\documentclass{book}
%\usepackage[letterpaper, portrait, margin=1cm]{geometry}
%\usepackage[letterpaper, bindingoffset=0.2in, left=1in,right=1in,top=.5in,bottom=.5in,footskip=.25in,marginparwidth=5em]{geometry}
\usepackage[letterpaper, left=1in,right=1in,top=.5in,bottom=.5in,footskip=.25in,marginparwidth=1cm]{geometry}
% ---------------------
% mini-table-of-content
% ---------------------
\usepackage{minitoc}
\setcounter{minitocdepth}{1}
\setlength{\mtcindent}{24pt}
\setcounter{secnumdepth}{-2}
%\renewcommand{\mtcfont}{\small\rm}
%\renewcommand{\mtcSfont}{\small\bf}
%\usepackage{setspace}
%\usepackage{tocloft}
%\setlength\cftparskip{-1.2pt}
%\setlength\cftbeforesecskip{1.3pt}
%\setlength\cftaftertoctitleskip{2pt}
%\renewcommand{\cftsecafterpnum}{\hspace*{02.0em}}
%\renewcommand{\cftsubsecafterpnum}{\hspace*{02.0em}}

% ---------------------------
% Chinese Characters Packages
% ---------------------------
\usepackage{fontspec} 
\usepackage{xeCJK}
\setmainfont{Times}
\setCJKmainfont{BiauKai}
\newfontfamily\sblgoodhebrew{SBL BibLit}[Script=Hebrew,Contextuals=Alternate]
\newfontfamily\sblgoodgreek{SBL BibLit}[Script=Greek,Contextuals=Alternate]

\usepackage{ifpdf,cite,algorithmic,url,tikz}
\usepackage[cmex10]{amsmath}

% ---------------------------
% Hebrew Characters Packages
% ---------------------------
\usepackage{polyglossia}
\setmainfont{Times New Roman}

% -------
% General
% -------
\usepackage{multicol}
\usepackage{multirow}
\usepackage{color,colortbl}
\usepackage{xparse}
\usepackage{pbox}
\usepackage{stackengine}
\usepackage{titlesec}% http://ctan.org/pkg/titlesec
\usepackage{tabularx}
\usepackage{xltabular}
\usepackage{titlesec}
\usepackage{makecell}
\newcommand{\sectionbreak}{\clearpage}

\author{
  Editor, Michael Chan\\
  \texttt{michaelchan\_wahyan@yahoo.com.hk}
}
\usepackage{tocloft}

\usepackage{hyperref}
\hypersetup{
    colorlinks=true, % set true if you want colored links
    linktoc   =all , % set to all if you want both sections and subsections linked
    linkcolor =blue, % choose some color if you want links to stand out
}

% ----------
% Afterword
% ----------
\usepackage{marginnote}
\usepackage{sectsty}
\usepackage{ragged2e}
\usepackage{lineno}
\usepackage{xcolor}
\usepackage{paracol}

\begin{document}

\clearpage
%% temporary titles
% command to provide stretchy vertical space in proportion
\newcommand\nbvspace[1][3]{\vspace*{\stretch{#1}}}
% allow some slack to avoid under/overfull boxes
\newcommand\nbstretchyspace{\spaceskip0.5em plus 0.25em minus 0.25em}
% To improve spacing on titlepages
\newcommand{\nbtitlestretch}{\spaceskip0.6em}
\pagestyle{empty}
\begin{center}
\bfseries
\nbvspace[1]
\Huge
{%\nbtitlestretch
\Large
\textbf{港九培靈研經會講章 \\
       Hong Kong Bible Conference 1928-2007
       }}

\nbvspace[1]

{\large
Editor: Michael\\
\texttt{michaelchan\_wahyan@yahoo.com.hk}
}

\nbvspace[1]

{\large
Revision: \texttt{v1.3}\\
Last Update: \today
}


\vfill
\begin{tikzpicture}
    %remove comment for OT cover%\node (0,0) [opacity=0.03]{\includegraphics[width=15cm]{../bible_out/ot_frontcover.png}} ;
    %remove comment for NT cover%\node (0,0) [opacity=0.03]{\includegraphics[width=15cm]{../bible_out/christ_on_cross.png}} ;
    %remove comment for Bible cover%\node (0,0) [xshift=0.8cm, yshift=+2cm, opacity=0.03]{\includegraphics[width=10cm]{./christ_on_cross.png}} ;
    %remove comment for Bible cover%\node (0,0) [              yshift=-2cm, opacity=0.03]{\includegraphics[width=14cm]{./ot_frontcover.png}} ;
\end{tikzpicture}
\vfill

\end{center}

\newpage

\setcounter{tocdepth}{0}
\dominitoc
\begin{multicols}{3}
\addtocontents{toc}{\protect\hypertarget{toc}{}}
\tableofcontents
\end{multicols}

\large
%\twocolumn

% the color definition syntax is as follow:
% \definecolor{name}{system}{definition}
% example: a mono-channel color can be defined as
%          \definecolor{Gray}{gray}{0.9}
% example: an rgb-3-channel color can be defined as
%          \definecolor{LightCyan}{rgb}{0.88,1,1}
%          \definecolor{pink}{rgb}{0.68,0,0.68}

\definecolor{CUV1LightRed}{rgb}{1,0.75,0.75}     % for CUV1
\definecolor{LZZVLightGray}{rgb}{0.9,0.9,0.9}    % for LZZ
\definecolor{KJVVLightGreen}{rgb}{0.75,1,0.85}   % for KJV
\definecolor{CUV2LightYellow}{rgb}{1,1,0.75}     % for CUV2
\definecolor{CNVVLightBrown}{rgb}{1,0.85,0.7}    % for CNV
\definecolor{NRSVLightBlue}{rgb}{0.75,1,1}       % for NRSV
\definecolor{WENLLightPurple}{rgb}{0.95,0.85,0.9}% for WENL
\definecolor{TCV19PaleGreen}{rgb}{0.85,1,0.95}   % for TCV19
\definecolor{MSGVLightWhite}{rgb}{0.98,0.98,0.98}% for MSGV
\definecolor{NETSLightRed}{rgb}{1,0.75,0.75}     % for NETS
\definecolor{JPS1917LightYellow}{rgb}{1,1,0.75}  % for JPS1917
\definecolor{SBLGNTPaleRed}{rgb}{1,0.85,0.80}    % for SBLGNT

\section{目錄}\label{sec:toc}
{ \scriptsize


\begin{xltabular}{\textwidth}{|p{0.08\textwidth} p{0.07\textwidth} p{0.25\textwidth}|p{0.15\textwidth} p{0.35\textwidth}|}
\hline
屆別 & 講號 & 經卷參照 & 講員 & 講題 \\
\hline
\hline
第1屆 & 第1講 &  & 古約翰 & \hyperref[sec:1322]{第一講 復興} \\
第1屆 & 第2講 &  & 古約翰 & \hyperref[sec:1326]{第二講 私藏數分的價銀} \\
第1屆 & 第3講 & 耶利米書4:3-4:4 & 古約翰 & \hyperref[sec:1327]{第三講 要開墾你們的荒地} \\
第1屆 & 第4講 &  & 古約翰 & \hyperref[sec:1329]{第四講 從墮落的地方起來} \\
第1屆 & 第5講 & 路加福音11:10-11:13 & 古約翰 & \hyperref[sec:1330]{第五講 禱告有效的秘訣} \\
第1屆 & 第6講 & 雅各書5:13-5:18 & 古約翰 & \hyperref[sec:1331]{第六講 有效力的禱告} \\
第1屆 & 第7講 & 以弗所書5:18-5:21 & 古約翰 & \hyperref[sec:1328]{第七講 要被聖靈充滿} \\
第1屆 & 第1講 & 路加福音9:18-9:36 & 吳子坤 & \hyperref[sec:1319]{開會詞} \\
第1屆 & 第1講 & 彼得前書2:4-2:7 & 王載 & \hyperref[sec:1351]{第七講 寶貝的耶穌} \\
第1屆 & 第2講 & 希伯來書10:4-10:7 & 王載 & \hyperref[sec:1352]{第八講 「我來了」} \\
第1屆 & 第3講 & 約翰福音10:22-10:30 & 王載 & \hyperref[sec:1353]{第九講 猶疑} \\
第1屆 & 第4講 & 箴言9:6-9:6 & 王載 & \hyperref[sec:1354]{第十講 五種愚人} \\
第1屆 & 第1講 & 使徒行傳2:30-2:36 & 祈理平 & \hyperref[sec:1345]{第一講 主釘十字架與復活} \\
第1屆 & 第2講 &  & 祈理平 & \hyperref[sec:1346]{第二講 靈洗} \\
第1屆 & 第3講 & 使徒行傳2:14-2:22 & 祈理平 & \hyperref[sec:1347]{第三講 耶穌再臨} \\
第1屆 & 第1講 &  & 趙柳塘 & \hyperref[sec:1338]{第一講　概論分三} \\
第1屆 & 第2講 & 加拉太書1:1-2:21 & 趙柳塘 & \hyperref[sec:1339]{第二講} \\
第1屆 & 第3講 &  & 趙柳塘 & \hyperref[sec:1340]{第三講} \\
第1屆 & 第4講 & 加拉太書2:15-4:7 & 趙柳塘 & \hyperref[sec:1341]{第四講} \\
第1屆 & 第5講 &  & 趙柳塘 & \hyperref[sec:1342]{第五講} \\
第1屆 & 第6講 & 加拉太書5:5-5:6 & 趙柳塘 & \hyperref[sec:1343]{第六講} \\
第1屆 & 第7講 &  & 趙柳塘 & \hyperref[sec:1344]{第七講} \\
第1屆 & 第1講 & 列王記下4:1-4:7 & 高樂弼 & \hyperref[sec:1348]{第四講 空瓶} \\
第1屆 & 第2講 & 馬太福音15:21-15:28 & 高樂弼 & \hyperref[sec:1349]{第五講 迦南婦} \\
第1屆 & 第3講 & 使徒行傳20:17-20:35 & 高樂弼 & \hyperref[sec:1350]{第六講 禱告與讀經} \\
第1屆 & 第1講 & 路加福音9:27-9:36 & 黃原素 & \hyperref[sec:1320]{開會講道} \\
第1屆 & 第1講 & 利未記1:1-7:10 & 黃原素 & \hyperref[sec:1334]{第一講　五祭之概論} \\
第1屆 & 第2講 & 何西阿書6:1-6:6 & 黃原素 & \hyperref[sec:1321]{閉會講道} \\
第1屆 & 第2講 & 利未記1:1-1:17 & 黃原素 & \hyperref[sec:1335]{第二講　燔祭之研究} \\
第1屆 & 第3講 & 利未記2:1-2:16 & 黃原素 & \hyperref[sec:1336]{第三講　素祭之研究} \\
第1屆 & 第4講 & 利未記3:1-3:17 & 黃原素 & \hyperref[sec:1337]{第四講　平安祭之研究} \\
第1屆 & 第5講 & 利未記4:1-4:35 & 黃原素 & \hyperref[sec:1333]{第五講　贖罪祭之研究} \\
第1屆 & 第6講 & 利未記5:1-5:19 & 黃原素 & \hyperref[sec:1332]{第六講　贖愆祭之研究} \\
\hline
\hline
第40屆 & 第1講 & 耶利米書1:2-1:3 & 吳勇 & \hyperref[sec:1534]{第一講 今日的時代} \\
第40屆 & 第2講 & 耶利米書2:9-2:11 & 吳勇 & \hyperref[sec:1535]{第二講 哀哭} \\
第40屆 & 第3講 & 耶利米書7:1-7:10 & 吳勇 & \hyperref[sec:1540]{第三講 三件傷心的事} \\
第40屆 & 第4講 & 耶利米書1:5-1:9 & 吳勇 & \hyperref[sec:1541]{第四講 呼召耶利米} \\
第40屆 & 第5講 & 創世記12:9-12:20 & 吳勇 & \hyperref[sec:1543]{第五講 復興的亞伯拉罕} \\
第40屆 & 第6講 & 創世記4:17-4:24 & 吳勇 & \hyperref[sec:1542]{第六講 人子近了} \\
第40屆 & 第7講 & 使徒行傳1:9-1:11 & 吳勇 & \hyperref[sec:1544]{第七講 主怎樣來?} \\
第40屆 & 第8講 & 民數記6:1-6:8 & 吳勇 & \hyperref[sec:1545]{第八講 請主來} \\
第40屆 & 第9講 & 使徒行傳26:28-26:28 & 吳勇 & \hyperref[sec:1546]{第九講 基督徒的分辨} \\
第40屆 & 第10講 & 約翰福音15:1-15:16 & 吳勇 & \hyperref[sec:1547]{第十講 住在基督裡} \\
第40屆 & 第11講 & 啟示錄2:1-2:4 & 吳勇 & \hyperref[sec:1548]{第十一講 末世教會} \\
第40屆 & 第12講 & 馬太福音18:15-18:20 & 吳勇 & \hyperref[sec:1539]{第十二講 教會的路} \\
第40屆 & 第13講 & 約翰福音7:37-7:38 & 吳勇 & \hyperref[sec:1538]{第十三講 生命的江河} \\
第40屆 & 第14講 & 以賽亞書66:7-66:8 & 吳勇 & \hyperref[sec:1537]{第十四講 神的傷心} \\
第40屆 & 第15講 & 馬太福音7:13-7:14 & 吳勇 & \hyperref[sec:1536]{第十五講 主的道路} \\
第40屆 & 第16講 & 哥林多前書2:14-2:15 & 吳勇 & \hyperref[sec:1533]{第十六講 跟隨耶穌} \\
第40屆 & 第1講 & 出埃及記17:8-17:16 & 吳明節 & \hyperref[sec:1508]{第一講 旗開得勝的摩西} \\
第40屆 & 第1講 & 馬太福音2:1-2:18 & 吳明節 & \hyperref[sec:1532]{第一講 逼迫逃難} \\
第40屆 & 第2講 & 馬太福音4:1-4:11 & 吳明節 & \hyperref[sec:1531]{第二講 被魔試探} \\
第40屆 & 第2講 & 撒母耳記上28:15-28:19 & 吳明節 & \hyperref[sec:1509]{第二講 失敗的掃羅} \\
第40屆 & 第3講 & 馬太福音13:53-13:58 & 吳明節 & \hyperref[sec:1530]{第三講 被人厭棄} \\
第40屆 & 第3講 &  & 吳明節 & \hyperref[sec:1510]{第三講 撒勒法的寡婦} \\
第40屆 & 第4講 & 馬太福音10:2-10:2 & 吳明節 & \hyperref[sec:1529]{第四講 感情的傷害} \\
第40屆 & 第4講 & 列王記上19:1-19:18 & 吳明節 & \hyperref[sec:1511]{第四講 灰心絕望的先知以利亞} \\
第40屆 & 第5講 & 列王記下2:1-2:15 & 吳明節 & \hyperref[sec:1512]{第五講 以利沙靈程的四個步驟} \\
第40屆 & 第5講 & 馬太福音27:33-27:34 & 吳明節 & \hyperref[sec:1528]{第五講 被釘十架} \\
第40屆 & 第6講 & 尼希米記1:1-1:4 & 吳明節 & \hyperref[sec:1513]{第六講 修建聖城的尼希米} \\
第40屆 & 第6講 & 馬太福音25:41-25:46 & 吳明節 & \hyperref[sec:1527]{第六講 常為弟兄而受苦} \\
第40屆 & 第7講 &  & 吳明節 & \hyperref[sec:1507]{第七講 背道的約拿} \\
第40屆 & 第1講 &  & 楊濬哲 & \hyperref[sec:1461]{序} \\
第40屆 & 第1講 &  & 鄭德音 & \hyperref[sec:1463]{港九培靈研經大會四十週年閉會詞} \\
第40屆 & 第1講 & 何西阿書1:1-1:1 & 鮑會園 & \hyperref[sec:1515]{第一講 神的話降臨} \\
第40屆 & 第2講 & 何西阿書1:4-1:5 & 鮑會園 & \hyperref[sec:1516]{第二講 耶斯列的罪} \\
第40屆 & 第3講 & 何西阿書2:1-2:23 & 鮑會園 & \hyperref[sec:1517]{第三講 神的管教} \\
第40屆 & 第4講 & 何西阿書3:1-3:5 & 鮑會園 & \hyperref[sec:1518]{第四講 完全的愛} \\
第40屆 & 第5講 & 何西阿書4:1-4:6 & 鮑會園 & \hyperref[sec:1519]{第五講 生命的知識} \\
第40屆 & 第6講 & 何西阿書4:1-4:19 & 鮑會園 & \hyperref[sec:1520]{第六講 神的爭辯} \\
第40屆 & 第7講 & 何西阿書5:1-6:11 & 鮑會園 & \hyperref[sec:1521]{第七講 悔改的呼召} \\
第40屆 & 第8講 & 何西阿書7:1-7:16 & 鮑會園 & \hyperref[sec:1526]{第八講 要心裡思想} \\
第40屆 & 第9講 & 何西阿書8:1-9:17 & 鮑會園 & \hyperref[sec:1524]{第九講 以色列人的偶像} \\
第40屆 & 第10講 & 何西阿書10:1-10:15 & 鮑會園 & \hyperref[sec:1525]{第十講 心懷二意} \\
第40屆 & 第11講 & 何西阿書11:1-12:1 & 鮑會園 & \hyperref[sec:1523]{第十一講 慈繩愛索} \\
第40屆 & 第12講 & 何西阿書12:1-13:16 & 鮑會園 & \hyperref[sec:1522]{第十二講 當歸向你的神} \\
第40屆 & 第13講 & 何西阿書14:1-14:9 & 鮑會園 & \hyperref[sec:1514]{第十三講 豐盛的生命} \\
第40屆 & 第1講 &  & 黃原素 & \hyperref[sec:1462]{港九培靈研經大會四十週年致詞} \\
\hline
\hline
第41屆 & 第1講 & 以弗所書1:3-1:14 & 孫維廉 & \hyperref[sec:1549]{第一講 天上屬靈的福氣(七福氣)} \\
第41屆 & 第1講 &  & 孫維廉 & \hyperref[sec:1471]{第一講 七個揭開(啟示錄)} \\
第41屆 & 第2講 & 啟示錄1:9-1:18 & 孫維廉 & \hyperref[sec:1472]{第二講 第二個揭幕 ── 在異象中的耶穌} \\
第41屆 & 第2講 & 以弗所書1:4-1:4 & 孫維廉 & \hyperref[sec:1465]{第二講 第一福氣 ── 神揀選我們} \\
第41屆 & 第3講 & 以弗所書1:5-1:5 & 孫維廉 & \hyperref[sec:1466]{第三講 第二福氣 ── 得兒子的名分} \\
第41屆 & 第3講 & 啟示錄1:19-1:20 & 孫維廉 & \hyperref[sec:1473]{第三講 第三個揭幕 ── 教會的主} \\
第41屆 & 第4講 & 以弗所書1:6-1:6 & 孫維廉 & \hyperref[sec:1467]{第四講 第三福氣 ── 被接納的恩典} \\
第41屆 & 第4講 & 啟示錄5:1-5:14 & 孫維廉 & \hyperref[sec:1474]{第四講 第四個揭幕 ── 揭開七印的獅子} \\
第41屆 & 第5講 & 以弗所書1:7-1:8 & 孫維廉 & \hyperref[sec:1468]{第五講 第四福氣 ── 赦免我們的罪} \\
第41屆 & 第5講 & 啟示錄8:1-8:6 & 孫維廉 & \hyperref[sec:1475]{第五講 第五個揭幕 ── 信徒禱告的效力} \\
第41屆 & 第6講 & 以弗所書1:9-1:10 & 孫維廉 & \hyperref[sec:1469]{第六講 第五福氣 ── 知道神的美意} \\
第41屆 & 第6講 & 啟示錄19:11-19:16 & 孫維廉 & \hyperref[sec:1476]{第六講 第六個揭幕 ── 再來的主} \\
第41屆 & 第7講 & 以弗所書1:11-1:12 & 孫維廉 & \hyperref[sec:1464]{第七講 第六福氣 ── 得神的基業} \\
第41屆 & 第7講 &  & 孫維廉 & \hyperref[sec:1470]{第七講 第七個揭幕 ── 祂是阿拉法,俄梅戛} \\
第41屆 & 第8講 & 以弗所書1:13-1:14 & 孫維廉 & \hyperref[sec:1325]{第八講 第七福氣 ── 聖靈為印記} \\
第41屆 & 第1講 &  & 林道亮 & \hyperref[sec:1506]{第一講 得永遠的生命} \\
第41屆 & 第2講 &  & 林道亮 & \hyperref[sec:1505]{第二講 人像神} \\
第41屆 & 第3講 &  & 林道亮 & \hyperref[sec:1504]{第三講 信耶穌與彼此相愛} \\
第41屆 & 第4講 &  & 林道亮 & \hyperref[sec:1503]{第四講 永火的審判} \\
第41屆 & 第5講 &  & 林道亮 & \hyperref[sec:1502]{第五講 進窄門} \\
第41屆 & 第6講 &  & 林道亮 & \hyperref[sec:1501]{第六講 進窄門得永生} \\
第41屆 & 第7講 &  & 林道亮 & \hyperref[sec:1500]{第七講 愛的兩面} \\
第41屆 & 第8講 &  & 林道亮 & \hyperref[sec:1499]{第八講 愛心的實行 ── 領人歸主} \\
第41屆 & 第9講 &  & 林道亮 & \hyperref[sec:1498]{第九講 認罪與平安} \\
第41屆 & 第10講 &  & 林道亮 & \hyperref[sec:1497]{第十講 更清楚認識罪} \\
第41屆 & 第11講 &  & 林道亮 & \hyperref[sec:1496]{第十一講 認識神的大名} \\
第41屆 & 第12講 &  & 林道亮 & \hyperref[sec:1495]{第十二講 以馬內利} \\
第41屆 & 第13講 &  & 林道亮 & \hyperref[sec:1494]{第十三講 持守末後的救恩} \\
第41屆 & 第14講 &  & 林道亮 & \hyperref[sec:1493]{第十四講 天國子民的資格} \\
第41屆 & 第15講 &  & 林道亮 & \hyperref[sec:1492]{第十五講 天國子民的福份} \\
第41屆 & 第16講 &  & 林道亮 & \hyperref[sec:1491]{第十六講 主的新命令} \\
第41屆 & 第1講 &  & 楊濬哲 & \hyperref[sec:1324]{序} \\
第41屆 & 第1講 & 列王記上1:1-22:53 & 胡恩德 & \hyperref[sec:1478]{第一講 失敗與希望} \\
第41屆 & 第2講 & 列王記上1:5-1:6 & 胡恩德 & \hyperref[sec:1479]{第二講 從神而來的權柄} \\
第41屆 & 第3講 & 列王記上2:1-2:9 & 胡恩德 & \hyperref[sec:1480]{第三講 公義的王權} \\
第41屆 & 第4講 &  & 胡恩德 & \hyperref[sec:1481]{第四講 智慧與豐富} \\
第41屆 & 第5講 &  & 胡恩德 & \hyperref[sec:1482]{第五講 耶和華神的殿} \\
第41屆 & 第6講 & 列王記上6:23-6:25 & 胡恩德 & \hyperref[sec:1483]{第六講 神的同在} \\
第41屆 & 第7講 & 列王記上10:1-10:10 & 胡恩德 & \hyperref[sec:1484]{第七講 名聲與失敗} \\
第41屆 & 第8講 & 列王記上12:1-12:33 & 胡恩德 & \hyperref[sec:1485]{第八講 分裂的國度} \\
第41屆 & 第9講 & 列王記上13:1-13:34 & 胡恩德 & \hyperref[sec:1486]{第九講 神的見證人} \\
第41屆 & 第10講 & 列王記上15:29-15:30 & 胡恩德 & \hyperref[sec:1487]{第十講 必按公義審判我們} \\
第41屆 & 第11講 & 列王記上17:1-17:24 & 胡恩德 & \hyperref[sec:1488]{第十一講 黑暗背道中神仍作工} \\
第41屆 & 第12講 & 列王記上17:17-17:24 & 胡恩德 & \hyperref[sec:1489]{第十二講 禱告的事奉} \\
第41屆 & 第13講 & 列王記上19:9-19:18 & 胡恩德 & \hyperref[sec:1490]{第十三講 為萬軍耶和華神大發熱心} \\
第41屆 & 第14講 & 列王記上20:1-20:6 & 胡恩德 & \hyperref[sec:1477]{第十四講 向人施恩的神} \\
\hline
\hline
第42屆 & 第1講 &  & 楊濬哲 & \hyperref[sec:1419]{序} \\
第42屆 & 第1講 &  & 滕近輝 & \hyperref[sec:1420]{第一講 七個呼聲} \\
第42屆 & 第2講 &  & 滕近輝 & \hyperref[sec:1422]{第二講 七福} \\
第42屆 & 第3講 &  & 滕近輝 & \hyperref[sec:1423]{第三講 基督的七對名字} \\
第42屆 & 第4講 &  & 滕近輝 & \hyperref[sec:1424]{第四講 新耶路撒冷的七種完全} \\
第42屆 & 第5講 &  & 滕近輝 & \hyperref[sec:1425]{第五講 七個對比} \\
第42屆 & 第6講 &  & 滕近輝 & \hyperref[sec:1426]{第六講 主耶穌的七對異象} \\
第42屆 & 第7講 &  & 滕近輝 & \hyperref[sec:1427]{第七講 七新} \\
第42屆 & 第8講 &  & 滕近輝 & \hyperref[sec:1428]{第八講 撒旦作為的七象} \\
第42屆 & 第9講 &  & 滕近輝 & \hyperref[sec:1429]{第九講 撒旦的七個名字} \\
第42屆 & 第10講 &  & 滕近輝 & \hyperref[sec:1430]{第十講 七個數字} \\
第42屆 & 第11講 &  & 滕近輝 & \hyperref[sec:1431]{第十一講 七讚} \\
第42屆 & 第12講 &  & 滕近輝 & \hyperref[sec:1421]{第十二講 天上七次發聲讚} \\
第42屆 & 第1講 & 創世記3:1-3:13 & 王永信 & \hyperref[sec:1446]{第一講 你在那裡?} \\
第42屆 & 第2講 & 帖撒羅尼迦前書2:3-2:10 & 王永信 & \hyperref[sec:1447]{第二講 信徒應有的三種表現} \\
第42屆 & 第3講 & 創世記4:1-4:11 & 王永信 & \hyperref[sec:1448]{第三講 你的兄弟在那裡?} \\
第42屆 & 第4講 & 創世記28:10-28:22 & 王永信 & \hyperref[sec:1449]{第四講 雅各屬靈上的轉機} \\
第42屆 & 第5講 & 出埃及記17:8-17:16 & 王永信 & \hyperref[sec:1450]{第五講 爭戰的兩支軍隊} \\
第42屆 & 第6講 & 以斯拉記1:1-1:3 & 王永信 & \hyperref[sec:1451]{第六講 為萬軍之耶和華重建聖殿} \\
第42屆 & 第7講 & 尼希米記1:1-1:4 & 王永信 & \hyperref[sec:1452]{第七講 重建聖城} \\
第42屆 & 第8講 & 羅馬書1:14-1:16 & 王永信 & \hyperref[sec:1453]{第八講 欠債} \\
第42屆 & 第9講 & 馬太福音6:24-6:23 & 王永信 & \hyperref[sec:1460]{第九講 教會面前的兩條路} \\
第42屆 & 第10講 & 以弗所書2:4-2:7 & 王永信 & \hyperref[sec:1454]{第十講 不徹底的教會} \\
第42屆 & 第11講 & 馬太福音12:30-12:4 & 王永信 & \hyperref[sec:1455]{第十一講 中立式的教會} \\
第42屆 & 第12講 & 士師記13:1-13:5 & 王永信 & \hyperref[sec:1456]{第十二講 失去能力的教會} \\
第42屆 & 第13講 & 士師記17:1-17:13 & 王永信 & \hyperref[sec:1457]{第十三講 假的教會} \\
第42屆 & 第14講 & 撒母耳記下12:1-12:10 & 王永信 & \hyperref[sec:1458]{第十四講 悔改的教會} \\
第42屆 & 第15講 & 撒母耳記上4:1-4:11 & 王永信 & \hyperref[sec:1459]{第十五講 以便以謝感恩的教會} \\
第42屆 & 第16講 & 歷代志上28:1-28:5 & 王永信 & \hyperref[sec:1445]{第十六講 建造教會} \\
第42屆 & 第1講 & 以西結書3:17-3:21 & 黃聿侯 & \hyperref[sec:1439]{第一講 領人歸主的責任} \\
第42屆 & 第1講 & 馬太福音16:13-16:20 & 黃聿侯 & \hyperref[sec:1550]{第一講 基督是永生神的兒子} \\
第42屆 & 第2講 & 馬太福音16:13-16:20 & 黃聿侯 & \hyperref[sec:1433]{第二講 人說人子是誰} \\
第42屆 & 第2講 & 詩篇122:1-122:9 & 黃聿侯 & \hyperref[sec:1440]{第二講 怎樣領人赴會} \\
第42屆 & 第3講 & 馬太福音16:13-16:20 & 黃聿侯 & \hyperref[sec:1434]{第三講 教會建造在磐石上} \\
第42屆 & 第3講 & 詩篇119:89-119:96 & 黃聿侯 & \hyperref[sec:1441]{第三講 略談讀經} \\
第42屆 & 第4講 & 馬太福音16:13-16:27 & 黃聿侯 & \hyperref[sec:1435]{第四講 生活在教會中的信仰} \\
第42屆 & 第4講 & 路加福音11:1-11:13 & 黃聿侯 & \hyperref[sec:1442]{第四講 學習禱告} \\
第42屆 & 第5講 & 羅馬書12:2-12:31 & 黃聿侯 & \hyperref[sec:1443]{第五講 肢體的生活} \\
第42屆 & 第5講 & 馬太福音16:13-16:20 & 黃聿侯 & \hyperref[sec:1436]{第五講 得勝陰間的權柄} \\
第42屆 & 第6講 & 馬太福音25:14-25:30 & 黃聿侯 & \hyperref[sec:1444]{第六講 做忠心有見識的管家} \\
第42屆 & 第6講 & 馬太福音16:13-16:27 & 黃聿侯 & \hyperref[sec:1437]{第六講 面向十字架的信仰} \\
第42屆 & 第7講 & 馬太福音16:13-16:27 & 黃聿侯 & \hyperref[sec:1432]{第七講 背負十字架的能力} \\
第42屆 & 第7講 & 約翰福音12:20-12:36 & 黃聿侯 & \hyperref[sec:1438]{第七講 至死盡忠} \\
\hline
\hline
第43屆 & 第1講 & 約翰福音20:19-20:23 & 周主培 & \hyperref[sec:1409]{第一講 將往何處?} \\
第43屆 & 第2講 & 羅馬書3:22-3:13 & 周主培 & \hyperref[sec:1410]{第二講 並沒有分別} \\
第43屆 & 第3講 & 馬太福音12:18-12:21 & 周主培 & \hyperref[sec:1411]{第三講 壓傷的蘆葦將殘的燈火} \\
第43屆 & 第4講 & 路加福音19:10-19:10 & 周主培 & \hyperref[sec:1412]{第四講 失喪的人} \\
第43屆 & 第5講 & 彼得前書4:17-4:6 & 周主培 & \hyperref[sec:1413]{第五講 求神鑑察} \\
第43屆 & 第6講 & 羅馬書1:16-1:16 & 周主培 & \hyperref[sec:1414]{第六講 福音本是神的大能} \\
第43屆 & 第7講 & 箴言23:26-23:26 & 周主培 & \hyperref[sec:1415]{第七講 將心歸神} \\
第43屆 & 第8講 & 啟示錄1:8-1:20 & 周主培 & \hyperref[sec:1416]{第八講 「沒有異象民就滅亡」} \\
第43屆 & 第9講 & 約翰福音21:15-21:19 & 周主培 & \hyperref[sec:1417]{第九講 奇妙的愛} \\
第43屆 & 第10講 & 詩篇16:1-16:11 & 周主培 & \hyperref[sec:1418]{第十講 靠主喜樂} \\
第43屆 & 第1講 &  & 楊濬哲 & \hyperref[sec:1382]{序} \\
第43屆 & 第1講 &  & 祈柏 & \hyperref[sec:1395]{第一講 我們活潑的盼望} \\
第43屆 & 第2講 & 彼得前書1:13-2:10 & 祈柏 & \hyperref[sec:1396]{第二講 聖潔生命之供養} \\
第43屆 & 第3講 & 彼得前書2:11-2:25 & 祈柏 & \hyperref[sec:1397]{第三講 蒙召受苦} \\
第43屆 & 第4講 & 彼得前書3:1-3:22 & 祈柏 & \hyperref[sec:1398]{第四講 家庭生活} \\
第43屆 & 第5講 & 彼得前書3:8-3:20 & 祈柏 & \hyperref[sec:1399]{第五講 積極的聖潔生活} \\
第43屆 & 第6講 & 彼得前書4:12-5:11 & 祈柏 & \hyperref[sec:1400]{第六講 為信仰而活} \\
第43屆 & 第1講 &  & 鮑會園 & \hyperref[sec:1401]{第一講 復興的預備──以斯拉記全貌} \\
第43屆 & 第2講 & 以斯拉記1:1-1:4 & 鮑會園 & \hyperref[sec:1402]{第二講 復興的啟示} \\
第43屆 & 第3講 & 以斯拉記1:1-1:11 & 鮑會園 & \hyperref[sec:1403]{第三講 復興的呼召} \\
第43屆 & 第4講 & 以斯拉記2:68-3:13 & 鮑會園 & \hyperref[sec:1404]{第四講 順服的喜樂} \\
第43屆 & 第5講 & 以斯拉記4:1-5:2 & 鮑會園 & \hyperref[sec:1405]{第五講 竭力爭戰} \\
第43屆 & 第6講 & 以斯拉記1:1-6:22 & 鮑會園 & \hyperref[sec:1406]{第六講 得勝有餘} \\
第43屆 & 第7講 & 以斯拉記8:1-8:36 & 鮑會園 & \hyperref[sec:1407]{第七講 持守見證} \\
第43屆 & 第8講 & 以斯拉記9:1-10:17 & 鮑會園 & \hyperref[sec:1408]{第八講 重新得力} \\
\hline
\hline
第44屆 & 第1講 & 箴言29:18-29:18 & 吉兆頎 & \hyperref[sec:1364]{第一講　關於異象外的意義} \\
第44屆 & 第3講 & 箴言29:18-29:18 & 吉兆頎 & \hyperref[sec:1365]{第三講　異象的功用} \\
第44屆 & 第4講 & 耶利米書23:16-23:16 & 吉兆頎 & \hyperref[sec:1366]{第四講　異象的鑑別} \\
第44屆 & 第5講 & 以西結書1:4-1:14 & 吉兆頎 & \hyperref[sec:1367]{第五講　異象的教訓} \\
第44屆 & 第6講 & 以西結書1:15-1:21 & 吉兆頎 & \hyperref[sec:1368]{第六講　活物的行動} \\
第44屆 & 第7講 & 以西結書1:22-1:28 & 吉兆頎 & \hyperref[sec:1370]{第七講　活物的環境} \\
第44屆 & 第8講 & 以西結書40:1-40:9 & 吉兆頎 & \hyperref[sec:1369]{第八講　聖殿的異象} \\
第44屆 & 第9講 &  & 吉兆頎 & \hyperref[sec:1371]{第九講　續「聖殿的異象」} \\
第44屆 & 第1講 &  & 楊濬哲 & \hyperref[sec:1355]{序} \\
第44屆 & 第1講 & 路得記1:1-1:22 & 費述凱 & \hyperref[sec:1372]{第一講　路得信心的選擇} \\
第44屆 & 第1講 & 希伯來書1:1-1:14 & 費述凱 & \hyperref[sec:1377]{第一講　你我心目中的耶穌,是怎樣的一位耶穌呢?} \\
第44屆 & 第2講 & 路得記2:1-2:23 & 費述凱 & \hyperref[sec:1374]{第二講　路得信心的試驗} \\
第44屆 & 第2講 & 希伯來書3:12-3:19 & 費述凱 & \hyperref[sec:1378]{第二講　當時教會的信徒,是怎樣的一些信徒呢?} \\
第44屆 & 第3講 & 路得記3:1-3:18 & 費述凱 & \hyperref[sec:1375]{第三講　路得信心的順服} \\
第44屆 & 第3講 & 希伯來書5:8-5:14 & 費述凱 & \hyperref[sec:1379]{第三講　永遠得救的問題} \\
第44屆 & 第4講 & 路得記4:1-4:22 & 費述凱 & \hyperref[sec:1373]{第四講　路得信心的賞賜} \\
第44屆 & 第4講 & 希伯來書7:1-7:28 & 費述凱 & \hyperref[sec:1380]{第四講　彌賽亞也是祭司嗎?} \\
第44屆 & 第5講 & 希伯來書10:1-10:18 & 費述凱 & \hyperref[sec:1381]{第五講　耶穌只一次獻上身體夠了嗎?} \\
第44屆 & 第6講 & 希伯來書11:1-11:40 & 費述凱 & \hyperref[sec:1376]{第六講　信主的定義和三種鼓勵} \\
第44屆 & 第1講 & 馬太福音11:28-11:30 & 陳終道 & \hyperref[sec:1323]{第一講　基督的性格} \\
第44屆 & 第2講 & 馬可福音6:34-6:34 & 陳終道 & \hyperref[sec:1356]{第二講　基督的工作} \\
第44屆 & 第3講 &  & 陳終道 & \hyperref[sec:1357]{第三講　基督的禱告} \\
第44屆 & 第4講 &  & 陳終道 & \hyperref[sec:1358]{第四講　基督的愛} \\
第44屆 & 第5講 &  & 陳終道 & \hyperref[sec:1359]{第五講　基督所責備及稱許的人} \\
第44屆 & 第6講 &  & 陳終道 & \hyperref[sec:1360]{第六講　基督所發的問題} \\
第44屆 & 第7講 &  & 陳終道 & \hyperref[sec:1361]{第七講　基督與瑪門} \\
第44屆 & 第8講 & 馬太福音15:17-15:20 & 陳終道 & \hyperref[sec:1362]{第八講　從四福音裡看舊約的基督} \\
第44屆 & 第9講 & 約翰福音14:1-14:31 & 陳終道 & \hyperref[sec:1363]{第九講　基督再來} \\
\hline
\hline
第45屆 & 第1講 &  & 吳勇 & \hyperref[sec:1260]{第一講 耶穌的寶血} \\
第45屆 & 第2講 & 啟示錄4:6-4:8 & 吳勇 & \hyperref[sec:1262]{第二講 生命的異象} \\
第45屆 & 第3講 & 撒母耳記上21:1-21:15 & 吳勇 & \hyperref[sec:1263]{第三講 軟弱和恢復} \\
第45屆 & 第4講 & 傳道書11:9-11:10 & 吳勇 & \hyperref[sec:1266]{第四講 青年與問題} \\
第45屆 & 第5講 & 但以理書2:31-2:35 & 吳勇 & \hyperref[sec:1268]{第五講 基督再來} \\
第45屆 & 第6講 & 民數記14:20-14:24 & 吳勇 & \hyperref[sec:1270]{第六講 被召者多選上者少} \\
第45屆 & 第7講 &  & 吳勇 & \hyperref[sec:1272]{第七講 建造聖城耶路撒冷} \\
第45屆 & 第8講 &  & 吳勇 & \hyperref[sec:1273]{第八講 會幕完工了} \\
第45屆 & 第9講 &  & 吳勇 & \hyperref[sec:1275]{第九講 同心合意興旺福音} \\
第45屆 & 第10講 &  & 吳勇 & \hyperref[sec:1277]{第十講 住在耶路撒冷的人} \\
第45屆 & 第1講 & 但以理書2:1-2:49 & 唐佑之 & \hyperref[sec:1246]{第一講 環境與經歷} \\
第45屆 & 第2講 & 但以理書2:1-2:49 & 唐佑之 & \hyperref[sec:1247]{第二講 夢幻或啟示} \\
第45屆 & 第3講 & 但以理書2:1-3:30 & 唐佑之 & \hyperref[sec:1248]{第三講 承受與擺脫} \\
第45屆 & 第4講 & 但以理書7:1-7:28 & 唐佑之 & \hyperref[sec:1250]{第四講 時間及永恆} \\
第45屆 & 第5講 & 但以理書8:1-8:27 & 唐佑之 & \hyperref[sec:1252]{第五講 異象及史實} \\
第45屆 & 第6講 & 但以理書9:1-9:27 & 唐佑之 & \hyperref[sec:1253]{第六講 蒙愛與期望} \\
第45屆 & 第7講 & 但以理書6:1-6:28 & 唐佑之 & \hyperref[sec:1255]{第七講 猛戰與天使} \\
第45屆 & 第8講 & 但以理書10:1-10:21 & 唐佑之 & \hyperref[sec:1258]{第八講 爭戰與得力} \\
第45屆 & 第9講 & 但以理書11:1-12:13 & 唐佑之 & \hyperref[sec:1259]{第九講 星光與永福} \\
第45屆 & 第1講 & 使徒行傳1:9-1:11 & 胡恩德 & \hyperref[sec:1225]{第一講 主再來的重要性} \\
第45屆 & 第2講 &  & 胡恩德 & \hyperref[sec:1227]{第二講 盼望主再來的態度} \\
第45屆 & 第3講 & 以弗所書1:15-1:18 & 胡恩德 & \hyperref[sec:1230]{第三講 盼望主再來的福氣} \\
第45屆 & 第4講 & 以弗所書1:13-1:14 & 胡恩德 & \hyperref[sec:1232]{第四講 得贖的日子} \\
第45屆 & 第5講 &  & 胡恩德 & \hyperref[sec:1234]{第五講 基督徒的審判} \\
第45屆 & 第6講 & 馬太福音24:1-25:46 & 胡恩德 & \hyperref[sec:1235]{第六講 主再來的初步豫兆} \\
第45屆 & 第7講 & 馬太福音24:15-24:31 & 胡恩德 & \hyperref[sec:1241]{第七講 主再來與有關猶太人的預兆} \\
第45屆 & 第8講 & 馬太福音24:15-24:16 & 胡恩德 & \hyperref[sec:1242]{第八講 敵基督顯現與主再來} \\
第45屆 & 第9講 &  & 胡恩德 & \hyperref[sec:1244]{第九講 星光與永福} \\
\hline
\hline
第46屆 & 第1講 & 雅歌1:1-1:4 & 于力工 & \hyperref[sec:1197]{第一講 基督是我們的愛人} \\
第46屆 & 第2講 &  & 于力工 & \hyperref[sec:1198]{第二講 基督是我們的牧人} \\
第46屆 & 第3講 & 雅歌1:13-1:14 & 于力工 & \hyperref[sec:1199]{第三講 基督是我們的榮耀} \\
第46屆 & 第4講 & 雅歌5:10-5:16 & 于力工 & \hyperref[sec:1201]{第四講 基督是我們的主} \\
第46屆 & 第5講 & 雅歌4:7-4:15 & 于力工 & \hyperref[sec:1203]{第五講 基督是我們的欣賞} \\
第46屆 & 第6講 & 雅歌2:2-2:6 & 于力工 & \hyperref[sec:1204]{第六講 基督是我們的旌旗} \\
第46屆 & 第7講 & 以賽亞書5:1-5:3 & 于力工 & \hyperref[sec:1286]{第七講 基督是我們的園主} \\
第46屆 & 第8講 & 雅歌5:15-5:23 & 于力工 & \hyperref[sec:1206]{第八講 基督是我們的朋友} \\
第46屆 & 第9講 & 雅歌1:1-8:14 & 于力工 & \hyperref[sec:1207]{第九講 基督是我們的一切} \\
第46屆 & 第1講 & 約翰福音1:35-1:39 & 周主培 & \hyperref[sec:1208]{第一講 你們尋找甚麼} \\
第46屆 & 第2講 & 哥林多前書2:1-2:5 & 周主培 & \hyperref[sec:1210]{第二講 耶穌基督} \\
第46屆 & 第3講 &  & 周主培 & \hyperref[sec:1212]{第三講 祂釘十架} \\
第46屆 & 第4講 & 以賽亞書6:1-6:13 & 周主培 & \hyperref[sec:1213]{第四講 光中見主} \\
第46屆 & 第5講 &  & 周主培 & \hyperref[sec:1214]{第五講 黎明日出} \\
第46屆 & 第6講 & 以賽亞書32:15-32:15 & 周主培 & \hyperref[sec:1215]{第六講 曠野樹林} \\
第46屆 & 第7講 & 約翰福音9:4-9:4 & 周主培 & \hyperref[sec:1217]{第七講 當前急務} \\
第46屆 & 第8講 & 馬太福音12:20-12:20 & 周主培 & \hyperref[sec:1218]{第八講 殘燈不滅} \\
第46屆 & 第9講 & 約翰福音12:1-12:11 & 周主培 & \hyperref[sec:1219]{第九講 模範教會} \\
第46屆 & 第10講 & 約書亞記24:14-24:28 & 周主培 & \hyperref[sec:1221]{第十講 全家事奉} \\
第46屆 & 第1講 &  & 鮑會園 & \hyperref[sec:1184]{第一講 知識的根基} \\
第46屆 & 第2講 & 歌羅西書1:9-1:14 & 鮑會園 & \hyperref[sec:1185]{第二講 真正的知識} \\
第46屆 & 第3講 & 歌羅西書1:14-1:14 & 鮑會園 & \hyperref[sec:1186]{第三講 無比的基督} \\
第46屆 & 第4講 & 歌羅西書1:23-1:23 & 鮑會園 & \hyperref[sec:1187]{第四講 蒙福的職份} \\
第46屆 & 第5講 &  & 鮑會園 & \hyperref[sec:1188]{第五講 有意義的生活} \\
第46屆 & 第6講 & 歌羅西書2:8-2:23 & 鮑會園 & \hyperref[sec:1189]{第六講 以基督誇勝} \\
第46屆 & 第7講 & 歌羅西書3:1-3:17 & 鮑會園 & \hyperref[sec:1190]{第七講 新人的生活} \\
第46屆 & 第8講 & 歌羅西書3:18-3:18 & 鮑會園 & \hyperref[sec:1191]{第八講 新人的責任} \\
第46屆 & 第9講 & 歌羅西書4:2-4:18 & 鮑會園 & \hyperref[sec:1192]{第九講 保羅最後的勸告} \\
\hline
\hline
第47屆 & 第1講 & 以弗所書3:3-3:6 & 包忠傑 & \hyperref[sec:1131]{第一講 基督的奧秘} \\
第47屆 & 第2講 & 彌迦書5:2-5:27 & 包忠傑 & \hyperref[sec:1134]{第二講 舊約中的基督} \\
第47屆 & 第3講 & 腓立比書2:5-2:8 & 包忠傑 & \hyperref[sec:1136]{第三講 道成肉身的基督} \\
第47屆 & 第4講 & 提多書2:11-2:15 & 包忠傑 & \hyperref[sec:1137]{第四講 基督的恩典} \\
第47屆 & 第5講 & 希伯來書9:11-9:28 & 包忠傑 & \hyperref[sec:1140]{第五講 基督的寶血} \\
第47屆 & 第6講 & 路加福音9:22-9:26 & 包忠傑 & \hyperref[sec:1142]{第六講 基督的十字架} \\
第47屆 & 第7講 & 路加福音24:13-24:46 & 包忠傑 & \hyperref[sec:1143]{第七講 基督的復活與升天} \\
第47屆 & 第8講 & 馬可福音11:1-11:10 & 包忠傑 & \hyperref[sec:1144]{第八講 再來之王} \\
第47屆 & 第1講 &  & 滕近輝 & \hyperref[sec:1161]{第一講 聖靈好像火一樣} \\
第47屆 & 第2講 &  & 滕近輝 & \hyperref[sec:1163]{第二講 靈火的能力在以利亞身上的具體表現} \\
第47屆 & 第3講 &  & 滕近輝 & \hyperref[sec:1164]{第三講 聖靈好像耶和華的手} \\
第47屆 & 第4講 &  & 滕近輝 & \hyperref[sec:1165]{第四講 聖靈的手降在以西結身上} \\
第47屆 & 第5講 &  & 滕近輝 & \hyperref[sec:1166]{第五講 聖靈的作為} \\
第47屆 & 第6講 & 馬太福音13:31-13:33 & 滕近輝 & \hyperref[sec:1167]{第六講 新酵與舊酵} \\
第47屆 & 第7講 &  & 滕近輝 & \hyperref[sec:1168]{第七講 聖靈的油也是能力的油} \\
第47屆 & 第8講 &  & 滕近輝 & \hyperref[sec:1169]{第八講 聖靈完全的能力} \\
第47屆 & 第9講 &  & 滕近輝 & \hyperref[sec:1170]{第九講 三個關於聖靈的比喻} \\
第47屆 & 第10講 & 以西結書1:12-1:21 & 滕近輝 & \hyperref[sec:1171]{第十講 聖靈的充滿} \\
第47屆 & 第1講 & 哥林多後書1:1-1:11 & 鮑會園 & \hyperref[sec:1148]{第一講 苦難與安慰} \\
第47屆 & 第2講 & 哥林多後書1:12-1:12 & 鮑會園 & \hyperref[sec:1150]{第二講 生命的記號} \\
第47屆 & 第3講 & 哥林多後書3:1-4:18 & 鮑會園 & \hyperref[sec:1152]{第三講 新約的執事} \\
第47屆 & 第4講 & 哥林多後書5:1-5:21 & 鮑會園 & \hyperref[sec:1154]{第四講 基督的使者} \\
第47屆 & 第5講 & 哥林多後書6:1-7:16 & 鮑會園 & \hyperref[sec:1156]{第五講 天父的兒女} \\
第47屆 & 第6講 & 哥林多後書8:1-9:15 & 鮑會園 & \hyperref[sec:1157]{第六講 樂意的奉獻} \\
第47屆 & 第7講 & 哥林多後書10:1-10:18 & 鮑會園 & \hyperref[sec:1158]{第七講 屬靈的權柄} \\
第47屆 & 第8講 & 哥林多後書12:1-13:14 & 鮑會園 & \hyperref[sec:1160]{第八講 足夠的恩典} \\
\hline
\hline
第48屆 & 第1講 &  & 史祈生 & \hyperref[sec:1059]{第一講 生命的印證(一)} \\
第48屆 & 第2講 & 馬太福音5:4-5:4 & 史祈生 & \hyperref[sec:1060]{第二講 生命的印證(二)} \\
第48屆 & 第3講 & 馬太福音5:5-5:5 & 史祈生 & \hyperref[sec:1061]{第三講 生命的印證(三)} \\
第48屆 & 第4講 & 馬太福音2:1-6:34 & 史祈生 & \hyperref[sec:1062]{第四講 生命的成長(一)} \\
第48屆 & 第5講 & 馬太福音5:7-5:7 & 史祈生 & \hyperref[sec:1095]{第五講 生命的成長(二)} \\
第48屆 & 第6講 & 馬太福音5:8-5:8 & 史祈生 & \hyperref[sec:1097]{第六講 生命的成長(三)} \\
第48屆 & 第7講 & 馬太福音5:9-5:9 & 史祈生 & \hyperref[sec:1099]{第七講 生命的操練(一)} \\
第48屆 & 第8講 & 馬太福音5:10-5:12 & 史祈生 & \hyperref[sec:1100]{第八講 生命的操練(二,三)} \\
第48屆 & 第1講 & 列王記上18:12-18:40 & 孫德生 & \hyperref[sec:1101]{第一講 聖靈的火} \\
第48屆 & 第2講 & 約翰福音14:15-14:18 & 孫德生 & \hyperref[sec:1103]{第二講 聖靈是誰} \\
第48屆 & 第3講 & 使徒行傳2:1-2:8 & 孫德生 & \hyperref[sec:1104]{第三講 聖靈降臨} \\
第48屆 & 第4講 & 以弗所書4:25-4:32 & 孫德生 & \hyperref[sec:1106]{第四講 罪} \\
第48屆 & 第5講 & 以弗所書5:15-6:9 & 孫德生 & \hyperref[sec:1108]{第五講 聖靈充滿} \\
第48屆 & 第6講 & 加拉太書5:16-5:23 & 孫德生 & \hyperref[sec:1110]{第六講 聖靈所結的果子} \\
第48屆 & 第7講 & 路加福音11:1-11:13 & 孫德生 & \hyperref[sec:1111]{第七講 聖靈與禱告} \\
第48屆 & 第8講 & 使徒行傳11:19-11:30 & 孫德生 & \hyperref[sec:1126]{第八講 傳福音給萬民} \\
第48屆 & 第9講 & 約書亞記14:6-14:14 & 孫德生 & \hyperref[sec:1128]{第九講 迦勒專心跟從耶和華} \\
第48屆 & 第10講 & 創世記45:1-45:15 & 孫德生 & \hyperref[sec:1129]{第十講 高山與深淵} \\
第48屆 & 第1講 &  & 張子華 & \hyperref[sec:1051]{第一講 神的兩個方法} \\
第48屆 & 第2講 & 哈巴谷書1:12-2:10 & 張子華 & \hyperref[sec:1052]{第二講 義人必因信得生} \\
第48屆 & 第3講 & 哈巴谷書3:1-3:19 & 張子華 & \hyperref[sec:1053]{第三講 苦難中的喜樂} \\
第48屆 & 第4講 & 瑪拉基書1:1-1:5 & 張子華 & \hyperref[sec:1054]{第四講 懷疑神的愛} \\
第48屆 & 第5講 & 瑪拉基書1:6-1:6 & 張子華 & \hyperref[sec:1055]{第五講 藐視神的名} \\
第48屆 & 第6講 & 瑪拉基書2:17-2:17 & 張子華 & \hyperref[sec:1056]{第六講 懷疑神的公義} \\
第48屆 & 第7講 & 瑪拉基書3:8-3:12 & 張子華 & \hyperref[sec:1057]{第七講 奪取神的供物} \\
第48屆 & 第8講 & 瑪拉基書2:10-2:16 & 張子華 & \hyperref[sec:1058]{第八講 兩個問題} \\
\hline
\hline
第49屆 & 第1講 &  & 張子華 & \hyperref[sec:777]{第一講 順服神權柄的重要性} \\
第49屆 & 第2講 &  & 張子華 & \hyperref[sec:778]{第二講 活在神話語旨意中的和諧生活} \\
第49屆 & 第3講 & 啟示錄2:1-2:4 & 張子華 & \hyperref[sec:779]{第三講 與神和諧關係破壞的原因－ 失去起初的愛心} \\
第49屆 & 第4講 & 撒母耳記上15:1-15:23 & 張子華 & \hyperref[sec:780]{第四講 與神和諧的關係－ 順服神} \\
第49屆 & 第5講 & 馬太福音5:21-5:25 & 張子華 & \hyperref[sec:781]{第五講 與人和諧關係－ 如何解憎恨,不和} \\
第49屆 & 第6講 & 以弗所書2:10-2:10 & 張子華 & \hyperref[sec:782]{第六講 接受你自己} \\
第49屆 & 第7講 & 以弗所書6:1-6:3 & 張子華 & \hyperref[sec:783]{第七講 家庭的和諧－ 向父母服權} \\
第49屆 & 第8講 & 以弗所書6:4-6:4 & 張子華 & \hyperref[sec:784]{第八講 家庭的和諧－ 父母的責任} \\
第49屆 & 第9講 & 創世記1:26-1:27 & 張子華 & \hyperref[sec:785]{第九講 建立夫妻和諧生活} \\
第49屆 & 第1講 &  & 戴紹曾 & \hyperref[sec:786]{第一講　初期教會的耶穌} \\
第49屆 & 第2講 &  & 戴紹曾 & \hyperref[sec:787]{第二講　聖靈和祂的工作} \\
第49屆 & 第3講 &  & 戴紹曾 & \hyperref[sec:788]{第三講　被聖靈充滿的司提反} \\
第49屆 & 第4講 &  & 戴紹曾 & \hyperref[sec:789]{第四講　巴拿巴完全的奉獻} \\
第49屆 & 第5講 &  & 戴紹曾 & \hyperref[sec:790]{第五講　巴拿巴衷心饒恕他的弟兄} \\
第49屆 & 第6講 &  & 戴紹曾 & \hyperref[sec:791]{第六講　巴拿巴的事奉} \\
第49屆 & 第7講 &  & 戴紹曾 & \hyperref[sec:792]{第七講　巴拿巴的見證} \\
第49屆 & 第8講 &  & 戴紹曾 & \hyperref[sec:793]{第八講　尋求神的旨意} \\
第49屆 & 第9講 &  & 戴紹曾 & \hyperref[sec:794]{第九講　初期教會增長的情形} \\
第49屆 & 第10講 &  & 戴紹曾 & \hyperref[sec:795]{第十講　為主受苦} \\
第49屆 & 第1講 &  & 楊濬哲 & \hyperref[sec:1385]{序} \\
第49屆 & 第1講 &  & 艾理德 & \hyperref[sec:768]{第一講 當心你的生活} \\
第49屆 & 第2講 & 約翰一書1:1-1:4 & 艾理德 & \hyperref[sec:769]{第二講 與神與弟兄相交} \\
第49屆 & 第3講 & 約翰一書1:5-2:29 & 艾理德 & \hyperref[sec:770]{第三講 與基督教不能分開的倫理宣告} \\
第49屆 & 第4講 & 約翰一書1:5-2:6 & 艾理德 & \hyperref[sec:771]{第四講 與基督教不能分開的倫理宣告(續上)} \\
第49屆 & 第5講 & 約翰一書2:7-2:27 & 艾理德 & \hyperref[sec:772]{第五講 生活的證據} \\
第49屆 & 第6講 & 約翰一書2:18-2:27 & 艾理德 & \hyperref[sec:773]{第六講 生活的試驗} \\
第49屆 & 第7講 & 約翰一書2:28-2:29 & 艾理德 & \hyperref[sec:774]{第七講 殷勤去作神的工} \\
第49屆 & 第8講 & 約翰一書3:1-4:21 & 艾理德 & \hyperref[sec:775]{第八講 神兒女的權利和義務} \\
第49屆 & 第9講 & 約翰一書3:18-3:24 & 艾理德 & \hyperref[sec:776]{第九講 愛裡沒有懼怕} \\
\hline
\hline
第50屆 & 第1講 & 路加福音24:49-24:5 & 楊濬哲 & \hyperref[sec:736]{第一講 為甚麼我們要被聖靈充滿} \\
第50屆 & 第2講 &  & 楊濬哲 & \hyperref[sec:738]{第二講 甚麼是被聖靈充滿} \\
第50屆 & 第3講 &  & 楊濬哲 & \hyperref[sec:739]{第三講 被聖靈充滿的途徑} \\
第50屆 & 第4講 &  & 楊濬哲 & \hyperref[sec:740]{第四講 被聖靈充滿的果效} \\
第50屆 & 第1講 & 使徒行傳17:24-17:25 & 鄭德音 & \hyperref[sec:731]{第一講 神是誰?} \\
第50屆 & 第2講 & 創世記3:1-3:24 & 鄭德音 & \hyperref[sec:732]{第二講 「人的罪惡」與「神的救贖」} \\
第50屆 & 第3講 & 以賽亞書43:1-43:8 & 鄭德音 & \hyperref[sec:734]{第三講 神的救贖工作} \\
第50屆 & 第1講 &  & 陳終道 & \hyperref[sec:714]{第一講 聖經奇妙的影響力} \\
第50屆 & 第2講 & 詩篇119:89-119:17 & 陳終道 & \hyperref[sec:715]{第二講 聖經的可信} \\
第50屆 & 第3講 & 以賽亞書50:4-50:4 & 陳終道 & \hyperref[sec:716]{第三講讀聖經應有的態度} \\
第50屆 & 第4講 & 申命記6:4-6:9 & 陳終道 & \hyperref[sec:717]{第四講 靈修讀經和熟讀聖經} \\
第50屆 & 第5講 & 詩篇1:1-1:6 & 陳終道 & \hyperref[sec:718]{第五講 精細讀經} \\
第50屆 & 第6講 & 提摩太後書2:15-2:15 & 陳終道 & \hyperref[sec:720]{第六講 如何避免誤解聖經} \\
第50屆 & 第7講 & 彼得後書1:20-1:21 & 陳終道 & \hyperref[sec:721]{第七講 查經可能遇到的難題} \\
第50屆 & 第8講 & 歌羅西書2:16-2:17 & 陳終道 & \hyperref[sec:723]{第八講 聖經的預表和靈意} \\
第50屆 & 第9講 & 路加福音24:27-24:27 & 陳終道 & \hyperref[sec:724]{第九講 聖經的一貫性} \\
第50屆 & 第1講 & 歷代志下20:1-20:12 & 高文 & \hyperref[sec:750]{第一講　承認罪過} \\
第50屆 & 第2講 & 馬太福音18:15-18:20 & 高文 & \hyperref[sec:751]{第二講　活在和諧禱告生活中} \\
第50屆 & 第3講 & 希伯來書9:11-9:12 & 高文 & \hyperref[sec:752]{第三講　得主寶血潔淨} \\
第50屆 & 第4講 & 約翰福音14:12-14:19 & 高文 & \hyperref[sec:754]{第四講　聖靈充滿的生活} \\
第50屆 & 第5講 & 約翰福音17:17-17:26 & 高文 & \hyperref[sec:755]{第五講　事奉永活的神} \\
第50屆 & 第6講 & 路加福音19:1-19:10 & 高文 & \hyperref[sec:757]{第六講　得著新生命} \\
第50屆 & 第7講 & 羅馬書8:14-8:18 & 高文 & \hyperref[sec:758]{第七講　活在神全美的旨意裡} \\
第50屆 & 第8講 & 以弗所書5:25-5:33 & 高文 & \hyperref[sec:759]{第八講　活在聖潔生活中} \\
第50屆 & 第9講 & 約翰一書3:1-3:3 & 高文 & \hyperref[sec:761]{第九講　活在主再來的盼望中} \\
第50屆 & 第10講 & 啟示錄12:11-12:11 & 高文 & \hyperref[sec:762]{第十講　得勝的生活} \\
\hline
\hline
第51屆 & 第1講 & 撒迦利亞書1:1-1:17 & 唐佑之 & \hyperref[sec:663]{第一講　四匹馬的異象} \\
第51屆 & 第2講 & 撒迦利亞書1:18-2:13 & 唐佑之 & \hyperref[sec:674]{第二講 四角與準繩的異象} \\
第51屆 & 第3講 & 撒迦利亞書3:1-3:10 & 唐佑之 & \hyperref[sec:675]{第三講 約書亞的異象} \\
第51屆 & 第4講 & 撒迦利亞書4:1-4:2 & 唐佑之 & \hyperref[sec:676]{第四講 教會更新的果效} \\
第51屆 & 第5講 & 撒迦利亞書6:9-6:13 & 唐佑之 & \hyperref[sec:677]{第五講 教會應走的路線} \\
第51屆 & 第6講 & 撒迦利亞書7:1-7:7 & 唐佑之 & \hyperref[sec:678]{第六講 教會復興的真義} \\
第51屆 & 第7講 & 撒迦利亞書10:1-10:7 & 唐佑之 & \hyperref[sec:679]{第七講 好牧人} \\
第51屆 & 第8講 & 撒迦利亞書13:1-13:1 & 唐佑之 & \hyperref[sec:680]{第八講 泉源和活水} \\
第51屆 & 第1講 &  & 張慕皚 & \hyperref[sec:682]{第一講 認識我們的時代} \\
第51屆 & 第2講 &  & 張慕皚 & \hyperref[sec:683]{第二講 屬靈生命的更新} \\
第51屆 & 第3講 &  & 張慕皚 & \hyperref[sec:685]{第三講 禱告生活的更新} \\
第51屆 & 第4講 & 彼得前書4:7-4:8 & 張慕皚 & \hyperref[sec:686]{第四講 奉獻生活的更新} \\
第51屆 & 第5講 &  & 張慕皚 & \hyperref[sec:688]{第五講 彼此相愛的更新} \\
第51屆 & 第6講 &  & 張慕皚 & \hyperref[sec:690]{第六講 崇拜生活的更新} \\
第51屆 & 第7講 &  & 張慕皚 & \hyperref[sec:691]{第七講 人價值觀的更新} \\
第51屆 & 第8講 &  & 張慕皚 & \hyperref[sec:692]{第八講 時代使命感的更新} \\
第51屆 & 第1講 &  & 楊濬哲 & \hyperref[sec:1383]{序} \\
第51屆 & 第1講 & 出埃及記24:9-24:10 & 滕近輝 & \hyperref[sec:693]{第一講 屬靈的境界} \\
第51屆 & 第2講 & 申命記34:1-34:3 & 滕近輝 & \hyperref[sec:696]{第二講 愛慕屬靈的福份 ― 信心的探險} \\
第51屆 & 第3講 &  & 滕近輝 & \hyperref[sec:697]{第三講 舉手得勝 ― 何烈山上的復興} \\
第51屆 & 第4講 &  & 滕近輝 & \hyperref[sec:698]{第四講 完全的奉獻} \\
第51屆 & 第5講 &  & 滕近輝 & \hyperref[sec:700]{第五講 各各他山上的經歷} \\
第51屆 & 第6講 &  & 滕近輝 & \hyperref[sec:701]{第六講 全心全意歸向神} \\
第51屆 & 第7講 &  & 滕近輝 & \hyperref[sec:703]{第七講 看見基督的榮耀} \\
第51屆 & 第8講 &  & 滕近輝 & \hyperref[sec:704]{第八講 合神心意的教會接受大使命} \\
第51屆 & 第9講 &  & 滕近輝 & \hyperref[sec:706]{第九講 儆醒預備等候主再來} \\
\hline
\hline
第52屆 & 第1講 &  & 史保羅 & \hyperref[sec:647]{第一講　愛是恆久忍耐又有恩慈} \\
第52屆 & 第2講 & 哥林多前書13:4-13:5 & 史保羅 & \hyperref[sec:649]{第二講　愛是不嫉妒不自誇不張狂} \\
第52屆 & 第3講 & 哥林多前書13:5-13:5 & 史保羅 & \hyperref[sec:651]{第三講　愛是不作害羞的事不求自己的益處不輕易發怒} \\
第52屆 & 第4講 & 哥林多前書13:5-13:6 & 史保羅 & \hyperref[sec:652]{第四講　愛是不計算人的惡不喜歡不義只喜歡真理} \\
第52屆 & 第5講 & 哥林多前書13:1-13:13 & 史保羅 & \hyperref[sec:654]{第五講　凡事包容,相信,盼望,忍耐} \\
第52屆 & 第6講 & 哥林多前書13:1-13:3 & 史保羅 & \hyperref[sec:656]{第六講　無比的愛} \\
第52屆 & 第7講 & 哥林多前書13:8-13:13 & 史保羅 & \hyperref[sec:657]{第七講　愛的增長} \\
第52屆 & 第8講 &  & 史保羅 & \hyperref[sec:659]{第八講　當代神的聲音} \\
第52屆 & 第9講 &  & 史保羅 & \hyperref[sec:661]{第九講　神的呼召} \\
第52屆 & 第10講 &  & 史保羅 & \hyperref[sec:662]{第十講　神的旨意} \\
第52屆 & 第1講 &  & 楊濬哲 & \hyperref[sec:1384]{序} \\
第52屆 & 第1講 &  & 焦源濂 & \hyperref[sec:614]{第一講 約伯記簡介} \\
第52屆 & 第2講 & 約伯記1:6-1:22 & 焦源濂 & \hyperref[sec:618]{第二講　撒旦對約伯的第一次攻擊} \\
第52屆 & 第3講 & 約伯記2:1-3:26 & 焦源濂 & \hyperref[sec:620]{第三講　撒旦對約伯的第二次攻擊} \\
第52屆 & 第4講 & 約伯記4:1-4:21 & 焦源濂 & \hyperref[sec:621]{第四講　失敗的安慰者-以利法(一)} \\
第52屆 & 第5講 & 約伯記8:1-8:22 & 焦源濂 & \hyperref[sec:623]{第五講　失敗的安慰者-以利法(二)} \\
第52屆 & 第6講 & 約伯記11:1-31:20 & 焦源濂 & \hyperref[sec:624]{第六講　失敗的安慰者-比勒達} \\
第52屆 & 第7講 & 約伯記11:1-11:12 & 焦源濂 & \hyperref[sec:626]{第七講　失敗的安慰者-瑣法} \\
第52屆 & 第8講 & 約伯記38:1-42:17 & 焦源濂 & \hyperref[sec:628]{第八講　疑團盡消 皆大歡喜} \\
第52屆 & 第1講 &  & 薛玉光 & \hyperref[sec:632]{第一講　認識獨一的真神} \\
第52屆 & 第2講 & 約翰福音17:1-17:3 & 薛玉光 & \hyperref[sec:633]{第二講　認識唯一的救主耶穌} \\
第52屆 & 第3講 & 約翰福音16:7-16:14 & 薛玉光 & \hyperref[sec:634]{第三講　認識聖靈} \\
第52屆 & 第4講 & 耶利米書32:36-32:39 & 薛玉光 & \hyperref[sec:636]{第四講　信服永不改變的聖經} \\
第52屆 & 第5講 & 耶利米書2:36-2:36 & 薛玉光 & \hyperref[sec:638]{第五講　神學信仰的合一} \\
第52屆 & 第6講 & 約翰一書1:1-1:4 & 薛玉光 & \hyperref[sec:639]{第六講　信徒生命的合一} \\
第52屆 & 第7講 & 帖撒羅尼迦前書1:9-1:10 & 薛玉光 & \hyperref[sec:641]{第七講　教會事奉的合一} \\
第52屆 & 第8講 & 馬太福音28:16-28:20 & 薛玉光 & \hyperref[sec:642]{第八講　福音工場的合一} \\
\hline
\hline
第53屆 & 第1講 & 使徒行傳1:1-1:11 & 戴紹曾 & \hyperref[sec:591]{第一講 聖靈與差傳事工} \\
第53屆 & 第2講 & 使徒行傳4:23-4:31 & 戴紹曾 & \hyperref[sec:593]{第二講 禱告與差傳} \\
第53屆 & 第3講 &  & 戴紹曾 & \hyperref[sec:595]{第三講 神的話與差傳} \\
第53屆 & 第4講 & 使徒行傳8:1-8:8 & 戴紹曾 & \hyperref[sec:596]{第四講 基督徒與差傳} \\
第53屆 & 第5講 & 使徒行傳13:1-13:3 & 戴紹曾 & \hyperref[sec:603]{第五講 教會與差傳} \\
第53屆 & 第6講 & 使徒行傳14:19-14:28 & 戴紹曾 & \hyperref[sec:604]{第六講 差傳與事奉} \\
第53屆 & 第7講 & 使徒行傳18:1-18:4 & 戴紹曾 & \hyperref[sec:605]{第七講 差傳工作與家庭} \\
第53屆 & 第8講 & 腓立比書3:7-3:16 & 戴紹曾 & \hyperref[sec:608]{第八講 差傳工作的原動力} \\
第53屆 & 第9講 & 使徒行傳16:1-16:5 & 戴紹曾 & \hyperref[sec:610]{第九講 差傳工作的訓練和裝備} \\
第53屆 & 第10講 &  & 戴紹曾 & \hyperref[sec:611]{第十講 差傳與戴氏家族(略述)} \\
第53屆 & 第1講 & 出埃及記2:23-2:25 & 林克己 & \hyperref[sec:574]{第一講 蒙神呼召的神蹟} \\
第53屆 & 第2講 & 出埃及記15:22-15:26 & 林克己 & \hyperref[sec:576]{第二講 苦水變甜的神蹟} \\
第53屆 & 第3講 & 創世記11:31-11:32 & 林克己 & \hyperref[sec:578]{第三講 力上加力的神蹟} \\
第53屆 & 第4講 & 列王記上18:20-18:46 & 林克己 & \hyperref[sec:583]{第四講 降火下雨的神蹟} \\
第53屆 & 第5講 & 哈巴谷書3:16-3:19 & 林克己 & \hyperref[sec:585]{第五講 靠主喜樂的神蹟} \\
第53屆 & 第6講 & 馬可福音6:30-6:44 & 林克己 & \hyperref[sec:586]{第六講 五餅二魚的神蹟} \\
第53屆 & 第7講 & 馬太福音8:23-8:27 & 林克己 & \hyperref[sec:588]{第七講 耶穌履海的神蹟} \\
第53屆 & 第8講 & 羅馬書8:1-8:2 & 林克己 & \hyperref[sec:590]{第八講 得勝有餘的神蹟} \\
第53屆 & 第1講 &  & 楊濬哲 & \hyperref[sec:1386]{序} \\
第53屆 & 第1講 &  & 鮑會園 & \hyperref[sec:555]{第一講 本書簡介及教會的分爭結黨} \\
第53屆 & 第2講 & 哥林多前書2:14-3:7 & 鮑會園 & \hyperref[sec:558]{第二講 三等人} \\
第53屆 & 第3講 & 哥林多前書3:10-3:17 & 鮑會園 & \hyperref[sec:559]{第三講 屬靈生命的建造} \\
第53屆 & 第4講 & 哥林多前書4:1-4:13 & 鮑會園 & \hyperref[sec:565]{第四講 錯誤的思想} \\
第53屆 & 第5講 & 哥林多前書5:1-5:13 & 鮑會園 & \hyperref[sec:567]{第五講 淫亂的問題} \\
第53屆 & 第6講 & 哥林多前書6:12-6:20 & 鮑會園 & \hyperref[sec:569]{第六講 基督徒生活的原則} \\
第53屆 & 第7講 & 哥林多前書7:1-7:17 & 鮑會園 & \hyperref[sec:570]{第七講 婚姻的生活} \\
第53屆 & 第8講 & 哥林多前書9:1-9:27 & 鮑會園 & \hyperref[sec:571]{第八講 屬靈的權柄} \\
\hline
\hline
第54屆 & 第1講 & 以西結書2:1-3:11 & 唐佑之 & \hyperref[sec:545]{第一講 苦澀與甘甜} \\
第54屆 & 第2講 & 以西結書9:3-9:19 & 唐佑之 & \hyperref[sec:546]{第二講 離愁與復合} \\
第54屆 & 第3講 &  & 唐佑之 & \hyperref[sec:547]{第三講虛 謊與真理} \\
第54屆 & 第4講 & 以西結書15:1-15:8 & 唐佑之 & \hyperref[sec:548]{第四講 脆弱與堅韌} \\
第54屆 & 第5講 & 以西結書18:1-18:4 & 唐佑之 & \hyperref[sec:549]{第五講 生命的定律} \\
第54屆 & 第6講 & 以西結書33:1-33:9 & 唐佑之 & \hyperref[sec:550]{第六講 守望的吶喊} \\
第54屆 & 第7講 & 以西結書34:1-34:6 & 唐佑之 & \hyperref[sec:551]{第七講 牧者的哀歌} \\
第54屆 & 第8講 & 以西結書37:1-37:10 & 唐佑之 & \hyperref[sec:552]{第八講 復興的景象} \\
第54屆 & 第9講 & 以西結書47:1-47:12 & 唐佑之 & \hyperref[sec:553]{第九講 增長的影響} \\
第54屆 & 第10講 & 以西結書1:1-1:3 & 唐佑之 & \hyperref[sec:554]{第十講 歷史的進程} \\
第54屆 & 第1講 &  & 張慕皚 & \hyperref[sec:536]{第一講 教會–基督的身體} \\
第54屆 & 第2講 &  & 張慕皚 & \hyperref[sec:537]{第二講 教會–基督的新婦} \\
第54屆 & 第3講 &  & 張慕皚 & \hyperref[sec:538]{第三講 教會–教會的牧養(一)} \\
第54屆 & 第4講 &  & 張慕皚 & \hyperref[sec:540]{第四講 教會的牧養(二)} \\
第54屆 & 第5講 &  & 張慕皚 & \hyperref[sec:541]{第五講 教會的見證} \\
第54屆 & 第6講 &  & 張慕皚 & \hyperref[sec:542]{第六講 天國的策略} \\
第54屆 & 第7講 &  & 張慕皚 & \hyperref[sec:543]{第七講 教會–聖靈的居所} \\
第54屆 & 第8講 &  & 張慕皚 & \hyperref[sec:544]{第八講 教會的使命} \\
第54屆 & 第1講 &  & 馬有藻 & \hyperref[sec:528]{第一講 馬太福音導論} \\
第54屆 & 第2講 &  & 馬有藻 & \hyperref[sec:529]{第二講 以色列的王來了} \\
第54屆 & 第3講 &  & 馬有藻 & \hyperref[sec:530]{第三講 以色列的王來了(續講一)} \\
第54屆 & 第4講 & 馬太福音5:17-7:29 & 馬有藻 & \hyperref[sec:531]{第四講 以色列的國度來了(續講二)} \\
第54屆 & 第5講 &  & 馬有藻 & \hyperref[sec:532]{第五講 以色列的王被棄絕了} \\
第54屆 & 第6講 & 馬太福音15:1-15:39 & 馬有藻 & \hyperref[sec:533]{第六講 以色列的王被棄絕了(續講一)} \\
第54屆 & 第7講 & 馬太福音19:27-20:16 & 馬有藻 & \hyperref[sec:534]{第七講 以色列的國度被奪去了} \\
第54屆 & 第8講 &  & 馬有藻 & \hyperref[sec:535]{第八講 以色列的國度被奪去了(續講)} \\
\hline
\hline
第55屆 & 第1講 & 傳道書1:1-1:11 & 劉承業 & \hyperref[sec:501]{第一講 導論與引言} \\
第55屆 & 第2講 & 傳道書1:12-2:26 & 劉承業 & \hyperref[sec:502]{第二講　體驗人生的虛空} \\
第55屆 & 第3講 & 傳道書3:1-5:7 & 劉承業 & \hyperref[sec:503]{第三講　觀察人生的虛空} \\
第55屆 & 第4講 & 傳道書5:8-6:12 & 劉承業 & \hyperref[sec:504]{第四講　福祿壽仍屬虛空} \\
第55屆 & 第5講 & 傳道書7:1-7:29 & 劉承業 & \hyperref[sec:505]{第五講　智慧人的體驗} \\
第55屆 & 第6講 & 傳道書8:1-9:12 & 劉承業 & \hyperref[sec:506]{第六講　智慧人接納神的主權} \\
第55屆 & 第7講 & 傳道書9:3-10:20 & 劉承業 & \hyperref[sec:507]{第七講　智慧面面觀} \\
第55屆 & 第8講 & 傳道書11:1-12:7 & 劉承業 & \hyperref[sec:508]{第八講　當未雨綢繆} \\
第55屆 & 第9講 & 傳道書12:8-12:14 & 劉承業 & \hyperref[sec:509]{第九講　總結 敬畏神的人生} \\
第55屆 & 第1講 & 啟示錄2:1-3:22 & 史保羅 & \hyperref[sec:518]{第一講　教會何時需要復興} \\
第55屆 & 第2講 & 歷代志下7:12-7:14 & 史保羅 & \hyperref[sec:519]{第二講　教會如何可以復興} \\
第55屆 & 第3講 &  & 史保羅 & \hyperref[sec:520]{第三講　教會復興帶來甚麼} \\
第55屆 & 第4講 &  & 史保羅 & \hyperref[sec:521]{第四講　聖靈與信徒} \\
第55屆 & 第5講 &  & 史保羅 & \hyperref[sec:522]{第五講　為何未被聖靈充滿} \\
第55屆 & 第6講 &  & 史保羅 & \hyperref[sec:523]{第六講　聖靈在信徒生命中的工作} \\
第55屆 & 第7講 &  & 史保羅 & \hyperref[sec:524]{第七講　如何引致聖靈擔憂} \\
第55屆 & 第8講 &  & 史保羅 & \hyperref[sec:525]{第八講　夜半呼聲} \\
第55屆 & 第9講 & 以西結書3:17-3:19 & 史保羅 & \hyperref[sec:526]{第九講　對冷淡信徒應負的責任} \\
第55屆 & 第10講 & 箴言22:15-22:15 & 史保羅 & \hyperref[sec:527]{第十講　如何建立基督化的家庭} \\
第55屆 & 第1講 & 哥林多前書14:40-14:40 & 張子華 & \hyperref[sec:510]{第一講　神對教會的呼召} \\
第55屆 & 第2講 & 哥林多前書1:10-1:17 & 張子華 & \hyperref[sec:511]{第二講　教會的紛爭} \\
第55屆 & 第3講 & 哥林多前書6:1-6:11 & 張子華 & \hyperref[sec:512]{第三講　教會與法院} \\
第55屆 & 第4講 & 哥林多前書1:18-4:5 & 張子華 & \hyperref[sec:539]{第四講　神的智慧} \\
第55屆 & 第5講 & 哥林多前書5:1-5:13 & 張子華 & \hyperref[sec:513]{第五講　教會在倫理上的破產與管教} \\
第55屆 & 第6講 & 哥林多前書7:1-7:40 & 張子華 & \hyperref[sec:514]{第六講　基督徒的獨身和婚姻} \\
第55屆 & 第7講 & 哥林多前書10:1-10:13 & 張子華 & \hyperref[sec:515]{第七講　基督徒在恩典中所享受的自由} \\
第55屆 & 第8講 & 哥林多前書8:1-8:13 & 張子華 & \hyperref[sec:516]{第八講　基督的自由和社交生活} \\
第55屆 & 第9講 & 哥林多前書11:1-11:34 & 張子華 & \hyperref[sec:517]{第九講　敬虔的崇拜生活} \\
\hline
\hline
第56屆 & 第1講 & 俄巴底亞書1:1-1:1 & 林克己 & \hyperref[sec:368]{第一講　耶和華的呼召(一節)} \\
第56屆 & 第2講 & 俄巴底亞書1:1-1:2 & 林克己 & \hyperref[sec:370]{第二講　耶和華的主權} \\
第56屆 & 第3講 & 俄巴底亞書1:3-1:4 & 林克己 & \hyperref[sec:372]{第三講　耶和華的偉大} \\
第56屆 & 第4講 & 俄巴底亞書1:5-1:6 & 林克己 & \hyperref[sec:373]{第四講　耶和華的提醒} \\
第56屆 & 第5講 & 俄巴底亞書1:7-1:7 & 林克己 & \hyperref[sec:375]{第五講　耶和華的警告} \\
第56屆 & 第6講 & 俄巴底亞書1:8-1:9 & 林克己 & \hyperref[sec:376]{第六講　耶和華的智慧} \\
第56屆 & 第7講 & 俄巴底亞書1:10-1:14 & 林克己 & \hyperref[sec:378]{第七講　耶和華的責備} \\
第56屆 & 第8講 & 俄巴底亞書1:15-1:16 & 林克己 & \hyperref[sec:379]{第八講　耶和華的審判} \\
第56屆 & 第9講 & 俄巴底亞書1:17-1:21 & 林克己 & \hyperref[sec:381]{第九講　耶和華的國度} \\
第56屆 & 第1講 &  & 楊濬哲 & \hyperref[sec:1387]{序} \\
第56屆 & 第1講 & 啟示錄4:1-4:8 & 高文 & \hyperref[sec:402]{第一講　神的尊榮} \\
第56屆 & 第2講 & 啟示錄4:9-4:10 & 高文 & \hyperref[sec:404]{第二講　金的冠冕} \\
第56屆 & 第3講 & 啟示錄5:1-5:10 & 高文 & \hyperref[sec:406]{第三講　寶血買贖} \\
第56屆 & 第4講 & 啟示錄5:11-5:14 & 高文 & \hyperref[sec:407]{第四講　讚美的巔峰} \\
第56屆 & 第5講 & 啟示錄6:9-6:11 & 高文 & \hyperref[sec:409]{第五講　苦難為何延續} \\
第56屆 & 第6講 & 啟示錄7:9-7:10 & 高文 & \hyperref[sec:411]{第六講　神的拯救} \\
第56屆 & 第7講 & 啟示錄11:15-11:18 & 高文 & \hyperref[sec:412]{第七講　審判的確實} \\
第56屆 & 第8講 & 啟示錄12:10-12:12 & 高文 & \hyperref[sec:413]{第八講　得勝者的凱歌} \\
第56屆 & 第9講 & 啟示錄15:2-15:4 & 高文 & \hyperref[sec:415]{第九講　萬民的聚集} \\
第56屆 & 第10講 & 啟示錄19:1-19:9 & 高文 & \hyperref[sec:416]{第十講　羔羊的婚筵} \\
第56屆 & 第1講 &  & 麥希真 & \hyperref[sec:382]{第一講　信神,信人,自信?} \\
第56屆 & 第2講 &  & 麥希真 & \hyperref[sec:384]{第二講　一個平凡基督徒的禱告生活} \\
第56屆 & 第3講 &  & 麥希真 & \hyperref[sec:386]{第三講　如何讀經,而且實行出來} \\
第56屆 & 第4講 &  & 麥希真 & \hyperref[sec:388]{第四講　今天如何能消除緊張} \\
第56屆 & 第5講 &  & 麥希真 & \hyperref[sec:389]{第五講　明天如何能保持堅忍} \\
第56屆 & 第6講 &  & 麥希真 & \hyperref[sec:391]{第六講　最有效的細胞小組 -- 家} \\
第56屆 & 第7講 &  & 麥希真 & \hyperref[sec:392]{第七講　金錢讓我們正視它} \\
第56屆 & 第8講 &  & 麥希真 & \hyperref[sec:396]{第八講　我們承擔了一個沉重的任務} \\
第56屆 & 第9講 &  & 麥希真 & \hyperref[sec:398]{第九講　成年人,還是青年人?} \\
\hline
\hline
第57屆 & 第1講 & 詩篇103:5-103:2 & 滕近輝 & \hyperref[sec:491]{第一講　詩篇中七幅復興的圖畫} \\
第57屆 & 第2講 &  & 滕近輝 & \hyperref[sec:492]{第二講　詩人奇妙的經歷} \\
第57屆 & 第3講 & 詩篇90:6-90:17 & 滕近輝 & \hyperref[sec:493]{第三講　詩篇中關於事奉的信息} \\
第57屆 & 第4講 & 詩篇119:1-119:176 & 滕近輝 & \hyperref[sec:494]{第四講　詩篇中神言的反合性功效} \\
第57屆 & 第5講 & 詩篇22:1-22:1 & 滕近輝 & \hyperref[sec:495]{第五講　詩入與神的深交} \\
第57屆 & 第6講 & 詩篇18:28-18:2 & 滕近輝 & \hyperref[sec:496]{第六講　詩篇中關於禱告的比喻} \\
第57屆 & 第7講 & 詩篇19:1-19:14 & 滕近輝 & \hyperref[sec:497]{第七講　神的四種言語} \\
第57屆 & 第8講 & 詩篇106:5-106:19 & 滕近輝 & \hyperref[sec:498]{第八講　詩人的喜樂} \\
第57屆 & 第9講 & 詩篇84:1-84:2 & 滕近輝 & \hyperref[sec:500]{第九講　熱愛聖殿聖城的詩篇} \\
第57屆 & 第10講 & 詩篇122:1-122:9 & 滕近輝 & \hyperref[sec:499]{第十講　熱愛聖殿聖城的詩篇} \\
第57屆 & 第1講 & 羅馬書8:1-8:4 & 艾榮 & \hyperref[sec:482]{第一講　如今就不被定罪了} \\
第57屆 & 第2講 & 羅馬書8:1-8:4 & 艾榮 & \hyperref[sec:483]{第二講　公義的實踐} \\
第57屆 & 第3講 & 羅馬書8:5-8:8 & 艾榮 & \hyperref[sec:484]{第三講　公義與我們的頭腦} \\
第57屆 & 第4講 & 羅馬書8:9-8:13 & 艾榮 & \hyperref[sec:485]{第四講　煩惱的身體} \\
第57屆 & 第5講 & 羅馬書8:14-8:18 & 艾榮 & \hyperref[sec:486]{第五講　神的兒女} \\
第57屆 & 第6講 & 羅馬書8:19-8:25 & 艾榮 & \hyperref[sec:487]{第六講　誠懇的期望} \\
第57屆 & 第7講 & 羅馬書8:26-8:27 & 艾榮 & \hyperref[sec:488]{第七講　聖靈的時代} \\
第57屆 & 第8講 & 羅馬書8:28-8:30 & 艾榮 & \hyperref[sec:489]{第八講　神為人的計劃} \\
第57屆 & 第9講 & 羅馬書8:31-8:39 & 艾榮 & \hyperref[sec:490]{第九講　得勝有餘} \\
第57屆 & 第1講 &  & 鮑會園 & \hyperref[sec:472]{第一講　摩利沙人彌迦──時代的先知} \\
第57屆 & 第2講 & 彌迦書1:1-1:7 & 鮑會園 & \hyperref[sec:473]{第二講　信息的權威} \\
第57屆 & 第3講 & 彌迦書1:8-1:16 & 鮑會園 & \hyperref[sec:475]{第三講　擊打和醫治} \\
第57屆 & 第4講 & 彌迦書2:1-2:13 & 鮑會園 & \hyperref[sec:476]{第四講　罪中的大罪} \\
第57屆 & 第5講 & 彌迦書3:1-3:12 & 鮑會園 & \hyperref[sec:477]{第五講　假先知的路} \\
第57屆 & 第6講 & 彌迦書4:1-5:4 & 鮑會園 & \hyperref[sec:478]{第六講　榮耀的盼望} \\
第57屆 & 第7講 & 彌迦書5:5-6:5 & 鮑會園 & \hyperref[sec:479]{第七講　王的被拒絕} \\
第57屆 & 第8講 & 彌迦書6:6-6:16 & 鮑會園 & \hyperref[sec:480]{第八講　清楚的準則} \\
第57屆 & 第9講 & 彌迦書7:1-7:20 & 鮑會園 & \hyperref[sec:481]{第九講　奮興的道路} \\
\hline
\hline
第58屆 & 第1講 &  & 吳勇 & \hyperref[sec:463]{第一講 教會的破口} \\
第58屆 & 第2講 & 瑪拉基書3:2-3:2 & 吳勇 & \hyperref[sec:464]{第二講 沒落的教會} \\
第58屆 & 第3講 &  & 吳勇 & \hyperref[sec:465]{第三講 隨從世俗的教會} \\
第58屆 & 第4講 &  & 吳勇 & \hyperref[sec:767]{第四講 教會的大使命} \\
第58屆 & 第5講 &  & 吳勇 & \hyperref[sec:467]{第五講 教會的人才} \\
第58屆 & 第6講 &  & 吳勇 & \hyperref[sec:468]{第六講 聖靈的果子} \\
第58屆 & 第7講 &  & 吳勇 & \hyperref[sec:469]{第七講 聖靈的能力} \\
第58屆 & 第8講 &  & 吳勇 & \hyperref[sec:470]{第八講 充滿聖靈} \\
第58屆 & 第9講 &  & 吳勇 & \hyperref[sec:471]{第九講 聖靈的工作} \\
第58屆 & 第1講 & 提摩太後書3:1-3:17 & 張慕皚 & \hyperref[sec:455]{第一講 基督教「時間觀」} \\
第58屆 & 第2講 & 創世記1:26-1:28 & 張慕皚 & \hyperref[sec:456]{第二講 基督教「工作觀」} \\
第58屆 & 第3講 & 申命記5:12-5:15 & 張慕皚 & \hyperref[sec:457]{第三講 基督教「休閒觀」} \\
第58屆 & 第4講 & 提摩太前書6:6-6:10 & 張慕皚 & \hyperref[sec:458]{第四講 基督教「財富觀」} \\
第58屆 & 第5講 & 歷代志下1:1-1:17 & 張慕皚 & \hyperref[sec:459]{第五講 基督教「成功觀」} \\
第58屆 & 第6講 & 以弗所書5:22-5:22 & 張慕皚 & \hyperref[sec:460]{第六講 基督教「家庭觀」} \\
第58屆 & 第7講 & 羅馬書13:1-13:7 & 張慕皚 & \hyperref[sec:461]{第七講 基督教「政府觀」} \\
第58屆 & 第8講 & 約翰福音10:10-10:10 & 張慕皚 & \hyperref[sec:462]{第八講 基督教「豐盛人生觀」} \\
第58屆 & 第1講 & 路得記1:1-1:5 & 焦源濂 & \hyperref[sec:447]{第一講　領袖的鑑戒} \\
第58屆 & 第2講 & 路得記1:1-1:15 & 焦源濂 & \hyperref[sec:448]{第二講　好人的災難} \\
第58屆 & 第3講 & 路得記1:5-1:18 & 焦源濂 & \hyperref[sec:449]{第三講 神愛中的人} \\
第58屆 & 第4講 & 路得記1:8-1:18 & 焦源濂 & \hyperref[sec:450]{第四講 愛的跟隨} \\
第58屆 & 第5講 & 路得記1:15-1:18 & 焦源濂 & \hyperref[sec:451]{第五講 神的引導} \\
第58屆 & 第6講 & 路得記1:15-1:18 & 焦源濂 & \hyperref[sec:452]{第六講 希望的中心} \\
第58屆 & 第7講 & 路得記1:22-2:7 & 焦源濂 & \hyperref[sec:453]{第七講 安息的尋求} \\
第58屆 & 第8講 & 路得記2:16-2:21 & 焦源濂 & \hyperref[sec:454]{第八講 愛的真諦} \\
\hline
\hline
第59屆 & 第1講 &  & 戴紹曾 & \hyperref[sec:1269]{第一講 亞伯拉罕} \\
第59屆 & 第2講 &  & 戴紹曾 & \hyperref[sec:1271]{第二講 約瑟} \\
第59屆 & 第3講 &  & 戴紹曾 & \hyperref[sec:1274]{第三講 摩西的生平} \\
第59屆 & 第4講 &  & 戴紹曾 & \hyperref[sec:1276]{第四講 摩西的見證} \\
第59屆 & 第5講 &  & 戴紹曾 & \hyperref[sec:1278]{第五講 約書亞和迦勒} \\
第59屆 & 第6講 &  & 戴紹曾 & \hyperref[sec:1279]{第六講 撒母耳} \\
第59屆 & 第7講 &  & 戴紹曾 & \hyperref[sec:1280]{第七講 大衛} \\
第59屆 & 第8講 &  & 戴紹曾 & \hyperref[sec:1281]{第八講 以斯帖} \\
第59屆 & 第9講 &  & 戴紹曾 & \hyperref[sec:1282]{第九講 但以理} \\
第59屆 & 第10講 &  & 戴紹曾 & \hyperref[sec:1283]{第十講 尼希米} \\
第59屆 & 第1講 &  & 楊濬哲 & \hyperref[sec:1388]{序} \\
第59屆 & 第1講 &  & 許道良 & \hyperref[sec:1236]{第一講 似是綑綁} \\
第59屆 & 第2講 &  & 許道良 & \hyperref[sec:1237]{第二講 落在地裡的一粒麥子} \\
第59屆 & 第3講 &  & 許道良 & \hyperref[sec:1238]{第三講 主人如此 何況僕人} \\
第59屆 & 第4講 &  & 許道良 & \hyperref[sec:1239]{第四講 沒有間隔的羊圈} \\
第59屆 & 第5講 &  & 許道良 & \hyperref[sec:1240]{第五講 葡萄樹上的枝子} \\
第59屆 & 第6講 &  & 許道良 & \hyperref[sec:1243]{第六講 沒有選擇的命令} \\
第59屆 & 第7講 &  & 許道良 & \hyperref[sec:1245]{第七講 肩負使命的客旅} \\
第59屆 & 第8講 &  & 許道良 & \hyperref[sec:1249]{第八講 愛的十字架} \\
第59屆 & 第1講 &  & 金新宇 & \hyperref[sec:1251]{第一講 神榮耀的統治} \\
第59屆 & 第2講 &  & 金新宇 & \hyperref[sec:1254]{第二講 耶穌基督是王} \\
第59屆 & 第3講 &  & 金新宇 & \hyperref[sec:1256]{第三講 先求祂的國 — 國度的重要性} \\
第59屆 & 第4講 &  & 金新宇 & \hyperref[sec:1257]{第四講 神的國 — 現在和將來} \\
第59屆 & 第5講 &  & 金新宇 & \hyperref[sec:1261]{第五講 國度與教會} \\
第59屆 & 第6講 &  & 金新宇 & \hyperref[sec:1264]{第六講 怎樣進神的國} \\
第59屆 & 第7講 &  & 金新宇 & \hyperref[sec:1265]{第七講 誰是天國的子民} \\
第59屆 & 第8講 &  & 金新宇 & \hyperref[sec:1267]{第八講 願你的國降臨 — 國度與差傳} \\
\hline
\hline
第60屆 & 第1講 & 約書亞記1:1-1:9 & 劉承業 & \hyperref[sec:1193]{第一講 得勝的先決條件} \\
第60屆 & 第2講 & 約書亞記2:1-2:14 & 劉承業 & \hyperref[sec:1194]{第二講 信心的榜樣} \\
第60屆 & 第3講 & 約書亞記3:1-3:11 & 劉承業 & \hyperref[sec:1195]{第三講 憑信的戰爭} \\
第60屆 & 第4講 & 約書亞記5:1-5:9 & 劉承業 & \hyperref[sec:1196]{第四講 神要作元帥} \\
第60屆 & 第5講 & 約書亞記7:1-7:9 & 劉承業 & \hyperref[sec:1200]{第五講 聖潔的重要} \\
第60屆 & 第6講 & 以弗所書6:14-6:20 & 劉承業 & \hyperref[sec:1202]{第六講 屬靈爭戰的策略} \\
第60屆 & 第7講 & 約書亞記11:16-11:23 & 劉承業 & \hyperref[sec:1205]{第七講 攻守的並重} \\
第60屆 & 第8講 &  & 劉承業 & \hyperref[sec:1209]{第八講 領袖的質素} \\
第60屆 & 第9講 &  & 劉承業 & \hyperref[sec:1211]{第九講 同心與配搭} \\
第60屆 & 第1講 & 以賽亞書40:1-40:10 & 唐佑之 & \hyperref[sec:1216]{第一講 安慰} \\
第60屆 & 第2講 & 以賽亞書41:1-41:29 & 唐佑之 & \hyperref[sec:1220]{第二講 重新得力} \\
第60屆 & 第3講 & 以賽亞書43:1-43:28 & 唐佑之 & \hyperref[sec:1222]{第三講 見證} \\
第60屆 & 第4講 & 以賽亞書44:1-49:26 & 唐佑之 & \hyperref[sec:1223]{第四講 裝備} \\
第60屆 & 第5講 & 以賽亞書50:1-50:11 & 唐佑之 & \hyperref[sec:1224]{第五講 待命} \\
第60屆 & 第6講 & 以賽亞書51:1-52:15 & 唐佑之 & \hyperref[sec:1226]{第六講 更新} \\
第60屆 & 第7講 & 以賽亞書53:1-53:12 & 唐佑之 & \hyperref[sec:1228]{第七講 承受} \\
第60屆 & 第8講 & 以賽亞書53:1-53:12 & 唐佑之 & \hyperref[sec:1229]{第八講 捨上} \\
第60屆 & 第9講 & 以賽亞書54:1-54:17 & 唐佑之 & \hyperref[sec:1231]{第九講 展望} \\
第60屆 & 第10講 & 以賽亞書55:1-55:13 & 唐佑之 & \hyperref[sec:1233]{第十講 歌頌} \\
第60屆 & 第1講 &  & 楊濬哲 & \hyperref[sec:1389]{序} \\
第60屆 & 第1講 & 創世記4:1-4:8 & 陳終道 & \hyperref[sec:1172]{第一講 第一個殉道士亞伯} \\
第60屆 & 第2講 & 創世記5:28-5:29 & 陳終道 & \hyperref[sec:1173]{第二講 劃時代的挪亞} \\
第60屆 & 第3講 & 創世記37:1-37:36 & 陳終道 & \hyperref[sec:1174]{第三講 隨處得勝的約瑟} \\
第60屆 & 第4講 & 民數記12:1-12:9 & 陳終道 & \hyperref[sec:1175]{第四講 神所尊重的摩西} \\
第60屆 & 第5講 & 以賽亞書6:1-6:13 & 陳終道 & \hyperref[sec:1176]{第五講 看見異象的以賽亞} \\
第60屆 & 第6講 & 馬太福音3:1-3:17 & 陳終道 & \hyperref[sec:1177]{第六講 嫉惡如仇的施洗約翰} \\
第60屆 & 第7講 & 使徒行傳6:3-6:6 & 陳終道 & \hyperref[sec:1178]{第七講 勇敢殉道的司提反} \\
第60屆 & 第8講 & 使徒行傳9:1-9:9 & 陳終道 & \hyperref[sec:1179]{第八講 外邦使徒保羅} \\
第60屆 & 第9講 & 路加福音1:38-1:12 & 陳終道 & \hyperref[sec:1180]{第九講 永垂千古的平凡人} \\
\hline
\hline
第61屆 & 第1講 & 使徒行傳1:14-1:14 & 劉銳光 & \hyperref[sec:1124]{第一講 同心} \\
第61屆 & 第2講 & 使徒行傳2:43-2:47 & 劉銳光 & \hyperref[sec:1125]{第二講 相愛} \\
第61屆 & 第3講 & 使徒行傳1:14-1:15 & 劉銳光 & \hyperref[sec:1127]{第三講 禱告} \\
第61屆 & 第4講 & 使徒行傳1:8-1:33 & 劉銳光 & \hyperref[sec:1130]{第四講 能力} \\
第61屆 & 第5講 & 使徒行傳1:21-1:22 & 劉銳光 & \hyperref[sec:1132]{第五講 見證} \\
第61屆 & 第6講 & 使徒行傳2:46-2:47 & 劉銳光 & \hyperref[sec:1133]{第六講 喜樂} \\
第61屆 & 第7講 & 使徒行傳9:20-9:22 & 劉銳光 & \hyperref[sec:1135]{第七講 信仰純正} \\
第61屆 & 第8講 & 使徒行傳8:25-8:7 & 劉銳光 & \hyperref[sec:1138]{第八講 傳福音} \\
第61屆 & 第9講 & 使徒行傳1:2-1:2 & 劉銳光 & \hyperref[sec:1139]{第九講 聖靈} \\
第61屆 & 第1講 &  & 張子華 & \hyperref[sec:1115]{第一講 導論} \\
第61屆 & 第2講 & 羅馬書4:21-4:4 & 張子華 & \hyperref[sec:1116]{第二講 亞伯拉罕的信心 (一)} \\
第61屆 & 第3講 &  & 張子華 & \hyperref[sec:1117]{第三講 亞伯拉罕的信心 (二)} \\
第61屆 & 第4講 &  & 張子華 & \hyperref[sec:1118]{第四講 亞伯拉罕的信心 (三)} \\
第61屆 & 第5講 &  & 張子華 & \hyperref[sec:1119]{第五講 亞伯拉罕的信心 (四)} \\
第61屆 & 第6講 &  & 張子華 & \hyperref[sec:1120]{第六講 以撒的順服} \\
第61屆 & 第7講 & 創世記32:22-32:30 & 張子華 & \hyperref[sec:1121]{第七講 雅各的真我作為} \\
第61屆 & 第8講 &  & 張子華 & \hyperref[sec:1122]{第八講 約瑟的升高} \\
第61屆 & 第9講 &  & 張子華 & \hyperref[sec:1123]{第九講 約瑟的完美生命} \\
第61屆 & 第1講 & 哥林多前書1:1-1:7 & 韓力生 & \hyperref[sec:1141]{第一講 請看新約教會 (哥林多前書概覽)} \\
第61屆 & 第2講 & 哥林多前書1:10-1:31 & 韓力生 & \hyperref[sec:1145]{第二講 人如何尋找神 (上)} \\
第61屆 & 第3講 & 哥林多前書2:1-2:16 & 韓力生 & \hyperref[sec:1146]{第三講 人如何尋找神 (下)} \\
第61屆 & 第4講 & 哥林多前書3:1-3:23 & 韓力生 & \hyperref[sec:1147]{第四講 建造的代價} \\
第61屆 & 第5講 & 哥林多前書4:1-4:21 & 韓力生 & \hyperref[sec:1149]{第五講 明確的判斷 — 動機} \\
第61屆 & 第6講 & 哥林多前書5:1-5:13 & 韓力生 & \hyperref[sec:1151]{第六講 明確的判斷 — 行為} \\
第61屆 & 第7講 & 哥林多前書6:1-6:20 & 韓力生 & \hyperref[sec:1153]{第七講 明確的判斷 — 爭執} \\
第61屆 & 第8講 & 哥林多前書7:1-7:40 & 韓力生 & \hyperref[sec:1155]{第八講 獨身或結婚 — 神的呼召} \\
第61屆 & 第9講 & 哥林多前書12:1-12:31 & 韓力生 & \hyperref[sec:1159]{第九講 恩賜的分配及規則} \\
第61屆 & 第10講 & 哥林多前書13:1-13:13 & 韓力生 & \hyperref[sec:1162]{第十講 神聖的潤滑劑} \\
\hline
\hline
第62屆 & 第1講 &  & 劉少康 & \hyperref[sec:1085]{第一講 重建的時機 — 尼希米的呼籲} \\
第62屆 & 第2講 &  & 劉少康 & \hyperref[sec:1086]{第二講 見證的重建 — 尼希米的領導} \\
第62屆 & 第3講 &  & 劉少康 & \hyperref[sec:1087]{第三講 敬拜的重建 — 以斯拉的行動} \\
第62屆 & 第4講 &  & 劉少康 & \hyperref[sec:1088]{第四講 心靈的重建 — 以斯拉的心志} \\
第62屆 & 第5講 &  & 劉少康 & \hyperref[sec:1089]{第五講 生活的重建 — 以斯拉的責備} \\
第62屆 & 第6講 &  & 劉少康 & \hyperref[sec:1090]{第六講 委身的重建 — 以斯帖的榜樣} \\
第62屆 & 第7講 &  & 劉少康 & \hyperref[sec:1091]{第七講 優次的重建 — 哈該的信息} \\
第62屆 & 第8講 &  & 劉少康 & \hyperref[sec:1092]{第八講 信心的重建 — 撒迦利亞的挑戰} \\
第62屆 & 第9講 &  & 劉少康 & \hyperref[sec:1093]{第九講 愛心的重建 — 瑪拉基的勸勉} \\
第62屆 & 第1講 & 詩篇51:1-51:19 & 杜偉力 & \hyperref[sec:1094]{第一講 往下扎根} \\
第62屆 & 第2講 &  & 杜偉力 & \hyperref[sec:1096]{第二講 往下扎根} \\
第62屆 & 第3講 &  & 杜偉力 & \hyperref[sec:1098]{第三講 往下扎根} \\
第62屆 & 第4講 &  & 杜偉力 & \hyperref[sec:1102]{第四講 往下扎根} \\
第62屆 & 第5講 &  & 杜偉力 & \hyperref[sec:1105]{第五講 向上結果} \\
第62屆 & 第6講 &  & 杜偉力 & \hyperref[sec:1107]{第六講 向上結果} \\
第62屆 & 第7講 &  & 杜偉力 & \hyperref[sec:1109]{第七講 向上結果} \\
第62屆 & 第8講 &  & 杜偉力 & \hyperref[sec:1112]{第八講 向上結果} \\
第62屆 & 第9講 &  & 杜偉力 & \hyperref[sec:1113]{第九講 向上結果} \\
第62屆 & 第10講 & 詩篇51:1-51:19 & 杜偉力 & \hyperref[sec:1114]{第十講 向上結果} \\
第62屆 & 第1講 &  & 楊濬哲 & \hyperref[sec:1284]{序} \\
第62屆 & 第1講 & 腓立比書1:9-1:19 & 鮑會園 & \hyperref[sec:1076]{第一講 現實的智慧──分辨是非的禱告} \\
第62屆 & 第2講 & 腓立比書1:20-1:30 & 鮑會園 & \hyperref[sec:1077]{第二講 最好的選擇──活著就是基督} \\
第62屆 & 第3講 & 腓立比書2:1-2:4 & 鮑會園 & \hyperref[sec:1078]{第三講 努力的方向──同心合意的事奉} \\
第62屆 & 第4講 & 腓立比書2:5-2:11 & 鮑會園 & \hyperref[sec:1079]{第四講 完全的榜樣──以基督的心為心} \\
第62屆 & 第5講 & 腓立比書2:12-2:18 & 鮑會園 & \hyperref[sec:1080]{第五講 工作的目的──作成得救的工夫} \\
第62屆 & 第6講 & 腓立比書3:1-3:11 & 鮑會園 & \hyperref[sec:1081]{第六講 最高的願望 — 認識基督和祂的大能} \\
第62屆 & 第7講 & 腓立比書3:12-3:16 & 鮑會園 & \hyperref[sec:1082]{第七講 追求的目標 — 凡事作完全人} \\
第62屆 & 第8講 & 腓立比書3:17-4:3 & 鮑會園 & \hyperref[sec:1083]{第八講 生活的盼望 — 作天上的國民} \\
第62屆 & 第9講 & 腓立比書4:4-4:13 & 鮑會園 & \hyperref[sec:1084]{第九講 堅守的動力 — 靠主站立得穩} \\
\hline
\hline
第63屆 & 第1講 & 以弗所書1:9-1:9 & 劉東崑 & \hyperref[sec:1035]{第一講 神為何創造宇宙萬物} \\
第63屆 & 第2講 & 以弗所書3:2-3:5 & 劉東崑 & \hyperref[sec:1036]{第二講 神如何創造宇宙萬物} \\
第63屆 & 第3講 & 以弗所書3:9-3:11 & 劉東崑 & \hyperref[sec:1037]{第三講 神對宇宙萬有作如何安排} \\
第63屆 & 第4講 & 以賽亞書43:7-43:7 & 劉東崑 & \hyperref[sec:1038]{第四講 神為何創造人類} \\
第63屆 & 第5講 &  & 劉東崑 & \hyperref[sec:1039]{第五講 神如何創造人類} \\
第63屆 & 第6講 & 以弗所書1:9-1:11 & 劉東崑 & \hyperref[sec:1040]{第六講 神在人身上有預定的美意} \\
第63屆 & 第7講 & 馬太福音22:14-22:14 & 劉東崑 & \hyperref[sec:1041]{第七講 神如何成就祂的美意} \\
第63屆 & 第8講 & 傳道書3:1-3:2 & 劉東崑 & \hyperref[sec:1042]{第八講 神安排時間,祂如何藉暫時成就永遠} \\
第63屆 & 第9講 & 以弗所書3:2-3:11 & 劉東崑 & \hyperref[sec:1043]{第九講 人如何經歷時間進入永遠} \\
第63屆 & 第1講 &  & 楊濬哲 & \hyperref[sec:1392]{序} \\
第63屆 & 第1講 &  & 麥希真 & \hyperref[sec:1066]{第一講 基督的原初} \\
第63屆 & 第2講 &  & 麥希真 & \hyperref[sec:1067]{第二講 基督的門徒} \\
第63屆 & 第3講 &  & 麥希真 & \hyperref[sec:1068]{第三講 基督的天國} \\
第63屆 & 第4講 &  & 麥希真 & \hyperref[sec:1069]{第四講 基督的教會} \\
第63屆 & 第5講 &  & 麥希真 & \hyperref[sec:1070]{第五講 基督的教導} \\
第63屆 & 第6講 &  & 麥希真 & \hyperref[sec:1071]{第六講 基督的愛顧} \\
第63屆 & 第7講 &  & 麥希真 & \hyperref[sec:1072]{第七講 基督的審判} \\
第63屆 & 第8講 &  & 麥希真 & \hyperref[sec:1073]{第八講 基督的再來} \\
第63屆 & 第9講 &  & 麥希真 & \hyperref[sec:1074]{第九講 基督的受死} \\
第63屆 & 第10講 &  & 麥希真 & \hyperref[sec:1075]{第十講 基督的復活} \\
第63屆 & 第1講 & 哥林多後書5:9-5:10 & 黃子嘉 & \hyperref[sec:1044]{第一講 盼望的生活} \\
第63屆 & 第2講 & 希伯來書11:39-11:39 & 黃子嘉 & \hyperref[sec:1045]{第二講 信心的生活} \\
第63屆 & 第3講 & 約翰一書3:16-3:23 & 黃子嘉 & \hyperref[sec:1046]{第三講 愛心的生活} \\
第63屆 & 第4講 & 羅馬書1:7-1:7 & 黃子嘉 & \hyperref[sec:1047]{第四講 成聖的生活} \\
第63屆 & 第5講 & 約翰福音16:33-16:39 & 黃子嘉 & \hyperref[sec:1048]{第五講 得勝的生活} \\
第63屆 & 第6講 & 羅馬書12:1-12:2 & 黃子嘉 & \hyperref[sec:1049]{第六講 奉獻的生活} \\
第63屆 & 第7講 & 羅馬書12:1-12:8 & 黃子嘉 & \hyperref[sec:1063]{第七講 身體的生活} \\
第63屆 & 第8講 & 歌羅西書1:9-1:10 & 黃子嘉 & \hyperref[sec:1064]{第八講 結果子的生活} \\
第63屆 & 第9講 & 馬太福音12:18-12:40 & 黃子嘉 & \hyperref[sec:1065]{第九講 宣教的生活} \\
\hline
\hline
第64屆 & 第1講 & 馬太福音5:1-5:3 & 張慕皚 & \hyperref[sec:1017]{第一講 虛心 — 靈命成長的起點} \\
第64屆 & 第2講 & 馬太福音5:4-5:4 & 張慕皚 & \hyperref[sec:1018]{第二講 哀慟 — 洞悉人性的真面目} \\
第64屆 & 第3講 & 馬太福音5:5-5:5 & 張慕皚 & \hyperref[sec:1019]{第三講 溫柔 — 合乎中度的自我形象} \\
第64屆 & 第4講 & 馬太福音5:6-5:6 & 張慕皚 & \hyperref[sec:1020]{第四講 飢渴慕義 — 尋求活在神的旨意中} \\
第64屆 & 第5講 & 馬太福音5:2-5:2 & 張慕皚 & \hyperref[sec:1021]{第五講 憐恤 — 愛的實踐} \\
第64屆 & 第6講 & 馬太福音5:8-5:8 & 張慕皚 & \hyperref[sec:1022]{第六講 清心 — 邁向屬靈最高境界} \\
第64屆 & 第7講 & 馬太福音5:9-5:9 & 張慕皚 & \hyperref[sec:1023]{第七講 使人和睦 — 跟隨基督的腳蹤行} \\
第64屆 & 第8講 & 馬太福音5:10-5:12 & 張慕皚 & \hyperref[sec:1024]{第八講 為義受逼迫 — 背十字架跟隨主} \\
第64屆 & 第1講 & 路加福音2:25-2:38 & 戴紹曾 & \hyperref[sec:1025]{第一講 誠心盼望得見耶穌} \\
第64屆 & 第2講 & 馬太福音2:1-2:12 & 戴紹曾 & \hyperref[sec:1026]{第二講 博士遠來拜見耶穌} \\
第64屆 & 第3講 & 約翰福音1:29-1:34 & 戴紹曾 & \hyperref[sec:1027]{第三講 先鋒引人候見耶穌} \\
第64屆 & 第4講 & 約翰福音3:1-3:21 & 戴紹曾 & \hyperref[sec:1028]{第四講 夜裡先生訪見耶穌} \\
第64屆 & 第5講 & 約翰福音4:27-4:42 & 戴紹曾 & \hyperref[sec:1029]{第五講 井旁婦人遇見耶穌} \\
第64屆 & 第6講 & 馬可福音9:1-9:8 & 戴紹曾 & \hyperref[sec:1030]{第六講 變像山上只見耶穌} \\
第64屆 & 第7講 & 馬可福音10:46-10:52 & 戴紹曾 & \hyperref[sec:1031]{第七講 討飯瞎子看見耶穌} \\
第64屆 & 第8講 & 路加福音19:1-19:10 & 戴紹曾 & \hyperref[sec:1032]{第八講 樹上撒該想見耶穌} \\
第64屆 & 第9講 & 約翰福音12:20-12:36 & 戴紹曾 & \hyperref[sec:1033]{第九講 希利尼人求見耶穌} \\
第64屆 & 第10講 & 約翰福音20:1-20:23 & 戴紹曾 & \hyperref[sec:1034]{第十講 復活清晨去見耶穌} \\
第64屆 & 第1講 &  & 楊濬哲 & \hyperref[sec:1391]{序} \\
第64屆 & 第1講 &  & 許道良 & \hyperref[sec:1009]{第一講 曙光遙望} \\
第64屆 & 第2講 & 彼得前書1:16-1:16 & 許道良 & \hyperref[sec:1010]{第二講 手潔心清} \\
第64屆 & 第3講 &  & 許道良 & \hyperref[sec:1011]{第三講 擇善棄惡} \\
第64屆 & 第4講 &  & 許道良 & \hyperref[sec:1012]{第四講 贏取汝心} \\
第64屆 & 第5講 &  & 許道良 & \hyperref[sec:1013]{第五講 為義受苦} \\
第64屆 & 第6講 &  & 許道良 & \hyperref[sec:1014]{第六講 責無旁貸} \\
第64屆 & 第7講 &  & 許道良 & \hyperref[sec:1015]{第七講 善牧楷模} \\
第64屆 & 第8講 &  & 許道良 & \hyperref[sec:1016]{第八講 步步為營} \\
\hline
\hline
第65屆 & 第1講 & 出埃及記1:1-2:25 & 周永健 & \hyperref[sec:983]{第一講 神顧念我們──我下來是要救他們} \\
第65屆 & 第2講 & 出埃及記3:1-6:30 & 周永健 & \hyperref[sec:984]{第二講 神呼召我們事奉祂──我要打發你們去見法老} \\
第65屆 & 第3講 & 出埃及記7:1-12:51 & 周永健 & \hyperref[sec:985]{第三講 神用大能的手行奇事──埃及人就知道我是耶和華} \\
第65屆 & 第4講 & 出埃及記14:1-15:27 & 周永健 & \hyperref[sec:986]{第四講 神使我們得勝──耶和華必為你們爭戰} \\
第65屆 & 第5講 & 出埃及記19:1-19:25 & 周永健 & \hyperref[sec:987]{第五講 神將我們分別為聖──如鷹將你們背在翅膀上帶來歸我} \\
第65屆 & 第6講 & 出埃及記15:1-17:16 & 周永健 & \hyperref[sec:988]{第六講 神體恤我們的軟弱──耶和華是在我們中間不是} \\
第65屆 & 第7講 & 出埃及記25:1-25:40 & 周永健 & \hyperref[sec:989]{第七講 教導我們──我從天上和你們說話了} \\
第65屆 & 第8講 & 出埃及記25:1-31:18 & 周永健 & \hyperref[sec:990]{第八講 神與我們同在──使我可以住在他們中間} \\
第65屆 & 第1講 & 啟示錄1:9-1:20 & 張慕皚 & \hyperref[sec:991]{第一講 教會 — 末世的金燈臺} \\
第65屆 & 第2講 & 啟示錄2:1-2:7 & 張慕皚 & \hyperref[sec:992]{第二講 以弗所教會 — 愛的光輝} \\
第65屆 & 第3講 & 啟示錄2:8-2:11 & 張慕皚 & \hyperref[sec:993]{第三講 士每拿教會 — 苦難的光輝} \\
第65屆 & 第4講 & 啟示錄2:12-2:17 & 張慕皚 & \hyperref[sec:994]{第四講 別迦摩教會 — 聖潔生活的光輝} \\
第65屆 & 第5講 & 啟示錄2:18-2:29 & 張慕皚 & \hyperref[sec:995]{第五講 推雅推喇教會 — 職業的光輝} \\
第65屆 & 第6講 & 啟示錄3:1-3:6 & 張慕皚 & \hyperref[sec:996]{第六講 撒狄教會 — 生命更新的光輝} \\
第65屆 & 第7講 & 啟示錄3:7-3:13 & 張慕皚 & \hyperref[sec:997]{第七講 非拉鐵非教會 — 福音的光輝 (一)} \\
第65屆 & 第8講 & 啟示錄3:7-3:13 & 張慕皚 & \hyperref[sec:998]{第八講 非拉鐵非教會 — 福音的光輝 (二)} \\
第65屆 & 第9講 & 啟示錄3:14-3:22 & 張慕皚 & \hyperref[sec:999]{第九講 老底嘉教會 — 長久興旺的光輝 (一)} \\
第65屆 & 第10講 & 啟示錄3:14-3:22 & 張慕皚 & \hyperref[sec:1000]{第十講 老底嘉教會 — 長久興旺的光輝 (二)} \\
第65屆 & 第1講 &  & 楊濬哲 & \hyperref[sec:1390]{序} \\
第65屆 & 第1講 & 路加福音15:11-15:32 & 褚永華 & \hyperref[sec:1001]{第一講 浪子的比喻} \\
第65屆 & 第2講 & 馬可福音4:1-4:20 & 褚永華 & \hyperref[sec:1002]{第二講 撒種的比喻} \\
第65屆 & 第3講 & 路加福音16:19-16:31 & 褚永華 & \hyperref[sec:1003]{第三講 財主與拉撒路的比喻} \\
第65屆 & 第4講 & 路加福音10:25-10:37 & 褚永華 & \hyperref[sec:1004]{第四講 好撒瑪利亞人的比喻} \\
第65屆 & 第5講 & 馬太福音25:1-25:13 & 褚永華 & \hyperref[sec:1005]{第五講 十個童女的比喻} \\
第65屆 & 第6講 & 馬太福音25:14-25:30 & 褚永華 & \hyperref[sec:1006]{第六講 交銀僕人的比喻} \\
第65屆 & 第7講 & 馬可福音11:12-11:21 & 褚永華 & \hyperref[sec:1007]{第七講 無花果樹的比喻} \\
第65屆 & 第8講 & 馬太福音20:1-20:16 & 褚永華 & \hyperref[sec:1008]{第八講 葡萄園的比喻} \\
\hline
\hline
第66屆 & 第1講 &  & 沈保羅 & \hyperref[sec:973]{第一講 祂到底是誰} \\
第66屆 & 第2講 &  & 沈保羅 & \hyperref[sec:974]{第二講 雙重的考慮} \\
第66屆 & 第3講 &  & 沈保羅 & \hyperref[sec:975]{第三講 死亡,苦難,害怕} \\
第66屆 & 第4講 &  & 沈保羅 & \hyperref[sec:976]{第四講 那時 : 誤會,挫折,拒絕} \\
第66屆 & 第5講 &  & 沈保羅 & \hyperref[sec:977]{第五講 這是我的愛子} \\
第66屆 & 第6講 &  & 沈保羅 & \hyperref[sec:978]{第六講 祂的人生} \\
第66屆 & 第7講 &  & 沈保羅 & \hyperref[sec:979]{第七講 在那裡吃} \\
第66屆 & 第8講 &  & 沈保羅 & \hyperref[sec:980]{第八講 戰勝驚恐} \\
第66屆 & 第9講 &  & 沈保羅 & \hyperref[sec:981]{第九講 地平線的那一邊} \\
第66屆 & 第10講 &  & 沈保羅 & \hyperref[sec:982]{第十講 愛的比較} \\
第66屆 & 第1講 & 啟示錄4:1-5:14 & 蘇穎智 & \hyperref[sec:954]{第一講 在地如在天的敬拜} \\
第66屆 & 第2講 & 啟示錄6:1-7:17 & 蘇穎智 & \hyperref[sec:955]{第二講 浩劫當前 — 六印之災} \\
第66屆 & 第3講 & 啟示錄8:1-9:21 & 蘇穎智 & \hyperref[sec:956]{第三講 禍不單行 — 六號之災} \\
第66屆 & 第4講 & 啟示錄10:1-11:19 & 蘇穎智 & \hyperref[sec:957]{第四講 能帶給人出路的人} \\
第66屆 & 第5講 & 啟示錄12:1-13:18 & 蘇穎智 & \hyperref[sec:959]{第五講 神選民的仇敵} \\
第66屆 & 第6講 & 啟示錄14:1-16:21 & 蘇穎智 & \hyperref[sec:960]{第六講 窮途末路的地球 — 七碗之災} \\
第66屆 & 第7講 & 啟示錄17:1-19:21 & 蘇穎智 & \hyperref[sec:961]{第七講 必有後「福! 覆!」的大巴比倫} \\
第66屆 & 第8講 & 啟示錄19:1-20:15 & 蘇穎智 & \hyperref[sec:962]{第八講 大結局 — 有人快樂有人愁} \\
第66屆 & 第9講 & 啟示錄21:1-22:21 & 蘇穎智 & \hyperref[sec:963]{第九講 更美的家} \\
第66屆 & 第1講 & 彼得後書1:3-1:3 & 韋理信 & \hyperref[sec:964]{第一講 在動盪轉變的社會中需要一個至善至美的標準} \\
第66屆 & 第2講 & 彼得後書1:4-1:4 & 韋理信 & \hyperref[sec:965]{第二講 至善至美的標準之基礎 — 神寶貴極大的應許} \\
第66屆 & 第3講 & 提摩太後書1:5-1:5 & 韋理信 & \hyperref[sec:966]{第三講 在信心德行上至善至美的標準} \\
第66屆 & 第4講 & 彼得後書1:5-1:5 & 韋理信 & \hyperref[sec:967]{第四講 在知識上至善至美的標準} \\
第66屆 & 第5講 & 彼得後書1:6-1:6 & 韋理信 & \hyperref[sec:968]{第五講 在節制上至善至美的標準} \\
第66屆 & 第6講 & 彼得後書1:6-1:6 & 韋理信 & \hyperref[sec:969]{第六講 在忍耐上至善至美的標準} \\
第66屆 & 第7講 & 彼得後書1:6-1:6 & 韋理信 & \hyperref[sec:970]{第七講 在虔敬上至善至美的標準} \\
第66屆 & 第8講 & 彼得後書1:7-1:7 & 韋理信 & \hyperref[sec:971]{第八講 在愛弟兄和愛眾人上至善至美的標準} \\
第66屆 & 第9講 & 彼得後書1:8-1:11 & 韋理信 & \hyperref[sec:972]{第九講 如何達到至善至美的標準} \\
\hline
\hline
第67屆 & 第1講 &  & 滕近輝 & \hyperref[sec:945]{第一講 國度真理的重要性} \\
第67屆 & 第2講 &  & 滕近輝 & \hyperref[sec:946]{第二講 舊約的國度} \\
第67屆 & 第3講 &  & 滕近輝 & \hyperref[sec:947]{第三講 國度的奧秘} \\
第67屆 & 第4講 &  & 滕近輝 & \hyperref[sec:948]{第四講 國度的原則} \\
第67屆 & 第5講 &  & 滕近輝 & \hyperref[sec:949]{第五講 國度的擴展} \\
第67屆 & 第6講 &  & 滕近輝 & \hyperref[sec:950]{第六講 國度的勝利} \\
第67屆 & 第7講 &  & 滕近輝 & \hyperref[sec:951]{第七講 國度的價值觀} \\
第67屆 & 第8講 &  & 滕近輝 & \hyperref[sec:952]{第八講 國度的本體與外延} \\
第67屆 & 第9講 &  & 滕近輝 & \hyperref[sec:953]{第九講 國度與得救的異同} \\
第67屆 & 第1講 &  & 焦源濂 & \hyperref[sec:936]{第一講 受苦之前的約伯 — 幸福的真諦} \\
第67屆 & 第2講 &  & 焦源濂 & \hyperref[sec:937]{第二講 大禍臨頭的約伯 — 靈界的奧秘} \\
第67屆 & 第3講 &  & 焦源濂 & \hyperref[sec:938]{第三講 面對災難的約伯 — 真人的真像 (一)} \\
第67屆 & 第4講 &  & 焦源濂 & \hyperref[sec:939]{第四講 面對災難的約伯 — 真人的真像 (二)} \\
第67屆 & 第5講 &  & 焦源濂 & \hyperref[sec:940]{第五講 面對災難的約伯 — 真人的真像 (三)} \\
第67屆 & 第6講 &  & 焦源濂 & \hyperref[sec:941]{第六講 面對災難的約伯 — 真人的真像 (四)} \\
第67屆 & 第7講 &  & 焦源濂 & \hyperref[sec:942]{第七講 面對災難的約伯 — 家庭的危機} \\
第67屆 & 第8講 &  & 焦源濂 & \hyperref[sec:943]{第八講 知友陪伴的約伯 — 肢體生活的榜樣} \\
第67屆 & 第9講 &  & 焦源濂 & \hyperref[sec:944]{第九講 知友陪伴的約伯 — 友誼的珍貴} \\
第67屆 & 第1講 & 羅馬書1:1-1:17 & 鮑會園 & \hyperref[sec:927]{第一講 福音的託付} \\
第67屆 & 第2講 & 羅馬書1:18-2:15 & 鮑會園 & \hyperref[sec:928]{第二講 世人的需要} \\
第67屆 & 第3講 & 羅馬書3:21-3:31 & 鮑會園 & \hyperref[sec:929]{第三講 救恩的完成} \\
第67屆 & 第4講 & 羅馬書5:1-5:11 & 鮑會園 & \hyperref[sec:930]{第四講 救恩的功效} \\
第67屆 & 第5講 & 羅馬書6:1-6:23 & 鮑會園 & \hyperref[sec:931]{第五講 救恩的結果} \\
第67屆 & 第6講 & 羅馬書7:1-7:25 & 鮑會園 & \hyperref[sec:932]{第六講 實際的經驗} \\
第67屆 & 第7講 & 羅馬書8:1-8:25 & 鮑會園 & \hyperref[sec:933]{第七講 勝利的秘訣} \\
第67屆 & 第8講 & 羅馬書12:1-12:21 & 鮑會園 & \hyperref[sec:934]{第八講 生活的動力} \\
第67屆 & 第9講 & 羅馬書15:20-15:33 & 鮑會園 & \hyperref[sec:935]{第九講 教會的使命} \\
\hline
\hline
第68屆 & 第1講 & 約翰福音17:3-17:3 & 吳勇 & \hyperref[sec:911]{第一講 不把神當神} \\
第68屆 & 第2講 & 以西結書10:18-10:3 & 吳勇 & \hyperref[sec:912]{第二講 離開與回來} \\
第68屆 & 第3講 & 出埃及記40:34-40:34 & 吳勇 & \hyperref[sec:913]{第三講 三個願望} \\
第68屆 & 第4講 & 創世記25:22-25:22 & 吳勇 & \hyperref[sec:914]{第四講 屬肉體與屬靈} \\
第68屆 & 第5講 & 民數記17:8-17:8 & 吳勇 & \hyperref[sec:915]{第五講 生命的守護者} \\
第68屆 & 第6講 & 以西結書47:5-47:5 & 吳勇 & \hyperref[sec:916]{第六講 火與水的異象} \\
第68屆 & 第7講 & 腓立比書2:6-2:8 & 吳勇 & \hyperref[sec:917]{第七講 兩個安息} \\
第68屆 & 第8講 & 馬可福音3:21-3:15 & 吳勇 & \hyperref[sec:918]{第八講 約翰的補網} \\
第68屆 & 第9講 & 撒母耳記上2:29-2:22 & 吳勇 & \hyperref[sec:919]{第九講 復興的器皿} \\
第68屆 & 第10講 & 使徒行傳13:22-13:22 & 吳勇 & \hyperref[sec:920]{第十講 人與路} \\
第68屆 & 第1講 & 約翰福音4:43-4:45 & 歐達 & \hyperref[sec:902]{第一講 先知的險境} \\
第68屆 & 第2講 & 詩篇24:1-24:3 & 歐達 & \hyperref[sec:903]{第二講 與主默契} \\
第68屆 & 第3講 & 詩篇15:1-15:5 & 歐達 & \hyperref[sec:904]{第三講 正直忠誠} \\
第68屆 & 第4講 & 約翰福音4:19-4:24 & 歐達 & \hyperref[sec:905]{第四講 融和合一} \\
第68屆 & 第5講 & 哥林多前書12:1-12:7 & 歐達 & \hyperref[sec:906]{第五講 竭誠事主} \\
第68屆 & 第6講 & 詩篇119:9-119:9 & 歐達 & \hyperref[sec:907]{第六講 潔淨純全} \\
第68屆 & 第7講 & 帖撒羅尼迦前書1:6-1:10 & 歐達 & \hyperref[sec:908]{第七講 門徒樣式} \\
第68屆 & 第8講 & 哈巴谷書1:5-1:5 & 歐達 & \hyperref[sec:909]{第八講 福音使命} \\
第68屆 & 第9講 &  & 歐達 & \hyperref[sec:910]{第九講 牧者的禱告} \\
第68屆 & 第1講 &  & 蕭壽華 & \hyperref[sec:893]{第一講 跟從耶穌到底是甚麼意思?} \\
第68屆 & 第2講 &  & 蕭壽華 & \hyperref[sec:894]{第二講 天國子民的品格與見證} \\
第68屆 & 第3講 &  & 蕭壽華 & \hyperref[sec:895]{第三講 超越文士的義(一)} \\
第68屆 & 第4講 &  & 蕭壽華 & \hyperref[sec:896]{第四講 超越文士的義 (二)} \\
第68屆 & 第5講 &  & 蕭壽華 & \hyperref[sec:897]{第五講 天國子民隱藏的操練} \\
第68屆 & 第6講 &  & 蕭壽華 & \hyperref[sec:898]{第六講 天國子民的人生目標} \\
第68屆 & 第7講 &  & 蕭壽華 & \hyperref[sec:899]{第七講 天國子民的「眼光」} \\
第68屆 & 第8講 &  & 蕭壽華 & \hyperref[sec:900]{第八講 天國子民的關係} \\
第68屆 & 第9講 &  & 蕭壽華 & \hyperref[sec:901]{第九講 天國子民的委身} \\
\hline
\hline
第69屆 & 第1講 &  & 周永健 & \hyperref[sec:874]{第一講 確認主權 經歷信實} \\
第69屆 & 第2講 &  & 周永健 & \hyperref[sec:875]{第二講 持定身份 聖潔自守} \\
第69屆 & 第3講 &  & 周永健 & \hyperref[sec:876]{第三講 恆常禱告 認罪祈求} \\
第69屆 & 第4講 &  & 周永健 & \hyperref[sec:877]{第四講 不拜別神 專心靠主} \\
第69屆 & 第5講 &  & 周永健 & \hyperref[sec:878]{第五講 謙卑自省 莫虧主恩} \\
第69屆 & 第6講 &  & 周永健 & \hyperref[sec:879]{第六講 回歸國土 同心建造} \\
第69屆 & 第7講 &  & 周永健 & \hyperref[sec:880]{第七講 合力互助 共度危難} \\
第69屆 & 第8講 &  & 周永健 & \hyperref[sec:881]{第八講 聆聽主道 樂意遵行} \\
第69屆 & 第1講 &  & 陳黔開 & \hyperref[sec:882]{第一講 教會 — 你既愛之,又…} \\
第69屆 & 第2講 &  & 陳黔開 & \hyperref[sec:883]{第二講 敬拜 — 參與或投入} \\
第69屆 & 第3講 &  & 陳黔開 & \hyperref[sec:884]{第三講 佈道 — 使萬民作門徒} \\
第69屆 & 第4講 &  & 陳黔開 & \hyperref[sec:885]{第四講 教導 — 學無止境} \\
第69屆 & 第5講 &  & 陳黔開 & \hyperref[sec:886]{第五講 禱告—誰曾與神款款深談呢?} \\
第69屆 & 第6講 &  & 陳黔開 & \hyperref[sec:887]{第六講 恩賜 — 你尚待發掘的是甚麼?} \\
第69屆 & 第7講 &  & 陳黔開 & \hyperref[sec:888]{第七講 肢體 — 你屬誰?} \\
第69屆 & 第8講 &  & 陳黔開 & \hyperref[sec:889]{第八講 受死 — 有益?死地?} \\
第69屆 & 第1講 & 羅馬書1:1-1:32 & 黃子嘉 & \hyperref[sec:890]{第一講 福音是神彰顯榮耀與奧秘} \\
第69屆 & 第2講 & 羅馬書1:1-1:32 & 黃子嘉 & \hyperref[sec:891]{第二講 福音是神救人大能與大恩} \\
第69屆 & 第3講 & 羅馬書1:1-1:32 & 黃子嘉 & \hyperref[sec:892]{第三講 福音是神賜人因信而稱義也賜人平安喜樂} \\
第69屆 & 第4講 & 羅馬書6:1-6:23 & 黃子嘉 & \hyperref[sec:921]{第四講 福音是神使人稱義與成聖} \\
第69屆 & 第5講 & 羅馬書8:1-8:39 & 黃子嘉 & \hyperref[sec:922]{第五講 福音是神賜人榮美與榮耀} \\
第69屆 & 第6講 & 羅馬書12:1-12:21 & 黃子嘉 & \hyperref[sec:923]{第六講 福音是神要人奉獻作活祭} \\
第69屆 & 第7講 & 羅馬書12:1-12:21 & 黃子嘉 & \hyperref[sec:924]{第七講 福音是神要建立榮耀教會} \\
第69屆 & 第8講 & 羅馬書12:1-12:21 & 黃子嘉 & \hyperref[sec:925]{第八講 福音是神要建立美善社會} \\
第69屆 & 第9講 & 羅馬書1:1-1:32 & 黃子嘉 & \hyperref[sec:926]{第九講 福音是神要託付榮耀使命} \\
\hline
\hline
第70屆 & 第1講 &  & 曾立華 & \hyperref[sec:726]{第一講 為追求誠信正直感的禱告} \\
第70屆 & 第2講 & 腓立比書1:12-1:20 & 曾立華 & \hyperref[sec:727]{第二講 在困境中勇於表達誠信正直的品格} \\
第70屆 & 第3講 & 腓立比書1:21-1:30 & 曾立華 & \hyperref[sec:728]{第三講 恰當的正直生活形態} \\
第70屆 & 第4講 & 腓立比書2:1-2:11 & 曾立華 & \hyperref[sec:729]{第四講 同心培育基督謙卑正直的氣質} \\
第70屆 & 第5講 & 腓立比書2:12-2:30 & 曾立華 & \hyperref[sec:730]{第五講 誠信正直的領袖} \\
第70屆 & 第6講 & 腓立比書3:1-3:16 & 曾立華 & \hyperref[sec:733]{第六講 獲得誠信正直質素的秘訣} \\
第70屆 & 第7講 & 腓立比書3:17-4:1 & 曾立華 & \hyperref[sec:735]{第七講 破壞誠信正直的因素} \\
第70屆 & 第8講 & 腓立比書4:2-4:23 & 曾立華 & \hyperref[sec:737]{第八講 整全誠信正直的生活} \\
第70屆 & 第1講 &  & 林三綱 & \hyperref[sec:856]{第一講 所羅門 (耶底底亞) 生平} \\
第70屆 & 第2講 &  & 林三綱 & \hyperref[sec:857]{第二講 靈命經歷(1) 兩個心願(上)} \\
第70屆 & 第3講 &  & 林三綱 & \hyperref[sec:858]{第三講 靈命經歷(1) 兩個心願(下)} \\
第70屆 & 第4講 &  & 林三綱 & \hyperref[sec:859]{第四講 靈命經歷 (2) 躥山越嶺} \\
第70屆 & 第5講 &  & 林三綱 & \hyperref[sec:860]{第五講 靈命經歷 (3) 夜間的華轎} \\
第70屆 & 第6講 &  & 林三綱 & \hyperref[sec:861]{第六講 靈命經歷 (4) 患難與榮美} \\
第70屆 & 第7講 &  & 林三綱 & \hyperref[sec:862]{第七講 治國建殿} \\
第70屆 & 第8講 &  & 林三綱 & \hyperref[sec:863]{第八講 獻殿} \\
第70屆 & 第1講 & 創世記37:1-37:36 & 焦源濂 & \hyperref[sec:864]{第一講 一個與弟兄迥別的青年} \\
第70屆 & 第2講 & 創世記37:1-37:36 & 焦源濂 & \hyperref[sec:865]{第二講 一個「作夢」的聖徒} \\
第70屆 & 第3講 & 創世記37:1-37:36 & 焦源濂 & \hyperref[sec:866]{第三講 人間辛酸備嘗} \\
第70屆 & 第4講 & 創世記39:1-40:23 & 焦源濂 & \hyperref[sec:867]{第四講 受大的打擊,更深的扎根} \\
第70屆 & 第5講 & 創世記41:1-41:57 & 焦源濂 & \hyperref[sec:868]{第五講 突然青雲直上} \\
第70屆 & 第6講 & 創世記41:1-42:38 & 焦源濂 & \hyperref[sec:869]{第六講 舊夢重提} \\
第70屆 & 第7講 & 創世記42:1-43:34 & 焦源濂 & \hyperref[sec:870]{第七講 卻沒有一個人能看透他} \\
第70屆 & 第8講 & 創世記44:1-45:28 & 焦源濂 & \hyperref[sec:871]{第八講 一群真弟兄} \\
第70屆 & 第9講 & 創世記47:1-47:31 & 焦源濂 & \hyperref[sec:872]{第九講 甚麼是真的成就?} \\
第70屆 & 第10講 & 創世記50:1-50:26 & 焦源濂 & \hyperref[sec:873]{第十講 神的意思原是好的} \\
\hline
\hline
第71屆 & 第1講 &  & 劉承業 & \hyperref[sec:741]{第一講 詩人為何要讚美?} \\
第71屆 & 第2講 &  & 劉承業 & \hyperref[sec:742]{第二講 「讚美詩」的讚美} \\
第71屆 & 第3講 &  & 劉承業 & \hyperref[sec:743]{第三講 「感恩詩」的讚美} \\
第71屆 & 第4講 &  & 劉承業 & \hyperref[sec:744]{第四講 「個人哀歌」的讚美} \\
第71屆 & 第5講 &  & 劉承業 & \hyperref[sec:745]{第五講 「團體哀歌」的讚美} \\
第71屆 & 第6講 &  & 劉承業 & \hyperref[sec:746]{第六講 苦難中的讚美} \\
第71屆 & 第7講 &  & 劉承業 & \hyperref[sec:747]{第七講 逆境中的讚美} \\
第71屆 & 第8講 &  & 劉承業 & \hyperref[sec:748]{第八講 如何實踐讚美} \\
第71屆 & 第1講 & 約翰福音4:39-4:42 & 滕近輝 & \hyperref[sec:24]{第一講 真救主:七證} \\
第71屆 & 第2講 &  & 滕近輝 & \hyperref[sec:23]{第二講 真葡萄樹} \\
第71屆 & 第3講 & 約翰福音8:31-8:36 & 滕近輝 & \hyperref[sec:22]{第三講 真自由} \\
第71屆 & 第4講 & 約翰福音1:43-1:46 & 滕近輝 & \hyperref[sec:21]{第四講 真以色列人} \\
第71屆 & 第5講 & 約翰福音1:1-1:14 & 滕近輝 & \hyperref[sec:20]{第五講 真光} \\
第71屆 & 第6講 & 約翰福音8:31-8:31 & 滕近輝 & \hyperref[sec:19]{第六講 真門徒} \\
第71屆 & 第7講 & 約翰福音12:1-12:8 & 滕近輝 & \hyperref[sec:18]{第七講 真愛主} \\
第71屆 & 第8講 & 約翰福音4:23-4:24 & 滕近輝 & \hyperref[sec:17]{第八講 真敬拜 -- 聖禮的內在化} \\
第71屆 & 第9講 & 約翰福音7:18-7:18 & 滕近輝 & \hyperref[sec:16]{第九講 真事奉} \\
第71屆 & 第10講 & 約翰福音16:13-16:13 & 滕近輝 & \hyperref[sec:9]{第十講 真屬靈} \\
第71屆 & 第1講 &  & 盧家駇 & \hyperref[sec:749]{第一講 貪愛世界} \\
第71屆 & 第2講 &  & 盧家駇 & \hyperref[sec:753]{第二講 離棄愛心} \\
第71屆 & 第3講 &  & 盧家駇 & \hyperref[sec:756]{第三講 取死身體} \\
第71屆 & 第4講 &  & 盧家駇 & \hyperref[sec:760]{第四講 忘掉委身} \\
第71屆 & 第5講 &  & 盧家駇 & \hyperref[sec:763]{第五講 隨流失去} \\
第71屆 & 第6講 &  & 盧家駇 & \hyperref[sec:764]{第六講 缺乏目標} \\
第71屆 & 第7講 &  & 盧家駇 & \hyperref[sec:765]{第七講 不傳福音} \\
第71屆 & 第8講 &  & 盧家駇 & \hyperref[sec:766]{第八講 抗拒使命} \\
\hline
\hline
第72屆 & 第1講 & 馬太福音24:1-24:51 & 張慕皚 & \hyperref[sec:643]{第一講 內虛外暴的末世文化} \\
第72屆 & 第2講 & 馬太福音24:1-24:51 & 張慕皚 & \hyperref[sec:644]{第二講 在利己主義中邁向滅亡} \\
第72屆 & 第3講 & 但以理書1:1-1:21 & 張慕皚 & \hyperref[sec:645]{第三講 根基毀壞,義人還能作甚麼?} \\
第72屆 & 第4講 & 馬太福音24:1-24:51 & 張慕皚 & \hyperref[sec:646]{第四講 敵基督與全球化的災難} \\
第72屆 & 第5講 & 以弗所書6:1-6:24 & 張慕皚 & \hyperref[sec:648]{第五講 咆哮覓食的末世惡魔} \\
第72屆 & 第6講 & 以弗所書6:1-6:24 & 張慕皚 & \hyperref[sec:650]{第六講 面對惡魔,勝算在握} \\
第72屆 & 第7講 & 以弗所書6:1-6:24 & 張慕皚 & \hyperref[sec:653]{第七講 攜械披甲,嚴陣以待} \\
第72屆 & 第8講 & 馬太福音25:1-25:46 & 張慕皚 & \hyperref[sec:655]{第八講 深夜中的守望者} \\
第72屆 & 第9講 & 以賽亞書52:1-52:15 & 張慕皚 & \hyperref[sec:658]{第九講 離開巴比倫大城的呼召} \\
第72屆 & 第10講 & 馬太福音24:1-24:51 & 張慕皚 & \hyperref[sec:660]{第十講 主耶穌何時再來?} \\
第72屆 & 第1講 &  & 蘇穎睿 & \hyperref[sec:606]{第一講 新與變} \\
第72屆 & 第2講 &  & 蘇穎睿 & \hyperref[sec:607]{第二講 見到就信?未必!} \\
第72屆 & 第3講 &  & 蘇穎睿 & \hyperref[sec:609]{第三講 跳下去!等陣先!} \\
第72屆 & 第4講 &  & 蘇穎睿 & \hyperref[sec:612]{第四講 封建枷鎖} \\
第72屆 & 第5講 &  & 蘇穎睿 & \hyperref[sec:613]{第五講 民以食為先} \\
第72屆 & 第6講 &  & 蘇穎睿 & \hyperref[sec:615]{第六講 是我,不要怕!} \\
第72屆 & 第7講 &  & 蘇穎睿 & \hyperref[sec:616]{第七講 重見光明} \\
第72屆 & 第8講 &  & 蘇穎睿 & \hyperref[sec:617]{第八講 等待難捱的日子} \\
第72屆 & 第9講 &  & 蘇穎睿 & \hyperref[sec:619]{第九講 生離死別} \\
第72屆 & 第1講 &  & 陳濟民 & \hyperref[sec:622]{第一講 耶穌的榮耀} \\
第72屆 & 第2講 &  & 陳濟民 & \hyperref[sec:625]{第二講 生命的賜與者} \\
第72屆 & 第3講 &  & 陳濟民 & \hyperref[sec:627]{第三講 從耶穌的「不孝」談起} \\
第72屆 & 第4講 &  & 陳濟民 & \hyperref[sec:629]{第四講 恨死了} \\
第72屆 & 第5講 &  & 陳濟民 & \hyperref[sec:630]{第五講 至死的愛} \\
第72屆 & 第6講 &  & 陳濟民 & \hyperref[sec:631]{第六講 知與信} \\
第72屆 & 第7講 &  & 陳濟民 & \hyperref[sec:635]{第七講 言與行} \\
第72屆 & 第8講 &  & 陳濟民 & \hyperref[sec:637]{第八講 怎有可能超越} \\
第72屆 & 第9講 &  & 陳濟民 & \hyperref[sec:640]{第九講 世人怎能信} \\
\hline
\hline
第73屆 & 第1講 &  & 于力工 & \hyperref[sec:587]{第一講 由眾到一讀經法:用全部《聖經》來讀一段話} \\
第73屆 & 第2講 &  & 于力工 & \hyperref[sec:589]{第二講 心理讀經法:用人的天性來讀主的話} \\
第73屆 & 第3講 &  & 于力工 & \hyperref[sec:592]{第三講 性格讀經法:用人的性格來讀主的話} \\
第73屆 & 第4講 &  & 于力工 & \hyperref[sec:594]{第四講 禱讀讀經法:用經文變成禱告的話語} \\
第73屆 & 第5講 &  & 于力工 & \hyperref[sec:597]{第五講 創意讀經法:用已存之事物加上新的功用} \\
第73屆 & 第6講 &  & 于力工 & \hyperref[sec:598]{第六講 由一到眾讀經法:用一處經文來看全部《聖經》的說明} \\
第73屆 & 第7講 &  & 于力工 & \hyperref[sec:599]{第七講 靈修讀經法:用靈修的基本方法來讀主的話} \\
第73屆 & 第8講 &  & 于力工 & \hyperref[sec:600]{第八講 邏輯讀經法:推理讀經,因果讀法} \\
第73屆 & 第9講 &  & 于力工 & \hyperref[sec:601]{第九講 三一讀經法:讀出神的性情來} \\
第73屆 & 第10講 &  & 于力工 & \hyperref[sec:602]{第十講 靈進讀經法:用靈命進深之系統來讀經} \\
第73屆 & 第1講 &  & 李思敬 & \hyperref[sec:556]{第一講 我父在天} \\
第73屆 & 第2講 &  & 李思敬 & \hyperref[sec:557]{第二講 願爾名聖,爾國臨格} \\
第73屆 & 第3講 &  & 李思敬 & \hyperref[sec:560]{第三講 爾旨得成,在地若天} \\
第73屆 & 第4講 &  & 李思敬 & \hyperref[sec:561]{第四講 所需之糧,今日賜我} \\
第73屆 & 第5講 &  & 李思敬 & \hyperref[sec:562]{第五講 我免人負,求免我負} \\
第73屆 & 第6講 &  & 李思敬 & \hyperref[sec:563]{第六講 俾勿我試,拯我出惡} \\
第73屆 & 第7講 &  & 李思敬 & \hyperref[sec:564]{第七講 以國,權,榮,皆爾所有} \\
第73屆 & 第8講 &  & 李思敬 & \hyperref[sec:566]{第八講 爰及世世,誠心所願} \\
第73屆 & 第9講 &  & 李思敬 & \hyperref[sec:568]{第九講 凡求者得也} \\
第73屆 & 第1講 &  & 樓恩德 & \hyperref[sec:572]{第一講 那人獨居不好} \\
第73屆 & 第2講 &  & 樓恩德 & \hyperref[sec:573]{第二講 罪就伏在門前} \\
第73屆 & 第3講 &  & 樓恩德 & \hyperref[sec:575]{第三講 漸漸挪移帳棚} \\
第73屆 & 第4講 &  & 樓恩德 & \hyperref[sec:577]{第四講 從暗笑到喜笑} \\
第73屆 & 第5講 &  & 樓恩德 & \hyperref[sec:579]{第五講 神必親自預備} \\
第73屆 & 第6講 &  & 樓恩德 & \hyperref[sec:580]{第六講 必在這地昌盛} \\
第73屆 & 第7講 &  & 樓恩德 & \hyperref[sec:581]{第七講 從不知到知道} \\
第73屆 & 第8講 &  & 樓恩德 & \hyperref[sec:582]{第八講 不然我就死了} \\
第73屆 & 第9講 &  & 樓恩德 & \hyperref[sec:584]{第九講 因信留下遺命} \\
\hline
\hline
第74屆 & 第1講 & 以弗所書4:1-4:1 & 戴紹曾 & \hyperref[sec:434]{第一講 蒙召成為見證} \\
第74屆 & 第2講 & 彼得前書3:9-3:9 & 戴紹曾 & \hyperref[sec:435]{第二講 蒙召承受福氣} \\
第74屆 & 第3講 & 路加福音5:1-5:13 & 戴紹曾 & \hyperref[sec:439]{第三講 蒙召得以自由} \\
第74屆 & 第4講 & 哥林多前書1:15-1:16 & 戴紹曾 & \hyperref[sec:440]{第四講 蒙召成為聖潔} \\
第74屆 & 第5講 & 哥林多前書1:9-1:9 & 戴紹曾 & \hyperref[sec:441]{第五講 蒙召與主聯合} \\
第74屆 & 第6講 & 歌羅西書3:15-3:15 & 戴紹曾 & \hyperref[sec:442]{第六講 蒙召得享平安} \\
第74屆 & 第7講 & 彼得前書2:9-2:10 & 戴紹曾 & \hyperref[sec:443]{第七講 蒙召宣揚基督} \\
第74屆 & 第8講 & 彼得前書2:18-2:23 & 戴紹曾 & \hyperref[sec:444]{第八講 蒙召為主受苦} \\
第74屆 & 第9講 & 彼得前書5:10-5:10 & 戴紹曾 & \hyperref[sec:445]{第九講 蒙召得享榮耀} \\
第74屆 & 第10講 & 彼得後書1:10-1:10 & 戴紹曾 & \hyperref[sec:446]{第十講 蒙召殷勤堅定} \\
第74屆 & 第1講 & 路加福音4:14-4:21 & 江耀全 & \hyperref[sec:419]{第一講 愛的福音} \\
第74屆 & 第2講 & 路加福音1:46-1:56 & 江耀全 & \hyperref[sec:420]{第二講 愛的順服} \\
第74屆 & 第3講 & 路加福音7:36-7:50 & 江耀全 & \hyperref[sec:421]{第三講 愛的恩典} \\
第74屆 & 第4講 & 路加福音10:25-10:37 & 江耀全 & \hyperref[sec:422]{第四講 愛的承擔} \\
第74屆 & 第5講 & 路加福音10:38-10:42 & 江耀全 & \hyperref[sec:423]{第五講 愛的服侍} \\
第74屆 & 第6講 & 路加福音15:11-15:32 & 江耀全 & \hyperref[sec:424]{第六講 愛的寬宏} \\
第74屆 & 第7講 & 路加福音17:11-17:19 & 江耀全 & \hyperref[sec:425]{第七講 愛的感恩} \\
第74屆 & 第8講 & 路加福音19:1-19:10 & 江耀全 & \hyperref[sec:426]{第八講 愛的相遇} \\
第74屆 & 第9講 & 路加福音24:13-24:35 & 江耀全 & \hyperref[sec:427]{第九講 愛的同在} \\
第74屆 & 第1講 & 哥林多後書4:7-4:11 & 溫偉耀 & \hyperref[sec:428]{第一講 軟弱與剛強 -- 寶貝與瓦器} \\
第74屆 & 第2講 & 哥林多後書12:1-12:10 & 溫偉耀 & \hyperref[sec:429]{第二講 從高峰到深淵 -- 「三層天」與「一根刺」} \\
第74屆 & 第3講 & 哥林多後書2:14-3:6 & 溫偉耀 & \hyperref[sec:430]{第三講 向人推薦自己 -- 怎樣給人留下深刻的印象?} \\
第74屆 & 第4講 & 哥林多後書11:21-11:33 & 溫偉耀 & \hyperref[sec:431]{第四講 誇口 -- 怎樣評價自己?} \\
第74屆 & 第5講 & 哥林多後書1:3-1:10 & 溫偉耀 & \hyperref[sec:432]{第五講 火浴鳳凰 -- 跟苦難交個朋友} \\
第74屆 & 第6講 & 哥林多後書3:4-3:11 & 溫偉耀 & \hyperref[sec:433]{第六講 榮光的職事 -- 上帝要怎樣的人去事奉祂?} \\
第74屆 & 第7講 & 哥林多後書9:6-9:15 & 溫偉耀 & \hyperref[sec:436]{第七講 「零風險」投資 -- 捐輸的屬靈策略與回報} \\
第74屆 & 第8講 & 哥林多後書5:14-5:21 & 溫偉耀 & \hyperref[sec:437]{第八講 新造的人 -- 活出轉化了的屬靈生命} \\
第74屆 & 第9講 & 哥林多後書5:1-5:10 & 溫偉耀 & \hyperref[sec:438]{第九講 今生與永恆 -- 灑脫而進取的生命情操} \\
\hline
\hline
第75屆 & 第1講 & 詩篇1:1-1:6 & 焦源濂 & \hyperref[sec:399]{第一分題:有福之人--第一講 蒙福的源頭} \\
第75屆 & 第2講 & 詩篇1:1-1:6 & 焦源濂 & \hyperref[sec:400]{第一分題:有福之人--第二講:蒙福的途徑} \\
第75屆 & 第3講 & 詩篇1:1-1:6 & 焦源濂 & \hyperref[sec:401]{第一分題:有福之人--第三講 蒙福的見證} \\
第75屆 & 第4講 & 詩篇1:1-1:6 & 焦源濂 & \hyperref[sec:403]{第一分題:有福之人--第四講 蒙福的為人} \\
第75屆 & 第5講 & 詩篇1:1-1:6 & 焦源濂 & \hyperref[sec:405]{第一分題:有福之人--第五講 蒙福的保証} \\
第75屆 & 第6講 & 詩篇23:1-23:6 & 焦源濂 & \hyperref[sec:408]{第二分題:詩的人生--第六講 人生之詩的秘訣} \\
第75屆 & 第7講 & 詩篇23:1-23:6 & 焦源濂 & \hyperref[sec:410]{第二分題:詩的人生--第七講 無憂無慮的人生} \\
第75屆 & 第8講 & 詩篇23:1-23:6 & 焦源濂 & \hyperref[sec:414]{第二分題:詩的人生--第八講 目標正確的人生} \\
第75屆 & 第9講 & 詩篇23:1-23:6 & 焦源濂 & \hyperref[sec:417]{第二分題:詩的人生--第九講 不畏艱苦的人生} \\
第75屆 & 第10講 & 詩篇23:1-23:6 & 焦源濂 & \hyperref[sec:418]{第二分題:詩的人生--第十講 經歷美好的人生} \\
第75屆 & 第1講 & 路加福音4:1-4:44 & 蔡元雲 & \hyperref[sec:383]{第一講 同進曠野:廿一世紀的試探} \\
第75屆 & 第2講 & 路加福音6:1-6:49 & 蔡元雲 & \hyperref[sec:385]{第二講 同登高山:廿一世紀的文化救贖} \\
第75屆 & 第3講 & 路加福音9:1-9:62 & 蔡元雲 & \hyperref[sec:387]{第三講 同背十架:廿一世紀的苦難} \\
第75屆 & 第4講 & 路加福音10:1-10:42 & 蔡元雲 & \hyperref[sec:390]{第四講 同闖禁區:廿一世紀的種族衝突} \\
第75屆 & 第5講 & 路加福音10:1-11:54 & 蔡元雲 & \hyperref[sec:393]{第五講 同動同靜:廿一世紀的屬靈操練} \\
第75屆 & 第6講 & 路加福音15:1-15:32 & 蔡元雲 & \hyperref[sec:394]{第六講 同尋迷羊:廿一世紀的浪子} \\
第75屆 & 第7講 & 路加福音23:1-24:53 & 蔡元雲 & \hyperref[sec:395]{第七講 同死同活:廿一世紀的能力危機} \\
第75屆 & 第8講 & 路加福音24:1-24:53 & 蔡元雲 & \hyperref[sec:397]{第八講 同被差遣:廿一世紀的普世宣教} \\
第75屆 & 第1講 & 約書亞記1:1-1:18 & 蕭壽華 & \hyperref[sec:365]{第一講 只要剛強壯膽} \\
第75屆 & 第2講 & 約書亞記2:1-2:24 & 蕭壽華 & \hyperref[sec:366]{第二講 尋求真實的信心} \\
第75屆 & 第3講 & 約書亞記3:1-3:17 & 蕭壽華 & \hyperref[sec:367]{第三講 神在江河中開路} \\
第75屆 & 第4講 & 約書亞記5:1-5:15 & 蕭壽華 & \hyperref[sec:369]{第四講 調整與神的關係} \\
第75屆 & 第5講 & 約書亞記6:1-6:21 & 蕭壽華 & \hyperref[sec:371]{第五講 信靠順服中的得勝} \\
第75屆 & 第6講 & 約書亞記7:1-7:26 & 蕭壽華 & \hyperref[sec:374]{第六講 不容罪叫教會傾覆} \\
第75屆 & 第7講 & 約書亞記8:30-8:35 & 蕭壽華 & \hyperref[sec:377]{第七講 修築祭壇奉獻自己} \\
第75屆 & 第8講 & 約書亞記14:6-14:15 & 蕭壽華 & \hyperref[sec:380]{第八講 專心信靠跟隨神 - 迦勒的見證} \\
\hline
\hline
第76屆 & 第1講 & 箴言1:1-1:7 & 李思敬 & \hyperref[sec:338]{第一講 入門與定向} \\
第76屆 & 第2講 & 箴言1:8-1:33 & 李思敬 & \hyperref[sec:339]{第二講 動靜的抉擇} \\
第76屆 & 第3講 & 箴言2:1-2:22 & 李思敬 & \hyperref[sec:340]{第三講 敬虔還須正直} \\
第76屆 & 第4講 & 箴言3:1-3:35 & 李思敬 & \hyperref[sec:341]{第四講 說「不」的勇氣} \\
第76屆 & 第5講 & 箴言4:1-5:6 & 李思敬 & \hyperref[sec:342]{第五講 小心思想防線} \\
第76屆 & 第6講 & 箴言5:7-6:11 & 李思敬 & \hyperref[sec:343]{第六講 離惡莫猶豫} \\
第76屆 & 第7講 & 箴言6:12-7:23 & 李思敬 & \hyperref[sec:344]{第七講 還姦淫真面目} \\
第76屆 & 第8講 & 箴言7:24-9:18 & 李思敬 & \hyperref[sec:345]{第八講 明辨生與死} \\
第76屆 & 第1講 & 哥林多後書5:11-5:21 & 羅祖澄 & \hyperref[sec:346]{第一講 基督徒的新生命} \\
第76屆 & 第2講 & 帖撒羅尼迦前書1:1-1:3 & 羅祖澄 & \hyperref[sec:347]{第二講 新生命的標誌} \\
第76屆 & 第3講 & 詩篇127:1-127:5 & 羅祖澄 & \hyperref[sec:348]{第三講 生命的重整} \\
第76屆 & 第4講 & 腓立比書1:1-1:11 & 羅祖澄 & \hyperref[sec:349]{第四講 建立生命中的戰友} \\
第76屆 & 第5講 & 傳道書2:1-2:11 & 羅祖澄 & \hyperref[sec:350]{第五講 超越生命的虛空} \\
第76屆 & 第6講 & 約翰福音20:19-20:23 & 羅祖澄 & \hyperref[sec:351]{第六講 勝過生命的困鎖} \\
第76屆 & 第7講 & 出埃及記17:8-17:12 & 羅祖澄 & \hyperref[sec:352]{第七講 得力的生命} \\
第76屆 & 第8講 & 以賽亞書6:1-6:10 & 羅祖澄 & \hyperref[sec:353]{第八講 甘受差遣的生命} \\
第76屆 & 第1講 & 腓立比書2:1-2:11 & 高文 & \hyperref[sec:354]{第一講 耶穌基督是主} \\
第76屆 & 第2講 & 馬太福音7:13-7:14 & 高文 & \hyperref[sec:355]{第二講 選擇耶穌的窄路} \\
第76屆 & 第3講 & 路加福音19:1-19:10 & 高文 & \hyperref[sec:356]{第三講 耶穌的邀請} \\
第76屆 & 第4講 & 約翰福音4:31-4:42 & 高文 & \hyperref[sec:357]{第四講 為耶穌收割} \\
第76屆 & 第5講 & 馬太福音28:16-28:20 & 高文 & \hyperref[sec:358]{第五講 耶穌的大使命} \\
第76屆 & 第6講 & 約翰福音21:15-21:24 & 高文 & \hyperref[sec:359]{第六講 耶穌發出的決定性問題} \\
第76屆 & 第7講 & 約翰福音14:15-14:17 & 高文 & \hyperref[sec:360]{第七講 耶穌的同在} \\
第76屆 & 第8講 & 路加福音24:44-24:49 & 高文 & \hyperref[sec:362]{第八講 耶穌所賜的復興} \\
第76屆 & 第9講 & 啟示錄12:10-12:12 & 高文 & \hyperref[sec:363]{第九講 在耶穌的權柄裡得勝} \\
第76屆 & 第10講 & 啟示錄19:1-19:9 & 高文 & \hyperref[sec:364]{第十講 與耶穌同赴婚筵} \\
\hline
\hline
第77屆 & 第1講 & 出埃及記2:11-3:6 & 吳獻章 & \hyperref[sec:319]{第一講 永恆的託付──摩西} \\
第77屆 & 第2講 & 撒母耳記上16:1-16:13 & 吳獻章 & \hyperref[sec:320]{第二講 臥虎藏龍} \\
第77屆 & 第3講 & 撒母耳記下15:13-15:8 & 吳獻章 & \hyperref[sec:321]{第三講 權柄的失落──看父親的背影} \\
第77屆 & 第4講 & 詩篇23:1-23:6 & 吳獻章 & \hyperref[sec:322]{第四講 被牧之羊} \\
第77屆 & 第5講 & 以賽亞書6:1-6:13 & 吳獻章 & \hyperref[sec:323]{第五講 危險的任務} \\
第77屆 & 第6講 & 約翰福音1:6-1:7 & 吳獻章 & \hyperref[sec:324]{第六講 影響力的本質:名嘴背後} \\
第77屆 & 第7講 & 馬太福音4:1-4:11 & 吳獻章 & \hyperref[sec:325]{第七講 危險的試探} \\
第77屆 & 第8講 & 馬可福音15:37-16:8 & 吳獻章 & \hyperref[sec:326]{第八講 基督的震撼教育} \\
第77屆 & 第9講 & 使徒行傳16:1-16:10 & 吳獻章 & \hyperref[sec:327]{第九講 異象．團隊．委身──聖靈．人與宣教} \\
第77屆 & 第1講 & 啟示錄2:1-2:7 & 盧家駇 & \hyperref[sec:328]{第一講 起初的愛心} \\
第77屆 & 第2講 & 以弗所書2:11-2:13 & 盧家駇 & \hyperref[sec:329]{第二講 我愛我主教會} \\
第77屆 & 第3講 & 羅馬書7:21-7:25 & 盧家駇 & \hyperref[sec:330]{第三講 除去幽暗} \\
第77屆 & 第4講 & 羅馬書12:1-12:11 & 盧家駇 & \hyperref[sec:331]{第四講 心裡火熱} \\
第77屆 & 第5講 & 使徒行傳4:36-4:37 & 盧家駇 & \hyperref[sec:332]{第五講 完全奉獻} \\
第77屆 & 第6講 & 約翰福音13:31-13:35 & 盧家駇 & \hyperref[sec:333]{第六講 新命令} \\
第77屆 & 第7講 & 腓立比書4:10-4:14 & 盧家駇 & \hyperref[sec:334]{第七講 笑傲逆境} \\
第77屆 & 第8講 & 羅馬書9:1-9:3 & 盧家駇 & \hyperref[sec:335]{第八講 骨肉之親} \\
第77屆 & 第9講 & 馬太福音28:19-28:20 & 盧家駇 & \hyperref[sec:336]{第九講 遍傳地極} \\
第77屆 & 第10講 & 帖撒羅尼迦前書4:13-4:18 & 盧家駇 & \hyperref[sec:337]{第十講 候主再臨} \\
第77屆 & 第1講 & 馬太福音1:18-1:25 & 鮑維均 & \hyperref[sec:310]{第一講 神與我們同在} \\
第77屆 & 第2講 & 馬太福音2:1-2:12 & 鮑維均 & \hyperref[sec:311]{第二講 猶太人的王} \\
第77屆 & 第3講 & 馬太福音4:1-4:11 & 鮑維均 & \hyperref[sec:312]{第三講 曠野之旅} \\
第77屆 & 第4講 & 馬太福音6:5-6:13 & 鮑維均 & \hyperref[sec:313]{第四講 願你的國降臨} \\
第77屆 & 第5講 & 馬太福音9:18-9:26 & 鮑維均 & \hyperref[sec:314]{第五講 來跟從我} \\
第77屆 & 第6講 & 馬太福音13:1-13:23 & 鮑維均 & \hyperref[sec:315]{第六講 看卻看不見} \\
第77屆 & 第7講 & 馬太福音18:21-18:35 & 鮑維均 & \hyperref[sec:316]{第七講 七十個七次} \\
第77屆 & 第8講 & 馬太福音19:1-19:12 & 鮑維均 & \hyperref[sec:317]{第八講 二人成為一體} \\
第77屆 & 第9講 & 馬太福音27:1-27:10 & 鮑維均 & \hyperref[sec:318]{第九講 無辜之人的血} \\
\hline
\hline
第78屆 & 第1講 & 以弗所書4:1-4:3 & 唐崇明 & \hyperref[sec:707]{第一講 甘為愛奴獻一生} \\
第78屆 & 第2講 & 腓立比書3:9-3:12 & 唐崇明 & \hyperref[sec:708]{第二講 畢生所冀只為神} \\
第78屆 & 第3講 & 路加福音2:41-2:52 & 唐崇明 & \hyperref[sec:709]{第三講 父事為念定終生} \\
第78屆 & 第4講 & 哥林多前書13:1-13:13 & 唐崇明 & \hyperref[sec:710]{第四講 愛得真誠愛得純} \\
第78屆 & 第5講 & 耶利米書15:19-15:21 & 唐崇明 & \hyperref[sec:711]{第五講 心灰意冷點心燈} \\
第78屆 & 第6講 & 耶利米哀歌3:19-3:27 & 唐崇明 & \hyperref[sec:712]{第六講 信實廣大全是恩} \\
第78屆 & 第7講 & 以賽亞書40:27-40:31 & 唐崇明 & \hyperref[sec:713]{第七講 展翅上騰破雲層} \\
第78屆 & 第8講 & 出埃及記17:8-17:16 & 唐崇明 & \hyperref[sec:719]{第八講 旗開得勝奔前程} \\
第78屆 & 第9講 & 提摩太後書2:15-2:16 & 唐崇明 & \hyperref[sec:722]{第九講 無愧工人逾千軍} \\
第78屆 & 第10講 & 約翰福音21:15-21:19 & 唐崇明 & \hyperref[sec:725]{第十講 重新立志顧羊群} \\
第78屆 & 第1講 & 創世記1:1-1:31 & 楊錫鏘 & \hyperref[sec:664]{第一講 創造的主權} \\
第78屆 & 第2講 & 創世記1:1-1:10 & 楊錫鏘 & \hyperref[sec:665]{第二講 創造的設計(一)} \\
第78屆 & 第3講 & 創世記1:11-1:31 & 楊錫鏘 & \hyperref[sec:666]{第三講 創造的設計(二)} \\
第78屆 & 第4講 & 創世記2:1-2:3 & 楊錫鏘 & \hyperref[sec:667]{第四講 創造的意義} \\
第78屆 & 第5講 & 創世記2:4-2:17 & 楊錫鏘 & \hyperref[sec:668]{第五講 創造的禮物(一)} \\
第78屆 & 第6講 & 創世記2:18-2:25 & 楊錫鏘 & \hyperref[sec:669]{第六講 創造的禮物(二)} \\
第78屆 & 第7講 & 創世記3:1-3:7 & 楊錫鏘 & \hyperref[sec:670]{第七講 受造物的悖逆} \\
第78屆 & 第8講 & 創世記3:7-3:19 & 楊錫鏘 & \hyperref[sec:671]{第八講 受造物的審判} \\
第78屆 & 第9講 & 創世記3:8-3:24 & 楊錫鏘 & \hyperref[sec:672]{第九講 受造物的蒙恩} \\
第78屆 & 第1講 & 約翰一書1:1-1:4 & 郭文池 & \hyperref[sec:673]{第一講 重尋相交生活的豐盛} \\
第78屆 & 第2講 & 約翰一書1:5-2:2 & 郭文池 & \hyperref[sec:681]{第二講 重尋聖潔生活的豐盛} \\
第78屆 & 第3講 & 約翰一書2:3-2:17 & 郭文池 & \hyperref[sec:684]{第三講 重尋光明生活的豐盛} \\
第78屆 & 第4講 & 約翰一書2:18-2:29 & 郭文池 & \hyperref[sec:687]{第四講 重尋信心生活的豐盛} \\
第78屆 & 第5講 & 約翰一書3:1-3:10 & 郭文池 & \hyperref[sec:689]{第五講 重尋無罪生活的豐盛} \\
第78屆 & 第6講 & 約翰一書3:11-4:6 & 郭文池 & \hyperref[sec:694]{第六講 重尋靈裡生活的豐盛} \\
第78屆 & 第7講 & 約翰一書4:7-5:4 & 郭文池 & \hyperref[sec:695]{第七講 重尋相愛生活的豐盛} \\
第78屆 & 第8講 & 約翰一書5:5-5:12 & 郭文池 & \hyperref[sec:699]{第八講 重尋救恩生活的豐盛} \\
第78屆 & 第9講 & 約翰一書5:13-5:21 & 郭文池 & \hyperref[sec:702]{第九講 重尋平安生活的豐盛} \\
\hline
\hline
第79屆 & 第1講 & 馬可福音4:1-4:20 & 奧思邦 & \hyperref[sec:291]{第一講 比喻目的:回應呼召} \\
第79屆 & 第2講 & 路加福音15:11-15:32 & 奧思邦 & \hyperref[sec:292]{第二講 神恩浩瀚,不可思議} \\
第79屆 & 第3講 & 路加福音12:13-12:21 & 奧思邦 & \hyperref[sec:293]{第三講 無知的人,罪有應得} \\
第79屆 & 第4講 & 馬太福音22:1-22:14 & 奧思邦 & \hyperref[sec:294]{第四講 神發邀請,切勿推延} \\
第79屆 & 第5講 & 約翰福音10:1-10:18 & 奧思邦 & \hyperref[sec:295]{第五講 善牧耶穌,唯一門路} \\
第79屆 & 第6講 & 路加福音16:1-16:13 & 奧思邦 & \hyperref[sec:296]{第六講 盡用資源,全為基督} \\
第79屆 & 第7講 & 馬太福音18:21-18:35 & 奧思邦 & \hyperref[sec:297]{第七講 豈可或缺:彼此饒恕} \\
第79屆 & 第8講 & 馬太福音25:14-25:30 & 奧思邦 & \hyperref[sec:298]{第八講 信徒責任──才幹歸神} \\
第79屆 & 第9講 & 路加福音11:1-11:13 & 奧思邦 & \hyperref[sec:299]{第九講 信徒生活──禱告核心} \\
第79屆 & 第1講 & 瑪拉基書1:1-1:5 & 張慕皚 & \hyperref[sec:300]{第一講 回歸真愛的豐盛中} \\
第79屆 & 第2講 & 瑪拉基書1:6-1:14 & 張慕皚 & \hyperref[sec:301]{第二講 回歸真敬拜的豐盛中} \\
第79屆 & 第3講 & 瑪拉基書2:1-2:9 & 張慕皚 & \hyperref[sec:302]{第三講 回歸行道的豐盛中} \\
第79屆 & 第4講 & 瑪拉基書2:10-2:16 & 張慕皚 & \hyperref[sec:303]{第四講 回歸貞潔生活的豐盛中(一)} \\
第79屆 & 第5講 & 瑪拉基書2:10-2:16 & 張慕皚 & \hyperref[sec:304]{第五講 回歸貞潔生活的豐盛中(二)} \\
第79屆 & 第6講 & 瑪拉基書2:17-3:6 & 張慕皚 & \hyperref[sec:305]{第六講 回歸真信心的豐盛中} \\
第79屆 & 第7講 & 瑪拉基書3:7-3:12 & 張慕皚 & \hyperref[sec:306]{第七講 回歸物質生活的豐盛中} \\
第79屆 & 第8講 & 瑪拉基書3:13-3:15 & 張慕皚 & \hyperref[sec:307]{第八講 回歸真屬靈的生命中(一)} \\
第79屆 & 第9講 & 瑪拉基書3:16-3:18 & 張慕皚 & \hyperref[sec:308]{第九講 回歸真屬靈的豐盛中(二)} \\
第79屆 & 第10講 & 瑪拉基書4:1-4:6 & 張慕皚 & \hyperref[sec:309]{第十講 回歸真盼望的豐盛中} \\
第79屆 & 第1講 &  & 李寶珠 & \hyperref[sec:282]{第一講 主題介紹──前因後果} \\
第79屆 & 第2講 &  & 李寶珠 & \hyperref[sec:283]{第二講 信的根據} \\
第79屆 & 第3講 &  & 李寶珠 & \hyperref[sec:284]{第三講 信的証據} \\
第79屆 & 第4講 &  & 李寶珠 & \hyperref[sec:285]{第四講 信的理據} \\
第79屆 & 第5講 &  & 李寶珠 & \hyperref[sec:286]{第五講 愛的模式} \\
第79屆 & 第6講 &  & 李寶珠 & \hyperref[sec:287]{第六講 愛的方式} \\
第79屆 & 第7講 &  & 李寶珠 & \hyperref[sec:288]{第七講 望的焦點} \\
第79屆 & 第8講 &  & 李寶珠 & \hyperref[sec:289]{第八講 望的眼光} \\
第79屆 & 第9講 &  & 李寶珠 & \hyperref[sec:290]{第九講 望的視力} \\
\end{xltabular}

}


\newpage

\section{第1屆~港九培靈會~古約翰~第一講}
\label{sec:1322}
\textbf{復興}
\newline
\newline
連結: \href{https://www.hkbibleconference.org/session-message/view/1322}{\texttt{https://www.hkbibleconference.org/session-message/view/1322}}
\newline
\newline
\hyperref[sec:code]{< < < PREV SERMON < < <}
~
\hyperref[sec:toc]{[返目錄]}
~
\hyperref[sec:1326]{> > > NEXT SERMON > > >}
\newline
\newline
我們讀了使徒行傳第一章一至十一節那段經文,便發生一條復興的問題,我今要對各位講的,也是這復興的問題,但我所講的,總不出乎聖經之外.我們知真的復興,就是聖靈得人而在人的心中握權.教會歷史中最大的復興,就是耶路撒冷初期的復興.我們要得復興的模範,也要回到耶路撒冷初期的教會.當耶路撒冷教會未復興之前,便有一百二十人充滿聖靈,後由這些人為主作證,而引導多人歸到主的羊欄,使教會由此大大的復興起來.今我們要知,主的意旨,不獨望耶路撒冷的人得復興,亦望廣州的人得同樣的復興.我們若要榮耀主,那就須得勝靈的充滿,好像昔日在耶路撒冷得勝靈充滿的人一樣.若我們不得勝靈充滿,則我們不獨不能代表主,且要受魔鬼的擾害.想主耶穌初到曠野時,因先充滿聖靈,故能拒試探而勝魔鬼；若我們不充滿聖靈,則必為魔鬼所勝.然我們若不真誠悔改,則必不能得真的復興；若不得勝靈充滿而得真的復興,則我們還不及世俗的人,這因我們既不能救人,且反害人.中國教會不進步,就是太多這樣不得勝靈充滿的人!從前在北方有一教會請我往他們的教會主理一禮拜復興會集,我答沒有這麼多空閒的時候,只應允一日,因在我的工作範圍中有許多人我要向他們作工的.隨後他們再來請求,說他們的教會是失敗的教會,若不復興則完全沒有望了.同時另有一人對我說:『這教會充滿了偽善的人,不如不去,任他跌倒就是了.』但我覺那跌倒了的教會,正急需復興.於是往作工數日,乃見中西的信徒到台前認罪,原來他們的教會被魔鬼掌權於他們當中,致令中西領袖時起紛爭.但雖如此,然聖靈在會中顯出,使人大大悔改認罪.可見聖靈激動人心,好像潮水之湧舟,其力非言語所能形容.我一次搭船往天津,因水退落,船乃擱淺,不能移動；後數時,甚至小漁船也同樣的不能移動半步.在這樣的情形下,出盡馬力也不能移動,只有等候潮水湧進,潮水一湧進,那船就自然而然的可以任意移動了.然今日的教會,正像那擱淺的船,若不得勝靈充滿,則不能有什麼進步；若得勝靈充滿,則必能有真的復興.
\newline
\newline
(一)真的復興,使人知罪自責.這知罪自責,並不是向非信徒說,乃是向信徙說.在復興會中,神的靈在人的心責備人.神使聖靈來世,是為保守人；倘人不順服聖靈,則聖靈離而魔鬼來.在真的復興會中,不獨聖靈責備人,且使人的隱罪顯明出來.人若有隱罪,則當在人前認出來；人若認了,則主必赦免.尤其是教會中的傳道,教員,及一切的領袖,若是犯了罪的,則更當快快承認,因為這些人犯了罪使聖靈懷憂及眾人跌倒.倘人受聖靈的責備而知己曾得罪兄弟,則當到兄弟處或在人前認罪；因為為眾人所知的罪,故當在眾人前承認.在北方教會有一會吏,因當初開復興會時他沒有空閒來赴會,及開到第四天便見這會吏痛哭禱告,他向眾認罪,說己的罪就是使母的靈魂落了地獄.因他曾痛罵他的妻子,他的母見他如此,便說,你既是這樣,我就不信你所信的耶穌了.他的母既不信耶穌,故他說:『我使母的靈魂落了地獄』.
\newline
\newline
(二)真的復興乃復回順服的心.主囑說:『人若愛我,就當遵守我的誡命.』而今我們試想一己曾否遵守主的誡命.神又囑人當守聖日,因為這日為一奇妙的日,乃紀念主為我們而捨己而釋放我們,使我們脫離罪的權勢之日.可惜多人犯主日而背主的誡命.主有許多誡命叫人遵守,但未重生的人則不能守.然聖靈既來指示我們當如何遵守,則我們當受聖靈的引導.在滿洲有一會吏,是管理樑木,他暗暗把樑木偷來,建自己的屋.有一晚聖靈感動他的心,在會中有數百會眾無一不痛心自責.最奇妙的就是那些沒有來堂敘會的人也在家裡同受聖靈的責備.就是那位偷木的會吏雖沒有赴會,但他雖在家中也受聖靈的感動而痛心自責.他在那夜覺得他不能過夜便要死了；他想把己的罪狀寫出,但因手不能寫,乃使子代筆寫出己曾偷了若干的樑木,並囑子到會中向眾讀出.但到了次早,他還未死,他於是到會中認罪.
\newline
\newline
(三)真的復興乃回復信上帝的心.許多人把上帝認作木偶；若聖靈在我們中施感動,則我們定然知上帝不是木偶而是生且活的神.前在武昌會集中,我聞一人禱告,他說:『我覺得神在我們的會中,顯明出來,好像可以用手摸到』.
\newline
\newline
(四)真的復興乃回復信聖經的心.聖經乃是神默示的.得勝靈充滿的人,則必知這經為靈的糧食；又知經的言語,為兩刃的利劍.
\newline
\newline
(五)真的復興乃回復祈禱的心.主耶穌最大力量,就是他的祈禱.在聖經中有許多大祈禱的人,這些大祈禱的人,都是能作大工的人.因此,復興會若是缺少了祈禱的人,那就必然不能得復興.今就廣州培靈會來說,若是各位不負祈禱的責任,則我們所望的也不能成.
\newline
\newline
(六)真的復興乃回復愛神愛人的心.所以有人說,要見仁愛的顯著,須到真復興的教會裡.我想這個會集,若有聖靈充滿,則這培靈會必震動全球；因為廣東人多僑居全球；若你們既受了聖靈充滿,則必發出愛兄弟的而心往傳道.聖靈若在人心得全權,則必打敗人愛世的心.教會所以失敗,在受社會轉移.聖經明說:「人若愛世界愛父的心就不在他裡面了.」門徒尚未受聖靈充滿時,便互相爭大,今日教會裡的人,也是這樣.若聖靈在這些人的心得了權,則必打敗人的爭大的心.
\newline
\newline
(七)真的復興,乃多人得救.你看五旬節後的門徒,他們一得復興,便因而有數千人得救,後來保羅在以弗所傳道,也得多人歸主；在高麗一地因得勝靈充滿,不滿兩月,便有二千人得救,及一年,便有五萬人棄偶像而事奉上帝.今我們若得勝靈充滿,也必同樣的感動多人歸主.各位,你們願不願見真的復興?倘若你們沒有這渴慕復興的心,主必不使你們得有復興.神的恩賜,不是賤的.所以你要納最高的代價.這最高的代價,就是犧牲自己,若肯犧牲自己來求主的教會復興,主必使大的復興顯出.
\newline
\newline
(八)真的復興,乃是以上帝為大.聖經說:「非以勢,非以力,乃以主之靈,事始能成.」由這經可知我們要成大事,則必須以上帝為大.前在美國紐約省曾開一復興會,當在會期內一日,斐利先生入一工廠,這廠有許多女子在那裡做工,他的心裡便有禱告的重擔,入時見了二位女子譏笑他,他便站著禱告,求聖靈感化她們；過一刻間,便見許多做工的女子大受感動而痛哭認罪.後來工廠司理,見這樣的情況,他雖未信耶穌,但他也覺主道的重要,便立命全廠停工而開會講道.於是在這次會集中有多人得救.這樣,我們若得了真的復興,必以上帝為大；若以上帝為大而靠聖靈,亦必像斐利先生一樣的引多人得救.
\newline
\newline


\newpage

\section{第1屆~港九培靈會~古約翰~第二講}
\label{sec:1326}
\textbf{私藏數分的價銀}
\newline
\newline
連結: \href{https://www.hkbibleconference.org/session-message/view/1326}{\texttt{https://www.hkbibleconference.org/session-message/view/1326}}
\newline
\newline
\hyperref[sec:1322]{< < < PREV SERMON < < <}
~
\hyperref[sec:toc]{[返目錄]}
~
\hyperref[sec:1327]{> > > NEXT SERMON > > >}
\newline
\newline
我們讀了使徒行傳第五章第一至第二這兩節經文,便知道亞拿尼亞同他的妻子撒非喇的罪狀就是自欺.他們以為自己是好人,其實,他們實在不是好人.他們因為貪名,便中了魔鬼的詭計.他們見巴拿巴因賣了自己的田地,把價銀盡獻而在教會中得了名,便生心貪得同樣的名,於是夫妻同謀,賣了田產,把價銀私自留下數分,只把其餘的獻與教會.論到他們獻一半價銀與教會,這本是一件好事；可惜他們說謊,不說是一半的價銀,而說是盡獻,因此,他們便作成了一件不好的事.他們誤會當時的教會,是彼得主理；卻不知是聖靈主理.聖靈降臨這個世界,是為使人榮耀主,今見教會中有這樣求自己榮耀的人,便為教會懷憂而施嚴厲的處置.因此,亞拿尼亞夫婦,便得一不好的結局.這樣看來,我們真要千萬儆醒,不要步他們的後塵而欺聖靈.因為,雖然人只能見我的外貌,卻不能見我的內心,但聖靈則能洞察人心的隱微.你看猶大衛十二使徒之一,他有份事主,人都以為他是主的好門徒.誰知他常懷貪心,雖見主行大奇能,但他也不決心照神旨行,而為主的真門徒.他若照神旨而行,則不懷藏貪念了.他既收藏幾分,故後來魔鬼在猶大心裡得了權,猶大乃喪心賣主.各位,猶大這樣的跌倒,我們若不小心儆醒,則不免做了猶大一樣的人.我們為主僕的人,當時常順服主；因為我們的戰爭,不是與屬血肉的,乃是與屬靈的魔鬼.我們勿藏仇敵在心,勿收藏價銀,試看掃羅,他因收藏數分價銀而失敗.初時,神甚是祝福他,故膏立他為以色列的王.後來神囑他滅盡亞瑪力人,且儆醒他勿在異邦人中污神的名,誰想掃羅奉命雖滅亞瑪力人並他們所有的,但他卻貪上好的牛,羊,牛犢羊羔,並一切美物不肯滅絕.神見掃羅這樣不順服,便對撒母耳說:我立掃羅為王,我後悔了；因為他轉去不跟從我,不遵守我的命令.撒母耳聽了,便甚憂愁,終夜哀求神；但神不聽,這是因掃羅已犯了罪.事後掃羅對撒母耳說:耶和華的命令我已遵守了.撒母耳說:我耳中聽見有羊叫,牛鳴,是從哪裡來呢?掃羅說:這是百姓從亞瑪力人那裡帶來的,因為他們愛惜上好的牛羊,要獻與耶和華你的神；其餘的我們都滅盡了.撒母耳說:耶和華喜悅燔祭和平安祭,豈如喜悅人聽從他的話呢?聽命勝於獻祭,順從勝於公羊的脂油.悖逆的罪,與行邪術之罪相等；頑梗的罪,與拜虛神和偶像的罪相同；你既厭棄耶和華的命令,耶和華也厭棄你作王.今我們是在主內為王為祭司的人,神已命我們要滅絕心內的亞瑪力人,這心內的亞瑪力人,就是我們心內的驕傲與妒忌,但不知我們有沒有順服神的命令.說到驕傲與妒忌.也許有人說,這惡性是由父母遺傳來的,我們實在沒有法子把他滅絕.不錯,我們實在沒法,但我告訴你們,你們若真靠主的主便能助你滅盡一切的驕傲與妒忌.甚望各位,勿藉詞沒有方法而寬恕自己,而收藏數分價銀.而今你們有沒有不良的思念?有沒有隱藏,亞瑪力人牛羊的貪心?神說:貪心與拜偶像同罪,故我們不應隱藏我們的貪心,應當把他滅絕.願望我們儆醒；不然,魔鬼一入心內,我們便軟弱如水了.
\newline
\newline
我們知神使聖靈來世,是為助我們高舉主；因為神的意旨,是願我們榮耀主.在五旬節前的門徒,多是歸榮自己,彼此爭大；後來,到了聖靈來執權,便把一切的榮耀歸於主.
\newline
\newline
我們知神能在昔日怎樣察看門徒的心,今日也能怎樣察看我們的心.
\newline
\newline
在山東有一位很有聲望的傳道人,他因對神的心改變了,神便不使用他.有一次,神遣我往一地主理奮興會,以上那位傳道人也在會中聽道受感而悔改,他悔改從前盜竊了神的榮耀.這位傳道人,只因想歸榮自己,故聖靈不與他同在,他便沒有能力傳道.由此可知,我們若凡事歸榮主,那就得勝靈同在；若得勝靈同在,則我們必能說出真理,使人知罪悔改.當高麗教會大奮興時,有人說,那時個個信徒,都是禮拜堂；因為那時的信徒,個個都傳道,甚至世人見他們這樣熱切的傳道,不耐煩而逼迫他們.我想人若肯獻身與主,則主必用之而震動這個世界.不知我們當中,有沒有像高麗信徒這般熱心的人?假若我們是不傳主道的啞吧信徒,我們就好像亞拿尼亞夫婦收藏數分價銀.惟是我們要傳主道的人,應當多讀聖經.前十七年,北方譯出官話聖經,我自此官話譯本出版到如今,在這十七年間,我已讀了五十二次；不單讀官話聖經,還讀英文聖經.然我們傳主道的人,不單要多讀聖經,還要以靈祈禱.若信徒欲傳主道而不以靈祈禱,這就是收藏數分價銀的人.不知我們有沒有這樣收藏數分價銀的人.各位要知,我們的工作,就是祈禱讀經.有一次我為河南一教會請一傳道人,不過兩月,那教會又請我把他辭退.我問他會不會講道,他們答會；又問講得好不好,他們又答好.我說:既會講道,又講得好,你們為何要辭退他呢?他們回答說:他雖是會講道又講得好,但是他說不講道,只貪睡覺.他天天晏起,起後到市上賣數本福音書便回來晚餐,餐後又是睡覺；因此,甚望把他辭退.我們想想,這傳道人是何等的懶惰,他真是亞拿尼亞一樣收藏數分價銀的人.
\newline
\newline
主說:「人若愛我,就必遵守我的誡命」.各位,主豈未告訴你,七日中當休息一日麼?想神在古時用他的大能大力,拯救以色列人出那受專制痛苦的埃及,故以色列人在每七日中遵命留一日以紀念此事.神特意要以色列人這樣；若以色列人遵命守之,那就必然得福.然在以色列人時,是律法時代；惟是我們今日,是在恩典時代.我們今日所守的,是主的日；這日是紀念主使我們,脫離魔鬼的專制痛苦.這樣看來,前之拯救是小；今之拯救是大但未知,在廣州究竟有幾許信徒守主日.我在這二十年內,主理過許多奮興會,我覺有許多信徒都好像亞拿尼亞那樣收藏數分價銀.然自五旬節奮興以後,凡是奮興的,必然尊敬地守主日.在高麗教會奮興時,每當主日,必見許多商店的門,掛著那主日休息的牌.許多人都以為守主日是很難的事；但真奮興的人,就沒有這樣的感覺.在廿年前,台灣之北,有一長老會的教會,會中有一姓林的長老,是一植茶的富人.那地也有許多信徒植茶且知茶當採的時期.那些茶一到當採擇的時期便要採,若是及時不採,則茶便失了風味,成為次等的茶.有一次,茶當採的時期恰是主日.那教會中有許多植茶的信徒,都是收藏數分價銀的人,雖是遇著主日,也必往採；惟是那姓林的長老則不往採而與家人同守主日.那時許多世人便議論這位長老,說他被他的神欺負了,但這位長老卻安然處之,因他愛神過於愛錢.後遲數禮拜,便有茶商到那地購茶.那些在主日採的茶；都得了上價；惟那長老的茶,因及時不採,成為次貨,所以茶商不購；因此世人更加譏笑他,甚至那些有名無實的信徒,也同樣的譏笑,惟這位長老,卻不因此而憂愁.誰想再遲兩禮拜,那些茶商又再到那地購茶,原來因各地的茶失收,故逼得再回,連那些次貨也用高價購買；因此,這位長老的茶價反多得四百元.因此,前譏笑他的人,反羞愧起來.各位,我們若真心愛主而守主的日,則主必然祝福我們.
\newline
\newline
真奮興的人,當十輸其一；因為我們是主用重價贖來的人.我們既是在主內得了恩典；那就當實行那「白白得來的,當白白施出」的教訓.十輸其一,是神在律法時代施出的命令.倘若中國信徒,能夠效法以色列人十輸其一,那就中國教會,不致這樣的貧窮而仰賴外國的金錢了.然我們要知,就是凡是收藏數分價銀的人,神必定懲責他.若是順服神命而十輸其一,則不致像那些不得神的祝福而耗失許多錢財的人一樣.今我們處在這個恩典時代,對於捐輸,當超過十輸其一.想五旬節時的教會,沒有一點欠缺,就是因沒有收藏數分價銀的人存在.從前高麗有一教會的教友,對我說及他們願供奉十幾位傳道人的款項.我說:你們的力這樣微薄,又怎能實行呢?今你們說這樣的話,是不是染了神經病呢?他們流淚答我說:主為我們犧牲一切,難道我們不當犧牲一切為主麼.各位,這樣熱烈的愛,我們也當有呢.願我們都有這熱烈的愛來愛主,不獨獻金錢與主,且獻兒女,使之事主.當像亞伯拉罕,只一子也獻與神.然我們要知,神只一愛子,為救世人緣故,而賜給世人.神給我六子女,倘此子女獻身來中國傳道,則我以此為最快樂.
\newline
\newline
有一次我在山東一教會聚會,見那會的傳道與長老,竟每一人願以子女獻與神.我想你們中有子能為總統,則你們必定喜樂,但我們要知,做總統的其實還不及做耶穌的代表.想我初信道後,便喜為一政治家；但我父母則喜我為一傳道人.主雖用重價贖我,但我收藏數分價銀,所以不願為傳道人,後因聖靈感動我,那時我雖獻身,但只願在加拿大傳道,卻不願來中國,我當時這樣順從自己,亦仍是收藏數分價銀的人.及後有一夜我聽知中國十分需要主道,我才決志來中國,到如今有四十多年了.
\newline
\newline


\newpage

\section{第1屆~港九培靈會~古約翰~第三講}
\label{sec:1327}
\textbf{要開墾你們的荒地}
\newline
\newline
連結: \href{https://www.hkbibleconference.org/session-message/view/1327}{\texttt{https://www.hkbibleconference.org/session-message/view/1327}}
\newline
\newline
\hyperref[sec:1326]{< < < PREV SERMON < < <}
~
\hyperref[sec:toc]{[返目錄]}
~
\hyperref[sec:1329]{> > > NEXT SERMON > > >}
\newline
\newline
耶利米書 4:3-4:4
\newline
\begin{longtable}{cl}
\hline
\hline
章節 & 經文 (和合本修訂版)\\
\hline
4:3 & \begin{tabularx}{0.7\textwidth}{X} 耶和華對猶大人和耶路撒冷人如此說：「你們要為自己開墾荒地，不要撒種在荊棘裡。 \end{tabularx} \\ \\ \relax
4:4 & \begin{tabularx}{0.7\textwidth}{X} 猶大人和耶路撒冷的居民哪，你們當自行割禮，歸耶和華，將你們心裡的污穢除掉；免得我的憤怒因你們的惡行發作，如火燃起，甚至無人能熄滅！」 \end{tabularx} \\ \\
[1ex]
\hline
\hline
\end{longtable}
在耶利米四章三至四節說:「耶和華對猶大和耶路撒冷的人如此說,要開墾你們的荒地,不要撒種在荊棘中.猶大人和耶路撒冷的居民啊,你們當自行割禮,歸耶和華,將心裡的污穢除掉；恐怕我的忿怒,因你們的惡行發作,如火著起,甚至無人能以熄滅!」這兩節經所說就是神儆醒他的子民的話.這裡所說的「要開墾你們的荒地」,這荒地究竟是說什麼呢?我前在山西時,也曾見過山西在一千八百十七年那次的大饑荒,死了許多人；及遲卅年,我又再過山西,便見許多的房屋傾倒,田生叢草,這田地前本種有禾麥,惟因無人料理,便成了荒蕪的地.人在這荒地播種,則必徒然,若要得好收穫,則惟有墾地與除荒草.神從前對以色列人說,「要開墾你們的荒地」,這就是指他們的心田.他們從前的心田,本是好的,但因而今變壞了,神就不喜悅他們的心田荒蕪,因此就命他們快把心的荒田開墾.神想人得復興,就是想人除去心田的荒草.人需要復興,就是證明人的現在,不像從前那樣的好,人所以不像從前那樣的好,是因不扺擋魔鬼而任他人心內隨播草種.今晚要說數種荒草,也許說不出你的荒草,但聖靈則能夠逐件指出來,若你得知了你的荒草是什麼,你便當受聖靈的引導把他除去.
\newline
\newline
戀慕世界的荒草　我們能不能免魔鬼的擾呢?這答話定然是一不字,因為主耶穌在世時,尚且不能免,所以我們自然是不能免的.從前主受試探時,魔鬼帶主上了一座最高的山,將世上的萬國,與萬國的榮華,都指給主看而對主說,你若俯伏拜我,我就把這一切都賜給你.想那時魔鬼雖是把這個的榮華誘主,但主充滿了聖靈,所以魔鬼不能入據主的心.可是,人與主是大有分別的；因為人是屬血氣的,所以每每墮入魔鬼的網羅.
\newline
\newline
昔日主在西乃山成立教會時,用了許多的證據,證明自己是神.及到後來,他遣他的獨生子來到這個世界,這個世界的人,因不認識他的獨生子而釘之於十字架.他們為甚麼不認識主耶穌就是神的獨生子呢?這就是因他們受了魔鬼的誘惑而戀慕世界.在五旬節時,主把他的教會建立在真理的上；因此,教會成為堅固的教會.可是,到了千年以後,教會便給空間與魔鬼；於是信徒乃受魔鬼的擾,而戀慕世界.教會處在這個末世,本當敬虔度日,等候主耶穌再臨；惟惜他們不是這樣,他們以為自己甚是富厚,卻不知自己的赤貧.若是教會能夠常常專心仰慕主,則魔鬼自然無法可施.
\newline
\newline
在河南省有一信徒,我以為他是個可靠的人.有一日,他來對我說,我想在北方開設教會,我回答他說,這是甚好的事.他又說,我將開一煤礦,把所得的利,十輸其一,用來供給教會,可是,而今那煤礦因官封禁不能開,因此,連教會也不能開,而今我想請你助我成功,給你的名片一張與我見官.我隨答他說:兄弟啊,你當千萬小心,不要被魔鬼使你陷於煤礦裡,這事我實不能助你的.他說:你既不助我,就證明你是不願教會發達的人.我說:我不是用人的方法來發達教會,乃是用神的法.我總不答允他的請求,他便去了,但過了兩禮拜,他又來對我說,礦已開了但被官押了數位兄弟去,你若不救則這數位兄弟必被斬首.我說我不肯助你做這樣的事.他說:你這樣欠缺慈心,見兄弟快要被斬首了還不往救.我說你不知被押中有一是魔鬼的工具麼?你與世人合做而求神的祝福麼?我又不答允他的請求,他便往求一神父,那神父即寫了一信與當地政府要求即時釋放那些被押的人.唉,在教會中,不少這樣受魔鬼的擾而戀慕這個世界的信徒,若在座中有人有這戀慕世界的惡草的,那就當立刻把他除掉.
\newline
\newline
驕傲的荒草:驕傲實在是害人的物,故許多人都是因驕傲而國破家亡身受苦痛.你們記得否?尼布甲尼撒王為什麼受神的責備,變為牛一樣的吃草七年呢?原來就是因他的驕傲.幸他在這受罰的七年中學得謙卑,故信神而得神的祝福.這驕傲的荒草,多種在從前的法利賽人中,故他們在殿堂中祈禱,竟居然自稱為義.然這荒草,在今日的教會中,也未嘗沒有.我們知教會的紛爭,多是由於驕傲,實教會既有驕傲的人,則教會必因而受惡的影響.
\newline
\newline
在河南省有一信徒,他初本是一可惡的罪人,後得主把他從罪惡中拯救出來,於是他就成為神的兒女.許多人見這樣的人也會改變,不禁因他而詫異.後遲一年,魔鬼便誘他驕傲；因此,那教會便受他的影響而不得進步.有一次,我探見他,他的母親是一位年將七十的老人,她對我說:你不要對我講耶穌,這個因我這樣的兒子,我就不喜歡聽了.他方與妻相爭罵,終底也不與妻和好,故他的妻說:信了耶穌的人是這樣的,我就不信耶穌了.唉,可惜她們不知使他成為這樣的人,是魔鬼的工作.我們有沒有這驕傲的荒草呢?若我有的,那就要快快把他除掉.
\newline
\newline
不良的品性的荒草:這不好的品性,已叫多人失敗.在一千九百零八年,有一教會開一奮興會,到會的人數,不過五十多人,這因為他們的學校卻不停課,以為讀書還重要過聽道.後數日,他們得勝靈感動,乃使學生赴會聽道,在主日那日,全堂都座滿了.那日我講道講到一半,我見得魔鬼也在堂作工,我見一西婦因性情不良與一姊妹失和而使此會大受阻礙.見了後我便請二人禱告,但二人的禱告,無一是有力的,於是我即宣佈在我們的會中,有人阻礙聖靈.
\newline
\newline
遲十月後,他們又再開復興會,在開會後第二晚,我便見那不良性情的西婦向從前失和的姊妹認錯.她認錯後,她的子亦因她把不良性情除掉而歸信耶穌.我們試想,她若不把不良的性情除掉,則她的子亦因她而不得救.各位,若你的心有這不良的性情,你也當立刻除掉.
\newline
\newline
妒忌的荒草:從前魔鬼把這妒忌的荒草撒在約瑟的兄長們中,他們便生妒忌而要把弟弟謀殺了,又撒在但以理的同事中,但以理的同事便心生妒忌,而要把但以理陷害了；又撒在猶大的領袖中,他們便生妒忌主耶穌的心,終而把主釘於十字架了.在一千九百零五年有一地有一教會開一復興會,神在會中所用的人,是別的地方的人,於是該地的領袖,便生妒忌的心,而使教會不進步,他們寧使人的靈魂沉淪,卻不把妒忌除掉.
\newline
\newline
以惡態度責人:我想做領袖的,做父母的,有時有責人的本份,但當責人的時候,有不可忘記的,就是以愛來責人,我信主耶穌責人,都是以愛.
\newline
\newline
在河南有一傳道人,他很愛教友,甚至願捨命為教友.有一次我聽見他說,「我甚傷心,一晚我求神用我,若不用我,則我寧願死了,」你看他因這樣的愛教友,真是難得的人.惟惜他雖是這樣的愛教友,但當他責人的時候,態度是很不好的,因此,信徒向我請求把他調往別地.後來這人得了一位牧師的忠告說；「當你責人的時候,你若用鏡把你面貌端詳一下,我信你一見了自己那樣惡形狀,你連自己也必憎恨自己起來.」尚幸他得了這個教訓,他便把惡態度除掉,終底成為一個很好的人.
\newline
\newline
在座為人家姑的姊妹啊,當你責備媳婦的時候,當拿鏡來照一照己的容貌,是怎樣的.從前在河南省有一售書的兄弟,他們夫婦本甚相愛,每逢她的家姑出外,則他們喜樂地度日,但她的家姑一回,則她在凡事上都是受苦的,因此她終因受苦不起而自殺了.
\newline
\newline
在你們中,若有人用惡心或惡的態度來責人的,則我望你快快把這個惡的態度除掉﹗
\newline
\newline
(一)毀謗人　你既說人的好處,則勿說人的壞處.有一傳教的醫生,他每逢有新來華的西教士,則必說人的是非.這人這樣的說人的是非,不是神叫他這樣,乃是魔鬼叫他的.這人既常執人的短處來非議人,故令他人難受.在我國有一俗語說,你若毀謗那狗,不如把他殺了還好,因一說他不好,則人人也踐踏他.試想,毀謗狗已不好了,若較狗更高貴的人我們竟毀謗他,這豈不是更不好麼?
\newline
\newline
從前有一猶太人,一日,有一鄰人對他說及一婦人怎樣的作卑賤的事,這猶太人本應對他的鄰人說,你快離開我罷!但他不把鄰人逐退,卻靜耳而聽,且隨聽隨說那婦人的行為.隨後有一晚他才確知那婦人是個好人,他立即找一位拉比,問該怎樣的做才好.拉比答他說,你明早拿一雞來見我.次早,他即拿一雞往見拉比,拉比說你當拿雞前行,但行一步須拔一雞毛,及到拔盡時才可回來見我.他於是照所吩咐的做去,把雞毛拔盡了便回來見拉比,拉比說,你當把所拔出的雞毛,一一拾回.他於是照吩咐的去拾,但拾來拾去只拾回三條,他乃對拉比說,請你不要責備我,我拾了許久只拾得這些,他說了這話,拉比便乘勢教訓他說,發出毀謗人的話則易,但收回則難了.
\newline
\newline
請想,我們有沒有這毀謗人的惡草,若是有的則我們當快快的除掉他罷!
\newline
\newline
(二)輕躁　聖經說憑你的話,定你為義或不義.這樣看來,故我們要謹慎我們的言語.在河南省有一兄弟,甚好口才,他所說的詼諧,能夠令全堂的人笑倒,但惜他為人輕躁,故我雖需要傳道的人,但我因他的輕躁總不敢用他.
\newline
\newline
各位,有沒有這輕躁的惡草,若是有的,我望你快快除掉了罷!
\newline
\newline
(三)謊話　我們是信徒,也許沒有說大話的惡草；但俗語有說:「無直接的大話而有間接的大話.」又北方有一俗語,說大話有彎彎曲曲的大話.我們要知,大話是由受了魔鬼的擾而來的.從前我曾說了一句間接的大話而受神責備.當時因我在滿洲,主理復興會,有一人,向我問及一位姊妹,我本知那姊妹甚愛華人,就是為華人捨命也是願意的；但她也有給人說閒話的短處.當時我只以她的短處回答那人,沒有說及她的好處.這樣我所說的話,豈不是成了一幅半頭半身而為人所不願看的相麼.當我說了這人的那夜我便照常講道,講後,請人禱告,總覺全無能力我心甚憂,因前日已見主已動工,今夜竟是這樣,我便對神說,求你指明這個阻礙在哪裡,神說:阻礙就在你,我即求主赦罪,決意再對那問我的人說及那半未說的.及我認罪後,聖靈便作工.從這裡看來就知我們的大話,是如何的阻礙教會進步.
\newline
\newline
各位,若有這惡草,請你快快除掉了罷!
\newline
\newline


\newpage

\section{第1屆~港九培靈會~古約翰~第四講}
\label{sec:1329}
\textbf{從墮落的地方起來}
\newline
\newline
連結: \href{https://www.hkbibleconference.org/session-message/view/1329}{\texttt{https://www.hkbibleconference.org/session-message/view/1329}}
\newline
\newline
\hyperref[sec:1327]{< < < PREV SERMON < < <}
~
\hyperref[sec:toc]{[返目錄]}
~
\hyperref[sec:1330]{> > > NEXT SERMON > > >}
\newline
\newline
我們讀了「啟示錄一章一至七節」便知以弗所教會是一個完全的教會,但照主如火焰的眼目看來,就有許多未稱主心的欠缺,有一件事是要受主責備的,就是他們「把起初的愛心離棄了,」離棄了起初的愛心主總覺他們無好處,因為這個「起初的愛心」沒有別的美德可以彼得上來作個替代.主升上後已打發他應許的聖靈來臨,住在信徒裡面,所以我們應當結出仁愛的靈果,發出仁愛的靈光才是.聖經上記著說:「我若能說萬人的方言,並天使的話語,卻沒有愛,我就成了鳴的鑼,響的鈸一般；有全備的信,叫我能彀移山,卻沒有愛,也就算不得甚麼；我若將所有的賙濟窮人,又捨己身叫人焚燒,卻沒有愛,仍然與我無益......」主是要我常有那初時的愛,甚望各位不要推卻說不能恆久存著那「起初的愛心」；要求聖靈為我們保存那初時的愛,因為忠心愛主的人,必然得著聖靈保存他的愛心能彀長久.從前本仁約翰曾見一異象,看見一人不住的倒水入熾旺的火中而火仍燃燒不息,我們都知水是可能滅火的,但今竟然不滅,那就希奇了,誰知火之不滅,因為後面有一人不斷地把油傾入,所以那火可以常燃不滅.我們在這裡就該得著靈意的教訓；倒水入火的人可以表喻魔鬼,因為牠是常常不住的要撲滅信徒靈裡之火熱的；那個傾油入火的人,可以表喻聖靈,因為聖靈不斷地把他的恩膏傾注入我們靈裡,使靈裡的火更要熾旺起來,現在聖靈已到來了,我們若是給他充滿我們,他就常常保有我們愛主的愛,由此可知我們在主的面前,沒理由推諉我們無初時的愛了.
\newline
\newline
初時的愛是甚麼?
\newline
\newline
(一)有初時的愛就是始終跟從耶穌的
\newline
\newline
施洗約翰當耶穌由水而上的時候他就為主作證,後來他的門徒和猶大人因辯潔禮來問他,他答他們說:「他必興旺,我必衰微.」可見施洗約翰此時仍有初時的愛了,耶穌的門徒撇下一切來跟從耶穌到底,可見他們初時的愛了,保羅丟棄萬事,奉主名往各處傳道,榮耀主的名字,這豈不是保羅起初的愛心嗎?八十年前在加拿大有位神道大學的學生,畢業後獻身往南洋傳道,但那處的人,是會食人的,果然不夠兩年他母親便收到音信,說他的兒子在南洋已經被野人食了.那時他的弟弟正在神道學校畢業回來,幫忙他的母親耕田,他的母親是個寡婦,看見母親得了哥哥的信息,就接過來詳細的讀給母親聽聽；讀了,「心裡自己說道,倘若我去南洋代哥哥作工就好了.」便將信給還母親,但不敢把心事向母親說出來.母親聽了大兒子已為土人食了,就說:「我的細兒子,能夠代替我的大兒子再去南洋傳道,引導那些食人的土人歸主就好了.」說完就注望著細子的眼兒,知道他的細子是願意去的,後來他們的細子去到那裡,剛剛三年,又被那裡的土人所食.」在這裡可見他母子們起初的愛心了.
\newline
\newline
(二)有初時的愛就是始終高舉耶穌的
\newline
\newline
凡是有初時的愛之人,必定高舉耶穌的,施洗約翰對門生說:「看哪!上帝的羔羊......」他早已知道高舉耶穌,是會失去自己的地位的,但他願意「耶穌必興旺,自己必衰微.」從前主的門徒能高舉耶穌,所以得了許多人來歸主.保羅寫信給哥林多的教會兄弟說:「我已定意在你們中間,不知道別的,只知道耶穌基督並他釘十字架.」我們若是仍有初時的愛,我們的口是不會啞的,是高舉耶穌而常常傳講他的十架的.聖靈來了不單是門徒講耶穌,凡聽道而信主的人,也是會講耶穌的.廣州市信徒這麼多,但為主作工的那麼少,是不是初時的愛消滅了麼?昨天已說過高麗的教會得了奮興之後,他們家家的人都信了主,各家皆變成福音堂,天天傳揚主的道理.所以我們有初時之愛,這愛必然鞭策我們,多多傳講主的道理,教會因此就會得勝經的奮興了.當一九零七年的時候,在日本有一個大奮興會,起初是因為有幾位西國傳道人得聞英國某地得了一個大奮興,於是他們為著自己的教會,懇切祈求得著奮興,不久就有一位傳道人竟然被聖靈充滿,無論往何處作工,總是得著奮興的.後來有一位監獄長,請了這位傳道人到一個監裡去講道,這監裡之犯人多是犯了殺人罪的,聽道數天,就有數百人信主,獄卒也有許多得救.其中有一位反對不願聽道,當日俄交戰的時候那獄卒因為捉了一俄人,日皇給他一個獎章,他更驕傲起來,那些得救的獄卒向他講道,他反要怒罵他們.有一次在監獄裡開會佈道時,這獄卒也來赴會,聽道後就大受靈感悔改認罪,即在這天領洗歸主,過了三天,他竟然引了四十獄卒歸主.以上所說的都是有初時的愛,能夠高舉基督的.
\newline
\newline
(三)有初時的愛就是不願使聖靈懷憂的
\newline
\newline
保羅到以弗所傳講基督,主耶穌的名從此就在那裡尊大,平素行邪術的,也有許多人把書拿來,堆積在家人面前焚燒,他們計算書價便知道共合五萬塊錢.他們知道要與魔鬼脫離關係,這是令聖靈喜悅的,我們有初時的愛,就像一本活的聖經一樣,世人是不會找聖經看的,倘若我們仍有初時的愛,便能夠活活地把聖經的道理行出來,別人就因我而識主,以致得救；倘若我行魔之行,別人必因我而傾跌,聖靈因我而憂傷,你看以弗所的信徒,他們將所有令聖靈懷憂的書籍焚燒淨盡,這就是他們的初愛了.
\newline
\newline
(四)有初時的愛就是能夠犧牲為主的
\newline
\newline
我們讀經知道五旬節的教會有欠缺否?高麗的教會得了奮興之後,有欠缺否?從前在北印度有一位信徒,有數百銀,不信銀行暗埋土中,從此心裡常常為著這些土中的銀掛念擔憂,埋了數天,更恐怕他人偷竊,屢屢半夜起來遷移所埋的銀,另埋他處,一連幾個月都是這樣的在夜半中勞苦.有一次開大奮興會的時候,他就得了聖靈之感,立刻掘出所埋的銀盡獻歸主,教會司庫請問存記何人捐款,他就說不用書名,書則辱我極甚,祇書收銀若干便妥,這位信徒從此得回他初時的愛了.我們若有初時的愛,也能夠這樣為主犧牲的.有一次我在山東開奮興會,有一位會吏受了感動,情願將四個兒子全獻給主,求主接納為傳道人,這個會吏的初愛就復原了,倘若我們初時的愛復回,不單能夠捨棄一切歸主,並且生命也可全獻歸主的,我有一次由滿洲回到北京來,有一會堂請我說滿洲的奮興情形,時有一中學女生,禱告說:感謝主把靈恩降下滿洲,求天父把靈恩降下此城,因為人民之心好像往日滿洲人心那樣乾旱......」我聽了她的禱告,心裡很受感動,看他的面貌好像天使一般,因為她充滿了聖靈.再遲十個月那堂又請我到那裡講道,我講道之時覺得會中有阻礙,後來我問該堂的牧者,他說是會吏不喜歡你來,因為你喜歡人悔改認罪,所以他們不獨自己不來,連別人也要攔阻他,因為庚子年當皇帝逃走的時候,這些會吏們快快入到宮裡去偷了許多財寶.我到堂講道之時,那皇帝經已回到宮裡,他們以為認罪必為皇帝知悉而生命難保,所以他們的心仍是剛硬著.到了奮興會最後那一晚,我又請會中的兄弟姊妹懇切的禱告天父,但是禱告總無效力,我心很是憂愁,恰巧那位中學女生又開聲流著淚禱告說:「父呀!我知道我會的領袖和教友,又知阻擋此次奮興會而不肯使你旨意成就者是誰；父呀!你是否要得著一個犧牲,然後奮興我們的教會嗎?若然是的,我願意作這個犧牲,若有益於你的教會我願主塗抹我的名落地獄,這個犧牲能夠除去教會的阻礙,我甘心樂意的犧牲獻給主.」這就是這位女學生初愛的犧牲了.初時的教會,總是滿有初時的愛的,我記得庚子年的時候,在直隸省有父子二人,為拳匪捉去,後來對他們說,若不肯拜菩薩,則生埋土中.他們說:耶穌為我捨命,我亦不怕死,為耶穌捨命.那些匪徒就把他生埋土中,其子年方十四,亦同埋在父旁痛哭流涕,匪就對子說,你肯拜菩薩就釋放你.那小子說,我決不違背所信的道,寧與父同死,同得的賞賜.匪就大罵他說:番鬼第二也有此膽志,乃掘坭埋及兩父子的口說,快快悔改,不然,一刻鐘便要死了,他們說,死也不悔.這就是有了初時之愛,甘心樂意為主的犧牲,我們有沒有這樣的犧牲呢?你的初愛是不是離棄了呢?聖經記著說:「所以應當回想你是從哪裡墮落的,並要悔改,行起初所行的事.若你不悔改,我就臨到你那裡,把你的燈臺從原處挪去.」(啟二之五節)這話是「那右手拿著七星,在七個金燈臺中間行走的耶穌向以弗所教會所說的,也是向著今日廣州的教會來說的,願「離棄起初的愛心之教會」快快悔改,從墮落的地方起來,免得主臨到你那裡,把你的燈臺從原處挪去,阿們.
\newline
\newline


\newpage

\section{第1屆~港九培靈會~古約翰~第五講}
\label{sec:1330}
\textbf{禱告有效的秘訣}
\newline
\newline
連結: \href{https://www.hkbibleconference.org/session-message/view/1330}{\texttt{https://www.hkbibleconference.org/session-message/view/1330}}
\newline
\newline
\hyperref[sec:1329]{< < < PREV SERMON < < <}
~
\hyperref[sec:toc]{[返目錄]}
~
\hyperref[sec:1331]{> > > NEXT SERMON > > >}
\newline
\newline
路加福音 11:10-11:13
\newline
\begin{longtable}{cl}
\hline
\hline
章節 & 經文 (和合本修訂版)\\
\hline
11:10 & \begin{tabularx}{0.7\textwidth}{X} 因為凡祈求的，就得著；尋找的，就找到；叩門的，就給他開門。 \end{tabularx} \\ \\ \relax
11:11 & \begin{tabularx}{0.7\textwidth}{X} 你們中間作父親的，誰有兒子求魚，反拿蛇當魚給他呢？ \end{tabularx} \\ \\ \relax
11:12 & \begin{tabularx}{0.7\textwidth}{X} 求雞蛋，反給他蠍子呢？ \end{tabularx} \\ \\ \relax
11:13 & \begin{tabularx}{0.7\textwidth}{X} 你們雖然不好，尚且知道拿好東西給兒女，何況天父，他豈不更要把聖靈賜給求他的人嗎？」 \end{tabularx} \\ \\
[1ex]
\hline
\hline
\end{longtable}
古今以來,少有人的禱告能像使徒那樣的有能力.惟今願各位的禱告能像使徒一樣,所以特把禱告有效的秘訣向各位講述:
\newline
\newline
(一)虔敬　虔敬的心,是禱告的人所必須有的.可是我有時見人的禱告,反不樂起來；這因為有些人把自己與神列為同等,更有些人好像還要高神一等；他們不知神是萬王之王,人是受造的物,這樣,受造的人,那配與那位造物的神同等.你看主耶穌每逢禱告,尚且必恭必敬,他雖配與神同等,還自卑不已,更何況我們呢.我們研究聖經,都知在禱告有功效的人,都必有虔敬的心.如但以理雖為巴比崙國的大人物,但他禱告神時,卻雙膝跪地,謹存謙卑的態度虔敬的心.又如亞伯拉罕也謙謙卑卑的代所多瑪向神禱告,把己比作灰塵.
\newline
\newline
(二)認罪　我們向神禱告須先認罪；因為我們都是虧缺神的榮耀的人.神若以人的行為來論定人,則無一是公義的人；所以我們惟靠主血才能見神.如昔日的以色列人,無一人能入至聖所,惟大祭司靠著血纔能入去.今主用己血為人開一生且新的路,使凡靠他血的人,都能進到神前.回想當時的法利賽人,真是善頌善禱,可是他們的禱告總不得著功效.為什麼呢?這是因為他們總不承認己罪.惟先知以賽亞,本是聖潔的人,但他一見神便自嘆息而認己是污穢不堪的人.但以理亦一個聖潔的人,當他禱告時,也認己罪及眾民的罪.今我們若藏惡,則神定然不聽我們的禱告.
\newline
\newline
我有一次在滿洲一教會開奮興會,當未開會前,那教會裡有位西國的傳道人對我說:「此教會的領袖,沒有犯大罪的人,若是有的,連我也蒙羞了.」及開會到了第三日,聖靈便感動人心,人人都爭著禱告；惟是各領袖總不禱告,我便對以上所說那位傳道人說:「教會的阻礙,就在你們的領袖.」那位傳道人的妻,就回答我說:「不然,若是阻礙就在領袖,我們真是羞恥了.」但到次晚,我就見一位體格很強壯的人跪著禱告,此時我見人人都跪著禱告,這是最奇的事,因為那會堂是長老會的,平時禱告總是企著.那強壯的人對我說:「當我來赴會時,我早就決志不認罪；惟是聖靈感動我,叫我的心大大不安,甚至不能飲食.」他又認罪禱告說:「父呀!你知我是個犯第七誡的人,又從前有一會吏犯了大罪,我雖知道,卻不講與牧師知,這因那會吏送我一件皮袍,那皮袍就塞了我的口.」那時他禱告完了,就隨即把皮袍撕碎.原來那人不是個普通教友,卻是個傳道人,那晚因此人禱告認罪,於是全堂的人也受感動而各自傾心禱告.
\newline
\newline
(三)償還　犯罪是魔鬼的工作,若我們要反魔鬼的工作為神的工作,那就必須先毀壞魔鬼的工作.你看撒該,他在耶利哥作稅吏長時只知訛詐斂財；但當一悔改了就對主說,「我若訛詐了誰,就還他四倍」.他既肯認罪,又肯償還,所以主就回答他說:「你真是亞伯拉罕的子孫,今天救恩臨到你家了.」
\newline
\newline
有次我在滿洲一教會主理奮興會,在那教會中有一會吏因為貪財的緣故,便拿五百元與人合股做香燭的生意.有一教友知道這事,就隨處對人說,我們的會吏做那些不正當的生意,求得些不當得的錢.試想這教友不為會吏禱告,只說是非,這是應當的麼?誰知會集到了第二晚,我就見那會吏流淚禱告,願意取回五百元股本獻與教會；不特如此,還要罪己四百五十元,一併獻與教會；最後又求神祝福,且願意每年奉獻九十元.次日,以前說會吏是非的教友便與我到會吏的家裡一同跪下禱告,他且向會吏認罪,他這樣的做,也是償還.
\newline
\newline
在座各位信主的兄姊啊!主已發出「白白得來的要白白施出去」的命令了.他要你我傳福音我們若不去傳,那就多人因你我不傳而淪亡,後來主必和你我算賬.既是如此,所以我們當快快的償還福音的債.
\newline
\newline
(四)饒恕　我們要得主的赦免,就要先赦免他人,所以主禱文中有說:「你們饒恕人的過犯,你的天父也必饒恕你們的過犯；你們不饒恕人的過犯,你們的天父也必不饒恕你們的過犯.」有一次我在北京開會,有一傳道對我說:「求你代我禱告,因我不禱告已有三月了.這因為在庚子年時,拳匪把我的家人滅盡,我為傳道人,本不應報仇,但我有一友代我報仇,把殺我家的拳匪拿著,那匪有子,他想把二子也找來一併殺盡；但因拳匪的鄰人把二子收藏,所以找不著,我的友便把那拳匪殺了.此後我便移居到城外,我的友對我說,你若查出他的二子所在,我便來殺他.後來查了兩年就查著了,我便對一西牧師說知,我以為他會助我,誰知他卻對我說:「最好把他們帶來學校讀書,你供給他們的費用.」我聽了這些話,便很忿怒的回答他說,這事我怎能做,因為他們是我的仇人.無幾我的友人來了一信說:「你是孝子,我已為你做了一件事,但你一件還未做,就是那二子你還未查出,使我好把他們殺了；今我因你的事未了,使我一生不能回家.」當我一得此信,報仇之心更切,便快快通知他.唉,我自有此報仇之心,便失去禱告的心.今我在此數日會集中,受聖靈感動,定意不報仇,且助我仇人的二子入學.」
\newline
\newline
各位!我們從這事看來,人一有不饒恕他人的心,同時就失去禱告的心,所以我們要饒恕他人.
\newline
\newline
(五)信心　人若沒有信心,就不能得神的喜悅.我們從聖經而知有信的人,才能得著聖靈充滿.當日掃羅為以色列人的王時,與非利士人爭戰,那時非利士人的營中,來了一位勇猛的討戰人,這人就是哥利亞,在以色列人軍中,沒有一有信的人靠著上帝的大能,敢與他交戰.後來大衛因奉父命到營探望他的哥哥,因而見哥利亞向以色列人罵陣,他因有信,就靠上帝的大能而擊殺了威武絕倫的哥利亞.今在廣州也有許多罵陣的哥利亞,我們當像大衛一樣的信靠主而得勝也.
\newline
\newline
(六)一心　合千萬心為一心是件頂重要的事.比方珠江的水力強厚,所以能夠渡香港的船來到廣州,若把珠江的水力分而為十為百,那就沒有一船能來,這因為水力弱了.中國本是大國,照理應異常的富強,今竟這樣衰弱,其因由於不一心；國固然要合千萬心為一心,就是教會也要合千萬心為一心；若能合一,那就自然勝敵.我們都知兩牝山羊彼此相遇,就必相鬥,這是他們的特牲.但有時相遇卻不相鬥,如他們在橋板相遇時,照特性就必相鬥,惟是他們也有思想,知若相鬥就必相亡,於是一就自然而然的伏下,一就跨過,彼此就相安的過去.今我們若能學效他們,那就沒有紛爭了.惜我們少有人肯伏下,所以我們寧願教會分開,卻不肯降卑自己而使教會合一.想五旬節時所以得以聖靈降臨,由於門徒合心的禱告.今要培靈會得勝靈降臨,也要我們合心禱告.假若不得勝靈充滿,這不是神的錯,乃是我們不一心的錯.
\newline
\newline
(七)感謝和讚美　我們的禱告不單當求主,還要感謝和讚美主.大衛為王時,雖國政紛煩,但他卻能常讚美主；所以他所作的詩,多是讚美主之詞.又保羅雖擔負重多煩難的責任,但他仍能常讚美主.今日我們的責任,沒有像大衛保羅那樣的煩難,所以更該常讚美主.
\newline
\newline
(八)順服神旨　神要做我們的主,故我們的主權當奉給他,且當凡事只求他的意旨成就.然我們在順服中要凡事小心,不要中魔鬼的詭計；因為我們有時為父母或為親友禱告,但總不見主應允,到那時魔鬼就會對我們說,神不欲救你的父母親友了.各位!當你遇著這樣的聲音時,就該小心,更該加倍的順神旨.你看主耶穌一生是何等的順服神,就是死亡來到也不改變他對神的順服心,今我們無論遇什麼事,也當像他一樣.
\newline
\newline
(九)恆心　當主在世時,耶利哥城外有兩個瞎子坐在路旁,聽說是耶穌經過就喊著說:「主啊,大衛的子孫,可憐我們罷!」但當時眾人責備他們,不許他們作聲,他們卻越發喊著說:「主啊,大衛的子孫,可憐我們罷.」你看,他們是何等的恆心,我們知恆心的禱告,定不會落空.如慕勒牧師為他的友一世不斷的禱告,有兩友到他死時還未得救；但到他死後,他的兩友便悔改歸主.又如在英國有一婦人為丈夫禱告,一月又一月,到了十七個月還未見主應允；及到十八個月,有一日他的丈夫從外回家,徑直上樓去,婦以為丈夫洗手預備食餐,但許久還不落樓,婦就希奇起來,於是上樓窺探,見丈夫流淚的禱告,他就知神已答應他的禱告了.
\newline
\newline


\newpage

\section{第1屆~港九培靈會~古約翰~第六講}
\label{sec:1331}
\textbf{有效力的禱告}
\newline
\newline
連結: \href{https://www.hkbibleconference.org/session-message/view/1331}{\texttt{https://www.hkbibleconference.org/session-message/view/1331}}
\newline
\newline
\hyperref[sec:1330]{< < < PREV SERMON < < <}
~
\hyperref[sec:toc]{[返目錄]}
~
\hyperref[sec:1328]{> > > NEXT SERMON > > >}
\newline
\newline
雅各書 5:13-5:18
\newline
\begin{longtable}{cl}
\hline
\hline
章節 & 經文 (和合本修訂版)\\
\hline
5:13 & \begin{tabularx}{0.7\textwidth}{X} 你們中間若有人受苦，他該禱告；有人喜樂，他該歌頌。 \end{tabularx} \\ \\ \relax
5:14 & \begin{tabularx}{0.7\textwidth}{X} 你們中間若有人病了，他該請教會的長老們來為他禱告，奉主的名為他抹油。 \end{tabularx} \\ \\ \relax
5:15 & \begin{tabularx}{0.7\textwidth}{X} 出於信心的祈禱必能救那病人，主必叫他起來；他若犯了罪，也必蒙赦免。 \end{tabularx} \\ \\ \relax
5:16 & \begin{tabularx}{0.7\textwidth}{X} 所以，你們要彼此認罪，互相代求，使你們得醫治。義人祈禱所發的力量是大有功效的。 \end{tabularx} \\ \\ \relax
5:17 & \begin{tabularx}{0.7\textwidth}{X} 以利亞與我們是同樣性情的人，他懇切地祈求不要下雨，地上就三年六個月沒有下雨。 \end{tabularx} \\ \\ \relax
5:18 & \begin{tabularx}{0.7\textwidth}{X} 他又禱告，天就降下雨來，地就有了出產。 \end{tabularx} \\ \\
[1ex]
\hline
\hline
\end{longtable}
今晚的題目就是「有效力的禱告」,有效力的禱告,上帝的子民也曾禱告過:先知以利亞的禱告,也曾得了很大的功效,從前以色列人忘記他們的上帝去拜那異邦的巴力,上帝藉著這位先知挽回許多背逆的百姓,同被擄去巴比倫的但以理亦曾得著禱告的功效,得上帝依期釋放以色列人歸回故土；當日使徒從橄欖山回到耶路撒冷時的禱告也有這樣的大效,得著聖靈降臨奮興教會,內地會的創始者戴德生也同得著這禱告的實效,現在已有一千西牧,在華傳揚主道,並有數千華人牧者及傳道人同作主工,這個教會每人的需要,全是靠賴天父供給的,在歐戰期間,美國的加拿大雖是用了許多錢財去作戰費,但內地會未曾有所缺乏,這內地會就是給我們一個憑證,證明他們的禱告是有效的.
\newline
\newline
(一)有效的禱告,須要合上帝的旨意.在五旬節有一百二十人同心合意,憑著應許來禱告上帝,因為先知約珥曾預言說:「上帝說,在末後的日子,我要將我的靈澆灌凡有血氣的......」所以使徒們憑著這應許來禱告,後來彼得也為著這應許來作證.我們能夠按照主的應許來禱告,便是合乎神的意旨而得神的答允.當主耶穌未升天的時候,門徒問及復興以色列國的事,主答他們說:「父憑著自己的權柄,所定的時候日期,不是你們可知道的,但聖靈降臨你們身上你就必得著能力......」所以他們就執著這個應許在所住的樓房一心禱告主,主就把屬靈的恩雨下降他們的身上,為主名成就大工.高麗一八零七年的大奮興,是由合神旨的禱告而成的,因為當日那地的美英的牧者在高麗的京城會議時,恰巧有一位美國人由印度來將印度教會奮興的情形報會眾說:「印度此次之得大奮興實由美國有多位信徒代禱的能力,使兩月間得有六千人得救......」各人聽聞這個消息,心裡很是羞愧,因為高麗從來未有這個大工,於是有二十位西國的牧者興起說,印度教會既然得了主的奮興,高麗教會當然也得著主的奮興,他們既然有了這個篤信和切望,就定了一定的時間,日日為高麗得著大奮興來懇懇切切的禱告主,經過四個多月的工夫,聖靈就充滿了高麗的教會,一切工作就大勝從前了,在兩個月當中,得有二千人歸信主的道理,一年中有五萬人離棄偶像歸入基督名下.高麗的信徒執著上帝曾賜靈恩給印度教會來做應許禱告主,便得了聖靈的大奮興.
\newline
\newline
在廣州曾有這樣的大奮興否?兩個月中曾有二千人歸主否?廣東一省的人民,多過高麗全國,你想上帝愛廣東人多過愛高麗人麼?主願意把聖靈充滿我儕如高麗一樣,我們當有禱告的重擔,不單為本國教會禱告主,更當為他人為別國來禱告主,五旬節時的使徙得了聖靈便有代禱的重擔在心,為萬國來禱告.從前英國曾有一次很大的奮興會,兩月中有七萬人得救,從此這地方就有多人有禱告之重擔來為各國各地的教會來禱告主,各地的教會就因而奮興了.你切望廣東教會得著大奮興,當先負起代禱的重責.
\newline
\newline
(二)有效的禱告須要獨自與上帝交易清楚.禱告有效的人就是與神沒有隔膜的人,以利亞和但以理是個與神無隔膜的人,所以他們的禱告,得著了很大的功效.在五旬節中的一百二十人,個個是與神無隔膜的,所以他們就豐豐富富得了主的恩賜.今晚是培靈會的第六晚,倘若第一晚人與神交易清楚,聖靈早已丕顯在我們中間了,今晚是第六晚,聖靈尚未顯見在人心作工,這緣故,就是人與神仍未交易清楚.我曾與主交易清楚,從前在河南彰州與一位西國牧師失和,這次失和本來是那位牧者不合,但人有錯,每每推諉他人,這是人的常情,我們若與人不和而有所紛爭,自己亦當分負一半「不著」的責,不可完全諉過他人,有不和即當認罪.那時我正要往各處主理奮興會的講道,當我講道之時,聖靈的細聲責備我說:「你是偽善的人,你若與人不清楚,我必不與你同往.」我聽聞了這個警告後,心裡很是憂愁,因為聖靈不與我同工,就不會為主名成就甚麼了,但我的心仍是推諉,以為那位西國的牧者應當先來認錯,然後我與他和睦才是,誰知聖靈要我先去認錯,我說不是我錯,後去北京作工,頭一晚我主理祈禱會,聖靈的細聲又來責備我說:「你是偽善的,你未與兄弟和睦,不可作工.」那時要開會了,但是聖靈不斷地在我心擾我,於是唱詩之時心不安,禱告之時,也不曉怎樣禱告,後來講道的時候,聖靈又來責我,使我連講道也不會講了,我的口一面講,但我的心就一面戰爭,終底聖靈得勝了我,我就決心散會後即往與那位牧者認罪.我心一決,當晚的祈禱會中,看見一人在會中流淚,那夜我與神交易清楚,眾人之中多有流淚的,此後不但禱告有效,講道也有效了.所以我們應對神對人總要無虧負,無隔膜,主才用我們,主亦常常等候得著這樣的人,當以利亞的時候主要等用一人,就是以利亞這樣的一個人禱告,神就藉著這人的禱告,打破魔鬼的工作.深望中國有五百個像以利亞這樣的人,使中華早日歸主.在但以理的時候,上帝亦等候一人,後來得著但以理,藉他的禱告,使選民強敵的手下,釋放回去祖國.主耶穌升天之後,亦需一百二十個清楚的人,同心禱告,教會就得奮興.上帝今日亦需要廣東教會這樣的禱告人,來負起禱告的重責,你們願意做一個禱告的器皿任神使用麼?
\newline
\newline
神的旨意是要我們充滿聖靈,好作有效力的禱告.或有人說,我們不像以利亞為大先知,但聖經說以利亞是與人同性情,無有差別的,所不同之處,就是充滿主的聖靈,我們若然順服聖靈,聖靈就使用我們,我們的禱告就有大效.從前在愛爾蘭有個很軟弱的禱告器皿,這女子名安拿,她在學校六七個月,學ABC也學不成,醫生說她缺乏腦力,後來這個女子長大了一字不識,但她的靈仍然能夠感受聖靈而得救,她歸主之後,因不識字而禱告天父說:「天父呵!我不能讀你的書信(聖經)我是神之子女為何不能讀神的書呢?我不求主使我能讀別的書,獨求主使我能讀你的書,我心就滿足了......」後來神就答應她的禱告,她的禱告總是有效的.她曾服事一位醫生,這醫生的廚房有一井,久已乾涸無水,遠處之井則有水,當日有一人問她說:「你的天父既然是應允人的所求的,何不請他命這個乾涸的井有水呢?」那位女子安拿就去禱告天父說:「亞爸呵!求你令這個井有水,免得這人不信你而致落地獄.」那個人第二天早晨去取水,就見乾涸的井竟然有水了.她的禱告真是有效的禱告了.
\newline
\newline
主甚是喜歡我們進入他的禱告學校,我禱告無效是主不答允麼?或者是我有阻礙呢?彼得不認主他的禱告無效了；羅得往所多嗎戀慕貨財,他的禱告無效了；老底嘉的教會,雖自誇其富厚,但是禱告無效了；我們也是不冷不熱的,怎能得禱告的大效呢?我們要知禱告,是最重要的,主的門徒無別的可靠,惟可靠的是主.英國的莫勒先生,曾設立一間孤兒院,內有孤兒二千,理事的人員有一百五十個,全院共計二千一百五十人,每年皆無別的款支給,每日的需要,全賴上帝供養,一切皆由莫先生禱告天父而得,今晚食罷,明日不知怎樣了,但莫先生積數十年的經驗,從無一日有欠缺的,因為他的禱告是有效的禱告.禱告是有價值的,我們若知道禱告的價值,就會負責代他人禱告了,我們已先得著主的拯救,主願意我們引導他人得救,所以他人下了陰間,實是我們的責任,代人禱告的工夫,是不是可以輕看的呢?馬丁路得每日的工作很忙,他用兩個鐘頭的時間來禱告,工作加多,祈禱的時間同時加多.莫勒先生已得過五萬次有效的禱告,我們有懇切的心,來禱告天父,他必為我們成全的,我們人人皆有禱告的責任,停止祈禱是阻礙聖靈的工作,我們有祈禱的能力而不用他,難免主會向我們懲責.我一人肯代他人祈禱,他人又肯代他人祈禱,一而十,十而百,百而千,多人願負禱告的責,必有多人因此得救,多人肯代教會禱告必有許多教會因而奮興了.人的靈魂下獄是我的責,無靈由我引導而歸主,空手見父,是我們莫大的恥辱.以色列人不乘時機來榮耀主,至終大受主的責備,猶太國,地位居中,主甚願他廣傳救恩於萬國,但猶太人終如無花果樹之因無果而見詛,奪其恩而施之異邦.我們今日白白得受主的恩典,應該白白的施與別人,使多人得救,歸榮我們在天之父,阿們.
\newline
\newline


\newpage

\section{第1屆~港九培靈會~古約翰~第七講}
\label{sec:1328}
\textbf{要被聖靈充滿}
\newline
\newline
連結: \href{https://www.hkbibleconference.org/session-message/view/1328}{\texttt{https://www.hkbibleconference.org/session-message/view/1328}}
\newline
\newline
\hyperref[sec:1331]{< < < PREV SERMON < < <}
~
\hyperref[sec:toc]{[返目錄]}
~
\hyperref[sec:1319]{> > > NEXT SERMON > > >}
\newline
\newline
以弗所書 5:18-5:21
\newline
\begin{longtable}{cl}
\hline
\hline
章節 & 經文 (和合本修訂版)\\
\hline
5:18 & \begin{tabularx}{0.7\textwidth}{X} 不要醉酒，酒能使人放蕩；要被聖靈充滿。 \end{tabularx} \\ \\ \relax
5:19 & \begin{tabularx}{0.7\textwidth}{X} 要用詩篇、讚美詩、靈歌彼此對說，口唱心和地讚美主。 \end{tabularx} \\ \\ \relax
5:20 & \begin{tabularx}{0.7\textwidth}{X} 凡事要奉我們主耶穌基督的名常常感謝父神。 \end{tabularx} \\ \\ \relax
5:21 & \begin{tabularx}{0.7\textwidth}{X} 要存敬畏基督的心彼此順服。 \end{tabularx} \\ \\
[1ex]
\hline
\hline
\end{longtable}
在主歷一千九百年時,拳匪把他們的身心獻給魔鬼,向魔鬼禱告與獻祭.他們有信心,信魔鬼能護助他們,不為外國的炮所傷.他們既被魔鬼充滿,所作所為的盡是破壞的工作.由此可見魔鬼充滿人,乃叫人殺人.我們若被聖靈充滿,就大有能力去救人.被聖靈充滿的人,只有一目的,就是使他人得救.由此可見被聖靈充滿的重要,所以聖經囑咐人,要「被聖靈充滿.」
\newline
\newline
神願意以弗所教會被聖靈充滿,我信神也願意廣州教會也被聖靈充滿；這因為神願信徒都充滿聖靈.充滿聖靈,是信徒的本份.所以我們不當說我任意......這任意不得神的喜悅；因此我們只當任神的意旨成就在我們的身上.我們若抵擋聖靈,不任聖靈充滿,這就是罪.神應許賜聖靈給人,神所賜人的總是最好的.所以主在路加十一章十三節說:「你們雖然不好,尚且知道拿好東西給兒女,何況天父,豈不更將聖靈給求他的人麼?」然在神許多的應許中,沒有一樣應許像這應許說得如此清楚.今為父母的既把應許賜給子女,那為子女的就當快快的接納父母的賞賜.假若國王有所賞賜而國民卻不接受,這就是表示憎惡和輕視他的主了.而國王的賞賜,也許對於國民有害,惟神的賞賜則決不會有害他的兒女.我們讀以弗所五章十八節和路加十一章十三節,就知神應許賜聖靈給人,既是這麼的好,人若不接受,試想,這是何等的逆神呢?信徒為主血所買贖的人,本當接受聖靈；所以沒有一信徒當說我可任意醉酒,這因為醉酒就是得罪神.神不像那些予人難堪的人,他所命人做的人都必能做.若神以擔不起的重擔給人擔,那就難了；惟神所命的,盡是人所能行的.今神命神的兒女要被聖靈充滿,這事是人人所能的.神的兒女當未做工前,必須先充滿聖靈,主尚先充滿聖靈然後到曠野,所以他能得勝.各位,若你不充滿聖靈,則神決不使你到曠野與魔鬼交戰.主能行異跡,使多人榮耀神,也是全藉聖靈的能力.今神也要我們被聖靈充滿而榮歸他名.反是,則神名必因我們而受辱,茲將缺欠聖靈的弊害論列如下:
\newline
\newline
(一)信徒若缺欠聖靈,靈性就不能孕育.比方有兩夫婦生了一兒子,他們就以兒子為最寶,想為父母的人都是這樣的看重自己的兒子.到了十個月後,那兒子便會爬地了,那時他的父母便歡喜到沒有言語可以形容,你若持著萬元要換他們的兒子,他們必愛子而不愛銀；這因為他們在兒子身上的希望是超過萬元之上的.若是遲到十年再見此子還是爬地,不能行走,不能講話,只能吃奶,不能吃肉.今日在廣州有不發育的信徒,昔日在哥林多教會也有.所以保羅在哥林多前書三章一至三節說:「弟兄們,我從前對你說話,不能把你當作屬靈的,只得把你們當作屬肉體,在基督裡為嬰孩的.我是用奶餵你們,沒有用飯餵你們,因為從前你不能喫,就是如今還是不能.」惟是那些被聖靈充滿的,我必發育生長；且力上加力,一人可追千人,二人可追萬人.問廣州有幾許這樣的信徒,能使未信的人發生驚慌呢?
\newline
\newline
(二)信徒若缺欠聖靈,就不能進到成聖的地位.主已由死復活升天,但必再來此世.他離我們而去,我們就好像是寡婦；所以我們當仰望主再來.主願我們無瑕疵,到將來相見時不致羞愧.假若這裡有兄弟等候新娘到來,那新娘穿著污垢的衣服,身體數月也不沐浴,臭味觸鼻也要作嘔.試想這樣蒙不潔的女子,你樂意不樂意和他結婚.想我們若不被聖靈充滿,也像這被人厭惡的新娘.在巴西國有一例,就是王欲娶妻,就必在全國中求最美貌又最聖潔的處女.凡被選的眾女子,當未見王之前,照例先潔淨身體十二個月,六個月用沒藥油,六個月用香料和潔身之物,滿了日期,然後挨次進去見王.當時以斯帖要準備一年之久,才得被選為王后.你看,為王后尚且如此忍耐,何況我們為主耶穌的新婦,豈不當更聖潔麼?所以我們要等候主,就當被聖靈充滿,而進到成聖的地位.
\newline
\newline
(三)信徒若缺欠聖靈,就不能結果.主說:「不是你們揀選了我,是我揀選了你們,並且分派了你們,叫你們去結果子.」既是這樣,我們若只有葉而無果,這又如何能得主的喜悅呢?所以在主的園裡的,當結出果來.主已應許人能作較他更大的工作.然要作大工,那就必須先得勝靈的充滿.五旬節時,門徒一被聖靈充滿,一日就有三千人歸主,廣州有沒有這樣的景象呢?為什麼五旬節時能夠這樣,在今日的廣州又不能呢?啊!我們若謙卑在主前,主也必成就此大事於廣州.高麗國的信徒謙卑於主前,不過二月時間,便有五千人得救.各位能這樣的謙卑,主也必使用你救多人.望你們不要忘記了那不結果的無花果樹,那樹是預表以色列人,因他不結果,就要受主的咒詛.敢問廣州去年有多少得救的人呢?在座中也許有信主數十年還未引一人歸主.各位,你若被聖靈充滿,雖欲不結果,也不得了.
\newline
\newline
(四)信徒若缺欠聖靈,就必使神名受辱.有一日主對彼得說:「西門,西門,撒旦曾求得著你們,為要篩你們,像篩麥子一樣.」惟是彼得聞了卻自恃自己是磐石.但無幾時,此石就在祭司的院裡跌倒.彼得在院裡有機會作證而不證.在路加廿二章記著說:「他們在院子裡生了火,一同坐著；彼得也坐在他們中間.有一個使女看見彼得坐在火光裡,就定睛看他說,這個人素來也是同他一夥的.彼得卻不承認說,女子,我不認得他.過了不多時候,又有一個人看見他,說,你也是他們一黨的.彼得說,你這個人,我不是.約過了一點鐘,又有一個人極力的說,這個人實在是同他一夥的；因為他也是加利利人.彼得說,你這個人,我不曉得你說的是甚麼.唉,我們從上看來,就見此石且變為泥了!各位,彼得因為自恃而忽略了聖靈,但三次不認主,在那仇敵也得了話柄來譏主了,這因大門徒也不認他.請想我們,有沒有像彼得那樣的玷污主名.
\newline
\newline
(五)信徒若缺欠聖靈,就必任令多人沉淪.今神以聖靈賜給信徒,這是神給信徒最大的恩物.聖靈使用耶穌做成救贖的工夫；也使我們叫多人得救.想初期的信徒若不接受聖靈,則我們也遭受其害；今我們若不接受聖靈,則多人也遭受其害.如莫勒牧師,初為一平常的傳道人,有一禮拜會後,有二姊妹對莫勒說:「我們為你多多禱告,願聖靈充滿了你」.誰知莫勒聽了,很不悅意,以為自己已滿足作工的能力了.後幸他幡然覺悟,認識自己實在缺欠聖靈,從此就切切而求,約有四月之久.一日到銀行去,便在路上禱告,求神照著應許,使他聖靈充滿.在那時,主應允他的禱告,聖靈就滿注他,甚至求主減少,恐不能容.他一受了聖靈充滿,就與前大大的改變.在四十三年前,莫勒來到我所住的城開奮興會,雖使用最大的會堂,但座位也坐滿了,個個都受感流淚.莫勒當時宣佈說:「明晚會集,不許女界來,也不許男信徒來,若是要來,就必須帶未信的人同來.我有一醉酒的友人,我往請他赴會,初卻推辭,後見懇請,就與我同去.那晚所講的題目是「不要自欺,上帝是輕慢不得的；人種的是甚麼,收的也必是甚麼.」我聞了也十分驚懼.他所講的好像神審人的話一樣,人人也驚懼,我的友也戰慄；那晚便有三百人信主,我的友也在其中.後來我這位友人,當他臨死時寫一信給我,說:「我們在天上相見罷!」各位試想,莫勒若不接受聖靈充滿,只恃己力來作工,最多也不過得少數人信主而已,不能叫大多數人得福,他若不接受聖靈充滿,且神也必審判他,說:「我欲你被靈充滿,但你棄而不受,結果使多人沉淪.」你想,莫勒的罪是何等的大呢?又如賣主的猶大,他也有機會接受聖靈,但他不決志成就神旨,魔鬼就得使用他來賣主.我想當時釘主的猶大人落地獄時,必罵猶大說:「我們因你而落地獄,你若接受聖靈就不會賣他,你不賣他,我們就不會釘他.」你看,一不接受聖靈的猶大,不單害己,且害及多人,望各位接受聖靈充滿,不要失此良機.
\newline
\newline
(六)惟被聖靈充滿,就必盡心愛神愛人.所以被聖靈充滿的人,必為未信的人禱告；也必為不冷不熱的教會禱告.各位,你若被聖靈充滿,神必使你多結果子.在英國一千八百八十年時,有一青年牧師,他三十歲便死了,多人為他惋惜.然他雖年青而死,但他已引三萬人歸主.他平日所講的道,句句都是聖經,大有權能,使人聽了就必心裡發熱.又斐尼牧師曾說,有一患肺病的人,生有五子,但因為貧窮,所以異常困苦.後有一商人憐恤他,供給他家的所需,他甚感那商人的恩,見他還未得救,乃切切為他禱告,未幾,那商人悔改歸主.他又為三十多間教會的牧師禱告；他知一人往印度傳道,又為那人禱告.他把所代禱的人寫在日記內,及到他死後,在出殯日,他的妻就取他的日記給與斐尼牧師,斐尼就依著他代禱的教會往傳道,結果卒有一萬人信主.你看這樣的人,尚能使多人得救,何況你們呢.
\newline
\newline


\newpage

\section{第1屆~港九培靈會~吳子坤~開會詞}
\label{sec:1319}
\textbf{開會詞}
\newline
\newline
連結: \href{https://www.hkbibleconference.org/session-message/view/1319}{\texttt{https://www.hkbibleconference.org/session-message/view/1319}}
\newline
\newline
\hyperref[sec:1328]{< < < PREV SERMON < < <}
~
\hyperref[sec:toc]{[返目錄]}
~
\hyperref[sec:1351]{> > > NEXT SERMON > > >}
\newline
\newline
路加福音 9:18-9:36
\newline
\begin{longtable}{cl}
\hline
\hline
章節 & 經文 (和合本修訂版)\\
\hline
9:18 & \begin{tabularx}{0.7\textwidth}{X} 耶穌獨自禱告的時候，門徒也同他在那裡。耶穌問他們：「眾人說我是誰？」 \end{tabularx} \\ \\ \relax
9:19 & \begin{tabularx}{0.7\textwidth}{X} 他們回答：「是施洗的約翰；有人說是以利亞；還有人說是古時的一個先知又活了。」 \end{tabularx} \\ \\ \relax
9:20 & \begin{tabularx}{0.7\textwidth}{X} 耶穌問他們：「你們說我是誰？」彼得回答：「是神所立的基督。」 \end{tabularx} \\ \\ \relax
9:21 & \begin{tabularx}{0.7\textwidth}{X} 耶穌切切吩咐他們，命令他們不可把這事告訴任何人； \end{tabularx} \\ \\ \relax
9:22 & \begin{tabularx}{0.7\textwidth}{X} 又說：「人子必須受許多的苦，被長老、祭司長和文士棄絕，並且被殺，第三天復活。」 \end{tabularx} \\ \\ \relax
9:23 & \begin{tabularx}{0.7\textwidth}{X} 耶穌又對眾人說：「若有人要跟從我，就當捨己，天天背起自己的十字架來跟從我。 \end{tabularx} \\ \\ \relax
9:24 & \begin{tabularx}{0.7\textwidth}{X} 因為凡要救自己生命的，必喪失生命；凡為我喪失生命的，他必救自己的生命。 \end{tabularx} \\ \\ \relax
9:25 & \begin{tabularx}{0.7\textwidth}{X} 人就是賺得全世界，卻喪失了自己，或賠上自己，有甚麼益處呢？ \end{tabularx} \\ \\ \relax
9:26 & \begin{tabularx}{0.7\textwidth}{X} 凡把我和我的道當作可恥的，人子在自己的榮耀裡，和天父與聖天使的榮耀裡來臨的時候，也要把那人當作可恥的。 \end{tabularx} \\ \\ \relax
9:27 & \begin{tabularx}{0.7\textwidth}{X} 我實在告訴你們，站在這裡的，有人在沒經歷死亡以前，必定看見神的國。」 \end{tabularx} \\ \\ \relax
9:28 & \begin{tabularx}{0.7\textwidth}{X} 說了這些話以後約有八天，耶穌帶著彼得、約翰、雅各上山去禱告。 \end{tabularx} \\ \\ \relax
9:29 & \begin{tabularx}{0.7\textwidth}{X} 正禱告的時候，他的面貌改變了，衣服潔白放光。 \end{tabularx} \\ \\ \relax
9:30 & \begin{tabularx}{0.7\textwidth}{X} 忽然有摩西和以利亞兩個人同耶穌說話； \end{tabularx} \\ \\ \relax
9:31 & \begin{tabularx}{0.7\textwidth}{X} 他們在榮光裡顯現，談論耶穌去世的事，就是他在耶路撒冷將要完成的事。 \end{tabularx} \\ \\ \relax
9:32 & \begin{tabularx}{0.7\textwidth}{X} 彼得和他的同伴都打盹，但一清醒，就看見耶穌的榮光和與他一起站著的那兩個人。 \end{tabularx} \\ \\ \relax
9:33 & \begin{tabularx}{0.7\textwidth}{X} 二人正要和耶穌分離的時候，彼得對耶穌說：「老師，我們在這裡真好！我們來搭三座棚，一座為你，一座為摩西，一座為以利亞。」他卻不知道自己在說些甚麼。 \end{tabularx} \\ \\ \relax
9:34 & \begin{tabularx}{0.7\textwidth}{X} 說這些話的時候，有一朵雲彩來遮蓋他們；他們一進入雲彩就很懼怕。 \end{tabularx} \\ \\ \relax
9:35 & \begin{tabularx}{0.7\textwidth}{X} 有聲音從雲彩裡出來，說：「這是我的兒子，我所揀選的。你們要聽從他！」 \end{tabularx} \\ \\ \relax
9:36 & \begin{tabularx}{0.7\textwidth}{X} 聲音停止後，只見耶穌獨自一人。當那些日子，門徒保持沉默，不把所看見的事告訴任何人。 \end{tabularx} \\ \\
[1ex]
\hline
\hline
\end{longtable}
當今天培靈第二會開會之始,我回想既往,觀察現在,就有三件感謝主的事.追想去年第一會會集的情況,各赴會的兄姊,不因天氣酷熱而不前,卻能始終到會,這是我們要感謝主的第一件事.當第一會散會後,各同勞念念不忘第二會的會集,今雖當此不寧的環境中,竟得安然開會,這是我們要感謝主的第二件事.今當此天氣炎熱中,各同勞和各兄姊,暑也不避,同來此會,真是顧念靈性的事,過於肉體的事,這是我們要感謝主的第三件事.
\newline
\newline
其次在此開會之始,對于第二會會前籌備的經過,會期中的工作,及會後的結果的話,要藉此機會對各位略講.
\newline
\newline
(1)說到會前的籌備,我們因不敢靠人的力量來成就神的工作,所以屬人的工作,幾等於零,每會不過用數十分鐘會集,討論如何分配工作及聘請講員而已.然我們雖不靠人的力量,卻無依無靠的投靠神；所以從第一會散會到而今,我們除自己多多的禱告外,還要請求同心的兄姊代禱,感謝主,他感動了許多他的兒女,擔起這會禱告的責任.
\newline
\newline
(2)至於會期中的工作,不用說各位也知是為什麼了.今我國教會所需的就是靈性的培養所以本會的工作就注重這點.我想基督教入中華有一百廿餘年之久,所得的人不過三十餘萬.近年中華歸主的聲浪日日入我們的耳鼓,但各位試想中華究到何時才能歸主呢?從前用一百二十餘年,才得三十餘萬人歸主,照此推算,再經一千二百餘年,也不過得三百餘萬人；又再經一萬二千餘年,中華還未歸主呢!為什麼教會進步得如此遲滯呢?這實因信徒缺乏靈性的培養,故沒有靈力來引人歸主.今日教會對於工作的討論和組織,真是忙到不了.可是靈性的培養,就異常的缺乏,不然,中國教會,決不至如此衰落,如此沒有生產.這是我們同有的感覺,故本會為此而特殊注重靈性的培養.
\newline
\newline
(3)最後說到會終的結果如何,就人方面來說,責任固然在主講的人,但最大的還在赴會的兄姊.我們若要此會會集有良好的結果,那就要實行使徒的遺樣,「同心合意的恆切禱告」；(使一章十四節,這同心的禱告,就是好結果的因.主耶穌說:『若是你們中間有兩個人在地上,同心合意的求甚麼事,我在天上的父,必為他們成全.』(太十八章十九節)從這裡可見,主不貴人多少,只貴人同心；假若萬人禱告而萬心,這樣的禱告是沒有用的.我們雖三二人,若是契合一心托主名而求的,父神必為我們成全.所以本會的希望,就是希望赴會兄姊,契合一心而恆切禱告.求神的意旨,成就在我們的身上.
\newline
\newline


\newpage

\section{第1屆~港九培靈會~王載~第七講}
\label{sec:1351}
\textbf{寶貝的耶穌}
\newline
\newline
連結: \href{https://www.hkbibleconference.org/session-message/view/1351}{\texttt{https://www.hkbibleconference.org/session-message/view/1351}}
\newline
\newline
\hyperref[sec:1319]{< < < PREV SERMON < < <}
~
\hyperref[sec:toc]{[返目錄]}
~
\hyperref[sec:1352]{> > > NEXT SERMON > > >}
\newline
\newline
彼得前書 2:4-2:7
\newline
\begin{longtable}{cl}
\hline
\hline
章節 & 經文 (和合本修訂版)\\
\hline
2:4 & \begin{tabularx}{0.7\textwidth}{X} 要親近主，他是活石，雖然被人所丟棄，卻是神所揀選、所珍貴的。 \end{tabularx} \\ \\ \relax
2:5 & \begin{tabularx}{0.7\textwidth}{X} 你們作為活石，要被建造成屬靈的殿，成為聖潔的祭司，藉著耶穌基督獻上蒙神悅納的屬靈祭物。 \end{tabularx} \\ \\ \relax
2:6 & \begin{tabularx}{0.7\textwidth}{X} 因為經上說：「看哪，我把一塊石頭放在錫安—一塊蒙揀選、珍貴的房角石；信靠他的人必不蒙羞。」 \end{tabularx} \\ \\ \relax
2:7 & \begin{tabularx}{0.7\textwidth}{X} 所以，這石頭在你們信的人是珍貴的；在那不信的人卻有話說：「匠人所丟棄的石頭已作了房角的頭塊石頭。」 \end{tabularx} \\ \\
[1ex]
\hline
\hline
\end{longtable}
主耶穌說:「有人變賣了一切所有而買了一塊地；且他所賣去一塊更大的地,而買回一塊更小的地.為什麼呢?這因為他知道地雖小,但地裡卻藏有寶貝.各位,為什麼有人捨去一切而跟從耶穌呢?這也因為他知耶穌是寶貝.
\newline
\newline
然而,世人以什麼為寶貝呢?我見世人有以生命為寶貝.故一有侵害他生命的病來到,他就延請醫生診治,不惜錢財來買藥；為的,是使生命延長.說到世人以生命為寶貝一事,就是魔鬼也知道了.他在約伯書說:「以皮代皮,情願捨去一切所有的保全性命.」然而人的生命,不過數十年便了；惟耶穌是真的生命,是永遠的生命.他說:「復活在我,生命也在我.」今我告訴各位,耶穌是生命的主,你信耶穌,就得永生的福.想人生此世,不過數十年間,任是最長壽命的,終底也有一日離世；可是離世之後,又到何處去呢?惟我們信耶穌的人,就知死後是到一極樂的地,在那地沒有痛苦,沒有眼淚,也沒有魔鬼.試思,那地這麼的好,叫我們到那地的路是何等的好呢?我們若得了到那地的路,就必以那路為寶貝.主耶穌說:「我就是道路,真理,生命,人若不藉著我,沒有人能到父那裡去.」
\newline
\newline
其次,世人有以自由為寶貝,故有「不自由,毋寧死」之說.為奴的人,是最苦不過的人；惜人不知,犯罪則為罪奴；嗜賭則為賭奴；吸煙則為煙奴.然幸主耶穌降世,為釋放一切為奴的人.他說:「你們曉得真理,真理必叫你們得以自由.」許多不信主的人,他以為一信了主,便受束縛；誰知受主的束縛,才是真自由；不受,才是真束縛.啊!真理能釋放人一切的罪及迷信,我們說耶穌是寶貝,就因他是真理.
\newline
\newline
復次,世人有以糧食為寶貝,故有「民以食為天」這句話.人若餓了三天,那人就不得了.但主說:「不要為那能壞的食物勞力,要為那能存到永生的食物勞力.」又說:「我就是生命的餅,叫人喫了就不死.」然我們要喫這生命的餅麼?那就必須有信；這因為信主是沒有毒的,才會接受主.我從前聽說耶穌不可信,說一信了,眼就要被挖去.但而今卻有人問我,信了耶穌,是否要被人挖去心和眼?我就回答一聲「是」.那人聽了,便以為奇,我就再申釋說,不信耶穌的人的心是污穢的,但一信了耶穌,耶穌就把污穢的心挖去,換回一個純潔的新心；不信耶穌的人的眼,只想看戲和看淫書,但一信了耶穌,耶穌就把眼挖去,換回可愛的眼,如聖經說:「他的眼像鴿」；眼被改換了,便愛看聖經而不愛看淫書了.
\newline
\newline
復次,世人也有以水為寶貝.如遇早年,各地的人都為著水懷憂.惟主是生命的水,這水可白白取得.主耶穌說:「凡喝這水的,還要再喝；但是喝我們賜的水,就永遠不渴.」我們都知,渴是最苦的事,當在渴時,雖金也不以為寶了.所以人當渴時,要水而不要金.你看那些傷兵,所呼求的只是水.各位,我們最苦的是什麼呢?就是靈的渴.幸主是生命的水,他足以止我們的渴；所以耶穌是寶貝.
\newline
\newline
我願在此數日中,得勝靈感動我們的心,叫那些已信主的人,以主比萬事萬物為更寶貝；那些未信主的人就認識耶穌是寶貝而歸信他.彼得書只說,主在信的人中為寶貝.各位,我今不問你們信了主未,只問你們覺得主是不是寶貝?猶太人釘生命的主而要求釋放盜賊,這就因他們不以主為寶貝.猶大賣主,這也因他不以主為寶貝.惟保羅為掃羅時,就任意逼迫門徒,後因主一向他顯示,他就對主說:「主啊,我當作什麼?」後來他在腓立比書有說:「只是先前於我有益的,我現在因基督都當作有損的.不但如此,我也將萬事當作有損的,因為我以認識我主基督耶穌為至寶.我為他已經失去萬事,看作糞土,為要得著基督」.各位試思,當保羅未信主前,就以主為最污賤；及到信主後,就以主為至寶.從這裡可見,信能叫人的眼睛開亮,認識耶穌為至寶,所以你的信越大,就見主越寶；反是則見主越小.這不是主有所改變,不過人的信心改變而已.今我們當求主加大我們的信心,和有屬天之智.我們若有了,就必知所揀選.比方:把珍珠和餅來任兒童選擇,他就必捨珍珠而取餅；惟年大有智的人,就必以珍珠為寶.我願神賜你們有屬天之智而選擇耶穌；因為耶穌是我們永遠可愛的至寶.想人當有病時,就不以平昔所寶貝的為寶.如食物在病時,那就不為人所愛,如平時一樣了.主不是這樣的主,在我們身體康健時,固是寶,就在病痛時,也是寶；在我們生時固以主是寶,在死時就更以主是寶了.可惜世人,只知珠寶為寶；但是死亡到了,那些珍寶又有什麼用呢?在那時,不獨不能叫人得安慰,且叫人多流點淚.從前有一富厚的太太,她用一鐵箱藏了許多的珍寶.有一次她病重了,見醫生對家人耳語,又見家人神色蒼忙,她於是明白已病沒有望了.她便大聲痛哭,叫人取鐵箱來,抱著大哭；及到臨危時,兩手還緊握著那鐵箱.各位,你看那為平時所寶的,到死時就不寶了.想人在世得了寶,就會有賊來搶劫；所以刻刻小心防守.惟主耶穌的寶,就不必什襲藏之；且當多多傳開,使人同得這寶.世人要得珍寶,就必須出納代價.我前見一姊妹藏一鑽石戒指,價值一千元；我說,此寶我得不到,惟主耶穌,只用信心的手便可取得.世上的寶,今天得了,明天也許失去；惟得主的寶,就必永遠得著.從前滿清的皇帝死了,便用許多珍寶同葬,但前數月新聞紙宣傳,說許多皇陵已被盜掘了.從這裡看來,屬世的珍寶,生不能攜來,死不能帶去；惟主的寶,是永久不失的.我想世人彼此相見時就好,惟同居久了,就不好了,這因為見他的毛病了.惟主就不是這樣,初見是寶貝,久見則更覺寶貝.我信了主十年,越久越覺主是寶貝.我知在座有信主多年的人,必同有這樣的經驗.耶利米書九章廿三至廿四節說:「智慧人不要因他的智慧誇口,勇士不要因他的勇力誇口,財主不要因他的財物誇口；誇口的卻因他有聰明,認識我是耶和華.」這裡神就囑咐人,只當以認識他而誇口.各位,你若識盡世上的學問而未識主,那末,你仍是愚人；得著全世界的貨財而未得得主,仍是赤貧的人.今願各位就在今晚接納耶穌,因他是你的寶貝,是你的智慧.從前有一博士渡河,他問船夫說:「你識英文否?」答說,「不識.」又問:「識算學否?」答說,「只知二加三為五」.又問:「曉天文地理否?」答說,「這都不識.」博士聽還未了,便說,你真是愚人了.無幾,船便開到河中,忽風浪大作,舟便沉下,博士且沉且呼,那時船夫便發聲說:「博士先生,你識浮水否?」博士窘迫萬狀,不得已而向船夫求救.各位,將來有一時候,我們都要過那死河,惟有信主的人才能得救.今願各位接主入心,主是至寶,此寶並不是可見而不可得的寶.若是可見而不可得,那就與我們沒有關係.有時晚上,我與我的兒子閒遊,他就向我討月亮.各位,月亮是可見而不可得的；惟主不是這樣,主說,道離人不遠,就在人的心裡.各位,你若肯謙卑到主前,求主入你的心裡,就必得主而歡然回家了.
\newline
\newline


\newpage

\section{第1屆~港九培靈會~王載~第八講}
\label{sec:1352}
\textbf{「我來了」}
\newline
\newline
連結: \href{https://www.hkbibleconference.org/session-message/view/1352}{\texttt{https://www.hkbibleconference.org/session-message/view/1352}}
\newline
\newline
\hyperref[sec:1351]{< < < PREV SERMON < < <}
~
\hyperref[sec:toc]{[返目錄]}
~
\hyperref[sec:1353]{> > > NEXT SERMON > > >}
\newline
\newline
希伯來書 10:4-10:7
\newline
\begin{longtable}{cl}
\hline
\hline
章節 & 經文 (和合本修訂版)\\
\hline
10:4 & \begin{tabularx}{0.7\textwidth}{X} 因為公牛和山羊的血不能除罪。 \end{tabularx} \\ \\ \relax
10:5 & \begin{tabularx}{0.7\textwidth}{X} 所以，基督到世上來的時候，就說：「祭物和禮物不是你所要的，但你曾給我預備了身體。 \end{tabularx} \\ \\ \relax
10:6 & \begin{tabularx}{0.7\textwidth}{X} 燔祭和贖罪祭是你不喜歡的。 \end{tabularx} \\ \\ \relax
10:7 & \begin{tabularx}{0.7\textwidth}{X} 那時我說：看哪！我來了，我的事在經卷上已經記載了；神啊！我來為要照你的旨意行。」 \end{tabularx} \\ \\
[1ex]
\hline
\hline
\end{longtable}
「我來了」的我,就是主:這「我來了」的話是主對神講的.這三字若在常人口中說出來,就沒有什麼意義；惟出自主口,那就大不同了,我想凡是真信主的人,讀到這話,定然有感於心.有一天我讀此經,心受很大的感動；我覺神的愛,在這「我來了」這一事為更大.主對神說這句「我來了」的話,已聽過了一千九百二十八年了.想世上來了許多許多的人,但惟此我來了才與我們有鉅大的關係；所以世界用他降生之年為紀元.然主說這話,不是易的；因他早已看透世界的苦,就是猶太人用石擊他,猶大賣他,最終還以十字架待他,他也見了；他雖見了,但他還要說「我來了」；這「我來了」的聲音,為最慈悲的聲音,他來是為拯救我們；他捨天堂的榮耀而降到這奸惡的世界；他本來富足,為我們而成了貧窮；他本是神,今竟成為人子,所以他說我來了這句話是要付頂大的代價才能說的.主既為我們納了頂大的代價,所以我們可因信而得救.各位,此救恩是白白可得的；但我們要知救恩的成就不是容易的,是要主降世代死才能成就.從前大衛與非利士人交戰,那時大衛在山寨,非利士人的防營在伯利恆.大衛渴想說,甚願有人將伯利恆城門旁井裡的水,打來給我喝.那時有三個勇士聞了這話,就闖過非利士人的營盤,從伯利恆城門旁,井裡打水,拿來奉給大衛；他卻不肯喝,將水奠在耶和華面前說,我的神阿,這三個人冒死去打水；這水好像他們的血一般,我斷不敢喝.各位,此水在平時就沒有什麼價值,惟是三勇士冒死去取來的,這就有價值了.主說:「口渴的人也當來,願意的都可以白白取生命的水喝」各位,你渴麼?但你可到主這裡來,取生命的水喝；這水是主白白的給人喝的.不過喝這水人要知,這水是主用血換來的.
\newline
\newline
茲把(一)主怎樣來,(二)為什麼來,(三)為何人來,(四)為人作了什麼事諸問題論列如下:
\newline
\newline
(一)
\newline
\newline
(二)主為什麼來世.
\newline
\newline
(1)照神的意旨行事.　　前亞當因悖逆而背神命,於是因一人悖逆,眾人就成為罪人；耶穌為後亞當,他來就為照神的意旨行事,於是,因一人的順從,眾人也成為義了.
\newline
\newline
(2)應驗舊約預言.　　主若不來,那就舊約的預言,豈不是白寫的麼?主說:「你查考聖經,給我作見證的,就是這經.」這樣主若果真不來,那就詩篇廿三篇,以賽亞五十三章,真是白寫了.但主已來了,舊約的預言也一一應驗了.
\newline
\newline
(3)作更美好的祭品.　　考猶太人犯了罪,就牽牛或羊到神前,按手祭牲之上,把他宰殺了,以表明代死.但聖經說:「公牛和山羊的血,斷不能除罪,」所以主要到世上來,作贖罪的祭物.施洗約翰為主作證說:「看哪,這是上帝的羔羊,除去世人罪孽的」.今主釘十字架而流血,就是為人贖罪.我們要感謝主,主若不道成肉身,就不能釘十字架而流血,主若不是具有血肉的身體,羅馬的兵也不能用鎗刺他的.
\newline
\newline
(三)主為何人來.
\newline
\newline
(1)恩典與真理.
\newline
\newline
(2)生命.
\newline
\newline
(3)光.
\newline
\newline
(4)主來為人作了什麼事.
\newline
\newline
今惟願各位立刻信靠主,因他是可信靠的主；凡親就他的人,都不會失望而返.聽說英國人以得與英王談話為無上的榮耀.從前有一英人對人說:「今日英王和我談話」,言下有得色.後來人知道英王對他所說的話,就是「你快的滾開」這一句叱他速退的話.惟主雖為萬王之王,但他總不丟棄凡來親就他的人.各位,你若未得主,那就不是主丟棄你們；乃是你們丟棄主.主來世,就是為接人,惜人拒絕他；各位而今你接待他呢,還是拒絕他呢?這事關乎你的永生或是永死,願你開心接主；若主不在你心,則主與你沒有有什麼關係；他來了一千九百二十八年了,但他入你的心,有多少年呢?
\newline
\newline


\newpage

\section{第1屆~港九培靈會~王載~第九講}
\label{sec:1353}
\textbf{猶疑}
\newline
\newline
連結: \href{https://www.hkbibleconference.org/session-message/view/1353}{\texttt{https://www.hkbibleconference.org/session-message/view/1353}}
\newline
\newline
\hyperref[sec:1352]{< < < PREV SERMON < < <}
~
\hyperref[sec:toc]{[返目錄]}
~
\hyperref[sec:1354]{> > > NEXT SERMON > > >}
\newline
\newline
約翰福音 10:22-10:30
\newline
\begin{longtable}{cl}
\hline
\hline
章節 & 經文 (和合本修訂版)\\
\hline
10:22 & \begin{tabularx}{0.7\textwidth}{X} 那時正是冬天，在耶路撒冷有獻殿節。 \end{tabularx} \\ \\ \relax
10:23 & \begin{tabularx}{0.7\textwidth}{X} 耶穌在聖殿裡的所羅門廊下行走。 \end{tabularx} \\ \\ \relax
10:24 & \begin{tabularx}{0.7\textwidth}{X} 猶太人圍著他，對他說：「你讓我們猶豫不定到幾時呢？你若是基督，就明白地告訴我們。」 \end{tabularx} \\ \\ \relax
10:25 & \begin{tabularx}{0.7\textwidth}{X} 耶穌回答他們：「我已經告訴你們，你們卻不信。我奉我父的名所行的事可以為我作見證。 \end{tabularx} \\ \\ \relax
10:26 & \begin{tabularx}{0.7\textwidth}{X} 但是你們不信，因為你們不是我的羊。 \end{tabularx} \\ \\ \relax
10:27 & \begin{tabularx}{0.7\textwidth}{X} 我的羊聽我的聲音，我認識牠們，牠們也跟從我。 \end{tabularx} \\ \\ \relax
10:28 & \begin{tabularx}{0.7\textwidth}{X} 並且，我賜給他們永生；他們永不滅亡，誰也不能從我手裡把他們奪去。 \end{tabularx} \\ \\ \relax
10:29 & \begin{tabularx}{0.7\textwidth}{X} 我父所賜給我的比萬有都大，誰也不能從我父手裡把他們奪去。 \end{tabularx} \\ \\ \relax
10:30 & \begin{tabularx}{0.7\textwidth}{X} 我與父原為一。」 \end{tabularx} \\ \\
[1ex]
\hline
\hline
\end{longtable}
許多人都是這樣說:「我要親眼見著耶穌我才信耶穌,」從前我也是這樣的人,各位,信不在乎眼見,你若不信,就是親眼見了,你也不信.就如約翰十章廿四節所載的猶太人,他們親眼見耶穌,親耳聽耶穌,就該信了；可是他們也不信.他們圍著耶穌說:「你叫我們猶疑不定到幾時呢?」你想,他們自己猶疑而不信耶穌,反說主叫他們猶疑,這真是世人同有的口氣.有些人自己跌倒了,反怪責別人；自己生氣了,反說人惹他；自己不信主,就推某人或某事阻他；這些一切或還可勉強講去,惟說耶穌叫他不信,那就千講萬講也講不去了.唉!世人不信,就怪到上帝,那是人的一種通病.如亞當犯了罪,就怪神說,你所賜給我與我同居的女人,她把那樹上的果子給我,我就喫了.又如猶太人不信,就反怪耶穌叫他們猶疑.然聖經說耶穌是信心創始成終的主,他來是叫人信,不是叫人猶疑；所以猶太人所說叫他們猶疑的話,是無理的話.試想,人若對飯說,你叫我餓到幾時呢?對水說,你叫我渴到幾時呢?對光說,你叫我不見到幾時呢?這可以麼?唉!不是光叫人不見,飯叫人飢餓,水叫人口渴；乃是因人不信!所以耶穌回答他說:「我已經告訴你們,只是你們不信!」在這裡可見人的猶疑,不是從耶穌那裡來的.既是這樣,猶疑究從哪裡來的呢?這是我們同有的問話.茲把那三個猶疑的源頭,略論如下:
\newline
\newline
(1)魔鬼.
\newline
\newline
(2)人智.
\newline
\newline
(3)異端.
\newline
\newline
我們從上看來,就知猶疑的源頭就是魔鬼,人智,和異端.一日,我讀到這節經(約翰十章廿四節)心甚受感,我想主的言詞中,沒「若是」這樣猶疑不定的言詞,所告訴人的,都是實實在在的.主初在曠野被試探時,魔鬼便用你「若是」這樣的疑問詞來試探主；後主被釘時,眾人譏誚主說,你「若是」神的兒子,就從十字架上下來罷.唉!他們何等蒙昧,卻不知他真是神的兒子,人的救主.
\newline
\newline
各位,若耶穌的言語不可信,那末,世上就沒有一個人的言語是可信的了.耶穌說,所作的事,可為我作見證；我們若虛心查考主所作的事,就必深知耶穌真為我們的救主了,試思,有什麼人能夠像主?當他降生時,有天使為他唱歌說:「在至高之處,榮耀歸與上帝；在地上平安,喜悅歸與人.」降生後,就有三博士從遠道來拜他.再看他在世只三十三年,三十歲前前為木匠,他未曾進過大學；三十歲後只傳道三年,週遊行善,他既不是將軍,又不是帝王,也不是博士或財主；但他卻對門徒說:「你們往普天下傳福音給萬民聽,」又說:「天地要廢去,我的話卻不能廢去」；他說這些話的時候,還未有新聞紙,或留聲機；然他竟說傳往普天下,你看,他所說的話,而今是否傳遍?若當時我們聽到他說這些話,必定不當他是神子,就必定當他是患了神經病而大言不慚的人了.各位,我們當確知,耶穌確是神子.聖經已明證耶穌為神子,這是我所當信的；這因為世上沒有一本書能比聖經為更可信可靠的.當耶穌未來世前,舊約便已預言；隨後則應驗不爽.他做了許多奇事,這就可證明他為人的救主.從他所言和所做的事,就知世上沒有一人能及他.他在十字架被釘時,他不發怒,也不怨天尤人,反求神赦免殺他的人的罪.還有一頂大的證據,就是他的復活.羅馬書一章裡說:「他從死裡復活,以大能顯明是神的兒子.」主已復活了,所以主的墳墓是空的.他今在神的右邊為祭司.我深知他是活的,因他常允我的禱告,及指示我的一切.在聖經中有許多憑據,證明耶穌是神的兒子,人的救主；惜乎人竟不信他!
\newline
\newline
然而,不信耶穌的到底有什麼的結果呢?這是我們所當知的事.在約翰八章廿四節載著主的警告說:「你們若不信我是基督,必要死在罪中.」各位當記著,這「你若不信耶穌為救主,則必要死在罪中」的警告,是主對人最大的警告.想我若不信了耶穌,那就而今不知到了哪裡去了.我感謝他,他救我出罪中,所以我要永久的信他靠他.
\newline
\newline
說到不信,那就雖是熱心聽道也是無用.所以希伯來書說:「只是所聽的道與他們無益,因為他們沒有信心與所聽的道調和.」有人對我說,你所講的道真好,但他不信；這樣,這道與他無益.昨有人請我喫晚飯,假若我對不喫,只說真好真好,你想我是何等的愚人呢?同樣,聞道而不信道的人,實在也是一個可憐的愚人!然我們接聲則用耳,接光則用目,接飯則用口；照樣,我們接道則當用信心.各位!你若不信主為你的救主,則你必要死在罪中；你若聞道而不信道,則你的心必更冰冷.
\newline
\newline
唉,世人對於一切假話,就快快的信；惟叫他信真的話,那就千難萬難了!如在創世記裡記著說:雅各有十一個兒子,少子是約瑟,他愛約瑟,過於愛其他的兒子；因此約瑟的哥哥就妒他恨他,終底把他賣給了以實瑪利人?他的哥哥既然把他賣了,後就宰了一隻公山羊,把約瑟的那件彩衣染了血,就打發人送到他們的父親那裡說,請檢了這個,請認一認,是你兒子的外衣不是.他認得,就說,這是我兒子的外衣,有惡獸把他喫了,約瑟被撕碎了,撕碎了!你想,這豈不是一件假的事麼?但是雅各一見便信了,卻不猶疑,也不派人調查.事後過了五十年,他的兒子們又告訴雅各說:「約瑟還在,並且作埃及全地的宰相.」這是真的事,但是他卻不信了.想世人也是這樣,你告訴人說,人是神所創造的,人若信靠耶穌,就必得救,這是真確的事；可是他們不信!惟你若告訴人說,人是從猿猴進化來的,這是假的事；惟是他們就連連稱「是」.你若告訴人說,信耶穌的死後上天堂,他們就不信；若說人死好像狗一樣,他就信了.
\newline
\newline
今我願你們除去猶疑,信靠耶穌為你們的救主；不要說我要見了才信.比方我是銀行的行長,請你為司庫,但我不會對你說,我見你之時才信你,不見就不信了；若是這樣,這就表明我不信你了.
\newline
\newline
說到這裡,要講講信耶穌的好處.信耶穌的好處是什麼呢?就是罪惡得以赦免,並且得與神和好,前是與神為仇敵,今靠主血而得親近神；前為魔鬼的奴僕,今得為神的兒女；前在黑暗的勢力下,今得在光明的國裡；前在等候沉淪的地位,今因信主而得永生.什麼是永生呢?千年算得永生否?算不得!千萬年算得永生否?算不得!真的,就是萬萬年也算不得是永生的年限,因為永生是悠悠無限的.這永生真是寶貝!願你們都想念永生,而知永生是何等的寶貝；更願你們因信主而得永生.
\newline
\newline
各位,我們要看重神所賜的恩典,我為此而講了又講；為的,是深願你們得此最大的永生福分.今你們一決心歸信耶穌,就立刻脫離魔鬼勢力和轄制,而得為神的兒女,此福大得無比!此事要得無比!主耶穌說:「現在你們信麼?」這「現在」信主的問題,為你們目前最要的事.惜有人因家人未信,所以他也不信.回想我信道時,父母還未信.但我們不要等家人信了,自己然後信；當從我先信起.信主好像逃火,當然快快的逃.比方家裡起火,當然不會遲遲疑疑的呼集家人,排起隊來,然後慢慢的走.各位,當我們請客時,可以講客氣,但若走火,就用不著客氣了.信耶穌就是逃避地獄的火,當從你先信,立著得救的地位,然後引導父母和家人；這樣,你和你一家都可得救了!
\newline
\newline


\newpage

\section{第1屆~港九培靈會~王載~第十講}
\label{sec:1354}
\textbf{五種愚人}
\newline
\newline
連結: \href{https://www.hkbibleconference.org/session-message/view/1354}{\texttt{https://www.hkbibleconference.org/session-message/view/1354}}
\newline
\newline
\hyperref[sec:1353]{< < < PREV SERMON < < <}
~
\hyperref[sec:toc]{[返目錄]}
~
\hyperref[sec:1345]{> > > NEXT SERMON > > >}
\newline
\newline
箴言 9:6-9:6
\newline
\begin{longtable}{cl}
\hline
\hline
章節 & 經文 (和合本修訂版)\\
\hline
9:6 & \begin{tabularx}{0.7\textwidth}{X} 你們要離棄愚蒙，就得存活，並要走明智的道路。」 \end{tabularx} \\ \\
[1ex]
\hline
\hline
\end{longtable}
今試問各位,在你們中間,有有誰甘願做個愚人.我知你們中,定然沒有人甘願做；不特你們不願做,設有人說你們是愚人,你們定必怫然不悅了.這不特你們自己不願做愚人,若是你們中為丈夫的娶了愚的妻子,為妻子的嫁了愚的丈夫,你們也是不悅意的啊,真的,人人都不願意做愚人,可是許多人卻走在愚人的路上,反自以為智；這真是愚人中的愚人了.今特提出聖經所說那五種愚人,來向各位略講；為的,是願各位以之為鏡,而知所為人,茲分述如下:
\newline
\newline
(一)心裡說沒有神的人
\newline
\newline
(二)以行惡為戲耎的人
\newline
\newline
(三)只知為自己攢財的人
\newline
\newline
(四)聞道而不行的人
\newline
\newline
(五)不儆醒預備的人
\newline
\newline
在上所講的五種愚人,第一,就是那心裡說沒有神的人；第二,就是那以犯罪為戲耎的人；第三,就是那只知為自己積財的人；第四,就是那聽道而不行道的人；第五,就是不儆醒預備的人.各位你是不是愚人?若是,那末,在這五種愚人中,是哪一種的愚人?聖經說:「你們愚蒙人,要捨棄愚蒙,就得存活；並要走光明的道.」各位要知,神遣耶穌降世,使愚蒙的人成為智者；人若接受耶穌,他就成為一有智的人.所以聖經說:「敬畏耶和華,是智慧的開端；認識至聖者,便是聰明.」末了,惟願你們有智慧,能明白這事,肯思念你們的結局.」
\newline
\newline


\newpage

\section{第1屆~港九培靈會~祈理平~第一講}
\label{sec:1345}
\textbf{主釘十字架與復活}
\newline
\newline
連結: \href{https://www.hkbibleconference.org/session-message/view/1345}{\texttt{https://www.hkbibleconference.org/session-message/view/1345}}
\newline
\newline
\hyperref[sec:1354]{< < < PREV SERMON < < <}
~
\hyperref[sec:toc]{[返目錄]}
~
\hyperref[sec:1346]{> > > NEXT SERMON > > >}
\newline
\newline
使徒行傳 2:30-2:36
\newline
\begin{longtable}{cl}
\hline
\hline
章節 & 經文 (和合本修訂版)\\
\hline
2:30 & \begin{tabularx}{0.7\textwidth}{X} 既然大衛是先知，他知道神曾向他起誓，要從他的後裔中立一位坐在他的寶座上。 \end{tabularx} \\ \\ \relax
2:31 & \begin{tabularx}{0.7\textwidth}{X} 他預先看見了，就講論基督的復活，說：『他不被撇在陰間；他的肉身也不見朽壞。』 \end{tabularx} \\ \\ \relax
2:32 & \begin{tabularx}{0.7\textwidth}{X} 這耶穌，神已經使他復活了，我們都是這事的見證人。 \end{tabularx} \\ \\ \relax
2:33 & \begin{tabularx}{0.7\textwidth}{X} 他既被高舉在神的右邊，又從父受了所應許的聖靈，就把你們所看見所聽見的，澆灌下來。 \end{tabularx} \\ \\ \relax
2:34 & \begin{tabularx}{0.7\textwidth}{X} 大衛並沒有升到天上，但他自己說：『主對我主說：你坐在我的右邊， \end{tabularx} \\ \\ \relax
2:35 & \begin{tabularx}{0.7\textwidth}{X} 等我使你的仇敵作你的腳凳。』 \end{tabularx} \\ \\ \relax
2:36 & \begin{tabularx}{0.7\textwidth}{X} 故此，以色列全家當確實知道，你們釘在十字架上的這位耶穌，神已經立他為主，為基督了。」 \end{tabularx} \\ \\
[1ex]
\hline
\hline
\end{longtable}
今日所講的題目,就是見證人的見證大網.說到見證的大綱最要的有三,在使徒行傳二章卅至卅六節其中一段,就是見證大綱之一；也是傳道宣言之大綱,惟此大綱不是人人所喜聽的.雖然不為人所喜聽,但我們是神所選立的見證人,那就不當求悅於人,只要求悅於神；若得見悅於神,其餘就沒有問題了.
\newline
\newline
見證的大綱是什麼呢?就是主釘十字架與復活.今時與昔日有什麼分別呢?我信是大同小異.若是說有差別,這差別就在今日的人心較昔日的人心離神更遠.也許有人說,這流血贖罪的道,太不合時宜了.但我總不明那些人為什麼會說此道是不合時宜的道；也許他以為今時是科學昌明的世代,所以要勸我們順潮流,不要說主釘十字架與復活之道.各位,這樣的話,我實在不敢苟同.而今在美國也有些牧師不敢說此道；他們以為今日傳道要變更計劃,所以他們只用十五分鐘講道,講道後便開跳舞會；你看他們順潮流是順得何等十分呀!想他們原不是這樣的；初時也好像羅得一樣,初向所多瑪去,後便漸漸入了所多瑪城.回想我初來中國學話時,不明白漸漸這二字是什麼意義,後來才知一步進一步就是漸漸.而今主的見人順潮流,也是漸漸.撒旦用這個漸漸才能成就他的詭計；否則人就不上他的當了.
\newline
\newline
從前有一神學生問他的教授說:「主血實有赦罪的權否?」教授不答他,只在粉板上寫一廢字,他的意思以為主血是廢的.唉,這樣言論,隨處都有,我們總要小心,今日我願各位認定主流血贖罪之道,就是我們見證人的見證大綱.
\newline
\newline
或問主流血贖罪的救法,是從何而來的呢?這話我可直捷的回答說,此不是從人腦來,也不是從猶太人來,而是從神來的.我們若讀創世記三章,就知人的始祖在樂園犯了罪後,神就宰一羊,以羊的皮作衣服而蓋他們的羞.又讀創世記四章,就知有二兄弟獻祭,一草地的出產為供物獻給神,神不看中他的供；一宰羊而以羊的脂膏獻給神,神便看中.為什麼?這就是因一照神所設的法而獻,所以神就悅納.今我們既知主流血贖罪的救法,就是神所設立的救法,我們就當知所應取的態度了.我想對此道大約有三種人:一是遺忘了此道；二是故作不知此道；三是藐視此道.各位試想這樣的人,又怎能得到神的祝福呢?前有一鄉婦在醫院裡留醫,因而聞道信主,但惜醫生總不能醫好他:有一日他問醫生說:「我尚有幾許生命呢?」醫生回答他說:「尚有六個月」；他聽了便想返鄉,同房的人就勸阻他說:「你若回鄉,就只有三個月命」,他聽了就回答那同房的人說:「我寧願生三月而傳主流血贖罪之道；耶穌為我一生,難道我給回三月與他就不得麼.」今在座有這麼多傳道的人,我想問你,一年中有多少次講及此道.
\newline
\newline
我問為什麼要主流血,人才得救贖呢?原來這就是因人犯了罪.人既犯了罪,神不設別法救人,只要主被掛在木頭,流他的血；人靠著他的血才能得救.為什麼人要靠著主血才能得救呢?這問題我以為只可等到與神面對面時問神.照聖經所說,我們就知主若不流血,神就不以恩典給人；神向人所施的恩典,乃由主出.我感謝主,我因主的憐恤而得拯救,我若沒有了神的恩,而今我不知在哪裡.主若不流他的血,不把神的恩給人,那就敗壞的人,不知成了怎樣更敗壞的人.無怪印度有一烈士,曾有人請他演說,他只說「我今日得為這樣的人,全靠著耶穌」這一句話,說了便坐下.真的,我們敗壞的人,若要成為新人,不靠耶穌,還有誰可靠呢?
\newline
\newline
各位,你們若注重說主血的道,就必得人來聽；若說詼諧,那就你們真不及那些講古的人.我曾見一位傳道人講諧談,聽了叫我幾乎要嘔；惟是我也曾見一講古的人說諧談,和那傳道人比較起來,勝過萬萬.若是叫我講諧談,那就無暇了,若講主血的道,那就日夜也講,這因為主血大有能力.就我來說,我初也是魔鬼家裡的人,嗜好煙酒和各樣的惡.忽一日受感,我就獻身給主,當時我望天,也覺天為之新；其實不是天新,是祁理平新.感謝主,當我入堂時和出堂時前後真是兩人.
\newline
\newline
從前有一人曾對我說:「不知主血能不能潔淨我的罪.」說這話的人,家裡有九個偶像,他為巫有五十餘年.後來他歸了主,他又說:「我尋真道五十餘年也尋不到,而今才知主血的道」.他嗜吸雅片煙,當他煙癮發作時就入禱告室禱告,有一日之久,大呼主血救我,他就在這次懇切的禱告中得主釋放.又有一在北方的同事對我說:「在他那裡有一有三百萬家財的吸煙富人,因他不理生產,後來他的家財不知怎樣的竟散盡了,他處困苦中便把兒子買去,後來連妻子也買了,一切都賣了,不過只他一副不值錢的骨頭還未買去.他既到了山窮水盡的時候便投海自盡,後竟遇救,有一牧師見了他,便對他說:「耶穌能釋放你」,他聽了便接受了血的道,後三月便信主,不數月便重了數十斤.由此可見,主血的道是何等的有能力,惜許多傳道人只對人說當擇善,當行善,卻不對人說主流血救贖的道.
\newline
\newline


\newpage

\section{第1屆~港九培靈會~祈理平~第二講}
\label{sec:1346}
\textbf{靈洗}
\newline
\newline
連結: \href{https://www.hkbibleconference.org/session-message/view/1346}{\texttt{https://www.hkbibleconference.org/session-message/view/1346}}
\newline
\newline
\hyperref[sec:1345]{< < < PREV SERMON < < <}
~
\hyperref[sec:toc]{[返目錄]}
~
\hyperref[sec:1347]{> > > NEXT SERMON > > >}
\newline
\newline
昨講主的十字架,為見證之第一大綱,此為聖彼得所言；今所講的題目,是聖靈的施洗,此是第二條的大綱.
\newline
\newline
今所講的是彼得對於教會當受聖靈施洗的主張,未曉我們是否贊成.我信我們若是贊成第一條的；亦當贊成第二條的.我以為若不接納第二條的,那就第一條也可不必接納了.我正間禱告說:「求主在此數日間,復興主的工人」,主若允我們的禱告而復興我們,不知我們是否願意.
\newline
\newline
猶太人在幾千年的長久時期中,說望主臨,可是主降臨了,因為不合他心所接的,於是耶穌等了數千年猶太人,態度又變了.為甚麼他們等主這麼的久長,當主來了,他們終竟把主殺了,這真是奇事.而今我們求主,但我們不要像猶太人一樣,當日若對猶太人說,你必殺你所望的基督,他們必不信的,各位既然求主復興,望各位在此數日間,納復興當納的代價.
\newline
\newline
有一信者,求神叫他謙卑,神於是叫許多人攻擊他,他不明神為什麼這樣做,然神叫人攻擊他,就是要他納謙卑的代價.我怕會中有攻擊主,好像當時的人；查當時攻擊主的人,不是羅馬人,卻是猶太人.
\newline
\newline
昨晚古約翰牧師講道後,主的靈感動人心,當堂有人認罪,於是有些人,以為怪事,誰知這一極尋常的事.彼得在(徒二章三十九節)說:「因為這應許是給你你們,和你們的兒女,並一切在遠方的人,就是主我們的上帝所召來的.」這是什麼的應許呢?就是「必領受所賜的聖靈.」這應許不單是屬猶太人,也屬是在座的你我.在啟示錄說:「凡有耳......」這耳是指心；這就是說:「凡有心聽的就當聽.但惜沒多人沒有心聽,若說其他屬世的事,那就聽得津津有味.他們以為聖靈充滿是猶太人的事,卻不知從這「凡有耳......」的凡字,就可是關乎人人的事.而今我問你們曾否受了聖靈的施洗?或以為這樣的問是得罪了你們,其實這並不算得罪.當日保羅曾問施洗禮約翰的門徒說:「你們曾否受了聖靈.」各位,保羅這問,難道是得罪當日的信徒麼?當時那些門徒答保羅說:「沒有,也未曾聽見有聖靈賜下來」.我想而今也有些兄姊信道,雖是多年卻未聞有聖靈,此真奇之又奇.有人以為一信了主便受聖靈充滿,誰知這是不真的,因為使徒當未上耶路撒冷前,他們還未有靈充滿.有一次主對門徒說,你們當以「名在天上」為喜.請問信者的名在天國,主知否?這主定然知道.從這裡可見名在天國與受聖靈充滿是大有分別的:有以為信主得罪赦而得勝潔,就是聖靈充滿,其實聖潔與聖靈充滿也是有大分別的；有以為祈禱有力,能逐魔鬼,便是聖靈充滿,但門徒未得勝靈充滿前,也能逐魔,可見能逐魔與聖靈充滿也是大有分別的.
\newline
\newline
被聖靈充滿者是怎樣的:
\newline
\newline
(一)有非常之勇敢的
\newline
\newline
(二)能犧牲
\newline
\newline
(三)有大愛
\newline
\newline
(四)獻身是長久的
\newline
\newline
(五)認識自己在世為寄居的人
\newline
\newline
(六)仰望主再來
\newline
\newline
(七)可代人禱告
\newline
\newline
(八)愛讀聖經
\newline
\newline
(九)有能力為主作證
\newline
\newline
在英國有二位牧師,他們初是掘煤的工人,後得勝靈充滿,就大有能力,四出為主作證,今年四月十九在一最大的會堂講道,有一萬人聽他講道,有一千人受洗,有一日有三百人在會堂同得勝靈的醫治.
\newline
\newline
又在美國有一牧師,他是個博士,他已傳道十五年,自己還未信耶穌.他好說詼諧,討人歡喜,他說三分鐘便了.有一日,他聞有一奮興的會集,他也赴會,到會時,有人請他坐在前行的座位,他聽道後,便得勝靈之感動,主席問,誰信耶穌的血使罪得赦的就請舉手,於是眾人舉手,但他不舉.又再問,誰想聖靈赦罪,請前來俯伏禱告,那時他忘了博士的頭銜,便到台前禱告,於是他就受了聖靈之施洗.後來他在家人的面前作證說,自己是一個未重生；未受聖靈施洗的人,續著又說,他現在怎樣的得勝靈的施洗,最終便對會眾說,而今我們已得了聖靈的能幹來證道,前時會堂中的跳舞場要立刻取消了.有一天講道後,問誰願得勝靈施洗的,請前來,那日便有五百多人,得了聖靈施洗.
\newline
\newline
各位:你想,得了聖靈施洗的人,是何等的有力為主作證呢?現在我們要知得勝靈施洗的應許是屬你我的,所以我願我們都要求得勝靈的施洗,且為此事作見證,阿們.
\newline
\newline


\newpage

\section{第1屆~港九培靈會~祈理平~第三講}
\label{sec:1347}
\textbf{耶穌再臨}
\newline
\newline
連結: \href{https://www.hkbibleconference.org/session-message/view/1347}{\texttt{https://www.hkbibleconference.org/session-message/view/1347}}
\newline
\newline
\hyperref[sec:1346]{< < < PREV SERMON < < <}
~
\hyperref[sec:toc]{[返目錄]}
~
\hyperref[sec:1338]{> > > NEXT SERMON > > >}
\newline
\newline
使徒行傳 2:14-2:22
\newline
\begin{longtable}{cl}
\hline
\hline
章節 & 經文 (和合本修訂版)\\
\hline
2:14 & \begin{tabularx}{0.7\textwidth}{X} 彼得和十一個使徒站起來，他就高聲向眾人說：「猶太人和所有住在耶路撒冷的人哪，這件事你們要知道，要側耳聽我的話。 \end{tabularx} \\ \\ \relax
2:15 & \begin{tabularx}{0.7\textwidth}{X} 這些人並不像你們所想的喝醉了，因為現在才早晨九點鐘。 \end{tabularx} \\ \\ \relax
2:16 & \begin{tabularx}{0.7\textwidth}{X} 這正是藉著先知約珥所說的： \end{tabularx} \\ \\ \relax
2:17 & \begin{tabularx}{0.7\textwidth}{X} 『神說：在末後的日子，我要將我的靈澆灌凡血肉之軀的。你們的兒女要說預言；你們的少年要見異象；你們的老人要做異夢。 \end{tabularx} \\ \\ \relax
2:18 & \begin{tabularx}{0.7\textwidth}{X} 在那些日子，我要把我的靈澆灌，甚至給我的僕人和婢女，他們要說預言。 \end{tabularx} \\ \\ \relax
2:19 & \begin{tabularx}{0.7\textwidth}{X} 在天上，我要顯出奇事，在地下，我要顯出神蹟，有血，有火，有煙霧。 \end{tabularx} \\ \\ \relax
2:20 & \begin{tabularx}{0.7\textwidth}{X} 太陽要變為黑暗，月亮要變為血，這都在主大而光榮的日子未到以前。 \end{tabularx} \\ \\ \relax
2:21 & \begin{tabularx}{0.7\textwidth}{X} 那時，凡求告主名的都必得救。』 \end{tabularx} \\ \\ \relax
2:22 & \begin{tabularx}{0.7\textwidth}{X} 「以色列人哪，你們要聽我這些話：拿撒勒人耶穌就是神以異能、奇事、神蹟向你們證明出來的人，這些事是神藉著他在你們中間施行，正如你們自己知道的。 \end{tabularx} \\ \\
[1ex]
\hline
\hline
\end{longtable}
在新約中有三百二十次說及主再臨的事,為什麼許多教會領袖竟不講此問題?若說此事不可信,為什麼新約有如此多次說及此事?這是使徒的錯誤,還是我們的錯誤呢?有人說:「我不敢講,是因各人對於此事的見解不一.」唉!若只因此而不敢講,那就聖經的道,無一可講；因為眾人對聖經的見解,也是不一的.茲將主再臨諸問題略論如下:
\newline
\newline
主再臨的證據
\newline
\newline
(一)主的預言　主在約翰十四章一至三節預告門徒說:「你們心裡不要憂愁；你們信上帝,也當信我.在我父的家裡,有許多住處；若是沒有,我就早已告訴你們了；我去原是為你們預備地方去.我若去為你們預備了地方,就必再來接你們到我那裡去；我在那裡,叫你們也在那裡.」各位!主這樣的預言,這樣的應許,你信不信?在這裡主不說「或者」再來,是堅決的說,「就必」再來.你們以為主所說的話,可信不可信?若是不可信,那就聖經沒有一節可信；惟是這「就必再來」的話是可信的,故聖經每一節也可信.聖經的可信,因是受感於聖靈而寫的.
\newline
\newline
(二)天使的見證　當日使徒在橄欖山,望著主升天而去,他們的眼似乎不能轉回來；這因為他們的心在主那裡.信徒的心與主同去了,所以他們望主再來.各位兄姊,我們的心在哪裡呢?若是在世界,那就無怪沒有時候來讀經,聽道,祈禱了；各位同勞；你的心在那裡?若在耶穌那裡,那就千好萬好,不然,那就殆了!使徒心在主那裡,就定睛望天,望到天有聲回答他.他們只管望,不管望到幾何時,終底他們就得一個最有盼望的消息,就是「他怎樣往天上去,他還要怎樣來.」各位,天使報這個消息,有沒有錯呢?然我們就聖經而知天使所報的消息,沒有一個有錯的.如報給牧羊人的消息,是沒有錯,報給約瑟的消息也是沒有錯的；其餘如報給撒瑪利亞的消息,往看主的墳墓的婦人的消息,哥尼流的消息,保羅在船上遇險時的消息,也都是沒有錯的.今我們從各消息都是沒有錯的想來,就更信天使所報「他還要怎樣來」的消息,也是沒有錯的了；這「沒有錯」三字,盼望各位多多注重!
\newline
\newline
(三)使徒得的默示　我不解我們為什麼信使徒許多言,卻不信使徒從主而得的默示；初時我也不信,因為我嗜好吸煙.我既因己的罪惡而怕主再臨,我想主若再臨,我便不知怎樣了.我母是一受聖靈施洗的人,她常說:「主必快來」,故我每早必偵察她昨夜有沒有被主攝去.當時鄰人中有一人也常與我母談論主再臨的事,我實在厭棄他來說這些話；故我一見他從前門來,我就從後門去了.但而今我不怕了,不獨不怕,且最歡喜聽聞主再臨的聲音；因我知一切的罪,已為主血所遮蓋,所以我不用因罪而害怕了.各位,你害怕否?但無論你害怕不害怕,主必再臨,這是人所不能更改的；你若要不害怕,就只有站在主血之下.
\newline
\newline
主怎樣再臨
\newline
\newline
有人說主已再臨在我們的心,這話是背乎聖經的話,且也是與天使的,「他怎樣往天上去,他還要怎樣來」的見證相反.主升天之時是有可摸的形體；故使徒在啟示錄載著見主的釘痕.又有人說,人若死了,主就臨在那人處；不知人死是人去,不是主來.惟是聖經告訴我們說,主再臨時,是和著眾天使來的.如在猶大書中有說:「看哪!主帶著他的千萬聖者降臨,」你想,主再臨時的榮耀是何等的榮耀,惟惜多人不知!
\newline
\newline
主在何時再臨
\newline
\newline
從前我在美國時拾得一單張,說明主在某時某日再來；但我不信這樣的話,惟信主必再來.在提摩太後書三章所說的共有二十個主再臨的兆頭,今只說一二,以證明主再臨的時期是如何的近了.主說:「你看見這些兆頭,就知近及門了」.
\newline
\newline
(一)教會的狀況　各位試想今日的教會,熱心多呢,還是冷淡多呢?有人批評赴此培靈會的人,多是婦女；同時有人回答說,將來上天堂的也多是婦女.聖經曾說,聖靈充滿他的婢女,可知婦女近主多,得主的使用也多.今日的救會,是何等的冷淡,何等的忘神的恩.今日許多牧師,不說悔罪改過的道,只講迎合人心的什麼主義.今就歐美教會來說,也是一樣的冷淡；信徒為討一己的方便而用收音機聽道,卻不到禮拜堂；受洗也用相片來受洗,這真是教會衰落的證據.
\newline
\newline
(二)世界的狀況　最近美國報載,美國已給一約與英國,說此約成立,兩國便可和平.有人說,某某家可令世界和平；但主說,此世沒有和平.而今雖有什麼國際聯盟會,什麼和平會,年年不斷的會議,可是毒氣廠卻日日開工製造.你看,這是何等的假,若想世界得真和平,那真是難了,意大利國宣言說:「到一千九百五十三年時,我國便可出五百萬兵.」這是歡迎主來麼?不!世界而今預備一空前絕後的戰爭,故各位不要受人欺,你預備的是什麼呢?我就是預備接主,你呢?
\newline
\newline
(三)猶太國的狀況　主說:「你見無花果樹萌芽,就知夏時到了.」照樣,我們見猶太國的狀況,就知主快來了.從前有一國王問一書記說:「你怎能知有神」?書詞回答說:「我因猶太人而知」.這因為猶太人雖是國破家亡,散於天下,但他們仍得神祝福.而今每月約有三百猶太人回國.今猶太人已預備建築聖殿,且已訓練祭司.聖經預言橄欖山要分開,政府因信預言必應驗,故不許猶太人在山上建屋.
\newline
\newline
各位,你會否預備接主?你願否宣傳主再臨的問題?啊!主再來了,我就歡喜,因我有兩個兒女在那邊；一是十三個月,一是三歲,當三歲的臨去世時,我們捨不得他,他回答我們說:「容我到妹處!」各位,你有沒有親友在那邊等候你?
\newline
\newline


\newpage

\section{第1屆~港九培靈會~趙柳塘~第一講　概論分三}
\label{sec:1338}
\textbf{第一講　概論分三}
\newline
\newline
連結: \href{https://www.hkbibleconference.org/session-message/view/1338}{\texttt{https://www.hkbibleconference.org/session-message/view/1338}}
\newline
\newline
\hyperref[sec:1347]{< < < PREV SERMON < < <}
~
\hyperref[sec:toc]{[返目錄]}
~
\hyperref[sec:1339]{> > > NEXT SERMON > > >}
\newline
\newline
第一層　我們因何研究此書?　分三
\newline
\newline
(甲)因為迦拉太是一個很熱心而忽然退落了的教會──可作為我們的鑑戒
\newline
\newline
(乙)因為迦拉太教會在根本信仰上搖移了──可作我們的警告
\newline
\newline
(丙)因為保羅向此退落的教會仍有熱烈的愛──可作我們的法則
\newline
\newline
(甲)迦拉太教會乃保羅首生的兒子,第一次出門傳道的時候,便建設了迦拉太諸教會.由本書一章二節中寫信給「迦拉太的各教會」一語,可知迦拉太并非指一個教會,乃指迦拉太一省的諸教會.查(徒十三十四章)可知安提阿,以哥念,路司得,得比皆在此範圍之中,皆為保羅第一次出門傳道時所建立的教會.當時他們都是極熱心的疼愛保羅,甚至有信徒願意將他們自己的眼睛剜出來給保羅.(四之十四,十五)因保羅當日建立那些教會之時,確是由於聖靈的能力,確堅的信心,故在(三之一,三)有云:耶穌基督釘十字架已活活盡在他們眼前；他們又確已靠聖靈入門,而且已接受了聖靈,(三之二)并曾為主受了苦難.(三之四)本來跑得很好,因何快快離開(一之六)忽然由恩典的地位墜落了呢?(五之七,四)
\newline
\newline
如此便可見信徒之柔弱如羊,以及魔鬼的狠毒了；當保羅離開他們之時,便有外羊內狼的偽師入了他們之中,宣傳一種奇異的福音:(一之六七)假作熱心,間離迦人與保羅之愛,使迦人為他們熱心.(四之十七)可惜迦人不查,好奇心重,接受了那偽師異端的教訓,受了迷惑,便逐漸的退落了,這是何等可嘆可憐的事；因此保羅稱那些偽師們為窺探者,(二之四)攪擾者；(五之十)但保羅仍本其父母之愛心,熱熱烈烈的說:我小子阿；我為你們再受生產之苦,直等到基督成形在你們心裡.(四之十九).
\newline
\newline
領袖們:你看保羅於迦拉太教會,以聖靈生之,真理養之,為之切禱,每次出門又必探視之,(徒十四之二十一,二十二又十五之三十六)那信徒們一經試誘,尚且離道如此之易；況無保羅的愛心,靈力,殷勤,看顧教會的,又當如何?是以無怪今日教會冷淡者,如此之多!
\newline
\newline
(乙)看迦拉太教會他們的敗壞,并非完全的離棄了真道,不過在所信的救道上,又加上了割禮儀文與律法等事而已.(三之三,四之八至十一)又:(五之二至四)如此,不但失卻了「純全」二字,且於他們的根本信仰,也動搖了﹗所以保羅說:一點麵酵,能使全團都發起來.(五之九)那是何等危險的事；故主耶穌當日,亦嘗嚴重的警告信徒,當謹防法利賽人的酵.(可八之十五),以此,凡各使徒於偽師的教訓,莫不是嚴謹的防禦之.在本書(二之五)保羅說:於此,傳異端之人「我們就是一刻的工夫,也沒容讓順服他們」.又看那最仁慈的使徒約翰說:若有人到你們那裡不是傳這教訓,不要接他到家裡,也不要問他的安,(約翰二書十節)然而他們為什麼這樣嚴厲呢?因為那偽師是正道的仇敵,是羊群的豺狼,生死不兩立的:所以保羅說:如此,為要叫福音的真理,仍存在你們的中間.(二之五)
\newline
\newline
領袖們:容讓異端侵入,就是斷送群羊的生命,凡為善牧的,必要盡力爭辯之,防守之,(猶大書之三)信徒們容納異端,即就離棄正道；使他的信仰根本動搖,可是自尋死路!汝若不信,試一觀迦拉太教會,他們接納異端的結果,就是:──退步(四之十)──失福(十五)──與基督隔絕(五之四)──從恩典的地位墜落(五之四)!!
\newline
\newline
有云:一點砒霜攙入全盃白糖之中,那白糖便變了性,能使人死,一點異端入於正道之中,亦是如此.
\newline
\newline
(丙)迦拉太教會雖然退落了,失敗了,向保羅的愛情也斷絕了；但這個創立迦會的保羅怎麼樣呢?他說:小子阿,我為你們再受生產之苦,直等到基督成形在你們心裡.(四章十九節)那一句話中,包含無限的感情,與熱烈的血淚呢!這句話使吾們聯想到耶穌為耶路撒冷的痛哭,(太二十三之三十七至三十九節)與摩西在西乃山巔向耶和華捨棄一己之永生而代全以色列人之請願了；(出三十二之三十至三十四節)
\newline
\newline
領袖們:教會的工作,雖是世界一切工作中至高貴的,但也是至煩難的.吾們總要得耶穌保羅摩西那種捨己為人熱烈的愛情,然後方能成功.所以耶穌三次向彼得先要求他的「愛心」,然後乃以羊群的重責託之.(約二十一之十五至十七)我儕試一思念保羅的經歷,他又曾向哥林多的信徒說:我也甘心樂意為你們的靈魂費財費力,難道我越發愛你們,就越發少得你們的愛麼?(哥後十二章十五節)設使你遇見了那種的經歷,你當如何?──辭職呢,──求另調呢,──或是改業呢?──解決那問題的,還是在(本書四之十九).同勞們；有誰父母因子女之背逆而請辭父母的職份呢?善牧人,必不因教會的冷淡,困難而辭職的,只有因此多多為他們流淚祈禱,求主奮興之而已.吾們更當記著雅各的話說,我們弟兄們,你們中間若有失迷真道的,有人使他回轉；這人該知道叫一個罪人從迷途上轉回,便是救一個靈魂不死,並且遮蓋許多的罪(五之十九,二十)
\newline
\newline
第二層　我們如何研究此書?　分三
\newline
\newline
(甲)研究書信的普通方法是怎樣?
\newline
\newline
(乙)本書中所最當注意的問題是什麼?
\newline
\newline
(丙)本書的分解如何?
\newline
\newline
(甲)研究書信的普通方式──
\newline
\newline
(一)該教會所在地的地理歷史如何?
\newline
\newline
迦拉太是省名,也是一城的名,在小亞細亞.本書所說的迦拉太,是指著一省說的,故迦拉太教會是包括了安提阿,以哥念,路司得,得比一切的教會.
\newline
\newline
那人民尚武好戰,而性情更變無常,即如他們初時悅聽保羅,待之如同天使,願意將自己的眼目剜出來,送給保羅.(四之十三至十五)及那傳異端的偽師到來,一聽他們的教訓,即離棄了保羅而別從一福音.(一之六)
\newline
\newline
(二)該教會成立的歷史如何?
\newline
\newline
讀徒十三十四章,知迦拉太諸教會,都是保羅親手創立的.所以在書信中,他不特自稱使徒,也常叫迦拉太的信徒為小子們.(四之十九)而且書中的口氣,如同父親訓子一般.
\newline
\newline
(三)作書者誰,又作書之時與地如何?
\newline
\newline
本書是保羅在以弗所作的,約在主後五十七年.
\newline
\newline
此當注意,解經家於此一項常有出入.故亦有云本書乃保羅於五十八年時在哥林多所作.
\newline
\newline
(四)該書信具何目的?
\newline
\newline
保羅作這書,乃想迦拉太的信徒離棄異端,仍歸真道.(五之十又四之二十)
\newline
\newline
(五)全書之要義如何?
\newline
\newline
保羅在迦拉太書,說明得救之道,乃由於信與基督之恩,非由於行；與羅馬人書的分別,就在本書注重「非由於行」四字(二之十六,十七又五之四)
\newline
\newline
(六)全書如何分解?
\newline
\newline
拉太書分三段
\newline
\newline
第一段是事實,(一,二章)論保羅為使徒之憑據如何?
\newline
\newline
── 關於個人之行為即做僕人的應當怎樣?
\newline
\newline
第二段是理論,(三,四章)論得救之道,由信非由於行.
\newline
\newline
── 關於聖道之奧妙即得救的方法怎樣?
\newline
\newline
第三段是實踐,(五,六章)論信徒當如何實行主道?
\newline
\newline
── 關於信徒日日之行為,即得救了的人應當怎樣?
\newline
\newline
第三段　本書的幾個特點
\newline
\newline
(一)本書是保羅與家兄弟聯名的一封公函,我儕當要特別嚴重的來讀.(一之一,二)
\newline
\newline
(二)本書不是達於一個堂會的書信,乃是達於一地眾教會的書信,是眾信徒人人應讀的.(一之二)
\newline
\newline
(三)本書信的開端,完全沒有一句稱讚或感謝的話；因為一個失福入於異端的教會,只有令人心傷嘆息,有何感讚之可言呢?
\newline
\newline
(四)保羅咒詛偽師,「咒詛」二字原文「是阿拿提瑪」(一之八,九)這是希臘文中一個很嚴重不易輕用的字,意即「剪除於群眾之中,且永見絕於主神之前.」因為偽師傳異端,使人絕於基督,由恩典中墜落,且永絕於主前(五之四)其本身豈不應自受此刑麼?保羅如此咒詛他們,實不為過.
\newline
\newline
(五)保羅的靜修生活,惟在本書中略略知之,他信道復隱藏了三年在亞拉伯(一之十七,十八)十四年後乃明顯於眾,那是指著使徒行傳十伍章的事說的.
\newline
\newline
(六)保羅為人之正直公義,在指責彼得一事明顯了出來；(二之十一)保羅為要福音之真存在信徒之中,故盡力爭戰,不愧為一個忠誠的僕人(提前一之十二)(提後四之七,八)
\newline
\newline
(七)二章二十節可是信徒靈程上一句最重要最深妙的話,其意有四──
\newline
\newline
(甲)生命是由死而來,(釘十字架而活)
\newline
\newline
(乙)復活的生命,是屬天的是屬神的生命,(基督在我裡面活)
\newline
\newline
(丙)如此的生命,不是科學可能說明,乃是神秘的.信上帝的兒子而活,意即憑上帝子之信而活.
\newline
\newline
(丁)如此的生命,是由恩而得,屬神之賜,非人之力.(愛我,為我捨己)
\newline
\newline
(八)亞伯拉罕是第一個傳福音的人,他所傳的,也是因信稱義之道(三之八)可作耶穌宣傳救道之預影,比較約三之十六節之信他的......得永生)參看本書三之十四及注意八節中之「萬國都必因你得福」一語
\newline
\newline
(九)(四章十九節)是作主工的一節最要的言語,也是成功的一個秘訣.
\newline
\newline
(十)書信中惟此一書是保羅親手所寫的.(六之十一)由是,更顯出此書的特別重要,及他的價值,以及保羅向迦拉太人熱烈的愛.
\newline
\newline
(十一)在異邦教會的保羅書信中,惟在此一書有為以色列祝福的話.(六之十六)
\newline
\newline
(十二)本書最終的話是祝福,(六之十八)迦人雖然入於異端,失主之福,但保羅至終之希望,仍願主恩在他們的心裡.這是教訓吾們,凡信徒與教會無論敗壞至何境地,由主之恩,仍然有望,仍可代求,仍當望主施以特別的恩典.(約壹五之十六)
\newline
\newline


\newpage

\section{第1屆~港九培靈會~趙柳塘~第二講}
\label{sec:1339}
\textbf{第二講}
\newline
\newline
連結: \href{https://www.hkbibleconference.org/session-message/view/1339}{\texttt{https://www.hkbibleconference.org/session-message/view/1339}}
\newline
\newline
\hyperref[sec:1338]{< < < PREV SERMON < < <}
~
\hyperref[sec:toc]{[返目錄]}
~
\hyperref[sec:1340]{> > > NEXT SERMON > > >}
\newline
\newline
加拉太書 1:1-2:21
\newline
\begin{longtable}{cl}
\hline
\hline
章節 & 經文 (和合本修訂版)\\
\hline
1:1 & \begin{tabularx}{0.7\textwidth}{X} 我使徒保羅和所有跟我一起的弟兄，寫信給加拉太的眾教會。我作使徒不是由於人，也不是藉著人，而是藉著耶穌基督與使他從死人中復活的父神。 \end{tabularx} \\ \\ \relax
1:2 & \begin{tabularx}{0.7\textwidth}{X} 見上節 \end{tabularx} \\ \\ \relax
1:3 & \begin{tabularx}{0.7\textwidth}{X} 願恩惠、平安從我們的父神和主耶穌基督歸給你們！ \end{tabularx} \\ \\ \relax
1:4 & \begin{tabularx}{0.7\textwidth}{X} 基督照我們父神的旨意，為我們的罪捨己，要救我們脫離現今這罪惡的世代。 \end{tabularx} \\ \\ \relax
1:5 & \begin{tabularx}{0.7\textwidth}{X} 願榮耀歸給神，直到永永遠遠。阿們！ \end{tabularx} \\ \\ \relax
1:6 & \begin{tabularx}{0.7\textwidth}{X} 我很驚訝你們這麼快就離開那位藉著基督之恩呼召你們的神，而去隨從別的福音； \end{tabularx} \\ \\ \relax
1:7 & \begin{tabularx}{0.7\textwidth}{X} 其實並沒有另一個福音，不過有些人騷擾你們，要把基督的福音更改了。 \end{tabularx} \\ \\ \relax
1:8 & \begin{tabularx}{0.7\textwidth}{X} 但無論是我們或是天上來的使者，若傳福音給你們，與我們所傳給你們的不同，他該受詛咒！ \end{tabularx} \\ \\ \relax
1:9 & \begin{tabularx}{0.7\textwidth}{X} 我們已經說了，現在我再說，若有人傳福音給你們，與你們以往所領受的不同，他該受詛咒！ \end{tabularx} \\ \\ \relax
1:10 & \begin{tabularx}{0.7\textwidth}{X} 我現在是要得人的心，還是要得神的心呢？難道我在討人的喜歡嗎？我若仍舊想討人的喜歡，我就不是基督的僕人了。 \end{tabularx} \\ \\ \relax
1:11 & \begin{tabularx}{0.7\textwidth}{X} 弟兄們，我要你們知道，我所傳的福音不是按照人的意思； \end{tabularx} \\ \\ \relax
1:12 & \begin{tabularx}{0.7\textwidth}{X} 因為我不是從人領受的，也不是人教導我的，而是藉著耶穌基督的啟示而來。 \end{tabularx} \\ \\ \relax
1:13 & \begin{tabularx}{0.7\textwidth}{X} 你們聽說過從前我在猶太教中的行徑，我怎樣竭力壓迫殘害神的教會。 \end{tabularx} \\ \\ \relax
1:14 & \begin{tabularx}{0.7\textwidth}{X} 在猶太教中，我比本國許多同輩的人更激進，為我祖宗的傳統更熱心。 \end{tabularx} \\ \\ \relax
1:15 & \begin{tabularx}{0.7\textwidth}{X} 然而，那位把我從母腹裡分別出來、又施恩呼召我的神，既然樂意 \end{tabularx} \\ \\ \relax
1:16 & \begin{tabularx}{0.7\textwidth}{X} 把他兒子啟示在我心裡，讓我在外邦人中傳揚他，我就沒有跟有血有肉的人商量， \end{tabularx} \\ \\ \relax
1:17 & \begin{tabularx}{0.7\textwidth}{X} 也沒有上耶路撒冷去見那些比我先作使徒的，惟獨到阿拉伯去，後來又回到大馬士革。 \end{tabularx} \\ \\ \relax
1:18 & \begin{tabularx}{0.7\textwidth}{X} 過了三年，我才上耶路撒冷去見磯法，和他同住了十五天。 \end{tabularx} \\ \\ \relax
1:19 & \begin{tabularx}{0.7\textwidth}{X} 至於別的使徒，除了主的兄弟雅各，我都沒有見過。 \end{tabularx} \\ \\ \relax
1:20 & \begin{tabularx}{0.7\textwidth}{X} 我現在寫給你們的是在神面前說的，不說謊話。 \end{tabularx} \\ \\ \relax
1:21 & \begin{tabularx}{0.7\textwidth}{X} 以後我到了敘利亞和基利家一帶； \end{tabularx} \\ \\ \relax
1:22 & \begin{tabularx}{0.7\textwidth}{X} 那時，在基督裡的猶太各教會都沒有見過我的面。 \end{tabularx} \\ \\ \relax
1:23 & \begin{tabularx}{0.7\textwidth}{X} 不過他們聽說「那從前壓迫我們的，現在竟傳揚他原先所殘害的信仰」。 \end{tabularx} \\ \\ \relax
1:24 & \begin{tabularx}{0.7\textwidth}{X} 他們就為我的緣故歸榮耀給神。 \end{tabularx} \\ \\ \relax
2:1 & \begin{tabularx}{0.7\textwidth}{X} 過了十四年，我再上耶路撒冷去，巴拿巴同行，也帶了提多一起去。 \end{tabularx} \\ \\ \relax
2:2 & \begin{tabularx}{0.7\textwidth}{X} 我是奉了啟示上去的；我把在外邦人中所傳的福音對弟兄們說明，我是私下對那些有名望的人說的，免得我現在或是從前都徒然奔跑了。 \end{tabularx} \\ \\ \relax
2:3 & \begin{tabularx}{0.7\textwidth}{X} 但跟我同去的提多，雖是希臘人，也沒有勉強他受割禮； \end{tabularx} \\ \\ \relax
2:4 & \begin{tabularx}{0.7\textwidth}{X} 因為有偷著混進來的假弟兄，暗中窺探我們在基督耶穌裡擁有的自由，要使我們作奴隸， \end{tabularx} \\ \\ \relax
2:5 & \begin{tabularx}{0.7\textwidth}{X} 可是，為要使福音的真理仍存在你們中間，我們一點也沒有讓步順服他們。 \end{tabularx} \\ \\ \relax
2:6 & \begin{tabularx}{0.7\textwidth}{X} 至於那些有名望的，不論他們是何等人，都與我無關；神不以外貌取人。那些有名望的，並沒有加增我甚麼。 \end{tabularx} \\ \\ \relax
2:7 & \begin{tabularx}{0.7\textwidth}{X} 相反地，他們看見了主託付我傳福音給未受割禮的人，正如主託付彼得傳福音給受割禮的人； \end{tabularx} \\ \\ \relax
2:8 & \begin{tabularx}{0.7\textwidth}{X} 那感動彼得、叫他為受割禮的人作使徒的，也感動我，叫我為外邦人作使徒。 \end{tabularx} \\ \\ \relax
2:9 & \begin{tabularx}{0.7\textwidth}{X} 那些被認為是教會柱石的雅各、磯法、約翰知道神所賜給我的恩典，就跟我和巴拿巴握右手以示合作，同意我們往外邦人那裡去，他們往受割禮的人那裡去。 \end{tabularx} \\ \\ \relax
2:10 & \begin{tabularx}{0.7\textwidth}{X} 他們只要求我們記念窮人，這也是我一向熱心在做的。 \end{tabularx} \\ \\ \relax
2:11 & \begin{tabularx}{0.7\textwidth}{X} 後來，磯法到了安提阿，因為他有可責之處，我就當面反對他。 \end{tabularx} \\ \\ \relax
2:12 & \begin{tabularx}{0.7\textwidth}{X} 從雅各那裡來的人未到以前，他和外邦人一同吃飯，及至他們來到，他因怕奉割禮的人就退出，跟外邦人疏遠了。 \end{tabularx} \\ \\ \relax
2:13 & \begin{tabularx}{0.7\textwidth}{X} 其餘的猶太人也都隨著他裝假，甚至連巴拿巴也隨夥裝假。 \end{tabularx} \\ \\ \relax
2:14 & \begin{tabularx}{0.7\textwidth}{X} 但我一看見他們做得不對，與福音的真理不合，就在眾人面前對磯法說：「你既是猶太人，卻按照外邦人的樣子，不按照猶太人的樣子生活，怎麼能勉強外邦人按照猶太人的樣子生活呢？」 \end{tabularx} \\ \\ \relax
2:15 & \begin{tabularx}{0.7\textwidth}{X} 我們生來就是猶太人，不是外邦罪人； \end{tabularx} \\ \\ \relax
2:16 & \begin{tabularx}{0.7\textwidth}{X} 可是我們知道，人稱義不是因律法的行為，而是因信耶穌基督，我們也信了基督耶穌，為要使我們因信基督稱義，不因律法的行為稱義，因為，凡血肉之軀沒有一個能因律法的行為稱義。 \end{tabularx} \\ \\ \relax
2:17 & \begin{tabularx}{0.7\textwidth}{X} 我們若求在基督裡稱義，自己卻還被視為罪人，那麼，基督是罪的用人嗎？絕對不是！ \end{tabularx} \\ \\ \relax
2:18 & \begin{tabularx}{0.7\textwidth}{X} 如果我重新建造我所拆毀的，這就證明自己是違犯律法的人。 \end{tabularx} \\ \\ \relax
2:19 & \begin{tabularx}{0.7\textwidth}{X} 我因律法而向律法死了，使我可以向神活著。我已經與基督同釘十字架， \end{tabularx} \\ \\ \relax
2:20 & \begin{tabularx}{0.7\textwidth}{X} 現在活著的不再是我，乃是基督在我裡面活著；並且我如今在肉身活著，是因信神的兒子而活；他是愛我，為我捨己。 \end{tabularx} \\ \\ \relax
2:21 & \begin{tabularx}{0.7\textwidth}{X} 我不廢掉神的恩；如果義是藉著律法而獲得，那麼基督就白白死了。 \end{tabularx} \\ \\
[1ex]
\hline
\hline
\end{longtable}
總題
\newline
\newline
(一)保羅與他的職份
\newline
\newline
(二)保羅與他的宗教
\newline
\newline
(三)保羅與傳異端的偽師
\newline
\newline
(四)保羅與教會的同勞
\newline
\newline
第一層　保羅與他的職份　有三
\newline
\newline
(甲)他的職份怎樣得來?(一之一及十)
\newline
\newline
(乙)他所傳的道怎樣得來?(一之十一,十二)
\newline
\newline
(丙)保羅盡職是用何種的方法?(一之十五至十七又二之一,二)
\newline
\newline
(甲)保羅論到他的職份他提出兩個要點?(一之一)
\newline
\newline
(一)他說不是由於人.
\newline
\newline
(二)他說不是藉著人；如是,他將他的地位先確定了.因為他的職份是藉著耶穌基督及父上帝而得,故他作工(一)不必要得人心(二)亦不求人的喜悅.(一之十)他的作工,只向上帝與主耶穌基督負絕對的責任；無論迦拉太教會的信徒向他存何觀感,他只要問心無愧的向神盡責,來 ── 勸勉 ── 教訓 ── 督責.
\newline
\newline
同勞們:我儕為主作工,如不先確定我們的地位,只重看教會的聘請與選舉,或機關的委任,而不重看神的選立；我儕的工作,必定陷於求悅於人,順從人意的網羅,終至灰心辭職而改業!
\newline
\newline
(乙)進一步就看到保羅論他所傳的道(一之十一,十二)也有三個要點 ──
\newline
\newline
(一)不是出於人意的,
\newline
\newline
(二)不是從人領受的,
\newline
\newline
(三)不是由人教道的,── 乃是耶穌基督啟示來的!
\newline
\newline
這樣看來,保羅於人神兩方面,是分得狠清楚的.他的職份固然不是由於己意的選擇與人的委任；故此他所傳的道,也完全無 ── 人意,人力之工在於其中；乃是全由主耶穌的啟示而來,故可稱為 ── 神道,── 非人道,而能運行在聽者的心中,發出效用來.(帖前二之十三,西一之六)
\newline
\newline
同勞們:我們在講臺宣傳的言語,是由我們的腦汁中絞出來的呢?或是由報紙雜誌中搜羅出來的呢?還是由膝頭上祈禱中得勝靈的感動啟示,由聖經中發明出來的呢?(耶利米二之十三)有說:「離棄主這活水的泉源,為自己鑿出池子,是破裂不能有水的池子.那就是聖經明明的告訴我們說:當日的偽先知們,離棄主的教訓,猶如離棄了活水的泉源；自己的工作是無用的,猶如破裂不能存水的池子一般.而且他們輕輕忽忽的醫治主百姓的損傷,說平安了,平安了,其實沒有平安!(耶八之十一)這類的工人,今日教會中是很多的.然而聖經中還有一方面的教訓,就是親近主,祈禱主的人,他的心湧出美辭,他的話頭是快手筆.(詩四十五之一)以賽亞為此亦曾作証說:「主耶和華賜我受教者的舌頭,使我知道怎樣用語扶助疲乏的人,主每早晨提醒提醒我的耳朵,使我聽像受教者一樣」(五十章四節)同勞們:你講道時覺得枯乾麼?你勸人的言語覺得無力麼?你當早早起身,多多祈禱,然後看看功效如何!
\newline
\newline
(丙)今再看到保羅如何盡職及其成功的方法 ── 他既得了主的恩召及啟示,并接受了傳道異邦的使命(一之十五,十六)他就 ──
\newline
\newline
(一),不與屬血氣的商量,
\newline
\newline
(二),不去見那先他作使徒的人,── 惟獨往亞拉伯去.
\newline
\newline
那獨字的關係,是何等的重大；今日有人明明得了保羅一切的主所賜的機會,但至終仍不能在主的工場效力,是何緣故呢?因為他要合於己意,也要得人的贊同,然後方敢放膽的遵行主旨.試看耶和華說:我的意念非同你們的意念；我的道路,非同你們的道路；天怎樣高過地,照樣我的道路高過你們的道路；我的意念,高過你們的意念.惡人常離棄自己的道路,不義的人當除掉自己的意念,歸向耶和華；耶和華就必憐恤.(賽五十五章八,九,七節)如此樂意遵行主旨而得了好處的,我們可看在創世記二十二章所載亞伯拉罕的一段經歷.即當注意三節之清早二字.(又參看徒二十二章十九至二十一節)由是,可知主僕的成功,就在於得知主旨,即去遵行便是了.
\newline
\newline
第二層　論保羅與他的宗教 ── 由這方面,使吾們看明了上帝救人的方法.保羅本是一個十分固守猶太舊教的法利賽人,他自說:他在猶太教中,比他本國許多同歲的人,更有長進；為他祖宗的遺傳,更加熱心.你們聽見我從前在猶太教中所行的事,怎樣逼迫殘害上帝的教會；(本書一章十四,十三節)這樣的人,上帝竟選了他做使徒；託以最重最大之責任,你說奇不奇呢?(二之七,八及徒九之十五,十六)所以當日主叫亞拿尼亞去見保羅為他祈禱,使他得見,為他施洗之時；亞拿尼亞心中十分的懷疑.(徒九之十三)及後保羅到了大馬色為主作証,多人為之驚奇.(徒九之二十一)後來又到了耶路撒冷,想與門徒結交,他們卻都怕他,不信他是門徒.(徒九之二十六節)然而上帝竟無絲毫的思疑,選用了他,那是何等的愛呢!故此,他即將他祖宗所傳,自幼所信的宗教,拋棄了；專心盡力的來事奉以恩召他的主耶穌基督.而且他又向天下萬人自證說:在罪人中我是罪魁,然而我蒙了憐憫,是耶穌基督要在我這罪魁身上,顯明他一切的忍耐,給後來信他得永生的人作榜樣；(提前一章十五,十六節)他又說,基督的愛激勵了我,所以他樂意為主盡忠至死.(哥後五章十四節)
\newline
\newline
同勞們:我們今日所傳之耶穌,因何無力感動人.離棄他虛偽的敬拜,來事奉生且榮的上帝；亦不能使人窺破世界的虛榮,來接受確且實的永生呢?
\newline
\newline
試看五旬節時的彼得,一日傳道有三千人悔改信主.那三千人都是完全自由,毫無宗教,風俗,習慣等束縛的人麼?然則他們因何得有能力勝過了他們的宗教,風俗,習慣,而歸信救主耶穌呢?又看保羅在以弗所的工作,亦能使平素行邪術的人,絕大的犧牲而歸信救主.(使十九章十九節)
\newline
\newline
那都是耶穌基督的感力,使他們得了勝,吸引了他們歸於自己:吾人今再一讀保羅傳道之經歷,便知道他在大馬色道中之所以四肢投地,降服在於主前的,全在於──主之光,── 與主之聲.(徒九章三,四節)那就是使保羅脫離了他舊宗教束縛的能力!
\newline
\newline
同勞們:吾們有何方法使一切人意料想不到,人必推測,絕對不能信主耶穌之人歸信耶穌呢?亦惟靠 ──(一)聖靈之光,光照其心之黑暗.(徒二十六章十八節)──(二)上主之言,以擊破其心之頑梗；(耶二十三章二十九節)──(三)耶穌之愛,以鎔化其心之剛愎而已!保羅說:我們爭戰的兵器,本不是屬血氣的,乃是在上帝面前有能力,可以攻破堅固的營壘；將各樣的計謀,各樣攔阻人認識上帝的那些自高之事,一概攻破了.又將人所有的心意奪回,使他都順服基督.(哥後十之四,五)
\newline
\newline
第三層　論保羅與偽師 ── 這層乃論保羅怎樣對待宣傳異端之人.
\newline
\newline
第一,他說:這等人應當咒詛.(一之八,九)第二．他說:一刻功夫也沒有容讓順服他們.(二之五)第三．他說:他必擔當他的罪名.(必受其刑五之十)第四．他說:願這等人被斷絕(五之十二)
\newline
\newline
保羅因何要這等嚴厲的對待這等人呢? ── 第一:因為他們更改了基督的福音.(一之七)第二:因為他們攪擾信徒.(一之七；五之十)第三:因為他們暗中窺探,如保羅所說:有偷著引進來的假兄弟,私下窺探我們在基督耶穌裡的自由,要叫我們作奴僕.(二之四)他又說:那偷進人家牢籠無知婦女的,正是這等人,這些婦女擔負罪惡,被各樣的私慾引誘(提後三之六)第四:因他們離間信徒,想人為他的熱心.(四之十七)第五:因為他們阻礙信徒的進步叫人不順從真理(五之七)第六:因他們勉強信徒受割禮,得以誇口.(六之十二,十三)第七:因為他們使得救之人,與基督隔絕,從恩典中墜落.(五之四)
\newline
\newline
偽師教訓的流毒,有此種種；是凡真心愛主愛護信徒的僕人,應當如何呢?寬容他們,或是拒絕他們呢?今日異端充斥教會之中,是誰之過呢?耶穌說:你從世界上賜給我的人,我已將你的名顯明與他們,他們本是你的,你將他們賜給我,他們也遵守了你的道.我與他們同在的時候,因你所賜給我的名保守了他們,我也護衛了他們.其中除了那滅亡之子,沒有一個滅亡的.(約十七之六,十二)保羅說:我就是一刻的工夫也沒有容讓順服他們,為要叫福音的真理仍存在你們中間.(二之五)然而有人覺得如此,是太狹窄了!
\newline
\newline
同勞們:豺狼進入羊群,是要劫奪吞噬.善牧如何對付豺狼?吾們亦當如何對付一切宣傳異端,以及離棄正道教訓之人.
\newline
\newline
第四層　論保羅與教會的同勞 ── 好了,吾們可以離開黑暗的景象,進入光明之地了.因為偽師是我們的敵人,所以,不得已我說許多傷心的話；同勞是吾們的同志,由此吾們看出了保羅高貴的人格,必能給吾們許多良好的模範.
\newline
\newline
(甲)不結交權貴(一之十八至二十二又二之二,六)── 趨炎附勢,本是惡劣社會中一種污點.不料主耶穌聖潔清高的教會,不知不覺的竟也沾染了這種的惡習；各自擁護一個領袖,使紛爭結黨起來,(哥前一之十又三之三,四節)保羅於此一事,是為今日教會最要學效的.
\newline
\newline
吾們看那幾節聖經,便可知保羅是一個絕不出風頭的領袖.他得主絕大光榮的選召,他竟過了三年才上耶路撒冷去住了十五天；見一見彼得雅各便走了.(一之十八),我們可說他是一個絕無社交之人,對於今日許多教會領袖,他誠然有愧了!
\newline
\newline
再過了十四年,他又上耶路撒冷,還是遵著啟示,為著教會爭辯信仰的大問題而去的.(徒十五章)那時他已出門一次,建立了幾個教會,已經為主作了一點工夫的了.然而他上去,乃是背地裡向耶路撒冷教會的領袖們談談.在那會議中,主任說話的,是彼得與巴拿巴.(徒十五之七,十二)這種沉靜隱默的態度,令人是何等的可慕可愛呢?(太十二之十九)
\newline
\newline
(乙)各盡己職(二之六至八)── 保羅說:至於那些有名望的,不論他是何等人,都與我無干；上帝不以外貌取人,那些有名望的,並沒有加增我甚麼?反倒看見了主託我傳福音給那未受割禮的人,正如託彼得傳福音給那受割禮的人；那感動彼得叫他為受割禮之人作使徒的,也感動我叫我為外邦人作使徒.保明白了他的職份,他不忌不貪,惟忠心的盡職.今日教會中因何這麼多的紛爭呢?蓋因許多人不明白自己的職份,甘心守份的去作他應作之工而已.
\newline
\newline
(丙)忠心規友(二之十一至十四)── 此為今日教會中最大之欠缺,故論之略詳 ──
\newline
\newline
(一)保羅所規勸者何人? ── 彼得
\newline
\newline
彼得為誰呢?豈不是那使徒的首領,教會的柱石麼?他也會有錯誤麼?若他行錯了,誰人竟敢勸勉呢?保羅說:後來磯法到安提阿,因他有可責之處,我就當面抵擋他.(二之十一)在此有兩方面的教訓:(一),做領袖們的當知道,你亦有行錯,當受人勸勉之時.更當知你所處之地位特高,尤易使人見你之錯過,故凡為領袖之人,當如何謙卑,更當如何謹慎!(二),教會中實甚欠缺忠誠如保羅之人,如能有多人如保羅,主道必更榮顯光明於此暗世了!
\newline
\newline
(二)因何規友?── 因彼得行錯,其餘的猶太人也都隨著他裝假,甚至連巴拿巴也隨夥裝假.(二之十三)
\newline
\newline
領袖行錯,影響至大；因愛護教會之信徒,故不能不即時更正之.所以保羅說,我一看見他們行的不正,與福音的真理不合,就在眾人面前指正他.(二之十四)試看利未記四章所記贖罪祭之例之別 ──
\newline
\newline
犯罪者 祭品 節數
\newline
\newline
(三)如何規友.── 正人的過失,如醫士之用刀割症,是一件很要小心注意的事.
\newline
\newline
其法有三 ──
\newline
\newline
一,當有完全之忍耐溫柔與謙卑；
\newline
\newline
二,當日完全的愛心,
\newline
\newline
三,當有懇切之祈禱
\newline
\newline
保羅說:若有人偶然被過犯所勝,你們屬靈的人,就應當用溫柔的心把他挽回過來.(六之一)由此看來,正人的過失,本是屬靈的人的責任；又當用溫柔之心,方能有效.他又曾教訓提摩太說:主的僕要溫溫和和的待眾人,善於教導,存心忍耐,用溫柔勸戒那抵擋的人；或者上帝給他們悔改的心,可以明白真道,規勸行錯之人,如無忍耐與溫柔,是剛愎其心；使罪人深入網羅,永無脫罪自由之日,由忍耐,我儕以等待主所賜吾人以勸勉的機會,然後以溫柔的言語解說指正之,使之自悟.(箴二十五之十五)但吾們又當記著,吾們勸人的人,還要低著頭,曲著腰,謙卑在罪人之前,俯服在主的足下,方能有效.試一思想主耶穌如何方能洗門徒之足呢?站住呢,還是低頭伏腰呢?吾們因為未能完全自卑於主前人前,故永不能洗人之足,去人之污!
\newline
\newline
有了完全的忍耐與溫柔,又當加上完全的愛心.這是規勸那過失之人的主要能力.耶穌曾對門徒說:你們可以放心,我已經勝了世界.(約十六之三十三)耶穌不特勝了世界之苦難,也勝了世界的罪惡；並且將世界罪惡敵己之人,使之降服了自己,那纔是真正的勝了世界；然而耶穌究竟是用何種方法勝了世界? ── 就是他的愛.保羅說:惟有基督在我們還作罪人的時候,為我們死；上帝的愛,就在此向我們顯明了.(羅五之八)保羅己身可為見証,他說:在罪人中我是個罪魁,然而我蒙了憐憫；是因耶穌基督要在我這罪魁身上,顯明他一切的忍耐,給後來信他得永生的人作榜樣.(提前一之十五,十六)故此他為基督如此之愛激勵了,他便甘心情願的不再為自己活,乃為替他們死而後復活的主活；(哥後五之十五)是以拿破侖曾自嘆說:拿撒勒人耶穌勝過了我,因為拿破侖曾有兵百萬,稱雄世界﹗但失敗,被囚荒島之中.從之者若干?然而耶穌基督,以愛臨世,千數百年後,世界甘心為其捨命者,難以數計.
\newline
\newline
夫世界最惡毒最永久者,其為罪.然愛能以勝之.設吾人有完全之愛,未有不能感失足之兄姊回頭,與使罪人悔改也.如其不然,是在吾人之愛心不完全耳.如吾人已盡其完全之愛,而仍無功效,是其人之定命,是為沉淪之人也.吾憶某聖堂有一聯云 ──
\newline
\newline
宣戰只攻心,一愛能平群敵壘﹗援殊以手,重生況普萬民靈,
\newline
\newline
又嘗聞一軼事,亦可引用於此 ──
\newline
\newline
昔有一最強惡的強盜,被官兵捉住了.曾用了許多許多極嚴厲的刑罰審問他,仍不能他招認他的罪惡,以定案.一日,他的母親去監中看他,見他全身的皮肉都損傷了,他母親用手摸他的傷痕,他熱烈慈愛的淚,禁不住點點滴滴的弔了下來；她說:兒呀,認了罷；然後官長一問,便全認了.兄弟們:吾們因何不能使犯罪之兄姊認罪悔改呢?蓋因吾們熱烈慈愛之淚,尚未滴在他們罪痕之上罷!看保羅之規勸彼得,純全是出於愛彼得,與愛眾信徒們之心；是以彼得至終還是認保羅為他的愛兄.(彼後三之十五)(箴二十七之五,九又二十八之二十三)
\newline
\newline
領袖們,聖經說:我們如缺少了智慧,應當向主祈求.(雅一之五)有人有了過犯,你當如何規正他,你如切切祈禱主,主必指示你最好的方法.你若覺欠缺了忍耐,溫柔,謙卑,愛心,以致在規勸人的事上失敗了,你如祈禱主,主必賜你充足的忍耐,溫柔,謙卑,興愛心.因為主曾對他的僕人保羅說:主的恩典是彀保羅用的；如是,也是彀吾們用的!(哥後十二之九)
\newline
\newline


\newpage

\section{第1屆~港九培靈會~趙柳塘~第三講}
\label{sec:1340}
\textbf{第三講}
\newline
\newline
連結: \href{https://www.hkbibleconference.org/session-message/view/1340}{\texttt{https://www.hkbibleconference.org/session-message/view/1340}}
\newline
\newline
\hyperref[sec:1339]{< < < PREV SERMON < < <}
~
\hyperref[sec:toc]{[返目錄]}
~
\hyperref[sec:1341]{> > > NEXT SERMON > > >}
\newline
\newline
本書第一大段之分析
\newline
\newline
第一章
\newline
\newline
大旨 ── 保羅與福音及偽師與異端之比較
\newline
\newline
分段 ──
\newline
\newline
(一)本書之小引(一至五節)
\newline
\newline
(二)偽師與異端(六至十節)
\newline
\newline
(三)保羅之蒙召及其所傳福音之來歷(十一至二十四節)
\newline
\newline
第一段　本書之小引(一至五節)
\newline
\newline
(甲)恩惠平安(三節)諸書信問安中祝福之言,其次序皆如此(哥前一之三彼後一之二)先恩惠而後平安,世人如不在主恩典之下,必不能得平安；蓋平安之源,乃主基督.彼之名稱,為和平之君.(賽九之六,七,路二之十四)故保羅云:願恩惠平安從父上帝與我們的主耶穌基督歸與你們.(一之三)凡不主者,不特不得平安,且主之震怒常在其身上；(約三之三十六)試思吾人每當戰爭之時,其最需求者何?非平安乎.故平安實為人類第一之需求,不特戰爭時然,平常亦然.如人之身心,家庭不得平安,其能工作乎?即此一事,亦足以顯明主道之為人類之第一需求矣.耶穌曰:我留下平安給你們,我將我的平安賜給你們,我所賜的,不像世人所賜的.(約十四之二十七)兄姊們!吾主賜吾人如此屬靈屬天之完全平安,汝曾得享否?(太十一之二十九)若多求主恩,必多得主之平安矣.
\newline
\newline
(乙)(四節)在此表明主為吾人所成全之救贖,不特關於已過罪惡之得赦免,及未來身體之復活,以享永生.其最要者,為現在的,今生的,如此節所云:救我們脫離這罪惡的世代.曾蒙主血救贖的兄姊,主的救恩曾使汝之身心脫離這世代的罪惡,使汝不為此世的罪惡所攪所勝否?主耶穌至終為信徒向上主之祈禱說:我不求你叫他們離開世界,只求你保守他們脫離那惡者.(約十七之十五)脫離這罪惡的世代一語,尚有一義,即主之救恩,能救吾人免與此罪惡之世代同受審判,同遭沉亡,而作燭世之光.(腓二之十五)故吾人當聽此天來之聲曰:我民哪,你們要從那城出來,免得與他一同有罪,受他所受的災殃.(啟十八之四)
\newline
\newline
第二段　論偽師與異端(六至十節)
\newline
\newline
在本段中吾們應有三個問題
\newline
\newline
(甲)偽師為何要傳異端呢?
\newline
\newline
(乙)何為異端?
\newline
\newline
(丙)異端的結果怎麼樣?
\newline
\newline
異端的原始,是因「就人」與「求悅於人」.他要投合人,不惜將福音更改了!所以異端的產生,無論是由於科學的,哲學的,或心理學的,總總是由「人智」兩字而來.世人之智慧,是釘主於十字架；(哥前二之八)故異端也是將主再釘了十字架.如(歌羅西書二之二十至二十三)所說:異端似乎是有智慧,然而正用的時候,就都敗壞了,其實在克制肉體的情慾上,是毫無功效.
\newline
\newline
然而偽師因何要傳異端?還有一個原因,就是「環境」.「在提後四章三四節」保羅說:時候要到,人必厭棄純正的道理,並且掩耳不聽從真道；是以傳純正道真的人,是很難得人人喜悅的.由此,傳道欲避免「藐視」與「艱難」,遂將真道加以投合人心的言語,如在本書六章十二節保羅說的很明白,凡希圖外貌體面的人,都勉強你們受割禮,無非是怕自己為基督的十字架受逼迫,在五十一節保羅又說:弟兄們,我若仍舊傳割禮,為甚麼還受逼迫呢?若是這樣,那十字架討厭的地方就沒有了.(參看一之十)因此,傳道者在不知不覺之中就離棄了正道而入於異端；以致己受其害,而害他人.他們的結果 ── 就是被咒詛.(一之八)保羅說:無論是我們,是天上來的使者,如傳異端,皆受此刑；決不能因其學識,與地位而倖免.是以吾人處今日之教會,更當常存敬畏主之心而作工而施教,以免陷入魔之網羅.(箴二十九之二十五,哥後七之一)
\newline
\newline
第三段　保羅之蒙召及其所傳福音之來歷(十一至廿四節)
\newline
\newline
此段於第二講中已大概說過了,今可將福音二字略論一下,
\newline
\newline
由(一章六節),可知福音就是「基督之恩」,其有四義:
\newline
\newline
(一)因信基督而稱義,(二之十六)
\newline
\newline
(二)因受基督而受聖靈,(三之十四)
\newline
\newline
(三)因信基督而有新生命,(二之二十又四之六)
\newline
\newline
(四)因信基督而為上帝之後嗣,(四之七,三之二十六至二十九,又羅八之十六,十七)
\newline
\newline
歸納一句說:福音就是基督耶穌.如(馬可一之一)所說:上帝的兒子耶穌基督福音的起頭；由信耶穌基督,即能與耶穌基督同享一切在天,在地,屬身,屬靈,今世,來世之福樂,那就是福音的果效.讀(羅八之十七)請注意「和基督同作後嗣」一語,又看(哥前三之二十二)保羅說:或保羅或亞波羅,或磯法或世界,或生或死,或現今的事或將來的事,全是你們的!
\newline
\newline
但迦會中有人另傳一種福音,與他們所領受的不同.(一之九)他們在基督之外,又加上了割禮守日月節令等之儀文.於「信」之法,又添上了「行」之法,那偽師所傳的,就是 ── 「要得救,當信基督,亦當守律法.」因為這亦有片面的真理,所以那傳此異端之人,仍稱他們所傳的為「福音」；故保羅在(一之七)說那并不是福音,不過有些人攪擾你們,要把基督福音更改了.領袖們,吾人如不效法保羅及使徒們,由主所領受的完全傳於信徒:那恐怕吾們不堪為主的忠僕,所傳的也不是福音呀!(哥前十五之一至四,彼後一之十二至二十一,約一書一之一至四)
\newline
\newline
第二章
\newline
\newline
大旨 ── 論保羅之品格及其信仰
\newline
\newline
分段 ──
\newline
\newline
(一)保羅怎樣愛護信徒?(一至五節)
\newline
\newline
(二)保羅怎樣和處同勞?(六至十節)
\newline
\newline
(三)保羅怎樣忠心規友?(十一至十四節)
\newline
\newline
(四)保羅怎樣信仰基督?(十五至二十一節)
\newline
\newline
第一至第三段,在第二講中詳細論過了,茲僅擇一件前未論及之事論之如下 ──
\newline
\newline
(甲)保羅十四年後之上耶穌撒冷 ──
\newline
\newline
這是指保羅信主後之十四年而言,約在主後五十年；其事實見徒十五章.由該章二節,如此一班人之上耶路撒冷,是會眾的意思.然在本章二節,保羅則說:「是奉啟示而去的.」這兩方都是真的,眾人雖取決了叫這一班去,然而保羅也得了主的啟示；所以就他自己個人,可說是奉啟示而去.在此頗有教訓,教會今有一派,以為眾人取決,就是主旨；那是很危險的,你看釘耶穌在十字架,豈不是全體取決的麼?(太二十七之二十三至二十五)或以為那是未信主的人的意念,當然是有錯誤；比較教會自有不同.吾們試看當日耶路撒冷教會的領袖,如何叫四人與保羅同行潔禮,入殿守節,(徒二十一之二十三,二十四)以至發生絕大之風潮,由此保羅失其自由者四五年之久.此舉有無錯誤?即如彼得等之補選瑪提亞為使徒亦恐非主旨.是以吾人在群眾中,當追求徒十五之二十八所言之辦事精神,即「聖靈與群眾之同意.」更當愛慕安提阿教會之經歷,即群眾之意念行動,皆由聖靈之指導而發；專心禁食祈禱,以候主旨之啟示.(徒十三之二)所以個人言,亦當凡事直接祈禱上主.眾意如係合於主旨,汝必行之心安；否則求主,主必為汝開路,打消眾議,更變眾心,主必為汝成之,以免汝之行錯,徒勞無益.其他一方之錯誤,即在於獨自行動,不與眾商,此亦為不合,且其危險甚大.蓋汝即投身於一團體中,即當凡事與眾同心,方免紛亂.即或主有啟示亦當忍耐切禱,求主感動眾心,使眾明白,免失相愛聯絡之道,以致團體破裂.忍耐切禱,更可免誤以己意為主旨之危險,以致獨斷獨行.試觀本章八,九節保羅所云:主感彼為異邦使徒,眾領袖亦與之行右手相交之禮,為之祝福,與之合作；(十節)此為何等佳美之事,故遵行主旨者,未必一定要分別獨立,以致失於合一之道,以傷主心,而遺笑外人也.(羅十二之十八,詩一百三十三篇全)
\newline
\newline
保羅此次之上耶路撒冷,其關於後世教會者絕大即為爭得信仰之自由,將其列祖所不能負之軛,即 ── 割禮,儀文等由耶路撒冷教會眾領袖正式取決革除之.(徒十五之十,十一及二十八,二十九)又保羅此次之上耶路撒冷,既得安提約教會會眾之同意,又得勝靈之啟示,然亦帶提多同去,此事亦有關於辦事上之手續.蓋攜帶提多,即係攜帶憑證,亦為後代立證據.蓋提多為希利尼人,保羅引之歸信耶穌,然未嘗使之受割禮,此為保羅工作上之憑.又提多在此大會議中,仍無人強之受割禮,(三節)此又可為後代教會免異邦人受割禮之真實證據.是以吾人於教會之辦事上,亦當手續完備,證據確實,以免爭端.凡所經眾取決之事,亦當豋入記錄,以資參攷,而垂永久.
\newline
\newline
第四段　保羅怎樣信仰基督?(十五至二十一節)
\newline
\newline
本段聖經,本應歸入下段,因為那是屬於神道的論證,今讀這幾節聖經,便知道保羅的信仰如何了.
\newline
\newline
(十六節)本節已將第二大段的大旨標明了:就是稱義之道,是因信耶穌基督,──不是因行律法.故他毅然決然說:凡有血氣的,沒有一人因行律法稱義.這是全稱的斷定,世間並無其他救法,如徒四之十二所說:除他(耶穌)以外,別無拯救；因為在天下人間,沒有賜下別的名,我們可以靠著得救.到稱義的地位,祇有一條路,就是信之路.行之路,祇能引人到被咒之地位；(雅二之十,本書三之十至十一)保羅說:連我也信了基督耶穌,這是保羅將他本身的經歷,來作聖道的保證.我們若能於一切的教訓上,皆有如此經歷之實證,那教訓便成為至有能力的了.
\newline
\newline
(十七節)對於得救,是完全由信,不能有絲毫行之功在於其中之道,是一節很有力量證明,吾人試作圖以明之
\newline
\newline
罪人由信基督而得稱義(十六節)
\newline
\newline
(甲)　稱義 ─── 基督 ─── 罪人
\newline
\newline
然就本節(七節)之義分析之即──
\newline
\newline
(乙)──┐
\newline
\newline
若罪人因信賴基督,仍然至一罪人之地位,即可斷定基督是叫人犯罪的.如是基督的死,不僅是徒然,而且是引人犯罪的了:這樣看來,行之法,是永遠不能與信之法并存的,又可作圖以證明之 ──
\newline
\newline
(丁)　稱義 ── 行律法 ── 信基督 ── 罪人
\newline
\newline
丙圖是二－一節的解說；那就是人不必定由信基督,亦可由行律法達到稱義的地位.由是基督十架之功,并非獨一之救贖,如是其死徒然；
\newline
\newline
丁圖是證明罪人信基督,仍要再經行之法,乃能稱義；則基督流血救贖之功,為不完全.
\newline
\newline
故罪人欲得稱義,其唯一之法,即信基督.如甲圖.
\newline
\newline
(二十節)為本書第一之要節,在第一講中經已分析,又因其屬於第二大段故歸第四講再詳論之.
\newline
\newline


\newpage

\section{第1屆~港九培靈會~趙柳塘~第四講}
\label{sec:1341}
\textbf{第四講}
\newline
\newline
連結: \href{https://www.hkbibleconference.org/session-message/view/1341}{\texttt{https://www.hkbibleconference.org/session-message/view/1341}}
\newline
\newline
\hyperref[sec:1340]{< < < PREV SERMON < < <}
~
\hyperref[sec:toc]{[返目錄]}
~
\hyperref[sec:1342]{> > > NEXT SERMON > > >}
\newline
\newline
加拉太書 2:15-4:7
\newline
\begin{longtable}{cl}
\hline
\hline
章節 & 經文 (和合本修訂版)\\
\hline
2:15 & \begin{tabularx}{0.7\textwidth}{X} 我們生來就是猶太人，不是外邦罪人； \end{tabularx} \\ \\ \relax
2:16 & \begin{tabularx}{0.7\textwidth}{X} 可是我們知道，人稱義不是因律法的行為，而是因信耶穌基督，我們也信了基督耶穌，為要使我們因信基督稱義，不因律法的行為稱義，因為，凡血肉之軀沒有一個能因律法的行為稱義。 \end{tabularx} \\ \\ \relax
2:17 & \begin{tabularx}{0.7\textwidth}{X} 我們若求在基督裡稱義，自己卻還被視為罪人，那麼，基督是罪的用人嗎？絕對不是！ \end{tabularx} \\ \\ \relax
2:18 & \begin{tabularx}{0.7\textwidth}{X} 如果我重新建造我所拆毀的，這就證明自己是違犯律法的人。 \end{tabularx} \\ \\ \relax
2:19 & \begin{tabularx}{0.7\textwidth}{X} 我因律法而向律法死了，使我可以向神活著。我已經與基督同釘十字架， \end{tabularx} \\ \\ \relax
2:20 & \begin{tabularx}{0.7\textwidth}{X} 現在活著的不再是我，乃是基督在我裡面活著；並且我如今在肉身活著，是因信神的兒子而活；他是愛我，為我捨己。 \end{tabularx} \\ \\ \relax
2:21 & \begin{tabularx}{0.7\textwidth}{X} 我不廢掉神的恩；如果義是藉著律法而獲得，那麼基督就白白死了。 \end{tabularx} \\ \\ \relax
3:1 & \begin{tabularx}{0.7\textwidth}{X} 無知的加拉太人哪，耶穌基督釘十字架，已經活現在你們眼前，誰又迷惑了你們呢？ \end{tabularx} \\ \\ \relax
3:2 & \begin{tabularx}{0.7\textwidth}{X} 這是我惟一要問你們的：你們領受了聖靈，是因律法的行為或是因聽信福音呢？ \end{tabularx} \\ \\ \relax
3:3 & \begin{tabularx}{0.7\textwidth}{X} 你們既然以聖靈開始，如今竟要以肉身終結嗎？你們是這樣的無知嗎？ \end{tabularx} \\ \\ \relax
3:4 & \begin{tabularx}{0.7\textwidth}{X} 你們受這麼多的苦都是徒然的嗎？如果真是徒然的， \end{tabularx} \\ \\ \relax
3:5 & \begin{tabularx}{0.7\textwidth}{X} 那麼，神賜給你們聖靈，又在你們中間行異能，是因律法的行為或是因聽信福音呢？ \end{tabularx} \\ \\ \relax
3:6 & \begin{tabularx}{0.7\textwidth}{X} 正如亞伯拉罕「信了神，這就算他為義」。 \end{tabularx} \\ \\ \relax
3:7 & \begin{tabularx}{0.7\textwidth}{X} 所以，你們知道：有信心的人才是亞伯拉罕的子孫。 \end{tabularx} \\ \\ \relax
3:8 & \begin{tabularx}{0.7\textwidth}{X} 聖經既然預先看見神要使外邦人因信稱義，預先傳福音給亞伯拉罕，說：「萬國都必因你得福。」 \end{tabularx} \\ \\ \relax
3:9 & \begin{tabularx}{0.7\textwidth}{X} 可見，那有信心的人和有信心的亞伯拉罕一同得福。 \end{tabularx} \\ \\ \relax
3:10 & \begin{tabularx}{0.7\textwidth}{X} 凡出於律法的行為都是受詛咒的，因為經上記著：「凡不持守律法書上所記的一切而去行的，都是受詛咒的。」 \end{tabularx} \\ \\ \relax
3:11 & \begin{tabularx}{0.7\textwidth}{X} 沒有一個人靠著律法在神面前稱義，這是明顯的，因為經上說：「義人必因信得生。」 \end{tabularx} \\ \\ \relax
3:12 & \begin{tabularx}{0.7\textwidth}{X} 律法並不出於信，而是說：「行這些事的就必因此得生。」 \end{tabularx} \\ \\ \relax
3:13 & \begin{tabularx}{0.7\textwidth}{X} 既然基督為我們成了詛咒，就把我們從律法的詛咒中贖出來。因為經上記著：「凡掛在木頭上的都是受詛咒的。」 \end{tabularx} \\ \\ \relax
3:14 & \begin{tabularx}{0.7\textwidth}{X} 這是要使亞伯拉罕的福，因著基督耶穌臨到外邦人，使我們能因信得著所應許的聖靈。 \end{tabularx} \\ \\ \relax
3:15 & \begin{tabularx}{0.7\textwidth}{X} 弟兄們，我照著人的觀點說，人的遺囑一經確定，沒有人能廢棄或加增。 \end{tabularx} \\ \\ \relax
3:16 & \begin{tabularx}{0.7\textwidth}{X} 那些應許原是向亞伯拉罕和他後裔說的，並不是說「和眾後裔」，指許多人，而是說「和你那個後裔」，指一個人，就是基督。 \end{tabularx} \\ \\ \relax
3:17 & \begin{tabularx}{0.7\textwidth}{X} 我是這麼說，神預先所立的約不能被四百三十年以後的律法廢掉，使應許失效。 \end{tabularx} \\ \\ \relax
3:18 & \begin{tabularx}{0.7\textwidth}{X} 因為承受產業若是出於律法，就不再是出於應許；但神是憑著應許把產業賜給亞伯拉罕。 \end{tabularx} \\ \\ \relax
3:19 & \begin{tabularx}{0.7\textwidth}{X} 這樣說來，為甚麼要有律法呢？律法是為過犯的緣故而加上去的，等候那蒙應許的子孫來到才結束，是藉著天使經中保之手而設立的。 \end{tabularx} \\ \\ \relax
3:20 & \begin{tabularx}{0.7\textwidth}{X} 但中保本不是為單方設立的；神卻是一位。 \end{tabularx} \\ \\ \relax
3:21 & \begin{tabularx}{0.7\textwidth}{X} 這樣，律法是與神的應許對立嗎？絕對不是！如果律法的頒佈能使人得生命，義就誠然出於律法了。 \end{tabularx} \\ \\ \relax
3:22 & \begin{tabularx}{0.7\textwidth}{X} 但聖經把萬物都圈在罪裡，為要使因信耶穌基督而來的應許歸給信的人。 \end{tabularx} \\ \\ \relax
3:23 & \begin{tabularx}{0.7\textwidth}{X} 但這「信」還未來以前，我們被看守在律法之下，像被圈住，直到那將來的「信」顯明出來。 \end{tabularx} \\ \\ \relax
3:24 & \begin{tabularx}{0.7\textwidth}{X} 這樣，律法是我們的啟蒙教師，直到基督來了，好使我們因信稱義。 \end{tabularx} \\ \\ \relax
3:25 & \begin{tabularx}{0.7\textwidth}{X} 但這「信」既然來到，我們從此就不在啟蒙教師的手下了。 \end{tabularx} \\ \\ \relax
3:26 & \begin{tabularx}{0.7\textwidth}{X} 其實，你們藉著信，在基督耶穌裡都成為神的兒女。 \end{tabularx} \\ \\ \relax
3:27 & \begin{tabularx}{0.7\textwidth}{X} 你們凡受洗歸入基督的都披戴基督了： \end{tabularx} \\ \\ \relax
3:28 & \begin{tabularx}{0.7\textwidth}{X} 不再分猶太人或希臘人，不再分為奴的自主的，不再分男的女的，因為你們在基督耶穌裡都成為一了。 \end{tabularx} \\ \\ \relax
3:29 & \begin{tabularx}{0.7\textwidth}{X} 既然你們屬於基督，你們就是亞伯拉罕的子孫，是照著應許承受產業的了。 \end{tabularx} \\ \\ \relax
4:1 & \begin{tabularx}{0.7\textwidth}{X} 我說，雖然那承受產業的是整個產業的主人，但在未成年的時候卻與奴隸毫無分別， \end{tabularx} \\ \\ \relax
4:2 & \begin{tabularx}{0.7\textwidth}{X} 仍是在監護人和管家的手下，直等他父親預定的時候來到。 \end{tabularx} \\ \\ \relax
4:3 & \begin{tabularx}{0.7\textwidth}{X} 我們也是一樣，在未成年的時候，被世上粗淺的學說所奴役，也是如此。 \end{tabularx} \\ \\ \relax
4:4 & \begin{tabularx}{0.7\textwidth}{X} 等到時候成熟，神就差遣他的兒子，為女子所生，且生在律法之下， \end{tabularx} \\ \\ \relax
4:5 & \begin{tabularx}{0.7\textwidth}{X} 為要把律法之下的人贖出來，使我們獲得兒子的名分。 \end{tabularx} \\ \\ \relax
4:6 & \begin{tabularx}{0.7\textwidth}{X} 因為你們是兒子，神就差他兒子的靈進入我們的心，呼叫：「阿爸，父！」 \end{tabularx} \\ \\ \relax
4:7 & \begin{tabularx}{0.7\textwidth}{X} 可見，你不再是奴隸，而是兒子了，既然是兒子，就靠著神也成為後嗣了。 \end{tabularx} \\ \\
[1ex]
\hline
\hline
\end{longtable}
論救贖之大道(二章十五節至四章七節)
\newline
\newline
救恩之根基......耶穌基督
\newline
\newline
得救之方法......信耶穌基督
\newline
\newline
(一)由信耶穌基督而得稱義(二之十六,三之六至九)
\newline
\newline
(二)由信耶穌基督而得生命(二之二十)
\newline
\newline
(三)由信耶穌基督而受聖靈(三之二,十四)
\newline
\newline
(四)由信耶穌基督而為嗣子(三之二十六,二十九；四之六,七節)
\newline
\newline
第一層,在第三講已略論過了；今欲再詳論稱義,── 即得救之時期.
\newline
\newline
罪人信耶穌基督要經歷若干時間方能稱義得救呢?試看(羅十之九,十)說:你若口裡認耶穌為主,心裡信上帝叫他從死裡復活,就必得救.因為人心裡相信,就可以稱義；口裡承認,就可以得救.(在約翰五之二十四)耶穌說:我實實在在的告訴你們,那聽我的話,又信差我來者的,就有永生,不至於定罪,是已經出死入生了；
\newline
\newline
照這兩節聖經看來,人之「得救」與「出死入生,」并非要經歷一長久之時間,乃在於信主耶穌,接納之,承認之,便即時成就了.但救贖之完成即由稱義以至稱義(或稱成聖)而復活；(或肉身變化)乃要經歷時日,即由信之時以至主再臨之日.試讀羅五之一,二；凡人由信稱義以至稱義或成聖,又明明可見,有幾步的經歷:
\newline
\newline
一,稱義............(因信)(一節)
\newline
\newline
二,與上帝相和.........(因信)(一節)
\newline
\newline
三,進入恩典中.........(因信)(一節)
\newline
\newline
四,站在恩典中.........(因信)(一節)
\newline
\newline
五,歡歡喜喜的盼望神的榮耀.........(因信)(一節)
\newline
\newline
一二兩件可云是信徒得救的初步,三至五件是信徒前進的功夫；但可歸納一句說:一切都是因信而成就的,并無絲毫行為儀文等事在於其中,如(來十二之二)所云:我們當仰望為我們信心創始成終的耶穌.是以,吾們得救的工夫,是由信始,亦由信終；既由信耶穌基督而得稱義,即當進前到下層所論之事.
\newline
\newline
第二層　由信耶穌基督而得生命(二之二十)
\newline
\newline
(甲)這生命是由死得來的 ──
\newline
\newline
這一層,在第一講中已略說過了,但在此仍要詳論一下.在這一層,比前一層就大有進步了；稱義是完全由於上帝的恩,因為罪人由信耶穌而得救；他自己實在無「義」,不過上帝以耶穌之義,當作吾們的義.在這個「稱」字已是明顯了.(參看三之六)之「算為」二字.
\newline
\newline
到了第二這一步,并不是說「算為義,」乃是說信徒應當真確有耶穌基督復活豐盛的生命.這生命不是理想,乃是經歷；所以保羅在這一節之中有六個「我字.」這生命也不是「外表,」乃是「內心；」也不是改良或加增以前原有(為罪人時)的生命,乃是完完全全的一個屬天的新生命,(基督在我裡面活著)一句證明了.信徒在這裡兩步的經歷,於耶穌十字架的關係也大有分別 ──
\newline
\newline
1.稱義 ── 是在十字架下,仰望釘死復活的基督耶穌而得的.
\newline
\newline
2.新生命 ── 是在十字架上,與耶穌基督同死同復活而得的.
\newline
\newline
耶穌說一粒麥子不落在地裡死了,仍舊是一粒；若是死了,就結出許多子粒來.(約十二之二十四)這是證明生命與豐盛的生命,是由死而來；猶如本節(二之廿一節)保羅所云:我已經與基督同釘十架,現在活著的不再是我,乃是基督在我裡面活著.在(以弗所四之二十二至二十四)保羅說道,穿上新人,豈不說先要脫去舊人麼?
\newline
\newline
(乙)這新生命是由信得來的 ──
\newline
\newline
保羅說:我於今在肉身活著,是因信上帝的兒子而活,但就英文是說「由上帝子之信而活.」這樣說來,得救的人自己是完全無絲毫功勞份子的；因此,「信」亦不是吾們的,乃是「上帝子耶穌基督的信.」我儕得救之人,能常自覺主恩之如此完全深厚否?在(哥林多前書十五之十)保羅說:然而我今日成了何等人,是蒙上帝的恩纔成的；並且他所賜我的恩,不是徒然的,我比眾使徒格外勞苦,這原不是我,乃是上帝的恩與我同在.
\newline
\newline
(丙)這新生命也是由主的愛而來的 ──
\newline
\newline
本節第三段,保羅說:他是愛我,為我捨己.由此可知,如無主愛為我儕捨己,我儕的信也是徒然.是以救贖的恩,是由上帝之愛與耶穌之愛,(約三之十六)如無此愛,世人是絕望的.我儕傳道者,向淪亡之人,如無如此捨己之愛,何能使人生發信心,以接受主耶穌救贖之恩乎?約翰說:不是我們愛上帝,乃是上帝愛我們,差他的兒子,為我們的罪作了挽回祭,這就是愛了﹗(一書四之十)
\newline
\newline
第三段由信耶穌基督而受聖靈(三之二至十四) ──
\newline
\newline
保羅在(三之二)問得好,他說:我只要問你們這一件,你們受了聖靈,是因行律法,還是因聽信福音呢?故在上節,他隨即警告迦人說:你們既靠聖靈入門,如今還靠肉身成全麼?你們是這樣的無知麼?今日如此無知之人,教會中固屬不少；惟惜竟有人專立一會,以誘惑人,就是安息會即六日會,亦稱基督復臨安息會.那等就是靠聖靈入門,而靠肉身完全的.他門要一切由信稱義而得新生命的信徒,都要遵守禮拜六即猶太教之安息日,并諸多的儀文,然後才能得救.(我們當要防禦那等偽師!)
\newline
\newline
迦拉太教會的信徒,當然是由信福音而受了聖靈,認識上帝,更可說是被上帝所認識的.(四之九)且此聖靈進入其心,使之呼叫阿爸父.(四之六羅八之十六)此種宗教之經驗,迦會信徒許多人在聞信福音之時已經得著了,然亦有人還要歸回那懦弱無用的小學,情願再給他作奴僕；(四之九)真是無知,實又可憐；偽師之惡毒,實又可恨.
\newline
\newline
再讀(三之十四)可有兩層的意思 ──
\newline
\newline
(甲)叫亞伯拉罕的福,因基督耶穌可以臨到外邦人 ──
\newline
\newline
此段第一:是說到稱義之事,即如(三之八)說:上帝要叫外邦人因信稱義,就早已傳福音給亞伯拉罕,說,萬國都必因你得福.第二；說到嗣業,如(三之二十九)說:你們既屬乎基督,就是亞伯拉罕的後裔,是照著應許承受產業的了.(參看弗三之六)
\newline
\newline
由信耶穌基督所得著的稱義與嗣業,猶太人與異邦人是毫無分別.保羅在羅馬三之二十二也證明了,說:上帝的義,因信耶穌基督,加給一切相信的人,並沒有分別.而且由信耶穌基督所得著的,比較亞伯拉罕所得的不特無別,而且有加；如羅馬八之十七所說:和基督同作後嗣,我們將來的福,非特與亞伯拉罕同等；更與基督同等,那是何等奧妙奇特可感頌的事!
\newline
\newline
(乙)使我們因信得著所許的聖靈 ──
\newline
\newline
接受聖靈,井未嘗包括在亞伯拉罕應許之內,此應許惟獨由基督耶穌而得；四福音中充滿了這種論證.但其最要緊的,可說是(約翰七之三十八,三十九及一之三十三)諸君定要細讀之.吾們亦要記著耶穌的話說:我將實情告訴你們,我去是與你們有益的,我若不去,保惠師就不到你們這裡來；我若去,就差他來.(約十六之七)由是聖靈是因耶穌釘身,復活,升天而降世的.信徒可憑信心,即能接之居於己心.(約十四之十六,十七,多三之六,弗三之十七)
\newline
\newline
信徒何時方能接受聖靈呢?彼得說:你們各人要悔改,奉耶穌基督的名受洗,叫你們的罪得赦,就必領受所賜的聖靈.(徒二之三十八)兄姊們:你已悔改受洗了──(1)你曾確知罪已得赦否?(2)你已確知聖靈已入居汝心否? ── 汝當信耶穌基督,即能立刻得之!
\newline
\newline
第四段　由信耶穌基督而為嗣子(三之二十六,二十九；四之六,七)
\newline
\newline
子,或嗣子,是本書中一個極大的題目；由三之十五節至四章末節,可說是完全論證此題.試將此段聖經中之大意,提示如下──
\newline
\newline
(甲)信基督耶穌就是亞伯拉罕的後裔(三之二十九)
\newline
\newline
(乙)信基督耶穌乃是與應許之子 ── 以撒同等(四之二十八)
\newline
\newline
行律法乃是與血肉之子 ── 以實馬利同等(四之二十三)
\newline
\newline
(丙)信基督耶穌的都是上帝的兒子而嗣業(三之二十六；四之七)
\newline
\newline
本段是以亞伯拉罕為中心,論證基督內之人之地位如下 ──
\newline
\newline
(甲)信基督耶穌就是亞伯拉罕的後裔(三之二十九)
\newline
\newline
舊約的中心人物,最要的,是摩西與亞伯拉罕,至新約時代,猶太人常常誇耀的是「亞伯拉罕的子孫」一語.(太三之九,約八之三十三)他們自覺自己是亞伯拉罕的子孫,故其地位極高,而為上帝特選的族.今偽師之誘惑異邦信徒,或亦以此一事.然而保羅在此證明,世人之成為亞伯拉罕之後裔,並非由行律法,乃由信耶穌基督；而且亞伯拉罕之蒙選,亦非憑割禮,乃憑信心.(三之六)是以保羅在(羅四之九,十)說:亞伯拉罕因信稱義,是在受割禮之前.(二)割禮,是亞伯拉罕由信稱義的印證,這兩節書,說得極之明白,當要細讀之.
\newline
\newline
亞伯拉罕本身,還有一件重要的事,就是嗣業.亞伯拉罕在世時,主上帝曾親口向他應許諸多的福氣,而且一一關乎其後裔；故凡為其子孫之人,可算世人中至為有福之人.但信耶穌基督之人,保羅說:照著應許,承受產業；即亞伯拉罕之業,與猶太人同等.所以在下文保羅又論證 ──
\newline
\newline
(乙)信基督耶穌乃是與應許之子 ── 以撒同等.(四之二十八)
\newline
\newline
行律法乃是與血肉之子 ── 以實馬利同等.(四之二十三)
\newline
\newline
亞伯拉罕的兒子,也有兩等:一個是由上帝的應許而生的,就是以撒一人,惟獨他才能承繼上帝所應許之業.所以在(創世記十五之四,)耶和華曾親口向他說過了,你本身所生的,纔成為你的後嗣.此外還有一個,是由血肉所生的,就是以實馬利.在(創二十一之十)撒拉對亞伯拉罕說:你把這使女和他兒子趕出去,因為這使女的兒子,不可與我的兒子以撒一同承受產業.
\newline
\newline
再者,看亞伯拉罕於其將離世時,如何處理其家事,即知以撒與其他諸子之不同.(創二十五之五,六)說:亞伯拉罕將一切所有的都給了以撒,亞伯拉罕把財物分給他庶出的眾子,趁著自己還在世的時候,打發他們離開他的兒子以撒,往東方去了.
\newline
\newline
所以保羅在本書(四之二十一,二十二)說:你們這願意在律法以下的人,請告訴我,你們豈沒有聽見律法麼?因為律法上記著亞伯拉罕有兩個兒子；一個是使女生的,一個是自主之婦人生的.靠著律法的,那就是使女之子 ── 以實馬利,不是以撒；那是不能承繼嗣業的.故凡離棄信之法而從行之法者 ──(1)今生,即失去了一切屬靈之福.(四之十五)(2)將來,失去了一切屬天之嗣業,而且是受咒詛的,如保羅說:凡以行律法為本的,都是被咒詛的.(三之十)
\newline
\newline
(丙)信基督耶穌的都是上帝的兒子而嗣業(三之二十六,四之七)
\newline
\newline
保羅明明的說:信基督耶穌的都是上帝的兒子.既為兒子,上帝就差他兒子的靈進入你們的心,呼叫阿爸父,就靠著上帝為嗣.信徒們:我們為上帝的兒子,靠上帝而為後嗣,并不要客觀的憑證,世人的公認:乃是憑著主觀的憑證,就是 ── 主靈之在我心作證,那是無錯誤無更變不能移奪永遠的憑證.你信道以來,曾得著了麼?吾們既有這種有福的指望,就當如何潔淨自己,像他潔淨一樣,那就是約翰的勸勉.(一書三之三)保羅又說:我想現在的苦楚,若比起將來要顯於我們的榮耀,就不足介意了.又說:我們若盼望那所不見的,就必忍耐等候.(羅八之十八,二十五)
\newline
\newline
兄姊們:吾們緊記著約翰的話說:接待他的,就是信他名的人,他就賜他們權柄,作上帝的兒女.(一之十二)由此看來,作上帝的兒女,并非由於己之行為,乃由主之恩賜,主之權能,及吾人之(信)!
\newline
\newline


\newpage

\section{第1屆~港九培靈會~趙柳塘~第五講}
\label{sec:1342}
\textbf{第五講}
\newline
\newline
連結: \href{https://www.hkbibleconference.org/session-message/view/1342}{\texttt{https://www.hkbibleconference.org/session-message/view/1342}}
\newline
\newline
\hyperref[sec:1341]{< < < PREV SERMON < < <}
~
\hyperref[sec:toc]{[返目錄]}
~
\hyperref[sec:1343]{> > > NEXT SERMON > > >}
\newline
\newline
本書第二大段之分析
\newline
\newline
第三章
\newline
\newline
大旨 ── 論信之功效如何
\newline
\newline
分段 ──
\newline
\newline
(一)由於福音之信之功效如何(一至五節)
\newline
\newline
(二)亞伯拉罕信心之功效如何(六至十四節)
\newline
\newline
(三)應許與基督耶穌之關係如何(十五至二十節)
\newline
\newline
(四)律法與基督耶穌之關係如何(廿一至廿九節)
\newline
\newline
本段乃入於理論,即關於聖道之奧妙；由之可得知世人得救之方法如何,其要旨──即為得救之道由信,非由於行.第三章乃論其前半,即信之功效如何.茲分層研究之 ── 第一段由於福音之功效如何(一至五節)
\newline
\newline
本段保羅向迦拉太人有五次的質問,使其同憶信之一切好處,皆由於聽信福音而來.
\newline
\newline
(甲)質問一......因何離棄耶穌(一節)
\newline
\newline
(乙)質問二......如何得受聖靈(二節)
\newline
\newline
(丙)質問三......因何離棄聖靈之法(三節)
\newline
\newline
(丁)質問四......因何為耶穌受苦(四節)
\newline
\newline
(戊)質問五......我保羅之能力由何而得(五節)
\newline
\newline
保羅開端便說:無知的迦拉太人哪!(一節)其中表現無限的愛情與嘆惜.凡為父母之人,目睹其子女趨於敗壞死亡之途,誰能不生感嘆；況屬靈的父保羅呢?保羅傳福音與迦人,及眾教會,不獨在乎言語,也在乎權能和聖靈並充足的信心；(帖前一之五)不是靠著智慧委婉的言語,乃是用聖靈和大能的明證.(哥前二之四)猶如本章一節所說:耶穌基督釘十字架,已經活畫在你們眼前.所以他要同迦人說:誰迷惑了你們,使你們離棄了耶穌基督呢?這答案就是由於偽師的攪擾.你看保羅這樣真的工作,魔鬼尚能敗壞之,我們的工作,又當要如何小心去做呢?(哥前三之十至十三)偽徒們,當要如何的謹慎,去堅守他們所得的道呢?(帖後二之十五)
\newline
\newline
第二:保羅問他們在前如何接受聖靈?(二節)是在聽信我保羅所傳的福音而得,還是在跟從偽師遵守割禮之後呢?這是迦人據他們的經歷,可以回答保羅的.
\newline
\newline
第三:故此,保羅接著的進前向他們說:你們既靠聖靈入門,如今還靠肉身成全麼?(三節)照迦人的思想,以為他們現今能知前所未知,行前所未行；加以偽師之自誇,便感覺保羅前所傳者之不完全!自以為今是進步了,實乃歸回那懦弱無用的小學,(四之九至十一)從恩典中墜落了,(五之四)故保羅深深嘆惜說──你們是這樣的無知麼?
\newline
\newline
第四:保羅再將他們的一件經歷提醒之,使之覺悟；就是為耶穌基督福音之道,而曾受苦難.讀徒十三十四章,知保羅在迦拉太各地傳道,曾受了無數的困逼,信徒們當然也必與之同受困逼；所以他回歸安提阿之時,就堅固門徒的心,勸他們恆守所信的道.又說:我們進入上帝的國,必須經歷許多艱難!迦人既為福音受苦如此之多,今因欲避免逼迫與厭棄,(六之十二)便離棄了十字架純正福音之路,而從守割禮與儀文,豈不是前功盡廢麼?(四節)主耶穌基督的工場中,不知有多少是半路而廢,徒然受苦的.信徒們,亦不知有多少向來跑得很好,而一經試煉誘惑便中止了!是以保羅說:我只有一件事,就是忘記背後,努力面前的,向著標竿直跑,要得上帝在基督耶穌裡從上面召我來得的獎賞.(腓三之十三十四)保羅又說:我不以性命為念,也不看為寶貴,只要行完我的路程,成就我從主耶穌所領受的職事,證明上帝恩惠的福音.(徒二十之二十四)耶穌在十字架至終的話是:「成了」.(約十九之三十)我們離世見主之時,能否如主耶穌一樣的說成了?
\newline
\newline
第五:保羅再將他的經歷來向迦人作證,是以他問迦人說:那賜給你們聖靈又在你們中間行異能的,是因你們行律法呢?還是因你們聽信福音呢?今若以保羅之工作,經歷,比較偽師,則迦人應當即時覺悟；惟惜迦人仍然多人為偽師熱心,這因他們受了魔力的迷惑!你想這是何等可恨可怕的事呢?保羅向腓立比人說:凡是真實的,可敬的,公義的,清潔的,可愛的,有美名的；若有甚麼德行,若有甚麼稱讚,這些事你們都要思念!(四之八)耶穌說:你們當信我,我在父裡面,父在我裡面；即或不信,也當因我所作的事信我.(約十四之十一)由是可知我們傳道的經歷,於信徒信心的關係,是何等的重大!保羅工作的功效,既可為福音之信而作證；迦人亦可放心信從而無疑.
\newline
\newline
第二段　亞伯拉罕信心之功效如何(六至十四節)
\newline
\newline
本段是論證亞伯拉罕由上帝所得一切的好處,都是由信之法,并非由行律法之法；凡以行律法為本的,都是被咒詛的!(十節)其中最重要的一節,就是說,叫亞伯拉罕的福,因基督耶穌,可以臨到外邦人,使我們因信得著所應許的聖靈；(十四節)而且在基督裡的人所得的,比較亞伯拉罕更多,那就是「更得著所應許的聖靈.」
\newline
\newline
亞伯拉罕就是以色列所誇耀的一人,也是偽師所傳的一人,── 就是直接由上帝接受割禮者.他所憑的是什麼呢?豈不仍是信心麼.保羅說:亞伯拉罕信上帝,這就算為他的義.(六節)
\newline
\newline
在本段的論證,是很合於論理學的規律的.(選輯)
\newline
\newline
一,亞伯拉罕 ── 信上帝 ── 這就算為他的義(六節)
\newline
\newline
二,以信為本的人 ── 就是亞伯拉罕的子孫(七節)
\newline
\newline
三,以信為本的人 ── 和有信心的(亞伯拉罕) ── 一同得福稱義(九節)
\newline
\newline
保羅又說:上帝欲使後人因稱義,故先施行於亞伯拉罕之身,早傳福音給亞伯拉罕說,萬國都必因你得福.(八節)當注意「萬國」二字,即未受割禮之異人,亦包括於其內了.
\newline
\newline
後段保羅又論證律法:
\newline
\newline
1.凡以行律法為本的, ── 都是被咒詛的.(十節)
\newline
\newline
2.基督為我們受了咒詛, ── 救贖出我們脫離律法的咒詛.(十三節)
\newline
\newline
由第一義知萬人都被咒詛,因為,在羅馬(二十之十二)說:凡在律法以下犯了罪的,也必按律法受審判.(二之十四十五)說:沒有律法的外邦人,自己就是自己的律法；律法的功用刻在他們心裡,他們是非之心同作見證,並且他們的思念互相較量,或以為是或以為非.即如本章二十二節說聖經把眾人都圈在罪裡,使所應許的福,因信耶穌督歸給那信的人；因為基督耶穌,掛在木頭上,受了律法咒詛,(十三節)將律法下的人贖了出來.也本無罪而為罪人受了罪惡的報應；就是死也將良心審判之下的人救了出來.
\newline
\newline
在本段保羅亦將世人決定不能由行律法之法得以稱義,說得明白 ──
\newline
\newline
1.凡不照律法書上所記一切之事去行的,就被咒詛； ── 故沒有一個人靠著律法在上帝面前稱義.(十,十一節)
\newline
\newline
2.聖經說:義人必「因信」得生；由此可知上帝所賜世人得生之路 ── 是信.
\newline
\newline
3.律法說:行這些事的,就必因此活著；(十二節) ── 由是上帝的律法,使人得生的路, ── 是行.
\newline
\newline
4.由行律法之法,世人都死了, ── 然由信之法,世人可由死得生.即如羅馬五章所論,基督耶穌的恩典,是勝過了律法的公義.
\newline
\newline
第三段應許與基督耶穌之關係如何(十五至二十節)
\newline
\newline
本段乃論亞伯拉罕的應許,是成就在其子孫基督的身上.(十六節)茲分層略論如下:
\newline
\newline
(一)應許律法的比較
\newline
\newline
凡應許都是白賜的.譬如吾人應許以一物賜贈與人,乃是出於應許者的自動與樂意.即如亞伯拉罕所得之一切福澤,乃是由應許,即上帝樂意的施賜.當日猶太人不服信之法即不行白白稱義之法.然而,不知亞伯拉罕之法,亦即福音之法.所以保羅問說:我們的祖宗亞伯拉罕,憑著肉體得了甚麼呢?(羅四章一節)然而,上帝應許,亞伯拉罕和他後裔,必得承受世界,不是因律法乃是因信而得的義.(羅四章十三節)然而,上帝應許,亞伯拉罕和他後裔,必得承受世界,不是因律法乃是因信而得的義.(羅四章十三節)所以人得為後嗣,是本乎信.因此,就屬乎恩,叫應許定然歸給一切後裔；不但歸給那屬乎律法的,也歸給那效法亞伯拉罕之信的.(十六節)但猶太人不憑著信心求只憑著行為求；他們正跌在那絆腳石上.(羅九章三十二節)
\newline
\newline
律法乃是絕對的,不是樂意的施賜；乃是強權的命令,在(希伯來十二之十八至二十九)已描寫得明白.論到立律法時的情形,一句話包括了,就是恐懼戰兢(二十一節)凡在律法下,歷代之人都是如此,恐懼戰兢!
\newline
\newline
再看到律法與應許相隔時間,約有四百三十年,那就是由亞伯拉罕蒙召起計,至出伊及止；然後在前四百三十年之應許,未嘗因以後的律法而廢掉,叫應許歸於虛空,(七節)也不能以人意更改,即如在十五節說.人的文約若已經立定了,就沒有能廢棄或加增；這就是證明應許的效用,是更大更可靠過律法.今偽師們,叫迦拉太的信徒,離棄了基督耶穌,永不更改可靠的應許；來信靠那使人不能遵守,反受咒詛的律法,那是何等可憐可惜的事!
\newline
\newline
(二)應許之成就
\newline
\newline
主上帝所賜亞伯拉罕的應許,不只及其本身；乃是包括他的子孫.(十五節)應許的事件,不只承受產業,更要施福異邦萬民.(八節)然而在於何時及在誰的身上方成就這應許呢?保羅說:上帝並不是說.眾子孫,指著許多人,乃是你那一個子孫,指著一個人,就是基督.(十六節)這是證到(創世記十三之十五及十七之七,八)在這幾節文中,都是用一「裔」字,但在英文不用SEEDS而用SEED,其分別只有一個s.現經保羅說明了,方知那小小一個字母,有這重大的關係；使吾們讀經的人,當要如何小心注意,即一點,一畫,一字,一語,都不可忽略更改.由此一節,又可知耶穌基督,實為聖經中唯一之中心人物；創世之前,萬代之後,一切之事事物物,莫不由耶穌基督,而歸於耶穌.
\newline
\newline
論到應許,乃是上帝由恩與愛,直接賜與亞伯拉罕的；但律法乃為過犯添上的；(十九節)他的用處,就在保守眾人,以待應許的子孫(即耶穌基督)來到；而且律法是藉天使,託中保,經了兩種的手續,那又是個律法不及應許的憑據.
\newline
\newline
總之保羅要那跟從偽師,倚靠律法的信徒,看出律法與福音,行與信之分別,恢復他們始初的信.
\newline
\newline
第四段律法與基督耶穌之關係如何(二十一至二十九節)
\newline
\newline
耶穌曾說:莫想我來要廢掉律法和先知；我來不是要廢掉,乃是要成全.我實在告訴你們,就是到天地都廢去了,律法的一點一畫也不能廢去,都要成全.(太五之十七十八)由是在耶穌基督裡的人,不是廢掉律法,乃是因為耶穌基督,已為世人成全了律法；故得以自由,不再受律法的束縛.第一個證明是在(羅馬七章),保羅說；我們藉著基督的身體,在律法上也是死了；叫你們歸於別人,就是歸於那從死裡復活的,叫我們結果子給上帝.但我們既然在捆我們的律法上死了,現今就脫離了律法,叫我們服事主,要按著心靈的新樣,不按著儀文的舊樣.(四節六節)(第二)保羅說:耶穌要救律法下的人,因而生於律法之下,(本書四之四,五)而且親身掛在木頭,為我們受了咒詛,救贖出我們脫離律法的咒詛.(三之十三)這樣,我們信徒的自由,不是打破律法不守；律法的自由,乃是耶穌為我們受了咒詛,出了贖價而得自由；這是何等可誇的事!因何偽師們,尚要將迦人放在律法之下呢?在本段保羅推重律法,是更加確定了基督耶穌救恩的價值.律法的地位,到底是怎樣呢?
\newline
\newline
(一)律法定人罪使人歸基督(二十二節)
\newline
\newline
(二)律法保守選民以待主耶穌降臨(二十三節)
\newline
\newline
(三)律法是訓蒙師引我們(律法下之人)到基督(二十四節)
\newline
\newline
由這三義看來,律法并非阻人歸信基督耶穌,只有助人歸信耶穌；無論異邦,猶太人,皆該由律法而歸信耶穌,因何律法反阻礙律法下之人到耶穌處,是何故呢?這就是因律法下之人,未曾因律法而認識己罪,反以律法為誇,那就是律下之人,不能信主的一個大弊病.信徒們!你知否有人因多聽福音而更剛愎,其弊亦在知而不行.
\newline
\newline
由以上三義看來,律法確是一條引人到耶穌救恩的路；惟惜多人不行在其中,只是站在其中誇口,故律法之於他們,毫無益處.
\newline
\newline
但由上文看來,律法本是一種臨時暫行的東西,(二十五節)其至終之目的,乃是──基督的救恩.迦人既歸信了耶穌,已得為上帝之子,照應許已承受了產業,(二十六至二十九)豈非律法之目的已達?因何偽師,仍要引迦人再行求救之路呢?保羅一語揭破了,說他們偽師,不過想要迦人為他們熱心而已!(四之十七)
\newline
\newline
「律法既是導人歸信耶穌基督的,故耶穌來,即是成就了律法」.
\newline
\newline
以目的言,固屬如此；以其經歷言,亦莫不如此.
\newline
\newline
第四章
\newline
\newline
大旨 ─── 子,奴.
\newline
\newline
分段 ──
\newline
\newline
(一)世人如何成為上帝之子(一至七節)
\newline
\newline
(二)保羅之警告與勸勉(八至二十)
\newline
\newline
(三)子奴之論證......(二十一至二十一節)
\newline
\newline
本章之大旨,即為一子字；此亦為本書之一重要問題.凡在律法下者為奴；在基督內者方為子,即由基督之救贖,而得此為子之地位.茲論之如下:
\newline
\newline
第一段世人如何成為上帝之子(一至七節)
\newline
\newline
本段的關鍵,全在(四節),人類之希望,地位之升高,皆在此一節.
\newline
\newline
(甲)上帝差遣他的兒子降世之前, ── 世人的地位怎麼樣呢?即以選民而論,亦不過一,孩童；二,奴僕；(一節)三,在師傅和管家的手下；(二節)四,受管於於世俗小學之下.(三節)
\newline
\newline
── 時候滿足 ──
\newline
\newline
(乙)上帝差遣他的兒子降世之後, ── 世人的地位怎麼樣呢?
\newline
\newline
一,得於律法下贖出；二,得著兒子的名份；(五節)三,得接受上帝子耶穌之靈於心(六節)；四,能以呼叫阿爸父.(六節)
\newline
\newline
如此便脫離了奴僕之地位,而為嗣子.(七節)
\newline
\newline
第二段保羅之警告與勸勉(八至二十節)
\newline
\newline
保羅於救贖原理之證論將終,應該論到迦人本身之事,所以有此──
\newline
\newline
(一)警告(八至十一節)
\newline
\newline
保羅為何有此警告呢?在(十一節)已說明了.他說:我為你們害怕,惟恐我在你們身上是枉費了工夫.枉費了工夫,這是何等傷心的話?傳道人不願信徒們使他的工夫枉費,然而,我們能不使聖靈在吾們身上枉費工夫麼?
\newline
\newline
「認識」兩字,就是保羅在於本段警告的主旨其意有三:── 一,世人由認識上帝罪人的地位,到 ── 二,認識了上帝(稱義)地位,而到 ── 三,被上帝所認識(蒙悅)的地位.
\newline
\newline
然後怎樣呢? ── 退步了,還要歸回那懦弱無用的小學,情願再給他作奴僕,(九節)這真是可憐可惜阿!
\newline
\newline
第一級,是世人未得救前之情形 ── 「不認識上帝.」
\newline
\newline
然而,世人以這五個字是很平常而無關重要的,常常輕易的說了出來.孰不知人類心思之智愚,人格之尊卑,社會之益損,心靈之光暗,靈魂之生死,全繫於此一語.因為聖經說:敬畏耶和華是知識的開端.(箴一之七)在羅馬一章說到人格墮落的原因,就在「不認識上帝」與「故意棄絕上帝」.(十八至三十二節)所以說:他們既然故意不認識上帝,上帝就任憑他們存邪僻的心,行那些不合理的事.(二十八節)在(迦四之八)說不認識上帝的人,是給那些本來不是神的作奴僕就是偶像與私慾的奴僕.這等人是絕對沒有自由,罪惡不能自拔的；惟有真理方能使人得釋於罪,得以自由.(八之三－二)由此可知,唯一救世的方法,就是 ── 使人認識上帝, ── 明白真理.
\newline
\newline
第二級,是罪人由喪失沉淪的地位,到了稱義得救的地位,是由什麼方法呢?不過是──認識上帝四字而已.所以耶穌說:那聽我話,又信差我來者的,就有永生,不至於定罪,是已經出死入生了.(約五之二十四)約翰更直白的說:人有了上帝的兒子,就有生命；沒有上帝的兒子,就沒有生命.(壹書五之十二)
\newline
\newline
第三級,是上帝所認識,這是最為寶貴的,這級與前就大不相同了!
\newline
\newline
因為認識上帝,為一事；上帝所認識,又是一事.例如吾人由報紙雜誌等,認識了不少偉大的人物,可是偉大的人物,未嘗認識我；如能得偉大人物認識了我,我的地位,生活,便大大的不同了.信徒們,吾們認識上帝的人,便為此全能天地之主上帝所認識的,你常思想此語之寶貴否?感謝上帝!迦拉太的信徒,有許多是為上帝所認識了的；今因偽師的誘惑,便無知無識的,甘心情願,再歸回那懦弱無用的小學謹守日子,月份,節期,年分,割禮；為這一切的儀文,律法,作奴僕,(九十節)竟棄上帝子之地位而不顧；不特令保羅傷心,就令後世讀者,亦為之同聲一嘆!
\newline
\newline
警告完了,保羅又切心的
\newline
\newline
(二)勸勉(十二至二十節)其意有四:
\newline
\newline
1.勸眾當效己；(十二節)　2.勸眾勿失初心；(十三至十六節)　3.勸眾當防偽師；(十七,十八節)　4.保羅之切愛信徒.(十九,廿節)
\newline
\newline
(甲)勸眾當效己(十二節)
\newline
\newline
「你們要像我一樣.」這一句話,有幾多領袖們敢說?而又能說之無愧呢!然而,我們領袖當知,即或不言,信徒們亦自然要留心觀察效法你的；因為我們的地位,就是代表耶穌,如同耶穌代表上帝一般.(參看一之二十,哥後三之三,約十三之三十五)昔有一牧師說:我們任職之時,如肯向會眾宣佈說:「你們要像我一樣,」他就凡事不能不小心謹慎勉力.在十二節中還有一句重要的話,就是「我也像你們一樣,」這是保羅指到二章中所說的事,除去猶太人一切規例,自由與未受割禮的來往飲食.(參看哥前九之二十至二十三)(像你們一樣),是導人成聖,使人進步的一條要理.我們領袖,應該高高站著,凡事使人學效；然亦當降到要人同等,表現他的同情,然後方能令人觀感.若是高高的站著,叫人到你那裡來,是不能的!這就是先聖先賢所以失敗的原因.但耶穌基督,要為世人的救主,他便要降世為人,而且要做人類至卑下的一種人 ── 奴僕謙卑順服,甚而死於十架.(腓二之六至八)這樣世人中,不獨無人與他有隔膜,無人不能向他發生同情,而接受他的救恩；且使人勉勵的去效法他.
\newline
\newline
信徒們,耶穌要救世人,所以他到我們罪人的世界來；歐美的西教士要救世界黑暗的民族,也是到他們那裡去,學他們的語言,觀察他們的風俗習慣,住在他們中間,方能作那救人的工夫.今我們想救鄉村的人,當然要到鄉村去.引人成聖的人們!你們不只是應有區別聖潔的人生,亦當有與罪世表其憐憫的同情,這因為你們,不當與這污穢邪惡的社會隔絕,當與之交處,作拯救的工作,而不自染污穢.如保羅教訓腓立比人說,你們無可指摘,誠實無偽,在這彎曲悖謬的世代,作上帝無瑕疪的兒女；你們顯在這世代中,好像光明照耀.(二之十五)
\newline
\newline
信徒們!基督徒的聖潔,并不是佛教的出世自潔,深山修煉；乃是人世不污,使人仰慕基督教的救贖；不是像儒教見孺子之墮井而只生憐憫之心,乃是入井救人.所以耶穌不是在天上叫人上到他那裡領受救恩,乃是降臨世界,貢獻救恩.保羅這「我也像你們一樣」一語,是很能感動迦拉太人的.
\newline
\newline
(乙)勸眾勿失初心(十三至十六節)
\newline
\newline
始初的愛,是最為寶貴而又可紀念的.(十三十四節)是保羅使迦人追念其當日愛彼之心如何,所以他說:「你一點沒有虧負我,」(十二節)而且願意以己目代保羅之目,接待他如同上帝的使者,如同基督耶穌.然而,至於斯時,當日所誇之福已歸烏有了!并將那以真理告訴他們的保羅,視如仇敵了!迦人情感如此之變遷,大大教訓於吾人.
\newline
\newline
1. 「所誇的福」信徒熱心傳道,人以之為教會之福；信徒亦嘗以己之熱心為傳道人之福.然而保羅說信徒熱心愛傳道人,愛教會,那就是信徒所誇的福,其理何在?這因信徒有主愛在心,然後能愛傳道人,與教會.今日之信徒,如能明白此理,教會自然興盛；因為許多信徒,以為不熱心捐助,聚集,教會與傳道,必受了無限的損失與福氣；然而,保羅說,那是自己失福.
\newline
\newline
(丙)勸眾當防偽師(十七,十八)
\newline
\newline
偽師也是很熱心的講道,或是更為盡力的.然而他們的目的錯了；因為他們不懷好意,要離間信徒,使人為他們熱心.這樣,他們的熱心與盡力,皆為可詛的；因為這是利己的熱心與盡力.同勞們!我們作工的熱心與盡力,具何目的呢?如此之人,保羅勸信徒防之；今我們亦當常常自審,防備自己,免入了偽師之途.
\newline
\newline
所以保羅終結一句說:在善事上常用熱心待人,原是好的.(十八節)
\newline
\newline
(丁)保羅的切愛(十九,二十節)
\newline
\newline
本文在第一講中,第一層中已略略論過了.總之,傳道人於信徒的工作有三:就是生,義,教.始以福音之,維以主道養之,終以聖靈之言教之；此三法已盡,而信徒仍再退步離道如迦人者,吾人即當為之再受生產之苦,直到基督成形在其心中.吾人任此傳道之職,其能有此忍耐與愛心否?信徒們,你之靈程與生命,莫不日夜繫於傳道者之心；你知之否?你能使你之傳道者無憂慮否?在(二十節)保羅之心,戰爭無已,憂慮滿衷；因為他不知道迦人究竟如何,他巴不得改換口氣,那是他的切望,就是他於絕望中在信裡所生的希望.
\newline
\newline
第四段子奴之論證(二十一至三十一節)
\newline
\newline
本段是保羅引用一段舊約的聖經,來論證新約信徒之位置,即是以經解經之法,或者以為此法是太穿鑿附會了；然而,保羅亦應用此法.今有許多人對此法表示不滿,至於批評反對,那又何妨呢,本段無時間詳加討論,惟作兩表以明之:
\newline
\newline
兩婦之預表(4:24)
\newline
\newline
預表狀況
\newline
\newline
婢女　 　┌　舊約　　　　　　　  ┐　所生之子為奴(:││24)
\newline
\newline
主母　 　┌　新約　　　　　　　　┐　所生者為子
\newline
\newline
二子之預表(4:22)
\newline
\newline
狀況預表
\newline
\newline
┌　使女所生(:2)　┐　1. 舊人
\newline
\newline
┌　主婦所生(:23) ┐　1. 新人
\newline
\newline


\newpage

\section{第1屆~港九培靈會~趙柳塘~第六講}
\label{sec:1343}
\textbf{第六講}
\newline
\newline
連結: \href{https://www.hkbibleconference.org/session-message/view/1343}{\texttt{https://www.hkbibleconference.org/session-message/view/1343}}
\newline
\newline
\hyperref[sec:1342]{< < < PREV SERMON < < <}
~
\hyperref[sec:toc]{[返目錄]}
~
\hyperref[sec:1344]{> > > NEXT SERMON > > >}
\newline
\newline
加拉太書 5:5-5:6
\newline
\begin{longtable}{cl}
\hline
\hline
章節 & 經文 (和合本修訂版)\\
\hline
5:5 & \begin{tabularx}{0.7\textwidth}{X} 至於我們，我們是靠著聖靈，憑著信心，等候所盼望的義。 \end{tabularx} \\ \\ \relax
5:6 & \begin{tabularx}{0.7\textwidth}{X} 因為在基督耶穌裡，受割禮不受割禮都沒有功效，惟獨使人發出仁愛的信心才有功效。 \end{tabularx} \\ \\
[1ex]
\hline
\hline
\end{longtable}
大旨......信徒當如何實行主道(五六節)
\newline
\newline
那就是得救了的人,應當怎樣?是關於信徒日日的行為,屬於實踐的；故當詳細論之.因為此兩章的教訓甚多,且要攷本講專論第五章,第七講專論第六章.
\newline
\newline
第五章
\newline
\newline
大旨......得贖之人於一己之心靈應當如何
\newline
\newline
分段 ──
\newline
\newline
(一)當堅守基督所賜之自由(一至十五節)
\newline
\newline
(二)當賴聖靈以滅肉體之情慾(十六至二十六節)
\newline
\newline
第一段當堅守基督所賜之自由(一至十五節)
\newline
\newline
(一)基督所賜的自由怎樣
\newline
\newline
(甲)脫於律法(三之十三,四之五)
\newline
\newline
此意在前已詳論過了,但在此保羅再說:──
\newline
\newline
一,受割禮就與你們無益(二節)
\newline
\newline
二,受割禮是欠著行全律法的債(三節)
\newline
\newline
因為凡在律法以下犯了罪的,也必按律法受審判.(羅二之十二)信徒們,既賴基督耶穌得脫於律法,而今又要自置於律法之下,就要自負律法的審判.然而,「全律法」三字,誰能擔受得起呢?但今之言守安息日......等之人,不過重言舊約律法中的一二事,於全律法究竟是不顧的,那又何益?
\newline
\newline
三,靠律法稱義的,是與基督隔絕,從恩典中墜落了；(四節)引人守割禮行律法的偽師,不是助人前進,乃是使人墮落!
\newline
\newline
因此,保羅毅然決然的斷定說:在基督耶穌裡受割禮不受割禮,全無功效；惟獨使人生發仁愛的信心,纔有功效.(六節)信基督的人,不重在儀文的外貌；乃重在仁愛的信心.
\newline
\newline
以上是向法下的猶太人說的.還有一方面,就是基督所賜的自由,是:
\newline
\newline
(乙)脫於罪惡
\newline
\newline
(在一之四)保羅曾說:基督的救恩,不僅使吾人得以稱義,而且脫離這罪惡的世代,就是勝罪.世人自從亞當以來,人類都在罪惡的權勢之下(羅五之十四)受罪惡的管轄,不能不去犯罪.然而,因為賜生命聖靈的律,在基督耶穌裡釋放了我,使我脫離罪和死的律了!(羅八之二)所以說在基督耶穌裡的,就不定罪；這是說由於聖靈與耶穌之法,吾人可以脫於罪惡之權勢,得以自由,可不犯罪.「不定罪」之意,即此.(在二之二十)說:吾們既與基督同釘十架,便得基督在吾們裡面活著；有了基督耶穌如此之豐盛生命,罪惡的能力在於吾們的身便消滅了.由是可知基督徒的得勝,是在能勝罪的,并非避罪；是積極的,并非消極的.
\newline
\newline
(二)享此自由者應怎樣
\newline
\newline
1. 當堅立(一節)
\newline
\newline
保羅說:基督釋放了我們,叫我們得以自由；所以要「站立得穩」.「站立得穩」,是一句很重要的話.吾們為信徒,不貴於在得主恩之多；乃在於既得主恩而能保守.如隨得隨失,或是忽冷忽熱,那又有何用處呢?站立得穩於熱心前進之信徒,尤當注意；所以保羅訓哥林多人說:自己以為站得穩的,須要謹慎,免得跌倒.(前十之十二)(在十六之十三)又說:你們務要儆醒,在真道上站立得穩.信徒堅立之法,其要有三:── 1.避　2.追　3.深
\newline
\newline
1. 「避」,就是逃避罪惡:── 2.追,就是追求公義,信德,仁愛,和平；(提後二之二十二)── 3.深,就要深居於基督,而且要在他裡面生根.(西二之六七)
\newline
\newline
2. 勿縱慾(十三節)
\newline
\newline
保羅說:弟兄們,......吾們由耶穌流血而得救贖的人,不當再為自己活,乃當為替我們死而復活的主活.(哥後五之十四十五)可惜教會中,常有以基督所賜之自由,而作縱慾之機會,藉著自由,遮蓋惡毒.(彼前二之十六)他們的結局是怎樣呢?由(民十一章)以色列人在基博羅哈他瓦求食鵪鶉一事,便可得著狠要的教訓.以色列人求鵪鶉,已是縱慾；主大發怒,摩西亦不喜悅；(十節)至終主在怒中,應允了他們祈求,大賜鵪鶉,他們便大縱其慾,聖經說:百姓起來,終日終夜,並次日一整天捕取鵪鶁,至少的也取了十賀梅珥,為自己擺列在營的四圍.肉在他們牙齒之間,尚未嚼爛,耶和華的怒氣,就向他們發作,用最重的災殃,擊殺了他們.那地方便叫作基博羅哈他瓦,因為他們在那裡葬埋那起貪慾之心的人.(三十二至三十四節)
\newline
\newline
基博羅哈他瓦,意譯就是貪慾之人的墳墓,這豈不是一個很可紀念的地方麼!世界就是一個大大的情慾的墳墓,世人故當逃脫這世界情慾敗壞.然而有人,在主十架之後,仍不免還走歸基博羅哈他瓦去!猶大就是其中的人；他用貪財之心,為自己建造了一個又高又大,舉世可見的墳墓,向信徒日夜的警告.與舊約羅得妻的鹽柱,(創十九之二十六)遙遙相應；故保羅有(十六至二十六)滅慾的大文章.
\newline
\newline
3.相服役(十三至十五節)
\newline
\newline
保羅說:不要以自由當作放縱情慾的機會,總要用愛心互相服事.(十三節)這兩事是相對的,因為吾們得了自由的人如不將主所賜諸恩,用以互相服事,難免用以互相嫉忌,紛爭；觀哥林多教會之經歷,便是吾們的鑑戒.保羅說:他們在主裡面,凡事富足,口才知識都全備；在恩賜上,也沒有一樣不及別人的.(前一之四五,七,)然而,哥林多教會是怎樣的情景呢?論紛爭,結黨,在聖經中他們可算第一了!
\newline
\newline
(前一之十一又三之三,四)今日教會,不為主傳道,推廣福音,以愛事人；惟求一己會的完備與富足；故此腐化了!死海的水,為何鹹苦不能飲呢?因為牠惟知接受,而不轉施.加利利海,是有受有施,故他的水便不同了.所以保羅說:你們總要以愛心互相服事,這是一句真確竭力的囑託.(在十四節)他又說:全律法都包在愛人如己這一句話之內,如能以愛互相服事,便是盡了律法,何必還要守割禮呢.(十五節)是一句至終至要的警告,說:你們要謹慎,若相咬相吞,只怕要彼此消滅了.在此,吾們當注意彼此兩字,今日主之教會在中華如此之多,然如此衰落,恐怕就是這個原因罷!
\newline
\newline
第二段當賴聖靈以滅肉體之情慾(十六至二十六節)
\newline
\newline
第五章分兩段,前段重在一「守」,本段則重在一「賴」；此兩字即為信徒靈程之大要.因為信徒日日靈修的功夫,不外保守所得的,進取未得的.進取之法如何呢?保羅不是說以己的力量,乃是說要賴聖靈的能力；在開首一節便說:你們當順著聖靈而行,就不放縱肉體的情慾了.(十六節)因為阻礙吾們靈程前進的,就是在自己的情慾；這事解決了,一切也就解決.信徒若要戰勝情慾,就必須 ──
\newline
\newline
一,順著聖靈而行 ── 就不放縱肉體的情慾(十六節)
\newline
\newline
二,被聖靈引導 ── 就不在律法以下(十八節)
\newline
\newline
由上文看來,信徒得勝情慾,以及不受律法的攪擾,不外順著聖靈的引導.然得勝情慾的人又當怎樣呢?就當:── 賴聖靈而生 ── 遵靈而行(二十五節文譯)
\newline
\newline
你看,這是何等容易的事,然而,那可咒的偽師,因圖己利,不惜以他們自己,及其祖宗所不能負的軛,加在迦人的頸項上!(參看徒十五之十)!唉,這是何等的罪過!今吾們再將本段分析 ──
\newline
\newline
(一)論情慾
\newline
\newline
(二)論聖靈
\newline
\newline
(一)論情慾 ── (甲)性質　　(乙)類別
\newline
\newline
(甲)情慾的性質 ──
\newline
\newline
1. 其數無限量
\newline
\newline
2. 發展極迅速
\newline
\newline
3. 情形極黑暗
\newline
\newline
夫情慾之數,在(本書五之十九至二十一節)官譯為十五件,文譯則為十七；在羅馬一之二十九至三十二,官譯為二十件,文譯為二十二,其數較聖靈之數多至兩倍.然,至於今日的罪惡,則更為日日加多,須另造新名詞以名之方可.然聖靈之果不多,僅有九件,但以一果,亦足以消滅情慾之多數行為；這是光明得勝黑暗,恩典得勝罪惡的明證.
\newline
\newline
論到情慾的發展,真是迅速可怕.讀創世記一至四章,知上帝造人,而亞當犯罪,而該隱殺弟,以自驕,(十七節)復仇,(二十四節)等等的惡慾,於很短的時間,已滿載於第四章中了!到了十代,人類便淪亡,到了不可見,不可聞,不可救之地位,而有洪水滅世之慘.誠如保羅所說,一點麵酵,能使全團都發起來.(五之九)這樣,吾們應當如何謹慎的嚴防罪惡,而除清舊酵呢.(哥前五之六至八)
\newline
\newline
說到情慾情形的黑暗,保羅對以弗所教會說,在你們中間,連題都不可,方合聖徒的體統.(五之三)耶和華因所多瑪俄摩拉罪惡之深重,升聞於天,所以要用天火以除滅之.(創十八之二十)前之洪水滅亡,亦即此故.然而,今日社會之人,竟將一切污穢之事,用畫片影了出來,用筆寫了出來,隨地隨時,當眾談論,不以為恥；這就證明人在黑暗中就不自覺己的黑暗!然而,吾們光明之子,屬於白晝之人,就當行事光明,結光明之果,即良善,公義,誠實等事,那暗昧無益的事,要與人同行,倒要責備行這事的人.(弗五之八至十一參看帖前五之五至十一)
\newline
\newline
(乙)情慾的類別
\newline
\newline
在此所言的情慾,雖有十五件之多,然而可分之為四類:
\newline
\newline
(子)性慾 ── 如淫,污穢,邪蕩,(十九節)男女兩犯者.
\newline
\newline
(丑)靈慾 ── 拜偶像!邪術,異端.(二十節)
\newline
\newline
(寅)社交之障礙
\newline
\newline
一,仇恨,二,爭競,三,忌恨,四,惱怒,五,結黨,六,紛爭,七,嫉妒.(二十至二十一)
\newline
\newline
(卯)肉體之嗜好 ──
\newline
\newline
一,醉酒,二1,荒宴.(二十一節)
\newline
\newline
以上所說的,并非包括世間一切情慾,不過列舉顯而易見的以言之；但此四類,則可概括一切了!
\newline
\newline
聖經既說情慾的事,是顯而易見的；(十九節)那就是說,凡是屬於情慾之行,無不自知的.既自知,即當如何呢?在二十四節說:凡屬基督耶穌的人,是已經把肉體連肉體的邪情私慾,同釘在十架上了.信徒們,我們是不是將情慾的事一一都已釘在十架上呢?如未,你肯不肯如此行呢?保羅怎樣警告迦人呢,他說:我從前告訴你們,現在又告訴你們,行這樣事的人,必不能承受上帝的國.(二十一節)你確信麼?夏絓亞當犯罪,以致累及後世萬代,豈不是在乎一字的疑惑麼?上帝說:你喫的日子必定死；但蛇(魔鬼)對他說,不一定死.(創二之十七又三之四)今日多少信徒,徘徊於聖潔的路上,躊躇不前,豈亦不是懷疑於此一必字麼?他以為上帝未必如此嚴重.然而,汝特因留戀一情慾而不捨,以致不能承受上帝的國,能不遺恨永遠麼?
\newline
\newline
(二)論聖靈
\newline
\newline
一,聖靈的性質
\newline
\newline
二,聖靈的果子
\newline
\newline
三,聖靈而生的人
\newline
\newline
聖靈的性質怎樣呢?在十七節說:情慾和聖靈相爭,聖靈和情慾相爭；這兩個是彼此相敵,使你們不能作所願意作的.然而,吾們順著聖靈,情慾便制服了.在此可知,聖靈之大力,是在信徒尊重他,順從他,方能顯出.所以在二十四節,其重要亦在一釘字,那是一步樂意遵行的工夫.這樣,信徒之不能去慾從靈,只好自己負責.今日教會中不少由墳墓中復活出來了的拉撒路；然而,還是纔著手腳,裹著布,臉上包著手巾,這是何等虧缺主耶穌的榮耀,埋沒聖靈的能力呢?
\newline
\newline
再看到聖靈的果子,就是顯出聖靈的美麗；與哥前十二章所說,聖靈的恩賜不同,恩賜是顯明聖靈的能力.所以在哥前十三之七說的明白,聖靈在各人身上,叫各人得益處.然而,在本章二十三節則說,這樣的事,沒有律法禁止；就是說生命勝了死亡,榮耀遮掩了黑暗.而且二十二節論到聖靈,說是「結果；」論情慾,則說「事或行」.一是由於天然的,造化生命的；一是由於屬世的；能力血氣的.靈乃屬天；慾乃屬世.樹木之結果,那是何等自然的事呢?因此,順著靈而行的人,也是如此；亦不需你的力量工作,比較行割禮律法,又是一個相對的證明.
\newline
\newline
論到聖靈的果子,惟言九件,件件都是令人可愛慕可親近的.(二十二,二十三節)猶如一串珍珠,最好最大的一顆在頭,就是仁愛；最尾結束的就是節制.那有何教訓呢?我以為這樣聖靈的果實,雖是白白的施賜於人；但吾人亦當慎重使用之,節制在尾,或是這樣的意思.你看,哥林多教會,得著方言的恩賜,不善使用,以致教會混亂起來；這是何等可怕的事呢!請詳讀哥前十四章末尾,就論到聖靈而生的人,便是本段的終結,亦是本章的總結.其有三事:
\newline
\newline
(子)不要貪圖虛名
\newline
\newline
(丑)不要彼此惹氣
\newline
\newline
(寅)不要互相嫉妒
\newline
\newline
這三事是屬靈的人所當提防的.靠靈行事的人,有時會求榮於人,去貪圖虛名,求人的稱讚；魔鬼也曾用過此法去害耶穌.他說,你若是上帝的兒子,可以從這裡跳下去；因為經上記著說,主要為你吩咐他的使者,保護你,他們要用手托著你,免得你的腳碰在石頭上.(路四之九至十一)這就是靈性的驕傲一種的試惑,屬靈的人,入了那(重)網羅,是極不容易脫離的!你看,有許多自命聖潔屬靈的人,總不見他有何善果與助人的事顯出；這是因為他一點屬靈的驕傲,已拒絕人於千里之外了!
\newline
\newline
第二段當提防的是什麼呢? ──惹氣,這就平等相處,會發生的危險；凡是屬靈的人,要更為謙卑.聖經說:基督不悅己；又說:若是能行,總要盡力與眾人和睦.(羅十五之一至三又十二之十八,十六)屬靈之人,豈當互相激怒麼?
\newline
\newline
第三所當提防的是:── 不要互相嫉忌,這是在下對上說的.有的靈歷自覺不及人之時,會自然生一種嫉忌之心,議論批評人,這不是平常人會犯的,尤其是稍具有靈眼的人,最易犯的.女先知米利暗,大祭司亞倫,都曾犯過了!(民十二章)吾們更當何等的小心呢.保羅說:各人看別人比自己強；(二之三)這真是屬靈的人所當具的態度.
\newline
\newline


\newpage

\section{第1屆~港九培靈會~趙柳塘~第七講}
\label{sec:1344}
\textbf{第七講}
\newline
\newline
連結: \href{https://www.hkbibleconference.org/session-message/view/1344}{\texttt{https://www.hkbibleconference.org/session-message/view/1344}}
\newline
\newline
\hyperref[sec:1343]{< < < PREV SERMON < < <}
~
\hyperref[sec:toc]{[返目錄]}
~
\hyperref[sec:1348]{> > > NEXT SERMON > > >}
\newline
\newline
得贖之人於各方面應盡之責任如何(第六章)
\newline
\newline
第六章是接著前章討論聖靈的題目,前章是說到聖靈的性質與功能；本章就說到一個屬靈的人,他於各方面有何應盡的責任.現將全章分析一下,便顯明他的責任如何.
\newline
\newline
(一)當挽回失足之信徒(一至五節)
\newline
\newline
(二)當捐助教會之需用(六至十節)
\newline
\newline
(三)當釘世界於十字架(十一至十六節)
\newline
\newline
(四)本書之終結(十七,十八節)
\newline
\newline
第一段向教會兄姊之責任如何(一至五節)
\newline
\newline
屬靈的人,其第一的責任,就是要扶持那些軟弱失足的兄姊們.所以,在(一節)說:兄弟們,若有人偶然被過犯所勝,你們屬靈的人,就當用溫柔的心,把他挽回過來.你看,屬靈的人,主并不是要他口誇一切所得之福,乃是要他幫助別人；今日有許多自以為是屬靈的,便遠遠的站在一邊,只是議論別人的長短,輕視他人,那可是合於聖經的教訓麼?主使他多接受聖靈,豈不是要用他幫助別人?昔日有一個血漏的婦人,她一接近了耶穌,她十二年來,人所不能醫治的病,便立刻好了.(可五之二十五至二十九節)又有一個患痲瘋的人,他祈求耶穌,耶穌不特不責備他,反為他醫愈了他永不能好的瘋疾.(太八之一至四)你看,這兩個人都是律法所定為不潔的人,不許親近人的,因為他們能使人成為不潔；但他們放膽的親近耶穌,耶穌的能力便使他們的不潔成為潔淨了.故凡一切屬靈的,亦當如是,要使軟弱不聖潔的人因之能受感而成聖潔；這才是屬靈人的好處!這才是教會需求!所以屬靈人并不是古玩,架上的玩品,乃是室中之用具；吾們可明白這個意思麼?而且挽回罪人的功夫是聖靈,要屬他的人第一當盡的責任,不過他要小心罷了,所以一節之下接著說,又當自己小心,恐怕也被引誘.
\newline
\newline
本段又說到(擔子)
\newline
\newline
一,是各人都有一個擔子要自己擔的(五節)
\newline
\newline
二,各人有時有重的擔子那是互相擔當的(二節)
\newline
\newline
屬靈的人,對於軟弱犯罪的人,不要責備太甚；因為他確是擔當不起了.屬靈的人有能力,就當為他盡力,代他擔當.然而代擔是什麼意思呢?就是為他祈禱,同他戰爭,助他脫離誘惑.猶如天使拉著羅得的手和他妻子的手,並他兩個女兒的手,把他們領出(將亡的所多馬)來安置在城外一般.(創十九之十六如此),便算完全了基督的律法.真是屬靈的人,并不像那偽師,勉強信徒受割禮,守月節,使人負他們和他們祖宗所不能負的軛.
\newline
\newline
(三,四節)所以保羅再向屬靈人警告一句說:人若無有,自己還以為有,就是自欺了!各人應當察驗自己的行為,我很願今日教會的信徒領袖,尤其是自命為屬靈的信徒領袖,常常自己察驗,能無以無為有的自欺否?我常見著聽著有與我有關係的人退步失敗了,我是何等自愧呢!保羅說:有誰軟弱我不軟弱呢?有誰跌倒我不焦急呢?(哥後一十之二十九)願主賜今日一切屬靈人有保羅如此的心志,教會便大奮興了.昔日大奮興家斐尼有一段話說:凡人得主愛復興在他心中之時,他如見人不愛他所切愛的上帝,其心必痛苦；是凡深信之人,彼必勸勉之,劬勞為之祈禱,更以其信德之手扶持之至於主前,望主復興於其心,亦望舉世之人皆歸信主.
\newline
\newline
至終他說,他所誇的就專在自己不在別人,那又是對偽師說的,因為偽師是專以別人為誇口的呀!(十三節)
\newline
\newline
第二,保羅便說到教會的供給上,他說在道理上受教的,當把一切需用的供給施教的人,(六節)這是屬靈人應盡的責任.中華教會已成立一百二十多年了,信徒們已有四五十萬,然而,有幾多自給的宣教士?自立的教會?即此一端,也就可知今日中華教會之靈程了!本來金錢算不得什麼,但保羅亦曾說過,那是試驗信徒們愛心的實在.(哥後八之七八節)金錢是人們的血汗,凡能愛主救贖的血,寶貴主僕們勞力的汗的信徒,然後能將他血汗的金錢貢獻出來.今日中國的教會不特靈糧的供給日夕仍多仰賴西教士,就是傳教士的身糧不也終年靠著西信士的供給麼.在此說是在道理上受教的,那就是一自主立的原理,無論何人,受了何人的教,得了他的益助,就當有當盡的責任.所以在哥前九章保羅已說的詳細,在那裡他舉出三個引證:
\newline
\newline
一,兵士當受糧餉
\newline
\newline
二,園夫當食果子
\newline
\newline
三,牧養牛羊的當喫牛羊的奶(七節)
\newline
\newline
總結 ── 傳福音的當靠著福音養生(十四節)
\newline
\newline
我願屬靈的領袖,在施靈訓之時,切勿忘了這一供給施教者重要的教訓；亦望屬靈的教友們,切勿忽略了供給施教者這一件本份,切勿學效哥林多教會,一切靈福無缺,惟於捐助一事則大不如人；這就是今日中華教會,同有的弊病!惟主靈奮興他的教會,則能同時一洗此恥.
\newline
\newline
吾們再要注意「一切」二字,此不僅指到金錢,亦指著一切日用之品物,無論食的,著的,凡自己所需的,亦當念及施教者之欠缺,憑主之愛,聖靈之感而供給之.你看外國教會之能遣宣教士於天下萬國,豈不是由於信徒之能獻己麼?由此吾亦切望中華信徒,能由貢獻田產起首,以實行此道.(在七節)保羅第二次說「不要自欺」,就是要屬靈的信徒,省察於此第二方面的責任曾盡了未.然而,在此他為何說道,「上帝是輕慢不得的呢,」兄姊們,因為捐助的事,不僅是向人與教會,實是向上帝；所以在馬拉基書論到猶太人不盡捐助之責,便說,是奪取上帝的物,因而通國的人都受了咒詛(三之八至十)唉呀!上帝是輕慢不得的.然而,今日中華信徒多多犯了此罪而不覺!例如不盡份捐助啦,認捐而不繳納啦,以偽幣獻主啦,種種傷心的事,說之不盡,言之可恥!這一切不特失福,自傷其德,而且使主之教會衰弱不振,人多不能得救,所以在(七節)說人種的是甚麼,收的也是甚麼.今如不求聖靈除去此罪惡,亦是今日奮興之一大阻力.(八節)又是論到吾們捐助的精神.在腓立比保羅有說:有的傳基督,是出於嫉妒紛爭,也有的是出於好意.(一之十五)捐助之法,亦分靈慾兩種,如亞拿尼拿夫婦之捐助,便是屬慾的,聖靈所以將他除滅了；(徒五之一至十二)因為他們是因人的稱讚而捐助.這樣,吾人該當如何的小心呢?所以耶穌有施濟右手勿叫左手知道之訓.(太六之一至四)亞拿尼亞夫婦,就是順著情慾而撒的,故所收的是敗壞.故此今日中有許多自立教會,仍是紛爭污穢,不是像老底嘉之似富實貧,則如撒狄之名生實死!(啟三之一至十三)靈命如何呢?救人的實際如何呢?順實著情慾撒種,不過敗壞而已!是故吾們為主獻己捐助,更當靠著聖靈去作,方能得著靈益實效.
\newline
\newline
本段的下一半便論到教會的慈善事業,可分為「向眾人」與「向信徒」兩方面.(十節)今日中華教會中一切慈善事業,多是外國信徒的靈果；本地果子,能有多少?這一事也是教會所當注意,而為主道之榮者.因為聖靈的果子一種是叫恩慈,耶穌說:你們的光也當這樣照在人前,叫他們看見你們的好行為,便將榮耀歸給你們在天上的父.(太五之十六)光照暗黑社會的,豈不是慈善的事業麼?向信徒一家的人更當這樣一語,就是信徒的相助；這相助,比較在社會服役,施行慈善,尤為不可少之事.「有了機會」四字,是說到吾們行善,當隨時預備,尋機會去行,不等待人來找汝,乃是要去尋人；不要學效那下耶利哥,既見道旁被盜所傷之人竟望望而去之之祭司與利未人,乃要學效那善撒馬利亞人；(路十之二十五至三十七)因為今日社會之痛苦已達極點,其所能解除之者,惟在於世人所憎惡的敵人,撒馬利亞人──耶穌之教會而已.(九節)乃是安慰堅定行善者的話,說我們行善不可喪志,若不灰心,到了時候,就要收成.你看,世界行惡的精神,總不會疲倦的；是以吾人行善豈可灰心喪志,而為世界之惡所勝.耶穌說:我父作事直到於今,我也作事.(約五之十七)蓋上帝耶和華,創造地極的主,並不疲乏,也不困倦.(賽四十之二十八)凡等候耶和華的,必從新得力,他們必如鷹展翅上騰；他們奔跑,卻不困倦,行走卻不疲乏.(三十一節)吾們行善之人,豈可靠己之力以行；惟當仰望等候主,必能不喪志灰心,到時就必收成.(詩一百二十六之五,六)說:流淚撒種的,必歡呼收割；那帶種流淚出去的,必要歡歡樂樂的帶禾捆回來.
\newline
\newline
第三段當釘世界於十字架(十二至十六節)
\newline
\newline
本書至終一段的討論與勸勉,仍是十字架；可云十架始,(一之一至三)十架生,(二之二十)十架終.但偽師們,因要避免十架的逼迫,欲求人的喜悅,便宣傳割禮,強人受割.(十二,十三節)十字架明明是一條得生得榮全智之路；但在猶太人為絆腳石,在外邦人為愚拙,(哥一之二十三)竟捨之弗由,都淪亡了!偽師畏其難,投世人心之所悅,竟偏離了,故入於迷途之中；惟能與主同釘同死之人,方能與保羅同說,我斷不以別的誇口,只誇我們主耶穌基督的十字架.因這十字架,就我而論,世界已經釘在十字架上；就世界而論,我已經釘在十字架上.(十四節)
\newline
\newline
再者,(十四節)的下半,保羅所誇十字架,并不僅向著教會,更是向著世界.保羅說因這世界已經釘在十字架上,那有兩方面的意思:
\newline
\newline
一,吾們 ── 因世界不愛耶穌 ── 釘世界於十字架
\newline
\newline
二,世界 ── 因吾們信愛耶穌 ── 釘吾們於十字架
\newline
\newline
看保羅這兩句的次序,是很有教訓的.信徒要先釘世界於十架,世界然後釘信徒於十字架；就是信徒當先與世界區別,然後世界方釋放你自由.吾們應當記得,摩西要帶以色列民出伊及,法老豈不是一再留著他麼?及至摩西向法老說了連一蹄也不留下斷然決然的話,(出十之二十六)然後再一次首生者死之的大災,以色列人方出離了伊及；經再一次紅海的追逐,伊及人方「連一個沒有剩下」(出十四之二十八)以色列人便真正的自由了.所以信徒要因切愛耶穌,真正完全的將世界釘了十架,十架方真捨棄你,釘你於十架；否則她仍然愛戀牽引你,一有機會,便完全再得回你作他的奴僕.吾們惟靠十字架,方能與世界隔絕；凡己心能得勝了世界的,方能救人脫離這世界,免與世界淪亡.能如此者,受割不受割,都無關緊要,因為他已成為新造之人了!(參看十五節)故凡能釘世界於十字架的,於偽師割禮的誘惑也都勝過了；信徒心靈常受罪惡攪擾的,都因他未曾向世界死罷!
\newline
\newline
(十六節)為以色列祝福在總論已說過了.
\newline
\newline
第四段本書之終結(十七,十八節)
\newline
\newline
你看保羅只用兩個意思來結束本書.
\newline
\newline
1.就本身說 ── 從今以後人都不要攪擾我,因為我身上帶著耶穌的印記.(十七節)
\newline
\newline
2.就信徒說 ── 弟兄們,願我主耶穌基督的恩常在你們心裡,阿們!(十八節)
\newline
\newline
第二義在第一講已說過了,就事無論如何,總要信徒求恩,今將第一義略論一下,便告完結.
\newline
\newline
印記,或作「司提基瑪」本來這個字不是好意思,在當日的社會,這印記是施在奴僕,或被擄者的身上,用紅鐵烙之；但今保羅說,我身上負著主耶穌之司基提瑪,不是屬人的奴僕的印記,或屬國家的兵丁的印記,或戰敗被擄者的印記；乃是為主耶穌基督奴僕的印記,屬天國的兵丁的印記,為主聖靈所敗終身降服於主耶穌權下之印記,即聖靈所印證者.那是何等榮幸可誇的事!而然,照世界而論,保羅不過為一奴僕,被擄者,世界的污穢,萬物中的渣滓而已；(哥前四之十三)但保羅只誇主耶穌基督的十字架,以耶穌之印記為誇,比較偽師以人手所受的割禮為誇,有大分別.聖經說:偽師們以辱為榮,心在世事.(三之十九)然而,主之忠僕,似乎一無所有,卻是樣樣都有；(哥後六之九,十)似辱實榮.所以保羅至終警告一句說:從今以後,人都不要攪擾我,那是他已戰勝一切誘惑,艱難的決辭!為什麼他拒絕這一切的事呢?就是因他身上已帶著耶穌的印記.
\newline
\newline
兄弟姊妹們!你未能勝過試誘,在你未曾與敵盡力戰爭以至流血的地步.然國家的榮,是國民的血換來的；人的榮,是汗換與腦汁換的；惟基督的榮,是向內勝利而得的.這榮也為那與主同釘同死而同生之人所共有.因此我們更當記著,保羅在此所說的印記,并非意想的,乃是事實的.我們試想,曾否為主多受如此之苦?(參看哥後十一之六,二十三至三十二)
\newline
\newline
各講較講稿略加詳明,因演講時限於時間,未盡欲言,今既許以文字發表,故敢多借篇幅以貢獻.
\newline
\newline
編者草草書成,不無錯誤,或言論過激,順先進諸公,幸以教正!
\newline
\newline


\newpage

\section{第1屆~港九培靈會~高樂弼~第四講}
\label{sec:1348}
\textbf{空瓶}
\newline
\newline
連結: \href{https://www.hkbibleconference.org/session-message/view/1348}{\texttt{https://www.hkbibleconference.org/session-message/view/1348}}
\newline
\newline
\hyperref[sec:1344]{< < < PREV SERMON < < <}
~
\hyperref[sec:toc]{[返目錄]}
~
\hyperref[sec:1349]{> > > NEXT SERMON > > >}
\newline
\newline
列王記下 4:1-4:7
\newline
\begin{longtable}{cl}
\hline
\hline
章節 & 經文 (和合本修訂版)\\
\hline
4:1 & \begin{tabularx}{0.7\textwidth}{X} 有個先知門徒的妻子哀求以利沙說：「你的僕人，我丈夫死了，他敬畏耶和華是你所知道的。現在有債主來，要帶走我的兩個孩子給他作奴隸。」 \end{tabularx} \\ \\ \relax
4:2 & \begin{tabularx}{0.7\textwidth}{X} 以利沙對她說：「我可以為你做甚麼呢？告訴我，你家裡有甚麼？」她說：「婢女家中除了一瓶油之外，甚麼也沒有。」 \end{tabularx} \\ \\ \relax
4:3 & \begin{tabularx}{0.7\textwidth}{X} 以利沙說：「你到外面去向所有的鄰舍借器皿，要空的器皿，不要少借。 \end{tabularx} \\ \\ \relax
4:4 & \begin{tabularx}{0.7\textwidth}{X} 然後你回家，關上門，你和你兒子在裡面把油倒在所有的器皿裡，倒滿了就放在一邊。」 \end{tabularx} \\ \\ \relax
4:5 & \begin{tabularx}{0.7\textwidth}{X} 於是婦人離開以利沙去了。她關上門，把自己和兒子關在家裡。他們把器皿拿給她，她就倒油。 \end{tabularx} \\ \\ \relax
4:6 & \begin{tabularx}{0.7\textwidth}{X} 器皿都滿了，她對兒子說：「再給我拿器皿來。」兒子對她說：「沒有器皿了。」油就止住了。 \end{tabularx} \\ \\ \relax
4:7 & \begin{tabularx}{0.7\textwidth}{X} 婦人去告訴神人，神人說：「你去賣了油還債，你和你兩個兒子可以靠著所剩的過活。」 \end{tabularx} \\ \\
[1ex]
\hline
\hline
\end{longtable}
今以以利沙為寡婦行奇事的故事來講,也許各位以為這是早年的故事,與今日的我們沒有什麼關係.然神歡喜把這事放在聖經中,是不是叫我們多知一件故事呢?我以為不然.我信聖經中所載的事都含有屬靈的要訓.茲將經文中那三件關於寡婦的事論列如下:
\newline
\newline
(一)極大的欠缺　　此婦是一寡婦,他的丈夫本是敬虔的人；今丈夫死了且遺下兩個兒子,所以他的責任就因而重了.然而不但重了且又欠債,欠了又不能償還；加以那債主又不是良善的人,他的困苦萬狀可由此而知了.有一天,債主對他說,你若不還債,我就要取你的兩個兒子作奴僕.此寡婦處在這樣的絕境下,誰也會向他生一同情的心.然而此寡婦的債,究有什麼靈訓呢?原來此債,就是預表人的靈債.所謂靈債,不是指罪而言,因人的罪,已由主耶穌代贖了,我們靠主的救贖,就已得了自由.不過我們的罪雖已由主代償清了,但我們還有靈債.這裡所說的靈債不是別的,就是信徒應行而不行的事.這不行的事,就是信徒的靈債；既有靈債,就證明對神的標準尚未達到,或問神向人所定的標準是怎樣的呢?然我們若要得到這個答案,就當到聖經裡找.主說:「你要盡心,盡性,盡意,盡力,愛主你的上帝,」這就是神向人所定的標準.今請我們細想,我們是否已經照此標準行事,又已否滿足此標準的要求.假若我們還未照標準行事,和滿足標準的要求,這還未就是我們的久缺,我們的靈債.主又說:「我賜給你們一條新命令,乃是叫你們彼此相愛.」想我們都是要守主要我們所守的；但我們能否已照所想做去?若是不能,這「不能」就是我們的欠缺,我們的靈債了.又:彼得前書一章十五節說:「那召你們的既是聖潔,你們在一切所行的事上也要聖潔,」這聖潔的標準,我們是否已經達到?又馬太五章四十八節:「你們要完全,像你們的天父完全一樣.」各位,你對這節聖經發生怎樣的意念呢?神是否真是要人像他一樣?我知神實在想人像他一樣的完全；因為神不說謊.難道獨此節經是神用來作裝飾麼?不!神既這樣的說,就必實在有這樣的意思.各位,我們若省察己心,就必像那寡婦一樣的為債掛心.保羅說:「我所願意的善,我反不作；我所不願意的惡,我到去作.」他這樣的事與心違,就覺得心裡有個律挾持著他,所以他就發生「我真是苦啊,誰能救我脫離這取死的身體呢?」他對神既有這樣的呼叫,神就開他的靈眼,叫他知有一法,可以使他一生得勝,這法就是,主的靈在他的靈裡,他就得著自由.今所說的寡婦,他不是異邦人,乃是屬主的人.他不推諉說:「我是貧寒的寡婦,雖欠人債,也是沒有法的；而且許多人也同我一樣的欠債.」她不說這樣推諉的話來安慰自己,也不以欠債當一件不必掛心的小事,卻以欠債為一件可掛心的大事.然她知欠債是苦事,也知向神求救,由此看來,此寡婦可作一般要解決靈債的人的模樣.今我們要為一屬靈的人,必先知己屬靈的欠缺,我們若在這些愛神愛人和潔己的事上失敗了,就不當以為這是不必掛心的小事,當學效此寡婦,切切的向神求助,而為一凡事得勝的人.
\newline
\newline
(二)神補足人欠缺的美法　　此寡婦一向神呼求救助,神就答應她；所以凡覺己有欠缺的人也當效他.主僕以利沙問寡婦說:「我可以為你作甚麼呢?」各位試想,神到底用什麼方法來補足此婦人的欠缺?神雖能用多法,惟在此神卻用一特別的法,這法就是從婦人所有的而濟助她.所以以利沙問她說:「你家裡有甚麼?」她回答說:「婢女家中,除了一瓶油之外,沒有甚麼.」今我們從她的答話,就知她不注重那些油,也許她以為這微小,不用言了,誰知她不注重的很小油,神就用那些小的油,增多起來,使她足以還債而有餘.
\newline
\newline
然而這裡所說的「油」,有什麼靈訓呢?此靈訓對我們又有什麼關係呢?查在聖經所說的油,是預表聖靈,神在舊約時代,用油膏祭司,與主,以表受膏聖靈.今此婦得油,可比我們得勝靈而償清所欠的靈債.我們知此婦原已有油比如信徒有聖靈在心.聖經說:「人若沒有靈,那就不屬耶和華.」今我甚怕各位有聖靈而不知聖靈的價值,反以為不甚重要,可有可無.或有人只知聖靈的名,而不知聖靈能助人一生行得勝的路.感謝主,他有一寶貴的應許賜給我們,這應許可在聖經裡見得的.在馬太七章十一節說:「你們雖然不好,尚且知道拿好東西給兒女,何況你們在天上的父,豈不更把好東西給求他的人麼.」又在路加十一章十三節說:「你們雖然不好,尚知道拿好東西給兒女,何況天父,豈不更將聖靈給求他的人麼.」這兩節聖經,同一的意思.馬太所說的好的,就是路加所說的聖靈.神既然賜聖靈給人,人若得著,就得著一切的好處.今姑且說一比方,以顯明此意.比方今有一貧人到這裡求助,我就有二個方法可以助他.第一法是和他一同出外購備一切所需的東西；第二法是給錢與他,叫他自行購備.由此看來,他有錢,就能有一切,所以錢就是他所需的好物.神今已賜給聖靈與我們,使我們得補足一切的欠缺；且使我們獲得一切屬靈的好處.屬靈的好處是什麼呢?
\newline
\newline
(1)是愛.
\newline
\newline
(2)是聖潔.
\newline
\newline
(3)是知聖經的奧義.
\newline
\newline
(4)是有能力的禱告.
\newline
\newline
(5)是有能力的講道.
\newline
\newline
(6)是喜樂.
\newline
\newline
(三)接受的法.
\newline
\newline
(1)備空的瓶.
\newline
\newline
(2)遵命倒油.
\newline
\newline


\newpage

\section{第1屆~港九培靈會~高樂弼~第五講}
\label{sec:1349}
\textbf{迦南婦}
\newline
\newline
連結: \href{https://www.hkbibleconference.org/session-message/view/1349}{\texttt{https://www.hkbibleconference.org/session-message/view/1349}}
\newline
\newline
\hyperref[sec:1348]{< < < PREV SERMON < < <}
~
\hyperref[sec:toc]{[返目錄]}
~
\hyperref[sec:1350]{> > > NEXT SERMON > > >}
\newline
\newline
馬太福音 15:21-15:28
\newline
\begin{longtable}{cl}
\hline
\hline
章節 & 經文 (和合本修訂版)\\
\hline
15:21 & \begin{tabularx}{0.7\textwidth}{X} 耶穌離開那裡，退到推羅、西頓境內。 \end{tabularx} \\ \\ \relax
15:22 & \begin{tabularx}{0.7\textwidth}{X} 有一個迦南婦人從那地方出來，喊著說：「主啊，大衛之子，可憐我！我女兒被鬼纏得很苦。」 \end{tabularx} \\ \\ \relax
15:23 & \begin{tabularx}{0.7\textwidth}{X} 耶穌卻一言不答。門徒進前來，求他說：「這婦人在我們後頭喊叫，請打發她走吧。」 \end{tabularx} \\ \\ \relax
15:24 & \begin{tabularx}{0.7\textwidth}{X} 耶穌回答：「我奉差遣只到以色列家迷失的羊那裡去。」 \end{tabularx} \\ \\ \relax
15:25 & \begin{tabularx}{0.7\textwidth}{X} 那婦人來拜他，說：「主啊，幫幫我！」 \end{tabularx} \\ \\ \relax
15:26 & \begin{tabularx}{0.7\textwidth}{X} 他回答：「拿孩子的餅丟給小狗吃是不妥的。」 \end{tabularx} \\ \\ \relax
15:27 & \begin{tabularx}{0.7\textwidth}{X} 婦人說：「主啊，不錯，可是小狗也吃牠主人桌上掉下來的碎屑。」 \end{tabularx} \\ \\ \relax
15:28 & \begin{tabularx}{0.7\textwidth}{X} 於是耶穌回答她說：「婦人，你的信心很大！照你所要的成全你吧。」從那時起，她的女兒就好了。 \end{tabularx} \\ \\
[1ex]
\hline
\hline
\end{longtable}
今早所講的是一舊約的婦人,現在所講的卻是一新約的婦人.在馬太十五章廿一至廿八節,和馬可七章廿四至卅節,所載的雖同說一事,但其中有點分同別；望各位小心來聽,就明白那婦人一切的事了.
\newline
\newline
記載那婦人的事是聖經,這聖經是一本奇妙的書.無論怎麼的書讀了一次,就叫人不願再讀了,或再讀一次,就把書的內容已盡知無遺了；惟聖經就不同了,每次有每次的新教訓；從這裡就可證明聖經不是由人作成,乃是由神靈所感而寫成的.今講這迦南婦人的故事,這故事為我們熟知的.在字面就論,人人也知道了；惟字裡多未顯靈隱藏的教訓,就不為人所知.今想講在這故事所包含的幾個靈訓.我們當先知耶穌動作,然後再研究此婦人,看有什麼教訓.我們若研究耶穌,就得有二個要訓.
\newline
\newline
(一)耶穌的感動力.
\newline
\newline
我信在座中有的是領袖,但你們要想你似主還是似法利賽人；有沒有人追我們好像追主；我們是否像主忙個不了,也許我們的時間除主日聚集外,統是空閒；也許我們用好些計劃叫人要親近我們也不能親近.這樣的疏遠人,是好事,還是不好的事?你是主的工人,若常把時間來為自己,你就是不忠實的工人了!
\newline
\newline
在美國有許多富翁,聽說日日有多人向他們乞助；但我在美國時,就沒有人乞助於我.為什麼呢?這因為人知我沒有錢,故他們不來；惟是有錢的,那就常常有人來了.蜜蜂知那種花有可釀蜜的材料,就親近那種花.我們也當想,有沒有人想親近我.主名如香膏,其實人人都有一種氣味.不是香的,就是臭的.有人有罪的氣味,這味神不願聞；有人有情慾的氣味,有人有利己的氣味.我知有多人能為耶穌作工,他也熱心為主作工,但他所作的,就是為己；你只見他卻不見耶穌.願各位沒有這利己的氣味,而有耶穌的香味.保羅在哥林多後書二章二十四節說:「隨處有基督之香.」你們都知彼得約翰被提到公會裡去,眾人都看出他們原是沒有學問的小民；他們雖沒有學問,但他們有一種香氣,能夠使人知他們是耶穌的門徒.查他們能使人知他們是耶穌的門徒,乃因他們多與主同在.同樣,我們多讀經,多禱告,並同他們一樣的有耶穌的香味.主留存我們在這世界,就想我們為他作證,叫人從我們而認識主.不知人看見了我們,信主還是離棄主?我們知我們的行為能影響別人.
\newline
\newline
(二)耶穌的愛心.
\newline
\newline
我們真是欠缺愛人的心!無論你有何等的才幹,若缺少了愛,就不能救人的靈魂.在約翰第十章說有三種牧羊人,一是賊；一是工人；一是牧人.工人雖不是賊,但他牧羊不是為愛,乃是為錢,故狼來了就捨羊而逃；因他是工人,故他愛己過於愛羊.惟善牧者若遇了狼,就為羊捨命,因他是善牧,所以他愛羊過於愛己；羊聽聞他聲就跟從他.請你想,我們是工人,還是善牧?我們知主不即應婦人的求,不是為不愛她；是要為她做成大事,引她進入信德的更深處.我們從上看來,就知主有何等的感力,何等的愛了!
\newline
\newline
今再進而研究此迦南婦人,看有什麼教訓.我們當知,主此次為最後到推羅西頓境的一次；此地雖多人,惟只一人得主的恩典；她尋得主,便立即求主,她就成為尋主的人的模範.茲研究她求得所欲的所以然.
\newline
\newline
(一)因她有切心
\newline
\newline
此婦人知他女兒所欠缺的是什麼,故她不到耶穌面前說些無關痛癢的話,乃懇切的說,「主啊,大衛的子孫,可憐我,我的女兒被鬼附著甚苦.」今人傳道少說地獄的痛苦；願神感動我們,確知不信耶穌的定必落地獄.望神叫我們有如此婦的切心.在夜半求餅的教訓,就是恆心的就得著.所以我若不恆心切求,神就不應允所求；此婦一生只得此一次機會救他的女兒；故她一見主便立刻切切求主,且至求得.令我們可以日日到主前切心為信徒祈禱,但惜我們不肯!
\newline
\newline
(二)因她有自卑的心.
\newline
\newline
(三)因她有篤信的心.
\newline
\newline
今研究此婦人的故事,就知主怎樣的待他；且又知他怎樣勝過艱難.
\newline
\newline
(一)不睬他.
\newline
\newline
(二)叫他灰心.
\newline
\newline
(三)羞辱她.
\newline
\newline
從上看來,就可見此婦人的信心是非常之大；故耶穌見他這樣而非常的歡喜,故說；「婦人,你的信心是大的；照你所要的給你成全了罷.」各位,耶穌今日也如昔日一樣,故無論何人,我們若有此婦的切心,自卑心,篤信心來求主,主亦必成就我們的所求.
\newline
\newline
在座中若有還未得救的人,你也當自卑到神前,神必救你；望神激動已信主的人,而為那些未得救的人掛心,使他們終底得救.望主祝福作主工的工人,真能犧牲一切.我願常為中國祈禱,使中國教會復興.在此數年間,魔鬼想盡力撲滅主的教會,難道我們讓他得勝麼?然我們要感謝主,因神常用魔鬼的詭計,而成就他的意旨；我們中國的信徒若切求主,則主必利用魔鬼的詭計,使主道在中國流行,而拯救多人.不過我們要知,神在未拯救人之前,須先復興他的兒女.
\newline
\newline


\newpage

\section{第1屆~港九培靈會~高樂弼~第六講}
\label{sec:1350}
\textbf{禱告與讀經}
\newline
\newline
連結: \href{https://www.hkbibleconference.org/session-message/view/1350}{\texttt{https://www.hkbibleconference.org/session-message/view/1350}}
\newline
\newline
\hyperref[sec:1349]{< < < PREV SERMON < < <}
~
\hyperref[sec:toc]{[返目錄]}
~
\hyperref[sec:1320]{> > > NEXT SERMON > > >}
\newline
\newline
使徒行傳 20:17-20:35
\newline
\begin{longtable}{cl}
\hline
\hline
章節 & 經文 (和合本修訂版)\\
\hline
20:17 & \begin{tabularx}{0.7\textwidth}{X} 保羅從米利都打發人往以弗所去，請教會的長老來。 \end{tabularx} \\ \\ \relax
20:18 & \begin{tabularx}{0.7\textwidth}{X} 他們來了，保羅對他們說：「你們自己知道，自從我到亞細亞的第一天，我怎樣跟你們相處， \end{tabularx} \\ \\ \relax
20:19 & \begin{tabularx}{0.7\textwidth}{X} 怎樣凡事謙卑，以眼淚服侍主，又因猶太人的謀害經歷試煉。 \end{tabularx} \\ \\ \relax
20:20 & \begin{tabularx}{0.7\textwidth}{X} 你們也知道，凡對你們有益的，我沒有一樣隱瞞不說的，或在公眾面前，或在每一個人的家裡，我都教導你們， \end{tabularx} \\ \\ \relax
20:21 & \begin{tabularx}{0.7\textwidth}{X} 不論猶太人和希臘人，我都已證明他們當在神面前悔改，信靠我們的主耶穌。 \end{tabularx} \\ \\ \relax
20:22 & \begin{tabularx}{0.7\textwidth}{X} 現在我被聖靈催迫要往耶路撒冷去，雖然不知道在那裡會遭遇甚麼事， \end{tabularx} \\ \\ \relax
20:23 & \begin{tabularx}{0.7\textwidth}{X} 但知道聖靈在各城裡向我指證，說有捆鎖與患難等著我。 \end{tabularx} \\ \\ \relax
20:24 & \begin{tabularx}{0.7\textwidth}{X} 我卻不以性命為念，只要走完我的路程，完成我從主耶穌所領受的職分，為神恩典的福音作見證。 \end{tabularx} \\ \\ \relax
20:25 & \begin{tabularx}{0.7\textwidth}{X} 「我素常在你們中間到處傳講神的國；現在我知道，你們眾人以後不會再見到我的面了。 \end{tabularx} \\ \\ \relax
20:26 & \begin{tabularx}{0.7\textwidth}{X} 所以我今日向你們作證，你們中間無論何人死亡，罪不在我。 \end{tabularx} \\ \\ \relax
20:27 & \begin{tabularx}{0.7\textwidth}{X} 因為神一切的旨意，我並沒有退縮不傳給你們的。 \end{tabularx} \\ \\ \relax
20:28 & \begin{tabularx}{0.7\textwidth}{X} 聖靈立你們作全群的監督，你們就當為自己謹慎，也為全群謹慎，牧養神的教會，就是他用自己血所買來的。 \end{tabularx} \\ \\ \relax
20:29 & \begin{tabularx}{0.7\textwidth}{X} 我知道，在我離開以後必有兇暴的豺狼進入你們中間，不顧惜羊群。 \end{tabularx} \\ \\ \relax
20:30 & \begin{tabularx}{0.7\textwidth}{X} 就是你們中間也必有人起來，說悖謬的話，要引誘門徒跟從他們。 \end{tabularx} \\ \\ \relax
20:31 & \begin{tabularx}{0.7\textwidth}{X} 所以你們要警醒，記念我三年之久，晝夜不斷地流淚勸戒你們各人。 \end{tabularx} \\ \\ \relax
20:32 & \begin{tabularx}{0.7\textwidth}{X} 現在我把你們交託給神和他恩惠的道；這道能建立你們，使你們和一切成聖的人同得基業。 \end{tabularx} \\ \\ \relax
20:33 & \begin{tabularx}{0.7\textwidth}{X} 我未曾貪圖一個人的金、銀或衣服。 \end{tabularx} \\ \\ \relax
20:34 & \begin{tabularx}{0.7\textwidth}{X} 你們自己知道，我靠兩隻手工作來供給我和同工的需用。 \end{tabularx} \\ \\ \relax
20:35 & \begin{tabularx}{0.7\textwidth}{X} 我凡事給你們作榜樣，叫你們知道應當這樣勞苦，扶助軟弱的人，又當記念主耶穌的話，說：『施比受更為有福。』」 \end{tabularx} \\ \\
[1ex]
\hline
\hline
\end{longtable}
當我在禱告間,得勝靈感我,要我講使徒行傳廿章卅二節這節經.昨讀的經文,就是保羅對當時教會的臨別贈言.這贈言的大意,首說己前所做的事工,後說及將來教會要遇者的事.他知將來必有兇暴的豺狠來殘害信徒；故特預告將來不但在外有逼害,在內還有像豺狼一般的偽師,說悖謬的話來引誘人.因此,他就甚願信徒堅立主道上而不背教,在臨別時就以至要的話儆醒信徒.他在卅二節說及二事,一是將他們交託上帝,一是將他們交託恩惠的道.我以為在這樣的情勢下,無別的事較這事為更要了.
\newline
\newline
(一)禱告.
\newline
\newline
各位,我們有沒有利用機會,成就神的大事?我們既知禱告就是我們的福份,為什麼不快快的去做?雅各書說:「人若知道行善,卻不去行,這就是他的罪了.」所以撒母耳對以色列人說:「至於我,斷不停止為你們禱告,以致得罪耶和華.」我在這數月間,得與內地會全球的總幹事,作多次很長的禱告；我見他禱告的生活及態度的誠懇,我的心就受了許大的感動；他有重大的責任在肩上,雖每日有許多工作,但他每日都找數時來禱告,他好像大祭司站在神前,為眾百姓禱告；他每日為數百人按名禱告,為各地教會及各傳道事業禱告；我信他最大的工,不在管理內地會全世界的工作,而在他的禱告.今我問各位每日有多少時候禱告?也許你們回答我說,沒有時候；但我知你們每日總有些時候可用來禱告.甚望各位,效法保羅的樣式而多多為他人禱告.在此你們要知,魔鬼必用千方百法來阻人的禱告；惟你們當效使徒的人,專心禱告,事事靠主,那就必定勝過魔鬼.
\newline
\newline
(二)讀經
\newline
\newline
前數月我留在高麗,想各位早知在高麗的大奮興會中,有多人得救.在一千九百零七年時,有一大奮興會,初在一有平陽的城市起首.我曾到過那地,會見一位祁牧師,此牧師就在此會講道；當講道時,聖靈就降臨下來,會眾就大受感動.祁牧師是一位瞎眼的人,當他目未瞎前,已讀經廿二遍；所以聖經很熟,甚至成本也能念出.當日本虐待高麗人及逼迫教會時,祁牧師就被押入獄裡；他雖在獄裡但常思念聖經而多得安慰.他對我說:「神許我進入獄裡,叫我多想神的言語,我今覺越思想而越覺有力.」今向各位,特講神怎樣的祝福他的事.在一千九百零七年時,獨他在那地為牧師；到今那地已多開十七間禮拜堂,教友成為牧師的已有四十五人,成為會吏的六十五人.他是高麗牧師中一位熱心的人.高麗的領袖因喜愛禱告與讀經,連教友也被感染而同樣的喜愛禱告與讀經；所以高麗教會有這麼好的成績.有時開研經大會,各鄉間的人也來；那些貧窮的人,就在平時積蓄些錢作赴會的費用.當一天早上,我們到平陽時,他就問我,願不願見一最早的早禱會,我們去到時,就見有三百多人跪著禱告,他們覺得高麗當再有一次更大的復興；他們既有這樣的誠懇,教會就大得神的祝福.
\newline
\newline
在詩篇一篇二節有說,愛神的人,就必喜愛神的言語.又在詩篇一百一十九篇那篇中,有一百六十節,其中有廿節,說及聖經怎樣的可愛.望各位都像大衛一樣的愛聖經.我們想做強壯的信徒,就只有讀經.聖經教訓初信的人,當渴慕靈奶.大衛以聖經為足前的燈.今我們人人都需要聖靈的引導,使己不致蕩入迷途；但要得勝靈的引導,就必須讀經.這因神的言語有監察人心,使人辨別靈慾的能力；所以希伯來書四章十二節說；「上帝的道是活潑的,是有功效的,比一切兩刃的劍更快,甚至魂與靈,骨節與骨髓,都能刺入剖開,連心中的思念和主意,都能辨明.」又神的言語,有保守人免入迷途的力量.你看,主在曠野受魔鬼試探時,他雖是神的獨生子,但他還要用神的言語來破魔鬼的詭計；他也要心藏神的言語,好使試探來時,而得以不敗.今我們當效他,保己不致一遇試探,便跌倒了.各位要知,魔鬼既敢試探主,也必會來試探你我.有時他會搖動我們的信,叫我們懷疑自己不得救；但你無別的話可答魔鬼,惟有聖經能夠回答他.當魔鬼來誘惑你,你若是讀經的人,就知所分別,而認識他的一切詭計.大衛說:「主的言語藏在我心,就不至得罪主.」我們都知,聖經為聖靈的劍；所以我們講道時,當用聖經,惟聖經才能使人悔改.保羅囑咐提摩太說:「務要傳道,無論得時不得時.」我知今日有些牧師,不傳道而傳科學.今有些人譏誚我們,說我們傳太舊的道.聞王載先生在港時,也受人這樣的譏誚.各位:我們不要怕人的譏誚,只要怕真理；所以我們要堂堂正正,忠實地傳聖經.我們能夠這樣,神就定必祝福他的道,與及我們.我們不必用言語為聖經辯護,只當對人說；聖經就自能為己辯護.我知中國許多信徒願意中國得大的復興,但要多人得復興,就必須先使一己得復興.神說:「我使水充滿那乾涸之地.」他怎樣的充呢?他就是先叫我們渴,然後才以生命的水充滿我們；我們得充滿了,才能充滿他人.主說:「口渴的人也當來；願意的都可以白白的取生命的喝水」.
\newline
\newline
各位!我們若要得到復興,惟有禱告與讀經.
\newline
\newline


\newpage

\section{第1屆~港九培靈會~黃原素~開會講道}
\label{sec:1320}
\textbf{開會講道}
\newline
\newline
連結: \href{https://www.hkbibleconference.org/session-message/view/1320}{\texttt{https://www.hkbibleconference.org/session-message/view/1320}}
\newline
\newline
\hyperref[sec:1350]{< < < PREV SERMON < < <}
~
\hyperref[sec:toc]{[返目錄]}
~
\hyperref[sec:1334]{> > > NEXT SERMON > > >}
\newline
\newline
路加福音 9:27-9:36
\newline
\begin{longtable}{cl}
\hline
\hline
章節 & 經文 (和合本修訂版)\\
\hline
9:27 & \begin{tabularx}{0.7\textwidth}{X} 我實在告訴你們，站在這裡的，有人在沒經歷死亡以前，必定看見神的國。」 \end{tabularx} \\ \\ \relax
9:28 & \begin{tabularx}{0.7\textwidth}{X} 說了這些話以後約有八天，耶穌帶著彼得、約翰、雅各上山去禱告。 \end{tabularx} \\ \\ \relax
9:29 & \begin{tabularx}{0.7\textwidth}{X} 正禱告的時候，他的面貌改變了，衣服潔白放光。 \end{tabularx} \\ \\ \relax
9:30 & \begin{tabularx}{0.7\textwidth}{X} 忽然有摩西和以利亞兩個人同耶穌說話； \end{tabularx} \\ \\ \relax
9:31 & \begin{tabularx}{0.7\textwidth}{X} 他們在榮光裡顯現，談論耶穌去世的事，就是他在耶路撒冷將要完成的事。 \end{tabularx} \\ \\ \relax
9:32 & \begin{tabularx}{0.7\textwidth}{X} 彼得和他的同伴都打盹，但一清醒，就看見耶穌的榮光和與他一起站著的那兩個人。 \end{tabularx} \\ \\ \relax
9:33 & \begin{tabularx}{0.7\textwidth}{X} 二人正要和耶穌分離的時候，彼得對耶穌說：「老師，我們在這裡真好！我們來搭三座棚，一座為你，一座為摩西，一座為以利亞。」他卻不知道自己在說些甚麼。 \end{tabularx} \\ \\ \relax
9:34 & \begin{tabularx}{0.7\textwidth}{X} 說這些話的時候，有一朵雲彩來遮蓋他們；他們一進入雲彩就很懼怕。 \end{tabularx} \\ \\ \relax
9:35 & \begin{tabularx}{0.7\textwidth}{X} 有聲音從雲彩裡出來，說：「這是我的兒子，我所揀選的。你們要聽從他！」 \end{tabularx} \\ \\ \relax
9:36 & \begin{tabularx}{0.7\textwidth}{X} 聲音停止後，只見耶穌獨自一人。當那些日子，門徒保持沉默，不把所看見的事告訴任何人。 \end{tabularx} \\ \\
[1ex]
\hline
\hline
\end{longtable}
昨年散會日最後的講題是甚麼呢?諒各位尚能記憶,所講的,就是「戰」的問題.大意說,各位散會後,各走各路,各上各的戰場去,為主作戰,若不幸有失敗的,就第二屆之會集,少了些人了；現在感謝主!第二會已到了,見各位踴躍回來,可知各位在這一年中,打了美好的仗,為主作了多多的善工,現在喜樂回來赴會,我們應該同心讚頌主!
\newline
\newline
本屆所希望之點,在這經文中,可以概見,容我漸漸往下說來.我們知到耶穌和三個們徒,在山上的工作,實在是一件特別奮興的工作；門徒在這事以先,受了耶穌很多嚴厲剌激的警告,人情上說,不免會心意灰冷,耶穌明白這事,所以特地用素未彰顯的榮耀,給他們看看,以慰他們的灰心.耶穌在這山上的奮興中,他的目的是甚麼呢?不外兩個字的工作,就是見「和聞」的工作是了.這裡的見聞,不是他們素來所曾有的,是他們一生所未經過的,他們所見,是見主的榮；所聞,也是聞主的聲.見人,聞人,不會大有益,且或有大害,見主,聞主,是大大有益的.我們在這裡聚會十天,若是全然見人的作為,聞人的言論,於我們究有甚麼益處?所以我們十分希望在這十天的聚會,所見所聞的,都是全然由主而來的.各位若確實在心靈中見過了主,聞過了主的聲音,那就雖然閉會散去,但這樣的靈益,都會永遠隨著我們的了.他們三個主徒看見了甚麼呢?恐怕就是:
\newline
\newline
(一)見主的國.耶穌說:「站在這裡的,有人沒嘗死以前,必看見上帝的國.」(二十七節)耶穌說此,就是指他在山上變化的事,他們有摩西和以利亞同在,更加證明這是天國的情形,因這兩位真是天國的人.
\newline
\newline
教會,本是天國的代表,但可惜有名無實,我們見天國,或入天國,不必等到死了以後才會見會入,在今世的生活中,也應當見了天國,入了天國.比方說永生,我們不是等到死後才得永生,今日信主,即已入了永生.(約五章二十四節)各位,教會今日在實在可以代表天國嗎?世人見教會,就會感動他們信有天國嗎?天國的榮耀,平安,喜樂,權能,在教會可以顯出來嗎?我以為未必!所以我們切切的盼望我們在這十天的聚會裡,會將天國的事顯現出來,給我們見,給我們聞.我說過了,天國若確然實現,應有大榮耀,大平安,大喜樂,大權能.我們最怕在會集中,人被高舉,主被壓低.今日的教會多半是這樣的情形,注重人手問題,追求人的方法,不肯受靈的指導,所以主斷難得高舉；主既不能在教會中尊大,那榮耀,平安,喜樂,權能,也就必不能實現.世界在靈界中,只有兩個國度,就是上帝國和撒旦國.上帝國的尊榮,不能多大彰顯；而撒旦國的惡勢,反為積極擴張,我們試閉目一思,社會中那裡不有撒旦的作為呢?那裡不有他的旗幟呢?啊!讓我們切切祈禱,叫主的國快快顯現出來；
\newline
\newline
(二)見主的榮.耶穌是榮耀的,但門徒總沒見過,他們此次一見,以後信從耶穌的心,必大大加增.各位,「見主榮」這一句話,不可輕看的,我們常說為主作證,若我們沒有見過主,怎可以作證?這個證字意思就是,對於這事,我已見,已知,已嘗過,已明白曉得,現在我要為這作實證.所以報紙的戰事紀載,人常看為虛言,若有人由戰爭陣上回來報告,他的言論是有價值的.(彼後一章十六至十八節)我們傳道,不是講學,心靈也沒有經驗,斷不能作有能力的見證.各位啊!我們若要作個好證人,要自己先見主.還有一件應當注意的,就是必要先見主榮,然後可見己辱,以賽亞,不是一個大好的先知麼?因何會有「禍哉!我滅亡了,因為我是嘴唇不潔之人」的話呢?就是因他「眼見大君王萬軍之耶和華」呀.(六章一至五節)約伯未得主向他顯現以前,他自覺很公義,所以在他二十九至三十一這三章書中,有一百四十六個我字,就是說我如何公義,我如何慈愛,我如何得主的祝福,我地位如何的高上,總之,我......我......我......我......我......,這一個「我」字,就是約伯得罪主的根源.後來主向他顯現說話,他就向主說:「我從前風聞有你,現在親眼看見你,因此我厭惡自己,在塵土和爐灰中懊悔.」(四十二章五六節)美國宣信先生說,英文字母之「i」字,(我也)比較別的字母格外出色,所以每版文字中,最易奪人目的,就是各語中的「i」字.我們的工作也是如此,工作中充滿了「我」,這個「我」,就是聖靈的大仇敵.各位啊!你若見了主的榮,你就不會有這麼多的「我」了.所以盼望各位,人人都有這樣「碎我」或非我惟主的經驗,更望在這十天聚會中得這經驗.
\newline
\newline
現在讓我們說到「聞」的教訓,此次的聞,是非常的聞,是他們一生未曾聞過的聞,因為在雲裡說出來的,就是上帝親自發聲說的.請我們注意「雲」字!聖經上說:「有一朵雲彩來遮蓋他們,他們進入雲彩裡.」他們這樣的境地,真是好過不了,因為這是天的情形,不是地的景況.雲表明聖靈,我們切願此次聚會,都要像他們進入靈裡面,就是我們要被聖靈覆蓋,在聖靈中作事.各位要預備耳朵,聽由聖靈而來的聲音.我們心中常有兩種聲音,一由聖靈來的,一由仇敵來的；我們要辨別,我們要認識,若是聖靈的話,我們要聽,我們要側耳聽；我們是耶穌的羊,必會聽耶穌的聲音.比方無線電的聲碼,若不是電報司機人,就不會明白裡面的意思；但他們極會,聽得很清楚.上帝屢屢藉著聖靈向我們打無線電報,我們要學聽他的消息.各位,我們一生作信徒,若不會聽主的消息,真是不幸之極了.雲裡的聲音說甚麼呢?是說:「這是我的愛子,你們要聽他!」為甚麼在此千載一時的機會,上帝不向人多說幾句呢?只說「這是我的愛子,你們要聽他」這兩句普通的話呢?夠了,他們若實行聽從這兩句話,就夠了!頭一句,是為耶穌作證；後一句,是勸人們絕對的順從耶穌.我們若能絕對順服耶穌,豈不是夠了麼?盼望我們在這十天的聚會,人人失了自己,去完全順服耶穌,讓耶穌為我們會裡的主席,聖靈為我們會裡的總指揮.我們又想到他們的光景,真是令人可慕的.彼得說:「夫子,我們在這裡真好,」他有不願下山的意思.是的,確實好,但時候未到,山下有很大的工作等候他們去做.我們在這培靈會中,倘若真有天國的光景,也不是我們當常久居留的.我們得了主祝福之後,要回去為主作工.還有一事要注意的,他們在山上的光景,比較下山的光景,各位試一比較,真有天淵之別.為甚麼呢?這因有耶穌同在和無耶穌同在的分別.有耶穌同在,監獄就是天堂；沒有耶穌同在,王宮也是地獄.各位你們一定要散開,回去你的工場,不要畏懼你的環境惡劣,只怕沒有耶穌和你同在.最後還有一句話,請各位不要忘記本年第二會訓說:「主啊!求你在這些年間,復興你的工作；在這些年間,顯明出來.(哈三章一節)感謝主,中國各地教會已略有復興的火燒起,願我們為此祈求,使主的工作,真正在這數年間復興起來,阿們!
\newline
\newline


\newpage

\section{第1屆~港九培靈會~黃原素~第一講　五祭之概論}
\label{sec:1334}
\textbf{第一講　五祭之概論}
\newline
\newline
連結: \href{https://www.hkbibleconference.org/session-message/view/1334}{\texttt{https://www.hkbibleconference.org/session-message/view/1334}}
\newline
\newline
\hyperref[sec:1320]{< < < PREV SERMON < < <}
~
\hyperref[sec:toc]{[返目錄]}
~
\hyperref[sec:1321]{> > > NEXT SERMON > > >}
\newline
\newline
利未記 1:1-7:10
\newline
\begin{longtable}{cl}
\hline
\hline
章節 & 經文 (和合本修訂版)\\
\hline
1:1 & \begin{tabularx}{0.7\textwidth}{X} 耶和華從會幕中呼叫摩西，吩咐他說： \end{tabularx} \\ \\ \relax
1:2 & \begin{tabularx}{0.7\textwidth}{X} 「你要吩咐以色列人，對他們說：你們中間若有人要獻供物給耶和華，可以從牛群羊群中獻牲畜為供物。 \end{tabularx} \\ \\ \relax
1:3 & \begin{tabularx}{0.7\textwidth}{X} 「他的供物若以牛為燔祭，要獻一頭沒有殘疾的公牛，獻在會幕的門口，他就可以在耶和華面前蒙悅納。 \end{tabularx} \\ \\ \relax
1:4 & \begin{tabularx}{0.7\textwidth}{X} 他要按手在燔祭牲的頭上，為自己贖罪，就蒙悅納。 \end{tabularx} \\ \\ \relax
1:5 & \begin{tabularx}{0.7\textwidth}{X} 他要在耶和華面前宰公牛犢；亞倫子孫作祭司的要獻上血，把血灑在會幕門口壇的周圍。 \end{tabularx} \\ \\ \relax
1:6 & \begin{tabularx}{0.7\textwidth}{X} 他要剝去燔祭牲的皮，把燔祭牲切成塊。 \end{tabularx} \\ \\ \relax
1:7 & \begin{tabularx}{0.7\textwidth}{X} 亞倫祭司的子孫要在壇上生火，把柴擺在火上。 \end{tabularx} \\ \\ \relax
1:8 & \begin{tabularx}{0.7\textwidth}{X} 亞倫子孫作祭司的要把肉塊連頭和脂肪，擺在壇上燒著火的柴上。 \end{tabularx} \\ \\ \relax
1:9 & \begin{tabularx}{0.7\textwidth}{X} 燔祭牲的內臟與小腿要用水洗淨，祭司要把整隻全燒在壇上，當作燔祭，是獻給耶和華為馨香的火祭。 \end{tabularx} \\ \\ \relax
1:10 & \begin{tabularx}{0.7\textwidth}{X} 「人的供物若以綿羊或山羊為燔祭，要獻一隻沒有殘疾的公羊。 \end{tabularx} \\ \\ \relax
1:11 & \begin{tabularx}{0.7\textwidth}{X} 他要在壇的北邊，在耶和華面前宰羊；亞倫子孫作祭司的要把血灑在壇的周圍。 \end{tabularx} \\ \\ \relax
1:12 & \begin{tabularx}{0.7\textwidth}{X} 他要把燔祭牲切成塊，祭司就要把肉塊連頭和脂肪，擺在壇上燒著火的柴上。 \end{tabularx} \\ \\ \relax
1:13 & \begin{tabularx}{0.7\textwidth}{X} 內臟與小腿要用水洗淨，祭司要把整隻獻上，全燒在壇上。這是燔祭，是獻給耶和華為馨香的火祭。 \end{tabularx} \\ \\ \relax
1:14 & \begin{tabularx}{0.7\textwidth}{X} 「人獻給耶和華的供物若以鳥為燔祭，就要獻斑鳩或雛鴿為他的供物。 \end{tabularx} \\ \\ \relax
1:15 & \begin{tabularx}{0.7\textwidth}{X} 祭司要把鳥拿到壇前，扭斷牠的頭，把鳥燒在壇上，鳥的血要流在壇的旁邊； \end{tabularx} \\ \\ \relax
1:16 & \begin{tabularx}{0.7\textwidth}{X} 又要把鳥的嗉囊和裡面的髒物除掉，丟在壇東邊倒灰的地方。 \end{tabularx} \\ \\ \relax
1:17 & \begin{tabularx}{0.7\textwidth}{X} 他要拿著鳥的兩個翅膀，把鳥撕開，卻不可撕斷；祭司要把牠擺在壇上燒著火的柴上焚燒。這是燔祭，是獻給耶和華為馨香的火祭。」 \end{tabularx} \\ \\ \relax
2:1 & \begin{tabularx}{0.7\textwidth}{X} 「若有人獻素祭為供物給耶和華，就要獻細麵為供物，把油澆在上面，加上乳香， \end{tabularx} \\ \\ \relax
2:2 & \begin{tabularx}{0.7\textwidth}{X} 帶到亞倫子孫作祭司的那裡。祭司要從細麵中取出滿滿的一把，又取些油和所有的乳香，把這些作為紀念的燒在壇上，是獻給耶和華為馨香的火祭。 \end{tabularx} \\ \\ \relax
2:3 & \begin{tabularx}{0.7\textwidth}{X} 素祭所剩的要歸給亞倫和他的子孫；在獻給耶和華的火祭中，這是至聖的。 \end{tabularx} \\ \\ \relax
2:4 & \begin{tabularx}{0.7\textwidth}{X} 「若獻爐中烤的素祭為供物，要用調了油的無酵細麵餅，或抹了油的無酵薄餅。 \end{tabularx} \\ \\ \relax
2:5 & \begin{tabularx}{0.7\textwidth}{X} 若以鐵盤上的素祭為供物，就要用調了油的無酵細麵， \end{tabularx} \\ \\ \relax
2:6 & \begin{tabularx}{0.7\textwidth}{X} 分成小塊，澆上油；這是素祭。 \end{tabularx} \\ \\ \relax
2:7 & \begin{tabularx}{0.7\textwidth}{X} 若以煎鍋煎的素祭為供物，就要用油與細麵做成。 \end{tabularx} \\ \\ \relax
2:8 & \begin{tabularx}{0.7\textwidth}{X} 要把這樣做成的素祭帶到耶和華面前，拿給祭司，祭司要帶到壇前。 \end{tabularx} \\ \\ \relax
2:9 & \begin{tabularx}{0.7\textwidth}{X} 祭司要從素祭中取出作為紀念的燒在壇上，是獻給耶和華為馨香的火祭。 \end{tabularx} \\ \\ \relax
2:10 & \begin{tabularx}{0.7\textwidth}{X} 素祭所剩的要歸給亞倫和他的子孫；在獻給耶和華的火祭中，這是至聖的。 \end{tabularx} \\ \\ \relax
2:11 & \begin{tabularx}{0.7\textwidth}{X} 「凡獻給耶和華的素祭都不可以有酵，因為你們不可把任何的酵或蜜燒了，當作火祭獻給耶和華。 \end{tabularx} \\ \\ \relax
2:12 & \begin{tabularx}{0.7\textwidth}{X} 你們可以把這些獻給耶和華當作初熟的供物，但是不可獻在壇上作為馨香的祭。 \end{tabularx} \\ \\ \relax
2:13 & \begin{tabularx}{0.7\textwidth}{X} 凡獻為素祭的供物都要用鹽調和；在素祭中，不可缺少你與神立約的鹽。一切的供物都要加鹽獻上。 \end{tabularx} \\ \\ \relax
2:14 & \begin{tabularx}{0.7\textwidth}{X} 「你若獻初熟之物給耶和華為素祭，就要獻在火中烘過的新麥穗，就是磨碎的新穀物，當作初熟之物的素祭。 \end{tabularx} \\ \\ \relax
2:15 & \begin{tabularx}{0.7\textwidth}{X} 你要加上油和乳香；這是素祭。 \end{tabularx} \\ \\ \relax
2:16 & \begin{tabularx}{0.7\textwidth}{X} 祭司要把供物中作為紀念的，就是一些磨碎的新穀物和一些油，以及所有的乳香，都焚燒，是獻給耶和華的火祭。」 \end{tabularx} \\ \\ \relax
3:1 & \begin{tabularx}{0.7\textwidth}{X} 「人獻平安祭為供物，若是從牛群中獻，無論是公的母的，要用沒有殘疾的，獻在耶和華面前。 \end{tabularx} \\ \\ \relax
3:2 & \begin{tabularx}{0.7\textwidth}{X} 他要按手在供物的頭上，在會幕的門口宰了牠。亞倫子孫作祭司的，要把血灑在壇的周圍。 \end{tabularx} \\ \\ \relax
3:3 & \begin{tabularx}{0.7\textwidth}{X} 從平安祭中，他要把火祭獻給耶和華，就是包著內臟的脂肪和內臟上所有的脂肪， \end{tabularx} \\ \\ \relax
3:4 & \begin{tabularx}{0.7\textwidth}{X} 兩個腎和腎上的脂肪，即靠近腎兩旁的脂肪，以及肝上的網油，連同腎一起取下。 \end{tabularx} \\ \\ \relax
3:5 & \begin{tabularx}{0.7\textwidth}{X} 亞倫的子孫要把這些擺在燒著火的柴上，燒在壇的燔祭上，是獻給耶和華為馨香的火祭。 \end{tabularx} \\ \\ \relax
3:6 & \begin{tabularx}{0.7\textwidth}{X} 「人向耶和華獻平安祭為供物，若是從羊群中獻，無論是公的母的，要用沒有殘疾的。 \end{tabularx} \\ \\ \relax
3:7 & \begin{tabularx}{0.7\textwidth}{X} 若他獻一隻綿羊為供物，就要把牠獻在耶和華面前。 \end{tabularx} \\ \\ \relax
3:8 & \begin{tabularx}{0.7\textwidth}{X} 要按手在供物的頭上，在會幕前宰了牠。亞倫的子孫要把血灑在壇的周圍。 \end{tabularx} \\ \\ \relax
3:9 & \begin{tabularx}{0.7\textwidth}{X} 從平安祭中，他要取脂肪當作火祭獻給耶和華，就是靠近脊骨處取下的整條肥尾巴，包著內臟的脂肪和內臟上所有的脂肪， \end{tabularx} \\ \\ \relax
3:10 & \begin{tabularx}{0.7\textwidth}{X} 兩個腎和腎上的脂肪，即靠近腎兩旁的脂肪，以及肝上的網油，連同腎一起取下。 \end{tabularx} \\ \\ \relax
3:11 & \begin{tabularx}{0.7\textwidth}{X} 祭司要把這些燒在壇上，是獻給耶和華為食物的火祭。 \end{tabularx} \\ \\ \relax
3:12 & \begin{tabularx}{0.7\textwidth}{X} 「人的供物若是山羊，就要把牠獻在耶和華面前。 \end{tabularx} \\ \\ \relax
3:13 & \begin{tabularx}{0.7\textwidth}{X} 要按手在牠的頭上，在會幕前宰了牠。亞倫的子孫要把血灑在壇的周圍， \end{tabularx} \\ \\ \relax
3:14 & \begin{tabularx}{0.7\textwidth}{X} 又要從供物中把火祭獻給耶和華，就是包著內臟的脂肪和內臟上所有的脂肪， \end{tabularx} \\ \\ \relax
3:15 & \begin{tabularx}{0.7\textwidth}{X} 兩個腎和腎上的脂肪，即靠近腎兩旁的脂肪，以及肝上的網油，連同腎一起取下。 \end{tabularx} \\ \\ \relax
3:16 & \begin{tabularx}{0.7\textwidth}{X} 祭司要把這些燒在壇上，作為馨香火祭的食物；所有的脂肪都是耶和華的。 \end{tabularx} \\ \\ \relax
3:17 & \begin{tabularx}{0.7\textwidth}{X} 在你們一切的住處，脂肪和血都不可吃，這要成為你們世世代代永遠的定例。」 \end{tabularx} \\ \\ \relax
4:1 & \begin{tabularx}{0.7\textwidth}{X} 耶和華吩咐摩西說： \end{tabularx} \\ \\ \relax
4:2 & \begin{tabularx}{0.7\textwidth}{X} 「你要吩咐以色列人說：若有人無意中犯罪，在任何事上犯了一條耶和華所吩咐的禁令， \end{tabularx} \\ \\ \relax
4:3 & \begin{tabularx}{0.7\textwidth}{X} 或是受膏的祭司犯了罪，使百姓陷在罪裡，他就當為自己所犯的罪，把沒有殘疾的公牛犢獻給耶和華為贖罪祭。 \end{tabularx} \\ \\ \relax
4:4 & \begin{tabularx}{0.7\textwidth}{X} 他要把公牛牽到會幕的門口，在耶和華面前按手在牛的頭上，把牛宰於耶和華面前。 \end{tabularx} \\ \\ \relax
4:5 & \begin{tabularx}{0.7\textwidth}{X} 受膏的祭司要取些公牛的血，帶到會幕那裡。 \end{tabularx} \\ \\ \relax
4:6 & \begin{tabularx}{0.7\textwidth}{X} 祭司要把手指蘸在血中，在耶和華面前對著聖所的幔子彈血七次， \end{tabularx} \\ \\ \relax
4:7 & \begin{tabularx}{0.7\textwidth}{X} 又要把一些血抹在會幕內，耶和華面前香壇的四個翹角上，再把公牛其餘的血全倒在會幕門口燔祭壇的底座上； \end{tabularx} \\ \\ \relax
4:8 & \begin{tabularx}{0.7\textwidth}{X} 又要取出這頭贖罪祭公牛所有的脂肪，就是包著內臟的脂肪和內臟上所有的脂肪， \end{tabularx} \\ \\ \relax
4:9 & \begin{tabularx}{0.7\textwidth}{X} 兩個腎和腎上的脂肪，即靠近腎兩旁的脂肪，以及肝上的網油，連同腎一起取下， \end{tabularx} \\ \\ \relax
4:10 & \begin{tabularx}{0.7\textwidth}{X} 正如從平安祭的牛身上所取的，祭司要把這些燒在燔祭壇上。 \end{tabularx} \\ \\ \relax
4:11 & \begin{tabularx}{0.7\textwidth}{X} 但公牛的皮和所有的肉，以及頭、腿、內臟、糞， \end{tabularx} \\ \\ \relax
4:12 & \begin{tabularx}{0.7\textwidth}{X} 就是全公牛，要搬到營外清潔的地方倒灰之處，放在柴上用火焚燒。 \end{tabularx} \\ \\ \relax
4:13 & \begin{tabularx}{0.7\textwidth}{X} 「以色列全會眾若犯了錯，在任何事上犯了一條耶和華所吩咐的禁令，而有了罪，會眾看不出這隱藏的事； \end{tabularx} \\ \\ \relax
4:14 & \begin{tabularx}{0.7\textwidth}{X} 他們一知道犯了罪，就要獻一頭公牛犢為贖罪祭，牽牠到會幕前。 \end{tabularx} \\ \\ \relax
4:15 & \begin{tabularx}{0.7\textwidth}{X} 會眾的長老要在耶和華面前按手在公牛的頭上，把牛宰於耶和華面前。 \end{tabularx} \\ \\ \relax
4:16 & \begin{tabularx}{0.7\textwidth}{X} 受膏的祭司要取些公牛的血，帶到會幕那裡。 \end{tabularx} \\ \\ \relax
4:17 & \begin{tabularx}{0.7\textwidth}{X} 祭司要用手指蘸一些血，在耶和華面前對著幔子彈七次， \end{tabularx} \\ \\ \relax
4:18 & \begin{tabularx}{0.7\textwidth}{X} 又要把一些血抹在會幕內，耶和華面前壇的四個翹角上，再把其餘的血全倒在會幕門口燔祭壇的底座上。 \end{tabularx} \\ \\ \relax
4:19 & \begin{tabularx}{0.7\textwidth}{X} 他要取出公牛所有的脂肪，燒在壇上。 \end{tabularx} \\ \\ \relax
4:20 & \begin{tabularx}{0.7\textwidth}{X} 他要處理這牛，正如處理那頭贖罪祭的公牛一樣，他要如此去做。祭司要為他們贖罪，他們就蒙赦免。 \end{tabularx} \\ \\ \relax
4:21 & \begin{tabularx}{0.7\textwidth}{X} 他要把牛搬到營外燒了，像燒前一頭公牛一樣；這是會眾的贖罪祭。 \end{tabularx} \\ \\ \relax
4:22 & \begin{tabularx}{0.7\textwidth}{X} 「官長若犯罪，在任何事上無意中犯了一條耶和華－他的神所吩咐的禁令，而有了罪， \end{tabularx} \\ \\ \relax
4:23 & \begin{tabularx}{0.7\textwidth}{X} 他一知道自己犯了罪，就要牽一隻沒有殘疾的公山羊為供物。 \end{tabularx} \\ \\ \relax
4:24 & \begin{tabularx}{0.7\textwidth}{X} 他要按手在羊的頭上，在耶和華面前宰燔祭牲的地方把牠宰了；這是贖罪祭。 \end{tabularx} \\ \\ \relax
4:25 & \begin{tabularx}{0.7\textwidth}{X} 祭司要用手指蘸一些贖罪祭牲的血，抹在燔祭壇的四個翹角上，再把其餘的血倒在燔祭壇的底座上。 \end{tabularx} \\ \\ \relax
4:26 & \begin{tabularx}{0.7\textwidth}{X} 祭牲所有的脂肪都要燒在壇上，正如平安祭的脂肪一樣。祭司要為他的罪贖了他，他就蒙赦免。 \end{tabularx} \\ \\ \relax
4:27 & \begin{tabularx}{0.7\textwidth}{X} 「這地的百姓若有人無意中犯罪，在任何事上犯了一條耶和華所吩咐的禁令，而有了罪， \end{tabularx} \\ \\ \relax
4:28 & \begin{tabularx}{0.7\textwidth}{X} 他一知道自己犯了罪，就要為所犯的罪牽一隻沒有殘疾的母山羊為供物。 \end{tabularx} \\ \\ \relax
4:29 & \begin{tabularx}{0.7\textwidth}{X} 他要按手在贖罪祭牲的頭上，在燔祭牲的地方把牠宰了。 \end{tabularx} \\ \\ \relax
4:30 & \begin{tabularx}{0.7\textwidth}{X} 祭司要用手指蘸一些祭牲的血，抹在燔祭壇的四個翹角上，再把其餘的血全倒在壇的底座上； \end{tabularx} \\ \\ \relax
4:31 & \begin{tabularx}{0.7\textwidth}{X} 又要把祭牲所有的脂肪都取下，正如取平安祭牲的脂肪一樣。祭司要把脂肪燒在壇上，在耶和華面前作為馨香的祭。祭司要為他贖罪，他就蒙赦免。 \end{tabularx} \\ \\ \relax
4:32 & \begin{tabularx}{0.7\textwidth}{X} 「人若牽一隻綿羊為贖罪祭作供物，就要牽一隻沒有殘疾的母羊。 \end{tabularx} \\ \\ \relax
4:33 & \begin{tabularx}{0.7\textwidth}{X} 他要按手在贖罪祭牲的頭上，在宰燔祭牲的地方宰了牠，作為贖罪祭。 \end{tabularx} \\ \\ \relax
4:34 & \begin{tabularx}{0.7\textwidth}{X} 祭司要用手指蘸一些贖罪祭牲的血，抹在燔祭壇的四個翹角上，再把其餘的血全倒在壇的底座上； \end{tabularx} \\ \\ \relax
4:35 & \begin{tabularx}{0.7\textwidth}{X} 又要把祭牲所有的脂肪都取下，正如取平安祭的羊的脂肪一樣。祭司要按獻給耶和華火祭的條例，把脂肪燒在壇上。祭司要為他所犯的罪贖了他，他就蒙赦免。」 \end{tabularx} \\ \\ \relax
5:1 & \begin{tabularx}{0.7\textwidth}{X} 「若有人犯了罪，就是聽見了誓言，他本來可以作證，卻不把所看見、所知道的說出來，必須擔當他的罪孽。 \end{tabularx} \\ \\ \relax
5:2 & \begin{tabularx}{0.7\textwidth}{X} 若有人摸了任何不潔之物，無論是野獸的不潔屍體，家畜的不潔屍體，或是群聚動物的不潔屍體，他雖不察覺，也是不潔淨，就有罪了。 \end{tabularx} \\ \\ \relax
5:3 & \begin{tabularx}{0.7\textwidth}{X} 或是他摸了人的不潔之物，就是任何使人成為不潔的不潔之物，他雖不察覺，但一知道，就有罪了。 \end{tabularx} \\ \\ \relax
5:4 & \begin{tabularx}{0.7\textwidth}{X} 若有人隨口發誓，或出於惡意，或出於善意，這人無論在甚麼事上隨意發誓，雖不察覺，但一知道，就在這其中的一件事上有罪了。 \end{tabularx} \\ \\ \relax
5:5 & \begin{tabularx}{0.7\textwidth}{X} 當他在這其中的一件事上有罪的時候，就要承認所犯的罪， \end{tabularx} \\ \\ \relax
5:6 & \begin{tabularx}{0.7\textwidth}{X} 並要為所犯的罪，把他的贖愆祭牲，就是羊群中的一隻母綿羊或母山羊，獻給耶和華為贖罪祭，祭司要為他的罪贖了他。 \end{tabularx} \\ \\ \relax
5:7 & \begin{tabularx}{0.7\textwidth}{X} 「若他的力量不夠獻一隻綿羊，就要為所犯的罪，把兩隻斑鳩或是兩隻雛鴿獻給耶和華為贖愆祭：一隻作贖罪祭，一隻作燔祭。 \end{tabularx} \\ \\ \relax
5:8 & \begin{tabularx}{0.7\textwidth}{X} 他要把這些帶到祭司那裡，祭司就先把贖罪祭獻上，從鳥的頸項上扭斷牠的頭，但不把鳥撕斷。 \end{tabularx} \\ \\ \relax
5:9 & \begin{tabularx}{0.7\textwidth}{X} 祭司要把一些贖罪祭牲的血彈在祭壇的邊上，其餘的血要倒在壇的底座上；這是贖罪祭。 \end{tabularx} \\ \\ \relax
5:10 & \begin{tabularx}{0.7\textwidth}{X} 他要依照條例獻第二隻鳥為燔祭。祭司要為他所犯的罪贖了他，他就蒙赦免。 \end{tabularx} \\ \\ \relax
5:11 & \begin{tabularx}{0.7\textwidth}{X} 「他的力量若不夠獻兩隻斑鳩或兩隻雛鴿，就要為所犯的罪把供物，就是十分之一伊法細麵，獻上為贖罪祭；不可加上油，也不可加上乳香，因為這是贖罪祭。 \end{tabularx} \\ \\ \relax
5:12 & \begin{tabularx}{0.7\textwidth}{X} 他要把細麵帶到祭司那裡，祭司要取出滿滿的一把，作為紀念，按照獻火祭給耶和華的條例把它燒在壇上；這是贖罪祭。 \end{tabularx} \\ \\ \relax
5:13 & \begin{tabularx}{0.7\textwidth}{X} 至於他在這幾件事中所犯的任何罪，祭司要為他贖了，他就蒙赦免。剩下的都歸給祭司，和素祭一樣。」 \end{tabularx} \\ \\ \relax
5:14 & \begin{tabularx}{0.7\textwidth}{X} 耶和華吩咐摩西說： \end{tabularx} \\ \\ \relax
5:15 & \begin{tabularx}{0.7\textwidth}{X} 「若有人在耶和華的聖物上無意中犯了罪，有了過犯，就要獻羊群中一隻沒有殘疾的公綿羊給耶和華為贖愆祭，或依聖所的舍客勒所估定的銀子，作為贖愆祭。 \end{tabularx} \\ \\ \relax
5:16 & \begin{tabularx}{0.7\textwidth}{X} 他要為在聖物上的疏忽賠償，另外加五分之一，把這些都交給祭司。祭司要用贖愆祭的公綿羊為他贖罪，他就蒙赦免。 \end{tabularx} \\ \\ \relax
5:17 & \begin{tabularx}{0.7\textwidth}{X} 「若有人犯罪，在任何事上犯了一條耶和華所吩咐的禁令，他雖不察覺，仍算有罪，必須擔當自己的罪孽。 \end{tabularx} \\ \\ \relax
5:18 & \begin{tabularx}{0.7\textwidth}{X} 他要牽羊群中一隻沒有殘疾的公綿羊，或照你所估定的價值，給祭司作贖愆祭。祭司要為他贖他因不知道而無意中所犯的罪，他就蒙赦免。 \end{tabularx} \\ \\ \relax
5:19 & \begin{tabularx}{0.7\textwidth}{X} 這是贖愆祭；因他確實得罪了耶和華。」 \end{tabularx} \\ \\ \relax
6:1 & \begin{tabularx}{0.7\textwidth}{X} 耶和華吩咐摩西說： \end{tabularx} \\ \\ \relax
6:2 & \begin{tabularx}{0.7\textwidth}{X} 「若有人犯罪，得罪了耶和華，就是在鄰舍寄託他的東西或抵押品上行詭詐，或搶奪，或欺壓鄰舍， \end{tabularx} \\ \\ \relax
6:3 & \begin{tabularx}{0.7\textwidth}{X} 或是撿了失物行了詭詐，起了假誓，在人所做的任何事上犯了罪； \end{tabularx} \\ \\ \relax
6:4 & \begin{tabularx}{0.7\textwidth}{X} 他既犯了罪，有了過犯，就要歸還他所搶奪的，或是因欺壓所得的，或是別人寄託他的，或是他所撿到的失物， \end{tabularx} \\ \\ \relax
6:5 & \begin{tabularx}{0.7\textwidth}{X} 或是起假誓得來的任何東西，就要全數歸還，另外再加五分之一。在查出他有罪的日子，就要立刻賠還給原主。 \end{tabularx} \\ \\ \relax
6:6 & \begin{tabularx}{0.7\textwidth}{X} 他要獻羊群中一隻沒有殘疾的公綿羊，給耶和華為贖愆祭，或照你所估定的價值，給祭司作贖愆祭。 \end{tabularx} \\ \\ \relax
6:7 & \begin{tabularx}{0.7\textwidth}{X} 祭司要在耶和華面前為他贖罪；他無論做了甚麼事，以致有了罪，都必蒙赦免。」 \end{tabularx} \\ \\ \relax
6:8 & \begin{tabularx}{0.7\textwidth}{X} 耶和華吩咐摩西說： \end{tabularx} \\ \\ \relax
6:9 & \begin{tabularx}{0.7\textwidth}{X} 「你要吩咐亞倫和他的子孫說，燔祭的條例是這樣：燔祭要放在壇的底盤上，從晚上到天亮，壇上的火要不斷地燒著。 \end{tabularx} \\ \\ \relax
6:10 & \begin{tabularx}{0.7\textwidth}{X} 祭司要穿上細麻布衣服，又要把細麻布褲子穿在身上，把在壇上燒剩的燔祭灰收起來，放在壇的旁邊。 \end{tabularx} \\ \\ \relax
6:11 & \begin{tabularx}{0.7\textwidth}{X} 然後，他要脫去這衣服，穿上別的衣服，把灰拿到營外潔淨之處。 \end{tabularx} \\ \\ \relax
6:12 & \begin{tabularx}{0.7\textwidth}{X} 壇上的火要不斷地燒著，不可熄滅。每日早晨，祭司要在壇上燒柴，把燔祭擺在壇上，並燒平安祭牲的脂肪。 \end{tabularx} \\ \\ \relax
6:13 & \begin{tabularx}{0.7\textwidth}{X} 壇上的火要不斷地燒著，不可熄滅。」 \end{tabularx} \\ \\ \relax
6:14 & \begin{tabularx}{0.7\textwidth}{X} 「素祭的條例是這樣：亞倫的子孫要在壇前把這祭獻在耶和華面前。 \end{tabularx} \\ \\ \relax
6:15 & \begin{tabularx}{0.7\textwidth}{X} 祭司要從素祭中的細麵取出一把，再取些油和素祭上所有的乳香，把這些作為紀念的燒在壇上，是獻給耶和華為馨香的祭。 \end{tabularx} \\ \\ \relax
6:16 & \begin{tabularx}{0.7\textwidth}{X} 亞倫和他子孫要吃素祭剩下的；要在聖處吃這無酵餅，在會幕的院子裡吃。 \end{tabularx} \\ \\ \relax
6:17 & \begin{tabularx}{0.7\textwidth}{X} 烤餅不可加酵。這是從獻給我的火祭中歸給他們的一份；如贖罪祭和贖愆祭一樣，這份是至聖的。 \end{tabularx} \\ \\ \relax
6:18 & \begin{tabularx}{0.7\textwidth}{X} 亞倫子孫中的男丁都要吃，因為這是你們世世代代從獻給耶和華的火祭中，他們永遠應得的份。凡摸這些祭物的都要成為聖。」 \end{tabularx} \\ \\ \relax
6:19 & \begin{tabularx}{0.7\textwidth}{X} 耶和華吩咐摩西說： \end{tabularx} \\ \\ \relax
6:20 & \begin{tabularx}{0.7\textwidth}{X} 「這是亞倫受膏的日子，他和他的子孫所要獻給耶和華的供物：十分之一伊法細麵，如他們經常獻的素祭，早晨一半，晚上一半。 \end{tabularx} \\ \\ \relax
6:21 & \begin{tabularx}{0.7\textwidth}{X} 要在鐵盤上用油調和，調勻後，就拿去烤。素祭烤熟了要分成小塊，作為獻給耶和華馨香的祭。 \end{tabularx} \\ \\ \relax
6:22 & \begin{tabularx}{0.7\textwidth}{X} 亞倫子孫中接續他受膏為祭司的，要把這素祭獻上，全燒給耶和華。這是永遠的定例。 \end{tabularx} \\ \\ \relax
6:23 & \begin{tabularx}{0.7\textwidth}{X} 祭司一切的素祭要全部燒了，不可以吃。」 \end{tabularx} \\ \\ \relax
6:24 & \begin{tabularx}{0.7\textwidth}{X} 耶和華吩咐摩西說： \end{tabularx} \\ \\ \relax
6:25 & \begin{tabularx}{0.7\textwidth}{X} 「你要吩咐亞倫和他的子孫說，贖罪祭的條例是這樣：要在耶和華面前宰燔祭牲的地方宰贖罪祭牲；這是至聖的。 \end{tabularx} \\ \\ \relax
6:26 & \begin{tabularx}{0.7\textwidth}{X} 獻贖罪祭的祭司要吃這祭物；要在聖處，就是在會幕的院子裡吃。 \end{tabularx} \\ \\ \relax
6:27 & \begin{tabularx}{0.7\textwidth}{X} 凡摸這祭肉的都要成為聖；這祭牲的血若濺在衣服上，你要在聖處洗淨那濺到血的衣服。 \end{tabularx} \\ \\ \relax
6:28 & \begin{tabularx}{0.7\textwidth}{X} 煮這祭物的瓦器要打碎；若祭物是在銅器裡煮，要把這銅器擦淨，用水沖洗。 \end{tabularx} \\ \\ \relax
6:29 & \begin{tabularx}{0.7\textwidth}{X} 祭司中所有的男丁都可以吃；這是至聖的。 \end{tabularx} \\ \\ \relax
6:30 & \begin{tabularx}{0.7\textwidth}{X} 若將任何贖罪祭的血帶進會幕，為要在聖所贖罪，那肉就不可吃，要用火焚燒。」 \end{tabularx} \\ \\ \relax
7:1 & \begin{tabularx}{0.7\textwidth}{X} 「贖愆祭的條例是這樣：這祭是至聖的。 \end{tabularx} \\ \\ \relax
7:2 & \begin{tabularx}{0.7\textwidth}{X} 人在哪裡宰燔祭牲，也要在哪裡宰贖愆祭牲；其血，祭司要灑在壇的周圍。 \end{tabularx} \\ \\ \relax
7:3 & \begin{tabularx}{0.7\textwidth}{X} 祭司要獻上牠所有的脂肪，把肥尾巴和包著內臟的脂肪， \end{tabularx} \\ \\ \relax
7:4 & \begin{tabularx}{0.7\textwidth}{X} 兩個腎和腎上的脂肪，即靠近腎兩旁的脂肪，以及肝上的網油，連同腎一起取下。 \end{tabularx} \\ \\ \relax
7:5 & \begin{tabularx}{0.7\textwidth}{X} 祭司要把這些燒在壇上，獻給耶和華為火祭，作為贖愆祭。 \end{tabularx} \\ \\ \relax
7:6 & \begin{tabularx}{0.7\textwidth}{X} 祭司中所有的男丁都可以吃這祭物，要在聖處吃；這是至聖的。 \end{tabularx} \\ \\ \relax
7:7 & \begin{tabularx}{0.7\textwidth}{X} 贖罪祭怎樣，贖愆祭也是怎樣，都有一樣的條例，用贖愆祭贖罪的祭司要得這祭物。 \end{tabularx} \\ \\ \relax
7:8 & \begin{tabularx}{0.7\textwidth}{X} 獻燔祭的祭司，無論為誰獻，所獻燔祭牲的皮要歸給那祭司，那是他的。 \end{tabularx} \\ \\ \relax
7:9 & \begin{tabularx}{0.7\textwidth}{X} 任何素祭，無論是在爐中烤的，用煎鍋或鐵盤做成的，都要歸給獻祭的祭司。 \end{tabularx} \\ \\ \relax
7:10 & \begin{tabularx}{0.7\textwidth}{X} 任何素祭，無論是用油調和的，是乾的，都要歸亞倫的子孫，大家均分。」 \end{tabularx} \\ \\
[1ex]
\hline
\hline
\end{longtable}
舊約諸種祭禮,是與基督有極關係之預表,並對信徒有極重要之靈訓,我們研究五祭所得的結果就是:
\newline
\newline
┌ 預表......完全救贖的「耶穌基督」
\newline
\newline
第二課以後,每個論題下之論文,不外是解釋如何預表耶穌基督,及如何教訓信徒,此後不必每條再書如何預表,如何教訓,可將預表與教訓混合而言,請讀者留意.
\newline
\newline
出埃及記,利未記,民數記及舊約歷史書中之諸般祭禮,雖甚覺繁雜,然亦不過總歸此五種祭禮,惟獻祭的時候,事故,分量,有分別耳.
\newline
\newline
(壹)何為五祭?
\newline
\newline
(貳)五祭如何分類
\newline
\newline
(參)五祭預表之總意如何
\newline
\newline
聯合五種祭品之預表,可以看見一個完全之救贖者,耶穌基督,如人立於五面大鏡之中,可見其人之全身,我們研究耶穌,亦須由面面觀察,方可見耶穌之為完全神人.
\newline
\newline
(肆)五祭之次序如何?有神人兩方面的看法不同
\newline
\newline


\newpage

\section{第1屆~港九培靈會~黃原素~閉會講道}
\label{sec:1321}
\textbf{閉會講道}
\newline
\newline
連結: \href{https://www.hkbibleconference.org/session-message/view/1321}{\texttt{https://www.hkbibleconference.org/session-message/view/1321}}
\newline
\newline
\hyperref[sec:1334]{< < < PREV SERMON < < <}
~
\hyperref[sec:toc]{[返目錄]}
~
\hyperref[sec:1335]{> > > NEXT SERMON > > >}
\newline
\newline
何西阿書 6:1-6:6
\newline
\begin{longtable}{cl}
\hline
\hline
章節 & 經文 (和合本修訂版)\\
\hline
6:1 & \begin{tabularx}{0.7\textwidth}{X} 來，我們歸向耶和華吧！他撕裂我們，也必醫治；打傷我們，也必包紮。 \end{tabularx} \\ \\ \relax
6:2 & \begin{tabularx}{0.7\textwidth}{X} 過兩天他必使我們甦醒，第三天他必使我們興起，我們就在他面前得以存活。 \end{tabularx} \\ \\ \relax
6:3 & \begin{tabularx}{0.7\textwidth}{X} 我們要認識，要追求認識耶和華。他如黎明必然出現，他必臨到我們像甘霖，像滋潤土地的春雨。 \end{tabularx} \\ \\ \relax
6:4 & \begin{tabularx}{0.7\textwidth}{X} 以法蓮哪，我可以向你怎樣行呢？猶大啊，我可以向你怎樣做呢？因為你們的慈愛如同早晨的雲霧，又如速散的露水。 \end{tabularx} \\ \\ \relax
6:5 & \begin{tabularx}{0.7\textwidth}{X} 因此，我藉先知砍伐他們，以我口中的話殺戮他們；對你的審判如光發出。 \end{tabularx} \\ \\ \relax
6:6 & \begin{tabularx}{0.7\textwidth}{X} 我喜愛慈愛，不喜愛祭物；喜愛人認識神，勝於燔祭。 \end{tabularx} \\ \\
[1ex]
\hline
\hline
\end{longtable}
今我們來赴這培靈會,我知你們各人的心如同我的心一樣,就是想在這會中得勝靈重新充滿我們的心,好使我們能作主的真代表.正有一位兄弟禱告,求神使他能作主的真代表,他這樣禱告的話甚感我心.我們今來赴這會,不是代表他人,乃是代表自己.今來這會,是為己來,不是為他人來.然而,立這樣的心,可算是私心麼?不,不過是先求奮興自己然後使己能夠去奮興他人.
\newline
\newline
今年培靈會第三會會訓,就是何西阿六章三節.在這節聖經裡顯明有三件是以色列人當日所需要的,玆把他分述如下:
\newline
\newline
(一)是光
\newline
\newline
以色列人第一所需要的是晨光,是這必來的晨光.在這裡也許有人說,來赴這會不知到能不能得著這樣的福,恐怕主不會向我顯現,我不會向主得著這麼大的福.這樣麼?各位,你若是這麼說,敢問你來赴這會,是仰望人呢?還是仰望主呢?請你不要疑慮罷!因這節聖經明說,主必確現如晨光,這確現在昔是這樣的,在今也是這樣的.當日以色列人為甚麼不得著光呢?不是因光不確現,這是因他們退後與背光啊!故何西阿時代,就是以色列人的黑暗時代.由這想來,昔日以色列人得不著光因這樣,今日我們得不著光也因這樣.是以我們要小心,不要像以色列人這樣的退後背主啊!
\newline
\newline
(二)是雨
\newline
\newline
以色列人第二所需要的是甘雨.試問香港這數月來所缺乏的是甚麼呢?你們必回答說是大雨.這以色列所欠的既是雨,故結果就使他們的心靈乾涸.敢問你們在這數月來心裡有沒有雨?你們的心若同樣的欠雨,結果就必同以色列人一樣的乾涸!我曾向一位牧師問及他的靈況怎樣?他回答說,『沒有進步,也沒有退後,穩穩陣陣的.』阿!各位,這樣的『穩陣』你可願意得著嗎?我不信這不進不退也得穩陣的話,因為我信不進的就必退後,這可拿單車來作例證.各位,我們的心靈若不是潤澤的,這就必乾涸；決不會不潤澤的還會不乾涸的.敢問你的心靈是潤澤的呢,還是乾涸的呢?你若是像一株欠水的樹,就定然不會結果!
\newline
\newline
(三)是生命
\newline
\newline
以色列人不只欠『光』,欠『雨』,更還欠『生』命.我昨年在港見一位數十年的老傳道,他對我說,『我教會無不信主的人,所有的都是高舉聖經的；雖是如此,但惜總無生氣!』各位,以色列人當時也是崇拜神,(何六之六)但惜祇有外貌的崇拜,惟實際則不識神.我知教會中外貌的信徒則多,而具有屬靈生命的信徒則少!
\newline
\newline
各位,以色列人想得著『光』,『雨』,和『生命』,就該到那處呢?你們定然會答說,到神處.故他們既背光而向黑暗,若要再得神恩,就應該回歸神處.可惜以色列人要『光』,要『雨』,卻不肯歸神,不肯歸『光』,『雨』,與『生命』源頭的神!
\newline
\newline
何西阿五章十三節說:『以法蓮見自己有病,猶大見自己有傷,他們就打發人往亞述去見耶雷布王,他卻不能醫治你們,不能治好你們的傷.』這裡看來,可見以色列人棄能醫治他們的神,而求不能醫治他們的亞述王,他們的愚昧是何等的深!這真要使後世的人為他們深致慨嘆!兄弟們:我們有時也有靈病,然而要走去那處求醫呢?求埃及還是求神呢?從前有一位傳道人,心靈憂傷,這因他犯了不與人和的罪.一日有一位牧者問他,『你的容色與前不同,為甚麼呢?』那傳道回問他,『我該要甚麼才能得回我所失的呢?』可惜該牧者答他說,『你當去多多作工,』不知道作工也不能療其病,這正同求亞述一樣,在這世上,實無別的法子能療治我們的心病,祇有一法,這就是『來.』來作甚麼?就是『來歸.』來歸那處?試問到亞述處是歸麼?到伊及處是歸麼?不是!惟回到神處那才是歸處呢!因為只有神才能醫治我們.所以何西阿六章一節大聲的宣告說:『來罷,我們歸向耶和華；他撕裂我們,也必醫治；他打傷我們,也必纏裹.』
\newline
\newline
各位,我們要注意路加十五章所載的浪子,當他離父家時,行走的路途是灣曲的,遼遠的；到後來他醒悟起來的時候,他就說,『我要起來,到我父親那裡去,』那時他回家的路途是直捷的.請問,當浪子離家後,他的父親有沒有遷家呢?沒有.在這看來,就知我們的神不會遷家,遷家這事,祇有浪子才會!然而父為甚麼不遷家呢?我們可簡捷的回答說,因為等待浪子回家.他不但不遷家來等待,且更日夜倚門倚閭的企望啊!
\newline
\newline
又我們讀路得記,就見路得和著家姑拿俄米等,初住在伯利恆那裡,後因遭遇飢荒,就往摩押地那裡寄居.請問,摩押地是不是神的家?不是.後來拿俄米在摩押地,聽聞耶和華眷顧他自己的百姓,賜糧食過他們.他就和兩個兒婦起行,離開所住的摩押地,而回猶大地去.在這裡看來,可見拿俄米是有何等的智慧.因為他一知道了錯就回歸,祇是一心的回歸,不再走別的路途.想今日教會中都有很多的人離了神,雖不是盡離,但都成為浪子,在父家外漂流!因是這樣,故聖經說『來!』現在我們若還不聽從,而再走別路,那就有禍了!但起而歸的,那就必得回他所當得的福!
\newline
\newline
主在以賽亞一章十八節說:『你們來,我們彼此辯論.』今早劉心慈先生說,『這會不是議事會,』這言是真的；但從別方面看來,這也是一議事會,是神與人的議事會.故今請你們不要理會別人的事,當要與神理清你自己的罪.這樣,你們不當說,『神呀,我有很多事還未理妥,我不暇來辯.』須知神不是叫你來辯論他人的罪,叫你來是辯論自己的罪.然而我們要知,我們的罪,雖比朱砂更紅些,但主也能白他如雪.今我們雖願中國教會在這五年內得加倍的人數,但我們若不與神理清我們的罪,神就因我們的罪而不奮興他的教會.試看撒母耳七章一至二節的記載,就是說神的約櫃,在基列耶琳山中,有二十年之久,到了那時,以色列全家,才傾向耶和華.當這時候,以色列人,只是哀求,哀求,求神憐己,但他自己總不憐自己,祇是憐罪惡,而不除罪惡；像這樣的求主可憐,求到永遠都沒用的!當時撒母耳見他們這樣,就對以色列全家說,『你們若是一心歸向耶和華,就要把外邦的神,和亞斯他錄,從你們中間除掉.』各位,你們若滿屋滿心有偶像,那就不要想向主得甚麼福.在鬱林那處有兩位姊妹,同入教會,經有十多年,有一位甚得主祝福,滿心喜樂；有一位則常說主不聽他禱告,又不向著他祝福.有一日,我妻往探她,見她屋內有一香爐,她知這香爐被我妻看見了,就面都紅起來.各位,你心內有甚麼物件阻礙你歸神呢?有沒有呢?知不知呢?若是有的,那就願你不獨知而且要除去.在何西阿五章四節裡說:『他們所行的使他們不能歸向上帝,因有淫心在他們裡面,他們也不認識耶和華.』在這裡所說的淫心,就是指以色列人,離棄上帝,而歸向偶像.這淫就是屬靈的淫.因為這樣,撒母耳就囑咐他們先把偶像除去,若是不除,那就定然不得神的祝福.
\newline
\newline
各位,你知自己的病麼?猶大知以法蓮也知.但你知否為甚麼失敗?為甚麼離開神?若知,那就要快的除去偶像,而一心歸向神如浪子一樣的立心回家.今望各位都除去屬靈的淫心而來歸神.在何西阿五章十五節看來,就知神是怎樣的忍耐等待他的兒女自覺.
\newline
\newline
各位,今你有甚麼事或甚麼物阻擋你得那『光』,『雨』,與『生命』?要知,神不獨不阻擋你,更一路等待你,如伯利恆等待拿俄米,浪子的父等待浪子.我知神必等待你,不過你若到這會,不立志除去偶像,又不清罪,你就不要想在這會會集裡得福.
\newline
\newline
末了,惟願望你立志離暗就光,而得靈雨滿心!
\newline
\newline


\newpage

\section{第1屆~港九培靈會~黃原素~第二講　燔祭之研究}
\label{sec:1335}
\textbf{第二講　燔祭之研究}
\newline
\newline
連結: \href{https://www.hkbibleconference.org/session-message/view/1335}{\texttt{https://www.hkbibleconference.org/session-message/view/1335}}
\newline
\newline
\hyperref[sec:1321]{< < < PREV SERMON < < <}
~
\hyperref[sec:toc]{[返目錄]}
~
\hyperref[sec:1336]{> > > NEXT SERMON > > >}
\newline
\newline
利未記 1:1-1:17
\newline
\begin{longtable}{cl}
\hline
\hline
章節 & 經文 (和合本修訂版)\\
\hline
1:1 & \begin{tabularx}{0.7\textwidth}{X} 耶和華從會幕中呼叫摩西，吩咐他說： \end{tabularx} \\ \\ \relax
1:2 & \begin{tabularx}{0.7\textwidth}{X} 「你要吩咐以色列人，對他們說：你們中間若有人要獻供物給耶和華，可以從牛群羊群中獻牲畜為供物。 \end{tabularx} \\ \\ \relax
1:3 & \begin{tabularx}{0.7\textwidth}{X} 「他的供物若以牛為燔祭，要獻一頭沒有殘疾的公牛，獻在會幕的門口，他就可以在耶和華面前蒙悅納。 \end{tabularx} \\ \\ \relax
1:4 & \begin{tabularx}{0.7\textwidth}{X} 他要按手在燔祭牲的頭上，為自己贖罪，就蒙悅納。 \end{tabularx} \\ \\ \relax
1:5 & \begin{tabularx}{0.7\textwidth}{X} 他要在耶和華面前宰公牛犢；亞倫子孫作祭司的要獻上血，把血灑在會幕門口壇的周圍。 \end{tabularx} \\ \\ \relax
1:6 & \begin{tabularx}{0.7\textwidth}{X} 他要剝去燔祭牲的皮，把燔祭牲切成塊。 \end{tabularx} \\ \\ \relax
1:7 & \begin{tabularx}{0.7\textwidth}{X} 亞倫祭司的子孫要在壇上生火，把柴擺在火上。 \end{tabularx} \\ \\ \relax
1:8 & \begin{tabularx}{0.7\textwidth}{X} 亞倫子孫作祭司的要把肉塊連頭和脂肪，擺在壇上燒著火的柴上。 \end{tabularx} \\ \\ \relax
1:9 & \begin{tabularx}{0.7\textwidth}{X} 燔祭牲的內臟與小腿要用水洗淨，祭司要把整隻全燒在壇上，當作燔祭，是獻給耶和華為馨香的火祭。 \end{tabularx} \\ \\ \relax
1:10 & \begin{tabularx}{0.7\textwidth}{X} 「人的供物若以綿羊或山羊為燔祭，要獻一隻沒有殘疾的公羊。 \end{tabularx} \\ \\ \relax
1:11 & \begin{tabularx}{0.7\textwidth}{X} 他要在壇的北邊，在耶和華面前宰羊；亞倫子孫作祭司的要把血灑在壇的周圍。 \end{tabularx} \\ \\ \relax
1:12 & \begin{tabularx}{0.7\textwidth}{X} 他要把燔祭牲切成塊，祭司就要把肉塊連頭和脂肪，擺在壇上燒著火的柴上。 \end{tabularx} \\ \\ \relax
1:13 & \begin{tabularx}{0.7\textwidth}{X} 內臟與小腿要用水洗淨，祭司要把整隻獻上，全燒在壇上。這是燔祭，是獻給耶和華為馨香的火祭。 \end{tabularx} \\ \\ \relax
1:14 & \begin{tabularx}{0.7\textwidth}{X} 「人獻給耶和華的供物若以鳥為燔祭，就要獻斑鳩或雛鴿為他的供物。 \end{tabularx} \\ \\ \relax
1:15 & \begin{tabularx}{0.7\textwidth}{X} 祭司要把鳥拿到壇前，扭斷牠的頭，把鳥燒在壇上，鳥的血要流在壇的旁邊； \end{tabularx} \\ \\ \relax
1:16 & \begin{tabularx}{0.7\textwidth}{X} 又要把鳥的嗉囊和裡面的髒物除掉，丟在壇東邊倒灰的地方。 \end{tabularx} \\ \\ \relax
1:17 & \begin{tabularx}{0.7\textwidth}{X} 他要拿著鳥的兩個翅膀，把鳥撕開，卻不可撕斷；祭司要把牠擺在壇上燒著火的柴上焚燒。這是燔祭，是獻給耶和華為馨香的火祭。」 \end{tabularx} \\ \\
[1ex]
\hline
\hline
\end{longtable}
一, 此祭教訓信徒崇拜主最深的工夫.
\newline
\newline
二, 祭牲之一切完全焚燒.(一章九節)
\newline
\newline
三, 特表耶穌向神為救贖的完全獻己.
\newline
\newline
四, 祭牲須獻者,親自剝皮切塊.
\newline
\newline
一,公牛(三至九節)二,綿羊或山羊(十至十三節)三,斑鳩或雛鴿.(十四至十七節)所以有如此之等別者,是因人之貧富力量而設.(主耶穌誕後八日,其母只攜雛鴿或斑鳩入殿獻祭,可知主在地上之家庭,是甚貧窮的.)茲述其預表與靈訓如下:
\newline
\newline
(甲)牛,力大能役,耐苦順服,可表耶穌之畢生工作,乃役於人.(可十章四十五節)至終順服至死,甚至死於十字架.(腓二章五至八節)
\newline
\newline
(乙)綿羊,性良而潔,耶穌在世之為人,溫柔清潔,被詬不反抗,受苦而忍.(賽五十三章七節)
\newline
\newline
(丙)山羊,性野,並指罪惡,(太廿五章卅三節)表耶穌因救人而成為罪人,(賽五十三章十二節)「列在罪犯之中」他與二盜同釘,已應此言,又為我們「成為罪,」(哥後五章廿一節)「為咒詛.」(迦三章十三節)
\newline
\newline
(丁)斑鳩雛鴿,亦性馴而潔,卑弱堪憐(賽卅八章十四節,表耶穌為人,亦處於卑微可憫之地位,且為我們成為貧.(路九章五十八節)
\newline
\newline
以上四種祭品,無論大小,均須無疵無瑕,否則不蒙悅納,我們於此種祭牲之事上,得有三個心靈的教訓:
\newline
\newline
一, 獻己於主之權利,人人可有.主所以許用鳩鴿,是恐人無力獻牛羊,然若真心按照規則而獻上,鳩鴿蒙納,與牛羊無異,信徒獻己,不論地位,智愚,長幼,若真是由愛而獻,以順服而保守,必蒙主悅也.
\newline
\newline
二, 獻己之工作,不重事工之大小,乃重愛心之多少.有人獻己後,作大奮興家,得多人歸主,在教會作成偉大的工；有人獻己後,在事工上無甚特色,他時常保守順服,或作個人傳道,或暗中代人祈禱,或為主名忍受諸般痛苦,如此之人,其在主前之馨香,實不遜於為主作大事之人也.
\newline
\newline
三, 獻己之條件.須耐苦服役,順服至終,品性溫良,常居卑弱,忍受苦難,如此四種祭牲之性者,然後得主悅納.
\newline
\newline
「耶和華面前」就是獻一切祭禮之處,摩西以前時代,各地隨建坵壇,崇拜鬼魔,故特區別一處,指定為拜耶和華之所,即在會幕前之祭壇.(利十七章三,四,五,八,九節)當日定了律法,凡宰牲或獻祭,不到耶和華面前獻給的,要從民中剪除,其若是嚴厲大概在新約有以下的兩個表意:
\newline
\newline
一,表十字架是人類拜主的獨一祭壇.但凡不托十字架上救主之名之一切祈禱,都不蒙納,且被憎惡,此等祭拜,是奉事鬼魔的,不惟如此,今日新派不信十字架上代死的恩為人類挽回之根基者,他們一切奉事之工作,都算徒然.
\newline
\newline
二,真正的崇拜,必須在主前,方蒙悅納.人前顯出熱心的奉事,輕忽鑑隱的主,此乃人所常犯,真正根基的福澤,不是在人前而得,乃由我們獨對主時,在其前傾吐我心,作甜密心交之崇事,此方是屬靈之大福.
\newline
\newline
(甲)親身攜祭牲至會幕前.一章三節
\newline
\newline
一,欲免罪得救,須親自承認十字架上之主,而作求救之工,自己不為,無人可代.
\newline
\newline
二,獻己之工,亦無人可代行.
\newline
\newline
三,我們至主前,必須蒙了基督的聖義,方蒙悅納.
\newline
\newline
(乙)親手按於祭牲之上.(一章四節)
\newline
\newline
一,區別之作聖品.約書亞于會眾前受按手禮,以區別之為領袖,(申卅四章九節)新約教會之按手禮,即倣此意,其意即于眾前,以此重大責任託之,而祝其成功.
\newline
\newline
二,區別之以證其死罪.證者按手於咒詛聖名的犯罪者,然後會眾打死他(利廿四章十四節)但此燔祭之按手,不同此意,乃表其本人之罪,盡託於此無辜白受冤苦的無疵之牲,而定其死,人的罪,牲代之死,基督亦然,白受天下人類之罪,而致死於十字架.
\newline
\newline
三,按手者與此牲聯屬不分,基督之死,亦即我們之死,我們有罪,基督亦同受罪報；主召我們,是要我們與基督相連屬,(哥前一章七節)所以有信徒與基督同苦,同死,同葬,同起,同起,同升,同榮等靈理.總之,主在世如何,我們亦如何,(約一書四章十七節)主云,「從我,」又云「學我,」此皆欲我們與他實行此相連屬之道也.
\newline
\newline
(丙)親手殺牲於主前.(一章五節)
\newline
\newline
此則的靈表與靈訓,就是基督之死,是因我們的罪,如經云,「基督為我們成了咒詛」(迦三章十三節)雖當日殺基督的,是出於猶太人之主意,羅馬人之權,其實是因我們的罪,所以基督之死,是我有份殺他,但千祈一殺,不可再殺.(來六章六節)
\newline
\newline
(丁)親手剝皮切塊.(一章六節)
\newline
\newline
一,剝皮表耶穌之內部清潔,即其心靈中無一點瑕疵,可以在主前完全表白顯露,內裡誠實,得主喜愛,(詩五十一篇六節)訓我們行事為人,內外皆須在主前呈露無遺,我們獻己之工,很易生上一層皮,皮外是獻了於主,皮內(心)是仍為己有,或內心受不起主的鑑察,至獻己的路,漸生障阻,以致失敗,故真正獻己,必當剝皮.
\newline
\newline
二,切塊比剝皮更深一層,剝皮是鑑察內心的工作,切塊是鑑察性根的工作,內心之是非,是自己良心知到的,性根之是非,良心所不知的,切塊是受主分析得透徹,細微鑑隱我們在主前,曾否受過這種電鑑之工作,又能否忍受之,切塊用刀,切性根要用主的道；(來四章十二,十三節)因主道活潑有力,比兩刃利刀更鋒利,甚至人的魂與靈,骨節與骨髓,(即言行與思想)皆能刺入剖開,人心中的思念與主意,亦皆辨明,萬物無一不要顯露於主前,就是人們魂靈間的混雜,亦用此刀而切析清楚.此剝皮切塊之工,其最大的條件,是內外無疵,我們獻己之條件,是內外無罪.
\newline
\newline
(甲)洒血之禮.(一章五節)
\newline
\newline
一,洒血是祭司獨有之專責,表流血代贖之工,是耶穌一己之專責；世人雖有諸多改良,救濟等工作,但對於流血代贖之工,則不能加上一指.獻祭者,雖有多種手續,但一至洒血之工,即須停手,讓祭司代行.
\newline
\newline
二,血,是全祭禮最注重之一點,血表生命,即基督之生命,為人類之罪而傾注於十字架.除素祭特表耶穌之人格外,一切祭禮,若不有洒血之工,主之心總不滿足,而人之罪亦不能免.耶穌之外,世界諸多先知,先見,聖賢,智,哲,救濟人民,敬拜上主,惟主之心,總不知足；須有耶穌一次流血,上主之心,然後即刻滿足；人若肯信靠此血,則前雖有至大之罪,亦可即免除,不再追究.血的工,如是之重要,我們不可不注重.
\newline
\newline
三,洒血於壇之四圍,表普世人皆須賴耶穌之流血,然後可以得近主.此祭壇是主曾區別之為神人相接之處.我們心中天良的虧欠,己經洒去,此洒去,就是祭壇的工作；我們可以存著誠心與充足的信心,到神面前也.(來十章廿二節)
\newline
\newline
所洗的是臟與腿,臟是內部,表意念；腿是外部,表行為.耶穌之意念行為,完全清潔,將自己無瑕無疵之身,獻給上帝,(來九章十三節)訓我們當賴主的靈與道,(水)潔淨內心與行為,而常居於成聖之地位.賴水與靈,而使人成聖之道,可看(約三章五節又十七章十七節,弗五章廿六,廿七節.)
\newline
\newline
燔祭之焚,與贖罪祭之焚,原文有不同之意義；因燔祭之焚,是馨香上升,令主喜悅,贖罪祭之焚,是因犯罪而焚,欲即毀滅之情形.由此火焚之工,我們得以下的教訓:
\newline
\newline
一,擺在柴上.(八節)柴是木,可表十字架,耶穌在十字架上獻了自己,此擺字,有按著次序的教訓；耶穌由馬槽至十字架,是皆按著受苦的秩序即所謂十字架的路.在腓立比書(二章五至八節)就是一條按序受苦路,至終死在十字架.我們獻己的路,亦有一定的次序,經過諸多爭戰,忍苦羞辱,擊碎的工夫,然後可進到陳列主前為馨香的獻己.
\newline
\newline
二,焚獻一切.一切焚燒在壇上,乃蒙主悅納之唯一條件,即所謂完全奉獻,屬魂屬身各部的,一切完全被主操權,完全順服,獻牛羊為燔祭者,以此乃最後之手續.我們獻己,能保守完全的順服,就是再無別的更好方法.
\newline
\newline
三,馨香蒙納,若依以上手續而獻,便為馨香；有一部分不被火焚,該部分即為惡臭.基督在上帝前為馨香的靈祭面蒙悅納.(彼前二章五節,弗五章二節)火表聖靈,聖靈,我們內外的一,切須放在壇上,主的面前,受主的靈膏沐即以可保存不臭；我們最怕獻不蒙納之祭.以賽亞時代,猶太人殷勤獻祭守節,但因為手中有血(即罪)被主厭棄,(賽一章十一至十八節)故當盡除罪惡,免礙聖火.
\newline
\newline
一,獻鳥之工,較牛羊更煩苦.獻牛羊的,多由其本人親為之,惟鳥則否,祭司須為之折頸,擠血,裂身,焚燒,又要格外小心.猶太人謂獻鳥是祭司至難之工,此訓我們為主服役之人,對於那些心貧弱無力之信徒服役,更須細心代勞,體貼忍耐.
\newline
\newline
二,對無力者之特別待遇.牛羊要盡焚一切,鳥之嗉子及翎毛不焚；此或可表人之心靈高下不同,心靈貧弱者,或屬靈知識淺薄者,主待之多以慈愛憐憫,與心靈強壯,屬靈知識高深者有不同.
\newline
\newline
三,蒙納則一.牛羊馨香,鳥亦馨香,不重物之大小,重在人之盡心.據猶太人云,若獻鳩則須老鳩,獻鴿則須雛鴿,即人在席上以為如此方是美味者,即用之以獻主；但若有牛羊,則須獻牛羊,總之必須盡力誠心而後已.
\newline
\newline
亞倫子孫,乃歷代相繼為祭司者,他們雖是人,但已為上帝選升為祭司之族,其地位甚高,雖須為己罪先獻祭,但己常常作代人獻祭之工矣.訓我們為信徒不可常常在犯罪求赦之地位,應心靈堅壯,前進至有為人代禱之能力；因為主召我們,為「有王位之祭司」,(彼前一章九節)當隨時舉起聖潔之手,隨處為人禱告,如亞倫子孫之日日代人獻祭焉.
\newline
\newline
一,當常不息.燔祭是馨香之祭,由愛心而獻,愈多愈好,愈得主悅；故獻者甚多,常有終夜至旦,壇上之火不熄.訓我們奉獻之工,是須常常有進步,此路愈行愈窄,又須常常保存此奉獻之火,在一切生活上不停息,步步踏著「無我惟主」之條例,則必至有豐盛之生命也.
\newline
\newline
二,添柴.欲壇上之火不滅,須每日添柴,(十二節)此添柴之工,是祭司之責任.此可訓我們為教會之領袖,擔負祭司職責之人,須有不已的祈禱,不使教會祈禱之工停止,然後可以得多人歸主,及得信徒多行獻己之工.今日的教會,祈禱之火不常顯著,所以祈禱之靈不大傾注於教會,奮興之王,故不大顯著,豈不惜乎?
\newline
\newline
三,上帝之火.我們應知會幕前壇上,耶和華曾由天降下火,焚燒其上的祭牲,(利九章廿四節)此火或壇上第一次之火,(參代下七章一至三節)所羅門之獻祭亦如是,上帝既用天火燒起我們之祭,我們應負責留心,不可使之熄滅,應將上帝賜我們之恩賜,再如火挑旺起來.(提後一章六節)可惜過了不多時候,即有非聖之火獻於香壇,致有火由主前出,焚死祭司二人.新約五旬節,有聖靈之火由天降,惜乎過不多時,有在奉獻事上犯罪,致有聖靈在主前擊死二人,故我們當慎重上帝之火也.
\newline
\newline
主理祭禮之祭司,要穿白細麻衣服,及白細麻褲.白細麻表聖潔公義,教訓服役主家之人,當保守此條件,代人獻祭之工,自己至要聖潔,且未穿之先,用水洗身.(利十六章四節)基督至聖的,所以能為萬人之祭司；願我們作新約祭司的人,常常保守自己之衣潔白,且頭上不缺香膏.(道九章八節)
\newline
\newline
所何以要傾灰,免阻火焚,此乃祭司每晨之工；此灰,是由獻祭而來的,此灰留於壇內無益,且有阻礙.教訓我們在各種靈德上,如犧牲,虔敬,節制,熱心,獻奉等,皆甚易於不知不覺中,留有各種渣滓,應常省察,作傾除之工,使焚獻之火常燃.論到預表,此灰或可指耶穌已死之體,傾於營外潔處,即指主葬於城外花園.耶穌負罪之體,要葬埋,然後得有復生之生命；我們之罪,亦須清除,然後得勝靈之大能也.
\newline
\newline


\newpage

\section{第1屆~港九培靈會~黃原素~第三講　素祭之研究}
\label{sec:1336}
\textbf{第三講　素祭之研究}
\newline
\newline
連結: \href{https://www.hkbibleconference.org/session-message/view/1336}{\texttt{https://www.hkbibleconference.org/session-message/view/1336}}
\newline
\newline
\hyperref[sec:1335]{< < < PREV SERMON < < <}
~
\hyperref[sec:toc]{[返目錄]}
~
\hyperref[sec:1337]{> > > NEXT SERMON > > >}
\newline
\newline
利未記 2:1-2:16
\newline
\begin{longtable}{cl}
\hline
\hline
章節 & 經文 (和合本修訂版)\\
\hline
2:1 & \begin{tabularx}{0.7\textwidth}{X} 「若有人獻素祭為供物給耶和華，就要獻細麵為供物，把油澆在上面，加上乳香， \end{tabularx} \\ \\ \relax
2:2 & \begin{tabularx}{0.7\textwidth}{X} 帶到亞倫子孫作祭司的那裡。祭司要從細麵中取出滿滿的一把，又取些油和所有的乳香，把這些作為紀念的燒在壇上，是獻給耶和華為馨香的火祭。 \end{tabularx} \\ \\ \relax
2:3 & \begin{tabularx}{0.7\textwidth}{X} 素祭所剩的要歸給亞倫和他的子孫；在獻給耶和華的火祭中，這是至聖的。 \end{tabularx} \\ \\ \relax
2:4 & \begin{tabularx}{0.7\textwidth}{X} 「若獻爐中烤的素祭為供物，要用調了油的無酵細麵餅，或抹了油的無酵薄餅。 \end{tabularx} \\ \\ \relax
2:5 & \begin{tabularx}{0.7\textwidth}{X} 若以鐵盤上的素祭為供物，就要用調了油的無酵細麵， \end{tabularx} \\ \\ \relax
2:6 & \begin{tabularx}{0.7\textwidth}{X} 分成小塊，澆上油；這是素祭。 \end{tabularx} \\ \\ \relax
2:7 & \begin{tabularx}{0.7\textwidth}{X} 若以煎鍋煎的素祭為供物，就要用油與細麵做成。 \end{tabularx} \\ \\ \relax
2:8 & \begin{tabularx}{0.7\textwidth}{X} 要把這樣做成的素祭帶到耶和華面前，拿給祭司，祭司要帶到壇前。 \end{tabularx} \\ \\ \relax
2:9 & \begin{tabularx}{0.7\textwidth}{X} 祭司要從素祭中取出作為紀念的燒在壇上，是獻給耶和華為馨香的火祭。 \end{tabularx} \\ \\ \relax
2:10 & \begin{tabularx}{0.7\textwidth}{X} 素祭所剩的要歸給亞倫和他的子孫；在獻給耶和華的火祭中，這是至聖的。 \end{tabularx} \\ \\ \relax
2:11 & \begin{tabularx}{0.7\textwidth}{X} 「凡獻給耶和華的素祭都不可以有酵，因為你們不可把任何的酵或蜜燒了，當作火祭獻給耶和華。 \end{tabularx} \\ \\ \relax
2:12 & \begin{tabularx}{0.7\textwidth}{X} 你們可以把這些獻給耶和華當作初熟的供物，但是不可獻在壇上作為馨香的祭。 \end{tabularx} \\ \\ \relax
2:13 & \begin{tabularx}{0.7\textwidth}{X} 凡獻為素祭的供物都要用鹽調和；在素祭中，不可缺少你與神立約的鹽。一切的供物都要加鹽獻上。 \end{tabularx} \\ \\ \relax
2:14 & \begin{tabularx}{0.7\textwidth}{X} 「你若獻初熟之物給耶和華為素祭，就要獻在火中烘過的新麥穗，就是磨碎的新穀物，當作初熟之物的素祭。 \end{tabularx} \\ \\ \relax
2:15 & \begin{tabularx}{0.7\textwidth}{X} 你要加上油和乳香；這是素祭。 \end{tabularx} \\ \\ \relax
2:16 & \begin{tabularx}{0.7\textwidth}{X} 祭司要把供物中作為紀念的，就是一些磨碎的新穀物和一些油，以及所有的乳香，都焚燒，是獻給耶和華的火祭。」 \end{tabularx} \\ \\
[1ex]
\hline
\hline
\end{longtable}
一,此祭非贖罪所用.
\newline
\newline
二,此祭無血.
\newline
\newline
三,此祭之預表,重在耶穌之為人.
\newline
\newline
耶穌稱為「人子」,重看此人字,耶穌是一個完人,燔祭表耶穌獻己而死,以十字架為標的；素祭則全不關耶穌之死及流血等事,乃專指耶穌之如何為人,即所謂耶穌在世之「人格」.夫認識耶穌之神性固極重要,認識耶穌之人格,亦極重要也.經云,道成肉身居於我們中間,「道」與「肉身」聯成一句即指耶穌之全完神性,及完全人性.「我們見其榮」,若耶穌不成為肉身,我們焉能見之.論耶穌在世之人格,真是無可指摘,全聖全義,聖經稱之為「聖者」,(徒三章十四節)「義者」,(徒七章五十二節)我們可稱之為「完人.」
\newline
\newline
四,此祭不可單獨獻,必須配他祭同獻.(利廿三章十二,十三,十八節,民廿八章七至十五節),以上經節,屢見有配獻二字,由此想到該隱之祭,所以不蒙悅納者,因他單獨獻上素祭；當日曾受過上帝以血為祭之道(創三章廿一節)而不從,亞伯以血而獻,(希伯來十一章四節)評其有信,即信血能免人罪,故蒙悅納.今日教會有所謂新學派者,他們一切言論,思想,信仰,總是在耶穌之人格上立場,放棄耶穌之神性信仰.他們一切事工,均無血的根基,去救贖之工甚遠；既無血；即不能有膏,聖靈此等教會,是謂之「惟獻素祭」之教會,定不能蒙主悅納也.然根基信仰派的人,不是放棄耶穌之美備人格不談,但不單獨專談,乃配其神能而談,二者相比,當言其神性重于人格也.
\newline
\newline
甲,細麵.(一章一,四,五,六,七節)
\newline
\newline
一,表主耶穌品性之謙卑,溫柔.夫麵已幼滑無比.且加上細,可見耶穌在世為人,細密不苟；如家庭中之孝順,臨終託母,祝福小孩,執遺屑十二筐等事,俱可耗之,此教訓我們為人應處事細密,因為不矜細行,終累大德,在小事上不苟,在大事方不苟也.
\newline
\newline
二,表至耶穌之品性完備不偏.麵粉粒粒皆細,無大小,聖經內之名人,各有欠缺,各有輕重；惟有耶穌之品性,方方無缺,面面皆圓,故謂世無完人,惟耶穌耳.此訓我們行事,當盡合乎主,行諸善事,凡事悅主之意.(西一章十節)
\newline
\newline
三,表耶穌磨難中之美備人格.由一粒麥在地中生出,至成為餅,其聞經過許多風霜折磨擊打之事,且經火,然後成餅.耶穌在生三十餘年,由誕生至死,其間所歷,亦經過極大火煉之試,(腓二章五至八節)至終死在十字架,其犧牲忍苦之人格,終不變更,卒成為人之生命餅.(約六章卅五節)此訓我們當忍受苦,難在試煉中顯出喜樂,(雅一章廿四節)使我們經過折磨後,仍能卓立.(弗六章十三節)
\newline
\newline
乙,油.(二章一,四,五,六,七,十五節)油表聖靈
\newline
\newline
一,傾于細麵之上.(一節六節)
\newline
\newline
普通素祭,用麵粉一伊法十分之一,(約一升七合一油,一欣四分之一,(約二斥半)以斤半油,傾於一升七合麵粉上,油之分量,極為豐富,表耶穌受聖靈豐盛的膏沐(徒十章卅八節)上主以聖靈膏拿撒勒人耶穌,在約但汀上受聖靈降臨其上,歸拿撒勒故鄉,入會堂讀經,首一句即云,主之靈臨我,主之膏膏我.(路四章十八節)此訓我們當受聖靈豐盛之充滿,以作其工,方有大效.
\newline
\newline
二,和於麵中.(四,五,七節,六章廿一節)
\newline
\newline
此表耶穌不只受膏於聖靈,且充于聖靈,甚至浸入聖靈中,與聖靈混合,聯合不分；故其一生治病,逐鬼,講道,俱有聖靈之工在內,如麵粉與油調勻,二者成為一體.此訓我們一生之生活,當與聖靈混合,在靈裡作工,在靈裡祈禱,在靈裡生活；如此,在主前必為馨香之祭,蒙主悅納.
\newline
\newline
丙,乳香.(二章一節,二節,十五節)
\newline
\newline
一,完全焚燒不留.所以用乳香加於麵粉或麵餅之上者,取其馨香也.此表耶穌之功勞,及其向外之工作,在主前實為馨香之祭；素祭之麵,只取一握,焚之於壇上,乳香,則完全焚燒,此表耶穌在世之工作,無論大小,其榮不歸己,完全皆歸父.(約五章四十一節又八章五十節)此訓我們勿偷竊主之榮耀,(詩一百十五篇一節)當求榮主,勿求榮己,偷竊主榮,是主僕最大的危險.慕翟先生講道後,有人讚他曰,「先生,你講得真好,」慕先生答曰,「是,魔鬼剛才對我說了,」所以魔鬼不怕我會講,只怕我以榮盡歸於主.我們講道有力,是因上帝之祝福,人之代禱,自己與神交通,此皆非有所可誇的.有人在他桌上書云,「不可思想你自己,」因思己好,則生驕,思己不好則灰心,我們作了工,則盡託於主可也.
\newline
\newline
二,蒙主記憶.(二節九節)各種素祭,均有此語,我們在工所得之榮耀雖盡歸主,但我們工作已得主之記憶,主非不義,必不忘我之工,只恐我們之工,不是馨香,不蒙悅納耳.
\newline
\newline
丁,鹽.(二章十三節,民十八章十九節)
\newline
\newline
鹽之效用有二:
\newline
\newline
一,保存.　鹽之效用,於人極有關係,而人反輕之；耶穌來世作救世保世之工,而人反輕視不接.不獨耶穌保世,信徒亦然.(太五章十三節)敵基督者,久欲大肆惡毒,毀懷此世,奈有教會在,阻他不能大大發動,須教會被主接去空中,敵基督者然後可大行其惡也.(帖後二章六至八節)如是,信徒一日在世界,世界一日末能大壞,豈不是信徒能保世乎﹖
\newline
\newline
二,調和.　即使人有味之謂,主耶穌之道,最能使人為人有味,有希望,有安慰,不獨耶穌為然,信徒亦是,(西四章六節)使人調和,教人進天國,惟當至要勿失其鹹,因鹽失了鹹,絕無用處,惟被棄於外,任人踐踏耳.
\newline
\newline
麵油酒
\newline
\newline
一 羔伊法十分之一一欣四分之一一欣四分之一(四,五節)
\newline
\newline
一 牡綿羊伊法十分之二一欣三分之一一欣三分之一(六,七節)
\newline
\newline
一 牡犢伊法十分之三一欣四分之二一欣四分之二(八至十節)
\newline
\newline
一,奠禮.(民廿八章七節)
\newline
\newline
奠祭與素祭,同為他祭牲配之而獻,所奠者為酒,即當日由壓酒處醡出之葡萄汁,其色似人類之血,酒不焚於壇,須攜入聖所奠於主前,(民廿八章七節,出卅章九節)此可表耶穌在十架,為人類傾出其生命,如酒由壓酒處而出,(示十九章十五節)配獻祭牲之物,即麵,油,酒,三種,麵(餅)表耶穌之體,油表聖靈,酒表耶穌之血也.
\newline
\newline
二,素祭分量,隨祭牲大小而定.
\newline
\newline
燔祭之預表在獻己,素祭之物如麵,油,酒等,其預表重在在世之行為,由此可教訓我們獻己者,受聖靈之感動深,聖經之智識多,獻己之工作高明,則向世之行為,當然與之相稱,即品格愈須高潔,然後可以配得起.
\newline
\newline
甲,酵.
\newline
\newline
一,表罪.酵表罪者,散見於舊新約,其最著者,是除酵節,(出十二章十五節)以色列人在節內七日之第一日,要盡除一切酵,一點不能留,一點不能食,何以若是嚴肅乎；蓋酵表罪也,教會既得主流血救贖之恩,(逾越節)應該盡除罪惡,使成一個新教會,因為一點的酵,能發全體,(哥前五章六至八節,太十三章卅三節)故凡祭物之完全預表主基督者,即禁用之,以證主之完全聖義無罪.
\newline
\newline
二,表異端.主云,謹防法利賽及撒吐該之酵,此即表他們的異端,但凡與純正福音之道不符的各種教理,皆算是酵,能混亂正道,並敗壞正道,今日教會中學一種最可惡的酵,已敗壞了不少福音純正之事業,此酵是何,即所謂新派神學是也.
\newline
\newline
以上兩則,教訓我們當謹防一切酵毒,在奉事上不可存有些須之罪惡,與異端,務要盡除,俾成新團而後已.
\newline
\newline
乙,蜜.
\newline
\newline
一,蜜雖甜而易變酸.耶穌之言行,是永遠聖義馨香的,任遭何種擊打,皆可永久存留,與世名人之言行不同；世上名人之言行價值,要受時代支配,主之言行,是精練有力,永有價值.此訓我們在世之言行,當求真材實料,永有效力.
\newline
\newline
二,耶穌在世之生活,不是甜的生活.　其生活是鹹,是苦,無日不依著十字架之道路走,無時不抱著犧牲之觀念,讀以賽亞五十三章,可見耶穌之生活,是無半點甜味.此訓我們為徒者,在世不當求甜蜜之生活,即安逸之生活,為主之名,常勞苦以效我主焉.
\newline
\newline
甲,用生麵.(一節)用麵粉澆上油,加上乳香,不用過火炊,即可帶至祭司代獻.
\newline
\newline
乙,用爐烤.(四節)
\newline
\newline
丙,用鏊炊.(五節)
\newline
\newline
丁,用盤煎.(七節)
\newline
\newline
戊,用烘穗.(十四,十五節)
\newline
\newline
以上五種獻法,不外兩個意,就是經過火與磨,及按獻者之方便.
\newline
\newline
一,經過火磨　火與磨,皆苦極之工作,此表耶穌在世為人,歷盡此種火之苦難.以賽亞五十三章全載耶穌在世受苦之事,此中有云,「他多受痛苦,常經憂患」,又云,「耶和華定意將他壓傷,使他受痛苦」,有云耶穌在十字架上六小時,由上午九時至下午三時,試思耶穌在此六小時,何等折磨.我們知耶穌在世之美備人格,不是偶然僥倖的,乃經過火磨生活而來.保羅一,曉得和他一同受苦,效法他的死.耶穌火磨之工已完,我們今日在世亦有此火磨之工,以補其未盡受之苦.(西一章廿四節)
\newline
\newline
二,所以有如此多種獻法者,乃因按各人家中所有之器物.此訓我們對於善行一事,每個信徒應有之,不能推諉.最單簡的獻法,是用生麵,主亦納之；所以我們為信徒,每人皆受過主之靈恩,每人應有善行,所受者,各雖有不同,然但各可按所受而行,便可得主納也.
\newline
\newline
甲,祭司可食.(即亞倫子孫)(二章三節,十節,六章十六節)
\newline
\newline
祭司以上帝為業,靠聖所得食,此種業,算是至可靠的；上帝永不死,祭司之業也永不衰亡.照聖經所載,素祭之麵油中,僅取出一握,及油少許而焚於主,其餘盡歸祭司,及其子孫,全國人每日獻上素祭的品物必不少,而祭司所得者必極其豐足,此可訓代獻者之權利,何等豐厚.我們代人祈禱中代求之工,自己所得之靈益實多,愈多愈益,如是,我們奉事主者,在世之生活上,似已犧牲,其實在靈福上,有豐盛之得也.
\newline
\newline
乙,亞倫子孫之男子可食.(六章十八節)
\newline
\newline
平安祭,男女可食,此指主前之喜樂,此素祭所餘,要男丁方可食,在舊約的男子地位,確然有別；在新約對主之靈益言,則男女皆一也.男子性剛強,女子性弱,古今一例,故可表凡在主內心靈強壯者,可有份享主之素祭,即與主同表善行於世,其心靈軟弱者,則不能也.此訓我們在靈德工程上,當作一大丈夫,勿向撒旦罪惡世界投降也.」(提後一章七節)
\newline
\newline
丙,祭司為己所獻之素祭,須完全焚燒.(六章廿二,廿三節)
\newline
\newline
祭司所獻者,用麵伊法十分之一,(民十五章四節)即一俄梅珥.昔日以色列民在野取嗎哪,每人每日一俄梅珥,此乃一日之糧；祭司所獻一俄梅珥而全焚之,即訓為奉事主之人,應盡一日之精神能力,以作主之工也.
\newline
\newline


\newpage

\section{第1屆~港九培靈會~黃原素~第四講　平安祭之研究}
\label{sec:1337}
\textbf{第四講　平安祭之研究}
\newline
\newline
連結: \href{https://www.hkbibleconference.org/session-message/view/1337}{\texttt{https://www.hkbibleconference.org/session-message/view/1337}}
\newline
\newline
\hyperref[sec:1336]{< < < PREV SERMON < < <}
~
\hyperref[sec:toc]{[返目錄]}
~
\hyperref[sec:1333]{> > > NEXT SERMON > > >}
\newline
\newline
利未記 3:1-3:17
\newline
\begin{longtable}{cl}
\hline
\hline
章節 & 經文 (和合本修訂版)\\
\hline
3:1 & \begin{tabularx}{0.7\textwidth}{X} 「人獻平安祭為供物，若是從牛群中獻，無論是公的母的，要用沒有殘疾的，獻在耶和華面前。 \end{tabularx} \\ \\ \relax
3:2 & \begin{tabularx}{0.7\textwidth}{X} 他要按手在供物的頭上，在會幕的門口宰了牠。亞倫子孫作祭司的，要把血灑在壇的周圍。 \end{tabularx} \\ \\ \relax
3:3 & \begin{tabularx}{0.7\textwidth}{X} 從平安祭中，他要把火祭獻給耶和華，就是包著內臟的脂肪和內臟上所有的脂肪， \end{tabularx} \\ \\ \relax
3:4 & \begin{tabularx}{0.7\textwidth}{X} 兩個腎和腎上的脂肪，即靠近腎兩旁的脂肪，以及肝上的網油，連同腎一起取下。 \end{tabularx} \\ \\ \relax
3:5 & \begin{tabularx}{0.7\textwidth}{X} 亞倫的子孫要把這些擺在燒著火的柴上，燒在壇的燔祭上，是獻給耶和華為馨香的火祭。 \end{tabularx} \\ \\ \relax
3:6 & \begin{tabularx}{0.7\textwidth}{X} 「人向耶和華獻平安祭為供物，若是從羊群中獻，無論是公的母的，要用沒有殘疾的。 \end{tabularx} \\ \\ \relax
3:7 & \begin{tabularx}{0.7\textwidth}{X} 若他獻一隻綿羊為供物，就要把牠獻在耶和華面前。 \end{tabularx} \\ \\ \relax
3:8 & \begin{tabularx}{0.7\textwidth}{X} 要按手在供物的頭上，在會幕前宰了牠。亞倫的子孫要把血灑在壇的周圍。 \end{tabularx} \\ \\ \relax
3:9 & \begin{tabularx}{0.7\textwidth}{X} 從平安祭中，他要取脂肪當作火祭獻給耶和華，就是靠近脊骨處取下的整條肥尾巴，包著內臟的脂肪和內臟上所有的脂肪， \end{tabularx} \\ \\ \relax
3:10 & \begin{tabularx}{0.7\textwidth}{X} 兩個腎和腎上的脂肪，即靠近腎兩旁的脂肪，以及肝上的網油，連同腎一起取下。 \end{tabularx} \\ \\ \relax
3:11 & \begin{tabularx}{0.7\textwidth}{X} 祭司要把這些燒在壇上，是獻給耶和華為食物的火祭。 \end{tabularx} \\ \\ \relax
3:12 & \begin{tabularx}{0.7\textwidth}{X} 「人的供物若是山羊，就要把牠獻在耶和華面前。 \end{tabularx} \\ \\ \relax
3:13 & \begin{tabularx}{0.7\textwidth}{X} 要按手在牠的頭上，在會幕前宰了牠。亞倫的子孫要把血灑在壇的周圍， \end{tabularx} \\ \\ \relax
3:14 & \begin{tabularx}{0.7\textwidth}{X} 又要從供物中把火祭獻給耶和華，就是包著內臟的脂肪和內臟上所有的脂肪， \end{tabularx} \\ \\ \relax
3:15 & \begin{tabularx}{0.7\textwidth}{X} 兩個腎和腎上的脂肪，即靠近腎兩旁的脂肪，以及肝上的網油，連同腎一起取下。 \end{tabularx} \\ \\ \relax
3:16 & \begin{tabularx}{0.7\textwidth}{X} 祭司要把這些燒在壇上，作為馨香火祭的食物；所有的脂肪都是耶和華的。 \end{tabularx} \\ \\ \relax
3:17 & \begin{tabularx}{0.7\textwidth}{X} 在你們一切的住處，脂肪和血都不可吃，這要成為你們世世代代永遠的定例。」 \end{tabularx} \\ \\
[1ex]
\hline
\hline
\end{longtable}
甲,此祭之預表,是表耶穌之死,使神人兩方面復和.
\newline
\newline
夫燔祭是重指耶穌向神獻己,素祭是重指耶穌在人間為完全人,作人類之代表.平安祭之地位,是在中間.(見圖)我們賴其子之死,得與神復和.(羅五章十節)故平安祭最重之意,是人因耶穌而得與神和好,除去了種種隔膜,內心恆享平安,(賽九章六,七節,弗二章十四節)而又得為義,並獲與神心交.(約一書一章三節)
\newline
\newline
乙,諸祭中,此祭之牲歸主者至少至美.(三章三,四,九節)
\newline
\newline
用牛,(一節用羊羔),(七節用山羊),(十二節此與燔祭大致相同,惟獨無為,何以?大概有三個緣故:
\newline
\newline
一,因平安祭是神人兩方平安喜樂之祭,須神人共享祭牲,鳥太細小,人難分享.
\newline
\newline
二,因鳥太細小,甚難區別何部分歸主,何部分歸人,因鳥之脂油無多,全鳥不知何部分算最佳美.
\newline
\newline
三,因平安祭是為感謝,甘心,許願而獻,即因此三事而獻,則至少必須籌備一隻山羊,方可表其感謝甘心之誠意.
\newline
\newline
感謝,甘心,償願,此三個原因,均屬於自由不強,茲論其四個表訓如下:
\newline
\newline
一,為感謝而獻者,許用有酵之餅.(七章十三節)
\newline
\newline
酵,既表罪,與異端,何以許用酵為祭?因此祭是關於神人兩方,人之方面是有罪性,此祭亦有洒血之工,(十四節)可知有赦罪之工在內,故可用酵.(在利廿三章十七節)五旬節獻祭之二大餅,亦用酵,因此二餅,一指以色列,一指猶大,此外祭物中,未見何處許用酵也.」但此酵餅,非謂主視為馨香而享受之,不過陳列於其前,而供其鑑察耳.又此酵餅,僅是一個,此可訓我們之罪在主前,並不馨香,且常陳於主前,受主之審判；惟願我們之罪,僅一而不可再,如酵餅之僅有一個而已.我們之罪,當陳於主前,求主赦免,主必以仁慈寬容赦宥之.
\newline
\newline
二,此祭只獻極少之分於主,餘則盡歸於人食之.
\newline
\newline
受獻者,代獻者,獻者三方面均可同享,歡心樂意而食.我若能以至優至美之分歸於主,即我們之全心全愛獻於主,則必在主之內有心滿意足的屬靈喜樂與利益；主所賜回我之恩澤,豐盛有餘,使我們心滿意足,讚美主恩也.
\newline
\newline
又可表我們人類以耶穌為生命餅.(約六章五十三至五十八節)
\newline
\newline
三,感謝而獻者,其肉必食於是日；但為償願或甘心而獻者,其肉可留至第二日.(七章十五,十六節)
\newline
\newline
何以,因可表感謝出於熱誠,熱誠之心,當即日盡力即行之,恐怕遲至明日即消失了；又訓信徒感謝之心,當日日新鮮,有新鮮之讚美.由甘心與許願而出者,出於愛心,愛主之心,不可一日即了,應有依依不捨之情.但祭肉不出三日,出三日即不食,而用火熱之；(十七節)此表耶穌之體不能過於三日夜在地中,其身不能臭壞.(徒二章廿七節)
\newline
\newline
四,甘心而獻者,許用不完全之牲.(廿二章廿三節)
\newline
\newline
因感謝與償願而獻者,必用完全無疵之牲,因感謝出於熱誠,償願,既應許必償,已入於律法範圍.但甘心是本乎愛心,在此愛心上,主特許人有一點之自由,但當小心,主人是規定人因甘心獻者必須獻不完全牲,不過在此特用此一點之自由,以試驗人之愛心程度耳.慎哉,此實乃一大試驗之關頭也；若人家中有完全之牲而不獻,而特獻以不完全者,是乃褻瀆主,不得主之悅納.(一章十四節)
\newline
\newline
甲,與燔祭之相同點.無疵之牲,牽至幕前,按手其首,宰於門前,洒血於壇四圍.
\newline
\newline
乙,與燔祭之不同點.不論牡牝,不剝牲皮,不洗臟腿,不盡焚,不用鳥.此五件不同點之原因,見上文第二第三論可詳知.
\newline
\newline
丙,祭牲歸主之分如何?
\newline
\newline
一,脂油,(三章三,四,九節)「蓋臟之脂油」,「臟上所有之脂油」,「靠腰兩旁之脂油」,「肝上的網子」,「整肥的尾巴」,(指羊尾該地之羊尾壹十餘斤)此皆時屬乎脂油.牲畜之肥美,在乎脂,故脂,可謂至甘至美.
\newline
\newline
二,腰子,(三,四節)腰,或是當日牲畜肉食中算是美食未定,或指到情愛而言,亦未定,因腰是屬乎內部之物,如是或可謂之至情.
\newline
\newline
以上兩種,算是牲體中最優之分,主特要取此；主所求於人者,非次等的,乃上等的.如所云,初熟,首生,歸主；耶穌時,婦人以至真且至貴之香膏膏之；(約十二章一至七節)主所求於亞伯拉罕者,是他最寶過於生命的以撒,(創廿二章二節)云,「你的兒子,你的獨生子,你所愛的以撒」.此訓我們獻於主,當以最好的.燔祭重在盡,平安祭在最好.
\newline
\newline
一,要自潔,(七章十九,廿,廿一節)
\newline
\newline
因此乃聖潔之牲,食之者亦當潔；否則剪滅於民中,所謂剪滅,大抵是一,出會,二,處死.不潔而食,有此大結果,真是可怕；我們食主之聖餐,亦當有此感想.(哥前十一章廿七至廿九節)不潔而食,即藐視主之聖禮,干犯主之聖體矣.人不潔,必不能享主之筵席.
\newline
\newline
二,要喜樂,(申十二章十八節)
\newline
\newline
平安祭,是喜樂之祭,摩西未立此禮之前,古代人已有平安祭之筵席.(創卅一章五十四節,出十八章十二節)我們若有罪,未與神和好,心中失了平安,則斷不能在主前有屬靈滿足之喜樂；喜樂一失,諸般惡念舊性復來,此為靈程上一個極大之危機也.
\newline
\newline
三,要在主前,(申十二章四至七節十七,十八節)
\newline
\newline
何以要在主前?(若離太遠可以在本城申十二章廿一節)因為要記念此筵席之主是耶和華.喜樂由耶和華所賜,我之好處,不在主以外；(詩十六篇二節)我們欲得真正之喜樂,只在主前密室得之.
\newline
\newline
四,禁食脂與血.(三章十七節)
\newline
\newline
血,是生命；脂,是豐盛之生命.凡畜牲瘦弱,必先消失其脂.血是其生命之根本,如是血可表救贖,脂可表犧牲；大有影響於耶穌救世之工,且使人每飯不忘主之救恩,又教人當以己至優之點歸主.
\newline
\newline
一,搖胸(卅節)
\newline
\newline
「搖一搖」,祭司將牲之胸,在主前搖過一搖,即由主前接受之,歸回自己；下文「舉之」,恐亦大致是也.胸,表愛也,此胸搖後,全歸祭司子孫所食；訓我們的愛力,是由我們先以己之愛獻於主,然後由主前接受之也.祭司得胸,不獨自享,亦與眾子孫同享；我們主之愛,亦當同人分享主之愛也.
\newline
\newline
二,舉右腿(卅節)
\newline
\newline
右腿,即肩,表負擔之能力；亦訓我們當先將己之力獻於主,然後由主前取回回.大能力如摩西之杖,在主前擲之,後執回,是一條神杖,大有拯救之能.(出三章四節)
\newline
\newline
以上所言,胸,(表愛)肩,(表能力)俱為祭在主所得之分.由此我們回想大祭司之以弗得之胸牌及肩牌,(出廿八章七至廿一節)兩者俱有十二支派之名,祭司入覲主時即服之；可見作祭司的至大至要之責任,就是胸懷其民,肩負其民.我們為傳道者,不亦如是乎?但此兩職之能力,往何處得之?曰,惟有主前而已﹗
\newline
\newline


\newpage

\section{第1屆~港九培靈會~黃原素~第五講　贖罪祭之研究}
\label{sec:1333}
\textbf{第五講　贖罪祭之研究}
\newline
\newline
連結: \href{https://www.hkbibleconference.org/session-message/view/1333}{\texttt{https://www.hkbibleconference.org/session-message/view/1333}}
\newline
\newline
\hyperref[sec:1337]{< < < PREV SERMON < < <}
~
\hyperref[sec:toc]{[返目錄]}
~
\hyperref[sec:1332]{> > > NEXT SERMON > > >}
\newline
\newline
利未記 4:1-4:35
\newline
\begin{longtable}{cl}
\hline
\hline
章節 & 經文 (和合本修訂版)\\
\hline
4:1 & \begin{tabularx}{0.7\textwidth}{X} 耶和華吩咐摩西說： \end{tabularx} \\ \\ \relax
4:2 & \begin{tabularx}{0.7\textwidth}{X} 「你要吩咐以色列人說：若有人無意中犯罪，在任何事上犯了一條耶和華所吩咐的禁令， \end{tabularx} \\ \\ \relax
4:3 & \begin{tabularx}{0.7\textwidth}{X} 或是受膏的祭司犯了罪，使百姓陷在罪裡，他就當為自己所犯的罪，把沒有殘疾的公牛犢獻給耶和華為贖罪祭。 \end{tabularx} \\ \\ \relax
4:4 & \begin{tabularx}{0.7\textwidth}{X} 他要把公牛牽到會幕的門口，在耶和華面前按手在牛的頭上，把牛宰於耶和華面前。 \end{tabularx} \\ \\ \relax
4:5 & \begin{tabularx}{0.7\textwidth}{X} 受膏的祭司要取些公牛的血，帶到會幕那裡。 \end{tabularx} \\ \\ \relax
4:6 & \begin{tabularx}{0.7\textwidth}{X} 祭司要把手指蘸在血中，在耶和華面前對著聖所的幔子彈血七次， \end{tabularx} \\ \\ \relax
4:7 & \begin{tabularx}{0.7\textwidth}{X} 又要把一些血抹在會幕內，耶和華面前香壇的四個翹角上，再把公牛其餘的血全倒在會幕門口燔祭壇的底座上； \end{tabularx} \\ \\ \relax
4:8 & \begin{tabularx}{0.7\textwidth}{X} 又要取出這頭贖罪祭公牛所有的脂肪，就是包著內臟的脂肪和內臟上所有的脂肪， \end{tabularx} \\ \\ \relax
4:9 & \begin{tabularx}{0.7\textwidth}{X} 兩個腎和腎上的脂肪，即靠近腎兩旁的脂肪，以及肝上的網油，連同腎一起取下， \end{tabularx} \\ \\ \relax
4:10 & \begin{tabularx}{0.7\textwidth}{X} 正如從平安祭的牛身上所取的，祭司要把這些燒在燔祭壇上。 \end{tabularx} \\ \\ \relax
4:11 & \begin{tabularx}{0.7\textwidth}{X} 但公牛的皮和所有的肉，以及頭、腿、內臟、糞， \end{tabularx} \\ \\ \relax
4:12 & \begin{tabularx}{0.7\textwidth}{X} 就是全公牛，要搬到營外清潔的地方倒灰之處，放在柴上用火焚燒。 \end{tabularx} \\ \\ \relax
4:13 & \begin{tabularx}{0.7\textwidth}{X} 「以色列全會眾若犯了錯，在任何事上犯了一條耶和華所吩咐的禁令，而有了罪，會眾看不出這隱藏的事； \end{tabularx} \\ \\ \relax
4:14 & \begin{tabularx}{0.7\textwidth}{X} 他們一知道犯了罪，就要獻一頭公牛犢為贖罪祭，牽牠到會幕前。 \end{tabularx} \\ \\ \relax
4:15 & \begin{tabularx}{0.7\textwidth}{X} 會眾的長老要在耶和華面前按手在公牛的頭上，把牛宰於耶和華面前。 \end{tabularx} \\ \\ \relax
4:16 & \begin{tabularx}{0.7\textwidth}{X} 受膏的祭司要取些公牛的血，帶到會幕那裡。 \end{tabularx} \\ \\ \relax
4:17 & \begin{tabularx}{0.7\textwidth}{X} 祭司要用手指蘸一些血，在耶和華面前對著幔子彈七次， \end{tabularx} \\ \\ \relax
4:18 & \begin{tabularx}{0.7\textwidth}{X} 又要把一些血抹在會幕內，耶和華面前壇的四個翹角上，再把其餘的血全倒在會幕門口燔祭壇的底座上。 \end{tabularx} \\ \\ \relax
4:19 & \begin{tabularx}{0.7\textwidth}{X} 他要取出公牛所有的脂肪，燒在壇上。 \end{tabularx} \\ \\ \relax
4:20 & \begin{tabularx}{0.7\textwidth}{X} 他要處理這牛，正如處理那頭贖罪祭的公牛一樣，他要如此去做。祭司要為他們贖罪，他們就蒙赦免。 \end{tabularx} \\ \\ \relax
4:21 & \begin{tabularx}{0.7\textwidth}{X} 他要把牛搬到營外燒了，像燒前一頭公牛一樣；這是會眾的贖罪祭。 \end{tabularx} \\ \\ \relax
4:22 & \begin{tabularx}{0.7\textwidth}{X} 「官長若犯罪，在任何事上無意中犯了一條耶和華－他的神所吩咐的禁令，而有了罪， \end{tabularx} \\ \\ \relax
4:23 & \begin{tabularx}{0.7\textwidth}{X} 他一知道自己犯了罪，就要牽一隻沒有殘疾的公山羊為供物。 \end{tabularx} \\ \\ \relax
4:24 & \begin{tabularx}{0.7\textwidth}{X} 他要按手在羊的頭上，在耶和華面前宰燔祭牲的地方把牠宰了；這是贖罪祭。 \end{tabularx} \\ \\ \relax
4:25 & \begin{tabularx}{0.7\textwidth}{X} 祭司要用手指蘸一些贖罪祭牲的血，抹在燔祭壇的四個翹角上，再把其餘的血倒在燔祭壇的底座上。 \end{tabularx} \\ \\ \relax
4:26 & \begin{tabularx}{0.7\textwidth}{X} 祭牲所有的脂肪都要燒在壇上，正如平安祭的脂肪一樣。祭司要為他的罪贖了他，他就蒙赦免。 \end{tabularx} \\ \\ \relax
4:27 & \begin{tabularx}{0.7\textwidth}{X} 「這地的百姓若有人無意中犯罪，在任何事上犯了一條耶和華所吩咐的禁令，而有了罪， \end{tabularx} \\ \\ \relax
4:28 & \begin{tabularx}{0.7\textwidth}{X} 他一知道自己犯了罪，就要為所犯的罪牽一隻沒有殘疾的母山羊為供物。 \end{tabularx} \\ \\ \relax
4:29 & \begin{tabularx}{0.7\textwidth}{X} 他要按手在贖罪祭牲的頭上，在燔祭牲的地方把牠宰了。 \end{tabularx} \\ \\ \relax
4:30 & \begin{tabularx}{0.7\textwidth}{X} 祭司要用手指蘸一些祭牲的血，抹在燔祭壇的四個翹角上，再把其餘的血全倒在壇的底座上； \end{tabularx} \\ \\ \relax
4:31 & \begin{tabularx}{0.7\textwidth}{X} 又要把祭牲所有的脂肪都取下，正如取平安祭牲的脂肪一樣。祭司要把脂肪燒在壇上，在耶和華面前作為馨香的祭。祭司要為他贖罪，他就蒙赦免。 \end{tabularx} \\ \\ \relax
4:32 & \begin{tabularx}{0.7\textwidth}{X} 「人若牽一隻綿羊為贖罪祭作供物，就要牽一隻沒有殘疾的母羊。 \end{tabularx} \\ \\ \relax
4:33 & \begin{tabularx}{0.7\textwidth}{X} 他要按手在贖罪祭牲的頭上，在宰燔祭牲的地方宰了牠，作為贖罪祭。 \end{tabularx} \\ \\ \relax
4:34 & \begin{tabularx}{0.7\textwidth}{X} 祭司要用手指蘸一些贖罪祭牲的血，抹在燔祭壇的四個翹角上，再把其餘的血全倒在壇的底座上； \end{tabularx} \\ \\ \relax
4:35 & \begin{tabularx}{0.7\textwidth}{X} 又要把祭牲所有的脂肪都取下，正如取平安祭的羊的脂肪一樣。祭司要按獻給耶和華火祭的條例，把脂肪燒在壇上。祭司要為他所犯的罪贖了他，他就蒙赦免。」 \end{tabularx} \\ \\
[1ex]
\hline
\hline
\end{longtable}
一,此祭預表,重在耶穌挽贖人類之本罪.
\newline
\newline
二,此祭不是馨香祭,是免罪祭.
\newline
\newline
利未記一章一節有云,「主曉諭摩西說」,所說,即一二三章之燔祭,素祭,平安祭,之禮；至四章一節又再云,「主曉諭摩西說」,所說,即四五章之贖罪,贖愆,二之祭禮.由此而觀,可見馨香祭與罪祭,天然分兩大段也.
\newline
\newline
三,此祭不由隨意獻,是不可免而獻.獻之,不是令上帝享受馨香；乃是令上帝掩面不願觀.如銅蛇掛在竿上,是極不可看之事,耶穌掛在十字架,更是也.
\newline
\newline
四,獻者分等.(四章三節十三節廿二節廿七節)
\newline
\newline
五,極注重洒血之工.(四章五至七節,十六至十八節,廿五節三十節)
\newline
\newline
在當日是為贖人所犯之大罪而用,但預表就是耶穌之代死,贖我們之本罪.請注意四章二節之誤犯了一件罪之誤字,若云誤,則不能謂之本罪；豈知以色列人有摩西五經數百條之法律,口誦心思,若犯了律法所不許之罪,何得謂誤之?但聖經云誤,此可見罪之能力,有能令人生於罪中,而尚以為所犯之罪,是乃無心之誤.聖經證明人雖知神,不以神尊之,故上帝聽其謬妄,以致盈其罪惡之量,而聽候審判.嗚呼!世風日下,罪惡日盈,教育道德之力,斷不能挽斯世之邪惡；故上帝特遣其子降世,用血之秘能,然後堪以挽救人之罪性與惡慾也.
\newline
\newline
祭司(二至十二節), 會罪(十三至廿一節), 官長(廿二至廿六節), 平民(廿七至卅五節).
\newline
\newline
一,受膏祭司犯罪所獻之贖罪祭.(二至十二節)
\newline
\newline
二,會眾全體犯罪所獻之贖罪祭.(十三至二十一節)
\newline
\newline
三,官長犯罪所獻之贖罪祭.(廿二至廿六節)
\newline
\newline
四,平民犯罪所獻之贖罪祭.(廿七至卅五節)
\newline
\newline
甲,如何置備祭牲?(六章廿四至卅節)
\newline
\newline
贖罪祭之置備,最注意者,此牲是至聖的,(廿五節)別祭多注意無疵,無瑕,完全無殘疾.但此祭注意至聖,因表耶穌之身；若非至聖,則不能為天下人類贖罪.因至聖之故,則有以下數條例:
\newline
\newline
1. 要在主前宰燔祭之處,宰贖罪祭牲.(廿五節)
\newline
\newline
2. 可食之牲肉(即羊,官長與平民所獻者)要在聖處食.(廿六節)
\newline
\newline
3. 凡摸牲肉,亦必成聖.(廿七節)
\newline
\newline
4. 牲血偶染於衣,要在聖處洗之.(廿七節)
\newline
\newline
5. 烹牲肉之瓦器,打碎；銅器,要磨擦.(廿八節)
\newline
\newline
6. 祭司中男丁可食.(廿九節)
\newline
\newline
7. 其血攜入聖所的,(即牡犢,大祭司與會眾所獻者)其肉不可食,必用火焚之.(卅節)
\newline
\newline
乙,如何殺祭牲?
\newline
\newline
1. 自認滿於罪惡.
\newline
\newline
2. 認此牲是無疵無辜.
\newline
\newline
3. 認此牲之苦,是白為我而受.
\newline
\newline
4. 認此牲之死,是為我死.
\newline
\newline
5. 信我之罪,已為此牲擔負.
\newline
\newline
6. 信因此牲之代死,主必赦我罪.
\newline
\newline
四,血禮.
\newline
\newline
1. 彈於主前.(六章六節十七節)
\newline
\newline
2. 塗於香壇之四角.(七節十八節)
\newline
\newline
3. 塗於祭壇之四角.(廿五節三十節)
\newline
\newline
4. 餘血傾於祭壇之基.(七節十八節廿五節卅節)
\newline
\newline
五,焚脂.(及腰子)(八九十九,廿六,卅一節)
\newline
\newline
贖罪方面,重在流血,獻己方面,重在焚燒,此贖罪祭仍有馨香二字,(卅一節)此指耶穌獻己於主,行主之旨,為罪嘗死味,受諸般痛苦,為順服之祭；故耶穌之於上帝,誠為馨香親愛,而蒙納,所以脂腰盡焚.
\newline
\newline
六,焚牲之全體於營外.(四章十二節廿一節)
\newline
\newline
焚脂之焚,是馨香之焚,牲皮肉等之焚；(十二節廿一節)是可惡之焚,表主耶穌之身,滿了人類之罪,在主前視為可憎,所以在城外埋之,以應驗此贖罪牲之被焚於營外也.(太廿一章卅九節,約十九章十七,十八節)於我們心靈之教訓,是我主在邑外受苦,我們當出邑外就他,忍受他所受之凌辱；所以但凡從耶穌,皆要以此為標準,即須我們去就他,不是他來就我們,我們要效他的生活,不是要他來效我們的生活也.(來十三章十一至十三節)
\newline
\newline


\newpage

\section{第1屆~港九培靈會~黃原素~第六講　贖愆祭之研究}
\label{sec:1332}
\textbf{第六講　贖愆祭之研究}
\newline
\newline
連結: \href{https://www.hkbibleconference.org/session-message/view/1332}{\texttt{https://www.hkbibleconference.org/session-message/view/1332}}
\newline
\newline
\hyperref[sec:1333]{< < < PREV SERMON < < <}
~
\hyperref[sec:toc]{[返目錄]}
~
\hyperref[sec:1534]{> > > NEXT SERMON > > >}
\newline
\newline
利未記 5:1-5:19
\newline
\begin{longtable}{cl}
\hline
\hline
章節 & 經文 (和合本修訂版)\\
\hline
5:1 & \begin{tabularx}{0.7\textwidth}{X} 「若有人犯了罪，就是聽見了誓言，他本來可以作證，卻不把所看見、所知道的說出來，必須擔當他的罪孽。 \end{tabularx} \\ \\ \relax
5:2 & \begin{tabularx}{0.7\textwidth}{X} 若有人摸了任何不潔之物，無論是野獸的不潔屍體，家畜的不潔屍體，或是群聚動物的不潔屍體，他雖不察覺，也是不潔淨，就有罪了。 \end{tabularx} \\ \\ \relax
5:3 & \begin{tabularx}{0.7\textwidth}{X} 或是他摸了人的不潔之物，就是任何使人成為不潔的不潔之物，他雖不察覺，但一知道，就有罪了。 \end{tabularx} \\ \\ \relax
5:4 & \begin{tabularx}{0.7\textwidth}{X} 若有人隨口發誓，或出於惡意，或出於善意，這人無論在甚麼事上隨意發誓，雖不察覺，但一知道，就在這其中的一件事上有罪了。 \end{tabularx} \\ \\ \relax
5:5 & \begin{tabularx}{0.7\textwidth}{X} 當他在這其中的一件事上有罪的時候，就要承認所犯的罪， \end{tabularx} \\ \\ \relax
5:6 & \begin{tabularx}{0.7\textwidth}{X} 並要為所犯的罪，把他的贖愆祭牲，就是羊群中的一隻母綿羊或母山羊，獻給耶和華為贖罪祭，祭司要為他的罪贖了他。 \end{tabularx} \\ \\ \relax
5:7 & \begin{tabularx}{0.7\textwidth}{X} 「若他的力量不夠獻一隻綿羊，就要為所犯的罪，把兩隻斑鳩或是兩隻雛鴿獻給耶和華為贖愆祭：一隻作贖罪祭，一隻作燔祭。 \end{tabularx} \\ \\ \relax
5:8 & \begin{tabularx}{0.7\textwidth}{X} 他要把這些帶到祭司那裡，祭司就先把贖罪祭獻上，從鳥的頸項上扭斷牠的頭，但不把鳥撕斷。 \end{tabularx} \\ \\ \relax
5:9 & \begin{tabularx}{0.7\textwidth}{X} 祭司要把一些贖罪祭牲的血彈在祭壇的邊上，其餘的血要倒在壇的底座上；這是贖罪祭。 \end{tabularx} \\ \\ \relax
5:10 & \begin{tabularx}{0.7\textwidth}{X} 他要依照條例獻第二隻鳥為燔祭。祭司要為他所犯的罪贖了他，他就蒙赦免。 \end{tabularx} \\ \\ \relax
5:11 & \begin{tabularx}{0.7\textwidth}{X} 「他的力量若不夠獻兩隻斑鳩或兩隻雛鴿，就要為所犯的罪把供物，就是十分之一伊法細麵，獻上為贖罪祭；不可加上油，也不可加上乳香，因為這是贖罪祭。 \end{tabularx} \\ \\ \relax
5:12 & \begin{tabularx}{0.7\textwidth}{X} 他要把細麵帶到祭司那裡，祭司要取出滿滿的一把，作為紀念，按照獻火祭給耶和華的條例把它燒在壇上；這是贖罪祭。 \end{tabularx} \\ \\ \relax
5:13 & \begin{tabularx}{0.7\textwidth}{X} 至於他在這幾件事中所犯的任何罪，祭司要為他贖了，他就蒙赦免。剩下的都歸給祭司，和素祭一樣。」 \end{tabularx} \\ \\ \relax
5:14 & \begin{tabularx}{0.7\textwidth}{X} 耶和華吩咐摩西說： \end{tabularx} \\ \\ \relax
5:15 & \begin{tabularx}{0.7\textwidth}{X} 「若有人在耶和華的聖物上無意中犯了罪，有了過犯，就要獻羊群中一隻沒有殘疾的公綿羊給耶和華為贖愆祭，或依聖所的舍客勒所估定的銀子，作為贖愆祭。 \end{tabularx} \\ \\ \relax
5:16 & \begin{tabularx}{0.7\textwidth}{X} 他要為在聖物上的疏忽賠償，另外加五分之一，把這些都交給祭司。祭司要用贖愆祭的公綿羊為他贖罪，他就蒙赦免。 \end{tabularx} \\ \\ \relax
5:17 & \begin{tabularx}{0.7\textwidth}{X} 「若有人犯罪，在任何事上犯了一條耶和華所吩咐的禁令，他雖不察覺，仍算有罪，必須擔當自己的罪孽。 \end{tabularx} \\ \\ \relax
5:18 & \begin{tabularx}{0.7\textwidth}{X} 他要牽羊群中一隻沒有殘疾的公綿羊，或照你所估定的價值，給祭司作贖愆祭。祭司要為他贖他因不知道而無意中所犯的罪，他就蒙赦免。 \end{tabularx} \\ \\ \relax
5:19 & \begin{tabularx}{0.7\textwidth}{X} 這是贖愆祭；因他確實得罪了耶和華。」 \end{tabularx} \\ \\
[1ex]
\hline
\hline
\end{longtable}
(壹)贖愆祭之特點
\newline
\newline
一,預表耶穌之死,其血能除去我們日常所犯之愆尤.
\newline
\newline
二,向神向人,各有當盡之本分.
\newline
\newline
三,與贖罪祭之異點.
\newline
\newline
贖罪祭,經文指出犯之者何等人.(四章三節,十三節,廿二節,廿七節)即謂某人有罪；贖愆祭,經文指出所犯何罪(五章一至六節,十四至十九節,六章一至七節)即謂犯了何罪；根據此意,即可定贖罪祭是指贖人本罪.因為人得救,不必自陳犯了何罪,只認自己是罪人足矣,贖愆祭,是指赦除人日常所犯之罪愆,因為人若蒙赦免,須要一件一件罪惡承認,方可得赦而有平安,如是,人若按例獻了罪祭,則可云「我已蒙赦免了!」(即得救了)按例獻了贖愆祭,則上帝可云「背逆我的罪愆我已經除去了!」
\newline
\newline
(貳)贖愆祭之用處
\newline
\newline
其用處,就是贖人向神,向人,向己所犯之各種愆尤,即日常處己待人之諸般不義；表耶穌之流血,不獨為一次贖人類之本罪,且有隨時潔淨信者之愆尤.約翰一書一章九節與二章一節就是此意；民數十九章之紅牝牛,焚化成灰,用其灰作淨水,備為除各種不潔之用,此牝牛一次死,其潔除污穢之功效長久可用,此可表證耶穌之流血亦有此功效也.我們若云無罪就是自欺,以上帝為誑,所以我們無日無時不要靠賴我主流血功勞,蓋我一切愆尤,並用其寶血之能,保我不陷於罪也(約一書一章九節)
\newline
\newline
(參)罪愆之種類
\newline
\newline
甲,得罪自己之罪愆.(五章一至四節)
\newline
\newline
一,匿事不報.一節(得罪自己之良心,)(心罪,)我們信徒知見魔鬼陷害人,及耶穌之拯救人,我們匿而不報,是否有罪?何事得罪自己良心比此更甚?
\newline
\newline
二,摸了一切不潔.(二節三節)(自污,)(手足之罪,)即沾染世俗一切惡慾.
\newline
\newline
三,躁妄發誓,(四節)(自陷,)(口舌之罪)
\newline
\newline
乙,得罪神之罪愆.(五章十四至十九節)
\newline
\newline
一,在聖物上誤犯.(十五節)所謂在聖上犯罪,大扺是以下二事,一,擅取歸主之物二,應歸主之物而不歸主,如不納十分之一,(瑪三章八節)等事.此事舊約尚定為罪,何況新約?
\newline
\newline
二,誤犯誡命之一.(十七節)猶太人今有六百一十三條誡律,安息日行多一步,踏斷一柴,皆算是罪.誡律如是之多,有時犯了,尚不知；我們信徒雖不如是,但主之聖義愛律,不容易守,要常求主監察,日日靠其寶血之恩而保守；若有受主聖靈責備之事當悔改求赦.
\newline
\newline
丙,得罪他人之罪愆.(六章一至七節)
\newline
\newline
一,不忠於信託.(二節)傳道人須格外謹慎.
\newline
\newline
二,詭詐.(二節)社會無日無之,我們信徒要小心.
\newline
\newline
三,謀奪.(二節)欠債不願償,是不是屬於謀奪?
\newline
\newline
四,欺壓.(二節)舊教會曾犯此大罪.
\newline
\newline
五,假偽.(三節)所謊妄誓等．
\newline
\newline
(肆)如何獻贖愆祭?
\newline
\newline
一,祭物等級
\newline
\newline
用牝羊,或羊羔,或山羊,(六節)或斑鳩,雛鴿,(七節)或細麵,(十一節)或牡綿羊,(十六節)此皆因人之貧富力量而定；並因所犯之罪是向何方面而定.但最須注意的,是許用細麵為贖愆祭,此乃格外特例待遇,但須照以下三條例.一,用一伊法十分之一,即一日之糧,盡其全力而獻.二不傾油,表主耶穌在十字架死時,聖靈亦離之.三不加乳香,表耶穌在客西馬尼之求離苦杯,不蒙允許.
\newline
\newline
二,向神方面如何?
\newline
\newline
凡犯有以上各種罪愆之一,初不知道,後乃知,一知,即第一本份向主認,(五章五節,民五章七節)帶祭牲到主前,請祭司代獻贖愆之祭.
\newline
\newline
三,向受虧損者方面如何?
\newline
\newline
向受虧損者,要賠償,(六章五節)尚要加伍分之一而償；因此可以表明真誠悔改之意.倘若物主已死,則歸其親屬；若無親屬則照數加伍分之一歸於服事耶和華的司祭.(民五章五至八節)
\newline
\newline
以上兩方面向神獻祭認罪,是「挽」的工夫.向人賠償,是「還」的工夫.要挽,要還,此愆尤方可除去.此訓我們若有犯了各種之愆尤罪惡,應聽聖靈之指示,聖經之辨法,向兩方面清楚；否則只祈禱,亦不能除免心中的不平安.
\newline
\newline
(伍)贖愆祭之條例(七章一至七節)
\newline
\newline
一,要在殺燔祭牲處殺牲.(二節)
\newline
\newline
二,要洒血在祭壇之四圍.(二節)
\newline
\newline
三,要焚一切肥油及腰子(三,四,五節)
\newline
\newline
四,祭司男丁可食於聖所.(六節)
\newline
\newline
以上條例之預表與教訓見前.
\newline
\newline
五祭之總結
\newline
\newline
一,燔祭──耶穌為人向上帝完全盡獻己──上帝完全滿足而悅之
\newline
\newline
二,素祭──耶穌為人完全聖義──上帝與人完全滿足耶穌之為人
\newline
\newline
三,平安祭──耶穌為神人獻己使神人和好──罪由血得赦並因之而進至與主心交
\newline
\newline
四,贖罪祭──耶穌流血贖人本罪──人本罪得完全宥免
\newline
\newline
五,贖愆祭──耶穌流血潔人愆──人之愆尤盡得除去
\newline
\newline
總五祭之工作而觀是謂之完全救贖
\newline
\newline


\newpage

\section{第40屆~港九培靈會~吳勇~第一講}
\label{sec:1534}
\textbf{今日的時代}
\newline
\newline
連結: \href{https://www.hkbibleconference.org/session-message/view/1534}{\texttt{https://www.hkbibleconference.org/session-message/view/1534}}
\newline
\newline
\hyperref[sec:1332]{< < < PREV SERMON < < <}
~
\hyperref[sec:toc]{[返目錄]}
~
\hyperref[sec:1535]{> > > NEXT SERMON > > >}
\newline
\newline
耶利米書 1:2-1:3
\newline
\begin{longtable}{cl}
\hline
\hline
章節 & 經文 (和合本修訂版)\\
\hline
1:2 & \begin{tabularx}{0.7\textwidth}{X} 亞們的兒子猶大王約西亞在位第十三年，耶和華的話臨到耶利米。 \end{tabularx} \\ \\ \relax
1:3 & \begin{tabularx}{0.7\textwidth}{X} 從約西亞的兒子猶大王約雅敬在位的時候，直到約西亞的兒子猶大王西底家在位的末年，就是第十一年五月間耶路撒冷被擄時，耶和華的話也常臨到耶利米。 \end{tabularx} \\ \\
[1ex]
\hline
\hline
\end{longtable}
我們從聖經中看見耶利米的世代,與我們今日時代相同.耶利米在猶大王約西亞時作先知.神要他把祂的話向猶大百姓宣告,神把末後世代選民的光景顯給他們看,同時神也把末後世代要說的話講給我們聽.今天是末後的世代,正如路加福音廿一章所講到的兩件事.一是耶路撒冷要被外邦人踐踏,直到外邦人的日期滿了.很早以前,尼布甲尼撒王造了一個大像,這像頭是精金,腳是半鐵半泥的.我們觀察這個像,就知道今日的時代已經到了半鐵半泥的兩個時代,就是民主時代和極權政治的時代,這正是末後的外邦日期.二是無花果樹和各樣的樹發了芽.這是指甚麼說的呢?無花果樹是代表猶太國,發芽是指復國,猶太國在一九四八年已經復國了,同時世界上亦有許多弱小民族,在這同一期間建立自己的國家.
\newline
\newline
耶利米時代是末後的時代,同樣今天也是末後的世代；神在耶利米時代,叫選民看見了末後世代的光景,同樣神在今天也叫基督徒看見未後世代的光景.
\newline
\newline
現在我們要來思想兩件事:(一)耶利米時代的背景.(二)當時的屬靈光景.
\newline
\newline
當時的背景──耶利米出來做先知是在約西亞王在位十三年的時候,在此之前,以色列到民族已分為南北兩國,北國叫以色列國,南國叫猶大國.北國已經滅亡了一百多年,而南國仍存在.原因是猶大國還有一些敬虔的人,所以神保存了這個國家.以色列到國有十九個王,猶大國也有十九個王,但以色列到國沒有一個好王,而猶大國卻有像希西家,約坦,約西亞,烏西雅......這樣的好王,但最後還是要滅亡的.神不因為這個國家有了幾個敬虔的人,就一直容忍下去,神只不過給這個國家多一點機會和日子罷了!同樣,今日教會中亦有少數敬虔的人,所以神還留下教會的燈光,但是神不會一直容忍下去的,祂只是給予教會多一些機會和時間,若最後不悔改,神要吹熄教會的燈光,神要離開教會.
\newline
\newline
今日的世代與耶利米時代是一樣的,耶利米的一生經過五個王,其中只有一個是好王,其餘四個都是壞王.耶利米時代,好人少,壞人多.在啟示錄二章和三章講到七個神的教會,除了士每拿,非拉鐵非外,其餘的幾個均受神的責備和憎惡.它們之中好的少,今日的教會也如此.今日的社會好人少,壞人一大堆.主耶穌時代,吃餅的人多,跟隨主的人少.今日的教會屬靈的人少,屬世的人多.今日的信徒冷淡的多,熱心的人少.
\newline
\newline
耶利米時代,約哈斯,約雅敬,約雅斤,西底家都是壞王,而且老百姓也敗壞,所以這個國家從上到下都壞透了,牧養的人壞,受牧養的人也冷淡.正像今天教會,不但受帶領的人壞,帶領的人也壞,因為為他們冷淡,屬世.
\newline
\newline
列王記下廿三章廿五節說:「在約西亞以前,沒有王像他盡心盡性盡力的歸向耶和華,遵行摩西的律法；在他以後,也沒有興起一個王像他.」但壯志未酬身先死,而被埃及王所殺.這人真有摩西的心志,又忠心於神,不顧自己的得失.當時,雖有這樣好的人起來,做復興的工作,可惜這種復興像曇花一現就過了,原因全國的人,好的少,壞的多,從下到上都壞透了,故此不能有所作為.同樣,今天的教會也有一兩個屬靈的人起來,如約西亞一樣,也是無能為力的,很快就要消失的.耶利米時代的光景,也就是今天的光景.
\newline
\newline
屬靈的光景──那時代是一個背道的時代,神藉著耶利米的口,講出背道的事.他們把幼年的恩愛,婚姻的愛情都丟棄了,背道了.在啟示錄二章至三章裡講到七個教會,而老底嘉教會是不冷不熱的,神要把它從口中吐出去.其中第一個教會以弗所,它把起初的愛心離棄了,這是靈性墮落,軟弱,跌倒,冷淡的總因.若果一對夫婦,離棄起初的愛情,就會常吵架,冷淡,甚至動武,弄到家不安寧.
\newline
\newline
我們跟隨主,過了一段時間後,就與主疏遠了.有甚麼憑據呢?就是不像以前那樣親切,讀經不像以前渴慕,聚會不像以前熱心,工作沒有以前殷勤,這些現象就是與主的愛疏遠了.
\newline
\newline
以前,以色列人在曠野,在未曾耕種之地跟隨神,祂都記得,因為他們對主有愛情.對主有真正的愛,才能完全為祂.好像約伯對神的忠心,敬愛,他說:「賞賜的是耶和華,收取的也是耶和華.」(伯一21)他愛神並非因神給他許多的恩典.有一個青年,在戰場上與敵人作戰,不幸手打斷了,眼也盲了,他寫信給他的未婚妻說:「你收到這信就自由了,我們解除婚約了.」他得著回信說:「我的手作你的手,我的眼作你的眼.」乃是她愛他之故.又如路得與拿俄米的故事,拿俄米叫她回去找丈夫嫁人去,而她要跟隨拿俄米,她說:「不要催我回去不跟隨你,你往那裡去,我也往那裡去；你在那裡住宿,我也在那裡住宿；你的國就是我的國,你的神就是我的神；你在那裡死,我也在那裡死,也葬在那裡；除非死能使你我相離......」(得一16-17)她不怕得失,甚麼環境下都跟隨拿俄米.
\newline
\newline
今天我們跟隨主,要在任何環境中跟隨祂,不要道路平坦時才跟隨,崎嶇時就不跟隨,教會光景好才跟隨,不好就不跟隨,這樣做會受到神責備的.
\newline
\newline
神記得那時以色列歸耶和華為聖,即「那時」分別為聖,有幼年的愛,和婚姻的愛情,有起初的愛,分別出來為神,不但把時間分別出來為神,甚至連金錢,時間,整個人全都為神.今日的教會,各人把時間,金錢,自己,也不能分別為聖為神,乃是把起初的愛離棄了,背道了.
\newline
\newline
神揀選以色列民為祂的選民,要求他們成為獨居的民.但他背道了,離棄了神,他們的敬虔變成了外表的儀式,他注重獻祭,節期,和儀式.同樣,神今日要求我們是內心的敬拜,不是外表的儀式.今天我們有禮拜堂,但已經沒有實際了.正如耶利米責備當時的猶大人是背道的子民,今天,我們在神看來亦是背道的人.
\newline
\newline
耶利米責備他們說:「你的衣襟上有無辜窮人的血,你殺他們......」(耶一34)他們裡面有罪,裡面沒有感覺,是麻木了.好像患痳瘋的人,他們的皮膚是麻木的.有一個母親接到他兒子在戰場上受傷的電報消息,他傷得很重,生,死的成分各一半.她趕到那裡,司令官不准她進去看她的兒子,她苦苦要求,結果准了她,但要她不准對他說話,她進到醫院看見他的兒子被包著帶,她一聲不響,用手摸她的兒子,她兒子就說話了:「媽媽,你甚麼時候來的?」愛是敏感的.如果愛不存在,屬靈的感覺也失去了,麻木了.當時的猶大在神面前有罪,他們還說沒罪.
\newline
\newline
雖然他們背逆神,但是,神還是呼喚他們回來罷!他們卻一再的不理會.今日的光景也是一樣,神一再的呼喚我們回來,而我們卻不理會祂的呼聲.有一次,我到阿飛學校去講道.我看見一個廿一,二歲的青年,淪落為阿飛.因為他犯罪,被拉進這間學校,沒幾天就逃出來.他父母用車子送他回阿飛學校去,到了門口,他不下車,他父母拉他也不下,向他叩頭也不聽,甚至叩頭叩到流血也不受感動,他還說:「我看慣了,不要來這一套.」各位,我們對主也是這樣嗎?
\newline
\newline
今天我們與當時的猶大人一樣,屬靈的心麻木了,不聽神的呼喚,所以我們落在苦難當中.求神甦醒我們的心,不要再頑梗,不要做一個背道的人.
\newline
\newline


\newpage

\section{第40屆~港九培靈會~吳勇~第二講}
\label{sec:1535}
\textbf{哀哭}
\newline
\newline
連結: \href{https://www.hkbibleconference.org/session-message/view/1535}{\texttt{https://www.hkbibleconference.org/session-message/view/1535}}
\newline
\newline
\hyperref[sec:1534]{< < < PREV SERMON < < <}
~
\hyperref[sec:toc]{[返目錄]}
~
\hyperref[sec:1540]{> > > NEXT SERMON > > >}
\newline
\newline
耶利米書 2:9-2:11
\newline
\begin{longtable}{cl}
\hline
\hline
章節 & 經文 (和合本修訂版)\\
\hline
2:9 & \begin{tabularx}{0.7\textwidth}{X} 「我因此必與你們爭辯，也與你們的子孫爭辯。這是耶和華說的。 \end{tabularx} \\ \\ \relax
2:10 & \begin{tabularx}{0.7\textwidth}{X} 你們且渡到基提海島察看，派人往基達去留心查考，看可曾有過這樣的事。 \end{tabularx} \\ \\ \relax
2:11 & \begin{tabularx}{0.7\textwidth}{X} 豈有一國換了它的神明嗎？其實那不是神明！但我的百姓將他們的榮耀換了那無益的東西。 \end{tabularx} \\ \\
[1ex]
\hline
\hline
\end{longtable}
聖經曾記載耶利米為他的百姓哀哭,耶穌為耶路撒冷哀哭,保羅為他的肢體哀哭.求主也給我們流淚的眼睛,為教會,為失喪的靈魂,為罪惡的世界,為要臨到世界的苦難哀哭.
\newline
\newline
耶利米為何哀哭?猶大人拜了虛無的神,這是神最憎恨的事.神說「......留心查考,看曾有這樣的事沒有.豈有一國換了他的神麼,其實這不是神；但我的百姓,將他們的榮耀,換了那無益的神.」(耶二10-11)
\newline
\newline
耶利米時代,猶太人換了神,真神不拜,換假神來拜.今日世界百分之八十的人,心中是相信有神的.相信有神並不難,有兩個原因.(一)被理性說服,這是外在的原因強迫相信有神.我們居住的地球,本身除自轉外,還繞太陽轉,於是人類知道宇宙中有一位神存在.(二)內在的原因,人心覺得空虛,而追求一位永恆的神.有一次,我去看一個患了癌症的法官.起初我不知要向他說甚麼才好?後來我忽然對他說:「你知道你是一個罪人嗎?」他聽後望著我,我感到對他講了這句話很冒失.他在法庭上常指責犯人是一個罪人,但是,想不到今天反而有人在他面前指責他是一個罪人呢!他的態度起初很嚴肅,後來漸漸變為溫和了,他回答道:「我知道我是一個罪人.」我又對他說:「你知道自己是一個罪人,要帶罪去見神嗎?這是非常可怕的事.」他聽了之後,非常害怕,於是我便向他傳講神的救恩,他接受了,他也得救了.若果這時候你向他講學問,錢財,地位的事,對他一點都沒用,若你向他講神的事,他會不能不信神的.
\newline
\newline
今天雖然有百分之八十的人相信有神,但多數人是拜假神的.為甚麼拜假神,不拜真神呢?拜假神容易,不必用心靈和誠實來拜,不必有條件的拜,做君子也好,做小人也好,牠不會干涉你.你的生活清潔也好,污穢也好,牠不會來管你.可是,拜真神就不同了,主耶穌在約翰福音第四章對撒瑪利亞的婦人說:「神是個靈；所以拜祂的,必須用心靈和誠實拜祂.」你口稱祂為主,心也稱祂為主；口尊祂為王,心也尊祂為王；你外面所事奉的,也就是內心所事奉的.敬拜真神,祂有計劃要你參加,祂要萬人得救,不願一人沉淪,所以你要去傳福音給人聽.耶穌基督降世,為要拯救罪人,建立教會,以便能廣傳福音,牧養羊群,所以你不但是信徒,而且也是肢體,你要與人配搭.
\newline
\newline
主耶穌在世三十三年,被釘在十字架上,死了,復活,升天,祂不是死了就不存在了.祂要繼續的延長下去,要藉你的身體延長下去,所以你的身體就是祂的身體,讓祂可以繼續彰顯現每個世紀看,若果你不榮耀祂,祂就要干涉你,你羞辱祂,祂就要責備你.假神就不會這樣做.
\newline
\newline
耶利米哀哭,是因為他國家的猶大人把神換了假神來拜,拜假神不要有條件,內外不必一樣.牠不會責備人,牠也沒有甚麼計劃要你參加,拜牠很方便,沒有麻煩.拜真神就不同,要內外一樣.所以當時的猶大人拜方便的假神,把真神換假神來拜.今天我們很多人,雖然沒有把真神換假神來拜,但是卻把真神當假神來拜了,口與心不一樣.口稱祂為主,心卻不尊祂為主；口稱祂為王,心卻不尊祂為王.我外面與你和好,內心卻不與你和好.內心滿了污穢,虛假,和妒忌.今天的教會,雖然沒有把真神換假神來拜,但卻把真神當假神來拜,所以別人找不到見證,找不到果子.
\newline
\newline
猶大人假意歸向神
\newline
\newline
這是耶利米哀哭的第二件事.猶王約西亞是個好王,他遵行摩西的律法,聖經說:「凡猶大國和耶路撒冷所有交鬼的,行巫術的,與家中的神像,和偶像,並一切可憎之物,約西亞盡都除掉,成就了祭司希勒家在耶和華殿裡所得律法書上所寫的話.」(王下廿三24)耶利米書第四章裡講到神的忿怒不止息,但祂是一位慈悲,憐憫人的神,不輕易發怒的神,祂又是不長久責備的神.只要你回轉到祂那裡就夠了,像新約中浪子回到他父親那裡一樣,父親仍愛他.又如彼得三次不認主,主回頭看了他一眼,他就自責痛哭,主復活後,他就把主復活的事,向門徒報告.
\newline
\newline
既然約西亞王聽見神的律法,就起來除掉一切的偶像,交鬼的,行邪術的,何以神的怒氣不止息呢?乃是他們假意歸向神.表面上殿是修築起來了,內心的殿沒修起來.會眾守逾越節,表面是敬虔,內心並不敬虔.外面除去偶像,內心的偶像並沒除去.我們在聖經中找到一個原則,荒涼的原因是神離開了,復興的原因是神回來.但是,神究竟怎樣離開,怎樣回來.例如神不是從約書亞離開,就是他跪到死求到死,甚至哭到死,神也不能回來.因祂是從亞干離開的,所以只有從亞干身上回來.直至他們把亞干打死後,神就回來了,這是千古不變的原則.今天的人也是如此,只是表面上歸向神.
\newline
\newline
當時的猶大人在約西亞王領導下,所做的復興工作是表面的,並非從內心歸向神.正像今日的教會,有的只是波浪式的復興,一下子復興,起來一下子又下去,其原因是假意歸向神.今天有很多愛主的弟兄姊妹,看見今日的教會光景就不滿足,所以便禁食禱告,甚至晝夜禱告向神求復興.可是神的話說,荒涼的原因是因為罪.我們就來對付罪,表面的罪對付了,內心的罪卻不肯對付；不必付代價就對付,要付代價的就不對付.對人表面和睦,內心卻恨那人,這樣做就是假意歸向神.當時的猶大人是如此,今的教會也如此.所有的復興工作,就像曇花一現似的,很快就完了.
\newline
\newline
聖經說:「你們當在耶路撒冷的街上,跑來跑去,在寬闊處尋找,看看有一人行公義,求誠實沒有,若有我就赦免這城.」(耶五1)當時沒有一人行公義,求誠實,這是耶利米哀哭的第三件事.神說:「我在他們中間尋找一人重修牆垣,在我面前為這國站在破口防堵,使我不滅絕這國；卻找不著一個.」(結廿二30)「一個」不能按表面字句來解釋,因為在那個時代,除了耶利米之外,還有哈巴谷,西東雅等都是敬虔的人.「一個」的解釋應該是指敬虔的人,體貼神心的人很少而說的.主耶穌時代,吃餅的人很多,真正跟隨主的人就很少.今天聽道的人很多,行道的人也很少.傳福音是一場屬靈的爭戰,要把撒旦權勢下的人救出來,也許參加的人很多,但最後能夠站住的人卻很少.例如在教會中有一個弟兄出了事,指責批評的人很多,安慰代禱的人很少.教會今日受撒旦攻擊,正在苦難當中,教會中肢體願同享受的很多,願同甘苦的很少.主耶穌教導我們要同心合意興旺福音,願意在這工作上有份的人很多,但在傳福音時受到打擊就後退了,所以還能站在崗位上的人則非常之少.
\newline
\newline
一個人要從內心歸向神,才有希望.如尼希米回國重建城牆的事,百姓響應他,所以他的工作很快就做好.今天的教會要用心靈和誠實拜祂,才有希望.
\newline
\newline
神說祂尋找一個人,表明神的恩典浩大,其次是表明這個人很重要,神要藉著他把祂的恩典流出去.
\newline
\newline
當時連一個真心歸向神的人都沒有,所以耶利米哀哭了.同樣,在今天神也為這時代哀哭.求主幫助我們,叫我們歸向祂,是出於真意的,而不是出於假意的.
\newline
\newline


\newpage

\section{第40屆~港九培靈會~吳勇~第三講}
\label{sec:1540}
\textbf{三件傷心的事}
\newline
\newline
連結: \href{https://www.hkbibleconference.org/session-message/view/1540}{\texttt{https://www.hkbibleconference.org/session-message/view/1540}}
\newline
\newline
\hyperref[sec:1535]{< < < PREV SERMON < < <}
~
\hyperref[sec:toc]{[返目錄]}
~
\hyperref[sec:1541]{> > > NEXT SERMON > > >}
\newline
\newline
耶利米書 7:1-7:10
\newline
\begin{longtable}{cl}
\hline
\hline
章節 & 經文 (和合本修訂版)\\
\hline
7:1 & \begin{tabularx}{0.7\textwidth}{X} 耶和華的話臨到耶利米，說： \end{tabularx} \\ \\ \relax
7:2 & \begin{tabularx}{0.7\textwidth}{X} 「你當站在耶和華殿的門口，在那裡宣講這話說：所有從這些門進來敬拜耶和華的猶大人哪，當聽耶和華的話。 \end{tabularx} \\ \\ \relax
7:3 & \begin{tabularx}{0.7\textwidth}{X} 萬軍之耶和華－以色列的神如此說：你們要改正你們的所作所為，我就使你們仍然居住這地。 \end{tabularx} \\ \\ \relax
7:4 & \begin{tabularx}{0.7\textwidth}{X} 不要倚靠虛謊的話，說：『這是耶和華的殿，是耶和華的殿，是耶和華的殿！』 \end{tabularx} \\ \\ \relax
7:5 & \begin{tabularx}{0.7\textwidth}{X} 「你們若實在改正你們的所作所為，彼此誠然施行公平， \end{tabularx} \\ \\ \relax
7:6 & \begin{tabularx}{0.7\textwidth}{X} 不欺壓寄居的和孤兒寡婦，不在這地方流無辜人的血，也不隨從別神陷害自己， \end{tabularx} \\ \\ \relax
7:7 & \begin{tabularx}{0.7\textwidth}{X} 我就使你們仍然居住這地，就是我從古時所賜給你們祖先的地，從永遠到永遠。 \end{tabularx} \\ \\ \relax
7:8 & \begin{tabularx}{0.7\textwidth}{X} 「看哪，你們倚靠虛謊無益的話語。 \end{tabularx} \\ \\ \relax
7:9 & \begin{tabularx}{0.7\textwidth}{X} 你們豈可偷盜，殺害，姦淫，起假誓，向巴力燒香，隨從素不認識的別神， \end{tabularx} \\ \\ \relax
7:10 & \begin{tabularx}{0.7\textwidth}{X} 又來到這稱為我名下的殿，在我面前敬拜，說『我們平安無事』，為了要行這一切可憎的事呢？ \end{tabularx} \\ \\
[1ex]
\hline
\hline
\end{longtable}
現在我要講到耶利米時代,有三件可怕的事發生,而令到神的心非常的傷心.
\newline
\newline
(一)這不是耶和華的殿
\newline
\newline
那時,猶大人分為南北兩國,北國以色列國早在一百年前被亞述所滅,這時候,猶大國仍存在.但是,猶大國也是危機四伏,故要靠與別國聯盟才能生存,而如今別國也靠不住了.這時候猶大王約西亞看見有神的國家才興盛,沒有神的就衰亡,所以他就起來領導民眾求神,修聖殿.殿修好後,各處各方的人都來到耶路撒冷守逾越節,祭牲成千成萬,祭司穿著莊嚴的衣服,群眾唱歌,敲鑼擊鼓,大家都非常虔誠敬拜,列王記下廿二章廿二節說:「自從士師治理以色列人,和以色列王猶大王的時候,直到如今,實在沒有守過這樣的逾越節.」因此民眾便高興的叫喊,說這是耶和華的殿,以為殿修得好看,又有這麼多人守節,獻這麼多祭牲,祭司穿得很莊嚴,會眾的歌聲又那麼好聽,就說是耶和華的殿了.
\newline
\newline
可是,耶利米卻說這不是耶和華的殿,與他們的看法完全相反.甚麼才是耶和華的殿呢?所羅門王建好聖殿又把它奉獻給神,於是耶和華的榮光充滿這殿,有神的同在,才叫耶和華的殿.經上說:「豈不知你們的身子就是聖靈的殿；這聖靈是從神而來,住在你們裡頭的.」(林前六19)有神的同在,才叫神的殿.若沒有神的同在,不是神的殿.儘管殿做很美,儀式也莊嚴,祭牲很多,也不能叫做耶和華的殿.
\newline
\newline
在此情形之下,耶利米指示他們,怎樣讓神回到他們中間?要改正他們的行為,對鄰舍要施行公平,做公務員的要清潔待人,做生意的要誠實待客人,要有慈愛叫人得溫暖,穿不暖,吃不飽的人,你要看顧他.接著神又責備他們說:「看哪!你們倚靠虛謊無益的話.」(耶七8)他們偷盜,殺害,姦淫,起假誓,向巴力燒香,並隨從素不認識的別神.甚麼是偷盜呢?從金錢來說,當納的不歸給神,在神的眼光中這算為偷盜.又如禮拜日是聖日,不歸神用,而為自己用,神說:「你若在安息日掉轉你的腳步,在我聖日不以操作為喜樂,稱安息日為可喜樂的,稱耶和華的聖日為可尊重的；而且尊敬這日,不辦自己的私事,不隨自己的私意,不說自己的私話.」(賽五十八13)這一日當為神,我們用來為自己,這就是偷盜.神責備他們中間有殺害的事發生,殺害就是傷人,我們沒有用刀傷人,但因懷恨之故,就會用言語造謠,毀謗來傷人,這是殺害.
\newline
\newline
來到殿裡的人,又向殿敬拜的人,都是對鄰舍不公平的人,欺壓孤兒寡婦的人,偷盜殺害的人,這些人在其中,神怎能在他們當中呢?如果要讓神回到他們當中,就要完全改正行為,否則,雖把殿修得再好看,熱熱鬧鬧守節也不是神的殿.當時殿能給個人或家庭帶來祝福,平安,希望；但是沒有神的殿,徒有外表,沒有實際,絕不能給人帶來祝福,平安和希望.如果把不公平,沒憐憫,偷盜,殺害各樣的罪帶到神名下的殿,只有加增神的咒詛和忿怒.例如以利的兩個兒子,以為約櫃搬來了就有希望,豈知反招咒詛,就當他們出征打仗時,死在非利士人的手下.
\newline
\newline
教會的工作是從罪中救出一些人來,使這些人從黑暗進入光明,從污穢進到清潔,神因這班人就賜福給這個時代,如果這班人,只注重敬虔的儀式,沒有敬虔的實際,這樣的教會,不但不能帶給時代平安,希望,反招神的忿怒與咒詛.
\newline
\newline
(二)輕輕忽忽的醫治神的百姓
\newline
\newline
這是耶利米責備當時祭司的罪,輕忽醫治神性的罪.如果一個醫生這樣對待他的病人,不外有兩個原因:(一)不看人命為重要.(二)節省寶貴時間,多看一些病人,多賺一點錢.今日的教會,輕輕忽忽的帶人信主,不看靈魂為一些,如醫生不看病人為寶貴一樣.醫生輕忽醫治的結果會:(一)叫病人的病加重(二)可能叫病人死亡.(三)關係病人,也關係到自己.(四)影響病人的家庭.(五)影響自己的聲譽.
\newline
\newline
教會輕忽的醫治也有以下幾種結果:(一)隨便叫人信主,以為他參加了慕道班,不久就給他施洗,加入教會,就恭喜「得救」了.其實他還沒得救,沒平安,不是弟兄,乃是你造成的,這樣就使他今世痛苦.他未信前痛苦,信了更痛苦,因為他沒生命,你要求他過神的生活,是不行的.如一個猴子,你要求他過人的生活,這是不可能的事.隨便叫人加入教會,這是叫人死亡.(二)關係到主,這人不是「基督徒」之前,好壞與主無關,如今輕忽的醫治,掛了基督徒的名,犯罪作惡,主的名因他受羞辱,他有臭名,就關係到主.(三)關係到教會.我在歐洲時,看見教會沒有復興的氣氛,一片死寂.在挪威看見一間教會,它的旁邊都是墳墓,做死的工作當然是死,做活的工作當然是活.例如一粒麥子,接受水分陽光就長出來,若你對一塊石頭培植,永遠培不出生命來.(四)關係自己,一個傳道人,輕輕忽忽的醫治,不但關係別人,而且也關係到自己,名譽不好.這樣,我們的重擔就越來越大,內心越來越痛苦,信心也拖跨了.所以,今日有很多神的工人,從工場上退了出去,不但害了自己,而且也害了別人.你引他入教,使他成了沉淪之子,神說:「這惡必死在罪孽之中,我卻向你討他喪命的罪.」(結三18)
\newline
\newline
(三)對神如林中的獅子
\newline
\newline
聖經說:「我的產業向我如林中的獅子；他發聲攻擊我；因此我就恨惡他.」(耶十二8)「產業」是指猶大百姓,亦指今日的基督徒.耶穌是猶大支派的獅子.民數記中講到猶大支派的旗號是獅子,神要求選民有獅子的樣式.在啟示錄第四章裡講到四活物,其實「活物」與「生命樹」同一字根,有生命的人才有人的面,像牛犢,像飛鷹,像獅子.耶穌是我們的榜樣.獅子不是用來對付神的,而是要對付撒旦的.如主到聖殿去,就像獅子,趕出在聖殿裡做買賣的人.又如有一個牧師去探訪他的教友,那教友正在打麻將,等他打完後,他對牧師說:「對不起,久等了.」牧師答道:「不要緊!」耶穌是神的羔羊,我們是神的羊.約翰說主是好牧人,我們與主的關係是牧人與人,這裡說,「我的產業向我如林中的獅子」,這就不對了.當時的百姓攻擊神,抵抗神.
\newline
\newline
神不是因子民一有罪就離開,乃是先警戒,然後才刑罰他們.猶大人有罪,神用耶利米警戒他們達四十年之久,他們刑罰他,這警戒就停止.耶利米還是不停的警戒他們,他們要把他處死,兇得像獅子,不肯接納神的警戒,所以神就離開他們了.如掃羅就是一個很好的例子,神先警戒他在吉甲做偷取牛羊的事,後來他更加犯罪,神警戒他,他不聽,結果神就離開他.正如神所說的:「我離了我的殿宇,撇棄我的產業,將我心裡所親愛的,交在他仇敵手中.」(耶十二7)今日神對我們,也是先警戒,不聽才刑罰,交給仇敵.又如所羅門做王,順服神,便四境平安,後來他犯罪,娶外邦女子為妻,又拜偶像,結果四境的仇敵,起來與他作對.神用聖經感動他,看見自己的罪,多麼可怕.叫他出醜,叫他知罪,叫人看,為要叫他不敢驕傲.神是願意遮蓋人的罪惡的,但到了人不肯把罪放在主的寶血下時,祂就放在人的眼下.
\newline
\newline
當時猶大人對神如獅子,抵抗神,結果神沒辦法,才把他們交給仇敵,這是神最傷心的事.獅子是對撒旦的,不是對付神的.但剛好相反,對付神的是獅子,不是對付撒旦,結果,神就把他交給仇敵.今日我們有很多軟弱,失敗,叫神傷心.但沒有一件事比對神像獅子更傷心的.求主幫助我們接受這警戒.
\newline
\newline


\newpage

\section{第40屆~港九培靈會~吳勇~第四講}
\label{sec:1541}
\textbf{呼召耶利米}
\newline
\newline
連結: \href{https://www.hkbibleconference.org/session-message/view/1541}{\texttt{https://www.hkbibleconference.org/session-message/view/1541}}
\newline
\newline
\hyperref[sec:1540]{< < < PREV SERMON < < <}
~
\hyperref[sec:toc]{[返目錄]}
~
\hyperref[sec:1543]{> > > NEXT SERMON > > >}
\newline
\newline
耶利米書 1:5-1:9
\newline
\begin{longtable}{cl}
\hline
\hline
章節 & 經文 (和合本修訂版)\\
\hline
1:5 & \begin{tabularx}{0.7\textwidth}{X} 「我尚未將你造在母腹中，就已認識你；你未出母胎，我已將你分別為聖，派你作列國的先知。」 \end{tabularx} \\ \\ \relax
1:6 & \begin{tabularx}{0.7\textwidth}{X} 我就說：「唉！主耶和華，看哪，我不知道怎麼說，因為我年輕。」 \end{tabularx} \\ \\ \relax
1:7 & \begin{tabularx}{0.7\textwidth}{X} 耶和華對我說：「不要說『我年輕』，因為我差遣你到誰那裡去，你都要去；我吩咐你說甚麼話，你都要說。 \end{tabularx} \\ \\ \relax
1:8 & \begin{tabularx}{0.7\textwidth}{X} 你不要怕他們，因為我與你同在，要拯救你。這是耶和華說的。」 \end{tabularx} \\ \\ \relax
1:9 & \begin{tabularx}{0.7\textwidth}{X} 於是耶和華伸手按住我的口，對我說：「看哪，我已將我的話放在你口中。 \end{tabularx} \\ \\
[1ex]
\hline
\hline
\end{longtable}
那時代是敗壞,人心邪惡的時代,神便召耶利米出來工作,做神的器皿,做神的先知.今日也是末世的時代,神要召像耶利米這樣的人出來工作.
\newline
\newline
神甚麼時候呼召耶利米呢?聖經說:「我未將你造在腹中,我已曉得你,你未出母胎,我已分別你為聖；我已派你作列國的先知.」(耶一5)他未出母胎,神已派定他的工作了,可見耶利米作先知,是神預定的,神先是預知才有預定.神先預知耶利米可以用,所以然後才預定耶利米來用.神預定用耶利米,並非胡裡胡塗的預定,乃是預知才預定.由此可見,神要用一個人並非無目的.
\newline
\newline
神用耶利米有三個原因,茲述如下:
\newline
\newline
(一)他沒有結婚
\newline
\newline
神對他說:「你在這地方不可娶妻,生兒養女.」(耶十六1)是不是沒有結婚才能為神用呢?這要看神對某一人的要求而決定的,例如舊約的耶利米是沒有結婚的,又如何西阿先知是結過婚的.新約的保羅是沒有結過婚的,而彼得是結過婚的,由此可知不能解釋沒結婚才能為神用.耶利米沒有結婚,乃是神吩咐的,原因是神在那時代交給他一項工作,要他完成.他不結婚才能擔任神交託的工作和見證.「結婚」如何解釋呢?婚姻是神所預定的,是一種權利,婚姻是建立家庭,是一種享受.例如神看亞當獨居不好,就為他造一個配偶給他,要他建立家庭,享受婚姻的權利和快樂.神不是不許可人有應當的權利和享受,乃是為著時代的需要,好像耶利米一樣,他答應神的要求,為著神的工作,犧牲當有的權利和享受.這是耶利米的存心,所以他能被神使用.
\newline
\newline
今天是怎樣的時代?是一個非常的時代,社會是那樣的邪惡,人心是那樣的詭詐,不潔的事,日益增多,今日的時代正如路加福音第十七章講的挪亞的時代一樣,人心中所想的盡都是惡,地上滿了強暴,人心險惡,社會敗壞.今日的教會,一片荒涼,不能作光作鹽.不像使徒時代的教會能夠影響社會.但是,今日的教會,不但不能影響社會,反而受社會的影響呢!信仰不純正的教會,索性變成社會組織,純正的不甘心變質,就靠一些例定的聚會來拖延.今天,神盼望有人能答應祂的要求,犧牲自己的權利和享受,才有能力去做神的工作.
\newline
\newline
怎樣犧牲自己的權利和享受呢?例如睡眠是人的權利和享受.主為天國計劃,揀選門徒,祂就把這權利和享受放棄,整夜禱告.又如吃飯也是人的權利和享受,主知道神的計劃,要在地上設立天國,主要執行這計劃,就必須經歷大爭戰,所以祂便把四十天吃飯的權利和享受放棄了,才勝過撒旦的攻擊.我認識一個弟兄,他是教會的執事,他本來是一個科學家,在幾年內他帶領很多的台灣青年信主.神得著他,大大使用他.有一次有一個國際機構以重金聘請他,他為著神的工作,卻放棄了這個優厚的聘請.
\newline
\newline
今日為著傳福音救靈魂,為著建立教會推廣天國,我們常常發現傳福音傳不出去,建立教會建立不起來,因此,我們就以為人心剛硬,神的話格格不入,以為肢體間問題那樣多,所以配搭不起來,其實冥冥之中有一股看不見的反叛力量,就是神的仇敵,它在攔阻,在破壞.
\newline
\newline
神要人為著祂的工作,犧牲更多的權利和享受.讓神的能力大量傾倒在我們身上.我們要為神的緣故,要澈底奉獻,更多禁食,禱告.
\newline
\newline
(二)他忠心傳神的話
\newline
\newline
神對他說:「你要對這百姓說,耶和華如此說,看哪!我將生命的路,和滅亡的路,擺在你們面前.住在這城裡的,必遭刀劍,饑荒,瘟疫而死；但出去歸降圍困你們迦勒底人的,必得存活,要以自己的命為掠物.」(耶廿一8-9)神給他的信息,是叫他們向敵人投降.耶利米因這信息受苦,甚至死亡.但他還是要傳,他所說不吉利的話,人不歡迎,但他不討人的喜悅,乃討神的喜悅,所以他還得要傳.他為神的話受苦,五次被囚,有一次被囚在沒水而有淤泥的爛土中,受盡苦楚.以致後來太監捉他起來時,要用碎布和破爛衣服來墊住他的路膊,可見他是瘦得皮包骨,瘦骨如柴.
\newline
\newline
他為神的話,受人冤曲,被人打,說他是敵國的間諜,被人來羞辱,幾乎喪命,但他仍對神忠心.他是一個屬靈的勇士,雖然他受盡了痛苦,受盡了冤曲,到了太監以伯米勒把他從土牢中提出後,領他去見西底家王,他還是照樣傳,這一切苦沒有叫他喪志,甚而變節.正如本仁約翰一樣,官員勸他不傳福音,就釋放他,他卻說:「今天釋放,明天就去傳.」一切打擊,不能叫他喪膽,折磨不能叫他灰心.耶利米預言猶大人要滅亡,這是百姓不願聽,不肯聽的話,他所傳的沒有功效,但他還是要傳,因他對工作的態度,只顧忠心於神,不計較自己的成敗.
\newline
\newline
今天,神需要這種人.神的工作,一定有撒旦的攻擊,無中生有的冤枉,慘無可忍的痛苦.在哥林多前書九章裡講到保羅受了撒旦攻擊,也是同樣的方法,撒旦攻擊他不娶妻,因為他守獨身,用情慾來攻擊他.這是撒旦破壞主工人的聲譽,使人對他失信用.主工人失了信用,所傳信息不是全無效,便要打折扣.因為主工人與弟兄姊妹有分不開的關係,到了主工人失去帶領作用,一般羊群,就徬徨失措.今日需要耶利米這樣的人起來,才能站主得住.他雖受這麼大的打擊,但信息始終不變,今日有很多人受一點的打擊後.就愛心變了,信心變了,甚至服事的崗位也變了.耶利米是一個屬靈的勇士,他始終站在那裡,受毀謗站在那裡,受痛苦站在那裡,甚至知道工作沒果效也站在那裡.到了耶路撒冷陷落後,巴比倫護衛長尼布撒拉旦釋放他,且請他去巴比倫,他拒絕了,仍留在耶路撒冷,還是站在那裡,要與百姓同受苦難.好像李溫斯敦一樣,當他橫過非洲大陸腹地後,船長勸他回國,可得大榮譽,但他還是回去非洲,到死都站在那裡.
\newline
\newline
(三)他有火熱的心
\newline
\newline
耶利米為神工作的心火熱,他裡面有一股力量.世界上需要火的力量,飛機需要火來騰空,鋼鐵需要火來熔化,封閉的東西,需要火來打開.如以利亞在迦密山上,火從天上一降下,巴力的先知被殺,以色列人得復興,認識耶和華是神.今日的教會更需要火的力量,才可以把冷淡的信徒推動起來!
\newline
\newline
神的話是火的來源,聖經說:「我得著你的言語,就當食物吃了；你的言語,是我心中的歡喜快樂.」(耶十五16)神的話是力量,如以馬忤斯路上的門徒,主與他們同行,解經給他們聽,他們就心裡火熱.食物化成身體的一切,神的話是生命的一切,生命是生活的一切.耶利米把神的話當作食物吃了.我們讀到創世記對該隱說的話:「你的兄弟亞伯在那裡?」(創五9)我們就要省察,有沒有弟兄隨流失去,要我們去尋找.又當我們讀到約翰福音叫人要合而為一時,我們是否因肉體血氣的緣故與人隔閡,亮光不同與人嫌隙.我們的身子是神的殿,在他裡面有沒有牛羊鴿子和兌換銀錢的人?你要省察.你若把神的話化成生活的一切,就會火熱起來.從耶利米身上,我們就見到聖靈的靈力,他本來是一個軟弱的人,變成一個剛強的人,他從膽怯變成膽壯.當時環境很可怕,但他勝過了,雖責任很難擔當,他也負起來了,痛苦很難受他也忍受了,工作沒有果效他仍對神忠心,聖靈給他一切的力量來負起神交託給他的使命.
\newline
\newline
我們要向耶利米學習,在任何的環境和打擊下,要對神至死忠心,忍受一切的痛苦,不要怕困難,腳步要站穩,要犧牲自己的享受和權利,像耶利米那樣,才能被神使用,只有這樣,才能被神重用.
\newline
\newline


\newpage

\section{第40屆~港九培靈會~吳勇~第五講}
\label{sec:1543}
\textbf{復興的亞伯拉罕}
\newline
\newline
連結: \href{https://www.hkbibleconference.org/session-message/view/1543}{\texttt{https://www.hkbibleconference.org/session-message/view/1543}}
\newline
\newline
\hyperref[sec:1541]{< < < PREV SERMON < < <}
~
\hyperref[sec:toc]{[返目錄]}
~
\hyperref[sec:1542]{> > > NEXT SERMON > > >}
\newline
\newline
創世記 12:9-12:20
\newline
\begin{longtable}{cl}
\hline
\hline
章節 & 經文 (和合本修訂版)\\
\hline
12:9 & \begin{tabularx}{0.7\textwidth}{X} 後來亞伯蘭漸漸遷往尼革夫去。 \end{tabularx} \\ \\ \relax
12:10 & \begin{tabularx}{0.7\textwidth}{X} 那地遭遇饑荒。亞伯蘭因那地的饑荒嚴重，就下到埃及，要在那裡寄居。 \end{tabularx} \\ \\ \relax
12:11 & \begin{tabularx}{0.7\textwidth}{X} 將近埃及，他對妻子撒萊說：「看哪，我知道你是美貌的女人。 \end{tabularx} \\ \\ \relax
12:12 & \begin{tabularx}{0.7\textwidth}{X} 埃及人看見你會說：『這是他的妻子』，他們就會殺我，卻讓你活著。 \end{tabularx} \\ \\ \relax
12:13 & \begin{tabularx}{0.7\textwidth}{X} 所以，請你說你是我的妹妹，使我可以因你得平安，我的性命也因你存活。」 \end{tabularx} \\ \\ \relax
12:14 & \begin{tabularx}{0.7\textwidth}{X} 亞伯蘭到達埃及時，埃及人看見那女人極其美貌。 \end{tabularx} \\ \\ \relax
12:15 & \begin{tabularx}{0.7\textwidth}{X} 法老的臣僕看見了她，就在法老面前稱讚她。那女人就被帶進法老的宮中。 \end{tabularx} \\ \\ \relax
12:16 & \begin{tabularx}{0.7\textwidth}{X} 法老就因她厚待亞伯蘭，給了亞伯蘭許多牛、羊、公驢、奴僕、婢女、母驢、駱駝。 \end{tabularx} \\ \\ \relax
12:17 & \begin{tabularx}{0.7\textwidth}{X} 耶和華因亞伯蘭妻子撒萊的緣故，降大災擊打法老和他的全家。 \end{tabularx} \\ \\ \relax
12:18 & \begin{tabularx}{0.7\textwidth}{X} 法老召了亞伯蘭來，說：「你向我做的是甚麼事呢？為甚麼沒有告訴我她是你的妻子？ \end{tabularx} \\ \\ \relax
12:19 & \begin{tabularx}{0.7\textwidth}{X} 為甚麼說『她是我的妹妹』，以致我把她接來要作我的妻子呢？現在 ，看哪，你的妻子在這裡，帶她走吧！」 \end{tabularx} \\ \\ \relax
12:20 & \begin{tabularx}{0.7\textwidth}{X} 於是法老吩咐人把亞伯蘭和他妻子，以及他一切所有的都送走了。 \end{tabularx} \\ \\
[1ex]
\hline
\hline
\end{longtable}
我們走屬靈的路可能會跌倒,跌倒的情形有兩種:(一)跌倒沒起來.(二)跌倒又起來.亞干跌倒沒起來,大衛跌倒後起來了,新約裡有猶大跌倒沒有起來,彼得跌倒後起來了.現在讓我們來看跌倒的亞伯拉罕.
\newline
\newline
(一)亞伯拉罕的跌倒
\newline
\newline
他跌倒是漸漸遷往埃及的,正如羅得漸漸遷往所多瑪一樣,可見跌倒並非一下子就跌倒的,乃是漸漸跌倒的.撒旦在我們身上工作,不是一下子就把我們拉下世界的,乃是一步一步的.不是一下子離開神的,乃是一步一步的.如讀經,不會一下子就停止讀經,乃是一步一步的叫你不讀經,今日叫你忙,沒時間讀,明日叫你煩,使你讀不進去.今日叫你讀不出味道來,明天讀得厭煩,慢慢的就不讀經了.同樣,禱告不會一下子就不禱告,撒旦乃叫你一步一步的不禱告,今天叫你身體疲累,禱告不能長久,明天叫你心意紛亂而禱告不進去,這樣下去就會慢慢的不禱告了.
\newline
\newline
亞伯拉罕為何會跌倒?他因飢荒才下埃及去的,飢荒是苦難,他下埃及就跌倒了.今天,許多人因苦難而跌倒.迦南是神應許之地,亞伯拉罕走進迦南,是走在神的旨意裡面,怎會有飢荒?又像神吩咐以利亞去到基立溪旁,過了些日子,溪水乾了,這是神吩咐的地方,這也是神旨意的地方,怎會乾呢?這是許多人的問題.其實,神旨意的迦南會有飢荒的,神旨意的基立溪水會乾的,好像保羅想去庇推尼,但聖靈阻止他,而要他去馬其頓,然而他到了馬其頓的第一個城市,就被人拉去坐牢.他去馬其頓是神的旨意,有苦難,神會負責的.如以利亞來到基立溪旁,不數日溪水乾了,他是走在神旨意裡,所以神會負責,祂吩咐烏鴉在那裡供養他.迦南是神旨意的地方,祂也會負責.
\newline
\newline
去年我花了美金四百三十元去印度,是聖靈感動我去的.未去之前,有些西教士勸我不要去,因為印度政府不歡迎傳道人入境.神要我去,我就去了.當我到達印度之時,在海關就遇到印度移民官的為難,他們要我答應兩個他們提出的條件:1.明早九點去見移民廳長.2.在印度前後只准住一個禮拜.我都答應了,後來他們帶我到一間古老宏大的沒人住的神學院裡.到了半夜有人打門,我便很害怕,以為有流氓分子來搗蛋,他們見到我問起我的姓名後便說:「你把我們弄得很苦呀!」我想,我才第一天到印度,怎麼會有這種事呢?我真是百思莫解.接著其中的一個對我說:「你的手續是我給你辦的.我告訴你,明天不要去見移民廳長,而且你喜歡在印度住多久就住多久.」他說完,便用筆把海關移民官寫給我的兩個條件劃掉.我到印度神為我負責.
\newline
\newline
亞伯拉罕下埃及跌倒,不是犯罪跌倒,乃是對神失了依靠的信心,就跌倒了.神叫亞伯拉罕跌倒是有目的的,是要叫他認識自己.如彼得在雞叫之前,要三次不認主乃是主所預言的,主為何不妨止他呢?主讓他跌倒,是要他認識自己.因為彼得是一個充滿自己的人.
\newline
\newline
亞伯拉罕下到埃及,他怕自己的妻子因長得美,怕人殺了他奪他的妻子去,埃及人是沒有道德的,故此他就與妻子撒拉商量,不要說是夫妻關係,而說是兄妹關係.他說謊了,雅各也有這毛病,他不是以掃,裝在父親面前說是以掃.亞伯拉罕是個假裝,說謊的人.他為自己犧牲別人,如果說是夫婦,也許埃及王不要她,她是別人的妻子要來做甚麼,如果說撒拉美麗,當時帝王專政,還怕埃及沒有美女嗎?這是犧牲別人保護自己.
\newline
\newline
亞伯拉罕犧牲別人而謀自己的益處,他因著撒拉得了許多牛羊駱駝.有一次,我在外國講道,散會後有一個女子哭著來見我,她父母要把她嫁給一位年齡比她大兩倍的老頭子為妻,因他很有錢,她父母想得些錢來養老.這對父母是專為自己的利益打算.
\newline
\newline
亞伯拉罕的作為,真卑鄙,真下賤,於是法老召他來,當面責備他.有一個牧師訴訟在法官面前,法官指責他與世人一樣,怎能叫人信主呢!亞伯拉罕是羞恥的,下賤的,神叫他認識自己的光景.
\newline
\newline
(二)復興的亞伯拉罕
\newline
\newline
他下埃及,又從埃及回來,從那裡跌倒,便從那裡出來.他容讓羅得,羅得因牛羊的事與亞伯拉罕的牧人相爭,於是他便對羅得說:「你我不可相爭,你的牧人和我的牧人也不可相爭,因為我們是骨肉.」(創十三8)今天的社會有相爭,明爭暗鬥,為名爭,為利爭,因相爭的緣故使用手腕手段,陰謀詭計,所以人情就越來越淡薄,人心越來越兇險,社會變成烏煙瘴氣.今大的教會也有相爭,明爭暗鬥,雖沒有為錢財相爭,但有為意氣而爭,所以愛心就淡薄了,肢體疏遠了,教會變成烏煙瘴氣.
\newline
\newline
亞伯拉罕對羅得說,我們因是骨肉,不可相爭.我們是神的兒女,今天有許多相爭,有的人主張教會在千禧年被提,有人認為教會是在千禧年後被提,各執一理,不讓對方,爭個不休.
\newline
\newline
當時羅得要與亞伯拉罕相爭.如果亞伯拉罕用血氣與羅得相爭,事情就糟了,羅得屬肉體,亞伯拉罕對付了肉體,所以爭不起來.亞伯拉罕對他說:「遍地不都在你眼前嗎?請你離開我,你向左,我就向右,你向右,我就向左.」亞伯拉罕是沒有自己的人,他有的就是愛,他心中充滿神,有神就有愛.一個充滿愛的人,在教會中是講愛,不講理的.羅得是下輩,亞伯拉罕是長輩,若是講理羅得是沒理可言的,但亞伯拉罕對他不講理,只講愛.
\newline
\newline
愛是甚麼?為人捨己,願意吃虧之意.亞伯拉罕沒有肉體,所以別人的肉體引不動他的肉體.一個沒有自己的人頭腦冷靜,理智清楚,有大量的胸懷,不會用肉體與人相爭的.羅得在亞伯拉罕面前是何等的醜,亞伯拉罕是何等的美.
\newline
\newline
(三)亞伯拉罕拯救羅得
\newline
\newline
失敗的亞伯拉罕是卑賤的,現在他復興了,多麼美啊,羅得是一個無情無義的人,亞伯拉罕不計較他的惡,聖經說:「亞伯蘭聽見他姪兒被擄去,就率領他家裡生養的精練壯丁三百一十八人,直追到但......將被擄掠的一切財物奪回來,連他姪兒羅得和他的財物,以及婦女人民,也都奪回來.」(創十四14-16)雖然羅得待他不好,但他不理這些,還當他是骨肉,你待自己的肢體是這樣嗎?
\newline
\newline
亞伯拉罕受神對付後,心胸寬大了,在拯救羅得的事上打了屬靈的勝仗.當他打勝仗回來時,所多瑪王對他說:「你把人口給我,財物你自己拿去罷!他對所多瑪王說:「我已經向天地的,至高的耶和華起誓；凡是你的東西,就是一根線,一根鞋帶,我都不拿,......」(參看創十四22-23)他的心中除神之外,沒有別的所愛了.
\newline
\newline
亞伯拉罕打勝仗回來時,麥基洗德出來迎接他,對他說:「願天地的主,至高的神,賜福與亞伯蘭.」(創十四19)神在亞伯拉罕的身上,可稱為天主,又可稱為地的主,兩個都可以.甚麼叫地的主呢?神在天上,藉著亞伯拉罕把祂的旨意,榮耀彰顯出來.神的旨意能在你身上通行無阻嗎?神的榮耀能在你身上完全彰顯出來嗎?
\newline
\newline
神揀選亞伯拉罕做祂的器皿,從卑微變成尊貴,今天神也要在你身上變成像亞伯拉罕那樣貴重的器皿,使神的旨意和榮耀能在你身上彰顯出來.
\newline
\newline
盼望主在你身上,能成為天的主和地的主.
\newline
\newline


\newpage

\section{第40屆~港九培靈會~吳勇~第六講}
\label{sec:1542}
\textbf{人子近了}
\newline
\newline
連結: \href{https://www.hkbibleconference.org/session-message/view/1542}{\texttt{https://www.hkbibleconference.org/session-message/view/1542}}
\newline
\newline
\hyperref[sec:1543]{< < < PREV SERMON < < <}
~
\hyperref[sec:toc]{[返目錄]}
~
\hyperref[sec:1544]{> > > NEXT SERMON > > >}
\newline
\newline
創世記 4:17-4:24
\newline
\begin{longtable}{cl}
\hline
\hline
章節 & 經文 (和合本修訂版)\\
\hline
4:17 & \begin{tabularx}{0.7\textwidth}{X} 該隱與妻子同房，她就懷孕，生了以諾。該隱建造一座城，就照他兒子的名字稱那城為以諾。 \end{tabularx} \\ \\ \relax
4:18 & \begin{tabularx}{0.7\textwidth}{X} 以諾生以拿，以拿生米戶雅利，米戶雅利生瑪土撒利，瑪土撒利生拉麥。 \end{tabularx} \\ \\ \relax
4:19 & \begin{tabularx}{0.7\textwidth}{X} 拉麥娶了兩個妻子：一個名叫亞大，一個名叫洗拉。 \end{tabularx} \\ \\ \relax
4:20 & \begin{tabularx}{0.7\textwidth}{X} 亞大生雅八；雅八是住帳棚、牧養牲畜之人的祖師。 \end{tabularx} \\ \\ \relax
4:21 & \begin{tabularx}{0.7\textwidth}{X} 雅八的兄弟名叫猶八；他是所有彈琴吹簫之人的祖師。 \end{tabularx} \\ \\ \relax
4:22 & \begin{tabularx}{0.7\textwidth}{X} 洗拉又生了土八‧該隱；他是打造各樣銅器鐵器的工匠。土八‧該隱的妹妹是拿瑪。 \end{tabularx} \\ \\ \relax
4:23 & \begin{tabularx}{0.7\textwidth}{X} 拉麥對他兩個妻子說：亞大、洗拉啊，聽我的聲音；拉麥的妻子啊，側耳聽我的言語：大人傷我，我把他殺了；小孩損我，我把他害了。 \end{tabularx} \\ \\ \relax
4:24 & \begin{tabularx}{0.7\textwidth}{X} 若殺該隱，遭報七倍，殺拉麥的，必遭報七十七倍。塞特和以挪士 \end{tabularx} \\ \\
[1ex]
\hline
\hline
\end{longtable}
聖經中有講到主的降生,受難,復活,升天,再來五件事,但這五件事的中心是主再來.主再來是聖經預言的中心,神在每一個時代都有祂的說話.今日的時代,是黑夜已深,白晝將近的時代,正是主要來的日子.
\newline
\newline
主的門徒曾問及祂再來的預兆是甚麼?現在從三方面論述之:
\newline
\newline
(一)社會的預兆
\newline
\newline
聖經說:「挪亞的日子怎樣,人子降臨也要怎樣?」(太廿四37)今天的時代與挪亞的時代一樣的敗壞,挪亞時代被神用洪水毀滅了,在創世記第四章裡論到挪亞時代敗壞的六個光景.
\newline
\newline
1.建造城市,正如今日是城市的時代.五十年前,百分之七十的人是住在農村,但今天,卻是百分之七十的人住在城市.因此城市的人口增加,人心敗壞,思想盡都是惡.人口多,謀生不易.城市的人重利不重義,道德墮落,好像挪亞時代的人一樣.
\newline
\newline
2.拉麥娶了兩個妻(創四19).婚姻是神定的,一夫一妻制是神定的.那時代的人是如此,今日世代也如此,男的不是她的丈夫,女的不是他的妻子,因此就破壞了今日的社會,為滿足自己不正當的情慾,就變成了濫交的風氣,敗壞到極點了.
\newline
\newline
3.牧養牧畜(創四20).當時的人,勞碌辛苦,是為滿足自己的需要和享受.今天的人,希望明天比今天好,就用時間的三法制來換取:休息,工作,睡覺各用一日的三分之一.有的人操勞到患肺病,苦不堪言.甚至有的人還用身體和人格來換取享受,結果導致男盜女娼的社會,人格和道德敗壞到不可收拾.
\newline
\newline
4.猶八是彈琴吹蕭之人的祖師(創四21)那時代是喜好音樂的時代,好像今日的時代,人們從古典音樂到爵士音樂,從披頭士音樂到貓王音樂,到台上又叫又喊,真肉麻.無論是留聲機,收音機,電視,廣告,以及各種娛樂,無不用音樂來吸引人,挪亞時代的人嗜好音樂,今日的人也如此.
\newline
\newline
5.土八該隱是打造各種銅鐵利器的人(創四22).今天的時代,人類用銅,不鏽鋼來製造殺人武器了,發達到登峰造極了.從槍到大炮,從大炮到核子彈,飛機等都不鏽鋼來製造,非常犀利,比起挪亞時代,真有天淵之別.
\newline
\newline
6.殺人為可誇(創四23).經上說:「壯年人傷我,我把他殺了,少年人損我,我把他害了.」拉麥不但殺人,且以殺人為可誇.現代的人為債務殺人,為色情殺人,為爭吵殺人,好像張獻忠的七殺詩說的:「天造萬物以養人,人無一德以報神.殺殺殺殺殺殺殺.」今天的時代,真是隨便殺人的,有仇也殺,無仇也殺,簡直變成了一個殺人的世界.
\newline
\newline
挪亞的日子怎樣,人子的日子也要怎樣?
\newline
\newline
(二)教會的預兆
\newline
\newline
教會是主的身體,身體的目的是見證「裡人」,人的態度,動作都是介紹「裡人」.譬如你輕鬆活潑,是美國人,嚴肅莊重是英國人.
\newline
\newline
主的身體是藉著教會表現出來,主在地上建立的教會,撒旦很快的就來破壞.今日的教會,撒旦也來破壞,要他在地上不存在.使徒時代的教會,魔鬼用刀劍來殺害,但他失敗了,教會還是屹立不動.可是,今天的教會變質了.甚麼叫變質呢?在二百年前,人類不知牛奶會發酸,肉會發臭的原因,後來有巴斯德醫生研究出來,乃是細菌侵入的緣故.
\newline
\newline
撒旦侵入教會,教會就會變質.在約翰福音第四章裡,從主與撒瑪利亞婦人的談話就可以看出來,她是一個道德極敗壞的婦人.這婦人有禮拜,禮拜是宗教舉動.這婦人有宗教,宗教是講敬虔的,她的為人生活既如此敗壞污穢,怎談得上敬虔?宗教是不能沒有敬虔的,那麼,她的敬虔是甚麼敬虔呢?她說:「我們祖宗在這山上禮拜.」(約四20)祖宗代表遺傳.敬虔既不在人的生活上,而變成在一些遺傳的規條上,在一些地方,在一些儀文上,在一些節期上了.這是變質的宗教,也是墮落的宗教.
\newline
\newline
當時的法利賽人就是這種宗教的代表,他們被主責備.主責備他們說:「因為你們好像粉飾的墳墓,外面好看,裡面卻裝滿了死人的骨頭,和一切的污穢.」(太廿三27)法利賽人只是外表敬虔,他們有許多遺傳規條,因著主的門徒沒洗手吃飯,他們就來對耶穌說:「你的門徒為甚麼犯古人的遺傳呢?」(太十五2)他們在十字街頭禱告,以示敬虔,他們注重節期,注重遺傳,儀文,而不遵行神的旨意,行出來的是污穢敗壞,只有外表的敬虔,而沒內心的敬虔.主的降生,是在猶太教變了質的時候,同樣,今天的教會也與當時的猶太教一樣,由此可見,主要再來了.
\newline
\newline
(三)猶太人的預兆
\newline
\newline
經上記著說:「你們可以從無花果樹學個比方；當樹枝發嫩長葉的時候,......你們看見這一切的事,也該知道人子近了,正在門口了.」(太廿四32-33)無花果樹是預表猶太人.
\newline
\newline
當猶太人的南北兩個國家滅亡後,神就差遣耶穌降世了,主說:「我多次願意聚集你的兒女,好像母親把小雞聚集在翅膀底下,只是你們不願意.」(太廿三37)主來了,他們不接納祂,而且還把祂交在羅馬人手中釘十字架.按照羅馬的律法耶穌應該被釋放的,因為巡撫彼拉多也查不出祂有甚麼罪.換言之,祂是無罪的.問題在於文士和祭司聳動百姓,目的是要除滅耶穌.無知的猶太人釘耶穌在十字架,而且還承受殺耶穌的罪歸到他們子孫的身上,結果猶太民族,土地,國家都受咒詛.在主後七十年八月提多將軍攻入耶路撒冷,屠殺猶太人一百一十萬.後來有一人名叫巴米巴,起來反抗羅馬,結果羅馬派兵鎮壓,猶太人又被屠殺五十萬.以後歷世歷代猶太人受屠殺和驅逐,最後一次是希特勒的大屠殺,猶太人死六百萬.直到今天,以色列國共有猶太人二百七十萬,全世界共有一千多萬.
\newline
\newline
這個國家的民族,多次受屠殺,仍沒滅亡,所以有人說這民族優秀,民族意識堅強,領袖傑出,這樣說對嗎?雖亡國二千多年,但這民族仍然存在.其存在的原因是神在他們身上有一個永遠的計劃.猶太人在一九一三年開始發動復國,直至一九四八年五月十四日才正式復建以色列國.以西結書卅五章講到西珥與他們世代為仇,這件事從以阿的戰事上可以看出,以阿戰爭或衝突是世界大事的焦點,死海的寶藏值一千二百億美元,所以阿拉伯國要爭它.在一九六七年六月阿拉伯國家發動戰爭,十四國一億多人的聯盟國進攻二百七十萬人的以色列國,結果怎樣呢?阿拉伯國家被打敗,是甚麼原因呢?這是神的作為,要這民族為祂作見證.聖經說:「耶路撒冷要被外邦人踐踏,直到外邦人的日期滿了.」(路廿一24)去年的以阿戰爭,以色列直搗阿拉伯國家的心臟開羅,正如以賽亞書五十四章十七節所說的:「凡為攻擊你造成的器械,必不利用.」
\newline
\newline
以色列打了勝仗,收回了耶路撒冷舊城,到此「外邦人踐踏耶路撒冷的日期」是否滿足了呢?沒有,因為聖殿的地基上建了一座宏大的回教廟,「外邦人的日期滿足」是要到這個回教廟被拆毀,一定要拆毀,但時間還沒到.這間回教廟一旦拆毀,主就要回來!
\newline
\newline
主再來我是無法知道的,但是,我們從社會,教會,以及猶太人的許多預兆便可知道主快要再來了.所以,我們要警醒,等候祂再來.
\newline
\newline


\newpage

\section{第40屆~港九培靈會~吳勇~第七講}
\label{sec:1544}
\textbf{主怎樣來?}
\newline
\newline
連結: \href{https://www.hkbibleconference.org/session-message/view/1544}{\texttt{https://www.hkbibleconference.org/session-message/view/1544}}
\newline
\newline
\hyperref[sec:1542]{< < < PREV SERMON < < <}
~
\hyperref[sec:toc]{[返目錄]}
~
\hyperref[sec:1545]{> > > NEXT SERMON > > >}
\newline
\newline
使徒行傳 1:9-1:11
\newline
\begin{longtable}{cl}
\hline
\hline
章節 & 經文 (和合本修訂版)\\
\hline
1:9 & \begin{tabularx}{0.7\textwidth}{X} 說了這些話，他們正看的時候，他被接上升，有一朵雲彩從他們眼前把他接去。 \end{tabularx} \\ \\ \relax
1:10 & \begin{tabularx}{0.7\textwidth}{X} 他升上去的時候，他們定睛望天，看哪，有兩個人身穿白衣站在他們旁邊， \end{tabularx} \\ \\ \relax
1:11 & \begin{tabularx}{0.7\textwidth}{X} 說：「加利利人哪，你們為甚麼站著望天呢？這離開你們被接升天的耶穌，你們見他怎樣升上天去，他也要怎樣來臨。」 \end{tabularx} \\ \\
[1ex]
\hline
\hline
\end{longtable}
在使徒行傳一章裡講到主怎樣去,祂還是要怎樣來的,主再來分三個階段,茲述如下:
\newline
\newline
(一)主來到空中
\newline
\newline
這時候,我們就被提,如聖經所載,兩個人在田裡,取去一個,撇下一個.兩個女人推磨,取去一個,撇下一個.那一夜,兩個人在床上,要取去一個,撇下一個(參太廿四40-41,路十七34).有人被提,有人被撇下.在帖撒羅尼迦前書四章裡提到在基督裡死了的人要復活,活著還存留的人就要改變身體.聖經中記載兩種的聲音,一種是呼叫,只用過一次,就是主耶穌呼叫拉撒路從墳墓中出來.另一種是號筒聲,叫我們升到天上去,與主相會.這事只能意會,不能體會.好像主登山變像,衣裳放光,彼得看見就很快樂,不知要說甚麼才好.我們被提也是如此.
\newline
\newline
血肉之體不能承受神的國.今天是太空時代,以血肉之體不能上太空,因為太空沒有空氣,沒有壓力,如果上去的人不穿太空衣,即刻會死.我們自肉之體被提,就必須要改變身體.保羅說:「血肉之體,不能承受神的國；必朽壞的,不能承受不朽壞的.」(林前十五50)被提是霎時發生的事.現在的太空人升空,用火箭打上去,非常迅速.我們被提,不是自己去,乃是主的權能叫我們去.例如拍X光,拍好也不知!快到眼看不出.到那時被提的霎那間,世界各地均有很多人失蹤,失蹤的全是基督徒.
\newline
\newline
主到空中,第一件事是我們被提,第二件事眾人在基督台前要受審判.主來到空中,有幾種人要被撇下:1.不信的人.2.硬心的猶太人.3.教會中有名無實的人.
\newline
\newline
被提的基督徒在空中做甚麼?要受審判,審判要從神的家起首.這審判與得救無關,而與得賞有關.神要根據哥林多前書三章裡提及的兩種工程來審判祂的兒女,就是金,銀,寶石和草木禾楷的工程.這是基督徒的工作,前者有價值,後者無價值.保羅也有講到極重的榮耀,和極輕的榮耀.如輕的物質,放在水中,它浮起來,被人看見.重的物質,就沉在水中,不被人看同.草木禾楷的工作,雖被人看見,但沒有價值,不能存留到永世.金銀寶石的工作,有價值,雖不被人看見,但可存留永世.這裡還有一個解釋:今天我奉主的名工作,趕鬼,醫病,但主說,作惡的人離我去罷!今天的工作,分出兩個源頭來.當時,神叫亞伯拉罕的妻子生一個兒子,而她怎樣想呢?只要有一個兒子就行了.但是,惟有撒拉生的,神才接受,夏甲生的,神不承認.一個人為神作工,神要看是從那裡來的,若是從神來的,祂接受,從血氣來的,神不接受.從神來的是金銀寶石,從血氣來的是草木禾楷.你的工程在基督審判台前顯露的是金銀寶石,或是草木禾楷?又有人說,這種工程是對主忠心與善良而言,忠心是工作,善良是生活.人有偏差,不問私生活,只要有才幹就行了,一切無所謂.今日的人,沒有善良的心,不做慈善事業,不做屬靈的事,其實他們是敗壞了.然而,主要我們生活與工作一樣重要.在空中,在祂的審判台前,祂要審判我們的良善與忠心,要合格,才能得賞賜.
\newline
\newline
(二)主來到地上
\newline
\newline
主甚麼時候來到空中,人不知道,天使也不知,主自己也不知.就像賊來一樣,我們不知道.主來到空中,我們看不見,因為有人失蹤,所以就知道主已經來過.在撒迦利亞書十四章裡,即主降生前五百年,已預言到主來到地上要發生的事,是我們可以看見的:1.主的腳要站在耶路撒冷的橄欖山上.2.耶和華必作全地的主.在主的腳未踏在橄欖山的前一刻,耶路撒冷城要被攻陷.自從一九四八年以色列建國以來,阿拉伯國家常攻打以色列國.從一九五六年和一九六七年的兩次進攻以色列.就可看見阿拉伯國家與以色列國的仇恨之大.但是,阿拉伯國家反被以色列打敗,正如以西結書預言到他們不能攻下耶路撒冷一樣.
\newline
\newline
當主的腳一踏在橄欖山的時候,主要把攻擊以色列國的仇敵親自消滅,將他們的屍首丟給鷹吃.主要捉拿敵基督,假先知,撒旦,把他們扔在火湖裡.
\newline
\newline
主來到地上,有幾樣事要發生:
\newline
\newline
1.有聖者同來.猶大書說主要與千萬的聖者駕雲降臨,這些聖者是在天空審判,工程經火仍存留,被稱為忠心良善的人,他們要得國度的賞賜.他們要與主一同作王掌權.我相信掌權的程度有分別,在路加福音中講到賺十錠銀的人給他管十座城,其餘賺五錠銀的人就派他管五座城.神不偏待人,祂待人公平.
\newline
\newline
2.舊約的信徒復活.主來到空中,新約的信徒復活.正如經上所說:「從東到西,將有許多人來,在天國裡與亞伯拉罕,以撒,雅各,一同坐席.」(太八11)主來到地上,大衛帳棚重建.以色列國成為各國之首,耶路撒冷是世界的中心.以色列成為各國之首,並非猶太人的緣故,乃因耶穌是萬王之王.
\newline
\newline
3.萬物得贖.神造人管理萬物,因人犯罪,萬物受咒詛.當主來到地上的時候,萬物恢復原來榮耀,那時,「豺狼必與綿羊羔同居,豹子與山羊羔同臥；少壯獅子,與牛犢,並肥畜同群,小孩子要牽引他們.牛必與熊同食,牛犢必與小熊同臥；獅子必吃草與牛一樣,吃奶的孩子必玩耍在虺蛇的洞口,斷奶的嬰兒必按手在毒蛇的穴上.在我聖山的遍處,這一切都不傷人,不害物......」(賽十一6-9)
\newline
\newline
(三)主來到空中與地上之間
\newline
\newline
這就是神啟示給但以理的七十個七的預言,這乃是指神的聖民設的.其中兩個七和六十二個七都已經在歷史上應驗了.還有一個七正待應驗,這是末後的一個七,即教會和信徒被提的時期.
\newline
\newline
最末後的一個七分為兩半,即三年半.前三年半是災難時期,但並非大災難.後三年半是大災難時期.前三年半,世界的領袖堅定盟約,假仁假義,後來又背約,有人自稱為神,要人拜他,逼迫信徒.當時,地上掛名的基督徒,悔改信主了,他們不肯拜敵基督,也不拜獸,所以他們買不到食物,且要被殺,因為他們的頭上沒有印記.我們絕不可存僥倖的心理,現在不信主,到那時就太遲了.
\newline
\newline
後三年半是大災難的時期,耶路撒冷被大軍圍攻,且被攻下,或者你說不在中東,不要緊.但在啟示錄中講到末期的戰爭是非常厲害的,正如今日洲際飛彈之多,世界沒有一處是安全的地方,那時的人被戰爭所殺有三分之一,被瘟疫所殺有三分之一,情況非常恐怖,因為災禍非常之大,在災難之前凡被提的都是有福的.
\newline
\newline
主從空中到地上的階段,地上有七年的大災難.馬太福音廿四章裡講到那日「屍首在那裡,鷹也在那裡.」甚麼叫屍首呢?人的身體沒有了生命,死了,才叫屍首.今天在教會中有很多屍首,就是那些沒有生,掛名的基督.所以我們必須要在神的面前省察,自己是否是一個「屍首」?
\newline
\newline
我們在主再來時要被提,但是,被提與得救不同.你信神,就得救,神是一視同仁的.然而得神的賞賜就不同了.神是公義的,祂要根據我們的工作和生活決定的.在這裡我們必須要鄭重的思想兩個問題:第一,我們的工程是金銀寶石的呢,還是草木禾楷的呢?第二,我們是否是一個在教會裡掛名的基督徒?有一次,我去看一個牧師的兒子,他對我說他做了五十年掛名的基督徒.人可以騙得,但卻騙不過神.
\newline
\newline
你是那一種的人?請你自己深深的省察!
\newline
\newline


\newpage

\section{第40屆~港九培靈會~吳勇~第八講}
\label{sec:1545}
\textbf{請主來}
\newline
\newline
連結: \href{https://www.hkbibleconference.org/session-message/view/1545}{\texttt{https://www.hkbibleconference.org/session-message/view/1545}}
\newline
\newline
\hyperref[sec:1544]{< < < PREV SERMON < < <}
~
\hyperref[sec:toc]{[返目錄]}
~
\hyperref[sec:1546]{> > > NEXT SERMON > > >}
\newline
\newline
民數記 6:1-6:8
\newline
\begin{longtable}{cl}
\hline
\hline
章節 & 經文 (和合本修訂版)\\
\hline
6:1 & \begin{tabularx}{0.7\textwidth}{X} 耶和華吩咐摩西說： \end{tabularx} \\ \\ \relax
6:2 & \begin{tabularx}{0.7\textwidth}{X} 「你要吩咐以色列人，對他們說：無論男女，若許了特別的願，就是拿細耳人的願，願意離俗歸耶和華， \end{tabularx} \\ \\ \relax
6:3 & \begin{tabularx}{0.7\textwidth}{X} 他就要遠離清酒烈酒，也不可喝任何清酒烈酒做的醋；不可喝任何葡萄汁，也不可吃鮮葡萄和乾葡萄。 \end{tabularx} \\ \\ \relax
6:4 & \begin{tabularx}{0.7\textwidth}{X} 在一切離俗的日子，任何葡萄樹上所結的，甚至果核和果皮，都不可吃。 \end{tabularx} \\ \\ \relax
6:5 & \begin{tabularx}{0.7\textwidth}{X} 「在他一切許願離俗的日子，不可用剃刀剃頭。在離俗歸耶和華的日子未滿之前，他要成為聖，要任由頭上的髮綹生長。 \end{tabularx} \\ \\ \relax
6:6 & \begin{tabularx}{0.7\textwidth}{X} 在他一切離俗歸耶和華的日子，不可挨近死屍。 \end{tabularx} \\ \\ \relax
6:7 & \begin{tabularx}{0.7\textwidth}{X} 即使他的父母或兄弟姊妹死了，他也不可因他們使自己不潔淨，因為他頭上有離俗歸神的記號 。 \end{tabularx} \\ \\ \relax
6:8 & \begin{tabularx}{0.7\textwidth}{X} 在他一切離俗的日子，他是歸耶和華為聖的。 \end{tabularx} \\ \\
[1ex]
\hline
\hline
\end{longtable}
在舊約時代,大衛的兒子押沙龍反叛,聲勢浩大,但後來慢慢失勢,直到有一天押沙龍逃到樹林,因頭髮懸掛樹上,結果他被約押所殺.叛亂平息,大衛仍然留在城外,差人去見祭司撒督和亞比亞,說:「你們是我的弟兄,是我的骨肉；為甚麼在人後頭請王回來呢?」(撒下十九11-12)
\newline
\newline
上列的故事可以用來預表我們的主,我們是主的弟兄,教會是主的骨肉,正如夏娃是亞當的骨肉.大衛要他們請他回來,主也要我們請祂回來.關於請主回來的事,彼得後書三章十二節說:「切切仰望上帝的日子來到.」這句話英文譯作「催促祂的日子來到.」怎樣催促呢?
\newline
\newline
(一)一般人講的主再來
\newline
\newline
從時間上去看,有人說神六日造天地,第七日安息.可是在彼得書信裡,神看一日如千年,千年如一日.所以,六日即六千年.人類歷史從亞當到亞伯拉罕有二千年,從亞伯拉罕到耶穌有二年,又從耶穌到今年的一九六八年,總共是五千九百六十八年,意即尚有三十二年,就滿六千年,主就要再來.
\newline
\newline
有人從時事上去看,外邦日期越來越壞,現今雖然物質文明,享受越來越好,交通從福特汽車到噴射飛機,科學已從挖掘古代秘密到探測太空秘密,雖然科學非常進步,發達,但道德越來越壞.從青年去看,亞洲以阿飛做代表,歐洲以披頭四做代表,美洲則以頹廢派的嬉皮上做代表.這樣下去,下一代的社會簡直不堪設想.去年葛培理博士在日本,有一個國會議員問他:「今日美國嬉皮士很多,下一代的青年如何辦?將來的世界又如何辦?」這問題令葛培理不知所答.
\newline
\newline
現今外邦人的宗教越來越壞,人類最古之時是敬拜一神,後來變成多神,現今則變成無神.就以有神的人來說,敬拜神成了外面的空殼,有外表沒有實際.人壽保險是人的保障,宗教是犯罪的保障.
\newline
\newline
今日的政治越來越壞,第一次世界大戰完了,不久,二次世界大戰又重演了,結果教會,城市變成焦土,通貨膨脹帶來人類的貧窮,政治的鬥爭產生了極權主義者,在二次世界大戰後,很多國家陷落在鐵幕後面,人類失卻自由,人對神的盼望喪失了.戰爭的結果,製造了無數無家可歸的人,遺下數不清的寡婦孤兒.這些事並沒有叫人的理性恢復.第三次世界大戰,似乎指日可待,馬上就要發生似的.
\newline
\newline
現今的世代是無花果樹和各樹發芽的時代,無花果樹是替猶太民族,各種的樹是指世界各國民族.聖經說:「無花果樹發芽,其他的樹也發芽.」猶太國復國,其他國也復國了,例如在非洲,亞洲在二次大戰後,有許多新的小國建立.
\newline
\newline
(二)主要來,天要留祂
\newline
\newline
主再來,要到萬物復興的時候,祂才來,現在神把祂留下,這是先知的預言.神在聖經中把祂的真理擺了出來,在啟示錄中講人可增添或刪減.神已經將祂的聖言交給猶太人了,這是羅馬書告訴我們的.但是,猶太人怎樣持守神的聖言呢?在約珥書講到一件悽慘的事,他們將神的產業給蝗蟲吃,蝗蟲吃剩的,蝻子來吃,蝻子吃剩的,螞蚱來吃.他們把神的話給各種蟲吃,歷代以來,神的話受到多次的虧損,在第三世紀時,君士坦丁提倡政教合一,第五世紀時,天主教第一個教皇起來,吃了神的真理,教會沒有了生命,變成一套制度,儀式,和規條等.
\newline
\newline
真理一定要復興,神用祂的僕人,把損失埋藏了的真理再發掘出來.例如十五世紀的馬丁路德把因信稱義的真理發掘出來.又如雅各逃避哥哥以掃,走到曠野,在曠野睡覺,作一夢,神與他同在,起身後把油澆在石頭上,石頭代表人,油代表神,因而那地方叫伯特利,意即石頭澆油才是伯特利的材料,光是有石頭也不行.伯特利是代表教會.意即光是人不是教會的材料,光是神也不是教會的材料,而是要人裡頭有神才是教會的材料.
\newline
\newline
從歷史來看,真理發現了,還不夠,而要實行,這樣才叫真理復興.例如材料有了,屋要建起來,人才能進去住.真理發現了,就要去實行,這樣,主才再來.
\newline
\newline
今天教會由牧師負責,信徒奉獻金錢,用錢來請牧師,這不是教會真理.甚麼是教會真理呢?教會是身體,身體是肢體配搭起來的,這才叫教會真理.若這真理我們發現了,知道了,但我們只顧社會事業,而卻來不及照顧到教會的工作,所以很多人只肯做信徒,不肯做肢體,這樣不肯把真理實行出來,還不是復興真理,主要到真理復興才回來.
\newline
\newline
(三)神要器皿來復興真理
\newline
\newline
在聖經民數記第六章提到一種人,叫拿細耳人.他的心被神恩感,被神的愛激動,在神面前許願分別為聖歸給神終身做拿細耳人.舊約的拿細耳人就是新約羅馬書十二章第一節所說的「將身體獻上,當活祭」的人.換言之,舊約的拿細耳人即是新約的奉獻的人.他們是作神義的器皿.做拿細耳人有四件事:
\newline
\newline
1.清酒,濃酒都不喝.酒代表喜樂,不是神不叫我喜樂.問題是:喜樂可從神來,也可從魔鬼而來.所以一個奉獻的人,只可享受從神而來的喜樂.因為撒旦的喜樂,是叫人墮落,反叛神,這種喜樂可用金錢買得,亞干,猶大可為例.金錢又可買色情的快樂,如參孫,大衛是也.奉獻的人要拒絕從外面來的快樂,好像底馬一樣,為著貪愛世界就迷失了.
\newline
\newline
2.是留長頭髮的.聖經在哥林多前書十一章中講到女人留長頭髮是榮耀,男人留長頭髮是羞恥,而拿細耳人是男人.聖經講到羞辱有兩種:一是為罪之故,如參孫被非利士人戲弄.二是為神為人之故,如主釘十字架時,羅馬的兵丁戲弄祂,羞辱祂.故此,拿細耳人受羞辱,不是為自己,乃為神,為人.他所受的羞辱,他並不介意.
\newline
\newline
3.不可有污穢.葡萄不可吃,甚至連葡萄乾也不可吃.為甚麼呢?意即要澈底離開罪,不讓罪進入思想,五官,碰一下都不可.大衛犯罪就是讓罪進入五官.
\newline
\newline
4.不可摸死屍.死屍意即污穢.也不可摸死了家人的屍體,家人是代表親情.拿細耳人要住在聖潔中,好像參孫一樣,他聖潔時神與他同在.他犯罪時神離開了他.
\newline
\newline
今日的人,本來不犯罪的,但因親情之故,犯了罪.在尼希米記講到一件事,聖殿本來是祭司的,但因祭司以利亞實與亞捫族人多比雅結了親,就竟叫他在聖殿住下來.這個社會,尤其是中國人,更講究人情的.有一次,我兒子要考學校,有一位弟兄來見我,他是該校入學試的主考人,他叫我吩咐我兒子在試卷上作一個記號,他說一看見就通過.他在賣人情,我就責備他說:「你受主教導這麼久,連這個真理都不懂嗎?」
\newline
\newline
舊約時代有兩種人:一是先知,代神說話的人.二是士師,治理國家的人.今日教會亦須要有兩種人:一是傳福音的人,使神的兒女不致走入異端,迷途.二是治理教會的人,使神的教會不至荒涼,否則教會就要荒涼了.
\newline
\newline
主要來,神把祂留住,要到真理復興時,祂才來.所以神就會藉著祂的靈在你身上做工作,要用你做拿細耳人.神要用這種器皿,才能叫這個時代的真理得著復興.
\newline
\newline
你願意做神這種器皿嗎?
\newline
\newline


\newpage

\section{第40屆~港九培靈會~吳勇~第九講}
\label{sec:1546}
\textbf{基督徒的分辨}
\newline
\newline
連結: \href{https://www.hkbibleconference.org/session-message/view/1546}{\texttt{https://www.hkbibleconference.org/session-message/view/1546}}
\newline
\newline
\hyperref[sec:1545]{< < < PREV SERMON < < <}
~
\hyperref[sec:toc]{[返目錄]}
~
\hyperref[sec:1547]{> > > NEXT SERMON > > >}
\newline
\newline
使徒行傳 26:28-26:28
\newline
\begin{longtable}{cl}
\hline
\hline
章節 & 經文 (和合本修訂版)\\
\hline
26:28 & \begin{tabularx}{0.7\textwidth}{X} 亞基帕對保羅說：「你想稍微勸一勸就能說服我作基督徒了嗎？」 \end{tabularx} \\ \\
[1ex]
\hline
\hline
\end{longtable}
基督徒有兩種,一種是自己的,一種是神生的.貨物也有真假,鈔票也有真假,有時分得出,有時是分辨不出的.主的眼目如同火焰,我們瞞不過祂,祂能看見人的外面和內心.但是,魔鬼做敗壞的工作,主說,麥子和稗子到收割時要分出來的.人子在祂榮耀裡降臨時,要坐在榮耀的寶座上,萬民要聚集在祂面前,祂要把他們分辨出來,如牧羊人分出綿羊與山羊一般,把綿羊放在右邊,把山羊放在左邊；一是永生的,一是永死的.
\newline
\newline
現在我要與大家思想做的基督徒與神生的基督徒,兩者如何分辨出來.
\newline
\newline
(一)做的基督徒
\newline
\newline
基督徒是從信入門,但是,信也有程度的分別.雅各書說:「你信上帝只有一位,你信的不錯；鬼魔也信,卻是戰驚.」(雅二19)所以,做的基督徒與鬼魔的信是相同的.甚麼是鬼魔的信心呢?有一日,主到格拉森去,那裡有一個被鬼附著的人,看見耶穌,就拜祂,並且大聲說:至高神的兒子耶穌.鬼魔也信耶穌,相信的程度不同.做的基督徒與鬼魔的信心是一樣的,所以他就要留在定罪的地位上.甚麼是定罪的地位上呢?信耶穌,知道自己是一個罪人,但他沒為自己的罪懊悔,沒有從罪中出來.他未信前,他性情如何,信後還是如此.從前如何狡猾,惡毒,醉酒,好賭......現在還是一樣.
\newline
\newline
每個人做基督徒,必須知道要與世人不同,聖經說:「你們務要從他們中間出來,與他們分別.......」(林後六17)可是他們除了懂得一點敬拜儀式,如做禮拜,唱詩,禱告,和懂得一點聖經知識外,與世上的人沒有分別.比如一個女人,她說是基督徒,吃飽飯沒事做,就到鄰舍去,說長道短,批評論斷人,這種人與世上的人有何分別呢?又有一種人心胸狹小,又小心眼,別人講他一句話,他就受不了,雖然不當面計較,但卻在暗中報復.如果是一個男人,他說他是基督徒,社會上的人怎樣?他也怎樣,沒有分別,社會上的人講究面子,鋪張,排場,所以就虛偽應酬,沒有實際,與社會上的人毫無二致.今日的人講利不講義,為利損害親情,友誼,甚至自己的人格,道德也可以損害.
\newline
\newline
基督徒應當有靈與心同證為神的兒女,聖經說:「聖靈與我們的心同證,我們是上帝的兒女.」(羅八16)不必假設,有心證和靈證,這個「心」怎麼解釋呢?我們這個心,並非指跳動的心,乃是指感情的所在,一個人心證是神的兒女,他的感情就討厭恨惡自己和這個世界,他的意志願意從罪中出來,從世界分別出來.他的感情如何向著神呢?今日與人來往是用言語,同樣,與神的關係也是用言語,會像嬰孩愛奶一樣愛慕神的話,他的意志順服神的話,如亞伯拉罕順服神,願意獻以撒為祭一樣.又如馬利亞順從神,就是損害自己的利益也順從.如果一個基督徒是做的,對神就沒感情,他不會愛慕神的話,也不會遵行神的話.
\newline
\newline
甚麼是靈證呢?保羅說:「因為所賜給我們聖靈,將上帝的愛澆灌在我們心裡.」(羅五5)愛是恩典的解釋,神愛我們,賜恩典給我們.一個感恩的人,會常常感謝神的恩典.因著神賜給我們的是祂的獨生兒子耶穌基督,所以我們一生一世對神的態度就是要報恩.好像昔日大衛一樣,得約拿單的恩情,約拿單死後,大衛做了王,他就問:這個地方有約拿單的親人沒有?因為他要報恩,所以他就恩待約拿單的兒子米非波設,與他一同吃飯.靈證不是叫聖靈擔憂的,如果聖靈擔憂,心中就沒有平安.當你不應當行這事時,神的靈會反對你.當你妒忌時,心就會不平安.不是基督徒,他說謊不會心中不安的.一個沒有心證和靈證的基督徒是假的基督徒.
\newline
\newline
(二)神生的基督徒
\newline
\newline
聖經一開頭在創世記就講到生命樹.聖經結束,在啟示錄中也講到生命樹.所以生命是貫穿於聖經中.這生命是神的生命,永恆的生命,神要人得祂的生命,有相同的生命,才能與神相通,有永恆的生命,才能進入永恆.
\newline
\newline
生命是生長的,會生長才是生命,否則就不是生命.如把種子埋在土中,不久,它就發芽,長莖,長葉,這是生長.又如小貓生下來,不久眼開了,可以走動,這是生長.照樣,若是生的基督徒,也會生長.
\newline
\newline
甚麼叫生長?他的年紀與他的性格有關.如果一個三五歲的小孩愛哭愛鬧,到十幾歲還是如此,就不正常了.一個人生長到甚麼地步,他的志氣就怎樣.
\newline
\newline
基督徒的榜樣是誰呢?主耶穌是我們一生的榜樣.主耶穌來到世上,祂與哀哭的人同哀哭,表示同情人,而且祂還要幫助人.世人幫助了人,是要別人的報酬的,要得別人的代價.但是,主幫助人不要人的報酬,也不要人的代價,所以一個基督徒應當像主耶穌一樣.
\newline
\newline
一個基督徒生長到甚麼地步,他的力量和知識也就到同樣的地步.照樣,生的基督徒,長到基督的身量,讓基督本性榮美活在身上.從前邱吉爾檢閱海軍,有些軍艦上有人,有些沒有人,因為有偽裝,但邱吉爾是聰明人,他分辨得出那些艦有人無人.軍官問他怎樣分辨?他說,在艦上空有海鷗飛的就有人,否則就是沒有人.因為有人的軍艦有東西流出來,海鷗可以吃.一個基督徒有東西流出來,才是真基督徒的憑據.
\newline
\newline
約翰說神生的不犯罪,但有時還是會犯罪的,犯了罪怎麼辦呢?在神與人中間有一位中保耶穌,祂的血永遠有功效.凡被過犯偶勝的人,不致完全失望.我們要當心這句話,不是說有了一位中保負責,就可隨意犯罪了.有一次王明道先生說了一個比喻,有一個人到他朋友家去住,第二天早上他起來刮臉,不慎割傷了臉,他的朋友有一種藥,抹了很快就好.他很喜歡那瓶藥,那朋友把藥送了他,可是他回家後,刮臉時就亂劃,以為有特效藥抹傷口.試問:世間有這樣的人嗎?絕對沒有.而且,靈界也沒有這種人.
\newline
\newline
從神生的基督徒是不肯犯罪的,主看我們是祂的羊,偶跌泥裡,痛苦,呼叫,牠巴不得被救出來,重生的基督徒也是如此,偶然犯罪,痛苦不可言喻,盼望立即從罪中解救出來.聖經說:「藉著祂的靈,叫你們心裡的力量剛強起來.」(弗三16)從神生的仍能犯罪.未信主前,有個舊人,信主後,裡面有了一個新人.舊人是邪惡的,新人是聖潔的.因此,新舊兩人就相爭,新人是內裡的人,舊人是外面的人.這兩者,誰勝呢?舊人強,新人弱,舊人就得勝,於是你的靈便落在黑暗污穢中.如果新人得勝了,便進入光明聖潔中.我有個第七的兒子,身體強壯,他哥哥姐姐吃不下的東西,他拿去吃了,胃口很好,吃得多,身體就強壯.
\newline
\newline
神生的基督徒是要受神的管教的,在生命上成了聖潔.聖經說:「神又使祂成為我們的智慧,公義,聖潔,救贖.」(林前一30)生命是有他的本能的,豬喜愛污穢,羊喜歡乾淨.神又用真理教導我們聖潔,聖潔是神的生命.基督徒要逃避情慾的事,好像大衛,因他不肯逃避情慾,不肯聽神的管教,就與拔示巴犯了罪.如果有一種人,不肯聽神的管教或教導,又以犯罪,污穢的生活為快樂的,這種人你要當心.
\newline
\newline
從神生的基督徒,是願受神的管教的,他也不肯犯罪.做的基督徒就完全不同.所以,真假基督徒,可以在生命上將他分辨出來.請問:你是那一種基督徒?
\newline
\newline


\newpage

\section{第40屆~港九培靈會~吳勇~第十講}
\label{sec:1547}
\textbf{住在基督裡}
\newline
\newline
連結: \href{https://www.hkbibleconference.org/session-message/view/1547}{\texttt{https://www.hkbibleconference.org/session-message/view/1547}}
\newline
\newline
\hyperref[sec:1546]{< < < PREV SERMON < < <}
~
\hyperref[sec:toc]{[返目錄]}
~
\hyperref[sec:1548]{> > > NEXT SERMON > > >}
\newline
\newline
約翰福音 15:1-15:16
\newline
\begin{longtable}{cl}
\hline
\hline
章節 & 經文 (和合本修訂版)\\
\hline
15:1 & \begin{tabularx}{0.7\textwidth}{X} 「我就是真葡萄樹，我父是栽培的人。 \end{tabularx} \\ \\ \relax
15:2 & \begin{tabularx}{0.7\textwidth}{X} 凡屬我不結果子的枝子，他就剪掉；凡結果子的，他就修剪乾淨，使枝子結果子更多。 \end{tabularx} \\ \\ \relax
15:3 & \begin{tabularx}{0.7\textwidth}{X} 現在你們因我講給你們的道已經潔淨了。 \end{tabularx} \\ \\ \relax
15:4 & \begin{tabularx}{0.7\textwidth}{X} 你們要常在我裡面，我也常在你們裡面。枝子若不常在葡萄樹上，自己就不能結果子；你們若不常在我裡面，也是這樣。 \end{tabularx} \\ \\ \relax
15:5 & \begin{tabularx}{0.7\textwidth}{X} 我就是葡萄樹，你們是枝子。常在我裡面的，我也常在他裡面，這人就多結果子，因為離了我，你們就不能做甚麼。 \end{tabularx} \\ \\ \relax
15:6 & \begin{tabularx}{0.7\textwidth}{X} 人若不常在我裡面，就像枝子被丟在外面，枯乾了，人撿起來，扔進火裡燒了。 \end{tabularx} \\ \\ \relax
15:7 & \begin{tabularx}{0.7\textwidth}{X} 你們若常在我裡面，我的話也常在你們裡面，凡你們想要的，祈求，就給你們成全。 \end{tabularx} \\ \\ \relax
15:8 & \begin{tabularx}{0.7\textwidth}{X} 你們多結果子，我父就因此得榮耀，你們也就是我的門徒了。 \end{tabularx} \\ \\ \relax
15:9 & \begin{tabularx}{0.7\textwidth}{X} 我愛你們，正如父愛我一樣；你們要常在我的愛裡。 \end{tabularx} \\ \\ \relax
15:10 & \begin{tabularx}{0.7\textwidth}{X} 你們若遵守我的命令，就會常在我的愛裡，正如我遵守了我父的命令，常在他的愛裡。 \end{tabularx} \\ \\ \relax
15:11 & \begin{tabularx}{0.7\textwidth}{X} 「我已對你們說了這些事，是要讓我的喜樂存在你們心裡，並讓你們的喜樂得以滿足。 \end{tabularx} \\ \\ \relax
15:12 & \begin{tabularx}{0.7\textwidth}{X} 你們要彼此相愛，像我愛你們一樣，這是我的命令。 \end{tabularx} \\ \\ \relax
15:13 & \begin{tabularx}{0.7\textwidth}{X} 人為朋友捨命，人的愛心沒有比這個更大的了。 \end{tabularx} \\ \\ \relax
15:14 & \begin{tabularx}{0.7\textwidth}{X} 你們若遵行我所命令的，就是我的朋友。 \end{tabularx} \\ \\ \relax
15:15 & \begin{tabularx}{0.7\textwidth}{X} 以後我不再稱你們為僕人，因為僕人不知道主人所做的事；但我稱你們為朋友，因為我從我父所聽見的一切都已經讓你們知道了。 \end{tabularx} \\ \\ \relax
15:16 & \begin{tabularx}{0.7\textwidth}{X} 不是你們揀選了我，而是我揀選了你們，並且派你們去結果子，讓你們的果子得以長存，好使你們奉我的名，無論向父求甚麼，他會賜給你們。 \end{tabularx} \\ \\
[1ex]
\hline
\hline
\end{longtable}
在這幾節聖經裡有十二個「常」字,這「常」字應翻作「住」字.住在主裡是基督徒的生活秘訣,我們以浪子作喻,他住在家中,一切都豐富.不在家住,則貧窮.同樣,住在主裡的基督徒是豐富的,這是基督徒的生活秘訣.
\newline
\newline
有一次我領一個大學團契主辦的大會.散會後,有一個教授與我握手,發現我發高燒,他說明天的聚會若不能領則由他代替.第二天早上起來,我的高燒還沒退,我很難過.晚上我去講道之前,我跪下禱告,主就給我一節聖經,在以弗所書一章三節裡:「祂在基督裡,曾賜給我們天上各樣屬靈的福氣.」禱告完,站起身,一身輕鬆,當晚照樣去講道.戴德生時常追求得勝,一次他寫信給他的姐姐,告訴她他得著得勝的秘訣了,就是約翰福音第十五章的真理,樹在枝上,枝在樹上,但枝子必向樹求取水份,養份,同樣,住在主裡就會有水份,養份,可以得生,這就是基督徒的生活秘訣.
\newline
\newline
生活在主裡的基督徒有各樣的福氣.
\newline
\newline
(一)結果子
\newline
\newline
聖經說,凡屬祂不結果子的人,即是無靈的人,在加拉太書五章裡講到聖靈結的果子,果子是由靈結出來的.一個人屬耶穌,又不結果子,就是一個沒有靈的人,也是沒有生命的人,沒生命就不會有屬靈的生活了,當然沒有屬靈的果子結出來.即使做基督徒很久,常守禮拜,也是掛名的,實際上是死的.結局如何呢?正如聖經所說,不結果子的枝子要剪去,同樣,不結果子的基督徒,祂也要剪去.
\newline
\newline
屬主的人,要結起碼的果子.聖經說:「你們要結出果子來,與悔改的心相稱.」(太三5)即是說要結出悔改的果子來.現在我要講兩個比喻:第一個是浪子的比喻,他離家後,有一天他醒悟過來,悔改了,他起來,離開污穢,離開罪,追求神,沒有空虛了.第二個例子是撒該悔改歸主的事,他悔改後對主說:「我把所有的分一半給窮人,訛詐誰賠他四倍.」從前他愛錢財,現在棄掉錢財,他的心思都改變了.以前以罪為樂,現在以罪為痛苦.以前屬地,現在屬天了.他結了起碼的果子,就是悔改的果子,他起來離開罪,醒悟沒有神的空虛,現在起來追求神了,現在他看重屬靈的事了.
\newline
\newline
一個人結出起碼的果子,並不能令主滿足,祂要幫助我們,修理我們,目的是使我們能多結果子.主如何修理呢?祂以事情不如意,處處受挫折和打擊,打擊我們的驕傲,獨去獨來,武斷,叫我們尊重神,尊重別人.主有時叫人處處打擊我們暴躁的脾氣,而結出溫柔的果子.信徒雖然有平安喜樂,但也有悲傷痛苦.目的是叫我們多結果子,所以祂必須修理我們.如雅各一生不知經過多少次的修理,神用以掃,拉班來修理他,甚至到了晚年,神還用他兒子約瑟來修理他,所以他經過這多次的修理,就結出這麼多榮美的果子.
\newline
\newline
(二)祈求就給我們成就
\newline
\newline
聖經說,奉主的名祈求,必蒙成就,為甚麼有人奉主名求,不蒙應允呢?原因乃是為自己祈求.例如以利亞逃避耶洗別,在羅騰樹下求死的禱告神沒有應允他,因他正落在自己當中.又如西庇太的妻子為兩個兒子祈求一個坐在主左邊,一個坐在右邊,主說,「他不知所求的是甚麼?」
\newline
\newline
落在自己裡面的禱告,神不能成就.約翰壹書五章十四節說:「我們若照祂的旨意求甚麼,祂就聽我們.」不照祂的旨意求,祂就不聽.如有雷子稱號的約翰,求神降火下來燒滅不信的人,這不是神的旨意,主要人,非滅人,故此主不聽他的祈求.這就是妄求,求也沒用,主不會成就.
\newline
\newline
主說住在祂裡面,主的話就在我們裡面,住在主裡面的人,主就要在他裡面向他說話,比如你為一個病人禱告,若是你與某弟兄有怨恨,先去與他和好,除去攔阻,主才會聽你的禱告.又比如你跪下為一件工作禱告,主說,我呼喚你不聽,這樣你呼喚我也不聽.這樣,我們就應當省察了,有甚麼阻隔,求主光照出來.祂要幫助一個人,要你去做,去幫助一個人,主會幫助你,祂會成就你的禱告.主的工作要人同工,如水變酒的事,主要叫人挑水倒在水缸中,祂要人同工.如神要滅所多瑪時,神記起羅得一家,神要救他們,但要人同工,所以神就要亞伯拉罕為所多瑪城禱告.
\newline
\newline
(三)父愛我們
\newline
\newline
住在主裡面就是住在愛中,撒旦在我們裡面毫無作為.耶穌在約但河受洗後,神說:「這是我的愛子,我所喜悅的.」後來祂在曠野四十晝夜,很飢餓,又口渴,結果,撒旦就極力試探主,要祂與神有距離.又如大衛被掃羅追趕,幾乎無路可走,他以為這一生完了,那裡知道當時非利士人犯境,結果掃羅放下大衛的事.神愛大衛,神也愛我們,祂保護我們,叫我們忍受得住.聖經說:「誰能使我們與基督的愛隔絕呢?難道是患難麼,是困苦麼,是逼迫麼,是飢餓麼,是赤身露體麼,是危險麼,是刀劍麼.」(羅八35)有一年我到一個地方領會,剛好那裡的一間教會牧師辭了職,他們執事會開會,決定要再請一個傳道人,他們提出三個條件:1.要中年的.2.要曉得那裡的土話.3.要熟乎那裡的風土人情.因此,有個長老對我說,決定請我牧養他們的教會.這件事表面上我是答應了,但我還要問過主才作最後決定.後來回家後,跪下禱告求問主,主說:「若不愛我勝過愛自己的父母,妻子,兒女,弟兄,姐妹,和自己的性命,就不能作我的門徒.」(路十四26)因為在那地方有我父母,主不許可我去.所以我就回覆那長老,主不答應我去.可是,那長老就去叫父親來說服我,我的確很難做.我父親要我養他的老,而我對他說,我可以接他去台灣住,他又不願意.最後我對他說:「爸爸,對不起你,神不答應.」剛好,他手中拿著一碗飯,摔在地上,入房關門大哭,那時晚我去了北馬丁加奴領會.直到第五天,會畢後,我出來看見一排人穿著白衣服,站著不出聲,牧師交給我一張電報,是弟弟打來的,電報說:父喪出葬.我馬上趕回家去,我問及父親如何死的?傷心而死.今天我仍有不孝的罪.但在主裡面,祂把愛賜給我,又保護我.沒有甚麼可以叫我與主的愛隔絕的,祂賜各樣屬靈的福氣給我.祂引導我走義路,使我得勝.
\newline
\newline
在神的愛中,就一定要遵守祂的話,祂的話就是命令.我結婚不久,為生活之故,離家外出,走了兩天的路,在路上有一個人截住我說:「你是吳勇嗎?」我以為是要捕捉我的,原來他拿了一封電報給我,電報說:妻病重.結果,馬上掉回頭趕回去,兩天零一夜的路用一天時間走完,回到家,腳痛不能動,都爛光了.為何路上不覺得痛呢?因為有愛妻之心.同樣,我們今天要遵行神的命令,就要犧牲,付代價.如亞伯拉罕在山上獻以撒之事,的確不是容易的.他多年的盼望要毀在自己的手下,但他願意順服神,他內心有神的愛,所以他才樂意犧牲一切為神,遵行神的旨意和命令.
\newline
\newline
今天,神要叫你順服甚麼呢?要我們彼此相愛,愛不能愛的人.愛那些不好的,傷我們心的,說我們長短的人.神要我們多結果子,領人歸主.有一個傳道人的太太被邀去探望一個朋友的兒子,他正在醫院患腦膜炎留醫.這種病會傳染的.到了醫院,這病人的母親也不敢進去,還戴著口罩,傳道人的太太看見也怕了,但主對她說:「凡愛惜自己生命的,必要失喪生命......」於是她膽壯了,進去用嘴對著他的耳邊大聲傳福音,結果他得救了,結了果子.愛主就會勝過一切的懼怕,愛主的人就會得著主賜的各樣福氣.
\newline
\newline


\newpage

\section{第40屆~港九培靈會~吳勇~第十一講}
\label{sec:1548}
\textbf{末世教會}
\newline
\newline
連結: \href{https://www.hkbibleconference.org/session-message/view/1548}{\texttt{https://www.hkbibleconference.org/session-message/view/1548}}
\newline
\newline
\hyperref[sec:1547]{< < < PREV SERMON < < <}
~
\hyperref[sec:toc]{[返目錄]}
~
\hyperref[sec:1539]{> > > NEXT SERMON > > >}
\newline
\newline
啟示錄 2:1-2:4
\newline
\begin{longtable}{cl}
\hline
\hline
章節 & 經文 (和合本修訂版)\\
\hline
2:1 & \begin{tabularx}{0.7\textwidth}{X} 「你要寫信給以弗所教會的使者，說：『那右手拿著七顆星，在七個金燈臺中間行走的這樣說： \end{tabularx} \\ \\ \relax
2:2 & \begin{tabularx}{0.7\textwidth}{X} 我知道你的行為、勞碌、忍耐，也知道你不容忍惡人。你也曾察驗那自稱為使徒卻不是使徒的，看出他們是假的。 \end{tabularx} \\ \\ \relax
2:3 & \begin{tabularx}{0.7\textwidth}{X} 你能忍耐，曾為我的名勞苦而不困倦。 \end{tabularx} \\ \\ \relax
2:4 & \begin{tabularx}{0.7\textwidth}{X} 然而，有一件事我要責備你，就是你把起初的愛心拋棄了。 \end{tabularx} \\ \\
[1ex]
\hline
\hline
\end{longtable}
啟示錄第二和第三章講到七個教會,第一個是以弗所,最後一個是老底嘉.七個教會代表七個時代,老底嘉教會是最後的教會,代表末後教會的光景.有以弗所教會才會有老底嘉教會.
\newline
\newline
聖經說:「你要寫信給老底嘉教會的使者,說,那為阿們的,為誠信真實見證的,在神創造萬物之上為元首的.」(啟三14)講明自己的身份,自己介紹出來,自己是阿們,誠信真實的.阿們是誠實.誠信是真實.簡單的說,誠信與真實就是阿們,表明自己是絕對實在的.
\newline
\newline
有甚麼事是絕對實在的呢?神的應許是絕對實在,祂有赦罪的應許,人跌倒,犯罪,可以到主面前悔改認罪.主有看顧我們的應許,落在苦難中,神會看顧我們.不是憑感覺,乃憑聖經,是絕對實在的.
\newline
\newline
神的警告也是絕對實在的,我們應當放在心中.有一天我去訪一位做了五十年掛名的教友,不能分辨真假.神賜智慧給所羅門,叫他分辨妓女的兒子出來.馬太福音第七章講到好樹不能結壞果子,壞樹不能結好果子.我們從果子可分別樹,從生活可分辨人.
\newline
\newline
約書亞時代,攻取艾城之前,神吩咐不可取當滅之物,但其中有個士兵,名叫亞干,看見金錢就貪心,取了當滅之物.他以為人看不見,但神看見了.神的警戒是真實的,結果查出亞干來,被人用石頭打死在亞割谷.同樣神警戒教會,也是絕對實在的.祂創造萬物,祂在萬物之先,為萬物的元首.祂不但是時間內的神,而且也是時間外的神.祂對猶太說:「沒亞伯拉罕之先,就有了我.」又如保羅對腓立比的教會說:「祂那超乎萬名之上的名,叫一切在天上的,地上的,和地底下的,因耶穌的名,無不屈膝.」(腓二9-10)猶太人不認識祂,他們用拳頭打祂,傷害祂,吐唾祂,羞辱祂,最後還把祂殺害.今天的人對祂也如此,看基督與他們無關.
\newline
\newline
猶太人不接待祂,祂復活後,神立祂為基督,為主.祂是救主,有一切的權柄.祂吩咐的,我們不行.祂警戒的,我們不對付.這樣待祂,祂在我們身上就沒有權柄了.
\newline
\newline
主責備老底嘉教會,主說出祂的身分,祂的眼目如同火焰.今天,有很多人裡外不一樣的,外面是溫柔,裡面是殘暴；外面愛神,內心是愛自己,外面是君子,裡面是小人.如亞干就是這種人.
\newline
\newline
主說:「我知道你的行為,你也不冷也不熱；我巴不得你或冷或熱.」(啟三15)又說:「你既如溫水,也不冷也不熱,所以我必從我口中把你吐出去.」(啟三16)不冷不熱是溫的.主巴不得他或冷或熱,意即冷是不信,熱是真信.不冷不熱是信又是不信.你說他不信,他又有做禮拜,唱詩,捐錢,禱告.你說信嗎?他又與污穢的生活同居,說污穢的說話.這種人比不信的人更不如,因為不信的人,聽見主的話,覺得自己是個罪人,害怕了,他會悔改,信主,知耶穌是救世主.不冷不熱,是溫水的人,掛名的基督徒,他們屢次抗拒聖靈的感動.可是,主說要把他吐出去.
\newline
\newline
他們說發了財,很富足.其實,他們是痛苦的,貧窮的.但他們不知自己是貧窮,可憐的.為何說發財富足呢?用外面的東西代替,有大禮拜堂,有身分的人做執事,有大的唱詩班.質如何不管.只要信徒盡本分做禮拜,做一點禱告,讀一點聖經就夠了.為人如何?與神關係如何?一概不管.今日廿世紀的人,裝璜好,裡面賣甚麼東西不管.同樣,今日教會落在老底嘉教會的光景,不但自滿,而且還自欺.說貧窮,實在貧窮.甚麼都沒有,他還說有,不虛心,有惡心,不滿人的責備.
\newline
\newline
主勸老底嘉教會要買三樣的東西:
\newline
\newline
一,買火煉的金子
\newline
\newline
彼得書信說:「信心如同不能朽壞的金子.」(參彼一7)當時,教會失去了信心.教會屬靈上貧窮,要藉信心才能領取信心.他們失去了屬靈的恩賜,如治理,醫病,牧養,禱告的恩賜.奉主的名祈求,就必得著.失去信心,休想從主那裡得著甚麼.老底嘉教會失去信心,很可憐,他們是赤身的,即羞恥.按道理,神的兒女應感化社會,可是,社會卻感染了神的兒女.一個人貧窮,有金子就富足了.屬靈貧窮,信心就沒有了.
\newline
\newline
二,買白衣穿
\newline
\newline
用來遮蓋羞恥.聖經中講到兩種衣服,一是浪子的衣服,浪子回家時,父親把袍子給他穿,代表基督的義.到神的家,因信稱義.是白白得的,一信即得.一是信徒所穿的潔白細麻衣,這細麻衣就是聖徒所行的義.這就是約翰福音第一章所講的道成肉身,耶穌用自身來說明,把神的道表明出來.但凡事上在愛中不可虧欠.從前法蘭西斯離世前,叫他的學生,鄰居到他跟前來,要求他們饒恕,他說:「我在世界上有許多對不起人的事,我知道,愛有虧欠.」哥林多前書講到愛是恆久忍耐的.如果愛心不夠,基督的香氣在我們的身上不能發出來.當時老底嘉教會沒有信徒所穿的義,同樣,今天的教會也是如此.
\newline
\newline
三,買眼藥
\newline
\newline
老底嘉教會是瞎眼的,只看見現在,沒看見將來的榮耀.有一個人到非洲去傳道,年老退休了.他回國時與美國總統同船,船抵岸時,文武百官來迎迓總統回國,這個傳道人看見人家如此威武,心中就很不平,但他的妻子責備他,他便跪下禱告:「我一生為你,還落到如此光景.」但他心中立即有一個聲音對他說,你將來有榮耀了.很多人沒有看見將來的榮耀,如老底嘉教會一樣,只注重地,不注重屬天.
\newline
\newline
神為甚麼把以弗所教會放在第一,把老底嘉教會放在最末後呢?因以弗所教會造成老底嘉教會,它是他們不冷不熱,失卻信心的老底嘉教會的總因.以弗所教會的光景,把起初的愛心丟棄了.有如一個青年男子,交上一個青年女子.他的眼放在她身上,心也放在她身上.愛主的人,他的心,他的眼也放在主身上.
\newline
\newline
以弗所教會失去了愛,不追求主,放鬆了.今日的教會已經落到以弗所教會的光景.追求主放鬆,常是信主多年的人,例如很多信主多年的人,從前他們與主關係甘甜,親密,禱告句句動心,現在就沒有這種光景了.主說的話,不是當作對自己說的,以為是對別人說的,禱告也不是用心靈的交通,而是普通的習慣.放鬆屬靈的光景,老底嘉的光景就會臨到,不冷不熱,失去信心,變成貧窮.
\newline
\newline
放鬆屬靈的追求,剛信主的人不會發生,因他幼稚,甚麼都不知道,他有追求的心.但信主三,五年後,他老練了,甚麼都知道了,他就會放鬆屬靈的光景.
\newline
\newline
你要當心,屬靈的光景,對於你來說是非常重要,求神保守你不放鬆,保守你的屬靈光景,否則你就會落在老底嘉教會的光景.
\newline
\newline


\newpage

\section{第40屆~港九培靈會~吳勇~第十二講}
\label{sec:1539}
\textbf{教會的路}
\newline
\newline
連結: \href{https://www.hkbibleconference.org/session-message/view/1539}{\texttt{https://www.hkbibleconference.org/session-message/view/1539}}
\newline
\newline
\hyperref[sec:1548]{< < < PREV SERMON < < <}
~
\hyperref[sec:toc]{[返目錄]}
~
\hyperref[sec:1538]{> > > NEXT SERMON > > >}
\newline
\newline
馬太福音 18:15-18:20
\newline
\begin{longtable}{cl}
\hline
\hline
章節 & 經文 (和合本修訂版)\\
\hline
18:15 & \begin{tabularx}{0.7\textwidth}{X} 「若是你的弟兄得罪你，你要去，趁著只有他和你在一起的時候，指出他的錯來。他若聽你，你就贏得了你的弟兄； \end{tabularx} \\ \\ \relax
18:16 & \begin{tabularx}{0.7\textwidth}{X} 他若不聽，你就另外帶一個或兩個人同去，因為『任何指控都要憑兩個或三個證人的口述才能成立』。 \end{tabularx} \\ \\ \relax
18:17 & \begin{tabularx}{0.7\textwidth}{X} 他若是不聽他們，就去告訴教會；若是不聽教會，就把他看作外邦人和稅吏。 \end{tabularx} \\ \\ \relax
18:18 & \begin{tabularx}{0.7\textwidth}{X} 「我實在告訴你們，凡你們在地上所捆綁的，在天上也要捆綁；凡你們在地上所釋放的，在天上也要釋放。 \end{tabularx} \\ \\ \relax
18:19 & \begin{tabularx}{0.7\textwidth}{X} 我又實在告訴你們，若是你們中間有兩個人在地上同心合意地求甚麼事，我在天上的父必為他們成全。 \end{tabularx} \\ \\ \relax
18:20 & \begin{tabularx}{0.7\textwidth}{X} 因為，哪裡有兩三個人奉我的名聚會，哪裡就有我在他們中間。」 \end{tabularx} \\ \\
[1ex]
\hline
\hline
\end{longtable}
有一次,我在台灣與一弟兄交通,談及社會和教會的問題,就如社會來說,娛樂地方多,特別是夜總會,生意好,場場客滿,為甚麼呢?有人說是經濟繁榮的現象,進去的是有錢人.他們進去是刺激情慾,並非袋中有錢.其原因是精神苦悶,空虛,用刺激麻醉的方法,來解決精神的問題.今日的人走這條路,因此道德越來越下流,人的情慾如毒蛇.有一次,我到巴西,參觀一間毒蛇研究所.我們進去,見到滿地都是毒蛇,爬來爬去,我見到就混身發麻,腳也發抖,我再也不能行了,我害怕.所長說:「動物有一種性情,你惹他,他惹你,不惹他,他也不惹你.」他馬上試給我看,他踢牠一下,牠馬上跳過來咬他.同樣,情慾也是如此,你動他,他就咬你,如洪水泛濫,那時貞操,廉恥也不顧了,用情慾來刺激自己.以前中國女子,對丈夫不貞之事,不敢過問,逆來順受,把希望寄託到兒女的身上去.可是,今天的女子就不同了,不能忍受丈夫的不潔,於是就造成家庭和兒女的分裂問題,可是兒女的問題會牽連到社會的問題.
\newline
\newline
教會對內給信徒要培靈,靈培不起來,就沒生活見證,不信前怎樣做人,信徒還是如此.沒愛心,也沒有愛主的心,對世界,對自己不能捨.若基督徒對世界,對自己不能捨,那他就不能跑屬靈的道路.
\newline
\newline
教會對外佈道,無法進行,你講天堂不希奇,講地獄不怕,講罪就煩,講神的愛無動於衷.以前教會多少人現在還是照舊,人數不會增加.按照主對教會的使命,要做鹽,做光.教會要有功用,如鹽可防腐,防止道德的腐化.光是照亮黑暗的,黑暗代表痛苦,教會的功用是要除去人的痛苦.教會的功用發不出,因此,我與這個弟兄交談了很多這類的問題,他很灰心,很失望.社會如此,教會也是如此,接著他又問:「我們今日的路到底如何?」因此我就想到今天應走的路,就想到古代三個歷史.
\newline
\newline
一,尼希米的歷史
\newline
\newline
有人從耶路撒冷回來,告訴尼希米,耶路撒冷的城牆已經破爛不堪,門也燒了,以色列人遭大能,受凌辱.尼希米聽到後,要一心救以色列人,他的路是甚麼呢?他不是用人的辦法,乃用禱告來成就.
\newline
\newline
二,以斯帖的歷史
\newline
\newline
哈曼得到權柄,命令他的部下,在亞達月十三日要殺滅全國的猶太人,同時也要把他們的財產奪來佔為己有.以色列人見到這光景,以色列人活的路在那裡呢?聖經說,末底改叫以斯帖出來救以色列人,以斯帖要求全國的猶大人禁食禱告,求神開路.
\newline
\newline
三,希西家的歷史
\newline
\newline
亞述國十八萬大軍包圍猶大國,敵人勢力很大,沒有指望.但希西家不用政治,軍事,開會來解決,乃去找以賽亞,要他向神禱告.希西家的路就是禱告.
\newline
\newline
那弟兄說,社會,教會是如此,那麼今天的路是禱告.從尼希米,以斯帖,希西家三人身上,一切的路都是禱告.今天教會的路是禱告,不是辯論,不是研究教會的辦法.唯一的路是禱告.很多弟兄姊妹知道這條路,但教會的培靈,佈道都不行,對內對外沒有功效,結果有了懷疑,以為此路表不通.
\newline
\newline
今天教會的路仍是要禱告,以利亞為例.列王記上十七,十八章講到天不降雨,地上一片飢荒,長不出東西來.今天的人,他們的心靈也是如此,屬靈一片飢荒,長不出東西來.當時有兩種屬靈飢荒的人:一是俄巴底亞,他是亞哈王家宰.當亞哈王殺神的先知時,他熱心為神做事,他把神的先知五十個放一處,又將別外五十個放一處.表面看來,他很靈,很愛神.他起來保護神的先知.亞哈的王朝犯罪作惡,背逆神,怎不叫人痛心呢!但是,俄巴底亞不敢出來講一句話,如果他那樣做,職位要解除,頭也要搬家.他的一切做法是為己為利.他服事的不是神,而是亞哈,對神沒頁獻,屬靈的人是如此的嗎?二是亞哈王,一片飢荒時,他不是找神,而是找草.今天多人去找金錢,地位,宴樂,交情,不是去找神.亞哈時代敗壞,今天的教會裡,很多人也是如此.亞哈去找草,為牲畜,今日的世人,甚至是基督徒,是為自己,不是為別人.當亞哈見到以利亞時,他說:「使以色列遭災的是你嗎?」以利亞回答他:「是你,不是我.」他不承認,他抵擋,他把責任推到別人身上去,要推在以利亞身上去.
\newline
\newline
今日的時代,敗壞到如此程度,要挽救過來,這條路是甚麼呢?以利亞到迦密山,禱告後,火就從天降下來.又如昔日摩西在曠野看見著火的荊棘,荊棘代表以色列人,火代表神,這個異象表示神與以色列人同在.在五旬節時,火落在各人的頭上,他們禱告時,神就顯現,禱告有功效.
\newline
\newline
以利亞禱告,神降下火來,做了潔淨的工作,百姓喊說:「耶和華是神」.人心歸向神.這把火把巴力的先知盡被殺在基順溪旁,接著還有雨水降下,飢荒解除了.
\newline
\newline
今天,基督教注重講道,組織,不是注重禱告.這樣的教會,不會見到有火降下來,也不會見到神的顯現.所以,要禱告.要怎樣禱告呢?以利亞禱告之前,前先用十二塊石頭築壇,十二塊石頭代表十二支派,南北不分裂,原是南北國分裂的,現在放在一起,表示合一.因此,今日的禱告,要同心合意的禱告才有功效.詩篇一百三十三篇說:「看哪,弟兄和睦同居,是何等的善,何等的美.這好比那貴重的油,澆在亞倫的頭上......」用禱告除去不合一,污穢的罪,妒忌的罪.
\newline
\newline
聖經說:「凡你們在地上所捆綁的,在天上也要捆綁；凡你們在地上所釋放的,在天上也要釋放.」(太十八15)這是禱告的功效.在馬太十八節的「同心合意的求甚麼事」,原文是在「一切事情上,同心合意.」所以同心合意的意義並非指一個目標,乃在一切事上同心合意.例如開一個佈道會,目標是共同的,願萬人得救,不願一人沉淪,這是主的命令,為此事禱告,求主帶五百人聽道,二百人得救.但結果聽道的人只有一百人,得救的不到十人.為甚麼呢?不是在一個目標上同心合意.要在甚麼事上同心合意呢?1.在真理上要同心合意,聖經中有的就接納,沒有的就不接納.今日的人錯就錯在這裡,聖經中沒有還接納下來.2.感情上要同心合意,要有共同的目標,要有共同的步伐.感情不合一,步伐就不會整齊了.當心感情會受傷的.我有一位同工,他待我像約拿單待大衛一樣,有一次他去探訪一個有病的獨身姊妹,事後被人閒言.我知道後便對他說,對這種事要有智慧,不可給魔鬼留地步,他誤會我的意思,結果感情破裂了.
\newline
\newline
今日的教會,真理上和感情上都不合一,因此禱告就沒有功效.今天的教會要怎樣做才能合一呢?教會的人,都要起來同心合意,才能影響教會,不合一就不能影響教會.要起來挖溝,溝越大,越深,容量就大,神的恩典就會更大.今日帶領教會的人,對神,對人,對自己的問題挖出來,容量大,神恩典充滿更大.
\newline
\newline
以利亞獻祭時,將祭牲切成塊,表示謙卑對付罪.當拿單先知責備大衛的罪時,大衛謙卑的認罪悔改.一個謙卑的人才能對付罪.合一的工作要有謙卑的心.又如保羅與馬可同工,曾經分開,但後來又彼此同工,因為兩人都有謙卑的心,所以能合一.
\newline
\newline
今天教會的路是禱告,要禱告才通.並非為一個目標,乃是在一切事上同心合意禱告,才有功效,才能改變教會.我們要向神呼求,叫今天的教會能影響社會,不是教會受社會的影響,求神賜福給今天的教會.
\newline
\newline


\newpage

\section{第40屆~港九培靈會~吳勇~第十三講}
\label{sec:1538}
\textbf{生命的江河}
\newline
\newline
連結: \href{https://www.hkbibleconference.org/session-message/view/1538}{\texttt{https://www.hkbibleconference.org/session-message/view/1538}}
\newline
\newline
\hyperref[sec:1539]{< < < PREV SERMON < < <}
~
\hyperref[sec:toc]{[返目錄]}
~
\hyperref[sec:1537]{> > > NEXT SERMON > > >}
\newline
\newline
約翰福音 7:37-7:38
\newline
\begin{longtable}{cl}
\hline
\hline
章節 & 經文 (和合本修訂版)\\
\hline
7:37 & \begin{tabularx}{0.7\textwidth}{X} 節期的最後一天，就是最隆重的一天，耶穌站著，喊著說：「人若渴了，到我這裡來喝！ \end{tabularx} \\ \\ \relax
7:38 & \begin{tabularx}{0.7\textwidth}{X} 信我的人，就如經上所說：『從他腹中將流出活水的江河來。』」 \end{tabularx} \\ \\
[1ex]
\hline
\hline
\end{longtable}
主耶穌在約翰福音第七章卅八節說:「信我的人,就如經上所說,從他腹中要流出活水的江河來.」到底活水的江河從那裡流出來的?流出後的結果如何?
\newline
\newline
第一,活水的江河從那裡流出來的
\newline
\newline
在以西結書四十七章說這活水是從殿的門檻下流出的,又由殿的右邊,在祭壇的南邊往下流.保羅說:「豈不知你們的身子就是聖靈的殿麼!」(林前六19)神的靈住在裡頭.如何知道身體是神的殿呢?因神的靈住在他裡頭.主耶穌對撒瑪利亞婦人說:「神是個靈,所以拜祂的,必須用心靈和誠實拜祂.」(約四24)靈與靈才能相通.有神的靈住在裡頭,禱告,讀經,才能與神交通.我們常說殿即是聖殿,聖是分別出來,分別給神,所以我們要將身體分別出來歸給神,看為聖潔.
\newline
\newline
主的殿是禱告的殿,當主進入聖殿的時候,看見賣牛羊,賣鴿子,兌換銀錢的人,主就忿怒,把神的殿成了賊窩,於是把他們趕出殿去.因為神的殿是清潔的,不可有污穢.同樣,我們的身體也是清潔的,不可用言語,思想來污穢他,他是神的殿.
\newline
\newline
生命的江河從殿流出,也從壇下流出.有生命的實際,才能流出來.一個奉獻的人,要變成一個流通的管子,讓的權來管理,把聖潔,慈愛流露出來.這件事,從先知以賽亞身上便可以看出來.當烏西雅王死後,他去聖殿,見到撒拉弗有六個翅膀,用兩個翅膀遮臉,兩個翅膀遮腳,兩個翅膀飛翔.他清楚看見這個異象,明白一個奉獻給神的人亦應該如此.各翅膀的功用:遮面代表敬畏,怕得罪神.很多人對罪的看法,不與約瑟一樣.波提乏之妻引誘約瑟犯罪,約瑟不肯就範,怕得罪神,干犯神.可是,今天很多人怕犯罪,不是敬畏神,乃是怕名譽受損,怕輿論的攻擊.又有一些人怕犯罪得罪神,怕失去神的祝福.不犯罪,怕叫神傷心,才是敬畏神的人.
\newline
\newline
遮腳代表自己看不見,是謙卑.真服事神的人是謙卑的人.因撒拉弗的呼喊聲,以致門檻的根基也震動了,有這種人服事神,能力是大的.當時殿充滿了煙雲,煙雲代表神的恩典.以賽亞看到服事神的人應該與撒拉弗一樣.神對以賽亞說:「我可以差遣誰呢?誰肯為我們去呢?」(賽六8)以賽亞說:「我在這裡,請差遣我.」(賽六8)他再次把自己奉獻給神,神的榮耀充滿他,結果神的榮耀在他身上彰顯出來.
\newline
\newline
第二,活水的江河怎樣流的?
\newline
\newline
聖經說:「量了一千肘,使我趟過水,水到踝子骨.他又量了一千肘,使我趟過水,水便到膝.再量了一千肘,使我趟過水,水便到了腰.又量了一千肘,水便成了河,使我不能趟過,......(結四十七3-5)水到踝子骨,到膝,到腰,是越來越深.從踝子骨到腰,就看見一個真理,進入生命的豐盛.也不是一步登天的,而是一步一步的.今日,走屬靈的路沒有一下子就達到的.
\newline
\newline
踝子骨表示遵行神的話,無論得時不得時務要傳道,在甚麼環境下也要傳福音,這樣才叫做遵行神的話.主禱文說:「免了我們的債,如同我們免了人的債.」(太六12)債代表別人對你的虧欠,發現後,要免去別人的債,這樣才是對主誠實.若果還記住別人的債,就是對主不誠實.不誠實,就要認罪.
\newline
\newline
膝代表禱告的意思,跪下時膝碰到地上,用禱告追求進深.經上說:「同那清心禱告主的人追求公義,信德,仁愛,和平.」(提後二22)約書亞攻艾城時打敗仗,他到神面前求告,知他們中間有人取了當滅之物,就是亞干這個人.以賽亞看見自己嘴唇不潔,又住在嘴唇不潔的民中.詩人說:「因為在你那裡,有生命的源頭；在你的光中,我們必得見光.」(詩三十六9)
\newline
\newline
用禱告追求信心,經歷神的榮耀,信心才能增加,進步,以穆勒先生為例,他一生養幾千個孤兒,從不向人募損,只向神求.有一次,有個做生意的人對他說,「你把需要告訴我,我可以考慮,幫助你.」他回答:「我不告訴你,我只告訴神.」
\newline
\newline
腰代表服事.那時代的用人,用布綁著腰,是為工作方便.服事可以進深,當別人的難處作為你的難處,給你學習忍耐的功課,別人批評,攻擊你,你便可以學習饒恕人的功課.教會中有人做頭,爬在你前面,你遇到各種的困難,事情不順利,你就學習到等候的功課.或者你做事失敗,別人藐視你,你就學習不灰心.所以在服事中可以進深.
\newline
\newline
由此可知,遵行神的話,禱告,服事是達到屬靈進深的三條路.進深是甚麼?簡單的話說,是失去自己.從踝子骨到腰,水的阻力越來越大.雅各對耶穌說:「人要顯揚名聲,沒有暗處行事的；你如果行這些事,就當將自己顯明給世人看.」耶穌回答說:「我的時候還沒有到；你們的時候常是方便的.」(約七4,6)進深就是不方便；不進深的人,只有自己,可以隨意來去.追求進深就不行了,以前去夜總會,賭場,現在追求進深,不方便去了,去就不平安,你就不會去.以前敢說傷人的話,今大要去求人饒恕,才得平安.從前要穿甚,就穿甚麼,別人不能干涉,現在就不行了,穿的要榮耀神.由此可見,追求進深的人,是完全不方便的人.不追求的人,就很方便了.水到腰就沒有自己了,以神的主張作他的主張,以神的意念作他的意念,完全失去自己的人了.
\newline
\newline
生命的江河流出來,完全為神.生命的江河流出來後結果怎樣呢?
\newline
\newline
一,流入鹽海,使水變甜
\newline
\newline
鹽海是死海,流入不流出,故叫它為死海.在地平線下八百尺,氣候很熱.有一次我到那裡,風吹來,皮膚也會痛.海水很鹹,很苦,喝一兩口會叫人致命.人生有許多苦,一個接一個,如空虛的苦,環境變化的苦,人事打擊的苦.人生是一片的苦.生命的江河變苦為甜.在我來香港之前,有一個姊妹來找我,哭著對我說:「我在上海結婚,是牧師給我證婚的,現在我要與丈夫離婚,我與他結婚三十八年,但他與我結婚後的第九天就去尋花問柳,如今又在外面搞女人......」夫妻變仇敵.他要我去見她丈夫,當我見到她丈夫時,他無話說,他要求神饒恕,要我為他禱告,結果夫妻和好如初,我來香港前,他們還為我餞行呢!苦的家庭變成甜的家庭了.保羅說:「我真是苦啊,誰能救我脫離這取死的身體呢?」(羅七24)
\newline
\newline
二,治水所到之處,百物都必生活
\newline
\newline
活水把生命帶出來,活出來.保羅就是這樣的人,外邦人用「瘟疫」來形容他,因他身上有基督,所到之處有人信主.福音帶給人生命,如保羅用福音生了提摩太,使他成了有用的人.
\newline
\newline
三,叫兩岸的樹都結出果子來
\newline
\newline
環境不好結忍耐的果子,人毀謗時結溫柔的果子,如亞伯拉罕與羅得為牧羊之事,亞伯拉罕不與羅得相爭,他結愛心的果子來.約瑟說:「你們原是要害我,但神的意思是好的.」他結出和平的果子.常有果子結出,葉子可以治病.
\newline
\newline
可是,有兩種地方是得不到醫治的:一是窪濕之地.二是泥濘之地.他們雖不覺得自己甜,但也不覺得苦,不如江河好,也不像乾地那樣壞,他還是濕的,不像海洋豐富,但也不像死海一無所有.也不像猶大,比上不足,比下有餘,自為滿足,這種人沒有希望,不得醫治.
\newline
\newline
我希望你不像泥濘或窪濕的人,不以現在的屬靈光景為滿足.不是自義的人,而是努力追求進深,讓活水的江河從你生命中流出來.
\newline
\newline


\newpage

\section{第40屆~港九培靈會~吳勇~第十四講}
\label{sec:1537}
\textbf{神的傷心}
\newline
\newline
連結: \href{https://www.hkbibleconference.org/session-message/view/1537}{\texttt{https://www.hkbibleconference.org/session-message/view/1537}}
\newline
\newline
\hyperref[sec:1538]{< < < PREV SERMON < < <}
~
\hyperref[sec:toc]{[返目錄]}
~
\hyperref[sec:1536]{> > > NEXT SERMON > > >}
\newline
\newline
以賽亞書 66:7-66:8
\newline
\begin{longtable}{cl}
\hline
\hline
章節 & 經文 (和合本修訂版)\\
\hline
66:7 & \begin{tabularx}{0.7\textwidth}{X} 「錫安未曾陣痛就生產，疼痛尚未來到，就生出男孩。 \end{tabularx} \\ \\ \relax
66:8 & \begin{tabularx}{0.7\textwidth}{X} 國豈能一日而生？民豈能一時而產？但錫安一陣痛就生下兒女，這樣的事有誰聽見，有誰看見呢？ \end{tabularx} \\ \\
[1ex]
\hline
\hline
\end{longtable}
神的傷心藉阿摩司先知表達出來.他是耶羅波安王時代的先知.那時軍事非常成功,收復邊界,從哈馬口直到亞拉巴海.政治也成功,做王四十一年.經濟也成功,夏天有夏天的房屋,冬天有冬天的房屋,而且還有象牙房屋,那時的人享受富裕.
\newline
\newline
既然以色列國如此成功,何以神要悲哀,傷心呢?神不是從軍事,政治,經濟面去看的,而是從社會和宗教方面去看的.神看以色列國如一棵茂盛的大樹,外面青翠,裡面卻爛朽了.這樹雖好看,但不能持久,所以過了五十年該國就被亞述所滅.
\newline
\newline
(一)社會道德方面
\newline
\newline
他們沒有憐憫的心,他們為銀子賣了義人,為一雙鞋賣了窮人.他們對窮人頭上所蒙的灰也垂涎......當時的社會,敗壞到非常可怕,沒有一點憐憫的心.本來在舊約時代,神吩咐猶太人,收割麥子時,要留下一些,叫窮苦的,孤寡的人,可以去收拾,才有得吃.即是說,有的人要看顧沒有的人,如果今日的人能夠這樣行,這個社會就要減少很多犯罪的勾當了.因許多盜竊,劫匪,卑鄙的事,殺人,傷天害理的事,都是窮到要命才逼出來的.同時,也可以減少許多自殺的案件,報上我們常見貧病交加的人上吊自殺,走上末路.
\newline
\newline
那時代的人自私,犧牲別人來滿足自己,不但沒有憐憫的心,而且變成殘忍的心.更可悲的,父子同一女子行淫,這是淫亂,污穢到極點.今日的社會也是如此.如電影,報紙,廣告,雜誌等都盡量渲染色情的東西,以這些黃色的東西來吸引人,除滿足自己的利益外,便是以此來迎合人的需要,因為人的本性已經敗壞到極點,非常可怕,到達不可收拾的地步.
\newline
\newline
商業場所是講究唯利是圖的,加上人心敗壞,道德觀念不存在了,也沒有貞操的觀念了.所以男人顧不到對妻子貞潔.由於種種緣故,今天色情泛濫之盛,沒有一個時代可與之相比,不但有逼良為娼的事,而且也有丈夫逼妻為娼,父母逼女為娼.今天的社會,不但可用金錢買淫亂,甚至不用金錢也可得淫亂,有夫之婦搭上有婦之夫,有婦之夫搭上有夫之婦.阿摩司時代如此,今日的社會也如此.當時神憎恨猶太人的罪,今日神也憎恨這社會的罪.
\newline
\newline
還有,他們又在神的廟中,喝受罰之人的酒.這些人利用偶像的地位,向拜偶像的人搾取.這種事,今天到處可見,政府人員向百姓搾取,就以東南亞某個國家,海關人員向入口的人要錢.這種事敗壞自己的人格,必定要惹起神的震怒.
\newline
\newline
神責備巴珊的女人,祂責備男人,也責備女人,即是責備婦女的虛榮.世界上男人的成功是女人的緣故.女人被造的目的是甚麼?就如神造亞當後見他獨居不好,所以神就造夏娃來幫助他.對丈夫來說,妻子不但在道德上幫助,而且在事業上也得幫助.丈夫的缺點,妻子看得最清楚,糾正他的缺點.做妻子的,使丈夫成功時不驕傲,失敗時不灰心.但是,有很多男人失敗,也是女人的緣故,女人本是內助,如今卻成了內阻.神對巴珊的母牛說,當聽我的話,你們欺負貧寒的,壓碎窮乏的,對家主(即丈夫)說,拿酒來,我們喝罷.酒是刺激肉體的東西,女人為她們的肉體,好虛榮,講究打扮,裝飾,時髦.
\newline
\newline
當你走過商店時,有百分之八十是賣女人的東西,而且千變萬化.巴珊的女人,看見別的女人打扮好,也希望自己也如此.今天有的的男人,為了妻子的虛榮,便做出許多違背天良的事,今日社會很多的罪,是因女人的緣故.
\newline
\newline
(二)教會實際方面
\newline
\newline
他們不是不聽神的訓誨,而是聽了不實行.據說歐美國家參加主日禮拜的人有百分之四十.人數很多,但教會沒有影響社會,以致倫理墮落,道德敗壞,他們厭棄神的律例.聽了之後,不照著去行,所以一點作用也沒有.正如以西結書卅三章卅二節說的:「他們看你如善於奏樂聲幽雅之人所唱的雅歌；他們聽你的話卻不去行.」
\newline
\newline
當時的人把聽道當作聽戲,今天的信徒更是聽戲的專家,他們會分析牧師的講章內容如何!只聽不行.神要祂的兒女聽從祂,遵行祂,和順從祂.不順從的後果就是死,亞當為例.神吩咐他不可吃分別善惡樹上的果子,吃後就要死,死有三種:1.靈死.2.體死.3.屬靈生活的死.如一個得救的人,對神的話沒反應,對神的愛也沒反應,這簡直在死亡中.人與神交通是藉著神的話,到了不行神的話的時候,工具沒有了,我們就軟弱,工作沒有力量,爭戰更不能得勝,這種光景就是屬靈生活的死亡.一個重生了的人,還會死亡嗎?不是指靈的,死乃是指他屬靈生活的死亡.
\newline
\newline
所羅門曾被靈感動說:「我兒,要把你的心給我.」這是神要說的話,藉所羅門講出來.表明神要得人的心,佔有人的心.如果我們的心被神之外的甚麼得了,那甚麼就是偶像,無論是屬靈的事或是社會上的事也是如此解釋.例如一個女子寫情書:「你是我的偶像......」表明你得著,佔有她的心.若有一件物佔有你的心,得著你的心,它就是偶像.若有一個人佔有你的心,這個人就是你的偶像.有的人,他的心被金錢佔有了,他以為金錢萬能,口說主是他的一切,其實金錢才是他的一切.有的人被學問佔有他的心,口裡說主是他一切,其實他的心中只有學問.有的人以地位為他的一切,以地位為榮耀,他的心被地位佔了.以為有了它,人就會恭維他,討好他,不但有物質享受,也有精神享受.
\newline
\newline
使徒保羅說:「我為你們起的憤恨,原是神那樣的憤恨；因為我曾把你們許配一個丈夫,要把你們如同貞潔的童女,獻給基督.」(林後十一2)保羅憤恨他們不把基督當作主,把祂看作是奴隸.今天的人也是如此,不是差我們,而是我們差祂,使用祂.
\newline
\newline
神厭惡他們的節期,也不喜悅他們的嚴肅會.神為甚麼那麼說呢?因為那時的百姓一面犯罪,一面又守嚴肅會,一面犯罪,又一面守節期,神當然不會喜悅.今天的人也是如此,很多人一面跳舞,一面又做禮拜,一面犯罪,一面又作聖工,神是喜愛公平,公義的人.神警告我們說:「不要將你們的肢體獻給罪作不義的器具；倒要像從死裡復活的人,將自己獻給神:並將肢體作義的器具獻給神.」(羅六13)
\newline
\newline
他們還有一種罪,就是不為罪擔憂,不為約瑟擔憂.約瑟是預表基督.他年少時,被哥哥們妒忌,逼迫,賣他下埃及,而他在埃及做了宰相,後來迦南地有飢荒,約瑟因而救了他們.我們的主為天父所愛,祂到世上,猶太人把祂釘十字架,但祂卻得榮耀.
\newline
\newline
基督與教會是聯合成一體的,今日的教會有異端,把心靈的實際變成了形式,這是今日教會的苦難,昔日的以色列人不為約瑟的苦難擔憂,同樣今日的基督徒,也不為教會的苦難擔憂.
\newline
\newline
今天神的教會有了苦難,你有與基督表同情嗎?你有起來對付這些苦難嗎?你有為今日教會的苦難擔憂嗎?
\newline
\newline


\newpage

\section{第40屆~港九培靈會~吳勇~第十五講}
\label{sec:1536}
\textbf{主的道路}
\newline
\newline
連結: \href{https://www.hkbibleconference.org/session-message/view/1536}{\texttt{https://www.hkbibleconference.org/session-message/view/1536}}
\newline
\newline
\hyperref[sec:1537]{< < < PREV SERMON < < <}
~
\hyperref[sec:toc]{[返目錄]}
~
\hyperref[sec:1533]{> > > NEXT SERMON > > >}
\newline
\newline
馬太福音 7:13-7:14
\newline
\begin{longtable}{cl}
\hline
\hline
章節 & 經文 (和合本修訂版)\\
\hline
7:13 & \begin{tabularx}{0.7\textwidth}{X} 「你們要進窄門。因為通往滅亡的門是寬的，路是大的，進去的人也多； \end{tabularx} \\ \\ \relax
7:14 & \begin{tabularx}{0.7\textwidth}{X} 通往生命的門是窄的，路是小的，找到的人也少。」 \end{tabularx} \\ \\
[1ex]
\hline
\hline
\end{longtable}
在約翰福音十四章講到一件事,多馬問主:「我們不知道你往那裡去?怎麼知道那條路呢?」耶穌答道:「我就是道路,真理,生命.」主在世上卅三年,一初經過,所言所行都是真理.要明白真理,找到真理,才有生命.
\newline
\newline
有一次,我去看一個有學問的人,問他:「你有神的生命嗎?」他答:「怎樣才有呢?」我便用約翰壹書五章十二節那裡說的「人有了神的兒子就有生命」這句話對他講,神的兒子成了人,祂成了人有兩個目的:第一,要有一個身體,才能叫人看見.第二,有了身體才能作你的替身,替你釘十字架.他聽後就對我說:「這些道理我都聽過,我都懂了,但如何得生命呢?」接著我便把約翰三章十六節讀出給他聽,信耶穌就有生命,不信則沒有.比方,這裡有一滴水,有人說它有細菌,有人說它沒有細菌.從肉眼看它,沒有細菌,若是放在顯微鏡下看,就可看見細菌了.
\newline
\newline
主曾叫三個人復活:1.叫睚魯的女兒復活,起來後,給她東西吃.2.叫拿因城寡婦的兒子復活,復活後說話.3.叫拉撒路,復活後行走.我對這個有學問的人講到生命是會的,要接受神的糧.生命會說話,讚美神,見證神.生命會行走,行神的道.主是真理的道路,明白真理的道路後,就可得著生命.
\newline
\newline
到底主的道路是怎樣的呢?
\newline
\newline
(一)貧窮的路
\newline
\newline
有一個文士來,對主說,夫子,你無論往那裡去,我要跟從你.耶穌對他說,狐狸有洞,天空的飛鳥有窩,人子卻沒有枕頭的地方(太八19-20).很多人跟隨主,有錯誤的看法.如西庇太的兩個兒子,他們的母親在耶穌面前求,使她兒子一左一右的坐在耶穌兩旁,以為跟隨主有高官顯祿.那個文士要跟隨主,也有同樣的錯誤,以為跟隨主後可以富有享福,他的看法是屬世的.
\newline
\newline
主耶穌把這條貧窮的道路擺在他面前.我每次證婚時總對新娘說,你願嫁某人為妻嗎?你是否願意在他健康或軟弱時,富足或貧窮時愛他呢?如果她說,只在健康時,富足時才愛他,那我一定宣佈這婚證不成了,因為這樣的人,不能做好妻子,不能白頭偕老.
\newline
\newline
跟隨主是貧窮的,事實上許多跟隨主的人並不貧窮,而且還很富足.故這貧窮的意思,即是無論甚麼環境必跟隨到底.如殉道史記載波利甲被捕,若他否認耶穌即可得釋放,但他卻說:「我服事祂多年,不願放棄祂,生要跟隨祂,死也要跟隨祂.」結果被殺,為主殉道.可是,底馬因貪愛世界,離開了主,中途變節,不能跟隨主到底.還有比底馬更厲害的人,就如猶大,為卅兩銀子賣了耶穌,結果他失去了屬靈的福氣.
\newline
\newline
我們的人生並非一帆風順,一路平坦的,如果我們沒有在任何環境中都跟隨主到底的心願,到了一天可能像底馬一樣,甚至像猶大一樣.
\newline
\newline
跟隨主的人,就是完全聽從主的話.有一天,主要我們為祂的榮耀和工作,有所捨,有所撇,那時你若捨下不下,撇不下,你將會抵抗神的旨意,那樣就會失去將來的賞賜.所以主耶穌擺出這條路,要跟隨祂的人,在任何環境下,跟隨祂到底.
\newline
\newline
(二)艱難的路
\newline
\newline
進入神的國,艱難很多.聖經講到兩種的門路,一種門是寬大的,容易進去；另一種門是窄的,不易進去,非按規矩不能進去.一種是大路,容易走,一種是小路,不易走.屬世的路是大的,容易走,屬靈的路是窄的,不易走.以約瑟為例,他被哥哥們賣去埃及.約瑟在波提乏家當家宰,一天,波提乏的妻千方百計的引誘他犯罪.這時候,擺在他面前有兩條路,一是屬世的路,易走,只要順從她就行了.一是屬靈的路,不易走,要制死肉體的情慾,為神要分別為聖,結果約瑟下監坐牢,受了冤曲.
\newline
\newline
進入神的國不容易,有艱難,原因何在呢?聖經說:「我心裡柔和謙卑,你們當負我的軛,學我的樣式,這樣,你們心裡就得安息.」(太十一29)
\newline
\newline
我們當學習主的柔和,祂被罵不還口,受苦不說威脅的話.主是謙卑的,祂取了如僕的形像,成了人的樣式,意即祂原來是神,為我們成了人.神是不受限制的,為我們成了受限制的,祂在十字架上可以下來,但祂不下來.我們要學主的樣式,今天人與人之間有難處,不能相處,乃因沒有十字架.十字架有兩根木頭,一橫一直,橫與直是代表神與人的關係.當時的人加給主的羞辱,戲弄,傷害,是為人的緣故.今日我們要看人加給我們的羞辱,戲弄和傷害是為主的緣故.在德國,有一個教授,一天,有一個女人抱著一個孩子來找他,那時他正在上課.他問道:「你找誰?」她說:「找你!」又問:「你手上抱的......」她答:「是孩子,是你罪惡所生的孩子.」她把孩子拋在他手中就走出去了.他的學生們看見此情形都走光了,教授很痛心,又被別人罵.有一天,他忍不住,到森林中去,跪下來禱告.他的心平靜了,抱孩子回家,後來這婦人又回來抱回她的孩子去了.因為這個女人是有神經病的.這個教授被人戲弄,傷害,他不還口,他學了主的樣式.
\newline
\newline
(三)受苦的路
\newline
\newline
聖經說:「因基督也為你們受過苦,給你們留下榜樣,叫你們跟隨祂的腳蹤行.」(彼前二21)跟隨主就是祂走的路,我們也走,祂所受的,我們也受.在歌羅西書一章廿三節說:「補滿基督患難的缺欠.」這句話可用國父遺囑來說明,「革命尚未成功,同志仍須努力」,從前革命先烈流血捨命,如黃花崗七十二烈士.國父把滿清推翻,民國建立了,何以說未成功?因孫中山先生要建立民主國家,但內部軍閥盤據,所以土地沒有改革,平均地權未實現,這是說後繼人要繼續努力,所以說那時先烈的革命只完成一部分,還有一部分等後繼人去完成.從前先烈們拋頭顱流熱血才完成革命一部分,所以後的要繼續去完成.我們的主也只是完成天國基礎的一部分,還有一部分神要藉著我們去完成.主在受苦中完成天國的一部分,建設的部分要我們去完成.因此我們就要努力把福音傳遍天下,把人遷到愛子的國度中去.我們這樣做才叫補滿基督患難的缺欠.從前大利有個革命領袖名加里波特,有一次,他對千萬青年演講,鼓勵他們為祖國自由而參戰,但有一個膽小的青年問他,「先生,我參戰,得甚麼賞賜.」他不妥協的說:「或者傷,或者殘,或者甚至死亡,但你要知道,意大利因你的傷殘死亡,便得了自由.」這是苦難的道路.
\newline
\newline
撒旦是神的對頭,我們站在神一邊,我們也是撒旦的對頭.所以撒旦常用苦難來攻擊我們,使我們的道路不是通順,而是崎嶇的；不是安逸的,而是有苦難的.
\newline
\newline
十年前,南美洲有個未開化的奧卡族,曾有五個美國青年人入去傳福音.開始用飛機散下福音單張,直至他們沒有惡意便在他們村落降下飛機.晚上住宿村中,因受另一族人的煽動,結果五個青年人被奧卡族人所殺.後來,有個殉道者的姊妹決心進去,要傳福音給奧卡族,叫他們信耶穌.經過她的努力,今天的奧卡族有人信了主,也有了教會.去年在柏林召開的世界宣教大會,殺害殉道者的兇手也出席了該次的大會.今天,奧卡族改變了,不是人的力量,乃是福音的能力.
\newline
\newline
主擺在我們面前的道路,是貧窮的,艱難的,受苦的.你肯走這條路嗎?你有這個心志嗎?
\newline
\newline


\newpage

\section{第40屆~港九培靈會~吳勇~第十六講}
\label{sec:1533}
\textbf{跟隨耶穌}
\newline
\newline
連結: \href{https://www.hkbibleconference.org/session-message/view/1533}{\texttt{https://www.hkbibleconference.org/session-message/view/1533}}
\newline
\newline
\hyperref[sec:1536]{< < < PREV SERMON < < <}
~
\hyperref[sec:toc]{[返目錄]}
~
\hyperref[sec:1508]{> > > NEXT SERMON > > >}
\newline
\newline
哥林多前書 2:14-2:15
\newline
\begin{longtable}{cl}
\hline
\hline
章節 & 經文 (和合本修訂版)\\
\hline
2:14 & \begin{tabularx}{0.7\textwidth}{X} 然而，屬血氣的人不接受神的靈的事，他反倒以這為愚拙，並且他不能了解，因為這些事惟有屬靈的人才能領悟。 \end{tabularx} \\ \\ \relax
2:15 & \begin{tabularx}{0.7\textwidth}{X} 屬靈的人能看透萬事，卻沒有一人能看透他。 \end{tabularx} \\ \\
[1ex]
\hline
\hline
\end{longtable}
主在加利利呼召彼得,安得烈,約翰,雅各來硍隨祂.以後在馬太立第十章祂又說,要跟隨祂,就必須背起他的十字架來跟隨.前者是初步的跟隨,後者是進一步的跟隨.現在我們要來看甚麼是初步的跟隨?甚麼是進一步的跟隨?
\newline
\newline
(一)初步的跟隨
\newline
\newline
主初步的呼召,目的有兩個,給他們知道天國近了,主禱文說「願你的國降臨」,呼召人進祂的天國,做祂的國民.國民有兩種,一種是做的,一種是生的,兩者有很大的區別.做的沒有神的生命,不能進天國.生的有神的生命,可以進入神的國.馬太福音第五,六,七的三章聖經是講到天國的憲法,一定要有天國的生命,才能明白天國的憲法.人裡頭有神的靈才能明白屬靈的事.否則就不能明白.我與內人結婚,她先信主,很熱心,要我陪她去做禮拜,本來我不願意去的.因剛結婚,愛情深厚,故我就去了.那天的聚會,聖靈做工,聽道的人流淚哭了,一直喊著「烏老鼠」,我抬頭四圍觀望,連「烏老鼠」的影子都看不見,那裡來「烏老鼠」呢!散會後我與內人一起回家,我就把「烏老鼠」的事對她說,她聽後大大的責備我,我真是無知,原來他們是喊讚美主.福建話「讚美主」的俗語「呵腦主」,「呵腦主」與「烏老鼠」音相近,因為福建人讀烏為呵.
\newline
\newline
約翰壹書五章一節說:「凡愛生他之上帝的,也必愛從上帝生的.」如一個家庭,因愛父母,就愛父母所生的弟兄姊妹.有一次我去紐西蘭,未去之前,有一位傳道人對我說:「你在紐西蘭有朋友嗎?」他便把紐西蘭一個開浸信會書店的朋友地址告訴我,叫我去找他.當我到了紐西蘭後,適逢開會,所以下,中,上的所有旅店都住滿了人.我怎麼辦呢?我這時便想起傳道人的話,我決定去找開書店的人.找到他,快下班了.他先帶我去吃中國餐,是他請客.吃完飯又帶我去觀光.然後又週到的安排我住宿.他說的話使我快樂.因為他有神的生命,他便愛神所生的人.
\newline
\newline
主呼召門徒,叫他們要得人如得魚一般.「叫」字的英文是「造就」.造就甚麼呢?以前是得魚的,現在是要得人了.得魚是為自己,為家庭而活.得人是為別人而活,用口說耶穌,用身活耶穌.人家看見就羨慕,便要追求耶穌.祂呼召我們不但做天國的國民,而且祂要造就我們,叫我們為人而活.
\newline
\newline
主第一次召彼得,安得烈；第二次召約翰,雅各.兩次都是召兩個人.主派門徒出去作工,也是兩個兩個去的.神造亞當,見他獨居不好,就造配偶來幫助他.一個人不好,兩個就好了.在屬靈的路程上,一個軟弱,另一個可扶助；一個跌倒,另一個給他警戒；一個工作有錯誤,另一個幫他糾正；一個灰心,另一個鼓勵他；一個傷心,另一個給他安慰.兩個人的目的,就是如此!
\newline
\newline
主呼召的彼得,安得烈是做撒網的工作,而約翰,雅各則是做補網的工作.彼得向猶太人撒網,結果有三千人,五千人信主,向外邦人撒網使哥尼流一家人得救.使徒時代,神的道弄偏了,從心移到頭腦,但約翰是做補網的工作,他把神的道弄正,把網補好.他說:「論到起初原有的生命之道,就是我們所聽見所看見,親眼看過,親手摸過的.」(約壹一1)為何約翰這樣說呢?他要糾正他們頭腦的知識.所以,主要有撒網和補網的工人.今日的教會,也須要有這兩種人.
\newline
\newline
在大衛的時代,有人做葡萄園,牧羊的工作；有人做對內,也有人做對外的工作.一個做這,一個做那,叫神的工作發展,不受攔阻.主召彼得,安得烈,和主召約翰,雅各有相同的事發生,他們都是「立刻」跟隨主的.如果不是立刻順從,則會消滅聖靈的感動.有一次,我到監獄去講道,有一個犯人就是懊悔他消滅聖靈的感動.雖然他從小就聽過神的道,但學壞了,想發財,結果加入強盜集團.聖靈一直在他心中做感動的工作,可是卻消滅聖靈的感動,甚至還對神理論,「現在我兩手空空,生活沒有依靠,等我發了財之後,便將我的心獻給你.」後來因殺人被捕入獄,判無期徒刑.他的後果是消滅聖靈感動而招致的.今日若聖靈感動我們,就須立即行動,否則消滅聖靈感動,就要一生懊悔了.
\newline
\newline
彼得,安得烈,約翰,雅各跟隨主是捨了船,捨了網,也捨了父親.父親代表親情,所以他們是捨了物質親情來跟隨主的.以前有一個聖人,名叫法蘭西斯,出身貴族,他要去做痲瘋人的工作,遭家人反對.他父親要與他斷絕父子關係,不給他財產,若他要去做這種工作.他願意撇下物質和親情來做這工作,這是真正的跟隨主.
\newline
\newline
(二)進一步的跟隨主
\newline
\newline
這個跟隨主是背起自己的十字架來跟隨的.主要去耶路撒冷受死,彼得聽後說:「主啊!此事不會臨到你身上.」主答:「撒旦,退後邊去吧!」彼得是體貼人的意思,不是體貼神的意思,這裡有三件事攔阻他的.
\newline
\newline
第一,主所走的是受苦的路
\newline
\newline
彼得拉他,表明人不願走受苦的路.聖經說:「萬事互相效力.」神造就人是用各種方法的.屬靈的人是造就出來的.
\newline
\newline
第二,撒旦要破壞人的靈性
\newline
\newline
主說,撒旦退我後面去吧!撒旦要在我們身上爭地位,把神從我們心中驅逐出去.
\newline
\newline
第三,體貼人不體貼神的意思
\newline
\newline
例如主在客西馬尼園禱告,祂對神說:「然而要成就你的旨意,不要成就我的旨意.」以前上海有一個人,他要到英國劍橋大學去學醫,想不到神呼召他出來做救人靈魂的工作,於是屬世的工作和屬靈的工作,兩者在他心中交戰,聖靈感動他,他得勝了.當船到法國馬賽時,他對船長說要回去上海.當他回到上海時,家人不明白他的行徑,然而他卻要遵行神的旨意,而願犧牲自己的益處.
\newline
\newline
不捨己就得不到造就,而會變成撒旦的立足之地.主說要背起自己的十字架來跟隨祂.到底十字架是甚麼呢?十字架是被憎惡,輕視的意思.以前孫大信信主後,被家人趕出家門,被輕視,被憎恨.他父親給他辦喪事,當他死了.那天趕他走時,他吃的食物是下了毒藥的,在路上毒性發作,幸被人救起.十字架是棄絕之意.今日的人,寄望於物質,希望明天比今天好.今天的人,為了肉體的享受,就不理道德與貞操的觀念了.若有一個人是分別為聖,聖潔慈愛的,別人會把他當作傻瓜,沒有出息的人.
\newline
\newline
我年幼的時代是重男輕女的時代,我母親一連生了三個女兒,祖父很不高興,母親也為此受了很多冤屈.生到第四個,我出世了,祖父很高興,看為掌上明珠,甚得他寵愛,以致養成任性的性格.當我在十三歲時,父親罵了我一頓,我一頓腳就離家出走,不要父母,家庭了.我一走,母親就病倒,不吃,不睡,醫生說是心病,不是生病,醫藥沒用,最好是叫兒子回來.後來母親打聽到我在檳城的親戚家中,她來找我.起初我不睬她,她傷心極了,我良心發現,回到她身邊,是因著她的愛的感動.
\newline
\newline
從使徒時代到今天,有很多人跟隨主.今天,主要你跟隨祂,你要初步的跟隨主,還是要進一步的跟隨祂呢?
\newline
\newline


\newpage

\section{第40屆~港九培靈會~吳明節~第一講}
\label{sec:1508}
\textbf{旗開得勝的摩西}
\newline
\newline
連結: \href{https://www.hkbibleconference.org/session-message/view/1508}{\texttt{https://www.hkbibleconference.org/session-message/view/1508}}
\newline
\newline
\hyperref[sec:1533]{< < < PREV SERMON < < <}
~
\hyperref[sec:toc]{[返目錄]}
~
\hyperref[sec:1532]{> > > NEXT SERMON > > >}
\newline
\newline
出埃及記 17:8-17:16
\newline
\begin{longtable}{cl}
\hline
\hline
章節 & 經文 (和合本修訂版)\\
\hline
17:8 & \begin{tabularx}{0.7\textwidth}{X} 那時，亞瑪力來到利非訂，和以色列爭戰。 \end{tabularx} \\ \\ \relax
17:9 & \begin{tabularx}{0.7\textwidth}{X} 摩西對約書亞說：「你為我們選出人來，出去和亞瑪力爭戰。明天我要站在山頂上，手裡拿著神的杖。」 \end{tabularx} \\ \\ \relax
17:10 & \begin{tabularx}{0.7\textwidth}{X} 於是，約書亞照著摩西對他所說的話去做，和亞瑪力爭戰。摩西、亞倫和戶珥都上了山頂。 \end{tabularx} \\ \\ \relax
17:11 & \begin{tabularx}{0.7\textwidth}{X} 摩西何時舉手，以色列就得勝；何時垂手，亞瑪力就得勝。 \end{tabularx} \\ \\ \relax
17:12 & \begin{tabularx}{0.7\textwidth}{X} 但摩西的雙手沉重，他們就搬一塊石頭來放在他下面，他就坐在上面。亞倫與戶珥扶著他的手，一個在這邊，一個在那邊，他的手就穩住，直到日落。 \end{tabularx} \\ \\ \relax
17:13 & \begin{tabularx}{0.7\textwidth}{X} 約書亞用刀打敗了亞瑪力和他的百姓。 \end{tabularx} \\ \\ \relax
17:14 & \begin{tabularx}{0.7\textwidth}{X} 耶和華對摩西說：「你要把這事記錄在書上作紀念，又念給約書亞聽：我要把亞瑪力的名字從天下全然塗去。」 \end{tabularx} \\ \\ \relax
17:15 & \begin{tabularx}{0.7\textwidth}{X} 摩西築了一座壇，起名叫「耶和華尼西」。 \end{tabularx} \\ \\ \relax
17:16 & \begin{tabularx}{0.7\textwidth}{X} 他說：「我指著耶和華的寶座發誓，耶和華必世世代代和亞瑪力爭戰。」 \end{tabularx} \\ \\
[1ex]
\hline
\hline
\end{longtable}
摩西帶領以色列人與亞瑪力人爭戰,得勝了便築一座壇,起名叫「耶和華尼西」.就是「神是我的旌旗」.這是神幫助以色列人作戰.
\newline
\newline
在人類歷史上,神揀選以色列人是最奇妙的.神藉此民族來顯明祂的旨意.用各種方法來訓練,使以色列人成為神的精兵.當以色列人作戰時,神作他們的幫助,元帥,故此,凡以色列人倚靠神作戰的,必然得勝,反之,就必定失敗.
\newline
\newline
以色列人出埃及,過紅海,神將法老的軍兵投在海中,以色列人未曾接戰,兵不血刃,便渡過紅海.
\newline
\newline
故此,與亞瑪力人爭戰,可說是以色列人出埃及以來,第一次正式戰爭.對以色列人出埃及進迦南來說,這是一場很重要的戰爭,如失敗,今後在進迦南的路上,更顯得軟弱.故此,以色列人在這場戰爭中,必須得勝.
\newline
\newline
我聽過許多人解釋此段經文,說亞瑪力人是代表我們的肉體,基督徒常與肉體爭戰,此解釋,很能幫助我們.
\newline
\newline
但以色列人是全會眾與亞瑪力人爭戰,故此,如解釋為教會與罪惡爭戰更為恰當.摩西是以色列人的領袖,教會也有領袖和許多同工,故看此爭戰為全教會與罪惡爭戰,更能幫助我們.現在我要分三點解釋:
\newline
\newline
一,作戰時,摩西手中拿杖.把杖舉起,以色列人就得勝,手下垂,以色列人便失敗.因此這杖不能垂下,只能舉起,才能打勝仗.在此要明白,不是這杖有甚麼能力,乃是因為這是神的杖,其表意有二:
\newline
\newline
1. 代表神的權柄.摩西與仇敵爭戰,需要神賜他權柄.如沒有權柄,就不能與仇敵爭戰.詩篇八十九篇三十二節與以西結書十九章十一節,均有提及杖是權柄的預表.今天教會領袖也需要從神來的權柄.
\newline
\newline
2. 杖也是表神的能力.詩篇一百一十篇二節說到神有能力的杖.與仇敵爭戰,權柄與能力是都不能缺少的事.把神的權柄能力舉起來的時候,就得勝仇敵.不是以色列人的英勇,也不是以色列人是大能的勇士,是神的權柄,神的能力護衛他們.今天教會的領袖,聖工人員,以及全會眾,以及全會眾,與罪惡戰爭,同樣的不能缺少神的權柄與能力.
\newline
\newline
做教會聖工,同樣的要神的權柄,能力.這權柄是指福音說的.無論傳道,長老,執事,聖工人員,都需要神賜我們這福音的權柄.在哥林多前書九章十四至十八節便是說到神賜我們福音的權柄.
\newline
\newline
今天教會,聖工人員的失敗,不是缺少聰明,也不是缺少才幹,乃是看錯權柄.在教會作工,不是倚賴議案通過之權柄,乃是要倚靠神所賜福音的權柄.
\newline
\newline
能力,是聖靈充滿的能力.「但聖靈降臨在你們身上,你們就必得著能力.」(徒一8)作神的工作,不是屬聖靈的人,一定失敗.從事聖工,若不被聖靈充滿,一定無能力與撒旦作戰.
\newline
\newline
許多教會興旺,一定對福音有熱心,聖工人員一定過著屬靈的生活,有屬靈的生命.反過來說,如教會冷淡,一定是對福音冷淡,信徒不是屬靈的,乃是屬血氣的.
\newline
\newline
各位弟兄姊妹,教會在今天的世代,每天都是過著戰爭的生活,不但個人戰爭,而且是全會眾在戰爭,不戰爭就不能存在.
\newline
\newline
摩西舉杖,全會眾得勝,垂下手來全會眾便失敗,這告訴我們,作教會領袖的,不把神的福音權柄高舉不為功.
\newline
\newline
二,摩西有好同工.約書亞是大能的勇士,衝鋒陷陣,攻擊仇敵.摩西差派他,他勇往直前,不懷疑,服從神的權柄,接受神的能力,與摩西同心合意,一心一德的作戰.
\newline
\newline
當摩西舉杖,手下垂.摩西也有肉體的軟弱,但同工來幫助他.不但搬石頭來,而且一邊一個人,扶住摩西的手,從早到晚,於是摩西的手都舉起來了.就大大得勝,將仇敵消滅.
\newline
\newline
這事何等美好,亞倫和戶珥與摩西同心,扶助摩西.他們是扶助摩西舉杖,不是因摩西有肉體的軟弱,便輪流舉杖.摩西是神選召的,權柄,能力是神賜給摩西的.如果亞倫,戶珥要舉杖,便是奪權者.如果他們發生奪權,全會眾的的爭戰,必定屍骸遍野,慘重失敗.亞倫,戶珥不是從摩西手中奪權,乃是盡力協助摩西舉起神的權柄.
\newline
\newline
今天教會有三兩個同工,便常常發生困難.有人說,傳福音不難,同工才是困難.神選召同工,各有託付,各有責任,應同心合意,發展聖工,為何發生困難?這往往起因於奪權,人人要當第一,所以發生紛爭,同工紛爭,全教會必然失敗.
\newline
\newline
神帶領以色列人,給他們最好的領袖,且有這許多同工.神也恩待教會,給教會有好領袖,好同工.若能彼此同心,教會必然大大得勝,主將得救的人,天天加給他們.
\newline
\newline
三,摩西吩咐約書亞,從會眾選出能出戰之人,帶兵與亞瑪力人作戰.被選的人中,相信有些是不願去的,有些卻不被選上的,但有力多被選上的便勇往直前,與仇敵爭戰.
\newline
\newline
摩西與亞瑪力人爭戰,不是單靠約書亞,也不是單靠一二同工,乃號召全會眾,大批戰鬥員,投到戰場上去.這給我們看見,教會與罪惡爭戰,不是少數人的工作,是全教會信徒的責任.
\newline
\newline
當神選召我們,有人推諉逃避,說身體軟弱,但大多數願為神戰爭.今天教會有一好現象,很多教會特別訓練聖工人員,訓練信徒去作聖工.
\newline
\newline
人類社會是廣大的.人也多,三兩人作不成,必須全會眾大家同心合意,為神,為天國,為全人類的靈魂,投身作戰,作一戰鬥員吧!
\newline
\newline
摩西築了一座壇,起名叫「耶和華尼西」,意即「神是我的旌旗」.完全得勝,消滅仇敵.
\newline
\newline
諸位弟兄姊妹,今天是一個戰鬥的世代,讓我們都起來,為神爭戰,因為「耶和華是我們的旌旗」.
\newline
\newline


\newpage

\section{第40屆~港九培靈會~吳明節~第一講}
\label{sec:1532}
\textbf{逼迫逃難}
\newline
\newline
連結: \href{https://www.hkbibleconference.org/session-message/view/1532}{\texttt{https://www.hkbibleconference.org/session-message/view/1532}}
\newline
\newline
\hyperref[sec:1508]{< < < PREV SERMON < < <}
~
\hyperref[sec:toc]{[返目錄]}
~
\hyperref[sec:1531]{> > > NEXT SERMON > > >}
\newline
\newline
馬太福音 2:1-2:18
\newline
\begin{longtable}{cl}
\hline
\hline
章節 & 經文 (和合本修訂版)\\
\hline
2:1 & \begin{tabularx}{0.7\textwidth}{X} 在希律作王的時候，耶穌生在猶太的伯利恆。有幾個博學之士從東方來到耶路撒冷，說： \end{tabularx} \\ \\ \relax
2:2 & \begin{tabularx}{0.7\textwidth}{X} 「那生下來作猶太人之王的在哪裡？我們在東方看見他的星，特來拜他。」 \end{tabularx} \\ \\ \relax
2:3 & \begin{tabularx}{0.7\textwidth}{X} 希律王聽見了，就心裡不安；耶路撒冷全城的人也都不安。 \end{tabularx} \\ \\ \relax
2:4 & \begin{tabularx}{0.7\textwidth}{X} 他就召集了祭司長和民間的文士，問他們：「基督該生在哪裡？」 \end{tabularx} \\ \\ \relax
2:5 & \begin{tabularx}{0.7\textwidth}{X} 他們說：「在猶太的伯利恆。因為有先知記著： \end{tabularx} \\ \\ \relax
2:6 & \begin{tabularx}{0.7\textwidth}{X} 『猶大地的伯利恆啊，你在猶大諸城中並不是最小的；因為將來有一位統治者要從你那裡出來，牧養我以色列民。』」 \end{tabularx} \\ \\ \relax
2:7 & \begin{tabularx}{0.7\textwidth}{X} 於是，希律暗地裡召了博學之士來，查問那星是甚麼時候出現的， \end{tabularx} \\ \\ \relax
2:8 & \begin{tabularx}{0.7\textwidth}{X} 就派他們往伯利恆去，說：「你們去仔細尋訪那小孩子，找到了就來報信，我也好去拜他。」 \end{tabularx} \\ \\ \relax
2:9 & \begin{tabularx}{0.7\textwidth}{X} 他們聽了王的話就去了。忽然，在東方所看到的那顆星在前面引領他們，一直行到小孩子所在地方的上方就停住了。 \end{tabularx} \\ \\ \relax
2:10 & \begin{tabularx}{0.7\textwidth}{X} 他們看見那星，就非常歡喜； \end{tabularx} \\ \\ \relax
2:11 & \begin{tabularx}{0.7\textwidth}{X} 進了房子，看見小孩子和他母親馬利亞，就俯伏拜那小孩子，揭開寶盒，拿出黃金、乳香、沒藥，作為禮物獻給他。 \end{tabularx} \\ \\ \relax
2:12 & \begin{tabularx}{0.7\textwidth}{X} 因為在夢中得到主的指示，不要回去見希律，他們就從別的路回自己的家鄉去了。 \end{tabularx} \\ \\ \relax
2:13 & \begin{tabularx}{0.7\textwidth}{X} 他們走後，忽然主的使者在約瑟夢中向他顯現，說：「起來！帶著小孩子和他母親逃往埃及，住在那裡，等我的指示；因為希律要搜尋那小孩子來殺害他。」 \end{tabularx} \\ \\ \relax
2:14 & \begin{tabularx}{0.7\textwidth}{X} 約瑟就起來，連夜帶著小孩子和他母親往埃及去， \end{tabularx} \\ \\ \relax
2:15 & \begin{tabularx}{0.7\textwidth}{X} 住在那裡，直到希律死了。這是要應驗主藉先知所說的話：「我從埃及召我的兒子出來。」 \end{tabularx} \\ \\ \relax
2:16 & \begin{tabularx}{0.7\textwidth}{X} 希律見自己被博學之士愚弄，極其憤怒，差人將伯利恆城裡和四境所有的男孩，根據他向博學之士仔細查問到的時間，凡兩歲以內的，都殺盡了。 \end{tabularx} \\ \\ \relax
2:17 & \begin{tabularx}{0.7\textwidth}{X} 這就應驗了耶利米先知所說的話： \end{tabularx} \\ \\ \relax
2:18 & \begin{tabularx}{0.7\textwidth}{X} 「在拉瑪聽見號咷大哭的聲音，是拉結哭她兒女；她不肯受安慰，因為他們都不在了。」 \end{tabularx} \\ \\
[1ex]
\hline
\hline
\end{longtable}
世界上有幾個不容易解決的問題,無人能避免,我們又必須經歷.許多科學家,哲學家,教育家,他們都在設法決這些問題,結果都失望了.這些問題第一就是罪惡的問題,第二就是苦難的問題,第三就是死亡的問題.這些問題每天都圍繞著我們,我們必須對它有所認識.由於時間關係,只能簡單地研究人生苦難問題.
\newline
\newline
從馬太福音,可以看見耶穌基督所受的苦難.耶穌所受的苦難與每一位弟兄姊妹都有密切的關係,或都說,耶穌所受的苦難,與全世界的人類都有關係的.無論老少,貧富貴賤,每一個人都被苦難圍繞,不能逃脫.世人都喜歡快樂,不願意痛苦,所以每個人都要解決苦難問題.擺在我們前面的問題,一切的苦難,只有在耶穌基督裡面才獲得解決,因為在耶穌基督裡面,我們享受祂為我們成就的一切恩典.
\newline
\newline
從馬太福音二章一至十八節的經文中,看見耶穌基督嬰孩的時候就受逼迫逃難,逼迫祂的人就是當時的希律王,在這段經文中可以看見三種人,這三種人與耶穌逃難的事情有密切的關係:
\newline
\newline
第一種人──博士,「有幾個博士從東方來到耶路撒冷......」這裡說有幾個博士,我不知道到底有幾個?他們是聰明人,有學問,是知識份子的代表.他們從東方看見主的星,行走了很遠的路,來尋找主耶穌,找到之後,就俯伏下拜,並獻上黃金,乳香,沒藥.博士看見嬰孩耶穌之後,在夢中被主指示,不要回去見希律,就從別路回本地去了.這些博士他們隨著星尋求主,對主有好的信仰,在行動上已表明出來了,這是值得稱讚的.他們聰明,有學問,是無可否認的；他們雖然有長處,有敬畏主的心,又為主獻上禮物,但他們也有缺點的.這些博士雖然起初發現了主的星,是有神的啟示,但後來仍然靠自己的聰明智慧研究出來的.或者他們研究舊約聖經,知道這個人生下來是作猶太人的王,按照古時人的傳說,一顆星是代表一個人的,靠著自己的查考和推斷,隨著這星,來到耶路撒冷,就看不見那星了.他們便十分張揚地問:「那生下來作猶太人之王的在那裡；我們在東方看見他的星,特來拜他.」這問題傳開後,就發生了嚴重的後果了.猶太人和希律王都不安,以希律王懷恨在心,把伯利恆城裡和四境所有的兩歲以下的男孩都殺盡了.博士若不在耶路撒冷問這個問題,只要他們暗暗地尋找主,一定會找到的.
\newline
\newline
希律王聽見了東方博士查問關於耶穌降生在何處,就心裡不安,暗暗的召了博士來,他們在希律王的權柄地位下,而他們又是愛面子的,雖然希律王暗中召他們,他們沒有拒絕,還是去了.這些博士,可作為一切聰明人的代表,他們歸主之後,若肯謙卑下來,主是會大大使用他們的.相反地,他們仍然倚靠自己的聰明,智慧和才幹的話,無論他們怎樣肯為主作工,結果都是會引起教會的混亂與不妥的.教會中十分注重青年和小孩的工作,因為他們年輕,心靈純潔無邪,比較容易接受主的話.因此教會中時常行兒童佈道會,少年夏令會,青年退修會等等,許多愛主的弟兄姊妹為這項工作奉獻,參加兒童工作,但很少看見有學問的人,博士,參加這項工作的,甚至未見過有博士做主日學教員的；做主日學校教員的多半是青年團契的青年人,他們的學問雖然比不上博士,但他們的熱心是博士所不及的.
\newline
\newline
第二種人──希律王,希律是一切以自己為中心的人,他是王,權力都在他手中,當他聽見博士說:「那生下來作猶太人之王的在那裡；我們在東方看見他的星,特來拜他.」希律想,我既是王,將來豈不是奪取我的地位和權柄?我便不能唯我獨尊了.他雖然是王,但是十分愚蠢的,他召齊了祭司長和民間的文士來,問他們說:「基督當生在何處?」他們回答說:「在猶太的伯利恆；因為有先知記著說「猶大地的伯利恆啊,你在猶大諸城中,並不是最小的；因為將有一位君王,要從你那裡出來,牧養我以色列民.」這是舊約的預言,是彌迦書所講的,他們將這事告訴希律王.希律便暗暗的召了博士來,叫他們去尋找,尋到了,就來報信,我也好去拜祂.其實他並不是要去拜主,乃是他的詭計,想要殺害耶穌.結果將伯利恆城內及四境的兩歲以內的小孩都殺死了,他以為小孩都殺了,那生下來作王的那小孩也必死了.
\newline
\newline
希律既然相信神啟示先知所講的話,將來有一位救主要作猶太人之王的,祂要拯救這個世界.他雖然相信先知所講的話,但還是將小孩殺了,他沒有謙卑下來,他是以自我為中心的.在今天的教會中,也有許多信徒,他們並非不信上帝的啟示和先知所講的話；他們也信聖經,他們一切都信,但不能謙卑下來,離開他們的寶座,權柄,面子,地位,及自己的榮耀,所以他們還是作愚蠢的事,還是以自我為中心.求主幫助我們,巴不得我們都能像保羅一樣的說:「我活著是為基督.」
\newline
\newline
第三種人──百姓,就是昔日的猶太人,他們聽見一個生下來作猶太人之王的,就心裡不安.猶太人都清楚知道,這是舊約的先知所預言的.當他們失敗之後,神應許給他們一位救主,就是他們所仰望的彌賽亞.當時猶太人被羅馬統治,沒有自己的國家,所以他們仰望有一位救主,救他們脫離羅馬的統治.但是當他們聽見有一位新的王要來了,心中就非常不安,他們恐怕這位王來了,就必然會同羅馬帝國發生爭戰,這樣一來,就要妨礙他們安定的生活了.他們忘記了以賽亞書九章所說的話:「有一嬰孩為我們而生,有一子賜給我們,政權必擔在他的肩頭上；他的名稱為奇妙,策士,全能的神,永在的父,和平的君.他的政權與平安必加增無窮；他必在大衛的寶座上治理他的國......!」先知彌迦說:這位王來了,是要牧養我以色列民.祂來了是要復興以色列國,並不是引起爭戰的,因為祂是和平之君.
\newline
\newline
猶太人被羅馬統治,羅馬帝國統治的手段十分高明,給予猶太人宗教自由,他們可以在聖殿中自由敬拜,尊重他們的風俗習慣,所當時的法利賽人,祭司,文士,長老等都仍然被人尊重,仍然有相當地位.有時候羅馬王召他們來詢問.那些人都發了財,生活優裕；他們不願意因為有一位新的王來了,引起戰爭,影響他們安居樂業的生活,這是他們錯誤的想法.從歷史上來看,世界上最會賺錢的是猶太人,最愛錢的也是猶太人.在耶路撒冷的猶太人都發了財,所以聽見有新王興起,就心中不安了.昔日的猶太人可以作為今日貪愛世界之人的代表,在教會中,貪愛世界的信徒,若要他認罪悔改是很難的,除非主的靈進入他的心,抓緊他,不然,是不易悔改的.那些貪愛世界的信徒,當他們聽見神的僕人勸他們將錢財奉獻給神使用的時候,心中不安了,因為世界太可愛,太美麗了.昔日的猶太人聽見彌賽亞來了,心中不安了.今天的信徒聽見主再來了,是否滿心歡喜快樂地說:主啊!我願你來.若你捨不得這個可愛的世界,你聽見主再來了,心裡就不安了.若你是愛主的兒女,你必像聰明的童女一樣,儆醒等候主來.那麼,我們必歡喜快樂地回答說:「主耶穌啊,我願你來.」
\newline
\newline


\newpage

\section{第40屆~港九培靈會~吳明節~第二講}
\label{sec:1531}
\textbf{被魔試探}
\newline
\newline
連結: \href{https://www.hkbibleconference.org/session-message/view/1531}{\texttt{https://www.hkbibleconference.org/session-message/view/1531}}
\newline
\newline
\hyperref[sec:1532]{< < < PREV SERMON < < <}
~
\hyperref[sec:toc]{[返目錄]}
~
\hyperref[sec:1509]{> > > NEXT SERMON > > >}
\newline
\newline
馬太福音 4:1-4:11
\newline
\begin{longtable}{cl}
\hline
\hline
章節 & 經文 (和合本修訂版)\\
\hline
4:1 & \begin{tabularx}{0.7\textwidth}{X} 當時，耶穌被聖靈引到曠野，受魔鬼的試探。 \end{tabularx} \\ \\ \relax
4:2 & \begin{tabularx}{0.7\textwidth}{X} 他禁食四十晝夜，後來就餓了。 \end{tabularx} \\ \\ \relax
4:3 & \begin{tabularx}{0.7\textwidth}{X} 那試探者進前來對他說：「你若是神的兒子，叫這些石頭變成食物吧。」 \end{tabularx} \\ \\ \relax
4:4 & \begin{tabularx}{0.7\textwidth}{X} 耶穌卻回答說：「經上記著：『人活著，不是單靠食物，乃是靠神口裡所出的一切話。』」 \end{tabularx} \\ \\ \relax
4:5 & \begin{tabularx}{0.7\textwidth}{X} 魔鬼就帶他進了聖城，叫他站在聖殿頂上， \end{tabularx} \\ \\ \relax
4:6 & \begin{tabularx}{0.7\textwidth}{X} 對他說：「你若是神的兒子，就跳下去！因為經上記著：『主要為你命令他的使者，用手托住你，免得你的腳碰在石頭上。』」 \end{tabularx} \\ \\ \relax
4:7 & \begin{tabularx}{0.7\textwidth}{X} 耶穌對他說：「經上又記著：『不可試探主—你的神。』」 \end{tabularx} \\ \\ \relax
4:8 & \begin{tabularx}{0.7\textwidth}{X} 魔鬼又帶他上了一座很高的山，將世上的萬國和萬國的榮華都指給他看， \end{tabularx} \\ \\ \relax
4:9 & \begin{tabularx}{0.7\textwidth}{X} 對他說：「你若俯伏拜我，我就把這一切賜給你。」 \end{tabularx} \\ \\ \relax
4:10 & \begin{tabularx}{0.7\textwidth}{X} 耶穌說：「撒但，退去！因為經上記著：『要拜主—你的神，惟獨事奉他。』」 \end{tabularx} \\ \\ \relax
4:11 & \begin{tabularx}{0.7\textwidth}{X} 於是，魔鬼離開了耶穌，立刻有天使來伺候他。 \end{tabularx} \\ \\
[1ex]
\hline
\hline
\end{longtable}
在主耶穌遭受的苦難中,魔鬼的試探也是非常難受的,試探是魔鬼使用的最可怕的武器,千千萬萬的信徒曾在試探之下失敗跌倒；幾乎每一個信徒都有這個經驗.當我們忍受試探時,痛苦是非常大的,雖然我們藉著信心勝過它,但仍然每日都有試探臨到我們.主耶穌教導門徒禱告說:「......不叫我們遇見試探.」試探卻每時每刻跟著我們.我們不能避免試探,但可以藉著主耶穌勝過試探.主耶穌是神的──兒子,祂就是真神,也免不了被魔鬼試探,何況我們呢?
\newline
\newline
有人說:主耶穌是神的兒子,是神,在受試探之時當然不會跌倒了.這是不對的.有時候我們忽略了耶穌基督是真神；另一方面也忽略了魔鬼試探的厲害.有時候我們會這樣想:神是全能的,為甚麼不禁止魔鬼試探祂的兒女呢?自從魔鬼犯罪之後,神將它拘禁起來,在黑暗裡等候大審判.但神未將它放入硫磺火湖之前,魔鬼仍有它的權力.因此神准許它活動；准許它試探神的兒女.
\newline
\newline
受試探有兩個作用,第一個作用是要看看屬神的人是否忠心,是否愛祂.第二個作用是要叫魔鬼知道,神的兒女是可以勝試探的.耶穌基督是道成肉身住在我們中間.既是道成肉身,就與我們眾人一樣,所以必須遭遇試探.這是魔鬼對神的兒子的試探,我們承受了基督的救恩之後,成為神的兒女,也同樣受到神的兒女地位的試探.魔鬼說:「......你若是神的兒子......」耶穌是神的兒子,這是毫無懷疑的,但魔鬼說:「......你若是......」比方,我們是神的兒女,這是毫無疑問的,但魔鬼問我們是否神的兒女呢?你若軟弱,信心就動搖了.
\newline
\newline
(一)對神兒子生活的試探
\newline
\newline
耶穌禁食四十晝夜之後,後來就餓了,魔鬼來了,對祂說:「你若是神的兒子,可以吩咐這些石頭變成食物.」魔鬼的意思說:你看,你現在餓了,沒有東西吃,神是不會叫祂的兒子挨餓的,這裡沒有可吃的東西,你可以吩咐這些石頭變成餅,就可以解決飢餓的問題了.若不能把石頭變成餅充飢,就證明你不是神的兒子了,因為神的兒子是有權柄的.這試探十分厲害的.
\newline
\newline
各位弟兄姊妹:當我們生活困難,不能解決的時候,魔鬼就進前來說:「你不是神的兒女嗎?為何要挨餓呢?」我們就會說:「我做了神的兒女,還要挨餓,生活還是艱苦,為何還要做神的兒女呢?」就這樣為了生活,衣,食,住,行的問題,整日忙忙碌碌,為了生活,漸漸地對神冷淡了,連做禮拜都不去了,這樣神兒子的地位就動搖了.我們也遭遇主耶穌昔日的試探,但是那一個基督徒能把石頭變餅呢?我們沒有辦法,信心就動搖了,但不用怕,耶穌基督為我們受過試探,我們倚靠祂的得勝,成為我們今天的得勝.
\newline
\newline
耶穌是神的兒子,祂回答魔鬼說:「人活著不是單靠食物.......」「不」!絕對不聽人的說話,不管人的話說得如何好聽,理由如何充份,也不能聽人的話,總是「不」聽.記得我們是神的兒女,無論任何問題臨到我們,對神的信仰和神兒女的地位總不要動搖.最好的辦法就是「不理」,「不聽」.
\newline
\newline
神的兒子能否把石頭變餅?耶穌可以把水變酒；能叫死人復活；把石頭變餅並非難事,完全做得到,但祂不這樣做.因為把石頭變食物不是神的意思,在全部聖經中找不到神叫人把石頭變餅的話.由此可見石頭變餅不是神的意思.耶穌可以把石頭變餅,但魔鬼叫祂變,祂就不變了.在此,我們要特別注意一件事,若不是出於神的意思,任何生活困難用不正當的辦法來解決,都是出於魔鬼的試探,我們必須靠主勝過它.
\newline
\newline
人活在世上有更高的價值和意義,我們做神的兒女,活在世上不單為了吃喝快樂,是要遵行神的旨意,這是屬靈的事.把石頭變餅充飢,是屬肉體的事.屬靈的事比較屬肉體的事更有價值,雖然活在世上生活艱苦,但在靈裡面富足的,故此我們有喜樂.有些人忽略了這點,所以跌倒了.主耶穌十分清楚神的意思,祂堅持神的話:「人活著不是單靠食物.......」祂得勝了,這是偉大的得勝.若有人縱然挨餓,仍舊遵行神的話,不至動搖對神的信心.感謝主,他已經得勝了.
\newline
\newline
魔鬼叫耶穌把石頭變餅,雖然可變,但本質仍是石頭,若吃了石頭,一定發生問題.吃了石頭心就剛硬,變成灰心；心腸變硬了.一個信徒若用不正當的方法來解決生活問題,就是吃喝自己的罪,心腸就一天比一天剛硬,不會悔改了.所以我們要拒絕這件事.
\newline
\newline
(二)對神兒子宗教信仰的試探
\newline
\newline
魔鬼把耶穌帶到聖殿的頂上,對祂說:「你若是神的兒子,可以跳下去；因為經上記著說,「主為你吩咐祂的使者,用手托著你,免得你的腳碰在石頭上.」魔鬼試探耶穌,叫祂懷疑神的話是否可靠,聖經的話是否可靠,你是神的兒子嗎?你可以試試從殿頂跳下去,神會吩咐天使用手托著你的.這是宗教信仰的試探,這試探是非常可怕的.神的兒子,應進入聖殿獻祭和敬拜,但魔鬼叫祂站在殿頂上,跳下去,讓人看見.這樣一跳,必定頭破血流.神不會派天使托住祂的,因為不是出於神的意思,是出於魔鬼的意思.有些人上了魔鬼的當,在教會中沒有好好事奉神,偏要出風頭,行人所不能行的.我們必須知道,神的兒女是溫柔謙卑的,要與神有甜蜜交通的,不是要誇張叫人看見.神的兒子若進入聖殿,必是由石階一級一級上去；神的兒子是走路的,不會在殿頂跳樓的.若真的跳樓,神經就不正常了.
\newline
\newline
耶穌如果從殿頂跳下,必不會死,因祂是神的兒子.但祂不跳,因為神從來沒有叫祂跳樓.在信仰上,宗教上,跳樓絕非正道,故此主拒絕了.詩篇九十一篇論到一個人愛主,倚靠主的話,主要在他所行的路上保護他.一個人做了神的兒女,聽神的話,守神的誡命,神吩咐天使用手托住他,這是對的.但魔鬼濫用神的話,這是我們應該分別清楚的.對生活的試探,我們容易明白；最難明白的就是假先知,假傳道的話.他們同樣是講聖經,甚至是你從來未聽過的,他們也可以行神蹟,魔鬼也一樣會講聖經,但聖經的意思不是這樣的.許多人解經解得好聽,但意思是不對的.所以我們若不好好讀經,過屬靈的生活,對魔鬼的試探,就不容易得勝了.
\newline
\newline
(三)對神兒子權柄地位的試探
\newline
\newline
魔鬼又帶耶穌上了一座最高的山,將世上的萬國,與萬國的榮華,都指給祂看,對祂說:「你若俯伏拜我,我就把這一切都賜給你.」耶穌是神的兒子,是要承受神的工作和神的產業的,祂有崇高的地位與權柄,祂是天國產業的繼承者.魔鬼卻試探祂說:萬國與萬國的榮華都在我手中,我要給誰就給誰,只要拜我,我就把這一切都給你,你可以作世界的王,你作了王有最高的地位權柄,對人類不是有很大的貢獻嗎......主耶穌說:「撒旦退去罷；因為經上記著說,『要拜主你的神,單要事奉他.』」地位權柄不是最重要的,單單敬拜事奉神,才是最重要的.魔鬼失敗就離開耶穌了.
\newline
\newline
在這裡我們要思想一件事,魔鬼的試探是厲害的,多少弟兄姊妹在這試探中失敗跌倒,求主幫助我們勝過這試探.魔鬼試探我們的時候,也是要將萬國和萬國榮華給我們,若基督徒都從魔鬼手中得到最大的權柄,大家都作世界的王,勢必彼此攻擊,互相殘殺了.若中了魔鬼的詭計,在教會中必起紛爭了.
\newline
\newline
主耶穌得勝試探完全是引用聖經,靠神的話得勝.弟兄姊妹,我們為何不熱心虔誠研究聖經呢?世界上沒有甚麼書本比聖經寶貴,不讀聖經的人,常常會失敗的.
\newline
\newline
魔鬼失敗之後,耶穌勝過試探,有天使來伺候祂,祂開始作神的工作.所以,我們若不勝過試探,怎能作工呢?怎能傳福音呢?有些人為主作工沒有力量,是因為沒有天使與他同在,他和魔鬼沒有斷絕關係,所以無能力作工.希伯來書二章十八節說:「他自己既然被試探而受苦,就能搭救被試探的人.」希伯來書四章十五至十六節又說:「因我們的大祭司,並非不能體恤我們的軟弱；他也曾凡事受過試探,與我們一樣,只是他沒有犯罪.所以我們只管坦然無懼的,來到施恩的寶座前,為要得憐恤,蒙恩惠作隨時的幫助.」我們的大祭司耶穌,祂體恤我們的軟弱,祂凡事受過試探,而且都得勝了.我們只要倚靠祂,祂必幫助我們,使我們能勝過試探.
\newline
\newline


\newpage

\section{第40屆~港九培靈會~吳明節~第二講}
\label{sec:1509}
\textbf{失敗的掃羅}
\newline
\newline
連結: \href{https://www.hkbibleconference.org/session-message/view/1509}{\texttt{https://www.hkbibleconference.org/session-message/view/1509}}
\newline
\newline
\hyperref[sec:1531]{< < < PREV SERMON < < <}
~
\hyperref[sec:toc]{[返目錄]}
~
\hyperref[sec:1530]{> > > NEXT SERMON > > >}
\newline
\newline
撒母耳記上 28:15-28:19
\newline
\begin{longtable}{cl}
\hline
\hline
章節 & 經文 (和合本修訂版)\\
\hline
28:15 & \begin{tabularx}{0.7\textwidth}{X} 撒母耳對掃羅說：「你為甚麼攪擾我，招我上來呢？」掃羅說：「我十分為難，因為非利士人攻擊我，神離開我，不再藉先知或夢回答我。因此請你上來，好指示我應當怎樣做。」 \end{tabularx} \\ \\ \relax
28:16 & \begin{tabularx}{0.7\textwidth}{X} 撒母耳說：「耶和華已經離開你，與你為敵，你何必問我呢？ \end{tabularx} \\ \\ \relax
28:17 & \begin{tabularx}{0.7\textwidth}{X} 耶和華照他藉我所說的話為他自己實現了。耶和華已經從你手裡奪去國權，賜給別人，就是大衛。 \end{tabularx} \\ \\ \relax
28:18 & \begin{tabularx}{0.7\textwidth}{X} 因為你沒有聽從耶和華的話，沒有執行他對亞瑪力人的惱怒，所以今日耶和華向你做這事。 \end{tabularx} \\ \\ \relax
28:19 & \begin{tabularx}{0.7\textwidth}{X} 耶和華也必將你和以色列交在非利士人手裡。明日你和你兒子們必與我在一處了；耶和華也必將以色列的軍兵交在非利士人手裡。」 \end{tabularx} \\ \\
[1ex]
\hline
\hline
\end{longtable}
這段經文記載掃羅請女巫招撒母耳上來,今天未有時間講述關於女巫的事,但依據撒母耳所說的這一番話,便可知掃羅失敗的原因.
\newline
\newline
撒母耳所說的有下列三點原因:1. 你沒有聽從耶和華的命令.2. 你沒有滅絕亞瑪力人.3. 我已經從你手裡奪去國權,交給大衛；又將你的眾子交在非利士人手裡.
\newline
\newline
掃羅是以色列第一個君王,他被選立完全是出於神的恩典.以色列人要求撒母耳立王治理他們,撒母耳順從神的命令,掃羅便被膏立為以色列王,他應該是撒母耳的同工,是神的忠心僕人.但從聖經來看,掃羅失敗了,而且失敗得很慘.我們常常指這是對教會領袖一個教訓,其實對每個基督徒也是很有教訓,關係到每個基督徒的成功與失敗.
\newline
\newline
一,撒母耳記上十三章八至十四節的這段經文裡,說到當時非利士人來攻打以色列人,掃羅正在吉甲,撒母耳曾對他說過,過了七日便來獻祭,然後才出兵與非利士人交戰.那知到了第七天,撒母耳還未到,百姓因怕非利士人,便散去,且躲起來,掃羅那時也驚惶失措,竟然自己上去獻祭.獻祭剛完畢,撒母耳到了,問他說:「你作的是甚麼事呢?」掃羅說:「因為我見百姓離開我散去,你也不照所定的日期來到,而且非利士人聚集在密抹.因為我恐怕沒有禱告耶和華,非利士人下到吉甲攻擊我,我就勉強獻上燔祭.」撒母耳對掃羅說:「你作了糊塗事了,沒有遵守耶和華你神所吩咐你的命令,若遵守,耶和華必在以色列中堅立你的王位,直到永遠.現在你的王位必不長久,耶和華已經尋著一個合祂心意的人,立他作百姓的君,因為你沒有遵守耶和華所指示你的.」
\newline
\newline
撒母耳第七天便到,可見他沒有失約.掃羅的罪,是沒有聽從神的命令,獻祭是祭司的職份,是亞倫後代的職責.掃羅不是祭司,也不是亞倫的後代,無職責獻祭.獻祭是神藉祭司與人交通,是聖潔的,獻祭有許多規定,有許多條文,掃羅不是不明白,他獻祭表明不把神的命令放在眼中,他要藉自己的地位去奪祭司的職權,是很大的罪.
\newline
\newline
神設立祂的僕人,各有各的職份,各人應站在自己的本份之內,神沒有託付,不應去干涉別人職份內的事.掃羅應該懂得,可是他去做不應做的事,奪取神設立祭司的職份,違背神的命令.掃羅應作的本分,是作王的事情,不應獻祭；他離開王的地位,去作祭司的工作,這是大罪.
\newline
\newline
魔鬼從那裡而來?牠就是離開了本位的天使.祂造天使是為服事人的,但他離開本位,驕傲自大,就墮落成為魔鬼,神要用鐵鍊將他鎖住,等候大日的審判.
\newline
\newline
神將責任託付給我們,是要我們在本位上忠心,完成神的託付.
\newline
\newline
基督徒是神的兒子,應有做兒子的本份.守本份,忠心為神工作.今天教會有許多混亂的事發生,聖工人員有熱心,有忠誠,可是他們有高位,有大權柄,自以為能當神作工,無論大小事都去做,但一個人的精力和時間有限.而且神的工作,不是交給一個人去做的,乃是全會眾均有責任,任何人不能把一切工作由自己去做.一個牧師忠心做牧養工作,神已悅納；至於其他聖工人員在神所託付的範圍內,忠心工作,也蒙神悅納.每個弟兄姊妹,能表現愛教會,就忠心神所託付的工作,配合起來,教會一定興旺.如果你干涉我,我干涉你,甚至干涉牧師講道,規定牧師講道,如掃羅奪去祭司的權柄,自己去獻祭.求主保守眾弟兄姊妹,在教會守本份,忠心工作,使神得榮耀,教會得興旺.
\newline
\newline
二,掃羅第二個失敗的原因是記在撒母耳記上十五章一至卅一節.撒母耳告訴掃羅,耶和華差遣我膏你為以色列王,是治理他的百姓.以色列人出埃及的時候,怎樣與亞瑪力人爭戰,怎樣滅絕亞瑪力人,你也要去攻擊亞瑪力人,滅盡他們所有的,不可顧惜他們,無論人口,牛羊和一切牲畜都要盡行殺死.可是掃羅將亞瑪力打敗,殺了他們的百姓,卻將亞瑪力王留下不殺,且留下許多肥美的牛羊,不肯滅絕,還對撒母耳說:「耶和華的命令,我已遵守了.」撒母耳說:「我耳中聽見有羊叫,牛鳴,是從那裡來的呢?」掃羅還假惺惺作態說:「這是百姓從亞瑪力王那裡帶來的,因為他們愛惜上好的牛羊,要獻與耶和華你的神.」撒母耳說:「你住口吧,聽命勝於獻祭,順從勝於公羊的脂油.......你既厭棄耶和華的命令,耶和華也厭棄你作王.」掃羅不但想將責任推到百姓身上,且要求撒母耳和他一同回國去,好讓人看見他和大先知一起,便高抬他,擁護他,這是一種權術,且是一種詭詐,可見掃羅貪愛世界過於愛神,而且內裡不誠實.假如真是是百姓留下牲畜,他不去禁止,等於不遵守神的道,讓百姓去犯罪,這樣的領袖,是可悲的,神將百姓託付給他治理,他不禁止百姓去犯罪,是多麼失職.他留下亞瑪力王,一定是同情他,他濫用感情公然違抗神的命令,這是多麼愚拙的事情啊!
\newline
\newline
今天教會也接受了神的託付,將仇敵消滅,如果我們愛世界,愛財物,對神一定不忠心.有許多人寧願濫用感情,違背神的道.教會裡有許多有地位的人,知道信徒犯罪,為了感情,不但不去責備他,不勸勉他,不攔阻他,且表同情,實在放棄了神的命令.
\newline
\newline
在今天的教會裡,還有一更大的危險.就是隨大眾的意見,開會討論,便是民主.在政治上,社會上民主的作風是好的.但教會應該由神當家,不是信徒當家,就算全世界信徒的意見,也不能代表神,不能超越神的道.因為教會裡高抬民主作風,於是遠離了神的道.今天世界有許多大會,幾十個國家,幾千人聯合開會,一次又一次的討論,表面看來是好的,但在教會來看,是與神的道不同.追隨大眾意見,難免教會變為社會化,容易失去見證.求主保守,救我們脫離世界,神的主意是第一的,最高的,絕不可違犯.猶太人和官長都大聲說:「釘祂十字架,釘祂十字架.」這是眾意,難道主耶穌真是該釘十字架嗎?這事已經過去,但千萬不要讓這樣的事情再在教會裡發生.
\newline
\newline
三,掃羅第三個失敗的原因是記載在撒母耳記上十六章十四至廿三節.耶和華的靈離開掃羅,惡魔便降在掃羅身上.從前有神的靈在掃羅身上,他便是大能的勇士,是以色列的王,與列國爭戰,簡直是戰無不勝,攻無不克.但神的靈離開他,惡魔就降在他身上.在此情形下,他應立刻向神認罪悔改,求神收回祂的怒氣,這是掃羅唯一可走的道路.但他不如此作,卻叫大衛來彈琴,惡魔就離開他,這樣作也不錯,可是後來卻喜歡人稱讚他,而妒忌大衛,一次再次的想刺殺大衛.連兒子約拿單因與大衛友好也要殺.後來更找女巫,交鬼,問命運,其結果多麼可怕,對國家,對大臣,對兒女,對生活均錯亂了.故到撒母耳記上的末後一章,是記載他在戰場上失敗,為仇敵所殺,被梟首示眾,以色列第一個王,結果如斯,能不令人搖首歎息!
\newline
\newline
各位弟兄姊妹,當留心查察,在我們身上的是神的靈,抑是魔鬼的靈,心懷疑忌,生活便會錯亂.求主幫助我們,以掃羅的失敗,作我們前車之鑑.我們是神所選立的,將拯救靈魂的責任託付我們,我們當盡忠竭力,完成神的使命,求主幫助,請低頭禱告!
\newline
\newline


\newpage

\section{第40屆~港九培靈會~吳明節~第三講}
\label{sec:1530}
\textbf{被人厭棄}
\newline
\newline
連結: \href{https://www.hkbibleconference.org/session-message/view/1530}{\texttt{https://www.hkbibleconference.org/session-message/view/1530}}
\newline
\newline
\hyperref[sec:1509]{< < < PREV SERMON < < <}
~
\hyperref[sec:toc]{[返目錄]}
~
\hyperref[sec:1510]{> > > NEXT SERMON > > >}
\newline
\newline
馬太福音 13:53-13:58
\newline
\begin{longtable}{cl}
\hline
\hline
章節 & 經文 (和合本修訂版)\\
\hline
13:53 & \begin{tabularx}{0.7\textwidth}{X} 耶穌說完了這些比喻，就離開那裡， \end{tabularx} \\ \\ \relax
13:54 & \begin{tabularx}{0.7\textwidth}{X} 來到自己的家鄉，在會堂裡教導人，以致他們都很驚奇，說：「這人哪來這樣的智慧和異能呢？ \end{tabularx} \\ \\ \relax
13:55 & \begin{tabularx}{0.7\textwidth}{X} 這不是那木匠的兒子嗎？他母親不是叫馬利亞嗎？他兄弟們不是叫雅各、約瑟、西門、猶大嗎？ \end{tabularx} \\ \\ \relax
13:56 & \begin{tabularx}{0.7\textwidth}{X} 他姊妹們不是都在我們這裡嗎？他這一切是從哪裡來的呢？」 \end{tabularx} \\ \\ \relax
13:57 & \begin{tabularx}{0.7\textwidth}{X} 他們就厭棄他。耶穌對他們說：「先知除了在本鄉和自己的家之外，沒有不被尊敬的。」 \end{tabularx} \\ \\ \relax
13:58 & \begin{tabularx}{0.7\textwidth}{X} 耶穌因為他們不信，沒有在那裡行很多異能。 \end{tabularx} \\ \\
[1ex]
\hline
\hline
\end{longtable}
耶穌回到自己的家鄉,自己的人不接待祂.家鄉的人對祂認識非常清楚,耶穌的身世大家都明白,說祂是木匠的兒子,所以厭棄祂.耶穌不但被拿撒勒人厭棄；全猶太地的人都厭棄祂,當祂被彼拉多審判的時候,猶太人大聲疾呼說:「釘他十字架......」他們要求釋放殺人兇手巴拉巴,也不要耶穌.有時我們想,耶穌在本國的工作是否完全失敗?這是一個很大的奧祕.以弗所書三章四至六節說:「你們念了,就能曉得我深知基督的奧祕；這奧祕在以前的世代,沒有叫人知道,像如今藉著聖靈啟示他的聖使徒和先知一樣；這奧祕就是外邦人在基督耶穌裡,藉著福音,得以同為後嗣,同為一體,同蒙應許.」這奧祕是甚麼?就是外邦人與猶太人一樣得救.保羅在羅馬書十一章十一至十二節中,特別提到一個奧祕,就是猶太人厭棄耶穌,救恩就臨到外邦.外邦人在神的恩約上是無分的,在猶太人看來,神是以色列人的神,不是外邦人的神.舊約的先知,好像拿著望遠鏡在高山上觀望,看見耶穌基督道成肉身,降生在伯利恆；看見耶穌稱為拿撒勒人,又被本族的人厭棄；釘十字架,又從死裡復活；將來還要再來.關於耶穌之事看得十分清楚,但有一件事沒有看見,就是神的奧祕；就是保羅所說的基督的奧祕.耶穌被釘十字架之後,又從死裡復活,活在世上四十天之久,向門徒顯現,並差遣聖靈來,建立教會,這是神沒有啟示先知的.也就是說因猶太人厭棄耶穌,救恩就臨到外邦.外邦人組織了教會,教會就成了基督的身體.
\newline
\newline
今天,我們能坐在禮拜堂聽福音,都是因為猶太人將耶穌厭棄,而我們卻得著了.保羅說:我們是野橄欖樹,接在真橄欖上.因猶太人厭棄耶穌,使外邦人蒙恩得救.由此看見,神的恩典何等大.耶穌被猶太人厭棄,他們看不起耶穌,這種情形照人看來對耶穌不利的；神卻藉著這件事把救恩普及世界.我們若對耶穌有厭棄的心,神的恩典就臨到另外的人了.你若不接受主耶穌的呼召,耶穌會另外選召別人代替你的工作.
\newline
\newline
猶太人為何厭棄耶穌?有三方面值得思想的:
\newline
\newline
一,猶太人厭棄耶穌,只從耶穌是人那方面看,沒有看耶穌是神那一方面──若看耶穌是人那一方面,一定會厭棄祂,因為他們如此看法,就會將耶穌和其他的救主一起比較；以為耶穌不過如老子,孔子,釋迦牟尼,穆罕默德一樣吧了.彼得,雅各和約翰跟耶穌上了高山,耶穌登山之後變了形像；忽然有摩西以利亞向他們顯現,同耶穌說話.彼得對耶穌說,我們在這裡真好,你若願意,我就在這裡搭三座棚,一座為你,一座為摩西,一座為以利亞.彼得看耶穌不過像摩西和以利亞一樣,所以忽然有一朵雲彩遮蓋他們,且有聲音從雲彩裡出來說:「這是我的愛子,我所喜悅的,你們要聽他.」耶穌是一位真人,又是真神,猶太人厭棄祂,只是看見祂是人那一方面,所以他們說:這不是木匠的兒子麼?他母親不是馬利亞麼?這人從那裡有這一切的事呢.以為他是一個普通的人,算不得甚麼,故此厭棄祂.這些硬心的人一方面是這麼說.另一方面又說:這人為何能行異能呢?他們看見耶穌所行的奇事,又聽見耶穌所說的話非常有智慧,歷史上所有的人物沒有一個能超過耶穌.但他們還是不信祂,厭棄祂.他們心中驕傲,這是十分可惜的事.
\newline
\newline
在基督徒當中,有些人非常愛耶穌,但有些人卻把耶穌看成和世人差不多,沒有把祂當作神.忽略了耶穌是神,這信仰就不純正了.我們應當相信祂是人又是真神,祂是人類生命的主,祂是掌管一切的主宰.我們雖然有困難,有苦難,有疾病,但算不得甚麼,我們有一位至高的神,耶穌基督是我們的救主,我們倚靠祂,就能勝過一切.
\newline
\newline
許多時候,我們沒有將耶穌當作真神,只把祂當作一個大人物來崇拜,這也是不純正的信仰.雅歌書第五章說:新婦睡了,良人回來,她說我脫了衣裳怎能再穿上呢?我洗了腳怎能再沾污呢?當他起來開門的時候,她的良人已轉身走了.她走到街上問耶路撒冷的女子,看見她的良人沒有.耶路撒冷的女子說:你的良人比別人的良人有何強處?女子說:我的良人白而且紅,超乎萬人之上.他的頭像至精的金子,他的頭髮厚密纍垂,黑如烏鴉......我的良人無人彼得上祂,祂全然可愛,這就是我們的主耶穌.當我們對人作見證的時候,只說到自己的好處,自己怎樣改變,忽略了講主的榮美,把主耶穌放在一邊,這是不對的.各位弟兄姊妹,我們不要忘記,我們的主是人又是神,祂是無人能比的,不然,我們為何要信祂?
\newline
\newline
二,猶太人厭棄耶穌,是因為他們傳統的觀念錯誤了──猶太人的觀念是尊重權柄和錢財的.在今天教會中也有同樣的事情,重視有錢的人,重視有地位的人.有錢有地位的人若加入教會,悔改信主,奉獻給神,一定使別人得著幫助,這是一件好的事情.另一方面,我們不能因此而厭棄貧窮人,或冷落貧窮的弟兄姊妹,這樣做是錯誤的,是不符合聖經的.求主保守我們的教會,不會有這種錯誤的事情發生.
\newline
\newline
當時猶太人輕看耶穌,不過是木匠的兒子,是一個窮小子,雖然說話有智慧,又能行異能,但到底沒有地位和權柄.故此,沒有尊重祂.今天的信徒,對耶穌不敢如此,但對待窮苦的弟兄就很難說了,要知道厭棄一個小弟兄,就是厭棄主耶穌,主不會喜悅這種事的.
\newline
\newline
三,猶太人厭棄耶穌,是因為他們不明白聖經──在舊約聖經中,預言到耶穌要來,他們的彌賽亞要來,祂來的時候,帶著天使駕雲降臨,坐在耶路撒冷的寶座,管轄列國.許多先知都提到彌賽亞來是何等威風.但是以賽亞書五十三章的說法就不同了:「他被藐視,被人厭棄,多受痛苦,常經憂患.他被藐視,好像被人掩面不看的一樣；我們也不尊重他......他被欺壓,在受苦的時候卻不開口；他像羊羔被牽到宰殺之地,又像羊在剪毛的人手下無聲......」猶太人沒有注意這些,只注意彌賽亞來的威嚴.但主耶穌降生在馬槽,並非這樣威風,所以猶太人厭棄他,也因為他門不明白聖經.
\newline
\newline
猶大衛何賣耶穌?就是因為他懷疑耶穌不是彌賽亞,同時也是不明白聖經,這樣,就將他跟從了三年的夫子出賣了.
\newline
\newline
弟兄姊妹,今天我們知道,耶穌基督是神,又是人.我們知道是從聖經中查考出來的,所以我們要天天讀聖經.昔日的猶太人,常常查考舊約,還不明白,如果我們一天不肯讀一章聖經,又怎能明白呢?
\newline
\newline
馬太福音廿三章卅七節說:「耶路撒冷啊,耶路撒冷啊,你常殺害先知,又用石頭打死那奉差遣到你這裡來的人；我多次願意召集你的兒女,好像母親把小雞聚集在翅膀底下,只是你們不願意.」弟兄姊妹,在我們中間若有人未真正接受耶穌基督的,請你聽聽耶穌的呼聲.猶太人不願意,不信,厭棄了耶穌,所以耶穌說:「你們的家要變成荒場......」主後七十年,羅馬將軍提多,帶兵攻入耶路撒冷,一百二十萬人被殺,城牆被毀,這是猶太人厭棄耶穌的罪所受的咒詛.耶穌愛我們,拯救我們,若我們不願意,就免不了災禍.我們的主已經伸開雙手說:「我是你們的避難所,山寨,我是你們的救主,到我這裡來吧!」各位朋友,願你現在就接受祂.
\newline
\newline


\newpage

\section{第40屆~港九培靈會~吳明節~第三講}
\label{sec:1510}
\textbf{撒勒法的寡婦}
\newline
\newline
連結: \href{https://www.hkbibleconference.org/session-message/view/1510}{\texttt{https://www.hkbibleconference.org/session-message/view/1510}}
\newline
\newline
\hyperref[sec:1530]{< < < PREV SERMON < < <}
~
\hyperref[sec:toc]{[返目錄]}
~
\hyperref[sec:1529]{> > > NEXT SERMON > > >}
\newline
\newline
在聖經裡面,我們見到神有許多同工,有出名的,也有不出名的；有些是大人物,有些是小人物.今天要看一位無名的小人物,就是供養先知以利亞的撒勒法寡婦.聖經沒有記載她的名字,她的身世我們也無法稽考,只知道她是個寡婦,與她的兒子相依過活.
\newline
\newline
當亞哈作以色列王的時候,與王后耶洗別行耶和華眼中看為惡的事,事奉敬拜巴力,把全國人民陷在罪裡,離棄真神,國中滿了強暴的事.當時先知以利亞奉差遣往見亞哈王說:「我指著所事奉永生耶和華以色列的神起誓,這幾年我若不禱告,必不降露不下雨.」神的話又臨到以利亞,吩咐他藏在基立溪旁,要喝那溪裡的水,並吩咐烏鴉在那裡供養他.直到溪水乾了,神又叫他起來,去投靠西頓撒勒法的寡婦.
\newline
\newline
以下我們要從這寡婦身上,看看她幾種美德.
\newline
\newline
(一)她以愛心接待神人
\newline
\newline
諺云:寡婦門前是非多.她是個寡婦,以一個寡婦的身份接待一個陌生的男子,是難免有閒言的.然而她認識以利亞是個聖潔的神人,而用愛心來接待他.聖經說:愛裡是沒有懼怕的,愛既完全,就把懼怕除去.(約壹四18)許多人不肯作愛心的事,總是因懼怕,怕人批評,怕人講閒話......,以至把愛心丟掉,應做的事,應幫助的人,反而沒有做到.她憑愛心接待先知,這是值得我們效法的.因為在人生的過程中,未必能夠一帆風順,我們是否為了傳統和習慣,而放棄了作愛心的事,幫助別人或受助於人呢?這寡婦如果不肯接待以利亞,她就不能與神同工,她的事蹟,也不能記錄於聖經上了.她雖然沒有名字記在聖經上,但她的名字卻已記在天上的生命冊上了.她憑愛心接待神人,所以能成為神的同工.不管環境如何,我們總不要忘記,用愛心接待客旅.亞伯拉罕因為接待客旅,在不知不覺間,他竟然接待了耶和華的使者.(來十三2)主也教訓我們凡接待小子中一個的,就是接待祂,太廿五章審判山羊和綿羊的事,也有這樣的教訓,因為凡作在小子一個的身上,就是做在主的身上了.我們說愛主,也當愛自己的弟兄,若說愛主而不愛弟兄,這愛是虛假的；人若不能愛看得見的弟兄,怎能愛看不見的主呢?寡婦所作的,甚為美好!因為她以愛心接待神人,就是接待主了.
\newline
\newline
(二)她以養生之物供養神人
\newline
\newline
在飢荒的日子裡,一把麵粉和一點油,也很寶貴.她母子兩人,預備吃了這僅有的糧食,以後只好等死!那時以利亞來向她求餅,在情理上實在是不對的.但她肯聽從神人的話,為他先作餅,把自己在患難中養生的僅有物也全所有奉獻了!真是難能可貴!她的奉獻與主所稱讚那奉獻兩個小錢的窮寡婦,情形是相同的,因為她們所奉獻的,不是有餘的,乃是連養生之物也獻上了.正如舊約中的燔祭,要把祭牲完全焚燒,我們今日也應當把我們的生命不保留地奉獻.五祭中的素祭,是表示信徒在生活中的奉獻.神鑑察我們不是我們奉獻的有多少,乃是在主面前,我們的存心如何?亞拿尼亞和撒非喇也奉獻,但他們卻私自留下幾分,存心欺哄聖靈,所以死在彼得腳前.摩西領以色列人出埃及,要求法老准許他們離開,法老跟他們討價還價,那時只准男人出去,卻要留下婦孺和牲畜；但摩西卻要全數出埃及,連一蹄也不留下.我們獻給神也要全數的歸祂,不能一點留給法老.
\newline
\newline
寡婦的一把麵粉和一點油,原來算不得甚麼,但首先奉獻,就成了神蹟了!不但養活了先知,連她母子兩人,也靠著這一點食物維持,直到天降雨來的日子.
\newline
\newline
如果我們肯奉獻,完全的奉獻,全所有奉獻；神決不叫我們餓死,祂必負責,因為祂是信實的神.我們為甚麼會窮呢?是否我們不肯奉獻,而攔阻了神施恩的手施恩給我們呢?
\newline
\newline
(三)她以信心將死了的兒子交託給神人
\newline
\newline
寡婦的兒子後來因病得甚重,氣如游絲.婦人就來見以利亞,對他說:「神人哪!我與你何干?你竟到我這裡來,使上帝想念我的罪,以致我的兒子死呢?」神人來到她家的時候,也許她總以為自己好,自己了不得,當她接待以利亞後,深深認識到以利亞實在是一位聖潔的神人,深感自己不配,不潔.
\newline
\newline
婦人的兒子忽然死去,也有神的美意,多少時候,一個愈愛神的人,遭受到的試煉也多.婦人只有這一個兒子,她不因兒子去世而發怨言,只安靜的對以利亞說:「神人啊!我與你何干?」我做錯了甚麼事,使神想起我的罪來呢?「我的罪」,她以為兒子的死是她的罪.神因你的聖潔,更想到我的罪.一個人如果在苦難中,能夠想念自己以往在神面前的罪,就必在苦難中獲益.
\newline
\newline
當彼得蒙召的時候,他看見那一網圈住了許多的魚,就驚慌起來,對主說:「主啊!離開我,我是個罪人.」後來他竟然蒙召作一個得人的漁夫.
\newline
\newline
這婦人是大有信心的,她信心的眼睛看見了神的作為,她不但承認自己的罪,而且以信心將兒子交給神人.主耶穌說:「你若能信,就必看見神的榮耀.」主又說:「當以利亞的時候,天閉塞了三年零六個月,遍地有大飢荒,那時,以色列中有許多寡婦,以利亞並沒有奉差遣往他們一個人那裡去；只奉差遣往西頓的撒勒法一個寡婦那裡去.」(路四25-26)她是大有信心的婦人,主提到她,記念她,因為她能在苦難中接待神人,與神同工,神就叫她的兒子,由死裡復活過來.
\newline
\newline


\newpage

\section{第40屆~港九培靈會~吳明節~第四講}
\label{sec:1529}
\textbf{感情的傷害}
\newline
\newline
連結: \href{https://www.hkbibleconference.org/session-message/view/1529}{\texttt{https://www.hkbibleconference.org/session-message/view/1529}}
\newline
\newline
\hyperref[sec:1510]{< < < PREV SERMON < < <}
~
\hyperref[sec:toc]{[返目錄]}
~
\hyperref[sec:1511]{> > > NEXT SERMON > > >}
\newline
\newline
馬太福音 10:2-10:2
\newline
\begin{longtable}{cl}
\hline
\hline
章節 & 經文 (和合本修訂版)\\
\hline
10:2 & \begin{tabularx}{0.7\textwidth}{X} 這十二使徒的名字如下：頭一個叫西門（又稱彼得），還有他弟弟安得烈，西庇太的兒子雅各和雅各的弟弟約翰， \end{tabularx} \\ \\
[1ex]
\hline
\hline
\end{longtable}
「這十二使徒的名,頭一個叫西門,又稱彼得......」這裡說到西門彼得是十二使徒中的頭一個,他在當時的教會中地位甚高,他有許多長處,非常尊敬主,按遺傳說彼得是倒釘十字架而死的.天主教說彼得是最大的使徒.彼得在五旬節之前是非常自信的,喜歡講話,在使徒中常發表代表性的話語.耶穌看著他,許多時候,耶穌只帶著彼得,雅各,約翰三個人.有一次,耶穌到了該撒利亞腓立比的境內,就問門徒說,人說我人子是誰.他們說,有人說是施洗約翰,有人說是以利亞,又有人說是耶利米或是先知裡的一位.耶穌說,你們說我是誰.西門彼得回答說,你是基督,是永生神的兒子.耶穌對他說,西門巴約拿,你是有福的(參太十六章).彼得對耶穌的看法,乃是在天上的父所指示的.主耶穌說:「......彼得,我要把我的教會建造在這磐石上；陰間的權柄不能勝過他.我要把天國的鑰匙給你,凡你在地上所捆綁的,在天上也要捆綁；凡你在地上釋放的,在天上也要釋放.」(太十六18-19)可見彼得在耶穌的眼中是一個很好的使徒.
\newline
\newline
當我們研究聖經的時候,看見彼得對耶穌的愛不是純潔的愛,乃是情感的愛.不信的人沒主耶穌的情感,只有人的情感,像丈夫愛妻子,妻子愛丈夫；父親愛兒子,兒子愛父親.但這種愛不是聖經所說的神的愛.基督徒愛主不是用自己原有的感情愛主的,乃是用天父的愛來愛主.一般人是先有感情然後才有愛,這種愛是很容易變的.比方一個青年男子,愛上一個女子,發生了情感,深深相愛,就要結婚了,若不能結婚,就會造成為情犧牲的慘劇.但是結婚之後,慢慢地情感就變動了,兩人的感情不能維持下去了,從前願意為它犧牲的愛也不存在了.這是世界愛,父親愛兒子,因為兒子是他所生的,若對兒子沒有感情,生下來就把他丟了.世界的愛是非常薄弱的,感情的維持是不能長久的,是自私自利的,所以很容易破裂的.感情破裂了,愛就沒有了.所以,我們所追求的愛不是這種愛.我們所追求的是高尚的愛,也就是基督的愛.當我們悔改之後,神的愛就住在我們心中,神的愛藉著聖靈住在我們裡面,正如以弗所書三章十七節所說:「基督因你們的信,住在你們心裡,叫他們的愛心有根有基.」這經文何等寶貴,同時也說明只有神的愛所產生的感情,才是真實的.這種愛才能愛別人,愛窮人,愛仇敵.在這裡有許多我們不認識的人,因為有天父和耶穌基督住在心中,有神的愛在心中.所以主內的弟兄姊妹是非常有感情的,這情感因為有主的愛在我們心中而發出的,這種愛是永不改變的.外邦人的愛是先有情感,然後才產生愛,基督徒是先有神的愛,然後才產生情感的.情感可以破裂,愛是不會破裂的.弟兄姊妹,我們為何不信耶穌?為何不感謝讚美主呢?
\newline
\newline
彼得是最好的基督徒,但在他未被聖靈充滿之前,他是一個情感作用的人.「那時候彼得進前來,對耶穌說,主啊,我弟兄得罪我,我要饒恕他幾次呢?到七次可以麼.」彼得以為饒恕弟兄七次已經夠了,但是耶穌說,不是七次,乃是七十個七次(無限的饒恕).有一次西庇太的兒子的母親,同他兩個兒子(雅各,約翰)上前來拜耶穌,求耶穌在祂的國裡,讓她的兒子一個坐在主的右邊,一個坐在主的左邊.其餘十個門徒聽見了,就惱恕他們兄弟二人(參太廿章).這十個門徒當中,有彼得在內,可見弟兄中有一次過錯,彼得就惱怒他,不原諒他們了.
\newline
\newline
耶穌和門徒吃完了最後的晚餐,就出來往橄欖山去.那時耶穌對他們說,今晚你們為我的緣故,都要跌倒.彼得說,眾人雖然為你的緣故跌倒,我卻永不跌倒.但當耶穌被捉拿的時候,門徒都離開了,彼得也逃跑了.
\newline
\newline
彼得對耶穌:「看哪,我們已經撇下所有的跟從你,將來我們要得甚麼呢?」(太十九27)彼得撇下的是甚麼呢?馬太福音四章廿節說:「他們就立刻捨了網,跟從了他.」他們不過撇下了漁網,不是撇下一切.耶穌被釘死之後,彼得對其他門徒說:「我打魚去.」他們說,我們也和你同去.由此見得,恐怕彼得的漁網都沒有撇下呢?或是已撇下了,又隨時可以拿回來用.可見彼得愛耶穌是有目的的,他愛耶穌不是從神而來的愛.彼得愛耶穌,是情感的愛,不是屬靈的愛.有三件事可以證明.
\newline
\newline
第一個證明──體貼人的意思,不是體貼神的意思
\newline
\newline
耶穌對祂的門徒說:祂必須上耶路撒冷去,受長老祭司文士長許多的苦,並且被殺,第三日復活.彼得就拉著他,勸他說:「主啊,萬不可如此,這事必不臨到你身上.」(太十六22)耶穌轉過來,對彼得說:「撒旦退我後邊去吧,你是絆我腳的；因為你不體貼神的意思,只體貼人的意思.」(太十六 23-24)你看彼得是何等的愛耶穌,但耶穌卻說,撒旦退我後邊去吧!這是人意思的愛,不單不是愛耶穌,反而絆耶穌的腳.所以,凡是體貼人意思的,不是體貼神意思的,都是絆耶穌的.神願意萬人得救,不願意有一人沉淪,世人都在魔鬼的手下受轄制,神要每一個祂的兒女享受自由的快樂.世人無辦法解決這難處,神要找一個人來代替世人受苦,所以叫祂的兒子來到世界,代替世人受苦,使世上的人都脫離這苦難.彼得愛耶穌,不願耶穌受苦.彼得體貼夫子的肉體,卻沒有看見世上的苦難.
\newline
\newline
我們若明白神的意思,就知道耶穌捨去生命,是為了拯救許多人的生命.因此,有不少宣教師,願意到荒蕪之地為主作工,傳揚福音,把生命獻上放在祭壇上.但是有不少基督徒,雖然加入教會已經不少日子了,他只會體貼自己,在外邦人面前一點見證都沒有.世人都是逃避痛苦的,只是體貼自己親愛的人,就不管神的心意如何了.
\newline
\newline
第二個明證──三次不認主.
\newline
\newline
耶穌對門徒說,今夜你們為我的緣故,都要跌倒；彼得說,眾人雖然為你的緣故跌倒,我卻永不跌倒.耶穌說,我實在告訴你,今夜雞叫以先,你要三次不認我.彼得說,我就是必須和你同死,也總不能不認你.耶穌愛彼得,彼得愛耶穌,這是真實的,但這愛不是從神而來的.所以當耶穌被捉拿之後,就三次不認主了.「彼得在外院子坐著,有一個使女前來說,你素來也是同那加利利人耶穌一夥的.彼得在眾人面前卻不承認,說:我不知道你說的是甚麼?既出去,到了門口,又有一個使女看見他,就對那裡的人說,這個人也是同拿撒勒人耶穌一夥的.彼得又不承認,並且起誓說,我不認識那個人.過了不多的時候,旁邊站著的人前來,對彼得說,你真是他們一黨的,你的口音把你露出來了.彼得就發咒起誓說,我不認得那個人,立時雞就叫了.」(太廿六69-74)你看,情感的愛變動得那麼快,這和猶大有多大分別?分別在彼得悔改了,而猶大卻沒有悔改.
\newline
\newline
第三個證明──動刀.
\newline
\newline
當耶穌被捉拿的時候,彼得伸手拔出刀來,將大祭司的僕人砍了一刀,削掉了他一個耳朵.耶穌對他說,收刀入鞘吧,凡動刀的,必死在刀下.彼得動刀因為他愛耶穌,他曾說過願意與主同死,現在有人來捉拿主耶穌,於是他就拿刀出來,和別人拚命了.看來這個門徒很不錯,但是他的行動完全是出於情感和血氣,不是從神而來的愛心.
\newline
\newline
各位,請問你有沒有血氣激動的經驗呢?耶穌對彼得說:「你想我不能求我父,現在為我差遣十二營多天使來麼.若是這樣,經上所說,事情必須如此的話,怎樣應驗呢?」(太廿六53-54)但是耶穌沒有要求,祂體貼神的意思,因為父要祂受苦,為了拯救世人.
\newline
\newline
我們時常求主保護,但是彼得用自己的刀來保護主,你說可笑不可笑.聖經說:「從來沒有人能抵擋真理.」我們應該靠真理保護,靠真理得勝才對.讓我們省察自己愛主的心,是出於情感的愛還是出於天父的愛.彼得經過五旬節之後,被聖靈充滿,一次講道便有三千人悔改.為何他有這麼大的能力?就是因為他認識了真理,明白了天父的愛的結果.
\newline
\newline
弟兄姊妹:讓我們用彼得的事作為鑑戒,提醒自己,用屬靈的愛去愛主,及愛弟兄姊妹.
\newline
\newline


\newpage

\section{第40屆~港九培靈會~吳明節~第四講}
\label{sec:1511}
\textbf{灰心絕望的先知以利亞}
\newline
\newline
連結: \href{https://www.hkbibleconference.org/session-message/view/1511}{\texttt{https://www.hkbibleconference.org/session-message/view/1511}}
\newline
\newline
\hyperref[sec:1529]{< < < PREV SERMON < < <}
~
\hyperref[sec:toc]{[返目錄]}
~
\hyperref[sec:1512]{> > > NEXT SERMON > > >}
\newline
\newline
列王記上 19:1-19:18
\newline
\begin{longtable}{cl}
\hline
\hline
章節 & 經文 (和合本修訂版)\\
\hline
19:1 & \begin{tabularx}{0.7\textwidth}{X} 亞哈把以利亞一切所做的和他用刀殺眾先知的事都告訴耶洗別。 \end{tabularx} \\ \\ \relax
19:2 & \begin{tabularx}{0.7\textwidth}{X} 耶洗別就派使者到以利亞那裡，說：「明日約這時候，我若不使你的性命像那些人的性命一樣，願神明重重懲罰我。」 \end{tabularx} \\ \\ \relax
19:3 & \begin{tabularx}{0.7\textwidth}{X} 以利亞害怕，就起來逃命，到了猶大的別是巴，把僕人留在那裡。 \end{tabularx} \\ \\ \relax
19:4 & \begin{tabularx}{0.7\textwidth}{X} 他自己在曠野走了一日的路程，來到一棵羅騰樹下，就坐在那裡求死，說：「耶和華啊，現在夠了！求你取我的性命吧，因為我不比我的祖先好。」 \end{tabularx} \\ \\ \relax
19:5 & \begin{tabularx}{0.7\textwidth}{X} 他躺在羅騰樹下睡著了。看哪，有一個天使拍他，對他說：「起來吃吧！」 \end{tabularx} \\ \\ \relax
19:6 & \begin{tabularx}{0.7\textwidth}{X} 他觀看，看哪，頭旁有燒熱的石頭烤的餅和一壺水，他就吃了喝了，又再躺下。 \end{tabularx} \\ \\ \relax
19:7 & \begin{tabularx}{0.7\textwidth}{X} 耶和華的使者回來，第二次拍他，說：「起來吃吧！因為你要走的路很遠。」 \end{tabularx} \\ \\ \relax
19:8 & \begin{tabularx}{0.7\textwidth}{X} 他就起來吃了喝了，仗著這飲食的力走了四十晝夜，到了神的山，就是何烈山。 \end{tabularx} \\ \\ \relax
19:9 & \begin{tabularx}{0.7\textwidth}{X} 他在那裡進了一個洞，在洞中過夜。看哪，耶和華的話臨到他，說：「以利亞，你在這裡做甚麼？」 \end{tabularx} \\ \\ \relax
19:10 & \begin{tabularx}{0.7\textwidth}{X} 他說：「我為耶和華－萬軍之神大發熱心，因為以色列人背棄了你的約，毀壞了你的壇，用刀殺了你的先知，只剩下我一人，他們還要追殺我。」 \end{tabularx} \\ \\ \relax
19:11 & \begin{tabularx}{0.7\textwidth}{X} 耶和華說：「你出來站在山上，在耶和華面前。」看哪，耶和華從那裡經過。在耶和華面前有烈風大作，山崩石裂，耶和華卻不在風中；風後有地震，耶和華也不在其中； \end{tabularx} \\ \\ \relax
19:12 & \begin{tabularx}{0.7\textwidth}{X} 地震後有火，耶和華也不在火中；火以後，有輕微細小的聲音。 \end{tabularx} \\ \\ \relax
19:13 & \begin{tabularx}{0.7\textwidth}{X} 以利亞聽見，就用外衣蒙臉，出來站在洞口。聽啊，有聲音向他說：「以利亞，你在這裡做甚麼？」 \end{tabularx} \\ \\ \relax
19:14 & \begin{tabularx}{0.7\textwidth}{X} 他說：「我為耶和華－萬軍之神大發熱心，因為以色列人背棄了你的約，毀壞了你的壇，用刀殺了你的先知，只剩下我一人，他們還要追殺我。」 \end{tabularx} \\ \\ \relax
19:15 & \begin{tabularx}{0.7\textwidth}{X} 耶和華對他說：「去吧，從原路回去，往大馬士革的曠野去。到了那裡，你要膏哈薛作亞蘭王， \end{tabularx} \\ \\ \relax
19:16 & \begin{tabularx}{0.7\textwidth}{X} 又膏寧示的孫子耶戶作以色列王，並膏亞伯‧米何拉人沙法的兒子以利沙作先知接續你。 \end{tabularx} \\ \\ \relax
19:17 & \begin{tabularx}{0.7\textwidth}{X} 將來逃過哈薛之刀的，必被耶戶所殺；逃過耶戶之刀的，必被以利沙所殺。 \end{tabularx} \\ \\ \relax
19:18 & \begin{tabularx}{0.7\textwidth}{X} 但我在以色列中留下七千人，是未曾向巴力屈膝，未曾親吻巴力的。」 \end{tabularx} \\ \\
[1ex]
\hline
\hline
\end{longtable}
以利亞是舊約時代的大先知,他曾做過轟轟烈烈的大事,他是大有信心的祈禱人,曾禱告求天閉塞不下雨三年零六個月,又曾使撒勒法寡婦的兒子復活,與巴力先知迦密山上決鬥.在許多方面,都能使人認識他是一位神所重用的先知.如今卻因為王后耶洗別一句恐嚇的話,而灰心絕望:耶洗別差遣人去見以利亞說:「明日約在這時候,我若不使你的性命像那些人的性命一樣,願神明重重降罰與我.」以利亞看見這光景,就起來逃命,到了猶大別是巴,將僕人留在那裡.自己在曠野走了一日的路程,來到一棵羅騰樹下,就坐在那裡求死,說:「耶和華啊!罷了,求你取我的性命,因為我不勝於我的列祖.」從以利亞的遭遇看來,可見人總有軟弱,當他的靈程到達高峰時,即會降下來；這也是每一個信徒必有的現象.當我們熱心服事主時,很容易會灰心跌倒!
\newline
\newline
(一)他把神忘記了
\newline
\newline
他忘記了神,所以怕人.當他離開迦密山後,耶洗別就差人去見他.以他殺了巴力先知,也必被耶洗別所殺.他聽見了心裡害怕,所以起來逃跑,因為跑路困乏無力,又飢又渴,以至身體軟弱,灰心絕望.以利亞為何會落到這樣的光景呢?正因為他只想到自己,卻忘記了神.他應該知道自己的身份和地位,神揀選他作先知,使用他與他同在,所以他就能作人所不能作的事.
\newline
\newline
神是萬有的主宰,世上一切有權,有勢,有地位的人,都不能與他相比.可惜以利亞忘記了神,他只見到惡劣的環境,皇后耶洗別要殺他.過去他曾勇敢地站在亞哈面前,亞哈說:「使以色列遭災的就是你麼?」以利亞說:「使以色列遭災的不是我,乃是你,和你父；因為你們離棄耶和華的誡命,去隨從巴力.」以利亞當時有這勇敢,因神的能力在他身上彰顯,他在迦密山能戰勝巴力先知,他能英勇無退縮,因為神與他同在,任何惡劣環境他都不怕!但看環境看人,他的信心就失去了!
\newline
\newline
當馬丁路得改教時,神將重大的任務託付他,他面對教皇和無數的主教,國王,當時他實在有點膽怯灰心.一日回家,見妻子穿裡孝袍,路得問,是誰死了?回答說:神死了!路得斥責妻子說,神是永活的,那會死.妻答:神既是永活的,你又何灰心喪志到此地步呢?路得聽後,恍然大悟,立即起而改教.
\newline
\newline
(二)他失去了愛同胞的心
\newline
\newline
神問以利亞說,你在此作甚麼,他答「我為萬軍之耶和華大發熱心,因為以色列人背棄了你的約,毀壞了你的壇,用刀殺了你的先知,只剩下我一個人,他們還要索我的命.」(王上十九10)他逃跑是因為耶洗別尋索他的命,但他卻把這責任歸咎在百姓身上.他失去了愛同胞的心,他的工作對象就沒有了.他灰心絕望,他逃跑,但不是跑到橄欖山,加略山,乃是跑到何烈山,躲在洞中.摩西曾在這山上接受律法,與神面對面.他也在此山等候神,看神如何對付他的同胞.摩西為愛同胞,求神赦免他們不信背逆製金牛犢的罪,在山上禁食四十晝夜禱告,甚至在生命冊上塗抹他的名,他也在所不惜.保羅為愛同胞,大有憂愁,甚至與基督分離,他也願意.以利亞既失去了愛同胞的心,還有甚麼工作可作呢?
\newline
\newline
(三)他以為只有自己為神孤軍奮鬥
\newline
\newline
他看不見別人也為神作工,以為只有自己為耶和華大發熱心.別人都不同情他,也看不起他.神卻說,我留下了七千人是未向巴力屈膝的.俄巴底亦曾救了一百個先知,把他們藏在洞裡.以利亞親身聽聞了,誰知此刻他竟然向神說:「只剩下我一人,他們還要尋索我的命.」
\newline
\newline
一個神的僕人,若只知有自己,是十分危險的,因為忘記了別人,更忘記了神.
\newline
\newline
不要以為自己在講台上了不起,卻應當知道在台下有不少支持你的人,為你代禱!否則你的工作就毫無能力,毫無效果.不但不能感動人,更毫無價值了!所以不要灰心失望,因為有不少人正為你代禱!
\newline
\newline
以利亞竟然失敗到如此光景
\newline
\newline
1. 他逃跑不再為耶和華戰鬥,他跑到無人之地,跑到神審判之地.
\newline
\newline
2. 他把僕人撇下,只顧自己逃命,不顧別人死活.
\newline
\newline
3. 他在羅騰樹下求死,連他的祈禱也成了犯罪.
\newline
\newline
千萬人面臨死亡,他卻逃命躲下來,何其自私.他懼怕神審問,又向神抱怨!如此的祈禱,也成為罪.神卻愛他,記念他,前來親近他,與他說話.山崩石碎,烈風大作,風後地震,地震後有火,神卻不在其中.以利亞到那裡去躲避神的面呢?耶和華用微小的聲音向他說話,他不願為神作工,神就另外揀選了以利沙接續他的職任,從此以利亞工作也就完了.但願我們看以利亞的失敗,作為鑑戒,從失敗的地方靠著主的恩典剛強起來!不要以為非你不可,因為神能夠興起別人來代替你!
\newline
\newline


\newpage

\section{第40屆~港九培靈會~吳明節~第五講}
\label{sec:1512}
\textbf{以利沙靈程的四個步驟}
\newline
\newline
連結: \href{https://www.hkbibleconference.org/session-message/view/1512}{\texttt{https://www.hkbibleconference.org/session-message/view/1512}}
\newline
\newline
\hyperref[sec:1511]{< < < PREV SERMON < < <}
~
\hyperref[sec:toc]{[返目錄]}
~
\hyperref[sec:1528]{> > > NEXT SERMON > > >}
\newline
\newline
列王記下 2:1-2:15
\newline
\begin{longtable}{cl}
\hline
\hline
章節 & 經文 (和合本修訂版)\\
\hline
2:1 & \begin{tabularx}{0.7\textwidth}{X} 耶和華要用旋風接以利亞升天的時候，以利亞與以利沙從吉甲往前行。 \end{tabularx} \\ \\ \relax
2:2 & \begin{tabularx}{0.7\textwidth}{X} 以利亞對以利沙說：「耶和華差遣我往伯特利去，你可以留在這裡。」以利沙說：「我指著永生的耶和華，又指著你的性命起誓，我必不離開你。」於是二人下到伯特利。 \end{tabularx} \\ \\ \relax
2:3 & \begin{tabularx}{0.7\textwidth}{X} 在伯特利的先知的門徒出來，到以利沙那裡，對他說：「耶和華今日要接你的師父離開你，你知不知道？」他說：「我知道，你們不要作聲。」 \end{tabularx} \\ \\ \relax
2:4 & \begin{tabularx}{0.7\textwidth}{X} 以利亞對以利沙說：「耶和華差遣我往耶利哥去，你可以留在這裡。」以利沙說：「我指著永生的耶和華，又指著你的性命起誓，我必不離開你。」於是二人到了耶利哥。 \end{tabularx} \\ \\ \relax
2:5 & \begin{tabularx}{0.7\textwidth}{X} 在耶利哥的先知的門徒來靠近以利沙，對他說：「耶和華今日要接你的師父離開你，你知不知道？」他說：「我知道，你們不要作聲。」 \end{tabularx} \\ \\ \relax
2:6 & \begin{tabularx}{0.7\textwidth}{X} 以利亞對以利沙說：「耶和華差遣我往約旦河去，你可以留在這裡。」以利沙說：「我指著永生的耶和華，又指著你的性命起誓，我必不離開你。」於是二人一同往前行。 \end{tabularx} \\ \\ \relax
2:7 & \begin{tabularx}{0.7\textwidth}{X} 有五十個先知的門徒同去，遠遠地站在他們對面；他們二人在約旦河邊站住。 \end{tabularx} \\ \\ \relax
2:8 & \begin{tabularx}{0.7\textwidth}{X} 以利亞捲起自己的外衣，用來打水，水就左右分開，二人走乾地過去。 \end{tabularx} \\ \\ \relax
2:9 & \begin{tabularx}{0.7\textwidth}{X} 過去之後，以利亞對以利沙說：「我未被接去離開你以前，你要我為你做甚麼，只管求。」以利沙說：「願感動你的靈雙倍感動我。」 \end{tabularx} \\ \\ \relax
2:10 & \begin{tabularx}{0.7\textwidth}{X} 以利亞說：「你求的是一件難事。我被接去離開你的時候，你若看見我，就必得著；若不然，就得不著了。」 \end{tabularx} \\ \\ \relax
2:11 & \begin{tabularx}{0.7\textwidth}{X} 他們邊走邊說話的時候，看哪，有火馬和火焰車出現，把二人隔開，以利亞就乘旋風升天去了。 \end{tabularx} \\ \\ \relax
2:12 & \begin{tabularx}{0.7\textwidth}{X} 以利沙看見，就呼叫說：「我父啊！我父啊！以色列的戰車騎兵啊！」以利沙不再看見他的時候，就把自己的衣服撕為兩片。 \end{tabularx} \\ \\ \relax
2:13 & \begin{tabularx}{0.7\textwidth}{X} 他拾起以利亞身上掉下來的外衣，回去站在約旦河邊。 \end{tabularx} \\ \\ \relax
2:14 & \begin{tabularx}{0.7\textwidth}{X} 他用以利亞身上掉下來的外衣打水，說：「耶和華－以利亞的神在哪裡呢？」打水之後，水也左右分開，以利沙就過去了。 \end{tabularx} \\ \\ \relax
2:15 & \begin{tabularx}{0.7\textwidth}{X} 在耶利哥的先知的門徒從對面看見他，說：「感動以利亞的靈臨到以利沙身上了。」他們就來迎接他，俯伏於地，向他下拜， \end{tabularx} \\ \\
[1ex]
\hline
\hline
\end{longtable}
我們看過以利亞的生活,他是神所重用的先知,但他也會灰心絕望,逃到何烈山躲起來,在羅騰樹下求死.在這種情況下,神就不能再用他了,神吩咐他膏以利沙代替他的職份.各位!一個神的工人,倘若灰心至極點,神就無法再使用他了；因為我們對於神的工作的態度,不能以為非我不行,沒有我就無人能作,須知神能興起你,也能興起別人!正如末底改對以斯帖說:「此時你若閉口不言,猶大人必從別處得解脫,蒙拯救,你和你父家,必然滅亡,焉知你得了王后的位份,不是為現今的機會麼?」(斯四14)所以我們無論遭遇甚麼,總不要灰心,神不能用灰心的人,神要我們剛強壯膽作大丈夫.在本段經文內,已見到以利沙起來代替了以利亞,這是今日為神工作者,一個嚴重的警告.
\newline
\newline
以利沙在未開始工作時,神對他要施以嚴格的訓練,我們要作基督精兵,若不經過嚴格訓練,就不能作基督的精兵.
\newline
\newline
現在看以利沙靈裡的四個步驟:
\newline
\newline
(一)要從吉甲前往
\newline
\newline
以利亞說我要離開吉甲,往伯特利去.舊約有三個吉甲,這裡所說的吉甲,在以法蓮地.以利亞與以利沙從吉甲前往.「吉甲」這名字是「輥」的意思.(書五9)約書亞要把埃及人給的羞辱輥出去,這裡所記叫我們清楚知道,以利亞帶領以利沙經聖靈裡的第一步,是要把一切罪惡羞辱除去.關於這一點,的確是我們聖工人員必須要對付清楚的.一個假悔改的人,不能得神使用,所以要把罪和羞辱從我們身上輥出去!有人容易灰心跌倒,因為罪還未從他身上輥出去,所以無法為神工作.腓三19提到有些人是十架的仇敵,他們以羞辱為榮耀,專以地上的事為念.我們若仍有羞辱的罪,就應當除去,方能奔跑靈裡.
\newline
\newline
(二)要到伯特利
\newline
\newline
以利亞對以利沙說:「耶和華差遣我往伯特利去,你可以在這裡等候.」以利沙說:「我指著永生的耶和華,又敢在你面前起誓,我必不離開你.」於是二人下到伯特利.按伯特利,乃是「神的殿,天的門」之意.(參創十八16-19)以利亞帶以利沙來到伯特利,是叫以利沙想起雅各曾在此地遇過神,也在此地與神交通.以利沙要服事神,為神工作,是不能單靠跟隨師傅的,他必須與神有交通,與神面對面；一個聖工人員,若不明白神的旨意,與神交通,支取神的能力,他的工作就毫無效果,只靠自己和倚靠人,是必然失敗的!故此以利亞要把以利沙帶到伯特利去.多少時候我們工作失敗,讀經沒有亮光,禱告沒有能力,生活不能得勝,就是因為我們未能與神有密切的交通.伯特利不特是雅各與神交通的所在,更是神應許賜福與雅各之地.神要與他同在,他無論往那裡去,神必保佑他,領他歸回這地,總不離棄他,直到成全向他所應許的.倘若我們能到伯特利,神也必使我們得著所應許之福.神若與我們同在,就無人能傷害我們,也沒有難題是祂不能為我們解決的.我軟弱,祂會扶持我；我憂傷,祂會安慰我......我們不但要從吉甲前往,更要來到伯特利──神的殿,天的門.
\newline
\newline
(三)往耶利哥去
\newline
\newline
這個城的名我們很熟識,耶利哥是有香氣之意,但它名不符實,因為實在是很臭.以利亞與以利沙同到耶利哥,在靈程上有特別的意義.耶利哥城在歷史上,有許多寶貴的事蹟.像以色列人,曾在到達耶利哥同守逾越節,為感恩紀念神拯救他們出埃及過紅海.他門抵達此地後時,吃到迦南的土產,嗎哪就止住了.他們把耶利哥城圍繞了十三次,城就不攻自陷了.以利亞為何要以利沙來到耶利哥呢?到耶利哥要吃乾糧,嗎哪是代表靈奶；到了耶利哥要打仗,必須吃乾糧,不能再吃奶了.作先知的,不僅是要離開罪惡,得神同工,更要吃乾糧,要長進,要經得起苦難.神呼召我們,要在救人的事工上,與祂同在.約書亞在耶利哥曾先遇到耶和華軍隊的元帥.他跟從主就能戰勝仇敵.我們若能緊緊的跟從主,忍受苦難,就必能打勝仗.正如保羅說:「感謝上帝,常帥領我們在基督裡誇勝,並藉著我們在各處顯揚那因認識基督而有的香氣.因為我們在上帝面前,無論在得救的人身上,或滅亡的人身上,都有基督馨香之氣,在這等人,就作了死的香氣叫他死；在那等人,就作了活的香氣叫他活.」(林後二14-16)
\newline
\newline
(四)往約但河去
\newline
\newline
以利沙要繼續以利亞的工作,他要求加倍的靈.他不求地位,金錢,尊貴,享受,乃求加倍的靈感動.他清楚知道,為神工作,不是倚靠勢力,不是倚靠才能,乃是倚靠耶和華的靈,方能成事.如果沒有神的靈感動,就甚都不能作.約但河的名字是水窪之地,下行之河.約但河處於較低之地勢,較地中海水平低一百十餘尺.以利沙求加倍的靈感動,他必須降卑自己,到低窪之地.他肯降卑自己,然後能得到加倍的靈感動,接受了以利亞的外衣,用以打水；約但河的水就分開,讓他過去.日後他所作的工作,遠勝於以利亞,他行了偉大的神蹟,歷代君王均稱他為父.因他得了加倍的靈充滿,也實在為神所重用!
\newline
\newline


\newpage

\section{第40屆~港九培靈會~吳明節~第五講}
\label{sec:1528}
\textbf{被釘十架}
\newline
\newline
連結: \href{https://www.hkbibleconference.org/session-message/view/1528}{\texttt{https://www.hkbibleconference.org/session-message/view/1528}}
\newline
\newline
\hyperref[sec:1512]{< < < PREV SERMON < < <}
~
\hyperref[sec:toc]{[返目錄]}
~
\hyperref[sec:1513]{> > > NEXT SERMON > > >}
\newline
\newline
馬太福音 27:33-27:34
\newline
\begin{longtable}{cl}
\hline
\hline
章節 & 經文 (和合本修訂版)\\
\hline
27:33 & \begin{tabularx}{0.7\textwidth}{X} 他們到了一個地方，名叫各各他，就是「髑髏地」。 \end{tabularx} \\ \\ \relax
27:34 & \begin{tabularx}{0.7\textwidth}{X} 士兵拿苦膽調和的酒給耶穌喝。他嘗了，不肯喝。 \end{tabularx} \\ \\
[1ex]
\hline
\hline
\end{longtable}
耶穌被釘十字架,這題目在教會中聽得很多,我們也常常唱耶穌被釘十架的詩歌.但是,耶穌被釘十字架,是一個奧祕的真理,並非每人都能懂的.保羅說:「因為十字架的道理,在那滅亡的人為愚拙；在我們得救的人卻為神的大能.」(林前一18)哥林多屬希臘國,希臘人的哲學思想是非常高超的,他們提倡真善美的哲學.故此保羅傳福音至哥林多的時候,他們說:神是真善美的總和,怎能把祂釘死在十字架上呢?保羅說:「我不知道別的,只知道耶穌基督並祂釘十字架.」我們常聽十字架的道理,不易明白.當耶穌被釘十字架的時候,那裡有兵丁,文士,長老,祭司等,許多人在譏誚祂,證明那些人都不懂耶穌釘十架的事.保羅說:「我們講的,乃是從前所隱藏,上帝奧祕的智慧,就是上帝在萬世之前,預定使我們得榮耀的.這智慧世上有權有位的人沒有一個知道的；他們若知道,就不把榮耀的主釘在十字架上了；」(林前二7-8)可見當時聰明人也不明白,魔鬼也不明白,耶穌是神的兒子,魔鬼曾試探祂,但失敗了.當耶穌在客西馬尼園祈禱的時候,把自己獻上,為全世界的人作贖罪祭.祂流汗如血,非常驚恐,祈禱一次,看看門徒都睡了.祂再去禱告說:「我父啊,這杯若不能離開我,必要我喝,就願你的旨意成全.」又來見他們睡著了,那時神就打發天使來加添他的力量.第三次禱告之後,羅馬的兵丁已經來捉拿祂了.祂若要逃避是來得及的,但神不喜歡祂逃避,因為神在祂身上有重要的任務,要祂釘死在十架上.魔鬼那時把一個意念放在祂心裡:「你是無罪的,為何要替這些罪人釘十字架呢?他們也不愛你.門徒都不理你了,仇敵來捉你了,你為何替他們死呢?」故此耶穌心裡驚恐害怕.神的兒子,無罪的耶穌,全能的主宰,也驚恐害怕.但祂得勝了,為代替全人類的罪,把自己交在神的手中,這道理魔鬼是不明白的.當時的人也不明白.
\newline
\newline
當時,被釘十字架的人,要把他的罪狀寫出來的,所以彼拉多要把耶穌的罪狀寫出來,上面寫著「猶太人的王」.當時猶太人被羅馬王統治,派彼拉多作總督,若是多了一個王,就是背叛羅馬王的統治.「猶太人的王」,這幾個字,是用三種文字寫的,拉丁文,希臘文和希伯來文.拉丁文是代表羅馬權勢；希臘文是代表哲學,智慧的權威；希伯來文是代表最高的宗教勢力.十字架上寫著「猶太人的王」,這罪狀是猶太人見證祂是猶太人的王.政治權勢,哲學思想,宗教權威都一同證明祂的罪狀.世上的人都憑著自己的思想來批評耶穌的不對,這是不合理的.世界的人看這是耶穌的罪狀,在耶穌的周圍有許多人譏誚祂,兵丁分祂的衣服,這是代表暴力的思想.有些過路的人也譏誚耶穌,這班人與耶穌有甚麼關係呢?這些人不同情便算了,為何還要笑祂呢?這是一切閒雜人的代表,他們吃飽了沒有甚麼正經事做,便批評教會,批評傳道,批評耶穌.今天,不信的人也在批評耶穌和教會,這完全是錯誤的.還有一種人就是官府,是代表政府的權勢,也在譏誚耶穌.有許多地方對信耶穌的人很不客氣,並加以逼害.還有一種人就是強盜,一個已經悔改了,另一個釘在十字架上仍然有心情來批評耶穌.有些人在世上正陷在罪的痛苦中,自己不悔改,還要譏誚基督徒.由此可見,我們所處的環境是十分困難的,罪惡是多麼可怕,基督徒所處的環境是多麼危險?當時圍著耶穌的人,也有不少已經悔改的了,並且為主作見證.耶穌在十字架上擔當世人的罪,祂還禱告求神赦免他們.
\newline
\newline
還有那些祭司長,文士,長老,是猶太人的中保,就是在聖殿中為犯罪的人獻祭贖罪的.這些人卻把耶穌交在外邦人手中,還站在旁邊,看祂釘十字架.這些人明白律法,有地位,每天查考聖經,但是不明白耶穌就是基督.這樣讀聖經有甚麼用呢?聖經說:「祭司長和文士並長老,也是這樣戲弄祂,說:祂救了別人,不能救自己.祂是以色列的王,現在可以從十字架上下來,我們就信祂.」(太廿七41-42)他們承認耶穌救了別人,但還要釘在十字架上.這些人真是壞到極處了.耶穌是猶太人的王,是猶大支派的獅子,大衛的後裔,但猶太人不明白.當時以色列已亡國,猶太國仍存留,這就是他們自己的王,這些人不信,心中詭詐到極.
\newline
\newline
(一)耶穌釘十字架成就了和睦
\newline
\newline
保羅說:「那時,你們與基督無關,在以色列國民以外,在所應許的諸約上是局外人,並且活在世上沒有指望,沒有上帝.你們從前遠離上帝的人,如今卻在基督耶穌裡,靠著他的血,已經得親近了.因他使我們和睦,將兩下合而為一,拆毀了中間隔斷的牆；而且以自己的身體,廢掉冤仇,就是那記在律法上的規條；為要將兩下,藉著自己造成一個新人,如此便成就了和睦；既在十字架上滅了冤仇,便藉著十字架,使兩下合為一體,與上帝和好了.」(弗二12-16)現今有些人說:那些不信的外邦人,也有他們的神,這是不對的.耶穌基督被釘十字架,廢掉了冤仇,成就和睦.人和神因為罪的緣故,成了仇敵；人與人之間也結了仇恨.直到現在,幾千年來,人不斷地彼此嫉忌,互相殘殺,是因為仇恨仍然存在之故.猶太人和外邦人不和睦,猶太人說,神是他們的,與你們外邦人無分,所以,人與人之間為了信神的緣故,彼此不和,互相攻擊.耶穌基督釘死在十字架上,使神與人和好,藉著十字架,人可以和神交通.今天,坐在我們中間有各省的人,有不同的習慣和不同的言語,但是十分融洽,一些也不隔膜,這是十字架的功用,彼此和睦了.
\newline
\newline
(二)耶穌的十字架消除了罪案
\newline
\newline
「又塗抹了在律例上所寫,攻擊我們有礙於我們的字據,把他撤去,釘在十字架上；既將一切的掌權的擄來,明顯給眾人看,就仗著十字架誇勝.」(西二 14-16)在我們的良心之中,也有不少罪案,將來審判的時候,都要在神的面前顯明出來,那是十分難受的,因為我們所犯的罪甚多.如今,耶穌基督釘十字架,成了罪案,祂的血已經塗抹了我們的罪.我們的罪雖多,但耶穌的血已遮蓋了,罪案已經消除了.十字架是一個武器,把我們的罪案消滅了,神再不記念我們過去的過犯了.
\newline
\newline
(三)耶穌的十字架消除了仇恨
\newline
\newline
「兒女既同有血肉之體,他也照樣成了血肉之體；特要藉著死,敗壞那掌死權的,就是魔鬼；並要釋放那些一生因怕死而為奴僕的人.」(來二14-15)死亡威脅著每一個人,包括肉體的死亡,靈魂的死亡,與將來永遠的死亡,這三種死亡都不能逃脫.因為死亡的權柄在魔鬼的手中.感謝主,祂死了,就把死權破壞了.耶穌是永遠不死的,我們相信祂,死亡已廢除了.耶穌將全世界人類的仇恨背在自己身上,釘死在十字架上,將永遠的生命賜給我們.我們雖然是罪人,但已得著釋放了.你若不信,就不能逃罪.
\newline
\newline
死是可怕的,何況是身死,靈死,永死呢.神的兒子,聖潔的主宰,無罪的主.為我們受了這一切,擔當了一切.使我們明白耶穌釘十字架的奧祕,神還要我們將這個十字架背起來,為全人類消除仇恨,獻上自己,你肯不肯呢?把十字架的真理在你的生命中表現出來.主耶穌說:「不背著他的十字架跟從我的,也不配作我的門徒.」(太十38)
\newline
\newline


\newpage

\section{第40屆~港九培靈會~吳明節~第六講}
\label{sec:1513}
\textbf{修建聖城的尼希米}
\newline
\newline
連結: \href{https://www.hkbibleconference.org/session-message/view/1513}{\texttt{https://www.hkbibleconference.org/session-message/view/1513}}
\newline
\newline
\hyperref[sec:1528]{< < < PREV SERMON < < <}
~
\hyperref[sec:toc]{[返目錄]}
~
\hyperref[sec:1527]{> > > NEXT SERMON > > >}
\newline
\newline
尼希米記 1:1-1:4
\newline
\begin{longtable}{cl}
\hline
\hline
章節 & 經文 (和合本修訂版)\\
\hline
1:1 & \begin{tabularx}{0.7\textwidth}{X} 哈迦利亞的兒子尼希米的言語如下：亞達薛西王二十年基斯流月，我在書珊城堡中。 \end{tabularx} \\ \\ \relax
1:2 & \begin{tabularx}{0.7\textwidth}{X} 那時，我有一個兄弟哈拿尼，同幾個人從猶大來。我問他們那些被擄歸回、剩下殘存的猶太人和耶路撒冷的情況。 \end{tabularx} \\ \\ \relax
1:3 & \begin{tabularx}{0.7\textwidth}{X} 他們對我說：「那些被擄歸回剩下的餘民在猶大省那裡遭大難，受凌辱；耶路撒冷的城牆被拆毀，城門被火焚燒。」 \end{tabularx} \\ \\ \relax
1:4 & \begin{tabularx}{0.7\textwidth}{X} 我聽見這話，就坐下哭泣，悲哀幾日，在天上的神面前禁食祈禱， \end{tabularx} \\ \\
[1ex]
\hline
\hline
\end{longtable}
以斯拉與尼希米這兩卷書,在聖經中是十分寶貴的；尤其在現代,我們身居海外,更易於聯想到尼希米當日的背景.他們被擄寄居外邦,生活不安定,思鄉和掛念祖國的心情是在所不免的.想到耶路撒冷,那裡有聖城,有神的聖殿,是他們宗教的中心,敬拜神的所在,亦為政治的中心.大衛曾在那裡作王,亦為民族集合的中心.耶路撒冷被佔被毀,實在是以色列民族的奇恥大辱,為了復國國家民族,在歷史上以色列人曾不斷與外邦人流血爭戰.耶路撒冷雖是個小城,但在以色列人心目中,卻非常看重.他們愛慕耶路撒冷,歌頌耶路撒冷,也哀慟耶路撒冷.從尼希米這本書裡,我們可以看出以色列人是如何英勇地,修建耶路撒冷聖城.
\newline
\newline
(一)坐下哭泣,悲哀祈禱(一1-4)
\newline
\newline
尼希米在書珊城的宮中,聽見那些被擄歸回的猶大人,在猶大省遭大能,受凌辱,並且耶路撒冷的城牆拆毀,城門被火焚燒；他就坐下哭泣,悲哀祈禱.按人情說,他有崇高的地位,生活優裕,應該無憂無慮才是.但他愛國心切,因為聽見同胞遭難,城牆被毀,故此悲哀不已.請問我們有沒有愛同胞的心,有沒有迫切愛人靈魂的心呢?願主的愛激勵我們,也能效法尼希米,有迫切愛國的心!
\newline
\newline
(二)我默禱天上的神,神施恩的手就幫助我(二1-9)
\newline
\newline
尼希米聽見同胞遭受患難,城牆被毀,他不但傷心難過,禁食祈禱,並且還見諸行動.他在王面前侍立,面帶愁容；王問他甚麼緣故,他就默禱天上的神.求主差遣他回到猶大去.人先有勇敢信心,神就在人心中動工；正如尼希米有了回國的決心,神就在亞達薛西王的心中感動他.並允許大河西的省長,准他經過,又發詔書,吩咐把王囿的木材給他使用.人若愛神,有心為神工作；神必步步帶領,處處引導.
\newline
\newline
(三)我們起來建造,奮勇作這善工(二17-18)
\newline
\newline
尼希米回到耶路撒冷後,他就起來號召百姓一同建造城牆.第三章記載他們分工合作,奮勇修建,結果五十二天就建成了.感謝神!叫我們明白到同心合力作工,是何等的重要.今日在神的家裡,也有許多工作正等待我們去做,當教會呼召我們起來作工時,我們能否同心合作,完成主的聖工哩!
\newline
\newline
(四)一手作工,一手拿兵器(四1-4,17-18,六1-9)
\newline
\newline
尼希米遭受強大的敵人參巴拉,多比雅,和亞拉伯人基善,以及其餘的仇敵所陷害和恐嚇.說他們謀反,要作王,意思要使尼希米手軟,以致工作不能完成.尼希米就禱告,求神堅固他的手.並且他還勉勵貴冑,官長,和其餘的人,不要怕他們,當記念主是大而可畏的,你們要為弟兄,兒女,妻子,家產爭戰；他與官長都站在百姓的後邊,使百姓一手作工,一手拿兵器；修造的人,都腰間佩刀修造,吹角的人在他旁邊.作領袖的人,如能做百姓的後盾,百姓就得鼓勵,奮勇作工.尼希米還有一個應付仇敵的好戰略,就是不求和,不妥協.一面抵抗,一面建造,趕緊的造,步步為營的造,所以在五十二天內就完工了.弟兄姊妹!我們要抵擋仇敵魔鬼撒旦,也必須教會內長老,執事,同工,會友,大家一齊同心倚靠主,就必然能勝過牠了!
\newline
\newline
(五)像我這樣的人豈要逃跑呢(六10-14)
\newline
\newline
尼希米不但有外面的仇敵,更有內部的仇敵,就是參巴拉,多比雅賄買了示瑪雅,要尼希米依從他們,叫他懼怕停工.尼希米卻說:像我這樣的人,豈要逃跑呢?豈能進入殿裡保全命呢?他不伯死,也不停工,他清楚知道他們的話不是從神來的.他不灰心,也不絕望!當仇敵見到他們的工程迅速完成便懼怕,愁眉不展.因為這工作完成,乃是出於神.假使我們發覺教會內部有動亂,我們也應當有堅定不移的信心去勝過!
\newline
\newline
(六)心裡籌劃,盡力贖回弟兄(五1-12)
\newline
\newline
猶大人回國的百姓,貧苦者居多,他們為要得糧度命,向官長求助.官長就要他們典賣田園,使兒女作僕婢.因此百姓和他們的妻子,大大呼號,埋怨弟兄不該如此.尼希米就挺身起來,斥責貴冑和官長,召開大會攻擊他們.必須將人口和財物歸還百姓.
\newline
\newline
在他們的族譜上,凡尋不出自己名字的,都算為不潔,不准供祭司的職位.(七61-65)這些人就是閒雜人,混入進來的,要與他們絕交.
\newline
\newline
猶大人中有許多日久漸與外邦人通婚的,以至他們的兒女,一半說猶大的話,一半說各族的話.尼希米就斥責他們,拔去他們的鬍子,要他們棄絕外邦的女子.
\newline
\newline
在一個屬神的家庭內,若與外邦人通婚,必難以持久的,我們必須慎重從事.
\newline
\newline
(七)召開民眾大會,宣讀律法書(八1-8)
\newline
\newline
尼希米不但修造城牆,而且還要復興他們的宗教生活.他們對神的祭祀敬拜,已荒廢日久,宗教生活,已低落到極點,又與雜族人通婚.尼希米見這光景,就與文士以斯拉起來,召開民眾大會,宣讀神的律法書,吩咐百姓遵守,並舉行盛大的住棚節(八14)又叫百姓各按恩賜多多的奉獻.因為尼希米發覺許多人不肯奉獻,以至聖殿的銀庫空了,祭司生活不能維持,紛紛回到家園.他們在宗教生活上,實在破產了!故尼希米要起來開奮興大會,使他們覺悟悔改.今日港九教會的信徒,倘若不愛主,也不肯熱心奉獻,神的僕人就要受苦,願主的愛再激勵我們,肯樂意為主奉獻吧!
\newline
\newline
尼希米見許多人,不守安息日,他就斥責他們,並且吩咐在安息日要把城門關閉,不准出入；這樣就杜絕了那些在安息日趕著進城做買賣的人.以後商人和販賣貨物的,就不敢冒犯了.
\newline
\newline
尼希米能夠體會神的心意,他誠然是一個合乎神心意的好同工.
\newline
\newline


\newpage

\section{第40屆~港九培靈會~吳明節~第六講}
\label{sec:1527}
\textbf{常為弟兄而受苦}
\newline
\newline
連結: \href{https://www.hkbibleconference.org/session-message/view/1527}{\texttt{https://www.hkbibleconference.org/session-message/view/1527}}
\newline
\newline
\hyperref[sec:1513]{< < < PREV SERMON < < <}
~
\hyperref[sec:toc]{[返目錄]}
~
\hyperref[sec:1507]{> > > NEXT SERMON > > >}
\newline
\newline
馬太福音 25:41-25:46
\newline
\begin{longtable}{cl}
\hline
\hline
章節 & 經文 (和合本修訂版)\\
\hline
25:41 & \begin{tabularx}{0.7\textwidth}{X} 「王又要向那左邊的說：『你們這被詛咒的人，離開我！進入那為魔鬼和他的使者所預備的永火裡去！ \end{tabularx} \\ \\ \relax
25:42 & \begin{tabularx}{0.7\textwidth}{X} 因為我餓了，你們沒有給我吃；渴了，你們沒有給我喝； \end{tabularx} \\ \\ \relax
25:43 & \begin{tabularx}{0.7\textwidth}{X} 我流浪在外，你們沒有留我住；我赤身露體，你們沒有給我穿；我病了，我在監獄裡，你們沒有來看顧我。』 \end{tabularx} \\ \\ \relax
25:44 & \begin{tabularx}{0.7\textwidth}{X} 他們也要回答：『主啊，我們甚麼時候見你餓了，或渴了，或流浪在外，或赤身露體，或病了，或在監獄裡，沒有伺候你呢？』 \end{tabularx} \\ \\ \relax
25:45 & \begin{tabularx}{0.7\textwidth}{X} 王要回答：『我實在告訴你們，這些事你們沒有做在任何一個最小的弟兄身上，就是沒有做在我身上了。』 \end{tabularx} \\ \\ \relax
25:46 & \begin{tabularx}{0.7\textwidth}{X} 這些人要往永刑裡去；那些義人要往永生裡去。」 \end{tabularx} \\ \\
[1ex]
\hline
\hline
\end{longtable}
主耶穌說:「飢渴慕義的人有福了,因為他們必得飽足.」感謝主,我們若能在祂的面前領受祂的話,就得飽足.主所受的苦,不是為別人受的,乃是為我們每一個人受的.想起主的受苦,我們的信心就堅定,愛心就增長了.
\newline
\newline
主耶穌曾經說過一個比喻:「祂又要向那左邊的說,你們這被咒詛的人,離開我!進入那為魔鬼和他的使者所預備的永火裡去!因為我餓了,你們不給我吃；渴了,你們不給我喝；我作客旅,你們不留我住；我赤身露體,你們不給我穿；我病了,我在監裡,你們不來看顧我.他們也要回答說,主啊,我們什麼時候見你餓了,或渴了,或作客旅,或赤身露體,或病了,或在監裡,不伺候你吃.王回答說,我實在告訴你們,這些事你們既不作在我這弟兄中一個最小的身上,就是不作在我身上了.這些人要往永刑裡去；那些人要往永生裡去.」(太廿五41-46)我們知道耶穌所講的比喻,是要說明一件事實.這固比喻是說到將來的審判,王是代表他自己.耶穌說:「這些事你們若不作在我這弟兄中一個最小的身上,就是不作在我身上了.」我時常想,我們得救之後,作了神的兒女,神愛我們,神願意我們得著最好的,祂帶領我們一步一步前進.開始時,我們像嬰孩,吃奶,不會吃乾糧,也不能作甚麼.慢慢長大成為小孩,但仍十分幼稚,常常跌倒的.慢慢長大成人,像基督的身量.像保羅所說的「我活著就是基督」,那麼完全了.但我們不能達到這地步,主一步一步帶領我們,教我們如何向標準努力.每一個人都有自己所愛的,愛甚麼?愛的對象是甚麼,這值得研究了.我們生活在世上,每人所愛的不同,信徒與非信徒所愛的不同,讓我們看看,基督徒所愛的是甚麼?
\newline
\newline
(一)基督徒應愛世界
\newline
\newline
聖經不是說過「不要愛世界和世界上的事」麼?這是約翰一書所說的話.向下唸去就是:「就像肉體的情慾,眼目的情慾,並今生的驕傲,都不是從父來的,乃是從世界來的.」聖經叫我們不要愛世界,是不要愛世界上的情慾,犯罪的事.世界其實是可愛的,每一樣花草,飛鳥,樹木......都彰顯神創造之奇妙.主耶穌說:「你想野地裡的百合,怎麼長起來；他也不勞苦,也不紡線；然而我告訴你們,就是所羅門極榮華的時候,他所穿戴的,還不如這花一朵呢.」(太六 28-29)野地的百合,飛鳥,神也愛他們,也給他們裝飾.神所愛的一切東西,都是可愛.神不愛的是世界的罪惡,我們生活在世界上,工作在世界上,生命也在這世界上,怎能不關心世界的一切呢?世界實在是美麗的,每逢假日,青年人都聯群結隊到郊外旅行,或到海濱游泳,欣賞大自然的景色,因為神所創造的一切都是美好的.只因為戰爭摧殘這世界,這是罪惡所造成的,罪惡不是我們所愛的.世界無論怎樣好,但都要過去的,聖經說:這世界將來都要用火焚燒.世界的問題,根本是罪惡的問題,罪惡問題若不解決,世界是無法解決的.當耶穌坐在白色大寶座的時候,罪惡仍然存在；直到撒旦被投入硫磺火湖裡之時,罪仍存在,要想將世界的罪除去是不可能的.我們一方面在世上享受神所創造的自然,我們愛世界的一切事物,一切的人,要努力傳福音,叫世人脫離罪,勝過罪,不受罪的轄制,這就是愛世界,幫助這世界了.
\newline
\newline
(二)生命比世界更寶貴
\newline
\newline
耶穌說:「人若賺得全世界,賠上自己的生命,有什麼益處呢；人還能拿什麼換生命呢?」(太十六26)若將世界給你,你將生命給我,這樣是沒有人肯換的.所以,世界不是最貴重的.生命有兩部分,就是肉體的生命和靈魂的生命.肉體若沒有生命就死亡了,爛了；靈性若沒有生命就沉睡了.主耶穌說:「豈不知你們的身體就是聖靈的殿麼?」許多人說,信耶穌上天堂.不但是靈魂上天堂,肉體也上天堂.所以不能犯罪摧毀肉體的.有些人很容易自殺,豈不知我們的身體是耶穌用重價買來的麼?我們要在身體上榮耀神,就要有好的靈命,也應愛自己的身體,因為生命比世界更寶貴的.
\newline
\newline
(三)以耶穌為至寶
\newline
\newline
保羅說:「......我也將萬事當作有損的,因我以認識我主基督耶穌為至寶.我為他已經丟棄萬事,看作糞土,為要得著基督」(腓三8)比生命更寶貴的就是耶穌基督,耶穌是至寶,有了耶穌才能保守我們的生命；有了上帝的兒子基督就有生命.我們的生命是寶貴的,若將我們的生命和基督的生命比較,就要說選擇基督了.故此有許多人為基督的事工,捨棄自己的生命,為了主的福音而殉道.
\newline
\newline
有一位姊妹,三十多歲了,有一個孩子十分可愛.丈福音為信主的緣故,被官捉去殺死了.她還有一個父親,三個人相依為命.這位姊妹因為傳揚基督,被官捉去,用繩子縛住,吊在獅子坑的洞中.官兵對她說:「你還信耶穌麼?你只要說不信,我們就放你回去,你若信,我們就把繩子砍斷,掉你下去餵獅子.」這位姊妹說:「我以耶穌基督為至寶.」官兵又叫她的父親抱著她的孩子來,父親老淚縱橫,勸她否認主......姊妹說:「我愛你們,不願你們孤單,但是,殺身體不能殺靈魂的,我仍以耶穌基督為至寶,我不能否認主,你們回去吧!」最後,她終於被掉在獅子坑中,被獅子撕碎了.這是實在的事,這是真正的信仰,許多人為了主撇下了生命,故能叫千千萬萬人悔改,一代一代傳主的福音,藉著耶穌的血,一直傳到現在,我們基督徒才能享受主的恩典.生命誠然寶貴,但比不上基督寶貴.
\newline
\newline
(四)愛弟兄就是愛耶穌
\newline
\newline
「所以,你在祭壇上獻禮物的時候,若想起弟兄向你懷怨,就把禮物留在壇前,先去同弟兄和好,然後來獻禮物.」(太五23-24)這裡有一個寶貴的教訓,就是弟兄和睦.我們的生命固然寶貴,但比生命更寶貴的就是基督.弟兄姊妹也是寶貴的.為基督可以捨棄生命,為弟兄姊妹也可以捨棄生命.故此,最寶貴的是弟兄了.若弟兄跌倒,或喪失了生命,我們卻上天堂,他們要下地獄,我們不願意的.我們願意捨棄自己的生命,換取弟兄的生命,這是耶穌基督的真理.
\newline
\newline
保羅說:「我是大有憂愁,心裡時常傷痛；為我弟兄,我骨肉之親,就是自己被咒詛,與基督分離,我也願意.」(羅九2-3)保羅為耶穌基督已丟棄萬事,看作糞土.後來保羅為基督的緣故,在羅馬被殺.他一切都可以丟棄,不能捨棄耶穌.他又說,為了骨肉之親,弟兄姊妹得救,我保羅自己被咒詛,與基督分離也願意.只要同胞能悔改.這就說明了只有愛弟兄才是愛基督.我們有七萬萬同胞,七萬萬骨肉之親,他們敗壞,沉淪,不認識耶穌,未接受救恩.你有沒有為他們的靈魂心裡憂傷呢?
\newline
\newline
耶穌說:「因為我餓了,你們不給我吃；渴了,你們不給我喝；我作客旅,你們不留我住；我赤身露體,你們不給我穿；我病了,我在監裡,你們不來看顧我.......我實在告訴你們,這些事你們既不作在我這弟兄中一個最小的身上,就是不作在我身上了.這些人要往永刑裡去；那些義人要往永生裡去.」(太廿五 41-46)這段經文說明了真正的信徒,若愛主耶穌,就必須表現在愛弟兄姊妹上面.有不少弟兄姊妹,對於教會中貧窮的弟兄姊妹,或老病弱者,沒有好好幫助和看顧.我們若能愛那些弟兄姊妹,就是愛主了.這才是主所喜歡的.求主教導我們懂得怎樣愛主,因而愛我們的弟兄姊妹.
\newline
\newline


\newpage

\section{第40屆~港九培靈會~吳明節~第七講}
\label{sec:1507}
\textbf{背道的約拿}
\newline
\newline
連結: \href{https://www.hkbibleconference.org/session-message/view/1507}{\texttt{https://www.hkbibleconference.org/session-message/view/1507}}
\newline
\newline
\hyperref[sec:1527]{< < < PREV SERMON < < <}
~
\hyperref[sec:toc]{[返目錄]}
~
\hyperref[sec:1461]{> > > NEXT SERMON > > >}
\newline
\newline
約拿是個特別的先知,他的名字有「鴿子」之意.鴿子的性情是馴良的,可是他適得其反,列王記下十四章廿五節記載,他在以色列王耶羅波安第二在位的時候作先知,他也曾做過許多工作；可是從約拿書看來,他所做的工作實在令人失望!從他的工作與行為看來,可知他是不順從的先知,他是個背道逃跑的先知,在神面前是不蒙悅納的.主耶穌在馬太福音內曾提到他:「約拿三日三夜在大魚肚腹中,人子也要這樣三日三夜在地裡頭.當審判的時候,尼尼微人要起來定這世代的罪,因為尼尼微人聽了約拿所傳的,就悔改了,看哪!在這裡有一人比約拿更大.」(太十二40-41)
\newline
\newline
許多人讀到本書時,以為本書是寓意,不是事實,但你若留心看一章至四章,雖短短的數章書,就知道這不是寓意,乃是實實在在的歷史記載,而且每章都是真實的.
\newline
\newline
(一)神命令約拿去尼尼微(一1-3)
\newline
\newline
尼尼微城罪大惡極上達於天,神就差遣約拿往尼尼微去警告他們.尼尼微是在中東,約拿不去尼尼微,卻逃到他施去.他施即今日的西班牙,他不去東方卻逃至西方,這可看出他的背道行為來.為何他要躲避耶和華呢?四章二節說明他躲避的理由,他說:「耶和華啊!我在本國的時候,豈不是這樣說麼?我知道你是有恩典,有憐憫的神,不輕易發怒,有豐盛的慈愛,並且後悔不降所說的災,所以我急速逃往他施去.」他明白神有慈愛有恩典,也不輕易發怒,既知如此,此就應該好好地與神同工,聽從神的差遣,警告尼尼微人及早悔改,轉離所行的惡道.但他卻以為既有恩惠慈愛,且後悔不降所說的災,我又何必叫他們悔改呢?這表明他實在不明白神的旨意,不明白神的德性!
\newline
\newline
(二)神要追討約拿背道的罪(一4)
\newline
\newline
約拿躲避耶和華,逃往他施,他乘船航行,但耶和華使海中起大風浪,海就狂風大作,甚至船幾乎破壞.眾人就只好將船上的貨物拋在海中,為要使船輕些.神要呼召人,人若不聽從,那麼他所走的路,就必不亨通；他必受到神的攔阻.約拿此刻要借船逃跑,豈料全船的人就要因他而遭害.(一4-10)這時,全船的人都非常驚慌；因一個人的不順從,就連累到別人!後來船主發現約拿躲在底艙,躺臥,沉睡.於是船主叫他起來,求告他的神.這個先知,一點也不顧及別人......!這時船上的人,就向他發問了.
\newline
\newline
1. 你這沉睡的人哪!為何這樣呢?起來求告你的神,或者神顧念我們,使我們不至滅亡!約拿已知道,起這風浪是因他的緣故,所以他實在無言可答!
\newline
\newline
2. 來罷!我們製籤,看看這災禍臨到我們是因誰的緣故,於是他們掣籤,掣出約拿來!你看如果我們今日不為神工作,我們逃跑,要走自己的道路；倘若危險災禍臨到世人,這是誰的緣故呢?試想,我們有沒有責任呢?能否卸下我們的責任呢?我們是否貪愛世界,對別人的得救漠不關心呢?他們若死亡責任在誰身上呢?無怪保羅說:「我傳福音是不得已的,若不傳福音,我便有禍了.」(林前九16)
\newline
\newline
3. 你以何事為業?那些人這樣問約拿,約拿能不慚愧嗎?約拿是個先知,是個傳道人,神的話語信息,要藉著他作出口,但他竟然躲在艙底沉睡.他不能為神好好地工作,就必羞辱神的名了.比如我是個牧師,人若問我,你以何事為業,眼見群羊失散,神家荒涼,我將何以作答呀!
\newline
\newline
4. 你從那裡來?約拿是從以色列來的,但船是去他施的.他施不是神指示他要去的地方,到他施大可以發財,走自己的路,做自己的事.但卻不是神所差派要去的,如今約拿被船上的人所考問,他心中自然有說不盡的痛苦!
\newline
\newline
5. 你是那一國那一族的人?他清楚自己是以色列國的人,是神特選的民族,他應該隨時榮耀神!今日不信主的人,也會問我們,「你是那一國那一族的人?」我們如何作答呢?能否回答說:我是大天國的子民,我是屬乎神的呢?我們有沒有這美好的見證呀!
\newline
\newline
6. 我們當向你如何行?他們不想陷害約拿,卻向約拿如此發問,因為他們已知約拿是個先知,為了尊重,他只好讓他自己回答,好使海浪平靜,但當時海浪卻越發翻騰!約拿就對他們說:「你們將我抬起來,拋在海中,海就平靜了；我知道你們遭這大風,是因我的緣故.」約拿說這話的意思乃是,風浪既由我起,丟我下海,風浪即止,平安無事了!一個人倘若錯了,若能自覺,對整個團體,實在是有大大影響的.這時船上的人,只好把約拿拋在海裡,風的狂浪就平息了.耶和華安排了一條大魚吞了約拿,他就在魚腹中渡過了三日三夜.
\newline
\newline
(三)約拿在魚腹中悔改(二1-10)
\newline
\newline
約拿在魚腹中,四圍黑暗,又被海草纏繞著,魚肚又腥又臭,十分難過!這時他就開聲向神禱告,「求告耶和華他的神.」第二章全章都是他的禱文,從他的禱文中,可知當時他是如何的迫切,在神面前傾心吐意悔改,並向神許願.神就應允他的禱告,吩咐大魚把他吐在旱地上.
\newline
\newline
(四)尼尼微全城人畜都悔改
\newline
\newline
第三章記載耶和華的話,二次臨到約拿,叫他起來,往尼尼微大城去.向尼尼微城的人,大聲宣告說:「再等四十日,往尼尼微必傾覆了!」出乎意外的,尼尼微全城的人,聽見這話,他們都悔改了!尼尼微人信服神,從最大的到至小的,都穿麻衣,連王也下了寶座,脫下朝服,披上麻布,坐在灰中.王又命人遍告通城的人,禁食禱告,甚至牲畜也不吃草不喝水,在神面前自卑,切切求告神!各人回頭離開所行的惡道,丟棄手中的強暴.於是神鑑察他們的行為,見他們離開惡道,祂就後悔,不把所說的災禍降與他們了!
\newline
\newline
(五)約拿不合理的發怒
\newline
\newline
神赦免了尼尼微全城的人,這事約拿大大不喜悅,且甚發怒.那時他向神提出理由說:「我知道你是有恩典,大有憐憫的神,不輕易發怒,有豐盛的慈愛,並且後悔不降所說的災,所以我急速逃往他施去.」(四2)他以為早知神會後悔不降所說的災,故此不願去傳道.並且因此求死,以為死了比活著還好.神就問他,你這樣發怒,合乎理麼?
\newline
\newline
他出到城外,在東邊為自己搭了一座棚,坐在棚下,要看那城究竟如何?
\newline
\newline
這時神就安排一棵蔑麻,使它生長高過約拿,影兒遮蓋約拿的頭,救他脫離苦楚,約拿因此大大喜樂.但次日黎明,神卻安排了一條蟲子,咬這蔑麻,以致枯槁.日頭出來,神又安排炎熱的東風,日頭曝曬約拿的頭,使他發昏,他就為自己求死!說我死了比活著還好.
\newline
\newline
神所用的器皿,都順服神,就從神的命令吩咐.如:大海起風浪,大魚吞約拿,尼尼微王下寶座,外邦罪人悔改,一棵小蔑麻,一條小毛蟲,樣樣東西無不聽從神的話；惟獨約拿硬心硬到底!
\newline
\newline
藉著約拿的故事,求主憐憫我們,好叫我們知道人的靈魂比一切更寶貴!尼尼微城有十二萬多小孩的靈魂,神要愛惜,何況香港有成千成萬的孩童,他們的靈魂,豈非也是神所寶貴的嗎?讓我們都起來,作忠心傳福音的使者,不要作背道的先知約拿!(朱建磯記)
\newline
\newline


\newpage

\section{第40屆~港九培靈會~楊濬哲~序}
\label{sec:1461}
\textbf{序}
\newline
\newline
連結: \href{https://www.hkbibleconference.org/session-message/view/1461}{\texttt{https://www.hkbibleconference.org/session-message/view/1461}}
\newline
\newline
\hyperref[sec:1507]{< < < PREV SERMON < < <}
~
\hyperref[sec:toc]{[返目錄]}
~
\hyperref[sec:1463]{> > > NEXT SERMON > > >}
\newline
\newline
港九培靈會自一九二八年創辦迄今,匆匆已屆四十週年.回顧過去四十年中,除了日軍佔領時期外,大會均按年舉行.在這漫長的歲月中,上帝所給我們的恩典,真如以色列人出埃及,經曠野,有雲柱,有火柱,有天來糧食,有意外保護,使我們日日享受上帝的同在和更新的恩典.緬懷既往,只有俯伏敬拜,感謝讚美.
\newline
\newline
去年香港受到騷動影響,一時人心惶惶,但不旋踵,雨過天青,香港已回復安定,經濟再趨繁榮,今年大會又於安定中如期舉行,這一切都是上帝的大恩.
\newline
\newline
過去的四十年,已在主恩中平安過去,未來的日子如何,我們無法逆睹,但知昨日帶領我們的上帝,未來的日子仍然引我領我,因此我們只有舉起仰望的眼睛,求主憐憫,讓本會諸同仁,同心同德,藉著本會的工作,能夠對弟兄姊妹有供應,對教會有貢獻,務使佳果纍纍,榮神益人,有厚望焉.
\newline
\newline


\newpage

\section{第40屆~港九培靈會~鄭德音~港九培靈研經大會四十週年閉會詞}
\label{sec:1463}
\textbf{港九培靈研經大會四十週年閉會詞}
\newline
\newline
連結: \href{https://www.hkbibleconference.org/session-message/view/1463}{\texttt{https://www.hkbibleconference.org/session-message/view/1463}}
\newline
\newline
\hyperref[sec:1461]{< < < PREV SERMON < < <}
~
\hyperref[sec:toc]{[返目錄]}
~
\hyperref[sec:1515]{> > > NEXT SERMON > > >}
\newline
\newline
(一)
\newline
\newline
在路加福音講到不結果子的無花果樹,暫時留著它,加上土,多施肥料,或許會結果子,不然,才把它砍去.我們聽道,不結果子,是沒有用.甚至你聽了四十年,八十年的道也沒有用.在使徒行傳第七章裡講到腓利這個人,他是我們要學習的榜樣,他向路上的太監傳福音.
\newline
\newline
我們聽了道,還要去傳道.太監回去後他就傳道,今天亞比西尼亞仍然是信主的國家,乃是這個太監的功勞.
\newline
\newline
今天主須要教會起來傳福音,我們得了福音後,就要去傳,向家庭傳,向親友傳,向朋友傳.否則就是不結果子,不傳福音的無花果樹了.
\newline
\newline
若果主現在召你回去,你有甚麼果子?有多少個靈魂向主交代呢?
\newline
\newline
培靈會四十年過去了,聽了主的道,要效法太監謙卑聽福音,也效法太監忠心去傳福音,帶領更多的人信耶穌.
\newline
\newline
(二)
\newline
\newline
馬太福音第五章至第七章的三章聖經中,主親自主持一個培靈會,祂在山上講長篇的道理.我最後的臨別贈言是在馬太福音第七章十三節:「你們要防備假先知.」
\newline
\newline
聽道後,要照所聽的去行,才有永生,否則就要滅亡.永生永死兩種人的分別,在馬太福音七章中講得很清楚,聽了道去行就得永生,聽後不去行,就要滅亡.聖經又講到兩種樹,好樹與壞樹,好樹結好果子,壞樹結壞果子.主說,天國的道,照著去行,便可進入神的國,否則就不能.主說,惡人,離開我去吧!全備的道,愚人照著去行就要得智慧,聰明人不照著去行,也要變成愚拙了.照著主的道去行,他的人生是建立在磐石上,不照著去行,他的人生就是建立在沙土上.
\newline
\newline
各位聽了很多道,聽了之後,有無去行呢?主說,聽了道去行的,才是永生.主第二次的臨別贈言,祂說,要往普天下去傳福音,你們聽了道後,有無向家人,親戚,朋友,未認識的人傳福音呢?
\newline
\newline
主勉勵我們聽道後要行道,主三年工作後,向祂的同胞臨別贈言,不結果子的無花果樹,暫時保留,鬆土,加肥,到下一年若仍不結果子才把它砍下.這是主對選民的警告,同樣,也是對我們警告,不結果子也要被砍去.神對我們有一個目的:
\newline
\newline
一,神栽培我們,大家都是一樣的.
\newline
\newline
二,神放我們在教會中,神要用道栽培我們.
\newline
\newline
三,神要我們結出悔改的果子來,神要果子,我們有果子交給祂嗎?
\newline
\newline
四,神要你聽了道之後,立刻去行道.
\newline
\newline
五,神的忍耐並非是永遠的.
\newline
\newline
信主後一次得救就永遠得救,得救的信徒不行道,不結果子,在這世界上便失去見證,失去屬靈的福分,而且,將來在主臺前還要受審判.主要你聽道後,不但要行道,而且應該殷勤去行道.
\newline
\newline


\newpage

\section{第40屆~港九培靈會~鮑會園~第一講}
\label{sec:1515}
\textbf{神的話降臨}
\newline
\newline
連結: \href{https://www.hkbibleconference.org/session-message/view/1515}{\texttt{https://www.hkbibleconference.org/session-message/view/1515}}
\newline
\newline
\hyperref[sec:1463]{< < < PREV SERMON < < <}
~
\hyperref[sec:toc]{[返目錄]}
~
\hyperref[sec:1516]{> > > NEXT SERMON > > >}
\newline
\newline
何西阿書 1:1-1:1
\newline
\begin{longtable}{cl}
\hline
\hline
章節 & 經文 (和合本修訂版)\\
\hline
1:1 & \begin{tabularx}{0.7\textwidth}{X} 當烏西雅、約坦、亞哈斯、希西家作猶大王，約阿施的兒子耶羅波安作以色列王的時候，耶和華的話臨到備利的兒子何西阿。 \end{tabularx} \\ \\
[1ex]
\hline
\hline
\end{longtable}
當耶羅波安作以色列王的時候,神的話臨到何西阿,所以本書都是耶和華臨到的話.
\newline
\newline
耶羅波安作王是北國最強盛的時代,外無仇敵,內無困難,生活相當安定,人民生活水準一天一天的提高.過去雖然曾有亞述國入侵的事,但已忘記得一乾二淨.
\newline
\newline
因百姓生活安定,以致不再有聽神的話的心,故神藉何西阿賜下信息.百姓離開神,神特賜下祂的話來提醒百姓.雖然這是二千七百至二千八百年前的事,但拿來與現今時代比較一下,總覺得神藉何西阿所講的信息也是對我們所講的信息.我們一年前困難,似乎已經過去了,生活逐漸安定,不再用倚靠外力來幫助了,也漸漸不聽神的話了.
\newline
\newline
何西阿所傳的信息的特性
\newline
\newline
1. 這話一臨到何西阿身上,他立刻遵從,立刻去實行.神的話有權威,人應當遵從.祂從神權威的話,可以表示一個人對屬靈長進的態度.
\newline
\newline
舊約時代,對神的話,絕對順服,不論何事,神怎樣說,就怎樣遵從.
\newline
\newline
在民數記三十六章,記載神命摩西分配產業,西羅非哈沒有兒子,只有女兒,於是問題來了,他來到摩西面前,摩西也來到神面前,神說怎樣解決,百姓就怎樣解決,一點也沒有爭論.
\newline
\newline
今天教會有了問題,如何解決,不是到人的面前,問問是否合法,乃是要求問,看神的話如何說?
\newline
\newline
摩西年老,工作未完成,神怎樣安排新的領袖領會眾進迦南,不是叫人投票選舉,而是摩西到神面前來說:「願耶和華萬人之靈的神,立一個人治理會眾.可以在他們面前出入,也可以引導他們,免得耶和華的會眾如同沒有牧人的羊群一般.」於是耶和華對摩西說:「嫩的兒子約書亞,是心中有聖靈的,你將他領來按手在他頭上.」(民廿七15-18)於是會眾便照耶和華藉摩西所說的話,無人爭論.
\newline
\newline
亞伯拉罕得到神的賜福,是遵從神的命令把兒子獻上作為燔祭.看來有有點不合理,為甚麼將活人作祭牲.但亞伯拉罕不爭辯,以神的話為權威.今天若要神賜福,也是要遵從神的話,以神的話為權威.
\newline
\newline
2. 第二個性質,是要看神的話怎樣在何西阿身上出現.不是他心中忽然有感動,或是有了亮光,是神的話臨到他.他從未想過神的話會臨到,所以不是依靠環境,也不是依靠心中的感動,而是客觀的事實.若只是心中的感動,或被稱為亮光,無論這意思多好,也不能代替神的話.
\newline
\newline
神的話也是早賜給我們了,就是在聖經裡面.若是從神來的亮光,是神的感動,與聖經一定沒有矛盾.無論亮光多真確,感動多合理,若與聖經所說不符,就不是神的話.
\newline
\newline
3. 第三個特性.從神來的話,不是人研究的結果,不是人類文化的高峰.從神來的話,永遠不是從人做出發點.在神的恩典下,人類文化雖有進步,但人的話永遠不能代替神的地位.
\newline
\newline
神的話怎樣臨到何西阿
\newline
\newline
1. 神的話只臨到何西阿,不是臨到眾先知,也不是臨到以色列會眾.
\newline
\newline
神的話臨到人,要成為福氣,成為力量,成為幫助,但神的話是永遠臨到個人的.神的話不能臨到社會,因為社會,甚麼都是以大團體作出發點.神所看重的是每個人的生命,每個人的生活.神拯救的目標,也是每個人的生命.
\newline
\newline
就算我們明白神的話臨到何西阿有甚麼意義,或對耶羅波安有甚麼意義,也不能造就我的生命,我們要進一步明白神向我講甚麼?我們常常說,這篇道理講得好,但講得好與我有甚麼關係?比如你說這畫家畫的圖畫很好,這是畫家的成就,與我是沒有關係的.如果要我個人得著造就,就要說,神的話是臨到我,與我有關係.
\newline
\newline
我可以從聖經明白主是我們的力量,是我們的安慰,不錯,主耶穌的確是.當你有困難時,有痛苦時,感覺孤單時,主耶穌確是你的安慰,確是你的力量,與我的生命有直接的關係.
\newline
\newline
2. 神的話臨到我們身上,這「臨到」是一個非常特別的字,是很有力的字.
\newline
\newline
在撒母耳記上十六章和十九章記載掃羅違背神,神的靈離開他,惡魔便臨到他身上,他就失去了能力,不能管制自己的思想,也不能管制自己的行動,心靈被魔鬼掌管著,佔據著,需要有神的靈的能力的大衛來奏琴,為他驅去惡魔.神的話「臨到」何西阿,也是抓住何西阿的心,控制著整個何西阿,勝過何西阿的生活.
\newline
\newline
今天神的話有沒有臨到我?控制我的思想,勝過我的生活.
\newline
\newline
使徒行傳第二章,記載彼得講道時,人聽了覺得扎心,這是聖靈的能力,人不能不受感動,也不能不接受救恩.
\newline
\newline
耶利米說神的話臨到我,我便心裡覺得有燒著的火(耶廿9).神的能力勝過我,控制我,如火燒在我心裡.
\newline
\newline
主復活那天晚上,有兩個門徒從耶路撒冷往以馬忤斯,談論到舊約神的說話,心裡就覺得火熱.心被感動,被神的能力掌管,被神的能力勝過,從此,神的能力佔據了他們的心.
\newline
\newline
但神的話臨到何西阿,何西阿也付上了代價.對何西阿來說,這代價太大,太難,也太痛苦.照人的本性,實在太難接受.因為神要他娶淫婦為妻.這婦人曾使何西阿在生活上太痛苦,令何西阿太羞辱,但神要藉此作教導以色列人的工具.何西阿沒有與神分辯說:「神阿!我是被人尊敬的先知,怎可與這淫婦結婚!我一生的幸福就此斷喪嗎?我會因此失去見證,失去作工的能力.」他也沒有像約拿立刻逃走,雖然他心中或許起了極大的掙扎,但他沒有違背神的話,他順服了,從此成為神重用的人.神的話臨到我們,我們也許要付上代價.
\newline
\newline
或今天,或以後,神的話臨臨你,要你付上代價,才能得到他的祝福,你肯嗎?
\newline
\newline
神的話臨到有何結果?
\newline
\newline
1. 他領受了神的話,娶這淫婦為妻,被人棄絕,失去人的尊嚴,成為人譏笑的對象.今天你如接受神的話,也要付上這麼大的代價.
\newline
\newline
大有學問的人一樣地固執說,此書真是神的話嗎?把自己的生活放在這狹窄的信仰上,多麼愚蠢.但保羅說:「我不以福音為恥,這福音本是神的大能,要救一切相信的.」(羅一16)他知道向人作見證,會被人輕視,但他願領受神的話.
\newline
\newline
2. 成為神重用的器皿.在當時悖逆的世代,何西阿成為一盞明燈.把神的旨意傳開,使人認識神.神祝福他,使用他.今天我們若領受神的話,神也祝福我們,使用我們.
\newline
\newline
在人面前似乎沒有甚麼成就,但能完成神旨意,就是最大的成就.因為「人若賺得全世界,賠上自己的靈魂,有甚麼益處呢?」
\newline
\newline
然而當神的話臨到我們,我們怎樣面對神的話呢?
\newline
\newline


\newpage

\section{第40屆~港九培靈會~鮑會園~第二講}
\label{sec:1516}
\textbf{耶斯列的罪}
\newline
\newline
連結: \href{https://www.hkbibleconference.org/session-message/view/1516}{\texttt{https://www.hkbibleconference.org/session-message/view/1516}}
\newline
\newline
\hyperref[sec:1515]{< < < PREV SERMON < < <}
~
\hyperref[sec:toc]{[返目錄]}
~
\hyperref[sec:1517]{> > > NEXT SERMON > > >}
\newline
\newline
何西阿書 1:4-1:5
\newline
\begin{longtable}{cl}
\hline
\hline
章節 & 經文 (和合本修訂版)\\
\hline
1:4 & \begin{tabularx}{0.7\textwidth}{X} 耶和華對何西阿說：「給他起名叫耶斯列；因為再過片時，我要懲罰耶戶家在耶斯列流人血的罪，也必終結以色列家的王朝。 \end{tabularx} \\ \\ \relax
1:5 & \begin{tabularx}{0.7\textwidth}{X} 到那日，我必在耶斯列平原折斷以色列的弓。」 \end{tabularx} \\ \\
[1ex]
\hline
\hline
\end{longtable}
神的話臨到何西阿,要他娶淫婦為妻,他順從了.他與這淫婦結婚,是代表神與以色列人的關係.他們生的兒子,起名叫「耶斯列」.神說:「因為再過片時,我必討耶戶家在耶斯列殺人流血的罪.」(四節)這名字說明神要向祂的百姓施行審判.
\newline
\newline
神興起耶羅波安,給他十支派的以色列人,但他犯了大罪,帶領百姓違背神,拜金牛犢.後來當亞哈作了以色列王,他是一個大惡人,他不單獻祭與巴力,而且把獻祭與神的事,完全破壞,神便興起耶戶家來刑罰亞哈.在耶斯列殺了亞哈,也殺了耶洗別和他們的七十個兒女.耶戶殺了亞哈全家,是神熱心,執行神的旨意.
\newline
\newline
但神對他說:「我必討耶戶在耶斯列殺人流血的罪.」神既命令耶戶殺亞哈,為甚麼又刑罰耶戶?
\newline
\newline
聖經告訴我們,當耶戶作了以色列王後,也照樣領百姓去拜偶像.耶戶為了亞哈拜偶像去執行神的旨意,如今自己又領百姓拜偶像,更顯出耶戶的罪了.他既知拜偶像是罪,還去拜偶像,豈不是明知故犯嗎?神要刑罰他,有何不對?
\newline
\newline
耶戶幾方面的罪:
\newline
\newline
1. 他奉神的命令,審判亞哈拜偶像的罪,但他自己卻去犯此罪.他明知神憎恨此事,但他屬靈的生命卻是麻木的,忘卻了神的話,正是一個沒有以神的話與生活發生關係的人的鑑戒.
\newline
\newline
小孩子犯錯,母親處罰他,這不是仇恨,而是要孩子學習功課.我們犯罪,神已管教,但不久再犯,乃是沒有藉此刑罰來學習功課.
\newline
\newline
去年神要香港教會學習一件功課,但現今平安了,許多信徒便忘記了這個功課.以為這是偶然,與我們的罪無關.因為香港教會快要沉睡了,神藉去年的事來甦醒教會,我們不可以神的管教,作為不知道.
\newline
\newline
2. 神吩咐耶戶去審判亞哈,定亞哈為有罪.但他定別人的罪,卻不定自己的罪.別人犯罪要受刑罰,我犯罪不一定要受刑罰.定別人之罪易,定自己的罪難.主耶穌說:「先去掉自己眼中的梁木,然後才能看得清楚,去掉你弟兄眼中的刺.」(太七5)最好要定別人的罪以先,先定自己的罪.
\newline
\newline
3. 耶戶殺了亞哈全家,是神的命令.但表面上去服從還不夠,就是說表面上是做神的工作,裡面有自私的動機；做神的工,是為了另外的目的.主耶穌說:「當那日必有許多人對我說,主阿,主阿!我不是奉你的名傳道,奉你的名趕鬼,奉你的名行許多異能麼?我就明明的告訴他們說,我從來不認識你們,這些作惡的人,離開我去吧!」請三思此言.
\newline
\newline
拿答,亞比戶在神面獻上凡火,是耶和華沒有吩咐他們的,結果受神刑罰,死在神面前.
\newline
\newline
亞拿尼亞,撒非喇,變賣田產,卻私自留下幾分,也是被神刑罰而死.
\newline
\newline
看見別人奉獻,自己不奉獻,恐怕被人譏笑自己不愛主,不愛教會,不屬靈,於是我亦奉獻.可是這是在神的命令中,加上另外的目的,是要受神審判的.
\newline
\newline
巴蘭為神工作,卻一方面想得到巴勒的好處,這也是出於錯誤的動機.
\newline
\newline
有一位弟兄,工作能力好,表現得很熱心,曾得到好幾個學位.他起初是讀英國文學的,快得博士學位之時,轉讀神學.如今畢業十多年了,前兩年看見他時,仍是一個小職員,他雖是神學畢業的,但神不用他.可見做聖工的人,不可有錯誤的動機.
\newline
\newline
4. 耶戶奉神的命殺了亞哈全家,以為對自己有很大的好處.因為亞哈是王,殺了亞哈,自己便可以作王了.
\newline
\newline
做神的工,先衡量對自己有沒有好處.從對自己有沒有利益來做出發點,有利益便努力去做,沒有利益便馬虎了事.若順服神有利益,便順服,無利益便不順服.在屬靈的事上沒有目標,事奉神是為了自己的好處.
\newline
\newline
屬靈生活無目標,無標準,無原則,可有可無.知道應該愛主,應該追求,但要付代價,要對付自己,或是須等候主,要在人面前承認罪,這就太困難了.不是不想作,但實在不方便了.有好處,便可以追求；有方便,才有屬靈的表現.
\newline
\newline
我有一位友人想赴美謀生,但領取入境證很難,有人教他說去旅行,到了那裡,才作另外的打算.於是有人說:「說一句謊話不要緊,是為了方便呀!」但我們是否每遇一件困難的事,便說一次謊話呢?
\newline
\newline
有一位信徒,本是很愛主的.有一次他去某一國家,想帶一點違禁品,便使用黑錢來賄賂.這等於把見證收藏起來.
\newline
\newline
有人為了一點利益,隨時放棄基督徒的信仰,神不能不對付罪,神對耶戶怎樣,我們也不能例外.
\newline
\newline
無屬靈原則的生活,以利己為生活原則,其結果定是悲慘的.
\newline
\newline
過去神重用的僕人,後來屬靈生活上無原則,無標準,神便放棄他,像這樣的人,比比皆是.
\newline
\newline
今天在此有沒有這樣的信徒,過去熱心事奉神,愛神,肯為神犧牲,如今因有生活上的罪,神放棄了他,且受神的刑罰,過痛苦的生活,有何救法?
\newline
\newline
感謝神!耶戶的罪可怕,神審判已定,但神有豐盛的恩典.請聽神的話:「然而以色列的人數必如海沙,不可量,不可數,從前在甚麼地方對他們說,你們不是我的子民,將來在那裡必對他們說,你們是永生神的兒子.」(一10)以色列雖犯罪,但神仍向他們施恩,只要他們悔改,回到神面前.神已將福氣預備好,等人去領取,只要你肯脫離罪.耶戶在耶斯列的罪,顯明神的審判的可怕,亦藉此事顯明神的恩典,只要百姓肯離罪,神仍要施恩.
\newline
\newline
願神藉耶戶在耶斯列的罪,成為我們的鑑戒與蒙恩.
\newline
\newline


\newpage

\section{第40屆~港九培靈會~鮑會園~第三講}
\label{sec:1517}
\textbf{神的管教}
\newline
\newline
連結: \href{https://www.hkbibleconference.org/session-message/view/1517}{\texttt{https://www.hkbibleconference.org/session-message/view/1517}}
\newline
\newline
\hyperref[sec:1516]{< < < PREV SERMON < < <}
~
\hyperref[sec:toc]{[返目錄]}
~
\hyperref[sec:1518]{> > > NEXT SERMON > > >}
\newline
\newline
何西阿書 2:1-2:23
\newline
\begin{longtable}{cl}
\hline
\hline
章節 & 經文 (和合本修訂版)\\
\hline
2:1 & \begin{tabularx}{0.7\textwidth}{X} 你們要稱你們的眾弟兄為阿米，稱你們的眾姊妹為路哈瑪。 \end{tabularx} \\ \\ \relax
2:2 & \begin{tabularx}{0.7\textwidth}{X} 要跟你們的母親理論，理論，---因為她不是我的妻子，我也不是她的丈夫---叫她除掉臉上的淫相和胸間的淫態， \end{tabularx} \\ \\ \relax
2:3 & \begin{tabularx}{0.7\textwidth}{X} 免得我剝光她，使她赤身，如剛出生的時候一樣，使她如曠野，如乾旱之地，乾渴而死。 \end{tabularx} \\ \\ \relax
2:4 & \begin{tabularx}{0.7\textwidth}{X} 我必不憐憫她的兒女，因為他們是從淫亂生的兒女。 \end{tabularx} \\ \\ \relax
2:5 & \begin{tabularx}{0.7\textwidth}{X} 他們的母親行了淫亂，懷他們的做了可羞恥的事；因為她說：「我要跟隨我所愛的，我的餅、水、羊毛、麻、油、酒，都是他們給的。」 \end{tabularx} \\ \\ \relax
2:6 & \begin{tabularx}{0.7\textwidth}{X} 因此，看哪，我要用荊棘堵塞她的道，築牆擋住她，使她找不著路； \end{tabularx} \\ \\ \relax
2:7 & \begin{tabularx}{0.7\textwidth}{X} 以致她追隨所愛的人，卻追不上，尋找他們，卻尋不著，就說：「我要回到前夫那裡去，因我那時比現在還好。」 \end{tabularx} \\ \\ \relax
2:8 & \begin{tabularx}{0.7\textwidth}{X} 她不知道是我給她五穀、新酒和新的油，又加添她的金銀；他們卻用來供奉巴力。 \end{tabularx} \\ \\ \relax
2:9 & \begin{tabularx}{0.7\textwidth}{X} 因此，我要在收割的日子收回我的五穀，在當令的季節收回我的新酒，我要奪回她用以遮體的羊毛和麻。 \end{tabularx} \\ \\ \relax
2:10 & \begin{tabularx}{0.7\textwidth}{X} 如今我必在她所愛的人眼前顯露她的羞恥，無人能救她脫離我的手。 \end{tabularx} \\ \\ \relax
2:11 & \begin{tabularx}{0.7\textwidth}{X} 我必使她的宴樂、節期、初一、安息日，她一切的盛會都止息。 \end{tabularx} \\ \\ \relax
2:12 & \begin{tabularx}{0.7\textwidth}{X} 我要毀壞她的葡萄樹和無花果樹，就是她所說「我所愛的給我為賞賜」的；我要使它們變為荒林，為野地的走獸所吞吃。 \end{tabularx} \\ \\ \relax
2:13 & \begin{tabularx}{0.7\textwidth}{X} 我要懲罰她素日給諸巴力燒香的罪；那時她佩戴耳環和珠寶，跟隨她所愛的，卻忘記我。這是耶和華說的。 \end{tabularx} \\ \\ \relax
2:14 & \begin{tabularx}{0.7\textwidth}{X} 因此，看哪，我要誘導她，領她到曠野，我要說動她的心。 \end{tabularx} \\ \\ \relax
2:15 & \begin{tabularx}{0.7\textwidth}{X} 在那裡，我必賜她葡萄園，又賜她亞割谷作為指望的門。她必在那裡回應，像在年輕時從埃及地上來的時候一樣。 \end{tabularx} \\ \\ \relax
2:16 & \begin{tabularx}{0.7\textwidth}{X} 那日你必稱呼我伊施，不再稱呼我巴力。這是耶和華說的。 \end{tabularx} \\ \\ \relax
2:17 & \begin{tabularx}{0.7\textwidth}{X} 因為我必從她口中除掉諸巴力的名號，不再有人提這名號。 \end{tabularx} \\ \\ \relax
2:18 & \begin{tabularx}{0.7\textwidth}{X} 當那日，我必為我的百姓，與野地的走獸、天空的飛鳥和地上爬行的動物立約；又要在國中折斷弓和刀，止息戰爭，使他們安然躺臥。 \end{tabularx} \\ \\ \relax
2:19 & \begin{tabularx}{0.7\textwidth}{X} 我必聘你永遠歸我為妻，以公義、公平、慈愛、憐憫聘你歸我； \end{tabularx} \\ \\ \relax
2:20 & \begin{tabularx}{0.7\textwidth}{X} 也以信實聘你歸我，你就必認識耶和華。 \end{tabularx} \\ \\ \relax
2:21 & \begin{tabularx}{0.7\textwidth}{X} 耶和華說：那日我必應允，我必應允天，天必應允地， \end{tabularx} \\ \\ \relax
2:22 & \begin{tabularx}{0.7\textwidth}{X} 地必應允五穀、新酒和新的油；這些都必應允在耶斯列身上。 \end{tabularx} \\ \\ \relax
2:23 & \begin{tabularx}{0.7\textwidth}{X} 我為自己必將她種在這地。我必憐憫羅‧路哈瑪；對羅‧阿米說：「你是我的子民」；他必說：「我的神。」 \end{tabularx} \\ \\
[1ex]
\hline
\hline
\end{longtable}
以色列人不但犯罪離開神,且堅持著抵擋神.他們犯罪不是偶然的,而是經常的.經過神的提醒,勸勉,責備,但仍不肯歸給神.
\newline
\newline
「你們要與你們的母親大大的爭辯.」「爭辯」是一個很特別的字,可譯作「你們應繼續向母親苦苦的勸她.」換言之,神多次的勸勉他們,他們仍不聽從.
\newline
\newline
神對不肯歸從之百姓,仍繼續的勸勉,且苦苦的勸勉,為的是勸他們悔改.但剛硬的百姓,不想順從.人可以不順從,神不把人當作機器,不是把工作方法定好,指定人去做.神要我選擇順從神的道路,但給我們自由選擇順從或不順從.但若選擇不順從之後,神的刑罰必臨到,神是公義的神,萬不以有罪為無罪.我們可以不順從,但無法逃避不順從的結果.但神仍揀選以色列人,作神的百姓.他們雖背逆,離開神,神必降刑罰與他們.但神對屬祂的人的刑罰與對不屬祂的人的刑罰,往往不相同.
\newline
\newline
從舊約的以賽亞,耶利米,以西結等先知書均可看見,神刑罰過以色列人,也刑罰過埃及,刑罰過亞述人,亞蘭人.神刑罰外邦人,目的是消滅.但神刑罰百姓,目的是管教.管教的刑罰與毀滅的刑罰有極大的分別.
\newline
\newline
這章經文裡,看見神刑罰以色列人,是因為他們自取其咎,在生活上,是明知故犯.
\newline
\newline
有時你冷淡,不追求,生活不好,神似乎不刑罰,也許神的刑罰未到,但早晚必臨到.
\newline
\newline
有一次慕迪講到神的審判,他說:「不論何時犯罪,神必刑罰.」有一人前來對慕迪說:「傳講耶穌是騙人的,我從來沒有敬畏神,但我萬事順遂.」慕迪說:「神與你結賬之日,不一定是在你收成的日子.」不錯,神結賬的日子,也許不是在今天,明天,但一定臨到我們.現在要看看神如何將管教臨到以色列人.
\newline
\newline
(一)生活失了意義,失去真正的目標
\newline
\newline
這是多麼可憐的生活.外表生活好,生活舒適,但這是沒有意義的生活,沒有目標,是可憐的生活.
\newline
\newline
北平菜式有一款叫填鴨.把鴨餵飽,放在一小籠裡,不能自由行動.吃的雖然是好的食物,但單吃不運動,很快長得肥而大.但養鴨的人,目的只是為宰烹而已.可見沒有意義,沒有目標的生活,多麼可憐.
\newline
\newline
「因為他說,我要隨從所愛的,我的餅,水,羊毛,麻,油,酒,都是他們給的.」(5節)換句話說,他們所追求的,所愛的,是餅與水,因他們給我羊毛,麻,油和酒,所以我隨從他.這是為肉體而追求,從他們追求的方向,便可以看出他們真正的生活是甚麼.
\newline
\newline
要得肉體滿足,是以色列人追求的目標.雖然可能得著餅,但生活的目標卻失去了.
\newline
\newline
認識主耶穌的人,知道生活多麼有價值.但真正生活的目標是甚麼?我們今天生活追求的是甚麼呢?屬靈的生活是永遠的,今天的生活是暫時的.你一年中,有多少時候是為屬世的生活安排,又用多少時間去為屬靈的生活安排呢?
\newline
\newline
俗語說:「哀莫大於心死.」追求屬靈上進的心沒有了,只為餅與水,油與酒,羊毛與麻.追求更高更屬靈的心已不存在,只要得餅就覺得滿足,這是以色列人的生活情況,多麼可憐!
\newline
\newline
(二)生活充滿攔阻
\newline
\newline
「因此,我必用荊棘堵塞他的道,築牆擋住他,使他找不著路.」就是說,神使他的生活中,充滿攔阻,充滿困難,條條路均堵塞,沒有路可通到神那裡.無法可解決困難,自己對自己也完全灰心.
\newline
\newline
今天我們為何在生活上也有如此多的困難,可能是因罪的緣故,生活給困難阻擋著,無法衝過.因為將生活放在肉體的環境裡,看重肉體,注重物質,當然無法勝過困難.神使他們的物質,成為他們的審判,成為他們的灰心,失望.
\newline
\newline
(三)沒有工作的果效
\newline
\newline
「他必追隨所愛的,卻追不上.他必找尋他們,卻尋不見.」(7節)他們不是不追求,不尋找?但雖然努力追求,努力尋找,卻得不到果效.
\newline
\newline
我們單努力追求,努力工作,為主作見證.但作見證,不能發生果效.因為靈裡貧窮,冷淡.努力追求錯了方向.
\newline
\newline
要勝過罪惡,但不能勝過；努力作見證,但得不著果效；愈努力愈沒有果效.怎辦?終於灰心喪志.
\newline
\newline
所以一天到晚在忙,愈忙愈亂,在主前被記念的究竟有多少?內心是冰冷的,與主交通有攔阻,祈禱不通,屬靈生命沒有長進,沒有渴慕主的道的心.於是不敢讓人知們我們屬靈生活的光景,用作主工來掩飾與主交通有了攔阻.這是根本的錯誤,愈尋找,愈不能尋見,這是神第三個刑罰.
\newline
\newline
(四)失去分辨是非的能力
\newline
\newline
「他不知道是我給他五穀,新酒和油.又加增他們的金銀,他卻以此供奉巴力.」(8節)神賜給他們的,他們反用來供奉巴力,以為是巴力所賜.他們在靈裡失去了分辨是非的能力.明明是神所賜的,滿有價值的,卻以為是巴力所賜的,變為沒有價值的.明知放縱肉體是神所不喜悅的,但仍追求肉體的滿足,在屬靈裡失掉了分辨是非的能力.以賽亞先知說:「禍哉,那些稱惡為善,稱善為惡,以暗為光,以光為暗,以苦為甜,以甜為苦的人.」(賽五20)失去了屬靈辨別力,是多麼可憐的生活.
\newline
\newline
不久前,我在巴士上,看見一老人正在掏錢買票,那司機把車一開一停,以致這老人跌在車上.司機的服務精神已不好,售票員還在說:「不喜歡乘巴士,不如去乘的士啦!」他們明知是錯,還以為對.
\newline
\newline
有人得罪我,我便要報復,這因為辨別是非的心失去了.
\newline
\newline
教會工作走錯了方向,失了生活的目標,不是一時犯錯,是整個屬靈分辨力喪失了.
\newline
\newline
神不勉強人走祂的路,人沒有分辨屬靈的亮光,神不感動他,不提醒他.主耶穌說:「你裡頭的光若黑暗了,那黑暗何等大呢!」(太六23)小孩子不能分辨是非,不足為奇,若是二十歲,或是三十歲,還不會分辨是非,那是多可憐!前幾天我到新界去,看見有兩個基督徒在談話,正談論著下次馬票開彩的時間,談論得興高采烈.世俗的事,會分辨,屬靈的事卻沒有分辨的亮光.
\newline
\newline
(五)在靈裡有極大的貧窮
\newline
\newline
「因此,到了收割的日子,出酒的時候,我必將我的新酒五穀收回,也必將他應當遮體的羊毛和麻奪回來.如今我必在他所愛的眼前,顯露他的醜態,必無人救他脫離我手.」神要將他的空虛,貧窮,軟弱,可憐,顯露出來.缺乏得不到滿足,痛苦得不到安慰.他們雖追求五穀新酒,但肚腹卻得不到飽足,他們赤身露體,神要將他的罪顯露出來.離開主的痛苦,想是人人都有的經驗,內心想得到安慰,卻得不到安慰.
\newline
\newline
在舞場,電影院裡尋找刺激一下,或許忘記了一時的虛空,但離開舞場,電影院之後,心靈更覺不安,更是空虛.多工作,可能多得錢.但心中沒有平安,晚上不能入寐.千萬富翁自殺,就是因為金錢不能滿足他的心.有學問的人怎樣,前幾天有一位有學問的人奢談,世上沒有神,但他說:「人人以為我有很大的成就,但我的心靈是澈底的空虛.」你看多可憐!
\newline
\newline
神給我們選擇順服與不順服,但若選擇不順服,祂要管教.但若肯悔改,神的態度改變了.請聽:
\newline
\newline
「耶和華說,那日你必稱呼我伊施,不再稱呼我巴力.」(16節)
\newline
\newline
「我必聘你永遠歸我為妻,以仁義,公平,慈愛,憐憫,聘你歸我.」(19節)
\newline
\newline
「耶和華說,那日我必應允天,天必應允地；地必應允五穀,新酒和油,這些必應允耶斯列民.」(21,22節)
\newline
\newline
神正等待我們悔改歸回,順從祂.神願意賜福給他的百姓,多過百姓渴望神的賜福!「回頭吧,回頭吧!何必死亡呢?」
\newline
\newline


\newpage

\section{第40屆~港九培靈會~鮑會園~第四講}
\label{sec:1518}
\textbf{完全的愛}
\newline
\newline
連結: \href{https://www.hkbibleconference.org/session-message/view/1518}{\texttt{https://www.hkbibleconference.org/session-message/view/1518}}
\newline
\newline
\hyperref[sec:1517]{< < < PREV SERMON < < <}
~
\hyperref[sec:toc]{[返目錄]}
~
\hyperref[sec:1519]{> > > NEXT SERMON > > >}
\newline
\newline
何西阿書 3:1-3:5
\newline
\begin{longtable}{cl}
\hline
\hline
章節 & 經文 (和合本修訂版)\\
\hline
3:1 & \begin{tabularx}{0.7\textwidth}{X} 耶和華又對我說：「你去愛那情人所愛卻犯姦淫的婦人，正如耶和華愛那偏向別神、喜愛葡萄餅的以色列人。」 \end{tabularx} \\ \\ \relax
3:2 & \begin{tabularx}{0.7\textwidth}{X} 於是我用十五舍客勒銀子和一賀梅珥半大麥買她歸我。 \end{tabularx} \\ \\ \relax
3:3 & \begin{tabularx}{0.7\textwidth}{X} 我對她說：「你當多日與我同住，不可行淫，不可歸與別人，我對你也一樣。」 \end{tabularx} \\ \\ \relax
3:4 & \begin{tabularx}{0.7\textwidth}{X} 因為以色列人必多日過著無君王，無領袖，無祭祀，無柱像，無以弗得，無家中神像的生活。 \end{tabularx} \\ \\ \relax
3:5 & \begin{tabularx}{0.7\textwidth}{X} 後來以色列人必歸回，尋求耶和華---他們的神和他們的王大衛。在末後的日子，他們必敬畏耶和華，領受他的恩惠。 \end{tabularx} \\ \\
[1ex]
\hline
\hline
\end{longtable}
「耶和華對我說,你再去愛一個淫婦,就是他情人所愛的；好像以色列人,雖然偏向別神,喜愛葡萄餅,耶和華還是愛他們.」這節經文似乎有點難以解釋.但所指的淫婦,顯然是指第一章神要何西阿所娶的妻.因為這婦人離開了丈夫,不知丈夫的愛.本來在愛中與丈夫結合,但現在被別的東西吸引去了,離開了丈夫.至今仍未回頭,心中仍有所愛的,忘記了丈夫,所以神吩咐何西阿再娶這婦人為妻.
\newline
\newline
按人來看,這怎麼可能.如說把一些好處給她,幫助她也行,但再娶她為妻,實在有困難.但何西阿仍然去做,何西阿作這事,是表明神對百姓的愛.何西阿的生活,是神與百姓的預表.百姓心中麻木了,不明白神的愛.現在神藉何西阿把對百姓的愛顯露出來.
\newline
\newline
世人也不明白神的大愛,神也是要藉屬祂的人的生活,在世人面前把神的愛彰顯出來.信徒在世人面前是主耶穌的代表,你看,主將何等重大的責任放在我們身上!今天是一個悖逆的世代,信主的人在世人面前要將主的愛彰顯出來,你看,我們身份何等高貴.我們不要辜負主的使命,也不可辜負這高貴的身份.
\newline
\newline
現在看看何西阿怎樣將神的愛彰顯出來.
\newline
\newline
1. 神已經揀選百姓,拯救百姓,用全愛來愛百姓.但百姓仍然離開神,還說他們的好處,是外邦偶像給他們的.可見人不可愛,但神是愛,這也是人愛與神愛的分別.神愛以色列百姓,不是他們有何可愛.
\newline
\newline
有主的愛在我們當中,就可以在世人面前將主的愛彰顯出來.不但愛那些對我好的人,也愛那些對我不好的人；愛朋友,也愛仇敵.要像何西阿一樣,愛那不可愛的人.
\newline
\newline
神的愛是超乎自然的愛,不受別人愛的影響,是我裡面有神完全的愛.
\newline
\newline
基督徒的愛不是被動式的愛,不是反應式的愛.你愛我,我就愛你；你對我好,我當然對你好；你愈不愛我,我更愈不愛你.你可愛時我愛你,不可愛時便不愛你.基督徒有神的愛,才能彼此相愛.
\newline
\newline
基督徒生活方式,也不受外力控制.比如因同工說了我一句話,我心中便不快,這是被同工控制.我的心極其難過,因為某弟兄譭謗我,這樣,我的生活便為弟兄所控制.他說我好,我很感謝；說我不好,也很感謝,這就是不受外力控制.
\newline
\newline
如果有神超然的愛存在我裡面,我愛人不是用反應式去愛人；我的愛,勝過別人的譭謗,也勝過別人的輕視.
\newline
\newline
神愛以色列百姓,不是因為他們多好,多可愛.今天我們的愛,也是因為有這亮光.這是我們最高的標準.主耶穌說:「你們要完全,像你們的天父完全一樣.」「因為他叫日頭照好人,也照歹人；降雨給義人,也給不義的人.」何西阿愛那不可愛的妻子,我們也要用盡方法愛那些不愛我的人.
\newline
\newline
2. 是真正的愛.神用全心愛百姓,這是真正的愛.沒有人說:「不管她愛不愛我,只要她知道我愛她便算.」因為不是單方面將愛表明便算,因為這不是真愛.
\newline
\newline
真正的愛,二人同心,才得到滿足,愛達不到對象,多麼痛苦!母親愛兒子,兒子愛母親,但永不能見面,這是多麼痛苦.
\newline
\newline
「我對他說,你當多日為我獨居,不可行淫,不可歸別人為妻,我向你也必這樣.」(3節)神這樣全心愛百姓,不叫他們與祂連在一起,要獨居,顯明神的愛更偉大.
\newline
\newline
這是給百姓有機會反省,知道自己的錯誤,是為百姓的好處.但在神方面,愛是得不到滿足的,是極難忍受的.按神的愛要與百姓結合,但為百姓的好處,寧願為他們獨居,神的愛能勝過一切.神對百姓的愛沒有改變,但要他們獨居,反省自己,巴不得此時他心中有愛神的感覺,在愛上如何跌倒,應回過頭來,便如何起來.
\newline
\newline
3. 百姓對神的愛,有甚麼反應?這章聖經給我們看見,我們與神恢復愛,有兩個很顯明的步驟.請你查考一下,究竟你在那一個步驟上:
\newline
\newline
i. 第一個步驟是:「以色列人也必多日獨居,無君王,無首領,無祭祀,無柱像,無以弗得,無家中的神像.」(4節)
\newline
\newline
他們是屬神的兒女,是神所愛的人,與神有法律上的關係.但他們獨居,無祭司,意即沒有與神交通.有法律上的關係,但無享受與神交通的關係.無君王,意即神在他們身上無權柄,這是多可憐的生活.如果這是基督徒的生活,他是一個多可憐的基督徒呀!
\newline
\newline
基督徒想進電影院,進舞場,但想想若給人看見,多難為情；但與神的交通也沒有.屬世與屬靈的享受均失掉了.罪中享樂失掉了,屬靈滿足也沒有了,這樣的基督徒,生活多可憐!
\newline
\newline
「你們是世上的鹽,鹽若失了味,怎能叫他再鹹呢?以後無用,不過丟在外面被人踐踏.」因為鹽是鹽,不是糖呀!
\newline
\newline
「你們是世上的光,城造在山上,是不能隱藏的.」若不能發光,是多麼可憐呀!
\newline
\newline
與不信的人在一起,不能作基督的見證；與屬靈的人在一起,又不敢貪戀罪中的享樂,這是多麼可憐!
\newline
\newline
「我使用銀子十五舍客勒,大麥一賀梅耳半,買他歸我.」(2節)主在十字架流血,將我們從罪中買回來,這雖然在法律上有了關係,但失了屬靈的快樂,沒有享受屬靈的交通,若能獨居省察,悔改認罪,便可達第二個階段.
\newline
\newline
ii. 「後來以色列人必歸回,尋求他們的神耶和華和他們的王大衛.在末後的日子,必以敬畏的心歸向耶和華,領受他的恩惠.」(5節)以色列百姓因明白了自己所需要的,也明白了自己的軟弱就歸回祂.省察之後,他們知道無享受,無喜樂,無滿足,便思想要歸回.浪子耗盡了所有,生活窘迫,便想起家父的糧食的豐富.於是起來,離開遠方的可憐地位,回到家裡,重新享受父愛.
\newline
\newline
但願我們省察之後,便起來,回到天父溫暖的家,享受天父為我們預備的慈愛,真愛和大愛.
\newline
\newline


\newpage

\section{第40屆~港九培靈會~鮑會園~第五講}
\label{sec:1519}
\textbf{生命的知識}
\newline
\newline
連結: \href{https://www.hkbibleconference.org/session-message/view/1519}{\texttt{https://www.hkbibleconference.org/session-message/view/1519}}
\newline
\newline
\hyperref[sec:1518]{< < < PREV SERMON < < <}
~
\hyperref[sec:toc]{[返目錄]}
~
\hyperref[sec:1520]{> > > NEXT SERMON > > >}
\newline
\newline
何西阿書 4:1-4:6
\newline
\begin{longtable}{cl}
\hline
\hline
章節 & 經文 (和合本修訂版)\\
\hline
4:1 & \begin{tabularx}{0.7\textwidth}{X} 以色列人哪，當聽耶和華的話。耶和華指控這地的居民，因為在這地上無誠信，無慈愛，無人認識神； \end{tabularx} \\ \\ \relax
4:2 & \begin{tabularx}{0.7\textwidth}{X} 惟起誓、欺騙、殺害、偷盜、姦淫、殘暴、流血又流血。 \end{tabularx} \\ \\ \relax
4:3 & \begin{tabularx}{0.7\textwidth}{X} 因此，這地悲哀，其上的居民、野地的走獸、天空的飛鳥都日趨衰微，海中的魚也必消滅。 \end{tabularx} \\ \\ \relax
4:4 & \begin{tabularx}{0.7\textwidth}{X} 然而，人都不必爭辯，也不必指責。你的百姓與抗拒祭司的人一樣。 \end{tabularx} \\ \\ \relax
4:5 & \begin{tabularx}{0.7\textwidth}{X} 日間你必跌倒，夜間先知也要與你一同跌倒；我要滅絕你的母親。 \end{tabularx} \\ \\ \relax
4:6 & \begin{tabularx}{0.7\textwidth}{X} 我的百姓因無知識而滅亡。你拋棄知識，我也必拋棄你，使你不再作我的祭司。你既忘了你神的律法，我也必忘記你的兒女。 \end{tabularx} \\ \\
[1ex]
\hline
\hline
\end{longtable}
「以色列哪,你們當聽耶和華的話.耶和華與這地的居民爭辯,因這地上無誠實,無良善,無人認識神.」百姓當中充滿罪,若說外邦人不認識神,這也難怪,但連以色列人不認識神,那就奇怪了.
\newline
\newline
何西阿時,以色列人在各處仍有祭壇,熱心敬拜,但神說:「我的民因無知識而滅亡.」(6節)為甚麼說百姓沒有知識.這「知識」的意義是甚麼?聖經裡所說的知識,有各種不同的意義,每字代表一種知道的方法,用這種方法才可得到一種知識.
\newline
\newline
在這裡何西阿所用的「知識」,是一個特別的字.先談談有關知識的問題.
\newline
\newline
有些知識是從別人領受,比方我們未見過秦始皇,但認識秦始皇的知識,是別人間接傳給我們的.
\newline
\newline
前幾天伊朗有地震,但香港沒有人親眼看見,因為這是無線電傳給我們知道的.這間接的知道,有些是歷史的事實,有些是無線電傳真.又如聖經裡也有許多歷史的事實,亞伯拉罕蒙神揀選；以色列人出埃及,主降生在伯利恆城裡,被釘在各各他山的十字架上,這些知識,我們都從間接領受而來的.
\newline
\newline
第二種知識,如數學方程式,不能單憑別人傳授,還需自己運用頭腦才明白.
\newline
\newline
哲學原理也一樣,只背誦一些書籍是不夠的,還需自己去體驗.這是代表思想方式.
\newline
\newline
第三種是用這種知識的方法,去得到別種知識.我到泰國,得到很好的招待.吃一種叫榴連的生果,初時覺得味道很不好,但勉強吃一些,漸漸不覺得難吃.
\newline
\newline
後來有人問我榴連是怎樣味道的,我雖然可以解釋說,又甜又軟,像軟糖.但你不明白我所說的,便再問,但我無法再作介紹,只有請你嘗嘗,你才明白這味道兒.這種知識是無法傳授的,用盡方法也不行,只有親身經歷才知道.
\newline
\newline
可以用別人傳達這知識給你,記下,思想,就算想通,但未親身經歷,仍無法得這知識,這就是第六節所說的「知識」.以色列人,不認識神,不明白神,可能聽過神的話,思想過,且有高深的教育,但未親身經歷耶和華.
\newline
\newline
正如先知以賽亞說:「牛認識主人,驢認識主人的槽,以色列卻不認識我.」(賽一3)
\newline
\newline
有人可能知道真理,知道關於主耶穌的事,或教會的信條,但未親身經歷主耶穌.主耶穌在我裡面,不是活的救主,不能改變我的生命.雖然明白,但未得生命的造就.
\newline
\newline
有人會講道,但屬靈生命枯乾,明白真理,但不明白主耶穌.
\newline
\newline
今天我們到底認識主的知識如何?生命有沒有得著造就.
\newline
\newline
有時明白神的真理太多,但真理不能感動他們,這比不信的人更可憐.知道真理,但真理卻成為定我們的罪.
\newline
\newline
為甚麼以色列百姓缺乏這知識,主要有兩個原因.
\newline
\newline
1. 缺乏對神有真正的認識,只將重點放在表面上.如只注重獻祭,或獻祭的地方.在敬拜上,節期上,他們的宗教生活只是一些儀式.
\newline
\newline
今天我們來赴培靈會是為了甚麼?我聽過有人說,八天我每堂都到齊,沒有一堂缺席過.他就是人來,我也來,且來得更殷勤.
\newline
\newline
有人讀聖經是為安慰一下良心,麻木一下生活不好的良心.單有外表儀式,用儀式代替與神的交通,所以「我的民因無知識而滅亡.」
\newline
\newline
2. 缺乏對神有真正的認識,是因他們拒絕屬靈的知識.愈聽神的話,應愈接受.愈聽主的聲音,應愈清楚,愈順服.但今天在猶疑下,明天更難明白,後天便會成為重擔.
\newline
\newline
「你棄掉知識,我也棄掉你.」(6節中)因百姓不認識神,他們棄掉知識.一次棄掉,再次棄掉,沒有屬靈的感覺.赤足走路不慣,遇有沙石,痛苦不堪.但有人隨意走在沙石上,不覺痛苦,因他習慣了,硬皮生厚了.這硬皮,是死皮,沒有知覺的.屬靈的事也如此,因一次,再次拒絕,屬靈的感覺便麻木了,不再領略神的知識,這是不順從最可怕的結果.
\newline
\newline
久犯罪,便不覺得罪是罪；多次不順從,便不覺得不順從的難過,因為與神的關係麻木了.
\newline
\newline
有一對夫婦,感情十分融洽,得罪對方,即道歉認錯.但不久,情形不同了,感情漸破裂,甚至公開對罵,故意說些使對方難聽的話.其實這些話,對人對己都沒有好處.問她為何講這些話,她說目的叫他難過,她說目的叫他刺心,因為他們的心已麻木了.
\newline
\newline
基督徒有時也是這樣,知道這些事是主不喜悅的,但做了也不覺得難過.
\newline
\newline
拒絕知識另一個結果是:「起假誓,不踐前言,殺害,偷盜,姦淫,行強暴,殺人流血,接連不斷.」(2節)屬靈生命,永遠不能靜止的,要向前走.讀經,祈禱,敬拜仍不能叫屬靈生命進展.不錯,若沒有讀經,祈禱和敬拜,生命更敗壞,但不是外表做做便算.起假誓,不進步,結果走上更悲慘的道路.
\newline
\newline
怎樣才得到認識神的知識?
\newline
\newline
「人若立志遵著他的旨意行,就必曉得這教訓或是出於神,或是憑著自己說的.」(約七17)順服神,必定有一個遵從主命令的心.如果在屬靈上有真知識,必要遵行主的命令.屬世的事,先知道才遵守,先明白了才遵守.但屬靈的事,先有了一個遵守主命的心,才能與主有更深切的交通.
\newline
\newline
怎樣才算真明白主的旨意?
\newline
\newline
1. 在約翰福音十四章裡:「有了我的命令又遵守的,這人就是愛我的；愛我的必蒙我父愛他,我也要愛他,並且要向他顯現.」(21節)「人若愛我,就必遵守我的道；我父也必愛他,並且我們要到他那裡去,與他同住.」(23節)以上的「顯現」與「同住」；「命令」與「道」有甚麼分別?
\newline
\newline
比方,有一天,有位媽媽要外出,吩咐她的小孩不要到外邊去同那些野孩子玩耍.她的孩子果然為了答應過媽媽,所以沒有到門外去玩耍,但他想出一個很好的方法,乃是招請那些孩子到家裡玩耍.他說媽媽只說不要出去,但沒有說不可以讓他們進來.各位,這小孩果真遵守了媽媽的命令了嗎?他遵命了,但不照媽媽的旨意行事.
\newline
\newline
有人可以將聖經列成條文,一點不違背,但可能不按天父的心去行,可能遵守命令,但沒有體貼父的旨意.
\newline
\newline
比如說,聖經那處說不可賭博?不可看電影?主雖然未說過,但我們都知道,若是做了,便是違背主的命令.因為這與聖潔的生活不符,會攔阻屬靈生命的長進.
\newline
\newline
2. 「清心的人有福了,因為他們必得見神.」(太五8)這「見」字也是一個特別的字,不單肉眼見,乃從靈裡見而接受.認識神,要有清潔的心,和聖潔的生活.
\newline
\newline
聖潔包含兩方面的意義:
\newline
\newline
不吸煙,不飲酒,不賭博,不看電影,這是是消極的聖潔.
\newline
\newline
清心的「清」字是積極的意思.要認識神,要沒有虛假,沒有別的目的.
\newline
\newline
3. 安靜等待神.「你們要休息,要知道我是神.」(詩四十10)「休息」可譯作「安靜」.世上不論甚麼東西,對找都失去吸引力,便能安靜在神面前.
\newline
\newline
每天用半小時把別的事放下,而安靜在神面前多麼難.不但要「安靜」,還要「敞開」我們的心,才能真正認識主.
\newline
\newline
4. 要有一個真正相信主的心.相信主耶穌,藉著主耶穌,才能認識神.「因為在天下人間,沒有賜下別的名.」「因為那已經立好的根基,就是耶穌基督,此外沒有人能立別的根基.」(林前三11)不藉著主,不能親近神.
\newline
\newline
在追求認識神的道,也許有困難,因有肉體的吸引,但主耶穌基督能在一切環境上幫助我們,得勝有餘.
\newline
\newline


\newpage

\section{第40屆~港九培靈會~鮑會園~第六講}
\label{sec:1520}
\textbf{神的爭辯}
\newline
\newline
連結: \href{https://www.hkbibleconference.org/session-message/view/1520}{\texttt{https://www.hkbibleconference.org/session-message/view/1520}}
\newline
\newline
\hyperref[sec:1519]{< < < PREV SERMON < < <}
~
\hyperref[sec:toc]{[返目錄]}
~
\hyperref[sec:1521]{> > > NEXT SERMON > > >}
\newline
\newline
何西阿書 4:1-4:19
\newline
\begin{longtable}{cl}
\hline
\hline
章節 & 經文 (和合本修訂版)\\
\hline
4:1 & \begin{tabularx}{0.7\textwidth}{X} 以色列人哪，當聽耶和華的話。耶和華指控這地的居民，因為在這地上無誠信，無慈愛，無人認識神； \end{tabularx} \\ \\ \relax
4:2 & \begin{tabularx}{0.7\textwidth}{X} 惟起誓、欺騙、殺害、偷盜、姦淫、殘暴、流血又流血。 \end{tabularx} \\ \\ \relax
4:3 & \begin{tabularx}{0.7\textwidth}{X} 因此，這地悲哀，其上的居民、野地的走獸、天空的飛鳥都日趨衰微，海中的魚也必消滅。 \end{tabularx} \\ \\ \relax
4:4 & \begin{tabularx}{0.7\textwidth}{X} 然而，人都不必爭辯，也不必指責。你的百姓與抗拒祭司的人一樣。 \end{tabularx} \\ \\ \relax
4:5 & \begin{tabularx}{0.7\textwidth}{X} 日間你必跌倒，夜間先知也要與你一同跌倒；我要滅絕你的母親。 \end{tabularx} \\ \\ \relax
4:6 & \begin{tabularx}{0.7\textwidth}{X} 我的百姓因無知識而滅亡。你拋棄知識，我也必拋棄你，使你不再作我的祭司。你既忘了你神的律法，我也必忘記你的兒女。 \end{tabularx} \\ \\ \relax
4:7 & \begin{tabularx}{0.7\textwidth}{X} 祭司越發增多，就越發得罪我；我必使他們的榮耀變為羞辱。 \end{tabularx} \\ \\ \relax
4:8 & \begin{tabularx}{0.7\textwidth}{X} 他們吞吃我百姓的贖罪祭，滿心願意我的子民犯罪。 \end{tabularx} \\ \\ \relax
4:9 & \begin{tabularx}{0.7\textwidth}{X} 將來百姓所受的，祭司也必承受；我必因他們所行的懲罰他們，照他們所做的報應他們。 \end{tabularx} \\ \\ \relax
4:10 & \begin{tabularx}{0.7\textwidth}{X} 他們吃，卻不得飽足；行淫，卻不繁衍；因為他們離棄耶和華，常行 \end{tabularx} \\ \\ \relax
4:11 & \begin{tabularx}{0.7\textwidth}{X} 淫亂。酒和新酒奪去人的心。 \end{tabularx} \\ \\ \relax
4:12 & \begin{tabularx}{0.7\textwidth}{X} 我的百姓求問木偶，以為木杖能指示他們；淫亂的心使他們失迷，以致行淫離棄他們的神， \end{tabularx} \\ \\ \relax
4:13 & \begin{tabularx}{0.7\textwidth}{X} 在各山頂獻祭，在各高岡上燒香，在橡樹、楊樹、大樹之下，因為那裡樹影美好。所以，你們的女兒行淫，你們的媳婦犯姦淫。 \end{tabularx} \\ \\ \relax
4:14 & \begin{tabularx}{0.7\textwidth}{X} 我不因你們的女兒行淫或你們的媳婦犯姦淫懲罰她們；因為人自己轉去與娼妓同居，與神廟娼妓一同獻祭。這無知的百姓必致傾倒。 \end{tabularx} \\ \\ \relax
4:15 & \begin{tabularx}{0.7\textwidth}{X} 以色列啊，你雖然行淫，猶大卻不可犯罪；不要往吉甲去，不要上到伯‧亞文，也不要指著永生的耶和華起誓。 \end{tabularx} \\ \\ \relax
4:16 & \begin{tabularx}{0.7\textwidth}{X} 以色列倔強，猶如倔強的母牛；現在耶和華能牧放他們，如在寬闊之地牧放羔羊嗎？ \end{tabularx} \\ \\ \relax
4:17 & \begin{tabularx}{0.7\textwidth}{X} 以法蓮親近偶像，任憑他吧！ \end{tabularx} \\ \\ \relax
4:18 & \begin{tabularx}{0.7\textwidth}{X} 他們喝完了酒，荒淫無度，他們的官長甚愛羞恥的事。 \end{tabularx} \\ \\ \relax
4:19 & \begin{tabularx}{0.7\textwidth}{X} 風把他們捲在翅膀裡，他們必因所獻的祭蒙羞。 \end{tabularx} \\ \\
[1ex]
\hline
\hline
\end{longtable}
「耶和華與這地的居民爭辯,因這地上無誠實,無良善,無人認識神.但起假誓,不踐前言.」我們分析全本何西阿書,可分為兩段,頭一段是一至三章,四至十四章是另一段.第一段神責備祂的百姓,說及他們消極的罪,應行不去行,應良善沒有良善.
\newline
\newline
第四章才進一步責備百姓積極的罪.不但無良善,無誠實,而且殺害,偷盜,姦淫,行強暴,殺人流血,接連不斷.這也是今天教會的景況.傳信息不易,因為要責備罪,我是香港教會一份子,教會有罪,每個信徒都有責任.所以我預備這信息的時候,心中很是痛苦,因為神指責百姓的罪,也是指責今天香港教會的罪.
\newline
\newline
但感謝神,神不但責備罪,他為我們開了一條出路.昨天我說過,屬靈的生命不能靜止,不能在真空裡生活.說因沒有時間,就不追求.不對付罪,就是為罪留下機會.開始時,不過是無誠實,無良善,結果是殺害,偷盜,姦淫,行強暴.開始時不過是不熱心,結果,行為完全與主相違背.所以屬靈的事不能靜止,如逆水行舟,不進則退.
\newline
\newline
為此,神要與百姓爭辯.「爭辯」可譯作「控訴」.過去,神曾與百姓立約,神與百姓均有遵守的義務.
\newline
\newline
以色列百姓若遵守與神所立的約,就有福.若不遵守,神便要控訴百姓,要在律法上爭辯.
\newline
\newline
以色列百姓是神所揀選的,是神愛的對象.本來與神合而為一,但今時卻站在敵對的地位上.今天我們的生活是否也是與神站在敵對地位,抑或站在神的一邊?
\newline
\newline
神提出要控訴他們.因為:
\newline
\newline
(一)藉著別人的罪,從中取利
\newline
\newline
「他們吃我民的贖罪祭,滿心願意我民犯罪.」(8節)「贖」,「祭」字在原文是沒有的.所以可譯為「他們吃我民的罪.」這是當時祭司們生活的情況,靠百姓犯罪,他們從中得利.他們使百姓犯罪,但又責備他們,抓住他們的罪,去從中得利,有甚麼事比這事更慘酷,更殘忍!
\newline
\newline
教會歷史上也有,犯了罪的人可買贖罪票,教會便可從中得利,以致發達.今天教會也有一些這樣的事,有信徒犯了罪,因送禮給我,我怎能責備他的罪.
\newline
\newline
有人出賣學位,使人犯罪,助人掩罪,便可從中得利.
\newline
\newline
(二)在罪中放縱自己
\newline
\newline
「他們吃,卻不得飽.」(10節)「吃」字在文法上是一個特別的字.應譯為「他們繼續吃,但吃不飽.」愈吃愈不飽,下文說他們種的是東風,吃的也是東風.吃東西不能飽足,便想多吃一點,放縱自己的生活.這是描寫人在罪中的情況,得不到滿足,仍然追求；不飽足,繼續犯罪.但第一次所犯的罪,易對付,第二次所犯的罪更大.他們這樣放縱自己,卻得不到飽足.
\newline
\newline
(三)過著雙重的生活
\newline
\newline
「你們的兒女淫亂,你們的新婦行淫,我卻不懲罰他們；因為你們自己離群與娼妓同居,與妓女一同獻祭.」(10節)他們犯罪也和你一樣.男人犯罪不算罪,女人犯罪才算罪,這是不行的.你責備別人的罪,自己卻犯這罪.
\newline
\newline
有人對待同樣的罪,如果是朋友所犯便不出聲,是仇敵所犯便責備.容忍自己的罪,卻去責備別人的罪.生活上對自己有一個標準,對別人又另一個標準.在人前有一個標準,在神前又有另一個標準.
\newline
\newline
教會事工受攔阻,也是因有雙重的生活標準.今天教會對待傳道人的看法,一個讀完了高中,又讀四年神學的傳道人,所得的薪金不及一個幼稚園教師.還說,傳道人應當吃苦的.但事奉神的人應當吃苦,教會的弟兄姊妹,為甚麼不吃苦呢?
\newline
\newline
不久前有一立教會領袖說,以傳道人同樣薪金去請一位司機,也請不到.可見不是青年人不肯奉獻,乃是一個小學教員的薪金兩倍於一個傳道人,於是青年人對奉獻道路便裹足不前.請問弟兄姊妹過著舒適的生活,要傳道人過痛苦的生活,這是合理不合理.這是教會的雙重生活的標準.
\newline
\newline
(四)倔強不順從
\newline
\newline
「以色列倔強,猶如倔強的母牛.」這倔強,不是性情的倔強,乃如一隻母牛,你將軛放在他們身上,他們要掙脫,不肯向著神所指示去行,不願受限制.
\newline
\newline
神說:「現在要放他們,如同放羊羔在寬闊之地.」他們要掙脫神的軛,神就將他們放在寬闊處,即曠野之地,那兒無保護,無安穩.神藉先知所傳的信息,像畫下界線,你要越出範圍,便是無保障.神也給我們每人一條界線,我們若安心在神範圍下生活,那是福氣,安穩無憂.
\newline
\newline
我不會唱歌,若為此埋怨神,聽見歌聲,便心中難過.這是越出神的範圍,自尋煩惱.有人看見別人有成就,便說神不公平,這也是越出神的界線.
\newline
\newline
神將傳道拯救靈魂的責任給我,但若覺得不滿足,結果連傳道責任也不盡責.因為如果越過神的範圍,我們便落在痛苦中,毫無保障.
\newline
\newline
(五)拜偶像的生活
\newline
\newline
「我的民求問木偶,以為木杖能指示他們；因為他們的淫心,使他們失迷,他們就行淫離棄神,不守約束,在各山頂,各高崗的橡樹,楊樹,栗樹之下,獻祭燒香.」(12,13節)當日百姓有神的祭壇,也守耶和華的節期,表面按摩西的律法去行.但當時迦南地還有許多偶像,百姓的心,又不願離開這些偶像,表面敬拜耶和華,但心卻追隨偶像.
\newline
\newline
按真正的敬拜,不是一兩小時的事,而是內心所事奉的神,這才是真正屬靈的生活.
\newline
\newline
事奉神,又事奉瑪門；事奉學問,也事奉政治勢力,必一事無成.於是神給百姓兩個選擇.
\newline
\newline
1. 以色列百姓家中沒有偶像,但仍是一個拜偶像的人,他們不聽神的勸導,所以「以法蓮親近偶像,任憑他吧!」(17節)神不勉強百姓,他不順從,神就任憑他.
\newline
\newline
2. 神要百姓「回轉吧!回轉吧!何必死亡呢!」神要他們思想,再回轉過來歸向神.
\newline
\newline
這兩條路也擺在我們面前,應選擇那條呢?
\newline
\newline


\newpage

\section{第40屆~港九培靈會~鮑會園~第七講}
\label{sec:1521}
\textbf{悔改的呼召}
\newline
\newline
連結: \href{https://www.hkbibleconference.org/session-message/view/1521}{\texttt{https://www.hkbibleconference.org/session-message/view/1521}}
\newline
\newline
\hyperref[sec:1520]{< < < PREV SERMON < < <}
~
\hyperref[sec:toc]{[返目錄]}
~
\hyperref[sec:1526]{> > > NEXT SERMON > > >}
\newline
\newline
何西阿書 5:1-6:11
\newline
\begin{longtable}{cl}
\hline
\hline
章節 & 經文 (和合本修訂版)\\
\hline
5:1 & \begin{tabularx}{0.7\textwidth}{X} 眾祭司啊，要聽這話！以色列家啊，要留心聽！王室啊，要側耳而聽！審判將臨到你們，因你們在米斯巴如羅網，在他泊山如張開的網。 \end{tabularx} \\ \\ \relax
5:2 & \begin{tabularx}{0.7\textwidth}{X} 這些悖逆的人大行殺戮，我要斥責他們眾人。 \end{tabularx} \\ \\ \relax
5:3 & \begin{tabularx}{0.7\textwidth}{X} 至於我，我認識以法蓮，以色列不能向我隱藏。以法蓮哪，現在你竟然行淫，以色列竟然被污辱。 \end{tabularx} \\ \\ \relax
5:4 & \begin{tabularx}{0.7\textwidth}{X} 他們所做的使他們不能歸向神，因有淫亂的心在他們裡面；他們不認識耶和華。 \end{tabularx} \\ \\ \relax
5:5 & \begin{tabularx}{0.7\textwidth}{X} 以色列的驕傲使自己臉面無光；以色列和以法蓮必因自己的罪孽跌倒，猶大也必與他們一同跌倒。 \end{tabularx} \\ \\ \relax
5:6 & \begin{tabularx}{0.7\textwidth}{X} 他們牽著牛羊去尋求耶和華，卻尋不著；因他已轉去離開他們。 \end{tabularx} \\ \\ \relax
5:7 & \begin{tabularx}{0.7\textwidth}{X} 他們不忠於耶和華，生了私生子。現在新月必吞滅他們和他們的地業。 \end{tabularx} \\ \\ \relax
5:8 & \begin{tabularx}{0.7\textwidth}{X} 你們當在基比亞吹角，在拉瑪吹號，在伯‧亞文發出警報；便雅憫哪，留意你的背後！ \end{tabularx} \\ \\ \relax
5:9 & \begin{tabularx}{0.7\textwidth}{X} 到了懲罰的日子，以法蓮必變為廢墟；我在以色列眾支派中，已指示將來必成的事。 \end{tabularx} \\ \\ \relax
5:10 & \begin{tabularx}{0.7\textwidth}{X} 猶大的領袖如同挪移地界的人，我必把我的憤怒如水傾倒在他們身上。 \end{tabularx} \\ \\ \relax
5:11 & \begin{tabularx}{0.7\textwidth}{X} 以法蓮因喜愛遵從荒謬的命令，就受欺壓，在審判中被壓碎。 \end{tabularx} \\ \\ \relax
5:12 & \begin{tabularx}{0.7\textwidth}{X} 我對以法蓮竟如蛀蟲，向猶大家竟如朽爛。 \end{tabularx} \\ \\ \relax
5:13 & \begin{tabularx}{0.7\textwidth}{X} 以法蓮見自己有病，猶大見自己有傷，以法蓮就前往亞述，差遣人去見大王；他卻不能醫治你們，不能治好你們的傷。 \end{tabularx} \\ \\ \relax
5:14 & \begin{tabularx}{0.7\textwidth}{X} 我必向以法蓮如獅子，向猶大家如少壯獅子。我要撕裂，並且離去，我必奪去，無人搭救。 \end{tabularx} \\ \\ \relax
5:15 & \begin{tabularx}{0.7\textwidth}{X} 我要去，我要回到原處，等他們自覺有罪，尋求我的面；急難時他們必切切尋求我。 \end{tabularx} \\ \\ \relax
6:1 & \begin{tabularx}{0.7\textwidth}{X} 來，我們歸向耶和華吧！他撕裂我們，也必醫治；打傷我們，也必包紮。 \end{tabularx} \\ \\ \relax
6:2 & \begin{tabularx}{0.7\textwidth}{X} 過兩天他必使我們甦醒，第三天他必使我們興起，我們就在他面前得以存活。 \end{tabularx} \\ \\ \relax
6:3 & \begin{tabularx}{0.7\textwidth}{X} 我們要認識，要追求認識耶和華。他如黎明必然出現，他必臨到我們像甘霖，像滋潤土地的春雨。 \end{tabularx} \\ \\ \relax
6:4 & \begin{tabularx}{0.7\textwidth}{X} 以法蓮哪，我可以向你怎樣行呢？猶大啊，我可以向你怎樣做呢？因為你們的慈愛如同早晨的雲霧，又如速散的露水。 \end{tabularx} \\ \\ \relax
6:5 & \begin{tabularx}{0.7\textwidth}{X} 因此，我藉先知砍伐他們，以我口中的話殺戮他們；對你的審判如光發出。 \end{tabularx} \\ \\ \relax
6:6 & \begin{tabularx}{0.7\textwidth}{X} 我喜愛慈愛，不喜愛祭物；喜愛人認識神，勝於燔祭。 \end{tabularx} \\ \\ \relax
6:7 & \begin{tabularx}{0.7\textwidth}{X} 他們卻如亞當背約，在那裡向我行詭詐。 \end{tabularx} \\ \\ \relax
6:8 & \begin{tabularx}{0.7\textwidth}{X} 基列是作惡之人的城，被血沾染。 \end{tabularx} \\ \\ \relax
6:9 & \begin{tabularx}{0.7\textwidth}{X} 成群的祭司如強盜埋伏等候，在示劍的路上殺戮，行了邪惡。 \end{tabularx} \\ \\ \relax
6:10 & \begin{tabularx}{0.7\textwidth}{X} 在以色列家我看見可憎的事，在以法蓮那裡有淫行，以色列被污辱了。 \end{tabularx} \\ \\ \relax
6:11 & \begin{tabularx}{0.7\textwidth}{X} 猶大啊，我使被擄之民歸回的時候，必有為你所預備的豐收。 \end{tabularx} \\ \\
[1ex]
\hline
\hline
\end{longtable}
本書悔改的話語,一而再的向百姓呼召.「來吧!我們歸向耶和華,他撕裂我們,也必醫治；他打傷我們,也必纏裹.過兩天,他必使我們甦醒,第三天,他必使我們興起,我們就在他面前得以存活.」(六1-2)照此段經文,百姓似乎有悔改的心,但我們仍未能照字面解釋.因「主說:以法蓮哪,我可向你怎樣行呢?猶大啊,我可向你怎樣呢?因為你們的良善,如同早晨的雲霧,又如速散的甘露.因此,我藉先知砍伐他們,以我口中的話殺戮他們；我施行的審判如光發出.」(六4-5)可見百姓的悔改是表面的,是虛假的,所以神要更厲害的責備他.
\newline
\newline
看看他們虛假的悔改
\newline
\newline
(Ａ)藉自己的功勞代替內心的悔改
\newline
\newline
神已宣佈他們的罪,但他們想在神面前行賄賂.神啊!我給你一些好處,不要審判我吧!「以色列的驕傲,當面見證自己.故此,以色列和以法蓮必因自己的罪孽跌倒；猶大也必與他們一同跌倒.他們必牽著牛羊去尋求耶和華,卻尋不見；他已經轉去離開他們.」(五5-6)神既宣佈百姓的審判,他們應當悔改,但他們說,將牛羊獻給神,但內心違背神,獻祭行動與內心不同,所以神不接納.
\newline
\newline
我若生活上有罪孽,可以向神說,我多些奉獻,你可不責備我嗎?這是不可以的,一切奉獻,不能代替心中真正的悔改.百姓以為奉獻牛羊就行,但牛羊不能帶他們到神面前.有悔改的行動,但無悔改的心,神是不悅納的.
\newline
\newline
(Ｂ)遮蓋自己的困難
\newline
\newline
「以法蓮見自己有病,猶大見自己有傷,他們就打發人往亞述去見耶雷布王,他卻不能醫治你們,不能治好你們的傷.我必向以法蓮如獅子,向猶大家如少壯獅子.我必撕裂而去,我要奪去,無人搭救.」(五13-14)百姓生活有危險,國家有難,不悔改所行的,卻去求助他國.他們不查出危難的原因是甚麼.
\newline
\newline
教會工作沒有見證,信徒生命沒有長進,不能結出屬靈的果子,雖有嚴密的組織,也不能解決問題.
\newline
\newline
心靈寂寞,與主的交通中斷.想得一時的滿足,追求自己喜悅的事.滿足肉體的享樂.
\newline
\newline
有一位弟兄,因有困難,以致影響青年團的氣氛,其實一個困難解決,又有新的困難.有人以為純正信仰教會未得到應得的果效.是因世界的聲勢,異端,新神學派的聲勢太大,其實真正的困難,不是這些原因,是因為神的靈已不在我們中間,失去屬靈的作戰能力.
\newline
\newline
(Ｃ)覺得自己的罪大太難對付
\newline
\newline
不是沒有悔改的心,也不是不恨惡罪惡.只是他們的悔改是暫時的,出現如同早晨的雲霧.巴勒斯坦氣候乾燥,白日海風吹來,潮濕成一層雲霧,看起來很大,但太陽一出,便不見了.他們不能保持悔改的心,永遠不能得到神的悅納.
\newline
\newline
(Ｄ)用信心做口號,來遮掩缺乏的信心
\newline
\newline
「過兩天他必使我們甦醒.第三天他必使我們興起.」他們說,我有軟弱,神不祝福.但神是不輕易發怒的,過兩天神會使我們甦醒,用不著擔心.神的教會永不失敗,到了時候,神必賜福,他們用神的恩慈,信實,權柄來作藉口.他們又說,我知有罪,但我已是得救的人,聖靈重生了我,我一生在神的家,還怕甚麼.他們表面似乎很有信心,其實是用來遮掩剛硬的心,藉神的恩慈來放縱自己.
\newline
\newline
真正的悔改
\newline
\newline
真正的悔改,首先對罪有所認識.甚麼是罪?兇殺,偷盜,貪心等是罪,當然是,但不是罪的基本意義.罪不是外表的,乃在內心與神的關係,不可把罪看得太輕,太淺,或用行動來遮掩罪.
\newline
\newline
比如丈夫得罪妻子,拿一塊糖給妻子,請她不發脾氣.恐怕因此使妻子更大發脾氣,因為這樣不能解決困難,而且加深了問題.如要解決困難,捨此並無別法,就是丈夫要向妻子認錯.同樣的我們犯罪得罪神,只有向神認罪,才能得到饒恕.
\newline
\newline
要知道罪是攔阻我們與神交通,要認識罪的可怕,就要認識神是多偉大,多聖潔,這樣便可認識在神面前,罪是多可怕,多可惡.
\newline
\newline
以賽亞看見神的偉大與聖潔,便說:「禍哉,我滅亡了.」
\newline
\newline
認識罪的性質,就是神多偉大,罪是多可怕,為罪要有痛悔的心.
\newline
\newline
遊子在外面想起母親不喜歡罪惡,便不敢犯罪.我們若知道神的心難過,就能恨惡罪惡.
\newline
\newline
認罪與道歉,亦有不同.比如我踏了你的腳,我便說:「對不起.」我道歉後,就沒有責任.你若不饒恕我,我還可以定你的罪.但認罪不是如此,認罪後,不是將責任推卸給神,而是仍然願負起罪的一切後果.
\newline
\newline
故真正的悔改,先有謙卑的心,且願為罪負責.不但有決心離罪,且願為罪付任何代價.不遮掩,不推卸,最後用行動將神的話行出來,不是單單聽道,還要行道.
\newline
\newline
神說,轉回吧,何必死亡呢?不推卸,不遮掩,用誠心在神面前,神正等待將福氣賜給我們.神要賜福給我們的心,比我們要得福的心更大.願神賜福大眾!
\newline
\newline


\newpage

\section{第40屆~港九培靈會~鮑會園~第八講}
\label{sec:1526}
\textbf{要心裡思想}
\newline
\newline
連結: \href{https://www.hkbibleconference.org/session-message/view/1526}{\texttt{https://www.hkbibleconference.org/session-message/view/1526}}
\newline
\newline
\hyperref[sec:1521]{< < < PREV SERMON < < <}
~
\hyperref[sec:toc]{[返目錄]}
~
\hyperref[sec:1524]{> > > NEXT SERMON > > >}
\newline
\newline
何西阿書 7:1-7:16
\newline
\begin{longtable}{cl}
\hline
\hline
章節 & 經文 (和合本修訂版)\\
\hline
7:1 & \begin{tabularx}{0.7\textwidth}{X} 我正要醫治以色列的時候，以法蓮的罪孽和撒瑪利亞的邪惡就顯露出來。他們行事虛謊，內有賊人入侵，外有群盜劫掠。 \end{tabularx} \\ \\ \relax
7:2 & \begin{tabularx}{0.7\textwidth}{X} 他們以為我不在意他們一切的惡行；現在，他們所做的在我面前纏繞他們。 \end{tabularx} \\ \\ \relax
7:3 & \begin{tabularx}{0.7\textwidth}{X} 他們行惡使君王歡喜，說謊使官長快樂。 \end{tabularx} \\ \\ \relax
7:4 & \begin{tabularx}{0.7\textwidth}{X} 他們全都犯姦淫，如同烤熱的火爐，師傅在揉麵到發麵時暫時停止煽火。 \end{tabularx} \\ \\ \relax
7:5 & \begin{tabularx}{0.7\textwidth}{X} 在我們君王宴樂的日子，官長因酒的烈性而生病，王與褻慢的人握手。 \end{tabularx} \\ \\ \relax
7:6 & \begin{tabularx}{0.7\textwidth}{X} 他們臨近，心裡如火爐一般，他們等待，如烤餅的整夜睡覺，到了早晨如火焰熊熊。 \end{tabularx} \\ \\ \relax
7:7 & \begin{tabularx}{0.7\textwidth}{X} 他們全都熱如火爐，吞滅他們的審判官。他們的君王都仆倒，他們中間無一人求告我。 \end{tabularx} \\ \\ \relax
7:8 & \begin{tabularx}{0.7\textwidth}{X} 以法蓮混居在萬民中，以法蓮是沒有翻過的餅。 \end{tabularx} \\ \\ \relax
7:9 & \begin{tabularx}{0.7\textwidth}{X} 外邦人消耗他的力量，他卻不知道；頭髮斑白，他也不覺得。 \end{tabularx} \\ \\ \relax
7:10 & \begin{tabularx}{0.7\textwidth}{X} 以色列的驕傲使自己臉面無光。他們雖遭遇這一切，仍不歸向耶和華－他們的神，也不尋求他。 \end{tabularx} \\ \\ \relax
7:11 & \begin{tabularx}{0.7\textwidth}{X} 以法蓮好像鴿子愚蠢無知，他們求告埃及，投奔亞述。 \end{tabularx} \\ \\ \relax
7:12 & \begin{tabularx}{0.7\textwidth}{X} 他們去的時候，我要把我的網撒在他們身上；我要捕獲他們如同空中的鳥。我必按他們會眾所聽到的懲罰他們。 \end{tabularx} \\ \\ \relax
7:13 & \begin{tabularx}{0.7\textwidth}{X} 他們因離棄我，必定有禍；因違背我，必遭毀滅。我雖想要救贖他們，他們卻向我說謊。 \end{tabularx} \\ \\ \relax
7:14 & \begin{tabularx}{0.7\textwidth}{X} 他們在床上呼號，卻不誠心哀求我；他們為求五穀新酒而聚集，卻背叛我。 \end{tabularx} \\ \\ \relax
7:15 & \begin{tabularx}{0.7\textwidth}{X} 我雖管教他們，堅固他們的膀臂，他們卻圖謀邪惡抗拒我。 \end{tabularx} \\ \\ \relax
7:16 & \begin{tabularx}{0.7\textwidth}{X} 他們歸向，但不是歸向至上者；終究必如鬆弛的弓。他們的領袖必因舌頭的狂傲倒在刀下，這在埃及地必成為人的笑柄。 \end{tabularx} \\ \\
[1ex]
\hline
\hline
\end{longtable}
何西阿書的分段為二,一至三章是用生活預表神對百姓的大愛,神吩咐他們所做的事,清楚的均有領表意義.
\newline
\newline
四至十四章按一至三章的註解,將整個預表詳細解釋,也是每個百姓需要學習的功課.
\newline
\newline
一至三章是藉事實表明神的愛,四至十四章是神賜給百姓直接的教訓.
\newline
\newline
當時百姓過著安定生活,國中無困難,有能幹的君王──耶羅波安第二世,忙於建設國家,恢復過去所失去的聲望,也忙於生產,以謀恢復經濟的雄厚力量,故此就疏於敬拜神.在這時百姓因沒有時間上去獻祭,就把牛羊送去便算了事.在這樣情形下,神要向百姓施教訓.
\newline
\newline
他們生活太忙,工作擔子太重,沒有時間思想,又不敢思想.在神前所犯的罪,若思想起來,便要負責,怎麼辦才好?他們在生活上,工作上太認真,所認定的是工作和生活,不是神.將整個心放在工作和生活上,沒有時間思想神.心中充滿生活和工作,沒有時間再思想屬靈的生命,沒有時間思想與主的關係.
\newline
\newline
於是神要他們思想,生命的目標是甚麼?生活的重點是甚麼?
\newline
\newline
今天在這社會裡,許多東西充滿著我們的心,生活太忙,屬靈的心便麻木了,沒有時間過問屬靈的事,所以神也要我們心裡思想.
\newline
\newline
神提出要百姓思想的事
\newline
\newline
1. 「以法蓮與列邦人攙雜；以法蓮是沒有翻過的餅.」(8節)以色列百姓的生活如此,確要思想.
\newline
\newline
以色列百姓做餅,把餅放入爐中,忘記取出,一面已燒焦,另一面卻還是生粉.他們的生活也如此,一方面非常認真,一方面卻異常疏忽；一面熱如火,一面冷像冰；一面熟識,一面忘記.
\newline
\newline
有人讀聖經,入神學院研究,也是一樣,真理認識多,但生命得造就太少.明白真理,但屬靈的生命卻是冷冰的.有些被稱為神學家,但心裡從未與主發生關係,只是頭腦知識而已,未履行為真理應付出的生活.
\newline
\newline
有人追求屬靈的事很熱心,但忽略了屬靈的根基,也是像一個未翻過的餅.
\newline
\newline
有一姊妹,過去很熱心,後來結婚了,比前更火熱,我以為是屬靈的長進,可是久而久之,我發覺她的熱心是從情感來的,從肉體來的.
\newline
\newline
很多異端信仰,錯誤教訓,就是這樣出來,他們把重點,只偏重一面.我們要有屬靈的追求,也要有真理為根基；要神的教訓,但必須與生活發生關係.
\newline
\newline
2. 被外邦人同化,失去見證能力.「以法蓮與外邦人攙雜.」當時在以色列人中,有許多亞蘭人,迦南人,亞述人雜居,以色列人效法外邦人的生活習慣,效法他們的敬拜方式.
\newline
\newline
我們雖然生活在這世界裡,但我們不屬這個世界,需要與外邦人分別.基督徒在世界中生活,是要將屬靈的生命力量彰顯出來,影響別人,將他們從黑暗中領來歸向神.但如果未有影響他們,反而被他們同化,便失去見證的能力,忘記基督徒在世上是光是鹽.教會在世上與世人一同生活,不是要人知道我們是一個大團體,也不是要人知道我們裡面有很多有學問,有金錢的人,而是讓他們知道我們生活中有屬靈的能力,影響著他們,使他們不得不歸服神.
\newline
\newline
以法蓮要思想生活的目標是甚麼?生活的意義是甚麼?
\newline
\newline
3. 屬靈生命麻木.「頭髮斑白,他也不覺得.」(9節下)這是年紀衰老的象徵,但自己還不覺得.自己的軟弱,也不知道.自己屬靈的失敗,沒有長進,愛心減少,但心靈麻木,沒有感覺.屬靈生活衰老,這是多麼可憐的現象.
\newline
\newline
今天神要我們思想,屬靈的真實情況如何?過去熱心,愛主,看見人言行不符,心中難過；看見電影廣告,心如破碎.但今天沒有感覺.過去每天若不靈修,不祈禱,心如乾旱的河.現在卻因工作忙,無聚會,無靈修,無禱告,卻處之泰然,這是屬靈生命已經衰老了.過去若缺乏愛心,心中便不平安；不饒恕人,心中也不平安；口不慎言,心靈難過,如今全不覺得,因為屬靈生命已經衰老了,失去感覺.
\newline
\newline
4. 「他們並不誠心哀求我,乃在床上呼號.」(14節)不是認罪,乃是發怨言.禱告,不是出於誠心.充滿屬靈偏見,無屬靈生命實質.有屬靈的外表,沒有屬靈的能力.叫人知道自己有祈禱,有讀經,整個生命目標是在人前,只是叫人知道.主耶穌責備這樣的人說:「按名是活的,其實是死的.」他們祈禱,為求五穀新酒,為物質而追求,為生活滿足而追求.
\newline
\newline
在這方面,神要我們思想,有兩方面:
\newline
\newline
1. 「他們中間無一人求告我.」(7節)神刑罰他們,他們不聽求告.知道自己屬靈生活太差,明知自己有需要,但不敢求告.今天要他們思想,神刑罰他們的目的,不是叫他們滅亡,也不是叫他們離開神.乃是要叫他們知道,從那裡跌倒,便要從那裡起來.
\newline
\newline
主耶穌說:「康健的人用不著醫生,有病的人才用得著.」又說:「我來不是召義人,乃是召罪人悔改.」
\newline
\newline
2. 「他們歸向,卻不歸向至上者,他們如翻背的弓.」(16節)「翻背」可譯作「詭詐」.即不誠實.射箭原來是要射別人,卻反射了自己.這原是毀滅敵人的工具,卻毀滅了自己.
\newline
\newline
我們如翻背的弓,主怎能使用.主叫他作見證,但他失去作見證的能力.主叫他領人歸主,但他卻成了歸主的人的絆腳石.如翻背的弓,不但沒有好處,反而破壞聖工.
\newline
\newline
過去為主重用,有能力,有果效,別人知道我是基督徒.但如今工作無果效,不能引人歸主,無喜樂,與主交通無甜蜜.過去為主工作,心中快樂.今則成為重擔.過去在屬靈的道路上追求,愈走愈快,愈走愈輕.今則屬靈的追求,成為重負.因為這是翻背的弓,不能為主使用.
\newline
\newline
主說,你心裡要思想,要記念.在甚麼地方跌倒,就要在甚麼地方起來.誰要回轉歸向主,便要思想,要記念.
\newline
\newline


\newpage

\section{第40屆~港九培靈會~鮑會園~第九講}
\label{sec:1524}
\textbf{以色列人的偶像}
\newline
\newline
連結: \href{https://www.hkbibleconference.org/session-message/view/1524}{\texttt{https://www.hkbibleconference.org/session-message/view/1524}}
\newline
\newline
\hyperref[sec:1526]{< < < PREV SERMON < < <}
~
\hyperref[sec:toc]{[返目錄]}
~
\hyperref[sec:1525]{> > > NEXT SERMON > > >}
\newline
\newline
何西阿書 8:1-9:17
\newline
\begin{longtable}{cl}
\hline
\hline
章節 & 經文 (和合本修訂版)\\
\hline
8:1 & \begin{tabularx}{0.7\textwidth}{X} 你用口吹角吧！敵人如鷹攻打耶和華的家；因為他們違背了我的約，干犯了我的律法。 \end{tabularx} \\ \\ \relax
8:2 & \begin{tabularx}{0.7\textwidth}{X} 他們必呼求我：「我的神啊，我們以色列認識你了。」 \end{tabularx} \\ \\ \relax
8:3 & \begin{tabularx}{0.7\textwidth}{X} 以色列丟棄良善；仇敵必追逼他。 \end{tabularx} \\ \\ \relax
8:4 & \begin{tabularx}{0.7\textwidth}{X} 他們立君王，並非出於我；立官長，我卻不知道。他們用金銀為自己製造偶像，以致被剪除。 \end{tabularx} \\ \\ \relax
8:5 & \begin{tabularx}{0.7\textwidth}{X} 撒瑪利亞啊，耶和華已拋棄你的牛犢；我的怒氣向拜牛犢的人發作。他們要到幾時方能無罪呢？ \end{tabularx} \\ \\ \relax
8:6 & \begin{tabularx}{0.7\textwidth}{X} 因這牛犢是出於以色列，是匠人所造的，並不是神。撒瑪利亞的牛犢必被打碎。 \end{tabularx} \\ \\ \relax
8:7 & \begin{tabularx}{0.7\textwidth}{X} 他們所栽種的是風，所收割的是暴風；禾稼不長穗，無以製成麵粉；即便製成，外邦人也必吞吃它。 \end{tabularx} \\ \\ \relax
8:8 & \begin{tabularx}{0.7\textwidth}{X} 以色列被吞吃，如今在列國中像人所不喜愛的器皿。 \end{tabularx} \\ \\ \relax
8:9 & \begin{tabularx}{0.7\textwidth}{X} 他們投奔亞述如獨行的野驢。以法蓮雇用情人， \end{tabularx} \\ \\ \relax
8:10 & \begin{tabularx}{0.7\textwidth}{X} 他們雇用列國；如今我要聚集他們，他們必因君王和官長所加的重擔開始衰微。 \end{tabularx} \\ \\ \relax
8:11 & \begin{tabularx}{0.7\textwidth}{X} 以法蓮為贖罪增添許多祭壇，這些祭壇卻使他犯罪。 \end{tabularx} \\ \\ \relax
8:12 & \begin{tabularx}{0.7\textwidth}{X} 我為他寫了許多條律法，他卻以為與他毫無關係。 \end{tabularx} \\ \\ \relax
8:13 & \begin{tabularx}{0.7\textwidth}{X} 他們獻祭物作為給我的供物，卻自食其肉，耶和華並不悅納他們。現在他必記起他們的罪孽，懲罰他們的罪惡；他們必返回埃及。 \end{tabularx} \\ \\ \relax
8:14 & \begin{tabularx}{0.7\textwidth}{X} 以色列忘記造他的主，建造宮殿，猶大增添許多堅固的城；我卻要降火在他的城鎮，吞滅其堡壘。 \end{tabularx} \\ \\ \relax
9:1 & \begin{tabularx}{0.7\textwidth}{X} 以色列啊，不要歡喜，像萬民一樣快樂；因為你行淫離棄你的神，喜愛各禾場上賣淫所得的賞金。 \end{tabularx} \\ \\ \relax
9:2 & \begin{tabularx}{0.7\textwidth}{X} 禾場和壓酒池都不足以餵養他們，它的新酒也必缺乏。 \end{tabularx} \\ \\ \relax
9:3 & \begin{tabularx}{0.7\textwidth}{X} 他們必不得住耶和華的地；以法蓮卻要返回埃及，在亞述吃不潔淨的食物。 \end{tabularx} \\ \\ \relax
9:4 & \begin{tabularx}{0.7\textwidth}{X} 他們必不得向耶和華獻澆酒祭，所獻的祭也不蒙悅納。他們的祭物如居喪者的食物，凡吃的必使自己玷污；因為他們的食物只為自己的口腹，必不得入耶和華的殿。 \end{tabularx} \\ \\ \relax
9:5 & \begin{tabularx}{0.7\textwidth}{X} 到盛會的日子，在耶和華的節期，你們要怎樣行呢？ \end{tabularx} \\ \\ \relax
9:6 & \begin{tabularx}{0.7\textwidth}{X} 看哪，他們要逃避災難；埃及人要收殮他們，摩弗人要埋葬他們。蒺藜盤踞他們貴重的銀器，荊棘必佔據他們的帳棚。 \end{tabularx} \\ \\ \relax
9:7 & \begin{tabularx}{0.7\textwidth}{X} 降罰的日子近了，報應的時候已經來到。以色列必知道，先知愚昧，受靈感動的人狂妄，皆因你多多作惡，大懷怨恨。 \end{tabularx} \\ \\ \relax
9:8 & \begin{tabularx}{0.7\textwidth}{X} 以法蓮替我的神守望；至於先知，他所到之處都有捕鳥人的羅網，在他神的家中也遭人懷恨。 \end{tabularx} \\ \\ \relax
9:9 & \begin{tabularx}{0.7\textwidth}{X} 他們深深敗壞，如在基比亞的日子一樣。耶和華必記起他們的罪孽，懲罰他們的罪惡。 \end{tabularx} \\ \\ \relax
9:10 & \begin{tabularx}{0.7\textwidth}{X} 我發現以色列，如在曠野的葡萄；我看見你們的祖先，如春季無花果樹上初熟的果子。他們卻來到巴力‧毗珥，獻上自己做羞恥的事，成為可憎惡的，與他們所愛的一樣。 \end{tabularx} \\ \\ \relax
9:11 & \begin{tabularx}{0.7\textwidth}{X} 以法蓮，他們的榮耀如鳥飛去，必不生產，不懷胎，不成孕； \end{tabularx} \\ \\ \relax
9:12 & \begin{tabularx}{0.7\textwidth}{X} 他們縱然將兒女養大，我卻要使他們喪子，一個也不留。我離棄他們，他們就有禍了。 \end{tabularx} \\ \\ \relax
9:13 & \begin{tabularx}{0.7\textwidth}{X} 我看以法蓮如推羅栽於美地。以法蓮卻要將自己的兒女帶出來，交給行殺戮的人。 \end{tabularx} \\ \\ \relax
9:14 & \begin{tabularx}{0.7\textwidth}{X} 耶和華啊，求你加給他們，加給他們甚麼呢？要使他們懷孕流產，乳房枯乾。 \end{tabularx} \\ \\ \relax
9:15 & \begin{tabularx}{0.7\textwidth}{X} 因他們在吉甲的一切惡事，我在那裡憎惡他們。因他們所行的惡，我必把他們趕出我的殿，不再愛他們；他們的領袖都是悖逆的。 \end{tabularx} \\ \\ \relax
9:16 & \begin{tabularx}{0.7\textwidth}{X} 以法蓮受擊打，其根枯乾，不能結果，即或生產，我也要殺他們所生的愛子。 \end{tabularx} \\ \\ \relax
9:17 & \begin{tabularx}{0.7\textwidth}{X} 我的神必棄絕他們，因為他們不聽從他；他們必飄流在列國中。 \end{tabularx} \\ \\
[1ex]
\hline
\hline
\end{longtable}
以色列人分為南北兩國,北國稱以色列,南國稱為猶大.北國從此轉離耶和華,直到滅亡.神一再指責他們的罪,他們始終不能脫離耶羅波安帶領他們拜金牛犢的罪.他們敬拜神是真誠的,但走錯方向,拜了偶像還不知是偶像,犯罪不知是在罪中生活,得罪神也不知道是得罪神,還以為是敬拜神.
\newline
\newline
第九章也是責備他們拜偶像的罪.
\newline
\newline
他們敬拜偶像的情形
\newline
\newline
1. 拜偶像出於無知.如同使徒行傳第十七章的雅典人一樣,敬拜神但不認識神,所以他們敬拜的是未識之神,所有的神,甚至從來未聽過之神都敬拜.
\newline
\newline
可見人心或許不認識神,但敬拜神的心是存在的.因為人是按神的形像,樣式而造的,故不能離開神而生存.
\newline
\newline
「他們必呼叫我說,我的神阿!我們以色列認識你了.」(八2)這句話的重點在「我」字和「你」字.他們口裡說認識神,心裡卻與神未建立關係.
\newline
\newline
若明白整個基督的信仰,和一切信條,但未與基督發生關係,也不是敬拜.
\newline
\newline
2. 為方便而造成的偶像.他們離開耶路撒冷起立北國之時,在政治上分開,所以回到耶路撒冷聖殿敬拜神不方便,耶羅波安便想出一個方便的方法,就把樹立拜牛犢的祭壇.於是他們轉離耶和華,直到亡國.
\newline
\newline
若為環境方便,而改變敬拜方式,這與聖經沒有衝突,這是錯誤的說法.
\newline
\newline
記得日本入侵華北時,我剛在那兒,日本與中國基督教成立所謂「華北基督教團」,代表所有基督教可在政府中有發言權.日本人所拜的是神道教,我們所拜的是真神.日本人提議互相參加敬拜,那時有人提出反對,有人則說為了環境,委屈一點吧,我們仍可傳福音呀!試問改變了敬拜方式,是不是真誠的敬拜.
\newline
\newline
有位神學畢業生,要到某國傳道,因領取入境證有困難.有人說,可用遊客身份申請.當時我反對,試問若果辦事人員問到那裡做甚麼?你回答說是為傳道抑回答是為遊歷?若說謊,是神所不許的.
\newline
\newline
比如說有人要在書局偷一本聖經去送給一個人,希望救他的靈魂,你說神會許可你做嗎?
\newline
\newline
隨機應變,應付環境的敬拜方式,一定不是神所喜悅的.
\newline
\newline
3. 高舉自己而立偶像.「他們投奔亞述,如同獨行的野驢,以法蓮賄買朋黨.他們雖在列邦中賄買人,現在我卻要聚集徵罪他們.」(八9-10)神興起亞述刑罰他們,利用亞蘭人攻擊他們,是要他們對付罪.但他們卻用錢去買朋友來幫助,不澈底明瞭困難的原因,只應付目前的困難.
\newline
\newline
共產主義以思想解決困難；人文主義叫人相信自己,以解決困難；納粹主義卻以武力來解決困難.這都是人想出來應付目前困難的方法.
\newline
\newline
4. 明明認識神,但另立標準.明知人類中的罪惡,死亡等問題無法解決.所以不得不承認宇宙中有權柄,不得不承認有神的掌管著.但為了應付目前困難,自己立一個評判標準.「以法蓮增添祭壇取罪.」(11節)「以色列忘記造他的主,建造宮殿.」(14節)以色列百姓知道神的尊貴,知道屬靈的事的重要,但今日情況,無法得主的喜悅,神無法賜福.便用另外的方法,增添祭壇,以為增加一個祭壇,顯出我更愛主.
\newline
\newline
或許有人說,教會增添教友,不是顯得我更愛主嗎?我的事奉工作更成功嗎?過去一間教會,今有分堂；過去是小教會,今建大教會,別人會說我更成功.可是他忘記造他的主,主重要?還是聖殿重要?不是憑外表來衡量屬靈的成就.
\newline
\newline
今天教會所需要的不是更大的禮拜堂,而是需要有愛主的弟兄姊妹；不是加添教友,而是要加添熱心事奉.
\newline
\newline
神若果祝福教會,效果多,教友多,我們應感謝主.但不可用數目代替屬靈生活的果子,以免增添教友,成為我們的偶像.
\newline
\newline
真正敬拜神的方法是甚麼?
\newline
\newline
1. 如果要真敬拜神,不能按私意,要按神的話語.神的律法包括一切,不要說,聖經沒有說不可買馬票,便可以買馬票.因為神的律法包括一切屬靈的生活.
\newline
\newline
神的話是摩西所當遵守的,也是二十世紀教會所當遵守的.不是知識少,年紀老,小孩子才遵守；我有學問便不用遵守.
\newline
\newline
2. 要出於真誠的心.屬靈的事是永遠建在心靈中的,不是建在外表上.以色列百姓只獻牛羊,不顧自己生活的糜爛,是不可以的.雖然傳福音,行異能,但若沒有屬靈生活是與主無關的.我們不能用口親近神,心卻遠離神.
\newline
\newline
3. 要出於由靈裡愛主的心.不是出於愛主的心,不能作事奉的工,主不要心裡與祂持相反的人事奉祂.
\newline
\newline
主要的不是效果,不是外表,主耶穌要的是內心愛主的人事奉祂.
\newline
\newline
神用基甸所帶的三百人,不用起初的三萬二千人,免得他們說是我有能力戰勝敵人.
\newline
\newline
4. 要真正從內心承認耶穌為主.只口裡承認主的人,不是以主耶穌為主.所以真正敬拜的人,必定以耶穌為主,不單是個人的救主,也是生活的主,家庭的主,事業上的主.因為「那真正拜父的,要用心靈和誠實拜祂.因為父要這樣的人拜祂.」(約四24)
\newline
\newline


\newpage

\section{第40屆~港九培靈會~鮑會園~第十講}
\label{sec:1525}
\textbf{心懷二意}
\newline
\newline
連結: \href{https://www.hkbibleconference.org/session-message/view/1525}{\texttt{https://www.hkbibleconference.org/session-message/view/1525}}
\newline
\newline
\hyperref[sec:1524]{< < < PREV SERMON < < <}
~
\hyperref[sec:toc]{[返目錄]}
~
\hyperref[sec:1523]{> > > NEXT SERMON > > >}
\newline
\newline
何西阿書 10:1-10:15
\newline
\begin{longtable}{cl}
\hline
\hline
章節 & 經文 (和合本修訂版)\\
\hline
10:1 & \begin{tabularx}{0.7\textwidth}{X} 以色列是茂盛的葡萄樹，結果繁多。果子越多，就越增添祭壇；土地越肥美，就越建造美麗的柱像。 \end{tabularx} \\ \\ \relax
10:2 & \begin{tabularx}{0.7\textwidth}{X} 他們心懷二意，現今要定為有罪。耶和華必拆毀他們的祭壇，粉碎他們的柱像。 \end{tabularx} \\ \\ \relax
10:3 & \begin{tabularx}{0.7\textwidth}{X} 現在他們要說：「我們沒有王；因為我們不敬畏耶和華，王又能為我們做甚麼呢？」 \end{tabularx} \\ \\ \relax
10:4 & \begin{tabularx}{0.7\textwidth}{X} 他們講空話，以假誓立約；因此，懲罰如苦菜滋生在田間的犁溝中。 \end{tabularx} \\ \\ \relax
10:5 & \begin{tabularx}{0.7\textwidth}{X} 撒瑪利亞的居民必因伯‧亞文的牛犢驚恐；它的百姓為它悲哀，它的祭司為它戰兢，因為榮耀已經離開它。 \end{tabularx} \\ \\ \relax
10:6 & \begin{tabularx}{0.7\textwidth}{X} 人必將牛犢帶到亞述，當作禮物獻給大王。以法蓮必蒙羞，以色列必因自己的計謀慚愧。 \end{tabularx} \\ \\ \relax
10:7 & \begin{tabularx}{0.7\textwidth}{X} 撒瑪利亞的王要滅亡，如水面上的泡沫一般。 \end{tabularx} \\ \\ \relax
10:8 & \begin{tabularx}{0.7\textwidth}{X} 亞文的丘壇，以色列犯罪的地方必毀壞，荊棘和蒺藜必長在他們的祭壇上。他們要向大山說：遮蓋我們！向小山說：倒在我們身上！ \end{tabularx} \\ \\ \relax
10:9 & \begin{tabularx}{0.7\textwidth}{X} 以色列啊，你從基比亞的日子以來就時常犯罪，他們仍停留在那裡。攻擊罪孽之輩的戰事豈不會臨到基比亞嗎？ \end{tabularx} \\ \\ \relax
10:10 & \begin{tabularx}{0.7\textwidth}{X} 我必隨己意懲罰他們，他們為雙重的罪所纏；萬民必聚集攻擊他們。 \end{tabularx} \\ \\ \relax
10:11 & \begin{tabularx}{0.7\textwidth}{X} 以法蓮是馴良的母牛犢，喜愛踹穀，我要將軛套在牠肥美的頸項上，我要使以法蓮被套住；猶大必耕田，雅各必耙地。 \end{tabularx} \\ \\ \relax
10:12 & \begin{tabularx}{0.7\textwidth}{X} 你們要為自己栽種公義，收割慈愛。你們要開墾荒地，現今正是尋求耶和華的時候；等他臨到，公義必如雨降給你們。 \end{tabularx} \\ \\ \relax
10:13 & \begin{tabularx}{0.7\textwidth}{X} 你們耕種奸惡，收割罪孽，吃的是謊言的果實。因你倚靠自己的行為，仰賴你眾多的勇士， \end{tabularx} \\ \\ \relax
10:14 & \begin{tabularx}{0.7\textwidth}{X} 所以在你百姓中必掀起鬧鬨，你一切的堡壘必被拆毀，就如沙勒幔在爭戰的日子拆毀伯‧亞比勒，將城中的母子一同摔碎。 \end{tabularx} \\ \\ \relax
10:15 & \begin{tabularx}{0.7\textwidth}{X} 伯特利啊，因你們的大惡，你們必遭遇如此。黎明來臨，以色列的王必全然滅絕。 \end{tabularx} \\ \\
[1ex]
\hline
\hline
\end{longtable}
神在這章聖經發出最嚴重的警告,以致百姓不敢不思想,為甚麼神的大審判要臨到我.心中一定不平安,他們想到要為罪悔改,但又捨不得過去的罪中生活.
\newline
\newline
為著將臨到的可怕的審判,他們思想要悔改,但又捨不得罪中的生活.不如不要神,但他們又不敢離開神,心靈在極度矛盾之中.神說:「他們心懷二意,現今要定為有罪.」
\newline
\newline
我們不能說,我已經信主,我不能不要主；但另方面看世界也可愛.一面抓住耶穌,一面又捨不得罪中的享受.
\newline
\newline
又愛神,又愛瑪門,今天許多人過著這樣心懷二意的生活.不敢離開神,又不肯放棄罪中生活,在神面前也要受責備,神的大審判也要臨到他們.
\newline
\newline
用三個不同象徵來形容百姓的生活
\newline
\newline
1. 「他們為立約說謊言,起假誓；因此,災罰如苦菜滋生在田間的犁溝中.」(四節)他們在對付罪的事上不專心.「苦菜」,原文是毒草.在文學著作上所稱,即哲學家蘇格拉底自殺時所吃的毒草,他吃此毒物,全身麻木,不久便死去.這是一種毒藥,但不是敵人撒在田間的,而且自己將它栽在泥溝中.不是一時軟弱犯罪,而是自己用心計劃去犯罪.並且栽培,澆灌,使罪在心中長成為一股力量.「如苦菜滋生在田間的犁溝中.」是用心栽培罪,用法安排罪.有一愛好罪的心,難怪不能對付.這是他們用心培養,可說是他們心血的結晶.
\newline
\newline
這毒物雖然危險,但想到是自己用許多心血栽培出來的,不能放棄.或許暫時收起來,有機會便放出來.
\newline
\newline
雖然為罪在神面前痛哭過,但對付罪時,自己又留一條出路.知道這罪在神面前不蒙悅納,但這罪對自己有好處,便又捨不得放棄.
\newline
\newline
不是不恨這罪,也不是不覺羞恥,只是為了這罪對自己也有好處,便捨不得離罪.
\newline
\newline
有一位弟兄,想奉獻讀神學,因不夠資格,便做一假證件,這間學院不接納,另一間學院也不接納.他對我說,禱告之後,決心認錯.以後我遇見他,問他怎樣處置那證件,他還收藏起來,有時找工作做或者需要.既已認錯,為甚麼還留下,真是心懷二意.
\newline
\newline
有一位弟兄在一救濟團體作事,常常接到有些不記名字的奉獻,便將它收起來,據為己有.在一個特別聚會中,他受感動,起來認錯,恨惡自己.但不願繼續工作,提請辭職,卻私自配了一條信箱鑰匙.暫時離開這犯罪的機會,但以後有需要時,仍可有機會.清醒時知錯,但又不肯完全斷絕.許多人滅亡,就是因為「心懷二意」.
\newline
\newline
2. 沒有專心跟從神.「如水面的沫子一樣」(七節)「沫子」原文是「刨柴」,放在水中,隨水飄流.不能忍受一點波浪,自己無力抵擋.自己想去的方向,波浪把他帶到相反的方向去.沒有穩當的基礎,未有保持的精神,不能始終不渝.
\newline
\newline
今日可能對付罪,他日環境改變,便不能對付罪.今天在培靈會裡對付罪,有追求,有熱心.但散會後回家,因家人無人追求,便也不追求.在熱心的環境上熱心,但在不好的環境上怎能熱心.
\newline
\newline
比如初打針之時有功效,過了一個時期,藥力便消失了.培靈會有如打針,但散會後,漸漸藥力便消失.
\newline
\newline
不是沒有熱心,乃是要靠環境來栽培這熱心,我的追求如環境的轉變,如水沫隨波逐流一樣.
\newline
\newline
以色列人不單對付罪不專心,連追求上也不專心.
\newline
\newline
3. 事奉上不專心.「以法蓮是馴良的母牛犢,喜愛踹穀.」(11節)神不是要馴良,是要耕種.但以色列人踹穀可,耕種則不可.「牛在場上踹穀的時候,不可籠住他的嚼.」(林前九9)踹穀雖辛苦,但可自由吃穀.但耕田便不一樣,因為耕田不可吃泥土.踹穀的時候肯熱心,努力,認真.但耕田得不到收穫,便吃不飽.不是不肯為主工作,有收穫才去為主工作.這就是踹穀可,耕田則不可.
\newline
\newline
做禮拜,有奶粉,有米派就來,做傳道,沒有人稱讚便不幹.
\newline
\newline
有一教會,經濟好,教友多,請牧師薪金高.有一次老牧師死了,登報請新牧師,有二百多人來應徵,均說是蒙神引導,難道真的神引導二百多個牧師到一間教會工作嗎?
\newline
\newline
有人神學畢業了,被派到鄉村工作,便不喜悅,說我的學院聲望大,教授學問多,我是在此畢業的,不應派我去鄉村工作.不是鄉村不好,而是教會小,薪金低.
\newline
\newline
做醫生來傳道則可,但放下醫生不做,去傳道則不可.以法蓮是馴良的母牛,踹穀可,耕種則不可.
\newline
\newline
神勸告百姓
\newline
\newline
「你們要為自己栽種公義,就能收割慈愛.現今正是尋求耶和華的時候；你們要開墾荒地,等他臨到,使公義如雨降在你們身上.」(12節)將地開墾,神可能降下秋雨春雨.未開墾荒地,降雨無用.
\newline
\newline
心靈未預備好,神怎能祝福.
\newline
\newline
荒地看來沒有盼望,但一開墾,神便降下甘霖,叫地生長植物.能長出美麗的花,能長成大莊稼.
\newline
\newline
但有人以為我的荒地,一向沒有開,別人也不開,我何必開.而且這地太硬,開墾不易,太辛苦.也許這心因剛硬,因裡面有神的敵人,在罪中仍覺得滿足,或是自己無力開,開不動.但神說,你若不開,永遠是荒地,不能作莊稼,神的福氣不能臨到你.
\newline
\newline
就算開墾地,對付罪,祈禱,讀經,有誰知道.我若去派單張,教主日學,講道,則有人稱讚.
\newline
\newline
不管甚麼,要開懇荒地,神的恩雨才能沛降.
\newline
\newline
你心懷二意,要到幾時呢?要專心對付罪,專心追求,也要專心事奉,要開墾荒地,神必將大福賜給你.
\newline
\newline


\newpage

\section{第40屆~港九培靈會~鮑會園~第十一講}
\label{sec:1523}
\textbf{慈繩愛索}
\newline
\newline
連結: \href{https://www.hkbibleconference.org/session-message/view/1523}{\texttt{https://www.hkbibleconference.org/session-message/view/1523}}
\newline
\newline
\hyperref[sec:1525]{< < < PREV SERMON < < <}
~
\hyperref[sec:toc]{[返目錄]}
~
\hyperref[sec:1522]{> > > NEXT SERMON > > >}
\newline
\newline
何西阿書 11:1-12:1
\newline
\begin{longtable}{cl}
\hline
\hline
章節 & 經文 (和合本修訂版)\\
\hline
11:1 & \begin{tabularx}{0.7\textwidth}{X} 以色列年幼的時候，我愛他，就從埃及召我的兒子出來。 \end{tabularx} \\ \\ \relax
11:2 & \begin{tabularx}{0.7\textwidth}{X} 先知越是呼喚他們，他們越是遠離，向諸巴力獻祭，為雕刻的偶像燒香。 \end{tabularx} \\ \\ \relax
11:3 & \begin{tabularx}{0.7\textwidth}{X} 我曾教導以法蓮行走，我用膀臂抱起他們，他們卻不知道是我醫治他們。 \end{tabularx} \\ \\ \relax
11:4 & \begin{tabularx}{0.7\textwidth}{X} 我用慈繩愛索牽引他們；我待他們如人鬆開牛兩腮旁邊的軛，彎下身來餵養他們。 \end{tabularx} \\ \\ \relax
11:5 & \begin{tabularx}{0.7\textwidth}{X} 他們不必返回埃及地；然而亞述人要作他們的王，因他們不肯歸向我。 \end{tabularx} \\ \\ \relax
11:6 & \begin{tabularx}{0.7\textwidth}{X} 刀劍必臨到他們的城鎮，毀壞門閂，吞滅眾人，都因他們自己的計謀。 \end{tabularx} \\ \\ \relax
11:7 & \begin{tabularx}{0.7\textwidth}{X} 我的百姓偏要背離我，他們雖向至高者呼求，他卻不抬舉他們。 \end{tabularx} \\ \\ \relax
11:8 & \begin{tabularx}{0.7\textwidth}{X} 以法蓮哪，我怎能捨棄你？以色列啊，我怎能棄絕你？我怎能使你如押瑪？怎能使你如洗扁？我回心轉意，我的憐憫燃了起來。 \end{tabularx} \\ \\ \relax
11:9 & \begin{tabularx}{0.7\textwidth}{X} 我必不發猛烈的怒氣，也不再毀滅以法蓮。因我是神，並非世人，是你們中間的聖者；我必不在怒中臨到你們。 \end{tabularx} \\ \\ \relax
11:10 & \begin{tabularx}{0.7\textwidth}{X} 耶和華如獅子吼叫，他的兒女必跟隨他。他一吼叫，他們就從西方戰兢而來。 \end{tabularx} \\ \\ \relax
11:11 & \begin{tabularx}{0.7\textwidth}{X} 他們必如雀鳥從埃及戰兢而來，又如鴿子從亞述地來到。我必使他們住自己的房屋；這是耶和華說的。 \end{tabularx} \\ \\ \relax
11:12 & \begin{tabularx}{0.7\textwidth}{X} 以法蓮用謊言圍繞我，以色列家用詭計環繞我；猶大卻仍與神同行，向聖者忠心。 \end{tabularx} \\ \\ \relax
12:1 & \begin{tabularx}{0.7\textwidth}{X} 以法蓮以風為食物，終日追逐東風，增添虛謊和殘暴，與亞述立約，也把油送到埃及。 \end{tabularx} \\ \\
[1ex]
\hline
\hline
\end{longtable}
本書第一段中心是第三章,述說神對百姓無比的大愛.然後從第四章到十四章,解釋前三章預表的意義,前後銜接,先後呼應.神用「慈繩愛索」牽引他們.何西阿傳揚神對百姓的大愛,所以有人稱何西阿為愛的先知.
\newline
\newline
全書雖充滿責備的信息,但均以愛做出發點.百姓犯罪,神因愛他,才責備他.
\newline
\newline
神對罪不輕易放過,我們可以自由選擇罪惡的道路,但無自由選擇逃避審判的道路.
\newline
\newline
神審判外邦國度,施行刑罰,各有不同.神審判埃及,如今仍存在.但刑罰亞述,亞蘭,已完全毀滅,不留痕跡.在神的恩中不歸向神,神要向他施刑罰.
\newline
\newline
百姓犯罪,神刑罰他們,但又將他們從巴比倫帶回來.他們再犯罪,神便使耶路撒冷全毀,但神應許他們復興.所以神刑罰百姓,不是毀滅,乃是管教,使他們在神面前,領會神的愛.神極愛百姓,但為要管教,便向他們施行刑罰,神的心是多麼痛苦.神管教百姓,為的是使百姓離開錯路歸向神.如果仍剛硬著心,不順從神學功課,麻木不領會,這是多麼愚拙的事.
\newline
\newline
神藉環境向香港教會講話,要香港教會學功課.去年的事便是神藉環境來管教我們.
\newline
\newline
何西阿書講到神的審判,但主要還是述說神的愛.神的審判目的,是要甦醒我們的心,是神要完成祂的旨意的工具.如果站在神的方面來看,我們便得著造就.
\newline
\newline
現在看看神怎樣向他們施行大愛
\newline
\newline
「我用慈繩愛索牽引他們；我待他們如人放鬆牛的兩腮夾板,把糧食放在他們面前.」(4節)第三章說及丈夫對妻子的愛,本章是說父親對兒女的愛,可見本書所說愛的程度,和愛的結果.
\newline
\newline
(Ａ)神愛的性質
\newline
\newline
1. 「慈繩愛索」這是父親對兒女的愛,關心對方的愛,真正看到對方的需要.正如雅各從巴旦亞蘭回來.途遇以掃,送上禮物.以掃要護送他們.但雅各對他說:「我主知道孩子們年幼嬌嫩,牛羊也正在乳養的時候,若是催趕一天,群畜都必死了.求我主在僕人前頭走,我要量著在我面前群畜和孩子的力量,慢慢的前行.」(創卅三13-14)神也知道我們的軟弱,也按著我們各人的需要來帶領我們,不會將過重的擔子放在我們身上.祂用愛,一步一步的領我們走.我們能走快,神就領我們走快些.父母一切都為兒女來計劃.但他們做錯事,父母會責罰他們,但父母的心不是因責罰兒女而快樂,反覺得非常痛苦.不過,所有責罰都是為兒女的好處.
\newline
\newline
2. 愛是神的揀選.神揀選以色列為他的百姓,不是因為百姓有甚麼好處.按勢力,神可揀選巴比倫；按豐富,神可揀選埃及；按軍事神可揀選亞述和亞蘭；按經濟,神可揀選腓尼基.可見神揀選以色列百姓,不是他們有甚麼可取的地方.神愛他們,不是他們有甚麼可愛,這才顯出神奇妙的愛,我們才被愛上.今天我們若明白神的真愛,才能改變兄弟姊妹間的一切問題.神吩咐我們彼此相愛,不是我們那一方面有可愛之處,乃是神的愛在我們中間彰顯.
\newline
\newline
神的愛,沒有先計算標準,不管愛的對象如何,還是愛他.因此,有神愛的力量,在我們裡面推動我們,我們不得不愛我們的弟兄姊妹.
\newline
\newline
神的愛,沒有為自己的好處,如果為神自己的好處,一定不揀選以色列.神的愛不受環境支配,今天我們需要這樣的愛.
\newline
\newline
(Ｂ)神愛的表現(愛的結果)
\newline
\newline
1. 「以色列年幼的時候,我愛他,就從埃及召出我的兒子來.」(十一1)神愛百姓,將他們從捆縛中釋放出來,愛使他們離開為奴的生活.神的愛使我們從罪中得釋放,得自由.神用慈繩愛索,代替了罪的捆綁,代替了為奴的捆綁.
\newline
\newline
「慈繩愛索」的愛,需要對方內心的反應.愛臨到一個人身上,這人心裡定有反應,願將自己放在這愛的力量影響之下.
\newline
\newline
愛如堅強的鎖,沒有甚麼力量可以把他打破.逼迫,勢力,貧窮,痛苦均不能對這把愛鎖有絲毫的摧毀.神的愛臨到百姓身上,使他們離開為奴的地位,給他們新自由.可是他們內心沒有甘心順從這愛的捆綁,這是百姓當時的情形.
\newline
\newline
他們離開了在埃及為奴的生活,但裡面的勢力未有離開.神可以助我們打破外力,但神不勉強我們接納祂的愛.若強迫我們接受,就不是真愛.
\newline
\newline
手放開,石下墜,不能顯出石的特點,只顯出地心吸力.金,銀,鑽石均一樣下墜,無法顯出價值.因為它們無揀選的餘地,只往下墜一途.照樣神不給我們違背的機會,就顯不出順從,顯不出真愛神.
\newline
\newline
神的愛提醒我,像輕輕牽引我一下,使我有思想,發出對這輕輕牽引有何反應.神的「慈繩愛索」輕輕牽引我心,如果我甘心順從,這愛在我裡面,便成為一股堅強不能打破的力量.
\newline
\newline
2. 但若不順從這愛的牽引,不接納這愛的提醒,可能有以下的兩個結果.
\newline
\newline
(甲)「他們必不歸回埃及地,亞述人卻要作他們的王.」(十一5)百姓離開了埃及的捆綁,但未能得到永遠的自由.神的恩典已向他們顯現,但他們未對愛的吸引有適當的反應.因他們脫離埃及的捆綁,但不肯順從神的愛的吸引,便又落亞述捆綁之中.若一個基督徒離開罪的情形說:「污鬼離了人身,就在無水之地過來過去,尋求安歇之處；既尋不著,便說,我要回到我所出來的屋裡去；到了,就看見裡面打掃乾淨,修飾好了,便去另帶了七個比自己更惡的鬼來,都進去住在那裡.那人末後的景況比先前更不好了.」(路十一26)
\newline
\newline
(乙)「以法蓮吃風,且追趕東風,時常增添虛謊和強暴,與亞述立約,把油送到埃及.」(十二1)風是易得的,有時風可充滿我,但不能使我滿足.風有力,但不能培植生命.我們可盡量吃風,但風對我們生命無價值.風可不斷吹來,但不能叫生命長進.吃風,追風,心靈都不能得到力量,滿足和長進.有人信主很久,知道真理,但內心仍是空虛.得救贖,但心中得不到喜樂,安慰,沒有享受與神甜蜜的交通.整個屬靈的生命,好像吃風,心靈永遠空虛.
\newline
\newline
可是神仍然用「慈繩愛索」輕輕牽引他,提醒他,要他有反應,要他歸回.神已將豐富的糧食放在他們面前,等他們享受.
\newline
\newline
我真願我的內心也有一個接受神愛的反應,接納神「慈繩愛索」的牽引,使我能享受神的豐富,心靈裡得著神的肥甘.
\newline
\newline


\newpage

\section{第40屆~港九培靈會~鮑會園~第十二講}
\label{sec:1522}
\textbf{當歸向你的神}
\newline
\newline
連結: \href{https://www.hkbibleconference.org/session-message/view/1522}{\texttt{https://www.hkbibleconference.org/session-message/view/1522}}
\newline
\newline
\hyperref[sec:1523]{< < < PREV SERMON < < <}
~
\hyperref[sec:toc]{[返目錄]}
~
\hyperref[sec:1514]{> > > NEXT SERMON > > >}
\newline
\newline
何西阿書 12:1-13:16
\newline
\begin{longtable}{cl}
\hline
\hline
章節 & 經文 (和合本修訂版)\\
\hline
12:1 & \begin{tabularx}{0.7\textwidth}{X} 以法蓮以風為食物，終日追逐東風，增添虛謊和殘暴，與亞述立約，也把油送到埃及。 \end{tabularx} \\ \\ \relax
12:2 & \begin{tabularx}{0.7\textwidth}{X} 耶和華指控猶大，要照雅各所行的懲罰他，按他所做的報應他。 \end{tabularx} \\ \\ \relax
12:3 & \begin{tabularx}{0.7\textwidth}{X} 他在腹中抓住哥哥的腳跟，壯年的時候與神角力， \end{tabularx} \\ \\ \relax
12:4 & \begin{tabularx}{0.7\textwidth}{X} 他與天使角力，並且得勝。他曾哀哭，懇求施恩。在伯特利遇見耶和華，耶和華在那裡吩咐我們， \end{tabularx} \\ \\ \relax
12:5 & \begin{tabularx}{0.7\textwidth}{X} 耶和華是萬軍之神，耶和華是他可記念的名。 \end{tabularx} \\ \\ \relax
12:6 & \begin{tabularx}{0.7\textwidth}{X} 所以你當歸向你的神，謹守慈愛和公平，常常等候你的神。 \end{tabularx} \\ \\ \relax
12:7 & \begin{tabularx}{0.7\textwidth}{X} 商人手持詭詐的天平，喜愛欺壓。 \end{tabularx} \\ \\ \relax
12:8 & \begin{tabularx}{0.7\textwidth}{X} 以法蓮說：我果然富有，得了財寶；我所勞碌得來的一切，人必找不到我有甚麼可算為有罪的惡。 \end{tabularx} \\ \\ \relax
12:9 & \begin{tabularx}{0.7\textwidth}{X} 自從你出埃及地以來，我就是耶和華－你的神；我必使你再住帳棚，如同節期的日子一樣。 \end{tabularx} \\ \\ \relax
12:10 & \begin{tabularx}{0.7\textwidth}{X} 我已吩咐眾先知，又增加異象，藉先知設比喻。 \end{tabularx} \\ \\ \relax
12:11 & \begin{tabularx}{0.7\textwidth}{X} 基列沒有罪孽嗎？他們誠然是虛假的，在吉甲獻牛犢為祭；他們的祭壇如同田間犁溝中的亂堆。 \end{tabularx} \\ \\ \relax
12:12 & \begin{tabularx}{0.7\textwidth}{X} 從前雅各逃到亞蘭地，以色列為娶妻子工作，為娶妻子而牧放。 \end{tabularx} \\ \\ \relax
12:13 & \begin{tabularx}{0.7\textwidth}{X} 後來耶和華藉先知領以色列從埃及上來，也藉先知看顧他們。 \end{tabularx} \\ \\ \relax
12:14 & \begin{tabularx}{0.7\textwidth}{X} 然而以法蓮大大惹動主怒，他所流的血必歸到他身上。主必使他的羞辱歸還給他。 \end{tabularx} \\ \\ \relax
13:1 & \begin{tabularx}{0.7\textwidth}{X} 從前以法蓮說話，人都戰兢，他在以色列中居處高位；但他因巴力犯罪就死了。 \end{tabularx} \\ \\ \relax
13:2 & \begin{tabularx}{0.7\textwidth}{X} 如今他們罪上加罪，為自己鑄造偶像，憑自己的聰明用銀子造偶像，全都是匠人所製的。論到它，有話說：獻祭的人都要親吻牛犢。 \end{tabularx} \\ \\ \relax
13:3 & \begin{tabularx}{0.7\textwidth}{X} 因此，他們必如早晨的雲霧，又如速散的露水，如被狂風吹離禾場的糠秕，又如煙囪冒出的煙。 \end{tabularx} \\ \\ \relax
13:4 & \begin{tabularx}{0.7\textwidth}{X} 自從你出埃及地以來，我就是耶和華－你的神；除了我神以外，你不認識別的，在我以外，並沒有救主。 \end{tabularx} \\ \\ \relax
13:5 & \begin{tabularx}{0.7\textwidth}{X} 我曾在曠野，就是那乾旱之地認識你。 \end{tabularx} \\ \\ \relax
13:6 & \begin{tabularx}{0.7\textwidth}{X} 他們得到餵養，就飽足；既得飽足，就心高氣傲，因而忘記了我。 \end{tabularx} \\ \\ \relax
13:7 & \begin{tabularx}{0.7\textwidth}{X} 因此我向他們如同獅子，又如豹伏在道旁。 \end{tabularx} \\ \\ \relax
13:8 & \begin{tabularx}{0.7\textwidth}{X} 我如失去小熊的母熊，攻擊他們，撕裂他們的胸膛。在那裡我必如母獅吞吃他們，如野獸撕開他們。 \end{tabularx} \\ \\ \relax
13:9 & \begin{tabularx}{0.7\textwidth}{X} 以色列啊，你自取滅亡了 ，因為我才是你的幫助。 \end{tabularx} \\ \\ \relax
13:10 & \begin{tabularx}{0.7\textwidth}{X} 現在，你的王在哪裡呢？讓他在你的各城中拯救你吧！你曾說「給我立君王和官長」，那些治理你的又在哪裡呢？ \end{tabularx} \\ \\ \relax
13:11 & \begin{tabularx}{0.7\textwidth}{X} 我在怒氣中將王賜給你，又在烈怒中將王廢去。 \end{tabularx} \\ \\ \relax
13:12 & \begin{tabularx}{0.7\textwidth}{X} 以法蓮的罪孽被捲起來，他的罪惡被收藏起來。 \end{tabularx} \\ \\ \relax
13:13 & \begin{tabularx}{0.7\textwidth}{X} 產婦的疼痛必臨到他身上；他是無智慧之子，如同臨盆時未出現的胎兒。 \end{tabularx} \\ \\ \relax
13:14 & \begin{tabularx}{0.7\textwidth}{X} 我必救贖他們脫離陰間，救贖他們脫離死亡。死亡啊，你的災害在哪裡？陰間哪，你的毀滅在哪裡？憐憫必從我眼前消逝。 \end{tabularx} \\ \\ \relax
13:15 & \begin{tabularx}{0.7\textwidth}{X} 他在弟兄中雖然旺盛，卻有東風颳來，就是耶和華的風從曠野上來。他的泉源必乾涸，他的源頭必枯竭，這風必奪走他所積蓄的一切寶物。 \end{tabularx} \\ \\ \relax
13:16 & \begin{tabularx}{0.7\textwidth}{X} 撒瑪利亞要擔當罪孽，因為背叛自己的神。他們必倒在刀下，嬰孩必被摔碎，孕婦必被剖開。 \end{tabularx} \\ \\
[1ex]
\hline
\hline
\end{longtable}
「從前以法蓮說話,人都戰兢,他在以色列中居處高位；但他在事奉巴力的事上犯罪就死了.」(十三1)那時,以色列百姓有敬拜神的心,尊敬神,神也尊敬他們,與他們同在,賜福他們.
\newline
\newline
以法蓮是大國,被人尊敬,但不久,約在北國建國後一百年,百姓因事奉巴力轉離神,不能不被神所棄絕.有善始,但不能保持善終,多可憐!
\newline
\newline
有人起初有追求,跟隨主覺得輕快,但一天比一天沉重,跟主不再是喜事,不再是輕省.是主改變嗎?不是的.
\newline
\newline
「所以你當歸向神.」意即歸向你最初所敬拜,所認識,所仰望的神,歸向你最初與神的關係.神就使你有高位,說話也使人戰兢,因與神有交通,為神所使用.
\newline
\newline
今天你聽聞神呼召你歸向祂,有甚麼感動?你能恢復起初與神的交通,神也能恢復最初給你的福氣.
\newline
\newline
(一)歸向一位雅各曾與祂較力的神
\newline
\newline
雅各剛從哈蘭回來,他用盡自己的方法,詭計,成就大事.空手出去,如今有妻兒,牲畜成群回來.雅各心想,我是多麼成功.憑我的力量,方法,計劃,才有今日的成功.那天晚上,妻兒,牲畜已過了河,只剩下他一人,能安靜,放下一切,神的使者──即神自己親自向雅各顯現.雅各或許以為是以掃?是迦南人?他想,無論是誰,我雅各不是輕易給人欺侮的,便與他較力,而且得勝.可是那人在此時將他的大腿窩,摸了一把,雅各的大腿窩就扭了.大腿窩是一切行動之力的所在,如今被神毀壞了.雅各所恃的,所倚靠的,雅各知道是神毀壞的.現今,他知道在神面前,不能再靠自己的力量.於是,雅各立刻要求神祝福,神是賜福的神,是掌管一切的神.如今他以神為他的倚靠,為他的幫助.
\newline
\newline
當神的話臨到我們的時候,我們便知到自己的敗壞,污穢,軟弱.但我們沒有力量離開敗壞,也無能力離開污穢,軟弱.神的話使我傷心難過,就在此,聖靈感動我,使我看見主為我預備的盼望,我在失敗中說:「主阿!救我!」這樣,像有一股力量湧進我心,整個生命改變了,前路充滿盼望,對主有火熱的心,主能幫助我一切,使我經驗主的大能.可惜,這是十年,二十年或多年前的事,現在心懷二意,像有甚麼在吸引我,攔阻我,令我看不見初時的亮光.不過,只要不再倚靠自己,向神投降,順從神的感動,神會照樣像初時那樣賜福的!
\newline
\newline
(二)不要忘記神
\newline
\newline
1. 「與天使較力,並且得勝,哭泣懇求,在伯特利遇見耶和華.耶和華萬軍之上帝在那裡曉諭我們以色列人；耶和華是他可記念的名.」(十二4-5)在創二十八章記載雅各逃難時,來到伯特利,在夢中看見天開了.雅各醒來便說:「這是神的殿,也是天的門.」清早起來,便立石作柱子,澆油在上面.後從哈蘭回迦南,不願到伯特利,而來到示劍.但神對他說,起來,往伯特利去.伯特利是神的殿,天的門.雅各曾經在那裡敬拜神,但現在已名成利就,不用倚靠神了.過去為神築壇,向神許願已忘記,連敬拜神也忘記.
\newline
\newline
有位弟兄,初信主時,每天用一小時半祈禱,讀經.後來因做工,工作忙,多時不來聚會.最近結婚了,家務漸多.他說:「過去一年來,每天最多用十分鐘讀經,祈禱.」如繼續下去,可能更少.最初有祭壇,生活也為祭壇.現今連敬拜也沒有,生活失敗.
\newline
\newline
今天,你我的生活中,有沒有祭壇,家庭的祭壇是不是荒廢了.在香港社會裡,人人生活忙碌,維持家庭中有敬拜多麼困難阿!
\newline
\newline
雅各為妻兒和牲畜事忙,那有家庭敬拜?回到迦南,他們貪愛示劍舒適的生活,可是他唯一的女兒,落在示劍人中,使以色列家蒙羞.這罪應歸誰身上?是底拿嗎?她小時常常聽道,但現今因家中已沒有祭壇了,忘記敬拜了.底拿跌倒,犯罪,誰應負責?
\newline
\newline
基督徒兒女犯罪,誰負責?今日青年人犯罪,誰負責?
\newline
\newline
歸向神吧!恢復家庭祭壇吧!使兒女知道甚麼是罪,過聖潔的生活.
\newline
\newline
(三)神呼叫雅各回伯特利
\newline
\newline
雅各就對家中的人說:「你們要除掉你們中間的外邦神,也要自潔,更換衣裳.」(創卅五2)
\newline
\newline
伯特利的神是聖潔的,所以雅各回伯特利,先潔淨自己,恢復家庭祭壇.
\newline
\newline
雅各過去在哈蘭生活,與拉班一同生活,衣食住行等習慣與哈蘭人一樣,在雅各帳幕內有哈蘭地的偶像.
\newline
\newline
基督徒住在親友當中,一切生活習慣都與他們相同,因為他們不信主.在家不能讀經祈禱,因為家中無人信主,大家圍著電視.一切生活都給不信主的人的生活習慣圈住,難怪失去了見證.
\newline
\newline
有一弟兄,自己信主,但家人不信主.他很愛主,因全家用不誠實的方法申請巴西居住.他說如自己不說謊,全家不能去.所以他在家中無法有見證.
\newline
\newline
若要恢復在神面前有見證.唯一方法,毀掉哈蘭的偶像,一切生活要聖潔,因伯特利的神,是聖潔的神.
\newline
\newline
(四)完全離開埃及
\newline
\newline
「自從你自埃及地以來,我就是耶和華你的神.」(十三9)神領百姓出埃及,是顯出神愛百姓的奇妙.救他門脫離為奴的生活,也從罪中生活得釋放,進入神恩典中得自由.
\newline
\newline
今天,若有人仍未信主,仍在埃及,被罪惡捆綁.現在給你一個機會,當歸向帶你出埃及的神.
\newline
\newline
百姓出埃及,進入曠野,仍未忘懷埃及地的黃瓜韭菜,與及鍋中的肉,以神所賜的嗎哪太簡單.要順從神,歸向神,便要放棄在埃及的生活和肉體在罪中的享受.
\newline
\newline
在曠野,神吩咐百姓造會幕,神要與他們同在.但他們不知神的同在,不能享受與神的交通.
\newline
\newline
領你出埃及到曠野的神,不但叫你離開罪中的享受,而且曠野不是最終的目的.曠野只是過渡時期.前面有流奶與蜜的美地.
\newline
\newline
今天神賜福香港教會,但最終目的,是要我們享受與祂交通,進入那甜蜜的享受.
\newline
\newline
神領百姓進入迦南,從敵人手中,搶回產業.神今天也要我們將人的靈魂仇敵手中搶回來.
\newline
\newline
神吩咐百姓歸向伯特利的神,願你也歸向這神.
\newline
\newline


\newpage

\section{第40屆~港九培靈會~鮑會園~第十三講}
\label{sec:1514}
\textbf{豐盛的生命}
\newline
\newline
連結: \href{https://www.hkbibleconference.org/session-message/view/1514}{\texttt{https://www.hkbibleconference.org/session-message/view/1514}}
\newline
\newline
\hyperref[sec:1522]{< < < PREV SERMON < < <}
~
\hyperref[sec:toc]{[返目錄]}
~
\hyperref[sec:1462]{> > > NEXT SERMON > > >}
\newline
\newline
何西阿書 14:1-14:9
\newline
\begin{longtable}{cl}
\hline
\hline
章節 & 經文 (和合本修訂版)\\
\hline
14:1 & \begin{tabularx}{0.7\textwidth}{X} 以色列啊，你要歸向耶和華－你的神，你因自己的罪孽跌倒了。 \end{tabularx} \\ \\ \relax
14:2 & \begin{tabularx}{0.7\textwidth}{X} 當歸向耶和華，用言語向他說：「求你除盡罪孽，悅納善行，我們就用嘴唇的祭代替牛犢獻上。 \end{tabularx} \\ \\ \relax
14:3 & \begin{tabularx}{0.7\textwidth}{X} 亞述不能救我們，我們不再騎馬，也不再對我們手所造的偶像說：『你是我們的神』；孤兒在你那裡得蒙憐憫。」 \end{tabularx} \\ \\ \relax
14:4 & \begin{tabularx}{0.7\textwidth}{X} 我必醫治他們背道的病，甘心愛他們，因為我向他們所發的怒氣已轉消。 \end{tabularx} \\ \\ \relax
14:5 & \begin{tabularx}{0.7\textwidth}{X} 我必向以色列如甘露；他必如百合花開放，如黎巴嫩的樹扎根。 \end{tabularx} \\ \\ \relax
14:6 & \begin{tabularx}{0.7\textwidth}{X} 他的嫩枝必延伸，他的榮華如橄欖樹，香氣如黎巴嫩的香柏樹。 \end{tabularx} \\ \\ \relax
14:7 & \begin{tabularx}{0.7\textwidth}{X} 曾住在他蔭下的必歸回，使五穀生長，他們要發旺如葡萄樹，他的名氣如黎巴嫩的酒。 \end{tabularx} \\ \\ \relax
14:8 & \begin{tabularx}{0.7\textwidth}{X} 以法蓮說：「我與偶像有何相干？」我應允他，顧念他：我如青翠的松樹，你的果實從我而來。 \end{tabularx} \\ \\ \relax
14:9 & \begin{tabularx}{0.7\textwidth}{X} 智慧人必明白這些事，聰明人必知道這一切。耶和華的道是正直的，義人行在其中，罪人卻在其上跌倒。 \end{tabularx} \\ \\
[1ex]
\hline
\hline
\end{longtable}
神藉先知責備百姓,乃要他們悔改離開罪惡歸向神.這章聖經是神對百姓極寶貴的應許.這章聖經如同兩人對話.先知代神向百姓說話,然後又轉過來,代百姓向神表露他們的心.百姓立刻有了反應,神也立刻有了反應.
\newline
\newline
神的話臨到百姓,若他們沒有反應,神的話便停止；若百姓有反應,神立刻給予應許.
\newline
\newline
神為百姓預備有豐盛的生命,但百姓必須按神所定的方法去得.
\newline
\newline
「以色列阿,你要歸向耶和華你的神.」(1節)這是很寶貴的呼召.
\newline
\newline
百姓落在困難中,是因他們的罪,不能將責任推卸給別人.今日教會太冷淡,不能說因社會太邪惡.有人遭遇太困難,便怨天尤人.若做錯了事,要承擔錯誤,要有今是昨非的態度,才有真正悔改歸向神的心.同時我們要承認無力對付罪,因為若靠自己的標準去愛主,靠自己的力量掙扎,失敗時便灰心喪志.
\newline
\newline
所以我們要說:「願你吸引我,我就快跑跟隨你.」(歌一4)主就立刻給我們力量.
\newline
\newline
求你除淨罪孽,悅納善行.
\newline
\newline
這樣,我們就把咀唇的祭代替牛犢獻上.我們不向亞述求救,不騎埃及的馬,也不再對我們手所造的說,你是我們的神,因為孤兒在你耶和華那裡得蒙憐憫.」(2至3節)這是百姓的反應.
\newline
\newline
1. 「不可向亞述求救.」過去有困難,危險,便去靠亞述,求別人的力量來幫助自己,靠別人來扶持.感謝神,現在他們不再靠亞述.
\newline
\newline
2. 「不騎埃及的馬.」表示他們不靠自己工具的力量.埃及人以馬為打仗的工具,過去以色列百姓以為亞述,亞蘭的軍隊我不怕,因為我有埃及的馬.有人以為我的軟弱,失敗,是因工具不夠,如今已知道靠馬──工具也不能得勝.
\newline
\newline
教會不興旺,有人歸咎經濟力量不雄厚.有了金錢,靈裡沒有力量也是假的.
\newline
\newline
3. 「不再對我們手所造的說,你是我的神.」(3號)即不再靠偶像.過去信仰偶像,今日已知它們不能幫助我們.
\newline
\newline
本來以色列百姓,從少便認識耶和華是神.可是入迦南地之後,迦南人有許多偶像,以色列人便以為迦南人興旺,是偶像的力量.故以色列的百姓雖仍在殿裡敬拜神,心卻追隨他們的偶像.猶太歷史家說:「當時百姓敬拜的是耶和華,但事奉的卻是迦南的偶像.」
\newline
\newline
今日許多人敬拜主耶穌,可是心中卻事奉金錢.
\newline
\newline
百姓所得的結果
\newline
\newline
「把嘴唇獻上」(2節).就是將身體獻上,當作活祭,如羅馬書十二章一節所說.嘴唇是代表心聲.過去百姓將身體獻給罪,作罪惡的奴.如今卻將身體上給神,作義的奴僕.神立刻就給予他們寶貴的應許.
\newline
\newline
1. 「我必醫治他們背道的病,甘心愛他們.因為我的怒氣向他們轉消.我心向以色列如甘露,他必如百合花開放,如利巴嫩的樹木扎根.他的枝條必延長,他的榮華如橄欖樹.」(4至6節上)換句話說,神在一切物質的需要上,豐富的賜福給他們.
\newline
\newline
2. 「他的香氣如橄欖樹.曾住在他蔭下的必歸回,發旺如五穀,開花如葡萄樹.他的香氣如利巴嫩的酒.」(6-7)
\newline
\newline
過去,以色列百姓在罪中生活,不能榮耀神,不能發香,令神受羞辱.今天他們悔改歸向神,神應許使他們發香.正如保羅說:「感謝神,常帥領我們在基督裡誇勝,並藉著我們在各處顯揚那因認識基督而有的香氣.」(林後二14)不錯,在以色列人的歷史上,似乎未看見以色列人得著這福氣,原來先知所看見的是主耶穌來了,藉著祂的靈在信徒裡面.誰帶有基督的靈,誰便帶著基督的香氣.
\newline
\newline
不信主的人,已經定罪,是滅亡的人,所以帶著死氣.信徒能帶著香氣,是因有基督的香氣在我們裡面.
\newline
\newline
保羅西拉雖被囚監牢中,仍唱詩讚美神,這就是從生命中所發出的香氣.保羅在耶路撒冷被捕,以至被囚在羅馬,但他在猶太人面前,外邦人面前,均發出基督的香氣.保羅帶著捆綁到羅馬去,途遇颶風,因有基督的香氣發出,主便將全船的人的生命賜給他.
\newline
\newline
主耶穌為神發出更大的香氣.主耶穌受洗的時候,從水裡上來,天上有聲音說:「這是我的愛子,是我所喜悅的.」還有,在世人面前,主耶穌甘心放在罪人的地位,甚至死在十字架上,為了拯救世人,因此,主耶穌整個生命都為神發出香氣.
\newline
\newline
主耶穌在客西馬尼園禱告,心裡憂愁,極其傷痛,有天使來加他的力量,他順服神,飲盡苦杯,為神發出極大的香氣.所以我們要發出香氣,必須完全順服.為神發出香氣,需要付代價,需要順服,你肯嗎?
\newline
\newline
有人想得神的祝福,但不肯付代價,不肯對付罪.要發出基督的香氣,但覺得付代價太大,多可惜!
\newline
\newline
有一位弟兄,在神學院受造就,有恩賜,有追求.但有一女朋友,對主冷淡,愛慕世界.有人勸他離開她,但他覺得太痛苦,這代價太大,不肯離開,多可惜.後來畢業了,兩人結了婚.兩年前,我看見他,他一直沒有出來服事主,只在工廠裡工作.我與他交談,他心中極其痛苦.
\newline
\newline
順服要犧牲,要放下自己的意見,有痛苦.但不順服,不能為主發出香氣.順服要代價,其實不順服也要付出代價,可能所付出的代價更大.
\newline
\newline
「誰是智慧人,可以明白這些事；誰是通達人,可以知道這一切.因為,耶和華的道是正直的；義人必在其中行走,罪人卻在其上跌倒.」(9節)
\newline
\newline


\newpage

\section{第40屆~港九培靈會~黃原素~港九培靈研經大會四十週年致詞}
\label{sec:1462}
\textbf{港九培靈研經大會四十週年致詞}
\newline
\newline
連結: \href{https://www.hkbibleconference.org/session-message/view/1462}{\texttt{https://www.hkbibleconference.org/session-message/view/1462}}
\newline
\newline
\hyperref[sec:1514]{< < < PREV SERMON < < <}
~
\hyperref[sec:toc]{[返目錄]}
~
\hyperref[sec:1549]{> > > NEXT SERMON > > >}
\newline
\newline
港九培靈研經大會今年已四十年了,在此我要講幾句感恩的話.一九二六年,廣州有一個宣教師夏令會,我曾去參加一,二次,很得造就,會場屬靈氣氛很好.會畢,我就有一個感想,若是每年有這種固定的聚會就好了,可以使教會的職員和信徒得著靈性的造就.後來我便和趙柳塘牧師,程文熙牧師,羅聖愛先生等,一同禱告,商量怎樣做.大家討論的結果,同意用「培靈會」的名稱,注重聖經真理,培養信徒的靈性為目的,後來又將這名稱改為「培靈研經大會.」
\newline
\newline
第一屆的廣州培靈研經大會是在一九二七年舉行的,當時的講員是王載義先生.許多人聽了道,就愛慕聖經.自此之後,一共開了七年,因為廣州淪陷在日軍手裡,所以該聚會就暫時停頓,光復後開了幾次,但又停了.
\newline
\newline
香港的培靈研經大會是怎樣開辦的呢?以前,有位楊少泉醫生到廣州赴會,他建議在香港也需要有這樣的聚會,於是有張吉盛先生,李求恩牧師等組織了香港培靈研經大會委辦會,因此香港的培靈研經大會就誕生了,但是講員必須先到廣州主講,待廣州的聚會完畢後,才到到香港來領會.換句話說,廣州的培靈研經大會是香港培靈研經大會的前身.每一年的聚會,神都很使用,以致叫很多人的靈性得著幫助和長進.
\newline
\newline
香港培靈研經大會到今天是四十年,我們滿意了嗎?除了向神感恩外,我們還未滿意,我們還要繼續的追求長進,因為今天的需要更大,逼迫更多,所以就更需要有這個培靈研經大會.我們盼望來聽道的信徒,聽了道得著復興,得著力量,回到自己的教會去,謙卑服事主,為主作見證.
\newline
\newline
有人說培靈研經大會開到五十年就更好了,主若未回來,我們當然會開到五十年,但是我希望不必開到五十年,在這十年內,主就回來,教會被提,這是我的心願.聖經中有三百一十八次是講到主再來,叫我們要警醒等候主的再來.
\newline
\newline
今天,我們要向神感恩,也要向神祈求,求主賜福今年的三位講員,叫他們講出的信息,能使我們得造就與幫助.
\newline
\newline


\newpage

\section{第41屆~港九培靈會~孫維廉~第一講}
\label{sec:1549}
\textbf{天上屬靈的福氣(七福氣)}
\newline
\newline
連結: \href{https://www.hkbibleconference.org/session-message/view/1549}{\texttt{https://www.hkbibleconference.org/session-message/view/1549}}
\newline
\newline
\hyperref[sec:1462]{< < < PREV SERMON < < <}
~
\hyperref[sec:toc]{[返目錄]}
~
\hyperref[sec:1471]{> > > NEXT SERMON > > >}
\newline
\newline
以弗所書 1:3-1:14
\newline
\begin{longtable}{cl}
\hline
\hline
章節 & 經文 (和合本修訂版)\\
\hline
1:3 & \begin{tabularx}{0.7\textwidth}{X} 願頌讚歸給我們主耶穌基督的父神。他在基督裡曾把天上各樣屬靈的福氣賜給我們。 \end{tabularx} \\ \\ \relax
1:4 & \begin{tabularx}{0.7\textwidth}{X} 因為他從創世以前，在基督裡揀選了我們，使我們在他面前成為聖潔，沒有瑕疵，滿有愛心。 \end{tabularx} \\ \\ \relax
1:5 & \begin{tabularx}{0.7\textwidth}{X} 他按著自己旨意所喜悅的，預定我們藉著耶穌基督得兒子的名分， \end{tabularx} \\ \\ \relax
1:6 & \begin{tabularx}{0.7\textwidth}{X} 使他榮耀的恩典得到稱讚；這恩典是他在愛子裡白白賜給我們的。 \end{tabularx} \\ \\ \relax
1:7 & \begin{tabularx}{0.7\textwidth}{X} 我們藉著這愛子的血得蒙救贖，過犯得以赦免，這是照他豐富的恩典， \end{tabularx} \\ \\ \relax
1:8 & \begin{tabularx}{0.7\textwidth}{X} 充充足足地賞給我們的。他以諸般的智慧聰明， \end{tabularx} \\ \\ \relax
1:9 & \begin{tabularx}{0.7\textwidth}{X} 照自己在基督裡所立定的美意，使我們知道他旨意的奧祕， \end{tabularx} \\ \\ \relax
1:10 & \begin{tabularx}{0.7\textwidth}{X} 要照著所安排的，在時機成熟的時候，使天上、地上、一切所有的，都在基督裡面同歸於一。 \end{tabularx} \\ \\ \relax
1:11 & \begin{tabularx}{0.7\textwidth}{X} 我們也在他裡面得了基業；這原是那位隨己意行萬事的神照著自己的旨意所預定的， \end{tabularx} \\ \\ \relax
1:12 & \begin{tabularx}{0.7\textwidth}{X} 為要使我們，這些首先把希望寄託在基督裡的人，頌讚他的榮耀。 \end{tabularx} \\ \\ \relax
1:13 & \begin{tabularx}{0.7\textwidth}{X} 在基督裡你們聽見真理的道，就是那使你們得救的福音，你們也信了他，就受了所應許的聖靈為印記。 \end{tabularx} \\ \\ \relax
1:14 & \begin{tabularx}{0.7\textwidth}{X} 這聖靈是我們得基業的憑據，直等到神的子民得救贖，使他的榮耀得到稱讚。 \end{tabularx} \\ \\
[1ex]
\hline
\hline
\end{longtable}
今日,我請各位朋友一同幻想自己在以弗所城.那裡雖然沒有汽車,飛機等新設備,但它卻是亞細亞省的首都,人口二十二萬五千.以前是一個大城,為羅馬帝國交通的中心點,宗教氣氛也很濃厚,他們敬拜女神亞底米.亞底米的廟宇,以大理石建造,為世界七大奇觀之一,人們拜女神,行為十分不潔,我們看聖經,便知道他們用魔術,又講鬼,故有好些人被鬼附著.又有人因賣亞底米銀像賺了很多錢.他們有一個可容納二萬四千人的戲園.又有一個圓形的劇場,有人和獸在那裡相搏鬥.人們喜歡看見流血殘忍的事.由此可見以弗所是一個有迷信,容易犯罪,多賺錢,多享受,充滿屬世快樂的城市.以弗所與今日大城市的情形沒有多大差別.保羅寫信給以弗所教會,你想像你是他們當中的一位信徒,保羅寫這封信的時候,你就要留心看,當心思想.因為你以前都是一個拜偶像的人,你和以弗所其他人的生活是一樣的,思想是屬世的,根本沒有想過天堂,地獄.保羅和一些同工來了,講及生命之道,起初在會堂中講,後來又到學校中,有二年多的時間,研究深奧的道理,心裡所想的,生活所表現的都有改變.今天我們的題目便是「天上屬靈的福氣」.
\newline
\newline
主對尼哥底母說了一句重要的話,記載在約三:12──「我對你們說地上的事,你們尚且不信,若說天上的事,如何能信呢?」
\newline
\newline
世界是我們的起點,我們生活在世界.祂造亞當希望他得永生,故有一棵生命樹,希望人得永生,有超然的生命,要離開世界,住在天堂.耶穌說了上面那番話,意思是人非信不可；不論懂不懂,都要憑信心接受.
\newline
\newline
神賜給我們屬地的福氣,有空氣,朋友等,可惜因著罪的緣故,都變成咒詛.因為屬世的福氣,令人敗壞；當有享受時,便連看聖經的時間都沒有了.但神有新的恩典賜給屬祂的人,地上的卻會過去,我們新的福氣是屬天的,在基督裡的,為新的族類,我為基督的後裔,在耶穌裡.我們快樂的,現在得著不能想像的大恩典.
\newline
\newline
當太空人從月球看地球,他們說很美麗,無奈不能用說話來形容.我們讀本段經文時,亦有同感.這是豐富寶貴的福氣,如泉源一樣.我們福氣很多,如多人在泉源邊汲水,取之不盡,因為實在太多.
\newline
\newline
保羅是一個聰明人,大有學問,但仍不能將這些福氣形容出來.今天只告訴你們有這些福氣,叫你們羨慕這七個福氣.求主給我們得祂的基業,又有聖靈為印記,同蒙天上的恩典.
\newline
\newline


\newpage

\section{第41屆~港九培靈會~孫維廉~第一講}
\label{sec:1471}
\textbf{七個揭開(啟示錄)}
\newline
\newline
連結: \href{https://www.hkbibleconference.org/session-message/view/1471}{\texttt{https://www.hkbibleconference.org/session-message/view/1471}}
\newline
\newline
\hyperref[sec:1549]{< < < PREV SERMON < < <}
~
\hyperref[sec:toc]{[返目錄]}
~
\hyperref[sec:1472]{> > > NEXT SERMON > > >}
\newline
\newline
啟示錄是「揭幕」的意思.這幾天我們要講主耶穌的七個揭幕.單講耶穌而不講獸和假先知.啟示錄是最後的一本書,亦是最重要,最寶貴的；若缺少此書,聖經就不完全.主耶穌是啟示錄的鑰匙,讀了使我們更深認識主.
\newline
\newline
約翰寫信給亞西亞的七個教會,一開始就介紹主耶穌基督這位三位一體的神,就是聖父,聖子,聖靈.約翰在啟示錄一章四,五節說:「神是昔在今在以後永在的,神和祂寶座前的七靈,並那誠實作見證的,從死裡首先復活,為世上元首的耶穌基督.」這裡說那昔在今在以後永在的是指聖父；寶座前的七靈是指聖靈；那誠實作見證的便是耶穌基督.在創世記,我們已見三一神的合作,創世記第一章,一至三節,第一節可見聖父──「起初神創造天地」.第二節可見聖靈──「神的靈運行在水面上」.第三節可見聖子──「神說要有光,就有了光」.這是那稱為道的耶穌.廿六節,我們見三一神開會,這是象徵的話:是一位三元的神照著祂們的形像做人.
\newline
\newline
在這最後的一卷書,約翰先介紹三一的神,最重要是救主耶穌；祂愛我們,以自己的血使我們離開罪惡,後來還從死裡復活,將來仍要再來.啟一:7──「看哪,祂駕雲降臨,眾目要看見祂,連刺祂的人也要看見祂,地上的萬族都要因祂哀哭.這話是真實的.阿們.」主耶穌升天後,有兩位穿白衣的人對門徒們說話:「加利利人哪,你們為甚麼站著望天呢!這離開你們被接升天的耶穌,你們見祂怎樣往天上去,他還要怎樣來.」我們不知祂何時回來,但時候近了.有人歡迎主來,等待那榮耀的一天,應該預備好.一剎那間,我們便得到那不死的身體.有人會懼怕,因為他們刺過主,主來對他們是可怕的.這裡介紹我們認識主.第八節,神說:「我是阿拉法,我是俄梅戛,是昔在今在以後永在的全能者」.這是指我們的救主耶穌.有一位有權威的解經家說這裡表明耶穌是神.
\newline
\newline
第一個揭幕──衪是起初
\newline
\newline
這裡耶穌說祂自己是「阿拉法,是俄梅戛,是昔在今在以後永在的全能者」,是初是終.17節──「我一看見,就仆倒在祂腳前,像死了一樣.祂用右手按著我說,不要懼怕.我是首先的,我是末後的.」廿一:6──「祂又對我說,都成了.我是阿拉法,我是俄梅戛,我是初,我是終......」廿二:13──「我是阿拉法,我是俄梅戛,我是首先的,我是末後的,我是初,我是終.」這原來是神自己的稱呼,但主耶穌用來稱呼自己.阿拉法,俄梅戛是希臘文首末的兩個字母.主用這個象徵自己是有意思的.我們最先學字母,但學後一生受用.雖然字母數目不多,卻產生幾十萬字.我們寫作,談話都缺不了這些字母.
\newline
\newline
主耶穌是簡單的,孩子也會領會,長大後一生學主耶穌都學不完.我們像在大海邊,得著的不過是幾杯水.主說祂是根本,我們要認識祂.愛因斯坦是猶太人,對主耶穌十分欽佩說:「拿撒勒的耶穌十分偉大.」
\newline
\newline
第一世紀時,信徒彼此問安時說:「耶穌是主」.他們所說的主是代表耶和華,意思說耶穌是耶和華.主是昔在,今在,永在,這道理非常重要.
\newline
\newline
(一)創造的初
\newline
\newline
歌羅西書一16-17「因為萬有都是靠祂造的,無論是天上的,地上的,能看見的,不能看見的,或是有位的,主治的,執政的,掌權的,一切都是藉著祂造的,又是為祂造的.祂在萬有之先.萬有也靠祂而立.」
\newline
\newline
在第五世紀時,有人說耶穌是被造的,比天使高,比父神低.其實祂是神所生,並非被造,是獨生子,萬有都由祂而生,因此祂是初,將來宇宙的完結,亦是由祂而成.
\newline
\newline
(二)復活的初
\newline
\newline
歌羅西書一章18節「祂也是教會全體之首,祂是元始,是從死裡首先復生的,使祂可以在凡事上居首位.」
\newline
\newline
生命由主而來,祂是首先復活的.拉撒路的復活是暫時的,因為後來他又死去,但主復活乃永遠活著,是教會之始.林後五:17──「若有人在基督裡,他就是新造的人；舊事已過,都變成新的了.」在基督裡有新的人種,信徒是由祂開始,祂是初,叫我們成為新造的人.
\newline
\newline
(三)我是首先的,是末後的
\newline
\newline
打仗時,前鋒和後裔都是十分危險的；我們的主是前是後.主耶穌來世界,曾嘗過貧窮的生活和人生的滋味.希伯來書說祂為人人嘗過死味,看四福音便知我們的主嘗到底,祂在客西馬尼園禱告,願主的旨意成就,後來祂又嘗到復活的滋味.主因著前面的榮耀就輕看這一切,祂是在我們前面作先鋒.
\newline
\newline


\newpage

\section{第41屆~港九培靈會~孫維廉~第二講}
\label{sec:1472}
\textbf{第二個揭幕 ── 在異象中的耶穌}
\newline
\newline
連結: \href{https://www.hkbibleconference.org/session-message/view/1472}{\texttt{https://www.hkbibleconference.org/session-message/view/1472}}
\newline
\newline
\hyperref[sec:1471]{< < < PREV SERMON < < <}
~
\hyperref[sec:toc]{[返目錄]}
~
\hyperref[sec:1465]{> > > NEXT SERMON > > >}
\newline
\newline
啟示錄 1:9-1:18
\newline
\begin{longtable}{cl}
\hline
\hline
章節 & 經文 (和合本修訂版)\\
\hline
1:9 & \begin{tabularx}{0.7\textwidth}{X} 我—約翰就是你們的弟兄，在耶穌裡和你們一同在患難、國度、忍耐裡有份的，為神的道，並為給耶穌作的見證，曾在那名叫拔摩的海島上。 \end{tabularx} \\ \\ \relax
1:10 & \begin{tabularx}{0.7\textwidth}{X} 有一主日我被聖靈感動，聽見在我後面有大聲音如吹號， \end{tabularx} \\ \\ \relax
1:11 & \begin{tabularx}{0.7\textwidth}{X} 說：「把你所看見的寫在書上，寄給以弗所、士每拿、別迦摩、推雅推喇、撒狄、非拉鐵非、老底嘉那七個教會。」 \end{tabularx} \\ \\ \relax
1:12 & \begin{tabularx}{0.7\textwidth}{X} 我轉過身來要看看是誰的聲音在跟我說話。我一轉過來，看見了七個金燈臺； \end{tabularx} \\ \\ \relax
1:13 & \begin{tabularx}{0.7\textwidth}{X} 在燈臺中間有一位好像人子的，身穿垂到腳的長袍，胸間束著金帶。 \end{tabularx} \\ \\ \relax
1:14 & \begin{tabularx}{0.7\textwidth}{X} 他的頭與髮皆白，如白羊毛，如雪；他的眼睛好像火焰， \end{tabularx} \\ \\ \relax
1:15 & \begin{tabularx}{0.7\textwidth}{X} 雙腳好像在爐中鍛鍊得發亮的銅，聲音好像眾水的聲音。 \end{tabularx} \\ \\ \relax
1:16 & \begin{tabularx}{0.7\textwidth}{X} 他右手拿著七顆星，從他口中吐出一把兩刃的利劍，面貌好像烈日放光。 \end{tabularx} \\ \\ \relax
1:17 & \begin{tabularx}{0.7\textwidth}{X} 我看見了他，就仆倒在他腳前，像死人一樣。他用右手按著我說：「不要怕。我是首先的，是末後的， \end{tabularx} \\ \\ \relax
1:18 & \begin{tabularx}{0.7\textwidth}{X} 又是永活的。我曾死過，看哪，我是活著的，直到永永遠遠；並且我拿著死亡和陰間的鑰匙。 \end{tabularx} \\ \\
[1ex]
\hline
\hline
\end{longtable}
這段聖經描寫主耶穌基督的榮耀.凡是超凡的,必須象徵才能明白.這裡說從主口中出來一把兩刃的利劍,是象徵的.神是永遠的道,故祂的話語有能力刺透人心.祂的榮耀像日頭一樣,口不能述說.大科學家碰到奇妙的事也要象徵,電子只不過是光速,一秒行十八萬六千英里,是不能描寫出來的,故必須象徵.
\newline
\newline
約翰見到救主,沒話可說,故以兩刃的利劍,太陽的烈光,鍛鍊過的銅,雪白的頭髮來形容主.我父親由瑞典到中國,夏天常在園中禱告,仰天說「上帝的兒子呀」,如看見祂一樣.因為他不大清楚中國的四聲,故講道困難,剛巧那時天氣又熱,他仍盡力講,因他裡面見神的榮耀,但說不出.這卻給我一個深刻的印象,我知他見了主的榮耀,故我信耶穌是真的.
\newline
\newline
約翰說看見一位像人子.約翰寫這書是在主升天後六十多年,他跟隨主三年半,看見主的死和復活,升天；又冒險跟隨主六十多年.那時他是一個老人,被囚在拔摩海島上,忽然看見主,是多年所跟從,滿有榮耀的耶穌,故他仆倒如死一般.今天我們要看四件事,是主耶穌自己說的.第十八節──「又是那存活的我曾死過,現在又活了,直活到永永遠遠.並且拿著死亡和陰間的鑰匙.」這節經文有四個意思.
\newline
\newline
(一)我活著,又是那存活的
\newline
\newline
青年人有一大問題,就是人生有甚麼意義?有人覺得人生沒有意義,故有兩條路,盡量享受或自殺.但主說祂是存活的.認識主耶穌的人,就知道人生的意義.我以前曾經絕望,等我認識主,便明白人生的意義.
\newline
\newline
我們現在再看看耶穌,起先是最卑微的,為馬利亞所生,這是「大哉敬虔的奧秘,神在人間顯現」.祂十二歲時,祂父母在聖殿中找到祂時,祂說:「你們豈不知道我要以我父的事為念麼?」祂雖然是神的兒子,但仍順服在世上的父母,回到拿撒勒.有人說耶穌的父親約瑟早死,於是家庭的責任便落在耶穌身上,順服馬利亞至三十歲才出來傳道,醫病,不住的勞苦作工.當我們看見主耶穌如何度過祂的一生,就知道人生的意義如何.祂的一生可以說是奉獻給神和人.有人整天大聲叫嚷,只想到自己及自己的利益,不想到神,又不做事榮耀神,這是一個失敗的人生.耶穌三十三歲便死了,似乎太早一點,但祂死得與別人不同,怪不得百夫長說「這真是一個義人,是神的兒子」.他驚奇主的死,因為祂的生命和人的生命不同.
\newline
\newline
(二)我曾死過
\newline
\newline
主死在十字架上,藉著死勝過死亡,這是十分合乎邏輯.科學家有一個防毒蛇咬的方法.他們先捉毒蛇,把它的頭放在杯的邊緣,取出毒液後,注射到馬的身上,取出反毒素.這便是以蛇的毒來勝過蛇毒,聖經說主藉著死勝過撒旦.人活著怕死,但死了就不怕它.我們雖然不懂,但知主死在十架上,便叫我們勝過死亡.這真是福音.
\newline
\newline
(三)現在又活了,直活到永永遠遠
\newline
\newline
我們不敢說這話,但主能,因祂從死裡復活,永不再死.世人常說「萬歲」,那不過是空談,只有主才是萬歲,直活到永遠.當救主再來時,我從死裡復活,身體改變,我便敢說:「現在又活了,直活到永永遠遠」.將來復活時,身體要變為不朽壞的.「一粒麥子,落在地裡死了,就結出許多子粒來」,沒有植物學家能解釋此事,這是神賜復活的大能.
\newline
\newline
(四)並且拿著死亡和陰間的鑰匙
\newline
\newline
世人最怕的便是死亡和陰間,但基督徒不怕這事.我們曾住過一所旅店,店主的父親半夜死了,請和尚替他念經.但我們不怕,因主耶穌拿著死亡和陰間的鑰匙.保羅說他曾到過三層天,但因為太美麗,太好,他竟說不出話來.十字架上的囚犯對主說:願你得國降臨的時候紀念我.主耶穌說,今日你便要和我同在樂園中.因祂有鑰匙,故能說此話.有一位人子,現在在天上,高過諸天.當太空人第一次旅行,我看聖經,十分高興,因為我們的主,一早已旅行過,且到父神家中,有一天我們也要去,而且勝過死亡,永遠活著.
\newline
\newline


\newpage

\section{第41屆~港九培靈會~孫維廉~第二講}
\label{sec:1465}
\textbf{第一福氣 ── 神揀選我們}
\newline
\newline
連結: \href{https://www.hkbibleconference.org/session-message/view/1465}{\texttt{https://www.hkbibleconference.org/session-message/view/1465}}
\newline
\newline
\hyperref[sec:1472]{< < < PREV SERMON < < <}
~
\hyperref[sec:toc]{[返目錄]}
~
\hyperref[sec:1466]{> > > NEXT SERMON > > >}
\newline
\newline
以弗所書 1:4-1:4
\newline
\begin{longtable}{cl}
\hline
\hline
章節 & 經文 (和合本修訂版)\\
\hline
1:4 & \begin{tabularx}{0.7\textwidth}{X} 因為他從創世以前，在基督裡揀選了我們，使我們在他面前成為聖潔，沒有瑕疵，滿有愛心。 \end{tabularx} \\ \\
[1ex]
\hline
\hline
\end{longtable}
神在創世以前,就預定揀選我們,這福氣實在太大了.
\newline
\newline
神的話安定在天,甚麼都會改變,神的話永不改變.
\newline
\newline
去年聖誕我聽廣播,聽三太空人誦讀創世記第一章,為全世界人祝福,我的心十分興奮.後來他們又讀詩篇第八篇三,四節.這表明神的話,無論何時何地,就是到月球去,一樣是人類的需要.
\newline
\newline
我講道喜歡用比方.以弗所人吃五穀魚肉,今天我人一樣吃五穀魚肉,沒有分別.這是基本的需要,數千年沒有改變.上帝的話對我們來說,也是如此.
\newline
\newline
這世界不住改變,但主的話沒有改變,主的愛沒有改變,人類犯罪的本性也沒有改變.只有物質的表面有改變.
\newline
\newline
現在讓我們思想神在基督裡賜給我們的第一福氣──祂預定揀選我們.
\newline
\newline
這福氣實在太大,我們思想它,能夠挑旺我們的信心,愛主的心和盼望.
\newline
\newline
第一,揀選的事實──這事太奇妙,我們無法推測,無法了解.為甚麼祂看中了我,揀選了我?有人不肯相信,這因為神的愛太偉大,超過我們所能夠明白的.
\newline
\newline
其實,人類太渺小,我們不能明白的事情太多,何況屬神的事?
\newline
\newline
可是我深相信,因為這是神的啟示,並且我們非得到它不可.
\newline
\newline
有許多證據,可以證明,(1)主有主權,我們在神手中,神有權柄,無人能抵禦.(2)主有聰明智慧,祂作事無人能攔阻.
\newline
\newline
我每次想到神的揀選,覺得太奇妙.世上有千千萬萬人,比較我好,主為何特地揀選我.我的智慧不多,能力不多,神卻揀選那被人輕視的,有如哥林多前書第一章廿六至卅一節所說的,這實在奇妙.
\newline
\newline
第二,揀選的條件──這條件乃是「在基督裡」.
\newline
\newline
世人都在亞當裡,自小就習慣犯罪.我五歲時就會說謊.偷吃卻說沒有,天然的說謊.打架,偷竊,生下就犯罪.
\newline
\newline
罪的結局是死,個個人都逃不了死.
\newline
\newline
美國有人死了,用法冷藏,預備有一日科學家可以醫治他的病.但人一定要死,因為這是罪的結果,無法避免.生命在主,只有主能夠叫我們「不至於死」,那是在基督裡賜給我們生命,使我們勝過死亡.好像樹枝與樹幹,肢體與身體的關係一樣.
\newline
\newline
第三,揀選的時間──叫我們聽了發昏,原來是在萬世以前.
\newline
\newline
宇宙的年齡應在一百億年以前.但時間在上帝眼中,都是「現在」,永遠是現在.神就在創世以前預定揀選我們.
\newline
\newline
我們不是宿命論,相信命運支配,我們乃相信一生的道路和日子,都在上帝的主權和能力統治下面.
\newline
\newline
挪亞造方舟,預定得救,其他不信的,不肯進入方舟,他們不得救.
\newline
\newline
第四,揀選的目的──神揀選我們,有祂高尚的目的,乃要我們「聖潔沒有瑕疵」.
\newline
\newline
有人很想到天堂,但天天犯罪,放縱情慾,如何能入天堂,入了一定也住不慣.
\newline
\newline
我看見很多信徒,日日前進得勝,也看見不少信徒落後失敗.神揀選我們,乃要我們聖潔沒有瑕疵.
\newline
\newline
聖經中有很多人看見主的榮耀,如以賽亞,但以理等,將來我們也要站在榮耀中,面對面見我們的主.到那日我們能坦然站立否?
\newline
\newline
發條必須用鐵多方煆煉,再加化學藥品煆煉,一斤可值六萬元.上帝也是如此煆煉我們,造就我們,要我們聖潔,耳,目,口,身,都聖潔沒有瑕疵,這是揀選我們的目的.
\newline
\newline


\newpage

\section{第41屆~港九培靈會~孫維廉~第三講}
\label{sec:1466}
\textbf{第二福氣 ── 得兒子的名分}
\newline
\newline
連結: \href{https://www.hkbibleconference.org/session-message/view/1466}{\texttt{https://www.hkbibleconference.org/session-message/view/1466}}
\newline
\newline
\hyperref[sec:1465]{< < < PREV SERMON < < <}
~
\hyperref[sec:toc]{[返目錄]}
~
\hyperref[sec:1473]{> > > NEXT SERMON > > >}
\newline
\newline
以弗所書 1:5-1:5
\newline
\begin{longtable}{cl}
\hline
\hline
章節 & 經文 (和合本修訂版)\\
\hline
1:5 & \begin{tabularx}{0.7\textwidth}{X} 他按著自己旨意所喜悅的，預定我們藉著耶穌基督得兒子的名分， \end{tabularx} \\ \\
[1ex]
\hline
\hline
\end{longtable}
得兒子名分是一種無始無終的福氣,是每個基督徒所共同經驗的；口說不清楚,是神豫定我們得著的.「豫定」兩字,原文作神所揀選的.神立我們作兒子,並非收留；祂要我們作祂的兒子才滿足.古時羅馬人,孩子需請老師教導,到十八九歲,身量夠長時,才配做兒子.於是父親便擺筵席,請朋友來,除了這個兒子外,大家都穿上大衣.父親在朋友面前,正式認他為兒子後,他才可穿上袍子.保羅所用的話就是這個意思.若非神的揀選,我們豈能有此福分.我們可以別的聖經來解釋此點.
\newline
\newline
「我想現在的苦楚,若比起將來要顯於我們的榮耀,就不足介意了.受造之物,切望等候神的眾子顯出來.因為受造之物服在虛空之下,不是自己願意,乃是因那叫他如此的.但受造之物仍然指望脫離敗壞的轄制,得享神兒女自由的榮耀.我們知道一切受造之物,一同歎息勞苦,直到如今.不但如此,就是我們這有聖靈初結果子的,也是自己心裡歎息,等候得著兒子的名分,乃是我們的身體得贖.」(羅八:18-23)
\newline
\newline
這是將來主回來的日子,我們得著兒子的名分,那時將會得到一個新的身體,是榮耀,不死,像主的身體.保羅在此告訴我們一個奧秘,受造之物切望等候神的眾子顯現出來,他們現在都生活在虛空之下,到那時便要得著釋放.全宇宙現在正在等待中,包括被造之物在內.因罪的緣故,把他們害了,叫他們受苦.我們等待著神的兒子顯現,像以賽亞先知所說的,那日,獅子可以和羊羔在一起.主再來的日子,我們要和祂一同顯現,你信這事嗎?天地可廢去,神的話卻不廢去.各位,這是將來的事,我們要盼望等待.
\newline
\newline
我們現在要講一些過去的事.羅馬人認的兒子,必須是自己的親生兒子.以弗所書雖沒有提及重生的事,但從別的書卷,便知道必須重生才能得兒子的名分.我們得神兒子的名份,有兩個憑據.第一,我們有神的生命,就是永生.我們原有的生命,是短暫的,但神生命是無始無終的.聖經說,信子的人有永生.我們既有神的生命,稱為祂的兒子,就在神的性情上有份.以後,我們所愛好的,思想的,都與前不同,我們只羨慕神所羨慕的.
\newline
\newline
在一千九百年前,有一件大事發生了.就是天上的神,取了人的形像,來到人間,從處女馬利亞生下來.我們現在所用的公元,就是以主降生作紀元.就是反對基督耶穌的人,也不能反對公元的年份.故主來便將世界的歷史改變了.保羅說:大哉!敬虔的奧秘.還有另一個奧秘,就是若有人重生,原文作從上而生,他就有新的生命及性情,其實,每個基督徒得救都是一個神蹟.有一個無神主義者,朋友邀他的女兒參加查經班,他一同去時,另一太太問「認識主否?」他便晝夜思想.他雖然有美麗的外貌,但他得不到心裡的平安.後來,他決志信主.他的問題雖多,但卻有心追求,故漸漸長進,成為一位敬虔的基督徒.每一個人得救,都是一個敬虔的奧秘.神的生命,在人身上工作,人有時仍會失敗,會犯罪,但事後他會認罪悔改,繼續向前走,仍為神的兒子.
\newline
\newline
讓我們談談現在的事.我們應當向著標竿直跑,準備那輝煌榮耀的將來.兒子非長大不可,我曾見有憂傷的父母,他們雖有漂亮的孩子,可惜他們身體不發育,腦子不好,不聽話,叫他的父母難過.基督徒亦有不長進的危機,故保羅曾寫信給三間教會.加拉太教會,原本很熱心,但因有人在他們當中傳講新的道理,說除信耶穌外,還要守律法,才能得救成聖,他們便改變了.保羅寫信勉勵他們說,信耶穌不但得救,還能使人完全.哥林多教會,分門別類,有屬保羅的,屬阿皮羅的,屬彼得的,但保羅說,只有主為我們死在十架上,為我們的罪流血,復活,我們都是基督的門徒.既然我們都屬基督,就應信靠祂.希伯來書是保羅寫給猶太信徒的,他們完全信主耶穌,但有人說必須不離開猶太教,因他們有美麗的聖殿,有祭司,不用受逼迫,但保羅說,不能因這些而離開基督,否則,將完全滅亡.基督徒要小心,面前有極大的榮耀,故我們必須以聖經為標準,行為合乎真道.不要被異端的風吹去,不要像嬰孩沒有主意,要長大成人,更深認識主.讓聖靈工作.主必快來,讓我們帶著得贖的身體,與主面對面,得兒子的名分.
\newline
\newline


\newpage

\section{第41屆~港九培靈會~孫維廉~第三講}
\label{sec:1473}
\textbf{第三個揭幕 ── 教會的主}
\newline
\newline
連結: \href{https://www.hkbibleconference.org/session-message/view/1473}{\texttt{https://www.hkbibleconference.org/session-message/view/1473}}
\newline
\newline
\hyperref[sec:1466]{< < < PREV SERMON < < <}
~
\hyperref[sec:toc]{[返目錄]}
~
\hyperref[sec:1467]{> > > NEXT SERMON > > >}
\newline
\newline
啟示錄 1:19-1:20
\newline
\begin{longtable}{cl}
\hline
\hline
章節 & 經文 (和合本修訂版)\\
\hline
1:19 & \begin{tabularx}{0.7\textwidth}{X} 所以，你要把所看見的事、現在的事和以後將發生的事，都寫下來。 \end{tabularx} \\ \\ \relax
1:20 & \begin{tabularx}{0.7\textwidth}{X} 至於你所看見、在我右手中的七顆星和那七個金燈臺的奧祕就是：七顆星是七個教會的使者，七個燈臺是七個教會。」 \end{tabularx} \\ \\
[1ex]
\hline
\hline
\end{longtable}
今天,站在神的立場來看,世界上最重要的組織便是教會.一千九百多年來,神注意自己的教會,我們雖然不同宗派,但不要緊,最重要的便是我們都是祂的門徒,蒙過祂的恩,是祂的肢體.集合起來便是主的身體；祂是元首,一切都由祂指揮.教會能存到今天,是一件奇妙的事.撒旦藉著很多人想毀滅教會,但不成功,因為我們天上的主是元首.就形體來說,主不在世上,因祂已在上帝的右邊,我們看不見祂.但在靈裡來說:我見祂,愛祂並尊敬祂.世人說主耶穌是一位偉大的救主,但在一千九百多年前已死了,只是精神猶存.我們的看法卻不同,我們深信主已帶著身體升天,但主也充滿教會中.祂說:若有兩三個人奉主名聚會,主便在他們當中.這當然是指靈裡說的.世界雖然忘記主,不理會主,然而,主卻常在教會中顯現,我們深知祂在我們中間.在使徒行傳中,我們見信徒和主常有密切的交通.主叫他們出發,賜福給他們,主是活潑,長活在他們中間.現在,我們要看七封信,他們寫得十分有次序.
\newline
\newline
(一)啟二:1
\newline
\newline
「你要寫信給以弗所教會的使者,說,那右手拿著七星,在七個金燈台中間行走的」.
\newline
\newline
我們要注意主的手和腳,都因被釘,留有釘痕,而且主的手和腳亦沒有休息過.祂雖在父神的右邊,但祂不分晝夜,手裡拿著七星,就是教會.祂的腳亦不停的在教會中巡行.教會是燈台,祂在燈台中走來走去,為我們加油；為我們剔亮,使燈光照得更清楚.若教會失去起初的愛心,教會一定退步.像在家庭中,孩子睡著,但母親仍走來走去,看看孩子們如何,主亦如此在教會中行走,你聽見主的腳步聲沒有?
\newline
\newline
(二)啟二:8
\newline
\newline
「你要寫信給士每拿教會的使者說,那首先的,末後的,死過又活的.」
\newline
\newline
在這裡,主耶穌要我們想兩件事,祂提醒教會,祂是神,是首先的,又是末後的.祂又提醒我們,祂是人,曾經死過,但現在又活過來.因為祂是神,故不會死,無始無終；但祂亦是人,有可能死,是有始有終的.我們的主永遠是人子,是死而又活,無始無終的.一千九百年來,教會曾為此事多次爭辯.在二,三世紀時,有人說他們信耶穌是神,非人,只不過是一個顯像.今天,有人說耶穌只是一個模範,崇高的人,不是神.士每拿教會是一個受逼迫的教會,主安慰他們,主說你們要為我的緣故受逼迫,甚至於死.主提醒他們說不要怕,祂原是神,有人的身體,跟你們一樣受逼迫,為人所恨所罵,甚至處死,但祂第三日復活.人若為主受逼迫,將來也要復活,與祂一同享受永生.
\newline
\newline
(三)啟二:12
\newline
\newline
「你要寫信給別迦摩教會的使者說,那有兩刃利劍的說.」
\newline
\newline
主耶穌說:「從祂口中有兩刃的利劍出來」,這是一所失敗的教會,有撒旦的座位,犯了巴蘭的錯誤,在道理上退步,主要以口中兩刃的利劍來對付.教會常在受逼迫的時候火熱,不受攻擊時便冷淡.教會冷淡時主便藉著祂的僕人,不客氣的責備,把兩刃的利劍放在信徒的心中,將人心中的情形顯露出來.葛培理到處傳道,都說「神如此說」,「聖經說」.每次教會復興,都因著聖經,別的方法都不成功.早二百多年前,英國十分黑暗,教會亦冷淡,衛斯理每日騎馬八九十英哩,講三次道,只是講主耶穌,起初人以臭蛋擲他,教會亦不歡迎他.他在教會的墳場,找一塊高的石碑,站在上面講道,十多年後,英國便改變,因此就得救.今日我們覺得有兩刃利劍在我們心中麼?覺得自己失敗冷淡麼?
\newline
\newline
(四)啟二:18,23
\newline
\newline
「你要寫信給推雅推喇教會的使者,說,那眼目如火焰,腳像光明銅的神之子,說......」,「我又要殺死他的黨類,叫眾教會知道,我是那察看人肺腑心腸的.並要照你們的行為報應你們各人.」
\newline
\newline
在這一段,主顯現給我們,「眼如火焰,監察人心肺腑」令人戰慄.這是一所最壞的教會.舊約的耶洗別,父親拜巴力,嫁給亞哈王,把推羅西頓的壞風俗帶來,有幾百個巴力祭司,只有以利亞先知敢對付他們.最後皇后從窗掉下死了,除了手,腳和頭骨,都給狗吃掉.是個非常可怕的教會,神知道她需要奮興.多年前,我曾親眼看見陝西教會的復興,聖靈大大的工作,很多人為自己的罪流淚,以致地氈都濕了.有人不想進入教會,但不能,因神來了,非認罪不可,主的眼如同火焰.印尼不久以前的大復興,有幾十萬人歸主.有人說教會的時代完了,現在是一個新的時代,但我不信,因我們的神是活潑的主,祂以口中的利劍及眼目的火焰監察我,這是祂的恩典,救我們的方法.
\newline
\newline
(五)啟三:1
\newline
\newline
「你要寫信給撒狄教會的使者,說,那有神的七靈,和七星的說,我知道你的行為,按名你是活的,其實是死的.」
\newline
\newline
這裡說到有七星,七靈,七台.當教會受逼迫或冷淡,主便以這「三七」來工作.啟五:6說,有羔羊,像被殺過,是神的七靈,奉差遣往普天下去.主的計劃就是將福音傳到地極,在各處地方,都有祂所揀選的人,等待聖靈.聖靈來到時,便幫助人,這是主的方法.有人說教會的宗派多,但主的工作乃然前進.南美有很多部落和新畿內亞等地,都有很多人未聞福音,聖靈感動青年人冒險入森林中傳福音,又把福音譯成文字,教那些土人.有聖靈的感動,使者便去傳道,在黑暗中成就祂的美意,主的工作就不停的進前.
\newline
\newline
(六)啟三:7
\newline
\newline
「你要寫信給非拉鐵非教會的使者說,那聖潔,真實,拿著大衛的鑰匙,開了就沒有人能關,關了就沒有人能開的,說:」
\newline
\newline
這是一所最理想的教會,十分敬虔,又結果子.有敞開的門來傳福音,主不開門,就無人能進去.
\newline
\newline
教會有二青年,雖知道往新畿內亞傳道有被殺的危險,但仍等待主開門,門開了他們便進去,土人把他們殺了,以為沒有人再去,但主卻感動幾十位青年人,他們自告奮勇的去,去把福音傳出去,有整部落信主,把他們原有的偶像,弓箭都燒掉,不再殺人,轉而事奉真神.主為我們敞開的門,非入不可,非引人信主不能.
\newline
\newline
(七)啟三:14
\newline
\newline
「你要寫信給老底嘉教會的使者,說,那為阿們的,為誠信真實見證的,在神創造萬物之上為元首的,說:」
\newline
\newline
主是信實的見證人,第一次來世是為真理作見證.祂和我們見證的方法不同,每次祂都是說我實實在在的告訴你們,從來都是「是就說是,不是就說不是」.有人說主的道理是相對的,但耶穌說恨人的就是殺人,思想上動淫念的便是犯罪.人若不重生就不能進神的國.主耶穌的道理是絕對的,但猶太人不接納祂,將祂交給彼拉多時,祂說祂是為真理作見證的,不是屬世界的,乃被釘十架.祂被猶太人棄絕,但神接納祂,把祂升到天上,坐在祂的右邊.
\newline
\newline
第七個教會是可憐的,祂拒絕主在門外,不冷不熱.有些教會,主被關在門外,教會外貌雖好看,但只有敬虔的外貌,而沒有敬虔的實意.
\newline
\newline
我們若接納主,祂就坐在我們的寶座上.若剛硬,祂有兩刃的利劍和眼目的火焰.我們要跟從主到底,將來得主的獎賞.
\newline
\newline


\newpage

\section{第41屆~港九培靈會~孫維廉~第四講}
\label{sec:1467}
\textbf{第三福氣 ── 被接納的恩典}
\newline
\newline
連結: \href{https://www.hkbibleconference.org/session-message/view/1467}{\texttt{https://www.hkbibleconference.org/session-message/view/1467}}
\newline
\newline
\hyperref[sec:1473]{< < < PREV SERMON < < <}
~
\hyperref[sec:toc]{[返目錄]}
~
\hyperref[sec:1474]{> > > NEXT SERMON > > >}
\newline
\newline
以弗所書 1:6-1:6
\newline
\begin{longtable}{cl}
\hline
\hline
章節 & 經文 (和合本修訂版)\\
\hline
1:6 & \begin{tabularx}{0.7\textwidth}{X} 使他榮耀的恩典得到稱讚；這恩典是他在愛子裡白白賜給我們的。 \end{tabularx} \\ \\
[1ex]
\hline
\hline
\end{longtable}
這節經文有幾個譯本.英文有被神接納的意思,新的英文聖經譯作在神裡面被歡迎.又有聖經譯作神以恩典對待我們.中文聖經說:恩典在愛子裡賜給我們,這裡有接納,歡迎,恩待的意思.
\newline
\newline
(一)接納的需要
\newline
\newline
心理學家告訴我們,人最需要便是被接納.有心理病的人就因為不被社會接納,故感覺到非常孤單,這樣就足夠引神經損害.這問題與我們得救有關.神在創世之前,就揀選了我們,豫定得兒子的名分.我們原是卑賤,不配,由塵土而造的.科學家說我們體質的化學成份和泥土一樣.我們是神用塵土所造的,將來必歸於塵土.我們是一無所有的.神一切的來源,萬物都是藉祂造的.故保羅說:我們一切所有的,都是神所賜.我們本是罪人,但因著神的揀選,我們成了祂的兒子.三千年前,大衛想到這問題,「我觀看你指頭所造的天,並你所陳設的月亮星宿,便說,人算甚麼?你竟顧念他.世人算甚麼?你竟眷顧他.」(詩八3-4)美國太空人雖然做了歷史空前的事,但當他們回地球時,他們說:「人算甚麼?你竟顧念他」.當人考察大自然時,亦不禁說:「人算甚麼」?
\newline
\newline
(二)接納的方法
\newline
\newline
神歡迎我們,主耶穌以「浪子回頭」的比喻(路十五:11-32)來描寫天父對待我們的心.雖然浪子將家業揮霍掉,但當他回頭時,他父親不顧他的污穢,仍抱著他親嘴,這是神接納我們的樣式,天父以祂的大愛及恩典傾倒在我們身上.
\newline
\newline
(三)接納的程度
\newline
\newline
接納客人,朋友及親屬,有不同的程度.神是在愛子裡接納我們.美國南北戰爭的時候,有一件動人的事,有一個富家的獨生子和一個貧窮的青年,在戰場上結為好朋友.那有錢的青年不幸死亡,他托那窮青年帶給家人一封信,請他們待他如對待他自己一樣,他當然可以得到一切豐富的享受.神是在愛子裡接納我們,又將豐富的恩典加在我們的身上.
\newline
\newline
主耶穌在世時,他是一個真正的人,也是神的兒子.祂為馬利亞所生而為人子,他得父神的愛心,為祂所接納,所滿意的.當耶穌三十歲受洗時,天開了,聖靈降臨,天父說:「這是我的愛子,我所喜悅的.」天門為祂而開.祂那時是完全的,無罪的,卻為我們的緣故,揀選了十字架的道路.一次祂禱告說:「父啊,願榮耀歸與你的名」.上帝說:「我已榮耀了我的名」.上帝曾三次從天上說話,表明耶穌是祂的愛子.祂因背負我們的罪,以致死在十架上時,神掩面不看,天黑地震,主耶穌說:我的神,我的神,為甚麼離棄我?後來死而埋葬,祂完成救贖的工作後,三天祂從死裡復活.四十日後,又被接升天,坐於父神之右邊.
\newline
\newline
弗一:21──「遠超一切執政的,掌權的,有能的,主治的和一切有名的.不但是今世的,連來世的也都超過了.」祂在至高之處,永遠被接納的.我們雖然不是神兒子,但主仍在愛子裡接納我們,如接納愛子一樣.我們現在帶著這必死的身體,在外表看不到神接納我們.但我們憑信心到主面前來,主就歡迎我們.
\newline
\newline
人信主後便完全改變.我在鄉下講道時,有一位窮人在城門的洞居住.得救後,他的人生便改變了.我到別處佈道時,他為我拿皮包.我們二人在主面前沒有多大的分別.我們都是得救,被主接納的人,但將來誰的榮耀大就不知道,可能他要比我大.
\newline
\newline
神的愛接納我們.故信主的人便有快樂.
\newline
\newline


\newpage

\section{第41屆~港九培靈會~孫維廉~第四講}
\label{sec:1474}
\textbf{第四個揭幕 ── 揭開七印的獅子}
\newline
\newline
連結: \href{https://www.hkbibleconference.org/session-message/view/1474}{\texttt{https://www.hkbibleconference.org/session-message/view/1474}}
\newline
\newline
\hyperref[sec:1467]{< < < PREV SERMON < < <}
~
\hyperref[sec:toc]{[返目錄]}
~
\hyperref[sec:1468]{> > > NEXT SERMON > > >}
\newline
\newline
啟示錄 5:1-5:14
\newline
\begin{longtable}{cl}
\hline
\hline
章節 & 經文 (和合本修訂版)\\
\hline
5:1 & \begin{tabularx}{0.7\textwidth}{X} 我看見坐在寶座那位的右手中有書卷，正反面都寫著字，用七個印密封著。 \end{tabularx} \\ \\ \relax
5:2 & \begin{tabularx}{0.7\textwidth}{X} 我又看見一位大力的天使大聲宣告說：「有誰配展開那書卷，揭開那七個印呢？」 \end{tabularx} \\ \\ \relax
5:3 & \begin{tabularx}{0.7\textwidth}{X} 在天上、地上、地底下，沒有人能展開、能閱覽那書卷。 \end{tabularx} \\ \\ \relax
5:4 & \begin{tabularx}{0.7\textwidth}{X} 因為沒有人配展開、閱覽那書卷，我就大哭。 \end{tabularx} \\ \\ \relax
5:5 & \begin{tabularx}{0.7\textwidth}{X} 長老中有一位對我說：「不要哭。看哪，猶大支派中的獅子，大衛的根，他已得勝，能展開那書卷，揭開那七個印。」 \end{tabularx} \\ \\ \relax
5:6 & \begin{tabularx}{0.7\textwidth}{X} 我又看見寶座和四個活物，以及長老之中有羔羊站著，像是被殺的，有七個角七隻眼睛，就是神的七靈，奉差遣往普天下去的。 \end{tabularx} \\ \\ \relax
5:7 & \begin{tabularx}{0.7\textwidth}{X} 這羔羊前來，從坐在寶座上那位的右手中拿了書卷。 \end{tabularx} \\ \\ \relax
5:8 & \begin{tabularx}{0.7\textwidth}{X} 他一拿了書卷，四活物和二十四位長老就俯伏在羔羊面前，各拿著琴和盛滿了香的金爐；這香就是眾聖徒的祈禱。 \end{tabularx} \\ \\ \relax
5:9 & \begin{tabularx}{0.7\textwidth}{X} 他們唱新歌，說：「你配拿書卷，配揭開它的七印；因為你曾被殺，用自己的血從各支派、各語言、各民族、各邦國中買了人來，使他們歸於神， \end{tabularx} \\ \\ \relax
5:10 & \begin{tabularx}{0.7\textwidth}{X} 又使他們成為國民和祭司，歸於我們的神；他們將在地上執掌王權。」 \end{tabularx} \\ \\ \relax
5:11 & \begin{tabularx}{0.7\textwidth}{X} 我又觀看，我聽見寶座和活物及長老的周圍有許多天使的聲音；他們的數目有千千萬萬， \end{tabularx} \\ \\ \relax
5:12 & \begin{tabularx}{0.7\textwidth}{X} 大聲說：「被殺的羔羊配得權能、豐富、智慧、力量、尊貴、榮耀、頌讚。」 \end{tabularx} \\ \\ \relax
5:13 & \begin{tabularx}{0.7\textwidth}{X} 我又聽見在天上、地上、地底下、滄海裡和天地間一切所有被造之物，都說：「願頌讚、尊貴、榮耀、權勢，都歸給坐在寶座上的那位和羔羊，直到永永遠遠！」 \end{tabularx} \\ \\ \relax
5:14 & \begin{tabularx}{0.7\textwidth}{X} 四活物就說：「阿們！」眾長老也俯伏敬拜。 \end{tabularx} \\ \\
[1ex]
\hline
\hline
\end{longtable}
啟示錄第五章──「我看見坐寶座的右手中有書卷,裡外都寫著字,用七印封嚴了.我又看見一位大力的天使,大聲宣傳說,有誰配展開那書卷,揭開那七印呢?在天上,地上,地底下,沒有能展開能觀看那書卷的.因為沒有配展開,配觀看那書的,我就大哭.長老中有一位對我說,不要哭.看哪,猶大支派中的獅子,大衛的根,他已得勝,能展開那書卷,揭開那七印.我又看見寶座與四活物並長老之中,有羔羊站立,像是被殺過的,有七角七眼,就是神的七靈,奉差遣往普天下去的.這羔羊前來,從坐寶座的右手裡拿了書卷.他既拿了書卷,四活物和二十四位長老,就俯伏在羔羊面前,各拿著琴,和盛滿了香的金爐.這香就是眾聖徒的祈禱.他們唱新歌,說,你配拿書卷,揭開七印.因為你曾被殺,用自己的血從各族各方,各民各國中買了人來,叫他們歸於神,又叫他們成為國民,作祭司,歸於神.在地上執掌王權.我又看見,且聽見,寶座與四活物並長老的周圍,有許多天使的聲音.他們的數目有千千萬萬.大聲說,曾被殺的羔羊,是配得權柄,豐富,智慧,能力,尊貴,榮耀,頌讚的.我又聽見,在天上,地上,地底下,滄海裡和天地間一切所有被造之物,都說,但願頌讚,尊貴,榮耀,權勢,都歸給坐寶座的和羔羊,直到永永遠遠.四活物就說,阿們,眾長老也俯伏敬拜.」
\newline
\newline
這段經文中的話多是象徵性的.第四章約翰親自到天上,親眼看見神的榮耀.舊約時,摩西想看見神,但不能,神將他放在磐石穴中,要祂過後,才看見神的背.(出卅三20-23)當摩西從山上下來,面上發光,因為他看見神的榮耀.以賽亞在聖殿見主的榮耀時；便說:「禍哉,我滅亡了.」但以理在河邊,以西結在巴比倫曾見神的榮耀.他們站立不住.我們若要看見神,只有在耶穌基督身上看見,因為神的榮耀從耶穌表現出來.
\newline
\newline
在第五章,我們看到一位像獅子又像羔羊的.約翰說:神手中拿著一卷羊皮的書卷,有七印封著,是非常重要的.似乎和解決世界問題的事情有關.所以要在天上,地上和地底下找一位來開這七印.這裡有三個意思.
\newline
\newline
(一)沒有一個
\newline
\newline
起初他們找不到一位配揭開七印,於是約翰大哭.他一生事奉主,看見撒旦的權勢,人們犯罪墮落,如何能解決這個問題呢?基督徒要關心這個世界.有些人到地獄,犯罪一點不傷心,這是失敗.世上有很多人想解決這個問題,但他們沒有能力.美國競選總統,有一位州長說,美國不需要總統,只需要一位彌賽亞來對付世界的問題.
\newline
\newline
(二)有一位
\newline
\newline
長老中有一位告訴約翰,叫他不要哭,因為有一位是猶大的獅子,他能解決這問題.創世記四十九章8-12節.「猶大阿,你弟兄們必讚美你,你手必捏住仇敵的頸項,你父親的兒子們必向你下拜.猶大是個小獅子.我兒阿,你抓了食便上去,你屈下身去,臥如公獅,蹲如母獅,誰敢惹你.......」雅各有十二個兒子,流便最大,但神卻揀選了猶大,要從他支派中找一個王.雅各的預言便是指此而言.
\newline
\newline
幾百年後,神差先知往伯利恆,要為祂預備一位王.耶西大孩子很俊美,但祂沒有看上,祂只揀選牧羊的大衛,他雖然是年青,卻敢和獅子摔交,有小獅子的膽量；少壯時,又能以小石塊打倒巨人歌利亞.多年來又帶以色列軍隊打敗敵人,他很聰明,是一隻老獅子.但大衛終於年老死去,他不能長作獅子,雅各的預言沒有完全臨到大委身上.
\newline
\newline
雅各的預言歸到耶穌基督的身上,祂是永遠的獅子,又是王.這表明祂是神.每次聖經提到耶穌都有兩方面.第一第二章提及祂是神又是人；第五章提及祂是獅子,又是羔羊.看四福音,祂能叫風平浪靜,可見祂是神.但他在船上睡著,休息,表明祂是人,因為神不會打盹.祂是大衛的根,因為祂是神,祂是大衛的枝,能開花,結果.從創世記到啟示錄是一貫的,主令我們驚奇,使我們愛祂,拜祂.
\newline
\newline
(三)唯一的那位
\newline
\newline
約翰沒有看見獅子,因為獅子是神,不能看見.神在人裡面,故不見獅子,只見羔羊.羔羊像被殺,說出耶穌為人,祂被釘,在手,腳及肋旁留下死的痕跡,耶穌好像被殺,但祂復活了.
\newline
\newline
本來獅子與羔羊不能一同存在,但神是聖潔,公義,祂的本性是獅子亦是羔羊.天上的神取了人的身體,使本來無限的成為有限,充滿了宇宙.這是神救人的奧秘,為我們信仰的根本.能打開那七印的便是那為我們被殺的那位,將來祂要掌管世界.神在過去幾千年中,要我們明白這道理.啟示錄有二十八次提及羔羊.創世記第四章,曾提及羔羊.始祖犯罪,被逐出樂園後,若不藉著血就不能與主親近.該隱和神發生衝突,因為他不肯以羔羊來獻祭.
\newline
\newline
創世記二十二章,有奇妙的事.亞伯拉罕年老得子,但神要他將以撒獻上為祭,這是象徵基督十字架的道理.另一方面,以撒是長子,長子代表家庭,若不以血代贖便會死.但在最緊急時,神呼叫亞伯拉罕不要殺以撒,祂已為他預備了羔羊為祭.這羔羊代替了以撒．以色列家的人都要以羔羊的血染在門楣上,以免長子被殺.可見救贖非用血不可.
\newline
\newline
主耶穌到了三十歲,便離開馬利亞,出去傳道.到了約但河,施洗約翰說有一位在祂以後來的,他不配和祂解鞋帶,是神的兒子.施洗約翰在門徒面前,指著耶穌說:「這是神的羔羊,背負世人罪孽的」.祂正是那位能代負我罪的羔羊.
\newline
\newline
美國心理學家著書指出精神病的來源,大多數有心病的人都覺得自己有罪,需要有人替他除罪,贖罪.衛斯理及他弟弟講救贖的道理,因為除此無別的方法把罪除去,使人得釋放.約翰從此,知耶穌是羔羊.第五章說到羔羊的愛如此愛我們,救贖了我們.啟五9-10「他們唱新歌說,你配拿書卷,配揭開七印.因為你曾被殺,用自己的血從各族,各方,各民,各國中買了人來,叫他們歸於神.」這裡包括我們在內.主救我們,使我們有平安,有永生,有得救的把握.四年前,我回美探老朋友,他住在養老院中,他雖然學識少,但他卻滿面笑容,他確知主救了他.
\newline
\newline
世界上宗教多,人生觀亦不同,但有一件事,除耶穌外,無一人能叫我們得救和平安.在一個廣播節目中,有一位牧師到芝加哥,他在火車站,拿著錄音機,分別問人有關得救,神及死後的問題,很多人答不來,但遇到基督徒時,就有答案.因為他們有得救,赦罪的憑據.這是羔羊的愛.
\newline
\newline
啟示錄六章15-17節:「地上的君王,臣宰,將軍,富戶,壯士,和一切為奴的,自主的,都藏在山洞和巖石穴裡.向山和巖石說,倒在我們身上吧,把我們藏起來,躲避坐寶座者的面目,和羔羊的忿怒因為他們忿怒的大日到了,誰能站得住呢?」主在世上也有忿怒,曾怒目視人,以繩為鞭,兩次潔淨聖殿.將來更可怕的,便是這拒絕,神的世界要面對羔羊的忿怒.
\newline
\newline
啟示錄七章15-17節:「所以他們在神寶座前,晝夜在祂殿中事奉祂,坐寶座的要用帳幕覆庇他們.他們不再飢,不再渴.日頭和炎熱,也必不傷害他們.因為寶座中的羔羊必牧養他們,領他們到生命水的泉源.神也必擦去他們一切的眼淚.」將來天堂千萬信徒跟從羔羊,一切的眼淚都要擦乾,只有喜樂,這是羔羊的眷顧.沒有不能解決的問題.主說「小子呀,女兒呀!放心吧,你的罪赦了.」這是羔羊所說的.
\newline
\newline


\newpage

\section{第41屆~港九培靈會~孫維廉~第五講}
\label{sec:1468}
\textbf{第四福氣 ── 赦免我們的罪}
\newline
\newline
連結: \href{https://www.hkbibleconference.org/session-message/view/1468}{\texttt{https://www.hkbibleconference.org/session-message/view/1468}}
\newline
\newline
\hyperref[sec:1474]{< < < PREV SERMON < < <}
~
\hyperref[sec:toc]{[返目錄]}
~
\hyperref[sec:1475]{> > > NEXT SERMON > > >}
\newline
\newline
以弗所書 1:7-1:8
\newline
\begin{longtable}{cl}
\hline
\hline
章節 & 經文 (和合本修訂版)\\
\hline
1:7 & \begin{tabularx}{0.7\textwidth}{X} 我們藉著這愛子的血得蒙救贖，過犯得以赦免，這是照他豐富的恩典， \end{tabularx} \\ \\ \relax
1:8 & \begin{tabularx}{0.7\textwidth}{X} 充充足足地賞給我們的。他以諸般的智慧聰明， \end{tabularx} \\ \\
[1ex]
\hline
\hline
\end{longtable}
我們講到父在愛子裡接納我們,但罪的問題若不解決,便不能得到.為了解決這問題,神又給我們另一個恩典──赦免我們的罪.「我們藉這愛子的血,得蒙救贖,過犯得以赦免,乃是照祂豐富的恩典.」
\newline
\newline
(一)赦免的需要
\newline
\newline
每一個人,甚至基督徒都有求神赦罪的需要.戴德生是神忠心的僕人,中國內地會的創辦人,他每晚都在神面前,求神赦免他每日所犯的罪.人犯罪可分三方面來看:
\newline
\newline
(1)我們本身是罪人,故我們常犯罪.雖然大衛王愛上帝,但他仍犯大罪.故他寫詩篇第五十一篇,省察過去,求神赦罪.他說:「我是在罪孽裡生的.在我母親懷胎的時候,就有了罪」.我們是罪人,故孩子不用人教,便會說謊,自私,貪心等.
\newline
\newline
(2)人自己揀選犯罪.有時,我們不服從神,不要神管治我們,不以聖經為標準.我們自己揀選道路,自己作主,結果犯罪變壞.有人年老時,回顧一生就傷心,痛悔.因他們所揀選的成了罪.
\newline
\newline
(3)罪的奴僕.起初,人犯罪以為是一種享受,但結果卻被罪轄制,而為罪的奴僕.
\newline
\newline
無論我們所犯的罪是體面的或是難看的,它們的根源都是一樣,所以結果也是一樣.基督徒要離開罪,才能與主同享天上的福樂.
\newline
\newline
世人多不肯認罪.他們認為若把環境改好,人自然會變好.只要各人生活一樣,就沒有罪的存在.他們怕提起「罪」,否認罪的存在.他們以為罪只不過是小毛病,或者把罪歸於別人,像亞當犯罪就推諉夏娃.人亦常找藉口,或自己想辦法來掩飾自己的過犯.有人賺很多錢,但虧負了別人,於是便做善事,或苦待自己.來彌補自己的過錯.但他們除不了罪,心中沒有平安.
\newline
\newline
幾十年前,有位老先生住在廟宇中,每日只喝水及吃松子,但心中沒有平安.有位牧師問他,你得到了平安沒有?便和他同讀約翰三章十六節「神愛世人,甚至將祂的獨生子,賜給他們,叫一切信祂的,不至滅亡,反得永生.」又告訴他天上的神已為我們預備了赦罪之恩.便離開拜偶像的生活,信主得救.世人都需神親自來救我們.
\newline
\newline
有一位在大學讀書的美國小姐,他不信神存在,一切只不過是偶然而成,沒有道德,只有相對.有一天,他犯了罪,雖然他不認罪,以為沒有律法,但他失去了心中的平安,覺得心裡有重擔,他懷疑自己的人生觀錯了,最後他說,恨不得在宇宙間有一位可赦免她的罪.
\newline
\newline
(二)赦免的方法
\newline
\newline
這是一個十分寶貴的方法,祂藉愛子的血,為我們的罪,釘在十字架上.祂心破碎,為我們流盡全身的血.流血是代表捨命.這恩典是七個恩典中最當中的一個.我們單效法主是不能得救.
\newline
\newline
美國報紙,曾報導英國倫敦大學哲學系的主任祖博士,他和哲學家羅素及歷史學家 H. G. WELL 常談世界問題.他們是人民主義者,故稱人能自救,並不接納人是罪人的道理.第一次大戰後,大受打擊,覺得人不甚好,但仍希望人能找到一些自救的方法,二十年後,看見第二次大戰及納粹黨行為,知道自己以前想錯了,人畢竟是敗壞,不能自救的.雖然科學,哲學及教育很發達,但人仍然是敗壞墮落的.他覺得勝經所講的對,便信主,極力為主作見證.神救贖我們用了很大的代價,祂以寶血及那無限的生命來換取我們,故祂能赦免我們的罪.當癱子被放到主面前時,祂看見他內心的重擔及沒有平安,所以對他說:「放心,你的罪赦了」.這是一句最寶貴的話.有一個女人犯了很多罪,但當她認罪,主說:「女兒,放心吧,你的罪赦了」.赦了,我們的罪就像雲彩消散.有很多人本來憂愁,但信主後,因罪得赦,便很快樂,與主和好.
\newline
\newline
(三)赦免的程度
\newline
\newline
「這恩典是神用諸般的智慧聰明,充充足足賞給我們的.」(弗一8)
\newline
\newline
神赦免我們的罪非一件簡單的事.記得以前在鄉下佈道時,有一個青年人問我為何要浪費時間,金錢,天天講神愛世人.若講一句赦罪的話,豈不簡單嗎?但這宇宙是道德的宇宙,而神是道德的根基,是公義,是光.因此,祂不能違背本性,只說赦免我們.若如此行,祂便失去公義的本性.神的十字架是一項十分聰明的恩典,祂擔當我們的罪,出代價,乃願我們得救.聰明的父母,會用智慧來管教兒女,並不放縱他們.神的赦免也有二個條件,(1)叫我們看見祂公義,聖潔的本性.(2)世人必要認罪,主才赦免.我們必須有此二條件,才能得到那永遠的福氣.
\newline
\newline
到天堂見主時,我們能見祂釘痕的手及肋旁的槍傷.我若有機會與主說話,就必感謝主,因祂赦免我的罪,又改變了我,我們要常感謝祂,因祂的恩典豐富偉大.
\newline
\newline
基督徒的危險,就是犯了罪便隱藏起來,神曾應許我們說:「祂是信實,公義,我們若認罪,祂就必赦免我們,洗淨我們一切的不義」.當我們一犯罪,就立刻求主赦免,這是非常要緊.若要維持屬靈的生活,就要做到這一點.若將我們的失敗和過犯藏在心內,沒有和主算賬,就成為我們心中的包袱.有好些信徒的失敗,便是犯小罪沒清賬,故越積越多,而成為與主間的障礙物.故我們要每日認罪,靠主得勝.
\newline
\newline


\newpage

\section{第41屆~港九培靈會~孫維廉~第五講}
\label{sec:1475}
\textbf{第五個揭幕 ── 信徒禱告的效力}
\newline
\newline
連結: \href{https://www.hkbibleconference.org/session-message/view/1475}{\texttt{https://www.hkbibleconference.org/session-message/view/1475}}
\newline
\newline
\hyperref[sec:1468]{< < < PREV SERMON < < <}
~
\hyperref[sec:toc]{[返目錄]}
~
\hyperref[sec:1469]{> > > NEXT SERMON > > >}
\newline
\newline
啟示錄 8:1-8:6
\newline
\begin{longtable}{cl}
\hline
\hline
章節 & 經文 (和合本修訂版)\\
\hline
8:1 & \begin{tabularx}{0.7\textwidth}{X} 羔羊揭開第七個印的時候，天上寂靜約有半小時。 \end{tabularx} \\ \\ \relax
8:2 & \begin{tabularx}{0.7\textwidth}{X} 我看見那站在神面前的七位天使，有七枝號賜給他們。 \end{tabularx} \\ \\ \relax
8:3 & \begin{tabularx}{0.7\textwidth}{X} 另有一位天使拿著金香爐來，站在祭壇旁邊；有許多香賜給他，要和眾聖徒的祈禱一同獻在寶座前的金壇上。 \end{tabularx} \\ \\ \relax
8:4 & \begin{tabularx}{0.7\textwidth}{X} 那香的煙和眾聖徒的祈禱從天使的手中一同升到神面前。 \end{tabularx} \\ \\ \relax
8:5 & \begin{tabularx}{0.7\textwidth}{X} 天使拿著香爐，盛滿了壇上的火，倒在地上；就有雷轟、響聲、閃電、地震。 \end{tabularx} \\ \\ \relax
8:6 & \begin{tabularx}{0.7\textwidth}{X} 拿著七枝號筒的七位天使預備好要吹號。 \end{tabularx} \\ \\
[1ex]
\hline
\hline
\end{longtable}
讀經:啟示錄八章1-6節:「羔羊揭開第七印的時候,天上寂靜約有二刻.我看見那站在神面前的七位天使,有七枝號賜給他們.另有一位天使拿著金香爐,來站在祭壇旁邊,有許多香賜給他,要和眾聖徒的祈禱,一同獻在寶座前的金壇上.那香的煙,和眾聖徒的祈禱,從天使的手中一同升到神面前.天使拿著香爐,盛滿了壇上的火,倒在地上.隨有雷轟,大聲,閃電,地震.拿著七枝號的七位天使.就豫備要吹.」
\newline
\newline
這一段經文:好像說一位天使,並不是說到神.但解經家以為這是說我們的主.這裡最主要是禱告的信息.今天教會最缺少禱告,世界的情形亦十分嚴重.我們必須知道信徒禱告的力量和效果十分大,但需要有忍耐,亦要明白禱告的原則.
\newline
\newline
(一)天上的寂靜
\newline
\newline
寂靜本來有休息的意思,但這裡帶有等待之意.天上有等待,因為將有重要的事情發生.天上寂靜無聲,為要等待神在世上完成祂的美意.
\newline
\newline
另一種寂靜的等待,是在戰爭之前；在未到開戰的時候,戰場是寂靜的,所謂暴風雨前的寂靜.這裡天上的寂靜亦有此意.
\newline
\newline
(二)天使
\newline
\newline
這裡說到另外有一位天使,拿著金香爐.這位天使是指我們的救主耶穌.在聖經,天使有好幾個解釋,最簡單的便是天上屬靈的,無罪,圍繞著寶座服侍主的天使.希伯來書的作者說他賜福給我們.在客西馬尼園,他賜力量主耶穌的身體,當主在曠野四十晝夜,受試探後,天使便來服侍祂.他亦眷顧愛主的僕人.在第二章及三章,有七星在主手中,有七教會的使者.天使亦是耶和華的使者,曾向摩西,亞伯拉罕,大衛等顯現,祂是未降世為人子的耶穌.當人見祂時,便說看見耶和華神.
\newline
\newline
主是我們的大祭司,在我們天父的右邊為我們禱告.我們的名字刻在祂心版和手上,這裡提到世上的聖徒有三件事.
\newline
\newline
受苦的聖徒──啟六:9-11「揭開第五印的時候,我看見在祭壇底下,有為神的道,並為作見證被殺之人的靈魂.大聲喊著說,聖潔真實的主啊,你不審判住在地上的人給我們伸流血的冤,要等到幾時呢?於是有白衣賜給他們各人.有話對他們說,還要安息片時,等著一同作僕人的,和他們的弟兄,也像他們被殺,滿足了數目.」
\newline
\newline
信徒在地上受苦是一千九百年來的事實,被鞭打,被迫害,被殺戮,這是受苦的信徒.
\newline
\newline
受印記的信徒──啟七:1-3「以後我看見四位天使站在地的四角,執掌地上四方的風,叫風不吹在地上,海上和樹上.我又看見另有一位天使,從日出之地上來,拿著永生神的印.他就向那得著權柄能傷害地和海的四位天使,大聲喊著說,地與海並樹木,你們不可傷害,等我們印了我們,神眾僕人的額.」
\newline
\newline
奇妙的神,把信徒印上了印記,又應許受印記的信徒,雖受苦,但不用怕,因為人能傷害我們,不能殺害我們的靈魂,因為連我們的頭髮,我們的天父也曾數過.有一天,我們要復活,並得著榮美.
\newline
\newline
在地上禱告的信徒──啟八:3-4「另有一位天使拿著金香爐,來站在祭壇旁邊.有許多香賜給他,要和眾聖徒的祈禱一同獻在寶座前的金壇上.那香的煙,和眾聖徒的祈禱,從天使手中一同升到神面前.
\newline
\newline
我有一位朋友每日禱告,至為主而死.我知有一天我會和他相見.神應許我們,天上有一位為我們禱告.我們到天堂,全憑這一位大祭司.
\newline
\newline
(三)香
\newline
\newline
這是象徵的話.神吩咐摩西要做一種特別的香料,只在聖所用.若大祭司擅用他種香料,他必被處死.祭司進聖所時,要帶著血,另一手拿香爐.這裡預表我們的主耶穌帶著自己的血進去,這亦是象徵的.
\newline
\newline
主耶穌完全滿足神的心,以致神曾打開天,對耶穌說:「這是我的愛子,我所喜悅的」.香料便有這個意思.主耶穌是完全的,聖潔,公義,又沒有瑕疵.我愛主的名,如雅歌所說:如倒出來的香膏.
\newline
\newline
(四)禱告
\newline
\newline
天上寂靜,因為地上要受審判,天使吹號,因為神的國近了,信徒要與主同作王.有些人祈禱的範圍太狹窄,只為自己,我們應該像主耶穌教我們禱告時說:「願人都尊你的名為聖,願你的國降臨」.這是一個大的禱告,因它包括整個世界.我們實在需要神的國降臨.我很喜歡旅行禱告,在美國時,為亞洲教會禱告.在亞洲時為美國的教會禱告,由中部至東西岸,願神祝福他們.
\newline
\newline
我們禱告,沒有主便無效.我們常加一句:奉主的名求.我們這樣說是有理由的,因主耶穌是我們的大祭司,將我們的禱告和主耶穌的美德,合在一起,像把香爐的香加上去,一同升到神的面前.
\newline
\newline
(五)地上有雷轟,閃電,地震
\newline
\newline
這是說禱告有效力.神最愛聽的禱告,便是救恩臨到世上,奮興臨到信徒.故信徒為這事熱心祈禱,必定有效驗.我曾為奮興祈禱,又親眼看見復興.有一位老姊妹,本來不識字,但信主後,不但會讀經,更會唱讚美詩.後來有饑荒,但他的臉仍像天使一樣.這是復興的效果.在三百年前,有一位青年在印第安人中作工,雖然天氣寒冷,但他仍迫切禱告至流汗,後來人心便改變過來,接受耶穌.這是禱告的效果.
\newline
\newline
將來這些事必會臨到,若世人不悔改,必有審判,求主仍降恩典,有一天,祂必回來.我們一方面要警告人,另一方面我們必須禱告.我們常用主禱文,但我們可曾想到它的意思,是十分真實的.願主的國降臨.若主不來,我們便沒有盼望.將來主要再來,世上的國要永遠成為祂的國.
\newline
\newline


\newpage

\section{第41屆~港九培靈會~孫維廉~第六講}
\label{sec:1469}
\textbf{第五福氣 ── 知道神的美意}
\newline
\newline
連結: \href{https://www.hkbibleconference.org/session-message/view/1469}{\texttt{https://www.hkbibleconference.org/session-message/view/1469}}
\newline
\newline
\hyperref[sec:1475]{< < < PREV SERMON < < <}
~
\hyperref[sec:toc]{[返目錄]}
~
\hyperref[sec:1476]{> > > NEXT SERMON > > >}
\newline
\newline
以弗所書 1:9-1:10
\newline
\begin{longtable}{cl}
\hline
\hline
章節 & 經文 (和合本修訂版)\\
\hline
1:9 & \begin{tabularx}{0.7\textwidth}{X} 照自己在基督裡所立定的美意，使我們知道他旨意的奧祕， \end{tabularx} \\ \\ \relax
1:10 & \begin{tabularx}{0.7\textwidth}{X} 要照著所安排的，在時機成熟的時候，使天上、地上、一切所有的，都在基督裡面同歸於一。 \end{tabularx} \\ \\
[1ex]
\hline
\hline
\end{longtable}
「都是照祂自己所豫定的美意,叫我們知道祂旨意的奧秘,要照所安排的,在日期滿足的時候,叫天上地上一切所有的,都在基督裡面,同歸於一.」(弗一9-10)
\newline
\newline
我們首先要知道神是有計劃的.我們從「照祂自己所豫定的」,「照所安排的」,知道創造天地宇宙的神是有計劃的,有祂自己的美意.世人都有自己的計劃,成功了他們便快樂.但這些都是小規模的.我們天父的計劃卻是偉大而美好,要祂眾子們明白祂的計劃,知道祂旨意的奧秘.神的旨意就是奧秘,有許多青年人問「為何有宇宙」?「為何有我」?連人生的短暫也是一個奧秘.若果沒有聖經,我們就得不著答案.神藉著聖經告訴我們,將來我們必會完全明白.神創造世界,見到一切所造的都好；但不好的,就是撒旦來了,它使人犯罪,令我們進入滅亡.罪在那裡,那裡就有敗壞.有一點可以證明的,就是一個健康的人,他犯了法,可能隱瞞沒有人知,但他的精漸漸崩潰起來,可能引致神經病.有一位老師寫信請葛培理牧師幫忙,使說若早五年,他可能不會寫信來,因為那時他一無所缺,但他現在內心感覺空虛,他不想活下去,醫生們亦無法醫治他.心內的空虛,乃是因為失敗,和神的關係脫了節.人犯了罪,整個宇宙都籠罩在罪惡裡.連下等動物也看見人的敗壞,怕了人類,這就是罪的結果.不過天父已為我們豫備了解決的方法,記載在羅馬書第八章十八至廿五節.「被造之物都在嘆息勞苦」.這都因為我們和基督脫節的緣故.我們的手脫了節,不能說沒有相干；因為會覺得痛苦,當人和神脫節,就會犯罪,覺得痛苦.
\newline
\newline
神有一個美滿的計劃──「要照所安排的,在日期滿足的時候,使天上地上一切所有的,都在基督裡面同歸於一.」(弗一10)神的計劃就是在日期滿足的時候,使我們同歸於一.人在世界是短暫的,但神卻是永遠的.祂從不急忙,到了恰當的時候,神就會成就祂的計劃,如加拉太書四章四節所說的一樣.主是一切的中心,祂也是人類,歷史的中心.在創世前有一段永遠的時間,稱為亙古；將來也有一段永遠的時間,就是主國降臨之時.而主卻在以兩段永遠的當中,成為歷史的中心.但以理書說主快來,但沒有人知彌賽亞何時來.到了時候,世界黑暗,宗教很多,但卻不能解決人的問題,人渴望主來,那時,主就來了.主未來是因為時間未到,我們還未把福音傳遍天下,但今天有許多電台都在傳播神的話,由此可見主必快來.神好像一個輪子的中心,一切都在祂裡面,到時,我們便在主內同歸於一.
\newline
\newline
很多科學家都想知道將來如何,故他們飛昇月球,希望尋找到宇宙的來源,以便推測將來情形如何.但聖經告訴我們,將來撒旦必被打敗,被鎖起來,不能再在世界作惡,那時主的國度要降臨,再沒有戰爭混亂的事,也沒有危險或緊張的生活,只有主賜的平安.聖經說所有的兵器有一天都要改成和平的用途,以賽亞書乃舊約的福音,它說將來有一個黃金時代,是天地要聯合,人與人,人與獸都要和好.那是我們所羨慕的一天,充滿著神豐滿的恩典.
\newline
\newline
主的福音乃和平的福音,能給人心有平安.有一個人家中有問題,夫婦不和,差不多要離婚,但那位太太聽了福音後,便信主,跟著就有改變,家庭就重新有幸福,有平安.從前有一個船主及船員,在外地犯法,被逐出境,漂到一個小島上,與島上的土女通婚,後因自相殘殺,只剩一船員生還.在一偶然機會,找到一本聖經,看後,他的生命起了改變,他信主悔改.從此便和自己的太太及其他船員的太太和孩子們同讀神的話.使這島成為一個最理想的社會.這是一個比較大規模的例證,後來有人到那島,寫書說及那島.又有土人未信主,未聞福音,傳道人被他們殺了,但仍有人去傳音,終於他們改變了,全族人都信主,燒去他們的箭和符.他們沒有聖經,沒有文字,乃由別人教他們,他們結繩以知何時為星期日,又學背經文及唱詩等.
\newline
\newline
將來主回來時,世界就會改變,在主裡面便得救,這是聖經上的真理,主必快來,讓我們天天一同盼望,等候主.
\newline
\newline


\newpage

\section{第41屆~港九培靈會~孫維廉~第六講}
\label{sec:1476}
\textbf{第六個揭幕 ── 再來的主}
\newline
\newline
連結: \href{https://www.hkbibleconference.org/session-message/view/1476}{\texttt{https://www.hkbibleconference.org/session-message/view/1476}}
\newline
\newline
\hyperref[sec:1469]{< < < PREV SERMON < < <}
~
\hyperref[sec:toc]{[返目錄]}
~
\hyperref[sec:1464]{> > > NEXT SERMON > > >}
\newline
\newline
啟示錄 19:11-19:16
\newline
\begin{longtable}{cl}
\hline
\hline
章節 & 經文 (和合本修訂版)\\
\hline
19:11 & \begin{tabularx}{0.7\textwidth}{X} 後來我看見天開了。有一匹白馬，騎在馬上的稱為「誠信」、「真實」，他審判和爭戰都憑著公義。 \end{tabularx} \\ \\ \relax
19:12 & \begin{tabularx}{0.7\textwidth}{X} 他的眼睛如火焰，頭上戴著許多冠冕；他身上寫著一個名字，除了他自己沒有人知道。 \end{tabularx} \\ \\ \relax
19:13 & \begin{tabularx}{0.7\textwidth}{X} 他穿著浸過血的衣服；他的名稱為「神之道」。 \end{tabularx} \\ \\ \relax
19:14 & \begin{tabularx}{0.7\textwidth}{X} 眾天軍都騎著白馬，穿著又白又潔淨的細麻衣跟隨他。 \end{tabularx} \\ \\ \relax
19:15 & \begin{tabularx}{0.7\textwidth}{X} 有利劍從他口中出來，用來擊打列國。他要用鐵杖管轄他們，並且要踹全能神烈怒的醡酒池。 \end{tabularx} \\ \\ \relax
19:16 & \begin{tabularx}{0.7\textwidth}{X} 在他衣服和大腿上寫著「萬王之王，萬主之主」的名號。 \end{tabularx} \\ \\
[1ex]
\hline
\hline
\end{longtable}
這裡我們看見天開,世界和天上有一個敞開的交通.
\newline
\newline
今天我們看見地上也是寂靜無聲的,因為有重大的事情即將來到.在前文,我們見巴比倫被毀滅,這代表世界在背叛神,又有災難臨到世界,但在天上卻有最大的快樂.救主耶穌,就是神的羔羊,救贖了自己的百姓,當人數滿足時,羔羊便要婚娶,在天上有大筵蓆.這是宇宙中最快樂的一天,因為救贖的計劃完成了.以弗所書五章說:神將教會賜給基督,像新婦穿著白色的衣服,聖潔無有瑕疵.但那時世界仍有問題未解決,等候基督從天降臨.
\newline
\newline
主第一次來世界,被釘在十架之上,從早上九時至十二時,有群眾嘈雜的聲音,釘釘的聲音,但主安靜忍受,又有許多人說嘲笑的話.直到正午,天忽然黑了,只有人低哼之聲,主說了幾句話:「我渴了」,「我的神,我的神,為甚麼離棄我?」大自然都寂靜無聲,因宇宙的主死在十架之上.
\newline
\newline
主第二次來亦如此,我們要注意騎白馬的人,祂是主耶穌基督,祂在聖經裡有二百多個名字.世上得勝的人都騎馬,象徵凱旋得勝.預言中亦提過,早三千年,大衛作詩在詩篇四十五篇1-7節.「我心裡湧出美辭,我論到我為王作的事.我的舌頭是快手筆.你比世人更美,在你嘴裡滿有恩惠,所以,神賜福給你,直到永遠,大能者啊,願你腰間佩刀,大有榮耀和威嚴.為真理,謙卑,公義,赫然坐車前往,無不得勝,你的右手必顯明可畏的事.你的箭鋒快,射中王敵之心,萬民仆倒在你以下.神阿,你的寶座是永永遠遠的,你的國權是正直的.你喜愛公義,恨惡罪惡,所以神,就是你的神,用喜樂油膏你,勝過膏你的同伴.」這一段前面好像說一個主,後面又似乎在說一個神.這位騎白馬是一位創造天地的主.祂騎驢駒進入耶路撒冷,應驗了先知撒迦利亞的話:「耶路撒冷啊,你們要高興,因為王來了,非常的謙卑,騎驢進來」.古時,要登位的王必要騎驢,摩西叫他們不可騎馬,因為外邦人都是騎馬,故他們騎驢.有人不喜歡主再來,亦不愛聽主再來.你愛主再來嗎?若不愛,可能你心中有麻煩,我們可以來試驗一下自己.
\newline
\newline
主第一次來,人不歡迎,希律派兵迫害,文士聖經雖熟,知基督生在伯利恆,(彌迦書五章二節),但說完就忘記,只有東方的三博士去找.我們應當羨慕主再來.
\newline
\newline
再來的主祂是誠信,真實和公義.世界實在須要一個這樣的人.主再來的工作就是審判,祂眼光如火焰,踹神的酒醡,就是審判之意.主再來亦有戰爭,祂的衣服濺滿了血,這不是十字架的血,這是毀滅仇敵的血.祂再來亦要轄管,祂有鐵杖擊打列國,這審鄰非來不可.
\newline
\newline
我們喜歡知道一些關乎耶穌的事,就是祂的名.祂有一個名字是沒有人知的.名字在聖經內是本性的意思.耶穌的名字有救主,先知,道路等意思.有一些事我們永遠不會知道,就是到天堂也不會完全知道的.世上的人,我們也不能完全明白.耶穌實在太豐富,我們到天堂亦不能看清楚耶穌.另一個名字就是神的道.約翰福音第一章說,「太初有道,道與神同在,道就是神」.道就是道路,但約翰所說的道亦是話.耶穌將神的奧妙偉大顯明出來.祂從開始就將神說出來.我們若認識耶穌就認識神,祂將永生的道顯出來,祂是萬王之王,萬主之主.
\newline
\newline
騎馬來的人並非單獨的來,有天上的眾軍陪同,他們身穿白衣,全都騎馬,他們都是公義,良善,無罪,有千千萬萬.這是一個進侵.常有人想到太空的侵犯.啟示錄有兩個侵犯,一個從天上來,一個從地底來.啟九1-3──「第五位天使吹號,我就看見一個星從天落到地上.有無底坑的鑰匙賜給他.他開了無底坑,便有煙從坑裡往上看,好像大火爐的煙.日頭和天空,都因這煙昏暗了.有蝗蟲從煙中出來飛到地上,有能力賜給他們,好像地上蠍子的能力一樣」.蝗蟲是從無底坑上來,邪靈要來侵犯這個世界.主第一次來時,邪靈曾多次侵犯,被主驅逐.這是的而且確的事.約翰說末世時亦有邪靈從地裡起來侵犯世界,這情形已經有了.世界的道德墮落到何種地步呢?衣服裝飾,美術都是奇怪的.英國有一個美術家,常畫新潮派的畫,信主後便改畫神所造自然的東西.因為邪靈侵犯,故世界不安,神讓邪靈短時侵佔,最後要從天上進侵.看帖撒羅尼迦後書二章便會明白.人不接受真理,神要任他們受欺騙.日子到了,我們應該儆醒.天上萬軍要來,除掉罪惡,使世界變成和平,變成一個新的世界.故此我們要豫備自己,盡力工作,除清罪惡,努力傳福音.過去有好些歐美的青年出去傳道,現在亦有不少亞洲的青年「自告奮勇」地為主作工.工作等主來,便有賞賜,免得空手見主.
\newline
\newline
有一個青年人,鋼琴彈得很好,她受感動要去非洲傳道,但不願意.牧師告訴她,當主叫彼得宰了污穢的物吃時,他說「主啊,不能.」其實他說「主」,就不可說「不能」；因為他是僕人.他用一張紙寫上「主」和「不能」二字,叫她仔細思想,然後畫去其中一個.終於,這位姊妹畫去「不能」二字,順服著往非洲傳道.我們的主,不久必再來,祂是我們的主,每個基督徒不能對主說「不可」.願我們對主說:「主啊,我肯,願照你的旨意行」.
\newline
\newline


\newpage

\section{第41屆~港九培靈會~孫維廉~第七講}
\label{sec:1464}
\textbf{第六福氣 ── 得神的基業}
\newline
\newline
連結: \href{https://www.hkbibleconference.org/session-message/view/1464}{\texttt{https://www.hkbibleconference.org/session-message/view/1464}}
\newline
\newline
\hyperref[sec:1476]{< < < PREV SERMON < < <}
~
\hyperref[sec:toc]{[返目錄]}
~
\hyperref[sec:1470]{> > > NEXT SERMON > > >}
\newline
\newline
以弗所書 1:11-1:12
\newline
\begin{longtable}{cl}
\hline
\hline
章節 & 經文 (和合本修訂版)\\
\hline
1:11 & \begin{tabularx}{0.7\textwidth}{X} 我們也在他裡面得了基業；這原是那位隨己意行萬事的神照著自己的旨意所預定的， \end{tabularx} \\ \\ \relax
1:12 & \begin{tabularx}{0.7\textwidth}{X} 為要使我們，這些首先把希望寄託在基督裡的人，頌讚他的榮耀。 \end{tabularx} \\ \\
[1ex]
\hline
\hline
\end{longtable}
「我們也在祂裡面得了基業,那原是那位隨已意行作萬事的,照著祂旨意所豫定的.叫祂的榮耀,從我們這首先在基督裡有盼望的人,可以得著稱讚.」(弗一11-12)這裡說我們要得神的基業,讓我們看二件事,就是得基業的當然性及基業的特點.
\newline
\newline
(一)得基業的當然性
\newline
\newline
父母的基業當然是由子女承受.我們有神兒子的名份,當然可以受神的產業及得神的恩典.神的產業和人的產業是有分別的.人的基業要父母死後才能得到,但我們的天父不會死,而祂的產業亦是取之不盡的.主耶穌所說「浪子回頭」的比喻,說到浪子分了家產後才離家,但當他回家時,他的父親仍有家產,似乎是用之不盡,但主卻沒有作任何解釋.主又以五餅二魚餵飽五千人,不計婦孺,我想就是多來一萬人仍是沒有問題的,因天父是無限的.
\newline
\newline
我們既是神的後嗣,就可得主的產業,這基業乃為我們保存,將來才會得著.保羅在羅馬書八章32節說:「神既不愛惜自己的兒子為我們眾人捨了,豐不也把萬物和祂一同白白的賜給我們麼?」祂將祂自己所愛的兒子賜給我們,其他的就算不得甚麼.例如一對將結婚的青年,男子買了一串貴重的項鏈送給他女朋友,他斷不會對她說,我只送你項鏈,盒子你不能拿去.
\newline
\newline
(二)基業的特點
\newline
\newline
一切的福氣都是在基督耶穌裡面,凡不在主裡的,都會毀壞,是暫時的,但在主內的卻要永遠長存.有一位朋友,是美國鋼鐵大王的兒子,他有很多產業,但他看到世上一切都屬於暫時的,他住在聖經學院,駕駛的也是普通的汽車,後來愛上一位貧窮的小姐,結婚時帶她回家見他的父母,她才知道他乃富家子弟,雖覺不配,但那男子說世上的財富並不重要,因一切都要過去.
\newline
\newline
我們神的基業包括一切,像保羅對哥林多信徒所說的:一切都是屬你們的.有一個牧師到城裡講道,由一位會友接待他,他帶牧師到各處遊覽,看到一個美麗的花園,宏大的圖書館,他說是他父親的,起初,那牧師以為他神經不正常,事實上,他乃相信保羅的話──一切都是你們的.
\newline
\newline
溫柔的人要承受地土,當我父親由瑞典到中國傳道時,他原來很會做生意.在海上,他想到若要傳福音,就不能再做生意.心中有些怏怏.這是撒旦的攻擊.但他讀到詩篇三十七篇第九節:「惟有等候耶和華的必承受地土」.心中就大得安慰.溫柔的人,未必現在就承受地土,乃要等候耶和華.神曾應許亞伯拉罕,向東向西向南向北所看到的地方,將來都要歸給他,成為他的產業；但他仍住在帳棚度過一生.妻子死了,沒有墳地,墳地主人因為他是客人,不能講價,又把他虧待了.不錯,溫柔人要承受地土,只不過是時候未到.今天我們仍要等候,到時候便要承受地土.
\newline
\newline


\newpage

\section{第41屆~港九培靈會~孫維廉~第七講}
\label{sec:1470}
\textbf{第七個揭幕 ── 祂是阿拉法,俄梅戛}
\newline
\newline
連結: \href{https://www.hkbibleconference.org/session-message/view/1470}{\texttt{https://www.hkbibleconference.org/session-message/view/1470}}
\newline
\newline
\hyperref[sec:1464]{< < < PREV SERMON < < <}
~
\hyperref[sec:toc]{[返目錄]}
~
\hyperref[sec:1325]{> > > NEXT SERMON > > >}
\newline
\newline
啟廿一6「祂又對我說,都成了.我是阿拉法,我是俄梅戛,我是初,我是終.我要將生命泉的水白白賜給那口渴的人喝.」啟廿二:13「我是阿拉法,我是俄梅夏,我是首先的,我是末後的,我是初,我是終.」
\newline
\newline
聖經從開始到末了,都是講及主耶穌,祂是初,也是終.
\newline
\newline
(一)耶穌是救恩的阿拉法,俄梅戛
\newline
\newline
聖經是一本合乎理性的書,從始到末了都有次序的.創世記說到世人犯罪,全世界都表明人是墮落的,本性是壞的.本書第二十一及二十二章,我們可見一個理想的境界,因為耶穌來了,祂說了三個字:「都成了」.祂完全了救贖,世界便恢復原狀.起初,神咒詛地,地便生出荊棘,後來神咒詛人,若不守法,就要被咒詛.但這裡只有福氣.得救的人都身穿白衣服,被神潔淨了.從前在樂園中有生命樹,但主現在說祂是生命,凡到祂那裡的,都有永生.祂也有生命的水泉,因為地被咒詛,故有沙漠,這是罪的結果.人雖曾被趕出樂園,但現在可進入神的國,新耶路撒冷的門為我們而開.由此可見耶穌是救恩的阿拉法和俄梅戛.
\newline
\newline
(二)主是歷史的阿拉法,俄梅戛
\newline
\newline
啟二十二:16「我耶穌差遣我的使者,為眾教會將這些事向你們證明.我是大衛的根,又是他的後裔.我是明亮的晨星.」
\newline
\newline
這節聖經是關係歷史.美國有二位歷史家,最近寫完幾本世界歷史,寫了四十年.最後一本有結論,但我並不滿意.它的結論太悲觀,因他們從歷史覺得死了便完了,終於世界會變成一個死的世界.我有這本聖經,知道神在歷史中有一個結論,祂是有計劃的.
\newline
\newline
祂是大衛的根,說明祂是大衛的神.祂在世界上選了一小民族,便是以色列民；又尋找大衛,來完成祂的美意.詩篇二十二篇,大衛已預言到我們的救主要再來,描寫到主死在十架上.神和大衛立約,將來有一個後裔永遠做王.耶穌是大衛的後裔,祂亦是大衛的主,有一天,要坐在神的右邊,這要應驗先知的話.
\newline
\newline
主若不來,歷史便沒有結論,世界屬主,一切的盼望亦屬於祂,將來祂要來.主是晨星,晨星令人快樂,而且非常明亮.一切的權柄亦在耶穌的手中.世界上有很多人都想起來掌握權柄.拿破崙曾試過,結果失敗,在海島中,他說:他所有只是暫時的,因為是靠自己的力量,但主耶穌所有的,卻是永遠,因為是靠著愛.
\newline
\newline
(三)耶穌是聖經的阿拉法,俄梅戛
\newline
\newline
啟十九:10「......因為預言中的靈意,乃是為耶穌作見證.」啟廿二:6-10「天使又對我說,這些話是真實可信的,主就是眾先知被感之靈的神,差遣祂的使者,將那必要快成的事指示祂僕人.看哪,我必快來,凡遵守這書上預言的有福了.這些事是我約翰所聽見所看見的.我既聽見看見了,就在指示我的天使腳前俯伏要拜他.他對我說,千萬不可.我與你,和你的弟兄眾先知,並那些守這書上言語的人,同是作僕人的.你要敬拜神.他又對我說,不可封了這書上的預言,因為日期近了.」
\newline
\newline
看了這段經文,可知耶穌是聖經的鑰匙.我在日本曾做教會的司庫,只有一根鎖匙.有一天坐人力車時失去了.後來出高價尋找,那車夫便拿來還我.因為那條鎖匙給了他也沒有用,但於我卻是要緊的.
\newline
\newline
沒有鑰匙,看聖經便沒有意思.若認識主,聖經就給你打開了.聖經是為耶穌作見證,祂是聖經的阿拉法及俄梅戛,故我們必須看完全本聖經才有意義.以前馬克吐溫有一本書寫一半便沒有繼續,令人看了,發生很多問題.世界上一切的故事都沒有完結,但聖經就不同,因主是初,亦是終.將來我們到那裡理想的新世界中,耶和華便住在人間,那裡沒有流淚和死亡.我們現在便要等待和羨慕那一天.有一天,我們便要看見一個非常完滿的結束.
\newline
\newline
在第二十及二十一節,全都講及主耶穌,有三件事:(1)最後的應許,(2)最後的禱告,(3)最後的祝福.
\newline
\newline
很多人都喜歡爭說最後的話,但這裡是耶穌說最後的話.主在耶路撒冷的假日中,利未人拿水,記念摩西擊打磐石,向神獻祭.主耶穌站起來說:「人若渴了,可到這裡來喝,信我的人,就如聖經上所說,從他腹中要流出活水的江河來」.天地都可廢去,但神的話不能廢去,一直存到永遠.祂既說再來,就必再來.我們若等到祂來才預備,便來不及,因為祂來時有如閃電.我們必須預備好,盼望那榮耀的日子來臨.
\newline
\newline


\newpage

\section{第41屆~港九培靈會~孫維廉~第八講}
\label{sec:1325}
\textbf{第七福氣 ── 聖靈為印記}
\newline
\newline
連結: \href{https://www.hkbibleconference.org/session-message/view/1325}{\texttt{https://www.hkbibleconference.org/session-message/view/1325}}
\newline
\newline
\hyperref[sec:1470]{< < < PREV SERMON < < <}
~
\hyperref[sec:toc]{[返目錄]}
~
\hyperref[sec:1506]{> > > NEXT SERMON > > >}
\newline
\newline
以弗所書 1:13-1:14
\newline
\begin{longtable}{cl}
\hline
\hline
章節 & 經文 (和合本修訂版)\\
\hline
1:13 & \begin{tabularx}{0.7\textwidth}{X} 在基督裡你們聽見真理的道，就是那使你們得救的福音，你們也信了他，就受了所應許的聖靈為印記。 \end{tabularx} \\ \\ \relax
1:14 & \begin{tabularx}{0.7\textwidth}{X} 這聖靈是我們得基業的憑據，直等到神的子民得救贖，使他的榮耀得到稱讚。 \end{tabularx} \\ \\
[1ex]
\hline
\hline
\end{longtable}
聖靈為印記的福氣,是非常的寶貝.在這段聖經內,我們可見神的合作,神乃三位一體的神.我雖傳道多年,仍不能詳細解釋；但我卻深信不疑.科學家說立體有四面,第四面就是時間,雖然他們用數學或其他方法來證明,但我仍不明白,不過我卻相信.神的智慧比我們的智慧高,我們雖然不明白,但我們仍可看見三位合一的神是合一的.
\newline
\newline
我們一切福氣的範圍,包括我們剛說過六個福氣,都是從天父來的.聖子便是範圍,一切的福都是在愛子裡來的；故我們只誇耶穌,因恩賜是祂帶來的；離開了祂就不能得甚麼.
\newline
\newline
第七個福才提到聖靈,祂是我們得福的保證.有了產業,就必須有保證.本章十三節「你們既聽見真理的道,就是那叫你們得救的福音,也信了基督,既然信祂,就受了所應許的聖靈為印記.」這裡三個最重要的字便是「聽」,「信」,「受」.有了這三樣便得救.這是十分簡單.有人以為太簡單,但神要每人都得救,故把救恩放在最低的水平上,連孩子聽了都有機會接受.若你不得救就要怪自己了.救恩簡單,為要叫我們得福.有人聽福音,但不信,以為不合現實.我曾舉行佈道會,有青年說不夠現實.後來問他們的牧師才知他們需要救濟品.於是,第二天,我便用很厲害的話責備他們.因為他們只看見現在所需要的,而看不到天上的永久.那晚,那青年便信主.有人只聽而不信,以為那只是一些超然的事.
\newline
\newline
聖靈在兩方面都有見證.
\newline
\newline
(一)從神的觀點──是印記
\newline
\newline
每一個屬主的人都有印記,故神知道誰是屬乎祂的.祂以聖靈為印記,故不會失掉或毀滅.當人信主時,聖靈便入人心中.得到聖靈便有盼望,思想行為便改變,與前不同.
\newline
\newline
主不但是我們的產業,我們也是祂的產業.可能有人不會明白.當一對男女結婚時,誰屬誰呢?我們和神的關係亦是如此,我們能享受祂的恩典,祂叫我們有價值,主把最無用的變為最有用的.像保羅說,他本是一個罪魁,但主改變他,這是一個奧秘,以致他能在獄中寫信與聖徒們彼此勉勵.
\newline
\newline
(二)從我們的觀點──憑據
\newline
\newline
「這聖靈,是我們得基業的憑據,直等到神子民被贖,使祂的榮耀得著稱讚.」(弗一14)
\newline
\newline
溫柔的人要承受地土,我們將來往太空旅行亦不需要買票,到時候,有聖靈為證為質.我們要有聖靈為證,才能承受神的國為產業.時候快到了,我們要忍耐等候,將來我們要一同歡喜快樂,同受產業,但今天就必須聽,信,受.
\newline
\newline


\newpage

\section{第41屆~港九培靈會~林道亮~第一講}
\label{sec:1506}
\textbf{得永遠的生命}
\newline
\newline
連結: \href{https://www.hkbibleconference.org/session-message/view/1506}{\texttt{https://www.hkbibleconference.org/session-message/view/1506}}
\newline
\newline
\hyperref[sec:1325]{< < < PREV SERMON < < <}
~
\hyperref[sec:toc]{[返目錄]}
~
\hyperref[sec:1505]{> > > NEXT SERMON > > >}
\newline
\newline
三十一年前的今天,我第一次在這裡領培靈會.在這三十一年當中,無論是個人,國家,社會都有了一個極大的變動.三十一年前,安樂園的創辦人張吉盛先生還健在,承他老人家熱誠招待,我永久不會忘記.那時楊少泉牙醫師替我翻譯,這位老人家愛主熱忱,而今,這些主內的長輩都回到主那裡享受他們的安息.我還記得三十一年前的今天晚上,這裡的正堂和閣樓擠滿了人,連講台上,台階上也滿坐了人,較之今晚,令我有無限感慨,這是一幅圖畫,告訴我們說:教會在這末世是越過越不成樣子了!
\newline
\newline
目前在社會上,雖然科學和知識發達,可是人生卻比以前更虛空.人們對於生命是甚麼,努力地研究,可是對於怎樣生活,卻茫無所知!就是教會也是如此,他們以為舊的信仰不合時代,應另找新的,卻又不知道新的在那裡!事實,真理無新舊,他們放下真理而尋找真理,結果等於緣木求魚,陷在困惑中!這次我們所討論的總題是「信,望,愛」.因為我覺得；今天大多數信徒,他們的信仰有了問題,他們得救後沒有盼望,他們要想彼此相愛,卻愛不起來!
\newline
\newline
不久前,洛杉磯(Los Angeles)有一位牧師打電話問我一個問題,在加州 U. C. L. A. 大學,大學生的座談會裡,有一位學生問他:「信耶穌為要作好人,信其他的宗教也是要作好人,並且其他宗教確有比信耶穌更好的人,信耶穌和其他宗教有甚麼不同呢?」我就對他說:「要解答這個問題,並非三言兩語就能了事,簡單的答覆是,信耶穌不是為要作好人!如果信耶穌單是為作好人,那就不信耶穌也可以!一切宗教都要人作好人,我們中國人有了孔孟的倫理也夠了.我再說:信耶穌能作好人,卻不是為作好人!」
\newline
\newline
一次,一位青年問我一個問題:「我們勸人信耶穌,為了要上天堂,他們覺得這樣,說法似乎太自私!可是不這樣說,又怎樣勸人信耶穌呢?」我說:「請你看聖經,聖經在那一處告訴我們信耶穌為要上天堂呢?」嚴格地說來,信耶穌上天堂這句話,是信徒們受了佛教的影響.聖經只說信耶穌得永生,現在就得永生,不必等到離世的時候!就在這點上,我們看見今天許多基督徒的信仰錯誤了.信是我們的基礎,如果我們信錯了,基礎壞了,還有甚麼前途呢?據我觀察,目前一半以上基督徒,很可能都不是得救重生的.願主憐憫!
\newline
\newline
在真的基督徒中,就是那些確實重生得救的,又因為他們沒有盼望,沒法在生命中長進.一位神學畢業的弟兄,有一次告訴我,他太太向他埋怨:「你已經重生得救了,你辛苦服事主上天堂,不辛苦服事主也一樣上天堂,那你何必要這樣辛苦呢?」我們的舊天性總是喜歡偷懶,若我已經確實得救了,清清楚楚地重生了,「天上的子民證」我已拿到了,我何必還要殷勤服事主,辛辛苦苦地服事主呢?許多基督徒信了耶穌後,毫無盼望；他們得救了,卻沒有盼望在他們前面鼓勵他們前進!事實我們的主給我們極大的盼望,可是我們不知道!弟兄姊妹們:你信了主,你有甚麼盼望嗎?若是你沒有,你的生命當然沒法長進,你作基督徒沒有喜樂,今天教會的烏煙瘴氣,就是因為傳道人和信徒們沒有了盼望.雖然聖經裡充滿了盼望,末世的教會卻聽不見盼望的信息.
\newline
\newline
我們都知道弟兄姊妹們彼此相愛的重要,因為主耶穌離世之前,特給我們這條新的命令,說:「你們若彼此相愛,眾人因此就認出你們是我的門徒了.」說起彼此相愛,目前的教會可說是根本談不上也見不到!許多弟兄姊妹在教會中,彼此從未認識尊姓大名.富有的弟兄姊妹們看見貧苦的弟兄姊妹,視若無睹,連毫毛也不照顧一下.照我所知,香港教會中不少有錢的信徒,甚至放重利給貧苦的弟兄姊妹；他們不知道神給他們有錢,為要他們顧念貧苦的弟兄姊妹哩?他們還一天到晚在那裡說愛主,這是撒謊!神說:「不愛他所見的弟兄,就不能愛沒有看見的神.」願富有的弟兄三思之!
\newline
\newline
今天晚上讓我們注意「信」的問題,因為這是最基礎的問題.如果一個信徒信得不對,他就沒有盼望,也就不會彼此相愛；因為這三件的次序是一定的:先要有信,才有望,然後才能彼此相愛.如果信得不確實,就沒法望得對,也就永遠沒法愛得起.願主憐憫我們,藉著這次培靈會,叫我們信得正,望得對,愛得起.
\newline
\newline
現在讓我們討論為甚麼要信耶穌?是否為上天堂呢?若真的單為上天堂,基督教真的太自私了.目前信徒的自私,有大部份是教會訓練出來的,因為教會裡一直傳說信耶穌能得這樣,能得那樣,不知不覺養信徒有貪得無饜的心,求主憐憫我們!我們為甚麼要信?第一個答覆是:因為我們是人,不是動物,所以要信.
\newline
\newline
人是否動物?
\newline
\newline
人不是動物,因為神造人和造別的動物不同；傳道書三章十一節:「神造萬物,各按其時成為美好,又將永生安置在世人心裡,然而,神從始至終的作為,人不能參透」.神造人與造動物不同.神造動物的時候,說有就有了；但神造人的時候,卻把一樣寶貝放入人裡面.這寶貝是甚麼呢?中文譯作「永生」,但本節小字卻說:「永生原文作永遠.」本節原文說得很清楚,神造人時是把「永遠」放在人裡面,不是「永生」,因為永生不能造成的,乃是要人信神才能得到的.當我們被造時,神在我們裡面放入一件寶貝叫「永遠」.你們記得嗎?神造人時:「神就照著自己的形像造人.」甚麼是神的形像呢?神是永遠的,神把祂自己的「永遠」放在我們裡面,叫我們像祂,這節聖經也告訴我們,為甚麼神要把這個「永遠」放在我們裡面,叫我們像祂?我不大懂得譯中文聖經的人為甚麼把下半句的開始譯為「然而」,這兩字在原文是三合一的一個字,應譯作「除此以外」,或「沒有那個......就」.神造人時把「永遠」放在人裡面,沒有那個,神從始至終的作為,人就不能參透.換句話說:人應該明白神,因為人有「永遠」.神將「永遠」放在我們裡面,目的就是要我們參透神的作為.這就是人與動物的不同點:動物是沒有宗教性的,或說是沒有靈覺的,人是有的.你有沒有看見猴子媽媽帶著小猴子到猴子禮拜堂作禮拜呢?我從沒有見過一隻老虎牧師講道.動物沒有認識神的本能,人卻有這本能.據調查:世上沒有一個民族是沒有宗教思想的.我仍記得李文斯頓(Livingstone)和他岳父的辯論.李氏一生是在非洲傳道,他的岳父是一個對宗教不太認真的人.他們有一次在倫敦開會,他的岳父強調著說,據他所知,世上有些民族是沒有宗教的.他說完以後,李文斯頓立刻站起來說:「對不起,據我所知道,上沒有一種民族是沒有宗教的.」雖然在非洲有些民族,我們的確看不見他們的偶像,這是因為他們迷信,以為自己的神若給別人看見了,那神所賜的福便會跑到別人家裡去的!所以他們便把他們的神藏在帳棚裡的隱密處.總之,人和動物不同,因為人有宗教性,能感覺神的存在也能明白神的事情.
\newline
\newline
在下文第十二至十四節又告訴我們,這「永遠」應使人知道神的二方面:十二,十三節知道神的恩賜──神的愛；十四節知道神的公義──人要敬畏.這樣,人與動物之間有了鴻溝:人有永遠,動物沒有永遠；人有宗教思想,動物沒有宗教思想；人有良知良能,動物沒有良知良能.最可惜的是今天有些傳道人也會說:「人死後一切便完了,正像一盞燈滅了一般.」這分明是他們不知道人和動物的不同.人有永遠,永不會消滅的!
\newline
\newline
人為甚麼像動物?
\newline
\newline
人的墮落,就是人墮落到動物的地位.當神造人後,把人放在神和動物的中間,目的是要人接受祂的生命像祂.撒旦用動物的特性──吃──試探人,叫人墮落到動物的地位.正像彼得後書所說:「但這些人好像沒有靈性,生來就是畜類,以備捉拿宰殺的」(二13),(參詩四十九:10,12,20；九十四:8；猶10)你們記得但以理在巴比倫王伯沙撒元年所見的一個異象嗎?神把動物當作國度的象徵,啟示但以理,因為在神的眼光中,今天的邦國不是別的,不過是野獸而已.那和野獸不同的「永遠」,已在人內裡失去了功用,雖然神的永能和神性是明明可知的,人就不願意明白神,也不願意尋求神,反而偏離正路,一同變為無用；結果神放棄了人,任憑人為所欲為.神放棄了的,魔鬼就接收過去,把人的心眼弄瞎了,叫人無法見靈光,(林後四4)因為神造人時賦給人選擇的自由,人可以選牛扒和豬扒作他的晚餐,神從不干涉.若人定要離棄神,神只有任憑人!因此,人在神前不再是正常的人,而是像動物.我們常說:「人心不古,獸慾橫行」,這真是神眼光中看人的寫照.以往的年日中,人們受了道德倫理的管束,獸性沒法盡量發揮；現在道德倫理已無能為力,人類的獸性便氾濫著社會的每個角落了!
\newline
\newline
謝謝神,祂不讓人類這樣悲哀下去,永遠沉淪在動物的地位,卻要把人和動物的不同點恢復過來,叫人能知道神的事情.因為除非那在人內裡的「永遠」的功能恢復,人就無法成為神所盼望的人.凡願意成為神所盼望的人,或說凡願意恢復那「永遠」的功能的,就必須信耶穌.我相信弟兄姊妹們都會背約翰福音三章十六節:「神愛世人,甚至將祂的獨生子賜給他們,叫一切信祂的,不至滅亡,反得永生.」
\newline
\newline
人怎能成為神所盼望的人?
\newline
\newline
我們都知道人「永遠」,但那「永遠」的正常功用已失去了.自從罪進入了這世界,死就到處橫行；甚麼地方有了罪,甚麼地方就有死.人有罪,人「永遠的功能,就向神死了.聖經所說向神死了.有它自己的意思.我們以為死就是斷氣,心跳停止,脈息全無.但聖經裡許多的「死」它不是斷氣的意思,乃是指和環境脫離了關係.例如我現在和弟兄姊妹有關係,我在這裡講,你們在那裡聽；我講得對,你會接受；我講得不對,你會起反感.萬一我突然間倒下去死了,我與你們之間的關係立刻斷絕了.你們可憐我,我不覺得；你們憎恨我,我也不知曉,我和環境毫無關係了.我們向神死了,就是說我們和神的環境脫離了關係.神的愛,我不欣賞；神的怒,我也不在乎.人成為動物,因為人與神的關係斷絕了.所以除非把人和神的關係恢復,人就無法成為正常的人.要恢復我們與神的關係,我們就要得永生.
\newline
\newline
甚麼是永生呢?剛才我說:「那一段聖經說信耶穌可以上天堂?」或若你們會說:「約翰福音三章十六節豈不是說信耶穌得永生嗎?得永生就是上天堂啊!」不錯,有時在喪事禮拜中,我們會看見有些輓聯或花圈上會寫著「某某人永生」,意思說某某人上了天堂.可是在聖經裡的「永生」卻沒有這意思.「永生」就是永遠的生命.世上有植物的生命,動物的生命,人的生命,但都不是永遠的.山上有些大樹至多也不過二三千年壽命而已.據說烏龜和白鶴能活到五百年,但也不是永遠的.只有神的生命是永遠的；所以「永生」就是神的生命.你信耶穌所得到的是一件禮物,就是神自己的生命.因為惟有神自己的生命,才能把你失去了的「永遠」的功用恢復過來.神把「永遠」放在人們裡面,目的就是要人們接受祂的永生.所以現在除非你接受耶穌的生命,你就無法在神前成為正常的人!
\newline
\newline
弟兄姊妹們,你們信了耶穌,你們是怎樣信的?信耶穌是要有消極和積極兩方面的.耶穌說:「你們當悔改,信福音.」消極方面是要悔改,積極方面是要信福音.我剛才說有許多人信耶穌信錯了,因為有許多人信耶穌,卻從未悔改過.一個沒有悔改的信,就不是信!舊約時代的先知們,新約時代的使徒們,都強調悔改的重要性.耶穌在世傳道時,一開始就說:「當悔改,信福音.」可惜末世的教會忽略了這一點,誰要受洗,來者不拒,不管他有沒有悔改信福音.因此,今天教會中充滿了許多死的教友.教會應是神的教會,但現在充滿了許多像野獸的教友.朋友們,你們要悔改呀!
\newline
\newline
末了,我要作我自己的見證來結束今晚的信息.我從小是在一個聖公會傳道人的家裡長大,我有一個很好的父親,當我六七歲時就教我讀聖經.每天早晚二次,他會讓我讀聖經中我所認識的字,不認識的字,他便替我讀.這樣,到了我十二,三歲的時候,他有時帶我去鄉村佈道.我識得怎樣唱詩,怎樣禱告,也懂得怎樣講道.我還記得在十五六歲的時候,當父親出門旅行,我便代替他領禮拜講道.所以教會裡的弟兄姊妹們都稱讚我,說我是小牧師,真好!我自己也真的以為了不得,那裡知道,那時的我仍舊是一個死的我哩!我和神沒有個人的接觸,我向神仍是脫離了關係.
\newline
\newline
十九歲那年,當我在師範學校肄業的時候,學校附近開了一個奮興會,因為是在晚上舉行,所以我便去聽道.有一晚,講道的是一個美國宣教士,用英文講「罪」的題目,由一位中國弟兄傳譯.很希奇,神藉著他的信息叫我看見自己是一個罪人.我從前以為自己還不錯,是個相當標準的基督徒,能唱詩,能禱告,也能講道給別人聽.那晚,神恩待我,讓我看見自己是個罪人!我才知道,為甚麼許多時候,我想抵抗罪,卻沒有能力.
\newline
\newline
我繼續聽下去,散會以前,他開始呼召,叫一切自覺有罪的都到台前去認罪接受主.有許多人上前去了,但我不敢!雖然我心中想上去,我沒有勇氣!最後,他又說:「你們當中若有不願上來的,請回家去,好好地在神面前算一算罪賬,請耶穌進入你心中.」當晚我就回到宿舍裡,在神面前同神算賬.可是好像同時有一個聲音說:「你在世上這麼多年,那裡認得清楚你的罪啊!」謝謝神,祂堅固了我的心,到底我開始認我的罪了.當我越認罪,我覺得越輕鬆,且有一種亮光進入我心.到了後來我真覺得「罪重擔皆脫落」,「喜樂潮溢我魂,似海濤之滾滾!」從那時起,聖經成為一部新的書；我所信的耶穌不再是理論,卻是真實的生命；我也知道甚麼叫信耶穌.感謝神!
\newline
\newline
祂沒有讓我在十九歲以前去世,否則我會永遠沉淪!朋友:悔改是消極的,接受生命是積極的,你有沒有悔改信福音呢?
\newline
\newline
不要等待,現在正是悅納的時候,現在正是拯救的日子,悔改罷!
\newline
\newline


\newpage

\section{第41屆~港九培靈會~林道亮~第二講}
\label{sec:1505}
\textbf{人像神}
\newline
\newline
連結: \href{https://www.hkbibleconference.org/session-message/view/1505}{\texttt{https://www.hkbibleconference.org/session-message/view/1505}}
\newline
\newline
\hyperref[sec:1506]{< < < PREV SERMON < < <}
~
\hyperref[sec:toc]{[返目錄]}
~
\hyperref[sec:1504]{> > > NEXT SERMON > > >}
\newline
\newline
昨晚我們看見人不是動物,因為神造人不像造動物.神給人有「永遠」,動物沒有「永遠」.可是人在神前像動物,因為那在人內裡「永遠」的功用已經失去了.人與動物不同,就是因為人有「永遠」,「永遠」的功能一失去,人與動物就沒有多大區別了!因此,世上有許多人不認識神的存在,渾渾噩噩過了一世!謝謝神,藉著主耶穌基督的永生,叫人能成為神盼望的人,就是成為神的兒子.我昨晚已沉重地說過,信耶穌不是要做好人,也不是為要上天堂,那末信耶穌是為甚麼呢?信耶穌是為要作神的兒子.如果一個人不肯作神的兒子,僅作個好人,在神眼中他仍是個罪人.因為神的命令是叫人信神的兒子耶穌基督的名.不遵守這命令的,就是在道德倫理上完全,他仍是個罪人；正像耶穌所說:「為罪,是因為他們不信我!」這樣,信耶穌就是作神的「兒子」.我不說作神的「兒女」,因為我們大家在基督裡都只有「神兒子的生命」,沒有「神女兒的生命」.因此某一教會稱姊妹們為女弟兄,是很有意義的.
\newline
\newline
今晚我們要討論人為甚麼要信耶穌的第二點,就是「人像神」.人之所以為人,一方面是因為人不是動物,另一方面也是因為人像神.前年我在美國加州北部的一個浸信會領奮興會,討論了人像神的問題後,有一個年老姊妹對我說,她在教會五六十年了,還不知道耶穌是怎樣救她的,那晚她才知道耶穌的救恩是這樣奇妙在她裡面作工.那位姊妹有了救恩的經驗,卻沒有救恩的知識,所以生命不容易長進.今天有許多弟兄姊妹靈命不能長進,或許也是因為根基不好,不知道所信的是甚麼?當神造人的時候,祂放一件寶貝在人內裡,這寶貝名叫「生命之氣」.這生命之氣是從神出來的,不是神的靈,卻是神使人有生命的一種原素.換句話說,人有了軀體以後,仍是沒有生命,等神把生命之氣放進後,人才起了心理的活動──魂.創世記第一章告訴我們:神造人像祂.甚麼地方像祂呢?當然不是我們的頭或腳像祂,聖經所說的像祂,是指那從地裡面出來的生命原素.那生命原素能起情感,思想和意志的作用.當生命之氣和人的身體一接觸,情感,思想和意志的作用就表現出來.這三種作用,就是心理學家所說的心理作用,也是聖經所說的魂的作用.事實,心理學就是魂學.你們知道心理學這名稱是二個希拉字合成的:第一個是「魂」,第二個是「學」.所以心理學實際就是魂學.心理的分析,雖然新近的心理學家各有主張,但大體上仍不能超越康德 KANT 的三折學──情感,思想,意志.情感,大家都知道是甚麼.例如現在天氣這麼熱,你感覺到辛苦,這就是情感的作用.今天上午,若有人問你今晚去不去培靈會,你的情感會起作用──歡喜或不歡喜培靈會.接著,你的思想起了作用,想想從前藉著培靈會蒙過恩,你應該去；可是同時,又因天氣那麼熱,你的情感不高興去.這樣你就有了戰爭,去不去呢?末後,你讓意志起作用,決定來培靈會,所以你終於來了.總而言之,人在每天生活中,離不了情感,思想和意志的作用.
\newline
\newline
人的靈從神而來
\newline
\newline
我們都不能否認,這些作用在人內裡,已成為無韁之馬,辦事只問感情歡喜不歡喜,不理真理對不對.
\newline
\newline
為甚麼人的心理會成為無韁之馬?原因是生命之氣本身的作用已失去權威.生命之氣除了發生情感,思想和意志的作用外,它本身亦起作用.照原定的次序,它本身的作用應該是管理人們的情感,思想和意志.那麼它本身的作用是甚麼呢?請看約伯記三十二章八節:「但在人裡面有靈,全能者的氣,使人有聰明.」(伯卅二8)甚麼是人裡面的靈呢?那就是全能者的氣.這「氣」和創二7的「氣」同一字.那全能者的氣有甚麼功能呢?它的功能是使人有聰明.這是指那一種聰明呢?物理上的聰明,還是化學上的聰明呢?請看詩篇一一九篇就能知道.本篇詩篇所提「悟性」二字和約伯記的「聰明」同字.例如:「求你賜我悟性,我便遵守你的律法」(三十四節)「求你賜我悟性,可以學習你的命令.」(七十三節)「求你賜我悟性,使我得知你的法度.」(一二五節)「求你賜我悟性,我就活了.」(一四四節)換句話說,人裡面的靈,就是全能者的氣,能使人有聰明和悟性,明白神的律法,命令和法度.這樣,生命之氣的第一個功能,就是令人明白屬神的事情,我們姑且稱它為屬靈的悟性或靈學.第二種功能是記在箴言二十章二十七節.本節的「靈」字和「生命之氣」的「氣」字是同字的.因為翻譯中文的人覺得若譯成「人的氣是耶和華的燈」有些不通,便把「氣」字譯作「靈」.這裡指神給人的生命之氣,就是耶和華的燈,它的功能是鑑察人的心腹.換句話說,生命之氣在我們裡面有監察的功能,或說是良心的作用,就是孟子所說的「是非之心」.在全部聖經裡清清楚楚能找到生命之氣本身的功能,就是這兩件:屬靈的悟性和是非之心.屬靈悟性的功能應是知道神的事情；良心應是依屬靈悟性所知道的來斷定是非.神藉著人的靈,一方面叫人們知道祂的旨意,另一方面叫人們去執行祂的旨意,使人們的情感,思想和意志都受靈的管轄,不致成為無韁之馬.
\newline
\newline
為甚麼人會不像神?
\newline
\newline
可是人的心理為甚麼會像快行的獨峰駝,狂奔亂走呢?聖經告訴我們:這是因為我們的悟性和良心被污穢了.提多書一章十五節說:「在污穢不信的人,甚麼都不潔淨,連心地和天良,也都污穢了.」這裡「心地」二字,在新約中除了幾次外,都是指屬靈的悟性說的.林前二章十六節「主的心」和「基督的心」的「心」字和這裡的「心地」是同字的.在本節清楚告訴我們:不信的人,他們的悟性和天良都污穢了.這樣,我們知道:人的心理成為無韁之馬,是因為那管理心理的屬靈悟性和天良污穢了,失去了功能!但這並不是說人的悟性和天良消滅了.雖然有時我們會說:某人真沒有良心.其實,只有動物才沒有良心的存在.
\newline
\newline
人的悟性和良心雖然失去功用,但它們的本體仍舊存在,這是神的恩典.因為「污穢」二字的本意原指「塗上黑漆」而言.我們若把燈泡塗上了黑漆,燈光就發不出來了；人的悟性和良心被污穢了,也是如此.它們的權威既然無從執行,各地土司就要稱皇稱帝了.
\newline
\newline
我們知道,神的恩典分為二種:普通的思典和拯救的思典.雖然人類敗壞,神的普通恩典仍和人們悟性和良心同在,所以大多數人們仍有「有神」的觀念和「是非」的辨別.我從前在國內作個人佈道的時候,常問一些老先生宇宙中有否神的存在,他們多數會回答說:「神當然是有的,不然的話,人是從那裡來的呢?」就是良心也是如此,它雖然污穢了,但它仍有些作用.社會和國家的法律,就是從人不正常的良心所產生的.正因這個原故,有些法律和神旨是格格不相入的.
\newline
\newline
神怎樣使人再像祂
\newline
\newline
第一講時,我提到神造人時給人「永遠」.這「永遠」好像燈泡,要電進入後才能發光.我們的悟性和良心要有永遠的生命才能起屬靈的作用.那裡有這永生呢?或說誰有這永遠的生命呢?你沒有,我也沒有；莫罕默德沒有,釋迦牟尼也沒有.因為這世界的生命都是暫時的；永遠的生命是在神那裡.可是神那麼大,人們那麼小,祂怎能進入人們裡面呢?神就讓自己道成肉身,成為你我一樣大小的人.雖然如此,祂仍不能進入你我裡面.所以祂必須死.不錯,祂死是為我們的罪,但若單為罪,祂死已夠了,何必要復活呢?祂復活為要藉著復活產生新生命,正如祂所說的:「一粒麥子不落在地裡死了,仍舊是一粒,若是死了,就結出許多粒子來.」(約十二24)這新生命合你的需要,也合我的需要,一進到我們裡面,那悟性和良心的功用就開始恢復.這恢復有如瞎眼得看見,正如使徒行傳二十六章十八節說:「我差你到他們那裡去,要叫他們的眼睛得開.」眼睛一開,就能「從黑暗中歸向光明,從撒旦權下歸向神.」這「眼睛得開」當然是指人們靈裡的眼睛,就是悟性和良心的功能.悟性和良心的功能恢復了,人就能看見神的愛,神的義,像神了.
\newline
\newline
弟兄姊妹們,你們信了耶穌,我請問你:「你有沒有看見耶穌?耶穌有多高?是肥或瘦?」你若說:「我不知道!」那你所信的是迷信了.若你說:「我雖然沒見耶穌,但我知道我所信的是誰!」我相信你所說是真的,因為你在屬靈的悟性中確已看見了耶穌.有些弟兄姊妹,在蒙恩的時候看見了神的愛,會像小孩子一般對神說:「哦,神呀!你真愛我,我不配!」也有些弟兄姊妹,在得救的時候會打自己,說自己是該死的.人都愛自己的,為甚麼會打自己呢?這是因為他們良心的正常功用恢復了,發出神公義的審判來.
\newline
\newline
所以弟兄姊妹們,信穌不是隨便的事,你一定要接受永遠的生命,就是耶穌的生命.若沒有這永遠的生命,你的悟性與良心的功用就無法恢復,你的情感,思想和意志也就永遠沒有管束.你若有了永生,你的情感,思想和意志都會改變,因為你裡面換了新主人,你會面目一新!弟兄姊妹們:請你自問一下,你裡面的靈覺和良心的正常功用有沒有恢復?如果還未恢復,請原諒我說一句嚴重的話；恐怕你還沒有得救哩!耶穌為你死,為你復活,他的目的是要祂的生命給你,要接受這生命!因為經說:
\newline
\newline
「凡接待祂的,就是信祂名的人,祂就賜他們權柄,作神的兒女.」(約一12)
\newline
\newline


\newpage

\section{第41屆~港九培靈會~林道亮~第三講}
\label{sec:1504}
\textbf{信耶穌與彼此相愛}
\newline
\newline
連結: \href{https://www.hkbibleconference.org/session-message/view/1504}{\texttt{https://www.hkbibleconference.org/session-message/view/1504}}
\newline
\newline
\hyperref[sec:1505]{< < < PREV SERMON < < <}
~
\hyperref[sec:toc]{[返目錄]}
~
\hyperref[sec:1503]{> > > NEXT SERMON > > >}
\newline
\newline
有人問:「我們可以不信耶穌嗎?」我的回答是「不可以」.因為聖經這樣說:「神的命令就是叫我信祂兒子耶穌基督的名,並且照祂所賜給我們的命令彼此相愛.」(約壹三章二十三節)這裡告訴我們,神有一個命令,這命令就是叫我們信祂兒子耶穌基督的名.我們都知道,不聽神的話,或說違背神的命令就是罪.所以除非我們故意得罪神,不然的話,我們就不能不信耶穌.
\newline
\newline
今晚我們的討論分為二點:第一,神的命令有了二次的修改,第二,神的命令有二方面.
\newline
\newline
神的命令有了二次的修改:
\newline
\newline
首先我們要看神怎樣修改祂的命令.神修改祂的命令,不是祂命令的原則有了問題需要改正,乃是因環境需要有所更改而已.好像現在在修改中的英文聖經,不是真理有甚麼錯誤需要改正,乃是因為所用的古典文,不合現代的需要.神修改祂的命令,不是原則上有甚麼變動,乃是為了人的需要不同,神降低祂自己來就人.聖經是神的啟示,神的啟示是前進的:原則不改變,細則因時代的不同須要更改.例如:在舊約神的名字是耶和華,在新約我們看見了耶穌；在舊約有安息日,在新約有主日.安息日會和耶和華見證人的錯誤,就在不明瞭神啟示的前進性.不錯,在舊約的時代是守安息日的,但到了新約；主日是我們紀念主的日子.我在這裡要順便提一提,主日不是給我們安息的日子!有許多教會把主日當作安息日,這是一件錯誤的事.你或說:「那未我在主日安息可不可以?」當然可以!但要看你是為誰安息.若說今天是主日,你特別以安息來紀念主,願主賜福給你；你以教主日學,崇拜主來紀念主,主也要賜福給你.若說主日是給你休息的日子,我不相信主會賜福給你!這是主的日子,不是你的日子.你若在這日一切為主,主會賜福給你.我相信許多信徒在主日不很蒙福,是因為不明瞭主日的真諦.我再說:主日是主的日子,不是我們的日子.我們若把這日尊主為聖,主為我們所預備每主日的福會臨到我們.現在讓我們看神怎樣修改祂的命令.
\newline
\newline
你們知道,最初神的誡命是在出埃及記二十章.當我年青的時候,每次聖餐禮拜的時候曾唸十條誡.我雖然唸,可是不明白為甚麼十條誡裡沒有叫我們信耶穌?感謝主的恩典,後來我的聖經知識加增了,我才知道這十條誡不是我們的,乃是為舊約選民的,到了新約,神已把祂的命令修改了.
\newline
\newline
第一次修改命令是在福音書裡.耶穌對少年人說:「你要盡心,盡性,盡意愛主你的神.這是誡命中的第一,且是最大的.其次也相倣,就是要愛人如己.」(太廿二37-39)這樣,在四福音的時候,耶穌已把十誡歸納為二誡:第一要盡心,盡性,盡意愛主你的上帝!第二,就是要愛人如己.命令的原則沒有改變,細則有了更動,叫神的命令比較以前的更能適合任何人的需要.可是在這修改的命令中「信神的兒子主耶穌基督的名」仍舊沒有放入去.這是因為耶穌當時還沒有釘死而復活,贖罪的救恩還沒有成功,所以無從放入.
\newline
\newline
第二次修改是在約翰壹書第三章廿三節.這裡說「神的命令就是叫我們信祂兒子耶穌基督的名,且以他所賜給我們的命令彼此相愛」.在福音書裡,耶穌把十誡歸納為二,本節聖靈把二誡納為一,因為本節「命令」二字在原文是單數.換句話說,到約翰壹書時候,神已把十誡歸納為一,所以若有人問你信耶穌要守多少誡命.你應說只有一個命令,就是要信也要愛.
\newline
\newline
神的命令有二方面
\newline
\newline
現在,我們要看神命令的二方面.我們已說過,這裡只有一個命令,可是有二方面.第一方面是講到「信」的問題我們為甚麼要信?聖經告訴我們,因為這是神的命令:神要我們信祂兒子,耶穌,基督的名!弟兄姊妹們:請問你們所信的耶穌是怎樣的耶穌呢?你說你所信的耶穌是從前來到世界,生活在世上三十多年,為我們死在十字架上,後來復活升天,祂真好,所以我相信祂.或說祂是世上超奇的人,祂是完人之範,所以我相信他.如果我們信耶穌不過如此,照聖經所說,我們仍是滅亡者!不久以前,我在美國和一位神學教談起趕鬼的問題.他說當日耶穌所趕的鬼,無非是些神經錯亂的人.照他說,耶穌也是一位心理治療家,所以他識得怎樣醫治那些人.我就問他曉不曉得鬼附和神經錯亂的區別.他說沒有此種經驗.我對他說:「感謝主,我倒有些小經驗.」我就告訴他被鬼附和神經錯亂的區別,也告訴他可以用甚麼方法來試驗他們.他覺得很希奇,他說:「真有這種事情嗎?」你看他可憐不可憐?他還在美國的神學當教務長呢.他甚至於不信耶穌基督的升天是事實.這樣的人,自己不信耶穌是神,怎能造就學生呢?
\newline
\newline
第一我們所信的耶穌是神的兒子.曾經有人問過這麼一個滑稽的笑話,他說:「你們說耶穌是神的兒子,那麼神的太太是誰呢?」我聽了真的為他啼笑皆非,想不到一個受過教育的人竟會出一個這麼滑稽的問題,我們要知道神的兒子的意思,先要明瞭父子的關係是甚麼.比如說一天你家裡有許多孩子,他們長得一樣高,穿一樣的衣服,甚麼都是一樣的,連說話的聲音也是一樣,但他們當中只有三個是你的兒子,你就不會叫所有的孩子都是你的兒子.同時,除了三個孩子外,所有的孩子也不會叫你作爸爸.他們一樣肥,一樣瘦,一樣黑,一樣白,為甚麼他們不都叫你作父親呢?你也為甚麼不能叫他們都是你的兒子?你會說,這很簡單,因為那三個是你生的,那三個孩子的生命是從你身上分出來的,所以有了父子的關係.甚麼叫作「生」?「生」就是生命的分裂,兒子的生命是從父親身上分裂出來的.所以父子的關係是生命的關係,兒子的生命就是父親自己的生命.當我們明白這關係,就沒有人會問神的太太是誰了.耶穌是神的兒子,就是說耶穌基督就是神自己,分身降世.我們所信的耶穌,就是不折不扣的神.猶太人明白這意思,所以當耶穌說神是祂父時,「猶太人越發想要殺祂,因祂不但犯了安息日,並且稱神為他的父,將自己和神當作平等.」(約五章18節)我們所信的耶穌就是神,祂是不折不扣的神,也是不折不扣的人,萬物是藉著祂有的,萬物是靠著祂存在的,萬物有一天也要歸在祂裡面.(羅十一36)這是我們所信的耶穌,祂是神的兒子.
\newline
\newline
第二,我們所信的耶穌是耶穌.我不知道弟兄姊妹們能否告訴我,耶穌這名稱是甚麼意思?然,有些弟兄姊妹們雖未明白這名稱的意義,卻已得了耶穌這名稱所包含福份.可是要我們的信仰堅固,能幫助別人,我們一定要明白這名稱.耶穌這名稱是二個字合併而成的:「耶」是「耶和華」的縮寫,「穌」是「拯救」的意思.「穌」在聖經中是一個很普遍的字,叫一個不完全的人得完全,叫被綑綁的釋放,都可用這個「穌」字.所以「耶穌」的意思是「耶和華的拯救」,或「耶和華是救主.」
\newline
\newline
當我們知道耶穌是「耶和華的拯救」,舊約對我們的重要性就不言自明了.因為耶穌既是舊約裡的那位耶和華成為我們的拯救,叫我們的不完全,成為完全,受綑綁的得釋放,我們若要透澈認識耶穌,就必要認識在舊約裡的那位大有能力,大有拯救的耶和華.信耶穌不特罪得赦免,也得脫離罪的權勢.因為我們所相信的耶穌是舊約的耶和華,祂不特救以色列的長子免死亡,也救他們脫離法老的奴役.我們的罪藉著祂的死得拯救,我們未來的軟弱,會藉著他的能力得拯救.在美國有好些青年,信了耶穌不肯受浸禮,我問他們的原因,他們說怕不能堅持到底,他們不知道耶穌的拯救是點,也是線,信祂,不但得救時罪得赦免──點,得救後一生也能得祂的拯救──線.為這原故,天使特地把取名叫「耶穌」的原因向約瑟說明了.「你要給祂起名叫耶穌,因祂要將自己的百姓從罪惡裡救出來.」(太一 21)耶穌之所以為「耶」,就是祂要把信祂的人從罪惡的權勢和綑綁之下救出來.我們已信了耶穌,我們相信祂會救我們到底,因為祂是為我們信心創始成終的耶穌.沒有一個人能說祂信了耶穌,一定能站立得穩,一定能打美好的勝仗.可是他能說,靠主的恩典和保守,他一定會過得勝的生活.哈利路亞!祂有能力,因祂的名是耶穌.「穌」字也有「保存」的意思,所以你若怕不能保存,耶穌能保存.耶穌真是我們奇妙的救主.
\newline
\newline
第三,我們所相信的耶穌是基督.基督是譯音,如果根據中文的譯音,那是毫無意義的.今晚我們沒有時間詳細研究「基督」這名銜,我只能簡單地提一提.首先,我們要知道:耶穌與基督,有不同的意義.經上說:「我已經與基督同釘十字架,現在活著的,不再是我,乃是基督在我裡面活著.」(加二20)我們相信耶穌為我們受苦,復活,升天,現在在上帝右邊,聽我們的禱告,保守我們；可是耶穌仍是耶穌,我們仍是我們,神所要的「以馬內利」仍無法實現.我們必須相信耶穌是基督,基督活在我們裡面,天人才能合一,神在永世的計劃才能實現.當我們信耶穌的時候,請耶穌進入我們裡面,這位耶穌是基督,第一,祂是先知,所以祂在我們裡面,藉著聖靈的聲音教訓我們；祂是我們的祭司,藉著聖靈用說不出來的歎息為我們禱告；祂也是我們的君王,在我們裡面治理我們.事實,我們稱為基督徒,就是因為基督住在我們裡面.
\newline
\newline
耶穌是我們的基督,祂現在在我們裡面作王,有一天,他會降臨實現祂的國度,祂要在宇宙中作王.弟兄姊妹們,你們所信的耶穌是否這麼一位基督呢?
\newline
\newline
第二方面是講到「愛」的問題.有人會說,這裡沒有清楚叫我們愛神呀!在這裡,我們可以看見神的智慧.祂確實沒有清楚地叫我們愛祂,因為本節只說「且照祂所賜給我們的命令彼此相愛」.請弟兄姊妹們注意,這裡有兩個字是從前福音書裡的命令所沒有的,那就是「彼此」,指信主的和信主的之間「彼此」要相愛.至於為甚麼聖靈只注重「彼此」,而忽略了愛神呢?請看約翰壹書四章二十節就會知曉:「人若說,我愛神,卻恨他的弟兄,就是說謊話的.不愛他所看見的弟兄,就不能愛沒有看見的神.」又第五章一節:「凡信耶穌是基督的,都是從神而生.凡愛生祂之神的,也必愛從神生的.」
\newline
\newline
我們的神真是奇妙!祂懂得我們的心,也識得我們的詭詐.祂不要我們愛祂,卻要我們信祂；但信祂後,祂要我們愛弟兄姊妹.好像祂說:「你們若愛我,就必愛你們的弟兄姊妹；你們若不愛弟兄姊妹,就是不愛我!」為甚麼我說祂識得我們的詭詐呢?你記得法利賽人嗎?他們常說只要愛神,便用不著愛人.當時有些作兒女的,甚至把供養父母的錢送到聖殿去,就不負責供養.所以耶穌斥責他們是假冒為善的人.(太五5-9)神看見人心的詭詐,就不叫人愛祂,而叫人愛弟兄姊妹們.請我們留意:信神的兒子耶穌基督的名和彼此相愛是一個命令,不是二個.你如果覺得信耶穌基督是重要的,那末愛弟兄姊妹們也是同樣重要.毋怪聖靈藉著約翰說得很清楚:「愛弟兄,就曉得是已經出死入生了.沒有愛心的,仍住在死中!」換句話說,除非沒有信主,若信了主,我們就沒法不愛弟兄姊妹,否則,我們在神前仍是犯法的.
\newline
\newline
再者,弟兄姊妹們,你們知道信和望有一天會停止,愛卻要進入永世裡.我們若不在永世裡,我們就不知道何為相愛!願主恩待我們,叫我們看見這命令的重要性,信祂也彼此相愛.
\newline
\newline


\newpage

\section{第41屆~港九培靈會~林道亮~第四講}
\label{sec:1503}
\textbf{永火的審判}
\newline
\newline
連結: \href{https://www.hkbibleconference.org/session-message/view/1503}{\texttt{https://www.hkbibleconference.org/session-message/view/1503}}
\newline
\newline
\hyperref[sec:1504]{< < < PREV SERMON < < <}
~
\hyperref[sec:toc]{[返目錄]}
~
\hyperref[sec:1502]{> > > NEXT SERMON > > >}
\newline
\newline
照前幾次所討論的,有人以為我不信天堂地獄.我不是不信,但聖經沒有告訴我們信耶穌是為上天堂,聖經只說不信耶穌的會受永遠的審判!這永遠的審判,因中文沒有更合適的字,便姑且把它譯成地獄.這「永遠的審判」是事實.為甚麼我們要信耶穌,因為有永火!
\newline
\newline
第一,我們怎知道有永火?有許多事實不一定要親身經驗的.在學校裡寫文章,特別是現代的格式,引證有權威人的話語是不可少的.當我在大學研究院裡教授的時候,我喜歡叫學生寫論文代替考試,每學期每課我要他們寫五六篇來代替大考.他們交來的論文,除了我要看它的結構內容以外,我也留意它的引證是否有權威.若那引證不夠權威的話,我便叫他們再做,再去埋頭在圖書室裡,直等他們找到有權威的引證為止.所以有權威的言論,可以代替我們親身的經驗.有些人在美國有權威,有些人在中國有權威,或者有些人在香港和九龍有權威.我們知否有一個人,他有天上的權威,也有地上的權威!聖經說:「耶穌進前來,對他們說,天上地下所有的權柄都賜給我了」(太廿八18),這「權柄」二字應譯作「權威」.耶穌說:「天上地下所有的權柄都賜給我了」.那他是最有權威的一位了.他是神,祂能看見我們所看不見的,也能聽到我們所聽不到的,讓我們來看祂對永遠審判的言論怎樣.在四福音中,耶穌至少十次以上說及「地獄」二字,有權威的言論說了一次已夠了,何況耶穌說了十次以上呢?無論他是誰,若說沒有永火的審判,他就是撒旦!因為撒旦就是「反對」的意思,耶穌說有地獄,有永火,他說沒有,豈不是反對神嗎!他反對耶穌,他就是撒旦!
\newline
\newline
弟兄姊妹們,在這末世的時候,教會需要人起來為真理作見證.讓我們相信耶穌所說的一切,因為惟有祂的話是最有權威的.事實,若沒有永火,耶穌又何必為我們受死呢?他死就要救我們脫離永遠的沉淪,免去神永遠的義怒!總之,永遠的審判是事實,沒有人能推翻,也沒有人可取消!因此,我們要傳福音,要苦口婆心去勸人悔改信耶穌,因為永火是事實!聖經告訴我們:神的話安定在天,他說有就有；他命立,就立.天地要廢去,他的話一點一劃也不能更改.他說有地獄,就是有地獄!
\newline
\newline
第二,為甚麼要有永火?聖經告訴我們:神是愛,(約壹四8)便不應該有永火了!神是愛,但我們不要忘記神也是烈火!(來十二29)神是一位啟示者,他要把自己啟示出來.人既是照著祂的形像造的,所以人也會表現自己.我們說話,寫字,衣著等都在那裡表現自己.現在不少青年人,打扮得不三不四,這就告訴我們,他們裡面不三不四.「有諸內形諸外」,本是中國的古訓.一位畫家如果為錢而畫圖,他不是個上流的畫家.上流的畫家不是為錢而畫,而是為表現他的爇術天才；你看了他的畫,就知道他是怎樣的人.中國人很注重寫字,以前家父常對我說:「寫字是第一碗菜!」意思說,別人看見你的字,就知道你的為人了.
\newline
\newline
神是一位啟示者,祂要啟示祂的愛,也要啟示祂的公義；祂要我們認識祂是天父,也要我們認識他是宇宙的管理者,按公義行事.全部聖經告訴我們,祂一方面是慈愛,一方面是公義.這樣的啟示,不但在聖經中,在自然中也看得見:迅雷暴雨表彰神的威嚴,薰風時雨表彰神的慈愛.神造人的目的,是要啟示祂自己給我們,好叫我們表彰祂,榮耀祂.甚麼是榮耀祂呢?是高大的禮拜堂嗎?是多人參加的聚會嗎?不是的!榮耀神的真意是把神的本體顯大,讓別人在我們身上看見神的永能和神性.因為神的本體就是榮耀.(來一3)榮耀神,就是把神表彰出來.神既是慈愛的,祂要人彰顯祂的慈愛；祂也是公義的,祂也要有受造之物彰顯祂的公義.這樣,有人會說,神既需要祂的受造之物彰顯祂的公義,那末有一部份人進入永火是免不了的!神又怎能責怪都些進入永刑的呢?但聖經告訴我們,神不必要人們彰顯祂的公義.換句話說,這永火並非是為人預備的!那是為誰預備的呢?聖經說:「王又要向那左邊的說,你們這被咒詛的人,離開我,進入那為魔鬼和他的使者所預備的永火裡去」.太十五41),那永火是為魔鬼和他的使者所預備的,人可不必參與.可是魔鬼不單喜歡自己墮落,也喜歡人參與他的背逆,他正引誘更多人和他一齊去神為他所預備的地方.耶穌的救恩正好相反,祂來拯救人,為要叫人彰顯神的慈愛,正如聖經所說:「神愛世人,甚至將祂的獨生子賜給他們,叫一切信祂的,不至滅亡,反得永生」.讚美主!但人不是機器,如果人不願意彰顯神的慈愛,祂不勉強人.神給人選擇的自由,祂不會剝奪人的自由,因為祂不是「出乎反乎」的神!
\newline
\newline
第三,永火到底是怎樣的?永火是真的火嗎?我說不是!那你會覺得希奇,為甚麼我一方面說有永火,一方面又說沒有火.請留意:我的確說了有永火,卻沒有火!人的知識是根據比較而來的.屬靈的一切,有物質可以比較的,我們才會明白；沒有物質可以比較的,我們就無從知曉.火是屬於物質世界的,在永世裡沒有物質的一切.不過永遠審判的痛苦,差不多像火燒一樣難受,聖靈只好用火來描寫它.那麼永火到底是甚麼呢?聖經說:「你發怒的時候,要使他們如在炎熱的火爐中.耶和華要在他的震怒中吞滅他們.那火要把他們燒盡了」.(詩廿一9)這裡聖靈把火描寫耶和華的震怒!(參詩八十九46；結三十八；啟十四10).
\newline
\newline
這樣,那永火,並不是我們所能見的火,乃是神永遠的忿怒.形容常不及實際,火字也不能表達神震怒的嚴重性.雖然我們沒有嘗過永火的滋味,永火一定是比火燒更難當；更痛苦!物質的火,一會兒就把物質消滅了；永火是不滅的,是永遠的!(可九44-48)這裡,有人會反對說:「神是慈愛的,人即使進入永火受苦,神最後還會動慈心把人救出來的!」我也願望如此.但這話是騙人的.因為說這話的人不認識神性!請記清楚:神不特是慈愛的,公義的,也是永遠的:祂慈愛和公義的啟示也是永遠的!也不能讓他的慈愛永遠彰顯,卻讓他的公義彰顯一時.祂不能歪曲祂的神性!朋友們:你有選擇的自由,你若要彰顯神的慈愛,主耶穌已為你預備了天天的救恩,他會幫助你成全；如果你要選擇神的公義,神也沒辦法,因為祂不會剝奪他所給你的權利──這權利叫「願意」(啟廿二17；羅二 7-11)你願意甚麼?你願意彰顯神的慈愛,還是神的公義?請記得:神的慈愛是永遠彰顯的,神的公義也是永遠彰顯的,祂由得你選擇!
\newline
\newline
第四,進到永火情形.進到永火是有先後的.啟示錄十九章二十節:「那獸被擒拿,那在獸面前曾行奇事,迷惑受獸印記,和拜獸之人的假先知,也與獸同被擒拿.他們兩個就活活的被扔在燒著硫磺的火湖裡.」這是頭一批.這裡火上加上「硫磺」二字,描寫永火是比尋常的火更可怕.啟示錄二十章十節:「那迷惑他們的魔鬼,被扔在硫磺的火湖裡,就是獸和假先知所在的地方.他們必晝夜受痛苦,直到永永遠遠.」這是第二批,說明永火是永永遠遠的!啟示錄二十章十五節:「若有人名字沒有記在生命冊上,他就被扔在火湖裡.」這是第三批,指一切掛名的基督徒和不信主的人,啟示錄二十一章八節:「惟有膽怯的,不信的,可憎的,殺人的,淫亂的,行邪術的,拜偶像的,和一切說謊話的,他們的分就在燒著硫磺的火湖裡.這是第二次的死.」這是第一個總結,說永火就是第二次的死──永遠沉淪!
\newline
\newline
曾有人問我:「人是那一部份進入永火呢?是魂,是靈,還是全人呢?」我不回答,只把馬太十章二十八節指給他看:「那殺身體不能殺靈魂的,不要怕他們.惟有能把身體和靈魂都滅在地獄裡的,正要怕他.」無論信主的或不信主的都要復活,因為在永世裡的是全人.換句話說,彰顯祂的公義,神要你整個位格,也要你整個位格彰顯祂的慈愛,所以凡今天不留意靈魂問題的人,到那日會「哀哭切齒」!(太十三42-50)他們到那時會後悔,自己痛恨自己,他們會恨自己的愚昧,恨自己的頑梗.他們會說:「以前太太勸我悔改信耶穌,我不接受；先生勸我接受耶穌為救主,我不聽；兒女流淚勸我,我說他們迷信；傳道人勸我,我把他趕出去；有人送聖經給我,我把它丟在地上,我為甚麼會這樣愚昧無知!今天到了這樣地步,我已沒有辦法,只有在永遠的審判裡,永永遠遠彰顯神的公義啊!」他們會哀哭,他們會切齒!
\newline
\newline
朋友們:為了永火,今晚我請求你不要去彰顯神的公義；為了你的靈魂,我請求你來彰顯神的慈愛!耶穌知道那永火是不容易擔當的,所以祂甘心情願為你死.你知道釘十字架是痛苦之事,但我們的主情願在十架上擔當我們的罪,為要我們能免去永遠的審判.朋友:單是這一點,你如果不信主的話,到那天,你就夠痛恨自己,痛責自己了!
\newline
\newline
弟兄姊妹們,你們已信了耶穌,我還要請問你們是否真信耶穌?你信耶穌是否按聖經的手續呢?甚麼手續?耶穌說:「當悔改,信福音!」你有沒有悔改過?有沒有求耶穌的寶血潔淨你,然後求耶穌進入你的心?聖經說:「凡接待祂的,就是信祂名的人,祂就賜他們權柄,作神的兒女.」(約一12)現在就是悅納的時候,現在是拯救的日子,你有沒有接待耶穌在心裡?不要等到明天,因為明天不一定是你的啊!
\newline
\newline


\newpage

\section{第41屆~港九培靈會~林道亮~第五講}
\label{sec:1502}
\textbf{進窄門}
\newline
\newline
連結: \href{https://www.hkbibleconference.org/session-message/view/1502}{\texttt{https://www.hkbibleconference.org/session-message/view/1502}}
\newline
\newline
\hyperref[sec:1503]{< < < PREV SERMON < < <}
~
\hyperref[sec:toc]{[返目錄]}
~
\hyperref[sec:1501]{> > > NEXT SERMON > > >}
\newline
\newline
今晚我要講「望」.
\newline
\newline
今天教會有許多稱為信徒的,說不定有一半以上仍未得救!那得救了的,也很少數過屬靈的生活.原因很多,其中之一,就是沒有前進的推動力推動他們.換句話說,信主之後應「盼望」!不知為甚麼在這末世中,「盼望」的真理竟聽不見.所以弟兄姊妹信主之後,不知道該怎樣前進.
\newline
\newline
美國有些教會有幾千教友,主日崇拜人數不及一半.問他們為甚麼不去崇拜?他們說牧師未講以先,他們已知道他要講甚麼了!牧師說來說去是要人信耶穌,他們已經信了耶穌,再也得不到甚麼,那他們何必去禮拜堂呢?這樣說法,當然是「自我原諒」而已!主日到禮拜堂,第一是敬拜神,第二是得些屬靈的糧食.作基督徒要知所先後,才不致自欺!同時,教會方面也應該注意:傳福音是緊要的,但得救的人也需要乾糧,單是靈奶並不能滿足他們的要求!
\newline
\newline
盼望是重要的
\newline
\newline
學生入學,若不盼望畢業,便不會用功.作父母的,在兒女身上沒有盼望,也不會化心血.盼望是一種推動力.一個作家盼望他的作品流傳,便不分晝夜地努力著作.同樣的,弟兄姊妹們,為甚麼我們愛主愛不起,前進進不成?因為沒有盼望!我再問:各位信主後有甚麼盼望嗎?若說沒有,當然你愛主愛不起來,熱心也熱不起來,因為你沒有推動力!今天有許多信徒,只要知道自己是大有盼望的,便會前進了.
\newline
\newline
盼望是在天上的
\newline
\newline
那末這盼望到底是怎樣的呢?歌羅西書一章四至五節說:「因聽見你們在基督耶穌裡的信心,並向眾聖徒的愛心,是為那給你們存在天上的盼望.這盼望就是你們從前在福音真理的道上所聽見的.」這裡告訴我們,那盼望是信基督耶穌並愛眾聖徒的因,是神為我們保存在天上的,我們還未得到,如何才能得到呢?「只要你們在所信的道上恆心,根基穩固,堅定不移,不至被動失去福音的盼望.」(西一23)這盼望名為福音的盼望,不單人有份,凡受造的也有份(請看本節的小字說).福音不單是赦罪和得到耶穌的生命,且有盼望在它裡面.這盼望是神為人存留在天上的,人能否得到,是要看人努力的情形,是否會在所信的道上恆心,根基穩固,堅定不移!神知道,在世上有許多事物會引動我們,要我們離間福音的盼望!那些被引動離開了的,雖然他們信了福音,得了生命,卻得不到福音的盼望!他們蒙召,卻不得選上!
\newline
\newline
盼望是活潑和可誇的
\newline
\newline
「願頌讚歸與我們主耶穌基督的父上帝,他曾照自己的大憐憫,藉耶穌基督從死裡復活,重生了我們,叫我們有活潑的盼望,可以得著不能朽壞,不能衰殘,為你們存留在天上的基業.」(彼前一3-4)這裡告訴我們,那盼望是活潑的.活潑是盼望的性質,有了這盼望,我們的靈命就會活潑,沒有盼望,我們的靈命就會死氣沉沉!這盼望是:我們的基業.不在地上,乃存留在天上的.希伯來書告訴我們另一種性質:「但基督為兒子,治理神的家.我們若將盼望的膽量和可誇(直譯),堅持到底,便是他的家了.」(來三6)我們沒有可誇的,除了指主誇口外,主為我們所存留的盼望,也是我們的誇口.撒旦有時會引誘我們,叫我們不必為主吃苦,為人應聰明些,多找些錢財,多建些房屋.我們應該說:「感謝主,他為我存留在天上的基業,比世上一切的更美!」這是我們的可誇,也是我們的膽量.我們為了盡忠於主,人若殺害我們,我們既有這盼望,還怕甚麼!死亡不過是門,經過這門就承受了我們的永生.司提反有這盼望,所以他有這勇氣.古今教會多少殉道者,因有這盼望,視死如歸!弟兄姊妹們,今天我們為甚麼沒有勇氣,貪愛世界,就是因為我們沒有他們所有的盼望,為甚麼不少傳道人貪愛錢財,走巴蘭的道路,因為他們的心眼瞎了,看不見這榮耀的盼望.為人總是要有盼望的,侚們既然看不見永世的盼望,只好在這邪惡的世代裡找盼望.願主光照我們看見那榮耀的盼望.
\newline
\newline
盼望所應許的永生
\newline
\newline
提多書一章二節告訴我們:「盼望那無謊言的神,在萬古之先所應許的永生.」不是別的,乃是在萬古之先所應許的永生.這裡或有人會問:「你不是說信耶穌便得永生了嗎?前幾晚你不是強調著說:信耶穌就是讓耶穌進入心裡,有了祂,就有了永生嗎?為甚麼現在又說要盼望得永生?那麼信耶穌到底是已經得到了永生呢?還是仍要盼望呢?」這的確是一個問題!弟兄姊妹們:聖經初步的真理是很簡單的,可是進步的真理,就不如你想像的那麼簡單.有人把聖經看得太簡單,不肯化工夫,不肯學習.他們以為聖經是任何人都可以自通的.我差不多化了五十年的時間在聖經裡,我仍無法測度聖經是多麼長闊高深!進步的真理,不是普通徒能明白的,否則,神為甚麼賜給教會教師的恩賜呢?恩賜意思是屬天的才幹,不是人人都有的,乃是神隨己意所賜的.神在教會設立教師,為要他們把深奧的真理教導祂兒女.別聽人說聖經是可以無師自通的,這是懶惰人的說法.願神在末世多多興起教師的人才,教導祂兒女們明白聖經.比如「永生」二字,有的人說指天堂,有的說指高超的人格.二十年前,我遇見我們教會的會督,他說:「你以為信耶穌就得永生嗎?」我說:「永生並不是道德高超的意思,永生指永遠的生命,就是耶穌的生命!」聖經清楚地說:「有了神的兒子,就有生命.」他卻把它解為高尚的人格!我再說:「永生就是永遠的生!」
\newline
\newline
解釋聖經要根據上下文,因為古代文字常是一字數意,我們不能囫圇吞棗地來解釋.例如大學裡的「親親而仁民」,第一個「親」是動詞,第二個「親」是名詞,我們不能作同樣的解釋.我說這話,希望我們讀經不要大意,免得誤解聖經,把人引到歧途上去.永生原來不止一個意思,它可指我們在信主時已得到的禮物,也可指我們還沒有達到的目的地.
\newline
\newline
耶穌說:「凡為我的名撇下房屋,或是弟兄,姐妹,父親,母親,兒女,田地的,必要得著百倍,並且承受永生.」(太十九29)這裡清楚地告訴我們,有一種永生是需要付代價才能承受的.換句話說,聖經明說有二種永生:一種是信耶穌已得到的,另一種是要努力才能承受的.那麼這種要努力才能得到的到底何所指呢?「耶穌說,我實在告訴你們,你們這些跟從我的人,到復興的時候,人子坐在祂榮耀的寶座上,你們也要坐在十二個寶座上,審判以色列十二個支派.」(太十九28)這樣,根據上下文,這裡的「永生」就是與基督同坐寶座一同作王!在信主時已得到的永生,就是主耶穌的生命,是我們信主的基礎；日後與基督一同作王的永生,是我們的盼望.在這裡,我給你們提醒,凡聖經中說得永生是有條件的,都是與基督同作王的永生；凡是沒有條件的,就是基督的生命.當然,在聖經中也有不少次數提到永生是句括二方面的.
\newline
\newline
這是我們的「前言」,現在讓我們開始討論正題；我們望甚麼?耶穌說:「你們要進窄門.因為引到滅亡,那門是寬的,路是大的,進去的人也多.引到永生,那門是窄的,路是小的,找著的人也少.」(太七13-14)現在我們把這節聖經分作下列四點來討論:第一,誰要進窄問.第二,為甚麼要進窄門.第三,甚麼是小路.第四,甚麼是窄門.
\newline
\newline
第一,進要進窄門
\newline
\newline
我們常見一幅窄門和寬門的圖畫,那裡畫著進窄門的是上天堂,進寬門的是入地獄.這樣的說法,真理是對的,若有人引證這節聖經來解釋這圖畫,那就錯了的,因為這節聖經不是對未信之人說的.我要問你們,本節的「你們」是指誰說的?我已說過,看聖經不可斷章取義,我們要知道「你們」指誰,一定要根據上下文來定.馬太福音五章一節說:「耶穌看見這許多的人,就上了山,既已坐下,門徒到他跟前來.他就開口教訓他們.」這樣,這三章聖經是對門徒說的,不是對未信主之人說的.難怪下文繼續說:「你們是世上的鹽」；「你們是世上的光」；又說:「你們的光也當這樣照在人前,叫他們看見你們的好行為,便將榮耀歸給你們在天上的父.」他們既是世上的光和鹽,天上的神既是他們的父,這「你們」當然指已信主之人,因為未信主的人,天上的神不能作他們的父.再看第七章:「你們祈求,就給你們.」(七節)這是指會禱告的人.「你們在天上的父.」(十一節)這是指神的兒女.「你們要進窄門.」(十三節)也當然是指神的兒女!
\newline
\newline
弟兄姊妹們,你們有沒有注意:你信主以後,前面有一個窄門是你要進去的.你別以為信了耶穌,就一了百了!你的兒子被他的學校接受了,便是開始進了教育的窄門和小路,他要努力前進,才能學成；否則,他不能畢業.當我看見三四歲的小孩時,常對他們說:「要趁機會享受你這金時代!」因為他們一到了五六歲,就要進幼稚園的窄門了,以後一直要走小學,中學,大學,研究院的小路,這窄門真窄,小路真小呀!但不這樣掙扎,你就不能畢業!當你報名入學,你就成了某校的學生,但不能為畢業生.你要畢業就要努力.同樣的,信了耶穌就是進了神的學校,蒙召了,能否選上,要看你在主裡長進的情形而定!弟兄姊妹們,別以為光是傳道人要進神學,事實上,每個基督徒都在那裡進神的學校!雖然我們能否畢業,要看我們的訓練如何而定,但有一件是定規了的；我們的主斷不會給那些未受過訓練的人在他的國度裡有份!
\newline
\newline
第二,為甚麼要進窄門?
\newline
\newline
弟兄姊妹們,在你們前面有窄路小路,你有沒有進去行走呢?你若不進窄門,就一定要進寬門,因為在我們前面沒有第三扇門!窄門或寬門,我們有自由可以選擇；可是我們沒有自由選擇不進任何一門!
\newline
\newline
所以目前只有二種基督徒,我所說的是真基督徒,一種是進窄門,一種是進寬門.為甚麼我們要進窄門?消極方面,因為寬門是引到滅亡的!那你會說:「這樣,作基督徒還會滅亡嗎?」我回答是；作基督徒,信主到底,必不至滅亡!可是聖經中「滅亡」二字是從另一個字根轉出來的,這字根可有幾個意思,例如馬太記載；「門徒看見,就很不喜悅,說:何用這樣枉費呢?」(太廿六8)這裡「枉費」二字與「滅亡」是同一字根的.本節既是對基督徒說的,所以能譯作「枉費」就更合適；寬門大路是引人到枉費裡去的!任何基督徒都有自由進寬門走大路,因為神給我們選擇的自由,可是我們選大路,我們就要枉費神的一切了!年前,有一母親向我哭訴說:「你知道我另有一個兒子,我花了一生的心血在他身上,我受了不知道多少的苦,叫他吃得好,穿得好,希望他成人,誰知他會成為這麼一個敗類,真枉費了我一生的心血!」她說完後,就痛哭!弟兄姊妹們,今日有不少基督徒,令神為他們痛哭!你從前多麼可憐,神拯救了你,恩待了你,不論在屬靈上,物質上,神都豐豐富富地賜給你,今天,你有了一點地位,賺了幾文錢,你就把神等在腦後,忘記了恩待你的神,甚至在人前否認祂,你這麼忘恩負義,枉費了神的心血,你的良心何在!可是這種沒良心的基督徒偏偏多的是,因為主說:「那門是寬的,路是大的,進去的人也多!」
\newline
\newline
願神憐憫我們,不讓我們作忘恩負義的人,就是枉費了神心血的人.
\newline
\newline
「蒙召的人多,選上的人少!」願神保守我們,叫我們作選上的人!
\newline
\newline


\newpage

\section{第41屆~港九培靈會~林道亮~第六講}
\label{sec:1501}
\textbf{進窄門得永生}
\newline
\newline
連結: \href{https://www.hkbibleconference.org/session-message/view/1501}{\texttt{https://www.hkbibleconference.org/session-message/view/1501}}
\newline
\newline
\hyperref[sec:1502]{< < < PREV SERMON < < <}
~
\hyperref[sec:toc]{[返目錄]}
~
\hyperref[sec:1500]{> > > NEXT SERMON > > >}
\newline
\newline
昨晚我們討論了為甚麼要進窄門的消極方面.
\newline
\newline
積極方面,窄門小路是引到永生.這裡的永生不是指一件禮物,乃是指一個目的地.講留意本節「永」字旁邊的三小點,這是說,在原文並無「永」字,乃是譯者加上去,為要令讀者容易明白些.這種加法,有時加得很好,也有時會加壞了,讀經的需要留意.這裡加得很對,因為「生命」和「永生」是通用的.例如:「信子的人有永生.不信子的人得不著永生,上帝的震怒常在他身上.」(約三36)(參約五24；六53－54；約壹五11)第一個「永生」和第二個「永生」是通用的.
\newline
\newline
第三,甚麼是小路?
\newline
\newline
「凡稱呼我主阿,主阿的人,不能都進天國.惟獨遵行我天父旨意的人,纔能進去.當那日必有許多人對我說,主阿,主阿,我們不是奉你的名傳道,奉你的名趕鬼,奉你的名行許多異能麼.我就明明的告訴他們說,我從來不認識你們,你們這些作惡的人,離開我去罷.」(太七21-23)這裡告訴我們；到那日,有許多稱呼「主阿主阿」的人,主會不認識他們!為甚麼呢?因為他們不遵行天父的旨意.國度只有遵行天父旨意的人才能進去,信了耶穌而不遵行天父旨意的人,不能進國度!我希望有一天全世界的傳道人會看見這幾節聖經!他們的確是有恩賜的傳道人,因為連鬼都聽他們的話!他們不單能趕鬼,且能奉主名行異能和奇事!可是他們不能進國度.因為主「從來不認識」他們.難道他們奉主名趕鬼,主還不知道嗎?這裡不認識是不滿意的意思,主不滿意他們在祂面前的生活,所以他們在神眼前是「作惡的人」.「作惡」二字原文是「不守法的」意思.神不讓那些不守法的人在祂公義的國度有份,因為未來的國度是不能讓那些不守法,自私自利的人進去的.他們奉主名傳道,趕鬼,行異能和奇事,當然是得救的,但不能進國度,因為他們不守法.換而言之,天國有天國的憲法,凡不守憲法的,就不能與主同作王!
\newline
\newline
現在我們知道要走小路就要守法的意思.那這天國的憲法在那裡呢?是在馬太福音五,六,七章內.這三章是天國子民的憲法,非天國子民不在其內!
\newline
\newline
這三章可說是聖經中的寶石,非常榮耀的三章,可惜常常被人誤解,以為內容是對任何人說的.昨晚我已提過,這三章的「你們」是指信徒說的,因為「不要與惡人作對.有人打你的右臉,連左臉也轉過來由他打.」(太五39)是不能對任何人說的,也沒有人能辦得到的.如果是對任何人說的,那末當有賊在你客聽裡偷東西的時候,你就要拱手迎他到你臥室再偷些了!當一個國家被敵人侵佔了第一座城,也就該請敵人佔據第二座城了!但是這些經文只對信徒們說的,所以如果有賊要入屋偷　,信徒應設法自衛,因為天國的憲法只對天國的子民有效,對天國以外的人毫無關係!相反的,如果有一位弟兄或姊妹打你的右臉,你若要作天國的好公民,就應該對他說:「弟兄,如果打了我的右臉能令你舒服的話,請你也打我的左臉吧.」若他是個真信主的,你這樣說了,他還會打你嗎?恐怕他要哭了.
\newline
\newline
記得有一次,有一弟兄對我發生誤會,用電話直罵了我半個鐘頭,我一直一聲不響.等他罵完了,我向他說:「還有嗎?」他說:「沒有了!」我說:「那麼再見罷!」過了幾天,那弟兄來見我,他有些不好意思地對我說:「那天你贏了!」我說:「我一直由你罵,怎會我贏了呢?」他說:「就是因為你不出聲,所以你贏了!」
\newline
\newline
這是天國的憲法,要走小路的人,一定要常讀這三章聖經,照上面說的行事為人.弟兄窮乏時向你借錢,你若向弟兄收重利,你便沒有守這憲法,天國也就沒你份了!以前,有時弟兄向我借錢,我從來不肯借的.如果有弟兄姊妹有需要的話,叫他們把錢拿去用就是了,不必還我.總之,這三章是天的憲法,當弟兄與弟兄之間有紛爭的時候,若能想起這三章的話,就會放下自己,彼此相讓,教會裡的問題會由此大減!
\newline
\newline
記得去年我在印尼椰城的時候,為他們查馬太福音第五章,會後,有二位執事一起去向另一位年老的執事認罪,那三位執事很快樂,全體信徒也得快樂.弟兄姊妹們,你們當然可以照自己的意思行事為人,但你若要進國度的話,那就不能由你自主了!你要走小路,守天國的憲法才行!
\newline
\newline
當我們實行馬太福音第五,六,七章的教訓,就會發覺困難重重!可是你們知道嗎?美國那些打疊球的人,他們在社會上是光榮的一群.但他們受訓時所吃的苦,恐非你們所知了!他們睡覺受限制,食物受限制,生活要循規蹈矩,克苦耐勞,為要訓練他們成為好的球員.凡是社會上專門人才,都是受過嚴格訓練的.聖經說:「如果我們和他一同受苦,也必和他一同得榮耀!」(羅八17)
\newline
\newline
弟兄姊妹們,別以為信主後就一定萬事興隆!如果主看得起你,或許你的道路相當崎嶇,叫你受訓練.總之,走小路與主同苦,是引到與主同王的永生!
\newline
\newline
第四,甚麼是窄門?
\newline
\newline
走小路是要進了窄門後才行的.甚麼是窄門呢?前年我在印尼的時候,有一次聚會完了,有一位執事對我說,他從不知道要進窄門,若果知道的話,他早就進去了!他以前以為窄門是專為傳道人的,以後才知道窄門也是為每一個基督徒的.窄門是點,小路是線.那「點」是在甚麼地方呢?羅馬人書一至八章是論生命,九至十一章是插入論神對以色列人的旨意,十二章開始是論生活開始時有一窄門,就是「將身體獻上,當作活祭,是聖潔的,是神所喜悅的.你們要如此事奉是理所當然的.不要效法這個世界.只要心意更新而變化,叫你察驗何為上帝的善良,純全可喜悅的旨意.」
\newline
\newline
這裡叫我們不要效法這世界,要我們裡面的靈覺更新而變化.基督徒的變化不是從外面變進去的,乃是由裡面變出來的.所以要靈覺更新才能變化,亦惟有靈覺新的人,才能察驗神的旨意.因為知道神的旨意和講英文或寫字很相同,都要經過長期練習才可以的.靈覺漸漸更新,神旨也漸漸明白了.總之,明白神旨,不會如隕石從天上跌下來,乃要化時間學習的.所以打開聖經用手指一指,當作神的旨意,是最迷信的.以前我見過一個禮拜堂,分為兩派,互相攻擊.當我去勸他們的時候,其中一派的人對我說,今早他祈禱後,打開聖經,指著列王記下第二十五章,尼布甲尼撒王亡以色列國的記載,所以他知道,是神要他們把對方趕出去.神的聖經被人用作占卜了!神不肯胡亂讓人知道祂的旨意,所以要人先進窄門,不效法這世界,心意能更新而變化,才給人知道祂的旨意.因為一個人若不立定心志,神不放心把祂的旨意啟示他.
\newline
\newline
窄門,就是這裡所說的奉獻.有許多人以為奉獻自己.奉獻並不是說你今天奉獻了,明天就要去讀神學.神的國度裡不要自願兵,他要自己抽調的.奉獻是你的名分,選不選作傳道人是神的事.今天教會裡太多自己奉獻的人,神從未有選召他們,也沒有給他們能力和亮光,所以他們的信息只像一皮袋的水,用一些時就完了.美國教會,有許多牧師只能在一堂作四,五年牧師,因為四,五年後他們的講章用完了,他們便要換禮拜堂.因為在神學裡所學的有限,聖靈又沒有給他們那活水的泉源他們沒法在一處長住下去.許多西方宣教士由於一時情感的激動,以為求中國傳教多屬偉大.等到到了工場,左右碰釘,困難重重後,便灰心喪志,得過且過,良心內疚!所以奉獻了的弟兄姊妹們,千萬不要立時入神學.除非神選召你,千萬別作傳道人!凡是青年人按立作傳道時,我常是為他們流淚,因為在他們前面只有二條路；一是忠心為主,那就有眼淚,痛苦,患難,困難等等著他；一是敷衍過日子,貪愛世界,不愛主,那就有一天會站在主前受審判!
\newline
\newline
若有人說,如果他奉獻了,萬一神拿去他一切所有的,那怎麼好呢?這可說是杞人憂天!記得在三十年前,我在一個奉令會中向很多青年人講奉獻的問題.後來有一個中年人在那裡哭,我問他為甚麼哭?他說,一方面他想奉獻,另一方面他又怕奉獻一切以後神會叫他去拉車,他拉不動怎麼辦?我幾乎笑了出來,我說:「弟兄,如果你的兒子來到你面前,願意從今以後聽你的話,作你所吩咐的事,你會不會對你的兒子說:『把你一切都給我,你去拉黃包車吧!』你想慈愛的天父會這樣兇殘待你嗎?」終於他奉獻了,主也很賜福給他.
\newline
\newline
弟兄姊妹,若你信主多時,還未奉獻過的話,那你仍未進窄門哩!我們奉獻一些金錢,一些時間,不能滿足主的心,因主所要的是我們的心.主要你同掌王權,要你先進窄門.你若未把心奉獻給主,今日是你的機會,讓我們奉獻罷!
\newline
\newline


\newpage

\section{第41屆~港九培靈會~林道亮~第七講}
\label{sec:1500}
\textbf{愛的兩面}
\newline
\newline
連結: \href{https://www.hkbibleconference.org/session-message/view/1500}{\texttt{https://www.hkbibleconference.org/session-message/view/1500}}
\newline
\newline
\hyperref[sec:1501]{< < < PREV SERMON < < <}
~
\hyperref[sec:toc]{[返目錄]}
~
\hyperref[sec:1499]{> > > NEXT SERMON > > >}
\newline
\newline
我們化了四晚討論「信」的問題,因為「信」是基礎,基礎若不好,我們就無法討論下去!
\newline
\newline
昨晚我們討論「望」的問題.我們信耶穌不是盼望得永生,因為我們信耶穌時已得了永生；我們也不是盼望上天堂,因為甚麼時候我們有了神的生命,天堂就是我們的家,好像一個人到外國去一作了公民,那國家一切的權利都是他的了.那未我們望甚麼?望得另一方面的永生,那就是與基督同作王.為要補充一下,今晚想和各位再查一查聖經,免得有人仍未明白與基督作王的重要性.
\newline
\newline
你們知道神造人的目的是甚麼嗎?人作事都有目的,神造人也有祂的目的,並且這目的記載在聖經裡,但不知為甚麼許多人讀經時總是看不見.現在我們要查幾處聖經,希望我們能認識這真理.
\newline
\newline
第一處是創世記一章二十八節,這裡說神造人是要人治理萬物.可惜自從始祖墮落後,神要人同王的計劃就受了阻礙.可是神是全能的,祂所定的計劃或會遲延,但永不取消.人雖然敗壞,神仍有辦法.
\newline
\newline
第二處是在詩篇八篇四至八節記載著:「人算甚麼,你竟顧念他:世人算甚麼,你竟眷顧他.你叫他比天使微小一點,並賜他榮耀尊貴為冠冕.你派他管理你手所造的,使萬物,就是一切的牛羊,田野的獸,空中的鳥,海裡的魚,凡經行海道的,都服在他的腳下.」這裡再次說出了神造人的目的,就是要人管理萬物；但看不出神怎樣恢復祂的目的.要知道神怎樣恢復祂的目的,一定要讀希伯來書二章六至十節.這裡記著說:「但有人在經上某處證明說:『人算甚麼,你竟顧念他,世人算甚麼,你竟眷顧他.你叫他比天使微小一點,賜他榮耀尊貴為冠冕,並將你手所造的都派他管理.叫萬物都服在他腳下.』既叫萬物都服他,就沒剩下一樣不服他的.只是如今我們還不見萬物都服.惟獨那成為比天使小一點的耶穌,因為受的苦,就得了尊貴榮耀為冠冕,叫他因著上帝的恩,為人人嘗了死味.原來那為萬物所屬,為萬物所本的,要領許多的兒子進榮耀裡去,使救他們的元帥,因受苦得以完全,本是合宜的.」這樣,祂是藉耶穌恢復他的目的.耶穌為人人嘗了死味,為要領許多兒子進入榮耀裡,成就神永遠的計劃──要人與祂同王.所以我們被拯救只是開始,不是目的.我們被拯救後,神要訓練我們,使我們有資格與祂同王才是目的.
\newline
\newline
第三處要看神的目的有沒有實現.啟示錄二十二章五節說:「不再有黑夜.他們也不用燈光日光,因為主上帝要光照他們.他們要作王,直到永永遠遠.」 ──目的實現了!這裡告訴我們,信徒將與基督同王.這事實不單是在千禧年內,也是在永世內,因為聖經中重複的「永遠」,都是指永世而言.這樣,全部聖經告訴我們,神造人的目的是要人永遠與祂一同作王.這同王也就是福音的盼望；(西一23)但這盼望會被引動而失去,所以聖經多次要我們「持定」或「堅持」這盼望.(參來三6,14；六18；十23等)持定不是容易的一件事,以結果被召的多,選上的少；得救的多,國度有份的少!一個能傳道,趕鬼,行異能的傳道人,在神的國度內不一定有份!有份在國度的人一定要守法,若不守法,有一天,耶穌會對他們說:「我從來不認識你們,你們這些不守法的人(直譯),離開我去吧!」所以弟兄姊妹們,不是光說:「主阿!主阿!」就可以進國度,我們一定要按照馬太福音第五,六,七章的教訓去行事為人,才可以進去!因為這三章是特為信主的人預備的,不信的人是沒有份的.我希望這幾句補充的話,能加增各位的盼望,去打那美好的勝仗.
\newline
\newline
今天晚上我們要討論「愛」的問題.弟兄姊妹們都知道,在教會中彼此相愛的重要.但愛是在「信,望」之後的,如果信得不好,就望不起來,也無從彼此相愛.信要信得不錯,望要望得確實,才能彼此相愛!我們知道:信,望都會過去的,惟有愛永存,因為將來在天國裡信和望都用不著了,惟有愛充滿著每一角落,如同永充滿著海洋一般!所以我們今日若不學習彼此相愛,日後即使神讓我們進入國度,我們也會覺得格格不相入!屬靈的事,許多人以為是天上掉下來的禮物；其實除了救恩以外,都是要人去實行才能得到的.信要實行,望要實行,愛也要實行才能得到.智慧的神,從未教導祂的兒女作個懶惰人!
\newline
\newline
愛的反面
\newline
\newline
以前我們已討論過,神只給我們一條命令,但這命令有二方面；一是要信神兒子耶穌的名,一是要彼此相愛.(約壹三23)今晚我們要討論愛的二方面── 反面和正面.許多時候我們只懂得愛的正面,不知道它的反面,所以有時把愛的正反面倒置了也不知道.換句話說,我們若不明白愛的反面,我們的愛就很難純粹.甚麼是愛的反面呢?愛的反面也就是罪的正面.
\newline
\newline
甚麼是罪呢?吸煙喝酒嗎?謊話打罵嗎?這些不過是罪的表顯而已,不是罪的本質.我們若不明瞭罪的本質,就很難切實相愛!所以要和大家討論罪的本質.
\newline
\newline
我們知道,每一樣物件都有它的本質.一把劍,它本質是鋼；一張台,它的本質是木.我們基督徒若想愛神愛人,一定要先明瞭愛的反面──罪的本質是甚麼.我們平時要知道某一物件的本質,就用分析的方法去分析它.做學生的都知道,化學老師會把水的本質分成二份氫氣一份氧氣,告訴我們說,水的原素二氫化氧.今晚我們也用分析的方法來研究罪的本質.
\newline
\newline
現在我要問各位,聖經那一處記載著宇宙中第一件罪?這是在以賽亞書十四章十二至十四節.十二節說有一顆明亮之星墜落了,十三,十四節說它墜落的原因.它的原因共有五句話:(一)我要昇到天上；(二)我要高舉我的寶座在神眾星以上；(三)我要坐在聚會的山上,在北方的極處；(四)我要昇到高雲之上；(五)我要與至上者同等.現在讓我們來分析這五句話:它們的動詞和前置詞的短句,除了一次外,都是各不相同的；可是它們的主詞和副動詞卻每一句相同,那就是「我要」!這樣,這些「我要」就是宇宙間第一件罪的本質.撒旦所犯的並不是了不得的大罪,他沒有在天上放火,聖經也沒有說他殺了任何天使,他的罪不過是「我要」.甚麼是「我要」呢?「我要」的意思就是甚麼都以自己為中心.神恩待撒旦,給他尊榮,讓他管理一些星球,他卻不肯順服神,專以滿足自己的慾望為前提,所以神在聖經中宣告說:「你有罪了.」(結二十八26)這是宇宙中第一次的審判,也是撒旦之所以成為撒旦的總因!這樣,「我要」或說「為自己」就是宇宙間第一件罪的本質.我們往往以「自己」為中心,滿不在乎,不知道「自己」就是罪啊!
\newline
\newline
現在讓我們再看地球上的第一件罪,這是在創世記第三章.我不知道各位讀這章有沒有稀奇:亞當不過吃了一隻果子,神就要這樣對待他,連他的子孫也受累,這不是太嚴厲嗎?到底亞當吃的是甚麼果子呢?聖經說:「只是分別善惡樹上的果子,你不可喫,因為你喫的日子必定死.」(創二17)這樣,亞當吃的是分別善惡的果子.可是無論如何只是一個果子,有甚麼大不了呢!如果有母親買了水果回家,叫孩子們別吃,留到飯後才吃,但在她出去後,有一個頑皮孩子偷吃了一個,母親發覺後,最多訓責那孩子一頓,或體罰他幾下,決不至把他趕出家門吧!可是神卻把亞當趕出伊甸,這是為甚麼呢?這是告訴我們,神注重罪的原則過於行為.吃分別善惡樹上的果子,在物質上,不過吃了一個不許吃的果子,但在原則上,卻和宇宙間第一件罪同樣地嚴重.神要亞當以祂的善為善；可是亞當不順服,要以自己看為善的為善,自己看為惡的為惡.換句話說,亞當不願以創造他的神為神,他要以自己為神!所以地球上第一件罪的本質也是「自己」!
\newline
\newline
我們再看地球上第二件罪.創世記第四章告訴我們,一個哥哥為了嫉妒殺死弟弟.甚麼是嫉妒呢?嫉妒是當一人看見別人比自己強,自己不肯上進,反而用手段把別人拖下去,免得別人比自己強.這樣,嫉妒的本質不是「自己」是甚麼呢!
\newline
\newline
我們若細心地分析聖經中每一件罪,便會發現罪的本質是「自己」,這就是神所最憎惡的!以色列人犯了很多罪,使神不得不執行公義,把以色列國和猶太國都滅亡了.他們所犯形形色色的罪,神把它們只歸納為兩點:「因為我的百姓,作了兩件惡事,就是離棄我這活水的泉源,為自己鑿出池子,是破裂不能存水的池子.」(耶二23)這兩點,事實上只是一件罪的兩方面:消極方面,他們離棄了神；積極方面,他們甚麼都為自己.以色列人忘記了領他們出埃及和賜他們迦南地的神,甚麼都憑自己,至終國破民亡!不單以色列人如此,就是全世界也是如此.聖經說:「我們都如羊走迷,各人偏行己路；耶和華使我們眾人的罪孽都歸在祂身上.」(賽五十三6)甚麼是全世界的罪呢?就是各人偏行己路.總之,全部聖經中罪的本質就是「自己」.以色列人的罪如此,全世界的罪也是如此!
\newline
\newline
弟兄姊妹們:我們若要愛神愛人,我們必要先明白愛的反面.無論你怎樣屬靈,若在神前以自己為大前提,你犯了罪!你奉獻金錢,若希望自己的大名登在報上；你唱詩,若希望別人讚美你的聲音；你講道,若希望別人誇獎你的口才,你不是在那裡榮耀神,乃是在那裡犯罪啊!願你立刻停止這種邪念.為甚麼我們要以金錢來買罪呢?或以詩歌和講道來犯罪呢?請記住:一切以自己為中心的事,無論你怎樣原諒自己,神從不以你為無罪!
\newline
\newline
愛的正面
\newline
\newline
昨晚有許多人奉獻,也有許多人表示已經奉獻了,謝謝主!奉獻就是不再為自己活,乃是為替我們死而復活的主活.這本是聖靈果子的根──沒有自己.
\newline
\newline
說起聖靈果子,有些人常傳聖靈果子的道,可是仔細觀察他們的生活,便會發覺他們所結的是罪(為自己)的果子.聖靈的果子就是裡面沒有「自己」的果子.例如聖靈第一個果子是仁愛.甚麼是愛呢?這是常用的一個字,許多人會講,會寫,也有許多基督徒會說:「主啊,我愛你!」但當你仔細觀察,便會發覺,許多人根本不知何為愛!我們幫助人,可能是出於愛,也可能不是出於愛.愛需要動作,可是動作不能代表愛的真實性.這是一個謎,世上有許多書籍提及它,但從未有一書給我們清楚的提示.感謝主,它在聖經裡說得清清楚楚!
\newline
\newline
在聖經裡,我們找到愛兩方面的定義,就是主觀方面和客觀方面,今晚我們只能討論主觀方面.這方面的定義是記載在約翰壹書三章十六節,我希望青年的弟兄姊妹特別留意,因為你們正在談愛的階段.這裡說:「主為我們捨命,我們從此就知道何為愛；我們也當為弟兄捨命.」單是「主為我們捨命,我們從此就知道何為愛.」我們很容易了解,因為主給我們榜樣,叫我們知道；愛就是為人捨命!可是有了下面一句:「我們也當為弟兄捨命,」就有問題了.我們都只有一條命,所以也只能捨一次命；這樣,難道我們至多只能愛弟兄一次嗎?在這裡,我們覺得需要原文的幫忙了.在原文「捨命」的「捨」,原意是「放下」；「命」,原是「魂」.所以在原文本節是:「主為我們放下魂,我們從此就知道何為愛,我們也當為弟兄放下自己.」換句話說:甚麼是愛的正面呢?愛的正面就是放下自己.
\newline
\newline
有時,有些青年人對我說:「他(或她)很愛我,但他(或她)不肯同我去禮拜堂!」我就說:「若是他真愛你,別說上禮拜堂,就是下監獄,他也甘心同你去啊!」若有人說愛你,可是一點也不肯放下自己,他是在那裡撒謊,別信他!曾有一對年輕夫婦,上街買了一個美麗的花瓶,回家後,因為太太要把花瓶放在鋼琴上,丈夫要把它放在寫字台上,結果就爭吵起來.可是西方人,就是在爭吵時,也是彼此稱呼「親愛的」.他們彼此不斷地說:「親愛的!」同時他們一點也不肯放下自己!這是愛嗎?還是欺騙呢?
\newline
\newline
我們愛神,就當放下自己,否則便是欺騙神了!耶穌在世時能愛神,沒有罪,祂是有秘訣的,那秘訣就是沒有自己.祂說:「我實實在在的告訴你們,子憑著自己不能作甚麼,惟有看見父所作的,子纔能作；父所作的事,子也照樣作.」(約五19)祂又說:「我憑著自己不能作甚麼；我怎麼聽見,就怎麼審判；我的審判也是公平的；因為我不求自己的意思,只求那差我來者的意思.」(約五30)在祂告訴眾人祂是生命糧之後,祂就說:「因為我從天上降下來,不是要按自己的意思行,乃是要按那差我來者的意思行.」(約六36)當祂去客西馬尼祈禱前,祂說:「但要叫世人知道我愛父,並且父怎樣吩咐我,我就怎樣行.」(約十四 31)耶穌在世,祂的行事為人,都是以愛神為中心,所以他能說:「你們中間,誰能指證我有罪呢?」(約八36)
\newline
\newline
弟兄姊妹們:我們要愛神,一定先要知道神所憎惡的是甚麼?神所憎惡的既是我們的自己,我們就當放下自己.從今天起,一切都為神而活,也照神的意思而作.神若把你放在某公司裡,你要在那裡向神盡忠.就是作服事人的,也要為主的緣故服事人.作父母的,應為神的緣故養育兒女,藉著祈禱把他們交託給神,神會負責；若專為自己而養育兒女,神不會賜福.作學生的,讀書不應為自己的前途而用功,應是為榮耀主而努力,神必賜福.
\newline
\newline
弟兄們:你們早晨起來梳頭,帶領呔,有否想到是為誰而作?姊妹們:你們有時塗口紅,有否想過是為誰而塗.以前在廣西有一位信主的姊妹,說因自己的臉特別黃,為了不使別人以為信主的人餓得臉也黃了,所以塗上口紅,我說:「姊妹:塗你的口紅吧,願主賜福你!」總之,我們信主的,已不是自己的人,乃是主用寶血買去了的,所以凡事不應為自己活.為自己的,便沒有愛,不能進神的國度,因神的國度裡是沒有自己的!願神憐憫我們,叫我們憎惡自己,愛神更深!
\newline
\newline


\newpage

\section{第41屆~港九培靈會~林道亮~第八講}
\label{sec:1499}
\textbf{愛心的實行 ── 領人歸主}
\newline
\newline
連結: \href{https://www.hkbibleconference.org/session-message/view/1499}{\texttt{https://www.hkbibleconference.org/session-message/view/1499}}
\newline
\newline
\hyperref[sec:1500]{< < < PREV SERMON < < <}
~
\hyperref[sec:toc]{[返目錄]}
~
\hyperref[sec:1498]{> > > NEXT SERMON > > >}
\newline
\newline
前幾晚我們討論原則上的問題,今晚我們要討論技術的問題.昨晚我們要討論到愛的正反面；今晚我們要討論怎樣去實行愛.
\newline
\newline
有不少弟兄姊妹想幫助別人信耶穌,蒙恩典,但不知應怎樣開始,怎樣進行,和怎樣結束.個人見證是一個領人歸主的好方法,你教會裡的牧師可能很忙,沒有時間多做個人佈道的工作,弟兄姊妹應該負起這責任來.我們常聽見別人說:來我們教會聽道吧!或說:信耶穌可以得福氣,蒙恩典!這樣說當然是可以的,我們若懂得個人佈道的方法,工作就不致這麼皮毛了!經上說:「有智慧的必能得人.」(箴十一30)求主憐憫我們,使我們在救人技術上有聰明智慧,能領多人歸主.
\newline
\newline
在每次開始個人佈道的時候,你總要找一些彼此有興趣或相同的事談一談.談了一會,你不妨對他發一個問題,例如:「你知道甚麼是信耶穌嗎?」他或者會說:「信耶穌就是去禮拜堂做禮拜.」你說:「我是問甚麼是信耶穌,不是信耶穌要作甚麼.」這樣你再問一二句,直到他說:「我不曉得.」你就要乘機說:「讓我告訴你好嗎?」如果他說:「好!」你要感謝主,因為祂已為你開了門.如果他佢絕,那你別講下去.千萬不要勉強他,要留下第二次的機會.我所提的不過是例子,方法可以不同,原則都是千篇一律的一把他的思想引到耶穌身上.如果他肯聽你的話,那你就要一步一步的引導他到耶穌那裡.下文幾個步驟都是十分緊要的,如果你能靠著主的恩典,採用以下幾個步驟,就不難領人歸主了.
\newline
\newline
在未討論步驟以前,你的信心需要調整一下.
\newline
\newline
第一,你要相信聖經是神的話.聖經是神的,它的本身有能力,因為經上說:「聖經都是神所默示的,於教訓,督責,使人歸正,教導人學義,都是有益的.」(提後三16)所以個人佈道時,一定要帶聖經.聖經是你的武器,只要你知道怎樣使用,一定是有功效的.千萬別和他論聖經的可靠或不可靠,你只要把聖經拿出來,讀給他聽,也在可能範圍內叫他自己讀,就會有好結果.譬如有賊來你家,你手裡有把寶劍,你會不會對賊辯論你寶劍的來源而不舉起劍來指著他呢?神的話是兩刃的利劍,用不著你辯護,你拿起來使用就是了.
\newline
\newline
第二,你要信聖經本身是有生命的.聖經是神所默示的,「默示」二字原意是「吹過氣的」所以聖經是神吹過氣的.甚麼是吹過氣呢?你們記得在聖經中神曾經向一件受造物吹過氣,產生了奇妙的結果嗎?那就是創世記二章七節的人,神向他吹了氣後,他就成為有靈的活人.神的氣是有生命的,只要我們把神的話播散出去,神的話自己會向下扎根向上展枝!
\newline
\newline
第三,要相信聖靈會與你同工,聖靈是真理的靈,祂要為真理作見證.你若把真理說出去,聖靈要負責.反過來說,若你所說的不合聖經,聖靈便不能與你同工了,因為祂是忌邪的神!記清楚這幾點後,你便可以開始個人佈道工作了.
\newline
\newline
弟兄姊妹們:我們都不願意空手去見天父.我們已討論過,得救是白白的,但進國度,除了神恩外,也要憑成績.我們不單要愛主,也要救人.怎樣救人呢?要有以下幾個步驟?
\newline
\newline
第一,你要用各種方法使他知道自己是個罪人,不然的話,他不會悔改,也不會接受基督救恩.這幾晚你們已很清楚知道,信耶穌是要悔改,信福音,因為這是信的兩方面:消極方面是要悔改,積極方面是要福音.在美國我常對弟兄姊妹們說,因為這是信的兩方面:消極方面是要悔改,積極方面是要福音.在美國我常對弟兄姊妹們說,葛培理博士傳福音有一大缺點,就是不注重悔改.他只是叫人立志,卻不從悔改開始,因此工作很難有永久性的功效.無論是誰傳福音,若不是從悔改開始,請記得:這工作一定只是曇花一現!所以,當你作個人佈道的時候,要記清楚,在你面前的是個罪人,你要用各種說法,使他知道自己是個罪人,肯承自己的罪.第二,你要使他知道,耶穌有妙法拯救他.第三,你要領他親自與主面對面.做到了這三步,通常地來說,這人就會蒙恩得救.
\newline
\newline
第一步,採用羅馬書三章二十三節:「因為世人都犯了罪,虧缺了上帝的榮耀.」這一節聖經最好請你的朋友自己唸,或者你和他一起唸；如果光是你唸給他聽,效果就不會像他念的好.唸完後,你不妨問他:「是誰有了罪?」他會說:「是世人.」你問他:「這些人是否包括我?」他會說:「是的.」那你可以進一步同他討論以下的問題:「為甚麼我們會成為罪人呢?我們有否說過一句謊言,或偷過一次東西呢?因為一句謊言或十句謊言,同樣能叫人成為說謊者,正如偷一次是賊,偷十次亦是賊,毋怪聖經說:『世人都犯了罪!』因為人在世上至少總有一二次的罪過呀!」這時,你最好能作自己的見證,說到自己以前如何得罪神,得罪人,但神不計算你過往的惡,怎樣在你身上施行拯救.你說這一切,無非是要他知道自己是罪人.聖靈到世上來,本來是要人為罪,為義,為審判,自己責備自己,為罪乃是祂工作的第一步,所以你使人知罪,便是和聖靈同工了.
\newline
\newline
第二步,採用羅馬六章二十三節:「因為罪的工價乃是死.惟有神的恩賜,在我們的主基督耶穌裡,乃是永生.」在這裡,你要問他罪的工價是甚麼?他一定會說是死.這時你要留心解釋這個「死」字,就是與環境脫離關係的意思.比如有人的眼睛瞎了,他和日月脫離了關係,看不見陽光和月亮的美,卻不能因此而說沒有陽光和月亮,不過因為他的眼睛瞎了,他和這一切的環境斷絕了關係而已!照樣,罪的工價也是叫人和屬神的環境脫離了關係,許多人不信神的存在,就是因為和神的環境脫離了關係所致.這是罪的結果,叫人向神死,不認識神.神的永能和神性在宇宙中是明明可見的,但人受了罪的蒙蔽,就不認識神,向神死了!雖然如此,但神仍是憐憫的神,人雖然不認識神,祂仍設法要人認識祂,叫人和祂的關係恢復.祂用甚麼方法呢?祂送人一件禮物,(恩賜二字在原文就是禮物)叫人接受祂的禮物,就能脫離罪,認識神.但這禮物是要人去得的,就如耶穌說:「凡勞苦擔重擔的人,可以到我這裡來!」我們去那裡呢?這禮物在那裡呢?它是在我們的主基督耶穌裡,名為永生.只要我們去基督耶穌裡,我們就能得那稱為永生的禮物.
\newline
\newline
第三步,在這裡有一問題,就是耶穌在那裡?永生是神給我們的,卻要我們在耶穌裡得,可耶穌在那裡呢?你可請他看啟示錄三章二十節:「看哪,我站在門外叩門；若有見我聲音就開門的,我要進到他那裡去,我與他,他與我一同坐席.」你可向他解釋如下:「看哪」二字,在聖經中是叫人「留意下文」,在這裡,耶穌叫我們留意,祂不是在天上,也不是在北極,乃是在你的心外叩門!你以往是否感覺到人生的空虛,世事的煩惱,有時甚至會覺得生不如死呢?這些或許是主耶穌在你心外的叩門,祂叫你看見人生的短促和虛幻,使你覺到永世的重要.我們每早起來梳洗,進早餐,上班,勞碌奔波,天天如此,到頭來仍免不了一死,人生實在太無意義啊!耶穌來,是要叫無意義的人生成為有意義的.祂說凡聽見祂聲音就開門的,祂要進到他那裡去與他一同坐席.請注意本節「我要」二個字,這就是說耶穌已經立定主意,如果你肯打開心門,祂一定要進到你裡面來,和你彼此有交通.祂已定規了,你肯不肯立定主意打開你的心門呢?
\newline
\newline
現在到了緊要關頭.你要相信聖經是神的話有生命能力,你要相信聖靈是與你同工,你問他好不好現在就打開他的心門.他或者會說:「我不知怎樣打開心門!」你就說:「好不好我和你一同祈禱?」當你同他祈禱時,你一定要從幫助他認罪開始.你可以這樣祈禱說:「主阿,求你赦免我的罪,饒恕我一切的過犯；以前一切錯誤的地方,願你用寶血潔淨我；願你給我恩典,使我到你面前承認自己的罪.求主耶穌基督現在進到我裡面,賜我生命和能力,與我同在.」你祈禱時,應該你說一句,叫他跟著你說一句,把你自己當作他,使他認罪,邀請耶穌進入他心內.是否這樣就算完工了呢?如果這樣就完工,那你還是未把他領到主耶穌面前,或許會前功盡廢哩!你一定要領他到他自己開口祈禱的地步,這是很重要的!他或會說他不會祈禱,但你一定要鼓勵他開口,多少時候,聖靈就在這一剎那作工,當他一開口時,眼淚就落下來,承認自己是罪人!這是聖靈的工作,多少時候,一個從未祈禱過的人,他會禱告得很屬靈,他知道怎樣接受耶穌進入他內心.在他祈禱接受耶穌以後,你要恭喜他；如果他有聖經的話,請他在聖經的第一頁寫下某月某日接受了耶穌作紀念.
\newline
\newline
弟兄姊妹們:這三步是非常簡單,如果你們肯常常練習的話,不久,你們便會成為打魚的老手!各位:請問你有沒有領過一個人信主呢?假如你從未有過,那將來你見主耶穌的時候,你怎樣回答呢?耶穌會對你說:「我為你到世上受苦,為你流血,為你死在十字架上,為你由死復活；我拯救你,赦免了你一切的罪,為甚麼你不白白的得來白白的捨去呢?」弟兄姊妹們:我們要趁著白日作工啊!黑夜將到,就沒有人能作工了!
\newline
\newline
現在我要補充幾句.如果有人在禱告後仍有疑惑,他說沒有感覺到耶穌進入他裡面,這是因為他用理智來代替了信心的緣故,你要加強他的信心.你可對他說:「如果我在你門外叩門,定要進來,當你開了門,我會不會掉頭走了?如果我真的走了,那我是甚麼人呢?」他若說:「是撒謊者!」那就最好沒有了.那說:「耶穌說祂在你心外叩門,祂要進來,如果你開了門,祂會走了嗎?如果他真的離開你走了,祂是不是一個撒謊者?耶穌會撒謊嗎?你感覺到或未感覺到,沒有多大分別,你一開門,祂已進入了!」有一次,有一位弟兄在加省二埠已信了主,他的太太仍未信,我就去探訪他們,按著以上三步和她談道,以後我們一起祈禱.在她自己祈禱完後,她對我說,她沒有感覺到她丈夫所感覺的經驗.我便把以上的解釋對她說了,她聽後,貶貶眼睛說:「呵!我知道了,我打開心門時,耶穌已進入我心了,覺不覺得是沒有關係的,要緊的,是相信耶穌的應許,祂是信實的耶穌,我為甚麼會疑惑祂呢!」自此以後,她有了改變,夫妻二人熱心愛主,服事主,今天作教會的柱石.
\newline
\newline
弟兄姊妹們:愛是要有動作的!主耶穌拯救我們,就是要我們去愛主也愛人.有時我們單愛人,有時我們單愛主,惟有作個人佈道的工作是愛主,愛人,同時也愛己!如果你能引導一個人信了耶穌,你的信心會加增,你與主的交通會更親密,你的靈命也會長進.教會的發達,不能完全靠牧師傳道努力,乃要靠弟兄姊妹們同心合力,興旺福音!
\newline
\newline
親愛的讀者,當你讀到這裡,我願你立個志願,從今天起,到明年今天止,你願否至少領一人歸主.一年有三百六十五天,這樣漫長的日子,你先祈禱求神為你預備一位對象,然後採用這 ABC,相信你必能領人歸主吧!請你低頭,祈禱,立願.願主賜福給你!
\newline
\newline


\newpage

\section{第41屆~港九培靈會~林道亮~第九講}
\label{sec:1498}
\textbf{認罪與平安}
\newline
\newline
連結: \href{https://www.hkbibleconference.org/session-message/view/1498}{\texttt{https://www.hkbibleconference.org/session-message/view/1498}}
\newline
\newline
\hyperref[sec:1499]{< < < PREV SERMON < < <}
~
\hyperref[sec:toc]{[返目錄]}
~
\hyperref[sec:1497]{> > > NEXT SERMON > > >}
\newline
\newline
這次我向各位傳的信息是很基礎的；雖然是基礎的,但希望對各位是新的信息.這裡有許多基督徒老不長進,原因是在信仰上有了問題.也有一些弟兄姊妹們,雖然有了靈命,但因沒有盼望推動他們前進,他們靈命無法長進,盼望是一種推動力,任何人作一件事情都需要它,沒有它就沒有成功!我們作基督徒的,若沒有這一種推動力,信心就無法增長,愛心也無法發展!
\newline
\newline
說起愛心,我實在覺得港九教會弟兄姊妹的彼此相愛很差,這當然有其他特殊的原因,但信仰的根基不好,一定是其中之一!有人自以為已有了天上公民的證書,相愛也好,不相愛也好,橫豎都可以進天堂,何必花精神去愛弟兄姊妹呢?所以我這次要向各位傳講的仍是基礎的信息.希望靠主的恩典,我們的信仰能夠糾正,推動力能獲得,彼此相愛能實行.耶穌說:「你們若有彼此相愛的心,眾人因此就認出你們是我的門徒了.」一個不相愛的教會是不能代表基督的!
\newline
\newline
今晚,我們要繼續討論為甚麼要信耶穌.簡單地說:因為我們需要屬天的平安.我們常聽人說:「信耶穌得平安!」耶穌復活後的第一句話,就是「願你們平安!」耶穌受死之前,他應許說:「我留下平安給你們!」我不知廣東如何,在廣西的信徒見面時不說「早」或「晚」,卻說:「平安!平安!」這是最好的問候.我們感覺到,今天有許多人沒有平安,生活非常苦悶,在家庭裡孩子不聽話,在學校裡學生不守規則,青年人無禮和剛愎自用,老年人夜郎自大,小夫小妻愛情破裂,老夫老妻不能和睦偕老,甚至教會也充滿著嫉妒紛爭,對真理不是不及,就是太過；此外,國際上沒有永久的友誼,也沒有永久的仇根,只有永久的利益!這一切都令人心惶惶,得不到平安!人們既是這樣,在基督裡有沒有救法呢?感謝神,有!因為耶穌給我們的救恩,不但是一次的救恩,也是天天的救恩.祂不特一次救我們脫離極大的死亡,也天天救我們脫離恐懼.願主憐憫,叫我們今晚能看見自己的真相,藉著主耶穌的寶血得著拯救.現在讓我們分作以下三點來討論:
\newline
\newline
誰得不到平安
\newline
\newline
經上說:「惡人必不得平安!」(賽五十七21)仔細看這節經文,這是神親口說的.祂說:「沒有平安給惡人們!」(直譯)因為平安是祂的,(參約十四 27)祂不能讓惡人浪費祂的平安,所以祂把平安限制起來,不給他們.那末誰是惡人呢!凡得不到平安的都是惡人!這樣,你或會說:「假如我信了耶穌後不得平安,那我也是惡人嗎?」聖經中惡人不一定指外邦人,特別是本節,讀了上文幾節,明知是指以色列人中遠離神的人.可拉,大坍,亞比蘭反對摩西(民十六 26),以色列人中不肯事奉神的(瑪三18),離棄神的(詩一一九53),都看為惡人,所以詩篇九篇十七節說:「惡人就是忘記神的」人!(「外邦人」原意「人民」或「國家」)這樣,在神的眼中誰是惡人呢?就是忘記了神的人.你若是信徒,我請問你有沒有忘記神呢?你若說怎麼知道我有沒有忘記神?我請問你每天有沒有禱告?用飯前有沒有謝謝神?有許多作父母的知道教導孩子,在別人給他們禮物時要說謝謝；可是有多少父母享受了神的恩典,有說過謝謝嗎?如果你沒有忘記神,你每天就不會忘記禱告和讀經,主日不會忘記敬拜神,行事為人也不會目中無神!末世的時候,信徒對主的愛心漸漸冷卻,對靈修也可有可無,對神漸漸忘記,毋怪許多信徒心中沒有屬天的平安；平安是神的,是客觀的,一定要親近神才能支取.所以忘記了神,等於拋棄了平安.當以色列人忘記神的時候,出入的人,不得平安；但當他們盡意尋求祂,祂就賜福他們四境平安(代下十五3,5,15).
\newline
\newline
惡人,就是忘記神的人,必不得平安!為甚麼他們得不到平安呢?因為他們「被自己手所作的纏住了!」(詩九16)惡人得不到平安,一方面是因神不肯把平安給他,另一方面也是因為自作孽,受了綑綁.
\newline
\newline
為甚麼得不到平安
\newline
\newline
本段經文說:「惟獨惡人,好像翻騰的海,不得平靜,其中的水,常湧出污穢和淤泥來.」(賽五十七20)這裡告訴我們:惡人之所以不得平安,是因他們有一個翻騰的海.為甚麼他們的海會翻騰呢?是因海底有污穢和淤泥.那海底的污穢和淤泥是那裡來的呢?又為甚麼它們會湧上來呢?現在各位請在腦海裡畫三個圈:外面是一個大圈,中間的小些,裡面的最小.這三個圈,第一圈代表我們的身體,第二圈代表我們的心理,第三圈代表我們的靈裡.在心理方面,我們再三個圈:情感是最外面的一圈,思想是中間的一圈,意志是最裡面的一圈.在靈裡方面,我把它分為二個圈,外面的是良心,裡面的是靈覺.這就是全人的情形.
\newline
\newline
一個未信耶穌的人,他裡面的靈覺是不合格的(羅一28),所以對於神的事情,他無法知道.雖然如此,靈覺(或說宗教性)仍是靈覺,人們對宗教的事情仍有些少感覺,但因受了罪的敗壞,這些感覺已不再正確了!靈覺的功用既已不再正確,良心的功用也就各是其是,各非其非!雖然如此,因為神普通的恩典,神仍讓人們的良心起作用.為此,社會上仍有法律的存在,人也仍有一部份判斷是非惡的辨別力.良心的外面是意志,這是你的「自我」.當你說:「我......我.......」大都是指著你的意志說的.因它是「心理」政府的總統,沒有它的簽字,法令是不生效力的!
\newline
\newline
現在讓我們再畫一圖,在二條平行的垂直線當中,我們用橫線把它分為三格:第一格稱為意識界,第二格潛意識界,第三格無意識界.意識界存放著有問即答的知識,潛意識界儲在著經過聯想而獲得的知識,無意識界卻埋藏著遺忘殆盡的知識.
\newline
\newline
我們生活在社會中,天天和環境接觸,得到不少印象.這些印象可分為二種:一種是有害的印象,一種是無害的印象.當一種無害的印象臨到我們的時候,因為是無害的,便漸漸進入我們的潛意識界和無意識界,成為我們的常識和經驗.有許多事,當我們一碰時,就知道怎樣應付或辦理,差不多可以不必用思想,這些常識和經驗,是從以往許多印象積聚形成的.但當一種有害印象臨到時,因為是有害的,我們不願意它存在,就會想方法忘記它.在這裡,我們常有一錯覺:以為凡是我忘記了的,便真的把丟在腦後了!不曉得任何印象一進入我們裡面,便再也不能出去了!我們的心思有如照相機裡面的菲林,外界的印象一進入菲林,便再也不能出去了.我們稱為忘記了,無非是我們用意志把它壓下去!但因為這是有害的印象,我們的良心不肯接受它,意志要往下壓,良心要向上反,這麼一來,我們便需要一部份的精力去抑壓它.如果一個人有很多有害印象的話,他就要費很多精力去鎮壓它們!這樣一來,他便會讀書沒有精神,辦事提不起勁,甚至晚上睡不著覺,他那裡還會有平安啊!他的精力都消耗在鎮壓那些叛徒上,那裡還有能力應付日常生活和工作!不特如此,每當意志一疏怠時,那在潛意識界的壞印象就會偷偷地跑到他意識界來；那在無意識界的有害印象,雖然不能跑上來,卻能發出泡泡到他意識界來,叫他困惱!有時我們看見水底有泡泡上來,我們知道池底一定有腐爛的東西,我們的腦海也是如此!許多孩子們,父母無論說甚麼,他們都會加以反坑,這或許是由於他們在過去的時日中受了些有害的印象所發的泡泡而起的!這些印象,不一定是父母所給的,有不少或許是由他們自己幻想出來的.綜之,若我們的腦海底(無意識界)埋藏著大量有害印象,日常便會有大量的泡泡到海面上(意識界)來成為波浪,形成翻騰的海,不得平靜!
\newline
\newline
有一次,美國新澤西有一個人去見醫生,說自己因為一天到晚忙碌,家人勸他飲些酒.他工作得很辛苦,似乎飲酒不過是辛勤的補償而已.到了後來,他的生活沒有別的,就是工作和飲酒.但後來檢查的結果,才知道他工作這樣辛苦,飲酒這樣多,並不是互為因果.這二件都是果,他的因是另有所屬的!那因令他不敢靜下來,因為一靜下來,他便會全身緊張.所以他便日夜藉工作和飲酒,使心思忙碌,頭腦麻木,不讓那因的泡泡湧上來,翻騰出污穢和淤泥來!
\newline
\newline
弟兄姊妹們,如果你內心覺得不不安,一定是有原因的!我們常有一個錯誤思想,以為認罪是非信徒的事,我們信了主的,便不需要認罪了.事實,約翰壹書所寫關於認罪的事,不是為不信的人,卻是為已信的人!我們內心若不平安,必有某種原因.我們裡面若有一種衝動是我們所不能控制的,也必有一原由.我們不可用「忘記它」來自己欺騙自己,一定要把它查出來,交給主.主是奇妙的,只要我們查得出,交得出,祂就會搬走這印象,叫這印象的能力在我們心裡完全消失了.
\newline
\newline
幾年前,我在大學研究院教書的候,有時在上課時順便提及這些事情.有一位加拿大同學聽見後,他來和我談及他的一個朋友,本來她是念醫學的,且是一個很好的基督徒,但以後非常灰心喪志,書也不能唸了,人漸漸瘦下去.她已看過許多醫生,但沒有一個能幫助她,那同學問我肯不肯幫她.我說如果她肯來的話,我願意和她談談,但我不能保證能幫助她.不久,她真的由加拿大來到美國的東南部,我便和她談了大約二小時,我發現她的污穢和淤泥,並她聖經的誤解,就把聖經打開,請她自己唸,和她一同祈禱,感謝主,她完全得了釋放,高興地回加拿大去,一直在醫院裡很努力地服事主.
\newline
\newline
弟兄姊妹們:我舉這些例子,是要各位明白,雖然耶穌肯赦免人的罪,但他要人把罪交給祂.你若不承認你的罪,祂就無從赦免.我們都知道:神造了食物,我們不能說,神已造了食物,那我們不吃也會飽的.我們一定要吃,(當然仍要靠主的恩典),這食物才能成為我們的營養；否則,食物永遠不會供給我們身體的需要.但對屬靈的事,有些信徒反有這種愚昧的思想,以為甚麼耶穌都會擔當,是的,耶穌會擔當；但我們的主作事則的,我們行事為人一定要照祂的規則才行.我們讀聖經,為要知道祂的規則,實行祂的法度.在祂許多規則中,有一條是帶應許的規則:我們若繼續認自己的眾罪,(直譯)主是聖潔的,公義的,必要赦免我們的眾罪,洗淨我們一切的不義.
\newline
\newline
去年我在印尼的時候,有一位牧師要我去看一位姊妹.她以前很熱心愛主,但那時身體很衰弱,連去禮拜堂也沒有力氣.我見到她,問她許多問題,後來才發現她的病根.原來她有一個兒子,起初想奉獻讀神學,但因為她和她丈夫一同反對,那青年就跑到美國去做生意了.現在那青年生意做得不錯,對主卻完全忘記了,這使她非常不平安!她告訴我,她一睡覺的時候,便想到她孩子；一想著孩子,便睡不著覺!這樣,她的病情無非是「內疚無已」!她感覺到,她的孩子離開主完全是她的責任!她很有錢,但錢不能給她平安!我向她建議以下三點:(一)叫她捐助一個獎學金,栽培一個青年服事主,代替她的兒子.(二)叫她寫信向她的兒子認錯,請他原諒.(三)勉力去教會服事主,因為她的軟弱不是真的身體上的軟弱,乃是靈裡的軟弱.她當時完全答應了.這就是說我們要認眾罪──一件一件地承認.
\newline
\newline
弟兄姊妹們:這翻騰的海,不特叫我們頭腦不平靜,也能實際地叫我們身體成殘疾.目前有許多身體的缺陷,可歸因於海底湧上來的泡泡!在歐戰時,有一個英國兵跳進了德國兵的戰壕裡,殺死了三個德國兵,他的上司見到,當然讚賞他的勇敢.誰知在上司走後,他竟發起抖來；過了些時,他成了個啞吧!大家都不知道是為甚麼,連他自己也不明白.後來才查出,原來當他跳入那戰壕的時候,那三個德國兵已經受了傷,他真忍心地把他們打死來邀功!我們都知道,照國際公法,受傷的敵人是不許打死的.他打死了三個受傷的兵,得到了讚賞,但他內心起了衝突.他怕萬一給別人知道,不但沒有獎賞,而且還要受罪.於是詭詐的肉體向他提議:既然除你自己以外沒有人知道,只要你不說別人就無法知道了!他接受了這泡泡,便成了啞吧.
\newline
\newline
弟兄姊妹們:我們多少時候得不著心理和生理上的平安,是因為在我們的腦海裡有著翻騰,在意識界中有著波浪!這些翻騰和波浪是因為在無意識界裡有污穢和淤泥.那些污穢和淤泥不能用忘記來取消它們,卻能被寶血洗淨.這就是為甚麼聖經強調地勸勉我們要認罪.有些人討厭基督教整天叫人認罪,他們不知道,繼續地,細細地認罪,是我們支取天天救恩的必需手續!
\newline
\newline
怎樣能得到平安
\newline
\newline
在本段經文共有二次提及「我要醫治他」(賽五十七18-19).換句話說,不平安就是疾病,等於我們身體發肺癆,或毒瘤一樣嚴重,我們要求主醫治才是.如果人有重病,又不當心的話,他的病便會越來越嚴重．不平安也是如此,我們若不及早解決,不平安會越來越不平安!有人曾問我,基督徒會不會神經病?我說:「如果他常有不平安,又不肯解決,有一天會到神經病的地步!」
\newline
\newline
我們需要主的醫治,他是大醫生.許多醫學家心理學家不能醫治的病,祂都能醫治.祂怎樣醫治呢?祂搬走我們的病根.「赦免」二字在新舊約原是「搬走」的意思.當我們求主赦免,主不單赦免我們,也實際地把我們的罪搬了出去,這就是救恩的奇妙.我們若要祂醫治,我們必要交給祂.就是我們去普通醫生那裡,也要詳細告訴他病狀；否則,病治不好是我們的責任.主願意醫治我們,只要我們告訴祂,祂會負責.我們應該怎樣告訴祂呢?本段經文又說:「因為那至高至上,永遠長存,名為聖者的如此說,我住在至高至聖的所在,也與心靈痛悔謙卑的人同居,要使謙卑人的靈甦醒,也使痛悔人的心甦醒.」這樣,我們應該「痛悔」地告訴祂.
\newline
\newline
「痛悔」原來意思是「粉碎」.神歡喜完整的奉獻,卻悅納粉碎的悔改!我們要主的醫治主要我們把不平安的原由一五一十,詳詳細細地告訴祂.千萬別光說:「主阿,赦免我的罪!赦免我一切的罪!」我們常說,禱告有如鋪路,我們鋪到那裡,主的恩典也進到那裡.認罪亦是如此,我們認到那裡,主的赦免也到那裡.若我們肯詳詳細細地承認自己的罪,(越細越好)那主的恩典也會詳詳細細地臨到我們.如果我們內心不平安,同時又找不出那淤泥來,我們只有求主藉著祂的靈啟示我們,叫我們能找得出來.靈是叫我們「想起」的靈,祂會幫助我們發現那些被抑壓在我們無意識界裡的污穢.祂能引導我們思想,使我們想起.當我們一發現那些有害的印象,便要向主清清楚楚地承認出來,交給祂,讓祂搬走.
\newline
\newline
弟兄姊妹們:多少時候我們不能長進,我們不能完全享受神的平安,就是因為海底裡淤泥和污作祟,我們一定要對付它們.謝謝主,祂的救恩是天天的救恩!
\newline
\newline


\newpage

\section{第41屆~港九培靈會~林道亮~第十講}
\label{sec:1497}
\textbf{更清楚認識罪}
\newline
\newline
連結: \href{https://www.hkbibleconference.org/session-message/view/1497}{\texttt{https://www.hkbibleconference.org/session-message/view/1497}}
\newline
\newline
\hyperref[sec:1498]{< < < PREV SERMON < < <}
~
\hyperref[sec:toc]{[返目錄]}
~
\hyperref[sec:1496]{> > > NEXT SERMON > > >}
\newline
\newline
昨晚和大家討論,為甚麼要信耶穌,因為我們要得屬天的平安.我們常聽人說；平安即是福!真的,無論你多麼富有,沒有平安有甚麼用處呢!如果我們作信徒的沒有平安,我們就不能彰顯主耶穌基督；如果我們內心常時翻騰,我們就不能使神的榮耀在我們身上彰現.昨天晚上我們討論的是比較主觀的,今天晚上我們要討論客觀的.
\newline
\newline
我們為甚麼信耶穌?因為罪的可怕!我說了這句話,恐怕有人會說:「罪有甚麼可怕呢?」如果我們明白聖經,我們會知道,罪是超乎我們想像地可怕!許多時候,就是因為我們不認識罪的可怕,對耶穌的救恩我們不欣賞,對受罪綑綁的靈魂我們不同情,甚至當罪侵入我們自身,家庭和教會的時候,我們也不知道!罪的確可怕,願主叫我們洞悉罪的真相,知道如何依靠主,不得罪主!
\newline
\newline
中文的「罪」字,實際上是一幅可怕的圖畫,可惜許多人不留意!罪字分為兩半,上半是個四字,下半是個非字.上半的四字原來是個網字,如网,下半的非是為非作歹的非,整幅字畫暗示我們,天網恢恢,疏而不漏；為非作歹,難逃天網!凡罪一定有審判,別想逃避!天網雖然疏,但這疏卻不致使你漏網,希特勒,墨索里尼等都是好例子!雖然中國的古聖賢者知道罪的可怕,可是他們的見解到底不是屬靈的看法.神在聖經中以屬靈的眼光來告訴我們罪的真相,說罪是人格化的!有人或以為這不過是神特別要形容罪的可怕吧!我們知道:神絕不會把無的當作有.神的話語是真理,真理一定要事實和理論相符的,不然的話,便成為虛謊!
\newline
\newline
神不特在聖經裡給我們看見罪是人格化的,祂也告訴我們,有一天,罪會實際地罪成肉身!那看不見的罪,竟會成為人形,你看多麼可怕啊!罪成肉身將在那大罪人身上應驗,他是沉淪之子,會在末世顯靈出來(帖後二3).主耶穌曾道成肉身,日後,罪也要成為肉身,行各樣的異能神蹟,和一切虛假的奇事.所以罪在世上的活躍和可怕是超乎普通人所想像的!舊約時候以色列人有逃城,你們都知道這是基督的預表.基督為甚麼要作我們的逃城呢?其中原因之一,是為了罪的可怕!願主憐憫,今天晚上讓我們看罪的真面目,叫我們知道如何隱藏在主裡,不受罪的轄制,也不讓罪侵犯我們!
\newline
\newline
弟兄姊妹們能不能告訴我,聖經中甚麼地方第一次提到「罪」字?請我們留意:凡是在聖經中第一次提到的一件事或一個字,都是十分重要的.比如聖經中第一次提到主日(約二十19-23),你若仔細去讀,你會發現,主為主日所預備的福份有多少!第一次的罪字記載在創世記四章七節.創世記二章七節論到人的被造；三章七節論到人第一次失敗；四章七節是第一次提到「罪」字.這裡說:「你若行得好,豈不蒙悅納；你若行得不好,罪就伏在門前；他必戀慕你,你卻要制伏他.」這樣,當聖經第一次提到「罪」字,神就說罪是人格化的──「他必戀慕你.」「戀慕」二字在舊約共用過三次,指妻子戀丈夫(創三16),或男情人戀慕女情人(歌七10).罪是多麼可怕,她竟有豐富的情感來戀慕人們去犯罪!今晚讓我們研究這第一個「罪」字；罪的性別,性情,動作和含義.
\newline
\newline
罪的性別
\newline
\newline
舊約文字男女性的分別,和中英文男女性的分別不同,凡是像母親的,溫柔的,軟弱的,屬下的和抽象的,舊約文字都用女性來表明.這裡的罪字是個女性的,由此可知罪的詭詐!普通人以為罪是可怕的,醜陋的,殊不知她的外觀是美貌的,柔弱的,令人容易接近的.正如魔鬼,普通人以為是紅髮綠眼,十分難看,殊不知他是宇宙中最美麗的.雖然如此,我猜想他現在一定是走樣了,因為這麼多年來,他內心積存著恨毒,那恨毒能使他的美麗變成醜陋!一個人內心的恨毒,能使他的心理和生理都受影響.如果你們當中有懷恨人的,請記得,受虧損的不是你所恨的人,乃是你自己啊!願主幫助你作聰明的人,把你的恨毒交給主,讓主搬出去,饒恕人是得快樂的秘訣,也是美容的秘方!
\newline
\newline
為甚麼罪要女性呢?與其說是天然,毋寧說是為方便任務起見,因為經上說:「總要趁著還有今日,天天彼此相勸,免得你們中間,有人被罪迷惑,心裡就剛硬了.」罪的任務是迷惑人,那未花容月貌當然比粗手大足方便得多!我作孩子的時候,鄰居告訴我水塘裡有水鬼.我問他怎麼知道?他說:水鬼常在水塘裡變成一朵很可愛的花,但當你伸手去採時,那花會漸漸移遠離開你,讓你失丟平衡,跌入水中,他便用泥塞住你的口鼻,叫你窒息而死!自從我聽了這話,以後每逢看見水塘裡有美麗的花,便嚇得連忙走開.這些可能是迷信,但罪的本領確是如此.她常花技招展地招呼你,引誘你,要你到她那裡去,應許你享受一切.許多無知的人聽信她的花言巧語,就上了當.在箴言書裡有一段話,我們平常以為是描寫淫婦而已,其實是罪迷惑人的寫照:
\newline
\newline
「我曾在我房屋的窗戶內,從我窗檽之間,往外觀看,見愚蒙人內,少年人中,分明有一個無知的少年人,從街上經過,走近淫婦的巷口,直往通到他家的路去,在黃昏,或晚上,或半夜,或黑暗之中.看哪,有一個婦人來迎接他,是妓女的打扮,有詭詐的心思.這婦人喧嚷不守約束,在家裡停不住腳,有時在街市上,有時在寬闊處,或在各巷口蹲伏,拉住那少年人,與他親嘴,臉無羞恥對他說,平安祭在我這裡；今日纔還了我所許的願.因此我出來迎接你,懇切求見你的面,恰巧遇見了你,......淫婦用許多巧言誘他隨從,用諂媚的嘴逼他同行.少年人立刻跟隨他,好像牛往宰殺之地,又像愚昧人帶鎖鍊,去受刑罰；直等箭穿他的肝,如同雀鳥急入網羅,卻不知是自喪己命.眾子阿,現在要聽從我,留心聽我口中的話.你的心,不可偏向淫婦的道；不要入他的迷途.因為被他傷害仆倒的不少；被他殺戮的而且甚多.他的家是在陰間之路,下到死亡之宮.」(七6-27)在這裡,神告訴我們,罪在世上好像一個淫婦,她會用各種花言巧語來引誘我們,特別是無知的人；因為聖經告訴我們,在愚昧人中有一個無知的少年人!我們不願意做個無知的人,可是我們若聽從罪的引誘,聖經就說我們是無知的人!
\newline
\newline
在這裡,我要和你們談一談試探的真相.甚麼是試探呢?試探就是說,用錯誤的手,達到正當的目的.例如在耶穌受試探這事上,耶穌飢餓的時候,找些食物來充飢,本是合法的需要,目的完全正當,並沒有甚麼錯.可是用甚麼手段來找食物,能叫整個行為成為合法或不合法的!魔鬼知道這關頭,他就給主耶穌一個正當的目的,和一個不正當的手,要主耶穌注意目的,(以石頭變成食物充飢)忽略手段.(在神旨外擅自妄動)耶穌知道,神所要的是目的和手段同樣合法,不然的話,就是得罪神,所以耶穌拒絕他的建議.
\newline
\newline
弟兄姊妹們:目前有許多人提倡「為達目的,不擇手段的論調,我們千萬不要受欺騙.「只要目的正當,手段不正當無關要緊」,是從撒旦來的禍音,也就是撒旦試探的內容,我們信主的人,行事為人不但要目的正當,手段也要正當.凡是目的正當,而手段不正當的,主從來不會算他為無罪!今天有許多教會不注意這點,他們滿以為只要目的正當便可以了,手段不正當主會原諒!他們不擇手段去賺錢,還為自己辯護說:「我又不是為自己賺錢,為主,為教會,怕甚麼!」他們不知道,如果手段不正當,他們不是榮耀主,乃是羞辱主!他們也會用各種手段引人到教會裡來,(甚至開跳舞會),給他們受洗,騙他們說,他們已經信了耶穌,不叫他們認罪悔改,接受耶穌為救主,及自圓其說:「我們怎能知道一個人重生得救啊!」總之,甜言蜜語的罪不特在社會中有,在教會中也有,有智慧的人卻從不上她的圈套!
\newline
\newline
罪的性情
\newline
\newline
罪的性別是女的,外表似乎弱不禁風,可是她的內裡,或說本身,卻是強而有力的.在約裡,大多數字是由三個字母作字根的,凡是字根當中那字母重複的,就是表示那字是強而有力的意思,罪字就是這麼強而有力的字.不特如此,「罪就伏在門前」的「伏」字是分詞──有名詞的形式,可作形容詞用,也有動詞的動作.照希伯來文文法,若主位是女性,那形容她的字和動詞也應該是女性,但這裡卻不是如此.如果聖經不是靈寫的話,我一定會說摩西的文法有了錯誤.「罪」字是女性,「伏」字也應該是女性才對,但「伏」字是男性,難道真的摩西用錯了文法嗎?摩西學了埃及人一切的學問,而這裡的文法很簡單,他斷不至用錯文法!這裡是聖靈的特別文法,為要告訴我們,罪的外表是女性,行動是男性!當罪接近我們的時候,她是一個女姓；當我們一旦被她抓住以後,她是一個強而有力的男姓.我常想埃及的獅身女首石像和罪很相似.在埃及的金字塔旁有一座獅身女首石像,從遠處看它,很像一個女人,在近處看便是一隻獅子,罪也是如此!
\newline
\newline
在希伯來文文法,男性是表示勇敢,危險,大力,野蠻等等.所以罪的外觀雖然是溫柔,軟弱,實際卻是野蠻,大力!這種事實,今天在社會上隨處可以找到例子.美國有許多青年人,被引誘吸食毒物以後,可以看見從未見過的事物,感覺到感覺以外的幻境,真的有飄飄欲仙的味兒!於是他們繼續吸食,不久,就無法脫離了!等到無錢購買毒物時,甚麼事也會去作,偷得到就偷,搶得到就搶,名譽不顧,性命也不顧!這就是罪的本色,正如經上記著說:「雖然惡在他口中甘甜,他藏它在舌下,愛戀不捨,含在口中；他的食物卻成為虺蛇的毒!」(伯二十12-14講者重譯).
\newline
\newline
罪的動作
\newline
\newline
罪有外表,內容,也有動作.她的動作也可以「伏」字上見到.這「伏」字是名詞也是動詞.當名詞用時,在阿拉伯文,是獅子的綽號:當動詞用時,它的動作是繼續不斷的.罪是沒有假期的,她會晝夜戀慕我們,她會伏在門口,等在門後,乘機捉住我們!這「伏」字也有它的語態,它的語態是自動的.我們常以為犯罪是我們自己作得主的事,自己要犯就犯,不犯就不犯；但聖經說罪是自動的,不需要你去找她,她會來找你!(民卅二23)罪不但是自動,且會找機會,如經上說:「然而罪趁著機會」,「因為罪趁著機會」.(羅七8-11)我們都知道機會的重要.在希臘有個機會的神話說:有一次,有一人在田野散步,忽然看見迎面來了個怪人.那人除了額前有一朿頭髮外,頭上牛山濯濯,腿上反有翅膀,走起路來如飛.這人就上前問他尊姓大名?他說名叫「機會」；問他為甚麼只在額前有頭髮,他說因為他是機會,如果有人要抓他,就必須在他前面抓他,一旦他過去了,便再也不能抓住他；問他為甚麼腿上有翅膀,他說因為他要走得快,叫人不很容易抓住他!罪很懂得機會的重要,她常伏在門前,在那裡等機會,待你一出門,她就可以跳在你身上,抓住你!所以我們要保守自己在基督裡.凡在基督裡的,她就奈何不得了!如經上所說:「我又賜給他們永生,他們永不滅亡,誰也不能從我手裡把他們奪去,我父把羊賜給我,他比萬有都大,誰也不能從我父手裡把他的奪去.」(約十28-29)
\newline
\newline
總而言之,罪的行動是連續性的,你靜坐時她會來,你作工時他也會來；她能把罪行如電影放映給你看,她也能把罪中之樂播音給你聽；只有在基督裡,她就無法施她技倆?願主憐憫,叫我們常住在主裡面.
\newline
\newline
罪的含義
\newline
\newline
罪包含著救恩,說起來你可能難以置信.這字在舊約有一一六次譯作「贖罪祭」,三次「刑罰」,三次「除罪」,其餘都是「罪」.這是告訴我們；甚麼地方有了罪,甚麼地方就有刑罰,同時甚麼地方也有贖罪祭.神的恩惠慈愛,不特一生一世隨著大衛,也隨著一切犯罪者!只要我們肯認罪,無論多大的罪,多久的罪,都能得赦免.相反地,罪的刑罰也跟著罪!甚麼場合有了罪,甚麼場合就有罪的刑罰,千萬別以為可以僥倖,神的忍耐是無限的,卻不是無底的,千萬不耍玩烈火!(來十二 29)
\newline
\newline
當神對該隱說這話的時候,該隱一定知道這字一切的含義.他知道神所說的罪,也指贖罪祭說的.雖然他行得不好,贖罪祭卻在那裡等候他,戀慕他,甘心情願解決他「不好」,只要他肯接受.可惜該隱讓這千古難逢的機會溜之大吉!弟兄姊妹們:悔罪是有機會的,別讓機會失去,成為一失足成千古恨!(啟二21-22)
\newline
\newline
弟兄姊妹們:屍首在那裡,老鷹也在那裡；罪在那裡,罪的刑罰也在那裡；可是神的恩典也在那裡.罪的工價本是死,但罪也被神的恩典跟著,你要罪的刑罰,還是要神的恩典呢?神給我們有選擇的自由,神不勉強我們.若你覺得心中不安,良心在那裡控告你,不要輕忽良心聲音,散會後,找個安靜的場合,把良心帶到你意識裡的罪,一一認清楚,交給主,讓主搬出去.不要等待,明天或許不是你的!
\newline
\newline


\newpage

\section{第41屆~港九培靈會~林道亮~第十一講}
\label{sec:1496}
\textbf{認識神的大名}
\newline
\newline
連結: \href{https://www.hkbibleconference.org/session-message/view/1496}{\texttt{https://www.hkbibleconference.org/session-message/view/1496}}
\newline
\newline
\hyperref[sec:1497]{< < < PREV SERMON < < <}
~
\hyperref[sec:toc]{[返目錄]}
~
\hyperref[sec:1495]{> > > NEXT SERMON > > >}
\newline
\newline
我們討論過:「神的命令就是叫我們信祂兒子耶穌基督的名」,我們所信的是神的兒子,是耶穌,是基督,也因信祂而要彼此相愛.今晚,我們要進一步討論:神的兒子既是神自己,那麼神自己是誰呢?「耶穌」既是「耶和華的拯救」,耶和華是誰呢?換句話說:我們為甚麼要信?因為我們的神是神,也是耶和華!
\newline
\newline
神
\newline
\newline
我們都知道,聖經一開始就提到神,「起初神創造天地.」神這個字(共三式),在舊約提了二千八百多次,可知它的重要性了.我們個人所信的是神,教會所敬拜的是神；聖經是神的話,宇宙是神所創造的,那麼我們所信的神是怎樣的呢?
\newline
\newline
神字在中文,是由一個縮寫的名詞和一個動詞合成的.左邊是縮寫的「神」字,右邊是「申」字,引申的意思.照說文,神的意思就是「天神引申萬物」.這樣在創世記第一章用神字再合適也沒有了,因為這一章告訴我們,神怎樣創造萬物.當我們每次提到「神」字的時候,讓我們記得,宇宙萬物是祂創造的啊!
\newline
\newline
神字在有些聖經譯作上帝.上帝這名稱,在四書五經裡約提過一百五十次,那麼上帝是甚麼意思呢?根據說文:「帝者,諦也,謂天上審判者也.」從前有皇帝管理人民,宇宙中也有一位審判者審判萬民.因他是高高在上,所以稱為上帝.這樣,說到創造萬物,用神字比較合適；說到管理萬物,上帝就比神字更貼切.無論如何,神太偉大,廣大,地上沒有一個名稱能把祂表現得完全!
\newline
\newline
那麼,聖經中的神到底是何所指呢?普通研究舊約的人,把這神字歸納成以下三個意義:
\newline
\newline
(一)極有能力,如經上記著說:「我手中原有能力害你,只是你父親的上帝昨夜對我說,你要小心,不可與雅各說好說歹.」(創卅一29；參箴三27；彌二1)
\newline
\newline
在這裡,「能力」二字原是與神同字.這名詞在舊約中有時是複數,有時是單數.單數時,若是個專有名詞,就是神；若是個普通名詞,便是能力.例如「我是伯特利的神」(創卅一13),「築一座壇給神」(創卅五1-3),這二神字和能力是同字的.當希伯來人要表示一個名稱偉大或無限時,常把單數變成複數.所以舊約神字十分之九是複數字.如果有人對你說,希伯來人是信多神的,因為神字在舊約是複數.說這話的人,不過顯他的無知而已,他不知道希伯來文複數的意義有十多種啊!
\newline
\newline
我們所信的神是大有能力的,祂創造宇宙,創造奇妙的人.我們也知道,祂肯施恩給我們,那麼為甚麼我們這麼軟弱可憐呢?我常說,真理是有客觀和主觀兩方面的,客觀能成為我們主觀的有多少,是在乎我們能支取多少而定.例如,宇宙間充滿了空氣,肺大的便多吸些,肺小的便少吸些.神是大有能力的,為甚麼我們不能得祂的大能力呢?是因我們的信心小.為甚麼我們的信心不夠大呢?恐怕是因我們對神的認識不夠多吧!神是大有能力的,祂能創造宇宙,祂能興國亡國,也能使亡了二千年的以色列國家重建,在祂沒有難成的事,這是我們所信的神!
\newline
\newline
(二)極其偉大.神字在舊約有時也作形容詞用,好像我們有時把天字當作形容詞說:「天大的事有我擔當!」在聖經中,例如:「他們便起行前往,上帝使那周圍城邑的人都甚驚懼,就不追趕雅各的眾子了.」(創卅五5)這裡的「甚」字原是「神」字,意思是「極其」!「約拿便照耶和華的話起來,往尼尼微去,這尼尼微是極大的城,有三日的路程.」(拿三3)這裡的「極」字也是「神」字,形容「大」字.在中文有一次把「神」譯作「高」字:「你的公義,好像高山,」(詩卅六6)原是「好像神的山」.(英文譯作偉大.)新約也有一次把「神」譯作「非凡」.司提反說:「那時,摩西生下來,俊美非凡......」(徒七20)「俊美非凡」原是「神美」.這樣,神字第二個意思,就是極其偉大,或說至大.我們所信的神是極其偉大的,可惜我們常把祂縮到極小!我們都知道馬利亞在以色列女子中是最蒙福的,她有甚麼特長作我們主的母親呢?是因她的「我心尊主為大」(路一46-47).這句話中文譯得太文雅些,原意是「我魂把主放大」.這是馬利亞的秘訣,她知道怎樣在她情感,心思,意志裡把主放大,所以她的靈能在神她的救主裡為樂.我們不能在神裡享受喜樂,是因我們把神縮得太小,沒有把神放大!願主憐憫我們,叫我們能把神放大,就能以神為樂.請放心,我們把神無論怎樣放大,斷不致把神放得太大!祂是充滿萬有的,毋怪大衛說:「我若升到天上,你在那裡；我若在陰間下榻,你也在那裡!」(詩一三九8)
\newline
\newline
(三)極其可畏.神本名為「可畏者」,如雅各對拉班所說:「但願亞伯拉罕的神,和拿鶴的神,就是他們父親的神,在你我中間判斷；雅各就指著他父親以撒所敬畏的神起誓.」(創卅一53；參上章四十二節,神字旁亦應有三點,惜譯者遺忘.)在中文聖經裡,凡一字旁邊有三點,指該字並非原文,乃譯者為使讀者容易明白加上去的.有時加上去確使原意更明顯,有時反叫原意更暗昧,讀者要留意.本節應是「雅各就指著他父親以撒的可畏者起誓.」這「可畏者」並非和「神」同字,也不是和「神」字同根；但神字本身有「畏懼」的意義.我們知道亞拉伯人稱神作「亞拉」,「亞拉」這字和「神」字同根,由動詞「畏懼」轉出來的.這樣,「神」本有「畏懼」的意思,所以聖經用另一字稱祂為「可畏者」.有時我們單注意神的慈愛,忘記了神的可怕,甚至有人說:「在舊約我們要怕神,在新約我們只要愛神.」這樣的話,是一知半解的說法,最容易使人上當!讓我們查一查新約,是否我們不需要怕神?「你們既稱那不偏待人,按各人行為審判人的主為父,就當存敬畏的心,度你們在世寄居的日子.」(彼前一17)本節「當存敬畏的心」應是「在敬畏裡」.「所以我們既得了不能震動的國,就當感恩,照神所喜悅的,用虔誠敬畏的心事奉上帝.」為甚麼要敬虔呢?「因為我們的神是烈火.」(來十二28)神是我們的慈父,也是背逆之子的烈火,惟有智慧之子知道,愛父和怕神應雙方並重!
\newline
\newline
以上是根據字根所歸納的三個意義.現在讓我們查一查這三個意義是否合聖經的定義.在申命記記著說:「因為耶和華你們的神他是萬神之神,萬主之主,至大的神,大有能力,大而可畏,不以貌取人,也不受賄賂.」(十17)這裡告訴我們:神是至大,大有能力,大而可畏.在尼希米書也有同樣的語法:「我們的神阿,你是至大,至能,至可畏,守約施慈愛的神.」(九32)這裡告訴我們,我們的神是「三至」；至大,至能,至可畏.
\newline
\newline
末世的教會不像教會,就是因為忘記了神的三「至」.神是至大的,教會卻把祂縮小到甚麼都無祂份!神家裡的事情,甚麼都憑善惡果的知識來斷定是非,把神的法令丟在腦後,甚至把祂關在門外,讓祂在門外叩門!(啟三20)神是至能的,教會卻以為祂甚麼不能:祂不能行神跡,祂不能聽祈禱,祂是垂死的老人；甚至有人說祂已死!教會這樣褻瀆神,怎能希望祂顯大能呢!神也是至可畏的,教會卻對祂視若無睹,造成末世危險的日子,叫人專顧自己,貪愛錢財......違背父母,忘恩負義,無親情,不解怨......賣主賣友,任意妄為,自高自大,愛宴樂不愛神,有敬虔的外貌,卻背了敬虔的實意!(提後三1-5)這等人,神叫我們都要躲開,怎能希望祂自己和這等人同在呢!毋怪末世教會越過越死,許多傳道人的眼睛也越過越瞎!他們走遍洋海陸地,勾引一個人入教；既入了教,卻使他作地獄之子,願主憐憫我們,叫我們知道為甚麼要信?並所信的是誰?我們的神是至大,至能,至可畏,所以我們要信!
\newline
\newline
耶和華
\newline
\newline
神既是至大無邊,至能無比,也是至可畏無雙,要給祂自己一個適當名字,實非易事!雖然這樣,祂還是設法給我們知道祂的尊名.當摩西請問祂大名的時候,祂對摩西說:「我是自有永有的,又說,你要對以色列人這樣說,那自有的打發我到你們這裡來.」(出三14)
\newline
\newline
這裡中文又是譯得太文雅,反叫原意不明.當摩西問神叫甚麼名字時,神說:「我是自有永有,」直譯應是「我是就是我是」；他又說:「那自有的,」直譯應是「那我是」.這樣,神的大名就是「我是」.
\newline
\newline
「我是」是神的大名,卻不是「耶和華」的意思!所以若有人問你「耶和華」的意思,你說「我是」,你就錯了.因為「我是」和「耶和華」,雖然是同一動詞,卻不同身；「我是」是第一身,耶和華是第三身.當神稱祂自己時當然應說「我是」,當我們稱祂時便要說「他是」了.所以耶和華的意思是「他是」.這意思大衛很懂得,因為他曾說:「耶和華是(他是)我的巖石,我的山寨,我的救主,我的磐石,我所投靠的；他是我的盾牌,是拯救我的角,是我的高臺.」(詩十八 2)這樣,耶和華的大名很像一本支票簿已簽上了存款人的名,只要我們在「他是」下面填上我們的需要去支取,我們一切的需要都能獲得.所以到神面前來的人,必須要有這支票簿,正如希伯來書所說:「人非有信,就不能得神的喜悅；因為到神面前來的人,必須信有神,且信他賞賜那尋求他的人.」(十一6)本節的「信有神」原文是「信他是」.所以到神面前來的人,必須信他是!這是我們祈禱的秘訣,我們要信「他是......」,不疑惑祂能否幫助我們.耶和華甚麼都有,甚麼場合都在!你要能力嗎?祂是你的能力；你要智慧嗎?祂是你的智慧；你要平安,喜樂,自由嗎?祂是你的一切!讓「我們只管坦然無懼的,來到施恩的寶座前,為要得憐憫,蒙恩惠作隨時的幫助.」(來四16)
\newline
\newline
我常感覺到,神為自己取名「祂是」,真是一件絕頂智慧,深入淺出的啟示.因為人類文字中,除此以外,實無其他能啟示祂的真相.祂若說自己名叫「慈愛」,祂就沒有啟示祂的公義:祂若說自己是智慧和能力,祂就沒有告訴人祂的聖潔和信實.所以祂必須是「祂是」,才能啟示祂一切,也能讓祂救贖的人支取一切.
\newline
\newline
祂的大名是「他是」,不是「他曾是」,也不是「他將是」.因為在祂沒有過去和將來,只有現在,正如希伯來書告訴我們:「耶穌基督,昨日今日一直到永遠是一樣的.」(十三8)「他是」昨日今日一直到永遠一樣的偉大,一樣有能力,一樣可畏.祂在舊約時興邦亡邦,今天祂是；祂在新約時醫病趕鬼,今日他也能.在祂沒有改變,也沒有轉動的影兒,所以我要信祂!
\newline
\newline
我們討論聖經文字的時候,我們要注意兩方面的意義:一是文字的意義,一是聖經的意義.有時一個字兩方面的意義是相同的,有時文字有文字的意義,聖經又有聖經的意義.當我們研究神字,我們發現兩方面的意義是相等的；但「耶和華」的兩方面意義卻不是相等,乃是相成的.耶和華的字義是「他是」,聖經中的定義是把「他是」甚麼摘要地舉出來.耶和華說:「耶和華是有憐憫,有恩典的神,不輕易發怒,並有豐盛的慈愛和誠實.為千萬人存留慈愛,赦免罪孽,過犯,和罪惡；萬不以有罪的為無罪,必追討他的罪,自父及子,直到三四代.」(出卅四6-7)
\newline
\newline
這裡清楚告訴我們,耶和華是怎樣的神.當我們研究「神」的時候,我們看不見神的全面,現在藉著「耶和華」,神就把祂另一面啟示給我們.在這兩節聖經裡,神啟示了三點關於耶和華的特性；一,耶和華是慈愛的；二,耶和華是信實的；三,耶和華是公義的.
\newline
\newline
憐憫和恩典是神愛兩方面的表顯；憐憫是消極方面,我有罪,祂赦免我；恩典是積極方面；我有軟弱,祂剛強我.憐憫是拯救我脫離死亡,恩典是賜給我新生命.因為祂憐憫我,祂把我的苦杯倒盡；因為祂恩待我,祂叫我的福杯滿溢.我們一而再,再而三地能得到神的憐憫和恩典,是由於祂不輕易發怒.這一切都因為祂的名是耶和華,祂是慈愛的神.
\newline
\newline
耶和華不但是慈愛的神,也是「有豐富的慈愛和誠實」的神.這裡「慈愛」二字,譯文不很準確.這字在舊約共用了二四七次,中文常譯作「慈愛」實際是「共同責任,休戚相關,利害與共」的意思.所以本字常和「誠實」或「信實」(參創廿四27,書二14,詩八十九2,24,33,49)或「守約」(參申七 9,代下六14,尼一5)同用,這裡譯作共同責任,諒和原意相近些.耶和華是富於共同責任和誠實的,祂說了怎樣,就是怎樣.天地要廢去,祂的信實卻存到萬代.我們縱然失信,祂仍是可信的；因為祂是耶和華,祂不能背乎自己!可惜末世的教會不負責任,說是認識神,行事卻和祂相背；說是事奉主,為人卻以自己為中心,毋怪耶穌寫信給老底嘉教會時,說自己是「那為阿們的,為誠信真實見證的」,反映教會在末世時代的不誠實,說假阿們和作假見證,那麼我們為甚麼要信!不是因為教會的可靠,也不是因為信徒的可愛,乃是因為耶和華向我們的信實和可靠.耶和華是應當讚美的!
\newline
\newline
最後,也是最重要的,是祂「不以有罪的為無罪,必追討他的罪,自父及子,直到三四代」.我們常注意的慈愛,忽略了神的公義,以致我們的信心失去了平衡.我們有時會錯想,以為神既是慈愛的,對於我們一些思想上的錯誤,生活上的不正軌,祂一定滿不在乎,姑息寬容的.這是我們自己欺騙自己,祂的訓示是:「若有一文錢沒有還清,你斷不能從那裡出來!」相反地,我們若認自己的罪,無論我們的罪多大,多深,祂是信實的,是公義的,必要赦免我們的罪,洗淨我們一切的不義!
\newline
\newline
弟兄姊妹們:你聽到這裡,你的良心有否指責你所隱藏的罪呢?不要學取巧,不要想僥倖,我們騙不了神,祂斷不以有罪為無罪!請你在聚會後找個安靜的地方,好好地同神算一算賬,不要遲延,因為明天不一定是你的!請聽主對推雅推喇教會說的話:「我曾給他悔改的機會,他卻不肯悔改他的淫行.看哪,我要叫他病臥在床,......同受大患難!」(啟二21-22)
\newline
\newline


\newpage

\section{第41屆~港九培靈會~林道亮~第十二講}
\label{sec:1495}
\textbf{以馬內利}
\newline
\newline
連結: \href{https://www.hkbibleconference.org/session-message/view/1495}{\texttt{https://www.hkbibleconference.org/session-message/view/1495}}
\newline
\newline
\hyperref[sec:1496]{< < < PREV SERMON < < <}
~
\hyperref[sec:toc]{[返目錄]}
~
\hyperref[sec:1494]{> > > NEXT SERMON > > >}
\newline
\newline
我們今晚仍是討論為甚麼要信耶穌,答案是:為了「以馬內利」.
\newline
\newline
以馬內利是耶穌的另一個名字,在未研究這名字以前,讓我們先研究甚麼是宗教.今天我收到一些問題,總括起來就是說:世上有些人立心行善,不作虧心事,他們和基督徒有甚麼分別?這樣的問題,和說世上所有的宗教都是好的；宗教都是叫人為善,基督教也不過如此,那何必一定要信基督教呢?
\newline
\newline
宗教二字是由外國轉來的.在辭源中有以下的定義:「以神道設教,設立誡約,使人信仰者也.」在辭海中的定義是:「有所宗,以為教者也.」以上的定義,根據中文「宗教」二字來解釋,也說得過去；宗就是崇拜意思,教就是教訓.這樣,佛教,道教,基督教......都有個別的崇拜和教訓,世上也就真的有許多宗教了.但我們要問:這是否宗教原來的意義?宗教的原意是「綁回去」,意思是把「人和神聯繫在一起」.所以在任何宗教裡,沒有這種「把人和神聯繫在一起」的本質和方法的,便不能稱為宗教.這樣,儒教便不能稱為宗教.雖然後期儒教有人提倡「天人合一」,只不過是摸不著邊際的理論而已.嚴格地說來,原始的佛教根本是無神的,更談不到神人的聯繫了!總之,世上真正的宗教根本很少.就是那些有「神人聯繫在一起」成份的宗教,若不是成份不純粹,就是方法不正確.結果,芸芸眾生,正像聖經所說的「困苦流離,如同羊沒有牧人一般!」這就是耶穌為甚麼要降生──設立一個真宗教!
\newline
\newline
在馬太福音一章一至十八節內,我們看見神揀選一個民族,建立一個國度,維持一個宗教,都未成功.直到耶穌基督降世,神所想建立的民族,國度和宗教才得實現.在本章二十一節說:「她將要生一個兒子,你要給他起名叫耶穌:因他要將自己的百姓從罪惡裡救出來.」這是說耶穌要建立一個新的民族.在二十二節又說:「這一切的事成就,是要應驗主藉先知所說的話說:『必有童女,懷孕生子,人要稱他的名為以馬內利』」.這就是說神要藉著耶穌創立真的宗教.「以馬內利」是一個三合一的名稱:「以馬」就是同在,「內」就是我們,「利」就是神.原意就是「神與我們同在」.這就是宗教的本質,神要藉著耶穌成就祂創世的計劃──神人合一!這也是我們信耶穌的原因,因為除祂以外,在天上人間,沒有賜下別名,我們可以與神和好,和神同在!現在讓我們分四點討論如左:
\newline
\newline
神要以馬內利
\newline
\newline
亞當被造後,神就與人同在,和人有交通.但從罪入了世界後,死(與神隔絕)也來了.自此以後,罪在那裡,人與神的隔膜也在那裡,無法與神同在.雖然如此,神要與人同在的計劃卻從未改變.當神救以色列人出埃及後,神叫他們造會幕.祂說:「又當為我造聖所,使我可以住在他們中間.」(出廿五8)神要他造會幕的目的,是要和他們同在.神在會幕中向他們顯現,他們在會幕敬拜神,藉著會幕神可以和他們同在.會幕是主耶穌的預表,所以今天我們也在主耶穌裡與神同在.以後,以色列人立了國,神就叫所羅門王造聖殿.當他為聖殿行奉獻禮時,神的榮耀充滿聖殿,如列王記上所說:「祭司從聖所出來的時候,有雲充滿耶和華的殿；甚至祭司不能站立供職,因為耶和華的榮光充滿了殿.」(八10-11)
\newline
\newline
猶太人知道,這榮光就是證實神與他們同在了.可惜以後他們演成了一種迷信,以為耶路撒冷有聖殿,聖殿有神的同在,耶路撒冷永遠不會滅亡.他們卻不知道,在他們犯罪的時候,神的榮光離開了他們.讓我們來看一看神榮光離開他們的情形.
\newline
\newline
在以西結書第八章時,雖然他們無惡不作,神仍停留在殿中,如以西結所說:「......在神的異象中,帶我到耶路撒冷朝北的內院門口,在那裡有觸動主怒偶像的坐位,就是惹動忌邪的.誰知在那裡有以色列神的榮耀,形狀與我在平原所見的一樣.」(三4)但是因為殿內充滿著觸動主,惹動主的污穢,所以在以西結十章時,耶和華就開始離開聖殿,向上升.可是祂仍不忍立時離開,所以「耶和華的榮耀從基路伯那裡上升,停在門檻以上.」(4)當祂停在門檻上時,祂仍讓祂榮光充滿殿的院宇!雖然祂不忍離開以色列人,終因他們所行一切「大可憎的事」(結八6,13,15),祂從門檻上升,停在半空,在基路伯以上.「耶和華的榮耀從殿的門檻那裡出去,停在基路伯以上.」(18)那時基路伯停在殿的東門口,神在他們以上,說明神仍是不忍離開那為祂自己名所建造的殿.祂真是有憐憫,有恩典的神,不輕易發怒的耶和華.就是在祂離開耶路撒冷後,祂仍停在橄欖山上不走!「耶和華的榮耀從城中上升,停在城東的那座山上.」(結十一23)橄欖山實在是眾山中最有福的,在這山上,神子曾為耶城哀哭(路十九41),為我們升天(路廿四50,徒一12),並日後為我們再來(亞十四4),聖潔的神,沒法和可憎可厭之物妥協,在忍無可忍之後,終於離地上升.「靈將我舉起,在異象中藉著上帝的靈,將我帶進迦勒底地,到被擄的人那裡；我所見異象就離我上升去了.」(結十一24)
\newline
\newline
神如何渴慕要與人同在!可是人偏要行可憎可厭的事惹動主怒,叫神沒法與人同在!末世教會更是變本加厲,把主逐出教會,讓祂在漫漫長夜裡,受盡風霜!(參歌五2,啟三20)請問各位:神有否離開你的教會?教會若沒有神同在,那教會就不是宗教,她的存在不過等候主審判而已!
\newline
\newline
神離開以色列人後,新舊約間靜默了四百年,沒有異象,沒有啟示,直到道成了肉身,神才再與人同在.祂住在我們中間,充充滿滿有恩典有真理,叫一切信祂的,愛祂的,可以和神同住(約十四23)當祂升天前,祂應許和我們同在,直到世界的末了(太廿八20).這是真宗教,神所要的「以馬內利」.有一天這神人聯繫在一起,要實現,正如啟示錄所記:「我聽見有大聲音從寶座出來說,看哪,神的帳幕在人間；祂要與人同住,他們要作祂的子民,神要親自與他們同在,作他們的神.」(廿一3)這樣,神建立宗教是為國度,宗教是國度的前身,也是以馬內利的前半；宗教和國度併在一起,才是整個以馬內利,就是神所要的.
\newline
\newline
人要以馬內利
\newline
\newline
人和動物不同就是因為人有宗教性.雖然人的宗教性敗壞,不能了解神的事情,可是因了神的恩典,宗教性起碼的功能仍舊存在.由於那個功能,人類對宗教仍有以下四點感想:
\newline
\newline
一,宇宙中確有一個偉大的權力.雖然人們給祂各種不同名稱,這權力的存在是無可置疑的.
\newline
\newline
二,人類的禍福是根據人與那權力的關係而定.如果關係好,人會得福；如果不好,人會受禍.
\newline
\newline
三,今日世界的痛苦,人心的恐慌,不幸的災禍,都由於人和這權力關係不正常.起了衝突.
\newline
\newline
四,有沒有方法可以使人與那權力的關係正常,衝突消弭.
\newline
\newline
這樣說來,人的本性就是想以馬內利,想與神調整關係,想和神親近.這就是有些天然宗教的出發點,可惜因為它們的方法不正確,結果等於緣木求魚!
\newline
\newline
事實,不但人要以馬內利,就是一切受造之物都在盼望以馬內利.經上記著說:「受造之物,切望等候神的眾子顯出來......指望脫離敗壞的轄制,得享神兒女自由的榮耀.」(羅八19-21)「因為萬有都是本於他,倚靠他,歸於他.願榮耀歸給他,直到永遠.阿們!」(羅十一36)這就是以馬內利!
\newline
\newline
如何活出以馬內利
\newline
\newline
信主的人,有基督在內心作主,這就是以馬內利──神和我們同在.神是至大,至能,至可怕的.我們有這麼一位神在裡面,那末我們也應該是至大,至能,至可怕的.但往往因為沒有信心的緣故,我們不能把神的至大至能表現出來.今晚我要向各位作三十多年前第一次趕鬼的見證.有一個夏天,我由學校回到家鄉,附近一個內地會的禮拜堂,有一位長老來請我下午禮拜講道.我仍記得,講題是「城造在山上是不能隱藏的.」講完後,我自己覺得還講得不差.正要散會的時候,那請我來講道的長老上前來謝謝我,同時他要我為一位被鬼附了的姊妹祈禱趕鬼.我聽了之後,吃了一驚!雖然我以前跟過父親去趕鬼,自己卻從沒有嘗試過.可是剛才自己講道講得那麼起勁,現在如果說不能趕鬼,豈不是太難為情!我到底答應了他.我請他把那姊妹帶到講台前面.我問他附的是甚麼鬼,他說是啞巴鬼,這姊妹已有二年未說過話.我知道這事相當棘手,只好看神肯不肯施施憐憫了!我就請會眾一起跪下來替這可憐的姊妹祈禱,我自己當然很懇切地禱告.祈禱後,我便站起來仿效我父親的手法,用手指著那姊妹說:「魔鬼,我奉主耶穌基督的名叫你出來,」可是她一聲也不響,我知道祈禱未曾生效,心裡著忙起來!但除了禱告外有甚麼辦法呢?於是我再請會眾一起祈禱,希望神會憐憫我們.在二次祈禱後,我用同樣手法,同樣的話語,那姊妹不特不出聲,反用眼瞪著我!我那時真是上天無路,入地無門!我幾乎哭出來!可是除了祈禱外,有甚麼辦法呢?我請大家再次禱告,我自己可說和神摔交一樣地祈禱.我流淚,我求神顯出祂的榮耀,不然的話,我怎能傳道呢?正在懇切祈禱的時候,我忽然覺得心中起了一股信心,覺得神已聽了祈禱.我便起來,用手指指著地說:「撒旦,我奉主耶穌基督的名叫妳一定要出來!」我緊張得連手指也發抖,當我在說「一定要」的時候,我不自覺地用手指推了她的額頭一下,就在這一推之下,她對我說:「先生!謝謝你!」那時,整個會眾的快樂與奮興,真是非筆墨所能形容的!讚美主,就是這樣,祂給我第一次趕鬼的經驗.雖然我不配,我無知,祂仍恩待我.縱然我失信,祂仍是可信的.哈利路亞!弟兄姊妹們:在你我裡面的神是一樣的,祂從不偏待人.讓我們用信心來表彰祂的內在,因為祂所賜給我們的,不是膽怯的靈,乃是剛強,仁愛謹守的靈.(提後一7)
\newline
\newline
目前有些基督徒,因為先天或後天的缺陷,常常有自卑感.自卑感的實情是:一個人內心的深處覺得自己不如人,可是在外表上不肯示弱,偏要相反地用驕傲和自大來自誇.這樣的人,目前在社會中越來越多.可是基督徒卻不應有自卑感,因為我們裡面有至大,至能,至可畏的神,願主憐憫我們,叫我們看見這真理的真實性,不然的話,教會裡的問題會更多!
\newline
\newline
在我們裡面的,不特是神!祂也是耶和華.祂是慈愛的,所以我們要儘量憐憫人,恩待人.別人得罪自己要饒恕,弟兄姊妹的需要我們儘量幫助.祂也是負責的,誠實的.末世教會就是不誠實,不負責.社會上充滿了虛謊,毫無責任心,我們一定要誠實,一言既出,駟馬難追.行事為人,寧人負我,毋我負人.不然的話,我們就不能彰顯住在我們裡面的耶和華,那真理的靈也就不能在我們裡面運行了.我們作基督徒的,同時要公義.許多人把神當作歡喜菩薩,笑面佛,忘記了神也是烈火!也有人以為作基督徒應該笑口常開,即使別人有大逆不道,我們也應該說不要緊.弟兄姊妹們:我們是慈愛的,但也是公義的.我們應該不對的就是不對,對的就是對,為真理作見證.教會是真理的柱石和根基,我們應該建立真理,支持真理.對異端,對不法,「未罵先笑,未打先倒」,不是真基督徒的作風；不然的話,基督也不會被釘在十字架上了.今天教會之所以到這地步,就是因為許多弟兄姊妹不敢為真理說話!我們要為主盡忠.耶和華既斷不以有罪的為無罪!我們自己有罪一定要承認,不承認便得不到赦免.如果弟兄姊妹有錯,我們一定要憑愛心直說.不要各人自掃門前雪,因為一個亞干的罪能令全以色列失敗啊!
\newline
\newline
如何得著以馬內利
\newline
\newline
現在,我們要看怎樣才能得著以馬內利?我們要有以馬內利一定要注意客觀和主觀兩方面.客觀方面請看詩篇六十八篇十八節；「我已經升上高天,擄掠仇敵,你在人間,就是在悖逆的人間,受了供獻,叫耶和華上帝可以與他們同住.」這裡「供獻」二字,中文譯者又譯得太文雅,原是「禮物」.主耶穌升上高天的時候接受了禮物,有了這禮物,才能叫神與人同住,叫以馬內利成為事實.這禮物是甚麼呢?羅馬書告訴我們,這禮物是永生.「因為罪的工價乃是死.惟有上帝的恩賜(原文禮物),在我們的主基督耶穌裡,乃是永生.」(六23)永生就是那禮物,沒有那永生,以馬內利就沒法實現,神和人就沒法合一.
\newline
\newline
以上所說的是客觀方面,或說是主為我們所預備的.現在讓我們討論主觀方面,或說怎樣才能叫主的禮物(生命)成為我們的.使徒約翰告訴我們說:「人有了神兒子就有生命,沒有神兒子就沒有生命.」(約壹五12)我們必須接受基督才能有生命!換句話說,我們要有以馬內利(神人合一),必須接受以馬內利(基督).「凡接待祂的,就是信祂名的人,祂就賜他們權柄,作神的兒女.」(約一12)
\newline
\newline
弟兄姊妹們:我感覺到有責任要問你；自己有永生嗎?以馬內利在你裡面有否成為事實呢?你有沒有詳細地認過你的罪,接受耶穌到你內心?如果你沒有,那你還沒有主的生命,你還未曾以馬內利!一天,你會後悔說:「麥秋已過,夏令已完,我們還未得救!」(耶八20)神要以馬內利,你心中的深處也要以馬內利,以馬內利一直在那裡等候你,你肯否打開心門接受以馬內利呢?耶穌說:「到我這裡來的,我總不丟棄他!」(約六37)
\newline
\newline


\newpage

\section{第41屆~港九培靈會~林道亮~第十三講}
\label{sec:1494}
\textbf{持守末後的救恩}
\newline
\newline
連結: \href{https://www.hkbibleconference.org/session-message/view/1494}{\texttt{https://www.hkbibleconference.org/session-message/view/1494}}
\newline
\newline
\hyperref[sec:1495]{< < < PREV SERMON < < <}
~
\hyperref[sec:toc]{[返目錄]}
~
\hyperref[sec:1493]{> > > NEXT SERMON > > >}
\newline
\newline
今天晚上我們要討論「望」的問題,願主恩待我們,叫我們看見我們所望的是甚麼.在這末世的教會裡,許多信徒不知道自己的盼望是甚麼,所以生命沒法長進,愛心沒法加增!現在請讀以下一段聖經:
\newline
\newline
「耶穌往耶路撒冷去,在所經過的各城各鄉教訓人.有一個人問他說:「主阿,得救的人少麼?」耶穌對眾人說:「你們要努力進窄門；我告訴你們,將來有許多人想要進去,卻是不能!及至家主起來關了門,你們站在外面叩門說:「主阿,給我們開門；」他就回答說:「我不認識你們,不曉得你們是那裡來的!」那時,你們要說:「我們在你面前喫過喝過,你也在我們的街上教訓過人.」他要說:「我告訴你們,我不曉得你們是那裡來的；你們這一切作惡的人,離開我去罷.」你們要看見亞伯拉罕,以撒,雅各和眾先知,都在上帝的國裡,你們卻被趕到外面；在那裡必要哀哭切齒了.從東,從西,從南,從北,將有人來,在上帝的國裡坐席.只是有在後的將要在前,有在前的將要在後.」(路十三22-30)
\newline
\newline
本段經文告訴我們:有人問主耶穌,得救的人少嗎?主耶穌不回答他多或少,只是鼓勵和警告他們,要他們努力進窄門.在這裡有一個問題；得救是白白的,還是要努力爭取的?今晚希望藉著這問題,來討論盼望甚麼,怎樣盼望.
\newline
\newline
盼望甚麼
\newline
\newline
我們都知道救恩是白白的,舊約這樣說(賽五十五1),新約這樣說(羅三24,啟廿二17),那坐寶座的也這樣說:「我是阿拉法,我是俄梅戛,我是初,我是終；我要將生命泉的水白白賜給那口渴的人喝.」(啟二十一6)可是本段經文卻清楚告訴我們要努力,並且是主親口告訴我們的!難道我主會自相矛盾?斷無此事,我們的主從不出乎反乎!在這裡,或許有人說:「是否努力二字另有別解呢?」好罷,讓我們先查一查努力這二字.
\newline
\newline
努力二字在新約共用過七次,共有七種不同的繙譯.在七次中耶穌用過二次,其餘的都在保羅書信中,耶穌除本文外,在祂受審於彼拉多前時也用過一次；「耶穌回答說,我的國不屬這世界,我的國若屬這世界,我的臣僕必要爭戰,使我不至於被交給猶太人,只是我的國不屬這世界.」(約十八36)這裡「爭戰」二字是與「努力」同字的.那麼「努力進窄門」也是「爭戰進窄門」了.保羅在提摩太前後書各用過一次.他說:「你要為真道打那美好的仗,持定永生.你為此被召,也在許多見證人面前,已經作了那美好的見證.」(提前六12)他又說:「那美好的仗我已經打過了,當跑的路我已經跑盡了,所信的道我我已經守住了.」(提後四7)在這二節經文裡,「打仗」二字是與「努力」同字的.這樣,耶穌是說:「你們要打仗進窄門!」打仗可不是容易的事!要受訓練,要有兵器,要吃苦耐勞,要機警靈巧；對領袖要服從,對同志要互助；不然的話,就永無打勝仗的日子!在歌羅西書保羅用過二次,中文譯為「盡心竭力」和「竭力」(西一29,四 12).此外,在哥林多前書一次,中文譯作「較力爭勝」(林前九25).這樣,得救真的這麼難嗎?怎麼又說白白的呢?得救到底是白白的呢?還是要打仗的呢?要解答這問題,「努力」二字既然不能給我們亮光,我們試在「得救」二字上找出路.
\newline
\newline
細讀新約,我們知道得救是有三方面的,正如保羅所說:「他曾救我們脫離那極大的死亡,現在仍要救我們,並且我們指望他將來還要救我們.」(林後一 10)我們的主是昔在今在永在的救主.祂在已往救我們,現在天天救我們,將來還要救我們.已往的救恩是白賜的,正如保羅說:「神救了我們,以聖召召我們,不是按我們的行為,乃是按他的旨意和恩典；這恩典是萬古之先,在基督耶穌裡賜給我們的.」(提後一9)現在的救恩是本乎恩的,也藉著祈禱.保羅曾說「弟兄們,我還有話說,請你們為我禱告,好叫主的道理快快行開,得著榮耀,正如在你們中間一樣；也叫我們脫離無理之惡人的手,因為人不都是有信心.」(帖後三 1-2)保羅在世時,有許多無理的惡人橫行不法,保羅希望從他們手中得拯救.怎樣才能得拯救呢?「請你們為我們禱告.」祈禱,代禱是天天得救的秘訣.所以要天天得拯救,單是信不夠,還要禱告才行.禱告的實際,並非人軟弱的表示,乃是屬天的手續,能把客觀的福份成為主觀的.神的應許和拯救都是客觀的,我們要應用祈禱的手續,把神的一切成為我們的.這手續不特保羅知道,大衛也很認識它的重要性.(撒下七章)當神應許了他國度和寶座之後,他立即為神所應許的一切禱告,證實神的應許.將來的救恩是救我們進入祂的國度.在提摩太後書記著說:「主必救我脫離諸般的兇惡,也必救我進他的天國.」(四18)怎樣才能被救入天國呢?就是本章第七節所說的:「那美好的仗我已經打過了；當跑的路我已經跑盡了；所信的道我已經守住了；從此以後,有公義的冠冕為我存留......」保羅是一個得救的人,他為甚麼要打美好的仗呢?因為他盼望主會救他進入天國.討論到這裡我相信各位都已明白:耶穌叫我們努力的救恩不是過去的,乃是救人未來的天國.所以我們要爭戰,要較力爭勝,要盡心竭力,才能進入天國,或說得末世的救恩；否則,我們會被關在門外,因為有一天家主會起來關門的!
\newline
\newline
關在門外確是個嚴重的問題,雖然有許多基督徒置若罔聞!記得當日挪亞造方舟的時候,洪水未來以前,挪亞苦口婆心地勸人進方舟,別人反譏笑他是神經病；等到洪水真的來了,人們就來在方舟外邊叩門!可是已太遲了,挪亞愛莫能助,毫無辦法,因為方舟的門是耶和華關的；所以連造方舟的工匠們也不得進入!弟兄姊妹們,請記得:耶和華的慈愛是無盡的,但不是無底的；到了時候,施恩時期一滿,那就一秒鐘也不再放寬了!
\newline
\newline
在先知以西結的時候,以色列人的背逆和頑梗已到了極點,神在他們身上的恩典也到了底,所以神說,即使有挪亞,但以理,約伯在其中,他們只能自己的性命,連他們的兒女們也救不了!相信在以色列亡國的時候,仍有不少義人向神呼求,但神不聽他們的禱告,因為施恩的時候已過去了!得救有機會,選上也有機會,機會一去不再來!
\newline
\newline
近年我有一感覺,似乎神在末世的時候聽人的禱告不像往昔那麼容易.神的作為也不像往昔那麼多方彰顯,所以甚至有人說聖經所說的神蹟異能是不可靠的,因為他們沒有看見這種事實.我自己的經驗也告訴我,許多時候,除非懇切不斷地祈禱,甚至禁食祈禱,很難得神的垂聽.我想；這難道就是家主快要起來關門的先聲嗎?目前有許多弟兄姊妹覺得要過敬虔的生活比以前更不容易,這也是否快要關門的預告呢?無論如何,聖經告訴我們,有一天家主一定要起來關門的.那時天要像銅,地要像鐵,人們「雖然號哭切求,卻得不著門路」!(來十二17)
\newline
\newline
怎樣盼望
\newline
\newline
神給我們福音的盼望(西一23)有失去的可能,正像主對非拉鐵非教會所說的:「我必快來,你要持守你所有的,免得人奪去你的冠冕.」(啟三11)國度原已經分配給我們,至於我們能否實際地承繼,要看我們能否持守.如果我們不能持守,就不能承繼這一份「存留在天上的基業」!(彼前一4)教會中有跌倒退後的信徒是免不了的,不然,聖經不會有「持守」或「堅持」等勸告.願主憐憫我們,別的信徒或會退去,不再和主同行,我們卻要對主說:「主啊,你有永生之道,我們還跟從誰呢?」(約六68)我們要持守,要較力爭勝地持守,不讓我們失去福音的盼望才是!如果我們要不失去這盼望,我們一定要留意以下二方面:
\newline
\newline
(一)消極方面──聚會,聽道不一定能持守盼望.本段經文告訴我們,在主否認那些人以後,他們向主辯護說:「我們在你面前喫過喝過.」這是說他們曾和主有過交通.有些信徒喜歡參加教會各種聚會,尤不肯放過教會聚餐的機會.每次舉行聚餐時必到,特別聚會時更活躍.這些原是好的,但不一定能持定盼望；因為這一切可能完全是為自己,不是為主!在我們的慾望中有合群慾.有人在社會上找不到這種慾望的出路,到教會來找到了,就在教會裡滿足他的慾望,他還以為是熱心愛主哩!這裡或許有人會問:「那末愛主和滿足合群慾有甚麼分別呢?」它們的分別是在以誰為中心.若他真的愛主,他一定在凡事上讓主居首位；如果他是為了滿足自己的合群慾,那他一定會凡事以自己為前提.這樣的人,他遲早會離開教會或團體,甚至會破壞教會,那裡談得到持守的問題!其實,如果一個人真愛主,除非是主對他不起,他怎會離開主所愛的教會或團體呢?為了這個緣故,教會中有不少交誼會,聚餐會等等,不特不能叫主得榮耀,反成為得罪主的機會,也有不少以鬧意見或紛爭收場!
\newline
\newline
那些人繼續為自己辯護說:「你也在我們的街上教訓過人.」這是說聽道不一定能持定盼望.有人常常聽道,甚麼聚會都參加,看來很虔誠,其實卻如保羅所說的:「常常學習,終久不能明白真道.」(提後三7)
\newline
\newline
從前上海有一位姊妹,常常赴宋尚節博士的聚會,每次當宋博士呼召說:「你們中間有誰以前賭錢,現在立志不再賭的,舉起手來!」她每次都照舉如儀,可是每次回家後也照賭如儀!真可說得常常聽道,終久不能明白真道!
\newline
\newline
我傳講福音的時候,總是強調信耶穌不是為作好人.信耶穌能作也應作好人,但不是為作好人!會後,常有人會對我說:「牧師,你的道講得真好,我們信耶穌為要作好人.」有不少信徒聽了幾十年的道,連基本的理都不清楚,這樣的信徒,他怎樣不失去盼望呢?因為神的國度如果給他進去,一定會給他敗壞了的!
\newline
\newline
(二)積極方面──行義,守法才能成全盼望,因為主的國度是重質不重量的.本段經文告訴我們,主拒絕那些人的第一原由,是因主不認識他們.主說:「我不認識你們,不曉得你們是那裡來的!」那末誰是主所認識的人呢?保羅說:「然而上帝堅固的根基立住了,上面有這印記說,主認識誰是他的人；又說,凡稱呼主名的人,總要離開不義.」(提後二19)這裡說主所認識的人,必要離開不義.「義」字在聖經中有二個意思:一是宣告無罪(稱義),一很近古文觀止的定義:「行而宜之」.例如撒迦利亞和以利沙伯:「他們二人,在上帝面前都是義人,遵行主的一切誡命禮儀,沒有可指摘的」.(路一6)這裡說他們稱為義人,是因他們在主的誡命禮儀上行而宜之.簡單地說,要給主認識,或說要進入國度,總要離間一切行而不宜的,力行一切行而宜的.
\newline
\newline
主拒絕他們的第二原因,是因他們作惡.主說:「你們這一切作惡的人,離開我去罷!」「作惡」二字原意是「不法」,或說「不守法」.這二字和馬太七章的「作惡」同字(23).那段聖經明說,將來連許多大佈道家,神醫家都不能進國度,因為他們不守法.在這裡,或有人問:「既然這些大佈道家不守法,為甚麼主仍用他們領人得救呢?」我們知道,使人歸主,不一定因為傳道人的好生命,多少時候,完全是神為了祂自己的緣故,應用傳道人的恩賜,施行拯救.例如你有一園丁,技術高明,常把你的花園整理得井井有條,可是他有令人難忘的脾氣.你為了不易找人接替,只好暫時忍氣用他.我們服事主有甚麼成就,可能也是這種情形.我們千萬不可自欺,免得我們傳福音給別人,自己反被主不認識,關在門外,哀哭切齒!
\newline
\newline
弟兄姊妹們:我們盼望甚麼?我們盼望神在萬古前所應許的永生(多一2),就是那不能朽壞,不能玷污,不能衰殘,為我存留在天上的基業.(後前一4)這基業是我們的賞賜,(西三24)也是我們的冠冕.(提後四8)但這冠冕會被人奪去,(啟三12)這福音的盼望會失去.(西一23)所以我們要努力,較力爭勝地持守,將起初確實的信心堅持到底,(來三14)要給主認識,要棄絕一切不義,凡事行而宜之,守法,聽命.這樣,我們這個信蒙神能力保守的人,必能得著所預備到末世要顯現的救恩.(彼前一5)「但願使人有盼望的神,因信,將諸般的喜樂平安充滿你們的心,使你們藉著聖靈的能力,大有盼望!」(羅十五 15)啊們!
\newline
\newline


\newpage

\section{第41屆~港九培靈會~林道亮~第十四講}
\label{sec:1493}
\textbf{天國子民的資格}
\newline
\newline
連結: \href{https://www.hkbibleconference.org/session-message/view/1493}{\texttt{https://www.hkbibleconference.org/session-message/view/1493}}
\newline
\newline
\hyperref[sec:1494]{< < < PREV SERMON < < <}
~
\hyperref[sec:toc]{[返目錄]}
~
\hyperref[sec:1492]{> > > NEXT SERMON > > >}
\newline
\newline
國家有憲法,基督徒也有天國的憲法.昨晚我們討論,要持守福音的盼望,除了神的保守外,要行義,要守法.那末我們要行甚麼義,守甚麼法呢?天國的義和法,記載在馬太福音第五至第七章裡,凡天國的子民都要履行這幾章,非天國子民是沒有份的.
\newline
\newline
這三章的目的,有些像以色列人的誡命,律例,典章.神把這些賜給以色列人,要摩西教訓他們,使他們在未來國度裡可以實行.(申五31)在基督徒未進入國度之前,主也給我們這三章,叫我們練習,使我們配在未來國度裡與祂同王.這三章內所說的,是教導我們怎樣對待弟兄姊妹,神和自己.願主恩待我們,叫我們能努力學習,配進主的國度!
\newline
\newline
今天晚上,我們只能討論這憲法的引言,或說是天國子民的資格,就是我們平常所說的八福.八福原是主為我們所預備國度裡的福份,只有夠資格的人才能享有.這不是八樣福,乃是一福的八方面.一切盼望進國度的人,都應靠主恩實踐力行!請記得:「靠著那加給我力量的,凡事都能作」.(腓四13)不要,怕只要信!
\newline
\newline
天國子民的資格
\newline
\newline
(一)虛心的人有福了.「虛心」二字,我們平時說是謙虛的意思,或說要倒空自己.其實原意並不這樣.「心」字原是「靈」字,「虛」字原是「貧窮」的意思,耶穌實際是說:「靈裡貧窮的人有福了!」(參英文本)在新約中共有二字可譯作貧窮:一是不太有錢,但還過得日子；一是窮得非討飯過不得日子.這裡的「虛」是第二個字.所以在路加福音這字有兩次譯作「討飯的」:「又有一個討飯的,名叫拉撒路」,「後來那討飯的死了.」(路十六20-22)這樣,本句應是:「靈裡非討飯過不得日子的人有福了!」靈是不吃飯的,靈怎能討飯呢?靈所要的是屬靈的福份,單是主有,也只能向主要.所以本節原意是「凡靈裡非向主求告就不敢動的人有福了!」這就是主耶穌所說:「......我沒有一件事,是憑著自己作的.」(約八28)這樣的人,雖然生活在地上,天卻已在他內裡掌權,天國(天的統治權)已是他的.
\newline
\newline
我們喜歡誇耀自己.賺了錢是自己本事,書讀得好是自己用功,兒女是自己生的,房子是自己買的,甚麼都以為靠自己才有今天.但是一個靈裡貧窮的人,深知一切都是主的恩典.有了主的恩典.他能作這樣,得那樣；沒有主的恩典,他就不敢動.他真感覺到天在他身上掌權.記得我在伯瓊斯大學神學院教書時,在談話時常常說:「靠主的恩典!」學生們以為這是我的口頭禪罷了.後來有一天,一位同學來見我說:「林博士,你常說靠主的恩典,起初我們還以為是你的口頭禪,後來才知道你說的真心如此!」我們是神的兒子,我們的一切都是神的恩典.弟兄姊妹們:你們是否正在受天國的訓練?是否凡事都會先告訴神?不論你用飯,上學,工作,買賣,甚至理髮等等,都會先告訴神,求神賜福,引導,看顧呢?這樣作不是神經病,卻是真正天國子民的訓練.我們若現在不順服天上統治權,將來在天國裡怎能掌權呢?那非求告主過不得日子的靈有福了,因為天國已是他們的了!
\newline
\newline
(二)哀慟的人有福了.不是叫我們信主的人整天悲傷哭泣.如果我們整天拉長了臉,用淚洗臉,又怎能領人歸主呢?「哀慟」二字,原意指悲哀的心境說的,如傳道書所說:「智慧人的心,在遭喪之家；愚昧人的心,在快樂之家.」(七4)
\newline
\newline
在我國的孝禮上,如果父喪,為兒子的要帶孝,不理髮,不修面,親友中有喜慶他也不去參加,這就是這裡「遭喪之家」的意思.聖經說:「智慧人的心,在遭喪之家.」這是甚麼意思呢?在屬靈上說,主耶穌時刻與我們同在；但在人方面說,我們的主已是離開我們到遠方去了,他要得國回來,我們的心對他怎樣呢?有些丈夫出門遠行,太太會寫信告訴他,她晚上睡不著覺,用飯也食不知味；別人請她出去玩,她也沒有興趣!這一切都因她親愛的丈夫不在她身邊,叫她無心享受.信的結束,她常會說:「親愛的呀,你早點回來吧!」弟兄姊妹們,我們親愛的主已到遠方去了,他仍未回來.他留我們在這裡,我們的心是否在祂那裡呢?若我們的心真的在那裡,世界的虛榮,我會無心享受；別人對名利的追逐,我會不感興趣.因為我的喜樂,享受,是在我親愛的主那裡.我現在願意孤孤單單地等候,直到我親愛的得國回來!那時他會看我如同裝飾整齊的新婦!毋怪約翰說:「主耶穌呀,我願你來!」(啟廿二20)
\newline
\newline
弟兄姊妹們:我們是天國的子民,我們的盼望不在這世上.當然,你若愛世界,愛娛樂,不愛主,你仍有這自由；但你若要在天國有份,請你記得主曾說:「學生不能高過先生,僕人不能高過主人!」(太十24)這裡所說的學生先生和主人僕人並不是普遍的說法,乃是專指信徒與耶穌之間的關係.我們在這世上是寄居的,我們的心應在遭喪之家!我們的主背十字架得國,難道我們可坐花轎進天國嗎?讓我們現在與他一同受苦,將來也與他一同作王!
\newline
\newline
(三)溫柔的人有福了.甚麼是溫柔呢?是走路時低下頭,說話很細聲嗎?是「未罵先笑,未打先倒」嗎?這些不是聖經裡的溫柔呀!馬太記載了耶穌溫柔的榜樣說:「要對錫安居民說,看哪,你的王來到你這裡,是溫柔的,又騎著驢,就是騎著驢駒子.」主耶穌是王,他可大張威勢的進耶路撒冷；但他有那當得,寧不得；有那權利,寧放棄,甘願騎著驢駒進城,這就是溫柔.溫柔簡單的意義,就是肯為主的緣故,把當得的寧放棄而不得.創世記十三章告訴我們亞伯拉罕的溫柔.依東方人的規矩,作叔父的當有優先權揀選牧場,但亞伯拉罕甘願讓姪兒羅得先揀,並且說:「你向左,我就向右；你向右,我就向左!」這是他的溫柔.結果,神對他說:「凡你所看見的一切地,我都要賜給你和你的後裔,直到永遠.」(創十三15)溫柔的人必承受地土!
\newline
\newline
我幼時唸過一課書叫「大樹將軍」.說到中國古時的一位將軍,在每次戰爭中,他總是身先士卒,但在論功行賞時,他總是到大樹下站著,讓別人去得獎賞,自己情願不要,這是溫柔!這位將軍,比較今天有許多教會領袖,有風頭可出的時候,爭先恐後地站在前面,等到要盡義務的時候,便退避三舍,真不可同日而語!今天在教會中,有不少弟兄姊妹歡喜掛名當甚麼長,甚麼主任,甚麼委員,可是他們的聖經知識,連舊約幾卷,新約幾卷都不知道,教會工作十次中九次缺席,還美其名為老資格!(抑老腐敗呢?)一個國家的腐敗,在於愛名利的人太多,盡責任的人太少.我常對美國人說,美國已大非昔比了!三十年前,我第一次去西方,看見美國政府人員辦事的忠勤和快捷,令我五體投地.我心裡說:「他們國家的富強是理所當然的啊!」那時,我回觀自己政府的官員,又懶又貪,只知道奉承,不知道責任!例如我申請出國護照,半年還未發下來.等我託長官查問,他們連我的申請書都找不到!以後總算找到了,不久護照也發下來了.如我不托長官去查問,我相信他們永遠也不會發給我!可是美國近來也漸漸走下坡,若不是主繼續施恩,美國的國運恐也不會太長吧!不特邦國如此,今天教會也是一樣.所以除非主大興溫柔的人才,教會的前途非常悲觀!
\newline
\newline
願主憐憫我們,叫我們看穿近處的虛榮,看清主國度的榮耀和要求,讓我們為主多作從來沒有人知道的工作,學習溫柔的功課,去承受地土──天上的基業.
\newline
\newline
(四)飢渴慕義的人有福了.昨晚我已說過,「義」字有二個意義:在舊約和四福音裡的意義是「行而宜之」,在保羅書信的意義是「宣告無罪」.「行而宜之」是對已得救的人說的.我們應該渴望行得合宜,如同飢餓的人渴慕食物.這樣的人,求就得著,尋必尋見,叩門必得開門,因為他們渴望遵行神的旨意.
\newline
\newline
弟兄姊妹們:請記得我們研讀聖經不是為得知識,乃是為明白神的旨意,識得是非之分,知道怎樣去行而宜之.昨晚我們討論過,在國度裡,主只認識那些「行而宜之」的信徒,可見義對國度的重要!我們要渴慕義,學習義,追求義,因為惟有義要居在新天新地之中.(彼後三13)今晚有許多兄姐來自很遠的地方,我相信你們是有飢渴慕義的心才來的.但同時我也相信,在你們中間,有些人對於赴培靈會已成為一種習慣!他們為了安慰自己良心,或想得些功德來赴會,對於真理有否聽進去,根本不在乎.這樣沒有飢渴慕義心的人,無論他赴多少會,不特得不到造就,反會越聽,耳聽越發沉,眼睛越昏迷!
\newline
\newline
馬可福音第四章是討論「聽」的問題.我們平時總是注重講道的講得好不好,但這章聖經是說聽道的聽得好不好!第二十四節是本章的中心,耶穌說:「你們所聽的要留心,你們用甚麼量器(來)量,(給人二字未見於原文)也必用甚麼量器量給你們......」(24)這裡告訴我們,神量給我們的福,是根據我們帶到神前的量器.我記得在上海有一位愛主的姊妹,她每次崇拜之前會祈禱,求主讓祂的話進入她心裡.後來她自己作證說,她沒有一次得不到主的光照的!我聽了以後,很受感動.毋怪那姊妹如同鴿子一般,令人一見就得造就.她每次到主前都存著飢渴慕義的心,主是信實的,當然祂不會使她失望.總之,你帶了甚麼量器來量,主也必照甚麼量器量給你.你若帶著批評的量器來,你會得著煩惱回去；帶著無所謂的量器來,空空如也回去；帶著安慰良心的量器來,良心得些安慰回去.照樣,你若帶了虛心領教的量器來,你必會滿載而歸.
\newline
\newline
(五)憐恤人的人有福了.憐恤人是信徒必須的品格.我們信仰純正的信徒,常批評那些傳社會福音的人,只注重救濟,不注重救恩；只顧念人的貧窮,不理人屬靈的需要.他們這種作風,當然是捨本逐末,前後倒置!但我們往往太過忽略了憐恤人的重要.有些兄弟姊妹們很愛教會,教會要錢建造房屋,他會奉獻幾千幾萬元,若有貧苦的弟兄需要救濟時,他連幾塊錢也不肯施予,這是沒有憐恤啊!主說:「我喜愛良善,不喜愛祭祀；喜愛認識上帝,勝於燔祭.」(何六6)「良善」小字「或作憐恤」.根據新約主的引用,這裡應譯作憐恤(太九13).主不喜愛我們的祭祀,祂要我們憐恤貧窮的兄姐.教會裡若有窮苦的兄姐們,我們要設法幫助,不應該讓他們到外面去找幫助.如果張家有人需要錢,張家的人不肯幫忙自己的人,反要他向李家求助,這是張家的羞恥!照樣,若主的兒女中有需要,主家中的人不肯幫助,反要他向外邦人求助,這也是主的羞恥,經上說:「人若不看顧親屬,就是背了真道,比不信的人還不好；不看顧自己家裡的人,更是如此.」(提前五8)
\newline
\newline
為了時間關係,今晚只能到此為止,其餘,若主今晚不回來,留待明晚討論,以上都是天國子民的資格,凡願意在主國度裡有份的,希望不把這些資格當作耳邊風才好.主說:「我告訴你們:你們的義,若不勝於文士和法利賽人的義,斷不能進天國!」(太五20)
\newline
\newline


\newpage

\section{第41屆~港九培靈會~林道亮~第十五講}
\label{sec:1492}
\textbf{天國子民的福份}
\newline
\newline
連結: \href{https://www.hkbibleconference.org/session-message/view/1492}{\texttt{https://www.hkbibleconference.org/session-message/view/1492}}
\newline
\newline
\hyperref[sec:1493]{< < < PREV SERMON < < <}
~
\hyperref[sec:toc]{[返目錄]}
~
\hyperref[sec:1491]{> > > NEXT SERMON > > >}
\newline
\newline
我們已作了天國的子民,一定要履行天國的憲法,叫我們有屬天的資格,日後可以和基督同王.昨天晚上,我們已經討論了五個資格；今天晚上,我們由第六資格開始.
\newline
\newline
(六)清心的人有福了.若我們要知道清心是甚麼意思,一定要先知道這裡的「心」字是何所指.「心」在聖經中所指的甚多；有時指思想,有時指情感,有時指意志,有時指良心,也有時指我們一切的精神作用.所以研究聖經中「心」字的時候,我們要特別留意,一定要細看上下文,有時甚至要新舊約對照,才能清楚.要知道這裡清心的意義,我們一定要找其他的經文幫忙.約翰告訴我們說；「我們的心若責備我們,神比我們的心大,一切事沒有不知道的.親愛的弟兄阿,我們的心若不責備我們,就可以向神坦然無懼了.」(約壹三20-21)甚麼人不敢見神呢?愛心責備的人.所以這裡說:「我們的心若不責備我們,就可以向神坦然無懼了.」神是無所不在的,惟有當人心有內疚時,人心懼怕神,才見不到神!一個孩子作了錯事,心有內疚,就怕見爸爸的面.如果他毫無自責,心中坦然,那裡會怕見爸爸!照樣,若我們沒有錯處,良心不自責,我們就能見神了.所以「心」在這裡是指良心的功用說的.如果一個人的良心有了問題,他禱告會不通,聽道會乏味,對傳道人會討厭,對愛主的弟兄姊妹會避見,那裡還敢見神啊!毋怪作教會領袖的,一定「要存清潔的良心」.(提前三9)
\newline
\newline
這樣,清心是指清潔的良心說的.如果我們的良心不清潔,我們就不能坦然無懼地見神.謝謝主,在我們得救時,我們的良心已被主的寶血洗淨(來九 14)；我們承認了罪,我們的良心也不再控告我們了.(來十2)從那時起,良心的功用恢復,一有不義的地方,良心就會責備我們.相信我這樣說的時候,在座有人的良心正在那裡責備你,說你有那些事是不應該作的!每次當良心斥責我們的時候,我們便不能和主面對面了!我們一定要再一次求主寶血洗淨,千萬不要隱藏自己的罪,或說讓它去!凡未到主面前認清的一切罪,只有兩條路:一是把它認清,一是用意志把罪壓下去.(普通人所說的忘記它)凡壓下去的罪,永遠是潛伏著,我們不特要消耗許多精力去抑壓它,它也會令我們生病!各位,千萬別把罪抑壓住,要把罪交給耶穌,祂會潔淨你,使你成為清心的人.願主憐憫我們,叫我們都能像保羅所說的:「弟兄們:我在神面前行事為人,都是憑著純善的(原文有此)良心,直到今日.」(徒廿三1)
\newline
\newline
(七)使人和睦的人有福了.我們知道,基督是宇宙中的和事佬.祂不特藉著十字架,拆毀了中間的牆,將以色列人和外邦人合而為一,成就了和睦；且藉著十字架,使兩下歸為一體,與神和好,(弗二14-16)連著萬有,無論是地上的,天上的,都與神和好了!(西一20)所以祂也要我們去作和事佬.我們有時看見教會中弟兄姊妹發生爭執,不特沒有盡力去勸解,反會存著幸災樂禍的心去批評,或甚至鼓動!這樣的人,不能稱為神的兒子,不要希望主認識他,也不要希望在天國裡有份!
\newline
\newline
我們勸人信耶穌,是叫人與神和好,免得神的震怒降在他們身上；但我們往往忽略了弟兄姊妹間的和睦.當他們中間發生意見的時候,我們會抱著「各人自掃門前雪,休管他人瓦上霜」的態度.各位,若你不想「選上」,你可以有你的自由；如你想要在主的國度裡有份,那你就沒有這自由,你一定要負起和事佬的責任.若弟兄中有二人發生誤會,你能不能請他們一同到家裡來,拉著他們的手,請他們看在主的面上和你面上彼此饒恕呢?在真摯的愛裡,我相信鐵石的心腸也會軟化了!
\newline
\newline
我少年時遭遇了一件永不忘記的事.我有一個淘氣的弟弟,有一天,他被媽媽斥責了幾句以後,心中很不服氣,就用腳拼命地踢著一扇門出氣!我覺得他這種態度很不應該,就自持著大哥的身份,去賞了他幾個耳光.事後,我的良心有些不安.為要安慰我的良心,我便去把這事告訴了媽媽.我想:她就算不稱讚我,也不至於會責備我吧!她聽了後,半晌不出聲,然後流著淚對我說:「亮呀,兄弟和同看娘面,千朵桃花一樹開!」從此以後,我們兩兄弟非常相愛.弟兄姊妹們,請記得:教會裡的紛爭,仇恨,也會叫在天上的慈父傷心流淚啊!願主幫助我們,叫我們肯負起和事佬的責任來!如果我們當中彼此有仇恨的,應該彼此饒恕,如主饒恕我們一樣.我們在世都是客人,今天不知道明天事,何必要賭氣呢?
\newline
\newline
(八)為義受逼迫的人有福了.為真理受逼迫,或為行在神眼中看為正的事受毀謗,是不能避免的!因為經上記著說:「凡立志在基督耶穌裡敬虔度日的,都要受逼迫!」(提後三2)這樣,除非我們不敬虔度日,否則,逼迫是在那裡等候我們!這種逼迫,不特是敬虔的標誌,也是向主是否盡忠的試金石.凡經得起試煉的,不特天國是他們的,且有大賞!(彼前一7)我們的主為要強調這資格,再用兩節經文來敘述它.(十一,十二節)所以基督徒為主名受逼迫的毀謗,甚至殉道,倒要歡喜；不特歡喜,也要快樂,因為神為他預備了先知的賞賜!綜之,為主名受逼迫或殉道是神的高級恩典,只有那些配得的才能享有.所以凡為主名受逼迫,甚至殉道者,當自引以為榮,因為這證明你們是選上的.「你們是與基督一同受苦,使你們在祂榮耀顯現的時候,也可以歡喜快樂!」(彼前四13)
\newline
\newline
我再說,弟兄姊妹們:為主受苦,甚至為主殉道,是神特殊的恩典!司提反是個好例子,教會歷史中許多殉道的偉人,也都可以為此作見證.希望凡在火煉中度日的兄姐們能剛強壯膽,知道這至暫至輕的苦楚,要為你們成就極重無比永遠的榮耀!(林後四17)基督升天後,坐在神右邊,但當殉道者回天家時,祂會站起來歡迎,這是何等的「稱讚,榮耀和尊重」啊!(徒七55)
\newline
\newline
天國子民的福份
\newline
\newline
我們都知道,本段經文普通稱為八福篇,因為除了末一資格用過兩次福字外,其餘每一資格都有一個福字.雖然福字用了九次,福的說明卻只有七種,因為第一和第八福同是「天國是他們的」.所以謂其說八福,毋寧說七福更準確些.其實這七福,無非是天國福份七方面的描寫而已.因為在神眼中,惟有得進天國的人才算是真有福的!天國好比一株樹,七福是樹上的果子.現在讓我們嚐一嚐這些果子的滋味!
\newline
\newline
(一)天國是他們的.我們都知道,當我們信主的時候,我們便與基督一同活過來,一同坐在天上.(弗二5)但我們能否與基督同王,要看我們能否讓天國的權柄,在我們生活中施行；能否在一切逼迫患難中,打那美好的仗.如果我們能的話,那末主說,天國已是我們的.我們用不著擔心我們能否得那不能朽壞,不能玷污,不能衰殘,為我們存留在天上的基業.(彼前一4)相反地,我們相信:「我們這因信蒙神能力保守的人,必能得著所預備,到末世要顯現的救恩.」(彼前一5)或說「永不衰殘的榮耀冠冕」.(彼前五4)
\newline
\newline
(二)他們必得安慰.在世為主過「遭喪之家」生活的人,他們所得的安慰是永世的,正如保羅說:「但願我們主耶穌基督,和那愛我們,開恩將永遠的安慰,並美好的盼望,賜給我們的父神.」(帖後二16)這永遠的安慰和美好的盼望,當然都是指著天國說的.因為惟有在天國裡的一切,是永遠的,也是美好的.這安慰和盼望是要我們努力去得,但同時我們要牢記:天國的一切,仍是由於父神為誰預備,就賜給誰!(太二十23)我們的態度是:「作完了一切所吩咐的,只當說:我們是無用的僕人,所作的本是我們應份作的.」(路十七10)
\newline
\newline
(三)他們必承受地土.當我年青的時候,每次讀到這裡,總是百思不得其解!「地土」到底何所指呢?謝謝主,一天讀到以賽亞書說:「你哀求的時候,讓你所聚集的拯救你罷；風要把他們颳散,一口氣要把他們都吹去；但那投靠我的必得地土,必承受我的聖山為業.」(賽十七13)我才知道,這地土原指神行政中心說的.因為本節的「必得地土」,等於「必承受神的聖山為業」.神的聖山既在舊約常指神行政中心,所以「得地土」就是「承受神行政的中心」.今天肯為主當得寧不得的,日後要在主行政中心掌權.
\newline
\newline
(四)他們必得飽足.凡是飢渴慕義的人,在今世一定得不到飽足.因為越是渴慕行而宜之的人,越會感覺到,作惡的和迷惑人的,越久越惡.不特社會上不法的事天天加增,就是教會也一過一天離經叛道.義人住在其中,看見聽見他們不法的事,他的義心就天天傷痛,那裡還會飽足!感謝主,祂有應許:「但我們照他的應許,盼望新天新地,有義居在其中.」(彼後三13)主應許我們將有新天新地,有義長居其中.這新天新地將充滿了「行而宜之」!甚麼都以神的尺為準則,甚麼都以神的旨意為依歸,毫與今天的不法,不德,不義.義人在其中將豐富地享受一切他所渴慕的.換句話說,惟有在天國裡,義人才能飽足!
\newline
\newline
(五)他們必蒙憐恤.他們將蒙憐恤呢?主耶穌曾告訴我們,那些憐恤人的,就是見弟兄饑了給他吃,渴了給他喝,作客旅留他住,赤身露體給他穿,病了或坐監去看他的人,他們是蒙天父賜福的.他們的蒙憐恤是承受天國,正如馬太所記:「於是王要向那右邊的說,你們這蒙我父賜福的,可來承受那創世以來為你們所豫備的國.」(太廿五34)這個不是臨時計劃的,也不是猶太人拒絕基督後才補充的,乃是神在創世以來預備給義人的,就是憐恤人的人.
\newline
\newline
(六)他們必得見神,我們信主的人,一定已會見了神,否則我們就沒法信耶穌!那末這裡所說,清心的人必得見神,到底是甚麼意思呢?在新約裡,「見神」可有二個意思:一是今世的,一是在國度裡的.今世的,覺神同在；在國度裡的,實際與神同在,正如啟示錄所記:「他的僕人都要事奉他,也要見他的面；他的名字必寫在他們的額上......他們要作王,直到永永遠遠.」(啟廿二4)在國度裡,神要親與他們同在,作他們的神.惟有今天清心的人就是行事為人憑著清潔良心的信徒,在國度裡才能與基督同王,見祂的面,永遠屬祂；這不特在千禧年時,也在永世裡,因為作王是永永遠遠的!
\newline
\newline
(七)他們必稱為神的兒子.我們信主後已經是神的兒子了,為甚麼這裡又說必稱為神的兒子?這問題,主耶穌已經代替我們答覆.祂說:「得勝的,必承受這些為業,我要作他的上帝,他要作我的兒子.」(啟廿一7)這樣,這裡所說的兒子,是指那些受國度為業的得勝者.我們不是信了主,有了主的生命就夠了,我們還要努力去承受天上的基業.我們能否有資格去承受它,要看我們在世上的造就如何.若我們今天肯效法我們的主作和事佬,我們就會作祂的承繼的兒子.總之,得神生命的兒子和承神基業的兒子是二,不是一!
\newline
\newline
以上七福,不過是描寫天國福份的各方面,總括起來,只有一個福,就是與基督同王,承受神在創世以前為我們所預備的基業,實現神在萬古前所應許的永生.弟兄姊妹們:我再說,得救是白白的,進神的國度是要資格的!這資格是訓練成的,經過訓練後,有人及格,也有人不合格!所以經上說:「被召的人多,選上的人少!」
\newline
\newline
弟兄姊妹們:我相信你們大多數已經得救,希望你們不叫主失望,努力進窄門,跑小路,接受天的統治,過「遭喪之家」的生活,凡事重義務輕權利,渴慕神旨得成,向眾人行善,(向信徒一家的人更當這樣)存清潔的良心,作神和人,人和人之間的和事佬,以歡喜快樂的心忍受逼迫,甚至殉道.賜平安的神,必親自使你們全然成聖,保守你們的靈,與魂,與身子,在我主耶穌基督降臨的時候,完全無可指摘!
\newline
\newline
各位,做人難,做基督徒更難,要與主同王更難上加難!但不要怕,只要信,因主的恩典總是夠我們用的.我們只要努力向標桿直跑,祂會扶持我們.我們不敢靠自己,但我們要盡自己所能的.相信那位在我們心裡動了善工的,必會完成祂的工作在我們身上,因為他是信實的,祂的信實在到萬代!哈利路亞!
\newline
\newline


\newpage

\section{第41屆~港九培靈會~林道亮~第十六講}
\label{sec:1491}
\textbf{主的新命令}
\newline
\newline
連結: \href{https://www.hkbibleconference.org/session-message/view/1491}{\texttt{https://www.hkbibleconference.org/session-message/view/1491}}
\newline
\newline
\hyperref[sec:1492]{< < < PREV SERMON < < <}
~
\hyperref[sec:toc]{[返目錄]}
~
\hyperref[sec:1324]{> > > NEXT SERMON > > >}
\newline
\newline
今天晚上,是這次培靈會最後一個晚上,我要和各位討論愛的問題.在這末世,人心不古,世態炎涼的社會中,要找純真的要是何等難啊!不特年青人的愛,是自欺自迷；連夫妻間,只存著利益的動機,而沒有真摯的自我犧牲!教會在這邪惡的世代中,許多稱為信徒的,也只知專愛自己,貪愛錢財,不愛父母,愛娛樂不愛神.他們常為聖經中一個名詞,一種措辭；宗派中的某種成規,某種儀式,唇槍舌劍地地鬧得天翻地覆,卻把基督徒的精華──彼此相愛──完全忘記.這正如經上所記:「但命令的總歸就是愛,這愛是從清潔的心和無虧的良心,無偽的信心,生出來的.有人偏離這些,反去講虛浮的話,想要作教法師,卻不明白自己所講說的,所論定的.」(提前一5-7)這些辯論和爭執,大都無關真理的原則,也不能發明神在信上所立的章程,反把出自清潔的心,和無虧的良心,無偽的信心的愛偏離了!這原因,除了社會的影響以外,也因教會中加多了講虛浮話的人.這些人要作教法師,卻又沒有真材實料.有些稱為聖工人員的,自己還沒有重生,怎能明白自己所講的呢?自己沒有遇見神,怎能叫人遇見神呢?可是作了教法師,總要有些教材才是,否則怎能騙得到麵包呢?所以只好採些世俗的虛談,和那些敵真道似是而非的說話來自欺欺人,造成教會烏煙瘴氣的局面!
\newline
\newline
我們都知道,「彼此相愛」是基督的遺囑,也是我主的臨別贈言,它的重要性不言可知.祂說:「我賜給你們一條新命令,乃是叫你們彼此相愛；我怎樣愛你們,你們也要怎樣相愛.你們若有彼此相愛的心,眾人因此就認出你們是我的門徒了.」(約十三34-35)這是一條新命令!全部聖經有許多神的命令,惟有這一條,我們的主特別指出是「新」命令；並且是永遠常新的.(約壹二7-8)為甚麼這命令是新的呢?這命令新到甚麼程度?為甚麼要這新命令?為要回答這些問題,讓我把這新命令分成以下四點來討論.
\newline
\newline
命令的來源是新的──「我賜給你們一條新命令」.
\newline
\newline
我們都知道,全部聖經可以用「愛」字來總括,說到神怎樣愛人,人也應該怎樣愛神,而又彼此相愛.所以彼此相愛本是舊命令,為甚麼主說是一條新命令呢?這回答,我們可以在這「新」字裡找到.
\newline
\newline
在新約聖經共有四字譯作「新」字.三字用在耶穌回答約翰門徒的詰問,一字說耶穌為我們了一條又新又活的路(來十20).當約翰門徒詰問耶穌,為甚麼祂的門徒不禁食?祂回答說:「沒有人把新布補在舊衣服上,因為所補上的,反帶壞了那衣服,破的就更大了.也沒有人把新酒裝在舊皮袋裡:若是這樣,皮袋就裂開,酒漏出來,連皮袋也壞了；惟獨把新酒裝在新皮袋裡,兩樣就都保全了.」(太九16-17)這裡新布的「新」字,原意指「未經磨光細」,或說「未漂過」,所以「新布」實指「生布」.新酒和新皮袋各用了不同的「新」字,因為新酒和新皮袋的「新」是不同的.新酒和舊酒雖然是時代不同.質料卻相同；新皮袋與舊皮袋卻是質料不同.舊的皮袋漸漸腐蝕陳舊,不能再用,必須換新的,本命令的「新」字,是和新皮袋的「新」同字的.這是告訴我們說:以前的命令已經陳舊不合用了,現在需要一個新命令,是裡面到外邊完全新的,並不是由舊皮袋改成的.
\newline
\newline
不特內容和外觀是新的,來源更是新的.本節經文說:「我賜給你們一條新命令.」這命令是「我」賜給的.這命令的本源是在那位為我們道成肉身,死而復活的主裡面.它不是由良心產生,不是由社會的需要構成,也不是從舊約發展出來的,乃是我們的主新頌佈的命令.我們都知道；凡不守主命令的,就是罪人!請注意:彼此相愛既是主的命令,凡不彼此相愛的,也就是背逆之子!我們知道「不信耶穌」就是罪,卻不知「不彼此相愛」也是罪!這是祂的命令,所以我們必須遵守,沒有選擇的餘地.換句話說,就是我們不想彼此相愛,也要彼此相愛,除非我們甘願被主關在天國門外!同時,這既是祂的命令,所以我們彼此相愛的動機,不是因為弟兄們的可愛,也不是為了滿足個人的需要,乃是為了主的緣故!
\newline
\newline
命令的範圍是新的──「乃是叫你們彼此相愛」.
\newline
\newline
當耶穌在世時,祂為了時代的需要,訂正了舊約十誡為二誡:「你要盡心,盡性,盡意,愛主你的神......其次也相倣,就是要愛人如己.」(太廿二 37-39)但第二條範圍仍太廣大,有時叫人無從開始；因為要逢人都愛如己,在實際上不可能性很大!耶穌知道門徒的軟弱,祂要他們更實際,就把範圍縮小,由愛人如己成為「你們彼此相愛」.因為祂知道:除非弟兄姊妹們能夠先彼此相愛,便永遠也愛不起眾人!凡事由小至大,由點成線,愛人也是如此.毋怪彼得說:「有了虔敬,又要加上愛弟兄的心,有了愛弟兄的心,又要加上愛眾人的心.」(彼後一7)
\newline
\newline
我們都知道,愛朋友和愛自己人不同.自己人,無論如何總是自己人.自己孩子,就算不聽話,作父母的還是他的父母,但外人就不同了.主要我們對弟兄姊妹們的彼此相愛,如同自己人一樣.為此,聖經稱教會為「神的家」,(提三15)信徒為「信徒一家的人」,(加六10)或「神家裡的人」.(弗二19)既然是自己人,有甚麼不可直言無諱,或不能體貼饒恕呢?又為甚麼不情願受欺,不情願吃虧呢?請聽主的警告:「你們不饒恕人的過犯,你們的天父也必不饒恕你們的過犯!」(太五15)相反地,經上應許說:「我們若彼此相愛,神就住在我們裡面,愛祂的心在我們裡面得以完全了.」(約壹四12)
\newline
\newline
末世的時候,不少信徒確注意「相愛」二字,卻忽略了「彼此」.他們常埋怨教會不關懷他們,弟兄姊妹沒有愛心!可是他們自己對教會呢?連「拔一毛而利天下」也不肯,一舉手之勞的幫忙也不願!相反地,教會中有利益時,他們會爭先恐後；有義務時,卻袖手旁觀,反說別人沒有思想!主知道我們的詭詐,所以祂沒有叫我們只要有愛心,乃是叫我們彼此相愛.換句話說,祂告訴我們；愛是雙方的,不是單程的；愛是傳染的,不是隔離的.所以,如果你覺得教會沒有愛心,請你自問,你為弟兄姊妹出過多少力,化過多少錢,費過多少時,幫過多少忙,祈過多少禱,甚至流過多少淚!若說沒有多少,那麼被埋怨的不是教會,乃是你自己!請注意神的話:「不要自欺,神是輕慢不得的；人種的是甚麼,收的也是甚麼!」(加六7)願神憐憫,叫我們在彼此相愛上,能常以為虧欠,不讓我們利用「相愛」二字,攻他人,或發展自我!
\newline
\newline
話雖如此,要實際履行彼此相愛,確非容易!特別在這末世的時候,不法的事增多,許多人的愛心漸漸冷淡了,(太廿四12)好人總是吃虧,好事常是上當!例如在美國,以前搭順風車很容易,因為大都車主樂意幫助無車階級.近來因為有些不法之徒,竟趁這機會行劫,大多車主就不再給人方便了.不但社會上不法,教會中亦有許多不法之徒,忘恩負義,賣主賣友.無論你怎樣用愛心多方幫助他,培植他,他不特不感激,反會轉過來咬你一口!為這緣故,許多人的愛心漸漸冷談了,但請記得:彼此相愛是主的命令,有一天我們要向主交賬.人會忘記我們,虧待我們,主卻永不忘記我們.經上記著說:「我們行善不可喪志；若不灰心,到了時候,就要收成.」(加六9)
\newline
\newline
命令的要求是新的──「我怎樣愛你們,你們也要怎樣相愛」.
\newline
\newline
上文我們已討論過,為甚麼我們要彼此相愛,因為這是主的命令,也因為弟兄姊妹是自己人.在這裡,我們再加一件,我相愛是因為主愛我們.祂怎樣愛我們呢?以賽亞告訴我們,祂為我們被藐視,被人厭棄,多受痛苦,常經憂患；為我們擔當憂患,背負痛苦；為我們的過犯受害,罪孽壓傷；為我們的平安受刑罰,醫治受鞭傷......因為祂這樣愛我們,所以祂也要求我們要這樣彼此相愛,為弟兄作這個那個,因為愛是有動作的.
\newline
\newline
動作是要緊,可是只有動作而沒有動機,那動作會成為假冒為善.我們要效法主愛的動作,更要明瞭我們的存心.我主所最愛的門徒告訴我們說:「主為我們捨命,我們從此就知道何為愛；我們也當為弟兄捨命.」(約壹三16)在香港那一邊,我已經提過,這節經文直譯應是:「主為我們放下魂整個來說,是指人的自己；分開來說,是指人內裡,情感,思想,意志的效用,一生為他人放下自己的情感.祂常為了他人的軟弱和需要,作祂自己所不想作的事情.例如在他休息時講道給五千人聽；叫彼得去海邊釣魚,得銀付稅；特別在客西馬尼亞園的祈禱,和在十字架為仇敵求赦免,都是給我們榜樣和要求,要我們放下自己的喜好,願意......去愛神家裡的人.祂也常為他人放下自己的思想,去思想他人所思想的.例如祂讓彼得在水面行走；特別是在十字架上萬分痛苦時,想起祂母親的未來,都是叫我們知道,應該放下自己的思想,計劃,意見來愛弟兄姊妹們.祂的意志也是如此,當一個迦南婦人哀求時,祂就例外地幫助她；在祂復活後,祂還要從以馬杵斯往前行,卻給革流巴和他的朋友強留住,這一切都包含在「我怎樣愛你們,你們也要怎樣想愛」的要求裡.祂要我們放下自己的定規,志趣,目的,來愛信徒一家的人.總之,祂要我們彼此相愛到祂愛我們的程度.願主赦免我們,也幫助我們,叫我們能滿足祂的要求.請記得:在祂沒有難成的事!
\newline
\newline
命令的後果是新的──「眾人因此就認出你們是我的門徒了」.
\newline
\newline
初期的教會,常給外邦人稱譽說:「看哪,他們是何等彼此相愛啊!」相反地,今天的教會常給人批評說:「看哪,信耶穌的和不信耶穌的有甚麼不同啊!」彼此相愛是耶穌門徒的標誌,沒有那個,怎能和不信主的有分別呢!特別在末世的時候,信徒若沒有相愛,就沒有見證；沒有相愛,就沒法領人歸主.
\newline
\newline
在啟示錄第二,三章裡,主只對一個教會說:「看哪,我在你面前給你一個敞開的門,是無人能關的.」(啟三8)主給那一個教會開傳道的門呢?非拉鐵非教會.「非拉」就是愛的意思,「鐵非」就是弟兄.換句話說,主只給兄弟相愛的教,開無人能關的門,這樣,要廣傳福音有功效,教會一定要先彼此相愛.不然的話,就是走遍洋海陸地,領了一個人到教會,難保他不會因教會的炎涼而離棄主!
\newline
\newline
許多人以為愛只是內心的,不知道愛是有行動的.沒有行動的愛,是不完全的愛.因為純真的愛是不能坐守家園的,它也不能被人佔為私有.它會像光,不住地發揮它的光熱.好像意大利格言說:「愛比千把火炬更熱!」純真的愛既然要發揮,一發揮就會產生後果,這是很自然地被吸引到主面前來,正如約翰所說:「彼來沒有人見過神；我們若彼此相愛,神就住在我們裡面.」(約壹四12)
\newline
\newline
願主憐憫我們,叫我們明白這新命令的重要性.我們彼此相愛時,不特是因為主先愛我們,也是因為我們是祂家裡的人.祂為我們放下自己的情感,思想,意志,我們也當為愛弟兄放下這一切.我們若能這樣行,他人就會看見基督的美,來歸向基督,使主的名得著榮耀.願一切榮耀,權能都歸給坐寶座的和羔羊,直到永永遠遠,啊們!
\newline
\newline


\newpage

\section{第41屆~港九培靈會~楊濬哲~序}
\label{sec:1324}
\textbf{序}
\newline
\newline
連結: \href{https://www.hkbibleconference.org/session-message/view/1324}{\texttt{https://www.hkbibleconference.org/session-message/view/1324}}
\newline
\newline
\hyperref[sec:1491]{< < < PREV SERMON < < <}
~
\hyperref[sec:toc]{[返目錄]}
~
\hyperref[sec:1478]{> > > NEXT SERMON > > >}
\newline
\newline
一九六九年香港人心安定,市面恢復繁榮,教會工作也照常發展.本會工作蒙主恩佑,得以向前邁上一步；這應當深深感謝主恩者.
\newline
\newline
今年講員,主講奮興會者,為林道亮牧師.林牧師旅美多年,為世界有名原文專家,對於真理造就深湛,蒙林牧師不辭勞苦,遠道西來,所帶給聽眾的信息,深入淺出,會眾得益良多.
\newline
\newline
主講查經會者,為台灣協同會孫維廉牧師.孫牧師早歲在西北一帶,多年為主辛勞.迨河山變色,始移旌寶島,雖已花甲高齡,但忠心為主,數十年如一日.此次主講以弗所書及啟示錄,正是合時信息,聽眾深為感動.
\newline
\newline
主講培靈會者為本港胡恩德先生.胡先生擔任本會講員已非一次.今年信息特別取自列王記,從歷史中認識神的救恩和公義.
\newline
\newline
我們感謝神,給本會預備這幾位講員,不但信仰純正,靈歷豐富；且大有能力,帶給聽眾極大的祝福.
\newline
\newline
本會記錄員,歷年來冒炎暑與講員同工作,至深欽佩.惟今年因各位工作過忙,無法分身,乃由戎子由,譚雅芳,吳天樂諸弟妹,分擔謄記.感謝上帝,在祂的家中親自興起人來,學習事奉.
\newline
\newline
願主親自賜福帶領本會當前的工作.
\newline
\newline


\newpage

\section{第41屆~港九培靈會~胡恩德~第一講}
\label{sec:1478}
\textbf{失敗與希望}
\newline
\newline
連結: \href{https://www.hkbibleconference.org/session-message/view/1478}{\texttt{https://www.hkbibleconference.org/session-message/view/1478}}
\newline
\newline
\hyperref[sec:1324]{< < < PREV SERMON < < <}
~
\hyperref[sec:toc]{[返目錄]}
~
\hyperref[sec:1479]{> > > NEXT SERMON > > >}
\newline
\newline
列王記上 1:1-22:53
\newline
\begin{longtable}{cl}
\hline
\hline
章節 & 經文 (和合本修訂版)\\
\hline
1:1 & \begin{tabularx}{0.7\textwidth}{X} 大衛王年紀老邁，雖然蓋著外袍，仍不夠暖和。 \end{tabularx} \\ \\ \relax
1:2 & \begin{tabularx}{0.7\textwidth}{X} 臣僕對他說：「不如為我主我王找一個年輕的少女，侍立在王面前，照顧王，睡在王的懷中，好使我主我王得暖。」 \end{tabularx} \\ \\ \relax
1:3 & \begin{tabularx}{0.7\textwidth}{X} 於是他們在以色列全境尋找美貌的少女，找到了一個書念女子亞比煞，帶到王那裡。 \end{tabularx} \\ \\ \relax
1:4 & \begin{tabularx}{0.7\textwidth}{X} 這少女極其美貌，她照顧王，伺候王，王卻沒有與她親近。 \end{tabularx} \\ \\ \relax
1:5 & \begin{tabularx}{0.7\textwidth}{X} 那時，哈及的兒子亞多尼雅妄自尊大，說：「我要作王」，就為自己預備座車、騎兵，又派五十人在他前頭奔跑。 \end{tabularx} \\ \\ \relax
1:6 & \begin{tabularx}{0.7\textwidth}{X} 他父親從來沒有責怪他，說：「你為何這麼做？」他非常俊美，生在押沙龍之後。 \end{tabularx} \\ \\ \relax
1:7 & \begin{tabularx}{0.7\textwidth}{X} 亞多尼雅與洗魯雅的兒子約押和亞比亞他祭司商議；他們就順從亞多尼雅，幫助他。 \end{tabularx} \\ \\ \relax
1:8 & \begin{tabularx}{0.7\textwidth}{X} 但撒督祭司、耶何耶大的兒子比拿雅、拿單先知、示每、利以，以及大衛自己的勇士都不順從亞多尼雅。 \end{tabularx} \\ \\ \relax
1:9 & \begin{tabularx}{0.7\textwidth}{X} 亞多尼雅在隱‧羅結旁瑣希列磐石那裡獻牛羊、肥犢為祭，請了他的眾兄弟，就是王的眾兒子，以及所有作王臣僕的猶大人。 \end{tabularx} \\ \\ \relax
1:10 & \begin{tabularx}{0.7\textwidth}{X} 但他沒有邀請拿單先知、比拿雅和勇士們，以及他的弟弟所羅門。 \end{tabularx} \\ \\ \relax
1:11 & \begin{tabularx}{0.7\textwidth}{X} 拿單對所羅門的母親拔示巴說：「哈及的兒子亞多尼雅作王了，你沒有聽見嗎？我們的主大衛卻不知道。 \end{tabularx} \\ \\ \relax
1:12 & \begin{tabularx}{0.7\textwidth}{X} 現在，來，我給你出個主意，好保全你和你兒子所羅門的性命。 \end{tabularx} \\ \\ \relax
1:13 & \begin{tabularx}{0.7\textwidth}{X} 你去，進到大衛王那裡，對他說：『我主我王啊，你不是曾向使女起誓說：你兒子所羅門必接續我作王，他必坐在我的王位上嗎？亞多尼雅怎麼作了王呢？』 \end{tabularx} \\ \\ \relax
1:14 & \begin{tabularx}{0.7\textwidth}{X} 看哪，你還在那裡與王說話的時候，我會隨後進去，證實你的話。」 \end{tabularx} \\ \\ \relax
1:15 & \begin{tabularx}{0.7\textwidth}{X} 拔示巴進入內室，到王那裡。那時，王很老了，書念女子亞比煞正伺候著王。 \end{tabularx} \\ \\ \relax
1:16 & \begin{tabularx}{0.7\textwidth}{X} 拔示巴向王屈身下拜，王說：「你要甚麼？」 \end{tabularx} \\ \\ \relax
1:17 & \begin{tabularx}{0.7\textwidth}{X} 她對王說：「我主啊，你曾向使女指著耶和華－你的神起誓：『你兒子所羅門必接續我作王，他必坐在我的王位上。』 \end{tabularx} \\ \\ \relax
1:18 & \begin{tabularx}{0.7\textwidth}{X} 現在，看哪，亞多尼雅作王了，你，我主我王卻不知道。 \end{tabularx} \\ \\ \relax
1:19 & \begin{tabularx}{0.7\textwidth}{X} 他獻許多牛羊、肥犢為祭，請了王的眾兒子和亞比亞他祭司，以及約押元帥，他卻沒有請王的僕人所羅門。 \end{tabularx} \\ \\ \relax
1:20 & \begin{tabularx}{0.7\textwidth}{X} 但你，我主我王啊，以色列眾人的眼目都仰望你，等你告訴他們，在我主我王之後誰坐你的王位。 \end{tabularx} \\ \\ \relax
1:21 & \begin{tabularx}{0.7\textwidth}{X} 若不然，我主我王與祖先同睡的時候，我和我兒子所羅門必列為罪犯了。」 \end{tabularx} \\ \\ \relax
1:22 & \begin{tabularx}{0.7\textwidth}{X} 看哪，拔示巴還與王說話的時候，拿單先知也進來了。 \end{tabularx} \\ \\ \relax
1:23 & \begin{tabularx}{0.7\textwidth}{X} 有人奏告王說：「看哪，拿單先知來了。」拿單進到王面前，臉伏於地，向王叩拜。 \end{tabularx} \\ \\ \relax
1:24 & \begin{tabularx}{0.7\textwidth}{X} 拿單說：「我主我王，你果真說過『亞多尼雅必接續我作王，他要坐在我的王位上』嗎？ \end{tabularx} \\ \\ \relax
1:25 & \begin{tabularx}{0.7\textwidth}{X} 他今日下去，獻了許多牛羊、肥犢為祭，請了王的眾兒子和軍官們，以及亞比亞他祭司；看哪，他們正在亞多尼雅面前吃喝，說：『亞多尼雅王萬歲！』 \end{tabularx} \\ \\ \relax
1:26 & \begin{tabularx}{0.7\textwidth}{X} 至於我，就是你的僕人，和撒督祭司、耶何耶大的兒子比拿雅、王的僕人所羅門，他都沒有請。 \end{tabularx} \\ \\ \relax
1:27 & \begin{tabularx}{0.7\textwidth}{X} 這事果真出於我主我王嗎？王卻沒有告訴僕人們，在我主我王之後誰坐你的王位。」 \end{tabularx} \\ \\ \relax
1:28 & \begin{tabularx}{0.7\textwidth}{X} 大衛王回答說：「召拔示巴到我這裡來。」拔示巴就來，站在王面前。 \end{tabularx} \\ \\ \relax
1:29 & \begin{tabularx}{0.7\textwidth}{X} 王起誓說：「我指著救我性命脫離一切苦難的永生的耶和華起誓。 \end{tabularx} \\ \\ \relax
1:30 & \begin{tabularx}{0.7\textwidth}{X} 我既然指著耶和華－以色列的神向你起誓說：你兒子所羅門必接續我作王，他必繼承我坐在我的王位上，我今日必這樣做。」 \end{tabularx} \\ \\ \relax
1:31 & \begin{tabularx}{0.7\textwidth}{X} 於是，拔示巴屈身，臉伏於地，向王叩拜，說：「我主大衛王萬歲！」 \end{tabularx} \\ \\ \relax
1:32 & \begin{tabularx}{0.7\textwidth}{X} 大衛王又說：「召撒督祭司、拿單先知、耶何耶大的兒子比拿雅到我這裡來！」他們就都來到王面前。 \end{tabularx} \\ \\ \relax
1:33 & \begin{tabularx}{0.7\textwidth}{X} 王對他們說：「要帶領你們主的僕人，讓我兒子所羅門騎我自己的騾子，送他下到基訓。 \end{tabularx} \\ \\ \relax
1:34 & \begin{tabularx}{0.7\textwidth}{X} 在那裡，撒督祭司和拿單先知要膏他作以色列的王；你們也要吹角，說：『所羅門王萬歲！』 \end{tabularx} \\ \\ \relax
1:35 & \begin{tabularx}{0.7\textwidth}{X} 你們要跟隨他上來，使他坐在我的王位上，他要接續我作王。我已立他作以色列和猶大的君王。」 \end{tabularx} \\ \\ \relax
1:36 & \begin{tabularx}{0.7\textwidth}{X} 耶何耶大的兒子比拿雅回應王說：「阿們！願耶和華－我主我王的神這樣說。 \end{tabularx} \\ \\ \relax
1:37 & \begin{tabularx}{0.7\textwidth}{X} 耶和華怎樣與我主我王同在，願他照樣與所羅門同在，使他的王位比我主大衛王的王位更大。」 \end{tabularx} \\ \\ \relax
1:38 & \begin{tabularx}{0.7\textwidth}{X} 於是，撒督祭司、拿單先知、耶何耶大的兒子比拿雅，以及基利提人和比利提人都下去，讓所羅門騎上大衛王的騾子，送他到基訓。 \end{tabularx} \\ \\ \relax
1:39 & \begin{tabularx}{0.7\textwidth}{X} 撒督祭司從帳幕中取了盛膏油的角來，膏所羅門。他們就吹角，眾百姓都說：「所羅門王萬歲！」 \end{tabularx} \\ \\ \relax
1:40 & \begin{tabularx}{0.7\textwidth}{X} 眾百姓跟隨他上來，吹著笛，大大歡呼，地被他們的聲音震裂。 \end{tabularx} \\ \\ \relax
1:41 & \begin{tabularx}{0.7\textwidth}{X} 亞多尼雅和所有的賓客剛吃完，聽見這聲音；約押聽見角聲就說：「城中為何有這響聲呢？」 \end{tabularx} \\ \\ \relax
1:42 & \begin{tabularx}{0.7\textwidth}{X} 他正說話的時候，看哪，亞比亞他祭司的兒子約拿單來了。亞多尼雅說：「進來吧！你是個賢明的人，必是來報好消息的。」 \end{tabularx} \\ \\ \relax
1:43 & \begin{tabularx}{0.7\textwidth}{X} 約拿單回答亞多尼雅說：「我們的主大衛王已經立所羅門為王了！ \end{tabularx} \\ \\ \relax
1:44 & \begin{tabularx}{0.7\textwidth}{X} 王派撒督祭司、拿單先知、耶何耶大的兒子比拿雅，以及基利提人和比利提人和所羅門一起去，叫他騎上王的騾子。 \end{tabularx} \\ \\ \relax
1:45 & \begin{tabularx}{0.7\textwidth}{X} 撒督祭司和拿單先知已經在基訓膏他作王了。他們從那裡歡呼著上來，城都震動，這就是你們所聽見的聲音。 \end{tabularx} \\ \\ \relax
1:46 & \begin{tabularx}{0.7\textwidth}{X} 所羅門也已經登上國度的王位了。 \end{tabularx} \\ \\ \relax
1:47 & \begin{tabularx}{0.7\textwidth}{X} 王的臣僕也來為我們的主大衛王祝福，說：『願神使所羅門的名比你的名更尊榮，使他的王位比你的王位更大。』王在床上屈身敬拜， \end{tabularx} \\ \\ \relax
1:48 & \begin{tabularx}{0.7\textwidth}{X} 王也這樣說：『耶和華－以色列的神是應當稱頌的，因他今日賞賜一個人坐在我的王位上，我也親眼看見了。』」 \end{tabularx} \\ \\ \relax
1:49 & \begin{tabularx}{0.7\textwidth}{X} 亞多尼雅所有的賓客都戰兢，起來，各走各路去了。 \end{tabularx} \\ \\ \relax
1:50 & \begin{tabularx}{0.7\textwidth}{X} 亞多尼雅懼怕所羅門，就起來，去抓住祭壇的翹角。 \end{tabularx} \\ \\ \relax
1:51 & \begin{tabularx}{0.7\textwidth}{X} 有人告訴所羅門說：「看哪，亞多尼雅懼怕所羅門王。看哪，他抓住祭壇的翹角，說：『願所羅門王先向我起誓，必不用刀殺死僕人。』」 \end{tabularx} \\ \\ \relax
1:52 & \begin{tabularx}{0.7\textwidth}{X} 所羅門說：「他若作賢明的人，連一根頭髮也不致落在地上；他若作惡，必要死亡。」 \end{tabularx} \\ \\ \relax
1:53 & \begin{tabularx}{0.7\textwidth}{X} 於是所羅門王派人叫亞多尼雅從壇上下來，他就來向所羅門王下拜。所羅門對他說：「你回家去吧！」 \end{tabularx} \\ \\ \relax
2:1 & \begin{tabularx}{0.7\textwidth}{X} 大衛的死期臨近了，就吩咐他兒子所羅門說： \end{tabularx} \\ \\ \relax
2:2 & \begin{tabularx}{0.7\textwidth}{X} 「我要走世人必走的路了。你當剛強，作大丈夫， \end{tabularx} \\ \\ \relax
2:3 & \begin{tabularx}{0.7\textwidth}{X} 遵守耶和華－你神所吩咐的，照著摩西律法上所寫的行耶和華的道，謹守他的律例、誡命、典章、法度，好讓你無論做甚麼，不拘往何處去，盡都亨通。 \end{tabularx} \\ \\ \relax
2:4 & \begin{tabularx}{0.7\textwidth}{X} 耶和華必成就他所說關於我的話，說：『你的子孫若謹慎自己的行為，盡心盡意憑信實行在我面前，就不斷有人坐以色列的王位。』 \end{tabularx} \\ \\ \relax
2:5 & \begin{tabularx}{0.7\textwidth}{X} 你也知道洗魯雅的兒子約押向我所做的事，他對付以色列的兩個元帥，尼珥的兒子押尼珥和益帖的兒子亞瑪撒，殺了他們。他在太平之時，如同戰爭一般，流這二人的血，把這戰爭的血染了他腰間束的帶和腳上穿的鞋。 \end{tabularx} \\ \\ \relax
2:6 & \begin{tabularx}{0.7\textwidth}{X} 所以你要照你的智慧去做，不讓他白髮安然下陰間。 \end{tabularx} \\ \\ \relax
2:7 & \begin{tabularx}{0.7\textwidth}{X} 你當恩待基列人巴西萊的眾兒子，請他們常與你同席吃飯，因為我躲避你哥哥押沙龍的時候，他們親近我。 \end{tabularx} \\ \\ \relax
2:8 & \begin{tabularx}{0.7\textwidth}{X} 看哪，在你這裡有來自巴戶琳的便雅憫人，基拉的兒子示每。我到瑪哈念去的那日，他用狠毒的言語咒罵我。後來他卻下約旦河迎接我，我就指著耶和華向他起誓說：『我必不用刀殺死你。』 \end{tabularx} \\ \\ \relax
2:9 & \begin{tabularx}{0.7\textwidth}{X} 但現在你不要以他為無罪。你是有智慧的人，必知道怎樣待他，使他白髮流血下陰間。」 \end{tabularx} \\ \\ \relax
2:10 & \begin{tabularx}{0.7\textwidth}{X} 大衛與他祖先同睡，葬在大衛城。 \end{tabularx} \\ \\ \relax
2:11 & \begin{tabularx}{0.7\textwidth}{X} 大衛作以色列王四十年：在希伯崙作王七年，在耶路撒冷作王三十三年。 \end{tabularx} \\ \\ \relax
2:12 & \begin{tabularx}{0.7\textwidth}{X} 所羅門坐他父親大衛的王位，他的國度非常穩固。 \end{tabularx} \\ \\ \relax
2:13 & \begin{tabularx}{0.7\textwidth}{X} 哈及的兒子亞多尼雅到所羅門的母親拔示巴那裡，拔示巴問他說：「你是為平安來的嗎？」他說：「為平安來的。」 \end{tabularx} \\ \\ \relax
2:14 & \begin{tabularx}{0.7\textwidth}{X} 他又說：「我有話對你說。」拔示巴說：「你說吧。」 \end{tabularx} \\ \\ \relax
2:15 & \begin{tabularx}{0.7\textwidth}{X} 亞多尼雅說：「你知道這國原是歸我的，全以色列也都期望我作王。然而，這國反歸了我兄弟，因這國歸了他是出乎耶和華。 \end{tabularx} \\ \\ \relax
2:16 & \begin{tabularx}{0.7\textwidth}{X} 現在我有一件事求你，請你不要推辭。」拔示巴對他說：「你說吧。」 \end{tabularx} \\ \\ \relax
2:17 & \begin{tabularx}{0.7\textwidth}{X} 他說：「求你請所羅門王把書念女子亞比煞賜我為妻，因他必不拒絕你。」 \end{tabularx} \\ \\ \relax
2:18 & \begin{tabularx}{0.7\textwidth}{X} 拔示巴說：「好，我必為你對王提說。」 \end{tabularx} \\ \\ \relax
2:19 & \begin{tabularx}{0.7\textwidth}{X} 於是，拔示巴來到所羅門王那裡，要為亞多尼雅說話。王起來迎接，向她下拜，然後坐在自己的位上，又為王的母親設一座位，她就坐在王的右邊。 \end{tabularx} \\ \\ \relax
2:20 & \begin{tabularx}{0.7\textwidth}{X} 拔示巴說：「我要向你提出一個小小的請求，請你不要回絕我。」王對她說：「母親，請提出來，我必不回絕你。」 \end{tabularx} \\ \\ \relax
2:21 & \begin{tabularx}{0.7\textwidth}{X} 拔示巴說：「請你把書念女子亞比煞賜給你哥哥亞多尼雅為妻。」 \end{tabularx} \\ \\ \relax
2:22 & \begin{tabularx}{0.7\textwidth}{X} 所羅門王回答母親說：「為何替亞多尼雅求書念女子亞比煞呢？可以為他求王國吧！他是我的兄長，不但為他，也為亞比亞他祭司和洗魯雅的兒子約押求吧！」 \end{tabularx} \\ \\ \relax
2:23 & \begin{tabularx}{0.7\textwidth}{X} 所羅門王指著耶和華起誓說：「亞多尼雅講這話是自己送命，不然，願神重重懲罰我。 \end{tabularx} \\ \\ \relax
2:24 & \begin{tabularx}{0.7\textwidth}{X} 耶和華堅立我，使我坐在父親大衛的王位上，照著他所應許的為我建立家室；現在我指著永生的耶和華起誓，亞多尼雅今日必被處死。」 \end{tabularx} \\ \\ \relax
2:25 & \begin{tabularx}{0.7\textwidth}{X} 於是所羅門王派耶何耶大的兒子比拿雅去擊殺亞多尼雅，他就死了。 \end{tabularx} \\ \\ \relax
2:26 & \begin{tabularx}{0.7\textwidth}{X} 王對亞比亞他祭司說：「你回亞拿突歸自己的田地去吧！你本是該死的，但因你在我父親大衛面前抬過主耶和華的約櫃，又與我父親同受一切苦難，所以我今日不殺死你。」 \end{tabularx} \\ \\ \relax
2:27 & \begin{tabularx}{0.7\textwidth}{X} 所羅門就革除亞比亞他，不讓他作耶和華的祭司。這就應驗了耶和華在示羅論以利家所說的話。 \end{tabularx} \\ \\ \relax
2:28 & \begin{tabularx}{0.7\textwidth}{X} 雖然約押沒有擁護押沙龍，卻擁護了亞多尼雅；這消息傳到約押那裡，他就逃到耶和華的帳幕，抓住祭壇的翹角。 \end{tabularx} \\ \\ \relax
2:29 & \begin{tabularx}{0.7\textwidth}{X} 有人告訴所羅門王：「約押逃到耶和華的帳幕，看哪，他在祭壇的旁邊。」所羅門就派耶何耶大的兒子比拿雅，說：「去，殺了他。」 \end{tabularx} \\ \\ \relax
2:30 & \begin{tabularx}{0.7\textwidth}{X} 比拿雅來到耶和華的帳幕，對約押說：「王這樣吩咐：『你出來吧！』」他說：「不，我要死在這裡。」比拿雅就去回覆王，說：「約押這樣說，他這樣回答我。」 \end{tabularx} \\ \\ \relax
2:31 & \begin{tabularx}{0.7\textwidth}{X} 王對他說：「你可以照著他的話去做，殺了他，把他葬了，好叫約押流無辜人血的罪不歸在我和我的父家。 \end{tabularx} \\ \\ \relax
2:32 & \begin{tabularx}{0.7\textwidth}{X} 耶和華必使約押的血歸到他自己頭上，因為他擊殺兩個比他又公義又良善的人，就是尼珥的兒子以色列的元帥押尼珥和益帖的兒子猶大的元帥亞瑪撒，用刀殺了他們，我父親大衛卻不知道。 \end{tabularx} \\ \\ \relax
2:33 & \begin{tabularx}{0.7\textwidth}{X} 這二人的血必歸到約押和他後裔頭上，直到永遠；惟有大衛和他的後裔，以及他的家與王位，必從耶和華那裡得平安，直到永遠。」 \end{tabularx} \\ \\ \relax
2:34 & \begin{tabularx}{0.7\textwidth}{X} 於是耶何耶大的兒子比拿雅上去，擊殺約押，殺死他，把他葬在曠野約押自己的家裡。 \end{tabularx} \\ \\ \relax
2:35 & \begin{tabularx}{0.7\textwidth}{X} 王就立耶何耶大的兒子比拿雅作元帥，代替約押，又使撒督祭司代替亞比亞他。 \end{tabularx} \\ \\ \relax
2:36 & \begin{tabularx}{0.7\textwidth}{X} 王派人召示每來，對他說：「你要在耶路撒冷為自己建造房屋，住在那裡，不可從那裡出來到任何地方去。 \end{tabularx} \\ \\ \relax
2:37 & \begin{tabularx}{0.7\textwidth}{X} 你當確實知道，你何日出來過汲淪溪，就必定死！你的血必歸到自己頭上。」 \end{tabularx} \\ \\ \relax
2:38 & \begin{tabularx}{0.7\textwidth}{X} 示每對王說：「這話很好！我主我王怎樣說，僕人必照樣做。」於是示每住在耶路撒冷許多日子。 \end{tabularx} \\ \\ \relax
2:39 & \begin{tabularx}{0.7\textwidth}{X} 過了三年，示每的兩個奴僕逃到瑪迦的兒子迦特王亞吉那裡去。有人告訴示每說：「看哪，你的奴僕在迦特。」 \end{tabularx} \\ \\ \relax
2:40 & \begin{tabularx}{0.7\textwidth}{X} 示每起來，備上驢，往迦特到亞吉那裡去找他的奴僕，從迦特帶他的奴僕回來。 \end{tabularx} \\ \\ \relax
2:41 & \begin{tabularx}{0.7\textwidth}{X} 有人告訴所羅門：「示每出耶路撒冷到迦特去，又回來了。」 \end{tabularx} \\ \\ \relax
2:42 & \begin{tabularx}{0.7\textwidth}{X} 王就派人召示每來，對他說：「我豈不是叫你指著耶和華起誓，並且警告你說『你當確實知道，你何日出來到任何地方去，就必定死』嗎？你也對我說：『這話很好，我必聽從。』 \end{tabularx} \\ \\ \relax
2:43 & \begin{tabularx}{0.7\textwidth}{X} 你為何不遵守你對耶和華的誓言和我吩咐你的命令呢？」 \end{tabularx} \\ \\ \relax
2:44 & \begin{tabularx}{0.7\textwidth}{X} 王又對示每說：「你向我父親大衛所做的一切惡事，你自己心裡都知道，耶和華必使你的罪惡歸到你自己的頭上。 \end{tabularx} \\ \\ \relax
2:45 & \begin{tabularx}{0.7\textwidth}{X} 但所羅門王必蒙福，大衛的王位必在耶和華面前堅立，直到永遠。」 \end{tabularx} \\ \\ \relax
2:46 & \begin{tabularx}{0.7\textwidth}{X} 於是王吩咐耶何耶大的兒子比拿雅，他就出去，擊殺示每，示每就死了。這樣，國度在所羅門的手中鞏固了。 \end{tabularx} \\ \\ \relax
3:1 & \begin{tabularx}{0.7\textwidth}{X} 所羅門與埃及王法老結親，娶了法老的女兒，接她進入大衛城，直等到建完了自己的宮和耶和華的殿，以及耶路撒冷周圍的城牆。 \end{tabularx} \\ \\ \relax
3:2 & \begin{tabularx}{0.7\textwidth}{X} 當那些日子，百姓仍在丘壇獻祭，因為還沒有為耶和華的名建殿。 \end{tabularx} \\ \\ \relax
3:3 & \begin{tabularx}{0.7\textwidth}{X} 所羅門愛耶和華，遵行他父親大衛的律例，只是還在丘壇獻祭燒香。 \end{tabularx} \\ \\ \relax
3:4 & \begin{tabularx}{0.7\textwidth}{X} 所羅門王到基遍，在那裡獻祭，因為基遍有極大的丘壇。所羅門在那壇上獻了一千祭牲為燔祭。 \end{tabularx} \\ \\ \relax
3:5 & \begin{tabularx}{0.7\textwidth}{X} 在基遍，耶和華夜間在夢中向所羅門顯現；神說：「你願我賜你甚麼，你可以求。」 \end{tabularx} \\ \\ \relax
3:6 & \begin{tabularx}{0.7\textwidth}{X} 所羅門說：「你曾向你僕人我父親大衛大施慈愛，因為他用忠信、公義、正直的心行在你面前。你又為他存留大慈愛，賜他一個兒子坐在他的王位上，正如今日一樣。 \end{tabularx} \\ \\ \relax
3:7 & \begin{tabularx}{0.7\textwidth}{X} 現在，耶和華－我的神啊，你使僕人接續我父親大衛作王；但我是幼小的孩子，不知道應當怎樣出入。 \end{tabularx} \\ \\ \relax
3:8 & \begin{tabularx}{0.7\textwidth}{X} 僕人住在你揀選的百姓中，這百姓之多，多得不可點，不可算。 \end{tabularx} \\ \\ \relax
3:9 & \begin{tabularx}{0.7\textwidth}{X} 所以求你賜僕人善於了解的心，可以判斷你的百姓，辨別是非。不然，誰能判斷你這麼多的百姓呢？」 \end{tabularx} \\ \\ \relax
3:10 & \begin{tabularx}{0.7\textwidth}{X} 所羅門因為求這事，就蒙主喜悅。 \end{tabularx} \\ \\ \relax
3:11 & \begin{tabularx}{0.7\textwidth}{X} 神對他說：「你既然求這事，不為自己求壽、求富，也不求滅絕你仇敵的性命，只求能明辨，可以聽訟， \end{tabularx} \\ \\ \relax
3:12 & \begin{tabularx}{0.7\textwidth}{X} 看哪，我會照你的話去做，看哪，我會賜你智慧和明辨的心，在你以前沒有像你的，在你以後也沒有興起像你的。 \end{tabularx} \\ \\ \relax
3:13 & \begin{tabularx}{0.7\textwidth}{X} 你沒有求的，我也賜給你，就是富足、尊榮，使你在世一切的日子，列王中沒有一個能比你的。 \end{tabularx} \\ \\ \relax
3:14 & \begin{tabularx}{0.7\textwidth}{X} 你若遵行我的道，謹守我的律例、誡命，正如你父親大衛所行的，我必使你長壽。」 \end{tabularx} \\ \\ \relax
3:15 & \begin{tabularx}{0.7\textwidth}{X} 所羅門醒了，看哪，是個夢。他就來到耶路撒冷，站在耶和華的約櫃前，獻燔祭和平安祭，又為眾臣僕擺設宴席。 \end{tabularx} \\ \\ \relax
3:16 & \begin{tabularx}{0.7\textwidth}{X} 那時，有兩個妓女來，站在王面前。 \end{tabularx} \\ \\ \relax
3:17 & \begin{tabularx}{0.7\textwidth}{X} 一個婦人說：「我主啊，我和這婦人同住一屋。她在屋子裡的時候，我生了一個孩子。 \end{tabularx} \\ \\ \relax
3:18 & \begin{tabularx}{0.7\textwidth}{X} 我生了以後第三天，這婦人也生了。我們是一起的，屋子裡除了我們二人之外，再沒有別人在屋子裡。 \end{tabularx} \\ \\ \relax
3:19 & \begin{tabularx}{0.7\textwidth}{X} 夜間，這婦人的兒子死了，因為她壓在她的兒子身上。 \end{tabularx} \\ \\ \relax
3:20 & \begin{tabularx}{0.7\textwidth}{X} 她半夜起來，趁你使女睡著的時候，從我旁邊把我兒子抱走，放在她懷裡，又把她死的兒子放在我懷裡。 \end{tabularx} \\ \\ \relax
3:21 & \begin{tabularx}{0.7\textwidth}{X} 清早，我起來要給我的兒子吃奶，看哪，他死了；早晨我仔細察看他，看哪，他不是我所生的兒子。」 \end{tabularx} \\ \\ \relax
3:22 & \begin{tabularx}{0.7\textwidth}{X} 另一個婦人說：「不！我的兒子是活的，你的兒子是死的。」但這一個說：「不！你的兒子是死的，我的兒子是活的。」她們就在王面前爭吵。 \end{tabularx} \\ \\ \relax
3:23 & \begin{tabularx}{0.7\textwidth}{X} 王說：「這婦人說：『這是我的兒子，他是活的，你的兒子是死的。』那婦人說：『不！你的兒子是死的，我的兒子是活的。』」 \end{tabularx} \\ \\ \relax
3:24 & \begin{tabularx}{0.7\textwidth}{X} 王就說：「給我拿刀來！」人就把刀拿到王面前來。 \end{tabularx} \\ \\ \relax
3:25 & \begin{tabularx}{0.7\textwidth}{X} 王說：「把活孩子劈成兩半，一半給這婦人，一半給那婦人。」 \end{tabularx} \\ \\ \relax
3:26 & \begin{tabularx}{0.7\textwidth}{X} 活孩子的母親為自己的兒子心急如焚，對王說：「求我主把活孩子給那婦人吧，萬不可殺死他！」那婦人說：「這孩子也不歸我，也不歸你，你們就劈了吧！」 \end{tabularx} \\ \\ \relax
3:27 & \begin{tabularx}{0.7\textwidth}{X} 王回應說：「把活孩子給這婦人，萬不可殺死他，因為這婦人是他的母親。」 \end{tabularx} \\ \\ \relax
3:28 & \begin{tabularx}{0.7\textwidth}{X} 全以色列聽見王這樣判斷，就都敬畏王，因為他們看見他心中有神的智慧，能夠斷案。 \end{tabularx} \\ \\ \relax
4:1 & \begin{tabularx}{0.7\textwidth}{X} 所羅門作全以色列的王。 \end{tabularx} \\ \\ \relax
4:2 & \begin{tabularx}{0.7\textwidth}{X} 這些是他的官員：撒督的兒子亞撒利雅作祭司， \end{tabularx} \\ \\ \relax
4:3 & \begin{tabularx}{0.7\textwidth}{X} 示沙的兩個兒子以利何烈、亞希亞作書記，亞希律的兒子約沙法作史官， \end{tabularx} \\ \\ \relax
4:4 & \begin{tabularx}{0.7\textwidth}{X} 耶何耶大的兒子比拿雅作元帥，撒督和亞比亞他作祭司， \end{tabularx} \\ \\ \relax
4:5 & \begin{tabularx}{0.7\textwidth}{X} 拿單的兒子亞撒利雅作宰相，拿單的兒子撒布得作祭司和王的顧問， \end{tabularx} \\ \\ \relax
4:6 & \begin{tabularx}{0.7\textwidth}{X} 亞希煞作管家，亞比大的兒子亞多尼蘭掌管服勞役的工人。 \end{tabularx} \\ \\ \relax
4:7 & \begin{tabularx}{0.7\textwidth}{X} 所羅門在全以色列有十二個官員，供給王和王室的食物，每年各人供給一個月。 \end{tabularx} \\ \\ \relax
4:8 & \begin{tabularx}{0.7\textwidth}{X} 這些是他們的名字：在以法蓮山區有便‧戶珥； \end{tabularx} \\ \\ \relax
4:9 & \begin{tabularx}{0.7\textwidth}{X} 在瑪迦斯、沙賓、伯‧示麥、以倫‧伯‧哈南有便‧底甲； \end{tabularx} \\ \\ \relax
4:10 & \begin{tabularx}{0.7\textwidth}{X} 在亞魯泊有便‧希悉，他管理梭哥和希弗全地； \end{tabularx} \\ \\ \relax
4:11 & \begin{tabularx}{0.7\textwidth}{X} 在多珥山岡有便‧亞比拿達，他娶了所羅門的女兒她法為妻； \end{tabularx} \\ \\ \relax
4:12 & \begin{tabularx}{0.7\textwidth}{X} 在他納和米吉多，以及靠近撒拉他拿、耶斯列下邊的伯‧善全地，從伯‧善到亞伯‧米何拉直到約緬的另一邊有亞希律的兒子巴拿； \end{tabularx} \\ \\ \relax
4:13 & \begin{tabularx}{0.7\textwidth}{X} 在基列的拉末有便‧基別，他管理在基列的瑪拿西子孫睚珥的城鎮，巴珊的亞珥歌伯地的六十座大城，各有城牆和銅閂； \end{tabularx} \\ \\ \relax
4:14 & \begin{tabularx}{0.7\textwidth}{X} 在瑪哈念有易多的兒子亞希拿達； \end{tabularx} \\ \\ \relax
4:15 & \begin{tabularx}{0.7\textwidth}{X} 在拿弗他利有亞希瑪斯，他也娶了所羅門的一個女兒巴實抹為妻； \end{tabularx} \\ \\ \relax
4:16 & \begin{tabularx}{0.7\textwidth}{X} 在亞設和亞祿有戶篩的兒子巴拿； \end{tabularx} \\ \\ \relax
4:17 & \begin{tabularx}{0.7\textwidth}{X} 在以薩迦有帕路亞的兒子約沙法； \end{tabularx} \\ \\ \relax
4:18 & \begin{tabularx}{0.7\textwidth}{X} 在便雅憫有以拉的兒子示每； \end{tabularx} \\ \\ \relax
4:19 & \begin{tabularx}{0.7\textwidth}{X} 在基列地，就是亞摩利王西宏和巴珊王噩之地，有烏利的兒子基別，他一個官員管理這地 。 \end{tabularx} \\ \\ \relax
4:20 & \begin{tabularx}{0.7\textwidth}{X} 猶大人和以色列人如同海邊的沙那樣多，都吃喝快樂。 \end{tabularx} \\ \\ \relax
4:21 & \begin{tabularx}{0.7\textwidth}{X} 所羅門統治諸國，從大河到非利士地，直到埃及的邊界。所羅門在世的日子，這些國都向他進貢，服事他。 \end{tabularx} \\ \\ \relax
4:22 & \begin{tabularx}{0.7\textwidth}{X} 所羅門每日所用的食物：三十歌珥細麵，六十歌珥粗麵， \end{tabularx} \\ \\ \relax
4:23 & \begin{tabularx}{0.7\textwidth}{X} 十頭肥牛，二十頭草場的牛，一百隻羊，還有鹿、羚羊、麃子，以及肥禽。 \end{tabularx} \\ \\ \relax
4:24 & \begin{tabularx}{0.7\textwidth}{X} 所羅門管理整個大河西邊，從提弗薩直到迦薩，以及大河西邊的諸王，屬他的四境盡都平安。 \end{tabularx} \\ \\ \relax
4:25 & \begin{tabularx}{0.7\textwidth}{X} 所羅門在世的日子，從但到別是巴，猶大和以色列各人都在自己的葡萄樹下和無花果樹下安然居住。 \end{tabularx} \\ \\ \relax
4:26 & \begin{tabularx}{0.7\textwidth}{X} 所羅門擁有給戰車用的四萬個馬棚，還有一萬二千名騎兵。 \end{tabularx} \\ \\ \relax
4:27 & \begin{tabularx}{0.7\textwidth}{X} 這些官員各按自己的月份供給所羅門王，以及一切與他同席之人的食物，一無所缺。 \end{tabularx} \\ \\ \relax
4:28 & \begin{tabularx}{0.7\textwidth}{X} 他們各按其分，把給馬與快馬吃的大麥和乾草送到指定的地方去。 \end{tabularx} \\ \\ \relax
4:29 & \begin{tabularx}{0.7\textwidth}{X} 神賜給所羅門極大的智慧和聰明，以及寬闊的心，如同海邊的沙。 \end{tabularx} \\ \\ \relax
4:30 & \begin{tabularx}{0.7\textwidth}{X} 所羅門的智慧超過所有東方人的智慧，和埃及人一切的智慧。 \end{tabularx} \\ \\ \relax
4:31 & \begin{tabularx}{0.7\textwidth}{X} 他的智慧勝過萬人，勝過以斯拉人以探，以及瑪曷的兒子希幔、甲各、達大。他的名聲傳遍四圍的列國。 \end{tabularx} \\ \\ \relax
4:32 & \begin{tabularx}{0.7\textwidth}{X} 他作箴言三千句，詩歌一千零五首。 \end{tabularx} \\ \\ \relax
4:33 & \begin{tabularx}{0.7\textwidth}{X} 他講論草木，從黎巴嫩的香柏樹直到牆上長的牛膝草，又講論飛禽、走獸、爬行動物和魚類。 \end{tabularx} \\ \\ \relax
4:34 & \begin{tabularx}{0.7\textwidth}{X} 地上凡曾聽過他智慧的君王，都派人來；萬民都有人來聽所羅門的智慧。 \end{tabularx} \\ \\ \relax
5:1 & \begin{tabularx}{0.7\textwidth}{X} 推羅王希蘭是大衛平生的好友。希蘭聽見以色列人膏所羅門接續他父親作王，就派臣僕到他那裡。 \end{tabularx} \\ \\ \relax
5:2 & \begin{tabularx}{0.7\textwidth}{X} 所羅門也派人到希蘭那裡，說： \end{tabularx} \\ \\ \relax
5:3 & \begin{tabularx}{0.7\textwidth}{X} 「你知道我父親大衛因四圍的戰爭，不能為耶和華－他神的名建殿，直等到耶和華使仇敵都服在他腳下。 \end{tabularx} \\ \\ \relax
5:4 & \begin{tabularx}{0.7\textwidth}{X} 現在耶和華－我的神使我四圍太平，沒有仇敵，沒有災禍。 \end{tabularx} \\ \\ \relax
5:5 & \begin{tabularx}{0.7\textwidth}{X} 看哪，我吩咐要為耶和華－我神的名建殿，是照耶和華向我父親大衛說的：『我必使你兒子接續你，坐你的王位，他必為我的名建殿。』 \end{tabularx} \\ \\ \relax
5:6 & \begin{tabularx}{0.7\textwidth}{X} 現在，請吩咐人在黎巴嫩為我砍伐香柏木，我的僕人必幫助你的僕人。至於你僕人的工錢，我必照你所定的給你。你知道，在我們中間沒有人像西頓人那樣擅長砍伐樹木。」 \end{tabularx} \\ \\ \relax
5:7 & \begin{tabularx}{0.7\textwidth}{X} 希蘭聽見所羅門的話，就很高興，說：「今日耶和華是應當稱頌的，因為他賜給大衛一個有智慧的兒子，治理這眾多的百姓。」 \end{tabularx} \\ \\ \relax
5:8 & \begin{tabularx}{0.7\textwidth}{X} 希蘭送信給所羅門，說：「你派人向我所提的那事，我已聽見了；論到香柏木和松木，我必照你一切的心願去做。 \end{tabularx} \\ \\ \relax
5:9 & \begin{tabularx}{0.7\textwidth}{X} 我的僕人必把這木料從黎巴嫩運到海裡，我會把它們紮成筏子浮在海上，運到你告訴我的地方，在那裡拆開，你就可以收取；你也要照我的心願做，把食物給我的家。」 \end{tabularx} \\ \\ \relax
5:10 & \begin{tabularx}{0.7\textwidth}{X} 於是希蘭照所羅門的心願，給他香柏木和松木； \end{tabularx} \\ \\ \relax
5:11 & \begin{tabularx}{0.7\textwidth}{X} 所羅門給希蘭二萬歌珥麥子，二十歌珥搗成的油，作他家的食物。所羅門每年都是這樣給希蘭。 \end{tabularx} \\ \\ \relax
5:12 & \begin{tabularx}{0.7\textwidth}{X} 耶和華照著所應許的賜智慧給所羅門。希蘭與所羅門和平相處，二人彼此立約。 \end{tabularx} \\ \\ \relax
5:13 & \begin{tabularx}{0.7\textwidth}{X} 所羅門王從全以色列挑取服勞役的人，徵來的人有三萬， \end{tabularx} \\ \\ \relax
5:14 & \begin{tabularx}{0.7\textwidth}{X} 派他們輪流每月一萬人上黎巴嫩去；一個月在黎巴嫩，兩個月在家裡。亞多尼蘭管理他們。 \end{tabularx} \\ \\ \relax
5:15 & \begin{tabularx}{0.7\textwidth}{X} 所羅門有七萬扛抬的，八萬在山上鑿石頭的。 \end{tabularx} \\ \\ \relax
5:16 & \begin{tabularx}{0.7\textwidth}{X} 此外，所羅門有三千三百個監督工作的官長，監管百姓做工。 \end{tabularx} \\ \\ \relax
5:17 & \begin{tabularx}{0.7\textwidth}{X} 王下令，他們就鑿出又大又貴重的石頭來，用以立殿的根基。 \end{tabularx} \\ \\ \relax
5:18 & \begin{tabularx}{0.7\textwidth}{X} 所羅門的工匠和希蘭的工匠，以及迦巴勒人，把石頭鑿好，預備了木料和石頭來建殿。 \end{tabularx} \\ \\ \relax
6:1 & \begin{tabularx}{0.7\textwidth}{X} 以色列人出埃及地後四百八十年，所羅門作以色列王第四年西弗月，就是二月，他開工建造耶和華的殿。 \end{tabularx} \\ \\ \relax
6:2 & \begin{tabularx}{0.7\textwidth}{X} 所羅門王為耶和華所建的殿，長六十肘，寬二十肘，高三十肘。 \end{tabularx} \\ \\ \relax
6:3 & \begin{tabularx}{0.7\textwidth}{X} 殿的正堂前走廊長二十肘，與殿的寬度一樣，殿前寬十肘； \end{tabularx} \\ \\ \relax
6:4 & \begin{tabularx}{0.7\textwidth}{X} 他為殿做了有框嵌壁式的窗戶。 \end{tabularx} \\ \\ \relax
6:5 & \begin{tabularx}{0.7\textwidth}{X} 靠著殿牆，圍著外殿和內殿的牆，周圍建造了廂房； \end{tabularx} \\ \\ \relax
6:6 & \begin{tabularx}{0.7\textwidth}{X} 下層寬五肘，中層寬六肘，第三層寬七肘。他在殿牆的周圍造坎，免得梁木插入殿牆裡。 \end{tabularx} \\ \\ \relax
6:7 & \begin{tabularx}{0.7\textwidth}{X} 殿是用山中鑿成的石頭建的，所以建殿的時候，鎚子、斧子和別樣鐵器的響聲都沒有聽見。 \end{tabularx} \\ \\ \relax
6:8 & \begin{tabularx}{0.7\textwidth}{X} 在殿右邊當中的廂房有門，可以從螺旋梯上到中層，再從中層上到第三層。 \end{tabularx} \\ \\ \relax
6:9 & \begin{tabularx}{0.7\textwidth}{X} 所羅門完成殿的建造。他用香柏木作梁木和橫板，遮蓋殿頂。 \end{tabularx} \\ \\ \relax
6:10 & \begin{tabularx}{0.7\textwidth}{X} 靠著整個殿所造的廂房，每層高五肘，香柏木的梁板擱在殿的牆坎上。 \end{tabularx} \\ \\ \relax
6:11 & \begin{tabularx}{0.7\textwidth}{X} 耶和華的話臨到所羅門，說： \end{tabularx} \\ \\ \relax
6:12 & \begin{tabularx}{0.7\textwidth}{X} 「論到你所建的這殿，你若遵行我的律例，謹守我的典章，遵從我的一切誡命，行在其中，我必向你應驗我所應許你父親大衛的話。 \end{tabularx} \\ \\ \relax
6:13 & \begin{tabularx}{0.7\textwidth}{X} 我必住在以色列人中間，並不丟棄我的百姓以色列。」 \end{tabularx} \\ \\ \relax
6:14 & \begin{tabularx}{0.7\textwidth}{X} 所羅門完成殿的建造。 \end{tabularx} \\ \\ \relax
6:15 & \begin{tabularx}{0.7\textwidth}{X} 他用香柏木板建造殿的內牆，從殿的地到牆頂都貼上木板，又用松木板鋪地。 \end{tabularx} \\ \\ \relax
6:16 & \begin{tabularx}{0.7\textwidth}{X} 他在殿的後部建了一間內殿，長二十肘，從地到牆用香柏木板，作為至聖所。 \end{tabularx} \\ \\ \relax
6:17 & \begin{tabularx}{0.7\textwidth}{X} 殿，就在內殿的前面，長四十肘。 \end{tabularx} \\ \\ \relax
6:18 & \begin{tabularx}{0.7\textwidth}{X} 殿裡一點石頭都不顯露，一概用香柏木遮蔽；香柏木上刻著野瓜和綻開的花。 \end{tabularx} \\ \\ \relax
6:19 & \begin{tabularx}{0.7\textwidth}{X} 他在殿的中間預備內殿，在那裡安放耶和華的約櫃。 \end{tabularx} \\ \\ \relax
6:20 & \begin{tabularx}{0.7\textwidth}{X} 內殿長二十肘，寬二十肘，高二十肘，都貼上純金。他又用香柏木做壇。 \end{tabularx} \\ \\ \relax
6:21 & \begin{tabularx}{0.7\textwidth}{X} 所羅門用純金貼殿內，又用金鏈子掛在內殿前，內殿也貼上金子。 \end{tabularx} \\ \\ \relax
6:22 & \begin{tabularx}{0.7\textwidth}{X} 整個殿都貼上金子，直到貼滿；內殿前的整個壇，也都包上金子。 \end{tabularx} \\ \\ \relax
6:23 & \begin{tabularx}{0.7\textwidth}{X} 他在內殿裡用橄欖木做兩個基路伯，各高十肘。 \end{tabularx} \\ \\ \relax
6:24 & \begin{tabularx}{0.7\textwidth}{X} 這基路伯的一個翅膀長五肘，另一個翅膀長五肘，從一個翅膀尖到另一個翅膀尖共有十肘； \end{tabularx} \\ \\ \relax
6:25 & \begin{tabularx}{0.7\textwidth}{X} 第二個基路伯也是十肘；兩個基路伯的尺寸、形狀都一樣。 \end{tabularx} \\ \\ \relax
6:26 & \begin{tabularx}{0.7\textwidth}{X} 這一個基路伯高十肘，第二個基路伯也是如此。 \end{tabularx} \\ \\ \relax
6:27 & \begin{tabularx}{0.7\textwidth}{X} 他把兩個基路伯安在內殿中間。基路伯的翅膀是張開的，這基路伯的一個翅膀挨著這邊的牆，第二個基路伯的一個翅膀挨著那邊的牆，向內的兩個翅膀在殿中間彼此相接。 \end{tabularx} \\ \\ \relax
6:28 & \begin{tabularx}{0.7\textwidth}{X} 二基路伯都包上金子。 \end{tabularx} \\ \\ \relax
6:29 & \begin{tabularx}{0.7\textwidth}{X} 殿周圍的牆上全都刻著基路伯、棕樹和綻開的花，內外都是如此。 \end{tabularx} \\ \\ \relax
6:30 & \begin{tabularx}{0.7\textwidth}{X} 殿的地板都貼上金子，內外都是如此。 \end{tabularx} \\ \\ \relax
6:31 & \begin{tabularx}{0.7\textwidth}{X} 他用橄欖木製造內殿的入口、門楣和五邊形的門柱。 \end{tabularx} \\ \\ \relax
6:32 & \begin{tabularx}{0.7\textwidth}{X} 在橄欖木做的兩門扇上刻著基路伯、棕樹和綻開的花，都貼上金子。基路伯和棕樹上也灑上金子。 \end{tabularx} \\ \\ \relax
6:33 & \begin{tabularx}{0.7\textwidth}{X} 他又為外殿的入口，用橄欖木製造門柱，是四邊形的。 \end{tabularx} \\ \\ \relax
6:34 & \begin{tabularx}{0.7\textwidth}{X} 他用松木做兩扇門。這一扇有兩葉摺疊，第二扇也有兩葉摺疊。 \end{tabularx} \\ \\ \relax
6:35 & \begin{tabularx}{0.7\textwidth}{X} 上面刻著基路伯、棕樹和綻開的花，雕刻物都均勻地貼上金子。 \end{tabularx} \\ \\ \relax
6:36 & \begin{tabularx}{0.7\textwidth}{X} 他又用三層鑿成的石頭、香柏木一層建造內院。 \end{tabularx} \\ \\ \relax
6:37 & \begin{tabularx}{0.7\textwidth}{X} 所羅門在位第四年西弗月，立了耶和華殿的根基。 \end{tabularx} \\ \\ \relax
6:38 & \begin{tabularx}{0.7\textwidth}{X} 到十一年布勒月，就是八月，殿和一切屬殿的都按著樣式造成。他建殿共用了七年。 \end{tabularx} \\ \\ \relax
7:1 & \begin{tabularx}{0.7\textwidth}{X} 所羅門為自己建造宮殿，十三年方才建成整座宮殿。 \end{tabularx} \\ \\ \relax
7:2 & \begin{tabularx}{0.7\textwidth}{X} 他建造黎巴嫩林宮，長一百肘，寬五十肘，高三十肘，有四行香柏木柱，柱上有香柏木橫梁； \end{tabularx} \\ \\ \relax
7:3 & \begin{tabularx}{0.7\textwidth}{X} 廂房以上覆蓋著香柏木，在四十五根柱子之上，每行十五根。 \end{tabularx} \\ \\ \relax
7:4 & \begin{tabularx}{0.7\textwidth}{X} 窗戶有三排，三排的窗與窗相對。 \end{tabularx} \\ \\ \relax
7:5 & \begin{tabularx}{0.7\textwidth}{X} 所有的門和門柱都有四方形的框，共有三行，彼此相對。 \end{tabularx} \\ \\ \relax
7:6 & \begin{tabularx}{0.7\textwidth}{X} 他建造有柱子的廳，長五十肘，寬三十肘。在這前面有走廊，前面有柱子和頂蓋。 \end{tabularx} \\ \\ \relax
7:7 & \begin{tabularx}{0.7\textwidth}{X} 他又建造一個有座位的廳，就是審判廳，他在那裡審判；這廳的地板從這邊到那邊都鋪上香柏木。 \end{tabularx} \\ \\ \relax
7:8 & \begin{tabularx}{0.7\textwidth}{X} 廳後面的院內有所羅門自己住的宮殿，都用同樣的建造方式。所羅門又為所娶法老的女兒建造一座宮，建造方式與這廳一樣。 \end{tabularx} \\ \\ \relax
7:9 & \begin{tabularx}{0.7\textwidth}{X} 建造這一切所用的石頭都是貴重的，按著尺寸鑿成，用鋸子裡外鋸齊；從根基直到房檐，從外頭直到大院，都是如此。 \end{tabularx} \\ \\ \relax
7:10 & \begin{tabularx}{0.7\textwidth}{X} 根基是貴重的大石頭，有長十肘的，有長八肘的； \end{tabularx} \\ \\ \relax
7:11 & \begin{tabularx}{0.7\textwidth}{X} 上面有香柏木和按著尺寸鑿成的貴重石頭。 \end{tabularx} \\ \\ \relax
7:12 & \begin{tabularx}{0.7\textwidth}{X} 大院周圍有鑿成的石頭三層、香柏木板一層，都照耶和華殿的內院和殿的走廊的樣式。 \end{tabularx} \\ \\ \relax
7:13 & \begin{tabularx}{0.7\textwidth}{X} 所羅門王派人從推羅把戶蘭接來。 \end{tabularx} \\ \\ \relax
7:14 & \begin{tabularx}{0.7\textwidth}{X} 他是拿弗他利支派中一個寡婦的兒子，父親是推羅人，是作銅匠的。戶蘭滿有智慧、聰明、技能，善作各樣的銅器。他來到所羅門王那裡，為王做一切的工。 \end{tabularx} \\ \\ \relax
7:15 & \begin{tabularx}{0.7\textwidth}{X} 戶蘭製造兩根銅柱，一根高十八肘，第二根柱子用繩子量，周圍是十二肘； \end{tabularx} \\ \\ \relax
7:16 & \begin{tabularx}{0.7\textwidth}{X} 他做了兩個柱頂安在柱上，是用銅鑄造的，一個柱頂高五肘，第二個柱頂也高五肘。 \end{tabularx} \\ \\ \relax
7:17 & \begin{tabularx}{0.7\textwidth}{X} 柱子頂上有裝飾的網子和編成的鏈子，一個柱頂有七個，第二個柱頂也有七個。 \end{tabularx} \\ \\ \relax
7:18 & \begin{tabularx}{0.7\textwidth}{X} 他做了柱子，第一根柱子的柱頂上，周圍有兩行網子在柱子上面遮蓋柱頂，第二根柱頂也是這樣做。 \end{tabularx} \\ \\ \relax
7:19 & \begin{tabularx}{0.7\textwidth}{X} 走廊柱子頂上的柱頂高四肘，刻著百合花。 \end{tabularx} \\ \\ \relax
7:20 & \begin{tabularx}{0.7\textwidth}{X} 兩根柱子上面有柱頂，柱頂靠近網子的圓凸面上，有石榴的行列環繞著，共二百個，第二個柱頂也是如此。 \end{tabularx} \\ \\ \relax
7:21 & \begin{tabularx}{0.7\textwidth}{X} 他把兩根柱子立在殿的走廊前：右邊立一根，起名叫雅斤；左邊立一根，起名叫波阿斯。 \end{tabularx} \\ \\ \relax
7:22 & \begin{tabularx}{0.7\textwidth}{X} 柱頂上刻著百合花。這樣，柱子的工程就完畢了。 \end{tabularx} \\ \\ \relax
7:23 & \begin{tabularx}{0.7\textwidth}{X} 他又鑄一個銅海，周圍是圓的，直徑十肘，高五肘，用繩子量周圍是三十肘。 \end{tabularx} \\ \\ \relax
7:24 & \begin{tabularx}{0.7\textwidth}{X} 銅海邊緣下面的周圍有野瓜的形狀，每肘十個，共兩行，繞著銅海，是造銅海的時候鑄上去的。 \end{tabularx} \\ \\ \relax
7:25 & \begin{tabularx}{0.7\textwidth}{X} 銅海安在十二頭銅牛上：三頭向北，三頭向西，三頭向南，三頭向東。銅海安在牛上，牛尾都向內。 \end{tabularx} \\ \\ \relax
7:26 & \begin{tabularx}{0.7\textwidth}{X} 銅海厚一掌，邊如杯邊，像百合花，容量是二千罷特。 \end{tabularx} \\ \\ \relax
7:27 & \begin{tabularx}{0.7\textwidth}{X} 他用銅製造十個盆座，每座長四肘，寬四肘，高三肘。 \end{tabularx} \\ \\ \relax
7:28 & \begin{tabularx}{0.7\textwidth}{X} 銅座的造法是這樣：周圍各有嵌邊，嵌邊裝在框架中。 \end{tabularx} \\ \\ \relax
7:29 & \begin{tabularx}{0.7\textwidth}{X} 裝在框架中的嵌邊上有獅子和牛，以及基路伯。框架上有小座，獅子和牛的上面和下面有錘成的花紋浮雕。 \end{tabularx} \\ \\ \relax
7:30 & \begin{tabularx}{0.7\textwidth}{X} 每座有四個銅輪和銅軸，它有四個支架在盆以下，這些支架是鑄成的，各邊都有花紋。 \end{tabularx} \\ \\ \relax
7:31 & \begin{tabularx}{0.7\textwidth}{X} 它的口在柱頂裡，向上高一肘，口是圓的，做法如座一樣，直徑是一肘半，口上也有雕工。嵌邊是方形的，不是圓的。 \end{tabularx} \\ \\ \relax
7:32 & \begin{tabularx}{0.7\textwidth}{X} 四個輪子在嵌邊以下，輪軸與座相連，每輪高一肘半。 \end{tabularx} \\ \\ \relax
7:33 & \begin{tabularx}{0.7\textwidth}{X} 輪的樣式如同車的輪子；軸、輞、輻、轂都是鑄成的。 \end{tabularx} \\ \\ \relax
7:34 & \begin{tabularx}{0.7\textwidth}{X} 每個座四邊有四個盆形的支架，這些支架是與座從一整塊鑄成的。 \end{tabularx} \\ \\ \relax
7:35 & \begin{tabularx}{0.7\textwidth}{X} 座頂有圓架，高半肘；座頂有支柱和嵌邊，是與座從一整塊鑄成的。 \end{tabularx} \\ \\ \relax
7:36 & \begin{tabularx}{0.7\textwidth}{X} 他在支柱和嵌邊上，每個空處刻上基路伯、獅子和棕樹，周圍有花紋。 \end{tabularx} \\ \\ \relax
7:37 & \begin{tabularx}{0.7\textwidth}{X} 他按照這樣的做法造了十個盆座，它們的鑄法、尺寸、樣式全都相同。 \end{tabularx} \\ \\ \relax
7:38 & \begin{tabularx}{0.7\textwidth}{X} 他又造十個銅盆，每盆的容量四十罷特，直徑四肘。在十個座上，每座安設一盆。 \end{tabularx} \\ \\ \relax
7:39 & \begin{tabularx}{0.7\textwidth}{X} 他把五個安置在殿的右邊，五個安置在殿的左邊，又把銅海安置在殿的右旁，在東南邊。 \end{tabularx} \\ \\ \relax
7:40 & \begin{tabularx}{0.7\textwidth}{X} 戶蘭又造了盆、鏟子和盤子。這樣，戶蘭為所羅門王做完了耶和華殿一切的工： \end{tabularx} \\ \\ \relax
7:41 & \begin{tabularx}{0.7\textwidth}{X} 兩根柱子和柱子頂上兩個如碗的柱頂，以及蓋著如碗柱頂的兩個網子； \end{tabularx} \\ \\ \relax
7:42 & \begin{tabularx}{0.7\textwidth}{X} 四百個石榴，安在兩個網子上，每網兩行石榴，蓋著柱子上面兩個如碗的柱頂； \end{tabularx} \\ \\ \relax
7:43 & \begin{tabularx}{0.7\textwidth}{X} 十個盆座和其上的十個盆； \end{tabularx} \\ \\ \relax
7:44 & \begin{tabularx}{0.7\textwidth}{X} 銅海和其下的十二頭牛； \end{tabularx} \\ \\ \relax
7:45 & \begin{tabularx}{0.7\textwidth}{X} 盆、鏟子、盤子。戶蘭給所羅門王為耶和華殿造的這一切器皿都是用光亮的銅， \end{tabularx} \\ \\ \relax
7:46 & \begin{tabularx}{0.7\textwidth}{X} 是王在約旦平原、疏割和撒拉但中間的泥巴地鑄成的。 \end{tabularx} \\ \\ \relax
7:47 & \begin{tabularx}{0.7\textwidth}{X} 所羅門允許這一切器皿不過秤，因為所用的銅太多，重量無法計算。 \end{tabularx} \\ \\ \relax
7:48 & \begin{tabularx}{0.7\textwidth}{X} 所羅門又為耶和華的殿造了各樣的器皿：金壇和獻供餅的金供桌； \end{tabularx} \\ \\ \relax
7:49 & \begin{tabularx}{0.7\textwidth}{X} 內殿前的純金燈臺，右邊五個，左邊五個，以及其上的花、燈盞、燈剪，都是金的； \end{tabularx} \\ \\ \relax
7:50 & \begin{tabularx}{0.7\textwidth}{X} 純金的杯、鉗子、盤子、勺子 、火盆，以及聖殿的最裡面，就是至聖所的門樞和外殿的門樞，都是金的。 \end{tabularx} \\ \\ \relax
7:51 & \begin{tabularx}{0.7\textwidth}{X} 所羅門王做完了耶和華殿一切的工，就把他父親大衛分別為聖的金銀和器皿都帶來，放在耶和華殿的庫房裡。 \end{tabularx} \\ \\ \relax
8:1 & \begin{tabularx}{0.7\textwidth}{X} 那時，所羅門召集以色列的長老、各支派的領袖和以色列人的族長到耶路撒冷，所羅門王那裡，要把耶和華的約櫃從大衛城，就是錫安，接上來。 \end{tabularx} \\ \\ \relax
8:2 & \begin{tabularx}{0.7\textwidth}{X} 以他念月，就是七月，在節期時，所有的以色列人都聚集到所羅門王那裡。 \end{tabularx} \\ \\ \relax
8:3 & \begin{tabularx}{0.7\textwidth}{X} 以色列眾長老一來到，祭司就抬起約櫃。 \end{tabularx} \\ \\ \relax
8:4 & \begin{tabularx}{0.7\textwidth}{X} 祭司和利未人將耶和華的約櫃請上來，又把會幕和會幕一切的聖器皿都帶上來。 \end{tabularx} \\ \\ \relax
8:5 & \begin{tabularx}{0.7\textwidth}{X} 所羅門王和聚集到他那裡的以色列全會眾一同在約櫃前獻牛羊為祭，多得不可勝數，無法計算。 \end{tabularx} \\ \\ \relax
8:6 & \begin{tabularx}{0.7\textwidth}{X} 祭司將耶和華的約櫃請進內殿，就是至聖所，安置在兩個基路伯的翅膀底下約櫃的地方。 \end{tabularx} \\ \\ \relax
8:7 & \begin{tabularx}{0.7\textwidth}{X} 基路伯張開翅膀在約櫃上面的地方，從上面遮住約櫃和抬櫃的槓。 \end{tabularx} \\ \\ \relax
8:8 & \begin{tabularx}{0.7\textwidth}{X} 這槓很長，從內殿前的聖所可以看見槓頭，從外面卻看不見。這槓直到今日還在那裡。 \end{tabularx} \\ \\ \relax
8:9 & \begin{tabularx}{0.7\textwidth}{X} 約櫃裡沒有別的，只有兩塊石版，就是以色列人出埃及地，耶和華與他們立約的時候，摩西在何烈山放在那裡的。 \end{tabularx} \\ \\ \relax
8:10 & \begin{tabularx}{0.7\textwidth}{X} 祭司從聖所出來的時候，有雲充滿耶和華的殿， \end{tabularx} \\ \\ \relax
8:11 & \begin{tabularx}{0.7\textwidth}{X} 祭司因雲彩的緣故不能站立供職，因為耶和華的榮光充滿了耶和華的殿。 \end{tabularx} \\ \\ \relax
8:12 & \begin{tabularx}{0.7\textwidth}{X} 那時，所羅門說：「耶和華曾說要住在幽暗之處。 \end{tabularx} \\ \\ \relax
8:13 & \begin{tabularx}{0.7\textwidth}{X} 我的確為你建了一座雄偉的殿宇，作為你永遠居住的地方。」 \end{tabularx} \\ \\ \relax
8:14 & \begin{tabularx}{0.7\textwidth}{X} 王轉過臉來為以色列全會眾祝福，以色列全會眾都站立。 \end{tabularx} \\ \\ \relax
8:15 & \begin{tabularx}{0.7\textwidth}{X} 所羅門說：「耶和華－以色列的神是應當稱頌的！因他親口向我父大衛應許的，也親手成就了；他曾說： \end{tabularx} \\ \\ \relax
8:16 & \begin{tabularx}{0.7\textwidth}{X} 『自從那日我領我百姓以色列出埃及以來，我未曾在以色列各支派中選擇一城，在那裡為我的名建造殿宇，但我揀選大衛治理我的百姓以色列。』 \end{tabularx} \\ \\ \relax
8:17 & \begin{tabularx}{0.7\textwidth}{X} 我父大衛的心意是要為耶和華－以色列神的名建殿。 \end{tabularx} \\ \\ \relax
8:18 & \begin{tabularx}{0.7\textwidth}{X} 耶和華卻對我父大衛說：『你有心為我的名建殿，這心意是好的； \end{tabularx} \\ \\ \relax
8:19 & \begin{tabularx}{0.7\textwidth}{X} 但你不可建殿，惟有你親生的兒子才可為我的名建殿。』 \end{tabularx} \\ \\ \relax
8:20 & \begin{tabularx}{0.7\textwidth}{X} 現在耶和華實現了他所應許的話，使我接續我父大衛坐以色列的王位，正如耶和華所說的，我也為耶和華－以色列神的名建造了這殿。 \end{tabularx} \\ \\ \relax
8:21 & \begin{tabularx}{0.7\textwidth}{X} 我也在那裡為約櫃預備一處。約櫃那裡有耶和華的約，就是他領我們列祖出埃及地的時候，與他們所立的約。」 \end{tabularx} \\ \\ \relax
8:22 & \begin{tabularx}{0.7\textwidth}{X} 所羅門當著以色列全會眾，站在耶和華的壇前，向天舉手， \end{tabularx} \\ \\ \relax
8:23 & \begin{tabularx}{0.7\textwidth}{X} 說：「耶和華－以色列的神啊，天上地下沒有神明可與你相比！你向那些盡心行在你面前的僕人守約施慈愛， \end{tabularx} \\ \\ \relax
8:24 & \begin{tabularx}{0.7\textwidth}{X} 這約是你向你僕人大衛守的，是你應許他的。你親口應許，親手成就，正如今日一樣。 \end{tabularx} \\ \\ \relax
8:25 & \begin{tabularx}{0.7\textwidth}{X} 耶和華－以色列的神啊，你向你僕人我父大衛應許說：『你的子孫若謹慎自己的行為，在我面前行事像你所行的一樣，就不斷有人在我面前坐以色列的王位。』現在求你信守這話。 \end{tabularx} \\ \\ \relax
8:26 & \begin{tabularx}{0.7\textwidth}{X} 以色列的神啊，現在求你成就向你僕人我父大衛所應許的話。 \end{tabularx} \\ \\ \relax
8:27 & \begin{tabularx}{0.7\textwidth}{X} 「神果真住在地上嗎？看哪，天和天上的天尚且不足容納你，何況我所建的這殿呢？ \end{tabularx} \\ \\ \relax
8:28 & \begin{tabularx}{0.7\textwidth}{X} 惟求耶和華－我的神垂顧僕人的禱告祈求，俯聽僕人今日在你面前的祈禱呼求。 \end{tabularx} \\ \\ \relax
8:29 & \begin{tabularx}{0.7\textwidth}{X} 願你的眼目晝夜看顧這殿，就是你說要作為你名的居所；求你垂聽禱告，你僕人向此處的禱告。 \end{tabularx} \\ \\ \relax
8:30 & \begin{tabularx}{0.7\textwidth}{X} 你僕人和你百姓以色列向此處祈禱的時候，求你在你天上的居所垂聽，垂聽而赦免。 \end{tabularx} \\ \\ \relax
8:31 & \begin{tabularx}{0.7\textwidth}{X} 「人若得罪鄰舍，有人強迫他，要他起誓，他來到這殿，在你的壇前起誓， \end{tabularx} \\ \\ \relax
8:32 & \begin{tabularx}{0.7\textwidth}{X} 求你在天上垂聽、處理，向你的僕人施行審判，定惡人有罪，照他所行的報應在他頭上；定義人為義，照他的義賞賜他。 \end{tabularx} \\ \\ \relax
8:33 & \begin{tabularx}{0.7\textwidth}{X} 「你的百姓以色列若得罪你，敗在仇敵面前，卻又歸向你，宣認你的名，在這殿裡向你祈求禱告， \end{tabularx} \\ \\ \relax
8:34 & \begin{tabularx}{0.7\textwidth}{X} 求你在天上垂聽，赦免你百姓以色列的罪，使他們歸回你賜給他們列祖的地。 \end{tabularx} \\ \\ \relax
8:35 & \begin{tabularx}{0.7\textwidth}{X} 「你的百姓若得罪了你，你使天閉塞不下雨；他們若向此處禱告，宣認你的名，因你的懲罰而離開他們的罪， \end{tabularx} \\ \\ \relax
8:36 & \begin{tabularx}{0.7\textwidth}{X} 求你在天上垂聽，赦免你僕人你百姓以色列的罪，將當行的善道教導他們，並降雨在你的地，就是你賜給你百姓為業之地。 \end{tabularx} \\ \\ \relax
8:37 & \begin{tabularx}{0.7\textwidth}{X} 「這地若有饑荒、瘟疫、焚風、霉爛、蝗蟲、螞蚱，或有仇敵圍困這地的城門，無論遭遇甚麼災禍疾病， \end{tabularx} \\ \\ \relax
8:38 & \begin{tabularx}{0.7\textwidth}{X} 你的百姓以色列，或眾人或一人，內心知道有禍，向這殿舉手，無論祈求甚麼，禱告甚麼， \end{tabularx} \\ \\ \relax
8:39 & \begin{tabularx}{0.7\textwidth}{X} 求你在天上你的居所垂聽、赦免、處理。因為你知道人心，惟有你知道世人的心，求你照各人所行的一切待他們， \end{tabularx} \\ \\ \relax
8:40 & \begin{tabularx}{0.7\textwidth}{X} 使他們在你賜給我們列祖的土地上一生一世敬畏你。 \end{tabularx} \\ \\ \relax
8:41 & \begin{tabularx}{0.7\textwidth}{X} 「論到不屬你百姓以色列的外邦人，若為你的名從遠方而來， \end{tabularx} \\ \\ \relax
8:42 & \begin{tabularx}{0.7\textwidth}{X} 他們因聽見你的大名和大能的手，以及伸出來的膀臂，來向這殿禱告， \end{tabularx} \\ \\ \relax
8:43 & \begin{tabularx}{0.7\textwidth}{X} 求你在天上你的居所垂聽，照著外邦人向你所求的一切而行，使地上萬民都認識你的名，敬畏你，像你的百姓以色列一樣，又使他們知道我所建造的是稱為你名下的殿。 \end{tabularx} \\ \\ \relax
8:44 & \begin{tabularx}{0.7\textwidth}{X} 「你的百姓若奉你的派遣出去，無論往何處與仇敵爭戰，他們若向耶和華所選擇的城，以及我為你名所建造的這殿禱告， \end{tabularx} \\ \\ \relax
8:45 & \begin{tabularx}{0.7\textwidth}{X} 求你在天上垂聽他們的禱告祈求，為他們伸張正義。 \end{tabularx} \\ \\ \relax
8:46 & \begin{tabularx}{0.7\textwidth}{X} 「你的百姓若得罪你，因為沒有人不犯罪，你向他們發怒，把他們交在仇敵面前，擄他們的人把他們帶到仇敵之地，或遠或近， \end{tabularx} \\ \\ \relax
8:47 & \begin{tabularx}{0.7\textwidth}{X} 他們若在被擄之地那裡回心轉意，在擄掠者之地悔改，向你懇求說：『我們有罪了，我們悖逆了，我們作惡了』； \end{tabularx} \\ \\ \relax
8:48 & \begin{tabularx}{0.7\textwidth}{X} 他們若在擄他們的仇敵之地盡心盡性歸向你，又向自己的地，就是你賜給他們列祖的地和你所選擇的城，以及我為你名所建造的這殿禱告， \end{tabularx} \\ \\ \relax
8:49 & \begin{tabularx}{0.7\textwidth}{X} 求你在天上你的居所垂聽他們的禱告祈求，為他們伸張正義， \end{tabularx} \\ \\ \relax
8:50 & \begin{tabularx}{0.7\textwidth}{X} 饒恕得罪你的子民，赦免他們向你所犯一切的過犯，使他們在擄他們的人面前蒙憐憫。 \end{tabularx} \\ \\ \relax
8:51 & \begin{tabularx}{0.7\textwidth}{X} 因為他們是你的子民，你的產業，是你從埃及，從鐵爐中領出來的。 \end{tabularx} \\ \\ \relax
8:52 & \begin{tabularx}{0.7\textwidth}{X} 願你的眼目看顧僕人和你百姓以色列的祈求；他們無論何時向你呼求，願你垂聽。 \end{tabularx} \\ \\ \relax
8:53 & \begin{tabularx}{0.7\textwidth}{X} 主耶和華啊，你將他們從地上萬民中分別出來作你的產業，是照著你領我們列祖出埃及的時候，藉你僕人摩西所應許的。」 \end{tabularx} \\ \\ \relax
8:54 & \begin{tabularx}{0.7\textwidth}{X} 所羅門在耶和華的壇前屈膝跪著，向天舉手；他在耶和華面前禱告祈求完畢的時候，就起來， \end{tabularx} \\ \\ \relax
8:55 & \begin{tabularx}{0.7\textwidth}{X} 站著，大聲為以色列全會眾祝福，說： \end{tabularx} \\ \\ \relax
8:56 & \begin{tabularx}{0.7\textwidth}{X} 「耶和華是應當稱頌的！因為他照著一切所應許的賜平安給他的百姓以色列，凡藉他僕人摩西應許賜福的話，一句都沒有落空。 \end{tabularx} \\ \\ \relax
8:57 & \begin{tabularx}{0.7\textwidth}{X} 願耶和華－我們的神與我們同在，像與我們列祖同在一樣，不撇下我們，不丟棄我們， \end{tabularx} \\ \\ \relax
8:58 & \begin{tabularx}{0.7\textwidth}{X} 使我們的心歸向他，遵行他一切的道，謹守他吩咐我們列祖的誡命、律例、典章。 \end{tabularx} \\ \\ \relax
8:59 & \begin{tabularx}{0.7\textwidth}{X} 願我在耶和華面前祈求的這些話，晝夜靠近耶和華－我們的神，好讓他每日為他僕人和他百姓以色列伸張正義， \end{tabularx} \\ \\ \relax
8:60 & \begin{tabularx}{0.7\textwidth}{X} 使地上的萬民都知道惟獨耶和華是神，沒有別的了。 \end{tabularx} \\ \\ \relax
8:61 & \begin{tabularx}{0.7\textwidth}{X} 所以你們當向耶和華－我們的神存純正的心，遵行他的律例，謹守他的誡命，如同今日一樣。」 \end{tabularx} \\ \\ \relax
8:62 & \begin{tabularx}{0.7\textwidth}{X} 王和全以色列一同在耶和華面前獻祭。 \end{tabularx} \\ \\ \relax
8:63 & \begin{tabularx}{0.7\textwidth}{X} 所羅門向耶和華獻平安祭，二萬二千頭牛，十二萬隻羊。這樣，王和全以色列為耶和華的殿行了奉獻之禮。 \end{tabularx} \\ \\ \relax
8:64 & \begin{tabularx}{0.7\textwidth}{X} 當日，王因耶和華殿前的銅壇太小，容不下燔祭、素祭和平安祭牲的脂肪，就將耶和華殿前院子的中間分別為聖，在那裡獻燔祭、素祭和平安祭牲的脂肪。 \end{tabularx} \\ \\ \relax
8:65 & \begin{tabularx}{0.7\textwidth}{X} 那時所羅門守節，從哈馬口直到埃及溪谷的以色列眾人都與他同在一起，成了一個盛大的會，在耶和華－我們的神面前七日又七日，共十四日。 \end{tabularx} \\ \\ \relax
8:66 & \begin{tabularx}{0.7\textwidth}{X} 第八日，王遣散百姓；他們都為王祝福。他們為耶和華向他僕人大衛和他百姓以色列所施的一切恩惠都心中喜樂，愉快地各回自己的帳棚去了。 \end{tabularx} \\ \\ \relax
9:1 & \begin{tabularx}{0.7\textwidth}{X} 所羅門建造耶和華的殿和王宮，以及一切所想要建造的都完畢了， \end{tabularx} \\ \\ \relax
9:2 & \begin{tabularx}{0.7\textwidth}{X} 耶和華第二次向所羅門顯現，如先前在基遍向他顯現一樣。 \end{tabularx} \\ \\ \relax
9:3 & \begin{tabularx}{0.7\textwidth}{X} 耶和華對他說：「我已聽了你在我面前的禱告和祈求，將你所建的這殿分別為聖，使我的名永遠立在那裡；我的眼、我的心也必時常在那裡。 \end{tabularx} \\ \\ \relax
9:4 & \begin{tabularx}{0.7\textwidth}{X} 你若以純正的心和正直行在我面前，效法你父大衛所行的，遵行我一切所吩咐你的，謹守我的律例典章， \end{tabularx} \\ \\ \relax
9:5 & \begin{tabularx}{0.7\textwidth}{X} 我就必堅固你在以色列國度的王位，直到永遠，正如我應許你父大衛說：『你的子孫必不斷有人坐以色列的王位。』 \end{tabularx} \\ \\ \relax
9:6 & \begin{tabularx}{0.7\textwidth}{X} 倘若你們和你們的子孫轉去不跟從我，不守我擺在你們面前的誡命律例，去事奉別神，敬拜它們， \end{tabularx} \\ \\ \relax
9:7 & \begin{tabularx}{0.7\textwidth}{X} 我就必把以色列從我賜給他們的地上剪除，也必從我面前捨棄那為我名所分別為聖的殿，使以色列在萬民中成為笑柄，被人譏誚。 \end{tabularx} \\ \\ \relax
9:8 & \begin{tabularx}{0.7\textwidth}{X} 這殿雖然崇高，將來凡經過的人必驚訝，嗤笑，說：『耶和華為何向這地和這殿如此行呢？』 \end{tabularx} \\ \\ \relax
9:9 & \begin{tabularx}{0.7\textwidth}{X} 人必說：『因為此地的人離棄領他們祖先出埃及地的耶和華－他們的神，去親近別神，敬拜事奉它們，所以耶和華使這一切災禍臨到他們。』」 \end{tabularx} \\ \\ \relax
9:10 & \begin{tabularx}{0.7\textwidth}{X} 所羅門建造耶和華殿和王宮這兩座殿宇，用了二十年才完成。 \end{tabularx} \\ \\ \relax
9:11 & \begin{tabularx}{0.7\textwidth}{X} 推羅王希蘭曾照所羅門所要的資助他香柏木、松木和金子，所羅門王就把加利利地的二十座城給了希蘭。 \end{tabularx} \\ \\ \relax
9:12 & \begin{tabularx}{0.7\textwidth}{X} 希蘭從推羅出來，察看所羅門給他的城鎮，看不順眼， \end{tabularx} \\ \\ \relax
9:13 & \begin{tabularx}{0.7\textwidth}{X} 就說：「我兄啊，你給我的是甚麼城鎮呢？」他就給這些城鎮起名叫迦步勒地，直到今日。 \end{tabularx} \\ \\ \relax
9:14 & \begin{tabularx}{0.7\textwidth}{X} 希蘭曾給所羅門一百二十他連得金子。 \end{tabularx} \\ \\ \relax
9:15 & \begin{tabularx}{0.7\textwidth}{X} 所羅門王挑取服勞役的工人，為要建造耶和華的殿、自己的宮、米羅、耶路撒冷的城牆、夏瑣、米吉多和基色。 \end{tabularx} \\ \\ \relax
9:16 & \begin{tabularx}{0.7\textwidth}{X} 先前埃及王法老上來攻取基色，用火焚燒，殺了城內居住的迦南人，把城賜給他的女兒，就是所羅門的妻子，作為嫁妝。 \end{tabularx} \\ \\ \relax
9:17 & \begin{tabularx}{0.7\textwidth}{X} 所羅門建造基色、下伯‧和崙、 \end{tabularx} \\ \\ \relax
9:18 & \begin{tabularx}{0.7\textwidth}{X} 巴拉，和位於境內曠野的達莫。 \end{tabularx} \\ \\ \relax
9:19 & \begin{tabularx}{0.7\textwidth}{X} 所羅門建造一切的儲貨城、戰車城、戰馬城，以及他所想要建造的，在耶路撒冷、黎巴嫩和自己治理全國中的一切建設。 \end{tabularx} \\ \\ \relax
9:20 & \begin{tabularx}{0.7\textwidth}{X} 至於所有剩下的百姓，不屬以色列人的亞摩利人、赫人、比利洗人、希未人、耶布斯人， \end{tabularx} \\ \\ \relax
9:21 & \begin{tabularx}{0.7\textwidth}{X} 那些以色列人在當地不能滅盡的人，所羅門徵召他們剩下的後代作服勞役的奴僕，直到今日。 \end{tabularx} \\ \\ \relax
9:22 & \begin{tabularx}{0.7\textwidth}{X} 惟有以色列人，所羅門不使他們作奴僕，而是作他的戰士、臣僕、官長、軍官、戰車長、騎兵長。 \end{tabularx} \\ \\ \relax
9:23 & \begin{tabularx}{0.7\textwidth}{X} 這些是所羅門工程的五百五十個監工，他們在百姓中監管作工的人。 \end{tabularx} \\ \\ \relax
9:24 & \begin{tabularx}{0.7\textwidth}{X} 法老的女兒從大衛城上到所羅門為她建造的宮裡。那時，所羅門才建造米羅。 \end{tabularx} \\ \\ \relax
9:25 & \begin{tabularx}{0.7\textwidth}{X} 所羅門每年三次在他為耶和華所築的壇上獻燔祭和平安祭，又在耶和華面前的壇上燒香。這樣，他完成了建殿。 \end{tabularx} \\ \\ \relax
9:26 & \begin{tabularx}{0.7\textwidth}{X} 所羅門王在以東地紅海邊，靠近以祿的以旬‧迦別製造船隻。 \end{tabularx} \\ \\ \relax
9:27 & \begin{tabularx}{0.7\textwidth}{X} 希蘭派他的僕人，就是熟悉航海的船員，與所羅門的僕人一同坐船航海。 \end{tabularx} \\ \\ \relax
9:28 & \begin{tabularx}{0.7\textwidth}{X} 他們到了俄斐，從那裡得了四百二十他連得金子，運到所羅門王那裡。 \end{tabularx} \\ \\ \relax
10:1 & \begin{tabularx}{0.7\textwidth}{X} 示巴女王聽見所羅門因耶和華的名所得的名聲，就來要用難題考問所羅門。 \end{tabularx} \\ \\ \relax
10:2 & \begin{tabularx}{0.7\textwidth}{X} 她帶著很多的隨從來到耶路撒冷，有駱駝馱著香料、極多金子和寶石。她來到所羅門那裡，向他提出心中所有的問題。 \end{tabularx} \\ \\ \relax
10:3 & \begin{tabularx}{0.7\textwidth}{X} 所羅門回答了她所有的問題，沒有一個問題太難，王不能向她解答的。 \end{tabularx} \\ \\ \relax
10:4 & \begin{tabularx}{0.7\textwidth}{X} 示巴女王看見所羅門一切的智慧，和他所建造的宮殿， \end{tabularx} \\ \\ \relax
10:5 & \begin{tabularx}{0.7\textwidth}{X} 席上的食物，坐著的群臣，侍立的僕人，他們的服裝，和他的司酒長，以及他在耶和華殿裡所獻的燔祭，就詫異得神不守舍。 \end{tabularx} \\ \\ \relax
10:6 & \begin{tabularx}{0.7\textwidth}{X} 她對王說：「我在本國所聽到的話，論到你的事和你的智慧是真的！ \end{tabularx} \\ \\ \relax
10:7 & \begin{tabularx}{0.7\textwidth}{X} 我本來不信那些話，及至我來親眼看見了，看哪，人所告訴我的還不到一半，你的智慧和你的福分超過我所聽見的傳聞。 \end{tabularx} \\ \\ \relax
10:8 & \begin{tabularx}{0.7\textwidth}{X} 你的人是有福的！你這些僕人常侍立在你面前、聽你智慧的話是有福的！ \end{tabularx} \\ \\ \relax
10:9 & \begin{tabularx}{0.7\textwidth}{X} 耶和華－你的神是應當稱頌的！他喜愛你，使你坐以色列的王位，因為他永遠愛以色列，所以立你作王，使你秉公行義。」 \end{tabularx} \\ \\ \relax
10:10 & \begin{tabularx}{0.7\textwidth}{X} 於是，示巴女王把一百二十他連得金子、極多的香料和寶石送給所羅門王；送來的香料，從來沒有像示巴女王送給他的那麼多。 \end{tabularx} \\ \\ \relax
10:11 & \begin{tabularx}{0.7\textwidth}{X} 希蘭的船隻也從俄斐運了金子來，又從俄斐運了許多檀香木和寶石來。 \end{tabularx} \\ \\ \relax
10:12 & \begin{tabularx}{0.7\textwidth}{X} 王用檀香木為耶和華的殿和王宮做欄杆，又為歌唱的人做琴瑟。以後再沒有這樣的檀香木運進來，也再沒有人見過，直到如今。 \end{tabularx} \\ \\ \relax
10:13 & \begin{tabularx}{0.7\textwidth}{X} 所羅門王除了照自己的厚意餽贈示巴女王之外，凡她所提出的一切要求，所羅門王都送給她。於是女王和她臣僕轉回，到本國去了。 \end{tabularx} \\ \\ \relax
10:14 & \begin{tabularx}{0.7\textwidth}{X} 所羅門每年所得的金子，重六百六十六他連得； \end{tabularx} \\ \\ \relax
10:15 & \begin{tabularx}{0.7\textwidth}{X} 另外還有來自商人和做生意的商品，以及阿拉伯諸王和各地省長的。 \end{tabularx} \\ \\ \relax
10:16 & \begin{tabularx}{0.7\textwidth}{X} 所羅門王用錘出來的金子打成二百面盾牌，每面盾牌用六百舍客勒金子； \end{tabularx} \\ \\ \relax
10:17 & \begin{tabularx}{0.7\textwidth}{X} 又用錘出來的金子打成三百面小盾牌，每面小盾牌用三彌那金子。王把它們放在黎巴嫩林宮裡。 \end{tabularx} \\ \\ \relax
10:18 & \begin{tabularx}{0.7\textwidth}{X} 王又製造一個大的象牙寶座，包上純金。 \end{tabularx} \\ \\ \relax
10:19 & \begin{tabularx}{0.7\textwidth}{X} 寶座有六層臺階，座的後背是圓的，座位之處兩旁有扶手，靠近扶手有兩隻獅子站立。 \end{tabularx} \\ \\ \relax
10:20 & \begin{tabularx}{0.7\textwidth}{X} 六層臺階上有十二隻獅子站立，分站左邊和右邊；任何國度都沒有這樣做的。 \end{tabularx} \\ \\ \relax
10:21 & \begin{tabularx}{0.7\textwidth}{X} 所羅門王一切的飲器都是金的，黎巴嫩林宮裡所有的器皿都是純金的。在所羅門的日子，銀子算不了甚麼。 \end{tabularx} \\ \\ \relax
10:22 & \begin{tabularx}{0.7\textwidth}{X} 王有他施船隻與希蘭的船隻一同航海，他施船隻每三年一次把金、銀、象牙、猿猴、孔雀運回來。 \end{tabularx} \\ \\ \relax
10:23 & \begin{tabularx}{0.7\textwidth}{X} 所羅門王的財寶與智慧勝過地上的眾王。 \end{tabularx} \\ \\ \relax
10:24 & \begin{tabularx}{0.7\textwidth}{X} 全地都求見所羅門的面，要聽神放在他心裡的智慧。 \end{tabularx} \\ \\ \relax
10:25 & \begin{tabularx}{0.7\textwidth}{X} 他們各帶貢物，就是銀器、金器、衣服、兵器、香料、馬、騾子，每年都有一定的數量。 \end{tabularx} \\ \\ \relax
10:26 & \begin{tabularx}{0.7\textwidth}{X} 所羅門聚集戰車騎兵；他有一千四百輛戰車，一萬二千名騎兵，安置在屯車城，在耶路撒冷的王那裡。 \end{tabularx} \\ \\ \relax
10:27 & \begin{tabularx}{0.7\textwidth}{X} 王在耶路撒冷使銀子多如石頭，香柏木多如謝非拉的桑樹。 \end{tabularx} \\ \\ \relax
10:28 & \begin{tabularx}{0.7\textwidth}{X} 所羅門的馬是從埃及和科威運來的，是王的商人按著定價從科威買來的。 \end{tabularx} \\ \\ \relax
10:29 & \begin{tabularx}{0.7\textwidth}{X} 從埃及進口的戰車，每輛六百舍客勒銀子，馬每匹一百五十舍客勒；赫人眾王和亞蘭諸王的戰車和馬，也是經由他們的手出口的。 \end{tabularx} \\ \\ \relax
11:1 & \begin{tabularx}{0.7\textwidth}{X} 所羅門王在法老的女兒之外，又寵愛許多外邦女子，就是摩押女子、亞捫女子、以東女子、西頓女子、赫人女子。 \end{tabularx} \\ \\ \relax
11:2 & \begin{tabularx}{0.7\textwidth}{X} 論到這些國的人，耶和華曾吩咐以色列人說：「你們不可跟他們通婚，他們也不可跟你們在一起，因為他們一定會誘惑你們的心去隨從他們的神明。」所羅門卻為了愛，緊緊跟從他們。 \end{tabularx} \\ \\ \relax
11:3 & \begin{tabularx}{0.7\textwidth}{X} 所羅門娶七百個公主，三百個妃嬪。這些妻妾誘惑他的心。 \end{tabularx} \\ \\ \relax
11:4 & \begin{tabularx}{0.7\textwidth}{X} 所羅門年老的時候，他的妻妾誘惑他的心去隨從別神，不像他父親大衛以純正的心順服耶和華－他的神。 \end{tabularx} \\ \\ \relax
11:5 & \begin{tabularx}{0.7\textwidth}{X} 所羅門隨從西頓人的女神亞斯她錄和亞捫人可憎的米勒公。 \end{tabularx} \\ \\ \relax
11:6 & \begin{tabularx}{0.7\textwidth}{X} 所羅門行耶和華眼中看為惡的事，不像他父親大衛專心順從耶和華。 \end{tabularx} \\ \\ \relax
11:7 & \begin{tabularx}{0.7\textwidth}{X} 那時，所羅門為摩押可憎的基抹和亞捫人可憎的摩洛，在耶路撒冷對面的山上建造丘壇。 \end{tabularx} \\ \\ \relax
11:8 & \begin{tabularx}{0.7\textwidth}{X} 他為所有的妻妾，就是那些向自己神明燒香獻祭的外邦女子，也是這樣做。 \end{tabularx} \\ \\ \relax
11:9 & \begin{tabularx}{0.7\textwidth}{X} 耶和華向所羅門發怒，因為他的心偏離了向他顯現兩次的耶和華－以色列的神。 \end{tabularx} \\ \\ \relax
11:10 & \begin{tabularx}{0.7\textwidth}{X} 耶和華曾吩咐他這件事，不可隨從別神，他卻沒有遵守耶和華所吩咐的。 \end{tabularx} \\ \\ \relax
11:11 & \begin{tabularx}{0.7\textwidth}{X} 所以耶和華對他說：「你既然是這樣，不遵守我所吩咐你守的約和律例，我必定把國度撕裂離開你，將它賜給你的大臣。 \end{tabularx} \\ \\ \relax
11:12 & \begin{tabularx}{0.7\textwidth}{X} 然而，因你父親大衛的緣故，我不在你的日子行這事，而要從你兒子的手中撕裂這國。 \end{tabularx} \\ \\ \relax
11:13 & \begin{tabularx}{0.7\textwidth}{X} 只是我不撕裂全國，卻要因我僕人大衛和我所選擇的耶路撒冷，保留一個支派給你的兒子。」 \end{tabularx} \\ \\ \relax
11:14 & \begin{tabularx}{0.7\textwidth}{X} 耶和華使以東人哈達興起，作所羅門的敵人；他是以東王的後裔。 \end{tabularx} \\ \\ \relax
11:15 & \begin{tabularx}{0.7\textwidth}{X} 大衛在以東的時候，約押元帥上去埋葬陣亡的人，殺了以東所有的男丁。 \end{tabularx} \\ \\ \relax
11:16 & \begin{tabularx}{0.7\textwidth}{X} 約押和以色列眾人在以東住了六個月，直到把以東的男丁盡都剪除。 \end{tabularx} \\ \\ \relax
11:17 & \begin{tabularx}{0.7\textwidth}{X} 那時哈達還是幼童；他和他父親的臣僕，以及幾個以東人逃往埃及。 \end{tabularx} \\ \\ \relax
11:18 & \begin{tabularx}{0.7\textwidth}{X} 他們從米甸起行，到了巴蘭，再從巴蘭帶著幾個人來到埃及，到埃及王法老那裡。法老給他房屋，吩咐給他糧食，又把地賜給他。 \end{tabularx} \\ \\ \relax
11:19 & \begin{tabularx}{0.7\textwidth}{X} 哈達在法老眼前大蒙恩寵，法老就把王后答比匿的妹妹嫁給他。 \end{tabularx} \\ \\ \relax
11:20 & \begin{tabularx}{0.7\textwidth}{X} 答比匿的妹妹給哈達生了一個兒子，叫基努拔。答比匿使基努拔在法老的宮裡斷奶，基努拔就與法老的眾子一同住在法老的宮裡。 \end{tabularx} \\ \\ \relax
11:21 & \begin{tabularx}{0.7\textwidth}{X} 哈達在埃及聽見大衛與他祖先同睡，約押元帥也死了，就對法老說：「請你讓我走，我要回本國去。」 \end{tabularx} \\ \\ \relax
11:22 & \begin{tabularx}{0.7\textwidth}{X} 法老對他說：「你在我這裡有甚麼缺乏？看哪，你竟想要回你本國去！」他說：「我沒有缺乏甚麼，只是懇求王准我回去。」 \end{tabularx} \\ \\ \relax
11:23 & \begin{tabularx}{0.7\textwidth}{X} 神又使以利亞大的兒子利遜興起，作所羅門的敵人。他曾逃避主人瑣巴王哈大底謝。 \end{tabularx} \\ \\ \relax
11:24 & \begin{tabularx}{0.7\textwidth}{X} 大衛擊殺瑣巴人的時候，利遜召集了一群人，自己作他們的領袖。他們往大馬士革，住在那裡，在大馬士革建立王國。 \end{tabularx} \\ \\ \relax
11:25 & \begin{tabularx}{0.7\textwidth}{X} 所羅門活著的時候，除了哈達為患之外，利遜也作以色列的敵人。他憎恨以色列，作了亞蘭人的王。 \end{tabularx} \\ \\ \relax
11:26 & \begin{tabularx}{0.7\textwidth}{X} 尼八的兒子耶羅波安也舉起手來攻擊王。他是所羅門的臣僕，以法蓮支派的洗利達人；他母親是個寡婦，名叫洗魯阿。 \end{tabularx} \\ \\ \relax
11:27 & \begin{tabularx}{0.7\textwidth}{X} 他舉手攻擊王是因先前所羅門建造米羅，修補他父親大衛城缺口的這件事。 \end{tabularx} \\ \\ \relax
11:28 & \begin{tabularx}{0.7\textwidth}{X} 耶羅波安是個大有才能的人。所羅門見這青年殷勤，就派他監管約瑟家所有服勞役的工人。 \end{tabularx} \\ \\ \relax
11:29 & \begin{tabularx}{0.7\textwidth}{X} 那時，耶羅波安出了耶路撒冷，示羅人亞希雅先知在路上遇見他；亞希雅身上穿著一件新衣。田野中只有他們二人，沒有其他的人。 \end{tabularx} \\ \\ \relax
11:30 & \begin{tabularx}{0.7\textwidth}{X} 亞希雅拿起穿在自己身上的新衣，把它撕成十二片， \end{tabularx} \\ \\ \relax
11:31 & \begin{tabularx}{0.7\textwidth}{X} 對耶羅波安說：「你可以拿十片。耶和華－以色列的神如此說：『看哪，我必從所羅門手裡撕裂這國，把十個支派賜給你。 \end{tabularx} \\ \\ \relax
11:32 & \begin{tabularx}{0.7\textwidth}{X} 我因我僕人大衛和我在以色列眾支派中所選擇的耶路撒冷城的緣故，仍為所羅門留一個支派。 \end{tabularx} \\ \\ \relax
11:33 & \begin{tabularx}{0.7\textwidth}{X} 因為他們離棄我，敬拜西頓人的女神亞斯她錄、摩押的神明基抹和亞捫人的神明米勒公，沒有像他父親大衛一樣遵從我的道，行我眼中看為正的事，守我的律例典章。 \end{tabularx} \\ \\ \relax
11:34 & \begin{tabularx}{0.7\textwidth}{X} 但我不從他手裡奪走整個國家，卻使他在活著的日子作君王，是因我所揀選的僕人大衛遵守我的誡命律例。 \end{tabularx} \\ \\ \relax
11:35 & \begin{tabularx}{0.7\textwidth}{X} 我必從他兒子手裡將王國奪走，賜給你十個支派， \end{tabularx} \\ \\ \relax
11:36 & \begin{tabularx}{0.7\textwidth}{X} 只留一個支派給他的兒子，使我僕人大衛在我所選擇立我名的耶路撒冷城那裡，在我面前常有燈光。 \end{tabularx} \\ \\ \relax
11:37 & \begin{tabularx}{0.7\textwidth}{X} 我選你，使你照你心裡一切所願的作王，成為以色列的王。 \end{tabularx} \\ \\ \relax
11:38 & \begin{tabularx}{0.7\textwidth}{X} 你若聽從我一切所吩咐你的，遵行我的道，行我眼中看為正的事，謹守我的律例誡命，像我僕人大衛所行的，我就與你同在，為你立堅固的家，像我為大衛所立的一樣，將以色列賜給你。 \end{tabularx} \\ \\ \relax
11:39 & \begin{tabularx}{0.7\textwidth}{X} 我必因這事使大衛的後裔遭受患難，但不是永遠的。』」 \end{tabularx} \\ \\ \relax
11:40 & \begin{tabularx}{0.7\textwidth}{X} 所羅門想要殺耶羅波安，耶羅波安起身逃往埃及。他到了埃及王示撒那裡，就住在埃及，直到所羅門死了。 \end{tabularx} \\ \\ \relax
11:41 & \begin{tabularx}{0.7\textwidth}{X} 所羅門其餘的事，凡他所做的和他的智慧，不都寫在《所羅門記》上嗎？ \end{tabularx} \\ \\ \relax
11:42 & \begin{tabularx}{0.7\textwidth}{X} 所羅門在耶路撒冷作全以色列的王四十年。 \end{tabularx} \\ \\ \relax
11:43 & \begin{tabularx}{0.7\textwidth}{X} 所羅門與他祖先同睡，葬在他父親大衛的城裡，他兒子羅波安接續他作王。 \end{tabularx} \\ \\ \relax
12:1 & \begin{tabularx}{0.7\textwidth}{X} 羅波安往示劍去，因以色列眾人都到了示劍，要立他作王。 \end{tabularx} \\ \\ \relax
12:2 & \begin{tabularx}{0.7\textwidth}{X} 尼八的兒子耶羅波安先前躲避所羅門王，逃往埃及，住在那裡。他還在埃及，聽見了這事， \end{tabularx} \\ \\ \relax
12:3 & \begin{tabularx}{0.7\textwidth}{X} 以色列人派人去請他來。耶羅波安就和以色列全會眾來，與羅波安談話，說： \end{tabularx} \\ \\ \relax
12:4 & \begin{tabularx}{0.7\textwidth}{X} 「你父親使我們負重軛，現在求你減輕你父親所加給我們的苦工和重軛，我們就服事你。」 \end{tabularx} \\ \\ \relax
12:5 & \begin{tabularx}{0.7\textwidth}{X} 羅波安對他們說：「你們走吧，過三天再來見我。」百姓就走了。 \end{tabularx} \\ \\ \relax
12:6 & \begin{tabularx}{0.7\textwidth}{X} 羅波安的父親所羅門在世的日子，有侍立在他面前的長者，羅波安王和他們商議，說：「你們出個主意，好把話帶回給這百姓。」 \end{tabularx} \\ \\ \relax
12:7 & \begin{tabularx}{0.7\textwidth}{X} 他們對他說：「現在王若像僕人一樣服事這百姓，用好話回覆他們，他們就永遠作王的僕人了。」 \end{tabularx} \\ \\ \relax
12:8 & \begin{tabularx}{0.7\textwidth}{X} 王不採納長者給他出的主意，卻和那些與他一同長大、在他面前侍立的年輕人商議。 \end{tabularx} \\ \\ \relax
12:9 & \begin{tabularx}{0.7\textwidth}{X} 他對他們說：「這百姓對我說：『你父親使我們負重軛，求你減輕一些。』你們出個甚麼主意，我們好把話帶回給他們。」 \end{tabularx} \\ \\ \relax
12:10 & \begin{tabularx}{0.7\textwidth}{X} 那些與他一同長大的年輕人對他說：「這百姓對王說：『你父親使我們負重軛，求你給我們減輕一些。』王要對他們如此說：『我的小指頭比我父親的腰還粗呢！ \end{tabularx} \\ \\ \relax
12:11 & \begin{tabularx}{0.7\textwidth}{X} 我父親使你們負重軛，現在我必使你們負更重的軛！我父親用鞭子懲罰你們，我要用蠍子懲罰你們！』」 \end{tabularx} \\ \\ \relax
12:12 & \begin{tabularx}{0.7\textwidth}{X} 耶羅波安和眾百姓遵照王所說「你們第三天再來見我」的話，第三天來到羅波安那裡。 \end{tabularx} \\ \\ \relax
12:13 & \begin{tabularx}{0.7\textwidth}{X} 王嚴厲地回答百姓，不採納長者給他出的主意。 \end{tabularx} \\ \\ \relax
12:14 & \begin{tabularx}{0.7\textwidth}{X} 他照著年輕人所出的主意對他們說：「我父親使你們負重軛，我必使你們負更重的軛！我父親用鞭子懲罰你們，我卻要用蠍子懲罰你們！」 \end{tabularx} \\ \\ \relax
12:15 & \begin{tabularx}{0.7\textwidth}{X} 王不依從百姓，因這事件是出於耶和華，為要應驗耶和華藉示羅人亞希雅對尼八的兒子耶羅波安所說的話。 \end{tabularx} \\ \\ \relax
12:16 & \begin{tabularx}{0.7\textwidth}{X} 以色列眾人見王不依從他們，百姓就回話給王，說：「我們在大衛中有甚麼份呢？我們在耶西的兒子中沒有產業！以色列啊，回你的帳棚去吧！大衛啊，現在你顧自己的家吧！」於是，以色列人都回自己的帳棚去了； \end{tabularx} \\ \\ \relax
12:17 & \begin{tabularx}{0.7\textwidth}{X} 至於住猶大城鎮的以色列人，羅波安仍作他們的王。 \end{tabularx} \\ \\ \relax
12:18 & \begin{tabularx}{0.7\textwidth}{X} 羅波安王派監管勞役的亞多蘭去，以色列眾人用石頭打他，他就死了。羅波安王急忙上車，逃回耶路撒冷去了。 \end{tabularx} \\ \\ \relax
12:19 & \begin{tabularx}{0.7\textwidth}{X} 這樣，以色列背叛大衛家，直到今日。 \end{tabularx} \\ \\ \relax
12:20 & \begin{tabularx}{0.7\textwidth}{X} 以色列眾人聽見耶羅波安回來了，就派人去請他到會眾那裡，立他作全以色列的王。除了猶大支派，沒有跟從大衛家的。 \end{tabularx} \\ \\ \relax
12:21 & \begin{tabularx}{0.7\textwidth}{X} 羅波安來到耶路撒冷，召集了猶大全家和便雅憫支派的人共十八萬，都是精選的戰士，要與以色列家打仗，好將王國奪回，歸所羅門的兒子羅波安。 \end{tabularx} \\ \\ \relax
12:22 & \begin{tabularx}{0.7\textwidth}{X} 但神的話臨到神人示瑪雅，說： \end{tabularx} \\ \\ \relax
12:23 & \begin{tabularx}{0.7\textwidth}{X} 「你去告訴所羅門的兒子猶大王羅波安，猶大和便雅憫全家，以及其餘的百姓，說： \end{tabularx} \\ \\ \relax
12:24 & \begin{tabularx}{0.7\textwidth}{X} 『耶和華如此說：你們不可上去與你們的弟兄以色列人打仗。你們各自回家去吧！因為這事是出於我。』」眾人就聽從耶和華的話，遵照耶和華的話回去了。 \end{tabularx} \\ \\ \relax
12:25 & \begin{tabularx}{0.7\textwidth}{X} 耶羅波安在以法蓮山區建了示劍，住在其中，又從示劍出去，建了毗努伊勒。 \end{tabularx} \\ \\ \relax
12:26 & \begin{tabularx}{0.7\textwidth}{X} 耶羅波安心裡說：「現在，這國恐怕仍會歸大衛家； \end{tabularx} \\ \\ \relax
12:27 & \begin{tabularx}{0.7\textwidth}{X} 這百姓若上耶路撒冷去，在耶和華的殿裡獻祭，他們的心必歸向他們的主猶大王羅波安。他們會殺了我，仍歸猶大王羅波安。」 \end{tabularx} \\ \\ \relax
12:28 & \begin{tabularx}{0.7\textwidth}{X} 耶羅波安王就籌劃，鑄造了兩個金牛犢，對眾百姓說：「你們上耶路撒冷去實在夠久了。以色列啊，看哪，這是領你出埃及地的神明。」 \end{tabularx} \\ \\ \relax
12:29 & \begin{tabularx}{0.7\textwidth}{X} 他把一個安置在伯特利，另一個安置在但。 \end{tabularx} \\ \\ \relax
12:30 & \begin{tabularx}{0.7\textwidth}{X} 這事使百姓陷入罪裡，因為他們甚至到但去拜那牛犢。 \end{tabularx} \\ \\ \relax
12:31 & \begin{tabularx}{0.7\textwidth}{X} 耶羅波安在一些丘壇建神殿，立不屬利未人的平民百姓為祭司。 \end{tabularx} \\ \\ \relax
12:32 & \begin{tabularx}{0.7\textwidth}{X} 耶羅波安定八月十五日為節期，像在猶大的節期一樣，自己上壇獻祭。他在伯特利這樣做，向他所鑄的牛犢獻祭，又把他所立丘壇的祭司安置在伯特利。 \end{tabularx} \\ \\ \relax
12:33 & \begin{tabularx}{0.7\textwidth}{X} 他在八月十五日，就是他自己心中所定的月份，在伯特利上到自己所造的祭壇；他為以色列人定了一個節期，親自上壇燒香。 \end{tabularx} \\ \\ \relax
13:1 & \begin{tabularx}{0.7\textwidth}{X} 看哪，有一個神人遵照耶和華的話從猶大來到伯特利。耶羅波安正站在壇旁燒香； \end{tabularx} \\ \\ \relax
13:2 & \begin{tabularx}{0.7\textwidth}{X} 神人遵照耶和華的話向壇呼叫，說：「壇哪，壇哪！耶和華如此說：『看哪，大衛家必生一個兒子，名叫約西亞，他必將在你上面燒香的丘壇祭司，宰殺在你上面，人的骨頭也必燒在你上面。』」 \end{tabularx} \\ \\ \relax
13:3 & \begin{tabularx}{0.7\textwidth}{X} 當日，神人設個預兆，說：「這是耶和華說的預兆：『看哪，這壇必破裂，壇上的灰必傾倒出來。』」 \end{tabularx} \\ \\ \relax
13:4 & \begin{tabularx}{0.7\textwidth}{X} 耶羅波安王聽見神人向伯特利的壇呼叫的話，就從壇上伸手，說：「拿住他！」王向神人所伸的手卻萎縮了，不能彎回。 \end{tabularx} \\ \\ \relax
13:5 & \begin{tabularx}{0.7\textwidth}{X} 壇也破裂了，壇上的灰傾倒出來，正如神人遵照耶和華的話所設的預兆。 \end{tabularx} \\ \\ \relax
13:6 & \begin{tabularx}{0.7\textwidth}{X} 王對神人說：「請你為我禱告，向耶和華－你的神懇求恩惠，使我的手復原。」於是神人向耶和華懇求，王的手就復原了，如平常一樣。 \end{tabularx} \\ \\ \relax
13:7 & \begin{tabularx}{0.7\textwidth}{X} 王對神人說：「請你跟我回宮，讓你恢復心力，我必給你賞賜。」 \end{tabularx} \\ \\ \relax
13:8 & \begin{tabularx}{0.7\textwidth}{X} 神人對王說：「你就是把你一半的王宮給我，我也不跟你進去，也不在這地方吃飯喝水， \end{tabularx} \\ \\ \relax
13:9 & \begin{tabularx}{0.7\textwidth}{X} 因為耶和華的話這樣吩咐我說：『不可吃飯喝水，也不可從你去的原路回來。』」 \end{tabularx} \\ \\ \relax
13:10 & \begin{tabularx}{0.7\textwidth}{X} 於是神人從別的路回去，不從他到伯特利來的原路回去。 \end{tabularx} \\ \\ \relax
13:11 & \begin{tabularx}{0.7\textwidth}{X} 有一個老先知住在伯特利，他的兒子來，把神人當日在伯特利所做的一切事和他向王所說的話，都告訴了父親。 \end{tabularx} \\ \\ \relax
13:12 & \begin{tabularx}{0.7\textwidth}{X} 父親對他們說：「神人從哪條路去了呢？」他的兒子都看到從猶大來的神人所去的路。 \end{tabularx} \\ \\ \relax
13:13 & \begin{tabularx}{0.7\textwidth}{X} 老先知吩咐兒子說：「你們為我備驢。」他們備好了驢，他就騎上， \end{tabularx} \\ \\ \relax
13:14 & \begin{tabularx}{0.7\textwidth}{X} 去追神人，遇見神人坐在橡樹底下，就對他說：「你是不是從猶大來的神人？」他說：「是我。」 \end{tabularx} \\ \\ \relax
13:15 & \begin{tabularx}{0.7\textwidth}{X} 老先知對他說：「請你跟我一起回家吃飯。」 \end{tabularx} \\ \\ \relax
13:16 & \begin{tabularx}{0.7\textwidth}{X} 神人說：「我不能跟你回去，與你同行，也不能在這地方跟你一起吃飯喝水， \end{tabularx} \\ \\ \relax
13:17 & \begin{tabularx}{0.7\textwidth}{X} 因為有耶和華的話吩咐我說：『你在那裡不可吃飯喝水，也不可從你去的原路回來。』」 \end{tabularx} \\ \\ \relax
13:18 & \begin{tabularx}{0.7\textwidth}{X} 老先知對他說：「我也是先知，和你一樣。有天使遵照耶和華的話對我說：『你去帶他一同回你的家，給他吃飯喝水。』」老先知在欺騙他。 \end{tabularx} \\ \\ \relax
13:19 & \begin{tabularx}{0.7\textwidth}{X} 於是神人跟老先知回去，在他家裡吃飯喝水。 \end{tabularx} \\ \\ \relax
13:20 & \begin{tabularx}{0.7\textwidth}{X} 他們坐席的時候，耶和華的話臨到那帶神人回來的先知， \end{tabularx} \\ \\ \relax
13:21 & \begin{tabularx}{0.7\textwidth}{X} 他就對從猶大來的神人宣告說：「耶和華如此說：『你既違背耶和華的指示，不遵守耶和華－你神的命令， \end{tabularx} \\ \\ \relax
13:22 & \begin{tabularx}{0.7\textwidth}{X} 反倒回來，在耶和華禁止你吃飯喝水的地方吃了飯喝了水，因此你的屍體必不得葬在你祖先的墳墓裡。』」 \end{tabularx} \\ \\ \relax
13:23 & \begin{tabularx}{0.7\textwidth}{X} 神人吃喝完了，老先知為他帶回來的先知備驢。 \end{tabularx} \\ \\ \relax
13:24 & \begin{tabularx}{0.7\textwidth}{X} 神人就去了，在路上有隻獅子遇見他，把他咬死。他的屍體倒在路上，驢站在屍體旁邊，獅子也站在屍體旁邊。 \end{tabularx} \\ \\ \relax
13:25 & \begin{tabularx}{0.7\textwidth}{X} 看哪，有人經過，看見屍體倒在路上，獅子站在屍體旁邊，就來到老先知所住的城裡述說這事。 \end{tabularx} \\ \\ \relax
13:26 & \begin{tabularx}{0.7\textwidth}{X} 那帶神人回來的先知聽見了，就說：「這是那違背了耶和華指示的神人，所以耶和華把他交給獅子；獅子撕裂他，咬死他，正如耶和華對他說的話。」 \end{tabularx} \\ \\ \relax
13:27 & \begin{tabularx}{0.7\textwidth}{X} 老先知吩咐他兒子說：「你們為我備驢。」他們就備了驢。 \end{tabularx} \\ \\ \relax
13:28 & \begin{tabularx}{0.7\textwidth}{X} 他去了，發現神人的屍體倒在路上，驢和獅子站在屍體旁邊，獅子卻沒有吃屍體，也沒有撕裂驢。 \end{tabularx} \\ \\ \relax
13:29 & \begin{tabularx}{0.7\textwidth}{X} 老先知把神人的屍體抬起，馱在驢上，帶回自己的城裡，要為他哀哭，為他安葬。 \end{tabularx} \\ \\ \relax
13:30 & \begin{tabularx}{0.7\textwidth}{X} 老先知把屍體葬在自己的墳裡，為他哀哭，說：「哀哉！我的弟兄啊！」 \end{tabularx} \\ \\ \relax
13:31 & \begin{tabularx}{0.7\textwidth}{X} 安葬之後，老先知對他兒子說：「我死了，你們要把我葬在神人所葬的墳裡，使我的屍骨在他的屍骨旁邊， \end{tabularx} \\ \\ \relax
13:32 & \begin{tabularx}{0.7\textwidth}{X} 因為他遵照耶和華的話，指著伯特利的壇和撒瑪利亞各城丘壇神殿所宣告的話必定應驗。」 \end{tabularx} \\ \\ \relax
13:33 & \begin{tabularx}{0.7\textwidth}{X} 這事以後，耶羅波安仍不離開他的惡道，立平民百姓為丘壇的祭司；凡願意的，他都分別為聖，立為丘壇的祭司。 \end{tabularx} \\ \\ \relax
13:34 & \begin{tabularx}{0.7\textwidth}{X} 這事使耶羅波安的家陷入罪裡，甚至他的家被剪除，從地面上消滅了。 \end{tabularx} \\ \\ \relax
14:1 & \begin{tabularx}{0.7\textwidth}{X} 那時，耶羅波安的兒子亞比雅病了。 \end{tabularx} \\ \\ \relax
14:2 & \begin{tabularx}{0.7\textwidth}{X} 耶羅波安對他的妻子說：「你起來改裝，使人認不出你是耶羅波安的妻子。你往示羅去，看哪，那裡有先知亞希雅，他曾告訴我說，你必作這百姓的王。 \end{tabularx} \\ \\ \relax
14:3 & \begin{tabularx}{0.7\textwidth}{X} 現在你手裡要帶十個餅、幾個薄餅和一瓶蜜到他那裡去，他必告訴你，孩子會怎樣。」 \end{tabularx} \\ \\ \relax
14:4 & \begin{tabularx}{0.7\textwidth}{X} 耶羅波安的妻子就照樣做，起身往示羅去，到了亞希雅的家。亞希雅因年紀老邁，兩眼發直，不能看見。 \end{tabularx} \\ \\ \relax
14:5 & \begin{tabularx}{0.7\textwidth}{X} 耶和華對亞希雅說：「看哪，耶羅波安的妻子來問你她兒子的事，因她兒子病了，你當如此如此告訴她。她進來的時候會扮成別的婦人。」 \end{tabularx} \\ \\ \relax
14:6 & \begin{tabularx}{0.7\textwidth}{X} 她剛進門，亞希雅聽見她的腳步聲，就說：「耶羅波安的妻子，進來吧！你為何扮成別的婦人呢？我奉差遣將凶信告訴你。 \end{tabularx} \\ \\ \relax
14:7 & \begin{tabularx}{0.7\textwidth}{X} 你回去告訴耶羅波安說：『耶和華－以色列的神如此說：我從百姓中提拔了你，立你作我百姓以色列的君王， \end{tabularx} \\ \\ \relax
14:8 & \begin{tabularx}{0.7\textwidth}{X} 將大衛家的國撕裂，賜給你，你卻不效法我僕人大衛，遵守我的誡命，全心順從我，行我眼中看為正的事。 \end{tabularx} \\ \\ \relax
14:9 & \begin{tabularx}{0.7\textwidth}{X} 你反倒行惡，比在你之前所有的人更嚴重；你離開了我，為自己立了別神，鑄了偶像，惹我發怒，將我丟在背後。 \end{tabularx} \\ \\ \relax
14:10 & \begin{tabularx}{0.7\textwidth}{X} 因此，看哪，我必使災禍臨到耶羅波安的家，把屬耶羅波安的男丁，無論是奴役的、自由的，都從以色列中剪除。我必除滅耶羅波安的家，如同人掃除糞土，直到消滅。 \end{tabularx} \\ \\ \relax
14:11 & \begin{tabularx}{0.7\textwidth}{X} 凡屬耶羅波安的人，死在城中的必被狗吃，死在田野的必被空中的鳥吃。這是耶和華說的。』 \end{tabularx} \\ \\ \relax
14:12 & \begin{tabularx}{0.7\textwidth}{X} 你起身回家去吧！你的腳一進城，孩子就死了。 \end{tabularx} \\ \\ \relax
14:13 & \begin{tabularx}{0.7\textwidth}{X} 以色列眾人必為他哀哭，為他安葬。凡屬耶羅波安的人，只有他可以葬入墳墓，因為在耶羅波安的家中，只有他向耶和華－以色列的神表現出好的行為。 \end{tabularx} \\ \\ \relax
14:14 & \begin{tabularx}{0.7\textwidth}{X} 耶和華必另立一王治理以色列，這一天，他必剪除耶羅波安的家；甚麼時候呢？現在就是了。 \end{tabularx} \\ \\ \relax
14:15 & \begin{tabularx}{0.7\textwidth}{X} 耶和華必擊打以色列，使他們搖動，像水中的蘆葦一樣，又將他們從耶和華賜給他們列祖的美地上拔出來，分散在大河那邊，因為他們造了亞舍拉，惹耶和華發怒。 \end{tabularx} \\ \\ \relax
14:16 & \begin{tabularx}{0.7\textwidth}{X} 因耶羅波安所犯的罪，又因他使以色列陷入罪裡，耶和華必將以色列交出來。」 \end{tabularx} \\ \\ \relax
14:17 & \begin{tabularx}{0.7\textwidth}{X} 耶羅波安的妻子起身回去，到了得撒，剛到門檻，孩子就死了。 \end{tabularx} \\ \\ \relax
14:18 & \begin{tabularx}{0.7\textwidth}{X} 以色列眾人為他安葬，為他哀哭，正如耶和華藉他僕人亞希雅先知所說的話。 \end{tabularx} \\ \\ \relax
14:19 & \begin{tabularx}{0.7\textwidth}{X} 耶羅波安其餘的事，他怎樣打仗，怎樣作王，看哪，都寫在《以色列諸王記》上。 \end{tabularx} \\ \\ \relax
14:20 & \begin{tabularx}{0.7\textwidth}{X} 耶羅波安作王二十二年，就與他祖先同睡，他兒子拿答接續他作王。 \end{tabularx} \\ \\ \relax
14:21 & \begin{tabularx}{0.7\textwidth}{X} 所羅門的兒子羅波安作猶大王。他登基的時候年四十一歲，在耶路撒冷，就是耶和華從以色列眾支派中所選擇立他名的城，作王十七年。羅波安的母親名叫拿瑪，是亞捫人。 \end{tabularx} \\ \\ \relax
14:22 & \begin{tabularx}{0.7\textwidth}{X} 猶大人行耶和華眼中看為惡的事，以所犯的罪惹動他的妒忌，比他們的祖先所犯的一切更嚴重。 \end{tabularx} \\ \\ \relax
14:23 & \begin{tabularx}{0.7\textwidth}{X} 因為他們在各高岡上，各青翠樹下築丘壇，立柱像和亞舍拉。 \end{tabularx} \\ \\ \relax
14:24 & \begin{tabularx}{0.7\textwidth}{X} 國中也有男的廟妓。他們效法耶和華在以色列人面前所趕出的外邦人，行一切可憎惡的事。 \end{tabularx} \\ \\ \relax
14:25 & \begin{tabularx}{0.7\textwidth}{X} 羅波安王第五年，埃及王示撒上來攻打耶路撒冷， \end{tabularx} \\ \\ \relax
14:26 & \begin{tabularx}{0.7\textwidth}{X} 奪了耶和華殿和王宮裡的寶物，盡都帶走，又奪走所羅門製造的一切金盾牌。 \end{tabularx} \\ \\ \relax
14:27 & \begin{tabularx}{0.7\textwidth}{X} 羅波安王製造銅盾牌代替那些金盾牌，交給看守王宮宮門的護衛長看管。 \end{tabularx} \\ \\ \relax
14:28 & \begin{tabularx}{0.7\textwidth}{X} 每逢王進耶和華的殿，護衛兵就舉起這些盾牌；隨後仍將盾牌送回護衛室。 \end{tabularx} \\ \\ \relax
14:29 & \begin{tabularx}{0.7\textwidth}{X} 羅波安其餘的事，凡他所做的，不都寫在《猶大列王記》上嗎？ \end{tabularx} \\ \\ \relax
14:30 & \begin{tabularx}{0.7\textwidth}{X} 羅波安與耶羅波安時常交戰。 \end{tabularx} \\ \\ \relax
14:31 & \begin{tabularx}{0.7\textwidth}{X} 羅波安與他祖先同睡，與他祖先同葬在大衛城。他母親名叫拿瑪，是亞捫人，他兒子亞比央接續他作王。 \end{tabularx} \\ \\ \relax
15:1 & \begin{tabularx}{0.7\textwidth}{X} 尼八的兒子耶羅波安王十八年，亞比央登基作猶大王， \end{tabularx} \\ \\ \relax
15:2 & \begin{tabularx}{0.7\textwidth}{X} 在耶路撒冷作王三年。他母親名叫瑪迦，是押沙龍的女兒。 \end{tabularx} \\ \\ \relax
15:3 & \begin{tabularx}{0.7\textwidth}{X} 亞比央行他父親從前所犯一切的罪，他的心不像他曾祖父大衛以純正的心順服耶和華－他的神。 \end{tabularx} \\ \\ \relax
15:4 & \begin{tabularx}{0.7\textwidth}{X} 然而耶和華－他的神因大衛的緣故，仍使大衛在耶路撒冷有燈光，立他兒子接續他作王，又堅立耶路撒冷。 \end{tabularx} \\ \\ \relax
15:5 & \begin{tabularx}{0.7\textwidth}{X} 因為大衛除了赫人烏利亞那件事，都行耶和華眼中看為正的事，一生沒有違背耶和華一切所吩咐的。 \end{tabularx} \\ \\ \relax
15:6 & \begin{tabularx}{0.7\textwidth}{X} 羅波安在世的日子常與耶羅波安交戰。 \end{tabularx} \\ \\ \relax
15:7 & \begin{tabularx}{0.7\textwidth}{X} 亞比央其餘的事，凡他所做的，不都寫在《猶大列王記》上嗎？亞比央常與耶羅波安交戰。 \end{tabularx} \\ \\ \relax
15:8 & \begin{tabularx}{0.7\textwidth}{X} 亞比央與他祖先同睡，葬在大衛城，他兒子亞撒接續他作王。 \end{tabularx} \\ \\ \relax
15:9 & \begin{tabularx}{0.7\textwidth}{X} 以色列王耶羅波安第二十年，亞撒登基作猶大王， \end{tabularx} \\ \\ \relax
15:10 & \begin{tabularx}{0.7\textwidth}{X} 在耶路撒冷作王四十一年。他祖母名叫瑪迦，是押沙龍的女兒。 \end{tabularx} \\ \\ \relax
15:11 & \begin{tabularx}{0.7\textwidth}{X} 亞撒效法他的高祖父大衛行耶和華眼中看為正的事， \end{tabularx} \\ \\ \relax
15:12 & \begin{tabularx}{0.7\textwidth}{X} 從國中除去男的廟妓，又除掉他祖先所造的一切偶像。 \end{tabularx} \\ \\ \relax
15:13 & \begin{tabularx}{0.7\textwidth}{X} 他甚至廢了他祖母瑪迦太后的位，因瑪迦造了可憎的亞舍拉。亞撒砍下她的偶像，在汲淪溪邊燒了， \end{tabularx} \\ \\ \relax
15:14 & \begin{tabularx}{0.7\textwidth}{X} 只是丘壇還沒有廢去。亞撒一生向耶和華存純正的心。 \end{tabularx} \\ \\ \relax
15:15 & \begin{tabularx}{0.7\textwidth}{X} 亞撒將他父親所分別為聖與自己所分別為聖的金銀和器皿都奉到耶和華的殿裡。 \end{tabularx} \\ \\ \relax
15:16 & \begin{tabularx}{0.7\textwidth}{X} 亞撒和以色列王巴沙在世的日子常常交戰。 \end{tabularx} \\ \\ \relax
15:17 & \begin{tabularx}{0.7\textwidth}{X} 以色列王巴沙上來攻擊猶大，修築拉瑪，不許人從猶大王亞撒那裡出入。 \end{tabularx} \\ \\ \relax
15:18 & \begin{tabularx}{0.7\textwidth}{X} 於是亞撒把耶和華殿和王宮府庫裡所剩下的金銀都交在他臣僕手中，派他們到住在大馬士革的亞蘭王，就是希旬的孫子，他伯利門的兒子便‧哈達那裡去，說： \end{tabularx} \\ \\ \relax
15:19 & \begin{tabularx}{0.7\textwidth}{X} 「你父曾與我父立約，我與你也要這樣立約。看哪，我把金銀送給你作禮物，請你廢掉你與以色列王巴沙所立的約，使他從我這裡撤退。」 \end{tabularx} \\ \\ \relax
15:20 & \begin{tabularx}{0.7\textwidth}{X} 便‧哈達聽從了亞撒王，就派遣他的軍官去攻打以色列的城鎮，攻下了以雲、但、亞伯‧伯‧瑪迦、全基尼烈、拿弗他利全地。 \end{tabularx} \\ \\ \relax
15:21 & \begin{tabularx}{0.7\textwidth}{X} 巴沙聽見了，就停工不修築拉瑪，仍住在得撒。 \end{tabularx} \\ \\ \relax
15:22 & \begin{tabularx}{0.7\textwidth}{X} 於是亞撒王向猶大眾人宣佈，不准任何人推辭，吩咐他們運走巴沙修築拉瑪所用的石頭和木料。亞撒王用它們來修築便雅憫的迦巴和米斯巴。 \end{tabularx} \\ \\ \relax
15:23 & \begin{tabularx}{0.7\textwidth}{X} 亞撒其餘的事，他英勇的事蹟，凡他所做的，以及他所建築的城鎮，不都寫在《猶大列王記》上嗎？只是亞撒年老的時候患有腳疾。 \end{tabularx} \\ \\ \relax
15:24 & \begin{tabularx}{0.7\textwidth}{X} 亞撒與他祖先同睡，與他祖先同葬在大衛城，他兒子約沙法接續他作王。 \end{tabularx} \\ \\ \relax
15:25 & \begin{tabularx}{0.7\textwidth}{X} 猶大王亞撒第二年，耶羅波安的兒子拿答登基作以色列王二年， \end{tabularx} \\ \\ \relax
15:26 & \begin{tabularx}{0.7\textwidth}{X} 拿答行耶和華眼中看為惡的事，行他父親所行的道，犯他父親使以色列陷入罪裡的那罪。 \end{tabularx} \\ \\ \relax
15:27 & \begin{tabularx}{0.7\textwidth}{X} 以薩迦人亞希雅的兒子巴沙背叛拿答，在非利士人的基比頓殺了他，那時拿答和以色列眾人正圍困基比頓。 \end{tabularx} \\ \\ \relax
15:28 & \begin{tabularx}{0.7\textwidth}{X} 猶大王亞撒第三年，巴沙殺了拿答，篡了他的位。 \end{tabularx} \\ \\ \relax
15:29 & \begin{tabularx}{0.7\textwidth}{X} 巴沙一作王就殺了耶羅波安全家，耶羅波安家凡有氣息的，一個也沒有留下，都殺滅了，正如耶和華藉他僕人示羅人亞希雅所說的話。 \end{tabularx} \\ \\ \relax
15:30 & \begin{tabularx}{0.7\textwidth}{X} 這是因為耶羅波安所犯的罪，他使以色列陷入罪裡，激怒了耶和華－以色列的神。 \end{tabularx} \\ \\ \relax
15:31 & \begin{tabularx}{0.7\textwidth}{X} 拿答其餘的事，凡他所做的，不都寫在《以色列諸王記》上嗎？ \end{tabularx} \\ \\ \relax
15:32 & \begin{tabularx}{0.7\textwidth}{X} 亞撒和以色列王巴沙在世的日子常常交戰。 \end{tabularx} \\ \\ \relax
15:33 & \begin{tabularx}{0.7\textwidth}{X} 猶大王亞撒第三年，亞希雅的兒子巴沙在得撒登基，作全以色列的王二十四年。 \end{tabularx} \\ \\ \relax
15:34 & \begin{tabularx}{0.7\textwidth}{X} 他行耶和華眼中看為惡的事，行耶羅波安所行的道，犯他使以色列陷入罪裡的那罪。 \end{tabularx} \\ \\ \relax
16:1 & \begin{tabularx}{0.7\textwidth}{X} 耶和華的話臨到哈拿尼的兒子耶戶，責備巴沙說： \end{tabularx} \\ \\ \relax
16:2 & \begin{tabularx}{0.7\textwidth}{X} 「我既從塵埃中提拔你，立你作我百姓以色列的君王，你竟行耶羅波安所行的道，使我的百姓以色列陷入罪裡，以他們的罪惹我發怒， \end{tabularx} \\ \\ \relax
16:3 & \begin{tabularx}{0.7\textwidth}{X} 看哪，我必除盡巴沙和他的家，使你的家像尼八的兒子耶羅波安的家一樣。 \end{tabularx} \\ \\ \relax
16:4 & \begin{tabularx}{0.7\textwidth}{X} 凡屬巴沙的人，死在城中的必被狗吃，死在田野的必被空中的鳥吃。」 \end{tabularx} \\ \\ \relax
16:5 & \begin{tabularx}{0.7\textwidth}{X} 巴沙其餘的事，凡他所做的和他英勇的事蹟，不都寫在《以色列諸王記》上嗎？ \end{tabularx} \\ \\ \relax
16:6 & \begin{tabularx}{0.7\textwidth}{X} 巴沙與他祖先同睡，葬在得撒，他兒子以拉接續他作王。 \end{tabularx} \\ \\ \relax
16:7 & \begin{tabularx}{0.7\textwidth}{X} 耶和華的話臨到哈拿尼的兒子耶戶先知，責備巴沙和他的家，因他行耶和華眼中看為惡的一切事，以他手所做的惹耶和華發怒，像耶羅波安的家一樣，又因他殺了耶羅波安全家。 \end{tabularx} \\ \\ \relax
16:8 & \begin{tabularx}{0.7\textwidth}{X} 猶大王亞撒第二十六年，巴沙的兒子以拉在得撒登基，作以色列王二年。 \end{tabularx} \\ \\ \relax
16:9 & \begin{tabularx}{0.7\textwidth}{X} 他的大臣心利，就是管理他一半戰車的軍官背叛他。當他在得撒，在王宮的管家亞雜家裡喝醉的時候， \end{tabularx} \\ \\ \relax
16:10 & \begin{tabularx}{0.7\textwidth}{X} 心利進去擊殺他，把他殺死，篡了他的位。這是猶大王亞撒第二十七年的事。 \end{tabularx} \\ \\ \relax
16:11 & \begin{tabularx}{0.7\textwidth}{X} 心利一坐上王位就殺了巴沙全家，連他的親屬和朋友，一個男丁也沒有留下。 \end{tabularx} \\ \\ \relax
16:12 & \begin{tabularx}{0.7\textwidth}{X} 心利滅絕巴沙全家，正如耶和華藉耶戶先知責備巴沙的話。 \end{tabularx} \\ \\ \relax
16:13 & \begin{tabularx}{0.7\textwidth}{X} 這是因為巴沙和他兒子以拉的一切罪，就是他們使以色列陷入罪裡的那罪，以虛無的神明惹耶和華－以色列的神發怒。 \end{tabularx} \\ \\ \relax
16:14 & \begin{tabularx}{0.7\textwidth}{X} 以拉其餘的事，凡他所做的，不都寫在《以色列諸王記》上嗎？ \end{tabularx} \\ \\ \relax
16:15 & \begin{tabularx}{0.7\textwidth}{X} 猶大王亞撒第二十七年，心利在得撒作王七日。那時軍兵正安營圍攻非利士人的基比頓。 \end{tabularx} \\ \\ \relax
16:16 & \begin{tabularx}{0.7\textwidth}{X} 軍兵在營中聽說心利已經背叛，殺了王，以色列眾人當日就在營中立暗利元帥作以色列王。 \end{tabularx} \\ \\ \relax
16:17 & \begin{tabularx}{0.7\textwidth}{X} 暗利率領以色列眾人，從基比頓上去，圍困得撒。 \end{tabularx} \\ \\ \relax
16:18 & \begin{tabularx}{0.7\textwidth}{X} 心利見城被攻陷，就進了王宮的堡壘，放火焚燒宮殿，自焚而死。 \end{tabularx} \\ \\ \relax
16:19 & \begin{tabularx}{0.7\textwidth}{X} 這是因為他犯罪，行耶和華眼中看為惡的事，行耶羅波安所行的道，犯他使以色列陷入罪裡的那罪。 \end{tabularx} \\ \\ \relax
16:20 & \begin{tabularx}{0.7\textwidth}{X} 心利其餘的事和他背叛的事，不都寫在《以色列諸王記》上嗎？ \end{tabularx} \\ \\ \relax
16:21 & \begin{tabularx}{0.7\textwidth}{X} 那時，以色列百姓分為兩半：一半隨從基納的兒子提比尼，要擁立他作王；另一半隨從暗利。 \end{tabularx} \\ \\ \relax
16:22 & \begin{tabularx}{0.7\textwidth}{X} 但隨從暗利的百姓勝過隨從基納兒子提比尼的百姓。提比尼死了，暗利就作了王。 \end{tabularx} \\ \\ \relax
16:23 & \begin{tabularx}{0.7\textwidth}{X} 猶大王亞撒第三十一年，暗利登基作以色列王十二年；他在得撒作王六年。 \end{tabularx} \\ \\ \relax
16:24 & \begin{tabularx}{0.7\textwidth}{X} 暗利用二他連得銀子向撒瑪買了撒瑪利亞山，在山上建城，按著山的原主撒瑪的名，給所建的城起名叫撒瑪利亞。 \end{tabularx} \\ \\ \relax
16:25 & \begin{tabularx}{0.7\textwidth}{X} 暗利行耶和華眼中看為惡的事，比他以前所有的王作惡更嚴重。 \end{tabularx} \\ \\ \relax
16:26 & \begin{tabularx}{0.7\textwidth}{X} 因為他行了尼八的兒子耶羅波安所行的道，犯他使以色列陷入罪裡的那罪，以虛無的神明惹耶和華－以色列的神發怒。 \end{tabularx} \\ \\ \relax
16:27 & \begin{tabularx}{0.7\textwidth}{X} 暗利其餘的事，他所做的和所顯出的英勇事蹟，不都寫在《以色列諸王記》上嗎？ \end{tabularx} \\ \\ \relax
16:28 & \begin{tabularx}{0.7\textwidth}{X} 暗利與他祖先同睡，葬在撒瑪利亞，他兒子亞哈接續他作王。 \end{tabularx} \\ \\ \relax
16:29 & \begin{tabularx}{0.7\textwidth}{X} 猶大王亞撒第三十八年，暗利的兒子亞哈登基作以色列王。暗利的兒子亞哈在撒瑪利亞作以色列王二十二年。 \end{tabularx} \\ \\ \relax
16:30 & \begin{tabularx}{0.7\textwidth}{X} 暗利的兒子亞哈行耶和華眼中看為惡的事，比他以前所有的王更嚴重。 \end{tabularx} \\ \\ \relax
16:31 & \begin{tabularx}{0.7\textwidth}{X} 他犯了尼八的兒子耶羅波安所犯的罪，還當作是小事，又娶了西頓王謁巴力的女兒耶洗別為妻，去事奉巴力，敬拜它， \end{tabularx} \\ \\ \relax
16:32 & \begin{tabularx}{0.7\textwidth}{X} 又在撒瑪利亞建巴力廟，在廟裡為巴力築壇。 \end{tabularx} \\ \\ \relax
16:33 & \begin{tabularx}{0.7\textwidth}{X} 亞哈又造亞舍拉，他所做的惹耶和華－以色列的神發怒，比他以前所有的以色列王更嚴重。 \end{tabularx} \\ \\ \relax
16:34 & \begin{tabularx}{0.7\textwidth}{X} 亞哈的日子，伯特利人希伊勒重修耶利哥。立根基的時候，他喪了長子亞比蘭；安門的時候，他喪了幼子西割，正如耶和華藉嫩的兒子約書亞所說的話。 \end{tabularx} \\ \\ \relax
17:1 & \begin{tabularx}{0.7\textwidth}{X} 住在基列的提斯比人以利亞對亞哈說：「我指著所事奉永生的耶和華－以色列的神起誓，這幾年我若不禱告，必不降露水，也不下雨。」 \end{tabularx} \\ \\ \relax
17:2 & \begin{tabularx}{0.7\textwidth}{X} 耶和華的話臨到以利亞，說： \end{tabularx} \\ \\ \relax
17:3 & \begin{tabularx}{0.7\textwidth}{X} 「你離開這裡往東去，躲在約旦河東邊的基立溪旁。 \end{tabularx} \\ \\ \relax
17:4 & \begin{tabularx}{0.7\textwidth}{X} 你要喝那溪裡的水，我已吩咐烏鴉在那裡供養你。」 \end{tabularx} \\ \\ \relax
17:5 & \begin{tabularx}{0.7\textwidth}{X} 於是以利亞去了，他遵照耶和華的話做，去住在約旦河東的基立溪旁。 \end{tabularx} \\ \\ \relax
17:6 & \begin{tabularx}{0.7\textwidth}{X} 烏鴉早上給他叼餅和肉來，晚上也有餅和肉，他又喝溪裡的水。 \end{tabularx} \\ \\ \relax
17:7 & \begin{tabularx}{0.7\textwidth}{X} 過了些日子溪水乾了，因為雨沒有下在地上。 \end{tabularx} \\ \\ \relax
17:8 & \begin{tabularx}{0.7\textwidth}{X} 耶和華的話臨到他，說： \end{tabularx} \\ \\ \relax
17:9 & \begin{tabularx}{0.7\textwidth}{X} 「你起身到西頓的撒勒法去，住在那裡，看哪，我已吩咐那裡的一個寡婦供養你。」 \end{tabularx} \\ \\ \relax
17:10 & \begin{tabularx}{0.7\textwidth}{X} 以利亞就起身往撒勒法去。他到了城門，看哪，有一個寡婦在那裡撿柴。以利亞呼喚她說：「請你用器皿取點水來給我喝。」 \end{tabularx} \\ \\ \relax
17:11 & \begin{tabularx}{0.7\textwidth}{X} 她去取水的時候，以利亞又呼喚她說：「請你手裡也拿點餅來給我。」 \end{tabularx} \\ \\ \relax
17:12 & \begin{tabularx}{0.7\textwidth}{X} 她說：「我指著永生的耶和華－你的神起誓，我沒有餅，罈內只有一把麵，瓶裡只有一點油。看哪，我去找兩根柴，帶回家為我和我兒子做餅。我們吃了，就等死吧！」 \end{tabularx} \\ \\ \relax
17:13 & \begin{tabularx}{0.7\textwidth}{X} 以利亞對她說：「不要怕！你去照你所說的做吧！只要先為我做一個小餅，拿來給我，然後為你和你的兒子做餅； \end{tabularx} \\ \\ \relax
17:14 & \begin{tabularx}{0.7\textwidth}{X} 因為耶和華－以色列的神如此說：『罈內的麵必不用盡，瓶裡的油必不短缺，直到耶和華使雨降在地上的日子。』」 \end{tabularx} \\ \\ \relax
17:15 & \begin{tabularx}{0.7\textwidth}{X} 婦人就照以利亞的話去做。她和以利亞，以及她家中的人，吃了許多日子。 \end{tabularx} \\ \\ \relax
17:16 & \begin{tabularx}{0.7\textwidth}{X} 罈內的麵果然沒有用盡，瓶裡的油也不短缺，正如耶和華藉以利亞所說的話。 \end{tabularx} \\ \\ \relax
17:17 & \begin{tabularx}{0.7\textwidth}{X} 這事以後，那婦人，就是那家的女主人，她的兒子病了，病得很重，甚至沒有氣息。 \end{tabularx} \\ \\ \relax
17:18 & \begin{tabularx}{0.7\textwidth}{X} 婦人對以利亞說：「神人哪，我跟你有甚麼關係，你竟到我這裡來，使神記起我的罪，以致我的兒子死了呢？」 \end{tabularx} \\ \\ \relax
17:19 & \begin{tabularx}{0.7\textwidth}{X} 以利亞對她說：「把你兒子交給我。」以利亞就從婦人懷中接過孩子來，抱到他所住的頂樓，放在自己的床上。 \end{tabularx} \\ \\ \relax
17:20 & \begin{tabularx}{0.7\textwidth}{X} 他求告耶和華說：「耶和華－我的神啊，我寄居在這寡婦的家裡，你卻降禍於她，使她的兒子死了嗎？」 \end{tabularx} \\ \\ \relax
17:21 & \begin{tabularx}{0.7\textwidth}{X} 以利亞三次伏在孩子的身上，求告耶和華說：「耶和華－我的神啊，求你使這孩子的生命歸回給他吧！」 \end{tabularx} \\ \\ \relax
17:22 & \begin{tabularx}{0.7\textwidth}{X} 耶和華聽了以利亞的呼求，孩子的生命歸回給他，他就活了。 \end{tabularx} \\ \\ \relax
17:23 & \begin{tabularx}{0.7\textwidth}{X} 以利亞把孩子從樓上抱下來，進了房間交給他母親，說：「看，你的兒子活了！」 \end{tabularx} \\ \\ \relax
17:24 & \begin{tabularx}{0.7\textwidth}{X} 婦人對以利亞說：「現在我知道你是神人，耶和華藉你口所說的話是真的。」 \end{tabularx} \\ \\ \relax
18:1 & \begin{tabularx}{0.7\textwidth}{X} 過了許多日子，到了第三年，耶和華的話臨到以利亞，說：「你去，讓亞哈看見你，我要降雨在地面上。」 \end{tabularx} \\ \\ \relax
18:2 & \begin{tabularx}{0.7\textwidth}{X} 以利亞就去，要讓亞哈見到他。那時，撒瑪利亞的饑荒非常嚴重。 \end{tabularx} \\ \\ \relax
18:3 & \begin{tabularx}{0.7\textwidth}{X} 亞哈召來他的管家俄巴底。俄巴底非常敬畏耶和華。 \end{tabularx} \\ \\ \relax
18:4 & \begin{tabularx}{0.7\textwidth}{X} 耶洗別殺耶和華先知的時候，俄巴底把一百個先知藏了，每五十人藏在一個洞裡，拿餅和水供養他們。 \end{tabularx} \\ \\ \relax
18:5 & \begin{tabularx}{0.7\textwidth}{X} 亞哈對俄巴底說：「我們要走遍這地，到一切水泉旁和一切溪邊，或者能找到青草，可以救活馬和騾子，免得喪失一些牲畜。」 \end{tabularx} \\ \\ \relax
18:6 & \begin{tabularx}{0.7\textwidth}{X} 於是二人分地巡查，亞哈獨自走一路，俄巴底獨自走另一路。 \end{tabularx} \\ \\ \relax
18:7 & \begin{tabularx}{0.7\textwidth}{X} 俄巴底在路上時，看哪，以利亞遇見他。俄巴底認出他來，就臉伏於地，說：「你是我主以利亞嗎？」 \end{tabularx} \\ \\ \relax
18:8 & \begin{tabularx}{0.7\textwidth}{X} 以利亞對他說：「我是。你去，告訴你主人說：『看哪，以利亞在這裡。』」 \end{tabularx} \\ \\ \relax
18:9 & \begin{tabularx}{0.7\textwidth}{X} 俄巴底說：「僕人犯了甚麼罪，你竟要把我交在亞哈手裡，使他殺我呢？ \end{tabularx} \\ \\ \relax
18:10 & \begin{tabularx}{0.7\textwidth}{X} 我指著永生的耶和華－你的神起誓，無論哪一邦哪一國，我主都派人去找你。若他們說：『不在這裡』，他就叫那邦那國的人起誓說，他們實在找不到你。 \end{tabularx} \\ \\ \relax
18:11 & \begin{tabularx}{0.7\textwidth}{X} 現在你說：『你去告訴你主人說，看哪，以利亞在這裡』； \end{tabularx} \\ \\ \relax
18:12 & \begin{tabularx}{0.7\textwidth}{X} 恐怕我一離開你，耶和華的靈就把你提到我所不知道的地方去。這樣，我去告訴亞哈，他若找不到你，就必殺我。僕人是自幼敬畏耶和華的。 \end{tabularx} \\ \\ \relax
18:13 & \begin{tabularx}{0.7\textwidth}{X} 耶洗別殺耶和華先知的時候，我把耶和華的一百個先知藏了，每五十人藏在一個洞裡，拿餅和水供養他們，難道沒有人把我做的這事告訴我主嗎？ \end{tabularx} \\ \\ \relax
18:14 & \begin{tabularx}{0.7\textwidth}{X} 現在你說：『你去告訴你主人說，看哪，以利亞在這裡』，他一定會殺我。」 \end{tabularx} \\ \\ \relax
18:15 & \begin{tabularx}{0.7\textwidth}{X} 以利亞說：「我指著所事奉永生的萬軍之耶和華起誓，我今日要讓亞哈見到我。」 \end{tabularx} \\ \\ \relax
18:16 & \begin{tabularx}{0.7\textwidth}{X} 於是俄巴底去迎見亞哈，告訴他這事。亞哈就去見以利亞。 \end{tabularx} \\ \\ \relax
18:17 & \begin{tabularx}{0.7\textwidth}{X} 亞哈見了以利亞，就說：「真的是你嗎？你這使以色列遭殃的人！」 \end{tabularx} \\ \\ \relax
18:18 & \begin{tabularx}{0.7\textwidth}{X} 以利亞說：「使以色列遭殃的不是我，而是你和你的父家，因為你們離棄耶和華的誡命，去隨從巴力。 \end{tabularx} \\ \\ \relax
18:19 & \begin{tabularx}{0.7\textwidth}{X} 現在你要派人去召集以色列眾人，以及耶洗別所供養的四百五十個巴力的先知和四百個亞舍拉的先知，叫他們都上迦密山到我這裡來。」 \end{tabularx} \\ \\ \relax
18:20 & \begin{tabularx}{0.7\textwidth}{X} 亞哈就派人到以色列眾人那裡，召集先知上迦密山。 \end{tabularx} \\ \\ \relax
18:21 & \begin{tabularx}{0.7\textwidth}{X} 以利亞近前來對眾百姓說：「你們心持二意要到幾時呢？如果耶和華是神，就當順從耶和華；如果是巴力，就當順從巴力。」百姓一言不答。 \end{tabularx} \\ \\ \relax
18:22 & \begin{tabularx}{0.7\textwidth}{X} 以利亞對百姓說：「作耶和華先知的只剩下我一個；巴力的先知卻有四百五十人。 \end{tabularx} \\ \\ \relax
18:23 & \begin{tabularx}{0.7\textwidth}{X} 請給我們兩頭牛犢，巴力的先知可以為自己挑選一頭牛犢，切成小塊，放在柴上，不要點火；我也預備一頭牛犢放在柴上，也不點火。 \end{tabularx} \\ \\ \relax
18:24 & \begin{tabularx}{0.7\textwidth}{X} 你們求告你們神明的名，我也求告耶和華的名。那應允禱告降火的就是神。」眾百姓回答說：「好主意。」 \end{tabularx} \\ \\ \relax
18:25 & \begin{tabularx}{0.7\textwidth}{X} 以利亞對巴力的先知說：「因為你們人多，先挑選一頭牛犢，預備好了，求告你們神明的名，卻不要點火。」 \end{tabularx} \\ \\ \relax
18:26 & \begin{tabularx}{0.7\textwidth}{X} 他們把所給他們的牛犢預備好了，從早晨到中午，求告巴力的名說：「巴力啊，求你應允我們！」卻沒有聲音，也沒有回應。他們就在所築的壇四圍蹦跳。 \end{tabularx} \\ \\ \relax
18:27 & \begin{tabularx}{0.7\textwidth}{X} 到了正午，以利亞嘲笑他們，說：「大聲求告吧！因為它是神明，它或許在默想，或許正忙著，或許在路上，或許在睡覺，它該醒過來了。」 \end{tabularx} \\ \\ \relax
18:28 & \begin{tabularx}{0.7\textwidth}{X} 他們大聲求告，按著他們的儀式，用刀槍刺割自己，直到渾身流血。 \end{tabularx} \\ \\ \relax
18:29 & \begin{tabularx}{0.7\textwidth}{X} 中午過去了，他們狂呼亂叫，直到獻晚祭的時候，卻沒有聲音，沒有回應的，也沒有理睬的。 \end{tabularx} \\ \\ \relax
18:30 & \begin{tabularx}{0.7\textwidth}{X} 以利亞對眾百姓說：「你們到我這裡來。」眾百姓就到他那裡，他把那已經毀壞了的耶和華的壇修好。 \end{tabularx} \\ \\ \relax
18:31 & \begin{tabularx}{0.7\textwidth}{X} 以利亞按照雅各子孫支派的數目，取了十二塊石頭；耶和華的話曾臨到雅各，說：「你的名要叫以色列。」 \end{tabularx} \\ \\ \relax
18:32 & \begin{tabularx}{0.7\textwidth}{X} 以利亞用這些石頭為耶和華的名築一座壇，在壇的四圍挖溝，可容納二細亞穀種。 \end{tabularx} \\ \\ \relax
18:33 & \begin{tabularx}{0.7\textwidth}{X} 他又在壇上擺好了柴，把牛犢切成小塊放在柴上，說：「你們用四個桶盛滿水，倒在燔祭和柴上。」 \end{tabularx} \\ \\ \relax
18:34 & \begin{tabularx}{0.7\textwidth}{X} 他又說：「倒第二次。」他們就倒第二次。他又說：「倒第三次。」他們就倒第三次。 \end{tabularx} \\ \\ \relax
18:35 & \begin{tabularx}{0.7\textwidth}{X} 水流到壇的四圍，溝裡也滿了水。 \end{tabularx} \\ \\ \relax
18:36 & \begin{tabularx}{0.7\textwidth}{X} 到了獻晚祭的時候，先知以利亞近前來，說：「耶和華－亞伯拉罕、以撒、以色列的神啊，求你今日使人知道你是以色列的神，我是你的僕人，我遵照你的話做這一切事。 \end{tabularx} \\ \\ \relax
18:37 & \begin{tabularx}{0.7\textwidth}{X} 求你應允我，耶和華啊，應允我，使這百姓知道你－耶和華是神，是你叫他們回心轉意的。」 \end{tabularx} \\ \\ \relax
18:38 & \begin{tabularx}{0.7\textwidth}{X} 於是，耶和華降下火來，燒盡燔祭、木柴、石頭、塵土，又燒乾了溝裡的水。 \end{tabularx} \\ \\ \relax
18:39 & \begin{tabularx}{0.7\textwidth}{X} 眾百姓看見了，就臉伏於地，說：「耶和華是神！耶和華是神！」 \end{tabularx} \\ \\ \relax
18:40 & \begin{tabularx}{0.7\textwidth}{X} 以利亞對他們說：「拿住巴力的先知，不讓任何人逃走！」眾人就拿住他們。以利亞帶他們到基順河邊，在那裡殺了他們。 \end{tabularx} \\ \\ \relax
18:41 & \begin{tabularx}{0.7\textwidth}{X} 以利亞對亞哈說：「你現在可以上去吃喝，因為有暴雨的響聲了。」 \end{tabularx} \\ \\ \relax
18:42 & \begin{tabularx}{0.7\textwidth}{X} 亞哈就上去吃喝。以利亞上了迦密山頂，屈身在地，把臉伏在兩膝之中。 \end{tabularx} \\ \\ \relax
18:43 & \begin{tabularx}{0.7\textwidth}{X} 他對僕人說：「你上去，向海觀看。」僕人就上去觀看，說：「沒有甚麼。」以利亞說：「你再去。」如此七次。 \end{tabularx} \\ \\ \relax
18:44 & \begin{tabularx}{0.7\textwidth}{X} 第七次，僕人說：「看哪，有一小片雲從海裡上來，好像人的手掌那麼大。」以利亞說：「你上去告訴亞哈，當套車下去，免得被雨阻擋。」 \end{tabularx} \\ \\ \relax
18:45 & \begin{tabularx}{0.7\textwidth}{X} 霎時間，天因風雲黑暗，降下大雨。亞哈就坐上車，往耶斯列去了。 \end{tabularx} \\ \\ \relax
18:46 & \begin{tabularx}{0.7\textwidth}{X} 耶和華的手按在以利亞身上，他就束上腰，奔在亞哈前頭，一路到耶斯列。 \end{tabularx} \\ \\ \relax
19:1 & \begin{tabularx}{0.7\textwidth}{X} 亞哈把以利亞一切所做的和他用刀殺眾先知的事都告訴耶洗別。 \end{tabularx} \\ \\ \relax
19:2 & \begin{tabularx}{0.7\textwidth}{X} 耶洗別就派使者到以利亞那裡，說：「明日約這時候，我若不使你的性命像那些人的性命一樣，願神明重重懲罰我。」 \end{tabularx} \\ \\ \relax
19:3 & \begin{tabularx}{0.7\textwidth}{X} 以利亞害怕，就起來逃命，到了猶大的別是巴，把僕人留在那裡。 \end{tabularx} \\ \\ \relax
19:4 & \begin{tabularx}{0.7\textwidth}{X} 他自己在曠野走了一日的路程，來到一棵羅騰樹下，就坐在那裡求死，說：「耶和華啊，現在夠了！求你取我的性命吧，因為我不比我的祖先好。」 \end{tabularx} \\ \\ \relax
19:5 & \begin{tabularx}{0.7\textwidth}{X} 他躺在羅騰樹下睡著了。看哪，有一個天使拍他，對他說：「起來吃吧！」 \end{tabularx} \\ \\ \relax
19:6 & \begin{tabularx}{0.7\textwidth}{X} 他觀看，看哪，頭旁有燒熱的石頭烤的餅和一壺水，他就吃了喝了，又再躺下。 \end{tabularx} \\ \\ \relax
19:7 & \begin{tabularx}{0.7\textwidth}{X} 耶和華的使者回來，第二次拍他，說：「起來吃吧！因為你要走的路很遠。」 \end{tabularx} \\ \\ \relax
19:8 & \begin{tabularx}{0.7\textwidth}{X} 他就起來吃了喝了，仗著這飲食的力走了四十晝夜，到了神的山，就是何烈山。 \end{tabularx} \\ \\ \relax
19:9 & \begin{tabularx}{0.7\textwidth}{X} 他在那裡進了一個洞，在洞中過夜。看哪，耶和華的話臨到他，說：「以利亞，你在這裡做甚麼？」 \end{tabularx} \\ \\ \relax
19:10 & \begin{tabularx}{0.7\textwidth}{X} 他說：「我為耶和華－萬軍之神大發熱心，因為以色列人背棄了你的約，毀壞了你的壇，用刀殺了你的先知，只剩下我一人，他們還要追殺我。」 \end{tabularx} \\ \\ \relax
19:11 & \begin{tabularx}{0.7\textwidth}{X} 耶和華說：「你出來站在山上，在耶和華面前。」看哪，耶和華從那裡經過。在耶和華面前有烈風大作，山崩石裂，耶和華卻不在風中；風後有地震，耶和華也不在其中； \end{tabularx} \\ \\ \relax
19:12 & \begin{tabularx}{0.7\textwidth}{X} 地震後有火，耶和華也不在火中；火以後，有輕微細小的聲音。 \end{tabularx} \\ \\ \relax
19:13 & \begin{tabularx}{0.7\textwidth}{X} 以利亞聽見，就用外衣蒙臉，出來站在洞口。聽啊，有聲音向他說：「以利亞，你在這裡做甚麼？」 \end{tabularx} \\ \\ \relax
19:14 & \begin{tabularx}{0.7\textwidth}{X} 他說：「我為耶和華－萬軍之神大發熱心，因為以色列人背棄了你的約，毀壞了你的壇，用刀殺了你的先知，只剩下我一人，他們還要追殺我。」 \end{tabularx} \\ \\ \relax
19:15 & \begin{tabularx}{0.7\textwidth}{X} 耶和華對他說：「去吧，從原路回去，往大馬士革的曠野去。到了那裡，你要膏哈薛作亞蘭王， \end{tabularx} \\ \\ \relax
19:16 & \begin{tabularx}{0.7\textwidth}{X} 又膏寧示的孫子耶戶作以色列王，並膏亞伯‧米何拉人沙法的兒子以利沙作先知接續你。 \end{tabularx} \\ \\ \relax
19:17 & \begin{tabularx}{0.7\textwidth}{X} 將來逃過哈薛之刀的，必被耶戶所殺；逃過耶戶之刀的，必被以利沙所殺。 \end{tabularx} \\ \\ \relax
19:18 & \begin{tabularx}{0.7\textwidth}{X} 但我在以色列中留下七千人，是未曾向巴力屈膝，未曾親吻巴力的。」 \end{tabularx} \\ \\ \relax
19:19 & \begin{tabularx}{0.7\textwidth}{X} 於是，以利亞離開那裡走了，遇見沙法的兒子以利沙；他正在耕田，在他前頭有十二對牛，自己趕著第十二對。以利亞經過他，把自己的外衣搭在他身上。 \end{tabularx} \\ \\ \relax
19:20 & \begin{tabularx}{0.7\textwidth}{X} 以利沙就離開牛，跑到以利亞那裡，說：「請你讓我先與父母吻別，然後我就跟隨你。」以利亞對他說：「因我對你所做的事，你去吧，然後回來。」 \end{tabularx} \\ \\ \relax
19:21 & \begin{tabularx}{0.7\textwidth}{X} 以利沙離開他回去，宰了一對牛，用套牛的器具煮肉給百姓吃，隨後就起身跟隨以利亞，服事他。 \end{tabularx} \\ \\ \relax
20:1 & \begin{tabularx}{0.7\textwidth}{X} 亞蘭王便‧哈達召集他的全軍，率領三十二個王，帶著馬和戰車，上來圍困撒瑪利亞，要攻打它。 \end{tabularx} \\ \\ \relax
20:2 & \begin{tabularx}{0.7\textwidth}{X} 他派使者進城到以色列王亞哈那裡，對他說：「便‧哈達如此說： \end{tabularx} \\ \\ \relax
20:3 & \begin{tabularx}{0.7\textwidth}{X} 『你的金銀都要歸我，你妻妾兒女中最美的也要歸我。』」 \end{tabularx} \\ \\ \relax
20:4 & \begin{tabularx}{0.7\textwidth}{X} 以色列王回答說：「我主我王啊，就照著你的話，我和我所有的都歸你。」 \end{tabularx} \\ \\ \relax
20:5 & \begin{tabularx}{0.7\textwidth}{X} 使者又來說：「便‧哈達如此說：『我已派人到你那裡，要你把你的金銀、妻妾、兒女都歸我。』 \end{tabularx} \\ \\ \relax
20:6 & \begin{tabularx}{0.7\textwidth}{X} 但明日約在這時候，我還要派臣僕到你那裡，搜查你的家和你僕人的家，你眼中一切所喜愛的都由他們的手拿走。」 \end{tabularx} \\ \\ \relax
20:7 & \begin{tabularx}{0.7\textwidth}{X} 以色列王召了國內所有的長老來，說：「你們要知道，看，這人是來找麻煩的！他派人到我這裡來，要我的妻妾、兒女和金銀，我並沒有拒絕他。」 \end{tabularx} \\ \\ \relax
20:8 & \begin{tabularx}{0.7\textwidth}{X} 所有的長老和眾百姓對王說：「不要聽從他，也不要答應他。」 \end{tabularx} \\ \\ \relax
20:9 & \begin{tabularx}{0.7\textwidth}{X} 以色列王對便‧哈達的使者說：「你們告訴我主我王說：『王頭一次派人向僕人所要的一切，僕人都依從，但這事我不能依從。』」使者就去回覆便‧哈達。 \end{tabularx} \\ \\ \relax
20:10 & \begin{tabularx}{0.7\textwidth}{X} 便‧哈達又派人到亞哈那裡，說：「撒瑪利亞的塵土若足夠跟從我的軍兵每人手拿一把，願神明重重懲罰我！」 \end{tabularx} \\ \\ \relax
20:11 & \begin{tabularx}{0.7\textwidth}{X} 以色列王回答說：「你們告訴他說，『剛束上腰帶的，不要像已卸下的那樣誇口。』」 \end{tabularx} \\ \\ \relax
20:12 & \begin{tabularx}{0.7\textwidth}{X} 便‧哈達和諸王正在帳幕裡喝酒，聽見這話，就對他臣僕說：「擺陣吧！」他們就擺陣攻城。 \end{tabularx} \\ \\ \relax
20:13 & \begin{tabularx}{0.7\textwidth}{X} 看哪，一個先知靠近以色列王亞哈，說：「耶和華如此說：『這一大群人你看見了嗎？看哪，今日我必把他們交在你手裡，你就知道我是耶和華。』」 \end{tabularx} \\ \\ \relax
20:14 & \begin{tabularx}{0.7\textwidth}{X} 亞哈說：「藉著誰呢？」他說：「耶和華如此說：『藉著跟從省長的年輕人。』」亞哈說：「誰要開戰呢？」他說：「你！」 \end{tabularx} \\ \\ \relax
20:15 & \begin{tabularx}{0.7\textwidth}{X} 於是亞哈數點跟從省長的年輕人，共二百三十二名，然後又數點以色列的眾軍兵，共七千名。 \end{tabularx} \\ \\ \relax
20:16 & \begin{tabularx}{0.7\textwidth}{X} 中午，他們出了城；便‧哈達和幫助他的三十二個王正在帳幕裡暢飲。 \end{tabularx} \\ \\ \relax
20:17 & \begin{tabularx}{0.7\textwidth}{X} 跟從省長的年輕人先出城。便‧哈達派人去，他們回報說：「有人從撒瑪利亞出來了。」 \end{tabularx} \\ \\ \relax
20:18 & \begin{tabularx}{0.7\textwidth}{X} 他說：「他們若為求和出來，要活捉他們，若為打仗出來，也要活捉他們。」 \end{tabularx} \\ \\ \relax
20:19 & \begin{tabularx}{0.7\textwidth}{X} 跟從省長的年輕人，和跟隨他們的軍兵，都出了城， \end{tabularx} \\ \\ \relax
20:20 & \begin{tabularx}{0.7\textwidth}{X} 各人遇見敵人就擊殺。亞蘭人逃跑，以色列人追趕他們；亞蘭王便‧哈達騎著馬和騎兵一同逃跑。 \end{tabularx} \\ \\ \relax
20:21 & \begin{tabularx}{0.7\textwidth}{X} 以色列王出城攻擊馬和戰車，大大擊殺亞蘭人。 \end{tabularx} \\ \\ \relax
20:22 & \begin{tabularx}{0.7\textwidth}{X} 那先知靠近以色列王，對他說：「去吧，你當自強，看清楚，也要知道你所要做的事，因為再過一年，亞蘭王會上來攻擊你。」 \end{tabularx} \\ \\ \relax
20:23 & \begin{tabularx}{0.7\textwidth}{X} 亞蘭王的臣僕對他說：「他們的神是山神，所以他們勝過我們。但在平原與他們打仗，我們一定勝過他們。 \end{tabularx} \\ \\ \relax
20:24 & \begin{tabularx}{0.7\textwidth}{X} 王當做這樣的事，把諸王革去，派軍官代替他們， \end{tabularx} \\ \\ \relax
20:25 & \begin{tabularx}{0.7\textwidth}{X} 又照著王喪失軍兵的數目，再招募一支軍隊，馬補馬，車補車。然後在平原與他們打仗，我們一定勝過他們。」王就聽臣僕的話，照樣去做。 \end{tabularx} \\ \\ \relax
20:26 & \begin{tabularx}{0.7\textwidth}{X} 過了一年，便‧哈達果然召集亞蘭人上亞弗去，要與以色列人打仗。 \end{tabularx} \\ \\ \relax
20:27 & \begin{tabularx}{0.7\textwidth}{X} 以色列人也召集軍兵，預備食物，去迎戰亞蘭人。 以色列人對著他們安營，好像兩小群的山羊；亞蘭人卻佈滿了地面。 \end{tabularx} \\ \\ \relax
20:28 & \begin{tabularx}{0.7\textwidth}{X} 有神人靠近，對以色列王說：「耶和華如此說：『亞蘭人既說我－耶和華是山神，不是平原之神，我必將這一大群人全都交在你手中，你們就知道我是耶和華。』」 \end{tabularx} \\ \\ \relax
20:29 & \begin{tabularx}{0.7\textwidth}{X} 以色列人與亞蘭人相對安營七日，到第七日兩軍開戰。那一日以色列人殺了亞蘭的十萬步兵， \end{tabularx} \\ \\ \relax
20:30 & \begin{tabularx}{0.7\textwidth}{X} 其餘的都逃向亞弗，到了城裡，城牆倒塌，壓死了剩下的二萬七千人。便‧哈達也逃入城內，藏在嚴密的內室裡。 \end{tabularx} \\ \\ \relax
20:31 & \begin{tabularx}{0.7\textwidth}{X} 他的臣僕對他說：「看哪，我們聽說以色列家的王都是仁慈的王；讓我們腰束麻布，頭套繩索，出去到以色列王那裡，也許他會存留王的性命。」 \end{tabularx} \\ \\ \relax
20:32 & \begin{tabularx}{0.7\textwidth}{X} 於是他們腰束麻布，頭套繩索，來到以色列王那裡，說：「王的僕人便‧哈達說：『求王饒我一命。』」亞哈說：「他還活著嗎？他是我的兄弟。」 \end{tabularx} \\ \\ \relax
20:33 & \begin{tabularx}{0.7\textwidth}{X} 這些人正在探測吉凶，就立即抓住他的話說：「便‧哈達是王的兄弟！」王說：「你們去請他來。」便‧哈達出來到王那裡，王就請他上車。 \end{tabularx} \\ \\ \relax
20:34 & \begin{tabularx}{0.7\textwidth}{X} 便‧哈達對王說：「我父從你父那裡所奪的城鎮，我必歸還給你。你可以在大馬士革為你自己設立街市，像我父在撒瑪利亞所設立的一樣。」亞哈說：「我照此立約，放你回去。」王就與他立約，放了他。 \end{tabularx} \\ \\ \relax
20:35 & \begin{tabularx}{0.7\textwidth}{X} 有一個人是先知的門徒，遵照耶和華的話對他同伴說：「你打我吧！」那人不肯打他。 \end{tabularx} \\ \\ \relax
20:36 & \begin{tabularx}{0.7\textwidth}{X} 他就對那人說：「你既不聽從耶和華的話，看哪，你一離開我，必有獅子咬死你。」那人一離開他，果然遇見獅子，把他咬死了。 \end{tabularx} \\ \\ \relax
20:37 & \begin{tabularx}{0.7\textwidth}{X} 先知的門徒又遇見一個人，對他說：「你打我吧！」那人就打他，把他打傷。 \end{tabularx} \\ \\ \relax
20:38 & \begin{tabularx}{0.7\textwidth}{X} 那先知就去了，用頭巾蒙眼，改了裝，在路旁等候王。 \end{tabularx} \\ \\ \relax
20:39 & \begin{tabularx}{0.7\textwidth}{X} 王從那裡經過，他向王呼叫說：「僕人出戰的時候，看哪，有人轉過來，帶了一個人到我這裡來，說：『你要看守這人，若他真的失蹤了，你的性命必代替他的性命，否則，你就要交出一他連得銀子來。』 \end{tabularx} \\ \\ \relax
20:40 & \begin{tabularx}{0.7\textwidth}{X} 僕人正在到處忙碌的時候，那人就不見了。」以色列王對他說：「你自己決定了，就必照樣判你。」 \end{tabularx} \\ \\ \relax
20:41 & \begin{tabularx}{0.7\textwidth}{X} 他急忙除掉蒙眼的頭巾，以色列王就認出他是一個先知。 \end{tabularx} \\ \\ \relax
20:42 & \begin{tabularx}{0.7\textwidth}{X} 他對王說：「耶和華如此說：『因你把我決定要消滅的人從你手中放走，所以你的命必代替他的命，你的百姓必代替他的百姓。』」 \end{tabularx} \\ \\ \relax
20:43 & \begin{tabularx}{0.7\textwidth}{X} 於是以色列王生氣，憂悶地回撒瑪利亞，到自己的宮去了。 \end{tabularx} \\ \\ \relax
21:1 & \begin{tabularx}{0.7\textwidth}{X} 這些事以後，又有一事。耶斯列人拿伯在耶斯列有一個葡萄園，靠近撒瑪利亞，亞哈王的宮。 \end{tabularx} \\ \\ \relax
21:2 & \begin{tabularx}{0.7\textwidth}{X} 亞哈對拿伯說：「把你的葡萄園給我作菜園，因為它靠近我的宮，我就把更好的葡萄園換給你。你若要銀子，我就按著價錢給你。」 \end{tabularx} \\ \\ \relax
21:3 & \begin{tabularx}{0.7\textwidth}{X} 拿伯對亞哈說：「耶和華不准我把我祖先留下的產業給你。」 \end{tabularx} \\ \\ \relax
21:4 & \begin{tabularx}{0.7\textwidth}{X} 亞哈因耶斯列人拿伯說「我不把我祖先留下的產業給你」，就生氣，憂悶地回宮，躺在床上，臉轉向內，也不吃飯。 \end{tabularx} \\ \\ \relax
21:5 & \begin{tabularx}{0.7\textwidth}{X} 耶洗別王后來對他說：「你為甚麼心裡這樣生氣，不吃飯呢？」 \end{tabularx} \\ \\ \relax
21:6 & \begin{tabularx}{0.7\textwidth}{X} 他對王后說：「我向耶斯列人拿伯說：『把你的葡萄園按價錢賣給我，或是你願意，我可以把別的葡萄園換給你。』他卻說：『我不把我的葡萄園給你。』」 \end{tabularx} \\ \\ \relax
21:7 & \begin{tabularx}{0.7\textwidth}{X} 耶洗別王后對王說：「你現在是不是治理以色列國呢？只管起來，心裡暢暢快快地吃飯，我會把耶斯列人拿伯的葡萄園給你。」 \end{tabularx} \\ \\ \relax
21:8 & \begin{tabularx}{0.7\textwidth}{X} 於是王后以亞哈的名義寫信，蓋上王的印，把信送給那些與拿伯同城居住的長老和貴族。 \end{tabularx} \\ \\ \relax
21:9 & \begin{tabularx}{0.7\textwidth}{X} 她在信上寫著說：「你們當宣告禁食，叫拿伯坐在百姓的高位上， \end{tabularx} \\ \\ \relax
21:10 & \begin{tabularx}{0.7\textwidth}{X} 又叫兩個無賴坐在拿伯對面，作證告他說：『你詛咒了神和王。』然後把他拉出去用石頭打死。」 \end{tabularx} \\ \\ \relax
21:11 & \begin{tabularx}{0.7\textwidth}{X} 那些與拿伯同城居住的長老和貴族，照耶洗別送給他們的信去做。正如她送的信上所寫， \end{tabularx} \\ \\ \relax
21:12 & \begin{tabularx}{0.7\textwidth}{X} 他們宣告禁食，叫拿伯坐在百姓的高位上。 \end{tabularx} \\ \\ \relax
21:13 & \begin{tabularx}{0.7\textwidth}{X} 有兩個無賴來，坐在拿伯對面。無賴當著百姓作證告他說：「拿伯詛咒神和王了！」眾人就把他拉到城外，用石頭打他，他就死了。 \end{tabularx} \\ \\ \relax
21:14 & \begin{tabularx}{0.7\textwidth}{X} 於是他們派人到耶洗別那裡，說：「拿伯被石頭打死了。」 \end{tabularx} \\ \\ \relax
21:15 & \begin{tabularx}{0.7\textwidth}{X} 耶洗別聽見拿伯被石頭打死，就對亞哈說：「你起來，去取得耶斯列人拿伯不肯出價賣給你的葡萄園吧！因為拿伯不在了，他已經死了。」 \end{tabularx} \\ \\ \relax
21:16 & \begin{tabularx}{0.7\textwidth}{X} 亞哈聽見拿伯死了，就起來，下去要取得耶斯列人拿伯的葡萄園。 \end{tabularx} \\ \\ \relax
21:17 & \begin{tabularx}{0.7\textwidth}{X} 耶和華的話臨到提斯比人以利亞，說： \end{tabularx} \\ \\ \relax
21:18 & \begin{tabularx}{0.7\textwidth}{X} 「你起來，去見在撒瑪利亞的以色列王亞哈。看哪，他下去要取得拿伯的葡萄園，他正在那園裡。 \end{tabularx} \\ \\ \relax
21:19 & \begin{tabularx}{0.7\textwidth}{X} 你要對他說：『耶和華如此說：你殺了人，還要取得他的產業嗎？』又要對他說：『耶和華如此說：狗在何處舔拿伯的血，狗也必在何處舔你的血。』」 \end{tabularx} \\ \\ \relax
21:20 & \begin{tabularx}{0.7\textwidth}{X} 亞哈對以利亞說：「我的仇敵啊，你找到我了嗎？」他說：「我找到你了。因為你出賣自己，行了耶和華眼中看為惡的事。 \end{tabularx} \\ \\ \relax
21:21 & \begin{tabularx}{0.7\textwidth}{X} 耶和華說：『看哪，我必使災禍臨到你，把你除滅。以色列中凡屬亞哈的男丁，無論是奴役的、自由的，我都要剪除。 \end{tabularx} \\ \\ \relax
21:22 & \begin{tabularx}{0.7\textwidth}{X} 我必使你的家像尼八的兒子耶羅波安的家，又像亞希雅的兒子巴沙的家，因為你惹我發怒，又使以色列陷入罪裡。』 \end{tabularx} \\ \\ \relax
21:23 & \begin{tabularx}{0.7\textwidth}{X} 論到耶洗別，耶和華說：『狗必在耶斯列的城郭吃耶洗別。 \end{tabularx} \\ \\ \relax
21:24 & \begin{tabularx}{0.7\textwidth}{X} 凡屬亞哈的人，死在城中的必被狗吃，死在田野的必被空中的鳥吃。』」 \end{tabularx} \\ \\ \relax
21:25 & \begin{tabularx}{0.7\textwidth}{X} （只是從來沒有像亞哈的，因他受耶洗別王后的唆使，出賣自己，行了耶和華眼中看為惡的事。 \end{tabularx} \\ \\ \relax
21:26 & \begin{tabularx}{0.7\textwidth}{X} 他行了最可憎的事，隨從偶像，正如耶和華在以色列人面前趕出的亞摩利人所行的一切。） \end{tabularx} \\ \\ \relax
21:27 & \begin{tabularx}{0.7\textwidth}{X} 亞哈聽見這些話，就撕裂衣服，禁食，貼身穿著麻布，也睡在麻布上，沮喪地走來走去。 \end{tabularx} \\ \\ \relax
21:28 & \begin{tabularx}{0.7\textwidth}{X} 耶和華的話臨到提斯比人以利亞，說： \end{tabularx} \\ \\ \relax
21:29 & \begin{tabularx}{0.7\textwidth}{X} 「亞哈在我面前這樣謙卑，你看見了嗎？因為他在我面前謙卑，所以在他的日子，我不降這禍；到他兒子的時候，我必降這禍於他的家。」 \end{tabularx} \\ \\ \relax
22:1 & \begin{tabularx}{0.7\textwidth}{X} 亞蘭和以色列之間連續三年沒有戰爭。 \end{tabularx} \\ \\ \relax
22:2 & \begin{tabularx}{0.7\textwidth}{X} 到了第三年，猶大王約沙法下去見以色列王。 \end{tabularx} \\ \\ \relax
22:3 & \begin{tabularx}{0.7\textwidth}{X} 以色列王對臣僕說：「你們不知道基列的拉末是屬我們的嗎？我們豈可不採取行動，把它從亞蘭王手裡奪回來呢？」 \end{tabularx} \\ \\ \relax
22:4 & \begin{tabularx}{0.7\textwidth}{X} 亞哈問約沙法說：「你肯同我去攻打基列的拉末嗎？」約沙法對以色列王說：「你我不分彼此，我的軍隊就是你的軍隊，我的馬就是你的馬。」 \end{tabularx} \\ \\ \relax
22:5 & \begin{tabularx}{0.7\textwidth}{X} 約沙法對以色列王說：「請你先求問耶和華的話。」 \end{tabularx} \\ \\ \relax
22:6 & \begin{tabularx}{0.7\textwidth}{X} 於是以色列王召集先知，約有四百人，問他們說：「我可以上去攻打基列的拉末嗎？還是不要上去呢？」他們說：「可以上去，因為主必將那城交在王的手裡。」 \end{tabularx} \\ \\ \relax
22:7 & \begin{tabularx}{0.7\textwidth}{X} 約沙法說：「這裡還有沒有耶和華的先知，我們好求問他呢？」 \end{tabularx} \\ \\ \relax
22:8 & \begin{tabularx}{0.7\textwidth}{X} 以色列王對約沙法說：「還有一個人，是音拉的兒子米該雅，我們可以託他求問耶和華。只是我真的很恨他，因為他對我說預言，從不說吉言，總是說凶信。」約沙法說：「請王不要這麼說。」 \end{tabularx} \\ \\ \relax
22:9 & \begin{tabularx}{0.7\textwidth}{X} 以色列王召了一個官員來，說：「你快去，把音拉的兒子米該雅召來。」 \end{tabularx} \\ \\ \relax
22:10 & \begin{tabularx}{0.7\textwidth}{X} 以色列王和猶大王約沙法在撒瑪利亞城門前的禾場，各穿朝服，坐在寶座上，所有的先知都在他們面前說預言。 \end{tabularx} \\ \\ \relax
22:11 & \begin{tabularx}{0.7\textwidth}{X} 基拿拿的兒子西底家造了鐵角，說：「耶和華如此說：『你要用這些角牴觸亞蘭人，直到將他們滅盡。』」 \end{tabularx} \\ \\ \relax
22:12 & \begin{tabularx}{0.7\textwidth}{X} 所有的先知也都這樣預言說：「可以上基列的拉末去，必然得勝，因為耶和華必將那城交在王的手中。」 \end{tabularx} \\ \\ \relax
22:13 & \begin{tabularx}{0.7\textwidth}{X} 那去召米該雅的使者對他說：「看哪，眾先知都異口同聲向王說吉言，你也跟他們說一樣的話，說吉言吧！」 \end{tabularx} \\ \\ \relax
22:14 & \begin{tabularx}{0.7\textwidth}{X} 米該雅說：「我指著永生的耶和華起誓，耶和華向我說甚麼，我就說甚麼。」 \end{tabularx} \\ \\ \relax
22:15 & \begin{tabularx}{0.7\textwidth}{X} 米該雅來到王那裡，王問他：「米該雅，我們可以上去攻打基列的拉末嗎？還是不要上去呢？」他對王說：「你可以上去，必然得勝，耶和華必將那城交在王的手中。」 \end{tabularx} \\ \\ \relax
22:16 & \begin{tabularx}{0.7\textwidth}{X} 王對他說：「我要你發誓多少次，你才會奉耶和華的名向我說實話呢？」 \end{tabularx} \\ \\ \relax
22:17 & \begin{tabularx}{0.7\textwidth}{X} 米該雅說：「我看見以色列眾人散佈在山上，如同沒有牧人的羊群一般。耶和華說：『這些人沒有主人，他們可以平安地各自回家去。』」 \end{tabularx} \\ \\ \relax
22:18 & \begin{tabularx}{0.7\textwidth}{X} 以色列王對約沙法說：「我豈沒有告訴你，這人對我說預言，從不說吉言，只說凶信嗎？」 \end{tabularx} \\ \\ \relax
22:19 & \begin{tabularx}{0.7\textwidth}{X} 米該雅說：「因此你要聽耶和華的話！我看見耶和華坐在寶座上，天上的萬軍侍立在他左右。 \end{tabularx} \\ \\ \relax
22:20 & \begin{tabularx}{0.7\textwidth}{X} 耶和華說：『誰去引誘亞哈上基列的拉末去陣亡呢？』這個這樣說，那個那樣說。 \end{tabularx} \\ \\ \relax
22:21 & \begin{tabularx}{0.7\textwidth}{X} 隨後有一個靈出來，站在耶和華面前，說：『我去引誘他。』 \end{tabularx} \\ \\ \relax
22:22 & \begin{tabularx}{0.7\textwidth}{X} 耶和華問他：『用甚麼方法呢？』他說：『我要出去，在他眾先知的口中成為謊言的靈。』耶和華說：『這樣，你去引誘他，必能成功。你出去，照樣做吧！』 \end{tabularx} \\ \\ \relax
22:23 & \begin{tabularx}{0.7\textwidth}{X} 現在，看哪，耶和華使謊言的靈入了你所有的這些先知的口，並且耶和華已經宣告要降禍於你。」 \end{tabularx} \\ \\ \relax
22:24 & \begin{tabularx}{0.7\textwidth}{X} 基拿拿的兒子西底家前來，打米該雅一巴掌，說：「耶和華的靈從哪裡離開我向你說話呢？」 \end{tabularx} \\ \\ \relax
22:25 & \begin{tabularx}{0.7\textwidth}{X} 米該雅說：「看哪，你進入嚴密的內室躲藏的那日，就必看見。」 \end{tabularx} \\ \\ \relax
22:26 & \begin{tabularx}{0.7\textwidth}{X} 以色列王說：「把米該雅帶走，交回給亞們市長和約阿施王子， \end{tabularx} \\ \\ \relax
22:27 & \begin{tabularx}{0.7\textwidth}{X} 你們要說：『王如此說，把這個人關在監獄裡，使他受苦，吃不飽喝不足，直等到我平安回來。』」 \end{tabularx} \\ \\ \relax
22:28 & \begin{tabularx}{0.7\textwidth}{X} 米該雅說：「你若真的能平安回來，那就是耶和華沒有藉我說這話了。」他又說：「眾百姓啊，你們都要聽！」 \end{tabularx} \\ \\ \relax
22:29 & \begin{tabularx}{0.7\textwidth}{X} 以色列王和猶大王約沙法上基列的拉末去。 \end{tabularx} \\ \\ \relax
22:30 & \begin{tabularx}{0.7\textwidth}{X} 以色列王對約沙法說：「我要改裝上陣，你可以仍穿王袍。」以色列王就改裝上陣去了。 \end{tabularx} \\ \\ \relax
22:31 & \begin{tabularx}{0.7\textwidth}{X} 亞蘭王吩咐他的三十二個戰車長說：「你們不要與他們的大將或小兵交戰，只要單單攻擊以色列王。」 \end{tabularx} \\ \\ \relax
22:32 & \begin{tabularx}{0.7\textwidth}{X} 那些戰車長看見約沙法就說：「這一定是以色列王！」他們轉過去與他交戰，約沙法就呼喊起來。 \end{tabularx} \\ \\ \relax
22:33 & \begin{tabularx}{0.7\textwidth}{X} 戰車長見他不是以色列王，就轉身不追他了。 \end{tabularx} \\ \\ \relax
22:34 & \begin{tabularx}{0.7\textwidth}{X} 有一人開弓，並不知情，箭恰巧射入以色列王鎧甲的縫裡。王對駕車的說：「我受重傷了，你掉過車來，載我離開戰場！」 \end{tabularx} \\ \\ \relax
22:35 & \begin{tabularx}{0.7\textwidth}{X} 那日，戰況越來越猛，有人扶著王站在戰車上，面對亞蘭人。到了傍晚，王就死了，血從傷處流入車底。 \end{tabularx} \\ \\ \relax
22:36 & \begin{tabularx}{0.7\textwidth}{X} 約在日落的時候，有喊聲傳遍軍中，說：「大家各歸本城，各歸本地吧！」 \end{tabularx} \\ \\ \relax
22:37 & \begin{tabularx}{0.7\textwidth}{X} 王死了，人把他送到撒瑪利亞，葬在撒瑪利亞。 \end{tabularx} \\ \\ \relax
22:38 & \begin{tabularx}{0.7\textwidth}{X} 他們在撒瑪利亞的水池旁洗他的車，有狗來舔他的血，有妓女在那裡洗澡，正如耶和華所說的話。 \end{tabularx} \\ \\ \relax
22:39 & \begin{tabularx}{0.7\textwidth}{X} 亞哈其餘的事，凡他所做的、他所修造的象牙宮和所建築的一切城鎮，不都寫在《以色列諸王記》上嗎？ \end{tabularx} \\ \\ \relax
22:40 & \begin{tabularx}{0.7\textwidth}{X} 亞哈與他祖先同睡，他兒子亞哈謝接續他作王。 \end{tabularx} \\ \\ \relax
22:41 & \begin{tabularx}{0.7\textwidth}{X} 以色列王亞哈第四年，亞撒的兒子約沙法登基作猶大王。 \end{tabularx} \\ \\ \relax
22:42 & \begin{tabularx}{0.7\textwidth}{X} 約沙法登基的時候年三十五歲，在耶路撒冷作王二十五年。他母親名叫阿蘇巴，是示利希的女兒。 \end{tabularx} \\ \\ \relax
22:43 & \begin{tabularx}{0.7\textwidth}{X} 約沙法效法他父親亞撒所行的道，不偏離左右，行耶和華眼中看為正的事。只是丘壇還沒有廢去，百姓仍在那裡獻祭燒香。 \end{tabularx} \\ \\ \relax
22:44 & \begin{tabularx}{0.7\textwidth}{X} 約沙法與以色列王和平相處。 \end{tabularx} \\ \\ \relax
22:45 & \begin{tabularx}{0.7\textwidth}{X} 約沙法其餘的事和他所行的英勇事蹟，以及他的戰役，不都寫在《猶大列王記》上嗎？ \end{tabularx} \\ \\ \relax
22:46 & \begin{tabularx}{0.7\textwidth}{X} 約沙法把他父親亞撒的日子所剩下男的廟妓都從國中除去了。 \end{tabularx} \\ \\ \relax
22:47 & \begin{tabularx}{0.7\textwidth}{X} 那時以東沒有立王，由總督治理。 \end{tabularx} \\ \\ \relax
22:48 & \begin{tabularx}{0.7\textwidth}{X} 約沙法造了他施船隻，要往俄斐去，把金子運來，卻沒有啟航，因為船在以旬‧迦別毀壞了。 \end{tabularx} \\ \\ \relax
22:49 & \begin{tabularx}{0.7\textwidth}{X} 亞哈的兒子亞哈謝對約沙法說：「讓我的僕人和你的僕人坐船同去吧！」約沙法卻不肯。 \end{tabularx} \\ \\ \relax
22:50 & \begin{tabularx}{0.7\textwidth}{X} 約沙法與他祖先同睡，與他祖先同葬在大衛城，他兒子約蘭接續他作王。 \end{tabularx} \\ \\ \relax
22:51 & \begin{tabularx}{0.7\textwidth}{X} 猶大王約沙法第十七年，亞哈的兒子亞哈謝在撒瑪利亞登基作以色列王；他作以色列王二年。 \end{tabularx} \\ \\ \relax
22:52 & \begin{tabularx}{0.7\textwidth}{X} 他行耶和華眼中看為惡的事，行他父母的道，又行尼八的兒子耶羅波安的道，使以色列陷入罪裡。 \end{tabularx} \\ \\ \relax
22:53 & \begin{tabularx}{0.7\textwidth}{X} 他事奉巴力，敬拜它，惹耶和華－以色列的神發怒，正如他父親一切所行的。 \end{tabularx} \\ \\
[1ex]
\hline
\hline
\end{longtable}
本書命名列王記,內容乃是論及關乎以色列猶大諸王的史事,有些人以為這些歷史事實非常瑣碎乏味,使人有不耐煩之感!然而我們要從這些歷史書中找出中心教訓,因而得著幫助.列王記上下明顯地告訴我們關乎兩方面的教訓:
\newline
\newline
(1)人與人的方法永遠是失敗的.
\newline
\newline
(2)神一直作主,並沒有離棄以色列子民,神對他們仍有盼望.
\newline
\newline
人是失敗:
\newline
\newline
從所羅門王往下閱讀,直至本書之末,均看到人的失敗,雖然以色列在萬國中被選作神的百姓,我們也明顯地看見有神的教導,同在及律法在他們中間.「會幕」便是一個很好的證明,會幕內有約櫃,其上有神的榮光,表明神與他們相近,又聖殿建成時,神的榮耀充滿聖殿,也清楚顯明了神的同在.除此以外,他們有神的律法,指示他們處理一切的事,甚至對待雀巢的愛心,何謂聖潔或污穢?衛生的問題,在律法內神一一詳言.神教導他們為使他們蒙福.他們有神的同在,可以隨時尋求祂,可惜他們仍是失敗,從他們在曠野飄流四十年的記載中,可清楚看見他們的失敗,十次試探神,惹動神怒.士師時代也是如此背逆,就是那被稱為特別的士師參孫,雖有神特別的能力但仍終失敗.
\newline
\newline
列王記告訴我們:人性就是人性,失敗總是失敗,當我們細察列王記歷史時,雖然某些王似乎不錯,但仍有失敗.例如:烏雅王,他驕倣自大,這驕倣不是向人,乃是向神,神劃定了作王的界限,不許王在聖殿內處理聖事,因這些事另有先知祭司處理,但他竟不遵照神的命令,私自進入聖殿,以致做成嚴重的後果,在以後的日子中,無法再恢復他的王族.
\newline
\newline
又如希西家王,聖經明言在以色列中沒有一個王像他如此盡心依靠神,他靠著神的能力不用一箭一刀,亞述軍隊,一晚間死了十八萬五千人,希西家倚靠神在聖殿中禱告,盡心歸向神,又領導百姓盡心歸向神,雖然他父親阿哈斯是猶大王中最壞的一個,但他卻不受他父親的影響,他如此盡心倚靠主,真是難能可貴.怎料當我們往下閱讀時,我們曉得他仍舊是失敗,原來他得到一病,神說他這病必要死,要他立遺囑,但他卻在神前強求,因而多得十五年壽命,這消息還遠播四方,甚至巴比倫國,差派使者到耶路撒冷把禮物送與希西家,聖經記載希西家喜歡見使者,為何喜歡?此乃驕傲之心,當別人稱頌,送禮物時驕傲便發作了,使者要求觀看國中寶藏,希西家王便把他自有的,國有的,寶庫的金子,銀子,香料,膏油和他武庫的一切軍器,並他所有的財寶都給他們看,其後先知以賽亞代表神斥責希西家說為何如此作?並向他宣告耶和華所定的日子將到,這一切都要被擄到巴比倫,連他所生的後裔都要到巴比倫,這位專心靠主的希西家王也遭失敗,也是因著他的驕倣.
\newline
\newline
希西家王多十五年壽命所生的兒子嗎哪西,完全不行父親的教訓,把外來的罪惡帶進子民中,耶和華清楚指明,就是因著嗎哪西及其子阿們的罪孽,使猶大國位無從挽回,神必定要除淨耶路撒冷的一切罪惡,將罪孽倒出去.直至阿們作王兩年,便有人起來把他殺死.
\newline
\newline
約西亞王八歲登基作王,他行耶和華眼中看為正的事,不偏左右,約西亞可算是一空前的王,他空前盡心盡性歸向神,在父親阿們,祖父嗎哪西大背道後作王,光景在是敗壞不堪,但他仍是盡心去做,除去一切的偶像,殺盡一切異端的祭司,甚至進到以色列國中除去偶像,這王可算是有希望吧!總算能達到神的旨意及目的吧!但他不到四十歲便死了,耶利米為他作哀歌,他死乃因法老王尼哥從埃及上來要攻擊靠近伯拉河的迦基米施,約西亞王出去干涉,法老王派使者見約西亞王說:這次他們是要攻擊與他爭戰之家,而且此事又有神的吩咐,神的同在,約西亞不聽,便改裝上陣,但被弓箭手射中受重傷而死,耶利米先知當時為他作哀歌,為何神不保守他?此也是顯明他的失敗,有些人以為約西亞既聽聞法老王出去爭戰是得神的旨意,他應先到聖殿求問神,昔日神既能使用騾說話,現在也能用異邦王向他表明祂的旨意.一個屬神的人,非靠自己政治的認識,政治的想法,乃是凡事要謙卑求求問主.約書亞也是如此失敗,有些人曾假裝從遠方來,用方法欺騙約書亞,其實他們來自近鄰,是神命定要滅的民族,但約書亞憑自己的見識,眼光來判斷這事,並沒有求問耶和華,因而有錯誤的做法.箴言書明言「要專心仰賴耶和華,不可依靠自己的聰明,在一切的事上都要認識祂,他必指引你的路.」
\newline
\newline
列王記上下兩卷,從所羅門王開始至西底家作王,均著重地表明一件事,就是「人是失敗」的,特別在西底家被擄以後,其侄約雅斤作王,神命定耶利米先知寫明他是沒有子孫,沒有後人,雖然他作王只是幾個月的時間,但神極其不滿意他.
\newline
\newline
神是希望
\newline
\newline
從諸王的史事中,我們看到沒有一個王是合神的心意,可以帶領以色列民達到神的水準,正如耶利米書中所言的:「耶和華要得著萬民的尊敬」,但感謝神,在羅馬書十五12約我們看見神是希望:「但願使人有盼望的神,因信將諸般的喜樂平安充滿你們的心,使你們藉著聖靈的能力,大有盼望」.「使人有盼望的神」,在英譯本解作「希望的神」,神本身便是希望的神,事情雖至破碎不堪,在人及天使眼中看為沒有希望,但在神仍有希望.
\newline
\newline
我們從以色列人出埃及的歷史中可常見這些情形,當他們行至前無去路,後有追兵時,已到絕望的境界,沒有盼望；但神卻說向前行,向海行,神蹟便出現了.甚至耶穌基督的降世也是使沒有希望的人來到主的面前,因著主有希望.耶穌能使死人復活,乃是表明他是有希望的主.有一次,有一個少年的官,問主應作何事才有永生,「做何事得永生」?問題根本便是錯,此不是主的主張,舊的眼光是憑自己守律法,新的方法是靠信,這人守了一切的律法,便以為自己是好人,但主向他解釋得永生是要變賣一切來跟從主,那人貪戀財物,便憂憂愁愁地走了.律法本身沒有希望,唯有認識耶穌才有希望.所以主跟著對他們說,財主進入天國比駱駝穿過針更難,其實財主不能得救,窮人也是不能,得救是在律法以外,故主又說:「在人不能,在神卻凡事都能」.
\newline
\newline
羅馬書告訴我們:「但如今神的義在律法以外顯明出來,因信耶穌加給一切相信的人,並沒有分別」.在人是不能,在神卻沒有不能,只要我們有信心,祂便把自己的兒子白白的賜給我們,所以神是使人有盼望.人是絕望的,但神卻有希望.
\newline
\newline
我們不要單因列王記題及人的失敗而灰心,因為當我們再往下閱讀歷代志及其他經文時,我們便得著希望.雖然大衛本身的子孫沒有希望,但耶西的根必另發生一條枝條,舊約聖經不但一次兩次把樹枝苗裔來形容耶穌；他是大衛的後人生的,但卻不是屬大衛血統的男子所生的,神為大衛興起一位特別的後裔,就是神的兒子耶穌基督,如此使大衛的家有希望,以色列人有希望,世界的人也得著希望.
\newline
\newline
概論
\newline
\newline
以色列人君王的制度是屬人的方法,從撒母耳記上便可見百姓要求撒母耳為他們立一王,撒母耳知道這事不是神所喜悅的就非常難過,但神告訴撒母耳,以色列民從出埃及開始便不要他,他們立王的原因是要有一人領導他們,正如列國一樣,其實在神的百姓,國度中間就是世人的行列,他們把人的方法搬進去,要使以色列國變成列國一樣.
\newline
\newline
這事是神所許可的意旨,不是神所命定的旨意,神要在他們有王的地方中顯出神的作為,他們要求一眼見的王,如列國一般,但他們的方法並不成功,反而亡在異國外邦人的手中,今天的教會也是如此,胡作胡為,神在聖經中把原則,真理,方法,無論是教會行政,個人生活都告訴我們,但人不願意遵照而行,反而越覺乏味,想引用外面世人的方法幫助教會向外人作見證,但這是註定失敗的,唯有用聖經的方法,原則,專心依靠主才有能力.今天香港的教會用外面的方法吸引青年人,其實若用世界的方法一定不能與之敵對,當你用世界方法的時候,世界是主人,那麼一定不能勝過他,例如用跳舞,打麻雀,這些人的方法來,吸引別人到禮拜堂是沒有可能的.
\newline
\newline
今天教會若再用人的方法,不用神的方法,難怪傳福音的「能」就不見了,必遭失敗,因為一切「能」的來源是在乎神,只要把世界的事加進去,「能」便不存在,反而使教會的信仰沖淡.例如英國有一對夫婦到訪某教會,到門前只見寫著:「本教會因XX太太死了便關門.」把聖經的信仰與外面事物相混和,正如砂,穀與米混在一起,不是純正的,教會若不以主的能力,主的靈奶為滿足,把外面的事搬進去,教會的能力必定消失.正如保羅對帖撒羅尼迦信徒說:「不在乎言語,乃在乎聖靈的權能及充足的信.」又對哥林多教會那些引誘信徒離開真道的人說:「我願意知道他們的權能,不是言語.」
\newline
\newline
若教會所用的方法不是神聖靈的方法,聖經的方法,是人的方法,便失去能力,難怪教會無力,信徒軟弱,我們的方法失敗,不少的教會仿照一主教所寫「對神誠實」這書的見解,願意向世人看齊.向世人看齊,為的是怕世人不接納而反對,但如此向他們看齊,世人也是不得滿足,因為世人向我們所要求的,不是稍為有一些妥協便夠,乃是要我們的信仰改至零度.
\newline
\newline
求主幫助我們站立得穩,人是失敗,教會也是失敗,但主永不失敗,我們要行祂的路,用祂的方法.
\newline
\newline


\newpage

\section{第41屆~港九培靈會~胡恩德~第二講}
\label{sec:1479}
\textbf{從神而來的權柄}
\newline
\newline
連結: \href{https://www.hkbibleconference.org/session-message/view/1479}{\texttt{https://www.hkbibleconference.org/session-message/view/1479}}
\newline
\newline
\hyperref[sec:1478]{< < < PREV SERMON < < <}
~
\hyperref[sec:toc]{[返目錄]}
~
\hyperref[sec:1480]{> > > NEXT SERMON > > >}
\newline
\newline
列王記上 1:5-1:6
\newline
\begin{longtable}{cl}
\hline
\hline
章節 & 經文 (和合本修訂版)\\
\hline
1:5 & \begin{tabularx}{0.7\textwidth}{X} 那時，哈及的兒子亞多尼雅妄自尊大，說：「我要作王」，就為自己預備座車、騎兵，又派五十人在他前頭奔跑。 \end{tabularx} \\ \\ \relax
1:6 & \begin{tabularx}{0.7\textwidth}{X} 他父親從來沒有責怪他，說：「你為何這麼做？」他非常俊美，生在押沙龍之後。 \end{tabularx} \\ \\
[1ex]
\hline
\hline
\end{longtable}
關乎列王記上卷的結構及分析:全書二十二章完整地分成兩半,內容如下:
\newline
\newline
(一)第一章至第十一章──國度分裂以前,所羅門王國.
\newline
\newline
(二)第十二章至第二十二章──國度分裂後,南(猶大國)北(以色列國)諸王的時代.
\newline
\newline
由第一章至第十一章經文的要旨:
\newline
\newline
第一章──所羅門王的王權.
\newline
\newline
第二章──所羅門王的公義.
\newline
\newline
第三章──所羅門王的智慧.
\newline
\newline
第四章──所羅門王的豐富.
\newline
\newline
第五章至第九章──所羅門王的建設.
\newline
\newline
第十章──所羅門王的聲譽.
\newline
\newline
第十一章──所羅門王的失足.
\newline
\newline
所羅門作王的權柄是從神而來,從撒母耳記中,我們可以知道神已命定所羅門作王,承繼大衛的王位,在列王記上第一章中看見所羅門的哥哥亞多尼雅要自尊作王,大宴政府中重要的人物及宗教上最高的領袖,所羅門的母親拔示巴從拿單口中聞得此事,覲見大衛王與他理論一番,先知拿單也隨同指證這事,大衛因而起誓要立所羅門作王,然後吩咐祭司撒督及其他人等在基訓立所羅門為王,萬民歡呼吹角,所羅門便坐在大衛的寶座上為王,大衛在床上敬拜稱謝神,因神使他親眼看到有人坐在他的寶座上.
\newline
\newline
所羅門的王權不是出於人乃是出於神,神曾明言他作王,正如昔日大衛作王一樣,現在神已應許所羅門作王,雖然他聽到亞多尼雅自尊作王,但他並沒有用方法,權力為自己奪回王權,乃是等候神的旨意完成,因此他得國的王權非靠己力而得,非靠爭權而來.正如詩篇所言,高舉非從東,從西,從南而來,唯有耶和華定了誰升高,誰降低.恐怕今天不少的基督徒沒有看到這點,倘若我們看到的話,絕不會為自己在教會上的地位,名聲而爭鬥,如果是主所定的,神必叫我們升高,不必自己來爭奪.所羅門之王國如此高,王權如此大,這都不是自己爭回來的,乃是神所安排,我們不要為自己爭取,更不要向神有所強求,否則我們就不認識「高舉」是從神而來,也不認識神是活著的主,管理天地萬物及我們一切的環境,甚至與我們為敵的,也在他管理之內.
\newline
\newline
當然我們不是停下來甚麼也不幹,最重要的還是倚靠神,我們若認識神是管理萬物的主宰,他命所羅門作王,所羅門是不能不作王的；正如大衛作王一樣.神又命耶穌不但作王且作基督,基督的地位,職份比王權更高,是神所賜給人最高的位份,但耶穌作此職份並不是用爭取的方法奪得的,所謂「基督」的位份是有神一切的權柄.所羅門坐在大的寶座上作王,是表明他有大衛作王的一切權柄,照樣耶穌作基督便能行使神一切的權柄.啟示錄描寫那位創造萬物宇宙的神,他的權柄,聖潔,威嚴是極其大,他手中拿著一書卷,內外都寫滿了字,在天上地底下的沒有一人配得展開這書卷,只有一位,就是大衛的根,猶大的獅子能得勝配展書卷,他就是曾被殺過的羔羊耶穌基督,手裡拿著的書是代表神所有的權力,國度,藉著此書授與那接受者的手中,耶穌配受神一切的權柄,豐富,他有這地位及被稱為基督職份,不是靠私人的鬥爭,力量爭取得來的,乃是耶和華所命定,因他完全服在神大能的手中,正如以賽亞書五十三章所形容的:耶和華定意將他壓傷,耶和華又以他為贖罪祭,他服在神的手中,任由神處理,因此腓立比書第二章描寫耶穌本有神的形像,但他謙卑自己成為人,又存心順服,以至於死,且死在十字架上,耶穌乃是完全服在神手中,直至落到最卑微的地步.神的兒子為人嘗死味,擔當了我們一切的刑罰,但至終神必升他為至高,又把一切超乎萬名之上的名賜給他,使萬民無不屈膝,無不口稱耶穌為主.神把這一個如此尊貴的名字給了基督,不是由於他爭取而來,相反地,正因他是完全服在神的手中.
\newline
\newline
所羅門作王與亞多尼雅自立為王兩事完全相反,亞多尼雅自尊說我必作王,而所羅門自始至終卻一言不發,完全讓神藉人為他完成一切的事.
\newline
\newline
偉大的使徒保羅,他得能力的源頭,得勝的秘訣,其中有一個是在哥林多後書第十二章,「因我甚麼時候軟弱,甚麼時候便剛強了.」保羅他為基督的緣故,以軟弱,凌辱等等為可喜樂的,因他甚麼時候軟弱,甚麼時候便剛強了.就是在那些世人以為是軟弱之中,卻顯明了在神前有能力而蒙福,所以得力不在乎我們自己,乃在乎我們肯服在神面前與否?神有旨意叫我們升高,祂就必叫我們升高,我們既是神的兒女,當然有神的旨意使我們升高,因為凡接待他的,就是信他名的人,他就賜他們權柄作神的兒女,神定意施恩與我們如同兒女一般,甚至犧牲自己也把最好的給我們,正如詩篇第八十四篇所說的:神沒有留下一樣好處,不給那些行動正直的人.我們是神的兒子,又是後嗣,羅馬書告訴我們,後嗣是承受一切的產業,神已命定我們與耶穌基督一同作後嗣,預定我們被高舉,我們應當有一種靈裡的高舉,屬靈裡的能力,權柄及榮耀,這並不是世人眼所見的,乃是伏在神前所得到的權柄.
\newline
\newline
浸信會有一位被主使用的牧師,他與神甚親近,也有好些著作,雖然一個靈性如此高超的人,但有時也有難處,因為凡婦人所生的就有難處,況且作為基督徒是神命定要受難的；很奇怪,當他有難處時便找一位名叫摩西的黑人交通,談話後心裡便覺得暢快,摩西他有神的榮耀,光彩及能力,就算一個神所使用的大僕人也要從他身上得安慰,這種使人得安慰的恩賜是好的,但也不能自己強求而得,強求時恐怕也要如亞多尼雅一般,自尊說我要作主,不是耶穌那種柔和,謙卑的心腸,唯有柔和,謙卑才是得力量的方法.
\newline
\newline
亞多尼雅與所羅門是互相對照的,神揀選所羅門王而棄絕亞多尼雅,其中的原因也有多少人的因素,這事與神揀選大衛及廢棄掃羅相仿.神揀選大衛作王時,事先並沒有告訴撒母耳誰是被選立的王,只吩咐他到耶西家中選立一王,當耶西眾子一個一個在他跟前行過時,首先看到耶西之長子,他的體態非常吸引撒母耳,撒母耳以為這人是王,但耶和華告訴他:人是看外貌,但神是看內心.掃羅作王是神順從百姓的意思,眾民因掃羅的外貌及身材便歡喜他作王,但他並不是神所設立及喜悅的王,神從他手中奪回國權,另立大衛為王,並且為他見證,證明他是合神心意的王,凡事遵照神旨而行.
\newline
\newline
亞多尼雅自尊要作王,但所羅門並非如此,神所鑑察的是人的內心.列王記上第三章記載所羅門獻上一千牲畜,神在夢中向他顯現,問他要求甚麼,他便求神給他智慧可以判斷眾多的民,所羅門並不因自己作了王而驕傲,反而知道自己地位,責任重大,又知道工作不太容易處理及判斷,故求神幫助.自高自大的人是不合乎神所使用的,只能使教會產生更多分爭結黨的事.希伯來書十二章保羅勸勉信徒要謹慎,不可失去主的恩典,又恐怕有毒更新出,這些毒根就是驕傲,當毒根在教會生長時,神的靈便不能作工.神選立所羅門也有多少人的因素在內,就是因為他有謙卑及肯仰賴耶和華的心.
\newline
\newline
箴言書第三章第六節如此勸勉我們說:「你要專心仰賴耶和華,不可倚靠自己的聰明,在你一切所行的事上都要認定他,他必指引你的路.」假若今天把我們各人的生平研究一番,究竟我們是屬所羅門還是亞多尼雅呢?很多時候,人的作品免不了是隱惡揚善,但聖經是十分誠實的,它會按我們所行的事,一一寫在神的冊子上.我們究竟是高舉自己還是讓神施恩?
\newline
\newline
我們得救是本乎恩,「恩」的解釋就是神在愛心中,將好處賜與那些不配受的人,今天我們相信耶穌為我們死,接納他作救主,我們便得到所羅門的地位並不是亞多尼雅的地位.聖經告訴我們:世上沒有義人,連一個也沒有,我們要謙卑在神前,用信心無倚無靠地接受主所成就的救恩,我們從前是死在罪惡過犯之中,行事為人隨從現今的風俗,順服空中掌權者的首領和魔鬼,又隨從自己肉體和心中所喜好的去行,包括自己的情慾,本為可怒之子,但今天我們卻得救,這得救是本乎恩,也因著信,並不是出於自己,免得有人自誇,假若我們自誇,便如亞多尼雅自誇一樣,求主幫助我們不以自己為聰明,肯服從他的旨意,恐怕我們開始得好,肯謙卑,但跑下去時便驕傲失敗.
\newline
\newline


\newpage

\section{第41屆~港九培靈會~胡恩德~第三講}
\label{sec:1480}
\textbf{公義的王權}
\newline
\newline
連結: \href{https://www.hkbibleconference.org/session-message/view/1480}{\texttt{https://www.hkbibleconference.org/session-message/view/1480}}
\newline
\newline
\hyperref[sec:1479]{< < < PREV SERMON < < <}
~
\hyperref[sec:toc]{[返目錄]}
~
\hyperref[sec:1481]{> > > NEXT SERMON > > >}
\newline
\newline
列王記上 2:1-2:9
\newline
\begin{longtable}{cl}
\hline
\hline
章節 & 經文 (和合本修訂版)\\
\hline
2:1 & \begin{tabularx}{0.7\textwidth}{X} 大衛的死期臨近了，就吩咐他兒子所羅門說： \end{tabularx} \\ \\ \relax
2:2 & \begin{tabularx}{0.7\textwidth}{X} 「我要走世人必走的路了。你當剛強，作大丈夫， \end{tabularx} \\ \\ \relax
2:3 & \begin{tabularx}{0.7\textwidth}{X} 遵守耶和華－你神所吩咐的，照著摩西律法上所寫的行耶和華的道，謹守他的律例、誡命、典章、法度，好讓你無論做甚麼，不拘往何處去，盡都亨通。 \end{tabularx} \\ \\ \relax
2:4 & \begin{tabularx}{0.7\textwidth}{X} 耶和華必成就他所說關於我的話，說：『你的子孫若謹慎自己的行為，盡心盡意憑信實行在我面前，就不斷有人坐以色列的王位。』 \end{tabularx} \\ \\ \relax
2:5 & \begin{tabularx}{0.7\textwidth}{X} 你也知道洗魯雅的兒子約押向我所做的事，他對付以色列的兩個元帥，尼珥的兒子押尼珥和益帖的兒子亞瑪撒，殺了他們。他在太平之時，如同戰爭一般，流這二人的血，把這戰爭的血染了他腰間束的帶和腳上穿的鞋。 \end{tabularx} \\ \\ \relax
2:6 & \begin{tabularx}{0.7\textwidth}{X} 所以你要照你的智慧去做，不讓他白髮安然下陰間。 \end{tabularx} \\ \\ \relax
2:7 & \begin{tabularx}{0.7\textwidth}{X} 你當恩待基列人巴西萊的眾兒子，請他們常與你同席吃飯，因為我躲避你哥哥押沙龍的時候，他們親近我。 \end{tabularx} \\ \\ \relax
2:8 & \begin{tabularx}{0.7\textwidth}{X} 看哪，在你這裡有來自巴戶琳的便雅憫人，基拉的兒子示每。我到瑪哈念去的那日，他用狠毒的言語咒罵我。後來他卻下約旦河迎接我，我就指著耶和華向他起誓說：『我必不用刀殺死你。』 \end{tabularx} \\ \\ \relax
2:9 & \begin{tabularx}{0.7\textwidth}{X} 但現在你不要以他為無罪。你是有智慧的人，必知道怎樣待他，使他白髮流血下陰間。」 \end{tabularx} \\ \\
[1ex]
\hline
\hline
\end{longtable}
所羅門的王權明顯是出於神,雖有人爭奪,但他仍照神所命定的坐在寶座上.所羅門為王,一方面有人的因素,神要看人的內心是否向他謙卑?亞多尼雅自尊作王,但終失敗,神按他所定的立所羅門為王.另一方面也有神的因素,神命先知拿單用膏油膏立所羅門為王.表明神一切的恩賜,權柄給他.膏油是從神的聖所中取出,證明是從神而來,有了權柄他作王便不能被廢棄.這裡,使我們想到主耶穌基督,祂的職份是神特別給與祂的.祂能行使神一切的權柄,正如希伯來書第一章,作者引用舊約的話「神阿,你的寶座是永永遠遠,你的國權是正直的」.基督作王乃因祂在父神面前存順服謙卑的心,並非如亞多尼雅用手段,他的王權是從神而來,因此他的國,他的寶座是永遠的.
\newline
\newline
自古以來,不少人想起來推翻耶穌的王權,當耶穌出世時,希律王便要殺他,直至現今仍不住有這些的事發生.詩篇第二篇描寫預言中的那位彌賽亞,但也有不少人以為,詩篇第二篇寫作的背景可能與所羅門為王的事有關連,所羅門如何被亞多尼雅起來推翻,基督也是如此；所羅門如何被神所堅立,基督也是一樣.所羅門差不多成為耶穌的縮影.詩篇第二篇指明,「你是我的兒子,我今日生你,你求我我就將列國賜你為基業,將地極賜你為國度」.神把他立在聖山上使他得一切的國權,雖然如此,但詩人又預言到「列國的爭鬥,世上的君王一齊起來,臣宰一同商議」,現今這時代,我們極容易看到這情形,各國聯合商討事情,直至末後有一天世上所有的領袖,都要一同商議反對神,敵視基督,然而詩人在詩篇第二篇中又告訴我們,雖然世上的人都起來反抗神,「但那坐在天上的必發笑,主必嗤笑他們說,我已經立我的君在錫安我的聖山上,受膏者說我要傳聖旨,耶和華曾對我說,你是我的兒子,我今日生你.」神必要主耶穌基督作王,雖然全世界的人集合起來反對他,動員全球的人力物力對抗他,要把耶穌的王權打倒,但我們必能看見坐在天上的發笑,因為耶穌的寶座是無法搖動的.
\newline
\newline
現今世界的人,甚至稱為基督教圈子以內的人,反叛的勢力相當大,他們要把耶穌的王權挪開,特別所謂「神死了」的理論,以為聖經中的那位神已經不存在,這位神可以跟著我們思想轉變而改變,由於時間改變而思想改變,思想改變則神的觀念也改變.這是信仰上基本的錯誤,因為神並不是我們思想上的產物,如同偶像無神的論調一般.神的名字稱為耶和華,意思表明他有自由又是自顯的,我們不能單靠我們的思想來認識神,況且我們的思想也並不十分可靠,人的思想可以顛倒曲直,正如以賽亞先知所說,「禍哉那些以光為暗,以暗為光,以甜為苦,以苦為甜的人」.人類可以用他們思想搜集許多的理論來為自己辯護,例如:世界大戰時有一德國的猶太人十分聰明,精通法律而專任辯護律師,任何一件案情只放在他的手中,他都能轉敗為勝.可見現今科學昌明,人的思想已達到把一切是非顛倒的程度.
\newline
\newline
神絕對不能用人的思想來尋找他,認識他,因為唯有神把自己啟示顯明出來,才能使人認識.雖然昔日人類被創造是按神的形像,但正如羅馬書一19所言:「神的事情,人所能知道的,原顯明在人心裡.」又因罪惡在人心中所行的破壞,甚至我們的良心也喪盡,我們的思想敗壞到沒有明白的」,我們的意志敗壞到「沒有尋求神的」,我們不會尋求神,只會尋求世界,我們對神的認識有如將殘的燈火,加以人糊塗錯誤的思想,便產生奇怪的宗教,如果人必要憑自己的亮光去尋求神這必然是失敗的,因為要認識神必須要神自己啟示出來.耶穌降世成人,這就是神顯現出來啟示與人,正如保羅所傳的福音,非由人教,乃是由神的啟示.舊約三十九卷,其中有一個特式,就是有關「耶和華說」,「神說」,「耶和華曉諭說」等的句語差不多出現一二千次,聖經寫作的人聲明,聖經的話是由於啟示,出自神,並非出自自己.啟示的意思就是以前是封密我們看不到的,但神願意把它顯露出來使我們明白,這不是靠人自己思想所得的.
\newline
\newline
聖經明載:當耶穌降臨時必有離道及大背道之事發生,正如詩篇第二篇所預言的,世上一切的君王都反對主,現今那些說「神死了」的人,其實也是極力反對主,與無神的理論只有一線之差,帶人到滅亡的地步.雖然人在教會中行使一切的破壞,但基督作王,他的王權是不能推翻的,因為耶和華已立他的君在錫安聖山上,正如昔日亞多尼雅推翻王朝,但到底所羅門仍是作王.在社會裡我們可以看見反叛的事越來越多,要推翻耶穌的王權,但我們要靠主站立得穩,認清楚基督的王權是堅定不移的,是永不改變的.
\newline
\newline
第一章裡,亞多尼雅要自立為王,有一班人跟隨他,但另有一班不從,他們要行在神旨中擁護神的受膏者所羅門.今日有不少基督徒也不肯服從聖經的權威,不要耶穌基督完全作王,對於真理只是部份接受,其他則任意而行,現在神所要求尋找的就是那些認識耶穌基督是神所設立的受膏者,知道他是基督,在神手中有一切的權柄,我們都要在他手中,雖然這類人的數目並不會太多,但神仍要尋求那些認識基督是神所設立,是與神有同等權威的人子.
\newline
\newline
很多基督徒走屬靈的道路,在開始時已走差了,基礎裡已有問題,只信耶穌是救主,能使人得福得永生,而罪又得赦免,與主同在好得無比,而不肯服在神前,自己仍要掌權,拒絕基督,不讓他作王.信主的人無心服事主,只想得福.我們必須認識當我們信基督時他便是我們的主,魔鬼常在我們心中擾亂,怕我們完全以基督為王,但我們該完完全全在祂的手中,以致我們得看神完全的能力.現今教會周圍的風氣就是敵對基督,但我們不可與時代同流合污,真理並不是相對的,神的啟示並不分時代性的,乃是直接給與罪人.
\newline
\newline
耶穌在住棚節最後的一天,高聲地說:「人若渴了,可以到我這裡來喝,信我的人就如經上所說,從他腹中要流出活水的江河來.」耶穌用水代表他所講的恩典及永生,也指聖靈.水是我們所需要的,水不受時代所限制.很可惜,現今多少青年基督徒,太多依照自己的意思而行,不認識耶穌的真理是啟示性的,耶穌的權柄並不因我們的時代不同而有所改變及取消,我們應一心一意服在神權柄之下.
\newline
\newline
第二章論及所羅門的公義.神的寶座不能沒有公義,大衛作王時他秉公行義,他臨終時曾吩咐所羅門行公義.公義就是正確的事,行在神旨意,神的法度律例之內.王上二章二三節大衛勸勉所羅門:「當剛強作大丈夫,遵守耶和華你神所吩咐的,照著摩西律法上所寫的行主的道,謹守他的律例誡命典章法度,這樣你無論作甚麼事,不拘往何處去,盡都亨通.」大衛勉勵所羅門要作蒙福的王,必先作一些公義的事,因為神是公義的神,我們應謹守他的誡命律例典章.
\newline
\newline
約押兩次犯了謀殺以色列元帥的罪,約押是一非常能幹的元帥,大衛生前未處理他,現在大衛覺得這事留下來,並不滿足神的心,因流人血的,他的血應被流,若把事情留下不對付,便留下污點不能滿足神的心,因此大衛要所羅門處理這事.
\newline
\newline
大衛逃難時示每咒詛大衛,大衛當時的境況中需要寬大,所以他曾發誓不殺示每,但示每究竟不是真心悔改愛神所立的王,後來大衛把這人心裡的不誠實試驗出來,故大衛又把他交與所羅門處理.
\newline
\newline
巴西萊,大衛逃避押沙龍時,得到巴西萊的好待,現在大衛吩咐所羅門也要好待他的兒女,如有好處合神的心意,便得到神公義的賞賜.
\newline
\newline
大衛是如此公義的王,他命所羅門作王也要如此,大衛不但處理國事的態度是公義,就是個人生活也是如此,他恨惡喜愛公義,正如詩篇一百零一篇所說的:「我要用智慧行完全的道,你幾時到我這裡來呢,我要存完全的心,行在我家中.邪僻的事我都不擺在我眼前,悖逆人所作的事,我甚恨惡,不容沾我身上.彎曲的心思,我心遠離,一切的惡人我不認識.在暗中讒謗他鄰居的,我必將他滅絕,眼目高傲,心裡驕縱的,我必不容他.我眼看國中誠實人,叫他們與我同住,行為完全的,他要伺候我,行詭詐的,必不得住在我家裡,說謊話的必不得立在我眼前,我每日早晨,要滅絕國中所有的惡人,好把一切作孽的從耶和華的城裡剪除.」大衛本性是追求良善,他憎恨罪惡的心有如對獅子一樣,不容許存留在他國中.
\newline
\newline
我們所信的神是公義的一位,我們所領受的是聖靈聖經,我們被主寶血洗淨稱為聖徒,我們必須要憎恨罪惡,不義的態度要除開.
\newline
\newline
當雅各臨終前,他為他的兒子祝福及說預言,西緬及利未受咒詛,雅各說:我的靈哪不要與他們同謀,我的心哪不要與他們聯絡.利未的名意是聯合,西緬的名意是聽見,他在預言中說不與他們聯合,又不願對他們有所聽見,可見雅各雖然是年紀甚老,對罪惡仍不疏忽.羅得是屬世基督徒的代表,但他在所多瑪城中看見罪惡的事,心裡天天傷痛,我們在罪中生活是否仍有難過的感覺,不願與惡同謀的心呢?
\newline
\newline
今天的教會有多少人真正憎恨不義的錢財,我們對公義的態度如何?對色情與主違背的事又如何處理呢?我們有否像大衛公義的態度?對罪惡的看法是跟隨還是推卻?我們作為神聖潔的兒女應該有神聖潔公義的心腸,罪惡就是罪惡,人雖忘記,不願對付,但神必不會忘記,必照我們所行的報應我們,只要我們肯承認自己的罪,神是信實的是公義的,必赦免我們的罪,洗淨我們一切的不義.我們既有王的地位也該有王的公義,羅馬書第二章告訴我們,我們要作公義,作用神給我們的權柄,所以我們要熟讀聖經,明白聖經,真正追求神的公義.
\newline
\newline


\newpage

\section{第41屆~港九培靈會~胡恩德~第四講}
\label{sec:1481}
\textbf{智慧與豐富}
\newline
\newline
連結: \href{https://www.hkbibleconference.org/session-message/view/1481}{\texttt{https://www.hkbibleconference.org/session-message/view/1481}}
\newline
\newline
\hyperref[sec:1480]{< < < PREV SERMON < < <}
~
\hyperref[sec:toc]{[返目錄]}
~
\hyperref[sec:1482]{> > > NEXT SERMON > > >}
\newline
\newline
超凡的智慧
\newline
\newline
所羅門所得的智慧,可說勝過當代一切所有的,他在神面前所求的是一種可聽訟,審判真假是非的智慧,但他所得的不但如此,是勝過萬人,而且聖經具體地把他的智慧記下來:「他作箴言三千句,又寫詩歌一百零五首」,這些是關乎文學,道德上的智慧.又有植物動物學方面的,由高大的香柏樹直至牆上微小的牛膝草；各類飛禽走獸,昆蟲,水族也有述說,他的智慧使天下列邦都派人去觀看聆聽,神賜他一顆廣大的心,能盛載多物,似乎這些都是超過人所能有,也是超過所羅門向神所求的,正如聖經所說:「神能照著運行在我們心裡的大力,充充足足的成就一切超過我們所求所想的.」神不但能聽我們的禱告,而且超過我們所求所想的,向我們施恩.
\newline
\newline
神賜人智慧,而且又豐富地向我們施恩,但可惜現今不少人單單追求科學上的知識,而缺少尋找神與人之間正當的關係.例如把牛角裝上了利刀,牛身裝滿了鐵甲及極利害的釘子,使它亂撞,非徒無益反而有害.人類現在正追求如此的知識,但心靈裡反而貧窮,沒有智慧弄好人與人間的關係,因為根本上人沒有追求與神之間的關係,人雖有很多知識,但這些知識並沒有教人如何成為一個與神聯絡得好的人,反使人成為比野獸更可怕的動物,回復以前森林野獸的生活,只注重科學上的技術,科學上利害的武器.
\newline
\newline
有一動物園,其中養了些猿類,在門前面寫著說:「單單這種動物便能使整個世界絕了種.」原來那字後是放了一座鏡,使看的人在鏡的反影中看見自己的樣子,意思表明單單就是人類已能有足夠的能力,把整個世界毀滅.
\newline
\newline
人是非常可怕的動物,人若失去神的形像,生命,樣式是十分可怕的.今天世界的人只追求屬世的學問知識,偏重一面,使人變成一個十分本事的猛獸.人能殺死世上其他一切的動物,因為人沒有追求更要緊的智慧,甚至現今的科學能把人完全改變,毀滅,使地面一切的事都改變.耶穌曾說在他未回來之前,必有不法之事增多,不法的定義就是沒有秩序的,不認識天地間有一大家應遵守的秩序,否認有永恆的真理,以為世上的事物都是相對的,他們想推翻神的律法及他自己,正如但以理書裡所說的,主再來之前有敵基督統治世界,改變律法及節期.人現今有一不法的心,說沒有一永恆的律法,留下地步給敵基督胡作胡為.
\newline
\newline
人的思想已到混亂不堪的地步,現今世界是極其顛狂,危險,人現今要追求智慧,但卻離開了真正的智慧,不認識「敬畏耶和華是智慧的開端,認識至聖者便是聰明.」所羅門所求的智慧,是使能在子民中執行處理國事,他不是為自己求富貴,求長壽,也不是為自己的好處,乃是好好地行使神的旨意,所以他所得的智慧是遠超過人與人之間,及人與神之間的關係.
\newline
\newline
現今我們應藉著聖經,追求與神和好的智慧,並不是單追求地上的學問.不少人心靈貧乏,沒有智慧,他們只憑地上的智慧行事,結果害了自己,屬神的智慧,是神默示與人,藉聖經寫明出來,我們要迫切尋求這些亮光,使我們到完全的地步,因為神的話能使人有智慧,眼目明亮,人心快樂,我們必因祂的亮光而得幫助及滿足的喜樂.
\newline
\newline
雅各書一章五節,雅各告訴我們:「你們中間若有缺少智慧的,應當求那厚賜與眾人,也不斥責人的神,主必賜給他.」我們要追求智慧,這並不是指世界學問,處事為人的智慧,乃是從天上來的智慧:「先是清潔,後是和平,溫良,柔順......」等等.神指明我們世人在祂面前是罪人,沒有明白的,沒有悟性,智慧,我們對神的事及與神的關係均是無知,只要我們在祂面前才能認識神是至聖者,他能使人有平安,有智慧.
\newline
\newline
現今教會對於聖道的書籍,沒有注重精義,生活上也沒有屬靈的智慧來信靠神,從表面看過去似乎有道理,有信心,也有追求,但實行起來卻缺乏了真正的智慧,我們必須認識自己的軟弱,來到神面前求教導:「若中間缺少智慧,當求那厚賜眾人的神」.我們要在主面前靠著聖經的話,追求真正的智慧,又曉得尋找永遠價值的事物,因為惟有遵行神旨意的才永遠長存.
\newline
\newline
培靈會第一屆在廣州舉行時,「求你開我的眼睛,使我看出你律法中的奇妙」(詩一百二十九篇十八節),此語成為該會的標語,我們事奉神的人,更應求主賜給我們解經的靈,能按著正意分解真理的道.
\newline
\newline
無窮的豐富
\newline
\newline
所羅門每日所用的食物,我們從第四章的記載中略知一二,「細麵三十歌珥,粗麵六十歌珥,肥牛十隻,草場的牛十隻,羊一百隻,還有鹿羚羊,麃子,並肥禽......各按各月供給所羅門王,並一切與他同席之人的食物,一無所缺.」這些都是表明所羅門除了有智慧外,還有無窮盡的豐富.不但在食物方面如此豐富,在第十章裡我們看到所羅門每年所得的金子有六百六十六他連得,所以他宮中一切的器皿都是用金製成,銀在耶路撒冷城的價值及份量如石頭一般,全國都能享受這樣的豐富,國泰民安十分太平.
\newline
\newline
神賜給所羅門有如此的智慧及豐富,他自己也是付上相當代價,在基遍獻祭時把一千的祭牲獻上,又在聖殿落成時向耶和華獻平安祭時,用牛二萬二千,羊十二萬作為祭品.我們今日站在神面前也應有所獻,然而我們十分感謝神,就是有一位大祭司耶穌已替我們獻與神,他所獻的乃將他自己整個心身靈獻上,不但為我們贖罪而且更是為我們成為一個馨香的禮物送給神,把這一切都歸在我們名下而獻,故比千千萬萬的祭品更寶貴.
\newline
\newline
我們今天有神寶貝的應許,就是:「你們若奉我的名無論求甚麼,我必成就.」「到如今你們求甚麼就必得著.」我們可以隨時隨地到主面前求,他必按他的豐富厚賜與我們,況且在我們未求以前,神已經把他的豐富賜給我們,因為有耶穌已為我們獻上祭物.
\newline
\newline
神曾把天上各樣屬靈的福氣賜給我們,我們在神面前正如一個蒙愛的兒女,保羅寫以弗所書時,寫出神各樣的豐富,又求神開他們的眼睛,使他們真知道屬神的人有何等大的盼望,神賜給我們的恩是十分豐富,我們應多追求屬天的豐富,不以世上的事物為滿足,因為有不少的人就算得著世界的事物,心靈仍覺非常空虛.
\newline
\newline
許多時候,我們心靈感到貧乏,沒有屬靈得勝的力量,又不能感動人來到神面前,不要忘記:神是豐富的一位,他願意向我們施恩,在我們未求以先已賜給我們了,就是在我們相信的時候已接受他各樣豐富的恩惠,而且超過我們所求和所想的.
\newline
\newline
羅馬書八章三十一節告訴我們,神若幫助我們,誰能敵擋我們呢?神站在我們一邊,誰能扺擋我們呢?我們一切的貧乏,罪惡,軟弱也包括在內.神既不愛惜自己的兒子,為我們眾人死了,豈不也把萬物和他一同白白的賜給我們嗎?神願意將他的豐富賜給我們,我們不要自怨貧乏,所羅門王昔日能得著神舊約方式的豐富,今日我們也能憑著應許得著神新約方式的豐富,因為我們已有極大的祭獻上,我們已有極大的應許.
\newline
\newline


\newpage

\section{第41屆~港九培靈會~胡恩德~第五講}
\label{sec:1482}
\textbf{耶和華神的殿}
\newline
\newline
連結: \href{https://www.hkbibleconference.org/session-message/view/1482}{\texttt{https://www.hkbibleconference.org/session-message/view/1482}}
\newline
\newline
\hyperref[sec:1481]{< < < PREV SERMON < < <}
~
\hyperref[sec:toc]{[返目錄]}
~
\hyperref[sec:1483]{> > > NEXT SERMON > > >}
\newline
\newline
一個蒙大恩的僕人,又是作王的所羅門,我們必能在他的身上看見神許多特別的恩惠.最初神給他一個應許就是:「你要甚麼可以求」,所羅門只求智慧,目的是要服事神,服事人,神所給他的智慧乃是過於他所求的,包括其他種種的智慧.在他建造的事上,也把神的恩典表彰出來,首先在他父親大委身上,神彰顯能力,福份,但大衛並非作建設的工作,只是負責爭戰,神使他勝過國內國外一切的仇敵,使國的基礎堅固,直至所羅門登基作王,便把殿建起來,他的工作是負責建設,神在這兩個王身上彰顯祂的榮耀及豐富,神要利用這些奇妙的能力,恩典吸引其他的民族歸向祂,神在這兩個王身上彰顯祂的榮耀及豐富,神要利用這些奇妙的能力,恩典吸引其他的民族歸向祂,神使用以色列整個民成為彰顯神榮耀的一個器皿,使別人看見而認識神.
\newline
\newline
神揀選教會,為要叫他施恩的心得以滿足,他有極大的恩惠,正如啟示錄所寫的,在人子耶穌的右手中手持一書卷,其中有了神一切的榮耀,能力,美德,權柄.在教會屬神的人身上,神要施恩,但卻找不到真正合宜能接受他恩的人.耶穌是教會的元首,基督要把祂充滿萬有的充滿教會,神揀選教會來成就他的計劃,又使教會得著他的榮耀和豐富,在地面彰顥出來,使不屬教會的人也能相信接受那位救主耶穌基督.
\newline
\newline
在神的計劃裡,揀選以色列,顯出祂的榮耀,豐富及能力.人不明白便問為何神選以色列民族,以為是神偏心.神愛世人,他選以色列是為要吸引萬民,神曾經應許以色列,他們若肯遵守神的話便得福,不順從便是禍.這些是舊約時代的福,是代表暫時榮耀的福份.新約時代的福,乃是更深一層的,是神最終最高的計劃.我們相信祂,服從祂,我們就會享受來生的福,這是新約時代的福,是深入的,長遠的,將來的,靈界的.
\newline
\newline
神揀選我們不單用以滿足他施恩的愛心,也要藉教會彰顯祂的榮耀,就是來生的榮耀.神不是藉外表的事物來爭勝,神藉先知耶利米責備約雅敬王:「難道你作王是在乎建宮殿爭勝麼?」現今的教會以為多建禮拜堂便是榮耀神,以外面的體面,外面那一套來榮耀神,這是十分錯誤,難道保羅斬首而死,彼得釘十字架,耶穌受別人恥辱而死這不是榮耀神嗎?真正榮耀神的心,不在乎外表的修飾,乃在乎愛人的心,寬容的態度,確實的信心等,我們在教會中要謙卑,愛人,持守主道,尊重主話,自然外界的人就能看見,把榮耀歸與神.
\newline
\newline
所羅門所建的殿為要表彰主,這殿與以西結在異象中所看見的殿作比較,便看到以西結所見的殿是表明神的聖潔,使污穢之物與人隔離；所羅門的殿乃是注重神豐富那方面.聖殿製作的材料是最上乘的,木用香柏木,是當時最貴重的木材,聖所及至聖所的器具用精金包上,殿內所用的石頭乃是用鉅鋸鋸齊,是王派幾萬人上山鑿成的,他們揀選又寶貴又大塊的石塊原裝運來,可見工程浩大,國中豐富的情形.
\newline
\newline
第七章裡我們看見會幕之前面,放有一大銅盤,院門口有一銅祭壇,銅盤不是放正在殿前的門口,乃是放在靠右的那邊,銅盤稱為銅海,外型巨大而重,祭司可在其間沐浴.在聖殿之前放有兩枝銅柱,柱頂有裝飾物稱之為鋼子,聖所內的台不但只有一個,而且每個都是用金製成的,金是反光的物體,所謂金碧煇煌,用以代表榮耀及豐富,由此可見所羅門的豐富是無法計算的.
\newline
\newline
我們靈性上若有神的恩典,我們必不會貧窮缺乏,我們要追求屬靈的豐富,例如豐富的交通,詩人的經歷,「神若向他閉口不言,他便要與死人一般」,可見與主交不通是一件非常苦的事.又要追求認識神的豐富,我們一切的好處都在於認識神的事上.耶穌死以前向父神禱告說:「認識你獨一的真神及耶穌基督就有永生」,這認識不是頭腦上的認識,正如以弗所書題到的「真知道」,深入的認識.我們得救以後也要追求有所認識,認識耶穌為至寶.神有很多的豐富,只要我們認識他就有豐富.保羅追求認識基督,要效法和祂一同死,他在哥林多傳道時,自己作工,他的賞賜就是傳福音給別人時,能使別人白白得福音,他以別人白白的領受為自己的賞賜.
\newline
\newline
我們中間不少是心裡貧乏,我們不能愛人,特別那些與我們為敵的,我們也沒有寬容的忍耐及力量,因為真正的力量是從神而來,神的殿是如此豐富,此乃表明神屬靈裡的豐富,我們必因祂聖殿裡肥美的膏油而得飽足.
\newline
\newline
銅盤是說明得救以後的人,每日生活中的罪惡都要洗淨,祂把我們的罪踏在祂的腳下,又投在深海裡,我們得救以後雖有失敗跌倒,但因有銅盤為我們預備,只要我們肯向主承認自己的罪,神是信實的,是公義的,必要赦免我們的罪,洗淨我們一切的不義,神豐富洗罪的恩典,一生的罪,每日的罪祂都能替我們洗淨.
\newline
\newline


\newpage

\section{第41屆~港九培靈會~胡恩德~第六講}
\label{sec:1483}
\textbf{神的同在}
\newline
\newline
連結: \href{https://www.hkbibleconference.org/session-message/view/1483}{\texttt{https://www.hkbibleconference.org/session-message/view/1483}}
\newline
\newline
\hyperref[sec:1482]{< < < PREV SERMON < < <}
~
\hyperref[sec:toc]{[返目錄]}
~
\hyperref[sec:1484]{> > > NEXT SERMON > > >}
\newline
\newline
列王記上 6:23-6:25
\newline
\begin{longtable}{cl}
\hline
\hline
章節 & 經文 (和合本修訂版)\\
\hline
6:23 & \begin{tabularx}{0.7\textwidth}{X} 他在內殿裡用橄欖木做兩個基路伯，各高十肘。 \end{tabularx} \\ \\ \relax
6:24 & \begin{tabularx}{0.7\textwidth}{X} 這基路伯的一個翅膀長五肘，另一個翅膀長五肘，從一個翅膀尖到另一個翅膀尖共有十肘； \end{tabularx} \\ \\ \relax
6:25 & \begin{tabularx}{0.7\textwidth}{X} 第二個基路伯也是十肘；兩個基路伯的尺寸、形狀都一樣。 \end{tabularx} \\ \\
[1ex]
\hline
\hline
\end{longtable}
列王記上第六章,神要在其中顯明祂是豐富的,在所羅門王建殿的事上可見到神有無限的豐富.另有一件可表彰神豐富的,就是殿的本身,殿前所羅門立兩枝銅柱,極其華美,放在殿中,左邊稱為雅斤,右邊稱為波阿斯.雅斤意為「祂必堅立」,波阿斯意為「在祂裡面有能力」,兩銅柱均是堅立之能力的意思,這麼大的銅往以色列人從未見過,用它放在殿前乃是表明堅立的能力,保守的能力.
\newline
\newline
神表彰他的豐富就是屬祂的能力,很多時候基督徒說他甚軟弱,若我們真的從心裡覺得十分軟弱的話,我們就要在聖殿建築裡得著安慰；就算我們甚是軟弱,只要我們來到神豐富的能力裡求主幫助,便可克服.最難過的,是我們自己軟弱卻不自知,我們需要發現自己的軟弱,使我們重新得回神的力量.神的力量是可以拿出來的,銅柱不是藏在殿內,乃是放在殿前,只要人肯來到殿前便可以看見,祂能建立保守我們,我們隱藏在祂裡面,神的能力就會包圍我們,又透過我們傾福下來.
\newline
\newline
所羅門的殿是代表神的豐富,大銅盤代表神洗罪的恩是何等寬大,我們任何一個基督徒都有祭司的身份,可以隨時得到洗淨,祭壇之後有洗濯的海放在那裡,使我們得潔淨.
\newline
\newline
聖殿內滿了基路伯,至聖所是大祭司每年進去一次獻祭的地方,在內可看見基路伯,在他翼底下是神的約櫃,櫃上有兩金基路伯俯視約櫃的蓋,就是施座,這是摩西時代所造的,那金的基路伯是所羅門時代加製,但也是照神所吩咐的去做,除此以外,其他牆壁上也刻滿了基路伯的形像,甚至會幕內的幕也有基路伯.基路伯是一種高級的靈界,比天使的地位更高,他們能親近神到一個相當近的地位,在神的寶座以下及周圍,基督伯之所有乃是表明有神的同在,殿內有基路伯乃是表明殿裡隨處也有神的同在.同時基路伯的形像是提醒人,而不是叫人來敬拜他,基路伯的腳像牛,有四種面:人,鷹,牛,獅等,這些不是人間的東西,因而使人產生敬畏的心,神本身威嚴,我們在神前當存教畏的心.聖經一方面強調神是可愛,但我們在神前要存敬畏的心,正如詩篇裡所說的,「存戰兢而快樂」,「因主而樂」.對主的態度是裡面存著敬畏戰兢的心,因主是可畏嚴肅的,我們作神子民的也應有威嚴.我們在神面前受審判時是戰慄恐懼的,因為審判刑罰要來,使肉體的人不能承受,耶穌作人時把神的威嚴隱藏起來,只剩下一部份屬靈的威嚴,因他在世與人交通,但在他死而復活以後便得回一切屬神的威嚴.
\newline
\newline
在聖殿裡有如此多雕刻的基路伯,使我們認識到神的威嚴,我們在祂面前不能隨便,因為有神同在.教會也有神同在,聚會也有神同在,我們每一個信徒合起來便成了聖殿,神的靈要住在其中.教會的大團體有主住在當中；我們若有兩三個人奉神的名聚會,祂也必在我們中間.我們不能憑自己的思想處理教會的行政,事務,乃是遵守聖經所言的去辦理,不可把屬肉體的東西帶進教會裡.
\newline
\newline
保羅寫信給哥林多教會,告訴他們他不以其中的信徒為屬靈,乃是屬肉體的,因他們常常分爭結黨,好勝非常,豐不知神的靈住在我們當中,若有人要毀壞神的殿,他也必受毀壞.若我們在教會中弄是非,便是毀壞神的威嚴了.在舊約時代有神同在,新約時代也有神同在,就是在教會聚集中與我們同在,因此我們所做的每一件事動機都要非常正確.
\newline
\newline
神偉大的僕人哥登先生,有一件事是在他傳道生活早期發生的,這事影響他傳道的心志,以成為有名的傳道人.有一次歌登在夢中夢見自己在講道,當時座位差不多快要滿,有一個十分平常的人走進來,找來找去找不到座位,最後他直行上前,找到一個座位坐下,坐下以後他定睛地望著哥登先生,使哥登感到十分託異,散會後,哥登便走下來希望與這一個人談話,怎料那人已不見了,只找到與他鄰座的那位,哥登從那人的口中知道那遲到的聽者就是拿撒勒人耶穌,而且他不但這一個聚會才來,以後每主日的聚會他都會來,哥登先生聽到這些話,便從夢中醒來.
\newline
\newline
聖經告訴我們:每一個聚會耶穌都會來參加,且留心聆聽,那麼我們在聚會之中,是否馬馬虎虎隨便就算,還是要尊主為大,討他的喜悅呢!
\newline
\newline
所羅門殿內滿有基路伯,目的是提醒我們殿內有主的同在,正如教會及教會的聚會有主同在一樣.以前耶穌與門徒一同生活在一起,他們感到異常的快樂.但當主死後他們非常難過,以為失去主的同在,以前一切需用主都知道,為他們安排妥當,使他們一無所缺,但他們卻不明白.雖然耶穌已離開,聖經仍然給我們應許,就是我們可求父,父就另賜給我們一位保惠師與我們同在,「何惠師」的名意是與你並行的,來幫助,安慰,供應你.
\newline
\newline
神把祂獨生的兒子耶穌基督賜給我們,後來又賜給我們真理的靈.聖靈接替了耶穌的職任,他與我們同在,但這位聖靈能明白我們一切的苦情,他能在我們裡面幫助我們,不是在旁邊表面上的工作,正如彼得在五旬節聖靈降臨那天所說的:所以你們各人應當悔改,奉主耶穌的名受洗,使你們的罪得赦,就必領受所賜的聖靈.
\newline
\newline
我們每一個信徒都成為神的殿,神能夠而且是樂意住在我們當中,每當我們信主以後罪得赦免,我們便有基督的生命成為我們的生命,他住在我們裡面為要幫助我們,成為有榮耀的盼望.
\newline
\newline


\newpage

\section{第41屆~港九培靈會~胡恩德~第七講}
\label{sec:1484}
\textbf{名聲與失敗}
\newline
\newline
連結: \href{https://www.hkbibleconference.org/session-message/view/1484}{\texttt{https://www.hkbibleconference.org/session-message/view/1484}}
\newline
\newline
\hyperref[sec:1483]{< < < PREV SERMON < < <}
~
\hyperref[sec:toc]{[返目錄]}
~
\hyperref[sec:1485]{> > > NEXT SERMON > > >}
\newline
\newline
列王記上 10:1-10:10
\newline
\begin{longtable}{cl}
\hline
\hline
章節 & 經文 (和合本修訂版)\\
\hline
10:1 & \begin{tabularx}{0.7\textwidth}{X} 示巴女王聽見所羅門因耶和華的名所得的名聲，就來要用難題考問所羅門。 \end{tabularx} \\ \\ \relax
10:2 & \begin{tabularx}{0.7\textwidth}{X} 她帶著很多的隨從來到耶路撒冷，有駱駝馱著香料、極多金子和寶石。她來到所羅門那裡，向他提出心中所有的問題。 \end{tabularx} \\ \\ \relax
10:3 & \begin{tabularx}{0.7\textwidth}{X} 所羅門回答了她所有的問題，沒有一個問題太難，王不能向她解答的。 \end{tabularx} \\ \\ \relax
10:4 & \begin{tabularx}{0.7\textwidth}{X} 示巴女王看見所羅門一切的智慧，和他所建造的宮殿， \end{tabularx} \\ \\ \relax
10:5 & \begin{tabularx}{0.7\textwidth}{X} 席上的食物，坐著的群臣，侍立的僕人，他們的服裝，和他的司酒長，以及他在耶和華殿裡所獻的燔祭，就詫異得神不守舍。 \end{tabularx} \\ \\ \relax
10:6 & \begin{tabularx}{0.7\textwidth}{X} 她對王說：「我在本國所聽到的話，論到你的事和你的智慧是真的！ \end{tabularx} \\ \\ \relax
10:7 & \begin{tabularx}{0.7\textwidth}{X} 我本來不信那些話，及至我來親眼看見了，看哪，人所告訴我的還不到一半，你的智慧和你的福分超過我所聽見的傳聞。 \end{tabularx} \\ \\ \relax
10:8 & \begin{tabularx}{0.7\textwidth}{X} 你的人是有福的！你這些僕人常侍立在你面前、聽你智慧的話是有福的！ \end{tabularx} \\ \\ \relax
10:9 & \begin{tabularx}{0.7\textwidth}{X} 耶和華－你的神是應當稱頌的！他喜愛你，使你坐以色列的王位，因為他永遠愛以色列，所以立你作王，使你秉公行義。」 \end{tabularx} \\ \\ \relax
10:10 & \begin{tabularx}{0.7\textwidth}{X} 於是，示巴女王把一百二十他連得金子、極多的香料和寶石送給所羅門王；送來的香料，從來沒有像示巴女王送給他的那麼多。 \end{tabularx} \\ \\
[1ex]
\hline
\hline
\end{longtable}
實際的名聲
\newline
\newline
所羅門的名聲遠播四方,甚至在阿拉伯以南的示巴女王,也不怕長途跋涉到耶路撒冷來見所羅門王,為要看看他的豐富與榮華.看後,示巴女王深深覺得所羅門的實過於他的名聲.可見所羅門的名聲,豐富,智慧,不單是外表的榮耀,乃是非常實際的.
\newline
\newline
啟示錄裡給老底嘉教會的信,乃是責備老底嘉教會,表現過於實際,自己以為發財富足,而實際是空無所有.故此神要從口中把他們吐出來,他們是:「沒有雨的雲彩,秋天沒有果子的樹,死而又死,連根被拔出來.
\newline
\newline
老底嘉教會情形如此,今天教會情形也是一樣:某些人自以為在知識上有認識,對信知識造詣相當高,便開始對聖經大膽地批評,而且在批評時加上「身為二十世紀的人」,「作為原子時代的人」等等的字眼,他們卻不知道在主的道上是平等的,無論甚麼地位,學識,在審判台前,並沒有分別.我們應存謙卑的心,接受神真實的話──聖經.正如人在風浪中一樣,不能說因我是大學生,研究院的博士,哲學博士,體育家等便以為不會暈浪,這並不根據外在的因素,乃是與內在實質上有關,倘若我們還不自知,以世上一切事物為誇口,便會落到老底嘉教會的光景中,自以為富足,得勝,然而在耶穌看來卻是十分貧乏.在得救的事上,我們不能把世上所得的搬進去,神所看的並非是屬世的好處,乃是自己屬靈方面的造就,不能自以為了不起,誇大歸於自己的實際,這類人主要教導他們,向神買火煉的精金.
\newline
\newline
神賜所羅門的豐富是十分實際的,這也是與所羅門本身有關,因他首先認識自己無有,神才賜他真正的豐足,過於他所求的.我們靈裡貧乏,是因為我們沒有饑渴的態度和謙卑的心.我們對聖經,不能只當為白紙黑字,講時自己也要在其中領受,不能言中無物,這樣是不能實際供應人內心的需要.在我們的經歷裡,倘若從來沒有在主面前被主對付,未藉著環境被主打倒,不認識自己的敗壞軟弱,因此不能透徹地認識主,所得著的豐富也是十分小,很多基督徒在得救開始時,沒有好好追求,不認識應付一切事情的是主,因此便常常頂撞主,抵擋主,傷主的心.我們要多多依靠主,多得主的豐富,如此才能有實際的豐富供應別人.從這裡我們便看出所羅門之所以能領受如此豐富,乃是與他謙卑的心有密切的關係.
\newline
\newline
所羅門的實際過於他的名聲,甚至在他以後,沒有一人能與他相比,直至末後猶大的獅子耶穌降生為止,耶穌基督在世時,曾題及所羅門王說:「示巴女王從南方來要聽所羅門智慧的話,看哪,這裡有一位比他更大的」.耶穌的豐富與舊約的有所不同,他所有的乃是更深的豐富,耶穌藉著他對別人的愛心,助人的恩典,表明他的豐富,這些都是所羅門沒有的.耶穌的恩典使跟從他的人一生不會忘記,特別是西庇太的兒子約翰,他為主作見證說；「道成肉身,住在我們中間,充充滿滿的有恩典有真理.」耶穌的豐富何等廣大.初傳道時,「耶和華用膏膏祂,叫祂傳好信息給謙卑的人,差遣祂醫好傷心的人,報告被擄的得釋放,被囚的出監牢」.所羅門的實際過於他的名聲,耶穌比他更豐富.直到今日,祂有很多摸不著的恩典,正如詩人謂:「你把你的話顯為大,過於你自己.」耶穌的豐富是何等長闊高深,我們不能用思想來作判斷.
\newline
\newline
示巴女王用許多難題來問所羅門,所羅門對答如流,以致示巴女王十分佩服.耶穌的智慧比所羅門更高更大,只要我們把難處帶到主面前,祂能幫助我們.示巴女王從遠方來見所羅門,除了要聽他智慧的話語外,她還把豐富的禮物送與所羅門；所羅門照他所得的厚厚還給她.我們來到主面前,祂能照我們的缺乏賜給我們,而且富富有餘,超過我們所想的,使我們的福杯滿溢.
\newline
\newline
我們最難度過的一關,就是要站在一位絕對公義執行律法的神面前,我們生出便是一個罪人,沒法行義,也不能行神的道,滿足神聖潔的要求,但感謝神只要我們接受耶穌的救恩,我們便能得稱為義人,靠著主經過了神的審判,主既能幫助我們經過神的審判,試問我們生活上還有甚麼是主不能為我們解決的呢?
\newline
\newline
最後的失敗
\newline
\newline
所羅門失敗跌倒,以致神的責罰臨到他,從他手中把國奪回,但因耶路撒冷是神所揀選的,又因大衛的緣故所以神施恩保留猶大及便雅憫支派,其他十個支派賜給耶羅波安,所羅門失敗的原因,就是他隨從偶像假神,為外邦妃嬪建造可憎惡的神廟,效法她們敬拜偶像.
\newline
\newline
曾經兩次得到神的顯現及堅立的所羅門,仍遭失敗,我們要小心.所羅門跌倒是慢慢墮落,不知結局是如此利害,直至神把他的國奪回.我們常憑著自己一兩次屬靈特殊的經歷,便以為了不起,結果遭遇失敗而不自知,不少人正如羅得一樣,把帳棚漸漸搬進城內,結果遭遇失敗.當所羅門初作王時,他對付亞多尼雅的祭司亞比亞他,亞比亞他以前曾在神面前抬過約櫃,又與大衛一同受過苦,但他失敗,被所羅門把他祭司的職份開除.
\newline
\newline
列王記由始至終就是告訴我們,人是失敗的,我們要加倍小心,不要自以為能得勝,能站立得穩,我們多多依靠主,免得成為得救行列中失敗的一個,一個傳道人最可惜的就是半途而廢,失節不能盡忠,當我們看見別人犯罪跌倒之時,我們不可自欺,反而加倍小心,認識我們自己是不能的,要依靠主才能站立得穩.a
\newline
\newline


\newpage

\section{第41屆~港九培靈會~胡恩德~第八講}
\label{sec:1485}
\textbf{分裂的國度}
\newline
\newline
連結: \href{https://www.hkbibleconference.org/session-message/view/1485}{\texttt{https://www.hkbibleconference.org/session-message/view/1485}}
\newline
\newline
\hyperref[sec:1484]{< < < PREV SERMON < < <}
~
\hyperref[sec:toc]{[返目錄]}
~
\hyperref[sec:1486]{> > > NEXT SERMON > > >}
\newline
\newline
列王記上 12:1-12:33
\newline
\begin{longtable}{cl}
\hline
\hline
章節 & 經文 (和合本修訂版)\\
\hline
12:1 & \begin{tabularx}{0.7\textwidth}{X} 羅波安往示劍去，因以色列眾人都到了示劍，要立他作王。 \end{tabularx} \\ \\ \relax
12:2 & \begin{tabularx}{0.7\textwidth}{X} 尼八的兒子耶羅波安先前躲避所羅門王，逃往埃及，住在那裡。他還在埃及，聽見了這事， \end{tabularx} \\ \\ \relax
12:3 & \begin{tabularx}{0.7\textwidth}{X} 以色列人派人去請他來。耶羅波安就和以色列全會眾來，與羅波安談話，說： \end{tabularx} \\ \\ \relax
12:4 & \begin{tabularx}{0.7\textwidth}{X} 「你父親使我們負重軛，現在求你減輕你父親所加給我們的苦工和重軛，我們就服事你。」 \end{tabularx} \\ \\ \relax
12:5 & \begin{tabularx}{0.7\textwidth}{X} 羅波安對他們說：「你們走吧，過三天再來見我。」百姓就走了。 \end{tabularx} \\ \\ \relax
12:6 & \begin{tabularx}{0.7\textwidth}{X} 羅波安的父親所羅門在世的日子，有侍立在他面前的長者，羅波安王和他們商議，說：「你們出個主意，好把話帶回給這百姓。」 \end{tabularx} \\ \\ \relax
12:7 & \begin{tabularx}{0.7\textwidth}{X} 他們對他說：「現在王若像僕人一樣服事這百姓，用好話回覆他們，他們就永遠作王的僕人了。」 \end{tabularx} \\ \\ \relax
12:8 & \begin{tabularx}{0.7\textwidth}{X} 王不採納長者給他出的主意，卻和那些與他一同長大、在他面前侍立的年輕人商議。 \end{tabularx} \\ \\ \relax
12:9 & \begin{tabularx}{0.7\textwidth}{X} 他對他們說：「這百姓對我說：『你父親使我們負重軛，求你減輕一些。』你們出個甚麼主意，我們好把話帶回給他們。」 \end{tabularx} \\ \\ \relax
12:10 & \begin{tabularx}{0.7\textwidth}{X} 那些與他一同長大的年輕人對他說：「這百姓對王說：『你父親使我們負重軛，求你給我們減輕一些。』王要對他們如此說：『我的小指頭比我父親的腰還粗呢！ \end{tabularx} \\ \\ \relax
12:11 & \begin{tabularx}{0.7\textwidth}{X} 我父親使你們負重軛，現在我必使你們負更重的軛！我父親用鞭子懲罰你們，我要用蠍子懲罰你們！』」 \end{tabularx} \\ \\ \relax
12:12 & \begin{tabularx}{0.7\textwidth}{X} 耶羅波安和眾百姓遵照王所說「你們第三天再來見我」的話，第三天來到羅波安那裡。 \end{tabularx} \\ \\ \relax
12:13 & \begin{tabularx}{0.7\textwidth}{X} 王嚴厲地回答百姓，不採納長者給他出的主意。 \end{tabularx} \\ \\ \relax
12:14 & \begin{tabularx}{0.7\textwidth}{X} 他照著年輕人所出的主意對他們說：「我父親使你們負重軛，我必使你們負更重的軛！我父親用鞭子懲罰你們，我卻要用蠍子懲罰你們！」 \end{tabularx} \\ \\ \relax
12:15 & \begin{tabularx}{0.7\textwidth}{X} 王不依從百姓，因這事件是出於耶和華，為要應驗耶和華藉示羅人亞希雅對尼八的兒子耶羅波安所說的話。 \end{tabularx} \\ \\ \relax
12:16 & \begin{tabularx}{0.7\textwidth}{X} 以色列眾人見王不依從他們，百姓就回話給王，說：「我們在大衛中有甚麼份呢？我們在耶西的兒子中沒有產業！以色列啊，回你的帳棚去吧！大衛啊，現在你顧自己的家吧！」於是，以色列人都回自己的帳棚去了； \end{tabularx} \\ \\ \relax
12:17 & \begin{tabularx}{0.7\textwidth}{X} 至於住猶大城鎮的以色列人，羅波安仍作他們的王。 \end{tabularx} \\ \\ \relax
12:18 & \begin{tabularx}{0.7\textwidth}{X} 羅波安王派監管勞役的亞多蘭去，以色列眾人用石頭打他，他就死了。羅波安王急忙上車，逃回耶路撒冷去了。 \end{tabularx} \\ \\ \relax
12:19 & \begin{tabularx}{0.7\textwidth}{X} 這樣，以色列背叛大衛家，直到今日。 \end{tabularx} \\ \\ \relax
12:20 & \begin{tabularx}{0.7\textwidth}{X} 以色列眾人聽見耶羅波安回來了，就派人去請他到會眾那裡，立他作全以色列的王。除了猶大支派，沒有跟從大衛家的。 \end{tabularx} \\ \\ \relax
12:21 & \begin{tabularx}{0.7\textwidth}{X} 羅波安來到耶路撒冷，召集了猶大全家和便雅憫支派的人共十八萬，都是精選的戰士，要與以色列家打仗，好將王國奪回，歸所羅門的兒子羅波安。 \end{tabularx} \\ \\ \relax
12:22 & \begin{tabularx}{0.7\textwidth}{X} 但神的話臨到神人示瑪雅，說： \end{tabularx} \\ \\ \relax
12:23 & \begin{tabularx}{0.7\textwidth}{X} 「你去告訴所羅門的兒子猶大王羅波安，猶大和便雅憫全家，以及其餘的百姓，說： \end{tabularx} \\ \\ \relax
12:24 & \begin{tabularx}{0.7\textwidth}{X} 『耶和華如此說：你們不可上去與你們的弟兄以色列人打仗。你們各自回家去吧！因為這事是出於我。』」眾人就聽從耶和華的話，遵照耶和華的話回去了。 \end{tabularx} \\ \\ \relax
12:25 & \begin{tabularx}{0.7\textwidth}{X} 耶羅波安在以法蓮山區建了示劍，住在其中，又從示劍出去，建了毗努伊勒。 \end{tabularx} \\ \\ \relax
12:26 & \begin{tabularx}{0.7\textwidth}{X} 耶羅波安心裡說：「現在，這國恐怕仍會歸大衛家； \end{tabularx} \\ \\ \relax
12:27 & \begin{tabularx}{0.7\textwidth}{X} 這百姓若上耶路撒冷去，在耶和華的殿裡獻祭，他們的心必歸向他們的主猶大王羅波安。他們會殺了我，仍歸猶大王羅波安。」 \end{tabularx} \\ \\ \relax
12:28 & \begin{tabularx}{0.7\textwidth}{X} 耶羅波安王就籌劃，鑄造了兩個金牛犢，對眾百姓說：「你們上耶路撒冷去實在夠久了。以色列啊，看哪，這是領你出埃及地的神明。」 \end{tabularx} \\ \\ \relax
12:29 & \begin{tabularx}{0.7\textwidth}{X} 他把一個安置在伯特利，另一個安置在但。 \end{tabularx} \\ \\ \relax
12:30 & \begin{tabularx}{0.7\textwidth}{X} 這事使百姓陷入罪裡，因為他們甚至到但去拜那牛犢。 \end{tabularx} \\ \\ \relax
12:31 & \begin{tabularx}{0.7\textwidth}{X} 耶羅波安在一些丘壇建神殿，立不屬利未人的平民百姓為祭司。 \end{tabularx} \\ \\ \relax
12:32 & \begin{tabularx}{0.7\textwidth}{X} 耶羅波安定八月十五日為節期，像在猶大的節期一樣，自己上壇獻祭。他在伯特利這樣做，向他所鑄的牛犢獻祭，又把他所立丘壇的祭司安置在伯特利。 \end{tabularx} \\ \\ \relax
12:33 & \begin{tabularx}{0.7\textwidth}{X} 他在八月十五日，就是他自己心中所定的月份，在伯特利上到自己所造的祭壇；他為以色列人定了一個節期，親自上壇燒香。 \end{tabularx} \\ \\
[1ex]
\hline
\hline
\end{longtable}
列王記上半題及人的失敗,所羅門開始時有謙卑的心,尋求神的智慧,服事主的百姓,所以神賜給他,遠超他所求.所羅門開始作王時表現很好,但他漸漸失敗,有好的開始,卻以失敗結束.
\newline
\newline
所羅門隨從異邦偶像而遭失敗,這失敗乃是由於不照神的教訓而行,與外邦不認識神的女子結合.婚姻使兩人的感情達到火熱的地步.始祖亞當在這情感上失敗,所羅門也是失敗在火熱裡,他如此作並非不明道理,他寫詩篇,箴言,傳道書,還有雅歌書.所羅門有屬靈的實際,也有很好的基礎,一個如此實際的人,到底還是失敗,失敗於異性屬肉體的吸引,這失敗連他自己也是想不到的.他父親大衛也是在這方面失敗,而自己卻步上父親所踐踏的途徑上.
\newline
\newline
這條路正是魔鬼的出路,所羅門在此失敗招惹了神的責罰,甚至神要從他手中把國奪回,掃羅以前失敗,是因他不遵照神旨意行事,私自獻祭,又留下當滅絕的牛群羊群,他不是全心全意服主的道,所以神把掃羅的國奪回.所羅門的國神沒有收回,乃是因他父親大衛的原故.
\newline
\newline
教會蒙大恩,有主豐富的能力同在,我們應當小心,教會的領袖不可因驕傲而跌倒,否則神的恩典很快便要過去了.所羅門初期肯謙卑在神面前,所以他有神滿足的恩典,後來因為好色出了毛病,就成為他的弱點,給撒旦留下地步,甚至國度因此傾覆.在列王記十二章裡給我們看到,所羅門的國度要受審判.我們個人在神面前也是如此,若自己造成一個破口,撒旦便從那裡進來,把整個人都奪去,只剩下一個空殼也不自知,還以為與以前一樣.基督徒蒙恩時也是我們最要防堵之時,常見有人身上帶有神的大能,而他們又是認真地歸向主的,只因小小的毛病造成小小的破口,便失去被主使用的恩典.
\newline
\newline
有一西教士勉勵一傳道人,要他在傳道的事工上小心三件事:女色,金錢,名譽.後來這傳道人果真失敗,就是在名譽那方面跌倒.好名和享受,都不是我們所應該追求的,我們乃要追求真正愛主,把我們一生放在更有意義的工作上.一個傳道人若有愛世界的心,那麼他的工作,說話便不能起作用,果效也不能在人心裡.
\newline
\newline
當大衛的的國度極其強盛時,他吩咐全國數點百姓,他不知道這是神所不喜悅的,他的元帥約押也勸他不要如此行,結果那次數點百姓死了幾萬人,因為大衛乃是要顯出自己的榮耀.我們若到一個靈性很高的地步中,我們更要謙卑,自以為站立得穩的更當謹守,不可依靠自己的學問,才幹,應隨時隨地求問主,以他為首.
\newline
\newline
神許可把刑罰加在大衛家,當然與所羅門王末後所犯的罪惡有極大的關係,但也有第二個因素,就是羅波安自己不好,不體恤民間的苦情；以法蓮人耶羅波安到羅波安面前,求他減輕百姓的工作,因他們甚是困苦,羅波安叫他過了三天再來見王.
\newline
\newline
羅波安在此事上顯出他的愚昧,他應利用三天的時間好好求問神,或是想念神施與他父親的一切恩典,好待他的百姓,我們很多時候失敗,恐怕也是因著缺乏禱告.
\newline
\newline
耶羅波安反叛得到北方十個支派,便在示劍作王,稱為以色列國,他作主以後深恐百姓再返南國,因以色列人一年有三個節期必要上到耶路撒冷朝拜神,所以他便造了兩隻金牛,一隻放在國境之北但的地方,一隻放在國境之南伯特利的地方,又自己定下節期命百姓向金牛獻祭.這事使以色列國陷在罪惡中直至國家滅亡,這件罪惡透過了北國的王朝,它的罪甚重,如一條毒根在他們中間.耶羅波安對神有一不信的惡念,神既把國度放在他手中,自然神會保守不讓子民再返南國.
\newline
\newline
今天不少人憑著自己的思想製造某些宗教,自稱為基督教,但不是根據聖經而行,這些流行的新派神學使教會失去為神作見證的力量,這些行動正如耶羅波安自金牛一樣,以別樣事物代替了真神,他們如此作,不但害了自己,也是敗壞了教會.
\newline
\newline


\newpage

\section{第41屆~港九培靈會~胡恩德~第九講}
\label{sec:1486}
\textbf{神的見證人}
\newline
\newline
連結: \href{https://www.hkbibleconference.org/session-message/view/1486}{\texttt{https://www.hkbibleconference.org/session-message/view/1486}}
\newline
\newline
\hyperref[sec:1485]{< < < PREV SERMON < < <}
~
\hyperref[sec:toc]{[返目錄]}
~
\hyperref[sec:1487]{> > > NEXT SERMON > > >}
\newline
\newline
列王記上 13:1-13:34
\newline
\begin{longtable}{cl}
\hline
\hline
章節 & 經文 (和合本修訂版)\\
\hline
13:1 & \begin{tabularx}{0.7\textwidth}{X} 看哪，有一個神人遵照耶和華的話從猶大來到伯特利。耶羅波安正站在壇旁燒香； \end{tabularx} \\ \\ \relax
13:2 & \begin{tabularx}{0.7\textwidth}{X} 神人遵照耶和華的話向壇呼叫，說：「壇哪，壇哪！耶和華如此說：『看哪，大衛家必生一個兒子，名叫約西亞，他必將在你上面燒香的丘壇祭司，宰殺在你上面，人的骨頭也必燒在你上面。』」 \end{tabularx} \\ \\ \relax
13:3 & \begin{tabularx}{0.7\textwidth}{X} 當日，神人設個預兆，說：「這是耶和華說的預兆：『看哪，這壇必破裂，壇上的灰必傾倒出來。』」 \end{tabularx} \\ \\ \relax
13:4 & \begin{tabularx}{0.7\textwidth}{X} 耶羅波安王聽見神人向伯特利的壇呼叫的話，就從壇上伸手，說：「拿住他！」王向神人所伸的手卻萎縮了，不能彎回。 \end{tabularx} \\ \\ \relax
13:5 & \begin{tabularx}{0.7\textwidth}{X} 壇也破裂了，壇上的灰傾倒出來，正如神人遵照耶和華的話所設的預兆。 \end{tabularx} \\ \\ \relax
13:6 & \begin{tabularx}{0.7\textwidth}{X} 王對神人說：「請你為我禱告，向耶和華－你的神懇求恩惠，使我的手復原。」於是神人向耶和華懇求，王的手就復原了，如平常一樣。 \end{tabularx} \\ \\ \relax
13:7 & \begin{tabularx}{0.7\textwidth}{X} 王對神人說：「請你跟我回宮，讓你恢復心力，我必給你賞賜。」 \end{tabularx} \\ \\ \relax
13:8 & \begin{tabularx}{0.7\textwidth}{X} 神人對王說：「你就是把你一半的王宮給我，我也不跟你進去，也不在這地方吃飯喝水， \end{tabularx} \\ \\ \relax
13:9 & \begin{tabularx}{0.7\textwidth}{X} 因為耶和華的話這樣吩咐我說：『不可吃飯喝水，也不可從你去的原路回來。』」 \end{tabularx} \\ \\ \relax
13:10 & \begin{tabularx}{0.7\textwidth}{X} 於是神人從別的路回去，不從他到伯特利來的原路回去。 \end{tabularx} \\ \\ \relax
13:11 & \begin{tabularx}{0.7\textwidth}{X} 有一個老先知住在伯特利，他的兒子來，把神人當日在伯特利所做的一切事和他向王所說的話，都告訴了父親。 \end{tabularx} \\ \\ \relax
13:12 & \begin{tabularx}{0.7\textwidth}{X} 父親對他們說：「神人從哪條路去了呢？」他的兒子都看到從猶大來的神人所去的路。 \end{tabularx} \\ \\ \relax
13:13 & \begin{tabularx}{0.7\textwidth}{X} 老先知吩咐兒子說：「你們為我備驢。」他們備好了驢，他就騎上， \end{tabularx} \\ \\ \relax
13:14 & \begin{tabularx}{0.7\textwidth}{X} 去追神人，遇見神人坐在橡樹底下，就對他說：「你是不是從猶大來的神人？」他說：「是我。」 \end{tabularx} \\ \\ \relax
13:15 & \begin{tabularx}{0.7\textwidth}{X} 老先知對他說：「請你跟我一起回家吃飯。」 \end{tabularx} \\ \\ \relax
13:16 & \begin{tabularx}{0.7\textwidth}{X} 神人說：「我不能跟你回去，與你同行，也不能在這地方跟你一起吃飯喝水， \end{tabularx} \\ \\ \relax
13:17 & \begin{tabularx}{0.7\textwidth}{X} 因為有耶和華的話吩咐我說：『你在那裡不可吃飯喝水，也不可從你去的原路回來。』」 \end{tabularx} \\ \\ \relax
13:18 & \begin{tabularx}{0.7\textwidth}{X} 老先知對他說：「我也是先知，和你一樣。有天使遵照耶和華的話對我說：『你去帶他一同回你的家，給他吃飯喝水。』」老先知在欺騙他。 \end{tabularx} \\ \\ \relax
13:19 & \begin{tabularx}{0.7\textwidth}{X} 於是神人跟老先知回去，在他家裡吃飯喝水。 \end{tabularx} \\ \\ \relax
13:20 & \begin{tabularx}{0.7\textwidth}{X} 他們坐席的時候，耶和華的話臨到那帶神人回來的先知， \end{tabularx} \\ \\ \relax
13:21 & \begin{tabularx}{0.7\textwidth}{X} 他就對從猶大來的神人宣告說：「耶和華如此說：『你既違背耶和華的指示，不遵守耶和華－你神的命令， \end{tabularx} \\ \\ \relax
13:22 & \begin{tabularx}{0.7\textwidth}{X} 反倒回來，在耶和華禁止你吃飯喝水的地方吃了飯喝了水，因此你的屍體必不得葬在你祖先的墳墓裡。』」 \end{tabularx} \\ \\ \relax
13:23 & \begin{tabularx}{0.7\textwidth}{X} 神人吃喝完了，老先知為他帶回來的先知備驢。 \end{tabularx} \\ \\ \relax
13:24 & \begin{tabularx}{0.7\textwidth}{X} 神人就去了，在路上有隻獅子遇見他，把他咬死。他的屍體倒在路上，驢站在屍體旁邊，獅子也站在屍體旁邊。 \end{tabularx} \\ \\ \relax
13:25 & \begin{tabularx}{0.7\textwidth}{X} 看哪，有人經過，看見屍體倒在路上，獅子站在屍體旁邊，就來到老先知所住的城裡述說這事。 \end{tabularx} \\ \\ \relax
13:26 & \begin{tabularx}{0.7\textwidth}{X} 那帶神人回來的先知聽見了，就說：「這是那違背了耶和華指示的神人，所以耶和華把他交給獅子；獅子撕裂他，咬死他，正如耶和華對他說的話。」 \end{tabularx} \\ \\ \relax
13:27 & \begin{tabularx}{0.7\textwidth}{X} 老先知吩咐他兒子說：「你們為我備驢。」他們就備了驢。 \end{tabularx} \\ \\ \relax
13:28 & \begin{tabularx}{0.7\textwidth}{X} 他去了，發現神人的屍體倒在路上，驢和獅子站在屍體旁邊，獅子卻沒有吃屍體，也沒有撕裂驢。 \end{tabularx} \\ \\ \relax
13:29 & \begin{tabularx}{0.7\textwidth}{X} 老先知把神人的屍體抬起，馱在驢上，帶回自己的城裡，要為他哀哭，為他安葬。 \end{tabularx} \\ \\ \relax
13:30 & \begin{tabularx}{0.7\textwidth}{X} 老先知把屍體葬在自己的墳裡，為他哀哭，說：「哀哉！我的弟兄啊！」 \end{tabularx} \\ \\ \relax
13:31 & \begin{tabularx}{0.7\textwidth}{X} 安葬之後，老先知對他兒子說：「我死了，你們要把我葬在神人所葬的墳裡，使我的屍骨在他的屍骨旁邊， \end{tabularx} \\ \\ \relax
13:32 & \begin{tabularx}{0.7\textwidth}{X} 因為他遵照耶和華的話，指著伯特利的壇和撒瑪利亞各城丘壇神殿所宣告的話必定應驗。」 \end{tabularx} \\ \\ \relax
13:33 & \begin{tabularx}{0.7\textwidth}{X} 這事以後，耶羅波安仍不離開他的惡道，立平民百姓為丘壇的祭司；凡願意的，他都分別為聖，立為丘壇的祭司。 \end{tabularx} \\ \\ \relax
13:34 & \begin{tabularx}{0.7\textwidth}{X} 這事使耶羅波安的家陷入罪裡，甚至他的家被剪除，從地面上消滅了。 \end{tabularx} \\ \\
[1ex]
\hline
\hline
\end{longtable}
在主看來全世界的人都沒有希望.挪亞洪水的刑罰,雖然如此重,但人不肯轉回；洪水以後,他們便集體行動,處事與神旨意相違,神要他們分散,神也承認人集團的力量大,他們就用這力量來反對神,迅速展開拜偶像的事.現今考古學已發現亞拉伯罕時代拜偶像極其利害,製出種種假神,假道理,雖然他們知道掌管洪水是一位真正創造萬物之主,但他們的心仍未歸向神.但無論如何,神不會因人的悖逆而離棄人類,他對我們有一個愛的心腸,正如他對以法蓮說:「我怎能棄絕你們呢」?神愛他們,如果他們肯順服主,神便賜給他們豐盛榮耀的福氣,使他們在列邦中為神作見證,引領他們歸向神.
\newline
\newline
然而撒旦卻要在這事上作破壞的工作,他把以色列民與世人的界線抹去,引誘以色列人,使他們與異邦沒有分別,以致外邦人不能從他們身上看見神的作為.以色列人進入迦南便違背神旨,與迦南人混雜在一起,彼此結親,因此拜偶像之物便進入他們家中.士師時代幾次的大墮落,神又施恩使他們幾次得著復興,這便反映當時他們的情形.後來他們還效法鄰邦要求立王,不以神為他們的王,但神仍然在如此惡劣的環境中行事,使祂的目的達到.所以在他們立第二任王時,神選出合祂心意的大衛王,他是一位遵行神旨蒙福的王,他在惡劣的民族中及其他非常凶悍的民族裡,表明他是順服主,依靠主的.所羅門初期,他能依靠主得到主豐富的智慧,神在他當中顯出祂莫大的恩惠,使他的名聲傳至中東各國,但魔鬼在身上作工,使極多外邦女子進到他宮中,又把他們所事奉的假神偶像帶進去,因而他的國便分裂.他離開神,因而招致失敗,一方面是神刑罰他；另一方面也是與他本身有關,天地間有因果的定律存在,所羅門是自吃其果.
\newline
\newline
國分裂以後,他兒子羅波安繼位,國家的命運可謂每下愈況,直至國家亡.特別北國,他們沒有出過一位賢君,每一位君王都拜耶羅波安所立的那金牛,這金牛從始至終沒有倒下來,魔鬼放這件東西在那裡,一直緊抓人心,直至他們滅亡為止.這十個支派分散各地,不知所蹤,神報應他們,因為他們破壞了神在他們身上作見證的祈望.列王記給我們看見人失敗,把神的見證破壞,但實際上神並不失敗,這也是神的計劃,神的智慧.
\newline
\newline
神的子民好像被消滅,不能為祂作見證,但我們從保羅的解釋裡便看見神的計劃,原來亞伯拉罕肉體所生的子孫不都是真以色列人,神另外揀選了一個小團體,要他們個人在神面前作個別忠心的見證.
\newline
\newline
列王記上卷可見人的失敗,他們不能站在異邦地為神作見證,反而在異邦的假神面前低頭下拜燒香,築祭壇.雖然滿篇都是失敗,但其中仍可看見神的作為.在列王記上十二章以後,便特別看見「神人」出現,這些是先知,那時代特別給他們起名叫「神人」,神人就是在這人身上有神的作為,說話帶有神能力的表現,他在人面前可以代表神.從這些神人裡我們便可知道神並沒有完全離開這個民族,神沒有完全撇下祂的見證,正如保羅所說:凡事都為你們,神每一件事目的都是要使我們得福.
\newline
\newline
在這些失敗的歷史中,神人的出現便成了書卷中非常突出的一個主題.在列王記上第十三章有一個不知名的神人他是一個與神關係失敗的神人,然而他已完成神托負他的工作.在大背道時,失去見證的能力,圈內的圈便突出來,那些真正為神作見證的人,雖然數目微小,但他們要起來代替整個民族的見證.整體的見證打破,小數的見證仍然存在.現今教會的情形也相仿,教會離開聖經的真理,使我們感到異常失望,雖然如此,罪惡滿盈,在其中仍有人起來為神作見證.
\newline
\newline
最近印度前途非常暗淡,本年印度政府的措施,其中通過一件大案件；就是任何人不得改變其宗教,此案件目的是要對付天主教及基督教,免至印度教的人,改變其信仰,凡改變別人宗教信仰者要受處罰或罰款,已有浸會的牧師及傳道人被政府控訴,又有天主教神父被囚預備處分.而且在一兩年之內,印度的傳教士要完全離開印度.如此看來,印度仍否有希望?神的見證是否完畢?雖然外表看來人要離開那裡,但神的工作仍會存留,正如神在中古時代保守了那些地下教會的情形一樣.
\newline
\newline
列王記第十三章開始,我們便看見神的先知如何起來為主作見證,這對我們今日教會甚有意思,按表面看,現今的教會似乎被外面勢力掩蓋,另一方面教會裡有發酵的東西,錯誤離教的道理,使人無心向主,藐視聖經,但神仍有他奇妙的作為.
\newline
\newline
神的心意是要在以色列中作王,又要在王的身上作工,使眾民都能認識神.以色列諸王應作神的僕人帶領子民,但他們失敗,直至亡國,也不能達到神的心意,所以神要在王室以外另立一些見證人,神永不失敗,祂必須維持祂的見證,人雖離開祂但祂沒有退出.神可以指著約伯向撒旦誇口,雖然撒旦攻擊約伯,約伯仍不犯罪咒詛神；神也可以用指頭指著但以理,使但以理成為他的見證,因為他能謹守神的律法,行在神的道中不偏左右；現在這大背道黑暗的時代中,並非整個民族失敗,神仍然可以指著以利亞向撒旦誇口,而且除了以利亞一人以外,祂仍留下七個千人未向巴力屈膝,我們是否能成為神的據點,給神的指頭指著我們向撒旦誇口.我們不但要在信仰上嚴格,甚至生活也需要嚴謹地合乎神的標準,遵照神的旨意而行.
\newline
\newline
第十三章裡所記載的神人要到伯特利去,因為耶羅波安替他的偶像獻祭,神把當說的話交託與神人,這些話甚難出口,但他仍要傳講.作為一個神人,是要為神作見證,不能隨從今世的風俗,潮流,要站在神面前作見證,使人認識神,世界的人輕看我們,藐視我們,正如保羅說:「猶大人是要神蹟,希利尼人是求智慧,我們卻是傳釘十字架的基督,在猶太人為絆腳石,在外邦人為愚拙......」保羅傳講神的道,靠神勝過一切.神人靠著主大膽把一切當說的講出來:「向壇呼叫說:壇哪,壇哪,耶和華如此說,大衛家必生一個兒子名叫約西亞,他必將邱壇的祭司,就是在你上面燒香的,殺在你上面,人的骨頭也必燒在你上面.」保羅對腓立比教會說:「有人行事與十字架為敵,我從前告訴你們,現在又流淚告訴你們,這些人的結局是沉淪」.我們能否像這神人一樣,在反對基督的人面前宣佈刑罰,給他們一個實在的警告.神驗中我們,將福音托付我們,我們便要向神有忠心的表現,作為一個神人,忠心是重要素質之一.
\newline
\newline
神人宣佈刑罰以後,耶羅波安便伸手命人把他捉下來,怎他伸出的手枯乾了,神人為他禱告後才復原,王請他吃飯,他不肯進去,以後城裡有一個老先知,說假話欺騙他,使他進城吃飯.這神人又不自己求問神,便順老先知的命進城去了,結果他失敗見噬於獅,這神人的失敗有神的作用在他身上,原來神要在以色列民中,藉著多方面的教導,表明神是活的神,叫他們能順服神.
\newline
\newline
在大背道墮落時,神給他們有重要的神磧,作為,希望挽回他們,但耶羅波安仍然犯罪,把自己的家及以色列家陷在罪惡滅亡中.我們應尊重主的話,聽祂的道而行,新約時代比較少用威嚴的方法來教導我們,乃用基督替死的大愛來激勵我們,我們應更加行在其中,持守神的道,叫祂能向撒旦指明我們是屬他的人,榮耀祂的聖名.
\newline
\newline


\newpage

\section{第41屆~港九培靈會~胡恩德~第十講}
\label{sec:1487}
\textbf{必按公義審判我們}
\newline
\newline
連結: \href{https://www.hkbibleconference.org/session-message/view/1487}{\texttt{https://www.hkbibleconference.org/session-message/view/1487}}
\newline
\newline
\hyperref[sec:1486]{< < < PREV SERMON < < <}
~
\hyperref[sec:toc]{[返目錄]}
~
\hyperref[sec:1488]{> > > NEXT SERMON > > >}
\newline
\newline
列王記上 15:29-15:30
\newline
\begin{longtable}{cl}
\hline
\hline
章節 & 經文 (和合本修訂版)\\
\hline
15:29 & \begin{tabularx}{0.7\textwidth}{X} 巴沙一作王就殺了耶羅波安全家，耶羅波安家凡有氣息的，一個也沒有留下，都殺滅了，正如耶和華藉他僕人示羅人亞希雅所說的話。 \end{tabularx} \\ \\ \relax
15:30 & \begin{tabularx}{0.7\textwidth}{X} 這是因為耶羅波安所犯的罪，他使以色列陷入罪裡，激怒了耶和華－以色列的神。 \end{tabularx} \\ \\
[1ex]
\hline
\hline
\end{longtable}
列王記把分國以後的猶大及以色列歷史並行而寫,在這幾段的經文中特別側重北方以色列國的情形.他們越久越得罪神,直至第四朝暗利作王,他得罪神比以前更甚,後來其子亞哈作王,亞哈在神面前行惡尤比列祖更甚,一代不如一代,所以他們的家及以色列國都得到不良的結局.他們以前三朝,得罪神,神把他們除滅.
\newline
\newline
亞哈是第四朝,比以前的王更壞,他們效法異族的風俗,本來他們的責任是要吸引別人到主面前,怎料反被別人吸引,以前他們種種的失敗,便成為我們現在的借鏡.
\newline
\newline
保羅寫信給哥林多信徒流淚告訴他們:「濫交是敗壞善行的」,隨便與不信的人建立密切的友誼,很可能敗壞自己的善行.信與不信的人不能同負一軛,我們信主的人不是要做隱士,但我們要小心不可向他們同流合污,教會裡很多失敗,就是與世界分得不太清楚.
\newline
\newline
「教會」這字在新約原文的意義,是被呼召出來的一個團體,是神在世人當中召出來的.我們被召出來,仍要回去作見證,但我們的身份必須要確立,別人以我們為奇怪,是免不了的,因為根本上我們與他們完全不同,我們有主的生命,聖靈,行事當然與他們有分別.
\newline
\newline
「暗利用二他連得銀子,向撒瑪買了撒瑪利亞山,在山上造城......」,在以色列人未進入迦南以前,神曾經吩咐摩西在他們進去以後,把土地好好的分給眾民,又用抽籤分地以示公道,神命定地是永不能賣,地界也不能移,成為他們長久的產業,但暗利不重視聖經的話,任憑自己的意思去行,買了撒瑪利亞山地.正如今日教會中有不少領袖,他們不尊重主的話行事.亞哈,貪愛拿伯的葡萄園,欲得此園,但拿伯對他說:我敬畏耶和華,萬不敢以先祖留下來的產業賣給你.
\newline
\newline
摩西在申命記十八章,傳神的命令說:「若有人作王應把律法書抄寫一份,平生頌讀」,但暗利,亞哈只憑自己的想法去做,因此聖經除了寫明他拜金牛的罪惡外,還以他不遵行神的律法,不聽神的話為大罪,使他不配作神子民的領袖.現今這個世代,人需要尊重神的話.
\newline
\newline
耶穌也行主的道,遵行神的話,祂認為神的話是不能改變的,「經上的話是不能廢掉」.當魔鬼來試探耶穌時,耶穌對他說「經上記著說:人活著不是單靠食物」；「經上記著說:不可試探主你的神；」耶穌用聖經的話應付仇敵魔鬼.他也是行在聖經中,以聖經為他生活的原則,完全服在聖經下面.「經上記著說」,原文只得一字,其意思初時沒有人知曉,直到近二三百年也不知道,考古學家在猶大地發掘到一些古物,就是羅馬時代用希臘文所寫的紙章,其中有一稿件有這「經上記著說」的字,原來當時的百姓有此習俗,若有兩人到攻府那裡立約,把約寫好以後由政府官員閱讀一遍,兩者同意後便在他們所簽的約上寫上「經上記著說」這個字,然後將約分成兩半,各執一邊都有這個字作印監,這字是代表政府所確定的,不能更改,神的話有神的權威,也是不能更改的.
\newline
\newline
今天作為基督徒的應當跟隨主,牢記聖經中的話,照聖經的法則來生活行事為人.有些人讀神的話,只用欣賞的態度,我們讀經聽道,乃是從神面前領受祂寶貴的真理,不能隨便,暗利不讀,不聽,不行聖經的話,這便成了他一生大惡中的一個.
\newline
\newline
暗利的兒子亞哈作王,他除了敬拜金牛,娶外邦女子為妻,這些惡事外,他還作了一件聖經明載是神所不歡喜的事,他的部下不敬拜神,與亞哈一樣,「...... 伯特利人希伊勒重修耶利哥城,立根基的時候,喪了長子亞比蘭,安門的時候,喪了幼子西割,正如耶和華藉嫩的兒子約書亞所說的話.」上面背道,下面也是背道,上面的人不尊重聖經,下面的人也不尊重聖經,結果國家越來越衰退,離主越來越遠.
\newline
\newline
以色列人輕看神的話語,又違背主的道,神應把他們放在一邊棄絕他們,自古至今,不少的人起來反對神,甚至逼害那些信主的猶太人,雖然環境如此惡劣,但神仍沒有棄絕他們,因為神的揀選是沒有後悔,堅定不移的,他們違背神的約,所以神便照祂所定的法則來刑罰他們,但他卻沒有離棄他們,只是用愛心來干涉他們.
\newline
\newline
在整個情勢十分惡劣下面,第十七章突然出現了一位人物,聖經並沒有言及他的來歷,只是把他稍作介紹:「基列寄居的提斯比人以利亞,對亞哈說我指著所事奉永生耶和華以色列的神起誓,這幾年我若不禱告,必不降露不下雨.」現今神的子民犯罪離棄神,神不丟棄他們,反而興起先知來,藉著祂的僕人向他們說話,希望他們能如浪子回頭,神對他子民真有極大的恩情.
\newline
\newline
我們的神現今對待我們也是如此恩愛,他永不離開我們,不斷用祂的僕人傳出祂的信息感動我們得著復興,我們因著祂有盼望.我們從以利亞的身上可以得著無限的安慰,看出神仍然關心我們,為要使我們得著復興,祂願意在末日的世代裡向我們說話,好像昔日與亞哈王說話一樣.
\newline
\newline


\newpage

\section{第41屆~港九培靈會~胡恩德~第十一講}
\label{sec:1488}
\textbf{黑暗背道中神仍作工}
\newline
\newline
連結: \href{https://www.hkbibleconference.org/session-message/view/1488}{\texttt{https://www.hkbibleconference.org/session-message/view/1488}}
\newline
\newline
\hyperref[sec:1487]{< < < PREV SERMON < < <}
~
\hyperref[sec:toc]{[返目錄]}
~
\hyperref[sec:1489]{> > > NEXT SERMON > > >}
\newline
\newline
列王記上 17:1-17:24
\newline
\begin{longtable}{cl}
\hline
\hline
章節 & 經文 (和合本修訂版)\\
\hline
17:1 & \begin{tabularx}{0.7\textwidth}{X} 住在基列的提斯比人以利亞對亞哈說：「我指著所事奉永生的耶和華－以色列的神起誓，這幾年我若不禱告，必不降露水，也不下雨。」 \end{tabularx} \\ \\ \relax
17:2 & \begin{tabularx}{0.7\textwidth}{X} 耶和華的話臨到以利亞，說： \end{tabularx} \\ \\ \relax
17:3 & \begin{tabularx}{0.7\textwidth}{X} 「你離開這裡往東去，躲在約旦河東邊的基立溪旁。 \end{tabularx} \\ \\ \relax
17:4 & \begin{tabularx}{0.7\textwidth}{X} 你要喝那溪裡的水，我已吩咐烏鴉在那裡供養你。」 \end{tabularx} \\ \\ \relax
17:5 & \begin{tabularx}{0.7\textwidth}{X} 於是以利亞去了，他遵照耶和華的話做，去住在約旦河東的基立溪旁。 \end{tabularx} \\ \\ \relax
17:6 & \begin{tabularx}{0.7\textwidth}{X} 烏鴉早上給他叼餅和肉來，晚上也有餅和肉，他又喝溪裡的水。 \end{tabularx} \\ \\ \relax
17:7 & \begin{tabularx}{0.7\textwidth}{X} 過了些日子溪水乾了，因為雨沒有下在地上。 \end{tabularx} \\ \\ \relax
17:8 & \begin{tabularx}{0.7\textwidth}{X} 耶和華的話臨到他，說： \end{tabularx} \\ \\ \relax
17:9 & \begin{tabularx}{0.7\textwidth}{X} 「你起身到西頓的撒勒法去，住在那裡，看哪，我已吩咐那裡的一個寡婦供養你。」 \end{tabularx} \\ \\ \relax
17:10 & \begin{tabularx}{0.7\textwidth}{X} 以利亞就起身往撒勒法去。他到了城門，看哪，有一個寡婦在那裡撿柴。以利亞呼喚她說：「請你用器皿取點水來給我喝。」 \end{tabularx} \\ \\ \relax
17:11 & \begin{tabularx}{0.7\textwidth}{X} 她去取水的時候，以利亞又呼喚她說：「請你手裡也拿點餅來給我。」 \end{tabularx} \\ \\ \relax
17:12 & \begin{tabularx}{0.7\textwidth}{X} 她說：「我指著永生的耶和華－你的神起誓，我沒有餅，罈內只有一把麵，瓶裡只有一點油。看哪，我去找兩根柴，帶回家為我和我兒子做餅。我們吃了，就等死吧！」 \end{tabularx} \\ \\ \relax
17:13 & \begin{tabularx}{0.7\textwidth}{X} 以利亞對她說：「不要怕！你去照你所說的做吧！只要先為我做一個小餅，拿來給我，然後為你和你的兒子做餅； \end{tabularx} \\ \\ \relax
17:14 & \begin{tabularx}{0.7\textwidth}{X} 因為耶和華－以色列的神如此說：『罈內的麵必不用盡，瓶裡的油必不短缺，直到耶和華使雨降在地上的日子。』」 \end{tabularx} \\ \\ \relax
17:15 & \begin{tabularx}{0.7\textwidth}{X} 婦人就照以利亞的話去做。她和以利亞，以及她家中的人，吃了許多日子。 \end{tabularx} \\ \\ \relax
17:16 & \begin{tabularx}{0.7\textwidth}{X} 罈內的麵果然沒有用盡，瓶裡的油也不短缺，正如耶和華藉以利亞所說的話。 \end{tabularx} \\ \\ \relax
17:17 & \begin{tabularx}{0.7\textwidth}{X} 這事以後，那婦人，就是那家的女主人，她的兒子病了，病得很重，甚至沒有氣息。 \end{tabularx} \\ \\ \relax
17:18 & \begin{tabularx}{0.7\textwidth}{X} 婦人對以利亞說：「神人哪，我跟你有甚麼關係，你竟到我這裡來，使神記起我的罪，以致我的兒子死了呢？」 \end{tabularx} \\ \\ \relax
17:19 & \begin{tabularx}{0.7\textwidth}{X} 以利亞對她說：「把你兒子交給我。」以利亞就從婦人懷中接過孩子來，抱到他所住的頂樓，放在自己的床上。 \end{tabularx} \\ \\ \relax
17:20 & \begin{tabularx}{0.7\textwidth}{X} 他求告耶和華說：「耶和華－我的神啊，我寄居在這寡婦的家裡，你卻降禍於她，使她的兒子死了嗎？」 \end{tabularx} \\ \\ \relax
17:21 & \begin{tabularx}{0.7\textwidth}{X} 以利亞三次伏在孩子的身上，求告耶和華說：「耶和華－我的神啊，求你使這孩子的生命歸回給他吧！」 \end{tabularx} \\ \\ \relax
17:22 & \begin{tabularx}{0.7\textwidth}{X} 耶和華聽了以利亞的呼求，孩子的生命歸回給他，他就活了。 \end{tabularx} \\ \\ \relax
17:23 & \begin{tabularx}{0.7\textwidth}{X} 以利亞把孩子從樓上抱下來，進了房間交給他母親，說：「看，你的兒子活了！」 \end{tabularx} \\ \\ \relax
17:24 & \begin{tabularx}{0.7\textwidth}{X} 婦人對以利亞說：「現在我知道你是神人，耶和華藉你口所說的話是真的。」 \end{tabularx} \\ \\
[1ex]
\hline
\hline
\end{longtable}
北方的以色列國,由亞哈王帶頭離開神,敬拜巴力,特別王后耶洗別尤其可惡,引進巴力要全國敬拜,又殺神的先知,拆毀神的祭壇,但神在這大背道的日子中仍作工,保存一些先知不被她所殺,就是那少數忠心為主作證的人.
\newline
\newline
在這時期,背道反對主相當利害,以別神代替神,國政敗壞,此乃離棄神的結果,但他們仍不肯悔改,正如現今的人一樣.以往不信的人以偶像代替獨一的真神,現在則以「無神論」,「唯物論」,「進化論」來代替聖經.又把人抬高,萬事以人為首,人高過一切,人代替神.現今背道的事也滲入教會之中.以為「神」是思想的產品.反對神的勢力也滲透入教會之中.但這並非表明神是失敗.神從不失敗,祂已選召一班小數的人,這些人自古至今從未動搖過,保羅對提摩太說.當時有些錯誤的道理在教會中發酵,這些是世俗的虛談,屬世虛渺的主張,他們的話像毒瘡,越爛越大,敗壞了好些人的信心.教會裡很多人的信心受影響跌倒.表面看去好像主失敗了,然而聖經的話安慰我們說:「神穩固的根基立住了,在上有這印記說,主認識誰是屬祂的人」.這是不能動搖的根基,神對教會的計劃,是祈望那班小數忠心為祂作見證的人.無論教會的情勢如何,神的作為仍要彰顯.現今雖然很多地方受到反基督勢力的壓迫,然而神在那裡保留一些屬神的人.聖經明言:在背道的日子,似乎沒有希望,但在上有穩固的根基,就是神認識誰是屬祂的,神的眼遍察全地,要幫助那些向他心存誠實的人.
\newline
\newline
以西結作先知的時代,北方以色列國已淪亡,南面猶大國已有部份被擄到巴比倫,神讓他在被擄的民中作神的僕人,傳神的信息,為神作見證,又留下耶利米在猶大國耶路撒冷那裡作先知,對未被擄的民說話.有一次先知以西結在異象中被神帶回猶大本國的聖殿中,他看見牆上畫滿了種種污穢之物,如走獸爬蟲等,又看見許多祭司在那裡燒香,在殿外向東拜日頭,又在北面院外有婦女手執樹枝向假神哭泣,這乃是巴力先知敬拜的儀式,他們在作可惡的事.以西結看完以後,便聽見神大聲呼叫,命天使行遍全地察看,凡為此城的罪惡哭泣的,便在他們額上印上記號,又吩咐另一位天使隨從出去,把那些額上沒有印記的人殺掉,雖然當時可謂是大背道之日,耶路撒冷由大至小都犯罪,但神認識誰是屬祂的人.
\newline
\newline
亞哈時代的敗壞甚多,王后耶洗別強迫百姓拜偶像,又殺先知,好像要把神的見證挪去,但神仍然找出祂的見證人提斯比人以利亞來,為祂作有力的見證.他出來所講的話十分沉重,他宣告亞哈狀及刑罰,「......我指著所事奉永生耶和華以色列的神起誓,這幾年我若不禱告,必不降露不下雨.」
\newline
\newline
神現在要說話,在全國背道時,祂表明自己審判的行動,公義的行動.未施行刑罰以先,先差屬祂的人向首領宣告,使他們知道在他們中間,仍有一個神人,他有神的權柄,神差遣以利亞去見亞哈,目的是要亞哈知道,神與他甚近,就是在他眼前,神未曾被他們推翻,他的寶座是在天上,他的寶座永永遠遠,天地要滅沒,神卻要長存.魔鬼總希望把主推翻,使整個教會的信仰失掉,但這是沒有可能的事.有很多青年人被世界所纏繞,以致信仰失掉,熱心失落,為神作見證的力量也失落,只用外表的事來掩飾自己,香港在一九二七至一九三七時多人蒙恩被主復興,但相隔數十年,竟有人在此路上失敗,不能專一事奉主,把世界混入純正的信仰中,教會已滿了反對神的道理.
\newline
\newline
在這時候,神要使用祂那圈內之小圈,誠誠實實心裡愛主的人,為祂作見證,要使世人知道神仍在人間.今日神要尋找像以利亞及七千未向巴力屈膝的人,忠心地在末世為祂作見證,當我們細心觀察現今的局勢時,我們便曉得萬物的結局近了,很多力量要集合起來,成為集體反對神的力量,我們再不能如此放鬆,把心放在世界上,否則我們便不認識現在緊張的情勢,沒有勇士的心為主作見證.我們無論作何事,要想念到神救我們,我們人生的目的是否以神為我們的主?末世時代裡,神需要我們奮勇起來為祂工作,我們不可再為自己打算,隨從自己喜好的去行,因為我們所面對的時代,是風暴的時代,初時以為平安無事,恐怕死亡的陰影隨即來到,我們面臨這一個比以利亞更利害的爭戰時代,我們的責任相當重大,應隨時作準備,不要像以利亞一樣.今天神也要使用這樣忠心的人為祂的真理作見證,我們要好好地預備自己.
\newline
\newline
魔鬼在各方面利用人的知識學說來反對主,促進它的勢力,以罪中之樂來吸引人,魔鬼已用他空前厲害的方法來吸引我們,攻擊我們,使我們不知不覺中便陷入它的陷阱.神要找出英勇的精兵像以利亞一樣,起來為祂的真理作戰.
\newline
\newline
以利亞與主的關係聯合得非常好,他能靠著主的名,禱告下雨就下雨,不降雨就不降雨,他用禱告來行使神的旨意及權柄,證明他是與神同心,同工,同行.神的心意在聖經裡藉著啟示和文字表達出來,但倘若一個人不願讀聖經,也不追求主的話,恐怕他必不能與主一同享受,假若一個人肯與神同心便能與神同在,我們才有主的權柄運用在我們身上.昔日彼得有權柄吩咐那瘸腳的人起身行走,他的權柄是來自拿撒勒人耶穌,因他與主聯合便能奉祂的名行事,很可惜今天不少的教會,金也有,銀也有,但沒有主的權柄.
\newline
\newline
一個人能與主同行,便能與祂一同使用權柄,正如以利亞一樣,以利亞開始說的那句話,已把他的身份表明出來,他是事奉永生耶和華神的人.若我們一心一意服事主,被主的愛激勵,我們必得主的使用,不但只是事奉,生活上也是如此.「事奉」英文直接的意思:就是我侍立在主人的面前,接受主的支配.如果我們真的這樣行,我們必被主使用,成為祂的精兵,在末世裡,使人知道有神僕人的存在,這是我們重要的任務.
\newline
\newline


\newpage

\section{第41屆~港九培靈會~胡恩德~第十二講}
\label{sec:1489}
\textbf{禱告的事奉}
\newline
\newline
連結: \href{https://www.hkbibleconference.org/session-message/view/1489}{\texttt{https://www.hkbibleconference.org/session-message/view/1489}}
\newline
\newline
\hyperref[sec:1488]{< < < PREV SERMON < < <}
~
\hyperref[sec:toc]{[返目錄]}
~
\hyperref[sec:1490]{> > > NEXT SERMON > > >}
\newline
\newline
列王記上 17:17-17:24
\newline
\begin{longtable}{cl}
\hline
\hline
章節 & 經文 (和合本修訂版)\\
\hline
17:17 & \begin{tabularx}{0.7\textwidth}{X} 這事以後，那婦人，就是那家的女主人，她的兒子病了，病得很重，甚至沒有氣息。 \end{tabularx} \\ \\ \relax
17:18 & \begin{tabularx}{0.7\textwidth}{X} 婦人對以利亞說：「神人哪，我跟你有甚麼關係，你竟到我這裡來，使神記起我的罪，以致我的兒子死了呢？」 \end{tabularx} \\ \\ \relax
17:19 & \begin{tabularx}{0.7\textwidth}{X} 以利亞對她說：「把你兒子交給我。」以利亞就從婦人懷中接過孩子來，抱到他所住的頂樓，放在自己的床上。 \end{tabularx} \\ \\ \relax
17:20 & \begin{tabularx}{0.7\textwidth}{X} 他求告耶和華說：「耶和華－我的神啊，我寄居在這寡婦的家裡，你卻降禍於她，使她的兒子死了嗎？」 \end{tabularx} \\ \\ \relax
17:21 & \begin{tabularx}{0.7\textwidth}{X} 以利亞三次伏在孩子的身上，求告耶和華說：「耶和華－我的神啊，求你使這孩子的生命歸回給他吧！」 \end{tabularx} \\ \\ \relax
17:22 & \begin{tabularx}{0.7\textwidth}{X} 耶和華聽了以利亞的呼求，孩子的生命歸回給他，他就活了。 \end{tabularx} \\ \\ \relax
17:23 & \begin{tabularx}{0.7\textwidth}{X} 以利亞把孩子從樓上抱下來，進了房間交給他母親，說：「看，你的兒子活了！」 \end{tabularx} \\ \\ \relax
17:24 & \begin{tabularx}{0.7\textwidth}{X} 婦人對以利亞說：「現在我知道你是神人，耶和華藉你口所說的話是真的。」 \end{tabularx} \\ \\
[1ex]
\hline
\hline
\end{longtable}
以利亞的時代,是非常黑暗,大背道的日子.好像現今教會,外面不信的人,反對主極其利害,他們又把反對的理論帶進教會,使無知信徒離開真理.撒旦在最後的日子中,用盡祂的力量反對主,罪惡勢力加增,使許多基督徒跌倒.現在除了不信的力量外,還有很多迫害基督徒的力量,使基督徒改變信仰.撒旦藉著一個人來統治世界,此人保羅稱他為「不法者」,「沉淪之子」,「大罪人」,也就是那末世的敵基督.他有一種空前的力量,撒旦在末世時代要利用他來行事.神的教會只剩下一小群,從外表看,教會好像面臨失敗,但我們可以放心,神永遠不會失敗,神必用屬祂的人起來為祂作見證.
\newline
\newline
以西結在被擄的猶太人中為神傳言,神曾對他說:凡神吩咐他說的話,他都要說出來,子民的額比鋼鐵還硬,雖在荊棘蒺藜中仍不要驚怕,專心傳講,信息傳出以後,聽或不聽不要緊,但要他們知道在他們中間有耶和華的先知.為神作見證,說明神並非退出,其實全世界都在祂統治之內,神把這些見證人放在地面,使他們曉得神離他們不遠,就在他們的旁邊.雖然人悖逆不肯聽,但神仍作見證表明祂並未退出世界.
\newline
\newline
當黑暗勢力佔據整個時代,神也用祂的僕人作有力量的見證,這些人不一定很多,可能有些是作隱藏的工作,但神要藉著這少數人起來為祂作見證.神今天在這末後的日子中,也要如此尋找屬祂的人為祂作見證,正如以利亞一人為主作見證一樣.
\newline
\newline
這些見證人的資格,並非只靠說話作見證,乃與生活有密切的關係.我們在神面前如何生活,是否有力量?是否尊敬主的道?否則稍被人對付便遭失敗.以利亞是一位事奉永生耶和華神的人,他站立在神面前,尊重神的道,順服祂的旨意,接受神的吩咐及支配.今天我們作基督徒的,若不專一事奉神,我們的生活及見證必定失敗.我們若不能專一事奉神,漸漸我們就會變成服事自己.人若以自己為生活行事的中心,不以神為中心,就會忘記自己是被主救回來,要離棄偶像,又要事奉神等候主從天降臨.恐怕不少基督徒對於神只是頭腦上的認識,不是真真經歷主是一位又真又活的神,不知道神就是在我們旁邊離我們不遠,我們的工作,生活是有一位主人來管理的.
\newline
\newline
我們屬神的人要查考聖經,知道甚麼是祂所喜悅的,無論是個人或是家庭都要遵祂旨意而行我們所處的地位是僕人的地位,應完全在祂手中,祂在我們身上有支配的權柄,我們應該奉獻給主,全然作祂的僕人,一個肯事奉永生耶和華的神的人,才有能力為祂作見證.
\newline
\newline
以利亞的時代,有不少人在偶像宗教墮落中倒下去,站立不住,政府最高的元首要全國的人起來敬拜神,但以利亞卻站穩.以利亞能忠心事奉神,因為他曾見過主,認識神是如此有力量可尊敬的一位,所以他放膽一人走到亞哈面前宣佈他當得重重的刑罰.以利亞有所見有所知,使他不得不發出警告.當以利亞親近神,祂便把當傳的話給他,使他能直接作大無畏的見證.
\newline
\newline
以利亞是一個事奉神的人,他事奉的生活離不開禱告的生活.我們要過一個事奉的生活,就要與禱告的生活發生密切的關係.以利亞能夠忠心事奉,乃因他在禱告上花了不少的工夫.以利亞的成功與他禱告有關,若不是一個勇於禱告的人,絕不能獲得如此能力和成功.
\newline
\newline
雅各書第五章告訴我們,以利亞與我們有一樣性情的人,他懇切禱告求不下雨,雨就不下了三年零六個月,他這樣做並非靠一時之衝動,乃是根據神的話.神在摩西時曾經明說,他們若離開神,神便使池乾旱,以利亞看見子民如此背逆不道,他便求神對付他們,使他們覺悟.後來又藉著禱告,雨又再次降在地面上.以利亞的禱告是十分謙卑,十分懇切的,這不是外表上的功夫,乃是心裡屈膝,非常熱切地禱告,正如雅各書形容他懇切禱告.現今教會的禱告是否在乎外表,沒有誠懇?一個人若不是生活在神前,清楚對付罪惡,就沒有神的靈感動帶領,若與主不接近,就沒有禱告的心.以利亞懇切禱告,與主同心,這樣才能成為他忠心有能力的禱告者及見證人.我們必須從心裡誠實懇切的禱告神.在王下第十七章以利亞使寡婦的兒子復活,乃是伏在孩子的身上,三次禱告主,並不灰心.以利亞是一個禱告的人,我們曾否有禱告的心,為著神的事情神的百姓屈膝,好像把自己的靈魂都傾倒下來一樣.
\newline
\newline
聖經提到有一個人名叫以巴弗,他與保羅一同坐監,他不是傳道人,卻周圍向別人傳講福音,在歌羅西書保羅對歌羅西教會的人稱讚他,說到有他們的人以巴弗,以巴弗為他們痛苦,常常為他們竭力禱告,使他們能在神旨意上站立得穩.我們是否也把自己交出來作禱告的器皿,作禱告的事奉,我們禱告的生活是否相差很遠.主耶穌是一位十分曉得禱告的人,他禱告至四更天,甚至是通宵的禱告.如果我們研究以往一百年神偉大的作為,歸根結底就是有很多人肯付上禱告的代價,這件事給我們體驗到禱告與工作的關係是極其明顯的,一個人必須過著認真禱告的生活,才能站立在神前忠心地事奉祂.
\newline
\newline
在迦密山上試驗真假的神,以利亞以一人的力量扺擋八百五十名假先知,神的先知與巴力的先知較力爭勝,以利亞十分誠懇地禱告神,便得蒙神的應允,一個習慣以禱告事奉神的人,他能以禱告與神同工,又能得蒙神的應允,因為他的禱告是有效的,有能力的,唯有這樣的人才能得到的神的使用.
\newline
\newline
我們自己要一心一意追求認識神,除去一切與神與人之間的罪惡,使自己得蒙神的應允.現今的世代,神的工作一代不如一代,恐怕是缺少禱告的人,我們只能發出緊迫的呼聲,就是要向普世發問:「禱告的人在那裡?」
\newline
\newline


\newpage

\section{第41屆~港九培靈會~胡恩德~第十三講}
\label{sec:1490}
\textbf{為萬軍耶和華神大發熱心}
\newline
\newline
連結: \href{https://www.hkbibleconference.org/session-message/view/1490}{\texttt{https://www.hkbibleconference.org/session-message/view/1490}}
\newline
\newline
\hyperref[sec:1489]{< < < PREV SERMON < < <}
~
\hyperref[sec:toc]{[返目錄]}
~
\hyperref[sec:1477]{> > > NEXT SERMON > > >}
\newline
\newline
列王記上 19:9-19:18
\newline
\begin{longtable}{cl}
\hline
\hline
章節 & 經文 (和合本修訂版)\\
\hline
19:9 & \begin{tabularx}{0.7\textwidth}{X} 他在那裡進了一個洞，在洞中過夜。看哪，耶和華的話臨到他，說：「以利亞，你在這裡做甚麼？」 \end{tabularx} \\ \\ \relax
19:10 & \begin{tabularx}{0.7\textwidth}{X} 他說：「我為耶和華－萬軍之神大發熱心，因為以色列人背棄了你的約，毀壞了你的壇，用刀殺了你的先知，只剩下我一人，他們還要追殺我。」 \end{tabularx} \\ \\ \relax
19:11 & \begin{tabularx}{0.7\textwidth}{X} 耶和華說：「你出來站在山上，在耶和華面前。」看哪，耶和華從那裡經過。在耶和華面前有烈風大作，山崩石裂，耶和華卻不在風中；風後有地震，耶和華也不在其中； \end{tabularx} \\ \\ \relax
19:12 & \begin{tabularx}{0.7\textwidth}{X} 地震後有火，耶和華也不在火中；火以後，有輕微細小的聲音。 \end{tabularx} \\ \\ \relax
19:13 & \begin{tabularx}{0.7\textwidth}{X} 以利亞聽見，就用外衣蒙臉，出來站在洞口。聽啊，有聲音向他說：「以利亞，你在這裡做甚麼？」 \end{tabularx} \\ \\ \relax
19:14 & \begin{tabularx}{0.7\textwidth}{X} 他說：「我為耶和華－萬軍之神大發熱心，因為以色列人背棄了你的約，毀壞了你的壇，用刀殺了你的先知，只剩下我一人，他們還要追殺我。」 \end{tabularx} \\ \\ \relax
19:15 & \begin{tabularx}{0.7\textwidth}{X} 耶和華對他說：「去吧，從原路回去，往大馬士革的曠野去。到了那裡，你要膏哈薛作亞蘭王， \end{tabularx} \\ \\ \relax
19:16 & \begin{tabularx}{0.7\textwidth}{X} 又膏寧示的孫子耶戶作以色列王，並膏亞伯‧米何拉人沙法的兒子以利沙作先知接續你。 \end{tabularx} \\ \\ \relax
19:17 & \begin{tabularx}{0.7\textwidth}{X} 將來逃過哈薛之刀的，必被耶戶所殺；逃過耶戶之刀的，必被以利沙所殺。 \end{tabularx} \\ \\ \relax
19:18 & \begin{tabularx}{0.7\textwidth}{X} 但我在以色列中留下七千人，是未曾向巴力屈膝，未曾親吻巴力的。」 \end{tabularx} \\ \\
[1ex]
\hline
\hline
\end{longtable}
在大黑暗大背道的時期,神像快被人推翻；人不要神就以為沒有神.這時神要為自己說話,天地是屬祂的,他一直在說話,雖然無言無語,也無聲音可聽,但祂的聲音卻傳遍天下,使整個世界滿了見證.在這時代,人另想一個道理,另立一位神,要把上帝推翻,但上帝雖有慈愛,卻不會包庇罪惡,容忍錯誤,正如耶穌到世上來時,約翰為他作見證,證明祂是道成肉身,住在我們中間,充充滿滿的有恩典有真理,恩典是愛的流露,賜與,但其中也有真理所充滿.耶穌行在神的真理,律法,原則之上,人要推翻神,但神不會錯誤登寶座作王,耶穌來是為神的恩典真理作見證.
\newline
\newline
那時代神興起的見證人,表明祂並沒有退位.以利亞在當時工作,有幾分像士師時代的參孫一樣,是在黑暗時代裡獨行的人,沒有其他同工,單獨一人作戰,直至末後神才告訴以利亞,在他以外還隱藏了七千人為神作見證.
\newline
\newline
一個為神作見證的人,必是在神面前侍立,並且事奉.若不是在神前聽祂的說話,必不能為祂作見證；單靠自己的理智來分析很容易會錯誤的,人當靠主的道,才能認識神,唯有肯侍立在神面前,才能與主親近,知道神的計劃和目的.耶穌曾對門徒如此說,他設立十二個人,要他們常和自己同在,也要差他們去傳道,這次序:先是與他親近,後是出去傳道,一個要傳講耶穌救恩的人,必定要先與主同在,否則怎樣傳講呢?很容易變成空洞的言語,言中無物.一個基督徒需要真知道神,心中的眼睛又要照明,知道神在蒙召的人身上的指望是何等浩大.以利亞真知道神的情形後,才能真正事奉神,為神作見證.
\newline
\newline
一個為主作見證的人,他必定是一個與主十分接近的人,有了這兩件事後,我們還要追求第三件事,正如以利亞兩次題到的:「我為耶和華萬軍之神大發熱心.」若沒有熱心是不可能為主作見證的.以利亞心裡火熱才能為主作工.今天的教會缺乏了心靈上的熱情,以致顯出沒「能」,冷冷淡淡的樣子,又沒有聖靈的作用,不但如此,反而又為自己辯論,批評那些有能力的人,只不過由於情感作用所驅使,他們卻不知道以利亞為主作見證,是為主大發熱心,這種熱心是正確的.
\newline
\newline
不少人起來,把基督教情感的部份挪去,他們以為信仰絕對是理智的,但我們試從基督教基本要道的一節經文裡看看,約翰福音三章十六節告訴我們,「神愛世人」,「愛」是情感作用的一個表現,理智在旁邊只作輔導作用,使我們的愛不會有錯誤,基督教基本的道理是關乎情感的,又如聖靈降臨,教會成立那天,聖靈進入人心,使他們為罪非常難過扎心,他們恆心遵守主的道,又為主作見證,信主的人數天天加增,這班初期教會的信徒也是帶有情感.耶穌去世之前告訴門徒,將來有保惠師聖靈來,他來了使人為罪自己責備自己,這責備之中也是感情上的作用.
\newline
\newline
我們信主是全心全意的尋求,在情感,意志,心靈上,以最深的心來信.愛是屬情感的事,理智便不是愛了,我們要愛神,愛人,不能只有理智而沒有情感,否則只屬理論方面,並不是實際的.
\newline
\newline
一個為主作見證的人,心裡需要火熱,耶穌責備老底嘉教會,他們對神的心是不冷也是不熱的,主要我們有一顆火熱的心,我們應常常省察自己是否侍立在神前?是否真正事奉神?神在我們心中的地位又如何?
\newline
\newline
以利亞有大熱心,他的熱心是受到一件事影響的,就是他看見以色列人背棄神的約,毀壞了神的祭壇,以利亞的熱心除了與神親近外,還由他看見這些悖逆的光景,使他心裡痛苦非常,因而火熱起來,不顧自己的孤單,及情況的惡劣,走到亞哈王面前宣告刑罰,指出以色列全國背棄主的約,違背主的律法,主的道,他們爭取敬拜新的神,新的道理,新的儀式,他們可能以為摩西時代的道理太陳舊不合時代的需要,他們所注重的是新與舊,而不理真或假.他們不知道聖經的話都是神藉著啟示曉論給世人的,出於神的啟示不會有改變,我們只要持守主的道,因為聖經的話並不是人類思想中的產物,也不因時間所限制而有所改變.以利亞見到神的選民背棄神的約,心裡便痛苦,覺得十分難堪,他有神的心,為著神置一切於不顧.我們在大黑暗的時代中,能否為神作大事?我們看見這些背道的事是否心裡也是傷痛,為他們難過?
\newline
\newline
列王記上十八章,我們發現有另外一個神的見證人出現,這人就是亞哈的家宰俄巴底,他在亞哈的宮中作事,神用這個人在亞哈的家裡成為見證人,可見神的工作是十分希奇,又是十分得勝的.列王記給我們看到當人離棄神,神就要在其中顯出祂特殊的工作.俄巴底不是神人,也不是先知,但他卻安靜專心地在他工作的崗位上事奉神,為神作見證,他在生活上表明是一個不怕人,惟獨敬畏耶和華的人,他沒有拜他家主所設立的偶像,也沒有在巴力面前屈膝,他知道自己不可偏離主道,因為神有公義,可畏,當約瑟的主母,引誘約瑟作壞事時,約瑟表明他敬畏神,萬不可作大惡得罪神.
\newline
\newline
俄巴底在亞哈的宮中,他十分敬畏神,他遵行神的旨意,王后要殺神的先知,他便把一百個神的先知隱藏起來,供給他們餅和水,這是神另外的一種見證人,不用聲音乃用行動,他特別要在人面前見證神.
\newline
\newline
列王記的歷史裡表明先知在神面前,其中要學習一個重要的功課就是,神差遣你作甚麼便要作甚麼,差遣你傳達甚麼話就說甚麼話,學習完全敬畏神,尊主的權,順主的命,傳主的話,有忠心的表現.
\newline
\newline
末世大背道之時,神需要這樣的人為祂作見證,求主引導我們的心使我們能敬畏神,為神大發熱心,忠心事奉祂.
\newline
\newline


\newpage

\section{第41屆~港九培靈會~胡恩德~第十四講}
\label{sec:1477}
\textbf{向人施恩的神}
\newline
\newline
連結: \href{https://www.hkbibleconference.org/session-message/view/1477}{\texttt{https://www.hkbibleconference.org/session-message/view/1477}}
\newline
\newline
\hyperref[sec:1490]{< < < PREV SERMON < < <}
~
\hyperref[sec:toc]{[返目錄]}
~
\hyperref[sec:1419]{> > > NEXT SERMON > > >}
\newline
\newline
列王記上 20:1-20:6
\newline
\begin{longtable}{cl}
\hline
\hline
章節 & 經文 (和合本修訂版)\\
\hline
20:1 & \begin{tabularx}{0.7\textwidth}{X} 亞蘭王便‧哈達召集他的全軍，率領三十二個王，帶著馬和戰車，上來圍困撒瑪利亞，要攻打它。 \end{tabularx} \\ \\ \relax
20:2 & \begin{tabularx}{0.7\textwidth}{X} 他派使者進城到以色列王亞哈那裡，對他說：「便‧哈達如此說： \end{tabularx} \\ \\ \relax
20:3 & \begin{tabularx}{0.7\textwidth}{X} 『你的金銀都要歸我，你妻妾兒女中最美的也要歸我。』」 \end{tabularx} \\ \\ \relax
20:4 & \begin{tabularx}{0.7\textwidth}{X} 以色列王回答說：「我主我王啊，就照著你的話，我和我所有的都歸你。」 \end{tabularx} \\ \\ \relax
20:5 & \begin{tabularx}{0.7\textwidth}{X} 使者又來說：「便‧哈達如此說：『我已派人到你那裡，要你把你的金銀、妻妾、兒女都歸我。』 \end{tabularx} \\ \\ \relax
20:6 & \begin{tabularx}{0.7\textwidth}{X} 但明日約在這時候，我還要派臣僕到你那裡，搜查你的家和你僕人的家，你眼中一切所喜愛的都由他們的手拿走。」 \end{tabularx} \\ \\
[1ex]
\hline
\hline
\end{longtable}
十個支派所組成的北國,他們的思想中,雖仍有神的名字,但自從耶羅波安起,便一直離開獨一的真神,敬拜耶羅波安所鑄的金牛.這罪惡不能離開他們的國,以致他們得到惡結果,國中產生混亂,失去權柄,屬天的福氣.亞哈是一個惡王,他的王后是外邦女子,她把外邦巴力假神帶進來,使全國都拜巴力,把國陷在罪惡之中.
\newline
\newline
以利亞在天旱三年零六個月後出來,見證神是天地間獨一的真神,又在他們面前殺了巴力的先知,以利亞滿心以為在他們面前行過如此大的神蹟,他們必轉回歸向神,怎料竟沒有人悔改,心意挽回,把巴力從他們國中除去.亞哈有勇氣行惡,沒有勇氣行善,神在這章聖經的歷史裡,表明祂要用特別的憐憫慈愛給他,這也是神最後的一步.神用公義,威嚴,責備,人反而硬心不怕,因為很多時候人心表現出撒旦的性格,撒旦的名義就是無論大小事情,都是敵對神.我們從前也是敵對神,現今神用愛心的力量感化我們,這力量比責罰的力量更大,本來這是亞哈不配受的恩,但神要施恩給他.
\newline
\newline
亞蘭國的軍隊攻打猶大以色列國,亞哈的京城撒瑪利亞被圍困,他們只得七千個士兵,而亞蘭的軍隊卻有十二萬七千人,所以他們不敢出去抗戰,只能閉關自守,在這光景中他們心靈必定十分不安,沒有辦法可想.這時神的恩典來到,藉著先知,向他們傳恩典的話,耶和華如此說,今天必要把這十二萬七千人交在他們手中.神的恩典是偉大的,亞哈聽到這些話似乎沒有甚麼信心,最後他還是聽從神的話,親自帶領二百三十名跟隨省長的少年人及七千民眾出去爭戰.
\newline
\newline
雖然亞蘭國的軍兵眾多,然而以色列人因有神的幫助便得勝,因為得勝乃在乎耶和華,不在自己的武力,智慧,乃在乎執掌權柄的主宰.
\newline
\newline
現今這個時候,我們也可以看見神的恩典活潑地工作,即如發生不久的以色列六日戰爭,在其中顯然有神蹟出現,有一個在羅省報館專門負責採訪工作的基督徒,他發表自己對以色列六日戰爭的看法,他指出有七件事證明這絕對是出於神蹟,在此處不能多題,其中有一事是戰略上的神蹟,以色列軍方的飛機,每天來回七百多次,不用修理,這些飛機能發揮百分之九十多的效力,這件事實在是令人難以致信的.
\newline
\newline
一個如此得罪神的亞哈,竟然在神面前蒙受如此大勝利,在整場戰爭中顯明了極大的神蹟,恩惠,憐憫.神使他們勝而又勝,至終獲大勝利.當戰爭開始時,亞蘭的軍隊遍滿了地面,以色列人的軍隊如軟弱的羊羔,但神卻使軟弱的成為剛強,少數勝了多數,可見神的恩典夠用,現今主將祂的恩典給我們,不知多麼深厚,我們本是不配得的,因我們如同以弗所書第二章所形容的,我們是死在罪惡過犯中,生活行為不是服從神,只是隨從現今的潮流,心中所喜好的去行,本來為可怒之子,我們就要在神前惹祂極大的烈怒,這樣我們不配受主的恩,只能接受審判.
\newline
\newline
神把白白的恩典加給我們,即如空氣,陽光,雨露.這都是神的恩典,而且還把救恩加給我們,這完全是出自恩典而不是行為,我們得恩典是白白的,但其中有一條件就是叫我們認識耶和華,是天地宇宙的管理者,一切都是屬祂的.
\newline
\newline
我們得恩典以後,是否仍然帶著自己的老我,舊脾氣,原則,觀念?我們若認真在主的光中察看自己時,我們便知道自己實在的情況,有事情臨到時便求主,沒有事時便自尊為大,我們在教會中多少時候是妄自行事,不承認神是活在教會之中,沒有迎合神的心,不是討神的喜悅,神向人施恩的目的,是使人知道耶和華就是獨一的真神.昔日神一次兩次把恩典給亞哈,今天神也樂意向我們施恩,聖經告訴我們:「我們若認自己的罪,神是信實的,公義的,必要赦免我們的罪,洗淨我們一切的不義.」「祂兒子的血為我們作了挽回祭」.「我們有一位中保就是那義者耶穌.」「耶和華是我的牧者,......祂使我的靈魂甦.」神有豐富的恩典,祂定規要救我們,祂的恩浩大,無論我們如何敗壞,只要我們自己的罪,祂有赦免之恩,使我們認識耶和華是神.
\newline
\newline
亞哈雖受神多恩,但他好像是一個木頭人,神兩次向他施恩仍然不知自己的罪,後來因他與神的關係疏遠不曉得神的心意,不知道亞蘭國王便哈達是神命定要殺的,反稱這人為他的兄弟,且與他立約又放他去了.所以神差遣先知對亞哈說:「耶和華如此說,因你將我定要滅絕的人放去,你的命就必代替他的命,你的民也必代替他的民.」亞哈雖承受神如此大恩惠,然而他不能達神施恩的目的,也沒有迎合神的心意.
\newline
\newline
在末後的日子,亞哈還做了兩件極愚昧的事.他因為貪愛拿伯的葡萄園,要拿伯把園賣給他,但拿伯因為敬畏耶和華,他不敢以先人的產業出賣.拿伯提到敬畏神的事,亞哈應當醒悟,但他只追求自己的慾望,悶悶不樂,最後王后耶洗別便用計謀把拿伯害死,佔據了他的葡萄園,神便差遣先知以利亞見亞哈,且宣佈他的結局:「凡屬亞哈的人,死在城中的,必被狗吃,死在田野的,必被空中的鳥吃.」亞哈聽見,心中非常駭怕,便撕裂衣服,神因他自卑便再向他施恩,告訴他當他還在世時不將此災降下直至他死後,可見神有無限豐富的恩典,只要有機會祂便樂意施恩.
\newline
\newline
亞哈穿麻衣撕裂衣服,並不是真正悔改,他對於自己基本的罪,如拜偶像卻沒有放棄,他只是懼怕神的刑罰而已,神用愛作挽回祭是祂最後的一個方法,若不接受刑罰便要來到.在第二十二章裡我們看見亞哈死了,這是神定意要殺他的,猶大王約沙法與亞哈預備上去攻打亞蘭國,希望得回基列的末拉,而神的先知米該雅已預言他不能上去,因為此去必然失敗,但亞哈不聽,改裝上陣,結果陣亡.
\newline
\newline
神用愛心教導我們,若不接受便有刑罰來到,教會時代是蒙恩的時代,空前表現神的恩典,我們能得到天上屬靈的福氣及產業,這實在是蒙大恩的時代,但也是最危險的時代,因為使人蒙恩是神最後的一步,人若不接受便要滅亡,不接受這恩典便有大審判臨到,我們個人如此,神向那些跌倒,不接受主恩的人有嚴厲的態度.
\newline
\newline
主要快來,世界要過去,我們若再貪戀世俗物質是最愚拙的事,應當一心事奉主,順祂的命而行,因著將來永遠無比的榮耀而努力,為主作見證.
\newline
\newline


\newpage

\section{第42屆~港九培靈會~楊濬哲~序}
\label{sec:1419}
\textbf{序}
\newline
\newline
連結: \href{https://www.hkbibleconference.org/session-message/view/1419}{\texttt{https://www.hkbibleconference.org/session-message/view/1419}}
\newline
\newline
\hyperref[sec:1477]{< < < PREV SERMON < < <}
~
\hyperref[sec:toc]{[返目錄]}
~
\hyperref[sec:1420]{> > > NEXT SERMON > > >}
\newline
\newline
一九七○年八月,港九培靈研經會,舉行第四十二屆大會,適兄弟應赴東南雅各地領會；未能與會,深感遺憾!幸得委辦工作人員分勞,使大會能如期舉行,圓滿閉會.
\newline
\newline
返港後,得悉大會經過情況,極為欣慰!問謹題一二,以誌主恩.
\newline
\newline
本屆講員:滕近輝牧師,黃聿侯先生,王永信牧師,他們都是主所重用的僕人,曾在各地領會,復興教會.
\newline
\newline
本屆會期,港方會場適逢鄰近之電話公司大廈,趕工打樁,嘈雜之聲,不絕於耳；使講者聽者,均感吃力.後因天雨,場地不能開工,環境就寧靜下來了.赴會的,有人說我們寧願下雨不方便,勝於吵耳聽不見.
\newline
\newline
黃聿侯先生父親病危,母親來電摧他回家.他經過懇切禱告,確知主旨,他就仍然留下來.及後再收電訊,噩耗傳來,得悉父親離世了,他仍帶著沉痛之心,繼續領會,使會眾深為感動!本會委辦會,曾拍電至星洲,代表培靈會會眾,向黃先生親屬致唁.其時兄弟適在星洲領會,亦親至黃家慰問.
\newline
\newline
香港會期最後兩天,適逢颶風襲港,豪雨淹道；講員,工作人員,會眾,仍不畏阻撓!冒雨涉水赴會,愛主之心不因環境受阻,真令人感動!
\newline
\newline
在信息方面,滕近輝牧師講解啟示錄中的「七」字；黃聿侯先生講基督徒的信仰與生活；王永牧師講解信徒與教會的復興.三位講員所講的,都是針對教會與信徒現實的需要.今已將紀錄刊印成冊,廉價出售,使大家同分享主恩.
\newline
\newline
在四十二屆大會閉會時,委辦使徒安活醫生夫婦,舉家遷往美國.他倆對培靈會工作,供獻至大!尤其是在招待方面,願主記念.
\newline
\newline
書記何道彰牧師,三十年來,忠心於培靈直至一九七○年十一月二十日,蒙主寵召,息其勞苦,實在令人感懷與景仰!
\newline
\newline
多年來蒙港九循道公會,借用堂址,舉行大會.在人事上,得到他們很大的幫助,謹綴數言,以誌不忘.
\newline
\newline
願神恩待培靈會工作,復興教會,直到主再來!阿們.
\newline
\newline


\newpage

\section{第42屆~港九培靈會~滕近輝~第一講}
\label{sec:1420}
\textbf{七個呼聲}
\newline
\newline
連結: \href{https://www.hkbibleconference.org/session-message/view/1420}{\texttt{https://www.hkbibleconference.org/session-message/view/1420}}
\newline
\newline
\hyperref[sec:1419]{< < < PREV SERMON < < <}
~
\hyperref[sec:toc]{[返目錄]}
~
\hyperref[sec:1422]{> > > NEXT SERMON > > >}
\newline
\newline
啟示錄是一本末世的書,內容異常豐富,可以從許多不同的角度來查攷.昔日宋尚節博士在初得救時曾用四十種不同的角度來查攷聖經,從聖經中得到無限豐富.
\newline
\newline
今天我要講的是七個呼聲,這都是末世的呼聲.此時是末世,我們實在需要這七個呼聲!
\newline
\newline
這七個呼聲,都用「看哪!」來表示.「看哪!」明顯的有兩個重要的意義!其一是叫人特別注意；其二是用靈眼看,用信心的眼睛看,看出屬靈的異象.箴言廿九章18節「沒有異象,民就放肆.」(英文聖經譯作「沒有異象,民就滅亡.」)「放肆」意即我行我素,隨意放縱.但若有異象,心中有所領受,便放棄自己的道路,按著神的異象行事.
\newline
\newline
啟示錄中十次用「看哪」,但是其中數次意義相同,所以可以歸納為七種呼聲.
\newline
\newline
第一個呼聲──主來呼聲
\newline
\newline
「看哪!他駕雲降臨,眾目要看見他,地上萬族都要因他哀哭,這話是真實的,阿們!」(一7).另廿二7及廿二12.
\newline
\newline
當主耶穌升天時,有天使對門徒說:「加利利人哪,你們為甚麼站著望天呢!這離開你們被接升天的耶穌,你們見他怎樣往天上去,他還要怎樣來.」(徒一 11).門徒將此話藏在心裡.在新約廿七卷聖經中,只有四卷沒有提及主再來的事:約翰貳書,參書,腓利門書,因篇幅太少故未提及；加拉太書因為專講「稱義」真理,故也沒有提及.還有兩卷書間接提到,這就是以弗所書及羅馬書,此二書內均說到身體得贖,而身體得贖就是主再來時所成全的事.
\newline
\newline
「眾目要看見他.」新約有些地方說主再來時像賊一樣,當然不是眾目要看見他.為何有上述兩種不同的說法?原來主再來分兩步驟；
\newline
\newline
(1)提取──即將信徒接去,在空中與主相遇,是隱藏的,世人看不見.
\newline
\newline
(2)主與被提的信徒降臨在地上建立國度──這次降臨,眾目都要看見.
\newline
\newline
正如舊約關於彌賽亞的預言是綜合性的,一方面說祂是受苦的僕人,另一方面又說祂是大有榮耀的君王.這並不是互相矛盾,乃是預言主兩次來世:一次是卑微的,一次是榮耀的.正如從遠處看兩個山峰,像是連在一起,但近處看,兩峰距離卻是很大.舊約先知是從遠處看；新約乃是從近處看.
\newline
\newline
「連刺他的人也要看見他.」使人想起撒迦利亞書十二章十節的話:「他們必仰望我,就是他們所扎的.」這個「我」字,就是神自稱.誰能扎耶和華?誰能扎天地的主宰?主耶穌是道成肉身,扎主耶穌,就是扎耶和華.猶太人一天不接受基督是彌賽亞,就一天不明白此節之真理.我們不只看見基督在十字架上,也看見神在基督裡面在十字架上為我們而死.
\newline
\newline
「他駕雲降臨.」雲彩象徵神的榮耀,這是說主在榮耀中降臨.
\newline
\newline
「萬族都要因他哀哭.」主在榮耀中降臨,這是好的,為何萬族要為他哀哭呢?這是因為審判的緣故.主第一次來,是落在人手中,人任意待他；主第二次來,是審判者的身份.你要細心思想,在這哀哭的人群中,你在其中嗎?抑或與主一同從榮耀中降臨?
\newline
\newline
第二個呼聲──直接審判的呼聲
\newline
\newline
「看哪!我要叫她病臥在床,那些與她行淫的人,若不悔改所行的,我也要叫他們同受大患難.」(啟二22).
\newline
\newline
這裡叫我們注意一件事,就是在推雅推喇教會裡面,耶洗別所造成的問題.耶洗別這名字,很明顯的不是真名,乃是象徵的名稱,正如上文巴蘭與安提帕,二名之用法一樣.巴蘭是舊約裡面的一個人,安提帕在原文的意思,是「反對一切」──反對一切與神旨意相對的事物.在推雅推喇教會裡面,有一個領袖所做的事情與舊約的耶洗別所做的相同.主在這裡要祂的教會注意一件事,這也是今日我們不能不注意的一件事,這是與信仰有關的一件事.舊約的耶洗別帶領神的百姓離開正確信仰的道路,去敬拜巴力.神就興起先知以利亞來,帶領那些向神忠心的百姓站立得穩.現在推雅推喇教會裡面,發生了同樣的情形,這個領袖所作所為的與耶洗別一樣,帶領神的百姓離開了信仰正確的路徑.請注意,在主致七教會的書信裡面,有七次提到關乎不正確信仰的警告,我們不能不注意這件事.當日的教會需要這警告,今日的教會更加需要；如果這警告是發自他人,我們倒可以不注意,但這是主親自對祂教會所說的話,我們就不能不注意了.在這裡另外有一件事值得我們注意的:在這裡把信仰上不正確的路線與淫亂連在一起講.以當日的歷史背景來說,在推雅推喇城裡的神廟內,有各種宗教的節期,在節期中經常進行淫亂的事.這就是說,不正確仰的路線常造成道德上的鬆弛,聖潔的標準跟著降低了.凡是離開聖經信仰的教會,在這聖潔的道理上,一定放鬆.還有一件事,是值我們注意的,就是主繼續警告說:祂給耶別悔改的機會,但是她不肯悔改,所以主使她病臥在床.她的罪是在床上的,受神審判的地方也是在床上.人種的是甚麼,收的也是甚麼.公義的神,有公義的審判.在舊約裡有許多這樣的例證:雅各的事,是多麼的明顯,他欺騙了他的父親,他也受到舅父的欺騙,和他兒子的欺騙；連在他婚姻的事上,也受到欺騙.神是輕慢不得的!
\newline
\newline
在這經文裡面,我們看見有三件事:(1)主警告我們不可離開信仰的正確的路線,(2)信仰上不正確的路線,一定忽視成聖的道理,(3)人種的是甚麼,收的也是甚麼.這呼聲不但值得我們的注意,更值得我們大聲疾呼.
\newline
\newline
第三個呼聲──傳福音的呼聲
\newline
\newline
「看哪!我在你面前,給你一個敞開的門,是無人能關的.」(啟三8).這是非常重要的一個呼聲,這是一個鼓勵的呼聲,叫我們想起了我們的責任.主為我們開了一個傳福音的門,我們的主開了門,沒有人能關；主關了門就沒有人能開.祂拿著大衛的鑰匙,天上地下的權柄,都在祂手中.今天在未世的時候,佈道的工作來到最後一個階段,而最後一個階段總是最艱鉅的.有些時是難見到效果的,因為世上的罪惡增加,人的物質觀念加強.黎明之前,是最黑暗的一個階段,所以在這階段裡傳福音,一定有許多許多的困難!在這一段的工作當中,需要更大的屬靈力量,方能把工作完成.當我們面對最後一里路,工作最大困難的時候,我們有何反應?我們是因為困難而後退呢,還是因更大的困難而更加追求呢?我們面對最後一段的工作,發現我們原來的能力不夠用了,原來的方法,好像沒有果效,怎辦?是停止呢?後退呢?灰心呢?抑因有更大的需要而向主祈求呢?讓我們今天聽主這個呼聲:「看哪!我給你一個敞開的門!」在今天最艱苦的階段中,鑰匙仍然在主的手裡；祂開的門,仍然是無人能關的.我們聽到這個聲音,我們的心就得安慰.弟兄姊妹們!我說句老實話,如果不是這位莊稼的主呼召我們,我們早已不想在這禾場上工作了!如果不是這位莊稼的主管理禾場,今天我們都做不下去了!我們本身的力量很少,我們能夠做的很有限.但我們力量的來源,我們的盼望,我們的鼓勵,全在乎主,祂手中有大衛的鑰匙,祂開了無人能關,關了無人能開.這是我們要用信心接受的一個信息,如果我們心裡沒有這個保證,我們就甚麼也不能做!最後的一里路,黎明前最黑暗的一個階段,我們把眼睛完全注定在主身上.這樣我們才能夠不灰心,不退後,不喪志!我們可以藉教會歷史證明出來,主開的門,實在無人能關.在教會歷史裡面,有多少多少的時候,按人來看,這門是銅牆鐵壁,誰都不能開的,但是主打開了緊閉之門,在福音傳入愛爾蘭之前,有一位主的僕人,叫做巴提客 St. Patrick,他在一個小山上禱告了四十天,求神開傳道的門,讓福音進入愛爾蘭,後來門就開了.今天我們仍然要接受這寶貴的呼聲.
\newline
\newline
「主說:非拉鐵非教會只有一點力量」,這句話到底是責備呢?還是稱讚呢?他們的教會按人看來,力量不大,但是他們以信服的心,盡上了這一點力量,結果主應許她說:「看哪!我給你一個敞開的門!」因為他們盡了他們的力量,門就為他們開了!感謝主,一點的力量,加上主,就是大大的資源了.一點的力量加上主,就等於成功!一點力量加上主,就等於敞開的門.撒勒法寡婦的瓶裡面有一點麵和油,一次就會吃完了,但加上了神,就可以足夠一千二百六十天之用.非拉鐵非教會的祕訣,就是她那一點力量肯拿出來用.今天的香港每一個教會,都有自己的一點點的力量,讓我們把這一點點的力量加上主,主就要在香港,在東南亞,為我們打開一個敞開的門!
\newline
\newline
第四個呼聲──由冷淡中復興的呼聲
\newline
\newline
「看哪!我站在門外叩門,若有聽見我的聲音就開門的；我就要進到他那裡去,我與他,他與我,一同坐席.」(啟三20).主在老底嘉教會的門外,發出這呼聲.為甚麼主在門外?有三件事把主推到門外去.這三事的力量很大,主在門外針對這三件事,發出呼聲:
\newline
\newline
一,注重物質,他們以為自己發了財,甚麼都不缺.他們讓物質代替了主的地位.在物質豐富的時候,就不需要主,把主擠到門外了!他們看重物質過於主.
\newline
\newline
二,不冷不熱.不冷不熱充滿了教會,對主半心半意.一個向主一心一意的人,斷不玉於不冷不熱；半心半意的結果,一定是不冷不熱.主在這樣的教會裡,就無所作為!祂在教會裡面沒有地位,就等於把祂趕到外面去了.
\newline
\newline
三,自滿自足.他們認為自己甚麼都有,他們沒有需要感.一個教會沒有需要感,就不會要主.我們之所以需要主,正因為祂能把我所需要的給我們.我們追求,正因為我們有需要.自滿的教會,主在裡面就沒有地位了!主就被趕到外面去.
\newline
\newline
主針對這三種情形說話了三句話:
\newline
\newline
(1)祂是在「萬有之上為元首的」(三15),這是祂原來的地位,祂在萬有之上.被造的一切物質,都在主之下；主的地位,應該在一切一切之上.這是祂應該有的地位,祂向老底嘉教會介紹祂自己.
\newline
\newline
祂是誰呢?祂乃是在萬物之上為元首的.祂超過物質的一切一切,這是老底嘉教會所不認識的,不了解的,不領會的,主就要他們如此認識祂自己.弟兄姊妹!今天我們是否需要有這個認識?我的主耶穌是誰?祂是在萬有之上,萬物之上為元首的.元首是甚麼?是第一位,首位,我的主是居首位的主,祂不應在第二位,不應該有任何的事物,在祂之上.今天我們的教會,是否需要主這樣自我介紹?在這物質主義的時代裡面,我們要認識主,仍然要讓祂居首位.如果主在我們的一生中,在我們的家庭中,在我們的事業中居首位,我們這班人就能成全神很多很多的事工了.主就能藉著我們來發揮祂的力量.
\newline
\newline
(2)「誠信真實見證的,稱為阿們的.」阿們是甚麼意思?香港教會習慣用「誠心所願」這句話來結束祈禱,這是很有意思的.我並非要提倡「誠心所願」,我是提倡用「阿們」的,因為全世界的基督徒同在一起祈禱的時候,只有廣東的基督徒說「誠心所願」.但我覺得「誠心所願」四個字是非常寶貴,正是「阿們」的意義.誠誠實實,全心全意,這是主的名字,主是全心全意的向著我們,這是祂一致的表現.祂是「誠信真實見證者」,祂為父的慈愛公義作見證,祂怎樣作見證?祂用生命作見證!祂見證父,直到死在十字架上,這樣的一顆心,正與不冷不熱完全相反.不冷不熱的心,不願上十字架.一個不冷不熱的心,永遠不能真說阿們.我們把老底嘉的不冷不熱,與信真實的見證者互相比較,馬上就能體會這個信息.我真是不冷不熱的基督徒!我真是不冷不熱的傳道人!我們在主的旁邊,真是與祂不相配!我們要接受這個信息,這個呼聲!
\newline
\newline
第五個呼聲──勝利的呼聲
\newline
\newline
「看哪!猶大支派中的獅子,大衛的根,祂已得勝,能以展開那書卷,揭開那七印.」(啟五4)這個呼聲真是寶貴.「獅子」是強壯的,有能力的,是獸中之王.當約翰聽見這個呼聲的時候,他就週圍觀望,他看見甚麼呢?他是否看見獅子?不是!他看見羔羊,而且是被殺的羔羊.希奇之希奇!猶大的獅子就是羔羊!被殺的羔羊就是猶大支派的獅子!祂的能力正在於他的被殺.在神奇妙的作為之中,被殺的羔羊變成了得勝的猶大支派的獅子.感謝神!基督的能力正在祂的軟弱裡.林後十三4)基督軟弱,我們也是這樣同祂軟弱.主的能力是在祂軟弱裡面,祂因順服父神而軟弱.祂在客西馬尼園的時候,曾經說:難道我不能求父為我差遣十二營多的天使來幫助我嗎?但跟著說:如果這樣,那些指著我說的預言,怎能成全呢?祂順服父的旨意而軟弱,被人捉拿,被人綑綁,被人處死!但感謝神!祂的能力就從祂的軟弱裡面產生出來.祂伸出祂的手－有釘痕的手,「好像被殺過的」表明祂有被殺的記號,那就是手中的釘痕.神伸出祂釘痕的手,接過書卷.馬上就有讚美發出!廿四位長老的讚美!四活物的讚美!得救者的讚美!天上,地下,地底下,一切受造之物同聲的讚美!甚麼時候發出讚美?就是當羔羊伸出祂釘痕的手接受書卷時發出的.祂能揭開七印,換言之,即神在這書卷裡的計劃,能夠完成.誰使神的計劃成全?被殺的羔羊.感謝主,這有釘痕的手掌管一切,藉著這隻釘痕的手,神的計劃能夠完全全的成全.在第一章我們看見,當約翰看見主的榮耀而俯伏在地的時候,主就用有釘痕的手撫摸他(一17,20),主的有釘痕之手掌握著陰間的鑰匙；這手掌握七星:七教會的領袖.主曾將這隻有釘痕的手,給門徒看,說:願你們平安,隨即差遣他們.請注意:主先給他們看見釘痕的手,然後差遣他們.看見在前,被差遣在後.門徒看見釘痕的手心中就有保證,知道祂是死而復活的主,勝過了死亡,罪惡,撒旦!這是我的主!我知道祂是誰!然後他們就歡喜接受主的差遣.主說:父怎樣差遣我,我也照樣差遣你們,如果你對一件事的成功的可能性有很大的懷疑,你一定不肯作那件事；如果某某人差遣你作一事,而你對於這人缺乏信心,你肯不肯接受他的差遣和委託呢?你當然不願接受.但是如果我們好像門徒一樣,知道這位主是復活的主,勝過撒旦,罪惡,死亡的得勝之主,我們就會說:「主阿!差遣我!我接受你的差遣和你的託付.」
\newline
\newline
請注意:主向他們一連說了兩次「願你們平安」(約廿19).主向門徒顯現的時候,說「願你們平安」；繼而在給他們看過釘痕的手之後,又第二次說:「願你們平安!」.第一次說的平安,可以算是問候,因為他們害怕,關上了門,主就向他們顯現說,願你們平安.第二次的意義更深長更寶貝,他們心中充滿了平安,是因為他們得了保證；第一次平安是主的復活所造成的客觀平安,第二次的平安是由保證而得的主觀的平安.主觀的平安加上了客觀的平安,就變成了平安的平安!這是主所賜的平安.看哪!猶太支派的獅子,祂已經得勝了,祂配展開書卷,完成了神的計劃!
\newline
\newline
第六個呼聲──儆醒的呼聲
\newline
\newline
「看哪!我來像賊一樣,那儆醒看守衣服,免得赤身而行,叫人見他羞恥的,有福了!」(啟十六15).相信大家都知道「衣服」的意義－表明好行為,即信徒的義.我們信主後,得到了主的義,但自己亦應當行義.這是信徒的衣服.我們應當保守這義,免得赤身羞恥.
\newline
\newline
請注意上文十三節的話:「有三個污穢的靈,好像青蛙......出去到普天下眾王那裡.」污穢的靈在末世的時代中,迷惑人作污穢的事,把許多污穢的思想放在人的心中.末世的時代乃是一個放縱情慾的時代,色情泛濫在社會中,在這污穢的時代裡面,我們聽見這個呼聲,就應當儆醒保守聖潔!
\newline
\newline
第七個呼聲－更新與成全的呼聲
\newline
\newline
「看哪!神的帳幕在人間,他要與人同住,他們要作祂的子民；神要親自與他們同在,作他們的神.」(啟廿一3)「看哪,我將一切更新了!」(廿二5)在這裡我們看見神的計劃完成.看哪!神在舊約裡所說的應驗了,在舊約時代,神藉著會幕來表明祂的同在；在新約時代神要藉教會與人同在,更要藉著聖靈與人同在.看哪!在新耶路撒冷城中,以馬內利的意義就完全實現,神的帳幕要在人間,主耶穌降生的時候,乃是以馬內利的開始,但在新耶路撒冷中是以馬內利的成全了!人神之間再無隔膜,再無距離,同住在一起,恢復了人神原來的情況.我們渴想那日子來到,那時我們就心滿意足了!
\newline
\newline


\newpage

\section{第42屆~港九培靈會~滕近輝~第二講}
\label{sec:1422}
\textbf{七福}
\newline
\newline
連結: \href{https://www.hkbibleconference.org/session-message/view/1422}{\texttt{https://www.hkbibleconference.org/session-message/view/1422}}
\newline
\newline
\hyperref[sec:1420]{< < < PREV SERMON < < <}
~
\hyperref[sec:toc]{[返目錄]}
~
\hyperref[sec:1423]{> > > NEXT SERMON > > >}
\newline
\newline
(一)「念這書上預言的,和那些聽見又遵守其中所記載的,都是有福的,因為日期近了.」(啟一3)
\newline
\newline
根據這處經文,我們知道信徒應該研讀啟示錄,這樣就會儆醒預備等候主來,並且見主時不致羞愧,反得獎賞,這實在是大福氣.
\newline
\newline
(二)「我聽見從天上有聲音說,你要寫下,從今以後,在主裡面而死的人有福了!聖靈說,是的,他們息了自己的勞苦,作工的果效也隨著他們.」(啟十四13).
\newline
\newline
這裡有一個很重要的問題,是平常很少提到的.此處的經節我們常在安息禮拜的時候用.但很少注意到「從今以後」四個字的意義.「從今以後」到底指甚麼時候說呢?解經家有幾個解法:
\newline
\newline
(1)這四個字,是連在上文,而非連在下文的.意即「我聽見從天上有聲音說,從今以後你要寫下......」但大多數希臘文的學者,都認為此四字是連在下文的,是指死時,而不是指約翰寫的時候.
\newline
\newline
(2)「從今以後」乃是指大災難中後三年半開始以後.大家都知道七年的大災難,係以後三年半更為厲害的,在此時期開始時即去世,可免去後三年半之大災難.
\newline
\newline
(3)此節是指那些在後三年半中殉道的信徒而言.他們的死不是平常的死,乃是犧牲的死；他們所受的,乃是最難最大的試鍊.故此,他們是特別有福的!我認為後者是正解.
\newline
\newline
(三)「在頭一次復活有分的有福了,聖潔了,第二次的死在他們身上沒有權柄,他們必作上帝和基督的祭司,並要與基督一同作王一千年.」(啟廿6)這一個福和上面的一個福有密切的聯繫,因為前者殉道的聖徒,他們必要復活,與主一同作王.到底是誰復活呢?就是那些為神之道被斬者的靈魂,和那沒有拜過獸的獸像,也沒有在額上和手上受過他印記之人的靈魂.他們乃是後三年半被殺的,他們要與主一同作王.他們忠心得勝,至拜獸像,他們有福了!能夠與主同作王.我們不要單求得救就算數,我們必須求與主一同作王.(提後二12).「我們若能忍耐,也必和祂一同作王.」忍耐是資格,有這資格的人,纔是有福的,因為他們必作王.
\newline
\newline
(四)「天使吩咐我說,你要寫上,凡被請赴羔羊婚筵的有福了!又對我說,神真實的話.」(啟十九9).到底誰是「被請赴羔羊婚筵的」呢?解經家的意見,各有不同.有人認為新婦是指教會,被請者是猶太人,因為教會為主,猶太人為副.有人認為兩者所指是相同的,因為教會是由信主的猶太人與外邦人組成的.教會包括各時代聖徒的整體.教會不可分開,不分猶太人或外邦人；基督沒有兩個身體,凡屬基督的,都在基督裡面,祂的身體是分不開的.在本書廿一章三節,新耶路撒冷也是用新婦來形容,新婦只有一個,並沒有兩個的.
\newline
\newline
如果我們是得救的人,我們就是基督的新婦,我們就要與主合而為一.那麼我們在世上的生活應該如何呢?我們是否以新婦的貞潔的資格迎見新郎呢?新郎最怕新婦不貞,如果一個得救的人去作信徒不應作的事,就會叫基督傷心.我們若能守住貞潔的身份就有福了!
\newline
\newline
(五)「看哪!我來像賊一樣,那儆醒,看守衣服,免得赤身而行,叫人看見他羞恥的,有福了!」(啟十六15)這一節已在七個呼聲之內提到,在此不詳述.
\newline
\newline
(六)「那些洗淨自己衣服的有福了!可得權柄到生命樣那裡,也能從門進城.」(啟廿二14)這衣服是甚麼呢?啟十九8說「就蒙恩得穿光明潔白的細麻衣,這細麻衣就是聖徒所行的義.」原文無「所行」二字,只有「聖徒的義」.這「義」字的,原文是多數的,故譯為所行的義.到底此義字何所指?行為之義或承受基督之義?
\newline
\newline
新婦在婚宴中所穿之義袍,當然是基督之義所製成,這是不容異議的,但是為什麼此處的義字是多數呢?這個問題並不難回答.信徒之義在基本上是承受基督之義,但是在這裡的基礎上,建立自己的義行.這也就是保羅所說「戰兢作成自己得救的工夫」同樣的意義.
\newline
\newline
(七)「我必快來,凡遵守這書上預言的有福了.」(啟廿二7).七福的首末二福都是特別與遵守啟示錄之預言有關.這加倍的注重,更令我們重視主的再來.
\newline
\newline


\newpage

\section{第42屆~港九培靈會~滕近輝~第三講}
\label{sec:1423}
\textbf{基督的七對名字}
\newline
\newline
連結: \href{https://www.hkbibleconference.org/session-message/view/1423}{\texttt{https://www.hkbibleconference.org/session-message/view/1423}}
\newline
\newline
\hyperref[sec:1422]{< < < PREV SERMON < < <}
~
\hyperref[sec:toc]{[返目錄]}
~
\hyperref[sec:1424]{> > > NEXT SERMON > > >}
\newline
\newline
啟示錄中主耶穌的七對名字,都是很寶貴的,讓我們一一查攷.
\newline
\newline
(一)啟廿二13「我是阿拉法,我是俄梅戛；我是首先的,我是末後的；我是初,我是終.」啟示錄第一章和這末一章都說到這名字,這一對名字確是寶貴.主耶穌是最先的,也是最後的,一切都在他手中；沒有一事在他之前,沒有一事在他之後；若有一事在他之前,或在他之後,他就不能掌管.
\newline
\newline
神的名稱為「自有永有的」.聖經文的意義是「我是我所是」.祂的存在並無「他因」,祂是一切的根源.一切原因和結果都在他裡面.這「是」字是現在式,無過去,無將來,是自有永有,時間空間都包括在裡面.
\newline
\newline
三一頌說:「神是萬福之源.」他是一切福樂的源頭,一切福樂都包在他裡面.祂所開始的祂必成全.
\newline
\newline
他是始,他是終.一切都是他掌管,宇宙的歸宿,救恩的完成,都在他掌中.領會這名字的,心中有最寶貴的保證.
\newline
\newline
(二)「那為阿們的,為誠信真實見證的.」(啟三14).啟示錄十九章十一節:「祂稱為誠信真實」.
\newline
\newline
「阿們」是對我們而言,「誠信真實見證者」是對神而言,形成一對.
\newline
\newline
哥林多後書一章廿四節說:「神的應許,不論有多少,在基督都是的,所以藉著他也都是實在的,叫神因我們得榮耀.」小字註明「實在」的原文是「阿們」.
\newline
\newline
「阿們」的意義如下:
\newline
\newline
(1)阿們是確定之意,永不改變,不是而又非,不搖動,所有的應許在基督裡面都是確定的.火論新約,舊約,都成全在基督裡.
\newline
\newline
(2)撒旦常常要否定神的計劃,要阻擋,要破壞.撒旦二字之意,就是敵對者,牠的工作,牠的立場是敵對神；在宇宙間造成敵對勢力,但牠永不能否定神的計劃.
\newline
\newline
神是沒有似是而非的,祂所決定計劃,必成全到永遠.
\newline
\newline
「誠信真實見證者」指祂是神的誠信真實的見證人.祂見證神的愛,見證神的真理,見證神的榮耀與公義.為了作見證祂死在十字架上.
\newline
\newline
(三)「祂的名稱為神之道.」(啟十九13).「我又看新城耶路撒冷由神那裡天而降,豫備好了就如新婦裝飾整齊,等候丈夫.」
\newline
\newline
「神的道」,表明主耶穌與神合一.「太初有道,道與神同在,道就是神.」「丈夫」則表明主耶穌與教會合一,主耶穌是丈夫,我們是新婦.因此,主耶穌與我們的合一,是愛的合一.此二名又形成寶貴的一對.
\newline
\newline
主耶穌與神合一也是愛的合一,因為祂是神懷裡的獨生子.
\newline
\newline
我想到此意真是感謝讚美神,主耶穌是神,也是人；是神人間的中保,我們有此中保,大可以放心.他是神,為人的樣式,最後造成合一,即人在基督裡與神合一.
\newline
\newline
保羅說藉著基督,我們與神和好,這是大合一的第一步.罪惡使我們與神隔絕,耶穌基督卻把這隔絕除去,把人與神間的鴻溝除去,使人與神,神與人大合一,成就了神的計劃.
\newline
\newline
世上沒有一種宗教,把人的地位抬舉得如此高.許多人以為基督教將人壓低,又以為人文主義才是給人崇高的地位.其實聖經說人是照著神的形像造的,而且最後還與神合一.聖經說人裡面有神的形像,人神合一,這地位豈不是最崇高嗎?
\newline
\newline
不過,人在被抬舉之先,必須被否定,被破碎.破碎之後方得建立,這是神的意念高過我們的意念.
\newline
\newline
保羅在羅馬書十一章末段說:「深哉,神豐富的智慧和知識,他的判斷,何其難測.誰知道主的心,誰作過他的謀士呢!」保羅被神奇妙的作為陶醉,忍不住高聲讚美!
\newline
\newline
(四)「在祂衣服和大腿上,有名寫著說,萬王之王,萬主之主.」(啟十九16)
\newline
\newline
「在神創造萬物之上為元首.」(啟三14)
\newline
\newline
「萬王之王,萬主之主」與「元首」是主的第四對名字.
\newline
\newline
「萬王之王,萬主之主」是對人而言,「元首」是對宇宙一切被造之物而言.
\newline
\newline
我們不但要認識主耶穌是先知與祭司,也要認識他「是首生的,在一切被造的以先.」「在凡事上居首位」,「萬王之王」.
\newline
\newline
今天我們是否認他是主,是我們心中的王?
\newline
\newline
(五)「明亮的晨星.」(廿二16)「用鐵杖轄萬國的.」這是主耶穌的第五對名字.
\newline
\newline
「明亮的晨星」表示盼望,而「鐵杖」表示審判.黑夜已深,白畫將近,黎明是晨星出現的時候.黎明前是最黑暗的.今日的政治與道德,都是最黑暗的.看這黑夜已深,心裡沈重,你的盼望在何處?是在主身上嗎?是盼望主快再來嗎?明亮的晨星,將黑暗驅走,將光明帶來.
\newline
\newline
但主「用鐵杖轄管萬國.」對不合神旨意的人,主要用公義審判他.
\newline
\newline
「鐵」是強力的意思,主按真理,按事實來審判,不是按人的看法與標準來審判.
\newline
\newline
甚麼是真理?神救恩的原則和人生的真道,合成真理.主將照出人心裡的隱情,和事實的真相.
\newline
\newline
(六)「猶大支派中的獅子」和「被殺的羔羊」是主第六對名字.基督是得勝者,是獅子.撒旦是紅龍,紅表流血,殘忍,龍表示神祕能力,但猶大支派中的獅子已得勝牠!將牠踏在下.
\newline
\newline
「羔羊」是軟弱的名字,耶穌像羔羊在剪毛的人手下無聲.
\newline
\newline
宇宙是無聲之聲,傳揚神的榮耀；十字是無聲之聲,傳揚神的大愛.基督的能力藏在軟弱裡面.
\newline
\newline
(七)「耶穌基督」(啟一1)耶穌是道成肉身的名字.基督是神格的名字,合成一對.
\newline
\newline
「你要給他起名叫耶穌.」這是人的名字；他是受膏者,是神所設立的基督.他不只是人,也是神.祂第一次來,顯為耶穌；第二次來,顯為基督,成全救恩.
\newline
\newline


\newpage

\section{第42屆~港九培靈會~滕近輝~第四講}
\label{sec:1424}
\textbf{新耶路撒冷的七種完全}
\newline
\newline
連結: \href{https://www.hkbibleconference.org/session-message/view/1424}{\texttt{https://www.hkbibleconference.org/session-message/view/1424}}
\newline
\newline
\hyperref[sec:1423]{< < < PREV SERMON < < <}
~
\hyperref[sec:toc]{[返目錄]}
~
\hyperref[sec:1425]{> > > NEXT SERMON > > >}
\newline
\newline
新耶路撒冷的七種完全,每種有二方面,可說是七對完全.
\newline
\newline
今天,我們未得完全,連救恩也未到完全的地步；但啟示錄末後兩章,給我們看見完全的完全.
\newline
\newline
(一)完全的平安與完全的快樂
\newline
\newline
A,完全的平安:「我又看見一個新天新地,因為先前的天地已經過去了,海也不再有了.」(啟廿一2).本來新天新地與新耶路撒冷有所不同,新天新地的範圍較廣,新耶路撒冷是其中的一部分.但是在屬靈的實質上,二者是相同的.
\newline
\newline
「海也不再有了.」為何特特提到無海?
\newline
\newline
聖經說的海,是指政治變動如海中波浪翻騰.如但以理書七章一節所載,有四個大獸從海中上來,意即四大帝國從地中海的政局變化中產生出來.無海即不會有波浪翻騰,一切平安穩定.今日世界到處不平安,但新耶路撒冷是平安穩定不動搖的國度,我們羨慕這平安.
\newline
\newline
過去兩年香港有不少人往美國,但現在已有多人陸續回來.過去從台灣去美國的與加拿大的留學生,大多數不返回台灣,但最近一年中有四百留學生回到台灣,因為美加遭遇經濟不景氣,謀職甚難.那裡有平安?
\newline
\newline
真平安在新耶路撒冷.撒冷本是平安之意,可是到今天為止,地上的耶路撒冷雖名為平安,卻不平安；因為人不認識主,平安之王不在人的心中.
\newline
\newline
新耶路撒冷才有真平安,因為那時認識耶和華的知識充滿遍地.
\newline
\newline
B,完全的快樂:「神擦去他們的眼淚,不再有死亡,也不再有悲哀,哭號,疼痛,因為以前的事都過去了.」(啟廿一4).先前的悲哀過去了,眼淚被擦乾.今天誰無流淚?這是流淚的世界,連快樂時也流淚,所謂高興到流淚,可見在這世界上連快樂也不能與流淚分開.世上沒有純快樂.快樂時總攙有悲哀的成份在內,即所謂樂極生悲.我們不能不承認,人是生存在痛苦的世界中,唯有神能擦乾我們的眼淚.那時才有完完全全的快樂.
\newline
\newline
(二)完全的愛與完全的合一
\newline
\newline
A,完全的愛:「我又看見聖城新耶路撒冷由上帝那裡從天而降,預備好了,就如新婦妝飾整齊,等候丈夫.」(啟廿一2).新耶路撒冷最顯著特點,是「就如新婦」,表示與基督之間的愛的結合.在新耶路撒冷,人人都有愛主的心,心中若沒有愛,住在新耶路撒冷,是無味的；住在這城內的人,充滿基督的愛,也充滿愛基督的心.
\newline
\newline
B,完全的合一(啟廿一12-21).這段經文有幾個數目,很有意義,表達完全合一.
\newline
\newline
「十二」表示完全,其特點是三乘四,三是神數目,四是受造者的數目,包括人在內所以十二是表示神與人的完全.
\newline
\newline
「一百四十四」.「十二」既是完全的數目,十二乘十二是一百四十四,故此,一百四十四表示完全的完全.在新耶路撒冷內,有神與人的合一,也有人與人的合一.今天人與人間未能完全和諧,人與神之間仍有隔斷之牆.但在新耶路撒冷城內,人與神,人與人都合一.新約時代屬神的人與舊約時代屬神的人之間,也沒有劃分.因此在這段經文中提到舊約十二支派和新約十二使徒.
\newline
\newline
在十七節說:一百四十四是「人的尺寸」,也是「天使的尺寸.」意思就是說,人與天使也完全合一.
\newline
\newline
這是一幅和諧的景象,人與神,人與人,人與天使都合一,這是多麼動人的一幅景象.
\newline
\newline
「長闊高都是四千里」(十六節).長闊高是三度,各度皆四千里,也即是三乘四千,等於一萬二千里,其基本意義仍然是十二,代表完全.長闊是面積,是平面的,表示得救的是此種理想的實現.我們每天所禱告的:「願你的旨意行在地,上如同行在天上.」這時完全實現了.在新耶路撒冷的人,絕對遵行神的旨意.
\newline
\newline
(三)完全的榮耀和完全的聖潔
\newline
\newline
A,完全的榮耀:「城中有上帝的榮耀,城的光輝如同極貴的寶石,好像碧玉,明如水晶.」(啟廿一11).「那城內又不用日月光照,因為神的榮耀光照,又有羔羊城的燈.」(啟廿一23)在25節末句說:「在那裡原沒有黑夜.」這表示完全光明.神榮耀的光充滿城內,黑暗權勢全表粉碎.「明亮如水晶」(廿二1),「不再有黑夜」(廿二5),整個新耶路撒冷,充滿光明,這是完全的榮耀.
\newline
\newline
今天的世界黑暗掌權,連屬主的人,心中也有黑暗.但在新耶路撒冷,充滿神的榮耀光明,是光明掌權.
\newline
\newline
B,完全的聖潔.新耶路冷的榮耀「如同極貴的寶石,好像碧玉」(十一節),其城牆也是碧玉造的(十八節).碧玉代表聖潔,神是聖潔的,神榮耀的光輝正是聖潔的光輝.在啟示錄第四章三節形容神好像碧玊,新耶路撒冷城的第一根基也是碧玉.神是聖潔的,只有聖潔的人才可進城,所以「凡不潔淨的」,總不得進耶城,只有名字寫在羔生冊上的纔得進去(啟廿一25)
\newline
\newline
(四)完全的敬拜與完全的事奉
\newline
\newline
A,完全的敬拜:「主神全能者和羔羊,為城的殿.」(啟廿一22)殿是神人相會之處,如「伯特利」所代表的意義.古時會幕是神人相會之所,今時因主耶穌降世,就不需要會幕.因有主耶穌道成肉身,住在我們中間,在主裡神人便可以相會,所以主說祂就是聖殿(約二19).
\newline
\newline
神自己為城的殿,表示我們的敬拜是完全的.今天的敬拜還未完全,因為有人身在禮堂內,心卻在堂外.但在新耶路撒冷,是心靈和誠實的敬拜.
\newline
\newline
B,完全的事奉:「在城裡有神和羔羊的寶座,他的僕人都要事奉他.」(啟廿二3)今天有事奉的工作,將來在新耶路撒冷也有事奉的工作.有人說,離世是進入更高事奉.那是可享受的事奉,是安息的事奉,完全的事奉.
\newline
\newline
(五)完全的同在與完全的了解
\newline
\newline
A,完全的同在:「我聽見有大聲音從寶座出來說,看哪,神的帳幕在人間,他要與人同住,他們要作他的子民,上帝要親自與他們同在作他們的神.」(啟廿一3).神的帳幕在人間,表示神與人完全同在.主耶穌降生,是「以馬內利」的開始,神藉著祂的兒子住在我們中間.但是,到了這節經文所指之時,「以馬內利」的意義便要完成,就是完全的同在.
\newline
\newline
B,完全的了解:「也要見他的面,他的名字必寫在的額上.」(啟廿二4).今不我們看不見神的面,因為從來沒有人看見神.「我們如今彷彿對著鏡子觀看,模糊不清,到那時就全知道,如同主知道我一樣.」(林前十三12)全知道,表示完全了解.
\newline
\newline
(六)完全的生命與完全的靈果
\newline
\newline
A,完全的生命:天使又指示我在城內街道當中一道生命水的河,明亮如水晶,從神和羔羊的寶座流出來.」(啟廿二1).在新耶路撒冷有生命河,表示完全的生命.今天,生命是在成長中,發展中.到那時生命河水充滿我們,使我們有豐富的,完全的生命.
\newline
\newline
B,完全的靈果:「在河這邊與那邊有生命樹,結十二樣果子,每月都結果子,樹上的葉子乃為醫治萬民.」(啟廿二2).這是屬靈的果子,正是加拉太書所指的果子.每日都有果子,所以是完全的靈果.
\newline
\newline
(七)完全的福分與完全的權柄
\newline
\newline
A,完全的福份:「以後再沒有咒詛.」(啟廿二3)沒有咒詛,只有祝福.今天仍有咒詛的影響在我們身上,身體得贖,咒詛也被完全除去.
\newline
\newline
B,完全的權柄:「他們要作王,直到永永遠遠.」(啟廿二5)
\newline
\newline


\newpage

\section{第42屆~港九培靈會~滕近輝~第五講}
\label{sec:1425}
\textbf{七個對比}
\newline
\newline
連結: \href{https://www.hkbibleconference.org/session-message/view/1425}{\texttt{https://www.hkbibleconference.org/session-message/view/1425}}
\newline
\newline
\hyperref[sec:1424]{< < < PREV SERMON < < <}
~
\hyperref[sec:toc]{[返目錄]}
~
\hyperref[sec:1426]{> > > NEXT SERMON > > >}
\newline
\newline
(一)神的印與獸的印
\newline
\newline
A,神的印:「我又看見另一位天使從日出之地上來,拿著永生神的印,他就向那得著權柄能傷害地和海的四位天使,大聲喊著說:地與海並樹木,你們不可傷害,等我們印了我們神眾僕人的額.」(啟七2-3)受神印的人有十四萬四千,這印是屬神號.教會的名冊上有你的名,但你身上若沒有神的印記,仍不是屬神.
\newline
\newline
但此處所指受的人與三章十二節所指「得勝的......我又要將神的名和我神城的名,並我的新名都寫在他上面.」有所不同.那是「得勝者」所獲得的名號.每個得救的人,都受了神的印記,但是得勝者才能獲得那些名號.
\newline
\newline
在以弗所一章十三至十四節提及受聖靈的印記,這印記是得基業的憑據,這印記也是表示屬神的記號,每一個得救的人都有這印記.
\newline
\newline
在啟示錄內有兩次提到十四萬四千人,即七章和十四章,這兩處的十四萬四千人,所指不同.第七章所記的十四萬四千人是額上有神印記的人,但十四章的十四萬四千人,額上寫著主的名和神的名,所以這是得救與得勝之不同.
\newline
\newline
B,獸的印:「他又叫眾人,無論大小貧富,自主的,為奴的,都在右手上,或是在額上,受一個印記.」(啟十三16)
\newline
\newline
受了獸的印記的人,表示他是屬獸的.在本書第十五章二節提到獸和獸的像並他名字數目,與及第十三章十七節所提獸印,獸名.可見屬獸之人身上也有獸的三個記號.所以受神印記和受獸印記是顯明對比.
\newline
\newline
印記是在額上,表示被人容易看見.
\newline
\newline
(二)兩個婦人的異象
\newline
\newline
A,「天上現出大異象來,有一個婦人,身披日頭,腳踏月亮,頭戴十二星的冠冕.他懷了孕,在生產的艱難中疼痛呼叫.」(啟十二1-2)在舊約記載約瑟的夢中,日月代表以色列的先祖,十二星是代表以色列的十二支派.故此這婦人明顯的代表以色列人.懷孕所生的孩子,其重要的任務,是用鐵杖管轄列國,這孩子就是耶穌.也是指耶穌的身體:教會.
\newline
\newline
生產艱難,表示在準備基督降世的過程中,經過許多艱難與痛苦.
\newline
\newline
B,大淫婦(啟十七1-6).這婦人被稱為大淫婦,代表撒旦的作為和計劃.這大淫婦騎在朱紅色的獸上,她名叫「大巴比倫」(五節)作世上淫婦騎在朱經色的獸上,她名叫「大巴比倫」(五節)作世上淫婦和一切可憎之物的母.
\newline
\newline
在第十七章末節:「你所看見的那女人,就是管轄地上眾王的大城.」故這淫婦與十八章所稱的大巴比倫城是三而一的.包括下列四特點:
\newline
\newline
(1)是淫亂的,包括世上一切情慾的罪惡.
\newline
\newline
(2)與君王有關,君王是他利用的工具,用來統治世界.
\newline
\newline
(3)遍佈有褻瀆的名號.其意義是高舉自己,與神為敵.
\newline
\newline
(4)喝醉了聖徒的血,和為耶穌作見證之人的血(六節),是指他逼迫教會,殺害信徒.這都是撒旦的作為和計劃.這大巴比倫最後的結果是被毀滅.
\newline
\newline
(三)羔羊與龍
\newline
\newline
A,羔羊:在啟示錄中,提到主耶穌時,多稱之為羔羊,表明他是馴良順服,贖罪的.
\newline
\newline
B,龍:是撒旦在啟示錄中主要的表現.在啟示錄中撒旦有七個名字.
\newline
\newline
(1)大龍──代表奧祕的權柄與能力.
\newline
\newline
(2)古蛇──代表其狡猾.
\newline
\newline
(3)魔鬼──代表其兇暴.
\newline
\newline
(4)撒旦──代表與神為敵.
\newline
\newline
(5)迷惑者──引誘人遠離神.
\newline
\newline
(6)亞巴頓──意即沈淪.
\newline
\newline
(7)亞玻倫──意即毀滅.
\newline
\newline
魔鬼的目的,是要人沈淪,毀滅,是要破壞人的信仰,破壞人的幸福.
\newline
\newline
主耶穌的名字是羔羊,魔鬼的名是龍,但羔羊要勝過龍.
\newline
\newline
(四)三位一體的對比
\newline
\newline
A,神的三位一很清楚的記在本書第一章,即神,基督與聖靈.
\newline
\newline
B,「我所看見的獸,形狀像豹,腳像熊的腳,口像獅子的口,那龍將自己的能力,座位和大權柄,都給了他.」(啟十三2)「我又看見另有一個獸從地中上來,有兩角如同羔羊,說話好像龍.他在頭一個獸面前,施行頭一個獸所有的權柄,並且叫他和住在地上的人,拜那死傷醫好的頭一個獸.」(啟十三 11-12)
\newline
\newline
這兩處經文都提到獸的三位一體.
\newline
\newline
龍是撒旦本身,是第一位,獸是第二位.獸又分為二:第一位是政治的領袖,第二位是宗教領袖.第一獸是從海中上來(十三1),海是代表政治的變化翻騰.第二獸是從地中上來,地比較穩定,所以代表宗教界.宗教的領袖為第一獸作見證,意即宗教勢力與政治勢力聯合起來,最後龍便要操縱一切.
\newline
\newline
第三位是污穢的靈,「我又看見三個污穢的靈,好像青蛙,從龍口,獸口並假先知的口中出來.」(啟十六13)
\newline
\newline
由此可知撒旦的三位一體.其目的是摹倣神.撒旦的戰術就是以假亂真.在摩西時代,術士在法老面前,摹倣摩西以杖變作蛇,使人混亂,使人敬畏神的心減少.
\newline
\newline
(五)新耶路撒冷城與大巴比倫城
\newline
\newline
A,新耶路撒冷,是基督的新婦,與基督合而為一.
\newline
\newline
B,大巴比倫是鬼魔的住處和各樣污穢之靈之巢穴,和各樣污穢可憎之雀鳥的巢穴.可以說,與神為敵之一切都包括在內.巴比倫的意思是混亂,此名是從巴別塔而來.當時人要高舉自己,傳揚自己的名,因而建造巴別塔.
\newline
\newline
今日,無論是家庭,政治,社會,道德都在混亂狀態之中,其根源就是如此.甚麼時候人高抬自己敵擋神,就產生混亂.巴比倫王尼布甲尼撒王是其最高象徵.
\newline
\newline
(六)基督與敵基督
\newline
\newline
這是啟示錄中最明顯的對比.前文已加以解釋.
\newline
\newline
(七)二次收割
\newline
\newline
A,在啟示錄十四章14-16所載是人子收割莊稼,是屬主的人被收割,是熟麥子的收割,收割好了,將麥子放在倉裡.
\newline
\newline
B,葡萄的收割(17-20)天使把鐮刀扔在地上,收取了地上的葡萄,丟在神忿怒的大酒醡中,有血從酒醡裡流出來.這是不屬神的人的大收割.
\newline
\newline
麥子收放在倉裡,葡萄收割放在神的酒醡中；所以一個是永生的收割,一個是審判的收割.
\newline
\newline
從這七個對比,可見撒旦的作為與神針鋒相對.
\newline
\newline


\newpage

\section{第42屆~港九培靈會~滕近輝~第六講}
\label{sec:1426}
\textbf{主耶穌的七對異象}
\newline
\newline
連結: \href{https://www.hkbibleconference.org/session-message/view/1426}{\texttt{https://www.hkbibleconference.org/session-message/view/1426}}
\newline
\newline
\hyperref[sec:1425]{< < < PREV SERMON < < <}
~
\hyperref[sec:toc]{[返目錄]}
~
\hyperref[sec:1427]{> > > NEXT SERMON > > >}
\newline
\newline
(一)榮耀的異象
\newline
\newline
甲,啟一12-16.這異象各要點正是針對七教會需要,解釋如左:
\newline
\newline
以弗所教會失去起初的愛心,所以燈臺的被挪去.因此主耶穌說:「我是那右手拿著七星,在七個金燈臺中間行走的.」(二1)這話是從第一章主耶穌榮耀的異像引述的,正針對以弗所教會的需要而發.主在「燈台中行走」鑑察,祂對於教會的情形異常清楚.同時主也供給燈台(教會)的需要,如剪燈蕊,加油等是.以弗所教會的燈台開始發暗,失去起初的愛,失去力量.今天我們是否也如此?
\newline
\newline
士每拿教會受大逼迫,受大患難十日.有解經家以為十日是指十次大逼迫,所以士每拿教會需要「首先的,末後的,死過又活的」(二8)主耶穌,因為祂曾受大苦難,也曾受大逼迫,但死過又活了,能加給他們在苦難中的力量.這正是主榮耀異象的要點之一.(一17,18)
\newline
\newline
別迦摩教會,容讓撒旦居在其中,在真理上有偏差,所以主說:「我是那有兩刃利劍的」(二12)兩刃之劍代表主的道,能救他們脫離異端,脫離背道的信仰.這正是一16的意義.
\newline
\newline
推雅推拉的教會,有淫亂,有耶洗別教導他們行姦淫,所以主耶穌說他是「眼目如火焰,腳像光明銅的神之子.」(二18)眼目如火焰,能洞察一切,這正是一14,15所顯明的.其腳掌所到之處都為光明之處.
\newline
\newline
撒狄教會按名是活的,其實是死的.所以主耶穌說他是那有神的七靈和七星的(三1),叫那按名是活的,其實是死的教會再活過來.這正是一16所說的.
\newline
\newline
非拉鐵非教會有一點力量.卻肯為傳音而用.主耶穌說:「我是那聖潔真實拿著大衛的鑰匙,開了就沒有人能關,關了就沒有人能開.」(三8)這正是一18所說的.
\newline
\newline
老底嘉教會不冷不熱,注重物質多過注重靈性.所以主耶穌說:「我是那為阿們的,為誠信真實見證的,在神創造萬物之上為元首的.」(三15).主在萬物之上,所以信徒要愛主多過愛物質.這正是一17所說的「首先的」之意.
\newline
\newline
從上述,可見主耶穌按教會之需要,引用此榮耀異象中不同之要點.
\newline
\newline
乙,啟十1-3.在此段經文用「另有一位大力的天使」,代表主耶穌.祂「披著雲彩」,雲彩代表神的榮耀.例如雲彩遮蓋摩西所建的會幕,有雲充滿所羅門所建之殿,這表示被神的榮耀所充滿.
\newline
\newline
「虹」是恩典的記號.在第四章3節說神寶座有虹圍著.
\newline
\newline
「臉面像日頭」,也是表示神的榮耀,「如同烈日放光」.「兩腳像火柱」,表示主的光明與聖潔.「右腳踏海,左腳踏地」,表示海陸都是屬主的.
\newline
\newline
所以這位大力的天使,明顯是指主耶穌而言.在聖經中,天使的意思有兩種用法.其一是天使,其二是負有使命者.希伯來書三章一節,稱主耶穌為使者.天使可譯作使者.
\newline
\newline
今天這罪惡的世界,真需要主的同在－如雲彩充滿聖殿；今日是黑暗的時代,需要主如烈日放光照耀我們；今日世界充滿迷惑,需要主如同火柱引導我們.
\newline
\newline
(二)勝利的異象
\newline
\newline
甲,啟五4-7.約翰看見因為沒有配展開,配觀看那書卷的,就大哭.立時有一個呼聲:「不要哭,看哪,猶大支派中的獅子,大衛的根,他已得勝,能以展開那書卷,揭開那七印.」這時約翰舉目觀看,但他所看見的不是獅子,乃是被殺過的羔羊.羔羊原是與獅子相反,但這羔羊正是得勝的獅子.勝是因被殺,他的死是勝利－藉著死敗壞那掌死權的魔鬼.這是羔羊得勝的異象.
\newline
\newline
乙,「我觀看,見天開了,有一匹白馬,騎在馬上的,稱為誠信真實.」(十九11).「他的名稱為神之道.」(13節).「萬王之王,萬主之主.」(16節).這也是得勝的異象,基督是萬王之王,萬主之主,「有利劍從他口中出來,可以擊殺列國,他必用鐵杖轄管他們.」(15節).「他頭上戴許多冠冕.」(12節).
\newline
\newline
這是一對得勝的異象:前者是指主第一次降生,走十架的道路,但被殺是得勝；後者是主耶穌第二次降臨時的勝利.
\newline
\newline
主再來分為兩階段,第一階段是信徒被提,是世人看不見的；第二階段是主與被提者一同降臨,建立千年國度.可以說,被提是主再來的準備.主第一次來是藉著死得勝；第二次再來是藉著神的神性與能力得勝.
\newline
\newline
(三)拯救之主的異象
\newline
\newline
甲,啟七9-17「此後,我觀看,見有許多的人,沒有人能數過來,是從各國各族各民各方來的.」他們圍繞羔羊的寶座.這些數不過來的人,是得救的人,被羔羊寶血洗淨的人,他們歡呼說:「願救恩歸與坐在寶座上我們的神,也歸與羔羊.」
\newline
\newline
二千年前主親口說過:「天國的福音,要傳遍天下,對萬民作見證,然後末期才來到.」(太廿四14)今天這預言快要成全了.
\newline
\newline
這沒有人能數過來的極大的群眾,你也在其中嗎?
\newline
\newline
約翰壹書五章13節說:「我將這些話傳給你們信奉神兒子之名的人,要叫你們知道自己有永生.」所以每個基督徒,必須知道自己有永生.知道自已得救的人,真是有福.
\newline
\newline
乙,十四1-5這裡羔羊與十四萬四千人同在,這些人是從地上買來的,所以也是得救的人.得救之後,忠心跟隨主,所以他們也是得勝的人.
\newline
\newline
(四)審判之主的異象
\newline
\newline
甲,啟十四14-20.這段經文明顯的指出主耶穌是審判之主,莊稼收割,將麥收在倉裡,麥子是指得救的人.但在第十七節所記是收割葡萄,放在神忿怒的大酒醡中,這是指基督的審判.
\newline
\newline
有人得永生,有人受刑罰,主是輕慢不得的,人種的甚麼,收的是甚麼.蒙恩得救的重畏的事,故此得救的人要多多感謝主.
\newline
\newline
乙,誰坐在白色的寶座上?主耶穌說:「父不審判甚麼人,乃將審判的事全交與子.」(約五22).故此,坐在寶座上的是聖子.
\newline
\newline
「我又看見死了的人,無論大小,都站在寶座前,案卷展開了,並且另有一卷展開,就是生命冊,死了的人都憑著這些案卷所記載的,照他所行的受審判.」(啟廿12).
\newline
\newline
(五)主耶穌在新耶路撒冷城為燈為殿的異象
\newline
\newline
「我未見城內有殿,因主神全能者,和羔羊為城的殿.那城內又不用日月光照,因有神的榮耀光照,又有羔羊為城的燈.」(啟廿一22-23)這是約翰所見新耶路撒冷的異象,兩次提到羔羊.
\newline
\newline
甲,主耶穌是城的燈.這形像異常寶貴,不用日月光照,因羔羊發光,新耶路撒冷充滿光明.
\newline
\newline
今天主在我們心中為光,他也是世界的光,但「光來到世間,世人因自己的行為是惡的,不愛光倒受黑暗......凡作惡的便恨光,並不來就光.」(約三19-20)主第一次來,人不愛就光,但在新耶路撒冷,不用日月光照,因為羔羊是光.
\newline
\newline
乙,主為城的殿,約翰在新耶路撒冷,找不到殿,因羔羊是殿,表示在耶穌基督裡面,有完全的敬拜.
\newline
\newline
殿是神人相會之處,主也曾說過:「你們拆毀這殿,我三日內要再建立起來」.這是指他的身體說的,主以他的身體為殿.在新耶路撒冷時,無有形之殿,羔羊為城的殿,我們敬拜,成全在他裡面,藉著他與神相會.
\newline
\newline
(六)主復活的異象
\newline
\newline
甲,「婦人生了一個男孩子,是將來要用鐵杖管轄萬國的,......被提神寶座那裡去了.」(啟十二5).這婦人是指以色列人,(見上文),婦人所生的孩子,有兩個解釋.
\newline
\newline
(1)這孩子是指基督,因祂是從以色列人出來.
\newline
\newline
(2)這孩子是指教會,因教會是基督的身體,與基督是合一的.這孩子要用鐵杖轄管萬國,這話可用在基督身上,也可以用在教會身上.正如第二章27節可用在主耶穌身,上也可以用在得勝的信徒身上.
\newline
\newline
「他的孩子被提到神寶座那裡去了.」(啟十二5)這表示主耶穌的復活,也表示教會被提.
\newline
\newline
這孩子出生時,龍要吞吃他.龍是指撒旦.這與基督降生時的情形很吻合.主一降生,撒旦便想毀滅他.同時也可用在教會身上,在教會早期歷史中,教會受十次大逼迫.但基督勝過撒旦由死復活,基督的教會也同樣勝過撒旦,被提升天.
\newline
\newline
乙,啟八3-5.這段經文提到「另有一位天使」,他不是普通的天使,是指主耶穌基督而言.他在祭壇旁邊,手拿著金香爐.金香爐是表祈禱,祂的祈禱和眾聖徒的祈禱,一同獻在神面前.聖經說復活的主耶穌在神的右邊,為眾聖徒祈求,是代禱者.所以這是一對復活之主的異象.
\newline
\newline
眾聖徒的祈禱,加上基督的代禱,會使地上發生特別的情形,即神計劃的進行.神的計劃與我們的祈禱有密切的關係.神重視我們的祈禱,遠超過我們所能了解的.
\newline
\newline
主要我們與他同工,可見主實在高舉了我們.
\newline
\newline
(七)主耶穌作王的異象
\newline
\newline
甲,「我又看見幾個寶座,也有坐在上面的,並有審判的權柄賜給他們.」(啟廿4上)為甚麼說到幾個寶座?除了神的寶座和羔羊的寶座,還有甚麼人的寶座?這裡繼續說:「我又看見那些因為給耶穌作見證,並為神之道被斬者的靈魂,和那沒有拜過獸與獸像,也沒有在額上和手上受過他印記之人的靈魂,他們都復活了,與基督一同作王一千年.」(4節下)故此這裡也有得勝信徒的寶座.
\newline
\newline
乙,「在城裡有神和羔羊的寶座,他的僕人都要事奉他.」(啟廿二3)
\newline
\newline
這裡所說的是在新耶路撒冷的寶座.而第廿一章所說的是地上作王的寶座.
\newline
\newline
這異象帶給我們極大的安慰與鼓勵.
\newline
\newline


\newpage

\section{第42屆~港九培靈會~滕近輝~第七講}
\label{sec:1427}
\textbf{七新}
\newline
\newline
連結: \href{https://www.hkbibleconference.org/session-message/view/1427}{\texttt{https://www.hkbibleconference.org/session-message/view/1427}}
\newline
\newline
\hyperref[sec:1426]{< < < PREV SERMON < < <}
~
\hyperref[sec:toc]{[返目錄]}
~
\hyperref[sec:1428]{> > > NEXT SERMON > > >}
\newline
\newline
(一)信徒的新名
\newline
\newline
「並賜他一塊白石,石上寫上新名,除了那領受的以外,沒有人能認識.」(啟二17)
\newline
\newline
按當時風俗,主人請客時,要為貴賓留座位,位前放一白石,石上寫著名字,表明這是貴客.被主稱為貴賓的人,是特別蒙恩的人,是得勝者.所以新名表示新地位,與基督有新的關係,也是代表新的恩典.他與主這些屬靈的關係,只有他自己才了解.所以這裡說:「除了那領受的人以外,沒有人能認識.」(二17)
\newline
\newline
「列國必見你的公義,列王必見你的榮耀,你必得新名的稱呼,是耶和華親口所起的.」(賽六十二2)甚麼新名?「你必不再稱為撇棄的,你的地也不再稱為荒涼的,你卻要稱為我所喜悅的,你的地稱為有夫之婦,因為耶和華喜悅你,你的地也必歸他.」(賽六十二4)這是說以色列人得神喜悅,蒙神賜福,與神合一,稱為有屬靈果子之民族,是有夫之婦,不是被丟棄的,而是有屬靈的生產.
\newline
\newline
「有夫之婦」,原文是 BEULAH 這是一個極寶貴的名稱.
\newline
\newline
以色列人蒙恩而得新名.這叫我們想起基甸也改了名字,他勝過巴力,故稱為耶路巴力,意即讓巴力與他辯論.他有勝過偶像之經歷,你我若勝過偶像,神也會給我們一個新名.
\newline
\newline
彼得也得了新名,他原來很軟弱,三次不認主,後來卻剛強起來,三次為主坐監.主給他新名磯法,是磐石之意.他未有得勝經驗之先即得此名.這恩典對他必有鼓勵的作用.他想到主既給我這新名字,我怎能軟弱而令主傷心失望!後來這新名字終在他身上得了成全.
\newline
\newline
雅各在得勝時,神給他一個新名──以色列.他與神較力,被神扭了大腿窩──受神對付此時,神說:「你要叫做以色列.」
\newline
\newline
你若勝過一切攔阻,為主付上代價,你也必得新名.新名是給與得勝信徒的.
\newline
\newline
(二)主的新名
\newline
\newline
「得勝的......並我的新名,都寫在他上面.」(啟三12下)這新名不是信徒的新名,而是主的新名.這新名是寫在得勝信徒身上的.這得勝的人對主有新的了解,主對他有了新的意義.我們認識耶穌為救主；認識祂為主；管理我一生,又認識祂為我的最好朋友；最後了解他是新郎,有愛的結合.靈程一步步長進的時候,對主的認識亦一步步加深.
\newline
\newline
(三)得救的新歌
\newline
\newline
「他們唱新歌說,你配拿書卷,配揭開七印,因為你曾被殺,用自己的血從各族各方,各民各國中,買了人來,叫他們歸於神.」(啟五9)這群得救的人所唱的新歌,是新約的救贖之歌.
\newline
\newline
(四)得勝的新歌
\newline
\newline
「他們在寶座前,並在四活物和眾長老前唱歌,彷彿是新歌,除了從地上買來的那十四萬四千人以外,沒有人能學這歌.」(啟十四3)這也是新歌,但這是得勝的信徒所唱的.這十四萬四千人是得勝的信徒,所以沒有得勝經歷的人,沒有人能學這歌.
\newline
\newline
這些得勝的信徒有幾特點:
\newline
\newline
(1)是童身.(4節上)表明聖潔.
\newline
\newline
(2)跟隨羔羊.(4節中)與主有密切關係.
\newline
\newline
(3)初熟的果子.(4節下)初熟的果子,在聖經中有七種意義:
\newline
\newline
1. 基督的復活,成為睡了之人初熟的果子.(林前十五20)
\newline
\newline
2. 在神的救恩大計劃中,祂的兒女先得救,成為萬物初熟的果子,將來萬物得贖(雅一18).這包括整個大宇宙.
\newline
\newline
3. 在基督徒中,使徒時代的基督徒,成為初熟的果子.(帖後二13英譯)
\newline
\newline
4. 在某一個地區之初信主者.(羅十六4)
\newline
\newline
5. 我們的靈先得救,將來身體得贖.所以我們的靈得救是初熟的果子.(羅八23)
\newline
\newline
6. 在雅各書一章18節中,原文有「一種初熟的果子之意」,含有不同程度之意.
\newline
\newline
7. 啟示錄十四章之十四萬四千人,稱為初熟的果子.即完全成熟之初果.
\newline
\newline
(五)新婦
\newline
\newline
「我們要歡喜快樂,將榮耀歸給他,因為羔羊婚娶的時候到了,新婦也自己預備好了.(啟十九7)這是指教會說的.榮耀中的教會是基督的新婦.這是無形的教會,也是基督的身體」.今天地上的教會中有麥子也有稗子,但在榮耀中的教會,只有麥子.在地上的教會是不純全的聯合,但榮耀中的教會,是得救旳人與基督完全給合.那時愛得完全.
\newline
\newline
在新約前廿六卷書中,看見到教會的開始與發展,但在此看見教會的成全.
\newline
\newline
「新婦也預備好了.」怎樣預備?「就蒙恩得穿光明潔白的細麻衣,這細麻衣,就是聖徒所行的義.」(8節).新婦的裝飾,是光明潔白的.「聖徒所行的義」在原文與英譯本都沒有「所行」二字,原文是「就是聖徒的義」.中文加上「所行的」,因為這「義」字是多數式,所以必係指信徒所行的一件一件的義.這是信徒在因信稱義的基礎上所行的義.
\newline
\newline
新婦的裝飾,是基督的義,加上信徒藉基督的義所表顯於行為上的義,多麼寶貴的裝飾!
\newline
\newline
(六)新耶路撒冷
\newline
\newline
「我又看見聖城新耶路撒冷由上帝那裡從天而降.」(啟廿一2).新耶路撒冷是舊耶路撒冷的成全.耶路撒冷是平安之意.但舊耶路撒冷沒有平安,新耶路撒冷才有.在新耶路撒冷有七種完全.(見另文).
\newline
\newline
(七)新天新地
\newline
\newline
這名稱與新耶路撒冷不同,到底有何分別?
\newline
\newline
新天新地,彼得稱之為「神的日子」(彼後三12-13).聖經告訴我們將來有三個日子:
\newline
\newline
(1)基督的日子,在保羅書信中常見,是屬主的人得榮耀,得賞賜的日子.
\newline
\newline
(2)主的日子,是審判患難的日子.
\newline
\newline
主的日子對未重生得救的人說的,是神的審判.
\newline
\newline
(3)神的日子,重點回到神身上,「基督既將一切執政的,掌權的,有能的,都毀滅了,就把國交與父神.」(林前十五24)
\newline
\newline
神的國度有五階段與形態
\newline
\newline
(1)神在以色列人中掌權.
\newline
\newline
(2)主耶穌說:「看哪!天國在你們心裡.」可譯作「在你們中間.」
\newline
\newline
(3)基督在教會中掌權,但仍不完全.
\newline
\newline
(4)千年國度,基督作王.但國度後,撒旦暫時被釋放,又與神為敵.
\newline
\newline
(5)基督將王權交給神,這就完成.
\newline
\newline
一步一步的走向完全,王權交全神手中.在新天新地中,一切都完全了.
\newline
\newline


\newpage

\section{第42屆~港九培靈會~滕近輝~第八講}
\label{sec:1428}
\textbf{撒旦作為的七象}
\newline
\newline
連結: \href{https://www.hkbibleconference.org/session-message/view/1428}{\texttt{https://www.hkbibleconference.org/session-message/view/1428}}
\newline
\newline
\hyperref[sec:1427]{< < < PREV SERMON < < <}
~
\hyperref[sec:toc]{[返目錄]}
~
\hyperref[sec:1429]{> > > NEXT SERMON > > >}
\newline
\newline
(一)煙中的蝗虫.(啟九1-6)
\newline
\newline
開了無底坑,便有煙從坑裡往上冒,因而天昏地暗,這也是撒旦的作為,因為他是黑暗之王.他使人心黑暗,社會黑暗,一片黑暗籠罩大地.
\newline
\newline
「無底坑」,顧名思義是無底的.使我們想到聖經中提及不知足的有三樣,其中一樣是陰間.陰間開著大門,進來的人愈多愈好,永遠填不滿.人犯罪墮落是無止境的,無底的.人的痛苦也是無底的.
\newline
\newline
「無底坑」代表陰間,代表罪的權勢,總括來說,是代表撒旦權勢.
\newline
\newline
無底坑內出來的蝗虫非人間實物,乃是屬靈的象徵,有四個理由:
\newline
\newline
(1)這些蝗虫不傷害青草,真蝗虫是傷害青草的.
\newline
\newline
(2)獨傷害額上沒有神印記的人.真蝗虫不會分別有沒有神印記的人.可見蝗虫是指不屬神的人在黑暗的權勢下犯罪受苦.
\newline
\newline
(3)不准害死,只准使人受痛苦.這是與肉身無關的痛苦,是心靈的痛苦,是黑暗權勢所造成的痛苦.
\newline
\newline
(4)真蝗不能從無底坑內飛出.
\newline
\newline
(二)馬.(啟九17-21)
\newline
\newline
馬,象徵戰爭,是破壞性的.馬的胸前有甲如火,如紫瑪瑙,並硫磺.馬的頭像獅子頭,有火,有煙,有硫磺從馬的口中出來.平常的馬無這些情形.
\newline
\newline
馬的能力在口裡和尾巴上.尾巴上有頭,表示前後皆有指揮──在撒旦指揮之下,進行其破壞工作.
\newline
\newline
同時這種情形也與今日的戰爭武器相似,所以可能是末世戰爭武器的象徵.
\newline
\newline
(三)紅龍.(啟十二3)
\newline
\newline
龍表神秘能力,紅表流血,所以龍代表一種殘害人類的能力.
\newline
\newline
聖經告訴我們,這龍分三步往下墮落:
\newline
\newline
(1)撒旦本來是光明的天使,後來墮落了.第一步墮落到「空中」.聖經說有三層天,第一層就是肉眼所能見之天(空中),第二層即代表撒旦掌權之處,是超物質的領域.第三層是神的寶座所在之處,保羅曾被提到這裡聽見隱秘的話語.
\newline
\newline
(2)撒旦從空中被趕逐到地上,迫害猶太人.
\newline
\newline
(3)墮落到無底坑裡.
\newline
\newline
(四)獸
\newline
\newline
(A)十角七頭的獸.(啟十三1)
\newline
\newline
這獸是敵基督,七頭代表敵對神的七個王,也即是埃及,亞述,巴比倫,波斯,希臘,羅馬,與及接受基督教復興的羅馬國.有一個似乎受了死傷,那死傷卻醫好了.表示該國表面上接受基督,但不變其敵對神之本質.在第十七章十節說,五位已傾倒了,約翰在當時看來,確實有五個王,即埃及,亞述,巴比倫,波斯,希臘已經傾倒了,表示已經過去了.
\newline
\newline
十角是末世十個王,我們現在還不知道是那十國,現在還未出現,這十國在出現時與神為敵,這十國將在舊羅馬帝國的版圖內出現,請留意時局的發展.
\newline
\newline
(B)另有一個獸.(啟十三11)
\newline
\newline
有兩角如同羔羊,但說話卻像龍.這是撒旦最厲害的戰略.
\newline
\newline
(五)青蛙.(啟十六13)
\newline
\newline
是污穢的靈,情慾就是撒旦的工具.
\newline
\newline
(六)大淫婦.(啟十七)
\newline
\newline
大淫婦代表宗教上的敵基督.
\newline
\newline
「坐在眾水之上.」眾水代表多國多方(15節).
\newline
\newline
「騎在朱紅色的獸上,那獸有七頭十角,遍體有褻瀆的名號.」(3節下)代表宗教上的敵基督與政治上的敵基督合作.
\newline
\newline
(七)大巴比倫堿.(啟十八)
\newline
\newline
「大巴比倫」是撒旦工作的表號.「巴比倫」這字是從「巴別」而來,巴別代表恩膏舉自己敵對神而發生混亂.這正是末世的情形；國家,社會,道德,政治均在混亂的情況之中.
\newline
\newline
從上述的七象,我們可以了解撒旦的作為.
\newline
\newline


\newpage

\section{第42屆~港九培靈會~滕近輝~第九講}
\label{sec:1429}
\textbf{撒旦的七個名字}
\newline
\newline
連結: \href{https://www.hkbibleconference.org/session-message/view/1429}{\texttt{https://www.hkbibleconference.org/session-message/view/1429}}
\newline
\newline
\hyperref[sec:1428]{< < < PREV SERMON < < <}
~
\hyperref[sec:toc]{[返目錄]}
~
\hyperref[sec:1430]{> > > NEXT SERMON > > >}
\newline
\newline
撒旦的七個名字正與撒旦作為的七象互相呼應:
\newline
\newline
(一)亞巴頓.(啟九11)
\newline
\newline
這無底坑的使者就是撒旦,這是歷代許多解經家所同意的.
\newline
\newline
亞巴頓,是希伯來文,意即沉淪,滅亡,這名很符合撒旦的工作.
\newline
\newline
此名正與無底坑上冒之煙中之蝗虫的意義相合,皆代表黑暗破壞的權勢.
\newline
\newline
(二)亞玻倫.(啟九11下)
\newline
\newline
意即毀滅者.這可代表撒旦的特性,牠的作為是毀滅人的.此名正與第十七節內代表戰爭和毀滅的馬互相呼應.
\newline
\newline
在人類歷史之中「五年一小戰,三年一大戰」這雖然不是精確的統計,但使人不難相信撒旦的作為是戰爭.
\newline
\newline
(三)大龍.(啟十二9)
\newline
\newline
又叫紅龍,紅表流血,龍表神秘的大能力,是殘忍流血的能力.
\newline
\newline
(四)古蛇.(啟十二9)
\newline
\newline
蛇表狡猾.此名正與那兩角如同羔羊的獸互相應合.(十三11)裝作基督,其實他是假基督.若將羔羊解作良善,撒旦就是羊皮披在狼身上,使人不會提防他.
\newline
\newline
在但以理書上敵基督的預象,乃是「乘人不備」就毀滅人(但八25),這是它的拿手好戲.古蛇正符合他狡猾的形像.
\newline
\newline
(五)魔鬼.(啟十二9)
\newline
\newline
這名字代表牠是窮兇極惡的,正與十三章那七頭十角的獸的形像.互相應和.
\newline
\newline
(六)撒旦.(啟十二9)
\newline
\newline
原文意即敵對者──與神為敵.此名正與十七章所載之大淫婦又名大巴比倫的形像互相呼應,因巴比倫之基本意義是與神為敵.
\newline
\newline
(七)迷惑者.(啟十二9)
\newline
\newline
此名正與十六章13節所說青蛙的意義相合,極其污穢,撒旦用情慾去迷惑人.
\newline
\newline
撒旦這七個名字,使我們認識他的作為.撒旦不喜歡人讀啟示錄,因為本書將它的面具揭下來,顯出其本相.聖經記載撒旦的十二個名字,啟示錄佔其中七個,其他五個名字如左:
\newline
\newline
「那惡者.」(太十三19)他來把人心裡的道奪去.
\newline
\newline
「說謊人之父.」(約八44)他自己說謊,也引人說謊.今天很多人的哲學是「假話說得多,便變成真話.」故此,在今天混亂的世代裡,要用主的話做試金石,不憑感覺,只憑主的話,主的話如兩刃利劍.
\newline
\newline
「世界的神.」(林後四4)是世界人所崇拜的.
\newline
\newline
「世界的王.」(約十四30,十六11)世人的思想被牠操縱.
\newline
\newline
「試探者.」(太四3)試探是側面攻擊.撒旦常常避開正面,從不提防的角度去試探人.
\newline
\newline
如拔鐵釘,在旁邊打打敲敲,然後一拔便拔出來.那試探人的便是要把我們從神的救恩中,團契中,愛心中拔出.
\newline
\newline


\newpage

\section{第42屆~港九培靈會~滕近輝~第十講}
\label{sec:1430}
\textbf{七個數字}
\newline
\newline
連結: \href{https://www.hkbibleconference.org/session-message/view/1430}{\texttt{https://www.hkbibleconference.org/session-message/view/1430}}
\newline
\newline
\hyperref[sec:1429]{< < < PREV SERMON < < <}
~
\hyperref[sec:toc]{[返目錄]}
~
\hyperref[sec:1431]{> > > NEXT SERMON > > >}
\newline
\newline
(一)善意的七個數字
\newline
\newline
啟廿一11-20這段經文是描寫新耶路撒冷情況,其中常有數字出現,藏有屬靈深意.
\newline
\newline
(一)十二(12,14,21)是完全之意,是神人和諧之完全.十二支派的名字在十二個門上,十二使徒的名字刻在城的十二根基上.一切得救蒙恩的人都包括在其中.
\newline
\newline
(二)一百四十四(17).這是十二乘十二的積,是完全的完全之意.人與人之間,人與神之間,人與天使之間都有完全合一的關係.恢復神起初創造的情況－宇宙的大和諧.
\newline
\newline
(三)「長闊高都是四千里」,(16)長闊都是四千里,是很容易想像的.但高也是四千里,真是不容易想像.這與上文所說城牆高一百四十四肘是有矛盾的.可見完全是象徵的意思,而非實際的尺寸.
\newline
\newline
有人說城牆連山的高度才有四千里,因為第十節曾提到山.但詳細看,那是天使帶約翰所到之高山,指示他看新耶路撒冷,並不是說新耶路撒冷在高山之上.
\newline
\newline
長寬高是三度,各長四千里,也等於三乘四,即一萬二千里.
\newline
\newline
在聖經中,一千表示成全,十二乘十二再乘一千,就是完全的完全.
\newline
\newline
十字架上的呼聲,「成了!」但今天救恩還未完成,要等到身體得贖後進入耶路撒冷才是完全的完全.
\newline
\newline
四千里高,也就是地與天相連之意,中國人的「人天合一」的觀念,這時成全了.聖經的真理實在是中個人理想的成全.
\newline
\newline
主耶穌教我們禱告說:「願你的旨意行在地上,如同行在天上.」這禱告,今天還未完成,你我生活中,多時不遵行神的旨意.但在新耶路撒冷中就得以成全.
\newline
\newline
這城將天地連起來,不分地上與天上,這異象,真是聖靈所感動而顯明約給約翰看的.
\newline
\newline
(四)一千二百六十天(啟十一3).
\newline
\newline
我們詳細查考啟示錄與但以理書時,就發現,一千二百六十天與三年半有不同的重點,三年半是指大災難,敵基督橫行,耀武揚威的日子,三年半與四十二個月的意思與重點相同.
\newline
\newline
而一千二百六十天的重點,是神在大災難期間看顧他的百姓.十二章六節將此意顯明.
\newline
\newline
在此時期內,有兩個見證人,要為主作見證一千二百六十天,這是最後,最有力,最寶貴的見證.雖然,此時敵基督勢力猖狂,但這兩個見證人,得神保守,看顧與幫助,他們工作未完,是不會死的.有許多主的僕人,使女,就是因為得此寶貴經文,而有力量,有安慰.是的,神的計劃在我們身下未完成之前,我們是不會死的.
\newline
\newline
有人說,耶路撒冷由被回教徒佔領,直到得解放是一千二百六十年,在此時期猶太人到處飄流,但在一千二百六十年期間,神卻保守看顧他們.
\newline
\newline
(五)十四萬四千人(啟七4)
\newline
\newline
十四萬四千,是十二乘十二,再乘以一千.十二乘十二是完全的完全；十千是完成的數目,所以十二乘十二,再乘以一千,就是完全的完全之完成.
\newline
\newline
完全與完成二者之間有分別.比如今天已有完全的救恩,但將來連身體也得贖之時,救恩才完成.
\newline
\newline
所以十四萬四千的異象,使我們提早看見最後的完成.
\newline
\newline
啟示錄內有兩特點:
\newline
\newline
(1)地上景況與天上景況常交替出現,地上的景況是充滿痛苦與罪惡:敵基督橫行,惡勢力遍地.但天上的景況充滿榮耀與喜樂,寶座上有神與羔羊.我們若學會了用天上的亮光來看地上一切,我們的思想就不會被地上的景況統治.
\newline
\newline
(2)在啟示錄內我們先看見最後的完成.第四章是預言的開始.但一開始就看見完成的景況,然後一步一步的顯明如何完成.
\newline
\newline
十四萬四千代表滿足的數目,表示得救的人,一個不少,一個不多,正按神的計劃成全.如果我們有十四萬四千位朋友,相信很多會忘記.但每一個得救的人,神都不會忘記,神都看顧保守.
\newline
\newline
昔日祭司的胸牌與肩上有以色列支派的名字,天天擔當他們的軟弱,也天天將他們放在心上,記念他們,這正是主耶穌的預表.誰也不能將他們從主的手中奪去.
\newline
\newline
(六)二十四(四4).本書多次提到二十四.
\newline
\newline
廿四是十二的雙倍.他們代表所有屬神的人,此節的冠冕是表示得神的賞賜.啟示錄裡提到兩種冠冕,原文有兩個不同的字:一是君王的冠冕,其次是得勝者的冠冕.這二十四位長老的冠冕是得勝者的冠冕.他們在神面前不住的讚美神.
\newline
\newline
(七)千年國度(啟二十章四節).第一次復活的人與基督一同作王一千年.
\newline
\newline
「我又看見幾個寶座,也有坐在上面的,並有審判的權柄賜給他們,我又看見那些因為給耶穌作見證,並為神之道被斬者之靈魂,和那沒有拜過獸與獸像,也沒有在額上和手上受過他印記之人的靈魂,他們都復活了,與基督一同作王一千年.(啟二十4)根據這段經文,這些與基督作王一千年的,是在大災難被殺的聖徒,加上由空中與主一同降臨的信徒.
\newline
\newline
有些奧祕,當我們祈禱,等候,思想,主就解開；但有些是神不讓我們早知道,我們就不必早知道.
\newline
\newline
這千年國度是一個完成.神用六天創造天地,第七天休息.在舊約,神多次向以色列人說到榮耀的盼望,在這一千年中,神對以色列人的應許就完全成就.
\newline
\newline
(二)惡意的七個數字
\newline
\newline
(一)十角七頭的獸.(啟十三1)
\newline
\newline
這七頭十角的角獸是敵基督,七與十皆表示完全,這就表示敵基督集一切與神為敵萬物之大成.敵基督有五種意義:
\newline
\newline
(1)廣義:凡與基督為敵之領袖皆稱為敵基督,例如「現在已經有好些敵基督出來了.」(約壹二18)約翰時代已有好些敵基督出來了.
\newline
\newline
(2)歷史上的意義:「又是七位王,五位已經傾倒了,一位還在,一位還沒有來到.」(啟十七10)七頭代表亞述,埃及,波斯,希臘,羅馬等七大帝國.在歷史上這是很長的過程.
\newline
\newline
(3)末世的意義.七角十角之獸的「本身」是在末世要出現的政權.「那先前有,如今沒有的獸,就是第八位.」(十七11).
\newline
\newline
(4)集體意義.十角代表十王,同時存在,有集體之意.
\newline
\newline
(5)個人意義.「並有那大罪人,就是沉淪之子,顯靈出來.」(帖從二3)可見有一人要做總代表.
\newline
\newline
有關這獸的七頭中,有一個似乎受傷,那死傷卻醫好了.這也是敵基督摹倣基督的一件事:基督曾藉著死敗壞掌死權的魔鬼,後來卻復活過來.敵基督在此點照樣作.當羅馬帝國接受基督教為國教,好像給撒旦一個打擊,但卻造成政教合一,破壞教會的純潔,也破壞教會的真正工作.這獸藉死得著權勢,正如基督藉死得勢.撒旦藉表面的失敗,完成最大的詭計.藉著接納基督教為國教,撒旦的作為進入教會內,破壞教會的純潔.
\newline
\newline
(二)六百六十六(啟十三18)
\newline
\newline
這數字主要有兩種解釋:
\newline
\newline
(1)此數字代表尼祿王.尼祿王的名字本是拉丁文,但用希伯來文書寫時,其字母所代表之數字加起來總和是六百六十六.
\newline
\newline
(2)此數字代表羅馬帝國,因為希臘文中的「拉丁」一字,其字母所代表數目之總和是六百六十六.
\newline
\newline
所以六百六十六代表末世羅馬帝國復興,這是敵基督的勢力.有關羅馬帝國復興的說法有二:
\newline
\newline
(1)羅馬帝國的特性在末世復興.
\newline
\newline
(2)羅馬舊版圖上興起一政權.
\newline
\newline
我覺得這兩說法可以合起來,它是一個敵對神的極大的勢力,同時也有政治上的形式.
\newline
\newline
就六六六本身來講,「六」少於七,故代表人的不完全,人有虧欠故落在神審判之下.
\newline
\newline
六百六十六是人發展到最高峰.人不完全,但有智力,能運用從所領受之才幹,但已變質,造成敵對神的一切一切.
\newline
\newline
正如但以理所見的巨像,「極其光耀,形狀甚是可怕.」末世敵基督要將人的成就集合起來敵對神.
\newline
\newline
(三)四十二個月(啟十三5).即三年半.
\newline
\newline
上文已經說過,一千二百六十天的數字含有好的意義,代表神在患難中看願意屬祂的人,例如兩個見證人傳道一千二百六十天；神的百姓在曠野被神養活一千二百六十天.
\newline
\newline
但四十二個月含有不好意義,代表敵基督做成痛苦與毀滅的一段時間.
\newline
\newline
在十三章看見敵基督得權柄,任意而行四十二個月,殺害屬神的人.
\newline
\newline
「他們要踐踏聖城四十二個月.」(啟十一2)這也是敵基督的作為.在此敵基督作為下,計算時間的單位是月,不能一年一年的數,須一個月一個月的數.因為事情發生很多,用一年做單位,就不準確,須用一個月一個數.
\newline
\newline
後三年半,是敵基督最猖狂之時,主耶穌說過:「若不減少那日子,凡有血氣的,總沒有一個得救的,只是為選民,那日子必減少了.」可見神仍有恩典.
\newline
\newline
(四)六百里(十四20).十四章是說到兩個大收割,十四節至十六節,是收割麥子,把好的麥子收在倉裡.十七至二十節是收割葡萄,放在神忿怒的大酒醡中,那酒醡踹在城外,就有血從酒醡裡流出來,高到馬的嚼環,遠有六百里.「六百里」的原文是一千六百士他迪,等於六百里,也等於二百英哩.
\newline
\newline
一千六百有兩種解釋:
\newline
\newline
(1)巴勒司坦地從南至北共有一千六百士他迪.這是代表全巴勒司坦的大戰爭.
\newline
\newline
(2)一千六百是四十乘四十.「四」字是代表地的四方,包括全地,換句話說,是代表完全的審判.
\newline
\newline
(五)二萬萬(啟九16).亦有兩種解釋:
\newline
\newline
(1)代表「多」之意.在五章十一節說天使有千千萬萬,英文譯本是「十千乘十千」.這裡的二萬萬是多而又多.可見這次戰爭動員人數之多.
\newline
\newline
(2)特別要注意:「他們的數目我聖見了.」這是著重的說法.我相信這話很有意義.約翰時代的戰爭一定不會有這麼多人參加,但現代戰爭動員二萬萬人,並非不可能.
\newline
\newline
(六)三分之一(啟八7-13)多次出現.如果把這些三分之一加在一起,超過全數,所以不是實數,而是表示小部份.這些苦難有範圍限定.患難過程中,有神管制.
\newline
\newline
(七)五個月(啟九5,10).有不同解釋:
\newline
\newline
(1)從前猶太國的蝗蟲災,最長為五個月,這是有季節性的.但這些蝗蟲若是真蝗蟲,便符合這解釋.不過這些蝗蟲是從無底坑的煙中出來的,當然不是真蝗蟲.(見上文)
\newline
\newline
(2)「五」是一半之意.十表完全,五是十之一半,意思是說,這災難的痛苦雖大,但比起將來的痛苦只是一半而已.可見將來的痛苦深重,實不能用筆墨形容.
\newline
\newline
(3)五有恩典之意,例如大衛用五光滑的石子,擊殺哥利亞,在神恩典下,半力便夠了.保羅說:「在教會中,用五句教導人的話,強如說萬句方言.」
\newline
\newline
第(2)(3)合在一起,就是說,今天罪惡的痛苦,不過是一半的痛苦.其中仍有神的恩典.
\newline
\newline


\newpage

\section{第42屆~港九培靈會~滕近輝~第十一講}
\label{sec:1431}
\textbf{七讚}
\newline
\newline
連結: \href{https://www.hkbibleconference.org/session-message/view/1431}{\texttt{https://www.hkbibleconference.org/session-message/view/1431}}
\newline
\newline
\hyperref[sec:1430]{< < < PREV SERMON < < <}
~
\hyperref[sec:toc]{[返目錄]}
~
\hyperref[sec:1421]{> > > NEXT SERMON > > >}
\newline
\newline
啟示錄有七讚,這七次讚美是因約翰見到天上的榮耀景象而發的.讚美是屬天的,只有屬天的聖徒纔能發出讚美來,屬地的信徒實在發不出心靈的讚美!為甚麼我們有時失去讚美呢?因為我們沒有仰望主.我們只看環境,看自己,看地上的事,因此,我們會灰心喪膽.主的話一次一次的提醒我們,要思念天上的事,不要思念地上的事.思念天上的事,就有讚美.
\newline
\newline
(一)(啟四8-11).在本段經文中,當將來的事展開時,約翰就見到天上寶座的榮耀,隨即聽見四活物和廿四位長老的讚美.四活物是一種特別的天使,即基路伯,代表被救贖的人性.自從亞當夏娃被逐出樂園以後,就有基路伯把守生命樹的道路.可見四活物與重得生命有關.民數記第二章記載百姓安營,有四面旗幟懸掛在營的四角,上面有四活物的形象,故此這四活物顯然與神的百姓有關.新約的四福音,是代表到主耶穌的四方面,祂是王,僕人,人子,神子.而主耶穌正代表原有的人性,亦即得救之人性.本段經文說四活物在寶座中與神一同作王,故四活物與得救的人有關.他們是讚美者.廿四位長老是代表新約和舊約中屬神的人.他們的讚美包括兩重點:
\newline
\newline
(1)「聖哉,聖哉,聖哉!」他們認識神的聖潔完全,有此認識就自然產生敬拜和讚美.我們對於好人,尚且尊敬仰慕,何況最完全的神呢?人所拜的如何,他也必如何!我們所拜的是聖潔的神,我們必趨於聖潔.
\newline
\newline
(2)神創造萬物都有其美好的旨意在內,神所作所為都有祂的定旨,安排,和計劃.最後也必然會成全.今天我們所看見一切混亂,污穢,罪惡之事都要過去.想到這裡,我們就不能不讚美了!主叫我們再看見祂的計劃必要完成.無一事能叫神失敗,但撒旦至終必要失敗!但願我們能夠在今世混亂中,看見神的成功與得勝而發出讚美!
\newline
\newline
(二)天使的讚美!(啟五8-14).天上地下一切被造之物都在讚美,千千萬萬的天使加入大唱詩班,將來萬物得贖,所以一同讚美神了!
\newline
\newline
今日因人的罪影響到一切被造之物,不能發出讚美的聲音；但將來這一切將要進入神的榮耀中,他們已被救贖,進入大讚美的行列裡:他們唱的是新歌,他們大聲讚美說:「你配拿書卷,配揭開七印,因為你曾被殺,用自己的血,從各族,各方,各民,各國中買了人來,叫他們歸於神.」在這段經文中,我們看見兩個對比；
\newline
\newline
(1)約翰大哭,因為沒有配展開那書卷的人.
\newline
\newline
(2)勝利歡呼:「看哪,猶太支派中的獅子,大衛的根,祂已得勝!」哀哭變為讚美.
\newline
\newline
保羅誇自己的軟弱,因為主的恩在他的軟弱上顯得完全,使他能不斷的發出感謝和讚美來.
\newline
\newline
(三)啟七9-16「我觀看見有許多的人,沒有人能數過來,是從各國,各族,各民,各方來的,站在寶座和羔羊面前,身穿毛衣,手拿棕樹枝；大聲喊著說,願救恩歸與坐在寶座上我們的神,也歸與羔羊.」他們為何讚美呢?因為坐寶座的,要用帳幕覆庇他們,他們不再飢,不再渴,日頭和炎熱,也必不傷害他們.因為寶座中的羔羊,必牧養他們,領他們到他們生命的泉源,神必擦去他們一切的眼淚!這是在神面前蒙福的情況,這些數不盡的群眾,手拿棕樹枝,從他們心中湧出感謝和讚美來.
\newline
\newline
(四)十四萬四千人讚美(啟十四1-3)這歌誰能唱呢?除了地上的十四萬四千人外,無人能唱.世上的歌雖有的難唱,但肯花時間去學,就斷沒有不可能唱的.但是這歌無人能學.因為這是他們的生活見證之歌.他們具有五個條件:(1)童身表示聖潔,林後十一章,說信徒如同貞潔的童女,獻給基督.(2)是從人間買來的,(3)口無謊言,(4)無瑕疵.(5)跟從羔羊.我們在地下如何,將來在天上也如何.在地下能跟從主,將來在天上必仍跟從主.有一位姊妹對我說,我們事奉主不應為了獎賞；對的,但主愛我們,將來把獎賞賜給的.神對亞伯拉罕說:「我要大大的賞賜你」英譯聖經說:「我就是你的獎賞.」誠然,主是我們杯中的分.我們不是為了得獎而愛主,但主卻歡喜將獎賞賜給我們.我們看重恩典,但更看重賜恩的主.
\newline
\newline
(五)那些勝了獸和獸像者(啟十五2-4)他們站在玻璃海上,拿著上帝的琴,唱神僕人摩西的歌,和羔羊的歌說:「主神!全能者啊!你的作為大哉,奇哉!萬世的王啊!你的道途義哉!誠哉!」這些人乃是為主殉道被殺者,他們曾受過迫害痛苦,他們得勝了!現在歡呼讚美!
\newline
\newline
(六)祭壇下的受冤者(啟十六4-5,7)這裡的讚美與十五章的讚美是一樣的.祭壇下的靈魂,看見天使倒神大怒的碗在江河和眾水裡,他們就發出讚美的聲音說:「主神!全能者啊,你的判斷,義哉!誠哉!」神的兒在地下受苦,受逼迫,神必記念,施行審判刑罰,到了時候,必為他的兒女伸冤.
\newline
\newline
(七)大眾的總讚美(啟十九1-8)在天上的群眾大聲說哈利路亞,救恩,榮耀,權能,都屬乎我們的神,祂的判斷是誠實公義的.這一個讚美是包括天上的廿四位長老,四活物,從寶座出來的聲音,群眾的聲音,大雷的聲音,包括宇宙間的一切聲音.啟示錄乃是一本讚美的書,我們的感情,意志,都應當讚美祂.尤其在這末世的日子,對主當敬拜感恩.(西四3)在患難中也當讚美,張開信心之口讚美神!
\newline
\newline
但願我們都能充滿著天上的讚美!
\newline
\newline


\newpage

\section{第42屆~港九培靈會~滕近輝~第十二講}
\label{sec:1421}
\textbf{天上七次發聲讚}
\newline
\newline
連結: \href{https://www.hkbibleconference.org/session-message/view/1421}{\texttt{https://www.hkbibleconference.org/session-message/view/1421}}
\newline
\newline
\hyperref[sec:1431]{< < < PREV SERMON < < <}
~
\hyperref[sec:toc]{[返目錄]}
~
\hyperref[sec:1446]{> > > NEXT SERMON > > >}
\newline
\newline
天上的聲音,一定是重要的聲音,屬主的人要留心領受.
\newline
\newline
(一)「此後,我聽見好像,群眾在天上大聲說:「哈利路亞,救恩,榮耀,權能,都屬乎的神.」(啟十九1)
\newline
\newline
這是多麼令人興奮的聲音,為什麼要讚美耶和華?「因他判斷了那用淫行敗壞世界的大淫婦,並且向淫婦討流僕人血的罪,給他們伸冤.」
\newline
\newline
所以這天上的聲音是勝利的聲音,大巴比倫的權勢與成就是可驚的,但從終點的眼光看,她要落在神公義的審判之下.
\newline
\newline
(二)「我又聽見從天上有聲音說,我的民啊,你要從那城出來,免得與他一同有罪,受他所受的災殃.」(啟十八4)這聲音要一切屬神的人離開巴比倫；屬神的人是屬新耶路撒冷的,不可住在巴比倫城中,這天上的聲音呼叫我們出來!出來!
\newline
\newline
很多屬神的人住在大巴比倫城內,在與神為敵的權勢下,行事為人,與神為敵,與十字架為敵.
\newline
\newline
耶利米書有兩個基本信息:
\newline
\newline
(1)回來──呼喚離開主的人回到神的面前.耶利米三次說:「回來!回來!回來!」
\newline
\newline
(2)出來──呼喚那些住在與神為敵的大城中的人出來!耶利米也有三次這樣的呼聲.(耶三12,14,22,五十8,五十一45,6)
\newline
\newline
「回來!」「出來!」是流淚的先知耶利米的信息,啟示錄最後的信息也是「我的民哪!你要從那城出來!」如果你遠離主,主要你回來.若你生活在大巴比倫中,愛世界多過愛主,主要你出來.
\newline
\newline
(三)「我聽見天上有聲音說:『你要寫下,從今以後,在主要面而死的人,有福了』聖靈說:『是的,他們息了自己的勞苦,作工的效果也隨著他們』.」(十四13).這節聖經給我很深刻的印象,先父去世時,在一幅輓帳上正寫著這節聖經.他是在講台上講道時腦充血死的,發生事情之後,昏迷了一週,在次主日晚忽唱了一首詩便離世,這情景我永不能忘記.
\newline
\newline
不錯,這節經文可以用在一切離世信徒身上,但這節經文有特別所指.「從今以後」就是指後三年半的時期,這是聖徒大受逼迫的時期.為主犧牲的人有福了.正如二十章六節所說:「在頭一次復活的人有福了,聖潔了,第二次的死,在他們身上沒有權柄,他們必作神和基督的祭司,並要與基督一同作王一千年.」他們息了工作的勞苦,作工的果效也隨著他們.」
\newline
\newline
為主付代價的人有福了,天上有聲音安慰他們.
\newline
\newline
(四)「我聽見天上有大聲音說,我神的救恩,能力,國度,並他基督的權柄,現在都來到了,因為那在我們神面前晝夜控告我們弟兄的已經被摔下去了.」(啟十二10).
\newline
\newline
這聲音是在龍被摔在地上時發出的(9節).甚麼時候龍被摔到地上?請看第五節「婦人生了一個男孩子,是將來要用鐵杖轄管萬國的,她的孩子被提到上帝那裡去了.」這婦人是指以色列人,孩子是基督,也是基督的身體－教會.
\newline
\newline
孩子被提,一方面是指主復活的事,一方面也指教會被提的事.所以教會被提時龍被摔在地上.
\newline
\newline
教會被提是在七年大災難之前,這是我從前所相信的.但我多年查攷,現在相信信徒是在七年災難之中被提,即是啟示錄十,十一,十二章所談到的時候.
\newline
\newline
其他四次天上的聲音,均與主耶穌再來有關.(啟十4,8,十一15,十二10)
\newline
\newline
結論:「證明這事的說,是了,我必也來.阿們!主耶穌啊,我願你來!」(啟廿二20)
\newline
\newline
主在啟示錄正文最後一句,證明祂再來是真的,好像蓋上印.我們的反應是「阿們!主耶穌啊!我願你來!」
\newline
\newline
在座的眾兄妹,都預備好了嗎?等候主來!
\newline
\newline


\newpage

\section{第42屆~港九培靈會~王永信~第一講}
\label{sec:1446}
\textbf{你在那裡?}
\newline
\newline
連結: \href{https://www.hkbibleconference.org/session-message/view/1446}{\texttt{https://www.hkbibleconference.org/session-message/view/1446}}
\newline
\newline
\hyperref[sec:1421]{< < < PREV SERMON < < <}
~
\hyperref[sec:toc]{[返目錄]}
~
\hyperref[sec:1447]{> > > NEXT SERMON > > >}
\newline
\newline
創世記 3:1-3:13
\newline
\begin{longtable}{cl}
\hline
\hline
章節 & 經文 (和合本修訂版)\\
\hline
3:1 & \begin{tabularx}{0.7\textwidth}{X} 耶和華神所造的，惟有蛇比田野一切的走獸更狡猾。蛇對女人說：「神豈是真說，你們不可吃園中任何樹上所出的嗎？」 \end{tabularx} \\ \\ \relax
3:2 & \begin{tabularx}{0.7\textwidth}{X} 女人對蛇說：「園中樹上的果子，我們都可以吃； \end{tabularx} \\ \\ \relax
3:3 & \begin{tabularx}{0.7\textwidth}{X} 只是園子中間那棵樹的果子，神曾說：『你們不可吃，也不可摸，免得你們死。』」 \end{tabularx} \\ \\ \relax
3:4 & \begin{tabularx}{0.7\textwidth}{X} 蛇對女人說：「你們不一定死； \end{tabularx} \\ \\ \relax
3:5 & \begin{tabularx}{0.7\textwidth}{X} 因為神知道，你們吃的日子眼睛就開了，你們就像神一樣知道善惡。」 \end{tabularx} \\ \\ \relax
3:6 & \begin{tabularx}{0.7\textwidth}{X} 於是女人見那棵樹好作食物，又悅人的眼目，那樹令人喜愛，能使人有智慧，她就摘下果子吃了，又給了與她一起的丈夫，他也吃了。 \end{tabularx} \\ \\ \relax
3:7 & \begin{tabularx}{0.7\textwidth}{X} 他們二人的眼睛就開了，知道自己赤身露體，就編織無花果樹的葉子，為自己做成裙子。 \end{tabularx} \\ \\ \relax
3:8 & \begin{tabularx}{0.7\textwidth}{X} 天起了涼風，那人和他妻子聽見耶和華神在園中來回行走的聲音，就藏在園裡的樹木中，躲避耶和華神的面。 \end{tabularx} \\ \\ \relax
3:9 & \begin{tabularx}{0.7\textwidth}{X} 耶和華神呼喚那人，對他說：「你在哪裡？」 \end{tabularx} \\ \\ \relax
3:10 & \begin{tabularx}{0.7\textwidth}{X} 他說：「我在園中聽見你的聲音，我就害怕；因為我赤身露體，我就藏了起來。」 \end{tabularx} \\ \\ \relax
3:11 & \begin{tabularx}{0.7\textwidth}{X} 耶和華神說：「誰告訴你，你是赤身露體呢？莫非你吃了那樹上所出的，就是我吩咐你不可吃的嗎？」 \end{tabularx} \\ \\ \relax
3:12 & \begin{tabularx}{0.7\textwidth}{X} 那人說：「你賜給我、與我一起的女人，是她把那樹上所出的給我，我就吃了。」 \end{tabularx} \\ \\ \relax
3:13 & \begin{tabularx}{0.7\textwidth}{X} 耶和華神對女人說：「你怎麼會做這種事呢？」女人說：「那蛇引誘我，我就吃了。」 \end{tabularx} \\ \\
[1ex]
\hline
\hline
\end{longtable}
香港是一個繁榮的都市,我每次到香港來都感到香港－在物質方面都有顯著的進步.然而叫我感到詫異的,就是為甚麼香港在傳福音的事工上,和信徒在靈性上,幾十年來依然故我未見長進?香港的教會為何仍未得著復興?這個問題是與我們每一個信徒有切身關係的,所以我們要從人類的首頁歷史來看看.
\newline
\newline
創世記第三章告訴我們一件重要的史實－就是人類如何墮落.神創造了亞當和夏娃,神的心意是要他們掌權,成就神旨；但他們犯罪墮落,從此神不能在他們身上得著榮耀了.始祖原來是可以與神自由交通的,但犯罪後便失去這種特權,且要受神的刑罰.世人也是如此,因為甘願走魔鬼的道路,寧願向魔鬼投降.不過,最令神難過的,莫過於祂自己的兒女,也走貪愛世界的道路了.亞當夏娃正用自己的自由意志,要選擇離開神,走自己的道路.當他們離開神後,在他們身上發現有兩種感覺:
\newline
\newline
(一)兩種感覺
\newline
\newline
(A)「羞恥」
\newline
\newline
罪使他們感到自己是赤身露體,又使他們良心受責備；本來他們是屬赤子之心,但犯罪墮落後,便發覺自己的本相因而感到羞恥.罪帶給人第一個感覺就是羞恥,知道自己是犯罪得罪神.在這情形下,人自然的反應就是用自己的方法來遮蓋自己的罪.他們拿無花果樹的葉子,為自己編作裙子,為要隱藏自己在神面前的羞恥,但人怎能用葉子來遮蓋神的忿怒呢?我們何嘗不是如此?常把自己的外表裝得十分敬虔與熱心,但這一切都是虛偽的,裡面卻不敢直接面對神,因為其中有很多罪惡未經對付,內心走世界的路,對神不順服,扺抗神,敵對神!若我們只有外表的敬虔,這就等於亞當夏娃以無花果樹葉子所編成的裙子,也就是人用方法來偽裝自己.不過,這在神光照下是無所遁形的.神如何對付亞當,今日也要如何對付我們,因為亞當的失敗,也是我們的失敗；時至今日,世人所編的仍是亞當所編的裙子,因為世人用盡方法來遮掩自己的錯過.
\newline
\newline
(B)「畏懼」
\newline
\newline
亞當夏娃是神所造的,他們本性原是無罪的.神將他們安放在伊甸園中居住,每當神到園中探問他們時,他們有如兒女之於父親一般的親切.但現今卻完全不同了,他們心裡出了毛病,所以當神到伊甸園,他們便產生懼怕的心,躲藏起來不敢見神.「藏在園裡的樹木中,躲避耶和華神的面.」這也是人的方法,樹怎能使人逃避神的面呢?世人用盡自己的方法來逃避神的面,令人感到歎息的,就是神的兒女也起來躲避神哩!神要我們熱心服事他,我們卻對這些事避而不談,躲在自己的知識,才幹,金錢,假裝的事奉後面,而逃避那真正的事奉和基督徒的本份.
\newline
\newline
亞當犯罪,神沒有放過他,起來尋找他；今天神也不會放過我們,也要起來尋找我們.耶和華對那人呼喚說:「你在那裡?」這句話不但是問「你在什麼地方?」也是問「你現在的境況如何?」「你生活的態度如何?」今天神也要向眾教會發出這一個問題:
\newline
\newline
(二)兩種立場
\newline
\newline
(A)你的信仰的立場在那裡?
\newline
\newline
教會在今日面臨一個極嚴重的挑戰,就是信仰變了質,信徒的生活每況愈下.歐美的教會比東方的教會尤甚,歐洲的教會又比美國的教會更為敗壞；教會的信徒走了樣,變了質.他們否認耶穌的神的兒子－否認耶穌是童貞女馬利亞所生；對於祂被釘死,以至埋葬,復活,升天的事產生不少疑團；他們將聖經的真實性和價值貶低.事實告訴我們,遠東的教會,信仰純正,但這現象是否跟時間先後而會改變呢?有一次我與一位西教士偶然談及這問題,他指出:美洲教會比遠東教會發展得早,他們信仰變質也早；歐洲教會比美洲教會發展得更早,他們的信仰變質也更早.遠東教會的信仰能否保持不變,這要待時間來證明了.
\newline
\newline
神要向我們遠東的教會發問:我們信仰的立場究竟在那裡?因為信仰是我們的根基,如果我們的教會夫去了純正的信仰,一切就都要完了.目前新神學的勢力已滲入到遠東的教會來,在台灣已有不少神學院變了質,受新派思想所影響；他們不是訓練神的僕人,而是練魔鬼的僕人.各位親愛的弟兄姊妹,我們寧可夫去一切物質有價值的東西,甚至有形教會的建築物,也不能失去教會純正信仰的立場.我們是否在這信仰動蕩的世代中,仍能站立得住,保守純正的信仰?這是現今神對他兒女們最關切的一個問題.
\newline
\newline
(B)你的生活立場如何?
\newline
\newline
神不但要知道我們信仰的立場在那裡?祂也要向我們發問:我們的信仰是否與生活打成一片?曾否活出主的道,活出基督徒的樣式來?世人之所以不看重教會,青年人之所以遠離教會；正因為他們發覺教會的信徒言行不一致,信徒的生活沒有見證,以至世人無法從這沒有見證的生活裡認識神!我們若想世人尊重教會,不能單靠純正的信仰,因為信仰的純正並不能感動人,並不能使人認罪悔改相信；唯有彼此相愛,充滿愛的生活,才可以表現出「基督徒」的言行是充滿著基督馨香之氣.
\newline
\newline
今天我們的信仰如何?生活的立場如何?對世界的態度又如何?我們應當深深自省,曾經有不少相當熱心的基督徒,蒙神賜福與他們,使他們的環境富裕,但可惜每當我們在物質富裕以後,靈性便開始退步了,因為人若貪愛世界,愛父的心,就不在他裡面了.在台灣的教會,以一情形來說主日只有早堂或午堂的崇拜,而沒有晚堂的聚會.因為台灣禮拜日晚間的電視節目特別精彩.所以人就不來教堂了.教堂與電視展開爭奪的時候,常是教堂失敗,電視得勝的,多麼的可憐!人在利用自由意志選擇時,寧可選擇了電視的節目,而放棄了教堂的聚會,從此便可以看出教堂前途是何等的暗淡了!
\newline
\newline
各位!究竟我們站在那裡?當我們面臨選擇的時候,我們是站在神的一邊,還是站在世界的一邊?從前以色列人在西乃的曠野,摩西一人獨自上山去領受十誡.百姓卻在山下跪拜金牛犢.神就吩咐摩西立即下山,摩西看見,便發怒將兩法版摔碎於地.他向以色列會眾發問:
\newline
\newline
「凡屬耶和華的都要到我這裡來.」按原文譯,意即:「凡站在耶和華一邊的,都要到我這裡來.」那一天利未支派的子孫站在摩西的一邊,表明他是屬耶和華的.
\newline
\newline
請問我們屬靈的態度如何?我們是否分別為聖,是否得勝地站立在神的一邊,抑仍在金錢上走多馬的路,追求世界的享受,對未信主的人漠不關心；若然,難怪神也要向我們發問,因為我們有很多的才幹都為世界用.當神呼召我們出來工作時,我們用自己的方法逃避神,用自己的方法編織葉子的裙子來遮蓋自己；然而神絕不放棄每一個屬祂的兒女,他急迫地,接二連三地向我們發問:「你信仰的立場在那裡?生活的立場在那裡?信仰與生活有否打成一片?曾否把所信的道實行出來?教會今天仍未復興你有何重責?」
\newline
\newline
神把這些問題放在我們面前,我們該好好禱告,求神首先復興我們.教會今天不能復興,正因為我們不體貼神的心意；不能慎重其事,馬馬虎虎地作基督徒,言行不一致.如果我們肯在神面前付出一點代價,為主多受一點苦,生活上有所改變,讓所聽的道實行出來；中國的教會才能重新振作,起來向外傳福音,得著聖靈的能力,又得著從天上而來的復興.
\newline
\newline


\newpage

\section{第42屆~港九培靈會~王永信~第二講}
\label{sec:1447}
\textbf{信徒應有的三種表現}
\newline
\newline
連結: \href{https://www.hkbibleconference.org/session-message/view/1447}{\texttt{https://www.hkbibleconference.org/session-message/view/1447}}
\newline
\newline
\hyperref[sec:1446]{< < < PREV SERMON < < <}
~
\hyperref[sec:toc]{[返目錄]}
~
\hyperref[sec:1448]{> > > NEXT SERMON > > >}
\newline
\newline
帖撒羅尼迦前書 2:3-2:10
\newline
\begin{longtable}{cl}
\hline
\hline
章節 & 經文 (和合本修訂版)\\
\hline
2:3 & \begin{tabularx}{0.7\textwidth}{X} 我們的勸勉不是出於錯誤，也不是出於污穢，也不是用詭詐。 \end{tabularx} \\ \\ \relax
2:4 & \begin{tabularx}{0.7\textwidth}{X} 但神既然認定我們經得起考驗，把福音託付我們，我們就照著傳講，不是要討人喜歡，而是要討那考驗我們的心的神喜歡。 \end{tabularx} \\ \\ \relax
2:5 & \begin{tabularx}{0.7\textwidth}{X} 因為我們從來沒有用過諂媚的話，這是你們知道的，也沒有藏著貪心，這是神可以作證的。 \end{tabularx} \\ \\ \relax
2:6 & \begin{tabularx}{0.7\textwidth}{X} 我們作為基督的使徒，雖然可以受人尊重，卻沒有向你們或向別人求榮耀， \end{tabularx} \\ \\ \relax
2:7 & \begin{tabularx}{0.7\textwidth}{X} 反而在你們當中心存溫柔，如同母親哺乳自己的孩子。 \end{tabularx} \\ \\ \relax
2:8 & \begin{tabularx}{0.7\textwidth}{X} 既然我們這樣愛你們，不但樂意將神的福音給你們，連自己的性命也樂意給你們，因為你們是我們所疼愛的。 \end{tabularx} \\ \\ \relax
2:9 & \begin{tabularx}{0.7\textwidth}{X} 弟兄們，你們記念我們的辛苦勞碌，晝夜做工，傳神的福音給你們，免得你們任何人受累。 \end{tabularx} \\ \\ \relax
2:10 & \begin{tabularx}{0.7\textwidth}{X} 我們對你們信主的人是何等聖潔、正直、無可指責，這有你們作證，也有神作證。 \end{tabularx} \\ \\
[1ex]
\hline
\hline
\end{longtable}
帖前二3-10是保羅寫給帖撒羅尼迦教會的信,從字裡行間,看出保羅如何愛帖撒羅尼迦教會的信徒,以及保羅如何在主裡作了他們的好榜樣,表明一個基督徒應有的三個態度:
\newline
\newline
(一)單純的愛心
\newline
\newline
很多基督徒「愛」神,因為神賜給他們許多特別的恩惠,這種「愛」不是神所喜悅的,因其中含有討價還價的成份；神給我恩惠纔愛神,若遇困苦艱難,就離開了神.許多基督徒順境亨通時滿口都是讚美的話,但當神向他有特別的要求,要他受點苦時；他便感到難過,心中懷疑為何屬神的人也要受苦,因而向神發出怨言.他們心中對神有一種錯覺,他們所要信的神,最好像聖誕老人一樣,只帶給人類恩惠與快樂,卻不要使人受苦.
\newline
\newline
神要求我們對祂的愛,是從表面進到祂裡面的,從貪圖外面的恩典而進到施恩者的裡面,不是為恩典而愛神,乃是為愛神而愛神.自己對主的愛心,純為主的愛所吸引,跟隨主也單為了主自己.宣信博士,是一位神所重用的僕人.他所寫的聖詩－主自己,正指出這個真理:「前要的是福祉,今要的是主.前我要得感覺,今要主言語.前我切慕恩賜,今要賜恩主.前我尋求醫治,今要主自己.」這正是他個人經歷到主自己的寶貴,從外面的認識要進了一步,對主自己有新的感受.我們也該有這樣的認識,當我們初信主的時候,或許是為了主的恩典而愛祂,但久而久之,我們應該從愛的外表進到愛的中心了.
\newline
\newline
你愛神究竟有多少份量呢?外表熱心的人並不一定是裡面有愛的人,若從他們所作的事中,仔細觀察,就很難看得到有多少真正愛成份了!保羅在哥林多前書第十三章中,詳談「愛」的真理,他題到:人若能說萬人的方言並天使的話語,又有先知講道之能,也明白各樣的奧祕,各樣的知識,又有全備的信能夠移山,只要在其中沒「愛」的存在,這一切都算不得甚麼.甚至自己將所有的財產賑濟窮人,就算身體焚燒,只是所作的事動機不是為了愛,這一切仍與我們無益.既然如此,我們必須反省,究竟我們在神面前一切事奉有多少是出於愛的成份?因為「愛」成了一切事奉的根基.
\newline
\newline
以前有一位神的僕人,他非常殷勤地牧養教會,所以這教會在當地成了表表者.有一天,他發了一個夢.他夢到有一位不速之客奪門而進,他一手持天平,另一手拿著許多類似化學藥品的東西,那人對他說:「你的熱心有多少?」這位牧師聽見異常快慰.他以為有許多人都想要知道他的熱心有多少?所以特派人來訪問.於是他從身上取出自己的熱心來,放在天平上,恰巧有一百磅,牧師看見,便洋洋自得,露出滿足的微笑,他以為自己有完全的熱心.這位奇異的客人,把這一百磅的熱心拿去,打成粉碎,再加上一些化學藥品,於是把它分析成七堆物質,放在牧師面前,然後嘆氣地對他說:「你想要知道你一百磅熱心內有什麼東西嗎?看－這堆東西是代表你個人的野心,佔百份之二十五,而貪圖別人的稱讚佔百份之二十一,宗教上的驕傲百份之十七,恩賜上的驕倣佔百份之十六,喜愛弄權柄佔百份之十四,這五堆合起來一共是九十三磅,其餘這兩堆就是:對神的愛心佔百份之四,對人的愛心佔百份之三.牧師,但願主拯救你.」這位牧師從夢中醒過來,他領悟到神藉夢賜給他的教訓.他在教友面前作見證:「前我求神救我脫離地獄,今則求神救我脫離自己.」
\newline
\newline
一個事奉主的人,常以為自己的工作是為神大發熱心,但可能這熱心內充滿了其他的雜質或動機.主也曾明言,到了審判那天,必有許多人對他說:我不是奉你的名傳道,奉你的名趕鬼,又行許多異能嗎?但主卻會對他們說:「我不認識你.」這些人在事奉主的動機上有了毛病.表面的事奉是不能得主的喜悅.今天我們愛主是否被主的愛所激勵,因愛主而愛主.正如約伯說:「祂雖殺我,我仍要愛祂.」(參伯六9-10)蓋恩夫人又說:「我寧可撰擇地獄的路,而不願違背神」究竟我們今天對神的愛如何?
\newline
\newline
(二)聖潔的生活
\newline
\newline
在這混亂罪惡的世代中,我們的生活必須聖潔,更像恩主.當然一個基督徒偶然也會為罪惡所勝,但聖靈居住在我們心中,隨時指責我們的不是,直至我們認罪以後,才能重獲平安與喜樂；因為一個人帶著罪惡,是絕不能有效地事奉主.在舊約的聖殿裡一切服事神的人都要聖潔,甚至那抬器皿的祭司也要聖潔.
\newline
\newline
人常易於容讓罪惡,甚至怕提及罪惡,能不犯「大罪」已感滿足了；但在神眼中看來,罪無大小之別,聖潔對於事奉神的人是必備的條件.人只看外貌,但神不像人看人,祂要察看我們的內心,因為祂看重人內心過於人手所作的工；只要我們的生活在神面前有毛病,我們所作的都算是草木禾楷,毫無價值.
\newline
\newline
登山寶訓的八福篇,從其中可看出神十分重視人裡面的問題.被稱為有福的八種人,除了憐恤人的人和使人和睦這兩種人外,其餘的六種:虛心的人,哀慟的人,溫柔的人,飢渴慕義的人,清心的人和為義受逼迫的人；他們蒙福的憑據,都是關乎自己本身的問題,並非從工作事奉中所帶來的結果.所以我們該追求聖潔,把信仰活在我們生命中；因為教會與信徒的生活是十目所視,十手所指的.比方:尖沙嘴鐘樓上所掛的大鐘,與我們手上所戴的錶,在功用上是一樣；但兩者之影響力顯然有不同.手上的錶所指的時間錯了,只影響一個人,但鐘樓的時間錯了則影響群眾；教會的言行所帶來的影響力,正如鐘樓之於群眾.
\newline
\newline
(三)殷勤的事奉
\newline
\newline
常見有人為著追求世上的物質,寧願浪費許多時間,可是在追求靈性長進的事上,一點代價也不肯付.有人幹自己的工作顯得十分殷勤,為何竟在親近神的事上顯得如此懶惰?許多信徒都犯了屬靈的懶惰,不願找時間與神交通；蔬忽自己的靈性,荒廢聖經寶貴的亮光及真理；單憑每主日所聽的道理來造就自己,試問我們如何能在基督裡長進?教會又從何得著復興?今天的教會有屬靈的大飢荒,傳道人有責,信徒也有責.信徒只盼望傳道人用他們消化後的食物來餵養自己,正如那些未長大成人,不能吃乾糧還要別人來餵養的嬰孩一樣.
\newline
\newline
在聖經裡,給我們看見,神選召一個人出來為他工作前,這個人必須在自己的崗位上忠心,絕不是躺在那裡等候人來分派他們工作.例如:摩西──正當在曠野牧羊時,神的呼召便臨到他,把拯救以色列民的重責放在他肩頭上.基甸──正在酒醡裡打麥時,耶和華的使者便向他顯現,把救民之使命託付他.當以利亞遇到以利沙的時候,以利沙正在田裡耕地,他便從田裡上來跟從了以利亞.使徒們的被召,都是在海邊撒網時遇見耶穌.還有許多的事例,都是給我們指明一件事:我們若要合乎主用,必須首先作一殷勤工作者.我們每天最少用二十分鐘靈修,讀經禱告,讓神的話進入我們心中,充滿我們整個心靈.滿有新的力量和恩典的,又讓所默想的道實行在我們生活之中.如是者,我們便發覺,我們的生命與以前顯然有別,完全被神的靈所充滿；正如活水江河一樣往外湧流,直接影響自己,又影響教會.
\newline
\newline
你盼望自己得著復興,教會也得著復興嗎?倘若你立定心志,讓單純的愛心,聖潔的生活,殷勤的事奉,充滿你整個生命.我敢說:復興的火必定臨到你,接著教會也能藉你而蒙恩.
\newline
\newline


\newpage

\section{第42屆~港九培靈會~王永信~第三講}
\label{sec:1448}
\textbf{你的兄弟在那裡?}
\newline
\newline
連結: \href{https://www.hkbibleconference.org/session-message/view/1448}{\texttt{https://www.hkbibleconference.org/session-message/view/1448}}
\newline
\newline
\hyperref[sec:1447]{< < < PREV SERMON < < <}
~
\hyperref[sec:toc]{[返目錄]}
~
\hyperref[sec:1449]{> > > NEXT SERMON > > >}
\newline
\newline
創世記 4:1-4:11
\newline
\begin{longtable}{cl}
\hline
\hline
章節 & 經文 (和合本修訂版)\\
\hline
4:1 & \begin{tabularx}{0.7\textwidth}{X} 那人和他妻子夏娃同房，夏娃就懷孕，生了該隱 ，她說：「我靠耶和華得了一個男的。」 \end{tabularx} \\ \\ \relax
4:2 & \begin{tabularx}{0.7\textwidth}{X} 她又生了該隱的弟弟亞伯。亞伯是牧羊的；該隱是耕地的。 \end{tabularx} \\ \\ \relax
4:3 & \begin{tabularx}{0.7\textwidth}{X} 過了一些日子，該隱拿地裡的出產為供物獻給耶和華； \end{tabularx} \\ \\ \relax
4:4 & \begin{tabularx}{0.7\textwidth}{X} 亞伯也把他羊群中頭生的和羊的脂肪獻上。耶和華看中了亞伯和他的供物， \end{tabularx} \\ \\ \relax
4:5 & \begin{tabularx}{0.7\textwidth}{X} 卻看不中該隱和他的供物。該隱就非常生氣，沉下臉來。 \end{tabularx} \\ \\ \relax
4:6 & \begin{tabularx}{0.7\textwidth}{X} 耶和華對該隱說：「你為甚麼生氣呢？你為甚麼沉下臉來呢？ \end{tabularx} \\ \\ \relax
4:7 & \begin{tabularx}{0.7\textwidth}{X} 你若做得對，豈不仰起頭來嗎？你若做得不對，罪就伏在門前。它想要控制你，你卻要制伏它。」 \end{tabularx} \\ \\ \relax
4:8 & \begin{tabularx}{0.7\textwidth}{X} 該隱與他弟弟亞伯說話 。 二人正在田間時，該隱起來攻擊他弟弟亞伯，把他殺了。 \end{tabularx} \\ \\ \relax
4:9 & \begin{tabularx}{0.7\textwidth}{X} 耶和華對該隱說：「你弟弟亞伯在哪裡？」他說：「我不知道！我豈是看守我弟弟的嗎？」 \end{tabularx} \\ \\ \relax
4:10 & \begin{tabularx}{0.7\textwidth}{X} 耶和華說：「你做了甚麼事呢？你弟弟血的聲音從地裡向我哀號。 \end{tabularx} \\ \\ \relax
4:11 & \begin{tabularx}{0.7\textwidth}{X} 現在你必從這地受詛咒，這地開了口，從你手裡接受你弟弟的血。 \end{tabularx} \\ \\
[1ex]
\hline
\hline
\end{longtable}
神曾向墮落以後的亞當發出第一個問題:「你在那裡?」接著祂又問犯了謀殺罪的該隱發出第二個問題:「你的兄弟在那裡?」這兩問題－前者是關係自己,後者是關係別人；前者是對內,後者是對外.
\newline
\newline
創世記第四章,記載有關人類第一宗謀殺案.在遠處有一塊青草地,其上躺著一個滿身鮮血的青年,他名叫亞伯；另外的一邊安放著一個祭壇,旁邊站著一個怒容滿面的青年,他的名叫該隱.這兩個青年是人類的第二代,他們也是同父同母的兄弟.該隱是兄,以耕種為生；亞伯是弟,以牧羊度日.有一天,他們各持不同的祭品,一同到神前獻祭；該隱所獻的是地裡的出產,亞伯卻將自己羊群中頭生的獻上.神悅納了那預表基督的羊羔,而不悅納那代表人手所作成的農產；因為前者是神所命定,後者是人為的方法.今天我們對神的奉獻與事奉也該檢討,我們曾否走了該隱的路?是否因已用錢財事奉主而感到滿足,卻疏忽了靈性上的追求?事奉神靠十架的能力,還是靠個人的力量?該隱看見神悅納了亞伯所獻的祭,心中甚不喜悅；臉色也轉變了,甚至產生恨的心理,殺人的意念,於是起來把弟弟殺死了.嫉妒成恨,可見一個人的恨心是十分兇惡可怕的.
\newline
\newline
謀殺案發生以後,神便來到犯罪的該隱面前,正如昔日亞當夏娃犯罪後的情形一樣神也同樣地對待以色列人,無論他們何時走差了路,神的手便伸出,審判便要臨到；這乃是必然的定律,因為神要向犯罪的人類顯現.今天的教會,無論何時變了質,起了紛爭,神也要出來干涉,審問.因為神並不偏待人,昔日神審判以色列民的手,今天也必然臨到眾教會.耶和華來到該隱該面前,向他發問:「你的兄弟在那裡?」這是神向犯罪的人類發出的第二個問題,與第一個問題:「你在那裡?」稍有不同.第一個問題的重心,是放在自己的身上,是向內的,其中的意思包括了:「你這個人怎樣?」「你在什麼地方?」「你裡面的生活情況如何?」但第二個問題的重心是放在別人身上,是向外的,其中表明了:「你的兄弟在什麼地方?」「他的情況如何?」你對他是否負責?別人的靈魂曾否關心?向外傳福音的負擔又怎樣?現今神也用這個嚴重的問題問教會,祂會對我們說:「我親愛的兒女,你的兄弟在那裡?你是否關心你家庭,學校,辦事處那些未得救人的靈魂?你對於港九,以及全世界人的靈魂之生死問題的態度又如何?」我們成為基督徒,個人努力追求聖潔,求靈性的進步,獨善其身是不夠的,這只完成了一半；另一半是要因著我們向外產生傳福音的負擔,顧及他人靈魂問題而得完成.因為神向今天教會的要求,不是單顧自己,保持現狀,此實違背屬靈的原則和教訓.
\newline
\newline
當主耶穌離世升天時,他曾吩咐門徒說:「但聖靈降臨在你們身上,你們就必得著的能力,並要在耶路撒冷,猶太全地和撒瑪利亞,直到地極,作我見證.」耶穌命令門徒出去傳福音,把福音向外擴展；直至全世界都有得聞福音的機會,這也是現今教會重要的異象.普遍來說中國的教會缺乏更遠的眼光,沒有顧及海外的需要,固步自封.更有不少教會依賴西方差會,在物質上的供應,而感到自滿自足；今天中國的教會信徒應一同起來,向海外傳福音.因為現今的中國人已遍佈全球,我們當趁著這末世時代,就是西方教會離棄真道,靈性走向下坡時；中個的信徒向普世傳福音,負起差派傳教士的重責.
\newline
\newline
該隱在這愛弟兄的事上作了壞榜樣,他不但沒有愛顧自己的兄弟,反而將他殺死了!難怪神要來到他面前,向他發問:「你的兄弟在那裡?」但該隱卻回答說:
\newline
\newline
(一)「......我不知道......」
\newline
\newline
「我不知道」,這話原文意義是:這事我管不著,我不知道,為何要問我.我們雖沒有明顯地向神強嘴,但我們又何嘗關心別人的靈魂?祇以為自己在傳福音的事上無關,故在態度上實與該隱無異.
\newline
\newline
(二)「......我豈是看守我的弟兄嗎?」
\newline
\newline
該隱面對面向神撒謊,昔日亞當犯了罪仍未敢在神前撒謊,只是推諉責任；現在該隱在犯罪的事上更進一步,由此可知人類到了第二代比第一代更墮落.該隱先向神撒謊要向神反問,為要堵住神的口.我們也常在自己的態度上如此頂撞神,每當教會請我們出來事奉時,我們故意推辭,以為這並非關己之責任.我們的所為正像該隱的態度,也行在該隱的路中.
\newline
\newline
「我豈是看守我的弟兄嗎?」該隱膽敢向神強嘴,簡直是傷透神的心.神創造諸世界,又在其上造了人類,目的為要人類掌權,但可惜人類從始祖亞當就開始墮落了；他的後代也照樣失敗,不能完成神的心意.直到挪亞時代人的罪惡滿盈,終於神要用洪水毀滅整個世界,卻保守了挪亞一家.但挪亞的後裔仍是犯罪,神又要在吾珥城另選出亞伯拉罕一族；盼望祂的子民能以神的心為心,但神的願望仍是落了空.直至新約時代,神差遣祂獨生的兒子耶穌基督,纔完成祂的心意.今天神又揀選了祂的教會,藉著教會為祂作見證；現今祂所關切的就是教會,能否成就祂的旨意,關懷別人的靈魔?愛顧自己的同胞?
\newline
\newline
舊約時代的該隱作了我們的警戒,把自己的弟兄殺死了；而新約時代的安得烈,則成了我們的榜樣.聖經告訴我們:當安得烈遇見耶穌以後,他先找著自己的哥哥西門彼得,對他說:我們遇見彌賽亞了,於是領彼得去見耶穌.神的心意是要我們每一個信徒,都成為一個歡喜領人歸主的安得烈,把我們兄弟,失喪的靈魂帶到神面前.
\newline
\newline
神首先要向我們發問:「你在那裡?」因為祂先要了解我們個人的情形如何?接著又向我們發問:「你的兄弟在那裡?」因為祂也願知道我們對外的情形又如何?今天中國教會應作的幾件事:
\newline
\newline
(1)培植熱心事奉主的青年.
\newline
\newline
(2)為受訓以後的青年預備機會,把他們領到事奉中.
\newline
\newline
(3)負起差派傳教士到遠方傳福音的責任.
\newline
\newline
今天有許多青年不肯奉獻傳福音,因為他們看見教會裡,傳道人們受了許多不必要的苦.教會刻薄傳道人,不肯供應他們一切生活上合理的需要.教會應該自省,負起培植青年傳教士的重責,又要鼓勵他們,差派他們出去傳福音.
\newline
\newline
在馬來亞有一所教會,雖然他們聚會的人數並不太多,但他們肯負起差派工人往外傳道的工作.開始那年只能派出兩三位,但每年他們都增加差派的數目,如是者,直至現今他們已經差派了二十多位同時間的傳道人往外地去了.現今的局勢給我們看見,西方的傳教士要退去了；中國的教會開始有力量,理應起來取而代之,訓練中國的信徒東南亞一帶傳福音.
\newline
\newline
「你的兄弟在那裡?」這是神向每一位信徒發問的一個問題,我們今天能否安然面對這問題呢?使徒保羅他能坦然無懼的對神說,我已成了弟兄的看守者,你所交託我的事情我已經成了,那美好的仗我已經打過了,當跑的路我已經跑盡了,所信的道我已經守住了.他的回答簡單而有力量,因為他已用盡自己的一生,在傳福音的事工上勞苦.
\newline
\newline
我們究竟如何對待我們的兄弟?倘若我們對他們置諸不理,漠不關心,就等於把他們殺死了.今天神要向我們每一個信徒發問:「你的兄弟在那裡?」我們究竟要走該隱的路?安得烈的路?還是保羅的路呢?
\newline
\newline


\newpage

\section{第42屆~港九培靈會~王永信~第四講}
\label{sec:1449}
\textbf{雅各屬靈上的轉機}
\newline
\newline
連結: \href{https://www.hkbibleconference.org/session-message/view/1449}{\texttt{https://www.hkbibleconference.org/session-message/view/1449}}
\newline
\newline
\hyperref[sec:1448]{< < < PREV SERMON < < <}
~
\hyperref[sec:toc]{[返目錄]}
~
\hyperref[sec:1450]{> > > NEXT SERMON > > >}
\newline
\newline
創世記 28:10-28:22
\newline
\begin{longtable}{cl}
\hline
\hline
章節 & 經文 (和合本修訂版)\\
\hline
28:10 & \begin{tabularx}{0.7\textwidth}{X} 雅各離開別是巴，往哈蘭去。 \end{tabularx} \\ \\ \relax
28:11 & \begin{tabularx}{0.7\textwidth}{X} 到了一個地方，因為已經日落，就在那裡過夜。他拾起那地方的一塊石頭枕在頭下，就躺在那地方。 \end{tabularx} \\ \\ \relax
28:12 & \begin{tabularx}{0.7\textwidth}{X} 他做夢，看哪，一個梯子立在地上，梯子的頂端直伸到天；看哪，神的使者在梯子上，上去下來。 \end{tabularx} \\ \\ \relax
28:13 & \begin{tabularx}{0.7\textwidth}{X} 看哪，耶和華站在梯子上面，說：「我是耶和華—你祖父亞伯拉罕的神，以撒的神。你現在躺臥之地，我要將它賜給你和你的後裔。 \end{tabularx} \\ \\ \relax
28:14 & \begin{tabularx}{0.7\textwidth}{X} 你的後裔必像地上的塵沙，必向東西南北開展；地上萬族必因你和你的後裔得福。 \end{tabularx} \\ \\ \relax
28:15 & \begin{tabularx}{0.7\textwidth}{X} 看哪，我必與你同在，無論你往哪裡去，我必保佑你，領你歸回這地。我總不離棄你，直到我實現了對你所說的話。」 \end{tabularx} \\ \\ \relax
28:16 & \begin{tabularx}{0.7\textwidth}{X} 雅各睡醒了，說：「耶和華真的在這裡，我竟不知道！」 \end{tabularx} \\ \\ \relax
28:17 & \begin{tabularx}{0.7\textwidth}{X} 他就懼怕，說：「這地方何等可畏！這不是別的，是神的殿，是天的門。」 \end{tabularx} \\ \\ \relax
28:18 & \begin{tabularx}{0.7\textwidth}{X} 雅各清早起來，拿起枕在頭下的石頭，立作柱子，澆油在上面。 \end{tabularx} \\ \\ \relax
28:19 & \begin{tabularx}{0.7\textwidth}{X} 他給那地方起名叫伯特利；那地方原先名叫路斯。 \end{tabularx} \\ \\ \relax
28:20 & \begin{tabularx}{0.7\textwidth}{X} 雅各許願說：「神若與我同在，在我所行的路上保佑我，給我食物吃，衣服穿， \end{tabularx} \\ \\ \relax
28:21 & \begin{tabularx}{0.7\textwidth}{X} 使我平平安安回到我父親的家，我就必以耶和華為我的神。 \end{tabularx} \\ \\ \relax
28:22 & \begin{tabularx}{0.7\textwidth}{X} 我所立為柱子的這塊石頭必作神的殿；凡你所賜給我的，我必將十分之一獻給你。」 \end{tabularx} \\ \\
[1ex]
\hline
\hline
\end{longtable}
亞伯拉罕,以撒,雅各是以色列民族的列祖.亞伯拉罕和以撒時代,只一人代表愛神；但從雅各開始以色列整個民族有了新的發展,他領導整個民族,從一人的敬拜進而成為一族的敬拜.雅各所生的十二個兒子就是以色列民的十二支派.神先得著亞伯拉罕,再藉著他得到整個以色列民族,神為要將祂旨意傳遍世界,正如神藉耶穌一人的工作得到教會,又藉著教會在這世上為祂傳福音一樣.
\newline
\newline
雅各在神面前的地位是十分重要的.倘若神要使用以色列人成為祂施的器皿,那麼他們必須接受神的訓練,所以神要從領導他們的雅各開始.
\newline
\newline
雅各從小就很有才幹,又有計謀,敢作敢為,佔便宣而從不吃虧.有一天,當他長兄以掃從回野回來時已經累昏了；雅各便乘人之危,用一碗紅豆湯把他長子的名份奪過來.這件事雙方面都有錯誤,以掃為了解決一時肚腹之快慰,竟肯將自己上好長子的名份出賣；而雅各的錯,在於對人不忠誠,缺乏真正的愛心,行事詭詐.有不少的信徒也常走了以掃及雅各的路,常見有人因貪圖一時的安逸,甚至將自己屬靈的福氣也出賣了；愛宴樂多過愛神,寧願選擇電視的節目,也不參加教會晚間的聚會.這些宴樂,電視的節目,就是我們的紅豆湯,它能使我們失去屬靈的福份.
\newline
\newline
雅各不但一次欺騙長兄,甚至與母親結連,欺騙老父以撒,領受了從神而來的祝福.其後以掃獲悉,甚為惱怒,並蓄意於父親去世後把雅各殺死.雅各聞之懼怕萬分,在母親主使下,便起來逃命往舅父拉班的家去了.雅各在父家十分舒服,現在他要離家出走了,在荒山野嶺中奔馳；四顧無人,不能再見雙親的面,在曠野遇到野獸的危險,盜賊的危險.當時環境的可怕,心靈上的自責,使雅各恩想到這是自己犯罪所帶來的後果,今天我們蒙恩的信徒,在父家能享神一切榮耀的福份；但我們犯罪偏要在無泉源,無光線,無供應的曠野中飄流,以致教會失去屬靈的權柄.正如當日的以色列人偏行己路,離棄真神,心存偶像,以致失去所羅門時代聖殿的榮耀一樣.雅各不在父家,要在曠野過飄流的生活；不是父親不要他,乃是他欺騙父兄的結果.
\newline
\newline
雅各在曠野,孤苦零丁,就在那裡神仍顧念他；並沒有把他遺忘,在他最需要人來幫助時,神在夢中向他顯現.且對他說安慰的話:「他夢見一個梯子立在地上,梯子的頭頂著天,有神的使者在梯子上,上去下來；耶和華站在梯子以上說......我也與你同在,你無論往那裡去,我必保佑你,領你歸回這地；總不離棄你,直到我成全了向你所應許的.」神絕不放棄他所愛的每一個兒女,祂的眼睛也常看顧我們的需用.雅各在離家的日子裡,神不願他灰心失意,所以便在夢中向他顯現,應許將賜與列祖的福氣臨到他身上.他醒了以後,第一個感覺就是「畏懼」.他說:「這地方何等可畏,這不是別的,乃是神的殿,天的門.」雅各從小就在敬畏神的家庭中長大,可惜他對神的認識只屬表面化,在外表上與神親近.他雖然犯了罪,神仍向他顯現,憐憫他,使他覺得這地是可畏之地.他清早起來,就給那地方起名叫伯特利,就是神的殿,天的門的意思.接著他向神許願:
\newline
\newline
(1)神若與我同在,在我所行的路上保佑我.
\newline
\newline
(2)又給我食物吃,衣服穿.
\newline
\newline
(3)使我平平安安的回到我父親的家.
\newline
\newline
倘若神答允所求,他就必以耶和華為神,凡神所賜給他的,必將十分之一獻上.從雅各的許願裡,我們發覺他是何等厲害的一個人,單在第一個條件裡就出現了三次「我」字,他不是說我要跟隨神；乃是要神來跟隨他,他在那裡神也要在那裡.雅各不但需要神來跟隨他,他還求神供應他一切屬世的物質,就是衣食住行也包括在內.雅各永遠不能脫離自己,他正代表了現今不少的工利信徒,為求自己的益處,為餅而來的群眾；有利益就以神為神,否則需要加以考慮.難怪有許多信徒,當神要他受一點苦時,他便怨天尤人,甚至從此退去.
\newline
\newline
雅各與神立約:倘若神保守他,他回來以後凡神賜給他的,他就把十分之一獻與神.換言之,神若賜他一百萬,他就給神十萬,自己還留下九十萬.雅各連向神也想佔便宜.有不少信徒懷了這種態度,若神使他富足他便向神捐獻,只把所剩餘的獻與神,這是不正確奉獻的態度.奉獻是我們信徒當盡的本份,我們富足時需要奉獻,窮貧時也需要奉獻；很多信徒在奉獻的事上佔了神的便宜,舊約的標準是獻上十分之一,新約時代都是屬神的,除了養生以外,應把一切都獻上.現今有許多傳福音的工作需要大量金錢支持,我們是否仍在地上為自己積聚財寶.
\newline
\newline
雅各對神存有工利的態度,雖然他曾一度在伯特利面對面朝見神,但可惜他頭腦仍未改變.他到舅父拉班家中居住,繼續施用手段計謀,奪了舅父的財產,得了兩個妻子,十二個兒子,一個女兒和一切肥美的羊群牛群,然後還起來逃跑.但神要對付他,從前他只會佔人便宜,現在神的手要臨到他,使他自食其果.當他們到示劍時,他的女兒受人凌辱,自己的兒子為復仇而成了殺人的兇手,以致他們再不能在那裡居住.神在使用一個人之前,必須徹底對付他,將他「老我」破碎.雅各也不能例外,所以當他預備回鄉時,就在雅博渡口,遇到一個人來和他摔跤,直到黎明；這個人來為要把雅各摔下去,但雅各的天然人性十分堅強；最後那人見自己勝不過他,就將他的大腿窩摸了一把,他的大腿就瘸了.大腿窩是人體最有力量的地方,神在那裡把他摸了一把,使他的大腿瘸了,讓他留下一個受了對付的記號.使他知道自己最強的地方,在神面前也算不得什麼.雅各受了神完全的對付,便降服下來,自此,他人生就有了很大的轉變.今天神也要在我們身上動工,因為我們還存有許多對付的罪惡,「老我」仍未在神前降服.
\newline
\newline
創世記三十五章,告訴我們:雅各再一次回到伯特利,他在那裡做了兩件事很有意義的工作:
\newline
\newline
(一)築了一座壇
\newline
\newline
雅各自從遭遇毘努伊勒的經歷後,他的人生有了很大的轉機.後來他再次來到伯特利,這時他再不先為自己尋找安身之所了,他關心到神的事,在那裡他為神築了一座壇.試問我們對神教會是否關心呢?教會不復興我們的責任又如何?我們曾否攔阻了神與我們同行?主耶穌教導我們向神禱告說:「願你的旨意行在地上,如同行在天上.」神的旨意之所以不能自由運行,乃因受了人的攔阻.信徒為求自己的享受,貪圖自己的利益.而忘記神的工作；以致講台無供應,教會沒有能力.神在雅各身上施行管教,破碎雅各,使他悔改.以前他只顧工利,牛群,羊群；但現今他不再單顧自己的事,只顧如何討神的喜悅.他為邵和華築了一座壇,在其上為耶和華獻祭,自此以後,他生命的目的就是為神工作.
\newline
\newline
(二)給那地起名叫伊勒伯特利
\newline
\newline
伊勒伯特利名義就是伯特利之神,即殿裡面的神.雅各屬靈的生命有了長進,以前他只認識有形式的殿,現今他的認識更進了一步,乃認識殿裡面的神.他從敬拜神的地方,進到認識神的本身；從認識神的外表進到神的中心.又從敬拜的形式,進到敬拜的實際.神完全得著雅各,因他整個人生有了深度的改變；他以後的人生勝過他以前的人生,以後的事奉,也勝過以前的事奉.
\newline
\newline
究道我們靈性的光景又如何?現在停留在什麼地方?是否仍在曠野?伯特利?毘努伊勒或是在伊勒伯特利?
\newline
\newline


\newpage

\section{第42屆~港九培靈會~王永信~第五講}
\label{sec:1450}
\textbf{爭戰的兩支軍隊}
\newline
\newline
連結: \href{https://www.hkbibleconference.org/session-message/view/1450}{\texttt{https://www.hkbibleconference.org/session-message/view/1450}}
\newline
\newline
\hyperref[sec:1449]{< < < PREV SERMON < < <}
~
\hyperref[sec:toc]{[返目錄]}
~
\hyperref[sec:1451]{> > > NEXT SERMON > > >}
\newline
\newline
出埃及記 17:8-17:16
\newline
\begin{longtable}{cl}
\hline
\hline
章節 & 經文 (和合本修訂版)\\
\hline
17:8 & \begin{tabularx}{0.7\textwidth}{X} 那時，亞瑪力來到利非訂，和以色列爭戰。 \end{tabularx} \\ \\ \relax
17:9 & \begin{tabularx}{0.7\textwidth}{X} 摩西對約書亞說：「你為我們選出人來，出去和亞瑪力爭戰。明天我要站在山頂上，手裡拿著神的杖。」 \end{tabularx} \\ \\ \relax
17:10 & \begin{tabularx}{0.7\textwidth}{X} 於是，約書亞照著摩西對他所說的話去做，和亞瑪力爭戰。摩西、亞倫和戶珥都上了山頂。 \end{tabularx} \\ \\ \relax
17:11 & \begin{tabularx}{0.7\textwidth}{X} 摩西何時舉手，以色列就得勝；何時垂手，亞瑪力就得勝。 \end{tabularx} \\ \\ \relax
17:12 & \begin{tabularx}{0.7\textwidth}{X} 但摩西的雙手沉重，他們就搬一塊石頭來放在他下面，他就坐在上面。亞倫與戶珥扶著他的手，一個在這邊，一個在那邊，他的手就穩住，直到日落。 \end{tabularx} \\ \\ \relax
17:13 & \begin{tabularx}{0.7\textwidth}{X} 約書亞用刀打敗了亞瑪力和他的百姓。 \end{tabularx} \\ \\ \relax
17:14 & \begin{tabularx}{0.7\textwidth}{X} 耶和華對摩西說：「你要把這事記錄在書上作紀念，又念給約書亞聽：我要把亞瑪力的名字從天下全然塗去。」 \end{tabularx} \\ \\ \relax
17:15 & \begin{tabularx}{0.7\textwidth}{X} 摩西築了一座壇，起名叫「耶和華尼西」。 \end{tabularx} \\ \\ \relax
17:16 & \begin{tabularx}{0.7\textwidth}{X} 他說：「我指著耶和華的寶座發誓，耶和華必世世代代和亞瑪力爭戰。」 \end{tabularx} \\ \\
[1ex]
\hline
\hline
\end{longtable}
摩西是以色列人偉大的領袖,負起領他們出埃及進迦南的重責.當以色列人出埃及行經曠野時,曾遇到不少的難處.他們在利非訂安營,有亞瑪力人與他們爭戰.這是他們在曠野遭遇的第一次爭戰,以後還要經過無數次的爭戰,才能到達迦南.舊約的以色列人代表了新約的信徒,埃及是預表世界.今天神要把教會從世界中領出來,分別為聖,其中的過程也經歷了不少次的爭戰；以色列人的爭戰有勝有敗,教會也是如此.昔日以色列人的際遇如何?可以成為今天教會的寫照.
\newline
\newline
以色列人與亞瑪力人的爭戰,是他們爭戰的開始,他們未曾受過任何軍事上的訓練；又有婦孺,牛群,羊群隨行.但對方卻是精勇的戰士,曾受嚴格訓練,又有充足的武器和裝備,按屬世的眼光來看,這次爭戰以色列人必然失敗,警慌逃跑的；但透過神的眼目,事情並非如此悲觀,因為爭戰的原則,不在乎依靠人的努力,乃在乎依靠耶和華.當日他們靠著神的恩典,不慌不忙地出去應戰了.今天教會爭戰的態度也該如此,雖然在世上我們只不過是極少數的一群人,真不曉得如何與世爭戰；但神爭戰從不看人數之多少,乃是看我們站在那一邊?自己?世界?或是神的一邊?倘若我們站在神那邊,祂必然使我們勝而又勝.
\newline
\newline
而亞瑪力人的軍隊來了,摩西吩咐約書亞從眾民中選出人來,出去和亞瑪力人爭戰,不是每一個以色列人都可以出去,這支軍隊是經過挑選而組成的.約書亞要用他屬靈的眼光,從眾民中選出那些有信心的,肯付代價的,願意離開家庭的,不貪生怕死的,要組一支耶和華的軍隊.這也成了我們事奉神的原則,我們曾否為神爭戰甚至一切都放棄；否則我們臨陣便退縮,成事不足而敗事有餘.魔鬼從不怕普通的信徒,因為這些人不能起什麼作用；牠只怕十字架的精兵,不貪生怕死的.倘若教會有這類肯為主付上生命代價的信徒,這教會必然興旺.
\newline
\newline
約書亞把組成這隊軍隊的隊員選好以後,他便一馬當先,帶領以色列人出發,直搗戰場.正當那時,另有一支軍隊也同時出發,這支軍隊只有三個隊員,就是摩西,亞倫,戶珥所組成的.他們沒有到前線衝鋒陷陣與亞瑪力人面對面交戰,他們卻上了山頂,到另外一個禱告的戰場.在那裡,摩西舉起雙手,為以色列人在山下的爭戰而禱告,摩西的禱告顯得十分逼切,因為以色列人要與精銳的亞瑪力人爭戰.整個爭戰是與他們整個民族的前途與生命有關,千鈞一髮,勝敗就目前.摩西懇切的禱告,使我們聯想到主耶穌,祂在客西馬尼園中曾逼切為我們禱告.摩西的禱告只關乎以色列人,但耶穌的禱告是關乎全人類.
\newline
\newline
戰爭開始了.約書亞帶領眾民英勇作戰,雖然當時有千萬的仇敵,但他們毫不駭怕,只憑著信心與神一同爭戰.摩西不斷在山頂上為他們的戰爭禱告,摩西何時舉起雙手；以色列人就得勝,何時垂手亞瑪力人就得勝.因此摩西的禱告便成為他們爭戰勝敗的關鍵,可見禱告在爭戰上的重要性.因在以色列人裡面,有一位肯付禱告代價的摩西,便使全民族有爭戰的力量.摩西禱告的手也會疲倦發沉的,亞倫和戶珥就成了他的幫手.他們把石頭搬來,放在摩西的手以下,亞倫與戶珥扶著他的手,一個在這邊,一個在那邊.他的手就穩住了,一直到日落的時候,約書亞用刀殺了亞瑪力王和他的百姓.這是何等激烈的一幅圖畫,摩西,亞倫,戶珥三人都是以色列民的首領；在這次戰爭中,他們全家合作,無分彼此,因此以色列人那天便獲得完全的勝利.
\newline
\newline
從這件事實裡,神要現今的會在其中學習三樣要緊的教訓:
\newline
\newline
(一)領袖的合作
\newline
\newline
教會的領袖們在處理教會的事務時,若充份表現出主內的合作精神,神必定賜福.當以色列人與亞瑪力人爭戰時,在山上作戰的豈不是有三人嗎?亞倫既是神所選立的大祭司,當然可以代表百姓到神前禱告.但當他看見摩西開始向神禱告,而神又悅納他所祈求；藉著他的禱告,以色列民就獲大勝；他便寧可作摩西的助手,而不在這事奉上爭為首.而摩西也沒有因神垂聽禱告,而顯出驕傲之心,因聖經記載他比眾人謙和.戶珥也沒有因著自己成為首領,要在山上負起禱告的爭戰,而自以為了不起.他們三個以色列民的首領完全在神的事工上合一,彼此合作不爭為首.神的教會應在事奉和禱告上同心合意,不凡心鬥角也不存自私之心,神的恩典必然臨到.
\newline
\newline
(二)兩代的合作
\newline
\newline
在這爭戰的兩支軍隊中,帶領以色列的首領們在山上展開禱告爭戰的是摩西,而帶領以色列民在山下與亞瑪力人爭戰的是約書亞.摩西是年長的,而約書亞是年少的,他們這年老和年少的兩代,在這爭戰中也充份表現他們的合作.一個到山上去,一個到平地去；各盡其才,約書亞年少力強,所以當摩西派他在眾民中,選出人來出去與亞瑪力人爭戰時,他毫無異議,對首領的權柄完全順服,他沒有因摩西離開危難的戰場,到山頂上安逸而不肯順服.摩西也給我們成年人作了好榜樣,他在爭戰時上了山頂禱告,不是為了貪生怕死,乃是有聖靈特別的帶領.因他年事已高,不宜在戰場中追殺,但他屬靈的經歷比年輕人深；可以為以色列民及約書亞代求,成為他們的後盾.
\newline
\newline
現今各地教會有一潮流,就是青年與成年分家,兩代不能彼此合作,青年人寧願另立團契而不到教會聚集.倘若今天的教會得不到青年,明天的教會就沒有指望了.因為青年是將來教會的柱石.在這種潮流裡,雙方面也該反省；青年人重理想,每當理想不能達到時便產生反抗；他們看見成年人的敗壞,言行不一致,便不願到教會來,更不肯順服教會的權柄.教會也應自省,因教會有很多方面是需要改善,純正的信仰當然不能改；但工作的方式,傳福音的方法,是應該更新而變化的.教會是神的家,在一個完整的家庭裡,需要有祖父母,父母親,兒女,孫子,有老有少；否則成了不平衡的現象,老年人愛顧年青的,又作他們的榜樣,少年人彼此順服合作,家庭才會有快樂.教會裡年長的,可以在事奉上負起領導的責任,年輕的彼此合作,一同興旺福音,如此神就會使我們在爭戰上有大大的勝利.
\newline
\newline
(三)全體的合作
\newline
\newline
當亞瑪力人來與以色列爭戰時,全體以色列人一同起來爭戰,沒有一人發怨言,也沒有一人逃跑；萬人一心,在那裡為爭戰而努力,所以他們便得了全體的勝利.爭戰以後,摩西築了一座壇,起名叫耶和華尼西,意思就是耶和華是我的旌旗.這次戰爭的勝利,帶給他們整個民族無限的光榮及歡欣,這一切都是基於他們全體合作帶來的成果.
\newline
\newline
末世的教會,也有無數的亞瑪力人來圍繞,他們來為要與教會爭戰,這就是包括了一切不信的勢力,反對神的集團.教會應該為復興的爭戰而努力,不是個人的力量,乃是全體總動員,只靠幾個人的力量實在獨力難支.教會裡每一個徒應各盡其才,運用各人不同的恩賜,或在禱告上,講道上,領人歸主,捐獻,探訪,種種事奉上.信徒除了主日崇拜的聚會以外,我們還該有其他事奉的機會,神賜給我們各人的才幹,就正如活水江河,用之不歇,直湧到永生.如果不肯工作,就會閒暇,活水的江河很快便會變成死水.但願信徒們起來多為主工作,因為我們越多工作,神就越賜繁盛的恩典.
\newline
\newline


\newpage

\section{第42屆~港九培靈會~王永信~第六講}
\label{sec:1451}
\textbf{為萬軍之耶和華重建聖殿}
\newline
\newline
連結: \href{https://www.hkbibleconference.org/session-message/view/1451}{\texttt{https://www.hkbibleconference.org/session-message/view/1451}}
\newline
\newline
\hyperref[sec:1450]{< < < PREV SERMON < < <}
~
\hyperref[sec:toc]{[返目錄]}
~
\hyperref[sec:1452]{> > > NEXT SERMON > > >}
\newline
\newline
以斯拉記 1:1-1:3
\newline
\begin{longtable}{cl}
\hline
\hline
章節 & 經文 (和合本修訂版)\\
\hline
1:1 & \begin{tabularx}{0.7\textwidth}{X} 波斯王居魯士元年，耶和華為要應驗藉耶利米的口所說的話，就激發波斯王居魯士的心，使他下詔書通告全國，說： \end{tabularx} \\ \\ \relax
1:2 & \begin{tabularx}{0.7\textwidth}{X} 「波斯王居魯士如此說：耶和華天上的神已將地上萬國賜給我，又委派我在猶大的耶路撒冷為他建造殿宇。 \end{tabularx} \\ \\ \relax
1:3 & \begin{tabularx}{0.7\textwidth}{X} 你們中間凡作他子民的，可以上猶大的耶路撒冷去，重建耶和華－以色列神的殿，他是在耶路撒冷的神；願神與這人同在。 \end{tabularx} \\ \\
[1ex]
\hline
\hline
\end{longtable}
以色列人是神所揀選的民族,神對他們有特別的愛情；然而他們離棄真神,以外邦偶像為神,因此神的審判與刑罰便臨到他們身上.神使用外邦人成為管教以色列的刑杖,讓他們被擄到外邦地,作異族的奴隸,受他們轄制,單從他們所寫的詩歌中,便可知道他們心中痛苦的境況:「我們曾在巴比倫的河邊坐下,一追想錫安就哭了；我們把琴掛在那裡的柳樹上,因為在那裡擄掠我們的,要我們唱歌......我們怎能在外邦唱耶和華的歌呢?......」(詩一百三十七篇)當以色列人在異地受苦時,他們就追想到昔日耶路撒冷聖殿的偉大,大衛所羅門時代的榮美,現在因他們離棄神,便受神管教的刑罰,他們要在外邦地為奴七十年之久,直到期滿了他們才能回到聖地－自己的家鄉.以色列人的遭遇,正可代表了今天教會的境況；什麼時候教會離開了神,向神不忠,就會被撒旦的權勢所擄,變成世俗化,在那裡教會就要受許多的苦.
\newline
\newline
以色列人先後經過七十年飄流的生活後,神便藉外邦王古列的手,准許他們返回猶大國；當時帶領子民回去的領袖是所羅巴伯.他帶領以色列民重回故土時,心裡第一件預備要做的事,就是將已毀壞的耶路撒冷聖殿重建.因為聖殿經過歷年來爭戰的破壞,劫掠的損害,其中一切的用具都變成支離破碎.但聖殿是神的家,它的榮耀就成了以色列民族的榮耀；它的羞辱也成了民族的羞辱,所以子民歸回時都懷著一顆熱情,盡力把聖殿重建.
\newline
\newline
被擄的民族回到猶大地,各歸各族找回自己的家譜,因為家譜在猶大人的觀念中是十分重要的.他們每一個人都有自己的家譜,可以追溯到他們的祖先亞伯拉罕的一代.但在被擄歸回的民中,聖經告訴我們有三家的人,他們不能指明自己的宗族譜系,證明他們是否以色列人.他們就是第來雅的子孫,多比雅的子孫,尼哥大的子孫.這三家人在族譜之中,尋查自己的譜系卻尋不著；因此算為不潔,不准供祭司的職任.可憐的三家族人,他們外表似乎是以色列人,滲雜在以色列民中,現在回到聖地了,但他們卻沒有找到自己的族譜；不能證明自己是真以色列人,因此不能參與他們建殿的工作.今天有不少人只是傳統式的到禮拜堂聚集,宗教對他們只不過形式化的節目；這種人十分危險,到末世審判時,主會對他們說:我不認識你們.因為他們未曾真正重生得救,生命冊上沒有他們的名字,以致他們沒有屬靈的家譜.我們今天需要自省,究竟我們是否真正重生得救?我們的名字有沒有記在生命冊上?我們是否被主所認識的?否則我們的光景正如那沒有族譜的三家以色列人一樣.
\newline
\newline
千萬的以色列人聚集在耶路撒冷,預備動工重建聖殿,雖然他們人數眾多,但工作起來卻似一人.彼此同心合意,沒有紛爭,不分宗派,不分長幼；努力在那裡事奉神,建造神的殿.但魔鬼看見便紅了眼睛,牠最駭的就是人同心向神；牠就必定要在旁邊找機會作破壞的工作.魔鬼殷勤不懶惰,牠常在那裡忍耐等候機會來攔阻神的工作.正當他們建殿時,旁邊的人便起來攻擊他們,使他們的手發軟,又向他們說一些假意的話:「請容我們與你們一同建造,因為我們尋求你們的神,與你們一樣......」從外表看過去,這些人的心意真不容易拒絕,因為魔鬼也喜歡裝作光明的天使；把一些糖衣包著的毒藥給屬神的人吃,使吃了的人中毒受傷甚至死亡.魔鬼的詭計雖然非常厲害,但牠對教會所施的方法,不外乎兩種:一是施壓力控制,逼迫教會；另一是暗地工作,使教會的信仰變質.使信徒靈性冷淡,以假亂真,混雜神的道.但我們必須要持守神的真道,在信仰上有根有基,否則對於道理之真假便難於分辨.
\newline
\newline
昔日那些人假意來親善他們,希望參與建殿的工作.他們用心不善,有意混雜其中作破壞份子,倘若以色列人沒有分辨善惡的智慧；容讓這班人在一起,如此建造的工作真不堪設想了.幸而當時的首領所羅巴伯有屬靈的眼光,他看出信與不信原不相配,不能在一起工作；因而拒絕了這個不速之客,也是除去一個工作上的危機.今天的教會也有不少類似的情形發生,我們常因某人有金錢,地位,才幹,權勢,便沒有顧及他信仰的問題；把他們拉到教會中負重責,這樣有如教會進行慢性自殺,這些人來了,教會的危機便產生,從此在教會中掌權的,不再是神而是人了.在教會中事奉神的弟兄姊妹,他們必須要有充足的信心,在教會內外有好名聲,又被聖靈充滿,如此,教會才能真正讓神有工作的機會.
\newline
\newline
這班有意破壞神工作的人遭受拒絕後,仍不甘心,他們又控訴於亞達薛西王.王不明白真相,於是傳旨意,強迫他們停工.各民在這壓迫下,對於建造聖殿的熱心,便漸漸冷卻,在敵人阻止停工的情況中,各民便各歸各城去了.當初他們是懷著一顆熱心,為耶和華神的殿工作,但當遇到難處便向後退了.教會的情形也是如此,教會中的事奉從來經不起考驗,試驗來了便向惡勢力低頭,一切的勇敢便不翼而飛.但神絕不能容讓這種冷淡的情況繼續往下去,神眷顧祂自己的百姓,為他們焦急.於是神就動工,興起先知哈該及撒迦利亞,成為祂的出口；對以色列民說勸勉的話及責備的話,向他們指明雖然罪惡的勢力是大,但神仍掌王權.所以無論遇到任何困難的環境,不能妥協,應該繼續為主奮勇前進,至終也必能勝利,因為神被稱為得勝的神.
\newline
\newline
先知哈該替神向以色列人傳言,責備又勸勉這班軟弱的子民,應當起來為建造耶和華聖殿的工作而努力,因為當時的以色列民竟說:「建造耶和華殿的時候尚未來到.」這句話是妥協,軟弱,沒有信心的話,百姓受到仇敵的攪擾,不能再繼續工作,以為建造殿宇的時候仍未來到；他們被惡劣的環境所征服,就把神的工作耽擱了.所以神差遣先知對他們說:「這殿仍然荒涼,你們自己還住天花皮的房屋麼?」子民放下建殿的工作,各歸各城,為自己建造精美的房屋,他們不但建造,而且把它修飾得十分雅觀,裝有天花板.(古時的猶太人不是每一所房屋都可以裝有天花板.)他們寧可花錢財,時間,在建造自己的房屋,而不多想及神的事,使神的殿仍然荒涼.我們也常把屬世的應酬,個人的享受看得十分重要,卻把神的事情耽延了,神的家荒涼無人關心,我們也毫無負擔,良心可算是麻木了.先知有見及此,便用幾個比喻來提醒他們,責備他們:
\newline
\newline
(1)「吃不得飽,喝卻不得足」──這指到屬靈上的飢餓,完全得不到靈糧的供應,於是產生了屬靈的飢荒,使到自己受痛苦.
\newline
\newline
(2)「穿衣服卻不得暖」──靈性的溫度是冰冷的,沒有溫暖,熱心,與活力,因為在裡面毫無聖靈的工作,只是形式化的宗教.
\newline
\newline
(3)「得工錢的,將工錢裝在破漏的囊中」──袋中有了破洞,錢便從洞裡漏出來.這代表教會發生了破洞,神的恩典便不能繼續存留,因著這些漏洞,神給我們許多教訓也是保留不住.
\newline
\newline
先知哈該用責備的話叫子民自我省察,曾否為著建造自己天花板的房屋；便忽略了神的工作,使屬神的工作受了枕延.先知接著又用勸勉的話,引領他們走回正路,終於神的殿再一次動工而且完成了,耶路撒冷便得到一次大的復興.我們也該如此省察,求神給我們教會有新的轉機,重新被建造起來,除去其中一切漏洞,成為火熱的家.
\newline
\newline


\newpage

\section{第42屆~港九培靈會~王永信~第七講}
\label{sec:1452}
\textbf{重建聖城}
\newline
\newline
連結: \href{https://www.hkbibleconference.org/session-message/view/1452}{\texttt{https://www.hkbibleconference.org/session-message/view/1452}}
\newline
\newline
\hyperref[sec:1451]{< < < PREV SERMON < < <}
~
\hyperref[sec:toc]{[返目錄]}
~
\hyperref[sec:1453]{> > > NEXT SERMON > > >}
\newline
\newline
尼希米記 1:1-1:4
\newline
\begin{longtable}{cl}
\hline
\hline
章節 & 經文 (和合本修訂版)\\
\hline
1:1 & \begin{tabularx}{0.7\textwidth}{X} 哈迦利亞的兒子尼希米的言語如下：亞達薛西王二十年基斯流月，我在書珊城堡中。 \end{tabularx} \\ \\ \relax
1:2 & \begin{tabularx}{0.7\textwidth}{X} 那時，我有一個兄弟哈拿尼，同幾個人從猶大來。我問他們那些被擄歸回、剩下殘存的猶太人和耶路撒冷的情況。 \end{tabularx} \\ \\ \relax
1:3 & \begin{tabularx}{0.7\textwidth}{X} 他們對我說：「那些被擄歸回剩下的餘民在猶大省那裡遭大難，受凌辱；耶路撒冷的城牆被拆毀，城門被火焚燒。」 \end{tabularx} \\ \\ \relax
1:4 & \begin{tabularx}{0.7\textwidth}{X} 我聽見這話，就坐下哭泣，悲哀幾日，在天上的神面前禁食祈禱， \end{tabularx} \\ \\
[1ex]
\hline
\hline
\end{longtable}
為神重建聖殿的工作,在極度艱難與工作中完成了.但以色列的聖城仍是破爛荒蕪,也需要他們合作重建.城在聖殿以外,聖殿在城以內；聖殿先重建,聖城便繼續完成；尼希米被神使用,成為一個建城的使者.
\newline
\newline
尼希米在書珊城的宮中,在亞達薛西王面前當酒政須得到王極度的信任.而王他每天有好幾次能面對面見王,比朝中一切文武官員的機會都多；因此他在王前必有一切美好的享受,又為王所寵愛.不過尼希米並未用著自己有安逸的享受,任大官之職,便將自己國家民族遺忘了.有人把耶路撒冷遭大能,受凌辱,城牆拆毀,城門被火焚燒的情形告訴尼希米；他聽見這消息,就坐下哭泣,悲哀了好幾日.向神禁食禱告,心中為著國民受痛苦,城被毀而極其難過；他沒有苟且偷安,身雖在異鄉,心仍切想神的殿神的城.今日我們中國的僑胞散居世界各地,我們在這享有信仰自由,物質豐富的香港裡；曾否也想念到那些落在水深火熱當中的同胞?更以他們的痛苦為苦?有如置身其間一樣.還是自顧自己,苟且偷安?
\newline
\newline
有一天尼希米面帶愁容,遞酒奉給王,王察覺了.王問他為何面帶愁容?尼希米於是禱告神,求神幫助他,也感動他,使他說出合宜的話來.尼希米雖站在王前,心卻想念神；口雖對王說話,心卻向神禱告.神便把當講的話藉著他,向亞達薛西王說明.他在王面前口吐真情,說明自己國家荒涼,聖城被火焚燒,人民受苦；雖身居皇宮,然心卻異常迫切.我們的生活會否像尼希米一樣,時刻與神有交通,常過著禱告的生活,心中時刻想念屬天的事.我們應該培養自己時刻過著禱告的生活,因為這生活能使我們常與神在一起,又隨時得到祂能力的幫助.因著尼希米默禱天上的神,神就把當說的話賜給他,王聽了又感動了他的心.亞達薛西王便傳旨意下來,下令讓尼希米回國,而且還派了馬兵護送他,又另外賜給他重建城牆所用的一切器材.尼希米就與百姓,順利回國.
\newline
\newline
尼希米回到耶路撒冷,當晚他不驚動仇敵,夜間起來查察聖城；察看城門被火焚燒,城牆被拆毀.後來他就召了猶大的領袖來,鼓勵他們為重建耶路撒冷聖城齊心努力.因為他們光有聖殿是不夠的,他們還需要有城牆.因為聖城有兩種特別的功用:
\newline
\newline
(1)守衛功用
\newline
\newline
城牆是用石頭砌成,是一件相當穩固及高大的工程.敵人若要毀壞聖殿,必先要毀壞城牆；因此城牆之於聖殿便有守衛的作用了.今天教會也需要有屬靈的城牆負起保衛的功用；我們為了保衛教會純正的信仰,聖潔,不容認罪惡入侵,就必須築上保衛的城牆.
\newline
\newline
(2)分別功用
\newline
\newline
沒有城牆的存在,就沒有範圍的界限,也不能有內外的分別；城嬙就可以把國家分成內外兩部份,在城牆的一邊是屬城外,另一邊是屬城內.信徒及教會的生活也該如此,與世界應有分別為聖的界限.這個分別不拘於外表衣服或是髮型,乃裡面的分別,在世而不屬世.活出屬天的樣式,不被罪惡所玷污,過一聖潔的生活；以基督的心為心,接受祂愛的生命,自然我們便與世人有所不同,這就是聖城存在的作用了.
\newline
\newline
尼希米對子民說:「來吧!我們重建耶路撒冷的城牆.」神與尼希米同工,所以當他們呼喚眾民時,他們對於這個呼喚有所響應,一同起來建造,於是重建聖城的工作又開始動工.尼希米是一位甚懂心理學的人,他領導眾民建城是另有方法,他規定:「眾祭司各在自己的房屋對面修造.」(尼三28)如此的建造,就能達成幾個目的.
\newline
\newline
(1)每一家建造自己房屋對面的城牆,如此,各人必然盡心去做,因為這段城牆面對自己,直接影響自己的安危；若建造時隨便敷衍,敵人便乘虛進攻,自己便首當其衝了.
\newline
\newline
(2)建造運動是以「家」成為單位.如此,每一家都有自己應盡的本份,就是修造自己房屋對面的城牆；合家老幼一起工作,在建造期間,沒有一個是閒人,或是躲藏起來不作工的.
\newline
\newline
(3)把建造城牆的工作,分與每一個家庭修造,倘若其中有一段城牆毀壞了；就算其他各部份十分堅固,也是失去城牆的效用.因為只需有一個漏洞,就足以使整座城都受到威脅,因此彼此間也是關心,合力一同建造.
\newline
\newline
教會的建造也該效法聖城的建造,教會不應有太多閒雜人等,每一份子都有自己事奉的崗位,好使教會沒有缺口,若有一個細小的漏洞,魔鬼就能乘虛而入,影響整個教會.例如教會有某一位信徒犯了一件罪,他仍繼續犯；牧師知道不去規勸他,也沒有其他人指責他的罪.那麼其他的信徒就相繼效尤了!又如果牧師三鍼其口,那麼,信徒就更任意妄為了.一個細小漏洞,足以使教會帶來無法收拾的局面,這都是魔鬼破壞教會工作的方法之一.
\newline
\newline
他們重建耶路撒冷聖城的工作,是相當艱鉅的,環境的惡劣足以使他們灰心喪膽了.但他們仍在尼希米領導下盡心竭力,而且更們發現到有兩種值得我們效法的人:
\newline
\newline
(1)一手工作一手拿兵器
\newline
\newline
他們一方面作工,一方面爭戰.因為有許多敵人乘虛而入,要逼使他們停工,所以他們要常作準備,手持兵器.信徒的生活就是爭戰式的生活,事奉也是爭戰式的事奉.每當我們為主工作時,我們都會感到有另一種力量反抗我們,又有聲音困擾我們.我們越多事奉主,魔鬼便越要引誘我們,把我們從一切事奉當中拖下來,所以我們該隨時隨地準備與牠爭戰.
\newline
\newline
(2)吹角的人
\newline
\newline
這些吹角的人,他們在那裡儆醒等候,因為他們成了守望者,當心看守營地,看看有沒有人入侵!那些建造聖城的人固然是重要,他們從早到夜,從天亮直到星宿出現,不斷地工作.但這些吹角的人也是不可缺少的,每當他們看到敵人入侵時,便向眾民提出警告,使各人準備作戰,教會也是需要這種守望的吹角者,就是在靈裡儆醒的人；他們留心觀察教會將遇到的危機,又提醒各人努力防備.
\newline
\newline
尼希米帶領子民努力建造,不分日夜,甚至他們睡覺休息的時候,也不脫衣服,為的是常作準備,隨時工作準備爭戰；他們不貪圖自己個人的舒適,殷勤地盡心服事主.尼希米也在他們當中成了美好的領導者,各方面為他們留下好榜樣,他沒有為自己存留薪俸,也沒有顧到自己的事情,又不求自己的好處,努力工作,因此眾民願意全力以赴.雖然他們沒有好的工具,只靠一雙手,就在極度惡劣的情形下；祇用短短的五十二天時間,便把城牆修造完畢,因為神施恩的手幫助他們.工作的成就不在乎人力量的多少,乃在乎人是否完全順服神.
\newline
\newline
倘若我們今天肯順服神,以神的事為念,常常儆醒,隨時準備,教會也能在短期內得到很大的改變,有新的轉機,更有復興的事實.教會之所以得不到復興,並非神不願意,乃是我們不肯為著復興教會付上代價,但願神憐憫我們.
\newline
\newline


\newpage

\section{第42屆~港九培靈會~王永信~第八講}
\label{sec:1453}
\textbf{欠債}
\newline
\newline
連結: \href{https://www.hkbibleconference.org/session-message/view/1453}{\texttt{https://www.hkbibleconference.org/session-message/view/1453}}
\newline
\newline
\hyperref[sec:1452]{< < < PREV SERMON < < <}
~
\hyperref[sec:toc]{[返目錄]}
~
\hyperref[sec:1460]{> > > NEXT SERMON > > >}
\newline
\newline
羅馬書 1:14-1:16
\newline
\begin{longtable}{cl}
\hline
\hline
章節 & 經文 (和合本修訂版)\\
\hline
1:14 & \begin{tabularx}{0.7\textwidth}{X} 無論是希臘人、未開化的人、聰明人、愚拙人，我都欠他們的債， \end{tabularx} \\ \\ \relax
1:15 & \begin{tabularx}{0.7\textwidth}{X} 所以願意盡我的力量把福音也傳給你們在羅馬的人。 \end{tabularx} \\ \\ \relax
1:16 & \begin{tabularx}{0.7\textwidth}{X} 我不以福音為恥；這福音本是神的大能，要救一切相信的，先是猶太人，後是希臘人。 \end{tabularx} \\ \\
[1ex]
\hline
\hline
\end{longtable}
教會是神的家,所以每一個信徒都該關心自己家裡的事.我們對於建立教會,護衛教會,甚至奮興教會都有重大責任.各人當存火熱的心事奉神,否則便欠了神的債.
\newline
\newline
什麼力量能激發一個人願意放棄一己的工作,起來終身事奉神呢?這個人在他為主工作之前,他必須從心裡先產生一顆愛神愛人的「熱心」,眼睛又看到這傳福音的需要,有千萬人在罪中滅亡的「異象」；然後他整個人才能產生出傳福音的「負擔」；最後便有所「行動」毅然撇下一切跟隨主.舉例來說:一小孩在馬路玩耍,剛巧有一輛汽車從橫街衝出,孩子不及走避,司機也不及煞掣,就在這千鈞一髮,極度危險時有人挺身而出,把這小孩搶救過來.或許你對這小孩子從未認識,而你竟肯救他,這正因為你看見這孩子處境的可怕,然後才產生出拯救他的負擔來.
\newline
\newline
保羅在羅馬書第一章指出:「無論是希利尼人,化外人,聰明人,愚拙人,我都欠他們的債.」保羅明言無論是誰,他都欠了他們的債,也就是普世的債.欠債,換言之就必定有債主存在；被人討債是十分難過的事,心中時刻有一種壓力的存在,覺得自己欠了別人的東西.保羅正活在這種欠債的負擔裡,以致使他終夜為著失喪者的靈魂禱告,又獻身努力在傳福音的事上,把福音傳給普世的人.一個人生存在世,便有機會嘗試欠別人債的滋味；父母欠兒女們的債,兒女又欠父母的債；夫妻之間也彼此欠債,朋友之間也欠了義氣和愛心的債.倘若一個人毫無欠債的感覺,他就是失去了責任感,他的頭腦也許不健全.例如一個小孩子,他毫不感到自己欠了別人什麼東西,他只會覺得別人在虧欠了他因為他還未懂得人性,不知道對待父母朋友要盡孝,盡忠.及至他長大成人,他纔開始感到自己在各方面都欠了別人的債.從這事例來看,我們可以曉得:一個人欠債的感覺,是與個人成熟的程度成正比的；一個人越成熟越感到自己欠了別人的債,因為越長成他的責任也越來越多.
\newline
\newline
基督徒的生命也是如此,一個初信信徒正如一個嬰孩,他不會有欠債的感覺.但當他的靈性長大成人以後,他裡面欠債的感覺則越發加增；以前從沒有想到要向別人傳福音,現今因著屬靈的生命成熟了,於是對失喪人的靈魂的負擔也開始擴展.欠債的感覺越來越重,欠債的範圍也越來越廣；不但關心自己的家人,朋友,甚至對那些從不認識的人,也顯出關懷的態度.深感他們流離失所,如同羊沒有牧人,又如一班坐在死蔭裡失喪的人,倘若我不向他們傳福音便有禍了.
\newline
\newline
保羅給我們作了很好的榜樣,當他蒙主呼召出去傳道時,可謂已盡竭力,甚至生命也不顧.雖然如此,他竟然感到自己有如一個大罪人,欠了別人許多福音的債；不但欠了他周圍的人,甚至欠了普世人的債.相反地,有不少基督徒,靈性低沉,他對於事奉的工作從來不幹,也沒有想過在傳福音的工作上有責.正因著兩者靈性深度不同,所以彼此的感覺與負擔也是不同.很可惜,現今有不少信徒,只會作聽道的人,從不把所聽聞的真理實行在生活當中,難怪心中毫無傳福音的負擔.因他心裡的活水江河,已經成了死水；只曉得領受恩典,而不會把恩典與別人分享.這樣的人永遠不能得到神的喜悅和榮耀.
\newline
\newline
教會的工作,可以用「運動會」來作個比方,在每一個運動會舉行之際,我們不外發覺在整個會場當中,只有四種不同的人存在:
\newline
\newline
(1)運動員──這些健兒祇參加運動會,各項節目的比賽,他們不做其他的工作；專心致志,用全副精神在運動場上賽跑,比賽.他們在比賽過程中,必須忘記其他一切的東西.這種人就是在教會裡用全時間事奉神的傳道人,牧師.他們為神的工作盡忠,專心以祈禱傳道為事.
\newline
\newline
(2)事務人員──他們是負責管理整個會場的雜務,例如:租借地方,管理交通,和擔任各項運動比賽的評判.這種人在整個運動會裡也是十分重要的；單有運動員,而沒有其他事務人員,整個會場就處理得不妥善.他們可比作在教會作執事,幹事的人.雖然不是全時間事奉,但他們卻協助教會各項工作的進行,與傳道人一同在教會中同心努力,使教會的事工得以興盛.
\newline
\newline
(3)啦啦隊──這班人在運動比賽進行時搖旗吶喊,負起鼓勵的責任；為賽跑中的運動員「加油」,運動員疲倦走不動時,為他們鼓舞,使他們重新得力.這班人在教會裡,專作守望,鼓勵和代禱告的工作.因著他們流淚的禱告,以致許多疲倦中全時間事奉神的傳道人們,再次得著屬天的氣力,繼續奔跑前面的路.
\newline
\newline
(4)觀眾──他們只是作旁觀者,隨便觀望,他們絕不是賽跑的人,也不是事務人員,或是工作人員.他們對於整個運動會毫無責任可言.倘若教會裡大多數的信徒,都是觀眾,只知享受神恩,而絕不願意為主作半點工作.而只有幾個人在運動場上拚命,這是不能起什麼大作用的!教會裡每一個信徒必須有事奉的崗位,沒有一個是閒人或是觀眾,神才能恩待教會,賜下力量和恩典,神復興的火也才臨到地面上的教會.
\newline
\newline
神在這末世指望教會的,就是教會能起來盡心事奉祂,使祂的旨意早日成就.神的旨意不能自由運行,無非因為教會的耽延,每當神呼召人出來事奉祂的時候,常遭受到人的逃避,所以祂的旨意便難於成就.普通來說,神的呼召分為兩大類:
\newline
\newline
(1)包括每一位重生得救的信徒
\newline
\newline
自從我們信主那天開始,這呼召便臨到我們身上,就是神盼望祂的兒女能在神所賜予的崗位上服事神,為福音而作見證.無論是工程師,醫生,教員,學生,家庭主婦......都應該在自己的工作崗位上盡忠；否則我們便違背了神的呼召,欠了福音的債.
\newline
\newline
(2)特別是蒙召少數的一群人
\newline
\newline
神在所有的信徒中,特別選召一班少數的人,使他們沒有其他的顧慮；用全時間來為主傳福音,又為福音的事工而努力.
\newline
\newline
我們必須把這兩種呼召弄清楚,神對於每位信徒都有第一種呼召.但我們教會傳福音的工作,也需要有第二種呼召的人出來負責.今天教會需要大量佈道,奮興,培靈的人才,文字工作人才,牧養教會的人才.因為教會的圈子裡現有的人才質和量都不夠,教會人才的缺乏；不是因為沒有人才,乃是許多人才寧可為世用,為己用,而不肯為主用.倘若今天的教會仍然靜坐不動,不多幾時教會的工作,就完全被人群所忘記；故此教會應有更多奮勇的戰士起來善用神所賜予的才幹,努力為主作工.把傳道的範圍擴大,帶到人群當中；不再閉關自守,這樣教會才會受到廣大群眾的注視,帶領更多的人歸向基督.
\newline
\newline
保羅深深感到自己欠了全世界人類福音的債,他不以福音為恥,情願盡力將福音傳至普天下.請問我們心中的感受如何呢?是否對於傳福音毫無負擔,願主感動我們,能樂意為主付上時間,金錢與力量,償還我們福音的債,與人同得福音的好處.
\newline
\newline


\newpage

\section{第42屆~港九培靈會~王永信~第九講}
\label{sec:1460}
\textbf{教會面前的兩條路}
\newline
\newline
連結: \href{https://www.hkbibleconference.org/session-message/view/1460}{\texttt{https://www.hkbibleconference.org/session-message/view/1460}}
\newline
\newline
\hyperref[sec:1453]{< < < PREV SERMON < < <}
~
\hyperref[sec:toc]{[返目錄]}
~
\hyperref[sec:1454]{> > > NEXT SERMON > > >}
\newline
\newline
「教會」的名義在舊約時代已經開始.自從亞當夏娃犯罪墮落以後,神特別選出亞伯拉罕一族,又藉著他一人得了以撒,雅各及以色列人的十二支派,進而得著普世.新約時代的情形亦相仿,神藉著祂兒子耶穌基督得到教會,為畏把祂旨意傳遍普天下.舊約時代的以色列民正代表了今天的教會,昔日神如何使用以色列人,今日也要使用教會.起初的亞當失敗了不能成就神旨,所以到了新約神就差耶穌到世上來,成為末後的亞當.首先的亞當的失敗,要在末後亞當身上成就；末後的亞當是神所喜悅的,藉著祂產生了教會.神指望新約教會能代表祂,因神一切工作的範圍不出乎教會,所以歷世歷代教會的工作是十分重要.為此,教會必須順服神,否則神的旨意因著人的弱軟便受到耽延.教會什麼時候軟弱了,發生問題,神的手便伸出,要在其中加以安慰,鼓勵,教導,藉著祂的僕人向教會發出警告的挑戰.舉些例子來說:
\newline
\newline
約書亞記二十四章:神藉祂的僕人約書亞,向以色列人發出一個挑戰,因為他們將要面臨一個相當惡劣的形勢.約書亞曾領導他們奮勇交戰,得迦南美地為業.現在約書亞年事已高,離世日子快到,不能繼續領導他們.約書亞想知道當他離世以後,他們能否站立得住；是否仍以耶和華為神,正如摩西,約書亞猶在之時一樣.現今的世代,神也興起祂的僕人,復興教會,信徒得到復興後,教會的事工就得以擴展.但當主的僕人離開以後,教會的情況又漸漸轉弱,冷淡之風再度侵入教會,異端四處興起.昔日以色列人就是面臨這種危機,因此神要藉約書亞向以色列人說話,將兩條路擺在他們面前.讓他們決定到底事奉耶和華,還是事奉假神?約書亞又在眾民面前表明己的立場:「至於我和我家我們必定事奉耶和華.」他要成為眾民的榜樣,對神至死忠心,但以色列民又如何?究竟他們選那一個?
\newline
\newline
出埃及記三十二章:摩西帶領以色列民至西乃曠野,他上山領受法板,百姓卻在山下拜金牛犢.摩西下山發怒摔碎法板,對百姓說:「凡屬耶和華的,都要到我這裡來.」按原文意即「凡站在耶和華那邊的都要到我這裡來.」那天利未支派的子孫都到摩西那裡.營內百姓只有兩條路的選擇,到底要選擇耶和華還是金牛犢?選擇真神還是假神?選擇迦南地還是埃及地?
\newline
\newline
列王記上十八章:迦密山上以色列人同樣面臨兩條路的選擇,以利亞時代,以色列離棄耶和華犯罪拜巴力,以異邦的神巴力亞舍拉為神.先知以利亞在迦密山向他們挑戰,警告以色列民,究竟他們心持兩意要到幾時?若耶和華是神,眾民就該歸順他；若巴力是神,就跟從巴力.就在那一天,以利亞用神蹟挽回以色列人的心,使他們伏在地上大聲喊叫:「耶和華是神,耶和華是神.」
\newline
\newline
約翰福音二十一章:提比哩亞的海邊,主耶穌又將這個激發性的挑戰放在彼得面前,主三次向他發問:「你愛我比這些更深嗎?」不錯,彼得跟隨主已有三年多,他曾認識耶穌是基督,是神的兒子；親眼看過主行了許多神蹟奇事,但現今主快要離開他了；他需要起來取代主耶穌傳福音的職事.就在這重大的關頭裡,主耶穌不能不讓彼得有個更清楚的選擇,到底他愛主是否比這些更多?
\newline
\newline
神也盼望知道我們對神的愛如何?我們口裡說愛神,這個愛到底到那一個程度呢?現今是一個非常的時代,一切事物都在急速的轉變中,神要向眾教會挑戰,「到底你選擇人的路,還是神的路?跟隨祂還是跟隨世界?愛神還是愛自己?現今的社會除非真正愛神,否則沒有一個教會能站立得住.因為魔鬼已加速進行混亂的工作,教會的信仰變了質；靈性方面大破產,對罪惡的判斷貶了值,黑白是非不能分明.教會的名雖在,但實際的生活已失去聖經的準繩!除非教會有一個新的轉機,否則前途真不堪設想.究竟我們選擇那一條?人的道路?抑是神的道路呢?
\newline
\newline
自從神造人以後,神就賜給人有自由的意志,有自由選擇的權利.從伊甸園開始,人便有了第一個選擇,但可惜始祖在第一個選擇裡,選了人的路；以至全人類墮落在罪中,後來神又賜給人律法,人對於律法可以選擇遵守或是違背；在恩典時代裡,人又可以用自由意志選擇救恩,或是神的呼召；也有權利可以拒絕.由此可見人的選擇是如何重要,決定得「永生」或是「永死」,決定成為神的僕人,或是在罪中墮落.
\newline
\newline
放在教會面前的兩條路,教會必須有所選擇,不能缺乏明確的態度,得過且過,否則教會就會漸漸被世界所吞噬.究竟教會要選擇那一條路?神的路還是人的路?神的路是窄路,需付代價,有愛心,討神喜悅,保守信仰,屬靈的合一,普世的佈道.但人的路是寬路,走的時候是容易的,得過且過,討人喜悅,放鬆信仰,人神合一,只顧自己.兩條路完全不同,有很大的區別；教會若要得著復興必須走窄路,肯付代價；而神的僕人又殷勤不懶惰,以神為首,這是唯一復興的門徑.教會單有外表的建築物是不夠的,我們心須把神接到我們中間來；使神的愛在信徒心裡燃燒起來,從外表的敬拜而進到敬拜的中心.
\newline
\newline
教會常常軟弱得不到神的復興,不是神不願意賜福,乃是我們教會本身發生了問題,我們必須要注意:
\newline
\newline
(1)選擇持守真道
\newline
\newline
現今教會裡有一個潮流,冠以美名稱之為「教會合一.」這實在是魔鬼迷惑教會的方法,教會合一運動裡面竟把天主教,希臘正教,不信派,拉進信仰純正的圈子裡.簡直是混亂真道,混雜神的教會,這運動普世漫延；許多不明真相的教會也被拉到這運動中.何謂「合一」?聖經告訴我們的原則是:一主,一信,一神；信仰相同,異象相同的就在主裡合一；互相幫助,不分界限,宗教,教會,在主內同工,同禱,這纔是真正合一的真諦.可惜如今信仰相同的反不合一,互相勾心鬥角,「該合的不合,不該合的倒合了.」我們要在這混亂的局勢中必須選擇神的路,持守真道,在屬靈上合一.
\newline
\newline
(2)選擇熱心傳福音
\newline
\newline
在傳福音的事上不能「保持現狀」,否則不進則退.今天人口數字不斷膨脹,教會保持不動便是落伍.究竟教會在傳福音的事上該選擇那一條路?是否感到救人靈魂有責?還是閉門自守,自滿自足?
\newline
\newline
(3)選擇討神喜悅
\newline
\newline
討神喜悅和討人喜悅,兩者不可兼得時,我們究竟選擇那一樣?中國人歡喜講面子,所以在屬靈方面也帶有人情,這是十分錯誤的.我們寧可得罪人,卻不能違反真理,要讓主在我們生命中居首位.
\newline
\newline
神將教會面前的兩條路放在我們眼前,究竟我們要選擇那一條?是否選擇以弗所教會的路?事奉神而沒有愛心；還是選擇別迦摩教會的路,失去純正的信仰.選擇推雅推拉教會的路,與世俗聯合,還是選擇老底嘉教會的路,自滿自足,不冷不熱,失去向普世傳福音的熱心.
\newline
\newline
約書亞要以色列民選擇,而且還要他們「今天」便撰擇,立刻對神有所表示.現今神把這兩條路也放在我們面前,看我們選擇那一條?
\newline
\newline


\newpage

\section{第42屆~港九培靈會~王永信~第十講}
\label{sec:1454}
\textbf{不徹底的教會}
\newline
\newline
連結: \href{https://www.hkbibleconference.org/session-message/view/1454}{\texttt{https://www.hkbibleconference.org/session-message/view/1454}}
\newline
\newline
\hyperref[sec:1460]{< < < PREV SERMON < < <}
~
\hyperref[sec:toc]{[返目錄]}
~
\hyperref[sec:1455]{> > > NEXT SERMON > > >}
\newline
\newline
以弗所書 2:4-2:7
\newline
\begin{longtable}{cl}
\hline
\hline
章節 & 經文 (和合本修訂版)\\
\hline
2:4 & \begin{tabularx}{0.7\textwidth}{X} 然而，神有豐富的憐憫，因著他愛我們的大愛， \end{tabularx} \\ \\ \relax
2:5 & \begin{tabularx}{0.7\textwidth}{X} 竟在我們因過犯而死了的時候，使我們與基督一同活過來—可見你們得救是本乎恩— \end{tabularx} \\ \\ \relax
2:6 & \begin{tabularx}{0.7\textwidth}{X} 他又使我們在基督耶穌裡與他一同復活，一同坐在天上， \end{tabularx} \\ \\ \relax
2:7 & \begin{tabularx}{0.7\textwidth}{X} 為要把他極豐富的恩典，就是他在基督耶穌裡向我們所施的恩慈，顯明給後來的世代。 \end{tabularx} \\ \\
[1ex]
\hline
\hline
\end{longtable}
教會的爭戰,有不徹底的表現,主因有二,就是不順服與小信.以彼得為例,當他看見主在海面上行走時,他也想嘗試；於是照主的話到海面上行走.可惜一看環境,因怕風浪,就沉下海裡去.主責備他小信,因他的信心不徹底,虎頭蛇尾,有始無終.
\newline
\newline
我們的信心不但要有始,也當成終,因我們的主足為我們信心創始成終的.所以要以主耶穌基督為根基.
\newline
\newline
在約書亞將離世時,曾囑咐百姓要繼續爭戰；因為還有許多未得之地,就是那流奶與蜜之地.他們在約書亞帶領之下,已有多次的戰蹟勝而又勝；但還有許多未得之地,需要他們全力以赴,把住在迦南地的異族完全趕出去.教會在世上,時刻都充滿屬靈的爭戰；沒有一個人例外,從我們蒙恩得救那天開始,便加入這爭戰的行列；在爭戰中不是得勝就是失敗,不能有中立的態度.我們戰役的結果,神都有記錄在天上.當末日審判的時候,案卷要在神面前展開,請問你能否坦然無懼見主?這爭戰的記錄便成為重要的關鍵.
\newline
\newline
士師記第一章,記載了以色列人在迦南地對爭戰各種的態度,這記錄可以作為今天教會爭戰的寫照和鑑戒.
\newline
\newline
(一)猶大的爭戰──不徹底式
\newline
\newline
以色列人的領袖約書亞死了以後,他們便起來求問神,誰應先去與迦南人爭戰,神吩咐猶大支派應先上去.猶大支派便出去奮勇交戰,他們開始是相當成功的；神也幫他們戰勝了許多敵人,他們又把山地的居民趕出去,這種表現是十分可嘉的.但聖經記載說:「耶和華與猶大同在,猶大就趕出山地的居民,只是不能趕出平原的居民,因為他們有鐵車.」(士一19)他們的爭戰只完成了一半,趕出山地的居民,而不能趕出平原的居民.因為平原的居民有鐵車,所以他們便懼怕,不敢把那平原的居民趕出.猶大支派的軍隊就在敵人鐵車前失去了信心,以致不能完成神的使命,打了一場不徹底的仗.
\newline
\newline
在這末世時代裡,圍繞著教會四周的,都是可怕的鐵車.它們在恐嚇我們,使我們喪失勇氣,失去信心.這些鐵車包含了世界勢力,反對神的運動,魔鬼的力量,經濟的壓迫,和一切的難處.當教會看見這些可怕的鐵車時,便裹足不前,因而使神的工作受到攔阻.有人曾研究,彼拉多為何在猶太人面前屈服,他指出彼拉多最大的毛病就是膽怯恐懼；他明知耶穌是一個義人,妻子也曾相勸,但他懼怕猶太人的勢力,便在這惡勢力的鐵車下屈服了.在教會裡也常看見這種情形,每當一個有權有勢的人做了某些錯事,看見的人也不敢直言指責；這無非是懼怕他們那些無形的鐵車,他們寧願使教會受虧損,受辱,而對惡勢力妥協,卻不敢向那些鐵車挑戰.
\newline
\newline
昔日猶大支派的人,倘若繼續憑信心在平原上爭戰,深信神必能幫助他們；使他們把平原的居民及鐵車一同趕出去.但他看鐵車比神還大,竟忘了神是大而可畏的；以為神不能得勝鐵車,因此在鐵車前失敗.我們為什麼軟弱和失敗,恐怕在我們心裡,也充滿了無數的鐵車,以致成為我們靈性上的威脅.須知道我們得救所獲得的恩典是白白的,但得救以後所過的爭戰的生活,就需要付出代價了.進攻迦南需要付代價,得迦南為業需要付更大的代價!要經過爭戰,流血,痛苦,纔能獲得完整的迦南美地.
\newline
\newline
聖經告訴我們:「當我們死在過犯中的時候,神既有豐富的憐憫,因祂愛我們的大愛,便叫我們與基督活過來；他又叫我們與基督一同復活,一同坐在天上.」(弗二4-6),這是何等確實的應許,就是教會要成為神的子民,得享屬天的地位,但很可惜許多信徒只知道有這屬天的地位,從未享受實際屬天的生活.雖己蒙恩成為神的兒女,「心」是聖靈的殿,但其中仍有許多罪惡不能趕出,正如昔日的以色列人一樣.他們雖有得迦南地為業的應許,但仍未佔領迦南地；敵人仍未完全剪除.雖進了迦南,卻未完全得迦南,更未能完全享受流奶與蜜之迦南.因為他們不徹底的爭戰,不完整的事奉；有始無終的信心,只趕出一半的敵人,他們在爭戰上又吃了一次敗仗.
\newline
\newline
(二)便雅憫的爭戰──與世相和
\newline
\newline
「便雅憫人沒有趕出,住在耶路撒冷的耶布斯人；耶布斯人仍在耶路撒冷,與便雅憫人同住,直到今天.」(士一21)聖經這裡又指出便雅憫人在爭戰上的軟弱.便雅憫人所得到的地方是相當重要的,他們所居住的地方包括了耶路撒冷；就是他們的聖城,但便雅憫竟然沒有把住在耶路撒冷的耶布斯人趕出,容讓耶布斯人與他們一起居住.舊約時代神曾多次警告以色列人,命令他們要把耶布斯人完全的趕出,可知耶布斯人敗壞的度已到極處,所以神纔多次吩咐把他們毀滅.以色列應當與他們勢不兩立,合力把他們趕出迦南地,可是他們竟能容忍,與世相和,違背神的命令.
\newline
\newline
今天的教會也與他們相仿教會該把異族從聖潔之所趕出去,這些異族就是代表了世界的罪惡.神要我們嚴厲對付罪惡,把罪惡趕出去!但教會卻容認罪惡,講台不敢提到罪,恐怕得罪人,人就不再來做禮拜.教會的經濟,就沒有人支持了；這種形勢是逼使教會與罪惡妥協,容讓罪惡.這種態度就等於便雅憫人與耶布斯人,和平共處的態度一樣.甚麼時候教會容認罪惡,與罪惡和平共存,面對罪惡而妥協,甚麼時候便給了魔鬼有立足之地.現今魔鬼的勢力,正在整個教會中蔓延；權勢也越來越大,教會因此而墮落了!被魔鬼所控制,從此,聖靈的力量便離開了教會.
\newline
\newline
可憐的便雅憫人,不把耶布斯人從他們中間趕出去!這些耶布斯人就是淫亂的人種,拜偶像的人種,這些孽種正可以玷污以色列人,影響他們整個民族.當便雅憫人與耶布斯人和平共存後,耶布斯人便成了便雅憫人的羅網；他們受痛苦,正因為他們不徹底執行神的命令,趕出耶布斯人.
\newline
\newline
(三)但人的爭戰──反被轄制
\newline
\newline
「亞摩利人強迫但人住在山地,不容他們下到平原.」(士一34)猶大支派的爭戰只作了一半,便雅憫人支派爭戰卻與敵方和平共處,而但支派的情況更加可憐.他們不但不能把敵人趕出,又不能彼此容忍；他們要受到亞摩利人的轄制,只許他們住在山上受苦,不能到平原中居住.
\newline
\newline
神的兒女也常受到外邦人的轄制,世界勢力的擾亂,不能在暗世成為明燈照燿,把生命的道表明出來.我們不能只因個人信仰純正,而感到自豪滿足,因神向我們的要求,除此以外還需要我們在生活上聖潔,得勝而有能力.很多信徒一生只追求自己的愛好,為自己打算,從不想到該為神的工作而努力,在主面前站立起來,成神一個合神心意的人.
\newline
\newline
今天的教會是否因為四周的鐵車滿佈,而感到失望懼怕?使我們喪失信心而不敢勇往直前呢?我們在神面前的生活如何?曾否靠主得勝一切?還是繼續追求自己的愛好?我們到了迦南,是否已得迦南為業?在屬天的位置,曾否嘗到與基督一同活過來,又一同坐在天上?有那些實際屬天的經歷?猶大的爭戰是有始無終,便雅憫的爭戰是與罪惡妥協；但支的爭戰是受世界轄制,遭受完全的失敗.今天你的爭戰又如何?
\newline
\newline


\newpage

\section{第42屆~港九培靈會~王永信~第十一講}
\label{sec:1455}
\textbf{中立式的教會}
\newline
\newline
連結: \href{https://www.hkbibleconference.org/session-message/view/1455}{\texttt{https://www.hkbibleconference.org/session-message/view/1455}}
\newline
\newline
\hyperref[sec:1454]{< < < PREV SERMON < < <}
~
\hyperref[sec:toc]{[返目錄]}
~
\hyperref[sec:1456]{> > > NEXT SERMON > > >}
\newline
\newline
士師時代可稱為以色列人的黑暗時代,他們的生活正像一個惡性的循環.當他們在居於迦南地後,久而久之,便開始行耶和華眼中看為惡的事,離棄真神.神便讓懲罰臨到他們身上,藉著外邦人的手苦待他們,就在極度痛苦之中,他們又想到真神,重新悔改向神呼求.神的憐憫終於又臨他們身上,特為他們興起士師,拯救他們脫離仇敵的手.但當神再度賜福他們,使他們重新過著安居的生活後,他們又再次離開神.如此,便成了一個惡性的循環.神多次的拯救他們,但可惜以色列民又多次的墮落離棄神.
\newline
\newline
不少信徒的生活也是如此,普遍來說:信徒靈性冷淡都是由於屬世亨通,世上享受太多所致.當神的打擊管教臨到時,他們又哀求神；稍為歸向神,神便施恩,因此靈裡再得到一次的復興.但當情勢好轉以後,他們又再離開神追求世界,正如在惡性循環中一樣.許多信徒不肯到教會,除非等到教會被關起來,不能隨便做禮拜時,纔會感到教會的可貴；到不能自由傳福音時,就會領悟到傳福音的可貴.等到聖經被人焚燒不得再閱讀時,纔會覺得勝經是神的話,是神賜給人的恩物.信徒為何不能早一點了解這個真理,在平安穩妥的日子中熱心事奉神?
\newline
\newline
以色列人又行了耶和華眼中看為惡的事,因此耶和華把他們交在迦南王耶賓的手中二十年之久,在外邦人手下受苦,在那裡已經流了許多眼淚,受了許多痛苦.當他們呼求時,神為他們興起一位女士師底波拉,拯救他們脫離仇敵的手.他們的仇敵是相當勇猛的,迦南人的首領耶賓手下有一名大將,他名叫西西拉,這人善征慣戰,並且他們又有大量的兵器,九百輛的鐵車.這個形勢足能使以色列人膽怯後退,但底波拉及其助手西拉,剛強壯膽,靠著神的恩典和信心,照著神的話去行,帶兵出去打仗.聖經告訴我們:「耶和華使西西拉和他一切車輛全軍潰亂.......」雖然仇敵如此勇猛,只要他們肯依靠神的幫助,他們的爭戰終必勝利.
\newline
\newline
我們事奉神的態度也當如此,不管魔鬼的勢力如何浩大,難處如何眾多；但教會絕不能屈服.因為教會的頭是耶穌基督,祂不能失敗,因他是完全的勝利者.教會必須緊緊跟隨這位元帥,絕不可分開,有如身體與頭相連在一起一樣.如此我們可以從主那裡得享一切的恩典,在戰爭裡靠主得勝且有餘.多少時候身體不能得勝,並不是因為頭的緣故,乃因身體發生了毛病；元帥在前頭走,可惜教會卻留在後頭.
\newline
\newline
以色列人打了一場完全的勝仗,雖然面對兇猛的仇敵,九百輛的鐵車,按人的眼光看來,真是銳不可擋,但是他們都要在以色列軍隊面前失敗.因為底波拉信靠神,跟隨神前往打仗,他的助手巴拉又跟隨底波拉,以色列民跟隨著他們兩位首領,因著這一連串的跟隨,以色列人跟隨神,神便使他們勝利.
\newline
\newline
他們在這次戰爭裡,認識耶和華是他們唯一的拯救者,又學習到全體一致,跟隨耶和華爭戰的功課.所以當爭戰完畢,底波拉便作詩歌,一方面讚美神恩,另一方面指責以色列人從前的不是.他們之所以受苦,乃因離棄耶和華,為自己選擇新神,因此爭戰的事就要臨到城門口.在屬靈上的原則,也必如此,什麼時候人另選新神而離棄真神,神便要藉著患難來管教他們.昔日的美國曾以基督教立國,所以明顯地給我們看見,神在一切的事上特別賜福,又使他們成為一強盛的國家.但現今美國人開始離棄神,不追求神,許多的教會牧師起來宣傳一些不純正福音,離開了列祖的信仰；以學問,金錢,武器及其他的物件為「新神」.他們今日仍然強盛,可算是神的特別的恩典；倘若他們仍不肯悔改,神就要把恩典收回,強盛之國也要滅亡!
\newline
\newline
世人所選擇的新神,是以金錢為神.世人敬拜財神,心中所思想的,盡都是屬世的；他們關心一切物質的東西,遠超過真神.因此生活充滿痛苦,毫無平安與喜樂；人在甚麼時候有了新神,就不能得到神的賜福.神的審判也要臨到世人,直至人肯向神降服,認罪悔改,否則就不能脫離神的手.
\newline
\newline
底波拉在他所作的詩歌中說:「耶和華的使者說:應當咒詛米羅斯,大大咒詛其中的居民；因為他們不來幫助耶和華,不來幫助耶和華攻擊勇士.」(士五 23)舊約裡常看見的「耶和華的使者」不是一個普通的天使,這名詞可以與神自己互相並用；正統的神學家相信,舊約時代耶和華的使者,正是新約時代道成肉身以前的耶穌基督.祂是三位一體真神的第二位,與無始無終的神同等.這次耶和華使者來,吩咐以色列人起來咒詛米羅斯,和那地的民.關於「米羅斯」的資料,在聖經的記載並不太多,但米羅斯在這裡受過咒詛以後,再沒有看見這名詞在聖經中出現.因為他們被神咒詛後,這地方的人都要被神完全消滅.
\newline
\newline
神不容易咒詛一個人,除非這是十分嚴重的事,米羅斯的居民沒有犯大罪,他們的罪狀是:「因為他們不來幫助耶和華,不來幫助耶和華攻擊勇士.」就是當底波拉及巴拉四處呼召沒有響應,沒有幫助耶和華出去打仗.米羅斯沒有犯行動上的罪,但他們的罪狀正是犯了「沒有行動的罪」；在開始爭戰時,神需要所有的人一起來與仇敵交戰,為神勞苦,流血犧牲.但米羅斯的人卻無動於中,在神需要人有所行動時,米羅斯的居民卻享受自己安靜的生活,而不肯出來幫助耶和華.這是他們犯了無可赦免的罪,他們是標準中立式教會的代表,雖然他們沒有犯外表的罪,卻犯了按兵不動的罪.
\newline
\newline
現今神極需要大量的工人,為福音而爭戰,搶救人的靈魂.倘若信徒們對這種呼聲充耳不聞,在生活上仍追求個人的快樂,按兵不動,這樣的人就是走了米羅斯的路,要受到神的咒詛,因為他們沒有出來幫助耶和華.耶和華也需要人幫助的嗎?祂豈不是那無所不能的一位,祂可以打發十二營天使下來打仗,或是在天上用力量消滅仇敵.但神有一個原則,就是祂不歡喜單獨作戰,祂願意人與祂合作,透過教會而爭戰.若教會不肯為主作戰,神旨便受到攔阻,雖然我們沒有站在魔鬼那邊,也沒有走世路,但只要我們站在中間路線；不積極走天路,我們便會受到神的咒詛.正如老底嘉教會的處境,不冷不熱,主要從口中將他吐棄,又將他的燈台挪移.從此他再不見生命之道,再不能稱為教會,只不過是一個宗教上的組織而已.
\newline
\newline
昔日主曾斥責耶路撒冷城,路旁的一棵無花果樹,使它立刻枯乾.這棵無花果樹不是蟲蛀,乃是需要它結果子的時候,而找不到果子；所以它便受到主的咒詛,枯乾而死亡.今天的教會也是在這危急存亡之際,一切新派的信仰,對聖經懷疑的態度直接影響教會.在神的眼看來,這些教會是沒有價值,沒有用的,是個半死的教會.
\newline
\newline
教會情況的好與壞,就要看我們生活的態度如何?米羅斯的居民受到咒詛,從此不被聖經所記念.因為他們在神需要幫助時,竟然按兵不動；他們雖然沒有殺人,但他們不救人就等於殺人,此乃米羅斯犯罪之嚴重性.但願神恩待今天的教會,不再走米羅斯的路,就在神審判我們之前,切實悔改；響應神救人的呼聲,參加搶救靈魂行列的隊伍.
\newline
\newline


\newpage

\section{第42屆~港九培靈會~王永信~第十二講}
\label{sec:1456}
\textbf{失去能力的教會}
\newline
\newline
連結: \href{https://www.hkbibleconference.org/session-message/view/1456}{\texttt{https://www.hkbibleconference.org/session-message/view/1456}}
\newline
\newline
\hyperref[sec:1455]{< < < PREV SERMON < < <}
~
\hyperref[sec:toc]{[返目錄]}
~
\hyperref[sec:1457]{> > > NEXT SERMON > > >}
\newline
\newline
士師記 13:1-13:5
\newline
\begin{longtable}{cl}
\hline
\hline
章節 & 經文 (和合本修訂版)\\
\hline
13:1 & \begin{tabularx}{0.7\textwidth}{X} 以色列人又行耶和華眼中看為惡的事，耶和華將他們交在非利士人手中四十年。 \end{tabularx} \\ \\ \relax
13:2 & \begin{tabularx}{0.7\textwidth}{X} 那時，有一個但支派的瑣拉人，名叫瑪挪亞。他的妻子不懷孕，不生育。 \end{tabularx} \\ \\ \relax
13:3 & \begin{tabularx}{0.7\textwidth}{X} 耶和華的使者向那婦人顯現，對她說：「看哪，以前你不懷孕，不生育，如今你必懷孕生一個兒子。 \end{tabularx} \\ \\ \relax
13:4 & \begin{tabularx}{0.7\textwidth}{X} 現在你要謹慎，清酒烈酒都不可喝，任何不潔之物都不可吃， \end{tabularx} \\ \\ \relax
13:5 & \begin{tabularx}{0.7\textwidth}{X} 看哪，你必懷孕，生一個兒子。不可用剃刀剃他的頭，因為這孩子一出母胎就歸給神作拿細耳人。他必開始拯救以色列脫離非利士人的手。」 \end{tabularx} \\ \\
[1ex]
\hline
\hline
\end{longtable}
教會若失去了能力,燈臺被挪開了；從外表看來,雖仍是一所教會,但裡面就失去屬靈的供應.從此再沒有聖靈的力量,活潑的生命,和完全的愛.雖仍有宗教的活動,而只不過是徒有外表,卻失去了教會的實際.
\newline
\newline
以色列人又行耶和華眼中看為惡的事,所以神把他們交在兇殘的非利士人手中,讓外邦人管轄他們達四十年之久.四十這數目在聖經中常用,它的意義是「完全」,正代表了他們受了完全的轄制與審判.當神刑罰他們時,他們又開始想到神,神的恩典終於又臨到他們身上.神差遣使者向瑪挪亞顯現,預言他將要得一個兒子,這孩子要分別為聖,成為拿細耳人,專心事奉神.在舊約時代拿細耳人是一種特別的人,他們許了特別的願,又遵守好些規則,不用剃頭刀剃頭髮,清酒濃酒都不喝.因為喝酒代表放縱,身為一個拿細耳人,他的生命是受神的管制,一生要走在神的路中.他裡面有管制,外面有分別為聖的記號.我們新約每一個信徒都是拿細耳人,在我們身上應當受到聖靈的管制,走十字架的道路；外面應當有成聖與世分別的記號,基督的生命就自然在我們身上流露了.這不是個人的造作,乃是被十字架所破碎；受神對付的記號,所謂有諸內而形諸外.
\newline
\newline
參孫在神的應許中誕生了,他父母按照耶和華所說的,把他分別為聖,按照拿細耳人的規矩把他養大.正當那時,以色列人受盡非利士人的壓迫,他們極其痛苦,他們理當起來反抗,靠神的大能拯教自己脫離捆綁；但他們卻失去這種勇氣,安於現狀,甘願為奴.參孫知道自己身為拿細耳人,神在他身上有特別的使命,所以他不甘願自己民族受人欺凌,於是立定心意,要拯救以色列人脫離非利士人的轄制.參孫既有這樣的負擔,以色列人應一同起來殺敵,但可惜當參孫出去和非利士人打仗時,竟有三千猶大人聯同來到參孫面前,埋怨地說:「非利士人轄制我們,你不知道嗎?你向我們行的是甚麼呢?......我們下來是要捆綁你,將你交在非利士人手中.......」神差遣參孫作士師為要拯救他們,他們竟把參孫交在非利士人手中,置他於死地.耶和華的靈大大感動參孫,他臂的繩就好像火燒的麻一樣,他的綁繩都從他手上脫落下來,他又拿起一塊未乾的驢腮骨,擊殺了一千非利士人.
\newline
\newline
現今教會也常落在這光景,受盡世界的轄制,變成世俗化,失去一切反抗的力量；可惜教會卻滿於現狀,安於不冷不熱的情況中,從不感到難過.在這種軟弱不堪的情況下,要求神復興,神施恩的手便臨到教會.可惜復興還未來到,人倒懼怕；雙手把這一切祝福放走.中國前著名佈道家──宋尚節博士,他指責罪惡太厲害,信徒因此得復興；但教會卻把神的僕人捆綁起來,攔阻他工作,到處攻擊他,甚至把他交在仇敵的手裡.我們的主耶穌,他也是當日猶太人的拯救者,但那些猶太人卻把他交在仇敵手中,置祂於死地.教會現今的光景也是如此.
\newline
\newline
在士師時代,參孫的際遇與其他的士師有所不同.其他的士師興起時,以色列人便一同起來與他並肩作戰；但參孫自始至終都是獨行俠,赤手空拳地爭戰,沒有任何人支持與幫助.參孫另外一個與眾不同的特點,就是他有神奇的力氣,能以一敵千.這一個大能的勇士,倘若他能自始至終順服神的帶領,不放鬆自己,必能成為偉大的拯救者.但可惜參孫失敗了,聖經用一句令人非常難過的話來形容他:「耶和華已經離開了他!」他本來被聖靈大大感動,甚至用七條未乾的青繩捆綁著他,他都能弄斷；但現在耶和華的能力離開他了,從此,在他只留下了一個外殼,毫無用處了!
\newline
\newline
參孫是一個大能勇士,因著他的心意墮落,不聖潔,放鬆自己；眼睛便開始犯罪,甚至身體各部份也跟著犯罪；在裡面「放鬆」的人,外面也必定「放鬆」,他整個人在不知不覺中就陷在罪裡.任何一個人墮落都是從心裡開始,然後身體各部份纔相繼.底馬之所以離開保羅,不是一朝一夕的事,他能成為保羅的同工實不容易；但因為他的心離開了神,眼裡看見世界的誘惑,心思便發生傾慕,整個人便跟著去了.保羅很悲嘆地說:「底馬貪愛現今的世界,往帖撒羅尼迦去了.......」
\newline
\newline
身為拿細耳人絕不能沾染任何不潔,倘若心裡不潔,外表也會不潔；因著裡面沒有管制,外面的記號也改變了.裡面是好,外面也必須是好；裡面是壞,外面也必然是壞.參孫首先犯罪,後來頭髮也被剪了.因為他裡面失去聖靈的能力,外面便失去特別的能力.可憐神離開了他,他仍不知道!他犯了姦淫的罪,往妓女那裡去,受了大利拉的誘惑；大利拉有美麗迷人的魅力,表面上她深愛參孫,但心裡卻要害他的命.從大利拉的表現,可以正視魔鬼,她會裝作光明的天使,她是包著糖衣的毒藥；外表很好看,裡面卻充滿了詭計,使人防不勝防.倘若沒有大利拉,參孫的頭髮不會被剪,眼睛不至被挖；更不至關在仇敵營中推磨,受盡凌辱和譏笑.
\newline
\newline
參孫享有神特別的恩典,特別的能力；但裡面卻不能持守,不能分辨,以致根基不穩固,甚至失去了能力!神的恩賜種類很多,包括了頭腦的.身體的；但這一切恩賜都要以屬靈為根基,否則對己對人都有損害.正如一把鋒利無比的剪刀,他落在成年人的手中,用途甚廣；但在一個小孩子手裡,就可能造成損害.我們屬神的人,裡面應當有分別的力量,曉得分辨是非.神的恩賜用在一個有屬靈份量的信徒身上,這恩賜有如耕農的工具,又如戰爭的武器.反過來說,一個人沒有屬靈的根基,妄用神的恩賜；這一切的力量,就可能成為自己或別人的網羅.參孫就是在這事上失敗了,受了魔鬼的引誘,因著追求情慾而失去拿細耳人的記號.他寧可選擇大利拉而放棄拿細耳的職份；正如以掃寧可選擇紅豆湯而失去長子名份一樣.教會也常見這種情形,信徒為追求一點物質,貪圖一時安逸,便放棄了主給我們的大恩典.比如在參加教會聚集前,剛剛有精采的電視節目；那麼信徒便往往選擇了電視的節目,而放棄教會的聚會了.「電視」恐怕就成了我們的「大利拉」與「紅豆湯」了.
\newline
\newline
「耶和華已經離開了他」,可憐的參孫自己還不知道.昔日以弗所教會失去起初的愛心,它仍是教會；狄撒教會名生實死,老底嘉教會不冷不熱,但教會的名還在；而教會的實質就不見了.外表的形式可能繼續存在,直至考驗來到,才能分辨真假.屬靈的仇敵非利士人來了,起來壓迫教會,教會站立得住與否?全憑裡面的根基如何?一棵高大的樹能使多人在其蔭下乘涼；有美麗的外體,英雄的體態；但有一天樹忽然倒下來,原來裡面被蟲蛀通了,不能再繼續生存.
\newline
\newline
神沒有忘記參孫,在他受苦時,神特別記念他；這時候,參孫也想念神了.他後悔自己所犯的罪,求神再賜給所收回的恩賜；神答允所求,於是他便重新得力.他在非利士人的廟裡,寧可喪失自己的生命,仍要顯出神的榮耀,那日他所殺的非利士人比生前還多.
\newline
\newline
許多年青的弟兄姊妹,有恩賜可以為主用,但被世界的誘惑成了「大利拉」；「紅豆湯」成了心中的攔阻,毀壞了神在你身上的恩賜和計劃,但願神的話感動我們的心,使我們悔改,復興作成祂所使用的器皿.
\newline
\newline


\newpage

\section{第42屆~港九培靈會~王永信~第十三講}
\label{sec:1457}
\textbf{假的教會}
\newline
\newline
連結: \href{https://www.hkbibleconference.org/session-message/view/1457}{\texttt{https://www.hkbibleconference.org/session-message/view/1457}}
\newline
\newline
\hyperref[sec:1456]{< < < PREV SERMON < < <}
~
\hyperref[sec:toc]{[返目錄]}
~
\hyperref[sec:1458]{> > > NEXT SERMON > > >}
\newline
\newline
士師記 17:1-17:13
\newline
\begin{longtable}{cl}
\hline
\hline
章節 & 經文 (和合本修訂版)\\
\hline
17:1 & \begin{tabularx}{0.7\textwidth}{X} 以法蓮山區有一個人，名叫米迦。 \end{tabularx} \\ \\ \relax
17:2 & \begin{tabularx}{0.7\textwidth}{X} 他對母親說：「你的一千一百塊銀子被人拿走了，為此你發咒起誓，也說給我聽。看哪，銀子在我這裡，是我拿的。」他母親說：「願我兒蒙耶和華賜福！」 \end{tabularx} \\ \\ \relax
17:3 & \begin{tabularx}{0.7\textwidth}{X} 米迦把這一千一百塊銀子還他母親。他母親說：「我把這銀子分別為聖，親手獻給耶和華，為我兒子造一尊雕刻的像，以及一尊鑄成的像。現在我把銀子交給你。」 \end{tabularx} \\ \\ \relax
17:4 & \begin{tabularx}{0.7\textwidth}{X} 米迦把銀子還他母親，他母親把二百塊銀子交給銀匠，去造一尊雕刻的像，以及一尊鑄成的像，安置在米迦的房子裡。 \end{tabularx} \\ \\ \relax
17:5 & \begin{tabularx}{0.7\textwidth}{X} 米迦這個人有了神堂，又造了以弗得和家中的神像，派他的一個兒子作祭司。 \end{tabularx} \\ \\ \relax
17:6 & \begin{tabularx}{0.7\textwidth}{X} 那時，以色列中沒有王，各人照自己眼中看為對的去做。 \end{tabularx} \\ \\ \relax
17:7 & \begin{tabularx}{0.7\textwidth}{X} 猶大的伯利恆有一個年輕人，是猶大族的人。他是利未人，寄居在那裡。 \end{tabularx} \\ \\ \relax
17:8 & \begin{tabularx}{0.7\textwidth}{X} 這人離開猶大的伯利恆城，要找一個可住的地方。他來到以法蓮山區米迦的家，還要往前走。 \end{tabularx} \\ \\ \relax
17:9 & \begin{tabularx}{0.7\textwidth}{X} 米迦對他說：「你從哪裡來？」他說：「我從猶大的伯利恆來。我是利未人，要找一個可住的地方。」 \end{tabularx} \\ \\ \relax
17:10 & \begin{tabularx}{0.7\textwidth}{X} 米迦說：「你就住在我這裡吧！我以你為父為祭司，每年給你十塊銀子和一套衣服，以及生活所需的食物。」利未人就來了。 \end{tabularx} \\ \\ \relax
17:11 & \begin{tabularx}{0.7\textwidth}{X} 利未人願意和這人同住；他待這年輕人如自己的兒子一樣。 \end{tabularx} \\ \\ \relax
17:12 & \begin{tabularx}{0.7\textwidth}{X} 米迦授這年輕的利未人祭司的職任，他就住在米迦的家裡。 \end{tabularx} \\ \\ \relax
17:13 & \begin{tabularx}{0.7\textwidth}{X} 米迦說：「現在我知道耶和華必恩待我，因為我有利未人作我的祭司。」 \end{tabularx} \\ \\
[1ex]
\hline
\hline
\end{longtable}
士師時代的歷史告訴我們,以色列的民族真是一代不一代,在神面前越來越墮落.最初他們的爭戰是一半的,有如不徹底的教會；其後他們又往後退,成為一個中立式的教會,對於神的爭戰漠不關心響應.進而再成為失去能力的教會,浪費神恩,失去聖靈所賜的力量.可是他們仍要往下去,達到極點,就是士師記十七章所記載的,他們成了一個假的教會.
\newline
\newline
(一)米迦的墮落(士十七1-6)
\newline
\newline
以法蓮山地有一個青年名叫米迦.他偷了母親一千一百舍客勒銀子,其後母親發覺錢被人偷去,因此十分著急；便咒詛那偷錢的人,米迦聽聞異常懼怕.由於舊約時代一切祝福或咒詛的話,都要應驗在人的身上；米迦就把錢歸還母親,又向母親認自己的罪.這位母親不但沒有教導她的兒子,反而對他說:「我兒啊,願耶和華賜福給你.」母親接過錢來分出一部份獻給耶和華,用來雕刻一個銀像,安置在米迦的屋內.這位母親的心意是十分矛盾,一方面祈求耶和華賜福給他的兒子,另一方面又為兒子做偶像.她的宗教以是形式化的宗教；她對神的認識,也是頭腦上的認識.她不是真心相信,甚至所行的都是違背神的事；口裡親近神,心思意念卻與神相離.米迦和他的母親一同犯罪,他們把最重要的十誡也忘記了.為自己來雕刻偶像,犯罪得罪神,母親年紀老,可能一時糊塗；但兒子不應如此,彼此在真理上胡作妄為.
\newline
\newline
米迦接受了母親所送的偶像,便在家中找一間房間,改裝成為神堂.把他們所有的偶像安置在那神堂,在這個屬神的屋子裡,不但有一個偶像,而且還有許多的偶像.正如許多信徒的心中,也充滿了偶像,佔有了神的地位一樣.米迦的家裡除了偶像以外,還製造了以弗得和家中的神像,以弗得是祭司的衣服；他不是利未人,只不過是以法蓮人.祭司之職與他無關,這祭司的衣服也不能擅自製造,聖所只能讓祭司進去.但米迦竟為自己製造了以弗得,又分派他的兒子作祭司；原文意義就是使自己兒子成聖,米迦可謂妄自尊大,在家中擅立聖職.當時以色列人的會幕是放在示羅,百姓需要齊集示羅,在會幕內敬拜神.人在這事上不能隨便,但米迦竟膽妄自製造以弗得,又立兒子作祭司.這個家庭,不但老母親不明事理,兒子也是不明,甚至孫子更是如此；妄任祭司之職,他們三代一同犯罪得罪神.
\newline
\newline
米迦的名意是「像神」,他的為人也是如此；他只像神,而不屬神.他的家庭設有神堂,以弗得,及祭司；外表好像是教會,但內裡卻不是.孫子任祭司之職,但實則卻不是利未人,不能擔任此職.因此這教會正好比作現今末世的假教會,只有外表的形式,而缺乏了實際.正因為當時以色列中沒有王,各人任意而行,在混亂之中,便產生了這教會.假的教會不但失去力量,而且還不能稱之為教會；因為其中一切屬神的人墮落了,他們口中充滿耶和華,心中卻充滿了偶像.
\newline
\newline
(二)利末人的墮落(士十七7-13)
\newline
\newline
舊約時代的利未人是代表了事奉神的人,他們在教會中替神傳言,假若他們事主不忠,教會也會有不忠誠的表現.聖經告訴我們:「......猶大伯利恆有一個少年人,是猶大族的利未人......這人離開猶大伯利恆城,要尋找一個可住的地方.......」這個利未人要找安身之所,表明他的生活發生了問題；本來利未人不是靠自己工作,來維持生活的,他們是分別為聖,特別事奉神的人.神命定其他支派的人需要負責他們的生活,現今這利未人的情形,正表明了以色列人忽略了利未人,在生活上所需要的供應.
\newline
\newline
甚麼時候教會冷淡,信徒靈性退步,他們對神的工作也會冷淡,對神僕人的需要及教會的需要,也缺乏了負擔.信徒對神的奉獻是我們的權利和義務,舊約的律法最少要把十分之一獻上；新約時代更要把標準提高,倘若我們在金錢奉獻上不忠,就是欠了神的債.今日的教會有許多工作都需要發展,例如:海外傳福音的工作,文字的工作等等；但因著經濟的缺乏,以致工作不能作全面性的展開.我們的信仰需要言行一致,不但聽道也需要行道；理論加上行動便實行出來.我們口裡說愛神愛人,但眼見許多失喪的靈魂走向滅亡之路而不顧,反為自己積藏財寶.由於教會對傳道人不合理的待遇與供應,使不少青年足不前,不肯獻身事奉主.甚麼時候教會的信徒靈性低沉,教會就會忘記神僕人的需要.
\newline
\newline
這個猶大族的利未人要離開猶大伯利恆,為要找一個可住的地方,便來到以法蓮山地,走到米迦的家.這利未人在極度困難的生活中,也開始墮落了；他忘記用信心事奉神,他忘記利未人一切的產業是耶和華.他開始尋找人的辦法,他到了米迦的家；在那裡,他必定看見米迦的家中有許許多多的神像.他本身是利未人,有聖潔祭司的職位；現今他看見米迦犯如此大罪,竟不出一句責備的話,只顧為自己找安居之所.米迦請他住在他的家,並以每年十舍客勒銀子,一套衣服,和度日的食物為酬聘那人為祭司.聖經告訴我們:「......利未人就進了他的家,利未人情願與那人同住.......」這是何等羞恥的事,一個利未人竟留在一所充滿偶像的家中,擔任祭司之職.他不是為服事神,乃是為尋找可住的地方；解決了每日的衣,食,住,行,他只注重物質方面的享受.米迦在真理方面作了一個糊塗人,他以為自己家裡有會幕的設備,又有利未人來作祭司之職,便組成一個真正的教會,必然蒙神的賜福.其實這一切都是外表,出於裝假的,他所擁有的教會只不過是一間假的教會.
\newline
\newline
啟示錄第十三章裡,記載兩獸的預言.第一隻從海中上來的獸代表了敵基督,另一隻從地裡上來的獸,在外表看來像羊羔,像基督,又像教會,但他說話卻像龍,龍就是代表魔鬼,也是代表了末世時代,一切不信派的假基督,和將來那些實行普世合一的宗教集團.這末世假的教會,外表的一切,與真的教會無異；有十字架,又有講壇,但裡面卻是魔鬼的思想,充滿了假神及偶像.現今是末世的時代,假的教會越來越顯靈了.
\newline
\newline
米迦這個名意果真應驗了,他實在是「像神」而不「屬神」.家裡佈置得像教會,但不是真教會；他的家庭極其可憐,老少三代一同在神面前裝假,失落,違背神,失去熱心.教會是神的家,在家裡有老,有少,才算是正常的現象,正如米迦一樣；這家可以一同起來真心服事神,卻也可以一同違背神.倘若他們全體,無論老少,在神的事上全心全意事奉神,神必然賜福.摩西在昔日帶領以色列人出埃及,在路途中行經利非訂,在那裡他們要與非利士人爭戰；當日以色列人中分成兩支軍隊,一支由摩西帶領到山上為以色列人的爭戰而禱告,展開禱告的戰場.另一支由約書亞帶領以色列民,在平地上與敵人交戰,因為以色列全家合一,年老的摩西與年少的約書亞完全合作,因此神賜福給他們,使他們得著完全的勝利.現今教會年長的與年少的兩代常產生隔膜,有分家的現象,這並不能推展教會的聖工,我們應當彼此謙卑,彼此尊重,同心合意興旺神的家,神的工作才能早日完成.
\newline
\newline
我們要求主救我們脫離假的教會,與假的事奉,不能再繼續在表面上裝假；各人應存謙卑的心,無論老少在教會中一同推進神的聖工.
\newline
\newline


\newpage

\section{第42屆~港九培靈會~王永信~第十四講}
\label{sec:1458}
\textbf{悔改的教會}
\newline
\newline
連結: \href{https://www.hkbibleconference.org/session-message/view/1458}{\texttt{https://www.hkbibleconference.org/session-message/view/1458}}
\newline
\newline
\hyperref[sec:1457]{< < < PREV SERMON < < <}
~
\hyperref[sec:toc]{[返目錄]}
~
\hyperref[sec:1459]{> > > NEXT SERMON > > >}
\newline
\newline
撒母耳記下 12:1-12:10
\newline
\begin{longtable}{cl}
\hline
\hline
章節 & 經文 (和合本修訂版)\\
\hline
12:1 & \begin{tabularx}{0.7\textwidth}{X} 耶和華差遣拿單到大衛那裡。拿單到了他那裡，對他說：「在一座城裡有兩個人，一個是富翁，一個是窮人。 \end{tabularx} \\ \\ \relax
12:2 & \begin{tabularx}{0.7\textwidth}{X} 富翁有極多的牛群羊群； \end{tabularx} \\ \\ \relax
12:3 & \begin{tabularx}{0.7\textwidth}{X} 窮人除了所買來養活的一隻小母羊之外，一無所有。小羊在他家裡和他兒女一同長大，吃他所吃的，喝他所喝的，睡在他懷中，在他看來如同女兒一樣。 \end{tabularx} \\ \\ \relax
12:4 & \begin{tabularx}{0.7\textwidth}{X} 有一客人來到這富翁那裡，富翁捨不得從自己的牛群羊群中取一隻招待來到他那裡的旅客，卻取了窮人的小母羊，招待來到他那裡的人。」 \end{tabularx} \\ \\ \relax
12:5 & \begin{tabularx}{0.7\textwidth}{X} 大衛就非常惱怒那人，對拿單說：「我指著永生的耶和華起誓，做這事的人該死！ \end{tabularx} \\ \\ \relax
12:6 & \begin{tabularx}{0.7\textwidth}{X} 他必須償還小母羊四倍，因為他做這事，沒有憐憫的心。」 \end{tabularx} \\ \\ \relax
12:7 & \begin{tabularx}{0.7\textwidth}{X} 拿單對大衛說：「你就是那人！耶和華－以色列的神如此說：『我膏你作以色列的王，我救你脫離掃羅的手； \end{tabularx} \\ \\ \relax
12:8 & \begin{tabularx}{0.7\textwidth}{X} 我將你主人的家業賜給你，將你主人的妃嬪交在你懷裡，又將以色列和猶大家賜給你；若還嫌少，我也會如此這般加倍賜給你。 \end{tabularx} \\ \\ \relax
12:9 & \begin{tabularx}{0.7\textwidth}{X} 你為甚麼藐視耶和華的命令，做他眼中看為惡的事呢？你用刀擊殺赫人烏利亞，又娶了他的妻子為妻，借亞捫人的刀殺死他。 \end{tabularx} \\ \\ \relax
12:10 & \begin{tabularx}{0.7\textwidth}{X} 現在刀劍必永不離開你的家，因你藐視我，娶了赫人烏利亞的妻子為妻。』 \end{tabularx} \\ \\
[1ex]
\hline
\hline
\end{longtable}
歷代以來,教會追求復興,門徑只有一個－徹底的認罪和悔改,除此以外並無別法了.要教會復興如初期時代一樣,我們就要為此付上代價；就是在大背道犯罪的事上,徹底的痛悔.
\newline
\newline
大衛犯罪以後悔改的經歷,成了我們好的榜樣.他雖然身為一國之君,又為神所愛,是一位合神心意的王；但他犯罪的史實也被記載在聖經內,可見神是公義不偏待人的.因著大衛外邊開始犯罪,眼睛受到罪的誘惑,心中便產生犯罪的念頭,接著實行去犯罪；每當一個人內裡的心思放鬆,外表也跟著放鬆便有所行動.大衛首先犯姦淫的罪,繼而犯謀殺的罪,可見罪惡連累的可怕.人不能因自己認識神,為神所愛,恩賜又多,而斷言自我再不犯罪；這一切外面的情況,都不能成為我們不犯罪的保障.
\newline
\newline
一日,太陽平西,大衛從床上起來,在王宮的平頂上遊行.大衛放鬆自己每天的生活,不治國政,睡至黃昏,醒過來時便思想壞事,因為他裡面的心思放鬆,外面的行為也放鬆.大衛在這國泰民安的社會裡,就放縱自己的私慾,奪人妻殺人夫,他在神面前犯了三件大罪；就是淫亂,殺人和硬心,因為他犯罪以後,並沒有立刻悔改.一年多過去,神藉先知拿單揭發他的罪行,他纔知道悔改,可見他硬心到甚麼程度.大衛是一個深識神的人,竟然在犯大罪以後；能賄賂良心渡過了一年多的生活,真是不可思議.
\newline
\newline
我們可以猜想,大衛在這年多以來的生活情況是如何:
\newline
\newline
(一)雙重的生活
\newline
\newline
大衛作了兩面人,他犯罪以後,仍然粉飾太平在宮中生活.他是過著雙重的生活,來掩飾自己在百姓面前所行的大惡.他在百姓面前是一副面孔,背著百姓卻又是另一副面孔.因為他身為一國之君,又是國中屬靈的領袖,他要在屬靈的事上作榜樣,所以他要在百姓面前裝作敬虔的外表,但在個人的私生活裡卻是充滿罪惡.我們可以痛斥大衛的假冒為善,雖然在我們生活裡,沒有犯這相同的罪,但也有很多類似的事例發生.例如明知那是神不喜悅的事仍堅持去做,對教會的事奉,態度錯誤,個人的生活充滿羞恥的事不敢見神.這些東西存留在我們生命中,日復一日,年復一年,仍然未徹底對付；有時被良心責備不安,但稍後隨即忘記了.
\newline
\newline
約翰壹書一9說:「我們若認自己的罪,神是信實的,是公義的；必要赦免我們的罪,洗淨我們一切的不義.」這節經文我們常歡喜用於佈道聚會中,請那些決信主的朋友悔改認罪,接納耶穌基督為救主；但這節經文對於信主的人也有同樣的價值,有時信主的人犯了罪更難認罪,心中更剛硬.這經文對世人,對教會,徹底向神認罪,都是十分適合的；除非我們脫離罪的生活,纔能完全順服神.
\newline
\newline
(二)良心剛硬
\newline
\newline
大衛犯罪一年多仍不肯悔改,久而久之,他的心越變越硬,甚至趨於麻木.他安於逸樂,絕不為自己的罪而憂傷；他拒絕聖靈的感動,逃避認罪的機會.他可說已經到了心死的地步.古語云:「哀莫大於心死,而身死次之.」大衛以為他能逃避神的審判,但神絕不放過任何一個犯罪的人；神要向犯罪的人討罪,除非他悔改認罪.倘若不悔改,就不能從罪中出來,也不能從罪中得釋放!從聖經裡,我們看見很多人犯罪以後,認罪的態度,可作我們的借鏡.
\newline
\newline
法老的認罪.神要法老釋放以色列人,讓他們出埃及.法老王卻反對,所以神容讓十大災害臨到埃及.法老王不知所措,便認罪懇求摩西為他禱告,允准以色列人出埃及.他雖認罪,但態度不夠徹底,只不過是由於畏懼之心產生,並且是另有目的的.
\newline
\newline
巴蘭的認罪.巴蘭說預言要咒詛神的選民,神的使者攔阻他.他就向神認罪懺悔,但他認罪的態度只屬表面化,因他貪心的程度比認罪更多,因此受到神的懲罰.
\newline
\newline
亞干的認罪.他的認罪不是由於悔過,乃因偷窺當滅之物被人揭破；逼於無奈而認罪,這被逼性的認罪,不能蒙神悅納.
\newline
\newline
掃羅的認罪.他逼迫大衛,後來又後悔,也曾多次向大衛認罪；但認罪不夠徹底,也不持久.他感到自己不該殺大衛,但當那感動過去,仇恨的心又起.
\newline
\newline
浪子的認罪.浪子回頭,他的認罪是十分撤底的；他在墜落失敗時,回到父家,向父親認罪說:我得罪了天,又得罪了你.他不但肯認罪,而且又肯謙卑,用行動證明他有悔改的表現.他認罪是從心裡痛悔以往的不是,他悔改是用行動來悔改,回到父家.所以他的悔改是徹底的,有深度的,也是持久的.
\newline
\newline
如今看大衛的悔改.神藉先知拿單向大衛忠告,指責大衛一年以來隱藏在心裡的罪惡；大衛聽見了,立刻悔改認罪.他說:「我有罪了,我得罪了耶和華.......」可能這句話在他心中潛伏了多時,但因私慾與高位,不能啟齒向神認罪.現在他果真悔改了,他在神面前俯伏禁食,不聽音樂,極其自卑痛悔.神就顧念大衛,因義人雖七次跌倒,也必起來.神差遣先知來指責他的罪惡,表明神未放棄他.這一次的經歷,使大衛不但認識神,更是認識自己；只不過是人,人有罪惡的本性,從裡面到外面是一無可取.如果大衛未曾經歷到犯罪的痛苦,可能他永遠不會認識自己.也許他想,像我這樣的人還會犯姦淫,謀殺人的罪嗎?神常讓激打,磨煉臨到我們身上,目的是讓我們更認識自己.除非人認識自己敗壞的本性,否則是不會在神面前謙卑的.
\newline
\newline
由以上的例子可以發現,我們表面上皮毛的悔改,心裡悔而不痛,在神面前是一無價值的.在這裡祇提出真正悔改所包括的四種意義:
\newline
\newline
(1)承認犯罪的事實──人犯罪以後歡喜用一些理由來掩飾,欺騙自己的良心.例如某些人不承自己是驕傲,只說這是個性強一點.這樣的辯駁,就否定了犯罪的事實.我們犯罪是一件一件犯的,認罪也該一件一件認；不能籠統的認,要承認自己個人犯罪的事實,必須向神認又要向人認,要付出代價來.
\newline
\newline
(2)認識罪惡本質是邪惡的──人常被罪迷惑,因為罪常以美麗的姿態出現；披有悅目的外衣,如糖衣的毒藥.我們不可被外表所朦蔽,要正視罪正如毒蛇烈火一樣.
\newline
\newline
(3)認識犯罪是直接得罪神──我們犯任何一件罪都與神有關,必要受神的刑罰,在神面前要向神負責任.倘若我們不認罪,我們就不能從那罪中出來；未信主的人認罪可得救,信主的人認罪可得福.信主的人犯罪而不認是可憎的,等於把主重釘十字架.
\newline
\newline
(4)認識人有罪的本性──我們不但痛悔外表的罪行,更要認識自己裡面的罪性,要注意犯罪的傾向與動機.不但在神面前求神赦免,所犯自私自利的罪,也要赦免這個犯自私自利罪的人.這樣我們才能在神面前有長進,把罪與罪人一同交在神的手,求神造清潔的心重新賜下正直的靈.
\newline
\newline
我們必須在神前認識自己的軟弱,把這歡喜犯罪的人交給主；求主用十字架來對付,又洗淨一切所認的罪.能這雙管齊下,我們對罪才能根本的解決.我們的悔改就有意義了,我們的人生才有轉變,我們的教會,才能得到復興.
\newline
\newline


\newpage

\section{第42屆~港九培靈會~王永信~第十五講}
\label{sec:1459}
\textbf{以便以謝感恩的教會}
\newline
\newline
連結: \href{https://www.hkbibleconference.org/session-message/view/1459}{\texttt{https://www.hkbibleconference.org/session-message/view/1459}}
\newline
\newline
\hyperref[sec:1458]{< < < PREV SERMON < < <}
~
\hyperref[sec:toc]{[返目錄]}
~
\hyperref[sec:1445]{> > > NEXT SERMON > > >}
\newline
\newline
撒母耳記上 4:1-4:11
\newline
\begin{longtable}{cl}
\hline
\hline
章節 & 經文 (和合本修訂版)\\
\hline
4:1 & \begin{tabularx}{0.7\textwidth}{X} 撒母耳的話傳遍了全以色列。以色列人出去與非利士人打仗，安營在以便‧以謝，非利士人安營在亞弗。 \end{tabularx} \\ \\ \relax
4:2 & \begin{tabularx}{0.7\textwidth}{X} 非利士人向以色列人擺陣。兩軍交戰的時候，以色列人敗在非利士人面前；非利士人在戰場上殺了他們約四千人。 \end{tabularx} \\ \\ \relax
4:3 & \begin{tabularx}{0.7\textwidth}{X} 百姓回到營裡，以色列的長老說：「耶和華今日為何使我們敗在非利士人面前呢？我們要將耶和華的約櫃從示羅抬到我們這裡來，好讓他來到我們中間，救我們脫離敵人的手掌。」 \end{tabularx} \\ \\ \relax
4:4 & \begin{tabularx}{0.7\textwidth}{X} 於是百姓派人到示羅，從那裡將坐在二基路伯上萬軍之耶和華的約櫃抬來。以利的兩個兒子何弗尼、非尼哈也與神的約櫃同來。 \end{tabularx} \\ \\ \relax
4:5 & \begin{tabularx}{0.7\textwidth}{X} 耶和華的約櫃到了營中，全以色列就大聲歡呼，連地都震動。 \end{tabularx} \\ \\ \relax
4:6 & \begin{tabularx}{0.7\textwidth}{X} 非利士人聽見歡呼的聲音，就說：「為何希伯來人在營裡這麼大聲歡呼呢？」他們知道耶和華的約櫃到了營中。 \end{tabularx} \\ \\ \relax
4:7 & \begin{tabularx}{0.7\textwidth}{X} 非利士人就懼怕，說：「有神明到了他們營中。」又說：「我們有禍了！從來不曾有這樣的事。 \end{tabularx} \\ \\ \relax
4:8 & \begin{tabularx}{0.7\textwidth}{X} 我們有禍了！誰能救我們脫離這些大能之神明的手呢？從前在曠野用各樣災禍擊打埃及人的，就是這些神明。 \end{tabularx} \\ \\ \relax
4:9 & \begin{tabularx}{0.7\textwidth}{X} 非利士人哪，要剛強，要作大丈夫，免得作希伯來人的奴僕，如同他們作你們的奴僕一樣。你們要作大丈夫，與他們爭戰。」 \end{tabularx} \\ \\ \relax
4:10 & \begin{tabularx}{0.7\textwidth}{X} 非利士人進攻，以色列人敗了，各往自己的家逃跑。被殺的人很多，以色列倒下的步兵有三萬。 \end{tabularx} \\ \\ \relax
4:11 & \begin{tabularx}{0.7\textwidth}{X} 神的約櫃被擄去，以利的兩個兒子何弗尼、非尼哈也都被殺了。 \end{tabularx} \\ \\
[1ex]
\hline
\hline
\end{longtable}
士師時代的以色列人,處於無政府狀態,因為他們國中沒有王,各人任意而行,在民間只見一片混亂與爭戰.
\newline
\newline
以色列人又與非利士人交戰,兩軍對峙的時候,以色列人竟敗在非利士人面前.百姓回到營內,以色列的長老彼此對說:耶和華今日為何使我們敗在非利士人面前呢?我們不如將耶和華的約櫃,從示羅抬到我們這裡來,好在我們中間救我們脫離敵人的手.耶和華的約櫃到了營中,以色列眾人就大聲呼叫,地便震動；非利士人聽見懼怕起來.以色列人與非利士人再次交戰,以色列人又敗了,被殺的甚多；神的約櫃被擄去,祭司以利兩個兒子也都被殺了.眾民急忙向老祭司以利報訊,以利聽聞之後就從他的位上往後跌倒,在門旁折斷頸項而死.以利的兒婦懷孕將到產期,他聽見神的約櫃被擄去,公公和丈夫都死了；就猛然疼痛,曲身生產,產了一個男孩,他給孩子起名叫以迦得.說:「榮耀離開以色列了.」
\newline
\newline
從這件事來看,可見以色列人的宗教觀念,是何等形式化.在他們信仰中,以為自己犯罪,與在耶和華以外另立別神是沒相關的,他們以為遇事只要隨時把約櫃搬出來,就可以鎮壓敵人,叫他們知難而退了.這實在是捨本求末之法.以色列人犯罪事奉偶像,神就離開他們；甚至連約櫃也讓非利士人擄去!神的榮耀歸開以色列實在是一件悲悽不過的事.
\newline
\newline
年青的撒母耳,被神呼召,繼承祭司以利之職,起來作以色列人屬靈的領袖.他苦心教導以色列人,周遊全國向他們傳講律法,召開培靈會,使百姓離開罪惡,悔改歸正.他又訓練先知門徒,開設先知學校,他殷勤作工達二十年之久,忠心致意教導以色列民,奇蹟便發生了,這是士師時代從未發生過的.「......過了二十年,以色列全家傾向耶和華.」(撒上七2)靠著撒母耳一人的努力,使全以色列家歸向真神,這工作實不容易,所以撒母耳對以色列全家說:「們若一心歸順耶和華,就要把外邦的神和亞斯他錄,從你們中間除掉；專心歸向耶和華,單單的事奉他.......」撒母耳要眾民作一個決定,倘若要接受耶和華是真神,就要全心歸向他,不能一方面愛神,另外一方面又愛巴力,假神,世界,瑪門.這教訓從古至今都是如此,人的心不能分為二；不是愛這個惡那個,就是重這個輕那個.在基督的生命中,有多少份量是全心愛神的呢?可能在我們一生中,從未全心愛過神,我們常讓其他的東西作主.請問你每早晨起來,睜開眼睛,第一個思想是否想到神呢?一個人如果夠在工作以先朝見神,那麼他一天的生活就會充滿快樂與能力.活水江河便自然而來,任何的試煉不會使我們喪膽.我們把主接到心裡,自然會成聖,並且成為神的勇士.不但重生,更要長大成人；不但有長進,更在真道上自建,成為一個有活潑生命的基督徒.
\newline
\newline
撒母耳告訴以色人要一心愛主,把神活活的接到心裡,他們的事奉才有果效.以色列人就尊著他的話去行,棄丟一切的偶像；那天在以色列眾民中,就有了復興的轉機.撒母耳為著這次的復興,足用了二十年悠長的歲月來預備,纔有美好的成果.
\newline
\newline
撒母耳在米斯巴召開大會,吩咐眾民都來參加.這是以色列全國的奮興培靈會,講員就是撒母耳.撒母耳教導他們,為他們禱告耶和華.眾民就聚集在米斯巴,打水澆在耶和華面前,當日禁食說:「我們得罪耶和華.」倒水表示將在神面前倒空之意,眾民就在那裡痛苦流涕彼此認罪.他們向列祖的神禱告,把痛悔的心,將來的盼望,都澆在神面前.他們又禁食,在神面前自卑,克苦己心.所以有大復興臨到.復興是由禱告與悔改產生的,因為禱告和悔改是復興的美好準備.撒母耳在米斯巴審判以色列人,他能徹底對付罪,以色列全家纔會歸向耶和華.
\newline
\newline
撒母耳為以色列人作了完全的工作:
\newline
\newline
(1)用二十年作準備工作,用心良苦.開奮興會,帶領眾民認罪,克苦己心,使以色列全家傾向耶和華.
\newline
\newline
(2)審判眾民,嚴厲對付罪惡,又把罪從他們中間滾出去.
\newline
\newline
當他們在米斯巴召開大會時,非利士人又上來攻擊他們.非利士人是撒旦的工具,常做破壞的工作,從不歡喜看見屬神的人團結在一起,教會團結了,牠便要起來施行詭計,使教會分化；使教會長輩與晚輩不和,彼此紛爭不能團結.牠要浪費我們的力量,分化我們,這是牠攻擊我們最厲害的工具.非利士人擺好隊伍預備要攻擊以色列人,百姓十分恐懼,但撒母耳非常鎮定,他在這時獻上一隻小羊羔,表示感恩,和求神保護.結果神用大雷的聲音,驚亂非利士人,他們就敗在以色列人面前.
\newline
\newline
二十年前,以色列人在以便以謝,慘敗於非利士人的手中；但二十年後的今天,他們兩軍又在此相遇,再決雌雄.這次他們卻轉敗為勝,他們的勝利並不是因為經過二十年的時光,刀槍戰略都進步了；戰爭勝敗的關鍵乃在乎他們是否心中有偶像?是否一心歸向耶和華.他們在爭戰中得了勝利,撒母耶便將一塊石頭,立在米斯巴和善中間,給石頭起名叫以便以謝,說:「到如今耶和華都幫助我們!」或許有人會發問:撒母耳向神感恩讚美,神是以便以謝的神,難道他忘記了二十年前,以色列人慘敗的情景嗎?當時祭司以利和他兒子的慘死,約櫃被擄到外邦,百姓被殺的有三十萬.這悽涼慘痛的往事,也算是神一直幫助以色列人嗎?神是以便以謝的神,因為聖經告訴我們:「萬事都互相效力,叫愛神的人得益處.」「萬事」就是包括了溫暖的及冷酷的,失敗的及得勝的,痛苦的及快樂的......,我們得勝的時候是「效力」,失敗的時候也是「效力」,無論我們的境況如何?他都會幫助我們,使我們藉那件臨到我們身上的事得著益處.
\newline
\newline
一個基督徒爭戰的生活,充滿了許多的失敗,常常犯罪,甚至不好意思到神面前認罪.正如保羅說:「我所願意的善,我反不作；我所不願意的惡,我倒去作......我真的苦啊,誰能救我脫離這取死的身體呢?」神拯救我們成為祂的兒女,但我們雖是重生了,卻還未長大成人；稱義了卻還未成聖,因而產生恐懼之心.倘若我們這樣感覺,這纔是正常,為當人感到自己的生活不像樣時,纔會竭力追求聖潔.肯追求然後有長進,主允准失敗臨到我們身上,經過失敗以後,我們就能認識自己.
\newline
\newline
神不容讓我們一天到晚都是失敗,他會給我們鼓勵,幫助我們得勝,我們也不能在一切的工作上完全亨通；神使我們成功,但也讓我們有時失敗,他使用這兩種方法使我們在雙重對付下,得到平衡的生長.不會太驕傲,也不會太自卑,因為兩種極端都不能使我們事奉神.神的慈愛和公義是平衡的,因此我們雖然失敗也得到赦免,我們若遠離神,祂的管教便臨到我們,所以我們必須倚靠祂常過得勝生活.
\newline
\newline
以便以謝的經歷對於撒母耳來說,完全是真實的.倘若二十年前以色列人未遭受失敗,打擊,和痛苦,他們的心仍未肯察覺己罪,不肯離棄偶像歸向耶和華.因著他一人努力事奉,帶領會眾謙卑悔改,便得到以便以謝寶貴的經歷.我們理當在一切失敗軟弱後,用行動悔改,徹底認罪；這樣神也能使我們教會,得著以便以謝的經歷.
\newline
\newline


\newpage

\section{第42屆~港九培靈會~王永信~第十六講}
\label{sec:1445}
\textbf{建造教會}
\newline
\newline
連結: \href{https://www.hkbibleconference.org/session-message/view/1445}{\texttt{https://www.hkbibleconference.org/session-message/view/1445}}
\newline
\newline
\hyperref[sec:1459]{< < < PREV SERMON < < <}
~
\hyperref[sec:toc]{[返目錄]}
~
\hyperref[sec:1439]{> > > NEXT SERMON > > >}
\newline
\newline
歷代志上 28:1-28:5
\newline
\begin{longtable}{cl}
\hline
\hline
章節 & 經文 (和合本修訂版)\\
\hline
28:1 & \begin{tabularx}{0.7\textwidth}{X} 大衛召集以色列所有的領袖，各支派的領袖、輪班服事王的官長、千夫長、百夫長、掌管王和王子一切產業牲畜的、宮廷官員、勇士，和所有大能的勇士，都到耶路撒冷來。 \end{tabularx} \\ \\ \relax
28:2 & \begin{tabularx}{0.7\textwidth}{X} 大衛王站起來，說：「我的弟兄，我的百姓啊，請聽我說！我心裡本想建造殿宇，安放耶和華的約櫃，作為我們神的腳凳，並且我已經預備了建造的材料。 \end{tabularx} \\ \\ \relax
28:3 & \begin{tabularx}{0.7\textwidth}{X} 只是神對我說：『你不可為我的名建造殿宇，因你是戰士，流了人的血。』 \end{tabularx} \\ \\ \relax
28:4 & \begin{tabularx}{0.7\textwidth}{X} 然而，耶和華－以色列的神在我父的全家揀選我作以色列的王，直到永遠。因他揀選猶大為領袖，在猶大家中揀選我父家，在我父的眾兒子裡喜悅我，立我作全以色列的王。 \end{tabularx} \\ \\ \relax
28:5 & \begin{tabularx}{0.7\textwidth}{X} 耶和華賜我許多兒子，在我兒子中揀選我兒子所羅門坐耶和華國度的王位，治理以色列。 \end{tabularx} \\ \\
[1ex]
\hline
\hline
\end{longtable}
無論在甚麼時候,教會前面都擺著兩條路；就是神的路,和人的路.教會可以自由地選擇其一,若要追求復興,就需要選擇神的路.選擇神的路,就是脫離不徹底的教會,中立式的教會,失去能力的教會,和假的教會.更必須互相認罪,成為一個認罪的教會,以便以謝的教會.讓神拆毀,拆毀以後,把教會重建起來,因為我們不但要在消極方面感恩,更要在積極方面努力建造.
\newline
\newline
以色列人要建造聖殿,神的居所,在信仰方面來說,這是他們整個民族的光榮；因為聖殿是他們向神獻祭與朝見神的地方.在大衛心靈裡,深感為耶和華建造聖殿的重要,也願負起這艱巨的工作.但神卻不把這重大的責任交付大衛的手.因為自從大衛犯罪以後,以色列國外有禍患,內有恐懼.鄰近的摩押人,阿們人,亞蘭人與他為敵!在大委身上的光彩消失了,而且他兒子押沙龍叛變,謀朝篡位.是後愛子又被人所弒,老年喪子,他心裡充滿了悲傷及痛苦.他回憶過去,曾奪人妻殺人夫,大逆不道的罪惡,滿手沾染了流人血的血腥；這一切都應驗了神藉著先知拿單向他宣告的刑罰:「刀劍必不離開你家!」
\newline
\newline
歷代志上二十八章,我們在這裡可以看見大衛悲傷的晚年,稍為有點改變；好像雲彩被挪去陽光再次透進來一樣.大衛開始恢復他以往,屬靈生活的活力及權柄.他招聚了全以色列各支派的首領,和輪班服事王的軍長,與千夫長,百夫長,掌管王與王子產業牲畜的；並太監以及大能的勇士,都到耶路撒冷來,召開全國性的大會.這聚會是十分重要的,因為大衛要藉著這個機會,向全國領導者宣告,他兒子所羅門王將要繼承他的王位,及代替他完成建造耶和華殿宇的心願.大衛作王快到四十年了,現今是他作王的末年；在他眾子之中,誰是將來繼承王位作以色列之君的,這是眾民極關切的問題.倘若眾子相爭,恐怕就會鬧出大戰.所以大衛必須在他離位以前,向眾民有清楚的交代,使他們知道所羅門不但將要負起國家領袖的地位,而且又成為他們屬靈的領袖.
\newline
\newline
「我兒所羅門哪!你當認識耶和華你的父神,誠心樂意的事奉祂.」這句出自慈父肺腑之言,把他心底裡的意思都說出來；也是大衛人生所經歷到確實的一句話.在他整個人生裡,他發覺有兩件事是身為一個屬靈的領袖不能不知的；首先就是認識神,然後心樂意事奉神.大衛知道所羅門所站的地位是十分重要的,就是肩負承上繼下的重責,正如昔日以撒的地位一樣.在以撒以上有亞伯拉罕,以下有雅各,他自己站在中間.神賜給亞伯拉罕一切福份及應許都是依靠以撒傳授下去；否則神賜給以色列民一切的福份及應許從此便中斷,無法傳授與雅各的十二支派及以後的人類了.所羅門就是站在這中間位置的一個人,在他以上有大衛王,在他以後則有大衛的後裔,耶西的根耶穌基督.因為從所羅門直至耶穌之間,再沒有一個君王有如此偉大的表現.大衛的子孫,根苗,這最偉大的應許,都要經過所羅門而得著應驗.倘若所羅門不小心,在他那裡便要中斷,可見他的位置是十分重要.所以當大衛立所羅門成為國家的領袖及屬靈的領袖時,他有許多說話需要向這位年輕的新王提醒:
\newline
\newline
(一)「你要認識神」
\newline
\newline
這句話是值得我們去思想的.我們是否真正認識神,又是被神所認識的呢?不錯,當我們蒙恩得救時,我們已認識到祂是我們的救主,天父.但不少信徒自得救以後,他們對祂的認識仍然很生疏,與神交通未成為密友.僅是馬虎地來到祂面前,未在心靈上認識祂；與祂的關係未曾深切了解,在生活上也未曾努力尋求他.我們得救是白白的,主已將救恩大功作成,只要我們接受便可以獲得；但認識神是需要我們付代價,努力在聖經的亮光中發掘,因為我們尋找神,正如尋找隱藏的珍寶一樣.
\newline
\newline
身為一個屬靈領袖的所羅門,他的責任可算十分重大；若不深入認識神,是無法完成神的託負.昔日建造耶和華神的殿宇,完全是根據神所指定的樣式,一點也不能有錯誤.我們今天建造教會也是這樣,更需要認識神的心意,因為這屬靈的殿比物質的殿更重要；責任更重大,建造更困難.聖殿可分成兩部份的建造:
\newline
\newline
(1)建造神的殿──建立教會
\newline
\newline
這是每一個基督徒的責任,各盡其職,彼此在教會裡發展自己的恩賜,互相服事,這樣就是為神建造教會.
\newline
\newline
(2)建造裡面的殿──聖靈的殿
\newline
\newline
我們每一個基督徒的身體就是神的殿,這裡面的殿就是教會殿的根基.每一個信徒裡面的殿建造得好,神的殿就能有更好的建造.若神的殿需要有能力,我們個人裡面的殿就先要建造來；不能存有污穢,應當完全聖潔.因著我們每一個人都有這如此重大的責任及託負,就是為神建造教會,及建造自己裡面的殿.所以我們更該竭力追求認識神,讓我們更深認識祂兩個屬性:
\newline
\newline
1. 聖潔的神
\newline
\newline
祂是完全聖潔,遠離罪惡,不能被騙的一位.我們不能一方面犯罪,另一方面又事奉祂；也不能過著向罪惡妥協的生活,又被聖靈所膏抹.因為污穢的器皿,不合乎祂使用；我們若不自潔,神就不能賜福使用我們,除非我們求神赦免一切的罪.
\newline
\newline
2. 愛的神
\newline
\newline
他是充滿愛的一位,他知道人會犯罪,所以他樂意赦免人的罪孽.祇要我們禱告認罪以後,他就使我們重新得力,再一次湧出活水江河來.年老的大衛王,雖然以前他曾一度犯罪墮落,失去屬靈的能力；但因著他完全痛悔己罪,神的能力和拯救就再一次臨到他.完全赦免他一切罪惡,使他再能帶著權柄來說勸勉全國首領的話.聖經說義人雖七次跌倒,神仍使他起來.
\newline
\newline
(二)「你要誠心樂意事奉神」
\newline
\newline
我們不但要認識神是我們的密友,而且還有第二步,就是誠心樂意事奉祂,僅僅滿足在認識神,只屬消極性的追求；我們也要做積極性的事,就是誠心樂意事奉神.否則認識神所帶來的復興過去以後,生活還是沒有改變,這一切都歸徒然.心中認識神後,神賜我們靈裡得到復興,因著這一點的復興,又起來事奉神.如此,纔是正確合適的態度.
\newline
\newline
服事神的工作包括很多,除了探訪,傳道,禱告以外,在生活中也有很多機會為主工作.例如:在一個不信主的家庭中,能否為主緣故顯出百般的忍耐,當閱讀一本污穢的書籍時,能否為主緣故立刻放下?或有人故意冒犯你,能否為主不發脾氣?公司中工作,是否為主緣故寧願損失那些發不義之財的機會.事奉主的方式很多,按著神的帶領去行,這就是事奉主了.
\newline
\newline
大衛王和所羅門王,可代表年少和年老的兩代.兩代能彼此和睦,毫無隔膜,彼此看顧相愛.今天的教會也該以此為榜樣,不能因老少之間的意見不同而分家.這種兩代分家的現象,近年來,在各地教會都有發生,好像成了潮流一樣.要彼此在主裡合一,同心合意興旺福音,為神的家盡忠.今天神需要許多信徒,在他們的崗位上事奉神,在生活裡活出基督的樣式.將才幹不但為世用,也是獻身為主用.神更需要有人對祂的呼召有響應,當神呼召信徒出來全時間事奉神時,他們應該起來加入祂的行列中.
\newline
\newline
但願神今天向我們說話,正如昔日向大衛和所羅門說話一樣.使我們終身不忘記:「你要認識神」,「你要誠心樂意事奉神」.
\newline
\newline


\newpage

\section{第42屆~港九培靈會~黃聿侯~第一講}
\label{sec:1439}
\textbf{領人歸主的責任}
\newline
\newline
連結: \href{https://www.hkbibleconference.org/session-message/view/1439}{\texttt{https://www.hkbibleconference.org/session-message/view/1439}}
\newline
\newline
\hyperref[sec:1445]{< < < PREV SERMON < < <}
~
\hyperref[sec:toc]{[返目錄]}
~
\hyperref[sec:1550]{> > > NEXT SERMON > > >}
\newline
\newline
以西結書 3:17-3:21
\newline
\begin{longtable}{cl}
\hline
\hline
章節 & 經文 (和合本修訂版)\\
\hline
3:17 & \begin{tabularx}{0.7\textwidth}{X} 「人子啊，我立你作以色列家的守望者，所以你要聽我口中的話，替我警戒他們。 \end{tabularx} \\ \\ \relax
3:18 & \begin{tabularx}{0.7\textwidth}{X} 我何時指著惡人說：『他必要死』；你若不警戒他，也不勸告他，使他離開惡行，拯救他的性命，這惡人必死在罪孽之中；我卻要從你手裡討他的血債。 \end{tabularx} \\ \\ \relax
3:19 & \begin{tabularx}{0.7\textwidth}{X} 倘若你警戒惡人，他仍不轉離罪惡，也不離開惡行，他必死在罪孽之中，你卻救了自己的命。 \end{tabularx} \\ \\ \relax
3:20 & \begin{tabularx}{0.7\textwidth}{X} 但是義人若轉離他的義而作惡，我要把絆腳石放在他面前，他必死亡；因你沒有警戒他，他必死在罪中，他素來所行的義不被記念；我卻要從你手裡討他的血債。 \end{tabularx} \\ \\ \relax
3:21 & \begin{tabularx}{0.7\textwidth}{X} 倘若你警戒義人，使他不犯罪，他就不犯罪；他因領受警戒就必存活，你也救了自己的命。」 \end{tabularx} \\ \\
[1ex]
\hline
\hline
\end{longtable}
基督徒的生活,是本乎信仰,因為信仰帶給我們生命－接受了救恩,就有生命,就有力量,有力量就有生活.神的原則就是這樣,祂頒下的十誡前一半是論到信仰,後一半是論到生活,在新約主耶穌論到誡命中最重要的乃是:「要盡心,盡性,盡意,盡力,愛神,其次要愛人如己.」保羅寫羅馬書指出神的忿怒,從天上顯明在一切不虔不義的人身上,可見沒有信仰便沒有生活.
\newline
\newline
我相信赴會的各位,多數是已經重生得救的,心中也盼望別人得救；正如家中添了丁一樣,全家都充滿喜樂,這是很自然的流露.倘若你信了主,而你的家人竟不知道的,恐怕你的生命就有問題,因為你活不出見證來.我們的見證若靠咀巴講是沒有用的,必須在生活中活出來吸引人,以致感動人歸主使榮耀歸給神!
\newline
\newline
剛才看過的經文論到以西結,他在被擄到巴比倫的猶大人中作先知,巴比倫是外邦拜偶像的地方,是他們的敵國.他就在這不自由,受逼迫的環境中被神揀選作先知.因此我們想到在自由環境裡而不事奉神,不願走屬靈的道路,至終必要後悔.參路加十六章財主和拉撒路的故事,財主在陰間,哀求亞伯拉罕打發人去他家中傳道,可惜他的後悔已經來不及了.如今我們雖不會因局面驟然轉變,不能傳福音而後悔；卻會因環境上的轉變－例如搬遷,疾病,死亡,而失去傳福音的機會.求主幫助我們,要把握時機為祂工作.保羅說道:「無論是希利尼人,化外人,聰明人,愚拙人,我都欠他他們的倩,若不傳福音我們便有禍了.」所以我們若不趁機會向人傳福音,就不但對人有虧欠,也對神有虧欠,對曾經向我們傳過道的人也有虧欠.因為希伯來書十三章告訴我們:「從前引導你們,傳神之道給你們的人,你們要想念他們,效法他們的信心,留心看他們為人的結局.」主已把責任交付我們,我們若不傳福音,就有禍了.
\newline
\newline
近代社會生活趨於商業化,連傳福音工作也趨於商業化.初期的教會,工作開始時只有很少的人,他們得不著人的同情,也沒有人幫助,工作孤單,四面受敵,他們能向主盡忠,主就將得救的人賜給他們.五旬節後,信主的人多起來了,就需要有行政的組織；但日子久了,起初的愛心和熱心,就漸漸消失,行政組織成為屬靈的累贅.以至日後的教會一開始工作,便進行組織,換言之,以組織來開始工作.須知歸信主的人,當他們信了主後,在主的帶領下,便會合而為一.所以不論個人去傳道或由教會派人出去傳道,動機都應該是為愛主；不應為了組織而工作,也不應為維持組織而工作.
\newline
\newline
有時我們會太注重人的需要,認為人落在可憐的光景,犯罪作惡,因此我們要拯救他們；不錯,這動機是對的,但只不過是一部份.若是我們單為這些人的需要而傳福音,就會失去了工作的意義.猶太人昔日拒絕主耶穌,認為他們不需要福音,他們看保羅如同渣滓,希臘人以為自己有哲學,所以拒絕福音.人根本不知道自己的需要,即使知道,也不肯接受拯救,這是可嘆的.我們務要記得,傳福音是為滿足主的心,並不是為工作而工作,也不是為工作的對象,而傳福音.若然,動機就錯,屬靈道路就不正確了.從前在英國有一道河,常常有人在那裡跳河自殺.人們看見這情形,認為有防止自殺的需要,便廣泛宣傳.結果,社會人士同情,捐助,買了一艘船整天在河上巡邏拯救.這個開始的動機原來不錯,但因為有了組織,捐錢也多起來,有救生船的設備,就沒有人跳河自殺了.於是整個組織就耽心起來,認為這樣下去便無可宣傳,不能再向人捐錢!組織機構也不可能存在了,他們就巴不得有人再跳河,好讓他們有工作可做,可以維持這個組織.所以我們傳福音的動機必須要正確.多少傳福音的機構變了質,因為純正信仰不受人歡迎,原來信仰純正的教會也不傳純正的信仰了.耶穌是神的兒子,是童女所生,被釘在十字架上,為我們流血贖罪,這些信息既是不受歡迎的,便不想再傳.主的復活,再來,審判活人,死人,更不想傳了.但工作不能不做,產業不能不保留,組織不能不存在；於是將它變成社會工作的一部份,裡面變了質沒有耶穌了.在另一個情況下,是為了填寫報告而工作,盡量多開新的工作；藉此可以多寫報告,收獲更多捐款.雖然,傳福音是純正的,可惜內部精神已變質不純正.以前的世代是以假神亂真神,現在的世代以無神亂有神,教會變了,本身也不知道!因此必須求主憐憫幫助,讓我們傳福音不僅是為了適應人的需要,更是為了滿足主的心.
\newline
\newline
我們為主作見證,不一定得到對方的心,只管講主在我們身上的作為；對方受感動與否是另一回事.當然講台上是要顧及聽眾的需要,要有供應,餵養,而最重要的是要把神的心意傳講出來.耶穌講道時常說:神如此說,在講台上是神在那裡發言,而不是表現我們自己.我們為主工作,必須要堅守對主的信仰；不要討好個人或是團體.所以我們應該時常檢討,對主的心是否真誠?若有一些差錯,就要立刻悔過,矯正.
\newline
\newline
論到為主傳道的方法,個人傳道比集會傳道更為收效.因為講台上有些信息,是一般性的,聽的人對所傳的也許會遺漏,或有疑問而又沒有機會發問；若個人對個人傳道,就方便很多了,因為能夠直接而深入.有人認為個人沒有機會為主作見證,這話是講不通的；登台講道的機會或許不易,可是對人作見證的機會隨時都有.只要我們是向主忠心,目標對著未得救的靈魂,神是會給我們力量的.假使我們對少數人不盡心盡力,祂就不會給我們有更大的機會.因此我們要把握機會,把得救,贖罪的道理,告訴鄰居,親友,也可以把福音單張派發給他們.我們要以主的事為念,常常記掛著應該為祂而工作；不要厭煩,不要怕難為情,即使受到不禮貌的待遇,也當忍受.派發一份單張似乎是很微小的事,但也得求主帶領.有一次我去派發單張,對方不但不要,還吐口沬在上面,我把它撿起來拭淨,讀了一遍；在路旁有人看見,就認為那吐口沫的人太無理了,我沒有不好的影響,反而使那人心中也受到感動.
\newline
\newline
我們為主工作,需要有極高度的忍受,因為不經過痛苦的,就不曉得去安慰人.有一次我手中的單張,派到最後的一張,本來想進入禮拜堂聚會了；但心中有感動要再派,就在那時候來了一個人,他接了單張進去,因此,便信了主.他原是碼頭工人,因患肺病得愈,而答謝醫生,醫生說你報答耶穌好了.於是他到禮拜堂想進去聽道,可是多次到禮拜堂,在門外徘徊,都沒有人理會他,現在竟因一張單張引他進禮拜堂而信了主.這固然是醫生的工作果效,也可說是我們對往來的人,慇勤態度的反應.傳道的工作如同接力賽跑,有人為他們禱告,有人帶領他們,有人向他們傳信息,終有一天他們會相信接受福音的.
\newline
\newline
傳福音的對像是罪人,對於罪人我們有錯覺的想法,以為罪人總是窮兇極惡的,事實上,罪人不一定生出一副兇惡相,就算是善良面孔的也是罪人.因為人一生下來就有了罪,人的外表是不能遮蓋罪的.許多人不喜歡提到罪,不滿意傳道人指責罪；魯迅曾寫過一篇文章,說到一家人生下一個孩子,大家圍攏著看這孩子,稱讚他.一個說他會長壽,一個說他會做官,一個說他會很聰明,另一個卻說他終有一天要死.做父母的聽了就不高興,原來說誠實話會得開罪人,虛偽恭維的話卻賺得人的喜歡.
\newline
\newline
我們是否怕得罪人,就不講罪呢?罪的事實是不能否認的.罪的代價乃是死,唯有到主面前來纔可以得赦免得拯救.我們若不趁機會傳福音,恐怕機會一過,便後悔莫及了,環境是隨時都會改變的,若不拚命為主工作,領人歸主,我們怎能坦然見主,向祂交帳呢?
\newline
\newline


\newpage

\section{第42屆~港九培靈會~黃聿侯~第一講}
\label{sec:1550}
\textbf{基督是永生神的兒子}
\newline
\newline
連結: \href{https://www.hkbibleconference.org/session-message/view/1550}{\texttt{https://www.hkbibleconference.org/session-message/view/1550}}
\newline
\newline
\hyperref[sec:1439]{< < < PREV SERMON < < <}
~
\hyperref[sec:toc]{[返目錄]}
~
\hyperref[sec:1433]{> > > NEXT SERMON > > >}
\newline
\newline
馬太福音 16:13-16:20
\newline
\begin{longtable}{cl}
\hline
\hline
章節 & 經文 (和合本修訂版)\\
\hline
16:13 & \begin{tabularx}{0.7\textwidth}{X} 耶穌到了凱撒利亞．腓立比的境內，就問門徒：「人們說人子是誰？」 \end{tabularx} \\ \\ \relax
16:14 & \begin{tabularx}{0.7\textwidth}{X} 他們說：「有人說是施洗的約翰；有人說是以利亞；又有人說是耶利米或是先知中的一位。」 \end{tabularx} \\ \\ \relax
16:15 & \begin{tabularx}{0.7\textwidth}{X} 耶穌問他們：「你們說我是誰？」 \end{tabularx} \\ \\ \relax
16:16 & \begin{tabularx}{0.7\textwidth}{X} 西門‧彼得回答說：「你是基督，是永生神的兒子。」 \end{tabularx} \\ \\ \relax
16:17 & \begin{tabularx}{0.7\textwidth}{X} 耶穌回答他說：「約拿的兒子西門，你是有福的！因為這不是屬血肉的啟示你的，而是我在天上的父啟示的。 \end{tabularx} \\ \\ \relax
16:18 & \begin{tabularx}{0.7\textwidth}{X} 我還告訴你，你是彼得，我要把我的教會建造在這磐石上，陰間的權柄不能勝過它。 \end{tabularx} \\ \\ \relax
16:19 & \begin{tabularx}{0.7\textwidth}{X} 我要把天國的鑰匙給你，凡你在地上所捆綁的，在天上也要捆綁；凡你在地上所釋放的，在天上也要釋放。」 \end{tabularx} \\ \\ \relax
16:20 & \begin{tabularx}{0.7\textwidth}{X} 當時，耶穌囑咐門徒不可對任何人說他是基督。 \end{tabularx} \\ \\
[1ex]
\hline
\hline
\end{longtable}
「基督是永生神的兒子」這句話是我們的信仰根基,初期教會的信仰也是建立在這句話上,使徒行傳第八2-4腓利向埃提阿伯太監傳道,太監就相信耶穌基督是神的兒子.現今在末世時候,異端興起,雖然有人自以為能持守純正,但因為世界逼迫教會,或教會在政府的支持下而工作,就會因此而漸漸變質.我曾經提醒一些青年信徒,閉起眼睛來講話,並不等於祈禱,甚至是失去與神交通.求主幫助,在這個世代回想當日的信仰,讓我們回到初期教會一樣,蒙受聖靈的光照,活出當日的光景,快快樂樂的把一切擺上,背起十字架而歌唱.
\newline
\newline
我們常說耶穌降生,祂是基督是永生神的兒子,但有時候我想應該反過來說:神的兒子降世,神立祂為基督,名為耶穌.基督是受膏者的意思,受膏是一種手續,目的是立祂為王,因此馬太二章引用舊約的話說:「將來有一位君王,要從那裡出來,牧養我以色列民.」(彌五2)與以賽亞書九章6節一樣的預言:「有一嬰孩為我們而生,有一子賜給我們,政權必擔在他的肩頭上.」耶穌的名,是要將百姓從罪惡裡救出來的意思,祂將要成為萬王之王,萬主之主.在當時的猶太人,都盼望有這樣一位的救主,因為當時他們所處的環境十分困難,由於違背神的旨意,結果被擄到巴比倫,被分散到世界各地,又被羅馬人所統治,非常可憐,故此盼望拯救更加迫切.按百姓的想法,當然盼望有一位能夠帶領他們復國的,所以按主降生的時候,博士來到耶路撒冷訪問,論到猶太人之王誕生,希律和全城的人卻不安,因為這位王既然誕生了,對於他的統治是十分嚴重的打擊.我們想到希律召齊祭司長文士問說:基督當生在何處?這句話包含很多意思,不但宗教方面,也關係政治方面,他們害怕猶太人會進行民族的革命.
\newline
\newline
關於耶路撒冷人之所以不安,我們應當研究路加所記載的報名上冊,就是現今所謂人口登記.凡舉行人口登記,正因為要應付非常局面.羅馬政府嚴查猶太人的戶籍,就可以控制他們.對於以色列人來說,是很大的侮辱,這事之前,曾有一班人因為反對報名上冊,有六千名法利賽人被釘十字架上.在宗教上神是不喜歡人數點祂的百姓,政治上猶太人也是反對人口登記的；因為這樣一來,他們會失去自由,會被徵兵役.這一次的人口登記,適藉博士又告訴他們有王出生,那麼他們會想到可能就會發生革命,而希律也可能要大規模屠殺了,故此耶路撒冷合城的人都不安──後來事實證明,凡兩歲以內的男孩果然被殺了.
\newline
\newline
「基督」兩個字很不容易釋明白,這不是個人的名乃是個職份,是神所膏立的.膏立是一種手續,目的乃要祂作王；裡面包含掌權,擔當政權,管理百姓的意思.在舊約先知的預言,往往把耶穌降生與祂再來會合在一起來講；因此耶穌降生的時候,人們就想到既然祂是擔當政權的,一定是會立刻做王.所以博士的問題是「那生下來作猶太之王的」,希律的問題卻是「基督當生在何處?」
\newline
\newline
我們的主耶穌,是在這個光景下降生的,而我們現今信耶穌,又是在甚麼氣氛下相信呢?為甚麼現今的教會,會缺少能力?我曾經到過印尼,一間禮拜堂有一二千信徒,聚會的時候坐得滿滿,但屋頂破爛也沒有錢修理,然而那裡的信徒,他們的住所多半很漂亮,好像哈該書說的:「這殿仍然荒涼,你們自己還住天花板的房屋.」那些基督徒想著到禮拜堂聚會是好的,可是不肯為做基督徒而付出一些代價.我又想到現今做傳道人,心中不是不迫切多得一些人歸主,然而人心剛硬.不過,歷史上從沒有說人心是柔軟的,只有主的僕人靠著主的力量纔能攻克人心,又增加自己的勇氣,更加願意付出力量來領人接受救主.
\newline
\newline
有些人以為把福音的標準降低,纔會使人容易接受,例如說認耶穌基督是神的兒子,信祂為我們釘死在十字架上,這樣的信仰太難,換過別的方法就會較易接受,但,如果人的信仰錯誤了,糾正就困難了.有人舉辦佈道會,覺得人很難接受永生的真理,因此改變語氣說:人信耶穌實際上是有好處的,如考試不及格是因為罪的攔阻,信主之後罪既除去一定及格的；這樣說就使人信主的觀念改變了,是為了考試及格而信主,若仍舊不及格,那人便埋怨了.也有人說信主便可保佑身體健康,疾病得醫治,生意發達,但這都是沒有聖經根據的.我要提醒大家,信主得救重生,肉體也得蒙救贖,是實在的.可是我們並不能以疾病得醫治為目的,因為信主之故而受辱,也是有可能的,像路加十九章記載的撒該一樣.可惜現今的教會,並不是那樣,怪不得傳福音沒有力量了.
\newline
\newline
主耶穌的降生,是為完成神的旨意,是要應驗舊約的預言.安得烈曾告訴彼得說:「我們遇見彌賽亞了.」(約一41)可見當時的百姓是多麼熱切盼望主的降生.不但如此,反對基督的希律也很注意祂的降生,因為魔鬼也知道聖經的預言必應驗,故此牠做更厲害的工作,牠要利用希律殺害主耶穌.主在世三十多年的過程中,魔鬼屢次利用猶太人的手傷害主；但一直不能成功,主被釘十字架,乃是主的時候到了,甘心捨命的.
\newline
\newline
我們所相信的耶穌基督,祂必再來,在教會中我們有沒有讓祂顯露呢?或許我們會想,主耶穌出現便好了,但要知道祂一出現,仇敵必然起來攻擊.祂降生的時候,希律就迫害祂:如果我們真要讓祂顯露,也必須為祂付出代價,抵擋不合真理的.我們生活在自由的環境裡,似乎看不見教會的逼迫,但從各方面看:哲學,新神學,物質的生活豈非要把基督徒的信仰改變麼?因此我們在這樣的壓力之下,並非強於猶太人在羅馬人的手下.因為我們是與空中掌權惡魔爭戰,若不得勝,便要敗於祂權下.但我們所相信的是擔當政權,統管萬有的主耶穌,在今天我們便要把權柄歸給祂,讓祂在我們的家庭,教會,個人中掌管權柄,承認祂永永遠遠作我們的主,我們的王.我們為這個信仰肯付出代價,肯認真對付,肯完全擺上,相信主必施恩,在我們身上顯露,滿有屬靈力量為祂作見證!
\newline
\newline


\newpage

\section{第42屆~港九培靈會~黃聿侯~第二講}
\label{sec:1433}
\textbf{人說人子是誰}
\newline
\newline
連結: \href{https://www.hkbibleconference.org/session-message/view/1433}{\texttt{https://www.hkbibleconference.org/session-message/view/1433}}
\newline
\newline
\hyperref[sec:1550]{< < < PREV SERMON < < <}
~
\hyperref[sec:toc]{[返目錄]}
~
\hyperref[sec:1440]{> > > NEXT SERMON > > >}
\newline
\newline
馬太福音 16:13-16:20
\newline
\begin{longtable}{cl}
\hline
\hline
章節 & 經文 (和合本修訂版)\\
\hline
16:13 & \begin{tabularx}{0.7\textwidth}{X} 耶穌到了凱撒利亞．腓立比的境內，就問門徒：「人們說人子是誰？」 \end{tabularx} \\ \\ \relax
16:14 & \begin{tabularx}{0.7\textwidth}{X} 他們說：「有人說是施洗的約翰；有人說是以利亞；又有人說是耶利米或是先知中的一位。」 \end{tabularx} \\ \\ \relax
16:15 & \begin{tabularx}{0.7\textwidth}{X} 耶穌問他們：「你們說我是誰？」 \end{tabularx} \\ \\ \relax
16:16 & \begin{tabularx}{0.7\textwidth}{X} 西門‧彼得回答說：「你是基督，是永生神的兒子。」 \end{tabularx} \\ \\ \relax
16:17 & \begin{tabularx}{0.7\textwidth}{X} 耶穌回答他說：「約拿的兒子西門，你是有福的！因為這不是屬血肉的啟示你的，而是我在天上的父啟示的。 \end{tabularx} \\ \\ \relax
16:18 & \begin{tabularx}{0.7\textwidth}{X} 我還告訴你，你是彼得，我要把我的教會建造在這磐石上，陰間的權柄不能勝過它。 \end{tabularx} \\ \\ \relax
16:19 & \begin{tabularx}{0.7\textwidth}{X} 我要把天國的鑰匙給你，凡你在地上所捆綁的，在天上也要捆綁；凡你在地上所釋放的，在天上也要釋放。」 \end{tabularx} \\ \\ \relax
16:20 & \begin{tabularx}{0.7\textwidth}{X} 當時，耶穌囑咐門徒不可對任何人說他是基督。 \end{tabularx} \\ \\
[1ex]
\hline
\hline
\end{longtable}
主在世上的時候,猶太人問施洗約翰他是否就是基督,他誠懇的回答不是；而他卻介紹耶穌出來,見證祂是神的羔羊,背負世人罪孽的.當時猶太人都渴望基督顯現,因為他們頭腦中,基督是神設立要祂作他們的王的；所以希律知道基督降生而大感不安,因為照預言說那生下來的基督,在以色列中要作王掌權.猶太人切望基督出現,因此曾有多人設法尋找祂,甚至有人假冒祂.在猶太歷史中曾有多次的暴亂,都是因有自稱為彌賽亞的出現.而施洗約翰很清楚自己不過是基督的先鋒,那位藉童貞女所生,名為耶穌的,纔是基督.所以他告訴他的門徒,他們便跟隨了祂,其中有安得烈,約翰等.
\newline
\newline
安得烈與主同住了一天,便發現施約翰介紹的那位真是基督,放是找著他的哥哥彼得,告訴他說我們遇見彌賽亞了(彌賽亞譒出來就是基督)不久,腓力找著拿但業,告訴他耶穌就是律法上所寫和先知所記的那一位.拿但業最初不相信,後來耶穌對他說腓力還沒有招呼你,你在無花果樹底上我就看見你了,主耶穌一講就講中了拿但業的心,因此他就承認耶穌是神的兒子,是以色列的王.
\newline
\newline
最初跟隨主的人,對祂都有很清楚的認識,因為當日的耶穌,日被人輕視的,在政治上也有很大危險；若承認祂是基督,就被視為革命份子,受到羅馬政府的監視,隨時有被拘捕的可能.希律王為要殺害耶穌,就連伯利恆兩歲以下的男童也殺死了.據歷史載,猶太人曾有兩次大規模的革命行動,一次有六千人被殺,另一次死了一千多人.當耶穌十二歲時候,曾有因為革命活動,而把所有木匠捉拿的事；照我猜想,耶穌家庭可能是製過十字架的.因為祂講的十字架不會是現今我們所欣賞,陳設,精美的那一種,而是隨時都把人釘上去的恐怖刑具.早期的門徒,都明白若跟隨耶穌便有這樣的遭遇.
\newline
\newline
在宗教上,耶穌自稱是神的兒子,但按猶太教義說,祂這樣自稱是褻瀆神的.神只有一位,怎可能有兒子?因此祂屢次被人用石頭打.現今宗教自由,我們可以信也可以不信；但當日的猶太人一生下來便是猶太教徒,若承認耶穌是神的兒子,便是違反了祖宗遺傳的信仰,要被人用石頭打死.所以當日的人相信主耶穌,是要付重大代價的,一信了祂,下一分鐘便要準備被釘十字架,被趕逐出會堂,被石頭打死,整個家庭被孤立.......然而他們的信仰非常堅定,所以見證有力量,可是現今我們的相信,卻為了要得主的好處,若得不著好處便發怨言,甚至主有好處給我們,不過延遲一點我們也不滿!祂若沒有屬世的好處給我們,便似乎要罷讀經,罷祈禱,罷奉獻,罷聚會,與祂對立.但當日的信徒,卻把自己的生命拿出來,任由人孤立,殺害!因為他們不注重外表的得著,不在乎主的祝福,卻在乎主自己.所信所傳者,就是神的兒子耶穌基督.
\newline
\newline
我們的主是豐富的,天地萬物都屬於祂,祂當然能夠給我們某些物質上的好處,然而,賺得全世界,賠上了自己的生命,有甚麼益處?我們所需要的應該不是主的祝福而是主自己,因為有了祂便有一切了.
\newline
\newline
馬太十六章載耶穌到該撒利亞腓立比境內,祂問門徒說:人說我人子是誰.他們中間,有人答:有人說是施洗的約翰,有人說是以利亞,又有人說是耶利米或是先知裡的一位；除了彼得,卻沒有人說祂是基督.這樣看來,門徒所得來的消息,比較一般猶太人所渴慕思想的基督都不如；因為博士報告猶太人的王降生,他們會引用經典證明耶穌是基督.對於施洗約翰的印象,我們想到分封王希律聽見耶穌所作的事,以為祂就是約翰復活.以利亞行神蹟,耶穌也行神蹟,耶利米忠心於工作,但見不到效果,耶穌也似乎是一樣──其實效果是在後來.摩西五經中早就說將來有先知興起像他,因此從摩西以後,猶太人日夕都盼望那先知的興起.
\newline
\newline
耶穌自降生至那時候已三十多年,不但與祂降生的時候那種氣氛不同,就算兩年前門徒跟隨主認耶穌是基督的光景也不同,那時候安得烈告訴彼得我們遇見彌賽亞了；約翰介紹門徒說這是神的羔羊,背負世人的罪孽的；拿但業說你是神的兒子是以色列的王.現在他們漸漸糢糊,只講聽自外人的話.只有彼得異於他們,他很突出的說:你是基督神的兒子.外人所講的怎樣?只要我們保持冷靜,是不會受影響的!可惜我們往往對教會灰心,只講外人之怎樣批評,自己卻無話可說,這是今天教會的致命傷,信仰到達這地步可以說壞到極.外人的批評是無可避免的,只怕他們不批評,我們便是與他們同化.但他們講他們的,我們應有基本的態度,是不能動搖的.約書亞臨終前,向以色列人鄭重表示:他們可以選擇所要事奉的神,至於我和我的家,我們是必定事奉耶和華.(書廿四15)他的態度堅決,不但沒有被百姓所拖垮,反之百姓回來跟著他.外面的輿論,正是我們信仰的考驗,我們應持定自己的意見,不要被外人的話所影響.
\newline
\newline
口快心直的彼得說:「是基督,是永生神的兒子.」耶穌便對他說:「西門巴約拿,你是有福的；因為這不是屬血肉的指示你的,乃是我在天上的父指示的.」主指示他,他纔會承認,斷不是他觀察而來.主的內心當時非常難過,因為外面的人對他估價低,以為祂是施洗的約翰,以利亞,耶利米或是先知裡的一位.但跟隨祂的人在兩年多以前就已經清楚,安得烈對彼得說:我們遇見彌賽亞了,腓力也曾對拿但業說:摩西在律法上所寫的和眾先知所記的那一位,我們遇見了.現今竟越來越糢糊!因此祂必須問,要他們清楚表示.我們的信仰是從聖靈來的,舊約時候靠信心接受預言,聖靈在人的心中作主,需要絕對順服.當日記下預言的人並沒有憑他們的經驗,與忠實的把所領受的寫下來,昔日的預言,後來都成為事實.現今我們比較當日有更多歷史資料,起碼耶穌那時候還未被釘未復活.所以彼得他們的信仰糢糊,主也諒解他們,沒有責備；如今我們卻有完整的聖經和聖靈的引導,比起前人實在有福多了.
\newline
\newline
信仰是聖靈的工作,與教條不同；教條在不信的人也可以照樣讀,因為它本身沒有能力.但信仰是有生命的,有能力的,它指導人的生活.我們對於主的信仰,不能在模稜兩可之間,不是用甚麼方法去研究；乃是藉祂的話來提供信仰,藉我們生活來順服見證.當日主對門徒所講的話,不一定立刻產生效果,好像彼得雖然有過承認祂是永生神的兒子的話,也似乎沒有效果,直到後來主耶穌復活升天,門徒被聖靈充滿,他們便開始傳耶穌是被立為主為基督,也對猶太人辨明祂是基督,萬膝都要跪拜,萬口都要稱頌.他們越傳越堅固.現今我們又經歷了千多年的歷史證實,聖靈不斷在人心中作主,因此我們的信仰是更加牢固,不管世人對祂怎樣看法,我們卻是按聖經所啟示聖靈所引導,在生活中見證祂的恩典；滿足主的心,也滿足我們內心的渴望.
\newline
\newline


\newpage

\section{第42屆~港九培靈會~黃聿侯~第二講}
\label{sec:1440}
\textbf{怎樣領人赴會}
\newline
\newline
連結: \href{https://www.hkbibleconference.org/session-message/view/1440}{\texttt{https://www.hkbibleconference.org/session-message/view/1440}}
\newline
\newline
\hyperref[sec:1433]{< < < PREV SERMON < < <}
~
\hyperref[sec:toc]{[返目錄]}
~
\hyperref[sec:1434]{> > > NEXT SERMON > > >}
\newline
\newline
詩篇 122:1-122:9
\newline
\begin{longtable}{cl}
\hline
\hline
章節 & 經文 (和合本修訂版)\\
\hline
122:1 & \begin{tabularx}{0.7\textwidth}{X} 我喜樂，因人對我說：「我們到耶和華的殿去。」 \end{tabularx} \\ \\ \relax
122:2 & \begin{tabularx}{0.7\textwidth}{X} 耶路撒冷啊，我們的腳站在你門內。 \end{tabularx} \\ \\ \relax
122:3 & \begin{tabularx}{0.7\textwidth}{X} 耶路撒冷被建造，如同連結整齊的一座城。 \end{tabularx} \\ \\ \relax
122:4 & \begin{tabularx}{0.7\textwidth}{X} 眾支派就是耶和華的支派，上那裡去，按以色列的法度頌揚耶和華的名。 \end{tabularx} \\ \\ \relax
122:5 & \begin{tabularx}{0.7\textwidth}{X} 他們在那裡設立審判的寶座，就是大衛家的寶座。 \end{tabularx} \\ \\ \relax
122:6 & \begin{tabularx}{0.7\textwidth}{X} 你們要為耶路撒冷求平安：「願愛你的人興旺！ \end{tabularx} \\ \\ \relax
122:7 & \begin{tabularx}{0.7\textwidth}{X} 願你城中有平安！願你宮內得平靜！」 \end{tabularx} \\ \\ \relax
122:8 & \begin{tabularx}{0.7\textwidth}{X} 為我弟兄和同伴的緣故，我要說：「願你平安！」 \end{tabularx} \\ \\ \relax
122:9 & \begin{tabularx}{0.7\textwidth}{X} 為耶和華－我們神殿的緣故，我要為你求福！ \end{tabularx} \\ \\
[1ex]
\hline
\hline
\end{longtable}
基督徒在地上有為主作見證的責任,不但自己相信,並且要讓人知道你是屬主的,帶領他們也相信.但,若只帶領人相信而不帶領他聚會,這樣就只做了一半.在他得救以後沒有加以教導,與主所教訓我們的:「凡我所吩咐你們的,都教訓他們遵守.」還未達到,還未曾把得救的人,造就成主的器皿,再去造就他人.
\newline
\newline
我相信在各人託負上,沒有不願意領人赴會的；但在帶領人去赴會的問題上有了問題,因此實行不來.這樣關係非常重要的,正如一個人患了病,若不及早治理,恐怕越來越嚴重,以致危害到他的生命.
\newline
\newline
我們領人聚會,必先領他去同往我們所要去的聚會,就如經常性的聚會.講道會帶他同去,祈禱會帶他同去,查經會也帶他同去,把聚會要看為生活的一部份.這是表明我們對於自己的聚會滿意,對於自己的聚會有信心.假如每次聚會四十分鐘,加上來回時間要三小時,一個禮拜赴三小時,一個禮拜赴三次會便要九個小時了.對於有工作的人,尤其在大專讀書的,這九個小時對他是有很大影響的,所以必須要使他在聚會中有所得著.覺得並不浪費,甚至認為再多用一點時間也有價值的.若然,他下次也喜歡跟你同去赴會.若赴會而無所得便難以有下次的機會了.詩人說:「人對我說.我們往耶和華的殿去,我就歡喜.」(詩一二二1)如果覺得聚會是一種享受,是一種福樂,便巴不得邀請人同去.許多信仰純正的人,很注意讀經祈禱,有愛心,但卻少有領人赴聚會的熱心.那些信仰不純的,只講自己好而反對別人的,這樣的人倒非常熱心,領人到他們的聚會.這是因為他們對自己的聚會,有信任心之故.另外有些人為了使他們聚會人數增多,會不擇手段去爭取；這是血氣的做法,也許一時很哄動,但到底是不能持久的.求主幫助我們能夠「把各人在基督裡完全的引到神面前」(西一28)當我們初信主的時候,領人赴會還有熱心,可是日子久了,便冷淡了,因為對聚會沒有新鮮的感覺的緣故.赴會變了循例式,如此怎能帶領人呢?倘若在教會裡沒有職位的話,便更閒散,去不去聚會也覺得無所謂,所以也不會帶領人.有人唱屬靈的高調,說赴會乃是看主,不是看人.在講道的時候,不把講員放在眼內,不理會講員講的是甚麼.對講員根本沒有靈裡的交道,屬靈的調子唱得越高,就越不會領人赴會了.
\newline
\newline
不能領人赴會的問題,既如上述.而我們怎樣去解決這些問題呢?有人說,人是屬血氣的,不喜歡走屬靈道路,所以就不願意參加聚會了.也有人說一天生活是那麼緊張,很難再安排時間去聚會.又有人認為外面的吸引太大,即如香港來說,報紙有很多家,而且篇幅又多；不但有國際新聞和社會新聞,特別吸引人的是長篇連載小說.喜歡小說的人天天追著那段小說看,到主日要赴會了,心中還思念著那回小說的下文.講員的信息和他頭腦中的小說起了衝突.另方面電視機和其他聲色的吸引,因此人就不去赴會了.
\newline
\newline
最大的問題,是到教會來沒有所得,這個問題除了我們個人靈性以外,不能不承認是由於教會對於恩賜的運用不夠分明.羅馬十二章和哥林多前書十二章,論恩賜都是專一的,當今二十世紀工業發達,科學昌明,人們都追求專業學問；在社會發展務要專業化,殊不知神在一千多年之前,已經說明專一的重要.這本聖經包括四十多位作者寫成,前後又歷時一千五百多年,所以纔能把真理透過他們表達.但現今的教會,只有一兩位傳道人,這樣,怎可以把聖經中的恩典,恩賜,恩惠都一一講明?一個教會有幾百個會友,卻只有一兩個傳道人,怎可能做好供應,滿足大家所需要?一個人的才能是有限的,不可能一人而兼具多種恩賜.這是可以想像得到的,假使一間學校,校長是他,校役也是他,怎可以把學校辦好?我們看哥林多書,羅馬書,至少見到八,九種恩賜,所以一間教會不可能單靠某一方面恩賜來維持,更不用說發展了.倘若只注重某一種恩賜,其他便會忽略；例如注重祈禱,聖經原則就會忽略了:注重傳信息,造就的功夫便不會持久了.
\newline
\newline
因此,巴不得主興起更多同工,有更多恩賜運用,而結合到實際問題.有人認為這樣講是多餘的,因為教會的經費不夠.以為供應一位傳道人已感困難,怎可以維持多位呢?但我們必須知道,經費的缺欠原因在那裡,當然教會不是要倚靠金錢；但作為教會中一份子是理當出錢出力的.以色列人在曠野時候,十二支派中的利未支派,是毋須做其他工作,而由十一支派供應的.到尼希米從巴比倫回來時,發覺利未人為了尋找生活,不能事奉神,因此他要各支派奉獻供應他們.這樣,傳道人是靠信徒供給,每個家庭都奉獻十分之一,十個家庭便可以供養一家傳道人；若連其他費用也算在內,二十家供應一家是毫無困難的.不錯,耶穌是勉勵跟從祂的人要背十字架,彼得,保羅也說要為主受苦.但,這是傳道人與主的關係,是他對主的心志.我們做信徒的不應該說,傳道人必須受苦,因為聖經教導我們:「在道理上受教的,當把一切需用的供給施教的人.」(加六6)我們理應供應他們的一切.有人認為供給傳道人是給他有中等的生活就夠了,這話並非全對,未知中等生活是用甚麼來作標準?在美國五百元的收入是中等,若拿來香港可能是高等了.所以應在所服事的人群中來定標準,把他與弟兄們生活拉成一致；恩賜的運用和彼此的隔膜,都易於溝通了.傳道人會視教會如家,信徒們到教會來也感到享受了,須知信徒們對教會是有奉獻責任的,我們不能把這責任放棄.
\newline
\newline
在教會裡需要有人做長老執事,分擔教會的工作.多少時選舉做長執的,以看他是否有錢來做標準；有錢的人一點也沒有錯,但只盼望他以金錢來支持工作是錯誤的.應該選取那些有才能的,讓他出來擔當傳道,彼此同心合意的把恩賜顯明,求主幫助我們,重視領人赴會,注意到幾方面:
\newline
\newline
(1)為所帶領的人禱告,我們對他說說,只是做了一半功夫；應該切切為他禱告,與他有較親近的接觸.這樣,邀請他赴會的時候,便有把握.自己的話語,像帶著權威,引起他對自己的信任,視我們如同他的師長,父母；他的家長,親屬對於我們就可以放心了.
\newline
\newline
(2)領他赴會的時間必須準確,約定了不應推翻取消,最好是提早一點帶領到會,以便介紹他與有關的人認識.若男女分座,更易使赴會的人與同性的接近有親切感.我們還要估計聚會進行中可能遇到的情形,例如聖餐的餅,杯,未受水禮的不應該領取；某事應該參加,某事不需要參加,都應該給他們先有心理上的準備.
\newline
\newline
(3)需要注意幫助赴會的人翻閱聖經.往往有人因為不會找聖經章節,下次便不願意赴會了.特別是在社會上有地位的知識份子；他找不到經節是多麼尷尬呢,我們應該注意幫助他.有陌生的人進來,我們要友善的招待他,不要光是招待相熟的,冷落那些陌生的,要使他進到聚會中感覺得有溫暖.假如那人來參加聚會幾次了,沒有聖經的,最好能買一本與聚會所用的相同版本聖經給他應用.
\newline
\newline
最後,帶領一個人來聚會,也要負責送他回去；尤其年青的學生,更應該這樣做.在路上我們可以和他談談聚會的領受；特別提出一點是自己想像中,認為他可能不明白的給他補課.這些都是比較具體可行的,願主施恩,讓我們這樣做,使神得著榮耀的,人得著幫助!
\newline
\newline


\newpage

\section{第42屆~港九培靈會~黃聿侯~第三講}
\label{sec:1434}
\textbf{教會建造在磐石上}
\newline
\newline
連結: \href{https://www.hkbibleconference.org/session-message/view/1434}{\texttt{https://www.hkbibleconference.org/session-message/view/1434}}
\newline
\newline
\hyperref[sec:1440]{< < < PREV SERMON < < <}
~
\hyperref[sec:toc]{[返目錄]}
~
\hyperref[sec:1441]{> > > NEXT SERMON > > >}
\newline
\newline
馬太福音 16:13-16:20
\newline
\begin{longtable}{cl}
\hline
\hline
章節 & 經文 (和合本修訂版)\\
\hline
16:13 & \begin{tabularx}{0.7\textwidth}{X} 耶穌到了凱撒利亞．腓立比的境內，就問門徒：「人們說人子是誰？」 \end{tabularx} \\ \\ \relax
16:14 & \begin{tabularx}{0.7\textwidth}{X} 他們說：「有人說是施洗的約翰；有人說是以利亞；又有人說是耶利米或是先知中的一位。」 \end{tabularx} \\ \\ \relax
16:15 & \begin{tabularx}{0.7\textwidth}{X} 耶穌問他們：「你們說我是誰？」 \end{tabularx} \\ \\ \relax
16:16 & \begin{tabularx}{0.7\textwidth}{X} 西門‧彼得回答說：「你是基督，是永生神的兒子。」 \end{tabularx} \\ \\ \relax
16:17 & \begin{tabularx}{0.7\textwidth}{X} 耶穌回答他說：「約拿的兒子西門，你是有福的！因為這不是屬血肉的啟示你的，而是我在天上的父啟示的。 \end{tabularx} \\ \\ \relax
16:18 & \begin{tabularx}{0.7\textwidth}{X} 我還告訴你，你是彼得，我要把我的教會建造在這磐石上，陰間的權柄不能勝過它。 \end{tabularx} \\ \\ \relax
16:19 & \begin{tabularx}{0.7\textwidth}{X} 我要把天國的鑰匙給你，凡你在地上所捆綁的，在天上也要捆綁；凡你在地上所釋放的，在天上也要釋放。」 \end{tabularx} \\ \\ \relax
16:20 & \begin{tabularx}{0.7\textwidth}{X} 當時，耶穌囑咐門徒不可對任何人說他是基督。 \end{tabularx} \\ \\
[1ex]
\hline
\hline
\end{longtable}
「我要把我的教會建造在這磐石上,陰間的權柄不能勝過他.」關於磐石的解釋,天主教說磐石就是彼得,這解釋錯誤很大.也有人解作耶穌自己,而按聖經上下文看,其正確意義,應該是形容教會要建造在穩固的基礎上面.這解釋關係我們的信仰很重要,現在分述如下:
\newline
\newline
天主教說希臘文磐石與彼得都同是一個字.因此18節說你是彼得,以後接著說我要把我的教會建造在這磐石上,便認為教會要建造在彼得身上,彼得就是教皇.但我們研究希臘文,彼得的名字,前面是有冠詞的,屬於陽性字；磐石是沒有冠詞的,屬於陰性字,所以這樣解釋是錯誤的.
\newline
\newline
天主教又說因為彼得認耶穌是永生神的兒子,主便給他三樣賞賜:
\newline
\newline
(1)教會建造在他的身上,他在那裡,教會就在那裡.
\newline
\newline
(2)把天國鑰匙給他,要開便開,要關便關.
\newline
\newline
(3)有捆綁或釋放的權柄,甚至後來的神父也有赦罪的權柄.
\newline
\newline
這是天主教的錯誤教義,認為人是靠行為得救,然而聖經告訴我們沒有人能夠靠行為得救的,若承認信仰也可以成為行為的一部份,接受救主也算為做了好事,這種說法是不通的.
\newline
\newline
基實,彼得在此承認耶穌是神的兒子,他不是第一個,因為早在兩年前,最少已經有兩三個人承認了,約翰一章35至51那一經文得清楚記載.若推想聖經以外的,博士朝見聖嬰的時候,希律和耶路撒冷合城的人都不安,便顯見那時已經有許多人認耶穌是永生神兒子的了.也許猶太人當時渴慕的基督是屬世化的,寄予政治期望的；由於祂不是力爭復國,甚至躲避人擁戴祂為王.因此人們纔對祂漸漸失望,連祂的先施洗約翰被下在監的時候,也對祂動搖,所以我們的信仰必須經得起考驗.
\newline
\newline
在彼得之前承認耶穌是基督的有施洗約翰,他介紹耶穌給門徒認識.因為在當時環境,大家對彌賽亞的觀念是要掌握地上政權的,若提及這稱謂,必須用一種代名詞,以免直接刺激羅馬統治者.約翰稱祂為神的羔羊,自己也答覆人所詢問,聲明他不是要來的基督,卻是另外一位.安得烈對彼得說我們遇見彌賽亞了,括號註明繙出來就是基督,這是後來約翰寫福音書的附註；而彌賽亞是舊約預言,說出來不會引起羅馬統治者當局太大反感.腓力對拿但業說是律法上所寫和眾先知所記的那一位.拿但業說拉比,你是神的兒子,以主耶穌為基督,就更為明顯了.所以天主教以為彼得認信有獎,那未免太天真了.
\newline
\newline
若說彼得就是教皇就更不通,我們根據歷史,羅馬教會成立之前,彼得來沒有去過那裡,教會向來都是以主耶穌為元首的,各人是肢體.只要信主得救重生,不管生在那裡都同為教會一肢體,甚至有沒有名字記在地上教會名冊也不緊要,有真的重生生命便成為教會的一份子.而天主教則重視教會的組織,認為有教皇纔有教會,所以把彼得抬舉作教皇.把他們自己奉獻給教會,乃是為教會組織而生活,並不是為了事奉主而奉獻.同時,這樣說法,在因信稱義的救恩是不合的,靠行為得救就不是靠聖靈更新了.
\newline
\newline
如今一般性的講法,以為磐石是指耶穌自己,這樣的解釋,可能是因為根據聖經串珠或者經文彙編而來.但串珠聖經和經文彙編對聖經的解釋並不絕對可靠,例如林前十章4節的靈磐石,是指當日摩西在曠野擊打的磐石,不可能說那石便是基督.況且這裡十六章講的磐石並不是比喻,乃是形容上文我的教會,要建造在像磐石一般的穩固基礎,與此類似的如馬太七房子建造在磐石上,不過那裡是明明的比喻.所以我們從上下文觀察,乃是形容教會的建造.
\newline
\newline
彼得承認耶穌是基督,也是我們信仰的基礎.初期教會傳道工作,就以這個信念來做標準,腓利對埃提阿伯太監講道,太監願意受洗的時候,也承認說:我信耶穌基督是神的兒子(徒八37)保羅寫信給以弗所教會說被建造在使和先知的根基上,從來使徒們和普通書信都沒有一句說過彼得是教皇的.總之彼得承認耶穌是基督,是永生神的兒子,耶穌稱許他,說要祂的教會建造在這磐石上.意思是這個信仰如同磐石的穩固,永不搖動.無論過去,現在,將來,都是要有這樣的信仰,祂是神的兒子,是神所立的王,祂擔當政權,掌握一切,並不是一位普通的革命領袖.我們儘管來自不同的教會,有不同的背景,但有同一的信仰,主耶穌不僅是拿撒勒人,神設立的救主,更是神的兒子；祂道成肉身住在世間,後來被釘死,復活,升天,還要再來.這一個信仰是絕對的,正如門徒所宣告:「除祂以外,別無拯救」(徒四12)「只有一位神,在神和人中間,只有一位中保,乃是降世為人的基督耶穌」(提前二4)
\newline
\newline
我們所信仰只有一位神,神人中間只有一位中保耶穌基督.而世人的信仰是分為有神和無神兩大範圍,並不是按種族或按地理,因為任何國度都有信與不信兩種人的.世上絕少人信無神,即使他們口頭強調無神而內心是很畏懼神的.在信神的人中間,有信泛神的,甚麼東西,甚麼地方都有神.有信獨一神的,例如猶太教,回教,希臘正教,天主教,基督教.但在所信獨一神裡面也有分別,是有救恩與沒有救恩,回教是不談救恩的,天主教對救恩很模糊；但我們相信有救恩,這救恩是從主耶穌而來,祂是神所設立的救主,在十字架上為我們成就了救贖.我們在主面前認罪悔改,祂赦罪之恩便臨到,獲得得救的生命；因此我們不但接受這救恩,並要把這救恩傳出去,使許多的人都像我們一樣蒙受主的恩典.倘若我們只是參加教會聚會而沒有認罪悔改,便得不著救恩了,與主耶穌仍是無份無關.所以基督徒的信仰要如同磐石,絕不受時間,人物,思想所動搖的.而在不動搖的信仰根基上面,便更有力量為主作見證,在世上為鹽為光榮耀主了.
\newline
\newline


\newpage

\section{第42屆~港九培靈會~黃聿侯~第三講}
\label{sec:1441}
\textbf{略談讀經}
\newline
\newline
連結: \href{https://www.hkbibleconference.org/session-message/view/1441}{\texttt{https://www.hkbibleconference.org/session-message/view/1441}}
\newline
\newline
\hyperref[sec:1434]{< < < PREV SERMON < < <}
~
\hyperref[sec:toc]{[返目錄]}
~
\hyperref[sec:1435]{> > > NEXT SERMON > > >}
\newline
\newline
詩篇 119:89-119:96
\newline
\begin{longtable}{cl}
\hline
\hline
章節 & 經文 (和合本修訂版)\\
\hline
119:89 & \begin{tabularx}{0.7\textwidth}{X} 耶和華啊，你的話安定在天，直到永遠。 \end{tabularx} \\ \\ \relax
119:90 & \begin{tabularx}{0.7\textwidth}{X} 你的信實存到萬代；你堅立了地，地就長存。 \end{tabularx} \\ \\ \relax
119:91 & \begin{tabularx}{0.7\textwidth}{X} 天地照你的典章存到今日；萬物都是你的僕役。 \end{tabularx} \\ \\ \relax
119:92 & \begin{tabularx}{0.7\textwidth}{X} 我若不以你的律法為樂，早就在苦難中滅絕了！ \end{tabularx} \\ \\ \relax
119:93 & \begin{tabularx}{0.7\textwidth}{X} 我永不忘記你的訓詞，因你用這訓詞將我救活。 \end{tabularx} \\ \\ \relax
119:94 & \begin{tabularx}{0.7\textwidth}{X} 我是屬你的，求你救我，因我尋求了你的訓詞。 \end{tabularx} \\ \\ \relax
119:95 & \begin{tabularx}{0.7\textwidth}{X} 惡人等著要滅絕我，我卻要揣摩你的法度。 \end{tabularx} \\ \\ \relax
119:96 & \begin{tabularx}{0.7\textwidth}{X} 我看萬事盡都有限，惟有你的命令極其寬廣。 \end{tabularx} \\ \\
[1ex]
\hline
\hline
\end{longtable}
詩篇裡最長的一篇是一一九節,共一七六篇,分為廿二段,每段八節.這篇最長的經文就是講論聖經本身的,它用不同的詞彙,如神的律法,法度,典章,命令,訓詞......來說明神的話語.多麼感謝主!我們人竟得著神把祂自己的話交給我們,意思就是祂把自己交給我們,正如我們最知心的人把最要緊的話,單獨告訴我們一樣.昔日保羅講過猶太人自誇有福氣,因為神把舊約交給他們.但,我們這些被猶太人不看在眼內的外邦人,神竟然把新舊的全部都給與；同時,感動人的聖靈又住在我們裡面,所以我們的福氣是比猶太人大得多了.猶太人很殷勤讀聖經－雖然只是聖經的一部份.我在耶路撒冷看那些拉比,見到他們在烈日蒸曬之下；仍坐在哭牆那邊讀得津津有味,我們豈不更應該殷勤讀經麼,基督徒的行事為人,要著他是否從聖經消化吸收而衡量的；若然,纔有知識,有生命,有力量.
\newline
\newline
我們的中華文化,就是孔子思想.若問甚麼是中國人?不應從種族或疆界來劃分,乃應該看他是否有中華文化的認識；如此英國蕭伯納也可看為中國人,因為他佩服中華文化,說世界第一等公民是有中華文化的人.同樣,甚麼人是基督徒?不應該絕對看他是受了水禮領受聖餐而論定,乃要看他有沒有主的話,在他生命中生活中流露而確定.
\newline
\newline
提摩太後書三章16節:「聖經都是神所默示的」意思是神用一口氣把它吹成.因為祂感動了不同時代,不同背景,不同階層的人物；歷時一千五百多年而寫成了這有系統有中心的聖經.保羅那時所指的聖經,還未把新約編選在內,但彼得寫彼得後書時,便認定使徒書信也是神所默示的,它的地位價值,與舊約聖經同列.我們研讀這最長的一篇詩篇,看見神如何教導我們去遵行順服,以致把祂傳揚.因為在這個過程裡,是包括研究分析,若光是在頭腦中作為知識探討；而不是去遵行於生活上,永遠得不著主的心,也失去祂感動人寫成聖經的目的.
\newline
\newline
舊約時代的先知,他們是非常謹慎遵行神話語的,神怎樣吩咐,他們便怎樣做.例如耶利米,他本身是一位愛國人士,但當神的話臨到,要他告訴百姓悔改,否則便要被仇敵擄去,倚靠埃及定必失敗.他被同胞誤會,以為他是奸細,擾亂人心.你含冤受屈,向神忠心,只有接受神的靈所感動,遵行祂的付託,就不顧慮自己.保羅說:「現在活著不再是我了,乃是基督在我裡面活著.」整個人要讓主的話在他生活中顯明.神的話語不但在個人身上有重大力量,在國家也大有影響.主前六零六年的三幾年之間,在希西家王之後的嗎哪西,敬拜偶像,聖殿荒涼.後來繼位的約西亞王重新修葺聖殿,纔發覺有神的律法書在裡面；於是宣告遵行神的話語,使這個國家能夠一度復興過來.可見聖經對個人,對家庭,對國族,對教會都產生極大力量,是我們不容置疑的.
\newline
\newline
人類的軟弱,雖然一度復興,不久又會故態復明.但感謝主的恩典,甚麼時候渴慕遵行祂話語的,便可以重新得力.昔日百姓被擄到巴比倫,便非常渴想神的說話,如今有些不自由的地區,已沒有聖經可看,但仍靠著以前記憶的經節得安慰.可惜我們在有自由讀聖經的環境,卻對於讀經不感興趣,不受感動,我們會推說生活緊張沒有時間去讀,其實我們每天花在看報祇,看其他刊物或浪費於別的事上,所佔的時間非常多.基督徒的靈性測驗,要看他每天用多少時間讀聖經,清晨起來,是先看報紙還是先看聖經?往往我們不是利用精力旺盛的時間來讀,卻在筋疲力竭之餘,勉強讀一段來求心理上的安慰.須知「時間就是生命.」消失一段時間便是消失一部份生命.尤其過了四十歲的人,把時間浪費,簡直就像讓毒瘤在那裡發展,多麼令人警惕!
\newline
\newline
讀經究竟容易嗎,有主的靈帶領引導,是容易的；沒有罪的阻隔,遵行也不困難.若我們要研究到它寫作背景文法等等,恐怕一個相當聰明的人,用相當適宜的方法,就算是二三十年,也不可能懂得很多.事實上沒有一個人對聖經完全懂得的.聖經的寫作過程是複雜的,地理包括非常廣,巴勒斯坦,埃及,巴比倫,希臘,羅馬都在其內；當時世界的變動情形,也隱隱約約在聖經上有線索.聖經原來的文字很舊,猶太人思想又非常深刻,而離我們現今時間,地理都相當遙遠,要把它仔細研究,更不是只讀一本聖經就可以.但研究這些與否都不重要,卻要重視神自己話語,祂的話語安定在天,永不改變.尤其祂的預言,太奇妙了!聖經中許多預言,至今已經成為歷史事實,未來的將要陸續見諸事實.我們可以說現在是站立在預言之上,因為時間一分一秒的過去,神的預言就是在那裡實現,祂的應許就像在我們腳下幾成事實,所以我們的福氣是多麼大!
\newline
\newline
人算甚麼,不過是泥土,像蟲類一樣的軟弱無用.並且因為犯罪,是罪該萬死的,然而神竟然將祂的話語託付,並且成為事實.耶穌是神的兒子,而道成肉身與我們同在,我們一信了主就有了道－神的話.我們接受祂,這道便佔據了我整個人,把我這個人都消化在祂裡面.到祂顯現的時候,我們必要像祂,這福氣是很難形容盡致的.我們不過是無用的瓦器,但寶貝放在裡面；它的能力,它的價值就完全不同.我們這沒有用的身體,聖靈竟然也住在裡面,祂在我們裡面並不是分開乃是溶合；當我們順服祂越多,溶化的量越大.換言之,我們順服主,主的話就把我們吞進去,所以我們平常讀經,順服的態度是非常重要的.猶太人以順服神為最重大的事,他們從巴比倫被擄歸回；以斯拉將律法書帶到會眾面前宣讀,他們聽了都說「阿們」.他所宣讀神的說話,不盡是勉勵,並且有許多責備,他們卻因這些話語而感動戰兢,樂意遵行.主在地上生活的時候,雖然滿有權柄能力,也是絕對順服神的說話,就如祂被捉拿的時候,祂的門徒拔刀相救；但祂說,「收刀入鞘罷,凡動刀的必死在刀下.你想我不能求父,現在為我差遣十二營多天使來麼；若是這樣,經上所說事情必須如此的話,怎麼應驗呢.」(太廿六)所以保羅在哥林多前書十五章說:「基督照聖經所說為我們的罪死了,尚且埋葬了,又照聖經所說,第三天復活了.」這是神的信實,也是祂話語不可變更的證明.
\newline
\newline
神向我們是這樣負責任,我們向祂順服是不用懷疑的；主的話語很明確,雖然深奧卻也易了領悟,許多人讀書都知道有精讀或略讀的方法,我們讀聖經的缺點是略讀得不夠,精讀又不全面.有人只是拿一兩節深入去鑽.由於全面的知識不廣；我們所得著的便零零星星,對於神的整個計劃,心意,氣魄都是很零碎的概念.盼望我們要從聖經創世記至啟示錄快快的讀一遍,放膽去讀,不用做筆記,也不用追求亮光.也許有人認為這樣略讀記憶不牢,不會的,我們如每日去讀,就能夠養活靈命,會不自覺的使靈性長進.若是那天未曾讀經,便惘惘然有所失,這就是靈性進步的狀態.這樣開始的時候真的似乎無所得.但習慣了,在談話中,思想中,以至講台上與講台下都會有了交通.在你裡面是有東西的,到了多讀之後,每卷書就會有印象,這些書卷內容特點是甚麼,神最重要的旨意要求是甚麼,都能夠明白的了.其他一般性的知識,時間,地點,記得與否,並非重要.只要記取其中重要教訓靠著主的力量去遵行便可以了.讀新舊約的時候都一樣,先要打好略讀的基礎,對書卷有整個概念,然後找一卷書,用一年時間去精讀,對於該書的內容段落大意都比較深刻的默記.例如讀羅馬書.一庄三章是講論罪,四五章是講因信稱義,六至八章講聖靈工作,九至十一章論猶太人的事,十二章以下是信徒生活,我們都要緊記著.但同時也不放棄略讀,假如讀經的時候,感覺上中文之間很難懂,不必先去找註釋,或用外文比較,或參關中文譯本的小字注.要不同角度去默想,會比較注釋得著為多.當有了得著的時候,你會非常歡喜快樂,更引起了興趣.聖經是非常奇妙的,裡面滿了安慰,鼓勵,教訓,一生受用不盡,間題是我們怎樣去領受,遵行.
\newline
\newline


\newpage

\section{第42屆~港九培靈會~黃聿侯~第四講}
\label{sec:1435}
\textbf{生活在教會中的信仰}
\newline
\newline
連結: \href{https://www.hkbibleconference.org/session-message/view/1435}{\texttt{https://www.hkbibleconference.org/session-message/view/1435}}
\newline
\newline
\hyperref[sec:1441]{< < < PREV SERMON < < <}
~
\hyperref[sec:toc]{[返目錄]}
~
\hyperref[sec:1442]{> > > NEXT SERMON > > >}
\newline
\newline
馬太福音 16:13-16:27
\newline
\begin{longtable}{cl}
\hline
\hline
章節 & 經文 (和合本修訂版)\\
\hline
16:13 & \begin{tabularx}{0.7\textwidth}{X} 耶穌到了凱撒利亞．腓立比的境內，就問門徒：「人們說人子是誰？」 \end{tabularx} \\ \\ \relax
16:14 & \begin{tabularx}{0.7\textwidth}{X} 他們說：「有人說是施洗的約翰；有人說是以利亞；又有人說是耶利米或是先知中的一位。」 \end{tabularx} \\ \\ \relax
16:15 & \begin{tabularx}{0.7\textwidth}{X} 耶穌問他們：「你們說我是誰？」 \end{tabularx} \\ \\ \relax
16:16 & \begin{tabularx}{0.7\textwidth}{X} 西門‧彼得回答說：「你是基督，是永生神的兒子。」 \end{tabularx} \\ \\ \relax
16:17 & \begin{tabularx}{0.7\textwidth}{X} 耶穌回答他說：「約拿的兒子西門，你是有福的！因為這不是屬血肉的啟示你的，而是我在天上的父啟示的。 \end{tabularx} \\ \\ \relax
16:18 & \begin{tabularx}{0.7\textwidth}{X} 我還告訴你，你是彼得，我要把我的教會建造在這磐石上，陰間的權柄不能勝過它。 \end{tabularx} \\ \\ \relax
16:19 & \begin{tabularx}{0.7\textwidth}{X} 我要把天國的鑰匙給你，凡你在地上所捆綁的，在天上也要捆綁；凡你在地上所釋放的，在天上也要釋放。」 \end{tabularx} \\ \\ \relax
16:20 & \begin{tabularx}{0.7\textwidth}{X} 當時，耶穌囑咐門徒不可對任何人說他是基督。 \end{tabularx} \\ \\ \relax
16:21 & \begin{tabularx}{0.7\textwidth}{X} 從那時起，耶穌才向門徒明說，他必須上耶路撒冷去，受長老、祭司長和文士許多的苦，並且被殺，第三天復活。 \end{tabularx} \\ \\ \relax
16:22 & \begin{tabularx}{0.7\textwidth}{X} 彼得就拉著他，責備他說：「主啊，千萬不可如此！這事絕不可臨到你身上。」 \end{tabularx} \\ \\ \relax
16:23 & \begin{tabularx}{0.7\textwidth}{X} 耶穌轉過來，對彼得說：「撒但，退到我後邊去！你是我的絆腳石，因為你不體會神的心意，而是體會人的意思。」 \end{tabularx} \\ \\ \relax
16:24 & \begin{tabularx}{0.7\textwidth}{X} 於是耶穌對門徒說：「若有人要跟從我，就當捨己，背起自己的十字架來跟從我。 \end{tabularx} \\ \\ \relax
16:25 & \begin{tabularx}{0.7\textwidth}{X} 因為凡要救自己生命的，要喪失生命；凡為我喪失生命的，要得著生命。 \end{tabularx} \\ \\ \relax
16:26 & \begin{tabularx}{0.7\textwidth}{X} 人若賺得全世界，賠上自己的生命，有甚麼益處呢？人還能拿甚麼換生命呢？ \end{tabularx} \\ \\ \relax
16:27 & \begin{tabularx}{0.7\textwidth}{X} 人子要在他父的榮耀裡與他的眾使者一起來臨，那時候，他要照各人的行為報應各人。 \end{tabularx} \\ \\
[1ex]
\hline
\hline
\end{longtable}
我們的信仰不是口頭講講的說話,也不是寫為信經的文字,或是神學論文,乃是活在教會中的生活.在我們信仰裡面,有兩件事似乎被認為是宗教儀式的:一件是受洗(受浸),一件是擘餅(聖餐)但在我們信仰生活裡,這兩件事是很重要的,英文稱之為兩個無聲的傳道人.信仰是要藉著生活表明,也藉著生活來傳達,信仰生活的高峰乃是愛和死.基督是神的兒子,祂在死的時候就表明了.教會傳揚主的道,不但需要有聲的傳道人,並需要無聲的傳道人.有聲的傳道,只是向聽眾講,聽的人不一定會共鳴；而無聲的傳道,卻會使人有很深印象,受很大感動.不曉得我們看見有人接受水禮,是否歡喜到淚也滴下；倘若見到別人信主,在神,人面前見證歸主也不受感動的,恐怕他的靈性也有問題!因為在他心目中,教會多一個人也無所謂,少一個人也無所謂.或者他對於教會救靈的工作,根本沒有甚麼負擔；事實上,有人歸主並非傳道牧師的事,乃是關係整個教會的事.我們同為肢體,應該有說不出的喜樂.求主幫助,讓我們看見人若接受水禮,就是撒旦的失敗,那人已經脫離的轄制,與主同死同葬同復活了.所以路加十五章告訴我們:一個罪人悔改,在神的使者面前,也是這樣為他歡喜；倘若我們看見人歸主而不喜樂,不受感動,怎能使教會見證有力量?
\newline
\newline
現今教會有許多有聲的傳道,信徒識得許多道理,彼此也有交通,可是沒有行動表示.這樣,與初期的教會相差很遠,初期教會是有行動而沒有聲響的,不像近今的教會有許多議案,議決了卻不去實行.教會是需要有行動的無聲傳道,聖靈在教會所作的工就是這樣.許多時候教會少人或沒有人接受水禮,所以沒有喜樂,而這責任乃是在我們自己,因為沒有人去傳道－無聲的傳道,是我們自己敗壞自己.神在我們身上寄予很大的盼望,讓我們擘餅紀念主也是一個很重要的教訓.我們不要把它視為宗教儀式之一,若然,不會渴慕主,只有一種神祕感而沒有享受交通的樂趣；即使參加擘餅的聚會,也覺得是重擔,對他自己沒有一點益處.甚至破壞了整個屬靈肢體的交通,這樣,靈性便成了嚴重病態了.
\newline
\newline
當我們思想洗禮或者說是浸禮時,是要思想施洗約翰所傳的洗禮是甚麼,主耶穌所受的洗是甚麼,我們所受的洗又是甚麼?約翰所傳的是悔改的洗禮.保羅傳道的時候,有一個人名叫阿波羅,他很有智慧,有口才,然而只曉得約翰的洗禮；所以亞居拉夫婦請他到家中來,與他交通勸勉,多少時候教會出現異端,是因為沒有注意到傳道人所講是否對,或者知道他不對也沒有幫助他糾正,這是非常遺憾的事.約翰傳的是悔改的洗,從保羅吩咐他們再受洗的事看來,便知道這洗禮是不完全的,我們不是以洗禮來表示悔改赦罪,乃是根據羅馬書六章所指出的,要與主同死,同埋葬,同復活.人必須在重生得救之後纔予以施洗的.所以洗禮與完備的救恩無關,受了洗禮並不一定得救,得救了並不是馬上就施洗；我們之受洗,是在人在神面前作見證,也向撒旦宣告不再為自己活,而是與主同活,今後成為新造的人,一舉一動有新生的樣式.
\newline
\newline
耶穌的受洗,是盡諸般的義.當時施洗約翰出來傳道施洗,知道祂就是永生神的羔羊,所以主請求受洗的時候,約翰曾攔阻祂.但主說暫且許我,因為我理當這樣盡諸般的義.當祂受洗從水裡上來,天上有聲音說,這是我的愛子,我所喜悅的.經上有三次記載天上有聲音見證祂(路九,約十二),後兩次是與祂受死有關的.因此祂受洗表示祂要經過死,埋葬,復活,祂的死是驚天動地的,而為主殉道的人也同樣震動天地.希伯來書十一章記載那些忠心於主見證的人,如同雲彩一般的圍著.若有人爭主而死,相信同樣天上和地下都知道.亞伯雖然死了,卻因信仍舊說話.這是無聲的傳道,無聲的行動,我們不必要求鳴鑼響鈸來宣傳,若有真實行動為愛主而表現,這就是得大的力量了.我們要決心與主同死同埋葬同復活,不是一天的決定；乃是要日日都如此過生活,在我們心中的聲音要與天上的聲音起共鳴,面對死亡也如司提反,保羅一樣的歌唱.
\newline
\newline
擘餅是神與人所立最末後的約.聖經載神曾與人立了許多的約,但人是軟弱失敗的；不能向神守完全的約,並且守也是暫時的,而神是昨日,今日,直到永遠都不改變的.因此,以多變的人與不變的神很難保守約的堅立.主耶穌來世上,神立祂作基督,這是最後一次的約,假如人仍不歸信祂,就只有自取沉淪了.所以希伯來書一章提到:神在古時藉著眾先知,多次,多方的曉諭列祖,就在這末世藉著祂兒子曉論我們.主與我們立的新約,是用祂自己的血!祂說:這是我立約的血,為多人流出來,使罪得赦的.祂吩咐人如此行,為的是紀念祂.又說:從今以後,我不再喝這葡萄汁,直到我在我父的國裡,同你們喝新那日子.這一個特別的約,在祂是絕對沒有失信的,因此當我們擘餅紀念祂的時候,應該想到,祂是用自己的血把我們的救贖過來的；同時也想到自己悔改得赦免便大大的喜樂；好像猶太人想起出及,雖然在幾百年後也仍舊充滿了喜樂一樣.這個約是那麼堅定不改變,我們這些不堪不配的人若要自己負責是不可能保守的,但感謝主!祂保守我們,使我們失敗了又起來.所以我們紀念祂是多麼喜樂!紀念主要有彼此相思的深厚感情,沒有一件東西比我們與主相親為更密切.主從天上降到世上來,雖然祂的愛,祂的死不是許多人能夠明白,但祂為愛我們而付出生命,不管我們反應如何,這就是愛.主知道一切,也眷顧我們無微不至.在我們祈求以先,祂已經知道,為何有時祂像沒有聽禱告似的,其實,祂所答應我們,給予我們的,都是照著祂最好的旨意；若是我們不了解祂,就反而增加我們心中的痛苦了.當我們想起某人離了世界,為他追思,可能會掉下眼淚,因為那人活在你的心中；若與那死者無關的人,即使死者死得再悲慘,他也不會有傷感的.那麼,我們紀念主擘餅,是覺得祂與我們無關?抑是覺得祂活在我們心中?若是覺得祂活在我們心中的話,怎可這樣冷淡?主耶穌的身體是為我們捨的,血是為我們流的,因此我們學習祂,為這信仰作見證,傳這得救的福音,要影響現今的世代,也要影響及於將來的世代!
\newline
\newline


\newpage

\section{第42屆~港九培靈會~黃聿侯~第四講}
\label{sec:1442}
\textbf{學習禱告}
\newline
\newline
連結: \href{https://www.hkbibleconference.org/session-message/view/1442}{\texttt{https://www.hkbibleconference.org/session-message/view/1442}}
\newline
\newline
\hyperref[sec:1435]{< < < PREV SERMON < < <}
~
\hyperref[sec:toc]{[返目錄]}
~
\hyperref[sec:1443]{> > > NEXT SERMON > > >}
\newline
\newline
路加福音 11:1-11:13
\newline
\begin{longtable}{cl}
\hline
\hline
章節 & 經文 (和合本修訂版)\\
\hline
11:1 & \begin{tabularx}{0.7\textwidth}{X} 耶穌在一個地方禱告。禱告完了，有個門徒對他說：「主啊，求你教導我們禱告，像約翰教導他的門徒一樣。」 \end{tabularx} \\ \\ \relax
11:2 & \begin{tabularx}{0.7\textwidth}{X} 耶穌對他們說：「你們禱告的時候，要說：『父啊，願人都尊你的名為聖；願你的國降臨； \end{tabularx} \\ \\ \relax
11:3 & \begin{tabularx}{0.7\textwidth}{X} 我們日用的飲食，天天賜給我們。 \end{tabularx} \\ \\ \relax
11:4 & \begin{tabularx}{0.7\textwidth}{X} 赦免我們的罪，因為我們也赦免凡虧欠我們的人。不叫我們陷入試探。』」 \end{tabularx} \\ \\ \relax
11:5 & \begin{tabularx}{0.7\textwidth}{X} 耶穌又對他們說：「你們中間誰有一個朋友半夜到他那裡去，對他說：『朋友！請借給我三個餅； \end{tabularx} \\ \\ \relax
11:6 & \begin{tabularx}{0.7\textwidth}{X} 因為我有一個朋友旅途中來到我這裡，我沒有東西招待他。』 \end{tabularx} \\ \\ \relax
11:7 & \begin{tabularx}{0.7\textwidth}{X} 那人在裡面回答：『不要打擾我，門已經關了，孩子們也同我在床上了，我不能起來給你。』 \end{tabularx} \\ \\ \relax
11:8 & \begin{tabularx}{0.7\textwidth}{X} 我告訴你們，雖不因他是朋友起來給他，也會因他不顧面子地直求，起來照他所需要的給他。 \end{tabularx} \\ \\ \relax
11:9 & \begin{tabularx}{0.7\textwidth}{X} 我又告訴你們，祈求，就給你們；尋找，就找到；叩門，就給你們開門。 \end{tabularx} \\ \\ \relax
11:10 & \begin{tabularx}{0.7\textwidth}{X} 因為凡祈求的，就得著；尋找的，就找到；叩門的，就給他開門。 \end{tabularx} \\ \\ \relax
11:11 & \begin{tabularx}{0.7\textwidth}{X} 你們中間作父親的，誰有兒子求魚，反拿蛇當魚給他呢？ \end{tabularx} \\ \\ \relax
11:12 & \begin{tabularx}{0.7\textwidth}{X} 求雞蛋，反給他蠍子呢？ \end{tabularx} \\ \\ \relax
11:13 & \begin{tabularx}{0.7\textwidth}{X} 你們雖然不好，尚且知道拿好東西給兒女，何況天父，他豈不更要把聖靈賜給求他的人嗎？」 \end{tabularx} \\ \\
[1ex]
\hline
\hline
\end{longtable}
禱告是一項最重大的事,是我們一輩子的功課,因此我們必須學習禱告,過禱告的生活.
\newline
\newline
路加十一章記載門徒求主教導他們禱告,可見他們多麼渴望像主耶穌一樣的禱告.學習是必須的,但要注意我們學習的動機,不是學別人禱告的時間長,也不是學別人講一些屬靈的話語；當求靈性生命的長進,求與主的交通緊密.有湛深的屬靈生命,那就不愁禱告時間不長,也不愁禱告的詞句不美了.最怕是屬靈的生命淺薄,而套用別人的話為自己的說話,這樣的禱告是空洞的；祈禱的人辛苦,和他一同祈禱的人也感到辛苦.為甚麼許多人不肯參加祈禱會,就是因為怕了那些禱告說話不實際的人；他們好像年青人學老年人口吻,言中無物.也有人怕來聚會是因為提出祈禱的,都是消極性的事,有疾病,有困難,就來禱告仰望神；得平安,順利時就不來向神感謝.信徒有急事,半夜也來敲傳道人的門要求為他代禱.但事情好轉了,當然不會在夜跑來告訴傳道人,甚至見了面也不再提；這些人有困難便帶給主,有快樂就不與主分享.所以我們要學習,禱告不是為渴慕祈禱,乃是渴慕與主有更好交道；禱告如同談話,我們通過禱告能夠更明白主,也讓主了解我們.這樣,我們與主的交道便有深度,能夠經常禱告,成了自然的,習慣的人心中有說不出的快樂.
\newline
\newline
聖經說:「你們雖然不好,尚且知道拿好東西給兒女,何況天父,豈不將聖靈給求祂的人麼.」耶穌講這些話的時候,聖靈還未降臨,門徒心中也沒有聖靈.而主耶穌與父神交通,從來沒有隔膜,所以很自然很快樂.門徒們知道禱告的好處,但行不出來；他們可能見到主時常禱告,有時是跪在那裡,有時是安靜坐著；或許是慢慢的行走看來禱告,然而他們總不能夠像主那樣的禱告.因為他們心中沒有聖靈,羅馬書說:「況且我們的軟弱有聖靈幫助,我們本不曉得當怎樣禱告,只是聖靈親自用說不出來的歎息,替我們禱告.」到現今我們是在主的恩典之下,有聖靈在心中,因此再不像昔日門徒的禱告沒有力量.當禱告的時候,主的靈與我們的靈一起交通,所以禱告是毫無困難的.
\newline
\newline
主教導我們禱告的幾項原則
\newline
\newline
為神自己和為祂的名祂的國度,能力,旨意而祈求.許多人渴慕主不是渴慕主自己,而是渴慕祂的好處,禱告祂從祂那裡得好處.當然,主會給我們,但記得主的話說:「你們要先求祂的國和祂的義,這些東西都要加給你們了.」若沒有神的國和神的義,就算得著全世界有甚麼益處?主在地上與門徒一起的時候,從世人的眼光來說,他們的背景是不錯的.馬太是稅吏,彼得,安得烈是富有的漁夫.我在迦百農見到一座會堂,據說對面的房子就是彼得的,面積可不小；但這些門徒寧願放下一切跟從主,甚至被逐出會堂,被釘在十字架上也在所不惜!論到主耶穌,物質豐裕固然談不上,在當時宗教和政治圈子裡都沒有地位；甚至百姓也不了解祂.但祂所選召的,門徒竟然這樣愛祂,他們只愛祂自己,而不是愛祂所給的東西.
\newline
\newline
「願人都尊你的名為聖」這節經文時常被人提起,「聖」是最高峰,是絕對的,反面就是邪惡,這是無可比擬的.世上有人相信無神,有人相信有神,在信有神的人中間,有信多神,有信一神.基督教和天主教,回教,希臘教都是相信獨一真神的,除祂以外無別神,所以祂是至聖至尊的.
\newline
\newline
禱告不是對神說一大堆詞句,乃是一種心志的表現.神向我們講話是藉著聖經,我們向祂表達心意是藉著祈禱.因此當我們說:「願人都尊祂的名為聖.」就是向神表達出自己心志,是一心尊崇實行的表示,可惜我們寶在缺少這樣的禱告.我們看主耶穌來到父神面前祈求,祂不但尋求神的旨意與神的帶領,並且很鄭重的宣告心志:「倘若可行,求你叫這杯離開我,然而不要照我的意思,只要照你的意思.」祂能夠這樣講也就這樣行,難怪門徒渴望像祂那樣的禱告；倘若我們一起聚會時有這樣禱告,就十分甜蜜了.
\newline
\newline
現今是末世,魔鬼的工作在靈界裡有很大力量.我們必須認識在屬靈的鬥爭中,除了祈禱以外再沒有別的辦法了.例如許多地方,已把福音的門戶都關閉了,我們真沒有辦法幫助他們,但－仍可以為他們禱告.在啟示錄說到,滿了香的金爐,這金爐就是眾徒的祈禱.」我們的禱告,有耶穌做我們的大祭司,像燒香一樣香氣上騰,直達到神那裡.每個人的禱告都是這樣的上達,如同啟示錄八章所描寫的,那禱告能夠震動天地.我們現今的焦慮,認為完全沒有辦法可想的事,到那時候,神必會為我們一一成就.猶太人被擄到巴比倫,經過七十年,聖殿已被毀；文化也被改變,所過的是迦勒底人的生活方式.這七十年間已經生產了第二,三代.似乎一切已成定局,不可能有希望歸回祖國.要想革命嗎?要想反攻嗎?一點盼望也沒有,然而神使用古列王,竟然叫他們歸回,這不是甚麼人的力量也做不到的,乃是禱告的果效.我們可以關起門來禱告,把不可能改變的形勢按著神的時候,可以改變過來.彼得被囚在監內,教會為他切切祈禱,監獄的門自動的打開.神就把他釋放,所以我們不要灰心失望,現在想到有甚麼難能的事?祈禱交託就行了.
\newline
\newline
禁食祈禱可以說是拼命的禱告,主在世上時常常禁食禱告.門徒趕不出污鬼,主說:這一類的鬼,若不禁食禱告,他就不出來.門徒也時常禁食禱告,我們看安提阿的教會的先知和教師,他們在那裡事奉主是禁食禱告.聖靈差巴拿巴和保羅去傳道,也禁食禱告.哥林多前書七章還提到夫妻間要禁食禱告,纔可以分房.可見禁食禱告在信徒生活中是時常實行的.在舊約中也有禁食的事,這是人與神之間非常自由的交道,是人向神迫切祈求的表示.哈曼要迫害猶大人,要把他們滅絕；末底改沒有辦法,只好求助於王后以斯帖.於是以斯帖和她的宮女,末底改並猶大人就為此事禁食三晝三夜.撒迦利亞先知,為了國難,他曾多次的呼籲全國百姓同心禁食.真誠的禁食是需要的,但切勿失去內在的聖靈能力,以致成為形式,這是我們所當注意的.
\newline
\newline
以賽亞書五十八章是一段關乎靈性的記載,我們讀到:「他們天天尋求我,樂意明白我的道,好像行義的國民,不離棄他們神的典章；向我求問公義的判語,喜悅親近我.」這是多麼令人欣賞的語啊!看下去,他們追問神:「我們禁食,你為何看不見呢......」神就指責他們:「你們禁食的日子,仍求利益,勒逼人為你作苦工；你們禁食,卻互相爭競......這樣的禁食,豈是我所揀選使人刻苦已心的日子麼......」(賽五十八3-9)所以我們禱告不是命令我們的主怎樣怎樣做,乃是讓祂的旨意在我們身上成就.我們本來是祂的兒女,應該順從祂,不能強祂來成全我的意思.我的家鄉福州,譯聖經時候把僕人稱為奴才,因為是奴才地位,所以沒有權力影響主人的.不過我們常常不自覺而指揮主耶穌,幸而祂心地寬廣,不會計較,只要我們肯禱告,神必應允,呼求時候,祂就垂聽.
\newline
\newline
祈禱為何不能達到神面前,不但是因我們的信心不足,最主要的是有東西在那裡阻擋著.那些東西並非我們不知道,若我們真的還未覺察,神會饒恕；但實際上那些阻擋著我們與神交通的罪,我們不肯放下,不肯對付!換言之,我們並不是真的以祂為聖,以祂為最尊?故此,禱告便受攔阻,巴不得我們除去與祂的隔膜,讓我們暢快地與祂交談；如同兒女倚在父母的懷中,盡情享受祂的大愛.這種快樂,比甚麼都好,這種情份,也唯有與主毫無隔膜的人,才能領會.
\newline
\newline


\newpage

\section{第42屆~港九培靈會~黃聿侯~第五講}
\label{sec:1443}
\textbf{肢體的生活}
\newline
\newline
連結: \href{https://www.hkbibleconference.org/session-message/view/1443}{\texttt{https://www.hkbibleconference.org/session-message/view/1443}}
\newline
\newline
\hyperref[sec:1442]{< < < PREV SERMON < < <}
~
\hyperref[sec:toc]{[返目錄]}
~
\hyperref[sec:1436]{> > > NEXT SERMON > > >}
\newline
\newline
羅馬書 12:2-12:31
\newline
\begin{longtable}{cl}
\hline
\hline
章節 & 經文 (和合本修訂版)\\
\hline
12:2 & \begin{tabularx}{0.7\textwidth}{X} 不要效法這個世界，只要心意更新而變化，叫你們察驗何為神的善良、純全、可喜悅的旨意。 \end{tabularx} \\ \\ \relax
12:3 & \begin{tabularx}{0.7\textwidth}{X} 我憑著所賜我的恩對你們每一位說：不要把自己看得太高，要照著神所分給各人的信心來衡量，看得合乎中道。 \end{tabularx} \\ \\ \relax
12:4 & \begin{tabularx}{0.7\textwidth}{X} 正如我們一個身子上有好些肢體，肢體也不都有一樣的用處。 \end{tabularx} \\ \\ \relax
12:5 & \begin{tabularx}{0.7\textwidth}{X} 這樣，我們許多人在基督裡是一個身體，互相聯絡作肢體。 \end{tabularx} \\ \\ \relax
12:6 & \begin{tabularx}{0.7\textwidth}{X} 按著所得的恩典，我們各有不同的恩賜：或說預言，要按著信心的程度說預言； \end{tabularx} \\ \\ \relax
12:7 & \begin{tabularx}{0.7\textwidth}{X} 或服事的，要專一服事；或教導的，要專一教導； \end{tabularx} \\ \\ \relax
12:8 & \begin{tabularx}{0.7\textwidth}{X} 或勸勉的，要專一勸勉；施捨的，要誠實；治理的，要殷勤；憐憫人的，要樂意。 \end{tabularx} \\ \\ \relax
12:9 & \begin{tabularx}{0.7\textwidth}{X} 愛，不可虛假；惡，要厭惡；善，要親近。 \end{tabularx} \\ \\ \relax
12:10 & \begin{tabularx}{0.7\textwidth}{X} 愛弟兄，要相親相愛；恭敬人，要彼此推讓； \end{tabularx} \\ \\ \relax
12:11 & \begin{tabularx}{0.7\textwidth}{X} 殷勤，不可懶惰。要靈裡火熱；常常服侍主。 \end{tabularx} \\ \\ \relax
12:12 & \begin{tabularx}{0.7\textwidth}{X} 在盼望中要喜樂；在患難中要忍耐；禱告要恆切。 \end{tabularx} \\ \\ \relax
12:13 & \begin{tabularx}{0.7\textwidth}{X} 聖徒有缺乏，要供給；異鄉客，要殷勤款待。 \end{tabularx} \\ \\ \relax
12:14 & \begin{tabularx}{0.7\textwidth}{X} 要祝福迫害你們的，要祝福，不可詛咒。 \end{tabularx} \\ \\ \relax
12:15 & \begin{tabularx}{0.7\textwidth}{X} 要與喜樂的人同樂；要與哀哭的人同哭。 \end{tabularx} \\ \\ \relax
12:16 & \begin{tabularx}{0.7\textwidth}{X} 要彼此同心，不要心高氣傲，倒要俯就卑微的人。不要自以為聰明。 \end{tabularx} \\ \\ \relax
12:17 & \begin{tabularx}{0.7\textwidth}{X} 不要以惡報惡，眾人以為美的事要留心去做。 \end{tabularx} \\ \\ \relax
12:18 & \begin{tabularx}{0.7\textwidth}{X} 若是可行，總要盡力與眾人和睦。 \end{tabularx} \\ \\ \relax
12:19 & \begin{tabularx}{0.7\textwidth}{X} 各位親愛的，不要自己伸冤，寧可給主的憤怒留地步，因為經上記著：「主說：『伸冤在我，我必報應。』」 \end{tabularx} \\ \\ \relax
12:20 & \begin{tabularx}{0.7\textwidth}{X} 不但如此，「你的仇敵若餓了，就給他吃；若渴了，就給他喝。因為你這樣做，就是把炭火堆在他的頭上。」 \end{tabularx} \\ \\ \relax
12:21 & \begin{tabularx}{0.7\textwidth}{X} 不要被惡所勝，反要以善勝惡。 \end{tabularx} \\ \\ \relax
13:1 & \begin{tabularx}{0.7\textwidth}{X} 在上有權柄的，人人要順服，因為沒有權柄不是來自神的。掌權的都是神所立的。 \end{tabularx} \\ \\ \relax
13:2 & \begin{tabularx}{0.7\textwidth}{X} 所以，抗拒掌權的就是抗拒神所立的；抗拒的人必自招審判。 \end{tabularx} \\ \\ \relax
13:3 & \begin{tabularx}{0.7\textwidth}{X} 作官的原不是要使行善的懼怕，而是要使作惡的懼怕。你願意不懼怕掌權的嗎？只要行善，你就可得他的稱讚； \end{tabularx} \\ \\ \relax
13:4 & \begin{tabularx}{0.7\textwidth}{X} 因為他是神的用人，是與你有益的。你若作惡，就該懼怕，因為他不是徒然佩劍；他是神的用人，為神的憤怒，報應作惡的。 \end{tabularx} \\ \\ \relax
13:5 & \begin{tabularx}{0.7\textwidth}{X} 所以，你們必須順服，不但是因神的憤怒，也是因著良心。 \end{tabularx} \\ \\ \relax
13:6 & \begin{tabularx}{0.7\textwidth}{X} 你們納糧也為這個緣故，因他們是神的僕役，專管這事。 \end{tabularx} \\ \\ \relax
13:7 & \begin{tabularx}{0.7\textwidth}{X} 凡人所當得的，就給他。當得糧的，給他納糧；當得稅的，給他上稅；當懼怕的，懼怕他；當恭敬的，恭敬他。 \end{tabularx} \\ \\ \relax
13:8 & \begin{tabularx}{0.7\textwidth}{X} 你們除了彼此相愛，對任何人都不可虧欠甚麼，因為那愛人的就成全了律法。 \end{tabularx} \\ \\ \relax
13:9 & \begin{tabularx}{0.7\textwidth}{X} 那不可姦淫，不可殺人，不可偷盜，不可貪婪，或別的誡命，都包括在「愛鄰如己」這一句話之內了。 \end{tabularx} \\ \\ \relax
13:10 & \begin{tabularx}{0.7\textwidth}{X} 愛是不對鄰人作惡，所以愛就成全了律法。 \end{tabularx} \\ \\ \relax
13:11 & \begin{tabularx}{0.7\textwidth}{X} 還有，你們要知道，現在正是該從睡夢中醒來的時候了；因為我們得救，現在比初信的時候更近了。 \end{tabularx} \\ \\ \relax
13:12 & \begin{tabularx}{0.7\textwidth}{X} 黑夜已深，白晝將近。所以我們該除去暗昧的行為，帶上光明的兵器。 \end{tabularx} \\ \\ \relax
13:13 & \begin{tabularx}{0.7\textwidth}{X} 行事為人要端正，好像在白晝行走。不可荒宴醉酒；不可好色淫蕩；不可紛爭嫉妒。 \end{tabularx} \\ \\ \relax
13:14 & \begin{tabularx}{0.7\textwidth}{X} 總要披戴主耶穌基督，不要只顧滿足肉體，去放縱私慾。 \end{tabularx} \\ \\ \relax
14:1 & \begin{tabularx}{0.7\textwidth}{X} 信心軟弱的，你們要接納，不同的意見，不要爭論。 \end{tabularx} \\ \\ \relax
14:2 & \begin{tabularx}{0.7\textwidth}{X} 有人信甚麼都可吃；但那軟弱的，只吃蔬菜。 \end{tabularx} \\ \\ \relax
14:3 & \begin{tabularx}{0.7\textwidth}{X} 吃的人不可輕看不吃的人；不吃的人也不可評斷吃的人，因為神已經接納他了。 \end{tabularx} \\ \\ \relax
14:4 & \begin{tabularx}{0.7\textwidth}{X} 你是誰，竟評斷別人的僕人呢？他或站立或跌倒，自有他的主人在，而且他也必會站立，因為主能使他站穩。 \end{tabularx} \\ \\ \relax
14:5 & \begin{tabularx}{0.7\textwidth}{X} 有人看這日比那日強；有人看日日都是一樣。只是各人要在自己的心意上堅定。 \end{tabularx} \\ \\ \relax
14:6 & \begin{tabularx}{0.7\textwidth}{X} 守日子的人是為主守的。吃的人是為主吃的，因他感謝神；不吃的人是為主不吃的，他也感謝神。 \end{tabularx} \\ \\ \relax
14:7 & \begin{tabularx}{0.7\textwidth}{X} 我們沒有一個人為自己而活，也沒有一個人為自己而死。 \end{tabularx} \\ \\ \relax
14:8 & \begin{tabularx}{0.7\textwidth}{X} 我們若活，是為主而活；我們若死，是為主而死。所以，我們或死或活總是主的人。 \end{tabularx} \\ \\ \relax
14:9 & \begin{tabularx}{0.7\textwidth}{X} 為此，基督死了，又活了，為要作死人和活人的主。 \end{tabularx} \\ \\ \relax
14:10 & \begin{tabularx}{0.7\textwidth}{X} 可是你，你為甚麼評斷弟兄呢？你又為甚麼輕看弟兄呢？因我們都要站在神的審判臺前。 \end{tabularx} \\ \\ \relax
14:11 & \begin{tabularx}{0.7\textwidth}{X} 經上寫著：「主說，我指著我的永生起誓：萬膝必向我跪拜；萬口必稱頌神。」 \end{tabularx} \\ \\ \relax
14:12 & \begin{tabularx}{0.7\textwidth}{X} 這樣看來，我們各人一定要把自己的事在神面前交代。 \end{tabularx} \\ \\ \relax
14:13 & \begin{tabularx}{0.7\textwidth}{X} 所以，我們不可再彼此評斷，寧可決意不給弟兄放置障礙或絆腳石。 \end{tabularx} \\ \\ \relax
14:14 & \begin{tabularx}{0.7\textwidth}{X} 我憑著主耶穌確知深信，凡物本來沒有不潔淨的，除非人以為不潔淨的，在他就不潔淨了。 \end{tabularx} \\ \\ \relax
14:15 & \begin{tabularx}{0.7\textwidth}{X} 你若因食物使弟兄憂愁，就不是按著愛心行事。基督已經為他死，你不可因你的食物使他敗壞。 \end{tabularx} \\ \\ \relax
14:16 & \begin{tabularx}{0.7\textwidth}{X} 所以，不可讓你們的善被人毀謗。 \end{tabularx} \\ \\ \relax
14:17 & \begin{tabularx}{0.7\textwidth}{X} 因為神的國不在乎飲食，而在乎公義、和平及聖靈中的喜樂。 \end{tabularx} \\ \\ \relax
14:18 & \begin{tabularx}{0.7\textwidth}{X} 凡這樣服侍基督的，就為神所喜悅，又為人所讚許。 \end{tabularx} \\ \\ \relax
14:19 & \begin{tabularx}{0.7\textwidth}{X} 所以，我們務要追求和平與彼此造就的事。 \end{tabularx} \\ \\ \relax
14:20 & \begin{tabularx}{0.7\textwidth}{X} 不可因食物毀壞神的工作。一切都是潔淨的，但有人因食物使人跌倒，這在他就是惡了。 \end{tabularx} \\ \\ \relax
14:21 & \begin{tabularx}{0.7\textwidth}{X} 無論是吃肉是喝酒，是甚麼別的事，使弟兄跌倒，一概不做，才是善的。 \end{tabularx} \\ \\ \relax
14:22 & \begin{tabularx}{0.7\textwidth}{X} 你有信心，就要在神面前持守。人能在自己以為可行的事上不自責就有福了。 \end{tabularx} \\ \\ \relax
14:23 & \begin{tabularx}{0.7\textwidth}{X} 若有人疑惑而吃的，就被定罪，因為他吃不是出於信心。凡不出於信心的都是罪。 \end{tabularx} \\ \\ \relax
15:1 & \begin{tabularx}{0.7\textwidth}{X} 我們堅強的人應該分擔不堅強的人的軟弱，不求自己的喜悅。 \end{tabularx} \\ \\ \relax
15:2 & \begin{tabularx}{0.7\textwidth}{X} 我們各人務必要讓鄰人喜悅，使他得益處，得造就。 \end{tabularx} \\ \\ \relax
15:3 & \begin{tabularx}{0.7\textwidth}{X} 因為基督也不求自己的喜悅，如經上所記：「辱罵你的人的辱罵都落在我身上。」 \end{tabularx} \\ \\ \relax
15:4 & \begin{tabularx}{0.7\textwidth}{X} 從前所寫的聖經都是為教導我們寫的，要使我們藉著忍耐和因聖經所生的安慰，得著盼望。 \end{tabularx} \\ \\ \relax
15:5 & \begin{tabularx}{0.7\textwidth}{X} 但願賜忍耐和安慰的神使你們彼此同心，效法基督耶穌， \end{tabularx} \\ \\ \relax
15:6 & \begin{tabularx}{0.7\textwidth}{X} 為使你們同心同聲榮耀我們主耶穌基督的父神！ \end{tabularx} \\ \\ \relax
15:7 & \begin{tabularx}{0.7\textwidth}{X} 所以，你們要彼此接納，如同基督接納你們一樣，歸榮耀給神。 \end{tabularx} \\ \\ \relax
15:8 & \begin{tabularx}{0.7\textwidth}{X} 我說，基督是為神真理作了受割禮的人的執事，要證實所應許列祖的話， \end{tabularx} \\ \\ \relax
15:9 & \begin{tabularx}{0.7\textwidth}{X} 並使外邦人，因他的憐憫，榮耀神。如經上所記：「因此，我要在外邦中稱頌你，歌頌你的名。」 \end{tabularx} \\ \\ \relax
15:10 & \begin{tabularx}{0.7\textwidth}{X} 又說：「外邦人哪，你們要與主的子民一同歡樂。」 \end{tabularx} \\ \\ \relax
15:11 & \begin{tabularx}{0.7\textwidth}{X} 又說：「列邦啊，你們要讚美主！萬民哪，你們都要頌讚他！」 \end{tabularx} \\ \\ \relax
15:12 & \begin{tabularx}{0.7\textwidth}{X} 又有以賽亞說：「將來有耶西的根，就是那興起來要治理列邦的；外邦人要仰望他。」 \end{tabularx} \\ \\ \relax
15:13 & \begin{tabularx}{0.7\textwidth}{X} 願賜盼望的神，因你們的信把各樣的喜樂、平安充滿你們的心，使你們藉著聖靈的能力大有盼望！ \end{tabularx} \\ \\ \relax
15:14 & \begin{tabularx}{0.7\textwidth}{X} 我的弟兄們，我本人也深信你們自己充滿良善，有各種豐富的知識，也能彼此勸戒。 \end{tabularx} \\ \\ \relax
15:15 & \begin{tabularx}{0.7\textwidth}{X} 但我更大膽寫信給你們，是要在一些事上提醒你們，我因神所賜我的恩， \end{tabularx} \\ \\ \relax
15:16 & \begin{tabularx}{0.7\textwidth}{X} 使我為外邦人作基督耶穌的僕役，作神福音的祭司，使所獻上的外邦人因著聖靈成為聖潔，可蒙悅納。 \end{tabularx} \\ \\ \relax
15:17 & \begin{tabularx}{0.7\textwidth}{X} 所以，有關神面前的事奉，我在基督耶穌裡是有可誇的。 \end{tabularx} \\ \\ \relax
15:18 & \begin{tabularx}{0.7\textwidth}{X} 除了基督藉我做的那些事，我甚麼都不敢提，只提他藉我的言語作為，用神蹟奇事的能力，並神的靈的能力，使外邦人順服； \end{tabularx} \\ \\ \relax
15:19 & \begin{tabularx}{0.7\textwidth}{X} 甚至我從耶路撒冷，直轉到以利哩古，到處傳了基督的福音。 \end{tabularx} \\ \\ \relax
15:20 & \begin{tabularx}{0.7\textwidth}{X} 這樣，我立了志向，不在基督的名已經傳揚過的地方傳福音，免得建造在別人的根基上； \end{tabularx} \\ \\ \relax
15:21 & \begin{tabularx}{0.7\textwidth}{X} 卻如經上所記：「未曾傳給他們的，他們必看見；未曾聽見過的事，他們要明白。」 \end{tabularx} \\ \\ \relax
15:22 & \begin{tabularx}{0.7\textwidth}{X} 因此我多次被攔阻，不能到你們那裡去。 \end{tabularx} \\ \\ \relax
15:23 & \begin{tabularx}{0.7\textwidth}{X} 但如今，在這一帶再沒有可傳的地方，而且這許多年來，我迫切想去你們那裡， \end{tabularx} \\ \\ \relax
15:24 & \begin{tabularx}{0.7\textwidth}{X} 盼望到西班牙去的時候經過，得見你們，先與你們彼此交往，心裡稍得滿足，然後蒙你們為我送行。 \end{tabularx} \\ \\ \relax
15:25 & \begin{tabularx}{0.7\textwidth}{X} 但如今我要到耶路撒冷去，供應聖徒的需要。 \end{tabularx} \\ \\ \relax
15:26 & \begin{tabularx}{0.7\textwidth}{X} 因為馬其頓和亞該亞人樂意湊出一些捐款給耶路撒冷聖徒中的窮人。 \end{tabularx} \\ \\ \relax
15:27 & \begin{tabularx}{0.7\textwidth}{X} 這固然是他們樂意的，其實也算是所欠的債；因為外邦人既然分享了他們靈性上的好處，就當把肉體上的需用供給他們。 \end{tabularx} \\ \\ \relax
15:28 & \begin{tabularx}{0.7\textwidth}{X} 等我辦完了這事，把這筆捐款交付給他們，我就要路過你們那裡，到西班牙去。 \end{tabularx} \\ \\ \relax
15:29 & \begin{tabularx}{0.7\textwidth}{X} 我也知道去你們那裡的時候，我將帶著基督豐盛的恩典去。 \end{tabularx} \\ \\ \relax
15:30 & \begin{tabularx}{0.7\textwidth}{X} 弟兄們，我藉著我們的主耶穌基督，又藉著聖靈的愛，勸你們與我一同竭力為我祈求神， \end{tabularx} \\ \\ \relax
15:31 & \begin{tabularx}{0.7\textwidth}{X} 使我脫離在猶太不順從的人，也讓我在耶路撒冷的事奉可蒙聖徒悅納， \end{tabularx} \\ \\ \relax
15:32 & \begin{tabularx}{0.7\textwidth}{X} 並使我照著神的旨意歡歡喜喜地到你們那裡，與你們同得安息。 \end{tabularx} \\ \\ \relax
15:33 & \begin{tabularx}{0.7\textwidth}{X} 願賜平安的神與你們眾人同在。阿們！ \end{tabularx} \\ \\ \relax
16:1 & \begin{tabularx}{0.7\textwidth}{X} 我對你們推薦我們的姊妹非比，她是堅革哩教會中的執事。 \end{tabularx} \\ \\ \relax
16:2 & \begin{tabularx}{0.7\textwidth}{X} 請你們在主裡用合乎聖徒的方式來接待她。她在任何事上需要你們幫助，你們就幫助她；因她素來幫助許多人，也幫助了我。 \end{tabularx} \\ \\ \relax
16:3 & \begin{tabularx}{0.7\textwidth}{X} 請向百基拉和亞居拉問安。他們在基督耶穌裡作我的同工， \end{tabularx} \\ \\ \relax
16:4 & \begin{tabularx}{0.7\textwidth}{X} 也為我的性命把自己的生死置之度外；不但我感謝他們，就是外邦的眾教會也感謝他們。 \end{tabularx} \\ \\ \relax
16:5 & \begin{tabularx}{0.7\textwidth}{X} 又向在他們家中的教會問安。向我所親愛的以拜尼土問安，他是亞細亞歸於基督的初結果子。 \end{tabularx} \\ \\ \relax
16:6 & \begin{tabularx}{0.7\textwidth}{X} 又向馬利亞問安，她為你們非常辛勞。 \end{tabularx} \\ \\ \relax
16:7 & \begin{tabularx}{0.7\textwidth}{X} 又向與我一同坐監的親戚安多尼古和猶尼亞問安，他們在使徒中是有名望的，也是比我先在基督裡的。 \end{tabularx} \\ \\ \relax
16:8 & \begin{tabularx}{0.7\textwidth}{X} 又向我在主裡面所親愛的暗伯利問安。 \end{tabularx} \\ \\ \relax
16:9 & \begin{tabularx}{0.7\textwidth}{X} 又向我們在基督裡的同工耳巴奴和我所親愛的士大古問安。 \end{tabularx} \\ \\ \relax
16:10 & \begin{tabularx}{0.7\textwidth}{X} 又向在基督裡經過考驗的亞比利問安。向亞利多布家裡的人問安。 \end{tabularx} \\ \\ \relax
16:11 & \begin{tabularx}{0.7\textwidth}{X} 又向我親戚希羅天問安。向拿其數家在主裡的人問安。 \end{tabularx} \\ \\ \relax
16:12 & \begin{tabularx}{0.7\textwidth}{X} 又向為主辛勞的土非拿和土富撒問安。向所親愛、為主非常辛勞的彼息問安。 \end{tabularx} \\ \\ \relax
16:13 & \begin{tabularx}{0.7\textwidth}{X} 又向在主裡蒙揀選的魯孚和他母親問安，他的母親就是我的母親。 \end{tabularx} \\ \\ \relax
16:14 & \begin{tabularx}{0.7\textwidth}{X} 又向亞遜其土、弗勒干、黑米、八羅巴、黑馬，和跟他們在一起的弟兄們問安。 \end{tabularx} \\ \\ \relax
16:15 & \begin{tabularx}{0.7\textwidth}{X} 又向非羅羅古和猶利亞，尼利亞和他姊妹，阿林巴和跟他們在一起的眾聖徒問安。 \end{tabularx} \\ \\ \relax
16:16 & \begin{tabularx}{0.7\textwidth}{X} 你們要以聖潔的吻彼此問安。基督的眾教會都向你們問安！ \end{tabularx} \\ \\ \relax
16:17 & \begin{tabularx}{0.7\textwidth}{X} 弟兄們，那些離間你們、使你們跌倒、違背所學之道的人，我勸你們要留意躲避他們。 \end{tabularx} \\ \\ \relax
16:18 & \begin{tabularx}{0.7\textwidth}{X} 因為這樣的人不服侍我們的主基督，只服侍自己的肚腹，用花言巧語誘惑老實人的心。 \end{tabularx} \\ \\ \relax
16:19 & \begin{tabularx}{0.7\textwidth}{X} 你們的順服已經傳於眾人，所以我為你們歡喜；但我願你們在善上聰明，在惡上愚拙。 \end{tabularx} \\ \\ \relax
16:20 & \begin{tabularx}{0.7\textwidth}{X} 那賜平安的神快要把撒但踐踏在你們腳下。願我們主耶穌基督的恩與你們同在！ \end{tabularx} \\ \\ \relax
16:21 & \begin{tabularx}{0.7\textwidth}{X} 我的同工提摩太，和我的親戚路求、耶孫、所西巴德，向你們問安。 \end{tabularx} \\ \\ \relax
16:22 & \begin{tabularx}{0.7\textwidth}{X} 我這代筆寫信的德提，在主裡向你們問安。 \end{tabularx} \\ \\ \relax
16:23 & \begin{tabularx}{0.7\textwidth}{X} 那接待我，也接待全教會的該猶，向你們問安。城裡的財務官以拉都和弟兄括土向你們問安。 \end{tabularx} \\ \\ \relax
16:24 & \begin{tabularx}{0.7\textwidth}{X} 願我們的主耶穌基督賜恩典給你們各位！阿們！ \end{tabularx} \\ \\ \relax
16:25 & \begin{tabularx}{0.7\textwidth}{X} 惟有神能照我所傳的福音和所講的耶穌基督，並照歷代以來隱藏的奧祕的啟示，堅固你們。 \end{tabularx} \\ \\ \relax
16:26 & \begin{tabularx}{0.7\textwidth}{X} 這奧祕如今顯示出來，而且按著永生神的命令，藉眾先知的書指示萬民，使他們因信而順服。 \end{tabularx} \\ \\ \relax
16:27 & \begin{tabularx}{0.7\textwidth}{X} 願榮耀，藉著耶穌基督，歸給獨一全智的神，直到永遠。阿們！ \end{tabularx} \\ \\
[1ex]
\hline
\hline
\end{longtable}
教會不是團體,但為了行政方便,有不完全的像團體一樣的組織；就如註冊,登記之類,或組織形式上的種種.教會的本質是肢體,以主耶穌為頭,弟兄姊妹為肢體；教會又是神的家,神是這個家的家長.一個身體有它的組織系統,主要在乎生命；可能組織系統有毛病,但生命依然存在.同樣教會是有生命的基督徒所組成,儘管組織可能不健全甚至破壞,生命仍然存在的.若外面有組織,而內裡沒有生命,便等於空零.因此教會不在乎組織,卻在乎肢體相愛相助.
\newline
\newline
保羅說:「我以神的慈悲勸你們,將身體獻上,當作活祭.」他並不是勸人奉獻做傳道,因為傳道人必先要奉獻；奉獻的不一定做傳道,這乃是叫我們每人把主權交出來,好討主的喜悅.獻為活祭並不是專為某些人而說,乃是為全體基督徒說的.我們聽見保羅以神的慈悲勸勉的話,是出於神的愛.舊約獻祭的祭的祭牲是完全的,絕無瑕疵的；而我們是個罪人,怎配奉獻呢?雖然因信稱了義,仍不敢說絕對沒有瑕疵.但感謝神的慈悲,祂愛子的血遮蓋我們,即使屢屢犯罪跌倒,也得蒙赦免.所以問題不是看我們配不配,而是看我們有沒有憑信心接受祂的慈悲憐憫?也許你曾有幾次幾乎把自己獻上,卻又放下了,請現在就記住,只要憑信心就可以成功了.
\newline
\newline
肢體是連於元首基督的,有元首纔有各樣的恩賜,叫各肢體在自己的崗位合一事奉元首.所以凡對自己奉獻有問題的,就不能與弟兄成為肢體.要過肢體生活,就應該把自己交出來,然後能連於元首基督.在教會裡面,是無分彼此的,只有語言的分別而無種族,階級的分別.猶太人,希利尼人(代表外邦人)都同飲於一位聖靈.但禱告,唱詩,見證要用語言；語言不相通的時候可以傳譯,再不然可以分開聚會.為主為奴的,都是主的兒女,所以根本沒有分別.正如同一個家庭,都從一個父親所生；不管那一個的地位如何,成就如何,在父親眼中都寶貴的,因為都是他的兒女.多少年來,我們都承認肢體生活只在於同一信仰；可是在實行方面,卻實行得很少.因此提出幾個問題來,我們可以彼此討論:
\newline
\newline
在教會裡面,有許多人在帶領某人信主和走屬靈道路上花了很大心血.那個被帶領的人,便被稱為屬靈的兒女,如同保羅之稱提多,提摩太一樣.我相信在保羅生平中,除了提摩太和提多以外,可能還有其他的人,也是他的屬靈兒女.保羅既然稱他們為兒子,可以想像他們對保羅是多麼順服,是具有多麼深厚的情感.感情的發生並不是單方面的,提摩太之離開家庭跟隨保羅,足見他受保羅影響之大.如果在屬靈的兒女以外,－又有屬世界的乾兒乾女,這樣就把世界的味道帶進教會來了.再下去就會產生不良影響,愛某人惡某人,把肢體分開了,屬靈的事是在心內的,不是掛在嘴唇或只有外表的動作.此外,對於肉身的兒女和屬靈的兒女,愛的態度也要注意.假如肉身的兒女,他們也是信主而肯走屬靈道路的,就應該兩方面都一樣看待.倘若肉身兒女不信或是並非真信的,這樣的兒女,是暫時的,在天家不會相見.屬靈兒女能與我們同心,同走十字架道路,是永遠的有價值的.我們會對於自己肉身的兒女安排他們的教育,前途；但對屬靈的兒女有沒有這樣為他們設想呢?應該盡為長者的責任,不但關懷他們的靈命,更應該體貼關心他們肉體的需要.若不然,帶領他信了主,他稱為神的兒女是對的；若我們稱他是肢體便不對了.因為對他並沒有盡上肢體相顧的本份.多少人離世的時候,在遺囑上寫明遺產是給肉身的兒女之外,就很少有人想到這些財產是主交託他管理的；若交與不認識主不愛主的兒女,豈不是把神的財產隨意浪費了嗎?所以應該思想交給神的家－教會,或者交給可靠的屬靈兒女才是.倘看他肯這樣做,表明他的屬靈生活是有見證的.對年青人來說,要尊敬感謝你們的屬靈父母,要供養他們像供養你們肉身的父母一樣.
\newline
\newline
許多基督徒有一種錯誤的觀念,以為傳道人供應我們的,不過是講台上或查經班上的話語；因此,為了表示報答也只用話語.須知話語是傳道人從神所領受的職事,他是用這些話語使你得生命,進而生命豐盛.同時,他的供應是盡他所能,我們豈可以薄薄的對他一番言語上的謝謝,或者用很少代價的物質表示.試想,自己對於父母怎樣供養?中個人尊師重道,有許多事物是帶給人很有幫助的,比方對於老師,認為「一天受教,終身為父」.考試榜上有名了,也應首先去叩謝老師,甚至老師行在前面,學生只能跟在後面走.總之,是敬重他教育自己成材.但現今我們對帶領我們的牧師,傳道,長老,執事,女傳道們,他們在我們的心中地位怎樣?也許認為他們不過是教會聘請的；我是這教會的會友,他是這教會的職員.這就錯了!教會既然不是社團組織,是神的家.那麼,他們就是這個家的家長.「從前引導你們,傳神之道給你們的人,你們要想念他們,效法他們的信心.」(來十三7)他們為我們多方禱告,看顧帶領我們豈不應該敬重和供給?肢體生活是多方面的,主愛我們,為我們而死,我們記念祂,願意為祂盡忠,難道不應該對肢體彼此相顧?
\newline
\newline
我不知道有多少人履行什一奉獻?其實這不過是起碼的奉獻標準.猶太人並不只這樣,這十分之一只是拿來給利未人,其他用途的還要奉獻.人若什一奉獻也不肯拿出來,他的肢體生活就不見得實在；倘若那人沒有把薪水固定的拿回家庭作用度,他對於家庭便是不負責,對教會也是一樣.求主向我們施恩,不是我們聽不見這樣奉獻的道理,卻是不肯實行.如果每個家庭都實行什一奉獻,二十個家庭就可以供給一家傳道人,那麼,那位傳道人生活便成為與一般弟兄同等.他在會教工作不會受窘,倘若他的生活待遇太高或太低,彼此之間的交通便不暢快,見面的禱告,問安不見得真誠實在.早期教會肢體間的生活,是彼此坦誠的,相愛相助的,請問現今我們如何?
\newline
\newline
保羅稱提摩太為真兒子,提摩太稱保羅為父老,這樣的教會生活,讓年長一輩的有很大安慰.不少人工作了一生,沒有人給他安慰,得不著溫暖.主對於為祂勞苦的人,肯定是有報答的；但在弟兄的身上,也應該給他們有同情有安慰有溫暖.若我們個別失去這肢體相愛的見證,也就影響整個教會.但願憑著主的慈悲帶領,有多少力量就做出多少來；不管別人怎樣,難道別人不信主你也不信?別人不傳福音你也不傳?主留下我們在此地是有很大作用的,應該趁著機會將肢體生活顯明,以主耶穌為教會的元首；向這世代的社會作見證,也在魔鬼面前有見證.
\newline
\newline


\newpage

\section{第42屆~港九培靈會~黃聿侯~第五講}
\label{sec:1436}
\textbf{得勝陰間的權柄}
\newline
\newline
連結: \href{https://www.hkbibleconference.org/session-message/view/1436}{\texttt{https://www.hkbibleconference.org/session-message/view/1436}}
\newline
\newline
\hyperref[sec:1443]{< < < PREV SERMON < < <}
~
\hyperref[sec:toc]{[返目錄]}
~
\hyperref[sec:1444]{> > > NEXT SERMON > > >}
\newline
\newline
馬太福音 16:13-16:20
\newline
\begin{longtable}{cl}
\hline
\hline
章節 & 經文 (和合本修訂版)\\
\hline
16:13 & \begin{tabularx}{0.7\textwidth}{X} 耶穌到了凱撒利亞．腓立比的境內，就問門徒：「人們說人子是誰？」 \end{tabularx} \\ \\ \relax
16:14 & \begin{tabularx}{0.7\textwidth}{X} 他們說：「有人說是施洗的約翰；有人說是以利亞；又有人說是耶利米或是先知中的一位。」 \end{tabularx} \\ \\ \relax
16:15 & \begin{tabularx}{0.7\textwidth}{X} 耶穌問他們：「你們說我是誰？」 \end{tabularx} \\ \\ \relax
16:16 & \begin{tabularx}{0.7\textwidth}{X} 西門‧彼得回答說：「你是基督，是永生神的兒子。」 \end{tabularx} \\ \\ \relax
16:17 & \begin{tabularx}{0.7\textwidth}{X} 耶穌回答他說：「約拿的兒子西門，你是有福的！因為這不是屬血肉的啟示你的，而是我在天上的父啟示的。 \end{tabularx} \\ \\ \relax
16:18 & \begin{tabularx}{0.7\textwidth}{X} 我還告訴你，你是彼得，我要把我的教會建造在這磐石上，陰間的權柄不能勝過它。 \end{tabularx} \\ \\ \relax
16:19 & \begin{tabularx}{0.7\textwidth}{X} 我要把天國的鑰匙給你，凡你在地上所捆綁的，在天上也要捆綁；凡你在地上所釋放的，在天上也要釋放。」 \end{tabularx} \\ \\ \relax
16:20 & \begin{tabularx}{0.7\textwidth}{X} 當時，耶穌囑咐門徒不可對任何人說他是基督。 \end{tabularx} \\ \\
[1ex]
\hline
\hline
\end{longtable}
基督已得勝陰間權柄的,我們是祂的門徒,也需要得勝陰間的權柄.
\newline
\newline
主對彼得說:「我還告訴你,你是彼得,我要把我的教會建造在這磐石上,陰間的權柄,不能勝過他.」小字的「權柄」原文作門,意思是陰間的門不能勝過他.所以下面講:「我要把天國的鑰匙給你,凡你在地上所捆綁的,在天上也要捆綁；凡你在地上所釋放的,在天上也要釋放.」
\newline
\newline
在這裡提到三件事:眼所不能見的「權柄」.它的力量超過人所能形容,權力有多大,乃要看賜權柄的是甚麼人.一個家長,他的權柄就是那個家；一個校長,他的權柄就是那所學校.我們的主,祂是萬王之王,宇宙間沒有一個人比祂更偉大,因此祂的權柄也最偉大.我們接待祂,祂就賜我們權柄,作神的兒女.(約一12)在主的權柄以外,另一個是掌握陰間死亡的權柄,那是罪惡的結局,但主已經勝過牠了.
\newline
\newline
「陰間」,是無底坑的門,永不滿足的,牠日日張大口吞噬,滿足牠是不可能的,只可能用勝過牠的方法.人是有罪性的,情慾驕傲各樣敗壞,越想滿足牠,便越感覺空虛,必須要靠主得勝,感謝主,祂賜天國鑰匙給彼得,也賜給與彼得同樣承認信仰的人.
\newline
\newline
「鑰匙」是可以關門也可以開門的.主的救恩是得勝陰間的權柄.保羅在(林前十五)歌唱:「死啊,你得勝的權勢在那裡?死啊,你的毒凡在那裡?死的毒凡就是罪,罪的權勢就是律法；感謝神!使我們藉著我們的主耶穌基督得勝.」主說:「我是道路,真理,生命,若不藉著我,就沒有人能到父那裡去.」這是祂為我們開了天堂的門,所以我們藉著祂的救恩能夠出死入生.
\newline
\newline
若我們單看馬太十六章,對於捆綁和釋放的問題還不容易明白,再看十八章15至35節,便見到所指是講饒恕.19節說兩個人在地上同心合意祈求,天父必成全,不會是指一般禱告.我在神面前領受的就是一個人得罪了我,我要赦免他,但恐怕他不知道我赦免了；而為了他的好處,免得他心中責備受捆綁,因此當面說過又一同祈禱.正如主教導我們禱告,我赦免了人,天父也赦免我.主的救恩是要我們過彼此相愛生活的,若是我們不肯饒恕,便是不肯相愛,等於救恩目的還未達到.我們的主已經勝過陰間,我們為何不肯饒恕,卻把陰間帶到地上來?我們活在地上不肯赦免,心中仍然是被罪所捆綁便沒有釋放.現今教會之所以沒有見證力量,也就是因為彼此不肯饒恕不肯赦免.我們既然蒙恩得救,與陰間權柄無關,過在地若天的生活；雖然人的軟弱和魔鬼的詭計,不免互相之間有得罪,但彼此饒恕,主的名便得榮耀了.
\newline
\newline
主耶穌勝過陰間的權柄,所以門徒開始傳道便強調祂從死裡復活,說明神已經立祂為基督.不管世人對祂的反應怎樣,但他們毫不動搖.而我們的主事實上是勝過死亡的,祂之從死裡復活,便堅固了門徒,一直忠心至死的傳道.傳說彼得是為主而倒釘十字架的,保羅是羅馬公民,不能夠把他釘十字架是砍頭而殉道的；他在提摩太後書四章唱出凱歌,澆奠的時候到了,是滿懷得勝信心.第一個殉道的司提反,人們圍著他咬牙切齒,但他卻像天使一樣的面貌.他沒有理會當前的環境,也沒有計較曾經與自己同伴的人在那裡,只是想到自己與主的關係,他看見神的榮耀,看見耶穌站在神的右邊.經上常說主是坐在神右邊的,但此處卻是說站起來.在司提反心中,不能不感動,自己算甚麼,或者釘死,或者打死,都不要緊,能夠見主的面與祂同在就是最大的滿足.
\newline
\newline
信仰可以分為兩方面:例如(來十一1)「信就是所望之事的實底,是未見之事的確據.」許多時我們會說現在還未有,將來纔成就.其實,神叫我們相信之日,就有那事存在,並不是要我們信了纔慢慢有.在神來說,那事早已存在,只因為人是屬肉體的,所以看不見,這是人和神之間時間,空間發生的困難.我們看有過去,現在,與未來,在神看是沒有這些的.我們相信耶穌降生,受死,復活,升天還要再來作王,似乎現在祂尚未作王；不過站在基督徒立場說,祂已經在我們心中作王了,並且祂將來必作萬王之王.舊約先知總是把祂降生與再來連在一起講的,祂必要再來,不管地上怎樣反對,即如當日法利賽人,撒都該人試探主行神蹟.雖然他們兩幫人本來不合作,現在竟然合作起來,他們聯合對付主,主卻不理會他們.所以在現今的日子,外面的人怎樣聯合攻擊主,攻擊教會攻擊我們,我們也不要失望.乃要學主的樣式,向主負責忠心到底,主得勝陰間的權柄,我們靠祂也得勝陰間權柄的.
\newline
\newline
我們蒙受了主赦罪之恩,在地上稱為基督徒－基督人,生命上是神的一部份,將來要改變形狀像主.所以我們要傳福音領人到教會,也在弟兄中間過相愛生活,不但有信仰,並且有生活表現；把生命自然的流露,成為活水的江河,幫助週圍的人都可以飽足.我們所過是在地若天的生活,對於「死」的問題也不應懼怕,因為我們是隨時預備好,早已解決了對弟兄的虧欠問題,也對主的託付盡了責任.怕死只不過是未信的時候所留下的印象,若是死了乃是改變當前的環境過一種新的生活,當我們面對死亡－仇敵,疾病,刀劍,都不逃避；卻像保羅,像司提反的一樣從心中發出勝利的歌唱.因為我們的信仰是堅固不動搖的,有勝過陰間的權柄,藉著這堅定的信仰,見證主,也幫助那些灰心失望軟弱的人!
\newline
\newline


\newpage

\section{第42屆~港九培靈會~黃聿侯~第六講}
\label{sec:1444}
\textbf{做忠心有見識的管家}
\newline
\newline
連結: \href{https://www.hkbibleconference.org/session-message/view/1444}{\texttt{https://www.hkbibleconference.org/session-message/view/1444}}
\newline
\newline
\hyperref[sec:1436]{< < < PREV SERMON < < <}
~
\hyperref[sec:toc]{[返目錄]}
~
\hyperref[sec:1437]{> > > NEXT SERMON > > >}
\newline
\newline
馬太福音 25:14-25:30
\newline
\begin{longtable}{cl}
\hline
\hline
章節 & 經文 (和合本修訂版)\\
\hline
25:14 & \begin{tabularx}{0.7\textwidth}{X} 「天國又好比一個人要出外遠行，就叫了僕人來，把他的家業交給他們。 \end{tabularx} \\ \\ \relax
25:15 & \begin{tabularx}{0.7\textwidth}{X} 他按著各人的才幹，給他們銀子：一個給了五千，一個給了二千，一個給了一千，就出外遠行去了。 \end{tabularx} \\ \\ \relax
25:16 & \begin{tabularx}{0.7\textwidth}{X} 那領五千的立刻拿去做買賣，另外賺了五千。 \end{tabularx} \\ \\ \relax
25:17 & \begin{tabularx}{0.7\textwidth}{X} 那領二千的也照樣另賺了二千。 \end{tabularx} \\ \\ \relax
25:18 & \begin{tabularx}{0.7\textwidth}{X} 但那領一千的去掘開地，把主人的銀子埋藏了。 \end{tabularx} \\ \\ \relax
25:19 & \begin{tabularx}{0.7\textwidth}{X} 過了許久，那些僕人的主人來了，和他們算賬。 \end{tabularx} \\ \\ \relax
25:20 & \begin{tabularx}{0.7\textwidth}{X} 那領五千的又帶著另外的五千來，說：『主啊，你交給我五千。請看，我又賺了五千。』 \end{tabularx} \\ \\ \relax
25:21 & \begin{tabularx}{0.7\textwidth}{X} 主人說：『好，你這又善良又忠心的僕人，你在少許的事上忠心，我要派你管理許多的事，進來享受你主人的快樂吧！』 \end{tabularx} \\ \\ \relax
25:22 & \begin{tabularx}{0.7\textwidth}{X} 那領二千的也進前來，說：『主啊，你交給我二千。請看，我又賺了二千。』 \end{tabularx} \\ \\ \relax
25:23 & \begin{tabularx}{0.7\textwidth}{X} 主人說：『好，你這又善良又忠心的僕人，你在少許的事上忠心，我要派你管理許多的事，進來享受你主人的快樂吧！』 \end{tabularx} \\ \\ \relax
25:24 & \begin{tabularx}{0.7\textwidth}{X} 那領一千的也進前來，說：『主啊，我知道你，你是個嚴厲的人：沒有種的地方也要收割，沒有播的地方也要收穫， \end{tabularx} \\ \\ \relax
25:25 & \begin{tabularx}{0.7\textwidth}{X} 我就害怕，去把你的一千銀子埋藏在地裡。請看，你的銀子在這裡。』 \end{tabularx} \\ \\ \relax
25:26 & \begin{tabularx}{0.7\textwidth}{X} 他的主人回答他說：『你這又惡又懶的僕人，你既知道我沒有種的地方也要收割，沒有播的地方也要收穫， \end{tabularx} \\ \\ \relax
25:27 & \begin{tabularx}{0.7\textwidth}{X} 就該把我的銀子放給兌換銀錢的人，到我來的時候可以連本帶利收回。 \end{tabularx} \\ \\ \relax
25:28 & \begin{tabularx}{0.7\textwidth}{X} 把他這一千奪過來，給那有一萬的。 \end{tabularx} \\ \\ \relax
25:29 & \begin{tabularx}{0.7\textwidth}{X} 因為凡有的，還要加給他，叫他有餘；沒有的，連他所有的也要奪過來。 \end{tabularx} \\ \\ \relax
25:30 & \begin{tabularx}{0.7\textwidth}{X} 把這無用的僕人丟在外面黑暗裡，在那裡他要哀哭切齒了。』」 \end{tabularx} \\ \\
[1ex]
\hline
\hline
\end{longtable}
上面的兩段經文,與救恩無關,只是對主的僕人而言.因為能夠做主僕人的,必然是重生得救的.正如馬太廿五章的十童女,也與救恩無關,只是叫人儆醒等候主再來,忠心於主的交託.同一章裡面講到作在弟兄中最小一個的身上,就是作在主的身上；說明儆醒,忠心都是為愛主,愛弟兄.有時我們會強調工作的重要,慢慢變為愛工作；但要注意,作工的目的是為愛主愛弟兄,而不是為愛那工作.所以主人給僕人五千,二千,一千,乃按他們才幹而給他們；路加所記載的也類似同樣教訓,要他們作忠心有見識的管家.彼得問這比喻是對我們說呢?還是為眾人.主並沒有立界限指定是誰,只說:「誰是忠心有見識的管家,為主人派他管理家裡的人；按時分糧給他們呢?主人來到,看見僕人這樣行,那僕人訧有福了.」祂的意思是指每一個人.保羅說:「所求於管家的,是要他有忠心.」(林前四2)他又在提多書說到監督,執事都是神的管家,因此可見我們都是神的管家.神的教會既然是神的家,我們受過肢體生活,若管家責任只落在少數人身上,像家庭只由父母管家,其他成員都不負責任.這樣,做管家的便很難,以後便沒有二代,父母離了世家庭便分散了.
\newline
\newline
在家庭中哥哥姐姐做管家,乃是父母所交託的,有時也分工給弟弟妹妹.做僕人使女的也如是,主人是按才幹把工作交託給僕人.總之,每個人對家庭都要盡責,這個家就會和諧興旺.從小習慣於管家,當你自己成家,就曉得自己管理家庭,不覺得吃力了.
\newline
\newline
教會也是如此,若管家光是傳道牧師,那傳道牧師便夠苦了.一個家長管理十人,八人,已感覺到困難.一個教會總有過百人以上,而且還不斷會發展,豈是少數人能理得好的呢?有人說,除傳道牧師之外,還有長老執事,但事實上還未夠的,必須人人都對這個家有責任感.否則,若有一兩個職員差派到別處,或許離了世；其他的人便茫茫然沒有方向,也因之而變質了.一位主的僕人,他說父母管教很嚴,誰看見有東西丟在地上,靠近的那個兄弟便要把他拾起來.不能推諉說不是我丟的.我見到有些少年人,在家裡吃了飯便走開,從不幫忙父母；他媽媽整天忙碌,他也不管,必須要從小就養成愛家庭幫助家庭工作的習慣.以後長大了纔會理家,以致應世.我們對於在家中管家的父母,要體會他們的難處,經常給他們安慰幫助；不是只在生日買點禮物,給他們一些錢就算事.教會也是一樣,倘若弟兄姊妹把來教會的當做客人,這個教會是可悲的.到教會聚會的,要把教會看作自己的家,千萬不要抱漠不相關的態度.聖經有沒有擺好,窗門有沒有關好；雖然小事,我們能夠關心,就表明愛主愛教會.經費情形如何,聚會人數如何,自己都不關心,以為這是長老,執事,牧師,傳道的事.當他們有一天離了職,誰去承繼?若選到你為長老執事,而你從來沒有過為教會工作的經驗,硬著頭皮去做也做不好的,因為你平日與教會的工作脫了節.在教會裡不可能人人都上講台,也不可能人人都有一個名目擔任職份；不過我們仍可以為愛教會,自然而然的工作.正如在家庭裡從來沒有計較,有甚麼名銜說幫助家務的.倘若教會裡人人都有責任感,大事小事都能夠通力合作,這教會一定興旺.第一次來過這教會的,由於被這裡的人彼此合作,關心的態度所影響,也就喜歡參加這裡的聚會.所以我們要習慣有服事主,服事人的表現,成為好管家.
\newline
\newline
我們之熱心作主的工,是羨慕聖工,而不是羨慕那些職份.凡是與神的家有關心,我們都樂意有一份,盡自己知識,才能去做.常常做,多多做,就越會得著主的賜福.既然做管家是為主而做,因此傳道人牧師不是由於人請自己,才做這工作.同樣,做長老執事不是為這職份而做.做長老執事的,因為一向都熱心服事主服事人,他的人生他的工作,都符合聖經中為長老執事的資格.作監督的也是一樣,是因為他的為人,他的忠心,見識,教會的人都看見了；所以請他牧養全群,不應今年選上便做,明年選不上便不做.在教會中事奉主是不應論名份職銜的.難道在家庭中做哥哥姐姐,也要弟弟妹妹選舉才做?我們在教會中服事主,不是因被分派,也不是靠人的鼓勵,更不是看人.我們有主的靈在裡面,有祂的話語在我們心中:因著祂的愛所感動,我們起來事奉,是理所當然的.
\newline
\newline
在家庭中管家有困難,要做有忠心有見識的管家更難,因為人事是最複雜的.中國文化重在人倫,做臣子的不管君王怎樣,也要忠心.做兒子的不管父親怎樣,也要孝順.而教會的情形比中國之專注重人倫更深一步.因為在人之中有得救與未得救的,得救的人裡面又分屬血氣的,屬肉體的,屬靈的,教會又是永生神的家.這個家是永存的,換言之,把天上的事搬到地上來.我們有責任把天上的榮耀向世人彰顯,所以做神的管家是比較困難.保羅為我們立了一個榜樣,他把我們認為難做的做出來.哥林多教會是他辛辛苦苦建立的,他織帳棚來自給,他為那些因信主受逼迫遭反對的,花了很多心血,想不到後來有人傳假道,進入來混亂他們；不但否認保羅是使徒,裡面還分黨結派,甚至對保羅人身攻擊,說他氣貌不揚言語粗俗.試想,像保羅對哥林多教會,如同父母養育子女一樣,竟然被攻擊,豈不心碎?但,他表示:「所求於管家的,是要有忠心,我被你們論斷；或被別人論斷,我卻以為極小的事,連我自己也不論斷自己.」這是做管家應有的態度,而教會裡面,追求和事奉比較熱心的,許多時會招來誤會批評.因為在一大群人中,靈性,背景經歷與主交通的深度,都各有不同.不可能絕對要求自己所說的,別人必須阿們.所以被人反對的時候,只要檢討自己,卻不要失望灰心.
\newline
\newline
保羅個人被人誤會批評,他完全接受,不申辯不解釋.因為既然別人有了成見,解釋就難!若爭取別人的同情,可能會事先就維護自己,而責備他人,反為不美.別人若不同情,若在這事情上加上甚麼,自己就會苦上加苦；惟有到主面前,單從祂那裡得力量得安慰.但若為了主的道,為了十字架的真理,非爭辯不可.這樣的爭辯不是為自己,是為了見證,主耶穌在地上的時候也是這樣.可惜,現今多少基督徒對人對事分不清.公私也分不清；因此教會內部所發的爭執,只是人事問題而非真理問題.我們不能推諉到信仰上去.這樣動機就會錯誤,巴不得我們像保羅一樣,個人的問題不予爭論.
\newline
\newline
還有,神的家和自己肉身的家要分明,尤其在教會中負有責任的,不要把教會的難處帶回家裡講論；若然,就會影響下一代不愛主,不肯為主工作,這是我們極重大的損失.在新加坡一位已離世的老牧師,他受了委屈,當場昏倒,但回到家中也不講出被欺負的內情,兒女勸他辭職他也不答應；所以在他過世之後,他的兩個兒子竟繼續承擔聖職.若我們把教會難處帶回家去宣揚,有意無意之間就領家人得罪主,帶來灰心失望,雖然自己還可以受得住的,妻子兒女卻受不住.或許當時自己的態度不好,語氣粗暴,留給家人的印象也不好,以後便沒有辦法改正了.在教會有職份的,是服事弟兄姊妹,他們難免有軟弱,我們要為他們禱告.個別約他們談話,不應把他們的事宣揚出去.有誰會對人說自己兒女不好的呢?屬靈的兒女也是自己用福音所生的,用主的道所栽培的；又怎可以把他們的弱點公開,而不好好地教導他們呢?所以應該在主的愛裡包容,忍耐他們.
\newline
\newline
至於應該怎樣做管家,聖經說要按時分糧,這有養育的意思在內的,養育就包括衣食住行.做父母的處處為兒女著想,因此我們在教會裡既是管家中的一份子,不但自己穿得好,吃得好,也要顧到全教會的需要.若自己在一群不太富足的弟兄中間,穿著得特別華麗,很容易使人難於接近,比較貧苦的弟兄就不肯來,這樣便失去弟兄了.我們雖然富足,也不該浪費,因為這些財物是主所交託管理,生活夠用就該知足了.應該節省下來,幫助貧寒的弟兄,尤其要為傳福音而獻上自己的財物.世人追求物慾,奢華宴樂,這不足為怪；但我們是不屬這世界的,在神家中都是一家人,怎可以不相顧呢?
\newline
\newline


\newpage

\section{第42屆~港九培靈會~黃聿侯~第六講}
\label{sec:1437}
\textbf{面向十字架的信仰}
\newline
\newline
連結: \href{https://www.hkbibleconference.org/session-message/view/1437}{\texttt{https://www.hkbibleconference.org/session-message/view/1437}}
\newline
\newline
\hyperref[sec:1444]{< < < PREV SERMON < < <}
~
\hyperref[sec:toc]{[返目錄]}
~
\hyperref[sec:1432]{> > > NEXT SERMON > > >}
\newline
\newline
馬太福音 16:13-16:27
\newline
\begin{longtable}{cl}
\hline
\hline
章節 & 經文 (和合本修訂版)\\
\hline
16:13 & \begin{tabularx}{0.7\textwidth}{X} 耶穌到了凱撒利亞．腓立比的境內，就問門徒：「人們說人子是誰？」 \end{tabularx} \\ \\ \relax
16:14 & \begin{tabularx}{0.7\textwidth}{X} 他們說：「有人說是施洗的約翰；有人說是以利亞；又有人說是耶利米或是先知中的一位。」 \end{tabularx} \\ \\ \relax
16:15 & \begin{tabularx}{0.7\textwidth}{X} 耶穌問他們：「你們說我是誰？」 \end{tabularx} \\ \\ \relax
16:16 & \begin{tabularx}{0.7\textwidth}{X} 西門‧彼得回答說：「你是基督，是永生神的兒子。」 \end{tabularx} \\ \\ \relax
16:17 & \begin{tabularx}{0.7\textwidth}{X} 耶穌回答他說：「約拿的兒子西門，你是有福的！因為這不是屬血肉的啟示你的，而是我在天上的父啟示的。 \end{tabularx} \\ \\ \relax
16:18 & \begin{tabularx}{0.7\textwidth}{X} 我還告訴你，你是彼得，我要把我的教會建造在這磐石上，陰間的權柄不能勝過它。 \end{tabularx} \\ \\ \relax
16:19 & \begin{tabularx}{0.7\textwidth}{X} 我要把天國的鑰匙給你，凡你在地上所捆綁的，在天上也要捆綁；凡你在地上所釋放的，在天上也要釋放。」 \end{tabularx} \\ \\ \relax
16:20 & \begin{tabularx}{0.7\textwidth}{X} 當時，耶穌囑咐門徒不可對任何人說他是基督。 \end{tabularx} \\ \\ \relax
16:21 & \begin{tabularx}{0.7\textwidth}{X} 從那時起，耶穌才向門徒明說，他必須上耶路撒冷去，受長老、祭司長和文士許多的苦，並且被殺，第三天復活。 \end{tabularx} \\ \\ \relax
16:22 & \begin{tabularx}{0.7\textwidth}{X} 彼得就拉著他，責備他說：「主啊，千萬不可如此！這事絕不可臨到你身上。」 \end{tabularx} \\ \\ \relax
16:23 & \begin{tabularx}{0.7\textwidth}{X} 耶穌轉過來，對彼得說：「撒但，退到我後邊去！你是我的絆腳石，因為你不體會神的心意，而是體會人的意思。」 \end{tabularx} \\ \\ \relax
16:24 & \begin{tabularx}{0.7\textwidth}{X} 於是耶穌對門徒說：「若有人要跟從我，就當捨己，背起自己的十字架來跟從我。 \end{tabularx} \\ \\ \relax
16:25 & \begin{tabularx}{0.7\textwidth}{X} 因為凡要救自己生命的，要喪失生命；凡為我喪失生命的，要得著生命。 \end{tabularx} \\ \\ \relax
16:26 & \begin{tabularx}{0.7\textwidth}{X} 人若賺得全世界，賠上自己的生命，有甚麼益處呢？人還能拿甚麼換生命呢？ \end{tabularx} \\ \\ \relax
16:27 & \begin{tabularx}{0.7\textwidth}{X} 人子要在他父的榮耀裡與他的眾使者一起來臨，那時候，他要照各人的行為報應各人。 \end{tabularx} \\ \\
[1ex]
\hline
\hline
\end{longtable}
我們的信仰,是面向十字架的,或者可以說面向死亡的.這樣的信仰不是口頭說說唱唱,乃是要付出極大代價的；當付出代價以至於死的時候,可能旁人並不明白,但,這是為主而死,即使受鞭打,刀劍,殺戮,或者勞碌了一生；總之是為主而背上十字架.
\newline
\newline
當主與門徒論信仰之後,接著讓門徒清楚知道,祂必須上耶路撒冷受許多的苦.因為主深知天父在祂身上有目的,不管人們對祂了解與否,祂仍是面向耶路撒冷走去.門徒恐懼,彼得害怕,其他門徒也一樣徬徨.這是門徒在五旬節聖靈降臨以前的情形,他們雖然有純正的信仰,可是還沒有聖靈的充滿,因此他們站立不住.在五旬節有了聖靈充滿之後,情形便大大不同,門徒站起來講道,三千,五千人悔改；一切情形都改變,凡物公用,擘餅紀念主.從前他們聽到主說:「若有人跟從我,就當捨己,北起他的十字架來跟從我.」他們就不知怎樣辦,好像那個少年財主,心中渴想跟從主,卻放不下那麼多的財產,結果憂憂愁愁地走了.但聖靈充滿之後,他們撇下一切也不感覺為難了.
\newline
\newline
主耶穌是明白一切的,所以在約翰福音說:「等真理的聖靈來了,他要引導你們明白一切的真理.」感謝主,現今有聖靈在我們身上工作,我們聽道,讀經便有所得著.因此當我們明白了信仰,便要為這信仰付出代價,盡基督徒的責任.求主幫助,給我們看見昔日的人怎樣為信仰付出代價；我們也學習他們一樣,為了信仰甘心背起十字架.
\newline
\newline
當日跟隨主的人,必被逐出會堂的.會堂制度在寫新約的時候,人人都明白是怎樣一回事,正如主耶穌降生,一提到彌迦書五章的預言,大家都了解一樣.會堂在某一方面來說,像現今的禮拜堂讓大家集中敬拜神,但會堂可以說是宗教,政治,文化一切活動的中心.那位管會堂的並非現今禮拜堂的堂役而是長老,負責人；若被趕逐出會堂,無異與這個社會脫離了關係.是非的判斷,產業的交易,甚至兒子所受基本教育的權利,都被剝奪了,不但個人被社會孤立,連整個家庭也同樣被孤立.我們看約翰九章主耶穌所醫好的那個瞎子,就是受到這樣的待遇.從約翰十二章我們又看到在官長中有好些信祂的,只因法利賽人的緣故就不承認,恐怕被趕出會堂.寫福音的人還註明,這是因他們愛人的榮耀,過於愛神的榮耀.意思很明顯,被趕出會堂是極其羞恥的事,在社會上再不易抬頭.各位弟兄姊妹,我們聽來,似乎與我們沒有多大的關係,但曾否想到,如今不少地方,基督教也受到極大逼迫,在社會上被人孤立,口糧也得不到配給,被人指為噬人的老虎,這就是因為他們相信耶穌是基督,是永生神的兒子.
\newline
\newline
耶穌基督為我們受死受苦,我們為他受些苦,根本算不得甚麼；所以在我們的信仰受到考驗的時候,務要堅持.我們所信的是救百姓從罪惡裡出來的主耶穌,祂必須是神的兒子纔能從死裡復活,不然我們就不能得救,沒有重生的生命.當我們有這樣堅強信仰的時候,不僅可以成為祂的門徒,並且祂還賜我們權柄作神的兒女.主耶穌在路加福音對跟隨祂的人說:「人到我這裡來,若不愛我勝過愛自己的父母,妻子,兒女,弟兄,姊妹和自己的生命,就不能作我的門徒；凡不背著自己十字架跟從我的,也不能作我的門徒.」祂是願意萬人得救,不願意有一個人沉淪的.但祂自己既然清楚神給祂的使命,故此也盼望跟隨祂的人,也能清清楚楚,祂毫不含糊的對大家講個明白,以至說:你們那一個要蓋一座樓,不先坐下算計花費,能蓋成不能呢......或是一個王,出去和別個王打仗,豈不先坐下酌量,能用一萬兵去敵那領二萬兵來攻打他的麼......意思是必須考慮清楚,免得後來失敗羞辱.
\newline
\newline
在約翰第六章記載,主耶穌用五餅二魚使許多人吃飽,便有許多人要跟從祂.主見他們的目的不對,便對他們說,昔日他們的祖宗在曠野吃過嗎哪,結果還是死了；必須要吃那從天上降下來的糧,纔可以不死!這天上來的糧就是祂的血祂的肉,許多猶太人聽了,都覺得為難,就離開了祂.祂就對十二個門徒說,你們也要去麼,西門彼得回答說:「主啊,你有永生之道,我們還歸從誰呢?」他們覺得主這樣愛他們,表示跟隨到底.主耶穌是教會的頭,信徒要彼此配合,不是要求人數多就可以,若勉強湊數只有帶來教會累贅.若真是肢體,必須一個一個都是有生命的,互相聯絡的.若沒有生命的混在裡面,倒不若讓他們離開還好.
\newline
\newline
主在地上選了十二個門徒,時常和祂在一起.但這些人跟隨了祂三年半,一個把祂出賣,一個不認祂,有兩個爭坐祂的左右.按人說這是極大的失敗,耶穌也看得很透澈,所以祂在約翰十七章的禱告非常痛苦,然而卻是為他們和以後傳道的人信祂的人來祈求；可見祂對我們的盼望是十分大,祂是憑著愛心信心而祝福,讓不可能的事變變為能.能夠為祂的榮耀而見證.當年的瞎子為信主的緣故,被逐出會堂.有人不敢公開承認主,也是怕被趕出會堂,可是祂那十二個門徒隨時有被趕逐的可能也不畏懼,是非常難得的.當日的耶穌是還未復活升天的,不過是拿撒勒人,在人看來似乎沒有成就；而彼得等,大概不是一個普通的漁夫,(我曾到過迦百農,他們現在正發掘彼得的房子,發現那房子與會堂有同樣大的面積.)他與兄弟肯放下家庭,冒政治,宗教,社會各方面的危險,卻去跟隨一位不為人所敬重的耶穌.當時的人不明白,而彼得他們竟緊緊跟隨,這就是主所禱告的:「凡你所賜給我的人,都到我這裡來；到我這裡來的,我總不丟棄他.」
\newline
\newline
在人的思想認識中,常會與靈性有矛盾的,但我們要在矛盾中越加掙扎倚靠主.向自己是死的,向神是活的.我們的主已經復活升天,且要再來作王.祂對於我們,和我們的家庭,教會都有說不盡的恩典.若是我們不肯背起十字架來跟從祂,就不但比當年的門徒不如,連那瞎子也比不上了.但願主給我們更大信心,讓我們愛祂勝過一切,隨時隨地都可以為祂的緣故而死.
\newline
\newline


\newpage

\section{第42屆~港九培靈會~黃聿侯~第七講}
\label{sec:1432}
\textbf{背負十字架的能力}
\newline
\newline
連結: \href{https://www.hkbibleconference.org/session-message/view/1432}{\texttt{https://www.hkbibleconference.org/session-message/view/1432}}
\newline
\newline
\hyperref[sec:1437]{< < < PREV SERMON < < <}
~
\hyperref[sec:toc]{[返目錄]}
~
\hyperref[sec:1438]{> > > NEXT SERMON > > >}
\newline
\newline
馬太福音 16:13-16:27
\newline
\begin{longtable}{cl}
\hline
\hline
章節 & 經文 (和合本修訂版)\\
\hline
16:13 & \begin{tabularx}{0.7\textwidth}{X} 耶穌到了凱撒利亞．腓立比的境內，就問門徒：「人們說人子是誰？」 \end{tabularx} \\ \\ \relax
16:14 & \begin{tabularx}{0.7\textwidth}{X} 他們說：「有人說是施洗的約翰；有人說是以利亞；又有人說是耶利米或是先知中的一位。」 \end{tabularx} \\ \\ \relax
16:15 & \begin{tabularx}{0.7\textwidth}{X} 耶穌問他們：「你們說我是誰？」 \end{tabularx} \\ \\ \relax
16:16 & \begin{tabularx}{0.7\textwidth}{X} 西門‧彼得回答說：「你是基督，是永生神的兒子。」 \end{tabularx} \\ \\ \relax
16:17 & \begin{tabularx}{0.7\textwidth}{X} 耶穌回答他說：「約拿的兒子西門，你是有福的！因為這不是屬血肉的啟示你的，而是我在天上的父啟示的。 \end{tabularx} \\ \\ \relax
16:18 & \begin{tabularx}{0.7\textwidth}{X} 我還告訴你，你是彼得，我要把我的教會建造在這磐石上，陰間的權柄不能勝過它。 \end{tabularx} \\ \\ \relax
16:19 & \begin{tabularx}{0.7\textwidth}{X} 我要把天國的鑰匙給你，凡你在地上所捆綁的，在天上也要捆綁；凡你在地上所釋放的，在天上也要釋放。」 \end{tabularx} \\ \\ \relax
16:20 & \begin{tabularx}{0.7\textwidth}{X} 當時，耶穌囑咐門徒不可對任何人說他是基督。 \end{tabularx} \\ \\ \relax
16:21 & \begin{tabularx}{0.7\textwidth}{X} 從那時起，耶穌才向門徒明說，他必須上耶路撒冷去，受長老、祭司長和文士許多的苦，並且被殺，第三天復活。 \end{tabularx} \\ \\ \relax
16:22 & \begin{tabularx}{0.7\textwidth}{X} 彼得就拉著他，責備他說：「主啊，千萬不可如此！這事絕不可臨到你身上。」 \end{tabularx} \\ \\ \relax
16:23 & \begin{tabularx}{0.7\textwidth}{X} 耶穌轉過來，對彼得說：「撒但，退到我後邊去！你是我的絆腳石，因為你不體會神的心意，而是體會人的意思。」 \end{tabularx} \\ \\ \relax
16:24 & \begin{tabularx}{0.7\textwidth}{X} 於是耶穌對門徒說：「若有人要跟從我，就當捨己，背起自己的十字架來跟從我。 \end{tabularx} \\ \\ \relax
16:25 & \begin{tabularx}{0.7\textwidth}{X} 因為凡要救自己生命的，要喪失生命；凡為我喪失生命的，要得著生命。 \end{tabularx} \\ \\ \relax
16:26 & \begin{tabularx}{0.7\textwidth}{X} 人若賺得全世界，賠上自己的生命，有甚麼益處呢？人還能拿甚麼換生命呢？ \end{tabularx} \\ \\ \relax
16:27 & \begin{tabularx}{0.7\textwidth}{X} 人子要在他父的榮耀裡與他的眾使者一起來臨，那時候，他要照各人的行為報應各人。 \end{tabularx} \\ \\
[1ex]
\hline
\hline
\end{longtable}
提起十字架,我們往往講不出聖經上所提的那十字架,因為我們讀和聽不大準確,領受也有差別.我前幾天曾與各位提及變質的問題,一個人慢慢變了質,往往是不自覺的,即是對十字架的觀念也如是.論到信仰的變質,一方面是外面的壓力,一方面是內部的腐化.由於外面的壓力太大,教會不敢傳講福音,而內部腐化是因為社會太討好教會,不但沒有干涉；反而用政治力量幫助教會,這樣的變質,就會令我們變了還不覺得；可以說別人害了我們,而我們還會謝謝他.也許他們不是立心破壞,只不過所做的不合神的旨意,便易走入了魔鬼的圈套,以致無可救藥.今日西方教會,多半就落在這光景裡,就個人而言,別人批評固然要忍受檢討,別人稱讚也不應過份歡喜；因為自己沒有那樣好處,而別人恭維你,使你得著不應得的稱讚,心中應感受到纔是.基督徒是應該誠誠實實的,是就說是,不是就說不是.主耶穌的話,永遠都是對的,無論責備,管教,我們都應該絕對信服.不但如此,別人對教會的態度無論是好是壞,對教會都應該有正確的認識.保羅幫助哥林多教會,卻被假先知攻擊；以致引起許多人對他誤會.但保羅說:「我被你們論斷,或被別人論斷,我都以為極小的事；連我也是不論斷自己,我雖不覺得自己有錯,卻也不能因此得以稱義,但判斷我的乃是主.」這是基督徒的人生觀,換言之,必須要認識自己.
\newline
\newline
現今一般人對十字架的印象,似乎很好,可能是因為兩方面而來的.
\newline
\newline
(一)耶穌釘在十字架上,人對十字架的觀念就改觀了.十字架本來非常醜惡,人聽見了也很恐懼,連提也不願提；但因為愛我們的主釘在十字架上,那可怕的十字架便變為可愛.第三世紀君士坦丁宣佈,以後再不准用十字架釘死囚犯,他不但看主耶穌為可愛,連十字架也可愛了.社會人士因為沒有讀聖經,對於十字架的本來面目不清楚,還值得原諒,但基督徒就不應該如此了.可愛的主被釘在可咒可詛恐怖的十字架上,祂救了我們,我們豈可把十字架看作可恥可怕的呢?
\newline
\newline
(二)因傳福音而影響人們對十字架的觀念.福音傳到那裡,那裡就有十字架的記號,無論髹漆的,木做的,金屬的,總是以它作為救恩的象徵,十字架的原來為刑具的作用已經改變了.所以基督徒對於十字架應有基本的認識,斷不能拿別的認識作為我們的信仰.我們與主有生命的關係,乃是藉著聖徒和聖靈的啟示而認識祂,正如保羅說:「我雖然然藉著外貌認過基督,如今卻不再這樣認識祂了.」十字架不是裝飾品,乃是殘酷的刑具.在各個不同的代,不同的社會裡,都會有不同形式的刑具.猶太人判處死刑不是釘十字架,而是用石頭打死的.聖經有許多這樣的例證,亞干是被石頭打死的,司提反也是死於石頭之下.用石頭打死有兩個意思,第一,這樣做表示經過公議,不是報私仇；約翰八章記載人們要用石頭打死那淫婦的事可為證.第二,是給眾人作警告,一個人被打之後,便成為一大堆石頭,讓看見這堆石頭的人知所警愓.
\newline
\newline
為甚麼主耶穌要被釘在十字架上?申命記廿一章的末後一段說明:「被掛的人是在神面前受咒詛的.」猶太人的敵人知道猶太人的觀念,所以設十字架來對付他們.以斯帖記所載的哈曼,曾設計了一個五丈高的木架,要把末底改掛在上面,但他自己卻被掛上去了.羅馬的統治者,他們要藉十字架來對付稱為基督徒的人,或罪大惡極的強盜,及與羅馬為敵的.猶太人處理自己內部的事,絕不會把同胞釘上十字架的；照樣,羅馬人,對付自己的公民也不施用十字架刑具的.傳說保羅只是被斬頭死,卻不是像彼得一樣的被掛,所以羅馬人看釘為字架是非常嚴重的刑罰.主前一○三,四年之間,猶太人起來革命,一次便有六千人被釘十字架的記錄；主前六年也有過八百人被釘,紀後六年大約在主十二歲時候,拿撒勒平原也有二千猶太人被釘的歷史.所以在主耶穌的思想裡,祂對十字架有很深的印象,祂不但說自己要被釘十字架,並且說跟從祂的人也要釘十字架.所以在當時跟隨祂的人,對於十字架的認識,不會把它視為如今裝飾品一樣的.當門徒承認祂是基督,是永生神的兒子之後,祂就指示門徒要上耶路撒冷受苦被殺；連跟從祂的,也要決定生死的問題.因為羅馬政府要以政治力量對付他們,撒都該人也附從羅馬政府.而法利賽人也因為主耶穌自稱為神的兒子,把祂看為褻瀆神的.所以猶太人絕不會同情耶穌和祂的門徒,他們正如羊進入狼群一樣.
\newline
\newline
多少時,我們會怪責彼得,以為他剛承認耶穌是基督,是永生神的兒子,但聽說耶穌要被釘十字架就攔阻耶穌,勸祂說:「主啊!千萬不可.」我相信若是我們生在那時代,也會像彼得一樣的說話.耶穌當然知道彼得的心情,按常情祂是會安慰他的,但祂對彼得說:「撒旦退我後邊去吧!你是絆我腳的,因為你不體貼神的意思,只體貼人的意思.」聽祂說這樣的話,會比甚麼都難過,那麼我們是否體貼神的意思,能否放下一切跟從祂呢?法利賽人反對耶穌,當然有他們的看法,他們是為了自己的利益打算.羅馬的巡撫恐怕生亂,當然不態容忍耶穌.但那些上來過節的人前幾天纔高呼和散那歸於大衛的子孫,現今卻又大喊釘祂十字架,真是令人費解了!這種情形,歷史上也常見的,因為人心就是那樣反覆無常.
\newline
\newline
神的兒子,本不應被咒詛,祂被釘在十字架上,是以無罪的,代替有罪的,擔當了世人的罪.保羅曾說:「在罪人中我是個罪魁,然而我蒙了憐憫.」保羅本來是反對主的,但在大馬色路上,被主光照後,明白了主的愛,便一生奉獻為福音而盡忠至死.我們對於十字架的觀感如何?是否對它模糊而使教會蒙受羞辱?抑像彼得,回頭堅固眾弟兄?許多人承認主有大恩,卻沒有把主的恩典流露.我們不必否定物質,社會地位,學問,事業等本身價值；但對於這些和主耶穌相比較,能否像保羅一樣,看萬事如糞土,以得著耶穌基督為至寶.若然在地上生活就有喜歡.十字架的道理,是日日否定自己,十字架誠然是可怕的,但現今變為可愛的.這不是由於我們自己,也不是從教育或遺傳而來,乃是主成就了救恩；一個可怕的人與可愛的主合在一起,便變為好的了.離了主,我們甚麼都不能作.保羅的工作,當然是主所呼召,但司提反的殉道,看見復活的主坐在神的右邊,會帶給保羅很深刻的印象.司提反被石頭打死的時候,似乎見不到效果；可是這效果,乃是深而遠的,以後的保羅天天肯冒死,就足以證明了.我們讀希伯來書講及信心的偉人,如同雲彩圍繞,們的見證,吸引著我們奔走前面的道路.為此我們跟隨主為主而死,不一定要像主釘十字架,像先賢被斬頭,打死,若是一生敬虔事奉,到老安然去世,也是為主.像路加二章那兩位敬虔的老人西面,西拿,見了救主便歡樂的甘願安然離世.彼得為拿撒勒的耶穌,肯把生命擺上,司提反,保羅看見復活的主,願意為主犧牲,我們今日比他們所知所經歷的還要多,豈不是更受當激勵,負起十字架跟從祂麼!
\newline
\newline


\newpage

\section{第42屆~港九培靈會~黃聿侯~第七講}
\label{sec:1438}
\textbf{至死盡忠}
\newline
\newline
連結: \href{https://www.hkbibleconference.org/session-message/view/1438}{\texttt{https://www.hkbibleconference.org/session-message/view/1438}}
\newline
\newline
\hyperref[sec:1432]{< < < PREV SERMON < < <}
~
\hyperref[sec:toc]{[返目錄]}
~
\hyperref[sec:1409]{> > > NEXT SERMON > > >}
\newline
\newline
約翰福音 12:20-12:36
\newline
\begin{longtable}{cl}
\hline
\hline
章節 & 經文 (和合本修訂版)\\
\hline
12:20 & \begin{tabularx}{0.7\textwidth}{X} 那時，上來過節禮拜的人中，有幾個希臘人。 \end{tabularx} \\ \\ \relax
12:21 & \begin{tabularx}{0.7\textwidth}{X} 他們來見加利利的伯賽大人腓力，請求他說：「先生，我們想見耶穌。」 \end{tabularx} \\ \\ \relax
12:22 & \begin{tabularx}{0.7\textwidth}{X} 腓力去告訴安得烈，然後安得烈同腓力去告訴耶穌。 \end{tabularx} \\ \\ \relax
12:23 & \begin{tabularx}{0.7\textwidth}{X} 耶穌回答他們說：「人子得榮耀的時候到了。 \end{tabularx} \\ \\ \relax
12:24 & \begin{tabularx}{0.7\textwidth}{X} 我實實在在地告訴你們，一粒麥子不落在地裡死了，仍舊是一粒；若是死了，就結出許多子粒來。 \end{tabularx} \\ \\ \relax
12:25 & \begin{tabularx}{0.7\textwidth}{X} 愛惜自己性命的，就喪失性命；那恨惡自己在這世上的性命的，要保全性命到永生。 \end{tabularx} \\ \\ \relax
12:26 & \begin{tabularx}{0.7\textwidth}{X} 若有人服事我，就當跟從我；我在哪裡，服事我的人也要在哪裡；若有人服事我，我父必尊重他。」 \end{tabularx} \\ \\ \relax
12:27 & \begin{tabularx}{0.7\textwidth}{X} 「我現在心裡憂愁，我說甚麼才好呢？說『父啊，救我脫離這時候』嗎？但我正是為這時候來的。 \end{tabularx} \\ \\ \relax
12:28 & \begin{tabularx}{0.7\textwidth}{X} 父啊，願你榮耀你的名！」於是有聲音從天上來，說：「我已經榮耀了我的名，還要再榮耀。」 \end{tabularx} \\ \\ \relax
12:29 & \begin{tabularx}{0.7\textwidth}{X} 站在旁邊的眾人聽見，就說：「打雷了。」另有的說：「有天使對他說話。」 \end{tabularx} \\ \\ \relax
12:30 & \begin{tabularx}{0.7\textwidth}{X} 耶穌回答說：「這聲音不是為我，而是為你們來的。 \end{tabularx} \\ \\ \relax
12:31 & \begin{tabularx}{0.7\textwidth}{X} 現在正是這世界受審判的時候；現在這世界的統治者要被趕出去。 \end{tabularx} \\ \\ \relax
12:32 & \begin{tabularx}{0.7\textwidth}{X} 我從地上被舉起來的時候，我要吸引萬人來歸我。」 \end{tabularx} \\ \\ \relax
12:33 & \begin{tabularx}{0.7\textwidth}{X} 耶穌這話是指自己將要怎樣死說的。 \end{tabularx} \\ \\ \relax
12:34 & \begin{tabularx}{0.7\textwidth}{X} 眾人就回答他：「我們聽見律法書上說，基督是永存的；你怎麼說，人子必須被舉起來呢？這人子是誰呢？」 \end{tabularx} \\ \\ \relax
12:35 & \begin{tabularx}{0.7\textwidth}{X} 耶穌對他們說：「光在你們中間為時不多了，應該趁著有光的時候行走，免得黑暗臨到你們；那在黑暗裡行走的，不知道往何處去。 \end{tabularx} \\ \\ \relax
12:36 & \begin{tabularx}{0.7\textwidth}{X} 你們趁著有光，要信從這光，使你們成為光明之子。」 \end{tabularx} \\ \\
[1ex]
\hline
\hline
\end{longtable}
這次大會我在香港方面所提及的,我們的信仰像磐石的堅固；因為主耶穌在地上受苦難,並且勝過死亡,祂是復活的.信祂的人也必復活,所以我們要持守所信.至死忠心.
\newline
\newline
約翰十一章記載主耶穌使拉撒路復活,許多不信的人因此就信了主.但反對主的猶太人便緊張了,耶路撒冷城地方不大,發生甚麼事情立刻就知道,何況祂工作已經兩年多；故此宗教方面的祭司長,政治方面的羅馬人,都商議想把主殺害.就在那時候,有一班希利尼人上來過節,找腓力要見耶穌,腓力又找到了安得烈帶他們同去見主.腓力,安得烈兩個名字是希臘化的,可能他們受希臘教育,懂得希臘語文.耶穌接見他們時,他們講甚麼,聖經有記載,但從主的回話看到:(1) 人子得榮耀的時候到了.(2) 一粒麥子不落在地裡死了,仍舊是一粒；若是死了,就結出許多子粒來.後來又說:「若有人服事我,就當跟從我,我在那裡,服事我的人也要在那裡；若有人服事我,我父必尊重祂.」或許希利尼人見到祂的時候,表示猶太人既不歡迎你,何不到我們那裡.主這樣表示,大概說明不是我離此地到你們那裡,而是你們要跟隨我到我那裡.這個「那裡」就是面向十字架,死亡.如此,便得到天父的尊重.今日事奉主的人,也當像昔日的希利尼人人要請主到我們那裡,適應我們的條件,環境做法.這樣的禱求,在不違背主的旨意,有受苦的心志,主或許會答允；但應該說,求主幫助使我們能夠與你同行,主到那裡,我們也到那裡.換言之,我們緊緊靠近主.我們以耶穌為中心,而不是以自己為中心,像希伯來書告訴我們的:耶穌以自己的血叫百姓成聖,在城門外受苦；照樣我們也當出到營外,就了祂去,忍受祂所受的凌辱.因為我們既以祂為我們的頭,便應該要受頭的指揮,與頭不能脫節的.
\newline
\newline
我們被天父的愛愛到不能自已,在靈性幼稚時期,會像小孩子和父母講話；要父母的愛他,甚至要父母答允了才滿足.到靈性長進了,就再不是這樣,乃每一次新近祂就覺得祂的愛太多太大.不是再不喜歡祂的愛,卻是禁不住內心有極大感受；巴不得拼了性命來報答祂,事實上我們生命經歷越深；便越覺得主的愛太豐富太奇妙,以致自己心中有難過的感覺.「尊重」是更進一步.父母長輩愛兒女,出於天然的生命,是人性自然的流露；但為父的經驗,知識,位份都遠越過兒女,怎可能「父」會「尊重」兒女呢?我們的天父是萬有之源,祂的偉大深厚不是用言語能形容的,但祂竟尊重我們.多少時我們的禱告,是祂按照祂的應許而蒙允,也許憑祂的愛而答允,也因為「尊重」我們而答允.既然如此,我們禱告了又不徹底實行,是多麼大的罪過!假如要開佈道會,天黑下起大雨來,我們為天時祈禱,感謝祂使天氣轉晴了；倘若我們不拼命講道,怎對得起祂.為此我們得蒙天父的尊重,為主盡忠至死,一點也不希奇.
\newline
\newline
耶穌說:一粒麥子不落在地裡死了,仍舊是一粒,若是死了,就結出許多子粒來.祂自己就是這樣的一粒落在地裡死的麥子,初期的使徒,非常軟弱,後來經驗到主的復活升天:因此大大剛強,十幾粒麥子便興起來了.司提反之死也結出許多子粒,最明顯的有後來稱為保羅的掃羅.當眾人對司提反咬牙切齒的時候,他看見天開了,人子站在神邊.這樣便使我們想到主對他的尊重,因為聖經屢屢說主是坐在神右邊的,現今對司提反竟像對長輩一般站起來.雖然司提反滿有智慧和聖靈能力,按照人看是了不起的人才；但主對他的尊重,不是因為他的智慧和能力,乃如詩人所說:「在耶和華眼中,看聖民之死,極為寶貴.」(詩一一六15)神所寶貴的,並非地上的金銀財寶,也不是人的聰明才幹,而是有為祂至死盡忠的態度.主這樣「尊重」司提反,難怪他面對死亡而毫不畏懼.相反,他的面貌好像天使,相信受他影響的除保羅之外,還大有其人.保羅原是法利賽人,他明白被掛在十字架上是被咒詛的,所以他極端反對耶穌之受門徒崇拜,加以迫害.但現在經歷了司提反從容就義的場面,不能不對主耶穌作重新估價.因此,在大馬色路上受光照的時候,他會問:「主啊,你是誰.」當他知道就是當日所逼迫的耶穌,是司提反甘於為祂受死的那一位,此後他就冒生命危險為主；奉獻生命在猶太人中,外邦人中傳福音雖受極大困苦與逼迫,仍不灰心不失望.直到最後,他講出成功的祕訣:「是知道所信的是誰,也深信祂能保全我所交付祂的.」他將近要為主殉道的時候,深信因為打過美好的仗,跑盡當跑的路,守住所信的道；必有公義的冠冕為他存留,而這冠冕也賜給凡愛慕主顯現的人.他這一個信念,可能回憶起當日司提反的光景.希臘文「冠冕」就是司提反.而為主殉道為主作見證的人,如同雲彩圍繞,這在前面是榮耀的,受父神和主所尊重的；因此面對十字架的凌辱苦難,也不介意了.
\newline
\newline
司提反是一粒落在地裡死了的麥子,保羅也是一粒落在地裡死了的麥子,都結出許多子粒來.主曾設比喻,麥子和稗子會同時在田間生長,要到收割的時候,才把麥子收藏；把稗子薅出,用火焚燒.這比喻說到基督徒有真有假,混合在一起；要到主再來的時候,必有所區別.我們是真的,還是假的；不但主再來的時候逃不了祂的審判,現在自己也不能掩飾,當受洗歸入主名之下,就表示與主同死；因此在我們的生命歷程中,遇到辱凌,逼迫,死亡也要堅守所信的道,若不是以死而後已的決心面對現實,則無以證明信仰的堅定.若不足以至死不渝的態度,則無以證明愛的恆久.求主幫助,讓我們看昔日司提反,保羅等為主殉道的楷摸；成為我們往後見證主的力量.
\newline
\newline
現今中國人之能夠有機會信道,是因為早期有人為中國傳福音流血流汗.在「庚子」事件以前,西教士來到中國傳道,便懷疑這福音種子在中華文化中能否生根.恐怕中國人受孔子思想影響太深,即使接受了這道,仍然不能在生命中有變化.可是庚子之亂不但西教士為此流血捨生,中國人也有因此犧牲生命的；便證明聖靈在人心中作了工,中國基督徒可以為主作見證.一八九五年即是庚子的五年前,我們福州的古田有一群倫敦會的教士,同心祈禱,求主感動差會肯拿出錢來印刷一種以羅馬字拼音的聖經,好供給那八萬不識漢字的人.不久來了一群排外主義的長毛(義和團)他們竟去把教士殺害,只逃脫一個.走脫的西教士就寫信報告事變經過,因而感動了一位被害教士的姊妹；把自己獻上,來華成全了譯聖經為土語之壯舉.還有其他被殺害的遺屬,也願意回中國來做傳道救人的工作.我每逢讀聖經便想起外祖母當時讀的土話本,是傳教士用鮮血換回來的.也想到幾家以我最先蒙恩信主,是外祖母為我家禱告了卅七年的成果.因此我相信像雲彩般見證人當中,不但有司提反,保羅,歷世歷代的殉道者；我的外祖母也在內,因為為她也是至死盡忠的使女.
\newline
\newline
至死盡忠,不一定要死於刀槍之下；能夠堅持信仰,不屈不撓,也與殉道者同樣得主的稱讚.一九五六年一月間,南美奧卡族來了五位青年白人,是為傳福音來的；但一開始工作就給土人殺害了!然而消息傳回美國,他們的遺屬不但不憎恨那些人,畏懼他們；反之,他們更願意繼承那些青年傳道者的遺志,冒險再去那裡.他們靠著主的恩典,聖靈的能力,把福音傳開了!曾經殺害他們親屬的兇手,也接受殉道者的繼承者洗禮了,葛培理在柏林開大會時候,那些奧卡族的土人,站在眾人面前作見證說,因為那些青年肯為主流血,今天他們才能成為福音的戰士.一粒麥子若不落地裡死了仍舊是一粒,若是死了,便出許多子粒來.主耶穌若是不經過死,我們不能得救,司提反若不是壯烈殉道,保羅不會成為主重用的工人.因此,在現今的世代,基督徒必須有至死悲心的心志,才可以得著公義的冠冕!
\newline
\newline


\newpage

\section{第43屆~港九培靈會~周主培~第一講}
\label{sec:1409}
\textbf{將往何處?}
\newline
\newline
連結: \href{https://www.hkbibleconference.org/session-message/view/1409}{\texttt{https://www.hkbibleconference.org/session-message/view/1409}}
\newline
\newline
\hyperref[sec:1438]{< < < PREV SERMON < < <}
~
\hyperref[sec:toc]{[返目錄]}
~
\hyperref[sec:1410]{> > > NEXT SERMON > > >}
\newline
\newline
約翰福音 20:19-20:23
\newline
\begin{longtable}{cl}
\hline
\hline
章節 & 經文 (和合本修訂版)\\
\hline
20:19 & \begin{tabularx}{0.7\textwidth}{X} 那日（就是七日的第一日）晚上，門徒因怕猶太人，所在的地方門都關了。耶穌來，站在當中，對他們說：「願你們平安！」 \end{tabularx} \\ \\ \relax
20:20 & \begin{tabularx}{0.7\textwidth}{X} 說了這話，他把手和肋旁給他們看。門徒一看見主就喜樂了。 \end{tabularx} \\ \\ \relax
20:21 & \begin{tabularx}{0.7\textwidth}{X} 於是耶穌又對他們說：「願你們平安！父怎樣差遣了我，我也照樣差遣你們。」 \end{tabularx} \\ \\ \relax
20:22 & \begin{tabularx}{0.7\textwidth}{X} 說了這話，他向他們吹一口氣，說：「領受聖靈吧！ \end{tabularx} \\ \\ \relax
20:23 & \begin{tabularx}{0.7\textwidth}{X} 你們赦免誰的罪，誰的罪就得赦免；你們不赦免誰的罪，誰的罪就不得赦免。」 \end{tabularx} \\ \\
[1ex]
\hline
\hline
\end{longtable}
今天晚上非常感謝神的恩典,在坐中多數是青年人.青年人不但是將來的領袖,也是今日的領袖.主耶穌特別關心青年人,祂在世上工作時對青年人特感興趣.耶穌基督在世上,不但傳道而且醫；祂曾叫大麻瘋得潔淨,瞎眼得看見.三次叫死了的人復活,而三次活過來的都是青年人.香港教會的前途在弟兄姊妹身上,整個東南亞的工作在你們身上.如神許可的話,中國的基督徒應繼續寫使徒行傳.那麼,今日青年就有份於這偉大的時代中了.
\newline
\newline
剛才所讀的經文:「那日是七日的頭一日晚上,門徒所在的地方因怕猶太人,門都關了.」這是一幅很適合現今世界的圖畫.那日是七日頭一日的晚上,沒有星星,沒有月亮,是恐怖黑暗的一晚.現在是黑暗掌權,黑夜已深的時候,黑暗充滿了大地.黑社會得勢,吸毒者眾；兇殺搶劫層出不窮,許多不可告人的事在黑暗中發生.光天化日之下整個世界都被黑暗籠罩.不但有水的污濁,也有空氣的污濁.但是,整個世界我們看見有光明的將來在那裡,然在黑暗中是看不出有任何的希望.請各位原諒我引用前聯合國秘書長韓瑪紹的話.在他未離世之前,他曾說過很悲慘的幾句話.他說「我看不見這世界會有光明和平的將來,我們經盡最大努力,但仍然到了一個不可收拾的光景.除非世界在最近年來有個屬靈的復興,否則這世界將會臨到萬劫不復的地步.」今天我們看見整個世界是黑暗,人類已邁進毀滅的途徑.不但聖經告訴我們末日即將來到,許多科學家有思想的人,都在想方法防備世界大局更趨惡化.在這黑暗的晚上,我們不知究竟要往那裡去.第一次世界大戰之後,海明威寫了本書「迷失的一代」,這本書叫每個青年人從夢中醒來.第二次大戰之後,青年人對存在主義思想模糊,人類究竟從何處來,往何處去?為何要存在這世界上?今天很多的青年為世界的紛亂心裡感到煩悶.我們看見舊的已倒下而新的卻未建立,這是一個真空時代.許多青年覺得前途茫茫,到底要往那裡去呢?四百萬人擠在這小島上,心裡都有說不出的苦悶.記得在兩年前,我離華盛頓到新加坡參加聚會,我問那裡的弟兄姊妹們新加坡情形如何?他們說「我們二百萬人住在這島上如同被關在籠裡.」這是我頭一個印像,後來我離開新加坡到印度尼西工,印尼是個島國.離開印尼到越南,聽見西貢的人說他們如處於孤島上.離開西貢到菲律賓,那地是個島國.離開菲島到香港,這裡也是個島.離開香港到台灣,那裡也是個島.離開台灣到琉球,也是個島.離開琉球到日本,日本更是島國.離日回美之前我到夏威夷,又是個島.在這許多島上,人們都覺迷茫,不知所從.所謂青年人打架,中年人打牌,老年人打坐.人心覺得有不可言喻的苦悶.有個青年人在他汽車後面寫著「不要跟著我因我已迷失了」.有次我和一個長髮青年談話,覺得他心裡有說不出的苦悶.他向我發個問題「第一次世界大戰的結果是甚麼?」還未等我回答,他自己已經作出答案「第一次世界大戰的結果就是第二次世界大戰」.第二次世界大戰至今,世界已經了五十多次戰爭了.
\newline
\newline
我們在黑暗中,心裡惶惑,不知將往何處!在黑暗中,我們懼怕.記得我小時,母親差我進房間拿東西,我不敢去,但母親說我是個男孩子一定要去,我只好邊走邊唱而去.我唱歌並非快樂,乃是藉以壯膽除怯.在黑暗中,人心慌惶不定,充滿了恐懼.記得我到耶加達時,有人勸告我不要把錶戴在手上,可能連錶帶手都會失去,所以我就把錶放入袋裡.到了西貢,講完道在回途中,陪我的人也對我說不要戴手錶.我說我已知道了,早就將錶放進口袋了.兩週前到菲律賓,有位牧師也告訴我別戴手錶,他還向我顯示他失錶傷手的疤痕.
\newline
\newline
到處人心充滿恐懼.我永遠不會忘記金路德博士被暗殺的那夜,在我住的那城,五百多處的大火著起來.我趕快駕車回家,因交通阻塞,車子只能慢駛.我見路旁有個婦人抱著孩子,雙眼望天,悲嘆地說「天啊,我們究竟要到怎樣的光景?」人心慌惶不定.幾年前香港很亂,許多人到加拿大去.先到溫哥華,再到滿地可.豈料那裡的學生暴動,只好搬到多倫多.世界無一處有平安,這是晚上的時候!
\newline
\newline
為何晚上有黑暗迷茫恐懼?聖經告訴我們:「因為罪進入世界,世界就被敗壞.」罪進入社會,社會被破壞.進入家庭,家庭被破壞.進入人心中,人心變成空虛黑暗,因為罪已經把整個世界破壞了.有人說,信耶穌的人總是說「人有罪」,可是我們自己沒有看見自己有罪.現在請姊妹們回答一簡單問題「怎樣能看見自己的臉?」「照鏡」.我再問「如果眼睛很明亮,鏡子也很清楚,這樣能否看見自己臉?」「能」.「今晚在黑暗裡請你們試試；眼睛明亮,鏡子清楚,能看見自己的臉嗎?」「不能.」為甚麼?因為沒有光.親愛的弟兄姊妹!基督耶穌的道不是人的幻想,乃是神的啟示.不是人自己的想像,乃是神的光照.光一進來就能看見自己的臉.光一進來,黑暗迷茫恐懼之夜就有了轉機.耶穌來萬事都要改變.聖經說「耶穌來站在他們當中,對他們說,願你們平安.」這是黑暗中的光明,是我們前途的希望.
\newline
\newline
今天有許多名義上為基督徒的,但實際上與耶穌並無關係.保羅說:「我知道我所信的是誰.」這是很重要的.請原諒我說幾句話,青年基督徒到了外國後就跌倒冷淡,為甚麼?因他們沒站起來.他們不是曾站在那裡教主日學；在詩班唱詩嗎,為何說沒站起來?他們站起來是因香港的密度太強,前後左右把他們擠得非站不可,如果有時不來聚會,親戚朋友紛紛質問因由；一旦到了國外,四面無人,站立不住,跌下來了.親愛的弟兄姊妹!你非要和耶穌同連係不可.請原諒我說一點自己的事.我有五個孩子,三個大的我已經為他們施水禮.未受禮之前我和他們談話說「孩子啊:我們之間是父子特別的關係.我們現在這國家,甚麼都有自由,你可自由揀選你的學業職業配偶,我都給你有自由選擇的權利.」有的人對事物沒有揀選,這是錯誤的,今天的揀選影響將來的結果.我對孩子說「甚麼你都可以自由選擇,只有一件事你無需也無法選擇的.誰是你的父母難道能揀選嗎?但你必定要揀選上帝做你的上帝.每個人都要與上帝發生個別的關係」.耶穌來,萬事都要改變；耶穌來,我們有喜樂.聖經說門徒看見耶穌就喜樂了.耶穌說「父怎樣差遣我,我也怎樣差遣你們.」耶穌來,罪都得赦免.
\newline
\newline
我們雖處在這廿世紀太空時代,但那古老不變的真理,仍然是耶穌基督降世,為要拯救罪人；這話是可信的,是十分可佩服的.今天,一切的思想都以「人」為中心.很多青年人以自己為中心.別以為人的力量大,其實是很有限.太大之物固然拿不起,太小之物也拿不起.太遠的看不見,太近的也看不見.聲音太高聽不見,太低也聽不見.我們眼所能及的光非常有限,太弱太強的都看不見.有件事你不能否認的,就是你不能用自己的手抓住頭髮,使自己的腳離地,你無論有多之力也不能將自己抱起.故此,人休想用自己的力量來救自己!但耶穌來,萬事都要改變.有很多人,甚至許多學生都說,你若能解答我們這些問題,我們就信耶穌了.我說,我沒法解回.他們問,耶穌能否解答?請原諒我說,耶穌也不能解答,或者說耶穌也不來解答；因為耶穌到世上,不是單為解答這些問題來的.人生充滿問題,今天解答了,明天又有問題,明天解答了後天還是有問題.比方在黑暗中,你帶我到你家,我每摸索到的東西就問這是甚麼,你一一為我解答了.後來電燈一亮,各物依舊,但已不再成問題了,因為有光便一目瞭然了.耶穌說「我是光,凡跟從我的,就不在黑暗裡走,必要得著生命的光」.人信了耶穌,不是說生命上的困難就沒有了,雖然這些事情還在,但已不成問題了.耶穌並沒有應許我們在人生道路上不遭遇危險,但祂應許永遠與我們同在.耶穌也沒應許我們在人生的海上不遇風浪,但祂應許我們必登彼岸.祂不是叫我們逃避困難,乃是叫我們有勇氣面對困難.當我們接受耶穌,我們有了新的希望,新的能力,新的道路,新的生活.親愛的弟兄姊妹!你有否讓耶穌在你生命中居首位；在你凡事上作主?單稱他為主還不夠,應讓祂在實際上稱為你的主.我們生活需要有中心,有目標,才有力量.耶穌說,你們要受聖靈,但聖靈降臨在你們身上,你們就必得著能力.有能力為主作見證,有能力為主而活.巴不得今晚主對我們說話.
\newline
\newline


\newpage

\section{第43屆~港九培靈會~周主培~第二講}
\label{sec:1410}
\textbf{並沒有分別}
\newline
\newline
連結: \href{https://www.hkbibleconference.org/session-message/view/1410}{\texttt{https://www.hkbibleconference.org/session-message/view/1410}}
\newline
\newline
\hyperref[sec:1409]{< < < PREV SERMON < < <}
~
\hyperref[sec:toc]{[返目錄]}
~
\hyperref[sec:1411]{> > > NEXT SERMON > > >}
\newline
\newline
保羅在羅馬書中有兩次說:「並沒有分別」.研究人類學者告訴我們,人在皮膚之下99.45\%是沒有分別的；就是說人是一樣的,所以不同的只是風俗習慣和言語.
\newline
\newline
我來到香港,盼望學粵語,否則像個啞吧.不過,有人告訴我,粵語有九個音,我一想國語僅四個音.那麼,聰明人才能講粵語,腦子簡單的只好講國語.
\newline
\newline
世界上各地的習慣,有很多不同的情形.有個印度朋友告訴我,印度某地方的人,習慣和我們大不相同,比方點頭稱「是」,我們是由上而下的點頭.但他們卻是由左而右的擺動,相傳北方有個老秀才,某次,他請多位西國宣教士吃飯,飯前他起立講話；我們都是基督徒,你們都會講中國話.不過各方面頗多不同之處,中國人宴客是先吃菜後喝湯,你們是先喝湯後吃菜.
\newline
\newline
中國人叫人「過來」是用手望前面招,你們是用手往後面擺.與人見面時,你們彼此是手；我們卻是拉自己的手,這樣又乾淨又衛生.他繼續說:還有一件事我最感奇怪的,就拉介紹自己的妻子,你們為甚麼說「這是我的外婦(原來是 wife)中國人是說「這是內人」或說「內子」.世界各也的風俗習慣,實在諸多不同.
\newline
\newline
今晚我們要思想聖經所說的「並沒有分別」.猶太人和希利尼人,並沒有分別.神的義,因信耶穌基督,加給一切相信的人,並沒有分別.
\newline
\newline
聖經最古老的一卷書約伯記說:「人為婦人所生,日子短少,多有患難；出來如花,又被割下.飛去如影,不能存留.」(伯十四1-2)摩西說:「我們一生的年日是七十歲,若是強壯的,可到八十歲.但其中所矜誇的,不過是勞苦愁煩,轉眼成空,我們便如飛而去.」(詩九十10)耶穌告訴我們說「在世上你們有苦難.」(約十六33)但耶穌應許說:「凡勞苦擔重擔的人,可以到我這裡來.我就使你們得安息.」(太十一28).
\newline
\newline
今天在這世上,到處充滿愁苦.嬰兒出世便哭,除了眼淚之外甚麼都不能帶來.人離世時甚麼都不能帶去,因為連眼淚也流乾了.人一生的歷史,是用眼淚所寫的；生下來時自己哭,離世時別人哭；所以在世的一生,就是從流淚到流淚.雖然世上有很多不同的言語,有時還須繙譯.但眼淚是一種毋須繙譯的言語；眼淚所包含的是心酸,勞苦,愁煩.
\newline
\newline
從未有雙眼不流過淚的.從未有口不嘆息過的；也從未有心靈不受過創傷的.世上充滿愁苦!窮者既苦,富者難道一定快樂嗎?窮人的苦是天天地挨去,而富人的苦,卻是令他長久受不了的.有次我到某地講道,有位同學陪我去,他告訴我那處是柯達的照像機,菲林的發明地,發明者名以斯曼柯達.他是個非常富有的人,因進款極多,以他的錢是以每秒鐘計算的.誰知當他在外國旅行時他竟自殺了!可見錢並不能解決人心裡的苦.
\newline
\newline
醜人固然覺得苦.但是紅顏何嘗不是多薄命.世界上最苦的職業,不是賣勞力而是賣笑；不願笑也得笑,那才真正的苦.
\newline
\newline
老年人苦,青年人也苦.特別在這時代,青年人是最苦悶的,不知究竟當往何處去；尋找生路,尋找出路,心裡有說不出的痛苦.在我住的地方有個青年人正在作工.我問他喜歡他的工作嗎,他說「沒有辦法.」人心裡的苦悶,非言語所能形容,為什麼?聖經說:因為人心空虛.人以為已得著了世界,沒得著時,盼望得著；得到後,想來也不過是如此.因為所得的不能滿足內心的空虛.我們眼睛,有眼睛的需要,耳朵有耳朵的需要；靈魂有靈魂的需要.但物質並不能滿足靈命的需要,正如畫餅不能充飢一樣.
\newline
\newline
世界到處都有暴動,因為沒有安息,有次我在菲律賓一個家庭聚會講道,聽眾都是在銀行界辦事的人.講完道,我坐在那裡,他們還在唱詩.有一個人過來問我怎能得著安息?他說:他的內心沒安息,進來了想出去,出去了又想進來；天天喝酒,吸煙,心裡不得著安息.感謝主!那天下午他接受了耶穌,他跪下來和我一同禱告,承認他的罪,心裡就得到真安息.
\newline
\newline
人最大的愁苦是什麼?希伯來書告訴我們:按著定命,人人都有一死,死後且有審判.死是人最大的愁苦.香港有多賣人參補品的店舖,可見香港人很注重補養身體,秋風一起大進補品.人都不願意死,死如同單行道,只能一方通行,無法倒轉回來.有的人心裡恐懼,不知死後到那裡去.雖然有人買了人壽保險,但也不能擔保不死.有句話說:狗比死了的獅子強.中國人說:死了一切完了.難怪很多人心裡有惡恨的惆悵和迷茫.中國詩人蘇東坡說「大江東去,浪濤金,千古風流人物古壘西邊,人道是,三國周郎赤壁.」在人的心中,都覺得人會過去,楚霸王死了,史太林死了,希特勒死了.當我在東京站在那監獄的門口,有人告訴我,東條曾下在此監.在太平洋戰爭時他是個不可一世的人物,也不免一死.人人都有一死,並沒有分別,黃泉路上無老少.蘇聯曾發明了一種藥簡名 AES,據說服此藥可活到一百五十歲,但結果那發明者六十五歲便死了.在這世界上,死是非常公平的,不管是否有名譽,有地位,有財產,有學問,並沒有分別.
\newline
\newline
為什麼人有愁苦有死亡呢?聖經告訴我們:「死是從罪來的,於是死便臨到眾人,因為眾人都犯了罪」.雖然人都不承認自己有罪,但人天天過著受罪的生活.聖經說:「世人都犯了罪,虧缺了神的榮耀.」(羅三23)人心比萬物都詭詐,壞到極處.(耶十七9)世界各地都有醫院,監獄,墳場.這就證明你不能否認在這世界上有病人,罪人,死人.聖經說人人都有罪.今天有許多未信耶穌的朋友們,或者要說「我並沒有坐牢為何要說我犯罪?我到底有何罪?」我作個比方:最近我在菲律賓,那邊宿務的芒果是最有名貴的.他們帶我去看芒果樹,我問說:「當小孩子吃完了芒果,把核扔在泥地上,過些日子,生枝長葉,這棵小樹是否叫芒果樹?」答「當然是」那麼,大的芒果樹已經開花結果叫它芒果樹,而那小棵剛開始生長的,你已經叫它芒果樹了.不是因為它開芒果花,結成芒果然後才叫它芒果樹；乃是因它是芒果樹才開芒果花結出芒果.好多時候,我們不是怕犯罪,乃是犯了罪怕人知道,所以才暗暗地犯罪.誰也不肯說「我犯了罪」.有的說:我是有道德的人.有的說:我是有學問的人.有的說,我是有財產的人.有的說,我是樂善好施的人.為何說我是罪人?我告訴你罪人還是罪人.請原諒我作個比方:你是中國人,我也是中國人.你是住在香港的中國人,我是住在外地的中國人.有英籍的中國人,有美籍的中國人,中國人還是中國人.從前的中國人,走路四方腳.說話之乎者也,住中國房子,吃中國飯.現在中國人講洋話,住洋樓,吃洋菜,坐洋車,穿洋服,中國人還是中國人.有名譽,有地位,有學問,有財產,罪人還是罪人.你是受過教育的罪人,你是樂善好施的罪人,罪人還是罪人.絕不因外表而改變內心.無論住九龍或香港,任你到了天涯海角,你若不接受主耶穌的寶血,不接受主耶穌在十字架上為我們所成就的救恩,罪人還是罪人.
\newline
\newline
神愛世人,不分男婦老幼,不論黑白黃棕各種族,並沒有分別.神愛世人,甚至將祂的獨生子賜給他們,叫一切信祂的不至滅亡,反得永生.並沒有分別.親愛的弟兄姊妹們!雖然這是一件很古老的事,但每逢我們思想到得救的事,我們的心如同充滿了火一樣,我們看見了神的愛.
\newline
\newline
親愛的朋友們!人活著不是單靠食物,人是靠著愛來生存.沒有愛的人生,是不完全的.為什麼有的人心裡充滿仇恨?因他們沒有經歷過愛的人生,沒有愛的人生是痛苦的.聖經告訴我們,神的愛是何等的長闊高深,包羅一切.神的愛叫祂兒子離開天上的榮耀.除了神的愛,沒有任何力量能叫耶穌降到地上；也沒有任何力量,能將祂從地上掛在十字架上.因著愛耶穌來到世上,尋找拯救失喪的人.
\newline
\newline
今天世上多多少少的人,好像壓傷的蘆葦,將殘的燈火.充滿各種的愁苦.人知道有一天必離開世界,知道自己有罪.神的愛解決我們一切的問題,能叫黑暗的心靈變成光明；能擦乾流淚的眼,止息嘆息的口,包裹創傷的心；叫下垂的手發酸的腿,再次挺起來,奔走前面的路程.
\newline
\newline


\newpage

\section{第43屆~港九培靈會~周主培~第三講}
\label{sec:1411}
\textbf{壓傷的蘆葦將殘的燈火}
\newline
\newline
連結: \href{https://www.hkbibleconference.org/session-message/view/1411}{\texttt{https://www.hkbibleconference.org/session-message/view/1411}}
\newline
\newline
\hyperref[sec:1410]{< < < PREV SERMON < < <}
~
\hyperref[sec:toc]{[返目錄]}
~
\hyperref[sec:1412]{> > > NEXT SERMON > > >}
\newline
\newline
馬太福音 12:18-12:21
\newline
\begin{longtable}{cl}
\hline
\hline
章節 & 經文 (和合本修訂版)\\
\hline
12:18 & \begin{tabularx}{0.7\textwidth}{X} 「看哪，我所揀選的僕人，我所親愛，心所喜悅的；我要將我的靈賜給他，他必將公理傳給外邦。 \end{tabularx} \\ \\ \relax
12:19 & \begin{tabularx}{0.7\textwidth}{X} 他不爭吵，不喧嚷，街上也沒有人聽見他的聲音。 \end{tabularx} \\ \\ \relax
12:20 & \begin{tabularx}{0.7\textwidth}{X} 壓傷的蘆葦，他不折斷，將殘的燈火，他不吹滅，直到他使公理得勝。 \end{tabularx} \\ \\ \relax
12:21 & \begin{tabularx}{0.7\textwidth}{X} 外邦人都要仰望他的名。」 \end{tabularx} \\ \\
[1ex]
\hline
\hline
\end{longtable}
這段經文,是論到我們的主耶穌基督,到世上來的使命:為要叫流淚的眼能夠擦乾,嘆息的口能夠止住,創傷的心得著包,下垂的手；發酸的腿能夠挺起來,奔跑前面的道路.
\newline
\newline
「壓傷的蘆葦祂不折斷」這是一幅很有意義而且充滿詩意的圖畫.平靜的河水,襯著藍天白雲,河邊蘆葦,臨風招展,優美至極!蘆葦春生秋枯,生長時外表堅固如竹,實質脆弱不堪.蘆葦如不被壓傷,略有一些用處；如果壓傷了,人就把它折斷.這是今天世界的光景.
\newline
\newline
雖然許多人不承認我們是被壓傷的蘆葦,不過聖經告訴我們,撒旦把全世界的人,都壓在牠手下.今天很多人以為他是自由的,可以任所欲為.可是有很多的人,卻被很小的罪所捆綁.犯罪時糊糊塗塗,犯了罪之後清清楚楚.說,下次我再也不犯了,一旦引誘試探來到卻又糊塗了.
\newline
\newline
今天多少人不知不覺中,被撒旦用罪惡的力量控制.離我住的城不遠,是江水和海水匯流處.海軍部有個碼頭在那裡,晚間燈光燦爛,照得螃蟹都爬出來了.有位弟兄在海軍部任事,他常常請我去捉蟹.有一晚我們捉了很多螃蟹,他告訴我螃蟹不能久藏,要即刻煮熟然後放入冰箱.我回家馬上煮了一大鍋開水,把螃蟹倒下去,結果全部爬了出來.第二天他問我螃蟹如何?我回答說「滿頭大汗,滿屋螃蟹」.我把煮蟹的過程告訴他.他聽完了,說「你錯了,應先將螃蟹倒在冷水裡,然後才點火,慢慢地熱起來,等牠們發覺時已太遲了,已爬不動了」.人犯罪也是如此,起初不覺苦,因為魔鬼把你置於安舒適的地方,你以為是享受,你以為這樣的生活太有意義了；殊不知下面的火點起來,你便逃不出來了.
\newline
\newline
人似被壓傷的蘆葦,從約翰福音中可看到有二人都是被壓傷的.一個是撒瑪利亞的婦人,她被罪孽所壓傷.不但別人知道她是個罪孽深重的女人,連她自己也知道.她為了避免別人看見,所以揀了一個沒有人打水的時候,出來打水.聖經告訴我們,她已經有五個丈夫,現有的還不是她的.這個女人正是代表今天的社會,今天世界的人崇拜性慾的神.性與色已達世界每個角落.月曆牌如果沒有女人的圖畫,就沒有銷路.打開報紙,很多廣告都用色情來吸引人,今天的人已到色情狂的地步了,正像挪亞時代一樣,人被色情所迷惑.魔鬼敗壞世界的方法,就是讓人沉迷在色情裡.挪亞時代,人終日所想的盡都是惡.今天人思想裡面除了色情之外沒有別的.羅馬將亡時,物質的享受到了最高峰,人就在色情中打滾.有許多人被色情所壓傷,他們得不著釋放的自由.耶利米書六章十五節就是說今天的光景:「他們行可憎的事知道慚愧麼,不然,他們毫無慚愧!也不知羞恥,因此他們必在仆倒的人中仆倒.我向他們討罪的時候,他們必致跌倒,這是耶和華說的」今天的人,已經不知何謂羞恥何謂慚愧.從前造的房屋和東西都是又厚又重,經久耐用.現在造的東西卻是又薄又輕.連人的道德也是到了又輕又薄的地步,人已經被罪孽所壓傷.
\newline
\newline
撒瑪利亞婦人,同時也為驕傲所壓傷.從她的說話可以看出,她認為她有打水的器具,有祖先,有宗教.今日人類就在驕傲的光景中,卻不知已被驕傲所壓傷.今天的人把自己當作神.在這末世時代,人所想的就是為了自己,所愛的,所體貼的也就是自己,以為自己了不起,結果被壓傷.
\newline
\newline
(聖經告訴我們另一人也是被社會的制度所壓傷,思想上受到別人的影響.)我曾經和一位不信有神的青年學生談道,感謝主恩典,最後他說「我現在明白了,如果我能知道一切關於神的事,那我就是神了.因祂是神,我是人,我承認有許多不的懂的.」那晚他接受了神.一點不錯,如果你能完全明白神的事情,那你就是神了.若你承認你不過是人,那有許多關於神的事,你就不能全明白.
\newline
\newline
尼哥底母是個有道德,有地位,有名譽,有學問,有年紀的法利賽人.可是他被社會制度所壓傷,他到耶穌面前.起初他以為自己完全知道,他說:「我們知道,你是由神那裡來作師傅的.」但主的回答卻是:「人若不重生,就不能見神的國.」照樣,現在也有許多基督徒,知道一點關於耶穌的事；但從未不知道耶穌自己,沒有和耶穌發生個別關係.有許多青年人 ,只是口說信耶穌,心卻不知所信是是誰.許多人在教會裡長大,但沒有真正的認識主；都被制度所壓傷,受別人所影響.
\newline
\newline
撒該願意看耶穌,卻看不見.人都以為是因他矮所以看不見,但聖經說他看不見是因人多的緣故,矮只是其次的原因.今天也有很多人沒有清楚認識主,卻在教會中自人,以為自己了不起,人人都不及他.
\newline
\newline
尼哥底母以為自己年紀高,當耶穌說到重生的時候.他說:「人已經老了,如何能重生呢?」他將神的道,看得太簡單太幼稚了.耶穌講的是天上的事,尼哥底母卻當作是地上的事,結果他聽來聽去不是聽不懂.這些人正是井底青蛙一樣,以為宇宙只有井口這麼大.聖經告訴我們,屬肉體的人,不能領會屬靈的事.據一位在菲洲內陸傳福音的西教士說:「那地的人從來不知道水會結冰,更不相信汽車能在路上開.有些人從未見過巴士,不明白巴士是什麼,對他們說會走路的房子可能還懂一點.」所以,耶穌的真理,不是人自己的想像,乃是神的啟示.當神啟示人的時候,人就能明白.人若未得神的啟示,就如同在黑暗中.
\newline
\newline
聖經說,世上的人不單像壓傷的蘆葦,也是將殘的燈火.從前猶太是用油來點燈的,油足時,燈光很亮,油將盡,燈光逐漸暗淡,烏煙愈增,於是人就把燈吹熄.
\newline
\newline
科學家告訴我們,人類已到了非常危險的光景.現在歐美很多大學有一門最吃香的科目,是空氣和水的污染.在工業發達的城市,再過幾十年後,可能人出門行路也要戴上防毒面具,這不是笑話.藍色多腦河的水已經由藍色變成黑色了.德國的萊茵河,已變成一條大溝渠了.香港九龍的人真有福,今日還有清水灣.
\newline
\newline
在這世界上,我們看見人類的毀滅,非但聖經告訴我們人類已到毀滅的邊緣.自去年始.有些地方已禁用 DDT,因為它殺害了無數的飛鳥和魚,虫.科學雖然給人類帶來許多方便,但也帶來許多煩惱.人好像將殘的燈火,過一分鐘少一分鐘,過一秒鐘少一秒鐘.人的一生是個減數,過多一天就少了一天.世界上很多有思想的人,都感到他們的生命在減縮中.海明威在他的一本著作中說:「我生活在真空管中,我像真空管一樣孤寂；我的電池已經乾了,我的電源已經枯竭了.」好像作詩的人摩西說:「主啊,你世世代代作我們的居所,諸山未曾生出,地與世界你未曾造成,亙古到永遠你是神.你使人歸於塵土,說,你們世人要歸回.在你看來千年如已過的昨日,又如夜更一更；你叫他們如水沖去,他們如睡一覺,早晨他們如生長的草,早晨生長,晚上割下枯乾.......」
\newline
\newline
人的一生,就像一盞將殘的燈火,已經毫無用處.耶穌曾經講過三個比喻:一百隻羊,失去一隻.有十塊錢,掉了一塊.有兩個兒子,失去了一個.一隻走迷的羊好像將殘的燈火,牠的生命面臨危險的境地.失去的錢,如果不找回來,它的價值就失去了.那失落的浪子,也好像將殘的燈火一樣.
\newline
\newline
耶穌到世上來,壓傷的蘆葦,祂不折斷；將殘的燈火,祂不吹滅.聖經告訴我們；全世界人的罪都歸在祂身上,連主也被罪所壓傷.因為祂被壓傷,所以祂能叫壓傷的蘆葦,再次挺起來.祂將命傾倒,以至於死,祂自己好像將殘的燈火一樣；人以為祂已經死了,已經完了.感謝主!祂復活了,因為祂的復活,帶給我們新的希望,新的能力,新的道路,看見新的異象,生活在世界上有新的盼望.
\newline
\newline
壓傷的蘆葦,祂不折斷,將殘的燈火,祂不吹滅.
\newline
\newline


\newpage

\section{第43屆~港九培靈會~周主培~第四講}
\label{sec:1412}
\textbf{失喪的人}
\newline
\newline
連結: \href{https://www.hkbibleconference.org/session-message/view/1412}{\texttt{https://www.hkbibleconference.org/session-message/view/1412}}
\newline
\newline
\hyperref[sec:1411]{< < < PREV SERMON < < <}
~
\hyperref[sec:toc]{[返目錄]}
~
\hyperref[sec:1413]{> > > NEXT SERMON > > >}
\newline
\newline
路加福音 19:10-19:10
\newline
\begin{longtable}{cl}
\hline
\hline
章節 & 經文 (和合本修訂版)\\
\hline
19:10 & \begin{tabularx}{0.7\textwidth}{X} 人子來是要尋找和拯救失喪的人。」 \end{tabularx} \\ \\
[1ex]
\hline
\hline
\end{longtable}
主耶穌親口宣佈祂的使命說:「人子來,為要尋找拯救失喪的人.」
\newline
\newline
保羅以他個人的經驗作見證說:「基督耶穌降世,為要拯救罪人.......」(提前一15)
\newline
\newline
「神愛世人,甚至將祂的獨生子賜給他們......」(約三16)以上幾節經文都提到「人」.
\newline
\newline
一九六八年十二月,有幾個人上了月球,在月球上走.美總統尼克遜那時有點得意忘形.他很神氣的說:「這是人類有史以來最偉大的成就.」次日,他的好朋友,也是神忠心的僕人──葛培理博士對他說:「總統先生,歷史最大的成就不是人類登月,乃是神降世.」人類登月是幾千年人的夢想,近數十年來花費了幾十億美金,動用了四萬多人力,把人送上月球.結果,僅帶回一些石塊,這就算是人的成就了.
\newline
\newline
神子降世乃是亙古之前神的計劃.祂將最重要的生命,為人傾倒出來.祂親自降世,為要將人從罪惡中救出來,進入神的永生.相形之下,人類登月不過是件渺小的事.歷史上最大的成就是──我主耶穌離開天上的榮耀,到這卑微污穢的塵世.
\newline
\newline
耶穌為何要到世上來?希伯來書十章五至七節告訴我們:「所以基督到世上來的時候,就說,神啊,祭物和禮物是你不願意的,你曾給我預備了身體.燔祭和贖罪祭是你不喜歡的.」那時我說:神阿!我來了為要照你的旨意行.」約翰一章十八節說:「從來沒有人看見神,只有父懷裡的獨生子將他表明出來.」約翰一書三章十八節說:「神的兒子顯現出來,為要除滅魔鬼的作為.」耶穌降世有很多目的,但有一目的我們不可否認,耶穌到世上,是為了人的緣故,就是為了像你我這樣的人.馬太一章廿一節告訴我們「......你要給他起名叫耶穌.因祂要將自己的百姓從罪惡裡救出來.」約翰三章十七節說:「神差祂的兒子降世,不是要定世人的罪,乃是叫世人因祂得救.」腓立比二章六至八節說「他本有神的形像......反倒虛己,取了奴僕的形像,成為人的樣式......就自己卑微,存心順服,以至於死,且死在十字架上.」所以我們看到耶穌到世上來之目的:(1)非為生,非為活,乃是死.(2)將神彰顯出來,叫神和人溝通關係.(3)為我們開了一條又新又活的路,領我們進入榮耀裡.(4)為要尋找拯救失喪的人.
\newline
\newline
人對於人本身的觀點:有一班人,他們以為人是猴子變的.我告訴弟兄姊妹,我從來沒見過猴子所變的人.可是我看見現在的人已經慢慢地變成......(恕我不講下去)我今天不和各位講進化論,不過略為一題).如果人是猴變的,那麼猴子一定還繼續地變,難道猴子認為已經夠了不必再變,試停止不變了嗎?如果一直地變,說不定變成似猴非猴,似人非人的樣子.有人說人和猴有許多相同之處,但不能因此說人是猴子變的,不過表明是同一主宰創造的.好像美國通用汽車製造廠造的汽車,大小種類很多,名稱各不同.你不能說這種汽車能進化成那種汽車,或說,低級的能進化成高級的.你只能說是同一廠家製造的.至今為止,猴子的血,不能輸於人體,而人的血,無論任何種族；只要血型相同,便可以彼此輸血.可惜!人偏要承認自己是猴子的子孫,但猴子卻不承認.猴子們似乎在談論說:「我們猴子中那兒有如此不肖的子孫,那兒有把小猴丟在一邊而自己去打麻將的猴媽媽；那兒有多下小猴和母猴,而顧著自己去跳舞的猴爸爸?我們猴子要吃新鮮的水果,那裡會把樹葉晒乾,然後放在嘴邊燒的呢?還有那些大麻；迷幻藥,我們猴子都不吃的.」
\newline
\newline
在這世界上,人越過越苦悶.剛才我看了一本書說的很好.我們這一代的人最痛苦,我們年青時受教要聽老年人的話.我們年紀大起來了,現在要聽青年人的話.或年青或年老都要聽人家的話,卻沒有人聽我們的話.這個世界實在變動的很厲害!
\newline
\newline
人究竟像甚麼呢?耶穌說:「人像迷失的羊,像失落的錢,像離開父親的浪子.」以賽亞書說:「人好像羊走迷了路,各人偏行己路.」
\newline
\newline
不久以前,我聽見一位心理學家分析今天人的光景.我覺得正是聖經裡所給我們的一切答案.(一)虛空 emptiness.人心虛空,得不到滿足.剛出世嬰孩,兩手抓得很緊,這表示說「我來撈世界」,世界有金子,銀子,車子,妻子,兒子,老了還有鬍子.(不過現在的青年人也有鬍子)我覺得人好像小貓,團團轉,拼命抓,抓不到,心煩燥；最後抓到的,只是自己一條尾巴.(二)無安息 restless.為什麼現在有很多暴動,很多人遊行.在我住的城裡,幾乎天天有人遊行.你是反對的,我是成贊成的,人的內心沒有安息.不久以前,我曾和一個人談話,他父親是銀行家,可以說,他有世界上的一切.但他說他的心沒有安息,內心痛苦,他常飲酒消愁,甚至服過迷幻藥,服後飄飄欲仙,清醒時,仍然痛苦在現實中.(三)無目標 aimlesss.雖然我們比祖先有更好的交通工具,但是卻沒有更好的去處.今天的人,力求迅速,最新式速度最快的飛機,每小時可飛二千里.可憐今天的人已經失去了目標.只求快速而失去了目標,結果就是毀滅.人心杳茫.前不見古人,後不見來者；念天地之悠悠,獨愴然而淚下.前途沒有目標,令人驚慌失掉.親愛的弟兄姊妹!親愛的朋友!今天人如羊走迷了路,前途茫茫.(四)孤獨 Ioneliness.有人說廿世紀人類最大的病疾是孤獨.你有否看見那些高樓大廈?那是人類最高孤單的紀念碑.一幢大廈,可能住了成千上萬的人,但是人與人卻沒有來往.怪不得二次大戰之後,存在主義這般流行,人心覺得孤單.親愛的朋友們!詩篇廿三篇說:「耶和華是我的牧者,我必不至缺乏.」我的內心不再空虛了,我什麼都夠了.「祂使我躺臥在青草地上,領我在可安歇的水邊.」我的內心有了安息,不再勞苦愁煩了.「為自己的名引導我走義路.」我們不再沒有目標,因祂像牧者引導祂的羊.「我雖然行過死蔭的幽谷,也不怕遭害,因為你與我同在.」有神同在,我們不再孤單了.
\newline
\newline
罪叫我們與神隔絕,因著罪內心空虛,因罪內心不得平安.聖經說:「惡人的心不得不安.」奧古斯丁曾說:神啊!為了你自己而創造我們,如果我們不回到你裡面；我們的內心,就永遠得不到安息.」主耶穌非常關心你,雖然你不關心自己.當人犯罪以後,人已失去應有的價值；失去了安息和自由,但可憐人自己卻不承認.有些統治階級,不把人當人,乃把人當工具.有些獨裁者,不把人當人,乃把人當炮灰.在醫藥界,卻把人當玩物.請原諒我說,今天在世界上,人已非人了.住在大城市的人,在政府眼光中如同一個號碼,因為現在電腦發達.但電腦不認姓氏男女,只認號碼.今天人已失去應有的價值!各位親愛的弟兄姊妹們!主耶穌,為了像你我這樣的人,離開天上的榮耀來到卑微的塵世,就為了要拯救你.在耶穌的比喻中,祂把人的價值漸漸提高.本來是百羊失一,再而十錢失一,後來二子失一.一次是百分之一,二次是十分之一,三次是二分之一.而且那回來的兒子,在父親的眼光中,是失而復得,死而復活的.耶穌說:「在猶太地方,一分錢可以買兩隻麻雀；兩分錢可買五隻.如果神不許可,連一個麻雀,也不會掉在地上.」每一個人都是耶穌所愛的.最後耶穌告訴我們,那失了一個兒子的父親,天天等候.他家雖有許多僕婢牛羊財產,但這一切都不能滿足他的心；除非失了的兒子再回來.多少時候,人自暴自棄,行屍走肉,怨天尤人；不自知在神眼光中,是非常重要的.人若賺得全世界,而賠上自己的生命,有甚麼益處呢?因為人的裡面,有個寶貴的靈魂.故此,耶穌為你釘在十字架上.雖然羊迷失了,錢失了價值,浪子失了榮耀.聖經告訴我們,當浪子回到家中,他父親把上好的袍子,把戒子,鞋子穿戴在兒子身上.
\newline
\newline
親愛的朋友們,我們真感奇妙,為何耶穌從天上降世又從地上被掛十字架上?祂本滿有神的形像,成了人的樣式,自己謙卑,存心順服,以至於死,且死在十字架上.祂在世上時曾說:「凡勞苦擔重擔的人,可以到我這裡來,我必使你們得享安息.」(太十一28)
\newline
\newline


\newpage

\section{第43屆~港九培靈會~周主培~第五講}
\label{sec:1413}
\textbf{求神鑑察}
\newline
\newline
連結: \href{https://www.hkbibleconference.org/session-message/view/1413}{\texttt{https://www.hkbibleconference.org/session-message/view/1413}}
\newline
\newline
\hyperref[sec:1412]{< < < PREV SERMON < < <}
~
\hyperref[sec:toc]{[返目錄]}
~
\hyperref[sec:1414]{> > > NEXT SERMON > > >}
\newline
\newline
我相信我們中間凡是愛主屬於主的,都有同樣的感動,今天的教會需要復興.
\newline
\newline
雖然有許多教會,好像士每拿一樣愛主,至死忠心；可是有好多教會,卻像以弗所一樣,他們失了起初的愛心.雖然有的教會像非拉鐵非一樣,他們愛主愛弟兄.可是我們不能否認,有很多教會像老底嘉一樣；不冷不熱,他們物質豐富,愛心卻冷淡了,還有很多的教會,按名是活的,其實是死的.所以不但未信主的人要悔改信福音,基督徒同樣要悔改.悔改不止一次,要常常在神面前求神鑑察.求神復興今天香港的教會,復興我們每一個人.如果願意未得救的人信主,那已信主的人應當悔改.若不是主的愛吸引,就沒有人到祂的面前.
\newline
\newline
有很多所謂基督徒者,竟站在禮拜堂門口攔阻了人進來.那些假冒為善,有名無實的基督徒,徒有基督徒的外表,而無基督在內心.(請原諒我用這句話告訴大家:我用眼淚來預備今晚的講章).巴不得勝靈在我們當中工作.我不敢說是對大家講的,也是對我自己講的.我在神面前說:「神啊!求你鑑察我,知道我的心思」.
\newline
\newline
多少時,我們在禱告中取巧；在求恩典的時候,就求神賜福,恩待帶領.認罪的時候,就禱告說:神啊!我們都是罪人,驕傲,欺騙,詭詐.......連別人都扯進去.「我」和「我們」完全是兩件事.你說你已認了罪,可是你沒有在神面前,個別的認識祂,我們只會講道給別人聽,禱告給別人聽；為別人認罪,卻沒有在神面前省察自己.
\newline
\newline
路加福音第三章,可以看到當時的社會處於四分五裂,天昏暗地的光景.由一至六節可以看出四種不同的代表.我們用四種不同的人物,來代表那時代的光景；同時,也代表今日的世界,社會和教會.
\newline
\newline
(1)該撒提庇留　　是羅馬的王.當耶穌降生時,羅馬的勢力很強大,可說天下是屬於羅馬的.那時,羅馬立國近二百年.據歷史家告訴我們,那時羅馬國的奴僕,超過人民的數目.因奴僕都是從外地擄來的,最窮之家也有奴僕,所以當時羅馬的物質享受,已達最高峰.當日他們強奢極慾,今日我們的情形豈不一樣嗎?港九到處是餐館酒家,隨時門庭若市；打開報紙,吃甚麼,穿甚麼,看甚麼都有.人在這世界要得到聽覺,視覺,嗅覺,各種感覺上的享受.這就是羅馬,該撒提庇留的代表時代.
\newline
\newline
(2)本丟彼拉多　　他是羅馬的巡撫,被派統治猶太地.他是個聰明的政客,投機耶巧,對上極事奉承.他步步升官,日日見財,官財相逢,他對下卻殘暴.為人沒有主見,好像蘆葦隨風飄搖.但他卻善於順風駛舵,當猶太人說你若釋放耶穌,你就不是該撒的中神時,他就和猶太人妥協.為了要維持個人的權力和地位,他竟向猶太人投降.在猶太人面前洗手,他的夫人對他說:「這義人的事,你一點不可管.」他懦弱以為因環境的逼迫,不得不這樣做.雖然聖經沒有明載,但歷史告訴我們:他的結局,是悲慘的,身敗名裂,被充軍而死.
\newline
\newline
(3)希律王　　他是分封的王,在民事方面他有權柄.耶穌說他是「狐狸」「詭詐的人」,他不擇手段,從他弟弟手裡奪了弟婦,佔為己有.為了這事,他殺了施洗約翰的頭.今日世界上,貪污兇殺的事很多,港九數月來已發生了幾百宗的兇殺案.在這時代,人心充滿仇恨兇殺,正如昔日的光景.
\newline
\newline
(4)亞那和該亞法　　他們是大祭司.一為岳父,一為女婿,二人狼狽為奸,把持宗教的權柄,勾結商人.容許賣牛羊鴿子,找換銀錢的人,在聖殿做買賣；以至神聖的殿,成為污穢的地方.
\newline
\newline
親愛的弟兄姊妹們!教會本應社會分別出來,但社會中有的,教會中都有.且披上宗教的外衣,令人撲朔迷離.無怪許多人反對基督教,許多人不願到禮拜堂,此其因也.所以我們必須悔改!我們不單要請未信主的人接受耶穌,我們應該讓主耶穌在我們心裡掌王權,戴冠冕,坐寶座,居首位,求主今晚鑑察我們.
\newline
\newline
約翰出來傳悔改的道,他在曠野呼喊:「預備主的道,修直祂的路.」願主耶穌也在我們中間作祂奇妙的工作.
\newline
\newline
親愛的弟兄姊妹!為可今天教會沒有復興?為何教會還是不冷不熱?(我已經為這事禱告了很久)因為我們把起初的愛離棄了.聖經告訴我們有四件事情,我們必須注意的.
\newline
\newline
(一)一切山窪都要填滿.山窪代表隱藏的罪.山窪是代表陷落離地面,遠遠看不見,而臨近才發覺的罪.多少時候,我們有隱而未現的罪,外表道貌岸然,內心卻充滿罪惡.我們不是怕犯罪,而是犯了罪怕人知道.有個在教會中很有名的人,不久之前自殺了,原來他是醉酒者.別人都不知,後來給人發現,他無臉生存下去.當我看到這新聞時,我心裡警惕自問,我們是為神工作的僕人,或作傳道,牧師,或作長老,執事；我們的思想和內心,是否充滿了罪惡?耶穌說:這是假冒為善的「法利賽人」,看來一表斯文,衣著高尚,佩戴經文,開口選美主!合口感謝主!但心靈裡卻充滿了罪惡.弟兄姊妹!我不願一一提及,我們在錢財上,女色上的失敗,在很多事上不能榮耀主.你自己知道,神也知道,連撒旦都知道,我們羞辱了主的名.除非我們悔改認罪在神面前,教會才能得復興.我們是否一同聚會,但從來不打招呼；雖一同事奉,卻彼此各有心病.隱藏的罪──山窪要填滿.
\newline
\newline
(二)大小山崗都要削平.山崗是高突出來的,就是我們心裡的驕傲.有人說,這是屬靈的驕傲.我覺得這話不合邏輯.屬靈的人不會驕傲,既驕傲就不屬靈.今天我能夠為主工作,完全是主的恩典,恩典從來不會叫人驕傲的.如果你知道一切所得的,都是主的恩典,有何驕傲可言?連我們的存留氣息生命都在乎祂,有何可驕傲呢?
\newline
\newline
人在艱難的時候,依靠神,仰望神；到了豐衣足食食時,說,耶和華是誰呢?當你有病時,你曾經禱告許願,等到病癒了,你就以為自己僥倖.當你卑微時,你求告神；得意時卻忘記了神.一個發怒的人,就失了自己.嫉妒的人就失去朋友.驕傲的人,就沒有神.我們是否法利賽人?嘗說:「神啊:我感謝你!我不像別人......我一個禮拜禁食兩次；凡我所得的,都捐上十分之一.......」在神面前自以為義呢?這是魔鬼的方法.聖經告訴我們,當初魔鬼犯罪,不過是從心裡起了驕傲,要高舉牠的寶座,要與至高者同等.
\newline
\newline
(三)彎彎曲曲的地方,要改為正直　　就是我們心內的詭詐.很多時候,我們是心口不一的.「主說,因為這百姓親近我,用嘴唇尊敬我,心卻遠離我.......」(賽廿九13).耶利米說:「他們的口,是與神相近,心卻與你遠離.」以西結說:「他們的口多顯愛情,心卻追隨財利.」這些也是今日的光景.口說:「主啊!我心愛你.」心裡想著三一三十一,逢三進一.甜言蜜語；心裡卻如猶大一樣.聖經說猶大是個賊,常取其中所存的.弟兄姊妹!不要作假冒為善的人,應當心口如一.有人說:基督徒撒謊最多的地方,是在禮拜堂裡.為許多人在那裡說:主啊!我心愛你.口唱:我要跟隨你.其實心卻不然.
\newline
\newline
雅各是個詭詐多端的人,他不但對人行詭詐,向神也耍花樣.他跪下來禱告說:「神阿,你賜給我的恩典,我真不配得；從前一根杖過河,現在成兩大隊.當他禱告完了,馬上打自己的算盤,把妻兒財產分成兩隊,將他最愛的放在最後,預備和他哥哥相打.所以神要在毗努伊勒對付他!
\newline
\newline
(四)高高低低的的道路要改為平坦　　基督徒也有「不信」的罪.我們的知識和信心,不能打成一片.知道是一件事,相信又是一件事.那些文士和法利賽人,他們知道基督當生在何處,可是他們沒有去敬拜主.很多人聽了主的話沒有遵行,為何?因為不信,相信才能順服.信甚麼?約翰四章告訴我們:「信耶穌所說的話.」許多所謂信徒的,沒有真正相信耶穌,沒有真誠跟隨基督,而且常常得罪神.聖經說我們有不信的惡心.勸別人信主,自己卻不接受.該亞法大祭司曾經為耶穌作見證(雖然他的話別有用意),當耶穌被釘十架前,他曾說:「獨不想一個人替百姓死,免得通國滅亡就是你們的益處.」(約十一50)他的見證是說給別人聽的,不是對自己講的,恐怕很多人講道也是如此.
\newline
\newline
我們若要看見神的救恩,一定要付上代價.一切山窪都要填滿,沒有隱藏的罪.大小山崗都要削平,不再驕傲.彎彎曲曲地方要改為正直,不再詭詐.高高低低的道路,要改為平坦,讓我們內心和行為；知識和信心,都互相調和.
\newline
\newline


\newpage

\section{第43屆~港九培靈會~周主培~第六講}
\label{sec:1414}
\textbf{福音本是神的大能}
\newline
\newline
連結: \href{https://www.hkbibleconference.org/session-message/view/1414}{\texttt{https://www.hkbibleconference.org/session-message/view/1414}}
\newline
\newline
\hyperref[sec:1413]{< < < PREV SERMON < < <}
~
\hyperref[sec:toc]{[返目錄]}
~
\hyperref[sec:1415]{> > > NEXT SERMON > > >}
\newline
\newline
羅馬書 1:16-1:16
\newline
\begin{longtable}{cl}
\hline
\hline
章節 & 經文 (和合本修訂版)\\
\hline
1:16 & \begin{tabularx}{0.7\textwidth}{X} 我不以福音為恥；這福音本是神的大能，要救一切相信的，先是猶太人，後是希臘人。 \end{tabularx} \\ \\
[1ex]
\hline
\hline
\end{longtable}
這一個時代是一個講能力的時代,一國工業的發達與否,與國家的力量是分不開的.古時用人力,後來用風力；繼之用水力,用電力,用原子能力,用核子能力.每個國家盡量擴展能力來源以求工業發達.第二次世界大戰後,香港逐漸成為工業城市.許多國家力求工業發展,因它給社會帶來繁榮.世界的工業的確日趨進步.但道德卻成了癱瘓的光景.人類可征服太空,但卻不能征服自己的情慾.人想盡辦法延長壽命,但卻無法改變生活.從前的人是滿了繭的雙手.現在的人滿了縐紋的額頰.在這世界,我們不能否認因工業發達造出各種機器,今日的機器靈巧像人一樣；但人卻慢慢地變成機器,人並無改變自己的生活.
\newline
\newline
海德公園是個任何人都可以講話的地方,演講的人可站在肥皂箱上發表,總會有人來聽的.某次,一位傳道人在那裡露天佈道,忽然有個急進的青年問傳道人說:「你講信耶穌可以得天上的平安快樂,死後可以得永生.現在站在我身邊的人,穿的是件破衣服,你能否給他件新衣服?」傳道人默默地感謝神,因神賜給他智慧作出如下的回答:「從他臉上我看出是個酗酒者,如果現在給了他新衣服,也不能擔保他永遠穿在身上.因他必將新衣服賣掉買酒喝.若他信耶穌,得到新的心；新的力量,生活自然也改變了.」
\newline
\newline
神的大能是改變人心的力量,給人有新的生命過新的生活,有新的力量走新的道路.今日人類一直關懷何謂左何謂右,而漠視何謂上何謂下.今日我們應當思想的,不僅是世上的事；我們應該相信神的力量,才能改變世界.曾有些學生問我,怎樣改造這世界?許多人對於這問題高談闊論.有的說:「要先改革政治.」但我們看見曾由君主改為民主,獨裁改為共和,還是不能達到改善的目的.政治的改革並不能改變人心,也不能改造世界.經濟系的學生說:「改造世界要從經濟著手,經濟發展,著人民衣食足；知榮辱,然後才能改變世界」.有許多罪惡,是有錢有勢力的人犯的.一個參議員駕駛車子,將自己的女秘書浸死.因他有錢有勢,結果若無其事,逍遙自在.在這個世界,偷一塊錢的是賊,偷百萬的倒是富翁.經濟不能改變人心.有人說:「教育普及,可以改變世界.」聖經告訴我們,「知識來,煩惱也來.」並不是說知識能代表煩惱,而知識不過是煩惱的說明書,知識不是悲劇而是悲劇的說明書.愚笨的人令人苦,聰明的人自己苦.知識越多,受教育越多,內心煩惱也越多.受過教育的人所犯的罪,是聰明的,令人不知不覺；沒有受過教育的,其中有些笨頭笨腦的,他們一犯罪就很容易被人揭穿.教育並不能改變人的內心,正如第二次大戰前,德國教育水準相當高；那時卻產生了一位名聞世界的獨裁者希特勒,和他的宣傳部長戈培爾,他是個具有博士學位的人.
\newline
\newline
如果要世界和平,每個人都應該放下武器；國家也不應製造重武器,世界和平才有盼望.我曾在日本住了三年,有一次我到早稻田大學講道.我一踏入他們的校園,就發現有許多標語,而且許多人圍在那裡看.我不懂日文,但標語上寫的是反對原爆,我會意是反對原子彈.我並不加以理會,一直行抵會場.經過介紹後,準備上台講道時,有個學生舉手說:「先生!我有問題要問「我們學校反對原子彈,每人都要簽名,你簽不簽?」這令我進退兩難.我唯有禱告,求主給我力量.然後我對他說:「在我未簽名之前,我必須知道為甚麼我要簽名.」他說:「這很簡單,為了反原子彈.」問我們:「為甚麼我們要反原子彈?」答:「就是反對戰爭.」我說:「我也反對戰爭.」但有一事我要問你:「在沒有原子彈以前,有沒有戰爭?」.他無言可答.我告訴他:「在沒有原子彈以前有飛機大砲,同樣有戰爭.」我繼續問他:「如果沒有槍砲是否沒有戰爭?」他默然不語.我說:「沒有槍砲也有刀劍.沒有刀劍也有竹竿木桿.沒有竹竿木桿,也有石頭,沒有石頭也有拳頭,還是可以打仗.」故此,我不簽名.他卻一言不發,無奈我何.我說:今天我來此所要講的,就是人的內心,怎樣才有平安?如困個人心內有平安,家才有平安,社會竹有平,國家才有平安,世界才有平安.除非接受那位和平之君,否則心裡永無平安.
\newline
\newline
世界怎能得到和平呢?人類怎能得到新的希望呢?政治,經濟,教育都不能解決世界和平問題.或有人說:「宗教可以改善世界.」但我告訴你,「不能」,甚至基督教也不能.第二次大戰後,日本興起了百餘種,大小不同的宗教.何謂宗教?就是人想尋求上帝,尋求到上帝那裡的途徑.這本來是好的,但是有錯.耶穌說:「我就是道路,真理,生命,若不藉著我,沒有人能到父那裡去.」(約十四6)
\newline
\newline
福音是什麼?羅馬書第一章告訴我們:一,福音是神所應許的(一節)二,福音是神的兒子(四節)三,福音是神的大能(十六節)四,福音是神的義(十七節)五,福音是神的事情(十九節)六,福音是神的永能和神性(廿九節)七,福音是神的榮耀(廿三節).
\newline
\newline
現在讓我對基督徒講幾句話.香港有四百萬人口,基督徒約有廿萬人,這心例實在太少了.香港有很多聚會的地方,但仍有許多人在黑暗中；原因是基督徒以福音為恥,不敢在人面前承認耶穌是救主.昨天午膳時,因為人多,只好和別人搭檯.我謝飯之後,見他們面面相觀.後來有人告訴我說:「他們也是基督徒.」基督徒若沒有清楚認識神的福音,便以福音為恥.保羅說他是神的使者,使者,換言之就是大使.如果一國的大使以自己為恥,那怎能辦外交呢,基督徒對於神的福音沒有清楚的了解,沒有真確的信心,便會以福音為恥.
\newline
\newline
福音不是空洞的理論,不是德的體系；福音是耶穌基督,本為神,後為人.道成肉身住在我們中間,充充滿滿有恩典,有真理.馬可一章一節說:「神的兒子,耶穌基督福音的起頭.」耶穌到世上來,不但傳福音,祂本身就是福音.有耶穌才有福音,耶穌到世上,將神的福音彰顯出來.由那裡彰顯神的能力?非常奇妙!神使兩根木頭成為十字架,顯出祂的能力.耶穌之生,死,復活彰顯了神的大能.感謝主!因著神的大能,主由童女馬利亞所生,被釘在十字架上,從死裡復活了.
\newline
\newline
我以為耶穌復活並不是神蹟.耶穌是神,如果一直在墳墓裡,那才是奇怪的事；我以為耶穌的死才是神蹟；永生神的兒子居然死了,為我們的罪甘心樂意死了.
\newline
\newline
聖經告訴我們,當人類犯罪時,神應許說:「女人的後裔要傷蛇的頭.」耶穌在十字架上已成就了這事,已經除滅了撒旦的作為,為我們開了一條又新又活的路.藉著耶穌基督的死,讓我們能夠親近神.
\newline
\newline
十字架成就了非常奇妙的事,(請原諒我說)是人所不能解決的事,也是神的難題.(可能我言之過分)神的慈愛和公義,在神心裡同等重要.神的公義,叫人和神隔絕了；神的慈愛,叫世上的罪人得救,因此,十字架成就了神的慈愛和公義.十字架是由一橫一直,兩根木頭做成的.十字架的交叉點,就是神慈愛他和公義的匯合.
\newline
\newline
聖經告訴我們,耶穌在十字架上流血受審,付出祂的贖價.請問,由何處贖我們出來,贖價付出給誰?是否付給撒旦?你想撒旦敢不敢接受?難道神與撒旦妥協?加拉太四章五節告訴我們,神將我們從律法之下贖出來.唯耶穌在十字架上流血,滿足了律法的要求,祂在十字架上說:「成了」,成了贖罪的救恩,使我們得著兒子的名分.叫一切在基督裡的人,不被定罪,成為新造的人,既然我們有這樣寶貴有能力的福音,為何我們不真正為主的福音齊心努力呢?
\newline
\newline
感謝主的恩典!神在香港興起了很多教會,特別對傳福音有負擔.中國人已往一直在接受的階段,很少派人到外面傳福音.我並不是說每個人都往外地福音,如果你不能在本地傳福音,難道能在外地傳福音嗎?這是不可能的,除非你像保羅說:「我不以福音為恥,因為這福音本是神的大能.」(羅一16)
\newline
\newline
我最近看一本福特先生所寫的書.他說,早在法國還在越南爭戰時,有個宣教士被共黨游擊隊擄去.因為他會講英語,青年軍官就跟他學習,二人就漸生感情.軍官對教士說:「幾個月來我跟你讀聖經,我發現其中真理,實在和共黨有很多相同之處；而且聖經中有很多,是我們所沒有的.我們認為共產主義,是世界上唯一重要的主義；所以我們寧願為黨而生,也甘願為黨而死；我覺得你們基督徒,獨惜沒有這樣的精神.」我看到這裡,心裡由衷地深受感動.
\newline
\newline
可惜基督徒沒有一顆,真正為主傳福音的心.路上行人絡繹,我不知道你們有否看到每個人裡面,都有個寶貴的靈魂?你是否以基督福音為恥?我們若沒有傳福音的靈,就不能像保羅「不以福音為恥.」不肯為福音而犧牲!如果香港每個基督徒,每年領一人歸主,不到廿五年,就可以領全港的人都信主了,但事情卻並不是如此樂觀的.基督徒自己認識主,也領人歸主；我相信香港教會必能得復興.如果每個基督徒真正不以福音為恥,我們就要看見神的能力彰顥.巴不得在坐每位都有這樣的負擔,有這樣的心志.對福音認識,對福音信任,知道福音就是耶穌基督,領人歸主,榮耀主名!
\newline
\newline


\newpage

\section{第43屆~港九培靈會~周主培~第七講}
\label{sec:1415}
\textbf{將心歸神}
\newline
\newline
連結: \href{https://www.hkbibleconference.org/session-message/view/1415}{\texttt{https://www.hkbibleconference.org/session-message/view/1415}}
\newline
\newline
\hyperref[sec:1414]{< < < PREV SERMON < < <}
~
\hyperref[sec:toc]{[返目錄]}
~
\hyperref[sec:1416]{> > > NEXT SERMON > > >}
\newline
\newline
箴言 23:26-23:26
\newline
\begin{longtable}{cl}
\hline
\hline
章節 & 經文 (和合本修訂版)\\
\hline
23:26 & \begin{tabularx}{0.7\textwidth}{X} 我兒啊，要將你的心歸我，你的眼目也要喜愛我的道路。 \end{tabularx} \\ \\
[1ex]
\hline
\hline
\end{longtable}
最近有些人問我:「你是怎樣信耶穌的?你的英文名字為甚麼叫摩西?」感謝主!這是我的好見證.
\newline
\newline
感謝神拯救了我的家庭.在我年幼的時候,我參加主日學,接受了主耶穌.之後,主日學老師教我應該為耶穌作見證.可是我很怕為耶穌作見證,老師給我約翰福音的單行本帶回家去.當我父親看見這書的時候,他非常生氣說:「一定是你把那洋教帶到家裡.」父親命我拿開,我就把書移到另一桌子.過了兩天我父親打開這本書來看,竟被裡面的話抓住了,他就把全書讀完.遂即問我「你的牧師在那裡?」那時我才八歲,我不敢帶父親到牧師那裡,因為怕他和牧師吵架,但我又不敢不帶他去.結果,我鼓起勇氣帶他去找牧師.他問牧師許多問題,以後每天都去和牧師談話.這樣!過了三個月,父親也接受耶穌,我們就受浸禮.我們信了耶穌之後,受著家族的逼迫.我們的大家庭非常保守,一向拜祖宗.為了我們信耶穌,族人開家庭會議,問我父親:「你們要耶穌呢?要祖先呢?」「若要耶穌,就必須離開家庭,無權接受祖傳的財產.如果要祖宗你就反教罷!」他們有數天的時間,給父親作考慮.感謝主!我父親說:「我寧願要耶穌.」於是雙親和我就被趕出族了.我們離了家,搬到別處.當晚,我們三人跪下禱告.我父親說:「上帝啊!如今我甚麼都沒有,僅有這個兒子,我願意將他奉獻給你!」.從那時開始,我的中文名字改了.本來是「祖」改成了「主」.在三個月之內,我父親把全部聖經看完一遍.他認為摩西把自己的百姓,從埃及為奴之地領出來.他對我說:「孩子啊!我要給你一個聖經名字叫摩西.」同時,他將箴言書廿三章26節的聖經節給我.家父已於今年二月離開世界到天父那裡了.我永不忘先父所給我的這節聖經── 「我兒!要將你的心歸我,你的眼目,也是喜悅我的道路.」以後我們就離開家鄉.我永遠不會忘記,我們跪在神的面前,父親把我奉獻給神.
\newline
\newline
馬太廿二章十四節,耶穌特別題到:「被召的人多,選上的人少.」由此,我們看到主耶穌心裡的願望.祂盼望我們不單作個被召的人,祂更渴望我們作個被選上的人,今天很多基督徒,願意遵行神的旨意,但很少的人能夠體會神的心意.
\newline
\newline
天上有聲音說:「這是我的愛子,我所喜悅的.」這裡給我們看到愛,和喜悅是兩件事.(請原諒我以自己的事作比喻)我有五個孩子,都是我所愛的.但其中有討我喜悅的；或有時討我喜悅,有時不討我喜悅的.巴不得我們基督徒,下個雄心大志,不僅蒙神所愛,也蒙神喜悅.
\newline
\newline
基督徒在今天,不肯努力進天國,不肯付代價在屬靈的事上.世界上的人,有不必努力,而承繼了巨量的財產；有人碰巧得了高的學位,有的男人並沒有極力追求而得賢妻.但是在屬靈的事上,如果不付上代價,就不能討主喜悅.
\newline
\newline
今晚我們不談得救的問題.我們要談怎樣蒙神喜悅?神說:「我兒,要將你的心歸我,你的眼目,也要喜悅我的道路.」所以,喜悅神的道路,是蒙神喜悅的一個妙法.
\newline
\newline
請大家一齊思想,為什麼耶穌說:「被召的人多,選上的人少.」耶穌在世上時,祂有時在露天,有時在海邊,有時在山上講道,常常聽眾擁擠.路加十二章一節告訴我們,有幾萬人聚集.被召的人多!當耶穌復活以後,最多的一次向五百人顯現.耶穌每次顯現時都對他們說:「你們要等候在耶路撒冷.」但等候在耶路撒冷的僅百廿人.這是耶穌親自囑咐的教會,參加禱告會的人,也不過佔四分一.至於今日的教會,更是望塵莫及了.耶穌在世上時,差遣門徒兩個兩個出去傳福音,一共差遣了七十人.耶穌自己揀選的門徒有十二個,其中有三個,無論主到何處都跟隨著.到了耶穌被釘在十字架上的時候,唯有主所愛的門徒約翰,站在十字架旁.被召的人多,選上的人少.真可憐!我們沒有讓主的心得著滿足.
\newline
\newline
保羅說:「如今常存的,有信,有望,有愛,其中最大的就是愛.」因為愛有永恆的價值,將來我們到了天上,我們所信的主已經看見了；我們所盼望的,已經得著了；但我們在永遠的愛上面,仍有信心,有盼望.親愛的弟兄姊妹阿!誠然!召的人多,選上的人少.
\newline
\newline
以色列百姓離開了又真又活的神,去敬拜偶像,離開了耶和華而拜巴力.他們輕看創造萬物的神而拜人手所造的；神的忿怒臨到他們,使他們的仇敵米甸人,像蝗虫一樣,來攻擊他們.當他們收割時,所有的都被奪走了.以色列人就哀求神,神聽他們的禱告；興起一人叫基甸的,來拯救他們.他不是壯士,而是很軟弱的人.神對他說「靠著我的能力,你可去拯救以色列人.「跟隨基甸的共有三萬二千人.神說:「人數太多,恐怕得了勝,以為是靠自己的力量,膽怯的可以叫他們回去.」神知道人的弱點,常是事未成就之前膽怯,事既成將功歸己.後來有二萬二千人回去,留下一萬人.神說人數還是太多,祂教基甸用喝水的方法試驗,結果有九千七百人跪下喝水,用手捧著餂水的有三百人.神對基甸說:「我要用這餂水的三百人拯救你們.」為甚麼那跪下喝水的人神不能用他們呢?有的說,跪下來喝水的表現他們不儆醒防敵.在我住處附近有幾個猶太人的會堂.我認識其中一位拉比,我常常有機會和他談到舊約聖經,我也把新約聖經介紹給他,我們彼此談得很投契.有次,我問他:「那餂水的三百人,神要用他們是否因他們儆醒?」他徵求我的看法.我說:「以上的看法我不大滿意.」一個真正的希伯來人是不隨便跪下的；唯有在萬軍之耶和華面前跪拜.而那九千七百人是慣拜偶像的,跪對於他們是不足輕重的事.所以那三百人才是真正敬畏神的人.」他說:「我非常同意你的看法.」
\newline
\newline
今天許多人為了主,而在險裡轉,為錢而無惡不作,且為錢不恥下跪.為錢基督徒變了質,失去在神面前的見證.
\newline
\newline
全本聖經,舊約有卅九卷,新約有廿七卷,共六十六卷書,除以斯帖和雅歌兩卷書之外,每卷都有題及神的名字.在以斯帖記,雖未看見神的名,卻看見神的手和祂的寶座,神仍然坐在寶座上.以斯帖記內容也是論到關於揀選的事.
\newline
\newline
有很多的童女被召入皇宮的女院,但被選為皇后的只有以斯帖.為何諸多人中唯有以斯帖被選上!當然是因為神坐在寶座上,神有祂的美意.但在人來講,以斯帖非同其他童女.聖經說她別無所求,凡看見她的都喜歡她.以斯帖進到皇宮,目的在於坐寶座戴皇后冠冕,她唯要榮耀的冠冕,除此她別無所求.
\newline
\newline
許多弟兄姊妹,年青時很愛主,跟隨主；逐漸愛心冷淡,轉向世界,跟隨瑪門,失去了原有最好的福份.
\newline
\newline
北美洲有個農夫告訴我,每逢冬天,他帶兩隻狗出去獵鹿,有時候獵到卻是隻兔子.在你人生道路上,有很多鹿,也有很多兔子,特別是青年人,起初追的是鹿,後來卻捨鹿而追子.青年人!你的學業,事業,婚姻也是這樣.今天有許多基督徒,在屬靈的道路上,沒有選上主最美最好的.
\newline
\newline
以色列人得罪了神,且遠離神,埋怨神,令神的心傷痛.神帶領他們埃及,過經海,經曠野,應許領他們進迦南,但以色列人懷不信的惡心,派了十二個人去窺探迦南地.結果有十個回來報告,那地的人又高又大,相形之下,自己如同蚱蜢.所以全會眾就埋怨神.於是神向以色列人發怒,使603548名二十歲以上的男丁倒斃在曠野,只有迦勒和約書亞二人說:「我們一定要去,因為神與我們同在.」神說:「因他另有一個心志,專一的跟從我」.今天教會裡,都是少數服從多數.請原諒我說:「多數人的意見未必對的.」耶穌在世時,多數人喊著說:「釘他十字架」.十字架的道路是極窄的,我們不能隨從群眾,在神的工作上,不是少數服從多數,乃是多數都當服從主.
\newline
\newline
今天基督徒與世界妥協,看風使舵,以為光與暗,善與惡,信與不信都可混在一起.很多時候,我們心懷二意,沒有專一的心跟從主.心懷二意的人,在他一所行的路上都沒有定見.親愛的弟兄姊妹!我們是神用重價買回來的,神盼望我們專一的跟從祂.
\newline
\newline
神說:「我兒,要將你的心歸我,你的眼目也要喜悅我的道路.」
\newline
\newline


\newpage

\section{第43屆~港九培靈會~周主培~第八講}
\label{sec:1416}
\textbf{「沒有異象民就滅亡」}
\newline
\newline
連結: \href{https://www.hkbibleconference.org/session-message/view/1416}{\texttt{https://www.hkbibleconference.org/session-message/view/1416}}
\newline
\newline
\hyperref[sec:1415]{< < < PREV SERMON < < <}
~
\hyperref[sec:toc]{[返目錄]}
~
\hyperref[sec:1417]{> > > NEXT SERMON > > >}
\newline
\newline
啟示錄 1:8-1:20
\newline
\begin{longtable}{cl}
\hline
\hline
章節 & 經文 (和合本修訂版)\\
\hline
1:8 & \begin{tabularx}{0.7\textwidth}{X} 主神說：「我是阿拉法，我是俄梅戛 ，是今在、昔在、以後永在的全能者。」 \end{tabularx} \\ \\ \relax
1:9 & \begin{tabularx}{0.7\textwidth}{X} 我—約翰就是你們的弟兄，在耶穌裡和你們一同在患難、國度、忍耐裡有份的，為神的道，並為給耶穌作的見證，曾在那名叫拔摩的海島上。 \end{tabularx} \\ \\ \relax
1:10 & \begin{tabularx}{0.7\textwidth}{X} 有一主日我被聖靈感動，聽見在我後面有大聲音如吹號， \end{tabularx} \\ \\ \relax
1:11 & \begin{tabularx}{0.7\textwidth}{X} 說：「把你所看見的寫在書上，寄給以弗所、士每拿、別迦摩、推雅推喇、撒狄、非拉鐵非、老底嘉那七個教會。」 \end{tabularx} \\ \\ \relax
1:12 & \begin{tabularx}{0.7\textwidth}{X} 我轉過身來要看看是誰的聲音在跟我說話。我一轉過來，看見了七個金燈臺； \end{tabularx} \\ \\ \relax
1:13 & \begin{tabularx}{0.7\textwidth}{X} 在燈臺中間有一位好像人子的，身穿垂到腳的長袍，胸間束著金帶。 \end{tabularx} \\ \\ \relax
1:14 & \begin{tabularx}{0.7\textwidth}{X} 他的頭與髮皆白，如白羊毛，如雪；他的眼睛好像火焰， \end{tabularx} \\ \\ \relax
1:15 & \begin{tabularx}{0.7\textwidth}{X} 雙腳好像在爐中鍛鍊得發亮的銅，聲音好像眾水的聲音。 \end{tabularx} \\ \\ \relax
1:16 & \begin{tabularx}{0.7\textwidth}{X} 他右手拿著七顆星，從他口中吐出一把兩刃的利劍，面貌好像烈日放光。 \end{tabularx} \\ \\ \relax
1:17 & \begin{tabularx}{0.7\textwidth}{X} 我看見了他，就仆倒在他腳前，像死人一樣。他用右手按著我說：「不要怕。我是首先的，是末後的， \end{tabularx} \\ \\ \relax
1:18 & \begin{tabularx}{0.7\textwidth}{X} 又是永活的。我曾死過，看哪，我是活著的，直到永永遠遠；並且我拿著死亡和陰間的鑰匙。 \end{tabularx} \\ \\ \relax
1:19 & \begin{tabularx}{0.7\textwidth}{X} 所以，你要把所看見的事、現在的事和以後將發生的事，都寫下來。 \end{tabularx} \\ \\ \relax
1:20 & \begin{tabularx}{0.7\textwidth}{X} 至於你所看見、在我右手中的七顆星和那七個金燈臺的奧祕就是：七顆星是七個教會的使者，七個燈臺是七個教會。」 \end{tabularx} \\ \\
[1ex]
\hline
\hline
\end{longtable}
「當主日我被聖靈感動.......」(啟一10)在今天主日晚上,巴不得主自己的靈感動我們的心,叫我們有個被提的靈到祂那裡,看見神在我們身上的要求和計劃.
\newline
\newline
當約翰孤單地在拔摩海島上,那時教會被逼迫,使徒四散.這正是我們現今的境況,我們不知將來的遭遇如何.約翰在孤島上,好像四顧無望,在這情形下,他看見了異象.
\newline
\newline
今日教會需要神給我們看見異象.我們常常聽而不聞,視而不見.並非神沒有異象向我們顯現.因為我們的眼目昏迷,所看見的是世界上的一切,正如彼得在風浪中快要沉下.請看今天的光景:惡勢力日趨澎漲,人心充滿苦悶；甚至基督徒亦覺前途茫茫,我們怨天尤人,凌空下淚!
\newline
\newline
世界前途究竟如何?有誰知道!在這樣光景中,我們真需要神給我們看見異象.所謂異象,就是見別人所未見的.我常對青年弟兄姊妹說:一個人的看見能決定他的揀選,而他的揀選將決定他的一生.如果沒有看見就沒有決定；沒有決定,就沒有方向,沒有方向,就沒有前途.我們需要看見──一個更新的看見.當我們看見樹上掉下來一個果子時,我們就把它拾起來吃,這是普通的反應.但牛頓看見樹上的果子掉下來,他就發現其中有一種力量；結果,他發明地心吸力的原理.當我們看見水煮開了就沖茶,但瓦特看見水開了,發現其中有蒸氣的力量；結果,他發明了蒸氣機.
\newline
\newline
在這動亂不安的時代裡,求主給我們看見異象.約翰在拔摩海島上看見異象,叫他下垂的手舉起來；發酸的腿,再有力量奔跑前程.約翰聽見天上有聲音說「我是阿拉法,我是俄梅戛,我是首先的,我是末後的,是昔在今在永在的全能者」.
\newline
\newline
約翰所看見的異象:──
\newline
\newline
(一)神仍然坐在寶座上.　　詩廿九篇說「洪水泛濫之時,耶和華坐著為王......直到永遠.」烏西雅坐在王的寶座上做王五十多年,當他崩的那一年,以賽亞說:「我看見主高高的坐在寶座上.」雖然世界會改變,人會改變,但我們的神永遠不改變.從摩西的禱告(詩九十),我們看見人的一生,是何等的短促,但是神從永遠一直到永遠的.歷史兩字,英文 History,照字譯中文是「他的故事」.神在歷史裡,神掌管歷史.
\newline
\newline
有時,我們抬頭問蒼天,為什麼神不看顧?為什麼神任憑世界到這樣的光景?弟兄姊妹!人所看見的不過僅是「一點」,而神所看的乃是「全面.」我有個特別的經驗.有次正是雨過天晴之際,我在天空中飛機裡,看見整個大圓圈的虹,是我有生以來所未曾見過的.平時我們在地面上所看見的虹,只有半個圈.由此可知,神所看的是全面的,人所看的不過是少部份.我還有個經驗.有一天,我在朋友辦公室的天台上(很高大廈的頂上)看遊行,居高臨下,看見全部的遊行隊伍.平時不上高處,只看見一隊過了又一隊,沒有法子可以看到整個的.我們在這世界所看見的也是如此,一朝又一朝,一代又一代；一個又一個過去了,好像遊行一樣.我們能看見的,只不過一部份.唯獨我們的神,祂是阿拉法,祂是俄梅戛,從頭到底,從首到尾,祂一目瞭然.如果我們知道神仍然坐在寶座上,當我們得意時,我們不會驕傲；失意時,我們也不垂頭喪氣了.
\newline
\newline
我從1949-1956年在印尼工作,開始時印尼剛獨立.因為各政黨勢力逐漸強大,教會漸受逼迫；困難,日見增加,傳福音的工作大受限制.到了一九六五年九月我們感到將有事情要發生,但不知何事.我們只有迫切向神禱告,感謝主!時局才轉危為安,,因此給全印尼教會帶來一次大復興.
\newline
\newline
(二)主耶穌站在教會當中.　　聖經說:「有一位好像人子,站在七燈台中間.」雖然當時的教會有很多軟弱和失敗,但主不離開他的教會.那時教會的環境非常黑暗,撒旦要毀滅神的工作；許多基督徒被迫害,四處失敗.除了士每拿和非拉鐵非以外,以弗所教會,失去起初的愛心.別迦摩教會,有撒旦的作為.推雅推喇的教會,有自稱為先知的婦人.撒狄的教會,按名是活的,其實是死的.老底嘉教會,不冷不熱.
\newline
\newline
許多弟兄姊妹,特別是青年人；一看見教會冷淡?失敗,就想離開自己的教會.曾有個青年,對一位神的僕人說:「我發現我的教會,非常冷淡倒退；我很想參加另外的教會,你能否介紹一個較好的教會?」那位神僕說:「你千萬別隨意離開你的教會,除非這教會的信仰完全崩潰了.如果你找一個完全的教會,恐怕因為你加入了那教會,而變成不完全了.」在地上我們找不到一個完全的教會,我們不能單單消極的離開教會,我們應為自己的教會禱告.不久以前,美國有一間衛理公會的大學,校方充滿問題,教員中爭權奪利.有四個青年人跪在神面前流淚禱告,求主復興他們的大學.一天早禱會,聖靈在他們中間作工,兩晝夜聚會沒有間斷；無人想要離開聚會,復興的火就燃起來了.不但這所大學得復興,其他的大學也得到復興.
\newline
\newline
教會雖然失敗,軟弱,沒有讓主作主作王,主仍然沒有離開祂的教會.因為主是教會的元首,甚至老底嘉教會,主也沒有離棄,主仍站在門外扣門.每一個教會中,都有得勝的人,主要呼召那得勝者.雖然我們沒有讓主作王,我們失敗,虧欠,我們把主關在門外；只看見人作主,人掌權,沒有主的地位,但主仍沒有放棄祂的教會,主仍在教會當中.
\newline
\newline
(三)約翰看見自己是神的奴僕.　　耶穌基督的真理,不是空洞的理論,不是僅外表上的承認；乃叫我們的內心,有親切的體會.「僕人」按原文英譯「奴隸」(slave).起初我對於奴隸的身份沒有特別的感覺.某次,我被赴黑人弟兄姊妹的聚會.黑人善唱,他們的歌聲真是感動我心.他們聚會,情形和我們完全不同.我每講一句,他們必有回音的反應.當我講到 slave,全會場忽靜默無聲,我繼續地講；有的人在點頭,老年者在流淚.我才明白他們的祖宗做過奴隸,他們知道甚麼是奴隸.從那時我真正學到甚麼,叫做承認主是我們的主.我們接受主耶穌做我們的主,可是沒有真正將祂當作主.無神論者說沒有神.聖經說,惡人的心說沒有神,愚頑人的心說沒有神.其實基督徒也可說,是無神論者,因為常常隨己意行事,沒有尊主為大.我們稱為基督的精兵,但我們常常不服從主命令.有一次彼得正禱告時,看見一個異象,有東西從天下來,有聲音說:「起來宰了吃.」彼得說:「主阿!」這是不可的,凡俗物和不潔淨的物,我從來沒吃過.」我覺得奇怪,稱祂為主,卻說「不是的,不...... 不......」所說的都是不,豈不矛盾嗎?今天有很多基督徒,稱祂為主,卻不遵行祂的旨意.稱祂是光卻不跟隨祂.稱祂是道路.卻不行在其上.稱祂是牧人,卻不聽祂的聲音,我們實在虧欠主!
\newline
\newline
最後,我要做個見證:當家父將我奉獻給神以後,他一直提醒我說:「你是屬於神的,你要為神工作.」但是我不肯,家母也不肯,因為我是獨生的兒子,是三房唯一的男丁.家母說:「別的事情都可以做,千萬可別做傳道；信主可以,愛主可以,但不可做傳道.」
\newline
\newline
當我十二歲那年,我患了傷寒病,死去了三小時,母親焦急,向神許願:「若孩子活過來,就讓他做傳道.」事後,母親一直對我說:「孩子啊!我們已將你奉獻給神了.」到了我年紀漸長,我看見教會的黑幕及一切情形.我說:「我不但不願做傳道,連耶穌我都不願相信了.主的愛在我身上,主管教我.抗戰時期,我患惡性虐疾,又死過去了.那次我清楚知道,離開我睡的床,背起包袱一直孤單地走；那是下雨的晚上,我覺得那條路很難走,走到河邊,等候要過擺渡,有個人對我說:「今晚沒船,你回去吧!」我就回來,神又把我救回來.母親對我說:「若有人不愛主,這人可咒可詛,你要服事祂.」感謝主!主步步帶領我.一九四九年,神差遣我到印尼去傳福音.我將這事告訴雙親,父親當然無問題,母親卻問我印尼在何處?我就向她說明.她說:「孩子啊!你忘記你是我獨生的兒子,在我未死以前,你絕不能離開我.我愛主,但是我也愛你!甚至愛你勝過我的性命.我不願你離開我,願你留在中國,在你講道時為你禱告.孩子,你再三考慮罷!」後來,我們一同禱告,過了數日,母親對我說:孩子啊!現在我學了一個功課,禱告改變了我的心；叫我順服神,不要照我的意思成就,要讓神的旨意成就.你去吧!我的禱告與你同在.你若親近神,你就等於跟我在一起；你若離開神,也等於離開了我」.
\newline
\newline
我們常常盼望禱告改變事物,事情順利,困難沒有.想不到禱告以後,黑雲依然滿佈,道路仍然崎嶇；但神給我們力量,叫我們有能力勝過患難.基督徒所信的神,不是草草讓祂改變環境來適合我們；乃是改變我們,來得勝環境.我學習了這個功課,將主當作主,祂要我們做甚麼呢?是要我們完完全全順服祂!遵祂而行.
\newline
\newline
我和母親分別了廿多年,我相信有一天能看見她；如果不是在這地上,就是在榮耀裡看見她；要謝謝她!因為她教我學會了功課,就是禱告改變我們的心,順服我們的主；遵主而行,要將主當作主.
\newline
\newline
親愛的弟兄姊妹!你是否看見這個異象──神仍然坐在寶座上,耶穌是教會的元首,我們是祂的奴隸.我們因愛心做祂的奴隸,因為祂用重價把我們贖回來,我們在身子上,應當榮耀祂.
\newline
\newline


\newpage

\section{第43屆~港九培靈會~周主培~第九講}
\label{sec:1417}
\textbf{奇妙的愛}
\newline
\newline
連結: \href{https://www.hkbibleconference.org/session-message/view/1417}{\texttt{https://www.hkbibleconference.org/session-message/view/1417}}
\newline
\newline
\hyperref[sec:1416]{< < < PREV SERMON < < <}
~
\hyperref[sec:toc]{[返目錄]}
~
\hyperref[sec:1418]{> > > NEXT SERMON > > >}
\newline
\newline
約翰福音 21:15-21:19
\newline
\begin{longtable}{cl}
\hline
\hline
章節 & 經文 (和合本修訂版)\\
\hline
21:15 & \begin{tabularx}{0.7\textwidth}{X} 他們吃完了早飯，耶穌對西門‧彼得說：「約翰的兒子西門，你愛我比這些更深嗎？」彼得對他說：「主啊，是的，你知道我愛你。」耶穌說：「你餵養我的小羊。」 \end{tabularx} \\ \\ \relax
21:16 & \begin{tabularx}{0.7\textwidth}{X} 耶穌第二次又對他說：「約翰的兒子西門，你愛我嗎？」彼得對他說：「主啊，是的，你知道我愛你。」耶穌說：「你牧養我的羊。」 \end{tabularx} \\ \\ \relax
21:17 & \begin{tabularx}{0.7\textwidth}{X} 耶穌第三次對他說：「約翰的兒子西門，你愛我嗎？」彼得因為耶穌第三次對他說「你愛我嗎」，就憂愁，對耶穌說：「主啊，你無所不知，你知道我愛你。」耶穌說：「你餵養我的羊。 \end{tabularx} \\ \\ \relax
21:18 & \begin{tabularx}{0.7\textwidth}{X} 我實實在在地告訴你，你年輕的時候，自己束上帶子，隨意往來；但年老的時候，你要伸出手來，別人要把你束上，帶你到不願意去的地方。」 \end{tabularx} \\ \\ \relax
21:19 & \begin{tabularx}{0.7\textwidth}{X} 耶穌說這話，是指彼得會怎樣死來榮耀神。說了這話，耶穌對他說：「你跟從我吧！」 \end{tabularx} \\ \\
[1ex]
\hline
\hline
\end{longtable}
耶穌在世界上時,常用問題引起人的注意.(約一38)記載；當兩個門徒跟隨耶穌時,第一個問題耶穌問:「你們要甚麼?」(或謂你們尋找甚麼?)耶穌知道人心裡有一種需要,所以祂到世上來,是要叫人心得滿足.感謝主!那兩個門徒得到滿足了,因為他們已經找到了耶穌基督.
\newline
\newline
不過,很奇妙的!在約翰福音最後一章裡面所記載的,不是人在尋找,而是主耶穌在尋找.主的愛實在太奇妙!本來我們的罪多過我們的頭髮.但耶穌的愛,神的愛,比罪還多.我覺得最奇妙的,是神要我們愛祂.神是創造萬有的主,祂甚麼都不缺少；但祂一直在等待,盼望,祂要求人愛祂.(請原諒我說)神有需要,要我們愛祂,這是件太奇妙的事.
\newline
\newline
今天晚上,我們願意看見我們的主──死而復活的主,為我捨身流血的主,祂孤孤單單地；在那一天清晨,祂站在海邊看見門徒正在打魚,祂心裡有個需要.當他們打魚回來,耶穌就問彼得:「約翰的兒子西門,你愛我比這些更深麼?」親愛的弟兄姊妹!祂配得我們的愛,因祂為我們捨身流血,祂為我們離開天上的寶座,祂用重價把我們買回來.
\newline
\newline
一個將軍可以有千萬的俘擄,但妻子只有一個.俘擄是用力量得來的,而妻子是用愛得來的.主耶穌可以勉強我們愛祂,但祂從來不勉強我們.祂柔聲和氣的問彼得:「你愛我比這些更深麼?」耶穌所問的那門徒,竟是三次不承認耶穌的,是曾經跌倒的彼得.他好像將殘的燈火,他的心裡再也沒有力量了.雖然耶穌復活以後也向彼得顯現；但在彼得心靈裡,好像沒有勇氣,沒有信,是否能得主喜悅.耶穌從未責備過彼得:「為何三次不承認我?」耶穌問彼得:「你愛我嗎?」這句話的裡面有赦免,已往的事耶穌不再記念.耶穌的話包含恩典,信任,耶穌問彼得:「你愛我比這些更深嗎?」
\newline
\newline
保羅告訴提摩太,在末世時代人所愛的有三件事:一,專顧自己(就是專愛自己).二,愛錢財.三,愛宴樂.(提後三1-5)
\newline
\newline
現今的世界,人已到了極端自私自利的光景,人都存在「個人主義」,所思的是「我」,所想的是「錢財」,在這生活緊張的社會也是難怪.但基督徒整天所思想的也是錢財,正像猶大一樣.耶穌告訴他「天」的事,猶大想「錢」的事.結果他用卅兩銀子將主賣了.今日人所愛的是宴樂,我看見菜館裡面都滿了人,戲院門口都排長龍；至於其他奢華宴樂的地方怎樣,我卻不得而知.有人說,因為現在社會,人都覺得心裡虛空,前途渺茫,所以揮金如土,力求心內的滿足.
\newline
\newline
今日有很多基督徒也是為自己,禱告的時候,求神賜福給我,求神保佑我,保佑我的家庭,兒女,生意,都是為「我」.我們讀經,查經,是為要多得靈糧,這是不錯的；但多少時候還是為自己.有時候我們來聽道,不是為在神話語裡尋求真理,乃是望真理能夠符合我.所聽的道如果和自己所想的一樣,就說講者講得不錯.如果有一天真理要改變你,真理要叫你俯伏在它之前又如何[呢?連耶穌在世時,有些人也說:「祂所講的話,誰能聽得進去呢?」
\newline
\newline
耶穌沒有問彼得別的問題,只問:「你愛我嗎?」你可以信任一個人,但未必愛他.可是你不能愛一個人而不信任他.比方把錢存在銀行,不一定是愛銀行或經理,只能說是信任這銀行.如果你愛一個人,甚麼都願意給,甚至連自己的命都願意捨.可以信而不愛,卻不能愛而不信.我們讀聖經,不是單盼望得著亮光,得著聖經的知識；最寶貴的,是得著神的話語.先知耶利米說:「我得著你的話語,就當食物吃了!你的話語是我心裡的歡喜快樂,因為我是成為你名下的人.」
\newline
\newline
兩年前我到台灣去,有機會和大學的同學在一起.有一次,我聽見下面有聲音喊,某某人呀,有你的信.」接著,上面那人就問:「是鋼筆寫的還是毛筆寫的?」我感到很奇怪!後來他們告訴我,如果毛筆寫的當然是「老子」來信,除非有特別的事成匯款.不然,不必急於取信.如果鋼筆寫的,當然是愛人的信,便急不及待了!弟兄姊妹!你將神的話當作鋼筆寫的,還是毛筆寫的呢?
\newline
\newline
感謝主的恩典!最近幾年,在美國和加拿大興起了許多查經班.每逄查考聖經,我們就想到一件事；不是我們查聖經,乃是讓聖經神的話來查考我們的心.為何有時我們的禱告會人很少,禱告的生活那麼平凡?這是和主有關係的.有一次我見到一對青年人,他們都曾參加數月前所舉行的夏冬會,這次男的告訴我,他們已結婚了,他說,再不結婚要破產了.原來他們在冬令會結識後,常常藉長途電話通話；一個在紐約,一個在加尼福尼亞,每月要花去三百多美元的電話費.和愛人談話,感覺時間過得特別快!如果你真是愛主,你禱告在祂面前,向祂傾心吐意,話語會綿綿不絕,覺得主是何等的甜美!正如一首詩歌說:「在花園裡,雖然黑暗已籠罩滿天,我仍捨不得離祂慈面,祂與我同行,又和我談心,並對我說,我是屬於祂.我們同在時,情境真美,沒有別人能領會.」親愛的弟兄姊妹!你有否這樣的心向著主?
\newline
\newline
聖經說:「人若愛我,就必遵守我的命令.」如果你真正愛主,你遵守主的命令並不難.比方為你所愛的人辦事,一點不怕辛苦,因為要完全討他喜悅.聖經記載一個人,雖然他的生活行為,有些我不太敢領教；但他的愛,真是令人五體投地,他的愛實在偉大!他名叫雅各,為了愛他的表妹拉結,情願服侍他舅父,含辛茹苦,在所不惜!他舅父要他服侍七年,然後才把女兒給他.偶有一次,雅各思家心切,又覺生活難苦,決意離去.但當他見到拉結,意念都打消了,辛苦也都忘記了!好容易七年如七天的過去了.他滿以為可以和所愛的人在一起,想不到結婚之夜,他發現是拉結的姊姊利亞.雅各很生氣,和他舅父交涉.他舅父卻說:「這是我們的規矩,大家姐應先出閣,如果要拉結,可以給你,你再服侍我七年.」為了愛,雅各忍受了十四年的等待.神也願意愛祂的人服事祂,如果我們服事主,是為著將來得賞賜,為了將來要戴冠冕,這樣的動機還不夠高超呢!有許多青年人到國外留學,其中有很多的太太,日夜勤勞工作,以幫補她丈夫的費用.所以當她丈夫畢業得了 Ph.D. 的學位,她也得了一個 P.H.T. 的學位(push husband through)按字意譯是「把丈夫推過去.」這實在很難得,我看見許多姊妹,辛辛苦苦為了要她丈夫成名,無一分薪水,無一點獎章,為的甚麼?為了愛.如果我們真正愛主,我們甚麼都願意給祂.愛,永遠都是給的.神愛我們,甚至將祂的獨生子賜「給」我們.如果你真正愛主,一定願意為主作見證；晝夜所思想的是主,所講的也是主.
\newline
\newline
有的人不犯罪是因怕神刑罰,這樣的動機不算高超.我們不犯罪,乃為免得勝靈擔憂.聖經說:「聖靈好像母親.」一般而論,母親較慈愛,父親屬嚴厲,我自己也有此經驗.感謝神!我有很好的父母,特別是母親.我小時很頑皮,記得母親只打過我一次,打了以後她哭了,自從那次再沒有打過我,並非我變好了.每次我做壞事,她替我瞞住,但她暗暗嘆息；我不悔改,她心裡很難過.一直等到她跪在神面前痛哭禱告的時候,我也跪在她旁邊求神饒恕我,聖靈用說不出來的嘆息,為我們禱告.舊約聖經說「聖靈好像鴿子哀鳴!」為甚麼聖靈嘆息?為甚麼我母親嘆息?她眼看我不對,又不願告訴我父親,她心裡很多苦衷.我們不犯罪的原因是免得主傷心,不是為著怕主鞭打我們!
\newline
\newline
耶穌問彼得:「你愛我嗎?」許多解經家告訴我們,希臘原文耶穌所問的愛,和彼得所說的愛不一樣的.現在我們暫不題希臘文的意思.我有個小小的見證:有次我和幾個外國同學談及愛的問題.他們問我中國人的愛是何意思?感謝神給我想到一個字,我說,愛在中國是 Love, L loyalty 忠誠,O obedience 順服,V vestal 貞潔,E eternity 永恆.中國人說:「臣忠於君.子順服父.夫妻之間貞潔,朋友就是永恆.」聖經是非常奇妙的,耶穌基督總是把人懂得的事對人講,因為天上的事,是人所不能了解的.聖經說:「耶和華坐著為王,直到永永遠遠.」又說:「你若至死忠心,我就把生命的冠冕賜給你.」多少時候,我們對神不忠心,我們一面愛神一面卻愛世界.聖經說:「若有人愛世界,愛父的心就不在他裡面了.」又說:「祂要做我們的父親,我們做祂的兒女.」子女對父親要順服.中國古時,父親對兒女有絕對的權柄.聖經說:「我們做兒女的在主裡面聽從父母.」聖經的話沒有改變,雖然現今世界道德已經墮落,夫妻之間已經失去貞操.保羅對哥林多教會說:「我對你們起了憤恨,好像上帝的憤恨一樣.因為我已經把你們許配了基督；像貞潔的童女一樣,可是你們的心偏向世界.」我們向著主應有純潔的心.耶穌說:「我不稱你們為僕人,我稱你們為朋友；因為我從父那裡所聽見的已經都告訴了你們.」(約十五15)聖經裡的愛,是耶穌基督問彼得說,你是否用神聖,忠貞,永遠不變的愛愛我?彼得說,我是用感情上的愛愛你.感情是浮淺的東西,容易改變.人的愛,有範圍,有條件,有時間性；可是神的愛是永遠不變的愛,神是要我們用這樣的愛愛祂.
\newline
\newline
最後,耶穌並沒有問說,約翰的兒子西門你愛不愛我的羊?耶穌說「你愛我嗎?你牧養我的羊.」特別在我們中間有許多做神工作的人,好多時候,我們的心受了委屈,我們會灰心喪膽.因為人不可愛,你給他吃一百個饅頭,他會忘記；你給他一個拳頭,他至死也記得.你對他千萬次好,他會忘記；你對他一次不好,他永不忘記.各位在工場上服事神的僕人,我恭恭敬敬地說:「主知道你的勞苦.」我們不是為著人的緣故服事主,乃是為主的緣故服事人.」耶穌並沒有說:「你愛不愛我的羊?」耶穌說:「你愛不愛我!」為了主的緣故,我們仰望那信心創始成終的主,因祂擺在前面的榮耀喜樂,就輕看羞辱.羞辱是從人來的,喜樂是從神來的.因著神給我們的喜樂我們便輕看人的羞辱,奔那擺在前面的路程.
\newline
\newline
今天晚上,主耶穌問我們說:「你愛我嗎?」雖然我們有虧欠,雖然我們有失敗,甚至像彼得一樣.猶大用三十塊錢把主賣了.彼得用三句話把他的老師出賣了.同是一樣的賣主,但耶穌赦免彼得.耶穌告訴彼得:「你愛我嗎?你牧養我的羊.」弟兄姊妹!無論你是主日學的老師,無論你在青年團契,無論你在婦女會,無論你在那一部份的工作上,你是否為主的緣故牧養主的羊.我們要看主自己,我們要忠心為主；在愛裡有信心,有忠心,有順服；在愛裡,我們當堅貞的向著主.
\newline
\newline


\newpage

\section{第43屆~港九培靈會~周主培~第十講}
\label{sec:1418}
\textbf{靠主喜樂}
\newline
\newline
連結: \href{https://www.hkbibleconference.org/session-message/view/1418}{\texttt{https://www.hkbibleconference.org/session-message/view/1418}}
\newline
\newline
\hyperref[sec:1417]{< < < PREV SERMON < < <}
~
\hyperref[sec:toc]{[返目錄]}
~
\hyperref[sec:1382]{> > > NEXT SERMON > > >}
\newline
\newline
詩篇 16:1-16:11
\newline
\begin{longtable}{cl}
\hline
\hline
章節 & 經文 (和合本修訂版)\\
\hline
16:1 & \begin{tabularx}{0.7\textwidth}{X} 神啊，求你保佑我，因為我投靠你。 \end{tabularx} \\ \\ \relax
16:2 & \begin{tabularx}{0.7\textwidth}{X} 我曾對耶和華說：「你是我的主，我的福氣惟獨從你而來。」 \end{tabularx} \\ \\ \relax
16:3 & \begin{tabularx}{0.7\textwidth}{X} 論到世上的聖民，他們是尊貴的人，是我最喜悅的。 \end{tabularx} \\ \\ \relax
16:4 & \begin{tabularx}{0.7\textwidth}{X} 追逐別神的，他們的愁苦必增加；他們所澆奠的血我不獻上，我嘴唇也不提別神的名號。 \end{tabularx} \\ \\ \relax
16:5 & \begin{tabularx}{0.7\textwidth}{X} 耶和華是我的產業，是我杯中的福分；我所得的，你為我持守。 \end{tabularx} \\ \\ \relax
16:6 & \begin{tabularx}{0.7\textwidth}{X} 用繩量給我的地界，坐落在佳美之處；我的產業實在美好。 \end{tabularx} \\ \\ \relax
16:7 & \begin{tabularx}{0.7\textwidth}{X} 我要稱頌那指引我的耶和華，在夜間我的心腸也指教我。 \end{tabularx} \\ \\ \relax
16:8 & \begin{tabularx}{0.7\textwidth}{X} 我讓耶和華常在我面前，因他在我右邊，我就不致動搖。 \end{tabularx} \\ \\ \relax
16:9 & \begin{tabularx}{0.7\textwidth}{X} 因此，我的心歡喜，我的靈快樂；我的肉身也要安然居住。 \end{tabularx} \\ \\ \relax
16:10 & \begin{tabularx}{0.7\textwidth}{X} 因為你必不將我的靈魂撇在陰間，也不讓你的聖者見地府。 \end{tabularx} \\ \\ \relax
16:11 & \begin{tabularx}{0.7\textwidth}{X} 你必將生命的道路指示我。在你面前有滿足的喜樂，在你右手中有永遠的福樂。 \end{tabularx} \\ \\
[1ex]
\hline
\hline
\end{longtable}
聖經是神所默示的,我相信從創世記至啟示錄,都是神自己的默示.按我個人的心得,神的默示或啟示有兩種:直接的啟示──在舊約裡常記載:「萬軍之耶和華如此說.」新約裡耶穌說:「我實實在在的告訴你們.」間接的啟示──神把祂的啟示放在人心,而後由人將神的心意述說出來.特別在詩篇裡有很多神間接的啟示,神的心意,藉著祂的兒女向我們顯明.「我的心哪,你曾對耶和華說,你是我的主!我的好處不在你以外.......我將耶和華常擺在我面前,因祂在我右邊,我便不至搖動.因此我的心歡喜,我的靈快樂.我的肉體也要安然居住.」......「我的心哪!你為何憂悶?為何在我裡面煩燥?應當仰望耶和華.」......我的心哪!你要稱頌耶和華,不可忘記祂的一切恩惠.」......「耶和華我的力量阿!我愛你.」......「誰能登耶和華的山?誰能居住你的帳棚?就是行為正直,作事公義,心裡說誠實話的人.」這些話並沒有「耶和華如此如此說」,乃是人由心中發出說:「耶和華是我的牧者,我必不至缺乏.」這班人就是已經認識神的人.神把啟示感動他,藉此顯明神的心意.
\newline
\newline
今天晚上我們要思想大衛對神的認識.
\newline
\newline
神願人蒙福,願人都走上蒙福的道路.無可否認,今日許多基督徒,心中充滿了愁苦.雖然主耶穌曾說:「在我裡面有平安,在世上有苦難.但你們可以放心,我已經勝過了世界.」(約十六33)為何基督徒沒有平安?雖然我們聽見主耶穌說「你們心裡不要憂愁,你們信神也當信我.在父家裡,有許多住處,若是沒有,我就早已告訴你們了.」為何今天基督徒心裡還是充滿了愁苦?我們豈非勸人信耶穌,說:「凡勞苦擔重擔的,可到主這裡來,祂必使你們得享安息.」為何我們還時常唉聲嘆氣?這裡第四節說:「以別神代替耶和華的,他們的愁苦必加增.」原文:「以別代替的.......」親愛的弟兄姊妹!我恐懼戰兢的,思想這節聖經,在神面前省察自己.巴不得勝靈的光,光照我們每一個人.有許多的愁苦是我們冤枉受的.神不願意我們受苦,而是我們自己出了毛病,人任意妄為,自高自大,偏行己路,自作聰明,以別的代替神.神說:「因為我的百姓,作了兩件惡事,就是離棄我這活水的泉源,為自己鑿出池子,是裂不能存水的池子」(耶二13)今天人用破裂不能存水的池子,來代替活水的泉源.彼得前書告訴我們:有的人是為基督的緣故受苦,有很多時候,人卻自討苦受.今天的人,常把前後顛倒,上下反覆.聖經教我們要思想天上的事,但我們所思想的,多是地上的事；我們常常把上好的福分失去了.聖經教我們:要先求神的國和祂的義,一切東西都要加給我們；要先與弟兄和好,然後才回來獻祭；要先除掉自己眼中的樑木,然後除他人眼中的刺.很多事情神所規定的,有祂的目的,有祂的原因.眼睛,耳朵,各有其用；以耳代眼,他的愁苦必加增.今晚是我們最後一次聚會.我心裡還有負擔.在坐中或有些朋友們,尚未接受耶穌作你的救主,原因就是你用別的來代替.聖經記載有些人,就是因為以別的代替,而致遭遇許多的痛苦:
\newline
\newline
1. 該隱　　當人類犯罪,失去了神的榮耀.神用羊羔的皮,遮蓋他們赤身露體的羞恥.按律法,若不流血罪就不得赦免.亞當兩個兒子,亞伯取羊群中頭生的獻祭；該隱以自己地裡的出產獻祭.結果神悅納了亞伯所獻的祭.今日世界上有許多的宗教,如果把基督耶穌的道,也列為宗教之一(其實不是宗教)那麼,世界上僅有兩種宗教.其一是人尋求神──用盡自己的力量,辦法,計劃和自己的宗教去尋求神.其二是神尋找人──基督耶穌降世為要拯救罪人,祂在十字架上已經為我們完成了救恩,這救恩只要我們接受便可以得救.
\newline
\newline
2. 亞伯拉罕　　當他七十五歲的時候,神叫他離開他的本鄉.並應許使他成為大國,賜福他的後裔,給他的子孫多如天上的星,海邊的沙.他等候了十年,仍舊膝下無子.所增加的不過是年歲而已.他實在急不及待了,當時他的妻子撒拉,自嘆身已老邁,生育無望.提議以使女夏甲代替她.在亞伯拉罕八十六歲那一年,生了以實瑪利,從此以後,他的心充滿的痛苦,十三年之久,神和他斷絕的交通；神沒有句話對他講,也沒有向他顯現.這十三年悠長的歲月,在亞伯拉罕人生的歷史上,猶如空白的紙張.今天有許多基督徒,讀經覺得乏味,讀來讀去都是「你們」「我們」「他們」.毫無心得,找不到亮光.曾經有位老太太對我說,她的禱告總是禱來禱去禱不通.我說:禱告不通是「塞住了」.我們的靈性常像囚在監牢一樣,因為和神的關係已經斷絕了.親愛的弟兄姊妹!你已有多久發現神和你交通斷絕了呢?亞伯拉罕違背神的命令,用自己的方法代替等候神,以致神向他隱藏.那十三年是亞伯拉罕畢生最痛苦的階段.到了他九九歲時,神才向他顯現.當他百歲時,撒拉給他生了一個兒子,使這對老夫妻喜出望外；得意洋洋.但是可憐的夏甲被利用的價值至此完了!主母要主人把他們母子倆趕出去.亞伯拉罕當然於心不忍,相處多年,恩情已如海洋之深重了.他只好向神申訴,希望神主持公道.但是神對他說:「凡撒拉對你說的話,你都該聽從.」(創廿一12)人就是這樣,有難處時才求問神.當時,撒拉要亞伯拉罕娶夏甲,正是他稱心如意的事；他並無求問神,便毫不猶豫地聽從了撒拉的話.這次,神還要他聽從撒拉的話.雖然他對夏甲母子依依不捨,亦無可奈何.弟兄姊妹!人的痛苦常是自取之咎!我們不能怨天尤人.亞伯拉罕終於把夏甲母子趕走了,這不但令他的心痛苦萬分,也給後代留下極大的禍根.至今中東危機是因亞伯拉罕造成的,阿拉伯人和以色列人成了世仇,真是一失足成千古恨!
\newline
\newline
3. 雅各　　是個最會欺騙的人.雖然他誠心要得神的福分,但他為求達到目的而不擇手段.他欺騙父兄,祝福雖然得了,可是卻因此而受苦.為了哥哥要殺他,不得不離鄉別母親,疲於奔命.經跋踄,歷艱辛,投靠他的舅父.誰碰著他,誰就被他欺騙；從他舅父手裡,他得了舅父兩個女兒和許多肥美的羊群.他空手而來,卻滿載而歸.到了他年老的時候,他最愛的兒子約瑟,被哥哥們賣了,他被騙以為約瑟給野獸吃了,他痛苦悲哀欲絕.人種的是甚麼,收的也是甚麼.
\newline
\newline
親愛的弟兄姊妹!請恭敬思想,我們年青時所犯的,有許多罪神已赦免了；可能有些罪,或將有一天,神要叫你自食其果.好像大衛犯了姦淫殺人之罪,報應就在他的家中.如果我們相信神,信神是掌管一切的,求神保守我們；使我們手潔心清,在神前,人前,撒旦前,作美好的見證.
\newline
\newline
4. 以色列人　　神領他們出埃及,過紅海到了西乃曠野時,神呼喚摩西上山.「百姓見摩西遲延不下山,就大家聚集到亞倫那裡；對他說,起來!為我們作神像,可以在我們面前引路.因為領我們出埃及地的那個摩西,我們不知道他遭了甚麼事.......鑄了一隻牛犢......他們就說,以色列阿!這是領他出埃及地的神.」(出卅二 1-4)他們以金牛犢代替神,他們的愁苦必加增.在這廿世紀太空時代,是個感覺的時代,看得見的,摸得著的,放在試管裡可以試驗的,才會領人相信.「信心」二字在我們的字典已經失落了.我們中間,不知有多少人還在拜金牛犢.「以眼見」代替「信心」,以「看得見的」代替「看不見的神.」以破裂不能存水的池子,代替神的計劃.
\newline
\newline
5. 烏撒　　當神的約櫃被非利士人擄去後,神的榮耀,能力,在非利士人中間彰顯,使他們不得不將約櫃送還猶大人.後來約櫃遷移到俄別以東的家,途中因為牛失前蹄,烏撒伸手扶住約櫃,致被神擊殺.聖經告訴我們,約櫃應由分別為聖的利未人摃抬.但他們卻用外邦人的方法,以牛車運載.非利士人這樣做,神任憑他們.但神的兒女這樣做,神就擊殺那用手扶約櫃的人.為何神叫牛失前蹄?當神的工作失去應有的樣子時,神寧可讓祂的約櫃摔碎.烏撒用手去扶,怪不得神擊殺他.英文有兩字 man 人,method 方法.神的工作不是用任何方法才能復興,乃是需要分別為聖的人.可惜今天教會所注重的是方法,忘記了神所使用的是怎樣的人.
\newline
\newline
每個公會,宗派,甚至團體,開始時都是很好的.開始的領袖,無論是馬丁路德,衛斯理約翰或其他神所用的僕人,他們都非常忠心向著神.他們好像發現了活水的泉源,用杯接著活水,告訴弟兄姊妹們杯裡有活水,我們一同來領受.他們繼續從活水的泉源領受,慢慢地杯裡的水乾了,杯子還在,裡面卻沒有水,人便開始崇拜那杯子.英文有兩字,movement 運動,monument 紀念碑.今天有很多 movement 成了 monument,已經沒有生命印證在其中,因為用了別的來代替.在我們工作的開始,真是看見神的榮耀,逐漸地,工作雖還在進行,但神的靈已經離開了.好像杯裡的水已經乾了,人便用各種的方法手段,吸引人來看這杯.
\newline
\newline
近幾年來,神賜福給東南亞的教會.但也不能否認在許多事上,叫神得不到榮耀,叫撒旦有羞辱我們的把柄.原因則「用別的來代替」,沒有用神的話,來遵行神的旨意；用神的方法來做神的工作,所以今日教會有愁苦.審判要從神的家開始.今天我們中間,有否假冒為善的；有否驕傲的心,詭詐的心?正如聖經所說:我們有一切的山窪,有大小的山崗,有彎彎曲曲的地方,有高高低低的道路.巴不得求神憐憫我們,不再以別的來代替,叫我們內心得到祂所賜的平安.求聖靈光照我們,巴不得勝靈的能力在我們內心,叫我們完完全全的向著主.
\newline
\newline


\newpage

\section{第43屆~港九培靈會~楊濬哲~序}
\label{sec:1382}
\textbf{序}
\newline
\newline
連結: \href{https://www.hkbibleconference.org/session-message/view/1382}{\texttt{https://www.hkbibleconference.org/session-message/view/1382}}
\newline
\newline
\hyperref[sec:1418]{< < < PREV SERMON < < <}
~
\hyperref[sec:toc]{[返目錄]}
~
\hyperref[sec:1395]{> > > NEXT SERMON > > >}
\newline
\newline
憶一九七○年,第四十二屆大會閉會感恩會時,王永信牧師對委辦會同人,致詞勉勵,盼望我們集中力量,辦好培靈會的工作.委辦會經過多次祈禱商討,深感港方大會,環境噪雜.會場除灣仔循道會外,未有較適中之地點；且委辦同人,工作羈身,未能全時間赴會,至感遺憾!雖然九龍循道會,每年借用場址,並且給我們許多方面之幫助,我們實在衷心感激!但為了能有較大之場地,以容納更多會眾起見；我們遂徵得九龍城浸信會同意,允借用堂址,舉辦第四十三屆大會.
\newline
\newline
林子豐博士對培靈會工作,關懷深切!可惜在四十三屆大會舉行前,他已安息主懷.但他在離世前,對培靈會借用堂址,早已有妥善安排.願主記念和賞賜他!
\newline
\newline
第四十三屆大會,轉移場地,是一個新的嘗試,誰也未能預料後果如何?感謝主!實在恩待培靈會.在末了的幾天聚會中,浸信會正堂,閣樓,副堂,均全部滿座,約二千五百多人；這種情形,是前所未見的.至於奉獻也超過了以往的數字.第廿二集培靈講道,還未閉會,二千本全數售罄!因此,本集要加印至三千本了.
\newline
\newline
本屆的三位講員:研經會由美北浸信宣道會,祁柏牧師主講彼得前書.他是一位神學院教授,對解經甚有心得.講道會由鮑會園牧師主講以斯拉記.鮑牧師是播道神學院的院長,該書真理對今日信徒,有重大關係.奮興會由美北福音使者協會,總幹事周主培牧師主講,他言詞懇切感人,激動了不少青年人獻身事主.
\newline
\newline
記錄的有蔣國仁弟兄,譚雅芳姊妹,和雷麗端姊妹.編輯的有朱建磯牧師,他們對本書的出版,盡了很大的努力!
\newline
\newline
我們更要感謝九龍城浸信會的同工,在會期間的幫助和招待部的熱誠招待,給與我們很大的鼓勵.
\newline
\newline
在此,我得介紹金新宇博士,歡迎他參加委辦會的工作,與我們同心事奉主!
\newline
\newline
晚堂奮興會的錄音帶,已翻印出售,對海外和本港的基督徒來說,這是一個好消息.
\newline
\newline
末了,願將一切榮耀獻給主,主若遲延再來,盼望培靈會也繼續舉辦下去.
\newline
\newline


\newpage

\section{第43屆~港九培靈會~祈柏~第一講}
\label{sec:1395}
\textbf{我們活潑的盼望}
\newline
\newline
連結: \href{https://www.hkbibleconference.org/session-message/view/1395}{\texttt{https://www.hkbibleconference.org/session-message/view/1395}}
\newline
\newline
\hyperref[sec:1382]{< < < PREV SERMON < < <}
~
\hyperref[sec:toc]{[返目錄]}
~
\hyperref[sec:1396]{> > > NEXT SERMON > > >}
\newline
\newline
彼得前書綱要
\newline
\newline
彼得前書之序言
\newline
\newline
Ａ　他們是一群試煉的子民
\newline
\newline
1.經歷多次試煉 一:6
\newline
\newline
2.為神的榮耀 一:7
\newline
\newline
Ｂ　他們是一群聖潔的子民
\newline
\newline
1.他們是恩典的兒女 一:13
\newline
\newline
2.不是屬世的兒子 一:14
\newline
\newline
Ｃ　他們是一群誠實的子民
\newline
\newline
1.他們避免作世界情慾糾纏之客旅 二:11
\newline
\newline
2.他們在不信者面前要誠實 二:12
\newline
\newline
Ｄ　他們是一群恩慈的子民
\newline
\newline
1.他們存慈憐的心 三:8
\newline
\newline
2.他們以善報惡 三:9
\newline
\newline
Ｅ　他們是一群喜樂的子民
\newline
\newline
1.他們在試煉中歡喜 四:11-12
\newline
\newline
2.他們等候喜樂之日 四:14
\newline
\newline
我們活潑的盼望
\newline
\newline
Ａ　盼望的本質 一:3-5
\newline
\newline
1.它是活潑的 一:3
\newline
\newline
2.它是永恆的 一:4
\newline
\newline
3.它是確實的 一:5
\newline
\newline
Ｂ　盼望的喜樂 一:6-9
\newline
\newline
1.在患難中  一:6-7
\newline
\newline
2.在基督裡 一:8
\newline
\newline
3.在靈魂的救恩  一:9
\newline
\newline
Ｃ　盼望的啟示 一:10-12
\newline
\newline
1.為眾先知所宣告 一:10
\newline
\newline
2.為聖靈所啟示 一:11
\newline
\newline
3.向信徒傳揚 一:12
\newline
\newline
聖潔生命之供養
\newline
\newline
彼前一:13-二:10
\newline
\newline
Ａ　我們是聖潔的子民 一:13-17
\newline
\newline
1.我們是祂順命的兒女 一:13-14
\newline
\newline
2.祂是我們聖潔之父 一:15-16
\newline
\newline
3.我們要以生活來敬畏祂 一:17
\newline
\newline
Ｂ　我們有聖潔的救贖 一:18-21
\newline
\newline
1.憑聖潔寶血得蒙救贖 一:18-19
\newline
\newline
2.由一位權能的神得蒙救贖 一:20
\newline
\newline
3.藉基督的復活得蒙救贖 一:21
\newline
\newline
Ｃ　我們以聖潔之道來飼養 一:22-二:3
\newline
\newline
1.這是潔淨之道 一:22-23
\newline
\newline
2.這是永恆之道 一:24-25
\newline
\newline
3.這是真誠之道 二:1-3
\newline
\newline
Ｄ　我們有聖潔的使命 二:4-10
\newline
\newline
1.我們是聖潔的祭司 二:4-5
\newline
\newline
2.我們是聖潔的國度 二:6-9
\newline
\newline
3.我們是神的子民 二:10
\newline
\newline
蒙召受苦
\newline
\newline
彼前二:11-25
\newline
\newline
Ａ　受苦的境況 二:11-14
\newline
\newline
1.那些與靈魂爭戰之事 二:11
\newline
\newline
2.那些毀謗基督徒精神的人 二:12
\newline
\newline
3.那些我們無法縱控的制度 二:13-14
\newline
\newline
Ｂ　受苦的裡面 二:15-20
\newline
\newline
1.這是神的旨意 二:15-17
\newline
\newline
2.這是為神的榮耀 二:18-19
\newline
\newline
3.這證實神的恩典 二:20
\newline
\newline
Ｃ　受苦的天職 二:21-25
\newline
\newline
1.基督的榜樣 二:21
\newline
\newline
2.祂爭取權利赴死 二:22-24
\newline
\newline
3.我們應有權利而活 二:25
\newline
\newline
為信仰而活
\newline
\newline
Ａ　在受苦中要喜樂 四:12-19
\newline
\newline
1.在受苦中不要驚訝 四:12-13
\newline
\newline
2.以基督徒身份確實去受苦 四:14-16
\newline
\newline
3.審判要從神的家起首 四:17-19
\newline
\newline
Ｂ　在事奉中要忠心 五:1-4
\newline
\newline
1.要作基督的見證 五:1
\newline
\newline
2.要作群羊的牧者 五:2-3
\newline
\newline
3.要配得面對主的顯現 五:4
\newline
\newline
Ｃ　不要憂慮 五:5-11
\newline
\newline
1.萬事要信賴主 五:5-7
\newline
\newline
2.務要謹守,儆醒 五:8-9
\newline
\newline
3.要靠神的幫助以致完全 五:10-11
\newline
\newline
註:本綱要乃祁柏牧師於培靈會查經會時,分送與會眾者,茲隨書刊登,以備作參攷之用.
\newline
\newline
─編者─
\newline
\newline


\newpage

\section{第43屆~港九培靈會~祈柏~第二講}
\label{sec:1396}
\textbf{聖潔生命之供養}
\newline
\newline
連結: \href{https://www.hkbibleconference.org/session-message/view/1396}{\texttt{https://www.hkbibleconference.org/session-message/view/1396}}
\newline
\newline
\hyperref[sec:1395]{< < < PREV SERMON < < <}
~
\hyperref[sec:toc]{[返目錄]}
~
\hyperref[sec:1397]{> > > NEXT SERMON > > >}
\newline
\newline
彼得前書 1:13-2:10
\newline
\begin{longtable}{cl}
\hline
\hline
章節 & 經文 (和合本修訂版)\\
\hline
1:13 & \begin{tabularx}{0.7\textwidth}{X} 所以，要準備好你們的心，謹慎自守，專心盼望耶穌基督顯現的時候帶給你們的恩惠。 \end{tabularx} \\ \\ \relax
1:14 & \begin{tabularx}{0.7\textwidth}{X} 作為順服的兒女，就不要效法從前蒙昧無知的時候那放縱私慾的樣子。 \end{tabularx} \\ \\ \relax
1:15 & \begin{tabularx}{0.7\textwidth}{X} 但那召你們的既是聖潔，你們在一切所行的事上也要聖潔； \end{tabularx} \\ \\ \relax
1:16 & \begin{tabularx}{0.7\textwidth}{X} 因為經上記著：「你們要成為聖，因為我是神聖的。」 \end{tabularx} \\ \\ \relax
1:17 & \begin{tabularx}{0.7\textwidth}{X} 既然你們稱那不偏待人、按各人行為審判人的主為父，就當存敬畏的心，度你們在世寄居的日子。 \end{tabularx} \\ \\ \relax
1:18 & \begin{tabularx}{0.7\textwidth}{X} 你們知道，你們得以從你們祖先傳下來虛妄的行為中救贖出來，不是靠著會朽壞的金銀等物， \end{tabularx} \\ \\ \relax
1:19 & \begin{tabularx}{0.7\textwidth}{X} 而是憑著基督的寶血，如同無瑕疵、無玷污的羔羊的血。 \end{tabularx} \\ \\ \relax
1:20 & \begin{tabularx}{0.7\textwidth}{X} 基督是神在創世以前所預知，而在這末世才為你們顯現的。 \end{tabularx} \\ \\ \relax
1:21 & \begin{tabularx}{0.7\textwidth}{X} 你們也因著他而信那使他從死人中復活、又給他榮耀的神，好讓你們的信心和盼望都在於神。 \end{tabularx} \\ \\ \relax
1:22 & \begin{tabularx}{0.7\textwidth}{X} 既然你們因順從真理而潔淨了自己的心靈，能真誠愛弟兄，就該以清潔的心彼此切實相愛。 \end{tabularx} \\ \\ \relax
1:23 & \begin{tabularx}{0.7\textwidth}{X} 你們蒙了重生，不是由於會朽壞的種子，而是由於不會朽壞的種子，是藉著神永活常存的道。 \end{tabularx} \\ \\ \relax
1:24 & \begin{tabularx}{0.7\textwidth}{X} 因為「凡血肉之軀的盡都如草，他的一切榮美像草上的花；草必枯乾，花必凋謝， \end{tabularx} \\ \\ \relax
1:25 & \begin{tabularx}{0.7\textwidth}{X} 惟有主的道永遠常存。」這話就是傳給你們的福音。 \end{tabularx} \\ \\ \relax
2:1 & \begin{tabularx}{0.7\textwidth}{X} 所以，你們要除去一切的惡毒，一切詭詐、假善、嫉妒，和一切毀謗的話。 \end{tabularx} \\ \\ \relax
2:2 & \begin{tabularx}{0.7\textwidth}{X} 要愛慕那純淨的靈奶，像初生的嬰孩愛慕奶一樣，好使你們藉著它成長，以致得救， \end{tabularx} \\ \\ \relax
2:3 & \begin{tabularx}{0.7\textwidth}{X} 因為你們已經嘗過主恩的滋味。 \end{tabularx} \\ \\ \relax
2:4 & \begin{tabularx}{0.7\textwidth}{X} 要親近主，他是活石，雖然被人所丟棄，卻是神所揀選、所珍貴的。 \end{tabularx} \\ \\ \relax
2:5 & \begin{tabularx}{0.7\textwidth}{X} 你們作為活石，要被建造成屬靈的殿，成為聖潔的祭司，藉著耶穌基督獻上蒙神悅納的屬靈祭物。 \end{tabularx} \\ \\ \relax
2:6 & \begin{tabularx}{0.7\textwidth}{X} 因為經上說：「看哪，我把一塊石頭放在錫安—一塊蒙揀選、珍貴的房角石；信靠他的人必不蒙羞。」 \end{tabularx} \\ \\ \relax
2:7 & \begin{tabularx}{0.7\textwidth}{X} 所以，這石頭在你們信的人是珍貴的；在那不信的人卻有話說：「匠人所丟棄的石頭已作了房角的頭塊石頭。」 \end{tabularx} \\ \\ \relax
2:8 & \begin{tabularx}{0.7\textwidth}{X} 又說：「作了絆腳的石頭，使人跌倒的磐石。」他們絆跌，因為不順從這道，這也是預定的。 \end{tabularx} \\ \\ \relax
2:9 & \begin{tabularx}{0.7\textwidth}{X} 不過，你們是被揀選的一族，是君尊的祭司，是神聖的國度，是屬神的子民，要使你們宣揚那召你們出黑暗入奇妙光明者的美德。 \end{tabularx} \\ \\ \relax
2:10 & \begin{tabularx}{0.7\textwidth}{X} 「你們從前不是子民，現在卻成了神的子民；從前未曾蒙憐憫，現在卻蒙了憐憫。」 \end{tabularx} \\ \\
[1ex]
\hline
\hline
\end{longtable}
我們有了活潑的盼望後,如何善後呢?主在世上時,每逢將真理講解後,便吩咐門徒照著去行.故使徒約翰說:「凡向他有這指望的,就潔淨自己,像他潔淨一樣.」(約壹三3)
\newline
\newline
我們若有了活潑的盼望,就真正相信主必再來,這信念推動我們過著聖潔的生活,彼得在此也同樣的提出過聖潔生活的四個原因:
\newline
\newline
(一)我們是聖潔的子民(一13至17)
\newline
\newline
我們為何要過聖潔的生活?「你們要聖潔,因為我是聖潔的.」可見因為我們的天父是聖潔,我們是聖潔的子民,所以要過聖潔的生活.
\newline
\newline
1. 我們是祂順命的兒女(13,14節)
\newline
\newline
我們為何要過聖潔的生活呢?在十四節提出了答案,「你們既作順命的兒女,就不要效法從前蒙昧無知的時候,那放縱私慾的樣子.」近代福音傳播工具異常發達,我們很喜歡模仿電視銀幕上明星們的舉動,但基督徒卻應以聖經為銀幕,以聖經所記載之義人為效法對象.
\newline
\newline
兒女們能作光宗耀祖的事,做父母的也覺得光榮；若兒女不肖,使家門羞辱,為父牧者也覺難堪.所以神的兒女們,要存心討天父的喜悅,使天父以我們為榮,使神看見主耶穌為我們死,所成就的救恩沒有白廢.
\newline
\newline
2. 祂是我們聖潔之父(15,16節)「那召你們的既是聖潔,你們在一切所行的事上也要聖潔.因為經上記著說:『你們要聖潔,因為我是聖潔的.』」簡言之有兩點,(一)是因天父是聖潔的,(二)是神的話勸告我們要過聖潔的生活.
\newline
\newline
3. 我們要以生活來敬畏祂(17節)「你們既稱那不偏待人,按各人行為審判人的主；就當存敬畏的心,度你們在世寄居的日子.」
\newline
\newline
在十八節還有一個有趣的詞語:「就是不偏待人的,按各人行為審判人的主」在社會的法律上,有時會出偏差,犯人可賂賄法官,警察,議員,使犯人不用受審,也不用定罪.有權威的便有影響力,但神是聖潔的,你不能賄賂祂,祂也不受政治勢力的影響,來改變祂的審判.故此,彼得指出我們要過聖潔的生活,這並非為了畏懼神,實在是為了愛祂之故.
\newline
\newline
十八節以下說到神親自為我們預備救恩,聖潔的神不願一人下入地獄,乃願人人得救:凡讀過彼得書的人,都知道這作事.
\newline
\newline
1. 憑聖潔寶血得蒙救贖(18至19)節)「知道你們得贖,脫去你們祖宗所傳流虛妄的行為,不是憑著能壞的金銀等物,乃是憑著基督的寶血,如同無瑕疵,無玷污的羔羊之血.」聖潔的恩主,卻為我們背負了罪孽,甘心為我們而死.主耶穌如此愛我們,我們也應當有愛主的心,為愛主而過聖潔的生活.正像作父母的必然喜歡子女過著聖潔純良的生活一樣.否則,子女就要受責打.若子女因為父母而過著良善生活,非但不會受責打,彼此之間就能夠愛心相通,過著和諧美滿的生活.
\newline
\newline
天父也不想用權威來迫使我們過聖潔的生活.祂是慈愛的,甚願我們能因為愛祂才遵守他的命令,過著聖潔的生活.
\newline
\newline
2. 由一位權能的神得蒙救贖.(20節)「基督在創世以前,是豫先被神知道的,卻在這末世,才為你們顯現.」在以弗所書一章四節說:「就是神從創立世界以前,在基督裡揀選了我們.」我們從這兩處聖經可以了解到神的心,知道祂在永恆中,創造了世界,也揀選了我們.獨惜人要奪取神的權柄,讓罪進入了世界,敗壞了世界.神原可以不創造世界,免得世界敗壞.不過神不改變祂的旨意和計劃,寧願預備耶穌,來救贖世界.神也早已預備迦略山,讓主耶穌為我們死在迦略山,因此我們便成了神愛的對象.
\newline
\newline
可見我們過聖潔的生活,不應是畏懼祂,也不是為了在律法上負責任,乃是因為愛祂.
\newline
\newline
3. 藉基督的復活得蒙救贖.(21節)」你們也因著祂,信那叫祂從死裡復活,又給祂榮耀的神；叫你們的信心和盼望,都在於神.」我們所唱的詩歌有「祂是我心所盼望.」不錯,我們不敢靠人,只靠神的羔羊.
\newline
\newline
(三)我們以聖潔之道來飼養(一22至二3)
\newline
\newline
我們要過著聖潔的生活,因為祂已供應我們聖潔的糧食.
\newline
\newline
1. 這是潔淨之道.(22至23節)「你們既因順從真理,潔淨了自己的心,以致愛弟兄沒有虛假,就當從心裡彼此切實相愛.你們蒙了重生,不是由於能壞的種子,乃是由於不能壞的種子,是藉著上帝活潑常存的道.」彼得和約翰同樣講及重生,因為惟有重生的人,才能進入神的國.肉體的生命是會過去的,當一個嬰孩生下來的時候,我們就想到,他要漸漸長大,會老,會死.在我們衰老時身體的各部份,都要退化,像眼矇,耳聾,牙齒脫落,白髮生.許多美國婦人,最怕白髮,她們把長出來的白髮染黑了,使人以為她還年輕.可是那些髮根卻不允合作,繼續長出來的還是白的.既有衰老的跡象,就該早些作永恆的打算了.
\newline
\newline
我們從母胎出生,是憑能壞的種子,但重生卻是由於不能壞的種子,這種子給我們永生.
\newline
\newline
神的道,一方面是不能朽壞的種子,給我們永生.另方面也是我們的靈糧.
\newline
\newline
2. 是永恆之道.(24至25節)「因為『凡有血氣的,盡都如草.他的美榮,都像草上的花,草必枯乾,花必凋謝.惟有主的道是永存的.』所傳給你們的福音,就是這道.」
\newline
\newline
3. 這是真誠之道.(二1至3)「所以你們既除去一切的惡毒,詭詐,並假善,嫉妒和一切毀謗的話.」神首要的要求,是要我們在祂面前做一個聖潔的人.到底聖潔的內容是甚麼?我們要弄清楚了,才能追求聖潔的生活.細看下面兩節經文便知道了,這兩節也是本書的鑰節:「就要愛慕那純淨的靈奶,像才生的嬰孩愛慕奶一樣,叫你們因此漸長,以致得救.你們若嘗過主恩的滋味,就必如此.」
\newline
\newline
保羅在提摩太後書三章十六至十七節也說過:「聖經是神所默示的,於教訓,督責,使人歸正,教導人學義,都是有益的,叫屬神的人得以完全,預備行各樣的善事.」當我們有不正當的行為,神的話不但指出我們的錯誤,而且教導我們行正當的事.神的話又無所不備,足以教導神的兒女們怎樣在世人面前行事為人.
\newline
\newline
我們要遵行神的道,必須先知道神的話如何說,且要常常學習.有人說,必須熟習神的話,神的話才能操縱我們的人生.
\newline
\newline
許多人有一個錯誤的想法,以為加入教會,受過水禮,也知道一些教義和禮節,便是一個基督徒.其實這只能叫他成為一個有宗教儀式的人而已,絕不能叫他成為真正的基督徒.
\newline
\newline
當我受聘在浸信會當牧職時,教會的執事便問我,你到教會來要做些甚麼工作?我回答說,首要的是要將浸信會的信徒,成為基督徒.因為有人加入教會,像加入其他團體作會員員一樣.
\newline
\newline
做了神聖潔的子民.有聖靈在我們心中,我們就能到不聖潔的民中為祂作見證.
\newline
\newline
(四)我們有聖潔的使命(二4至10)
\newline
\newline
我們怎樣在不信的人中,為神作聖潔的見證?
\newline
\newline
1. 我們是聖潔的祭司.(4至5節)「主乃活石,固然是被人所棄的,卻是被上帝所揀選所寶貴的.你們來到主面前,也就像活石,被建造成為靈宮,作聖潔的祭司,藉著耶穌基督奉獻神所悅納的靈祭.」彼得在這裡不祇說到一個活石(4節,而是說到許多活石.(5節),因此,我請各位細讀下面一段娙文:「聽見約翰的話,跟從耶穌的那兩個人,一個是西門彼得的兄弟安得烈,他先找著自己的哥哥西門,對他說:『我們遇見彌賽亞了.』(彌賽亞繙出來就是基督),於是領他去見耶穌.耶穌看著他說,你是約翰的兒子西門,你要稱為磯法(磯法繙出來,就是彼得)東方人為兒女起名喜歡表明他的性格,或表明對這孩子的期望.這裡提到彼得有兩個名字.一個是西門,意即不穩定,虛浮,無定見,無主張.細看彼得的性情確是如此,虛浮,衝動,暴躁,懦弱.曾被一使女詢問之下,竟矢口不認主.主耶穌深知他的性情,正是「西門」.另一個名是彼得,主耶穌要改變他的性情,由西門成為彼得.美國諺語有云:「兒女從父親取出來的.」中國也有諺語云:「餅印印出來的.」彼得正是大石(主耶穌)取出來的小石.西門變成彼得,彼得已不是朽木,經已變為活石了.
\newline
\newline
主乃大活石(4節),眾基督徒是小活石(5節).我與祁師母曾往參觀過開採石礦,鑿出來的石碩,用在各樣的建築上.
\newline
\newline
2. 我們是聖潔的國度.(9節)「惟有你們是被揀選的族類,是君尊的祭司,是聖潔的國度,是屬神的子民,要叫你們宣揚那召你們出黑暗入奇妙光明者的美德.」在神的眼中,我不是美國人,不是中國人,也不是白種人,不是黃種人,世上不論是那一種族,同是他所救贖的.各位,我們儘管外表不同,也可以被人看為合一的.
\newline
\newline
我們為什麼是君尊的祭司,是聖潔的國度,是屬神的子民?原因是要出黑暗,入奇妙光明.
\newline
\newline
我們以甚麼引以為榮?若能忘記地區性的優越感,不論中西人士,而成為神家中的人,這是何等美的事,實在是神所寶貝的百姓.
\newline
\newline
3. 我們是神的子民.(10節)「你們從前算不得子民,現在卻作了神的子民,從前未曾蒙憐恤,現在卻蒙了憐恤.」
\newline
\newline
我們既蒙了神的奇妙大恩,應在世人面前,過聖潔的生活,表明我們過去雖曾「被放逐」,又如「難民」,如今卻雙重歸屬神了.
\newline
\newline


\newpage

\section{第43屆~港九培靈會~祈柏~第三講}
\label{sec:1397}
\textbf{蒙召受苦}
\newline
\newline
連結: \href{https://www.hkbibleconference.org/session-message/view/1397}{\texttt{https://www.hkbibleconference.org/session-message/view/1397}}
\newline
\newline
\hyperref[sec:1396]{< < < PREV SERMON < < <}
~
\hyperref[sec:toc]{[返目錄]}
~
\hyperref[sec:1398]{> > > NEXT SERMON > > >}
\newline
\newline
彼得前書 2:11-2:25
\newline
\begin{longtable}{cl}
\hline
\hline
章節 & 經文 (和合本修訂版)\\
\hline
2:11 & \begin{tabularx}{0.7\textwidth}{X} 親愛的，你們是客旅，是寄居的，我勸你們要禁戒肉體的情慾；這情慾是與靈魂爭戰的。 \end{tabularx} \\ \\ \relax
2:12 & \begin{tabularx}{0.7\textwidth}{X} 你們在外邦人中要品行端正，好讓那些人，雖然毀謗你們是作惡的，會因看見你們的好行為而在鑒察的日子歸榮耀給神。 \end{tabularx} \\ \\ \relax
2:13 & \begin{tabularx}{0.7\textwidth}{X} 你們為主的緣故要順服人的一切制度，或是在上的君王， \end{tabularx} \\ \\ \relax
2:14 & \begin{tabularx}{0.7\textwidth}{X} 或是君王所派懲惡賞善的官員。 \end{tabularx} \\ \\ \relax
2:15 & \begin{tabularx}{0.7\textwidth}{X} 因為神的旨意原是要你們以行善來堵住糊塗無知人的口。 \end{tabularx} \\ \\ \relax
2:16 & \begin{tabularx}{0.7\textwidth}{X} 雖然你們是自由的，卻不可藉著自由遮蓋惡毒，總要作神的僕人。 \end{tabularx} \\ \\ \relax
2:17 & \begin{tabularx}{0.7\textwidth}{X} 務要尊重眾人；要敬愛教中的弟兄姊妹；要敬畏神；要尊敬君王。 \end{tabularx} \\ \\ \relax
2:18 & \begin{tabularx}{0.7\textwidth}{X} 你們作奴僕的，凡事要存敬畏的心順服主人；不但順服善良溫和的，就是乖僻的也要順服。 \end{tabularx} \\ \\ \relax
2:19 & \begin{tabularx}{0.7\textwidth}{X} 倘若你們為使良心對得起神，忍受冤屈的痛苦，這是可讚許的。 \end{tabularx} \\ \\ \relax
2:20 & \begin{tabularx}{0.7\textwidth}{X} 你們若因犯罪受責打而忍耐，有甚麼可稱讚的呢？但你們若因行善受苦而忍耐，這在神看來是可讚許的。 \end{tabularx} \\ \\ \relax
2:21 & \begin{tabularx}{0.7\textwidth}{X} 你們蒙召就是為此，因為基督也為你們受過苦，給你們留下榜樣，為要使你們跟隨他的腳蹤。 \end{tabularx} \\ \\ \relax
2:22 & \begin{tabularx}{0.7\textwidth}{X} 「他並沒有犯罪，口裡也沒有詭詐。」 \end{tabularx} \\ \\ \relax
2:23 & \begin{tabularx}{0.7\textwidth}{X} 他被辱罵不還口，受害也不說威嚇的話，只將自己交託給公義的審判者。 \end{tabularx} \\ \\ \relax
2:24 & \begin{tabularx}{0.7\textwidth}{X} 他被掛在木頭上，親身擔當了我們的罪，使我們既然在罪上死，就得以在義上活。因他受的鞭傷，你們得了醫治。 \end{tabularx} \\ \\ \relax
2:25 & \begin{tabularx}{0.7\textwidth}{X} 你們從前好像迷路的羊，如今卻歸回你們靈魂的牧人和監督了。 \end{tabularx} \\ \\
[1ex]
\hline
\hline
\end{longtable}
彼得在此段經文,提到如何將我們內在的聖潔,在世人面前表明出來.
\newline
\newline
(一)受苦的境況(二11-14)
\newline
\newline
1. 那些與靈魂爭戰之事(11節)「親愛的弟兄啊!你們是客旅,是寄居的,我勸你們要禁戒肉體的私慾,這私慾是與靈魂爭戰的.」
\newline
\newline
本節經文彼得對我們有三個稱呼:「親愛的弟兄」,「客旅」和「寄居的」.
\newline
\newline
在美國通稱那些歸主的人為「重生」,「得救」者.但許多人卻稱歸主為「蒙恩」,這名詞正好讓我們明白彼得用「親愛的」的意思,他好像說「你們是神所親愛的」.我們本來該受神的忿怒,反而得了神的恩惠,故此稱為親愛的.正如使徒約翰所說:「不是我們愛神,乃是神先愛我們.」
\newline
\newline
「客旅」,我曾到過歐洲,北菲等地,我懂得當地的語言文字所以我不覺得是異鄉.但來到東方的日本,香港,台灣等地,便覺得異鎯了,因為我不懂這些地方的語言文字.我們在世上雖是客旅,卻是天國的子民,所以要除掉世上的污穢,學習天國國民的生活.
\newline
\newline
過去,若為人不誠實,不良善,惡待人；現今要脫去這些行為,行事為人要與所信的福音相稱.是神的百姓,就該有屬天的生活.
\newline
\newline
2. 那些毀謗基督徒精神的人.(12節)「你們在外邦人中,應當品行端正,叫那些毀謗你們是作惡的,因為看見你們的好行為,便在鑑察的日子,歸榮耀給神.」人雖以惡待我,欺侮我,我都不能以惡報惡待他們.我們既蒙了主恩,就不能按律法,以眼還眼,以牙還牙了.因為主饒恕了我們,我們也當饒恕別人,主曾說:「你們願意人怎樣待你們,你們也要怎樣待人.」
\newline
\newline
不過,你雖立志為主而活,處處謙讓,世人可能不會欣賞你,反而趁機佔你便宜,到此時,或許有人會問,到底主知道我為祂而活嗎?在聖經裡,主耶穌沒有應許我們,若為主而活,便立刻報答我們.那麼,主報答在何時?請注意「鑑察的日子」,即是主再來之時.因此我們現在不是憑眼見,乃是憑信心等候.
\newline
\newline
「鑑察」原文有「查賬」之意,有一日,主要平衡結算,到那日,那些惡待你的人,豈不歸榮與天父.主耶穌說過:「你們的光,也當這樣照在人前,叫他們看見你們的好行為,便將榮耀歸給你們在天上的父.」(太五16)
\newline
\newline
3. 那些我們無法操縱的制度.(13至14節)「你們為主的緣故,要順服人的一切制度,或是在上的君王,或是君王所派,罰惡賞善的臣宰.」與我們相處之人,有各樣各式之人,那些有權柄管理我的人,有君王,臣宰,商店的老闆,機關的上司.我們不同意接受無政府主義,也不想在無政府狀態下任意妄為.「你們要為主的緣故,要順服人的一切制度.」主耶穌曾提出同樣意義的事,祂說「該撒的物當歸給該撒,神的物當歸給神.」若該撒的政府與神國沒有衝突.當然沒有問題,但若信仰上受了逼迫,也該站在主耶穌的一邊,忠心到底.或有人會問,為主忠心,要經歷這許多艱難,是否值得?我要忠誠告訴你,我們有活潑的盼望,到那日,主必為你伸冤,那些惡待你的人,要歸榮與天父.其實不需等到那日,就在今日,有人亦可為主受苦的緣故,使惡待他的人,受感歸主.這是一個更有力的見證.
\newline
\newline
彼得過去以為主耶穌會得金冠冕,但主所得的卻是荊棘冠冕；他以為主耶穌會成為彌賽亞(君王),自己便可為丞相；怎料他自己後來卻成了被放逐的,成為難民.故此彼得所講的是經歷,不是理論.
\newline
\newline
(二)受苦的裡面(二15-20)
\newline
\newline
1. 這是神旨意.(15至17節)
\newline
\newline
「因為神的旨意,原是要你們行善,可以堵住那糊塗無知人的口.」(15節)基督徒不一定是做錯了,才被人迫害.但因行善,也會被迫害,我們須忍辱負重,才能堵住了那些迫害者的口.
\newline
\newline
「你們雖是自由的,卻不可藉著自由遮蓋惡毒,總要作神的僕人.」(16節)我們無論做何事,雖有自由,但須持著謹守道德去使用自由.不要以惡報善,反要以恩慈相待.
\newline
\newline
「務要尊敬眾人,親愛教中的弟兄,敬畏上帝,尊敬君王.」(17節)這節總結以上所說.有些人我們要尊敬他,有些人我們要親愛他,但更要尊敬而又愛天上的父神.
\newline
\newline
2. 這是為神的榮耀(二18至19)
\newline
\newline
從這節開始,已進入研究本書信的高潮.我們為主受苦,不一定是因政府,或有權勢者迫害,可能由於人與人之間交往而來.
\newline
\newline
使徒時代的教會是彼此和睦的,今天的教會卻常有彼此嫉妒,分爭,惡待別人,霸佔公物等事發生.故此,我們要問,一個屬天的子民應該這樣麼?
\newline
\newline
「你們作僕人的,凡事要存敬畏的心順服主人,不但順服那良善溫和的,就是那乖僻的也要順服.」(18節)「敬畏」,不是懼怕,而是敬重.因為在神的眼中,那些有權勢的人,不只是你的官長,上司,老闆,而且看見基督也為他們死.有時我們想向官長,上司,老闆作見證,但他們不聽.但如果你以生活行為作見證,豈不更為有力嗎?
\newline
\newline
彼得知道神是至高無上的,連我們受苦,也在祂的掌握中,故此我們為主的緣故,行事為人,要本裡德行而行.
\newline
\newline
「倘若人為叫良心對得住神,就忍受冤屈的苦楚,這是可喜愛的.」(19節)神知道我們受苦,是否應當的.有些基督徒被辭退,或被革職,卻推諉說是為主受苦是不對的,其實是他做事不忠心.
\newline
\newline
3. 這證實神的恩典.(20節)「你們若因犯罪受責打,能忍耐,有甚麼可誇的呢!但你們若因行善受苦,能忍耐,這在神看是可喜愛的.」你騙我,我不知；我欺你,你不曉.但神知道人的內心,知道我們受苦是否該當的.
\newline
\newline
(三)受苦的天職(二21-25)
\newline
\newline
誰受苦,值得作為我們的榜樣呢?彼得在此引導我們到十字架下,仰望十字架上的基督.可見我們蒙召不但有活潑的盼望,蒙召也是為受苦.且在受苦,不隨從我們自己心中的反應,而專一順服神.
\newline
\newline
1. 基督是榜樣(21節)「你們蒙召原是為此,因基督也為你們受過苦,給你們留下榜樣,叫你們跟隨祂的腳蹤行.」
\newline
\newline
2. 祂爭取權利赴死(22至24節)「祂並沒有犯罪,口裡也沒有詭詐.祂被罵不還口,受害不說威嚇的話；只將自己交託那按公義審判人的主.」(22-23節)憑信心去做,有時覺得實在是困難,不過,就算困難,也要憑信心而作.
\newline
\newline
二十三節末句「只將自己交託那按公義審判人的主.」可見不是白白受苦,也見到在受苦中如何得勝.主耶穌被殘酷對待,沒有報復,只將自己交託那按公義審判人的主.
\newline
\newline
今天,除了主以外,沒有人被釘在十字架上,而會說:「父啊!赦免他們,因為他們所作的,他們不曉得的.」我們要學習的一件事,就是將自己交託那按公義審判人的主,因為他要按公義算賬,在「鑑察的日子」按公義審判.
\newline
\newline
或有人說,這太倫理化了,世界那有公義,能否提出一證據,證明有一天,神要按公義審判世界.我們這些疑問,雖屬幼稚,但神也憐憫我們,讓我們看看十字架上的主耶穌,「祂被掛在木頭上,親身擔當了我們的罪,使我們既然在罪上死,就得以在義上活.因他受的鞭傷,你們便得了醫治.」(24節)
\newline
\newline
3. 我們應有權利而活(25節)「你們從前好像迷路的羊,如今卻歸到你們靈魂的牧人監督了.」「監督」有在事務上查察之意,正如在商業上之會計,是查賬的.
\newline
\newline
我願告訴你們,我們可能在各地,受著不同的苦,但有一件是相同的,就是因為我們是為主的緣故,是在全能父神的手中.
\newline
\newline


\newpage

\section{第43屆~港九培靈會~祈柏~第四講}
\label{sec:1398}
\textbf{家庭生活}
\newline
\newline
連結: \href{https://www.hkbibleconference.org/session-message/view/1398}{\texttt{https://www.hkbibleconference.org/session-message/view/1398}}
\newline
\newline
\hyperref[sec:1397]{< < < PREV SERMON < < <}
~
\hyperref[sec:toc]{[返目錄]}
~
\hyperref[sec:1399]{> > > NEXT SERMON > > >}
\newline
\newline
彼得前書 3:1-3:22
\newline
\begin{longtable}{cl}
\hline
\hline
章節 & 經文 (和合本修訂版)\\
\hline
3:1 & \begin{tabularx}{0.7\textwidth}{X} 同樣，你們作妻子的，要順服自己的丈夫，這樣，即使有不信從道理的丈夫，也會因妻子的品行，並非言語，而感化過來， \end{tabularx} \\ \\ \relax
3:2 & \begin{tabularx}{0.7\textwidth}{X} 因為看見了你們敬虔純潔的品行。 \end{tabularx} \\ \\ \relax
3:3 & \begin{tabularx}{0.7\textwidth}{X} 你們不要藉外表來妝飾自己，如編頭髮，戴金飾，穿美麗的衣裳等， \end{tabularx} \\ \\ \relax
3:4 & \begin{tabularx}{0.7\textwidth}{X} 而要有蘊藏在人內心不衰退的美，以溫柔嫻靜的心妝飾自己；這在神面前是極寶貴的。 \end{tabularx} \\ \\ \relax
3:5 & \begin{tabularx}{0.7\textwidth}{X} 因為古時仰賴神的聖潔婦人正是以此為妝飾，順服自己的丈夫。 \end{tabularx} \\ \\ \relax
3:6 & \begin{tabularx}{0.7\textwidth}{X} 就如撒拉聽從亞伯拉罕，稱他為主。你們只要行善，不怕任何恐嚇，就成為撒拉的女兒了。 \end{tabularx} \\ \\ \relax
3:7 & \begin{tabularx}{0.7\textwidth}{X} 同樣，你們作丈夫的，要按情理跟妻子共同生活，體貼女性是比較軟弱的器皿；要尊重她，因為她也與你一同承受生命之恩。這樣，你們的禱告就不會受阻礙。 \end{tabularx} \\ \\ \relax
3:8 & \begin{tabularx}{0.7\textwidth}{X} 總而言之，你們都要同心，彼此體恤，相愛如弟兄，存憐憫和謙卑的心。 \end{tabularx} \\ \\ \relax
3:9 & \begin{tabularx}{0.7\textwidth}{X} 不要以惡報惡，以辱罵還辱罵，倒要祝福，因為你們正是為此蒙召的，好使你們承受福氣。 \end{tabularx} \\ \\ \relax
3:10 & \begin{tabularx}{0.7\textwidth}{X} 因為經上說：「凡要愛惜生命、享受好日子的人，要禁止舌頭不出惡言，嘴唇不說詭詐的話。 \end{tabularx} \\ \\ \relax
3:11 & \begin{tabularx}{0.7\textwidth}{X} 也要棄惡行善，尋求和睦，一心追求。 \end{tabularx} \\ \\ \relax
3:12 & \begin{tabularx}{0.7\textwidth}{X} 因為主的眼看顧義人，他的耳聽他們的祈禱；但主向行惡的人變臉。」 \end{tabularx} \\ \\ \relax
3:13 & \begin{tabularx}{0.7\textwidth}{X} 你們若熱心行善，有誰會害你們呢？ \end{tabularx} \\ \\ \relax
3:14 & \begin{tabularx}{0.7\textwidth}{X} 即使你們為義受苦，也是有福的。不要怕人的威嚇，也不要驚慌； \end{tabularx} \\ \\ \relax
3:15 & \begin{tabularx}{0.7\textwidth}{X} 只要心裡奉主基督為聖，尊他為主。有人問你們心中盼望的理由，要隨時準備答覆； \end{tabularx} \\ \\ \relax
3:16 & \begin{tabularx}{0.7\textwidth}{X} 不過，要以溫柔、敬畏的心回答。要存無虧的良心，使你們在何事上被毀謗，就在何事上使那些凌辱你們在基督裡有好品行的人自覺羞愧。 \end{tabularx} \\ \\ \relax
3:17 & \begin{tabularx}{0.7\textwidth}{X} 神的旨意若是要你們因行善受苦，這總比因行惡受苦好。 \end{tabularx} \\ \\ \relax
3:18 & \begin{tabularx}{0.7\textwidth}{X} 因為基督也曾一次為罪受苦 ，就是義的代替不義的，為要引領你們到神面前。在肉體裡，他被治死；但在靈裡，他復活了。 \end{tabularx} \\ \\ \relax
3:19 & \begin{tabularx}{0.7\textwidth}{X} 他藉這靈也曾去向那些在監獄裡的靈傳道， \end{tabularx} \\ \\ \relax
3:20 & \begin{tabularx}{0.7\textwidth}{X} 就是那些從前在挪亞預備方舟、神容忍等待的時候不信從的人。當時進入方舟，藉著水得救的不多，只有八個人。 \end{tabularx} \\ \\ \relax
3:21 & \begin{tabularx}{0.7\textwidth}{X} 這水所預表的洗禮，現在藉著耶穌基督的復活拯救你們，不是除掉肉體的污穢，而是向神懇求有無虧的良心。 \end{tabularx} \\ \\ \relax
3:22 & \begin{tabularx}{0.7\textwidth}{X} 耶穌已經到天上去，在神的右邊，眾天使、有權柄的、有權能的都服從了他。 \end{tabularx} \\ \\
[1ex]
\hline
\hline
\end{longtable}
第三章,彼得將說話的焦點,集中於家庭生活,尤其是夫妻關係上.一至七節,彼得所假設的問題,是妻子已相信主,而丈夫還未得救的.問題的關鍵,是如何使丈夫也如妻子一樣蒙神恩惠,得蒙拯救呢!
\newline
\newline
「你們作妻子的,要順服自己的丈夫,這樣,若有不信從道理的丈夫,他們雖然不聽道,也可以因妻子的品行被感化過來.」(一節)在這裡說明,即使丈夫沒有興趣聽道,不想認識神；但也可以因妻子良善的品行,而被感化信主.
\newline
\newline
有時做妻子的急於想丈夫信主,終日在丈夫面前,嚕囌的講見證,或將電台所播福音信息,吵耳的播放,丈夫下了班想在家中安靜,或閱讀一些報上新聞,已給吵耳的收音機聲浪紛擾,不得不走到外面去找尋娛樂.有時妻子對丈夫強用抗拒的方法,以為信主的家庭不應抽煙,喝酒；所以將家中的香煙丟掉,把酒倒光.丈夫看見,自然火光!試想這樣能引導丈夫歸主嗎?
\newline
\newline
經文中「不聽道」,可作「不說一句話」或解作不用多說話,只用你的貞潔的品行,便可使他感化.
\newline
\newline
第二節有「敬畏」兩字,「畏」不是懼怕,乃是「尊敬」與「敬重」之意,丈夫雖未信住,妻子也當敬重他.若不尊敬他,容易使他氣忿.
\newline
\newline
妻子信了主,以為家庭是自己的,不是丈夫的.所以丈夫要做甚麼事情,只得到外面去做.妻子渴望丈夫早日信主,卻忽略了要尊重丈夫.
\newline
\newline
究竟怎樣纔是敬重丈夫呢?首先我們要明白,他雖然未信主,但主耶穌的死也是為了他.其次,這家庭雖是你的,也是他的.雖然他所要做的事,未得你同意；但,這家也是他的!他有自由做這樣做那樣.
\newline
\newline
或者有人會這樣說,他在家中做的是主所不喜悅的事,會影響兒女們呀!可是你與他作對,反會使兒女同情他,結果,不是更影響了兒女嗎?
\newline
\newline
所以我們應該向兒女解釋,他們的父親雖未信主,但很顧家,愛家,乃是一個好父親.我們必須和睦相處,使家庭美滿幸福,快樂!西諺有云:「與其整天在他面前背誦約翰福音三章十六節,不如做好些朱古力餅給他吃.」中國古諺亦云:「不如把粥煲好些給他吃.」
\newline
\newline
你們不要以外面的辮頭髮,戴金飾,穿美衣,為裝飾,只要以裡面存著長久溫柔安靜的心為妝飾,這在神面前極寶貴的.」(三至四節)彼得並非說不用整理頭髮,不可戴珠寶.乃是說,這些都不是最緊要的；有外面的妝飾,較不如有溫柔的內心.他的話正加強貞潔的品行的重要.
\newline
\newline
「因為古時仰賴神的聖潔婦人,正是以此為妝飾,順服自己的丈夫.就如撒拉聽從亞伯拉罕,稱他為主.你們若行善,不因恐嚇而害怕,便是撒拉的女兒了.」(五至六節)彼得以撒拉敬重亞伯拉罕為例,勗勉信主的姊妹們.
\newline
\newline
「你們作丈夫的,也要按情理和妻子同住,因她比你軟弱,與你一同承受生命之恩的.所以要敬重她,這樣,便叫你們的禱告沒有阻礙.」(七節)現在,彼得把話題轉到信主的夫妻家庭了,彼得稱他們是一同承受生命之恩的.
\newline
\newline
保羅在以弗所書第五章,也詳細論及,信主的夫婦應如何相敬如賓.「不要醉酒,酒能使人放蕩,乃要被聖靈充滿.當用詩章,頌詞,靈歌,彼此對說,口唱心和的讚美主.凡事要奉我們主耶穌基督的名,常常感謝父神.又常存敬畏基督的心,彼此順服.」(弗五18-21)保羅提到一個被聖靈充滿的人,有三方面的證據.被聖靈充滿的人,是一個喜樂的人,常常以詩歌讚美主；又是一個感恩的人,常常感謝父神；更是一個在生活上守紀律的人,常存敬畏基督的心,彼此順服.
\newline
\newline
在(弗五22至33節),保羅把喜樂,感恩和彼此順服,應用到丈夫與妻子方面；六章一至四節是應用在父母兒女身上；六章五至九節是應用在僱主與僱員,即東主與勞工方面.
\newline
\newline
第一點,當丈夫想起妻子時,是否愛慕,抑是憎恨?廿一節說要「彼此順服」,丈夫不應以妻子為自己囊中物,應該彼此尊重,正如我們尊敬主一樣.
\newline
\newline
(弗五章廿二至廿四節)是對妻子說的,廿五節以下是對丈夫說的.為甚麼對丈夫說的話,用了好幾節經文,可見在神眼前,丈夫應負起領導妻子,領導家庭之責.
\newline
\newline
「因為丈夫是妻子的頭,如同基督是教會的頭,他是教會全體的救主.教會怎樣順服基督,妻子也要怎樣凡事順服丈夫.」(弗五23至24節)廿二節的「順服」與廿四節的「順服」兩字常被世人濫用,可惜連神的子女們也濫用它.不少男子常有強盛的男性優越感,以為自己像皇帝一樣,強妻子做自己要做的事,要她聽話,接受自己的主意.但聖經的「順服」,是有順,有服,有跟隨之意,是彼此的順服.廿四節的「順服」,有將自己放在對方之下,支持他,幫助你,不是作他絆腳石之意.
\newline
\newline
神為亞當造一個配偶幫助他.今天的時代是一個競爭的世代,丈夫實在需要妻子關懷與幫助,然後家庭有幸福.
\newline
\newline
「你們作妻子的,當順服自己的丈夫,......」到此只是一個逗號,還不是句號,這句還未完全,直到「如同順服主」.才是句號.多奇妙的聖經真理,妻子有權利向丈夫表明人類最高超的愛,但作丈夫的也有權利看見神對人崇高的愛；丈夫對於妻子之愛,有如神對教會之愛.一個丈夫應該關懷到妻子在屬靈上的福樂,正如關懷家庭在經濟上的福利一樣.
\newline
\newline
在廿三節更清楚說:「他又是教會全體的救主.」意即「供應者」.有人虐待妻子,待她如牲畜,這與聖經的教訓是相反的.婚後妻子固然屬於丈夫正如頭與身體,基督與教會一樣,所以作丈夫的,應該愛他的妻子；正如基督愛祂的教會一樣.
\newline
\newline
我們要怎樣實行「順服」與「尊敬」呢?保羅在這方面有很清楚的說明:「你們作丈夫的要愛你們的妻子,正如基督愛教會,為教會捨己.」許多人沒有遵照聖經教訓而行,好像只把禮物給了妻子,而沒有將自己給妻子一樣.有一次,約翰來對我說:「祁牧師,我已將一幢美麗的洋房,雅緻的傢俬給瑪利,但她仍不覺滿足!」但當我和瑪利談話的時候,我問她:「約翰是不是已把洋房傢俬給你?」她說:「是的!不過......」她哇一聲哭起來了!我就轉過來對約翰說:「你會不會對著電視機親嘴?會不會對著這些傢俬親嘴?你將這些給她是不夠的,她所要的是你自己.」約翰這時才恍然大悟.各位,倘若主耶穌基督把世界都給了我們,但不把祂自己給我們,我們是多麼可憐!因我們需要的,不是這些,乃是祂自己.
\newline
\newline
「要用水藉著道,將教會洗淨,成為聖潔.可以獻給自己,作個榮耀的教會,毫無玷污皺紋等類的病,乃是聖潔,沒有瑕疵的.」(弗五26-27)我在探訪會友家庭時,常見丈夫批評妻子,在妻子身上找過失,甚至公開的批評.世人可這樣做,但信主的丈夫無理由這樣作.主不喜歡這樣,試想那個丈夫沒有瑕疵呢?那個丈夫是完全的呢?能有資格指責別人的,只有主自己.但主的心意,實不想指責我們,也不欲宣揚我的醜惡,乃是叫我們在世人面前,能夠無玷污,可以堵住他們的口.神給我一個妻子,就是神給我一件美妙的禮物.你的朋友不肯來教會聽道,但可能到訪你的家,他若看見你們夫妻相敬相愛,這就是你們一個最好的見證,也是一個最好傳福音的機會.要用家庭的見證向他傳信息,硬心的人也會受感歸主的.
\newline
\newline
「丈夫也當照樣愛妻子,如同愛自己的身子.愛妻子,便是愛自己了.」(弗五28)讀四福音書,看見神藉著耶穌基督,如何為教會受辱,受苦,受死,所以丈夫也應該愛妻子,如同基督愛教會,為教會捨己.
\newline
\newline
「保養顧惜」(29節)這四字有保護,供養之意.作丈夫的如果這樣對妻子,他在世人面前就是個好榜樣.
\newline
\newline
「保養顧惜」,也是丈夫在婚禮時的諾言,我們要緊記!
\newline
\newline
一個人被聖靈充滿的人,有甚麼證據呢?看他們夫妻如何相待,這就有力的見證了,由此纔可使世人受感信主!
\newline
\newline


\newpage

\section{第43屆~港九培靈會~祈柏~第五講}
\label{sec:1399}
\textbf{積極的聖潔生活}
\newline
\newline
連結: \href{https://www.hkbibleconference.org/session-message/view/1399}{\texttt{https://www.hkbibleconference.org/session-message/view/1399}}
\newline
\newline
\hyperref[sec:1398]{< < < PREV SERMON < < <}
~
\hyperref[sec:toc]{[返目錄]}
~
\hyperref[sec:1400]{> > > NEXT SERMON > > >}
\newline
\newline
彼得前書 3:8-3:20
\newline
\begin{longtable}{cl}
\hline
\hline
章節 & 經文 (和合本修訂版)\\
\hline
3:8 & \begin{tabularx}{0.7\textwidth}{X} 總而言之，你們都要同心，彼此體恤，相愛如弟兄，存憐憫和謙卑的心。 \end{tabularx} \\ \\ \relax
3:9 & \begin{tabularx}{0.7\textwidth}{X} 不要以惡報惡，以辱罵還辱罵，倒要祝福，因為你們正是為此蒙召的，好使你們承受福氣。 \end{tabularx} \\ \\ \relax
3:10 & \begin{tabularx}{0.7\textwidth}{X} 因為經上說：「凡要愛惜生命、享受好日子的人，要禁止舌頭不出惡言，嘴唇不說詭詐的話。 \end{tabularx} \\ \\ \relax
3:11 & \begin{tabularx}{0.7\textwidth}{X} 也要棄惡行善，尋求和睦，一心追求。 \end{tabularx} \\ \\ \relax
3:12 & \begin{tabularx}{0.7\textwidth}{X} 因為主的眼看顧義人，他的耳聽他們的祈禱；但主向行惡的人變臉。」 \end{tabularx} \\ \\ \relax
3:13 & \begin{tabularx}{0.7\textwidth}{X} 你們若熱心行善，有誰會害你們呢？ \end{tabularx} \\ \\ \relax
3:14 & \begin{tabularx}{0.7\textwidth}{X} 即使你們為義受苦，也是有福的。不要怕人的威嚇，也不要驚慌； \end{tabularx} \\ \\ \relax
3:15 & \begin{tabularx}{0.7\textwidth}{X} 只要心裡奉主基督為聖，尊他為主。有人問你們心中盼望的理由，要隨時準備答覆； \end{tabularx} \\ \\ \relax
3:16 & \begin{tabularx}{0.7\textwidth}{X} 不過，要以溫柔、敬畏的心回答。要存無虧的良心，使你們在何事上被毀謗，就在何事上使那些凌辱你們在基督裡有好品行的人自覺羞愧。 \end{tabularx} \\ \\ \relax
3:17 & \begin{tabularx}{0.7\textwidth}{X} 神的旨意若是要你們因行善受苦，這總比因行惡受苦好。 \end{tabularx} \\ \\ \relax
3:18 & \begin{tabularx}{0.7\textwidth}{X} 因為基督也曾一次為罪受苦 ，就是義的代替不義的，為要引領你們到神面前。在肉體裡，他被治死；但在靈裡，他復活了。 \end{tabularx} \\ \\ \relax
3:19 & \begin{tabularx}{0.7\textwidth}{X} 他藉這靈也曾去向那些在監獄裡的靈傳道， \end{tabularx} \\ \\ \relax
3:20 & \begin{tabularx}{0.7\textwidth}{X} 就是那些從前在挪亞預備方舟、神容忍等待的時候不信從的人。當時進入方舟，藉著水得救的不多，只有八個人。 \end{tabularx} \\ \\
[1ex]
\hline
\hline
\end{longtable}
彼得前書三章八節以下,是彼得提出為主而活,要在積極方面,跟隨聖靈,隨祂引導而行.
\newline
\newline
「總而言之,你們都要同心,彼此體恤,相愛如弟兄,存謙卑憐憫的心.」(八節)基督徒的態度,要在生活上表彰出來.這裡所指的「同心」,並非說甚麼事情都要相同,我們對事情的看法,雖然見仁見智,但仍須有合一的心.
\newline
\newline
其次是「體恤」,我們要設身處地的為別人著想,同情別人,體恤別人.但可惜我們做事,常不理會別人的反應如何,就一意孤行.在日常生活中,應該想到有人在你以上,也有人在你以下.在你以下的,並不是他們的才能不及你；可能他未有機會施展才能與已,所以要體恤他.不過人心傾向總是要居人之上的, 地位優越的人,更應該體恤別人,因為不論何人,主耶穌都為他們死了.從十字架的觀點來看我們都是罪人,主卻憐恤了我們!
\newline
\newline
英文──「禮貌」二字,是從拉丁文來的,有善良,和敬重之意.如果別人知道我是一位教授,他會敬重我；若他知道我是一個苦力,他可能不尊重我了.但主耶穌卻不是這樣,祂對待我們,不論是苦力,或是教授,都是一樣的愛護.所以我們要彼此敬重,以禮相待.
\newline
\newline
我不知道在香港的基督徒,有沒有爭競,和報復的事.在聖經的觀點來說,這樣的行為就是罪,是生活上不聖潔的明證.
\newline
\newline
第九節彼得說得更率直,「不以惡報惡,以辱罵還辱罵,倒要祝福；因你們是為此蒙召,好叫你們承受福氣.」(九節)我們是蒙神救贖的人,在世人面前,不可起爭競.務要使別人藉你得福,因為像我們這些原是不配得福的人,也蒙了神的恩惠.
\newline
\newline
詩卅四12:「人若愛長壽,享美福,須要禁止舌頭不出惡言,嘴唇不說詭詐的話.」請注意「惡言」和「詭詐的話」.我們是為主而活,抑為自己而活?若為主而活就要「離惡行善,尋求和睦,一心追趕.」(十二節),不可存報復之心,總要以善報惡.
\newline
\newline
各位,不論你是何人,在神眼中,你都是祂所愛,所寶貴的,你總要遵祂的話而行.
\newline
\newline
「因為主的眼看顧義人,主的耳聽他們的祈禱!惟有行惡的人,主向他們變臉.」(十二節).
\newline
\newline
「你們若是熱心行善,有誰害你們呢?」主眼看顧義人,聽他們的禱告.
\newline
\newline
十二節所提的「義人」,是正直的,不彎曲,不欺騙,誠實之意,誰是「義」與「不義」.神必鑑察!雖然污穢的人也可以裝成聖潔,不義的人可以裝假稱義；但神是無所不知的,祂識透你的內心.祂亦按我們內心所存的報應我們.惟有在神面前正直的人,不彎曲的人:才能蒙福.
\newline
\newline
「你們既是為義受苦,也是有福的,不要怕人的威嚇,也不要驚慌.」(十四節)為善的人,不一定事事順利.我們有時反會被人佔便宜惡待.不過,「為義受苦,也是有福的.」
\newline
\newline
「只要心裡尊主基督為聖,有人問你們心中盼望的緣由,就要常作準備,以溫柔敬畏的心,回答各人.」(十五節)這是我喜愛的經節,彼得所講的不是理論,乃是他痛苦的經歷.他曾三次否認主,終於被主的愛感動回轉過來愛主.
\newline
\newline
「要心裡尊主為聖」拜偶像的人,尚且在廟宇中供奉他們的假神.尊假神為聖,我們在日常的生活裡,有沒有尊主為聖呢?有沒有讓主在我們心中作主,作王呢?主是我生活的中心呢?我是否讓主得著呢?我們當常作準備,以便回答各人.所以務要如嬰孩,愛慕靈奶,明白道理,就能回答眾人,而且可以說出盼望的緣由.回答別人的態度,不要發怒,不要反擊,要溫柔安靜的回答,雖然他們不禮貌對待我們,我們也要以溫柔回答他們.
\newline
\newline
「存著無虧的良心,叫你們在何事上被毀謗,就在何事上,可以叫那誣賴你們在基督裡有好品行的人,自覺羞愧.」(十六節)試將本節與二章十二節比較,就知道是平衡的,「你們在外邦人中,應當品行端正,叫那些毀謗你們是作惡的,因看見你們的好行為,便在鑑察的日子,歸榮耀給神.」別人雖然惡待我們,或攻擊我們,其實他們早已知道我們所作的是對的.事情往往是這樣,我們愈對,他們愈憎恨.
\newline
\newline
「神的旨意若是叫你們因行善受苦,總強如因行惡受苦.」(十七節)彼得有痛苦的經歷,但他希望我們能避免,不再重蹈他的覆轍.在路加福音廿二章記載,彼得三次不認主之前,主耶穌早已告訴他了.主清楚的說:「今日雞還沒有叫兩次以前,你要三次說不認得我.」在美國有這樣一句話說,牧師若在主日吃雞,這雞也在他的事奉上有份.不錯,儆醒彼得的雞,實在在事奉主工作上有份了.可惜,今天的傳道人不如這隻雞,因為傳道者的聲音,不能叫人想起主的話,也不能叫人為罪痛哭憂傷!
\newline
\newline
在主釘十字架到復活之一段時間內,彼得的心情,是很悲慘的,他不知道主能否看見他悔罪之心!馬可十六章七節記載,當主復活的早晨,祂對婦女們說:「你們可以去告訴祂的門徒和彼得.......」可見主并沒有忘記彼得.日後彼得寫彼得前書,對主復活後的關懷,真是滿心感激,不能盡言.
\newline
\newline
「因基督也曾一次為罪受苦,就是義的代替不義的,為要引我們到神面前.按著肉體說祂被治死,按著靈性說,祂復活了.」(十八節)這節經文與二章廿一至廿三節比較一下,也是平衡的.
\newline
\newline
「你們蒙召原是為此,因基督也為你們受過苦,給你們留下榜樣,叫你們跟隨他的腳蹤行.」(二21)「腳蹤」是指甚麼?三章十八節指出的是主為我們的罪受苦,以義代替不義.他雖戴上荊棘的冠冕,但他復活了；使我們不但為他的十字架,也為他的空墳墓而歡欣誇勝.
\newline
\newline
甚麼叫恩典呢?就是我不配得,而神白白賜給我的就叫做恩典.我們既白白得恩,也應當白白施恩,好叫別人也像我們一樣的蒙恩!主既降卑,叫我們得恩,我們也當自卑,叫別人蒙恩.
\newline
\newline
各位,我們高唱為主而活的詩歌,何等容易；但真正肯為主而活的人就少有了!但願為我們死而復活的主在我們裡面,加給我們力量去實行榮耀祂的名!
\newline
\newline


\newpage

\section{第43屆~港九培靈會~祈柏~第六講}
\label{sec:1400}
\textbf{為信仰而活}
\newline
\newline
連結: \href{https://www.hkbibleconference.org/session-message/view/1400}{\texttt{https://www.hkbibleconference.org/session-message/view/1400}}
\newline
\newline
\hyperref[sec:1399]{< < < PREV SERMON < < <}
~
\hyperref[sec:toc]{[返目錄]}
~
\hyperref[sec:1401]{> > > NEXT SERMON > > >}
\newline
\newline
彼得前書 4:12-5:11
\newline
\begin{longtable}{cl}
\hline
\hline
章節 & 經文 (和合本修訂版)\\
\hline
4:12 & \begin{tabularx}{0.7\textwidth}{X} 親愛的，有火一般的考驗臨到你們，不要奇怪，似乎是遭遇非常的事； \end{tabularx} \\ \\ \relax
4:13 & \begin{tabularx}{0.7\textwidth}{X} 倒要歡喜，因為你們是與基督一同受苦，使你們在他榮耀顯現的時候也可以歡喜快樂。 \end{tabularx} \\ \\ \relax
4:14 & \begin{tabularx}{0.7\textwidth}{X} 你們若為基督的名受辱罵是有福的，因為榮耀的靈，就是神的靈，在你們身上。 \end{tabularx} \\ \\ \relax
4:15 & \begin{tabularx}{0.7\textwidth}{X} 你們中間，不可有人因為殺人、偷竊、作惡、好管閒事而受苦。 \end{tabularx} \\ \\ \relax
4:16 & \begin{tabularx}{0.7\textwidth}{X} 若有人因是基督徒而受苦，不要引以為恥，倒要因這名而歸榮耀給神。 \end{tabularx} \\ \\ \relax
4:17 & \begin{tabularx}{0.7\textwidth}{X} 因為時候到了，審判要從神的家開始；若是先從我們開始，那麼，不信從神福音的人將有何等的結局呢？ \end{tabularx} \\ \\ \relax
4:18 & \begin{tabularx}{0.7\textwidth}{X} 「若是義人還僅僅得救，不虔敬和犯罪的人將有何地可站呢？」 \end{tabularx} \\ \\ \relax
4:19 & \begin{tabularx}{0.7\textwidth}{X} 所以，照神旨意受苦的人要一心為善，將自己的靈魂交給那信實的造物主。 \end{tabularx} \\ \\ \relax
5:1 & \begin{tabularx}{0.7\textwidth}{X} 所以，我這同作長老，作基督受苦的證人和分享將來所要顯現的榮耀的人，勉勵在你們中間的長老們： \end{tabularx} \\ \\ \relax
5:2 & \begin{tabularx}{0.7\textwidth}{X} 務要牧養在你們當中神的群羊，按著神的旨意照顧他們，不是出於勉強，而是出於甘心；也不是因為貪財，而是出於樂意。 \end{tabularx} \\ \\ \relax
5:3 & \begin{tabularx}{0.7\textwidth}{X} 不要轄制所託付你們的群羊，而是要作他們的榜樣。 \end{tabularx} \\ \\ \relax
5:4 & \begin{tabularx}{0.7\textwidth}{X} 到了大牧人顯現的時候，你們必得到那永不衰殘、榮耀的冠冕。 \end{tabularx} \\ \\ \relax
5:5 & \begin{tabularx}{0.7\textwidth}{X} 同樣，你們年輕的，要順服年長的。你們大家都要以謙卑當衣服穿上，彼此順服，因為「神抵擋驕傲的人，但賜恩給謙卑的人。」 \end{tabularx} \\ \\ \relax
5:6 & \begin{tabularx}{0.7\textwidth}{X} 所以，你們要謙卑服在神大能的手下，這樣，到了適當的時候，他必使你們升高。 \end{tabularx} \\ \\ \relax
5:7 & \begin{tabularx}{0.7\textwidth}{X} 你們要將一切的憂慮卸給神，因為他顧念你們。 \end{tabularx} \\ \\ \relax
5:8 & \begin{tabularx}{0.7\textwidth}{X} 務要謹慎，要警醒。因為你們的仇敵魔鬼，如同咆哮的獅子，走來走去，尋找可吞吃的人。 \end{tabularx} \\ \\ \relax
5:9 & \begin{tabularx}{0.7\textwidth}{X} 你們要用堅固的信心抵擋他，因為知道你們在世上的眾弟兄也正在經歷這樣的苦難。 \end{tabularx} \\ \\ \relax
5:10 & \begin{tabularx}{0.7\textwidth}{X} 那賜一切恩典的神曾在基督裡召了你們，得享他永遠的榮耀，在你們暫受苦難之後，必要親自成全你們，堅固你們，賜力量給你們，建立你們。 \end{tabularx} \\ \\ \relax
5:11 & \begin{tabularx}{0.7\textwidth}{X} 願權能歸給他，直到永永遠遠。阿們！ \end{tabularx} \\ \\
[1ex]
\hline
\hline
\end{longtable}
(註:本題因有連屬關係,故六至八合為一講)
\newline
\newline
我將四章一至十一節省略不提,原因該段經文是彼得重覆上文所說的話；所以由四章十二節以下,是論到為信仰而力行.
\newline
\newline
(一)在受苦中要喜樂(四12-19)
\newline
\newline
1. 在受苦中不要驚訝(四12-13)「親愛的弟兄阿!有火煉的試驗臨到你們,不要以為奇怪.(似乎是遭遇非常的事)」許多人以為愛主不夠,所以要受苦.彼得卻指出有時受苦,是咎由自取的.但是基本上是因為作基督徒的生活見證,對世人來說,是一種責備,故會觸怒他們.正如今天在有些地區中,為主作見證而受苦,並不是希奇的事,「倒要歡喜,因為你們是與基督一同受苦,使你們在他榮耀顯現的時候,也可以歡喜快樂.」(十三節)聖經也說過:「他因那擺在前面的喜樂,就輕看羞辱,忍受了十字架的苦難......」(來十二2)照樣,我們忍受苦難是因為主耶穌已戰勝了死亡,為了主寶貴的應許我們要歡欣!
\newline
\newline
2. 以基督徒的身份確實去受苦(四14-16)若要找出基督徒受苦的原因為何?彼得指出:「你們若為基督的名受辱罵,便是有福的,因為神榮耀的靈,常住在你們身上.」(十四節)神不會忘記我們為祂的名忠心,有一天,要按祂的公義來審判世界.
\newline
\newline
前面說過,有時基督徒受苦,是咎由自取的,「你們中間,卻不可有人,因為殺人,偷竊,作惡,好管閒事而受苦.」(十五節)信徒中,相信不會有殺人,偷竊的事；但說閒話,好管閒事總是有的.我們若稱為主名下的人,就不可作這些事.
\newline
\newline
「若為作基督徒受苦,卻不要羞恥,倒要因這名歸榮耀給神.」(十六節)「歸榮耀是說到不論遇福遇苦,別人可以從我們的反應而更尊敬我們,這就是歸榮耀與神.
\newline
\newline
3. 審判要從神的家起首(四17-19)「因為時候到了,審判要從神的家起首,若是先從我們起首,那不信從神福音的人,將有何等的結局呢!」(十七節)世人對我們不公平的待遇,不要以為神不伸張公義,有一天,恩典的時候要結束,神公義的審判便開始.那時「若是義人僅僅得救,那不虔敬和犯罪的人,將有何地可站呢?」(十八節)
\newline
\newline
「所以那照神旨意受苦的人,要一心為善,將自己靈魂交與那信實的造化之主.」(十九節)這節經文正如二章廿三節描述主耶穌的遭遇,「祂被罵不還口,受害不說威嚇的話,只將自己交託那按公義審判人的主.」(二23)我們是主的門徒,要跟隨主的腳蹤走.
\newline
\newline
(二)在事奉中要忠心(五1-4)
\newline
\newline
本章是對教會的負責人,教會的領袖們(長老們)而說的.
\newline
\newline
1. 要作基督的見證(一節)「我這作長老,作基督受苦的見證,同享後來所要顯現之榮耀的,勸你們中間與我同作長老的人.」彼得以自己的經歷,勸告教會的領袖們.
\newline
\newline
2. 要作群羊的牧者(二至三節)「務要牧養在你們中間神的群羊,按著神旨意照管他們,不是出於勉強,乃是出於甘心,也不是因為貪財,乃是出於樂意.也不是轄制所託付你們的﹐乃是作群羊的榜樣.」這兩節經文我們當熟讀記在心中,我們蒙召原是為此,這是對我們這些事奉神的人意義深長的囑咐.務要牧養在你們中間神的群羊」.牧人的本份是要餵養群羊.我個人相信,如在這幾天,我用解經式的講道,逐節逐字解,最能使聽道的人得神的話餵養.作神聖工的人,要曉得到神的士多房裡,看看有些甚麼食物,可作餵養飼料.神的士多房就是聖經,所以神的工人要熟讀聖經,學習找尋靈糧.記得在二章二至三節,要愛慕靈奶,像才生的嬰兒一樣,以至漸長,成為剛強的基督徒.而且「有人問你們心中盼望的緣由,就要常作準備,以溫柔敬畏的心,回答各人.」(三15)若熟悉的話,自然有適當的回答.五章二節再次著重「餵養神的群羊」,保羅說:「當用各樣的智慧,把基督的道理,豐豐富富的存在心裡......」(西三16)在使徒行傳二十章記載保羅將離開以弗所教會,也勸勉當日教會的領袖們:「聖靈立你們作全群的監督,你們就將為自己謹慎,也為全群謹慎,牧養神的教會,就是也用自己血所買來的.」今天教會有一個通病,就是靈奶的講章太多,乾糧的講章太少了!不錯,奶是嬰兒最好的食物,但我們不能到六十歲還吃奶過活呀!奶對嬰兒是足夠營養的食物,但長大後便要吃乾糧了,基督徒也要如此.今天教會出了問題,是吃奶過多,吃乾糧過少的緣故.
\newline
\newline
「按神的旨意,照管他們,不是出於勉強,乃是出於甘心」.若有傳道人,以為牧養是一份工作,與職業無異,不清楚是神的呼召,這實在是悲劇.若我到禮拜六晚才坐下來想明天要講甚麼道,那麼我應要離開神的聖工.因為每一個傳道人,平日要勤讀聖經,才可按神的旨意,照管教會,分糧與群羊.
\newline
\newline
說到「也不是因為貪財之事,」大家都知道,傳道人斷不會有「過多」之薪金,沒有財「可貪」,所以在此不贅.
\newline
\newline
三節是說到「我們不應作獨裁者」來轄管他們,若我每天早上在此只用口講道,下了講台不行道.不必等到主再來,已顯露出來了!所以不論在學校裡,在工作場所,辦公廳,要使別人看出我在生活所傳之道來.各位!若傳道有問題,就是因為我們未曾為主而活.
\newline
\newline
3. 要配得面對主的顯現(四節)「到了牧長顯現的時候,你們必得那永不衰殘的榮耀的冠冕.」當我們定睛仰望主,知道主會再來,便叫我們有穩定,有堅持下去的力量.彼得,保羅和約翰都說過同樣的話:「凡向他有指望的,就潔淨自己,像他潔淨一樣.」
\newline
\newline
(三)不要憂慮(五5-11)
\newline
\newline
1. 萬事要信賴主(五5)「你們年幼的,也要順服年長的.就是你們眾人,也當要以謙卑束腰,彼此順服,因為神阻擋驕傲的人,賜恩給謙卑的人.」這裡雖然開頭說:「你們年幼的,也要順服年長的.」但跟著說「你們眾人,也當以謙卑束腰,彼此順服」.這是否有矛盾?但仔細閱讀便知道不會有矛盾,這是教導年長的和年幼的和睦相處.在今天教會的事工上,年長的和年幼的,都有很大的距離.年長的,高高的坐在高位上,多年來不肯改變；他們寧願教會死去,也不肯接受新主張.年幼的卻主張多,不但要改變；而且要快速的改變,但不知到底變甚麼!年長的和年幼的,彼此堅持,不肯協調,教會怎能沒有衝突?
\newline
\newline
其實,我們各有弱點,也各有長處.年長的,信主久,與主同行了不少時候,確是知道的多,經驗多,在教會來說,是寶貴的.但年幼的當他們歸主時,主的能力在他們身上何等的大,能為他們成就大事.若年長的以經驗,年幼的以熱誠,二者配合,教會事工,相得益彰,是何等大的福氣.彼得意謂,年幼的要順服年長的經驗,年長的要順服年幼的熱誠；二者配合,在主領導之下,教會便有好見證,好福氣.在美國常說兩代之間,有一段很長的隔閡.但在聖經的亮光,光照之下,兩代之間的問題,年長的和年幼的都是出於驕傲.年長的說,甚麼我都知道；年幼的說,甚麼我都做得到.驕傲的人,神要阻擋他們,所以教會裡沒有新人得救,更沒有人在主恩中長進.若彼此息爭,攜手合作,就會帶來極大的祝福.
\newline
\newline
有時年長的懷疑年幼的是否真做得到,年幼的懷疑年長的是否有經驗,就難於合作,彼得在此為可能發生的問題提供答案:「你們要將一切的憂慮卸給神,因為他顧念你們.」(七節)在此將兩代啣接在主的十字架上,主關懷年幼的,也關懷年長的.我們將憂慮卸給祂,祂就將二者連結.不論年齡的差別,觀念的差別,思想的差別,抑或環境的差別,在主裡是可以合力同心的來事奉祂.
\newline
\newline
主對待我們,不是根據我們的年齡差別,也不是根據我們的種族不同,祂引導年齡不同的人,引導不同種族的人親近祂.只要仰望主,男女老幼都可以合得來,一同事奉主.
\newline
\newline
2. 務要謹守,儆醒(五8-9)「務要謹守,儆醒,因為你們的仇敵魔鬼,如同吼叫的獅子,遍地游行,尋找可吞吃的人.」(八節)如對中國人說,吼叫的老虎,別人或許更易領會.「你們要用堅固的信心抵擋他,因為知道你們在世上的眾弟兄,也是經歷這樣的苦難.」(九節)在研經會中,我喜歡用專題來討論魔鬼.我們要抵擋牠,但先要知道牠是甚麼?撒旦的工作,不但要我們變壞,酗酒,吸毒,撒旦的意思是「分開」,「挑撥」.牠挑撥離間,牠要我變良善,良善到不需要救主.牠不是侷促一隅,而是遍地游行,獵取食物,所以我們要謹守,儆醒.在第九節,提出唯一抵擋牠的方法,是在信仰上能夠站立得穩.要在信仰上站立得穩,要先知道信仰是甚麼?所以要多研究聖經,作個在信仰上站立得穩的好門徒.
\newline
\newline
3. 要靠神的幫助以致完全(五10-11)「那賜諸般恩典的神,曾在基督裡召你們,得享他永遠的榮耀,等你們暫受苦難之後,必要親自成全你們,堅固你們,賜力量給你們.願權能歸給他,直到永永遠遠,阿們!」這兩節經文,是彼得留下給我們偉大的祝福,他鼓勵我們,告訴我們有一位賜諸般恩典的神.首先,祂要成全我們,好能應付各方面的事.其次,是說要堅固我們,在神的真道上有好根基.正如保羅在羅馬書對信徒們說:「總沒有因不信,心裡起疑惑,反倒因信,心裡得堅固.」(羅四20)第三,祂賜力量給我們,使我們剛強,能夠抵擋魔鬼,在這個邪惡的世代,為主作美好的見證.這裡,還包括另一件事,使我們有目標.我們原本是一隻無舵之舟,隨波逐流,不知往那裡去.彼得勉勵我們,要在神的恩典中長進.尊主為聖,若有人問我盼望的緣由,就常作準備,以溫柔敬畏的心回答各人.我們不但到禮拜堂聚集時,認識神的話,甚至每日靈修,也可認識神的話.有了穩固的根基,然後繼續長進.彼得從前有過痛苦的經歷,三次不認主,顛沛流離；如今他回首過去,便對我們發出忠誠的勸告和熱切的盼望.
\newline
\newline
附彼得生平簡略
\newline
\newline
若有人問我,當我離開世界,要在墓誌銘上寫甚麼,我一定是說約翰一章卅七節「兩個門徒,聽見他的話,就跟從耶穌.」但願我們也有此榮幸,有人聽見我們講道,就跟從了耶穌.
\newline
\newline
先請看卅五和卅六節:「再次日,約翰同兩個門徒站在那裡,他見耶穌行走,就說,看哪,這是神的羔羊.」這兩個門徒中,有一個是安得烈,他的兄弟就是彼得.「聽見約翰的話,跟從耶穌的那兩個人,一個是西門彼得的兄弟安得烈,他先找著自己的哥哥西門,對他說,我們遇見彌賽亞了.(彌賽亞繙出來就是基督)於是領他去見耶穌.耶穌看著他說,你是約翰的兒子西門,你要稱為磯法.(磯法繙出來,就是彼得.)」(約一40-42)在此我願提醒作聖工的同工們,安得烈遇見了耶穌,立刻便找著自己的哥哥,帶他去見耶穌,這是作個人見證的最好方法.
\newline
\newline
「耶穌看著他」(四十二節),是「看透」,「看穿」之意.主知道他的衝動,驕傲,「西門」就是這意思.但主耶穌仍然揀選西門做自己的門徒.主不看他以前是甚麼樣子,主是要看他今後的改變.今天是西門不要緊,只要明天成為彼得便行了.
\newline
\newline
再請看馬太福音十六章十三節:「耶穌到了該撒利亞腓立比的境內,就問門徒說:『人說我人子是誰?』他們說:『有人說是施洗的約翰,有人說是以利亞,又有人說是耶利米,或是先知裡的一位.』耶穌說:『你們說我是誰?』這時主耶穌在世上的工作,已到了末期,快上迦略山,那時,有許多人跟從祂,也有許多人反對祂.現在主耶穌有興趣要知道人到底以為他是誰?但當祂提出問話之後,其他門徒還未有機會回答,彼得卻搶先回答說:「你是基督,是永生神的兒子.」主立即稱讚他說:「西門巴約拿,你是有福的.因為這不是屬血肉的指示你的,乃是我天上的父指示的.我還告訴你,你是彼得,我要把我的教會建造在這磐石上,陰間的權柄,不能勝過他.」(太十六17-18)有人以為主耶穌要把教會建在彼得身上,在希臘文「磐石」與「彼得」(石頭)是有分別的.主耶穌所說,乃要將教會建在彼得所認識這磐石之上,這磐石是指主耶穌基督.彼得說過:「主乃活石,毋然是被人所棄的.」(彼前二4)保羅也說過:「有基督自己為房角石.」(弗二20)
\newline
\newline
我們若相信彼得是磐石,那麼在馬太福音十六章廿三節,這磐石便要被粉碎了.「從此耶穌才指示門徒,他必須上耶路撒冷去,受長老祭司長文士許多的苦,並且被殺,第三日復活.彼得就拉著他,勸他說:『主阿!萬不可如此,這事必不臨到你身上.』」(21至22節)「他」原文有「責備」之意,你想石頭可以責備教會的磐石嗎?「耶穌轉過身來,對彼得說:『撒旦,退我後邊去罷!你是絆我腳的,因為你不體貼神的意思,只體貼人的意思.』(廿三節)主責備彼得是撒旦,當然撒旦不會是磐石.如果以為彼得是教會的頭,就不是順從神的權柄,仍順服撒旦的權柄.不錯,彼得當時好意護衛主,但他是想用人力去護衛,不是用神的能力來護衛.
\newline
\newline
現在請翻到路加福音廿二章,在卅二節裡有一個祈禱,可是這祈禱不蒙神應允.盼望我們若為朋友祈禱,不要得此結果.
\newline
\newline
這一段經文是彼得痛苦的經歷的高潮.主又說:『西門,西門,撒旦想要得著你,好篩你們,像篩麥子一樣.但我已經為你祈求,叫你不至於失去信心,你回頭以後,要堅固你的弟兄.』彼得說:『主阿!我就是同你下監,同你受死,也是甘心.』耶穌說:『彼得,我告訴你,今日雞還沒有叫,你要三次說不認得我.』」(路廿二31-34)有許多人,起先自恃,驕傲,硬頸,要等他自己的剛硬被粉碎了,才能蒙恩得救.彼得被粉碎了,他學習了一課,他憂傷,他痛苦.
\newline
\newline
彼得為何有此軟弱?他遠遠的跟著耶穌,還和大祭司的下人,一同在火堆前取煖.「他們拿住耶穌,把他帶到大祭司的宅裡,彼得遠遠的跟著.他們在院子裡生了火,一同坐著.彼得也坐在他們中間,有一個使女,看見彼得坐在火光裡,就定睛看他.說:「這個人也是同那人一夥的.」彼得卻不承認說:「女子!我不認得他.」過了不多時候,又有一個人看見他,說:「你也是他們一黨的.彼得說:「你這個人,我不是.」約過了一小時,又有一個人極力的說:「他實在是同那人一夥的,因為他也是加利利人.彼得說:「你這個人,我不曉得你說的是甚麼.正說話之間,雞就叫了.主轉過身來,看彼得,彼得便想起主對他所說的話,「今日雞叫以先,你要三次不認我.他就出去痛哭.」(路廿二54-62)當主的話提醒我們,若肯聽從,是可等的福氣,若不聽從,痛苦是何等的大.
\newline
\newline
最後,要看約翰福音廿一章,記得主復活的清晨,天使對婦女說:「你們可以去告訴祂的門徒和彼得」(可十六7),主特別關懷他.在約翰福音廿一章記載,當時有幾個門徒去打魚,彼得是他們的首領.當他們捨舟登岸,覺得很驚奇!主耶穌為他們生了火,又預備了早餐.根據希臘文的意思,主耶穌說:「孩子們,你們來吃.」這是何等親切的呼喚.他們吃完了早飯,一段很有意義的對話發生了.「耶穌對西門彼得說,約翰的兒子西門:『你愛我比這些更深麼?』彼得說:『主阿!是的,你知道我你(直譯,喜歡你).』耶穌對他說:『你餧養我的小羊.』耶穌第二次又對他說:『約翰的兒子西門,你愛我麼?』彼得說:『主阿!是的,你知道我愛你(直譯,我喜歡你.』耶穌說:『你牧養我的羊.』第三次對他說:『約翰的兒子西門,你愛我麼?(直譯,你只是喜歡我麼?)』彼得因為耶穌第三次對他說,你愛我麼(你只是喜歡我麼?)就憂愁.對耶穌說:『主阿!你是無所不知的,你知道我愛你(直譯,你知道我只喜歡你).』耶穌說:『你餧養我的羊.』
\newline
\newline
彼得回答耶穌說,主阿!我喜歡你.彼得喜歡耶耶穌,也喜歡自己,更喜歡許多東西.但主耶穌是要求彼得,不只喜歡耶穌,是要求他愛耶穌,為祂犧牲一切.當我們讀使徒行傳時,很喜歡見到彼得,後來確實不只喜歡耶穌,而且深愛耶穌,將一切奉獻給耶穌.
\newline
\newline
這次培靈會幾位講員,有共同願望,在培靈會開始的時候,你不過只喜歡耶穌,但培靈會結束時,盼望你成為一個全心全意愛主的人.
\newline
\newline


\newpage

\section{第43屆~港九培靈會~鮑會園~第一講}
\label{sec:1401}
\textbf{復興的預備──以斯拉記全貌}
\newline
\newline
連結: \href{https://www.hkbibleconference.org/session-message/view/1401}{\texttt{https://www.hkbibleconference.org/session-message/view/1401}}
\newline
\newline
\hyperref[sec:1400]{< < < PREV SERMON < < <}
~
\hyperref[sec:toc]{[返目錄]}
~
\hyperref[sec:1402]{> > > NEXT SERMON > > >}
\newline
\newline
以斯拉記,尼希米記和以斯帖記,差不多是同時代的作品,分成兩個主題.以斯拉,尼希米,記載從被擄之地歸回的百姓所作的工；以斯帖則記載仍留在巴比倫的以色列人所作的工.我們從以斯拉記中,知道百姓回去重建耶路撒冷的聖殿,從尼希米記又知道,他們還要建造城牆.以斯拉記和尼希米記這兩本書,在次序上也相當有意義.尼希米記是記載以色列人在外邦人中,重新恢復他們的見證；以斯拉是記載百姓回國後,重建聖殿的工作.聖殿是代表敬拜神,重靈裡與神的關係.以斯拉記在先,尼希米記在後,在他們實行真正有效的事奉前,必須先在靈裡恢復他們與神的關係.基督徒的事奉,不是一個工作﹐乃是一個生命的流露.若我們單憑表面努力去作,恐怕這工作只不過是個人努力所產生的效果,因為我們的生命若不是在神前被建造,一切的工作就成了枉然的勞苦.
\newline
\newline
研究本書前,讓我們先查考它的歷史背景.
\newline
\newline
‧主前六百多年,神的懲罰臨到以色列的百姓,藉外邦人把他們擄到巴比倫地.這只不過是神短暫的懲罰,因祂曾應許百姓,只要他們轉回專心尋求祂的面,祂必再施恩與他們.
\newline
\newline
‧七十年後,神的時候到了,神感動波斯王古列的心,下令准許百姓返回耶路撒冷,恢復他們敬拜神的生活.結果有五萬人在王的支援下,神的賜福中,願意返回耶路撒冷,決心重建聖殿.
\newline
\newline
‧過了兩三年,約主前五百三十年,百姓在建殿時受到當地許多人的攻擊.因為當人何時在神前真正追求時,便會遭遇到難處,所以當聖殿的根基剛肂好,便被迫停工了.百姓正心裡火熱努力工作,突然受到攔阻,盼望粉碎,心裡的力量也失掉!從哈該書中看見百姓建殿時,在靈裡過了一個黑暗的生活,聖殿未建成便為自己建天花板的房屋.
\newline
\newline
‧五,六年後,約主前五百一十四年,神差遣哈該及撒迦利亞兩位先知,返回耶路撒冷；用神的話教導百姓,結果他們便再次努力建殿.
\newline
\newline
‧五,六年後的主前五百一十年,聖殿建造完了,他們有了敬拜神的地方,百姓便以為達到目的；不再追求,因而對神的心又冷淡了.
\newline
\newline
‧又過了幾十年,約主前四百五十八年,神又興起祂的僕人以斯拉,從巴比倫地被差遣回去,以斯拉回到耶路撒冷,距聖殿被建已有六十多年之久,他主要的工作是用神的話教導百姓.以斯拉是一個特別受訓的文士,他回去用真理教導百姓,讓他們有真正的復興,以斯拉記第九章,十章記載這些事情.
\newline
\newline
‧再過幾十年,約主前四百四十五年,百姓在靈裡得到復興後；神的時候又到了,神再差遣尼希米及一些人回去,重建耶路撒冷的城牆.
\newline
\newline
耶路撒冷城牆得以被建,在百姓復興前,神當時做了許多預備的工作.很多時候預備的工作,往往比實際做出來的更重要.但因預備的工作不能看果效,沒有顯著的表現,所以不肯在預備工作上下工夫!現在讓我們看看神在各方面的預備:
\newline
\newline
(一)預備一個肯禱告的器皿──但以理
\newline
\newline
「波斯王古列元年,耶和華為要應驗藉耶利米口所說的話,就激動波斯王古列的心,使他下詔通告全國說.波斯王古列如此說:耶和華天上的神已將天下萬國賜給我,又囑咐我在猶大的耶路撒冷為他建造殿宇.」
\newline
\newline
這段經文是以斯拉記的開始(拉一1-2)也是歷代志下的結束(代下三十六22-23),兩本書記載同一件的事情,中間雖然相隔數十年.但神的計劃仍繼續不斷演變下去.實際上這兩節有關聖殿重建的經文,遠在但以理書中已有記載,(但九章)當但以理讀先知耶利米書時,看見神警告及應許百姓話.他們被擄七十年期滿,神要再施恩與百姓.但以理計算知道七十年的時間已期滿,現今該是神施恩之時了.他得到這個信息後,整個心受感動被信息緊扼著；而覺得非常沉重,以玫無法再做其他的事情,只有把自己一生放神.他知道現今國家處於困危的命運中,是因為自己和百姓的罪,他便代表百姓晝夜伏在神前認罪禱告.
\newline
\newline
禱告永遠是復興工作的開始,我們不能不為教會的需要禱告,何時看見神感動人真願意花時間禱告,這就是離復興之時不遠了.約翰衛斯理是英國教會的名人,有一次當他傳道時,有一個婦人對他說,她已為英國復興禱告了二十年,現在因看見神藉衛斯理所作的工,便知道神已應允了她禱告.這只是一個婦人的見證,我們不知道還有多少人已為英國的復興而禱告,以致做成了如此奇妙的復興.
\newline
\newline
聞說神復興了美國亞斯比尼大學.在復興前該大學有許多紛爭,情況有陷於完全分裂的狀態.幸而有四位同學,看見這情形心中非常難過,流淚在神面前禱告.在想不到時,神在那裡做了奇妙的工作；因著學校的復興,許多教會也得到復興.這是因為有人肯在神面前付代價,今天教會也極需要,這些肯為教會付上禱告代價的人.
\newline
\newline
(二)一個肯謙卑的器皿──設巴薩
\newline
\newline
以斯拉記第五章,記載有一個省長名叫設巴薩.他不是一知名人物,從第一章裡我們知道他是猶大的首領.此「首領」原文應譯作「王子」,換言之,他是王族,大衛直系的後裔.當百姓回耶路撒冷建殿時,波斯王把聖殿事奉的器皿交在設巴薩的手,回到耶路撒冷後便立他作當地的省長.由此我們可知他的工作是何等的重要,他有特別的身份,地位,資格和負起領導百姓復興之重任.但我們聖經的記載裡,知道他在人前並沒有了不起顯著的成就,也沒有得到人的稱讚,竟然一個重任便落在他身上,這正因為他是王的後裔,有真正王的身份,是神特別的百姓,有屬靈的生命,神就是要藉著這樣的人,去成就他的工作.今天的教會也需要這些,不是追求在人間作大事的人,來完成他的工作,正如神對巴錄說:「你要圖謀大事嗎?你不要圖謀.」神所使用的人,乃在乎他願否把為自己圖謀的雄心放下,交在神的手中成為合用的器皿.今天教會裡只注重人的成就,和宏大的表現；在教會許多事工上,都學了屬世的方法.但,我們應追求的,不是「大」,也不是單靠宣傳和組織,這不能使教會得著祝福,反而成為攔阻.正如設巴薩,他雖不是一個人材,但他肯謙卑放在神手中被神使用,神便把事奉神聖殿的器皿,放在他手裡,又把帶領百姓的責任交在他身上.
\newline
\newline
(三)一個不可少的工具──神的話
\newline
\newline
神要百姓在外邦地恢復他們的見證,返國重建城牆以前,先差遣以斯拉用真理教導百姓.如果沒有以斯拉的教導,我們真不敢想像百姓的光景如何了!也不敢盼望尼希米能帶領他們作如此的見證,更談不上能把耶路撒冷城牆重建起來.照樣我們若盼望神在我們中間作工,我們該讓神的話在心中動工,讓我們的生命建造在神的話語上.在神的話語上下工夫並不容易,不單是靈裡受了感動,心中有新的亮光,乃是真正明白神的話.可惜教會今天對於神的話有兩方面的極端:
\newline
\newline
(1)過份注重工作忽略靈裡的建造
\newline
\newline
有些人把一生的重點放在為主「忙」的事上,這是一個不可彌補的損失.他們以為越忙越好,結果忙到一個地步,甚麼事都辦不好.過份的忙碌便會產生虛空,注重工作,而忽略了靈裡的真實!如此容易使整個教會的工作,重心改在福利,醫療,教育和其他工作上,而忽略了以真理造就弟兄姊妹.
\newline
\newline
(2)追求屬靈經歷忽略真理工作
\newline
\newline
有些人覺得需要更多親近主,認為用時間與主親近是絕不可少的,但這追求若不建立在真理上,便很容易誤入岐途.在教會歷史上有許多追求屬靈的經驗,單追求靈裡的感動,而沒有讓感動建在真理上,以致所走的路有所偏差.在主後的第三,第四世紀,教會有一特別的名詞,有人被稱為「旗桿上的聖徒」,這些聖徒要追求屬靈境界,便產生避世的思想,為自己做了一枝高旗桿,在其上建造房子過著聖潔的生活.他們以為旗桿越高,越表明更與神親近,當時許多有價值的生命,便在旗桿上浪費了.我們追求真理需要建造在神的根基上.
\newline
\newline
在尼希米帶領百姓重建城牆前,神先帶領以斯拉用真理教導百姓,惟有用真理教導百姓才能使他們得到復興.甚麼時候我們遵行神的道,我們便走上復興的路上.
\newline
\newline


\newpage

\section{第43屆~港九培靈會~鮑會園~第二講}
\label{sec:1402}
\textbf{復興的啟示}
\newline
\newline
連結: \href{https://www.hkbibleconference.org/session-message/view/1402}{\texttt{https://www.hkbibleconference.org/session-message/view/1402}}
\newline
\newline
\hyperref[sec:1401]{< < < PREV SERMON < < <}
~
\hyperref[sec:toc]{[返目錄]}
~
\hyperref[sec:1403]{> > > NEXT SERMON > > >}
\newline
\newline
以斯拉記 1:1-1:4
\newline
\begin{longtable}{cl}
\hline
\hline
章節 & 經文 (和合本修訂版)\\
\hline
1:1 & \begin{tabularx}{0.7\textwidth}{X} 波斯王居魯士元年，耶和華為要應驗藉耶利米的口所說的話，就激發波斯王居魯士的心，使他下詔書通告全國，說： \end{tabularx} \\ \\ \relax
1:2 & \begin{tabularx}{0.7\textwidth}{X} 「波斯王居魯士如此說：耶和華天上的神已將地上萬國賜給我，又委派我在猶大的耶路撒冷為他建造殿宇。 \end{tabularx} \\ \\ \relax
1:3 & \begin{tabularx}{0.7\textwidth}{X} 你們中間凡作他子民的，可以上猶大的耶路撒冷去，重建耶和華－以色列神的殿，他是在耶路撒冷的神；願神與這人同在。 \end{tabularx} \\ \\ \relax
1:4 & \begin{tabularx}{0.7\textwidth}{X} 凡存留的人，無論寄居何處，那地的人要用金銀、財物、牲畜幫助他，還要為耶路撒冷神的殿甘心獻上禮物。」 \end{tabularx} \\ \\
[1ex]
\hline
\hline
\end{longtable}
以色列人在以斯拉領導下,經驗了一個真正的復興,但這個復興不是從以斯拉的頭腦裡想出來的,因為神從來沒有叫我們這些屬祂的人,為祂開始作甚麼工作.神的工作早已計劃好了,我們的責任是要順服神的帶領,若我們只用自己的思想,計劃來做事,恐怕做出來的,不過是我們的工作而已.
\newline
\newline
神按著祂的時候將祂的旨意啟示我們,多少時候我們不能完成神旨努力工作,可能由於我們不能分辨出神的旨意.所以今天讓我們思想一個關乎「啟示」的問題,藉著神將復興的啟示,賜給以色列百姓的事,更能明白神的旨意.
\newline
\newline
(一)啟示的對象
\newline
\newline
聖經告訴我們:「耶和華為要應驗藉耶利米口所說的話.」(拉一1)因此我們明白神整個的計劃並不是突然而來的.祂早已藉著僕人先知啟示了.神要作的工作永遠不會向神的僕人隱藏的.創世記十八章裡,神曾說:「我們要所作的事豈可瞞著亞伯拉罕呢?」,亞伯拉罕與神同行,與神的交通非常親密,以致神並沒有將毀滅所多瑪,蛾摩拉兩城的事向亞伯拉罕隱藏.在歷史上我們也曾自見神的旨意,照樣地向祂的僕人耶利米,但以理,及以後很多先知,使徒們顯明出來.
\newline
\newline
(二)領受啟示的條件
\newline
\newline
昔日神把啟示賜與眾先知們,今天神也能將祂的旨意啟示我們.只要我們具備有他們的條件,正如先知耶利米領受神的啟示時一樣.我們從耶利米書中,看見他幾個領受神旨意的條件.
\newline
\newline
(1)不顧一切順服神旨
\newline
\newline
耶利米所處的時代,是一個混亂,戰爭,極不容易處的時代,更是百姓落在一個迷失方向需要選擇的時代.一個這樣的時代,要領受神的旨意又遵行出來實不容易.耶利米之所以能夠領受神的旨意,因為他願意不顧一切,按著神的旨意去行；而且又願意為著遵行神的旨意,付上極大的代價.他曾被放在地牢裡,監獄中,最後因著遵行神旨又被擄到外邦地去；還被自己的同胞所迫害,終於為著遵行神的旨意而喪失了生命.耶利米就是這樣肯不顧一切,遵行神旨,置生死於度外的人.
\newline
\newline
(2)有屬靈智慧辨別真偽
\newline
\newline
耶利米有一個分辨是非的心,又有屬靈智慧.當時有許多假先知興起與他為敵,假如有五十個假先知與耶利米的意見不一,一個人的意見與五十個人的意見比較,雖然那些人佔多數,但耶利米仍然敢站在神的話語上,因他知道這是神的旨意.在混亂的時代裡,能夠單單按著神的話分辨那個是是神的旨意,實不容易,今天我們這個時代,照樣充滿了人的聲音,若我們要領受神的啟示,也必須在靈裡有分辨的智慧.按著神的話辨別真偽,如此雖有困難和危險在前,仍然可以照著神的話去遵行.因為耶利米有這樣的心志,神的啟示便臨到他.
\newline
\newline
神的啟示不但臨到耶利米,也激動了波斯王古列的心,使古列領受神的旨意.一個順服神的人,一個有預備心的人,他能成為一個特別的器血.但可惜許多時候神的旨意臨到人身上,有些人不肯順服,甚至屬神的人也不肯按神旨而行.雖然如此,但神的旨意也不會落空的,因為君王的心都是在神手中(箴廿一1),詩篇七十六篇十節告訴我們:「人的忿怒要成全神的榮美.」人對神旨意的不順服只會使自己受虧損,神仍然要照樣成就.因此我們看見神的旨意不單臨到那些順服他的人,也臨到那些不肯順服的人.因為神的旨意,要按祂所定的計劃臨到世人.
\newline
\newline
(三)啟示的內容
\newline
\newline
(1)認識耶和華是全地的主宰
\newline
\newline
古列是一個外邦的君王,不是敬畏神的人,但他領受這啟示時；不得不認耶和華是掌管一切啟示的主,一切都在神的手裡.古列對百姓說,神將天上地上一切賜給他,天上萬國賜給他,神有權柄賜與,因為他是掌管一切的主宰.神啟示的目的,是要叫世人知道祂有權柄作任何事.
\newline
\newline
(2)認識耶和華是啟示的目標
\newline
\newline
古列得到啟示後,他就讓百姓回耶路撒冷建造神的殿,這個啟示不單顯明耶和華是掌權天地萬物的主,而且表明一切事情的發生,都是為耶和華的緣故,耶和華是一切事物的目標.神賜下啟示,感動人作工,都是為使神的名彰顯.今天神在我們中間作工,將啟示感動我們,也是為著這一個使神的名彰顯的目的,這樣我們可以藉此測驗所領受的啟示是否從神而來.若這啟示的目標,不是單彰顯主的名,不是高舉主耶穌,那麼這啟示不可能從神而來.
\newline
\newline
(四)啟示的方法
\newline
\newline
在這混亂的時代裡,我們要領受神的啟示,必須要懂得如何分辨,這裡列舉幾種神啟示不同的方法:
\newline
\newline
(1) 藉著神自己的話
\newline
\newline
神的啟示永遠是藉著祂口中的話,就是聖經的真理而來,我們所領受的啟示不能與聖經分開,神要藉著聖經把啟示賜給我們,因此追求神啟示的範圍,也不能離開聖經以外.若所領受的真理與聖經不相符合,這便不是屬於神的.有很多人在這事上犯了錯誤,解釋聖經與上下文毫無關係,隨便把某處的意思解釋便以此為啟示,建立一些錯謬的教訓.例如:摩門教引證以賽亞書四章「在那日七個女人必拉住一個男人說,我們歸到你的名下.」他們就以此證明人可以多妻,這樣的解釋是完全錯誤的.今天有許多人解釋聖經只是按個人的喜好,一時的衝動,隨便引出一些錯誤的道理.為此,我們必須要花時間查聖經,按聖經的教訓,上下文的關係,聖經作者的目的而解釋,這樣才能真正明白聖經的教訓.
\newline
\newline
(2) 藉人心裡的感動
\newline
\newline
神將這啟示賜給百姓是藉耶利米的話,但神啟示賜給古列時是感動他的心.我們也可以藉著心裡的感動領受神的啟示,耶穌也曾說過,聖靈來為要領導我們進入真理(約十六13),但有不少的人,甚至是信仰純正的基督徒,只尋求心中新的感動,啟示的亮光,便以此為真理.我們該懂得按聖經的教訓和標準去衡量它,辨別這啟示是否從神而來,以下有個建議:
\newline
\newline
(a)「你們中間若有先知,或是作夢的起來向你顯個神蹟奇事,對你說我們去隨從你素來所不認識的別神,事奉他罷.他所顯的神蹟奇事,雖有應驗,你也不可聽那先知,或是那作夢之人的話......」(申十三1)
\newline
\newline
對於任何的啟示和教訓,我們都應該先發問:這教訓是使我親近神抑或親近人?更愛神或是更愛某一種理論立場.因此有人雖能行神蹟,但只要他的教訓使我們遠離神,這個啟示永遠不可能從神而來.
\newline
\newline
(b)「......若有人傳福音給你們,與你們所領受的不同,他就應當被咒詛.」(加一9)
\newline
\newline
加拉太人已經領受了真理,而且憑著福音得救了.若有另一種新的啟示福音傳來,便要測驗這啟示與從前領受的是否相同,倘若不同這便不是出於神.
\newline
\newline
(c)「......真理的聖靈來了......他要榮耀我,因為他將受於我的,告訴你們.」(約十六14)
\newline
\newline
耶穌表明聖靈的工作是要榮耀他,為祂作見證.「榮耀」的意思是指將原來的好處彰顯出來.任何的啟示,都是為要將主的榮美彰顯出來.聖靈將啟示賜給我們,目的要我們在新啟示中更看見主的榮美!所以若有任何啟示高舉聖靈過於主耶穌的,這啟示便有問題.舉例來說:講台上的燈光是很重要,但它的工作只不過為要將講台上的人顯出來,而不是使人欣賞它,否則,便失去它存在的目的.
\newline
\newline
聖靈來為要為主作見證,彰顯主的榮美,帶我們進入真理,這是祂工作的目的.
\newline
\newline
以上三處經文是衡量啟示真實性的標準.
\newline
\newline
這樣聖靈還有甚麼新的真理啟示呢?既已符合聖經的教訓,與從前領受的一樣；為可我們還需要領受聖靈的啟示呢?從聖經裡我們可以看到一個例子:當耶穌復活以後曾向以馬忤斯路上的兩個門徒顯現,他們邊走邊談論,門徒為著耶穌被釘死在十字架上心裡難過.耶穌為使他們明白自己,已應驗聖經的話從死裡復活,便與他們讀聖經,從摩西律法開始將凡與自己有關的事解釋清楚.讀的時候他們心裡火熱,後來便恍然大悟.主耶穌與門徒講論的並沒有超越聖經的範圍,可是從其中便領受了新的啟示.今日聖靈也要在人心中動工,藉著心中的感動,將啟示賜給我們.
\newline
\newline
(3) 藉事實的環境
\newline
\newline
古列得到神的啟示向百姓宣告,百姓便遵著啟示而行.神的啟示感動臨到,必在人心及世界上發生作用,我們要用這作用的果效,定這啟示的真確性.神的啟示臨到曾否進入心中?是否在我們生命中結出果子?因為神的話進入我們的生命,必在我們生命中有所行動,讓神的話在其中結出真實的果子來.
\newline
\newline


\newpage

\section{第43屆~港九培靈會~鮑會園~第三講}
\label{sec:1403}
\textbf{復興的呼召}
\newline
\newline
連結: \href{https://www.hkbibleconference.org/session-message/view/1403}{\texttt{https://www.hkbibleconference.org/session-message/view/1403}}
\newline
\newline
\hyperref[sec:1402]{< < < PREV SERMON < < <}
~
\hyperref[sec:toc]{[返目錄]}
~
\hyperref[sec:1404]{> > > NEXT SERMON > > >}
\newline
\newline
以斯拉記 1:1-1:11
\newline
\begin{longtable}{cl}
\hline
\hline
章節 & 經文 (和合本修訂版)\\
\hline
1:1 & \begin{tabularx}{0.7\textwidth}{X} 波斯王居魯士元年，耶和華為要應驗藉耶利米的口所說的話，就激發波斯王居魯士的心，使他下詔書通告全國，說： \end{tabularx} \\ \\ \relax
1:2 & \begin{tabularx}{0.7\textwidth}{X} 「波斯王居魯士如此說：耶和華天上的神已將地上萬國賜給我，又委派我在猶大的耶路撒冷為他建造殿宇。 \end{tabularx} \\ \\ \relax
1:3 & \begin{tabularx}{0.7\textwidth}{X} 你們中間凡作他子民的，可以上猶大的耶路撒冷去，重建耶和華－以色列神的殿，他是在耶路撒冷的神；願神與這人同在。 \end{tabularx} \\ \\ \relax
1:4 & \begin{tabularx}{0.7\textwidth}{X} 凡存留的人，無論寄居何處，那地的人要用金銀、財物、牲畜幫助他，還要為耶路撒冷神的殿甘心獻上禮物。」 \end{tabularx} \\ \\ \relax
1:5 & \begin{tabularx}{0.7\textwidth}{X} 於是，猶大和便雅憫的族長、祭司、利未人，凡是心被神感動的人都起來，要上耶路撒冷去建造耶和華的殿。 \end{tabularx} \\ \\ \relax
1:6 & \begin{tabularx}{0.7\textwidth}{X} 四圍所有的人都拿銀器、金子、財物、牲畜、珍寶支持他們，此外還有甘心獻的一切禮物。 \end{tabularx} \\ \\ \relax
1:7 & \begin{tabularx}{0.7\textwidth}{X} 居魯士王也把耶和華殿的器皿拿出來，這些器皿是尼布甲尼撒從耶路撒冷掠取，放在自己神明廟中的。 \end{tabularx} \\ \\ \relax
1:8 & \begin{tabularx}{0.7\textwidth}{X} 波斯王居魯士派米提利達司庫把這些器皿拿出來，點交給猶大的領袖設巴薩。 \end{tabularx} \\ \\ \relax
1:9 & \begin{tabularx}{0.7\textwidth}{X} 它們的數目如下：金盤三十個，銀盤一千個，刀二十九把， \end{tabularx} \\ \\ \relax
1:10 & \begin{tabularx}{0.7\textwidth}{X} 金碗三十個，備用銀碗四百一十個，其他器皿一千件。 \end{tabularx} \\ \\ \relax
1:11 & \begin{tabularx}{0.7\textwidth}{X} 金銀器皿共有五千四百件。被擄的人從巴比倫上耶路撒冷的時候，設巴薩把這一切都帶了上來。 \end{tabularx} \\ \\
[1ex]
\hline
\hline
\end{longtable}
神藉波斯王古列將復興的呼召,臨到以色列百姓,吩咐他們返回耶路撒冷,重建聖殿.被擄前百姓可以隨意自由地到聖殿敬拜事奉神,那時聖殿剛建好,所羅門為百姓祝福；使他們享受與神親近的滋味,心中充滿大喜樂.他們知道聖殿是敬拜神的地方,喜樂時可以讚美,困苦時可蒙安慰,悔改時可得赦免.但因百姓犯罪後受刑罰,被擄到巴比倫七十年,百姓不能敬拜事奉神.如果他們從沒有敬拜神的經驗,心裡便不會太難過；但他們因為曾嘗過親近神的滋味,有過親近神的經驗；知道愛神及享受過神的愛,如今卻在外邦地失去與神的交通,是十分痛苦的.如果我們曾得過一些寶貴的屬靈經驗,但如今卻失去喜樂,和與神的交通,以致禱告讀經不起勁,犯了罪卻失去了悔改的感覺的話.但願我們就能藉著今天的信息,重新建造心靈裡與神親近的殿.
\newline
\newline
(一)呼召的反應
\newline
\newline
神從不含糊,清楚地將復興的啟示放百姓面前,但百姓對這呼召,產生了兩種不同的反應:
\newline
\newline
一部份百姓聽了呼召,得到啟示,明白神的心意,思想過,衡量過,考慮過,卻不肯按照神旨回耶路撒冷去.
\newline
\newline
另一部份的百姓,在得到了同一呼召的時候,也經過想,衡量,考慮就決意回耶路撒冷去.
\newline
\newline
(二)不順服的原因
\newline
\newline
根據歷史記載被擄到巴比倫的百姓,有四十萬至六十萬之間的數目；但回到耶路撒冷去的至多只有五萬人,那表示有大部份的人不肯回去.其中的原因,聖經沒有明說,我們不敢評論.但猶太歷史家約瑟弗在猶太古史中曾說過:「當時在巴比倫的百姓,因不肯離棄在巴比倫的家業,而不肯回耶路撒冷去.」又記載說:「當時在巴比倫的百姓,因懼怕耶路撒冷嚴緊的生活而不肯回去.」這兩句話非常重要,顯露出當時留在巴比倫的百姓,正因這緣故不肯回耶路撒冷.
\newline
\newline
(1)不肯離棄外邦的家業
\newline
\newline
從以斯帖記可見留在巴比倫的猶太人,在被擄期間,努力為自己建立事業,使生活過得很安定.他們享受自己勞力的果效,所以神的呼喚臨到,他們也不肯回去,因為覺得付出的代價太大了.要回去,必須將家庭,朋友,產業,工作都放棄!單獨回到荒涼的耶路撒冷,離開安定的生活,而再重新創業,他們為了付出的代價太大,便不肯順服神的呼召.一個經過考慮而不肯順服的人,不是為了不知道要付上代價,而是為了怕付上代價,這是大悲劇!可憐許多基督徒也落在這境況中,不能按神的旨意而行.
\newline
\newline
(2)不肯過嚴緊的生活
\newline
\newline
約瑟弗告訢訴我們,因為猶太人怕嚴緊的生活,所以不肯回去耶路撒冷.可見猶太人的生活是嚴緊的必須按著神旨而行.可能某些事情別人可以作,他們卻不能；有些事別人可以不去作,而他們卻要作,他們的生活是應該與別人不同的.他們不能隨便吃喝,更不可犯罪,要持守神的律法,遵神旨而行.巴比倫人的生活卻相反,他們可以愛宴樂,任意妄行,放縱私慾.猶太人在巴比倫住了七十年,見慣了巴比倫人的生活,便覺得回去守嚴緊的生活實在太苦了.他們經過考慮之後,就決定不肯回去.過聖潔的生活,要討神喜悅的生活,一個蒙福與神有交通生活永遠是一個嚴緊的生活.很多人可能為了這原因,寧可過自由不受約束的生活,而畏懼走順服的道路.但有一件事我們需要注意的,就是不順服神的人,他們心裡永遠不平安；不服神的心,把我們一切的平安都奪去.
\newline
\newline
這些人選擇了自己歡喜的路,心中沒有平安；但他們又不甘心過著不平安的生活,於是用自的方法來填補心中的不平安.聖經告訴我們,這些不肯回去的猶太人拿金銀來幫助那些回去的人；從字裡行間,看出他們希望做一些事情,滿足神的心.於是捐些錢來幫助他們回去的路費,及建殿的所需.他們希望用這方法,去換取失掉的平安.但是這些身外屬物質的東西,絕不能代替個人內心與神的關係,也不能彌補在神面前的損失!
\newline
\newline
今天也有很多基督徒過著這樣的生活.一年前我曾在美一所華僑教會中講道,有一弟兄前來與我交通.經介紹後得知他本是與我同時蒙召,進入神學院受造就的,那時他很熱心很有恩賜.畢業後曾在教會事奉了一年,隨即往外國去,做了十多年生意,很有成就,是當地華僑中有名望的人,他談話中常提及在某某地方捐錢建立教會,某些工作也是由奉獻支持的.我聽了不禁使我驚訝,懷疑他是否為了這個原因,不聽從神的呼召；心裡沒有平安,不肯順服神走奉獻的道路,使想用金錢來彌補不順服神的內心痛苦.多奉獻錢財,多參加聚會這些原是好的,但若指望藉著這些來彌補心中的不順服,不肯對付的罪,擬將功贖罪的,這不是聖經的教訓.那些不肯順服神的人,想藉此法來賄賂自己的良心,不受責備,不至痛苦,如此下去良心就漸漸麻木了.
\newline
\newline
(三)順服的原因
\newline
\newline
那班肯真正順服復興呼召的百姓,勇敢地回去,是因為有幾個重要的原因:
\newline
\newline
(1)領袖領導
\newline
\newline
原來他們中間,有一班肯順服的領袖,猶大和便雅憫的族長,他們受神的激動.這些被神呼召的人,利未人,分別為聖的人,他們受神的感動,帶領百姓回去敬拜事奉.今天的教會,也正需要這樣的領袖們起來,在教會中起領導的作用.無論我們是義工,或是任何工作的負責人；也不可忽略我們的責任,在屬靈工作上負起領導的責任.耶穌說:「學生不能高過先生」,在屬靈的事上更是如此.如果領袖的人自己在屬靈的追求上沒有表現,生命中沒有對付,那麼便成了整個教會弟兄姊妹的阻攔,教會不能在屬靈上有長進.
\newline
\newline
猶太的首領設巴薩有一個特別的表現,王將五千多件聖殿器皿交在他的手,命他經管.一個重大的責任落在設巴薩的身上,他沒有推卻,樂意擔承.他沒有因為器皿的貴重而逃避責任.今天教會也極需要領袖們把擔子擔上,肯面對現實,肯負起責任,不容許教會中有罪惡的存在.
\newline
\newline
(2)渴慕神家
\newline
\newline
這班人被神激動以後,都要起來上耶路撒冷去.這可以代表他們心中有渴望回去的願望,他們看見耶路撒冷建殿的需要及價值.因著渴望的心,而甘心走上這道路.神也將渴慕的心放在我們裡面,使我們知道應走的,是甚麼道路.雖然知道有困難,但因著渴慕與神有交通,渴慕得復興,渴慕離開現今的地方,而遵行神復興的呼召.耶穌在登山寶訓裡說:「饑渴慕義的人有福了,因為他們必得飽足.」若果我們靈裡的追求如飢餓一般,我們必定到心靈上的滿足.
\newline
\newline
(3)獲得異象
\newline
\newline
百姓之所以不顧一切回到耶路撒冷,因為他們得到一個重新恢復敬拜神的異象.這個異象有如保羅所見馬其頓的異象,為這異象他不得不離開特羅亞往馬其頓去.百姓所領受的異象是具體的事實,在他心裡發生作用,感動,反應,這異象也成了他心中的負擔.於是不怕損失,不怕付代價,千辛萬苦回去,重新享受與神甜蜜的生活.
\newline
\newline
求主也將這個復興的啟示,異象放在我們心中.使我們不顧一切攔阻,重新建造自己生活的聖殿,恢復與神完全的交通.
\newline
\newline


\newpage

\section{第43屆~港九培靈會~鮑會園~第四講}
\label{sec:1404}
\textbf{順服的喜樂}
\newline
\newline
連結: \href{https://www.hkbibleconference.org/session-message/view/1404}{\texttt{https://www.hkbibleconference.org/session-message/view/1404}}
\newline
\newline
\hyperref[sec:1403]{< < < PREV SERMON < < <}
~
\hyperref[sec:toc]{[返目錄]}
~
\hyperref[sec:1405]{> > > NEXT SERMON > > >}
\newline
\newline
以斯拉記 2:68-3:13
\newline
\begin{longtable}{cl}
\hline
\hline
章節 & 經文 (和合本修訂版)\\
\hline
2:68 & \begin{tabularx}{0.7\textwidth}{X} 有些族長到了耶路撒冷耶和華的殿，為神的殿甘心獻上禮物，要在原有的根基上重新建造。 \end{tabularx} \\ \\ \relax
2:69 & \begin{tabularx}{0.7\textwidth}{X} 他們量力捐入工程的庫房，有六萬一千達利克金子，五千彌那銀子，以及一百件祭司的禮服。 \end{tabularx} \\ \\ \relax
2:70 & \begin{tabularx}{0.7\textwidth}{X} 於是祭司、利未人、百姓中的一些人、歌唱的、門口的守衛、殿役，各住在自己的城裡；以色列眾人都住在自己的城裡。 \end{tabularx} \\ \\ \relax
3:1 & \begin{tabularx}{0.7\textwidth}{X} 到了七月，以色列人住在自己的城裡；那時他們如同一人，聚集在耶路撒冷。 \end{tabularx} \\ \\ \relax
3:2 & \begin{tabularx}{0.7\textwidth}{X} 約薩達的兒子耶書亞和他的弟兄眾祭司，以及撒拉鐵的兒子所羅巴伯和他的弟兄，都起來建築以色列神的壇，要照神人摩西律法書上所寫的，在壇上獻燔祭。 \end{tabularx} \\ \\ \relax
3:3 & \begin{tabularx}{0.7\textwidth}{X} 他們在原有的根基上築壇，因為他們懼怕鄰邦民族，又在其上向耶和華早晚獻燔祭， \end{tabularx} \\ \\ \relax
3:4 & \begin{tabularx}{0.7\textwidth}{X} 並照律法書上所寫的守住棚節，按數照例每日獻所當獻的燔祭。 \end{tabularx} \\ \\ \relax
3:5 & \begin{tabularx}{0.7\textwidth}{X} 此後，他們獻常獻的燔祭，並在初一和耶和華一切分別為聖的節期獻祭，又向耶和華獻各人的甘心祭。 \end{tabularx} \\ \\ \relax
3:6 & \begin{tabularx}{0.7\textwidth}{X} 從七月初一起，雖然耶和華殿的根基尚未立定，他們開始向耶和華獻燔祭。 \end{tabularx} \\ \\ \relax
3:7 & \begin{tabularx}{0.7\textwidth}{X} 他們把銀子給石匠、木匠，把糧食、酒、油給西頓人、推羅人，好將香柏樹從黎巴嫩浮海運到約帕，是照波斯王居魯士所允准他們的。 \end{tabularx} \\ \\ \relax
3:8 & \begin{tabularx}{0.7\textwidth}{X} 他們到了耶路撒冷神殿的第二年，二月的時候，撒拉鐵的兒子所羅巴伯，約薩達的兒子耶書亞和其餘的弟兄，就是祭司和利未人，以及所有被擄歸回耶路撒冷的人，就開工建造；他們派二十歲以上的利未人，監督建造耶和華殿的工作。 \end{tabularx} \\ \\ \relax
3:9 & \begin{tabularx}{0.7\textwidth}{X} 於是何達威雅的後裔，就是耶書亞和他的子孫與弟兄、甲篾和他的子孫，他們和利未人希拿達的子孫與弟兄，都起來如同一人，監督那些在神殿裡做工的人。 \end{tabularx} \\ \\ \relax
3:10 & \begin{tabularx}{0.7\textwidth}{X} 工匠立耶和華殿根基的時候，祭司穿禮服吹號，利未人亞薩的子孫敲鈸，都照以色列王大衛親手所定的，站著讚美耶和華。 \end{tabularx} \\ \\ \relax
3:11 & \begin{tabularx}{0.7\textwidth}{X} 他們彼此唱和，讚美稱謝耶和華：「他本為善，他向以色列永施慈愛。」他們讚美耶和華的時候，眾百姓大聲呼喊，因為耶和華殿的根基已經立定。 \end{tabularx} \\ \\ \relax
3:12 & \begin{tabularx}{0.7\textwidth}{X} 然而有許多祭司、利未人和族長，就是見過先前那殿的老年人，現在親眼看見這殿立了根基，就大聲哭號，也有許多人大聲歡呼， \end{tabularx} \\ \\ \relax
3:13 & \begin{tabularx}{0.7\textwidth}{X} 百姓不能分辨歡呼的聲音或哭號的聲音，因為百姓大聲呼喊，聲音連遠處都可聽到。 \end{tabularx} \\ \\
[1ex]
\hline
\hline
\end{longtable}
以斯拉記第二章是篇很難讀的經文,全是人名數字,但細心閱讀時便看出一個重要的真理.百姓之所以能回到耶路撒冷,重新建造,再次為神作見證；這不是單屬於所羅巴伯,或耶書亞,省長設巴薩某一個人的功勞,乃是全體百姓合作起來的工作.所有屬神的人一齊被神感動,蒙神呼召,順服為神工作；也許他們中間有許多是不知名的人物,但因著同一的事奉便把工作做完了.
\newline
\newline
「以色列人彼此唱和讚美稱謝耶和華.」當百姓歸回時,他們也想不到迅速間能在神前發出讚美的聲音,何時只要人肯遵神旨意而行,他的生命便不能不發出讚美.
\newline
\newline
(一)讚美的因由
\newline
\newline
(1)如同一人的事奉
\newline
\newline
在第三章裡有三處經文原文裡,用不同的文字表達同一個意思.「如同一人」(一節),「他們都一同起來為神工作.」(九節),「眾民都大聲呼喊」(十節),因著他們同心合意,便發出讚美,他們同心合意走向耶路撒冷.在敬拜的事上,有同一心志,目的,敬拜與事奉的過程中沒有凡心鬥角.如此的敬拜使人心裡充滿快樂,便同心向神發出讚美.
\newline
\newline
(2)同心建造聖殿的根基
\newline
\newline
聖殿的根基代表與神的交通,重建聖殿便是恢復與神的交通.這交通是建立在同一的根基上,這就是主耶穌基督(林前三),這名以外再沒有根基,我們也不能靠任何人得救.聖殿的根基建好了,代表人心中有耶穌的地立,因為有這根基,人才能在其上建造敬拜神的生活.為著整個根基,主耶穌已建好了而讚美神的名.
\newline
\newline
(3)祭壇建好
\newline
\newline
他們建造以色列的祭壇,在壇上獻燔祭,這燔祭被焚燒,香氣達到神前,神的心得到滿足.因為神悅納百姓的敬拜事奉,為此百姓向神獻上讚美.我們向神的讚美不該是追求自己的生活滿足,乃要追求滿足神的心,這就是我們追求的目標.若以別的事為目標,只要稍遇困難攔阻便失去快樂.耶穌曾說:「我留下平安給你們.」(約十四)又將「喜樂賜給你們.」(約十五),當耶穌說這些話時,是在最後晚餐後.按人來說那時他要面對死亡,還有甚麼喜樂平安可以賜給我們呢?但因著他所行的是神的旨意,所以任何環境卻不能叫他失去喜樂,他整個喜樂平安,是建造在追求神旨意之上.百姓獻上燔祭,雖然殿未建成,仍有許多困難在他們面前；但他們仍高聲讚美神,因為他們的喜樂是建在神所悅納的事上.
\newline
\newline
(二)讚美的性質
\newline
\newline
「......他們讚美耶和華的時候,眾民大聲呼喊,因耶和華殿的根基已經立定；然而有許多祭司利未人,族長,就是見過舊殿的老年人；現在親眼看見立這殿的根基,便大聲哭號,也有許多人大聲歡呼.」(拉三11-12)
\newline
\newline
這喜樂不是完全的,乃是混雜的,因為那些沒有見過舊聖殿的人,他們在外邦地過了七十年的生活,曾聽聞敬拜神殿的榮耀,他們便產生渴慕之心；如今眼看殿的根基已立好了,可以敬拜神,因此便大聲呼喊.但年老的一輩,曾親眼見過所羅門時代所建的聖殿,在相比之下,覺得這殿太微小,太窮乏了!他們便大聲哭叫,換言之,他們有一個思想,就是因聖殿外表不夠美觀而心中憂傷.
\newline
\newline
有一件事是百姓未能領會的,神曾藉著先知哈該預言說:「萬軍之耶和華如此說,過不多時我必再一次震動天地,滄海,與旱地,我必震動萬國,萬國的珍寶必都運來.(意即萬國所羨慕的必來到),我就使這殿充滿了榮耀......這殿後來的榮耀,必大過先前的榮耀.......」神曾親自應許過,要賜福與後來的殿,殿的外型雖小,但主耶穌自己要進到裡面,屬靈的真實進到殿中,使其中充滿榮耀.
\newline
\newline
部份的百姓只看外面的微小,只追求外表的表現,因外表不夠美觀,便大聲哭號,心裡難過!他們忽略了靈裡的真實,而追求人看得見的表現；只求殿外表的偉大而忽略了神居住在裡面的寶貴,這是多麼可憐的現象.只要偉大的外殼,而失去最重要的實質.另一部份的人,不是回憶過去聖殿外表的偉大,他們只知道一件事,就是現在可以敬拜耶和華,因為已經有一個敬拜神的根基,在裡面可以親近神,便心滿意足了,毫不介意外表之美觀與否?所以他們雖然只有一所微小的聖殿仍能大聲讚美,他們已得到了屬靈的真實.
\newline
\newline
(三)讚美的重要性
\newline
\newline
(1)這是信心的見證
\newline
\newline
這時雖然聖殿仍未建妥,但因著親眼看見殿的根基,認為這根基已是足夠的保證,知道神的應許必然成就；所以當根基建好後,馬上便發出讚美,此乃信心的讚美.
\newline
\newline
「信心」是什麼?不信的人如此嘲笑著說:「信心是半夜裡在暗室中,找一隻黑貓.」意思是表明信心是一個不實際的東西,但聖經告訴我們信心是需要有證據的.使在(約二十36,31)表明,耶穌一生行了許多神蹟,但沒有完全記在福音書內；只把某些記載下來,為要使看見的人,信耶穌是基督是神的兒子；並且因信他的名,而得生命.以神蹟本身而言,耶穌並沒有利用它,直接地說我就是神的兒子.但當這些神蹟合在一起時,我們就曉得耶穌不能不是神的兒子了.這些證據與耶穌是神的兒子,有密切的關係.我們若憑著信心接受,信耶穌是基督,因這些證據而發出信心來,憑信心接受見證的價值.
\newline
\newline
照樣百姓看見聖殿的根基建好了,雖然這事表面上與神沒有什麼關係；但因著殿的根基被建,按照神的應許得以應驗.他們更知道神是信實慈愛的一位,因著這見證更認識神自己,這就是信心.憑著自己經驗過的事,更能信得過神,憑著自己所明白的真理,更相信神的話的可靠.以色列人就是因著這一個聖殿根基,被建的證據而信神,又讚美神.
\newline
\newline
(2)這是順服的見證
\newline
\newline
祭壇被建立,百姓便讚美耶和華.祭壇是順服的見證,過去百姓七十年在巴比倫過活,沒有祭壇的設立,也沒有機會獻祭.在巴比倫外邦地只能用口敬拜耶和華,因為沒有敬拜耶和華的壇,所以生活上便沒有特別的表現.「建立祭壇」這是敬拜耶和華的記號,這是生活上,動作上的記號；表示敬拜神的人,按神的方式生活.我們追求敬拜神,不單憑口來過復興敬拜的生活,乃是要在生活上敬拜記號；顯出我們是真基督徒,生活與別人不一樣.
\newline
\newline


\newpage

\section{第43屆~港九培靈會~鮑會園~第五講}
\label{sec:1405}
\textbf{竭力爭戰}
\newline
\newline
連結: \href{https://www.hkbibleconference.org/session-message/view/1405}{\texttt{https://www.hkbibleconference.org/session-message/view/1405}}
\newline
\newline
\hyperref[sec:1404]{< < < PREV SERMON < < <}
~
\hyperref[sec:toc]{[返目錄]}
~
\hyperref[sec:1406]{> > > NEXT SERMON > > >}
\newline
\newline
以斯拉記 4:1-5:2
\newline
\begin{longtable}{cl}
\hline
\hline
章節 & 經文 (和合本修訂版)\\
\hline
4:1 & \begin{tabularx}{0.7\textwidth}{X} 猶大和便雅憫的敵人聽說被擄歸回的人為耶和華－以色列的神建造殿宇， \end{tabularx} \\ \\ \relax
4:2 & \begin{tabularx}{0.7\textwidth}{X} 就去見所羅巴伯和族長，對他們說：「請讓我們與你們一同建造，因為我們也與你們一樣尋求你們的神。自從亞述王以撒‧哈頓帶我們上這地的日子以來，我們常向神獻祭。」 \end{tabularx} \\ \\ \relax
4:3 & \begin{tabularx}{0.7\textwidth}{X} 但所羅巴伯、耶書亞和其餘以色列的族長對他們說：「我們建造神的殿與你們無關，因為我們要照波斯王居魯士所吩咐的，自己為耶和華－以色列的神協力建造。」 \end{tabularx} \\ \\ \relax
4:4 & \begin{tabularx}{0.7\textwidth}{X} 那地的人就在猶大百姓建造的時候，使他們的手發軟，擾亂他們。 \end{tabularx} \\ \\ \relax
4:5 & \begin{tabularx}{0.7\textwidth}{X} 從波斯王居魯士年間，直到波斯王大流士在位的時候，那些人賄賂謀士，要破壞他們的計劃。 \end{tabularx} \\ \\ \relax
4:6 & \begin{tabularx}{0.7\textwidth}{X} 亞哈隨魯在位，他的國度剛開始的時候，他們上書控告猶大和耶路撒冷的居民。 \end{tabularx} \\ \\ \relax
4:7 & \begin{tabularx}{0.7\textwidth}{X} 亞達薛西年間，比施蘭、米特利達、他別和他們的同僚上書奏告波斯王亞達薛西。奏文是用亞蘭文寫的，以亞蘭文呈上。 \end{tabularx} \\ \\ \relax
4:8 & \begin{tabularx}{0.7\textwidth}{X} 利宏省長、伸帥書記也上奏亞達薛西王，控告耶路撒冷如下 \end{tabularx} \\ \\ \relax
4:9 & \begin{tabularx}{0.7\textwidth}{X} （那時，利宏省長、伸帥書記和他們其餘的同僚，法官、官員、軍官、波斯官員、亞基衛人、巴比倫人，和書珊迦人，就是以攔人， \end{tabularx} \\ \\ \relax
4:10 & \begin{tabularx}{0.7\textwidth}{X} 以及被亞斯那巴大人遷移、安置在撒瑪利亞城和大河西邊一帶地方其餘的人。現在 ， \end{tabularx} \\ \\ \relax
4:11 & \begin{tabularx}{0.7\textwidth}{X} 這是他們上奏亞達薛西王奏文的抄本）：「河西的臣僕上奏亞達薛西王，現在 \end{tabularx} \\ \\ \relax
4:12 & \begin{tabularx}{0.7\textwidth}{X} 請王知道，從王那裡上到我們這裡的猶太人，已經抵達耶路撒冷。他們正在重建這反叛惡劣的城，已經完成了城牆，正要修復根基。 \end{tabularx} \\ \\ \relax
4:13 & \begin{tabularx}{0.7\textwidth}{X} 如今請王知道，這城若再建造，城牆完工，他們就不再進貢、納糧、繳稅，王的國庫必受虧損。 \end{tabularx} \\ \\ \relax
4:14 & \begin{tabularx}{0.7\textwidth}{X} 如今，我們吃的鹽既然全是宮廷的鹽，就不忍見王吃虧，因此奏告於王， \end{tabularx} \\ \\ \relax
4:15 & \begin{tabularx}{0.7\textwidth}{X} 請王考察先王史籍，必會在史籍上查知這城是反叛的城，對列王和各省有害；自古以來，城中常有悖逆的事，因此這城曾被拆毀。 \end{tabularx} \\ \\ \relax
4:16 & \begin{tabularx}{0.7\textwidth}{X} 我們謹奏王知，這城若再建造，城牆完工，河西之地王就無份了。」 \end{tabularx} \\ \\ \relax
4:17 & \begin{tabularx}{0.7\textwidth}{X} 那時王諭覆利宏省長、伸帥書記和他們其餘的同僚，就是住撒瑪利亞和河西一帶地方的人，說：「願你們平安。現在 \end{tabularx} \\ \\ \relax
4:18 & \begin{tabularx}{0.7\textwidth}{X} 你們所呈給我們的奏本，已經清楚地在我面前讀了。 \end{tabularx} \\ \\ \relax
4:19 & \begin{tabularx}{0.7\textwidth}{X} 我已下令考查，得知這城自古以來果然背叛列王，其中常有反叛悖逆的事。 \end{tabularx} \\ \\ \relax
4:20 & \begin{tabularx}{0.7\textwidth}{X} 也曾有強大的君王治理耶路撒冷，統管河西全地，人就給他們進貢、納糧、繳稅。 \end{tabularx} \\ \\ \relax
4:21 & \begin{tabularx}{0.7\textwidth}{X} 現在你們要下令叫這些人停工，使這城不得建造，等到我再降旨。 \end{tabularx} \\ \\ \relax
4:22 & \begin{tabularx}{0.7\textwidth}{X} 你們當謹慎辦這事，不可遲延，何必讓損害加重，使王受虧損呢？」 \end{tabularx} \\ \\ \relax
4:23 & \begin{tabularx}{0.7\textwidth}{X} 亞達薛西王上諭的抄本在利宏和伸帥書記，以及他們的同僚面前宣讀，他們就急忙往耶路撒冷去見猶太人，用勢力和強權叫他們停工。 \end{tabularx} \\ \\ \relax
4:24 & \begin{tabularx}{0.7\textwidth}{X} 於是，在耶路撒冷神殿的工程就停止了，直停到波斯王大流士第二年。 \end{tabularx} \\ \\ \relax
5:1 & \begin{tabularx}{0.7\textwidth}{X} 那時，哈該先知和易多的孫子撒迦利亞，兩個先知奉以色列神的名向猶大和耶路撒冷的猶太人說預言。 \end{tabularx} \\ \\ \relax
5:2 & \begin{tabularx}{0.7\textwidth}{X} 於是撒拉鐵的兒子所羅巴伯和約薩達的兒子耶書亞起來，開始建造耶路撒冷神的殿，有神的先知在那裡幫助他們。 \end{tabularx} \\ \\
[1ex]
\hline
\hline
\end{longtable}
今天神在教會中尋找的人,就是那些肯不顧一切尋求神旨意而又肯遵行的人.人有許多的計劃和方法,但人終於受方法計劃的影響.因此眼光便偏差了,漸漸跟隨世俗.以為藉教育,政治,可以拯救國家；但真正能拯救國家的,乃是按神的話,又肯遵行神的道的人.神所揀選的一班百姓,他們不顧一切,回到耶路撒冷,祗按神旨作神工.當他們為自己建造天花板房屋時,他們不會遇到攻擊困難；但當他們遵行神的話作工時,困難便來到.撒旦不肯讓人毫無攔阻去完成神的工作,牠一定多方面攔阻攻擊；所以當他們建好聖殿的根基,預備蓋上蓋時,便被迫停工.在屬靈工作上永不會無戰事的,我們必須要認識戰爭中的敵人是誰?他們所用的方法又如何?最後百姓如何在戰爭中仍然站立得往?
\newline
\newline
(一)敵人是怎麼樣的?
\newline
\newline
以色列人開始建殿時,有撒瑪利亞人對他們說:「請容我們與你們一同建造,因為我們尋求你們的神,與你們一樣.......」(拉四2)敵人自稱與他們有共同敬拜,目標與工作都相同,彼此間毋須再分別了.但撒瑪利亞人是怎樣的?在(王下十七24)記載他們的背景,當以色列人被擄到亞述時,亞述王從巴比倫,古他,亞瓦,哈馬,和西法瓦音遷移人來安置在撒瑪利亞的城邑.這些人不是猶太人,也不是以色列人,是亞述王從別的地方把他們遷移來的.他們並不敬畏耶和華,所以神刑罰他們,差遣野獸獅子來傷害他們.他們看見野獸便害怕,因為知道這刑罰,是從耶和華來的!亞述王為此就差了一個祭司回去,好指教他們敬拜神的儀式.但這些撒瑪利亞人卻不接受真理,也無心信從真理,只憑外表建壇敬拜耶和華.
\newline
\newline
在(王下十七30)記載巴比倫人從那地方找到他們的神像,古他人照他們的神像,哈馬人照他們的神像；各人並不忘記自己的神,仍然做他們的神像.因此這些民又懼怕耶和華,又事奉他們的偶像；實際上心中敬拜的,乃是自己的偶像.按現今的話說,這些人既不是基督徒,又不是完全反對基督教；若他完全反對,我們還可以提防.但他們敬拜耶穌的名字,心裡卻事奉自己的偶像.他們口裡標榜好聽的話語,實際上卻走反對耶穌的路.這是敵人的真面目,我們應該小心防備.
\newline
\newline
(二)敵人的方法如何?
\newline
\newline
(1)妥協法
\newline
\newline
「他們說:請容我們與你們一同建造,因為我們尋求你們的神,與你們一樣.......」(拉四2)他們所用的方法非常可怕,先來到以色列人面前,叫百姓與他們妥協,一同事奉.曾有人說,信耶穌是可以的,但不要信得太怪誕,太極端,太保守了.這是可怕的事,倘若與仇敵妥協,就失去了見證.他們一面敬畏耶和華,一面卻又事奉自己的偶像,多少基督徒也像他們,在這地位上失敗了.也許妥協一點可把問題暫時解決.但須知道,凡與世俗為友的,便是與神為敵.我們需要有堅定的立場,因為妥協便是與世俗為友了.若與世俗為友便失去了為主作見證的力量.敵人祗希望我們能遷就,放下見證；但是唯有極端的見證,才能顯出我們的立場來.
\newline
\newline
這時,百姓在所羅巴伯,耶書亞和其他族長領導下,表示堅決的態度:「......我們建造神的殿,與你們無干,我們自己為耶和華以色列的神協力建造.」他們沒有說什麼,只認清楚自己的立場,既與別人無關,便繼續按神旨而行.
\newline
\newline
(2)賄賂法
\newline
\newline
敵人賄賂謀士,敗壞他們的計劃,控告耶路撒冷猶大的居民,對王上本說:「王該知道,從王那裡上到我們這裡的猶太人,已經到耶路撒冷；重建這反叛惡劣的城,築起根基,建造城牆.......」(拉四12)
\newline
\newline
仇敵知道必須用一個有效的方法,來攔阻他們,因此便捏造謊話控告百姓.百姓返回耶路撒冷,還沒有建造城牆,但他們卻向王誣告百姓建城牆；反叛國王,百姓便不得已要停工了.當時百姓若因自己的過失受攻擊,還容易忍受；但仇敵的擊只是謊話,便難於忍受了.正如主耶穌曾在八福篇裡說:「......逼迫你們,捏造各樣壞話毀謗你們,你們就有福了.......」(太五11),在我們的生活中,很容易在這事上失去見證；但我們不能因敵人用謊話控告,便降低地位同樣用謊話來對付他.波斯王受仇敵欺騙,強迫他們停工；他們受此攻擊,卻沒有用自己的方法抵擋,只把這事交與神.結果,後來神又打發先知撒迦利亞和哈該,勸勉他們繼續工作,殿便建成了.
\newline
\newline
(三)如何對付敵人?
\newline
\newline
百姓怎樣對付敵人?在這事上顯出他們靈裡所站的地位如何?他們無故被害,但沒有失去該有的見證.當敵人攻擊百姓時,百姓被迫停工.人往往因著工作沒有果效,便開始懷疑到這是否合乎神的旨意.一件工作是否合乎神旨,并不能立刻以果效來評判.不論果效多少?困難多少?仍不應灰心,因為工作的成敗,不在乎表面的效果.從歷史上可看見一些例證:
\newline
\newline
(1)保羅的異象
\newline
\newline
保羅在特羅亞看見異象,便離開特羅亞往馬其頓的腓立比去.到了腓立比,就立刻被人囚在獄中,這事是如何解釋呢?在監牢裡,是神的引導焉?行在神的旨意中,是不會錯的.保羅,西拉,雖在苦難中,仍然高歌讚美神.
\newline
\newline
(2)耶穌的際遇
\newline
\newline
主耶穌到世上工作,盡祂所能,栽培了十二個門徒.但想不到被釘十字架以前,竟然有一個門徒出賣自己,另一個否認自己,其他的都遠遠躲開了.祂所作的,似乎沒有甚麼果效,但耶穌並沒有埋怨,因祂知道有神的美意在其中.
\newline
\newline
(3)傳道的果子
\newline
\newline
內地會有一位傳教士被到青海,甘肅一帶工作.那地方的人很迷信,知識低,回教勢力又強,工作了二十年；只領了一個人歸主,好像工作失敗了.誰知神卻藉著這人信主,打開了青海傳福音的門,帶領許多人歸信主.
\newline
\newline
由此看來,工作有困難及攔阻,不多果效時；決不能按著工作外表的果效來批評工作的成就.百姓並不因這難處而灰心,因他們過去曾得到神清楚的啟示；知道神曾藉先知耶利米的口,預言到他們要回到耶路撒冷重建聖殿, 恢復向神的敬拜.現在建殿的事,雖然遭遇困難,有如黑暗中看不見亮光,更不清楚前面當走的路；但百姓沒有失望,徬惶.因為他們既然在光明中領受了神的話,又順服了遵神旨而行；深信神必定為他們成全建殿的工作.
\newline
\newline
百姓為我們留下了好榜樣,他們處境雖可怖,卻沒有懷疑在光明中,所領受的啟示,他們沒有用人的方法,反對壓迫自己的人,單單站在神前,不妥協,不屈服；結果神的時候到了,神便成全他們建殿的工作.照樣,我們需要重建生命中,與神交通的殿,撒旦也不放過我們；但只要認清楚這是神要我們作的,我們就不顧一切,遵神的話而行.
\newline
\newline


\newpage

\section{第43屆~港九培靈會~鮑會園~第六講}
\label{sec:1406}
\textbf{得勝有餘}
\newline
\newline
連結: \href{https://www.hkbibleconference.org/session-message/view/1406}{\texttt{https://www.hkbibleconference.org/session-message/view/1406}}
\newline
\newline
\hyperref[sec:1405]{< < < PREV SERMON < < <}
~
\hyperref[sec:toc]{[返目錄]}
~
\hyperref[sec:1407]{> > > NEXT SERMON > > >}
\newline
\newline
以斯拉記 1:1-6:22
\newline
\begin{longtable}{cl}
\hline
\hline
章節 & 經文 (和合本修訂版)\\
\hline
1:1 & \begin{tabularx}{0.7\textwidth}{X} 波斯王居魯士元年，耶和華為要應驗藉耶利米的口所說的話，就激發波斯王居魯士的心，使他下詔書通告全國，說： \end{tabularx} \\ \\ \relax
1:2 & \begin{tabularx}{0.7\textwidth}{X} 「波斯王居魯士如此說：耶和華天上的神已將地上萬國賜給我，又委派我在猶大的耶路撒冷為他建造殿宇。 \end{tabularx} \\ \\ \relax
1:3 & \begin{tabularx}{0.7\textwidth}{X} 你們中間凡作他子民的，可以上猶大的耶路撒冷去，重建耶和華－以色列神的殿，他是在耶路撒冷的神；願神與這人同在。 \end{tabularx} \\ \\ \relax
1:4 & \begin{tabularx}{0.7\textwidth}{X} 凡存留的人，無論寄居何處，那地的人要用金銀、財物、牲畜幫助他，還要為耶路撒冷神的殿甘心獻上禮物。」 \end{tabularx} \\ \\ \relax
1:5 & \begin{tabularx}{0.7\textwidth}{X} 於是，猶大和便雅憫的族長、祭司、利未人，凡是心被神感動的人都起來，要上耶路撒冷去建造耶和華的殿。 \end{tabularx} \\ \\ \relax
1:6 & \begin{tabularx}{0.7\textwidth}{X} 四圍所有的人都拿銀器、金子、財物、牲畜、珍寶支持他們，此外還有甘心獻的一切禮物。 \end{tabularx} \\ \\ \relax
1:7 & \begin{tabularx}{0.7\textwidth}{X} 居魯士王也把耶和華殿的器皿拿出來，這些器皿是尼布甲尼撒從耶路撒冷掠取，放在自己神明廟中的。 \end{tabularx} \\ \\ \relax
1:8 & \begin{tabularx}{0.7\textwidth}{X} 波斯王居魯士派米提利達司庫把這些器皿拿出來，點交給猶大的領袖設巴薩。 \end{tabularx} \\ \\ \relax
1:9 & \begin{tabularx}{0.7\textwidth}{X} 它們的數目如下：金盤三十個，銀盤一千個，刀二十九把， \end{tabularx} \\ \\ \relax
1:10 & \begin{tabularx}{0.7\textwidth}{X} 金碗三十個，備用銀碗四百一十個，其他器皿一千件。 \end{tabularx} \\ \\ \relax
1:11 & \begin{tabularx}{0.7\textwidth}{X} 金銀器皿共有五千四百件。被擄的人從巴比倫上耶路撒冷的時候，設巴薩把這一切都帶了上來。 \end{tabularx} \\ \\ \relax
2:1 & \begin{tabularx}{0.7\textwidth}{X} 這些是從被擄之地上來的省民，巴比倫王尼布甲尼撒把他們擄到巴比倫，他們重返耶路撒冷和猶大，各歸本城。 \end{tabularx} \\ \\ \relax
2:2 & \begin{tabularx}{0.7\textwidth}{X} 他們是同所羅巴伯、耶書亞、尼希米、西萊雅、利來雅、末底改、必珊、米斯拔、比革瓦伊、利宏、巴拿一起回來的。以色列百姓的人數如下： \end{tabularx} \\ \\ \relax
2:3 & \begin{tabularx}{0.7\textwidth}{X} 巴錄的子孫二千一百七十二名； \end{tabularx} \\ \\ \relax
2:4 & \begin{tabularx}{0.7\textwidth}{X} 示法提雅的子孫三百七十二名； \end{tabularx} \\ \\ \relax
2:5 & \begin{tabularx}{0.7\textwidth}{X} 亞拉的子孫七百七十五名； \end{tabularx} \\ \\ \relax
2:6 & \begin{tabularx}{0.7\textwidth}{X} 巴哈‧摩押的後裔，就是耶書亞和約押的子孫二千八百一十二名； \end{tabularx} \\ \\ \relax
2:7 & \begin{tabularx}{0.7\textwidth}{X} 以攔的子孫一千二百五十四名； \end{tabularx} \\ \\ \relax
2:8 & \begin{tabularx}{0.7\textwidth}{X} 薩土的子孫九百四十五名； \end{tabularx} \\ \\ \relax
2:9 & \begin{tabularx}{0.7\textwidth}{X} 薩改的子孫七百六十名； \end{tabularx} \\ \\ \relax
2:10 & \begin{tabularx}{0.7\textwidth}{X} 巴尼的子孫六百四十二名； \end{tabularx} \\ \\ \relax
2:11 & \begin{tabularx}{0.7\textwidth}{X} 比拜的子孫六百二十三名； \end{tabularx} \\ \\ \relax
2:12 & \begin{tabularx}{0.7\textwidth}{X} 押甲的子孫一千二百二十二名； \end{tabularx} \\ \\ \relax
2:13 & \begin{tabularx}{0.7\textwidth}{X} 亞多尼干的子孫六百六十六名； \end{tabularx} \\ \\ \relax
2:14 & \begin{tabularx}{0.7\textwidth}{X} 比革瓦伊的子孫二千零五十六名； \end{tabularx} \\ \\ \relax
2:15 & \begin{tabularx}{0.7\textwidth}{X} 亞丁的子孫四百五十四名； \end{tabularx} \\ \\ \relax
2:16 & \begin{tabularx}{0.7\textwidth}{X} 亞特的後裔，就是希西家的子孫九十八名； \end{tabularx} \\ \\ \relax
2:17 & \begin{tabularx}{0.7\textwidth}{X} 比賽的子孫三百二十三名； \end{tabularx} \\ \\ \relax
2:18 & \begin{tabularx}{0.7\textwidth}{X} 約拉的子孫一百一十二名； \end{tabularx} \\ \\ \relax
2:19 & \begin{tabularx}{0.7\textwidth}{X} 哈順的子孫二百二十三名； \end{tabularx} \\ \\ \relax
2:20 & \begin{tabularx}{0.7\textwidth}{X} 吉罷珥人九十五名； \end{tabularx} \\ \\ \relax
2:21 & \begin{tabularx}{0.7\textwidth}{X} 伯利恆人一百二十三名； \end{tabularx} \\ \\ \relax
2:22 & \begin{tabularx}{0.7\textwidth}{X} 尼陀法人五十六名； \end{tabularx} \\ \\ \relax
2:23 & \begin{tabularx}{0.7\textwidth}{X} 亞拿突人一百二十八名； \end{tabularx} \\ \\ \relax
2:24 & \begin{tabularx}{0.7\textwidth}{X} 亞斯瑪弗人四十二名； \end{tabularx} \\ \\ \relax
2:25 & \begin{tabularx}{0.7\textwidth}{X} 基列‧耶琳人、基非拉人、比錄人共七百四十三名； \end{tabularx} \\ \\ \relax
2:26 & \begin{tabularx}{0.7\textwidth}{X} 拉瑪人和迦巴人共六百二十一名； \end{tabularx} \\ \\ \relax
2:27 & \begin{tabularx}{0.7\textwidth}{X} 默瑪人一百二十二名； \end{tabularx} \\ \\ \relax
2:28 & \begin{tabularx}{0.7\textwidth}{X} 伯特利人和艾人共二百二十三名； \end{tabularx} \\ \\ \relax
2:29 & \begin{tabularx}{0.7\textwidth}{X} 尼波人五十二名； \end{tabularx} \\ \\ \relax
2:30 & \begin{tabularx}{0.7\textwidth}{X} 末必人一百五十六名； \end{tabularx} \\ \\ \relax
2:31 & \begin{tabularx}{0.7\textwidth}{X} 另一個以攔的子孫一千二百五十四名； \end{tabularx} \\ \\ \relax
2:32 & \begin{tabularx}{0.7\textwidth}{X} 哈琳的子孫三百二十名； \end{tabularx} \\ \\ \relax
2:33 & \begin{tabularx}{0.7\textwidth}{X} 羅德人、哈第人、阿挪人共七百二十五名； \end{tabularx} \\ \\ \relax
2:34 & \begin{tabularx}{0.7\textwidth}{X} 耶利哥人三百四十五名； \end{tabularx} \\ \\ \relax
2:35 & \begin{tabularx}{0.7\textwidth}{X} 西拿人三千六百三十名。 \end{tabularx} \\ \\ \relax
2:36 & \begin{tabularx}{0.7\textwidth}{X} 祭司：耶書亞家耶大雅的子孫九百七十三名； \end{tabularx} \\ \\ \relax
2:37 & \begin{tabularx}{0.7\textwidth}{X} 音麥的子孫一千零五十二名； \end{tabularx} \\ \\ \relax
2:38 & \begin{tabularx}{0.7\textwidth}{X} 巴施戶珥的子孫一千二百四十七名； \end{tabularx} \\ \\ \relax
2:39 & \begin{tabularx}{0.7\textwidth}{X} 哈琳的子孫一千零一十七名。 \end{tabularx} \\ \\ \relax
2:40 & \begin{tabularx}{0.7\textwidth}{X} 利未人：何達威雅的後裔，就是耶書亞和甲篾的子孫七十四名。 \end{tabularx} \\ \\ \relax
2:41 & \begin{tabularx}{0.7\textwidth}{X} 歌唱的：亞薩的子孫一百二十八名。 \end{tabularx} \\ \\ \relax
2:42 & \begin{tabularx}{0.7\textwidth}{X} 門口的守衛：沙龍的子孫、亞特的子孫、達們的子孫、亞谷的子孫、哈底大的子孫、朔拜的子孫，共一百三十九名。 \end{tabularx} \\ \\ \relax
2:43 & \begin{tabularx}{0.7\textwidth}{X} 殿役：西哈的子孫、哈蘇巴的子孫、答巴俄的子孫、 \end{tabularx} \\ \\ \relax
2:44 & \begin{tabularx}{0.7\textwidth}{X} 基綠的子孫、西亞的子孫、巴頓的子孫、 \end{tabularx} \\ \\ \relax
2:45 & \begin{tabularx}{0.7\textwidth}{X} 利巴拿的子孫、哈迦巴的子孫、亞谷的子孫、 \end{tabularx} \\ \\ \relax
2:46 & \begin{tabularx}{0.7\textwidth}{X} 哈甲的子孫、薩買的子孫、哈難的子孫、 \end{tabularx} \\ \\ \relax
2:47 & \begin{tabularx}{0.7\textwidth}{X} 吉德的子孫、迦哈的子孫、利亞雅的子孫、 \end{tabularx} \\ \\ \relax
2:48 & \begin{tabularx}{0.7\textwidth}{X} 利汛的子孫、尼哥大的子孫、迦散的子孫、 \end{tabularx} \\ \\ \relax
2:49 & \begin{tabularx}{0.7\textwidth}{X} 烏撒的子孫、巴西亞的子孫、比賽的子孫、 \end{tabularx} \\ \\ \relax
2:50 & \begin{tabularx}{0.7\textwidth}{X} 押拿的子孫、米烏寧的子孫、尼普心的子孫、 \end{tabularx} \\ \\ \relax
2:51 & \begin{tabularx}{0.7\textwidth}{X} 巴卜的子孫、哈古巴的子孫、哈忽的子孫、 \end{tabularx} \\ \\ \relax
2:52 & \begin{tabularx}{0.7\textwidth}{X} 巴洗律的子孫、米希大的子孫、哈沙的子孫、 \end{tabularx} \\ \\ \relax
2:53 & \begin{tabularx}{0.7\textwidth}{X} 巴柯的子孫、西西拉的子孫、答瑪的子孫、 \end{tabularx} \\ \\ \relax
2:54 & \begin{tabularx}{0.7\textwidth}{X} 尼細亞的子孫、哈提法的子孫。 \end{tabularx} \\ \\ \relax
2:55 & \begin{tabularx}{0.7\textwidth}{X} 所羅門僕人的後裔：瑣太的子孫、瑣斐列的子孫、比路大的子孫、 \end{tabularx} \\ \\ \relax
2:56 & \begin{tabularx}{0.7\textwidth}{X} 雅拉的子孫、達昆的子孫、吉德的子孫、 \end{tabularx} \\ \\ \relax
2:57 & \begin{tabularx}{0.7\textwidth}{X} 示法提雅的子孫、哈替的子孫、玻黑列‧哈斯巴音的子孫、亞米的子孫。 \end{tabularx} \\ \\ \relax
2:58 & \begin{tabularx}{0.7\textwidth}{X} 殿役和所羅門僕人的後裔共三百九十二名。 \end{tabularx} \\ \\ \relax
2:59 & \begin{tabularx}{0.7\textwidth}{X} 從特‧米拉、特‧哈薩、基綠、亞頓、音麥上來，不能證明他們的父系家族和後裔是否屬以色列的如下： \end{tabularx} \\ \\ \relax
2:60 & \begin{tabularx}{0.7\textwidth}{X} 第萊雅的子孫、多比雅的子孫、尼哥大的子孫，共六百五十二名。 \end{tabularx} \\ \\ \relax
2:61 & \begin{tabularx}{0.7\textwidth}{X} 祭司中，哈巴雅的子孫、哈哥斯的子孫、巴西萊的子孫，巴西萊因為娶了基列人巴西萊的女兒為妻，所以就以此為名。 \end{tabularx} \\ \\ \relax
2:62 & \begin{tabularx}{0.7\textwidth}{X} 這些人在族譜之中尋查自己的譜系，卻尋不著，因此算為不潔，不得作祭司。 \end{tabularx} \\ \\ \relax
2:63 & \begin{tabularx}{0.7\textwidth}{X} 省長對他們說，不可吃至聖的物，直到有會用烏陵和土明的祭司興起來。 \end{tabularx} \\ \\ \relax
2:64 & \begin{tabularx}{0.7\textwidth}{X} 全會眾共有四萬二千三百六十名。 \end{tabularx} \\ \\ \relax
2:65 & \begin{tabularx}{0.7\textwidth}{X} 此外，還有他們的僕婢七千三百三十七名，又有歌唱的男女二百名。 \end{tabularx} \\ \\ \relax
2:66 & \begin{tabularx}{0.7\textwidth}{X} 他們有七百三十六匹馬，二百四十五匹騾子， \end{tabularx} \\ \\ \relax
2:67 & \begin{tabularx}{0.7\textwidth}{X} 四百三十五匹駱駝，六千七百二十匹驢。 \end{tabularx} \\ \\ \relax
2:68 & \begin{tabularx}{0.7\textwidth}{X} 有些族長到了耶路撒冷耶和華的殿，為神的殿甘心獻上禮物，要在原有的根基上重新建造。 \end{tabularx} \\ \\ \relax
2:69 & \begin{tabularx}{0.7\textwidth}{X} 他們量力捐入工程的庫房，有六萬一千達利克金子，五千彌那銀子，以及一百件祭司的禮服。 \end{tabularx} \\ \\ \relax
2:70 & \begin{tabularx}{0.7\textwidth}{X} 於是祭司、利未人、百姓中的一些人、歌唱的、門口的守衛、殿役，各住在自己的城裡；以色列眾人都住在自己的城裡。 \end{tabularx} \\ \\ \relax
3:1 & \begin{tabularx}{0.7\textwidth}{X} 到了七月，以色列人住在自己的城裡；那時他們如同一人，聚集在耶路撒冷。 \end{tabularx} \\ \\ \relax
3:2 & \begin{tabularx}{0.7\textwidth}{X} 約薩達的兒子耶書亞和他的弟兄眾祭司，以及撒拉鐵的兒子所羅巴伯和他的弟兄，都起來建築以色列神的壇，要照神人摩西律法書上所寫的，在壇上獻燔祭。 \end{tabularx} \\ \\ \relax
3:3 & \begin{tabularx}{0.7\textwidth}{X} 他們在原有的根基上築壇，因為他們懼怕鄰邦民族，又在其上向耶和華早晚獻燔祭， \end{tabularx} \\ \\ \relax
3:4 & \begin{tabularx}{0.7\textwidth}{X} 並照律法書上所寫的守住棚節，按數照例每日獻所當獻的燔祭。 \end{tabularx} \\ \\ \relax
3:5 & \begin{tabularx}{0.7\textwidth}{X} 此後，他們獻常獻的燔祭，並在初一和耶和華一切分別為聖的節期獻祭，又向耶和華獻各人的甘心祭。 \end{tabularx} \\ \\ \relax
3:6 & \begin{tabularx}{0.7\textwidth}{X} 從七月初一起，雖然耶和華殿的根基尚未立定，他們開始向耶和華獻燔祭。 \end{tabularx} \\ \\ \relax
3:7 & \begin{tabularx}{0.7\textwidth}{X} 他們把銀子給石匠、木匠，把糧食、酒、油給西頓人、推羅人，好將香柏樹從黎巴嫩浮海運到約帕，是照波斯王居魯士所允准他們的。 \end{tabularx} \\ \\ \relax
3:8 & \begin{tabularx}{0.7\textwidth}{X} 他們到了耶路撒冷神殿的第二年，二月的時候，撒拉鐵的兒子所羅巴伯，約薩達的兒子耶書亞和其餘的弟兄，就是祭司和利未人，以及所有被擄歸回耶路撒冷的人，就開工建造；他們派二十歲以上的利未人，監督建造耶和華殿的工作。 \end{tabularx} \\ \\ \relax
3:9 & \begin{tabularx}{0.7\textwidth}{X} 於是何達威雅的後裔，就是耶書亞和他的子孫與弟兄、甲篾和他的子孫，他們和利未人希拿達的子孫與弟兄，都起來如同一人，監督那些在神殿裡做工的人。 \end{tabularx} \\ \\ \relax
3:10 & \begin{tabularx}{0.7\textwidth}{X} 工匠立耶和華殿根基的時候，祭司穿禮服吹號，利未人亞薩的子孫敲鈸，都照以色列王大衛親手所定的，站著讚美耶和華。 \end{tabularx} \\ \\ \relax
3:11 & \begin{tabularx}{0.7\textwidth}{X} 他們彼此唱和，讚美稱謝耶和華：「他本為善，他向以色列永施慈愛。」他們讚美耶和華的時候，眾百姓大聲呼喊，因為耶和華殿的根基已經立定。 \end{tabularx} \\ \\ \relax
3:12 & \begin{tabularx}{0.7\textwidth}{X} 然而有許多祭司、利未人和族長，就是見過先前那殿的老年人，現在親眼看見這殿立了根基，就大聲哭號，也有許多人大聲歡呼， \end{tabularx} \\ \\ \relax
3:13 & \begin{tabularx}{0.7\textwidth}{X} 百姓不能分辨歡呼的聲音或哭號的聲音，因為百姓大聲呼喊，聲音連遠處都可聽到。 \end{tabularx} \\ \\ \relax
4:1 & \begin{tabularx}{0.7\textwidth}{X} 猶大和便雅憫的敵人聽說被擄歸回的人為耶和華－以色列的神建造殿宇， \end{tabularx} \\ \\ \relax
4:2 & \begin{tabularx}{0.7\textwidth}{X} 就去見所羅巴伯和族長，對他們說：「請讓我們與你們一同建造，因為我們也與你們一樣尋求你們的神。自從亞述王以撒‧哈頓帶我們上這地的日子以來，我們常向神獻祭。」 \end{tabularx} \\ \\ \relax
4:3 & \begin{tabularx}{0.7\textwidth}{X} 但所羅巴伯、耶書亞和其餘以色列的族長對他們說：「我們建造神的殿與你們無關，因為我們要照波斯王居魯士所吩咐的，自己為耶和華－以色列的神協力建造。」 \end{tabularx} \\ \\ \relax
4:4 & \begin{tabularx}{0.7\textwidth}{X} 那地的人就在猶大百姓建造的時候，使他們的手發軟，擾亂他們。 \end{tabularx} \\ \\ \relax
4:5 & \begin{tabularx}{0.7\textwidth}{X} 從波斯王居魯士年間，直到波斯王大流士在位的時候，那些人賄賂謀士，要破壞他們的計劃。 \end{tabularx} \\ \\ \relax
4:6 & \begin{tabularx}{0.7\textwidth}{X} 亞哈隨魯在位，他的國度剛開始的時候，他們上書控告猶大和耶路撒冷的居民。 \end{tabularx} \\ \\ \relax
4:7 & \begin{tabularx}{0.7\textwidth}{X} 亞達薛西年間，比施蘭、米特利達、他別和他們的同僚上書奏告波斯王亞達薛西。奏文是用亞蘭文寫的，以亞蘭文呈上。 \end{tabularx} \\ \\ \relax
4:8 & \begin{tabularx}{0.7\textwidth}{X} 利宏省長、伸帥書記也上奏亞達薛西王，控告耶路撒冷如下 \end{tabularx} \\ \\ \relax
4:9 & \begin{tabularx}{0.7\textwidth}{X} （那時，利宏省長、伸帥書記和他們其餘的同僚，法官、官員、軍官、波斯官員、亞基衛人、巴比倫人，和書珊迦人，就是以攔人， \end{tabularx} \\ \\ \relax
4:10 & \begin{tabularx}{0.7\textwidth}{X} 以及被亞斯那巴大人遷移、安置在撒瑪利亞城和大河西邊一帶地方其餘的人。現在 ， \end{tabularx} \\ \\ \relax
4:11 & \begin{tabularx}{0.7\textwidth}{X} 這是他們上奏亞達薛西王奏文的抄本）：「河西的臣僕上奏亞達薛西王，現在 \end{tabularx} \\ \\ \relax
4:12 & \begin{tabularx}{0.7\textwidth}{X} 請王知道，從王那裡上到我們這裡的猶太人，已經抵達耶路撒冷。他們正在重建這反叛惡劣的城，已經完成了城牆，正要修復根基。 \end{tabularx} \\ \\ \relax
4:13 & \begin{tabularx}{0.7\textwidth}{X} 如今請王知道，這城若再建造，城牆完工，他們就不再進貢、納糧、繳稅，王的國庫必受虧損。 \end{tabularx} \\ \\ \relax
4:14 & \begin{tabularx}{0.7\textwidth}{X} 如今，我們吃的鹽既然全是宮廷的鹽，就不忍見王吃虧，因此奏告於王， \end{tabularx} \\ \\ \relax
4:15 & \begin{tabularx}{0.7\textwidth}{X} 請王考察先王史籍，必會在史籍上查知這城是反叛的城，對列王和各省有害；自古以來，城中常有悖逆的事，因此這城曾被拆毀。 \end{tabularx} \\ \\ \relax
4:16 & \begin{tabularx}{0.7\textwidth}{X} 我們謹奏王知，這城若再建造，城牆完工，河西之地王就無份了。」 \end{tabularx} \\ \\ \relax
4:17 & \begin{tabularx}{0.7\textwidth}{X} 那時王諭覆利宏省長、伸帥書記和他們其餘的同僚，就是住撒瑪利亞和河西一帶地方的人，說：「願你們平安。現在 \end{tabularx} \\ \\ \relax
4:18 & \begin{tabularx}{0.7\textwidth}{X} 你們所呈給我們的奏本，已經清楚地在我面前讀了。 \end{tabularx} \\ \\ \relax
4:19 & \begin{tabularx}{0.7\textwidth}{X} 我已下令考查，得知這城自古以來果然背叛列王，其中常有反叛悖逆的事。 \end{tabularx} \\ \\ \relax
4:20 & \begin{tabularx}{0.7\textwidth}{X} 也曾有強大的君王治理耶路撒冷，統管河西全地，人就給他們進貢、納糧、繳稅。 \end{tabularx} \\ \\ \relax
4:21 & \begin{tabularx}{0.7\textwidth}{X} 現在你們要下令叫這些人停工，使這城不得建造，等到我再降旨。 \end{tabularx} \\ \\ \relax
4:22 & \begin{tabularx}{0.7\textwidth}{X} 你們當謹慎辦這事，不可遲延，何必讓損害加重，使王受虧損呢？」 \end{tabularx} \\ \\ \relax
4:23 & \begin{tabularx}{0.7\textwidth}{X} 亞達薛西王上諭的抄本在利宏和伸帥書記，以及他們的同僚面前宣讀，他們就急忙往耶路撒冷去見猶太人，用勢力和強權叫他們停工。 \end{tabularx} \\ \\ \relax
4:24 & \begin{tabularx}{0.7\textwidth}{X} 於是，在耶路撒冷神殿的工程就停止了，直停到波斯王大流士第二年。 \end{tabularx} \\ \\ \relax
5:1 & \begin{tabularx}{0.7\textwidth}{X} 那時，哈該先知和易多的孫子撒迦利亞，兩個先知奉以色列神的名向猶大和耶路撒冷的猶太人說預言。 \end{tabularx} \\ \\ \relax
5:2 & \begin{tabularx}{0.7\textwidth}{X} 於是撒拉鐵的兒子所羅巴伯和約薩達的兒子耶書亞起來，開始建造耶路撒冷神的殿，有神的先知在那裡幫助他們。 \end{tabularx} \\ \\ \relax
5:3 & \begin{tabularx}{0.7\textwidth}{X} 當時河西的達乃總督和示他‧波斯乃，以及他們的同僚來對猶太人這樣說：「誰降旨讓你們建造這殿，完成這建築呢？」 \end{tabularx} \\ \\ \relax
5:4 & \begin{tabularx}{0.7\textwidth}{X} 於是我們告訴他們建造這建築物的人叫甚麼名字。 \end{tabularx} \\ \\ \relax
5:5 & \begin{tabularx}{0.7\textwidth}{X} 但神的眼目看顧猶太人的長老，以致沒有人叫他們停工，直到奏文上告大流士，得著他對這事的回諭。 \end{tabularx} \\ \\ \relax
5:6 & \begin{tabularx}{0.7\textwidth}{X} 這是河西的達乃總督和示他‧波斯乃，以及他們的同僚，就是住河西的官員，上書奏告大流士王的抄本， \end{tabularx} \\ \\ \relax
5:7 & \begin{tabularx}{0.7\textwidth}{X} 他們上書給王的奏文，其中寫著：「願大流士王諸事平安。 \end{tabularx} \\ \\ \relax
5:8 & \begin{tabularx}{0.7\textwidth}{X} 請王知道，我們往猶大省去，到了至大神的殿。這殿是用鑿成的石頭建造的，梁木插入牆內。這項工程進行迅速，在他們手中順利。 \end{tabularx} \\ \\ \relax
5:9 & \begin{tabularx}{0.7\textwidth}{X} 於是我們問那些長老，對他們這樣說：『誰降旨讓你們建造這殿，完成這建築呢？』 \end{tabularx} \\ \\ \relax
5:10 & \begin{tabularx}{0.7\textwidth}{X} 我們又問他們的名字，要記下他們領袖的名字，奏告於王。 \end{tabularx} \\ \\ \relax
5:11 & \begin{tabularx}{0.7\textwidth}{X} 他們這樣回答我們說：『我們是天和地之神的僕人，重建多年前所建造的殿，就是以色列一位偉大的君王建造完成的。 \end{tabularx} \\ \\ \relax
5:12 & \begin{tabularx}{0.7\textwidth}{X} 但因我們祖先惹天上的神發怒，神把他們交在迦勒底人巴比倫王尼布甲尼撒的手中，他就拆毀這殿，又把百姓擄到巴比倫。 \end{tabularx} \\ \\ \relax
5:13 & \begin{tabularx}{0.7\textwidth}{X} 然而巴比倫王居魯士元年，他降旨允准建造神的這殿。 \end{tabularx} \\ \\ \relax
5:14 & \begin{tabularx}{0.7\textwidth}{X} 神殿中的金銀器皿，就是尼布甲尼撒從耶路撒冷殿中掠取帶到巴比倫廟裡的，居魯士王從巴比倫廟裡取出來，交給派為省長，名叫設巴薩的， \end{tabularx} \\ \\ \relax
5:15 & \begin{tabularx}{0.7\textwidth}{X} 對他說：可以將這些器皿帶去，放在耶路撒冷的殿中，在原處建造神的殿。 \end{tabularx} \\ \\ \relax
5:16 & \begin{tabularx}{0.7\textwidth}{X} 於是那位設巴薩來建立耶路撒冷神殿的根基。但從那時直到如今，這殿尚未修建完畢。』 \end{tabularx} \\ \\ \relax
5:17 & \begin{tabularx}{0.7\textwidth}{X} 現在，王若以為好，請查閱巴比倫王的檔案庫，看居魯士王有沒有降旨允准在耶路撒冷建造神的殿。請降旨指示我們王對這件事的心意。」 \end{tabularx} \\ \\ \relax
6:1 & \begin{tabularx}{0.7\textwidth}{X} 於是大流士王降旨，要尋察典籍庫，就是在巴比倫藏檔案之處； \end{tabularx} \\ \\ \relax
6:2 & \begin{tabularx}{0.7\textwidth}{X} 在瑪代省亞馬他城的宮內尋得一卷，其中這樣寫著，「記錄如下： \end{tabularx} \\ \\ \relax
6:3 & \begin{tabularx}{0.7\textwidth}{X} 居魯士王元年，王降旨論到在耶路撒冷神的殿，要建造這殿作為獻祭之處，堅固它的根基。殿高六十肘，寬六十肘， \end{tabularx} \\ \\ \relax
6:4 & \begin{tabularx}{0.7\textwidth}{X} 要用三層鑿成的石頭，一層木頭，經費可出於王的庫房。 \end{tabularx} \\ \\ \relax
6:5 & \begin{tabularx}{0.7\textwidth}{X} 至於神殿的金銀器皿，就是尼布甲尼撒從耶路撒冷的殿中掠取帶到巴比倫的，必須歸還，帶回耶路撒冷的殿中，各按原處放在神的殿裡。」 \end{tabularx} \\ \\ \relax
6:6 & \begin{tabularx}{0.7\textwidth}{X} 「現在，河西的達乃總督和示他‧波斯乃，以及他們的同僚，就是住河西的官員，你們當遠離那裡。 \end{tabularx} \\ \\ \relax
6:7 & \begin{tabularx}{0.7\textwidth}{X} 不要攔阻這神殿的工作，任由猶太人的省長和長老在原處建造神的這殿。 \end{tabularx} \\ \\ \relax
6:8 & \begin{tabularx}{0.7\textwidth}{X} 我又降旨，吩咐你們為建造神的殿當向猶太人的長老這樣行：從王的財產中，由河西所繳納的貢銀，迅速支付這些人，免得工程停頓。 \end{tabularx} \\ \\ \relax
6:9 & \begin{tabularx}{0.7\textwidth}{X} 他們向天上的神獻燔祭所需用的公牛犢、公綿羊、小綿羊，以及麥子、鹽、酒、油，都要照耶路撒冷祭司的話，每日供給他們，不得有誤； \end{tabularx} \\ \\ \relax
6:10 & \begin{tabularx}{0.7\textwidth}{X} 好叫他們獻馨香的祭給天上的神，又為王和王眾子的壽命祈禱。 \end{tabularx} \\ \\ \relax
6:11 & \begin{tabularx}{0.7\textwidth}{X} 我再降旨，無論誰更改這命令，必從他房屋中拆出一根梁木，把他舉起，懸在其上，又使他的房屋為此成為糞堆。 \end{tabularx} \\ \\ \relax
6:12 & \begin{tabularx}{0.7\textwidth}{X} 任何王或百姓若伸手更改這命令，拆毀在耶路撒冷神的這殿，願那立他名在那裡的神將他們滅絕。我大流士降這諭旨，你們要速速遵行。」 \end{tabularx} \\ \\ \relax
6:13 & \begin{tabularx}{0.7\textwidth}{X} 於是，河西的達乃總督和示他‧波斯乃，以及他們的同僚，急速遵行大流士王所頒的命令。 \end{tabularx} \\ \\ \relax
6:14 & \begin{tabularx}{0.7\textwidth}{X} 猶太人的長老因哈該先知和易多的孫子撒迦利亞的預言，就建造這殿，凡事順利。他們遵照以色列神的命令和波斯王居魯士、大流士、亞達薛西的諭旨，建造完畢。 \end{tabularx} \\ \\ \relax
6:15 & \begin{tabularx}{0.7\textwidth}{X} 大流士王第六年，亞達月初三，這殿完工了。 \end{tabularx} \\ \\ \relax
6:16 & \begin{tabularx}{0.7\textwidth}{X} 以色列人、祭司和利未人，以及其餘被擄歸回的人都歡歡喜喜地為神的這殿行奉獻禮。 \end{tabularx} \\ \\ \relax
6:17 & \begin{tabularx}{0.7\textwidth}{X} 他們為這神殿的奉獻禮獻了一百頭公牛、二百隻公綿羊、四百隻小綿羊，又照以色列支派的數目獻十二隻公山羊，作以色列眾人的贖罪祭。 \end{tabularx} \\ \\ \relax
6:18 & \begin{tabularx}{0.7\textwidth}{X} 他們派祭司按著班次，利未人也按著班次在耶路撒冷事奉神，正如摩西律法書上所寫的。 \end{tabularx} \\ \\ \relax
6:19 & \begin{tabularx}{0.7\textwidth}{X} 正月十四日，被擄歸回的人守逾越節。 \end{tabularx} \\ \\ \relax
6:20 & \begin{tabularx}{0.7\textwidth}{X} 祭司和利未人一同自潔，他們全都潔淨了。利未人為被擄歸回的眾人和他們的弟兄眾祭司，並為自己宰逾越節的羔羊。 \end{tabularx} \\ \\ \relax
6:21 & \begin{tabularx}{0.7\textwidth}{X} 從被擄之地歸回的以色列人，並所有歸附他們、除掉這地外邦人的污穢、尋求耶和華－以色列神的人，都吃這羔羊。 \end{tabularx} \\ \\ \relax
6:22 & \begin{tabularx}{0.7\textwidth}{X} 他們歡歡喜喜地守除酵節七日，因為耶和華使他們歡喜。耶和華又使亞述王的心轉向他們，堅固他們的手，去做神－以色列神殿的工。 \end{tabularx} \\ \\
[1ex]
\hline
\hline
\end{longtable}
聖殿的根基雖已建好,但因敵人厲害的攻擊,百姓被迫停工；各人回去蓋自己天花板的房屋,而將建殿的工程拖延下去.忽略建造的工作,誠屬可怕的危機；正如一個基督徒,有生命的建造,才能作有力的見證.許多人停止了建造他們屬靈的殿宇,難怪只徒有外表,而靈裡的生命,卻非常枯乾軟弱.當以色列人停止建造聖殿時,他們的情形是相當可憐的；幸而有神的先知哈該和撒迦利亞,親自提醒責備他們,使他們重新恢復聖殿的建造.
\newline
\newline
(一)先知哈該指責他們:(該一1-11)
\newline
\newline
(1)要省察自己的行為
\newline
\newline
先知哈該和撒迦利亞,就是這時代工作者.因著哈該的督責,使百姓能恢復建殿,哈該到他們當中傳信息,主要就是說:「萬軍之耶和華如此說,你們要省察自己的行為.」先知重複這句「省察自己的行為」,就是提醒百姓,他們已失去當初的愛心,和當初建殿的動機,因而失去裡面推動的力量.表面看來,停止建殿的工作,是受到外界仇敵的攻擊,但真正的失敗乃是內在的軟弱.今天我們也該省察自己的行為,個人靈命不長進之原因何在?
\newline
\newline
(2)只顧天花板的房屋
\newline
\newline
先知哈該向百姓傳信息,是代表萬軍之耶和華指責他們.但百姓卻以「建造耶和華殿的時候尚未來到.」為藉口推辭；可是他們的自私,只顧自己,而不敢面對現實負起責任.我們何嘗不是如此,用忙碌的理由來推辭靈裡的追求；明明是自己應負的責任,反而找出理由來原諒自己.先知哈該責備百姓說:殿仍荒涼,你們還自己住天花板的房屋嗎?你們該省察自己的行為,看靈的生命失敗到一個甚麼境況?你們只關心自己肉體的喜好,享受；而失去靈裡的份量,不能分辨輕重.
\newline
\newline
(3)工作沒有果效
\newline
\newline
「你們撒的種多,收的卻少；你們吃卻不得飽,喝卻不得足；穿衣服卻不得暖,得工錢的,將工錢裝在破漏的囊中.」做工作沒有果效,這是靈裡衰落的表現,由於沒有關心自己,在神面前靈裡的建造,屬靈生命的建造便停止了.所以一切工作,都不蒙神的記念,作工也沒有真實的果效.
\newline
\newline
(4)禱告不蒙垂聽
\newline
\newline
「所以為你們的緣故,天就不降甘露,地也不出土產.」他們向神祈求,神卻不賜福,祈求也不蒙垂聽,屬靈的生命未建造得好,禱告也不會產生果效.
\newline
\newline
先知哈該毫無保留地,提醒及責備百姓,當神的信息臨到百姓時,他們有良好的反應；「......那時撒拉鐵的兒子所羅巴伯,和約撒答的兒子大祭司耶書亞,並剩下的百姓,都聽從耶和華他們神的話.......百姓也在耶和華面前存敬畏的心,耶和華的使者哈該奉耶和華的差遣對百姓說:「耶和華說,我與你們同在.」因這個責備的信息便成了百姓恢復重新建殿的力量,他們便照著神的吩咐省察自己的行為；改變自己的生活,便努力起來為神繼續建造.
\newline
\newline
(二)百姓悔改起來建造
\newline
\newline
(1)神前悔改對付己罪
\newline
\newline
當百姓在神前存敬畏的心,悔改認罪時,神便說:「我與你們同在.」神肯與他們同工,又有豐滿的福氣為他們預備,便立刻向他們施恩.「肯對付己罪」永遠是復興靈命的條件.百姓何時在神前有悔改的態度,神便等候著將福氣降下,賜福與百姓.今天神也等待著賜福與我們,只要我們願意作這樣的器皿,把生命毫無保留地放在神面前.
\newline
\newline
(2)轉移眼目仰望真神
\newline
\newline
「神的眼目看顧猶大的長老,以致總督等沒有叫他們停工.」(拉五5)敵人不能用任何方法來攻擊他們,因有耶和華眼目的看顧.他們知道如何仰望神看顧的眼目,這是信心的表現.他們本是最軟弱的百姓,沒有能力,財富,權柄和盼望；但他們在逼迫困難中,並沒有看惡劣的環境,反而將眼目轉向神的看顧!以致困難不能壓制他們,環境不能使他們失敗；真正的信心是透過困難,仍然看得見神.
\newline
\newline
(3)聽從神命遵行神道
\newline
\newline
「猶大長老因先知哈該,和易多的孫子撒迦利亞所說勸勉的話,就建造這殿,凡事亨通.他們遵著以色列神的命令和波斯王,古列,大利烏,亞達薛西的旨意,建造完畢.」聽從神命遵行神道,也是屬靈長進的基本因素,在混亂中,百姓沒有聽從其他的聲音；只聽從神向他們說的話,又照著遵行,結果他們在惡劣的環境裡得勝了.在現今混亂,充滿了各樣混雜聲音的時代裡；我們應該學習在混亂中,分辨神的聲音.我們必須在神話語上下功夫,有一個學習聽神的話的經驗,追求明白神話語的心；平時有聽覺的經驗,屬靈的耳朵便不會糢糊了.
\newline
\newline
(4)獻上一切在壇上
\newline
\newline
「以色列的祭司和利未人,並其餘被擄歸回的人,都歡歡喜喜的行奉獻神殿的禮.」他們把自己一切所有的都樂意獻在神前,當時他們所獻的為數不少；計有公牛一百隻,公綿羊二百隻,綿羊羔四百隻.雖然這個數目與所羅門時代聖殿奉獻的數目無與倫比.這些百姓剛從巴比倫回來,實在一無所有,連路上的糧食都是王所賜的,如今他們肯這樣獻上,實足以表現出他們已把自己所有的獻在神前.深信他們所獻的,比所羅門所獻的更蒙悅納.神所看重的,不是奉獻數目的多寡,而是看重奉獻的心,因為神曾說:「金子銀子都是我的.」
\newline
\newline
百姓接受先知的警告和提醒,徹底在神前痛悔,對付己罪.又將眼目轉移仰望真神.聽從神的命令,又順服遵行!更將一切都獻在神前,這真是奉獻生命的表現；若今天我們在靈裡有這樣的經歷,又聽到神的聲音,而肯獻上一切,就必蒙神所悅納.
\newline
\newline


\newpage

\section{第43屆~港九培靈會~鮑會園~第七講}
\label{sec:1407}
\textbf{持守見證}
\newline
\newline
連結: \href{https://www.hkbibleconference.org/session-message/view/1407}{\texttt{https://www.hkbibleconference.org/session-message/view/1407}}
\newline
\newline
\hyperref[sec:1406]{< < < PREV SERMON < < <}
~
\hyperref[sec:toc]{[返目錄]}
~
\hyperref[sec:1408]{> > > NEXT SERMON > > >}
\newline
\newline
以斯拉記 8:1-8:36
\newline
\begin{longtable}{cl}
\hline
\hline
章節 & 經文 (和合本修訂版)\\
\hline
8:1 & \begin{tabularx}{0.7\textwidth}{X} 這些是亞達薛西王在位的時候，同我從巴比倫上來的族長和他們的家譜： \end{tabularx} \\ \\ \relax
8:2 & \begin{tabularx}{0.7\textwidth}{X} 屬非尼哈的子孫有革順；屬以他瑪的子孫有但以理；屬大衛的子孫有哈突； \end{tabularx} \\ \\ \relax
8:3 & \begin{tabularx}{0.7\textwidth}{X} 屬示迦尼的子孫；屬巴錄的子孫有撒迦利亞，同著他按家譜計算，男丁一百五十人； \end{tabularx} \\ \\ \relax
8:4 & \begin{tabularx}{0.7\textwidth}{X} 屬巴哈‧摩押的子孫有西拉希雅的兒子以利約乃，同著他有男丁二百人； \end{tabularx} \\ \\ \relax
8:5 & \begin{tabularx}{0.7\textwidth}{X} 屬薩土的子孫有雅哈悉的兒子示迦尼，同著他有男丁三百人； \end{tabularx} \\ \\ \relax
8:6 & \begin{tabularx}{0.7\textwidth}{X} 屬亞丁的子孫有約拿單的兒子以別，同著他有男丁五十人； \end{tabularx} \\ \\ \relax
8:7 & \begin{tabularx}{0.7\textwidth}{X} 屬以攔的子孫有亞他利雅的兒子耶篩亞，同著他有男丁七十人； \end{tabularx} \\ \\ \relax
8:8 & \begin{tabularx}{0.7\textwidth}{X} 屬示法提雅的子孫有米迦勒的兒子西巴第雅，同著他有男丁八十人； \end{tabularx} \\ \\ \relax
8:9 & \begin{tabularx}{0.7\textwidth}{X} 屬約押的子孫有耶歇的兒子俄巴底亞，同著他有男丁二百一十八人； \end{tabularx} \\ \\ \relax
8:10 & \begin{tabularx}{0.7\textwidth}{X} 屬巴尼的子孫有約細斐的兒子示羅密，同著他有男丁一百六十人； \end{tabularx} \\ \\ \relax
8:11 & \begin{tabularx}{0.7\textwidth}{X} 屬比拜的子孫有比拜的兒子撒迦利亞，同著他有男丁二十八人； \end{tabularx} \\ \\ \relax
8:12 & \begin{tabularx}{0.7\textwidth}{X} 屬押甲的子孫有哈加坦的兒子約哈難，同著他有男丁一百一十人； \end{tabularx} \\ \\ \relax
8:13 & \begin{tabularx}{0.7\textwidth}{X} 屬亞多尼干的子孫，就是晚到的，他們的名字是以利法列、耶利、示瑪雅，同著他們有男丁六十人； \end{tabularx} \\ \\ \relax
8:14 & \begin{tabularx}{0.7\textwidth}{X} 屬比革瓦伊的子孫有烏太和撒刻，同著他們有男丁七十人。 \end{tabularx} \\ \\ \relax
8:15 & \begin{tabularx}{0.7\textwidth}{X} 我召集這些人在流入亞哈瓦的河旁邊，我們在那裡紮營三日。我查看百姓和祭司，發現並沒有利未人在那裡， \end{tabularx} \\ \\ \relax
8:16 & \begin{tabularx}{0.7\textwidth}{X} 就派人到以利以謝、亞列、示瑪雅、以利拿單、雅立、以利拿單、拿單、撒迦利亞、米書蘭等領袖，以及約雅立和以利拿單教師那裡。 \end{tabularx} \\ \\ \relax
8:17 & \begin{tabularx}{0.7\textwidth}{X} 我吩咐他們往迦西斐雅地方去見那裡的領袖易多，又告訴他們當向易多和他的弟兄，就是迦西斐雅那地方的殿役說甚麼話，好為我們神的殿帶事奉的人來。 \end{tabularx} \\ \\ \relax
8:18 & \begin{tabularx}{0.7\textwidth}{X} 蒙我們神施恩的手幫助我們，他們在以色列的曾孫，利未的孫子，抹利的後裔中帶了一個精明的人來，就是示利比，還有他的眾子與兄弟共十八人。 \end{tabularx} \\ \\ \relax
8:19 & \begin{tabularx}{0.7\textwidth}{X} 另外，還有哈沙比雅，同著他有米拉利的子孫耶篩亞，以及他的眾子和兄弟共二十人。 \end{tabularx} \\ \\ \relax
8:20 & \begin{tabularx}{0.7\textwidth}{X} 從前大衛和眾領袖派殿役服事利未人，現在從這殿役中也帶了二百二十人來，全都是按名指定的。 \end{tabularx} \\ \\ \relax
8:21 & \begin{tabularx}{0.7\textwidth}{X} 那時，我在亞哈瓦河邊宣告禁食，為要在我們神面前刻苦己心，求他使我們和我們的孩子，以及一切所有的，都得平坦的道路。 \end{tabularx} \\ \\ \relax
8:22 & \begin{tabularx}{0.7\textwidth}{X} 我以求王撥步兵騎兵幫助我們抵擋路上的仇敵為羞愧，因我們曾對王說：「我們神施恩的手必幫助凡尋求他的，但他的能力和憤怒必攻擊凡離棄他的。」 \end{tabularx} \\ \\ \relax
8:23 & \begin{tabularx}{0.7\textwidth}{X} 我們為此禁食祈求我們的神，他就應允我們。 \end{tabularx} \\ \\ \relax
8:24 & \begin{tabularx}{0.7\textwidth}{X} 我分派十二位祭司長，就是示利比、哈沙比雅和與他們一起的兄弟十人， \end{tabularx} \\ \\ \relax
8:25 & \begin{tabularx}{0.7\textwidth}{X} 把王和謀士、軍官，並在那裡的以色列眾人為我們神殿所獻的金銀和器皿，都秤了交給他們。 \end{tabularx} \\ \\ \relax
8:26 & \begin{tabularx}{0.7\textwidth}{X} 我秤了交在他們手中的有六百五十他連得銀子，一百他連得銀器，一百他連得金子， \end{tabularx} \\ \\ \relax
8:27 & \begin{tabularx}{0.7\textwidth}{X} 二十個金碗，值一千達利克，上等光亮的銅器皿兩個，珍貴如金。 \end{tabularx} \\ \\ \relax
8:28 & \begin{tabularx}{0.7\textwidth}{X} 我對他們說：「你們歸耶和華為聖，器皿也歸為聖；金銀是甘心獻給耶和華－你們列祖之神的。 \end{tabularx} \\ \\ \relax
8:29 & \begin{tabularx}{0.7\textwidth}{X} 你們要警醒看守，直到你們在祭司長和利未族長，以及以色列的各族長面前，在耶路撒冷耶和華殿的庫房內，把這些過了秤。」 \end{tabularx} \\ \\ \relax
8:30 & \begin{tabularx}{0.7\textwidth}{X} 於是，祭司和利未人把秤過的金銀和器皿接過來，要帶到耶路撒冷我們神的殿裡。 \end{tabularx} \\ \\ \relax
8:31 & \begin{tabularx}{0.7\textwidth}{X} 正月十二日，我們從亞哈瓦河邊起行，要往耶路撒冷去。我們神的手保佑我們，救我們脫離仇敵和路上埋伏之人的手。 \end{tabularx} \\ \\ \relax
8:32 & \begin{tabularx}{0.7\textwidth}{X} 我們到了耶路撒冷，在那裡住了三日。 \end{tabularx} \\ \\ \relax
8:33 & \begin{tabularx}{0.7\textwidth}{X} 第四日，金銀和器皿都在我們神的殿裡過了秤，交在烏利亞的兒子米利末祭司的手中。同著他的有非尼哈的兒子以利亞撒，還有利未人耶書亞的兒子約撒拔和賓內的兒子挪亞底。 \end{tabularx} \\ \\ \relax
8:34 & \begin{tabularx}{0.7\textwidth}{X} 那時，這一切都點過秤過了，重量全寫在冊上。 \end{tabularx} \\ \\ \relax
8:35 & \begin{tabularx}{0.7\textwidth}{X} 從被擄之地歸回的人向以色列的神獻燔祭，為以色列眾人獻十二頭公牛、九十六隻公綿羊、七十七隻小綿羊，又獻十二隻公山羊作贖罪祭，這些全都是獻給耶和華的燔祭。 \end{tabularx} \\ \\ \relax
8:36 & \begin{tabularx}{0.7\textwidth}{X} 被擄歸回的人把王的諭旨交給王的總督與河西的省長，他們就支助百姓和神的殿。 \end{tabularx} \\ \\
[1ex]
\hline
\hline
\end{longtable}
以斯拉歸回耶路撒冷後,為要預備百姓重建城牆.以色列人回國重建聖殿後六十年,那時,以斯拉便回國,負起教導百姓的責任.聖殿是百姓敬拜神的地方,代表靈裡與神的交通,聖殿被毀,人與神的交通便中斷了!現在聖殿得以重建,這是復興人與神的交通.百姓在靈裡得到復興,但城牆的建造,是代表百姓恢復在外邦地的見證.所以尼希米帶領百姓重建城牆的工作,也是相當重要.
\newline
\newline
以色列人恢復與神的交通敬拜後,要等六十多年才能恢復在外邦地的見證.因為神要等候使用一個器皿──以斯拉,起來幫助百姓恢復見證的工作,這個領導的人未興起,他們的工作就不能完全.照樣神今天在教會中,尋找一些能堵住破口的人,使用他們來工作.因為在教會裡敷衍塞責的基督徒太多了,我們應該在屬靈的事上有雄心大志,來防堵教會的破口.我們這班蒙恩的人,在神面前均有聖潔君尊祭司的職份,我們每一個都可以為神工作,問題乃在乎我們願否把自己交在神手中.
\newline
\newline
神使用以斯拉在當時代,做預備訓練的工作,成為神合用的器皿,因為在他身上有幾個基本的條件:
\newline
\newline
(一)屬神生命(拉七1-5)
\newline
\newline
以斯拉是西萊雅的兒子,西萊雅的祖宗是以利亞撒.他是大祭司亞倫的兒子,所以以斯拉直系是亞倫的後裔.利未人是神特別揀選的一個支派,以斯拉在支派中是特選的,又是神特別使用的人.他生命中與神所揀選的地位,有特殊的關係.他有一個屬神揀選的生命,又有一個特別的地位.在神所使用的百姓中,這是一個極重要的條件.事奉神的工作不是一種職業,乃是生命的關係；因上事奉神的人必須是亞倫的後裔,有君尊祭司的身份；又是神所揀選的族類百姓,也是天國的國民.今天教會中仍有一種錯誤的觀念,就是平信徒與受聖職的人有很大的分別.信徒以為那些人才有特別的權利為神工作,但聖經告訴我們,原來每一個屬神的人,都有君尊祭司的職份.沒有一個人能推諉自己的責任,事奉神的人與不事奉神的人的分別,乃在乎所領受之職,於耶穌及整個教會,每人都有直接的關係.
\newline
\newline
(二)精通律法(拉七6)
\newline
\newline
以斯拉從巴比倫上來,他是一位敏捷的文士,通達耶和華神賜與摩西的律法.「敏捷」原文是指軍人,巫術等人使用工具時,特別精銳敏捷.用武器的方法,與沒有受過訓練的人完全異同,因為他熟練之故.以斯拉是一位敏捷的文士,表明他在神的道上有過訓練；在真理上有過練習,明白如何使用真理；正如戰士之使用武器,音樂家之使用樂器一樣.這是教會需要注意的,要懂得應用神的話在自己的生活上,又要在神話語上下工夫.神的話是有力的工具,正如兩刃的利劍一般.要懂得如何運用,使他成為屬靈的武器.
\newline
\newline
(三)遵行真道(拉七10)
\newline
\newline
「以斯拉定志考究遵行耶和華的律法,又就律例典章教 訓以色列人.」以斯拉之所以能用真理教導以色列人,因為他先領受,又先遵行；如此便能在神話語上有功效,又能應用它來作戰.如果一個人自己沒有先在靈裡領受神的話,恐怕所傳的道是缺乏力量的,不能在別人心中發生果效.使徒行傳的作者路加也曾明說,他所記載的事,先是耶穌基督自己所行的,然後才記載祂所教導別人的.主耶穌也為我們留下好榜樣,先遵行神的話,然後才將神的話教導別人.
\newline
\newline
(四)經歷神恩(拉七27-28)
\newline
\newline
「以斯拉說:耶和華我們列祖的神,是應當稱頌的,因他使王起這心意,修飾耶路撒冷耶和華的殿；又在王和謀士,並大能的軍長面前施恩與我,因耶和華我神的手幫助我.......」以斯拉不但熟悉遵行神的話,而且在靈裡對神有真正的認識.整本聖經沒有題及以斯拉的出身,也沒有說到他工作前的詳情.他出來工作的感動是從神而來,因為他在靈裡與神有交通,經驗到神自己,便按神旨而行.靈裡的經驗,成為他生命行動的力量.道路雖然困難,但因著經驗過神恩的實在,便成了他的力量；不因遇到攔阻,便容易跌倒了.聖經教導我們,應當相信神,信主是從聽道而來,聽道是藉著神的話而來,這樣便成了信心的根基.
\newline
\newline
(五)訓練同工(拉八15-23)
\newline
\newline
以斯拉對神的工作有真正的負擔,又因在靈裡根基,使他知道自己不能承擔重任,需要預備更多人起來分擔為神工作.當他返回耶路撒冷後,卻發現當中沒有一個是利未人,也沒有神所揀選特別的百姓；更沒有一個在神的話上,受過訓練及有根基的人.於是他便到處尋找一些這樣的人,好與他同工；後來找到了二百六十名,便帶領他們一同回到耶路撒冷工作.以斯拉認識清楚,倘若要帶領百姓復興歸回,對神有認識,不能單憑自己的思想感覺.必須要建立在真道的根基上,而負起教導責任的人,更要在真理上受過訓練的.
\newline
\newline
今天的教會普遍來說,都忽略了為神訓練,預備合用的人材.當神興起祂工作以前,祂先要把負擔放在教會的弟兄姊妹身上.大家要一同起來工作,正如以斯拉回到耶路撒冷；先與同伴在整個耶路撒冷,做教導百姓的工作.鼓勵百姓起來事奉主,結果全城的百姓得著大復興；使城牆重建,百姓在外邦地作有力的見證.
\newline
\newline
現今中國大陸有七億人口,二十歲以下的青年沒有機會聽過耶穌的名字.倘若有一天大陸傳福音的門開了,而我們又沒有好好訓練的工人,試問怎能應付這龐大的需要呢?正如一九四六年時,麥克亞瑟將軍佔領日本.在他統治時,曾向美國許多傳福音的機構大聲疾呼,因為他發覺日本是一個成熟的工場.日本人將特別歡喜聽道理,但為了不夠工人的差派,以致福音的門被打開,而沒有工人進去.直至今天,一千個日本人中只有一位是基督徒,而且對於基督教傳教士的排斥也相當厲害,這是多麼可惜的事!求主叫我們有膽量選擇年青人接受訓練,直至傳福音的門打開時,有充足的工人,預備差派接受工作使命.
\newline
\newline
(六)言行一致(拉八21-23)
\newline
\newline
「......我求王撥步兵,馬兵,幫助我們抵擋路上的仇敵,本以為羞恥；因我曾對王說,我們神施恩的手必幫助一切尋求他的.......」以斯拉在王面前曾見證神施恩的手必幫助他,在太平盛世的巴比倫境內如此說並不困難,因為那裡沒有敵人.但現在處於曠野地,常在敵人面前又沒有王的步兵馬兵,更無人的幫助保護；單純仰望耶和華的幫助,實在談何容易!以斯拉認清楚了,沒有危險困難時對神有信心,到遇見困難時,便轉過來求人的幫助,是失見證的.必須言行一致,在生活上有見證.因為唯有言行一致的,才能顯出對神的信心,真正有信心的生活,是在困難試煉中信靠神的話.
\newline
\newline
因著以斯拉具備有這基本的條件,便成為神手中合用的器皿.在這時代,神也需要尋找一些人防堵破口,為神工作正如以斯拉一樣的人.
\newline
\newline


\newpage

\section{第43屆~港九培靈會~鮑會園~第八講}
\label{sec:1408}
\textbf{重新得力}
\newline
\newline
連結: \href{https://www.hkbibleconference.org/session-message/view/1408}{\texttt{https://www.hkbibleconference.org/session-message/view/1408}}
\newline
\newline
\hyperref[sec:1407]{< < < PREV SERMON < < <}
~
\hyperref[sec:toc]{[返目錄]}
~
\hyperref[sec:1364]{> > > NEXT SERMON > > >}
\newline
\newline
以斯拉記 9:1-10:17
\newline
\begin{longtable}{cl}
\hline
\hline
章節 & 經文 (和合本修訂版)\\
\hline
9:1 & \begin{tabularx}{0.7\textwidth}{X} 這些事完成以後，眾領袖來接近我，說：「以色列百姓、祭司和利未人沒有棄絕迦南人、赫人、比利洗人、耶布斯人、亞捫人、摩押人、埃及人和亞摩利人等列邦民族所行可憎的事。 \end{tabularx} \\ \\ \relax
9:2 & \begin{tabularx}{0.7\textwidth}{X} 因他們為自己和兒子娶了這些外邦女子，以致聖潔的種籽和列邦民族混雜，而且領袖和官長在這事上是罪魁。」 \end{tabularx} \\ \\ \relax
9:3 & \begin{tabularx}{0.7\textwidth}{X} 我一聽見這事，就撕裂衣服和外袍，拔了頭髮和鬍鬚，驚惶地坐著。 \end{tabularx} \\ \\ \relax
9:4 & \begin{tabularx}{0.7\textwidth}{X} 凡為以色列神言語戰兢的人，都因被擄歸回之人所犯的罪，聚集到我這裡來。我驚惶地坐著，直到獻晚祭的時候。 \end{tabularx} \\ \\ \relax
9:5 & \begin{tabularx}{0.7\textwidth}{X} 獻晚祭的時候我從愁煩中起來，穿著撕裂的衣服和外袍，雙膝跪下，向耶和華－我的神舉手， \end{tabularx} \\ \\ \relax
9:6 & \begin{tabularx}{0.7\textwidth}{X} 說：「我的神啊，我抱愧蒙羞，不敢向你－我的神仰面，因為我們的罪孽多到滅頂，我們的罪惡滔天。 \end{tabularx} \\ \\ \relax
9:7 & \begin{tabularx}{0.7\textwidth}{X} 從我們祖先的日子直到今日，我們的罪惡深重；因我們的罪孽，我們和君王、祭司都交在鄰國諸王的手中，被殺害，擄掠，搶奪，臉上蒙羞，正如今日的景況。 \end{tabularx} \\ \\ \relax
9:8 & \begin{tabularx}{0.7\textwidth}{X} 現在耶和華－我們的神暫且向我們施恩，為我們留下一些殘存之民，使我們如釘子釘在他的聖所，讓我們的神光照我們的眼目，使我們在受轄制之中稍微復興。 \end{tabularx} \\ \\ \relax
9:9 & \begin{tabularx}{0.7\textwidth}{X} 我們是奴僕，然而在受轄制之中，我們的神沒有丟棄我們，在波斯諸王面前向我們施恩，叫我們復興，能重建我們神的殿，修補毀壞之處，使我們在猶大和耶路撒冷有城牆。 \end{tabularx} \\ \\ \relax
9:10 & \begin{tabularx}{0.7\textwidth}{X} 「我們的神啊，既然如此，現在我們還有甚麼話可說呢？因為我們離棄了你的誡命， \end{tabularx} \\ \\ \relax
9:11 & \begin{tabularx}{0.7\textwidth}{X} 就是你藉你僕人眾先知所吩咐的，說：『你們要去得為業之地是污穢之地，因列邦民族的污穢和可憎的事，叫這地從這邊到那邊都充滿了污穢。 \end{tabularx} \\ \\ \relax
9:12 & \begin{tabularx}{0.7\textwidth}{X} 現在，不可把你們的女兒嫁給他們的兒子，也不可為你們的兒子娶他們的女兒，永不可求他們的平安和他們的利益，這樣你們就可以強盛，吃這地的美物，並把這地留給你們的子孫永遠為業。』 \end{tabularx} \\ \\ \relax
9:13 & \begin{tabularx}{0.7\textwidth}{X} 我們因自己的惡行和大罪，遭遇這一切的事，但你－我們的神懲罰我們輕於我們罪所當得的，又為我們留下這些殘存之民。 \end{tabularx} \\ \\ \relax
9:14 & \begin{tabularx}{0.7\textwidth}{X} 我們豈可再違背你的誡命，與行這些可憎之事的民族結親呢？若我們這樣行，你豈不向我們發怒，將我們滅絕，以致沒有一個餘民或殘存之民嗎？ \end{tabularx} \\ \\ \relax
9:15 & \begin{tabularx}{0.7\textwidth}{X} 耶和華－以色列的神啊，你是公義的，我們才能剩下這些殘存之民，正如今日的景況。看哪，我們在你面前有罪惡，因此無人能在你面前站立得住。」 \end{tabularx} \\ \\ \relax
10:1 & \begin{tabularx}{0.7\textwidth}{X} 以斯拉禱告，認罪，哭泣，俯伏在神殿前的時候，有以色列中的男女和孩童聚集到以斯拉那裡，成了一個盛大的會，百姓無不痛哭。 \end{tabularx} \\ \\ \relax
10:2 & \begin{tabularx}{0.7\textwidth}{X} 以攔的子孫，耶歇的兒子示迦尼對以斯拉說：「我們娶了這地的外邦女子，干犯了我們的神，然而現在以色列人在這事上還有指望。 \end{tabularx} \\ \\ \relax
10:3 & \begin{tabularx}{0.7\textwidth}{X} 現在，我們要與我們的神立約，送走所有的妻子和她們所生的，照著主和那些因我們神誡命戰兢之人所議定的，按律法去行。 \end{tabularx} \\ \\ \relax
10:4 & \begin{tabularx}{0.7\textwidth}{X} 起來，這是你當辦的事，我們必支持你，你當奮勇而行。」 \end{tabularx} \\ \\ \relax
10:5 & \begin{tabularx}{0.7\textwidth}{X} 以斯拉就起來，叫祭司長和利未人，以及以色列眾人起誓，要照這話去做；他們就起了誓。 \end{tabularx} \\ \\ \relax
10:6 & \begin{tabularx}{0.7\textwidth}{X} 以斯拉從神殿前起來，進入以利亞實的兒子約哈難的屋裡，到了那裡不吃飯，也不喝水，為被擄歸回之人所犯的罪悲傷。 \end{tabularx} \\ \\ \relax
10:7 & \begin{tabularx}{0.7\textwidth}{X} 他們通告猶大和耶路撒冷，叫所有被擄歸回的人聚集在耶路撒冷。 \end{tabularx} \\ \\ \relax
10:8 & \begin{tabularx}{0.7\textwidth}{X} 凡不遵照領袖和長老所議定，三日之內不來的，就必毀壞他所有的財產，把他從被擄歸回之人的會中開除。 \end{tabularx} \\ \\ \relax
10:9 & \begin{tabularx}{0.7\textwidth}{X} 於是，猶大和便雅憫眾人三日之內都聚集在耶路撒冷。那時是九月，那月的二十日，眾百姓坐在神殿前的廣場，因這事，又因下大雨，就都戰抖。 \end{tabularx} \\ \\ \relax
10:10 & \begin{tabularx}{0.7\textwidth}{X} 以斯拉祭司站起來，對他們說：「你們有罪了，因為你們娶了外邦女子，增添以色列的罪惡。 \end{tabularx} \\ \\ \relax
10:11 & \begin{tabularx}{0.7\textwidth}{X} 現在當向耶和華－你們列祖的神認罪，遵行他的旨意，離開這地的百姓和外邦女子。」 \end{tabularx} \\ \\ \relax
10:12 & \begin{tabularx}{0.7\textwidth}{X} 全會眾大聲回答說：「好！我們必照著你的話去做。 \end{tabularx} \\ \\ \relax
10:13 & \begin{tabularx}{0.7\textwidth}{X} 只是百姓眾多，又逢大雨的季節，我們沒有氣力站在外面；這也不是一兩天可以辦完的事，因我們在這事上犯了大罪。 \end{tabularx} \\ \\ \relax
10:14 & \begin{tabularx}{0.7\textwidth}{X} 讓我們的領袖代表全會眾留在那裡。我們城鎮中凡娶外邦女子的，當按所定的日期，會同本城的長老和審判官前來，直到辦完這事，神的烈怒轉離我們。」 \end{tabularx} \\ \\ \relax
10:15 & \begin{tabularx}{0.7\textwidth}{X} 惟有亞撒黑的兒子約拿單，特瓦的兒子雅哈謝反對這事，並有米書蘭和利未人沙比太支持他們。 \end{tabularx} \\ \\ \relax
10:16 & \begin{tabularx}{0.7\textwidth}{X} 被擄歸回的人就如此做了。以斯拉祭司按著父家指名選派一些族長。十月初一，他們一同坐下來查辦這事， \end{tabularx} \\ \\ \relax
10:17 & \begin{tabularx}{0.7\textwidth}{X} 到正月初一，才查清所有娶外邦女子的人數。 \end{tabularx} \\ \\
[1ex]
\hline
\hline
\end{longtable}
以色列人歸回重建聖殿,已大功告成了,他們可以重新恢復與神的交通和敬拜.雖然當中曾受過不少的攻擊和反對,但靠著神的恩典至終也能得勝.就在這時,在極大喜樂中,有一陰影蓋在他們身上；以斯拉發現百姓中間,有很多人犯大罪,更有領袖帶領百姓去犯罪.本來他為著這些肯離開巴比倫,歸回耶路撒冷,一同敬拜及追求的百姓,感到滿心歡喜快樂；但料想不到這些人靈裡的境況,恐怕還要比留在巴比倫的以色列人更壞!心中因百姓陷在罪惡中,而感到萬分失望.很多時候,我們在神前有一個心志,肯追求又肯付代價；一方面因著所付的代價,達到目標了.可是另一方面,又為撒旦留下一個極大的破口；一方面追求,而另一方面又疏忽了.若祗恢復與神的交通,而沒有儆醒,也不能保証生活蒙神的悅納.不要忘記撒旦是一個可怕的仇敵,牠知道我們那一方面最軟弱,最容易受攻擊；不要讓一方面的成功,遮蓋了我們的眼目.只要我們稍不留心和謹慎,便為撒旦所乘.
\newline
\newline
(一)罪惡的性質
\newline
\newline
他們知道以色列百姓,神的選民不該與外邦女子通婚,也不能將自己的女兒,給外邦人為妻,可是他們卻明知而故犯.這罪是很特別的,雖然在聖經裡沒有詳細記載,但猶大歷史家約瑟弗曾記載著說:「當百姓回到耶路撒冷時,他們知道應該向四週的人作見証；但他們以為外邦人太多,自己的人太少,因此便採用結親之法與外邦人在一起,向這些不認識神的人作見証.當時有多少人這樣做,我們不大清楚；但從這件事上給我們看出一個極重要屬靈的原則.雖然以色列人懷著一個為耶和華做見証的目的和動機,但實行時卻按照自己的方法去完成.這個目的是好的,可惜用了神不歡喜的方法去完成,為了達到目的而不擇手段.在屬靈的事上萬不可如此,屬靈的事情目的要緊.但達到目標的手段也是要緊,不能用人的詭詐的方法,完成神的工作.
\newline
\newline
在某處基督教機構裡,有一個同工作非常不誠實,常把東西私自拿去,他還解釋說:「偷一本聖經是為了拯救一個靈魂.」拯救人的靈魂目的是對的,可是這方法永遠不會蒙神的悅納.我們所做的工作是神賜福的,因此我們所用的方法,也該是神賜福的.
\newline
\newline
另一方面,「因他們為自己和兒子,娶了這些外邦女子為妻,以致聖潔的種類和這些國的民混雜......」(拉九2),這是罪的真正的性質,屬神的百姓和不屬神的人混雜了；以致屬神的人的性質失掉了,為神作見証的力量也沒有了!藉著混雜將屬神的人,與不屬神的人連在一起,失去應有的特徵.以色列人在混合的事上失去見証,今天的教會也需要有分別的見証；顯出我們與世人在個人生活上有分別,教會的事奉上,應有分別的見証來保守自己.
\newline
\newline
(二)勝罪的原因
\newline
\newline
1．敏銳的感覺
\newline
\newline
以色列人在這罪惡中,順利得到勝利,是因為在靈裡有非常敏銳的感覺.過去他們犯罪,要神刑罰他們,才在罪惡中甦醒過來.但這次犯大罪,並沒有經過神的刑罰,當自己發現到罪惡軟弱時,便在神前悔改；這是復興以後的生活.一個與神有交通的人仍然有可能落在罪中,但當落在罪中而知道悔改；這就是靈裡甦醒的見証,我們也需要在罪惡上有敏銳的感覺.
\newline
\newline
教會在初期時代,有夫婦亞拿尼亞和撒非喇欺哄聖靈,把奉獻的錢,私自留下幾份,彼得就叫他們立刻死亡!犯罪的人,決不能在教會中存在.因為當時的教會是聖潔的,對於罪惡有敏銳的感覺.今天的教會也需要對罪惡有敏銳的感覺,犯了罪不須要等神的警告,就應該立刻悔改.靈裡敏銳的感覺,能使我們在軟弱上,立刻悔改過來,不容許罪惡在我們生命中存在.
\newline
\newline
2．神前的痛悔
\newline
\newline
以斯拉發現百姓犯罪以後,便撕裂衣服和外袍,拔了頭髮和鬍鬚,驚懼憂悶而坐；表示自己為百姓的罪惡,心裡十分難過.他沒有向人發脾氣,只在神面前為著百姓和領袖們的罪痛哭流淚.因著他一人的悔改,所有的百姓都聚集在那裡一同悔改.以斯拉並沒有用責備,來改變人的行為；只是為他們的罪心裡悲哀痛苦,他便成了一位肯為罪悔改的領袖.藉著他的榜樣,把犯罪的百姓帶到神面前來.今天的教會也需要這樣的人,幫助那些犯罪的弟兄姊妹,帶領他們悔改.
\newline
\newline
3．真正的信心
\newline
\newline
本書第八章與第十章記載的事,顯出一個非常重要的真理.第八章記以斯拉有完全順服依靠神的心,在耶路撒冷危險的路上,有敵人來攻擊他；他本可以求波斯王派兵保護,但因著他曾表示相信依靠耶和華的可勝過一切困難.若求助於波斯王便失去見証,所以他藉禁食禱告仰望神,神就保護他在路上有平安.第十章載,當百姓知罪在神前悔改時,他們相信神能幫助他們勝過罪惡,他們就果真勝過了罪惡.神的作為往往與我們信靠神的心成正比,我們信神大,神就在我們當中顯大.倘若我們對神的觀念太少,信心太少；神就不能在我們中間行大事.主說:在人看是不可能,但在信的人凡事都能.
\newline
\newline
拉撒路死了四天了,耶穌來到他姊姊馬大,馬利亞的家；她們兩人在耶穌面前哭.他們知道耶穌愛拉撒路.也知道祂有權能；但拉撒路已經死了四天了還有甚麼辦法叫他復活呢?感謝主!祂使拉撒路從死裡復活,好叫他們曉得在主凡事都能,信心便能生長.
\newline
\newline
當時的百姓信得過神的偉大,因此神在他們中間行了大奇事.聖經說:「你要大大張口,我就給你充滿.」(詩八一10)所以我們該把信心放大,神就要在我們中間顯為大.不要用自己的信心,把神的作為限制了.
\newline
\newline
4．堅持的毅力(拉十13)
\newline
\newline
「會眾都大聲回答說,我們必照著你的話行,只是百姓眾多；又逢大雨的時令,我們不能站在外頭；這也不是一兩天辦完的事,因我們在這事上犯了大罪......」百姓的得勝,是因為他們有堅持的毅力.知道這件事的意義是十分重要,也相當難以處理.但卻願意在神前立約繼續不斷的去完成.這不是一時的熱心,他們為著對付一件事,肯用幾個月的時間去持守,保持始終如一的心.
\newline
\newline
許多人接受神話語的感動,心中火熱立志,但過不多時,這熱心便冷淡不致能持久.結果有始無終,漸漸失去在神面前的領受.我們應持守在神前的領受,直至神的應許在我們身上成就；讓我們始終如一,使領受成為力量,走完神要我們走的路程.
\newline
\newline


\newpage

\section{第44屆~港九培靈會~吉兆頎~第一講　關於異象外的意義}
\label{sec:1364}
\textbf{第一講　關於異象外的意義}
\newline
\newline
連結: \href{https://www.hkbibleconference.org/session-message/view/1364}{\texttt{https://www.hkbibleconference.org/session-message/view/1364}}
\newline
\newline
\hyperref[sec:1408]{< < < PREV SERMON < < <}
~
\hyperref[sec:toc]{[返目錄]}
~
\hyperref[sec:1365]{> > > NEXT SERMON > > >}
\newline
\newline
箴言 29:18-29:18
\newline
\begin{longtable}{cl}
\hline
\hline
章節 & 經文 (和合本修訂版)\\
\hline
29:18 & \begin{tabularx}{0.7\textwidth}{X} 沒有異象，民就放肆；惟遵守律法的，便為有福。 \end{tabularx} \\ \\
[1ex]
\hline
\hline
\end{longtable}
聖經告訴我們:「沒有異象,民就放肆；惟遵守律法的,便為有福.」「那人對我說:人子阿!凡我所指示你的,你都要用眼睛看,用耳聽；並要放在心上；我帶你到這裡來,特為要指示你,凡你所見的,你都要告訴以色列家.」這就是說,異象對信徒本身,是有重大關係的.若異象對我們沒有深切的影響,聖經就不會有這樣多的篇幅來記載,必沒有這麼多的人順服了.
\newline
\newline
異象的意義──就是神將未來的計劃,使有心追求的人,有負擔的人,給他們啟示,直到神的計劃完成.若沒有負擔,沒有追求的心願,決不能看見異象,因為這種人看不見神的心,摸不著神的心.故有異象之人,不論如何艱難,一定領受,一定接受負擔,接受神的託付擺上自己.以西結先知就有這樣的負擔,他是有追求,有心願的人.以西結雖在外邦人中間,神卻將未來的計劃啟示他.
\newline
\newline
所以「以西結」還有最大的異象,有最多的異象,＿有最準確的異象.這對我們來說,也是最有價值的異象,而且是非常重要的.因為以西結的處境,他的地位,他的身份,與我們一樣.他是祭司,我們也是祭司.彼得前書告訴我們,信徒也是祭司,而且是君尊的祭司.我們有一個任務,就是要將人帶到神面前.雖然我們現在所處的時代,是一個彎曲悖謬的世代,被魔鬼擄去,被罪惡擄去,但我們是君尊的祭司,如果我們願意在神面前承認我們的職份,願意對神說:「神啊!你怎樣將異象給以西結,也求你怎樣將異象給我；以西結怎樣完成你給他的異象,也求你讓我也能完成你的異象.」諸位!神不是不願意將異象給我們,問題是我們是否願意擔起做祭司的職份!你若有心願,有負擔,神就會將祂的計劃,祂的託付,向你顯明,使你盡祭司的職份.以西結的心願是最寶貴的,他雖是被擄之人,又是在被擄之人中間；但他不是沒有負擔的人,不是沒有追求的人,不是沒有領受的人,乃是一個能領受神未來計劃的人.後來以色列人回國,都應驗了神給以西結看見的異象.
\newline
\newline
摩西在埃及時,雖落在外邦人的手,但他有愛國的心志,他在選民中間有心願,有負擔,他不願意神的子民建立在魔鬼的國度裡,不願意神的兒女作魔鬼的工.他動了血氣,出了問題,但他的負擔卻未完成.他逃到曠野,有四十年之久,作牧人,受盡折磨.但他對神的負擔未有減少,沒有看輕也沒有退後.一直在神面前等候,仰望,順服.時候到了,有一天,神不得不把異象啟示他,亦不得不將計劃告訴他.因為他在神前有心願,有負擔,故神將異象啟示他!
\newline
\newline
(一)異象的重要
\newline
\newline
異象對我們是非常重要的,因為神將祂的心意向人啟示,是神將計劃向人披露.故箴言廿九章四節說:「沒有異象,民就放肆.」就是說,沒有異象,人就馬虎；沒有異象,人就隨便；沒有異象,人就放縱肉體,看輕神的工作,不重視神的託付.今天,在教會裡面,在神的家中,這種人太多了.對神的工作,可有可無,可做可不做.如果有異象的人,也就不能放鬆,也不能放下,至死也要完成.在以西結身上,就可以看見這樣情形,他至死忠心,一領受,就立即付諸行動.
\newline
\newline
摩西也是如此,他不願做王子,寧願與神的百姓一同受苦.一個大有學問,大有智慧,大有抱負,大有雄心的人；若不看重異象,神就要挑起他的負擔,調整他應走的道路,可見異象對我們何等重要.
\newline
\newline
還有,異象是對付魔鬼最好的利器,以西結一章三節說:「在迦勒底人之地,迦巴魯河邊,耶和華的話,特別臨到布西的兒子祭司以西結,耶和華的靈,降在他身上.」可見,異象是針對魔鬼的詭計,是勝過魔鬼最好的利器.一個信徒有了異象,就有屬靈的智慧,在神工作上,就能識破魔鬼的詭計.
\newline
\newline
神要揀選一個人,就是亞伯拉罕,他是迦勒底吾珥人.神要揀選他,魔鬼卻要破壞他,神要用他,魔鬼卻要奪去他.神選召亞伯拉罕出迦勒底的吾珥,為要建立祂的國,使以色列人成為祂的民；魔鬼卻要摩西成為牠的奴僕,成為牠的囚犯.但神要對付他,他在甚麼地方跌倒,就在甚麼地爬起；在甚麼地方軟弱,就在甚麼地方剛強.這是神的方法,是屬靈的原則.神要在迦勒底找一個忠心的人,神要藉著異象叫他看見魔鬼的詭計,神要將祂的計劃藉異象啟示他.
\newline
\newline
以西結聽見神的異象,就拚命去傳,拚命去講,也不顧一切,因為他識破魔鬼的詭計,他看見了異象,就將自己擺上.
\newline
\newline
今天神也要找能承擔異象的人,有託付的人.但可惜在神家中,太多人忽略異象!在被擄之地,還以為平平安安；在撒旦手上,仍以為平平安安.所以有人一離開主,就被撒旦擄去了.神要呼召你出來,到應許之地,活在祂的家裡；這是異象,對我們是重要的.神要保守有心願,有負擔的人；神要藉著他們復興祂的工作,拯救被擄的百姓.
\newline
\newline
以西結雖在被擄的人中,但沒有被潮流沖倒.「他在迦勒底人之地,迦巴魯河邊.」他在河邊,不是在河中,表示不被潮流沖倒.一個有心願,有負擔的人,是站在潮流的河邊；不被潮流沖去,不被潮流奪走的.
\newline
\newline
弟兄姊妹!如果你在屬靈的事上,感覺放鬆了,感覺隨便了；在神工作上,感覺棘手,這就是潮流.這潮流很利害,把許多人沖走了!許多人因貪愛世界,走巴蘭的道路,拜金牛犢；就走上世界的道路,被世界擄去了.神所保守的人,是有心願的人；愛神的人,就是不隨從今日潮流的人.
\newline
\newline
弟兄姊妹!請問你今日站在何處?如果你站在潮流中,就不能作神的中流砥柱.神要保守的人,乃是為神擺上的人,是服從神權柄的人.
\newline
\newline
保羅在亞基帕王面前受審時,他說:「我沒有違背從天上來的異象.」這就是說,保羅一直受神的權柄管治.他好像說,我今日雖站在你面前受審,也是面對異象,你不過是我異象的一部份,所順服的,乃是神的權柄.所以我們要將神的權柄,在教會裡,在我們家中,在兒女身上,在我們的工作場所,不折不扣的表現出來.
\newline
\newline
有異象的人,不但順服神直接的權柄,也順服那代表神的權柄.因為順服直接的權柄易,順服間接的權柄難.神是聖潔,公義,慈愛,信實,不改變的；故此我們必須服從直接的權柄,也要順服代表的權柄.比方在教會中有牧師,有長老,有執事,我們也當服從他們的權柄.因為一個有異象的人,是必然服從權柄的.沒有負擔,沒有異象的人,才會藉口抗拒.但有異象的人,不但服從直接的權柄,也服從間接的權柄.
\newline
\newline
每一個國家,都有兩處地方是人所不願意去的；一個是警察局,另一個是法院.尤其是香港人,更深知此中情形.假如我們請求港督撤銷警署和法院,行得通嗎?當然不可能的.因為港督若真的撤銷了,那麼,治安不是更糟了嗎?可見在教會內,一定要學會服從代表的權柄,亦即服從間接的權柄.
\newline
\newline
馬路當中有許多鐵欄,為了確保人車的安全,你不能說左右都沒有車子開來,跳過去又何妨.但若被警察看見,就會給你告到法院去,你還得要認罪受罰哩.你不能說跳欄杆是我的自由,因為你必須服從代表的權柄.
\newline
\newline
又如有幾個小孩子,不願意受父母的約束,便開會,決議開除父母,想另找尋合心意的好父母,這當然這是不能的.
\newline
\newline
所以在教會內,沒有異象,就不能服從神的權柄,也不能服從代表的權柄,更不能完成神所託負的.唯一方法,求神給你異象,能服從直接的權柄和間接的權柄；盡心竭力,去完成神的託付.不然,問題便出來了,被世界的潮流沖去了.
\newline
\newline
感謝神!有心願的人,他頭上的天不是銅的,不是鐵的,乃是敞開的；神的話臨到他們,如同臨到以西結一樣.
\newline
\newline


\newpage

\section{第44屆~港九培靈會~吉兆頎~第三講　異象的功用}
\label{sec:1365}
\textbf{第三講　異象的功用}
\newline
\newline
連結: \href{https://www.hkbibleconference.org/session-message/view/1365}{\texttt{https://www.hkbibleconference.org/session-message/view/1365}}
\newline
\newline
\hyperref[sec:1364]{< < < PREV SERMON < < <}
~
\hyperref[sec:toc]{[返目錄]}
~
\hyperref[sec:1366]{> > > NEXT SERMON > > >}
\newline
\newline
箴言 29:18-29:18
\newline
\begin{longtable}{cl}
\hline
\hline
章節 & 經文 (和合本修訂版)\\
\hline
29:18 & \begin{tabularx}{0.7\textwidth}{X} 沒有異象，民就放肆；惟遵守律法的，便為有福。 \end{tabularx} \\ \\
[1ex]
\hline
\hline
\end{longtable}
(1)要顯出人的光景
\newline
\newline
人有時自以為義,有時將壞的行為當作標準.但有了異象,就完全不同.因為異象第一個功用是顯出人的光景.人明白自己的光景,就將自己放在異象的天秤上.若不認識自己的光景,根本就沒有異象.沒有異象,就容易放鬆,可以隨意而行了.
\newline
\newline
摩西有本事打死了人,暗暗把他埋了,第二天又來,理直氣壯的去干涉以色列人,那知事敗,不得不逃亡到米甸曠野.當摩西看見了異象後,他就認識了自己說:我是甚麼人,怎能去見法老?我是拙口笨舌的,怎能去說服法老?
\newline
\newline
許多人喜歡引用摩西的話說:我是拙口笨舌的人,其實他卻是口若懸河的人,不過藉口推辭工作而已.許多人講屬靈的見證,但一句造就人的話也沒有.只會批評,論斷.弟兄姊妹,讓我們在神面前悔改過來吧!
\newline
\newline
以賽亞書,寫得洋洋大觀,究竟力量從何而來呢?乃是從異象而來的.因為以賽亞看見了,聽見了,也領受了,遵行了.他說:「禍哉!我滅亡了,我是嘴唇不潔的人,又住在嘴唇不潔的民中.」可見異象一來,光景就立刻顯明.你是活在神的光景中,抑或活在人的光景中呢?異象會叫你完全顯出來的.
\newline
\newline
沒有異象,神就無法用你,教會也無法進步,所以異象是顯出人的光景,你先要看見自己快將滅亡的光景,然後神才能用你!
\newline
\newline
保羅起初以逼迫教會,逼迫信徒為能事,以為是事奉上帝.他在人面前,在神面前都可以得意洋洋,大搖大擺,手拿文書,口吐恐嚇的話,門徒看見他,魂也會被嚇走離開了軀體.但他見了異象就悔改,為主所用,作外邦人的光.後來他也見證說:「在罪人中我是個罪魁,然而我蒙了憐恤!」可見,有了異象,人就大大的改變.誰去改變他,乃是異象使他改變.
\newline
\newline
保羅回到耶路撒冷,沒有人敢接待他,因為他曾逼害信主的人.可歎一個人失敗,眾人就定罪；但一個人悔改了,眾人仍不接納.我們不能在主的旨意之外去使用權柄的.
\newline
\newline
在耶路撒冷,只有一個人,就是巴拿巴,願意接待保羅,彼此得造就.於是他們在安提阿傳福音,建立教會.門徒稱為基督徒是從安提阿起首的；是因巴拿巴,保羅而來的.如果沒有巴拿巴的愛心,信心,不但保羅失去前途,主也失去了一個合用的器皿,豈不可惜嗎?
\newline
\newline
所以異象是要顯出人的光景,各位弟兄姊妹!你看見教會冷淡,失去起初愛心,基督徒去走世界的道路；你暗中流淚仍是不夠,要在明處有行動的配合,才可以挽回頹廢的局面.
\newline
\newline
(2)謹守神的引導與託付
\newline
\newline
「神說,一次兩次,世人卻不理會.人躺在床上沉睡的時候,上帝就用夢,和夜間的異象,開通他們的耳朵；將當受的教訓,印在他們心上.好叫人不從自己的謀算,不行驕傲的事,攔阻人不陷於坑裡,不死在刀下.」經文:(伯三三14-18)
\newline
\newline
現在與各位思想異象的功用的第二方面.聖經告訴我們,有異象的人與無異象的人,完全不同.有異象的人,是用神的方法去解決人的事；沒有異象的人,是用人的方法去解決人的事.有了異象,一定清楚神的託付,在屬靈的路上,事奉在工作上,不會茫無頭緒,無所適從,而且有條不紊,因為有異象印證了功用.
\newline
\newline
沒有異象,就會盲目奔走；有異象後,就謹守神的引導與託付.以賽亞看見異象後便說:「主啊!我在這裡,請差遣我.」沒有異象的時候,可以照自己的旨意隨便作工,有異象之後,便說:「主阿!我在這裡,請差遣我.」因為已清楚了神的託付.
\newline
\newline
有異象的人可以為主作工,沒有異象的人也可以為主作工；但工作的結果,工作的賞賜,大大不同.沒有異象的人,工作的果效不被神喜悅,也不被神嘉許.比方,我喜歡種花,我到花店去買花的種子,我看準了店主背著我的時候,偷了一包花種子.回家以後,照種的方法,把買的一包種子撒下土裡,也把偷來的那一包撒下土裡,生長起來是沒有分別的.花種沒有問題,但我這個人便有了問題.按屬靈的定律,道種沒有問題,但人有了問題,是人的關係.有了異象後,清楚知道如何工作,明白為誰作工,也知道如何作工,清楚神的託付.在神的家中,對付得好,調整得好,聖靈在裡面活得好,工作便認真.所以主耶穌教導門徒說,你們要進窄門,走小路.門窄路小,找著的人也少.但大路是通到滅亡去的,闊門也是把人引進滅亡的.
\newline
\newline
窄門小路,多一點,小一點都不引.而且是先入門,後走路.若先走路後進門,是異途同歸,先入門後走路,是永生之路.入闊門走大路,可任意而為；但入小門走窄路,就必須要方方正正的,受神拘束.
\newline
\newline
羅得和他的妻子從所多瑪出來後,妻子因貪愛世界,回頭一望便成了一根鹽柱.羅得後來生了二子,一是摩押人的始祖,一是亞捫人的始祖,他們與以色列人世世代代為仇.因為他違背了神的異象,沒有清楚神的託付.
\newline
\newline
保羅站在亞基帕王面前的時候,他說:「我沒有違背從天上來的異象.」他清楚神的託付,也守住神的託付.
\newline
\newline
感謝神!有異象的人,沒有一個不清楚神的託付的.彼得從來不到外邦人的家,不與外邦人交往,遵守猶太人的規條.但有了異象之後,他第一個走進外邦人中,且知道神是不偏待人的；因為他已清楚了神的託付,甚至至死忠心,永不改變.就算上十字架,也在所不辭.異象的功用,在彼得身上,在保羅身上是如此；在你我身上也是如此,都有同樣的功用.求主幫助我們也照樣清楚神的託付.
\newline
\newline
(3)打破人的傳統
\newline
\newline
人的傳統對人的捆綁很利害,特別是以色列人,他們有民族性的傳統,文化的傳統,道德的傳統,與及宗教的傳統.信主耶穌,仍無法衝破這些傳統.誰領保羅出這些傳統.他殺害信徒,逼迫教會,還以為為上帝發熱心；這是傳統,是牢不可破的傳統.誰能打破這些傳統,沒有人能,但異象到來時,就脫離傳統,也脫離捆綁.
\newline
\newline
彼得,保羅何等愛惜傳統,保守傳統.但異象一來,保羅就說,人得救不是靠傳統,也不是靠割禮；不是靠外表的規定,禮儀,因為律法是叫人死.
\newline
\newline
甚麼事令保羅衝破這些傳統,是異象令他衝破此傳統,主耶穌怎樣用異象引導他,保羅就照著遵行.
\newline
\newline
古時先賢,接受神的託付,與我們今人接受神的託付是一樣的,但果效不同；他們跟得上神的呼召,我們卻感覺困難.神給我們引導,但我們要思想,要考慮,要看看與傳統有無抵觸.如有,就寧願犧牲神的託付,仍要保存傳統；因此就失去為神工作的機會,失去為神工作的能力,這實在是很大的損失.
\newline
\newline
異象臨到保羅後,他和路加當夜就起行.在神家中的人,每人好像樂隊:有大小的樂器,有高低的樂器,有長短的樂器；發出不同的聲音,獨奏是不行的,合奏才好聽.但要隨著指揮的手,若不跟著指揮,奏出來的聲音就不和諧.
\newline
\newline
異象一來,傳統的觀念就沒有了.如果給傳統捆綁,讀經就有了問題,比如說,聖經說要謙卑,他們就把聖經裡所有的不要；只要謙卑,因為只戴了謙卑的眼鏡.如果用傳統的觀念來衡量聖經,非常危險.在台灣,有一次家庭禮拜,本來是我講道的,我後來轉請一位西國同工講道.這裡的家庭聚會有些不良的習慣,主人家常為來聚會的人預備一些水果,餅食等.那天,講道完畢,我進到廚房裡要看看這家的主婦,正在忙著搬茶點出來給客人吃.這主婦看見我,立刻說:「今天我丈夫在家,你不講道,你不講聖經說,丈夫要愛妻子.」我說:「聖經也說,妻子要順服丈夫.」我感謝神,我沒有照傳統觀念去行.
\newline
\newline
彼得也好,保羅也好,有異象,傳統就沒有了,這是異象有功用.
\newline
\newline
(4)管教神的兒女
\newline
\newline
在神的面前,神的兒女有異象的約束,所以聖經說:「沒有異象,民就放肆.」沒有異象,人就放縱；有了異象,就要受異象的約束,帶領,我們對異象忠心.
\newline
\newline
保羅來到哥林多,為道迫切,向猶太人證明耶穌是基督,他一直講,講到筋疲力竭,他想要離開,但神不讓他離開.主在異象中對保羅說:「不要怕,只管講!不要閉口,有我與你同在,必沒有人下手害你；因為在這城裡,我有許多百姓.」可惜我們沒有受異象的管治.
\newline
\newline
(a)管治人的腳
\newline
\newline
沒有異象的人走闊路,有異象的人走窄路；生活,動作,存留都要受神的管治.愈愛神,愈親近神,愈發覺神管教的重要；如果沒有神引導,就不能作甚麼.
\newline
\newline
保羅要去特羅亞去,聖靈禁止.保羅似乎真難以忍受,他可以說,神啊!你召我作外邦人的光,如今我不是向外邦人作見證嗎?但往那裡,到何處,你都不許.神卻要保羅到馬其頓去.馬其頓的異象,是監獄的異象,真是不易受.但保羅順從神的異象,他真不錯,他永不懷疑.神在異象中的感動,是不會有偏差的.
\newline
\newline
我們的主耶穌,在世上時,是順命接受管制的.主耶穌曾說:「我的時候,還沒有到；你們的時候,常是方便的.」因有管治,主就面對耶路撒冷,面對痛苦,死亡,苦杯也在所不辭；但他仍說,父所賜的杯,我怎能不喝?可見異象在神的兒女身上,是有管制的功用.
\newline
\newline
(b)管治人的說話
\newline
\newline
神的兒女講甚麼?對誰講?都有管治,沒有管治,就發生危險.如火,星星之火,就可以燎原.
\newline
\newline
聖經說,你要講有建設性,造就人的話.在話語上,要受管治.巴不得神幫助我們,使我的話語受約束.
\newline
\newline
有異象的管治,也一定有目的,這目的,一定能對人有好處.
\newline
\newline
(c)聖靈的管治
\newline
\newline
是管治我們所去的地方,有異象的人,神在他前面有帶領.
\newline
\newline
戴永明──戴德生的孫子.在抗戰時,他在西北設有一間神學院.學生畢業時,他將中國地圖放在地上,讓學生們禱告；然後求神引導他站在地圖那一位置,他就到那地方傳福音.其中有一女學生不站,她說:「這地圖上沒有我應去傳道的地方.」戴院長也不勉強她,她一畢業便結婚,別人就誤會他貪愛世界.
\newline
\newline
後來她的丈夫死了,遺下一個女兒.女兒長大後嫁一空軍丈夫,事後被派往台灣,姊妹也隨著女兒一家到了台灣,在台灣做了傳道.戴院長在台灣遇見了她,戴院長對她說:「現在我才明白,你以前不是因為貪愛世界；因為那時,中國地圖上是沒有台灣的.」她在神學院所得的異象,是我們看不見的；在若干年後,才顯明.她在高雄市一間循理會的教會作了女牧師.
\newline
\newline
神引導祂的兒女們,是很奇妙的.請問,神在你的身上的引導是甚麼?
\newline
\newline
我們對主的工作,漠不關心嗎?今日教會的荒涼,世人不信,我們能無動於中麼?
\newline
\newline
西國有一教會,徵求會友奉獻,若奉獻五十元,即可認養一孤兒.有一位年老半身不遂的姊妹,因為窮,拿不出五十元來.她就縫了一張被義賣,但沒有人要,令她很傷心.後來有一位教會領袖聽到了,就把她的義賣心事,向教會宣佈,結果捐款就源源而來.可見只要我們在神面前有異象,終必蒙福的.
\newline
\newline
今天神的工作,到處都需要有異象的人去幹；我們不要自欺,也不要退後,好好地遵守神給你的異象去做吧!
\newline
\newline


\newpage

\section{第44屆~港九培靈會~吉兆頎~第四講　異象的鑑別}
\label{sec:1366}
\textbf{第四講　異象的鑑別}
\newline
\newline
連結: \href{https://www.hkbibleconference.org/session-message/view/1366}{\texttt{https://www.hkbibleconference.org/session-message/view/1366}}
\newline
\newline
\hyperref[sec:1365]{< < < PREV SERMON < < <}
~
\hyperref[sec:toc]{[返目錄]}
~
\hyperref[sec:1367]{> > > NEXT SERMON > > >}
\newline
\newline
耶利米書 23:16-23:16
\newline
\begin{longtable}{cl}
\hline
\hline
章節 & 經文 (和合本修訂版)\\
\hline
23:16 & \begin{tabularx}{0.7\textwidth}{X} 萬軍之耶和華如此說：「你們不要聽這些先知向你們所說的預言。他們使你們成為虛無，所說的異象是出於自己的心，不是出於耶和華的口。 \end{tabularx} \\ \\
[1ex]
\hline
\hline
\end{longtable}
今天要思想有關「異象的鑑別」.我們為甚麼要知道異象的鑑別呢?因為魔鬼會裝作光明的天使,混亂主道,攔阻真理,攔阻信徒.所以耶利米警告當時的猶太人:「萬軍之耶和華如此說,這些先知向你們說豫言,你們不要聽他們的話；他們以空虛教訓你們,所說的異象,是出於自己的心,不是出於耶和華的口.」使徒約翰也告訴信徒:「親愛的弟兄阿!一切的靈,你們不可都信,總要試驗那些靈是出於上帝的不是,因為世上有許多假先知已經出來了.」
\newline
\newline
要鑑別屬靈的事,要注意源頭,是從天上來的,是從神來的抑或從地上,從人來.我們一定要清楚,若不清楚屬靈的事,就會被混亂.
\newline
\newline
(1)原則對照
\newline
\newline
異象有原則.鑑別異象,要以原則對照,原則的標準,是神給人的標準,是所有異象的見證,是所有異象的標準.
\newline
\newline
異象的原則:頭一件事是看見.不論是聖事或俗事,喜歡或不喜歡,必需知道所看見的是甚麼?
\newline
\newline
保羅所看見的,不一定是保羅所喜歡的；彼得所看見的,也不一定是彼得所喜歡的.以賽亞所看見的是聖潔的神,立刻便知道自己的污穢.
\newline
\newline
人所看見的,各有不同,我們不能效法人.因為神會按各人的屬靈光景,和屬靈需要,以及屬靈負擔來給人.
\newline
\newline
(2)聽見.每一個異象都有聲音,因為人的靈不能都通達,所以要有聲音.神藉異象將目的告訴人.神對保羅說:「起來,進城!」對摩西說:「將百姓領出來.」顯見各人的異象,是與各人所領受是有關係的.
\newline
\newline
聖經說:你聽後,存在心裡.那麼就當按神的託付,先領受然後引出來,因為異象是叫人動的.你可以面對神的工作,教會的託付按兵不動；但有了異象,你就不能不動了.
\newline
\newline
原則的對照,在鑑別異象時,是很要緊的.假先知傳出的異象,是出於自己的心,出於私慾,出於捏造.他們的異象,不是為榮耀神,高舉神,乃是為自己.
\newline
\newline
領受異象的人,一定要清楚屬靈的光景如何.若在神前強求異象,則不一定有異象.因為異象是神給人的,不是人人可以勉強神,不是人人可以控制神強賜的,因為人的道路,都由神所控制.
\newline
\newline
若異象不是出於神,那麼就算有異象,這異象定可怕.因為那異象要人做神不喜歡的事,把人的心攪擾紛亂.例如巴蘭,神用驢向他說話,不用異象,因他貪愛錢財,走不義的道路.這種人,若說有異象,也是他捏造的,是假設的,因為神斷不能將異象給他,神的靈也不支持他.
\newline
\newline
下面是幾件事可以對照個人屬靈的光景如何:
\newline
\newline
(A)請問你在神面前,生命是否有長進?與神有沒有生命上的關係,就要先讓祂的生命進入你裡面；這生命乃是復活的生命,是求不死的生命.人若將錄音帶種在花盆,花盆不能使它生出小錄音帶來.但將花的種籽放在花盆裡,就在小花生出來.為甚麼呢?因為錄音帶沒有生命,花種是有生命的.據報載,有一具二千年前的出土屍體,皮肉均完整.雖然完整,只不過是死的僵屍,決不能起來的.
\newline
\newline
主耶穌已復活了,因為他是生命.我們接受了主就與主有生命上的關係.
\newline
\newline
試問神的生命,在我們裡面,是靜止的還是長進的呢?你現今仍然誇張你的好處嗎?天天說自己不配,軟弱嗎?這就是生命沒有長進的象徵.若生命有長進,就必愛慕屬靈的事,拒絕壞的事,接受好的事,樂於與人分享.
\newline
\newline
(B)生活聖潔.若要辨別異象,就一定要在生活上與眾不同.別人隨著時代的潮流轉動,注重享受；但你就要遵守主的旨意,過聖潔的生活
\newline
\newline
聖經說:神造萬物,各從其類.我們是屬那一類的人,乃是屬有神生命的人.聖經又說:不可把兩樣的種子,種在你的葡萄園裡,又說:不可穿羊毛細麻兩樣攙雜料作的衣服.(參申廿二9-11),信與不信的,原不相配,不能同負一軛.正如婦女不可穿戴男子所穿戴的,男子也不可穿婦女的衣服；因為這樣行,都是耶和華你的神所憎惡的.(申廿二5)我們能否遵行神所吩咐的呢?問題是在乎你自己,看你能不能各從其類；是否願意為主而活,想不想過聖潔的生活而已.
\newline
\newline
以賽亞要火炭潔淨他,保羅要將一切看作糞土.所有得異象的人,在生活上,都要分別為聖.
\newline
\newline
今天的世界,信徒要過分別為聖的生活,有許多難處.但如果生命有長進,就沒有難處.屬靈的事,全於生活.比如吃飯,覺得牙齒辛苦,難道就不吃飯嗎?有人來聚會,因懼怕長老,牧師而來,這樣的崇拜有甚麼意義呢?還有人說,是為牧師來做禮拜的,難怪他們的生命不會長進.
\newline
\newline
你在神的家,有沒有工作的負擔呢?你得救後,一年,兩年,三年......八年,一直沒有工作負擔嗎?那就證明你的靈命,是枯乾的.這會令主傷心難過,使教會蒙受損失.各位!工作負擔有二:(一)是個人領受,應盡甚麼本份.(二)是由教會分派.不推辭,必可完成負擔.今天在神家中,沒有負擔的人太多,所以教會有軟弱,失去見證.信徒在家中,在教會內,在工作的地方,都應有屬靈的見證與負擔.
\newline
\newline
(C)眾人的引證.在神的家中,眾人的引證是很重要的.你是否愛神,不是照你心所想,眾人的觀點一定集中在有負擔的人身上.你若真心愛主,真有熱心,用不著口講,別人也知道.你有異象,是真是假,別人將你放在屬靈的生活上一稱,便可知道.
\newline
\newline
若在神前有異象的領受,眾人必會印證出來,因為過去所成就的,大家會為你感謝神；未來的工作,他們會用祈禱支持你,這就是眾人印證的重要性了.
\newline
\newline
(五)異象的引導
\newline
\newline
我們在神面前,有了異象,這異象在你身上有甚麼引導,是很重要的.
\newline
\newline
保羅在馬其頓所看見的異象,「聖靈既然禁止他們在亞西亞講道,他們就經過弗呂家,加拉太一帶地方.到了每西亞的邊界；他們想要往庇推尼去,耶穌的靈卻不許.他們就越過每西亞.下到特羅亞去.在夜間有異象現與保羅.有一個馬其頓人,站著求他說,請你過到馬其頓來幫助我們.保羅既看見這異象,我們隨即想要往馬其頓去,以為神召我們傳福音給那裡的人聽.」(徒十六6-10)這異象的引導,非常清楚.保羅是屬靈的,對異象一點不懷疑,沒有干擾,完成神的異象所託付的.
\newline
\newline
在異象中,保羅只看見一個男人,但他到了腓立比,卻看見一個女人,在河邊禱告.後來保羅又為了一個女孩子趕逐惡鬼竟然被捕入獄,卻因此而領導了獄卒一家信主.保羅完成馬其頓的異象,不是在平坦的路途上完成的,而是在艱苦中完成的.保羅被扣上木狗,受了捆綁,這是異象引導保羅所付出的代價.
\newline
\newline
各位,異象臨到你,引導你,你如何去完成呢?我們必須付出,異象所要我們付出的代價.監牢的門開了,保羅可以逃,但他不逃；可見,不是羅馬人捆綁他,而是異象捆綁他.保羅得釋放,也不是羅馬人釋放他,而是異象釋放他.腓立比的教會,是異象所成立的教會,是有異象的保羅所成立的.所以「從頭一天直到如今,他們是同心合意的興旺福音,保羅深信那在他們心裡動了善工的,必成全這工,直到耶穌基督的日子.」保羅受異象的引導,不惜付上代價,完成神的託付.由古到今,完成異象的人,無一人不完全順服.怪不得保羅說:「那美好的仗我已經打過了.當跑的路已經跑盡了,所信的道我已經守住了.從此以後,有公義的冠冕為我存留.」感謝神!這冠冕為保羅存留,但聖經說也為愛慕主顯現的人存留.
\newline
\newline
保羅面向耶路撒冷,亞迦布豫言保羅必被捆綁,門徒都勸保羅不要上耶路撒冷去,但保羅說:「你們為甚麼這樣痛哭,使我心碎呢!我為主耶穌的名,不但被人捆綁,就是死在耶路撒冷也是願意的.」(徒廿13-14)在神前保羅受異象的引導,是非常清楚的.雖然後來失去自由,但傳福音的機會是沒有減少過.所以他敢站在亞基帕王面前說:「我沒有違背從天上來的異象.」
\newline
\newline
各位,你有沒有從天上來的異象,有沒有打折扣?有沒有遵從?各人要向主負責,因為各人都要向主交賬,求主恩待我們,去完全順服主的異象引導；
\newline
\newline


\newpage

\section{第44屆~港九培靈會~吉兆頎~第五講　異象的教訓}
\label{sec:1367}
\textbf{第五講　異象的教訓}
\newline
\newline
連結: \href{https://www.hkbibleconference.org/session-message/view/1367}{\texttt{https://www.hkbibleconference.org/session-message/view/1367}}
\newline
\newline
\hyperref[sec:1366]{< < < PREV SERMON < < <}
~
\hyperref[sec:toc]{[返目錄]}
~
\hyperref[sec:1368]{> > > NEXT SERMON > > >}
\newline
\newline
以西結書 1:4-1:14
\newline
\begin{longtable}{cl}
\hline
\hline
章節 & 經文 (和合本修訂版)\\
\hline
1:4 & \begin{tabularx}{0.7\textwidth}{X} 我觀看，看哪，狂風從北方颳來，有一朵大雲閃爍著火，周圍有光輝，其中的火好像閃耀的金屬； \end{tabularx} \\ \\ \relax
1:5 & \begin{tabularx}{0.7\textwidth}{X} 又從其中顯出四個活物的形像。他們的形狀是這樣：有人的形像， \end{tabularx} \\ \\ \relax
1:6 & \begin{tabularx}{0.7\textwidth}{X} 各有四張臉，四個翅膀。 \end{tabularx} \\ \\ \relax
1:7 & \begin{tabularx}{0.7\textwidth}{X} 他們的腿是直的，腳掌好像牛犢的蹄，燦爛如磨亮的銅。 \end{tabularx} \\ \\ \relax
1:8 & \begin{tabularx}{0.7\textwidth}{X} 在四面的翅膀以下有人的手。這四個活物的臉和翅膀是這樣： \end{tabularx} \\ \\ \relax
1:9 & \begin{tabularx}{0.7\textwidth}{X} 翅膀彼此相接，行走時並不轉彎，各自往前直行。 \end{tabularx} \\ \\ \relax
1:10 & \begin{tabularx}{0.7\textwidth}{X} 至於臉的形像：四個活物各有人的臉，右面有獅子的臉，左面有牛的臉，也有鷹的臉； \end{tabularx} \\ \\ \relax
1:11 & \begin{tabularx}{0.7\textwidth}{X} 這就是他們的臉。他們的翅膀向上張開，各有兩個翅膀彼此相接，用另外兩個翅膀遮體。 \end{tabularx} \\ \\ \relax
1:12 & \begin{tabularx}{0.7\textwidth}{X} 他們各自往前直行。靈往哪裡去，他們就往哪裡去，行走時並不轉彎。 \end{tabularx} \\ \\ \relax
1:13 & \begin{tabularx}{0.7\textwidth}{X} 至於四活物的形像，就如燒著火炭的形狀，又如火把的形狀。有火在四活物中間來回移動，這火有光輝，從火中發出閃電。 \end{tabularx} \\ \\ \relax
1:14 & \begin{tabularx}{0.7\textwidth}{X} 這些活物往來奔走，好像電光一閃。 \end{tabularx} \\ \\
[1ex]
\hline
\hline
\end{longtable}
有人說,預言,異象,都是一樣東西,但絕對不是的.現在我要分別解釋如下:
\newline
\newline
(1)預言──必須照字面應驗,創世記三章十五節說:「我又要叫你和女人彼此為仇,你的後裔和女人的後裔,也彼此為仇；女人的後裔要傷你的頭,你要傷他的腳跟.」這是預言.到以賽亞時:「必有童女懷孕生子,給他起名叫以馬內利.」直到新約,天使對約瑟說:「她將要生一個兒子,你要給他起名叫耶穌,因他要將自己的百姓,從罪惡裡救出來.」這是預言,必照字面應驗.
\newline
\newline
(2)預表──不照字面解釋,乃按靈意解釋,一樣東西,可能包括幾個預表,一個人同時能預表,另一個人的幾方面.
\newline
\newline
如約瑟坐牢,預表主耶穌受苦；約瑟作宰相,預表主耶穌得榮耀.預表是注重靈意.
\newline
\newline
比方:信徒將來要被提是預言,但聖經裡有許多關於信徒被提的預表.如大災難前被提,以利亞和以諾被神接去,均未經大災難,是信徒大災難前被提的預表；挪亞經過洪水,一家八口得救,是預表選民經過大災難進入新世界.
\newline
\newline
以利亞被神接去,以利沙跟隨以利亞的腳蹤.路遇四十孩童,諷笑以利沙:「禿頭的,上去吧!」以利沙咒詛他們,有熊出來把他們吃掉.預表主將臨前三年半,就是四十二個月,有兩個獸出來,任意而行,這是預表,預表是多方面的.
\newline
\newline
(3)異象──是最具體,而且一定要符合原則,一定要按著各人的託付,屬靈的領受,而有所不同；故此,異象是對個人的.
\newline
\newline
(A)現在看以西結的異象,先看四活物形狀.
\newline
\newline
(1)像人的形狀.四活物有人的形像,且有四個臉,正表示方方正正的人,不是隨隨便便,馬馬虎虎,圓圓滑滑的.信徒在神前,在人前必須維持方方正正的生活.不論世人如何圓滑,信徒內外都要方正,但若沒有裡面的方正,就沒有外面的方正.
\newline
\newline
「活物」也表示有生命的,是方正的生命,是神的生命.主耶穌潔淨聖殿,祭司長和文士並長老進前來.問他說:「你仗著甚麼權柄作這些事,給你這樣權柄的是誰呢?」(太廿一23)但主耶穌要他們回答他,約翰的洗禮,是從天上來的,抑或是從人間來的?他們不敢回答.主耶穌便說:「我也不告訴你們,我仗著甚麼權柄作這些事.」主不彎曲,主是正直的.
\newline
\newline
活物第一個是人的臉,人的臉是代表這個人,印證這個人.
\newline
\newline
當人未老,臉已老時,便有話說這人「未老先衰」.但人已老,而臉面仍如同小孩.便又說,這人「童顏鶴髮」.人臉帶歡笑,你便問,有甚麼喜事?若人愁眉苦臉,你便問,有甚麼不如意的事?人的心裡如何,面上也如何.
\newline
\newline
我們在神的家中,不可忽略一件事,就是要好好做人.可惜,有人實在上了魔鬼的當,因為人生下來,魔鬼就不想我們好好的做人.今日世人有所謂牛年,豬年,羊年......等.說:在牛年生的便屬牛,在豬年生的便屬豬,在羊年生的便屬羊.......魔鬼要我們一生下來,就作牲畜.
\newline
\newline
信徒第一件緊要的是要會做人,不可中魔鬼的詭計.信徒不能台上,又在台下有不同的臉面；在家中是一個臉,在工作場上又是另一個臉.神對我們所要求的,要有正常人的臉；不然,就不能見證神.神要求我們要不變的臉,有喜樂的臉,有可愛的臉,有能表現裡面有神生命的臉.
\newline
\newline
我在台灣,每禮拜要到一村莊去,做家庭禮拜.這村莊百分之九十是天主教徒,只有很少的基督徒；但我們每週都在那村裡有聚會,所以有許多天主教徒認識我.有一天主教徒來對我說:「牧師!我要信耶穌.」我問他為甚麼要信耶穌?他說:「因為做人做夠了.」我說:「怎麼做人做夠了,便要信耶穌?難道信了耶穌不是做人嗎?」他說:「我家裡吵鬧得很,大人有大人吵；小孩有小孩吵,所以我做人作夠了.」各位!他把信耶穌看錯了,以為信主後,就不是作人.事實恰得相反,信主前,我們不會做人；不會做父母,也不會做兒女.信主後,纔真正會做人,這樣的人在神面前,才能滿足神的心.
\newline
\newline
(2)第二個臉是獅子的臉,你若會在神前做人；成為一個活人,成為一個有用的人,神就要你升高,要你知道將來在神的家中,神給你賞賜,給你寶座,給你冠冕,也給你榮耀.使我們信主的人,不再是普普通通,簡簡單單的人,乃成為特特別別的人.特特別別作甚麼人?獅子是獸中之王,表示今日不會作人,將來就不會作王；現在會作人,將來才會作王.甚麼人可與主同作王,一同坐寶座?就是得勝的人.獅子乃獸中之王,乃表示得勝者.
\newline
\newline
信徒在世上,雖被邪惡環繞!但神給我們生命,給我們性情,要過得勝的生活.勝過世界的情慾,與各樣的罪惡.若不能得勝,將來就不能與主同作王.今天我們有許多難處,乃是主給我們有得勝的機會.獅子有兇猛的臉,表示對各樣罪惡不留情面,不願意妥協.主今天叫我們不是逃避罪惡,乃是要做一個得勝罪惡的人.主潔淨聖殿時曾說:「我的殿必稱為禱告的殿,你們倒使他成為賊窩了.」主對罪惡,永不妥協,也永不低頭.所以要與主同作王,先要有獅子的臉,有勝過罪惡的生活.
\newline
\newline
(3)牛的臉,第三個是牛的臉,牛的臉是平衡獅子的臉.信徒想作王,是不錯的,主將來也要你作王.但現今你仍是人,仍在人群中,無論在何處,接觸的都是人.
\newline
\newline
牛的樣子,牛的性情,是忍耐,是犧牲,是服務.下面是一個故事:有一天,一隻牛遇見一隻豬,豬對牛說:「牛大哥,人對你都很尊敬,但對我就不客氣.其實我對人也有不少的貢獻,肉可給他們吃；皮和毛也可給他們用,為甚麼人總是罵我,從不感謝我；對貪睡的人,罵他像一隻豬；對污穢的人,也罵他像豬；對懶惰的人,也罵他像豬,但從不罵他像牛,這是何故?」牛回答說:「豬老弟；你對人有好處,是不錯的.但若與我相比,時間上就有些不同.你是死後對人有貢獻,我是活著的時候便對人有貢獻.」
\newline
\newline
各位!你聽見牛所說嗎?作基督徒的人,是服務的人；在活著的時候,就應有服務的精神,也應有犧牲的精神.人看見你有犧牲的表現,有忍耐的性情,就給人感覺,你真是個基督徒.
\newline
\newline
在台灣有一位姓謝的青年傳道人,他原是一個海員.在船上服務期間,與另一海員同宿一房間.謝弟兄便向他傳講耶穌,勸他信主.但他說:「我知道耶穌是一個大仁,大勇,大愛的人；你是信主的,你應該發揮基督的精神.」謝弟兄問:「怎樣發揮?」他說:「若你每天為我服務,早上為我舖床褥,晚上為我收床褥；我換下來的襪子為我洗......這樣,就是發揮基督的精神了.」謝弟兄聽了,就照他所說的話去做.但一年,兩年的過去,這同事仍不願意信主.後來謝弟兄調到另一船上去.調來與他同房間的,也是一個基督徒.他要求這基督徒為他服務,如要求謝弟兄一般,可是這個基督徒,不但不願意為他服務,反而大罵他一頓.他聽了,立即說:「我知道基督徒有好有歹,有忠心的,有虛假的.倘若我要信耶穌,就要作一個像謝弟兄似的基督徒.」各位,這人願意信主,是因為有一個是牛的臉的基督徒.
\newline
\newline
我們不可只是祈禱說:「主啊!我是你的僕人,你的使女,你的孩子,我要服事你!」但一回到家中,就回復本來的面目,要所有家裡的人都服事你.」
\newline
\newline
陶希聖先生,是一個很有學問的人.他說:「基督徒要忍人所不能忍,要耐人所不能耐.」這就是牛的樣子,牛的工作,也是牛的見證.
\newline
\newline
(4)鷹的臉.後面是老鷹的臉,鷹的臉,表示智慧.這臉是在後面,是看不見的,是隱藏的.
\newline
\newline
弟兄姊妹們!你在神的工作上,在神的家中；雖有人的形狀,有寶座的榮耀,有牛的精神；但如果沒有隱藏的智慧,所暴露的就都是肉體了.
\newline
\newline
一個有鷹臉的人,鷹是高飛的,不受捆綁,也不受攔阻,是超越一切的.一個將自己隱藏的人,也不被罪惡捆綁,不被虛偽捆綁,更不被名譽捆綁.今天教會裡,有人被罪惡捆綁,他有祈禱,有讀經,但常犯罪.還有一種人被尊貴捆綁,罪惡對他沒有辦法；但別人稱讚他,誇獎他,高舉他,這些東西便將他捆綁著了.
\newline
\newline
我們要學一個功課,就是要將自己隱藏,才能為主作工,才能真正活出基督的樣式.這智慧使我不受任何干擾,不受任何捆綁；對上直到神的寶座前,對下能直達到人間.
\newline
\newline
這活物的腿是直的,表明行動正直,一點不彎曲.牛蹄表明能分辨,會分辨所以行動能正直.牛蹄步行,站立,都屬穩健.
\newline
\newline
這活物有人的臉,但沒有人的腿(彎曲)；有獅的臉,但沒有獅的爪(殘忍)；有鷹的臉,但沒有鷹的爪(抓住世界),只有牛的腿.感謝神!有這四個臉和有牛的腿,在神面前,才能行動正直,叫神得榮耀,願我們實實在在像這四活物!
\newline
\newline


\newpage

\section{第44屆~港九培靈會~吉兆頎~第六講　活物的行動}
\label{sec:1368}
\textbf{第六講　活物的行動}
\newline
\newline
連結: \href{https://www.hkbibleconference.org/session-message/view/1368}{\texttt{https://www.hkbibleconference.org/session-message/view/1368}}
\newline
\newline
\hyperref[sec:1367]{< < < PREV SERMON < < <}
~
\hyperref[sec:toc]{[返目錄]}
~
\hyperref[sec:1370]{> > > NEXT SERMON > > >}
\newline
\newline
以西結書 1:15-1:21
\newline
\begin{longtable}{cl}
\hline
\hline
章節 & 經文 (和合本修訂版)\\
\hline
1:15 & \begin{tabularx}{0.7\textwidth}{X} 我觀看活物，看哪，有四張臉的活物旁邊各有一個輪子在地上。 \end{tabularx} \\ \\ \relax
1:16 & \begin{tabularx}{0.7\textwidth}{X} 輪子的形狀結構好像耀眼的水蒼玉。四輪都是一個樣式，形狀結構好像輪中套輪。 \end{tabularx} \\ \\ \relax
1:17 & \begin{tabularx}{0.7\textwidth}{X} 輪子行走的時候，向四方直行，行走時並不轉彎。 \end{tabularx} \\ \\ \relax
1:18 & \begin{tabularx}{0.7\textwidth}{X} 至於輪圈，高而可畏；四個輪圈周圍佈滿眼睛。 \end{tabularx} \\ \\ \relax
1:19 & \begin{tabularx}{0.7\textwidth}{X} 活物行走，輪子也在旁邊行走；活物離地上升，輪子也上升。 \end{tabularx} \\ \\ \relax
1:20 & \begin{tabularx}{0.7\textwidth}{X} 靈往哪裡去，活物就往哪裡去；輪子在活物旁邊上升，因為活物的靈在輪中。 \end{tabularx} \\ \\ \relax
1:21 & \begin{tabularx}{0.7\textwidth}{X} 活物行走，輪子也行走；活物站住，輪子也站住；活物離地上升，輪子也在旁邊上升，因為活物的靈在輪中。 \end{tabularx} \\ \\
[1ex]
\hline
\hline
\end{longtable}
昨天我講活物的形狀,今日續講活物的行動.
\newline
\newline
現在我先講活物的翅膀,然後看活物的行動.一章十五節前說到活物有四個翅膀,分為兩對,其功用有二:上兩翅彼此連接,下兩翅用來遮身.信主的人,不但會超越,且願意與別人來往；這秘訣就是能隱藏自己,不顯露我們的才幹,也不顯露人的肉體,遮住自己之後,才能與別人有交通,這是最高的原則.四活物雖然能騰空,雖然能到神的寶座前；但必須要隱藏自己,謙卑自己,約束自己.許多人在神面前屬靈的光景不壞,生活不差,見證也不低,但不能約束自己.愈屬靈,愈能愛主；愈是超越,愈是要約束自己.但願我們也有這四個翅膀,對別人,對自己都有益處.
\newline
\newline
(1)活物有形狀,也有行動.「我正觀看活物的時候,見活物的臉旁,各有一輪在地上.」正觀看活物的時候,最奇妙的,有輪子在活物的臉旁.這活物先有四個臉,才帶動輪子.主的教會在地上,正像輪子,有主的託付,有主的工作,有主的引導,有主的能力.說是主的甚麼,就是甚麼；說是主的尊貴,就是尊貴；說是主的榮耀,就是榮耀；說是主的屬天,就是屬天；說是主的豐富,就是豐富.這輪子,就是主的一切.臉面是代表這個人,說出這個人.能面對神的人,才能擔當得起神的託付.能將神表現出來,才能看見神,才能帶動神的託付,帶動神的教會,帶動神的負擔,這輪子完全由臉帶動的.
\newline
\newline
我們在主面前,如有屬靈的恩賜,好像活物之臉,在神面前是方方正正的.請問你這方正的臉,完全的恩賜,有沒有帶動主的工作?有沒有帶動主的託付?有沒有帶動主的吩咐?如果沒有,這臉就不能印證主的工作,不能印證主的吩咐,這臉不過是白白有的臉,徒然存在的臉.
\newline
\newline
一個基督徒在神家中,能愛主,能奉獻,能追求；但若果不能工作,就是美中不足,因為不能帶動這輪子.
\newline
\newline
作主的工作,不能說有亦可,沒有亦可；做亦可,不做亦可.活在神前,要表明活物的形狀,也要表現活物的行動.將活的自己的表明出來,你必定積極的去傳福音.保羅說:「所以我今日向你證明.你們中間無任何人死亡,罪不在我身上.因為神的旨意,我並沒有一樣避諱不傳給你們的.」(徒二十26-27)保羅傳福音非常積極,所以在神前,在人前都坦然無懼.怪不得他敢說:「那美好的仗我已經打過了,當跑的路我已經跑盡了,所有信的道我已經守住了；從此以後,有公義的冠冕為我存留.」
\newline
\newline
我們在神面前,不要靠自己,要靠神的力量；要用輪子的力量,去做神的工.有神的負擔,便忠心這負擔,但不可靠自己.若靠自己,就減輕這負擔,拆毀這負擔；若靠神的力量,就加重這負擔.不可看環境,不可看人,也不可看自己,不可看影響,只看在神前所領受的託付.
\newline
\newline
輪子要快,講究效力.主來的日子近了.如不努力,怎能在主面前交賬.求主幫助我們,像輪子一樣快速的,有力的,積極的為主工作.
\newline
\newline
(2)輪與活物相連.活物如何,輪子也如何,輪子怎樣,活物也怎樣.輪子順從活物,活物屬於輪子,輪子不能無活物,活物也不能無輪子.若活物沒有輪子,就失去平衡；輪子沒有活物,也失去見證.在神的家中,請問你有沒有與輪子相連接?有沒有與主的教會,與主的託付,與主的工作相連接?有沒有分不開的心志?有沒有分不開的見證?如果有的話,就是與輪子脫節,只是一隻活物而已,但失去見證,失去託付,失去了影響力.
\newline
\newline
神知道只有以西結能在被擄的人中間,帶動輪子,且能有效的前進.神需要以西結,因為以西結真能帶動輪子.我們在教會裡面,要有分不開的關係,要忠心於現今的託付.
\newline
\newline
有一個信徒,已經參加別個教會和工作,但每月仍將十分之一寄回母會.有一次我們在執事會中,曾提出這問題.各執事以為這弟兄做得好,應加以鼓勵,因有愛母會的心.可是我力反眾議,以為他身在那地的教會,應與那教會有切實的連繫,十分一應奉獻與那教會.若另有負擔,才將另一些錢寄來.因為我們離他遠,無法帶動他；那教會無連繫,也無法帶動他.結果,日子一久,他慢慢冷淡了,輪子慢慢不動了.故此,我們是不能與輪子脫了節的.
\newline
\newline
(3)輪比活物大.看起來很不相稱,但物是活的；輪子雖大,活物仍能把他帶動.
\newline
\newline
許多人覺得教會的工作太大,太重要.教會如活物,主要將託付給活物,主要將大的工作放在活物身上.主不將教會交給別人,只交給忠心的,活在神前的活物身上.因為除了活物以外,沒有東西能帶動這輪子,也沒有人有資格帶動這輪子.
\newline
\newline
(4)輪子的顏色,如水蒼玉.玉是屬天的顏色；是屬天的,是屬靈的,不是屬地的；不是跟從世界一同滅亡的.主的教會是屬天的,是聖潔的.主的教會只有一個顏色,是屬天的顏色.主的教會,只有一種性質,是屬靈的；只有屬天的人,屬靈的人有資格在教會內.誰能帶動輪子?乃是屬天的人,是屬靈的人,是活在神前的活物.
\newline
\newline
這顏色也說明主的教會是尊貴的.主耶穌在世上時,除拯救人外,無其他工作；神在地上除教會外,也設有其他工作.神的教會只有一種顏色,是屬靈的,屬天的,聖潔的,永遠的.保羅說,這教會是萬世以前隱藏的奧秘,直到主來.彼得揭開這奧秘,主將自己顯給眾人,就知道教會不是卑微的,不是屬地的,不是屬世界的組織；而是屬天的,屬靈的,屬神的引證.有許多人想把教會抹上其他顏色,但不論塗上甚麼顏色,還是徒勞無功.因為這顏色是主定的,若你要塗上其他顏色,主就要塗抹這人.所以我們在教會裡面,該當如此忠心.
\newline
\newline
(5)輪子的樣式是一個的.主的教會在地上,在宇宙中,也只有一個；卻分別在各地出現.教會雖有不同的名稱,不同的組織；但在主看來,仍是一個工作,一個目的,一個見證,一個教會.聖經說:「一主,一信,一洗,一神.」雖然有人想把教會變成多個,但屬神的教會,只有一個.所以在神的家中,只有一個負擔,一個目的,一個工作,一個方針；不是分門別類,不是貶抑別人,高抬自己.
\newline
\newline
有許多基督徒,不是活物的樣式,他們信主,只不過在水中拖一拖,這輪子是不動的.
\newline
\newline
(6)輪中套輪.許多輪子套起來,活物的工作,由小而大,由大而更大.先是耶路撒冷,漸次猶太全地；而至撒瑪利亞,直到地極.這是輪中套輪.
\newline
\newline
請問今天你在教會裡面光景如何?五年前如何?五年後又如何?是輪中套輪,抑或脫了一個,又脫了一個.
\newline
\newline
今天,主將工作交給你,是由小而大,抑由大而小.感謝主的恩典!有許多教會,先是一個教會；一個又一個,由一個變為許多個教會.一個小工作,被活物帶動,變成工作中還有工作；不是空空洞洞的輪子,而是輪中套輪.
\newline
\newline
彼得一動,輪子就套起來,三千人套上五千人.保羅一動,教會也套起來.一個,兩個,三個教會接續產生；一個,兩個,三個同工,接著手攜著手；一個,兩個,三個屬靈的兒子,相繼生產,這是輪中套輪.巴不得幫助我們能帶動輪子,輪中套輪；不然,教會便愈來愈縮小了.
\newline
\newline
(7)輪子可向上行走.由下而上,一直高升,不停的向上,一直彰顯主的教會,主的工作,在地上如何,將來在天上也如何.工作是從地上開始,地上的工作是我們的工作.天上的是神的工作.你在地上所捆綁的,將來在天上也要捆綁；你在地上釋放的,將來在天上也要釋放.神的工作是由地上帶動到天上.主耶穌說:「願你的旨意行在地上,如同行在天上.」
\newline
\newline
教會的工作,是屬靈的工作,完全用屬靈來帶動.我這人裡面如何,屬靈的負擔如何,工作的結果就如何.
\newline
\newline
神在宇宙中,有一中心,就是以色列；在以色列中是耶路撒冷.在耶路撒冷中有聖殿,聖殿的中心是至聖所；至聖所中有約櫃,約櫃中有施恩座,施恩座上有血；血到了神面前,罪便得赦免,這人就活過來.人的靈要活,要常常活,所以要常常進到施恩座前,讓主血潔淨.工作要有能力,禱告要有能力,生活要有能力,就要靠主寶血.
\newline
\newline
(8)最後,輪子轉向四方.活物在神前受引導,方向不同；但各盡其職,完成基督的身體.方向雖不同,但工作,目的則一.將來在神前的賞賜,也是一樣；所以要各盡其職,只要專心,盡職,輪子就不停止.
\newline
\newline
今天教會不敢向前行,反而開倒車,是上了魔鬼的當.傳福音,要努力,無論得時,不得時,總要專心.
\newline
\newline


\newpage

\section{第44屆~港九培靈會~吉兆頎~第七講　活物的環境}
\label{sec:1370}
\textbf{第七講　活物的環境}
\newline
\newline
連結: \href{https://www.hkbibleconference.org/session-message/view/1370}{\texttt{https://www.hkbibleconference.org/session-message/view/1370}}
\newline
\newline
\hyperref[sec:1368]{< < < PREV SERMON < < <}
~
\hyperref[sec:toc]{[返目錄]}
~
\hyperref[sec:1369]{> > > NEXT SERMON > > >}
\newline
\newline
以西結書 1:22-1:28
\newline
\begin{longtable}{cl}
\hline
\hline
章節 & 經文 (和合本修訂版)\\
\hline
1:22 & \begin{tabularx}{0.7\textwidth}{X} 活物的頭上面有穹蒼的形像，像耀眼驚人的水晶，鋪張在活物的頭頂上。 \end{tabularx} \\ \\ \relax
1:23 & \begin{tabularx}{0.7\textwidth}{X} 穹蒼之下，活物的翅膀伸直，彼此相對，每個活物用兩個翅膀遮住自己；每個活物用兩個翅膀遮住自己 ，就是自己的身體。 \end{tabularx} \\ \\ \relax
1:24 & \begin{tabularx}{0.7\textwidth}{X} 活物行走的時候，我聽見翅膀的響聲，像大水的聲音，像全能者的聲音，又像軍隊鬧鬨的聲音。活物站住的時候，翅膀垂下。 \end{tabularx} \\ \\ \relax
1:25 & \begin{tabularx}{0.7\textwidth}{X} 在他們頭上的穹蒼之上有聲音。他們站住的時候，翅膀垂下。 \end{tabularx} \\ \\ \relax
1:26 & \begin{tabularx}{0.7\textwidth}{X} 在他們頭上的穹蒼之上有寶座的形像，彷彿藍寶石的樣子；寶座的形像上方有彷彿人的樣子的形像。 \end{tabularx} \\ \\ \relax
1:27 & \begin{tabularx}{0.7\textwidth}{X} 我見他的腰以上有彷彿閃耀的金屬，周圍有彷彿火的形狀，又見他的腰以下有彷彿火的形狀，周圍也有光輝。 \end{tabularx} \\ \\ \relax
1:28 & \begin{tabularx}{0.7\textwidth}{X} 下雨的日子，雲中彩虹的形狀怎樣，周圍光輝的形狀也是怎樣。這就是耶和華榮耀形像的樣式，我一看見就臉伏於地。我又聽見一位說話者的聲音。 \end{tabularx} \\ \\
[1ex]
\hline
\hline
\end{longtable}
每個信徒都知道,環境對個人來說,是很重要的.無論是屬靈的,是屬世界的；環境的影響,環境的造就,是大有關係的.有好環境,就有好氣氛；對屬靈的培養,定有長進.如果環境不好,對屬靈的追求,處處都使我們打折扣.
\newline
\newline
我們雖然明白在神前領受的職份,曉得倚靠,但有時咬著牙根,面對現實,還不能完成神給我們的託付.
\newline
\newline
神給你屬靈的託付,你就不要妄想,不會遭遇困難；而且也永不可想,前途會一帆風順的.因為魔鬼會眼紅,會阻撓.但相反的,這樣,反而可以支取神的恩典.
\newline
\newline
以西結在非常惡劣的環境中,完成神所託付,沒有絲毫打折扣.神藉著異象向他表明的是明朗的,是澈底的,是造就的.我們是神的兒女,在此異象中,也能同得造就,這是我們查考異象最大的目的.
\newline
\newline
(1)穹蒼以下的環境.廿二節是說到穹蒼以上活物的光景.一是在人前,另一是在神前.一方面是向人負責,一方面是向神負責.一方面是向神領取分給人,這是對異象不可少的責任,是信徒在神前追求不可少的光景.許多人願意向神領受,但不願把從神領受的分給別人,不願意向教會,神的兒女負責.其實我們不但要負責,且不能打折扣.
\newline
\newline
(A)「活物的頭以上,有穹蒼的形像,看著像可畏的水晶,鋪張在活物的頭以上.」活物的頭通天.人在異象裡,顯出個人屬靈的光景.第一個要顯明的,究得上標準,頭能不能頂天.人犯罪後,天也變了顏色,天門關閉了,成為銅天,鐵天.但異象所見的天,是透明的天,明如水晶的天.
\newline
\newline
人失敗後,天門關閉了.人悔改認罪,追求,天也改變了,明亮了,與神的關係正常了.神可以看見人,人也可以看見神.人可以活在神前,神也可以活在人的中間.每一個信徒,我們的天,都應當是明亮的天.
\newline
\newline
以西結在神前一點不模糊,不混亂,清清楚楚.他傳出來的信息,活在神面前的光景,都是明如水晶.巴不得我們在神面前的讀經,祈禱,奉獻,奔走都是清清楚楚,無論結果如何,全無後悔.
\newline
\newline
(B)活物的翅膀能發聲.翅膀的聲音,不是普通的聲音,不是簡單的聲音,也不是混亂的聲音,都是特特別別的聲音.以西結難以找出足以形容的字句,所以他用「彷彿」,「像」等字眼,以西結在神前的光景太明亮了!這聲音,是神的恩典,沒有打折扣的聲音,不是口講,乃是行為,是遵行出來的聲音,這聲音足以引證以西結這個人.他領受的恩典,都是實實在在的,這是行為的聲音.我不是說不應向主講話,若要向主講話,應與所蒙的恩相配合.
\newline
\newline
我在台灣常常出外講道,因此我有很大的感想.台灣的教會,蒙恩在這裡,行為生活在這裡,聽道卻在那裡.台灣開培靈會很難,聚會時,到的人,一天少過一天.海外的教會,蒙恩在這裡,生活行為在這裡,聽道也在這裡,這是不同的地方.請記住,聽道不是為別人,聲音也不是為別人,永遠是為自己的.求主恩待我們,多發行為的聲音,少發高言大智的聲音.
\newline
\newline
有幾次門徒發聲,被主耶穌立即打回去.彼得說:「夫子,我們在這裡真好,可以搭三座棚,一座為你,一座為摩西,一座為以利亞.」但彼得說了,不知道說甚麼?這聲音馬上被打回去.
\newline
\newline
又一次雅各,約翰對主說:「主啊!你要我們吩咐火從天上降下來,燒滅他們,像以利亞所作的麼?」耶穌立刻責備他們:「你們的心如何,你們並不知道.」也立即被打回去,這聲音根本不該發.
\newline
\newline
再一次,西庇太兒子的母親來對耶穌說:「願你叫我這兩個兒子在你國裡,一個坐在你的右邊,一個坐在你的左邊.」耶穌回答說:「你們不知道所求的是甚麼?」
\newline
\newline
我們不可輕易發言,正(傳五2)說:「你在神面前,不可冒失開口.」感謝神!到了五旬節時,他們發聲,三千人,五千人出來；因為這聲音,是從神那裡出來的.「眾人聽見這話,覺得扎心,就向使徒說,我們當怎樣行?」彼得說:「你們各人要悔改,奉耶穌基督的名受洗.」感謝神!這聲音滿帶著神的權柄,滿帶著聖靈的能力.
\newline
\newline
(a)這聲音好像大水的聲音.大水的聲音,就是說出這聲音是配搭合一的聲音,不是混亂的,這聲音是大有力量的；能衝倒一切,如大水氾濫的聲音,無法抵擋,也無法防範.這是合而為一的聲音,眾口同聲.我們在神面前,若眾口同聲的讚美主,為神作見證；這樣的聲音,真是像大水的聲音,無法抵擋.能造就人,也能幫助人,可以征服人心,因為是同心合意所發的聲音.
\newline
\newline
(b)像神的聲音.滿有神的能力,尊貴,榮耀的聲音；是從神寶座來的聲音,是讚美神,歌頌神,榮耀神的聲音.
\newline
\newline
若要發聲,要帶有神的能力,帶有聖靈的能力；說神所要說的話,禁止說神所不要說的話.有許多基督徒,所發出的是猖狂的聲音,自高自大,隨便批評,定別人的罪；求主憐憫我們,好叫我們曉得如何發聲.
\newline
\newline
(c)像軍隊的聲音.軍隊所發的聲音:(1)是爭戰之聲,衝殺之聲.基督徒要與撒旦爭戰,要與環境爭戰,要面對環境爭戰.我們行為的聲音,就是爭戰的聲音.(2)得勝的聲音.我們在神前作工,不消說要靠神得勝；得勝後要發聲像軍隊得勝發聲一樣,沒有雜聲,聲音宏大.
\newline
\newline
神要我們像軍隊.以西結看見骸骨甚多,極其枯乾.不料,有響聲,有地震,骨與骨互相連絡.又見骸骨上有筋,也長了肉,又有皮遮蔽其上.後來,氣息進入骸骨,骸骨就活了,成了極大的軍隊,是耶和華的軍隊.神要求我們做他的軍隊,為神爭戰,為神發聲,要做得勝的聲音.(3)是快樂的聲音.軍隊是有紀律的,像個大家庭.得勝時,就發聲歡狂的聲音,快樂的聲音,是無法形容的喜樂之聲,因為這是無法形容的喜樂.在神面前,我們要發出感謝的聲音,讚美的聲音,快樂的聲音.
\newline
\newline
(d)穹蒼的聲音.怎樣聽穹蒼的聲音?必須靜下來,才能聽此聲音,這是回聲,這是反射之聲；是從神寶座來,是從穹蒼而來.
\newline
\newline
我們所發的聲音,有一回聲.你祈禱,立即就有見證出來；你祈禱,教會立即就有復興出來；你祈禱,那邊立即就有醫治出來,這是回聲.
\newline
\newline
會發聲的活物,秘訣在乎會聽神的聲音.主耶穌說,我的羊聽要我的聲音,就跟著我走.
\newline
\newline
安靜在主的腳前,常常聽見神寶座所發的聲音.為何他能聽見,因為他的天是敞開的.我們的天也是敞開的,我們的主對我們不是關閉的,是敞開的.
\newline
\newline
(2)穹蒼以上的環境.「在他們頭以上的穹蒼之上,有寶座的形像,彷彿藍寶石,在寶座形像以上,有彷彿人的形狀.」(結一26)
\newline
\newline
(a)看見了神的寶座.這寶座是活物發聲的根源,是活物發聲的原因.因為要讚美神,在祂所造的穹蒼之上讚美神.我們一讚美,就看見神的寶座,還看見了神的權柄,尊貴,榮耀.
\newline
\newline
你頭的上面有沒有寶座?受不受上面的約束?服不服上面的權柄,領不領上面的恩典?做基督徒的,不可任意妄為.要知道權柄是高高在天上,高高在穹蒼之上,在寶座之上.我們天天環繞在這寶座過日子,就不能不榮耀神.
\newline
\newline
台灣有一校園團契,造就了許多青年人.有一次夏令會,分組討論,每組七人,但每組放了八張椅子.他們說,這一張是主坐的,主在我們中間,參加我們的聚會.他們尊主為大,他們若照主的旨意行就好；若不照主的旨意行,就令主坐冷板凳了.
\newline
\newline
我常向青年人講笑,「你家中有多少人?」「有爸媽,有哥哥姊姊,有弟弟妹妹.」「我到你家中,可招待我嗎?」「一定招待.」「用甚麼招待?」「用水果,請你吃飯.」「宰了弟弟給我們吃好嗎?」「不能.」「為甚麼?」「是媽媽生的.」「豬也有媽媽生,牛也有媽媽生,為甚麼不可以宰?」「因為人命關天.」是的,因為他是人,我們要向主負責,主降世就是取了人的形狀,成為人的樣式.
\newline
\newline
(b)看見寶座上的人子.不消說,是做過人子的耶穌,是曾為我們受過苦的耶穌；是現在得榮耀的人,是得尊貴的人.
\newline
\newline
天上的寶座,有為人而預備的.為怎樣的人預備呢?旁邊有虹,表示守約的人.昔日神將虹給挪亞觀看.主坐在寶座有虹,是遵守與人所立的約.主耶穌向你守約,但你有向主守約?在受洗時,甚麼都答應牧師.水一乾,甚麼也沒有了,也不守約了.主是守約的,你不守約,主要審判你.
\newline
\newline
(c)腰中有火.神是烈火,不守約,要受審判,但與得救無關.每一個基督徒的生活,工作,都要交賬.只有兩種人,就是良善,忠心的人,要享受神的快樂；又惡又懶的人,要受審判.巴不得我們都是又良善,又忠心的人,守約的人.
\newline
\newline


\newpage

\section{第44屆~港九培靈會~吉兆頎~第八講　聖殿的異象}
\label{sec:1369}
\textbf{第八講　聖殿的異象}
\newline
\newline
連結: \href{https://www.hkbibleconference.org/session-message/view/1369}{\texttt{https://www.hkbibleconference.org/session-message/view/1369}}
\newline
\newline
\hyperref[sec:1370]{< < < PREV SERMON < < <}
~
\hyperref[sec:toc]{[返目錄]}
~
\hyperref[sec:1371]{> > > NEXT SERMON > > >}
\newline
\newline
以西結書 40:1-40:9
\newline
\begin{longtable}{cl}
\hline
\hline
章節 & 經文 (和合本修訂版)\\
\hline
40:1 & \begin{tabularx}{0.7\textwidth}{X} 我們被擄的第二十五年，耶路撒冷城攻破後十四年，正在年初，某月初十，就在那一天，耶和華的手按在我身上，把我帶到那裡。 \end{tabularx} \\ \\ \relax
40:2 & \begin{tabularx}{0.7\textwidth}{X} 在神的異象中，他帶我到以色列地，把我安置在一座極高的山上；在山的南邊有彷彿一座城的建築物。 \end{tabularx} \\ \\ \relax
40:3 & \begin{tabularx}{0.7\textwidth}{X} 他帶我到那裡，看哪，有一人面貌如銅，手拿麻繩和丈量的蘆葦竿，站在門口。 \end{tabularx} \\ \\ \relax
40:4 & \begin{tabularx}{0.7\textwidth}{X} 那人對我說：「人子啊，凡我所指示你的，你都要用眼看，用耳聽，並要放在心上。我帶你到這裡來，為要指示你；凡你所見的，都要告訴以色列家。」 \end{tabularx} \\ \\ \relax
40:5 & \begin{tabularx}{0.7\textwidth}{X} 看哪，殿外四圍有牆。那人手拿丈量的蘆葦竿，長六肘，每肘再加一掌。他量圍牆，寬一竿，高一竿。 \end{tabularx} \\ \\ \relax
40:6 & \begin{tabularx}{0.7\textwidth}{X} 他到了朝東的門，就上臺階，量這門的門檻，寬一竿；這門檻寬一竿。 \end{tabularx} \\ \\ \relax
40:7 & \begin{tabularx}{0.7\textwidth}{X} 又有守衛房，每間長一竿，寬一竿，守衛房之間相隔五肘。挨著通往殿之門走廊的門檻，一竿。 \end{tabularx} \\ \\ \relax
40:8 & \begin{tabularx}{0.7\textwidth}{X} 他量通往殿之門的走廊，一竿。 \end{tabularx} \\ \\ \relax
40:9 & \begin{tabularx}{0.7\textwidth}{X} 他量門的走廊，八肘；牆柱，二肘；門的走廊通往殿那裡。 \end{tabularx} \\ \\
[1ex]
\hline
\hline
\end{longtable}
以西結所看見的聖殿的異象,與另一先知所看見的聖殿的異象不同.以西結所看見的是特大的異象,因為他所看見的聖殿是特大的聖殿；這聖殿不止一個院子,而且有兩個院子.摩西所建的會幕,所羅門所建的聖殿,都只有一個院子；但以西結所看見的聖殿,是有兩個院子的.
\newline
\newline
這是說出神的家將來怎樣,我們在神的家裡也怎樣.
\newline
\newline
(A)本段經文一開始,提到的年代和年數日子,是很有意義,和很有價值的.因為從開始所看的異象,到聖殿的異象；從活物的異象,進到聖殿的異象,不是接續下來的.時間告訴我們,已經數十年了.以西結書第一章告訴我們,那時以西結已經卅歲,即在被擄後第五年.以西結現時年已近半百.從異象看見他作祭司已二十年.一個作祭司二十年的人,已經很有經驗,很老練了.神不能不把未來的計劃託付他.每個神的兒女,在神的家中,你的忠心怎樣,神的託付也怎樣.神要我們所經歷的,我們便要順服跟上這經驗.跟得上,就可以看見後面所走的路起伏不平；所以未來的路,也是起伏不平.在神寶座前遵守神的託付,就不怨天尤人,因為神要我們經歷.神要我們所經歷的,總是對我們有益,總要使用我們,讓主彰顯出來.
\newline
\newline
以西結首先看見的,是他自己的光景.凡是異象,一定說出自己的光景來,將自己的光景顯明.神說以西結已五十歲,作祭司已二十年.他在神面前有很重大的負擔,在此時,神又加重他的負擔；神就給他更大,更美的異象.
\newline
\newline
二十五年是說個人的光景,正月初十是說出異象的日期.這日期是神命令人唯一選擇的日子.正月初十是以色列人出挨及,神命令他們選擇羔羊之時.表示神給他的異象是在羔羊裡,是在基督裡,是在至高屬靈的境界裡,不是在巴比倫低窪之地.
\newline
\newline
這是一個好的開始,所以看見一個更大,更美屬靈的異象.把看見的異象的地方說出來,是在被擄的人中間；做被擄的人的祭司,二十五年環境仍沒有改變；但不同的,因為自己在神面前已經有了長進,神就讓他所領受的與人不同.
\newline
\newline
(B)「在神的異象中,帶我到以色列地,安置在至高的山上；在山上的南邊,有彷彿一座城建立.」見異象之地,是在以色列地,這是神帶他到以色列地,不但到以色列地,且是到以色列的高地.這異象叫我們看見,以西結不是在被擄之地,不是在被擄之人中,不作被擄之人.神將他提到迦南地,提到基督裡,提到屬靈之地,是神應許之地.與第一個異象不同,第一個異象,是在迦巴魯河邊,是在迦勒底之地,是失敗,滅亡之地,是在世界裡,被擄到世界裡.不愛主的人,不忠心的人,失落到巴比倫地去；完全到世界裡去,完全被迦巴魯河潮水衝走了.
\newline
\newline
各位弟兄姊妹!我們在神面前,蒙恩多久?作祭司多久?從前光景怎樣?如今光景怎樣?現在見證怎樣?對照一下,如果完全不同,就要多謝主；如果差不多,就要認罪悔改,不然,是非常可惜之事.
\newline
\newline
在以色列高地,是要平衡那低窪之地；那地是何等卑賤,高地是何等尊貴.在基督裡有長進的人,何等蒙福,應何等忠心為主作工；因為神恩待我們,將我們打入到祂兒子的模型裡,使我們能活出祂兒子的模樣,見證祂的兒子,榮耀祂的兒子.
\newline
\newline
以西結見異象,總脫離不了一件事,就是神的靈降臨在他身上；神的靈帶領他,充滿他,造就他,使用他.無論新舊約,我們必須知道,聖靈工作有兩方面的功用.一是為自己,一是為別人；一個為你的生命,一個為你的工作；一是降臨在你身上,一是充滿你.降臨在你身上,你就得著能力；充滿你,你就有生命,這樣,你便能活在神的面前.以西結看見第一個異象,神的靈降臨在他身上,到了這最後的異象；仍不能越過聖靈的工作,全是聖靈在他身上的工作.
\newline
\newline
巴不得我們外有聖靈能力,內有聖靈的生命；按著神的帶領,接受屬靈的建造,時時刻刻為神負擔；為神工作,忠心的將自己擺上.神不但給你異象,更傳這異象,成為別人的異象.
\newline
\newline
(C)解釋異象.以西結將所看見的異象解釋出來.他看見一座建造的城,這城就是耶路撒冷,因城內有聖殿,所以是耶路撒冷.
\newline
\newline
神要把聖殿建造在一座倒塌的城內,不像樣的城裡；就是說,神要復興這城,建造這城.
\newline
\newline
神把我們擺在神的家中,若有心願,有負擔,不是外面好不好,華麗不華麗.人在神前建造,一切都有,基督不在中間,一切都是徒然.
\newline
\newline
城中有殿,且有一個特別的人,是一個銅人,因為顏色如銅.這是一個受世界審判,又是審判人的人,銅表審判.進入院子後,就是祭壇,是銅的祭壇,用銅包裹.這些銅是有二百五十的可拉黨人,他們反叛摩西,死在耶和華面前；他們才拿的銅香爐,打成一塊銅塊,包在祭壇上,把香燒在壇上,然後將血帶到至聖所.
\newline
\newline
以西結看見的銅人,就是主耶穌；他受審判,是擔當世人的罪,受罪的審判.我們如今能稱為神的兒子,都是他為我們受了審判.這血要帶到至聖所,今天我們成了何等樣人,都因有主的血遮蓋了我們.
\newline
\newline
以色列人從被擄之地回來時,要到這城,就要經過這人；不經過這人,便沒有資格進城.今天活著的人,在神面前的人,有沒有經歷這人,若未經歷,結果是可怕的.
\newline
\newline
「我見殿四圍有牆,那人手拿量度的竿,長六肘,每肘是一肘零一掌,他用竿量牆,厚一竿,高一竿.」(六節)「他帶我到那裡,見有一人,顏色如銅,手拿麻繩,和量度的竿,站在門口.」(三節)
\newline
\newline
(1)人站在門口,誰也不能進去,要進去的人,他都要量一量,看你夠不夠資格進去.
\newline
\newline
有一處地方聚會,規定小孩不可進去,但有一小孩必定要進去,他說:「因為我的媽媽在裡面.」各位!這城將來怎樣進去,誰也不能保證,除非你說:「我認識那為我流血的人.」繩表示你的耐力?有沒有灰心?
\newline
\newline
(2)竹竿.竿上有尺寸.愛不愛主?忠心不忠心?有甚麼標準?許多人常以自己為標準,常以教會裡,最軟弱為標準.說自己怎樣的好,怎樣的屬靈.但我們是要拿神的標準來衡量.
\newline
\newline
六肘,每肘是一肘零一掌,六肘即六肘六掌.亞當是第六天造的,所以我們人是不夠標準,要加上主,成為七肘七掌,才能夠完全.加上主,才能幫助人,在地上做工,在天上也做工.
\newline
\newline
六肘六掌的標準是屬亞當的,是從亞當出來的.祂站在門口,你入,祂看見,你出,祂也看見；你在裡面,你在外面,他都看見；你對神,對人,儘在祂的眼內.
\newline
\newline
你進入教會,祂站在門口；你進入院子,祂也站在門口；你進入天門,祂也在門口,你不能超越他.我們在世上的年日,要多多儆醒,因主是輕慢不得的.
\newline
\newline
「他到了朝東的門,就上門的臺階,量門的這檻,寬一竿,又量門的那檻,寬一竿.」(六節)「其窗櫺和廊子,並雕刻的棕樹,與朝東的門,尺寸一樣,登七層臺階上到這門,前高有廊子.」(廿二節)「登七層臺階,上到這門......」(廿六節)以西結看見有臺階,有七層臺階.七層臺階之上,就是完完全全,聖聖潔潔的神的家.七層階之外,就是世界.我若要在神的家,就要上臺階,若要進入神的家,要登七層臺階.「七」是完全,是要經歷完全之意.許多信徒,說自己軟弱失敗,這是未經歷完全.若多批評,在外邦人則可,在教會裡,就要勒住自己的舌頭.
\newline
\newline
信徒在神的家,第一要想到這是登高,信徒的腳蹤是向上的.
\newline
\newline
要進神的家,要有次序.神的家是有次序的,外邦人不尊敬年長的,信徒不可以；外邦人可以不孝敬父母的,信徒則不可以.神的家有次序,不混亂；進入神的家,要學一個功課,就是守次序.
\newline
\newline


\newpage

\section{第44屆~港九培靈會~吉兆頎~第九講　續「聖殿的異象」}
\label{sec:1371}
\textbf{第九講　續「聖殿的異象」}
\newline
\newline
連結: \href{https://www.hkbibleconference.org/session-message/view/1371}{\texttt{https://www.hkbibleconference.org/session-message/view/1371}}
\newline
\newline
\hyperref[sec:1369]{< < < PREV SERMON < < <}
~
\hyperref[sec:toc]{[返目錄]}
~
\hyperref[sec:1355]{> > > NEXT SERMON > > >}
\newline
\newline
昨天我們看過上臺階,上院子.上了七層臺階,進門,到達院子裡面.這門有一很厚的牆,名「東門洞」.不像一度門,而像一個洞.因為牆又厚又深,所以不像門,而像一個洞.入東門洞,就進到外院,即從臺階進入第一外院.
\newline
\newline
進外院,不同入衛房,這是會幕與聖殿所設有的.但以西結所看見的異象,有兩重院子,可說是「恩上加恩」.不但經歷一個,又再經歷一個；不但享受一個,又再享受一個,這實在是神的恩典.
\newline
\newline
進洞,入院子後,到了衛房.甚麼叫衛房?是保衛聖殿的,是保衛進出聖殿的人.一入外院,就進到衛房,是進入神的恩典去,不但蒙神的恩典,亦蒙神的保守.許多人忽略,以為蒙神的恩典易,但求神保守難.其實賜恩典的主,亦是賜保守的主.蒙神的恩典,進入神的院子,神要保守你.詩人說:「你出你入,耶和華要保守你,從今時直到永遠.」這衛房,是完完全全為愛祂的人所預備的,進來的人,都蒙保守.衛房的尺寸與數目,恰到好處,是為人所預備的.衛房寬一竿,有六肘,共有六間；這說明神的保守,是完完全全的在人身上,施恩典的神,也是保守我們的神.
\newline
\newline
從衛房到鋪石地,不但有衛房,且有鋪石地.不是普通地,在神的院子裡有鋪石地.因人是神用土造的.但神從世界把人揀選出來,人也容易,回到世界去.只有蒙神恩典的人,亦蒙神保守他與世界隔絕,與世人有別.所以基督徒所站的地,有屬靈的立場.我們蒙恩得救成為神的兒女,站立的地是鋪石地.鋪石地唯一的用處是享受神.
\newline
\newline
人獻祭後,有些祭物歸給神,有些祭物歸給祭司；有些祭物歸給獻祭的人,但必須在鋪石地吃.所以要站在鋪石地享受神的豐富,就必須與世界有分別.說得嚴重一點,我們無法將祭物拿到別的地方去吃；我們要享受神的一切,享受神的豐富,享受神的託付,必須站在神兒女的立場上.若失去了立場,就不能見証神,無法進到神面前去享受神.
\newline
\newline
外院有衛房,有鋪石地,又分有外院,內院.經過神的保守,站穩個人的立場,享受神的恩典,但不能停留不動.一定要向前追求,因為前面還有臺階,還有經歷,才可進到內院,有臺階.上外院的台有七層,進內院的臺階有八層.上外院有七層臺階,是表完全進到神的豐富.進內院要上八層臺階,是表復活的數目,進入內院的人,有復活的生命,榮耀的生命,永遠的生命,主自己的生命.
\newline
\newline
上八層臺階,才可進入內院.作基督徒的人,不是有外面屬靈的術語就得；我們要追求,一定要注意屬靈的經歷,一定要步屬靈的步驟.追求如何,在神面前所領受也如何,不能跳越.一跳越,問題便出來了.米利暗,亞倫,可拉黨的人,他們跳越,問題就出來了!
\newline
\newline
你去到教會墳場嗎?有些墳墓修得底漂亮,為的是遮蓋死亡.墳前有十字架,有天使像,遮蓋了死亡.換句話說,死是羞辱,因為死是從罪來的,所以要將死遮蓋起來.
\newline
\newline
進入內院,中央是一個祭壇.進來敬拜神,靠近神,朝見神,一定要經過祭壇；因這祭壇上有祭物與祭物的血,有獻祭人的謙卑與悔改.犧牲是代表人的,是解決人的死.這樣,神才能接納.所以祭壇是在中央,神一切中心是祭壇；這祭壇就是十字架,神的兒子釘在上面.釘死之後,他的樣子是最醜的,甚麼人都不願看,連神也不願看；因為罪人的罪集中在他身上.他擔當了,他代替我們,拆毀了人與神中間阻隔的牆.所以神就接納我們,我們才可以坦然無懼的來到神面前,朝見神.
\newline
\newline
以西結所看見的壇,是一座講究的壇,不是普通的壇.內院是方方正正的,長與寬都是一百肘,壇在中央；這說明了神的愛,神的恩典,祂是不偏不倚的.
\newline
\newline
壇有幾種功用.這壇,上小下大,有四層；壇邊有臺階,下面有供臺,是放供品的.人獻祭後,在神面前永遠有功效.祭壇上有邊,所以祭肉不易掉下來,神要保守獻祭物的人,永遠保守他們.獻祭一方面是為贖罪,一方面是為感恩.一個感恩的人,才真正是蒙恩的人.
\newline
\newline
祭壇上小下大,非常穩固,愈上就愈小；表示愈追求的人,愈升高的人,愈屬靈的人,愈愛主的人,是最小最謙卑的人.聖經說,自卑的,神要將他升為至高；自高的,神要將他降為卑.
\newline
\newline
我們在神面前,若有獻祭,若有感恩,先要有謙卑.自高的人,與祭壇是背道而馳的.以色列人自高,不敬畏神,蒙了神的恩典而不感謝神；蒙神揀選但看不起外邦人,神就將他帶到外邦人中間去受苦.
\newline
\newline
但神已揀選一個人,是一個有追求的人,有忠心的人,願意倚靠神的人,就是以西結.
\newline
\newline
以色列人想蒙恩,想得救,想回去,必要悔改,必要謙卑,才能回去.所以七十年過後,回去的人,都是追求的,謙卑的,流淚的,悔改的.以斯拉,所羅巴伯和其他回國的人,都是謙卑的人,忠心的人,愛神的人,神要使用這些人.他們回去,他們一面作工,一面打仗；不受魔鬼攪擾,直到完成.
\newline
\newline
經過祭壇,然後到神的家.從內院進到聖所有十層臺階,這是說出人的經歷.進外院有七層臺階,進內院有八層臺階,進聖所有十層臺階；即共有廿五層,以西結被擄,此時恰巧是廿五年.
\newline
\newline
各位!你的經歷如何?你的力量也如何；你的追求如何,你的見証也如何.就是經過廿五層臺階,來到神面前.個人追求,何等重要!個人生活,個人追求,個人靈修,個人奉獻,在神家中,何等重要!
\newline
\newline
請問,你蒙恩了多少年?上一了多少層臺階?求神光照我們,讓我們想一想,倘若主今日來,我們是良善?是忠心?抑或是懶惰的僕人?
\newline
\newline
由內院到外殿,就是聖所,要上十層臺階.上了臺階,以西結看見一件東西,這東西是聖殿和會幕都沒有的.聖殿和會幕的聖所裡,有陳設餅桌子,桌上有陳設餅,還有燈臺和香壇.但這裡以西結所看見,是簡化了.「壇是木頭作的.他對我說,這是耶和華面前的桌子.」(結四十一22)以西結說看見一木壇,但又說是桌子,究竟是木壇,抑或是桌子,不如叫它做「壇桌」吧!壇是用來燒香的,桌子是用來擺餅的.燒香是為神,餅是為人,到此地步,可以說,神人已經合一了.
\newline
\newline
弟兄姊妹!桌上有你的供應,也有神的悅納的香壇.神看供應的人,等於悅納燒香的人.你有資格擺上供餅,神就滿足了.
\newline
\newline
以西結所看見的聖殿,比之會幕,比之所羅門所建的聖殿,一切都擴大了.但有一件東西,沒有擴大,那就是約櫃,因為約櫃是預表基督.外面都改變了,一切都會變遷,唯一不改變的就是基督.
\newline
\newline
到神前領受的人,應是發香氣的人.得供應的人,應是滿足神的人.就是說,滿足神的人,才是有資格供應的人.
\newline
\newline
這壇又叫內壇,四層的是叫外壇；從外壇到內壇.內壇外壇是分不開的,外壇是預表十字架.主耶穌說:「你們要背起十字架來跟隨我.」作神的兒女,一切不能越過十字架,因為十字架是對付舊人,彰顯新人.你要進入聖所燒香,先要點燈.燒香是預表禱告,點燈是預表神的話,若燒香不點燈,黑墨墨,你的禱告就有了問題；只會批評,自高如法利賽人.你燒香之前若點燈,有神的話作亮光,你的禱告就會如稅吏,自卑禱告.巴不得我們在神面前會點燈,然後燒香.
\newline
\newline
以西結最後將血帶到施恩寶座前.帶著血才可到神面前來,才可得神的恩典.誰活在聖殿裡,誰就可領受神一切的恩典.「萬事都互相效力,叫愛神的人得益處.」凡到至聖所的人,所看見的都是神的恩典,雖然困難,還是神的恩典；雖然損失,還是神的恩典；雖然失敗,還是神的恩典.一切的一切,都是神的恩典!
\newline
\newline


\newpage

\section{第44屆~港九培靈會~楊濬哲~序}
\label{sec:1355}
\textbf{序}
\newline
\newline
連結: \href{https://www.hkbibleconference.org/session-message/view/1355}{\texttt{https://www.hkbibleconference.org/session-message/view/1355}}
\newline
\newline
\hyperref[sec:1371]{< < < PREV SERMON < < <}
~
\hyperref[sec:toc]{[返目錄]}
~
\hyperref[sec:1372]{> > > NEXT SERMON > > >}
\newline
\newline
感謝主的恩典,第四十四屆港九培靈研經大會,已如期舉行,亦圓滿結束.會前的籌備,會期中的工作,以及善後,都是相當繁重吃力的.培靈會過去的成果,都有賴於會眾,主內同工,同道愛戴,合作,纔能眾志聖城.
\newline
\newline
年來港九地方不寧,治安惡劣,以致各教會聚會,(特別晚堂聚會)大受影響.倘非有飢渴慕義,和堅定不移的信心,是難以肯來參加聚會的.感謝主!培靈會未因此而受影響.反之,由新界港九各地來赴會的,大不乏人.大會期間,適海底隧道正式通車之期,故此,方便了許多由香港方面來赴會的同道.
\newline
\newline
本屆講員,有陳終道牧師.他是星加坡神學院的教授,聖經報的主編；對四福音中的耶穌基督,講解詳盡透澈.吉兆頎牧師,原為台中嘉義貴格會牧師,現為香港生命堂主任牧師；講解信徒的異象,喚醒了不少沉睡的信徒.費述凱牧師,是菲律濱聖經學院教授；主講希伯來書,對認識主耶穌與陳牧師互相呼應；並在路得記一書中,對信徒的信心與愛心,重新衡量.
\newline
\newline
本書蒙陳使華姊妹,蔣國仁弟兄,雷麗端姊妹三位記錄,朱建磯牧師編輯,何統雄先生設計封面.會期間又蒙九龍城浸信會的同工,同道,竭誠幫助熱心招待,統此衷心致謝!
\newline
\newline
第四十五屆大會籌備時,人事略有更動,黃仲凱牧師年老退休,副主席一職,由張有光牧師接替,而英文書記一職,則由金新宇博士接替.更慶幸的得九龍城浸信會主席,黃汝光博士,允參加委辦會工作,與我們同心事奉主!
\newline
\newline
在本書出版時,我已身在美國,然而我的心仍念念不忘,港九培靈會的工作.願主興起更多新血,跟隨前人的腳蹤,忠心事主!阿們.
\newline
\newline


\newpage

\section{第44屆~港九培靈會~費述凱~第一講　路得信心的選擇}
\label{sec:1372}
\textbf{第一講　路得信心的選擇}
\newline
\newline
連結: \href{https://www.hkbibleconference.org/session-message/view/1372}{\texttt{https://www.hkbibleconference.org/session-message/view/1372}}
\newline
\newline
\hyperref[sec:1355]{< < < PREV SERMON < < <}
~
\hyperref[sec:toc]{[返目錄]}
~
\hyperref[sec:1377]{> > > NEXT SERMON > > >}
\newline
\newline
路得記 1:1-1:22
\newline
\begin{longtable}{cl}
\hline
\hline
章節 & 經文 (和合本修訂版)\\
\hline
1:1 & \begin{tabularx}{0.7\textwidth}{X} 士師統治的時候，國中有饑荒。在猶大的伯利恆，有一個人帶著妻子和兩個兒子往摩押地去寄居。 \end{tabularx} \\ \\ \relax
1:2 & \begin{tabularx}{0.7\textwidth}{X} 這人名叫以利米勒，他的妻子名叫拿娥米；他兩個兒子，一個名叫瑪倫，一個名叫基連，都是猶大伯利恆的以法他人。他們到了摩押地，就住在那裡。 \end{tabularx} \\ \\ \relax
1:3 & \begin{tabularx}{0.7\textwidth}{X} 後來拿娥米的丈夫以利米勒死了，剩下她和兩個兒子。 \end{tabularx} \\ \\ \relax
1:4 & \begin{tabularx}{0.7\textwidth}{X} 兩個兒子娶了摩押女子，一個名叫俄珥巴，第二個名叫路得，在那裡住了約有十年。 \end{tabularx} \\ \\ \relax
1:5 & \begin{tabularx}{0.7\textwidth}{X} 瑪倫和基連二人也死了，剩下拿娥米，沒有丈夫，也沒有兒子。 \end{tabularx} \\ \\ \relax
1:6 & \begin{tabularx}{0.7\textwidth}{X} 拿娥米與兩個媳婦起身，要從摩押地回去，因為她在摩押地聽見耶和華眷顧自己的百姓，賜糧食給他們。 \end{tabularx} \\ \\ \relax
1:7 & \begin{tabularx}{0.7\textwidth}{X} 她和兩個媳婦就起行，離開所住的地方，上路回猶大地去。 \end{tabularx} \\ \\ \relax
1:8 & \begin{tabularx}{0.7\textwidth}{X} 拿娥米對兩個媳婦說：「你們各自回娘家去吧！願耶和華恩待你們，像你們待已故的人和我一樣。 \end{tabularx} \\ \\ \relax
1:9 & \begin{tabularx}{0.7\textwidth}{X} 願耶和華使你們各自在新的丈夫家中得歸宿！」於是拿娥米與她們親吻，她們就放聲大哭， \end{tabularx} \\ \\ \relax
1:10 & \begin{tabularx}{0.7\textwidth}{X} 對她說：「不，我們要與你一同回你的百姓那裡去。」 \end{tabularx} \\ \\ \relax
1:11 & \begin{tabularx}{0.7\textwidth}{X} 拿娥米說：「我的女兒啊，回去吧！為何要跟我去呢？我還能生兒子作你們的丈夫嗎？ \end{tabularx} \\ \\ \relax
1:12 & \begin{tabularx}{0.7\textwidth}{X} 我的女兒啊，回去吧！我年紀老了，不能再有丈夫。就算我還有希望，今夜有丈夫，而且也生了兒子， \end{tabularx} \\ \\ \relax
1:13 & \begin{tabularx}{0.7\textwidth}{X} 你們豈能等著他們長大呢？你們能守住自己不嫁人嗎？我的女兒啊，不要這樣。我比你們更苦，因為耶和華伸手擊打我。」 \end{tabularx} \\ \\ \relax
1:14 & \begin{tabularx}{0.7\textwidth}{X} 兩個媳婦又放聲大哭，俄珥巴與婆婆吻別，但是路得卻緊跟著拿娥米。 \end{tabularx} \\ \\ \relax
1:15 & \begin{tabularx}{0.7\textwidth}{X} 拿娥米說：「看哪，你嫂嫂已經回她的百姓和她的神明那裡去了，你也跟你嫂嫂回去吧！」 \end{tabularx} \\ \\ \relax
1:16 & \begin{tabularx}{0.7\textwidth}{X} 路得說：「不要勸我離開你，轉去不跟隨你。你往哪裡去，我也往哪裡去；你在哪裡住，我也在哪裡住；你的百姓就是我的百姓；你的神就是我的神。 \end{tabularx} \\ \\ \relax
1:17 & \begin{tabularx}{0.7\textwidth}{X} 你死在哪裡，我也死在哪裡，葬在哪裡。只有死能使你我分離；不然，願耶和華重重懲罰我！」 \end{tabularx} \\ \\ \relax
1:18 & \begin{tabularx}{0.7\textwidth}{X} 拿娥米見路得決意要跟自己去，就不再對她說甚麼了。 \end{tabularx} \\ \\ \relax
1:19 & \begin{tabularx}{0.7\textwidth}{X} 於是二人同行，來到伯利恆。她們到了伯利恆，全城因她們騷動起來。婦女們說：「這是拿娥米嗎？」 \end{tabularx} \\ \\ \relax
1:20 & \begin{tabularx}{0.7\textwidth}{X} 拿娥米對她們說：「不要叫我拿娥米，要叫我瑪拉，因為全能者使我受盡了苦。 \end{tabularx} \\ \\ \relax
1:21 & \begin{tabularx}{0.7\textwidth}{X} 我滿滿地出去，耶和華使我空空地回來。耶和華使我受苦，全能者降禍於我。你們為何還叫我拿娥米呢？」 \end{tabularx} \\ \\ \relax
1:22 & \begin{tabularx}{0.7\textwidth}{X} 拿娥米從摩押地回來了，她的媳婦摩押女子路得跟她在一起。她們到了伯利恆，正是開始收割大麥的時候。 \end{tabularx} \\ \\
[1ex]
\hline
\hline
\end{longtable}
路得記是舊約時代,一部愛情的故事.青年人特別喜歡這樣的題目.成年人也喜歡聽,盡在不言中.
\newline
\newline
舊約中的路得記和以斯帖記,這兩卷書是論到兩個女子的故事.可見聖經對於女子,頗為重視.
\newline
\newline
本篇故事發生於士師時代.那時以色列民敗壞至極,遠離耶和華,背叛神的命令,各行耶和華眼中看為惡的事.就在那不信的時代中,出現了一個女子,名叫路得,她滿有信心.
\newline
\newline
我們也和路得一樣,都是處於不信的時代中；所以盼望我們能從路得身上,得到寶貝的教訓.這卷書雖然僅有四章,卻具有許多值得我們學習的功課.
\newline
\newline
按理,以色列人有最好的機會,認識神,相信神；殊不知他們的信心反而小得可憐；他們懷疑神,不倚靠神.可是路得給我們外邦人,帶來很大的安慰.她也是外邦人,但她對神的信心,卻遠超過多數以色列人的信心.
\newline
\newline
路得記告訴我們:伯利恆有一戶人家,他們的信心失落；因遇飢荒,不曉得仰望倚靠神,反而離開神賜福之地,投靠外邦摩押地.
\newline
\newline
在聖經中,舊約時代的路得,可以作為新約教會的預表.神雖曾向以色列人顯出很多奇跡,為他們行過大事；他們卻快快忘記了神,轉去拜偶像.教會則不然,多數純正的基督徒,從信主至離世,一直愛主.所以教會,能使神得到莫大的安慰.雖然教會也有失敗,退步,冷淡的；但教會中實在有許多寶貝的基督徒,他們甚至不顧性命,為主殉道.
\newline
\newline
世界上熱心愛主的人,有千千萬萬萬,我們可能不知道；但神卻知道,他們給神極大的安慰.路得,這位姊妹給神的安慰,比伯利恆全城的以色列人更大,堪為教會極好的預表.
\newline
\newline
耶穌曾講過一個比方:(參太廿二章)某國王為他擺設婚筵,打發僕人去,請那些尊貴的人來赴席.豈知被請的人都以種種原因,托詞不來.主人打發僕人再次邀請,豈料竟遭毆打,殺害.於是王大發烈怒,吩咐僕人到街頭巷尾,無論男女老少,貧窮殘廢,孤苦零仃的,一概請來.王還為他們預備華麗的衣服穿上,大家喜氣洋洋共用盛筵.這班卑賤污穢殘缺醜陋的小民,一旦能進王宮,與王和王子同坐席,真是夢想不到的事.所以他們都感恩載德,畢生難忘.至於那班尊貴的客人,本是王的上賓,他們卻辜負了王的厚恩；雖然後來有小部分來參加,只不過是出於勉強,他們的臉上好像失去了光彩.這個故事,說明了神愛世人,甚至將祂的獨生子賜下.神為祂的選民擺設救恩的筵席,可惜他們拒絕了祂的大恩.於是恩典的門大開,讓我們這些本來與神的國度無份無關的外邦人,得以進到恩典門內.所以,神從我們這外邦人身上,所得的安慰和快樂,大大超過從祂的選民所得的榮耀和快樂.神雖然曾經為以色列人預備流奶與蜜之地,保守他們.為他們預備了好的領袖,使他們勝過仇敵,他們卻忘恩負義.
\newline
\newline
摩押是外邦,拜偶像,風俗極腐化之地.但「耶和華的眼目遍察全地,要顯大能幫助向祂心存誠實的人.」(代上十六9)神是看人的內心,不看人的外貌.祂看見摩押有一個聰明美麗的女子路得,(路得,是朋友的意思,也是美麗的意思)神看路得是個心口一致,誠實的女子.可以說,路得是舊約時代的拿但業.主耶穌論拿但業,說他是真的以色列人,心裡沒有詭詐.路得也是如此.神定意要幫助路得,使路得認識祂,相信祂,蒙祂的恩典.
\newline
\newline
伯利恆有一四口之家,因遭飢荒逃往摩押地,據解經家說,這家父子三人精於某種手藝.那時摩押王此需要這種人才,所以請他們一家人住在王宮的外院,替王作工.以利米勒是家長,這名字是主作王之意,或是耶和華作王之意.但是他卻自己作王,遇到飢荒,他就帶了全家逃到摩押.雖然他是以色列人,是神的選民,但他忘記了神的作為,忘記了他們的祖先,一百多萬人在曠野四十年.神養活他們,讓他們不至飢餓；難道十年的飢荒,神不能養活他一家四口嗎?以利米勒在神的眼光中失去了選民的價值,因為他失去了信心.一個沒有信心的基督徒是毫無價值的.
\newline
\newline
以利米勒非但失去信心,同時也失去了良心.他不但逃避飢荒到外邦之地,還給他的兩個兒子娶外邦女子為妻.基督徒中也有失去良心的.在菲律賓的華僑,有很多這樣的情形,甚至教會中的長老執事,將兒女與不信主的人婚嫁.為了對方是富商,利用自己兒女婚嫁,圖擴展生意,並沒有考慮到信與不信的問題.
\newline
\newline
神使萬事互相效力,為要叫愛祂的人得益處.後來,以利米勒和他兩個兒子都死了.神准(不是吩咐)以利米勒帶著妻兒逃往外邦,寧可犧牲祂的選民父子三人,為要得著外邦女子路得；叫路得認識祂,叫她得著永遠的福份.
\newline
\newline
拿俄米的名字,是甘甜的意思.她實在是一位很可愛的老婦人,她的脾氣甜如蜂蜜,她的言語中常常題到耶和華,她的口充滿耶和華的名.她說:「耶和華眷顧祂的百姓......耶和華降禍與我.......」她認識到一切出於耶和華.她的心裡充滿著耶和華,晝夜思想耶和華的話.她常用言語為耶和華作見證,述說耶和華的作為說,亞伯拉罕離開吾珥,聽見神的聲音,並大大賜福給他.她講以撒的故事,雅各的故事.約瑟他們如何的下埃及,作宰相,神如何聽他的禱告.神又帶領祂的百姓出埃及,過紅海,在曠野四十年,吃嗎哪.後來又帶他們過約但河,打敗他們的仇敵七族,終於到了流奶與蜜之地.
\newline
\newline
俄珥巴和路得,經常從她們婆婆拿俄米口中,聽見耶和華的作為,心裡漸漸地受感動,衷心敬愛她們的婆婆.特別是路得,她開始認識耶和華,覺悟自己原來所拜的是偶像,是骯髒的假神,唯獨耶和華是聖潔的真神.
\newline
\newline
有一天,拿俄米聽說神眷顧賜福自己的百姓.她想丈夫和兒子都死了,留下婆媳三人都是寡婦,不如自己回老家去.「拿俄米對兩個兒婦說,你們各回娘家去罷!願耶和華恩待你們,像你們恩待已死的人與我一樣,願耶和華使你們各在新夫家中得平安.」於是拿俄米與她們親嘴,她們就放聲而哭,說,不然,我們必與你回你本國去.(8-10)拿俄米實在是一位好婆婆,她和兩個媳婦之間,已經建立了良好的感情.「拿俄米說,我女兒們哪!回去罷,為何要跟我去呢?我還能生子作你們的丈夫麼,我女兒們哪!回去罷......我為你們的緣故,甚是愁苦,因為耶和華伸手攻擊我.」(11-13)「兩個兒婦又放聲而哭,俄珥巴與婆婆親嘴而別,只是路得捨不得拿俄米.」(14)這裡告訴我們,兩個年青的寡婦,終而作了最後的決定；俄珥巴與婆婆親嘴而別,回本國去拜她原來所拜的偶像,葬身在那永遠沉淪之地.唯路得因信心的選擇,毅然決意跟隨拿俄米.「路得說,不要催我回去不跟隨你,你往那裡去,我也往那裡去.你在那裡住宿,我也在那裡住宿.你的國就是我的國,你的神就是我的神；你在那裡死,我也在那裡死,也葬在那裡.除非死能使你我相離,不然,願耶和華重重的降罰與我.」(16-17)路得具有堅決的心意要跟隨耶和華,永不轉回.她情願離開尊貴富有的家族,拒絕從前所拜的假神,投入神選民的先烈中,進入神的國度,永遠不轉回.
\newline
\newline


\newpage

\section{第44屆~港九培靈會~費述凱~第一講　你我心目中的耶穌,是怎樣的一位耶穌呢?}
\label{sec:1377}
\textbf{第一講　你我心目中的耶穌,是怎樣的一位耶穌呢?}
\newline
\newline
連結: \href{https://www.hkbibleconference.org/session-message/view/1377}{\texttt{https://www.hkbibleconference.org/session-message/view/1377}}
\newline
\newline
\hyperref[sec:1372]{< < < PREV SERMON < < <}
~
\hyperref[sec:toc]{[返目錄]}
~
\hyperref[sec:1374]{> > > NEXT SERMON > > >}
\newline
\newline
希伯來書 1:1-1:14
\newline
\begin{longtable}{cl}
\hline
\hline
章節 & 經文 (和合本修訂版)\\
\hline
1:1 & \begin{tabularx}{0.7\textwidth}{X} 古時候，神藉著眾先知多次多方向列祖說話， \end{tabularx} \\ \\ \relax
1:2 & \begin{tabularx}{0.7\textwidth}{X} 末世，藉著他兒子向我們說話，又立他為承受萬有的，也藉著他創造宇宙。 \end{tabularx} \\ \\ \relax
1:3 & \begin{tabularx}{0.7\textwidth}{X} 他是神榮耀的光輝，是神本體的真像，常用他大能的命令托住萬有。他洗淨了人的罪，就坐在高天至大者的右邊。 \end{tabularx} \\ \\ \relax
1:4 & \begin{tabularx}{0.7\textwidth}{X} 他所承受的名比天使的名更尊貴，所以他遠比天使崇高。 \end{tabularx} \\ \\ \relax
1:5 & \begin{tabularx}{0.7\textwidth}{X} 神曾對哪一個天使說過：「你是我的兒子；我今日生了你」？又說過：「我要作他的父；他要作我的子」呢？ \end{tabularx} \\ \\ \relax
1:6 & \begin{tabularx}{0.7\textwidth}{X} 再者，神引領他長子進入世界的時候，說：「神的使者都要拜他。」 \end{tabularx} \\ \\ \relax
1:7 & \begin{tabularx}{0.7\textwidth}{X} 關於使者，他說：「神以風為使者，以火焰為僕役。」 \end{tabularx} \\ \\ \relax
1:8 & \begin{tabularx}{0.7\textwidth}{X} 關於子，他卻說：「神啊，你的寶座是永永遠遠的；你國度的權杖是正直的權杖。 \end{tabularx} \\ \\ \relax
1:9 & \begin{tabularx}{0.7\textwidth}{X} 你喜愛公義，恨惡罪惡；所以神，就是你的神，用喜樂油膏你，勝過膏你的同伴。」 \end{tabularx} \\ \\ \relax
1:10 & \begin{tabularx}{0.7\textwidth}{X} 他又說：「主啊，你起初立了地的根基，天也是你手所造的。 \end{tabularx} \\ \\ \relax
1:11 & \begin{tabularx}{0.7\textwidth}{X} 天地都會消滅，你卻長存；天地都會像衣服漸漸舊了； \end{tabularx} \\ \\ \relax
1:12 & \begin{tabularx}{0.7\textwidth}{X} 你要將天地捲起來，像捲一件外衣，天地像衣服都會改變。你卻永不改變；你的年數沒有窮盡。」 \end{tabularx} \\ \\ \relax
1:13 & \begin{tabularx}{0.7\textwidth}{X} 神曾對哪一個天使說：「你坐在我的右邊，等我使你的仇敵作你的腳凳」？ \end{tabularx} \\ \\ \relax
1:14 & \begin{tabularx}{0.7\textwidth}{X} 眾天使不都是事奉的靈，奉差遣為那將要承受救恩的人服務的嗎？ \end{tabularx} \\ \\
[1ex]
\hline
\hline
\end{longtable}
本屆是第四十四屆港九培靈研經大會.當大會開始第一屆時,我還不是一個基督徒,我還不知道怎樣接受主耶穌作我個人的救主.到了一九二九年八月,就是港九培靈研經大會第二屆開會那一年,神尋找了我這罪人,把我這迷失的羊領回；使我離開黑暗,進入光明,脫離了撒旦的權勢歸向神.
\newline
\newline
一九二九年八月四日,正是第二屆港九培靈研經大會開會的第四天,神開恩救了我這罪人.那天我還不知道兩年之後,神要差我到中國當傳教士；也不知道有個港九培靈研經大會的聚會,更不知道我將在這大會當講員.這些都是想不到而且測不透神的恩典.
\newline
\newline
你我心目中的耶穌到底是怎樣的一位耶穌呢?為何今晚先要問這問題呢?你我是否過著榮耀神有意義的生活呢?這完全在乎我們知不知道耶穌,也在乎我們認識耶穌到何等程度.
\newline
\newline
(約十四9)主耶穌以憂傷的聲音,對祂的門徒腓力說:「我與你同在這樣長久,你還不認識我嗎?」腓力對耶穌所發的問題,以及他對耶穌的態度和思想,都是錯誤的；他的一舉一動,都把耶穌看錯了.他跟隨耶穌兩三年,還是不認識耶穌.他以為耶穌不過是位普通的先知.他沒有認識耶穌是道成為人身的神,是神的兒子；是道路,真理,生命,是獨一的救主.人若非藉著祂,不能到父那裡去.腓力與耶穌,起居飲食工作,常在一起,他也常常看見耶穌行神蹟奇事,但他不認識耶穌是誰.這叫主耶穌非常的傷心和失望!
\newline
\newline
假使今天晚上,主親自向你我顯現,主要問:「我的僕人費述凱,我與你同在幾十年了,你還不認識我嗎?」
\newline
\newline
我事奉主耶穌已經四十二年了.是否至今我還不夠認識我的主耶穌?如果現在我還不夠認識祂,那麼,我的思想,我的訊息,我對耶穌的態度,我愛耶穌的心；甚至我事奉祂,都會受到影響.
\newline
\newline
認識耶穌,與我們的生活有非常密切的關係.你為人如何?在乎你的思想如何,你思想甚麼,就行出甚麼；習以為常,造成你的人格；你人格如何,都是從你的思想發出的.(箴廿三7)說「因為他心怎樣思量,他為人就是怎樣.......」
\newline
\newline
路加福音十九章記載撒該的故事.撒該是個稅吏,本來他不認識耶穌；但他曾風聞耶穌的事,有的說耶穌良善愛人,並能治病趕鬼.有的說耶穌為人不善,欺驅老百姓.撒該卻不知耶穌實際上為人如何.一天,他聽說耶穌來了,他就爬上桑樹,要看看耶穌是怎樣的人.耶穌來到了,抬頭一看對他說:「撒該快下來,今天我必住在你家裡.」這時候,撒該不但聽見耶穌的聲音,也面對面的見到耶穌.在耶穌未來之前,他對耶穌毫無印像,更不認識,所以耶穌對撒該的為人,並無影響力.當耶穌來到那剎那之間,撒該對耶穌的認識進步很快,他認識:
\newline
\newline
(一)耶穌是無所不能的神.因為他躲在桑樹上,耶穌一到便知,且知道他名叫撒該.
\newline
\newline
(二)耶穌愛罪人.撒該是個稅吏,稅吏是猶太人所憎恨的.他知道自己是個大罪人,耶穌竟然願意到罪人的家一同吃飯,不像那班法利賽人冷眼相看.
\newline
\newline
撒該從耶穌得到亮光,認識祂是無所不能的神的兒子,是罪人的救主.因此,撒該整個的人生就完全改變了.他再不欺侮人了,他也不再敲詐人的錢財了,他對人,對錢財,都有了新的態度,而這新的態度,帶來給他新的生活.
\newline
\newline
認識耶穌,終久會影響我們的生活.相反的,不認識耶穌,我們的人生會受到影響.所以今天晚上,我要特別強調每個基督徒,要多多追求認識耶穌基督.
\newline
\newline
你我對耶穌的思想到底怎樣,夠不夠清楚,夠不夠正確,夠不夠寬大?
\newline
\newline
藉著這六個晚上,盼望可以帶領諸位多多的認識希伯來書,所介紹給我們的主耶穌.
\newline
\newline
今天早晨,我參加了大會的第一堂聚會.我坐在那裡聽陳終道牧師講道,我越聽越覺得希奇,因為他的講的和我所預備的差不多.他多多強調主耶穌的人格,叫我們更深的認識耶穌基督；他引用馬太福音,論到耶穌的溫柔,論到耶穌如何的厭惡,假冒為善的法利賽人.又告訴我們,耶穌如何的注重實際,隱藏自己等品格上的優點.
\newline
\newline
我發現聖靈有同一的感動.陳牧師在新加坡,我在馬尼拉,卻不約而同的預備了偏重於一方面的講章.不過,聖靈雖然給我們同樣的負擔,但是給我們的題材卻不同.陳牧師講到四福音中所介紹給我們的耶穌；而我要講的是希伯來書,所介紹給我們的耶穌.
\newline
\newline
四福音中的耶穌基督是當時的人,在耶路撒冷,拿撒勒,迦百農,加利利,撒瑪利雅各地所看見和聽聞的耶穌.希伯來書的耶穌,是坐在天上所看見所愛的耶穌.兩書同是論到耶穌,使我們從兩方面更多的認識祂,好使我們認識耶穌的智慧,來改變我們的態度；認識祂的思想,以及整個的人生,來改變我們的人生.
\newline
\newline
寫希伯來書的,是位隱名的作者,他是一位使徒.聖靈幫助他,讓他看見舊約聖經有許多美麗的圖畫；就是那位釘十字架,從死裡復活,升上高天,坐在至高者右邊的耶穌.我們若從希伯來書看這位使徒,由舊約聖經上拿出來耶穌的圖畫,我們就會看見父神所看見的耶穌基督.
\newline
\newline
保羅在雅典,向當地的哲學家傳講福音說:「眾位雅典人哪!我看你們凡事很敬畏鬼神,我遊行的時候,觀看你們所敬拜的；遇見一座壇,上面寫著未識之神.你們所不認識而敬拜的,我現在告訴你們.」(徒十七22-23)當時猶太藉的信徒也是敬拜耶穌的,但其中有些人卻不認識所敬拜的耶穌基督,正像雅典人不認識所拜的神一樣.
\newline
\newline
希伯來書所題到的那些猶太信徒,他們認識耶穌不夠清楚；所以他們對耶穌的信心,愛心,熱誠,盼望,都受了影響.他們的信心混雜了懷疑,他們的熱誠漸漸的冷淡,他們的盼望也一天天的失去了；這都是因為他們對耶穌的認識,深度不夠的緣故.所以聖靈用這位隱名的使徒,寫希伯來書,將主耶穌基督從新的介紹,給當時的希伯來信徒.我相信他也要把主耶穌,重新的介紹給廿世紀的信徒.
\newline
\newline
今天晚上我要給你們三個問題,請各位回家一同追求認識耶穌.
\newline
\newline
現在讓我們先把希伯來書的內容,簡單的介紹一下:第一章告訴我們,耶穌比舊約時代的摩西,撒母耳,以利亞等先知更尊貴.又告訴我們,祂比舊約時代幫助以色列民族的天使也更加尊貴.起初希伯來的信徒是很看重先知；也很看重天使的.雖然他們大部分都受洗歸入基督,但他們都不知道耶穌比先知和天使更尊貴.因為他們有猶太教的背景,他們看耶穌不過是一個人,已經釘十字架死了,還不如天使不會死.無論如何,他們看先知和天使都比耶穌尊貴.他們真是不夠認識耶穌!當他們遇到迫逼時,他們便離開基督回猶太教去,盼望摩西可以救他們.聖靈對此無限感嘆,所以感動作者寫希伯來書,證明耶穌實在比舊約時代的任何先知天使更尊貴.我們也很需要這部書信來幫助我們更清楚的認識主耶穌.
\newline
\newline
我所要你做的功課如下:
\newline
\newline
(一)神加給祂的兒子耶穌三樣尊榮(是祂沒有加給先知和天使的)是甚麼?(一1-2)
\newline
\newline
(二)耶穌作了神的甚麼?(有兩樣,請看三節)
\newline
\newline
(三)耶穌所行的三件大事是甚麼?(有兩件是長遠作,天天做,不停止作的.有一件只做了一次就不必再做的.請看三節)
\newline
\newline
耶穌在世上的時候,常為門徒傷心；因為與他們同在長久,他們還不認識祂是誰.若他們果真認識,主就大得安慰了.有一次主問門徒說:「人說我人子是誰?你們說我是誰?」感謝主!彼得被聖靈感動,回答正確:「你是基督,是永生神的兒子.」彼得認識了耶穌後,雖然一度否認主,但他對過去的認識,永不會忘懷,永遠地牢記.主對彼得說:「西門彼得,你是有福的的,因為這不是屬血肉的指示你,乃是天上的父,指示你的.」
\newline
\newline
在這幾天的培靈會中,別人雖不會把主耶穌活活地介紹給你,但聖靈要把主耶穌啟示在我們心裡,叫我們真正的認識祂.並且愛祂,信靠祂,等候盼望祂回來.阿們!
\newline
\newline


\newpage

\section{第44屆~港九培靈會~費述凱~第二講　路得信心的試驗}
\label{sec:1374}
\textbf{第二講　路得信心的試驗}
\newline
\newline
連結: \href{https://www.hkbibleconference.org/session-message/view/1374}{\texttt{https://www.hkbibleconference.org/session-message/view/1374}}
\newline
\newline
\hyperref[sec:1377]{< < < PREV SERMON < < <}
~
\hyperref[sec:toc]{[返目錄]}
~
\hyperref[sec:1378]{> > > NEXT SERMON > > >}
\newline
\newline
路得記 2:1-2:23
\newline
\begin{longtable}{cl}
\hline
\hline
章節 & 經文 (和合本修訂版)\\
\hline
2:1 & \begin{tabularx}{0.7\textwidth}{X} 拿娥米有一個親戚，是她丈夫以利米勒本族的人，名叫波阿斯，是個大財主。 \end{tabularx} \\ \\ \relax
2:2 & \begin{tabularx}{0.7\textwidth}{X} 摩押女子路得對拿娥米說：「讓我到田裡去拾取麥穗，我在誰的眼中蒙恩，就跟在誰的身後。」拿娥米說：「女兒啊，你去吧。」 \end{tabularx} \\ \\ \relax
2:3 & \begin{tabularx}{0.7\textwidth}{X} 路得就去了。她來到田間，在收割的人身後拾取麥穗。她恰巧來到以利米勒本族的人波阿斯那塊田裡。 \end{tabularx} \\ \\ \relax
2:4 & \begin{tabularx}{0.7\textwidth}{X} 看哪，波阿斯正從伯利恆來，對收割的人說：「願耶和華與你們同在！」他們對他說：「願耶和華賜福給你！」 \end{tabularx} \\ \\ \relax
2:5 & \begin{tabularx}{0.7\textwidth}{X} 波阿斯對監督收割的僕人說：「那是誰家的女子？」 \end{tabularx} \\ \\ \relax
2:6 & \begin{tabularx}{0.7\textwidth}{X} 監督收割的僕人回答說：「她是摩押女子，跟隨拿娥米從摩押地回來的。 \end{tabularx} \\ \\ \relax
2:7 & \begin{tabularx}{0.7\textwidth}{X} 她說：『請你容許我拾取麥穗，在收割的人身後撿禾捆中掉落的麥穗。』她就來了，從早晨直到如今，除了在屋子裡坐一會兒，她都留在這裡。」 \end{tabularx} \\ \\ \relax
2:8 & \begin{tabularx}{0.7\textwidth}{X} 波阿斯對路得說：「女兒啊，聽我說，不要到別人田裡去拾取麥穗，也不要離開這裡，要緊跟著我的女僕們。 \end{tabularx} \\ \\ \relax
2:9 & \begin{tabularx}{0.7\textwidth}{X} 你要看好我的僕人正在哪塊田收割，就跟著女僕們去。我已經吩咐僕人不可侵犯你。你渴了，可以到水缸那裡喝僕人打來的水。」 \end{tabularx} \\ \\ \relax
2:10 & \begin{tabularx}{0.7\textwidth}{X} 路得就臉伏於地叩拜，對他說：「我既是外邦女子，怎麼會在你眼中蒙恩，使你這樣照顧我呢？」 \end{tabularx} \\ \\ \relax
2:11 & \begin{tabularx}{0.7\textwidth}{X} 波阿斯回答她說：「自從你丈夫死後，凡你向婆婆所行的，以及你離開父母和你的出生地，到素不相識的百姓中，這些事人都告訴我了。 \end{tabularx} \\ \\ \relax
2:12 & \begin{tabularx}{0.7\textwidth}{X} 願耶和華照你所行的報償你。你來投靠在耶和華－以色列神的翅膀下，願你滿得他的報償。」 \end{tabularx} \\ \\ \relax
2:13 & \begin{tabularx}{0.7\textwidth}{X} 路得說：「我主啊，願我在你眼前蒙恩。我雖然不及你的一個婢女，你還安慰我，對你的婢女說關心的話。」 \end{tabularx} \\ \\ \relax
2:14 & \begin{tabularx}{0.7\textwidth}{X} 吃飯的時候，波阿斯對路得說：「你到這裡來吃些餅，把你的一塊蘸在醋裡。」路得就在收割的人旁邊坐下。波阿斯把烘了的穗子遞給她。她吃飽了，還有剩餘的。 \end{tabularx} \\ \\ \relax
2:15 & \begin{tabularx}{0.7\textwidth}{X} 她又起來拾取麥穗，波阿斯吩咐僕人說：「她即使在禾捆中拾取麥穗，也不可羞辱她。 \end{tabularx} \\ \\ \relax
2:16 & \begin{tabularx}{0.7\textwidth}{X} 你們還要從捆裡抽一些出來，留給她拾取，不可責備她。」 \end{tabularx} \\ \\ \relax
2:17 & \begin{tabularx}{0.7\textwidth}{X} 這樣，路得在田間拾取麥穗，直到晚上。她把所拾取的麥穗打了約有一伊法的大麥。 \end{tabularx} \\ \\ \relax
2:18 & \begin{tabularx}{0.7\textwidth}{X} 路得把所拾取的帶進城去給婆婆看，又把她吃飽所剩的拿出來，給了婆婆。 \end{tabularx} \\ \\ \relax
2:19 & \begin{tabularx}{0.7\textwidth}{X} 婆婆問她說：「你今日在哪裡拾取麥穗？在哪裡做工呢？願那照顧你的得福。」路得告訴婆婆，她在誰那裡做工，說：「我今日在一個名叫波阿斯的人那裡做工。」 \end{tabularx} \\ \\ \relax
2:20 & \begin{tabularx}{0.7\textwidth}{X} 拿娥米對媳婦說：「願那人蒙耶和華賜福，因為他不斷地恩待活人死人。」拿娥米又對她說：「那人是我們本族的人，是一個可以贖我們產業的至親。」 \end{tabularx} \\ \\ \relax
2:21 & \begin{tabularx}{0.7\textwidth}{X} 摩押女子路得說：「他還對我說：『你要緊跟著我的僕人拾取麥穗，直到他們把我所有的莊稼收割完畢。』」 \end{tabularx} \\ \\ \relax
2:22 & \begin{tabularx}{0.7\textwidth}{X} 拿娥米對媳婦路得說：「女兒啊，你要跟著他的女僕出去，免得你在別人的田間受人騷擾。」 \end{tabularx} \\ \\ \relax
2:23 & \begin{tabularx}{0.7\textwidth}{X} 於是路得緊跟著波阿斯的女僕拾取麥穗，直到大麥和小麥收割完畢。路得仍與婆婆同住。 \end{tabularx} \\ \\
[1ex]
\hline
\hline
\end{longtable}
當拿俄米決意歸回本鄉,她和兩個媳婦辭別,我們可以想像到一幅動人的圖畫.他們婆媳三人站在摩押(拜偶像之地)和以色列(神所賜流奶與蜜之地)交界的山邊；各懷著沉重的心情.拿俄米為了兩年青守寡的媳婦著想,竭力催她們回去.俄珥巴終而回摩押地去了,路得和拿俄米仍留在那裡；雖經拿俄米再三催促,但路得決心要跟隨她,不肯轉回.
\newline
\newline
路得作了最後的決定,要跟隨耶和華,這完全是神的工作在她身上；並沒有人勸勉她,鼓勵她,催促她.她堅心的決定,完全是出於聖靈的催促,絲毫沒有人的成分在其中.
\newline
\newline
我在菲律賓工作廿一年,每年參加了很多的夏令會和冬令會.許多青年男女,他們在獻心會時奉獻自己；他們都是很好的青年,他們舉了手,站立起來；表示奉獻給主,要終身跟從主.起初,我們還擔心菲律賓聖經學院,可能不夠容納奉獻的青年就讀；但結果,竟然連一個了沒有!其最大原因,就是很多講師用了催請的方法,致他們迫於情面而舉手.當然不能說他們不誠心,但可惜出於一時的衝動,聚會散了之後,便忘記得一乾二淨.
\newline
\newline
路得所作的決定,實在令人欽佩!當時有很大的力量在催她回去；但唯有神,唯有聖靈,唯有主的話,唯有拿俄米以前向她所作的見證；在她心裡動工,這是主所動的善工.拿俄米沒有幫助路得決定,反而勸阻她的決定；再三催她回去,但她沒有消滅聖靈的感動.這樣的決定,才是有價值的,才能永遠長存.
\newline
\newline
路得的名,到底在主耶穌的家譜內.一個外邦女子,名字竟然列於主耶穌的家譜,存到永遠,多麼榮耀!這都是在於她最後的決定,她決定拒絕原來所拜的假神,揀選以色列的真神耶和華.
\newline
\newline
摩押的神名叫基抹.聖經說,基抹是摩押人可憎惡的神.是非常淫亂,敗壞,充滿邪靈的假神.路得堅決拒絕假神,選擇了又真又活的耶和華.她自己破釜沉舟,不留後路,因為她決意要跟隨耶和華,走向新的道路,永不轉回.我相信今晚在我們中間,有很多路得,已經決定,永不轉回!
\newline
\newline
路得不單選擇了一位真神,她也選擇了一個新的生活；她揀選了一個新的家庭,就是拿俄米小小貧寒在伯利恆的家.她拒絕了她的老家.她揀選了一個新的國度──以色列,拒絕了她的原籍──摩押,拒絕了榮華舒適的家.
\newline
\newline
相傳路得可能是摩押王室的公主,在大衛還未做王之前,掃羅屢次追殺大衛.大衛為策安全,他將雙親送到摩押王那裡請求庇護.由此推測,大衛家和摩押皇族必有親戚關係,因為大衛是路得的曾孫.
\newline
\newline
路得撇下了皇家,而情願揀選拿俄米孤寡貧窮可憐的家.她並不以此為恥,反引以為榮!路得清楚知道(羅十二2)神的旨意是善良純全可喜悅的.拿俄米的家,在人看來是可憐的,但實在滿有神的榮耀,平安,喜樂,聖潔和愛,這比皇家的東西寶貝得多.所以,她並不戀慕老家.皇宮雖然堂皇富麗,皇囿鳥語花香,皇室有的音樂,更有諂媚奉承的官吏僕婢；但其中所充滿,乃是野心,貪婪,妒嫉,誹謗,邪惡的穢史,污蹟.
\newline
\newline
路得那天在摩押邊界的山上,拒絕了一樣東西,也揀選了一樣東西.她對拿俄米說:「你在那裡死,我也在那裡死,也葬在那裡.」(17)路得拒絕了王室宏偉的墳山,而揀選了伯利恆城外牧羊地方小小的墳地.她這樣完全,絕對,不留後路的選擇,甚至天上的神都不能否認她的忠心.這種至死不移的選擇,在神面前是極其有價值的.
\newline
\newline
路得到了伯利恆以後,心中忽然有個疑問,而且疑問與日俱增；她又不敢告訴婆婆,恐怕婆婆憂慮過份!因為她們婆媳二人,赤手回到伯利恆.拿俄米呢,丈夫沒有了,兒子也沒有了,年紀又老又弱,未免心事重重.她們所住的,是以利米勒從前留下的老房子.住的方面雖然沒有問題.可是吃甚麼?用甚麼?這些都成問題.路得覺得自己年輕,應當負起家庭的責任.她心中有疑問,為了體貼婆婆,不敢張揚,自己暗暗流淚.我相信,她一定每天清早,跪下禱告她所信的神耶和華,將心中苦情告訴神.經濟問題,在她們家中固然嚴重,但路得心中感到另有一問題更加重大.她相信經濟問題耶和華必有預備；因為亞伯拉罕稱祂為耶和華以勒,在耶和華的山上必有預備.四十年之久神為整個以色列國負完全的責任,所以她具有極大的信心仰望耶和華.
\newline
\newline
路得心中的問題,是屬靈的問題.在各位心中有否屬靈的問題?抑或僅有經濟的問題?如果你信主以後,心裡只有經濟問題而沒有屬靈的問題；這樣的現象不太好,似乎夠不上一個信徒的資格.路得信主不久,心中已存在嚴重的屬靈問題.她說:「主啊!我揀選了你,我絕對要跟隨你,你是我的神,我要永遠跟隨你.我絕對不再轉回摩押地,我要追求認識你,我要向你表明,我已經撇下一切,跟隨拿俄米到你的國度來.但你未曾向我表示你是否要我.」摩西在律法上曾聲明,摩押人不可入耶和華的會籍直到十代.自摩西宣佈這命令至今最多兩代,還要等兩三百年嗎?我渴望得到一個憑據,表示你悅納我了.」路得每朝在橄欖樹下,為此事流淚禱告.她恆切的禱告,聖靈大大動了善工.
\newline
\newline
第二章記載路得出去找工作.我們中間,有沒有人正在找工作?這種味道,我也嚐過.害怕,找不到工作；盼望,找到工作；因為有了工作才有收入,才可維持生計.路得雖然有此需要,但她找工作目的,與一般人有所不同.她心中迫切的問題是「主是否要我?」她憑著信心出去找工作,目的在於找尋主的憑據；她盼望從耶和華手中,找到憑據,因她是個外邦人.如果耶和華聽她的禱告,帶領她到田間；在收割的人身後拾取麥穗,受到莊稼之主的許可；這就是憑據,就是證明耶和華悅納好.
\newline
\newline
這是路得信心的試驗.
\newline
\newline
請原諒我,現在我又要講到自己和我師母的事.原來我們彼此不相識,她是美國人,我是加拿大人.她生長在中國,因為他的父母在中國傳道多年.主帶領我到中國河南省內地會當傳教士,恰巧和她同一工場.當時中國的風俗守舊,男女間的界線劃分清楚；禮拜堂不但男女分座,中間還隔一道牆.那時的傳道人從不敢向女界著眼,只尊向男界講道.這可知當時保守的情形.至於男子向女子示愛,求婚,更是聞所未聞的事.所以教會對於男女間問題,也就格外嚴格.如不守中國規矩的西教士,立即勒令撤職回國.所以我只好循規蹈矩.那時我實在很需要媒人為我和對方談親事,奈何求助無人!我就寫了封一本正經的信,將我的感動和經驗告訴她；並盼望她經過禱告之後,如有同樣的感動,回信給我.信發出之後,一天天過去都不見回音；我心裡害怕,不安,因為我愛她,盼望得到一個憑據,教我知道不是「一面熱」.(這是北方的俗語,人在冷天睡在石板,上蓋厚皮襖,上熱下冷的意思.)我一直等了十七天,終於收到她的來信了；這時,我萬分緊張,手拿著信發抖,不敢拆開；連忙跑進房間關上門禱告,求主安靜我.她來信說,接到我的信時,正忙於青年夏令會,沒有時間為此事禱告；會畢才尋求主的旨意,最後主用聖經的話,給她一個憑據.看到利百加,當亞伯拉罕差僕人,去為他兒子以撒求親時的故事,她家長問她:「你願意跟隨這人去嗎?」利百加說:「我願意.」
\newline
\newline
路得出去找工作那天,她盼望得到憑據,給她清楚知道耶和華要她.第三節告訴我們:「......她恰巧到了以利米勒本族的人,波阿斯那塊田裡.」這是主的帶領.
\newline
\newline
波阿斯是個很愛主的以色列人,行年五十尚未婚娶,他的母親是喇合；他從小有好的家庭教育,他是當時以色列的紳士.
\newline
\newline
當以色列人過了約但河,要攻打耶利哥城之前,先派了兩個探子去窺探,他們到喇合的客店去打聽消息.喇合是第一個聽見以色列人過紅海的迦南人；她也聽到神用大能的手領以色列人出埃及,在曠野的四十年之久養活他們.她相信耶和華是真神,她和路得一樣有誠實的信心,所以神向她顯大能.她為兩探子保守秘密,幫助他們,用繩子從牆上縋下讓他們逃脫.她要求當他們來圍繞耶利哥城時,記念她和她的家.所以,她全家是唯一存留在耶利哥城；沒有被殺的迦南人,因為她虔誠熱愛耶和華.
\newline
\newline
路得入了波阿斯的田,他向她說了許多安慰的話.「波阿斯對路得說:女兒阿!聽我說,不要往別人田裡拾取麥穗；也不要離開這裡,要常與我的使女們在一處.」(9)波阿斯承認她是來投靠在耶和華翅膀下的,承認她是個真的信徒.所以路得聽了這些話,大得安慰.至此她清楚知道,耶和華悅納她了.
\newline
\newline
如果她沒有出去工作,在收割的人身後拾取麥穗,她就得不到這個安慰.
\newline
\newline
奉勸各位!牧師,傳道,好像是那些收割的人,他們領很多靈魂歸向主.靈魂好像麥穗,牧師,傳道會漏掉好多寶貝的靈魂,如同收割的人漏掉麥穗.基督徒們應學習路得,辛勤工作,在田間拾取,收割的人遺漏的麥穗.
\newline
\newline


\newpage

\section{第44屆~港九培靈會~費述凱~第二講　當時教會的信徒,是怎樣的一些信徒呢?}
\label{sec:1378}
\textbf{第二講　當時教會的信徒,是怎樣的一些信徒呢?}
\newline
\newline
連結: \href{https://www.hkbibleconference.org/session-message/view/1378}{\texttt{https://www.hkbibleconference.org/session-message/view/1378}}
\newline
\newline
\hyperref[sec:1374]{< < < PREV SERMON < < <}
~
\hyperref[sec:toc]{[返目錄]}
~
\hyperref[sec:1375]{> > > NEXT SERMON > > >}
\newline
\newline
希伯來書 3:12-3:19
\newline
\begin{longtable}{cl}
\hline
\hline
章節 & 經文 (和合本修訂版)\\
\hline
3:12 & \begin{tabularx}{0.7\textwidth}{X} 弟兄們，你們要謹慎，免得你們中間有人存著邪惡不信的心，離棄了永生的神。 \end{tabularx} \\ \\ \relax
3:13 & \begin{tabularx}{0.7\textwidth}{X} 總要趁著還有今日，天天彼此相勸，免得你們中間有人被罪迷惑，心腸剛硬了。 \end{tabularx} \\ \\ \relax
3:14 & \begin{tabularx}{0.7\textwidth}{X} 只要我們將起初確實的信心堅持到底，就在基督裡有份了。 \end{tabularx} \\ \\ \relax
3:15 & \begin{tabularx}{0.7\textwidth}{X} 經上說：「今日，你們若聽他的話，就不可硬著心，像在背叛之時。」 \end{tabularx} \\ \\ \relax
3:16 & \begin{tabularx}{0.7\textwidth}{X} 聽見他而又背叛他的是誰呢？豈不是跟著摩西從埃及出來的眾人嗎？ \end{tabularx} \\ \\ \relax
3:17 & \begin{tabularx}{0.7\textwidth}{X} 神向誰發怒四十年之久呢？豈不是那些犯罪而陳屍在曠野的人嗎？ \end{tabularx} \\ \\ \relax
3:18 & \begin{tabularx}{0.7\textwidth}{X} 他向誰起誓，不容他們進入他的安息呢？豈不是向那些不信從的人嗎？ \end{tabularx} \\ \\ \relax
3:19 & \begin{tabularx}{0.7\textwidth}{X} 這樣看來，他們不能進入安息是因為不信的緣故了。 \end{tabularx} \\ \\
[1ex]
\hline
\hline
\end{longtable}
上面的經文,略述當時教會信徒的生活,弊病和問題.因為那些信耶穌的猶太人,有諸多的問題；所以作者心裡,有特別重大的負擔,他盼望能夠更清楚,更圓滿地,將主耶穌介紹給他們；不僅希伯來書的作者,有這樣的負擔,彼得也同樣關心,那些猶太教的信徒.他們的情況較為特殊,和外邦信徒不同；他們本是猶太教,後來信了耶穌,所以問題更多了.
\newline
\newline
以色列人是神的選民,除了他們之外,無論是中國人,加拿大,任何國家的人；聖經統稱之為外邦人,所以我們都是外邦基督徒.
\newline
\newline
希伯來書的對象,是猶太的基督徒.本來他們相信摩西,遵行律法；後來聽了福音,看見使徒們所行的神跡奇事；他們就沒有辦法否認耶穌,是他們所盼望的彌賽亞.他們就受洗,成為基督徒了.
\newline
\newline
彼得常為猶太信徒講道.彼得書信是為那些分散在各地的猶太教基督徒寫的.彼得對他們說:「願恩惠平安,因你們認識神和我們主耶穌,多多的加給你們.」(彼後一2)彼得強調,認識神和主耶穌,能夠得到更多的好處,認識耶穌是多麼重要.他用認識耶穌,多得恩惠平安,為寫後書的開端.在本書結束時,他說:「你們要在我們救主,耶穌基督的恩典,和知識上有長進.」(彼得後三18)
\newline
\newline
認識耶穌,是彼得對猶太信徒最大的負擔.如果他們認識耶穌,一切問題自然也解決了；如果他們不認識耶穌,問題也就越久越多了.盼望你們能多多的認識耶穌,這也是我心裡最大的負擔.
\newline
\newline
希伯來書是一部側重高舉耶穌基督的書信,要將應得的榮耀,都歸給主耶穌.
\newline
\newline
本書作者在第一章,曾列舉耶穌諸般的榮耀.一至二節告訴我們,天父把三種極大的尊榮,加給祂的兒子耶穌,是祂從未加給任何先知和天使的.
\newline
\newline
(1)「就在這末世,藉著祂兒子曉諭我們」,藉著祂兒子,神就把自己完全的榮耀和美德表示出來,讓世界的人得以看見.所以主耶穌說:「人看見我,就是看見了父.」父把小部分的榮耀托付先知叫他們報告.這裡告訴我們,天上的父將祂全部的榮耀和美德,藉著主耶穌基督表現出來.絕對沒有任何的先知,或天使可以說:「人看見了我,就是看見了父.」唯獨耶穌有這樣的尊榮.
\newline
\newline
(2)「又早已立祂為承受萬有的.」神應許亞伯拉罕那流奶與蜜之地,要賜給他和他的後裔,作為永遠的產業.(整個以色列國,正如台灣那麼大,是個很小的國家.)阿剌伯人和埃及人,他們卻否認是神的安排.但他們能否向神提出抗議呢?神應許給耶穌基督的,是「承受萬有」,我們的主耶穌何其偉大,誰能向神發出抗議嗎?宇宙和其中所充滿的,卻屬於祂,祂是主人,因為神早已立祂為承受萬有的.
\newline
\newline
(3)「也曾藉著祂,創造諸世界」(約一3):「萬物是藉著祂造的；凡被造的,沒有一樣不是藉祂造的.」按原文「世界」可以譯成「時代」,主耶穌不但安排了律法時代,也安排了恩典時代,教會時代,將來還有千禧年的國度時代.主耶穌不但創造了萬有,祂還與天父安排時代.
\newline
\newline
我們的主耶穌真是偉大!那一位天使能與祂比較呢?
\newline
\newline
希伯來信徒腦海中的耶穌,實在太小了,是沒有榮耀的；難怪他們想像中的耶穌,對他們生活,不發生效力.如果我們心中的耶穌太小,我們的信心就小,愛心就小,熱誠更小.其實耶穌並不小,獨惜我們的心太窄小,不夠認識祂；我們的眼睛未睜開,還未看見祂真的榮耀而已.
\newline
\newline
第三節首句告訴我們:
\newline
\newline
(1)「祂是神榮耀所發的光輝」好比天父是太陽,耶穌是從太陽發出的光輝,從太陽帶來溫暖,亮光與生命.照樣,主耶穌時時刻刻把天父所有的,帶來供應我們.
\newline
\newline
(2)「是神本體的真像」世界各地都有大的佛廟,也有大禮拜堂.我是剛從菲律賓來的,菲律賓是天主教的國家,有很多天主教堂；內有馬利亞神像,耶穌神像,聖徒神像.從前我在中國,看見有許多佛廟,裡有觀音菩薩神像及其他種種神像；有些很巨大,小孩子看見就害怕.世界上的神像很多,有的是人畫的,有的是彫刻的,那些都非真的神像；宇宙之間,唯有我們的主耶穌,是神的真像.所以我不歡喜我的家有任何耶穌的像,難道人會畫神人會彫刻神嗎?因為這些都是錯誤的.唯有主耶穌是神本體的真像,多麼榮耀.那一位先知或天使敢說他是神的真像,他們都不過是神所造的而已.
\newline
\newline
第三節還告訴我們,主耶穌另一些的尊榮,祂所行的三件大事比祂在世上時所行的神跡奇事更大.祂在世上時,曾叫五千人吃飽,叫瞎子看見,瘸子行走......這裡說的耶穌做了三件更大的事,其中1,3兩件是時刻不停的做,永不下班的做.第2件只做一次,永遠有效.
\newline
\newline
(1)常用祂權能的命令托住萬有,太陽東升西沉,一年春,夏,秋,冬,四季輪轉,太空千萬大小星系,大大小小,千萬的宇宙,各行其道,毫無衝突；都是因祂常用權能的命令,托住萬有,何等偉大的耶穌!我們為祂而犧牲一切,甚至付上我們的性命為祂,也是值得的.如此偉大的耶穌,我們窄小的心地和頭腦怎能容納呢?我們必須放大我們的心田腦海,認識耶穌的偉大.祂若幫助我們,誰能抵擋我們呢?我們倚靠祂,永不失敗.如果我們認識耶穌,我們平凡的人生,將變成榮耀的人生.但願聖靈幫助,叫我們真正的認識四福音,和希伯來書,所介紹給我們的耶穌.
\newline
\newline
(2)「祂洗淨了人的罪」.這比祂托住萬有的工作更大更難,聖經說祂用指頭創造天地.但祂全身被掛在十字架上,流出寶血才能洗淨人的罪.這是耶穌最偉大的工作,是先知所不能作的；因為先知自己有罪,怎能替罪人死?這是天使不能完成的,因為天使沒有肉身,不過是服役之靈,不可能為人死.唯耶穌是無罪的人,祂才有資格擔當人的罪.祂是神,祂的生命,可以代替萬人的生命,祂的生命比萬人的生命更寶貴.祂獻上生命作萬人的贖價.祂是人,有資格為人死.祂是神,祂能叫萬人得救.耶穌釘在十字架上幾小時,最後,祂大聲喊說:「成了」,祂一次獻上祂的身體為祭,完成了永遠贖罪之功.感謝主!我們的耶穌,除祂以外,天下人間沒有賜下別名,我們可以靠著得救.
\newline
\newline
(3)「就坐在高天至大者的右邊」.這裡告訴我們,主耶穌掌握天父交托給祂,天上地下所有的權柄.(馬太廿八18-19)告訴我們,祂怎樣運用這權柄:「耶穌進前來,對他們說,天上地下所有的權柄,都賜給我了；所以你們要去,使萬民作我的門徒,奉父子聖靈我的名,給他們施洗.凡我所吩咐你們的,都教訓他們遵守；我就與你們同在,直到世界的末了.」耶穌運用最大的權柄,幫助我們把福音傳到天邊地極.我們的耶穌,手中有極大的權柄,用來幫助祂的教會,在這時代中把福音傳開.
\newline
\newline
今天晚上,我們僅思想希伯來書一章二至三節,但願聖靈藉著這兩節經文,好像在我們心中開了一個窗戶；讓我們看見從未看見的耶穌,叫我們多多認識祂.
\newline
\newline
盼望晚上大家回家,找個地方,自己跪下；打開聖經讀希伯來書一章二至三節,所介紹給我們的耶穌,再次瞻仰主耶穌的榮美.向祂承認,我們以往對祂的看法太平凡,太軟弱,對祂的認識不夠清楚.
\newline
\newline
我們認罪之後,可用保羅的禱告,作為自己的禱告:「主耶穌基督的神,求你將那賜人智慧和啟示的靈賞給我,使我真知道你.」
\newline
\newline
當時希伯來的信徒是多麼的軟弱,懷疑很多；甚至有的離開耶穌,回到摩西律法那裡去.這和廿世紀的基督徒差不多.因為不夠認識耶穌,信心不夠,聖潔不夠,愛人的心也不夠.我們對主的認識確實不夠!
\newline
\newline
所以我們在主面前,要承認我們的罪.
\newline
\newline
現在我要用何西阿書六章三節來作結束.「我們務要認識耶和華,竭力追求認識祂；祂出現確如晨光,祂必臨到我們像甘雨；像滋潤田地的春雨.」
\newline
\newline


\newpage

\section{第44屆~港九培靈會~費述凱~第三講　路得信心的順服}
\label{sec:1375}
\textbf{第三講　路得信心的順服}
\newline
\newline
連結: \href{https://www.hkbibleconference.org/session-message/view/1375}{\texttt{https://www.hkbibleconference.org/session-message/view/1375}}
\newline
\newline
\hyperref[sec:1378]{< < < PREV SERMON < < <}
~
\hyperref[sec:toc]{[返目錄]}
~
\hyperref[sec:1379]{> > > NEXT SERMON > > >}
\newline
\newline
路得記 3:1-3:18
\newline
\begin{longtable}{cl}
\hline
\hline
章節 & 經文 (和合本修訂版)\\
\hline
3:1 & \begin{tabularx}{0.7\textwidth}{X} 路得的婆婆拿娥米對她說：「女兒啊，我不該為你找個歸宿，使你享福嗎？ \end{tabularx} \\ \\ \relax
3:2 & \begin{tabularx}{0.7\textwidth}{X} 你與波阿斯的女僕常在一處，現在，波阿斯不是我們的親人嗎？看哪，他今夜將在禾場簸大麥。 \end{tabularx} \\ \\ \relax
3:3 & \begin{tabularx}{0.7\textwidth}{X} 你要沐浴抹膏，穿上外衣，下到禾場，一直到那人吃喝完了，都不要讓他認出你來。 \end{tabularx} \\ \\ \relax
3:4 & \begin{tabularx}{0.7\textwidth}{X} 他躺下的時候，你看準他躺臥的地方，就進去掀露他的腳，躺臥在那裡，他必告訴你所當做的事。」 \end{tabularx} \\ \\ \relax
3:5 & \begin{tabularx}{0.7\textwidth}{X} 路得說：「凡你所吩咐我的，我必遵行。」 \end{tabularx} \\ \\ \relax
3:6 & \begin{tabularx}{0.7\textwidth}{X} 路得就下到禾場，照她婆婆吩咐她的一切去做。 \end{tabularx} \\ \\ \relax
3:7 & \begin{tabularx}{0.7\textwidth}{X} 波阿斯吃喝完了，心情暢快，就去躺臥在麥堆旁邊。路得悄悄走來，掀露他的腳，躺臥在那裡。 \end{tabularx} \\ \\ \relax
3:8 & \begin{tabularx}{0.7\textwidth}{X} 到了半夜，那人驚醒，翻過身來，看哪，有個女子躺在他的腳旁。 \end{tabularx} \\ \\ \relax
3:9 & \begin{tabularx}{0.7\textwidth}{X} 他就說：「你是誰？」路得說：「我是你的使女路得。請你用你衣服的邊來遮蓋你的使女，因為你是可以贖我產業的至親。」 \end{tabularx} \\ \\ \relax
3:10 & \begin{tabularx}{0.7\textwidth}{X} 波阿斯說：「女兒啊，願你蒙耶和華賜福。你後來的忠誠比先前的更美，因為無論貧富的年輕人，你都沒有跟從。 \end{tabularx} \\ \\ \relax
3:11 & \begin{tabularx}{0.7\textwidth}{X} 女兒啊，現在不要懼怕，凡你所說的，我必為你做，因為我城裡的百姓都知道你是個賢德的女子。 \end{tabularx} \\ \\ \relax
3:12 & \begin{tabularx}{0.7\textwidth}{X} 現在，我的確是一個可以贖你產業的至親，可是還有一個人比我更親。 \end{tabularx} \\ \\ \relax
3:13 & \begin{tabularx}{0.7\textwidth}{X} 你今夜在這裡住宿，明早他若肯為你盡至親的本分，很好，就由他吧！倘若他不肯，我指著永生的耶和華起誓，我必為你盡上至親的本分。你只管躺到早晨。」 \end{tabularx} \\ \\ \relax
3:14 & \begin{tabularx}{0.7\textwidth}{X} 路得就在他腳旁躺到早晨，在人還無法彼此辨認的時候就起來了。波阿斯說：「不可讓人知道有女子到禾場來。」 \end{tabularx} \\ \\ \relax
3:15 & \begin{tabularx}{0.7\textwidth}{X} 他又對路得說：「把你所披的外衣拿來，握緊它。」她就握緊外衣，波阿斯量了六簸箕的大麥，幫路得扛上，他就進城去了。」 \end{tabularx} \\ \\ \relax
3:16 & \begin{tabularx}{0.7\textwidth}{X} 路得回到婆婆那裡，婆婆說：「女兒啊，怎麼樣了？」路得就把那人向她所做的一切都告訴了婆婆， \end{tabularx} \\ \\ \relax
3:17 & \begin{tabularx}{0.7\textwidth}{X} 又說：「那人給了我這六簸箕的大麥，對我說：『你不可空手回去見婆婆。』」 \end{tabularx} \\ \\ \relax
3:18 & \begin{tabularx}{0.7\textwidth}{X} 婆婆說：「女兒啊，等著吧，看這事結果如何，因為那人今日不辦妥這事，必不罷休。」 \end{tabularx} \\ \\
[1ex]
\hline
\hline
\end{longtable}
今天晚上,我們要特別注意「順服」兩個字.剛才所唱的詩歌「信靠順服」這是我們一生一世必須牢記的,信靠與順服,是基督徒生活中,最重要的功課.有了信靠,有了順服,我們的生活才有喜樂.
\newline
\newline
「順服」是人最難學的功課.雖然這兩個字,連我這個外國人也會寫,可是教人行出來,就不容易了.順服人的意思很難,順服神似乎更難.惟有路得,她甘心無條件的順服神.
\newline
\newline
中文新的聖經,我們常看到「信」字,如(約一11-12):「祂到自己的地方來,自己的人倒不接待祂.凡接待祂的,就是『信』祂名的人,祂就賜他們權柄,作神的兒女.」(約三16):「神愛世人,甚至將祂的獨生子賜給他們；叫一切『信』祂的,不至滅亡,反得永生.」(約五24):「我實實在在的告訴你們,那聽我話,又『信』差我來者的；就有永生......不至於定罪.」新約裡面充滿了「信」字.
\newline
\newline
我很喜歡研究中國字,原來這信字很有意思.左邊是個單企人,右邊是個言字；言字是由「,三口」合成.相信中國古時,有很多小國,每小國都有國王；所以全國有很多位國王,但皇帝只有一位.如果某國王很好,得到皇帝喜歡,就被召入京朝見皇帝；皇帝就用朱紅筆,在那國王的王字上面點了一點,王字就成了主字.所以那一點是代表「主」,就是三位一體的主口裡所出的話.「信」的定義就是:「一個軟弱的人完全倚靠三位一體的主口裡所出的話.」中國字實在很奇妙!你是否一個軟弱的人?這裡有許多主口裡所出的寶貴話語,可以幫助你.
\newline
\newline
當我讀神學的時候,我沒有錢,一切所需都是神為我預備的.雖然可免學費,但每月要繳膳費卅元.曾有一時,我僅有十五元,為了急需,我想把一些書賣給同學；但同學也窮得與我同病相憐.後來,我在馬太六章卅三節找到了神的話:「你們要先求祂的國,和祂的義,這些東西都要加給你們了.」所以我就抓住神的應許,天天為所差的十五元禱告.雖然禱告,可是口袋裡,仍然只有十五元,神試驗我非常厲害.到了最後的一天,我不敢到飯堂進膳,因為院長在那裡收膳費.我連忙跑到宿舍,唯一的盼望,是有人寄信給我內附支票.果然有封我的信,是我妹妹小學時候的老師寄來的；他是個基督徒,他很愛我.信的內容是:親愛的費述凱,不知道為何這一個多月來,每當我跪下禱告,就想你的名字.我就求問神說:「神啊!現在費述凱讀神學,是否你要我送錢給他；如果,求你指示我應給多少?現在奉主的名送上支票,請你用感謝主的心接受.」我打開支票(看,正是十五元,不多也不少,真奇妙!我這軟弱的神學生,信靠主的話,主果然為我成全了!
\newline
\newline
「信」的意思,也是「靠」的意思,而且包含順服的成分；正如我們剛才所唱的詩,「信靠順服,以外並無別法；若要得耶穌喜愛,惟有信而順服.」「信」,非單腦子贊成,點頭稱是.比方我對人說,某航空公司的飛機,速度快,又準時；食物好,服務週到,升降平穩.那人邊聽邊點頭,表示頭腦同意我的話,但他不一定敢乘搭那飛機,因害怕飛機在空中失事.他信而不靠,怕飛機靠不住.所以新約裡的信,要加上倚靠和順服,才是真正的信.
\newline
\newline
很多年前,在菲洲有位傳教士,想把新約譯成當地土語,但找不出合宜的字來譯信字.他只好把凡有「信」字的地方,留下了空白；經過兩年千思萬想,總是沒有辦法.他就把這事求告神!一天,有個土人替他從海口搬運了大批書信來到,汗流滿額,疲乏非常,那土人就往籐椅一躺,說:「啊!我這一百八十磅全身重量躺在這椅子上,多麼舒服呀!傳教士聽見了,恍然大悟!於是將所有「信」字的空白,都填上這一句話.如(約三16)「神愛世人,甚至將祂的獨生子賜給他們,叫一切將(全身躺在耶穌身上的人)；不至滅亡,反得永生.」
\newline
\newline
希臘文新約聖經,多半是用匹,斯,丟,偶(國音)四個音階.這是希臘文的「信」字,也有其他的字,不過這是最重要的.這個字不但有倚靠的意思,也有是順服的意思.如保羅在(提後一12)說:「......因為知道我所信的是誰?也深信祂能保全我所交付祂的,直到那日.」這個「信」字就是匹,斯,丟,偶,其意思義就是完全交托,不靠自己,將最重要的交托給祂；將自己的重擔,全都躺在祂身上.同一個希臘字也有順服的意思.請看(腓一29):「因為你們蒙恩,不但得以信服基督,並要為祂受苦.」中文把信的原文譯成「信服」,這兩字譯得很好,非常恰當,比(約三16)的「信」字,譯得更為圓滿.「信」,如果單單頭腦信,沒有用；看你服不服,便知道你信不信.如果「信」沒有「順」,那個信就有問題了.正如飛機,飛行迅速；但你不乘搭,對你仍然無用.耶穌的救恩預備好了,你如果單單頭腦相信,單聽美麗的聖經故事；你不倚靠,你不接受,你就滅亡.
\newline
\newline
中文的國語聖經,翻譯至今五十三年,在這些年間,中外各地言語的習慣,改變了很多.五十三年前譯聖經時,信字的含意和現在已經兩樣了.那時的「信」,是對朋友的態度,忠心信實對待朋友.信耶穌,對耶穌忠心信實,這是我們的信,充滿了意義.但可惜直到現在,已成了頭腦佩服,不是信靠,更不是信服了.
\newline
\newline
路得的信,因信而順服,具有這兩種的成分.她完全信靠神,她出去找工作,神為她開路.她恰巧到波阿斯的田工作,神感動波阿斯；歡迎祂,安慰她,她才曉得；不是單單她要耶和華,耶和華也要她了.所以,路得信心的試驗,有了很好的成果.
\newline
\newline
在第三章,我們看到路得需要學習一個新功課,是她在摩押地所沒有機會學到的.她是一位公主,經常有人侍候她,不必她去侍候別人.耶和華要她學習順服,這是一門最難的功課.
\newline
\newline
記得我兩個孩子小的時候,我吩咐他們做事情,他們不大甘願.可是因為爸爸大於他們,手也大,可以管教,所以只好勉強去做.但他們做,用自己的辦法,把自己的成分加入；不肯百分之百順服爸爸的話.順服看得見的父母,順服看得見的主人尚且不容易；要順服我們看不見的神,更是難乎其難了.路得卻學到了這個功課.
\newline
\newline
昨天晚上我已經告訴各位,在我應當結婚的時候,因遠離父母,只好請中國的郵局幫我做媒.路得的婆婆很關心媳婦的事.我想,拿俄米一定為路得的問題禱告多日了.以利米勒拿俄米有位至近的親屬,名叫波阿斯；他還沒有結婚,又富裕,又愛主；對待路得一向很好,對拿俄米也很客氣.但他從未表示他的心意.按理,他應自動出來盡至親的本份,負起拿俄米和路得生活的責任；將她們的產業贖回,免得本支派的產業流於外人之手.這是猶太人的規矩.
\newline
\newline
波阿斯之所以毫無表示,並不是因他缺少金錢,也不是缺少愛路得的心；既然如此,為何他不負至親之責,為何讓路得辛辛苦苦在他田裡,拾取麥穗維持生活呢?拿俄米為這些問題左思右想,不解其意.後來在禱告中,從神那裡得到了亮光.
\newline
\newline
波阿斯實在深愛路得,可是因自己的年齡和路得相差太遠.路得年輕貌美,盡可找位年青有為的如意郎君,過幸福的日子.波阿斯對路得雖愛慕在心,卻不敢顯露於形,更不敢有所表示.因為他不願將自己的快樂建築在路得的痛苦上.孰知路得自從那天到波阿斯的田地拾取麥穗,發覺波阿斯為人德高望重,暗地裡已經愛慕他了.
\newline
\newline
拿俄米蒙聖靈光照,從路得的言語和表情,看出祂對波阿斯一見鍾情了.拿俄米想:波阿斯是自己的至親.(按猶太人的規例,男子若沒有兒子而死了,應由至近親屬,娶死者的遺孀,生子歸死者名下,為死者留後.)拿俄米覺得應該挺身而出為這件大事作主,她就對路得說:女兒阿!我不當為你找個安身之處,使你享福麼!你與波阿斯的使女常在一處,波阿斯不是我們的親屬麼!他今夜在場上簸大麥.你要沐浴抹膏,換上衣服,下到場上,卻不要使那人認出你來.你等他吃喝完了,到他睡的時候,你看準他睡的地方；就進去掀開他腳上的被,你躺臥在那裡,他必告訴你所當作的事.路得說:「凡你所吩咐的,我必遵行.」(三1-5),於是路得完完全全地順服了.
\newline
\newline


\newpage

\section{第44屆~港九培靈會~費述凱~第三講　永遠得救的問題}
\label{sec:1379}
\textbf{第三講　永遠得救的問題}
\newline
\newline
連結: \href{https://www.hkbibleconference.org/session-message/view/1379}{\texttt{https://www.hkbibleconference.org/session-message/view/1379}}
\newline
\newline
\hyperref[sec:1375]{< < < PREV SERMON < < <}
~
\hyperref[sec:toc]{[返目錄]}
~
\hyperref[sec:1373]{> > > NEXT SERMON > > >}
\newline
\newline
希伯來書 5:8-5:14
\newline
\begin{longtable}{cl}
\hline
\hline
章節 & 經文 (和合本修訂版)\\
\hline
5:8 & \begin{tabularx}{0.7\textwidth}{X} 他雖然為兒子，還是因所受的苦難學了順從。 \end{tabularx} \\ \\ \relax
5:9 & \begin{tabularx}{0.7\textwidth}{X} 既然他得以完全，就為凡順從他的人成了永遠得救的根源， \end{tabularx} \\ \\ \relax
5:10 & \begin{tabularx}{0.7\textwidth}{X} 並蒙神照著麥基洗德的體系宣稱他為大祭司。 \end{tabularx} \\ \\ \relax
5:11 & \begin{tabularx}{0.7\textwidth}{X} 論到這事，我們有好些話要說，可是很難解釋，因為你們聽不進去。 \end{tabularx} \\ \\ \relax
5:12 & \begin{tabularx}{0.7\textwidth}{X} 按時間說，你們早該作教師了，誰知還需要有人再將神聖言基礎的要道教導你們；你們成了那需要吃奶、不能吃乾糧的人。 \end{tabularx} \\ \\ \relax
5:13 & \begin{tabularx}{0.7\textwidth}{X} 凡只能吃奶的，就不熟練仁義的道理，因為他是嬰孩。 \end{tabularx} \\ \\ \relax
5:14 & \begin{tabularx}{0.7\textwidth}{X} 惟獨長大成人的才能吃乾糧，他們的心竅因練習而靈活，能分辨善惡了。 \end{tabularx} \\ \\
[1ex]
\hline
\hline
\end{longtable}
這兩段經文,當我們讀的時候,會引起許多的問題,特別是六章一至九節.第九節末句有四個字是我們很難明白的,「近乎得救」這四個字是甚麼意思?好像教會中的信徒,將近得救還未得救；好像離天國很近,但是還沒有進去.他們是已經受洗加入教會的信徒,為何作者還說他們近乎得救呢?
\newline
\newline
第六章六節首句「若是離開道理」,這話是背道離教的意思.這句話也叫我們覺得希奇,究竟是不是指已經受洗的信徒說的?還是希伯來信徒中,有些曾蒙恩得救的信徒,後來離開了主,拒絕了主!是否救恩也會失去呢!今天得著救恩而明天就失去救恩呢?有這樣的事嗎?我們所得的救恩是永遠的,還是暫時的呢?這些問題非常嚴重,而且叫許多基督徒的心裡,引起了許多懷疑.我是個基督徒,我發脾氣是否還得救?我犯了大罪,神是否離棄我?救恩是否永遠的救恩?這些都是我們今天晚上要討論的問題.
\newline
\newline
如果我們對於得救還沒有把握,怎能做個快快樂樂的基督徒呢?怎能有勇氣大膽為耶穌作見證呢?怎會充滿喜業,充滿信心,和充滿盼望呢?怎敢開口勸別人信耶穌呢?絕對不會.所以關乎得救的問題,一定要尋求根本的解決,才能叫我們在真道上有長進.
\newline
\newline
請看五章九節:「祂既得以完全,就為凡順從祂的人,成了永遠得救的根源.」耶穌救了我們,祂所賜的救恩,是永遠得救的救恩.所以我們得救是永遠得救,這叫我們的心大得安慰.
\newline
\newline
耶穌親口說:「我的羊聽我的聲音,我也認識他們,他們也跟著我.」(約十27)你是不是耶穌的羊,你有沒有聽祂的聲音,有否跟隨祂?祂認不認識你,把你的名記在生命冊上?
\newline
\newline
耶穌說:「我又賜給他們永生,他們永不滅亡!誰也不能從我手裡把他們奪去.我父把羊賜給我,祂比萬有都大；誰也不能從我父手裡把他們奪去.」(約十 28-29)耶穌賜給我們永生,永不滅亡!祂所完成的救贖,是永遠的救恩,我們得救就是永遠得救.祂說:「誰也不能從我手裡,把他們奪去.......祂比萬有都大,誰也不能從我父手裡把他們奪去.」這是何等安慰的話.
\newline
\newline
如果你信了耶穌,把你的靈魂交托給祂,你自己卻沒有辦法把自己的靈魂拿出,撒旦更是沒有辦法；因為耶穌把得救的人,保守在祂的手中,同時也保守在天父的手中.基督徒是安全的,救恩是可靠的,因為神負責保守我們.
\newline
\newline
當我在加拿大讀神學時,有一晚,我參加某教會的禱告聚會.有位老弟兄流淚的禱告說:「主啊!我常抓住你,不然我真是不容易站著,恐怕我鬆了手,會掉到地獄裡面去...」他禱告得十分可憐.我讀過約翰十章廿七至廿九節,我心裡想:「老弟兄啊!不是你的手抓住神,乃是神的手保守你呀!」你的手會發酸,軟弱,一放鬆便會掉下地獄去.人的靈魂太寶貝,是用重價買來的；天父不讓人自己負責,乃由祂和耶穌親自負這重大的責任.耶穌親口說:「誰也不能從我手裡把他們奪去.誰也不能從我父手裡,把他們奪去.」故此,我們得救,是永遠的得救.
\newline
\newline
聖經記載,有次撒旦想要把摩西的屍體,從墳墓移出.撒旦對摩西的靈魂,卻沒有辦法下手,所以想利用他的屍體,放在重要的地點.讓以色列人看見,以為摩西回來了,便將他當偶像跪拜.神不許撒旦引誘以色列人,把摩西的屍體當偶像跪拜；神就打發米加勒(是位最大的天使)加以阻止.(參猶9)摩西的屍體,神尚且負責；我們的靈魂是主耶穌流寶血用重價買來的,神豈不永遠負責嗎?
\newline
\newline
信徒永遠得救,這是聖經中的真理.如果你有所懷疑,是因你對耶穌不夠認識,不清楚明白耶穌救贖的完備；以至給撒旦留地步,將懷疑的火箭射入你心中.
\newline
\newline
人對神容易偏向懷疑,而難轉向信靠.撒旦在今日常利用有大學問的人,把無神論的思想,向青年人灌注.
\newline
\newline
昨天我在青年會的電梯中,管理電梯的青年人,問我從那裡來的,我說從菲律賓來的.他問我是菲律賓的大使嗎?我說不是的,我不過是天國的一個小小的使者.他說:甚麼天國?那裡有天國?這話好像懷疑沒有天,只有地；沒有神,只有人.懷疑的根源是從樂園出來的.「夏娃對蛇說,園中樹上的果子我們可以吃；唯有園當中那棵樹上的果子,神曾說,你們不可吃,也不可以摸；免得你們死.」「蛇對夏娃說,神豈是真說不許你們吃園中所有樹上的果子麼?」(創三2-1)蛇的一句話,就叫夏娃懷疑神的話了.
\newline
\newline
如果撒旦對你說:「你信了耶穌真的得救嗎?真有永生嗎?真能白白的上天堂嗎?」你應該勇敢大膽以對,主耶穌曾親口說:「我又賜給他們永生,他們永不滅亡!誰也不能從我手裡,把他們奪去.我父把羊賜給我,祂比萬有都大,誰也不能從我父手裡把他們奪去.」耶穌不但救贖我們,祂也保守我們到永遠.
\newline
\newline
如果我們認識耶穌,了解祂完備的救恩,我們的頭腦充滿了耶穌基督的真理,智慧,我們就足以抵擋撒旦的火箭.以弗所六章十七節首句說:「戴上救恩的頭盔」頭盔是用來保護頭部的.我在馬尼拉日常以電單車代步,有時候還有師母在後面.我備有一頂很堅固的頭盔,以防萬一.因為頭都受傷最為危險,有了頭盔保護,可策安全.盼望我們都戴上救恩的頭盔,我們的頭腦若裝滿了救恩,充滿了認識耶和華的智慧,當撒旦把懷疑的火箭,射過來時便失去牠的力量了.
\newline
\newline
所以,我們「務要認識耶和華,竭力追求認識祂.」這樣,對自己得救,就毫無疑問了.只要我們追求認識耶和華,「祂出現確如晨光」黑夜了,晨光出現,這是天然的定律.我們只要追求認識耶和華,有個屬靈的定律,祂必向我們出現,「祂必臨到我們像甘雨,像滋潤田地的春雨.」
\newline
\newline
如果我們不追求,祂就不向我們出現.不扣門,祂就不給我們開門.不祈求,祂就不垂聽；不尋找,就找不著,不追求,就得不到；追求,一定得到.追求在乎我們,我們務要多多追求認識耶和華.
\newline
\newline
菲律賓聖經學院的師生們,除了許多工作之外,每年讀新約全書兩遍.我每天讀一百七十一節,一個三百六十五天可以把全部聖經讀完.我每天約用廿分鐘詳細讀經.讀耶和華的書,才能夠認識耶和華,追求認識耶和華是最重要的.
\newline
\newline
最近菲律賓聖經學院,有位從越南來的校友林王輝姊妹,她鼓勵全學院的同學和校友共廿七位,每位背讀一卷新約聖經.她說,因為常聽見基督徒聖經被沒收,萬一此事臨到菲律賓的基督徒,新約全書可以由廿七位弟兄姊妹,合力重寫出來.所以我們要竭力追求耶和華,將主的話多多藏在心裡.
\newline
\newline
永遠得救的道理,果然記在聖經上,但希伯來書六章三至五節,實在叫我們莫明其妙,這些基督徒是否得救的呢?他們原來已認識主.已經得救,為何又退步的呢?抑是慕道友尚未得救的呢?這有兩個可能:一是得救後又隨流失去了.其二是他們還站在救恩的門外往裡看.北方有句俗話說:「捧誰的碗,受誰的管；吃誰的飯,聽誰的話.」有的人站著觀望,盼望罪赦免,上天堂,有永生諸多好處.但又覺得信了耶穌,要受神的管教,實在不自由.所以停留在那裡慕道,一直站在門外至死,死了滅亡,錯過了得救的機會.
\newline
\newline
「論到那些已經蒙了光照,嘗過天恩滋味,又於聖靈有分,並嘗過神善道的滋味,覺悟來世權能的人.」(4-5)這裡描寫神在罪人身上所做的五種工作.例如我在未信耶穌之前三個月,有位老姊妹天天為我禱告,聖靈在我身上動工.假使聖靈不做這預備的工作,我無法悔改,我無法接受主,因為我們的眼目已經被撒旦弄瞎了.那時我看不見天堂的美麗,也看不見地獄的可怕.聖靈一動工,我覺悟到來世的權能,天堂是何等可羡慕的.我開始恐懼地獄的可怕,我若今日死去如何?我實在沒有把握,沒有盼望.我趕快跑到外面,雖然一片烏天黑地,我也無暇顧及,我跪下禱告.這樣,聖靈在罪人身上的預備工作,就臨到我的身上了.
\newline
\newline
古時候,長老會的神學家,把聖靈預備人心的工作,起了一個名詞叫「先來之恩」.在我未得救之前,神的恩典先臨到我身上,釋放我脫離撒旦.因為撒旦是迷惑普天下的,蒙蔽了我,不讓我信耶穌.聖靈來了,光照我,釋放我,教我分別天堂和地獄.聖靈這個工作是暫時的,在我的身上三個月.最後我就接受耶穌,出死入生,從撒旦的權下歸向神.
\newline
\newline
聖靈這些預備的工作,在希伯來信徒身上,並不是證明他們已經得救,乃是盼望他們得救.但他們猶豫不決,越遲延,越等候,他們的心越硬；最後,他們離棄道理,無法叫他們重新懊悔.他們把主耶穌重釘十字架.
\newline
\newline
許多年前,加拿大有個青年人,他的父親在大城市修理鐘表,這青年人也學會他父親的職業.但是他討厭這門工作；他覺得整天戴副放大鏡坐著修理鐘表,毫無意思.有天他對他父親說,要到西部溫哥華深山,去做砍伐樹木的工作,藉以鍛煉身體,呼吸新鮮的空氣.經他父親許可,他就離家去了.因他沒有經驗,砍伐一棵大樹倒下來時被壓重傷.三個伙伴用雪車渡過結冰的河,把他送往市區醫院求醫.經醫生診斷,必須鋸掉一條折骨破碎的腿,或可保全性命.但手術之前,需經傷者本人,或其家屬簽字.那時傷者因發高熱,已經三天昏迷不醒了,他的姓名,地址,家屬,一概不得而知；因此醫生無法施行手術.後來醫生忽然想起,藥商正推銷一種退熱的特效藥,所以就給他先注射該藥,果然退熱,人也清醒.醫生把傷勢嚴重和必須鋸腿的詳情告訴他,經他同意簽字,醫生便進行施手術.雖然來他缺了一條腿,但卻挽救了他的性命.
\newline
\newline
聖靈在一個罪人身上的「先來之恩」就像退熱的特效藥,可以退卻從撒旦來的熱,使我們頭腦清醒；讓我們有清楚的頭腦,自由的意志來選擇,接受主耶穌基督的救恩.
\newline
\newline


\newpage

\section{第44屆~港九培靈會~費述凱~第四講　路得信心的賞賜}
\label{sec:1373}
\textbf{第四講　路得信心的賞賜}
\newline
\newline
連結: \href{https://www.hkbibleconference.org/session-message/view/1373}{\texttt{https://www.hkbibleconference.org/session-message/view/1373}}
\newline
\newline
\hyperref[sec:1379]{< < < PREV SERMON < < <}
~
\hyperref[sec:toc]{[返目錄]}
~
\hyperref[sec:1380]{> > > NEXT SERMON > > >}
\newline
\newline
路得記 4:1-4:22
\newline
\begin{longtable}{cl}
\hline
\hline
章節 & 經文 (和合本修訂版)\\
\hline
4:1 & \begin{tabularx}{0.7\textwidth}{X} 波阿斯上到城門，坐在那裡，看哪，波阿斯所說那個可以贖產業的至親經過。波阿斯說：「某某先生，請你轉回來，坐在這裡。」他就轉回來坐下。 \end{tabularx} \\ \\ \relax
4:2 & \begin{tabularx}{0.7\textwidth}{X} 波阿斯又請了本城的十個長老來，對他們說：「請你們坐在這裡。」他們就都坐下。 \end{tabularx} \\ \\ \relax
4:3 & \begin{tabularx}{0.7\textwidth}{X} 波阿斯對那至親說：「從摩押地回來的拿娥米，現在要賣我們弟兄以利米勒的那塊地。 \end{tabularx} \\ \\ \relax
4:4 & \begin{tabularx}{0.7\textwidth}{X} 我想我應該向你說清楚：你可以買那塊地，當著在座的眾人和我百姓的長老面前，你若要贖就贖吧！倘若你不贖就告訴我，讓我知道，因為除了你以外，沒有人可以先贖，在你之後才輪到我。」那人說：「我要贖。」 \end{tabularx} \\ \\ \relax
4:5 & \begin{tabularx}{0.7\textwidth}{X} 波阿斯說：「你從拿娥米和摩押女子路得手中買這地的時候，也當買死人的妻子，使死人在產業上留名。」 \end{tabularx} \\ \\ \relax
4:6 & \begin{tabularx}{0.7\textwidth}{X} 那至親說：「這樣我就不能贖了，免得對我的產業有損。你儘管去贖我所當贖的吧，我不能贖了！」 \end{tabularx} \\ \\ \relax
4:7 & \begin{tabularx}{0.7\textwidth}{X} 從前，在以色列中要確認任何交易，無論是贖業或買賣，一方必須脫鞋給另一方。以色列中都以此為證。 \end{tabularx} \\ \\ \relax
4:8 & \begin{tabularx}{0.7\textwidth}{X} 那至親對波阿斯說：「你自己買吧！」於是把鞋脫了下來。 \end{tabularx} \\ \\ \relax
4:9 & \begin{tabularx}{0.7\textwidth}{X} 波阿斯對長老和所有在場的百姓說：「你們今日都是證人；凡屬以利米勒，以及基連和瑪倫的，我都從拿娥米手中買下來了。 \end{tabularx} \\ \\ \relax
4:10 & \begin{tabularx}{0.7\textwidth}{X} 我也娶瑪倫的妻子摩押女子路得，好讓死人可以在產業上留名，免得他的名在本族本鄉的城門中消失了。你們今日都是證人。」 \end{tabularx} \\ \\ \relax
4:11 & \begin{tabularx}{0.7\textwidth}{X} 在城門坐著的所有百姓和長老說：「我們都是證人。願耶和華使進你家的這女子，像建立以色列家的拉結和利亞二人一樣。又願你在以法他得亨通，在伯利恆有名聲。 \end{tabularx} \\ \\ \relax
4:12 & \begin{tabularx}{0.7\textwidth}{X} 願耶和華從這年輕女子賜你後裔，使你的家像她瑪從猶大所生法勒斯的家一樣。」 \end{tabularx} \\ \\ \relax
4:13 & \begin{tabularx}{0.7\textwidth}{X} 於是，波阿斯娶了路得為妻，與她同房。耶和華使她懷孕生了一個兒子。 \end{tabularx} \\ \\ \relax
4:14 & \begin{tabularx}{0.7\textwidth}{X} 婦女們對拿娥米說：「耶和華是應當稱頌的！因為他今日沒有使你斷絕可以贖產業的至親。願這孩子在以色列中得名聲。 \end{tabularx} \\ \\ \relax
4:15 & \begin{tabularx}{0.7\textwidth}{X} 他必振奮你的精神，奉養你的晚年，因為他是愛慕你的媳婦所生的。有這樣的媳婦，比有七個兒子更好！」 \end{tabularx} \\ \\ \relax
4:16 & \begin{tabularx}{0.7\textwidth}{X} 拿娥米接過孩子來，抱在懷中撫養他。 \end{tabularx} \\ \\ \relax
4:17 & \begin{tabularx}{0.7\textwidth}{X} 鄰居的婦人給孩子起名，說：「拿娥米得了一個孩子了！」她們就給他起名叫俄備得。俄備得是耶西的父親，是大衛的祖父。 \end{tabularx} \\ \\ \relax
4:18 & \begin{tabularx}{0.7\textwidth}{X} 這是法勒斯的後代：法勒斯生希斯崙； \end{tabularx} \\ \\ \relax
4:19 & \begin{tabularx}{0.7\textwidth}{X} 希斯崙生蘭；蘭生亞米拿達； \end{tabularx} \\ \\ \relax
4:20 & \begin{tabularx}{0.7\textwidth}{X} 亞米拿達生拿順；拿順生撒門； \end{tabularx} \\ \\ \relax
4:21 & \begin{tabularx}{0.7\textwidth}{X} 撒門生波阿斯；波阿斯生俄備得； \end{tabularx} \\ \\ \relax
4:22 & \begin{tabularx}{0.7\textwidth}{X} 俄備得生耶西；耶西生大衛。 \end{tabularx} \\ \\
[1ex]
\hline
\hline
\end{longtable}
有位親愛的弟兄,當我們上週查希伯來書；讀到何西阿書六章三節時,他受了感動,就把該節經文,分成四行作出美麗的字畫送給我,我心裡很受感動,現在特為公開多謝這位隱名的弟兄.我問他貴姓大名,他謙卑隱藏,因為他不想得地上的賞賜,他願意得將來天上的賞賜.
\newline
\newline
現在,請大家異口同聲來讀何西阿書六章三節:「我們務要認識耶和華,竭力追求認識祂.祂出現確如晨光,祂必臨到我們像甘雨,像滋潤田地的春雨.」這節聖經的重要訊息是:凡事必有因果,好像雨過天晴；乾旱之後降下甘雨,黑夜已過白晝來臨.我們若竭力追求認識耶和華,祂必快快出現,降福給我們.神的心是軟的,我們若向祂祈求,祂必賞賜我們.人心剛硬,常常消滅聖靈的感動；有的不肯作工領人歸主,有的不肯接受耶穌白白的救恩.
\newline
\newline
從前有位老秀才幫助我讀四書,他念一句「人之初,性本善」我也跟著他念,我把全部四書都學完了.四書內容論到因和果；有果必有其因,有因必產生其果.我們追求認識耶和華,這是因.所產生的果,就是祂要向我們顯現,祂要恩待我們.這是一定的,為何我們不追求,不多禱告呢?
\newline
\newline
路得記一章論到路得信心的選擇.二章論路得信心的試驗.三章論路得信心的順服.這些在耶和華的眼中都看為很寶貴的,祂的心大受感動,祂看見了路得對祂的態度,對祂的熱誠.祂看見了路得對祂的態度,對祂的懇切,和對祂的信心.
\newline
\newline
路得記一至三章可以說是「因」,從路得對耶和華的態度表現出來,如何仰望祂,渴望認識祂,相信祂,順服祂.她的誠心尋求,叫耶和華神受了感動.
\newline
\newline
路得記第四章是「果」論到路得從耶和華手中所得的賞賜.因路得的信心,生活在神面前；結果,神厚厚地賜恩典給她.
\newline
\newline
一,神使那一位至近的親屬,甘心放棄贖拿俄米的田地,和娶路得的權利.當波阿斯對那人說:「當贖那塊地的是你,其次是我；此外再沒有別人了,......」那人回答說:「我肯贖.」波阿斯說:「你從拿俄米手中買這地的時候,也當娶死人的妻摩押女子路得.使死人在產業上存留他的名.那人說:這樣我就不能贖了,恐怕與我的產業有礙,你可以贖我所當贖的,我不能贖了.」(四4-6)
\newline
\newline
二,長老和眾民都異口同聲的贊成.
\newline
\newline
三,波阿斯的心靈得到釋放.
\newline
\newline
四,波阿斯和路得結婚後,得了一個兒子,名字叫俄備得,這名字是事奉的意思.
\newline
\newline
五,拿俄米和路得經濟上的問題得解決,因有大財主波阿斯負一切的責任.
\newline
\newline
六,拿俄米心裡大得安慰.拿俄米的名字,本來是甜的意思,可是後來她的丈夫,和兩個兒子都死了,空空地從摩押回來,所以「拿俄米對他們說:不要叫我拿俄米,要叫我瑪拉,(是苦的意思)因為全能者使我受了大苦.」(20節)可是現在不苦了,她抱著孩子俄備得,滿心快樂,滿臉笑容,真是甜啊!「拿俄米就把孩子抱在懷中,作他的養母.鄰舍的婦人說,拿俄米得孩子了.......」(16-17)
\newline
\newline
七,路得是主耶穌肉身的先祖.請參看馬太一章五至六節「撒門從喇合氏生波阿斯,波阿斯從路得氏生俄備得,俄備得生耶西,耶西生大衛王.大衛從烏利亞的妻子生所羅門.」保羅說:「我主耶穌基督按肉體說是大衛的子孫.」主給路得最大的賞賜,就是賜她作為主耶穌肉身的先祖.
\newline
\newline
假使路得沒有選擇耶和華,沒有信靠耶和華,沒有順服耶和華,路得的生命,就不能從波阿斯開始,一代代的延續到主耶穌了.如果路得沒有選擇耶和華,沒有倚靠耶和華!那麼,伯利恆的馬槽,就永遠是空的,也就沒有救主聖誕的節日了.因為路得順服耶和華,跟從耶和華,從摩押地出來.要耶和華作她的神,她要作以色列的民；她才能成為主耶穌最重要的先祖.
\newline
\newline
路得最大的賞賜,是她的生命與主耶穌發生永永遠遠的關係.你是否與主耶穌,有生命上的關係呢?可能在這次培靈會之前,你不清楚自己是否得救?是否有永生?是否基督徒?心裡常常懷疑,但是在這次培靈會之間,神藉著某講師用祂的話,使你忽然明白；「凡接待祂的就是信祂名的人,祂就賜給他們權柄,作神的兒女.」(約一12)如果你已全心全意信祂,認為除祂以外沒有別的救法；那麼,你就是接待祂的,就是信祂名的人,就得到祂所賜的權柄,作神的兒女.你已經和主耶穌的生命上有了關係.
\newline
\newline
現在,我要給你們一個機會,如果在這次培靈會,神恩待你；把你的疑惑解開,叫你清楚明白你已經得救了,請你舉手表示.
\newline
\newline
感謝主!現在有很多人舉手,這是用手見證.我要用口作見證,為甚麼我要當傳教士呢?本來我希望成為一個航空工程師,我在我父親的農場工作,終日思想,如何可發明更快更有力量,更經濟的飛機發動機,是其他國家所沒有的.我想成功,想發財.我還為我名成利就之後作了一番計劃:一,為我母親蓋大洋樓.(因我母親住的是鄉村偏僻地方的小房子)二,幫助我七個妹妹和兩個弟弟讀大學.三,成家立室,和我的愛人一生一世,快快樂樂作夫妻.
\newline
\newline
神救我,我聽了福音,受了三個月的痛苦；聖靈叫我為罪,為義,為審判,自己責備自己.到了一九二九年八月四日,那時我未滿廿一歲,我把主耶穌接到我心裡,做個人的救主.我就開始相信祂,倚靠祂,順服祂.祂初感動我去讀聖經學校.我去時,把那些航空工程的書籍都帶去；我雖已經信主,但我並不打算改變原有的計劃,我要做一個飛機工程師.入學後,我聽老教授們講課,得到了很多幫助,並且很喜愛所修的課程；所以就漸漸忘記,那些航空工程的書籍.到了將近畢業的前些日子,我每天照常讀經.有一天,我讀到(馬可十章29-30)時,覺得這兩節經文的字特別大「耶穌說:我實實在在告訴你們,人為我和福音,撇下房屋,或是弟兄,姊妹,父母,兒女,田地；並且要受逼迫,在來世必得永生.」
\newline
\newline
我最愛我的母親,母親也最愛我.因為我父親在第一次世界大戰時陣亡,我是最大的孩子,常常幫助母親料理一切,所以我最不捨得撇下的是我的母親；其次是兩個弟弟,和七個妹妹.我還有個愛人,我不知道應把她列在何地位,有時我把她列在第一,可是後來我離開家鄉到中國去,我最想念的卻不是她,而是我的母親.母親是我最愛的人,其次,我愛我的老家,愛我的田地,我喜愛種地,下了種子,發芽,生長,成熟,出去收割.我們有牛,馬,豬,雞,真是樂也融融!我愛作航空工程師,我愛發財,愛一切物質的享受.
\newline
\newline
耶穌說,為祂和福音撇下,最愛的人和最愛的東西；就在今世得百倍,這句話正抓住我的心.百倍,實是了不起的數!那時候,加拿大的銀行,一百元存款,每年僅得四元利息,離一倍尚差很遠.但耶穌說,人若為祂和福音撇下一切所愛的；不但今世得百倍,且在來世必得永生.
\newline
\newline
我抓住了神的應許,安靜跪下在宿舍房裡,神開始對我說話:「孩子!現在有兩條路讓你選擇,你今年廿三歲,人平均生命是七十歲.假如你活到七十歲,這段時間你如何用法?如果你把四十七年的時間,用在你所愛的人和物方面,那就是消耗了你四十七年的光陰.比方有人白白送你四十七塊錢,你請了一些人,吃喝快樂,一下就花完了.相反地,你有四十七元,或者四十七美元,或是四十七塊黃金你不濫用?把它存於可靠的銀行,或投資在可靠的事業上,不但本錢依然存在,每年還可生利息.如果你投資在我這裡,利息高達百倍,而且可靠,因我是神,永不撒謊,難道你不肯嗎?難道你要耗盡四十七年的光陰,為你所愛的人和物質嗎?難道你要研究製造最好的飛機供空軍殺人嗎?願你將你四十七年的光陰,交在釘十架的耶穌手中；獻給最愛的耶穌基督；為最愛福音使用嗎?你從今開始,為我活下去,擔保你今生得百倍.你為我和福音撇下你所愛的母親,我必興起百個老姊妹愛你如你母親一樣.況且你的母親現在還不懂得為你禱告,我可以興起百個姊妹曉得為你禱告的,豈不更好嗎?」我說:「主啊!我願意將我餘下的光陰交在你手中.」
\newline
\newline
在這四十年中,神至少給我百個比我鄉下更好的家.我的老農村,最多不過四分之一方英里之大,如今神賜給我之地,是整個的東南亞,凡有中國人的地方,就是我撒福音種子之地.比較之下,我的老農莊如同一枚小郵票.放下一個,得到更好的百個,這非但不是犧牲；而且是大大得利.非常寶貴!
\newline
\newline
今天晚上是一九七二年培靈會的最後一晚,神要向你我挑戰.你還有多少年為祂而活?我四十年已經過了,餘下的年日,還是留在主耶穌手中,讓祂使用.你呢?是為自己,為所愛的人呢?或是翻過來對主說:「主啊!我雖愛我的父母,妻子,兒女,事業,但是我應當更愛耶穌.祂撇下一切拯救我,我應當讓別人有機會得救.主啊!我願意將我餘下的光陰交在你手中,求你一步步帶領我,我願意為你和福音撇下一切.」
\newline
\newline


\newpage

\section{第44屆~港九培靈會~費述凱~第四講　彌賽亞也是祭司嗎?}
\label{sec:1380}
\textbf{第四講　彌賽亞也是祭司嗎?}
\newline
\newline
連結: \href{https://www.hkbibleconference.org/session-message/view/1380}{\texttt{https://www.hkbibleconference.org/session-message/view/1380}}
\newline
\newline
\hyperref[sec:1373]{< < < PREV SERMON < < <}
~
\hyperref[sec:toc]{[返目錄]}
~
\hyperref[sec:1381]{> > > NEXT SERMON > > >}
\newline
\newline
希伯來書 7:1-7:28
\newline
\begin{longtable}{cl}
\hline
\hline
章節 & 經文 (和合本修訂版)\\
\hline
7:1 & \begin{tabularx}{0.7\textwidth}{X} 這麥基洗德就是撒冷王，是至高神的祭司。他在亞伯拉罕打敗諸王回來的時候迎接他，並給他祝福。 \end{tabularx} \\ \\ \relax
7:2 & \begin{tabularx}{0.7\textwidth}{X} 亞伯拉罕也將自己所得來的一切，取十分之一給他。他頭一個名字翻譯出來是「公義的王」，他又名「撒冷王」，是和平王的意思。 \end{tabularx} \\ \\ \relax
7:3 & \begin{tabularx}{0.7\textwidth}{X} 他無父、無母、無族譜、無生之始、無命之終，是與神的兒子相似，他永遠作祭司。 \end{tabularx} \\ \\ \relax
7:4 & \begin{tabularx}{0.7\textwidth}{X} 你們想一想，這個人多麼偉大啊！連先祖亞伯拉罕都拿戰利品的十分之一給他。 \end{tabularx} \\ \\ \relax
7:5 & \begin{tabularx}{0.7\textwidth}{X} 那得祭司職分的利未子孫，奉命照例向百姓取十分之一，這百姓是自己的弟兄，雖是從亞伯拉罕親身生的，還是照例取十分之一。 \end{tabularx} \\ \\ \relax
7:6 & \begin{tabularx}{0.7\textwidth}{X} 惟獨麥基洗德那不與他們同族譜的，從亞伯拉罕收取了十分之一，並且給蒙應許的亞伯拉罕祝福。 \end{tabularx} \\ \\ \relax
7:7 & \begin{tabularx}{0.7\textwidth}{X} 向來位分大的給位分小的祝福，這是無可爭議的。 \end{tabularx} \\ \\ \relax
7:8 & \begin{tabularx}{0.7\textwidth}{X} 在這事上，一方面，收取十分之一的都是必死的人；另一方面，收取十分之一的卻是那位被證實是活著的。 \end{tabularx} \\ \\ \relax
7:9 & \begin{tabularx}{0.7\textwidth}{X} 我們可以說，那接受十分之一的利未也是藉著亞伯拉罕納了十分之一， \end{tabularx} \\ \\ \relax
7:10 & \begin{tabularx}{0.7\textwidth}{X} 因為麥基洗德迎接亞伯拉罕的時候，利未還在他先祖的身體裡面。 \end{tabularx} \\ \\ \relax
7:11 & \begin{tabularx}{0.7\textwidth}{X} 那麼，如果百姓藉著利未人的祭司職任能達到完全—因為百姓是在這職分下領受律法的—為甚麼還需要按照麥基洗德的體系另外興起一位祭司，而不按照亞倫的體系呢？ \end{tabularx} \\ \\ \relax
7:12 & \begin{tabularx}{0.7\textwidth}{X} 既然祭司的職分已更改，律法也需要更改。 \end{tabularx} \\ \\ \relax
7:13 & \begin{tabularx}{0.7\textwidth}{X} 因為這些話所指的人本屬別的支派，那支派裡從來沒有一人在祭壇前事奉的。 \end{tabularx} \\ \\ \relax
7:14 & \begin{tabularx}{0.7\textwidth}{X} 很明顯地，我們的主是從猶大出來的；但關於這支派，摩西並沒有提到祭司。 \end{tabularx} \\ \\ \relax
7:15 & \begin{tabularx}{0.7\textwidth}{X} 倘若有另一位像麥基洗德的祭司興起來，我的話就更顯而易見了。 \end{tabularx} \\ \\ \relax
7:16 & \begin{tabularx}{0.7\textwidth}{X} 他成為祭司，並不是照屬肉身的條例，而是照無窮生命的大能。 \end{tabularx} \\ \\ \relax
7:17 & \begin{tabularx}{0.7\textwidth}{X} 因為有給他作見證的說：「你是照著麥基洗德的體系永遠為祭司。」 \end{tabularx} \\ \\ \relax
7:18 & \begin{tabularx}{0.7\textwidth}{X} 一方面，先前的誡命因軟弱無能而廢掉了， \end{tabularx} \\ \\ \relax
7:19 & \begin{tabularx}{0.7\textwidth}{X} （律法本來就不能成就甚麼）；另一方面，一個更好的指望被引進來，靠這指望，我們就可以親近神。 \end{tabularx} \\ \\ \relax
7:20 & \begin{tabularx}{0.7\textwidth}{X} 再者，耶穌成為祭司，並不是沒有神的誓言；其他的祭司被指派時並沒有這種誓言， \end{tabularx} \\ \\ \relax
7:21 & \begin{tabularx}{0.7\textwidth}{X} 只有耶穌是起誓立的，因為那位立他的對他說：「主起了誓，絕不改變。你是永遠為祭司。」 \end{tabularx} \\ \\ \relax
7:22 & \begin{tabularx}{0.7\textwidth}{X} 既是起誓立的，耶穌也作了更美之約的中保。 \end{tabularx} \\ \\ \relax
7:23 & \begin{tabularx}{0.7\textwidth}{X} 一方面，那些成為祭司的數目本來多，是因為受死亡限制不能長久留住。 \end{tabularx} \\ \\ \relax
7:24 & \begin{tabularx}{0.7\textwidth}{X} 另一方面，這位既是永遠留住的，他具有不可更換的祭司職任。 \end{tabularx} \\ \\ \relax
7:25 & \begin{tabularx}{0.7\textwidth}{X} 所以，凡靠著他進到神面前的人，他都能拯救到底，因為他長遠活著為他們祈求。 \end{tabularx} \\ \\ \relax
7:26 & \begin{tabularx}{0.7\textwidth}{X} 這樣一位聖潔、無邪惡、無玷污、遠離罪人、高過諸天的大祭司，對我們是最合適的； \end{tabularx} \\ \\ \relax
7:27 & \begin{tabularx}{0.7\textwidth}{X} 他不像那些大祭司，每日必須先為自己的罪，後為百姓的罪獻祭，因為他只一次將自己獻上就把這事成全了。 \end{tabularx} \\ \\ \relax
7:28 & \begin{tabularx}{0.7\textwidth}{X} 律法所立的大祭司本是有弱點的人，但在律法以後，神以起誓的話立了兒子為大祭司，成為完全，直到永遠。 \end{tabularx} \\ \\
[1ex]
\hline
\hline
\end{longtable}
起初的教會,有許多信主的猶太人,都懷著一個大問題:「彌賽亞也是祭司嗎?」他們知道彌賽亞將來是要作王的,但未聞祂將來也做祭司.他們所盼望的彌賽亞,將來不是做祭司,乃是個君王.
\newline
\newline
舊約先知,約在主前七百年,預言神將要差遣彌賽亞,救他們脫離外邦人的手.(參詩二1-9),神所差來的彌賽亞要掌大權,審判以色列的仇敵,祂要來拯救百姓脫離外邦人的手.本處經文只預言祂是大君王,並無題到祂做祭司.還有(詩七二1-12)以色列人極其羡慕他們的王快來,就是他們所盼望大有能力的君王；是勇敢的戰士,可救他們脫離羅馬帝國的壓迫.
\newline
\newline
當時羅馬軍兵用武力,殘害以色列人,稅吏極盡敲詐能事,壓迫他們!他們唯一的盼望,是神所應許的彌賽亞快快來臨.他們覺得所需要的,是位能救國家脫離羅馬帝國壓迫的君王,並非需要一位救主,來救他們脫離他們的罪.他們需要大君王救他們的國家,而無需大祭司救他們的靈魂.所以他們的盼望是屬於政治,而非屬靈的盼望.他們只感到處境痛苦,而不覺悟有罪.他們莫明神為何打發救主,而不打發君王來呢?他們看見耶穌行神跡奇事,又溫柔和善,只不過是個木匠的兒子,沒有權勢,沒有地位,沒有大衛般的勇敢,怎能作君王領導百姓爭戰復國.他們以為犯了罪,可到聖殿去,獻上牛羊或鴿子贖罪.聖殿裡已經有很多祭司,也有大祭司,可以為罪人獻祭贖罪；為何神再差一位祭司來,豈不是多餘嗎?他們僅知國家的危險.
\newline
\newline
耶穌與祂的門徒同在三年多,多次告訴他們說:「我來不是要受人服事,乃是要服事人；並且要捨命,作多人的贖價.」「捨命」「作贖價贖罪」,這是祭司的工作.
\newline
\newline
當耶穌死在十字架上時,祂的門徒傷心至極.雖然耶穌曾經一次又一次的告訴他們,祂來是為捨命贖罪來的,但他們不了解,不相信,甚至懷疑耶穌,若真是神所差來的彌賽亞,要救以色列國脫離外邦的手；要復興以色列國,為甚麼還至於被壞人釘十字架?既來救人,為何不能自救呢?
\newline
\newline
請看詩篇一一○篇,這裡可以聽見天父對祂兒子彌賽亞說話,大衛王把神的話報告給我們.1-3節告訴我們耶穌是君王,將來彌賽亞要來作王.5-7節說祂將來要作審判者,要刑罰罪人的.耶穌再來的日子,祂不但做王,祂也作審判者,第四節說祂「是照著麥基洗德的等次,永遠為祭司.」同一彌賽亞,而具有君王,審判者,祭司三大職務.
\newline
\newline
以色列人所注意的,是彌賽亞作君王,作審判者,卻忽略了祂作大祭司.因他們以為已經有了祭司,無此需要了.然而主耶穌是照著麥基洗德的等次,永遠為祭司；乃是特別的,非同他們原有的祭司.(普通祭司任職由卅歲至五十歲,大祭司由五十歲至離世)是暫時的,聖殿裡從來沒有一位祭司,永遠在那裡任職.這叫百姓很難受的,他們從各地來到耶路撒冷敬拜神,他們所認識的祭司,天天歡迎他們.百姓也把家鄉土產相贈,彼此道好；如同見了親愛的牧師,他們有甚麼困難也告訴他；他就為他們禱告,用神的話講給他們聽,他們就大得安慰回家.到了第二年再來的時候,想不到看不見他們認識的老祭司了；原來他已經退休了,回到山頂去,不在聖殿了.他們頗感不便,因為新祭司彼此不認識；他們的情形,新祭司也不了解.所以神看這制度不好,祂要立一位祭司,永遠作祭司；不但一生一世為你負責,而且認識你一直到永遠.神所興起永遠的祭司,就是主耶穌.
\newline
\newline
那些猶太教的拉比們,每逢安息日讀舊約聖經,給來禮拜的人聽,料必讀過詩篇一百十篇第四節；但他們卻不知道祂「是照著麥基洗德的等次永遠為祭司.」只知道彌賽亞來要救他們,為他們捨命,把祂自己獻上為贖罪祭.可能那班拉比們,只單方面的道理,沒有告訴他們彌賽亞來要做大祭司,作永遠的大祭司；永遠認識我們,為我們的靈魂負責,救我們到底.所以耶穌的門徒彼得,約翰,雅各,馬太,還未知道這真理.
\newline
\newline
(希伯來書八8-18)這裡表明舊約的制度不完全.原本律法放在聖殿的約櫃裡面,這裡說神要將祂的律法放在他們裡面,寫在他們心上.舊的律法是寫在兩塊石板上,現在神把祂的律法放在我們裡面；把祂的生命放在我們裡面,叫我們自動地聽祂的話,愛祂的話.以前,祂作他們的王,他們作祂的百姓.現在,神作我們的天父,我們作祂的兒女,這比他們更美好了!所以在詩篇一一○篇,神應許,祂的兒子到世上來要作君王；要作審判者,也要作祭司.
\newline
\newline
獻祭──為罪人獻上贖罪祭,是一個祭司最重要的工作.請大家思想一下,耶穌是祭司,祂要獻上甚麼為贖罪祭?祂能否到聖殿和那些祭司一起獻牛羊,鴿子,為贖罪祭?經上曾否題及耶穌進聖殿獻祭的事呢?沒有,因為唯有利未人亞倫的子孫,才有資格獻祭.耶穌是屬於猶大支派的,猶大支派出了很多王,從未出過祭司.故此,更加叫猶太人懷疑,也叫門徒懷疑.
\newline
\newline
神要廢棄舊約另立新約,要把祭司的制度完全更改.耶穌作祭司並不是照亞倫,照利未支派的等次.「耶和華起了誓,決不後悔,說,你是照著麥基洗德的等次,永遠為祭司.」這句話非常寶貝,充滿了靈意.耶穌作祭司,不是照著亞倫的等次,而是照著麥基洗德的等次；是永遠的祭司,獨一無二的祭司,祂永遠不退休,永遠不死；永遠無需新任的祭司,來取代祂的職位,可惜門徒沒有領會這個意義!
\newline
\newline
「何況基督藉著永遠的靈,將自己無瑕無疵獻給神,祂的血豈不更能洗淨你們的心；除去你們的死行,使你們事奉那永生神麼?」(九14)耶穌作祭司,祂獻上自己,無瑕無疵獻給神,作我們的贖罪祭.因為公羊的血,斷不能除罪,只能暫時遮蓋罪.所以施洗約翰說:「看哪,神的羔羊!」神的羔羊絕對除罪,永遠解決罪的問題.牛羊鴿子的血,只能暫時遮蓋罪,等到神的羔羊來了,才能除淨從亞當,一直遺留下來的罪.祭司所獻的祭,只有暫時的好處.神知道將來耶穌要替我們清還罪債,暫時用獻牛獻羊的辦法,代替基督獻上自己的功勞.
\newline
\newline
「我們憑這旨意,靠耶穌基督,只一次獻上祂的身體,就得以成聖.」(十10)九章14節說「獻上自己」,這裡更清楚的告訴我們「耶穌基督只一次獻上祂的身體,我們就「得以成聖.」這位永遠的祭司,自己捨身,自己流血,要真正除淨信靠祂的人的罪.
\newline
\newline
耶穌來,既不單作君王,作審判者,也來作祭司.祂死在十字架上,不但作祭司,更是作犧牲,獻上自己.因為他們心目中的彌賽亞是得勝的,除滅仇敵的,永遠不死的；要帶領他們恢復大衛和所羅門的黃金國度時代,要在萬國中為首.故此,耶穌雖預言祂將受死,門徒也不相信,因為與他們的觀念不合.
\newline
\newline
主深知祂是為死而來世的,祂也知道來世的使命是甚麼,所以擺在祂面前十字架的道路,祂也不怕!請看十二2:「仰望為我們信心創始成終的耶穌,祂因那擺在前面的喜樂；就輕看羞辱,忍受了十字架的苦難,便坐在神寶座的右邊.」因為祂知道死在十字架上可以完成救贖,可以搶救千千萬萬的靈魂免下地獄,將來可以在世上成立祂的天國,祂捨命可作多人的贖價.因此,祂坦然無懼上十字架,直到祂大聲喊說「成了」!世上的罪人,便可以從黑暗走向光明,可以從撒旦的權下歸向神.這樣的喜樂,叫祂忍受十字架的苦難.
\newline
\newline
「他們聚集的時候,問耶穌說,主阿,你復興以色列國,就在這時候嗎?」(徒一6)這是耶穌復活之後,門徒向主第一句的問話.由這一句話,顯明門徒尚不了解耶穌已完成了救贖；已經盡了祭司的本分,使世界萬民有得救的盼望.他們所關心的不是世界萬民,而是小小的以色列國.他們沒有想到彌賽亞作王之前,要先作祭司,救贖萬民；有了千千萬萬得救之民,王才可以建立祂的天國.所以「耶穌對他們說,父憑著自己的權柄,所定的時候日期,不是你們可以知道的.但聖靈降臨在你們身上,你們就必得著能力!並要在耶路撒冷,猶太全地,和撒瑪利亞,直到地極,作我的見證.」(徒一7-8)
\newline
\newline
耶穌吩咐門徒往普天下去傳福音,讓人有機會得救；等到得救的數目滿足,祂就要回來作王.感謝主!從那時開始,他們對主耶穌就有所認識了.他們有一百廿人聚集,開了十天的禱告會,聖靈降臨在他們身上；他們就得著能力,往普天下去傳福音,為君王尋找百姓,等祂回來掌王權.
\newline
\newline
所以,我們都要追求認識耶和華,不是單單認識祂一方面,而是更多方面的.
\newline
\newline
耶穌不但是君王,是審判者,也是祭司.
\newline
\newline


\newpage

\section{第44屆~港九培靈會~費述凱~第五講　耶穌只一次獻上身體夠了嗎?}
\label{sec:1381}
\textbf{第五講　耶穌只一次獻上身體夠了嗎?}
\newline
\newline
連結: \href{https://www.hkbibleconference.org/session-message/view/1381}{\texttt{https://www.hkbibleconference.org/session-message/view/1381}}
\newline
\newline
\hyperref[sec:1380]{< < < PREV SERMON < < <}
~
\hyperref[sec:toc]{[返目錄]}
~
\hyperref[sec:1376]{> > > NEXT SERMON > > >}
\newline
\newline
希伯來書 10:1-10:18
\newline
\begin{longtable}{cl}
\hline
\hline
章節 & 經文 (和合本修訂版)\\
\hline
10:1 & \begin{tabularx}{0.7\textwidth}{X} 既然律法只不過是未來美好事物的影子，不是本體的真像，就不能藉著每年常獻一樣的祭物，使那些進前來的人完全。 \end{tabularx} \\ \\ \relax
10:2 & \begin{tabularx}{0.7\textwidth}{X} 若不然，獻祭的事豈不早已停止了嗎？因為敬拜的人僅只一次潔淨，良心就不再覺得有罪了。 \end{tabularx} \\ \\ \relax
10:3 & \begin{tabularx}{0.7\textwidth}{X} 但是這些祭物使人每年都想起罪來， \end{tabularx} \\ \\ \relax
10:4 & \begin{tabularx}{0.7\textwidth}{X} 因為公牛和山羊的血不能除罪。 \end{tabularx} \\ \\ \relax
10:5 & \begin{tabularx}{0.7\textwidth}{X} 所以，基督到世上來的時候，就說：「祭物和禮物不是你所要的，但你曾給我預備了身體。 \end{tabularx} \\ \\ \relax
10:6 & \begin{tabularx}{0.7\textwidth}{X} 燔祭和贖罪祭是你不喜歡的。 \end{tabularx} \\ \\ \relax
10:7 & \begin{tabularx}{0.7\textwidth}{X} 那時我說：看哪！我來了，我的事在經卷上已經記載了；神啊！我來為要照你的旨意行。」 \end{tabularx} \\ \\ \relax
10:8 & \begin{tabularx}{0.7\textwidth}{X} 以上說：「祭物和禮物，以及燔祭和贖罪祭，不是你所要的，也不是你喜歡的。」這都是按著律法獻的。 \end{tabularx} \\ \\ \relax
10:9 & \begin{tabularx}{0.7\textwidth}{X} 他接著說：「看哪！我來了，為要照你的旨意行。」可見他除去在先的，為要立定在後的。 \end{tabularx} \\ \\ \relax
10:10 & \begin{tabularx}{0.7\textwidth}{X} 我們憑著這旨意，藉著耶穌基督，僅只一次獻上他的身體就得以成聖。 \end{tabularx} \\ \\ \relax
10:11 & \begin{tabularx}{0.7\textwidth}{X} 所有的祭司天天站著事奉神，屢次獻上一樣的祭物，這祭物永不能除罪。 \end{tabularx} \\ \\ \relax
10:12 & \begin{tabularx}{0.7\textwidth}{X} 但基督獻了一次永遠有效的贖罪祭，就坐在神的右邊， \end{tabularx} \\ \\ \relax
10:13 & \begin{tabularx}{0.7\textwidth}{X} 從此等候他的仇敵成為他的腳凳。 \end{tabularx} \\ \\ \relax
10:14 & \begin{tabularx}{0.7\textwidth}{X} 因為他僅只一次獻祭，就使那些得以成聖的人永遠完全。 \end{tabularx} \\ \\ \relax
10:15 & \begin{tabularx}{0.7\textwidth}{X} 聖靈也對我們作證，因為他說過： \end{tabularx} \\ \\ \relax
10:16 & \begin{tabularx}{0.7\textwidth}{X} 「主說：那些日子以後，我與他們所立的約是這樣的：我要將我的律法放在他們的心上，又要寫在他們的心思裡。」 \end{tabularx} \\ \\ \relax
10:17 & \begin{tabularx}{0.7\textwidth}{X} 並說：「他們的罪惡和他們的過犯，我絕不再記得。」 \end{tabularx} \\ \\ \relax
10:18 & \begin{tabularx}{0.7\textwidth}{X} 這些罪過既已蒙赦免，就不用再為罪獻祭了。 \end{tabularx} \\ \\
[1ex]
\hline
\hline
\end{longtable}
今晚來赴會的青年人很多.青年人的頭腦充滿著問題,而成年人則充滿了答案.但成年的人答案,不一定能解答青年人的問題.
\newline
\newline
當青年人竭力追求耶和華的時候,漸漸地,他們的問題就沒有了,因為從主耶穌那裡得到了答案,他們相信人類一切的難題,靠主才能得解決.當成年人竭力追求耶和華的時候,他們就發現自己的答案都是錯誤,膚淺,不澈底的；於是他們開始尋求更深,更難的答案.
\newline
\newline
請看八章八節(注意小字)「主指前約的缺欠說」這裡說明主不單指責以色列人犯了罪,同時也指出摩西所立的約有缺欠,特別對於祭司們獻祭的工作.
\newline
\newline
神是無所不知的,難道在摩西設立祭壇獻祭,未訓練祭司之先,不知道這些方法會失敗的嗎?既知道,為甚麼又安排亞倫當祭司,利未人當祭司；有聖殿,有聖壇呢?難道試驗了一千五百多年,才發現那些方法不完善嗎?不,神是無所不知的神,祂知道；但祂仍吩咐摩西設立會幕聖壇,吩咐所羅門建造聖殿；每年獻上千千萬萬的牛羊鴿子.
\newline
\newline
請看十4-9,這顯明神的兒子到世上來的時候,知道那些獻祭的事不對.天父俯視以色列人天天獻祭,祭物牲畜之血成河；心覺不快樂,不滿意.既然如此,當初為何要設立獻祭呢?關於這些問題,第十章有很好的答案.
\newline
\newline
一至三節告訴我們,雖然神對獻祭不滿意,不能達到祂最高的願望,不能除淨人的罪,不能救人到底；但卻能使以色列人,生發有罪感.如果不是每年一度到聖殿獻祭,他們的心將變成麻木不仁；終日醉生夢死,不怕神,不怕死,不怕審判.他們獻祭,自知有罪,該死,無辜小羊為他死了.獻祭激發他們有罪之感!使他們覺得亟需一位救主,並盼望這位救主來臨.
\newline
\newline
還有一用處,請看一節,這裡說祭物是將來美事的影兒.主耶穌將要死在十字架上完成救贖,福音要傳到天邊地極；凡信靠祂的人都可以蒙恩得救,罪完全得赦免,成為神的兒女,得永生,上天堂,這些是將來的美事.在摩西當時的立場看是將來,因為那時事情還未發生.所以這裡說......是將來美事的影兒.
\newline
\newline
我生長於加拿大,所住的地方,既無山無樹,是片大草地.當太陽平面時,我走在草地上；背後因落日而射,形成一個長影兒在前面.我的體高本是五呎十一吋,但影兒約七十呎.在我前面行走的人,相當遠的距離,未見我的真人,已見我的影兒了.
\newline
\newline
祭物是未來美事的影兒,他們獻羊為祭物,是預表神的羔羊,主耶穌到世上完成的救贖.赦罪之恩並不在祭物牲畜裡面,不過藉以預表和教訓.比如幼稚園的老師教學生,問他們五加五等於多少?抽象的數法,他們必莫明其妙!若把生果分別擺成兩排,有實物供他們逐一數點,他們才會明白.神訓練以色列民,也像幼稚園老師教學生一樣.要用實物授課,訓練他們明白,屬靈的真理和道路.
\newline
\newline
以色列人每年一度到聖殿獻祭,宰了牛羊,把血洒在祭壇上；最後把牲畜的肉也燒在祭壇上.祭司便對獻祭的人說:「願你平安,奉耶和華的名,你的罪赦免了!」
\newline
\newline
舊約是用希伯來文寫的.「赦免」含有兩個字意,一是遮蓋的意思,另一是打發它走的意思,走到那裡?罪都歸到耶穌身上.「我們都如羊走迷,各人偏行己路.耶和華使我們眾人的罪孽都歸在祂身上.」(賽五三6)罪得遮蓋,但遮蓋只是暫時的；心裡暫時得著平安,慢慢地良心上又覺得有罪了.不過,每年一次獻祭,可以提醒以色列人知罪,認罪；一年之中,心裡有些時的平安.
\newline
\newline
「及至時候滿足,神就差遣祂的兒子......要把律法以下的人贖出來,叫我們得著兒子的名分.」(加四4-5)在當時來說,時候未滿,救恩的日子還未來到,相距一千五百多年.這段悠久的時間,神為可憐犯罪的人安排了獻祭的方法,來解決罪的問題.但祭物只能暫時安於現狀,不能使良心永遠得著平安.雖不能說它有絕對的價值,然亦有其相當的用處.
\newline
\newline
當第一次世界大戰時,加拿大很多青年人從軍入伍,到歐洲與德國作戰,我的父親也在其中.那時候我才九歲.臨別之前,他帶我和小弟弟並兩個小妹妹同走一程.記得爸爸用手撫摸我的頭說:「孩子啊!你雖然還小,但是我離家以後,家裡又沒有佣人；一定有很多事情差你去做,你要好好地幫助媽媽,做個好孩子.」從此爸爸就和我們永別了.之後,媽媽每天打發我到伙食店去,她把所要的東西,列成一單,簽上名；讓我交給老板,老板就憑字條交物給我.老板不但沒有向我要錢,有時候他還另外送給我一包糖果；叫我帶回家交給媽媽,分給我和弟妹們吃.當時我覺得非常希奇,別人購物都付錢,我卻不必,憑一張紙條便可得回許多東西.到了我年紀漸長,才明白原來每到月底,陸軍總部發薪,父親寄來支票以供家用,所以母親就帶著支票去伙食店老板結賬.是父親未出發之前和老板商妥的.過了三四十年,我反覆思想,覺得這件事是個很好的比方.老板好比慈愛的天父,一切恩典都是從祂來的.家父好比主耶穌基督,創立世界以前祂已與天父商妥了.紙條好比以色列人所獻的祭物(牛羊鴿子等)紙條暫時可以使用；可以買到很多東西,但不能作為付清欠賬之用.則如獻祭,只能暫時遮蓋罪,從天父得到暫時赦罪之恩.不至滅亡,神的憤怒不至臨到他們,但終不能澈底除罪.所以,基督到世上來的時候,就說,神啊!祭物和禮物是你不願意的,你曾給我預備了身體.又說,我到世上,不是要叫人服事,乃是要服事人；並且要捨命,作多人的贖價.
\newline
\newline
獻祭可說是時候還未滿足的過渡辦法,耶穌既已到世上來,且死在十字架上,完成了救贖；人若仍守舊法到聖殿獻祭,神就不收納了.因為基督已獻了一次永遠的贖罪祭,祂一次獻祭,便叫那得以成聖的人永遠完全.
\newline
\newline
我在菲律賓時,每年要到市政府去領新的身份證,一連兩年,都是在同一地點的辦公室窗口處排隊等候,花了五角錢便可領到.到了第三年,我照以往兩年的經驗去做,想不到等了很久沒有結果.原來搬遷了.這好像以色列人,老是從利未祭司那裡買到赦罪之恩；耶穌已經了,若仍舊去找摩西和亞倫,是得不到救恩的.許多可憐的希伯來人,他們不夠認識耶穌,還守住猶太教的窗戶,盼望得平安.故此,希伯來書的作者,特別高舉耶穌基督是神的羔羊.用牛羊鴿子獻祭贖罪已成過去了,那不過只是影兒.
\newline
\newline
記得我小的時候,很欣賞馬戲班的表演.馬戲團未來之兩個月前,到處張貼廣告,多姿多采的圖片,富有吸引力,小孩子們都百看不厭.但是等到馬戲團正式來了,廣告便成了過去時的東西.注意力都轉到馬戲班了.
\newline
\newline
當時的希伯來人,他們認識不夠,正如馬戲來了,他們還在看那些破爛的廣告一樣.
\newline
\newline
感謝讚美主!後來他們漸漸多認識耶穌的寶貝；離開他的舊俗,真真正正以耶穌為獨一的救主,一切盼望在祂身上.
\newline
\newline


\newpage

\section{第44屆~港九培靈會~費述凱~第六講　信主的定義和三種鼓勵}
\label{sec:1376}
\textbf{第六講　信主的定義和三種鼓勵}
\newline
\newline
連結: \href{https://www.hkbibleconference.org/session-message/view/1376}{\texttt{https://www.hkbibleconference.org/session-message/view/1376}}
\newline
\newline
\hyperref[sec:1381]{< < < PREV SERMON < < <}
~
\hyperref[sec:toc]{[返目錄]}
~
\hyperref[sec:1323]{> > > NEXT SERMON > > >}
\newline
\newline
希伯來書 11:1-11:40
\newline
\begin{longtable}{cl}
\hline
\hline
章節 & 經文 (和合本修訂版)\\
\hline
11:1 & \begin{tabularx}{0.7\textwidth}{X} 信就是對所盼望之事有把握，對未見之事有確據。 \end{tabularx} \\ \\ \relax
11:2 & \begin{tabularx}{0.7\textwidth}{X} 古人因著這信獲得了讚許。 \end{tabularx} \\ \\ \relax
11:3 & \begin{tabularx}{0.7\textwidth}{X} 因著信，我們知道這宇宙是藉神的話造成的。這樣，看得見的是從看不見的造出來的。 \end{tabularx} \\ \\ \relax
11:4 & \begin{tabularx}{0.7\textwidth}{X} 因著信，亞伯獻祭給神比該隱所獻的更美，因此獲得了讚許為義人，神親自悅納了他的禮物。他雖然死了，卻因這信仍舊在說話。 \end{tabularx} \\ \\ \relax
11:5 & \begin{tabularx}{0.7\textwidth}{X} 因著信，以諾被接去，得以不見死，人也找不著他，因為神已經把他接去了；只是他被接去以前，已討得神的喜悅而蒙讚許。 \end{tabularx} \\ \\ \relax
11:6 & \begin{tabularx}{0.7\textwidth}{X} 沒有信，就不能討神的喜悅，因為到神面前來的人必須信有神，並且信他會賞賜尋求他的人。 \end{tabularx} \\ \\ \relax
11:7 & \begin{tabularx}{0.7\textwidth}{X} 因著信，挪亞既蒙神指示他未見的事，動了敬畏的心，造了方舟，使他全家得救。藉此他定了那世代的罪，自己也承受了那從信而來的義。 \end{tabularx} \\ \\ \relax
11:8 & \begin{tabularx}{0.7\textwidth}{X} 因著信，亞伯拉罕蒙召的時候就遵命出去，往將來要承受為基業的地方去；他出去的時候還不知往哪裡去。 \end{tabularx} \\ \\ \relax
11:9 & \begin{tabularx}{0.7\textwidth}{X} 因著信，他就在所應許之地作客，好像在異鄉，居住在帳棚裡，與蒙同一個應許的以撒和雅各一樣。 \end{tabularx} \\ \\ \relax
11:10 & \begin{tabularx}{0.7\textwidth}{X} 因為他等候著那座有根基的城，就是神所設計和建造的。 \end{tabularx} \\ \\ \relax
11:11 & \begin{tabularx}{0.7\textwidth}{X} 因著信，撒拉自己已過了生育的年齡還能懷孕，因為她認為應許她的那位是可信的； \end{tabularx} \\ \\ \relax
11:12 & \begin{tabularx}{0.7\textwidth}{X} 所以，從一個彷彿已死的人竟生出子孫，如同天上的星那樣眾多，海邊的沙那樣無數。 \end{tabularx} \\ \\ \relax
11:13 & \begin{tabularx}{0.7\textwidth}{X} 這些人都是存著信心死的，並沒有得著所應許的，卻從遠處觀望，且歡喜迎接。他們承認自己在地上是客旅，是寄居的。 \end{tabularx} \\ \\ \relax
11:14 & \begin{tabularx}{0.7\textwidth}{X} 說這樣話的人是表明自己要尋找一個家鄉。 \end{tabularx} \\ \\ \relax
11:15 & \begin{tabularx}{0.7\textwidth}{X} 他們若想念所離開的家鄉，還有回去的機會。 \end{tabularx} \\ \\ \relax
11:16 & \begin{tabularx}{0.7\textwidth}{X} 其實他們所羨慕的是一個更美的，就是在天上的家鄉。所以，神並不因他們稱他為神而覺得羞恥，因為他已經為他們預備了一座城。 \end{tabularx} \\ \\ \relax
11:17 & \begin{tabularx}{0.7\textwidth}{X} 因著信，亞伯拉罕被考驗的時候把以撒獻上，這就是那領受了應許的人甘心把自己獨生的兒子獻上。 \end{tabularx} \\ \\ \relax
11:18 & \begin{tabularx}{0.7\textwidth}{X} 論到這兒子，神曾說：「從以撒生的才要稱為你的後裔。」 \end{tabularx} \\ \\ \relax
11:19 & \begin{tabularx}{0.7\textwidth}{X} 他認為神甚至能使人從死人中復活，意味著他得回了他的兒子。 \end{tabularx} \\ \\ \relax
11:20 & \begin{tabularx}{0.7\textwidth}{X} 因著信，以撒指著將來的事給雅各、以掃祝福。 \end{tabularx} \\ \\ \relax
11:21 & \begin{tabularx}{0.7\textwidth}{X} 因著信，雅各臨死的時候給約瑟的兩個兒子個別祝福，扶著枴杖敬拜神。 \end{tabularx} \\ \\ \relax
11:22 & \begin{tabularx}{0.7\textwidth}{X} 因著信，約瑟臨終的時候提到以色列人將來要出埃及，並為自己的骸骨留下遺言。 \end{tabularx} \\ \\ \relax
11:23 & \begin{tabularx}{0.7\textwidth}{X} 因著信，摩西生下來，他的父母見他是個俊美的孩子，把他藏了三個月，並不怕王的命令。 \end{tabularx} \\ \\ \relax
11:24 & \begin{tabularx}{0.7\textwidth}{X} 因著信，摩西長大了不肯稱為法老女兒之子。 \end{tabularx} \\ \\ \relax
11:25 & \begin{tabularx}{0.7\textwidth}{X} 他寧可和神的百姓一同受苦，也不願在罪中享受片刻的歡樂。 \end{tabularx} \\ \\ \relax
11:26 & \begin{tabularx}{0.7\textwidth}{X} 他把為彌賽亞受凌辱看得比埃及的財物更寶貴，因為他想望所要得的賞賜。 \end{tabularx} \\ \\ \relax
11:27 & \begin{tabularx}{0.7\textwidth}{X} 因著信，他離開埃及，不怕王的憤怒，因為他恆心忍耐，如同看見那不能看見的神。 \end{tabularx} \\ \\ \relax
11:28 & \begin{tabularx}{0.7\textwidth}{X} 因著信，他設立逾越節，在門上灑血，免得那毀滅者加害以色列人的長子。 \end{tabularx} \\ \\ \relax
11:29 & \begin{tabularx}{0.7\textwidth}{X} 因著信，他們過紅海如行乾地；埃及人試著要過去就被淹沒了。 \end{tabularx} \\ \\ \relax
11:30 & \begin{tabularx}{0.7\textwidth}{X} 因著信，以色列人圍繞耶利哥城七日，城牆就倒塌了。 \end{tabularx} \\ \\ \relax
11:31 & \begin{tabularx}{0.7\textwidth}{X} 因著信，妓女喇合曾友善地接待探子，就沒有跟那些不順從的人一同滅亡。 \end{tabularx} \\ \\ \relax
11:32 & \begin{tabularx}{0.7\textwidth}{X} 我還要說甚麼呢？若要一一細說基甸、巴拉、參孫、耶弗他、大衛、撒母耳和眾先知的事，時間就不夠了。 \end{tabularx} \\ \\ \relax
11:33 & \begin{tabularx}{0.7\textwidth}{X} 他們藉著信，制伏了敵國，行了公義，得了應許，堵住了獅子的口， \end{tabularx} \\ \\ \relax
11:34 & \begin{tabularx}{0.7\textwidth}{X} 滅了烈火的威力，在鋒利的刀劍下逃生，從軟弱變為剛強，爭戰中顯出勇猛，打退外邦的全軍。 \end{tabularx} \\ \\ \relax
11:35 & \begin{tabularx}{0.7\textwidth}{X} 有些婦人得回從死人中復活的親人。又有人忍受嚴刑，拒絕被釋放，為要得著更美好的復活。 \end{tabularx} \\ \\ \relax
11:36 & \begin{tabularx}{0.7\textwidth}{X} 又有人忍受戲弄、鞭打、捆鎖、監禁、各等的磨煉； \end{tabularx} \\ \\ \relax
11:37 & \begin{tabularx}{0.7\textwidth}{X} 他們被石頭打死，被鋸鋸死，被刀殺，披著綿羊山羊的皮各處奔跑，受貧窮、患難、虐待。 \end{tabularx} \\ \\ \relax
11:38 & \begin{tabularx}{0.7\textwidth}{X} 這世界配不上他們，他們在曠野、山嶺、山洞、地穴，飄流無定。 \end{tabularx} \\ \\ \relax
11:39 & \begin{tabularx}{0.7\textwidth}{X} 這些人都是因信獲得了讚許，卻仍未得著所應許的， \end{tabularx} \\ \\ \relax
11:40 & \begin{tabularx}{0.7\textwidth}{X} 因為神給我們預備了更美好的事，若沒有我們，他們就不能達到完全。 \end{tabularx} \\ \\
[1ex]
\hline
\hline
\end{longtable}
「我們既有這許多的見證人,如同雲彩圍著我們,就當放下各樣的重擔,脫去容易纏累我們的罪,存心忍耐,奔那擺在我們前頭的路程.」(十二1)這裡說的是十一章所題的,約二三十位信心的英雄.他們如同雲彩圍著我們.
\newline
\newline
第十一章論及到舊約時代許多愛主的人,他們雖然經歷很多苦難,卻一直保持他們的信心.並詳述舊約時代的信徒,如何的倚靠神,一直仰望神,甚至死也不放棄他們的信仰.許多人雖然今生得不到所盼望的,但他們不灰心,不喪膽,實在有信心仍待來生得到.
\newline
\newline
第十二章所論到的就是那班信心英雄的見證.作者在十一章未報告那些見證之前,他先下了信心定義,信心是甚麼?信主的人是怎樣的?誰是信主的?
\newline
\newline
信心是甚麼?「信就是所望之事的實底,是未見之事的確據.古人在這信上得了美好的證據.我們因著信,就知道諸世界是藉神話造成的.......」(十一1-3)因他們有信心,神就公開的與他們同在,公開承認他們是祂的僕人,因著他們的信心,神就給他們美好的證據.
\newline
\newline
挪亞聽從神的命令造方舟,雖很多人取笑他說:「沒有海洋江河,大船何用之有?」但挪亞並不因此而停止造方舟.因為他得了神的啟示.
\newline
\newline
神說:除了挪亞一家之外,世人都敗壞了行為,我後悔造人在世界上.世人的盡頭到了,再過百廿年,我要開天上的窗戶,地的泉源；傾盆的大雨要降下四十晝夜,大山都被遮蓋；凡有氣息的,連人帶飛禽走獸一併毀滅.惟挪亞在我們面前是個義人,我要保全他和他一家人的生命.世上的動物,我也要保存,你要為我造方舟,你和你一家,并各樣動物,都要進入方舟,躲避洪水.挪亞因信神的啟示,就在地上造了方舟.
\newline
\newline
中國人的「船」字實在很有意思.「船」是由「舟八口」合成的.挪亞合家,就是挪亞夫婦和他們三個兒子三個媳婦,正是八口之家在方舟,得以促使生命.
\newline
\newline
挪亞相信神的話,他的八口之家,在方舟上得救.挪亞時代的人,不肯悔改,不離開假神；不歸向耶和華.結果,盡都滅亡.
\newline
\newline
挪亞在信上得了美的證據.他照著神的吩咐造成了方舟.到了神所定的日期滿了,神吩咐挪亞全家進入方舟；並吩咐要帶潔淨的動物七公七母,不潔淨的動物一公一母；空中的飛禽,走獸,昆蟲,都進入方舟.唯有背逆的人,都不肯進去,因為他們不信神的話.挪亞遵照神的吩咐準備好了,神就把方舟的門關閉了；接著雷電交加,淵源裂開,水勢浩大.不信神的人至此才信.他們見了「事實」,不得不信,而不是信神.他們看見危險臨到,才相信方舟才是安身之所；於是拚命打門,呼救,但方舟的門已關閉了,任何人再也進不去了,他們只有悲哀,哭號,等待淹斃!
\newline
\newline
挪亞相信神的話,沒有看見事實之先,就相信了.他相信神是信實可靠,是永不撒謊的神.挪亞具有真正的信心,所以神與他同在,悅納他,他因著信得了美好的證據.現在,一九七二年八月六日,九龍培靈會這一天「我們因著信,知道諸世界是藉神話造成的；這樣,所看見的,並不是從顯然之物造出來的.」(十一 3)
\newline
\newline
學生進學校讀書求知識,聽老師講課,就「知道」.基督徒讀聖經,神的話告訴我們,我們因著信就「知道」.假使大學教授所講的,與聖經相符合,我就相信,如果與聖經不符,請問你相信不相信?你是信教授,還是信聖經?我相信聖經,我知道諸世界是藉著神的話,從「無」造出「有」來.
\newline
\newline
近數年來,美國耗資四十億美元,派太空人登上月球；目的想從月球上的岩石,研究太空是如何形成的.太空人帶了一本聖經登月,那本聖經是特別影印而成的,全部聖經約三張郵票之大小.我莫明他們為何諸多費心,他們是否相信聖經的第一句「起初神創造天地」這話呢?如果他們不相信聖經,為何要帶聖經登月?如果他們相信聖經；為何要花費鉅資,登月研究些石塊呢?
\newline
\newline
聖經告訴我們「起初神創造天地.」基督徒因信就知道神創造天地萬物.這是信心,不是盲從,是聖經告訴我們,是有憑有據的,去相信神的話語是可靠的.
\newline
\newline
從前有個聰明的人,他希望研究科學,他對航空工程很有興趣,並且悉心研究,懂得飛機是靠著空氣的浮力,能夠在空中飛行的原理.有一天,他準備好了一切,擬乘飛機.當他看見了飛機,忽然懷疑起來了.他想,難道巨大沉重的飛機能浮起在空中飛行嗎?他就改變主意不乘飛機了.他只懂得原理,相信原則,卻不相信實際的飛機.他只頭腦相信原理,但他不肯靠飛機.有的人頭腦信耶穌,說,基督教的道理是好的,但卻不肯信靠耶穌.有的人在平安無事,一切順利的時候信靠耶穌；當環境改變,難處來到,他的心就搖動,疑惑,不信了.
\newline
\newline
亞伯拉罕將近百歲時,他仰望神,一點不疑惑.他的妻子撒拉也將近九十歲,這對年紀老邁的夫妻,按生理學,根本沒有生育的可能；但他們依靠永生的耶和華,生命的主,相信神的應許,「到了百歲的時候要得一個兒子」.這樣的信心,才是真正的信心.
\newline
\newline
我們在教會裡面多年,有沒有信心?配不配稱為信主的人?
\newline
\newline
英國有位出名的,信心偉人,名叫穆勒.他創辦了一間孤兒院.某次,在天氣嚴寒的一個早晨,院裡的小孩子們,都到膳堂預備吃早餐,桌上餐具陳列整齊如常,但是卻沒有麵包食物.全體照常謝飯:「感謝主!為我們預備今天豐富的飲食,奉主聖名,阿們.」禱告完畢,聽見有扣門的聲音,原來某麵包廠的送貨車,正在孤兒院門口；因車壞了沒法開動,滿載麵包；既不能及時送出,又不能回廠；所以只好把全車的麵包,都送給孤兒院.穆勒先生再次低頭:「謝謝天父宏恩,顧念兒女的需要；不早也不遲,賜下最好最新鮮的麵包.」這才是真正的信心!
\newline
\newline
我們是不是信徒?配不配稱為信主的人?
\newline
\newline
有位親愛的同工,昨晚向我提個很好的建議；在這次培靈會中講罪的問題.我覺得這意見很對,雖然來赴會的,多數是基督徒,但是,許多基督徒滿有問題；或有不好的習慣,一直改不過來；或懷恨不饒恕人,或愛世俗,愛虛榮,或背後批評人,假冒偽善的.......
\newline
\newline
一般基督徒很可能犯的最大的罪是甚麼?聖經告訴我們,那叫我們至於下地獄永遠死亡的罪,是最大的罪.是不得赦免的罪,就是拒絕聖靈對於耶穌的見證.聖靈感動我們的心,叫我們接受耶穌,作個人的救主.我們卻硬著心,消滅聖靈的感動,不肯接受.這是不可赦免的罪,因為除耶穌以外別無拯救；除了耶穌,沒有別的救法.耶穌說:「我實在告訴你們,世人一切的罪,和一切褻瀆的話,都可得赦免；凡一切褻瀆聖靈的,卻永不得赦免,乃要擔當永遠的罪.」(可三 28-29)「這話是因為他們說,他是被污鬼附著.」(可三30)
\newline
\newline
法利賽人由聖靈感動,看見耶穌行神跡奇事,他們的良心說:「可能祂是彌賽亞,是基督.」但他們違背良心又說:「他不是基督,是撒旦,是鬼王幫助他的.」他們把聖靈的工作歸功於魔鬼,褻瀆聖靈,褻瀆主耶穌；消滅聖靈的感動,拒絕耶穌的救恩.這是最大的罪!
\newline
\newline
我相信,今晚在坐中多數的弟兄姊妹,至每一位,永遠不會犯這罪吧!聖靈已經在你們心裡作見證,因著聖靈不斷地感動,你的心漸漸軟化了,到了最後,你抵擋不住說:「主耶穌啊!除你以外別無救法.」我廿一歲時說了上面這句話.聖靈感動我三個月,有位老姊妹一直為我禱告.起初我有自己的計劃,我想發財,所以屢次拒絕.那位老姊妹愈加懇切為我禱告,聖靈的感動也愈加迫切；最後,我不敢再拒絕了.我到禮拜堂去,我對主說:「主啊!我需要你的救恩,現在我願意,但我不知道應當怎樣信主.」當時,有位聖經學院的老院長,受了聖靈的感動,走到我身邊問說:「孩子啊!你是否基督徒?」我說「不是.」他又問「你願否做基督徒?」我心裡掙扎,後來我想,這是我的機會,我不敢再錯過這機會.我就求主赦免我的罪,我知道主為我死在十字架上；我不再消滅聖靈的感動,我出死入生得救了.次日,我向主人借了匹快馬跑回家.一進家門,就對我母親說:「我得救了.」我母親卻愕然莫明其妙!後來,我向家人作見證,我的老母親,弟弟,和五個妹妹,都一齊接受主了.
\newline
\newline
「耶和華如此說,依靠人血肉的膀臂,心中離棄耶和華的,那人有禍了.因他必像沙漠的杜松,不見福樂來到；卻要住曠野乾旱之處,無人居住的鹹地.倚靠耶和華,以耶和華為可靠的,那人有福了.他必像樹栽於水旁,在河邊扎根,炎熱來到,並不懼怕；葉子仍必青翠,乾旱之年毫無掛慮,而且結果不止.人心比萬物都詭詐,壞到極處,誰能識透呢?」(耶十七5-9)基督徒第二大罪是倚靠人,不倚靠神,倚靠物質,而不倚靠造物之主.羅馬書說:「他們敬拜受造之物,而不敬奉造物的主.」基督徒也是如此!
\newline
\newline
我服事主四十多年,我可為主作見證,祂實在是可靠向主.人不可靠,錢財也不可靠,正如北方有句俗話說:「靠誰誰跑,靠山山倒.」任何人,物,都不可靠.
\newline
\newline
基督徒發現奇妙的秘訣,就是堅心倚賴耶和華.保羅說:「我知道我所信的是誰.」多多知道我們的主,就更容易相信我們的主,不至於犯這大罪.懷疑神,是叫神傷心的罪.耶穌常對祂的門徒說:「你這小信的人啊,為何疑惑呢?」耶穌也同樣對我們說:「你這小信的人啊,多少時候我聽了你們的禱告；你為何還懷疑,猶豫呢?難道你們的信心死了嗎?」
\newline
\newline
我們要追求認識主,才能夠容易相信主.
\newline
\newline


\newpage

\section{第44屆~港九培靈會~陳終道~第一講　基督的性格}
\label{sec:1323}
\textbf{第一講　基督的性格}
\newline
\newline
連結: \href{https://www.hkbibleconference.org/session-message/view/1323}{\texttt{https://www.hkbibleconference.org/session-message/view/1323}}
\newline
\newline
\hyperref[sec:1376]{< < < PREV SERMON < < <}
~
\hyperref[sec:toc]{[返目錄]}
~
\hyperref[sec:1356]{> > > NEXT SERMON > > >}
\newline
\newline
馬太福音 11:28-11:30
\newline
\begin{longtable}{cl}
\hline
\hline
章節 & 經文 (和合本修訂版)\\
\hline
11:28 & \begin{tabularx}{0.7\textwidth}{X} 凡勞苦擔重擔的人都到我這裡來，我要使你們得安息。 \end{tabularx} \\ \\ \relax
11:29 & \begin{tabularx}{0.7\textwidth}{X} 我心裡柔和謙卑，你們當負我的軛，向我學習；這樣，你們的心靈就必得安息。 \end{tabularx} \\ \\ \relax
11:30 & \begin{tabularx}{0.7\textwidth}{X} 因為我的軛是容易的，我的擔子是輕省的。」 \end{tabularx} \\ \\
[1ex]
\hline
\hline
\end{longtable}
有一個家庭裡,父親與兒子很少見面,兒子清早上學,父親晚上才回家.有一天朋友請他們父子倆吃飯,並問這兒子,父親喜愛吃些甚麼,兒子說不知道；再問他父親愛吃鹹的,甜的,還是辣的,兒子又說不知道.這朋友就生氣的說:「你這作兒子的,怎麼總是不知道呢!難道父親是肥,還是瘦,你也不知道嗎?」以上是一個故事或是真有其事,我不敢肯定,不過今日許多基督徒對主基督的認識也是如此,只知有父子關係──與主生命上之關係──但卻停留在此關係上,對主毫無認識.今日的基督徒該更多的認識基督,更多的像基督；有祂的榮美,有祂的性格.從馬太福音中,我們可以看見主的性格:
\newline
\newline
(一)基督是溫柔的──(太十一28-30,廿一5.)
\newline
\newline
在太21章中記載主騎驢進京,基督是君主,是溫柔之君.
\newline
\newline
(太十一29)節中主說:「我心裡柔和謙卑.」基督不單是外表溫柔的騎驢,而且是心裡也一樣的溫和.一個人是否溫柔,不能單靠外表來判定,該看他被激動時是否會暴燥,抑或仍存喜樂的心感恩.
\newline
\newline
當基督自己說心裡柔和的時候,祂是在何等情況之下呢?我們查考(太11章)上文時就可以清楚知道,當時主基督正是在傳道毫無效果的情況之下.(太十一16至24)節中,就是當時的情形,主說:「我用什甚麼比這世代呢?好像孩童坐在街市上,招呼同伴說,我們向你們吹笛,你們不跳舞；我們向你們舉哀,你們不捶胸；約翰來了,也不吃,也不喝,人就說他是被鬼附著的.人子來了,也吃,也喝,人又說他是貪食好酒的人,是稅吏和罪人的朋友.但智慧之子,總以智慧為是.耶穌在諸城中行了許多異能,那些城的人終不悔改.就在那時候主責備他們說,哥拉汎哪,你有禍了!伯賽大阿,你有禍了!因為在你們中間所行的異能,若行在所多瑪,他還可以存到今日.但我告訴你們,當審判的日子,所多瑪所受的,比你還容易受呢.」主在此光景之下,祂灰心喪膽嗎?不,祂說:「天地的主,我感謝你!因為你將這些事,向聰明通達人就藏起來,向嬰孩就顯出來.」(太十一25.)在此光景下祂不但發出感謝,而且還說:「我心裡柔和謙卑.」可見基督的溫柔,是從心裡發出的.
\newline
\newline
當主被釘十架時,同釘的兩個強盜,其中一個在譏誚祂.祂是否因此而被激動?不,祂並沒有被激動,因為當另一強盜求主記念他時,祂回答說:「你要同我在樂園裡了.」從祂能如此地聽得清楚,並清楚地回答,可見祂當時並未有被激動.因為一個激動的人,對答是不會如此清楚確實的.
\newline
\newline
至高神的兒子來世被人輕看,這會使祂難過,但卻不動怒.精神雖痛苦,但仍然是溫柔而平靜的.這是基督的性格,溫柔而謙卑.
\newline
\newline
(二)基督是不爭競也不喧嚷的──(太十二18至21)
\newline
\newline
「祂不爭競,不喧嚷,街上也沒有人聽見祂的聲音.」
\newline
\newline
這好像是形容祂的動作,其實乃是基督的性格,難道主基督在街上沒有聲音嗎?祂在街上也有講道的,怎會在街上沒有人聽見祂的聲音呢?祂不爭競,不喧嚷,是描寫祂的性格是隱藏的；不愛出風頭,也不愛到處引人注目的.為甚麼當時的人們,提到主基督時總愛說:「這不是那木匠的兒子嗎?祂的兄弟不是在我們中間嗎?祂的妹妹們不就在這裡嗎?他們就嫌棄祂.」(可六3)為什麼他們只見到祂是木匠的兒子呢?乃因祂沒有那神氣十足的派頭,不像法利賽人,在衣服上繡上有經文.主基督是與普通人一樣,故世人憑外貌不認識祂.「智慧如同深水,惟明哲人才能汲引出來.」(箴廿5).主就像深水,祂並非愛顯露自己,連祂的兄弟也不了解祂.有一次過節的時候,祂的兄弟對祂說,如果人要顯揚名聲,就該趁這機會了.他們實在不明白耶穌!主所到的地方,有許多人跟從祂；但並非祂故意要顯揚自己,不過無法隱藏而已.如果我們有主豐盛的生命,就不必擔心別人會不注意你,也不必擔心你有無神氣的架子；因為自然就有人發現你,你要隱藏也不可能了.正如一個好醫生,也有真工夫,就算沒有掛牌行醫,總會有人找尋他求治病的.工夫好連門也不必裝修,不必用廣告宣傳,也會有人找上門來.我們的主是不做門面工夫的,街上也沒有人聽見祂的聲音；不吹號也不打鑼,更沒有替自己宣傳.雖然如此,但還是隱藏不住的!但願基督的生命,充滿在我們裡面,讓我們有基督的榮美,有基督的生命吸引人.
\newline
\newline
(三)基督是注意實際的──(太十二7至8.)
\newline
\newline
「我喜愛憐恤不喜愛祭祀」,主責備當時的人說:如果你們明白這話的意思,就不會將無罪的當作有罪了.憐恤是神的性格,是一種恩賜,也是一種品德.祭祀乃是宗教的禮儀,是叫人認識這一位應該敬畏的神.祭祀是告訴我們,神要刑罰罪惡,祭祀亦教導我們認識神,祂是滿有憐恤的.我們雖然有罪,但神要藉祭牲擔當我們的罪,故祭祀教導我們認識神的公義和神的憐恤.主耶穌基督就是神所打發來的救主:「看哪!神的羔羊,背負世人罪孽的.」這就是神憐恤的記號,教我們認識基督乃是救主,在整卷舊約中,祭祀之教導是教我們認識基督,乃是神為憐恤我們,而為我們預備的救主.但當時會幕內的人,不明白憐恤的真正意義,因從前猶太人已漸漸注重祭祀的禮儀,而忘記憐恤的意義.故神對以色列人說:「我喜愛憐恤,不喜愛祭祀.」我所喜愛的是實際的品德,而不是宗教之儀式,我要的是憐恤的性格,而不是裝模作樣的一套,所以基督是注重實際性格,非注重外表.基督教是有儀式,但儀式的地位,並不重於實際.
\newline
\newline
(太15章中.)有一次門徒吃飯時不洗手,法利賽人批評說:為甚麼要違反古人的遺傳,而用俗手吃飯?主說:你們為何因遺傳而違背神的誡命?神說當孝敬父母,而法利賽人說,我所供養父母的都已將之奉獻給神了.這藉口多麼好聽,他們就利用這作藉口而不須孝敬父母了.故神說:「這百姓嘴唇尊敬我,而心卻遠離我,拜我也是枉然.」主說這敬拜是假的,神是要心靈誠實敬拜祂.主耶穌注重內心之實際,而非外表之藉口.在(太9章)記載稅吏馬太請主耶穌吃飯,他們就批評主與罪人吃飯.這有甚麼關係呢?是否同罪人吃飯,就變成罪人呢?主與罪人吃飯,不但沒有成為罪人,而且更使他成為主的門徒.這可見法利賽人,只重外表而已.他們不與罪人吃飯,乃是要表示自己的清高；其實,與你吃飯後,能變成何種人才是最重要的.假如吃飯後仍然依舊,那就算多吃幾頓飯也是一樣的.法利賽人請主吃飯,而不與罪人吃飯,是要表示他們與主同等.主耶穌是注重實際的憐恤,見別人生病了,飢餓了,難道不給他們吃,不醫治他們而再留待明天嗎?所以主耶穌的性格是重實際而非裝成屬靈的外表.但願主的靈充滿我們,使我們有主的性格.
\newline
\newline
(四)基督是恨惡假冒為善的──(太23全章,參廿一12至14.)
\newline
\newline
(太23章)乃是主在世時最嚴厲之責備,責備法利賽人假冒為善.可見主乃誠信真實的主,祂雖溫柔,但恨惡罪,恨惡虛假,斥責虛假.祂溫柔,但並不容讓罪惡；祂是公義的,但沒有把公義作為暴躁生氣.主耶穌是公義的,但也是絕對溫柔的.試看當祂入聖殿時,見有賣牛羊的,兌換銀錢的......他就推翻桌子...... (太21章),又在(太23章13至31節)中說法利賽人的七禍,可見主基督之性格,非常憎惡假冒為善.
\newline
\newline
我們作基督徒的應記著,如果做一個真正的基督徒,就要真正的像主,要實實在在的,不可有私毫裝假.最實在的人就是最屬靈的人,因為這才最像主耶穌.
\newline
\newline
(五)基督是良善的──(太十九16,十16,十四27-31,十七27.)
\newline
\newline
在(太19章)中有一少年官見主,稱主為良善夫子.主曾教導我們說:我差你們去,如羊進入狼群,要靈巧像蛇,馴良像鴿子.我們的主雖然靈巧像蛇,滿有能力智慧；但祂是馴良的,充滿了良善的性格.
\newline
\newline
(太14章中),主在海面行走,門徒未看清楚,以為是鬼怪,後來清楚知道是主.彼得就請求也要在水面行走,但當他見了浪大而沉下,並喊叫主救他時；主就立刻拉住他,然後責備他.請注意主耶穌並非先責備他才拉起他,而是先救起了才責備他.主是良善的,門徒是無知,把主當作鬼怪,又沒有信心,但主的性情是良善的.
\newline
\newline
(太17章中),收稅的人向主收稅,主問彼得說:「世上君王誰收稅?向外人抑或向兒子?」彼得說向外人,意思是說主是萬有神的兒子,可以免稅了.但主對彼得說:「恐怕觸犯他們」,為甚麼?難道是怕嗎?是的,祂是怕別人不了解以至生氣,而至對主有懷疑,甚至使人跌倒.可知主有良善的性格.
\newline
\newline
我們今日是主的門徒,是否有主的性格?我們對人記恨在心?找機會報仇?別人有禍就快樂?我們是主的門徒,應該跟從主,是從主而生,在我裡面有主的性格及性情；祂的心,祂的靈,在我們裡面活著.今天查考聖經,不但讓我們多知,而且要多像祂,活著有基督的性情.有人說:我生來是如此,江山易改,品性難移.你如未信主,可以如此說；但你是信徒,有主在你裡面,又怎會不能改?有主的榮美在裡面,就應該活出來.
\newline
\newline


\newpage

\section{第44屆~港九培靈會~陳終道~第二講　基督的工作}
\label{sec:1356}
\textbf{第二講　基督的工作}
\newline
\newline
連結: \href{https://www.hkbibleconference.org/session-message/view/1356}{\texttt{https://www.hkbibleconference.org/session-message/view/1356}}
\newline
\newline
\hyperref[sec:1323]{< < < PREV SERMON < < <}
~
\hyperref[sec:toc]{[返目錄]}
~
\hyperref[sec:1357]{> > > NEXT SERMON > > >}
\newline
\newline
馬可福音 6:34-6:34
\newline
\begin{longtable}{cl}
\hline
\hline
章節 & 經文 (和合本修訂版)\\
\hline
6:34 & \begin{tabularx}{0.7\textwidth}{X} 耶穌出來，見有一大群的人，就憐憫他們，因為他們如同羊沒有牧人一般，於是開始教導他們許多事。 \end{tabularx} \\ \\
[1ex]
\hline
\hline
\end{longtable}
「耶穌出來,見有許多的人,就憐憫他們,因為他們如同羊沒有牧人一般,於是開口教訓他們許多道理.」(可六34.)
\newline
\newline
「只是坐在我左右,不是我可以賜的,乃是為誰預備的,就賜給誰.」(可十40).
\newline
\newline
在馬可福音中,論及主是僕人,僕人的本份是工作,(在可六34)的經文中,是耶穌看見許多人,如同羊沒有牧人一般,祂就憐憫他們.(又在可十45 )說:「人子來,並不是要受人的服事,乃是要服事人,並且要捨命,作多人的贖價.」這兩節聖經合起來,正如(約翰十14)的意思:就是「我是好牧人,好牧人為羊捨命.」我們該知道,一個牧人好與否,應看他是否願意為羊捨命,對工作是否忠心,甘願為羊捨命.現在我們看看主的工作如何:
\newline
\newline
(一)主工作的權能──(可一22)
\newline
\newline
「眾人很希奇祂的教訓,因為祂教訓他們,正像有權柄的人,不像文士.」
\newline
\newline
這些聽眾怎能聽出主的教訓有權柄?他們又怎會分辨出主的教訓和文士的不一樣?這些文士可說是研經專家,他們是有學識的,有考據的.聽眾聽了主所講的,和文士所講的之後,心裡會有一種敏銳的感覺,就是主所講的帶著權柄和能力,而文士卻沒有.因這些文士雖有學問,有考據,但他們所能給人的,只是知識；而耶穌基督卻有靈感,有啟示,有聖靈同在.祂所講的不單是知識,而且是信息,是生命的供應,這是主與文士不同之處.聖經不單是學術的書,雖然內裡有學術的價值,但不能單憑學術的眼光,或研究學術的態度,來領會聖經,更重要的是倚靠聖靈.假如聖經是人的著作,就用研究學術態度去研究已足夠,但聖經不是人的著作,祂基本是超然的啟示.所以如要明白聖經的話,就必須要在我們智慧之外,倚靠聖靈之啟示.故主所說的與文士所說的不一樣.主的工作有屬靈的權能,論及工作,就不能不談及權柄和能力.主在世上並沒有地位上的權柄,但祂卻有屬靈的權柄能力.
\newline
\newline
在(可二11)中,我們看見主醫治了一個癱子,他由四個人抬起來,而從房頂垂下的.主就對癱子說:「你的罪赦了.」旁邊的人及文士和法利賽人聽見了,他們心裡便起了議論說:「這人說了僭妄的話,除了神以外誰能赦罪呢?」耶穌問他們說:「我說你的罪赦了,或說起來,拿你的褥子行走,那一樣容易呢?」當然,在我們看來對癱子說:你的罪赦了是非常容易的；因為罪是否赦了,誰都不能清楚知道.但對癱子說,拿你的褥子行走,卻要他真的能起來行走,不然就會立刻顯出,祂的話缺乏能力.主知道他們心裡的意思,必以為主開空頭支票.故主要證明祂的權柄,特地對癱子說:「你的罪赦了!和拿你的褥子起來行走吧.」讓他們知道,兩樣都是容易的.當時,癱子真的起來行走,就可證明主的話有權柄,證明主叫他行走和赦罪都是一樣的實在.
\newline
\newline
在(可五25)中,那患十二年血漏的婦人,她擠在人群中找到了耶穌,她心裡想,只要摸著耶穌,她的病就會好了.當她真的用手摸著了耶穌衣裳的繸子時,耶穌登時覺得有能力從自己身上出去,這婦人血漏的源頭就立刻枯乾了.耶穌基督怎會有能力從身上出來?這能力是怎樣的能力?這並不是聲音,也不是叫喊,亦不是眼淚,而是能力.因此,這女人的病真的好了.這是生命的能力,是祂工作權能內所發出的能力.
\newline
\newline
(可十一)中,主來到耶路撒冷進入聖殿時,趕出殿裡作買賣的人,推倒兌換銀錢之人的桌子......將牛羊趕出.為甚麼那些人不反抗主耶穌?為甚麼他們不發怒?被主用鞭子趕打,怎麼又不還手?乃因主工作有權能,他們都服在祂的權威與能力之下,無人敢反抗.我們如果想在生活上,言語上,工作中能有權能,就要有下列三條件:
\newline
\newline
(a)我們的生活要有見證──有見證就有權能,使徒保羅對帖撒羅尼迦教會說:「我們福音傳到你們那裡,不獨在乎言語,也在乎權能和聖靈,並充足的信心；正如你們知道我們在你們那裡,為你們的緣故,是怎樣為人.」(帖前一5)
\newline
\newline
保羅把福音傳給帖撒羅尼迦的人,為甚麼說不獨在乎言語,而也在乎權能呢?因他在他們當中有美好的見證,正如他們知道保羅在他們中間時,如何的為人.
\newline
\newline
(b)我們所行的事要合乎真理──為甚麼耶穌鞭打那些賣牛羊的人,和推翻兌換銀錢人桌子時,他們不敢反抗?因主所行是在真理方面,我們所行如在真理中；則真理本身就是權威,真理就會使人良心發生共鳴作用.假如我們所說的話是真理,那麼雖然你沒有地位,或是貧窮,或被人反對；但別人良心裡,會發生共鳴,而承認你所說的是對的.因為真理本身,是帶著權威的,我們所行如合乎真理,那我們就有權威.
\newline
\newline
(c)對我們自己所說,所行的,應極熟識和清楚──例如我們傳福音,而自己對福音不清楚,不認識,則所說的話,就沒有能力.對自己所信的福音,要能貫澈通達,那你所說的話,就有權威.因為當你講述時是滿有把握的,並非模稜兩可；那就是說你所說的是極熟識和清楚時,就會有權威.我們的主,祂工作是帶著權能的.
\newline
\newline
(二)主工作的秘訣──(可六45至46)
\newline
\newline
「祂既辭別了他們,就往山上去禱告.」在別的福音書中,我們知道這時人們正要擁護主作王,甚至強迫祂作王.在(可一35)中亦說:主在清早天還沒亮時,就到曠野去禱告.主的工作為何會有權能?祂的秘訣就是與神維持著關係.今日的基督徒會說,我們現今是忙碌時代,與當時不同,其實當時也是一樣忙的,雖然我們今日生活緊張,但亦有許多方法,可把時間節省.例如今日聚會可坐車,但主當時卻要步行很遠的路；身體疲乏,是會影響精神的.聖經記載祂忙碌,甚至沒時間吃飯.
\newline
\newline
在(可六)中,上文說及主用五餅二魚行神蹟之前,祂就沒時間吃飯,就帶門徒退到曠野去歇一歇,但還未歇息,就有許多人來了.主見了他們如羊無牧人一般,就憐憫他們,開口教訓他們,然後才行神蹟.那麼,行完神蹟後,就該當休息了是嗎?但祂卻上山去禱告.
\newline
\newline
為何主在世時滿有權能去工作呢?因祂常禱告,因祂也是人,祂與我們一樣,站在人的地位,凡事受試探,只是祂沒有犯罪.主工作有能力之秘訣,就是無論在任何忙碌中,都維持著與神必需有的親密交通,故無論在何地方,或應付敵人,都帶有權柄和能力.
\newline
\newline
(三)主工作的目標──(可一38)
\newline
\newline
「耶穌對他們說,我們可以往別處去,到鄰近的鄉村,我也好在那裡傳道.因為我是為這事出來的.」
\newline
\newline
從這段經文,我們知道主很清楚來世的主要使命是為傳福音,叫人認識神為人所預備的救法.不過,為了傳福音祂也行神蹟,以證明祂的使命是從神那兒來的.祂知道工作主要目標是甚麼,雖然機會好,眾人都來找祂；但主好像故意放棄此機會,而要到別處去傳福音.我們常因工作機會好,而忘卻甚麼是神給我們的主要託付.主知道祂自己來世的使命是甚麼,故祂明白了工作的目標,就能在祂的工作中,分別出主要的和次要的.假如我們不明白主的託付,就不會分別；如不明白自己的工作目標,不知自己在神家中應作何事.若不知道神交給我自己的是何崗位,那麼我們的眼睛就會昏花,就會常忽略了本份的工作；甚至攪擾到別人,干預別人.那我們所作的,就不是神所要作的,於是就成了神家中,被廢棄的器皿!為甚麼主寧願損失兩千頭豬,來拯救一個人呢?在祂心中的計算,這人的靈魂價值不止兩千頭豬,當主用五餅二魚使五千人吃飽時,為甚麼要叫門徒收拾零碎呢?在這裡主有一個屬靈的評價,和屬靈的價值觀；應用時,祂願意花很大的代價；不應用時,祂儘量節省.這是今日教會所應學習的.我們缺乏了屬靈的價值觀,我們應有屬靈的眼光,去評價每一件事.
\newline
\newline
當彼得勸主不要去釘十架時(可八),主斥責彼得說:「撒旦,退我後面吧!」但另一方面,有些家屬說祂癲狂了,因祂忙到忘卻了吃飯,主根本沒有理會他們的說話,毫不在乎似的.從這比較,就可知主心目中,有一個標準；因彼得所說的話,與祂的使命衝突；祂看重使命,故嚴厲的責備彼得.
\newline
\newline
為何主稱讚窮寡婦?在(可12)中說她投下的兩個小錢,比財主所投的更多.主用甚麼眼光來評價這事呢?該知主的稱讚,並非單對寡婦；而是鼓勵那些像寡婦一般的奉獻,亦等於責備那些,像財主一般奉獻的人.試想,這豈不是得罪了財主嗎?如單靠寡婦這兩個小錢的奉獻,又怎麼行呢!不,祂另有標準來衡量這件事,因祂看重屬靈的價值,多於表面的價值.這種價值觀念,完全是因祂認清工作的目標.假如我們認清自己工作目標,認清神將我擺在祂家裡所應負之責任,和所應起之作用,那麼我們有我們的人生觀,我們就有對事情應有評價的看法；如此,才能作出有意義之事.
\newline
\newline
(可八2-3)和(可六)一樣,主在工作中有愛心,祂憐憫眾人,怎能讓眾人餓著回去呢?祂要讓他們吃飽.愛心是不講理由的,愛心是祂能作甚麼,就做在祂所愛的人身上.我們在工作上必需有愛主的心去作,(林前十三)說:我們如沒有愛,則所作的都沒有價值,就算能說萬人的方言,並天使的話語,卻沒有愛,就像鳴的鑼響的鈸一般.如所作的非因主而作,則在神面前不能存到永遠.當主在世時,祂實在用愛心來工作.
\newline
\newline
在(可十)中,少年的官問主該如何作才可以得永生?他自己說,他從小就遵守誡命,主聽了就愛他.不論這少年所說的是否屬實,但主卻覺得他有可愛的地方.我們有愛心否?跟從主三年多的彼得,竟三次不認主,這情況主早已預先知道.所以他說:「雞叫二次之先,你要三次不認我.」在彼得否認主時,主並沒有對彼得說話的機會,但卻預備了那雞讓祂叫兩次提醒彼得.這是主的愛,祂在工作裡滿有愛心.保羅工作之成功亦在此,他滿有基督的愛,在他的責備裡充滿了愛.他對哥林多的人說:你們學基督的,作師傅的雖有一萬,為父的卻不多.保羅以父的心,去責備哥林多人；所以他們雖受責備,但仍受感動.
\newline
\newline
(四)主工作的智慧──(可十二)中,可以見到主的智慧有三方面:
\newline
\newline
(a)有人試探主,可否納稅給該撒?主就拿上稅的銀錢問道,這像和這號是誰的?他們答說是該撒的,於是主說:「該撒的物應歸給該撒,神的物應歸給神.」眾人希奇祂的說話.顯明主在應付事情上有機智.
\newline
\newline
(b)祂對屬天的奧秘,有人所不知道的智慧.撒都該人不信復活,他們問主說,人若死了,沒有孩子,那麼照猶太人的風俗,他兄弟應娶他的妻,為哥哥生子立後.從前在我們這裡,有兄弟七人,先後都娶過這女子作妻,後來他們都死了；那麼,在復活時,這女人是誰的妻呢?主耶穌回答說:「你們都錯了,因你們不明白聖經,也不曉得神的大能；當復活時,人不娶,也不嫁,好像天上的使者一樣.」(太廿二29)意思是說人在復活後,就沒有夫妻關係的存在；而夫妻關係,只不過在今世存在.到復活時,是進到神的大家庭裡,我們是以神的大家庭,作為我們的觀念,而非以地上的小家庭作觀念.到那時我們觀念改了,人也復活了,根本就沒有這些問題存在了.今日的基督徒,若因妻死而娶繼室,上述的事,也等於主對你回答的問題了.
\newline
\newline
(c)有人問主誡命中那一條是最大的?舊約的律法與誡命,猶太人分成六百多條,問這許多條誡命中何者為最大呢?假如不通曉和不明白,那實在難以解答.主回答說:你要盡心,盡性,盡意,盡力愛主你的神；其次,就是愛人如己.這就是一切律法和先知道理的總綱.(太廿二37-40)可見主對真理,具有屬靈的智慧.這關乎屬天的奧秘事,需有聖靈的啟迪和通達真理,才能認識；那麼,也只有主才能回答了.當主十二歲時在聖殿一面聽,一面問；祂對狡猾的敵人,非常機智,正如祂所教訓我們的:「我差你們去,好像羊進入狼群；所以你們要靈巧像蛇,純良像鴿子.」祂的確純良像鴿子,但並非糊塗愚笨的.主要我們作誠實的人,即心地善良的人,但並非作傻子而被人騙!祂有充足的智慧,知道如何去應付敵人的詭計；但卻不用詭計去對付人,這就是主工作的智慧.
\newline
\newline


\newpage

\section{第44屆~港九培靈會~陳終道~第三講　基督的禱告}
\label{sec:1357}
\textbf{第三講　基督的禱告}
\newline
\newline
連結: \href{https://www.hkbibleconference.org/session-message/view/1357}{\texttt{https://www.hkbibleconference.org/session-message/view/1357}}
\newline
\newline
\hyperref[sec:1356]{< < < PREV SERMON < < <}
~
\hyperref[sec:toc]{[返目錄]}
~
\hyperref[sec:1358]{> > > NEXT SERMON > > >}
\newline
\newline
「眾百姓都受了洗,耶穌也受了洗；正禱告的時候,天就開了；聖靈降臨在祂身上,形狀彷彿鴿子.又有聲音從天上來說:你是我的愛子,我喜悅你.」(路三21-22)
\newline
\newline
耶穌為甚麼需要禱告?因祂站人的地位上.若按祂的神性而言,祂是不需要禱告的,祂禱告,就可證明祂謙卑；祂反倒虛己降卑為人,否則祂不需要禱告的.路加講及耶穌基督為人子,故多記祂的禱告,除有關於禱告之教訓及比喻之記載外,還單獨記了主的八個禱告,是別的福音書所沒有的:
\newline
\newline
(一)受浸時的禱告──(路三21-22)
\newline
\newline
經文記載主受洗時的禱告,在別的福音中雖亦記主受洗；但只有路加如此說:「正禱告的時候,天就開了,聖靈降臨在祂身上,形狀彷彿像鴿子,又有聲音從天上說:你是我的愛子,我喜悅你.」天何時為主而開?乃正禱告時,因何主的禱告會那麼容易叫天打開?那麼容易使聲音從天上來?又能使聖靈印證祂是神的愛子?我們該注意到上文說:主到施洗約翰前受洗,約翰說:我應受你的洗,為何反到這裡來,在別的福音書中我們可知主當時對約翰說:「我理當盡諸般的義.」這「義」字在小字上或譯作「禮」,這意思是宗教方面的禮儀.在認識耶穌基督非律法而來的義.保羅說:他認識耶穌基督非律法而來的義,乃是因信基督而來的義.所謂律法的義,即按律法的規矩而遵守.主說要盡諸般的義,所以要求約翰為祂施洗.有人問:為何耶穌要受洗?該知道約翰為祂施洗.有人問:為何耶穌要受洗?該知道約翰所傳的,乃悔改的洗,耶穌是不需要悔改的.那麼為甚麼要到約翰前受洗?我們應注意,耶穌之受洗並非悔改,乃是盡諸般的義.那麼律法的規矩及禮儀,是屬外表之儀式,但卻象徵屬靈之意義,那是次要的而非主要的；但主到世上來連次要的,附帶的也都去遵守,可見祂並非勉強,而是甘心樂意地降卑.主要的目的是為我們上十字架,否則祂不會在人看來是次要的也去遵守.故主對約翰說,理當盡諸般的義.這不但是主為我們生在法律之下,盡了律法下各種的禮儀；而且祂是非常甘心樂意降世,祂如此甘心降卑順服,就使到祂的禱告,使天為祂而開了.為甚麼我們的禱告常覺得,天像銅,地像鐵一般不易開呢?是否我們在神面前未甘心降卑,還未像我們的主耶穌?
\newline
\newline
(二)退到曠野的禱告──(路五16)
\newline
\newline
「有極多的人聚集來聽道,也指望醫治他們的病,耶穌卻退到曠野去禱告.」這是在許多人擁戴主,歡迎祂的時候；多人對祂有著指望的時候,主卻退到曠野去.當許多人在我們身上存著很大的期望,將你抬得又高又大的時候,我們是否像主退到神前禱告?為甚麼當時主要退到曠野禱告?祂並非無能力或心靈疲乏,而是正行神蹟的時候.祂的能力並無減少,祂乃要在許多人擁戴的氣氛當中,不容許這些擁戴祂的人,有絲毫摸到祂的心.祂要退到神面前,要向神有專一的忠誠和心志,完全不讓世界的虛榮及稱譽影響祂,故祂要退到曠野去.舊約中的但以理,他是在惡劣的環境中作宰相.在不尋常的政治場合裡,和對信仰不利的環境中；他怎能站立得住,怎能不失落他的信仰呢?他不但不失落,反能為神作美好的見證.在但以理六章中說:當大利烏王被人欺騙,而下命令,禁止人(除了王之外),不准向神或向人禱告時,但以理卻把窗戶打開,一日三次向神禱告,和素常一樣.他能在惡劣的環境中,維持他的見證；乃因他有素常的禱告,經常到神面前；有許多事,是世界虛榮的事,也會在教會之內；這就是人的虛榮和人的稱讚,這很容易使人的眼睛昏花了.使我們昏花的並非人的苦難或反對,也非人的指摘,而是人的稱許.特別是當你成功時,他們存著希望要在你身上,想得著些甚麼；這種光景,最易使人眼睛昏花.但主的眼睛並沒有昏花,而是退到曠野禱告.當我們看到別人成功和被別人注意時,我們的心會被搖動；羨慕這種光景,結果我們的腳步,就會走到這世界路上.貪圖虛浮的榮耀,這一切就可能會摸著我們的心.所以我們該效法主,在試探來到之先,在面臨試探之時,就退到曠野去.
\newline
\newline
(三)選立使徒的禱告──(路六12-13)
\newline
\newline
「那時,耶穌出去上山禱告,整夜禱告神；到了天亮,叫他的門徒來,就從他們中間挑選十二個人,稱他們為使徒.」
\newline
\newline
在別的福音書中雖也有提及揀選門徒,但路加卻特別記載,耶穌在山上整夜禱告後才選立門徒.那麼,禱告就是主選立使徒的開始,禱告是主建立教會的開始.這是主揀選同工的原則.當主整夜禱告後,才揀選十二門徒,有人問:主為何會選猶大呢?猶大是經主禱告後才揀選的,那麼豈不是選錯了嗎?我們會收錯了教友,但主卻不會錯的.祂經整夜禱告後也選上了猶大,那是甚麼意思?對猶大而言,主給他特別的機會和恩典,神願顯明祂的忿怒在剛硬的人身上,但也向這些人多給恩慈.所以對剛硬如猶大一般的人,主顯出更多的恩慈,給他更多的機會悔改!可惜猶大失去這機會.對主耶穌而言,經整夜禱告而選猶大,這猶大是將來賣祂的,那就是說,當選猶大時,祂已準備接受十架苦難,這證明主順服神的旨意,對這十字架之使命毫不拒絕.故祂選猶大在三年多工作上,不論祂所講所行的,對猶大而言是一種恩慈可感動他.對主而言,猶大乃是十架之記號,猶大是賣祂之徒,卻又常跟在祂身旁.祂從沒忘記這十架之使命,也從沒動搖過祂的心志；救主的禱告,可以顯明祂對神順服的心志.
\newline
\newline
(四)問人子是誰之前的禱告──(路九18)
\newline
\newline
「耶穌自己禱告的時候,門徒也同祂在那裡,耶穌問他們說:眾人說我是誰?」
\newline
\newline
當主耶穌問門徒人說人子是誰之前,祂曾禱告,這是路加特別記下的.主在世為人,和我們一樣有情感,有喜樂與悲哀和受試探,故在祂想知道別人如何批評祂之前,祂先禱告.為甚麼要向門徒說:「人說我人子是誰」呢?這話是問門徒的,乃是要他們知道,他們所跟從的耶穌是誰:別人可能多方的說話,但門徒就得弄清楚,他們所跟從的耶穌是怎樣的?祂在問之先,自己先禱告才問；乃是要自己心靈上先有準備.我們喜歡知道別人對我們的批評,別人說好時,我們很歡喜；別人說不好時,我們會傷心了!但主卻有準備,在未問之先禱告,可知主在世上常防守著祂的心,箴言書有說:「我兒,你要保守你心,勝過保守一切.」故當彼得勸主不要上十架時,主立刻斥責他!乃因祂早有準備,所以我們該效法主,在凡事上禱告,因別人說我們好就歡喜,說壞時就憂愁了.
\newline
\newline
(五)登山變相時之禱告──(路九28-29)
\newline
\newline
「說了這話以後,約有八天,耶穌帶著彼得,約翰,雅各上山去禱告,正禱告的時候,祂的面貌就改變了,衣服潔白放光.」
\newline
\newline
路加特別記載:主正禱告的時候,面貌就改變,真誠的禱告,能影響你的生命,影響你的為人,影響你的人生觀,影響你的待人的態度,至終更會影響你的面貌.有一個佈道家說:「一個人的面貌,四十以後就由自己負責.」意思即是說四十歲前是父母生的,但在幾十年後,你這人是誠實,是狡猾,是兇惡或溫和,都從面貌上表露無遺.所以你該為你的面貌負責任,「誠於中,則形於外」,我們裡面有些甚麼,都會表現出來.主禱告在山上變了形像,這種的改變,是一種啟示,是屬靈的教導,讓門徒認識到禱告真正的意義.使生命發生了改變,性情有改變,連面貌也改變,這才是真正的禱告.
\newline
\newline
(六)教導門徒禱告前之禱告──(路十一1)
\newline
\newline
「耶穌在一個地方禱告,禱告完了,有個門徒對他說:求主教導我們禱告,像約翰教導他的門徒一樣.」
\newline
\newline
有門徒到耶穌面前請耶穌教他們禱告,下文跟著記載的就是主禱文.為甚麼門徒要祂教禱告?乃因主自己是個禱告的人,祂禱告完了,門徒就請教祂.如果沒有人來問你如何禱告,那你不要怪別人而該怪自己.如果你是醫生,當然會有人找你醫病,但因你不是醫生,那怎會有人找你醫病呢!為甚麼許多屬靈的問題和屬靈的難處,別人不來找你呢?請你不要怪人,這是一種記號,是一種聲音,是告訴你,因你沒有可供給別人所需的；你不會解決問題,所以別人不將問題帶來求問你.你不能醫病,所以病人是不會來求診的.
\newline
\newline
(七)為彼得跌倒的代禱──(路廿二31-32)
\newline
\newline
主又說:「西門,西門,撒旦要得著你們,好篩你們,像篩麥子一樣.但我已經為你祈求,叫你不至於失去信心.你回頭以後,要堅固你的弟兄.」
\newline
\newline
「我已經為你祈求」,只有路加如此記載.在祂還未公開責備及指出彼得軟弱之先,祂已為他祈求,主早已知道彼得的危險,將會跌倒,故早已為他禱告了.祂的禱告是有功效的,因彼得雖跌倒,但會回頭的,所以祂警告彼得.如果主讓我們看見別人,有軟弱或危險,見到教會有漏洞,你作何感想呢?是給你批評的機會嗎?是表現你的眼光比別人清楚嗎?不,你該為所見到的禱告.別人還未知道的而你卻見到了,那你就該為挽回那人而禱告.沒有人能知道彼得三次不認主,但主早已知道,所以主先為他禱告,然後對他說:「我已為你祈求,叫你不至於跌倒,你回頭之後要堅固你的弟兄.」我們要感謝主!祂不但自己常禱告,也為我們禱告.當祂活在世上作人子時,祂是個忠心的代禱者.現在祂死而復活,坐在父神的右邊；是我們在高天上的祭司,我們更相信祂要為我們禱告.
\newline
\newline
(八)十字架上的禱告──(路廿三34)
\newline
\newline
「當下耶穌說:父啊!赦免他們,因為他們所作的,他們不曉得.」
\newline
\newline
這是路加獨有的記載.當主在被釘十架時的第一句話.就是為那些釘祂的人,而禱告的話.當你的生命在最危險時,要說的第一句話是甚麼呢?當然是把最要緊的話說出來.主被掛在十架上受苦時,卻說:「父啊赦免他們,因為他們所作的他們不曉得.」祂來世並非滅命,直至祂被釘十架上時,仍作救命的工夫,起碼因祂的禱告救了一個強盜.所以祂誠然是至死忠心,祂不是曾教訓過門徒嗎?彼得曾問祂說:「弟兄得罪我,我當饒恕他幾次呢?七次可以嗎?」主答:「不是七次,乃是七十個七次.」因為彼得並沒有饒恕弟兄,試想沒有一怎會有二,七是從一算上去的,如果弟兄得罪他第一次時已得著饒恕了,那麼,永遠是第一次,所以主說七十個七次.現在主自己上十字架,祂禱告求父赦免他們,由此足以證實主給門徒的饒恕教訓.感謝主,祂實在是喜歡饒恕人的主,當我們到祂面前要禱告的時候,是否真有主那樣饒恕人的心呢?不要忘記聖經的教訓,饒恕人要好像主一般.我們要效法主在世上時,凡事禱告,祂受浸時禱告,被人擁護時禱告,選立使徒前禱告,教人禱告前要禱告,祂要知道別人如何批評祂之前要禱告.祂的禱告,影響了祂的生命,改變祂面貌.預先為失敗的弟兄禱告,為敵人禱告,存饒恕人的心禱告,願主賜下禱告的靈給我們吧!
\newline
\newline


\newpage

\section{第44屆~港九培靈會~陳終道~第四講　基督的愛}
\label{sec:1358}
\textbf{第四講　基督的愛}
\newline
\newline
連結: \href{https://www.hkbibleconference.org/session-message/view/1358}{\texttt{https://www.hkbibleconference.org/session-message/view/1358}}
\newline
\newline
\hyperref[sec:1357]{< < < PREV SERMON < < <}
~
\hyperref[sec:toc]{[返目錄]}
~
\hyperref[sec:1359]{> > > NEXT SERMON > > >}
\newline
\newline
「神愛世人,甚至將祂的獨生子賜給他們,叫一切信祂的,不至滅亡,反得永生.」(約三16)
\newline
\newline
約翰福音將神的愛,表明得最清楚.這節經文就是將神的愛,向我們顯明和保證.我們本不曉得何為愛,是神把祂的獨生子賜給我們,為我們作挽回祭,我們才明白.主來世將神的愛表明清楚,在約翰福音中,可見到基督的愛有三方面:
\newline
\newline
(一)不計算的愛──保羅在林前曾說:「愛是不計算人的惡」,其實更不惜任何代價,主的愛就是如此.神將兒子賜給我們,神的愛是不計算任何代價的.
\newline
\newline
(a)在(約一46至57)中,有人介紹耶穌給拿但業,這是腓力給拿但業作的見證.他說:「在舊約先知所說的那一位,我們遇見了,就是拿撒勒人耶穌.」拿但業聽了之後說:「拿撒勒還能出甚麼好的麼?」這是句頂撞耶穌的話,但非常奇怪,主並不計較,反而稱讚他說:「這是個真以色列人,他心裡沒有詭詐的.」結果,拿但業跟從了主.主的愛並不計較我們的無知,不過祂要看我們內心的誠實.在約翰福音中,拿但業是唯一得主稱讚的人.這人卻是頂撞主的,為甚麼反而得主稱讚呢?因為主看出拿但業仍有可愛之處,因他心裡沒有詭詐.
\newline
\newline
(b)在(約五1-9)中,記載畢士大池旁的一個癱子,主過去問他說:「你要痊癒麼?」他回答說:「......沒有人把我放在池子裡.......」於是主就醫治他.在主未醫治他之前,是否知道那些猶太人,文士及法利賽人,反對祂安息日治病呢?是的,祂知道,但祂卻不計較敵人給祂的任何攻擊；愛使祂不計較並要付出代價!「愛」使祂忘記了自己,將遇到甚麼困難；「愛」使祂拯救了那躺臥三十八年的癱子.
\newline
\newline
(c)在(約十三章)中,記載主為門徒洗腳.在本章開始就說,主已知自己離世歸父的時候到了,祂愛世間屬自己的人,就愛到底,繼而又說她束起腰來,為門徒洗腳.那時,正是門徒爭論誰為大的時候.他們爭大,不單是主與他們坐席的時候,其實他們在路上時,已爭論誰大了.主怎樣教導他們呢?祂深知要上十字架,門徒卻以為祂要作王了.所以他們就爭著坐在主的左右邊.主並不計較一切,反而束起腰來替他們洗腳,然後教訓他們要謙卑,彼此服事,祂說「我是你們的夫子,尚且洗你們的腳,你們也當彼此洗腳.」留下榜樣,很有耐心地教導他們.
\newline
\newline
(d)在(約二十24至29)中,記載當主復活向門徒顯現時,多馬不在,故主特別為多馬顯現.因多馬不信主復活,他非看見主手上的釘痕,並用指頭探入祂的肋旁,就總不信,但主沒有計較多馬的無知,而向他顯現,並對多馬說:「伸過你的指頭來摸我的手,伸出你的手來,探入我的肋旁,不要疑惑,總要信.」多馬慚愧對主說:「我的主!我的神!」主的愛,真是不計算的愛.
\newline
\newline
在彼得的書信中告訴我們,「愛」能遮掩許多的罪,也許我們不大同意這說法!因我們不懂得何為愛.只有主耶穌,祂來為我們釘十架,我們才曉得何為愛.這十架的愛,開了我們的眼睛；這十架的愛,顯明了愛的真意義；這十架的愛,才能遮掩許多的罪.求主施恩讓我們明白,祂的愛是沒有計較的；但我們到祂面前,是否還有許多的計算呢?如果我們愛一個人,而整天和他計較；那麼,這愛就有問題了.假如你愛一個人,就巴不得使對方,從你得到愈多愈好.為何作母親的,常會餵兒餵得太飽呢?乃因母親怕兒子饑餓了.沒有一個母親會計較兒子吃了她多少的.如果你在一弟兄或姊妹身上,吝嗇計算對方得過你多少好處,那就不是愛了.因為愛是沒有計算的.
\newline
\newline
(二)忠實的愛──(約三1至15)
\newline
\newline
在約翰福音中,我們看到主的愛是忠實的,怎樣才是忠實的愛呢?就是誠實地把實際的需要告訴你.如愛不忠實,根本就不是愛.忠實的愛包括在實際上的給予.如有缺乏,祂不會裝作不知道,祂會實在地思想到你的需要,然後供應給你.忠實的愛,不但包括在物質上；也包括在靈性及品德方面.我們所缺,所不知的,祂會忠實相告.祂的愛不是溺愛,不是包容罪惡或坦護,不是結黨或私人的勾結；真正的愛不論對方是誰,都要以基督的愛好好勸導.在(約三章)中,主對尼哥底母說:「我實實在在的告訴你,你必須重生,人若不重生,就不能進神的國.」尼哥底母是宗教的領袖,是法利賽人,聖經讀了許多,道理也十分明白.如果與現今的人比較,就等於教會領袖或是長老；如果一個受人尊敬的長老,而你對他說:「你必須重生.」那麼他也許會生氣了,這句話並不是容易說出的.但當日主對尼哥母就如此說,因為祂要針對尼哥底母之需要而說,這才是真正的愛.
\newline
\newline
約(四1－26)中,耶穌對那撒瑪利亞婦人說?「你已有五個丈夫,現在有的不是你的.」假如現在真的有這麼一個女人在你面前,你是否敢於這樣對她說呢?而主竟如此對她說.感謝主!這婦人不但不生氣,反而受感動.她對人作見證說:「有一個人把我素來所行的都說出來,這人莫非就是基督嗎?」可見主給她最深的印象就是這句話,因這話剛好是打中她痛的地方,她就受感動,不單自己歸信了主；城裡許多人也因她而信主,這是真實的愛.
\newline
\newline
(約六27).當主用五餅二魚使五千人吃飽以後,主就離開他們渡到那邊去.許多人在那邊等候祂.主見這麼多的人就說:「你們不要為那壞的食物勞力,要為那不能朽壞的食物勞力.」祂明白這許多人來就近祂的目的,不是信祂,乃因五餅二魚得飽.他們聽了心裡有何反應呢!主的話嚴厲,會使人反感!從外表看來,他們是渴慕見主,豈知主卻針對他們心中的意念.主的話確實使他們難受,但卻是他們最大之需要.當我們被指摘時,該留心那指摘我們的人,他們必有極大的愛心的；因為他指摘我們時,必需冒得罪我們之危險,並且可能因此而結成了仇怨.許多人喜歡別人的稱讚,所以說好話是很容易的；但說責備的話,必定要再三思量,免得對方反面.如忠誠的勸勉,是出於愛,有人願意對你直言,你就該接受,這會叫你得益的.我們年輕時有父母和長輩的教導.年長後,立足在社會上,誰還管教我們呢?父母及長輩們,都不敢再對你說些甚麼,連朋友也是如此,到底這是福還是禍呢?那就要憑你自知之明了,有人一生錯到底,可惜!可惜!倘若有人願意在主裡,因基督的愛勸告你,或指出你的錯處,就算你不接受對方的勸告,也該接受他這份愛心,因為這是忠實的愛.內裡最寶貴的,是忠誠的勸告,人是較難看出自己的短處；故神常藉主內弟兄姊妹,及主僕之口勸告我們.我們應在主前領受這份愛心.這忠實的愛,與故意的挑剔是不同的.願我們明察!
\newline
\newline
在(約八1－11)中,記載一個正行淫的婦人,被人帶到主面前,他們對主說:這婦人是正行淫時被拿的.摩西在律法上吩咐我們用石頭打死她.主並不回答,只在地上畫字.這些人是存著恨惡罪惡的心嗎?不是的,他們是要婦人從此回頭嗎?不是的,他們根本沒有憐愛的心.他們有罪,還敢定人的罪,他們在主面前一個一個羞愧地走了.主對婦人說:沒有人定你的罪嗎?我也不定你的罪,從此不要再犯罪.主有真誠的愛心,因主並非害她,而是要把她從罪裡拯救出來.摩西對以色列人有很大的愛心,但也曾十分嚴厲地責備過他們.有人問在出埃及記三十二章中,摩西叫利未人殺以色列人；那時被殺的有三千,這是因以色列人造金牛犢之故.摩西沒有愛心嗎?不,(在出三十二31節)說,摩西上山向神禱告說,如果可行,他情願將他的名字從冊上塗抹.雖然摩西不可能用自己的靈魂贖以色列人的罪,但他有這樣的愛心,所以配得嚴厲責備以色列人；他責備,是因他深愛以色列人之故.有人以為不應該責備人；但也有人以為保羅和使徒何嘗不勸誡,責備人嗎!是否應該去責備人,視乎你責備時,是否出於愛心.如果你有像摩西一般深切愛同胞的心,那是可以的；因為你在責備中一定充滿了愛.出於恨和找藉口控告人,是不該的.求主施恩,讓我們能接受別人忠實的愛；但也要用忠實的愛,去愛別人.我們能體會到主何等地愛我們；使我們常受挫折,被人責備或輕看；這一切都是出於祂對我們忠實的愛.藉此,好使我們醒悟過來,從各樣錯誤中,回到祂的面前.
\newline
\newline
(三)最美的愛──愛會使你把最好的給別人.神將主耶穌給我們,那是將最好的,所以神也有權利向我們要求,把最好的獻上.在約翰福音,就可以看到這份愛.
\newline
\newline
(a)約九:1至3節:記載主醫治一個生來的瞎子的故事.
\newline
\newline
許多人會發出這樣的問題,神如果是慈愛的,為何叫他生下來就瞎眼呢?主所給的答案並非在身體上殘疾,而是如何去應付所遭遇的事.那生來的瞎子,因他瞎了,主才能在他身上顯出更大榮耀.祂用泥抹在他眼睛上,吩咐他去洗,一洗就能看見了,從此他為主作了美好的見證.著名的詩人芬妮克羅司比,她是一個瞎眼的詩人,因幼年患眼疾,被醫生錯用了藥而瞎了眼.我想先天瞎眼和後天瞎眼者相比,後者會更為痛苦.一個本來不瞎眼的,而被醫生弄瞎了眼,那是痛苦和不幸的事!芬妮本可以怨天尤人,抱怨上帝,咒罵醫生；但她並不如此,她安心領受神給她的遭遇.她雖瞎了眼,卻寫了許多詩歌.她如何能成為一個成功的女詩人呢?正因她是瞎了眼睛,才有此成果.神會揀選最好的給我們,問題是我們是否願意,是否領受而已.
\newline
\newline
約十一章中記拉撒路生病死了,在他未死之前,馬大和馬利亞打發人去對主說:「主呀!你所愛的人病了.」她們充滿了信心差人告訴耶穌,但卻眼巴巴地看著弟弟病重而死.而且死了後,耶穌還未來.她倆實在受到極大的打擊,直到弟弟死後四天,主耶穌才來.主雖擔擱了時間,到底還是為她倆的好處.因主將最好的給她們.主若早到,只說一句話便使拉撒路病癒,當然是可以的；這樣就無法看得見拉撒路從墳墓裡,復活的大神蹟了.這更高福份,只有神揀選有屬靈程度的人配得；因主曉得揀選最好的給我們.
\newline
\newline
舊約的約瑟被主母陷害而下獄,是因他在神面前持守他的純潔.他為主的緣故不願犯罪,一個青年能拒絕男女之試探,是多麼難得.一個好人約瑟,應得神的保祐才是,但卻反被捕下獄!這到底是神的愛?還是神不看顧?神實在是選擇最好的給他,試看!約瑟若不下獄,只不過在波提乏的家作一世的管家而已；若酒政沒有忘記約瑟而得釋放,也不過是仍再作管家而已.但神遲延聽他的禱告,到了時候才釋放,那就使他一躍而成全埃及的宰相了,顯見神實在將最好的給我們.
\newline
\newline
我們常緊緊抓著自己所要的,其實神才知道把甚麼是最好的給我們,祂的愛是沒有計算的.她既把獨生子白白賜給我們,豈不把萬物和祂一同的給我們.求主賜恩,讓我們不單是一個被主愛的人,也能深切地愛主；我們領受那最好的,也將最好的獻給祂!
\newline
\newline


\newpage

\section{第44屆~港九培靈會~陳終道~第五講　基督所責備及稱許的人}
\label{sec:1359}
\textbf{第五講　基督所責備及稱許的人}
\newline
\newline
連結: \href{https://www.hkbibleconference.org/session-message/view/1359}{\texttt{https://www.hkbibleconference.org/session-message/view/1359}}
\newline
\newline
\hyperref[sec:1358]{< < < PREV SERMON < < <}
~
\hyperref[sec:toc]{[返目錄]}
~
\hyperref[sec:1360]{> > > NEXT SERMON > > >}
\newline
\newline
我們是基督徒,相信沒有一個人,想成為被主責備的人,而是想得主的稱許.那麼,怎麼才可得主的稱許,而又不受責備呢?最簡單的方法,是看四福音內,主責備何種人,並稱許些怎樣的人；就可知所儆惕,和知所效法了.
\newline
\newline
「凡說話干犯人子的,還可得赦免,惟獨說話干犯聖靈的,今世來世總不得赦免.你們或以為樹好,果子也好；樹壞,果子也壞；因為看果子,就可以知道樹.毒蛇的種類,你們既是惡人,怎能說出好話來呢?因為心裡所充滿的,口裡就說出來.」(太十二32-34)
\newline
\newline
(一)毒蛇的種類──「你們既是惡人,怎能說出好話來呢!」這是句很厲害的責備說話.主在當時竟用這樣的話,責備法利賽人,究竟是甚麼緣故呢?這是因為他們犯了褻瀆聖靈的罪.許多信徒都擔心自己,會犯此褻瀆的罪!但主責備犯這罪的人,乃是假冒為善的法利賽人,說他們是毒蛇的種類.因為他們懂得道理,見過神蹟,聽過主講道,明白聖經.當主降生時,幾個博士去尋找耶穌,這些人都曉得先知的預言及基督降生,是在伯利恆.可惜他們卻沒有去尋找主,相信主.他們是明知而不信,故犯了褻瀆聖靈的罪.他們並非天國之子,他們是毒蛇的種類.主說:「你們看樹就可以知道他的果子了,樹好則果子也好；樹壞,則果子也壞.」他們為何要褻瀆聖靈呢?乃因這樹壞,果子也壞.主明說:「你們心裡所充滿的,口裡就說出來.」法利賽人所犯褻瀆的罪,非口頭上偶然之過失；而是裡面充滿了,口就說出來,是從毒蛇的生命中,流露出來的.故凡是褻瀆聖靈的人,儘管聽了道理,或者相當明白道理；但並沒有基督之生命在他裡面,因他們是毒蛇之種類.我們知道,主所說犯褻瀆聖靈和犯褻瀆聖靈的話是不同的.參看別的福音可知道.馬可說:「人所講一切的話可得赦免,唯獨褻瀆聖靈,永不得赦免.」一切褻瀆的話是否就包括褻瀆神的話呢?是的；因「瀆」字之詞意,是用在神方面,毀謗神就是「瀆」,但主把褻瀆聖靈,和褻瀆之話分開；因為褻瀆聖靈不單是在話語,非偶然之錯失；更是敵擋聖靈之態度,故意輕視神的作為之態度.這種人早就常消滅聖靈之感動,常抗拒聖靈之工作,他們才敢褻瀆聖靈.在此段上文說及主趕鬼,醫治一個被鬼附著又瞎,又啞的人；他們就說主靠鬼王別西卜趕鬼,其實他們早知道主,是靠神的能力趕鬼的.他們是毒蛇之種類,心裡所充滿的,口就說出來.主所責備的,乃是那些懂得許多道理,看過主的作為,但全不願受感動,而誠心去悔改的人.他們在教會裡已有相當時間,慕道已久,但總不決心相信；甚至有些人用屬靈的知識,以圖謀屬肉身的好處.這些人就成為聖經所說的假師傅,也是主所責備的,我們在世上會犯錯失,當聖靈感動我們心的時候,就切不可消滅聖靈的感動.雖然聖靈是永存在我們裡面,我們也實在是祂的兒女；主賜聖靈作我們的印記,直至見祂面時.但主既責備這種褻瀆聖靈的人,也就是那些常消滅聖靈感動的人,我們不該消滅聖靈的感動了.
\newline
\newline
(二)妄用屬靈權柄──主與門徒經撒瑪利亞上耶路撒冷去,撒瑪利亞人不接待她；因主面向耶路撒冷.(參路九54-56)這句話有雙重意思:一方面主是往耶路撒冷去,因撒瑪利亞人本與猶太人沒往來.主面向耶路撒冷,所以撒瑪利亞人不接待她.另一意思是:主面向耶路撒冷,乃是要去預備迎接,等待著祂的十字架；祂專心走此十架道路時,世人不歡迎祂.門徒見撒瑪利亞人不接待主,就生氣地說:主呀!你要我們吩咐火從天降下來消滅他們,像以利亞所作的麼!主就責備他們:你們的心如何?你們不知道,我來並非滅人的性命,乃是救人的性命.從此可知兩門徒存心如何,他們是要炫耀自己本領,你們不接待嗎,那麼就讓我們求主顯些本領給你們看看,讓你們知道利害.門徒是妄用了屬靈權柄,他們不知道屬靈之權柄,乃是為對付魔鬼；而非炫耀及表現自己比別人強.這樣的存心,實在是一種愛世界的虛榮心,是愛世界的另一種方式.當然,主是可以吩咐火從天降下,消滅他們；但主從不亂用屬靈權柄.祂不是曾用五餅二魚使五千人得飽嗎?但祂卻不願為自己,將石頭變成食物而充飢；祂曾為別人行過許多神蹟,但祂從不行一個神蹟使自己得飽.因祂不胡亂使用神所賜的權柄.如果神給我們屬靈的恩賜,我們該如何去使用?是用自己的名譽,地位,謀肉身的好處嗎?又如果神將我們放在祂家裡,給我們在聖工上有職份,給我們在教會上有地位,那我們如何去運用呢?是利用地位,權柄以炫耀自己嗎?若如此,則我們與這兩個門徒犯同一的毛病,主說:「你們的心如何,你們不知道.」
\newline
\newline
有一次有人帶小孩到主面前,門徒要攔阻.主責備門徒說:「讓小孩子到我這裡來,不要禁止他們.」為甚麼門徒要攔阻,不許這些人帶孩子來呢?因為他們覺得主是大人物,是奮興家,和小孩子講話是不相稱的；甚至連他們這班門徒也會丟臉!若是對大人物,對名人講道,我們作門徒的,就增加幾分光彩了.主要糾正他們的心.主在世上偉大之處,並非是祂能對大人物講道,而是祂能接待小孩子.我們對事常有錯誤之觀念,何為高貴,何為偉大,何事值得人尊敬,我們都不曉得分辨；而很易受到世界的潮流與及人的看法所影響.但主要責備我們,不應有這種的眼光,更不該利用屬靈的權柄,地位,恩賜以及才幹,為自己謀好處；而又炫耀自己的聲望與地位.所以主責備說:「你們的心如何,你們不知道.」
\newline
\newline
(三)不儆醒禱告──「你們不能儆醒片時麼?」(可十四37-38)這是主對門徒的責備,主在客西馬尼園禱告.門徒很疲乏,時已夜深,他們因困倦而在禱告時睡著了!主這樣責備他們合理嗎?困倦時睡了,本是平常的事,為何主要責備他們呢?請注意!主在這時候,何嘗不困倦呢?可能比他們更困倦,主尚且能儆醒,門徒卻不能,究竟何故?主也是人,和我們一樣,有血肉之體,也需要睡覺,為甚麼祂要儆醒禱告呢?因為主看見前面有最大的爭戰,有最大的危機,有一大黑暗勢力；前面擺有十字架,祂要面對這十字架,迎接這戰爭.祂看見這場戰爭關係及全人類,於是祂所有的困倦與睡意都消失了.可惜門徒卻沒有看見,他們怎能不睡呢!這就是客西馬尼園中,主耶穌與門徒分別之處了.主責備門徒不儆醒!按「儆醒」之意,是眼要明亮些,要看清現今是否應睡覺的時候?更叫我們曉得,我們所面對的戰爭,是何等厲害；我們對這世代所負之責任,又是何等的重大!現今是怎麼時代?是我們睡覺的時候嗎?先知以利沙責備他的僕人基哈西說:「這豈是受銀子,衣裳,置橄欖園,葡萄園,牛羊,僕婢的時候呢?」(王下五26)換言之,這豈是買房子享清福的時候呢?以利沙看到當時自己所負的使命是何等重大!基哈西只貪愛乃縵的財物,所以招咒詛,反而得了乃縵的大麻瘋.這是以利沙與基哈西所不同之處.
\newline
\newline
(四)百夫長的信心──百夫長為他的僕人而去求主,主正想去他家的時候,他就對主說:「主呀!不要勞動,你到我舍下,我不敢當；我也以為不配去見你,只要你說一句話,我的僕人就必好了.因為我在人的權下,也有兵在我以下,對這個說去,他就去,對那個說來,他就來.」(路七8)他相信只要主說一句話,他的僕人就會好的.主稱讚百夫長的信心.百夫長的信心有何特點呢?因他的信心能使他對屬靈的事情,有更深的領悟；能影響他待人接物的態度,而且影響他對神的看法.真的信心,誠然會影響到我們的人生觀!百夫長之所以有屬靈的領悟力,是因他的信心影響到他的生活.上文說,百夫長為僕人代禱,求主醫治.以他高貴的身份願降卑為僕人求,可知他深切愛僕人,對僕人尚且有愛心,從此亦可聯想他對妻子,他的兒女亦必定有愛心了.所以他的愛心影響到他的家庭生活.當時有幾個猶太人的長老,來替百夫長求主說:「你給他行這事,是他配得的,因他愛我們的百姓,給我們建造會堂.」(路七4-5)百夫長必不是以色列人,但他卻管理以色列人；他非大官,但他卻是百姓的父母官.猶太人的長老能為他作見證說:「他是配得的,因他愛我們的百姓,給我們建造會堂.」可見百夫長的信心,影響到他在工作上,職份上的忠心與愛心.所以他對神的事,有領悟力；因他的信心,而影響到他對神的態度.如果我們要得主的稱讚,我們的信心,也必須影響到我們的生活哩!
\newline
\newline
(五)迦南婦人──(參太十五28)關於這迦南的婦人的事,我們可能有些疑問,為何主會用這樣的態度對待她呢?好像故意與她為難,才答應她所求的.主起初不理睬她,後來又說些難聽的話:「不好把兒女的餅給狗吃.」好像是罵她是狗,然後才答應她所求.這樣看來,是否主輕看她呢?乃要試驗她的信心.這婦人之所以得到主的稱讚,是因為她的信心經得起試驗；如果不經這試驗,就不能將信心的特色,和可貴之處表現出來.例如神試驗亞伯拉罕,是否神要試驗了才清楚知道一切?祂早已知道一切了,但亞伯拉罕經此試驗,他信心之偉大就完全顯露出來,證明他實在是信心之父.(在太十五)先說迦南的婦人在主後面,主不理會她,然後對她說:「不要將兒女這些餅給狗吃.」然後婦人說:「主阿,不錯,但是狗也吃他主人桌子上掉下來的碎渣兒.」於是耶穌就稱讚她.假如將上文略去,而只記婦人來求主,主就立刻稱讚她的信心是大的,那麼這稱讚就不合理了!所以試驗並非為主,而是為這婦人；使她寶貴的信心,能顯露出來.如果主容許我們受試驗,那麼必定是給我們配受,而又是我們受得起的.正如小學三年級的學生,教師不會教他讀中三的課程.所以主若讓我們受試驗,必定先看出我們的信心有特點,而將我們信心的特點,顯露出來而已.
\newline
\newline
(六)一件美事──(太廿六10-13)我們讀此經文時,覺得很奇怪,為何馬利亞用香膏膏主的事,可以值得普天下傳福音的人,都要提及呢?主給這伯大尼的馬利亞如此高的稱讚,就是一件美事,普天以下,無論在甚麼地方傳福音,也要述說這女人所行的,作為紀念.
\newline
\newline
馬利亞膏主,和主來世之主要使命,有密切關係.主來世,是為人死,祂曾三番四次地告訴門徒,祂要上耶路撒冷,在文士,祭司長之手下受辱,甚至被殺,三日後復活,但門徒卻聽而不聞,又不明白,這對主來說,乃是一種痛苦.正如你有大的心志,卻沒有人明白,甚至你最好的朋友也不明白一般.如今主要為人死和復活,門徒好像完全不知道,也不明白；但馬利亞卻用香膏來膏主,這是表明她實在明白了.所以,主為祂作見證說:「她是為我安葬的事,把香膏預先澆在我身上.」(可十四8)意思是說:以後就算用香膏,要膏我也沒有機會了.馬利亞知道主要到何處去,更知祂要作的事；所以她不但用香膏膏主,更要讓主知道,在這裡,還有一個明白主心意的人,使主的心得到安慰!這也教導我們愛主是怎樣的一回事；不單是將最好的獻給祂,更是要明白主的心意如何.所以主稱讚她並吩咐在普天下,無論在甚麼地方傳這福音,也要述說這女所作的,以為紀念.
\newline
\newline
一斤名貴的真拿達香膏膏主,在馬利亞方面而言,因她知道主快要死而把握這最後機會,所以不惜付上任何代價.但在門徒而言,他們沒有預感,知道主的死,所以對馬利亞反感!馬利亞是個冷靜的人,她早已有心理上的準備,她並非預備人稱讚；而是等著別人的批評和議論.所以當猶大和其他門徒責備她時,她能默然無聲,靜候主的判斷.
\newline
\newline
假如你像馬利亞,把最美的香膏獻上,將最真誠的愛獻上,那麼你就該有馬利亞一般的心理準備；不是預備要得人的稱讚,乃是預備人的攻擊,為主受辱!你該像馬利亞一樣的安靜,等主來為你分訴!我們該知道,馬利亞的香膏何能顯出價值來呢?乃是在主還未死之前；如果在主死之後才獻上,那是太遲了!若等祂復活才帶到墳墓去膏祂,那就更遲了!如五餅二魚的神蹟,若在眾人都吃飽了,才拿出來,那就全無價值.所以,時間是要緊的,馬利亞能把握機會獻上最寶貴的.
\newline
\newline


\newpage

\section{第44屆~港九培靈會~陳終道~第六講　基督所發的問題}
\label{sec:1360}
\textbf{第六講　基督所發的問題}
\newline
\newline
連結: \href{https://www.hkbibleconference.org/session-message/view/1360}{\texttt{https://www.hkbibleconference.org/session-message/view/1360}}
\newline
\newline
\hyperref[sec:1359]{< < < PREV SERMON < < <}
~
\hyperref[sec:toc]{[返目錄]}
~
\hyperref[sec:1361]{> > > NEXT SERMON > > >}
\newline
\newline
(太十六14-15)中主問門徒說:「你們說我是誰?」
\newline
\newline
法利賽人和撒都該人,用許多問題來為難主,問應否納稅給該撒,又問人死後復活的情形,又問律法中那一條最大,主都給他們滿意的回答.然後反問他們一問題說:「你們知道基督,是大衛的子孫,但當大衛被聖靈感動時；他說:主對我主說,你坐在我的右邊,等我使你的仇敵作你的腳凳.基督既是大衛的後裔,按理他應該比大衛小,又怎會稱為主?」他們不回答.雖知主來世道成肉身,按肉身說祂是大衛的後裔,但按神性說,祂與神同等,祂是大衛的主.照樣,我們也來向主發問,而忘卻了回答主的問題.我們向主發的問題是複雜的:為何我會有這樣的遭遇,為何會得這病症,為何會娶得這麼一個壞媳婦,又為何會要侍候這麼一個婆婆.......但卻忘了主豈非用寶血買贖我們嗎?祂並不是叫我們為祂要受苦嗎?為何我們的心仍向世界?祂是為我們而死,為何我們仍犯罪呢?主今日向我們發問題,我們該如何回答?從聖經中看看主所發的問題,我們該如何回答呢?
\newline
\newline
(一)誰能使壽數多加一刻呢?(太六27)
\newline
\newline
這是主對我們發的問題,主要教我們有信心的過一個靠主的生活.主問說:「誰能使壽數多加一刻呢?」其實思慮只會使你壽數多減一刻,所以不必為衣食,生活等等來憂慮；只要先求祂的國和祂的義,這一切都加給你了.
\newline
\newline
我們常以為主給的恩典不夠用,其實是夠用的.主曾應許過,我的恩典夠你用,一天的難處一天當就夠了.你若將多天來的難處一天當,甚至多年的難處,都放在一天擔,那當然是不夠用,我們常忽略了主的儆告,我們只憑自己的思慮,不但壽數沒加增,反而減少；而且使我們人生充滿憂慮,不能榮耀主!
\newline
\newline
(二)人子是誰?(十六13-15)
\newline
\newline
主問門徒人說我人子是誰?門徒答:有人說是以利亞,有人說是耶利米,又有人說是施洗約翰復活了.但主問:「你們說我是誰?」主要清楚知道他們對自己的認識到何種程度.別人的議論並不重要,但跟從主的門徒,就應當對主有深切的認識.今日是科學時代,喜歡用客觀的態度來研究一件事,也留心別人的看法如何.不過主卻要我們對祂要有主觀的認識.聖經已將主介紹給我們,我們雖不像門徒那樣,從肉身上認識主；但可憑神的話去認識,最要緊的是在神的話語中,用心領受.我們不能靠參考書,也不能只根據別人所領受的,必須從聖經真理中清楚認識主.
\newline
\newline
「你們說人子是誰?」雖有人說是以利亞,是耶利米,或是先知裡的一位；但他們都不是耶穌,只不過在某一點上像耶穌而已.其實主耶穌乃集中所有歷史上,信心偉人之特點,祂是一位完全的救主.請問我們今日是否清楚認識,我們所跟從的主耶穌基督是誰呢?
\newline
\newline
(三)賺得全世界有甚麼益處?(可八36-37)
\newline
\newline
我們傳福音時常用此節經文,「生命」原文可譯作「靈魂」.上文是勸門徒要背十架跟從祂,然後才告訴門徒說:「人若賺得全世界,賠上自己的生命有甚麼益處?人還能拿甚麼來換生命呢?」主如此問門徒是何意?乃要勉勵門徒,使他們知道背十架跟從主,是值得的.人生中最有價值的工作,是拯救靈魂；無論作任何事,都不能賺全世界.如今主卻說賺一個人的靈魂,比賺全世界更有價值；我們跟從主,乃是要得靈魂,花時間及精神也是值得的,許多基督徒為自己的工作及為學業花精神；但為神的話語──聖經──而花時間,就以為不值得.所以主要喚醒我們,改正我們的人生觀,活著並非為肚腹,而是為更有價值的人生觀.因世界一切都要過去,唯有為靈魂付代價,才是有永遠的價值的.
\newline
\newline
(四)要我為你作甚麼?(可十36,51)
\newline
\newline
這兩處經文同是問一個問題,起初是問雅各和約翰,後來是問瞎子巴底買.雅各和約翰是門徒中最接近主的,三人中的兩個,他們見主這樣問,就求賜將來在主的國度裡,坐在主的左右.主說:你們不知你們所求的是甚麼?他們跟主已久,聽道又多,又為十二使徒之一；對主認識就該比任何人多,可惜所回答的卻不如瞎子!主以同一的問題問瞎子巴底買,他卻答說:「拉波尼,我要能看見.」他能按所需要的回答,他所要的,不是金錢,不是地位,不是其他一切.......他的回答正合乎主的心意,若能看見,那麼甚麼問題都可解決了.為甚麼瞎子比那兩們徒答得更好呢?因門徒仍有虛榮心,追求地位,抓緊機會,要坐主國度的左右.他們為了虛榮心,而錯過了大好的機會!巴不得我們被主問時,能曉得如何回答.我們需要屬靈的眼睛開得明亮,能看見現今時代的需要是甚麼?
\newline
\newline
(五)為何只看見別人眼中有刺?(路六41-42)
\newline
\newline
為何只看見別人眼中有刺,而不見自己眼中有樑木呢?為何你們的眼睛在別人身上如此精明,而對自己又如此昏暗呢?這是主對我們發出的問題.我們喜愛挑剔別人的不是,而作為對別人定罪的根據.若是如此,我們每天都可有挑剔別人的機會；但如果要找別人的長處,當然也可以找到的.主問為甚麼你只看見弟兄眼中有刺,而看不見自己眼中的樑木呢?因你並沒將自己放在主的光中；讓祂光照你,因你已落在黑暗中,所以看不見自己,何等可惜!
\newline
\newline
(六)那九個在那裡?(路十七17)
\newline
\newline
「潔淨了的不是十個嗎?那九個在那裡?」這是主所發的問題.也許他們有理由解釋,主不在固定的地方傳道,得醫治後,主已往別處了.希奇嗎?那一個卻能找到主!許多人遭遇疾病時,有困難時,就迫切求主,還請戚友,牧師代禱；但當病癒後,困難解決後,就不回頭謝恩了!請他到教會嗎?工作太忙,生活困難,沒有時間!砌詞推諉,無怪主問道:「那九個在那裡?」我們是個忘恩負義的人嗎?我們是那九個呢?還是那一個呢?願我們深自反省!
\newline
\newline
(七)你們有吃的沒有?(約廿一5)
\newline
\newline
主所發的問題,充滿了慈愛,當門徒整夜打魚而打不著甚麼時,主就向他們顯現,並問他們說:「小子,你們有吃的沒有?」這句話對已撇下世上職業而專心以祈禱,傳道為事的人,是特別有意義的.門徒本已撇下漁船,漁網,為甚麼現在又重操故業?是從前錯了?還是現在錯了呢?當然一般基督徒是帶著職業事奉主,站在基督徒的崗位上,盡他該盡的本份；但對一個獻身于於傳道的人,而再去做生意,則與奉獻的原意恰得相反了!彼得和其餘的門徒,當然有重操故業的自由；不過,終夜勞力卻一無所獲,這是他們收回奉獻的結果.
\newline
\newline
「小子,你們有吃的沒有?」從這句話可證明,曾與他們在一起的主,過去如何照顧他們生活上的需要；現在復活了,仍然依舊地照顧他們的需要.吃的,用的,祂都沒有忘記.當他們缺乏時,主早已在岸上為他們預備了炭火,有魚,有餅.主像過去活著時,如何看顧他們,並無分別.永不改變的主,祂永不忘記那些為祂勞苦傳道的人.
\newline
\newline
(八)你愛我比這些更深嗎?(約二十一15-17)
\newline
\newline
這是耶穌向彼得所發的問題,他不但問彼得,也是向我們發問:「你愛我,比這些更深嗎?」到底主所說的「這些」是指甚麼?從上文看,彼得聽主的指示,把網撒在船右邊,而獲得一百五十三條魚.對一個漁夫來說,的確是筆財富.對那些「沒得吃」的門徒來說,是寶貴的糧食.對那些整夜勞力,而一無所得的門徒來說,是一個奇妙的神蹟.你愛主比這屬世的財富更深麼?比屬靈奇妙的經歷更深麼?主問彼得實在愛打魚更多,還是愛祂更多?彼得毫不猶疑地答:「主啊,你知道我愛你.」但主不因此而感到滿足,祂再三發問,因為彼得在行動上未能表現更深的愛主!他不但沒有去牧養主的小羊,反而合伙重操故業!主曾呼召他撇下所有,作得人漁夫.現在他再拿起所撇下的,而說愛主,豈不食言!所以主要一而再地考問並對他說:「你餵養我的羊.」如果他能夠真心愛主比這些更深,就必須照著主所吩咐的去行.
\newline
\newline
也許有多少人也像彼得那樣說:「主啊!你知道我愛你.」但可惜行動卻與主的旨意相反!只會追求世界的虛榮,卻未顧念神家的需要,缺乏愛人靈魂的深切!未能顧念軟弱的弟兄.......整天所追求的是金錢,享受,逸樂!口裡說愛主,心卻與主遠離!我們當醒悟!回復起初愛主的心!
\newline
\newline
對於主所發的幾個問題,我們該怎麼回答呢?是否知道所信的是誰?是否天天信靠主的恩典,不為自己憂慮呢?是否已認識到人生最高價值與工作呢?眼睛是否已看見?是否像那九個忘恩負義者?是否相信永恆不變的主,供給我們一切所需,也體貼我們的軟弱,願主憐憫!教我們明白祂的心意,更了解祂的問題,回答得合適,得主喜悅.阿們!
\newline
\newline


\newpage

\section{第44屆~港九培靈會~陳終道~第七講　基督與瑪門}
\label{sec:1361}
\textbf{第七講　基督與瑪門}
\newline
\newline
連結: \href{https://www.hkbibleconference.org/session-message/view/1361}{\texttt{https://www.hkbibleconference.org/session-message/view/1361}}
\newline
\newline
\hyperref[sec:1360]{< < < PREV SERMON < < <}
~
\hyperref[sec:toc]{[返目錄]}
~
\hyperref[sec:1362]{> > > NEXT SERMON > > >}
\newline
\newline
(太六24及路十六13)上文是同一樣的經文,時間卻是不同；主在不同時間,而說出同一的信息.祂儆告信徒說一個僕人,不能服事兩個主人；不是愛這個,惡那個,就是重這個,輕那個；你們不能事奉神,又事奉瑪門.主為何一而再地論及這信息?可見信徒對基督與瑪門有重要的關係.主已明白宣告,在祂與金錢之間,是沒有並立的地位.如果要事奉基督,就不能事奉瑪門；若要事奉瑪門,就不能事奉基督.在基督與瑪門中間,該有一個選擇,不然就會重這個,輕那個.主要求我們的事奉,不是部份的,而是全部的；要整個的心,和全部的愛.在世上能與主爭奪,搶去我們的心的,就是錢財的力量(即瑪門的勢力).今天的世界,注重物質,金錢能換取物質；人愈重物質的享受,則愈重金錢.在世人眼中,金錢乃最有價值的東西,主亦不否認金錢的用處；不過,祂只說,不能事奉金錢而已.路加福音是一本多談及錢的書,以此書作根據,顯明人有兩方面的需要:一方面是錢,因我們活在世上,錢能使我們生活下去,這是關乎肉身方面.另一方面是需要禱告,禱告是維持我們與神的關係,能使我們在世上過著一個屬天的生活.所以路加福音多提及主的禱告,以及有關對錢財的教訓:
\newline
\newline
(一)基督與瑪門對立之事實:
\newline
\newline
(a)(路八14)是撒種的比喻.論四種的心田,第一種是種子落在路旁的比喻,這些人是未得救的.第二種落在淺土石頭裡的,這種人不過暫時相信.只有落在好土裡的,才是真正相信的人；他們能結實,有三十倍,六十倍和一百倍的.但有一種人最可憐的,就是種子落在荊棘裡,是表明神的話,並非沒有在他心中運行,作工或感動；而是受了攔阻,是荊棘的攔阻.荊棘在(路八)中解釋就是今生的思慮,錢財的宴樂擠住了.換言之,今生的憂慮,就是怎樣去賺錢；錢財的宴樂,就是怎樣去用錢.不論是賺錢或用錢,總之,這荊棘是關係錢的問題；忙於賺錢或忙於花錢,而使人不能接受耶穌.所以瑪門與基督是對立的.
\newline
\newline
(b)(路十二13-21)記載,無知的財主為自己的錢財打算,他說,我的出產那麼多,沒有地方收藏,怎麼辦呢?他便將倉房拆了,另蓋更大的；好收藏更多的糧食和財物,然後對靈魂說:「靈魂啊!現在你只管安安逸逸吃喝快樂罷!」但神對這無知的財主說:「無知的人呀!今夜我要你的靈魂,你所積蓄的要歸誰呢?」世人不肯接受主,不肯為自己的靈魂打算；乃因被錢財迷惑,沒看見靈魂的需要.這財主是一個典型的代表,以為有了錢就可滿足,其實人不是為錢而活的.有人以為搶劫的人,是因窮的緣故；果真如此,則窮人就是罪人了,愈窮則罪就愈大；相反的,有錢人就是好人,愈有錢則是最好的人了.這樣的解釋,不過是替有錢的人說話.其實在現今的世代中,真正犯大罪的人,不是窮的人而是有錢的人.因有些罪,倘若沒有錢,不能犯的.人以為金錢可以解決一切,也可解決道德問題,社會問題,以及心靈的空虛；像這無知的財主,忙忙碌碌,為自己建造倉庫收藏出產,但心靈卻空虛,又有甚麼用處呢?倘若神今夜要他的靈魂,那麼所預備的,又歸誰呢?
\newline
\newline
(c)(路十六19-31)記載一個財主的禱告,可以說乃是在全新約中最美麗的禱告.就是財主在陰間裡,哀求他祖亞伯拉罕,打發拉撒路用指頭尖蘸點水,涼涼他的舌頭.這禱告何等懇切謙卑,但卻不蒙神的悅納!因為太遲了!為何要到陰間才禱告呢?是因他有錢,為何輕忽了拉撒路,每天在他身旁所作的見證呢?也是因他有錢!有錢使他輕忽了機會,故瑪門和基督是對立的.
\newline
\newline
(d)(路十八18-30)記載,少年的官歡歡喜喜地來求主說:我該作甚麼事,才可以承受永生?這少年從小就遵守誡命,是否遵守誡命就可得永生呢?假若是,這少年應得永生了,何須再到主面前求?可知單遵守誡命,並不能得永生的.主並沒和他辯論這問題,只叫他知道纏累著他的是甚麼?攔阻他到神前不能得永生的是甚麼?這少年官是貪愛錢財的,所以主對他說:你要將一切所有的變賣分給窮人,還要來跟隨我.這少年就憂憂愁愁的去了,因為他的問題是在錢財上.甚麼使他離開主?甚麼使他滅亡?乃是錢財,故瑪門與基督本是對立的.
\newline
\newline
(二)基督對錢財的態度
\newline
\newline
在福音書中很少提及,主與門徒們的經濟來源,只在(路八1－3)中提及.當主受試探時,魔鬼叫祂將石頭變為食物.他說:人活著不是單靠食物,乃是靠神口裡所出的一切話.食物是為生活所需,也是為了錢財.主能抵擋試探,表明祂是靠神的話活著.許多基督徒和傳道人退後的真正原因,可能與食物有關,為了生活上的問題被世界拖去了.
\newline
\newline
在(路九58)記載,有人來跟從主,主說:「狐狸有洞,天空的飛鳥有窩；只是人子沒有枕頭的地方.」保羅說:「祂本來是富足,卻為你們而成了貧窮；叫你們因祂的貧窮可以成為富足.」(林後八9)主在世時,在物質及錢財上,都不會比我們好；有時我們會怨神給我們錢財過少,但你曾否知道,神的兒子來世比你更不如呢!祂的生活雖貧窮,卻沒表現出寒酸的樣子要人可憐!祂雖為我們成了貧窮,卻常責富足的人.祂被請赴筵席時,就教訓那爭坐位大的.在(路六24)祂說:「你們富足的人有禍了!因為你們受過你們的安慰.」難道富足人不可以得安慰嗎?不是的；這意思是說富足的人很易受人之奉承.他們少看見自己的缺少,就很易在說人稱讚奉承中被蒙蔽.主雖貧窮,但祂並非去奉承富人來說話；反而在真理的權威下,真正地將他們的缺點說出來.
\newline
\newline
(太十七24-27)有一次,有人來向主收稅,主問彼得世上的君王收稅,是向兒子呢?還是外人?彼得說:「當然是外人.」主說那麼兒子可免稅了!但為了使收稅人不至絆倒,就叫彼得去打魚,把魚釣上可得一塊錢作稅銀.主既沒有錢,為何不多釣些錢呢?可見主在世如何依靠神,祂的心從不在錢財上打算過,只是專心行神旨意.因為人的心在錢財上,則思想動態也在錢財上打算了.試看猶大見馬利亞用香膏膏主時就可知,猶大說香膏值三十兩銀子,可賙濟窮人；這道理多令人感動,是誰叫猶大說出這話的呢?是瑪門!是錢在他心中.有時你會掛著事奉基督的招牌,但所走的是瑪門的路!說話是根據瑪門而說,臉色也跟瑪門轉變!你是在瑪門的勢力下!主說:人要事奉祂,就不能事奉瑪門.
\newline
\newline
(三)基督怎樣教我們處理錢財──(路十六1至18節)
\newline
\newline
這裡說及一個不義的管家,說人告他浪費主人的錢財.主人要將他辭去,他在未離職交代前；將主人的債戶請來,問欠多少?欠一百簍油的寫五十,欠一百石麥的寫八十.到了離職後就往債戶家裡去.主人讚他作事聰明.這比喻的中心點,是說這不義的管家,有利用機會的聰明.他的不義有兩方面,一是浪費,不應用的他用了；一是肥己,應為主人用的,卻用在自己身上.我們不過是神的管家,神交錢財給你管,你就不可浪費錢財.許多人收五百不夠用,一千,二千也不夠用；不夠用的原因是浪費而非主少給他.肥己就是主給我們的錢財,有一部份是主要我們為祂而用的,但你卻為自己用了,那麼你就是不義的管家.試問那五餅二魚的神蹟,主是否先擘開了一份份才分給人吃呢?一萬人該花多少時間.該知這神蹟是在傳遞時發生的.主祝福了將它擘開遞給門徒,門徒再遞給眾人；這一擘一遞中,自己得了所需的,然後又遞給別人.神蹟在短快的時間中,使這麼多人吃飽,如果其中有一人不傳遞,則神蹟就停在他手裡；所以神給我們的錢,有一部份並非屬於我們的,要我們學習作祂的管家.要利用機會去用錢,這管家在主人辭退他時,手裡仍有權可運用主人的錢,於是慷慨用主人的錢,利用機會為自己打算.主在這比喻中教訓我們,趁今日還活著；手裡還有些屬於不義世代的錢財,該趕快使用它,為了結交朋友及拯救更多的靈魂,使我們有些財寶,積存在求存的帳幕裡.若不會利用機會,我們就比不義的管家更愚拙了!主在此比喻中為錢財起了一個名字,祂說:「錢財是別人的東西.」唯有用這錢財為主而活的和所作的,才是自己的,不然就是別人的了.神對無知的財主說:「無知的人呀!今夜我要你的靈魂,你所預備的,要歸誰呢?」趁著錢財在你手中時,好好地利用它,不要作一個不義且愚昧的管家.讓我們曉得活著非為錢財而活,是為事奉基督而活.我們需要錢財,但要善用錢財,為永存的事業而活.求主的話存在我們心裡,讓我們不事奉瑪門,而專心事奉主.
\newline
\newline


\newpage

\section{第44屆~港九培靈會~陳終道~第八講　從四福音裡看舊約的基督}
\label{sec:1362}
\textbf{第八講　從四福音裡看舊約的基督}
\newline
\newline
連結: \href{https://www.hkbibleconference.org/session-message/view/1362}{\texttt{https://www.hkbibleconference.org/session-message/view/1362}}
\newline
\newline
\hyperref[sec:1361]{< < < PREV SERMON < < <}
~
\hyperref[sec:toc]{[返目錄]}
~
\hyperref[sec:1363]{> > > NEXT SERMON > > >}
\newline
\newline
馬太福音 15:17-15:20
\newline
\begin{longtable}{cl}
\hline
\hline
章節 & 經文 (和合本修訂版)\\
\hline
15:17 & \begin{tabularx}{0.7\textwidth}{X} 難道你們不了解，凡進到口裡的，是經過肚子，又排入廁所嗎？ \end{tabularx} \\ \\ \relax
15:18 & \begin{tabularx}{0.7\textwidth}{X} 然而口裡出來的是出於心裡，這才玷污人。 \end{tabularx} \\ \\ \relax
15:19 & \begin{tabularx}{0.7\textwidth}{X} 因為出於心裡的有種種惡念，如兇殺、姦淫、淫亂、偷盜、偽證、毀謗。 \end{tabularx} \\ \\ \relax
15:20 & \begin{tabularx}{0.7\textwidth}{X} 這些才玷污人。至於不洗手吃飯，那並不玷污人。」 \end{tabularx} \\ \\
[1ex]
\hline
\hline
\end{longtable}
以上經文,可簡單的看見整個舊約的中心,是耶穌基督.四福音乃記載耶穌基督,來世的生活與行事；祂是舊約的律法和先知所預言的那一位.兩者在比較上而言,四福音和新約的書信較易了解,所記述的的是以主為中心,而舊約則較難了解,而又不大明顯.所以在我們心中,有時會混亂,到底我們以甚麼態度來讀舊約呢?對舊約之律法應否遵守呢?律法既是神的話語,神的話是不能廢的,又怎能不遵守?那麼是遵守一部份,還是全部呢?其實我們不可能遵守全部的,假如要遵守,則現在就該到耶路撒冷去獻祭了.既是不能,那麼是否該遵守一部份?但聖經的要求又非如此,因凡遵守全律法的,只在一條上跌倒；就是犯了眾條,這是雅各書說的.這可證明,神對我們遵守律法的要求是全部.主說:「莫想我來是廢掉律法,我來不是要廢掉,乃是要成全.」到底是誰成全?是我們還是基督?主如此說是要我,你成全嗎?抑或是祂成全呢?祂又說:你們的義若不勝過法利賽人和文士的義,斷不能進天國.」文士和法利賽人的義我們能勝過嗎?法利賽人一個禮拜禁食兩次,請問你一年有否禁食兩次,如沒有,則不能勝過法利賽人,則亦不能進天國了!到底是憑我們的義可以得救,抑或是接受耶穌基督的義?感謝主!不是我們成全律法,乃是主基督成全.我們並非憑守律法而來的義,可以稱義得救；乃因信基督而有的義,故這義是勝過文士和法利賽人；如沒有這義,則斷不能進天國.
\newline
\newline
路加廿四章記載主復活後,向以馬忤斯的兩個門徒顯現.兩門徒實在不明白主為何受害,主既是要來的彌賽亞,但為何會被人釘十字架,他們實在不明白.所以主當時就把摩西與眾先知所講的,凡是經上指著祂的話,都給這兩門徒講解.摩西與眾先知所講的中心題目是甚麼?就是耶穌基督,亦即整個舊約的中心是耶穌基督.
\newline
\newline
當腓力向拿但業介紹耶穌時,他說:「摩西律法上所寫的,和眾先知所記的那一位,我們遇見了.」這可見摩西律法上所寫的和眾先知所記的,都是向著「那一位」為中心而寫,故舊約和新約都是以耶穌基督為中心.
\newline
\newline
聖經被稱為新舊約全書,十分有意思.聖經是名稱,新舊約則將其主要內容表明出來,因內裡包括了兩個主要的約:一是舊約(出十九24)中,神與以色列人所立的約.(出二十四章)是神與以色列人立舊約的經過.摩西自己上山領受神的話,然後向百姓宣佈；百姓願意表示,然後摩西回覆神,再下山殺了羊,然後把血一半洒在約書上,一半洒在百姓身上；舊約就是如此成立了.那麼舊約的內容,就是許多的律法和條例,這都是表明和說明耶穌基督,為何舊約時代要人獻燔祭?因燔祭是屬耶穌基督的,祂把自己獻上,得蒙神的喜悅；故我們在祂裡面亦可靠祂將我們獻上,成為神所喜悅的祭.為何要獻素祭?素祭乃描寫基督在世上的生活,是聖潔的,是為我們受苦的,是神所悅納的.為何要獻平安祭,贖罪祭......?每一個祭都是預表耶穌基督.為何要守安息日?在舊約每個以色列人在安息日時不作工,是否藉此就可將安息日的意義完全表達清楚?其實這不過是外表守安息,心裡卻沒有守安息.按先知以賽亞書記載,我們可以知道,安息日不單不可作自己的工,亦不能隨自己的私意,講自己的私話,那才是真正守安息.故安息日的意思,乃是在耶穌基督裡,才有真正的安息.為何神用第七日安息?乃因祂創造的工作完了.以色列人將這日作為安息日,以紀念神如何帶領他們出埃及；但神在這日子,是表明我們在耶穌基督裡面,使我們成為新造的；祂從死裡復活,使我們因祂而與祂同死而復活.我們歸信祂時,祂就在新造裡,使我們造成,我們在祂裡面可以享受.為何舊約中不許人吃血?利未記十七章解釋說:血裡有生命可以贖罪.請問這血是甚麼血?可以有生命和可以贖罪呢?是牛羊的血嗎?如果說血內有生命,單憑血而言,你吃了使身體有生命,那麼,麵包不也可以維持我們的生命嗎?所以血內有生命和贖罪的意思?就是耶穌基督的血有生命.為何舊約不可吃血,乃是教導以色列人要尊重這快要來的耶穌基督的寶血,這是不可吃血的主要原因.在舊約時代有許多儀式,規條,人物,都是要表明那要來的耶穌基督,神是否偏待人而揀選以色列人呢?不；其實神揀選以色列人,不過是將他們作例子,證明人不能靠守律法得救.一方面為要成就亞伯拉罕的應許,另一方面以色列人被揀選出來,有神的訓練,祂差遣先知特別教訓他們,有祭司的制度,知道如何去事奉.假如世上有一民族,他們憑守律法稱義的話,那該是以色列人了.因為他們是受神特別訓練,有特別優越的條件,有特別的教導；但以色列人尚不能好好地,遵守神的律法,那麼那些未受過神教導的其他外邦人,就更不行了.故神選以色列人,乃證明了在祂自己之外,另尋求一個救法；這救法乃是要來的耶穌基督,故整個舊約乃是以耶穌基督為中心.為何會有祭司之制度,乃是證明不是世上每個人都能事奉神,只有蒙揀選的人才能事奉神；只有亞倫的子孫才可作祭司,只有我們這位屬天的大祭司耶穌基督,我們是祂子孫,我們可以事奉祂.故今天我們到祂面前,不是到舊約所表明的那光景去事奉神,因舊約是一種聚述,是描寫,是要表明以後來的那一位；正如希伯來書十章說:律法目的並非要人遵守,或靠守律法得救,只是要表明以後要來的而已.舊約不是有十誡嗎,內裡有十個「不」,我們的神是否充滿了「不」?這是神的目的嗎?神的「可以」在何處?乃是在耶穌基督裡,祂先將「不可」的告訴我們,然後又給我般們一條可行的路,這路就是耶穌基督.故整個舊約的目的是將我們帶到耶穌基督面前.保羅說:「我們能承擔職事,不是憑字句,乃是憑著精義.」保羅所說的並非教會的執事,乃是說從神所領受的職份,(林前三章),他能在神前領受屬靈的職份,可事奉他的主,他的神,非憑字句；如憑字句,保羅沒有資格作事奉神的人.因他雖是希伯來人,但他是屬便雅憫支派而不是屬利未支派的,他也不是亞倫子孫.按字句他不能事奉神,你知我也不能.保羅說憑「精義」,「精義」乃是舊約律法中所表明的那一個,那表明的是誰?就是主耶穌基督.今日我們在基督裡,就可承擔職事,有權利可事奉祂.因我們是被揀選的族類,有君尊祭司的職份；這身份在舊約中得不著,但在新約中,在耶穌基督裡,就可以得到的,故在新約耶穌基督的恩典下,可享受這福份.從舊約中的許多的應許和預表,可以看見耶穌基督,乃是舊約和新約的中心.神豈不是在咒詛蛇的時候應許說:「女人的後裔要傷蛇的頭,蛇要傷他的腳跟.」女人的後裔是指耶穌.亞當犯罪後,聖經就開始預表這位要來的耶穌基督.為何亞伯拉罕如此寶貴要得著以撒?乃因他看重神的應許,因神應許他子孫要如海邊的沙,天上的星那麼多.神應許說萬國要因他得福,如果他沒有兒子,那又怎能呢?所以他盼望得著以撒.在創二十二章18節中說:「萬國必因你的後裔得福.」(加三16)解釋這「後裔」並非指眾子孫,而是指那一位子孫,就是基督,故保羅在(加三章)中解釋了神的話.在(約八56)主對猶太人說:「你們的祖宗亞伯拉罕歡歡喜喜的仰望我的日子,既看見就快樂.」猶太人聽了就生氣說:「你還沒有五十歲怎會見到亞伯拉罕?」主告訴他們說:「未有亞伯拉罕以先,已經有了我.」猶太人就想用石頭打耶穌.亞伯拉罕所信和所盼望的是甚麼?主說他所盼望的是祂,亞伯拉罕是歡歡喜喜仰望他的日子.神豈不也應許大衛說:「他的寶座直到永永遠遠.」但我們知道所羅門的寶座並沒有到永遠,可是所羅門是預表基督,主才是真正坐寶座的王(太一1)中說及大衛的後裔,這「後裔」二字,原文指基督.耶穌基督是大衛的兒子,因所羅門預表基督.故在舊約中的一切好像是偶然地記載,其實並非如此,因歷史向來是挑選過才記載.聖經亦然,可以表明耶穌基督的才記下來,以及那能說明神救贖計劃的事也記下.故舊約的記載乃指向耶穌基督.亞當不就是祂的預表嗎?有許多人以為預表的講解,乃是亂講聖經,其實神對預言也有祂自己的解釋,是不能亂講的.亞當預表耶穌基督,乃是羅馬書中明說的(羅五14)林前十五章中也有提及,耶穌基督乃是末後的亞當.在(弗五章)中保羅說及基督與教會的關係,用夫妻相待來說明；說及丈夫如何愛妻子,妻子如何順服丈夫.到最後結論就說:我是指基督與教會而言,這是極大的奧秘.夫妻怎麼可成為一體?這是指基督與教會而言,是屬靈的奧秘；我們如何去領受,就用夫妻來表明.以夫妻之關係,幫助我們明白基督與教會之關係；但這關係仍是奧秘,這要我們到主面前才能完全明白.故亞當,以撒也是預表耶穌基督.摩西在申命記中,神對他說,我要在你弟兄中興起一位先知像你；這一位先知就是在(約一45)中摩西在律法上,和先知上所記載的那一位.施洗約翰在施洗時,有人問說:「你到底是誰?是否那先知呢?」意思就是說:他是否神所應許摩西所興起的「那先知」呢?這可證明耶穌基督就是舊約所預言要來的「那一位先知」.在(約三章)中,耶穌與尼哥底母論及重生的道理,後來引證以色列人在曠野舉銅蛇的事.祂說:摩西怎樣在曠野舉銅蛇,人子也照樣被舉起來,叫一切信祂的都得永生.摩西領以色列人在曠野中前行,他們心中所相信的是誰?他們仰望銅蛇,銅蛇的背後是說耶穌基督,這是耶穌自己講的,是聖經自己的解釋.大致上可說:所有預表皆指基督而言.有時我們用一件事象徵一種教訓,嚴格而言「預表」這名詞好像不大合適,我們可說是一種象徵.因聖經內的預表,都是在基督與教會內.在聖經內也見到許多信心的偉人,他們所信的是誰?乃是基督.我們已提過亞伯拉罕及以色列人,那摩西信的是甚麼?(希十一章)中說:摩西看為基督所受的凌辱,比埃及的財物更寶貴.難道摩西去到曠野不單是為了殺一個人,而是他心中有一盼望；就是盼望那一位要來的耶穌基督.所以他看為基督所受的凌辱,比埃及的財物更寶貴.
\newline
\newline
猶大書說,亞當的第七世孫以諾,曾預言耶穌基督要帶千萬的聖者降臨,在那時要審判那些存剛愎的心背叛的人,假師傅和教會中的假基督.那麼,亞當的第七世孫以諾的心中所信的是誰?他既預言耶穌基督的第二次降臨,則他心中所信的當然是耶穌基督.並非已來的,而是仍未來的耶穌基督.故猶大書中說,以諾與神同行,不單是信靠神而已,其實他還有傳道和說預言的.如按時間而言,他可稱為最早的先知.按聖經所記載的次序,創世記並無提及他說預言,但猶大書卻有記下.按時間來說,是比亞伯拉罕更早的,故整個舊約都是向著一個目的,就是描寫耶穌基督如何來.
\newline
\newline
今日我們在這裡面可以憑恩典得救,那麼我們和舊約的關係是甚麼?舊約對我們有何用處?律法的用處,並非要我們靠著它得救.在羅馬書和加拉太書說得很清楚,例如有病時照X光,可以知道我們身體有何病症,只可知是病得深或淺?但X光並不能醫我們的病.這就是律法的功用,它不能使我們得救,但卻使我們知道是個無盼望,罪惡深重,要滅亡的人；不能靠自己得救,而要在自己之外,找尋一救法.保羅在加拉太書中說:律法是我們訓蒙的師傅引導我們到耶穌基督那裡.我們明白自己有罪,不能自救,神要刑罰罪惡,那我們怎辦呢?神就為我們預備救法,就是耶穌基督.故整本舊約,無論是十誡,或獻祭之規矩,或神吩咐以色列人的各種生活規條；皆有一共同作用,就是必需要尋求自己以外的救法.
\newline
\newline
(雅各二10)說:「我們所遵守的,犯了一條,就是犯了眾條.」在上下文是提及凡是憑外貌看人的就是被律法定為犯法的,這就是犯了全律法,該被審判.按外貌待人,在十誡是無的,但卻是一樣,因全律法都有同一功用,必在律法外尋求救恩.保羅在(羅十4)說:律法的總結就是基督.使凡信祂的,都得稱義,整個律法有一個總目的,就是將基督介紹出來.律法將許多路都攔阻了,讓我們只走一條路,就是到耶穌基督面前.我們今日既已蒙恩,應在祂的恩典底下,將祂的形像和榮美活出來,把實際的屬靈活出.所有律法所要求及描寫,都是儀式的；而這儀式的目的,是要我們進入實際,這實際是基督.故耶穌基督來了,在新約之下的人,所要遵守及所要生活的,不是單憑字句,或做到儀式所要求的；而是要耶穌基督活在心中,讓生命實在活出來,才能滿足主的心意.舊約只要求把牛羊獻上,但在新約中,希伯來書已說:燔祭和贖罪祭是你不喜歡的；你為我預備了身體,為了甚麼?該知道單靠牛羊的血贖罪而心裡不悔改,人又不歸主；你身體又不分別為聖,難道神也會被你欺哄嗎?我犯了罪該怎辦呢?多捐些錢吧!難道神會悅納嗎?神的要求是如此嗎?舊約的教導是屬於較嚴謹的一類,不可這樣,不可那樣,但新約的要求,是要我們將心獻給祂.在恩典之下沒有「不可」,但其實神向我們的要求更多,且道德更高,祂不是要我們表面上去遵行,而是實際地遵行.非要我們付容易的代價,而是實際的代價.故今日活在新約之下的人有一種律法,就是無形的律法；是愛的律法,凡是違背愛的律法的,我們就不可作,有聖靈在我們心中,在凡事上教導我們怎樣行神的旨意.讓神教導我們活在新約的律法下,而不是舊約的律法下.
\newline
\newline


\newpage

\section{第44屆~港九培靈會~陳終道~第九講　基督再來}
\label{sec:1363}
\textbf{第九講　基督再來}
\newline
\newline
連結: \href{https://www.hkbibleconference.org/session-message/view/1363}{\texttt{https://www.hkbibleconference.org/session-message/view/1363}}
\newline
\newline
\hyperref[sec:1362]{< < < PREV SERMON < < <}
~
\hyperref[sec:toc]{[返目錄]}
~
\hyperref[sec:1260]{> > > NEXT SERMON > > >}
\newline
\newline
約翰福音 14:1-14:31
\newline
\begin{longtable}{cl}
\hline
\hline
章節 & 經文 (和合本修訂版)\\
\hline
14:1 & \begin{tabularx}{0.7\textwidth}{X} 「你們心裡不要憂愁；你們信神，也當信我。 \end{tabularx} \\ \\ \relax
14:2 & \begin{tabularx}{0.7\textwidth}{X} 在我父的家裡有許多住處；若是沒有，我就早已告訴你們了。我去原是為你們預備地方去。 \end{tabularx} \\ \\ \relax
14:3 & \begin{tabularx}{0.7\textwidth}{X} 我若去為你們預備了地方，就必再來接你們到我那裡去，我在哪裡，叫你們也在哪裡。 \end{tabularx} \\ \\ \relax
14:4 & \begin{tabularx}{0.7\textwidth}{X} 我往哪裡去，你們知道那條路。」 \end{tabularx} \\ \\ \relax
14:5 & \begin{tabularx}{0.7\textwidth}{X} 多馬對他說：「主啊，我們不知道你去哪裡，怎麼能知道那條路呢？」 \end{tabularx} \\ \\ \relax
14:6 & \begin{tabularx}{0.7\textwidth}{X} 耶穌對他說：「我就是道路、真理、生命；若不藉著我，沒有人能到父那裡去。 \end{tabularx} \\ \\ \relax
14:7 & \begin{tabularx}{0.7\textwidth}{X} 既然你們認識了我，也會認識我的父。從今以後，你們就認識他，並且已經看見他了。」 \end{tabularx} \\ \\ \relax
14:8 & \begin{tabularx}{0.7\textwidth}{X} 腓力對他說：「主啊，將父顯給我們看，我們就知足了。」 \end{tabularx} \\ \\ \relax
14:9 & \begin{tabularx}{0.7\textwidth}{X} 耶穌對他說：「腓力，我與你們在一起這麼久了，你還不認識我嗎？看見我的就是看見了父，你怎麼還說『將父顯給我們看』呢？ \end{tabularx} \\ \\ \relax
14:10 & \begin{tabularx}{0.7\textwidth}{X} 我在父裡面，父在我裡面，你不信嗎？我對你們所說的話不是憑著自己說的，而是住在我裡面的父在做他的工作。 \end{tabularx} \\ \\ \relax
14:11 & \begin{tabularx}{0.7\textwidth}{X} 你們要信我，我在父裡面，父在我裡面；即使不信，也要因我所做的工作信我。 \end{tabularx} \\ \\ \relax
14:12 & \begin{tabularx}{0.7\textwidth}{X} 我實實在在地告訴你們，我所做的工作，信我的人也要做，並且要做得比這些更大，因為我到父那裡去。 \end{tabularx} \\ \\ \relax
14:13 & \begin{tabularx}{0.7\textwidth}{X} 你們奉我的名無論求甚麼，我必成全，為了使父因兒子得榮耀。 \end{tabularx} \\ \\ \relax
14:14 & \begin{tabularx}{0.7\textwidth}{X} 你們若奉我的名向我求甚麼，我必成全。」 \end{tabularx} \\ \\ \relax
14:15 & \begin{tabularx}{0.7\textwidth}{X} 「你們若愛我，就會遵守我的命令。 \end{tabularx} \\ \\ \relax
14:16 & \begin{tabularx}{0.7\textwidth}{X} 我要求父，父就賜給你們另外一位保惠師，使他永遠與你們同在。 \end{tabularx} \\ \\ \relax
14:17 & \begin{tabularx}{0.7\textwidth}{X} 他就是真理的靈，是世人不能接受的。因為他們既看不見他，也不認識他；你們卻認識他，因他常與你們同在，也要在你們裡面。 \end{tabularx} \\ \\ \relax
14:18 & \begin{tabularx}{0.7\textwidth}{X} 我不會撇下你們為孤兒，我必到你們這裡來。 \end{tabularx} \\ \\ \relax
14:19 & \begin{tabularx}{0.7\textwidth}{X} 再過不久，世人不再看見我，你們卻會看見我，因為我活著，你們也要活著。 \end{tabularx} \\ \\ \relax
14:20 & \begin{tabularx}{0.7\textwidth}{X} 到那日，你們就會知道我在父裡面，你們在我裡面，我也在你們裡面。 \end{tabularx} \\ \\ \relax
14:21 & \begin{tabularx}{0.7\textwidth}{X} 有了我的命令而又遵守的人，就是愛我的；愛我的人，我父要愛他，我也要愛他，並且要親自向他顯現。」 \end{tabularx} \\ \\ \relax
14:22 & \begin{tabularx}{0.7\textwidth}{X} 猶大（不是加略人猶大）問耶穌：「主啊，為甚麼親自向我們顯現，而不向世人顯現呢？」 \end{tabularx} \\ \\ \relax
14:23 & \begin{tabularx}{0.7\textwidth}{X} 耶穌回答他說：「凡愛我的人就會遵守我的道，我父也會愛他，並且我們要到他那裡去，與他同住。 \end{tabularx} \\ \\ \relax
14:24 & \begin{tabularx}{0.7\textwidth}{X} 不愛我的人就不遵守我的道。你們所聽見的道不是我的，而是差我來之父的。 \end{tabularx} \\ \\ \relax
14:25 & \begin{tabularx}{0.7\textwidth}{X} 「我還與你們在一起的時候，已對你們說了這些事。 \end{tabularx} \\ \\ \relax
14:26 & \begin{tabularx}{0.7\textwidth}{X} 但保惠師，就是父因我的名所要差來的聖靈，他要把一切的事教導你們，並且要使你們想起我對你們所說的一切話。 \end{tabularx} \\ \\ \relax
14:27 & \begin{tabularx}{0.7\textwidth}{X} 我留下平安給你們，我把我的平安賜給你們。我所賜給你們的，不像世人所賜的。你們心裡不要憂愁，也不要膽怯。 \end{tabularx} \\ \\ \relax
14:28 & \begin{tabularx}{0.7\textwidth}{X} 你們聽見我對你們說過，我去了還要回到你們這裡來。你們若愛我，就會因我到父那裡去而喜樂，因為父比我大。 \end{tabularx} \\ \\ \relax
14:29 & \begin{tabularx}{0.7\textwidth}{X} 現在事情還沒有發生，我預先告訴你們，使你們在事情發生的時候會信。 \end{tabularx} \\ \\ \relax
14:30 & \begin{tabularx}{0.7\textwidth}{X} 我不再和你們多說了，因為這世界的統治者將到，他在我身上一無所能。 \end{tabularx} \\ \\ \relax
14:31 & \begin{tabularx}{0.7\textwidth}{X} 我這麼做是照著父命令我的，為了讓世人知道我愛父。起來，我們走吧！」 \end{tabularx} \\ \\
[1ex]
\hline
\hline
\end{longtable}
主耶穌曾應許門徒說:你們心裡不要憂愁,你們信神也當信我,並安慰門徒說:祂去是為他們預備地方去,並且祂也再來接他們到祂那裡去.(約十四章)
\newline
\newline
當主耶穌被帶到公會受審時,祭司長和全公會,尋找假見證,要控告祂,治死祂,他們叫祂起誓,考問祂,是神的兒子基督不是?主回答說:「你說的是,然而我告訴你們,後來你們要看見人子坐在那權能者的右邊,駕著天上的雲降臨.」(太二十六:64)當時大祭司撕裂衣服說:「祂說了僭妄的話,我們何必再用見證人.」他們就用這話作把柄,定耶穌的罪.主耶穌在敵人尋找把柄定祂罪時,就說出祂要再降臨,這是非常重要的事.我們確實相信,祂必定要再來!祂再來時是駕著天上的雲彩降臨,帶著榮耀的身體降臨的.我們信耶穌就可以得救,祂也就住在我們心中；並且給我們活潑的盼望,在祂駕雲榮耀降臨時有分.好些人把主再來的事說得抽象了,其實主自己說得十分具體,說明祂如何的降臨.
\newline
\newline
關於主再來的次序,大概可分成兩個主要的步驟:首先從父的右邊降臨到空中(帖前四章),帶著已睡的聖徒一同降臨(與他們的身體復合).祂派使者臨到空中,吹響號筒,信徒則被提到空中與主相遇；在空中有信徒在基督台前受審判,有羔羊的婚筵.在主耶穌從父右邊降臨到空中時,地上就開始末期的大災難.當地上的大災難非常嚴重時,猶太人紛紛回耶路撒冷聖地.在災期快結束時,猶太人被列國圍攻,只得呼求彌賽亞拯救!主基督在這時就從空中帶著眾聖徒降臨到地上,拯救猶太人.祂降臨時,是人人都可以看見的,連扎祂的人也看見祂.祂要結束災難,消滅敵人,審判列國,建立千禧年國度.在這千禧年國度裡,聖經描寫獅子與羊羔在一起,和平相處,地上再無爭戰,撒旦被捆綁丟在無底坑裡,直到千年國度結束,撒旦從無底抗裡釋放出來；就有最後的大戰爭,大戰最後撒旦被扔火湖裡.主耶穌要坐在白色大寶座上,施行最後審判,天地要在祂面前逃避,被琉璜火所毀滅!新天新地要自天降臨,神要與人同在,直到永遠.這是主耶穌再來的大致情形.
\newline
\newline
在馬太二十四章中,主親自詳述祂再來以前,和災難的情形.二十五章亦與主再來的預言有關.(太24章)可說是預言的總綱,如不明白這章聖經而去研究別的預言時,就難有一個綱領.(太廿四27-28),包括了災難之起頭及大災難.在第六節的末句中說:「只是末期沒有到」,這表示在第七節之前還不在末期之內；而在第七節後,就是新的開始.但在第十四節中卻又說:「天國的福音要傳遍天下,對萬民作見證,然後末期才來到.」那麼就可知道7至14節乃是災難的起頭.15至28節是大災難期,讀15節可見:「你們看見先知但以理所說的,那行毀壞可憎的,站在聖殿.16節又說:「那時」──即行可壞可憎的那時 ──,又19節「當那些日子」──指行毀壞可憎的日子──又21節「因那時必有大災難」,故「那時」及「那些日子」,就是大災難,此大災難如何的大?乃是空前絕後的,前所未有的；是災難中最後的一段時期的大災難.「末期」的時間,大概一點來說,就是七年的時間,但嚴格說:「則前三年半之前是災難的起頭,後三年半才是大災難.」
\newline
\newline
在(太二十四章7到14)主稱之為災難之起頭,這與啟示錄之前三年半的災難作比較,就發現非常的類似:
\newline
\newline
許多的爭戰(太24:7),(啟六3至4)第二印.
\newline
\newline
飢荒的災難(太24:7),(啟六8)第三印.
\newline
\newline
瘟疫的災難(路21:11),(啟六8)第四印.
\newline
\newline
地震的災難(太24:7),(路21:11),(啟六12)第六印.
\newline
\newline
近年來,世界各地地震甚多,或許現在科學發達,檢測地震的儀器比較靈敏；但事實上近年來地震愈來愈多.這些跡象,都顯明是大災期的災難的起頭.
\newline
\newline
日頭變黑像毛布,在末後大災難的七年中,日頭變黑不止一次.在(啟六章)揭第六印時說的日頭變黑該是較早的事,因所描寫的是日頭黑像布,是暫時的.因在以後的災禍裡,神還要擊打太陽,要使太陽失去三分之一的光亮,到(啟十六章)倒第四碗時,神把碗倒在日頭上,叫日頭能用火烤人.故在六章中的日頭變黑是暫時的,在(太廿四29-30)中,日頭又變黑,月亮也不放光；接著耶穌就從空中降臨地上.那就是說,當主耶穌從空中降臨時,還有一次日頭變黑,這些都是異象.(太廿四7-14)所說的天然界變化,與啟示錄所說的前三年半的情形是一樣.
\newline
\newline
還有在(太24:9)說:「那時有人要把你們陷在患難裡,也要殺害你們.」這是說到許多信徒將要殉道,與啟示錄中的第五印相似；就是說及許多被殺死的靈魂,正在呼求神為他們伸冤.
\newline
\newline
(太24:11)那時假先知要興起,迷惑許多人.那麼假先知和假基督與第一印相似.(啟六2)揭六印的第一印中,有一匹白馬出來,手拿著弓,好像基督,但卻不是基督.(啟19:12-23),因耶穌基督是戴著許多冠冕,身穿濺了血的衣服.故揭六印中的第一印不是基督,而是假基督,這是前三年半的事.
\newline
\newline
(太24:14),「天國的福音要傳遍天下,對萬民作見證,然後末期才來到.」這末期是指後三年半的大災難.天國福音傳遍天下,在(啟11章)中有兩個見證人在傳福音,他們傳到三年半時就被殺,他們被殺是後三年半之開始.這與(太廿四14)相反.許多人引用(太廿四14)作勸勉信徒去傳福音的經文.但我個人的見解是這福音要傳遍天下,末期才來到；所以這句話是在大災難末期才對.我們今天等候主再來,有很重要的原則,就是主隨時都可來.如果說天國福音要傳遍天下,主才再來,那豈不是還要等很久的時間嗎?其實,要福音傳遍,是極容易的事,只要有一人傳,就立刻可傳遍了.現今的電視和廣播就是一種預備,使大災難前三年半的兩個見證人；可在很短的時間,把福音傳給全世界,福音就立刻傳遍了.
\newline
\newline
(啟11:9),兩個見證人被殺時,屍首不准葬,陳放在那裡有三天半的時間.讓各民,各族,各方,各國的人都來觀看；三天半的時間,怎能讓全世界的人都來觀看呢?但今天有了人造衛星的轉播,就毫無問題了!故天國福音遍傳,該在災期的末期.不過我們必須努力傳福音,因為要使外邦人得救的數目添滿,然後主才再來.我們不知外邦人得救的數目何時才是添滿,所以基督徒有責任傳福音.
\newline
\newline
還有一節重要的經文就是(太廿四15),「那行毀壞可憎的,站在聖地.」這行毀壞可憎的是指敵基督,這經文與(但九27)對照,是說及末後的一個七.敵基督在一七之內,與許多人堅定盟約,在一七之半,就是前三年半時；還容許有祭祀的事存在.這時還未完全暴露他的真面目,仍容許猶太人獻祭,但到了一七之半──即後三年半,他就使祭祀與貢獻止息,也即毀了約,也就更厲害的逼迫神的選民.那麼,敵基督者就出來,故(太廿四15)說:那行毀壞可憎的,站在聖地；就是分別前三年半,與後三年半的中心點.這也是給我們對主再來的前三年半,與後三年半的災難所分別的標準.(15節之前是災難的起頭,15節之後是大災難),那麼,我們教會與基督徒是否在大災難前被提呢?可以肯定的說是在15節之前被提.15節之後是關係以色列人的,因為內裡說及他們逃難時不要守安息日.我們不是守安息日,是守主日；故我們被提後,在地上的是猶太人.但我們不能肯定教會被提,是在災難的起頭之時,(即前三年半)抑或快結束時的那一段時間被提.聖經未有肯定的說明,而在(啟4章)就開始不再提及教會在地上的事,可是,在啟7章就又說及教會仍在地上,故(啟3至7章)中間,我們不知何時被提,為此,我們更要隨時儆醒預備見主.
\newline
\newline
當教會被提時,地上的大災難更厲害,天然界有超然的變化,是一種明顯的記號,讓地上的人知道耶穌基督,快從空中降臨地上.故災難的起頭就是大災難的預兆.在大災難裡,凡經過災難的猶太人,大致都知道耶穌基督再來；因有很希奇的預兆給他們知道.可是我們今天等候主再來,並沒有任何特別之預兆,沒有任何變化,也沒有日頭變黑的現象,沒有大災難時那種嚴重的災禍──帶著神蹟性的災禍──.故不能計算主何時再來,那麼,有沒有其他的預兆呢?有的只是普通的預兆,是有歷史以來都有的,所以無法計算,在(太24:7)前所提及的說有人要迷惑我們,冒基督的名來,在信仰上有異端的發生；有打仗的風聲,這一類的事,是我們向來都聽到的；不過卻愈來愈多了,愈聽則消息愈壞,故我們知道主來的日子實在近了!當主耶穌說完要來的預兆時,就勸勉信徒要儆醒(太廿四44).
\newline
\newline
我們該如何等候主再來?就是要儆醒預備,在(太廿四32)說及幾個重要的比喻.在24章末段論及忠心的僕人,按時分糧給家中各人；但亦另有一惡僕,那麼忠心的與惡的有何分別呢?分別之處乃是惡僕說:我主人必來得遲,就動手打他的同伴,又和酒醉的人一同吃喝.主耶穌就說:在想不到的日子,不知道的時辰,那主人要來,重重的處治他.惡僕與忠僕分別在何處?惡僕並非不儆醒,也不是不忠心,而是他說:主人必來得遲,不需要那麼快忠心殷勤作工,那就是說待主人快來時才殷勤,才把所有的事都做得好,以表現忠心的樣子!一個忠心的僕人,他的忠心在何處呢?就是按時分糧,平常也如此,不論主人在家與否都按時作工；不管主人何時回來,他每天都是一樣.昨日忠心,今日亦忠心!這是最聰明的僕人,也是儆醒的意思.這也就是我們等主來的最好準備,不論主何時再來；我們都天天忠心,天天照常一樣地工作.
\newline
\newline
(太25章)中的十個童女比喻,有五個愚拙的,五個聰明的,到最後主勸勉信徒說:你們要儆醒,因為那日子,那時辰,你們不知道.結論是說明童女忠心地準備.十個童女,聰明的和愚拙的都打盹,但聰明的雖打盹而燈裡有油；那燈不是等新郎來才發亮,而是未來時已發亮,來了也發亮.而愚拙的則燈在新郎未來時已亮,但新郎來了卻不亮!意即他們的光在主來之前發亮而已.故儆醒的意思是要到主來了之後也要亮,不是虛假的稍為受聖靈的感動,也不是只有一些行為發出來；別人以為你是受聖靈感動了,而是要有真實生命所發出的行為.那種光能在基督台前一樣發亮的,這就是儆醒.該知道主耶穌並不是那麼容易欺騙的.你的光要經得起考驗,要認識到那充滿虛假的光明就是不儆醒.求主教我們作一個忠心儆醒的僕人,實實在在的活在主基督裡.這樣,不論祂何時來,我們都能經得起考驗,是聰明的童女,忠心的僕人；讓我們見主面時,都是歡歡喜喜的.阿們!
\newline
\newline


\newpage

\section{第45屆~港九培靈會~吳勇~第一講}
\label{sec:1260}
\textbf{耶穌的寶血}
\newline
\newline
連結: \href{https://www.hkbibleconference.org/session-message/view/1260}{\texttt{https://www.hkbibleconference.org/session-message/view/1260}}
\newline
\newline
\hyperref[sec:1363]{< < < PREV SERMON < < <}
~
\hyperref[sec:toc]{[返目錄]}
~
\hyperref[sec:1262]{> > > NEXT SERMON > > >}
\newline
\newline
(一九六八或一九六九之間)我曾被邀在港九培靈會講道.可惜今年我患心臟痳痹症,所以有點隱憂,但我深信主必加我力量,使我能夠領完這十天的聚會.時代不斷地轉變,越變越不像樣子了.這次培靈會我講的總題是「時代與異象」；願神引導我,所講的能適合各位的需要.
\newline
\newline
最近有一本書在國內外非常風行,書的內容是論到基督徒的正常生活.我現在所要題論到的是:怎樣才是一個正常的基督徒.今天晚上我要由傳福音開始,明晚將講得救蒙恩之人生活的樣式──生命的異象.後晚要提到我們生命時有軟弱,如何才能恢復.以後,要講基督的使命.
\newline
\newline
(出十二13)論血的記號.(徒廿25)論血要代價買的.(弗二13)論靠主寶血與神親近.(來九14)論血冼淨我們.(約壹一7-9)也是論到血洗淨我們的罪；但是兩處冼淨的對象不同.(弗九14)題到血冼淨我們的良心.(約壹一9)題到血洗淨我們的不義.(啟一5)把血和愛相題並論.(啟十二11)題到血可以勝過撒旦.
\newline
\newline
今晚在座中,可能基督徒居多,但將耶穌的福音,複習一次,也是需要的；因為福音是每個人的根基.在啟示錄二章曾提到主責備老底嘉教會；把起初的愛心離棄了,回想我們初信主時,對主有單純的愛；可是漸漸地就隨著時間而冷淡了,因此,靈性停頓而不長進.
\newline
\newline
上面七處的經文都論到血.茲分為兩方面來思想:
\newline
\newline
(一)血的顯明　一個人的血良好,這人必然身體健康正常；相反地,如果血不好,健康便有問題.
\newline
\newline
(二)血的功效　人體有紅白兩種血球.紅血球輸送氧氣到身體各器官,心臟需要氧氣,紅血球就輸送氧氣到心臟；腦要氧氣,它就輸送氧氣供腦部之需.白血球也有它的功用,當人體受到意外損傷時,空氣中的細菌,隨時可能由傷口侵入人體,而造成中血毒；白血球就奮力抵抗,把細菌消滅,保障我們身體,不受嚴重的破壞.所以,紅血球和白血球各有其功用.
\newline
\newline
現在我們從神的方面看,在舊約記載血的事很多.祭司把牛羊獻在祭壇上,又用血灑會幕,抹器皿.當祭司就職時,把血抹在耳朵上,抹在大拇指或腳趾上.血在聖經裡,佔了很重要的地位.從聖經來看,血有兩種的顯明:
\newline
\newline
一,血顯明神的大愛.(啟一5)把血和愛擺在一起,這血就是顯明神的大愛.
\newline
\newline
人與人之間,有很多方式可以表示愛的.丈夫對妻子表示愛,可在她生日的時候,送些她所喜歡的東西做禮物,以表示愛忱.丈夫對妻子的愛,如果單用言語或文字表明,恐怕不切實際,所以藉著一件東西,把他的愛表明出來.除了用禮物表明愛之外,朋友之間,當灰心失意有難處時,我們就向他表同情.如何表明對他的同情呢?就是與他同憂愁同悲傷,用關切來表示愛.人與人之間,可以用物質來表示愛,也可用心靈的關切來表示愛；但有一種更高超的表明,是用血來表示愛.在軍隊中,有時為了要對領袖表示盡忠,或對國家元首表示愛戴,就咬破自己的指頭寫血書,用血來表示愛.
\newline
\newline
神對人類,是用多種的方法,多少的次數來表示祂的愛.希伯來書一章一節告訴我們:「神既在古時藉著眾先知,多次,多方的曉諭列祖」神對人類表示愛,不單用一種方法,乃是用多種的方法；不單一次,乃是多次.羅馬書一章廿節:「......藉著所造之物,就可以曉得,叫人無可推諉.」所以神藉著受造之物,來對人表示愛.我們住在地上,有陽光,有空氣,有雨水,這些都是人所必需的.神藉著這些對人類表示愛,而且不止於此.我們單就猶太人來看,神賜給猶太人有特別的聰明,可以說,他們是一種奇異的民族.在美國大學裡,中國人算是夠聰明的民族,讀書的成績優異,但猶大民族常與競爭,他們似乎比中國人更聰明.世界上有很多奇異的人物,如科學家愛因斯坦,政治家華盛頓,都是正宗的猶太人.神不但把聰明賜給猶太人,也把財富賜給猶太人.世界上任何國家,其中最富有的,幾乎都是猶太人.可以說,猶太人是得天獨厚的民族.藉著神所賜給他們的聰明和財富,於一九四八年猶太國重新建立起來.
\newline
\newline
神對人類表示愛是藉著受造之物,而神對人類最高愛的表示,乃是藉著祂的兒子耶穌基督的寶血.「惟有基督在我們還作罪人的時候為我們死,神的愛就在此向我們顯明了.」(羅五8)
\newline
\newline
在挪威海港的地方,某次因風浪大作,船隻翻沉,事後岸上的人趕去救援!祗見遠處的沙灘上,好像有小東西移動.原來有個嬰孩尚生存於已死的母親懷中,他們連忙把嬰孩抱起；發現婦人的胸膛敞開,滿佈血跡.因為她被風浪侵襲,沒有食物充飢,以致沒有乳汁哺嬰,她只好用石塊把自己的乳頭割破,流鮮血代替乳汁來養活她的孩子.為了保全孩子的生命,這母親竟然犧牲自己的性命.切膚之愛真是叫人大受感動!
\newline
\newline
神曾經讓祂兒子的寶血,從十字架上流下,用血來表示祂最高超的愛.因此在啟示錄把愛和血擺在一起.如果有人問我:為何神讓祂兒子從天上來到地上?我給各位的答覆是:「神的愛.」為何神的兒子從寶座降到馬槽?也是「神的愛.」為何神的兒子在十字架上流血?更是「神的愛.」.如果有人問我:人的罪怎能得赦免?我給各位的答覆是:「神兒子的血.」怎樣才能上天堂?也是靠著「神兒子的血.」怎樣才能成為神的兒女?也是「神兒子的血.」神藉著祂兒子,將祂對人類的愛流露出來.
\newline
\newline
二,血顯明代價.(徒廿28):「就是祂用自己血所買來的.」平常我們買物,付出的代價越高,表明越有價值.我到過南美洲,常聞有綁票的事,他們所綁的多數是強國的大使；因為強國的大使,較小國的大使代價為高.人對人的看法,都是以地位財富,來衡量身價.但神卻不看人身份,不看人的地位,也不看人的財產.神對人的看法,乃是人的生命,比全世界還大.在四福音裡主耶穌說:「人若賺得全世界,賠上自己的生命,有什麼益處呢?」以人的看法,人不過是一身的血肉,骨頭,死了以後,便與草木同朽,如燈火熄滅,一切歸於無有.人對人的看法是卑賤的,然而神卻看人的裡面.神用塵土造人的身體,還把氣息吹到人的鼻孔裡,叫人成為有靈的活人；在神的眼目中,人除了有個暫時的身體之外,還有高貴的靈魂；就是除了有限的軀體,還有永恆的靈魂.神看人的價值是大的,所以祂為我們所付的代價是重的,用祂兒子的血來買我們；由此表明我們在神面前,有很高很重很大的價值.
\newline
\newline
耶穌的血有洗罪的功效　希伯來書說,血洗淨我們心,(原文作良心)約翰壹書說血洗淨我們一切的不義.人身體骯髒,可以用水洗；如洗不淨,可用肥皂洗.但是人的內心並非任何洗潔劑可以洗淨的.這裡所指的心,不是人體內輸送血液的心,也不是現代醫學上換心的那種心,乃是無形的「良心」.保羅被聖靈感動時,他說人的良心好像被熱鐵烙慣一樣.(批發木材的商人,把木編成筏,用熱鐵烙上各種記號,隨水流運送目的地.因為被熱鐵烙過的記號,雖經水飄流浸襲；永不磨滅,所以保羅以此為喻.)人犯了罪,良心便受一次傷,一次次的受傷,好像被熱鐵烙了印記,永遠存在.良心留下犯罪的印記,雖事過情遷,似乎記憶模糊不清,其實罪的印跡卻潛伏著；等到有一日,臨到生死邊緣時,罪都要浮現出來.
\newline
\newline
我曾去探望一位法官,他的妻子是個基督徒.他患了肺癌,情緒不安,呼吸困難,且發高熱.不宜與他詳談,只心簡單的話觸及他內心的感覺.我問他是否知道自己是個罪人?他忽猛然睜開眼睛,態度嚴肅；沒有想到身為法官,竟有人問他這樣的話.我只好站著靜觀其變只見他逐漸溫和回答:「我是個罪人.」接著我告訴他:「你要得永生嗎?罪不能把你帶進永生裡的,所以你感到懼怕和痛苦；只有耶穌的血,能洗淨人一切的罪.惟有耶穌的血能洗淨人的良心,良心得潔淨,心裡才有真正的平安.耶穌的血不但洗淨人的良心,還洗淨人的不義.」
\newline
\newline
啟示錄十九章說,將來,地和海要交出死人,墳墓要交出死人,陰間要交出死人.墳墓要交出死人,陰間要交出死人.死了的人都要站在白色大寶座前,羔羊要把案卷打開,按著各人的行為審判人.那時人生前所說的話,句句都要供出來；所作的事,件件都要顯出來.感謝神!耶穌的血不但能洗淨人的良心,還能洗淨人一切的不義.洗淨良心是關於今世的,叫人在有生之年心裡有平安.洗淨不義是關乎將來的,叫人進入永生.
\newline
\newline
三,耶穌的血能使我們勝過撒旦「弟兄勝過牠,是因羔羊的血......」(啟十二11)
\newline
\newline
一個建築商和朋友打牌,直至凌晨三時才回家,於途中小解,突然感到一陣寒風抖峭,使他毛骨悚然；回頭一看,有個獨眼怪人跟著他,使他覺得痛苦,不安,懼怕!他立即飛跑回家猛力扣門.他妻子啟門,見他面無人色,問發生何事,他驚魂不語,只用手指指著後面,但他妻子卻一無所見.結果他因驚慌過度,發高熱,被送入住院治療.其後他的鄰舍來請我去見他.我和醫生磋商,獲准單獨與病者會面,我問他:何處有蚊子?答:有污水之處就有蚊子.何處有蒼蠅?答:有垃圾的地方就有蒼蠅.我再問:何處有鬼?這個問題他卻無法作答.我告訴他,有罪惡的地方就有鬼,魔鬼專在罪中活動.我再說:你以為要打蚊蠅除滅牠們嗎?他說:此法不行,打了還會再生,惟有清除孳生之源,才能根絕.那麼,有罪之處就有鬼,這問題應該怎樣解決呢?他答:當然要除去罪.「信有罪嗎?」,他說:「人都有罪,多多少少有.」「既如此,你願意對付罪嗎?」他答願意.於是我就教他認罪.我把數十種的罪逐一宣讀,每說到他曾犯過的罪,他舉手說「耶穌的血洗淨我的罪.」當我唸到無可再唸時,我看見他漸漸入睡,我便回家了.次日凌晨四點多鐘,忽然樓下有一陣猛烈的扣門聲把我驚醒,原來是昨天住醫院那位楊先生,他經過認罪之後,和他糾纏的獨眼鬼不再存在了,所以他急不及待天明,趕來多謝神的恩典.那天正是主日,我和講員磋商,讓他見證耶穌寶血的功效.
\newline
\newline
四,靠耶穌的血得與神親近(弗二13)　「我們都如羊走迷,各人偏行己路.......」(賽五三6)今日世界的光景,為何每況俞下?為何諸多矛盾?為何一天比一天敗壞?為何充滿了痛苦?正因為「人如羊走迷,各人偏行己路.」以致低落到此地步.從前的人讀書,讀到中學畢業,能寫信,能做賬,便算了.如今不然,中學畢業,還要升大學,出國留學.可見世界進步了,要求大了.今日世界科學發達──人從地上研究到太空,從地球到了月球.雖然如此,但人越來越痛苦,世界越來越混亂,矛盾越來越深；因為罪叫人與神隔開,與神交通斷絕.神是人生命的源頭,人如果沒有神,就好像樹沒有根,慢慢地枯朽了；又好像水沒有源,漸漸地乾涸了.無論人的財富多厚,地位多高,學問多博,如果沒有神,其內心是黑暗和虛空的.人心剛硬!任憑神的呼喚都無動於中.
\newline
\newline
感謝神!藉著耶穌的血,藉著祂的大愛,叫我們的心化開.
\newline
\newline
五,血最後的功效是記號,出埃及記記載,神差使者殺埃及人的長子,如果那家有血的記號便可越過.所以血是以色列人活命的記號.感謝神的恩典!血也是我們蒙恩得救的記號.昔日血潑在門框上,今日血洒在我們的心版上,到了審判的日子,有血的記號就可越過.請問你,你內心裡,曾否有血的記號呢?
\newline
\newline


\newpage

\section{第45屆~港九培靈會~吳勇~第二講}
\label{sec:1262}
\textbf{生命的異象}
\newline
\newline
連結: \href{https://www.hkbibleconference.org/session-message/view/1262}{\texttt{https://www.hkbibleconference.org/session-message/view/1262}}
\newline
\newline
\hyperref[sec:1260]{< < < PREV SERMON < < <}
~
\hyperref[sec:toc]{[返目錄]}
~
\hyperref[sec:1263]{> > > NEXT SERMON > > >}
\newline
\newline
啟示錄 4:6-4:8
\newline
\begin{longtable}{cl}
\hline
\hline
章節 & 經文 (和合本修訂版)\\
\hline
4:6 & \begin{tabularx}{0.7\textwidth}{X} 寶座前有一個如同水晶的玻璃海。寶座的周圍，四邊有四個活物，遍體前後都長滿了眼睛。 \end{tabularx} \\ \\ \relax
4:7 & \begin{tabularx}{0.7\textwidth}{X} 第一個活物像獅子，第二個像牛犢，第三個的臉像人臉，第四個像飛鷹。 \end{tabularx} \\ \\ \relax
4:8 & \begin{tabularx}{0.7\textwidth}{X} 四個活物各有六個翅膀，遍體內外都長滿了眼睛。他們晝夜不住地說：「聖哉！聖哉！聖哉！主—全能的神；昔在、今在、以後永在！」 \end{tabularx} \\ \\
[1ex]
\hline
\hline
\end{longtable}
啟示錄四章6-8題到活物有三種特徵:(一)多眼,「前後遍體都滿了眼睛.」(二)多臉面,「第一個活物臉面像獅子,第二個臉面像牛犢,第三個臉面像人,第四個臉面像飛鷹.」(三)多翅膀,「四活物各有六個翅膀.」按以賽亞書六章二節記載:活物有六個翅膀；用兩個翅膀遮臉,兩個翅膀遮腳,兩個翅膀飛翔.這是「寶座中,和寶座周圍有四個活物」(啟四6中)的情形.「寶座前好像一個玻璃海如同水晶」(啟四6上)耶穌在世時,祂講道常用比喻,用人所明白的來解釋人所不能明白的事；也用人所能看見的,來解釋所看不見的.屬靈的東西是人的頭腦所不易領會的,也是肉眼所不易看見的.
\newline
\newline
神所創造的活物有三種:1.植物──如花草樹木 2.動物──如飛禽走獸 3.人類──人物　我想神所指的活物,並非植物,也非動物,乃是人物；不是指一般的大人物或小人物,而是指著有生命的人物.
\newline
\newline
人藉著福音就得著生命,因為聖經說:「人有了神的兒子就有生命.」那麼,這生命是什麼樣式的呢?約翰被聖靈感動時,他就寫了出來──應當有許多的眼睛,應當有許多的臉面,應當有許多的翅膀.就是說,一個有生命的人應有這些樣式.
\newline
\newline
多眼──預表多看見,表明一個有生命的人,眼睛是很透澈的；是非常深入的,也是極其完全的.聖經說:「前後遍體都滿了眼睛.」所以這許多眼睛是全面性的,預表我們看到將來──「耶穌再來」.如果我們前面有這樣的眼睛,我們就可以很透澈,很深入的,從聖經看到耶穌再來的事.有一次,門徒問耶穌:「你再來的時候有甚麼預兆呢?」聖經雖然沒有記載耶穌再來的年日,可是卻說到許多關於耶穌再來的預兆.
\newline
\newline
(一)但以理書二章尼布甲尼撒王,夢見一個巨大的像,頭是金的,膀臂是銀的,肚腹是銅的,腳趾是半鐵半泥的.整個像是代表外邦時期,到了腳趾已是末期；一半是鐵的光景,一半是泥的光景.我們生在這世代裡,看見過鐵的政治出現,也看見過泥的政治出現,這些事正應驗了那巨像的末期.
\newline
\newline
(二)馬太福音記載耶穌再來之前,國要攻打國,民要攻打民,多處必有地震,飢荒.(太廿四7)回溯第一次和第二次世界大戰,很明顯是國攻打國.隨後又看見南北韓北與南北越之戰,這是民攻打民.「多處必有地震飢荒.」近來很多地區都有地震,而且幅度較從前大,次數也多.至於飢荒,這是誰也意想得到的；因為現代科學發達,機器進步,產量大增.在工業社會中,呈現著顯著的現象,年青人紛紛湧到城市；留在鄉村的,都是年老的或體弱的農民；以致耕種操作,後繼無人.加以人口膨脹,所以多處必有飢荒的預兆已經開始了.泰國原是產米之地,現在竟要排隊購米；泰政府也禁止米糧出口了.
\newline
\newline
一個有生命活物的人,有許多的眼睛,我們從但以理書,看見了主再來的現象；也從馬太福音,看見主再來的現象.主必再來!基督徒應抱何種態度呢?
\newline
\newline
「用前面的眼睛看主耶穌再來」我們候車時,眼睛總朝著車開來的方向,主再來乃是從天而降的,我們也應當仰望著,等候主來.
\newline
\newline
「用後面的眼睛看現在」世界不過是暫時的,轉瞬即逝!有一個比喻說:有一隻狐狸聞到葡萄成熟放香,想跑進葡萄園大吃一頓.牠找到一個洞,僅能把頭伸入,身子太胖卻進不去,那麼只好挨餓三天,待消瘦之後才竄進去.牠在園內大吃特吃,盡享口福之樂,然後得意洋洋從速離開.那知吃得太飽肚子脹,只好再挨餓三天,待消瘦後,才能出去.這雖是簡單的比喻,正啟示我們,人在世界,好像這隻狐狸,空空地來,空空地離開.世界是暫的,空虛的!
\newline
\newline
「用內裡的眼睛看自己」基督徒最可怕的仇敵,不是魔鬼,不是世界,乃是自己.「自己」常叫我們在屬靈的道路上半途夭折.一次的脾氣,一次的貪心,一次的驕傲,都能叫我們的靈性受到虧損.我們對付自己,如同對付牆角上的蜘蛛網；今天掃乾淨,明天又出現.明天掃乾淨,後天又發現,真是應付不暇!今天撒了謊,就去對付撒謊；明天妒忌,就去對付妒忌；後天貪心,就去對付貪心,簡直一輩子都對付不了!聰明的人,會把蜘蛛找到並且打死牠,這樣才是治本之法.保羅說:「就世界而論,我已經釘死在十字架上.」保羅把自我釘死在十字架上.屬靈的毛病是從自己生出來的.我們裡面的眼睛看見以後,就不再體恤自己,愛惜自己；要照著聖經的吩咐,來對付自己.
\newline
\newline
「用外面的眼睛看所發生的事」世事有時發生在人的身上,有時發生在環境上.我們埋怨人,埋怨環境.如果有外面的眼睛,就可以從人或環境看穿過去.耶穌在世時,遇到風浪,祂就斥責風和浪.風浪本來是沒有位格的,但耶穌看穿在風浪的後面,有魔鬼在那裡作祟,耶穌斥責風浪,就是斥責風浪背後的魔鬼.還有一次,耶穌在該撒利亞境內,要上耶路撒冷去,受文士和祭司的害；彼得聽見,就對主說「這事不會臨到你身上!」耶穌說:「撒旦退我後面去吧!」耶穌有外面的眼睛,從彼得身上看見這話是出於他背後的魔鬼.所以我們如果有外面的眼睛,不單能看穿發生事情的人或環境,也能看見背後的東西；或是神的計劃,或是撒旦的詭計.神要藉著人,藉著環境,來修理我們,造就我們.
\newline
\newline
一個有生命活物的人,有許多的臉面.約翰被聖靈感動時,他看見活物有獅子,牛犢,人,飛鷹各個臉面.為什麼有這麼多的臉面呢?正如福音書共有四本,其中有很多重複的記載,既是這樣,有一本便夠了,為甚麼要有四本呢?因為每本福音書是表達耶穌的一方面,四本福音書才能表達耶穌的各方面.馬太所表達的耶穌是君王的樣式；馬可所表達耶穌是個僕人樣式；路加所表達的耶穌是人的樣式；約翰所表達的耶穌是神的樣式.一本福音書不夠表達完整的基督,四本福音書才能表達完整的基督.照樣一個臉面,不夠表明完整的基督,四種臉面才足以把完整的基督表明出來.
\newline
\newline
「人的臉面」耶穌基督是「神子」也是「人子」.耶穌有人的身份,祂應盡人的本分.有一次耶穌對彼得說:「你去釣魚,魚腹裡有一個銀錢,你可拿去納稅.」耶穌是人,祂有國民的身份,應盡國民的本份.當耶穌在十字架上時,對祂母親說:「看你的兒子」又對約翰說:「看你的母親.」祂雖將離人世,還是盡了人的本份,把母親交托給約翰.基督徒得救以後,神揀選我們做人,我們應當有人的本份.人的身份很多,男人在家庭中,有丈夫的身份,應盡丈夫的本份；又有兒子的身份,應盡兒子的本份來孝敬父母.我們在教會,有肢體的身份,應盡肢體的本分；把福音傳給別人,叫基督的國擴大.一個有生命的人,應當按他的身份,而盡其本分.
\newline
\newline
「獅子的臉面」耶穌基督在世時,祂實在有獅子的臉面,聖經說祂是「猶太支派的獅子」.有一次,耶穌進入聖殿,看見那裡有販賣牛羊鴿子的,有兌換銀錢的；耶穌大發義怒說:「我的殿是禱告的地方,你們卻把它變成賊窩了.」那時耶穌顯出獅子的臉面,用繩子當鞭把他們趕走,又推翻桌子,潔淨聖殿,獅子的臉面是對付撒旦和罪惡的.獅子是兇猛倔強的.一個有生命的活物,對罪應倔強兇猛.可惜有很多基督徒,對罪卻是軟弱的,甚至向罪低頭,和世人並沒有分別,所以不能博取世人的欽佩.
\newline
\newline
「牛犢的臉面」　牛有三種表現:──
\newline
\newline
甘心負軛　負軛是難當的,牛卻默然無聲地負軛在肩.人生於世,未免有許多難受難當的事會臨到,如被人誤會,毀滅,藐視.這些都是很難受難當的.在教會中,神的兒女有時受批評,被毀滅,遭藐視,就想找機會報復,以致常有紛爭的事發生.作為一個有生命的活物,應有牛的樣式,來承擔所負的軛.
\newline
\newline
自己勞苦別人享受　世界為何有許多紛爭呢?因為少數人要享受多數人的勞苦,把自己的幸福寄托在別人的勞苦上,所以天下永無寧日!如此世界還有什麼盼望呢?一個有生命的活物,是自己勞苦,讓別人來享受.但願這痛苦的世界,能從基督徒身上得到改觀!
\newline
\newline
最後犧牲被人宰殺　耶穌在世時說,基督徒好像鹽,又像光.我們知道,如果要讓鹽發出功用,必須把它化開；要讓光照亮,腊燭必須溶化,油必須燃燒.基督徒在地上,應有捨己的態度.
\newline
\newline
「飛鷹的臉面」　飛鷹能站在地上,特別有騰空的技能,所以很難在地上捉到牠.保羅在以弗所書,尤其在歌羅西書裡面,題到我們和基督一同復活,一同坐在天上.可惜我們自從信主至今,一直是坐在地上.有人對這段經文的解釋是──到了有一天(意思是將來永世裡)我們要和基督同坐在天上.弟兄姊妹!此段聖經並不是指著將來說的,乃是指著現在說的.請問!我們自信主之日開始,甚麼時候坐在天上呢?我常在機場看見很多母親送別兒女,她們臨別時說「孩子啊!你去了,我也跟著你去.」我也看見過很多做母親的,心情非常緊張；原來她的兒子在戰場上,前線緊張,後方也焦慮.耶穌說:「你的財寶在那裡,你的心也在那裡.」基督徒的寶貝是什麼?是「主耶穌」祂在天上,所以我們的心也應當在天上.世界雖然改變,地雖搖動,我們卻能堅定不移.飛鷹有騰空的樣式,就不至被地上拖住.
\newline
\newline
「許多翅膀」　1,兩翅膀遮臉──意思就是敬畏神,恐怕得罪神.基督徒之所言所行,能夠這樣的,有幾個呢?可惜多少人,非但不怕人,連神也不怕.基督徒若有敬畏神的心,教會就不至於低落到這樣的地步了.　2,兩翅遮身──意思是隱藏自己,掩蓋自己.馬太七章記載:有很多人對主說,主啊!主啊!我們不是奉你的名傳道,奉你的名趕鬼,奉你的名行異能神蹟麼?但主對他們說:「我從來不認識你們,你們這些作惡的人,離開我去吧!」
\newline
\newline
(林前三12)記載,有兩種工作,一是金,銀,寶石的工作:一是草,木,禾楷的工作.我想保羅被聖靈感動而寫這段聖經,絕不以傳道工作,為金銀寶石,以看門工作為草木禾楷的工作.金銀寶石是貴重的,草木禾楷毫無價值的.在保羅眼光中,「重」是貴重有價值的,「輕」是微不足道,毫無價值的.他說:「應當注重極重無比的榮耀,忍受至暫至輕的痛苦.」今天在神的家中,是否祗顯露博士的名銜,與人的奉獻呢?先知以賽亞對希西家王說:「你給巴比倫使者所看見的,都要被擄去.」這句話教訓我們,凡是顯露的東西,在天上都不能存留.可惜現在教會中,信徒的奉獻,和工作,都要顯露.顯露是不錯的,但顯露的東西在神眼目中,都是草木禾楷,沒有價值,不能存留.　3,兩翅飛翔──這是服事的意思.神要什麼樣式的人服事祂呢?就是兩翅遮臉,兩翅遮身的人；敬畏神能隱藏自己的人,才配得服事神.
\newline
\newline
「寶座和玻璃海」　這兩者是一致而分不開的,有寶座就有玻璃海；或者說,有玻璃海,就有寶座.玻璃海四面透明,全面敞開.基督徒有時生活在暗中,有時生活在光中.暗中的生活是秘密不可外揚的,光中生活是可以公開的.弟兄姊妹!如果你沒有玻璃海,你生活就不透明,你所作之事就見不得人.玻璃海表明聖潔,玻璃海的生活,才是聖潔的生活.我們若有玻璃海的生活,我們的裡面才有神的寶座；因為神要在聖潔的人裡面,設立寶座.寶座是代表掌權的意思,掌權是代表管治.今天基督徒的光景,沒有神的掌權和管治,因為他內心沒有神的寶座.弟兄姊妹!我並非要帶領你們回到律法時代,我們並非在律法之起來受神的管治.
\newline
\newline
我在南洋長大,自小喜歡吃辣椒,可以說是辣椒把我養大的；因為沒有辣椒,我就吃不下飯.到了我回國結婚以後,我的妻子不喜歡我吃辣椒,特別禁止我吃頂辣的辣椒.有一次,我要到菲律濱去領會,臨行時師母提我不要吃辣椒；我佯作聽不見,她卻鄭重其事的再三叮嚀.當我抵步講道時,提及我最喜歡吃辣椒,所以當天晚宴,教會特備各式各樣的辣椒款待我,使我垂涎欲滴,但不得不遵守諾言.同席有一個十二歲的孩子,他覺得很奇怪,便發出疑問:「吳長老你曾說最歡喜吃辣椒,現在怎麼不吃呢?」我說「老闆不讓我吃」他又問你的「老闆是誰?這麼厲害」我說是「吳師母」.如果管治我的是在律法之下；我既不在她身邊,她就管不得了我?然而這樣的管治,是基於夫婦的恩愛,所以無論如何都能生效.
\newline
\newline
主對我們的管治,不是用律法,乃是用十架的大愛；我們甘心情願來受祂的管治,被祂愛所感,而受祂的管治.讓主在我們心裡設立寶座,我們一生一世才有玻璃海的生活.
\newline
\newline
正常有生命的基督徒,應當有許多眼睛,許多臉面,許多翅膀.應當讓基督在你裡面設立寶座,做個讓基督可以管治的人.
\newline
\newline
耶穌出來傳福音的時候,祂說:「天國近了,你們應當悔改.」「悔改」就是從原來不受神管治的人,變成受神管治的人.
\newline
\newline
在座各位!你是神所治理的人嗎?你裡面有神的寶座嗎?若你裡面沒有寶座,你能否對主說:「主啊!我願意打開心門,讓你設立寶座在我心裡.」
\newline
\newline


\newpage

\section{第45屆~港九培靈會~吳勇~第三講}
\label{sec:1263}
\textbf{軟弱和恢復}
\newline
\newline
連結: \href{https://www.hkbibleconference.org/session-message/view/1263}{\texttt{https://www.hkbibleconference.org/session-message/view/1263}}
\newline
\newline
\hyperref[sec:1262]{< < < PREV SERMON < < <}
~
\hyperref[sec:toc]{[返目錄]}
~
\hyperref[sec:1266]{> > > NEXT SERMON > > >}
\newline
\newline
撒母耳記上 21:1-21:15
\newline
\begin{longtable}{cl}
\hline
\hline
章節 & 經文 (和合本修訂版)\\
\hline
21:1 & \begin{tabularx}{0.7\textwidth}{X} 大衛到了挪伯的亞希米勒祭司那裡，亞希米勒戰戰兢兢地出來迎接他，對他說：「你為甚麼獨自一人，沒有人跟隨你呢？」 \end{tabularx} \\ \\ \relax
21:2 & \begin{tabularx}{0.7\textwidth}{X} 大衛對亞希米勒祭司說：「王吩咐我一件事，對我說：『我差遣你，吩咐你的這件事，不可讓任何人知道。』因此我已告訴一些僕人到某處去。 \end{tabularx} \\ \\ \relax
21:3 & \begin{tabularx}{0.7\textwidth}{X} 現在你手中有甚麼？請你給我五個餅或是可以找到的食物。」 \end{tabularx} \\ \\ \relax
21:4 & \begin{tabularx}{0.7\textwidth}{X} 祭司對大衛說：「我手中沒有普通的餅，只有聖餅，只能給沒有親近婦人的年輕人。」 \end{tabularx} \\ \\ \relax
21:5 & \begin{tabularx}{0.7\textwidth}{X} 大衛回答祭司說：「我們確實沒有親近婦人，如同往常我出征的時候一樣。平常行路，僕人的身體都還分別為聖，何況今日豈不更使自己分別為聖嗎？」 \end{tabularx} \\ \\ \relax
21:6 & \begin{tabularx}{0.7\textwidth}{X} 祭司拿聖餅給他，因為在那裡沒有別的餅，只有那從耶和華面前撤下的供餅，就是換上熱餅的日子取下來的。 \end{tabularx} \\ \\ \relax
21:7 & \begin{tabularx}{0.7\textwidth}{X} 當日，有掃羅的一個臣僕在那裡，他留在耶和華的面前，名叫多益，是以東人，作掃羅的畜牧長。 \end{tabularx} \\ \\ \relax
21:8 & \begin{tabularx}{0.7\textwidth}{X} 大衛對亞希米勒說：「你手中有沒有槍或刀？因為王的事緊急，連刀劍兵器我都沒有帶。」 \end{tabularx} \\ \\ \relax
21:9 & \begin{tabularx}{0.7\textwidth}{X} 祭司說：「你在以拉谷所殺的非利士人歌利亞的那刀，看哪，裹在布中，放在以弗得後邊。你若要可以拿去，除此以外，再沒有別的了。」大衛說：「沒有甚麼可以跟它比的了！請你給我。」 \end{tabularx} \\ \\ \relax
21:10 & \begin{tabularx}{0.7\textwidth}{X} 那日大衛起來，躲避掃羅，逃到迦特王亞吉那裡。 \end{tabularx} \\ \\ \relax
21:11 & \begin{tabularx}{0.7\textwidth}{X} 亞吉的臣僕對他說：「這不是那地的國王大衛嗎？那裡的人跳舞唱和：『掃羅殺死千千，大衛殺死萬萬』，不是指著他說的嗎？」 \end{tabularx} \\ \\ \relax
21:12 & \begin{tabularx}{0.7\textwidth}{X} 大衛把這些話放在心裡，就很懼怕迦特王亞吉。 \end{tabularx} \\ \\ \relax
21:13 & \begin{tabularx}{0.7\textwidth}{X} 於是他在眾人眼前一反常態，在他們中間裝瘋作癲，在城門的門扇上胡寫亂畫，任由唾沫流在鬍子上。 \end{tabularx} \\ \\ \relax
21:14 & \begin{tabularx}{0.7\textwidth}{X} 亞吉對臣僕說：「看哪，你們看這人瘋了，為甚麼帶他到我這裡來呢？ \end{tabularx} \\ \\ \relax
21:15 & \begin{tabularx}{0.7\textwidth}{X} 我豈缺少瘋子，你們竟然帶這人到我面前瘋癲嗎？這個人可以進我的家嗎？」 \end{tabularx} \\ \\
[1ex]
\hline
\hline
\end{longtable}
大卻是個了不起的人物,他也有軟弱之處.當他投奔亞吉王時,亞吉的臣僕憶述大衛就是從前打死歌利亞,婦女們唱和「掃羅殺死千千,大衛殺死萬萬」的那一位.那時大衛十分懼怕,因為恐怕他們記起舊仇而對他不利；所以他就在亞吉面前「假裝瘋癲,在城門的門扇上胡寫亂畫,使唾沫流在鬍子上.」(撒上廿一13)
\newline
\newline
大衛的恢復(撒上廿二章)　大衛就離開那裡,逃到亞杜蘭洞.他的弟兄和他父親的全家聽見了,就都下到他那裡.凡受窘迫的,欠債的,心裡苦惱的,都聚集到大衛那裡,大衛就作他們的頭目.......大衛對摩押王說:求你容我父母搬來,住在你們這裡,等我知道神要為我怎樣行.......先知迦得對大衛說,你不要住在山寨；要往猶大地去.......」(1-5)從以上兩處經文,我們看見大衛的軟弱和恢復.
\newline
\newline
一個正常的基督徒,難免會軟弱,不過他會恢復過來的；不正常的信徒,就長久不能恢復.
\newline
\newline
軟弱可分為三方面:
\newline
\newline
(一)軟弱的光景從　(撒上廿一章),我們看見大衛因被亞吉王的臣僕揭穿他的身勢,就害怕,假裝瘋癲,正常的樣式就改變了.　1.舉動不正常　「在城門的門扇上胡寫亂畫.」寫和畫可以比作人的舉動；一個軟弱的人,雖然竭力地寫,拼命地畫,不過是胡寫亂畫；沒有意義,沒有目標,雖然幹得十分辛苦,卻徒勞無益,不見果效.　2.言語不正常　「唾沫流在鬍子上」.唾沫比作人的言語.雅各書說「要勒住舌頭」,但一個軟弱的人舌頭是勒不住的；因控制不住,言語隨便,滿口都是批評,論斷和攻擊人的話.
\newline
\newline
一個正常人的光景,他的舉動是有意義,有目標,也可說是有果效的.他的言語是能鼓勵,幫助,安慰和造就人的.到了軟弱的時候,一切都反常了.這是大衛軟弱的光景,也是信徒軟弱的光景.
\newline
\newline
(二)軟弱的結果　「亞吉對臣僕說,你們看,這人是瘋子；為什麼帶他到我這裡來呢?我豈缺少瘋子,你們帶這人來在我面前瘋癲嗎?這人豈可進我的家呢?」(撒上廿一14-15)這裡給我們看見,大衛軟弱以後受到迦特王亞吉的鄙視.照樣,基督徒軟弱也受到外邦人的鄙視.
\newline
\newline
按創世記記載,亞伯拉罕被人看他如王子一樣.王子是尊貴的意思.正常的基督徒,在世人眼中應有尊貴的樣式；因為我們原是從他們中間出來的；世人刮目相看,我們理應與世人有分別.
\newline
\newline
我在台灣曾到一個地方講道,在散會後乘搭火車回台北,半途機車發生故障,不能開行,搭客紛紛下車.那時的士乘機敲榨,到台北每人要八十元.我們寧願等著,希望能到搭到較便宜的車.忽然有一部大巴士從遠處開來,我就站在路中截車.有人問司機要車資多少,他說:「夜深了,大家都急於回家,不必計較,可隨意給就是了.」有人卻對司機說:「我們寧可先小人後君子,免得到時大家爭得面紅耳赤.」司機於是很客氣地說:「如果我要每位十元,算不算太多呢?」大家都認為這價錢很公道.車子到了某處,有個搭客準備下車,但偏找不到錢付車資.司機對他說「不要緊的,十塊錢於你於我都算不得什麼,夜深了,快回家吧!」那位搭客,實在覺很難為情,詢問司機的尊姓大名和地址,改日奉還.司機說:「不必,不必,你下車吧!」車中有位基督徒對我說,司機一定是個基督徒.他走過去問司機你是否基督徒?司機說「為什麼你這樣問我?」那人說:「因為我看你好像是基督徒.」司機說「基督徒就是基督徒,沒有好像基督徒的.」弟兄姊妹!我們在世人面前,應當讓基督的香氣發出來,叫人看我們是尊貴的,是與眾人有分別的.世人可行的我們不當行,世人可說的,我們不當說:我們應有尊貴的樣式,顯在世人之前.當大衛打歌利亞時,曾說:「你來攻擊我是靠著刀槍和銅戟；我來攻擊你是靠著萬軍之耶和華的名.」那時大衛在歌利亞面前,真有王子尊貴的樣式.後來他卻軟弱下去,在亞吉王面前,落到那樣可憐的地步,不但自己受羞辱,神的名也因他受虧損.「讀者文摘」曾載,有一個法國駐蘇聯的大使,中了對方所佈的美人計；當他與一個女人上了不名譽的勾當時,被女人的丈夫發現了!於是他們把握證據來威脅他,敲榨他.結果,不但自己受羞辱,還連累污辱了國家.
\newline
\newline
(三)軟弱的原因　大衛本來是個非常剛強勇猛的人.當他放羊的時候,遇見獅子和熊,啣走了他的羊；他便奮勇追趕,與野獸搏鬥,把羊從獅子口中搶救出來.當他出戰歌利亞時,也可以從他對歌利亞所說的話,看出他的剛強來.想不到日後在亞吉王的臣僕面前,竟然如此膽怯軟弱!參孫在人看來,也是個剛強的人,當人用繩子捆綁他；他把捆綁全掙斷,他的氣力實在驚人!但當他頭髮被剃去時,氣力竟然離開了他.被非利士人捆綁起來,還挖掉了眼睛,把他關在監牢裡.參孫由剛強變成軟弱,是因為是他的頭髮被剃掉,表明他干犯了神的聖潔,所以他就軟弱下來了.至於大衛軟弱的原因,聖經告訴我們:「大衛到了挪伯祭司亞希米勒那裡；亞希米勒戰戰兢兢的出來迎接他,問他說:「你為什麼獨自來,沒有人跟隨呢?大衛說,王吩咐我一件事,不要使人知道.」(撒上廿一1-4)其實,王並沒有吩咐他作一件,不要使人知道的事；是大衛用謊言捏造的.不真實的言語使大衛由剛強落到軟弱的地步.(該一6)末句:「得工錢的,將工錢裝在破漏的囊中.」工錢是工價.因勞力換取工價.工價是勞力的果效,但裝在破漏的囊中,結果都漏光了.許多基督徒,努力追求,得著屬靈豐盛的果效；因工作而得的果效；可惜果效常會漏了出去.從何處漏出呢?我想口是最易漏之處.許多基督徒拼命追求,靈裡卻毫無果效；最大原因,就是全部由口漏掉了!這是基督徒的毛病,也是傳道人的毛病.我們屬靈的經驗和大衛一樣,時而剛強,時而軟弱；有時倚靠神,有時倚靠自己.正常的基督徒雖然也會軟弱,但很快就會恢復過來.
\newline
\newline
大衛的恢復　「大衛對摩押王說,求你容我父母搬來,住在你們這裡,等我知道神要為我怎樣行.」(撒上廿二3)「等」是基督徒屬靈的功課.大衛恢復以後,看見一條屬靈的路,他學習「等」的功課,預備走上這條路.「等」是很難學習的功課,而且在等的過程中常常出了問題.亞伯拉罕七十五歲進入迦南地,神應許要賜他一個兒子.他等到八十五歲,仍沒有兒子；後來妻子撒拉建議他娶使女夏甲為妾.到他八十六歲時,夏甲生了兒子以實瑪利.從此,亞伯拉罕的家就失去了平安,惹上了無限的麻煩.這是因為他在「等」的方面失敗了!掃羅也在「等」的事上出問題.神藉著撒母耳吩咐掃羅,在吉甲的地方等候,他七天後回來獻祭.掃羅等到第七天,看見太陽已經西下,撒母耳還未到.他不再等了,就擅自獻祭；正當在那時,撒母耳就到了.撒母耳對他說「你作了糊塗事,神已另選合祂心意的人了.」神廢棄掃羅,是因他在「等」的事上失敗,神不能再使用他了.在屬靈的功課中,「等」是很重要的.約書亞攻陷耶利哥城,沒有用一刀一槍；因為攻陷耶利哥城的；不是以色列人的力量,乃是神的大能.希西家王被亞述十八萬大軍包圍,並未用一兵一卒,在一夜之間,亞述軍潰敗了；因為得勝屬乎是神,得勝是神的活動,不是希西家的活動.
\newline
\newline
神要屬祂的人放下自己的活動,等候祂的活動.各位同工該回想一下,可能你覺得從前講道,不如現在來得動人；但從前得人的果效,會多於現在.因為從前的工作,是聖靈作工,不是肉體作工.在教會,如果我們要發動一件事,雖然要計劃,也需要金錢,但不必去用人的方法；因為這是神的工作,不是人的工作,聖靈自然會負責,經費也必源源而來.
\newline
\newline
大衛恢復以後,他看見一條屬靈的路,他要學習走這條路；他停止自己的活動,完全讓神活動.大衛本有機會可以殺掃羅:一次在隱基底洞內.大衛入洞,發現掃羅在洞裡；掃羅卻不知道.第二次是在西弗的曠野,神叫掃羅和他的軍兵沉睡,甚至大衛走近他身邊都不察覺.兩次很好的機會可以殺掃羅,但大衛沒有這樣做,並且阻止跟隨他的人殺掃羅.當時掃羅追尋大衛的命,迫使大衛走投無路.既有機會,為何不殺掃羅?因他走上了屬靈的路,他知道掃羅被廢是神的旨意；自己必作王,也是神的旨意,他不必用血氣的手來幫助神.亞伯拉罕要得一個兒子,本是神的旨意,他錯用血氣的方法,就惹上了無限的痛苦.很多基督徒,僅知道神的旨意,而不等候神的時間；每逢有事,就靠著自己的聰明,憑著自己的衝動,沒有憑著聖靈的感動,也沒有等候神的安排.肉體的活動,不但在靈界無濟於事,反而給靈界製造麻煩.今日屬靈的境界,並不是沒有工作,為何這麼冷淡,這麼脆弱?因為其中的工作,出於人的活動,並不是神的活動；憑著人的聰明安排,並非出自聖靈的安排.大衛恢復以後,清楚看見屬靈的路；所以他要等神的吩咐,等神的啟示,等神的帶領.出於神的,在靈界中才有益處,否則毫無俾益.
\newline
\newline
大衛恢復的原因　「凡受窘迫的,欠債的,心裡苦惱的,都聚集到大衛那裡；大衛就作他們的頭目.」大衛既然作了他們的頭目,便負起責任來帶領他們.大衛之所以迅速恢復,是因著他負有責任在身.我用兩個很簡單的比方:小孩子肯用功讀書,是因為要爭取榮譽,榮譽迫使他用功.人為要享福而努力奮鬥.照樣,基督徒屬靈的恢復,責任是他恢復的原因,責任的需要叫他趕快恢復過來.從前摩西帶領了百多萬的以色列民眾,如果他沒有力量,沒有本事,怎能負此重任呢?摩西清楚知道,力量和本事均非他所有；然而責任迫使他四十晝夜在耶和華的聖山上,親近神,追求神.有責任在身,才會想到責任的重要.很多基督徒軟弱一直無法恢復,是因為缺乏了責任感,沒有責任感就沒有力量迫使他恢復.如果我們對主的工作有責任,就不敢軟弱下來.如果我們軟弱,怎能作群羊的榜樣呢?怎能挑擔重大的責任.我一旦軟弱,主的名便要因我受虧損.我們應當趕緊恢復!叫主的名從我身上恢復榮耀.大衛當時有責任,他挑擔家庭的責任,挑擔那批痛苦人的責任.如果我們對家庭有責任,對痛苦人有責任,我們就不敢軟弱.很多責任需要我們去挑擔,「家」是我們的責任,對家人應負牧養的責任.這樣可推展教會的工作,也可減輕教會的負擔.要挑擔社會的責任,貧窮的,苦惱的,我們要把他帶起來,以減少人間之痛苦.如果我們不這樣做,社會的人要恥笑基督徒是寄生蟲,是社會的毒瘤；他們說:這班人對社會不負責任,沒有貢獻,對社會無所俾益.有一次,內子去探望一個傷心的女人,這女人的丈夫在日本因患破傷風症死去,她傷心欲絕,覺得前途茫茫；自己既無本事,兩個孩子幼小,實在無法維持下去.她暗暗買了安眠藥,準備一死了之.當我內子探訪她時,她正哭泣,極力安慰她,她還是不停的哭.後來內子被她的哭所感動,也陪她同哭,足足哭了半天.在太陽下山之前,她忽安靜不來；拿出一包藥片交內子.說:「吳師母,我不用死了,因我現在找到人間的同情者.」如果她自殺了,兩年小孩子,不是要成為社會的累贅,給社會平添了負擔嗎?還有一次,我們到一間寺廟傳福音,有個大漢前來,表明他的身份是某大流氓集團的頭目.我們大吃一驚,以為他來找麻煩.沒想到他的態度很好；他告訴我們,他歸主之後,把流氓組織解散了.因他的恢復,社會減少了邪惡劫殺的事.所以基督徒要負責任,因為責任迫使我們恢復,責任叫我們不敢軟弱.可是有很多基督徒,對家庭,對社會,對肢體一概沒有責任,以為軟弱只關係自己,與別人無關,與神的家無關.他不想人從他身上得什麼,也不想神在他身上作什麼.無怪這樣的人很難恢復；偶有恢復,只不過是血氣的衝動和肉體的活動.這樣在神的家中沒有價值,反而叫神的家受虧損.
\newline
\newline
弟兄姊妹!我們要恢復,要真正的恢復,才能看見屬靈的路,走上屬靈的道路.
\newline
\newline


\newpage

\section{第45屆~港九培靈會~吳勇~第四講}
\label{sec:1266}
\textbf{青年與問題}
\newline
\newline
連結: \href{https://www.hkbibleconference.org/session-message/view/1266}{\texttt{https://www.hkbibleconference.org/session-message/view/1266}}
\newline
\newline
\hyperref[sec:1263]{< < < PREV SERMON < < <}
~
\hyperref[sec:toc]{[返目錄]}
~
\hyperref[sec:1268]{> > > NEXT SERMON > > >}
\newline
\newline
傳道書 11:9-11:10
\newline
\begin{longtable}{cl}
\hline
\hline
章節 & 經文 (和合本修訂版)\\
\hline
11:9 & \begin{tabularx}{0.7\textwidth}{X} 年輕人哪，你在年少時當快樂；在年輕時使你的心歡暢，做你心所願做的，看你眼所愛看的；卻要知道，為這一切，神必審問你。 \end{tabularx} \\ \\ \relax
11:10 & \begin{tabularx}{0.7\textwidth}{X} 所以，當從心中除掉愁煩，從肉體除去痛苦；因為年少和年輕之時，全是虛空。 \end{tabularx} \\ \\
[1ex]
\hline
\hline
\end{longtable}
在傳道書十一章9-10,對青年人用三種不同的稱呼:1.少年人2.幼年人3.一生的開端.我對於聖經的原文不大熟悉,但我曾看過關於此段聖經的解釋.主要的意思是題到青年人的可貴.從這段聖經,我們看到三件事:1.一般人的追求2.一般人的苦惱3.神所要爭取的對象.
\newline
\newline
人的問題很多,我們生活在問題的世界裡面.照一般人的看法,年紀越大問題越多.其實並不盡然.但一個人由小到老,至少會有六個一般性的問題:如下
\newline
\newline
教育
\newline
\newline
就業
\newline
\newline
婚姻
\newline
\newline
兒女
\newline
\newline
年老
\newline
\newline
死亡
\newline
\newline
我今年五十二歲,以上六個問題的大部份已經過了.我不能再讀書；因為我的頭腦和時間都趕不上了.所以教育的問題於我已成過去了.就業的問題也成為決定性了,因為我獻身傳道,一生不會再更改.婚姻的事也成為過去了,因為我在卅年前已成家了.兒女問題我也過去了,我已有七個兒女；(在外國人看來七個兒女是很多的,但以中國人來看卻是很普通的.)年老的問題,我也有了把握,因為信了主,一生前途都在神的手中,不必為年老的問題擔憂.至於死亡的問題,我更是解決了；我知道當我走完世界路程的時候,我要到那裡去.所以,六個問題在我來說都解決了.但在座中有很多人可能第一個問題來了,因他正在讀書.有的已經學成正在找工作.有的已有職業,正尋找對象.有的已結婚還沒有生兒育女.照我的看法,年紀大的,問題雖多,年少的問題也不少.
\newline
\newline
一般人所追求的是什麼?傳道書十一9「在幼年的日子,使你的心歡暢,行你心所願行的,看你眼所愛看的；卻要知道,為這一切的事,神必審問你.」這段聖經我們覺得奇怪,行心所願行的,看眼所愛看的,為何神要審問呢?人所追求的是眼所能看見的身外之物.人可分為兩類,一是普通的,人看為庸弱無能的.有的人以為人生短促,人生痛苦,出生時,自己哭；離世時,別人哭.人生無意義,無價值；因此,他為要保持這段短促的時間,便力謀享受.享受從何來?當然需要金錢來解決.衣食住行,無一不需金錢.衣著時髦講究,食有南樓北館,中菜西餐,樣樣需要金錢.現代城市中因汽車的威脅,工業的威脅,和人口擠迫的威脅,造成空氣污濁；人們都希望到海濱或山上居住,可是在在都需要金錢.至於行的方面,有海洋,有陸上,有空中的交通.航空可以節省時間,航海可以養心怡神.陸上的交通,可以看見城市建築的雄偉,也可看到鄉村風光的美麗.交通方式,任你選擇,但在在需要金錢.金錢不但給我們身外的享受,也給我們裡面的享受.有錢的人,身價提高,人也另眼相看,到處受人歡迎,成了奉承的偶像.人都以錢為可貴之物,所以攪盡腦汁去賺錢,賣盡面子去為錢,拼命追求的也是錢.另有一類的人,他們較有思想,天分較高,他們以為錢財不過是身外物,不值得花費太多的精神和時間去追求；他們注重立功,立業,立德；盼望這些存到永久,留到後代,給後人景仰,給後代紀念.天分較高有思想的青年,不易受金錢的誘惑；有遠見的人,以錢財為身外物,所以就轉向名譽,如飢似渴的追求.無論追求的方向是什麼,人的慾是無止境的!有了錢,還想有更多的錢；有了地位,還想有更高的地位.人的私慾永無止境,以致人與人之間,勾心鬥角,損人利己,自高自大,你爭我奪,相吞相咬；弄得世界一片混亂,非常痛苦,到了不可收拾的地步.
\newline
\newline
當從心中除掉愁煩9節作者體會到人心愁煩.使人愁煩的原因很多.人類為求進步,不斷努力思想,絞盡腦汁,拼命研究改良,把不良者淘汰,讓世界向著新的路上走.推進世界向新的路上走的原動力,是基於人的不滿現狀.從前的人,居於穴中,洞中,樹上；現代的人因為不滿現狀,才把舊有的推翻,改良.從前的人用手操作,如果安於現狀,就沒有機器代替人力,就無法大量生產了.古時出門要步行,如果滿於現狀,就沒有汽車飛機等交通工具了.因為不滿現狀,所以造福人群,促進社會文明進步,人類生活幸福.所以不滿現狀,是世界進步的因素.如果滿足現狀,世界就會停頓下來.以上是不滿現狀優點之一面.另一方面,如古時的婚姻,都是憑著父母之命,媒酌之言,那裡顧到情投意合,互相了解的問題.當一個女子臨婚嫁之前,在家要哭三天；因對方是高,是矮,是胖,是瘦,是好,是壞,無從知道,難免提心吊膽.那樣的婚姻,實在太封建,太落伍了.因而後來人就把它打倒推翻；男女要有自由的選擇,婚姻也要文明的結合.但今日的婚姻又如何呢?結也容易!離也容易!雙方同意,往律師樓辦理手續,就可各走各路了.進步的婚姻,竟落到這樣的境況,落到這樣的下場.這是人的苦惱!
\newline
\newline
知識的進步原是好的,但結果死亡的恐怖也隨之而來.從前的戰爭只是平面的戰爭,現代的戰爭變成立體的戰爭；從前是短距離的戰爭,現代是長距離的戰爭.從前爭戰的消滅是局部的,現在卻是全部的消滅.難道知識的進步不好嗎?為什麼帶來今日的局面?這也是人的苦惱!
\newline
\newline
我們覺得一切問題都關乎經濟,人如果能得到滿足,問題就沒有了.世界上有富人,也有窮人.到底富好呢或窮好呢?結果富則浮,窮則兇；這樣看來,富也不好,窮也不對.這也是人的苦惱!
\newline
\newline
人實在有很多苦惱,所以寫聖經者被聖靈感動時,寫出怎樣克除邪惡,就是克制邪情和惡慾.人是神創造的,神造人有情有欲,但並非邪情惡慾.如果沒有情,夫妻之間就平平淡淡,世界就多麼冷酷.有了情,人間才有一點溫暖,人與人才能彼此懷念.人也有慾,慾是一種需要；如果沒有慾,人就沒有需要,有好的音樂也不想聽；有好的藝術,也不想欣賞,有美味的食物,也不想吃.慾叫我們享受造物的奇妙.情是人與人之間的關係.欲是人與萬物的關係.神造人有情也有慾；但是今日人的情已變成邪情,慾已變成惡慾了!男女戀愛,如今卻變成亂愛.本來,男人應有女人,女人應有男人；這是聖經所許可的,也是道德法律所許可.但變成邪情以後,男人有女人,但女的並非男的妻子；女的有男的,但男的卻非女的丈夫.邪情落到此地步!這個世代和所多瑪,俄摩拉沒有兩樣!滿足個人的需要,是神所許可的,也是道德法律所許可的.用自己的勞力換取所需者,這是神許可的；努力用功爭取優良的成績,這是良心所許可的,滿足自己的需要,用勞力取得的代價,這是正確的,是神許可的.然而現今的慾,變成惡慾!比方我三年沒有吃過肉,不聞肉味,偶然看見隻雞,就順手一把捉回家宰了大吃大嚼.這樣用不正當的手段來滿足自己的需要,便是惡慾.邪情惡慾,理應出於不信的人之間,但是有許多基督徒卻在邪情惡慾上跌倒了,犯上了可怕的大罪!邪情惡慾因何而生?在二百多年前,人們不明白肉類為何會腐臭,牛奶為何會發酸,因為那時醫學和科學還未發明,後來有位醫生巴斯特發現其腐壞之故,是因細菌侵入變質所致.照樣,今日的光景,也因被邪靈侵入；以致情慾變質,成為邪情惡慾.這也是人的苦惱!
\newline
\newline
如何克服邪惡?許多在學校讀書的青年,看見社會敗落,痛心疾首,立下清白的心願要好好做人；才對得起自己的良心,對得起社會殊不知踏入社會以後,願有的志願失去了,心有餘而力不足了,站立不住了.這也是人的苦惱!
\newline
\newline
人的苦惱實在太多!現在人的光景是一種死的光景.耶穌曾說「讓死人埋葬他的死人.」有許多人腳雖能走,口雖能言,手雖能動；但在神看來,他們是死了的人.聖經記載:睚涯魯的女兒死了,做父親的為好傷心痛苦流淚,她沒有知覺.許多青年人走上邪惡的路,父母為他傷心的流淚和嘆息,他都看不見,也聽不見.這和死了的人有何分別?拿因城有個寡婦的兒子,死了以後,人把他抬往荒野,這與今天的青年有何分別!許多青年人,隨從今世的風俗,以為自己是時代的人物,沒有想到已經被抬到死亡的荒野去了.拉撒路死了,被葬在墳墓裡.在那漆黑陰濕的地方,他全無知覺.很多人落到下流污濁的地步,照樣毫無知覺.一班青年人,甚至一班普通的人,都有這種死的光景.
\newline
\newline
聖經給我們看見苦惱的問題,也給我們看見追求的問題.青年人應往那裡找尋方向呢?有人說:「青年人是迷失的一代!」應當怎樣做才對呢?朝那一方向走才正確呢?許多人迷失了方向,不知何去何從!
\newline
\newline
聖經如何描述青年人之可貴呢?第一描寫青年人的身子俊秀,體態美麗肌肉豐滿,姿容端莊.第二描寫青年人如初升的太陽,光茫萬丈；如水之源頭,年之春初,筍之嫩芽,花之開放,把年青人描寫得有力量.第三描寫青年人是後生可畏,來日方長,滿有朝氣,蓬蓬勃勃,景象萬千.從聖經的描述,可知年青的日子在神的眼光中極其寶貴.錦繡年華的青年時代,有著無瑕無疵的童心,有著漫長的路程.青年人容易教導,也容易帶領,所以成了撒旦與神爭奪的對象.
\newline
\newline
約瑟出來正是年青的時候.大衛出來也在年青的時候.耶穌騎驢進耶路撒冷,所騎的是隻驢駒子.從聖經觀點來看,年青時代在神面前實在寶貴.我們再從以斯帖記看見,末底改是年長者,以斯帖是年少者；末底改把以斯帖帶了出來,表明年長的支持年少的,提攜年青人被神所重用.
\newline
\newline
人類有許多苦惱的問題,實在無從找到答案.學問,科學都無法作答.惟有歸向神,牢牢抓住神.只有神才能解答我們一切的問題.但可惜很多基督徒沒有完全歸向神,只不過一半歸向神；一方面事奉神,一方面卻事奉瑪門.聖經告訴我們:「我們都如羊走迷,各人偏行己路.」(賽五三3)意思就是人離開了神的管治.世界為何會有諸多矛盾,痛苦,紛爭?就是因人離開神的緣故,世界才落到這樣可怕的境地.
\newline
\newline
當所羅門王在位的時候,有兩個妓女到他面前投訴.她們同住一處,各生一個兒子.有一晚夜,一個母親壓死了自己的兒子,卻換了別人的兒子.於是兩婦人爭吵起來,告到所羅門王那裡.神給所羅門王智慧來判斷這事；王吩咐人拿刀來,把活孩子劈成兩半,各得一半.活孩子的母親為自己的孩子心裡急痛,寧可求王把活孩子給那婦人,萬不可殺他!弟兄姊妹!一半的孩子,誰都不要,我們把一半給神,神怎能要呢?我們的嘴說歸向神,心卻遠離神；主日上午歸向神,下午卻遠離神!奉獻一點錢給神,生活卻遠離神；在禮拜堂歸向神,在家裡卻遠離神.
\newline
\newline
世人有痛苦,基督徒也照樣有痛苦,因為我們還沒有完全歸向神.所以人類的問題,無從得到答案.弟兄姊妹!邪情和惡慾是因參雜了邪靈所致的.路加福音記載,有人離開耶路撒冷,在路上被強盜打到半死.為何有這樣的遭遇?因為他離開耶路撒冷,給強盜有機可乘.有許多基督徒,因離開神的緣故,給邪靈有機可乘.故此,不信者有邪情惡慾,信者也有.如果基督徒到了這樣的地步,信和不信有何分別?
\newline
\newline
弟兄姊妹!惟有歸向神,一切問題才能得到解決.但神要我們完全歸向祂.求主今晚給我們實在的問題,也求神鑒察我們的心,我們歸向祂一半?是整體呢?求主光照我們,使我們在祂面前立下心志,完全歸向祂.
\newline
\newline


\newpage

\section{第45屆~港九培靈會~吳勇~第五講}
\label{sec:1268}
\textbf{基督再來}
\newline
\newline
連結: \href{https://www.hkbibleconference.org/session-message/view/1268}{\texttt{https://www.hkbibleconference.org/session-message/view/1268}}
\newline
\newline
\hyperref[sec:1266]{< < < PREV SERMON < < <}
~
\hyperref[sec:toc]{[返目錄]}
~
\hyperref[sec:1270]{> > > NEXT SERMON > > >}
\newline
\newline
但以理書 2:31-2:35
\newline
\begin{longtable}{cl}
\hline
\hline
章節 & 經文 (和合本修訂版)\\
\hline
2:31 & \begin{tabularx}{0.7\textwidth}{X} 「你，王啊，你正觀看，看哪，有一個很大的像，這像甚高，極其光耀，立在你面前，形狀非常可怕。 \end{tabularx} \\ \\ \relax
2:32 & \begin{tabularx}{0.7\textwidth}{X} 這像的頭是純金的，胸膛和膀臂是銀的，腹部和腰是銅的， \end{tabularx} \\ \\ \relax
2:33 & \begin{tabularx}{0.7\textwidth}{X} 腿是鐵的，腳是半鐵半泥的。 \end{tabularx} \\ \\ \relax
2:34 & \begin{tabularx}{0.7\textwidth}{X} 你正觀看，見有一塊非人手鑿出來的石頭打在它半鐵半泥的腳上，把腳砸碎； \end{tabularx} \\ \\ \relax
2:35 & \begin{tabularx}{0.7\textwidth}{X} 於是鐵、泥、銅、銀、金都一同砸得粉碎，如夏天禾場上的糠秕，被風吹散，無處可尋。打碎這像的石頭成了一座大山，覆蓋全地。 \end{tabularx} \\ \\
[1ex]
\hline
\hline
\end{longtable}
(但二31-35)看見一個又高又大的像.(太廿四28):「屍首在那裡,鷹也必聚在那裡.」
\newline
\newline
(太廿五1-4)題到十個童女,五個是聰明的,五個是愚拙的.(約廿一15):耶穌三次問彼得「你愛我比這些更深嗎?」(提後三1-5)；「末世必有危險的日子來到.」人「有敬虔的外貌,卻背了敬虔的實意.」(彼前四17)告訴我們「審判要從神的家起首.」(彼後三10):「主的日子要像賊來到一樣.」
\newline
\newline
耶穌基督再來的預兆:(一)可從但以理書所解釋的巨像看.(但二31-35)(二)可從挪亞的日子和羅得的日子看.(路十七26-29)(三)可從有敬虔的外貌,無敬虔的實意來看.(提後三1,5)
\newline
\newline
有一次,耶穌和門徒在一起,門徒問耶穌關於祂再來有何預兆?耶穌說到外邦的日期,這是指著但以理二章說的.去年十月我在(真理堂)曾講過此段聖經,是以政治觀點來講的.現在卻要從整個世界的景況來看.整個巨像的構造是金,銀,銅,鐵,和泥,上重下輕勢趨敗壞.
\newline
\newline
第一:我們看到世界的人越來越壞:
\newline
\newline
人的身體日趨敗壞.香煙尼古丁,酒精,鴉片,海洛英,梅毒細菌,(尤以國際梅毒細菌尤甚)各種毒素的侵襲.
\newline
\newline
人的思想越趨敗壞.有人主張人文主義思想,主張雙手萬能,人可以勝天.有人主張存在主義,積極的存在主義者,不信有將來,只信有現在；把全部的精力投資於「現在」.消極的存在主義者,當然不承認有將來,他只看見人生的短短幾十年；也不承認有現在對於一切都看作無所謂.人文主義也罷,存在主義也罷,攏統言之,他們是無神主義.
\newline
\newline
人的心越趨敗壞.人心包括情感和意志.如今的人,情感日益疏遠,冷酷,人的意志本應以良心和道德為標準；但如今卻不然,乃以利害為標準.
\newline
\newline
環境日益敗壞.住在香港的人,因為治安日差,失去安全感,時時提心吊膽；加以物價飛漲,生活艱難,實令人片刻難安!居住環境,空氣污染,連聲音也受污染.這些會關係到別人的情緒,因而影響人的心臟神經.
\newline
\newline
第二:世界情勢日趨敗壞　東南半島問題至今未能解決,仍在戰爭中.至於中東問題,以色列和阿拉伯仍世代為仇,衝突不已.
\newline
\newline
第三:整個世界面臨經濟崩潰地步.幣制不斷貶值,令人非常憂懼.世界光景日趨惡化.第一次世界大戰有八百萬人陣亡.第二次大戰四千萬人陣亡,一旦再發生世界大戰,其結果實不堪設想.總之,世界大勢日趨敗壞,人越來越壞,環境越來越壞,前途越來壞.
\newline
\newline
我們再看那大像最後的一段是鐵和泥雜在一起.世界各國的政治鐵泥不分,混在一起.
\newline
\newline
照著中國的封建思想來說,男女是授受不親；可是現代新潮的男女,無論是衣著,是頭髮,已弄至男女不分了!本來神創造萬物,是各從其類的,如今卻變成不倫不類了.
\newline
\newline
挪亞的日子　「挪亞的日子怎樣,人子的日子也要怎樣.那時的人又喫又喝,又娶又嫁.」(路十七26-27)今天的人的光景,又喫又喝,生的時候講究吃,死了也照樣吃.開張要吃,結束也照樣吃,無所不吃!「又」是再的意思,一個有身份的人,吃了又吃；甚至每晚三場五場應接不暇地吃.所以血壓高和糖尿症,已成了極普遍的都市病.「又娶又嫁」就是娶了再娶,嫁了再嫁；淫亂成了時尚!報紙,雜誌無一不渲染淫亂.許多女人的服裝,後面是空的,連前面也是空的,更有全面是空的.從聖經立場來看,拜偶像和淫亂,都是神最憎惡的.淫亂的風氣已到人不能再忍受的地步,神又怎能忍受呢?
\newline
\newline
羅得的日子　「又買又賣」(路十七28)關於這點,最好以股票來解釋.去年香港股票的買賣.到了瘋狂的地步.股票在正常的情況下來投資那是對的；可使社會事業發展,經濟流通,大家分享利潤.但如果變成了投機的買賣,就能造成三種創傷1.人心的創傷.人只注目在股票上,股票起落,心也隨之上落.2.神經創傷,股票的漲跌,影響到人的神經隨著緊張.3.性命創傷.我認識一個人,為了股票暴跌,從八樓跳下,可幸跌下時有帳棚擋著,命雖不絕,身體卻已受傷.羅得的日子「又耕種,又蓋造.」耕種可表農村社會,蓋造可表城市社會.現在農村的人,都湧進都市,所以田園無人耕種；農村破產,世界飢荒開始,情形日益嚴重.城市方面,因從農村湧入的人日增,所以人口擠迫,問題叢生,紛爭愈多,罪惡愈盛.人為了要滿足肉體需求,以至整日忙忙碌碌,又耕種,又蓋造.聖經作者,被聖靈啟示,寫出今日世界和社會的光景.
\newline
\newline
口是心非　「有敬虔的外貌,卻背了敬虔的實意」(提後三5)這是今日教會的光景.耶穌在馬太廿三章責備法利賽人,說他們能說不能行.「說」是言語,基督徒應該說屬靈的話語:但屬靈兩字,並非一種標語.文字是有內容的.有愛心,謙卑,不顧自己,這是屬靈的內容,非單用言語來說,而且是在生活上去履行.法利賽人所作的,是叫人看見,而裡面卻是空洞的.沒有實際的內容,只憑外表虛假的裝飾.他們在十字街口揚聲禱告,把經文繡上衣袖,裝作屬靈的模樣顯於人前,這是今日教會的光景.教會應如寶貝放在瓦器裡,然後才能發出莫大的能力來.如果瓦器裡面是空洞的,那就無可顯明了!這樣教會就會落到無能的地步.
\newline
\newline
我在十歲時,曾在學校參加演講比賽,講題和次序都是臨賽前抽籤決定的；豈料我抽的是第一號,所以最先登台.幸而開場白老師早已教過,因此,我能很熟練地說出:「校長!各位評判老師!各位同學!今天小弟代表本班來參加演講比賽......」說到這裡便停止的沒有下文了.我呆站在台上,發現台下的人,個個睜著大眼睛看我,急得我手足無措.下意識的用手一直拉自己的衣服,引至同學們哈哈大笑!而我就忍不住哭起來了.老師急忙把我拉了下台,那次我實在太不爭氣了!
\newline
\newline
有兩個青年男女,同在鄉村長大,一同上學,一同玩耍.他倆的感情很好,到好了大學畢業當論婚嫁的時候,因男的是個內向的人,他內裡的感情,無法用言語表達.有一天,他們一同出去遊玩,直到了一棵大樹下；二人心照不宣,男的理應先開口向女的求婚,但他沒有勇氣.那時,男的站在樹的一邊,女的坐在的樹的另一邊,二人黯然不語.後來女的抬頭面向男的望了一望,示意催促他開口；那裡想到男的一時慌張,下意識的舉動出來了,把大拇指放在胸口上,不停地轉動.女的看了,非常生氣,走到他跟前,頓足疾言說:「難道你只會這樣嗎?」男的表示說:「我不止會這樣,我還會從後面繞過來.」今日基督的教會,神的兒女,除了會這樣,會那樣,到底我們還會做什麼呢?除了會禮拜,會捐幾個錢,還會作什麼呢?
\newline
\newline
從世界的光景,社會的光景,和教會的光景這些預表,我們知道末世到了,耶穌快要再來了.
\newline
\newline
「主的日子要像賊來到一樣」(彼後三10)為何用賊作為比方呢?我曾講過,神常用人頭腦所能明白的來解釋不能明白的事.屬靈的事是人所不易理解的,耶穌自己說祂來要像賊一樣.耶穌再臨是來到空中,不是來地上.那時人看不見祂,好像賊來光顧,都是趁著人不知道的時候來的.既然看不見賊來,怎信賊來過呢?最少有兩種跡象,第一,非常凌亂,因為賊來過必翻箱倒履.第二,貴重之物失掉,因被賊偷走了.耶穌再來時,地上有苦難；還有,馬太廿四章說,那時取去一個撇下一個.什麼被取去?就是那些貴重的.那時,到處人找人,有的丈夫不見了；因很多人已被提到空中了.地上剩餘的,一種是猶太人,一種是不信的外邦人,還有一種是教會中沒有生命的信徒.為何猶太人會被撇下,不足為奇,因他們直到如今還不信耶穌是彌賽亞,不承認耶穌是神的兒子.外邦人被撇下,也不希奇；因為他們的心向著物質,只顧現在,不顧將來.但奇怪的是教會中竟然也有人被撇下!
\newline
\newline
「屍首在那裡,鷹也必聚在那裡.」(太廿四28)鷹代表刑罰和咒詛.屍首是無生命的,那裡沒有生命,那裡必有刑罰和咒詛.第三種的教會中人,可以說他們是沒有生命的人.猶大書說:是「沒有靈的人」.啟示錄說:「按名是活的,其實是死的.」人沒有生命,就沒有感覺,就沒有功用.教會中有許多沒有感覺的人,聽了道毫無反應；警戒的話他也不懼怕,對他講神的愛,他更無動於中.在教會中也可以見到許多沒有功用的人,來到教會,好像死屍擺在那裡；沒有動作,沒有反應,教會中多了他或少了他,都起不了作用.
\newline
\newline
感謝主的恩典!如果我們是有生命的人,當主再來時,祂就要接我們到祂那裡去.
\newline
\newline
審判要從神的家起首,將來被提的,是一切屬神的人.工作勤勞,靈命豐盛的被提.靈命軟弱,貧窮的人也照樣被提.或是有人要說「神的公義在哪裡?」聖經說:「審判要從神的家起首.」在空中基督的審判台前,祂要審判被提的人.馬太廿四章題到被提的真理,廿五章題到十個童女的真理.十個童女,五個是聰明的,五個是愚拙的.他們穿著一樣的衣服,手中都提著點著的燈；但其中最大的分別,就是那五個聰明的童女,除了一隻手拿燈之外,另一隻手還拿裝著油的器皿.這預表聰明的人,不單有生命,且有更豐盛的生命.愚拙的童女只有點亮的燈,器皿卻沒有裝著油,這表明地上有很多人,雖有生命,僅是幼小的生命,非但不能發出功用,且成為人的累贅.教會有些人,要向他一趟趟的請才來,一次次地勸才舉手.我想將來在審判台前,受虧損的,是那些一隻手沒有拿裝著油器皿的人,深信神是公義的.
\newline
\newline
更豐盛的生命是付代價追求得來的,付代價仰望神才有的.那麼,到底這種力量從哪里來呢?約翰廿一章告訴我們,耶穌三次問彼得:「你愛我,比這些更深嗎?」彼得所回答的,卻錯解了耶穌所問那愛的含義.耶穌所指是(約翰三16)的愛,這樣的愛是絕對的,是要犧牲的.耶穌問彼得:「你願意犧牲；你肯絕對的愛我嗎?」彼得用感情的愛來答覆耶穌.耶穌說:「你愛我比這些更深嗎?」「這些」我想是指:
\newline
\newline
餅和魚(代表物質)
\newline
\newline
炭火(代表舒服因為在猶太地方,清早很冷,有炭火覺得溫暖舒適.)
\newline
\newline
同伴.
\newline
\newline
耶穌問彼得的話,意思是說:「你愛我比餅和魚更深嗎?比炭火溫暖舒適的享受更深嗎?比你的同伴更深嗎?」我們靈性上常常存著這幾個問題.我們不能勝過物質的誘惑,不能勝過自己的享受；我們在孤單零丁的路上就不能走.我們看別人熱心才熱心,看別人服事才服事.在屬靈的路上需要有同伴才能向前走.人走屬靈的道路,必需有神的愛,才能勝過「這些」.耶穌所以這樣問彼得,因祂要把責任交付彼得,必須愛祂的才能受祂的託付.耶穌在十字架上時,把母親交給約翰,因祂知道約翰對祂的愛心如何.我們再從彼得身上來看,他給耶穌的愛,雖然不是耶穌所要的,但仍把責任交給他.五餅二魚在五千人中,本是微不足道的,無濟於事的食物,但交在耶穌手中,祂仍然接受.好像小孩子吃雞腿,媽媽問他要,他就撕下一點點送到媽媽咀裡.雖然只是全隻雞腿的微小部分,但媽媽也快樂地接受了.彼得給耶穌只是情感的愛,不是耶穌所渴望的,但祂也照樣接受了,所以才把責任交給彼得.許多神的兒女,莫說沒有神聖的愛；沒有犧牲的愛,沒有絕對的愛；他和神之間,連情感的愛也沒有.如果夫婦之間,沒有情感,這家庭還有盼望嗎?許多神的兒女與耶穌之間,連情感的愛都沒有；難怪對於神的話無動於中,毫無反應!靈性陷於全完停頓的地步,一直不長進,一直得不到復興；乃因他和神的情感斷絕了,以致造成現在教會的光景.假如夫妻之間問題多多,如果還有感情存在,則家庭尚有盼望.照樣,如果我們對神的情感還在,我們屬靈的盼望仍存在.可惜神的兒子與神的感情,已被世界完全霸佔了,剝奪淨盡；每逢來到神面前,毫無反應,空空而來,也空空而去!昔日如此,今日亦然!神的道在他心裡再不能生效了!
\newline
\newline
耶穌要再來,我們當往前走!
\newline
\newline
求主鑒察我們的心,是否與主失落了情感?一旦落到這樣的光景,是很可怕的.求主再次把十字架的大愛,刻劃在我們心中.感動我們,光照我們罷!
\newline
\newline


\newpage

\section{第45屆~港九培靈會~吳勇~第六講}
\label{sec:1270}
\textbf{被召者多選上者少}
\newline
\newline
連結: \href{https://www.hkbibleconference.org/session-message/view/1270}{\texttt{https://www.hkbibleconference.org/session-message/view/1270}}
\newline
\newline
\hyperref[sec:1268]{< < < PREV SERMON < < <}
~
\hyperref[sec:toc]{[返目錄]}
~
\hyperref[sec:1272]{> > > NEXT SERMON > > >}
\newline
\newline
民數記 14:20-14:24
\newline
\begin{longtable}{cl}
\hline
\hline
章節 & 經文 (和合本修訂版)\\
\hline
14:20 & \begin{tabularx}{0.7\textwidth}{X} 耶和華說：「我照著你的話赦免他們。 \end{tabularx} \\ \\ \relax
14:21 & \begin{tabularx}{0.7\textwidth}{X} 然而，我指著我的永生與遍地充滿了耶和華的榮耀起誓： \end{tabularx} \\ \\ \relax
14:22 & \begin{tabularx}{0.7\textwidth}{X} 這些人雖然都看過我的榮耀和我在埃及與曠野所行的神蹟，仍然這十次試探我，不聽從我的話， \end{tabularx} \\ \\ \relax
14:23 & \begin{tabularx}{0.7\textwidth}{X} 他們絕不能看見我向他們祖宗所起誓應許之地；凡藐視我的，一個也不得看見。 \end{tabularx} \\ \\ \relax
14:24 & \begin{tabularx}{0.7\textwidth}{X} 惟獨我的僕人迦勒，因他另有一個心志，專心跟從我，我要領他進入他所去過的那地；他的後裔必得那地為業。 \end{tabularx} \\ \\
[1ex]
\hline
\hline
\end{longtable}
(民十四20-24)我們看到約書亞和迦勒二人專心跟從神.(士七1-8)我們看到基甸和跟隨他的三百人.以斯帖記我們看到以斯帖.(耶十二7-8)題以色列百姓如同林中吼叫的獅子.(太廿二14)題到召的人很多,選上的人很少.如果某公司要請兩個職員,當廣告登出後,應徵的人,可能是幾個,以至幾十個,但錄取的僅有兩個.耶穌基督在世時,跟隨祂的人很多,但選上的僅十二人,其中再被選上的僅三人.
\newline
\newline
民十四章載,在以色列五六十萬的群眾中,被揀選的只有約書亞和迦勒.士師記記載,當時跟隨基甸的,有三萬二千人.神對基甸說,跟隨的人太多.所以基甸向他們宣告,膽怯的,可以回去；於是有三萬二千人回去,只剩下一萬,神說還是太多.到了最後,只留下三百人.以斯帖記記載,從一百廿七省無數的女子中,只選上以斯帖一人.由這三處聖經,我們思想到兩件事:
\newline
\newline
第一,揀選的需要
\newline
\newline
繼承前任　當時神揀選約書亞和迦勒,是因為摩西死了,需要有人繼承；否則以色列人只有停留在曠野,不能進入迦南了.所以,揀選這兩個後繼之人,是要完成摩西未完的工作.「死」有兩方面的解釋,一是身體的死；好像有很多主的僕人作工到了一個地步,體力漸衰,腦力也沒有了,在這樣情形之下；必須有人繼承,否則工作無法繼續.一是靈的死；例如以利當士師時,耶和華已不常有默示,意思是他與神的交通斷絕了；神對這位先知不能再有啟示了,所以神就揀選撒母耳來繼承以利.許多神的僕人工作到了一個地步,或因世界的緣故,或因罪惡的緣故,以致和神的交通斷絕；神對他再沒有啟示,如果後繼無人,神的工作就無法往前.
\newline
\newline
扭轉時代　窺探迦南地共有十二個探子,其中十個報告凶惡的信息,說迦南地的城如何高大,亞納族人何等強壯,我們站在他們面前,就如同蚱蜢一樣.以色列人聽了,便大喊大哭,以為摩西帶他們去送死；因此,他們想另立領袖帶領他們回埃及,這真是個倒轉歷史的時代!但約書亞和迦勒,肯站起來替神說話:「有耶和華同在,不要怕他們,他們不過是我們的食物.」(書十四9)因此當時的光景就扭轉過來.
\newline
\newline
現今的時代是個千變萬化的時代,從音樂和美術來看；古典樂已變為爵士樂,繼而變為搖擺樂,再變為披頭樂.現在最時尚的音樂,是將內裡的不安,痛苦的情緒,表達出來,所以這種音樂到處風行.現代的美術,如果你問:「所畫的是什麼?」答:「以視覺所感受的而命名.」意思乃是你看它像什麼它就是什麼.這豈非莫明其妙嗎?他說,就是莫明其妙.時代是莫明其妙,美術也莫明其妙.唉!這個時代,真是痛苦,不安,莫明其妙的時代!
\newline
\newline
至於教會,不熱心的人,對教會沒有責任感,以為來禮拜,就算盡了本分.就算熱心的人,也只不過捐點錢給教會,並沒有真誠的事奉.以致教會活不出榮美的見證,叫人羡慕和佩服,形成一片混亂和冷落!
\newline
\newline
如此社會!如此教會!神非要揀選人來扭轉不可.當時的以色列人,若不是有約書亞和迦勒起來,他們就把歷史倒轉了；因有約書亞和迦勒起來,才把情形及時扭轉.所以從民數記給我們看見,神揀選人的目的,除了繼承死的人之外,還要扭轉時代.
\newline
\newline
再看士師記,當時米甸人起來要攻擊神的百姓；剝奪他們的鐵器和糧,使他們膽怯軟弱無力交戰.
\newline
\newline
神的兒女,亦常處與打仗的狀態.傳福音是打仗,要把人從撒旦權下救出來.造就有兩個目的,一方面對付外面的世界,一方面對付裡面的肉體；神的兒女要對付內外,這就是打仗,與別人與自己都要打仗.魔鬼所用的方法,就是剝奪你的武器和糧食,使你沒有武器,使你身體軟弱,不能打仗；牠在人身上,從前使用這方法,如今亦然.撒旦對付傳道人的方法,一是關閉他的天,使他得不到天上來的啟示和感動,以致無糧可分,無信息可傳.一是破壞胃口,使他對屬靈糧食無法吸收.
\newline
\newline
基甸起來,要帶領以色列人去打仗.這場戰爭,並非是一人打的.好像摩西帶著以色列全軍,還有戶珥,亞倫,攻擊亞瑪力人一樣.基甸起來呼召,不是一人跟隨,乃是三百人跟隨他.今日的戰爭,既是猛烈之戰,需要很多人同心合力來爭戰.
\newline
\newline
以斯帖時代,以色列民在哈曼權下.有一天「哈曼對亞哈隨魯王說,有一種民族,散居在王各省的民中......與王無益；王若以為美,請下旨滅絕他們；我就捐一萬他連得銀子......於是王從自己手中摘下戒子給......哈曼......說,這銀子仍賜給你,這民也交給你,你可以隨意待他們.」(斯三8-11)哈曼就利用他的權力要把以色列民全然剪除,戮殺滅絕,並奪他們的財產為掠物.
\newline
\newline
今天神的兒女,也是落在撒旦的權下.(耶十二7)說:「將我心裡所親愛的,交在仇敵的手中.」或者我們不明白,為何心裡所親愛的,還忍心交在仇敵手中,但第八節解釋:「我的產業向我如林中的獅子,牠發聲攻擊我,因此我恨惡他.」聖經以獅子和羔羊描述耶穌:「猶大支派的獅子」「像羊在剪毛的人手下無聲.」耶穌在世時,凡指著祂所說的,也是指著我們說的,因祂是我們的標本,是我們的榜樣.祂既像獅子,羔羊,我們也應如此.獅子是向著撒旦而言,羊羔是向著神而言.耶穌把自己當作羔羊,軟弱交在神的手中,向著撒旦應當像獅子,向著神應當像羊羔.但今天神的兒女,適得其反,向著神像林中的獅子倔強,與神抵抗；向著撒旦像羊羔,軟弱地順從牠,被牠拉住我們的鼻子跟牠走.難怪神將祂心裡所親愛的,交在仇敵手中.昔日神把以色列民交在撒旦手中,今日也把祂的兒女交在撒旦手中,今昔相同.
\newline
\newline
神興起以色列,要把在撒旦權下的人呼喊出來,拯救出來.照樣,神給祂的兒女打擊,磨煉,目的要叫他們的心軟化,神興起人來向他們呼喊,叫他們從撒旦權下出來.從民數記看,揀選的目的為了繼往開來,為了扭轉時代.從士師記看,因有一場猛烈的戰爭,需要很多人參與.從以斯帖看,有許多人在撒旦的權下,神揀選人把他們呼喊出來,這些都是神揀選的需要.
\newline
\newline
第二,揀選的條件　民數記載,約書亞和迦勒專心跟從神.專心就是心不受任何事物所影響:
\newline
\newline
世界不能叫他變心　新約的底馬,因貪愛世界,離開保羅到帖撒羅尼迦去了,因世界叫他變了心.四福音題到猶大,用卅塊銀把耶穌賣了,世界叫他變了心.亞伯拉罕與羅得從米所波大米出來同走到迦南地,但他們的結局卻不同.神對亞伯拉罕說:「論福,我要賜大福給你,叫你的子孫要像天上星,和海邊的沙.」(創廿二17).至於羅得,結果家散人亡,身敗名裂,而且犯了亂倫之罪,留下孽種的後代.他何竟落到這地步呢?乃因他貪愛平原,把棚帳漸挪移靠近所多瑪,世界使他變了心.許多神的兒女怨像羅得一樣,世界使他變了心.
\newline
\newline
遇難不灰心　馬可本與保羅同工,出門傳道,可惜他半途而廢離開了保羅,因他遇到難處就灰心.許多人在工作開始時,非常熱烈,遇到難處,就心灰意冷!在神的工場上,我們面對仇敵撒旦,牠會製造種種難處,擺佈障礙,使我們心灰意冷.以致許多人變了質,半途而廢,這是因為不專心的緣故.
\newline
\newline
受打擊不灰心　我們看大衛王,當時掃羅妒忌他,到處追尋他的命.他如喪家之犬,四處逃亡,幾乎喪命；甚至國內無他立足之地.他便向國外尋找安身之處,遂投奔亞吉王,孰料亞吉王的臣僕妒忌他,不能容納他.雖然亞吉慕大衛之才想容納他；但為了要安撫自己的臣僕,只好忍痛辭退大衛.大衛無可奈何,又折回洗革拉,看見家園被毀,家眷失散,遭亞瑪力人洗劫.這時候,跟隨大衛的人,怨聲四起,因為他們跟隨大衛多年,至此一切皆空.甚至竟無立錐之地,激憤之餘,要用石頭打死大衛；大衛處此境地,所受打擊何其大!國內國外的人既不能容他,連跟隨的人,不但不諒解他的遭遇,反歸咎於他.聖經告訴我們:「大衛倚靠耶和華他的神心裡堅固.」(撒上卅6下半)
\newline
\newline
能專心者,即使世界也不能叫他變心,艱難也不能叫他灰心,打擊更不能叫他灰心.當時這些人被揀選,神為他們見證說:「因為他們專心跟從我.」在這世代中,專心的人才有用處.我們常看見有些人,一面要跟隨基督,一面手扶著犁頭向後著.神所要的,是專心的人；因為專心的人才能在任何境況中,始終屹立不動.
\newline
\newline
我們再從士師記來看揀選的條件.當時有三萬二千人跟隨基甸,後來膽怯,一批批退去.神指示基甸,用喝水的方法來試驗,最後留下三百人,這三百人用舌餂水,表明儆醒,是合格能用的.其餘有九千七百人,跪下捧水喝的,這是代表靈性鬆懈的,神不能錄用他們.因為他們跪下時,要把戰衣解開,把武器放下,才能方便下跪俯首喝水.
\newline
\newline
大衛做王時,約押帶兵上前線打仗,大衛在後方睡到日頭平西才起來,表明他靈性鬆懈；因而魔鬼有隙可乘,使大衛跌倒至嚴重的地步.彼得告訴我們:「魔鬼好像吼叫的獅子,遍地遊行,尋找可吞吃的人.」靈性鬆懈的人就被吞吃了.我們處於魔鬼極度活動的時代.在舊約時,撒旦是條蛇；在新約時,撒旦是條龍.蛇固然厲害,龍更厲害；可見撒旦作工,變本加厲.神的兒女,一旦鬆懈,撒旦便有機可乘來攻擊了.
\newline
\newline
以斯帖具備了被揀選的條件:
\newline
\newline
尊重長者　她沒有把自己的身份告人,這是末底改囑咐她的,可見這位年青人尊重年長者.許多家庭兩代之間,常有歧見,造成很大的難處.年輕的以為年長者固執,自高,持守傳統,誠然有很多缺點；但我們不能因噎廢食,歧視年長的；因年長者也有很多長處,很多經驗.不過他們因持守傳統,墨守成規；所以與年青人與他們成了兩代鴻溝,這邊過不去,那邊過不來.年長者雖有錯誤,但他們卻有豐富的經驗,可作為年青人的借鏡,可作前車之鑑,以免重蹈覆轍.如果年青人把年長好處抹煞；那麼,他們必會遭受許多無謂的虧損.
\newline
\newline
隱藏自己　以斯帖沒有把身份告人,表明她隱藏自己.我們事奉神,不在於工作如何,乃在於動機如何.主耶穌說:「一杯涼水也不能不得賞賜.」保羅說:「捨己身叫人焚燒,卻沒有愛,也算不得什麼.」難道人的身體還不及一杯涼水嗎?這是因為給人一杯涼水的動機是對的,而捨己身叫人焚燒的動機是不對的.我們服事神,動機若是為著自己,就算犧牲身體被燒,在神的眼目中,也沒有價值.如果小小一杯涼水,動機是為著神；人雖看為微少,但神卻看為可貴.從前的法利賽人,並非他們的教規教理不對,而是他們的動機不對,為要炫耀自己.
\newline
\newline
別無所求　以斯帖除了希該所派給她的分,她別無所求.大衛王逃避押沙龍時,米非波設因為瘸腿不利於行,所以留在耶路撒冷.米非波設的僕人洗巴在大衛面前說讒言,說他主人幸災樂禍,忘恩負義.大衛聽了,很是傷心,因為大衛為報答約拿單,所以待他的兒子米非波設如同自己的兒子.大衛聽了洗巴的讒言,信以為真,要把米非波設的財產,分一半給洗巴.到了大衛回耶路撒冷.米非波設下去迎接王的時候,王問他說,為何沒有同王去.經米非波設解釋,誤會才冰消.大衛提到他和洗巴均分地土的事,他說:「王既平平安安的回宮,就任憑洗巴都取去了也罷!」米非波設別無所求,惟獨要王.這樣的人是多麼可貴,他的心多麼單純.如果我們向主的心,不求名,不求利,甚至不求人的稱讚和看重,除基督以外,別無所求；這樣的心,在神面前何等可貴!但神的兒女如此者有幾?
\newline
\newline
甘願犧牲　以斯帖說:「我違例進去見王,我若死就死罷.」她為了全民族,個人的得失,在所不計.在哈曼權勢之下,作了驚天動地之大事.
\newline
\newline
軟弱變剛強　以斯帖也有她軟弱的時候,當末底改吩咐她去見王時,她說:「若不蒙召,擅入內院見王的,無論男女必被治死；除非王向他伸出金杖,不得存活.現在我沒有蒙召進去見王,已經三十日了.」(斯四11)她雖有軟弱,神還是用她,因她能及時從軟弱中剛強.人難免有時軟弱,只要肯從軟弱中起來,他仍有用.許多神的兒女,因為在軟弱中不能起來,以致失去他的功用.
\newline
\newline
求主給我們看見,被召的人多,選上的人少,被神選上是要具備條件的.約書亞和迦勒,專心跟從主.跟隨基甸最後留下的三百人,他們時刻儆醒.不敢鬆懈的以斯帖選上,她具備了很多的條件,其中最寶貴的,就是她能從軟弱中起來,所以才成為那時代被神重用的人.
\newline
\newline


\newpage

\section{第45屆~港九培靈會~吳勇~第七講}
\label{sec:1272}
\textbf{建造聖城耶路撒冷}
\newline
\newline
連結: \href{https://www.hkbibleconference.org/session-message/view/1272}{\texttt{https://www.hkbibleconference.org/session-message/view/1272}}
\newline
\newline
\hyperref[sec:1270]{< < < PREV SERMON < < <}
~
\hyperref[sec:toc]{[返目錄]}
~
\hyperref[sec:1273]{> > > NEXT SERMON > > >}
\newline
\newline
尼希米呼喊猶大人起來建造聖城(尼二17)第三章首段有六次論到其次某某人修造.第七章題到尼希米派人去看守城牆.第八章以斯拉講聖經.(啟廿一)論到城內的光景.廿二章論到城外的光景.從聖經目錄的次序看,約伯記在列王記和歷代志的下面,其實,約伯記是聖經較早的一卷書.尼希米記在聖經目錄上雖早於大小先知書,其實,尼希米記應排列於瑪拉基書之後,因為是同一時代.
\newline
\newline
聖經題到被擄歸回的猶大人有兩種建造.其一,是以斯拉記論建造聖殿.其二,是尼希米記論建造聖城.
\newline
\newline
神的榮耀在聖殿裡,神藉著聖殿對以色列民傳話,以色列民藉著聖殿向神敬拜.聖城是以色列民生活的所在,他們都住在城裡.聖城預表人的生活,聖殿是預表人的靈命.聖城代表外面的生活,也可以說是見證.我們有個很要緊的責任,就是把聖城建造起來,這聖城就是表明神的見證.我們每個人都有重大的責任,就是把耶穌基督,充充滿滿地表達給世人看.
\newline
\newline
聖城可作為疆界和保護之用.聖城內有精金的街道,碧玉的城牆,珍珠的城門.(啟廿一)城外有那些犬類,行邪術的,淫亂的,殺人的,拜偶像的,說謊的.可見城外城內,大有差別.內外之間,有城牆為界,以防城外那些罪惡混入.聖經常吩咐我們,要從他們中間出來,要與他們有分別.
\newline
\newline
一,人生目的有別　世人所注重的是地上的福氣,他們求名求利,生活患得患失.至於我們應以天上的福為重,與神交通,與肢體配搭,為神為人.因此,我們的人生是平安喜樂的.
\newline
\newline
二,生活樣式有別　世人在生活中,常以玩弄手法,玩弄手段為能事,人與人之間,只有虛假,沒有真實.基督徒應以聖經為準則,以愛己之心愛人,學習溫柔謙卑.處在這謊謬的世代,就算吃虧,也不能違背聖經的原則.這是基督徒應有的樣式,與不信者,應有清楚的界限.神說:「不可把兩樣的種子,種在你的葡萄園裡.不可并用牛驢耕地.不可穿羊毛細麻兩樣攙雜的衣服.」(參申廿二9-11).為何神這樣吩咐?因祂憎惡混雜.一旦混雜,見證就無效,價值就不存在,能力就失去了.
\newline
\newline
使徒行傳五章告訴我們,五旬節聖靈大大動工,教會非常興旺.當時亞拿尼亞夫婦,賣了田產,把價銀私自留下幾分,其餘的幾分,拿來放在使徒腳前.彼得指責他們說:「你不是欺哄人,是欺哄聖靈.」因此他們倒斃了!賣田產是表明犧牲,付代價；屬靈愛主,並非空洞的名詞,是要犧牲和付代價的.亞拿尼亞夫婦,他們羡慕屬靈的名聲而不肯犧牲,不願付上代價,他們運用說謊來朦蔽,希望達到其目的,用謊言謀取榮耀.這樣的事混入教會,教會就被破壞,虛假的榮耀進入教會,神的教會就被混雜,見證就失去,所以城牆應有界限,城內城外應劃分清楚,應小心防社會的東西混入教會.四福音載:主持聖殿的人,藉屬靈的事謀私人的利益,在那裡兌換銀錢,賣牛羊鴿子,耶穌來,把他們都趕走了.世人常以利為前提,社會上有錢的人,被稱為賢達,被推為首領.有很多教會,也效法世界!以為有錢的人才能作教會領袖這真是讓世界混入了教會.有很多教會以錢作為長老的資格,因他們認為經濟能解決一切的難處,有錢者在教會中,才能雄壯教會的聲勢.世人以手段作為對付人的方法.社會以奉承博得人的喜歡,以拉攏樹立個人的勢力.一般教會也效法這種樣式,遇事靠自己的聰明大家商量對策.如果教會也採用此法,沒有一同禱告倚靠神,落到如此地步,何等可憐!
\newline
\newline
尼希米見當時耶路撒冷城牆被拆毀,城門被焚燒,心裡難過.有許多教會,豈非如此,教會與社會已經沒有分治的界限了.尼希米的任務是要把城牆重建,設立看守的人,教會如同聖城,也應設立看守的人；因為城內的街道是精金造的,城牆碧玉造的,城門珍珠造的；精金代表價值,與榮耀,碧玉和寶石代表光輝.珍珠代表珍貴,精金成分純一；寶石和珍珠沒有瑕疵,不受城外外之物混污.照樣,教會也應純潔無瑕疵,無外界的污染,才能成榮耀光輝尊貴的教會.神造亞當之後,派他在伊甸園看守,這是意味到有些東西要混入.尼希米建造城牆之後,也派人看守,這也是意味到有仇敵要從外混入,城內的性質將受影響.神看純一不混是件非常嚴肅的事,祂不容許亞拿尼亞和撒非喇,把虛假的榮耀混入,也不容那些祭司以謀財利己為目的混進來.
\newline
\newline
彼得到安提阿和外邦人一同吃吃飯,後來,看見從耶路撒冷來的猶太人因懼怕這些奉割禮的人說話,就立即躲開；保羅看見,認為這事是不對的,把律法觀念混入神的救恩裡面；救恩的性質,就會改變.保羅認為事態嚴重,所以當面責備他,抵擋他,免得救恩的真理,因彼得的錯誤而變質.要看守城牆,防止仇敵侵入,免得影響純金的城牆變質,城牆是界限,城內才是屬於基督.不讓城外之物混入城內,這是教會的責任.城外是甚麼呢?是地上的政治,肉體的情感,人的天然性情,這些都不容混入教會；一旦混入教會,性質即變,教會價值即不存在.
\newline
\newline
尼希米建造城牆,是指示我們應負的責任.城牆是界限,界限代表分別.每一個神的兒女,都有看守城牆的重大使命,不容許世界的東西進入城牆裡面.
\newline
\newline
城牆另一個意思是保護:
\newline
\newline
一,保護神子民的利益　我們還未被提之前,外面有世界,裡面有肉體；直到了被提之日,身體改變,肉體即不存在；被提之日,身體離開地上,所以世界也不存在了.我們在地上之日,外面的世界和裡面的肉體,不斷地纏纏我們.
\newline
\newline
世界叫神的子民受虧損　從舊約聖經,我們看到以利沙的僕人基哈西,他的主人醫好乃縵的大痳瘋；乃縵備了許多禮物,向以利沙表示大大的感激,但以利沙拒收.他僕人在旁觀看,覺得這許多禮物主人不收,實在太可惜；所以暗地去追乃縵,騙取一他連得銀子和兩套衣服回來.以利沙問基哈西從那裡來,他還撒謊說:「僕人沒有往那裡去.」以利沙對他說:「乃縵的大痳瘋必沾染你,和你的後裔,直到永遠.」世界叫基哈西受了極大的虧損.新約的猶大,跟隨耶穌三年半,最後用卅塊錢把耶穌賣了給人釘十字架.猶大後來良心發現,把銀子丟在聖殿,出去吊死了.世界叫猶大受虧損,世界很容易叫人受虧損.許多神的兒女,當世界還未糾纏他的時候,在靈程上邁步前進；豈知偶然失足,世界迷惑了他,叫他受虧損.
\newline
\newline
肉體叫神的子民受虧損　摩西在曠野,百姓因沒有水喝而發怨言,摩西發怒擊打磐石兩下流出水來.結果神對摩西說:「你在以色列民面前不尊我為聖,所以你不得進入迦南.」肉體的脾氣叫摩西受虧損,大衛行淫污辱了拔示巴,還殺死她的丈夫；神要大衛賠四倍,肉體的情慾,叫大衛受虧損.
\newline
\newline
肉體的脾氣,肉體的情慾,虛假的榮耀,都能叫我們受虧損,那麼,用何方法來保護神子民的利益呢?
\newline
\newline
真理的教導　「誰能使我們與基督的愛隔絕呢?難道是患難麼,是逼迫麼,是飢餓麼,是赤身露體麼,是危險麼,是刀劍麼?」(羅八35)什麼都不能叫我們與基督的愛隔絕,可見愛有強烈的保護力量.我們愛主的心不當減少,也不當冷落,才不致失去保護靈性的力量.真理教導我們,「體貼肉體的就是死,體貼聖靈的乃是生命平安.」(羅八6)神的兒女與外邦人不同的,並非因我們有禮拜的儀式,有基督徒的名號,乃是因我們裡面有聖靈.(加五17):「情慾和聖靈相爭」有時我們受人的激動發脾氣,聖靈就指導我們趕緊制止.如果制止不住.聖靈就教導我們要趕快對付；我們若制止情慾,對付情慾,就是體貼聖靈；治死我們心中的惡行,叫我們生命俞加豐盛.裡面的生命豐盛,肉體就漸衰微,在我們身上無能為力了.
\newline
\newline
齊心合力　尼希米三章有六次說:其次某某人修造,這意思就是我們生活在一起,應當彼此照顧,互相配搭,這是及時的保護.從前長大麻瘋的乃縵,以利沙吩咐他到約但河去沐浴,他聽了不滿意,準備回國；後來他的僕人勸他說,先知若吩咐你做一件大事,你豈不作麼?何況說你沐浴得潔淨呢?因僕人及時的勸告,使他的錯誤挽回而得醫治.基督的教會許多人的軟弱要及時挽回,免得軟弱愈甚至於無可挽回.教會中有些人的不滿,若不及時挽回,就會像瘟疫蔓延,越傳越廣了.
\newline
\newline
城牆是一種見證,仇敵隨時要毀壞這個見證.使用變質的方法,叫一個真正的教會變成團體,成了社會活動之所,神的家變成客棧沒有愛和肢體的關係.教會本是禱告的殿,卻變成賊巢,變成謀利的地方；教會是神的居所,地變成魔鬼的巢穴,成了犯罪之處.教會變了質,見證全無價值,教會成了廢墟!魔鬼在教會中挑撥是非,引起紛爭,叫身體支離破碎.激動人的血氣,造成人的衝突；利用人的舒服,造成人的不滿；利用人的意見,引起人的爭端.魔鬼千方百計,目的要叫身體分裂,肢體破碎.教會遇此光景,如果沒有屬靈的分量,便索性任其公開表面化,公開的爭,公開的吵；如果有屬靈的分量,便覺得如此太無屬靈的體統.就把不滿的事,潛伏在心,表面好像相安無事,實則隨時可以爆發.魔鬼用敗壞的方法,使教會變質,用分離的手段,使身體破碎.
\newline
\newline
二,保護神自己的利益
\newline
\newline
用禱告保護神的見證　尼希米建造聖城時,用禱告來保護.尼希米記全卷,幾乎每章都有他的禱告.有一位姊妹,在禱告會中站起來禱告,因她話語過於累贅,拖著時間,一位弟兄實在忍不住去勸止她.那姊妹因失面子,而非常生氣,便大哭大鬧,大吵大叫,把聚會的氣氛,弄得一片混亂.會眾勸她不聽,又不能責她,因她已失去理智.那怎麼辦呢?有幾個屬靈的長者,馬上宣佈跪下懇切禱告,那位姊妹才受到控制,當眾認罪；因此,見證不但沒有破壞,反而更加榮美.
\newline
\newline
用聖潔保護神的見證　聖潔的教會,有聖靈的管制.弟兄姊妹與神之間,一有不對,即時起敬畏的心,好像聽見有微小的聲音說:「這是聖地,要把你的鞋脫下.」另外一方面的人看見了,就深覺神真是在他們中間.一個聖潔的教會,必能及時在神面前對付罪,不敢拖延；這樣,必叫一班不冷不熱的信徒復興過來.
\newline
\newline
每一個神的子民,都有責任把城牆建造起來；教會與社會應有清楚劃分的界線,叫神的子民,和神的利益都能得到很好的保護.尼希米說,城牆已經倒塌,城門已經被焚燒；是因他們列祖拜偶像,惹動神的怒氣而被擄到巴比倫去.(巴比倫是拜偶像的中心)後來神使他們從巴比倫歸回,他們受過教訓之後就把拜偶像的罪對付清楚.神在教會中,應有祂當神之地位和尊榮.若世界在教會中有地位,世界就變成偶像.在教會中佔了神的地位,許多人把它當作神來尊崇.教會亦以「傳統」取代神的地位,這些傳統沒有聖經為根據,是上一代留下來的；從前這一套,現在也是這一套,沒有想到神給人每天都是新的.有了一套傳統的東西,我們就不會再求新的引導,不求新的帶領了.老一套的傳統,幾乎把教會壓死.這套傳統,就成了教會的偶像.
\newline
\newline
修造城牆,修造的意思是把城牆倒塌的原因找出來,加以對付.并且在原有的根基上再造.現在教會為何會分裂,而且越分越多越亂呢?弟兄姊妹愛主是不錯的,屬靈程度也是不錯的.但他們對教會不滿,便另立爐灶,這是今日教會的大錯誤.尼希米建立城牆,不是在別的根基建立,乃是在原有的根基修造.各位弟兄姊妹!如果我們個個另立爐灶,教會前途就不可收拾了,教會的見證也就失去了.尼希米對猶大人說:「來罷!我們重建耶路撒冷的城牆,免得再受凌辱.」(尼二17)尼希米要猶大人起來修造城牆,把所受的羞辱洗掉；原有的根基受了傷,要重新建立起來.教會應走尼希米的路,把城牆被毀被燒的原因找出,一齊對付,重新修造.
\newline
\newline
求主給我們看見,建造耶路撒冷城牆,最急要的任務,就是要把界限樹立起來建造城牆,保護神子民的利益,保護神自己的利益.城牆的功用,也就是教會的功用.今天晚上要緊的信息,就是請弟兄姊妹要愛自己的教會,不要厭惡自己的教會.教會有軟弱,我們自己有責任,應當用禱告,配搭,聖潔,真理,教導,追求來保護她.另建爐灶的事,不是神要我們走的路,把病原找出來,才是神要我們走的路.很多人常批評自己的教會,論斷自己的教會,殊不知越批評,教會越冷落；越論斷,教會越損傷.我們應走的路,是要建立城牆,叫神的子民利益不受虧損,這是我們應走的路,是教會重大的問題.求主給我負擔的信息,與弟兄姊妹一同分享.
\newline
\newline
弟兄姊妹!愛你的教會,愛神的教會；因為教會是神的家,是基督的身體,是我們的本分.求主給我們有這樣的心志.
\newline
\newline


\newpage

\section{第45屆~港九培靈會~吳勇~第八講}
\label{sec:1273}
\textbf{會幕完工了}
\newline
\newline
連結: \href{https://www.hkbibleconference.org/session-message/view/1273}{\texttt{https://www.hkbibleconference.org/session-message/view/1273}}
\newline
\newline
\hyperref[sec:1272]{< < < PREV SERMON < < <}
~
\hyperref[sec:toc]{[返目錄]}
~
\hyperref[sec:1275]{> > > NEXT SERMON > > >}
\newline
\newline
(徒四16)論受膏者基督.(出卅二25-33)論以色列民拜金牛犢,摩西為他們禱告.(出卅六3-8)論百姓獻上禮物和工作.(出卅九32-33)論他們把做成的物件搬入會幕.(出四十1-4)論所作之物各盡其用,例如燈要點著,香要燃燒.(33-38)論工作完成,雲彩遮蓋會幕,耶和華的榮光充滿會幕.
\newline
\newline
摩西建造會幕完工,與神與人都有關係,現在我們要講與人的關係.
\newline
\newline
神的工作要用人當路來走,把祂的計劃藉人行出來.
\newline
\newline
摩西是神的好工人,一個好的工人是非常重要的.一個教會的強弱好壞,和工人是有密切關係的.好的工人,是磨煉出來的.摩西在埃及時,曾因兩個同胞爭吵,而為他們調解.豈料其中一人說「難道你要殺我,好像殺那埃及人一樣嗎?」摩西聽了很害怕,因為事敗了,想法老王必定會知道,也必不會放過他；所以他就急忙逃跑,到米甸曠野去.
\newline
\newline
磨煉有兩種:
\newline
\newline
一,神藉著人來磨煉　當時給摩西製造難處的人,不是外人,乃是他的同族.一個神的工人常常遇到麻煩,不是外邦人給的麻煩,乃是主內的肢體；不是無關係的給予麻煩,乃是有關係的人.摩西經過磨煉,才能成功,當他在曠野時,曾有大坍和亞比蘭結黨與他作對；還有哥哥亞倫和姊姊米利暗妒忌他.這些不是外人都是自己人,雖然他遭遇諸多難處,但他堅定不移,因為神要藉著人把他磨煉,直至成功.
\newline
\newline
二,神藉著環境來磨煉　摩西從小住在皇宮,突然逃到曠野去.由寶貴淪到貧窮,環境實在轉變得太厲害了!約瑟成為一個好工人,也是神用環境來磨煉成的.他在家裡,做人的兒子,在埃及做人的奴隸.在家裡穿彩色衣服,在埃及穿奴隸衣服.一下子受人寵愛,一下子被人鄙視,變化實在太厲害!環境極度的轉變,常使人站立不住.可是摩西一時被以色列人擁戴,一時被以色列人厭棄,一時人讚揚他,一時人咒罵他；他卻不以為然,不灰心也不喪膽.
\newline
\newline
一個好工人的產生,我想是從教導出來的.關於教導,在神學裡有很多種教導,有系統神學,有牧範學,有宗教比較學等等的教導.可是有一種教導,不是可從這些追求得到；乃是神親自教導,教他知道自己不能,知道神凡事都能.摩西本來在埃及,他實在對神有負擔,對以色列百姓有負擔.他以為憑著自己的熱情,年輕,埃及所學的學問,甚麼都能作；那裡想到憑這些一無所作,而且還要跑到埃及去.可見神不但要叫他的僕人知道,自己是無能的,還要知道神凡是都能的.當神揀選摩西的時候,他諸多推卻,說他不能勝任.因他已學會了「自己不能」的功課；然而他也學了「神凡事都能」的功課,所以有一次他對神說:「你給我們去,我們就去,你不給我們去,我們就不去.」他知道靠自己去是無用的,神帶他去才有用.神教導摩西,使他知道「自己不能」「神凡事都能」,這樣的工人,才不會倚靠自己,凡事都倚靠神,將神在他身上的大能完全顯明出來.
\newline
\newline
摩西是個好領袖　出埃及記卅二章載,摩西從山上下來時,看見百姓拜金牛犢,坐下喫喝,起來玩耍.他非常憤怒,就吩咐人就,你們把刀跨在腰間,各人殺他的弟兄,殺弟兄的意思,就是對付罪.表明不因手足之關係對罪苟息,對罪容讓.好的領袖對罪是絕不苟息容讓的.摩西聽見神說要滅絕百姓,就在神面前禱告,向神求情.聖經記載摩西的禱告,其中有一點一點的記號,有些解經家認為這是表明摩西向神禱告,因太過激動,話說不出來了.有些解經家認為是摩西禱告時,眼淚滴下來了.無論如何,這兩樣的意思,均足以表明他為以色列民禱告,是出於真誠的感情.摩西對神說「倘若你肯赦免」這表明一個好領袖為人代禱是不肯放鬆的.一直要禱告至神答應為止,否則不肯放下.我們常常為人禱告,神聽也好,不聽也可,都無所謂.但一個好的領袖,他的代禱是抓得緊緊的,不達目的,絕不放下.摩西對神說:「倘若你肯赦免他們的罪,......不然,求你從你所寫的冊上塗抹我的名.」由這話可知他肯犧牲自己.一個好的領袖不以自己的得失為念,乃以神和民的得失為念.
\newline
\newline
有一次我到外國,和一些學生談話,其中兩個人,批評倪柝聲和王明道蹧踏了神的恩賜,埋沒了神的恩賜.我對於這樣的批評甚覺驚訝.他們批判這兩位在內地的僕人,沒有顯出恩賜的功用,虧負了神,把神的恩賜蹧蹋了,把恩賜給埋沒了.我對他們說,作為一個神的工人,一生有兩件重大的事:一是工作,一是見證；我請他們批判一下,是工作重要,還是見證重要?有個學生迅速的回答:「當然見證重要.」我說:「這兩位神僕在內地,縱然沒有工作,但他們有沒有見證呢?」他們聽了纔醒悟過來.這兩位神的僕人留下了千古不朽的見證,他們以神的榮耀得失為念；以羊群的得失為念,而不以自己的得失為念.工作既屬重要,見證比工作更加重要!
\newline
\newline
會幕能完成,是與摩西有關係的.神找到這樣一個好工人,以色列百姓有一個這樣好的領袖.有了好的工人,兼有好的領袖,會幕才能完成.
\newline
\newline
好的以色列百姓　出埃及記卅六章三節論他們獻禮物,8節論他們獻工作.他們的奉獻多麼完全,不但奉獻禮物而且奉獻工作.在神的家中,有的人獻了禮物就不獻工作；獻了工作就不獻禮物,這是錯誤的.以色列百姓是把工作和禮物並獻,因為他們對神有認識.我們不但對耶穌要有認識,對基督也要有認識.耶穌二字的意思,就是要將自己的百姓從罪惡裡救出來.雖然神不能容許有罪變為無罪,但神能叫無罪的替我們成為罪.如果神能在地上找到一個無罪的,也只能代替一個人的罪,不能代替無數人的罪.這位要代替我們的罪的,必定要無罪的,也必定要無限的.在地上是找不到的,惟有神自己來解決；祂自己成了人,祂自己是無罪,可以代替人的罪.所以我們想到耶穌,就想到了神的愛,想到了神的恩.這樣的恩和愛,怎叫我們的心不受感動呢!基督二字的意思,使徒行傳說基督是受膏者.聖經上說有幾種人是受膏的:就是先知,君王,和祭司.先知受膏是受托於神,作神的出口.君王受膏是受托於神,治理神的百姓.祭司受膏是受托於神,管理聖殿的事.所謂受膏,就是受托託的意思.神把一個計劃托在基督身上,這計劃就是神的國,神的國就是神的教會.我們認識基督,知道神有大的計劃在祂裡面,我們因認識基督的緣故,就投身其中,在神的計劃裡一同有份.神的子民對耶穌基督認識清楚,才能獻上禮物並獻上工作,好像早時以色列百姓對神的認識透澈又清楚,所以他們把禮物和工作卻拿到會幕來.今天神的兒女,如果對主有清楚的認識,才能把物質和力量拿到教會來,使教會有物質和力量完成神的旨意.好的以色列百姓,一切都是為著會幕,會幕是預表教會.會幕是各種器具配搭起來的,教會是各肢體配搭起來的.會幕是祭司事奉神的地方,教會也是神子民事奉神的地方.當時的以色列百姓是好的百姓,因他們對會幕認識清楚,一切都是為著會幕.
\newline
\newline
出埃及記四十章載:燈台做好了,點上燈；香壇作好了,點上香；洗濯盆作好了,盛了水；這許多的器具,都是為著會幕使用,不是為了觀賞.教會裡面,我們每一人都是一個肢體,一個器皿.可惜今天教會的器皿,未能盡其功用,只不過擺著給觀賞而已.如果教會多數是這樣的人,教會怎會有見證,怎會有能力呢?你是個器皿,是個供人觀賞的,還是被神使用的呢?另一方面,祭壇做好了,搬入會幕；餅桌做好了,搬進會幕；燈台做好了,搬進會幕.一切器具做好了,放在自己的地方.置於外面,是沒有用處的,要搬入會幕,才可盡其功用.燈台代表發光,生活榮耀主的見證,這是為著教會.餅桌代表生命的糧,一個傳道人釋放生命的糧,這是為著教會.香壇代表禱告,禱告為著教會.洗濯盆代表聖潔,也是為著教會.當時以色列百姓作那些器具都是為著會幕,若不搬進會幕去,就沒有用場,沒有價值.我們當省察自己,所作的是為著什麼,我們屬靈的追求,屬靈的事奉,是為自己還是為教會?好的以色列百姓把器具做好,不是留在自己家裡,而是一一搬入會幕.
\newline
\newline
以色列百姓,他們把香壇做好,把燈台,餅桌也做好,不是把它們擱置,而把各器皿按位置擺列.「單獨」為著「整體」,這真理很要緊.並不是說,個人的追求不要緊,個人的追求也是很要緊的.不過個人的追求,若不是為著整體,在神面前是毫無價值的；因為神所要的,是整體的需要.耶利米作先知時,是非常屬靈的.他是個流淚的先知,既是屬靈為何要流淚呢?他個人是不錯的,但因整體不好.那時以色列民族遠離神,面臨困境危機,所以他一直為國人流淚.耶穌基督也為耶路撒冷而流淚,祂說:「耶路撒冷阿!你常殺害先知,又用石頭打死那奉差遣到你這裡來的人；我多次願意聚集你的兒女,好像母雞把小雞聚集翅膀下,只是你們不願意.」(路十二34)為什麼耶穌哭呢?因為祂雖然好,但整體卻不好,因而叫祂的心傷感痛苦.弟兄姊妹!你個人的屬靈,更應當為著整體而屬靈.基督的教會,單獨的屬靈可能有,但整體的屬靈在那裡呢?如果沒有整體的屬靈,單獨屬靈也算不得什麼.神所要的,乃是整體的屬靈.個人屬靈是不錯的,但神安排你的崗位；是叫你個人為著整體,與整體配搭,想不到中國的教會,只有單獨的表現,而沒有整體的表現,單獨表現雖然美麗,但神要的是整體的表現.
\newline
\newline
好的以色列百姓,一切都是為著會幕.有好的工人,有好的領袖,有好的以色列百姓,會幕纔能完成.會幕完了工,當時有雲彩遮蓋會幕,耶和華的榮光就充滿了會幕(出四十34)雲彩是代表恩典,天空有雲彩,甘雨即降下；大地便得到滋潤,百物重生了.因為會幕完了工,才有雲彩遮蓋.他們不是完了自己的工,不是完了性質相同的工.(出埃及記常有句話說:「是照耶和華所吩咐他的」,)他們乃是完了耶和華所吩咐的工.
\newline
\newline
弟兄姊妹!如果一件事重複的講,表明那件事是重要的事.聖經那件事講得多呢?摩西五經,論到會幕.列王和歷代,論耶和華的殿.四福音論神要建立祂的教會.保羅書信論教會的內容和真理.今日神的教會在地上,好像一片朦朧的霧,怎能叫人佩服,怎能叫人接受呢?
\newline
\newline
好的以色列百姓,建完了會幕之後,有雲彩遮蓋.當所羅門獻殿的時候,神的榮耀也充滿了殿.
\newline
\newline
「耶和華的榮光」榮光代表神的同在.以色列人在曠野,神與他們同在,他們就一無所缺.他們能力不缺,日用不缺,引導也不缺.感謝神的恩典!我們若照著神心意完成祂的工作,必有榮光充滿,教會什麼都不缺；人才不缺,恩賜不缺,神跡不缺,能力也不缺.教會有光輝,真正能夠吸引外面的人歸主.感謝神的恩典!會幕完了工,隨即有雲彩下來,隨即充滿了榮光.各位弟兄姊妹!從這短短的例子,我們可以明白會幕是何等的重要.會幕完了工,不是沒有原因的.神找到了好的工人,以色列百姓找到好的領袖,更有好的百姓配搭.他們對神有清楚的認識；一切都是為著會幕,會幕纔能建成.照樣,神也把會幕的工交給我們,我們的會幕完了工否?我們的會幕到底是何樣式?可惜!神的兒女,雖是器皿,但擱置一邊,沒有搬入會幕,沒有盡其功用.很多神的兒女,只顧個人,不顧整體,這樣,神的家怎能建立起來?我們沒有傳福音的力量；工作沒有推動的力量,原因是我們對教會沒有認識,以致教會如此脆弱.今日的世代,單獨的力量起不了作用,必須整體的力量才會發生效用.可是有那幾個弟兄姊妹,發現了整體的重要,關心神的家,投身於神的家呢?
\newline
\newline
求主藉著這信息,讓我們對教會有新的認識,好叫教會在現時代為神發出好的功用,叫神的雲彩降下來,榮光彰顯.這個責任在於神的工人,也在於神的兒女.神的工人不能推諉於神的百姓,神的百姓也不能推諉於神的工人.應彼此配合,然後才能建立起來.
\newline
\newline
求主藉此信息,使我們關心神的家,投身於神的家.
\newline
\newline
會幕能夠建成,我們的教會也應該建立起來!
\newline
\newline


\newpage

\section{第45屆~港九培靈會~吳勇~第九講}
\label{sec:1275}
\textbf{同心合意興旺福音}
\newline
\newline
連結: \href{https://www.hkbibleconference.org/session-message/view/1275}{\texttt{https://www.hkbibleconference.org/session-message/view/1275}}
\newline
\newline
\hyperref[sec:1273]{< < < PREV SERMON < < <}
~
\hyperref[sec:toc]{[返目錄]}
~
\hyperref[sec:1277]{> > > NEXT SERMON > > >}
\newline
\newline
「凡遵守命令的,必不經歷禍難；」(傳八5)「延長年日」(賽五三10)「你們受洗歸入基督的,都是披戴基督了.並不分猶太人,希利尼人,自主的,為奴的,或男或女；因為你們在基督耶穌裡,都成為一了.」(加三27-28)這裡題到三種沒有分別:不分國籍,不分階級,不分性別.「直到我們眾人在真道上同歸於一.」(弗四13)(腓立比書一章5,12,13,二22)都是論到興旺福音的事,教會的一切都應當為著福音.
\newline
\newline
為何要禮拜,是為著福音.為何要奉獻金錢?為何要配搭建立身體?都是為著興旺福音.保羅在腓立比書題到他為福音被捆鎖,題到多人因福音與他同勞.可見福音的工作是需要付代價的,而保羅為了福音的工作是需要付代價的,他實在付上了很大的代價.在哥林多後書保羅題到他曾經受很多的苦；被鞭打,被石頭打,遇陸上及海上的危險.他到馬其頓去,後來他在監牢裡；被人打至皮破血流.他在以弗所傳福音時,當地的人煽惑群眾暴動,幾乎要把他撕碎.保羅為著福音的緣故,實在付上了代價.
\newline
\newline
中國的福音是由外國的宣教士傳來的.那些宣教士也實在付上了代價.他們離別國境親人,遠路而來.從前保羅從耶路撒冷到羅馬,那時交通不便,沒有飛機,沒有輪船；需經長久時間才可達到.美國一個宣教士克爾威廉,他到印度去傳道,要航行五個月的水程.馬禮遜費了八個月的時候,才自英國到達廣州傳道.戴德生搭了一艘四百七十噸的輪船,經七個月水程才來到中國.外國的傳教士來中國傳道,實在付上了很多的代價.他們離別家園,要往人地生疏的異國,通常要半年才有家信,言語又不通,無法與當地人交往,備嘗孤獨之苦!每個國家都有它宗教背景,要使他們棄假歸真,實不容易.傳道工作是艱難的,困苦的!此外,還有性命的危險.當義和團得勢的時候,很多宣教士殉道被殺了.為了福音,是需要付上大代價的.
\newline
\newline
福音有種種的關係:
\newline
\newline
福音與別人的關係　耶穌把福音帶到撒瑪利亞.有個婦人,正午出去打水.在猶太地方,天氣酷熱,人們打水,都在清晨或傍晚的；但這婦人卻在正午出來打水；因為她是個不名譽的女人,避免露面見人.但耶穌把福音傳給她之後,她完全改變過來,她就跑到城裡見那些她平時所不願見的人.告訴他們:「有一人把我素來的事都說出來了.」福音能改變人!
\newline
\newline
在工業發達的時代裡,最易使人的情緒煩躁；為了解脫煩躁,人尋找刺激；尋求娛樂的刺激求解脫.更有甚者,是找尋色情的刺激,麻醉的刺激.許多人飲酒來麻醉自己,以毒品來麻醉自己.因為情緒不安煩躁,若不求解脫,就不能活下去.
\newline
\newline
福音與個人的關係　保羅說:「若不傳福音,我就有禍了.」(林前九16下半)有些基督徒,以為這話是神威脅我們,其實不然.(傳道書八5)說:「凡遵守命令的,必不經歷禍患.」反過來說:不遵守命令的,就必經歷禍患了.主說:「神叫日頭照好人,也照歹人.」照這話解釋,不單歹人遭遇禍患,連好人也會遭遇禍患.也可說,不信的人會遭遇禍患,信的人也會遭遇禍患.那麼,信與不信有何分別?其中是大有分別的.不信的人遭遇,沒有真神可倚靠,也沒有真神可求告.至於信者適得其反,真神是我們的倚靠,我們可以向祂求告庇佑.
\newline
\newline
「萬軍之耶和華說,我曾呼喚他們,他們不聽,將來他們呼求我,我也不聽.」(亞七13)我想這就是基督徒的禍患.沒有一件事比福音更大,神呼喚我們在福音上有份,我們卻視如飄風之過耳；到了我們真有需要,向神呼求時,祂也不聽了.(以弗所六15)說:「用平安的福音,當作預備走路的鞋穿在腳上.」鞋穿在腳上可以幫助走路,保護腳底不至受傷,也可防釘子刺傷,玻璃割傷.腳穿了鞋,免得沾染污穢,還可增加美觀.保羅用福音的鞋作預表,就是說,我們向人傳福音,好像鞋穿在腳上一樣.
\newline
\newline
福音與教會的關係　教會擔起傳福音的責任,如(以賽亞五三10)說:耶穌要延長年日,耶和華以祂為贖罪祭,把眾人的罪孽都歸在祂身上.意思就是說,耶穌被釘在十字架上,把命捨了,後來埋葬,復活,升天.聖經說祂要延長年日.「延長」可作「擴大」解釋.耶穌基督的年日,不但要延長,且要擴大.祂把這樣的責任交給教會,教會要擔當延長耶穌的責任.教會是個身體.神造了亞當,對他說:「要生養眾多,遍滿地面.」生養眾多是延長的意思.(路加三章)說:「亞當之後有塞特,塞特之後有以挪士,」一代代傳下去.亞當所負的,是「延長」之責,而且要遍滿地面,意思就是也要「擴大」.亞當的後裔,挪亞三個兒子,含到非洲地方,閃在亞洲地方,雅弗在歐洲地方；他們擔負擴大的責任.亞當的身體所擔負的工作,是世世代代,各處各方的工作.教會的身體也照樣要擔起延長和擴大的責任.如果這一代不傳與下一代,基督的年日怎能延長呢?如果此地不傳到彼地,基督怎能擴大呢?所以傳福音是神給教會的大使命.
\newline
\newline
從聖經看來,福音的工作,是一直要擴大的,耶穌說:「天國好比芥菜種.」芥菜種是要擴大的.耶穌在約翰福音十章說祂另外有羊,不但猶太人中間有祂的羊；外邦人中間也有祂的羊.耶穌復活以後,叫門徒要作見證,從耶路撒冷直到地極,這都是要擴大的.讓聖靈充滿在人的身上,叫人有說各地方言的恩賜,其目的也是要把福音擴大.後來,祂給保羅有馬其頓的異象,也是要保羅把福音擴大；把福音從歐洲擴大到亞洲去.由耶穌時代至使徒時代,神一直把這旨意顯給我們.耶穌不但要延長下去,還要擴大出去.擔負延長和擴大,是教會的責任,如果教會不擔起這責任,教會在神面前將如何交代?教會不只有人守禮拜,有人奉獻金錢,就可以滿足；教會的一切,當為著福音.今天的教會,不明白屬靈的定律.屬世的定律是越積越多；屬靈的定律卻是越分越多,越流越遠.從五餅二魚的神跡來看,在自己的手裡只夠供一人之需,分了出去,可供無數人之用.如以西結四七章所說,水越流越深,開始時在踝子骨,漸而到膝,到腰.屬靈和屬世的定律完全不同；屬世的,要抓往它,留著它；屬靈的,要放開它,要分出它,因而愈加豐富.在猶太地方,有個鹽海,又叫死海.也有個加利利海,又叫活海.活海裡有生物,海面常有雀鳥飛翔.死海則連鳥都不敢飛過,恐被海裡的蒸氣蒸倒下來.一死一活的海,同是約但河的水流下,同出一源,為何一死一活呢?原因很簡單；所以死是因水流下囤積起來.所以活是因水流下去以後,不斷又分流出去.教會屬靈光景的死,是因我們從神所得的,不肯向人分出去；結果,教會死氣沉沉.
\newline
\newline
怎樣興旺福音　保羅在腓立比書一章,教我們要同心合意,福音才能興旺.
\newline
\newline
同心合意必有聖靈同在(同工)從舊約多處可以看到:以利亞到迦密山上,把十二塊石頭砌在一起,他以十二塊石頭代表十二支派；因那時十二支派已經分開了,因羅波安壓迫百姓之故,分裂為南北兩國,就是猶太和以色列.以利亞把十二塊石頭砌疊在一起,是同心合意的預表.他把壇築好,向神禱告,就有火降下.可見同心合意,就有聖靈的工作.詩篇告訴我們:「看哪!弟兄和睦同居,是何等的善,何等的美.這好比那貴重的油,在亞倫的頭上流到鬍鬚,又流到他的衣襟.」(詩一三三1-2)和睦同居,也是同心合意.若同心合意,聖靈就澆灌下來.耶穌說:「地上若有兩三個人奉我的名聚會,我就在他們中間.」兩三個人也是同心合意；這樣,主就與我們同在.五旬節時,教會有聖靈同工；其中有兩個最顯著的現象,其一,是不斷推展出去,從耶路撒冷推展到撒瑪利亞,一直推展到整個亞細亞.其二,有聖靈同工,教會就有推展的能力.還有一現象,如果有聖靈同工,就有神蹟奇事出現.彼得,有神跡奇事隨著他.保羅,也有神蹟奇事隨著他.神蹟奇事足以證明神的恩道,叫一個頭腦很大的人,在祂面前,不能不俯伏下來.
\newline
\newline
有一件事我曾多次提過:人在身體強壯時,常會自高自傲；到了面臨死亡邊緣時,便沒有甚麼可誇可傲之處了.有個大學的校長,他患了嚴重的心臟病；當我向他傳福音時,他好像很聽話的孩子一樣.次日,他對我說他要受水洗,護士聽見了,連忙拉一下我的衣服,示意叫我離開病房.對我說:『吳長老!你千萬不要為他施浸,因為主診醫生吩咐要特別小心照料他,因他的心臟脆弱至極；若稍受刺激,恐隨時有性命危險.如果你把他浸在水裡,恐怕到時拉不上來呀!」我覺得她言之有理,萬一死在水裡,消息給報紙刊登出來,個人的名譽不要緊,可是教會的名譽就不得了.我正思索之際,校長在房裡叫我,他再三叮囑明晨為他施浸的事.那時,我不便對他說明.各位!難道他信耶穌,他願意接受,而不為他施浸,豈不矛盾嗎?那時我只好答應了.次日早晨,我很早起來禱告,因我極擔心會出事,我把這事交給神.到八點鐘,他的兒子陪他乘汽車,如約來到禮拜堂門口,他不讓兒子扶進去,他邊走邊喘息地進到裡面,其實極短的距離,卻走了很長的時間.後來他的兒子扶他下水.我沒有忘記,那時我禱告兩句話,一是為自己禱告,求神給我看見榮耀.一是為他禱告,求神給他經歷榮耀.接著,我問他:「你信耶穌是神的兒子嗎?答:「信」,我就把他浸在水裡了,隨即把他拉起來,我看見他面露笑容,當他兒子要扶他時,他卻把兒子推開,自己走動.他不需再回到醫院,而返東海大學,繼續執教.他就是東海大學的張啟東校長.教會同心合意,就有神蹟奇事,藉著神蹟奇事,能夠證明神的真道.
\newline
\newline
同心合意就有榮美的見證(加拉太書三章27-28)我們看見三種的「沒有分別」第一是不分國籍,不分猶太人或希利尼人.各國有不同的規矩,如中國人較為好客,宴客時,常常把筷子放在自己咀裡啜一下,然後夾菜敬賓,表示親切.如果請的是外國客人,真會把他嚇壞了；但如果同在主裡面,他就能受得起.西國人對於時間很寶貴,當你造訪,逾時還不告辭,他便會下逐客令的；他下逐客令,我們仍受得了.這些見證,都是很美的.第二是不分階級,就是不分貧富,在教會中,無論貧富,都交通得來,配搭得來.富的和貧的交通,不以為是遷就；貧的和富的配搭也並非高攀.這樣的見證多麼美,第三是不分性別,就是不分男女.並非不男不女,乃是表明我們在教會裡,男的女的,都一齊事奉,謹謹慎慎的,清清潔潔的,這個見證實在很美!所以一個教會同心合意,有榮美的見證,才能叫社會的人景仰和欽佩.
\newline
\newline
同心合意就有滿足的喜樂(腓立比書一至三章)論到喜樂.第四章說要常常喜樂,大大喜樂.為什麼有喜樂呢?第一章題到同心合意,第二章題到很多相同之處,意思是大家同心合意就有喜樂.這裡所謂相同(一樣),並非我們的個性相同,也非興趣一樣.那麼,到底有何相同,有何一樣呢?神所造的都不相同,才顯出祂的完全,才顯出祂的美麗.我們的個性,有急的,有慢的,有內向的,有外向的.如果大家一樣的急,或一樣的慢,那怎麼辦呢?如果大家都屬內向,有事大家都不講；這樣,人與人之間,難題就大了.人的興趣也不相同,有的喜歡遊山,有的喜歡玩水,有山有水,才顯出美麗.聖經說,我們的心相同,向著神的心相同,為著神的事工心要相同,這樣教會就有喜樂.我們在一起時,氣氛和諧.我們一起事奉,工作輕鬆,這樣的教會,內部有歌可唱,外面有力量可以發展.所以,同心合意,教會才有滿足的喜樂.
\newline
\newline
我們應在何事上同心合意呢?
\newline
\newline
一,在真道上同心合意　保羅在以弗所四章說「在真道上同歸於一.」真道是信仰的意思,我們是在信仰上要同心合意.
\newline
\newline
二,在感情上同歸於一　別讓感情把我傷,也不要傷感情.何謂感情把我傷?在尼希米記,有個人名叫多比雅,他外面與神有關,裡面卻與神無關.他按名是活的,其實是死的.多比雅是阿們族.是羅得和他女兒所生的後裔.這種人不能入耶和華的會,可是他卻住在聖殿裡,而且參於事奉.因他是當時大祭司以利亞實的親戚,為了感情的緣故,容許多比雅在聖殿,這就是感情把他傷了.今天中國的教會,給感情傷了的事,在在皆是,明知那人行為不對,明知他得罪神,但為了情面,不敢指摘.至於傷感情最易引起的由於口,和由於發生誤會.言語不小心和種種誤會,彼此傷了感情,感情一傷,目標雖是相同,步伐不能一致.興旺福音是我們共同的目標,為了感情受傷,步伐不能一致,福音不能傳開,神的家大受虧損.
\newline
\newline
三,在事奉上同心合意　尼希米建造城牆,聖經說:「某某人對著城牆修造.」這話何解?意思是我這段的修造和你那段的修造有關.教會也是如此,你的事奉和我的事奉有關,彼此的事奉要互相配搭；否則,身體就支離破碎,不能發生功用.在教會中,各個團契,息息相關,應一致行動,對付這邪惡的世界.
\newline
\newline
韓國教會,在若干年前發生了一件事:當時韓國政府擬印行一種新鈔票,票上印有和尚之像.教會風聞此事,即上書政府提出抗議；理由是鈔票既屬全國性使用,但全國百姓,有佛教徒,也有基督徒,所以不能以和尚代表所有的百姓.政府當局認為此乃政府之決策,教會無權干涉,故置之不理.後來教會再次上書,聲明若政府印行這種鈔票,全國教會將堅持拒用.結果,政府只好取消計劃.由此可見,整體的力量能夠影響國家和社會.弟兄姊妹!我們應在信仰上同心合意,在感情上同心合意,在事奉上同心合意,教會才有能力；否則,沒有聖靈同工,沒有榮美的見證,也沒有滿足的喜樂.
\newline
\newline
求主今天晚上,把「同心合意」四個字放在我們心裡,為著福音同心合意,好叫福音興旺起來!
\newline
\newline


\newpage

\section{第45屆~港九培靈會~吳勇~第十講}
\label{sec:1277}
\textbf{住在耶路撒冷的人}
\newline
\newline
連結: \href{https://www.hkbibleconference.org/session-message/view/1277}{\texttt{https://www.hkbibleconference.org/session-message/view/1277}}
\newline
\newline
\hyperref[sec:1275]{< < < PREV SERMON < < <}
~
\hyperref[sec:toc]{[返目錄]}
~
\hyperref[sec:1246]{> > > NEXT SERMON > > >}
\newline
\newline
(歷代志下卅六17-19)論到城被仇敵攻破,百姓被擄,被殺；聖殿被焚燒,因城被攻破,影響實在大.(尼三1,5,12,31)論到大祭司,貴胄,女人,銀匠,商人等五種人.(尼四20)論百姓聽見角聲就聚集.(十一1-3)論住在耶路撒冷城的人,他們都甘心樂意.(啟十四4)；「羔羊無論往那裡去,他們都跟隨祂.」
\newline
\newline
住在耶路撒冷的人,民眾為他們祝福,而且名字記錄在聖經上,可見住在耶路撒冷的事,是非常重要的.不但是民的需要,也是神的需要.
\newline
\newline
(一)住在耶路撒冷的人,負有扭轉時代的使命.
\newline
\newline
關於扭轉時代,前面曾略為題過.尼希米記和瑪拉基書都是在舊約的末了,將要進入新約的時代,有一班人就負起轉移的責任.如撒母耳當時所負的責任,是從士師時代轉入君王時代.撒母耳膏了掃羅,又膏了大衛為王.猶大人被擄到巴比倫去,那是被擄時代；神就興起但以理,以斯拉,尼希米,起來負扭轉的責任,從被擄轉到歸回的時代.除了聖經記載之外,我們看見在十五世紀時,神興起馬丁路得,他也是個扭轉時代的人；因為那時天主教把真理歪曲了.耶穌在十字架上已經完成了救恩,祂大聲喊說「成了」,祂的工作是毫無殘缺極其完全的,不需要加上什麼.天主教興起之後,主張要加上人的功德,要加上金錢,然後才能得救.神興起馬丁路得,把那歪曲的真理恢復過來.十七世紀是歐洲非常黑暗的時期,教會受了政治的影響,淪落黑暗的境地,教會極其荒涼；神就興起衛斯理約翰,把荒涼的教會扭轉過來.今日的時代,因受工業社會的影響,世人都忙忙碌碌,為物質,為享受而忙碌.因此,一般神的兒女,沒有時間和精神,為教會,更為神所用.另一方面,因世界各樣物質的迷惑,影響教會至於冷落.凡此都需要神用人來扭轉.住在耶路撒冷的人,當時所負的,是扭轉時代的使命.城是很重要的,城被攻破,民眾和聖殿均受影響；民眾有些被殺,有的被擄去巴比倫,他們把聖殿焚燒.所以,城一破,不但影響以色列民,也影響到耶和華的聖殿.
\newline
\newline
城代表教會,(日前已解釋過)城中有聖殿,教會中有神.城是用磚頭砌成的,教會是肢體配搭起來的.城是以列民生活之所,教會是信徒屬靈生活之所.城破,代表教會荒涼.城一破,民被擄,失去護衛,百姓無安全.照樣,教會荒涼,便沒有愛心牧養,沒有天上的信息餵養,弟兄姊妹沒有配搭操練,神的兒女就沒有安全感了.有些就被世界擄去,有些就被世界殺害了.尼布甲尼撒攻破耶路撒冷城,殺害了很多以色列百姓.很多神的兒女,也是被罪所殺,或被肉體所殺；好像以西結所見的異象「遍地都是骸骨.」教會中有許多尸體,乃是因教會荒涼,神的兒女被擄被殺所致.城破殿被毀,殿是敬拜的所在.聖經給我們看見,神與人之間的歷史,就是一個敬拜的歷史,從開始直到末了啟示錄,都是敬拜.一般教會的敬拜,不過是從屬世的興趣,轉向宗教的興趣而已.這種興趣的敬拜,是儀式的敬拜,是心不在焉的敬拜,沒有靈的敬拜.從前以色列百姓到聖殿,他們有敬畏的心,也有敬畏的態度.可是現在,信徒來到教會,大多沒有敬畏的心和敬畏的態度,卻說:「我們基督徒的敬拜,不像天主教只是儀式上,外表上,無心的敬拜.」豈知基督教的敬拜,如今已落到不但無心,連儀式也沒有了.城一破,殿立刻被毀!
\newline
\newline
(二)住在耶路撒冷城的人負有看守和供應的使命　城代表見證,城存在,見證也存在.住在耶路撒冷的人,作為我們的榜樣,他們一邊看守,一邊供應,以使命看顧見證,供應見證；以禱告看守見證,叫仇敵無機所乘.以肢體的關切看守見證,不叫一個弟兄或姊妹成為教會的破口.耶穌說:「凡有的,還要加給他,叫他有餘；凡沒有的連他所有的,也要奪去.」(太十三12)(彼得後書一5-7):「有了信心,又要加上德行......加上知識......加上節制......加上忍耐......加上虔誠......加上愛弟兄的心......又要加上愛眾人的人.」一一加上去.弟兄姊妹!我們應當為著見證,供應這見證,把自己的部份加進去,好叫見證更豐滿,更榮耀.如果自己的部份是火熱的,加上進去,教會就增加一分熱量.我們若有謙卑,供應進去,教會就加一分謙卑.我若有愛心,供應進去,教會就加一分愛心.想不到人沒有供應於教會,以致見證活不出來,見證榮耀不起來.這是今日教會的光景!
\newline
\newline
(三)住在耶路撒冷之人的見證　凡住在耶路撒冷的,名字都錄在聖經上,表明他們是神所記念的；因為體會神的心,與神同心,與神同行,神看他們如同寶貝.好像大衛的三個勇士,聽見大衛想喝伯利恆的井水,他們奮勇衝過敵人的營盤去取水.他們體會大衛的心,記念大衛的心,所以大衛看他們非常寶貝.城被焚燒了,城牆已倒塌了,神的心必為此傷痛.住在耶路撒冷的人,他們體會神的心,一同起來分擔神的傷心和痛苦,他們同心合力,趕快修造城牆,免得神的名繼續羞辱下去.這班人分擔神的感覺,供應神的需要,所以他們被神記念.神有莫大的能力,我們靠祂的權能行祂的旨意；但神需要人作祂的出口,藉著人,把祂的作為顯明出來.這班住在耶路撒冷城的人,他們作屬靈的先鋒,叫神的旨意不停頓下去,反而藉著他們行了出來.(這是從神方面講的)從人的方面來講,雖有很多人看見環境的艱難,自己的軟弱,致無法行出來,所以需要這班屬靈的先鋒；以振作他們的膽量,帶領他們的腳步.歷史告訴我們,以色列人在米甸人蹂躪之下,土產被毀壞,糧食被斷絕；他們想解脫現狀,卻沒有膽量.那時基甸作屬靈的先鋒,登高一呼,以色列百姓立即響應,起來跟隨他.如今我們的靈性也是一片荒涼,誰不想復興?誰不想改變現狀?但心有餘而力不足；沒有勇氣,沒有膽量,實在需要屬靈的先鋒!
\newline
\newline
屬靈先鋒者具有美好見證　尼希米三章告訴我們:當時到耶路撒冷去的,有大祭司,有貴胄,有女人,有銀匠,也有商人；雖是各色各樣不同身份的人,但他們都能夠互相配搭,一齊同心建立.為了城的緣故,把自己的身份丟腦後,屬世的資歷也丟在一邊；他們覺得所作的,是一件莊嚴有意義的事.這樣的見證多麼美!
\newline
\newline
「羔羊無論往那裡去,他們都跟隨祂.」這是服從管治的意思,這個問題非常重要,教會應有這樣的次序.當時的百姓都是頑強的,各有各的觀念,各有其觀點.在以斯拉時的工作,他們都是各行其事；因為人人都不願放棄自己的觀念,堅持著自己的意見.其中有些人更是自私自利.到了尼希米時代,這樣的光景過去了.他們各把地位讓給尼希米,而且認他為大家所當跟從的人,都照著尼希米吩咐他的去做.因此,那時工作進行極其順利,五十二天便完成了.有些猶大人,起初曾侵吞別人的田產,剝削他人的利益.後來他們跟隨尼希米,凡事聽從他的吩咐,把剝削來的,物歸原主.那時尼希米無論有何要求,他們都照辦.真是羔羊無論往那裡去,他們都跟隨.教會之中,如果有這樣的見證,是多麼美!牧人的帶領,羊群都跟隨；這不但有美好的見證,而且發出很大的力量.
\newline
\newline
從前猶大人在曠野,他們的行止,都是隨著角聲；角聲一響,他們都起來聚集.角聲是代表聖靈的聲音.耶穌在啟示錄說:「聖靈向眾教會所說的話,凡有耳的,就應當聽.」昔日神的百姓一切行動,完全受著角聲支配.今時教會的一切行動,也應當受著聖靈的支配.五旬節時代,因為受著聖靈的支配,所以每一天都是新的.使徒行傳第一章他們聚集在一起禱告.第二章聖靈降臨下來,他們所傳的道,叫人扎心.第三章瘸子起來行走,神蹟奇事出現.第四章約翰和彼得被捕,他們對官長說:「聽從你們,不聽從神,這在神面前,合理不合理你們自己酌量罷.」第五章亞拿尼亞和撒非喇,因為貪圖虛浮的榮耀,倒斃在使徒腳前.再繼續下去,執事被選出來了.每章都有新事發生.為什麼?因為神的恩典每天都是新的.可惜今日的教會,撇開聖靈的聲音,用著固定的規條,固定的計劃,限制神的作為,以致我們不倚靠聖靈,而倚靠自己的計劃,用了一套傳統的東西完全限制自己,沒有人等候聖靈；也沒有人仰望聖靈,自己的一套把自己壓死.當時的百姓用角聲來帶領,今日神對祂的教會,也是用聖靈帶領我們.可是教會憑著肉體的衝動,頭腦的聰明.其實處於這樣的世代,肉體的衝動,頭腦的聰明能作什麼呢?求主給我們看見這些人的見證,他們不但順從領袖的帶領,他們也順從聖靈的帶領,所以城牆才能完成.他們順從角聲,就是順從聖靈的引導,所以他們有豐滿榮美的見證.
\newline
\newline
這些住在耶路撒冷的人,他們獻上甘心的祭.他們住在那裡,並不是出於勉強,不是出於規條的限定,不是為著懼怕有不良的效果,也不是為著得報賞.他們只有一個心志,他們已經看見寶座的旨意,就投身在這旨意中；不顧自己的得失,好像路得馬丁改教的時候,他沒想到結果有個路得會.有很多青年人當他舉步的時候,就顧慮到第二步第三第四步,以致連一步都沒有跑上.從前神召亞伯拉罕,叫他離開本地,本族,父家,往神所指示的地方去.神並無預告他,迦南地的氣候風土人情如何?亞伯拉罕惟有遵從,沒有顧慮；因他知道一切神會負責,有神負責就夠了,無需考慮.
\newline
\newline
耶路撒冷是見證的中心,我們應當住在耶路撒冷.無可否認的,我們在家裡是見證,在學校是見證,在工作單位也是見證,任何地方都是見證.但耶路撒冷是見證的中心,教會是見證的中心,神要我們活在見證的中心；大家來關心這見證,負責這見證,維持這見證,好叫見證更加豐滿.大家來保護這見證,好叫見證不受虧損.大家來供應這見證,服事這見證.弟兄姊妹!我們是神用重價買來的,為的是買來使用,叫你站在見證上有用處.
\newline
\newline
住在耶路撒冷的人需要付上代價　最近,我領了一個少年佈道會,他們派了一個女孩子來招待我.第一天,她就被她父親打傷了頭,第二天,被她父親趕出去.她為了要供應這個見證,她付上了這些代價.兩星期前我在加拿大,有個護士關心於這個見證,她要供應這個見證,參與教會的服事；為此被她父親關起來,叫她沒有自由.為了見證,代價是要付的.當時住在耶路撒冷的人,他們要離開家園,離開他原來美好的環境,舒適的生活.如果他不肯付上這些代價,就不能住在耶路撒冷.有的人希望住在耶路撒冷,可是他們留戀自己的經營,愛慕自己進款豐富的職業,始終不願付上代價.當年的弟兄姊妹!不錯的,雖然住在耶路撒冷吃飯的錢,是許多人拿出來的,可是我們應當認定是從神那裡接取的.前幾天,我看見一個牧師的兒子,他告訴我:「我們離開教會已經廿多年,沒有教會的麻煩也已廿多年了.」許多青年人,不敢住在耶路撒冷,就是和這牧師的兒子一樣.固然教會有很多事情叫人寒心,叫人抬不起頭來；但是我們應當分擔主的苦難,許多青年人為此緣故,裹足不前,以致耶路撒冷無人居住.無人去看顧這個見證,去保護這個見證,以致耶路撒冷被仇敵蹂躪.教會落到如此地步,是因你不肯住在耶路撒冷；不負責任,沒有供應,沒有保護,以致落到荒涼的地步.對社會沒有功用,只跟著社會的尾巴走.弟兄姊妹!我們再看那班住在耶路撒冷的人,他們願付代價,願意犧牲,在住耶路撒冷,供應她,保護她.雖然這班人是少數,然而他們放棄一切來到見證的中心,他們遺言囑咐,叫我們要住在耶路撒冷,可是有多少人理會呢!
\newline
\newline
尼希米記十一章開始,有很多名字記載,讀經時常會把那些名字忽略了,以為沒有意思.但是人看為無意思的,神卻看為寶貴,因為那些人曾獻上甘心的祭.把自己獻上的,不顧自己利益的人；聖經會記念他們,神也記念他們,因他叫神的心得著安慰和滿足.他的名在神的腳前,到了有一天,你到天上,神要指給你看；這班住在耶路撒冷的人,是祂所欣賞,是祂最看為寶貴的人.
\newline
\newline
神叫我們住在耶路撒冷,就是要我們來供應祂的見證,來保護她的見證,一同來看顧祂的見證.如果耶路撒冷無人居住,城就被仇敵蹂躪,民將被擄,被殺,殿將被焚燒.
\newline
\newline
神在呼喚我們,住在耶路撒冷吧!
\newline
\newline


\newpage

\section{第45屆~港九培靈會~唐佑之~第一講}
\label{sec:1246}
\textbf{環境與經歷}
\newline
\newline
連結: \href{https://www.hkbibleconference.org/session-message/view/1246}{\texttt{https://www.hkbibleconference.org/session-message/view/1246}}
\newline
\newline
\hyperref[sec:1277]{< < < PREV SERMON < < <}
~
\hyperref[sec:toc]{[返目錄]}
~
\hyperref[sec:1247]{> > > NEXT SERMON > > >}
\newline
\newline
但以理書 2:1-2:49
\newline
\begin{longtable}{cl}
\hline
\hline
章節 & 經文 (和合本修訂版)\\
\hline
2:1 & \begin{tabularx}{0.7\textwidth}{X} 尼布甲尼撒在位第二年，他做了很多夢，心裡煩亂，不能睡覺。 \end{tabularx} \\ \\ \relax
2:2 & \begin{tabularx}{0.7\textwidth}{X} 王吩咐人將術士、巫師、行邪術的和迦勒底人召來，要他們把王的夢告訴王；他們就來，站在王面前。 \end{tabularx} \\ \\ \relax
2:3 & \begin{tabularx}{0.7\textwidth}{X} 王對他們說：「我做了一個夢，心裡煩亂，想要知道這是甚麼夢。」 \end{tabularx} \\ \\ \relax
2:4 & \begin{tabularx}{0.7\textwidth}{X} 迦勒底人用亞蘭話對王說：「願王萬歲！請將夢告訴僕人，我們就可以講解。」 \end{tabularx} \\ \\ \relax
2:5 & \begin{tabularx}{0.7\textwidth}{X} 王回答迦勒底人說：「這事我已決定，你們若不把夢和夢的解釋告訴我，就必被凌遲，你們的房屋必成糞堆； \end{tabularx} \\ \\ \relax
2:6 & \begin{tabularx}{0.7\textwidth}{X} 但你們若能說出這個夢和夢的解釋，就必從我得到禮物、賞賜和殊榮。現在，你們要把夢和夢的解釋告訴我。」 \end{tabularx} \\ \\ \relax
2:7 & \begin{tabularx}{0.7\textwidth}{X} 他們再一次回答說：「請王將夢告訴僕人，我們就可以講解。」 \end{tabularx} \\ \\ \relax
2:8 & \begin{tabularx}{0.7\textwidth}{X} 王回答說：「我確實知道你們是故意拖延，因為你們知道這事我已決定。 \end{tabularx} \\ \\ \relax
2:9 & \begin{tabularx}{0.7\textwidth}{X} 你們若不將夢告訴我，只有一個辦法對待你們；因為你們彼此串通，向我胡言亂語，要等候情勢改變。現在，你們要將夢告訴我，讓我知道你們真能為我解夢。」 \end{tabularx} \\ \\ \relax
2:10 & \begin{tabularx}{0.7\textwidth}{X} 迦勒底人回答王說：「世上沒有人能解釋王的事情；從來沒有君王、大臣、掌權者向術士、巫師，或迦勒底人問過這樣的事。 \end{tabularx} \\ \\ \relax
2:11 & \begin{tabularx}{0.7\textwidth}{X} 王所問的事很難，除了不與血肉之軀同住的神，沒有人能在王面前解釋。」 \end{tabularx} \\ \\ \relax
2:12 & \begin{tabularx}{0.7\textwidth}{X} 王因這事生氣，大大震怒，吩咐滅絕巴比倫所有的智慧人。 \end{tabularx} \\ \\ \relax
2:13 & \begin{tabularx}{0.7\textwidth}{X} 命令發出，智慧人將要被殺，人就尋找但以理和他的同伴，要殺他們。 \end{tabularx} \\ \\ \relax
2:14 & \begin{tabularx}{0.7\textwidth}{X} 王的護衛長亞略奉命去殺巴比倫的智慧人，但以理用婉言和智慧回應， \end{tabularx} \\ \\ \relax
2:15 & \begin{tabularx}{0.7\textwidth}{X} 向王的大臣亞略說：「王的命令為何這樣緊急呢？」亞略就把事情告訴但以理。 \end{tabularx} \\ \\ \relax
2:16 & \begin{tabularx}{0.7\textwidth}{X} 於是但以理進去求王寬限，好為王解夢。 \end{tabularx} \\ \\ \relax
2:17 & \begin{tabularx}{0.7\textwidth}{X} 但以理回到他的居所，把這事告訴他的同伴哈拿尼雅、米沙利、亞撒利雅， \end{tabularx} \\ \\ \relax
2:18 & \begin{tabularx}{0.7\textwidth}{X} 要他們祈求天上的神施憐憫，將這奧祕指明，免得但以理和他的同伴與巴比倫其餘的智慧人一同滅亡。 \end{tabularx} \\ \\ \relax
2:19 & \begin{tabularx}{0.7\textwidth}{X} 這奧祕就在夜間異象中顯明給但以理，但以理就稱頌天上的神。 \end{tabularx} \\ \\ \relax
2:20 & \begin{tabularx}{0.7\textwidth}{X} 但以理說：「神的名是應當稱頌的，從亙古直到永遠！因為智慧和能力都屬乎他。 \end{tabularx} \\ \\ \relax
2:21 & \begin{tabularx}{0.7\textwidth}{X} 他改變時間、季節，他廢王，立王；將智慧賜給智慧人，將知識賜給聰明人。 \end{tabularx} \\ \\ \relax
2:22 & \begin{tabularx}{0.7\textwidth}{X} 他顯明深奧隱祕的事，洞悉幽暗中的一切，光明也與他同住。 \end{tabularx} \\ \\ \relax
2:23 & \begin{tabularx}{0.7\textwidth}{X} 我列祖的神啊，我感謝你，讚美你，因你將智慧才能賜給我，我們所求問的現在你已指明給我，把王的事給我們指明。」 \end{tabularx} \\ \\ \relax
2:24 & \begin{tabularx}{0.7\textwidth}{X} 於是，但以理進到王所派滅絕巴比倫智慧人的亞略那裡去，對他這樣說：「不要滅絕巴比倫的智慧人，求你領我到王面前，我可以為王解夢。」 \end{tabularx} \\ \\ \relax
2:25 & \begin{tabularx}{0.7\textwidth}{X} 亞略就急忙領但以理到王面前，對王這樣說：「我在被擄的猶大人中找到一人，能將夢的解釋告訴王。」 \end{tabularx} \\ \\ \relax
2:26 & \begin{tabularx}{0.7\textwidth}{X} 王對那稱為伯提沙撒的但以理說：「你能將我所做的夢和夢的解釋告訴我嗎？」 \end{tabularx} \\ \\ \relax
2:27 & \begin{tabularx}{0.7\textwidth}{X} 但以理回答王說：「王所問的那奧祕，智慧人、巫師、術士、觀兆的都不能告訴王， \end{tabularx} \\ \\ \relax
2:28 & \begin{tabularx}{0.7\textwidth}{X} 只有那在天上的神能顯明奧祕。他已把日後將要發生的事指示尼布甲尼撒王。你在床上做的夢和你腦中的異象是這樣： \end{tabularx} \\ \\ \relax
2:29 & \begin{tabularx}{0.7\textwidth}{X} 你，王啊，你在床上所思想的是關乎日後的事，那顯明奧祕的主已把將來要發生的事指示你。 \end{tabularx} \\ \\ \relax
2:30 & \begin{tabularx}{0.7\textwidth}{X} 至於我，那奧祕顯明給我，並非因我智慧勝過一切活著的人，而是為了讓王知道夢的解釋，知道你心裡的意念。 \end{tabularx} \\ \\ \relax
2:31 & \begin{tabularx}{0.7\textwidth}{X} 「你，王啊，你正觀看，看哪，有一個很大的像，這像甚高，極其光耀，立在你面前，形狀非常可怕。 \end{tabularx} \\ \\ \relax
2:32 & \begin{tabularx}{0.7\textwidth}{X} 這像的頭是純金的，胸膛和膀臂是銀的，腹部和腰是銅的， \end{tabularx} \\ \\ \relax
2:33 & \begin{tabularx}{0.7\textwidth}{X} 腿是鐵的，腳是半鐵半泥的。 \end{tabularx} \\ \\ \relax
2:34 & \begin{tabularx}{0.7\textwidth}{X} 你正觀看，見有一塊非人手鑿出來的石頭打在它半鐵半泥的腳上，把腳砸碎； \end{tabularx} \\ \\ \relax
2:35 & \begin{tabularx}{0.7\textwidth}{X} 於是鐵、泥、銅、銀、金都一同砸得粉碎，如夏天禾場上的糠秕，被風吹散，無處可尋。打碎這像的石頭成了一座大山，覆蓋全地。 \end{tabularx} \\ \\ \relax
2:36 & \begin{tabularx}{0.7\textwidth}{X} 「這就是那夢；我們要在王面前講解那夢。 \end{tabularx} \\ \\ \relax
2:37 & \begin{tabularx}{0.7\textwidth}{X} 你，王啊，你是諸王之王。天上的神已將國度、權勢、能力、尊榮都賜給你。 \end{tabularx} \\ \\ \relax
2:38 & \begin{tabularx}{0.7\textwidth}{X} 世人和走獸，並天空的飛鳥，不論居住何處，他都交在你的手中，令你掌管這一切。你就是那金的頭。 \end{tabularx} \\ \\ \relax
2:39 & \begin{tabularx}{0.7\textwidth}{X} 在你以後必興起另一國，不及於你；又有第三國如銅，必掌管全地。 \end{tabularx} \\ \\ \relax
2:40 & \begin{tabularx}{0.7\textwidth}{X} 第四國必堅壯如鐵，就像鐵能打碎砸碎一切；鐵怎樣壓碎一切，那國也必照樣打碎壓碎。 \end{tabularx} \\ \\ \relax
2:41 & \begin{tabularx}{0.7\textwidth}{X} 你既看見像的腳和腳趾頭，一半是陶匠的泥，一半是鐵，那國將來也必分裂。你既看見鐵和泥攙雜，那國也必有鐵的力量。 \end{tabularx} \\ \\ \relax
2:42 & \begin{tabularx}{0.7\textwidth}{X} 那腳趾頭既是半鐵半泥，那國也必半強半弱。 \end{tabularx} \\ \\ \relax
2:43 & \begin{tabularx}{0.7\textwidth}{X} 你既看見鐵和泥攙雜，他們必有混雜的後裔，卻不能彼此相合，正如鐵和泥不能相合。 \end{tabularx} \\ \\ \relax
2:44 & \begin{tabularx}{0.7\textwidth}{X} 當諸王在位的時候，天上的神必另立一個永不敗壞的國度，這國度必不歸給其他百姓，卻要打碎滅絕所有的國度，存立到永遠。 \end{tabularx} \\ \\ \relax
2:45 & \begin{tabularx}{0.7\textwidth}{X} 你既看見非人手鑿出來的一塊石頭從山而出，打碎鐵、銅、泥、銀、金，那就是至大的神把將來要發生的事給王指明。這夢是確實的，這解釋也是準確的。」 \end{tabularx} \\ \\ \relax
2:46 & \begin{tabularx}{0.7\textwidth}{X} 當時，尼布甲尼撒王臉伏於地，向但以理下拜，並且吩咐人給他奉上供物和香。 \end{tabularx} \\ \\ \relax
2:47 & \begin{tabularx}{0.7\textwidth}{X} 王對但以理說：「你既能講明這奧祕，你們的神誠然是萬神之神、萬王之主，是奧祕的啟示者。」 \end{tabularx} \\ \\ \relax
2:48 & \begin{tabularx}{0.7\textwidth}{X} 於是王使但以理高升，賞賜他極多的禮物，派他管理巴比倫全省，又立他為總理，掌管巴比倫所有的智慧人。 \end{tabularx} \\ \\ \relax
2:49 & \begin{tabularx}{0.7\textwidth}{X} 但以理求王，王就派沙得拉、米煞、亞伯尼歌管理巴比倫省的事務，只是但以理仍在朝中侍立。 \end{tabularx} \\ \\
[1ex]
\hline
\hline
\end{longtable}
但以理書可說是舊約的啟示錄.在啟示錄中的記載,我們見到主之再來,神榮耀之國度降臨的榮耀；每次看到這卷書都使我們充滿希望,使我們信心堅定.
\newline
\newline
在啟示錄三百九十四節經文中,約有二百四十五節之多引用舊約.如果我們對舊約(特別是但以理書)沒有真正澈底的研究,就對啟示錄之了解仍未夠深刻；若多研究但以理書,就更能明白神的話語,更能體會神的心意.
\newline
\newline
神的話語非人所能說得完備,但神仍用如此偉大的真理啟示我們,教導屬祂的人；但願主寶血潔淨我們,使我們有屬靈的悟性,有屬靈的了解.神的國不是單在乎言語,更在乎聖靈的權能.神的話語解開,就會發出亮光,使我們能真正領受祂所給的恩典,有永遠的盼望.雖然世界敗壞,環境惡劣,我們又軟弱,但我們基督徒仍有活潑的盼望,我們不得不向神感謝.
\newline
\newline
但以理和啟示錄兩卷書,都是聖經啟示的信息.例如賽十三章及十四章,廿四至廿七章,以西結書一章及廿七至三十九章,撒迦利亞書九至十四,約珥書二及三章都是啟示的信息.
\newline
\newline
新約中除啟示錄外,普通書信有帖撒羅尼迦前後書,彼得前書一部份.四福音亦特別注重末世之信息,馬太廿四章,馬可十三章,路加廿一章皆為啟示之信息.
\newline
\newline
為什麼但以理書列在舊約的先知書,而其中信息內容與先知書不盡相同呢?中文聖經是根據希腊文舊約聖經的次序(亞歷山大的猶太人把舊約譯成希腊文後安排之目錄)我們當它是先知書,巴勒斯坦的聖經希伯來舊約將但以理書列在「著作」(即詩歌書及部分的歷史書),而不放在先知書中.因此,本書列在先知書之內,亦同時為啟示文學,因其中有特別之信息,比先知之信息更有深度.但以理是一位先知,先知乃神所差遣的,按當代之需要而傳合宜之信息.但以理是先知,先知是要站在時代的尖端；看到人所未看到的,傳出來也非人們所能易於明白的,傳出神的話語,有永存之價值.
\newline
\newline
但另一方面但以理與其他先知有所分別,因信息之內容是啟示文學.先知都秉承耶和華神命令,很直接,很清楚,並嚴厲的,帶責備的口吻指摘罪,叫人悔改.可是但以理的信息來源並沒有說明耶和華的吩咐,信息中沒有責備,也沒有嚴厲的口吻,卻有非常堅定的信心.今日,我們讀此卷書時,仍有不少的領受.常見其中有熱切的態度,並勸人悔改；神的恩典,神的復興,神的福份能引導我們,使我們能享受昔日神給以色列之榮耀.
\newline
\newline
在啟示文學中缺少樂觀的成份,我們覺得世上一切永無希望,因撒旦魔鬼之權勢使人感到失望,我們實在需要有榮耀的盼望.神之復興,似乎不能持久；直至有一天,世界和世界的情慾過去了,但遵行神旨意的人必永遠長存.因榮耀國度降臨時,那就真正在神前得著與主之恩惠與主福份.
\newline
\newline
在啟示文學中有許多象徵的預表,不容易明白及解釋,但求神幫助.正如科學家牛頓說:但以理書是聖經中最易明白的書卷,因其中不但將神的話說出,還藉天上的使者及地上之使者──神僕說給我們聽.並用各種不同的文字及正確的歷史攷據,使我們更易明白.願神在今次的培靈會中,藉聖靈光照我們,明白神的話語.
\newline
\newline
(但一1-4)「猶大王約雅敬在位第三年,巴比倫王尼布甲尼撒來到耶路撒冷,將城圍困；主將猶大王約雅敬,並神殿中器皿的幾分交付他手.他就把這器皿帶到示拿地,收入他神的廟裡,放在他神的庫中.王吩咐太監長亞施毗拿,從以色列人的宗室和貴胄中,帶個人來；就是年少沒有殘疾,相貌俊美,通達各樣學問,知識聰明具備；足能侍立在王宮裡的,要教他們迦勒底的文字言語.」
\newline
\newline
猶大王約雅敬在位第三年,據歷史攷據是主前六○五年或六○六年,巴比倫曾圍困擄掠猶大之耶路撒冷城共三次之多.第一次是六○五年,第二次是五九八年,第三次是五八七或五八六年.在第三次時,耶路撒冷才整個被毀,猶大完全敗亡；故但以理被巴比倫人擄去約在六○五年之時,他是當時的最優秀青年之一.當時神之殿受玷污,神的名受到羞辱.神是公義的神,以色列人悖逆了神(正如結廿四章中以色列人背了永約).以色列人是神的選民.為何神的公義會臨到他們?彼得前書告訴我們:神公義的審判要從神的家起首!今日教會需要復興,神的公義要表現出來.要你,我,教會──神的百姓,在神前悔改；若不悔改,神的審判便臨到,因神的審判要從神的家起首!
\newline
\newline
各位!今天我們研讀但以理書,必需要對時代有深刻的認識,現今是什麼世代?在此世代我們有何種地位?有何種功用?有何種責任?我們不單對時代有所認識,更要負起歷史上的責任.我們應知道,像但以理當時所認識神是統管萬有的神；祂是在歷史之中,也超乎歷史之上,在歷史中沒有任何一段事件,不是出於神的許可的.神更藉歷史表現祂救贖的工作,神要我們每個重生的人必須具有,時代的認識,歷史的責任.
\newline
\newline
當時但以理和他朋友之心靈感受,我們雖然未必完全明白；但知道他們都有許多的理想,有他們的抱負,對人生是充滿了希望與歡樂,多采多姿!可是當處在這種歷史裡面,他們就毫無辦法,不能不有一個嚴肅的態度.這種環境只有使他們悲觀,消極!因為國家敗亡了!他們失去了宗教的自由,又被擄到外國去；人生在這環境中受磨煉,神藉此造就他,潔淨他及揀選他,使用他.今日有許多人,其中包括了你和我,神也是如此將心思放在我們裡面；當然那時的歷史環境與現今我們所處的不同,但在另一方面看,亦有許多相同之處.西諺說:歷史如大磨石,磨坊中磨子不斷地轉動,轉得很慢卻磨得細.環境的磨練對一些人有很大的裨益,但對一些人可能有很大的危害.當時但以理及三位朋友在此遭遇中,當會有各種不同的感受及反應.也許有人能適應,但亦有人會厭惡而不肯順從環境；會在埋怨中順應環境,或是在無可奈何裡消極地去忍受.可是但以理及朋友都不是如此.他們是否識時務去順從環境嗎?還是以不滿之態度,強硬地抵擋嗎?或是消極地忍受去順從嗎?一切都不是,而是積極地承擔,無論如何都要高舉神!無論怎樣都堅持自己的信仰!無論什麼處境都要在外邦中作主的見證!使別人在他們身上看見神的榮耀；神的殿雖已毀壞,但在外邦仍可見到真正敬拜神的人.
\newline
\newline
當時巴比倫王之教育方法非常利害,他要使他們失去國家民族觀念和思想,要他們完全忘記自己到底是怎樣的人.在新約就並不注重這問題,但在舊約很重視；神為何使一些人成為一國,成為一民族然後才賜福他們.因為神所選的這批人,要成為傳福音的使者.神要使這批人從那時國家開始,一直成為福分傳遞.神要他們明白個人與團體間之關係,現今的人對這問題的認識很淺薄,對國家,對團體之精神亦不注意.許多人以為自己在教會可毫無責任,其實神為何設立教會?因神不單是賜福個人,亦藉個人之賜福而至教會,而至整個團體,神要使我們有團體之意識.這是舊約所注意的一種群體之精神.正如亞干犯罪,令全體受害,個人亦要向團體負責.個人在神前領受神的恩,神亦因此而賜福整個團體.我們需要與信徒彼此相交,相通之關係,需要神賜福我們的教會,我們的家庭及我們自己.神不但要我們向神負責,更要向教會,向屬靈團體,向弟兄姊妹們負責.但以理及朋友非常清楚的明白,可是照這環境方面來說,神甚至把給先知的權力都剝奪了,他們失去了他們的身份.神之兒女,神的子民,猶大國的百姓,他們身份都失去了.有許多基督徒在教會非常熱心,但當去到另一教會或海外時,慢慢就冷淡起來；因他們沒注意到自己的身份,他們的身份只是在一個小小的教會內；一個青年團契內,一個婦女會,一個很小的範圍內.失去身份的人是十分容易失敗的,但以理及三位朋友已失去他們應有的身份,那是最危險的時候,可是他們失敗了嗎?巴比倫王要使他們的歸屬感失去,要他們忘記自己是猶大人,忘記自己是神的選民；要他們對猶大人及以色列民族完全遺忘,毫無關係,歸屬感完全歸入巴比倫裡.許多人常說國家民族不甚要緊,只要我們在這裡安全就好了.然而一個人應有國家民族之觀念,不可只為獲得外國之國籍而沾沾自喜,其實你還不是中國人嗎?神要我們在這時代作中國人,該為此而感到驕傲.神給我們如此好的道德文化,要我們以宗教屬靈的真理來貫通.現今但以理及朋友是否因自己在巴比倫王宮中生活而驕傲呢?他們是在玩樂嗎?不,他們在此環境中,必需要照著神所賜的恩典,應發生作用,應有影響,應有見證.讓我們今日能在神前成為多方的見證人,巴比倫王不單要他們忘記自己的種族對宗教輕忽!連名字也要改為巴比倫的名字.可是他們沒忘記自己的身份,始終堅持自己的信仰；雖然當時巴比倫王對他們十分優待,吃王的膳和酒,有最舒服的環境,最遠大的前程,可是他們最尊貴的身份沒有失去.當一個人或幾個人在十分可怕的逼迫之下,為主而站立得穩是很困難的.但最困難的是勝過魔鬼「優先」的方法──最容易失敗跌倒,但以理及三位朋友卻立定心志決不玷污自己.因為要順服神,順服神的律法及命令,不肯失去自己身份,不肯忘卻自己的歸屬感；不肯向此世代低頭,甚至極微小的事,所以神豐盛之恩典使他們站立得穩.如三友在火yiu中,但以理在獅子坑,這一切的逼迫,他們都站立得穩；因他們從開始時就在最微小的事上,都始終為主而站立著.堅定之信心是一種拒絕,但卻不用強硬的抵抗只用溫柔的方法,正如求太監容他不玷污自己；用「求」的方法,客氣溫和有禮貌,這是屬靈的方法.這使太監等都因此而答應了,且不敢輕看他們.如真有決心,屬靈的工作才有果效.如有好的屬靈氣質,不但靈性身體都能作神的見證,每一方面都是神的見證人.
\newline
\newline
信心是要受磨煉的.一方面信心有確據,但另一方面並沒有,但以理不知將來會如何；可是卻知道神會幫助他,結果,但以理的智慧比巴比倫的術士們超過十倍.我們什麼都不必懼怕,最要緊的是需要祈禱,單是思想而不祈禱是不行的.今日之青年在學校讀書,得到許多知識；但知識是一件事,信仰又是另一件事.求主憐憫我們,能思想亦會禱告；在第一章中可見他常禱告,二章及九章也說及禱告.
\newline
\newline
到古列王元年但以理還在(但一:21),從巴比倫到波斯王朝,他數十年如一日.「但以理至今仍在」,他的見證,他的影響還在.他那為神發光,將神所傳的真正話語之能力還在,因他有屬靈的持久力.在最惡劣時代,他不單是不隨波逐流,反而成為中流砥柱.不單是成為神的器皿,充滿而又倒空；更成活水的江河,在河床上水流不絕.在世界的潮流中,不單不受影響；反而有新的活水流出,永流不息.愈流愈深,愈流愈廣,愈遠,愈長,直至今日我們仍在領受著.
\newline
\newline


\newpage

\section{第45屆~港九培靈會~唐佑之~第二講}
\label{sec:1247}
\textbf{夢幻或啟示}
\newline
\newline
連結: \href{https://www.hkbibleconference.org/session-message/view/1247}{\texttt{https://www.hkbibleconference.org/session-message/view/1247}}
\newline
\newline
\hyperref[sec:1246]{< < < PREV SERMON < < <}
~
\hyperref[sec:toc]{[返目錄]}
~
\hyperref[sec:1248]{> > > NEXT SERMON > > >}
\newline
\newline
但以理書 2:1-2:49
\newline
\begin{longtable}{cl}
\hline
\hline
章節 & 經文 (和合本修訂版)\\
\hline
2:1 & \begin{tabularx}{0.7\textwidth}{X} 尼布甲尼撒在位第二年，他做了很多夢，心裡煩亂，不能睡覺。 \end{tabularx} \\ \\ \relax
2:2 & \begin{tabularx}{0.7\textwidth}{X} 王吩咐人將術士、巫師、行邪術的和迦勒底人召來，要他們把王的夢告訴王；他們就來，站在王面前。 \end{tabularx} \\ \\ \relax
2:3 & \begin{tabularx}{0.7\textwidth}{X} 王對他們說：「我做了一個夢，心裡煩亂，想要知道這是甚麼夢。」 \end{tabularx} \\ \\ \relax
2:4 & \begin{tabularx}{0.7\textwidth}{X} 迦勒底人用亞蘭話對王說：「願王萬歲！請將夢告訴僕人，我們就可以講解。」 \end{tabularx} \\ \\ \relax
2:5 & \begin{tabularx}{0.7\textwidth}{X} 王回答迦勒底人說：「這事我已決定，你們若不把夢和夢的解釋告訴我，就必被凌遲，你們的房屋必成糞堆； \end{tabularx} \\ \\ \relax
2:6 & \begin{tabularx}{0.7\textwidth}{X} 但你們若能說出這個夢和夢的解釋，就必從我得到禮物、賞賜和殊榮。現在，你們要把夢和夢的解釋告訴我。」 \end{tabularx} \\ \\ \relax
2:7 & \begin{tabularx}{0.7\textwidth}{X} 他們再一次回答說：「請王將夢告訴僕人，我們就可以講解。」 \end{tabularx} \\ \\ \relax
2:8 & \begin{tabularx}{0.7\textwidth}{X} 王回答說：「我確實知道你們是故意拖延，因為你們知道這事我已決定。 \end{tabularx} \\ \\ \relax
2:9 & \begin{tabularx}{0.7\textwidth}{X} 你們若不將夢告訴我，只有一個辦法對待你們；因為你們彼此串通，向我胡言亂語，要等候情勢改變。現在，你們要將夢告訴我，讓我知道你們真能為我解夢。」 \end{tabularx} \\ \\ \relax
2:10 & \begin{tabularx}{0.7\textwidth}{X} 迦勒底人回答王說：「世上沒有人能解釋王的事情；從來沒有君王、大臣、掌權者向術士、巫師，或迦勒底人問過這樣的事。 \end{tabularx} \\ \\ \relax
2:11 & \begin{tabularx}{0.7\textwidth}{X} 王所問的事很難，除了不與血肉之軀同住的神，沒有人能在王面前解釋。」 \end{tabularx} \\ \\ \relax
2:12 & \begin{tabularx}{0.7\textwidth}{X} 王因這事生氣，大大震怒，吩咐滅絕巴比倫所有的智慧人。 \end{tabularx} \\ \\ \relax
2:13 & \begin{tabularx}{0.7\textwidth}{X} 命令發出，智慧人將要被殺，人就尋找但以理和他的同伴，要殺他們。 \end{tabularx} \\ \\ \relax
2:14 & \begin{tabularx}{0.7\textwidth}{X} 王的護衛長亞略奉命去殺巴比倫的智慧人，但以理用婉言和智慧回應， \end{tabularx} \\ \\ \relax
2:15 & \begin{tabularx}{0.7\textwidth}{X} 向王的大臣亞略說：「王的命令為何這樣緊急呢？」亞略就把事情告訴但以理。 \end{tabularx} \\ \\ \relax
2:16 & \begin{tabularx}{0.7\textwidth}{X} 於是但以理進去求王寬限，好為王解夢。 \end{tabularx} \\ \\ \relax
2:17 & \begin{tabularx}{0.7\textwidth}{X} 但以理回到他的居所，把這事告訴他的同伴哈拿尼雅、米沙利、亞撒利雅， \end{tabularx} \\ \\ \relax
2:18 & \begin{tabularx}{0.7\textwidth}{X} 要他們祈求天上的神施憐憫，將這奧祕指明，免得但以理和他的同伴與巴比倫其餘的智慧人一同滅亡。 \end{tabularx} \\ \\ \relax
2:19 & \begin{tabularx}{0.7\textwidth}{X} 這奧祕就在夜間異象中顯明給但以理，但以理就稱頌天上的神。 \end{tabularx} \\ \\ \relax
2:20 & \begin{tabularx}{0.7\textwidth}{X} 但以理說：「神的名是應當稱頌的，從亙古直到永遠！因為智慧和能力都屬乎他。 \end{tabularx} \\ \\ \relax
2:21 & \begin{tabularx}{0.7\textwidth}{X} 他改變時間、季節，他廢王，立王；將智慧賜給智慧人，將知識賜給聰明人。 \end{tabularx} \\ \\ \relax
2:22 & \begin{tabularx}{0.7\textwidth}{X} 他顯明深奧隱祕的事，洞悉幽暗中的一切，光明也與他同住。 \end{tabularx} \\ \\ \relax
2:23 & \begin{tabularx}{0.7\textwidth}{X} 我列祖的神啊，我感謝你，讚美你，因你將智慧才能賜給我，我們所求問的現在你已指明給我，把王的事給我們指明。」 \end{tabularx} \\ \\ \relax
2:24 & \begin{tabularx}{0.7\textwidth}{X} 於是，但以理進到王所派滅絕巴比倫智慧人的亞略那裡去，對他這樣說：「不要滅絕巴比倫的智慧人，求你領我到王面前，我可以為王解夢。」 \end{tabularx} \\ \\ \relax
2:25 & \begin{tabularx}{0.7\textwidth}{X} 亞略就急忙領但以理到王面前，對王這樣說：「我在被擄的猶大人中找到一人，能將夢的解釋告訴王。」 \end{tabularx} \\ \\ \relax
2:26 & \begin{tabularx}{0.7\textwidth}{X} 王對那稱為伯提沙撒的但以理說：「你能將我所做的夢和夢的解釋告訴我嗎？」 \end{tabularx} \\ \\ \relax
2:27 & \begin{tabularx}{0.7\textwidth}{X} 但以理回答王說：「王所問的那奧祕，智慧人、巫師、術士、觀兆的都不能告訴王， \end{tabularx} \\ \\ \relax
2:28 & \begin{tabularx}{0.7\textwidth}{X} 只有那在天上的神能顯明奧祕。他已把日後將要發生的事指示尼布甲尼撒王。你在床上做的夢和你腦中的異象是這樣： \end{tabularx} \\ \\ \relax
2:29 & \begin{tabularx}{0.7\textwidth}{X} 你，王啊，你在床上所思想的是關乎日後的事，那顯明奧祕的主已把將來要發生的事指示你。 \end{tabularx} \\ \\ \relax
2:30 & \begin{tabularx}{0.7\textwidth}{X} 至於我，那奧祕顯明給我，並非因我智慧勝過一切活著的人，而是為了讓王知道夢的解釋，知道你心裡的意念。 \end{tabularx} \\ \\ \relax
2:31 & \begin{tabularx}{0.7\textwidth}{X} 「你，王啊，你正觀看，看哪，有一個很大的像，這像甚高，極其光耀，立在你面前，形狀非常可怕。 \end{tabularx} \\ \\ \relax
2:32 & \begin{tabularx}{0.7\textwidth}{X} 這像的頭是純金的，胸膛和膀臂是銀的，腹部和腰是銅的， \end{tabularx} \\ \\ \relax
2:33 & \begin{tabularx}{0.7\textwidth}{X} 腿是鐵的，腳是半鐵半泥的。 \end{tabularx} \\ \\ \relax
2:34 & \begin{tabularx}{0.7\textwidth}{X} 你正觀看，見有一塊非人手鑿出來的石頭打在它半鐵半泥的腳上，把腳砸碎； \end{tabularx} \\ \\ \relax
2:35 & \begin{tabularx}{0.7\textwidth}{X} 於是鐵、泥、銅、銀、金都一同砸得粉碎，如夏天禾場上的糠秕，被風吹散，無處可尋。打碎這像的石頭成了一座大山，覆蓋全地。 \end{tabularx} \\ \\ \relax
2:36 & \begin{tabularx}{0.7\textwidth}{X} 「這就是那夢；我們要在王面前講解那夢。 \end{tabularx} \\ \\ \relax
2:37 & \begin{tabularx}{0.7\textwidth}{X} 你，王啊，你是諸王之王。天上的神已將國度、權勢、能力、尊榮都賜給你。 \end{tabularx} \\ \\ \relax
2:38 & \begin{tabularx}{0.7\textwidth}{X} 世人和走獸，並天空的飛鳥，不論居住何處，他都交在你的手中，令你掌管這一切。你就是那金的頭。 \end{tabularx} \\ \\ \relax
2:39 & \begin{tabularx}{0.7\textwidth}{X} 在你以後必興起另一國，不及於你；又有第三國如銅，必掌管全地。 \end{tabularx} \\ \\ \relax
2:40 & \begin{tabularx}{0.7\textwidth}{X} 第四國必堅壯如鐵，就像鐵能打碎砸碎一切；鐵怎樣壓碎一切，那國也必照樣打碎壓碎。 \end{tabularx} \\ \\ \relax
2:41 & \begin{tabularx}{0.7\textwidth}{X} 你既看見像的腳和腳趾頭，一半是陶匠的泥，一半是鐵，那國將來也必分裂。你既看見鐵和泥攙雜，那國也必有鐵的力量。 \end{tabularx} \\ \\ \relax
2:42 & \begin{tabularx}{0.7\textwidth}{X} 那腳趾頭既是半鐵半泥，那國也必半強半弱。 \end{tabularx} \\ \\ \relax
2:43 & \begin{tabularx}{0.7\textwidth}{X} 你既看見鐵和泥攙雜，他們必有混雜的後裔，卻不能彼此相合，正如鐵和泥不能相合。 \end{tabularx} \\ \\ \relax
2:44 & \begin{tabularx}{0.7\textwidth}{X} 當諸王在位的時候，天上的神必另立一個永不敗壞的國度，這國度必不歸給其他百姓，卻要打碎滅絕所有的國度，存立到永遠。 \end{tabularx} \\ \\ \relax
2:45 & \begin{tabularx}{0.7\textwidth}{X} 你既看見非人手鑿出來的一塊石頭從山而出，打碎鐵、銅、泥、銀、金，那就是至大的神把將來要發生的事給王指明。這夢是確實的，這解釋也是準確的。」 \end{tabularx} \\ \\ \relax
2:46 & \begin{tabularx}{0.7\textwidth}{X} 當時，尼布甲尼撒王臉伏於地，向但以理下拜，並且吩咐人給他奉上供物和香。 \end{tabularx} \\ \\ \relax
2:47 & \begin{tabularx}{0.7\textwidth}{X} 王對但以理說：「你既能講明這奧祕，你們的神誠然是萬神之神、萬王之主，是奧祕的啟示者。」 \end{tabularx} \\ \\ \relax
2:48 & \begin{tabularx}{0.7\textwidth}{X} 於是王使但以理高升，賞賜他極多的禮物，派他管理巴比倫全省，又立他為總理，掌管巴比倫所有的智慧人。 \end{tabularx} \\ \\ \relax
2:49 & \begin{tabularx}{0.7\textwidth}{X} 但以理求王，王就派沙得拉、米煞、亞伯尼歌管理巴比倫省的事務，只是但以理仍在朝中侍立。 \end{tabularx} \\ \\
[1ex]
\hline
\hline
\end{longtable}
但二1-16,神乃歷史之主.在歷史進行中,神的手在歷史的巨輪上,當歷史在進行時,神要彰顯祂的大能；神是在歷史之中,也在歷史之上.藉歷史將祂的救贖心意表明,當我們讀舊約時,就可發現神是藉歷史將祂自己表現出來.
\newline
\newline
歷史乃神之啟示,歷史是動的而不是靜的,神藉歷史的每一細節要彰顯表明祂偉大審判的公義作為.故歷史不單說明啟示,也是一直向前的,有目的的.雖有人說歷史是重複的,如一巨輪循環著,但另一方面是有方向的,直到歷史結束.
\newline
\newline
在此章中,可見歷史之屬靈哲學.當我們研究此古代歷史時,見先知的預言,整幅歷史圖畫,見神所有歷史中之心意.雖然一般人注意歷史之事實,這是縱面而非橫面.「縱」是從一個年代到另一個年代,但以理為王解夢,從「縱」的方面看此像,是從頭到腳要一件一件的發生.從「橫」的方面看,則還有另一種力量,那就是一塊石頭扔過來時；就見到另一力量,這並非人之力量,乃是神的能力.神更藉著歷史彰顯祂的作為,但也藉著這橫的力量作屬靈的解釋.
\newline
\newline
所以,從縱的力量可見歷史之事實,從橫的力量可見歷史之意義.我們試看當時之背景,尼布甲尼撒在位第二年,即但以理正在受訓時──這問題會引起許多經學家的問題.其實計算曆法不一樣,因巴比倫之算法那二年是二周年,乃有歷史之根據.但以理是在主前六○五年夏被擄,同年九月七日尼布甲尼撒之父死,於是尼布甲尼撒開始作王.故在六○五至六○四年之尼散月──即今之四五月──,可算是但以理受訓之年月.但另一方面,六○四至六○三年,乃是受訓第二年,在猶大人算是三年,巴比倫算是二年.可見第二章之背景,乃當但以理受訓完之後,尼布甲尼撒作王二週年之時所發生的一件大事.尼布甲尼撒王作夢後卻不知道夢的內容.現在的問題是:(一)到底此夢是否神之啟示?我們可說不是的,因為一個信奉異教的外邦人,一個有罪的人,神怎會向他啟示呢?這不過是野心家的夢境,因許多的思想顯明政治的野心.一個暴虐之君王,充滿了許多夢想,憂慮和煩惱,故在白天所思想的,晚上就在夢境出現.因這種夢實在太雜亂,所以夢醒時就不知所夢的是什麼?甚至忘卻,心裡就非常不安.那為什麼又有神的心意在其中?神讓他的心覺得煩雜.(二)這樣的人就需要神的僕人,需要神的信息,需要神的警戒的信息,公義的信息.今日世界上多少人心裡正在覺得煩惱,他們不知自己的夢想,幻想和理想是什麼.他們的人生好像夢幻的人生,似夢非夢,人生在極端的困惑與迷惘中；他們需要我們這些知道神心意,明白神話語的人.(三)人需要從神那裡得著能力與智慧,因為沒有神的能力與智慧,就無法知道神的心意更沒有神的信息.尼王作夢並非神給他啟示,所以有了此夢後,心裡就很煩惱.這裡我們可以見到幾種的人:
\newline
\newline
(一)尼布甲尼撒乃暴虐的罪人──當時尼王召巴比倫術士解夢之人,他們乃虛妄的弱者.在某一狀況下就振振有詞,頭頭是道,可是卻毫無內容；充滿虛妄,在此特殊情形之下,他們就表現出懦弱無用的樣子.因他們沒有神,他們完全是迷信.今日的世代不也是如此嗎?不是充滿了虛妄懦弱和兇惡強暴的的人嗎?這世代何等需要但以理和他的同伴,需要有同樣的信仰的你和我!需要神的工人,需要有屬靈智慧的敬虔者,需要忠心願被神使用的工人,和那些肯真正把自己完全擺上的神的兒女.但我們知道那兇暴的罪人,要把所有的術士們殺掉,因為罪人落在罪人的手中是可怕的；但義人是否也落在罪人的手裡?在外表會有這危險,甚至但以理和他的同伴,也幾乎被殺.可是神有永遠保守的能力,那就是說屬神的人是否有此信心,是否信心使之鎮定?但以理是有信心的人,據經上記載,他不但告訴王的護衛長,更到王前說他能解夢,當王心情不好時,任何事都會做出來,也許連但以理也殺掉；因但以理不過是個被擄者,在君王眼中是個奴隸,根本不值得重視；人不重視他,神卻重視他.各位!你我都是神的兒女,神保護我們好像保護眼中的瞳人.祂看重我們,因我們屬於祂!但以理有此信心,故他雄赳赳,氣昂昂；光明磊落地站在君王前,宣告他能解夢,他的生命在神手中,因為但以理有堅強的信念,他不懼怕.有一些人會想,他這樣會有危險；可是神的兒女有此能力.不單不逃避,反而有堅強的信心面對現實.不僅把壞的變好,也把惡劣的改善,這是神蹟,因神要在歷史中彰顯祂的作為,隨時改變現狀,但以理就將整個情勢挽救過來.
\newline
\newline
在聖經上常有這種情形,如果人沒有在神前真真得勝,他不能逃脫神的忿怒；他只不過是在神的審判火中,抽出的一根柴,僅僅的得救,抽出後只能放在一邊.有人在神前的態度就如在危險的火中被人將他救出來,可是卻沾沾自喜,覺得很滿足,但這不是你我作神兒女的態度.因我們並非在神審判火中抽出的一根柴,而是可經過水火,進入豐富之地.又知環境的火如試煉般臨到之時,不是僅僅抽出的一根柴；也不讓火滅掉,而把火成為福音火炬；高高舉起,使多人看見.但以理就有這樣的態度,因他知道仰望神,並向祂禱告.
\newline
\newline
(二)但以理乃禱告的人(但二17-18):「但以理回到他的居所,將這事告訴他的同伴哈拿尼雅,米沙利,亞撒利雅；要他們祈求天上的神施憐憫,將這奧秘的事指明,免得但以理和他的同伴,與巴比倫其餘的哲士,一同滅亡.」但以理是禱告的人,我們也禱告,可是我們的禱告生活如何?但以理知道在此情形下必需要禱告,他回去和與三個朋友一同禱告.該知道當我們遭遇困難時,並非開會商議對策；不是在研究實際情勢,而是在神前好好地同心合意地禱告.那是何等的寶貴!主在馬太十八章19節說:「地上有兩三個人同心合意奉祂的名禱告,祂必定垂聽.」這裡的涵義可用馬太十八章中的「同心合意」來說明:這好似「交響樂」 Symphony,樂器不同,但是奏出來的卻是和諧的.為什麼我們的禱告會會冷落而沉悶呢!實在因為個人在神前沒有好好的禱告,如果每人先有好的禱告,那麼一起禱告才有功效.有時我們不想參加禱告會,以為這禱告會不好,其實是我們自己本身有問題,願神憐憫我們先自己肯好好地禱告.
\newline
\newline
看這個儆醒禱告的人:(彼前四7)「萬物的結局近了,所以你們要謹慎自守,儆醒禱告.」但以理和他三個朋友真的全交主手裡,專心地禱告.哥林多前書七5「要專心禱告」,四2「恆切禱告」,歌羅西四12「竭力的祈求」,這樣禱告才有功效.當他們禱告完畢,心中充滿平安,這幾個都是禱告的人,當他們一起禱告時則功效更大,雅各書五:16「義人祈禱所發的力量,是大有功效的.」一個義人禱告有功效,何況四個呢!神在他們心中充滿喜樂,於是但以理就安然睡覺.
\newline
\newline
那時神指示他明白夜間的異象,神把祂的心意說出.今日你和我都需要如此,基督徒沒有平安常有雜夢,與世人一樣,那是因為缺少禱告.我們沒有主的平安在心中,故當事情發生時；就十分忙亂,東奔西跑,不肯安靜起來在神前禱告,直至沒辦法時才在對神說:「人的盡頭是你的開頭,求你助我!」可是心中仍充滿了雜亂,這怎能得到神的安息呢!「人的盡頭乃神之開頭」,這話是不錯的,因人在活動時神不能工作.人的辦法太多了,神的辦法就顯不出來.但我們為何要等到人盡頭時才讓神來開頭呢?神是否喜歡我們在盡頭時才讓我們依靠祂?不是的,神不喜歡我們屬肉體的工作,和我們自己的方法.神也不喜歡我們忙亂.主說:我們勞苦擔重擔的可以到祂面前.主呼喊我們:來吧,快來吧!不要到別的地方去,不要再想別的方法；因這是唯一的途徑,到主前可得安息.在神前真得安息,神願我們開始時就仰望祂.但以理與他三友始終一的仰望神,仰望那信心創始成終的主；讓我們不要到自己的盡頭,才給神開頭.該神讓在我們身上掌王權,今日世代需要有禱告之人,像但以理他們一樣.因現世代尼布甲尼撒王太多了,他們好像有財力,勢力,但卻都是罪人.他們用許多暴力及不正常方法,心中沒有平安；他們需要我們的幫助,我們有沒有去幫助他們嗎?我們敢去嗎?我們是否因他們有財有勢而覺自卑?不敢把福音傳出?不傳福音的人有禍了!為何我們要自卑!我們有萬王之王.人們需要幫助,他們在尋找卻尋不見,心中沒有平安,我們敢傳福音給他們嗎?有信息傳給他們嗎?我們沒有能力,沒有見證,也沒有智慧!該知道智慧是從神而來,能力也從神而來!當我們見但以理在歌頌耶和華神時,他說:「神的名是應當稱頌的,從亙古直到永遠,因為智慧能力都屬乎祂.」(但二20)
\newline
\newline
許多時候我們說了許多話,可惜都是廢話,我們應說些合宜的話.在教會中也很少有屬靈的話語,如果我們真被神的恩典充滿,就會有屬靈的話語發出；可惜我們與外邦人異教徒無異!去探訪時你能說些什麼話呢?和世人一般嗎?這是我們主要的失敗,這失敗成為莫大的損失,這世界就因我們的軟弱遭更大的損失.
\newline
\newline
我們沒有禱告,所以就失敗,沒有禱告的事奉也就沒有功效,於是感到工作困難重重!試看舊約先知約拿,他違背神,不肯到尼尼微傳道而到他施去,於是風浪大作.船上的水手是信假神的迷信者,他們的禱告無功效,但他們仍知求告他們的神；可是神的工人,神的兒女在何處呢!他──約拿──卻在艙中睡覺!當這世界充滿了危險恐懼時,在動蕩不安時,你我都還在睡覺嗎?教會沒有禱告,神的工人,神的兒女沒有禱告,世界還有什麼希望呢!當時在風浪中的船又如何?反而是船主叫約拿起來.他說:「你這沉睡的人阿!為何這樣呢?起來快求告你的神.」各位!你是否聽見這聲音?這是世界發出的聲音,叫我們這些沉睡的人快起來.假如但以理和他三友都是沉睡的人,那早已被殺了!幸而他們是儆醒禱告的人.我們多麼需要神的憐憫,多麼需要禱告；今日最缺乏的是禱告,我們都明白這個真理,都知道這屬靈的原則,但是就不肯禱告.
\newline
\newline
今日社會為什麼如此混亂不安,我們卻好像是個旁觀者；毫無關係似的,其實是有密切的關係.今日教會與社會的關係只有兩個可能,如果不是教會影響社會的話；那就是社會影響了教會.今日教會之見證何在?立場何在?你我到底作了些什麼?我們需要禱告.當我們禱告時,神就會大大施恩,因祂要復興祂的作為,復興祂的教會,復興祂的兒女!
\newline
\newline


\newpage

\section{第45屆~港九培靈會~唐佑之~第三講}
\label{sec:1248}
\textbf{承受與擺脫}
\newline
\newline
連結: \href{https://www.hkbibleconference.org/session-message/view/1248}{\texttt{https://www.hkbibleconference.org/session-message/view/1248}}
\newline
\newline
\hyperref[sec:1247]{< < < PREV SERMON < < <}
~
\hyperref[sec:toc]{[返目錄]}
~
\hyperref[sec:1250]{> > > NEXT SERMON > > >}
\newline
\newline
但以理書 2:1-3:30
\newline
\begin{longtable}{cl}
\hline
\hline
章節 & 經文 (和合本修訂版)\\
\hline
2:1 & \begin{tabularx}{0.7\textwidth}{X} 尼布甲尼撒在位第二年，他做了很多夢，心裡煩亂，不能睡覺。 \end{tabularx} \\ \\ \relax
2:2 & \begin{tabularx}{0.7\textwidth}{X} 王吩咐人將術士、巫師、行邪術的和迦勒底人召來，要他們把王的夢告訴王；他們就來，站在王面前。 \end{tabularx} \\ \\ \relax
2:3 & \begin{tabularx}{0.7\textwidth}{X} 王對他們說：「我做了一個夢，心裡煩亂，想要知道這是甚麼夢。」 \end{tabularx} \\ \\ \relax
2:4 & \begin{tabularx}{0.7\textwidth}{X} 迦勒底人用亞蘭話對王說：「願王萬歲！請將夢告訴僕人，我們就可以講解。」 \end{tabularx} \\ \\ \relax
2:5 & \begin{tabularx}{0.7\textwidth}{X} 王回答迦勒底人說：「這事我已決定，你們若不把夢和夢的解釋告訴我，就必被凌遲，你們的房屋必成糞堆； \end{tabularx} \\ \\ \relax
2:6 & \begin{tabularx}{0.7\textwidth}{X} 但你們若能說出這個夢和夢的解釋，就必從我得到禮物、賞賜和殊榮。現在，你們要把夢和夢的解釋告訴我。」 \end{tabularx} \\ \\ \relax
2:7 & \begin{tabularx}{0.7\textwidth}{X} 他們再一次回答說：「請王將夢告訴僕人，我們就可以講解。」 \end{tabularx} \\ \\ \relax
2:8 & \begin{tabularx}{0.7\textwidth}{X} 王回答說：「我確實知道你們是故意拖延，因為你們知道這事我已決定。 \end{tabularx} \\ \\ \relax
2:9 & \begin{tabularx}{0.7\textwidth}{X} 你們若不將夢告訴我，只有一個辦法對待你們；因為你們彼此串通，向我胡言亂語，要等候情勢改變。現在，你們要將夢告訴我，讓我知道你們真能為我解夢。」 \end{tabularx} \\ \\ \relax
2:10 & \begin{tabularx}{0.7\textwidth}{X} 迦勒底人回答王說：「世上沒有人能解釋王的事情；從來沒有君王、大臣、掌權者向術士、巫師，或迦勒底人問過這樣的事。 \end{tabularx} \\ \\ \relax
2:11 & \begin{tabularx}{0.7\textwidth}{X} 王所問的事很難，除了不與血肉之軀同住的神，沒有人能在王面前解釋。」 \end{tabularx} \\ \\ \relax
2:12 & \begin{tabularx}{0.7\textwidth}{X} 王因這事生氣，大大震怒，吩咐滅絕巴比倫所有的智慧人。 \end{tabularx} \\ \\ \relax
2:13 & \begin{tabularx}{0.7\textwidth}{X} 命令發出，智慧人將要被殺，人就尋找但以理和他的同伴，要殺他們。 \end{tabularx} \\ \\ \relax
2:14 & \begin{tabularx}{0.7\textwidth}{X} 王的護衛長亞略奉命去殺巴比倫的智慧人，但以理用婉言和智慧回應， \end{tabularx} \\ \\ \relax
2:15 & \begin{tabularx}{0.7\textwidth}{X} 向王的大臣亞略說：「王的命令為何這樣緊急呢？」亞略就把事情告訴但以理。 \end{tabularx} \\ \\ \relax
2:16 & \begin{tabularx}{0.7\textwidth}{X} 於是但以理進去求王寬限，好為王解夢。 \end{tabularx} \\ \\ \relax
2:17 & \begin{tabularx}{0.7\textwidth}{X} 但以理回到他的居所，把這事告訴他的同伴哈拿尼雅、米沙利、亞撒利雅， \end{tabularx} \\ \\ \relax
2:18 & \begin{tabularx}{0.7\textwidth}{X} 要他們祈求天上的神施憐憫，將這奧祕指明，免得但以理和他的同伴與巴比倫其餘的智慧人一同滅亡。 \end{tabularx} \\ \\ \relax
2:19 & \begin{tabularx}{0.7\textwidth}{X} 這奧祕就在夜間異象中顯明給但以理，但以理就稱頌天上的神。 \end{tabularx} \\ \\ \relax
2:20 & \begin{tabularx}{0.7\textwidth}{X} 但以理說：「神的名是應當稱頌的，從亙古直到永遠！因為智慧和能力都屬乎他。 \end{tabularx} \\ \\ \relax
2:21 & \begin{tabularx}{0.7\textwidth}{X} 他改變時間、季節，他廢王，立王；將智慧賜給智慧人，將知識賜給聰明人。 \end{tabularx} \\ \\ \relax
2:22 & \begin{tabularx}{0.7\textwidth}{X} 他顯明深奧隱祕的事，洞悉幽暗中的一切，光明也與他同住。 \end{tabularx} \\ \\ \relax
2:23 & \begin{tabularx}{0.7\textwidth}{X} 我列祖的神啊，我感謝你，讚美你，因你將智慧才能賜給我，我們所求問的現在你已指明給我，把王的事給我們指明。」 \end{tabularx} \\ \\ \relax
2:24 & \begin{tabularx}{0.7\textwidth}{X} 於是，但以理進到王所派滅絕巴比倫智慧人的亞略那裡去，對他這樣說：「不要滅絕巴比倫的智慧人，求你領我到王面前，我可以為王解夢。」 \end{tabularx} \\ \\ \relax
2:25 & \begin{tabularx}{0.7\textwidth}{X} 亞略就急忙領但以理到王面前，對王這樣說：「我在被擄的猶大人中找到一人，能將夢的解釋告訴王。」 \end{tabularx} \\ \\ \relax
2:26 & \begin{tabularx}{0.7\textwidth}{X} 王對那稱為伯提沙撒的但以理說：「你能將我所做的夢和夢的解釋告訴我嗎？」 \end{tabularx} \\ \\ \relax
2:27 & \begin{tabularx}{0.7\textwidth}{X} 但以理回答王說：「王所問的那奧祕，智慧人、巫師、術士、觀兆的都不能告訴王， \end{tabularx} \\ \\ \relax
2:28 & \begin{tabularx}{0.7\textwidth}{X} 只有那在天上的神能顯明奧祕。他已把日後將要發生的事指示尼布甲尼撒王。你在床上做的夢和你腦中的異象是這樣： \end{tabularx} \\ \\ \relax
2:29 & \begin{tabularx}{0.7\textwidth}{X} 你，王啊，你在床上所思想的是關乎日後的事，那顯明奧祕的主已把將來要發生的事指示你。 \end{tabularx} \\ \\ \relax
2:30 & \begin{tabularx}{0.7\textwidth}{X} 至於我，那奧祕顯明給我，並非因我智慧勝過一切活著的人，而是為了讓王知道夢的解釋，知道你心裡的意念。 \end{tabularx} \\ \\ \relax
2:31 & \begin{tabularx}{0.7\textwidth}{X} 「你，王啊，你正觀看，看哪，有一個很大的像，這像甚高，極其光耀，立在你面前，形狀非常可怕。 \end{tabularx} \\ \\ \relax
2:32 & \begin{tabularx}{0.7\textwidth}{X} 這像的頭是純金的，胸膛和膀臂是銀的，腹部和腰是銅的， \end{tabularx} \\ \\ \relax
2:33 & \begin{tabularx}{0.7\textwidth}{X} 腿是鐵的，腳是半鐵半泥的。 \end{tabularx} \\ \\ \relax
2:34 & \begin{tabularx}{0.7\textwidth}{X} 你正觀看，見有一塊非人手鑿出來的石頭打在它半鐵半泥的腳上，把腳砸碎； \end{tabularx} \\ \\ \relax
2:35 & \begin{tabularx}{0.7\textwidth}{X} 於是鐵、泥、銅、銀、金都一同砸得粉碎，如夏天禾場上的糠秕，被風吹散，無處可尋。打碎這像的石頭成了一座大山，覆蓋全地。 \end{tabularx} \\ \\ \relax
2:36 & \begin{tabularx}{0.7\textwidth}{X} 「這就是那夢；我們要在王面前講解那夢。 \end{tabularx} \\ \\ \relax
2:37 & \begin{tabularx}{0.7\textwidth}{X} 你，王啊，你是諸王之王。天上的神已將國度、權勢、能力、尊榮都賜給你。 \end{tabularx} \\ \\ \relax
2:38 & \begin{tabularx}{0.7\textwidth}{X} 世人和走獸，並天空的飛鳥，不論居住何處，他都交在你的手中，令你掌管這一切。你就是那金的頭。 \end{tabularx} \\ \\ \relax
2:39 & \begin{tabularx}{0.7\textwidth}{X} 在你以後必興起另一國，不及於你；又有第三國如銅，必掌管全地。 \end{tabularx} \\ \\ \relax
2:40 & \begin{tabularx}{0.7\textwidth}{X} 第四國必堅壯如鐵，就像鐵能打碎砸碎一切；鐵怎樣壓碎一切，那國也必照樣打碎壓碎。 \end{tabularx} \\ \\ \relax
2:41 & \begin{tabularx}{0.7\textwidth}{X} 你既看見像的腳和腳趾頭，一半是陶匠的泥，一半是鐵，那國將來也必分裂。你既看見鐵和泥攙雜，那國也必有鐵的力量。 \end{tabularx} \\ \\ \relax
2:42 & \begin{tabularx}{0.7\textwidth}{X} 那腳趾頭既是半鐵半泥，那國也必半強半弱。 \end{tabularx} \\ \\ \relax
2:43 & \begin{tabularx}{0.7\textwidth}{X} 你既看見鐵和泥攙雜，他們必有混雜的後裔，卻不能彼此相合，正如鐵和泥不能相合。 \end{tabularx} \\ \\ \relax
2:44 & \begin{tabularx}{0.7\textwidth}{X} 當諸王在位的時候，天上的神必另立一個永不敗壞的國度，這國度必不歸給其他百姓，卻要打碎滅絕所有的國度，存立到永遠。 \end{tabularx} \\ \\ \relax
2:45 & \begin{tabularx}{0.7\textwidth}{X} 你既看見非人手鑿出來的一塊石頭從山而出，打碎鐵、銅、泥、銀、金，那就是至大的神把將來要發生的事給王指明。這夢是確實的，這解釋也是準確的。」 \end{tabularx} \\ \\ \relax
2:46 & \begin{tabularx}{0.7\textwidth}{X} 當時，尼布甲尼撒王臉伏於地，向但以理下拜，並且吩咐人給他奉上供物和香。 \end{tabularx} \\ \\ \relax
2:47 & \begin{tabularx}{0.7\textwidth}{X} 王對但以理說：「你既能講明這奧祕，你們的神誠然是萬神之神、萬王之主，是奧祕的啟示者。」 \end{tabularx} \\ \\ \relax
2:48 & \begin{tabularx}{0.7\textwidth}{X} 於是王使但以理高升，賞賜他極多的禮物，派他管理巴比倫全省，又立他為總理，掌管巴比倫所有的智慧人。 \end{tabularx} \\ \\ \relax
2:49 & \begin{tabularx}{0.7\textwidth}{X} 但以理求王，王就派沙得拉、米煞、亞伯尼歌管理巴比倫省的事務，只是但以理仍在朝中侍立。 \end{tabularx} \\ \\ \relax
3:1 & \begin{tabularx}{0.7\textwidth}{X} 尼布甲尼撒王造了一個金像，高六十肘，寬六肘，立在巴比倫省的杜拉平原。 \end{tabularx} \\ \\ \relax
3:2 & \begin{tabularx}{0.7\textwidth}{X} 尼布甲尼撒王差人將總督、欽差、省長、參謀、財務、法官、地方官和各省的官員都召了來，為尼布甲尼撒王所立的像行開光禮。 \end{tabularx} \\ \\ \relax
3:3 & \begin{tabularx}{0.7\textwidth}{X} 於是總督、欽差、省長、參謀、財務、法官、地方官和各省的官員都聚集，站在尼布甲尼撒所立的像前，要為尼布甲尼撒王所立的像行開光禮。 \end{tabularx} \\ \\ \relax
3:4 & \begin{tabularx}{0.7\textwidth}{X} 那時傳令的大聲呼叫說：「各方、各國、各族的人哪，有命令傳給你們： \end{tabularx} \\ \\ \relax
3:5 & \begin{tabularx}{0.7\textwidth}{X} 你們一聽見角、號、琴、瑟、三角琴、鼓和各樣樂器的聲音，就當俯伏，拜尼布甲尼撒王所立的金像。 \end{tabularx} \\ \\ \relax
3:6 & \begin{tabularx}{0.7\textwidth}{X} 凡不俯伏下拜的，必立刻扔在烈火的窯中。」 \end{tabularx} \\ \\ \relax
3:7 & \begin{tabularx}{0.7\textwidth}{X} 因此百姓一聽見角、號、琴、瑟、三角琴和各樣樂器的聲音，各方、各國、各族的人就都俯伏，拜尼布甲尼撒王所立的金像。 \end{tabularx} \\ \\ \relax
3:8 & \begin{tabularx}{0.7\textwidth}{X} 在那時，有幾個迦勒底人進前來控告猶大人。 \end{tabularx} \\ \\ \relax
3:9 & \begin{tabularx}{0.7\textwidth}{X} 他們對尼布甲尼撒王說：「願王萬歲！ \end{tabularx} \\ \\ \relax
3:10 & \begin{tabularx}{0.7\textwidth}{X} 你，王啊，你曾降旨，凡聽見角、號、琴、瑟、三角琴、鼓和各樣樂器聲音的，都當俯伏拜這金像。 \end{tabularx} \\ \\ \relax
3:11 & \begin{tabularx}{0.7\textwidth}{X} 凡不俯伏下拜的，必扔在烈火的窯中。 \end{tabularx} \\ \\ \relax
3:12 & \begin{tabularx}{0.7\textwidth}{X} 現在有幾個猶大人，就是王所派管理巴比倫省事務的沙得拉、米煞、亞伯尼歌；王啊，這些人不理你的諭旨，不事奉你的神明，也不拜你所立的金像。」 \end{tabularx} \\ \\ \relax
3:13 & \begin{tabularx}{0.7\textwidth}{X} 當時，尼布甲尼撒大發烈怒，命令把沙得拉、米煞、亞伯尼歌帶過來；他們就把這幾個人帶到王面前。 \end{tabularx} \\ \\ \relax
3:14 & \begin{tabularx}{0.7\textwidth}{X} 尼布甲尼撒問他們說：「沙得拉、米煞、亞伯尼歌，你們不事奉我的神明，不拜我所立的金像，是真的嗎？ \end{tabularx} \\ \\ \relax
3:15 & \begin{tabularx}{0.7\textwidth}{X} 現在，你們若準備好，一聽見角、號、琴、瑟、三角琴、鼓和各樣樂器的聲音，就俯伏拜我所造的像；若不下拜，必立刻扔在烈火的窯中，有哪一個神明能救你們脫離我的手呢？」 \end{tabularx} \\ \\ \relax
3:16 & \begin{tabularx}{0.7\textwidth}{X} 沙得拉、米煞、亞伯尼歌對王說：「尼布甲尼撒啊，這件事我們不必回答你， \end{tabularx} \\ \\ \relax
3:17 & \begin{tabularx}{0.7\textwidth}{X} 即便如此，我們所事奉的神能將我們從烈火的窯中救出來。王啊，他必救我們脫離你的手； \end{tabularx} \\ \\ \relax
3:18 & \begin{tabularx}{0.7\textwidth}{X} 即或不然，王啊，你當知道，我們絕不事奉你的神明，也不拜你所立的金像。」 \end{tabularx} \\ \\ \relax
3:19 & \begin{tabularx}{0.7\textwidth}{X} 當時，尼布甲尼撒怒氣填胸，向沙得拉、米煞、亞伯尼歌變了臉色，命令把窯燒熱，比平常熱七倍； \end{tabularx} \\ \\ \relax
3:20 & \begin{tabularx}{0.7\textwidth}{X} 又命令他軍中的幾個壯士，把沙得拉、米煞、亞伯尼歌捆起來，扔在烈火的窯中。 \end{tabularx} \\ \\ \relax
3:21 & \begin{tabularx}{0.7\textwidth}{X} 這三人穿著內袍、外衣、頭巾和其他的衣服，被捆起來扔在烈火的窯中。 \end{tabularx} \\ \\ \relax
3:22 & \begin{tabularx}{0.7\textwidth}{X} 因為王的命令緊急，窯又非常熱，那抬沙得拉、米煞、亞伯尼歌的人都被火焰燒死。 \end{tabularx} \\ \\ \relax
3:23 & \begin{tabularx}{0.7\textwidth}{X} 但是這三個人，沙得拉、米煞、亞伯尼歌被捆綁著，掉進烈火的窯中。 \end{tabularx} \\ \\ \relax
3:24 & \begin{tabularx}{0.7\textwidth}{X} 那時，尼布甲尼撒王驚奇，急忙站起來，對謀士說：「我們捆起來扔在火裡的不是三個人嗎？」他們回答王說：「王啊，是的。」 \end{tabularx} \\ \\ \relax
3:25 & \begin{tabularx}{0.7\textwidth}{X} 王說：「看哪，我看見有四個人，並沒有捆綁，在火中行走，也沒有受傷；那第四個的相貌好像神明的兒子。」 \end{tabularx} \\ \\ \relax
3:26 & \begin{tabularx}{0.7\textwidth}{X} 於是尼布甲尼撒靠近烈火窯門，說：「至高神的僕人沙得拉、米煞、亞伯尼歌，出來，來吧！」沙得拉、米煞、亞伯尼歌就從火中出來。 \end{tabularx} \\ \\ \relax
3:27 & \begin{tabularx}{0.7\textwidth}{X} 那些總督、欽差、省長和王的謀士一同聚集來看這三個人，見火不能傷他們的身體，頭髮沒有燒焦，衣裳也沒有變色，都沒有火燒過的氣味。 \end{tabularx} \\ \\ \relax
3:28 & \begin{tabularx}{0.7\textwidth}{X} 尼布甲尼撒說：「沙得拉、米煞、亞伯尼歌的神是應當稱頌的！他差遣使者救護倚靠他的僕人，他們不遵王的命令，甚至捨身，在他們神以外不肯事奉敬拜別神。 \end{tabularx} \\ \\ \relax
3:29 & \begin{tabularx}{0.7\textwidth}{X} 現在我降旨，無論何方、何國、何族，凡有人毀謗沙得拉、米煞、亞伯尼歌的神，他必被凌遲，他的房屋必成糞堆，因為沒有別神能像這樣施行拯救。」 \end{tabularx} \\ \\ \relax
3:30 & \begin{tabularx}{0.7\textwidth}{X} 那時王在巴比倫省使沙得拉、米煞、亞伯尼歌高升。 \end{tabularx} \\ \\
[1ex]
\hline
\hline
\end{longtable}
(但二31-45)神在每一時代有祂的無限心意,並願將心意表明,使世人領悟；祂需要祂的僕人,將祂的話傳出來.如果沒有祂的僕人或祂合用之器皿,那麼祂的話語和心意,就沒法使人明白.
\newline
\newline
神為要使其他外邦人知道祂的公義審判,就必需祂的僕人傳祂的話.那麼,神是否向尼布甲尼撒王啟示呢?不是的,因為神不向罪人說話,但祂願將心意使罪人明白；知道祂的公義,並切實的悔改,所以尼布尼撒王的夢並非神的啟示.神不向尼布甲尼撒王啟示,而是向屬祂的人啟示；讓祂僕人再將祂的心意傳給世人,使他們接受救恩.但可惜每一個時代都很難找到合祂心意的人!
\newline
\newline
撒母耳時代,耶和華的言語稀少,不常有默示；因為找不到合祂心意的器皿,找不到合祂心意的但以理.各位!現今神也有許多的話語要讓世人明白,可是,合祂心意的人在那裡?有些人真的願把自己擺上,完全地奉獻,不保留,不猶疑,也不遲延；澈底地,始終如一地獻上,神要藉著這種人將祂的心意表達.所以從神的僕人來看,這夢是神的啟示,讓祂的僕人將祂公義的審判說出,神的啟示對於不信者是恐懼的,不安的,憂慮的.但對信靠祂的人卻是歡欣喜樂的,對不信的人是模糊的像在黑暗中.但對信的人,因是在光明中行有聖靈之光光照；就歡歡喜喜領受,因神要我們明白祂.
\newline
\newline
當時,但以理就是如此,所以能在王前解夢,因這夢是神的啟示.他解釋這個夢境,有一大像,這像很高大,很光耀且又很可怕；這像是象徵世界的權力,因世界之權力和屬靈之權力在爭戰中.世上權力看來高大──自高自大,很光耀──很體面,但卻可怕──因為是畸形的.注意這像,那是非常不正常的,因為是用不同的金屬造成:「這像的頭是精金,胸膛和膀臂是銀的,肚腹和腰是銅的,腿是鐵的,腳是半鐵半泥的.」(但二32-33),一個不同金屬造成的非常不正常的像.
\newline
\newline
(一)價值上的不正常──金子的價值是最貴的,其他的價值就每況愈下.世界亦是如此,初時來勢洶湧,但慢慢地價值一直減少.這不就是畸形的嗎?但在屬靈的事上則剛好相反,因神不但讓好的給我們；而且有更好的,甚至最好的也要賜給我們.
\newline
\newline
(二)重量上的不正常──金子是很重的,銀比不上金重,銅與鐵則更不如金.如研究它們的比重:則金是19.3,銀是10.5,銅8.5,鐵7.6,泥1.9,即泥的比重只有金的十分一,那可想像到這是上重下輕的像.試問這又怎能站得住呢,豈不是很易倒下來嗎?不錯,這根基太差,太脆弱,世界的事──屬世的事也是如此；完全沒根基,隨時會倒塌.
\newline
\newline
(三)硬度上的不正常──鐵和泥是最剛硬的,金是最軟的；那就是說愈來愈硬,這也是不正常的現象.但這就是屬世的東西,各位弟兄姊妹!我們該有屬靈的價值感和辨別力；要知什麼是有價值的和穩固的,什麼是沒價值和不可靠的.如果我們無屬靈的價值感,也就沒有安全感.所以當我們在神面前見這幅圖畫時,就何等需要神賜諸般的憐憫,使我們對歷史有更深的認識.
\newline
\newline
根據歷史的證據,一般的解釋是對的；金子是代表巴比倫,銀是指瑪代波斯,銅指希腊,鐵指羅馬.亦有人說,鐵是希腊,銅是波斯,銀是瑪代(瑪代和波斯分開),這都並不重要,不過我們見到當時的權力是如此:先是巴比倫,然後瑪代波斯,希腊,最後是羅馬,為何要從巴比倫開始呢?因那是外邦人的日子,(路廿一24),神的選民要敗亡,所以必須計算外邦人的日子.何時是外邦人日子完結?這是在主再來的那一天.其次,有兩隻腳,解經家說那是敘利亞和埃及,是世界的強權,亦可分東邊與西邊.至於十個腳趾我們在七和八章才研究.現今這世代也是包括在這一期間內,主必快來.
\newline
\newline
(但二34)「見一塊非人手所鑿出來的石頭」,這是一種感覺,好像一塊大石頭；其實這石頭大小並不重要,而能力卻很大.當這石頭擲過去時,就發出很大的能力,能把這麼大的像完全砸得粉碎.這是象徵著我們的主耶穌基督,祂雖是神；但祂要降世為人,祂是餘數的餘數.該知神審判人審判列國時,信靠祂的人就是餘數；當主被釘時,只剩下祂一個.但這生命的能力卻一直發展,主死後就復活.當我們讀到主從墳墓復活出來時,擋著墓前的大石頭被挪開,這石頭就成主復活的見證.生命的見證,象徵死亡的門打開.因主說:祂雖然死過,現在又活了,且活到永永遠遠,而且拿著陰間和死亡的鑰匙.死亡之門戶已經洞開,紀念的碑石已經豎立,表明耶穌基督得勝死亡.我們更知道信仰的基石也已奠定,因信靠主從死裡復活,才是有正確的信仰,且知道界限的石頭已安放,從此以後,生命與死亡已經分開,光明與黑暗分開,屬世與屬靈也分開.這石頭可象徵天國的磐石,主耶穌成為一塊非人手鑿成的石；生命的大能,能把世界的一切都砸爛.因祂已得勝,這非人手鑿成的石慢慢地擴充,石又怎能擴充呢?但見這石砸碎這像時,就變成一座大山,充滿天下.你我都是活石,要建立神的國,要將神的國建立起來充滿天下.彼前二章記說很清楚,我們都是活石,我們要與主一同建造.在十架以前,神只有一個兒子；神愛世人,甚至將祂的獨生子賜給世人,叫信靠祂的不至滅亡,反得永生.但在主死在十字架以後,神就有許多的兒女,現在在這兒也有許多神的兒女.信徒們,神給我們一個責任,一個重大的使命；就是把神的教會建立起來,把聖山建造起來,把神的國建造起來!這是何等興奮的事!我們需要能力,需要信心仰望主.
\newline
\newline
(但三1-8),二與三章的關係甚為密切,二章的像是在夢境中；但三章則是尼布甲尼撒王把像建立.或許他感到權勢上受到威脅,或許是他不肯輕易地放棄他的計謀,所以就自己建造一個像.不但頭是金的,甚至全身都是金的,雖然當時但以理的勸告曾使尼布甲尼撒王受到感動,甚至說萬神之神,萬王之王,可是他心中恐懼,神的啟示對不信的人是恐懼與不安,故要建一個金像(大概這像是木頭造而外層包金.)(以賽書四十:耶利米十).外強中乾,外表很輝煌,而裡面卻是醜陋,沒價值可言.且這像很奇怪,高六十肘,寬六肘,與一般人的身體比例不一樣.(六十肘即90呎,六肘即9呎),但這是巴比倫心目中的英雄造像,且「六」乃巴比倫人最喜歡的數字.那是說屬罪惡,屬世界的力量要想掙扎,要爭戰,要奮鬥,甚至用暴力以達目的.故王傳令所有的人,凡聽見樂聲都要跪拜,因這是他的方法；用暴力要得人的尊敬,用強制的命令要人順服；如有何人違反,他就嚴厲處置,甚至處死刑,這是惡者所發出的猙獰的狂笑.然而,用人的方法可以勝過一切嗎?信靠神的人在暴力下到底應有什麼態度呢?屬神的事與屬世的事到底不同；世上的人可用強制或暴力的方法,以求得外表的行動,至於內心如何就全不顧及.但神卻非這樣,祂從不強制,只以愛與勸說來感動.祂要我們甘心樂意到主面前將愛獻上,人雖常不尊重神,神卻尊重人；人雖常輕忽神,神卻從不輕忽人.尼布甲尼撒王用強暴的,強制的,壓迫的方法,果然有許多人去敬拜那「像」,可是但以理和他的三位朋友不肯,可見他們是真正屬於神的人.他們必受到控告,雖然這是事實,並非誣告,可是卻是個惡毒的控告.這「控告」在亞蘭文很有意思,是「撕碎」:野獸的動作,意要咬破至成塊,那就是野獸不是人了!這些迦勒底人簡直是野獸,是惡者,要到遍地去尋找可吞噬的人.他們的控告使尼布甲尼撒王震驚,直至他極其惱怒,因為人們不但控告這三位屬神的人,也控告尼布甲尼撒王.他們似乎說:「你這愚蠢的人,竟然重用這幾個猶大人；你這愚蠢的人,這些人根本不會效忠於你,他們的宗教給你這麼大的威脅,你的權柄在那裡?你的尊榮在那裡?」所以尼布甲尼撒王的心情非常不歡,為此,他一定要這三人跪拜.並且說:「如果不跪拜,那就立刻將你們扔在烈火的yiu中,有那一位神能救你們脫離我的手呢?」他好像是有必然得勝的把握.試問這幾個敬畏耶和華的人,除了妥協之外,怎可能逃避這危險呢!」
\newline
\newline
16至18節是三人信心的宣告,信心是無攷慮亦不需充份的解釋,亦不需要辯論.三人對王說:「無論在何景況中,我們絕不妥協,因信心使我們如此.」他們三人到底有否研究對付的辦法,或是想去等候神的旨意呢?一切都沒有,許多時候我們說尋求神的旨意,其實神的旨意早就清楚擺在前面；可是我們不順服,總是還在理論,想辦法逃避,想辦法推辭,想辦法的妥協.還要說:神啊!我不知道是否你的旨意,現在暫時拜一拜,我外表拜而內心不拜不也可以嗎,直到我清楚你的旨意時再算吧.這三人是否如此?信心是不等候憑據的,多少時候我們在等候憑據；真的信心不需要計較後果,也不求神蹟.他們並沒有求神救他們,他們只對尼布甲尼撒說:我們的神若救我們,祂就一定能救我們出來；否則也有祂的美意,無論如何只有一個答案,只有一個方向；就是完全順從神,不順從人.有誰的信心比這更偉大呢!這是完全澈底地將自己擺上,因主有祂的美意；當時王怒氣填胸,要把火加上七倍,其實這是最愚蠢的事.火增加七倍倒不如減七倍更好,那可以使神的僕人在火yiu中慢慢地燒死,不就更痛苦嗎?現在增七倍烈火,實在太猛烈了,甚至抬他們的人也都燒死.這些順服人而不順服神的人被燒死,神是不負這責任的.但那三個屬神的順服神的人,神負完全的責任.火yiu的路既是神要他們走,那就甘心樂意地走去,但神卻負祂僕人安全的責任.他們穿著褲子,內袍和外衣；這可能是以色列祭司所穿的衣袍,因在出埃及記廿九章說祭司所穿的是長袍.他們如此穿著,也表明在異邦人面前作見證──敬拜神的見證.這是何等大的信心!當火加上七倍時,他們更感到神的同在,因有人與他們同在火yiu中走；這是神,是道成肉身前的基督.當他們將自己完全交托神時,神就向他們施恩,他們沒向神求神蹟,但神給他們神蹟.沒求搭救,但神卻搭救.各位弟兄姊妹!我們何等需要神的恩典,我們如要得安慰,就必定會先有苦難；要保護就會有危險；要得勝,一定會有爭戰；要神蹟,就先會有試煉.有人說現今的世代沒神蹟,其實是有許多的神蹟；只是我們看不見,因為我們不禱告,不付代價,不為神作見證；神怎會將祂奇妙的能力彰顯出來呢!我們需要火般的試煉,我們需行走在火的行列.火對但以理三位朋友沒有任何作用,那像火一般的現世的力量,對基督徒也沒有作用.只要我們真正地仰望神,那麼火給我們的作用就是把我們捆綁燒掉.好似那三個人的經歷一樣.這是何等奇妙!世界有許多方法對付我們,但一切的方法都成為我們的祝福.這火的行列就是十字架的道路,我們需要走,因有主與我們同在並親自引導；只要看我們是否肯擺上,是否肯仰望神.
\newline
\newline
主在客西馬尼園禱告完畢時說:「起來,我們走吧!」我們需在這火的世代跟隨祂,只許向前,不許退後；只許成功,不許失敗,只許愛主,不許冷淡.若有人不愛主,那人是可咒可詛的,因為主已經近了.我們必需向前走,「起來,我們走吧!」.
\newline
\newline


\newpage

\section{第45屆~港九培靈會~唐佑之~第四講}
\label{sec:1250}
\textbf{時間及永恆}
\newline
\newline
連結: \href{https://www.hkbibleconference.org/session-message/view/1250}{\texttt{https://www.hkbibleconference.org/session-message/view/1250}}
\newline
\newline
\hyperref[sec:1248]{< < < PREV SERMON < < <}
~
\hyperref[sec:toc]{[返目錄]}
~
\hyperref[sec:1252]{> > > NEXT SERMON > > >}
\newline
\newline
但以理書 7:1-7:28
\newline
\begin{longtable}{cl}
\hline
\hline
章節 & 經文 (和合本修訂版)\\
\hline
7:1 & \begin{tabularx}{0.7\textwidth}{X} 巴比倫王伯沙撒元年，但以理在床上做夢，腦中看見異象，就記錄這夢，述說其中的大意。 \end{tabularx} \\ \\ \relax
7:2 & \begin{tabularx}{0.7\textwidth}{X} 但以理說：我在夜間的異象中觀看，看哪，天上有四風，突然颳在大海之上。 \end{tabularx} \\ \\ \relax
7:3 & \begin{tabularx}{0.7\textwidth}{X} 有四隻巨獸從海裡上來，牠們各不相同： \end{tabularx} \\ \\ \relax
7:4 & \begin{tabularx}{0.7\textwidth}{X} 頭一個像獅子，有鷹的翅膀；我正觀看的時候，牠的翅膀被拔去，牠從地上被扶起來，用兩腳站立，像人一樣，還給了牠人的心。 \end{tabularx} \\ \\ \relax
7:5 & \begin{tabularx}{0.7\textwidth}{X} 看哪，另有一獸如熊，就是第二獸，半身側立，口裡的牙齒中有三根獠牙。有人吩咐這獸說：「起來，吞吃許多的肉。」 \end{tabularx} \\ \\ \relax
7:6 & \begin{tabularx}{0.7\textwidth}{X} 其後，我觀看，看哪，另有一獸如豹，背上有四個鳥的翅膀；這獸有四個頭，還給了牠權柄。 \end{tabularx} \\ \\ \relax
7:7 & \begin{tabularx}{0.7\textwidth}{X} 其後，我在夜間的異象中觀看，看哪，第四獸可怕可懼，極其強壯，有大鐵牙，吞吃嚼碎，剩下的用腳踐踏。這獸與前面所有的獸不同，牠有十隻角。 \end{tabularx} \\ \\ \relax
7:8 & \begin{tabularx}{0.7\textwidth}{X} 我正思考這些角的時候，看哪，其中又長出另一隻小角；先前的角中有三隻角在它面前連根被拔出。看哪，這角有眼，像人的眼，有口說誇大的話。 \end{tabularx} \\ \\ \relax
7:9 & \begin{tabularx}{0.7\textwidth}{X} 我正觀看的時候，有寶座設立，上面坐著亙古常在者。他的衣服潔白如雪，頭髮如純淨的羊毛。寶座是火焰，其輪為烈火。 \end{tabularx} \\ \\ \relax
7:10 & \begin{tabularx}{0.7\textwidth}{X} 有火如河湧出，從他面前流出來；事奉他的有千千，在他面前侍立的有萬萬；他坐著要行審判，案卷都展開了。 \end{tabularx} \\ \\ \relax
7:11 & \begin{tabularx}{0.7\textwidth}{X} 於是我觀看，因這角說誇大的話，我正觀看的時候，那獸被殺，身體被毀，扔在火中焚燒。 \end{tabularx} \\ \\ \relax
7:12 & \begin{tabularx}{0.7\textwidth}{X} 其餘的獸，權柄都被奪去，生命卻得以延續，直到所定的時候和日期。 \end{tabularx} \\ \\ \relax
7:13 & \begin{tabularx}{0.7\textwidth}{X} 我在夜間的異象中觀看，看哪，有一位像人子的，駕著天上的雲而來，被領到亙古常在者面前。 \end{tabularx} \\ \\ \relax
7:14 & \begin{tabularx}{0.7\textwidth}{X} 他得了權柄、榮耀、國度，使各方、各國、各族的人都事奉他。他的權柄是永遠的，不能廢去，他的國度必不敗壞。 \end{tabularx} \\ \\ \relax
7:15 & \begin{tabularx}{0.7\textwidth}{X} 至於我－但以理，我的靈在我裡面憂傷，我腦中的異象使我驚惶。 \end{tabularx} \\ \\ \relax
7:16 & \begin{tabularx}{0.7\textwidth}{X} 我走近其中一位侍立者，問他這一切的實情。他就告訴我，使我知道這事的解釋： \end{tabularx} \\ \\ \relax
7:17 & \begin{tabularx}{0.7\textwidth}{X} 這四隻巨獸就是將要在世上興起的四個王。 \end{tabularx} \\ \\ \relax
7:18 & \begin{tabularx}{0.7\textwidth}{X} 然而，至高者的眾聖者必要得到這國度，並且擁有它，直到永遠，永永遠遠。 \end{tabularx} \\ \\ \relax
7:19 & \begin{tabularx}{0.7\textwidth}{X} 於是我想要更清楚知道第四獸的實情，牠與一切的獸不同，甚是可怕，有鐵牙銅爪，吞吃嚼碎，剩下的用腳踐踏； \end{tabularx} \\ \\ \relax
7:20 & \begin{tabularx}{0.7\textwidth}{X} 頭上有十隻角和那另長出的一角，三隻角在這角面前掉落；這角有眼，有口說誇大的話，形狀比牠的同類更強。 \end{tabularx} \\ \\ \relax
7:21 & \begin{tabularx}{0.7\textwidth}{X} 我觀看，這角與眾聖者爭戰，勝了他們， \end{tabularx} \\ \\ \relax
7:22 & \begin{tabularx}{0.7\textwidth}{X} 直到亙古常在者來到，為至高者的眾聖者伸冤，眾聖者得到國度的時候就到了。 \end{tabularx} \\ \\ \relax
7:23 & \begin{tabularx}{0.7\textwidth}{X} 那侍立者這樣說：第四獸就是世上要興起的第四國，與其他各國不同，它要併吞全地，並且踐踏嚼碎。 \end{tabularx} \\ \\ \relax
7:24 & \begin{tabularx}{0.7\textwidth}{X} 至於那十隻角，就是從這國中興起的十個王；後來又興起另一王，與先前的不相同，他要制伏三個王。 \end{tabularx} \\ \\ \relax
7:25 & \begin{tabularx}{0.7\textwidth}{X} 他說話抵擋至高者，折磨至高者的眾聖者，又改變節期和律法。眾聖者要交在他手中一年、兩年、又半年。 \end{tabularx} \\ \\ \relax
7:26 & \begin{tabularx}{0.7\textwidth}{X} 然而，他坐著要行審判；他的權柄要被奪去，毀壞，滅絕，一直到底。 \end{tabularx} \\ \\ \relax
7:27 & \begin{tabularx}{0.7\textwidth}{X} 國度、權柄和天下諸國的大權必賜給至高者的眾聖民。他的國是永遠的國，所有掌權的都必事奉他，順從他。 \end{tabularx} \\ \\ \relax
7:28 & \begin{tabularx}{0.7\textwidth}{X} 這事到此結束。我－但以理因這些念頭甚是驚惶，臉色也變了，卻將這事記在心裡。 \end{tabularx} \\ \\
[1ex]
\hline
\hline
\end{longtable}
(但七1-14),我們的神是歷史的主,在祂歷史之計劃中.歷史是有末期的,在這末期過去後；神就完全掌管這世界,賜下永遠的平安.
\newline
\newline
(這末期在路廿一24)說,乃是外邦人的日子.當神的選民被外邦人擄去之時,猶大國敗亡；外邦人日子,就從巴比倫時開始到現在,直等到主再來之日.我們生在外邦人日子的時代,我們有更大之歷史使命,所以我們何等需要神諸般的憐憫.在第二章中之像不夠清楚,描述也不夠完全,必須看第七章補充的解釋.從這章開始有四大異象,七章是第一異象,八章是第二異象,九章是第三異象,十至十二章是第四異象.這四大異像是但以理直接在神前領受,能將巴比倫王所見(二章)的像加以解說.現在但以理所見的異象,是由神的使者解釋,這觀點就不相同.二章中巴比倫王所見的巨像,是從人之眼光看世界的歷史；但在七章是以屬靈之眼光來看世界,非常醜陋,在七章中之獸就十分兇暴.我們需要聖靈帶領進入真理；不但進入且要深入,因聖靈能帶我們到屬靈之深處.尼布甲尼撒曾見神所指的像,亦曾感動；他知天上的神才是歷史之主,才是萬神之神,萬主之主.但他那罪惡之本性,和那不信之惡心,在裡面常發動,所以常失敗；正像他自己立金像以炫耀自己；用暴力強制人順服,非常狂傲.故在第四章中,當他狂傲時,神要他知自己不過是人,且是罪人；結果他精神錯亂像野獸般在吃草,得此嚴重教訓後才悔改；承認天上之神,特別在四章中,他在天上的神前低頭,於是尼布甲尼撒王乃在歷史中消失.在四章之後再不提及,他在主前五六○年已經死亡.他兒子 Amel-marduk 也是暴虐之君作王二年後,由他親戚 Neriglisser 纂位,得位四年後,由尼布甲尼撒王之女婿作王三年後,交他兒子,也就是伯沙撒(七1)他是尼布甲尼撒王的外孫.七章的開始是他的元年,八章是三年,五章則是他的末年；即主前五三九年,巴比倫整個王朝被波斯吞滅.故五章依次序說是應在八章之後,為何放在前面?因六章是大利烏王的時候,主要是一至六章說及個人及外邦之人之事；即以色列個人為神見證及外邦人君王個人的事.按次序先看七章,八章至五章,然後九章和六章.
\newline
\newline
但以理的三朋友為何以後就不再提及?現在已是伯沙撒王的時代,幾十年後之事了.但以理已屆老年,三位朋友也許被主接去；他們忠心的事奉,已得神厚厚的賞賜,他們工作及勞苦的果效必隨著他們.神將祂的工人接回去,可是卻將祂的工作存留著,交給忠心的僕人但以理.神在每個時代都選祂的工人,祂需要許多的工人,正如主說:「莊稼多,作工的人少.」無論工人有多少,都是不夠用的.但在另一方面看,神的工人並不少,不過真正合神心意的有多少?真正被神重用的器皿又有多少?祂心裡在焦急,一直在呼召,要你我站起來作時代的工人.我們有何理由推辭?實在不可遲延!因為只有一點點的時候,那是最嚴重的時間.神需要許多合祂使用的工人,神需要但以理,神需要那屬祂的信靠祂的所有兒女們.所以,但以理雖然老,而神仍用他；他年已老邁,可是在伯沙撒王時他仍活著,甚至大利烏王時,他還活著.(但一21),是否神難找合用的人,故而使忠心的老僕但以理不能離世而要存活著呢?這問題實在使但以理心裡難過!所以在當他見異象時,有兩次提到他心裡又懼怕又非常不安!(七章),因時代愈來愈嚴重,環境愈來愈惡劣,他實在真的無法在這光景中作工.正如主說:「父啊,求你救我脫離這個時候.」(十二)可是跟著又說:「我原是為這時候來的.」祂願意脫離這時代,但這是神的旨意和安排.我們在這時代,並非糊塗地虛度此生；也不是要我們內心不安,心靈軟弱.神要我們負應負的責任,明白神的話語,要見到異象中屬靈的意境,要看見別人所不能見到的；才能說別人不會說的話,而這些話是別人最需要明白的.
\newline
\newline
在七章與二章中的夢是相似的,因二章之人像主要有四部份,七章之獸也是四個,所以有許多相同點.這些獸是怎樣出來?當但以理在夢中見異象時；是見「天的四風陡起,颳在大海中」.「大海」就是世界,因世界像海一般,動蕩不安,變化莫測.在啟示錄中記載:當新天新地來臨,海就不再有了海島也沒有了.我們都生活在海島,覺得這是多麼動蕩.你我都是屬神的人,如果我們在這世界生活,就與世界共同沉浮；因我們沒有方向也沒有目標,隨著世界改變,所以應該在神前有堅定的信心.
\newline
\newline
「天起了風」,「風」在聖經中提及至少有一百五十次之多；舊約有九十次,新約有三十次,都是著重在神的能力.世界上的情勢不論如何改變,神的能力卻仍存在.雖然許多時候我們看不見,但並非因看不見而說不存在.我們正如此幼稚軟弱和膚淺,許多時我們只看見這世界權力的情勢,而不看見神的能力!先知以利沙有一次在神前禱告,求神開少年人的眼睛；我們多麼需要看見,要看見神的能力.因我們知道在這世界,特別這末世,在外邦人的日子裡,有兩種勢力在衝突；屬神的,與屬鬼魔的,屬靈的,與屬世的.我們在此兩種情形中爭戰,感到非常的困難和軟弱,我們實在須要神的憐憫.
\newline
\newline
在海中出來的第一個獸像獅子,獅為萬獸之王,有鷹的翅膀,鷹乃飛禽之王.這樣就正如但以理為尼布甲尼撒王所解說之人像一樣,那人像的頭是金的,是巴比倫,是有權力的.從政治之體系而言,乃是完全專制的政治,故其勢力甚大.這獸站立起來就像人一般,當時吸引許多人；這是「權力第一」的世界,許多人看重權力,勢力,正如第二章所說的.
\newline
\newline
又有一獸如熊一樣,我們知道熊就遠不如獅,因其動作不如獅；兇暴也不如獅,但卻仍是有力的；正如人像之膀臂及胸膛,這就是瑪代波斯的力量.在這兒說「旁跨而坐」,是表示它的力量是單方面不平衡的；試看瑪代波斯的政體並非專制政體,是分頭政治,勢力並不均衡.他們注意政治之組織,在管制上似乎相當成功；甚至口含三肋骨,那是指三個國家:埃及,敘利亞,呂彼亞.(以西結卅5)在以西結書提及當時外邦人很有力量,第一強權巴比倫之後,跟著是波斯,波斯起來,管制許多的國.
\newline
\newline
第三個像豹,正如尼布甲尼撒王所見的像中之銅腰,是指希腊.希腊的王中特別是亞力山大帝,非常強大,迷信武力,窮兵黷武,他侵略小亞細亞整個地方.在背有四翅膀,歷史告訴我們,亞力山大有四個大將軍.亞力山大三十一歲死後,權力就分由四將軍所掌握,也就是說,這獸有四個頭,都得了權柄.
\newline
\newline
第四個獸並沒有說像什麼,是因為無從描寫.我們想到啟示錄十三章說有一獸要上來,這獸有十個角,形狀像豹,腳像熊,口像獅,這是集合前三獸之部份；也正如尼布甲尼撒王所見的像的腳部份,腳有十指.十角是講羅馬,在初期教會至中世紀解經家有不同之解釋,可能不完全正確.因自羅馬到現在直至主再來時,是很長久的時間,會有許多變遷,實在很難說是指什麼,又那是指示什麼.因我們知道但以理所說的,乃是到波斯為止,而啟示錄乃是到羅馬帝國之時,這以後的就難以清楚解釋,在此見這獸有相同之點,那就是兇暴,殘忍,毀壞,這時候簡直零亂不堪.
\newline
\newline
此處與二章所記的也有不同之處,因二章中人像是整體,而這獸則是分開的.人像是固定的,雖是不同的金屬,但仍是整個的.可是在這裡有不同的獸,各具特點.二章的人像無動作,無生命；但在此是有動作,有獸的生命,故那兇暴的動作使人感到可怕.此處特別提到小的角及三角,據啟示錄的解釋,小角是指敵基督；三角亦是可能指惡者之使者,我們只可明白到這地步.在此混亂的世界,神仍坐著為王,巴比倫是專制之政治,希腊乃貴族政治,羅馬為帝國之政治.故有人想及人像的兩腳表示兩政體,有鐵有泥,泥代表民主政體,這民主政體如有規律就好,但如行得不好,就非常混亂.我們記得那人像被一非人手所鑿的石一打就全部砸碎.在此則不同,一獸過去另一獸就來,又有先後之別.在第二章中,石頭打去後成了大山,那是神的國,神自己作王.在詩篇九十至九十七篇有說:耶和華作王,世界就堅定.這神的國是神權政治,正如以色列在神前所立的國一樣,耶和華是他們的王,當以色列民要立王時,神對撒母耳說:「他們不是厭惡你,乃是厭惡耶和華神.」我們從舊約的以色列就想到新約的教會.在教會裡沒有別的頭,只有耶穌基督是我們的頭.神安定在天,當大水泛濫時,耶和華仍坐著為王.因為神管理萬有雖暫時容忍一些惡人的作為,終久要直接干預.世人實在經受不起罪之威脅,一直等待著好的政治出現,而世界的權力加厲.一獸比一獸更兇,不單是一般人受苦,神的兒女更加痛苦.有時我們被人辱罵說:「有那一個神能救你?」正如尼布甲尼撒王對三個義人所說的一樣.在詩篇說:「你們的神在那裡呢?」如真有神的話,神是否有能力?或說神是有能力但無慈愛,祂雖然可以將罪惡除去,但見祂沒有慈愛,祂不愛你.或說神是有慈愛而沒有能力,祂雖然愛世人,特別愛屬於祂的人；但沒有能力,愛莫能助.神是如此的嗎?不是的,祂是大有能力,滿有慈愛的.
\newline
\newline
「我觀看見有寶座」,上面坐著的是互古長存的.我們的神是昔在今在永在的神,祂是創始成終的神,祂坐著為王,祂不需站起來.希伯來一章說:「當主耶和華作成了救贖的工作就坐下.」「坐下」表明做成了工作.人在活動,人在蠢動,但耶和華坐著為王,我們可在主裡得安息.因祂如此向我們施恩,連衣服頭髮都潔白,那表明祂是聖潔的永恆的,人子也如火焰一樣.這是以西結一與十章說:這寶座如車一般有輪子的充滿火焰.詩九十七篇說:神在前面過的時候,火就降下,有神的公義的存在；神在何處,公義也隨著在何處.神行走在歷史當中,在人類中,社會中,教會中及家庭中；神是無所不在的神,是公義的神,因祂定要施行審判.神在歷史中,我們依靠祂,神一定要保守屬祂的,我們還怕什麼?在此特別提到,我在夜間的異象中觀看,見有一位像人子的,駕著天雲而來,被領到互古常在者面前,祂還要再到這世上來.天使曾對加利利人說,這位主怎樣被接到天上去,將來也照樣再來.(啟一7)「看啊!祂駕雲降臨,眾目要看見祂,連刺祂的人也要看見祂；地上的萬族都要因祂哀哭,這話是真的.」祂駕雲降臨,祂正來了,已在途中了,但我們仍未看見,是明天還是後天?現在已經來了,不過我們還未看見,因為歷史每一時刻,都要帶到歷史的結局.你我都有責任和使命,快領人蒙受主恩；還有一點點的時候,主就再來,當時但以理心裡甚難過並驚惶.七15及28節,「心中甚驚惶,臉色也改變.」他驚惶,乃因這世界的事放在他身上,福音的重任也放在他身上.這世界要滅亡,如再不起來,這世界會毀滅；我們不能再遲延,該熱切地為自己也為未得救的人禱告.讓得救的人數增加能使主心滿足,主愛已臨到世上.我們若不愛主,那是可咒可詛的,主再來的日子近了.
\newline
\newline


\newpage

\section{第45屆~港九培靈會~唐佑之~第五講}
\label{sec:1252}
\textbf{異象及史實}
\newline
\newline
連結: \href{https://www.hkbibleconference.org/session-message/view/1252}{\texttt{https://www.hkbibleconference.org/session-message/view/1252}}
\newline
\newline
\hyperref[sec:1250]{< < < PREV SERMON < < <}
~
\hyperref[sec:toc]{[返目錄]}
~
\hyperref[sec:1253]{> > > NEXT SERMON > > >}
\newline
\newline
但以理書 8:1-8:27
\newline
\begin{longtable}{cl}
\hline
\hline
章節 & 經文 (和合本修訂版)\\
\hline
8:1 & \begin{tabularx}{0.7\textwidth}{X} 伯沙撒王在位第三年，有異象向我－但以理顯現，是在先前所見的異象之後。 \end{tabularx} \\ \\ \relax
8:2 & \begin{tabularx}{0.7\textwidth}{X} 我在異象中觀看，見自己在以攔省書珊的城堡中；我在異象中又見自己在烏萊河邊。 \end{tabularx} \\ \\ \relax
8:3 & \begin{tabularx}{0.7\textwidth}{X} 我舉目觀看，看哪，有一隻公綿羊站在河邊，牠有兩隻角，這兩角都高，一角高過另一角，後長出來的比較高。 \end{tabularx} \\ \\ \relax
8:4 & \begin{tabularx}{0.7\textwidth}{X} 我見那公綿羊向西、向北、向南牴撞，沒有任何獸在牠面前站立得住，沒有能逃脫牠手的；牠任意而行，自高自大。 \end{tabularx} \\ \\ \relax
8:5 & \begin{tabularx}{0.7\textwidth}{X} 我正思想的時候，看哪，有一隻公山羊從西而來，遍行全地，腳不著地。這山羊兩眼當中有一隻顯眼的角。 \end{tabularx} \\ \\ \relax
8:6 & \begin{tabularx}{0.7\textwidth}{X} 牠往我先前所見、站在河邊、有雙角的公綿羊那裡，以猛烈的怒氣向牠直闖。 \end{tabularx} \\ \\ \relax
8:7 & \begin{tabularx}{0.7\textwidth}{X} 我見公山羊靠近公綿羊，向牠發怒，攻擊牠，折斷牠的兩角。公綿羊在公山羊面前站立不住；牠把公綿羊撞倒在地，用腳踐踏，沒有能救公綿羊脫離牠手的。 \end{tabularx} \\ \\ \relax
8:8 & \begin{tabularx}{0.7\textwidth}{X} 這公山羊長得極其高大，正強壯的時候，那大角折斷了，從角的下面向天的四方長出四隻顯眼的角來。 \end{tabularx} \\ \\ \relax
8:9 & \begin{tabularx}{0.7\textwidth}{X} 從四角中的一角又長出另一隻小角，向南、向東、向佳美之地，日漸壯大。 \end{tabularx} \\ \\ \relax
8:10 & \begin{tabularx}{0.7\textwidth}{X} 牠漸壯大，高及諸天萬象，把一些天象和星辰摔落在地，用腳踐踏。 \end{tabularx} \\ \\ \relax
8:11 & \begin{tabularx}{0.7\textwidth}{X} 牠自高自大，自以為高及萬象之君，牠除掉經常獻給君的祭，毀壞君的聖所。 \end{tabularx} \\ \\ \relax
8:12 & \begin{tabularx}{0.7\textwidth}{X} 因罪過的緣故，有軍隊和經常獻的祭交給牠。牠把真理拋在地上，任意而行，無往不利。 \end{tabularx} \\ \\ \relax
8:13 & \begin{tabularx}{0.7\textwidth}{X} 我聽見有一位聖者說話，又有一位聖者向那說話的聖者說：「這經常獻的祭、帶來荒涼的罪過、聖所與軍隊被踐踏的異象，要持續到幾時呢？」 \end{tabularx} \\ \\ \relax
8:14 & \begin{tabularx}{0.7\textwidth}{X} 他對我說：「要到二千三百日，聖所就必潔淨。」 \end{tabularx} \\ \\ \relax
8:15 & \begin{tabularx}{0.7\textwidth}{X} 我－但以理見了這異象，想要明白其中的意思。看哪，有一位形狀像人的站在我面前。 \end{tabularx} \\ \\ \relax
8:16 & \begin{tabularx}{0.7\textwidth}{X} 我聽見烏萊河中有人聲呼叫說：「加百列啊，要使這人明白這異象。」 \end{tabularx} \\ \\ \relax
8:17 & \begin{tabularx}{0.7\textwidth}{X} 他就來到我所站的地方。他一來，我就驚慌，臉伏於地。他對我說：「人子啊，你要明白，因為這是關乎末後時期的異象。」 \end{tabularx} \\ \\ \relax
8:18 & \begin{tabularx}{0.7\textwidth}{X} 他對我說話的時候，我在沉睡，臉伏於地。他就摸我，扶我站起來。 \end{tabularx} \\ \\ \relax
8:19 & \begin{tabularx}{0.7\textwidth}{X} 他說：「看哪，我要指示你惱怒結束的時候必成的事，因為這是關乎末後指定的時期。 \end{tabularx} \\ \\ \relax
8:20 & \begin{tabularx}{0.7\textwidth}{X} 你所看見那有雙角的公綿羊就是瑪代王和波斯王。 \end{tabularx} \\ \\ \relax
8:21 & \begin{tabularx}{0.7\textwidth}{X} 那公山羊就是希臘王；兩眼當中的大角就是第一個王。 \end{tabularx} \\ \\ \relax
8:22 & \begin{tabularx}{0.7\textwidth}{X} 至於角折斷了，又從角的下面長出四隻角，意思就是有四個國要從這國興起，只是權勢都不及它。 \end{tabularx} \\ \\ \relax
8:23 & \begin{tabularx}{0.7\textwidth}{X} 這四國末期，惡貫滿盈的時候，必有一王興起，面貌兇惡，詭計多端。 \end{tabularx} \\ \\ \relax
8:24 & \begin{tabularx}{0.7\textwidth}{X} 他的權柄極大，卻不是因自己的能力；他要施行驚人的毀滅，無往不利，任意而行，又要毀滅強有力的人和眾聖民。 \end{tabularx} \\ \\ \relax
8:25 & \begin{tabularx}{0.7\textwidth}{X} 他用權術使手中的詭計成功；他的心自高自大，趁人無備的時候毀滅多人。他又起來攻擊萬君之君，至終卻非因人的手而遭毀滅。 \end{tabularx} \\ \\ \relax
8:26 & \begin{tabularx}{0.7\textwidth}{X} 所說二千三百日的異象是真的，但你要將這異象封住，因為它關乎未來許多的日子。」 \end{tabularx} \\ \\ \relax
8:27 & \begin{tabularx}{0.7\textwidth}{X} 於是我－但以理昏倒，病了數日，然後起來辦理王的事務。我因這異象驚駭不已，但還是不能了解。 \end{tabularx} \\ \\
[1ex]
\hline
\hline
\end{longtable}
神是歷史之主,在歷史中沒有一件事是偶然發生的；不是一般人所說的命運,而是神的安排,神的計劃.因神在歷史中彰顯祂的作為,啟示祂的心意；祂的作為和心意就是救贖,因為祂要拯救這世界,免去極大的死亡及審判.祂是公義的,祂斷不以有罪的為無罪.我們這些蒙恩的人到祂面前說:主耶和華阿,你若究察罪孽,誰能站立得住呢?
\newline
\newline
但你有赦罪的恩,我們不過是蒙恩的罪人.我們看見整個世界被罪惡所佔有,是極激烈的爭戰,直到主來的日子.因此不但得勝,且永遠得勝；我們靠著那加給我們力量的主,不但凡事都能作；而且靠祂得勝,得勝有餘的.但以理在信心過程中,他有這樣的認識和體驗的見證,成為當代及歷史的信息.直至如今我們讀此書時,覺得是主在對我們說話,向我們施恩.
\newline
\newline
在但七章從始至終是說及四大異象,七章第一個,八章第二個,九章第三個,十至十二章第四個.這都是神直接向祂的僕人但以理啟示,且有神的使者──天使長為他解釋,那是何等奇妙奧秘的真理.在第七章許多是重複的,第八章見到的異象集中某部份,與七章有很大不同的重點；八章中的公綿羊公山羊,是集中在第二及第三部份.第一部份是巴比倫王國,金的頭是獅子,可是巴比倫王國很快就過去,為什麼?當時巴比倫國不是極強盛的嗎?但神的僕人,他站在時代的頂點上,他能看見人所不能看見的,所以能說出別人所不能說的話.各位!你我都需要有此屬靈的遠見及先見,假如真有屬靈的看見；才會在屬靈的生命上繼續長進,但願我們能看見祂要我們看見的.
\newline
\newline
但以理見異象時,中文聖經說:「我以為在以攔省書珊城.」(這「以為」兩字中文聖經在旁用小點點著,意即原文沒有.那即是說但以理以為到書珊城,其實並非真的到了這城,只是屬靈的感受.書珊城乃波斯的京城,因巴比倫王朝快亡時,波斯的強權即興起,神讓但以理看到這光景.在八章中伯沙撒王位第三年,再看五章就更清楚見到伯沙撒王末年,也像巴比倫一樣的過去了.屬世的一切都會改變,毀壞而過去,不會永存；祗有神的國才永不動搖,這是本書反覆的主題,我們該好好地把握著.
\newline
\newline
繼而我們看見人像中之第二及第三部份,亦可說是第二及第三獸.人像之第二部份是臂及胸膛,「我舉目觀看,見有雙角的公綿羊站在河邊.」(但八3),公綿羊乃波斯國之國號,(今考古家發現波斯古錢幣有此像.)波斯之前先有瑪代,其實是兩個國,後才併合為一,故兩個角中,其中一角比另一角高些,古列王比大利烏王力量更大.公綿羊出來好像很有力,但亦只不過二百多年,巴比倫王過去時主前五三九年.當第三次強權勢力希腊起來時約三三一年,希腊強權亦不過二百多年.可見世界情勢極不安定,大水泛濫之時,耶和華坐著為王,耶和華神仍是歷史之主.其他一切都會過去,沒永存的價值.
\newline
\newline
「我正思想時」,他在研究這勢力能持久嗎?不會的,因另有公山羊從西邊出來；這是希腊強權,遍行全地,腳都不沾塵土(八5),表示敏捷.亞力山大帝到處用武力征服國家,最後他甚至哭著說:他再沒有地方去征服了.這狂傲不可一世的樣子,表明當時他實在很有力量.這就是所說的「有一非常角」,他征服了許多地,今希腊附近的愛琴海,「愛琴」乃是公山羊的意思,公山羊象徵希腊國.
\newline
\newline
「四角中長出一小角」,小角是誰?我們知道亞歷山大帝有四個將軍,當亞歷山大帝死時,四個將軍就分裂,長出小角來是有一兇暴之將軍起來,向南向東,向榮美之地,漸漸成為強大(八9),向南是埃及,向東是亞述巴比倫的地區,是米所波大米的地區,向榮美之地是巴勒斯坦──神所應許之地.這小角大概就是指四將軍中之一名叫依比芬AntiochusIV(Epiphanes),是非常利害的人,提倡希利尼文化；把猶大人的宗教及生活傳統設法破壞消除；此兇暴褻瀆神的態度,甚為可怕,因為他逼迫神的子民.「他漸漸強大,高及天象,將些天象星宿拋落在地,用腳踐踏.」(但八10)這「星宿」有人說是指天使,又有說是以色列百姓,我以為他是指百姓,因他是迫害神的百姓.更可能是指神的使者,及神的僕人們,因他們在最惡劣的環境下,常領以色列人繼續敬拜神.他不但自高自大,且廢除聖殿之獻祭,且將聖所全部毀壞,並做各種褻瀆神的事!就是要將天上的神侮辱!
\newline
\newline
有人以為那可能是瑪客比的時代,那是在主前一三○年左右,那時代就非但以理時代所能明白的.聖經曾論及聖殿的三個階段:一是所羅門的聖殿,一是所羅巴伯的聖殿(是以色列人歸回後重建),又在新約中有希律建的聖殿,其實在歷史中還有瑪喀比(猶大政治獨立運動者)聖殿.因以色列人是神的選民仍要敬拜神.但希腊之暴君各種方法逼迫他們,且建造偶像,用豬來獻祭,破壞聖殿之聖潔,引誘青年人犯罪.雖然歷史不重演,但許多類似的事仍重複,因魔鬼的方法就是如此.現代在西方那種迷信的發展,比起東方真有過之而無不及；那是非常可怕的,且專門吸引青年人；在美國加省,各式各樣的宗教,甚至有稱為撒旦教會,用盡人最敗壞的方式來敬拜.但以理之預言已經提及,到底何時才能完結?實在無從知道,所以但以理極感痛苦,故此天使特別告訴他說:到二千三百日聖殿就被潔淨,這二千三百日是包括早祭與晚祭的次數,所以早與晚實際為一千一百五十日,有三年多的時間,亦即七章所記的一載,二載,半載.
\newline
\newline
(但八18)中記載「他與我說時,我便伏在地沉睡.」「沉睡」是整個昏過去,然後使者扶起後才接受神的啟示,那所付的代價又何等大!
\newline
\newline
當時但以理覺得恐懼,世界要滅亡了,到底以後如何?歷史負擔壓在但以理身上實非他所能擔當的；何況他想及神百姓將來之前途,在神公義彰顯時,又有誰能逃脫呢?
\newline
\newline
27節「於是但以理昏迷不醒,病了數日,然後起來辦理王的事務；我因這異象驚奇,卻無人能明白其中意思.」但以理雖不明白內容,可是知道是十分嚴重,他唯一能作的是忠心虔誠作應作的事.你我許多時也知神歷史中的作為,那我們能作些什麼呢?有時需要宣告審判之信息,有時需要禱告,有時我們做日常生活中十分平常的事,在神前忠心作應作的事,但以理就是如此,仍繼續辦理王的事,可是沒有多久這機會就失去,伯沙撒王遺忘了他,人是會忘記他,但神沒有忘記,時候到了他會出來為神作見證.當成功或興盛時,我們要謙卑；因為我們的一切,有什麼不是從神那兒領受的呢?既是領受,那又有什麼可誇的呢?總要謙謙卑卑地靠主的力量去作工.雖然有時我們被人忘卻,被忽略,甚至沒有機會見證神；但千萬不可自暴自棄,應在神面前好好地等候.時候到了,神會大大地使用我們,正如但以理一樣.
\newline
\newline
(但五1-12)中可見巴比倫王朝的敗亡,這是命運.世人之命運都在神手中,神使人升高,亦叫人降卑；神能使一個民族衰敗,亦能使一個民族興起.神恆久忍耐,又有恩慈,不輕易發怒.但當人不肯領受祂的恩慈時,祂忿怒的杯滿了；就必流溢出來,因神的公義必要顯露出來,因祂是公義的主.
\newline
\newline
伯沙撒王當時覺得很狂傲自滿,因他從未經過憂患,當然更不曉得何為屬靈的事,認識耶和華真神.他的外祖父尼尼甲尼撒王雖是暴虐之君王,但他的經歷使他知道一些公義的事,直至他屈服在神面前.但伯沙撒王生於安樂,且長於安樂；況且在他父王拿彼尼度的時代,巴比倫暫享安定.巴比倫的京城甚大,防禦好,城之周圍有十五英哩之防禦工事,城牆有三百五十呎,八十七呎闊,根本不可能有外邦軍隊攻破這永不敗亡的城.可惜這只不過是伯沙撒的想法,因當他最歡樂時；大利烏的軍隊已將巴比倫圍困.一個國家之敗亡並非剎那間的事,社會的敗壞及一個人的墮落,也不是一時的.一定有其敗壞的因素存在,當這些因素,發展下去,總有一天,那敗亡的事實就會突然的發生!這是必然的事實.當時伯沙撒的讌會是極不正常的,因臣僕的讌會不應有婦女參加,這暴虐的居王卻如此行.讌會中充滿了放蕩,甚至把從耶路撒冷擄來的,應分別為聖的聖殿的器皿,當作酒杯用；讚美偶像,罪惡的權勢似乎完全勝利,魔鬼在發出猙獰的狂笑!
\newline
\newline
可是,忽然間神的手就出現,神的手指在牆上寫字；剎那間,所有情況都改變,他們都十分害怕,不知該如何才好.於是太后(伯沙撒之母)將但以理介紹出來,可見當時但以理已被人遺忘；但神不但沒遺忘他,且叫他出來,神要他說話,要他責備.但以理到了巴比倫宮廷中不是已改了新的名字嗎,為何在此仍有他本來的名字?因為他的身份沒有失去.各位!你我的身份都不能失去,我們是屬神的人,神在我們身上都有一定的計劃,我們不能失去這尊貴的身份.人可輕看或藐視我們,但在必要時神要我們出來,請問我們都預備好了嗎?但以理隨時都預備好了,他聽到神的命令就挺身而出,仗義直言,把神公義的信息告訴不義的居王.
\newline
\newline
23至25節中,是解釋這文字「彌尼,彌尼,提客勒,鳥法珥新」,「彌尼」(有翻彌拿)意「數算」,「提客勒」(有翻舍克拉),意是「秤一秤」,「烏」意「以及」,「法珥新」(法珥新是現在式,毘勒斯是過去式),意「分開」,──「法珥新」之字根與「波斯」字音相同,即 P(h)-R-S.這是清楚說:要數算,要秤一秤,就要分開.
\newline
\newline
當神指頭在牆寫字時是指現在就要分裂,而但以理解說這已成事實,因神的話永不改變,安定在天,句句帶著能力.神叫我們數一數,秤一秤就顯出虧欠.「公義」裡字,本來是「天平」的意思,新約的啟示錄和舊約的但以理書,常有「數算」,這是啟示文學的特點；如一載,二載三載,又有二千三百日,及七十個七等；所以時間是在神的手中,公義的審判要來到.當晚伯沙撒殺,瑪代王大利烏就取了這國家.各位!你我都需有屬靈的深度才能看見神的計劃,我們要學習神給我們的方法；秤一秤這個世界,秤這個社會,秤我們的教會,還有秤我們的家庭和我們自己的心是怎樣.還要數算,求主教我們數一數我們的日子如何,好叫我們得著智慧的心.因為還有一點點的時間,那要來的就來,並不遲延,還有一點點的時間可作見證,這短短的時間我們能作些什麼?我們的心又如何?愛主嗎?還是愛自己愛世界?若有人不愛主,這人是可咒可詛的,因主來的時間快到了.
\newline
\newline


\newpage

\section{第45屆~港九培靈會~唐佑之~第六講}
\label{sec:1253}
\textbf{蒙愛與期望}
\newline
\newline
連結: \href{https://www.hkbibleconference.org/session-message/view/1253}{\texttt{https://www.hkbibleconference.org/session-message/view/1253}}
\newline
\newline
\hyperref[sec:1252]{< < < PREV SERMON < < <}
~
\hyperref[sec:toc]{[返目錄]}
~
\hyperref[sec:1255]{> > > NEXT SERMON > > >}
\newline
\newline
但以理書 9:1-9:27
\newline
\begin{longtable}{cl}
\hline
\hline
章節 & 經文 (和合本修訂版)\\
\hline
9:1 & \begin{tabularx}{0.7\textwidth}{X} 瑪代族亞哈隨魯的兒子大流士被立為王，統治迦勒底國元年， \end{tabularx} \\ \\ \relax
9:2 & \begin{tabularx}{0.7\textwidth}{X} 就是他在位第一年，我－但以理從書上得知，耶和華的話臨到耶利米先知，論耶路撒冷荒涼期滿的年數為七十年。 \end{tabularx} \\ \\ \relax
9:3 & \begin{tabularx}{0.7\textwidth}{X} 我面向主神，禁食，披麻蒙灰，懇切禱告祈求。 \end{tabularx} \\ \\ \relax
9:4 & \begin{tabularx}{0.7\textwidth}{X} 我向耶和華－我的神祈禱、認罪，說：「主啊，你是大而可畏的神，向愛主、守主誡命的人守約施慈愛。 \end{tabularx} \\ \\ \relax
9:5 & \begin{tabularx}{0.7\textwidth}{X} 我們犯罪作惡，行惡叛逆，偏離你的誡命典章， \end{tabularx} \\ \\ \relax
9:6 & \begin{tabularx}{0.7\textwidth}{X} 沒有聽從你僕人眾先知奉你的名向我們君王、官長、祖先和這地所有百姓所說的話。 \end{tabularx} \\ \\ \relax
9:7 & \begin{tabularx}{0.7\textwidth}{X} 主啊，你是公義的，但我們猶大人和耶路撒冷的居民，並你所趕到各國的以色列眾人，不論遠近，因為背叛了你，臉上蒙羞，正如今日一樣。 \end{tabularx} \\ \\ \relax
9:8 & \begin{tabularx}{0.7\textwidth}{X} 耶和華啊，我們和我們的君王、官長、祖先因得罪了你，臉上就都蒙羞。 \end{tabularx} \\ \\ \relax
9:9 & \begin{tabularx}{0.7\textwidth}{X} 主－我們的神是憐憫饒恕人的，我們卻違背了他， \end{tabularx} \\ \\ \relax
9:10 & \begin{tabularx}{0.7\textwidth}{X} 沒有聽從耶和華－我們神的話，沒有遵行他藉僕人眾先知向我們頒佈的律法。 \end{tabularx} \\ \\ \relax
9:11 & \begin{tabularx}{0.7\textwidth}{X} 以色列眾人都犯了你的律法，偏離、不聽從你的話；因此，你僕人摩西律法上所寫的詛咒和誓言傾倒在我們身上，因我們得罪了神。 \end{tabularx} \\ \\ \relax
9:12 & \begin{tabularx}{0.7\textwidth}{X} 神使大災禍臨到我們，實現了警戒我們和審判我們官長的話；原來耶路撒冷所遭遇的災禍是普天之下未曾有過的。 \end{tabularx} \\ \\ \relax
9:13 & \begin{tabularx}{0.7\textwidth}{X} 這一切災禍臨到我們，是照摩西律法上所寫的，我們卻沒有求耶和華－我們神的恩惠，使我們回轉離開罪孽，明白你的真理。 \end{tabularx} \\ \\ \relax
9:14 & \begin{tabularx}{0.7\textwidth}{X} 所以耶和華特意使這災禍臨到我們，耶和華－我們的神在他所行的事上都是公義的；我們並沒有聽從他的話。 \end{tabularx} \\ \\ \relax
9:15 & \begin{tabularx}{0.7\textwidth}{X} 主－我們的神啊，你曾用大能的手領你的子民出埃及地，使自己得了名聲，正如今日一樣，現在，我們犯了罪，作了惡。 \end{tabularx} \\ \\ \relax
9:16 & \begin{tabularx}{0.7\textwidth}{X} 主啊，求你按你豐盛的公義，使你的怒氣和憤怒轉離你的城耶路撒冷，就是你的聖山。因我們的罪惡和我們祖先的罪孽，耶路撒冷和你的子民被四圍的人羞辱。 \end{tabularx} \\ \\ \relax
9:17 & \begin{tabularx}{0.7\textwidth}{X} 我們的神啊，現在求你垂聽你僕人的祈禱懇求，為你自己的緣故使你的臉向荒涼的聖所發光。 \end{tabularx} \\ \\ \relax
9:18 & \begin{tabularx}{0.7\textwidth}{X} 我的神啊，求你側耳而聽，睜眼而看，眷顧我們那荒涼之地和稱為你名下的城。我們在你面前懇求，不是因自己的義，而是因你豐富的憐憫。 \end{tabularx} \\ \\ \relax
9:19 & \begin{tabularx}{0.7\textwidth}{X} 主啊，求你垂聽！主啊，求你赦免！主啊，求你側耳，求你實行！為你自己的緣故不要遲延。我的神啊，因這城和這民都是稱為你名下的。」 \end{tabularx} \\ \\ \relax
9:20 & \begin{tabularx}{0.7\textwidth}{X} 我正說話、禱告，承認我的罪和我百姓以色列的罪，為我神的聖山，在耶和華－我的神面前懇求； \end{tabularx} \\ \\ \relax
9:21 & \begin{tabularx}{0.7\textwidth}{X} 我正在禱告中說話，先前在異象中所見的那位加百列，約在獻晚祭的時候迅速飛到我這裡來。 \end{tabularx} \\ \\ \relax
9:22 & \begin{tabularx}{0.7\textwidth}{X} 他指教我說：「但以理啊，現在我來要使你有智慧，有聰明。 \end{tabularx} \\ \\ \relax
9:23 & \begin{tabularx}{0.7\textwidth}{X} 你剛開始懇求的時候，就有命令發出。現在我來告訴你，因你是蒙愛的；所以你要思想這事，明白這異象。 \end{tabularx} \\ \\ \relax
9:24 & \begin{tabularx}{0.7\textwidth}{X} 「為你百姓和你聖城，已經定了七十個七，要止住罪過，除淨罪惡，贖盡罪孽，引進永恆的公義，封住異象和預言，並膏至聖所。 \end{tabularx} \\ \\ \relax
9:25 & \begin{tabularx}{0.7\textwidth}{X} 你當知道，當明白，從發出命令恢復並重建耶路撒冷，直到受膏的君出現，必有七個七和六十二個七。耶路撒冷城連街帶濠都必在艱難中恢復並重建。 \end{tabularx} \\ \\ \relax
9:26 & \begin{tabularx}{0.7\textwidth}{X} 過了六十二個七，那受膏者被剪除，一無所有；必有一王的百姓來毀滅這城和聖所，它的結局必如洪水沖沒。必有戰爭，一直到末了，荒涼的事已經定了。 \end{tabularx} \\ \\ \relax
9:27 & \begin{tabularx}{0.7\textwidth}{X} 在一七之期，他必與許多人堅立盟約；一七之半，他必使獻祭與供獻止息。那施行毀滅的可憎之物必立在聖殿裡，直到所定的結局傾倒在那行毀滅者的身上。」 \end{tabularx} \\ \\
[1ex]
\hline
\hline
\end{longtable}
天使稱但以理為大蒙眷愛的人,誰才配你為大蒙眷愛的人?乃是一個肯禱告,真正信靠神而又將自己全交托神的人；真正全向真神支取神恩典,知神是信實的充滿慈愛的人；一個真正願將自己完全澆奠在主腳前的人.神要我們禱告,在此章書中,看見非常豐富的屬靈教訓；使我們能學習禱告的功課,好像但以理一樣,他真正將所有的心意向神傾吐.
\newline
\newline
當時的歷史是非常重要的時期,由巴比倫王朝過後,瑪代皮斯新強權正興起之時；在這交替的階段,但以理深感禱告的需要.但以理書不單說明將來末後之事,並說明在當今的日子.神的兒女們該做些什麼,這也說明末後事之奧秘,更進一步明白奧秘的經驗.不單是從神之啟示中有所領悟,更是在神面前有所獻呈.但以理書是特別著重禱告的書,例如在二章中但以理與三個朋友同心合意的禱告.六章中但以理個人的禱告生活,九章中但以理在神前之代求,十一章中但以理如何恆心持久的禱告.這一切都說明整卷書充滿了禱告的氣氛.當時大利烏作王,迦勒底已成為是瑪代波斯之一部份,時在主前五三九年,離但以理被擄六十多年.如果但以理是主前六○五年,則他已年老,他一生最寶貴的時間在外邦,無國家,無自由,無宗教的環境；在叛背神的暴虐君王下,在異教徒中,在政治暴力環境下.對許多青年而言,會隨波逐流,或勉強順應當前的環境,他卻好自為之,靠主始終站立得穩,始終成為時代的中流砥柱.什麼力量使他如此祂始終如一的依靠神?那是禱告的力量.祗有每天禱告,不住禱告,隨時多方地禱告；現今他更需要禱告,因他知道這是非常重要的歷史時期.
\newline
\newline
各位!我們是否對歷史負上責任?在聖經可以見到神所愛所使用的,總是對時代之轉變有非常敏銳感覺的人；是屬靈的覺得應該禱告,需真正懇切付代價禱告的人.但以理真正明白神的話和祂的心意,他說:在大利烏在位第一年,我但以理是從書上得知耶和華的話,(九2)從耶利米書中,耶利米先知在猶大敗亡時為神所立為列國的先知,因以色列歷史是神特別看重,所有外邦一切之演變皆與神之選民有關係,耶廿五章及廿九章皆提及耶路撒冷荒涼年數,以七十年為滿；不知但以理是從何算起,如果以他從被擄開始,七十年是從主前六○○五至五三九.可能另一算法更準確,就是耶路撒冷敗亡五八六或五八七年,一直到聖殿之重建五一五年,那麼算亦約為七十年.此乃神之作為及計劃,神叫百姓回耶路撒冷.神的話永遠可靠,這是但以理之信心,因他明白神之心意,並將所得的思想化為禱告.我們為何少禱告,禱告那麼膚淺,因為少得神的話語；缺少在神的話語方面下工夫,沒有屬靈的意識和思想,因而缺少禱告.屬神的人應該像但以理一般,真正在神的話語下工夫；真正得著神賜的話語,那時我們才願意禱告.那時我們必會逼切地懇求,因神的心意被我們摸著了.
\newline
\newline
禱告並非大聲呼叫就能表明心意,有時我們真想將內心的情緒,及思想等釋放出來,無法自制的.不過真實的禱告,乃是充滿神的話,真正在神的真道上造就自己；才知道如何在聖靈裡禱告,猶大書所說的在次序及果效都這樣.我們先要讓聖經充滿,才有聖靈充滿；但以理就是如此,試看他如何禱告:
\newline
\newline
(一)(但九3)「定意」的禱告:表示專心一意集中力量,很堅強地須要禱告.「我便禁食,披麻蒙灰,定意向主神祈禱懇求.」在福音書中記載主最後一次上耶路撒冷時,也是用這「定意」的字.雖然新舊約的用字不同,但意思卻是一樣,圖畫是一樣；確定的方向,定意向著神毅然決意,絕不猶疑的仰望神.
\newline
\newline
(二)懇切的祈求:他不單定意禱告,且懇切的祈求,迫切地求,正如主在客西馬尼園作付代價的禱告.假如我們能如此,則神的恩一定臨到我們,使多人蒙福.
\newline
\newline
(三)謙卑的禱告:他禁食披麻蒙灰,如此在神前痛悔的禱告,是非常澈底的在神前認罪悔改的態度.
\newline
\newline
(四)悟性的禱告:不是胡亂的思想,而是非常清楚地讓聖靈光照他；指示他,他在神前等候.他有屬靈的眼光可以看得見,心裡是何等地清楚到神前禱告.
\newline
\newline
我們讀此章(九章),是否注意其中有一特點,乃是別章所無的,在九:2,4,10,13,14,20幾節中,有提及耶和華的名.為什麼如此?因耶和華是以色列的神,是自有永有的神.與以色列民有聖約的關係,表明神信實的；我們人雖然失信,但神總不失信的,因神眷顧祂的百姓.人失敗了!但神永不失敗,信實的神,祂在歷史中要做成祂的工作.
\newline
\newline
但以理有此屬靈的思想及屬靈的了解,所以他在呼喚耶和華,耶和華......,神的名字好像倒出來的香膏韾香無比.我們在神前如提到祂的名字,心裡就歡喜,禱告就特別充實.
\newline
\newline
從14節後就沒有再提耶和華的名,從15至19節就提「主」神的名,表明另有一重點以表神的大能；提及耶和華時是表明神的信實與慈愛,提「主」是表明神的能力及權柄；再注意「你的城耶路撒冷是你的聖所,我們是你是子民；你要把你的臉光照在這荒涼城,求你要按著你的大仁大義.」在此禱告是完全集中於神,那是多麼寶貝的思想.所以他的禱告有功效,禱告中充滿了屬靈的思想又是何等的充實.我們不會禱告,也不知該如何禱告,求主教導我們如何禱告；感謝主!聖靈用說不出來的嘆息為我們代求.
\newline
\newline
(但九20-27),這是天使的答案,是最難解而又最寶貴之真理.迦百列天使對但以理說「你的話一出,神就發出命令.」你一禱告,神就出命令,我們如在神前真正禱告；神不能不垂聽,有什麼比禱告能力更大更有功效!因禱告的人必大蒙眷愛.不單要禱告,還要有屬靈思想,於是說及七十個七,為何叫七十個七?在猶大「七」乃完全之數,七十個七則更完全,因猶大人在摩西律法裡,神啟示以色列節期時；每七年有一安息年,七十個七就有七十個安息年.在安息年時地土不必種植,所有人不必勞苦.當以色列人猶大人違背神時,地土就荒涼,人就被擄在外邦；留下的也不能作什麼,共有七十年之久.人失敗,地就荒涼,但神永不失敗,在神計劃中有裡七十個安息年中；地土仍可安息,這七十個七分三期.
\newline
\newline
(一)第一個七.
\newline
\newline
(二)六十二個七.
\newline
\newline
(三)最後一個七.
\newline
\newline
許多研究聖經的人,愛研究這七十個七,有許多的寶貴意見,亦有許多不同的意見.但卻不可能有最後的結論；那並非神所講的不完全,而是人所領悟的不完全.我們人可能解釋錯誤,但神在歷史當中的時間,一點都沒錯.主耶穌說:你們的時候很方便,但我的時候還沒到.我們可以數算一下,但不一定準確.第一個七可以算得準,許多解經者皆認為對,共有四十九年.那是當以色列人敗亡時(主前五八七年)直等他們可重建耶路撒冷聖殿時(五二○年),約四十七至四十八年,約第七個七的時間.第二個是六十二個七,有許多不同之意見:大概在尼希米重建耶路撒冷城牆,約主前四四五年尼散月的初一日,直至受膏者被剪除時；即主耶穌進入耶路撒冷,承受十字架苦難之時.主後三十二年時,即尼散月初十日.約在十七萬三千八十天,如三十天為一個月；則為四八三年,即六十九個七.(六十二個七連前七個七,共六十九個七),這是一般保守派學者的算法,自由派的學認為這受膏者非指耶穌基督,而是指瑪喀比時的一位大祭司阿尼亞.若果如此算法,最末一個七在瑪客比被殺時就完了,故我們不能接受.我們相信此受膏者可能指君王,或祭司；此處必是指彌賽亞之預言,指主耶穌基督；這是我們所能見到的,只有到此為止.
\newline
\newline
最後一個七最難算,亦難作結論；因為不單是在屬靈方面上領悟力不夠,事實上歷史還在進展中.六十九個七已過去,現在正是在最後一個七內.是羅馬帝國開始直到主再來,在啟示錄提及,太廿四章,路廿一章及可十三章亦提及,可見最末一個七是很長的時期.
\newline
\newline
有人解釋,主後七十年時,耶路撒冷被提多將軍所毀滅,他們以為最後一個七已應驗；我們不能接受,這是猶太拉比的解釋.
\newline
\newline
另有人以為最後的七,在耶穌基督身上已經應驗.因主來世有三十三年半的時間,然後被釘十架.主復活後另外三年半,前一半是主道成肉身；另一半是主基督復活升天後教會之情況.這是教會時代,這也可能不正確.
\newline
\newline
那麼怎樣才正確?最後一個七之意義分成兩半,可能有更深之意思在內.為何六十九個七能以年代來計算?因有歷史事實存在.但若算遙遠之將來就很困難,正如站在山頭上,看見重山疊疊；好像每個山頭都連接一起,而其實則每個山與山之間,是有山有谷的,故不易計算最後一個七.可是一七之半的前半,我們能知道.在七章曾提一載二載半載,在八章也提二千三百日,有人以為有早晚祭,故應是一一五○日.又在啟十二章中提及一載二載半載,這可能是外邦人的日子.
\newline
\newline
你和我亦包括在頭一個七之半內,社會混亂,罪惡,魔鬼有很大的作為,神的兒女應如何儆醒禱告.另一個七之半,乃是主在(太廿四21)太廿四21之大災難,亦即耶卅17所說的雅各遭難的日子.願神憐憫我們,在這災難未來之前,知道在這時代我們所擔負的責任,正像先知一樣.從神領受話語,像寫廿一章內曾有從度瑪來的默示.這默示裡就聽見說:「守望的啊,夜裡如何?守望的啊,夜裡如何?」這是說夜間到底過了多少?早晨一定快來到,主一定來到,但在主未來之前,在大災難之前；天是那這幽暗,真是漫漫長夜!教會是金燈台,燈台只在黑夜裡照亮.我們都在晚上,無邊的黑暗中；守望的啊!夜裡如何?到底黑夜過了多少?為何黑夜是這麼的長?守望者回答說:早晨必快來到,榮耀的將來必來到.我們該存這樣的心,這樣的信念,並要這樣堅定的禱告.守望的一定要守望,因早晨必要來到,可是黑夜也會來到的.那就是在不信的人,將是永遠的黑夜!只一點點時間,那要來的就要來,祂並不遲延.所以我們每天仍要重複的說:若有人不愛主,那是可咒可詛的.我們愛祂,祂也必愛我們,我們無論求什麼,祂也必成全.我們都是大蒙眷愛的人,我們應該愛主,因主已經近了.
\newline
\newline


\newpage

\section{第45屆~港九培靈會~唐佑之~第七講}
\label{sec:1255}
\textbf{猛戰與天使}
\newline
\newline
連結: \href{https://www.hkbibleconference.org/session-message/view/1255}{\texttt{https://www.hkbibleconference.org/session-message/view/1255}}
\newline
\newline
\hyperref[sec:1253]{< < < PREV SERMON < < <}
~
\hyperref[sec:toc]{[返目錄]}
~
\hyperref[sec:1258]{> > > NEXT SERMON > > >}
\newline
\newline
但以理書 6:1-6:28
\newline
\begin{longtable}{cl}
\hline
\hline
章節 & 經文 (和合本修訂版)\\
\hline
6:1 & \begin{tabularx}{0.7\textwidth}{X} 大流士隨心所願，立了一百二十個總督，治理全國， \end{tabularx} \\ \\ \relax
6:2 & \begin{tabularx}{0.7\textwidth}{X} 又在他們以上立總長三人，但以理也在其中；使總督在他們三人面前呈報，免得王受虧損。 \end{tabularx} \\ \\ \relax
6:3 & \begin{tabularx}{0.7\textwidth}{X} 這但以理因有卓越的靈性，超乎其餘的總長和總督，王想立他治理全國。 \end{tabularx} \\ \\ \relax
6:4 & \begin{tabularx}{0.7\textwidth}{X} 那時，總長和總督在治國的事務上尋找但以理的把柄，為要控告他；只是找不到任何的把柄和過失，因他忠心辦事，毫無錯誤過失。 \end{tabularx} \\ \\ \relax
6:5 & \begin{tabularx}{0.7\textwidth}{X} 那些人就說：「我們要找但以理的把柄，若不從他神的律法中下手，就尋不著。」 \end{tabularx} \\ \\ \relax
6:6 & \begin{tabularx}{0.7\textwidth}{X} 於是，總長和總督紛紛聚集來見王，說：「大流士王萬歲！ \end{tabularx} \\ \\ \relax
6:7 & \begin{tabularx}{0.7\textwidth}{X} 國中的總長、欽差、總督、謀士和省長彼此商議，求王下旨，立一條禁令，三十天之內，不拘何人，若在王以外，或向神明或向人求甚麼，就必扔在獅子坑中。 \end{tabularx} \\ \\ \relax
6:8 & \begin{tabularx}{0.7\textwidth}{X} 王啊，現在求你立這禁令，在這文件上簽署，使它不能更改；照瑪代人和波斯人的例，絕不更動。」 \end{tabularx} \\ \\ \relax
6:9 & \begin{tabularx}{0.7\textwidth}{X} 於是大流士王在這禁令的文件上簽署。 \end{tabularx} \\ \\ \relax
6:10 & \begin{tabularx}{0.7\textwidth}{X} 但以理知道這文件已經簽署，就進自己的家，他家樓上的窗戶開向耶路撒冷。他一天三次，雙膝跪著，在他的神面前禱告感謝，像平常一樣。 \end{tabularx} \\ \\ \relax
6:11 & \begin{tabularx}{0.7\textwidth}{X} 於是，那些人紛紛聚集，發現但以理在他神面前祈禱懇求。 \end{tabularx} \\ \\ \relax
6:12 & \begin{tabularx}{0.7\textwidth}{X} 他們就進到王面前，向王提及禁令，說：「三十天之內不拘何人，若在王以外，或向神明或向人求甚麼，必被扔在獅子坑中，王不是在這禁令上簽署了嗎？」王回答說：「確有這事，照瑪代人和波斯人的例是不可更改的。」 \end{tabularx} \\ \\ \relax
6:13 & \begin{tabularx}{0.7\textwidth}{X} 他們對王說：「王啊，那被擄的猶大人但以理不理會你，也不遵守你簽署的禁令，竟一天三次祈禱。」 \end{tabularx} \\ \\ \relax
6:14 & \begin{tabularx}{0.7\textwidth}{X} 王聽見這話，就甚愁煩，一心要救但以理，直到日落的時候，他還在籌劃解救他。 \end{tabularx} \\ \\ \relax
6:15 & \begin{tabularx}{0.7\textwidth}{X} 那些人就紛紛聚集到王那裡，對王說：「王啊，當知道瑪代人和波斯人有例，凡王所立的禁令和律例都不可更改。」 \end{tabularx} \\ \\ \relax
6:16 & \begin{tabularx}{0.7\textwidth}{X} 於是王下令，人就把但以理帶來，扔在獅子坑中。王對但以理說：「你經常事奉的神，他必拯救你。」 \end{tabularx} \\ \\ \relax
6:17 & \begin{tabularx}{0.7\textwidth}{X} 有人搬來一塊石頭放在坑口，王用自己的璽和大臣的印，封閉那坑，使懲辦但以理的事絕不更改。 \end{tabularx} \\ \\ \relax
6:18 & \begin{tabularx}{0.7\textwidth}{X} 王回到宮裡，終夜禁食，不讓人帶樂器到他面前，他也失眠了。 \end{tabularx} \\ \\ \relax
6:19 & \begin{tabularx}{0.7\textwidth}{X} 次日黎明，王起來，急忙往獅子坑那裡去， \end{tabularx} \\ \\ \relax
6:20 & \begin{tabularx}{0.7\textwidth}{X} 臨近坑邊，哀聲呼叫但以理。王對但以理說：「永生神的僕人但以理啊，你經常事奉的神能救你脫離獅子嗎？」 \end{tabularx} \\ \\ \relax
6:21 & \begin{tabularx}{0.7\textwidth}{X} 但以理對王說：「願王萬歲！ \end{tabularx} \\ \\ \relax
6:22 & \begin{tabularx}{0.7\textwidth}{X} 我的神差遣使者封住獅子的口，叫獅子不傷我，因我在神面前無辜。王啊，在你面前我也沒有做過任何虧損的事。」 \end{tabularx} \\ \\ \relax
6:23 & \begin{tabularx}{0.7\textwidth}{X} 王因此就甚喜樂，吩咐把但以理從坑裡拉上來。於是但以理從坑裡被拉上來，身上毫無損傷，因為他信靠他的神。 \end{tabularx} \\ \\ \relax
6:24 & \begin{tabularx}{0.7\textwidth}{X} 王下令，把那些控告但以理的人和他們的妻子兒女都帶來，扔在獅子坑中。他們還沒有到坑底，獅子就制伏他們，咬碎他們的骨頭。 \end{tabularx} \\ \\ \relax
6:25 & \begin{tabularx}{0.7\textwidth}{X} 於是，大流士王傳旨給住在全地各方、各國、各族的人說：「願你們大享平安！ \end{tabularx} \\ \\ \relax
6:26 & \begin{tabularx}{0.7\textwidth}{X} 現在我降旨，我所統轄全國的人民，都要在但以理的神面前戰兢畏懼。因為他是活的神，永遠長存，他的國度永不敗壞，他的權柄永存無極！ \end{tabularx} \\ \\ \relax
6:27 & \begin{tabularx}{0.7\textwidth}{X} 他庇護，搭救，在天上地下施行神蹟奇事，救了但以理脫離獅子的口。」 \end{tabularx} \\ \\ \relax
6:28 & \begin{tabularx}{0.7\textwidth}{X} 如此，這但以理，當大流士在位的時候和波斯的居魯士在位的時候，大享亨通。 \end{tabularx} \\ \\
[1ex]
\hline
\hline
\end{longtable}
(但六1-11)在每一個時代內,神總尋找祂心意的人,那就是真正肯禱告及真把自己完全擺在主面前的人.神要我們成為當代的見證人,見證人須有雙重經歷:一方面是在幔子裡,另一方面是在營外.希伯來書對這兩方面的經歷十分注意,(來6章)中說,要進入幔子內,應與主有親密的靈交；真正了解主摸著祂的心意,這經歷很重要.還有來13章說:屬神的人,要出到營外,受祂所受之凌辱真正為主犧牲,這經歷與幔子內都很重要.
\newline
\newline
但以理六章是很熟悉的篇章,但以理這一經歷竟佔整整的一章,可見是何等重要.所以,雖然熟悉,但仍應思想,好讓我們得著新的亮光,以切實學習禱告的功課.
\newline
\newline
按時間而言,六章是在九章之後,因六章乃是在大利烏元年之時,相距時間雖不會很遠,但仍有先後之別.「大利烏乃歷史上最難解決的問題,因擄考古的發現及研究,在歷史上有三個大利烏.就是大利烏一世,二世,和三世.一世是在亞達薛西王及亞哈隨魯王之前,既在以前,必不是六章的大利烏.因九章中,大利烏乃亞哈隨魯王之子.又有人說大利烏二世是亞達薛西和亞哈隨魯之後,可是我們發現本章大利烏該在亞達薛西之前.大利烏三世則距此時更遠,應在主前三百多年前,所以這三個大利烏都不是但六章所說的.更有人說,大利烏乃是王的稱號,並非人名.(正如埃及王稱法老,羅馬王稱該撒)也有人說他就是古列,因大利烏的年代遲於以色列人回國之時.但在但以理書中說,大利烏王是瑪代人,古列是波斯人,那是有分別的.故根據學者考據而言,大利烏也許是將軍,可能其名是鳥格巴魯.他是第一個打倒巴比倫的,在主前五三九年十月初八日；十多天之後古列才到,然後委他管理加勒底國.大利鳥王治理時,立一百二十個總督；在上有三總長,但以理是其中的一個.為何如此重用他,是因他有美好的靈性.我們該在神前思想,怎樣才能作美好見證?先要有美好的靈性,因有美好的靈性,我們的才能方用得合式.神使用我們為主作見證,有些人靈性好,但才幹差；以致有人認為靈性與知識不能均衡,以為靈性好則知識一定低落.並不盡然,試看但以理和保羅,不也是有學問的人嗎?他們都被神大大的使用.我們不能自暴自棄,應靠主恩好自為之.大利烏看重但以理,因他作事認真又能幹,最重要的還是他生活聖潔,作事忠心又廉潔.可能其他官員是貪官污吏,對但以理之為人不滿意,於是用盡方法對付他.在六:6中,「總長和總督紛紛聚集來見王.」他們既不能找但以理的錯,故另用計謀,紛紛聚集!意即囂張熱鬧,像很興奮般對大利烏王說:「王啊!我們是想尊敬你的,因你是最偉大的王,我們擁護愛戴你.故在三十日內,不拘何人,若在王以外或向神或向人求甚麼,就必扔在獅子坑中.」王聽了當然興奮,因他作王不久,正需要權勢及安定的政治,需要幹部鞏固他的地位.這些人現今如此地擁護他,真是求之不得!他也許太興奮而竟上了圈套；因這些人中缺少了但以理王都沒有注意到.為何但以理不在一起?當時他好像什麼都不理會.世人常會如此作急速的決定,但後來卻又後悔.只有屬神的人眼睛才專注在主前,然後讓主引導,不慌張,不隨意,我們也該在主前平靜安穩等待,正如但以理一般.當但以理知道這消息,他仍很安靜,「但以理知道這禁令蓋了玉璽,就到自己家裡」(六10),在這情形之下但以理實在再無可選擇的餘地就回家去了,人之人格及靈性,不在乎行動(Action)面在乎反應,(Reaction),但以理所作的別人對他有了反應,可是他卻仍繼續下去,這表明他自己所作是對的.各位!當你的行動有人反對時,你該如何呢?是仍繼續下去?抑或去適應環境呢?於是我們的本性因此就顯出來了.我們信神,到教會聚會,在教會事奉,這是行動；但當遭遇患難,逼害,反對,辱罵時,你的反應如何?但以理的反應是絕不更改,絕不搖動的.寶貴不能淫,威武不能屈,始終如一,絕不妥協,絕不更改；照常一樣地禱告,因禱告是他的生活,他的反應就是照常禱告.我們平穩時健康時不肯禱告,當我們遭遇困難時有病時才去禱告,這種禱告如何能成為生活,怎會有功效,怎會蒙神悅納!又怎能讓神清楚答覆我們呢?平常不禱告,不讀神的話,當困難時或不知何去何從時才去尋求神；這又怎能明白神旨意!願神憐憫我們,教我們好好地禱告.但以理有好的禱告生活,好的讀經生活,他知道仰望神及呼求神.因專心尋求神就可尋見,他的窗戶向著耶路撒冷,乃是表示他心中所有的盼望.
\newline
\newline
(代下六34-39)說及假如神的百姓遭遇困難時,要好好地向耶路撒冷的聖所禱告,神就必允許.但以理一日三次禱告,並非故意要給人看見,只不過是出於自然.屬靈的事是繼續的,恆久的,始終如一,絕不更改；故他仍舊禱告,依舊在神前有仰望的心.當他禱告時,別人紛紛聚集,看但以理如何在神前禱告祈求.因他們亦知道但以理是不會改變的,他們囂張地到但以理面前,可是但以理並不因此而不再禱告；他依然禱告,始終如一,表裡一致.他的見證是自然的,他的見證並非說出來,也非做作出來.你我常很軟弱,不願作見證,也不敢作,沒時間作；就算有也是沒內容的,連自己也感到不自然的,所以見證就沒功效.許多人在台上說得振振有詞,頭頭是道,滔滔不絕,何等真實動聽；可是他們是否真正在生活上有所見證?但願我們多為傳道人禱告,讓他們能有表裡一致的見證,傳揚神的道.
\newline
\newline
如果生活失敗,見證亦必失敗,見證是自然的流露出來,像但以理一樣.試看一些假花,看來很美很完善,花瓣葉子都完整,甚至可以亂真那麼如何分別真假?太完全的毫無殘缺的該是假的,因真的總會在有些部份上有缺點.我們的見證時時也是如此,做得太好,不自然不得真實,但以理的見證是自然的,所以有功效.
\newline
\newline
(但六16-24)中,在那天晚上,王不能睡覺,因心裡沒有平安；可是但以理卻有平安,因他知道自己無辜,在主前亦沒有行過虧心之事.我們如何才得平安?是應在生活上聖潔,這樣無論在任何環境或遭遇任何困難,都可在神的前平靜安穩.我們的心和眼也不高大狂傲,正如斷奶的孩子在母親的懷中.相信但以理當晚全將自己交托神,他已年老,作了不少的事,在靈性上已到了爐火純青的程度；沒有不可交托的了,所以完全在神的愛中安息.次日早晨王到獅子坑邊哀聲呼喚但以理,他聽不見任何聲音,甚至連獅子的動作也聽不見,那多奇妙!神封住獅子的口,不但不咬但以理,連聲音也不能發出.王說:「你所事奉的神能救你脫離獅子嗎?」但以理回答說:「神已差祂的使者封住獅子的口,因我信靠我的神!」,由此可知,神僕及神的兒女就是如此地蒙神保守!
\newline
\newline
當但以理從獅子坑出來後,是否興奮地見證神的保守呢?不是的.當大利烏將那些控告他的人放在獅子坑裡咬死後,他幸災樂禍地說他們該死嗎?也沒有.一個人的人格考驗在乎他的反應.但以理的反應是絕不更改,他只回到家裡,仍向著窗戶,跪在神前禱告,他毫無改變,整個晚上在獅子坑裡的經驗沒使他改變.不論在幔子內或在營外,始終如一；在神前仰望,與神親近,他就是有這屬靈的生活.這是多麼的寶貴,他的見證又多麼的有效!
\newline
\newline
獅子坑現今在何處?早已淹沒在歷史的廢墟裡.但我們知道,今日的世界,不就正如一大獅子坑嗎,魔鬼如同吼叫的獅子,要吞喫我們.(提後四17)「主站在我旁邊,加給我力量,使福音被我盡都傳明,叫外邦人都聽見,我也從獅子口裡被救出來.」,保羅實在有這經驗.整個世界就是獅子坑,彼得書信說:萬物結局近了,你們要謹守儆醒禱告.彼前五8「你們要謹守儆醒,因為你們的仇敵魔鬼如同吼叫的獅子,要吞喫你們.」難道我們真的毫無感覺嗎?
\newline
\newline
當主在曠野試探的時候,有野獸亦有天使與祂在一起(可一13).主雖是神子,但祂降世成為人,祂的處境就是在野獸與天使之間.這也是你我的處境,我們有時像野獸般的性情,我們的情操何等卑下,不正常,不清潔.思想與言行會如禽獸一般,但另一方面,我們能得主的生命與主的恩典,就會像天使般那樣光明,美好,我們常在這兩者之間.司提反殉難之前,當他真正被聖靈充滿說出神的話,為逼害他的人禱告,他的面貌如天使一般.我們是像禽獸?抑或像天使?願神憐憫我們!尼布甲尼撒王貴為君王,在但四章中,他就變為野獸一樣.自從罪惡進入世界後,就不再是樂園而是曠野,人就在這曠野,常在野獸與天使之間.
\newline
\newline
但七章中,但以理豈不是看見獸上來嗎?可是,另又看見神的使者,有天使站在他旁邊解釋,並安慰剛強他.在此可見獅子坑裡的景況,一方面見有獅子野獸,另一方面亦有神的使者封住獅子口.各位弟兄姊妹!這是你我的處境,我們實在太危險了,因為到處都有獅子；但我們卻又很安全,因有神的使者與我們同在.只有仰望神的憐憫禱告,一直禱告,靠著主有說不盡的恩賜,及無限的能力幫助我們.
\newline
\newline


\newpage

\section{第45屆~港九培靈會~唐佑之~第八講}
\label{sec:1258}
\textbf{爭戰與得力}
\newline
\newline
連結: \href{https://www.hkbibleconference.org/session-message/view/1258}{\texttt{https://www.hkbibleconference.org/session-message/view/1258}}
\newline
\newline
\hyperref[sec:1255]{< < < PREV SERMON < < <}
~
\hyperref[sec:toc]{[返目錄]}
~
\hyperref[sec:1259]{> > > NEXT SERMON > > >}
\newline
\newline
但以理書 10:1-10:21
\newline
\begin{longtable}{cl}
\hline
\hline
章節 & 經文 (和合本修訂版)\\
\hline
10:1 & \begin{tabularx}{0.7\textwidth}{X} 波斯王居魯士第三年，有話指示那稱為伯提沙撒的但以理。這話是確實的，指著大戰爭；但以理明白這話，明白這異象。 \end{tabularx} \\ \\ \relax
10:2 & \begin{tabularx}{0.7\textwidth}{X} 那時，我－但以理悲傷了三個七日； \end{tabularx} \\ \\ \relax
10:3 & \begin{tabularx}{0.7\textwidth}{X} 美味我沒有吃，酒和肉沒有入我的口，也沒有用油抹我的身，直到滿了三個七日。 \end{tabularx} \\ \\ \relax
10:4 & \begin{tabularx}{0.7\textwidth}{X} 正月二十四日，我在大河，就是底格里斯河邊， \end{tabularx} \\ \\ \relax
10:5 & \begin{tabularx}{0.7\textwidth}{X} 舉目觀看，看哪，有一人身穿細麻衣，腰束烏法的純金腰帶。 \end{tabularx} \\ \\ \relax
10:6 & \begin{tabularx}{0.7\textwidth}{X} 他的身體如水蒼玉，面貌如閃電，眼目如火把，手臂和腳如明亮的銅，說話的聲音像眾人的聲音。 \end{tabularx} \\ \\ \relax
10:7 & \begin{tabularx}{0.7\textwidth}{X} 我－但以理一人看見這異象，跟我一起的人沒有看見，卻有極大的戰兢落在他們身上，他們就逃跑躲避， \end{tabularx} \\ \\ \relax
10:8 & \begin{tabularx}{0.7\textwidth}{X} 只剩下我一人。我看見這大異象就渾身無力，面容變色，毫無氣力。 \end{tabularx} \\ \\ \relax
10:9 & \begin{tabularx}{0.7\textwidth}{X} 我聽見他說話的聲音；一聽見他說話的聲音，我就沉睡，臉伏於地。 \end{tabularx} \\ \\ \relax
10:10 & \begin{tabularx}{0.7\textwidth}{X} 看哪，有一隻手摸我，使我膝蓋和手掌戰抖。 \end{tabularx} \\ \\ \relax
10:11 & \begin{tabularx}{0.7\textwidth}{X} 他對我說：「蒙愛的但以理啊，要思想我對你所說的話，只管站起來，因為我現在奉差遣來到你這裡。」他對我說這話，我就戰戰兢兢地站起來。 \end{tabularx} \\ \\ \relax
10:12 & \begin{tabularx}{0.7\textwidth}{X} 他說：「但以理啊，不要懼怕！因為自從第一日你立志要明白，又在你神面前刻苦自己，你的話已蒙應允；我就是因你的話而來。 \end{tabularx} \\ \\ \relax
10:13 & \begin{tabularx}{0.7\textwidth}{X} 但波斯國的領袖攔阻了我二十一天。看哪，天使長中的一位米迦勒來幫助我，因為我被留在波斯諸王那裡。 \end{tabularx} \\ \\ \relax
10:14 & \begin{tabularx}{0.7\textwidth}{X} 現在我來，要使你明白你百姓日後必遭遇的事，因為這異象關乎未來的日子。」 \end{tabularx} \\ \\ \relax
10:15 & \begin{tabularx}{0.7\textwidth}{X} 他向我這樣說，我就臉面朝地，啞口無聲。 \end{tabularx} \\ \\ \relax
10:16 & \begin{tabularx}{0.7\textwidth}{X} 看哪，有一位形狀像人的，摸我的嘴唇，我就開口說話，向那站在我面前的說：「我主啊，因這異象使我感到劇痛，毫無氣力。 \end{tabularx} \\ \\ \relax
10:17 & \begin{tabularx}{0.7\textwidth}{X} 我主的僕人怎能跟我主說話呢？我現在渾身無力，毫無氣息。」 \end{tabularx} \\ \\ \relax
10:18 & \begin{tabularx}{0.7\textwidth}{X} 有一位形狀像人的再一次摸我，使我有力量。 \end{tabularx} \\ \\ \relax
10:19 & \begin{tabularx}{0.7\textwidth}{X} 他說：「蒙愛的人哪，不要懼怕，願你平安！你要剛強！要剛強！」他一對我說話，我就覺得有力量，說：「我主請說，因你使我有力量。」 \end{tabularx} \\ \\ \relax
10:20 & \begin{tabularx}{0.7\textwidth}{X} 他說：「你知道我為甚麼到你這裡來嗎？現在我要回去與波斯的領袖爭戰，我去了之後，看哪，希臘的領袖必來。 \end{tabularx} \\ \\ \relax
10:21 & \begin{tabularx}{0.7\textwidth}{X} 但我要將那記錄在真理之書上的話告訴你。除了你們的天使米迦勒之外，沒有人幫助我抵擋他們。」 \end{tabularx} \\ \\
[1ex]
\hline
\hline
\end{longtable}
(但十1-9)中,這是但以理最後的異象──這異象從十章至十二章有長篇的記載.可說本書是最主要之啟示真理,除啟示錄之外,再沒有別處經文如此清楚.所以這信息甚為重要.神給祂的僕人見異象,在舊約中有不同的用字一看見──屬靈的看見.就是有屬靈的眼光及見解,能見別人所不能見的或別人所看不清的.
\newline
\newline
當我們真的看見時,心裡就會有特殊之感受,那我們才能向神仰望,等候和順服.另一是「默示」原意是「負擔」有重的負擔.是人所不能承擔,定要卸下,須有堅強之心意；這是一種艱苦之過程,在耶利米書中就有這樣特別的解釋.因他覺得環境困難,不知該如何做,甚至想放棄不傳神的話；可是卻又不能,因神的話在他裡面,像火在骨頭內一樣；無法閉塞得住,一定要說出來.有異象就有話語,有話語就有信息.箴言書說:沒有異象民就放肆.亦有說:沒有異象民就消滅,民就死亡.這異象的意思並非能非常之景象,而是神的話語及心意,心裡明白及順服.先知們不是單見特殊的經驗就算了,而是實在見到異象,且化為話語.因「神的話安定在天,永不改變.」「神的話沒有一句不帶著能力.」「神的話就是靈就是生命.」在我們裡面發生影響和果效.
\newline
\newline
現在我們看這最後的異象,不單見其中的圖畫,並要見其中的話語和信息.當時是古列王及大利烏王,這兩王的時間很難算；可能同時代的,亦可能同是一個人──有各不同的說法.但無論如何,聖經告訴我們,古列王發出命令叫以色列歸回.在以賽亞書中,我們發現神對以賽亞所說的話:古列乃是神的僕人,神所膏立的,那又怎樣解釋呢?因古列是外邦人,他是異教徒,但神為以色列人用這君王,讓神歷史之計劃能澈底地完全實行出來.根據歷代志下,以斯拉,尼希米所記之背景,古列發出命令後,以色列人歸回,這是應驗耶利米廿五至廿七章所說的預言.內裡記七十年以色列人歸回,耶路撒冷荒涼要恢復.現在他們已歸回,那時神向但以理顯異象,但以理的感受又如何?應該充滿了喜樂才對,因神的應許已實現,神的選民可回到神所應許之地,以色列民可恢復敬拜,那是多使人興奮的事!可是在當時,但以理卻悲傷三個七日,為什麼?因他看到當時之情況,神的時候到了,神的計劃要實現,神要大大地施恩典.神的怒氣不過是轉眼之間,他的恩典乃是一生之久,永遠長存.可惜以色列人沒預備好,許多不肯回去；為什麼?因為他們生活已安定了.巴比倫乃是個物質豐富之地,生活方便,謀生不難；所以他們再不願返回故鄉,寧願在這無國家無民族身份的地方,只要生活安定就行了!今日世人也如此,甚至基督徒亦然,只要有安定的生活,其他一切都並不重要.但神要他們回去,而他們又不願去,那是何等可悲!神為我們預備許多屬靈的恩典,而我們所要的卻只是物質生活.我們把靈性的生活放在次要的地位,以至許多的恩典及福份無法賜給我們,這是多麼大的損失.但以理見到這情形,心甚悲傷.為何但以理自己又不回去?是因他正作迦勒底作官,生活舒適而不願回去嗎?可能不是的,也可能這是他悲傷的緣故,他想回去但又不能回；正如摩西在毘斯迦山頂看見應許地而不能去.也許年紀老了不易走動(那時約有九十歲).他是何等希望與歸回的以色列人同行去建造聖殿.他見有許多以色列人不肯回去,也看見那些回去的人何等需要屬靈的帶領!神的時候已到,神的恩典已預備好,這種福份以色列人卻未預備去領受,那是多麼可惜,可悲及可憐!神應許許多恩典給我們,而我們不去接受；我們不願去,也不敢去.我們不肯付代價,不肯用信心的眼光去看祂,不敢伸出信心之手去領受.結果,就落在可憐可怕的光景中.
\newline
\newline
但以理是個屬靈的人,他心裡悲傷,因他見到許多別人所未能了解的事.在此時,他再沒別的辦法,他只有禱告,禱告是他的生活.當他在巴比倫時,他房子的窗戶向著耶路撒冷,他想望耶路撒冷,他等候更美的家鄉.在巴比倫城中有許多榮華美麗的景象他看不見；他所見的只是耶路撒冷,這是他心靈的眼.
\newline
\newline
看此處所記載的,他似乎不在城內,而是在希底結大河邊──希底結是巴比倫城外.可能城外有他居住的地方,當他有機會在城外時就禱告.他禱告的時間長達三個七日,即二十一日,一直在神前克苦己身地禱告.禱告時就是異象象中見到榮耀之主,因祂是無所不在的；是萬王之王,萬主之主,是管歷史的主.但以理見一人身穿細麻衣,腰束烏法精金帶,這是在啟一章中所描寫人子的圖畫.神的啟示何等奇妙,始終一樣地使人看見這榮譽的主,穿細麻衣,腰束金帶,因祂是位祭司.主為我們的大祭司,在天上為我們代求,那我們的禱告才有功效.
\newline
\newline
祂身體如水蒼玉,是透明的,是那麼榮耀美麗,面貌如閃電,眼目如火把,手和腳如光明的銅(但十6),這與啟一章所描寫的一樣,眼目明亮,看得那麼透澈,充滿公義與慈愛.「手腳如光明的銅」表明有力量,很潔淨；因銅是煉過的,所以光明.「說話如大眾之聲」,在啟示錄所說是「眾水的聲音」.表明聲音宏亮有力,而且是多種的聲音合起來；是和諧地合為一個聲音,一個啟示的信息.試看神的真理何等奇妙!人可從不同的角度,用不同的方法研究,而去解釋神的話但發現前後呈不矛盾,神的心意不會出爾反爾的.神的心意永不改變,祂有奇妙的心意要我們了解.
\newline
\newline
但以理見了異象以後就混身無力,甚至經受不住.我們該知道天上的異象,實在是地上的人所不能經受的.但我們也該知道,在神的愛中沒有懼怕,所以當讀下面經文時就可以知道,曾兩次提及「大蒙眷愛的但以理阿」.愛裡沒有懼怕,神藉祂的使者安慰祂的僕人.
\newline
\newline
(但十10-21)「大蒙眷愛的人阿!不要懼怕,願你平安,你總要堅強.」主要我們在祂面前有平安,也要堅強,從祂那裡得力量,才能聽祂的話.為什麼要用二十一日之時間禱告,這好像神的幫助及能力全未來到.經上說因上面有爭戰,這是屬靈的爭戰,我們禱告時神已立刻垂聽.天使已說明了當第一日專心求明白將來的事時,我就因你的言語而來；可是因有攔阻,那是波斯的魔君和希腊的魔君；他們是撒旦的使者要控制世界,凡不信神的都在撒旦控制之下.他要控制世上執政掌權的人,他們不擇手段以達強權的目的,這正好中了魔鬼的控制下,不過是魔鬼的奴僕而已.那我們這屬神的人又如何呢?竟在魔鬼的奴僕之下!受他們迫害,欺壓!我們何等需要神的憐憫,神能使我們得勝.我們要禱告,天上的爭戰一定會解決,上面的爭戰解決了才能解決下面的爭戰.我們在此也看見,這種爭戰甚至天上的使者也感困難,直至天使長中的一位米迦勒才能得勝.可見這是劇烈的爭戰,在羅八章及弗六章都有說及與空中屬靈氣的惡魔爭戰.是我們所不能得勝的,除非神的恩典及能力臨到我們.
\newline
\newline
我們需要禱告,禱告能使們得勝,但禱告就是爭戰；地上在爭戰時需要有天上爭戰,才能見神的能力顯出.我們是否有這種爭戰的經驗?假若沒有,則難有禱告的生活.試看主在客西馬尼園之禱告,是多麼艱苦的戰爭.
\newline
\newline
(林後十3-4)「因我們在血氣中行事,卻不憑著血氣爭戰.我們爭戰的兵器乃是在神面前有能力,可以攻破堅固的營壘.」
\newline
\newline
怎樣才算是不憑屬血氣的爭戰?怎樣的禱告才算不憑血氣呢?
\newline
\newline
(但十12)中是但以理的禱告,也說明了一種真正倚靠神的禱告,有三方面值得我們注意的:
\newline
\newline
(一)克苦己心──這表明情緒方面,這情緒要靠聖靈引導,帶到聖靈裡面,在聖靈裡面禱告.情緒有時可能在聖靈之外,那是魂,那是屬亞當的,是人的活動；這種情緒很危險,是不能體會到屬靈的經驗.所謂克苦己心,是將「己」在神前對付清楚；神才能在我們裡面作主作王,聖靈控制我們整個情緒.
\newline
\newline
保羅在(弗五18)說:「不要醉酒,乃要被聖靈充滿.」為何醉酒會與聖靈充滿作比較?因人喝酒後情緒會興奮,那是一種犯罪情緒的表現.當一個人將他的情緒真正放在主手中時,完全讓聖靈掌管及充滿時,也會興奮.但有禱告,但讚美感謝和歌唱；這就全不屬魂的,乃是在靈裡享受的恩惠,真正要憑神恩典禱告時；就該切實對付,克苦己心的在神面前.
\newline
\newline
(二)要求明白──求明白是有理性的,有理智的,理智未必與信心衝突.假如人是憑血氣,則那種理性是與信心有衝突；但我們的思想也全交給主,由主管理.那我們的思想就有屬靈的成分,有深度的能明神的事,能參透萬事.
\newline
\newline
(三)要專心──專心禱告,這是意志方面,由神管理我們,以祂的道路為我們的道路.因神的道路非同人的道路,神的意念非同人的意念.完全奉獻,讓主完全控制我們整個的人格,情緒,意志,理情,三方面都全屬神時,就可得神的能力.
\newline
\newline
以下的經文,我們可以看到天使曾三次使但以理得能力:
\newline
\newline
(一)第10節中說:「猛然有一手按在我身上,使我用膝和手掌,支持微起.」這裡本來所用的字句是:有一手按在我身上,搖我的膝,很用力地搖,經這一搖,但以理就剛強起來.我們需要主的力量震動我們,在此特別注重膝蓋,「膝蓋」是指禱告,我們要禱告才有力量.
\newline
\newline
(二)第15-16節,「有一位像人子摸我的嘴唇」,為什麼要摸嘴唇?乃是讓他開口說話,開口發問題,開口禱告,及開口傳信息,我們多麼需要神的恩惠來摸我們的嘴唇.
\newline
\newline
(三)第19節:「大蒙眷愛的人哪!不要懼怕,願你平安,你要堅強；他一說話,我便覺得有力量.」單是說話就有能力,因他已有力量聽神的話,有能力禱告.我們必需好好地在地上禱告,在禱告中爭戰,然後上面的爭戰才能幫助我們.
\newline
\newline
在出十七章中,當亞瑪力人在利非訂和以色列人爭戰.他們在非訂山谷下面打仗,那麼如何決勝負?勝敗的決定是在乎上面而非下面,山谷上面有亞倫,戶珥扶著摩西的手在禱告.摩西何時舉手禱告,則以色列人在下面就得勝；上下要配合,上面與下面都須經過爭戰.
\newline
\newline
我們既知道上面有主在堅持的爭戰,那麼在下面的我們又該如何?我們為何禱告無能力,也不見神復興的力量臨到?乃因沒有付代價,沒有真正在禱告中爭戰.各位!我們是基督的精兵,需要為主爭戰,靠主爭戰.
\newline
\newline
有一位傳道人曾說:要得著禱告的祝福,必需有血,有淚,有汗的禱告.也就是需付重大的代價,讓我們禱告,在神前竭力地爭戰；因只有一點點的時候,那要來的要來,祂不遲延.在這最短的時間裡,是最危險,最可怕,最容易跌倒的時候.如再不儆醒,那又如何到神面前交帳?若我們再不愛主,那實在是可咒可詛的主已經近了!
\newline
\newline


\newpage

\section{第45屆~港九培靈會~唐佑之~第九講}
\label{sec:1259}
\textbf{星光與永福}
\newline
\newline
連結: \href{https://www.hkbibleconference.org/session-message/view/1259}{\texttt{https://www.hkbibleconference.org/session-message/view/1259}}
\newline
\newline
\hyperref[sec:1258]{< < < PREV SERMON < < <}
~
\hyperref[sec:toc]{[返目錄]}
~
\hyperref[sec:1225]{> > > NEXT SERMON > > >}
\newline
\newline
但以理書 11:1-12:13
\newline
\begin{longtable}{cl}
\hline
\hline
章節 & 經文 (和合本修訂版)\\
\hline
11:1 & \begin{tabularx}{0.7\textwidth}{X} 「至於我，當瑪代的大流士元年，我曾起來扶助米迦勒，使他堅強。 \end{tabularx} \\ \\ \relax
11:2 & \begin{tabularx}{0.7\textwidth}{X} 現在我要指示你確實的事。」「看哪，波斯還有三個王要興起，第四王必富足遠勝諸王。他因富足成為強盛，就煽動各國攻擊希臘國。 \end{tabularx} \\ \\ \relax
11:3 & \begin{tabularx}{0.7\textwidth}{X} 必有一個勇敢的王興起，執掌大權，隨意而行。 \end{tabularx} \\ \\ \relax
11:4 & \begin{tabularx}{0.7\textwidth}{X} 他正興起的時候，他的國必瓦解，向天的四方裂開，卻不歸他的後裔，也不如他當年統治的權威；他的國必被拔出，歸給他後裔之外的人。 \end{tabularx} \\ \\ \relax
11:5 & \begin{tabularx}{0.7\textwidth}{X} 「南方的王必強盛，他的將帥中必有一個比他更強，執掌權柄，權柄甚大。 \end{tabularx} \\ \\ \relax
11:6 & \begin{tabularx}{0.7\textwidth}{X} 過了幾年，他們必結盟，南方王的女兒必來到北方王那裡，使約生效；但這女子不能保留實力，王的力量也未能存留。這女子、帶她來的、生她的和當時扶助她的必被殺害。 \end{tabularx} \\ \\ \relax
11:7 & \begin{tabularx}{0.7\textwidth}{X} 但從這女子的本家必另有一子接續王位，他要率領軍隊進入北方王的堡壘，攻擊他們，而且得勝， \end{tabularx} \\ \\ \relax
11:8 & \begin{tabularx}{0.7\textwidth}{X} 把他們的神像和鑄成的偶像，與金銀寶器都擄掠到埃及去。數年之內，他不去攻擊北方的王。 \end{tabularx} \\ \\ \relax
11:9 & \begin{tabularx}{0.7\textwidth}{X} 北方的王必侵入南方王的國土，但卻要撤回本地。 \end{tabularx} \\ \\ \relax
11:10 & \begin{tabularx}{0.7\textwidth}{X} 「北方王的兒子們必動干戈，招聚許多軍兵。他要前進，如洪水氾濫；要再度爭戰，直搗南方王的堡壘。 \end{tabularx} \\ \\ \relax
11:11 & \begin{tabularx}{0.7\textwidth}{X} 南方王必發烈怒，出來與北方王爭戰，擺列大軍；北方王的軍兵必敗在南方王的手下。 \end{tabularx} \\ \\ \relax
11:12 & \begin{tabularx}{0.7\textwidth}{X} 這大軍既被掃蕩，南方王的心就自高；他雖使萬人仆倒，卻不能保持勝利。 \end{tabularx} \\ \\ \relax
11:13 & \begin{tabularx}{0.7\textwidth}{X} 「北方王要再度擺列大軍，比先前更多。過了幾年，他必率領大軍，帶極多的裝備而來。 \end{tabularx} \\ \\ \relax
11:14 & \begin{tabularx}{0.7\textwidth}{X} 那時，必有許多人起來攻擊南方王，並且你百姓中的殘暴人要興起，應驗異象，他們卻要敗亡。 \end{tabularx} \\ \\ \relax
11:15 & \begin{tabularx}{0.7\textwidth}{X} 北方王必來建土堆攻取堅固城，南方的軍兵抵擋不住，就是精選的部隊也無力抵擋； \end{tabularx} \\ \\ \relax
11:16 & \begin{tabularx}{0.7\textwidth}{X} 前來攻擊南方王的必任意而行，無人在北方王面前站立得住。他要站在那佳美之地，用手施行毀滅。 \end{tabularx} \\ \\ \relax
11:17 & \begin{tabularx}{0.7\textwidth}{X} 「他必定意傾全國之力而來，與南方王訂約，把自己的女兒給南方王為妻，企圖敗壞他的國度。這計謀卻未得逞，自己也得不到好處。 \end{tabularx} \\ \\ \relax
11:18 & \begin{tabularx}{0.7\textwidth}{X} 其後北方王必轉頭，奪取許多海島。但有一將帥除掉北方王對人的羞辱，並且使羞辱歸到他自己身上。 \end{tabularx} \\ \\ \relax
11:19 & \begin{tabularx}{0.7\textwidth}{X} 他必轉頭回到本地的堡壘，卻要絆跌仆倒，歸於無有。 \end{tabularx} \\ \\ \relax
11:20 & \begin{tabularx}{0.7\textwidth}{X} 「那時，有一人興起接續他的王位，他為了王國的榮華，差官員橫征暴斂。這王過不多時就死了，不是因怒氣，也不是因戰役。」 \end{tabularx} \\ \\ \relax
11:21 & \begin{tabularx}{0.7\textwidth}{X} 「後來，有一個卑鄙的人興起接續他的王位，人未曾將國的尊榮給他，他卻趁人無備的時候前來，用詭詐奪取政權。 \end{tabularx} \\ \\ \relax
11:22 & \begin{tabularx}{0.7\textwidth}{X} 勢如洪水般的軍兵在他面前被沖沒，遭擊潰；立約的領袖也是如此。 \end{tabularx} \\ \\ \relax
11:23 & \begin{tabularx}{0.7\textwidth}{X} 他與人結盟之後，卻行詭詐。跟隨他的人雖不多，他卻日漸強盛。 \end{tabularx} \\ \\ \relax
11:24 & \begin{tabularx}{0.7\textwidth}{X} 他趁人無備的時候，來到國中極肥沃之地，做他祖宗和祖宗的祖宗未曾做過的事，瓜分擄物、掠物和財寶，又策劃進攻堡壘；然而這都是暫時的。 \end{tabularx} \\ \\ \relax
11:25 & \begin{tabularx}{0.7\textwidth}{X} 「他必奮勇向前，率領大軍攻擊南方王；南方王以極強的大軍迎戰，卻抵擋不住，因為有人設計謀害南方王。 \end{tabularx} \\ \\ \relax
11:26 & \begin{tabularx}{0.7\textwidth}{X} 吃王餉的使王敗壞，王的軍隊必被沖沒，仆倒被殺的甚多。 \end{tabularx} \\ \\ \relax
11:27 & \begin{tabularx}{0.7\textwidth}{X} 至於這二王，他們心懷惡計，同席吃飯卻彼此說謊，但計謀不成，因為結局要在指定的時期來到。 \end{tabularx} \\ \\ \relax
11:28 & \begin{tabularx}{0.7\textwidth}{X} 北方王必帶許多財寶回本地，但他的心反對聖約；他恣意橫行，回到本地。 \end{tabularx} \\ \\ \relax
11:29 & \begin{tabularx}{0.7\textwidth}{X} 「到了指定的時期，他必返回，侵入南方。這一次卻不像前一次， \end{tabularx} \\ \\ \relax
11:30 & \begin{tabularx}{0.7\textwidth}{X} 因為基提的戰船要來攻擊他，他就喪膽而退。他惱恨聖約，恣意橫行，要回來善待那些背棄聖約的人。 \end{tabularx} \\ \\ \relax
11:31 & \begin{tabularx}{0.7\textwidth}{X} 他要興兵，這兵必褻瀆聖所，就是堡壘，除掉經常獻的祭，設立那施行毀滅的可憎之物。 \end{tabularx} \\ \\ \relax
11:32 & \begin{tabularx}{0.7\textwidth}{X} 他必用巧言奉承違背聖約的惡人；惟獨認識神的子民必剛強行事。 \end{tabularx} \\ \\ \relax
11:33 & \begin{tabularx}{0.7\textwidth}{X} 民間的智慧人必訓誨許多人，然而在一段日子裡，他們必因刀劍、火燒、擄掠、搶奪而仆倒。 \end{tabularx} \\ \\ \relax
11:34 & \begin{tabularx}{0.7\textwidth}{X} 他們仆倒的時候，會得到少許援助，卻有許多人用詭詐加入他們。 \end{tabularx} \\ \\ \relax
11:35 & \begin{tabularx}{0.7\textwidth}{X} 智慧人中有些人仆倒，為要使他們受熬煉，成為潔淨、潔白，直到末了；因為還有一段日子才到所定的時期。 \end{tabularx} \\ \\ \relax
11:36 & \begin{tabularx}{0.7\textwidth}{X} 「王必任意而行，自高自大，超過所有的神明，又用荒謬的話攻擊萬神之神。他必行事亨通，直到主的憤怒結束，因為所定的事必然實現。 \end{tabularx} \\ \\ \relax
11:37 & \begin{tabularx}{0.7\textwidth}{X} 他不顧他祖宗的神明，也不顧婦女所仰慕的神明，任何神明他都不顧；因為他自大，高過一切， \end{tabularx} \\ \\ \relax
11:38 & \begin{tabularx}{0.7\textwidth}{X} 以敬奉堡壘的神明取而代之，用金、銀、寶石和珍寶敬奉他祖宗所不認識的神明。 \end{tabularx} \\ \\ \relax
11:39 & \begin{tabularx}{0.7\textwidth}{X} 他靠外邦神明的幫助，攻破最堅固的堡壘。凡承認他的，他要給他們許多尊榮，使他們管轄許多人，又分封土地作為報償。 \end{tabularx} \\ \\ \relax
11:40 & \begin{tabularx}{0.7\textwidth}{X} 「到末了，南方王要與北方王交戰。北方王要用戰車、騎兵和許多戰船，勢如暴風來攻擊他，又要侵入列國，如洪水氾濫。 \end{tabularx} \\ \\ \relax
11:41 & \begin{tabularx}{0.7\textwidth}{X} 他要侵入那佳美之地，許多國就被傾覆，但以東人、摩押人和大半的亞捫人必逃離他的手。 \end{tabularx} \\ \\ \relax
11:42 & \begin{tabularx}{0.7\textwidth}{X} 他要伸手攻擊列國，連埃及地也不得逃脫。 \end{tabularx} \\ \\ \relax
11:43 & \begin{tabularx}{0.7\textwidth}{X} 他要掌管埃及的金銀財寶和各樣珍寶，路比人和古實人都跟從他的腳步。 \end{tabularx} \\ \\ \relax
11:44 & \begin{tabularx}{0.7\textwidth}{X} 但從東方和北方必有消息傳來擾亂他，他就大發烈怒出去，要將許多人殺滅淨盡。 \end{tabularx} \\ \\ \relax
11:45 & \begin{tabularx}{0.7\textwidth}{X} 他要在海和榮美的聖山之間搭起王宮的帳幕；然而他的結局到了，無人能幫助他。」 \end{tabularx} \\ \\ \relax
12:1 & \begin{tabularx}{0.7\textwidth}{X} 「那時，保佑你百姓的天使長米迦勒必站起來，並且有大艱難，自從有國以來直到此時，未曾有過這樣的事。那時，你的百姓凡記錄在冊上的，必得拯救。 \end{tabularx} \\ \\ \relax
12:2 & \begin{tabularx}{0.7\textwidth}{X} 睡在地裡塵埃中的必有多人醒過來；其中有得永生的，有受羞辱永遠被憎惡的。 \end{tabularx} \\ \\ \relax
12:3 & \begin{tabularx}{0.7\textwidth}{X} 智慧人要發光，如同天上的光；那領許多人歸於義的必發光如星，直到永永遠遠。 \end{tabularx} \\ \\ \relax
12:4 & \begin{tabularx}{0.7\textwidth}{X} 但以理啊，你要隱藏這話，封閉這書，直到末時。必有許多人往來奔跑，知識就必增長。」 \end{tabularx} \\ \\ \relax
12:5 & \begin{tabularx}{0.7\textwidth}{X} 我－但以理觀看，看哪，另有兩個人站立：一個在河這邊，一個在河那邊。 \end{tabularx} \\ \\ \relax
12:6 & \begin{tabularx}{0.7\textwidth}{X} 其中一個對那在河水之上、穿細麻衣的說：「這奇異的事要到幾時才應驗呢？」 \end{tabularx} \\ \\ \relax
12:7 & \begin{tabularx}{0.7\textwidth}{X} 我聽見那在河水之上、穿細麻衣的，向天舉起左右手，指著那活到永遠的起誓說：「要到一年、兩年，又半年，粉碎聖民力量結束的時候，這一切的事就要應驗。」 \end{tabularx} \\ \\ \relax
12:8 & \begin{tabularx}{0.7\textwidth}{X} 我聽了卻不明白，就說：「我主啊，這些事的結局是怎樣呢？」 \end{tabularx} \\ \\ \relax
12:9 & \begin{tabularx}{0.7\textwidth}{X} 他說：「但以理，去吧！因為這話已經隱藏封閉，直到末時。 \end{tabularx} \\ \\ \relax
12:10 & \begin{tabularx}{0.7\textwidth}{X} 必有許多人使自己潔淨、潔白，且受熬煉；但惡人仍必行惡，沒有一個惡人明白，惟獨智慧人能明白。 \end{tabularx} \\ \\ \relax
12:11 & \begin{tabularx}{0.7\textwidth}{X} 從除掉經常獻的祭，設立那施行毀滅的可憎之物的時候起，必有一千二百九十日。 \end{tabularx} \\ \\ \relax
12:12 & \begin{tabularx}{0.7\textwidth}{X} 那等候，直到一千三百三十五日的有福了。 \end{tabularx} \\ \\ \relax
12:13 & \begin{tabularx}{0.7\textwidth}{X} 「至於你，你要去等候結局。你必安息，到了末期，你必起來，享受你的福分。」 \end{tabularx} \\ \\
[1ex]
\hline
\hline
\end{longtable}
但以理書中的十一及十二章乃是這卷書之重要預言.十一章歸納有關外邦人之預言.十二章乃是歸納以色列人──神的選民──之預言.也可說是有關教會,及有關你和我的話.
\newline
\newline
彼得後書說:聖經所有之預言,不能按人之私意解說；因為這不是出於人,乃是出於神.神的先知,僕人,神合用的器皿,藉神感動說出神的話；是從一個或許多個心靈而出,再進入我們心裡.要我們留意神的預言,反覆思想及應用在我們生活上,發出生命的果效及見證.正如燈照在黑暗的地方,成為見證.
\newline
\newline
在此我們看見神之預言及實際歷史之事實,預言與歷史在這卷書中都可見到.一部份對但以理來說是歷史,就是經過歷史之過程而更感到神之奇妙.因此,信神之預言,每句都不落空.雖然他未必能澈底的明白,但是神的話句句帶著能力,使你心得著安慰與喜樂.有一部份是預言,因在當時還未應驗；但對你我而言,則已應驗了.那我們可以更深切地了解,甚至比但以理更明白,這是何等的福份!
\newline
\newline
在十一章中可知那是波斯時代,是從大利烏王開始,波斯王朝算起.是但以理實際的經驗,是對歷史之發展觀察而得.在十一2中說:現在我將真實指示你,必有三王興起,還有第四王出來.此乃歷史考據,這第四王就是以斯帖記中的亞哈隨魯王,當時國勢甚為強盛.主前四百八十年曾發動大軍征服世界,特別在希腊強權興起時,他曾用二百六十四萬兵出征.正如在此所說,激動的攻希腊,結果失敗,因而國勢漸弱.約在主前333年亞歷山大起來將波斯王朝打敗,可見這是何等正確的歷史.在亞哈隨魯時代,那時但以理已不在了,不能看見預言的應驗.我們卻當清楚地明白,因從實際的歷史發展所以知道.神確實是在歷史之中,也是超乎歷史之上,更是藉歷史將祂永恆的計劃表明；因神在歷史之中有一願望和目的,就是願人相信祂得祂的救贖,讓神的國擴大.可是神在歷史中每一件事,都有祂的許可及工作.
\newline
\newline
以前我們曾提過,亞力山大大帝有四個將軍,後來他們彼此衝突,當時還有外邦的勢力──南方有埃及,北方有敘利亞.在(但十一1-9)節中,佔時約一百三十年.在20-32節中,則單提及希腊最利害兇暴的人,那就是依比芬將軍(Antiochus IV (Epiphanes),得後成為王.在20節說:「那時必有一人興起接續為王.」又21節「必有一個卑鄙的人興起接續為王.」這是約在主前一百六十多年,猶大國有一次叛變；就是瑪喀比──猶大政治獨立運動者──之叛變,將希腊強權壓下,但後來卻失敗了.
\newline
\newline
在十一章中常提及「定期」這一字.「到了定期事就了結.」(十一27末),「到了定期他必返回.」(十一29).
\newline
\newline
神有祂自己的時間,因神在歷史中；人以為可用武力或強權左右歷史,或創造歷史.其實人一切都是不可能的,因沒有一事不在神的手中.所以雖然敵神者不斷出現,在新約中提起敵基督,有人說是依比芬,因是逼害神子民及神的聖所.但往後仍是有敵基督的,就是但以理七章的小角及啟示錄之最後一獸.在(但十一36)說:「自高自大超過所有的神,又用奇異的話攻擊萬神之神.」這是何等可怕,又何等的可憐!惡者又怎能抵擋萬神之神呢!雖然他曾一度好像是凡事亨通,不可一世的樣子.可是,神的忿怒必定要臨到,正如(但十一36)說:「直到主之忿怒完畢,因為所定的事必然成就.」如參考啟示錄就可對但以理書更清楚了解.「忿怒」這字,在舊約最少提過二十二次,這字特別提到神之刑罰,及公義之審判,因神是公義之神祂公義之能力必然彰顯,我們屬祂的人,應該知道祂的公義,才能更了解祂的慈愛,我們更該曉得自己怎樣蒙恩,免去神公義的忿怒.如今又成為何等人,都是主的恩典.
\newline
\newline
在十一章特別指外邦人之歷史,故在第二章中說及的人像,就是從巴比倫開始的外邦人的日子,當時巴比倫已經過去,所以在十一章中,再不提及巴比倫.只說波斯的強權.
\newline
\newline
人像腳下有十指頭,但以理七章中有十個角,這十腳及十角到底是指什麼?至今有各種不同的推測,有些已在歷史應驗,而且現今正在這階段中,雖然我們現在不容易肯定地說清楚,但卻知道外邦人之日期快滿,因神使我們看見第十二章中更確定的預言.
\newline
\newline
在但以理十二章的開始,特別提及天使長米迦勒,參考啟示錄可知有劇烈爭戰,也是十章所說的劇戰.而天使長就站在神選民的立場,因神之選民是神所寵愛的,祂看為極其寶貴,故必保守他們.今天神也必保守你和我,因我們是屬祂的.各位!你是否看重自己的身份?最高的天使長要保護我們,因我們是大蒙眷愛的.祂如此愛護我們,我們又如何愛祂?主說:「愛祂的,祂必愛他.」「懇切尋求祂的必尋見.」感謝主祂賜我們如此大的恩典.天使長站起來充滿得勝能力,讓我們今天看見他的時候,就充滿歡悅和興奮.因靠著愛我們的主就能得勝,而且得勝有餘.
\newline
\newline
我們怎樣才可得勝?在(但十二1-4)可說是預言中最寶貴的話語,特別指末世時那全備的救恩.在第二節中說:「睡在塵埃中的必有多人復醒.」這乃是復活的道理,有人說:有新約復活的道理才顯得清楚,那是不錯的,因沒有主復活作初熟的果子.我們又怎知道復活呢?其實舊約中亦有經文說及復活,且說到以色列人復興,及預言教會復興,更說及蒙恩得救的人得復活.這寶貴而基本的道理,乃是要成為你我的經驗.
\newline
\newline
在賽廿六章中,我們也見到以賽亞受聖靈感動及充滿說:「睡在塵埃中的要興起歌唱」「塵埃」意是人本為塵土出來,亦要歸回塵土,人不過是塵土.所以舊約每說及「塵埃」時,乃是說死亡的朽壞墳墓.故以賽亞說在塵埃中的要起來,那不單是起來,更要充滿生命活力,而且還要歌唱!高聲地讚美神,這也是以色列民的復興.當以色列人敗亡時,外邦人以為信神的人已敗亡.他們說:「有何神可以將你們從危險痛苦中拯救出來呢!」正如尼布甲尼撒王對但以理三友所說狂傲的話.又在詩篇中說:巴比倫人叫他們歌唱,他們能唱出來嗎?他們已把豎琴掛在柳樹上.他們說:我們怎能在被擄之地歌唱呢.但神的百姓要歸回,要復興.故以賽亞充滿喜悅地說這預言.
\newline
\newline
在以西結三十七章中有一重要的預言,就是神的靈叫以西結看見,在平原上有許多屍體,並已成枯乾的骨頭.這些骨頭還能活嗎?在人看是不能,但當時大風吹來──這是指聖靈的能力,當聖靈的能力來時,不單是枯骨要連結起來,而且有肉.不單是死裡復活,且成強大之軍隊,雄赳赳,氣昂昂地,神的百姓如此地走在得勝者的行列.
\newline
\newline
各位親愛的弟兄姊妹!我們屬神的人,在生活上只不過受到一點兒的挫折,就垂頭喪氣嗎?那麼我們有什麼見證呢!讓我們抬起頭來,眾城門啊!把頭抬起來,把我們心裡的門窗抬起來；家裡的門戶抬起頭來,教會的門戶啊,要抬起頭來!因榮耀的王要進來,讓祂進入我們的心裡,家庭裡,教會裡!因主是得勝的,而且還要得勝!
\newline
\newline
結三十七及賽二十六章所載,不但是神選民復興,亦是教會復興.當神的靈在教會工作時,每一人都會起來成為強大的軍隊,在舊約中我們可看見個人復活的真理:
\newline
\newline
在詩十六10末中,神的聖者,他們的靈魂不會朽壞.
\newline
\newline
何十三14末說:「死亡阿!你的災害在那裡?陰間啊,你的毀滅在那裡?」林前十五55的記載:「死阿!你得勝的權勢在那裡?死阿,你的毒勾在那裡?」可能受這預言的影響.
\newline
\newline
復活的真理在但十二2就更清楚:「睡在塵埃中的必有多人復醒.」為何說多人?該譯作所有的人才對,原來這「多人」是因在大災難之前義人的復活,這是第一次復活.在啟示錄也清楚告訴我們.「其中有得永生的,有受羞辱永遠被憎惡的」(但十二2末),這裡告訴我們有些人得永生有些人被憎惡,我同意猶太「拉比」的意譯的加註:「必有多人復醒,那些義者得永生,其餘的人因沒得神的恩惠；因他們不信神,叛逆神,結果受羞辱,永遠被憎惡.」這些人當時沒有復活,祗有義者先復活,這也符合啟示錄所說的.在3節中所說的智慧人是指義人,他們不單是義人,更是使多人歸義的.
\newline
\newline
在但十二1中,提及大艱難,這是重要的真理,這大艱難在舊約有預言──申四30曾說及艱難的時候.耶三十7也清楚說明:「雅各遭難之日子」.在新約也有主耶穌的話語:在太廿四14-15中甚至引用但以理的話,可見這是神的靈所感動,大災難必要來,這是空前的大災.
\newline
\newline
在啟示錄與但以理兩卷書都清楚說明,當大災難來時神的兒女們,屬神的人,如果已經死了的要先復活,──睡在塵埃中的必有多人復醒.如若活著,神要將我們提到空中,那也就是說,信徒在大災難來臨之前被提.
\newline
\newline
「凡名錄在冊上的必得拯救」,(但十二1)也是啟示錄的「生命冊」.在舊約出埃及記三十二章亦已提及「生命冊」,那是當以色列人犯罪拜偶像,摩西向神禱告；求神拯救他們,寧可自己的名字在冊子上被塗抹.但神說:我不塗抹你的名,我要塗抹那不肯信靠神的人,生命冊是由神寫上,也是由神塗抹,我們不能作什麼.在詩篇中亦可看見同樣的話語:在生命冊的是神所紀念及眷顧的,因神所要拯救的是屬於神,是大蒙眷愛,這是全備的救恩,(羅八章)甚至我們身體也得贖.(林後一10):「祂曾救我們脫離極大的死亡,祂現在仍要救我們,並且指望祂將來還要救我們.」祂的救恩是完備的,讓我們戰戰兢兢地做成得救的工夫,這難道我們未得救嗎?不是的,乃是要保守我們救恩的能力.在3節中清楚的說:「智慧人必發光.」誰是智慧人?是信靠神敬畏神的人.在舊約智慧文學書中(箴言,約伯記,傳道書.)講及敬畏神是智慧的開端,故敬畏神的人有智慧；有明亮的眼睛,眼睛瞭亮全身就光明.我們有屬靈的見解及眼光,有清楚的目標及盼望,有聖潔的生活(約壹三3)中說:「凡向祂有這指望的,就潔淨自己,像祂潔淨一樣.」這樣的人必能發光,不是做光而是發光.很容易就將神見證流露出來.正如天上的光一樣,如穹蒼裡的光團與大氣這神氣對光有調節及發散作用.將光照耀到每個幽暗角落,這是你我的地位和功用,也是神把我們存留在世的目的.「那使多人歸義的」,這意思不單我們自己得著義,還要使別人得義.
\newline
\newline
神曾應許亞伯拉罕說:你的後裔要像天上的星海邊的砂一般,這「星」與「砂」是比作非常多的意思,可是卻有極大的分別.海邊的砂常被潮水沖擊,甚至被人用腳踐踏,如有人要建房屋在砂土上是不行的.因根基不穩固,但天上的星就不同；它在天上發出光芒,在最幽暗時它就發光.我們這些義人能領人歸主的話,就是如天上的星一樣照耀著；忍耐地照著,直到晨星出現之時.在林後及啟示錄有說:主是晨星宣告早晨的來臨,因這光發出來,直到黎明的時候.
\newline
\newline
「但以理阿,你要隱藏這樣,並封閉這書直到末時」(但十二4),為何要封閉?是不讓別人了解嗎?不是的,乃是要別人知道對主是要有信心.假如沒有信心,那麼,就算得著這些話語也不能發出功效.
\newline
\newline
「必有多人來往奔跑,知識就必增長」(但十二4).這經文有許多不同的解釋,有人說這是指末世的情況,「多人來往奔跑」是指交通非常發達,知識也甚增長.二十世紀是科學時代,交通及教育都發達.「來往奔跑」有小字寫著「切心研究之意」世人跑來跑去,那是何等的徬徨,他們不知何去何從,更不知要尋找甚麼,甚至當要尋找神的知識時也沒有機會.因為時候到了神的恩門已關,世人都要面對面的見神.「知識」這字,許多人如此翻釋(敘利亞及拉丁文譯本都作「知識」).但在七十士之譯本中,卻翻作「罪惡」,我覺得更為貼切.即就是說:許多人來來往往,不知作何事；整日奔波勞碌,結果,罪惡就更增長,這是末世的景況.
\newline
\newline
但以理又見有兩人,一個站在河這邊,一個站那一邊.其中一個問另一個人,(5節)然後這穿細麻布的(正如第十章的)再重複地解釋一載,二載,三載,三年的時間.(啟示錄亦有說及,並有說四十二個月).這些事一定要成就,但還有兩個數字在但以理書出現的就是:(十二11)中一二九○日；假如以三十日為一月,三年半應該是一二六○日,為何多出三十日?這乃是恩典的日子.我們當記得,挪亞造方舟時,曾勸人要信靠神,他傳義道有一二○年之久.但無人相信,因為當時的人們覺得挪亞是瘋狂的,不明白他所說的.結果,方舟造好,挪亞全家入了方舟,那時並未立刻下雨,而且再等了七天.還有七天的機會,神給人還有一點點的機會.在12節中說:等至一三三五日,那是再加增四十五日.神用很多的時間準備,卻用很短時間成全,祂用很長的時間一直叫人等著主來.正如主來世用三十年之時間在等候,但祂傳道只不過三年半.在那麼短的傳道時間中,最重要的,卻是最末的一個禮拜.我們常引用太廿四章,路廿一章,可十三章作為最重要的信息──主一定要再來.
\newline
\newline
在第8節中但以理說:「我不明白這事,主啊,這些事的結局是怎樣呢?」第9節主:「你只管去」,13節:「你且去等候結局」,12節:「等到一千三百三十五日的,那人便有福了.」
\newline
\newline
你只管走人生的道路,伸出信心的手來,讓主來牽引.我們要有忍耐,有信心,專心一意地等待,用禱告的心去等待.因等候神的必不羞愧,仰望神的必有光榮.
\newline
\newline
保羅曾說:當你在歡喜快樂的環境,一切都順利時,就該歡喜快樂高聲地讚美主,但如遇患難就要忍耐,且要等候.(羅十二12)有盼望就有喜樂,有患難就該忍耐,無論任何光景,總要好好地等候那結局；那就必得安息,你應走的路會走完,且該安息.但神的歷史並沒有結束；神的計劃,神的工作都沒結束；祂在末世之時要來,那時你還要起來享受你的福份.今日主亦對我們說:我們在世靠神的恩典走在錫安的大道上,力上加力,恩上加恩,福上加福.有一天走完人生的路程,要離世安息,與主同在真是好得無比.但如果主要我有存留生命等祂再來,也是主的恩典,無論如何總要好好地仰望主.但以理是大蒙眷愛的人,他已離世,他不會回來；可是等到有一天,我們上去可以看見他.只有一位祂去了還要再來,那是我們的救主基督.祂說:祂如何上去也如何來,祂一定再來,啟廿二7「看哪!我必再來.」20節說:「是的,我必快來.」我們該儆醒地等祂來,引多人歸向祂,見證祂.當晨星出現的時候,必有萬道光芒,宣告榮耀永恆的早晨,讓我們同聲說:「主啊!我願你來!」
\newline
\newline


\newpage

\section{第45屆~港九培靈會~胡恩德~第一講}
\label{sec:1225}
\textbf{主再來的重要性}
\newline
\newline
連結: \href{https://www.hkbibleconference.org/session-message/view/1225}{\texttt{https://www.hkbibleconference.org/session-message/view/1225}}
\newline
\newline
\hyperref[sec:1259]{< < < PREV SERMON < < <}
~
\hyperref[sec:toc]{[返目錄]}
~
\hyperref[sec:1227]{> > > NEXT SERMON > > >}
\newline
\newline
使徒行傳 1:9-1:11
\newline
\begin{longtable}{cl}
\hline
\hline
章節 & 經文 (和合本修訂版)\\
\hline
1:9 & \begin{tabularx}{0.7\textwidth}{X} 說了這些話，他們正看的時候，他被接上升，有一朵雲彩從他們眼前把他接去。 \end{tabularx} \\ \\ \relax
1:10 & \begin{tabularx}{0.7\textwidth}{X} 他升上去的時候，他們定睛望天，看哪，有兩個人身穿白衣站在他們旁邊， \end{tabularx} \\ \\ \relax
1:11 & \begin{tabularx}{0.7\textwidth}{X} 說：「加利利人哪，你們為甚麼站著望天呢？這離開你們被接升天的耶穌，你們見他怎樣升上天去，他也要怎樣來臨。」 \end{tabularx} \\ \\
[1ex]
\hline
\hline
\end{longtable}
主要再來,祂再來　的預言,乃是新約聖經主題之一.如果查攷聖經將有關主再來之經文刪去,則所餘下的經文就不多了,可知主再來這信息是何等的重要.
\newline
\newline
但今日的基督徒,對這問題反而不大重視；與初期教會之信徒相比,實相差甚遠!今日的基督徒覺得「死」比主再來更為重要,覺得「死」是一定要面臨的大問題,所以就多方的打算；準備──葬何處,死後穿何衣,火葬或是土葬.......其實人死了還能作什麼呢!何況死又不一定會臨到我們.在(林前十五51)說:「我們不是都要睡覺,乃是都要改變.」若主再來之時,我們仍活著,我們就毋須經過死亡了!
\newline
\newline
主再來必要成為事實,如果基督徒能夠實實在在的在主裡面,就不該日夕怕死等死了.我們很可能會在死時,等候主把我們的靈魂接去；但亦可能在祂再來之時,使我們身體改變而被接去.這是必定會臨到屬主的人身上的事.故此,我們對主再來就應該早作準備,更加重視了.
\newline
\newline
基督徒若太看重身後事,好像世俗人一樣,那就證明他們不是屬天的.他們不知道自己的身份,更不知道自己將來榮耀的福份.正如以掃因為看重屬地的便失去了屬靈的福份.尤其是在香港,多少人一味向世界追求!想得著世界.彼得說:「私慾是與我們與靈魂爭戰的.」(彼前二11)追求世界就屬乎私慾,以致我們寧可等死,卻不等候主再來.基督徒態度如此,實在對不起愛我們的主!要等候主再來,而且應該認為那是我們無上的福份；存這樣的盼望,才能得到主的喜悅!
\newline
\newline
我們對身後事,為甚麼會早作準備呢?是因為「死」必會來到的.「按著定命,人人都有一死.」(來九27)所以會對死看為重要.但對主再來,我們所知有限,亦不夠確切!其實主第一次來,我們已承認那是歷史上千真萬確的事實.(非信徒亦相信)那麼,主再來當然也是事實,且是主親口說的.(徒一11)載有兩位天使奉差遣對門徒說:「加利利人哪!你們為什麼站著望天呢,這離開你們被接升天的耶穌,你們見祂怎樣往天上去,祂還要怎樣來.」約翰在啟示錄也說主一定要駕雲降臨.(啟一7)主耶穌受審時,大祭司叫祂起誓答覆問題:「你是永生神的兒子不是?」主答:「你說是的,然而你們要看見人子坐在那權能者的右邊,駕著天上的雲降臨.」(太廿六63-64),這是主親口說的,且用祂的性命保證,支持這見證.因祂如此說,他才會被判死罪.人說謊是為免死,沒有人說謊而去叫人定自己罪的.祂何時再來當然我們不知道,何時何日,也祗有父知道.新約聖經是以主再來為主題,在廿七卷中有廿三卷是說主再來的,四福音書說及祂第一次來地上的工作和言行,以及祂再來的預言.如果一個基督徒不把主再來放在他的信仰裡,他的信仰就有問題了.
\newline
\newline
在約翰福音中雖較少提及主再來,但主仍親口說:「你們心裡不要憂愁,你們信神也當信我；在我父家裡有許多住處,若是沒有,我就早已告訴你們了.我去原是為你們預備地方,我若去為你們預備地方,就必再來接你們到我那裡去,我在那裡叫你們也在那裡.」(約十四1-3)又在約廿一章主答覆彼得問及約翰將來的問題時說:「我若要他等到我來的時候,與你何干?你跟從我罷!」這些經文都有提及主再來的預言的.
\newline
\newline
太廿四至廿五章,可十三章及八章末至九章初,路廿一章.十七和十二章內,這些經文都是在四福音中提及主再來的；聖經既然多提到這事,那麼教會及個人就更應當看重主再來了.今日的信徒對主再來的態度是很不正常的.有人怕主來,有人盼望祂遲來；甚至有人期望自己死了然後主才來!這些態度都是錯誤的.
\newline
\newline
初期教會的信徒十分盼望主再來,尤以使徒們,他們跟主三年多,主是他們的良師；他們的心已被主吸引,也有了感情,當然會盼望主再來.就算未見過主面的信徒,也有這盼望,正如(彼前一:8)說:「你們雖然沒見過祂,卻是愛祂.」那些曾見過主的人又怎會愛主呢?原來聖靈,會將主最高的愛,澆灌在我們心裡.我們有了愛主的心,自然就會盼望主再來了.故此人能盼望主來的最大動機乃出於愛；不單基督徒如此,一切受造之物也有這盼望,等候得贖的日子來到.(羅八22-23)雖然現在主已經與我們同在了,但我們總是感到美中不足,倘若等到主來了,我們都要改變成為榮耀的身體,整個靈,魂,體都與主合而為一.那是何等快樂的事阿!如果我們對主的愛有真的領悟,也嘗過主恩的滋味,那麼我們就不能不愛主了,我們的心必會傾向祂,願意跟從祂；也照祂所吩咐去行,並會盼望祂再來.
\newline
\newline
在啟示錄中提到羔羊的婚筵.這婚筵是指著主再來,與教會婚娶的日子.祂是新郎,教會是新婦,是一切蒙恩得救的人.這最完滿的結合,就是羔羊的婚筵.如果一個基督徒不愛慕主再來,也不切望主來；就等如一個女子定了婚,卻不願結婚一樣.這是否反常呢?主再來,就是教會與主婚娶之日,不切望這日子來臨的人,他應該要自責悔改!(羅五2)記載初期教會的信徒是何等切慕主來,可以作我們的榜樣,保羅說:「我們既因信稱義,就藉著我們的主耶穌基督,得與神相和；我們又藉著祂,因信得進入現在所站的這恩典中,並歡歡喜喜盼望神的榮耀.」求主改正我們之觀念,使我們常盼望主耶穌基督快來吧!
\newline
\newline


\newpage

\section{第45屆~港九培靈會~胡恩德~第二講}
\label{sec:1227}
\textbf{盼望主再來的態度}
\newline
\newline
連結: \href{https://www.hkbibleconference.org/session-message/view/1227}{\texttt{https://www.hkbibleconference.org/session-message/view/1227}}
\newline
\newline
\hyperref[sec:1225]{< < < PREV SERMON < < <}
~
\hyperref[sec:toc]{[返目錄]}
~
\hyperref[sec:1230]{> > > NEXT SERMON > > >}
\newline
\newline
主耶穌再來,乃是新約聖經中主題之一,主自己也說得很清楚,所以我們應看為重要,並該存這盼望.新約差不多每卷書的作者都提及主再來(只四卷是例外),可見主再來是何等重要.初期教會根據主的應許,一直盼望主再來,他們對於這盼望的態度,與今日的基督徒對此,實在差別很大.
\newline
\newline
今日的基督徒不大想到主來,只是想到自己會死；我們怕死,盡量把死推出去；怕主來,也盡量在思想上,把這意念推出去,這是非常的不正常的!是否初期教會的信徒,因為親近過主,與主有感情,而盼望主再來呢?不是的,其實他們也有不少人未見過主,他們卻有喜樂的心盼望主再來.聖經教訓我們要儆醒準備,並嚴重的警告我們；但也教訓我們,要用喜樂的態度等候主來.
\newline
\newline
(路廿一27-28)「那時,他們要看見人子,有能力,有大榮耀,駕雲降臨；一有這些事,你們就當挺身昂首,因為你們得贖的日子近了.」
\newline
\newline
在這段經文中,主警告世人說祂再來時的情境,祂吩咐門徒一有這些事,就該挺身昂首,應當歡喜快樂,因為得贖的日子近了.「死」雖然是可怕的事,但「死」對基督徒來說卻是一件得福的事,因我們要回到天家見主面.主再來更是我們有福的事,因我們的靈,魂,體,都要復活改變,永遠與主同在.主曾作比喻說:天國好比十個童女迎接新郎回來,這是歡喜而高興的事,我們也該歡迎祂再來.我們虔守主的晚餐,是因為主曾說:你們要如此行,為的是紀念我,直等至我來.主耶穌是教會的新郎,祂再來就是新郎要回來,這是教會歡喜快樂的日子.教會對新郎之情愛,是人間最強烈的愛,最喜樂的愛；這是教會喜樂之日,也是主最喜樂之日,因為祂喜歡看見這唯一的愛人──教會──.祂曾為教會受苦,這是祂愛教會的最高峰；祂再來時,就是愛教會強烈喜樂的最高峰,祂要歡迎教會回到祂懷抱.所以我們雖然沒有見過祂,卻是愛祂,如今雖不得看見,卻因信祂,就有說不出來滿有榮光的大喜樂!祂盼望我們因祂而喜樂,想及祂的名而心裡歡欣.有一詩歌:「耶穌你名何等甜蜜」,我們是否有此甜蜜呢?
\newline
\newline
(提多二13)說:「等候所盼望的福,並等候至大的神,和救主耶穌基督的榮耀顯現.」
\newline
\newline
讀這段經文,使我們看見:(一)我們等候主的福.(二)我們等候主自己.
\newline
\newline
我們今日盼望主再來,乃是盼望有福的日子來到,並非懼怕哀哭黑暗之日；所以要挺身昂首,要快樂等候至大的神,與主耶穌基督榮耀的顯現.保羅說:「現今我們是憑信心,並非憑眼見.」(林後五7)主榮耀的顯現,是能夠看見的,也是有福的.正如人間的榮耀,雖與自己無關,但我們亦會感到高興和榮耀.我們的主再來時是有榮耀的.猶大書說:「主帶著千萬的聖者降臨.」,有無比的榮耀,不單是光輝,更是榮耀和華美.正如地上的會幕,愈深入則愈華美；是用多種線織成的,配上基路伯；裡面還有燈台,有約櫃等,都是包金的.至於聖所和至聖所,一切都是用金裝飾的,所以聖所都充滿了華美!主再來時所顯現的華美,實非我們所能想像到的.試看天然界的日出日落,以及神所造的樹木,花草的華美,可知該是如何了.
\newline
\newline
(西三4):「基督是我們的生命,祂顯現的時候,你們也要與祂一同顯現在榮耀裡.」
\newline
\newline
這裡所說的榮耀,我們也有份；我們等候主得榮耀之日,也是眾聖徒與祂同得榮耀的日子,所以應該成為我們有福的盼望.
\newline
\newline
(提後四8):「從此以後,有公義的冠冕為我存留,就是按著公義審判的主,到了那日要賜給我的；不但賜給我,也賜給凡愛慕祂顯現的人.」
\newline
\newline
如果有人愛慕主顯現,主必會將公義之冠冕賜給他.這種人必然愛主生活必聖潔,也必是與主有密交.假如我們不喜歡主來,我們的生活必有問題,愛主的心必冷淡,看主是不值得愛的?求主赦免我們,也憐憫我們,好使我們有愛慕主的心,下功夫追求,使失去了的愛心,不願見主的心可復興過來,回頭愛主,而且愛慕主的顯現.
\newline
\newline
假如主來時,自己的人也不歡迎祂,試想祂心中會有何感受!(約一11)說:「祂到自己的地方來,自己的人倒不接待祂.」只有憎恨祂的人,妒忌祂的人才不歡迎祂,最後還會說:「除掉祂!釘祂十字架!」但祂第二次再來時我們就不能不歡迎祂了,因為我們是祂用血,用重價所買贖回來的.別人可以不歡迎,我們卻不能不歡迎祂.現在我們知主快要來了,卻仍冷冷淡淡地,不當為有福的盼望.試問,這種態度,又怎能不使主傷心呢?如再不改好這態度,那麼恐怕我們難免被撇下了.
\newline
\newline
有些人是怕主來審判,因為主的審判是嚴格的,是有賞罰的,我這樣的人看來是不能通過的了!主是那麼嚴厲公正!我又怎能逃得過祂的審判呢!瑪拉基書說:「祂來的日子,是如煉金之人的火,又如漂布的人的鹽.」(瑪三1)主來時任何的渣宰,剛愎,頑梗都不能存留,要完全打碎,我又怎能到主面前!我的思想如此污穢,我又怎能不怕?正如(太廿五41)中主說:「你們這被咒詛的人離開我,進入那為魔鬼和他的使者所預備的永火裡去.」那時我當如何!
\newline
\newline
但假如我們是個重生得救,已有生命的人,就不必存這懼怕的心了,試看(帖前五5-10):「你們都是光明之子,都是白晝之子；我們不是屬黑夜,也不是屬幽暗的,所以我們不要睡覺,像別人一樣,總要儆醒謹守.因為睡了的人是在夜間睡,醉了的人是在夜間醉；但我們既然屬乎白晝,就應當謹守,把信和愛當作護心鏡遮胸,把得救的盼望當作頭盔戴上,因為神不是豫定我們受刑,乃是預定我們藉著我們主耶穌基督得救.」祂替我們死,叫我們無論醒著睡著,都與祂同活.在世人來說,他們是絕不能逃脫,但保羅說我們是白晝光明之子；那屬黑夜的人對主來才會懼怕,所以我們應儆醒謹守做人(不是睡覺的人).因屬乎白晝,故當穿上白晝的服裝,把信和受當作護心鏡遮胸,把得救的盼望當作頭盔戴上.得救的盼望乃是我們的頭盔,主來時我們就完全得救(連身體也得救了).
\newline
\newline
所以,只管盼望,不要懼怕,神已預定我們得救.藉著我們的主得救,不是藉我們的善行得救.只管盼望,正如頭盔保護我們的要害處,這盼望也如戰場.當我們做人覺得灰心,困難而想放棄一切時；只要想到主快來,就會堅強起來而繼續下去.我們是得救的人,是永屬主的,我們不是預定受刑的人,我們是藉主的寶血買贖的.祂來時一定接我去,就算是我的生活不完全,也不至於沉淪.根據我們不沉淪這一點,我們就應該喜樂；也因這一點,而知主是愛我的.雖然祂會嚴厲的待我,但我已嘗到祂的愛,祂的恩情,我是祂所愛的,是祂所救的；所以歡迎祂,並愛祂顯現,求主幫助我們.
\newline
\newline


\newpage

\section{第45屆~港九培靈會~胡恩德~第三講}
\label{sec:1230}
\textbf{盼望主再來的福氣}
\newline
\newline
連結: \href{https://www.hkbibleconference.org/session-message/view/1230}{\texttt{https://www.hkbibleconference.org/session-message/view/1230}}
\newline
\newline
\hyperref[sec:1227]{< < < PREV SERMON < < <}
~
\hyperref[sec:toc]{[返目錄]}
~
\hyperref[sec:1232]{> > > NEXT SERMON > > >}
\newline
\newline
以弗所書 1:15-1:18
\newline
\begin{longtable}{cl}
\hline
\hline
章節 & 經文 (和合本修訂版)\\
\hline
1:15 & \begin{tabularx}{0.7\textwidth}{X} 因此，我既然聽見你們對主耶穌有信心，對眾聖徒有愛心， \end{tabularx} \\ \\ \relax
1:16 & \begin{tabularx}{0.7\textwidth}{X} 就不住地為你們感謝神，禱告的時候常常提到你們， \end{tabularx} \\ \\ \relax
1:17 & \begin{tabularx}{0.7\textwidth}{X} 求我們主耶穌基督的神，榮耀的父，把那賜人智慧和啟示的靈賜給你們，使你們真正認識他， \end{tabularx} \\ \\ \relax
1:18 & \begin{tabularx}{0.7\textwidth}{X} 照亮你們心中的眼睛，使你們知道他呼召你們來得的指望是甚麼，他在聖徒中所得榮耀的基業是何等豐盛， \end{tabularx} \\ \\
[1ex]
\hline
\hline
\end{longtable}
(弗一15-18)許多基督徒,對於主再來是漠不關心的,他們都存著不理會的態度.也有些是怕主再來的,其實我們應該以主再來為我們的福氣才對.在提多書二章裡說:我們等候所盼望的福,我們等候此福氣,並至大的神主耶穌基督的榮耀顯現.
\newline
\newline
在(弗一15-18)經文中,知以弗所教會相當的好,他們不但順從主,且親愛眾信徒；但保羅仍要他們認識一點,就是照明他們心中的眼睛.使他們知道主的恩召有何等指望(18節).求主開恩照他們,使他們知道前面有甚麼?這是當時保羅為他們所請求的.保羅那時坐在牢中,他看見主之榮耀與盼望,他自己覺得,也充滿了這種榮耀.而將來基督徒到主面前時,我們也同樣得著此榮耀,雖然如今我們所見的是有限,在約一書三章中說:「將來如何還未顯明,但我們知道,主若顯現,我們就必像祂,因為必得見真體.」可見主給我們的福氣何等大!這一點能使我們愛主更多,也使我們事奉祂及傳福就更努力去做.如果我們對主所賜的恩典看得清楚,並自己有所享受的有感動；那麼,這恩典就成為一種力量,使我們肯服事主,並去傳揚福音.神所重用的僕人慕迪先生,他翻閱聖經,要查看主的恩典,愈看越覺得主的恩典非常之大.在他心中顯明,甚至不能自制,他就跑到街上,碰見頭一個人就述說出來,那人當時就得救了!求主開我們的眼睛,好使我們也看見前面的盼望,成為力量,肯去服事主,肯去傳福音.我們更該求主開我們的眼睛,賜下啟示和智慧的靈.今天我們信了主,假如是實在地親近主；追求主的話,那麼你可發現在主裡的榮耀,福份,享受和喜樂；並非你在初信時所知道的,更非你未信時所能了解的.我們看到今生有福氣擺在前面有人對讀書有味道,有人對縫紉有味道,或對運動有味道,肯付代價竭力去追求；因為他覺得是享受,今生的尚且如此何況將來永遠的福份,又怎會輕易放過呢!保羅說你們要存這樣的心,就是追求渴慕的心,有時我們看自己似已有長進；但在主的恩典方面來看,則仍相差甚遠!連保羅在腓立比書中他對自己仍說:「要忘記背後,努力面前,向著標竿直跑.」這是追求的意思,我們大家都要往前直追.智慧的靈啟示我們,開我們的眼睛,好看見那擺在前面,就是主來時給我們的福份吧!本人在廿餘年來有機會領過不少夏令會,許多團契早已擬定了講題,這可以說有好也有不好.其實傳道人所領受的恐怕會比你所擬定的還好,但他給你限制了.他們所擬的題目也不外是愛主更深,主必快來,向下札根向上結果那一套,因為他們所領受的到底有限.講員能知你們所不知的,懂得你們所不懂的,較不如讓他們自定題目更好!現今許多人以為自己知得的太多,其實他還是不知.聖經你們能背多少呢?現今的事我們所知尚且有限.這何況是將來的事呢!有一個神所重用的僕人,一次讀以弗所書五章讀到基督愛教會,為教會捨己,突然有一個亮光,叫他想到主愛我們,即使全世界的人都得救,只有我一人沉淪,祂也願意來救我；如同牧人撇下九十九隻羊,去尋找那失落的羊一樣.他感到主的恩實在太大了,差不多要跳上屋頂,這是真真啟示的靈能使我們看見的.現今要追求的,尚且這樣多,何況將來的呢?所以保羅禱告說:求主開我們的眼,好使我們看見主再來帶給我們的福份!
\newline
\newline
(彼前一13)說:「所以要約束你們的心,謹慎自守,專心盼望耶穌基督顯現的時候,所帶來給你們的恩.」我們對將來的恩寶貝嗎?希罕嗎?這不是說我們死後才與主同在,那不是最完滿的,乃是指主來時使我們得榮耀的身體,百分之百的與主同在,並享主福的恩典.普通基督徒沒有想及此,若要想,只是想死後得福；那當然也是好的,正如保羅說:離開身體與主同在,就好得無比.我們教會有一弟兄,未死前已失聲,喉已破他對妻子打手勢；又以手向上指,及作彈琴狀忽然間他破喉而唱三一頌!這真莫名其妙!因為我們以往不習慣唱三一頌的,我也不明白他怎會懂得唱!他大聲唱,他已看見天上的大喜樂!可惜我們只想到現在的,卻忽略了將來的.但我們更應想主來,不要追求眼前的,地上的事,好像世人一樣.那些東西既然充滿了我們的心我們的心,就無法有餘地裝載將來的榮耀了.故彼得說:「你們要約束你們的心,謹慎自守,專心盼望；要專心,如讀書,駕駛,無論作何事我們尚且要專心,何況對主再來的事呢?所以我們要專心盼望；主再來時帶給我們的恩典.能專心者,才能認識祂的恩典.
\newline
\newline
(帖前四13-18)這段經文常在安息禮拜時誦讀,我們也會聽到.這就是主再來時帶給我們有復活的恩典.我們死,但還會復活,主叫我們有一個復活的身體；這是福氣,還不是第一的福氣,有賞賜,也不是第一的福氣,其實在17節末我們會和主永遠同在,才是第一最大的福氣.我們現在已經有主同在,但將來可以完全得著主.在外國許多青年人覺得生命空虛,但信主後就立刻覺得有主同在,充滿喜樂.現今我們有主之親近,有主之同在,我們就會有主之喜樂.有些教會的領袖喜歡為青年人尋找娛樂,我並非反對,如果正當的娛樂,本是好的；但如果像世人追求奢華,宴樂這就是聖經所反對的.愛宴樂不愛神,尤其羅馬人的荒宴這是聖經所定罪的.我們有了主的快樂,就不會愛世俗的娛樂了.倘若有人說,做基督徒就沒有快樂了,說這話的人,顯然是失體,也失禮了他為教會領袖的身份.他對主同在的快樂,知得太少了.主也曾對門徒說:「凡我所吩咐你們的,都教訓他們遵守,我就常與你們同在,直到世界的末了.」(太廿八20)這是指我們今生與主同在的快樂,但沒有說好得無比!因為好得無比的快樂,是將來在天上的.既然生前有主同在,我們也覺得是享受；何況將來在祂榮耀降臨時更認真的顯出,我們也認真認識主,享受這福份了.就如我們所愛的人,常在一處,這真是享受到同在一起的快樂了!保羅題到主再來的福氣中,看來這福氣應該列在第一.有一次本人在新界石崗上村,在車站附近的小舖,見一個看門的,原來裡面有一個小孩子.父母出去了,把小孩關在裡面,如果我有辦法開了門,又買雪糕,香蕉給他吃,叫他不要哭,可能他也不要,是否要母親帶雪糕回來呢?不,他也不是這樣想；是否要母親買玩具回來呢?可能是,但實在他所要的,乃是母親自己.有一個做母親的說,小孩子見了母親,就好像見了糖一樣；反過來說,母親看見他的兒子,也好像一塊糖一樣.
\newline
\newline
剛才我們所唱的詩歌是主自己,培靈會歷屆都會唱.到底我們懂不懂詩詞的意義呢?這詩的涵義是相當深的,作者宣信博士有此經歷,方能寫出.我們所追求的,是否主本身,不是追求能力,去佈道,阿!人的靈魂重要呀!主寶貝還是人的靈魂寶貝呢?如此說,是否要起革命?傳福音救人靈魂不重要嗎?不是的,要知道這一切的出發點,無論追求甚麼屬靈的事,讀經,禱告,傳道,一切應該是為了主自己,我為愛主而做,要主自己,以祂為榮!
\newline
\newline
有時我們跌倒,覺得一片烏黑,好像不見了主；正如雅歌書的新婦不見了良人一樣.從前有一位弟兄,困不高興人指責他,他們不參加聚會,禮拜天開汽車到沙灘,自己讀聖經,祈禱.他以為自己沒有錯,其實他已經錯了,因為他自大,不接受人指責.其後他返回聚會,我還記得他垂頭喪氣,不好意思,但他肯謙卑回來已算好了.他見證說:我雖在沙灘讀經,祈禱,但像死一般沒有平安!因為失去了主的同在,不少人可能也有這種經歷.故我們傳福音是為主,要追求是為主.要主有所得,那麼我們就覺得與主同在是好得無比了.將來那諸般福氣中最大的福氣是與主同在,主與我們一點間隔都沒有；因明明看見主的同在,這是何等大的福氣.許多青年到外國讀書,都會思家,而且歸心似箭,歡喜甚麼,歡喜是見到所愛的人同在；如果我們不曉得與主同在的福份,實在太可悲了,求主憐憫我們!
\newline
\newline


\newpage

\section{第45屆~港九培靈會~胡恩德~第四講}
\label{sec:1232}
\textbf{得贖的日子}
\newline
\newline
連結: \href{https://www.hkbibleconference.org/session-message/view/1232}{\texttt{https://www.hkbibleconference.org/session-message/view/1232}}
\newline
\newline
\hyperref[sec:1230]{< < < PREV SERMON < < <}
~
\hyperref[sec:toc]{[返目錄]}
~
\hyperref[sec:1234]{> > > NEXT SERMON > > >}
\newline
\newline
以弗所書 1:13-1:14
\newline
\begin{longtable}{cl}
\hline
\hline
章節 & 經文 (和合本修訂版)\\
\hline
1:13 & \begin{tabularx}{0.7\textwidth}{X} 在基督裡你們聽見真理的道，就是那使你們得救的福音，你們也信了他，就受了所應許的聖靈為印記。 \end{tabularx} \\ \\ \relax
1:14 & \begin{tabularx}{0.7\textwidth}{X} 這聖靈是我們得基業的憑據，直等到神的子民得救贖，使他的榮耀得到稱讚。 \end{tabularx} \\ \\
[1ex]
\hline
\hline
\end{longtable}
(弗一13-14,四30)這兩處經文都提及「得贖」和「聖靈」,這是憑據(印記),保證將來基督再來時我們得贖.在14節中說聖靈是我們得基業的憑據.這「憑據」意思是「質」,也是說聖靈作「質」.聖靈住在教會裡面,也住在人裡面.在(弗四30),聖靈亦稱神的印,或形象等；正如印記印在心中,所以我們屬於神是鐵定的,不變的.等主來時,我們有印記的被贖出來.為何說贖出呢?
\newline
\newline
因為亞當聽從撒旦的說話就得罪了神,於是死亡,憂愁,痛苦,憂慮,虛空就臨到.我們在惡者手下無法出來,我們得贖乃是主將我們救出,又付代價將我們贖出.在這大黑暗權勢下能贖出,並非易事；在天空執政掌權者之下,我們無法出來.即使在日常生活要脫去牠的控制,也不容易,何況那當然應臨到的死亡之下能脫離,且又得贖呢!如果神只差天使來贖我們,那是會失敗的；唯有神自己與神本身一樣的兒子,才能作此事.因祂在萬有之上,祂與神本有一樣的權力,且又與神有交通.但當祂在十架時,就因我們的罪而與神之交通斷絕了.主放下一切而負上這麼大的代價,才可以把我們救贖出來.祂有寶座,有權柄,在天地之上,是至高之主,是生命的源頭；這一切,祂卻因為救贖我們而放下了.本來這樣的犧牲,祂是絕對不能作的,也不應作的,但祂為我們作了.祂把我們從惡者手中,死亡權力之下,失望,虛空,勞碌,痛苦之下贖了出來,主的救贖是偉大的,而且是有步驟的.
\newline
\newline
(一)使我們的靈先得贖──我們是罪的奴僕,又是魔鬼的奴僕.我們犯罪痛苦,如今主肯救我們；那麼,我們只可說,我要以主為主,願主救我.這樣主就會立刻進入我們心中,使我們的靈得救贖.
\newline
\newline
主說:「我實實在在告訴你,那聽我話,又信那差我來者的,就有永生；不至於定罪,是已出死入生了.」(約五24)這是說我們的靈已脫離死亡,主生命的靈進入我們的靈裡,使我們與祂聯合,就是與基督成為一；這是靈得救贖,也是我們地位得救贖.聖靈成為「質」,神在無限恩典中,體諒我們信心的微小,就將祂的靈放在我們裡面成為「質」,以至神等到祂的工作完成時,就可得整個完滿全面的救贖,這是何等大的恩情!
\newline
\newline
(二)身體的得贖──(羅八19-23)
\newline
\newline
「靈」雖已得贖,但我們的身體仍然會有疾病,會死亡.故此,我們這有聖靈初結果子的,也心裡歎息等候,得著兒子的名份,乃是我們的身體得贖.另外,保羅在(林後五4)也有這樣的話說:「我們在這帳棚裡,歎息勞苦,并非願意脫下這個,乃是願意穿上那個；將叫這必死的,被生命吞滅了!」
\newline
\newline
「基督若在你們心裡,身體就因罪而死,心靈卻因義而活.」(羅八10)我們信主得永生,是立刻就可以得著,但並非指身體得贖而言,乃是指心靈得著.因為身體是仍然會死的,因有罪性存在；但心靈卻可因義而活,能與神交通並事奉主.我們在世界上,與其他受造之物,同在敗壞的轄制和虛空之下,我們仍要等候得贖的日子來到(弗四30)即(羅八21)說我們可以得享神兒女自由的榮耀.又(羅八23末)說等候兒子的名分,乃是我們的身體得贖.這是在主再來之時,救贖就完成了,再不為身體難過,我們身體已成榮耀的身體,這是最有福的.
\newline
\newline
(腓三20-21)說:「我們都是天上的國民,並且等候救主,就是主耶穌基督,從天上降臨；祂要按著那能叫萬有歸服自己的大能,將我們這卑賤的身體改變形狀,和祂自己榮耀的身體相似.」人的身體在每一個階段年齡都有威脅,如:幼年,青年,老年各有不同；但總是在各種的威脅中生活著.不過當身體得贖時,則一切的威脅都可消除了!
\newline
\newline
「弟兄們,我告訴你們說,血肉之體,不能承受神的國,必朽壞的,不能承受不朽壞的.」(林前十五50).這是說主來時號筒末次吹響時,我們未死而得主救贖的,我們的身體就在眨眼之間,號筒吹響時,立刻改變成為不能朽壞的身體了.在(帖四14)說:「主必親自從天降臨,有呼叫的聲音；和天使長的聲音,又有神的號吹響.那時我們這活著還存留的人,必定與復活的人被提到空中；在雲中與主相遇,與主永遠同在,你們當用這些話彼此安慰.」我們身體得救贖,是贖回給主自己,這成為主的快樂,也是我們的快樂.求主給我們愛與主同在的心.
\newline
\newline
我們身體得贖時與主同在,可以見到祂的真體.在(約壹三2)說:「親愛的弟兄阿!我們現在是神的兒女,將來如何?還未顯明；但我們知道主若顯現,我們必要像祂,因為必得見祂的真體.」主來時,我們全人進入祂救贖之恩典之下,亞當那時所失去的,我們可得回,主的恩典過於亞當所失去的.我們也不要忘記,亞當墮落時,不但身體變了,連環境也變了!萬物伏在虛空之下,受那惡者的轄制.所以要等到主的國降臨在地上,將我們所失去的得回.祂不但救贖我們全人,且連環境也救贖了.我們的信仰與世上的理想家不同,他們以為把世界的環境改好,那就有平安了.但數千年來,罪既掌權,何來平安呢?只見每況愈下愈來愈不平安.按聖經的預言,只說在主快臨之前,世上一切只有一片混亂；所以這些幻想的平安,妄想的平安,是基督徒所不應有的.主說在大災難前,有序幕性的災難,即多處有飢荒,瘟疫和地震等等,現今的日子可見這一切都應驗了.我們盼望主再來,是盼望祂的福份,在彼前一章中說:專心盼望主來的時候帶給你們的恩典.提多二章說等候盼望的福,基督徒應想到這些福份就會盼望主來,等候主來,我們就不要愛慕地上的,因地上的一切都會過去；祗要一心靠主,盼望得著主再來的福份.所以應存盼望喜樂的心等候主再來.
\newline
\newline


\newpage

\section{第45屆~港九培靈會~胡恩德~第五講}
\label{sec:1234}
\textbf{基督徒的審判}
\newline
\newline
連結: \href{https://www.hkbibleconference.org/session-message/view/1234}{\texttt{https://www.hkbibleconference.org/session-message/view/1234}}
\newline
\newline
\hyperref[sec:1232]{< < < PREV SERMON < < <}
~
\hyperref[sec:toc]{[返目錄]}
~
\hyperref[sec:1235]{> > > NEXT SERMON > > >}
\newline
\newline
在聖經中我們知道主再來乃是一件重要的事.特別在新約中,每卷書的作者都提及.主說:天地可廢,但主的話不能廢去.祂說祂再來是一件永不可改變的事.
\newline
\newline
許多基督徒對主再來的態度很不正常,我們應以愛慕,並歡迎的心,去仰望及等待祂；對主的再來,不應存懼怕,或不理會的心.主再來該成為我們心中之喜樂,因為這是有福的盼望,最大的福份,是我們與主永遠同在(帖前四17),我們的靈魂體都得贖,並且見到主榮耀的真體.身體不再是朽壞的,也再無疾病的痛苦,所以我們應放膽等候主再來.祂不是要我們受刑,乃是要我們得榮；使我們不論是睡了的,或醒著的,都可以與祂同活過來.「睡」的意思是死(帖前四章),但在五章之前所說則是要我們儆醒,放心等主來,這並非靠我們的生活；乃是靠主之救贖,救主之寶血得免沉淪.
\newline
\newline
主來既是要審判,那就該有嚴肅的另一方面,也就是說我們該加緊準備.我們肯去準備,就是表示我們有歡迎祂來的心意.主是我們所愛的,所以應該準備；最好的準備,是出於愛心.主會說祂來是要施行審判,這審判對愛主的人來說,是一個得賞賜的好機會.正如一個好學生,對考試的來臨,並非存懼怕的心；而是覺得那是一個機會,一個福份,因藉這機會可考取獎學金.所以一個正常愛主的基督徒,對主再來,該認定是個歡喜的日子.
\newline
\newline
「那美好的仗我已經打過了,當跑的路我已經跑盡了,所信的道我已經守住了；從此以後,有公義的冠冕為我存留.就是按著公義審判的主,到了那日要賜給我的；不但賜給我,也賜給凡愛慕祂顯現的人.」(提後四7-8)
\newline
\newline
當時保羅快要離世時說:「我現在被澆奠,我離世的時候到了.」(提後四6)他想及主來的時候要將冠冕賜給他,因他沒有中途背道,美好的仗已打過,工作又積極有進取；又按主之分派,將主道傳開,像攻破城池.攻破了帖撒羅尼迦,雅典......,雖然曾受過不少苦難,卻沒因此退後,灰心.他的助手提摩太,因受不了逼迫而軟弱,保羅自己雖在監牢中,卻反而安慰他說:「你要在耶穌基督的恩典上剛強起來.」最後他說:「美好的仗我打過了,當跑的路我跑盡了.......」信主的人有一條路是要跑的,就是屬靈的路,應一站一站的去跑完它.許多人中途停止了,毫無長進!但保羅並不如此,他有長進,他的長進包括靈性上,工作上和爭戰上,他一切都跑盡了.在保羅所書寫的十多卷書信中,從未見他寫過怕主再來；他一直步步與主聯系,所以他知道有公義冠冕為他存留.他死,不過是去戴公義冠冕而已.
\newline
\newline
有人認為冠冕有好幾種:即生命的冠冕,榮耀的冠冕,不朽壞的冠冕,公義的冠冕,喜樂的冠冕等.這都是永遠的,是實在的,可與主一同操王權,一同行審判,分享主基督的大權柄.這冠冕不單是給保羅,同時也給與主有正常交通,並愛慕祂顯現的人.
\newline
\newline
基督徒對將來之審判能有此態度,就是正常的；如果沒有,在他靈程上,許多屬靈的路可能他未走過,屬靈的福份還未得著.有許多屬靈的事物,是要去追求的.倘若沒有追求的心,就不會是個成功的基督徒了.正如老底家教會,自以為富足,禮拜堂華麗美觀,經濟又充足,甚至可差派人去傳道.會友人數增加,主日學可分有若干班數.但可惜的是自以為滿足,失去了追求和渴慕主道的心,對聖經沒有興趣.這正是老底家教會的情形.這種人不會愛慕主再來,就算會,也不過是騙自己而已.我們應有實際長進的生命.試問有多少人能說,主再來我就歡喜呢?或許有人望死亡,因為在世上正處在困苦中,「死」可以解脫了一切.但主再來你喜歡嗎?主的審判你喜歡嗎?主的審判是嚴厲的,主來之日如「煉金之人的火」,(瑪拉基三2-3),「漂布之人的鹼,他必坐下如煉淨銀子的必潔淨利未人.」「熬煉他們像金銀一樣.」那就是說,主再來時先接教會,然後與眾信徒到地上來.以色列人要經過大災難,在(太廿四21-22)中說:「主曾說因那時必有大災難,是創世以來所未有的,以後亦不會有的,若不減少那日子.......那時的以色列人逃到山上去,城內的不可到城外,在屋頂的.......」以色列人特別受災難,在耶利米書稱之為雅各遭難的時候.在但十二章說那時必有大艱難,那也是指以色列有患難.瑪拉基說主再來時如煉金之人的火,要將利未的子孫煉淨；藉大災難要將事奉神的人煉淨,其實是要整個民族煉淨.因為他們屢次犯罪總不能擺脫.所以神先使他們亡於亞述,再而猶大支派亡於巴比倫.後在巴比倫時讓他們歸回,然後預備主來；神如此利害地將他們連根拔出,他們應該悔改才對.
\newline
\newline
他們自主前五三六年歸回,直至二十世紀一九七三年；他們雖不像以前的日子拜偶像,但他們裡面腐化的情形,卻比前更甚!所以主耶穌說:「那人身上的鬼被趕出去了,但另外帶了七隻更惡的鬼來,那人以後境況就更不如前了!」以色列人也是如此,神的兒子來了,他們不但不接待,反而把祂釘在十字架上,如此狠毒,他們的罪惡經被擄後仍未除去,並且更加倍利害,所以主耶穌為他們哭.主升天後三十多年,即紀元七十年,羅馬軍隊圍攻耶路撒冷,猶太人雖經死守頑抗,但大勢已去而敗亡.
\newline
\newline
至主後一三○多年猶太人又起義,但未能成功.羅馬人將他們擄去,並將他們散開到天下各地.至二十世紀,他們的罪仍未清.今日神的時候到了,讓他們返回聖地,重建一個國家,這是一個奇異的神蹟!但,這並非因他們已悔改,他們內裡的罪還無法除掉!神會用到最後的辦法,就是煉金之人的火,在教會被提之後；大災難就來到,特別有以色列的一份被煉淨.他們的心知道主要來,果然主來了!他們就大哭悔改,然後承認接受主耶穌是基督是彌賽亞,以後就再不背道了.神與他們立下永約,永約是新的約生效的約.由此可見,主再來要我們站得住,就要利害地煉淨我們,可知這審判不易通過.我們本來應受審判後到地獄的；因我們是罪人,我們的思想,行為,動機都要表露出來.主說:「到時一切都要供出來,我們有甚麼辦法能救自己呢?神是完全的,祂的眼睛鑒察世人無微不至,祂不會將祂的律法標準降低.神是完全的,祂的要求也要完全,所以我們不能救自己.世人都伏在神審判之下,神的話一點一畫都不能廢棄；我們有甚麼辦法成為完全的呢?這是不可能的事.但神卻用大能者──祂的獨生子完成了這件我們不能的事.榮耀偉大的主死在十字架上,才能使我們得救.那就是說,我們得救,是主救我們要脫離那使我們下地獄的審判.」
\newline
\newline
主說:「我實實在在的告訴你們,信我的人有永生.」「叫一切信祂的都得永生.」這一切祂已做妥了,難道今日我們隨隨便便地過基督徒生活就行了嗎?不能的,彼得說:「我們為人應該何等敬虔.」基督徒不能以為對「滅亡」和「永生」的問題已解決了,已經出死入生了,隨隨便便地就算.我們應憑著祂的靈去追求上進,像保羅一般；跑完那應跑的路,打那美好的勝仗.基督徒另外有一個審判,是關乎工作賞賜的.試看雅各和保羅如何提醒我們:在(雅五8-9)「你們也當忍耐,堅固你們的心；因為主來的日子近了.」弟兄們,你們不要彼此埋怨,免得受審判,看哪!審判的主站在門前了.」這是對信主者而言在(羅十四10-12):「你這個人,為什麼論斷弟兄呢?又為什麼看輕弟兄呢?因我們都要站在神的台前.」經上寫著:「主說我憑著我的永生起誓,萬膝必向我跪拜,萬口必向我承認.這樣看來,我們必要將自己的事,在神面前說明.」又在(林後五10),說:「因為我們眾人,必要在基督台前顯露出來；叫各人按著本身所行的,或善或惡受報.」所以我們應何等地聖潔,何等地敬虔的度日以待主再來.
\newline
\newline


\newpage

\section{第45屆~港九培靈會~胡恩德~第六講}
\label{sec:1235}
\textbf{主再來的初步豫兆}
\newline
\newline
連結: \href{https://www.hkbibleconference.org/session-message/view/1235}{\texttt{https://www.hkbibleconference.org/session-message/view/1235}}
\newline
\newline
\hyperref[sec:1234]{< < < PREV SERMON < < <}
~
\hyperref[sec:toc]{[返目錄]}
~
\hyperref[sec:1241]{> > > NEXT SERMON > > >}
\newline
\newline
馬太福音 24:1-25:46
\newline
\begin{longtable}{cl}
\hline
\hline
章節 & 經文 (和合本修訂版)\\
\hline
24:1 & \begin{tabularx}{0.7\textwidth}{X} 耶穌出了聖殿，正離開的時候，門徒前來，把聖殿的建築指給他看。 \end{tabularx} \\ \\ \relax
24:2 & \begin{tabularx}{0.7\textwidth}{X} 耶穌回應他們說：「你們不是看見這一切嗎？我實在告訴你們，這裡將沒有一塊石頭會留在另一塊石頭上，而不被拆毀的。」 \end{tabularx} \\ \\ \relax
24:3 & \begin{tabularx}{0.7\textwidth}{X} 耶穌在橄欖山上坐著，門徒私下進前來問他：「請告訴我們，甚麼時候有這些事呢？你來臨和世代的終結有甚麼預兆呢？」 \end{tabularx} \\ \\ \relax
24:4 & \begin{tabularx}{0.7\textwidth}{X} 耶穌回答他們：「你們要謹慎，免得有人迷惑你們。 \end{tabularx} \\ \\ \relax
24:5 & \begin{tabularx}{0.7\textwidth}{X} 因為將有好些人冒我的名來，說『我是基督』，並且要迷惑許多人。 \end{tabularx} \\ \\ \relax
24:6 & \begin{tabularx}{0.7\textwidth}{X} 你們也將聽見打仗和打仗的風聲。注意，不要驚慌！因為這些事必須發生，但這還不是終結。 \end{tabularx} \\ \\ \relax
24:7 & \begin{tabularx}{0.7\textwidth}{X} 民要攻打民，國要攻打國，多處必有饑荒、地震。 \end{tabularx} \\ \\ \relax
24:8 & \begin{tabularx}{0.7\textwidth}{X} 這都是災難的起頭。 \end{tabularx} \\ \\ \relax
24:9 & \begin{tabularx}{0.7\textwidth}{X} 那時，人要使你們陷在患難裡，也要殺害你們；你們又要為我的名被萬民憎恨。 \end{tabularx} \\ \\ \relax
24:10 & \begin{tabularx}{0.7\textwidth}{X} 那時，會有許多人跌倒，也會彼此陷害，彼此憎恨； \end{tabularx} \\ \\ \relax
24:11 & \begin{tabularx}{0.7\textwidth}{X} 且有好些假先知起來，迷惑許多人。 \end{tabularx} \\ \\ \relax
24:12 & \begin{tabularx}{0.7\textwidth}{X} 因為不法的事增多，許多人的愛心漸漸冷淡了。 \end{tabularx} \\ \\ \relax
24:13 & \begin{tabularx}{0.7\textwidth}{X} 但堅忍到底的終必得救。 \end{tabularx} \\ \\ \relax
24:14 & \begin{tabularx}{0.7\textwidth}{X} 這天國的福音要傳遍天下，對萬民作見證，然後終結才來到。」 \end{tabularx} \\ \\ \relax
24:15 & \begin{tabularx}{0.7\textwidth}{X} 「當你們看見先知但以理所說的那『施行毀滅的褻瀆者』站在聖地（讀這經的人要會意）， \end{tabularx} \\ \\ \relax
24:16 & \begin{tabularx}{0.7\textwidth}{X} 那時，在猶太的，應當逃到山上； \end{tabularx} \\ \\ \relax
24:17 & \begin{tabularx}{0.7\textwidth}{X} 在屋頂上的，不要下來拿家裡的東西； \end{tabularx} \\ \\ \relax
24:18 & \begin{tabularx}{0.7\textwidth}{X} 在田裡的，不要回去取衣裳。 \end{tabularx} \\ \\ \relax
24:19 & \begin{tabularx}{0.7\textwidth}{X} 在那些日子，懷孕的和奶孩子的就苦了。 \end{tabularx} \\ \\ \relax
24:20 & \begin{tabularx}{0.7\textwidth}{X} 你們要祈求，好讓你們逃走的時候，不遇見冬天或安息日。 \end{tabularx} \\ \\ \relax
24:21 & \begin{tabularx}{0.7\textwidth}{X} 因為那時必有大災難，自從世界的起頭直到如今，從沒有這樣的災難，將來也不會有。 \end{tabularx} \\ \\ \relax
24:22 & \begin{tabularx}{0.7\textwidth}{X} 若不減少那些日子，凡血肉之軀的，就沒有一個能得救；可是為了選民，那些日子將減少。 \end{tabularx} \\ \\ \relax
24:23 & \begin{tabularx}{0.7\textwidth}{X} 那時，若有人對你們說：『看哪，基督在這裡！』或『在那裡！』你們不要信。 \end{tabularx} \\ \\ \relax
24:24 & \begin{tabularx}{0.7\textwidth}{X} 因為假基督和假先知將要起來，顯大神蹟、大奇事，如果可能，要把選民也迷惑了。 \end{tabularx} \\ \\ \relax
24:25 & \begin{tabularx}{0.7\textwidth}{X} 看哪，我已經預先告訴你們了。 \end{tabularx} \\ \\ \relax
24:26 & \begin{tabularx}{0.7\textwidth}{X} 若有人對你們說：『看哪，基督在曠野裡！』你們不要出去；或說：『看哪，基督在內室中！』你們不要信。 \end{tabularx} \\ \\ \relax
24:27 & \begin{tabularx}{0.7\textwidth}{X} 好像閃電從東邊發出，直照到西邊，人子來臨也要這樣。 \end{tabularx} \\ \\ \relax
24:28 & \begin{tabularx}{0.7\textwidth}{X} 屍首在哪裡，鷹也會聚在哪裡。」 \end{tabularx} \\ \\ \relax
24:29 & \begin{tabularx}{0.7\textwidth}{X} 「那些日子的災難一過去，太陽要變黑，月亮也不放光，眾星要從天上墜落，天上的萬象都要震動。 \end{tabularx} \\ \\ \relax
24:30 & \begin{tabularx}{0.7\textwidth}{X} 那時，人子的預兆要顯在天上，地上的萬族都要哀哭。他們要看見人子帶著能力和大榮耀，駕著天上的雲來臨。 \end{tabularx} \\ \\ \relax
24:31 & \begin{tabularx}{0.7\textwidth}{X} 他要差遣天使，用大聲的號筒，從四方，從天這邊直到天那邊，召集他的選民。」 \end{tabularx} \\ \\ \relax
24:32 & \begin{tabularx}{0.7\textwidth}{X} 「你們要從無花果樹學習功課：當樹枝發芽長葉的時候，你們就知道夏天近了。 \end{tabularx} \\ \\ \relax
24:33 & \begin{tabularx}{0.7\textwidth}{X} 同樣，當你們看見這一切，就知道那時候近了，就在門口了。 \end{tabularx} \\ \\ \relax
24:34 & \begin{tabularx}{0.7\textwidth}{X} 我實在告訴你們，這世代還沒有過去，這一切都要發生。 \end{tabularx} \\ \\ \relax
24:35 & \begin{tabularx}{0.7\textwidth}{X} 天地要廢去，我的話卻絕不廢去。」 \end{tabularx} \\ \\ \relax
24:36 & \begin{tabularx}{0.7\textwidth}{X} 「但那日子，那時辰，沒有人知道，連天上的天使也不知道，子也不知道，惟有父知道。 \end{tabularx} \\ \\ \relax
24:37 & \begin{tabularx}{0.7\textwidth}{X} 挪亞的日子怎樣，人子來臨也要怎樣。 \end{tabularx} \\ \\ \relax
24:38 & \begin{tabularx}{0.7\textwidth}{X} 在洪水以前的那些日子，人照常吃喝嫁娶，直到挪亞進方舟的那日， \end{tabularx} \\ \\ \relax
24:39 & \begin{tabularx}{0.7\textwidth}{X} 不知不覺洪水來了，把他們全都沖去。人子來臨也要這樣。 \end{tabularx} \\ \\ \relax
24:40 & \begin{tabularx}{0.7\textwidth}{X} 那時，兩個人在田裡，一個被接去，一個被撇下。 \end{tabularx} \\ \\ \relax
24:41 & \begin{tabularx}{0.7\textwidth}{X} 兩個女人推磨，一個被接去，一個被撇下。 \end{tabularx} \\ \\ \relax
24:42 & \begin{tabularx}{0.7\textwidth}{X} 所以，你們要警醒，因為不知道你們的主哪一天來到。 \end{tabularx} \\ \\ \relax
24:43 & \begin{tabularx}{0.7\textwidth}{X} 你們要知道，一家的主人若知道晚上甚麼時候有賊來，就必警醒，不讓賊挖穿房屋。 \end{tabularx} \\ \\ \relax
24:44 & \begin{tabularx}{0.7\textwidth}{X} 所以，你們也要預備，因為在你們想不到的時候，人子就來了。」 \end{tabularx} \\ \\ \relax
24:45 & \begin{tabularx}{0.7\textwidth}{X} 「那麼，誰是那忠心又精明的僕人，主人派他管理自己的家僕、按時分糧給他們的呢？ \end{tabularx} \\ \\ \relax
24:46 & \begin{tabularx}{0.7\textwidth}{X} 主人來到，看見僕人這樣做，那僕人就有福了。 \end{tabularx} \\ \\ \relax
24:47 & \begin{tabularx}{0.7\textwidth}{X} 我實在告訴你們，主人要派他管理所有的財產。 \end{tabularx} \\ \\ \relax
24:48 & \begin{tabularx}{0.7\textwidth}{X} 如果那惡僕心裡說：『我的主人會來得遲』， \end{tabularx} \\ \\ \relax
24:49 & \begin{tabularx}{0.7\textwidth}{X} 就動手打他的同伴，又和醉酒的人一同吃喝， \end{tabularx} \\ \\ \relax
24:50 & \begin{tabularx}{0.7\textwidth}{X} 在想不到的日子，不知道的時候，那僕人的主人要來， \end{tabularx} \\ \\ \relax
24:51 & \begin{tabularx}{0.7\textwidth}{X} 重重地懲罰他，定他和假冒為善的人同罪，在那裡他要哀哭切齒了。」 \end{tabularx} \\ \\ \relax
25:1 & \begin{tabularx}{0.7\textwidth}{X} 「那時，天國好比十個童女拿著燈出去迎接新郎。 \end{tabularx} \\ \\ \relax
25:2 & \begin{tabularx}{0.7\textwidth}{X} 其中有五個是愚拙的，五個是聰明的。 \end{tabularx} \\ \\ \relax
25:3 & \begin{tabularx}{0.7\textwidth}{X} 愚拙的拿著燈，卻沒有帶油； \end{tabularx} \\ \\ \relax
25:4 & \begin{tabularx}{0.7\textwidth}{X} 聰明的拿著燈，又盛了油在器皿裡。 \end{tabularx} \\ \\ \relax
25:5 & \begin{tabularx}{0.7\textwidth}{X} 新郎遲延的時候，她們都打盹，睡著了。 \end{tabularx} \\ \\ \relax
25:6 & \begin{tabularx}{0.7\textwidth}{X} 半夜有人喊：『看，新郎來了，你們出來迎接他。』 \end{tabularx} \\ \\ \relax
25:7 & \begin{tabularx}{0.7\textwidth}{X} 那些童女就都起來挑亮她們的燈。 \end{tabularx} \\ \\ \relax
25:8 & \begin{tabularx}{0.7\textwidth}{X} 愚拙的對聰明的說：『請分點油給我們，因為我們的燈要滅了。』 \end{tabularx} \\ \\ \relax
25:9 & \begin{tabularx}{0.7\textwidth}{X} 聰明的回答：『恐怕不夠你我用的；你們還是自己到賣油的那裡去買吧。』 \end{tabularx} \\ \\ \relax
25:10 & \begin{tabularx}{0.7\textwidth}{X} 她們去買的時候，新郎到了。那預備好了的，與他進去共赴婚宴，門就關了。 \end{tabularx} \\ \\ \relax
25:11 & \begin{tabularx}{0.7\textwidth}{X} 其餘的童女隨後也來了，說：『主啊，主啊，給我們開門！』 \end{tabularx} \\ \\ \relax
25:12 & \begin{tabularx}{0.7\textwidth}{X} 他卻回答：『我實在告訴你們，我不認識你們。』 \end{tabularx} \\ \\ \relax
25:13 & \begin{tabularx}{0.7\textwidth}{X} 所以，你們要警醒，因為那日子，那時辰，你們不知道。」 \end{tabularx} \\ \\ \relax
25:14 & \begin{tabularx}{0.7\textwidth}{X} 「天國又好比一個人要出外遠行，就叫了僕人來，把他的家業交給他們。 \end{tabularx} \\ \\ \relax
25:15 & \begin{tabularx}{0.7\textwidth}{X} 他按著各人的才幹，給他們銀子：一個給了五千，一個給了二千，一個給了一千，就出外遠行去了。 \end{tabularx} \\ \\ \relax
25:16 & \begin{tabularx}{0.7\textwidth}{X} 那領五千的立刻拿去做買賣，另外賺了五千。 \end{tabularx} \\ \\ \relax
25:17 & \begin{tabularx}{0.7\textwidth}{X} 那領二千的也照樣另賺了二千。 \end{tabularx} \\ \\ \relax
25:18 & \begin{tabularx}{0.7\textwidth}{X} 但那領一千的去掘開地，把主人的銀子埋藏了。 \end{tabularx} \\ \\ \relax
25:19 & \begin{tabularx}{0.7\textwidth}{X} 過了許久，那些僕人的主人來了，和他們算賬。 \end{tabularx} \\ \\ \relax
25:20 & \begin{tabularx}{0.7\textwidth}{X} 那領五千的又帶著另外的五千來，說：『主啊，你交給我五千。請看，我又賺了五千。』 \end{tabularx} \\ \\ \relax
25:21 & \begin{tabularx}{0.7\textwidth}{X} 主人說：『好，你這又善良又忠心的僕人，你在少許的事上忠心，我要派你管理許多的事，進來享受你主人的快樂吧！』 \end{tabularx} \\ \\ \relax
25:22 & \begin{tabularx}{0.7\textwidth}{X} 那領二千的也進前來，說：『主啊，你交給我二千。請看，我又賺了二千。』 \end{tabularx} \\ \\ \relax
25:23 & \begin{tabularx}{0.7\textwidth}{X} 主人說：『好，你這又善良又忠心的僕人，你在少許的事上忠心，我要派你管理許多的事，進來享受你主人的快樂吧！』 \end{tabularx} \\ \\ \relax
25:24 & \begin{tabularx}{0.7\textwidth}{X} 那領一千的也進前來，說：『主啊，我知道你，你是個嚴厲的人：沒有種的地方也要收割，沒有播的地方也要收穫， \end{tabularx} \\ \\ \relax
25:25 & \begin{tabularx}{0.7\textwidth}{X} 我就害怕，去把你的一千銀子埋藏在地裡。請看，你的銀子在這裡。』 \end{tabularx} \\ \\ \relax
25:26 & \begin{tabularx}{0.7\textwidth}{X} 他的主人回答他說：『你這又惡又懶的僕人，你既知道我沒有種的地方也要收割，沒有播的地方也要收穫， \end{tabularx} \\ \\ \relax
25:27 & \begin{tabularx}{0.7\textwidth}{X} 就該把我的銀子放給兌換銀錢的人，到我來的時候可以連本帶利收回。 \end{tabularx} \\ \\ \relax
25:28 & \begin{tabularx}{0.7\textwidth}{X} 把他這一千奪過來，給那有一萬的。 \end{tabularx} \\ \\ \relax
25:29 & \begin{tabularx}{0.7\textwidth}{X} 因為凡有的，還要加給他，叫他有餘；沒有的，連他所有的也要奪過來。 \end{tabularx} \\ \\ \relax
25:30 & \begin{tabularx}{0.7\textwidth}{X} 把這無用的僕人丟在外面黑暗裡，在那裡他要哀哭切齒了。』」 \end{tabularx} \\ \\ \relax
25:31 & \begin{tabularx}{0.7\textwidth}{X} 「當人子在他榮耀裡，同著眾天使來臨的時候，要坐在他榮耀的寶座上。 \end{tabularx} \\ \\ \relax
25:32 & \begin{tabularx}{0.7\textwidth}{X} 萬民都要聚集在他面前。他要把他們分別出來，好像牧人分別綿羊、山羊一般， \end{tabularx} \\ \\ \relax
25:33 & \begin{tabularx}{0.7\textwidth}{X} 把綿羊安置在右邊，山羊在左邊。 \end{tabularx} \\ \\ \relax
25:34 & \begin{tabularx}{0.7\textwidth}{X} 於是王要向他右邊的說：『你們這蒙我父賜福的，可來承受那創世以來為你們所預備的國。 \end{tabularx} \\ \\ \relax
25:35 & \begin{tabularx}{0.7\textwidth}{X} 因為我餓了，你們給我吃；渴了，你們給我喝；我流浪在外，你們留我住； \end{tabularx} \\ \\ \relax
25:36 & \begin{tabularx}{0.7\textwidth}{X} 我赤身露體，你們給我穿；我病了，你們看顧我；我在監獄裡，你們來看我。』 \end{tabularx} \\ \\ \relax
25:37 & \begin{tabularx}{0.7\textwidth}{X} 義人就回答：『主啊，我們甚麼時候見你餓了，給你吃；渴了，給你喝？ \end{tabularx} \\ \\ \relax
25:38 & \begin{tabularx}{0.7\textwidth}{X} 甚麼時候見你流浪在外，留你住；或是赤身露體，給你穿？ \end{tabularx} \\ \\ \relax
25:39 & \begin{tabularx}{0.7\textwidth}{X} 又甚麼時候見你病了，或是在監獄裡，來看你呢？』 \end{tabularx} \\ \\ \relax
25:40 & \begin{tabularx}{0.7\textwidth}{X} 王回答他們說：『我實在告訴你們，這些事你們做在我弟兄中一個最小的身上，就是做在我身上了。』 \end{tabularx} \\ \\ \relax
25:41 & \begin{tabularx}{0.7\textwidth}{X} 「王又要向那左邊的說：『你們這被詛咒的人，離開我！進入那為魔鬼和他的使者所預備的永火裡去！ \end{tabularx} \\ \\ \relax
25:42 & \begin{tabularx}{0.7\textwidth}{X} 因為我餓了，你們沒有給我吃；渴了，你們沒有給我喝； \end{tabularx} \\ \\ \relax
25:43 & \begin{tabularx}{0.7\textwidth}{X} 我流浪在外，你們沒有留我住；我赤身露體，你們沒有給我穿；我病了，我在監獄裡，你們沒有來看顧我。』 \end{tabularx} \\ \\ \relax
25:44 & \begin{tabularx}{0.7\textwidth}{X} 他們也要回答：『主啊，我們甚麼時候見你餓了，或渴了，或流浪在外，或赤身露體，或病了，或在監獄裡，沒有伺候你呢？』 \end{tabularx} \\ \\ \relax
25:45 & \begin{tabularx}{0.7\textwidth}{X} 王要回答：『我實在告訴你們，這些事你們沒有做在任何一個最小的弟兄身上，就是沒有做在我身上了。』 \end{tabularx} \\ \\ \relax
25:46 & \begin{tabularx}{0.7\textwidth}{X} 這些人要往永刑裡去；那些義人要往永生裡去。」 \end{tabularx} \\ \\
[1ex]
\hline
\hline
\end{longtable}


\newpage

\section{第45屆~港九培靈會~胡恩德~第七講}
\label{sec:1241}
\textbf{主再來與有關猶太人的預兆}
\newline
\newline
連結: \href{https://www.hkbibleconference.org/session-message/view/1241}{\texttt{https://www.hkbibleconference.org/session-message/view/1241}}
\newline
\newline
\hyperref[sec:1235]{< < < PREV SERMON < < <}
~
\hyperref[sec:toc]{[返目錄]}
~
\hyperref[sec:1242]{> > > NEXT SERMON > > >}
\newline
\newline
馬太福音 24:15-24:31
\newline
\begin{longtable}{cl}
\hline
\hline
章節 & 經文 (和合本修訂版)\\
\hline
24:15 & \begin{tabularx}{0.7\textwidth}{X} 「當你們看見先知但以理所說的那『施行毀滅的褻瀆者』站在聖地（讀這經的人要會意）， \end{tabularx} \\ \\ \relax
24:16 & \begin{tabularx}{0.7\textwidth}{X} 那時，在猶太的，應當逃到山上； \end{tabularx} \\ \\ \relax
24:17 & \begin{tabularx}{0.7\textwidth}{X} 在屋頂上的，不要下來拿家裡的東西； \end{tabularx} \\ \\ \relax
24:18 & \begin{tabularx}{0.7\textwidth}{X} 在田裡的，不要回去取衣裳。 \end{tabularx} \\ \\ \relax
24:19 & \begin{tabularx}{0.7\textwidth}{X} 在那些日子，懷孕的和奶孩子的就苦了。 \end{tabularx} \\ \\ \relax
24:20 & \begin{tabularx}{0.7\textwidth}{X} 你們要祈求，好讓你們逃走的時候，不遇見冬天或安息日。 \end{tabularx} \\ \\ \relax
24:21 & \begin{tabularx}{0.7\textwidth}{X} 因為那時必有大災難，自從世界的起頭直到如今，從沒有這樣的災難，將來也不會有。 \end{tabularx} \\ \\ \relax
24:22 & \begin{tabularx}{0.7\textwidth}{X} 若不減少那些日子，凡血肉之軀的，就沒有一個能得救；可是為了選民，那些日子將減少。 \end{tabularx} \\ \\ \relax
24:23 & \begin{tabularx}{0.7\textwidth}{X} 那時，若有人對你們說：『看哪，基督在這裡！』或『在那裡！』你們不要信。 \end{tabularx} \\ \\ \relax
24:24 & \begin{tabularx}{0.7\textwidth}{X} 因為假基督和假先知將要起來，顯大神蹟、大奇事，如果可能，要把選民也迷惑了。 \end{tabularx} \\ \\ \relax
24:25 & \begin{tabularx}{0.7\textwidth}{X} 看哪，我已經預先告訴你們了。 \end{tabularx} \\ \\ \relax
24:26 & \begin{tabularx}{0.7\textwidth}{X} 若有人對你們說：『看哪，基督在曠野裡！』你們不要出去；或說：『看哪，基督在內室中！』你們不要信。 \end{tabularx} \\ \\ \relax
24:27 & \begin{tabularx}{0.7\textwidth}{X} 好像閃電從東邊發出，直照到西邊，人子來臨也要這樣。 \end{tabularx} \\ \\ \relax
24:28 & \begin{tabularx}{0.7\textwidth}{X} 屍首在哪裡，鷹也會聚在哪裡。」 \end{tabularx} \\ \\ \relax
24:29 & \begin{tabularx}{0.7\textwidth}{X} 「那些日子的災難一過去，太陽要變黑，月亮也不放光，眾星要從天上墜落，天上的萬象都要震動。 \end{tabularx} \\ \\ \relax
24:30 & \begin{tabularx}{0.7\textwidth}{X} 那時，人子的預兆要顯在天上，地上的萬族都要哀哭。他們要看見人子帶著能力和大榮耀，駕著天上的雲來臨。 \end{tabularx} \\ \\ \relax
24:31 & \begin{tabularx}{0.7\textwidth}{X} 他要差遣天使，用大聲的號筒，從四方，從天這邊直到天那邊，召集他的選民。」 \end{tabularx} \\ \\
[1ex]
\hline
\hline
\end{longtable}
這段經文是門徒問主降臨,與世界的末了有什麼預兆,於是主就這樣回答他們,我的題目是關乎預兆的.現今的人看預兆有兩種的態度:
\newline
\newline
(一)聽預兆──只是聽而已,他們把預兆當作新奇的事去聽.
\newline
\newline
(二)怕聽預兆──聽了各種預兆,又見已一一地應驗了,就怕聽預兆.
\newline
\newline
預兆是會幫助我們的.我們應該喜樂,不必害怕.在第一至四世紀中,教會常有逼迫,那時他們只想到主快來的預兆,心裡就立刻興奮；因為一宿雖然有哭泣,早晨便必歡呼!可能他們有如保羅那樣,有患難,也有嘆息與勞苦；
\newline
\newline
願這世代快過去,更願主快來.不過,我們更應因愛主而願主快來!因祂是整個教會的愛人,故祂來我們就該歡喜.現今已有預兆給我們顯明了,那就該趁主來前,多作主工.
\newline
\newline
我們知道大災難之前奏的預兆,已逐漸臨到,只是末期還未到.主還說包括末世大逼迫,被萬民恨惡.基督徒應知逼迫來到就有戰爭,我們一定要靠主；因為逼迫一直至普遍化,而至成為世界性的時候,那麼主來的日子就更近了.在15及21節中都有「大災難」三個字,在耶利米書有說及；雅各遭難之時,男人像產難的婦人,用手捏腰,臉面都變青了.哀哉!那日為大,無日可比.(耶卅6-7)在但以理書十二章一節亦有提及:「必有大艱難,從有國以來,直到此時,沒有這樣的.」這是指以色列遭遇是空前絕後的災難.(太廿四21)「那時,必有大災難,從世界的起頭,直到如今,沒有這樣的災難,後來也必沒有.」啟示錄亦說:全世界皆落在大災難中,地變,海變,天變和人也變；國際中有許多變動,這一切都是大災難.有人為的,有神降的,亦有魔鬼給的.這災難的日子不過幾年,這災難有一重點是對付以色列人.故耶利米書說:雅各遭難之時,在大災中,神像是用量器量給世人,又特別量一份給以色列人.(耶卅15)因為他們之罪惡,在他們裡面始終沒有清除；神要最後一次熬煉他們,直至他們差不多全民族滅亡之時；主才特別顯出祂恩典臨到他們.唯有恩典能化人心,在他們無法抵受之時,即耶路撒冷被毀之時；以往阿拉伯人要將以色列人打退回地中海去,卻不能成事；但在大災難之後期,這事似乎應驗了.不但阿拉伯及四周的各國,還有遠方國家來攻,(聖經未有明說是誰?)那時,耶路撒冷就被攻了.當以色列人正絕望之時,主耶穌就突然降臨.(亞九9)「必有救主從錫安而出,要消除那國家一切的罪過.」主來向他們顯現「救了他們,然後他們必仰望那自己所扎過的,所恨過的人,就是他們的主彌賽亞.他們就放聲大哭,如喪獨生子,如喪長子,」(亞十二10)主說我們在耶路撒冷開一個洗罪的泉源,要洗淨雅各家的罪惡.(亞十三1)那時他們見主拯救他們,而不是向他們報復.主在恩典中向他們顯現,正如主向掃羅顯現一樣.以色列將在大災難後得主拯救,而消除了雅各家一切的過犯,他們就永不再離開主了.主將新心,新靈,放在他們心中.這大災難是為了以色列人,神最後要對付他們；這是大災難的目的,同時也是以色列人應得的報應.神在西乃山與他們所立的約.早已清楚說明:他們若不遵守神的道,神就要如何的報應他們.神的話在過去都應驗了,直到近世史亦如此記載.正如以色列民必在萬國中,被拋來拋去.雖然以色列人中,亦漸有悔改信主的青年,不過還未能說是全體的歸主,他們一直還是抵擋主.仍以釘主在十字架的那種態度,來拒絕主.神要等待,直至最後一次對付他們.今日他們雖散處世界各國,假如神降災的焦點是對付他們,以成全神的計劃；那麼散在各處,就不易對付了.所以主在(太廿四16)預言特別說明.「猶太」這名字,祂要先將他們集合在本地,才作最後的對付.撒迦利亞書說,主向他們顯現時是在橄欖山；如今他們已回國集合了,這是說明神最後對付以色列人的步驟快將實現了.主至今仍未再來,是因他們還未悔改.正如主耶穌第一次來時,他們要擁戴他作王.他們看見五餅二魚神蹟時,又要祂作王；可是他們不悔改,主又怎會作他們的王呢?直至今天,他們仍未悔改,反而反對基督教.這態度與從前一樣沒改變.主又怎會來,那麼神又如何將他們集合起來呢?是要他們等候大災難的日子來到,亦可說等待彌賽亞降臨.
\newline
\newline
教會是在大災難前先被接去的,所以我們不要再遲延,應快悔改.現在他們既已回本國,則我們亦應準備；不是準備受災,而是準備被提.有解經的人說,那時的大災難,是指主後七十年建那次的事重演.當時羅馬軍要攻入聖城,猶太人拼死保衛耶路撒冷,城內血流如溪河,情形極慘!查考聖經,則可知這大災難並非指這一次,而是雙關的.因在(太廿四29-30)說:「那些日子的災難一過去,日頭就變黑了,月亮也不放光,眾星要從天上墜落,天勢都要震動.那時人子的兆頭要顯在天上,地上的萬族都要哀哭；他們要看見人子,有能力,有大榮耀,駕著天上的雲降臨.」這裡說明災難「一過去」,馬上主耶穌就降臨,即是說兩件事是緊接著的,按本人的見解,在16-17節說:「那時,在猶太的,應當逃到山上,在房上的,不要下來拿家裡的東西.」那是說要急走,因事情來得急切,正像如今的爭戰一般,可見這並非指主後七十年之時.
\newline
\newline
今日猶太人已回本國,我們知大災難已快到,而且情形比以前更是厲害.現今地震,瘟疫,大旱......已不斷在最近發生.可知主快再來.主叫我們挺身昂首地等祂回來,讓我們也歡歡喜喜地迎接祂.當主來時,日,月,星辰都改變了,眾星要從天墜落,天勢都要震動.人會因海中波浪的響聲就慌張不定,海裡深淵泉源都變動(路廿一章).天地要被烈火焚燒,彼得說:這一切既如此銷化,我們為人就應何等敬虔,何等聖潔地切切盼望主的日子來到(彼後三10-11)我們應該輕看屬地的,屬物質的和暫時的；而看重那屬靈的,屬天和永遠的.
\newline
\newline


\newpage

\section{第45屆~港九培靈會~胡恩德~第八講}
\label{sec:1242}
\textbf{敵基督顯現與主再來}
\newline
\newline
連結: \href{https://www.hkbibleconference.org/session-message/view/1242}{\texttt{https://www.hkbibleconference.org/session-message/view/1242}}
\newline
\newline
\hyperref[sec:1241]{< < < PREV SERMON < < <}
~
\hyperref[sec:toc]{[返目錄]}
~
\hyperref[sec:1244]{> > > NEXT SERMON > > >}
\newline
\newline
馬太福音 24:15-24:16
\newline
\begin{longtable}{cl}
\hline
\hline
章節 & 經文 (和合本修訂版)\\
\hline
24:15 & \begin{tabularx}{0.7\textwidth}{X} 「當你們看見先知但以理所說的那『施行毀滅的褻瀆者』站在聖地（讀這經的人要會意）， \end{tabularx} \\ \\ \relax
24:16 & \begin{tabularx}{0.7\textwidth}{X} 那時，在猶太的，應當逃到山上； \end{tabularx} \\ \\
[1ex]
\hline
\hline
\end{longtable}
(徒二章)載,彼得在五旬節日向眾人見證主,他引證約珥預言,世界將有大事發生；這是指著主那大而可畏的日子而言.關於主來的日子,保羅稱之為「那日」,或是「主的日子」,「基督的日子.」.那大而可畏的日子,到底如何的大?可畏到甚麼程度呢?乃是有大能力顯出,使天變,地變,海變；地球和地上的社會,國際之間皆有大變動.
\newline
\newline
主來到地上,地上的萬族都要因祂哀哭!可知那是何等大的變動.摩西死,以色列人為他哀哭四十日；因摩西的死是一大的損失,不過這損失只關乎一個民族.地上萬族都要哀哭,乃是有重大的力量要臨到世界.
\newline
\newline
這力量第一是審判的力量,全世界都要伏地這審判的力量之下,是無人能逃脫或逃避的.隨著審判而來的是,這刑罰乃是從主而來的刑罰,是永遠的,是大而可畏的.所以人間有哀哭,有人躲到山洞地穴中叫喊說:山啊!倒在我身上吧!因怕見羔羊和坐寶座者的面目.主來是以威嚴,能力,審判而來,甚至於天象要倒塌下來.天上如此,地上也如此,是空前的,是極其可怕的!在(但九27太廿四15)說:(那行毀壞可憎的站在聖地,那時,在猶太的應當逃到山上.)所指的那人非從神而來,是有撒旦的力量,能統治全世界,他是空前的.但以理見那怪獸,同其他不同的在一起,這就是說,這國度可怕到極.現今人間的國度也要演變到這地步,那強大的武力是無從幻想的,而使世界所有聯合在他掌握中.當洪水來時,有大之能力在變動；將來主來時,亦發生很大的能力；是空前的,因創造之主親自來到.祂自己收拾這世界,自己對付惡人.在未到之前有地震,亦是大有力量的.今日神讓人發展人的力量,達到了最高峰；但人是惡人,是反叛的,所以所發展的,乃是空前敵擋神的力量.魔鬼也知道牠最後的時候到了,牠也盡所能地,把力量擺在敵基督身上.因此,有許多假基督,假先知,和邪靈出來,這都是魔鬼的力量.其中最厲害是敵基督,他發揮黑暗的撒旦力量,集合人的力量.人聯合起來的力量甚大.創世記中所記,人聯合起來建造巴別塔,要與神為敵,與神的旨意相反.神將人的口音變亂了,使人不能合作.人既不能合作,也就不能成功了.但在末世時,人與要神作一次最後之戰爭；因人在敵基督手下,人的力量也變了空前的大.這非人的力量,而是魔鬼的力量,將來必然發展到這地步.
\newline
\newline
敵基督行毀壞可憎的,在新舊約皆有記載.現今那敵基督者將要出來,(是神所許可的)在永遠的計劃裡面,我們會見到敵基督統治這世界.當敵基督出現時,會有下列之情況:
\newline
\newline
(一)交通方便──從來沒有一位統治者,能把全世界擺在自己的掌握中.但現今世界,地球好像縮小了；因為交通發達到高峰,數千里遙遠的路一剎那間即可到.電視,廣播可傳遞命令,傳遞主義,傳遞政治.這些條件,正適合由一個人出來統治全世界.所以現今乃是臨近敵基督出來之時.
\newline
\newline
(二)人的心已準備好──敵基督是與神對敵者,今日世人的心,愈來愈趨向反對神,這敵基督乃是個大無神論者.在(帖後二3-4)說:「人不拘用什麼法子,你們總不要被他誘惑,因為那日子以前;必有離道反教的事,並有那大罪人,就是沉淪之子,顯露出來.他是抵擋主,高抬自己.超過一切稱為神的,和一切受人敬拜的；甚至坐在神的殿裡,自稱是神.(但十一36-37)「王必任意而行,自高自大,超過所有的神.又用奇異的話,攻擊萬神之神.他必行事亨通,直到主的忿怒完畢,因為所定的事必然成就.他必不顧他列祖的神,也不顧婦女所羡慕的神；無論何神他都不顧,因為他必自大,高過一切.」現今人之無神風氣,已遍佈各地,成人,小孩都在懷疑.魔鬼使人心不想及神,高抬自己超過一切稱為神的.另做一神以適應現今人之思想及社會之情形,他們以為神是人想像中的,其實創造宇宙的主宰,祂將自己的一部份啟示給人.人離開了祂,就會失去了這啟示,但內裡總是知道有的.於是人就憑藉自己那虛妄墮落的思想,被罪惡破壞了的思想,驕傲的思想,去想出許多古怪的神,這是在人權力之下的.
\newline
\newline
他們什麼都不顧,只以人為大,不要神.以為人之科學及能幹已足,不需要神.人的心從未進步及改善過,人只有狂傲,放縱情慾,兇殺和殘忍,一天比一天厲害,人與人間紛爭結黨,互相攻擊謀害.這種精神及態度,乃是敵基督者所有,乃正如保羅說:末世危險的日子已臨近,是敵基督出現有大事發生之時.
\newline
\newline
(三)大團結──現今世界趨向於大團結,國際間亦如此,而這種聯合非聖經所說之聯合.這正是準備那統治者來臨,在(啟十一章)中可見:現今世界正有相當的準備,這敵基督的到來.世上的逼迫甚多,逼迫主的教會及基督徒.此乃敵基督來前之現象,是神與撒旦的戰爭,(撒旦以人為工具).(賽十二12)說:人與人戰爭,那時人比精金還少.故我們不要以為可在此安然居住,以為兒女長大了可安心享福.其實這思想不對的,我們應該要等主來,因就算我們是光明之子,也需要靠著主那非常偉大的力量來站穩.信仰方面需要站穩,而且要去得人.看現今的世界情況,可知主快再來了.
\newline
\newline
主駕雲降臨,是能力的降臨.空中也有能力的表現,現今太空也有能力,平時我們不覺得,其實細心研究,那對我們都有關係的.但將來能力更大,主身上的能力顯出.在啟示錄中提及神之寶座有雷轟閃電之聲.這可能是象徵,亦可能實在有,這雷轟的雷是聲音,我們可以聽到,閃電是在視覺中可見到,但另一件事是「能」.那是說神是有「能」不斷地發,現今天地仍不斷有「能」發出.祂不是揮動祂的手,而是用祂的言語.耶穌基督在地上表現祂的能力時；人們覺得驚奇,認為神蹟.當顯出祂的能力時,多是用祂的言語.其實神也是用祂的言語天地組織好,一直至如今,不過我們不覺得而已.但到主來時,是祂親自來,祂本身的「能」何等厲害!滿天滿地都是「能」.今日我們尋求主,是有救恩方面的能力,我們接觸並領受；我們心裡有信心,主的「能」就可在我們身上,叫我們服從的心,那麼「能」就有大.故歷代的信徒,他們身上能顯出一種屬靈的「能力」,至今仍有此「能力」將來更顯在自然界中,滿天地皆是.因而說:這是個大的日子.那日子要來到,也就是光明的日子.耶穌基督與新婦有大榮耀之時,我們該準備此大日子的來臨.
\newline
\newline
耶穌基督的大日子,固然是祂由死復活,(即今主日),但祂再來也是「主之大日子」這日子實在大!祂不單是自己由死裡復活,祂來了,使一切信祂的都由死裡復活.祂來要施行大審判,將榮耀賞賜祂的眾僕人,我們有沒有準備祂來臨呢?(詩廿九2)說:「以聖潔的裝飾敬拜耶和華.」我們要有聖潔的準備,在(啟十九8)中也說聖徒要穿細麻衣,即聖徒所行的義.今日我們不單是有神賜給我們的義,更要追求實際的行義.保羅在羅馬書六章中說「稱義」,即實際去行.「義」即是「對的」,也是說對的就做,不對的不去做.羅馬三章及彼前二章都說及這道理.假如我們違法的話,在政府來說,我們是不義；在神方面,我們更是不義.違背神的旨意就是不義和不對；所以我們身為聖徒,就應有聖徒的義,並且要行義.保羅說:「你們既然蒙召,行事為人就當與蒙召的恩相稱.」(弗四1)有了此恩給你,你行事為人就當與之平衡.我們應追求公義和聖潔.在(太廿四42,)說:「所以你們要儆醒,因為不知道你們的主,是那一天來到.」在(路廿一34):也說:「你們要謹慎!恐怕因貪食醉酒,並今生的思慮,累住你們的心,那日子就如同網羅,忽然臨到.」求主給我們有儆醒的心,常常準備,追求公義聖潔,有聖徒所行的義,等候主來.
\newline
\newline


\newpage

\section{第45屆~港九培靈會~胡恩德~第九講}
\label{sec:1244}
\textbf{星光與永福}
\newline
\newline
連結: \href{https://www.hkbibleconference.org/session-message/view/1244}{\texttt{https://www.hkbibleconference.org/session-message/view/1244}}
\newline
\newline
\hyperref[sec:1242]{< < < PREV SERMON < < <}
~
\hyperref[sec:toc]{[返目錄]}
~
\hyperref[sec:1197]{> > > NEXT SERMON > > >}
\newline
\newline
「凡遵守命令的,必不經歷禍難；」(傳八5)「延長年日」(賽五三10)「你們受洗歸入基督的,都是披戴基督了.並不分猶太人,希利尼人,自主的,為奴的,或男或女；因為你們在基督耶穌裡,都成為一了.」(加三27-28)這裡題到三種沒有分別:不分國籍,不分階級,不分性別.「直到我們眾人在真道上同歸於一.」(弗四13)(腓立比書一章5,12,13,二22)都是論到興旺福音的事,教會的一切都應當為著福音.
\newline
\newline
為何要禮拜,是為著福音.為何要奉獻金錢?為何要配搭建立身體?都是為著興旺福音.保羅在腓立比書題到他為福音被捆鎖,題到多人因福音與他同勞.可見福音的工作是需要付代價的,而保羅為了福音的工作是需要付代價的,他實在付上了很大的代價.在哥林多後書保羅題到他曾經受很多的苦；被鞭打,被石頭打,遇陸上及海上的危險.他到馬其頓去,後來他在監牢裡；被人打至皮破血流.他在以弗所傳福音時,當地的人煽惑群眾暴動,幾乎要把他撕碎.保羅為著福音的緣故,實在付上了代價.
\newline
\newline
中國的福音是由外國的宣教士傳來的.那些宣教士也實在付上了代價.他們離別國境親人,遠路而來.從前保羅從耶路撒冷到羅馬,那時交通不便,沒有飛機,沒有輪船；需經長久時間才可達到.美國一個宣教士克爾威廉,他到印度去傳道,要航行五個月的水程.馬禮遜費了八個月的時候,才自英國到達廣州傳道.戴德生搭了一艘四百七十噸的輪船,經七個月水程才來到中國.外國的傳教士來中國傳道,實在付上了很多的代價.他們離別家園,要往人地生疏的異國,通常要半年才有家信,言語又不通,無法與當地人交往,備嘗孤獨之苦!每個國家都有它宗教背景,要使他們棄假歸真,實不容易.傳道工作是艱難的,困苦的!此外,還有性命的危險.當義和團得勢的時候,很多宣教士殉道被殺了.為了福音,是需要付上大代價的.
\newline
\newline
福音有種種的關係:
\newline
\newline
福音與別人的關係　耶穌把福音帶到撒瑪利亞.有個婦人,正午出去打水.在猶太地方,天氣酷熱,人們打水,都在清晨或傍晚的；但這婦人卻在正午出來打水；因為她是個不名譽的女人,避免露面見人.但耶穌把福音傳給她之後,她完全改變過來,她就跑到城裡見那些她平時所不願見的人.告訴他們:「有一人把我素來的事都說出來了.」福音能改變人!
\newline
\newline
在工業發達的時代裡,最易使人的情緒煩躁；為瞭解脫煩躁,人尋找刺激；尋求娛樂的刺激求解脫.更有甚者,是找尋色情的刺激,麻醉的刺激.許多人飲酒來麻醉自己,以毒品來麻醉自己.因為情緒不安煩躁,若不求解脫,就不能活下去.
\newline
\newline
福音與個人的關係　保羅說:「若不傳福音,我就有禍了.」(林前九16下半)有些基督徒,以為這話是神威脅我們,其實不然.(傳道書八5)說:「凡遵守命令的,必不經歷禍患.」反過來說:不遵守命令的,就必經歷禍患了.主說:「神叫日頭照好人,也照歹人.」照這話解釋,不單歹人遭遇禍患,連好人也會遭遇禍患.也可說,不信的人會遭遇禍患,信的人也會遭遇禍患.那麼,信與不信有何分別?其中是大有分別的.不信的人遭遇,沒有真神可倚靠,也沒有真神可求告.至於信者適得其反,真神是我們的倚靠,我們可以向他求告庇佑.
\newline
\newline
「萬軍之耶和華說,我曾呼喚他們,他們不聽,將來他們呼求我,我也不聽.」(亞七13)我想這就是基督徒的禍患.沒有一件事比福音更大,神呼喚我們在福音上有份,我們卻視如飄風之過耳；到了我們真有需要,向神呼求時,他也不聽了.(以弗所六15)說:「用平安的福音,當作預備走路的鞋穿在腳上.」鞋穿在腳上可以幫助走路,保護腳底不至受傷,也可防釘子刺傷,玻璃割傷.腳穿了鞋,免得沾染污穢,還可增加美觀.保羅用福音的鞋作預表,就是說,我們向人傳福音,好像鞋穿在腳上一樣.
\newline
\newline
福音與教會的關係　教會擔起傳福音的責任,如(以賽亞五三10)說:耶穌要延長年日,耶和華以他為贖罪祭,把眾人的罪孽都歸在他身上.意思就是說,耶穌被釘在十字架上,把命舍了,後來埋葬,復活,升天.聖經說他要延長年日.「延長」可作「擴大」解釋.耶穌基督的年日,不但要延長,且要擴大.他把這樣的責任交給教會,教會要擔當延長耶穌的責任.教會是個身體.神造了亞當,對他說:「要生養眾多,遍滿地面.」生養眾多是延長的意思.(路加三章)說:「亞當之後有塞特,塞特之後有以挪士,」一代代傳下去.亞當所負的,是「延長」之責,而且要遍滿地面,意思就是也要「擴大」.亞當的後裔,挪亞三個兒子,含到非洲地方,閃在亞洲地方,雅弗在歐洲地方；他們擔負擴大的責任.亞當的身體所擔負的工作,是世世代代,各處各方的工作.教會的身體也照樣要擔起延長和擴大的責任.如果這一代不傳與下一代,基督的年日怎能延長呢?如果此地不傳到彼地,基督怎能擴大呢?所以傳福音是神給教會的大使命.
\newline
\newline
從聖經看來,福音的工作,是一直要擴大的,耶穌說:「天國好比芥菜種.」芥菜種是要擴大的.耶穌在約翰福音十章說他另外有羊,不但猶太人中間有他的羊；外邦人中間也有他的羊.耶穌復活以後,叫門徒要作見證,從耶路撒冷直到地極,這都是要擴大的.讓聖靈充滿在人的身上,叫人有說各地方言的恩賜,其目的也是要把福音擴大.後來,他給保羅有馬其頓的異象,也是要保羅把福音擴大；把福音從歐洲擴大到亞洲去.由耶穌時代至使徒時代,神一直把這旨意顯給我們.耶穌不但要延長下去,還要擴大出去.擔負延長和擴大,是教會的責任,如果教會不擔起這責任,教會在神面前將如何交代?教會不只有人守禮拜,有人奉獻金錢,就可以滿足；教會的一切,當為著福音.今天的教會,不明白屬靈的定律.屬世的定律是越積越多；屬靈的定律卻是越分越多,越流越遠.從五餅二魚的神蹟來看,在自己的手裡只夠供一人之需,分了出去,可供無數人之用.如以西結四七章所說,水越流越深,開始時在踝子骨,漸而到膝,到腰.屬靈和屬世的定律完全不同；屬世的,要抓往它,留著它；屬靈的,要放開它,要分出它,因而愈加豐富.在猶太地方,有個鹽海,又叫死海.也有個加利利海,又叫活海.活海裡有生物,海面常有雀鳥飛翔.死海則連鳥都不敢飛過,恐被海裡的蒸氣蒸倒下來.一死一活的海,同是約但河的水流下,同出一源,為何一死一活呢?原因很簡單；所以死是因水流下囤積起來.所以活是因水流下去以後,不斷又分流出去.教會屬靈光景的死,是因我們從神所得的,不肯向人分出去；結果,教會死氣沉沉.
\newline
\newline
怎樣興旺福音　保羅在腓立比書一章,教我們要同心合意,福音才能興旺.
\newline
\newline
同心合意必有聖靈同在(同工)從舊約多處可以看到:以利亞到迦密山上,把十二塊石頭砌在一起,他以十二塊石頭代表十二支派；因那時十二支派已經分開了,因羅波安壓迫百姓之故,分裂為南北兩國,就是猶太和以色列.以利亞把十二塊石頭砌疊在一起,是同心合意的預表.他把壇築好,向神禱告,就有火降下.可見同心合意,就有聖靈的工作.詩篇告訴我們:「看哪!弟兄和睦同居,是何等的善,何等的美.這好比那貴重的油,在亞倫的頭上流到鬍鬚,又流到他的衣襟.」(詩一三三1-2)和睦同居,也是同心合意.若同心合意,聖靈就澆灌下來.耶穌說:「地上若有兩三個人奉我的名聚會,我就在他們中間.」兩三個人也是同心合意；這樣,主就與我們同在.五旬節時,教會有聖靈同工；其中有兩個最顯著的現象,其一,是不斷推展出去,從耶路撒冷推展到撒瑪利亞,一直推展到整個亞細亞.其二,有聖靈同工,教會就有推展的能力.還有一現象,如果有聖靈同工,就有神蹟奇事出現.彼得,有神蹟奇事隨著他.保羅,也有神蹟奇事隨著他.神蹟奇事足以證明神的恩道,叫一個頭腦很大的人,在他面前,不能不俯伏下來.
\newline
\newline
有一件事我曾多次提過:人在身體強壯時,常會自高自傲；到了面臨死亡邊緣時,便沒有甚麼可誇可傲之處了.有個大學的校長,他患了嚴重的心臟病；當我向他傳福音時,他好像很聽話的孩子一樣.次日,他對我說他要受水洗,護士聽見了,連忙拉一下我的衣服,示意叫我離開病房.對我說:『吳長老!你千萬不要為他施浸,因為主診醫生吩咐要特別小心照料他,因他的心臟脆弱至極；若稍受刺激,恐隨時有性命危險.如果你把他浸在水裡,恐怕到時拉不上來呀!」我覺得她言之有理,萬一死在水裡,消息給報紙刊登出來,個人的名譽不要緊,可是教會的名譽就不得了.我正思索之際,校長在房裡叫我,他再三叮囑明晨為他施浸的事.那時,我不便對他說明.各位!難道他信耶穌,他願意接受,而不為他施浸,豈不矛盾嗎?那時我只好答應了.次日早晨,我很早起來禱告,因我極擔心會出事,我把這事交給神.到八點鐘,他的兒子陪他乘汽車,如約來到禮拜堂門口,他不讓兒子扶進去,他邊走邊喘息地進到裡面,其實極短的距離,卻走了很長的時間.後來他的兒子扶他下水.我沒有忘記,那時我禱告兩句話,一是為自己禱告,求神給我看見榮耀.一是為他禱告,求神給他經歷榮耀.接著,我問他:「你信耶穌是神的兒子嗎?答:「信」,我就把他浸在水裡了,隨即把他拉起來,我看見他面露笑容,當他兒子要扶他時,他卻把兒子推開,自己走動.他不需再回到醫院,而返東海大學,繼續執教.他就是東海大學的張啟東校長.教會同心合意,就有神蹟奇事,藉著神蹟奇事,能夠證明神的真道.
\newline
\newline
同心合意就有榮美的見證(加拉太書三章27-28)我們看見三種的「沒有分別」第一是不分國籍,不分猶太人或希利尼人.各國有不同的規矩,如中國人較為好客,宴客時,常常把筷子放在自己咀裡啜一下,然後夾菜敬賓,表示親切.如果請的是外國客人,真會把他嚇壞了；但如果同在主裡面,他就能受得起.西國人對於時間很寶貴,當你造訪,逾時還不告辭,他便會下逐客令的；他下逐客令,我們仍受得了.這些見證,都是很美的.第二是不分階級,就是不分貧富,在教會中,無論貧富,都交通得來,配搭得來.富的和貧的交通,不以為是遷就；貧的和富的配搭也並非高攀.這樣的見證多麼美,第三是不分性別,就是不分男女.並非不男不女,乃是表明我們在教會裡,男的女的,都一齊事奉,謹謹慎慎的,清清潔潔的,這個見證實在很美!所以一個教會同心合意,有榮美的見證,才能叫社會的人景仰和欽佩.
\newline
\newline
同心合意就有滿足的喜樂(腓立比書一至三章)論到喜樂.第四章說要常常喜樂,大大喜樂.為什麼有喜樂呢?第一章題到同心合意,第二章題到很多相同之處,意思是大家同心合意就有喜樂.這裡所謂相同(一樣),並非我們的個性相同,也非興趣一樣.那麼,到底有何相同,有何一樣呢?神所造的都不相同,才顯出他的完全,才顯出他的美麗.我們的個性,有急的,有慢的,有內向的,有外向的.如果大家一樣的急,或一樣的慢,那怎麼辦呢?如果大家都屬內向,有事大家都不講；這樣,人與人之間,難題就大了.人的興趣也不相同,有的喜歡遊山,有的喜歡玩水,有山有水,才顯出美麗.聖經說,我們的心相同,向著神的心相同,為著神的事工心要相同,這樣教會就有喜樂.我們在一起時,氣氛和諧.我們一起事奉,工作輕鬆,這樣的教會,內部有歌可唱,外面有力量可以發展.所以,同心合意,教會才有滿足的喜樂.
\newline
\newline
我們應在何事上同心合意呢?
\newline
\newline
一,在真道上同心合意　保羅在以弗所四章說「在真道上同歸於一.」真道是信仰的意思,我們是在信仰上要同心合意.
\newline
\newline
二,在感情上同歸於一　別讓感情把我傷,也不要傷感情.何謂感情把我傷?在尼希米記,有個人名叫多比雅,他外面與神有關,裡面卻與神無關.他按名是活的,其實是死的.多比雅是阿們族.是羅得和他女兒所生的後裔.這種人不能入耶和華的會,可是他卻住在聖殿裡,而且參於事奉.因他是當時大祭司以利亞實的親戚,為了感情的緣故,容許多比雅在聖殿,這就是感情把他傷了.今天中國的教會,給感情傷了的事,在在皆是,明知那人行為不對,明知他得罪神,但為了情面,不敢指摘.至於傷感情最易引起的由於口,和由於發生誤會.言語不小心和種種誤會,彼此傷了感情,感情一傷,目標雖是相同,步伐不能一致.興旺福音是我們共同的目標,為了感情受傷,步伐不能一致,福音不能傳開,神的家大受虧損.
\newline
\newline
三,在事奉上同心合意　尼希米建造城牆,聖經說:「某某人對著城牆修造.」這話何解?意思是我這段的修造和你那段的修造有關.教會也是如此,你的事奉和我的事奉有關,彼此的事奉要互相配搭；否則,身體就支離破碎,不能發生功用.在教會中,各個團契,息息相關,應一致行動,對付這邪惡的世界.
\newline
\newline
韓國教會,在若干年前發生了一件事:當時韓國政府擬印行一種新鈔票,票上印有和尚之像.教會風聞此事,即上書政府提出抗議；理由是鈔票既屬全國性使用,但全國百姓,有佛教徒,也有基督徒,所以不能以和尚代表所有的百姓.政府當局認為此乃政府之決策,教會無權干涉,故置之不理.後來教會再次上書,聲明若政府印行這種鈔票,全國教會將堅持拒用.結果,政府只好取消計畫.由此可見,整體的力量能夠影響國家和社會.弟兄姊妹!我們應在信仰上同心合意,在感情上同心合意,在事奉上同心合意,教會才有能力；否則,沒有聖靈同工,沒有榮美的見證,也沒有滿足的喜樂.
\newline
\newline
求主今天晚上,把「同心合意」四個字放在我們心裡,為著福音同心合意,好叫福音興旺起來!
\newline
\newline


\newpage

\section{第46屆~港九培靈會~于力工~第一講}
\label{sec:1197}
\textbf{基督是我們的愛人}
\newline
\newline
連結: \href{https://www.hkbibleconference.org/session-message/view/1197}{\texttt{https://www.hkbibleconference.org/session-message/view/1197}}
\newline
\newline
\hyperref[sec:1244]{< < < PREV SERMON < < <}
~
\hyperref[sec:toc]{[返目錄]}
~
\hyperref[sec:1198]{> > > NEXT SERMON > > >}
\newline
\newline
雅歌 1:1-1:4
\newline
\begin{longtable}{cl}
\hline
\hline
章節 & 經文 (和合本修訂版)\\
\hline
1:1 & \begin{tabularx}{0.7\textwidth}{X} 所羅門的雅歌。 \end{tabularx} \\ \\ \relax
1:2 & \begin{tabularx}{0.7\textwidth}{X} 願他用口與我親吻。你的愛情比酒更美， \end{tabularx} \\ \\ \relax
1:3 & \begin{tabularx}{0.7\textwidth}{X} 你的膏油馨香，你的名如傾瀉而出的香膏，所以童女都愛你。 \end{tabularx} \\ \\ \relax
1:4 & \begin{tabularx}{0.7\textwidth}{X} 願你吸引我跟隨你；讓我們快跑吧！王領我進入他的內室。我們必因你歡喜快樂，我們要思念你的愛情，勝似思念美酒。她們愛你是理所當然的。 \end{tabularx} \\ \\
[1ex]
\hline
\hline
\end{longtable}
「所羅門的歌,是歌中的雅歌.願他用口與我親嘴,因你的愛情比酒更美.你的膏油馨香,你的名如同倒出來的香膏.所以眾童女都愛你.願你吸引我,我們就快跑跟隨你.王帶我進了內室,我們必因你歡喜快樂.我們要稱讚你的愛情,勝似稱讚美酒,他們愛你是理所當然的.」
\newline
\newline
神造人有一個標準,這標準乃是祂自己,祂照著自己的形象造人.在一切被造之物中,惟有人接受了這至尊至榮的崇高地位.但自從人犯罪後,這標準被打破了!被毀壞了!感謝神,差遣耶穌來世,成全救恩；使我們從前在亞當裡所失去的,在基督裡得以恢復過來.回復原來的標準,有神的性情,有神的樣式.若有人要像父神,就應先認識基督.現今我們要從雅歌書中認識基督.
\newline
\newline
許多人認為雅歌書是一本談情說愛的書,為何會列入聖經內呢?若要明白這原因,就必須先要了解當時的背景.
\newline
\newline
雅歌書是所羅門在主前九世紀時寫的.這書已漸成為人們心中的安慰,並使人看見到神的心意,和著摸神的愛.在主前五世紀時,猶大人被擄到外邦,受盡痛苦,他們懷疑神,為何使他們受苦.當年有兩卷書,實在是使他們得安慰的:其一是約伯記──義人為神受苦.其次就是雅歌書,這書使他們看見神是他們的愛人.神既是他們的愛人,就決不會撇下他們,所以他們得安慰.神照樣愛我們,我們的痛苦乃是至輕至暫的；神要用永遠不變的愛愛我們.有誰能害我們和敵擋我們呢?讀此書就可使我們得安慰和滋潤!
\newline
\newline
有人覺得這卷書不易了解.但,若把它當作主說話的口吻去讀,就容易明白了.我們之與基督,正如良人之與新婦.基督與教會的愛,也是基督與我們的愛.
\newline
\newline
其次要明白作書時的背景,因為希伯來人男女間的愛情,他們親熱的態度與中國人相比,是有分別的.如親嘴擁抱是他們的習慣,我們卻並不如此.若能先了解,就不至引起誤會了.
\newline
\newline
「願他用口與我親嘴,因你的愛情比酒更美.」(歌一2)
\newline
\newline
「願你吸引我,我們就快跑跟隨你,王帶我進了內室......」(歌一4)
\newline
\newline
這兩節經文,是指我們與基督的兩種密切關係,──「親嘴」和「進入內室」──是相關的.沒有親嘴的關係,就絕對不能進入內室去,這種關係有重要的意思:
\newline
\newline
(一)祂是我們的生命氣息.(歌一2)
\newline
\newline
「口」,在創世記二章中曾說:當神造成了亞當之後,他仍毫無氣息；神就用口將生氣吹進亞當裡,於是亞當就成了有靈的活人.故「親嘴」,乃代表生命氣息進入亞當裡.「親嘴」也表示是我們與基督的關係.主乃是生命氣息的源頭,我們要得祂的供應,就必要在祂裡面,不斷地吸取祂一切的豐盛.吸取祂呼出來的生命,這是我們生命中重要的操練.因祂的口發出良言,口中的話如蜂房下滴的蜜甘甜.主口中的氣息能使我們生命活發,和感到滿足.我們常忽略了自己直接到主面前領受供應,若單是靠聽,支取供應,這是不正常的.是片面的.為甚麼我們不能直接向主吸取呢?應親自到主前,將裡面的問題向主呼出,然後吸入主供給的豐富.祂的話,沒有一句不帶著能力的；祂的話就是靈,是生命.主說:「我就是道路,真理,生命；若不藉著我,沒有人能到父那裡去.」生命的長進,屬靈的領受,裡面的安慰與甘甜,這一切都應直接地從主領受.
\newline
\newline
為甚麼主會成為我們的生命氣息呢?在(加二20)說:「我已經與基督同釘十字架,現在活著的,不再是我,乃是基督在我裡面活著.」「祂是愛我,為我捨己.」主為我們捨命,救我們脫離罪,為了成全救恩,為了我們的需要而將生命捨棄了!祂在十字架上所經歷的痛苦,為要使我們得醫治,得平安.唯有我們的主,才能成為我們的生命氣息.才能使我們生命得著滿足.因祂是生命的源頭,我們要繼續吸取祂生命的豐盛與需要.
\newline
\newline
(二)祂的膏油馨香──(歌一3)膏油可以代表:
\newline
\newline
代表醫治及滋潤.
\newline
\newline
a 得醫治的力量:「因祂受的鞭傷,我們得醫治.」(賽五三5).當主的膏油抹在我們身上時,我們就得醫治.
\newline
\newline
b 靈命得醫治:我們犯罪時,祂用膏油膏我們,我們就得赦免.
\newline
\newline
c 心靈得滋潤:我們在貧困中,缺乏時,祂就供應.有問題,祂就解決.在孤單時,祂作我良伴；經死蔭幽谷時,祂與我同在,這是恩膏的作用.主自己就是恩膏,因祂是先知,祭司,和君王.
\newline
\newline
代表按立──在舊約時代,作君王的需用膏油膏抹,作先知和作祭司的,需用膏油膏抹.當祭司在神前執行職務時,也要用膏油抹胸.
\newline
\newline
(三)祂有決定的權柄──祂有權柄,祂永遠作王.在洪水泛濫之時祂坐著為王.祂要管理,帶領,並引導我們走義路.所以我們接受這膏油時,該知道主是我們的王,祂會向我們說話；當我們禱告時,也聽到祂的聲音.
\newline
\newline
多少時候我們的禱告會落在單程交通的情況中,只有我們向主訴說,卻沒有機會讓主說話,所以聽不到主的聲音.你若聽見主說話時,就有平安與滿足的喜樂.若你犯罪時,祂就會責備你.在你裡面是否有先知的膏油?祂是祭司為我們代求,正如亞伯拉罕,他無論到何處,都築起祭壇,向神獻祭.假如恩膏在我們裡面,則無論在任何時間,任何工作中,都可憑著裡面的膏油,隨時築壇向主禱告.主基督是膏油的源頭,膏油在裡面就可以有權柄,你可奉祂的名叫撒旦退去；也可對世界說:「世界啊!你再沒有引誘我的力量了.」
\newline
\newline
(四)分別為聖
\newline
\newline
分別為聖的意思是,指這人,這事,這物永不再為己用,而單為主用.
\newline
\newline
基督是我們的吸引,是我們的方向和道路.我們一生要緊緊的跟隨主.主說:「你們要背起十字架來跟從我.」這是基督徒的道路,除此以外再沒有別的道路.保羅也說:「忘記背後,努力面前,向著標竿直跑.」因基督是我們的吸引,當祂吸引我們時,我們就快跑跟隨祂.
\newline
\newline
祂吸引我們到什麼地方去?乃是到「內室」裡去,(歌一4)內室是新郎新婦的地方,也是中國人所說的洞房.內室裡只有兩人,不能有第三者.這是我與主同在之地,也是我們的目標和方向.天天與主同在一起,魔鬼不能參在裡面,世界也不能在其中.若以會幕的名詞說,那是至聖所,也是保羅所說的三層天,以賽亞書所說的隱密處.
\newline
\newline
今日教會有一種普通現象,就是十八歲以上的青年人漸漸地減少,因為他們有男女間的約會.二十五歲至四十歲的人,也同樣漸少；因為他們婚後被家庭,兒女.......各樣問題纏擾而少到教會.我們今要進入內室,表明我與主單單同在,讓主與我說話.試問你今天是否讓家庭,世界和自己參入了你和基督之間的關係呢?求主恩待我們.
\newline
\newline


\newpage

\section{第46屆~港九培靈會~于力工~第二講}
\label{sec:1198}
\textbf{基督是我們的牧人}
\newline
\newline
連結: \href{https://www.hkbibleconference.org/session-message/view/1198}{\texttt{https://www.hkbibleconference.org/session-message/view/1198}}
\newline
\newline
\hyperref[sec:1197]{< < < PREV SERMON < < <}
~
\hyperref[sec:toc]{[返目錄]}
~
\hyperref[sec:1199]{> > > NEXT SERMON > > >}
\newline
\newline
我們讀雅歌書應注意的,就是這卷書並非一本愛情的故事.本書描寫良人與新婦從結識以至結婚,把心中的感覺,用不同的次序和與不同的人訴說出來；因此讀這卷書時,會容易發生困難.
\newline
\newline
從書中可見到書拉密女雖然知道良人愛她,她也愛良人.但在她心的深處,並不完全了解他,所以發生許多的問題和困難.這也是現今神的兒女們一般的困難.我們對主的認識常是片面的,判斷也是片面的；當我們遭遇問題時,就用片面的判斷而決定問題,以至引起錯誤!有人因此而跌倒,甚至懷疑所信的基督.
\newline
\newline
我們雖與主認識了,但卻不是完全的認識.正如在這三處經文中,書拉密女應認識她的良人是牧人,雖然她知道他是牧放群羊,但卻不知怎樣地牧養,這是她有所偏差的原因.現分三方面思想這問題:
\newline
\newline
(一)要跟隨主及尋求主的問題
\newline
\newline
「你這女子中極美麗的,你若不知道,只管跟隨羊群的腳蹤去.」(歌一8)這是耶路撒冷眾女所說的話.她們既與良人在一起,為甚麼叫她跟羊群的腳蹤去呢?因為書拉密女雖愛良人,但卻不一定天天跟隨著他,有時候離開了,又要尋找祂.祂問眾人,她的良人那裡去了?我們何嘗不是如此,走自己的路；不知道如何跟隨主.其實並非主離開我們,而是我們走了另一方向,找不到主；還以為主離開了我們,這是我們的錯誤!
\newline
\newline
基督徒夫妻間到底是誰順服誰呢?是順服妻子從神處所領受來的,抑是順服丈夫從神處所領受來的呢?按屬靈的原則及聖經教訓而言,丈夫乃一家之主,是頭,這並非就誰比較大；乃是要有屬靈的次序,總要一人順服另一人.由此看來,可以得到一個真理,就是我們是照自己的心意,判斷事情,去認識所要認識的；抑或是讓基督完全擺在我們面前,去認識我們所該去認識的呢.書拉密女若認識良人的話,就只管跟隨著他吧!
\newline
\newline
當主在世時,曾對門徒說:「我要到耶路撒冷,在那裡受迫害,被凌辱,被釘在十字架上.」當時彼得說:「主啊,你萬不可如此.」主耶穌說:「撒旦!退過後面去吧.」主為何要嚴厲斥責他呢?因為他不知道主的心意.他未摸到真理,並未真正認識基督；他只照自己的觀念,來決定耶穌基督的道路.所以主耶穌說:「若有人要跟從我,就當背起他的十字架來跟從我.」今天我們所走的道路,只管跟隨主,若照我們自己的看法,走自己的道路,持自己的意見,那麼到了一個時候,你就會對主說:「主啊!你在那裡?為何看不見你,為何不聽我的禱告!」其實這問題不在乎主,乃是在乎我們沒有跟從祂,沒有照祂的方向走,沒照祂的心意行.祂到那裡去,我也到那裡去；祂所指示我的,我也就去遵行.當效法主說:「神啊!祭物和禮物是你不願意的,你曾給我預備了身體,燔祭和贖罪祭是你不喜歡的,那時我說神呵!我來了,為要照你的旨意行.」(來十5-6)
\newline
\newline
各位神的兒女,今日的道路乃是單純的跟隨.因在我們的觀念裡,凡是亞當的子孫,都有自己的主張,自己的看法.在自己的看法裡就決定了自己的道路,於是我們就走錯了方向.
\newline
\newline
要找你的良人嗎?只管跟著羊群的腳蹤走去.這些羊如只隨著牧人走,就會走出一條道路.你若想知道主在何處,就先要察看羊群留下的腳蹤.一個跟隨主的人,也照樣會產生了腳蹤和一條道路,讓別人只需看見這腳蹤和道路,就知道牧人的所在.今天我們留下了什麼榜樣?什麼腳蹤?什麼道路?這是應該注意的,我們不單是要跟隨主,更要留下一條道路,讓這道路顯明時,別人就能立刻判斷出來,跟著這腳蹤走.所以我們對主之認識與了解,不單是對我們的生命,及靈命有關係,而且對別人的追求也有關係.如果我們走錯了路,也會把別人帶錯路,和走錯了方向.
\newline
\newline
我們的主在世上,祂一生所走的道路是十字架的道路,而祂留下了佳美的腳蹤.在希伯來書說:祂在城門外受苦,這樣,我們也當出到營外,就了祂去,受祂所受的凌辱.(來十三12-13),祂留下走十字架道路的蹤跡,今日我們也背著十字架出到城外去,一同經歷基督所經歷的.因我們這樣才可以知道主的所在.在雅歌書中,我們看到基督的樣式,也表示我們應經歷及應走的路.許多時候我們以為跟隨主,是要得主賜福和恩典,當遇到患難困苦,或會使我們不明白.其實主是我們的良人,祂的確愛我們,為我們捨命.我們跟隨主,要到城外去受凌辱.經上說:我們只管跟著羊群的腳蹤走.把你的山羊,牧放在牧人帳棚的旁邊.那是說只管將你的重擔放在主面前,祂在何處歇息,你就在何處停止.祂在何處發動,你就在何處行走,這樣,你的心才得著安慰,才能緊緊地跟隨主.跟隨主的道路,並非想像是一帆風順的道路,我們常因順境富裕而感恩,但卻從未為悲傷或缺乏,而感謝與讚美!更未能為所遭遇的,認定是變相的祝福.我們跟隨主必須要按著祂的標準來跟隨的.祂的帳棚在那裡,我們的山羊也在那裡；山羊乃是為獻燔祭用的,經過這火的焚燒才能發出香味.因此火祭就成為馨香祭,神也因我們所獻的而得喜悅.
\newline
\newline
(二)「良人屬我,我也屬祂,祂在百合花中牧放群羊.」──(歌二16)
\newline
\newline
在詩篇二十三篇中說:「耶和華是我的牧者,我必不至缺乏,祂使我躺臥在青草地上,領我在可安歇的水邊.」那麼為何在此說是牧放在百合花中?我們應知道,百合花不是生長於山頂,而是生在平原和山谷裡.山頂有陽光,山谷沒陽光,可能還是死蔭的幽谷,但卻長出百合花.當牧人牧放羊群時,會把群羊帶到山頂上,可能書拉密女喜歡一直住在山頂上,但良人卻不到山頂去,直到第四章才說及山頂的光景.可見主帶我們帶我們到山谷裡去,我們就要去.
\newline
\newline
百合花是代表得勝和聖潔之意,是經過死亡然後得勝,主常帶領我們經死蔭之幽谷.在幽谷裡,常有死亡的蔭影籠罩著,死亡壓在身上；使你感到死,不過只離開你一步.又像是快嘗到死的味道,這時候,你就會經歷到主的得勝.主在十字架上經歷了死亡,然後經歷復活.十架壓力是大的.死亡的權勢也在墳墓裡,但最後主說:「死啊!你的毒勾在那裡?死亡的權勢在那裡?」主向一切宣告得勝.在山谷中開出了百合花,是潔白的,聖潔的,得勝的!主並非單要我們經歷高山美麗的風景,而是也要我們在山谷中長出百合花.教我們身上彰顯出主的榮耀.感謝主!我們的主並非尋常的牧人!
\newline
\newline
(三)在園內牧放羊群
\newline
\newline
「我的良人下自己園中,到香花畦,在園內牧放群羊,採百合花.」(歌六2)
\newline
\newline
在這經文中,可知這女子失去了良人的蹤跡,但她知道良人下入自己園中,在園內牧放群羊.我們的主有祂自己的地方與範圍,祂要獨自計劃和主張,這是祂的自由.今日認識基督,跟隨祂的人,就應將自己的主張,自由......都要交在主手裡.祂到園子裡去.我們也要跟祂到園子裡.就算祂要關我們在病房裡,我們也該在病房裡.或要工場裡,婦女會,青年團契,無論在那裡只要是主的原則,我們就要順服.因為在園子裡,可以學到功課,可以使我們有長進.
\newline
\newline
有許多人,因對傳道人或長老不滿意,就會批評,繼而轉到別的教會去.這觀念是錯誤的!如果你逃避,換教會,並不能解決你的問題,因為你自己就有問題.我們該知道,無論主帶領到何處,我們就應到何處.若是祂要我們在祂園子裡,是祂所喜悅的地方,那麼我們就該歡喜.
\newline
\newline
在蓋恩夫人傳記中的最後一段說──當蓋恩夫人被關在巴士丁的監牢中時,就寫出以下的一首詩歌說:我好像一隻小鳥,自由飛翔在空中,但當小鳥被關在籠子裡,失去了自由時；為了神的緣故,只要你喜歡!我也喜歡.」如果今日主的原則是將你關在籠子裡,那麼,你也該說:為了主的緣故,我肯!為了主的緣故失去自由,我肯.為了主的緣故失去了健康,我肯我也喜歡；失去了財物,主啊我也喜歡.各位!你能唱首詩歌嗎?也許你今日不能唱,但要學習能唱,即使你未能走到此地步,仍然要一步一步的走,要走進這園子裡.
\newline
\newline
到香花畦在園內牧放群羊,在那兒採百合花.餵羊本是用草的,但為何會在這裡要採百合花呢?在前段是說山谷裡的百合花,在這段是說園子裡的百合花.在前段是在死亡蔭影裡,在此是在神的旨意主權裡.牧人要我們到這光景裡,為要我們裡面放出花朵來.並願有一天,新婦對新郎說:請你進到我的園子裡,好吃我佳美的果子.祂今日帶我們進入園子裡.要我們照著祂的規模,樣式,道路能長出百合花.可見屬靈的功課一步一步的學習,在不同的光景和不同的情形下,學習不同的功課,這是我們常遇的道路,在不同的惡劣環境裡要長出百合花.在困苦中要站立得穩,在優裕的環境裡能放出香氣；在困難中能打美好的勝仗,在舒暢的日子裡也要活出美好的樣式.舊約先知以利亞,能戰勝四百五十個假先知,但因一句說話而失敗逃命!所以自以為站立得穩的,就當謹慎,免得跌倒.無論處在何境況,都要長出百合花.
\newline
\newline
「我屬我的良人,我的良人也屬我」(歌六3)
\newline
\newline
「良人屬我,我也屬祂」(歌二16)
\newline
\newline
這兩節經文中一節是「良人屬我」,即以「我」為主.另一節是「我屬我的良人」,那是以「良人」為主.我屬良人.這是書拉密女經一切之後,在她真正進入園子後,有了純正的規模後,就知道了「我屬我的良人.」然後說:祂在百合花中牧放羊群.基督的立場是不變的,只有她自己在轉變.主今日向我們要求,要我們改變,從以「我」為中心,轉到以基督為中心.從前是基督屬我,現在是我屬基督.我們應進入這境地.多少時我只認識主一半,所以在生活上出了偏差,生命不能進步,在啟示錄十四章說:有十四萬四千人要跟隨羔羊,不論羔羊到何處,他們都跟隨祂.他們在基督榮耀的大場面出現,與基督出現一樣.因為他們知道基督是祭司,也是祭牲.可惜,許多人只知道基督是大祭司,而不知道也是所獻的祭牲；知是好牧人,而不知也是羔羊.知道基督君王,而知道主也是出城外的囚犯；知道祂是先知,而不知主是聽話順服的人.(詩四十7-8)說:我的事在經卷上記載了,我的神啊!我樂意照你的旨意行.我們認識基督是神,也是人,現今要跟隨祂,一方面享受祂所預備的一切,另一方面也接受祂為我們預備的一切,那麼才能開放出百合花,讓我們無論到何處,都跟隨著祂.
\newline
\newline


\newpage

\section{第46屆~港九培靈會~于力工~第三講}
\label{sec:1199}
\textbf{基督是我們的榮耀}
\newline
\newline
連結: \href{https://www.hkbibleconference.org/session-message/view/1199}{\texttt{https://www.hkbibleconference.org/session-message/view/1199}}
\newline
\newline
\hyperref[sec:1198]{< < < PREV SERMON < < <}
~
\hyperref[sec:toc]{[返目錄]}
~
\hyperref[sec:1201]{> > > NEXT SERMON > > >}
\newline
\newline
雅歌 1:13-1:14
\newline
\begin{longtable}{cl}
\hline
\hline
章節 & 經文 (和合本修訂版)\\
\hline
1:13 & \begin{tabularx}{0.7\textwidth}{X} 我的良人好像一袋沒藥，在我胸懷中。 \end{tabularx} \\ \\ \relax
1:14 & \begin{tabularx}{0.7\textwidth}{X} 我的良人好像一束鳳仙花，在隱‧基底的葡萄園中。 \end{tabularx} \\ \\
[1ex]
\hline
\hline
\end{longtable}
自亞當夏娃犯罪後,他們就赤身露體.他們犯罪之前,身上確有一件衣服；這衣服使他們能蔽體,不致露出羞恥來,這衣穿在身上,在使他們能與神交通,毫無懼怕!但犯罪後這衣服失落了,心中就產生了煩惱.為當他們被造時,有神的性情,有神的榮耀,犯罪後那榮耀就離開了,他們只有懼怕,羞恥,和戰慄!這榮耀在人的身上,是非常重要的,榮耀現已失去了,直等到基督再來時,祂必將這榮耀再歸到我們的身上.所以信主的人,並非僅僅作一個得救的基督徒.現要讓我們看看,基督的榮耀如何顯明,並如何停留在我們的身上.
\newline
\newline
(一)「王正坐席的時候,我的那噠香膏發出香味,我以我的良人為一沒藥,常在我懷中.」(歌一12-23.)
\newline
\newline
當哪噠膏香發出香氣時,對主的感覺和認識就不同了,這時會認識基督是一袋沒藥.在利未記一章記載:人獻過了燔祭後,在他裡面就發出香氣來.何謂燔祭?所謂燔祭就是當人犯了罪時,要經下列各種程序,然後罪才能得赦免.
\newline
\newline
要到神面前將祭物──牛,羊或鴿子獻上.並將手放在祭牲頭上,表示他所犯的罪,全歸於此祭牲,那麼罪才能得赦免.
\newline
\newline
當香氣發出時,書拉密女說:「我以我的良人為一袋沒藥.」「沒藥」是極為貴重的香料,這香料放在懷中,好好地保存.「沒藥」是用作膏屍體用的.人獻上燔祭後,他自己也被獻上作活祭,放出香氣來.保羅說:「弟兄們,我以神的慈悲勸你們,要將身體獻上當作活祭.」(羅十二1)基督徒認過罪後又再犯,雖然活在世界上,但表面那犯罪的性情和那愛好犯罪的心,還沒有塗上沒藥；那就是說要將自己放在主的面前,要一同經歷基督的死,與基督一同埋葬.故書拉密女以她的良人為一袋沒藥,放在懷中,這是說我讓基督的死在我裡面,我以基督的死為我的死.以基督為至寶,以祂的生命為我的生命；以祂的經歷為我的經歷,這是基督徒必需經歷的過程.
\newline
\newline
何謂基督的榮耀呢?主在(約17章)為門徒禱告的一段說:我已經得了榮耀,並將榮耀獻給父神.但在祂的祈禱中,祂求神再榮耀他.然後說:我的就是你的,你的也就是我的,我已經榮耀了你.那是甚麼意思呢?當基督要進入榮耀的門時,必須先經歷十字架,經歷死亡；然後將自己完全交給父神.人要進入榮耀,必須先經歷十架,才能說:「你的是我的,我的是你的.」可惜今天我們只會盼望神賜一切的福,但到奉獻的時候,卻只會說:「我的就是我的!」我們沒有榮耀歸給神,所以榮耀也不會臨到我們的身上了.
\newline
\newline
(約十二章記載)有幾個希利尼人對腓力說,我們願意見耶穌.主沒有答應立刻接見他們,反而說:「一粒麥子若不落在地上死了,就不能結出許多子粒來.」這話是指祂得榮耀而說的,如果我們要進入這榮耀的境地裡,就必需從這「死」的門進去.若不將基督的死當作我們的死,我們無法得著榮耀.所以書拉密女說:「我以我的良人為一袋沒藥,放在我的懷中.」各位!請問你今天有沒有將基督的死當作你的死呢?讓我們在主面前同一態度地向主說:「你既將一切都給了我,你為我捨命,主啊!今天我也將一切獻給你.」這一個祭在舊約利未記二章稱為素祭.就是說神給我許多的好處,現在要到神面前將神給的無論是物質,金錢,出產的獻給神作為素祭,也是馨香的祭.許多人誤會了羅馬書十二章一節所說的是作傳道的奉獻.其實那是指每一個人都應該如此奉獻.凡神兒女們都應奉獻,主既為我們死,我們從祂領受了恩典；就沒有理由再為自己活著了.故此應該把自己獻上,當作活祭,如此事奉,才能討神喜悅；所獻上的是火祭,即變成馨香的祭.
\newline
\newline
四福音都記載馬利亞用香膏膏主,當主耶穌坐席時,她把一玉瓶真哪噠香膏打破,將香膏傾倒在耶穌身上,香氣就充滿了那房子.這香氣並且充滿了全世界.因為耶穌說:無論你們在什麼地方傳福音,都要述說這女人所作的事,作個紀念.因為她在我身上所作的,是一件美事,是為我安葬預備的,這香氣就香遍了全世界了.所以我們裡面也應該有沒藥,塗抹我們的裡面,使我們經歷基督死的生活──獻在祭壇上.若我們能如此,就能對主說:你的是我的,我的也是你的了.我們到教會事奉,是否推辭很忙呢?是否常憂慮,失眠,吃喝失常呢?就是因為沒有沒藥的塗抹；把得失看得太重,以至心中煩惱不安......如果未經基督的死,仍然為自己而活,沒有沒藥在懷中,才會有這些問題的發生.人要進入榮耀,就必需先從死進入.
\newline
\newline
(二)「我以我的良人為一棵鳳仙花,在隱基底葡萄園中.」(歌一14)
\newline
\newline
鳳仙花是一棵的,不是一朵或一枝的,是整棵的.隱基底是在死海的旁邊,離耶利哥往南不遠,死海水是苦的,不能喝的,但隱基底卻有泉源.(大衛曾在此處躲避掃羅).隱基底是死海旁唯一肥美之地,是世外桃園.在那裡有葡萄園,有羊群,並各種植物出產.這表示主有復活的生命,豐富,榮耀的生命,這生命在我們的主身上.主是一袋沒藥,主是一棵鳳仙花,我們藉著主可進入那豐盛生命的境地.人若不經過死,就不會經歷復活.在羅馬書六章可清楚地知道,人若與主一同死,才可與祂一同復活,得到神的榮耀.我們在亞當失去的,現在可從基督裡得回來.請問你是否有豐盛的生命?榮耀是否在你裡面?要脫去容易纏累我們的罪.主的榮耀才能在我們身上彰顯.保羅說:凡稱呼主名的人,總要離開不義.(提後二19)要自潔脫離卑賤的事,才能作貴重的器皿.讓主的生命在我們這器皿裡彰顯他的豐富吧!那麼基督才能因我們得榮耀.
\newline
\newline
「我是沙崙的玫瑰花」(歌二1)
\newline
\newline
以前所講的是新婦書拉密女的感覺,現在是良人說的話:「我是沙崙的玫瑰花.」沙崙是個平原,靠近地中海不遠；此平原以前是迦得支派所得的地方,在迦密山旁.從約帕到迦密只有這平原,現在到這地方,是找不到玫瑰花的,但在古時的確是有的.我曾到過此地,雖然沒有玫瑰花,但再往深處仍可見到其他的花朵.沙崙的玫瑰放香,與別處不同；因為這地方是平原,也是戰爭之地.在以往的歷史,曾發生過多次戰爭,每次浩劫,必成為一片荒涼.但在這裡卻題到沙崙的玫瑰,即在戰爭被蹂躪之後,卻能開放美艷的花朵.這花朵會顯得特別的寶貴.照樣在基督徒的屬靈生命和生活中,也是個戰爭,這場戰爭中,我們如何作戰,如何可得勝,神就如何在我們身上得榮耀.基督來世,祂的一生,就是個屬靈的戰爭.這是不斷的戰爭,一生下來,魔鬼就與他打仗,人也與他打仗；魔鬼一直在祂身上攻擊祂.他在客西馬尼園,憂傷得幾乎至死,有天使加祂力量,這是一場屬靈的戰爭.祂上十字架時,是最後的戰爭.我們見到主被聖靈充滿,從約但河上來,聖靈將祂帶到曠野去,在那裡祂經歷了屬靈的爭戰.我們的一生,在不斷戰爭之下,正如在沙崙的平原中,作爭奪戰中的人.像約伯,他雖經歷痛苦,卻不知道這是一場屬靈的戰爭,是魔鬼向神的挑戰,是神著魔鬼去試煉約伯的.弟兄姊妹!我們屬主的人,當我們遭遇苦難時,不要失望,也不要灰心,因為這在靈界裡,可能是有一場戰爭.要經歷這場戰爭,玫瑰纔能發出香氣來.基督的香氣纔能在我們身上彰顯.
\newline
\newline
「我是谷中的百合花」.(歌二1)
\newline
\newline
百合花預表聖潔和潔白.當彼得寫書信時,他年已老；他說我當年在聖山時,我看見有我的主身體改變,面貌放光,身上的光如潔白的衣服一樣.馬太十七章及啟示錄一章也這樣的記載.什麼是基督的榮耀呢?乃是聖潔.我們若有聖潔的生活,就有神的榮耀在我們裡面.當祭壇下的靈魂喊著說:為什麼還不替我們伸冤?然後有聲音回答說:他們再等待片時,就有白衣賜給他們.白衣就是有基督的榮耀,基督的榮耀加在這些殉道者的身上,這是我們的盼望.所以我們要忘記背後,努力面前,向著標竿直跑.主是我們谷中的百合花,祂在艱難困苦中,為我們留下了榜樣.祂有聖潔的樣式,好叫跟從祂的人,得穿白衣有份於基督的榮耀裡.
\newline
\newline


\newpage

\section{第46屆~港九培靈會~于力工~第四講}
\label{sec:1201}
\textbf{基督是我們的主}
\newline
\newline
連結: \href{https://www.hkbibleconference.org/session-message/view/1201}{\texttt{https://www.hkbibleconference.org/session-message/view/1201}}
\newline
\newline
\hyperref[sec:1199]{< < < PREV SERMON < < <}
~
\hyperref[sec:toc]{[返目錄]}
~
\hyperref[sec:1203]{> > > NEXT SERMON > > >}
\newline
\newline
雅歌 5:10-5:16
\newline
\begin{longtable}{cl}
\hline
\hline
章節 & 經文 (和合本修訂版)\\
\hline
5:10 & \begin{tabularx}{0.7\textwidth}{X} 我的良人紅潤發亮，超乎萬人之上。 \end{tabularx} \\ \\ \relax
5:11 & \begin{tabularx}{0.7\textwidth}{X} 他的頭像千足的純金，他的髮綹卷曲，黑如烏鴉。 \end{tabularx} \\ \\ \relax
5:12 & \begin{tabularx}{0.7\textwidth}{X} 他的眼如溪水旁的鴿子，沐浴在奶中，安得合式。 \end{tabularx} \\ \\ \relax
5:13 & \begin{tabularx}{0.7\textwidth}{X} 他的兩頰如香花園，如香草臺；他的嘴唇像百合花，滴下沒藥汁。 \end{tabularx} \\ \\ \relax
5:14 & \begin{tabularx}{0.7\textwidth}{X} 他的雙手宛如金條，鑲嵌水蒼玉；他的身體如同雕刻的象牙，周圍鑲嵌藍寶石。 \end{tabularx} \\ \\ \relax
5:15 & \begin{tabularx}{0.7\textwidth}{X} 他的腿好比白玉石柱，安在精金座上；他的容貌如黎巴嫩，佳美如香柏樹。 \end{tabularx} \\ \\ \relax
5:16 & \begin{tabularx}{0.7\textwidth}{X} 他的口甘甜，他全然可愛。耶路撒冷的女子啊，這是我的良人，這是我的朋友。 \end{tabularx} \\ \\
[1ex]
\hline
\hline
\end{longtable}
假如我們要了解雅歌書,就必須從下列這幾個角來觀看:
\newline
\newline
(一)從三福音書看──馬太福音可見到基督是我們的王.馬可福音可見基督,如何在世上作僕人.路加福音則見到基督,在世上如何與人相處.
\newline
\newline
當他聽到聲音時,屬靈的眼睛開了.他所看見的基督,是與他蒙召時所看見的,和與他想像中的基督,是完全不同的.因他這時所見的是榮耀的基督,叫他不能不仆倒下來,軟弱下來；因他現在見到的是創造的主,是教會的頭,是掌握一切的神.祂要把一個新的世代,帶到地上來.這樣一位榮耀的主,使約翰的肉體倒下來,他的靈卻是清醒了.感謝主!我們何時見到榮耀的主,我們肉體就會軟弱,身體就仆倒在祂腳前；過去的觀念和經歷,甚至對過去基督的看法,都一齊放在主腳前.求主今天將祂榮耀的形像,向我們顯明.
\newline
\newline
基督是教會的頭,祂是在七燈台中行走的人子.約翰在拔摩海島得啟示,論到在異象中所見的基督──是面貌發光,眼目如同火焰的.祂憑手裡所拿著的七個星,和七金燈台的說話.這可見教會整個的光景,及神的兒女在神面前的正常表現.
\newline
\newline
(一)「我的良人白而且紅,超乎萬人之上.」(歌五10)
\newline
\newline
這句話使我們想到腓立比書二章六節中說:「他本有神的形像,不以自己與神同等為強奪的.」這「強奪」意思是「戀愛不捨」之意.祂與父是從亙古就在一起,一同創造這個世界,也設立了救恩的計劃.主耶穌曾說:我與父原為一.但祂不因此,而留在天上,反而到世上來；取了奴僕的形像,存心順服,以至於死,且死在十字架上.然後從死裡復活,神把祂升到至高之處,萬膝要向祂跪拜,萬口要稱頌祂.(歌五10)說「祂又白又紅」.這話不是單指健康而言,乃是指在祂裡面榮耀的表現；是超乎一切之上的,眾口都要頌揚祂,祂是榮耀的主!
\newline
\newline
(二)「他的頭像至精的金子,他的頭髮厚密纍垂,黑如烏鴉.」(歌五11)
\newline
\newline
按啟示錄一章說主的頭髮是白色的,而在(歌五12)卻說是黑色的,到底有沒有衝突呢?實在是沒有的,因為前者是榮耀的基督,後者是人性的基督.祂是人也是神,祂是教會的頭,是我們的主,在我們裡面,應佔有首位.查尼布甲尼撒王所夢見的異像,這像是用金屬造的,金的頭及各金屬造的身.後來他也照著夢中所見的大像,用精金仿照造成一像.他要改變神的計劃和次序,改變神所要用的東西,那像就成為偶像.按他所見的像,頭是用金子造的,表示有力量,有最大的榮耀.基督在教會的地位也必如此.金子乃代表神的性情,要在教會居首位,經將神的性情灌輸在教會裡面.但今日教會光景如何呢?難怪主在啟示錄中,要寫信給七教會.這七教會,各有各的問題,都失去了他們所應有的東西,神的性情就無法在教會裡面彰顯出來.末世的教會還能彰顯神的性情嗎?神的兒女們,有多少神的性情表現呢?在第一天的信息中,我已說過,神按著祂自己的形像造人,那形像是指神寶貴的性情,是比金子還美的性情,可惜人失去了標準!基督來世,祂要藉著死,將失去的性情帶回到我們裡面.我們在教會裡應該把這性情表現,讓主住在我們中間,充充滿滿的,有恩典,有真理；神一切的豐盛,都有形有體的居住在祂裡面.祂比宇宙更大,非我們所能了解的.感謝主!祂的性情要在我們裡面,看哪!我們的地位是何等的可愛,和何等的可羡慕!神的性情可用金子代表,這性情要流通在教會裡,也要流通在我們裡面.
\newline
\newline
(三)「他的眼如溪水旁鴿子的眼,用奶洗淨,安得合式」.(歌五12)
\newline
\newline
啟示錄中說,主的眼如火焰一般,是明亮的；雅歌書卻說祂眼,如同溪水旁的鴿子的眼.這兩處經文原是代表兩件事,也是一件事的兩方面.「眼如溪水旁鴿子的眼,是被洗淨的,代表什麼?」鴿子在聖經中代表聖靈,聖靈到我們中間,有兩種的表現:第一種是如鴿子一樣.當主受洗後,從水上來,天為祂開了,於是聖靈彷彿像鴿子降在祂身上.當聖靈澆灌在祂身上時,聖靈就把祂帶到曠野去.主順著那靈的引導,毫無意見.按人的看法,主來到世上已有三十年,祂可以對聖靈說:我該去耶路撒冷傳福音,告訴萬民悔改的道理.但祂並不如此,祂肯順服,因住在祂裡面的靈管理著祂.祂順服這管理,順服那住在祂的裡面的聖靈澆灌祂時,就像鴿子,祂就馴良像鴿子一般安安靜靜.照聖靈的引導,到曠野去.藉著溫柔,順服,倚靠,及完全的交托,去面對那屬靈的戰爭!這給我們看見,今日教會及神兒女們美麗的地方.我們讀雅歌書,可以明白基督對教會的關係.新婦也需要有鴿子的眼睛,要有屬靈的看見.要順服住在裡面的聖靈,這是基督留下給我們的模樣.主升天後,祂就差遣聖靈降臨,澆灌在門徒身上.聖靈像火焰的舌頭,分別落在各人的頭上,這是甚麼意思呢?聖靈來好像火一般,做潔淨,折毀,焚燒的工作.然後將能力帶來,他們的口就打開了,信息就被傳出去.現在是聖靈掌權的時代,教會一切的事,都在聖靈掌握中.神願望今日的教會,能順服聖靈,讓聖靈可掌權.試問我們的教會,在制度上,在傳統上,在人為的組織上,是誰在掌握呢?啟示錄說:「那眼目如同火焰的向眾教會說話.」我們的主坐在寶座上,要看我們的教會如何?假如有聖靈在我們教會中掌權,那麼我們應該馴良像鴿子,有鴿子的眼光看一切的事,順服,安靜,等候,謙卑和溫柔,在我們的生活中表現出來.所以聖靈澆灌我們,一方面是工作,一方面亦為生活.聖靈住在裡面,作內在的管理,對我們生活是重要的.
\newline
\newline
(四)「祂的兩腮如香花畦,如香草台.」(歌五13)
\newline
\newline
這裡說到基督的臉面.在啟示錄卻說,祂的面貌如同烈日放光.約翰見到異像時,正如彼得後來寫書信時所說的話:我們同祂在聖山的時候,曾親眼見過他的威榮.(參彼後一16-18)衣服潔白,面貌放光,那是屬天的美麗如香花畦,如香草台.請注意,在啟示錄中說祂的面貌卻如烈日放光.當摩西登西乃山,四十晝夜,下山的時候,他的面貌也發光,甚至以色列人不能見他的面,他要用帕子蒙上面.人見了神的光害怕?因為良心上有問題.良心被光照,責備,就懼怕!難怪亞當夏娃聽見神聲音時,就躲在樹叢中戰兢懼怕!在啟示錄二至三章,可見到教會已失去了立場,異端侵入向世界妥協；跟著世界走,學世界的樣式.當主光照時,教會應有覺悟,教會要悔改!信徒也要悔改!如果你怕這光,良心就在黑暗中.仍死在罪惡過犯中,良心已昏暗了(弗二1-2)求主使我們看完祂美麗的面孔,心中有喜樂,有安慰,有滋潤；那麼我們的良心才能坦然無懼.
\newline
\newline
(五)「他的嘴唇像百合花,且滴下沒藥汁.」(歌五13)
\newline
\newline
主的口能發出恩言.祂的言語像蜂房下滴的蜜一般,祂的話是我們腳前的明燈；祂的話安定在天,今日要藏在我們心中,可使我們不致得罪祂.祂的話像百合花,是在患難中長出來的,是潔淨的.是經過練淨的,是帶著能力的.祂不像文士和法利賽人,祂說話像有權柄的人.(可一22)教會有責任,要把主的話傳出來讓人聽見受感動.我們的口像百合花般,我們要作得勝的見證,說話要帶著能力,是美麗的,令人感到甘甜,滿足.一句話說得合宜,就像金蘋果落在銀網子裡.(箴廿五11),主的口滴下沒藥汁,因祂的口是曾經有過死亡的口,是被沒藥塗抹過的,滴出沒藥汁來.經過死亡,經過鍛煉的口,然後話語才有力量.故凡稱呼主名的人,總要離開不義.(提後二19)要說造就人的話,叫聽見的人得益處(弗四29)不合聖徒體統的話,連一句也不該說.我們的話語要滴下藥汁,要經過洗淨.成為清潔的言語.
\newline
\newline
這經文表示權柄.在啟示錄一至五章告訴我們,主右手中拿著七星,在七個金燈台中間行走；表示一切的奧秘和權柄,都在祂手中.當約翰在異像中看見七印封嚴的書卷時,就見印中定有非常奇妙的事.於是他說:誰能展開這書卷呢?他就大哭起來,長老中有一位說,唯有羔羊配開這書卷,因為一切的奧秘,權柄都在祂手中.神的一切計劃也在祂裡面,祂的手像金管,鑲嵌如水蒼玉；這表示祂的美麗,和神的一切豐盛.感謝主!教會在祂手中,教會的使者也在祂手中,神的兒女也在祂手中.詩篇告訴我們,祂用祂手中的巧妙,引導了他們,用祂心中的純正牧養他們.(詩七八72)祂的手是巧妙的手,恩惠的手；當祂的手被舉起來時,福氣就臨到我們.教會應有這樣的手,當舉起來時人就蒙福.當手舉起時,就有了影子,這影子就像沙漠中的涼亭,在乾旱之地,可讓人安歇.教會應像蔭涼之所,出時,可使人找到蔭涼之所,得到安息.我們都該彼此用愛心說誠實話,要彼此認罪,互相代求,然後才能使人感到教會是個涼亭,如此主才得榮耀.
\newline
\newline
(六)「他的身體如同雕刻的象牙,週圍鑲嵌藍寶石.」(歌五14)
\newline
\newline
這是代表基督曾經歷痛苦,經過雕刻,為我們受鞭傷,為的是要我們得平安.為我們經歷痛苦,是要使我們的罪得赦免,祂被人鞭打,戴上荊棘冠冕,被吐唾沫,經過許多的痛苦!如同經過雕刻的象牙一樣.在啟示錄四,五,六章中.當我們看見天上那幅圖畫時,寶座出現在人中間,有廿四位長老和眾天使,眾聖徒在唱歌.那寶座的面前有寶石,像天一般的藍色的寶石,那是表示敬拜的真實；教會也應將敬拜的真義,實際帶出來.這一位榮耀的基督,必需有榮耀的身體,就是今日的教會.祂要把教會帶到榮耀的境地,使每一位神的兒女,都要經歷這榮耀.
\newline
\newline


\newpage

\section{第46屆~港九培靈會~于力工~第五講}
\label{sec:1203}
\textbf{基督是我們的欣賞}
\newline
\newline
連結: \href{https://www.hkbibleconference.org/session-message/view/1203}{\texttt{https://www.hkbibleconference.org/session-message/view/1203}}
\newline
\newline
\hyperref[sec:1201]{< < < PREV SERMON < < <}
~
\hyperref[sec:toc]{[返目錄]}
~
\hyperref[sec:1204]{> > > NEXT SERMON > > >}
\newline
\newline
雅歌 4:7-4:15
\newline
\begin{longtable}{cl}
\hline
\hline
章節 & 經文 (和合本修訂版)\\
\hline
4:7 & \begin{tabularx}{0.7\textwidth}{X} 我的佳偶，你全然美麗，毫無瑕疵！ \end{tabularx} \\ \\ \relax
4:8 & \begin{tabularx}{0.7\textwidth}{X} 我的新娘，請你與我一同離開黎巴嫩，與我一同離開黎巴嫩。從亞瑪拿山巔，從示尼珥，就是黑門山頂，從獅子的洞，從豹子的山往下觀看。 \end{tabularx} \\ \\ \relax
4:9 & \begin{tabularx}{0.7\textwidth}{X} 我的妹子，我的新娘，你奪了我的心。你明眸一瞥，你頸項的鏈子，奪了我的心！ \end{tabularx} \\ \\ \relax
4:10 & \begin{tabularx}{0.7\textwidth}{X} 我的妹子，我的新娘，你的愛情何其美！你的愛情比酒甜美！你膏油的馨香勝過一切香料！ \end{tabularx} \\ \\ \relax
4:11 & \begin{tabularx}{0.7\textwidth}{X} 我的新娘，你的唇滴下蜂蜜，你的舌下有蜜，有奶。你衣服的香氣宛如黎巴嫩的芬芳。 \end{tabularx} \\ \\ \relax
4:12 & \begin{tabularx}{0.7\textwidth}{X} 我的妹子，我的新娘是上鎖的園子，是禁閉的園子，是封閉的泉源。 \end{tabularx} \\ \\ \relax
4:13 & \begin{tabularx}{0.7\textwidth}{X} 你園內所種的結了石榴，有佳美的果子，並鳳仙花與哪噠樹。 \end{tabularx} \\ \\ \relax
4:14 & \begin{tabularx}{0.7\textwidth}{X} 有哪噠和番紅花，香菖蒲和桂樹，並各樣乳香木、沒藥、沉香，與一切上等的香料。 \end{tabularx} \\ \\ \relax
4:15 & \begin{tabularx}{0.7\textwidth}{X} 你是園中的泉，活水的井，是從黎巴嫩湧流而下的溪水。 \end{tabularx} \\ \\
[1ex]
\hline
\hline
\end{longtable}
今天我們要從基督的眼光,來看教會和神的兒女們.從第一段經文中,論到新郎對新婦的認識,和稱讚.有些字句原來是新婦在新郎身上用的,現在反倒用在新婦身上.這可證明基督的一切豐滿和美德,可以在新婦身上彰顯出來.今天我講的題目,是基督是我們的欣賞.我們欣賞什麼呢?欣賞主的完全,因為只有祂才是真正的標準.我們是無法可以完全的.更無法可以在神面前說我已經完全了.
\newline
\newline
在(創十七1)說,亞伯拉罕年九十九歲時,耶和華向他顯現.對他說:「我是全能的神,你在我面前應當作完全人.」當神發出此要求之前,祂先說:我是完全的神,亞伯拉罕阿!你靠你自己,照你的標準,是無法作完人的.你該記得,我是全能的神,既然我是無所不能,就能使他這軟弱可憐,自己無善可陳的人得以完全.
\newline
\newline
(一)全然美麗
\newline
\newline
「我的佳偶,他全然美麗,毫無瑕疵.」(歌四7)感謝主!在良人的眼光中,按祂的標準,「你全然美麗,毫無瑕疵.」讚美主!祂判斷我們不是根據人的標準,乃是按祂的標準.這標準並不太高或太低,乃是照著神的心意而定.感謝主!祂稱讚教會全然美麗,也對我們說:他全然美麗,主欣賞我們便感到滿足,雖然我們不配,我們是他重價所買來的,是祂寶血洗淨的,是祂的道將我們潔淨.所以主會欣賞我們.主的欣賞一則是鼓勵我們,二則是儆戒我們.我們不能徒受主恩就放肆,或停留不動.必須忘記背後,努力面前,向著標竿直跑；好叫基督能得著我,我也能得著基督,這是我們今天應有的心願.(在以賽亞與以西結書均載)當主造基路伯時,祂說:「你是全然美麗.」這一位就是在神面前,侍立的天使長,神吩咐摩西造會幕,安置約櫃於至聖所,在至聖所裡的約櫃,是用皂莢木造的,一肘半高,二肘半長,外包精金.在約櫃上面,要用精金錘出一個施恩座,在施恩座旁邊,要用精金造兩個基路伯；是事奉神的天使,單侍立在施恩座前,用兩翅膀遮掩著臉,基路伯是對面的站著,神的聲音就從這裡發出.基路伯所佔的地位,是何其重要.人要到神面前,先要經會幕的外院,在那兒獻祭；然後帶著血經過洗濯盤,再進到聖所而到至聖所.既然基路伯佔有重要的地位,特別是受膏的基路伯,在神面前全然美麗.他有神寶座上一切的條件,他身上有一切寶石的樣式.讀啟示錄時可知那位坐寶座的,有藍寶石,紅寶石,和瑪瑙的樣子,但卻沒有提及樣子如何?當神造這天使長時,在他身上有一切寶石的樣子,他有神坐寶座的條件,才可以站立在神寶座前.教會是基督的新婦,應有寶座的條件,應披戴基督,行走在光明中.全然美麗,這句話,當年曾落在受膏的基路伯身上,但這基路伯在神的面前敗壞了!他心中起了驕傲,他說我要升到高天之上,要坐在寶座上要和神一樣,要像神一樣的發號施令.他墜落,成為魔鬼.這全然美麗的基路伯,竟然墜落與神為敵.這一件事應為我們的儆戒,主雖稱讚我們全然美麗,毫無瑕疵；我們卻要把這話奉獻給主,惟有主配得.若非主全然美麗,我們怎會有今天的美!若非主毫無瑕疵,我們怎能毫無瑕疵!主對我們的稱讚,一方面是鼓勵,一方面也是儆戒.
\newline
\newline
伊甸園是最好最美的地方,有神的同在,有神的賜福；但人卻可以從此墜落下去!我們不可不小心,人可在恩典中,思想自己的問題,以自己為中心.我們可能受試探,求主特別的保守我們；因為祂的稱讚表示祂有保守的能力,使我們永遠停留在於此.故主對我們的欣賞,要好好地注意；不可使欣賞成為虛空,求主恩待我們.
\newline
\newline
(二)與主同工同生活
\newline
\newline
「我的新婦,求你與我離開利巴嫩,與我一同離開利巴嫩,從亞嗎哪頂,從示尼珥與黑門頂.」(歌四8)
\newline
\newline
主對新婦有一要求,按利巴嫩,亞嗎哪,示尼珥和黑門,都是在加利利北部的高山的山峰.山上終年積雪,風景極優美,加利利的湖水,就是這從山頂上的雪溶了流下來的.祂說讓我們離開這些地方,在山頂上可享受一切,看到風景.在黑門山不必到山頂,只需往西邊看就可看見地中海.當年兩門徒離開耶路撒冷在以馬忤斯路上,遠遠就可見到約帕和地中海,並見北部的尼基多平原和黑門山.祂說我們走吧!別留戀在此地,讓我們回到了耶路撒冷去.那是作王的地方,是大君的城邑,有寶座,是放光的地方,是事奉神的地方.是神要我們看見的地方.照樣,我們的歸宿乃是耶路撒冷.這耶路撒冷是屬靈的耶路撒冷,有一天要從天上降下；這耶路撒冷也代表今日教會的實際.教會今天在地上,乃是天國的代表.你他是神的兒女,聚合在一起,因為在我們心中有天國.耶穌說:「天國在你心裡面.」你我就是這天國,我們應記著,那屬靈的表現.要從天上帶到地上來.今日的教會,有一天要將天上的耶路撒冷帶到地上；我們要永遠在耶路撒冷裡與基督一同作王.來吧!我們走去,向著耶路撒冷.按所羅門王在位時,如果是剪羊毛的時候,他也會下到剪羊毛的地方,人就特別為他預備帳棚,以今日的話來說,可稱之為行宮.我們在教會裡,是與基督一同作工,等到工作完畢後,就回到天上的耶路撒冷去.這是我們的盼望.主說,起來吧!我們向著這方向走去.不要留戀在我們的傳統裡,制度裡,個人的喜好裡!讓基督的喜好,成為我們的喜好,所以新郎對新婦說:「我們走吧!」讓我們一起過正常的生活.
\newline
\newline
(三)奪了我的心
\newline
\newline
「你奪了我的心,你用眼一看,用你項上的一條金鍊奪了我的心.」(歌四9)
\newline
\newline
這應該是新婦對新郎說的話,主實在是奪了我的心,有人曾用一名詞說:主的恩典臨到我們身上,就好像是不可抗拒的恩典,想推拒都不可能.主為什麼會愛我這個人!我有什麼長處!我比人有什麼強?然而主卻挑選了我,愛了我.當我們想到這裡的時候,真要深深的感謝祂,因祂的恩典是不可抗拒的!在此段經文說:「當你用眼一看時,你就奪了我的心；」這是主所說的話,這使主感到教會向祂的那種愛,真能令祂感到教會向祂的那種愛,真能令祂感到心中滿足.祂以我們的默念為甘甜,以我們獻上的脂油為滿足.在達爾比的「基督的滿足」一書中,可以見到基督如何滿足我們的心.在禱告時,感到主聽禱告的滋味甘甜,與主同在時,默想時,安靜等候時,心中會有說不出來的快樂.主今天思念我們時,祂坐在寶座上發出了微笑,我們也把主的心奪了.
\newline
\newline
金鍊是頸項上的鍊,按傳道書十二章上所說,金鍊代表頸,是表示一個人的意志.我們的意志要完全放在主的手中.保羅說:你們要立志作安靜人.當以斯帖帶那些被擄的人歸回時,他們就定意考查耶和華的律法.我們裡面應有一個定意,有一個決心,要把一切的意志放在祂身上,當這鍊子掛在我們頸上時,我的意志經過聖靈的火所潔淨.主的性情在我們裡面,就沒有自己的道路,須專一的跟隨主.惟有主能奪去我們的心,現在卻看見,當人跟隨主到一個地步時,與主愈接近時；奇妙的,就會將一切奪了過來放在他裡面.摩西在西乃山四十晝夜,與神面對面生活；當他下山時,面上有神的榮光發出來.
\newline
\newline
(四)要看我們的愛情
\newline
\newline
「你的愛情何其美,你的愛情比酒更美,你膏油的香氣勝過一切香品.」(歌四10)
\newline
\newline
新郎對新婦說:「你的愛情比酒更美.」這情不是在情感上的愛,當祂賜福你時,你心中就歡喜.當祂引導你時,你就快樂.當祂四面護衛你時,你會感到有安全.正如以色列人在曠野,有雲柱火柱作引導,飢餓時有嗎哪充飢,乾渴時有磐石出水,四十年之久,衣服沒有破,鞋子沒有穿爛.神賜這麼多的福份與恩典,有這麼多的福份與恩典,為甚麼他們不停留在曠野呢?神的應許地是迦南,他們必須進去得為業.兩個半支派的人,停留在約但河東岸,不願入迦南,已先後被塗抹了.神的恩典臨到我們,是為我們今天的需要.今日的飲食今日賜給我們,這不能當作是屬靈的標準.我們不能永遠留在培靈會裡,更不能停留在那十年前被聖靈充滿的光景裡.研究聖經的人,懂得勝靈工作方向的人,都明白恩賜乃是造就了我們,但我們必須要繼續往前走!不要把神賜我們的一點點福份,作為標準,是滿足於現況之下,我們愛主的心,要勝似酒.酒使人興奮刺激,但真正愛主的心,要脫離酒的光景；要進入真正愛的表面.「你膏油的香氣勝過一切的香氣.」請記著,死蒼蠅,使作香的膏油發出臭氣,故此不要那死蒼蠅敗壞了我們的香氣阿!
\newline
\newline
是要欣想諸般的香氣.恩賜是為我們放出香氣而用的,有難處時,就正是你放香氣的時候阿!
\newline
\newline
(五)「滴蜜,好像蜂房滴蜜,你的舌下有蜜有奶.」(歌四11)
\newline
\newline
現在基督向教會新婦說,你的嘴滴下蜜來,言語要甘甜,並經過鍛煉；發出來時,使人如聽見主的聲音一般.
\newline
\newline
宋尚節博士病重時,曾發出一封公開的信,我也有機會看過.信中給我有一感覺,覺得每一句話如珍珠.讀時另有滋味在其中,因為每句話都是由幾節聖經湊成的.有如下滴的蜜,流出奶.使你讀時,感到有供應,感到甘甜.
\newline
\newline
一個肯禱告親近主,他在主面前的光景,也必如此.我們要藉著別人的禱告,得到了鼓勵.要好好的下工夫與主親近,自然你也能滴蜜,流奶了.
\newline
\newline
(六)活水的泉源
\newline
\newline
「是關鎖的園,禁閉的井,封閉的泉源.」(歌四12)
\newline
\newline
在這園子裡有泉水,有果子,那是甚麼意思?乃是要叫主滿足.主要看在我們裡面能否流出活水的江河來?(約七35)主耶穌是那一個泉源,使園子結果子,滋潤別人.我們的泉源,是屬乎主的；我們的井,也是屬乎主的.把園奉獻單為祂,祂纔滿足.
\newline
\newline


\newpage

\section{第46屆~港九培靈會~于力工~第六講}
\label{sec:1204}
\textbf{基督是我們的旌旗}
\newline
\newline
連結: \href{https://www.hkbibleconference.org/session-message/view/1204}{\texttt{https://www.hkbibleconference.org/session-message/view/1204}}
\newline
\newline
\hyperref[sec:1203]{< < < PREV SERMON < < <}
~
\hyperref[sec:toc]{[返目錄]}
~
\hyperref[sec:1286]{> > > NEXT SERMON > > >}
\newline
\newline
雅歌 2:2-2:6
\newline
\begin{longtable}{cl}
\hline
\hline
章節 & 經文 (和合本修訂版)\\
\hline
2:2 & \begin{tabularx}{0.7\textwidth}{X} 我的佳偶在女子中，好像荊棘裡的百合花。 \end{tabularx} \\ \\ \relax
2:3 & \begin{tabularx}{0.7\textwidth}{X} 我的良人在男子中，如同蘋果樹在樹林裡。我歡歡喜喜坐在他的蔭下，嘗他果子的滋味，覺得甘甜。 \end{tabularx} \\ \\ \relax
2:4 & \begin{tabularx}{0.7\textwidth}{X} 他領我進入宴會廳，為我插上愛的旗幟。 \end{tabularx} \\ \\ \relax
2:5 & \begin{tabularx}{0.7\textwidth}{X} 請你們用葡萄餅增補我力，以蘋果暢快我的心，因我為愛而生病。 \end{tabularx} \\ \\ \relax
2:6 & \begin{tabularx}{0.7\textwidth}{X} 他的左手在我頭下，他的右手將我環抱。 \end{tabularx} \\ \\
[1ex]
\hline
\hline
\end{longtable}
在本書二章告訴我們,如何安息在主裡,新婦形容她愛人如園中的蘋果樹.在這樹下,她可以歡歡喜喜地,接受蘋果的滋味,和其中的甘甜.但想到始祖犯罪後,他們躲在林中,聽見了神的聲音,心中就懼怕.他們的眼不敢看神,人犯罪後,就失去了那蘋果樹下的安息.由生命樹下,而進入分別善惡的樹下.他們沉重地背負重擔,壓力甚重,這是人犯罪後的光景.耶穌說:「凡勞苦擔重擔的人,可以到我這裡來,我就使你們得安息.」人犯了罪,就背負重擔.
\newline
\newline
(一)律法的重擔──他們戰戰兢兢地守律法,怕守不住律法而不能得救,結果卻為神所撇棄.為什麼呢?因他們雖然行律法,嘴唇雖親近神,而心卻遠離神；所以神亦遠離他們.在他們身上有律法的擔子,因為守不住律法,所以緊緊抓住律法的條例.到了一個地步,這些條例就成為他們的重擔.他們真是苦!誰能救他們脫離這律法的重擔呢?他們在律法之下沒有自由,因為他們在禁果樹下過日子,所以沒有平安.
\newline
\newline
當初神在西乃山頒佈律法時,律法不是為他們得救而設的.按保羅說:他們過紅海,已經受了洗.(林前十1-2)他們是已被救贖的子民,已經得救,脫離了法老的管轄.亦已從埃及出來了,到了西乃山下,神纔給他們律法.不是要他們靠律法得救,乃是為幫助他們,要他們在主面前有享受,能過更完全的生活正如我們說「禮多人不怪」.禮貌是要教我們有一個好的,和正常的生活關係.神將律法頒下,也是要我們有正常的生活.當初亞當夏娃在伊甸園時,神對他們說:分別善惡樹上的果子不可吃,目的是要他們有正常生活；可是他們忽略了,而被引誘以至犯罪.人進入罪中時,就失去了平安.
\newline
\newline
(二)背負物質的重擔──人的心是沒有滿足的.神所應許給以色列人的迦南地,是流奶與蜜之地.但今天到聖地觀光,只見到處都是石頭.那麼流奶與蜜之地在何處?須知道以色列人犯了罪,地土就變為荒土.難怪主耶穌受試探時,第一個試探就是:「你若是神的兒子,就將這些石頭變為餅吧!」當主耶穌以五餅二魚使五千人吃飽時,那些人以為麵包問題解決了,從今以後不必再為吃的問題擔心了；於是他們都想擁護耶穌作王,他們就有了物質上的重擔.他們未嚐過生命的滋味,和蘋果的滋味,因此心中沒滿足也沒甘甜.當主耶穌看出他們心中的動機時,就遠遠地離開了.因主來不是要滿足人肉體之需要,而是滿足他們生命的需要.
\newline
\newline
(三)政治上的重擔──他們被羅馬所管治,失去自由；他們以為有耶穌基督作王,就可將仇敵滅絕.但主卻不是如此,而且說:「凡勞苦擔重擔的人,可以到我這裡來,就可以得安息.」所以我們可知真正的安息,是要從基督那裡領受.
\newline
\newline
(四)心靈裡的重擔──亞當夏娃犯罪時,就躲避不敢見神的面,為什麼?因在他們心靈裡有重擔.許多人心靈上有重擔,使他們煩亂.正如(詩四二5)說:「我的心啊!你為何憂悶?」這煩躁就是心裡的重擔,人被這重擔壓傷,他們就沒有安息.
\newline
\newline
在(歌二4)說:「他帶我入筵宴所,以愛為旗在我以上.」若有人要解決心裡的問題,心裡所有的重擔,就要到這旗幟之下.這是愛的旗幟,惟有愛才能滿足我們心裡一切的需要.許多人的問題,是因童年缺少了母愛和人的照顧,以至心裡無法滿足.神在舊約及新約中,都曾說及母親的地位,是何等崇高和重要.並且說兒女是耶和華的產業,所以母親的要愛自己的兒女.西諺有說:「搖搖籃的手,掌治天下.」一個作母親的,對孩子影響是大的,若能用愛去關心和注意,用正常的方法教導和帶領彼此間就可以得到滿足.嬰孩與母愛是分不開的,他需要母親的餵養與懷抱,他需要母親的溫暖.及至長大時,在他心中,仍覺得母親常在他身旁.基督的愛,也要這樣常在我們裡面,成為我們的力量,成為我們安息的泉源.只要我們向祂呼求,主的同在就向我們顯明.在主愛中生活,才能享到安息.愛的旗幟在我們上,我們就得到滿足,並享受到甘甜.
\newline
\newline
(歌六4):「我的佳偶阿!你美麗如得撒,秀麗如耶路撒冷.」得撒是個美麗的地方,耶路撒冷是大軍所在地,是敬畏神的地方.這極大的美麗,非言話或筆墨所能描述.這美是秀麗的美,也是神聖的美.這種美,我們是應該有的,應該有耶路撒冷的光景,在我們裡面.神所指定的地方,是神的保障,是神的座位中心.萬民都應回到此地來領受教訓.
\newline
\newline
「威武如展開旌旗的軍隊.」(歌六4).「那向外觀看如晨光發現,美麗如月亮,皎潔如日頭,威武如展開旌旗軍隊的是誰呢.」(歌六10).
\newline
\newline
在這裡說到威武的旌旗.我們應該安息,因這是我們裡面的操練,當我們真正享到這安息時,雖有人罵你,批評你,你都不會有任何的反應,因你裡面有了安息.正如關閉了的井,裡面的井水不會有波動的.基督徒裡面有安息,外面有威武,那是什麼一回事?基督是我們的旌旗,是我們的得勝,我們應有這屬靈的威風.這威風,並非我們的驕傲,知識及自我的誇獎,而是我們裡面的豐富,到了一個地步.我們親近神,魔鬼就離開逃跑了,基督就在我們心中作王,世界不再向我們起作用,這是屬靈的威武.保羅說:你們要作安靜的人.屬靈的威武是安靜,溫柔.主已留下了榜樣,祂的威武就是那溫柔,謙卑的樣式.這屬靈的戰爭裡,並不是那高聲的話語,或強烈的態度,而是我們在神前那謹守的心情.這旌旗,這威武,我們該如何表現出來?
\newline
\newline
(一)威武能使我們抵擋罪──任何時間有罪來到時,我們都能抵擋.肉體是我們最大的仇敵,魔鬼是我們的仇敵,的確是一大衛脅.但聖經上說:如果我們能親近神,魔鬼就會離開我們逃跑了(雅四7).世界是我們的仇敵,如果我們不去愛它,它就是死的.如果我們看世界是死的,那麼世界的一切就都不可愛了.所以我們最大的仇敵,乃是自己,保羅說:「我真是苦呵!誰能救我脫離這取死的身體呢?」(羅七24).我們常將肢體獻給罪,作不義的器具.(羅六13).所以要勝過肉體是最困難的,因肉體常被罪所利用,又牽引我們進入罪惡中.肉體內有犯罪的性情,就是情慾,這情慾在我們裡面發動時,就叫我們去犯罪.為這緣故,我們應有威武,應把這旗幟舉起；將主的名舉起.若要勝過自己,就千萬不能靠自己.要藉著禱告,操練神的同在.有時神不聽我們禱告,因神要我們學習功課；要我們養成一個多禱告的習慣,要看見自己的卑微.在旌旗之下,日頭照著,月亮的光也照下來；使我們看見自己.讓我們學習到如何禱告.及讚美,生命就也豐盛起來.能抵擋罪,勝過肉體.
\newline
\newline
(二)為真道而辯護──我們要站穩立場,守住我們的崗位.新婦之威武是在旌旗下,你我是軍隊,應有軍隊之風紀,和操練.應順服,聽命,照命令去行事.經過操練的軍隊,然後威風才發出來；在旌旗之下,只許得勝不許失敗.
\newline
\newline
(三)威武使我們靈裡剛強──我們不斷在屬靈的爭戰中,應披上全副軍裝,以便抵擋一切問題的來臨.十字架的道路是一場爭戰,這爭戰,切勿用肉體,要用安息去爭戰.我們倚靠主,是在他旗幟之下,在他得勝之下而向前行.
\newline
\newline
在歌羅西書三章十七節說:「無論作什麼,或說話,或行事,都要奉主耶穌的名,藉著祂感謝父神.」在此經文中告訴我們有兩個秘訣:
\newline
\newline
a 奉主的名──我們的一切生活都在主的旗幟下,這旗幟是什麼呢?一是愛旗,另一是軍旗,這軍旗有基督的名字的.古代作戰時,每一軍隊作戰都有其統領的名字底旗幟,今日我們作戰的軍旗,就有基督的名字.當旗幟一張開,魔鬼就逃跑,罪惡就離開.因為無論作什麼,或說話,或行事,都要奉耶穌的名.
\newline
\newline
b 感謝父神──無論在何時何環境都感謝父神,在富有時,貧窮時都應感謝父神.接受父神為你所安排的一切.那麼,你的威風就可以顯出來,撒旦要找我們的缺點和漏洞,當你埋怨或不滿時,撒旦就乘虛而入.但若我們處處感謝主,將主的名字旗幟高舉,則無論在何時主的名字都得著榮耀.
\newline
\newline


\newpage

\section{第46屆~港九培靈會~于力工~第七講}
\label{sec:1286}
\textbf{基督是我們的園主}
\newline
\newline
連結: \href{https://www.hkbibleconference.org/session-message/view/1286}{\texttt{https://www.hkbibleconference.org/session-message/view/1286}}
\newline
\newline
\hyperref[sec:1204]{< < < PREV SERMON < < <}
~
\hyperref[sec:toc]{[返目錄]}
~
\hyperref[sec:1206]{> > > NEXT SERMON > > >}
\newline
\newline
以賽亞書 5:1-5:3
\newline
\begin{longtable}{cl}
\hline
\hline
章節 & 經文 (和合本修訂版)\\
\hline
5:1 & \begin{tabularx}{0.7\textwidth}{X} 我要為我親愛的唱歌，我所愛的、他的葡萄園之歌。我親愛的有葡萄園在肥沃的山岡上。 \end{tabularx} \\ \\ \relax
5:2 & \begin{tabularx}{0.7\textwidth}{X} 他刨挖園子，清除石頭，栽種上等的葡萄樹，在園中蓋了一座樓，又鑿出酒池；指望它結葡萄，反倒結了野葡萄。 \end{tabularx} \\ \\ \relax
5:3 & \begin{tabularx}{0.7\textwidth}{X} 耶路撒冷的居民和猶大人哪，現在，請你們在我與我的葡萄園之間斷定是非。 \end{tabularx} \\ \\
[1ex]
\hline
\hline
\end{longtable}
在聖經中使我們看見神對救恩的計劃,包括了三件重要的事:
\newline
\newline
(一)在創世記二章中告訴我們,神為人安置了一個園子；要人在園子裡一面生活,一面工作.
\newline
\newline
按神的心意,原喜歡人住在祂所選的園子裡,一面工作生活,一面享受；但人卻失去了這園子.人因犯罪的緣故,就失去了工作和生活的地方；被逐到滿長荊棘蒺藜之地!於是人要汗流滿面才得糊口,這是第一件事.
\newline
\newline
為了要挽回失去的福份,回復伊甸園的生活；主耶穌就要釘在十字架上,流出祂的寶血；然後使我們回復亞當夏娃所失去的地位.按神的心意,不單如此,祂還要預備一個地方.這地方擴充得更大,更美麗,更榮耀,更有自由和享受；比以前更豐富,這就是將來的新耶路撒冷.這也是我們朝夕所盼望的,我們要永遠居住在其中,所以聖經從頭到末了,都講論及這三件事.在創一章說:「起初神創造天地,地是空虛混沌,淵面黑暗.」為何會黑暗,混沌呢?這與我們無關,因為是在亞當夏娃之前的事,原來新耶路撒冷降在地上時,城外那些行邪術的,有犬類,他們如何亦與我無關.
\newline
\newline
亞當夏娃失去了的園子,神就要重新把這個園子建立.在亞當的後裔中,找出一個人來,要使他們回到伊甸園的生活中,此人是亞伯.當亞伯還末將這種生活完全活出來時,就被哥哥殺了!於是神又再找到一個人,那是在洪水之前,神看中了以諾；以諾在地上與神同行三百年,最後就回到神那裡去.繼續神又看中了挪亞,挪亞與神同行,並作了大事,但神之心意仍未能展開.直到有一天,神在一個拜偶像的民族裡；一個邪惡,黑暗的世代裡；找到了一個人.這人就是亞伯拉罕.神呼召亞伯拉罕,叫他離開本地,本族,父家,往我所指示你的地方去,亞伯拉罕聽神的吩咐就去了.神說明我要賜福與你,我要與你立約.這約就是至高至大的神,全能全愛的神；降為人的地位,與他立約.要將當初原擬賜給亞當及夏娃的福份賜給他,只要他肯遵行神的話,就必得著.神與亞伯位罕,立了這約；為了此約,神必須將以色列人從埃及帶出來.到了西乃山神又與以色列人立了約.神要他們在列國中作君尊的祭司,代表祂生活在世上,回復亞當夏娃失去的一切,照祂心意去行.只要他們遵從神的話,神就要將一切的福份賜給他們.神吩咐他們不可拜偶像,不可隨從外邦人去事奉他們的神,百姓都齊聲說,阿們；神特別愛以色列人,在他們四圍圈上離笆,保護他們.先知以賽亞曾作歌說:親愛者所唱的葡萄園的歌.這葡萄園就是神給以色列人建造之地是迦南地,是應許之地.神要他們在這園子中,好好地工作和生活.只要肯聽從神的話,生活得像祭司一樣,就是他們的工作了.試問亞伯拉罕,以撒,雅各他們領了多少人歸向神?作了什麼工作?其實他們沒有什麼工作.亞伯位罕的一生,只表現了一個信息:他因著「信」神就稱他為義.他這一生所過的生活,就是表現信心的生活,這就是他的工作.許多時候我們忽略了一件事；常說自己無能,沒有恩賜,不會講道,不會工作,到底我一生是應該做些什麼呢?其實,神不是看我們的工作有多少,而是看我們的生活如何?若你的生活像主一樣,能活出祭司的身份來,活出君王而帶著尊貴；這就是你的工作了.你如果有好行為,外邦人看見了,就會將榮耀歸給神.不是我們能吸引人歸主,主曾說:當祂在地上被舉起的時候,就能吸引萬人歸向祂.所以基督若在我們身上,我們才能吸引人歸主.故生活就是工作.
\newline
\newline
以賽亞書五章提到,神又把他們安置了葡萄園.在這園子裡,他們又犯罪離開了神!各位能否看出這園子是什麼意思?園子是指他們的國家,神是園子裡的主人.經過了基督的十字架,基督的復活；神又另外安置一個園子,這就是羅馬書九,十,十一章所說的話.這三章說及以色列人犯罪,因此被神所丟棄,並且被趕出來,直趕到去巴比倫及分散天下各處.請注意:在主前七二二年時,十支派被亞述人擄去,從此十支派被分散在列邦.主前五八六年,僅剩下的猶大人,及利未支派的人,又被擄到巴利倫去,今日的猶大人,只有部分能真真正正知道他們是猶大支派和利未支派的後裔；但其他的支派,就不知到了那裡去了!顯然他們已從這園子中被趕出去.感謝神!祂從來沒有停止過工作,經過了基督的十字架之後,祂又建立了一個園子.在羅馬書十一章末段,保羅稱頌神的奇妙,和豐盛的恩典.祂是應該得著一切的榮耀和讚美的.祂的知識豐富,祂的判斷何其再難測,誰能知道主的心!萬有都是本於祂,倚靠祂.歸於祂,願榮耀歸給祂直到永遠!阿們.保羅在羅馬十二章說:所以我以神的慈悲勸你們,將身體獻上,作為活祭.你們如此事奉,乃是理所當然的,是聖潔的,是神所喜悅的.從此可見,神在這時代建立了另外一個園子,祂不稱之為葡萄園,而稱為教會.今日你我進入這園子裡,我們一同配搭,一群事奉,建立基督的身體.保羅曾說:「你們要在基督的裡面要建造.」彼此之間需建立,然後纔成為教會.保羅說:將身體獻上當作活祭.你我的心就是一個園子,這園子乃是經過主的寶血所買贖的.從以上的園子中,我們也看見三件事:
\newline
\newline
(歌一6)書拉密女說:「我同母的弟兄向我發怒,他們使我看守葡萄園；我自己的葡萄園卻沒有看守.」她被弟兄們強迫.看守別人的葡萄園；而忽略了自己的葡萄園.我們也要反省,是否集中力量在自己的葡萄園裡工作?好好地看守?這葡萄園是屬於你的,要專心看守.
\newline
\newline
(一)要專一──(羅十二6-18)說:作教導的,應專一教導,事奉的當專一事奉,治理的當專一治理......意思是每人自己所領受的恩賜,能在教會中善用它.可能你有許多種恩賜,但不應樣樣都由你而做.建立教會,是要彼此配搭,配搭乃是織在一起的意思；亦即彼此在一起建立之意.「門徒」二字有學習跟從的意思,我們不單作工就算；也要在別人身上學習,要安靜,順服,也要為別人的恩賜而感謝神.在自己的葡萄園中,要知道自己的本份,守住自己的崗位；那麼教會才有次序,能長進.肢體間才不至有論斷,批評,對付別人.求神憐憫我們,應用禱告的心,扶助有恩賜的人.我們有為牧師,長老,執事祈禱的嗎?雅各說,我們要彼此認罪,互相代求.這是非常美麗的操練,只有這樣才能將葡萄園看守得好,免去批評和拆毀.
\newline
\newline
(二)留心防避那毀壞葡萄園的小狐狸──因為我們的葡萄正在開花.(歌二15)葡萄正開花時,小狐狸最多.那就是說在教會中,在個人的生命裡,有許多漏洞,容易被小狐狸溜進來；傷葡萄的根,葡萄就沒有用了.求主教我們曉得如何防避.靠著園主一同將小狐狸捉拿,不讓牠進來,不要給魔鬼留地步.也不要給自己留犯罪的地步,免得有問題進入我們當中.問題來了；破壞也跟著來,難怪我們結不出果子來.我們常在小事上出問題,人若在小事上不忠心,神就無法將大事交他管理,如果我們盡心去作工,有一天神會憐憫我們說:你這又忠心又良善的僕人,我可以把更大的工作交給你.不然的話,神會說:你這又懶又惡的僕人,你到外面吧!所以我們要防備小狐狸.免得毀壞了果子.園子是屬園主的,我們不過是看守者.祂會負我們的責任,我們又何必憂心?何必緊張呢?到主面前抓緊祂的安息,將自己交給祂,祂會將我們從困難中帶領出來.
\newline
\newline
(三)任何環境都追求靈命長進──在(歌四13-五2),說及園內有果子和花,又說:北風呵!興起,南風呵!吹來,吹在我的園內,使其中的香氣發出來.願我的良人進入自己的園裡,喫祂佳美的果子.不單祂喫,連朋友也來喫；不單喫果子,也有香氣,有沒藥,乳香.有果,有花,有香料,這是豐富的園子.裡面泉源有井水,可以解渴.主要進入果子裡,北風可以吹來.雖然寒冷,但可以發出香氣,南風吹來；雖然很熱,但仍能結出果子來.無論在任何環境裡,都能產生各樣的果子.這表明基督的豐滿在我們裡面,一切的豐盛是有形有體,居住在我們裡面.兄姊們!我們要向前追求,不要停留.去享受園中的豐富,你會覺得甘甜滿足,生命也愈豐富.基督的豐富有多少?在歌羅西書中可見到最少十樣,現在讓我只提三樣:(一)豐富的知識.(二)豐富的信息,(三)豐富的盼望.
\newline
\newline
我們豐富的主,是我們今天所追求的,祂是何等豐富,是我們何等大的指望,別將自己看得太小,保羅說:這豐盛,在基督裡面,就是基督的奧秘.基督就是神的奧秘,為何說是神的奧秘?奧秘就是不能明白和不能完全了解的,也無辦法完全得著的.這奧秘可存留在心中,將基督放在心中,也住在基督裡面.我們要遵守祂的話.(約十五章).住在祂的話裡,就等於住在祂裡面了.
\newline
\newline
基督住在我裡面,我也住在基督裡面；這就是照著祂的話去行,將祂的話天天思想,也就等於住在基督裡面.當主來到我心中,看見有花朵,聞到了香氣,也喫到果子；主就能滿足.試問,我們能否滿足主的心呢?
\newline
\newline


\newpage

\section{第46屆~港九培靈會~于力工~第八講}
\label{sec:1206}
\textbf{基督是我們的朋友}
\newline
\newline
連結: \href{https://www.hkbibleconference.org/session-message/view/1206}{\texttt{https://www.hkbibleconference.org/session-message/view/1206}}
\newline
\newline
\hyperref[sec:1286]{< < < PREV SERMON < < <}
~
\hyperref[sec:toc]{[返目錄]}
~
\hyperref[sec:1207]{> > > NEXT SERMON > > >}
\newline
\newline
雅歌 5:15-5:23
\newline
\begin{longtable}{cl}
\hline
\hline
章節 & 經文 (和合本修訂版)\\
\hline
5:15 & \begin{tabularx}{0.7\textwidth}{X} 他的腿好比白玉石柱，安在精金座上；他的容貌如黎巴嫩，佳美如香柏樹。 \end{tabularx} \\ \\ \relax
5:16 & \begin{tabularx}{0.7\textwidth}{X} 他的口甘甜，他全然可愛。耶路撒冷的女子啊，這是我的良人，這是我的朋友。 \end{tabularx} \\ \\ \relax
6:1 & \begin{tabularx}{0.7\textwidth}{X} 你這女子中最美麗的，你的良人往何處去？你的良人轉向何處去了？我們好與你同去尋找他。 \end{tabularx} \\ \\ \relax
6:2 & \begin{tabularx}{0.7\textwidth}{X} 我的良人進入自己園中，到香花園，在園內放牧，採百合花。 \end{tabularx} \\ \\ \relax
6:3 & \begin{tabularx}{0.7\textwidth}{X} 我屬我的良人，我的良人屬我；他在百合花中放牧。 \end{tabularx} \\ \\ \relax
6:4 & \begin{tabularx}{0.7\textwidth}{X} 我的佳偶啊，你美麗如得撒，秀美如耶路撒冷，威武如展開旌旗的軍隊。 \end{tabularx} \\ \\ \relax
6:5 & \begin{tabularx}{0.7\textwidth}{X} 求你轉眼不要看我，因你的眼睛使我慌亂。你的頭髮如同一群山羊，從基列山下來。 \end{tabularx} \\ \\ \relax
6:6 & \begin{tabularx}{0.7\textwidth}{X} 你的牙齒如一群母羊，洗淨之後走上來，它們成對，沒有一顆是單獨的。 \end{tabularx} \\ \\ \relax
6:7 & \begin{tabularx}{0.7\textwidth}{X} 你的鬢角在面紗後，如同迸開的石榴。 \end{tabularx} \\ \\ \relax
6:8 & \begin{tabularx}{0.7\textwidth}{X} 雖有六十王后、八十妃嬪，並有無數的童女。 \end{tabularx} \\ \\ \relax
6:9 & \begin{tabularx}{0.7\textwidth}{X} 她是我獨一的鴿子、我完美的人兒，是她母親獨生的，是生養她的所寵愛的。女子見了都稱她有福，王后妃嬪見了也讚美她。 \end{tabularx} \\ \\ \relax
6:10 & \begin{tabularx}{0.7\textwidth}{X} 那俯視如晨曦、美麗如月亮、皎潔如太陽、威武如展開旌旗軍隊的是誰呢？ \end{tabularx} \\ \\ \relax
6:11 & \begin{tabularx}{0.7\textwidth}{X} 我下到堅果園，要看谷中青翠的植物，要看葡萄可曾發芽，石榴可曾放蕊； \end{tabularx} \\ \\ \relax
6:12 & \begin{tabularx}{0.7\textwidth}{X} 不知不覺，我彷彿坐在我百姓高官的戰車中。 \end{tabularx} \\ \\ \relax
6:13 & \begin{tabularx}{0.7\textwidth}{X} 回來，回來，書拉密的女子；回來，回來，我們要看你。你們為何要觀看書拉密的女子，像觀看兩隊人馬在跳舞呢？ \end{tabularx} \\ \\ \relax
7:1 & \begin{tabularx}{0.7\textwidth}{X} 尊貴的女子啊，你的腳在鞋中何等秀美！你的大腿圓潤，好像美玉，是巧匠的手做成的。 \end{tabularx} \\ \\ \relax
7:2 & \begin{tabularx}{0.7\textwidth}{X} 你的肚臍如圓杯，不缺調和的酒。你的肚子如一堆麥子，周圍有百合花。 \end{tabularx} \\ \\ \relax
7:3 & \begin{tabularx}{0.7\textwidth}{X} 你的兩乳好像一對小鹿，是母鹿雙生的。 \end{tabularx} \\ \\ \relax
7:4 & \begin{tabularx}{0.7\textwidth}{X} 你的頸項如象牙塔，你的眼睛像希實本、巴特‧拉併門旁的水池，你的鼻子彷彿朝向大馬士革的黎巴嫩塔。 \end{tabularx} \\ \\ \relax
7:5 & \begin{tabularx}{0.7\textwidth}{X} 你的頭在你身上好像迦密山 ，你頭上的髮呈紫色，王被這髮綹繫住了。 \end{tabularx} \\ \\ \relax
7:6 & \begin{tabularx}{0.7\textwidth}{X} 我親愛的，喜樂的女子啊，你何等美麗！何等令人喜悅！ \end{tabularx} \\ \\ \relax
7:7 & \begin{tabularx}{0.7\textwidth}{X} 你的身材好像棕樹，你的兩乳如同纍纍的果實。 \end{tabularx} \\ \\ \relax
7:8 & \begin{tabularx}{0.7\textwidth}{X} 我說：我要爬上棕樹，抓住枝子。願你的兩乳好像葡萄纍纍，願你鼻子的香氣如蘋果； \end{tabularx} \\ \\ \relax
7:9 & \begin{tabularx}{0.7\textwidth}{X} 你的上顎如美酒，直流入我良人的口裡，流入沉睡者的口中。 \end{tabularx} \\ \\ \relax
7:10 & \begin{tabularx}{0.7\textwidth}{X} 我屬我的良人，他也戀慕我。 \end{tabularx} \\ \\ \relax
7:11 & \begin{tabularx}{0.7\textwidth}{X} 來吧！我的良人，讓我們往田間去，在村莊住宿。 \end{tabularx} \\ \\ \relax
7:12 & \begin{tabularx}{0.7\textwidth}{X} 早晨讓我們起來往葡萄園去，看葡萄樹發芽沒有，花開了沒有，石榴放蕊沒有，在那裡我要將我的愛情給你。 \end{tabularx} \\ \\ \relax
7:13 & \begin{tabularx}{0.7\textwidth}{X} 曼陀羅草散發香味，在我們的門內有各樣新陳佳美的果子；我的良人，這都是我為你保存的。 \end{tabularx} \\ \\ \relax
8:1 & \begin{tabularx}{0.7\textwidth}{X} 惟願你像我的兄弟，像吃我母親奶的兄弟。我在外頭遇見你就與你親吻，誰也不輕看我。 \end{tabularx} \\ \\ \relax
8:2 & \begin{tabularx}{0.7\textwidth}{X} 我必引導你，領你進入我母親的家，她必教導我，我必使你喝石榴汁釀的香酒。 \end{tabularx} \\ \\ \relax
8:3 & \begin{tabularx}{0.7\textwidth}{X} 他的左手在我頭下，他的右手將我環抱。 \end{tabularx} \\ \\ \relax
8:4 & \begin{tabularx}{0.7\textwidth}{X} 耶路撒冷的女子啊，我囑咐你們，不要喚醒、不要挑動愛情，等它自發。 \end{tabularx} \\ \\ \relax
8:5 & \begin{tabularx}{0.7\textwidth}{X} 那靠著良人從曠野上來的是誰呢？在蘋果樹下，我叫醒了你；在那裡，你母親曾為了生你而陣痛，在那裡，生你的為你陣痛。 \end{tabularx} \\ \\ \relax
8:6 & \begin{tabularx}{0.7\textwidth}{X} 求你將我放在你心上如印記，帶在你臂上如戳記。因為愛情如死之堅強，熱戀如陰間之牢固，所發的光是火焰的光，是極其猛烈的火焰。 \end{tabularx} \\ \\ \relax
8:7 & \begin{tabularx}{0.7\textwidth}{X} 愛情，眾水不能熄滅，江河也不能淹沒。若有人拿家中所有的財寶要換愛情，就全被藐視。 \end{tabularx} \\ \\ \relax
8:8 & \begin{tabularx}{0.7\textwidth}{X} 我們有一小妹，她還沒有乳房，人來提親的日子，我們當為她怎麼辦呢？ \end{tabularx} \\ \\ \relax
8:9 & \begin{tabularx}{0.7\textwidth}{X} 她若是牆，我們要在其上建造銀塔；她若是門，我們要用香柏木板圍護她。 \end{tabularx} \\ \\ \relax
8:10 & \begin{tabularx}{0.7\textwidth}{X} 我是牆，我的兩乳像塔。那時，我在他眼中是找到平安的人。 \end{tabularx} \\ \\ \relax
8:11 & \begin{tabularx}{0.7\textwidth}{X} 所羅門在巴力‧哈們有一葡萄園，他將這葡萄園租給看守的人，每人為其中的果子要交一千銀子。 \end{tabularx} \\ \\ \relax
8:12 & \begin{tabularx}{0.7\textwidth}{X} 我有屬自己的葡萄園。所羅門哪，一千歸你，兩百歸看守果子的人。 \end{tabularx} \\ \\ \relax
8:13 & \begin{tabularx}{0.7\textwidth}{X} 你這住在園中的，同伴都要聽你的聲音，求你使我也得以聽見。 \end{tabularx} \\ \\ \relax
8:14 & \begin{tabularx}{0.7\textwidth}{X} 我的良人哪，求你快來！像羚羊，像小鹿，在香草山上。 \end{tabularx} \\ \\
[1ex]
\hline
\hline
\end{longtable}
雅歌書五章,論到對良人的認識.女子說:「我的良人,白而且紅,超乎萬人之上.」但對耶路撒冷的女子說:「這是我的良人,這是我的朋友.」按我們平常的看法,既是良人,又是朋友?到底表示什麼?要知道,這些耶路撒冷的女子,也包括皇宮中的女子,所羅門的妃嬪.我要告訴你,王是我的朋友,且不是普通的朋友.我們的神降到一個地位,使我感到祂是我們的朋友,而且祂又到另一個地步,對我們說:你是我的朋友,這關係非常重要.
\newline
\newline
(一)為什麼亞伯拉罕能稱為神的朋友?全能的神,創造了宇宙萬物,何以又謙卑到這麼一個地步,說:「亞伯拉罕,你是我的朋友!」各位有沒有感到神怎能稱人為朋友呢?假如我們有讀聖經,就應知道,神把愛祂的人都稱他們為朋友.以諾與神同行,有三百年之久.二人若不同心,又怎能同行?他們是好朋友.我們與主之間,有沒有這關係?當我們走路時,或工作時,正像在禱告時一樣?主若在我身旁,我就誠實,行事為人也坦白正直呢?也許我們會追問亞伯拉罕怎麼能成為神的朋友?神甚至曾說:「我們豈可隱瞞亞伯拉罕去作一件事呢.」(創十八17)當神要滅所多瑪和俄摩拉城時,神沒有向亞伯拉罕隱瞞,亞伯拉罕知道了後,就在神面前,用朋友的立場向神說話.他說:「神啊!若城中有五十個義人,你也要毀滅他們嗎?」他站在同等地位上與神談話.這是何等的甜密!所以我們在主面前,也應該這種談話.這才是真正禱告的生活,因為他可站在神的朋友的地位上與神談話.這不是報告,乃是在屬靈的操練裡,進入到談話的地步.禱告時將主的話默想,成為我們生命的供應,天天如此操練.進入禱告的生活裡,進入與神同在的境地裡.
\newline
\newline
(二)亞伯拉罕能稱為神的朋友,因他一方面有信心,一方面亦有行為.他不單信神,且信神的話,又照神的話去行.最明顯的例子就是神對他說:「亞伯拉罕,你將你所最愛的兒子,獻給我.」他沒有和神辯論,立時吩咐僕人預備驢,和柴,然後往神所指示的地方去.走了三天的路程,到了摩利亞山,就把僕人留下,帶以撒上山去.以撒說:火與柴都有了,但燔祭的羊羔在那裡?亞伯拉罕說:「我兒!耶和華必自己預備,作燔祭的羊羔.」當亞伯拉罕正要伸手殺他兒子時,耶和華的使者就攔阻他.叫他看見樹林中一隻公羊,就把那公羊殺了獻為燔祭,代替他的兒子.在這件事上,神對亞伯拉罕說:「你既行了這事,不留下你的兒子,就是你獨生的兒子,我便指著自己起誓說,論福,我要賜大福給你.」(創廿二17)亞伯拉罕能稱為神的朋友,因他肯順服神的話.在神面前願意放下一切；單純地順服,這是我們應學習的功課.順服後的亞伯拉罕,就帶來了一個大的果效.那就是:「原來在耶和華山上必有預備.」所以給那地方起名叫「耶和華以勒.」他的順服,將神的奇妙顯出,將神的新名字帶出來.從來沒有人知道耶和華以勒.但現在卻把神的應許顯明出來,好像從倉庫的寶貝帶出來一般,今天我們才會說耶和華以勒.
\newline
\newline
我們對順服的功課,應好好地學習.我們常反抗神,還以為照著真理行；所以不能把生命的豐富流露出來,反而失去了見証.
\newline
\newline
神的名字逐漸地帶出來,在創世記中就可以知道.因為沒有一個名字,是可以完全的代表神,例如耶和華,本來是救贖者的意思,但原文讀音方面,就是說沒辦法,將神的名字讀出來,就如一種呼求之意.當時希伯來人找不出一個名字,稱祂為神；所以就用這「啊!那一位至高至大至聖者啊!」感謝主!經過了屬靈的經歷,在神前面的順服.每學一功課時就說:「啊!原來神是全能者.」(以勒沙代)這全能者是剛強,有力,像大山一樣.之後,又見神的名字稱為奇妙,策士,和平之君,永在的父.你的生命如何?你就表現出你的神是如何.多少時從我們的談話中見証中,我們把神顯得太小太可憐了!不能作事也不聽禱告,把神限制得很小,這是我們的問題.我們要將神表現在人面前,祂是全能者,是奇妙,是策士,耶和華以勒.你如何順服,就如何的顯明.弟兄們!屬靈的道路,順服的道路,就是進入到朋友的地步.求主使我們看見這事實,不要限制神,因亞伯拉罕的行為是照神的話去行.今日教會弟兄姊妹們,許多是聽道專家,而不是行道專家.主曾說:凡聽見我的話又去行的,就好比一個聰明人把房子蓋在磐石上；雨淋,水沖,風吹,撞著那房子,房子總不會倒塌.我們不但要聽神的話,更要實行神的話.求主憐憫我們,祂實在給我們許多的福氣.也聽見了不少的真理.問題是肯不肯把所聽見的去操練?當我們在禱告會中,有幾多位能站起來禱告?當我們一起交通時,有多少人肯將主的話說出來分享?教會今天的貧乏,不是缺少信息,乃是缺少將信息拿出來操練.當我們聽到要愛主,要奉獻的時候,我們有沒有實行奉獻呢?我們是否只落在知識上?若只是落在知識上,就永遠無法長進了.一個人能認識基督,與其得著基督,是兩回事.如果你沒有去嘗試過,就不會知道有什麼果效.所以認識基督是要緊的,但得著基督更要緊；繼而效法基督就最要緊的.我們認識了祂,得著了祂.然後就要效法祂行事為人.許多時候,我們攔阻了人,並非我們沒有知識,而是沒有表現.到了一個時候,當人發覺你與別人有所不同的話,他才願意親近你,因為你的行為效法了基督.求主憐恤我們.我們要表現基督,無論任何場合,任何環境；你的一舉一動,一言一談都可以表現基督.因為別人看見了你的表現,就可經認識基督,所以我們要表現神是我們的朋友,祂實在是我們的朋友.
\newline
\newline
從良人的角度來看,在(歌五10-16)中,可以看見一幅圖畫.如果我們以祂為我的朋友,就可以看見不同的地方.「我的良人白而且紅,超乎萬人之上」(五10).誰能超乎萬人之上?可能是帝王,那麼帝王就是我們最尊貴的朋友.假如今天有一位總統對你說:我是你的朋友,那麼你的感覺如何?當然有所不同了.在十一節中說及祂的頭像至精的金子,乃是說基督是我們純潔的朋友,是鍛鍊過的.祂的頭髮方面,是表明祂是寶貝的朋友,因頭髮代表榮耀,美麗.「祂的眼睛如同鴿子」(十二節).這是代表正直的朋友.「嘴唇像百合花」(十三節).祂的口中有良言,是有良言真正的朋友,不是只會奉承的人.「祂的兩手像金管」(十四節).這是個援助的朋友,祂經過彫刻,是患難中的朋友,「祂的腿好像白玉石柱」(十五節).這是可靠的朋友.「祂的口極其甘甜」(十六節).這表示可親近的朋友.最後:「祂全然可愛」(十六節).是完全的朋友.用這不同的角度來看,可見是一位完全的主,完全的朋友.有了祂在身旁,作我們的朋友,實在感謝主!
\newline
\newline
這朋友的建立,在新約中告訴我們,這「建立」是在乎我們是否順服他,亞伯拉罕有信心,因有好的行為,照著祂的話去行,就成為神的朋友.
\newline
\newline


\newpage

\section{第46屆~港九培靈會~于力工~第九講}
\label{sec:1207}
\textbf{基督是我們的一切}
\newline
\newline
連結: \href{https://www.hkbibleconference.org/session-message/view/1207}{\texttt{https://www.hkbibleconference.org/session-message/view/1207}}
\newline
\newline
\hyperref[sec:1206]{< < < PREV SERMON < < <}
~
\hyperref[sec:toc]{[返目錄]}
~
\hyperref[sec:1208]{> > > NEXT SERMON > > >}
\newline
\newline
雅歌 1:1-8:14
\newline
\begin{longtable}{cl}
\hline
\hline
章節 & 經文 (和合本修訂版)\\
\hline
1:1 & \begin{tabularx}{0.7\textwidth}{X} 所羅門的雅歌。 \end{tabularx} \\ \\ \relax
1:2 & \begin{tabularx}{0.7\textwidth}{X} 願他用口與我親吻。你的愛情比酒更美， \end{tabularx} \\ \\ \relax
1:3 & \begin{tabularx}{0.7\textwidth}{X} 你的膏油馨香，你的名如傾瀉而出的香膏，所以童女都愛你。 \end{tabularx} \\ \\ \relax
1:4 & \begin{tabularx}{0.7\textwidth}{X} 願你吸引我跟隨你；讓我們快跑吧！王領我進入他的內室。我們必因你歡喜快樂，我們要思念你的愛情，勝似思念美酒。她們愛你是理所當然的。 \end{tabularx} \\ \\ \relax
1:5 & \begin{tabularx}{0.7\textwidth}{X} 耶路撒冷的女子啊，我雖然黑，卻是秀美，如同基達的帳棚，好像所羅門的幔子， \end{tabularx} \\ \\ \relax
1:6 & \begin{tabularx}{0.7\textwidth}{X} 不要因太陽把我曬黑了就瞪著我。我母親的兒子向我發怒，他們使我看守葡萄園；我自己的葡萄園我卻沒有看守。 \end{tabularx} \\ \\ \relax
1:7 & \begin{tabularx}{0.7\textwidth}{X} 我心所愛的啊，請告訴我，你在何處牧羊？正午在何處使羊歇臥？我何必像蒙著臉的女子在你同伴的羊群旁邊呢？ \end{tabularx} \\ \\ \relax
1:8 & \begin{tabularx}{0.7\textwidth}{X} 你這女子中最美麗的，你若不知道，只管跟隨羊群的腳蹤行，在牧人的帳棚邊，牧放你的小山羊。 \end{tabularx} \\ \\ \relax
1:9 & \begin{tabularx}{0.7\textwidth}{X} 我的佳偶，你好比法老戰車上的駿馬。 \end{tabularx} \\ \\ \relax
1:10 & \begin{tabularx}{0.7\textwidth}{X} 你的兩頰因髮辮而秀美，你的頸項因珠串而華麗。 \end{tabularx} \\ \\ \relax
1:11 & \begin{tabularx}{0.7\textwidth}{X} 我們要為你編上金鏈，鑲上銀飾。 \end{tabularx} \\ \\ \relax
1:12 & \begin{tabularx}{0.7\textwidth}{X} 王正坐席的時候，我的哪噠香膏散發香氣。 \end{tabularx} \\ \\ \relax
1:13 & \begin{tabularx}{0.7\textwidth}{X} 我的良人好像一袋沒藥，在我胸懷中。 \end{tabularx} \\ \\ \relax
1:14 & \begin{tabularx}{0.7\textwidth}{X} 我的良人好像一束鳳仙花，在隱‧基底的葡萄園中。 \end{tabularx} \\ \\ \relax
1:15 & \begin{tabularx}{0.7\textwidth}{X} 看哪，我的佳偶，你真美麗！看哪，你真美麗！你的眼睛是鴿子。 \end{tabularx} \\ \\ \relax
1:16 & \begin{tabularx}{0.7\textwidth}{X} 看哪，我的良人，你多英俊可愛！讓我們以青草為床榻， \end{tabularx} \\ \\ \relax
1:17 & \begin{tabularx}{0.7\textwidth}{X} 以香柏樹為房子的棟梁，以松樹作屋頂的椽木。 \end{tabularx} \\ \\ \relax
2:1 & \begin{tabularx}{0.7\textwidth}{X} 我是沙崙的玫瑰花，是谷中的百合花。 \end{tabularx} \\ \\ \relax
2:2 & \begin{tabularx}{0.7\textwidth}{X} 我的佳偶在女子中，好像荊棘裡的百合花。 \end{tabularx} \\ \\ \relax
2:3 & \begin{tabularx}{0.7\textwidth}{X} 我的良人在男子中，如同蘋果樹在樹林裡。我歡歡喜喜坐在他的蔭下，嘗他果子的滋味，覺得甘甜。 \end{tabularx} \\ \\ \relax
2:4 & \begin{tabularx}{0.7\textwidth}{X} 他領我進入宴會廳，為我插上愛的旗幟。 \end{tabularx} \\ \\ \relax
2:5 & \begin{tabularx}{0.7\textwidth}{X} 請你們用葡萄餅增補我力，以蘋果暢快我的心，因我為愛而生病。 \end{tabularx} \\ \\ \relax
2:6 & \begin{tabularx}{0.7\textwidth}{X} 他的左手在我頭下，他的右手將我環抱。 \end{tabularx} \\ \\ \relax
2:7 & \begin{tabularx}{0.7\textwidth}{X} 耶路撒冷的女子啊，我指著羚羊或田野的母鹿囑咐你們，不要喚醒，不要挑動愛情，等它自發。 \end{tabularx} \\ \\ \relax
2:8 & \begin{tabularx}{0.7\textwidth}{X} 聽啊！我良人的聲音，看哪！他穿山越嶺而來。 \end{tabularx} \\ \\ \relax
2:9 & \begin{tabularx}{0.7\textwidth}{X} 我的良人像羚羊，像小鹿。看哪，他站在我們的牆壁邊，從窗戶往裡觀看，從窗格子往裡窺探。 \end{tabularx} \\ \\ \relax
2:10 & \begin{tabularx}{0.7\textwidth}{X} 我的良人對我說：「我的佳偶，起來！我的美人，與我同去！ \end{tabularx} \\ \\ \relax
2:11 & \begin{tabularx}{0.7\textwidth}{X} 看哪，因為冬天已逝，雨水止住，已經過去了。 \end{tabularx} \\ \\ \relax
2:12 & \begin{tabularx}{0.7\textwidth}{X} 地上百花開放，歌唱的時候到了，斑鳩的聲音在我們境內也聽見了。 \end{tabularx} \\ \\ \relax
2:13 & \begin{tabularx}{0.7\textwidth}{X} 無花果樹的果子漸漸成熟，葡萄樹開花，散發香氣。我的佳偶，起來！我的美人，與我同去！ \end{tabularx} \\ \\ \relax
2:14 & \begin{tabularx}{0.7\textwidth}{X} 我的鴿子啊，你在磐石穴中，在陡巖的隱密處。求你容我得見你的面貌，求你容我得聽你的聲音；因你的聲音悅耳，你的容貌秀美。 \end{tabularx} \\ \\ \relax
2:15 & \begin{tabularx}{0.7\textwidth}{X} 請為我們擒拿狐狸，就是毀壞葡萄園的小狐狸，我們的葡萄正在開花。」 \end{tabularx} \\ \\ \relax
2:16 & \begin{tabularx}{0.7\textwidth}{X} 我的良人屬我，我也屬他，他在百合花中放牧。 \end{tabularx} \\ \\ \relax
2:17 & \begin{tabularx}{0.7\textwidth}{X} 我的良人哪，等到天起涼風、日影飛去的時候，願你歸回，像羚羊，像小鹿，在崎嶇的山上。 \end{tabularx} \\ \\ \relax
3:1 & \begin{tabularx}{0.7\textwidth}{X} 我夜間躺臥在床上，尋找我心所愛的；我尋找他，卻尋不著。 \end{tabularx} \\ \\ \relax
3:2 & \begin{tabularx}{0.7\textwidth}{X} 「我要起來，繞行城中，在街市上，在廣場上，尋找我心所愛的。」我尋找他，卻尋不著。 \end{tabularx} \\ \\ \relax
3:3 & \begin{tabularx}{0.7\textwidth}{X} 城中巡邏的守衛遇見我，「你們看見我心所愛的沒有？」 \end{tabularx} \\ \\ \relax
3:4 & \begin{tabularx}{0.7\textwidth}{X} 我剛離開他們，就遇見我心所愛的。我拉住他，不放他走，領他進入我母親的家，到懷我者的內室。 \end{tabularx} \\ \\ \relax
3:5 & \begin{tabularx}{0.7\textwidth}{X} 耶路撒冷的女子啊，我指著羚羊或田野的母鹿囑咐你們，不要喚醒，不要挑動愛情，等它自發。 \end{tabularx} \\ \\ \relax
3:6 & \begin{tabularx}{0.7\textwidth}{X} 那如煙柱從曠野上來，薰了沒藥、乳香，撲上商人各樣香粉的是誰呢？ \end{tabularx} \\ \\ \relax
3:7 & \begin{tabularx}{0.7\textwidth}{X} 看哪，是所羅門的轎，周圍有六十個勇士，都是以色列中的勇士。 \end{tabularx} \\ \\ \relax
3:8 & \begin{tabularx}{0.7\textwidth}{X} 他們的手都持刀，善於爭戰，各人腰間佩刀，防備夜間恐怖的攻擊。 \end{tabularx} \\ \\ \relax
3:9 & \begin{tabularx}{0.7\textwidth}{X} 所羅門王用黎巴嫩木為自己製作轎子。 \end{tabularx} \\ \\ \relax
3:10 & \begin{tabularx}{0.7\textwidth}{X} 轎柱是用銀做的，轎底是用金做的，坐墊是紫色的，其中所鋪的是耶路撒冷女子的愛情。 \end{tabularx} \\ \\ \relax
3:11 & \begin{tabularx}{0.7\textwidth}{X} 錫安的女子啊，你們要出去觀看所羅門王！他頭戴冠冕，就是在他結婚當天心中喜樂的時候，他母親給他戴上的。 \end{tabularx} \\ \\ \relax
4:1 & \begin{tabularx}{0.7\textwidth}{X} 看哪，我的佳偶，你真美麗！看哪，你真美麗！你的眼睛在面紗後好像鴿子。你的頭髮如同一群山羊，從基列山下來。 \end{tabularx} \\ \\ \relax
4:2 & \begin{tabularx}{0.7\textwidth}{X} 你的牙齒如新剪毛的一群母羊，洗淨之後走上來，它們成對，沒有一顆是單獨的。 \end{tabularx} \\ \\ \relax
4:3 & \begin{tabularx}{0.7\textwidth}{X} 你的唇好像一條朱紅線，你的嘴秀美。你的鬢角在面紗後，如同迸開的石榴。 \end{tabularx} \\ \\ \relax
4:4 & \begin{tabularx}{0.7\textwidth}{X} 你的頸項猶如大衛為收藏軍器而造的高塔，其上懸掛一千個盾牌，都是勇士的盾牌。 \end{tabularx} \\ \\ \relax
4:5 & \begin{tabularx}{0.7\textwidth}{X} 你的兩乳好像百合花中吃草的一對小鹿，是母鹿雙生的。 \end{tabularx} \\ \\ \relax
4:6 & \begin{tabularx}{0.7\textwidth}{X} 我要往沒藥山和乳香岡去，直到天起涼風、日影飛去的時候。 \end{tabularx} \\ \\ \relax
4:7 & \begin{tabularx}{0.7\textwidth}{X} 我的佳偶，你全然美麗，毫無瑕疵！ \end{tabularx} \\ \\ \relax
4:8 & \begin{tabularx}{0.7\textwidth}{X} 我的新娘，請你與我一同離開黎巴嫩，與我一同離開黎巴嫩。從亞瑪拿山巔，從示尼珥，就是黑門山頂，從獅子的洞，從豹子的山往下觀看。 \end{tabularx} \\ \\ \relax
4:9 & \begin{tabularx}{0.7\textwidth}{X} 我的妹子，我的新娘，你奪了我的心。你明眸一瞥，你頸項的鏈子，奪了我的心！ \end{tabularx} \\ \\ \relax
4:10 & \begin{tabularx}{0.7\textwidth}{X} 我的妹子，我的新娘，你的愛情何其美！你的愛情比酒甜美！你膏油的馨香勝過一切香料！ \end{tabularx} \\ \\ \relax
4:11 & \begin{tabularx}{0.7\textwidth}{X} 我的新娘，你的唇滴下蜂蜜，你的舌下有蜜，有奶。你衣服的香氣宛如黎巴嫩的芬芳。 \end{tabularx} \\ \\ \relax
4:12 & \begin{tabularx}{0.7\textwidth}{X} 我的妹子，我的新娘是上鎖的園子，是禁閉的園子，是封閉的泉源。 \end{tabularx} \\ \\ \relax
4:13 & \begin{tabularx}{0.7\textwidth}{X} 你園內所種的結了石榴，有佳美的果子，並鳳仙花與哪噠樹。 \end{tabularx} \\ \\ \relax
4:14 & \begin{tabularx}{0.7\textwidth}{X} 有哪噠和番紅花，香菖蒲和桂樹，並各樣乳香木、沒藥、沉香，與一切上等的香料。 \end{tabularx} \\ \\ \relax
4:15 & \begin{tabularx}{0.7\textwidth}{X} 你是園中的泉，活水的井，是從黎巴嫩湧流而下的溪水。 \end{tabularx} \\ \\ \relax
4:16 & \begin{tabularx}{0.7\textwidth}{X} 北風啊，興起！南風啊，吹來！吹在我的園內，使其中的香氣散發出來。願我的良人進入自己園裡，吃他佳美的果子。 \end{tabularx} \\ \\ \relax
5:1 & \begin{tabularx}{0.7\textwidth}{X} 我的妹子，我的新娘，我進入我的園中，採了我的沒藥和香料，吃了我的蜂房和蜂蜜，喝了我的酒和奶。我的朋友，請吃！我親愛的，請喝，多多地喝！ \end{tabularx} \\ \\ \relax
5:2 & \begin{tabularx}{0.7\textwidth}{X} 我身躺臥，我心卻醒。這是我良人的聲音；他敲門：「我的妹子，我的佳偶，我的鴿子，我完美的人兒，請你為我開門；因我的頭沾滿露水，我的髮被夜露滴濕。」 \end{tabularx} \\ \\ \relax
5:3 & \begin{tabularx}{0.7\textwidth}{X} 我脫了衣裳，怎能再穿上呢？我洗了腳，怎可再弄髒呢？ \end{tabularx} \\ \\ \relax
5:4 & \begin{tabularx}{0.7\textwidth}{X} 我的良人從門縫裡伸進他的手，我便因他動了心。 \end{tabularx} \\ \\ \relax
5:5 & \begin{tabularx}{0.7\textwidth}{X} 我起來，要為我的良人開門。我的兩手滴下沒藥，我的指頭有沒藥汁滴在門閂上。 \end{tabularx} \\ \\ \relax
5:6 & \begin{tabularx}{0.7\textwidth}{X} 我為我的良人開了門，我的良人卻已轉身走了。他說話的時候，我魂不守舍。我尋找他，竟尋不著，我呼叫他，他卻不回答。 \end{tabularx} \\ \\ \relax
5:7 & \begin{tabularx}{0.7\textwidth}{X} 城中巡邏的守衛遇見我，打了我，傷了我，看守城牆的人奪去我的披肩。 \end{tabularx} \\ \\ \relax
5:8 & \begin{tabularx}{0.7\textwidth}{X} 耶路撒冷的女子啊，我囑咐你們：若遇見我的良人，要告訴他，我為愛而生病。 \end{tabularx} \\ \\ \relax
5:9 & \begin{tabularx}{0.7\textwidth}{X} 你這女子中最美麗的，你的良人有甚麼勝過別的良人呢？你的良人有甚麼勝過別的良人，使你這樣囑咐我們？ \end{tabularx} \\ \\ \relax
5:10 & \begin{tabularx}{0.7\textwidth}{X} 我的良人紅潤發亮，超乎萬人之上。 \end{tabularx} \\ \\ \relax
5:11 & \begin{tabularx}{0.7\textwidth}{X} 他的頭像千足的純金，他的髮綹卷曲，黑如烏鴉。 \end{tabularx} \\ \\ \relax
5:12 & \begin{tabularx}{0.7\textwidth}{X} 他的眼如溪水旁的鴿子，沐浴在奶中，安得合式。 \end{tabularx} \\ \\ \relax
5:13 & \begin{tabularx}{0.7\textwidth}{X} 他的兩頰如香花園，如香草臺；他的嘴唇像百合花，滴下沒藥汁。 \end{tabularx} \\ \\ \relax
5:14 & \begin{tabularx}{0.7\textwidth}{X} 他的雙手宛如金條，鑲嵌水蒼玉；他的身體如同雕刻的象牙，周圍鑲嵌藍寶石。 \end{tabularx} \\ \\ \relax
5:15 & \begin{tabularx}{0.7\textwidth}{X} 他的腿好比白玉石柱，安在精金座上；他的容貌如黎巴嫩，佳美如香柏樹。 \end{tabularx} \\ \\ \relax
5:16 & \begin{tabularx}{0.7\textwidth}{X} 他的口甘甜，他全然可愛。耶路撒冷的女子啊，這是我的良人，這是我的朋友。 \end{tabularx} \\ \\ \relax
6:1 & \begin{tabularx}{0.7\textwidth}{X} 你這女子中最美麗的，你的良人往何處去？你的良人轉向何處去了？我們好與你同去尋找他。 \end{tabularx} \\ \\ \relax
6:2 & \begin{tabularx}{0.7\textwidth}{X} 我的良人進入自己園中，到香花園，在園內放牧，採百合花。 \end{tabularx} \\ \\ \relax
6:3 & \begin{tabularx}{0.7\textwidth}{X} 我屬我的良人，我的良人屬我；他在百合花中放牧。 \end{tabularx} \\ \\ \relax
6:4 & \begin{tabularx}{0.7\textwidth}{X} 我的佳偶啊，你美麗如得撒，秀美如耶路撒冷，威武如展開旌旗的軍隊。 \end{tabularx} \\ \\ \relax
6:5 & \begin{tabularx}{0.7\textwidth}{X} 求你轉眼不要看我，因你的眼睛使我慌亂。你的頭髮如同一群山羊，從基列山下來。 \end{tabularx} \\ \\ \relax
6:6 & \begin{tabularx}{0.7\textwidth}{X} 你的牙齒如一群母羊，洗淨之後走上來，它們成對，沒有一顆是單獨的。 \end{tabularx} \\ \\ \relax
6:7 & \begin{tabularx}{0.7\textwidth}{X} 你的鬢角在面紗後，如同迸開的石榴。 \end{tabularx} \\ \\ \relax
6:8 & \begin{tabularx}{0.7\textwidth}{X} 雖有六十王后、八十妃嬪，並有無數的童女。 \end{tabularx} \\ \\ \relax
6:9 & \begin{tabularx}{0.7\textwidth}{X} 她是我獨一的鴿子、我完美的人兒，是她母親獨生的，是生養她的所寵愛的。女子見了都稱她有福，王后妃嬪見了也讚美她。 \end{tabularx} \\ \\ \relax
6:10 & \begin{tabularx}{0.7\textwidth}{X} 那俯視如晨曦、美麗如月亮、皎潔如太陽、威武如展開旌旗軍隊的是誰呢？ \end{tabularx} \\ \\ \relax
6:11 & \begin{tabularx}{0.7\textwidth}{X} 我下到堅果園，要看谷中青翠的植物，要看葡萄可曾發芽，石榴可曾放蕊； \end{tabularx} \\ \\ \relax
6:12 & \begin{tabularx}{0.7\textwidth}{X} 不知不覺，我彷彿坐在我百姓高官的戰車中。 \end{tabularx} \\ \\ \relax
6:13 & \begin{tabularx}{0.7\textwidth}{X} 回來，回來，書拉密的女子；回來，回來，我們要看你。你們為何要觀看書拉密的女子，像觀看兩隊人馬在跳舞呢？ \end{tabularx} \\ \\ \relax
7:1 & \begin{tabularx}{0.7\textwidth}{X} 尊貴的女子啊，你的腳在鞋中何等秀美！你的大腿圓潤，好像美玉，是巧匠的手做成的。 \end{tabularx} \\ \\ \relax
7:2 & \begin{tabularx}{0.7\textwidth}{X} 你的肚臍如圓杯，不缺調和的酒。你的肚子如一堆麥子，周圍有百合花。 \end{tabularx} \\ \\ \relax
7:3 & \begin{tabularx}{0.7\textwidth}{X} 你的兩乳好像一對小鹿，是母鹿雙生的。 \end{tabularx} \\ \\ \relax
7:4 & \begin{tabularx}{0.7\textwidth}{X} 你的頸項如象牙塔，你的眼睛像希實本、巴特‧拉併門旁的水池，你的鼻子彷彿朝向大馬士革的黎巴嫩塔。 \end{tabularx} \\ \\ \relax
7:5 & \begin{tabularx}{0.7\textwidth}{X} 你的頭在你身上好像迦密山 ，你頭上的髮呈紫色，王被這髮綹繫住了。 \end{tabularx} \\ \\ \relax
7:6 & \begin{tabularx}{0.7\textwidth}{X} 我親愛的，喜樂的女子啊，你何等美麗！何等令人喜悅！ \end{tabularx} \\ \\ \relax
7:7 & \begin{tabularx}{0.7\textwidth}{X} 你的身材好像棕樹，你的兩乳如同纍纍的果實。 \end{tabularx} \\ \\ \relax
7:8 & \begin{tabularx}{0.7\textwidth}{X} 我說：我要爬上棕樹，抓住枝子。願你的兩乳好像葡萄纍纍，願你鼻子的香氣如蘋果； \end{tabularx} \\ \\ \relax
7:9 & \begin{tabularx}{0.7\textwidth}{X} 你的上顎如美酒，直流入我良人的口裡，流入沉睡者的口中。 \end{tabularx} \\ \\ \relax
7:10 & \begin{tabularx}{0.7\textwidth}{X} 我屬我的良人，他也戀慕我。 \end{tabularx} \\ \\ \relax
7:11 & \begin{tabularx}{0.7\textwidth}{X} 來吧！我的良人，讓我們往田間去，在村莊住宿。 \end{tabularx} \\ \\ \relax
7:12 & \begin{tabularx}{0.7\textwidth}{X} 早晨讓我們起來往葡萄園去，看葡萄樹發芽沒有，花開了沒有，石榴放蕊沒有，在那裡我要將我的愛情給你。 \end{tabularx} \\ \\ \relax
7:13 & \begin{tabularx}{0.7\textwidth}{X} 曼陀羅草散發香味，在我們的門內有各樣新陳佳美的果子；我的良人，這都是我為你保存的。 \end{tabularx} \\ \\ \relax
8:1 & \begin{tabularx}{0.7\textwidth}{X} 惟願你像我的兄弟，像吃我母親奶的兄弟。我在外頭遇見你就與你親吻，誰也不輕看我。 \end{tabularx} \\ \\ \relax
8:2 & \begin{tabularx}{0.7\textwidth}{X} 我必引導你，領你進入我母親的家，她必教導我，我必使你喝石榴汁釀的香酒。 \end{tabularx} \\ \\ \relax
8:3 & \begin{tabularx}{0.7\textwidth}{X} 他的左手在我頭下，他的右手將我環抱。 \end{tabularx} \\ \\ \relax
8:4 & \begin{tabularx}{0.7\textwidth}{X} 耶路撒冷的女子啊，我囑咐你們，不要喚醒、不要挑動愛情，等它自發。 \end{tabularx} \\ \\ \relax
8:5 & \begin{tabularx}{0.7\textwidth}{X} 那靠著良人從曠野上來的是誰呢？在蘋果樹下，我叫醒了你；在那裡，你母親曾為了生你而陣痛，在那裡，生你的為你陣痛。 \end{tabularx} \\ \\ \relax
8:6 & \begin{tabularx}{0.7\textwidth}{X} 求你將我放在你心上如印記，帶在你臂上如戳記。因為愛情如死之堅強，熱戀如陰間之牢固，所發的光是火焰的光，是極其猛烈的火焰。 \end{tabularx} \\ \\ \relax
8:7 & \begin{tabularx}{0.7\textwidth}{X} 愛情，眾水不能熄滅，江河也不能淹沒。若有人拿家中所有的財寶要換愛情，就全被藐視。 \end{tabularx} \\ \\ \relax
8:8 & \begin{tabularx}{0.7\textwidth}{X} 我們有一小妹，她還沒有乳房，人來提親的日子，我們當為她怎麼辦呢？ \end{tabularx} \\ \\ \relax
8:9 & \begin{tabularx}{0.7\textwidth}{X} 她若是牆，我們要在其上建造銀塔；她若是門，我們要用香柏木板圍護她。 \end{tabularx} \\ \\ \relax
8:10 & \begin{tabularx}{0.7\textwidth}{X} 我是牆，我的兩乳像塔。那時，我在他眼中是找到平安的人。 \end{tabularx} \\ \\ \relax
8:11 & \begin{tabularx}{0.7\textwidth}{X} 所羅門在巴力‧哈們有一葡萄園，他將這葡萄園租給看守的人，每人為其中的果子要交一千銀子。 \end{tabularx} \\ \\ \relax
8:12 & \begin{tabularx}{0.7\textwidth}{X} 我有屬自己的葡萄園。所羅門哪，一千歸你，兩百歸看守果子的人。 \end{tabularx} \\ \\ \relax
8:13 & \begin{tabularx}{0.7\textwidth}{X} 你這住在園中的，同伴都要聽你的聲音，求你使我也得以聽見。 \end{tabularx} \\ \\ \relax
8:14 & \begin{tabularx}{0.7\textwidth}{X} 我的良人哪，求你快來！像羚羊，像小鹿，在香草山上。 \end{tabularx} \\ \\
[1ex]
\hline
\hline
\end{longtable}
我們在過去八天的信息中,已看見基督的各方面.但雅歌書中所表現的基督,遠超過我們所能想能講的.這卷書亦有講及基督是我們的詩歌.是由我們心中所發出來的讚美,以及靈裡所發出的感覺.祂在我們靈裡作主,也在魂及肉體上作主.此書亦講及基督是我們的膏油,故成為我們的醫治,把我們生命中的問題解決.祂分別了我們,使我們有屬靈的權柄和地位.祂也是我們的兄弟,這表示我們與祂同受一個生命之恩.在教會中,祂是我們的頭,在我們作工時祂與我們一起.此書亦講及基督是我們的榮耀.在雅歌書中說:巴不得我的良人像我兄弟一樣.在人面前與我親嘴.這「親嘴」在希伯來人的風俗,是見面時的禮節,這使人看見他們之間的關係,這是他們的榮譽.
\newline
\newline
今天要講的是基督是我們的一切,所以要看整卷書.雅歌書並非是一連串不斷的故事,乃是講及書拉密女,與她的良人中間的一些主要的感覺.這感覺代表了生命之實際,也是我們與主之間的重要關係.在此關係中,可看見基督崇高的地位.既然基督是我們的一切,我們該讓祂在什麼地位上呢?神將祂放在最高的地位.當祂從死裡復活後,升上高天,得著一切應得的榮耀.將來祂再來時,父神要將一切的榮耀交給祂,國度交給祂,所有的人都要稱讚祂,讚美祂.父神既如此高舉祂,那麼,我們這一切跟隨主的人,生活在此世上,更要將主高舉!更應知主是我們的一切.若有此思想和操練,那麼我們的生命與生活,就能豐富起來.
\newline
\newline
基督既是我們的一切,首先我們應該認識主,認識祂的豐富.正如保羅在腓立比書中說要認識基督.在雅歌書中可見到,書拉密女不斷將她所認識的良人告訴我們,也告訴耶路撒冷的眾女子及兄弟們.這是我們今天應有的態度.我們應該認識主,當有人問及我們的主時,請問你如何回答?是否該用溫柔的話語,將主介紹給他們?在這卷書中我們看見幾件事:
\newline
\newline
(一)認識基督
\newline
\newline
Ａ)書拉密女以良人為最寶貴──保羅說:我以基督為至寶.要認識基督,就要像保羅一樣,把萬事當作有損.這裡有一個基本的教訓:人要認識基督,就應將萬事當作有損.若不將萬事當作有損,你對基督就沒有深刻的認識.你若當萬事為有益,就會倚靠物質,仰望別人,那就無法認識基督了.正如幾個瞎子摸大象,他們所知的都是其中一部份,是憑自己的心意和感覺；只憑一點點的知識來下判斷.那只不過是片面的,而不是完全的.我們到主面前,若用我們過去的一點點知識和經歷,就不能看見完全的基督.所以在本書中,新婦要認識新郎,是從多方面去經歷的.這經歷是每次都說及良人,或是良人說及對她的認識時；就照著她的認識在祂身上都發現了.例如當她說我的良人像一袋沒藥,良人也說,在我佳偶裡面有沒藥.可見他們的認識是合在一起的.人若貪愛世界,就沒辦法愛慕主.人是不能貪愛世界,又跟從主.保羅曾說:底馬貪愛現今世界,離開了我.(提後四10)
\newline
\newline
Ｂ)要離開原來的地方──人若抓住世界,就不配認識主基督,人要認識基督,應放下自己的偏見,並要離開了原來的地方,新郎對新婦說:讓我們離開利巴嫩,離開黑門山,往前看,那時就會有新的認識.
\newline
\newline
(二)得著基督──認識基督是好的,但還需得著基督
\newline
\newline
Ａ)尋找──新婦說我的良人你在何處:他說我臥躺在床上,尋找我心所愛的,尋找他卻尋不見.為什麼?因為她在夢中找基督.一個做夢的人是找不到的,特別是作白日夢的人,他不是思想神的話語.若要找著基督,該到祂所在的地方,才可找著.
\newline
\newline
Ｂ)將萬事當作糞土──保羅說:我將萬事當作糞土,為要得著基督.這是得著基督的附帶條件──把萬事當作糞土──把萬事看為有損還不夠,那只不過是放下觀念來認識基督,若要真正得著基督時,就要將萬事看作糞土.(歌二13)說:「無花果樹的果子漸漸成熟,葡萄樹開花放香；我的佳偶,我的美人,起來!與我同去.」這是說我們一同去吧.又(歌四7-8)節說:「我的佳偶,他全然美麗,毫無瑕疵.我的新婦,求你與我一同離開利巴嫩；從亞嗎哪頂,從示尼珥與黑門頂,從有獅子的洞,從有子豹的山,往下觀看.」我們應離開這一切境地,不可再留戀,世上一切的東西,都能將我們愛主的心奪去.現在不單是將萬事看為有損,更要當作糞土,這一切再不能吸引我們時,才可往前走.正如本書的一章中說:願你吸引我,我就快跑跟隨你.因世上一切若不放下,就不但不能愛主；且漸漸會變成小狐狸,把正在開花的葡萄樹也毀壞了.弟兄姊妹們,我們不但觀念要正確,更要去實行.所以有時神在我們身上作工,把我們最寶貴的東西取去,這是神要我們學一個功課.因這一切的東西不過是糞土,轉眼間都會過去.世界的一切,虛空的虛空,因為日光之下都虛空的,惟有日光之上的那一位主,才是真實的寶貴的.
\newline
\newline
Ｃ)捨棄自己──我們要得著基督,還在讓基督得著我們.甚麼時候我們能夠將萬事捨棄,甚麼時候基督就可以得著我.這雖然是不易明白,但聖經上是如此說:這位新婦失去了新郎,就去尋找,結果被城巡邏看守的人打傷了,奪去了披肩.但當她回到良人面前時,就立刻得著了良人.所以甚麼時候我們到主面前,甚麼時候我們就可以得著主,主也得著了我們,這樣才有密切的關係.感謝主!我們得著主,就要將自己放下,要捨棄自己；才能得著主,這是屬靈的秘訣.
\newline
\newline
(三)效法基督
\newline
\newline
在本書四章中,良人對女子所說的話和稱讚,其中基本的條件,也是我們對基督認識的條件.「我以我的良人為一袋沒藥,祂是鳳仙花,祂是沙崙的玫瑰花,是谷中的百合花.我在他的蘋果樹下,喫他的果子,我心中就有力量,我就歡喜.」當良人對女子說話時,他說:我進到你的園中,因在你的園子裡,有鳳仙花,有百合花,有沒藥又有乳香,有各樣的果子.這是什麼意思?這新婦裡面有了基督的一切,效法基督,進入到一個關係裡面.有了良人的性情；他裡面的性情,也就是良人所有的,這是我們應走的道路.一個正常的基督徒,乃是不斷的效法基督.我們得著主之後,應照著這樣去做.不能只是放在這一邊,而天天欣賞,這是不夠的.應照祂路上所留下的腳蹤向前走,祂溫柔,慈愛,憐憫,那我們也就該如此.他怎樣的施捨,我們也要紀念主的話,施比受更為有福.主如何走道路,我們也如何去走這道路.直等到一個時候,我們默想主,效法祂的作為和祂的榜樣；那麼基督就成形在我們裡面.我們要思想基督,天天默想基督,你的言語,動作,思想也都漸漸會像祂了.
\newline
\newline
研究近代人種學,語言學,心理學,社會學都會提及一件事,就是幾十年前在印度發現一對狼人.有兄弟兩人不知為何會在豺狼中,這狼就用奶餵養這兩小孩.及至被人捉到後,就觀察他們,見他們走路像狼一樣,且走得很快.他們之間的話沒人能明白,可惜這兩狼人不久就死去了.從這事上看來,他們與狼在一起,雖然他們的樣子還像人,但動作,語言及思想就像狼.因此,若人脫離世上一切卑賤的事,及容易纏累我們的罪,天天與主同在；常思念祂的話,默想祂的愛,我們也會像祂.盼望我們在生活中,禱告裡,思想上,都多思念主.祂說:「清心的人有福了.」又說:「凡勞苦擔重擔的人,可以到我這裡來,我就使你們得安息.」「我就是世界的光.」「人渴了可以到我這裡來.」「我的話就是靈,就是生命.」,我們把這些話語,深深印進了心坎裡.在新婦的心中是有沒藥,乳香,鳳仙花,百合花；這一切在良人裡面,也在我們裡面.祂的一切成為我們的一切,祂的豐滿成為我們的豐滿.
\newline
\newline
(四)活著就是基督
\newline
\newline
在(腓一21)保羅說:「因為我活著就是基督,我死了就有益處」.「我已經與基督同釘十字架,現在活著的,不再是我,乃是基督在我裡面活著.」(加二20)
\newline
\newline
我們不單是得著基督,效法基督,還要活著就是基督.別人看見我們時能說:「啊!基督在我們中間行走.」拿撒勒的基督,今天行在我的家中,當年走上各各他山上的基督,今天行走在教會裡.我有了這背負十字架的決心,我死了就有益處.為主的緣故,無論如何要讓人看見我就是基督.基督已升到天上,坐在父神的右邊；在地上代表祂的就是你和我,你我就是活在地上的基督.新婦對良人說:「讓我們去吧!讓我們在村莊住宿,可以到田間去或到園裡去.然後將我的愛情給你；因這一切在我裡面,有新陳佳美的果子.我的良人,這都是我為你所存留的.」基督的一切,既存留在我們裡面,那麼我們活著就是基督了.什麼是我們應不斷操練的呢?從今天開始直到見主面時,這一切都是該天天操練的.要天天認識基督,得著基督,效法基督,活著就是基督.那麼,當主坐筵席時,基督再來時,在婚筵中,就可以與他一坐席.然後祂真的可以對我們說:「我的佳偶,你全然美麗,是完全的人,是毫無瑕疵.」若能如此的操練,並得主的稱許,那麼我們就可以說:「我的良人哪,求你快來,如羚羊或小鹿在香草山上.」
\newline
\newline


\newpage

\section{第46屆~港九培靈會~周主培~第一講}
\label{sec:1208}
\textbf{你們尋找甚麼}
\newline
\newline
連結: \href{https://www.hkbibleconference.org/session-message/view/1208}{\texttt{https://www.hkbibleconference.org/session-message/view/1208}}
\newline
\newline
\hyperref[sec:1207]{< < < PREV SERMON < < <}
~
\hyperref[sec:toc]{[返目錄]}
~
\hyperref[sec:1210]{> > > NEXT SERMON > > >}
\newline
\newline
約翰福音 1:35-1:39
\newline
\begin{longtable}{cl}
\hline
\hline
章節 & 經文 (和合本修訂版)\\
\hline
1:35 & \begin{tabularx}{0.7\textwidth}{X} 又過了一天，約翰同兩個門徒站在那裡。 \end{tabularx} \\ \\ \relax
1:36 & \begin{tabularx}{0.7\textwidth}{X} 他見耶穌走過，就說：「看哪，神的羔羊！」 \end{tabularx} \\ \\ \relax
1:37 & \begin{tabularx}{0.7\textwidth}{X} 兩個門徒聽見他的話，就跟從了耶穌。 \end{tabularx} \\ \\ \relax
1:38 & \begin{tabularx}{0.7\textwidth}{X} 耶穌轉過身來，看見他們跟著，就問他們說：「你們要甚麼？」他們對他說：「拉比，你在哪裡住？」（「拉比」翻出來就是老師。） \end{tabularx} \\ \\ \relax
1:39 & \begin{tabularx}{0.7\textwidth}{X} 耶穌說：「你們來看。」他們就去看他在哪裡住。這一天他們就跟他同住；那時大約是下午四點鐘。 \end{tabularx} \\ \\
[1ex]
\hline
\hline
\end{longtable}
非常感謝神的恩典,經過三年,我又回到香港來；不但看見許多認識的弟兄姊妹,也見到不少初次見面的.感謝大會籌備的主席和同工們,給我機會為神傳話,特別是跟隨鮑會園牧師,和于力工牧師二位兄長,我實在從他們得很多的福份.我慎重要求各位,為我們在神面前禱告,同時也盼望各位在可能範圍內,千萬不要錯過早堂的聚會；萬一不能赴會,也請訂購該堂道的錄音帶,(我自己已經訂了)盼望藉此能夠幫助你.
\newline
\newline
耶穌問跟隨祂的兩個門徒說:「你們要什麼?」另一翻譯是:「你們尋找什麼?」這節聖經,是主耶穌在世上時提出的第一個問題.巴不得今晚第一次的聚會,主也在我們中間,向著我們發出這個問題.
\newline
\newline
世上的人,天天在尋找.不知找什麼?中國長江一帶,有許多的山；有一次,有兩個人站在山上觀看,看見長江上游船隻絡繹.其中一位較長有哲學思想者說:「那許多來來往往的船,非為名,則為利.」在香港,路上行人,也忙忙碌碌,有尋找學問的,也有尋找職業的,有尋找對象的,有尋找出路的,有爭名奪利的.九龍尖沙咀一個很高的建築物,上面懸著「名成利就食總督」的大廣告牌；這就是人心所想的,(幸虧現在是民主時代,言論自由,若是專制時代,誰敢連總督也吃了呢?)人都盼名成利就,但知識和財寶,並不能真正滿足人心.很多人極力尋找刺激；有人說,香港失業的人多,色情的事多,搶劫的事多；但據官方公佈,香港國泰民安,百姓安居樂業.據今年五月份工會的報告,香港有十二萬人半失業和失業,若一人代表四口之家.那麼,約近五十萬人處於生活困難之中.
\newline
\newline
狄更斯所寫的雙城記,開頭說:「這是最好,也是最壞的時代,是希望的春天,也是絕望的冬天；在我們前面樣樣都有,也是一無所有.」香港是東方明珠,是最美麗之地,也是最醜惡之地；有許多人如生活在天堂,也有千千萬萬人,如同過著地獄的生活.有些人覺得香港很熱鬧,但有人覺得很孤單,香港四面環水,但許多人卻覺得枯乾煩悶.從飛機上俯視,香港是綠洲?但卻如沙漠的乾旱,香港到處高樓大廈,有水坭森林之稱；在光天化日之下,卻有許多昏天黑地之處.在這樣光景之下,耶穌問:「你們尋找什麼?」
\newline
\newline
今晚,我們靠主恩典來思想,從約翰福音一至九章,我們看見人一步步地都在尋找,各種不同的人,都有其不同的需要.人心有很多渴望,約翰一章我們看到神的羔羊.約翰二章,耶穌以水變酒.第三章,有位高級的人物,夜裡見耶穌.第四章,犯罪的女人來見耶穌,在正午時,那撒瑪利亞婦人,和耶穌談論活水的問題.第五章,有個生病卅八年的人,躺臥在那裡.第六章,見到有多少飢餓的人,耶穌以五餅二魚給他們吃飽.第七章,耶穌站著高聲說:「人若渴了,可以到我這裡來喝.信我的人,就如經上所說,從他腹中要流出活水的江河來.」第八章,有個犯姦淫的婦人被捕,人要用石頭打死她.第九章,有個生來就瞎眼的人,他從未見過藍的天和白的雲,紅的玫瑰和綠的葉子；最可憐的,他沒有見過自己母親的臉,只能用手摸他慈母的臉.有個瞎眼的女孩,當她母親給她一個洋娃娃時,她一接過來,就摸了洋娃娃的眼睛；問她母親說:「洋娃娃的眼睛看得見嗎?」她自己看不見,巴不得她的洋娃娃能看得見.
\newline
\newline
人生的道路在乎自己揀選,如果眼看不見,只能接受他人所為你揀選的.一個瞎眼的人,人家領他,他就走；擺上食物在他的面前他就吃,給他衣服他就穿,沒有法子選.
\newline
\newline
親愛的同工們,親愛的弟兄姊妹!我們受了傳統的製作,別自以為有思想,其實你的思想,不過是別人灌輸的,你的思想都是前人所思想過的.有人說:「我不相信耶穌.」有人說:「聖經是錯的.」但你若問他有看過聖經嗎?他說:「稍微看過一點.」很多事情都是別人的思想放在你裡面.哦!今天有許許多多的人如瞎子一樣,黑白不分,是非不明；只講論左右的問題,忘記了上下的問題.我們並非單講向左,或向右的問題,乃是講人生出路在那裡?中國人有句話說:「不見盧山真面目,惟願生在此山中.」人雖身居其中,卻看不見盧山的真面目.很多時候,我們以局部的事代替整體,以有限的,衡量那無限的；而用暫時的,來思想永恆的；用物質來代替那永恆的生命.我相信各位若讀過那種書籍,就知道能力和物質可以交換,可「能」變成「物」,「物」可變成「能」.或說是無中生有,抽象變成實際,正如道教所說的無極到太極.其實,我們從聖經清楚看見；道成了肉身,住在我們中間,主耶穌基督將神顯明出來.
\newline
\newline
一個瞎眼的人,實在可憐!還把他當問題來研究.約翰第九章告訴我們,當門徒看見那瞎眼的人；他們就把他當作神學的問題,當作社會的問題來思想；研究到底是誰犯了罪,是這人呢?是他父母呢?我們常把人當作問題,說,這個問題兒童,那是個問題青年.耶穌卻不是這樣,祂沒有把人當問題來研究,乃把人當作人.耶穌說:「不是這人犯了罪,也不是他父母犯了罪；是要在他身上顯出神的作為來.」耶穌就用唾沫和泥,抹在瞎子的眼睛之上,吩咐他往西羅亞池子去玩,他就憑著信心去洗.什麼叫做信心?就是有客觀的事實,也有主觀的經驗.那瞎子聽了耶穌的話去洗,回頭就看見了.請問,他開了眼睛,第一先看到誰?他看見自己.世界上最難認識的就是自己,我們常看到別人,很少看到自己.弟兄姊妹!神把祂的光照在我們身上,叫我們先看見自己；在神面前,我們不再看見別人的短處,乃是先看見自己是何等樣的人.而以賽亞在神的面前,他看見了自己.約伯說:「我從前風聞有你,現在親眼看見你.」(伯四二5)因此,約伯在爐灰塵土中悔懊自己.彼得看見耶穌的時候,說:「主啊!請離開我,因我是個罪人.」主耶穌的道,主耶穌的光,叫我們看見自己,懂得自己.在神的光中,我們看見自己是罪人；當人認識自己是罪人的時候,他就會走上得救的道路.
\newline
\newline
現在讓我們再回到約翰福音一至九章裡面,請各位回家以後,細讀這幾章聖經.第八章記載一個正在犯姦淫而被捕的女人.第五章記載一個臥病卅八年的人.第四章記載一個大名鼎鼎的撒瑪利亞婦人.一題到她,我們知道她曾經有過五個丈夫.我覺得以上所題三個人,都很可憐,毫無疑問的,這三個都是罪人；但我們在第八章看見耶穌對犯姦淫而被捕的女人說:「我也不定你的罪,去罷!從此不要再犯罪.」在第四章我們看到撒瑪利亞婦人,當耶穌和她談話,門徒覺得很希奇,為甚麼耶穌和一個有罪的女人談話呢?至於第五章那臥病卅八年的人,我不知他患的何病,也不知他因何患病.但聖經告訴我們,耶穌對他說:「你已經痊愈了,不要再犯罪.」由此可見那人的病,是與罪有關的.第四章和第八章,很清楚給我們看見,有兩個犯罪的女人；但我卻不知道與她們有關的那些男人在那裡?這個世界,實在不太公平,容易錯怪人,將罪歸咎於女人,說女人是禍水；但我要問,當時那撒瑪利亞婦人的五個丈夫到那裡去呢?難道是她一人之錯嗎?那正在犯姦淫的女人,當時那男人那裡去呢?這個世界,我們所看到的,就是強大欺侮弱小.這個世界!偷一元有罪,偷百萬卻是富豪；殺一人坐監牢,殺萬人卻做將軍.這世界是不公平的世界!但主耶穌到世上來,祂體恤人的痛苦.
\newline
\newline
一個有罪的人,正像上述三個有罪的人,內心不平安,沒有喜樂；平安和喜樂是連在一起的.有罪的人,失去了自由.那撒瑪利亞婦人,她為避免見人,只好在正午──最炎熱無人打水的時候去打水.那臥病卅八年的人,能看,能聽,能說,但卻不能動,因他沒有了自由.許多人正像是這樣,被罪所捆綁；是真「夜裡思想天橋路,早晨起來磨豆腐,仍舊走老路.」他雖然思想很多,卻沒有力量去行.當人犯了罪,羞恥就進入；何處有罪,必定有羞恥,沒有平安.各位尚未接受耶穌的朋友們!你心裡有平安嗎?或者你家掛著「出入平安」的字句,但是「出」沒有平安,「入」也沒有平安；在外沒有平安,「入」也沒有平安；在外有麻煩,在家有爭吵.
\newline
\newline
我們每一個人都受遺傳和傳統的束縛,內心沒有平安；聖經告訴我們,除非與神和好,否則就得不著神的平安.如果以英文譯述意即:「必須與神和好,才能得著屬乎神之內的平安.」許多人雖然不承認自己有罪,但知道自己內心沒有平安.我們看見上述三人,一是罪人,一是犯人,一是病人.在這世界上,多少的犯人,多少的罪人,多少的病人,多少的死人；死人在墳墓裡,病人在醫院裡,犯人在監獄裡,可是罪人卻滿街跑.
\newline
\newline
當耶穌問門徒說:「你們要什麼?」他們卻回答說:「拉比,在那裡住?」真是答非所問.真奇妙!倆門徒為何跟隨耶穌!
\newline
\newline
第一次約翰說:「看啊!神的羔羊,除去世人罪孽的.」約翰說這話時,他想到了人的罪,這使我們想到耶穌能為我們作什麼.次日,約翰看見耶穌,就說:「看啊!這是神的羔羊.」約翰說到這裡就停住了.常常有人在未接受耶穌之前,開頭就問:「信了耶穌有何好處?」聖經說:「神的羔羊,除去世人罪孽的.」信了耶穌,罪得赦免.保羅說:「我在你們中間,我不知道別的,只知耶穌基督,並他釘十字架.」因祂釘十字架,我們的罪得了赦免.當我們第二次見到耶穌的時候,我們不是單講耶穌對我們有何好處,乃是因耶穌自己而愛耶穌.很多人信了耶穌,沒有認識清楚,僅以耶穌為救主,拯救我,赦免我的罪,那就夠了.我曾在日本工作三年,日本的青年非常渴慕和追求,他們聽道,在參加慕道班,毫無間斷；可是一直到了受水禮就告完畢.他們以為已經懂了,夠了,畢業了；沒有與耶穌自己發生關係.我們所想到的,不是單單耶穌是個大醫生,醫治我的病；不單因耶穌基督給我心裡有平安,乃是因真正的愛耶穌而跟隨祂,因為耶穌基督實在太奇妙.所以約翰第二次沒有提到人的罪孽,他只說:「看啊!神的羔羊,」 倆門徒就跟隨耶穌.
\newline
\newline
親愛的朋友們!親愛的同道們!或者你已經接受耶穌,你有否真正的讓耶穌在你生命中得著應得的地位?
\newline
\newline
我小時候,很喜歡聽趙世光牧師講道,我一直記得他曾講過,人生的一切就是○,○+○=○,無論多少個○,結果零,財產學問都等於零一樣,如果信了耶穌,就如零的前面加上一,三個○前面加上一就是1000,越多的前面加上一就數目越大.有了耶穌,你的生命就滿足了,就豐富了.最近,我自己在想,很多人有了耶穌,生命還是很空虛.今天在我們中間,已經有了耶穌的人,你是否還覺得空虛?是否你還在迷茫當中呢?原因在於你書寫的方法錯了.不錯,零加上一就有價值,但問題是你把一放在前面或後面.很多人剛到香港時,常禱告,讀經,聚會；現在房子,妻子,車子,金子,樣樣都有了,漸漸的把耶穌丟到後面去了.雖然你有耶穌,但所佔的地位很低,成分微乎其微.哦!我的心恐懼戰兢,我擔心每年培靈會赴會的人很多,等到散會後,各人還是老樣子.很多基督徒在香港,簡直如球掉在棉花堆中,不會再跳起來,有的人麻木了,有的人安逸的在那裡做夢.弟兄姊妹!你們究竟要什麼?尋找什麼?不單未信耶穌的人需要看見,基督徒也需要看見；沒有異象,民就滅亡.求你開我的眼睛,使我看見你律法中的奇妙.我們需要有屬靈的眼光.亞伯蘭和羅得,羅得看見所多瑪,摩蛾拉,就漸漸挪移帳棚,跟隨世界去了.
\newline
\newline
今天你要誰?你要什麼?當耶穌問你,你將如何回答?
\newline
\newline


\newpage

\section{第46屆~港九培靈會~周主培~第二講}
\label{sec:1210}
\textbf{耶穌基督}
\newline
\newline
連結: \href{https://www.hkbibleconference.org/session-message/view/1210}{\texttt{https://www.hkbibleconference.org/session-message/view/1210}}
\newline
\newline
\hyperref[sec:1208]{< < < PREV SERMON < < <}
~
\hyperref[sec:toc]{[返目錄]}
~
\hyperref[sec:1212]{> > > NEXT SERMON > > >}
\newline
\newline
哥林多前書 2:1-2:5
\newline
\begin{longtable}{cl}
\hline
\hline
章節 & 經文 (和合本修訂版)\\
\hline
2:1 & \begin{tabularx}{0.7\textwidth}{X} 弟兄們，從前我到你們那裡去，並沒有用高言大智對你們宣講神的奧祕。 \end{tabularx} \\ \\ \relax
2:2 & \begin{tabularx}{0.7\textwidth}{X} 因為我曾定了主意，在你們中間不知道別的，只知道耶穌基督並他釘十字架。 \end{tabularx} \\ \\ \relax
2:3 & \begin{tabularx}{0.7\textwidth}{X} 我在你們那裡時，又軟弱，又懼怕，又戰戰兢兢。 \end{tabularx} \\ \\ \relax
2:4 & \begin{tabularx}{0.7\textwidth}{X} 我說的話、講的道不是用委婉智慧的言語，而是以聖靈的大能來證明， \end{tabularx} \\ \\ \relax
2:5 & \begin{tabularx}{0.7\textwidth}{X} 為要使你們的信不靠著人的智慧，而是靠著神的大能。 \end{tabularx} \\ \\
[1ex]
\hline
\hline
\end{longtable}
兩年前,當我到東非洲工作時,回程中,航空公司為我安排有機會在羅馬停留幾天.我到希臘,盼望看看保羅當時工作的路線.我開始到哥林多,這是很重要的港口,現在相當繁華；保羅當日工作的地點,就在山丘的上面.我上到那裡,看見儘是斷牆殘壁,人事已非；因此,我們想到金恩吉爾教授的話,他說:「如果你跟這個時代的文化,與時代的文明結婚的話,你的下一代就成為寡婦.」意思就是一個傳道人,為要迎合人的心理,而講時代所需要的,人心所渴望的；以致忽略了最重要的信息,他不知不覺就失落了.雖然哥林多的舊城已成過去,我還記得保羅告訴哥林多的人說:「弟兄們!從前我到你們那裡去,又軟弱,又懼怕,又甚戰兢.我並沒有用高言大智,對你們宣傳神的奧秘.因為我曾定了主意,在你們中間不知道別的,只知道耶穌基督,並祂釘十字架.」保羅是位有學問,有智慧,有口才的人,但他並無用高言大智,講述神的奧秘；因為人的信不乎人的智慧,乃在乎神的大能.
\newline
\newline
我離開了希臘即往羅馬──永恆的城市,在此到處看見歷史.第一處我到 Catacomb,那是墳墓的地方,和普通墳墓不同.那裡的墳墓好像蜘蛛網,佈滿在地底下,其間有路可行,但是走來走去,還是在墳墓裡.據說.當時基督徒被迫害,一班基督徒在墳墓裡聚會,他們晚間點著臘燭,在那裡,說「祂已經復活了.」附近有個小教會名苦巴跌.據傳說,當時彼得在羅馬工作,(但有的人說彼得從未到過羅馬)那時基督教被迫害；有人勸彼得離開羅馬城為妙,於是他帶了一個門徒,在晨光稀微的時候,暗暗地要離開羅馬.忽然有一道大光向彼得照來,在光中彼得看見了主耶穌,彼得跪下說:「主啊!你往哪裡去?」彼得聽見有聲音對他說:「我要到羅馬去,因為我的門徒和兒子們在那裡被迫害.」彼得說:「主啊!你不要去,讓我回去.」說著,就起來,面向羅馬.他的門徒不知此事,問說:「主啊你往哪裡去?」彼得說:「回羅馬城去.」所以那教堂裡,掛著兩幅巨畫一是耶穌釘十字架,一是彼得倒釘十字架.據說當時彼得被捕,人要把他釘十字架,彼得說:「我不配像我的主那樣釘十字架,因為我曾三次否認主,你們應把我倒頭來釘罷.」
\newline
\newline
當時羅馬帝國軍力強盛,在中亞細亞和小亞細亞一帶,有很多羅馬的軍隊,保羅對羅馬人說:「我不以福音為恥,因為這福音本是神的大能,要救一切相信祂的人.」保羅膽敢對希臘人和羅馬人這樣地宣告,是因在他生命中有件最重要的事,就是──「耶穌基督並祂釘十字架.」這不但是他客觀上的信仰,也是他主觀的經驗；他面對面看見了耶穌,親耳聽見了耶穌的聲音,保羅不能不講.
\newline
\newline
今晚,我靠著主的恩典,和大家講一件最古老切實的事情,也就是上屆培靈會的主題；這個主題是魔鬼最不願我們講的.這個主題,曾經叫彼得因祂而三千人得救；這個主題,叫腓力曠野的路上,對太監講耶穌基督的事.
\newline
\newline
有篇著名的文章,告訴我們那無比的耶穌.雖然耶穌在世上時,從未唱過一首詩,但今天世界上有千千萬萬的詩,為歌頌祂而唱.雖然耶穌在世界上沒有留下隻字,但卻有千千萬萬的書本,是為祂而寫.耶穌曾當過木匠,但我們不知道祂曾作了何物,卻有千千萬萬的建築物是為耶穌的名而造.耶穌未曾訓練過一支軍隊,雖然十二門徒中曾有個賣祂的,也有個不承認祂的；但卻有千千萬萬的人,願靠著耶穌而活,也願為耶穌而死.這位耶穌,是我們所歌頌,是我們所傳揚的.
\newline
\newline
在我們中間,不知有那位,相信上帝而不相信耶穌?曾有青年人對我提出這個問題.說:「我相信有位創造天地的神,但是我不相信耶穌是我們的救主.」親愛的弟兄姊妹們!下午我臨時改變題目,我的心如同被火焚燒,我感到今晚必須講耶穌基督,並祂釘十字架.
\newline
\newline
兄弟的原籍是浙江紹興,是產酒勝地,如果各位不誤會我的意思；對我來講,傳耶穌是我的酒杯,不是我的飯碗.(一個嗜酒的人,可以不吃飯,而不能不喝酒的.)
\newline
\newline
司布真是位著名的講道家.有一次,當他講完道之後,一位聽眾問朋友說:「你覺得剛才所講的怎樣?」答:「司布真的人怎樣,我不知道；但我覺得他講的耶穌,非常有道理.」司布真說:「我現在不是勸你們加入教會,我乃是要求你們認識,我的主耶穌基督.」當彼得和約翰被迫害時,他說:「我們所知道的耶穌,我們不能不說.」我再次說,我所信的耶穌,不單是我父母我牧師的耶穌；而且是我個人和祂有關係的.靠主恩典,我認識祂,接受祂,跟隨祂,服事祂；如果有一天神願意的話,我願為祂而死.祂拯救了我,赦免了我的罪,醫治了我的病,祂救我的命脫離死亡.希斯尼爾曾經說；傳耶穌,就好像一個乞丐發現了某地施粥,立即告知其他的乞丐去取粥.我已經接受了耶穌,我再說:「傳耶穌不是我的飯碗,而是我的酒杯.」耶穌基督在我心靈裡,比世人更美.主啊!你的嘴裡滿有恩惠.除你以外在天上我還有誰呢?除你以外在地上我也沒有所愛的.你的名如同倒出來的香膏,所以眾童女就快跑跟隨你,女子啊!你要聽,你要想,不要記念你的名和你的父家,王就羡慕你的美貌.哦!我主耶穌基督,祂全然可愛.
\newline
\newline
有人說:「條條大路通羅馬.」我們可以說:「全本聖經66卷,卷卷講基督.」 全本聖經的主題是「耶穌基督.」創世記講耶穌是女人的後裔.出埃及記講耶穌是被殺的羔羊.利未記講耶穌是贖罪祭.民數記講耶穌是被擊碎的磐石.申命記講耶穌是先知.約書亞記講耶穌是軍隊的元帥.約伯記講我知道我的救贖主是誰.詩篇講耶穌是我們一切的一切.箴言講耶穌是我們的智慧.雅歌講耶穌是我們的良人,是我們是牧人.所有的先知書告訴我們,耶穌是和平的君王.四福音論耶穌是降世救人,復活的人子.使徒行傳告訴我們復活的耶穌基督.保羅,彼得,約翰的書信告訴我們,耶穌是永遠活著的主.啟示錄告訴我們那位再來的耶穌.馬太一章告訴我們,要給祂起名叫耶穌；因祂要將自己的百姓從罪惡裡救出來.使徒行傳(四12)告訴我們,除祂以外別無拯救；因為在天下人間,沒有賜下別的名,我們可以靠著得救.
\newline
\newline
我們在這世界上,充滿了許多的難題.佛教的思想說,人的難題生,老,病,死.聖經給我們看見,人的難題是罪惡的問題 ,我們看見人有罪,有憂愁,有死亡.
\newline
\newline
我在美國現住的是個很小的村莊,叫做樂園；就是蔡蘇娟小姐的地方.十幾年前,我剛到美國,有人問我要否到 Paradise 去?我一想,這是樂園,當然要去,而且急不及待想去.到了那裡,他帶我去見蔡蘇娟小姐；我發覺那裡有墳墓,也有醫生,覺得很奇怪.說,這不是樂園吧!樂園哪裡有墳墓有醫生呢?在這世界上,我們有很多問題.今天社會講福利,可是整個社會沒有安全.研究兒童心理兒童教育,可是我們看見不良少年何其多!香港很多新建的房屋,裡面卻充滿著許多破碎的家庭.男女講自由戀愛.可是離婚案卻與日俱增.各國致力和平談判.但一方面卻在準備武器.這個世界!犯罪的事觸目皆是.空氣,水,土地都有污染的問題.有一次,一位研究環境學的教授對我說:「污染的事,都是因人的罪行所造成.」我覺得很奇怪,罪的問題,本來是傳道人講的,如今連研究環境學的人也講人的罪.記得我少時,最欣賞 Blue Danube 一曲,藍色多腦河畔,多麼令人嚮往!那知現在該處,變成了骯髒的陰溝；因為工業污水都導向那裡流去.
\newline
\newline
我們所注意的,不應該是用何方法可以救這世界,我們所應當注意的,乃是誰才能救這世界?就是耶穌基督.許多人在研究,在思想,發現人的生命何等蒼白,空虛,幻滅,微小,人迷失了,孤獨,沉淪,荒謬,墮落.我們思想到存在主義,物質主義,人文主義,無神主義；問題研究了,但解決的方法在那裡?我來香港之前,先到越南的西貢,下了飛機,我首先向接機的人道賀說:「恭喜你們,現在沒有聽見炮聲了.」我旁邊的一位,黯然地指著報紙給我看「越南仍然有戰爭」,表面上已經和平,實際戰爭還未根本解決.這個世界好像將爆的火山,將沉的破船,垂死的病人一樣,用何方法來拯救這世界呢?聖經告訴我們,「惟有基督在我們還作罪人的時候為我們死,神的愛就在此向我們顯明了.」(約廿31):「但記這些事,要叫你們信耶穌是基督,是神的兒子,並且叫你們信了祂,可以因祂的名得生命.」保羅見證說:「基督耶穌降世,為要拯救罪人,在罪人中我是罪魁,然而我蒙了憐憫.」.
\newline
\newline
耶穌基督才能解決人的問題.耶穌到世上來不是單來傳福音.祂的本身就是福音,有了耶穌,我們才有福音.我們向失落的人說,人子來為要尋找拯救失喪的人；對存在主義的人說,耶穌基督昨日,今日,一直到永遠是一樣的；對荒謬虛無的人說,耶穌說:「我就是道路真理生命,若不藉著我,沒有人能到父那裡去.」對注重物質的人說:人活著,不是單靠食物,乃是靠神口裡所出的一切話.萬物是祂造的,凡被造的,沒有一樣不是藉著祂造的.對無神論者說,從來沒有人看見神,惟有父懷中的獨生子,將祂表明出來.哦!耶穌基督是我們一切問題的解答.最近有首很出名為青年人愛唱的英文詩歌:「沒有路就無處可去,沒有真理,就無從知道,沒有生命,就不能生存,簡直就不能活下去.」沒有耶穌基督,就沒有道路,我們就不知何往；沒有耶穌,世界就沒有真理,我們就不能明白.沒有耶穌,就沒有生命；我們就不能生活,簡直無活下去.
\newline
\newline
(約一1):「太初有道,道與神同在,道就是神.」這裡第一給我們看見的──耶穌是道.我認為中譯聖經勝過任何文字的翻譯.據英譯: In [the] beginning was the word 話,希臘文原意是:(Logical)中譯為合邏輯,意即合理,是經過頭腦思想的話,有條理有意思的話.思想看不見的東西,但從話表達出來.「從來沒有人看見神,只有在父懷裡的獨生子,將祂表明出來.」耶穌是神的話,耶穌將神表明出來.「道」字,中譯之妙,非任何國文字可比,道就是道理,道是路.老子的道德經說:「天生一,一生二,二生三,三生萬物.」聖經裡耶穌說:「我是道路,真理,生命,若不藉著我,沒有人能到父那裡去.」(約一1-5):「太初有道,道與神同在,道就是神,這道太初與神同在.萬物是藉著祂造的,凡被造的,沒有一樣不是藉著祂造的.生命在祂裡頭；這生命就是人的光.光照在黑暗裡,黑暗卻不接受光.」耶穌不但是道,且是創造天地萬物之主,我們不過是被造的人.
\newline
\newline
很多人不承認有神,不相信有耶穌.說:「因我沒有見過,」不能證明.蘇聯的太空人,到太空繞了一圈,居然大言不慚說:「那裡有上帝?我到太空根本沒有看見上帝.」葛培理說:「有一條蚯蚓從地裡竄出來,看了很久,說:我沒有看見赫魯曉夫,大概世上無此人吧!」這是何等愚蠢的話!世上雖看見天地萬物,卻不承認天地萬物是神創造的.
\newline
\newline
耶穌不但是生命,也是生命的光,是道成肉身的主；是神獨生的兒子,是神的羔羊──除去世人的罪孽的.
\newline
\newline
剛才我們講到世界上各種的問題,困人有罪之故,何處有罪,那處就沒有平安.心裡有罪,心就沒有平安.家裡有罪,家就沒有平安.社會有罪,社會就沒有平安.國家有罪,國家就沒有平安.世界有罪,整個世界充滿著死亡.怎能除去罪呢?這是很重要的問題.人想盡方法,以為科學發達能叫人類文明進步；但科學帶來空氣污染,水之污染.世界將到萬劫不復的地步!
\newline
\newline
用何方法除去人的罪呢?聖經說:「看那!神的羔羊,除去世人罪孽的.」許多人以為行善可以得救,可上天堂.如果人因罪坐牢,雖然在獄中表現良好,或許可以提早出獄；但不能因他現在的好,而把從前所犯的罪抵消.行善是人的本份.難道去作惡嗎?比如在路上開車,看見交通紅燈即停車；這並不是說明你是駕車能手,而是應份的事.你若不停,你就犯法了.並非人行善,罪便可以赦免.
\newline
\newline
今晚所要講的罪,不是單單你所犯的罪；或者你說,沒有殺人,搶劫,姦淫,偷盜.這些不過是罪的行為.最要者則罪裡面的性.人都有犯罪的可能性和傾向.
\newline
\newline
耶穌到世上來,祂是神的兒子,祂救我們脫離罪惡,故此,保羅說:「我在你們中間已經定了主意,我不知道別的；只知道耶穌基督,並祂釘十字架.」如果耶穌沒有釘十字架,耶穌不過只是宗教的教主,因祂釘十字架,成了人類的救主.如果十字架沒有耶穌,十字架只不過是羅馬最殘酷的刑具；有十字架並有耶穌,這十字架才是我們所敬仰的.因十字架曾經釘過耶穌,所以說,十字架永遠是我的榮耀.
\newline
\newline
人在世上無法可以除罪,惟仰望十字架上的耶穌,因祂受的鞭傷我們得醫治,因祂受的刑罰我們得平安.我們都如羊走迷,各人偏行己路,耶和華將眾人的罪都歸在祂身上,當我們接受耶穌的時候,我們的罪得著赦免,心得著平安,內心有喜樂.
\newline
\newline


\newpage

\section{第46屆~港九培靈會~周主培~第三講}
\label{sec:1212}
\textbf{祂釘十架}
\newline
\newline
連結: \href{https://www.hkbibleconference.org/session-message/view/1212}{\texttt{https://www.hkbibleconference.org/session-message/view/1212}}
\newline
\newline
\hyperref[sec:1210]{< < < PREV SERMON < < <}
~
\hyperref[sec:toc]{[返目錄]}
~
\hyperref[sec:1213]{> > > NEXT SERMON > > >}
\newline
\newline
澳門有一座古老的聖保羅禮拜堂,建築已逾百年,曾因海盜侵襲被焚,後座全部倒塌,僅留前座石牆,牆頂上十字架仍然屹立不動.一八五三年,英國的包靈爵士到香港任總督,曾在澳門海灣遇颱風吹襲,許多屋宇被毀,什物在海中飄流,惟聖保羅堂的十字架仍屹立,有感而作此詩:「年代久矣百物壞兮,十字寶架獨留存,先聖後聖讚其超奇,我也以此為榮尊.」這是首百廿多年前教會的詩歌,基督徒都能熟唱此詩,時至今日,十字架仍然奇妙莫測.
\newline
\newline
基督耶穌的道,有五件極其重要的事,就是:耶穌基督的降生,受死,復活,升天,再來.耶穌基督的降生是件極奇妙的事.耶穌基督的受死是件極痛苦的事.耶穌基督從死裡復活,是歷史上所未曾有的事.耶穌基督的升天,是件必然的事.耶穌基督的再來,是我們人類極大的盼望.常有青年弟兄和我談話,他們感到耶穌基督死而復活,是件不容易令人相信的事；他們認為人死了三天復活,是件太奇妙的事.現在我要問:「耶穌死而復活,是否一件奇妙的事?」(這問題三年前我曾在此問過)有人回答說:「奇妙!」但我說一點不奇妙,因祂是神,死而不復活才算奇妙.耶穌的降生,道成了肉身,是件奇妙的事,耶穌死了,是件最奇妙的事,「怎能如此主我的神,你為我死」你是神,是神的兒子,怎麼能死呢?如果有一天,你聽說周主培死了,一點都不奇妙,因他是人,遲早總要死的.我心裡始終覺得耶穌復活奇妙的事,祂是應該復活的,祂是永遠活著的.
\newline
\newline
當我想到耶穌的死,我心裡有說不出的感謝；祂是神竟然為我們死.所以保羅說:「我曾定了主意,在你們中間不知道別的；只知道耶穌基督,並祂釘十字架.」保羅有口才,能說許多其他的哲學,學問,教育,不過保羅說,他只知道一件最重要的事,就是耶穌基督,並祂釘十字架.
\newline
\newline
昨晚我們已經思想過耶穌基督,我不想用各的哲學,各種主義,來擾亂我們的思想,所以昨晚僅講耶穌基督.今晚我們要思想耶穌基督,並祂釘十字架.
\newline
\newline
十字架如果沒有耶穌,只不過是一種最殘酷的刑具.如果耶穌基督沒有釘在十字架上,祂的救法就沒有完成.一般人以為基督教若是單講罪,寶血,十字架,這是已經落伍了的思想；有的人認為這樣講,信心只不過是退了顏色的道理.當時保羅曾告訴哥林多教會的人,十字架的道理,在一般人的思想中,是愚拙道理.人以為神如此愚昧,要用十字架的道理,來告訴祂的救法.其實,人早就有覺罪的觀念.昔日猶太人知道內心有罪,以牛羊,鴿子的血獻祭贖罪.中國人也有這樣的思想,孔子說,獲罪於天,無所禱也.中國古代也有皇帝祭天的事.許多思想告訴我們,人的內心有「罪」的思想,有罪就不能到神面前.人要知道罪怎樣才能得赦免.以賽亞53章告訴我們:耶穌是神的羔羊,祂如羊牽到宰殺之地,在那裡代替了我們的罪.
\newline
\newline
十字架至少表明四件事
\newline
\newline
(一)表明人的罪,釘十字架是一件嚴重的事情,罪也是一件嚴重的事情；為了要解決罪的問題,神付了極大的代價.(來九22):「按著律法,凡物差不多都是用血潔淨的,若不流血,罪就不得赦免.」(林後五21):「神使那無罪的,替我們成為罪,」好叫我們從罪中出來.所以主耶穌說:「神啊!我來了,你已經為我預備了身體.」(路十九10):「人子來,為要尋找拯救失喪的人.」(太廿28):「人子來,不是要受人的服事,乃是要服事人；並且要捨命,作多人的贖價.」耶穌宣告說「好牧人為羊捨命.」為了人的罪,神差祂兒子到世上來.假如香港發生了一件事,非英國首相親臨不可,那必是件了不起的事.神可以差遣千萬天使,去作祂的工作；但為了人的罪,卻差祂的獨生子到世上來.魔鬼知道這事,所以當耶穌降生時,想藉著希律王的手殺害耶穌；魔鬼也千方百計地,攔阻耶穌上十字架去.耶穌從未犯過罪,祂在世時,未曾說過「請原諒我」「對不起」「我錯了」「我忘記了」；也沒有說過「差不多」「大概」「也許」「或者」諸如此類的話.耶穌說話都是:「我實實在在的告訴你們,」耶穌衷心切實的說「信子的人,就有永生.」耶穌說祂的話,就是靈,就是生命.話是人內心的表達,所謂之成於內,而形於外.我們的主耶穌從未犯過罪,連仇敵都找不出祂的把柄,審判祂的彼拉多,一而再,再而三的說:「我查不出祂的罪來.」哦!十字架給我們看見罪是件嚴重的事,不是因為耶穌有罪,乃是為著我們的罪.
\newline
\newline
(二)十字架是最痛苦的刑具聖經告訴我們,耶穌從釘十字架至斷氣,僅六個鐘頭.照一般情形,釘十字架至死要較長的時間(有的一天或兩天甚至有的三天)為什麼只經過這麼短的時間呢?耶穌的心實在受了極大的痛苦.祂在客西馬尼園被捉拿時,已是深夜的時候；加以長時間受審；一地又易一地,一官府又換一官府.到處被人戲弄,鞭打,辱罵,最後被定罪；人用荊棘冕戴在祂頭上,用紫袍袍穿在祂身上.人打祂,用布幪祂的眼問:「打你的是誰?」耶穌受盡了各樣的痛苦；最後,被兵丁押出去,背上沉重的十字架.越走越疲倦,力不能支,一次又一次倒下去.無情的羅馬丁兵,用皮鞭打祂,逼祂往前走,走了又倒,倒了挨打,打了勉強再走.到了羅馬兵丁沒法再逼祂走,便捉了一個從非洲來的,黑種人古利奈人西門,替耶穌背了十字架.我相信那時耶穌必定一路扶著西門,直到各各他山上.到了那裡,他們把十字架平放地上,用釘子把耶穌的雙手和雙腳釘在十字架,然後豎起.這樣全身的重量都集於被釘的部位,血一滴滴的流出來直到流盡.耶穌在十字架上說「我渴了!」但這並不是耶穌最覺痛苦的事.希伯來書說「祂為人人嘗了死味,」這死不單是身體的死,聖經告訴我們,這死是與神脫離關係.所以耶穌在十字架上時說:「我的神!我的神!為什麼離棄我?」親愛的弟兄姊妹!耶穌很少稱呼父為神,這裡耶穌說:「我的神,」何故?因當時祂站在罪人的地位,向創造萬物之主呼求.祂背負了世人的罪.叫祂最難過的,不是祂所流的血,也不是頭上荊棘的冠冕,而是祂感到神離棄了祂.罪的結果,就是人離開了神,因此,人覺得虛空,無目標,無安息,孤單,人心覺得沒有前途,沒有安慰,前途茫茫!你所遭遇的,你所感覺的,耶穌在十字架上都已嘗到了；耶穌在十字架嘗了死的味道.聖經告訴我們,人有兩次的死,一是身體的死,一是永遠的死亡,永遠的死和神斷絕了關係.有人間,神既愛我們,為何要造地獄呢?神造地獄並非為人預備,而是為魔鬼預備的.耶穌就是光,有祂就沒有黑暗.祂就是生命,死亡在祂面前變了顏色.祂就是平安的根源,雖然我們在世未免遭遇傷心的事,但我們絕不失望.神造地獄,絕非為了人的緣故；除非你拒絕祂,自甘到沒有神的地方去.哦,耶穌在十字架上,嘗到地獄的味道,一個真心信耶穌的人,不是將來才到天堂去,理應現在即在天堂；因為耶穌同在就是天堂,不是將來你才進天堂,將來更要永永遠遠和祂在一起.
\newline
\newline
耶穌在十字架上擔當了所有的人的罪.(羅五12):「罪是從一人入了世界,死又是從罪來的,於是死就臨到眾人,因為眾人都犯了罪.」有人說,這未免太不公平,亞當一人犯罪,以至眾人都成為罪人,親愛的朋友們!這是法律上的問題,不是道理上的問題.起初我不太懂,後來我在日本工作的期間；有位教授告訴我,當一九四五年,日本天皇投降了,於是所有日本人都成為亡國奴.這是非常公平的事,日本人不能說「我沒有打敗仗,為什麼我都投降呢?」這是法律的問題,不是你同意與否的問題.而亞當一人犯罪,所有的人都成了罪人.當日軍侵香港,每個香港人都要去領良民證,名義上叫良民證,其實就是亡國奴.弟兄姊妹!感謝主耶穌基督!雖然我們每人都有罪,內心都違背了律法；何謂律法?律法就是標準.今天的世界是無標準的世界.有的人說:「善惡沒有標準,高低沒有標準.」這在於尺度有問題,量度就沒有標準.神是我們的標準,神所規定的,如果人不達到標準就是罪人.世人的罪,不是單單從前才有的；有人說,如果耶穌今日到世上來,一定有許多人要跟隨祂,很多人相信祂.但是不一定,耶穌若今日到世上來,不單從前的猶太人會把祂釘十字架.現在世界不信耶穌的,都會把耶穌釘十字架.如果我是個尚未信耶穌的人,可能我也會把耶穌釘十字架；勿以為當時猶太人和羅馬兵丁那麼殘酷,他們把耶穌釘十字架；其實乃是因你我的罪,因人犯罪,以致耶穌在十字架上受痛苦.
\newline
\newline
(三)十字架是神的大能為何我們要講十字架──那古老的十字架,一千多年前的十字架,曾經拯救那悔改的強盜,曾經帶領千千萬萬的罪人進入光明.十字架的大能,若不流血,罪就不得赦免.(來十20):「是藉著祂給我們開了一條又新又活的路......」十字架就是為我們開了一條又新又活的路.有個小孩子進入禮拜堂時,指著掛在那裡的東西問媽媽: 「這是什麼?」他的說: 「你不知道嗎?這是「加」的記號.」很多人不信耶穌,是怕信了耶穌之後,行動不可隨便,有的人說,「信了耶穌,不可抽煙,不可喝酒,不可看電影,不可打牌......一連串的不可；等我戒除這些嗜好,然後來信耶穌.」耶穌基督的道理,並非不可做這,不可做那的道理.換言之,耶穌基督的道理不是「做」(do),也非「不做」(don and \#39;t do),乃是「做成了」(done).耶穌在十字架上說「成了」這才是十字架的道理.人靠自己無法行善,心想做好人卻沒有力量,惟十字架能改變人的心,能救人的靈魂.因為十字架已經擊敗了魔鬼的頭,在十字架前面,撒旦是失敗的仇敵,連死亡也被吞滅了.(來二14):「特要藉著死,敗壞那掌死權的,就是魔鬼；並要釋放那些一生因怕死,而為奴僕的人.」
\newline
\newline
十字架永遠是我的榮耀.十字架叫我們的罪得以洗淨；虧欠得著赦免,叫我們與神和好,叫我們與神和好,叫我們永遠稱為神的兒女.
\newline
\newline
(四)十字架表明神的大愛「神愛世人,甚至將祂的獨生子賜給他們；叫一切信祂的,不至滅亡,反得永生.」「神愛世人」英譯「神何等愛世上的人.」(我很喜歡英文的繙譯).記得有次我離家兩個多月,那時我的孩子還小,當我回來時,有一孩子用英文對我說「爸爸:「I am miss You 」,另一個孩子說「爸爸:I miss You So Much.」哦!「So」,神何等的,十分的愛我們,這愛非我們口舌所能述說.
\newline
\newline
十字架縱橫兩根木頭,將神的愛表明出來.我常常感覺,耶穌儘管可以從十字架下來,可以在未釘十字架以前；差遣千萬天使,救他離開那個世代.甚至可以差遣天使,拔出那些人的刀,把他們除滅!但天使只默默的看這件事情發生.你是神,居然能死!是因為你的愛.
\newline
\newline
親愛的朋友們!我在這裡大聲疾呼的告訴你,巴不得我全身是口,能夠清楚告訴你,神是何等的愛世人；如果世上惟你一人,耶穌基督也願意為你死.神為愛你之故,把祂的獨生子賜給你；祂知道你心裡的光景,知你內心的痛苦,知你被罪所壓傷,知你內心沒有自由.
\newline
\newline
十字架永遠是我們的榮耀,靠著十字架,人與神親近,與神和好了.耶穌基督曾一次受死,叫那義的代不義的；因祂的死,叫萬有與祂自己和好.
\newline
\newline
今晚我不傳別的,只傳耶穌基督並祂釘十字架；雖然這是一件古老的道理,卻是萬古常新的道理.千千萬萬罪人,因著十字架,得稱為神的兒女,從撒旦權下歸向神.跟隨十字架離開黑暗,歸向光明.因著十字架,今天晚上你回去時,將成為一個新造的人.「若有人在基督裡,他就是新造的人；舊事已過,都變成新的了.」因著十字架,你的罪得著赦免,心靈得著平安.
\newline
\newline


\newpage

\section{第46屆~港九培靈會~周主培~第四講}
\label{sec:1213}
\textbf{光中見主}
\newline
\newline
連結: \href{https://www.hkbibleconference.org/session-message/view/1213}{\texttt{https://www.hkbibleconference.org/session-message/view/1213}}
\newline
\newline
\hyperref[sec:1212]{< < < PREV SERMON < < <}
~
\hyperref[sec:toc]{[返目錄]}
~
\hyperref[sec:1214]{> > > NEXT SERMON > > >}
\newline
\newline
以賽亞書 6:1-6:13
\newline
\begin{longtable}{cl}
\hline
\hline
章節 & 經文 (和合本修訂版)\\
\hline
6:1 & \begin{tabularx}{0.7\textwidth}{X} 當烏西雅王崩的那年，我看見主坐在高高的寶座上。他的衣裳下襬遮滿聖殿。 \end{tabularx} \\ \\ \relax
6:2 & \begin{tabularx}{0.7\textwidth}{X} 上有撒拉弗侍立，各有六個翅膀：兩個翅膀遮臉，兩個翅膀遮腳，兩個翅膀飛翔， \end{tabularx} \\ \\ \relax
6:3 & \begin{tabularx}{0.7\textwidth}{X} 彼此呼喊說：「聖哉！聖哉！聖哉！萬軍之耶和華；他的榮光遍滿全地！」 \end{tabularx} \\ \\ \relax
6:4 & \begin{tabularx}{0.7\textwidth}{X} 因呼喊者的聲音，門檻的根基震動，殿裡充滿了煙雲。 \end{tabularx} \\ \\ \relax
6:5 & \begin{tabularx}{0.7\textwidth}{X} 那時我說：「禍哉！我滅亡了！因為我是嘴唇不潔的人，住在嘴唇不潔的民中，又因我親眼看見大君王－萬軍之耶和華。」 \end{tabularx} \\ \\ \relax
6:6 & \begin{tabularx}{0.7\textwidth}{X} 有一撒拉弗向我飛來，手裡拿著燒紅的炭，是用火鉗從壇上取下來的， \end{tabularx} \\ \\ \relax
6:7 & \begin{tabularx}{0.7\textwidth}{X} 用炭沾我的口，說：「看哪，這炭沾了你的嘴唇，你的罪孽便除掉，你的罪惡就赦免了。」 \end{tabularx} \\ \\ \relax
6:8 & \begin{tabularx}{0.7\textwidth}{X} 我聽見主的聲音說：「我可以差遣誰呢？誰肯為我們去呢？」我說：「我在這裡，請差遣我！」 \end{tabularx} \\ \\ \relax
6:9 & \begin{tabularx}{0.7\textwidth}{X} 他說：「你去告訴這百姓說：『你們聽了又聽，卻不明白；看了又看，卻不曉得。』 \end{tabularx} \\ \\ \relax
6:10 & \begin{tabularx}{0.7\textwidth}{X} 要使這百姓心蒙油脂，耳朵發沉，眼睛昏花；恐怕他們眼睛看見，耳朵聽見，心裡明白，回轉過來，就得醫治。」 \end{tabularx} \\ \\ \relax
6:11 & \begin{tabularx}{0.7\textwidth}{X} 我就說：「主啊，這到幾時為止呢？」他說：「直到城鎮荒涼，無人居住，房屋空無一人，土地極其荒蕪； \end{tabularx} \\ \\ \relax
6:12 & \begin{tabularx}{0.7\textwidth}{X} 耶和華將人遷到遠方，國內被撇棄的土地很多。 \end{tabularx} \\ \\ \relax
6:13 & \begin{tabularx}{0.7\textwidth}{X} 國內剩下的人若還有十分之一，也必被吞滅。然而如同大樹與橡樹，雖被砍伐，殘幹卻仍存留，聖潔的苗裔是它的殘幹。」 \end{tabularx} \\ \\
[1ex]
\hline
\hline
\end{longtable}
今天晚上,我們要思想,作為一個神的僕人,他的光景應當如何?
\newline
\newline
從以賽亞書第六章,或許你只看見一個神的僕人；其實這裡有兩個人的名字；一個是作王的烏西雅,一是作先知的以賽亞.
\newline
\newline
烏西雅似乎是個名不經傳的王,如果題到大衛和所羅門,我們都知道；但題到烏西雅,未免覺得生疏.
\newline
\newline
大衛作王四十年,這是一段相當長的時間.所羅門亦作王四十年.烏西雅也作王52年,(代下廿六3)這是很長的時間,可能是父,子,孫三代,都在烏西雅執政的王朝.一個作了52年之王的烏西雅,他的名字在以色列人中,當然家喻戶曉.歷代志下廿六章全章聖經,告訴我們關於烏西雅王的事.中國有詩云:「出師未折身先死,長壽英雄淚滿襟.」烏西雅為王之年,作了很多事情.
\newline
\newline
烏西雅是個尋求神的人,四節說,烏西雅行耶和華眼中為正事.他十六歲登基,就開始熱心追求討神的喜悅；他年青時就有愛主的心,有尋求神的心.
\newline
\newline
烏西雅是個軍事家(6-8)他和鄰國打仗,出去攻擊非利士人；他得勝且佔領許多領土,名聲直傳到埃及.
\newline
\newline
烏西雅是建築家(9節)他建造許多城樓.
\newline
\newline
烏西雅是水利家,(10節)他挖了許多井；又說他喜悅農事.(11-13節)說他對於戰略非常注意.(14-15)我們看見他對於軍事方面,工業方面,兵工方面非常注重,他無論對於內政,外交,政治,宗教,軍事,農業,水利,工業各方面,的確是一位傑出的人才,他在位五十二年,功績輝煌.可惜後來,16節告訴我們:「他既強盛,就心高氣傲,以致行事邪僻,干犯耶和華他的神,進耶和華的殿,要在香壇上燒香.」
\newline
\newline
雅各臨終時,扶著杖頭敬拜神.記得曾經有位神的僕人講過,(我在此也恐懼戰兢地講)一個人開始好壞,不能決定他的結局,必須蓋棺定論.求神憐憫!今天神保守我們,使用我們,我們毫無誇耀之處.神恩待我們,完全出於神的恩典,恩典絕不會叫人驕傲如果知道一切都是從神領受的,就無可值得驕傲.有首詩歌:「我一切所有,無非是接受,都是主賜下在我心中.故此我不驕傲,也不自誇,我是個罪人蒙恩典.」弟兄姊妹!我們若驕傲,將神的恩賜當作自己的才幹,這是最危險的.驕傲的人,無非是要叫人尊敬他,結果,反而叫人輕看他.記得成歸寄牧師曾說:「傳道人,驕傲不可有,傲骨不可無.」很多時候,我們失去神僕的骨格.我小的時候,家父常囑我要牢記:「虛心竹有低頭葉,傲骨梅無仰面花.」可惜這位神的僕人烏西雅,到了名成利就的時候,竟「心高氣傲,」結果一敗塗地.聖經給我們看見,天使長成為撒旦,就因為驕傲.神恩待我們,把我們置於自由的環境中,給我們豐衣足食；並非我們比別人好,乃是神的恩典.求神憐憫我們,教我學習雅各臨終之時,扶著杖頭敬拜神.
\newline
\newline
以賽亞一章一節告訴我們,以賽亞曾作烏西雅,約坦,亞哈斯,希西家等四代王朝的先知.當烏西雅王崩的那年,他的年紀尚輕；我相信那時他必因王崩寶座空,而覺得前途茫茫.烏西雅在位時,東征西討,國家強盛；一旦王崩,附庸之國,雖免紛組皈叛,國家的形勢,岌岌可危；百姓之心,也驚慌失措.以賽亞這時,經驗了他自己所說的話:「你們休要倚靠世人,他鼻孔裡不過有氣息；他在一切事上可算什麼呢?」很多時候,我們所講的,不過是頭腦的知識,並非生命的經歷.有的傳道人,平時講的頭頭是道；到了環境改變,卻把曾經講過的都忘記了.有個不會游泳的人,他就買了一本學習游泳的書,按照書本所示,天天在床上練習；到了他認為學習頗有心得,就正式去游泳.誰知一下水中,所學過的泳技全部忘記；手不會動,腳也不會動,只會在水中連呼救命.我們有了頭腦的知識,同時也需要實際的經驗；客觀的事實和主觀的經驗,應兼而有之.親愛的弟兄姊妹!許多基督徒只有客觀的事實,卻沒有生命的經驗；風平浪靜時,高聲歌頌讚美主.黑雲滿佈時,怨天尤人,臨風洒淚,心裡充滿痛苦.
\newline
\newline
以賽亞在烏雲滿佈,環境黑暗時,他抬頭仰望主.很多時候,基督徒從未有經驗,仰望為我們信心創始成終的耶穌.你所信的耶穌,可能你是父母的耶穌,牧師的耶穌；不是你自己所信的耶穌.保羅對提摩太說:「我知道我所信的是誰.」有很多基督徒,可能知道他所信的是什麼；但卻不知道他所信的是誰.
\newline
\newline
聖經說:「當洪水氾濫的時候,耶和華坐著為王,耶和華坐著為王,直到永永遠遠.」(詩廿九10)摩西的禱告:「主啊!你世世代代作我們的居所,諸山未曾生出,地與世界你未曾造成,從亙古到永遠你是神.」(詩九十1)這裡給我們看見,一個人對神有切實的信仰,所以他說:「我要向山舉目.」(一二一1)青年弟兄姊妹們!我們的思想並非要向左或向右,我們的思想應當向上.今日世界的風浪,有時衝我們向左,時而衝我們向右；你要留心聽,前面有聲音說:「這是正路,要行在其中.」(賽卅21)哦,以賽亞向上觀看!當你舉頭仰望主的時候,你就在光中看見了神,也看見了自己.當以賽亞看見自己,他說:「禍哉!我滅亡了,.......」(5)以賽亞第五章8,11,18,20,21,22,這六節聖經的開頭,以賽亞都是說「禍哉!」他責備人說:「你們這班人有禍了.」到了第六章,以賽亞說「禍哉!我滅亡了.」當人還沒有看見神之前,專指摘別人的不對.今日的世界,人互相責備,彼此指責.很多時候,我們用手指,指別人的不對,(特別是傳道人)請各位注意!當我們用手指指著別人的時候,有三隻指頭是指著自己的.
\newline
\newline
現在我要問:撒該要看耶穌,為什麼看不見?是因為矮嗎?路加第十九章第三節告訴我們,他所以看不見耶穌,是因人多的緣故.很多人不願意相信耶穌,(特別是青年人)雖然在教會的家庭,學校長大,卻不願意進入教會；因為他們看見很多的人,以致不得看見了耶穌；很多人跌倒,冷淡,就是因為看人,沒有看見自己.
\newline
\newline
以賽亞沒有看見神之前,他只看見神的時候,他說自己有了禍.
\newline
\newline
親愛的弟兄姊妹!不但尚未信主的人要悔改,基督徒要悔改,我恭敬的說, 我們傳道人要悔改.在這末世的時候,多少傳道人羞辱了主的名；打官司,上法庭,爭權奪利,勾心鬥角；吞滅人家的財產,欺侮孤兒寡婦.他們所作的,惟有神知道.以賽亞是傳道人,他開始傳道,他需要悔改.
\newline
\newline
人得救只有一次,但要常常在神面前悔改,說:「神啊!求你鑒察我,知道我的心思......看在我裡面,有什麼惡行沒有,引導我走永生的道路.」(詩一三九23-24)
\newline
\newline
我們盼望香港的教會復興,復興由個人開始.自己先在神面前悔改,承認自己的罪；求神光照,求人寬恕.神今日恩待香港的教會,非因我們的緣故,乃因祂自己的榮耀.神的確恩待我們,我們若不悔改,到了有一天,神的忿怒會臨到我們.
\newline
\newline
(賽六5):以賽亞說:「禍哉,我滅亡了!因為我是嘴唇不潔的人,又住在嘴唇不潔的民中；又因我眼見大君王萬軍之耶和華.」我們當在神面前悔改認罪,求神的聖靈光照我們.罪能把神的恩典攔阻,耶和華的膀臂並非縮短,祂的耳朵並非發沉,乃是我們的罪攔阻了祂.我們說要傳福音,要差宣教士出去；可是我們自己的行為,自己的光景,沒有在神面前悔改,以至神不施恩給我們.
\newline
\newline
惟願我們低頭在神面前,隨著聖靈的感動,求神光照我們,求神饒恕我們!
\newline
\newline


\newpage

\section{第46屆~港九培靈會~周主培~第五講}
\label{sec:1214}
\textbf{黎明日出}
\newline
\newline
連結: \href{https://www.hkbibleconference.org/session-message/view/1214}{\texttt{https://www.hkbibleconference.org/session-message/view/1214}}
\newline
\newline
\hyperref[sec:1213]{< < < PREV SERMON < < <}
~
\hyperref[sec:toc]{[返目錄]}
~
\hyperref[sec:1215]{> > > NEXT SERMON > > >}
\newline
\newline
昨天晚上,當我禱告時,我心裡有說不出來的感謝.我想到曾有人說:「將來我們到了天堂,會發現很多奇妙的事.」當然最奇妙的,就是看見我們奇妙的救主；還有一件奇妙的事,就是我們平時以為將再會在天上的人,而卻看不見了；以為不會在天上的人,反而看見了.像我這樣的人,居然也在神那裡,這是件更奇的事.我一面禱告,一面感謝,這確實是神的恩典!正如保羅說:「我今天成了何等樣的人,完全是神的恩典所造成.」聖經說:不是我們尋找神,乃是神尋找我們；不是我們先愛神,乃是神先愛我們.
\newline
\newline
當我們想到神的忍耐,寬容,特別在雅各的身上.有句話說,神的磨磨得很慢,但磨得很細.神一步步的在人身上工作,神等待.至於人,常常不能等待.如摩西,一出來就把個埃及人打死.神等待摩西,從一歲到四十歲；在皇宮裡,從四十歲到八十歲；在曠野中,神等待,忍耐,寬容.神所以如此,無非要我們悔改.親愛的弟兄姊妹!我們切勿忘記神的恩典,勿在神的恩典中墮落,勿因神的寬容忍耐,而任性去走自己的道路.
\newline
\newline
雅各是個很有才幹,有思想,聰明而且刻苦耐勞的人,他的一生,滿有神的恩典在他身上.雅各離開父母,家庭,到他舅父那裡.不辭勞苦,辛勤牧羊,白晝日晒,夜間露滴.經過了廿年,他思家心切,很想回去,惟恐他的哥哥,因為廿多年前種下的禍根,如今要自食其果.唉!人所種的是什麼,收的也是什麼.聖經告訴我們:他先差人去見以掃,那人回報說,他的哥哥帶了四百人迎面而來,使他心裡充滿恐懼,不知所措；他打算,安排,他跪下禱告,然後他進行一切.
\newline
\newline
查創世記卅二章,我們看見三件事,就是此次主題:「黑暗,黎明,日出.」13節:「當夜雅各在那裡住宿.......」21節:「......那夜雅各在隊中住宿.」22節:「他夜間起來.......」那是雅各渡著漫漫的長夜,雅各的生命中,一直都在黑暗裡,當他離開父家,疲於奔命,日已西沉,枕石而寐.有個傳道人形容雅各當時的情景說:「雅各夢見慈母流淚滴濕他的臉；醒來時,發覺原來是露水披面.」那時雅各實在太疲乏,他睡覺,醒而又睡；神向他顯現,對他說話,雅各夢寐初醒,說:「想不到神居然在這裡.」這是今日基督徒的光景.很多時候,我們離開了禮拜堂,離開了信徒,就離開了神.我常對青年人說,多時我們不是怕犯罪；而是犯了罪怕給人知道,所以暗暗的犯罪.雅各以為他離開了家,神仍留在家裡.我覺得基督徒禱告時,毋須說:「神啊!求你與我同在.」因為神從未離開我們,問題在於我們離開了神.正如光明白晝有太陽,黑雲滿佈時太陽依然在那裡.最要者,應常常知道神與我們同在,神永遠不離開我們.有時我個人去旅行,人家問「你是否一個人來呢?」我說,「不是的,神與我同在.」基督徒可能說:「我將耶和華擺在我面前,因此我的心很歡喜,我靈快樂,我的肉身也要安然居住.」如果我們知道神永遠不離開我們,那麼,我們的生活,就有很大的改變.
\newline
\newline
雅各一直在黑夜裡,他向神許願說:「神啊!你若賜我有吃的食物,有穿的衣服,帶領我平平安安的回來,我就以你作我的神.」這是何等幼稚的禱告!很多基督徒向神禱告說:「神啊!求你憐憫我,恩待我,保佑我,看顧我,和我的兒女,財產,生意.」日夜禱告的都是為已.請原諒我說:很多基督徒不過將耶穌當作菩薩來拜,從前拜財神,拜菩薩,現在拜耶穌；所求的不過是吃的,穿的,和保佑平安!難怪基督徒靈命不長進,僅求神的恩典,只要神的福分,從未想到神自己.難怪當試練臨到時,我們不知怎樣依靠主；當危險來到時,我們不知怎能得著平安.我們所注重的,不過是出入平安,有衣有食,過著極其很幼稚的生活,正在黑夜的光景.多時我們跪下禱告,清清楚楚的,起來作卻糊裡糊塗,禱告和行為,完全不一致,俗語所謂「小和尚念經有口無心.」(賽廿九13):「主說,因為這百姓親近我,用嘴唇尊敬我,心卻遠離我......」(結卅三31):「他們來到你這裡如同民來聚會,坐在你面前彷彿是我的民；他們聽你的話卻不去行,因為他們的口多顯愛情,心卻追隨財利.」這就是說,他們來到港九培靈會,如同民來聚會,他們坐在神面前,彷彿神的百姓；他們聽了神的話不去行,因為他們的口多顯愛情,心卻思想股票的起落.口多顯愛情,心卻追隨財利,這是我們今天的光景.我們的心常像小河,來潮就充滿了水；退潮時,心潮也低落了.當股票漲價,心高氣傲；股票跌價,心也掉落下來,臉也拉長了.人的心終日在打算盤,三一三十一,逢三進一；真是如以西結所說:「口多顯愛情,心卻追財利.」也正如以賽亞所說:「這百姓親近主,用嘴唇尊敬主,心卻遠離主.」
\newline
\newline
有個故事說,天上的天使養了一隻老虎,老虎漸長,從未吃過人的肉.有一天,老虎要求天使讓牠到地上去,因牠很想找人肉吃；經牠再三要求,天使才准許牠去,吩咐牠不可吃人,但老虎實在很想試人肉的味道.後來天使只好答應牠,不過有條件,就是不可吃信耶穌的人.老虎到了地上,在森林裡等了很久,才聽到有腳步聲,牠很高興,以為機會來了.那知走過來的人邊走邊唱讚美詩,牠大失所望,因來者是個基督徒.不料那人忽看見老虎,嚇得昏倒.老虎想,既然昏倒,不妨走近聞下人肉香味,牠從頭聞到腳,實在忍不住,於是一口又一口的吃了,僅留下兩片嘴唇,因為那人剛才在唱讚美詩.我不知道在座中有沒有人用嘴唇尊敬主,心卻遠離主.口多顯愛情,心卻追隨財利呢?
\newline
\newline
雅各跪下禱告時,頭頭是道,有條有理,但禱告完了,就照著自己的方法去做.基督徒有的依靠禱告,有的靠信心；我並非說禱告不好,也非說靠信心不對；但信心不是迷信,迷信是無對象無根據的.我們不是靠禱告,乃是當靠聽禱告的主.有的人禱告,就像留聲機的唱片,他的禱告,沒有真正的看見神,他所依靠的不是神自己,而單依靠自己的禱告.我盼望弟兄姊妹每禱告時,真正能夠看見我們的主.把我們每句禱詞,從心的深處發出.有一次,我禱告時,聖靈責備我是個無神論的人.無神論者沒有神只有自己.弟兄姊妹!你如何決定你的一生?如何運用你銀行的存款?如何選擇你的職業?如何選擇你的對象?你說:「我已經決定了.」雅各的禱告,正像自己寫好了一張支票,然後等神了簽字.一切都計劃好了,然後對神說:「求你拯救我.」雅各只想到怎樣與他哥哥和好,卻忘記了先與神和好；他一心想見他哥哥的面,卻忘記了；除非先與神和好,否則不能與人和好.如果與神的關係搞清楚,才能搞清楚與人的關係.今天有很多人常因小故而搬家改行業,無論如何,總是離不開人；有人之處,必有問題存在.
\newline
\newline
雅各孤單漫漫的長夜中,可能他在思想,怎樣面對他哥哥和所帶來的四百人大隊,當然寡不敵眾,最好,和哥哥個別摔跤,雅各正思想時,果然有人和他摔跤,黑暗中他不知對方何人,就用盡生平之力對付.弟兄姊妹!未知我們有誰正與神摔跤呢?你明知神的旨意,卻不願遵從,明知聖靈的感動,卻不肯順服；你與神摔跤已經很久了,你用盡方法來安慰了自己,賄賂自己的良心.
\newline
\newline
所有的罪都是得罪神,我們要當心,千萬不要在神面前行詭詐,不要有假冒偽善的心.我覺得一個詭詐的人,連神對他都沒辦法(可能我說得過分,當然神無所不能)一個假冒為善的人,不肯悔改,不肯承認自己的軟弱,一直與神摔跤.有很多人向神求平安；除非你與神和好,你就得不著神的平安.我們務要與神和好.
\newline
\newline
雅各與神摔跤,從黑暗直到黎明之前.照樣,這是你的黑夜,你在黑夜裡已經太久了,神在尋找你,神盼望你別與祂摔跤.今晚你來參加聚會,並非偶然,你自知心裡的光景,數年來你一直有件事不肯放下.有一人,一時候,一物你不肯放下；正像你把房間的門窗都打開,還保留一個小小的保險箱；你對神說,這是我私人的財產.
\newline
\newline
雅各一直與神摔跤,其實神可以輕易的勝過他；不過,很多時候,神給人機會.據游泳教授說,如果人被水溺,必拚命掙扎,要救他的人,必須善用方法,把他輕輕打昏,或暗暗從背後抱住他,否則反而會被他拉著一同沉下去.有首英文詩:「祂摸著我,教我的生命能夠完完全全的改變.」巴不得神今晚挨著我們,如同摸著雅各的大腿窩,這是人最強有力之處.當神摸著我們,我們就從黑暗進入黎明.巴不得神今晚摸著我們的心,轉動我們的心,好叫我們的心完全的歸向祂.
\newline
\newline
雅各被神摸著時,他覺得摸著他的非本常人,他就緊緊抓住,對那人說:「你不給我祝福,我就不容你去.」雅各從前抓住他的哥哥,抓住他的舅舅,從前抓住人,現在抓住神；他知道若沒有神,不但不能見他的哥哥,他自己也不能活下去.那人問他名叫什麼?說「名叫雅各.」(是抓的意思)他抓住哥哥的腿出母胎,他一直在抓,用自己的方法,詭詐拚命的抓.現在他知道名叫雅各,他才開始認識自己；好像那生來瞎眼的,照耶穌的吩咐,到西亞池一洗,他首先看見了自己.得著神的光照.才能夠認識自己.
\newline
\newline
黎明的時候,雅各面對面看見了神,這是件奇妙的事.聖經告訴我們,神對摩西說:「你不能看見我的面,你若看見我的面,你就要死.」當以賽亞看見神,他說「禍哉!我滅亡了.」當亞伯蘭看見了神,亞伯蘭死了,亞伯拉罕開始活過來.當雅各看見神,雅各死了,以色列活起來了.當約伯看見了神,他說:「我從前風聞有你,現在親眼看見你,因此我厭惡自己,在塵土和爐灰中懊悔.」當掃羅看見了從死裡復活的主,掃羅死了,保羅活出來了.見了神,就死了.
\newline
\newline
弟兄姊妹!我現在不問你信主有多少年,也不問你參加培靈會有多少次；我要問你曾否遇見過神?如果你真正的遇見了祂；你就死了,老我死了,自己死了,這是你的黎明.聖經上的條件和算法都很奇妙.聖經說:當你失掉時即你復得時,你死了時即你復活時,你放下時即得著時.當我們在神面前死了時,神的聖靈進入我們的裡面.「若有人在基督裡,他就是新造的人,舊事已過,都變成新的了.」(林後五17)哦!巴不得今晚在你是個轉機；當你遇見神,聖經說「日頭出來了.」雅各經過毘努伊勒,雖然他瘸了腿,走路失去舊日的風姿；但雅各已經成了王的兒子.(以色列的意思)聖經給我們看見一件極奇妙的事；人的盡頭就是神的起頭.雅各帶著瘸腿一步一步的走,看見迎面來了哥哥和四百人的大隊.他就俯伏在地,一連七次,到他哥哥面前.雅各的生命經過彫刻,經過神的工作,成為神的傑作,有神的痕跡在他身上.他們倆兄弟見面,擁抱大哭.照我的想像,以掃帶了四百人,在出發之前,一定對他們宣佈,在廿多年前如何吃他弟弟的虧.「有仇不報非大丈夫也,今日前來旨在報仇,非把弟弟打得稀糜爛碎不可.」如今竟然彼此擁抱親嘴大哭.那些人惟有面面相觀,可能還滴下同情之淚.這就是神的方法.
\newline
\newline
親愛的弟兄姊妹!一粒麥子若不落在地裡死了,仍舊是一粒；若是死了,就結出許多子粒來.神得勝的方法不同於人.雅各一向「抓」；但當他進到法老皇宮,他為法老祝福,出皇宮時又為法老祝福,正像個王太子一樣,彰顯他父神的榮耀.
\newline
\newline
弟兄姊妹!日頭已經出來了,千萬別再與神摔跤.求神憐憫我們!神知我們的光景；已經太久了,我們不肯順服神,不肯俯伏在神面前,難怪我們一直在黑夜中過生活.巴不得黑夜快快渡過,讓神的手摸著我們,轉動我們,好叫我們完完全全的向著祂.
\newline
\newline


\newpage

\section{第46屆~港九培靈會~周主培~第六講}
\label{sec:1215}
\textbf{曠野樹林}
\newline
\newline
連結: \href{https://www.hkbibleconference.org/session-message/view/1215}{\texttt{https://www.hkbibleconference.org/session-message/view/1215}}
\newline
\newline
\hyperref[sec:1214]{< < < PREV SERMON < < <}
~
\hyperref[sec:toc]{[返目錄]}
~
\hyperref[sec:1217]{> > > NEXT SERMON > > >}
\newline
\newline
以賽亞書 32:15-32:15
\newline
\begin{longtable}{cl}
\hline
\hline
章節 & 經文 (和合本修訂版)\\
\hline
32:15 & \begin{tabularx}{0.7\textwidth}{X} 等到聖靈從高處澆灌我們，曠野將變為田園，田園看似森林。 \end{tabularx} \\ \\
[1ex]
\hline
\hline
\end{longtable}


\newpage

\section{第46屆~港九培靈會~周主培~第七講}
\label{sec:1217}
\textbf{當前急務}
\newline
\newline
連結: \href{https://www.hkbibleconference.org/session-message/view/1217}{\texttt{https://www.hkbibleconference.org/session-message/view/1217}}
\newline
\newline
\hyperref[sec:1215]{< < < PREV SERMON < < <}
~
\hyperref[sec:toc]{[返目錄]}
~
\hyperref[sec:1218]{> > > NEXT SERMON > > >}
\newline
\newline
約翰福音 9:4-9:4
\newline
\begin{longtable}{cl}
\hline
\hline
章節 & 經文 (和合本修訂版)\\
\hline
9:4 & \begin{tabularx}{0.7\textwidth}{X} 趁著白日，我們必須做差我來的那位的工；黑夜來到，就沒有人能做工了。 \end{tabularx} \\ \\
[1ex]
\hline
\hline
\end{longtable}
約翰福音全卷,共記載了耶穌的七神蹟.耶穌在世上所行的神蹟,和所作的奇事極多；聖靈感動約翰記載了七個神蹟；每個神蹟,均有主耶穌基督最高之目的.耶穌使水變酒.使大臣的兒子復活.卅八年的病人得醫治.五餅二魚給五千人吃飽.耶穌在水面上走.使拉撒路復活.使瞎眼的得看見.每次耶穌行完神蹟後,繼之講道；祂有一個最高目的,祂不是單為神蹟而行神蹟.
\newline
\newline
(約九4):「趁著白日,我們必須作那差我來者的工；黑夜將到,就沒有人能作工了.」這段聖經可分成五段.(請原諒我用幾個英文字來講解更為清楚,讓聽者有較深刻的印象,特別是青年弟兄姊妹們).
\newline
\newline
「趁著白日」──機會(Opportunity)「我們」──合一(Unity)
\newline
\newline
「必須」──需要(Necessity)「作那差我來者的工」──優先(Priority)「黑夜將到,就沒有人能作工了」──迫切(Urgency)從這節聖經,神告訴我們要抓住機會,要同心合意,要認定何為必須作的事,何為最要緊的事,最迫切的關頭在那裡?
\newline
\newline
在工業社會裡,我們常聽到「時間比金錢更寶貴」的話,有人說,時間就是金錢.中國人的看法,時間較金錢更重要；所謂一寸光陰一寸金,寸金難買寸光陰；失去寸金有處尋,失去光陰無處覓.金錢可以存寄銀行,需用時支取.時間只能及時運用,卻無法存留.金錢在人的手中,多少各不同；可是神給於人的時間,一律晝夜廿四小時,無人可以多或少.如果存款在銀行,每月底有存款數目報告單；但是我們的時間,從來沒有一個報告說,還可活下去多久；無人可知他所餘下的光陰多少.事實來說,浪費時間,則浪費自己的生命.人許多人正在浪費自己生命的光景中,到了有一天,發現自己已經到了生命的盡頭.
\newline
\newline
弟兄姊妹!上一代的事情,我們無法負責任,無法去改變；下一代的時間,我們也無法去管理.我們能夠負責任的,惟現在的時候.我個人感覺,生活在這時代,青年人的時代,是現在的世代,(Now Generation)我們一切事情卻說 Now 現在.其實,「現在」根本沒有,過去的事,我們無法挽回；現在即迅速過去,一句話說了立即過去,一件事做了立即過去.我們不講過去,甚至也沒有現在,我們僅有的是不斷的將來,我們面對著無限的將來.有的人很可悲,常常談論過去,開口閉口都講過去；這是過去的人,過去的事,不能挽回,惟有留在我們的記憶中.趁著白日則趁著現在──不斷的將來.光陰如流水,既往下流,永遠不倒流回來.水雖不斷地流,但正在急流中的水,已非昨日的水了.聖經說,趁著年幼,應當記念造我們的主.
\newline
\newline
有個故事說,有個人看見個很古怪的人；光禿的頭,前面長出一辮子,他一把抓住那人的辮子.問「叫何名?」答「我名機會,入從窗,出從門,轉身一幌,你抓不到我.」
\newline
\newline
主說:「我父作事直到如今,我也作事.......」我們的主從不打盹；不過,請各位注意!除了神從不打盹,還有不打盹者；由培靈會開始,牠早到遲退,一直把守,撒旦繼續不斷地工作.
\newline
\newline
「趁著白日,我們必須作那差我來者的工.」這節聖經英文欽定本:「我必須作那差我來者的工」,不及中譯好；中譯「我們必須......」,太寶貝啦!可見工作時,不是我自己在工作；正如(林前三9):「我們是與神同工的.......」意思就是我們與神一同工作.真奇妙!沒有神我們什麼都不能做；神沒有我們,祂就不願意工作,祂願意等待,祂如今仍然在天上呼喊說:「我們能差遣誰呢?誰願為我們去呢?」神願意等待.神等待摩西,從他一歲至四十歲；從四十歲至八十歲.神一直等待,摩西不能等待,不能忍耐,神還是忍耐.
\newline
\newline
我們是與神同工的.有一大遊船載了許多遊客,大家在船上無所事事；有的看書,有的談話,有的休息.船上客廳有架鋼琴,一個很小的孩子上去,叮叮噹當地彈；大家覺得小孩子很有趣,都圍著鋼琴觀看,每彈一陣,大家都鼓掌；因此,小孩越彈越起勁.後來觀看的人漸少,掌聲也漸稀微,人群逐漸散開了.最後只剩一人,他覺得小孩很可愛,就抱著坐在他的膝上；彈奏美妙動聽的音樂,吸引了船上眾遊客欣賞.彈奏完畢,大家熱烈鼓掌,小孩得意洋洋,笑臉對彈琴的人說:「我們彈得很不錯.」很多時候,我們如同這小孩；神所作的工作,我們竟然以為是自己的本領.哦!如果沒有神,我們什麼都不能作.主說:「我們......,」祂將自己和我們置於同等地位.「我們」不但與神同工,「我們」也是人與人之間的同工；不是一個人的工作,乃是同心合意興旺福音.神在祂兒女身上的心願,祂願意有身體,有家,有城,有國度；教會是神的身體,我們互相作為肢體,在身體中千萬別有毒瘤.(毒瘤把所有的營養都集於自己)教會是神的家,非個人的地盤,千萬別有地盤主義.啟示錄廿一章說:「這是神的城.」我們萬勿為自己造巴別塔,要叫人知道我們的名聲.我們是為著神的國.非建立自己的王朝.親愛的弟兄姊妹!我們是與神同工的,千萬別唱獨腳戲,神所要的是個大交響樂團；有的吹小笛,有的吹大喇叭,各種樂器的聲音合奏起來.我們要聽見眾天軍天使,在天上歡呼.我們常有錯誤的觀點,以為牧師和傳道人,才是為神工作的,一切應由他們去做.有的人把人當作神,以為自己得救是因他的緣故；因他講道,我才得救,因他傳福音,我才悔改.好像有個人到茶樓飲茶,吃了許多東西；最後吃了一碟鳳爪,即結數,一看帳單,說,為何這麼多錢呢?我是吃了鳳爪才飽的,其他東西怎能算數呢?人得救,非單靠一人的工作,很多人為你禱告,很多人向你傳福音；正好有個神的僕人領奮興會,你得救了,你就歸功於他.常常教會信徒在奮興會之後,對本教會的牧師漸漸冷淡,以為自己的牧師不如那奮興家.我們是與神同工的,我們接力賽跑,每個人都有連帶關係,非單一人快跑就夠了.各人在自己跑的那段路上,都應盡自己的力量才對.當耶穌看見許多人沒東西吃,後來看見小孩手裡有五餅二魚,耶穌拿過來,禱告,擘開,遞給門徒,門徒再遞給眾人,眾人同心合意的工作；這是一幅多美的圖畫.那小孩子把所有的奉獻給主,主藉著他和門徒及眾人一齊工作,結果五千多人吃飽.
\newline
\newline
弟兄姊妹!今天在香港,有多少靈魂需要我們去傳福音,難道我們有空閒來你爭我奪,彼此批評,告狀,在報上登啟事,發新聞嗎?時候不多了,我們需要同心合意興旺福音.如果香港的教會,要得著神的恩典,我們每個工作者,必須悔改認罪在神面前,我們太虧欠神的名,太羞辱神的名,我們沒有在神面前同心合意.
\newline
\newline
「必須」我們一定要作,無其他辦法,神惟有用我們來傳福音.據傳說,當主耶穌升天時,千萬天使歡迎主,敬拜主.天使長問主說:「你在世上只有卅三年半,你到過的地方很少,接觸的人不多；怎能把救恩傳給全世界的人呢?」耶穌說:「當我離世時,我將使命告訴所有的門徒:『你們要往普天下去,傳福音給萬民聽.』我已把責任交給他們了.」天使長又問:「如果他們不作,你怎麼辦?」耶穌說:「如果他們不作,我也沒有別的辦法.」神可以差遣千萬天使到各處傳福音；但是神要我們每個人去傳福音.最近香港的教會,可說蒙神賜福；我看見很多下一代的青年都起來,在文字工作上,在傳福音的工作上,在許多事上為主工作.我們也看見有幾個教會,有的差派工人到各國各地傳福音.這種現象很好,但這並非屬於少數人的工作,而是人人都有責任.勿等到明年的培靈會或下次的佈道會我們再去為主作工；我們必須作工!保羅說:「責任已經托付我了.」中國所謂:受人之托.有人一日三省為人謀而不忠乎.我們有沒有忠心在神面前,神給我們的托付有沒有完成?保羅說:「我不以性命為念,也不看為寶貴,只要完成神所指望我走完的路程,完成神在耶穌基督裡所托付我的責任.弟兄姊妹!神不是單把福音的事交給幾位傳道人；每個基督徒,都有責任去傳福音領人歸主,這是一件必須的事.當主耶穌釘十架之前,特別在路加福音五,十三,十七至十九章,每節聖經都給我們看見,耶穌定意往耶路撒冷去.我們的主,祂知道前面有困難；有十字架,但他對父神說:「我來了,為要遵你的旨意行.」世上許多的激烈份子,革命份子,他們的頭顱可以拋,熱血可以流,但他們的使命一定要完成.今天的教會,有沒有人為主的緣故,離棄自己的家鄉,親人,到各處去傳福音?我曾經乘過一種玻璃底的船遊海,坐在船上,透過船底的玻璃,可以看見海裡大大小小各式樣的魚,於是我有所感想,很多人對傳福音漠不關心,對人的靈魂,正像坐在玻璃底的船上看海裡的魚.我們所需要的,乃是有人願意跳到海裡去捉魚,袖手旁觀是不夠的.英文有兩字:interest──興趣,involve──參與,很多人「興趣」.很少人「參與」,我有個朋友,他的體重二百多磅；他告訴我,他最喜歡運動.我問他喜歡何種運動,他說任何球類他都喜歡.正談時,他的太太聽見了,她對我說:「他整天坐在電視機前,看球賽,自己沒有拿過球.」他就是有興趣,但是沒有參與.救人的工作許多人都有興趣,但他們不肯參與工作.
\newline
\newline
人生到底為何?為吃飯嗎?為成家立業嗎?最後如何呢?死了,這是人生的目標嗎?有一次我和一個嘻皮士談道,我問他作何工作,他說已退休了；我很奇怪,為何才廿五歲已經退休了呢?他告訴我,他的父親和祖父一生的目標就是退休.當我工作之日就查問退休之後有何享受,他們工作的目標則為著退休.工作了廿卅年退休,總算目標達到.「至於我已經達到目標了.」這是美國一班青年的思想.
\newline
\newline
弟兄姊妹,何為你人生最重要的事?我們常常不知為何而活,活是為著吃,吃是為著活.人生就這吃活吃活地打圈子.耶穌說:「我是為著要遵行差我來者的旨意.」這是祂人生的目標.神已經清楚地為我們擺上前面的路,所以(來十二1):我們要奔那擺在我們前頭的路程.到了有一天,保羅說:當跑的路我已經跑盡了.你活著若是沒有目標,你就浪費了一生的時間；讓我再提醒你們,勿浪費時間!浪費時間則浪費自己的生命.
\newline
\newline
在教會中,常常有的傳道人或弟兄姊妹,心裡有很多困擾.願意工作的人,處處需要他.青年弟兄姊妹,願意服事主,詩班,主日學,團契,佈道團都要他；結果應接不暇,非常疲乏.他說:「我應該作什麼呢?」需要是件事,神的託付又是件事.你勿因需要而決定腳步,要因神的托付而規定你的方向,免得越幫越忙,吃力不討好,自找苦吃.我們應作那差我來者的工,在神面前禱告,認識清楚是神的托付,從神接受而後去做,做個專任的基督徒,向著目標前進!
\newline
\newline
「黑夜將到就沒有人能作工了」,很多時,有人以為這是傳道人常喊叫「狼來了」.不久以前,許多地方發生能源縮減問題；據云,照現在用油的情況,存油量不得使用超過四十年.東京有些地方,禁止汽車行駛,因空氣極度污染,行人戴面罩走路,許多大城市,都有空氣污染的現象.日本的魚類,因含水銀毒太重,不能再吃,吃了能形成人體變成矮小.在香港,吸毒,色情,情況嚴重.
\newline
\newline
在美國,前幾年,吸毒,性自由極泛濫.紐約一間著名音樂學校的基督徒團契,有一次請我們向他們同學傳福音.他們預先告訴我,目前美國報攤上有本甚流行的書,是關於人的問題,魔術問題,及神的問題.很多青年,傾向古代中國人迷於求仙求鬼,有一種東方神秘的宗教在他們青年的思想中,他們覺得吸毒不能解決問題,性慾也不過如此；未得到之先,以為得著便可滿足,哪知得到了也不過如此,現在他們相信鬼魔的道理.從前魔鬼戴著面具,現今竟敢公然露面,公開承認是魔鬼.有很多魔鬼禮拜堂,有的人穿了黑的衣服,稱為魔鬼教會的牧師,讓很多人崇拜.
\newline
\newline
弟兄姊妹!魔鬼自知時候不多了,猖獗在我們中間.今日香港雖有復興,青年人有愛主的心,教會對於差派的事極其注重,這固然是好現象.不過,並非僅差派人和奉獻金錢到遠處去傳福音,香港有五百多萬居民,如果你不能做個好的傳道人,到了別處也必定不能做好工作,你不要再欺哄自己吧!
\newline
\newline
在葡萄園裡,有的人早晨進去工作,有的中午,有的較遲；有的到了最後才進去工作,中國基督徒是最後的一批人在為神傳福音的事工上有份的.故此,我們不要再耽廷時日,因為黑夜將到,就沒有人能作工了.
\newline
\newline


\newpage

\section{第46屆~港九培靈會~周主培~第八講}
\label{sec:1218}
\textbf{殘燈不滅}
\newline
\newline
連結: \href{https://www.hkbibleconference.org/session-message/view/1218}{\texttt{https://www.hkbibleconference.org/session-message/view/1218}}
\newline
\newline
\hyperref[sec:1217]{< < < PREV SERMON < < <}
~
\hyperref[sec:toc]{[返目錄]}
~
\hyperref[sec:1219]{> > > NEXT SERMON > > >}
\newline
\newline
馬太福音 12:20-12:20
\newline
\begin{longtable}{cl}
\hline
\hline
章節 & 經文 (和合本修訂版)\\
\hline
12:20 & \begin{tabularx}{0.7\textwidth}{X} 壓傷的蘆葦，他不折斷，將殘的燈火，他不吹滅，直到他使公理得勝。 \end{tabularx} \\ \\
[1ex]
\hline
\hline
\end{longtable}
感謝神的恩典,不但給我有機會在聚會時,和弟兄姊妹一同思想神的話；有時也藉著個人的談話,能和弟兄姊妹在主裡面有團契,有交通.可惜時間匆促,未能與所需要的弟兄姊妹,一一個別談話.近數日來,當我和弟兄姊妹談話之中,我發現有個共同的問題,就是關於心內的軟弱,失敗,跌倒；如何地靠著主的恩典,把下垂的手,發酸的腿,再次的挺了起來,奔跑前面的道路；我們能夠站立得住,完全是神的恩典.恩典從不會叫人驕傲,不叫人自誇.
\newline
\newline
神在每個人身上,都有祂美好的計劃.許多時候,我們成功的經驗,得勝的見證,都能夠幫助人；而失敗的經驗,同樣也能幫助人.可是人常樂意講得勝的經歷,而不願意講失敗的見證.聖經記載很多失敗的人,叫我們看見,不在乎人的力量,乃全部靠那位得勝有餘的主.
\newline
\newline
今晚我們要思想一個出賣耶穌的門徒.當我們提到此事,就立即想到猶大；但今晚我不是要講猶大出賣耶穌,乃是講耶穌的另一個門徒,他用三句話出賣耶穌.何謂出賣?就是否認,在別人面前不承認素常最親愛的耶穌.我想基督徒們都知道那人就是彼得.現在,我們要從彼得失敗的經驗上,看他失敗的原因,及他如何恢復過來.但願這古老的故事,作為我們的借鏡,看見自己裡面的光景.
\newline
\newline
當我們失敗時,常責怪魔鬼,說是因牠的引誘以致犯罪,因祂的試探以致失敗.不錯,魔鬼的確是人失敗的引誘者；牠很聰明,常針對人的弱點加以引誘.牠曾引誘大衛犯姦淫；但大衛本身應負責任.多時我們錯怪了魔鬼,其實是自己替魔鬼留了地步.猶大用卅兩銀子賣耶穌,是猶大自己給魔鬼留了地步,聖經說:「猶大是掌管錢財的,他是個賊,常取其中所存的.」他常給魔鬼預備機會.撒旦知道我們裡頭的光景,神知道,我們自己也知道.
\newline
\newline
我們的主耶穌基督受魔鬼試探,當然是神的旨意,不過,魔鬼不是馬上試探耶穌的.耶穌在曠野四十晝夜,當祂肚子餓了；那試探人的就上前來,對耶穌說:「你若是神的兒子,可以叫石頭變成餅.」魔鬼一起在等待機會.當我們落在試探時,不要只怪魔鬼；要看自己裡面的光景如何,因為魔鬼是無孔不入的.有人以為亞當夏娃被造之初,在伊甸園中,無憂無慮,在那裡,四季如春,既無炎夏亦無寒冬；他們無須工作,不必流汗,無所事事,逍遙渡日；怎能犯罪呢?其實,這樣的觀念是錯誤的.(創二15):「耶和華神將那人安置在伊甸園,使他修理看守.」亞當和夏娃在伊甸園,並非沒有工作；他們的工作是修理,看守.神創造人,不是要讓他們在那裡「退休」「養老」,乃是要他們修理,看守；因有破壞,故要修理,因有仇敵,故要看守.撒旦能進到園裡,因亞當留了破口.人犯罪,常推諉是撒旦的引誘；其實,雅各說:「人被引誘受試探,乃是被私慾所牽引.」
\newline
\newline
為何彼得跌倒?(太廿六33):「彼得說:眾人雖然為你的緣故跌倒,我卻不跌倒.」彼得失敗最大原因,就是「我」在裡面.英文的「罪」字 SIN,「I」「我」在當中.「基督徒」英文字是 Christian,意則「Christ」I am Nothing 基督「我」是無有,惟有基督,我沒有了.(此非英文字典的解釋)基督徒失敗的根源,就是有「我」在裡面.許多基督徒常在不知不覺之中把自我顯露出來.記得我小時,聽過一淺陋的故事:在個池塘邊,有很多飛禽和動物.天氣漸冷,野鴨將南飛,牠們向烏龜辭行,烏龜問說:「你們用何方法到南方去呢?」野鴨說:「我們飛去的,」野鴨向烏龜盛讚南方景色美麗,氣溫溫暖,令烏龜十分嚮往；烏龜忽然想到辦法,叫兩隻野鴨分端抓住一根竹竿,自己咬緊竹竿中間,隨著野鴨南飛.他們先行試飛,非常成功,在地上看見的人,莫不讚賞方法巧妙!兩隻野鴨「哇哇」自鳴得意,烏龜也很高興；奈何未能向觀眾表明,牠是這巧妙方法的設計者.觀眾連聲讚好,說:「這麼好的辦法,不知是誰想的,真是聰明!」烏龜這回得意忘形,大聲應說「我......」遂即譁然墮地,粉身碎骨.
\newline
\newline
基督徒靈命上的危機,由於驕傲造成,聖經上諸如此類的記載甚多.以西結和以賽亞書告訴我們:美麗的天使長因為驕傲,墮落成為撒旦.蓋恩夫人有次寫信說:「親愛的弟兄,風聞你的靈性漸漸長進,但不是向上而是下沉.」請看,天秤重的一邊總是向下的；滿載貨的船沉重在海中；半桶的水,挑起來撲通作響,聲浪很大；滿桶的水,聲音反而很輕.許多基督徒靈性的光景,就像半桶水,聲浪很大,以為別人都跌倒,「我」不跌倒.哦!親愛的弟兄!巴不得我們在神面前,長大成人,滿有基督的身量.
\newline
\newline
耶穌基督,不以自己與神同等為強奪的,反倒虛已成了奴僕的樣子.耶穌生在世上卅三年,實在太奇妙!祂是神,但祂從未使人當祂是神,這實在很不容易.一個有勢力的人,很想叫人知道他的勢力；有學問的人,想叫人知道他是博士；有錢的人,是想叫人知道他是百萬,千萬,億萬,的富翁.耶穌是神,成為人,卅三年半中,在祂最危險的時候；祂可以差遣千萬天使,但祂沒有這樣.最後,人把祂釘十字架,說,你若是神的兒子,可以從十字架下來；但耶穌謙卑,溫柔,忍耐.我們的失敗,乃因驕傲,我常以此提醒自己；真正是謙卑的,自己覺得謙卑的人,就是他開始驕傲的時候.彼得的失敗,就在他驕傲的事上.
\newline
\newline
耶穌禱告,彼得睡覺.許多基督徒不禱告,在睡覺,所以靈性的光景日往下沉.屬靈的道路,不能走捷徑.有的人苯頭笨腦,也發了大財；有的人糊裡糊塗,卻娶了賢妻；有的人隨隨便便,讀書卻很不錯；但在屬靈的事上,絕不能如此.如果糊塗隨便,靈性便往下沉.在屬靈的事上,我們必須付上代價,須禱告,禱告就是爭戰.今天的教會,有計劃,開會,有奔跑,卻缺少禱告；我們用各種方法代替禱告,名為禱告會,人數很少,禱告的時間也不夠.以弗所書講到全副的軍裝,頭戴救恩的頭盔,用公義當作護心鏡遮胸,用真理當作帶子束腰,用平安的福音當作預備走路的鞋穿上,一手拿著信德的籐牌,一手拿聖靈的寶劍；就是神的道,靠著聖靈靈,隨時多方禱告祈求.親愛的弟兄姊妹!當我們與撒旦爭戰時,我們須付上禱告的代價.但以理一日三次跪在神面前禱告.我們的主耶穌,在工作之前,工作之後,通宵的禱告,不住的禱告.我們的主都需要禱告,何況我們呢?今天在我們當中,有多少弟兄姊妹；神知道,撒旦知道,自己也知道很久沒有好好的禱告.基督徒的爭戰,應當是一眼開,一眼閉,何意?乃是開著眼儆醒.
\newline
\newline
彼得遠遠的跟隨耶穌,這是他失敗因素之一.人離開主,非立刻離開,跌倒,也非馬上跌倒；舊約時代的羅得,聖經說他漸漸地挪移帳棚.人離開神,是逐漸的離開.求神憐憫我們!叫我們的心緊緊地跟隨主.
\newline
\newline
彼得跌倒,他如何給撒旦留地步呢?主禱文中有句「不叫我們遇見試探,」在小字是說「不叫我們進入試探,」遇見試探並非罪,有時試探來了,你可以拒絕,可以不接受.我們的主耶穌基督,祂也凡事受過試探,只是祂沒有犯罪.我們遇見試探,受試探,不是一件犯罪的事；但當心受魔鬼欺騙,好像我們的主曾經受過試探,但祂不接受.慕勒說:「你無法攔阻飛鳥經過你的頭上,可是你有權柄不讓鳥在你的頭上作窩.」哦,弟兄姊妹!彼得跌倒了,彼得失敗了,竟然三次否認耶穌.如果彼得承認主,可能他也被捕,受審判,被釘十字架,那時和耶穌同去；我相信耶穌必扶著他的肩膀,一步一步地同走各各他,彼得可幫助耶穌背負十字架.可惜彼得失去了最大的祝福.彼得不承認主,非但叫主的心得不到滿足,且叫主的心憂傷.「彼得發咒起誓的說,我不認得那個人.」哦,弟兄姊妹!多少時候,我們的失敗,叫主的心傷痛.
\newline
\newline
彼得怎樣回復過來呢?
\newline
\newline
(一)因主耶穌基督的代禱.(路廿二32):「耶穌說,我已經為你祈求.」我們那位慈悲忠信的大祭司,祂經常為我們禱告,我們今天能夠站立,乃因主為我們代禱.我相信很多人為這培靈會禱告.在我未離家以前,我有機會和蔡蘇娟小姐一同禱告,(就是那位暗室之後的作者)她說:「摩西,你去罷,我為你禱告,我每天為你禱告,我用禱告作你們的後盾.」我相信除她以外,還有許多的人,在世界各地,在香港,為這聚會禱告.神是聽禱告的神,況且我們的主耶穌基督,祂是慈悲忠信的大祭司,祂為我們禱告.
\newline
\newline
(二)因他想起耶穌的話,神的話是有能力的,當他記起耶穌的話,心裡就有力量.有一位主的僕人曾說:「聖經能使人離開罪惡,罪惡能使人離開聖經.」當人親近神的話時,想起神的話時；神的靈,神的力量在他裡面.
\newline
\newline
(三)因主耶穌基督回頭看彼得,有何力量能叫偶像退掉顏色?有何力量能叫我們脫離地的吸引?不是一切的教條,乃是主耶穌榮耀的顯現.當彼得被主一看,他的心是被主捉去了；耶穌並非以怒眼看彼得；耶穌一看,彼得想起祂的話來了.
\newline
\newline
願主的愛再次激勵我們!今晚各位來聚會,不是偶然的事,主為你禱告,用慈愛的心吸引你,用慈繩愛索牽引你,主願意你回到祂的面前.「雞叫」雞算不得什麼.神用著我們二人在祂面前,好像雞叫,有人聽了覺得討厭,厭煩.很多神的僕人,正像雞叫,叫你起來.神用祂自己的聲音,叫我們回到祂這裡來.
\newline
\newline
(路廿二62):「他就出去痛哭」這數字表示彼得悔改；當他看見主的眼睛,當他想起主的話；他聽見雞叫的時候,毫不遲延,就立刻,不等待的悔改.多時魔鬼說:「你再等一等,慢慢地.」我們的心正如聖經說:「好像被熱鐵烙慣.」開始時受感動,但不悔改.許多人,感而不動；悔而不改；如果是這樣,你每年來參加培靈會,對你無多大用處.但願我們的心,不至被熱鐵烙慣,能夠敏銳知道主的旨意,在我們身上,不要等待.下面接著說「出去」,就是離開罪惡.彼得不停在其中說,我還要向他們作見證.有的人本來吸毒,悔改後要到舊地去作見證,那是好的,不過自己要立得穩,不然還是及時離開,等你站穩後再回去作見證,這是聖經的原則.耶穌說:「凡勞苦擔重擔的人到我這裡來!」來了而後再去.有的人沒出來,還停留原處；好像那心裡有污鬼的人,當污鬼出去,另外七污鬼又進來.彼得本來和差役同坐在院子裡烤火；(太廿六)但他不再留在差役當中,他起來,出去,離開那地,出去痛哭.
\newline
\newline
弟兄姊妹!眼淚好像放大鏡一樣；透過眼淚,我們看見神在面前；眼淚把我們眼中的塵污都洗淨.我們需要流淚,為著我們的罪,為著人的靈魂,為著主的愛,我們都有理由流淚.我相信,彼得流的是悔改的眼淚,也是感恩的眼淚.他看見主的眼睛,看見自己的污穢,看見耶穌的慈愛,他就出去痛哭.
\newline
\newline
親愛的弟兄姊妹!我們在神面前都有虧欠和失敗,叫主傷心.我們今天來到主的面前,來到十字架前面,在各各他有赦免.雖然彼得像將殘的燈火,主耶穌卻沒有吹滅他；主加力量給他,教他回頭以後要堅固他的弟兄.
\newline
\newline


\newpage

\section{第46屆~港九培靈會~周主培~第九講}
\label{sec:1219}
\textbf{模範教會}
\newline
\newline
連結: \href{https://www.hkbibleconference.org/session-message/view/1219}{\texttt{https://www.hkbibleconference.org/session-message/view/1219}}
\newline
\newline
\hyperref[sec:1218]{< < < PREV SERMON < < <}
~
\hyperref[sec:toc]{[返目錄]}
~
\hyperref[sec:1221]{> > > NEXT SERMON > > >}
\newline
\newline
約翰福音 12:1-12:11
\newline
\begin{longtable}{cl}
\hline
\hline
章節 & 經文 (和合本修訂版)\\
\hline
12:1 & \begin{tabularx}{0.7\textwidth}{X} 逾越節前六天，耶穌來到伯大尼，就是他使拉撒路從死人中復活的地方。 \end{tabularx} \\ \\ \relax
12:2 & \begin{tabularx}{0.7\textwidth}{X} 有人在那裡為耶穌預備宴席；馬大伺候，拉撒路也在同耶穌坐席的人中間。 \end{tabularx} \\ \\ \relax
12:3 & \begin{tabularx}{0.7\textwidth}{X} 馬利亞拿著一斤極貴的純哪噠香膏，抹耶穌的腳，又用自己頭髮去擦，屋裡充滿了膏的香氣。 \end{tabularx} \\ \\ \relax
12:4 & \begin{tabularx}{0.7\textwidth}{X} 有一個門徒，就是那將要出賣耶穌的加略人猶大，說： \end{tabularx} \\ \\ \relax
12:5 & \begin{tabularx}{0.7\textwidth}{X} 「為甚麼不把這香膏賣三百個銀幣去賙濟窮人呢？」 \end{tabularx} \\ \\ \relax
12:6 & \begin{tabularx}{0.7\textwidth}{X} 他說這話，並不是關心窮人，而是因為他是個賊，又管錢囊，常偷取錢囊中所存的。 \end{tabularx} \\ \\ \relax
12:7 & \begin{tabularx}{0.7\textwidth}{X} 耶穌說：「由她吧！她這香膏本是為我的安葬之日留著的。 \end{tabularx} \\ \\ \relax
12:8 & \begin{tabularx}{0.7\textwidth}{X} 因為常有窮人和你們在一起，但是你們不常有我。」 \end{tabularx} \\ \\ \relax
12:9 & \begin{tabularx}{0.7\textwidth}{X} 有一大群猶太人知道耶穌在那裡，就來了，不但是為耶穌的緣故，也是要看耶穌使他從死人中復活的拉撒路。 \end{tabularx} \\ \\ \relax
12:10 & \begin{tabularx}{0.7\textwidth}{X} 於是眾祭司長商議連拉撒路也要殺了， \end{tabularx} \\ \\ \relax
12:11 & \begin{tabularx}{0.7\textwidth}{X} 因為有許多猶太人為了拉撒路的緣故，開始背離他們，信了耶穌。 \end{tabularx} \\ \\
[1ex]
\hline
\hline
\end{longtable}
近年來神自己開始了祂的工作,在許多城市中,我們看見許多青年弟兄姊妹們；或三五成群,或數十為伍；一同思想神的話,一同禱告,一同傳福音.這些弟兄姊妹的心裡有個問題:什麼叫做教會?我們是否教會?當我在北美各地工作時,我常遇見這樣的問題.
\newline
\newline
甚麼是教會?是否一所大禮拜堂；裡面且有一切的設備,才算是教會；或是只有兩三人在一起奉主的名聚會,就是教會呢?
\newline
\newline
一,教會是有主在其中的地方,如果沒有耶穌,雖有高大的禮堂,有一切好的設備,也不是教會.教會的先決條件,就是有主耶穌.有一次,英國女王到各城巡視,每到一城,城內各處均鳴鐘致敬；只有一城,女王到達時,沒有鐘聲；她的大臣覺得很失體面,大不滿意,找負責人交涉.那人說:「報告大臣,所以不打鐘有十個理由,第一因沒有鐘......」大臣說:「夠了,其他理由不必再提了.」一個教會最重要的,就是有耶穌基督在那裡.老底嘉雖為教會,但主耶穌卻站在門外叩門.多少名義上稱為教會的,實際上沒有主耶穌的地位,在他們中間沒有主耶穌的座位.教會中滿了一切的東西,但卻不見主耶穌.
\newline
\newline
二,教會是讓主安息的地方(約一11):「祂到自己的地方來,而自己的人倒不接待祂.」這叫主何等的傷心.我們的教會是否讓主得到安息；得到享受呢?很多時候,我們只顧自己的需要,不顧主耶穌的需要.
\newline
\newline
三,教會是榮耀主的地方(西一18):「祂也是教會全體之首；祂是元始......使祂可以在凡事上居首位.」在伯大尼地方,不但是有主耶穌基督的地方,是讓主常得安息的地方,也是榮耀主的地方.伯大尼地方,有人擺設筵席是為主的緣故,有人侍候也是為主的緣故,有人倒香膏抹耶穌是為主的緣故；所以教會是應該讓主得榮耀的地方,祂是那充滿萬有者所充滿的.
\newline
\newline
四,教會是一班從死裡復活的人聚集之處(約十二─末句):「就是祂叫拉撒路從死裡復活之處.」教會非社交之團體,非一班言論一致的人集合之處；乃是一班重生得救的人聚在一起,是耶穌基督所在之處,耶穌在每一個重生得救的人之心靈裡.所有心裡有耶穌的人,大家聚在一起,這就是教會.
\newline
\newline
教會裡彰顯榮耀的基督.基督徒要穿兩件衣服:一是全副的軍裝,一是上好的袍子,在家裡穿著上好的袍子,叫父親得著榮耀.在戰場上,我們應當穿起全副的軍裝.可惜今天有許多教會,在神的家中穿了全副的軍裝,弟兄姊妹彼此爭吵；沒有彰顯榮耀的基督,沒有叫主的心得安慰,仍然處於死未復活的光景.巴不得主耶穌藉祂復活的大能,在我們心裡彰顯祂的榮耀.如果今晚有那位還未重生得救的,你不要擔心你的生命快到盡頭,你應擔心你的生命從未開始.若是你的生命尚未開始,你應當擔心.若是你的生命未曾由死裡復活,未曾得著神的生命；你的心靈應當覺得懼怕,我們已到最後的機會了.
\newline
\newline
今天下午,當我們正要離開住處時,有人來敲門說:「請你下樓時通知接電話者,今晚我不接電話.」他邊說邊哭.後來我去敲他的門,聽見他在裡面哭,他開門,對我說:「今晚我不接任何電話,也不見任何人,我預備走了.」我為他很焦急,恐怕有不測之事發生.我對他說:「你要依靠神,要禱告.」他說:「我想神也無法拯救我,不過,請你放心,我不會自殺的,明晨再談吧!」我的心實在為他非常難過.這不過是我今天所見到的一人,相信如此之人在香港多著呢!他們流淚嘆息,心裡有創傷.
\newline
\newline
弟兄姊妹!教會就是需要把那些重生得救的人,從裡死復活的人,招集在神面前；不只是在外表,其實所需要的,乃是真真正正從死裡復活的人.
\newline
\newline
另一方面,我們看見許多弟兄姊妹,分工合作,同心合意.
\newline
\newline
教會有個錯誤的方針,以為單向某方面傳福音就夠了.有的向知識份子,有的向愚蠢的人,有的向富人,有向窮人的.很奇妙!教會如同神的院子,有高大的樹木,有美麗的花朵,有許多青草,各式各樣的人都有.我們不要單向某方面的人傳福音,不要以為中國人,應向同胞傳福音；如果中國人住在印尼,我們應向印尼人傳福音.住在日本,應該向日本人傳福音.住在美國或加拿大,應向當地的人傳福音.中國人所到之地,都有唐人街,但天堂卻沒有；如果你想在天堂找唐人街,你就必大失所望.
\newline
\newline
五,教會是包羅各式樣的人之團體(約十二1-11)這段聖經,給我們看見這個團體中有西門,有馬大,有拉撒路,有馬利亞,並且也有猶大.許多基督徒離開了自己的教會,他們到其他教會；因為覺得自己的教會,牧師,傳道人不好.有個基督徒請葛培理為他介紹一個最好的教會,因他覺得自己的教會不好.葛培理對他說:「世界上沒有完全的教會,萬一有,也因你加入而變成不完全.」不可從表面上看別人的教會好,許多大城市的基督徒,好像游擊隊,朝東暮西,這是很危險的.求主保守,讓我們所仰望的,非傳道人,非教會中的任何人,而是仰望為我們信心創始成終的耶穌.
\newline
\newline
我們看見教會中有西門,他本是長大痳瘋的.他有大房子,有錢財,他用錢財事奉主,這件事是不錯的,教會需要有些弟兄姊妹,把錢財奉獻在神面前；但很多時候,我們在神的家中,沒有做個好的管家.在這教會中,我們看見西門的錢財,也看見馬大的手,她在那裡事奉耶穌.很多人以為耶穌責備馬大,其實不是.耶穌的話是帶著嘆息,憐憫,及體諒的聲音說:「馬大!馬大!你為許多的事情忙碌,思慮,煩擾,但不可缺少的只有一件.」耶穌並沒有見怪她.當耶穌和十二門徒進到她家裡,難道一個身負家務的主婦,不手忙腳亂嗎?在這培靈會,我們看見有許多人坐在台上,還有許多無名的基督徒,他們在會前會期會後做了很多事情.我們需要馬大──需要弟兄姊妹為主勞碌.在這裡還有個拉撒路,他始終是一言不發；但有許多人為他的緣故回去信了耶穌.有的人天生沒有口才,很少發言；但他在神的面前,將復活的大能彰顯出來.電器之中,有無線電,電燈,電爐.我們需要無線電發聲音,電燈發光,但有時也需要電爐；雖然它無聲發出,但在它的裡面,滿有生命的能力.在我們的教會,有位姊妹,她不大說話；我們每個人都喜歡她,有事總去找她,當看見她時,心裡就有力量.我再次為蔡蘇娟女士感謝神!我和她住得很近,常有人去看她.她在暗房中已經住了四十多年,她今已八十歲了.人每當去看她,心裡就充滿能力；她並沒有對幾千幾百的人講道,可是她有個從死裡復活的生命,這生命滿有能力.我們在教會──神的家中,需要從死裡復活的生命；當人碰到你時,就覺得你裡面有一樣東西,是他所沒有的；當人接觸你之時,如同接觸到神的能力.
\newline
\newline
我們看見馬利亞,她把香膏倒在耶穌頭上,抹在耶穌的腳上,叫主的心得著滿足.馬太和馬可福音告訴我們,主耶穌有個命令,叫傳福音者無論往何處,都要述說這個女人所作的事,作個紀念.
\newline
\newline
何謂福音?福音是主耶穌所作的事.何謂奉獻?奉獻是人為主所作的事；這兩件事應連在一起.我們不但想到主為我們作了什麼,也要想到我們為主作什麼.主為我們釘十字架,為我們流血,為我們離開天上的榮耀；然而我們為主作了什麼?主的心願,不是叫我們單接受祂,為我們所做的工作在我們身上；更叫我們透過我們的身體,能夠在我們身上榮耀祂.
\newline
\newline
馬利亞的奉獻　我相信主一定預告馬利亞,關於將被釘十字架,被埋葬,三日復活了的事.馬利亞必為此反覆思想,應為主作些什麼,叫主的心得著滿足.耶穌說:「馬利亞是為我安葬之日作的.」馬利亞所作的,正合時宜；雖然有人在耶穌死後用香膏膏祂,但是已經太遲了.我永不忘記在上海,那時從東北有很多人逃難來；聚會時有位弟兄作見證,他邊講邊流淚,述說主如何恩待他,拯救他,保守他.他說他流淚,一方面是感恩地流淚,一方面是悔改地流淚.他原在東北有個大工廠,手下有許多工作人員；第一,他沒有把福音傳給他們,第二,當有力量奉獻時他不奉獻.現在要奉獻卻沒有錢財.他勸上海的弟兄姊妹,有機會為主工作時,應盡力去作；能夠奉獻時,要盡力奉獻.
\newline
\newline
今天在香港的弟兄姊妹,神給我們豐衣足食,正如末底改對以斯帖說:「焉知你得了王后的位分,不是為現今的機會嗎?」你可以服事主,可以奉獻,可以把自己交在主手中.否則,正如末底改對以斯帖說:「此時你若閉口不言,猶大人必從別處得解脫,蒙拯救,你和你父家,必至滅亡.」親愛的弟兄姊妹!神給我們有此機會,將祂自己的心願向我們啟示.最近幾年,我們看見許多教會,不論在北美,菲律賓,香港,有許多中國的弟兄姊妹奉獻錢財,為著國外的傳道.我心裡有個很大的負擔,覺得單奉獻錢財不夠,最要緊把你的心奉獻給主,莫以錢財代替我們的心.我們只注重國外傳道,然而在香港本地,我們的前後左右,有千萬人他們每天向著無神的地方前進.今晚我的心一直記念剛才和我講話的人,巴不得他能繼續存留世上,能聽見神自己的話；他對我說:「我不會自殺的,明天再談.」求神憐憫我們,開我們的眼睛,看見許多有需要的生命.耶穌說:「趁著白日,我們必須作那差我來者的工；黑夜將到,就沒有人能作工了.」
\newline
\newline
耶穌說馬利亞所作的是件美事,逾越節前六日,羔羊將被宰殺的時候,馬利亞作了件極有意義的事.她的奉獻是有心的奉獻,且是切切實實的奉獻.她用價值卅兩銀子的真哪噠的香膏,完全澆在耶穌的身上.哪噠香膏極其貴重,卅兩銀子在當時的世代,可用以買個奴僕.猶大就是用卅兩銀子把耶穌賣了.馬利亞用她所有的全部獻給主.很多時候,人要感嘆我們為耶穌所作的太過浪費.我讀中學時,有位級任待我非常好,常幫助我讀書,有時帶我出去散步.當我將要畢業時,老師問我畢業後擬作何事,我說要作傳道,他很驚奇地說:「為何要作傳道呢?」我就把我得救和奉獻的經過告訴他,他不斷搖頭嘆息；因為他認為是太可惜了.有的人以為到禮拜堂是給耶穌面子,覺得是一種浪費,為主工作也是浪費.他們沒有認識我們的主全然可愛,祂比世人更美,在祂口裡滿有恩惠；正如雅歌書說:「你的名如同倒出來的香膏,眾童女就快跑跟隨你.」人若對主有真正的認識,絕不至有浪費之感.當一個男子向一個女子求婚,把最好的鑽戒交在她手中,他認為全世界惟這女子是配.
\newline
\newline
我總是感到時間不夠,為主作的太少,每次為主去作總要好好地準備自己；因為我是代表我的君王,代表我的主,沒有一件東西對主是太好的.(請原諒我略題自己的事)我有五個孩子,三個已大學畢業了,其中一個已經結婚,兩個進神學.有一次我帶個兒子進城去,他問我能否順道到某處,我明白他的心,就帶他去.那家人對我們款待週到,他們問我的兒子畢業後擬作何工,我兒子對他們說要做傳道.因此事情遂有改變,現在那女子已經結婚了.他們認為做傳道是沒出息的.弟兄姊妹!多時我們認為奉獻給主是浪費的.記得戴德生說:「如果我有千條生命,我願意為中國人傳福音而擺上；如果有多少金鎊,我都願意為著在中國傳福音而花費.」
\newline
\newline
馬利亞的奉獻,是毫無保留的奉獻.她把香膏完全獻給主,而無留下一點為自己.據傳說,猶太女子,常預備香膏,為自己和新郎在結婚之日而用,可是馬利亞卻是為主的安葬獻上香膏.神非看我們奉獻多少,乃是看我們留下多少.聖經上記載那窮寡婦奉獻了兩個小錢,得著主的稱讚；因為她把所有養生的全部都獻上了,絲毫不為自己留下.許多時候,我們實在為自己留下太多,獻給主只不過少而又少,沒有把所有的獻在主的面前.馬利亞非但用香膏抹耶穌的腳,用眼淚哭在主面前,用頭髮擦主的腳,她完完全全把自己放在主的面前.
\newline
\newline
馬利亞的奉獻,叫主的心得著滿足,因彼得三次不承認主,所以耶穌孤單地上各各他；不過,我相信耶穌在各各他山上,並不孤單；因為馬利亞倒在祂身上的香膏,不但當時房子裡滿了香氣,一直今天,我們仍聞到香味.我相信當時耶穌的身上,香膏芳芬滿溢.在各各他山上,耶穌聞到香膏之香陣陣飄來,叫主的心得著滿足.
\newline
\newline
弟兄姊妹!我們不但要遵行主的旨意,更要體恤主的心意,叫我們能滿足主的心.我們的奉獻,最要緊的是叫主的心滿足,雖然別人不了解,不同情,只要主的心滿足.當猶大冷言冷語批評馬利亞時,她默然置之；但耶穌說:「不要難為她.」有時你被冤枉批評；在教會中為主所受的勞苦,無人欣賞,無人稱讚,反而有人說閒話.你不必難過,也不必說話,有主說話就夠了,主明白就夠了,主的心得滿足就夠了.
\newline
\newline
有主在的地方就是教會.你我心裡都有主,互相配合起來,成為一個身體；要彼此同心合意,才是真正的教會.
\newline
\newline
教會是主得安息的地方,我們願意給主有安息的地方嗎?教會是主得榮耀之處,當主得著榮耀時,我們的心也滿足了.啟示錄告訴我們,當四活物敬拜的時候,主的心得著滿足；他們把冠冕擺在主的腳前,讓主得榮耀,坐寶座,戴冠冕,掌王權.
\newline
\newline
求主憐憫我們!下個雄心大志,讓主的心得著滿足.
\newline
\newline


\newpage

\section{第46屆~港九培靈會~周主培~第十講}
\label{sec:1221}
\textbf{全家事奉}
\newline
\newline
連結: \href{https://www.hkbibleconference.org/session-message/view/1221}{\texttt{https://www.hkbibleconference.org/session-message/view/1221}}
\newline
\newline
\hyperref[sec:1219]{< < < PREV SERMON < < <}
~
\hyperref[sec:toc]{[返目錄]}
~
\hyperref[sec:1184]{> > > NEXT SERMON > > >}
\newline
\newline
約書亞記 24:14-24:28
\newline
\begin{longtable}{cl}
\hline
\hline
章節 & 經文 (和合本修訂版)\\
\hline
24:14 & \begin{tabularx}{0.7\textwidth}{X} 「現在你們要敬畏耶和華，誠心誠意事奉他，除掉你們列祖在大河那邊和在埃及事奉的神明，事奉耶和華。 \end{tabularx} \\ \\ \relax
24:15 & \begin{tabularx}{0.7\textwidth}{X} 若你們認為事奉耶和華不好，今日就可以選擇所要事奉的：是你們列祖在大河那邊所事奉的神明，或是你們所住這地亞摩利人的神明呢？至於我和我家，我們必定事奉耶和華。」 \end{tabularx} \\ \\ \relax
24:16 & \begin{tabularx}{0.7\textwidth}{X} 百姓回答說：「我們絕不離棄耶和華去事奉別神。 \end{tabularx} \\ \\ \relax
24:17 & \begin{tabularx}{0.7\textwidth}{X} 因為耶和華－我們的神曾領我們和我們的祖宗從埃及地為奴之家出來，在我們眼前行了那些大神蹟，並在我們所行的一切路上，和所經過的各民族中保護了我們。 \end{tabularx} \\ \\ \relax
24:18 & \begin{tabularx}{0.7\textwidth}{X} 耶和華又把各民族和住此地的亞摩利人都從我們面前趕出去。所以，我們也必事奉耶和華，因為他是我們的神。」 \end{tabularx} \\ \\ \relax
24:19 & \begin{tabularx}{0.7\textwidth}{X} 約書亞對百姓說：「你們不能事奉耶和華，因為他是神聖的神，是忌邪的神，必不赦免你們的過犯罪惡。 \end{tabularx} \\ \\ \relax
24:20 & \begin{tabularx}{0.7\textwidth}{X} 你們若離棄耶和華去事奉外邦的神明，耶和華在降福之後，必轉而降禍給你們，把你們滅絕。」 \end{tabularx} \\ \\ \relax
24:21 & \begin{tabularx}{0.7\textwidth}{X} 百姓對約書亞說：「不，我們要事奉耶和華。」 \end{tabularx} \\ \\ \relax
24:22 & \begin{tabularx}{0.7\textwidth}{X} 約書亞對百姓說：「你們選擇耶和華，要事奉他，你們自己作證吧！」他們說：「我們願意作證。」 \end{tabularx} \\ \\ \relax
24:23 & \begin{tabularx}{0.7\textwidth}{X} 「現在，你們要除掉你們中間外邦的神明，專心歸向耶和華---以色列的神。」 \end{tabularx} \\ \\ \relax
24:24 & \begin{tabularx}{0.7\textwidth}{X} 百姓對約書亞說：「我們必事奉耶和華---我們的神，聽從他的話。」 \end{tabularx} \\ \\ \relax
24:25 & \begin{tabularx}{0.7\textwidth}{X} 那日，約書亞就與百姓立約，在示劍為他們制定律例典章。 \end{tabularx} \\ \\ \relax
24:26 & \begin{tabularx}{0.7\textwidth}{X} 約書亞把這些話寫在神的律法書上，又拿一塊大石頭立在橡樹下耶和華聖所的旁邊。 \end{tabularx} \\ \\ \relax
24:27 & \begin{tabularx}{0.7\textwidth}{X} 約書亞對眾百姓說：「看哪，這石頭可以向我們作見證，因為它聽見了耶和華所吩咐我們的一切話；這石頭將向你們作見證，免得你們背叛你們的神。」 \end{tabularx} \\ \\ \relax
24:28 & \begin{tabularx}{0.7\textwidth}{X} 於是約書亞解散百姓，各自回到自己的地業去了。 \end{tabularx} \\ \\
[1ex]
\hline
\hline
\end{longtable}
聖經是神的啟示.啟示中；有神直接的啟示,也有神間接的啟示.在舊約聖經中,我們常常看見:「耶和華如此如此說.」在新約中,主耶穌常說:「我實實在在的告訴你們.」這是神直接的啟示.所謂神間接的啟示,乃是藉著人的經驗,和人心的感動說出話來,藉著人內心的體驗說出話來,如:「耶和華是我的牧者,我必不至缺乏；祂使我躺在青草地上,領我在可安歇的水邊.」這並非神怎樣說,乃是體驗到神恩典的人從他心裡說:「除了你以外,在天上我還有誰呢?除了你以外,在地上我也沒有可愛慕的.」「我的心哪!你要稱頌耶和華,不可忘記祂的一切恩惠；祂赦免你的一切罪孽,醫治你的一切疾病.祂救贖你的命脫離死亡,以仁愛和慈悲為你的冠冕.」「主啊!你世世代代作我們的居所.諸山未曾生出,地與世界你未曾造成,從亙古到永遠,你是神.」「我們一生的年日是七十歲,若是強壯可到八十歲；但其中所矜誇的,不過是勞苦愁煩；轉眼成空,我們便如飛而去.」「人為婦人所生,日子短少,多有愁苦,出來如花,飛去如影,不能存留.」這些話語,不是神直接說的,乃是從神那裡,有所領受,有所體驗的人,照神的心意向我們講的.在新約裡,保羅曾如此說；「我對其餘的人說,不是主說.」(林前七12)此乃神間接的啟示.當人透過神給他的遭遇,把客觀的信仰,成為他主觀的事實和經驗,這是件很重要的事情.
\newline
\newline
約書亞在他年紀老邁時說:「至於我,和我的家,我們必定事奉耶和華.」現在我們要思想三件事:
\newline
\newline
(一)人與神的關係.(二)全家的光景.(三)家庭的目標.
\newline
\newline
約書亞說:「至於我......」這是我個人在神面前的光景,不管別人如何.(請原諒我用基督教這幾字來講)基督教和神的關係,是極注重個人和神的關係；不但要注重個人修身,齊家,治國,平天下；基督徒更應想到個人與神的關係；我們的信仰,是個人的信仰.神惟有兒子,沒有孫子,沒有媳婦,也沒有女婿.你別想你父親是基督徒,是神的兒子,你就是神的孫子；別以為你丈夫是神的兒子,而你是神的媳婦；也別以為你妻子是基督徒,你就是神的女婿.基督徒應當直接和神自己發生關係,可惜教會中,我們聽道是為別人,認罪是為大家,禱告是為自己.有時我們聽道,心裡就覺得所聽的道；都是針對我的丈夫講的,如果他也在這裡聽多好呀.我們很聰明,認罪時都是說:「神啊!我們是罪人,我們都有驕傲,我們都有虧欠,我們都撒謊.......」所有的罪,都是「我們」犯的,簡直把別人都扯進去.至於禱告的時候,專求神賜福給我,恩待我,保佑我,賜福我的生意,兒女.......所有的祈求,都是為「我」.求神憐憫我們,叫「我」個人真正與神發生關係.
\newline
\newline
約書亞說:「至於......,」這是句非常勇敢的說話；雖然這是句孤單的講話,別人都不跟隨,「至於我」我願意跟隨主；雖然世界上大家都是「醉」了,只有一個人「醒」；都是「濁」的,只有一個「清」.弟兄姊妹!我們在這環境中,靠主恩典才能站立得住.一個大丈夫不是無所不為,大丈夫有所不為；能夠有所不為,才能有所為.雖然人都離開神,至於我,我願意走十字架的道路；這路是窄狹的,孤單的.人常常不了解你,不同情你.十字架的道路,只能孤單地走,那怎麼辦呢?有主同走,主在我裡面,我在主裡面,我獨自與耶穌同走.
\newline
\newline
弟兄姊妹!多數人的意見不一定是對的,當人高聲讚你的時候,就是你最危險之際.我們不要隨波逐流,死了的魚才隨波逐流；活的魚,乃是逆流而上.我們須在這時代中,彰顯神的榮耀,叫我們能夠作中流砥柱,逆流而上.
\newline
\newline
當以色列人將出埃及時,摩西一而再,再而三的要求法老,容百姓離開埃及,去事奉耶和華.約書亞離開埃及,為要事奉那位又真又活的神.約書亞到迦南地,當他年老時,他說:「至於我,和我的家,我們必定事奉耶和華.」他認定人生的目標,他向著標竿直跑,不偏左右,他知道神的心意.許多人跟隨主,不偏不移,緊緊的跟隨；一旦名成利就,卻慢慢的轉移了.約書亞站在神的一邊,無論別人如何,「至於我,和我的家,我們必定事奉耶和華.」有句英文話說:「一個人與神在一起,他就是大多數.」只要你與神站在一起,你就是絕大多數.但以理一人站在神面前,後來他被丟在獅子坑,他卻成了大多數.他的三個朋友,被丟在火yiu中；王看見其中有四個人,因為神與他們同在.只要站在神的一邊,你不要懼怕.
\newline
\newline
世人都站在寬闊的路上思想人生的意義.他們在那裡,只是存在而非存活；他們醉生夢死!施洗約翰在曠野,站在神的旁邊；大聲疾呼「天國近了,你們應當悔改.」挪亞出來傳道時,人都不接受,但他站在神的一邊.馬丁路德起來改教時,受了種種的迫害.衛斯理約翰,起來復興英國教會時,遭遇種種的攔阻.戴德生到中國傳福音時非常孤單.約翰加爾密在日內瓦城中,為了要改革宗教遇到許多的困難.宋尚節博士出來為主工作時,人都譏笑他.王明道先生站在神面前,這班人的臉都獅子一樣；他不看環境,不看方向,只迎著標竿直跑.親愛的弟兄姊妹!我們在神面前,是否這樣的決定?我們向著神宣告,向著撒旦警告,向著世人傳告說:「至於我,我們的家,我們必定事奉耶和華.」
\newline
\newline
有的基督徒好像蝙蝠一樣,他的樣子像老鼠,飛起來像鳥；即不像動物,又不像飛禽.很多基督徒,進禮拜堂像隻羊,出禮拜堂卻像頭狼；羊皮掛在大門上,假冒偽善不能進天堂.
\newline
\newline
數年前在印尼,聽說有個天主教徒,允諾等他所養的羊生了小羊,要奉獻一隻給他們的神父.有一天,他果然帶了一隻小羊,放在單車後面的籃裡,預備送去；到了半路,停下來,把單車和小羊置於茶室門口.他進去喝咖啡,在裡面他遇見一個朋友,當他朋友知道他有隻小羊,在門口是預備奉獻給天主教堂的神父,便偷偷地把家裡的一隻小狗換了小羊.等他喝完咖啡,遂即騎車趕路,毫不在意,到了教堂,神父出來相迎,十分謝謝這位好信徒,遵守諾言奉獻小羊一隻；不料打開籃子,發現卻是一隻小狗.神父很生氣,責他玩弄手段,存心欺騙,奉獻不忠.真是叫這教友如啞子吃黃蓮,垂頭喪氣地走了.在回程中,他再到那茶室喝咖啡,見到剛才那朋友,他怒氣沖沖,一言不發；朋友知情,急將小狗取出,小羊仍放入籃中.他回到家裡,指責他妻子糊塗從事,害他在神父面前丟臉,他妻子被責莫明其妙.他把籃打開要證實他妻子錯,那知抱出來卻是一隻小羊.弄得他啼笑皆非,說:「羊啊!你要做羊就做羊,做狗就做狗；為何一下子做羊,一下子做狗?」許多基督徒,他們的改變太大,在人面前,在教會中,在無人看見之處,表現各有不同.中國有兩個字「慎獨」,君子要慎獨.人常在大庭廣眾之前持守謹慎,但在獨處則易於放肆.
\newline
\newline
約書亞堅心決意說:「至於我,和我的家,我們必定事奉耶和華.」乃因他看見有一個人手裡有拔出來的刀.約書亞問他說:「你是幫助我們呢?是幫助我們敵人呢?」他回答說:「不是的,我來是要作耶和華軍隊的元帥.」(書五13-15)弟兄姊妹!我們常在禱告時有人習慣,未經思想地說:「主啊!求你幫助我.」這話不錯,但不只是幫助的問題,主是作我們的主.好多時,我們自己作主,以為求主幫助我們就夠了；以為自己能做,有主幫助更好.請記得!不是主來幫助我們,乃是主來引導我們.我們稱耶穌為救主,因祂救贖我們也作我們的主；可是我們並沒有在凡事上,讓主居首位.有的人僅一半讓主居首位.有的人在大事上讓主管理,小事自理.有的人小事交給主,如一般頭痛,咳嗽之類的小病,禱告求主醫治；但遭遇大病,就在非找醫生不可.我並非說,有病不需要醫生；乃是說,我們應在凡事上尊主為大.
\newline
\newline
人怎能服事主呢?必須面對面看見主,當你面對主面之後,你說:「從前我愛沉迷在繁華的夢裡,驕癡無知,舊事切莫重題,叫我從今之後,緊緊的跟隨你.」不要怕你的生命快到盡頭,應該擔心你的生命尚未開頭.我們應當有個新的開始.如何開始呢?就是要面對面看見神.神說:「誰看見我,他就要死.」老舊人死了,老我死了,肉體死了；在主的光中,老的,舊的都不能站立.
\newline
\newline
教會有個口號:「全家歸主」,聖經說:「當信主耶穌,你和你一家都必得救.」如果止於此,這是不夠,我們應當作到全家事奉.今天中國教會最大的危機,只想到得救的問題,得救以後,神將我包裝好,寫上「只為出口」這樣就夠了,這些貨專為上天堂.如果神救了我們,我們專為等死後上天堂,那不如早死早好.正像一個傳道人,想把自己的講章寫下來:有位信徒說:「牧師啊!我很喜歡看你的講章.」牧師說:「恐怕早哩,要等我過世以後才出版.」那位太太說:「巴不得越快越好.」弄得牧師有苦難言!
\newline
\newline
香港的教會,不要單講全家歸主,如果全家歸主,這是對的.全家歸主就是全所有的,都奉獻在主前；我不是只講錢財方面的,錢是身外物,單奉獻錢還不夠；全家都奉獻十分之一也不夠,這非神的心願.我再說:不是單單全家得救,乃要全家事奉,這是很重要的.教會不興旺,是因為許多基督徒在神面前,與神的工作脫了節.
\newline
\newline
有位朋友從印度的南部來,他說該地基督的道理非常廣傳.在村莊裡,每年有家庭燃燈的節期,全家發光；他們先到禮拜堂,天未黑便唱詩,他們聚會的時間很長,一齊唱詩,禱告,讚美；天黑了,他們繼續禱告唱詩,這時,有位長老來把當中的燈點亮,再繼續聚會.然後每家的家長,相繼從當中的燈火,引亮自己的燈.到散會時,每家的燈都點亮了,邊走邊唱詩,回到家裡,每人點亮一枝臘燭,開始家庭聚會.整個村莊,非常美麗.弟兄姊妹!我們的家,就是我們的教會；我們的教會,就是我們的家.傳道人們常常說:盼望弟兄姊妹,將教會成為我們的家；但從未說我們的家,成為教會.聖經說:「家中的教會」,每個教會應該是信徒的家,每個信徒的家都應該是神的教會.
\newline
\newline
約書亞說:「至於我,和我的家,我們必定事奉耶和華.」這句話使我心裡深受感動,同時也得到很大的安慰.當時約書亞很忙,從約書亞第一章開始,我們看見他東征西伐,到處奔跑,實在無暇顧及他的家,但神看顧他的家.我在這裡很慚愧的講一件事:我很少在家,特別是假期,多數在學生中間工作,每逢放假,有各種聚會,所以不能在家和孩子們一起.有一年的聖誕節,我立志要和孩子們一起過節,有個孩子,他感謝神說:「我們的浪父回家了.」
\newline
\newline
很多時候,我心裡很難過,(請原諒我講一點自己的事)十幾年前,我全家剛到美國；小孩子每到新的地方,就學當地的話吵架.在日本住了三年,他們學了日本話吵架；到了美國,又學了壞習慣.我心裡很難過!有一天早晨,家庭禮拜之後,我告訴他們說:「如果你們這樣,我寧可離開這裡,到非洲去；那是落伍未開化的地方.只要你們愛神,比住在美國好,比你們離開神為強.」我理直氣壯地說了這話.後來孩子們都上學去了,我在個人靈修時,神責備我說:「你到美國是為我還是為孩子呢?如果你到美國是為了孩子,你就到非洲去；如果是為我,我會照顧你的孩子.」於是,我跪下悔改,對神說:「神啊!我在這裡,非為孩子的緣故,乃是為你的緣故.」親愛的弟兄姊妹!我們為孩子打算是應該的；但應讓神居首位,如果你把孩子放在第一,沒有讓神居首位,當心有一天你要失去你的孩子.
\newline
\newline
記得家母曾寫信給我說:「孩子啊!感謝神,我們作了極好的決定,就是把你放在神的手裡；這樣,我們永遠得著你.」許多基督徒為父母的,失去了孩子,乃是因為沒有把孩子放在神的手中.人常望子成龍,這是沒錯的；但很危險!人常教訓孩子說「吃得苦中苦,方為人上人.」如果你的孩子在人之上,那麼誰在人之下呢?這是很危險的,千萬別盼望你的兒女作為人上人,乃要盼望你兒女作為神的兒女,做一個人就夠了.我曾對我的孩子說,你們要作什麼都好,我不勉強你們；如果你當個倒垃圾的,我也不怪你,只要你能打掃乾淨,盡基督徒所應做的,我都引以為榮.人都望子成龍,乘風破浪,今日世上的浪,實在太多了.
\newline
\newline
弟兄姊妹!你要把孩子交在神手中,求神憐憫我們!我們中間已經為父母的,應有如亞伯拉罕的心願,把以撒獻在祭壇上；像哈拿的心願,把撒母耳放在神手中；像摩西的父母,把兒女交在神手裡.未婚的弟兄姊妹!請記得,上一代的人,是望子成龍,養子防老,積穀防飢.我們千萬不要自私自利,我們要養兒女榮耀主,盼望兒女服事主,這是很重要的.
\newline
\newline
約書亞有非常確定的憑據說:「至於我,和我的家,我們必定事奉耶和華.」至於你,你的孩子在哪裡?有許多為父母的,夜不入寐,因他沒有把兒女帶到神的面前,他不知他的兒女在哪裡.約書亞卻留下服事神的榜樣.他知道他的孩子們,他了解他們,他不自私.以他的地位,他是個政治領袖,他很可以說:「至於我,和我的兒子,要繼續統治以色列國.」但他卻說:「至於我,和我的家,我們必定事奉耶和華.」有許多教會,成為私人的財產；許多的事不是屬於神的工作.我並非說,兒子不能繼任老子之職,問題不是因他是兒子,而是他夠資格,問題在於你有否個人的野心?
\newline
\newline
(路一74-75):「叫我們既從仇敵手中被救出來,就可以終身在祂面前,坦然無懼的用聖潔公義事奉祂.」這節聖經清楚告訴我們,一個人離開埃及,脫離法老轄制；就是得救重生的人,才能夠事奉神,才能夠來到神的面前.當我來到神面前時,我的心恐懼戰兢,每次站起來講道時,我覺得不是在弟兄姊妹面前講道,而是在神面前工作；所以我要用恭敬誠實的心,把這個時間交在神的手中.我們應當作個全時間的基督徒,絕不要作部份時間的基督徒.事奉神,並非每個人站在講台上,每個人出來做傳道.聖經說,我們不能事奉神,又事奉瑪門.難道我們不可做生意嗎?不,我們乃是需要基督徒的商人.當我的小兒子,讀法律的時候,我感謝神；我們需要基督徒的律師,如果神用著你,你可以在各方面作工；我們需要基督徒的醫生,工程師,店員,清潔工人倒垃圾者,電器技工,教授,法官......我們需要基督徒在每個角落.不是單站在台上講道的人,才是服事神；也不是每週兩小時主日聚會,來服事神.你若不服事神,你服事誰呢?不是服事瑪門,服事撒旦,就是服事自己的肚腹.
\newline
\newline
聖經說,我們都是服事神的.我常有機會對一班將出來做宣教士的弟兄姊妹講話；其中有的是教授,有的是醫生.我說:「你們千萬別以教授的身份,去當宣教士,以醫生的身份,去當宣教士,而應當以宣教士的身份,教書和行醫.每個弟兄姊妹,我們都是宣教士,都是服事神的.無論在公司做事,在學校做事,或在家庭負責家務,應以事奉神的心情去做這些事.」有個教會,是由信徒們輪流打掃禮拜堂；有次打掃的人做得非常乾淨,只有鋼琴的角落從來無人打掃.此人也置之不理,他回家去了；但心裡為此不安,他就再次去打掃乾淨.然後他跪下禱告說:「神啊!這個地方是為你打掃的；因為人看不見,惟你看見了.」弟兄姊妹!我們在任何環境服事神.教養兒女是為著神的緣故,賺錢是為神的緣故.基督徒應滲透到社會各部分；為主而活,靠主而活.
\newline
\newline
要敬畏神(亞廿四14)如果我們服事神不存敬畏的心,那就不是誠心誠意的服事.敬畏神,服事祂,這是榮耀體面的事情；當人服事萬王之王,這是極大的榮耀.
\newline
\newline
要誠心實意的事奉神　我們不能虛假的服事神,要誠心實意的事奉祂；這是我們真正服事神的態度.
\newline
\newline
要除去各樣的偶像,用聖潔公義事奉神前天我到某家,看見有個牌子寫著:「整潔為家庭之美.滿足為家庭之福.好客為家庭之樂.聖潔為家庭之冕.」一個家庭若沒有聖潔,就沒有聖潔的子女.求神憐憫我們!神並非看我們器皿的大小,貴重或卑賤.「人若自潔,脫離卑賤的事,就必作貴重的器皿；成為聖潔,合乎主用,......」(提後二21)弟兄姊妹!巴不得我們手潔心清的服事主.
\newline
\newline
要甘心專心服事神,約書亞甘心樂意,專心一意的服事神.一個將軍可以有千萬個的俘虜,但妻子只有一個.俘虜是用力量得來的,而妻子是用愛得來的.神用愛把我們得來,叫我們甘心的事奉祂,專心的服事祂.
\newline
\newline
叫神的心得到滿足,很多人樂意工作,目的叫自己快樂,他忘記了,神快樂更為重要.如果神快樂,我們也滿足了.
\newline
\newline
事奉最要緊的,就是:不但遵行神的旨意,也要體會神的心意.
\newline
\newline


\newpage

\section{第46屆~港九培靈會~鮑會園~第一講}
\label{sec:1184}
\textbf{知識的根基}
\newline
\newline
連結: \href{https://www.hkbibleconference.org/session-message/view/1184}{\texttt{https://www.hkbibleconference.org/session-message/view/1184}}
\newline
\newline
\hyperref[sec:1221]{< < < PREV SERMON < < <}
~
\hyperref[sec:toc]{[返目錄]}
~
\hyperref[sec:1185]{> > > NEXT SERMON > > >}
\newline
\newline


\newpage

\section{第46屆~港九培靈會~鮑會園~第二講}
\label{sec:1185}
\textbf{真正的知識}
\newline
\newline
連結: \href{https://www.hkbibleconference.org/session-message/view/1185}{\texttt{https://www.hkbibleconference.org/session-message/view/1185}}
\newline
\newline
\hyperref[sec:1184]{< < < PREV SERMON < < <}
~
\hyperref[sec:toc]{[返目錄]}
~
\hyperref[sec:1186]{> > > NEXT SERMON > > >}
\newline
\newline
歌羅西書 1:9-1:14
\newline
\begin{longtable}{cl}
\hline
\hline
章節 & 經文 (和合本修訂版)\\
\hline
1:9 & \begin{tabularx}{0.7\textwidth}{X} 因此，我們自從聽見的日子就不住地為你們禱告和祈求，願你們滿有一切屬靈的智慧和悟性，真正知道神的旨意， \end{tabularx} \\ \\ \relax
1:10 & \begin{tabularx}{0.7\textwidth}{X} 好使你們行事為人對得起主，凡事蒙他喜悅，在一切善事上結果子，對神的認識更有長進。 \end{tabularx} \\ \\ \relax
1:11 & \begin{tabularx}{0.7\textwidth}{X} 願你們從他榮耀的權能中，得以在一切事上力上加力，好使你們凡事歡歡喜喜地忍耐寬容， \end{tabularx} \\ \\ \relax
1:12 & \begin{tabularx}{0.7\textwidth}{X} 又感謝父，使你們配與眾聖徒在光明中分享基業。 \end{tabularx} \\ \\ \relax
1:13 & \begin{tabularx}{0.7\textwidth}{X} 他救了我們脫離黑暗的權勢，遷移到他愛子的國度裡。 \end{tabularx} \\ \\ \relax
1:14 & \begin{tabularx}{0.7\textwidth}{X} 藉著他的愛子，我們得蒙救贖，罪得赦免。 \end{tabularx} \\ \\
[1ex]
\hline
\hline
\end{longtable}
昨天已提過,保羅的每封書信,都是為著一個很實際的目的而寫的.歌羅西書亦不例外,因在那時可能有人傳異端的教訓；不單是傳講異端,而且在第二章中可以看見,已有許多人要遵守此教訓了.當時歌羅西教會,有很大的需要,保羅才寫此信.並且很直接地,為了解決重要的問題.歌羅西教會有許多人走上了異端之途,保羅一開始就要針對這問題；盼望能將這些人挽回過來.
\newline
\newline
剛才所看的這段經文,是保羅的禱告,這禱告是為了要挽回這班錯誤的弟兄而發出的.他不但祈禱給神聽,若是只給神聽,就不必寫給歌羅西教會了.另一方面乃是藉此禱告,教導歌羅西教徒明白真理；保羅所用的方法非常的美.他將真理告訴給他們,又藉真理改正他們思想上的錯誤.但他卻沒有直接地指正他們錯在何處,他先把他們帶到一個崇高的氣氛裡,帶領他們一同到神面前禱告；且在禱告中,他把所要說的真理直接地擺在他們面前.這種勸勉或指正的方法,給我們一個很好的榜樣.勸勉及指正別人的錯誤,是一件最難的事.保羅卻用很美妙的方法,作一個榜樣.歌羅西人犯了錯.保羅沒有直接指正他們的錯.沒去辯論或責備他們；還要帶領他們同到神面前一同禱告.我們需知道,和別人辯論是沒多大的好處,就算理由充足或很有力量,使對方無法辯駁；可能在理由與詞句上你可以得勝了,但實際上卻不能得到他的心.與別人講述真理,目的並非要勝過對方,而是要對方肯接受你所講的真理.真理不是憑你的詞句辯論可以得勝的,所以當我們向別人辯論真理的時候,可以善用護教學或衛道學,這些書籍對你很有幫助.但單憑著那些學問和理論,不能帶領人信耶穌.藉著護教學或衛道學,或許會將些難處除去了；但真正能將人帶到主面前的,乃是直接把主耶穌介紹給人.不是消極地與他辯論,而是積極地將他們的需要,擺在他們面前；保羅就是這樣作了我們的榜樣.
\newline
\newline
保羅為他們禱告的內容,是要叫他們明白真理,明白知識.因當時傳異端的人,非常注重知識,他們也高舉知識.保羅對他們說:你們看重知識,那是對的.現在我要把那真正的知識介紹給你們.這是最深奧最高尚的知識.保羅在此處所用「知識」一字很特別,這字不單是代表普通知識而已；而且還是兩字連在一起,可以說是高級的深奧和真正的知識.
\newline
\newline
(一)得真正知識的根基
\newline
\newline
「願你們在一切屬靈的智慧悟性上,滿心知道神的旨意.好叫你們行事為人對得起主,凡事蒙祂喜悅；在一切善事上結果子漸漸的多知道神.」(西一9-10)
\newline
\newline
照上述經文簡單一點來說,就是你們若要得真正的知識,就必需在善事上結出果子來,在行為上對得起主.這是聖經上所說得知識的特別方法,亦可說得真理的一方法,這是極重要的.在約翰福音七章裡,主耶穌講過一段特別的話:「人若立志遵著祂的旨意行,就必曉得這教訓是初於神,或是我憑著自己說的.」(約七17).在屬靈的知識和真理上,並不是憑著你的腦子,能了解多少:你若說把所有知識都明白了,得到很深奧的道理了,那麼自然就有長進了.聖經很忠誠地要把一切真理給我們,當我們知道了一部份真理後,就要按著我們所知道的真理去行.按著去行後,藉著我們靈裡的體驗,然後我們知道我們所行的是真理與否.你們若立志遵著真理去行,就曉得這教訓是出於神,或是出於我自己.所以首先必需有願意按著神的旨意行的心,對神旨意有根本上的順服,這是我們要明白真理的一個最基本根基.按人的知識而論,以為我明白了後就會知道,知道後才可以遵行.但聖經的真理卻將次序倒過來:我知道了,按所知的去遵行了,然後我才明白.我們若知道了一些真理,若不用經驗去體驗這真理,這真理就只不過是在你腦子裡思想,只是頭腦上的知識而已.我們用生活去體驗過,按著真理去遵行；然後可知此真理是否真的.我們所需要的,那是先知道所有的真理,乃是我已經知的真理,是否按著去行?若沒有按所知的真理去行,而想多知道真理,也許神不會讓我們知道得更多.為什麼有人讀聖經讀了一生之久,但在屬靈的真理上沒有更深的亮光?這是說在詞句上他知道了許多,在文法上也知道許多；但在靈裡的領受,卻好像是停止了!已知的真理不去行而想要再追求,那麼就算神再多給你,也沒用處!因為這真理在你身上不能發生作用,所以不必再給你了.人若立志遵著祂的旨意行,就可知道這教訓是出於神,或是出於自己.
\newline
\newline
(二)真理的主要基本內容
\newline
\newline
真理的主要基本內容,就是知道神的旨意.真正的知識,和最高尚的知識,無非是使我們知道神的旨意.
\newline
\newline
保羅為歌羅西人祈求,就是求神使他們知道神的旨意.在(一9)中說:「願你們在一切屬靈的智慧悟性上,滿心知道神的旨意.」這段經文中,有好幾處很難翻譯,這是其中之一句,原來的意思並非如此.原本的意思,保羅是這樣的:在你們智慧和悟性上,能夠知道神一切的旨意.我們今天要追求的,是神一切的旨意,或是神的每一個旨意.明白神在每一件事上,祂的心意是什麼?在每一件事上,祂的啟示是什麼?我們要知道一點神的旨意,可以說是容易的；或者要知道一點我喜歡的旨意,也比較容易.我容易領受的神的旨意,因這樣的事我容易做到?這是我喜歡的,對我環境來說是適合的.但保羅卻說不是單知道一些就算了?而是要知道神每一個旨意.神在每一件事上的旨意是甚麼?這旨意並非單指我個人的生活,我應當如何做?尋找何種職業?或去什麼地方?當然,這些我們該去追求.神也必在這些事上啟示我們,但還有比較更重要的,就是當我們明白真理時；需明白神在祂教會中的整個計劃.當我們追求神的旨意時,太容易將範圍縮得窄小了,我們只想到神的旨意,叫我自己做些什麼；或在某一件事上,神叫我做些什麼?免得我所做的脫節.但更重要的,是神在整個教會中的計劃是什麼?在今日這時代,神在整個屬祂的人身上的盼望是什麼?對整個人類的計劃是什麼?整個教會應走的方向又是什麼?整個福音的工作應如何去發展?整個教會應做些什麼?對教會工作主要的方向是什麼?許多時我們容易把那圈子,單放在自己身上,我們應把圈子放大些.神要在我身上啟示祂的主意,不但是限制在我的身邊,要把範圍推廣出去.對整個教會,整個香港的教會,和整個世界的教會,神在我們身上的心意是什麼?我們應走那一個方向,如果我們只想到自己,我,我個人；就永遠不張開眼睛,向遠的方面望了.我們應該知道神在我們身上的個人旨意,更應該注意到神對整個教會的計劃；這是保羅在這話中的重點.按當時歌羅西教會而言,保羅要他們不是單單注意所喜歡的知識,不是所聽到的所謂奧秘的知識,乃是神在整個的計劃,這是保羅在整段經文中,主要的目的.
\newline
\newline
(三)知識追求的目標和方法
\newline
\newline
「好叫你們行事為人對得起主,凡事蒙祂喜悅.」(西一10)
\newline
\newline
這句話很難翻譯「凡事蒙祂喜悅」,在整本新約中僅此處用了一次.雖然有好多處,說及追求神的喜悅,但與這裡的字眼是不同的.「凡事蒙祂喜悅」,有好幾個解經家,試圖把這個字意翻出來,有的說應翻作「在一切的事情上,追求神的旨意,過於喜歡我自己的旨意.」或譯作:「在所有神的事上,都要得祂的喜歡.」凡事蒙祂喜悅,這字的最基本意思,乃是說在一切的事上,你經過思想,經過選擇,經過攷慮,及衡量；那才擇選這樣做,你選擇這樣做的基本原因,是因這事是神所喜歡的.在這許多事上,你都可以行,每件事都可決定去做；但有些事你是知道這事是神喜歡的,有些事你知道是神所不喜歡的.在許多的事情中,你寧可把自己喜歡的旨意放下,來選擇神所喜歡的.這包括了一重要思想:就是你真正要做神所喜的事嗎?那不是你在做些什麼?你所做的也許是神所喜歡的,但你卻勉勉強強地去做,這不是神所喜歡了.雖或是神所喜歡,但你卻為了別的動機而去做,這不是神所喜歡的了.你要做神所喜歡的事,也應出於你自己的心,樂意做神所喜歡的,或者你喜歡做你計劃的事.你說為了一個原因可以把自己的意思放下,為了表示我是愛主的人,我可以把自己的意思放下.為了在人面前表示一個好的見証,可把自己的意思放下.為了勉勵人為激勵人,可以把自己意思放下.那我做乃是神的旨意,但做的時候心裡總是不大開心,為什麼別人不需要做而我要做?這樣的態度,那麼你所做的並不是神所喜歡的.保羅說及凡事要蒙祂喜悅,乃是叫你自己的意志先喜歡神的旨意.不是勉強的喜歡,追求神旨意；是而是你叫的意志,先樂意按著神的旨意行,然後才按著你所喜歡的去行.一個真正合乎神意旨的順服,永遠不應違背自己的喜歡；你會說既是順服,怎麼又會不違背自己的喜歡呢?這豈不是矛盾嗎?又要順服,又要不違背自己的喜歡；我覺得這就是特別的地方的關鍵就在於此,因真正的順服,是先叫內心真正願意按著神的旨意行,先叫我的意志,樂意的順服神的旨意了；當我按著神的旨意行時,神的旨意就變成我心裡最喜歡的了.這樣,我就自然喜歡按著神旨意去行,這是保羅說的「凡事蒙祂喜悅」的意思,如果我們真正得到這種知識,我們知識的根基,就建造在我們順服上.這知識的內容,是神完全的意旨,知識的方法,乃是凡事蒙祂喜悅.這是知識的意義.
\newline
\newline
(四)知識的表現
\newline
\newline
保羅說我們這樣追求知識,實在太難了!我怎能做得到呢!我還未完全脫離老我,雖然我願過百份之百的聖潔生活；但我活在世上時實在做不到!我的老我,舊人仍然存在,我怎能讓自己去百分之百的,喜歡神所喜歡的事呢!下面保羅對我們說:主要加添給我們力量.
\newline
\newline
「照祂榮耀的全能,得以在各樣的力上加力,好叫你們凡事歡歡喜喜的,忍耐寬容.」(西一11)
\newline
\newline
誠然!我們自己做不到,也沒有辦法對付老我.我雖然很願意,卻無法勉強他完全喜歡神所喜歡的事.
\newline
\newline
在這裡保羅說:「神要加力量給我們,要叫我們力上加力.」這是什麼意思?有幾件事應留意的.
\newline
\newline
保羅在此所用的動詞,是斷斷續續地,好像是繼續進行的動詞.這字非常重要,神要把祂的力量加給我們,是繼續不斷加給我們.神有豐富的力量,但不是一次全數給我們,而是繼續的加給.因我們的生活中,有繼續不斷的需要；何時需要,何時就有力量加上.正如血液在我們身體中,不斷地供應我們身體的需要.照樣我們在生活上順服神,做了一步,力量用完了；神就立刻加一新力量來.換句話說,就是在我們所走的每一步的道路上,都從神支取力量；不能說神領我走了第一步,我就靠自己走第二步.力量不能預支的,若今日神把全部力量都給了我,那麼我的生命就會毀滅或爆炸了!你的日子如何,你的力量也如何.按照你日子的需要,工程的需要；神就繼續不斷加力量給我們.這裡保羅所說的:「各樣的力加上力」,實在難翻譯,不過在這裡保羅卻引証出一個真理來,許多時候我們都會忽略了.
\newline
\newline
在此最少包括了兩個意思:(一)就是我們需要繼續不斷的仰望神,亦即等於繼續不斷地不再仰望自己.(二)是需要繼續不斷地,和神保持著應有的關係.燃料的油進入機器裡,若其中有一秒,或百份之一秒的時間油管塞住了；那麼機器就會因沒有燃料,所產的力量通過來而立刻停止了!當我們從神那裡不斷支取力量時,我們與神之間保持關係的油管,若是有一秒鐘停止了；被一些東西阻擋神,攔阻了,神的力量不能進去.我們靠著自己是不行的,神要給我們在各樣力上加力.
\newline
\newline
神力量來源有限制嗎?保羅說:「是照祂榮耀的權能.」(西一11).這意是照祂所顯示的光輝,榮耀；如同在施恩座上,祂所顯示的榮耀一樣.是無限的,與神的本性聯合一起的,也是說神給我們的力量是無限的.我們需要什麼力量,那麼就可有力量給你.保羅說:「各樣的力上加力.」是無限制,今天你需要能做儆醒禱告嗎?在你這軟弱的身體,你需要有喜樂?你要勝過某一個試探嗎?你需要在一個不快樂的環境中能忍耐嗎?你今天在缺之時你需要心能力去幫助人嗎?這一切,神都要給你力量,神有各種的力量.繼而保羅說:「所有的力量.」在所有的事情上,力加上力；不論你所需的力量是那一方面,神都能單單為你所需要愛的力量上,滿足你的需要.神的力量是無限的,但在實際的應用上有一個限制；這限制就是在於我們自己這面.那就是我們有沒有適合支取這力量的條件,有沒有支取這力量時的那種忍耐與倚靠?若沒有這條件,則無論如何總有了攔阻.這就是唯一限制神力量的原因.
\newline
\newline
神為什麼在各樣的事上,給我們力上加力,乃是為了要我們多結果子,或要我們作偉大的工作.不過保羅在此處卻不是論到這兩點,而是說神在各樣的事上,給我們力上加力.為的是「好叫我們凡事歡歡喜喜的忍耐寬容.」(西一11).神加給我們這許多力量,乃是叫我們在困難在痛苦中,能站立得穩.叫我們在需要時能支取力量,有忍耐寬容的心.或許你會有患難,逼迫或痛苦臨到,但這力量能幫助你,擔當得起.這樣就可以顯出我們在主面前,是否有真的知識.那就是說:你是否有保羅所說的真知識?就要看你在患難,逼迫和痛苦時,你那忍耐寬容的心有多少.在困苦試煉中,更能把人的本性顯露出來,在壓力之下,你對這事的反應如何?這樣自然可以把你生命的性格表露出來.
\newline
\newline
(五)知識的目的和果效
\newline
\newline
a 「祂救了我們脫離黑暗的權勢,把我們遷到祂愛子的國裡.」(西一13)
\newline
\newline
「遷到」這字在古典的希臘文中,常用這字來形容希臘國,在某處開拓一個殖民地.希臘國開拓殖民地與現今的情形不同.香港是英國殖民地,其中居民百份之九十八是中國人.希臘開拓殖民地是把希臘人搬到那地方去,而成一希臘城,這城一切生活與希臘一樣,住的是希臘人,說希臘話,用希臘錢；好像是開一個新的希臘國一般.這就是保羅在此所說的「遷」字.用今日的名詞可用「移民」為字代替.當你移民到美國做該國公民時,要在法官面前宣佈,今後是美國公民,同時亦宣佈今後不再是香港公民,脫離原來的國籍.我們乃是移到或遷到祂愛子的國裡.脫離那黑暗的權勢,進入了新的國度.若要進入這新的國度裡,必需與過去的一切一刀兩斷,不能處在兩個國籍中.保羅這句話,給我們一個很重要的責任?我們必需脫離黑暗的權勢,才能遷入愛子的國裡.
\newline
\newline
b 「我們在愛子裡得蒙救贖,罪過得以赦免.」(西一14)
\newline
\newline
主說天上地下所有的權柄都交在我手裡.祂不但是創造萬物的主,也是今日掌管萬物的主.一切都是祂的,我們在祂那裡得基業,得享受一切的好處?這是主在我們身上的旨意,和目的,我們要追求這真正的知識,非人思想裡所謂奧妙的知識.這種知識也許我們得不到,沒力量可行出來；但主會加添我們的力量,祂加力量最終的目的,乃是要我們脫離黑暗的權勢,遷入愛子的國裡；在祂裡面可以承受基業.這是非常崇高的思想,但這是個事實.
\newline
\newline


\newpage

\section{第46屆~港九培靈會~鮑會園~第三講}
\label{sec:1186}
\textbf{無比的基督}
\newline
\newline
連結: \href{https://www.hkbibleconference.org/session-message/view/1186}{\texttt{https://www.hkbibleconference.org/session-message/view/1186}}
\newline
\newline
\hyperref[sec:1185]{< < < PREV SERMON < < <}
~
\hyperref[sec:toc]{[返目錄]}
~
\hyperref[sec:1187]{> > > NEXT SERMON > > >}
\newline
\newline
歌羅西書 1:14-1:14
\newline
\begin{longtable}{cl}
\hline
\hline
章節 & 經文 (和合本修訂版)\\
\hline
1:14 & \begin{tabularx}{0.7\textwidth}{X} 藉著他的愛子，我們得蒙救贖，罪得赦免。 \end{tabularx} \\ \\
[1ex]
\hline
\hline
\end{longtable}
(西一14)可當作是保羅的禱告,不過由十五節開始,就是保羅要向歌羅西信徒,解釋他所要講的真理.在前面我們已見過保羅,曾向那些誤解真理的人,解釋和教訓,當他要改正他們的錯誤時,他就先帶領他們禱告.從今天的經文,可看見保羅又從另一角度來提醒他們.要提醒或指正人,在真理上的錯誤時,可用兩種方法:一種是直接指正他的錯誤,這是比較容易.
\newline
\newline
但這樣的指正,卻不容易得到預期的果效.雖然批評和指正是容易做到的事,不過也是極端危險的事.因為極度的批評可能把他基本信心的根基也搖動了.如果一個人對他基本的信心發生動搖,那麼他可能什麼都不相信了,所以我說直接的批評,有時會很危險,但卻比較容易.但這樣的指正,卻不容易得到預期的果效.雖然批評和指正是容易做到的事,不過也是極端危險的事.因為極度的批評可能把他基本信心的根基也搖動了.如果一個人對他基本的信心發生動搖,那麼他可能甚麼都不相信了,所以我說直接的批評,有時會很危險,但卻比較容易.常見許多人批評教會的錯誤,知識落伍,辦事方法也落伍；工作又是沒有果效,發展又不夠快,好像教會完全沒有價值了!這樣的批評是很容易的,但批評的果效可能非常嚴重.
\newline
\newline
另一種方法,就是不直接批評他所相信的錯誤,乃是把自己所認為對的,並所知道的真理,清楚擺在他前面；讓他看見真理多麼美好,然後在他心中可以作一比較.讓他知道這真理,實在比他自己所相信的異端好得多.這樣,他自然就會抓緊真理,而放棄他那錯誤的信仰.但這樣做是非常的不容易,消極的批評,永遠比積極的建議容易得多.保羅就是用這建議的辦法,他在向這些可能犯了錯誤的歌羅西人講話,沒有直接批評他們那一方面錯,只先把基督的美好擺在信徒面前；使他們看這位基督是多麼的美,多麼的完全.若是他們真知道基督的美,自然就會離開他們那錯誤的思想,而跟隨基督.把基督真正的榮美彰顯出來,積極地把好的擺在人前,這並非容易的工作；保羅作這工作何等地完全.下面是他說及基督的幾方面:
\newline
\newline
(一)基督是誰
\newline
\newline
「愛子是那不能看見之神的像,是首生的,是在一切被造的以先.」(西一15)
\newline
\newline
保羅在此解釋基督本性:
\newline
\newline
Ａ　那不能看見之神的像──這「像」與照片不一樣,照片只是外形的代表,但在此不單是外表的形像,更是基督的本性；也是神本性的代表,且是神本性的表現.在基督裡,就把神的本性完完全全的表現出來了?主耶穌在約翰十四章九節說,人看見了我,也就是看見了父.父的一切都在祂裡面.祂不單只是父的代表,更是父的顯現；那不是圖畫,也不是代表,乃是真正有形有體的表達出來.這是「像」的意義,是保羅所要表達的.主耶穌要把天父完完全全表達在人前,使人看見了祂,就好像看見了父一樣.主所代表,所有,所在生活出來的；一切都與父一樣.在第一講中我曾提及異端之錯誤思想,他們是說在神與人間有許多不同等級之被造之物；連基督也在內,他們說祂是最末後的,在本性上已離神很遠,但保羅卻否定此說,神的形像就在基督裡面.祂是完全與神一樣,完全將神表達出來.在下面的經文,保羅所用的詞句,就是把這意思顯明出來.
\newline
\newline
「因為父喜歡叫一切的豐盛,在他裡面居首位.」(西一19)父的一切是豐盛,神性,榮耀,智慧,能力,都完全地居住在主裡面.而且這「居住」之意,並非是暫時的兩三天,而是永遠地在家中居住的意思；是永久之居所,神之豐盛是永遠地,時常地,居住在主耶穌裡面,有些異端說:耶穌只不過是人,不過是個最完全的人,最聖潔的人.但他是人,又是基督,那該如何解釋?他們說基督是神的一個代表.當主耶穌工作時,基督就來到祂身上,像神附在人身上一般；但當主離世時,則基督這部份的神性,就離開了而歸到天上去.不過保羅卻不是這樣說,他說神一切的豐盛是永遠居住在基督裡,祂是神的形像.
\newline
\newline
Ｂ　基督是首生的──這節經文,許多人感到很難解釋.今日許多很純正信仰的教會,因此經節而成為異端.若祂是首生的,那我們是次生的,或第三,四生的,若根據此經文,則基督豈不是是我們的兄長了?若如此,則基督不就是和我們一樣子嗎?我們與神是不一樣的；但若是基督與我們一樣,則基督豈不就與神有分別了嗎?若是基督是首生,則祂與我們同一地位了.許多教會都為此而紛爭!「首生」這字實在難解釋,這字本身之意義,或許與我們的解法不同.這字在新約中只用過五次,在此乃是第一次,其次在西一18),「是從死裡首先復生的.」又(羅八29)「使他兒子在許多弟兄中作長子.」前面說及應效法祂兒子模樣,之後就說及在基督裡得榮耀.在(來一6)與(西一18)之用法是不同的.「神使長子到世上來的時候.」(來一6)(又在啟一5)「從死裡首先復活.」以上五處經文,若能詳細分析,就可發現每處都提及基督有特殊之處.他是從死裡首先復生的,但這並非說祂是在肉體上第一個復活.拉撒路比祂先復活,舊約中亦有許多人從死裡活過來.在此處是用這字,乃是祂從死裡復活的方式與其他的不同.祂要在眾弟兄中作長子,這是說祂不單是第一個出生就算了.試問舊約時代的人是否弟兄呢?耶穌出生時是比摩西和亞伯拉罕更早嗎?若按此時間,則祂不是長子了.故「首生」的意思乃是祂生的方式,與我們所說的生的方式不一樣；實際上保羅是用此字,來表達基督特別的地方；這不同的地方就按聖經的用法,這「首生」應包括幾個意義:
\newline
\newline
(1)基督乃創造萬物之主──我們不論是天上的,地上的,是使者,是人,都是神所造的；但基督就不是神所造的,是首生的.保羅在此就說出其中的分別:祂是生的,而我們是被造的.生與造,有極大之分別,洋娃娃,桌子......你可造出來；但卻不能造出一個兒子,用多少化學原料也造不出來.兒子是要生出來的,造的與生的基本上有所分別.所以基督與一切被造之物,絕對基本上有不同的地方.
\newline
\newline
(2)與神之本性合而為一──祂的本性與神全無分別,我們後裔,與我們在神面前完全一樣.我們的本性雖是罪人,同樣可與神發生關係；但在我們的後裔中,所有之事物就不是一樣.你可說猩猩聰明,但站在神面前,他與我們有基本上有分別.我們和我們的後裔,有同樣之本性,照樣基督是神的首生,祂的本性完全一樣.
\newline
\newline
(3)基督有承受神一切的權柄──記得以掃是因紅豆湯,而出賣長子的名份,賣了之後卻不甘心,又想將長子名份奪回.那何必如此麻煩!大哥,二哥,三哥又有何分別；在家中排行第幾有何關係呢?這是我們看法,但在神的看法並非如此,為何雅各與以掃會有此紛爭?他們都是伯亞拉罕的孫子.其實有很重要的分別,因長子有承受遺產的權利,主乃是神首生的.主在馬太廿八章說:天上地下一切所有的權柄,父已交給我了.而且在未復活之前,在約翰十三章也說:父已將萬有交在我手裡,祂是首生的,有承受神一切的權利,且有掌管一切的權利.這是權柄及地位的問題.
\newline
\newline
C　在一切被造以先──祂是在一切被造之物以前,有兩方面之意思:
\newline
\newline
(1)基督不是被造之物──在一切東西未造之前,基督已經存在.不論任何一切,都是在基督之後,那就是說基督不是被造之物.保羅在此繼續說基督是創造的主,那麼當然祂就不是被造之列.
\newline
\newline
(2)基督有永存的性格──在(箴八23-26)說:「從亙古,從太初,未有世界以前,我已被立；沒有深淵,沒有大水的泉源,我已生出.大山未曾奠定,小山未有之先,我已經生出.耶和華沒有創造大地,和田野,並世上的土質,我已生出.」箴言書說明基督與智慧的本性,從本段經文可見基督永存之性格.智慧在此解釋:在一切被造之先,智慧之先,甚至一切神的工作之先；這一切之先,乃是代表祂從永遠直到永遠存在的.這就是其中之真理,基督是永存的,沒有開始也沒有結尾.經上說:我是始,我是終.(啟一17)
\newline
\newline
(二)基督的工作
\newline
\newline
「因為萬有都是靠祂造的,無論是天上的,地上的,執政的,掌權的,一概都是藉著祂造的,又是為祂造的.」(西一16)
\newline
\newline
Ａ　基督是創造萬物的主──無論世上的一切物質,我們所見的東西,都是基督所造的.我們自己,我們的思想,理智,情感,我們這按著神的形像所造出來的人,也就是基督所造的.我們想及在我們以外的屬靈物質,雖然未見到那黃金街,碧玉城,這是基督所造的.在我們以外另有的屬靈被造之物:有天使,魔鬼,邪靈,這一切都是基督所造的.不論你想到在神以外的一切東西,都是基督所造的.這可見保羅在此說的,基督教的世界觀.因為這「萬有」一字是代表整個的宇宙,不論屬物質的,屬靈的,可見的,或看不見的,這一切都包括在內.本段經文中,保羅重複的用「萬有」這字,大概就是此原因.
\newline
\newline
既然一切都是基督所造,洛斯斯底派人所說的,神與人之間的單元,又如何解釋?一切都是基督所造,則當然基督就不是其中單元的一個,而是創造萬有之主宰.
\newline
\newline
Ｂ　基督是掌管萬有的──「萬有也靠祂而立.」(西一17)
\newline
\newline
萬有靠祂存在,並靠祂掌權.這是很重要的,有一東西掌管一切.各位可知原子彈是什麼嗎?乃是把最基本的物質單位就是原子,將之破裂,有能力發出來.今世上有這麼多的原子,若這原子都破裂了,那麼將如何?不會的,因為有基督在掌管一切,整個宇宙都在基督掌管之下.整個世界也在祂掌管之下,世上許多事發生我們都不明白,為何世上有這許多不公平的事?有些國家地土非常肥沃,但也有些國家非常的貧乏枯乾；有些民族生來就有優越的享受,但亦有一生在受著痛苦,是命運嗎?是其他的因素嗎?我們也許不會明白,但我們卻知道一件事情,基督在掌管一切.祂有原因,也有理由,不但世界是基督所掌管,人類的歷史也是在祂掌管之下.試看世界民族之演變,國家之興衰,又如我們中華民族的民運,我們不易明白,但我們知道整個歷史,在主的手裡掌管著.若不是主許可,沒有一件事可以發生.不單掌管人類歷史,亦掌管今天我們的生活環境,各人的生命,教會及其他一切；所以一切都不必掛慮及灰心,主是掌管一切的.
\newline
\newline
既然主是掌管一切的,然則我們有讓主來掌管我們嗎?這是一個問題,萬有都靠祂而立,那麼我呢?你呢?你的生活習慣,你做事的方法,主要掌管一切.
\newline
\newline
(三)基督是救贖萬有的主
\newline
\newline
「既然藉著祂十字架上所流的血,成就了和平,便藉著祂叫萬有,無論是地上的,天上的,都與自己和好了.」(西一20)
\newline
\newline
祂藉著十字架上所流的寶血,救贖了一切.我們知道保羅這整個的真理,是針對著歌羅西所有的異端而說的,這裡保羅簡單的提出了幾件事:
\newline
\newline
Ａ　藉主耶穌的寶血而得和好──因主為我們在十架流出寶血,而成就了和平.今日我們能達到神面前,並非因著洛斯底派人所說的奧秘知識,非靠著你知道得多少；或有何深奧思想,或有何奇妙特別之經驗,而是因主之寶血.
\newline
\newline
Ｂ　基督是道成肉身的──洛斯底派的人說,基督不是真正道成肉身.全能的神怎會成為我們這有罪的身體呢?如果祂沒有成為我們這有罪的人的身體,祂就不能為我們流血,為我們流血才能成就和平.故此,主耶穌為我們道成肉身,乃是絕對的事實.這方面我們需要特別注意:我們愛主,尊敬主,相信主,故每事都高舉主.主耶穌是神,我們都敬愛我們的主,我們的神.但有時我們太過高舉主時,把主的神性舉得非常高,而把主人性失去了,以為主是神就毫無軟弱.試問主耶穌為何會在船上睡覺呢?為何會口渴?疲乏,發怒呢?為何心裡會憂愁呢?祂的確是神,但祂的人性,卻也是百分之百的真實.
\newline
\newline
根據此經文,有人會有一另外的思想,如果基督在十字架上流了血,成就和平,萬有都與祂和好了,那豈非全世界的人都可以得救了嗎?故此,就發生了異端叫普救論,他們說全世界的人遲早都會得救的；因得救不是靠自己,而是靠主的寶血.主既叫萬有與祂和好,還有什麼不能得救呢?保羅恐怕有人有此種思想,故在此說明,基督是與我們和好；但要有一條件就是:「只要你們在所信的道上恆心,根基穩固,堅定不移.」(西一23).我們有信心,因著神的話,相信了這位為我們釘在十字架上流血的主,是極大的安慰,這也是一個保證.保羅所講的不是普救論.我們的確不錯的與神和好了,但能享受與祂和好的福份,就需要靠著我們對主的信心.這位如此奇妙的主,我們應該如何地尊敬祂,順服祂,愛祂,事奉祂!
\newline
\newline


\newpage

\section{第46屆~港九培靈會~鮑會園~第四講}
\label{sec:1187}
\textbf{蒙福的職份}
\newline
\newline
連結: \href{https://www.hkbibleconference.org/session-message/view/1187}{\texttt{https://www.hkbibleconference.org/session-message/view/1187}}
\newline
\newline
\hyperref[sec:1186]{< < < PREV SERMON < < <}
~
\hyperref[sec:toc]{[返目錄]}
~
\hyperref[sec:1188]{> > > NEXT SERMON > > >}
\newline
\newline
歌羅西書 1:23-1:23
\newline
\begin{longtable}{cl}
\hline
\hline
章節 & 經文 (和合本修訂版)\\
\hline
1:23 & \begin{tabularx}{0.7\textwidth}{X} 只要你們持守信仰，根基穩固，堅定不移，不致動搖，離開了你們從前所聽見的福音的盼望；這福音也是傳給天下一切被造之物的，我—保羅作了這福音的僕役。 \end{tabularx} \\ \\
[1ex]
\hline
\hline
\end{longtable}
(西一23),保羅說他也作了福音的執事,從廿四節開始,保羅就解釋他作了執事的工作.整段經文是解釋,他從神所領受的的托負和職份是什麼.在廿五節說:「我照神為你們所賜我的職份,作了教會的執事.」這是說我所做的工作,是從神所領的職份,「職份」這字是相當難解釋的,解經家有各不同的看法,甚至因此字而有完全不同的神學系統產生.
\newline
\newline
當然此字有時是指神為某時代,所有的特別工作方法,但在中文的翻譯十分正確.保羅在此並非說在神某一不同的職份所領受的是什麼.乃是說他領受的工作,並非出於自己的喜歡,或選擇；而是出於神按著祂那完全的計劃,選派我作此工.如今我所作的,不是出於我自己的意思,因此我不能按我自己的意思去作.保羅這種思想,就直接地說出,我們每一個基督徒應有的生活方式.若我們相信前面所講的,神今天不單是創造了我,且掌管了我們.神在我們每人的身上,都有一個完全的計劃.祂在我們身上有一目的,祂給我們有一責任；無論你今日的工作範圍是什麼,但在神面前都有一本份.這責任並非是我喜歡做什麼,或什麼事的是重要,而是今日神要我作些什麼.神給我的職份是什麼?若神給我的職份是傳道人,若我去做生意就錯了.照樣,若神今日給我的職份是做生意,我卻按自己的喜歡去做傳道人,那也是錯了.神既在我們身上有一清楚的計劃,那麼我們各人,就當分別負神給我們的職份.若用此種眼光來看基督徒的生活,那麼每一個基督徒,無論是作什麼,只要你按著神在你身上的計劃而行,你就是做神的工作了.今日許多人有一種觀念,認為某人纔是服事神的人,這意思是意味著說,有些人不是服事神的了.什麼人才是服事神的?讀神學,傳道,長老,執事......若這些人是服事神的人,其他的人又如何呢?他們不是服事神,那麼是服事誰呢?不是服事撒旦,或自己吧?我們必須認清楚,神在我們身上的計劃,如果我們真正按神的旨意而生活,這樣就是過著服事神的生活了.因此我們必須追求,在生活上如何來服事神!換言之,我們應當怎樣過基督徒的生活.今天太多的基督徒,把他們的生活分得太清楚了,只以為禮拜天的生活,才是真正有意義的生活.這態度只能在一觀點上看是對的,許多基督徒被這種態度所害了!有一個知名的傳道人(以前他是造鞋的),有人問他的職業是什麼?他說我的職業是為主作見證.我只不過是做鞋來維持生活,這種態度,也只能在一個觀點上看是很對的,但許多時因這種態度而使我們對整個生活,給人有錯誤的觀念.若我們認清楚了,今日我當教員,乃是神給我的呼召,公務員,政治家,文學家,經濟家,傳道人,這一切我們都認清楚是神給我們每人的呼召.並清楚了神在我身上一切的計劃,然則我們在工作時,就會有不同的觀念.我們會盡最大的努力,把神給我的職份做得最好.若每個基督徒對本身的職份,都如此看法,就應把作基督徒的態度和觀念,帶到整個生活上.如此,則基督徒生活的影響,就進入整個社會裡.試問你能否想像到,這是實際的情形呢!假如我是個政治家,作政治家不過是維持生活的方法；這工作與我做基督徒生活,完全是無關係,然則我作政治家就隨便馬虎地去應付,這樣的態度對嗎?保羅卻不是如此,他說我們在神前領受的職份,這職份不單是指傳道人領受,而且也是每一個基督徒所領受的.神在你上既有一個計劃,你的職份若不是在神的計劃中,那麼就該尋找神給你在祂計劃中的職份.神給每人都有職份,我們應該用什麼態度以盡我們職份呢?保羅告訴我們:
\newline
\newline
「我也為此勞苦,照著祂在我裡面運用的大能,盡心竭力.」(西一29)
\newline
\newline
保羅說我知我的職份是什麼?我就勞苦努力去作此工.「勞苦」是表示把你整個的精神,整個的力量放在這工作上,甚至流汗流血去作.換句話說就是毫無保留地,盡本身的最大努力任此職份.今天如果我在神前有所領受,就必須在生活上有好的見證.如果我在屬靈事上有責任,就應該把全副精神去完成這責任.基督徒在屬靈的生活上馬虎乃是常見之事,把生活化成幾部份；一部份是工作,一部份是基督徒.一週內禮拜三查經,禮拜六團契,主日崇拜時就作基督徒；其他的時間,就作另外的工作.這工作與基督徒沒有直接的關係,因此,把整個精神力量放在這工作的一面.這樣就把基督徒的生活,變成可有可無的裝飾品了!鮮花是可愛的,若把鮮花變成裝飾品,那麼就成為可有可無了!保羅覺得這種態度是錯誤的.為了完成這職份,就為此而受苦；把整個精神力量都集中起來,因此基督徒生活就成了必需的.我們在屬靈的事上看得很隨便,可能是生活的目標沒有認清楚；你也許知道這目標,但卻因其他的事,把他眼睛遮住了,因而失去了目標.你知道神在你身上有一個計劃,是非常美也是非常好的而你卻被許多東西遮蔽了視線,被週圍的一切吸引住了；所以你看不見神給你最終計劃!卻自滿足了.我們應該盡心竭力去追求,不要因很著雲彩的景象就滿足,更不要因看見現今的生活就滿足.應該認清楚真正的生活目標是什麼?要達到此目標,必需付代價,如果你真正認清這目標的價值,而肯付出代價去得它有何不值呢?
\newline
\newline
保羅又說追求屬靈的事,不是靠自己的眼光,知識和努力；而是要照著神在我裡面運用的大能.神應許給我們能力,這運用的大能,是指聖靈在我們裡面工作的力量.歌羅西書與以弗所書,很多方面是相像的,但在此卻完全不一樣.以弗所書聖靈的工作,講述得非常清楚,但歌羅西書,就沒有一次直接講及聖靈.他並非沒有聖靈的工作,而是間接的,從暗中的不同方向,把聖靈的工作表顯出來.他說在聖靈的追求上.所倚靠的,乃是祂在我們裡面運行的大能,要靠這能力；認清生活目標,就能作成神叫我們作的職份.在作成這職份時,應在生活的那方面有所表現呢?保羅說,領受職份包括三方面.
\newline
\newline
(一)補滿基督患難的缺欠職份
\newline
\newline
「現在我為你們受苦,倒覺歡樂,並且為基督的身體,就是為了教會,要在我肉身上補滿基督患難的缺欠.」(西一24)
\newline
\newline
這補滿缺欠,是補滿基督所受的苦,那是什麼的苦?當然不是十字架上完成教贖的苦；因這不是我們所有份的.而且這痛苦是百分之百的完全,毫無缺欠,不必補滿.現今所說的乃是代別人所受的痛苦.基督為我們受痛苦,保羅說我們要受的苦,乃是代替神的教會受的.怎樣才是為別人受苦,而補滿基督的痛苦呢?按保羅的經驗而言,當他往大馬色路上遇見了主時,主說對他說:「掃羅!掃羅!你為甚麼逼迫我?」其實保羅並沒有逼迫主,只不過逼迫基督徒而已.但基督徒乃是主的肢體,整個教會是主的身體；基督乃教會的頭,故基督的身體任何一部份受苦；也等於基督本身受苦.保羅逼迫基督徒,所以主說你為什麼逼迫我!我們是為了教會而受苦,藉著受苦,就完成了兩方面之工作.一方面藉著受苦,使我們生命長大.使徒曾說,藉著受苦可使我們更完全.另一方面,受苦可使我們脫離罪；痛苦是有此作用,能使我們生命漸漸長大.正如樹在山頂上,經過風吹雨打太陽晒；你會說這幼小的樹苗經大風吹它,對它豈不是極大的摧殘嗎?不,經過風雨後,你就會發現樹苗長大了.正因有這樣的摧殘,才使樹的根基越穩,也長得更茂盛.樹的生命經過了這摧殘,就有了抵擋摧殘的力量而生存.當然；若此樹還沒有生命的話,則無論如何的堅固；就算是一棵鐵樹,經過風吹雨打,也許第一天還很穩固,可是經過一段時間,這鐵也會生鏽而無用了.但如有生命的樹,經過摧殘,裡面有了抵擋的力量而生存著長大著.基督徒亦如此,為主緣故受了痛苦,靠主的力量而勝過；我們裡面就生出力量,使我們能站立得穩.
\newline
\newline
受苦的另一原因,我們對於今日的世界,非常的實際.若實際地看這社會,不作夢想,例如說我作基督徒得救了,奉獻作傳道人,為主犧牲了.那麼我這為主奉獻,又犧牲,又聖潔,樣樣都比人好!我傳的是純正的道理,是救人的福音；世界會歡迎我嗎?若你是存了此態度去工作,定會非常失望!雖然你態度很好,目的好,所傳的也好,但世人卻不接納你,反而毀謗你逼迫你!你為了主緣故,甚至會要犧牲你的性命.你該認清楚,世人對基督徒對屬靈的事,就是如此,他們認為屬靈的事是愚拙的,無價值的,他們會反對基督徒,逼迫為主工作的人.如果我們能認清楚了,存了這樣的態度,那麼就算遇見逼迫,反對,也不至於灰心失望.這就是羅馬書一章保羅所說的原因:「我不以福音為恥.」保羅好像的太消極了吧!為什麼他說不以福音為恥呢?應該說以他為榮耀才對,這原因就是保羅對這世界,這社會及人的心,認識得十分清楚.若我們有此實際的認識,就算有困難也不會灰心失望.主說:世人怎樣逼迫了我棄絕了我；照樣也要逼迫你們,棄絕你們!僕人不能大過主人,所以我們若為了主的的名而受逼迫,棄絕,也不至於灰心.相反的,若不信主的人,對我們基督徒的態度及見證,表示完全歡迎,沒有一種惹起別人反感的生活；每人都覺得基督徒的生活百分百是好的,那就該要小心,恐怕我們的見證有了問題.保羅在此提醒我們,應當受苦,不過並非說非受苦不可；或是越受苦越好,這觀念是錯誤的.乃是說,如果為了完成我們的職份,而需受苦,或要完成我的見證,而需要受苦；或要完成我的見證,而需要受苦；那就該因此而歡樂!我感謝主,祂給我職份,教我配藉著我的生活,而為主受苦.
\newline
\newline
(二)把奧秘顯明出來的職份
\newline
\newline
「這道理就是歷代所隱藏的奧秘,但如今向祂的聖徒顯明了.」(西一26)
\newline
\newline
當時歌羅西的一般異端說,有一奧秘這是真正的奧秘；這奧秘是從歷世歷代一直隱藏著的,可是今天不再隱藏了.如今已向眾聖徒顯明了,這是奧秘的真正意義.奧秘是過去神計劃裡,一個隱藏的真理,人無法明白,用盡方法也不能知道；但神的時候到了,祂把這真理啟示出來.
\newline
\newline
這奧秘是什麼呢?在(西一27)告訴我們,叫我們知道這奧秘,在外邦人中有何等的榮耀.基督在你們這歌羅西人心裡,在你們這不是亞伯拉罕後裔的外邦人底心裡；在你們這些過去,在神恩典並沒有份的人心裡,基督在你們裡面成了榮耀的盼望.我們傳揚祂,是傳揚這奧秘的意義；我們領受了這奧秘,這奧秘沒有辦法知道人,人不能用自己的方法去明白.但神啟示我們明白了,今日我們要藉著我們的生活,叫世人知道神的奧秘,並彰顯出來.這是我們的責任,藉我們的生活和工作,叫人知道外邦人,可以在基督耶穌裡蒙恩.基督成為我們榮耀的盼望.這奧秘在過去不能啟示出來,其原因乃是:
\newline
\newline
(1)在過去主耶穌還沒有來──這奧秘乃是基督在我們裡面成為榮耀之盼望.所以在基督未來之前,這奧秘不能啟示出來.
\newline
\newline
(2)在過去的人,靈裡未長大成熟──未長大成熟,所以神不能告訴他們.就算告訴了他們,他們也不懂；神若告訴舊約時代的以色列人說,我要叫香港的人蒙恩,這是不可能；他們不明白這事,神必需按人所能明白的,程度,將真理啟示.我們的屬靈程度長進一些,神就給更多一點真理我們；當我們屬靈的程度成熟了.神的時候到了.祂就把這奧秘啟示給我們.
\newline
\newline
(三)禱告的職份
\newline
\newline
「我願意你們曉得我為你們和老底嘉人,並一切沒有親自見面的人,是何等的盡心和竭力.」(西二1)
\newline
\newline
「盡心竭力」這四個字是特別的字,極難翻譯且是當時的圈子裡,成為一專門的學問.是指比武時,在場上互相較力掙扎之意,甚至人與獸鬥時拚命,在掙扎時所用的字.保羅在此說,我為了你們的緣故拚著命在掙扎,就如在戰場上遇了強大的敵人；拚命爭鬥以保存自己的生命.保羅當時是在羅馬坐監之時,可以顯明知道,他在肉體上不能與人較力；既如此,那麼就是靈裡的較力了,他聽到歌羅西人遇到困難,和各種試探和引誘,亦聽到了許多人已被引誘走上異端；他就如同遇見了敵人,面對面地與他較力一般.為你們的緣故,我在靈裡拚命地掙扎戰鬥,這是禱告的生活.他知道歌羅西人遇見了困難,走上了錯誤的道路,就拚命為他們禱告；這對我們是何等大的激勵.聽見別人走上異端,或是有何困難,我會如同遇見敵人,像較力般拼了命的去禱告,爭戰.保羅未到過歌羅西,在個人方面而言,他不認識那兒的基督徒,但他知道歌羅西的基督徒是主的肢體,知道他們如何如此情況,雖然在個人方面無得失,但卻是遇了敵人般在交戰,拚命的為他們禱告,這教我們何等的激勵!今日有些基督徒聽到別人有困難時,心裡毫無感動,靈裡也毫無負擔；聽見某教會有困難,或者有紛爭,或分裂,打官司告狀,或誰犯了罪,就心裡很興奮；向這邊說說,那邊談談,心裡毫無難過和為他禱告的心!這是否就是成為我們教會軟弱的原因,保羅不是如此,他不認識歌羅西的基督徒,沒有個人友誼的存在,卻因他們是基督徒,是主的肢體,保羅就立刻為他們禱告.好像是對敵人說,你死我活的一樣,如果我失敗,就等於那些基督徒失敗了,故拚命為他們禱告.我為這事情得勝就如同我自己的生命得保存一樣!為他們盡心竭力.而今日的世代,是多麼需要有保羅這樣的心,這實際上就是主所給基督徒的職份.我們有職份為主所受苦,有職份為這奧秘傳揚出去,亦有職份為眾信徒留下良好的祈禱的態度.
\newline
\newline


\newpage

\section{第46屆~港九培靈會~鮑會園~第五講}
\label{sec:1188}
\textbf{有意義的生活}
\newline
\newline
連結: \href{https://www.hkbibleconference.org/session-message/view/1188}{\texttt{https://www.hkbibleconference.org/session-message/view/1188}}
\newline
\newline
\hyperref[sec:1187]{< < < PREV SERMON < < <}
~
\hyperref[sec:toc]{[返目錄]}
~
\hyperref[sec:1189]{> > > NEXT SERMON > > >}
\newline
\newline
整段經文中包括一句話,這句話只有一個動詞,顯示出這是個命令.保羅吩咐歌羅西的基督徒說:「當遵祂而行.」這顯示出基督教的教訓,與當時異端的教訓有分別.異端是說要追求知識,追求奧秘,和超然的知識.保羅說你們應當有知識,然而單有知識是沒有幫助的?真正的目的乃是要遵祂而行.真正得到知識,是要行在生活上表現出來,所以今天的信息,可稱為有意義的生活.
\newline
\newline
單是明白許多深奧的真理,如果不能在生活上表現出來；所知的真理與我們就沒有什麼幫助.我們在這次培靈會得到了許多教訓,有人每年都參加培靈會,一天三堂很忠心地聽.為這樣的弟兄們,我們很感謝主,他們有一追求的心.但我仍覺得,單是有這聽的心是不夠的,假若我們每年都聽到許多的真理和教訓,而未能在我們的生活上發生作用；那麼我們所聽到的,對我們就毫無幫助了.正如雅各所說,聽道而不能行道一般.現今我們看看保羅如何吩咐我們.
\newline
\newline
「你既然接受了主基督耶穌,就當遵祂而行.」(西二6)
\newline
\newline
這節經文按中文的翻譯,未能將原意完全表達,也未能將其中的兩個連接詞譯出.這兩個連接詞的意思是非常重要的,那意思是:「你們如何接受了主基督耶穌,就當照樣遵祂而行.」然後才在第七節說及在主裡面生根建造,信心堅固.這樣,六節七節的句語就互相對照.把保羅所要說的真理表達出來了.如今我們思想一下:
\newline
\newline
(一)「你們如何接受了主基督耶穌.」(西二7)
\newline
\newline
首先是我們如何接受了主基督耶穌為主,這是表達每一個基督徒生命的根基,如果我們想過一個有意義的生活,就該先從這方面開始.必需從根本上,接受了主耶穌基督.保羅在這裡用的名詞非常完全,許多人曾經提出:在書信裡保羅每用及主耶穌的名稱時,都有他清楚的目的.例如在(加六17)說:「從今以後,人都不要攪擾我,因為我身上帶著耶穌的印記.」這裡沒用別的名稱而單用「耶穌」.這是代表耶穌的人性,耶穌在人間的生活.保羅說我為了耶穌的緣故,我身上帶著了特別的記號.但在別的地方,林後五章裡保羅說我們應如何地去追求,並如何的生活.這基本的動力,是有基督的愛激勵我們.在此單用「基督」而不用「耶穌」,因為基督的愛,是因祂在十字架上為我們受苦.「基督」乃是代表主耶穌為我們所做的工作.所以這節經文說:「我們接受了主基督耶穌.」那是說我們要接受這完全的主耶穌基督,亦需要接受基督,也要接受耶穌,接受這一位主.我們的主是耶穌,也是基督；我們不能忽略了,祂的性格上的任何一方面.當然一方面,是針對當時的教會中的異端而言,他們對基督有許多不同的看法,有人認為基督不過是神的一種顯現的方法,或認為基督不過是人；只是神的能力特別與祂同在而已.有人對耶穌的看法,不十分完全,所以保羅要針對他們.
\newline
\newline
保羅說,我們所接受的是主耶穌基督,但在另一方面,他所用的名詞,都有他本身的意義.首先「耶穌」這名詞是代表祂的人性,成為肉身的生活；代表祂在地上所為我們成就的.但主耶穌來到這世界,乃是為要把祂的百姓,從罪惡中拯救出來.我們應當記得,我們的主是耶穌；也是說祂乃是一個真正的人.如果我們忽略了,就會把主看得太虛無了.正如異端所說祂只不過是個顯現,不過我們知道主耶穌是童貞女所生的.
\newline
\newline
另一方面,有時因太尊敬主,認為祂非常之聖潔；無所不知,無所不能.尊敬祂各方面的好處,而結果就把主變成了一個超然的地位,差不多將祂的人性完全遺忘了.聖經上很清楚地有說及主的人性,祂會疲乏,飢餓,喜樂,忿怒,軟弱,掛慮,甚至有時有些事情祂不知道；這一切,都顯出主耶穌真正的是一個人.難道主耶穌是一個人,這件事是非常的重要嗎?是的,就是因為祂是人,祂才能同情體恤到我們的軟弱.因祂也曾凡事受過試探與我們一樣只是祂沒有犯罪.(來四15-16).這就可表明試探在祂身上有力量,祂可能被試探所勝；但因著祂對神的心,對神的態度,祂沒有被試探所勝.因此,祂能夠幫助我們在軟弱中的人.
\newline
\newline
又另一方面祂不單是耶穌,也是基督.那是代表祂真正的道成肉身,成了我們的彌賽亞；基督是代表祂的工作和職位.本來祂是神的兒子,但在永世的時候,祂不是彌賽亞.因人之需要,在神與主耶穌之同意下；主耶穌甘心順服神的旨意,謙卑自己而來到世上.為了要成就做彌賽亞的工作,所以做基督的第一件事,就是對神要完全的順服.因祂的順服,那才能來到世上成就救贖.祂說:「神啊!我來了!我來乃是按著你的旨意行.」祂存心謙卑以至於死.祂不單順服,且有真正謙卑,祂因有謙卑才能從神的地位變成人的地位,而道成肉身.不但是有順服與謙卑,更是在祂作彌賽亞這工作上,顯出了祂對人的大愛.因祂對我們的愛,甘心樂意地在十字架上為我們受苦,在十字架上擔當了我們罪的刑罰.所以能把我們從罪惡裡拯救出來,今天我們若接受了這位基督耶穌,就需要完全的接受祂.接受祂作我們的救主,接受祂在十字架上為我們成就的救贖,也接受祂是一位真正的一個完全的人,也是完全的神.要這樣地接受這位完全的主耶穌基督.
\newline
\newline
保羅在此又說:我們不單是接受主耶穌基督,更要用一特別的方式和態度去接受祂,這是很重要的.我們接受主耶穌基督,到我們的生命裡做什麼呢?乃是要接受祂作我們的主,按照聖經上「主」這字的用法,有幾種不同的意義:
\newline
\newline
先生或老師──這是普通的意義.祂收門徒,就可知祂有作先生的資格,祂又將真理啟示給我們知道,那祂也是我們的老師.
\newline
\newline
彌賽亞──聖經上有幾處是如此用法.祂是我們的主,我們的彌賽亞,祂來是為我們成全救贖.
\newline
\newline
權威與地位──這是保羅在這段經文的用法:「主」是代表祂的權威和地位,代表主與僕人之間的關係.
\newline
\newline
今天若我們接受了主,那就不單是接受祂作我們的老師,彌賽亞,救主.更是要接受祂在我們生命中掌管一切的主,祂應在我們身上有絕對的權柄.祂有權柄掌管我們的生活,換言之,我們若真正接受耶穌基督為主；就應對耶穌基督有一個完全順服的心.可否讓我們檢查一下我們的生活呢?我們真正接受耶穌基督為主嗎?我們是否在生活上尊祂為主?在生活中是否對祂完全的順服?生活的每一方面,都完全交在祂手裡?自己的家庭,工作,與人交往,喜好,心思意念和願望,這一切讓主耶穌來掌管?祂必需完全作我們的主,不可能是部份的.假若仍有部份生活留在自己手裡,那就等於沒有讓祂在我們生活中作主.要祂在我們生活中作主,可說是基督徒最困難的事情.保羅在此是用這個包括一切的名詞,祂作我們生活的主,不單是指生活的一方面.在生活的一方面順服主是較容易,我可把部份的時間奉獻主.一禮拜有七天,禮拜天完全奉獻給主；任何其他的事都不作,單作教會的工作.有這樣的基督徒生活,可說感謝主!但與主的要求距離仍太遠哩!主所要的並非一週內,只一天奉獻給主,由祂掌管；而其他六天由自己掌管,有權支配自己的生活.如果以主為主,祂就有權支配你七天的生活.你可將十分之一或十分之二奉獻給主,但仍應將其他的十分之七或八由主支配.主所要的是完全的交給祂掌管.今天有許多的基督徒有順服的心.但卻未完全的順服.只是在某些事上順服?而在其他事上仍是自己作主.若真正接受耶穌基督為主,就必需讓主來掌管整個的生活.不單是做工作,還要知道該怎樣去做.在生活上更要知道該怎樣去生活.外表的行動,內心的思想,全由主耶穌作主.如果你要尊稱主耶穌為主,就無論在何時,都要順服,倘若你仍不肯,那麼不要稱主耶穌為主了.保羅說我們的生活,應該在祂掌管之下.當然我們有資格選擇,是要耶穌作我們的主?還是仍然自己作主!
\newline
\newline
(二)就當照樣遵祂而行.(西二6)
\newline
\newline
「就當」這字極其重要,試想我們最初接受主時是怎樣接受的呢?是憑知識嗎?經驗嗎?不是的,我們不是憑任何方法,乃是憑著信心.我們相信了主,一切都該包括在信裡.我們知道所信的是誰,也經過思想去選擇,憑著信心接受了祂.保羅說,你們真的要過那有意義的生活嗎?你們的怎樣接受主耶穌的?那就要照樣地過生活.你是怎樣得救的?那你就怎樣地去生活.我們的生活並非靠自己的思想或計劃,我們得救既是憑著信；照樣,我們生活也需憑著信心.我們怎樣接受耶穌基督為主,就當照樣遵祂而行.我們不能憑著信心開始,而又憑自己的智慧,思想,計劃去完成.整個基督徒的生活,從始至終都需憑信心.你如何接受主耶穌基督為主,那你就照樣地遵祂而行.既憑信心接受主,也應憑信心過有意義的生活.航海人員有兩種計算船方位的方法,一種是憑星宿的位置計算；但單靠此法亦不完全,故另一種是憑整個航程中的方向.若要知自己的位置,就將從出發時,至現在的航程計算.向何方向走了若干哩,再轉方向又走了若干哩.不論向東,南,或北任何方向；就每次轉方向所走的都完全計算後,才可知自己在何位置.這種計算,不是憑著天上星宿的位置,而是憑整個航程的記錄.基督徒也應照此方法計算一下,不要單憑經驗,或成就；或別人的意見,或任何看得見摸得著的事情,要憑著聖經的話,告訴我們,指示我們的方向,這就是保羅所講的意義.你怎樣的接受了基督耶穌,照樣也要憑信心愛祂而生活.
\newline
\newline
如何遵祂而行?保羅在此也告訴了我們.這遵祂而行的「行」字乃是進行式.「生根」及「建造」「堅固」(七節).也是進行式.那就是說:藉著對主的信心,靠著祂生活的這心,是繼續不斷的進行方式.非在某特別的時間,特別關頭,或某一特別重要的事情；乃是在整個的生活裡,不論大小事,都經常不斷地遵祂而行,經常不斷地信靠依賴.在倚靠祂的事上,有幾種不同的態度:
\newline
\newline
(1)　生根──這字保羅表達了,我們與主之間的關係.正如種一棵樹一樣?樹種下就要生根,有了根,樹才能從地土裡得營養,得力量.照樣我們需在主裡生根,根要生到主裡面去,與地土保持繼續不斷地接觸.保持經常與主之間的交通,因而得營養得力量；若不能如此,則無法能遵祂而行.
\newline
\newline
(2)　建造──這字翻譯得不甚完善.此字之基本意義乃是打根基,並非指上面的建築物.在主裡面打了根基,把整個生命的根基打在主耶穌身上,則整個建築物自然是穩固.
\newline
\newline
(3)　信心堅固──此字亦十分難翻譯,是繼續不斷地,建造起來的意思.
\newline
\newline
我們如何建造?(西二7)保羅說:正如你們所領受的教訓,我們生根在主裡面；保持與主不斷地交通；然後建造在根基上.按著祂的教訓,和聖經上的真理,不但經常保持與主交通；更經常不斷地,在上面以真理建造自己.按主的真理建造,用時間不斷地在神的話語上用工夫.
\newline
\newline
「你們感謝的心也更增長.」
\newline
\newline
感謝也是我們靠祂遵祂而行的一方法,若經常對神有感恩的心,也就能建造相信的心.有感恩就等於沒忘記了主,也知道祂的可靠,無論在屬物質,或在屬靈的感恩時,就是你並沒忘記主賜給你的一切.所以感恩,也是我們遵祂而行的方法之一.
\newline
\newline
過一個意義的生活,根基必須建造起來.憑著信心相信祂,信祂是基督耶穌,然後在祂裡面生根；立下根基,繼續建造.遵行及追求祂的話,對祂存感恩的心.你們如何接受主耶穌基督為主,就當照樣遵祂而行.
\newline
\newline


\newpage

\section{第46屆~港九培靈會~鮑會園~第六講}
\label{sec:1189}
\textbf{以基督誇勝}
\newline
\newline
連結: \href{https://www.hkbibleconference.org/session-message/view/1189}{\texttt{https://www.hkbibleconference.org/session-message/view/1189}}
\newline
\newline
\hyperref[sec:1188]{< < < PREV SERMON < < <}
~
\hyperref[sec:toc]{[返目錄]}
~
\hyperref[sec:1190]{> > > NEXT SERMON > > >}
\newline
\newline
歌羅西書 2:8-2:23
\newline
\begin{longtable}{cl}
\hline
\hline
章節 & 經文 (和合本修訂版)\\
\hline
2:8 & \begin{tabularx}{0.7\textwidth}{X} 你們要謹慎，免得有人用他的哲學和虛空的廢話，不照著基督，而是照人間的傳統和世上粗淺的學說，把你們擄去。 \end{tabularx} \\ \\ \relax
2:9 & \begin{tabularx}{0.7\textwidth}{X} 因為神本性一切的豐盛都有形有體地居住在基督裡面； \end{tabularx} \\ \\ \relax
2:10 & \begin{tabularx}{0.7\textwidth}{X} 你們在他裡面也已經成為豐盛。他是所有執政掌權者的元首。 \end{tabularx} \\ \\ \relax
2:11 & \begin{tabularx}{0.7\textwidth}{X} 你們也在他裡面受了不是人手所行的割禮，而是使你們脫去肉體情慾的基督的割禮。 \end{tabularx} \\ \\ \relax
2:12 & \begin{tabularx}{0.7\textwidth}{X} 你們既受洗與他一同埋葬，也就在此禮上，因信那使他從死人中復活的神的作為跟他一同復活。 \end{tabularx} \\ \\ \relax
2:13 & \begin{tabularx}{0.7\textwidth}{X} 你們從前在過犯和未受割禮的肉體中死了，神卻赦免了你們一切的過犯，使你們與基督一同活過來， \end{tabularx} \\ \\ \relax
2:14 & \begin{tabularx}{0.7\textwidth}{X} 塗去了在律例上所寫、敵對我們、束縛我們的字據，把它撤去，釘在十字架上。 \end{tabularx} \\ \\ \relax
2:15 & \begin{tabularx}{0.7\textwidth}{X} 基督既將一切執政者、掌權者的權勢解除了，就在凱旋的行列中，將他們公開示眾，仗著十字架誇勝。 \end{tabularx} \\ \\ \relax
2:16 & \begin{tabularx}{0.7\textwidth}{X} 所以，不要讓任何人在飲食上，或節期、初一、安息日等事上評斷你們。 \end{tabularx} \\ \\ \relax
2:17 & \begin{tabularx}{0.7\textwidth}{X} 這些原是未來的事的影子，真體卻是屬基督的。 \end{tabularx} \\ \\ \relax
2:18 & \begin{tabularx}{0.7\textwidth}{X} 不要讓人藉著故作謙虛和敬拜天使奪去你們的獎賞。這等人拘泥在所見過的幻象，隨著自己的慾望無故地自高自大， \end{tabularx} \\ \\ \relax
2:19 & \begin{tabularx}{0.7\textwidth}{X} 不緊隨元首；其實，由於他全身藉著關節筋絡才得到滋養，互相聯絡，靠神所賜的成長而成長。 \end{tabularx} \\ \\ \relax
2:20 & \begin{tabularx}{0.7\textwidth}{X} 既然你們與基督同死而脫離了世上粗淺的學說，為甚麼仍像生活在世俗中一樣，去服從那「不可拿、不可嘗、不可摸」等類的規條呢？ \end{tabularx} \\ \\ \relax
2:21 & \begin{tabularx}{0.7\textwidth}{X} 見上節 \end{tabularx} \\ \\ \relax
2:22 & \begin{tabularx}{0.7\textwidth}{X} 這些都是根據人的命令和教導，論到這一切都是一經使用就都敗壞了。 \end{tabularx} \\ \\ \relax
2:23 & \begin{tabularx}{0.7\textwidth}{X} 這些規條使人徒有智慧之名，用私意崇拜，自表謙卑，苦待己身，其實在克制肉體的情慾上毫無功效。 \end{tabularx} \\ \\
[1ex]
\hline
\hline
\end{longtable}
保羅已告訴了我們,基督徒在主裡面所佔的地位,也講及我們真正生活的意義.如今保羅接著在此段經文中,向我們解釋,因著主耶穌為我們所成就的,我們在祂裡面已成就了一極大的工作.靠著祂我們今天可勝過一切!保羅在此段經文中,用及「在主裡」或「在基督裡」,有些是用「在祂裡面」有四至五次之多.這可顯出保羅所說,我們真正能有根基和能得著勝利,都是因為我們在主裡面；而並非因我們的知識如何的深,或知道許多人認為的奧秘,乃因我們的生命在基督裡.
\newline
\newline
在此保羅又說,我這樣地提醒他們,乃是要你們謹慎,免得有人用人的理學,或空虛的言語,把你們擄去.這「擄去」一字據解經家的看法,並不是單在靈裡的生命被擄去；乃是說把靈裡面智慧的能力搶去了.故此,說因人的理論的緣故,就失掉了分辨是非的能力,腦子裡那特別的能力被搶去了.因著異端和這些人的思想,你就被這些東西所掌管和控制,故不能分辨出真理,今天有許多人是如此的.那不能分辨真理的人,並非一定是沒有學問和知識或思想.有時或會是更有思想,和有高深的知識,但真正靈裡分辨是非的能力卻沒有了!因為是被這些人的理論所充滿,不能看出是對還是錯!自己認為對的,可能根本上是錯的,因此我們應留意保羅在此所提醒的.我們注重知識,需要學問,需要受造就,但卻也需要認清楚所得的知識與學問是什麼?若不小心,則可能我們的思想就這樣的被擄去了.有許多青年人,在教會裡很愛主,但中學或大學畢業後,有機會到外國讀書,本來是極熱心極愛主的,有見證,也熱心事奉；但在外國讀書一年數載後,就整個改變了!他知識更多,思想受的訓練也更深,但就是因人的理學,虛空的妄言,將他的思想搶去,失掉了那真正分辨是非的能力.保羅說我們若要在這方面得到勝利,就應如何地去追求.
\newline
\newline
(一)得勝的根基
\newline
\newline
「你們既受洗與祂一同埋葬,也就在此與祂一同復活.都因信那叫祂從死裡復活,神的功用.」(西二12)
\newline
\newline
保羅說,我們受洗,與主一同受死,一同埋葬；然後與主一同復活,我們應注意保羅用的「一同」這字.在(羅六5)中保羅給我們的教訓；我們要與主聯合,我們若在主死的形狀上與祂聯合,也要在祂復活的形狀上與祂聯合.保羅說及與主聯合這事是連接的,在死的形狀上與祂聯合,也必在復活上與祂聯合.若我們與祂一同死,也必照樣與祂一同復活,因若真正在生命上與主聯合,保羅說我們就藏在主裡面.若真正如此聯合了在主裡面,生命在主裡面,而不是外表的.也不是說要知識有多少,或在教會中的關係如何,工作如何,而且你的生命與主的關係如何?若生命真正隱藏在主裡,與主分不開,聯合在一起.那麼,試想主耶穌為我們在十字架上受死,死後埋葬,然後從死裡復活.若把死亡與埋葬看成一條線,基督本在死亡的一邊沒有受死,但為我們的罪的緣故,經過了死亡.祂受死了,受死後就如何?祂沒有留在死亡的那一邊.祂經過死,但在另一邊祂又復活過來了.今日我們不能想像主到死亡裡停留下來,沒經過死亡而復活,今天有一個已經成就的事實,我們無人能否認這事實,就是基督經過了死亡,然後從死亡裡復活過來.如果基督這樣經過死亡而復活,我們若在生命上,完全與主耶穌聯合在一起,隱藏在基督裡.我們能否想像主的一部份從死裡復活了,而另一部份卻仍留在死亡的一邊沒有復活.我們會如此想像嗎?主乃是整個的復活,不能留下一半或一部份沒有復活過來的,因此我們真正隱藏在主裡並生命上,真正與主聯合,則基督不會復活而把你留在死亡的那一邊.正如我的手是與我聯合在一起,是我的一部份,今天我來禮拜堂；斷不會把手留在家中,而我卻來了這裡.若是義手可能說會留下,但若是如此；則這手根本就不是我的,更不是與我聯合的.照樣我們若與主生命上聯合在一起,就絕不可能被主放在一邊而主復活.故基督徒真正相信了主耶穌,有主的生命,那在生命上必與基督一同復活.這是不可否認的事實,是永不會改變的,故保羅在(羅六3說:我們應當一舉一動有新生的樣式.生命已與主一同復活,就當過著復活生命應過的生活.若你是已復活的人,則保羅在羅馬書六章末與七章開始說:若我們真正地從死裡復活我們就脫離過去罪惡律法等一切的綑綁,這些都已在我們身上失去了效用,因我們不再在這裡面.正如犯貪污罪的葛柏要逃離香港,因為這樣在港犯了法,在香港的法律之下,香港的法律要按公理制裁他,但當他離開了香港,則香港法律的力量就不能達到制裁他,因英國不在香港法律統治之下,於是他就脫離了香港法律之勢力範圍.照樣,我們若已與主一同死亡,脫離了死亡在我們身上應有之勢力,從死裡出來了,又復活了,就進入另一範圍裡去生活.過去能叫我們死亡的律法,在我們身上已失掉了作用；若掌管我們的律法在我們身上失去了作用,則今天我們不必再受他的控制.保羅用此真理向我們解釋,所以我們和主已經同死又同復活了,就脫離過去那能叫我們死的勢力,進入新的生活裡.今天我們在主耶穌裡面,就已進入到新的生活範圍,這新生活範圍裡有何好處?保羅在此告訴我們.
\newline
\newline
(1)生命之豐盛.
\newline
\newline
「因為神本性一切的豐盛,都有形有體的居住在基督裡面；你們在祂裡面也得了豐盛.」(西二9-10)
\newline
\newline
這是我們與主同復活的第一樣好處.我們有一個豐盛的生命,或可說我們今天有「生命的豐盛.」在中文聖經有與此差不多的名詞,約翰福音十章說:「我來了是要叫羊得生命,並且得的更豐盛.」這「豐盛」與在此所說的不一樣.在此之「豐盛」一字甚難解釋,應可翻作完全,是完完全全的.我們今天在主裡,應有一完全的生命,因神的一切完滿,都在主耶穌裡.故我們若在耶穌基督裡,則我們的生命,也可以得著一切的完滿.這「完滿」在這段經文中,可包括下列各方面的完全.
\newline
\newline
(a)生命長進而至像主一般的完全──因神的一切都在主身上表現出來.神一切的豐盛,都在基督裡；我們也應將此一切的豐盛,進入我們裡面.那麼屬靈方面的生命,才是有長進和成熟的.
\newline
\newline
(b)屬靈的享受得完全──我們可想像保羅在此卷書中,所解釋屬靈生命的好處,這些好處,我們今天也有份,我們是否有這完全的生命?
\newline
\newline
(c)生命表現上的完全──在歌羅西三章中,保羅說及生命各方面的表現,屬靈生命表現上的完全,這完全是非常的完全的.
\newline
\newline
今天我們在主裡得到如何的完全呢?這是我們追求的.不是單一兩方面,而是追求最完全中的完全.我們的屬靈生命未長進到此程度,可是卻要在享受和表現上,追求達到,不盡本份又不肯付代價.正如使徒行傳八章中的西門想得彼得的能力,卻又不願付上彼得所付的代價.我們今天在許多屬靈事上,都可以看見這西門的例子,屬靈的事上,是沒有捷徑可走的.要在屬靈事上追求長進則應多到神面前,多禱告,讀神的話,在生命聖潔上追求,在順服上去追求,付上這代價.屬靈的事,應用屬靈的方法完成,神的標準,需用神的方法,才能達到.不能用自己的方法去完成神的目的.教會中今日亦如此,太著重表面,也太注重成功,沒有在神的話中所說的完全上追求.
\newline
\newline
(2)內心的平安
\newline
\newline
「你們從前在過犯,和未受割禮的肉體中死了,神赦免了你們一切過犯,便叫你們與基督一同活過來；又塗抹了在律例上所寫,攻擊我們有礙於我們的字據,把他撤去,釘在十字架上.」.(西二13-14)
\newline
\newline
這裡保羅用了當時習慣上的比喻,因釘十字架一定把罪名公報,寫在紙上釘在十字架上；好讓人知道為何罪而被釘.主被釘時,彼拉多找不出罪來,故寫上這是猶太人的王,這是存下的罪案.照樣我們每個得罪神的人,每人都有案底,一切罪都已寫上.啟示錄中說:在白色大寶座面前,神的書卷要打開,按著所寫的而定每人的罪；因這是罪案,無法逃避的.惟一逃避的方法是信主耶穌,因主的寶血已將信主的人的罪塗抹,都撤去了撕毀了,在主的十字上已替我們解決了.今日我們生命中並不完全,與不信主的人一樣.或有時比他們更壞,可是卻有一特別的地方,那就是我們在神面前,過去一切所犯的罪都被塗抹,故我們不必再為過去的罪而耽心.因已完全除掉,在十字架上已完全被毀壞了,今日神因主耶穌的緣故,不但赦免我們且將案底都除去；在神面前我們有完全的平安,不必再有任何掛慮.
\newline
\newline
(3)在世人面前誇勝
\newline
\newline
「既將一切執政的掌權的擄來,明顯給眾人看,就仗著十字架誇勝.」(西二15)
\newline
\newline
能在世人前顯出自己的好處,那並不是自己所能,而是主耶穌得來的.因此也能顯出神在我們身上的恩典,那就是──
\newline
\newline
(a)無案底.
\newline
\newline
(b)罪已得赦免.
\newline
\newline
(c)過得勝的生活.
\newline
\newline
(d)勝過世上的哲學.
\newline
\newline
(二)在何方面可誇勝──得勝的範圍是什麼.
\newline
\newline
「你們要謹慎,恐怕有人用他的理學,和虛空的妄言,不照著基督,乃照人間的遺傳,和世上的小學,就把你們擄去.」(西二8)
\newline
\newline
「理學」原意即今之「哲學」.這字有兩方面的意義:一方面是嚴格地在學術上所說的哲學,這是學問,許多人一生的研究它.保羅說要防避哲學,并非說反對哲學,因他對此也有研究,在此的意思,乃是要防避以自己為中心,武斷抵擋神的那種哲學.這個哲學保羅說要防避.另一意思就是例如今日的哲學博士,其實是物理學的哲學博士,又如人生哲學,是指人生基本的看法,這也是保羅所說的世上小學.這字甚難翻譯,呂鎮中牧師為譯這字而譯出幾個字意,以表達之:
\newline
\newline
a 世界心質之靈.
\newline
\newline
b 自然界的原質.
\newline
\newline
c 世俗初淺的宗教觀.
\newline
\newline
這都不易領會,但可表達其意,實際是代表人生最基本的看法,當時人生的基本看法,乃是星宿的迷情,對世俗事的迷信.也就是說,假若人相信觀念,就是將自己放在這思想上,由它控制,這是危險的；因這樣會將責任推卸在這些迷信上.結果也用這方法解決靈裡的問題.其實我們在神面前,就能勝過一切；自己不能負起責任,在生命上對付,則永遠不能有長進.今日在基督裡能誇勝,乃因在基督裡能勝過人的哲學.
\newline
\newline
(2)勝過律法的捆綁
\newline
\newline
在(西二16-19)保羅解釋說:我們拘泥那飲食上,節期,月朔,安息日的規條；其實這真正的是故意的謙虛,只不過是對別的東西的敬拜,對屬靈方面毫無幫助.今日仍有人有此觀念守此規條,以為不然神就不喜悅,也不能顯出是個基督徒.許多教會亦如此,當然,我們在生活上應有規條,但卻不能用規條以達到標準的方法,這是故意的謙虛.規條不能為我們成就任何事情,今天我們已有主的生命；生命的表現,決不是在表面上遵行些規條就算.
\newline
\newline
(3)勝過人迷信的行為
\newline
\newline
在(西二20-23)保羅說:服從那不可拿,不可嘗,不可摸等類的規條.這一切都不是舊約聖經中的,而是當時迷信的思想.當時許多基督徒覺得生活在世上,要隨從這種風俗,但保羅說我們不應隨從.照樣在今日的圈子裡,是否有這樣的迷信行為?過年過節,婚禮......等,許多都是迷信的遺傳,但今日的基督徒卻也隨從去做.
\newline
\newline
我們在主裡不單勝過人的哲學,勝過律法綑綁.也應勝過人迷信的行為,因為這一切對我們屬靈生命長進上毫無幫助.今天要追求的是在生命上與祂聯合,與祂同死同復活,生命才能豐盛及完全；能在人前為主誇勝,願主激勵我們.
\newline
\newline


\newpage

\section{第46屆~港九培靈會~鮑會園~第七講}
\label{sec:1190}
\textbf{新人的生活}
\newline
\newline
連結: \href{https://www.hkbibleconference.org/session-message/view/1190}{\texttt{https://www.hkbibleconference.org/session-message/view/1190}}
\newline
\newline
\hyperref[sec:1189]{< < < PREV SERMON < < <}
~
\hyperref[sec:toc]{[返目錄]}
~
\hyperref[sec:1191]{> > > NEXT SERMON > > >}
\newline
\newline
歌羅西書 3:1-3:17
\newline
\begin{longtable}{cl}
\hline
\hline
章節 & 經文 (和合本修訂版)\\
\hline
3:1 & \begin{tabularx}{0.7\textwidth}{X} 所以，既然你們已經與基督一同復活，就當求上面的事；那裡有基督，坐在神的右邊。 \end{tabularx} \\ \\ \relax
3:2 & \begin{tabularx}{0.7\textwidth}{X} 你們要思考上面的事，不要思考地上的事。 \end{tabularx} \\ \\ \relax
3:3 & \begin{tabularx}{0.7\textwidth}{X} 因為你們已經死了，你們的生命與基督一同藏在神裡面。 \end{tabularx} \\ \\ \relax
3:4 & \begin{tabularx}{0.7\textwidth}{X} 基督是你們的生命，他顯現的時候，你們也要與他一同在榮耀裡顯現。 \end{tabularx} \\ \\ \relax
3:5 & \begin{tabularx}{0.7\textwidth}{X} 所以，要治死你們在地上的肢體；就如淫亂、污穢、邪情、惡慾和貪婪—貪婪就是拜偶像。 \end{tabularx} \\ \\ \relax
3:6 & \begin{tabularx}{0.7\textwidth}{X} 因這些事，神的憤怒必臨到那些悖逆的人。 \end{tabularx} \\ \\ \relax
3:7 & \begin{tabularx}{0.7\textwidth}{X} 當你們在這些事中活著的時候，你們的行為也曾是這樣的。 \end{tabularx} \\ \\ \relax
3:8 & \begin{tabularx}{0.7\textwidth}{X} 但現在你們要棄絕這一切的事，就是惱恨、憤怒、惡毒、毀謗和口中污穢的言語。 \end{tabularx} \\ \\ \relax
3:9 & \begin{tabularx}{0.7\textwidth}{X} 不要彼此說謊，因為你們已經脫去舊人和舊人的行為， \end{tabularx} \\ \\ \relax
3:10 & \begin{tabularx}{0.7\textwidth}{X} 穿上了新人，這新人照著造他的主的形像在知識上不斷地更新。 \end{tabularx} \\ \\ \relax
3:11 & \begin{tabularx}{0.7\textwidth}{X} 在這事上並不分希臘人和猶太人，受割禮的和未受割禮的，未開化的人、西古提人、為奴的、自主的；惟獨基督是一切，又在一切之內。 \end{tabularx} \\ \\ \relax
3:12 & \begin{tabularx}{0.7\textwidth}{X} 所以，你們既是神的選民，聖潔、蒙愛的人，要穿上憐憫、恩慈、謙虛、溫柔和忍耐。 \end{tabularx} \\ \\ \relax
3:13 & \begin{tabularx}{0.7\textwidth}{X} 倘若這人與那人有嫌隙，總要彼此容忍，彼此饒恕；主怎樣饒恕了你們，你們也要怎樣饒恕人。 \end{tabularx} \\ \\ \relax
3:14 & \begin{tabularx}{0.7\textwidth}{X} 除此以外，還要穿上愛心，因為愛是貫通全德的。 \end{tabularx} \\ \\ \relax
3:15 & \begin{tabularx}{0.7\textwidth}{X} 你們要讓基督所賜的和平在你們心裡作主，也為此蒙召，歸為一體。你們還要存感謝的心。 \end{tabularx} \\ \\ \relax
3:16 & \begin{tabularx}{0.7\textwidth}{X} 當用各樣的智慧，把基督的道豐豐富富的存在心裡，用詩篇、讚美詩、靈歌，彼此教導，互相勸戒，以感恩的心歌頌神。 \end{tabularx} \\ \\ \relax
3:17 & \begin{tabularx}{0.7\textwidth}{X} 你們無論做甚麼，或說話或行事，都要奉主耶穌的名，藉著他感謝父神。 \end{tabularx} \\ \\
[1ex]
\hline
\hline
\end{longtable}
保羅告訴我們,我們已脫離律法的捆綁,不再在法律之下,亦不受迷信之思想所限制.但雖然勝過了,卻並非說可以放縱自己任意而生活,因為我們不在律法之下,但卻在基督的律法底下.我們雖然脫離了律法之捆綁,卻另有新的標準,基督的律就可作為我們的標準.保羅在這段經文說,基督徒應當追求的方向,基督徒的人生觀,往往會被人誤解.所以保羅要說明解釋,基督徒應有的生活方向.
\newline
\newline
(一)思念上面的事
\newline
\newline
「所以你們若真與基督一同復活,就當求在上面的事,那裡有基督坐在神的右邊.你們要思念上面的事,不要思念地上的事.」(西三1-2)
\newline
\newline
「思念」意思是整個的思想及心情放在這事上,若我們追求屬靈的事,愛屬靈的事,愛我們的主,追求親近主,應該從那方面著手呢?人的感情是不能勉強或命令的,但卻可以培養起來.我們可以控制思想及知識,使我們去思想及多知道所應知的事情.當多思想這件事時,自己就會有喜愛和渴望之心；如果我們培植我們的心意,培植我們對主的感情；那就應先從思想著手,你的思想要多讓主來充滿.充滿了主,自然對主的渴慕,就一天一天增長.若思想被其他的事或所充滿,就沒有機會思想主的事,屬靈的事,屬靈的心意就不能長進.我們不能命令我們情感,但有辦法控制思想；思想天上屬主的事,不要被屬世的事充滿.這並非說不願念地上的事,或不負責任,乃是叫我們的心,不要整天想著地上的事,要思念屬靈的,有價值的事.保羅說出幾個思念上面的原因.
\newline
\newline
(1)與主同死
\newline
\newline
「因為你們已經死了」(西三3)
\newline
\newline
「死了」是一次已經成了的動作,凡在主裡的人,一次已與主同死,不論你的經驗或感覺如何,這總是事實.我們應把與主同死看作一事實.死的生命怎能過活呢?基督徒與主同死了,已沒有非基督徒的生命,怎能過非基督徒的生活呢?已向世界死了,又怎能在世界活著呢?一個活在世界裡的基督徒,是指靈意活在世界裡.我們與主同死,又與主同復活.乃是已成事實的事.若沒有與主同復活,就尚未與主同死.若已與主同死,就不可能沒有與主同復活.主已從死裡復活,不會把你留在死裡.你已與主同經過死亡進入復活裡,又怎能將自己看作仍在死亡之前的階段呢.所以一個在完全罪裡生活的基督徒,實在是自相矛盾的名詞,這是不可能.若仍活在罪裡,完全思念屬地的事,那根本就沒有與主同死.今日屬主的人,已與主同死,我們就不能思念屬地的事.
\newline
\newline
(2)有新生命的原則
\newline
\newline
「你們的生命與基督一同藏在神裡面.」(西三3)
\newline
\newline
我們的生命,有完全不同的原則.若與(加二20)比較,就知保羅的意思.今日活著的不再是我,若今日我作基督徒,仍是讓自己活著,這根本就不是基督徒.如果是基督徒,就不能再讓自己活著,仍是要主在我裡面活著.如果每一件事,仍用自己的看法和標準去衡量,或以自己喜好為準；那麼基督又怎能活在你裡面呢?所以我們必須與主同死,讓基督活在我裡面,然後才會有新的生活原則.
\newline
\newline
(3)在主裡面有新盼望
\newline
\newline
「祂顯現的時候,你們也要與祂一同顯現在榮耀裡.」(西三4)
\newline
\newline
我們只知道主要再來,但卻未知道其餘另外的一個意義.這個意義就是我們都要改變,並要與祂永遠同在.要在那裡與祂同享受,在創世以前所有的榮耀.若基督再來與我無關,那就不是我們榮耀的盼望了.保羅說,我們與祂一同顯現,在榮耀與祂永遠同在.將有一天我們會在榮耀裡與祂同在.在今天我們應怎樣的生活,預備自己阿!我們是否與配主一同在榮耀裡,向世人顯現?如果我們的生活不配,那麼將來在榮耀裡時,是多麼難堪阿!我們真不配,沒有把自己預備好,那時恐怕不是榮耀,而是羞恥了.此外更有嚴重的後果,在國王婚筵的比喻稱:所請都不來,結果把街上的窮苦人都請來了.當坐席時,國王看見有一人沒穿禮服,他沒預備好,不配在此場合出現；只好丟在外面,哀哭切齒了!今天我們若不預備,那就不配在榮耀裡出現.因此保羅吩咐我們,不要思念地上的事,要思念上面的事.
\newline
\newline
(二)在地上生活應有的表現:
\newline
\newline
Ａ　要治死我們在地上的肢體
\newline
\newline
「所以要治死你們在地上的肢體,就如淫亂,污穢,邪情,惡慾和貪婪.」(西三5)
\newline
\newline
「治死」,「肢體」,這名詞當然是屬靈的意義而言,並非單是肉體的治死而已.正如耶穌說手犯了罪就把他砍掉,一雙眼睛犯罪就把他挖掉.但是難道把這犯了罪的手砍掉了,那罪就可解決了嗎?若如此則砍手的人就永不再犯罪了.事實上你的罪並非在手上,而是在你的生命裡.所以我們應治死肢體,而不是單單砍掉就算,應在靈裡將他治死,正如在肉體上將肢體砍掉一般!不論如何痛苦或艱難,都要如此；因為這是你的攔阻,你的罪一定要對付把他除去才行.
\newline
\newline
同時這「治死」一字,乃是繼續不斷地進行式.人若死了,肢體還活著嗎?還要繼續不斷地治死?這是要很現實地來看我們的生命:信主的人,都已有聖潔的生命,但不論神學的立場如何,看法如何,我們都承認在我表面有罪的力量存在.我活在肉體中一天,仍有罪的力量在我裡面.我們仍不能完全脫離這力量,而過著百分之百的聖潔生活.我們若要過聖潔的生活,仍需要經常地掙扎,經常靠主的恩典；與壓制著我那罪的力量.這就是保羅說要治死他,把他當死人看待,憑信心相信這肢體已死了.如貪心,妒忌,無愛心,壞脾氣,說謊......,常會發動使我們犯罪,我們要隨時把他治死.昨天對付了,今天仍要對付,隨時地將他治死.每次的治死,肉體的力量就會一次比一次軟弱,屬靈的生命就會一天比一天更長進.
\newline
\newline
Ｂ　脫去舊人穿上新人.(西三9-10)
\newline
\newline
脫去舊人穿上新人是非常重要,聖經常用這方式表現我們的生活.把生活比作衣服,意思十分重要,如何脫去舊人?如何穿上新人?舊人又是什麼?就是淫亂,污穢,邪情,惡慾,貪婪──即拜偶像──,以及惱恨,忿怒,惡毒,譭謗,並口中污穢的言語,這一切都要脫去；正如脫去舊衣服一般,然後穿上新衣服.脫舊衣服,是自己積極地把它除去,並非讓衣服在身上自己破爛,照樣除去舊人亦如此.要對付自己的老我──舊人──不能慢慢地對付,而是要積極地把它除掉,立刻除去舊人的生活；無論是邪情惡慾,或是貪婪說謊,或口中污穢的言語,都一次過地除去.因若不脫去舊衣服,就不能穿上新衣服,脫去舊人可分兩方面.
\newline
\newline
(a)除去行為上的罪
\newline
\newline
在(西三5)說及肉體的邪情私慾,今天許多人認為這些不是罪.世人的眼光以為這是自然的發展,或許有人說這是新潮,但新潮則與屬靈的事剛好相反.
\newline
\newline
(b)除去情感思想上的罪
\newline
\newline
在(西三8)說及精神上或心理上思想上所犯的罪:惱恨,忿怒,惡毒,毀謗並口中污的言語.我們口中是可說這種言語,污穢的故事,沒有根據的謠言；不造就人而只批評人的話語等,一概都要除去.
\newline
\newline
Ｃ　新人的生活
\newline
\newline
(a)愛控制一切(西三12-13)
\newline
\newline
我們應追求憐憫,恩慈,謙虛,溫柔,忍耐的心,彼此包容,彼此饒恕,這是對人應有之態度.但保羅又說:愛是聯絡全德的(西三14)這一切之外,還要存有愛,「愛」原來的意思乃是將一條帶子把腰束上.古時的人用帶子將那寬大的衣服束上,以便工作.也就是說「愛」好像一條帶子,能控制整件衣服.愛要作調劑的標準,無論是謙虛或忍耐其他,都以愛為出發點,以愛控制一切.
\newline
\newline
(b)主的平安在心中(西三15)
\newline
\newline
因我們在主裡過新人的生活,並知道這是主所悅納的,心靈在主裡就有完全的安息.
\newline
\newline
(c)把基督的道豐豐富富的存在心裡(西三16)
\newline
\newline
若心裡充滿了神的道,口裡所講的自然也就是神的道.心裡充滿什麼,口裡所發出來的也是什麼,這節經文與(西三2)互相呼應.「要思念上面的事」(西三12).「用神的話充滿我們」(西三16)若我們被神的話所充滿,則會用詩章,頌詞,靈歌,彼此教導,互相勸戒,心被恩感,歌頌神!
\newline
\newline


\newpage

\section{第46屆~港九培靈會~鮑會園~第八講}
\label{sec:1191}
\textbf{新人的責任}
\newline
\newline
連結: \href{https://www.hkbibleconference.org/session-message/view/1191}{\texttt{https://www.hkbibleconference.org/session-message/view/1191}}
\newline
\newline
\hyperref[sec:1190]{< < < PREV SERMON < < <}
~
\hyperref[sec:toc]{[返目錄]}
~
\hyperref[sec:1192]{> > > NEXT SERMON > > >}
\newline
\newline
歌羅西書 3:18-3:18
\newline
\begin{longtable}{cl}
\hline
\hline
章節 & 經文 (和合本修訂版)\\
\hline
3:18 & \begin{tabularx}{0.7\textwidth}{X} 你們作妻子的，要順服自己的丈夫，這在主裡面是合宜的。 \end{tabularx} \\ \\
[1ex]
\hline
\hline
\end{longtable}
保羅在前面說及基督徒生活應有之表現,無論在甚麼情況,都應有新人的表現.繼而說及基督徒與人的關係,特別在家庭中之生活；聖經特別注重關於家庭方面的教訓.因為在家庭中,才會將真正的自己顯露出來.我們在別人面前,很易會裝出使人喜悅的態度；但在家中,卻是自由的可隨便生活.在這裡保羅卻告訴我們,在家中的生活更為重要,故更應謹慎其次更重要的原因,家庭是我們在基本的生命裡,學習的地方.
\newline
\newline
聖經將我們將來要去的地方為天家,稱天上的神為我們的天父,我們是神的兒女.神把天上將來的關係,用地上家庭的關係來表達；將來我們要在天家的生活,在天家如何生活,今天就要在地上的家庭中學習.將來我們在天上,有怎樣生活的表現,今天在地上的家中就該預備好.今日地上的生活,乃是將來天家的準備；此家庭生活十分重要.家庭許多的關係,都是在屬靈關係的預表和代表.婚姻關係:保羅說丈夫與妻子之關係,如同基督與教會關係.整個家庭關係,都有屬靈之價值,因此聖經注重家庭生活的關係.
\newline
\newline
歌羅西書與以弗所書,其中有許多相似之處.可能歌羅西教會,亦讀以弗所書；以弗所書亦傳到歌羅西去.因此歌羅西讀過這兩卷書,則可更完全知保羅的教訓.因保羅在歌羅西書中基本之目的的及計劃之緣故,他寫到家庭關係方面十分簡單.好像把整個的教訓縮小起來,在以弗所書則更解釋,故當思想家庭關係時,就要看以弗所書中的教訓.在(弗五22-六9)都講同樣的教訓.若細心讀此段經文則更清楚,現分三段講述.
\newline
\newline
(一)夫妻間應有之教訓與責任.(西三18)妻子應順服丈夫,作丈夫的,當愛自己的妻子.在廿二節中,保羅又說我們當存敬畏基督的心來彼此順服.所以夫妻間的重要關係就是:
\newline
\newline
(1)應用愛心彼此順服──作丈夫應愛妻子,作妻子的應順服丈夫.有人對這問題走向了極端,以為丈夫應叫妻子順服,應管制她；在妻子身上有絕對的權柄.其實保羅並沒有了此意.他這命令是對妻子而言,那就是說作妻子的應從自己內心,有順服丈夫的心.而丈夫亦必須過著配得妻子順服的生活.照樣,作丈夫的也當愛妻子,此命令是對丈夫而言.丈夫應自己對妻子生出愛她的心,而妻子亦必須過著一種,能使丈夫愛她的生活才行.故夫妻間之生活,都要同樣的謹慎.
\newline
\newline
(2)夫妻聯合成為一體──這是基督教特別的寶貴教訓,是別的宗教所沒有提及的.丈夫與妻子二人成為一體,丈夫是妻子的頭,如同基督是教會的頭一樣.如丈夫是頭,妻子是身體.二人聯合成為一體,不必說這當然是指一夫一妻制度了,這是聖經的教訓.這才是完全的身體.所以頭與身體間就不應有別東西,在其中攔阻,應直接聯合成一體.但若頭與身體間,有別的東西在其中隔著,則這會成一奇怪的身體.照樣,夫妻間亦不能有別的東西參在其間；正如婆婆與媳婦或岳母的微妙關係.作老一輩的,不要參在年青的夫婦中間,引起各種問題；不要干涉在他們的生活方式裡,家庭間的問題,亦常會因此而發生,引起不必要的困擾.
\newline
\newline
(3)丈夫作全家的頭?在何事上丈夫是妻子的頭?據聖經的教訓如下:
\newline
\newline
(a)在屬靈生活上
\newline
\newline
在家庭屬靈生活中,丈夫需有領導的作用.若在家中作丈夫的,不能帶領整個家庭屬靈的生活,不能引導妻子及兒女一同敬拜神；則家中屬靈的生活必受攔阻.故作丈夫的,有非常重要的責任.特別在今日社會裡,許多青年受到社會生活一切的引誘,家庭屬靈的生活,就更為重要了,若作父親的作丈夫的,真能在家中屬靈的事上作頭,能經常在家一同禱告,讀經,及分享一切屬靈的福份；則深信對青年的一代有極大的影響.許多時候作父親就忽略了!
\newline
\newline
(b)在社會上
\newline
\newline
丈夫是妻子的頭,是家庭的頭；故在社會上,代表家庭的應該是丈夫.可以說今天全世界社會裡都是如此,這是神從起初就如此規定的.若有妻子要出頭,代替丈夫的地位,那麼這個家庭定會發生困難.但該明白這並非說,家中一切事情都是由丈夫管制,而是在於責任上,各人應分得清楚；負自己應負的責任,各有權柄處理自己的一切.保羅說家庭生活基本的關係,是彼此順服,彼此相愛.丈夫是妻子的頭,妻子是丈夫的身體.
\newline
\newline
若丈夫未信主又如何?妻子應順服他嗎?這是一個問題.
\newline
\newline
「你們作妻子的,要順服自己的丈夫,這樣若有不信從道理的丈夫,他們雖然不聽道,也可以因妻子的品行被感化過來.」(彼前三1)
\newline
\newline
作妻子的要順服丈夫,因他是丈夫而非因他是基督徒.不論他是否基督徒,都應盡自己本份；作妻子的藉著對丈夫的順服,就能顯出她溫柔的品行；顯出她信主的美好見證,藉此可以感動他信耶穌.若丈夫未信主,而妻子已信主；她愛主是本份的事,但倘若她太過份於做她認為愛主的事,而日夜到禮拜堂,探訪等工作,而忽略家務是不該的.愛主是對的,但要記著你作妻子的本份.若你是未婚的,則你可日夜地去作主工,但若是主帶領你結婚了,就應有婚姻裡的責任.因若日夜為去作主工,而忽略家庭的責任,則會在你那未信主的丈夫面前,失去了見證.照樣,信主的兄弟有位不信主的妻子,也應當用那應有的愛心去愛她.丈夫愛妻子,並非因她是基督徒才去愛她,而是因她是自己的妻子而愛她.不論對方如何,按聖經的教訓,丈夫應愛妻子,妻子應順服丈夫.
\newline
\newline
(二)父母與兒女間之關係.(西三20-21)
\newline
\newline
(a)在主裡要聽從父母
\newline
\newline
保羅說你們作兒女的,在主裡要聽從父母,這是理所當然的.(原文意思是「對的」,是「正確」的.)做兒女的,在主裡應順服父母,無論父母的情況如何,作兒女的都當順服.特別是信主的青年們,你父母若還未信主,則這方面的見證,就更重要.保羅並沒有說你們順服父母親,是因他們所說的話都是對的,而且是好的,是有經驗的有學問的；而是因為這是神的命令,故順服父母是應當的.中國對孝道素稱著於世界,但可惜近年來這孝道已愈來愈少,更是愈來愈不懂如何孝順父母.孝順父母,聽他們的話,不但在我們中國裡是傳統的美德,根本這就是神的命令.
\newline
\newline
保羅在此說,聽從父母要有一條件,乃是要在主裡聽從父母.也就是說,如果父母的命令與主裡有衝突的話,就當在主裡求主賜智慧應如何去做.你是信主的青年,父母不信主,若他們要你在神像前敬拜,那你當如何?求神教導給你智慧.這種命令則不是在主裡,而是與主有衝突.故保羅說我們應在主裡聽從父母.
\newline
\newline
(b)作父母的不要惹兒女的氣.
\newline
\newline
「惹兒女的氣」極難譯,意思是在兒女的思想發展上受到攔阻,因而生氣,這給作父母的有極重要的責任.我們在教養兒女,不要故意無理地惹兒女的氣,使兒女心靈裡受攔阻.作父母的對兒女應格外謹慎,要公平,要守信,要合理.
\newline
\newline
(三)主人與僕人應有之關係.(西三22-23)
\newline
\newline
(a)應盡心服事主人
\newline
\newline
保羅說作僕人的,應盡心服事你肉身的主人.因這是你的責任,並不因他是信主的,你才好好地服事他；而是你在這工作上,所領受的責任問題.應盡自己本份,忠心誠實地做,因為這是服事主,而不是服事人.在今日社會裡,作僕人的情況已漸改變而漸減少.不過這裡不單說作家庭裡的僕人,也許在生意上或工場上,我們都應有同樣的態度去作,為服事主而作.因為你是基督徒,今日你所站的崗位,就是神所安排的,應盡力忠心去作；不可利用其中,特別的機會和關係,去謀取不應有的利益.
\newline
\newline
(b)作主人的應用合理方法待僕人.
\newline
\newline
用合理的方法,待僕人,不可欺壓或不公平的待遇；不可利用自己優越的地位,用不合理的方法去待僕人.應以僕人所應得的去待他.因為我們的主在天上,有一天我們都要站在主的面前,為我們今天所領受的職份來交賬.
\newline
\newline


\newpage

\section{第46屆~港九培靈會~鮑會園~第九講}
\label{sec:1192}
\textbf{保羅最後的勸告}
\newline
\newline
連結: \href{https://www.hkbibleconference.org/session-message/view/1192}{\texttt{https://www.hkbibleconference.org/session-message/view/1192}}
\newline
\newline
\hyperref[sec:1191]{< < < PREV SERMON < < <}
~
\hyperref[sec:toc]{[返目錄]}
~
\hyperref[sec:1131]{> > > NEXT SERMON > > >}
\newline
\newline
歌羅西書 4:2-4:18
\newline
\begin{longtable}{cl}
\hline
\hline
章節 & 經文 (和合本修訂版)\\
\hline
4:2 & \begin{tabularx}{0.7\textwidth}{X} 你們要恆切禱告，在禱告中警醒感恩。 \end{tabularx} \\ \\ \relax
4:3 & \begin{tabularx}{0.7\textwidth}{X} 同時，也要為我們禱告，求神給我們開傳道的門，能宣講基督的奧祕， \end{tabularx} \\ \\ \relax
4:4 & \begin{tabularx}{0.7\textwidth}{X} 使我能按著所該說的話將這奧祕顯明出來，我為此而被捆鎖。 \end{tabularx} \\ \\ \relax
4:5 & \begin{tabularx}{0.7\textwidth}{X} 你們要把握時機，用智慧與外人來往。 \end{tabularx} \\ \\ \relax
4:6 & \begin{tabularx}{0.7\textwidth}{X} 你們的言談要時常帶著溫和，好像用鹽調味，讓你們知道該怎樣應對每一個人。 \end{tabularx} \\ \\ \relax
4:7 & \begin{tabularx}{0.7\textwidth}{X} 推基古是我親愛的弟兄，忠心的僕役，和我一同作主的僕人；他要把我一切的事都告訴你們。 \end{tabularx} \\ \\ \relax
4:8 & \begin{tabularx}{0.7\textwidth}{X} 我特意打發他到你們那裡去，好讓你們知道我們的情況，又讓他安慰你們的心。 \end{tabularx} \\ \\ \relax
4:9 & \begin{tabularx}{0.7\textwidth}{X} 我又打發一位親愛忠心的弟兄阿尼西謀同去；他也是你們那裡的人。他們會把這裡一切的事都告訴你們。 \end{tabularx} \\ \\ \relax
4:10 & \begin{tabularx}{0.7\textwidth}{X} 與我一同坐牢的亞里達古問候你們。巴拿巴的表弟馬可也問候你們。關於他，你們已經得到指示；他若到你們那裡，你們要接待他。 \end{tabularx} \\ \\ \relax
4:11 & \begin{tabularx}{0.7\textwidth}{X} 稱為猶士都的耶數也問候你們。奉割禮的人中，只有這三個人是為神的國與我作同工的，也是使我心裡得安慰的。 \end{tabularx} \\ \\ \relax
4:12 & \begin{tabularx}{0.7\textwidth}{X} 有一位你們那裡的人，作基督耶穌僕人的以巴弗問候你們。他禱告的時候常為你們竭力祈求，願你們能站穩而成熟，充分確信神一切的旨意。 \end{tabularx} \\ \\ \relax
4:13 & \begin{tabularx}{0.7\textwidth}{X} 他為你們、老底嘉和希拉坡里的弟兄多多勞苦，這是我可以為他作見證的。 \end{tabularx} \\ \\ \relax
4:14 & \begin{tabularx}{0.7\textwidth}{X} 親愛的醫生路加和底馬問候你們。 \end{tabularx} \\ \\ \relax
4:15 & \begin{tabularx}{0.7\textwidth}{X} 請問候老底嘉的弟兄以及寧法，和她家裡的教會。 \end{tabularx} \\ \\ \relax
4:16 & \begin{tabularx}{0.7\textwidth}{X} 你們宣讀了這書信，也要交給老底嘉的教會宣讀；你們也要宣讀從老底嘉轉來的書信。 \end{tabularx} \\ \\ \relax
4:17 & \begin{tabularx}{0.7\textwidth}{X} 你們要對亞基布說：「務要完成你從主所領受的職分。」 \end{tabularx} \\ \\ \relax
4:18 & \begin{tabularx}{0.7\textwidth}{X} 我—保羅親筆問候你們。要記念我在捆鎖中。願恩惠與你們同在！ \end{tabularx} \\ \\
[1ex]
\hline
\hline
\end{longtable}
簡單地說,歌羅西書是講及下列幾件事:
\newline
\newline
(一)我們所信的主耶穌是怎樣的一位主.
\newline
\newline
(二)主所成就的救贖如何.
\newline
\newline
(三)如何才能真知道這位主,及真知道祂為我們所成就的救贖.
\newline
\newline
我們得知以上三點知識,就不論在任何異端中,都能站立得住.若要在主裡真正站立得住,就必須有美好生活的表現.保羅將真理講完了,有關生活上的勸告也說了；最後他要用他自己的生活來見證.勸勉歌羅西的弟兄姊妹們,應如何過生活.保羅生活的見證十分的寶貴:
\newline
\newline
「你們要恆切禱告,在此儆醒感恩；也要為我們禱告,求神給我們開傳道的門,能以講基督的奧秘.」
\newline
\newline
保羅勸歌羅西信徒恆切禱告,然後又說也要為我們禱告,求神給我開傳道的門.按保羅的處境,怎能開一個傳道的門呢?在人看來,前面的門都已關了,不但工作門關了；而且更被關在牢獄中.保羅再沒有任何機會,可為主工作了,而且,當時他並不知要坐牢獄多久?甚至能否被釋放也成問題.然而保羅卻在此情況下,求神開傳道的門.試問還有可能的嗎?可是若從神的角度來看,傳道的門是永遠不能關閉的.如同使徒行傳中說,神的道不能被攔阻.只要有一個很順服神的器皿,神能使任何環境,打開傳道的門.保羅雖然在羅馬的監獄中,但卻成了他為主工作的好機會.在腓立比一章中,保羅從幾方面為我們講述.
\newline
\newline
「那在主裡的弟兄,多半因我受的捆鎖,就篤信不疑,越發放膽傳神的道.」(腓一14)而且又說在御營全軍都有人信耶穌.在人看來,保羅在監牢中,再沒有任何辦法可為主工作；但他卻藉在監牢的機會,把福音傳到御營軍裡.因他受的捆綁,使別的事奉主的人更加熱心和努力.另一方面,因他被囚在監裡,他知道歌羅西教會所遭遇的困難,而寫了這封極寶貴的歌羅西書.因著這卷書,使歷代的教會得無限的幫助.保羅在捆鎖中神仍為他打開傳道的門,使福音直接傳到御營軍裡.因他的忠心,使別人更加熱心,藉他亦使教會得幫助.我們要將這寶貴的真理,應用在我們的生活中.也許我們認為在生活中受限制,有缺欠及攔阻,正想為主工作；但覺得各方面的門都關著!我可用什麼方式事奉主呢?好像神在攔阻使我不能去作!我要為主去工作,但主卻叫我缺欠,身體健康有攔阻,或恩賜上有了限制,或環境的缺欠,許多方面都成為我們的攔阻.在這些攔阻的面前,我們可能會灰心失望!向神發怨言,向人埋怨,並且有許多不滿的表現!但在這種困難的環境裡,我們也可以藉著這些困難,來尋求神所要為我們打開的路；藉著困難神可以向我們施格外的恩典.在歷史上可以看出神如何在困難中,使用他的僕人.藉著保羅在監中,神使一個背叛的亞尼西母,認識耶穌.又在監牢裡,曾使獄卒認識主.許多許多這樣的事實,在教會歷史上都記載了.有一姐妹工作的範圍受到攔阻,只能被限制在一地方工作；但就在這一個地方,做了極大的事工.她把整本的聖經,用注音符號寫出來,使許多不識字的人,因此能讀聖經.有一愛主的姐妹,她眼睛忽然失明,在人看來是何等大的遺憾和損失!但神藉著她的損失完成神特別的計劃,她把整本聖經翻作盲人所用的字；因此世上許多盲人,因她的工作而得到祝福.有一姐妹有五小孩,在意外中五個小孩都死去了,在人看來是何等大的損失和痛苦；但她卻在痛苦中看見了傳福音的門,在緬甸為神作了奇妙的工作.她沒有因痛苦而向神發怨言,而是在痛苦中,問神有何機會可為神工作；因而開辦了孤兒院,完成了神的工作.在我們生活中,各有不同的困難,在困難中可以灰心失望.但另一方面,在困難中,也可以轉過來向神說:神啊!在這情況下你要我為你作何事?在人看為攔阻中,你為我開怎樣的傳道的門?保羅在監牢中,求神為他開傳道的門,神也因此而特別地使用祂.
\newline
\newline
(一)對外人應有愛心.(西四4-6)
\newline
\newline
「叫我按著所該說的話,將這奧秘發明出來.你們要愛惜光陰,用智慧與外人交往；你們的言語要常常帶著和氣,好像用鹽調和,就可知道該怎樣回答各人.」
\newline
\newline
保羅在前面曾對歌羅西信徒說:你們脫去舊人,過新人生活.應過著與世分離的生活,這是基督徒應有的生活.我們是聖徒,是與世人分別出來而分別為聖.但與世人的分別,只是在智方面及靈裡的生活目標上的分別；是生活動機的分別,而不是空間距離的分別.神不是要我們離開世界,到另一地方去過安舒的生活,我們經常要與世人接觸.現今看保羅如何教導我們與世人接觸應有的態度.
\newline
\newline
(1)要愛惜光陰──這字含有不要浪費光陰之意.但更重要的,是指我們應抓住我們所有的機會,不可放過任何機會為主工作.應愛惜這機會,神將我們放在不同的環境裡,都有各人所接觸的人.我所接觸不信主的人,可能其他的基督徒永遠未接觸過.我生活的環境,可能別的基督徒永沒機會去.那是說:在我們的生活環境中,我所有的特殊機會是別人所無.神既給我這特別的機會,就應抓緊和利用這機會去作見證.我們應該這樣地愛惜光陰,不可放過任何機會為主作見證,不可說以後或別人更有機會.
\newline
\newline
(2)言語應常帶和氣──言語常帶和氣,好像用鹽調和.這是極大和試探,我們愈熱情或看重於一件事上,心情也愈緊張,也愈急躁.特別我們有一真正火熱的心,愛人靈魂的心,看見別人未信主,走向滅亡的道路；我們的心會急迫.也許在這種情形之下,所說的話語沒帶著和氣.或是我們有真誠的愛心,向人講福音,認為所講的真理非常合理；但對方不肯接受!或用識誚的話相向我們,或找無理之問題和我們辯論；那時心裡忍受不了,於是話語中就失去了和氣.保羅說我們言語,要常帶著和氣,不要失去機會,把真正福音的溫柔美好處表達出來.
\newline
\newline
(3)知道如何回答各人──這裡極難翻譯,基本的意思,乃是說我們知道用什麼態度去待他.這態度的目的,是要藉此使他們接受福音.保羅在此有一好教訓:我們若要領人歸主,不是先把神最嚴勵及極重要的真理告訴他,而是先令他覺得神是位可信靠的神.為要使他信這位可靠的神,我們就應使他覺得我們是他可信任的朋友；然後才能與我們一同信這位可信的主耶穌.所以領人歸主,不是單是傳講福音而已.許多時候我們忽略了這些,平時與人沒有建立可信靠的關係,對人只是表示關心到靈魂,而未有關心到對方的整個人；所以未能真正地得著人心,這是保羅告訴我們的.
\newline
\newline
(二)對朋友及同工應有容忍的心.
\newline
\newline
Ａ　對同工的愛心如同兄弟.
\newline
\newline
(1)提摩太──保羅寫給歌羅西的信中,提及幾個朋友.例如他起初去信歌羅西教會說:「作基督耶穌使徒的保羅,和兄弟提摩太寫信給你們.」提摩太是個青年,是保羅的學生,是保羅帶領他信主的；是他屬靈的兒子,帶領他工作並教導他.但當寫這信時,他沒有說:我的門徒,我的兒子,而是在主裡與提摩太站在同等的地位.
\newline
\newline
(2)阿尼西母──「我又打發一位親愛忠心的兄弟阿尼西母同去.」(西四9)阿尼西母是腓利門的奴隸,歌羅西人都會知道,也許腓利門也是歌羅西人.阿尼西母從主人的家中走了,以為可得自由,在羅馬的法律中,這是極大的罪.他不但是逃走,而且可能還偷了主人一些錢財,逃到羅馬去.但在主的恩典中,在羅馬被帶進監牢裡,在監中遇見了保羅；因而得保羅的引領信了主,而且真正的悔改了.在腓利門書中可見保羅求腓利門,從新接納阿尼西母.現在他打發阿尼西母回腓利門家去.既回腓利門家,當然也就回歌羅西教會,(因腓利門大概是歌羅西人)阿尼西母今已信主,當然也回到歌羅西教會去.在此情況之下,保羅大可以說有一個悔改了的奴隸,阿尼西母要見你們.但保羅卻如此說:「我打發一位親愛忠心的兄弟.」現在他看阿尼西母是在主裡的親愛的兄弟.他見他悔改了,如何在主裡忠心,並順從主命樂意回到腓利門家中.
\newline
\newline
保羅不單是看到提摩太,在主裡所佔的地位及好處；同時也看見阿尼西母的好處.無論何人,保羅都看見他們特有的好處,這顯出對同工的愛心,待同工的態度.
\newline
\newline
Ｂ　對同工謙卑容忍饒恕及愛心.
\newline
\newline
「與我一同坐監的亞里達古問你們安,巴拿巴的表弟馬可也問你們安.」(說到這馬可,你們已經受了吩咐,他若到了你們那裡,你們就接待他.)(西四10)
\newline
\newline
在此提及巴拿巴及馬可,我們記得保羅與巴拿巴及馬可的關係,這在任何人身上都可說是極大的試探.保羅與巴拿巴是在第一次國外佈道時一同去的同工.在同工的工作上,有了些困難,就是巴拿巴的親戚馬可中途離開了他們；保羅對此事感到極失望.故在第二次出外佈道時,保羅與巴拿巴同去,但巴拿巴卻要帶馬可一同行；保羅卻不同意,因此保羅與巴拿巴分手了.這是不快樂的經驗,為了私人的感情,或個人的不同看法,結果發生了誤會.在誤會過了後,保羅自己再思想這問題,並見到馬可的改變.這改變是因為巴拿巴在此情況下沒有離棄他；且繼續用愛心帶領他,教導他；結果馬可成為神使用的僕人.保羅見及此,一定覺得當初自己堅持的意見是不對的.在律法上看這並不是錯,站在正義上也許不是錯,但站在愛心及人情上看,則當初的堅持是太嚴勵,知自己過去所做的不甚適當,故現在他的態度改變了.在改變以後,他在人的面前,也將此態度表明.他吩咐歌羅西人要接待馬可,因馬可是神的僕人,是親愛的弟兄.過去他所反對的人,現在他向人推薦,顯出他對錯誤有改正的心,有饒恕別人的心.在同工間彼此相處,這是何等重要的事!在有些情況下,站在真理上不能妥協；正義上不能改變,但有時在愛中會改變我們.在改變我們的觀點後,能在人前承認在某方面的觀點已改變；公開表達這意思,並非容易的事.
\newline
\newline
(三)因主的緣故勝過自己的情感
\newline
\newline
「我保羅親筆問你們安.」(西四18)
\newline
\newline
在前面保羅用幾種不同的方式,說及向他們問安的意思.現在又說親筆問你們安.保羅當時的處境,是在羅馬的監中,他從未與歌羅西信徒見過面；也未見歌羅西信徒,對他有什麼關心的跡象.甚至其他地方的基督徒,也沒表示關心他.保羅為了福音的緣故,為了主的緣故,落在這困難的境地裡；好像在這情況下,別人都將他忘記了.沒有人關心,在這境況中他的心,應變得何等的痛苦!甚至可用保羅在別卷書中的話說:在裡面可以生出一種苦毒的惡心.我在忠心為主工作,如今竟沒一人關心,既是如此,又何必關心你們!和關心別人呢?既然如此,那只好聽命了!不再理,也不再管,可以就此推搪下去.但保羅並不是如此,他沒在這不順利的環境下,或情感最容易發作之時,而任憑它發作.他仍然表示他與別的基督徒,在靈裡的關心,他沒有覺得自己可憐,或冤枉而忌恨而反抗.反而認為別人關心他,因此他也關心別人.「我保羅親筆問你們安.」這是何等難得的態度,而也是何等重要的態度.許多時候我們落在困難中,沒有人能幫助我們,或在特別情形下,沒法表示那種同情.這時,我們心靈裡合起了反抗,我如此忠心,努力,現在在此情況；無人同情,無人欣賞,無人稱讚,甚至別人離棄我.心裡的反抗及痛苦,這時可能憎恨任何的人.在同此情形下,為何不轉過來看主,也許沒有人明白,沒有人同情.別人不關懷也許別有原因,他不明白我們,也許有他的困難.因此我不向他們忌恨,我心裡沒有反抗；因我知主會明白我,在可能的範圍下,我仍對別人關心；仍然表示我在主裡與他的交通.
\newline
\newline
「你們要記念我的捆鎖.」現在我為主的緣故有困難,你們不要忘記我受捆鎖的原因,盼望藉此得激勵.我為福音受了捆鎖,歌羅西信徒也受了捆鎖,歌羅西信徒也受到激勵,亦更忠心的為福音而努力.
\newline
\newline


\newpage

\section{第47屆~港九培靈會~包忠傑~第一講}
\label{sec:1131}
\textbf{基督的奧秘}
\newline
\newline
連結: \href{https://www.hkbibleconference.org/session-message/view/1131}{\texttt{https://www.hkbibleconference.org/session-message/view/1131}}
\newline
\newline
\hyperref[sec:1192]{< < < PREV SERMON < < <}
~
\hyperref[sec:toc]{[返目錄]}
~
\hyperref[sec:1134]{> > > NEXT SERMON > > >}
\newline
\newline
以弗所書 3:3-3:6
\newline
\begin{longtable}{cl}
\hline
\hline
章節 & 經文 (和合本修訂版)\\
\hline
3:3 & \begin{tabularx}{0.7\textwidth}{X} 用啟示讓我知道福音的奧祕，正如我以前略略寫過的。 \end{tabularx} \\ \\ \relax
3:4 & \begin{tabularx}{0.7\textwidth}{X} 你們讀了，就會知道我深深了解基督的奧祕； \end{tabularx} \\ \\ \relax
3:5 & \begin{tabularx}{0.7\textwidth}{X} 這奧祕在以前的世代沒有讓人知道，像如今藉著聖靈向他的聖使徒和先知啟示一樣， \end{tabularx} \\ \\ \relax
3:6 & \begin{tabularx}{0.7\textwidth}{X} 就是外邦人在基督耶穌裡，藉著福音，得以同為後嗣，同為一體，同為蒙應許的人。 \end{tabularx} \\ \\
[1ex]
\hline
\hline
\end{longtable}
今年培靈研經會早堂的研經會,總題是「基督論」.這個題目十分重要,因為在末世的時代中,有許多異端興起,這些異端對於基督的位格,人性,神性,都有很大的影響.我們信主之人,對基督應有正確的認識,基督徒一生的計劃也應該以基督為中心.
\newline
\newline
關於外邦人得救及基督要來的事並非奧祕,在聖經中從舊約起,早就提到將來有一位彌賽亞要來,神向外邦人也有招聚.在新約中才論到教會的信息,使徒保羅稱之為「基督的奧祕」.意思就是說,我們外邦人與猶太人一同成為基督的身體,同為教會；而基督住在裡面,這是舊約中所沒有的信息,這啟示是神特別交給保羅的.
\newline
\newline
感謝主,我們在祂的恩典中有份,我們與神及主耶穌有密切的關係,這並非偶然的計劃,乃是從創世以前在神的計劃中已定規了.有人以為,人犯了罪,使神為難,神才為人類設立救恩.其實並非這樣,因為從起初神就知道一切,計劃一切了!
\newline
\newline
約翰福音十七章廿一至廿三節說:「使他們都合而為一,正如你父在我裡面,我在你裡面,使他們也在我們裡面；叫世人可以信你差了我來.你所賜給我的榮耀,我已賜給他們,使他們合而為一,像我們合而為一.我在他們裡面,你在我裡面,使他們完完全全的合而為一；叫世人知道你差了我來,也知道你愛他們,如同愛我一樣.」這就是最寶貴的關於信徒的信息,這信息將我們帶到父神裡面,與耶穌基督打成一片.至於基督如何使外邦人和猶太人合成新的一體,在保羅的書信中有所闡明:
\newline
\newline
「因我們是他身體上的肢體.為這個緣故,人要離開父母,與妻子連合,二人成為了一體；這是極大的奧祕,但我是指著基督和教會說的.」(弗五30-32)這樣密切的關係何等寶貴,創造天地萬物之主,拯救了我們罪人,又將我們與祂同在一起,這是保羅所說的教會的奧祕.
\newline
\newline
歌羅西書一章廿七節,說到基督在我們裡面的奧祕:「上帝願意叫他們知道,這奧祕在外邦人中有何等豐盛的榮耀；就是基督在你們心裡,成了有榮耀的盼望.」提前三章十六節提到敬虔的奧祕:「大哉,敬虔的奧祕,無人不以為然；就是神在肉體顯現,被聖靈稱義,被天使看見,被傳於外邦,被世人信服,被接在榮耀裡.」這是聖經中福音的中心,說到神如何將人改變過來,與主同享尊貴,榮耀.人因犯了罪離開了神,本來是沒有盼望的,但神藉著主耶穌基督將我們拯救過來,叫我們與祂的愛子耶穌一同享受神的愛.這也是教會中一項重要的教義,因為我們所傳的只有一位救主除祂以外,別無拯救!我們要將基督放在單獨的王位上,我們不能把基督與三萬萬個偶像同排一列；因為基督是創造宇宙之神的獨生兒子,祂是唯一的救主.我們在這個世界上渡過一生,如果能找到北極星就知道方向；主耶穌就是我們一生的北極星；祂引導我們一生的路程,我們的幸福都在祂裡面.
\newline
\newline
「因為父喜歡叫一切的豐盛,在他裡面居住.」這裡說到神對基督的計劃,一切都是為祂而造；一切的豐盛都在祂裡面,將來祂要作萬國的君主,祂為王超過世上一切的君主,祂要做審判官,審判全人類.這些道理對未重生的人來說是絆腳石,因為他們不承認耶穌是救主.而我們信主的人,要將一生的信仰立在主基督身上,只有祂是我們穩固可靠的根基.我們若有這種觀念,才有正確的途徑.雖然有許多人以為宗教越古老越好,儒家則宣稱他們是最老,最好的宗教(主前五一一年).又有人以為宗教越新越好,因此回教徒便宣稱他們的宗教是現代的宗教(主後五百年),他們說穆罕默德是最後的使者.但是我們所信的耶穌基督是首先的,又是末後的；祂是永生的神,所以祂又是最老的,最新的,最持久的.
\newline
\newline
箴言八章廿二至廿三節說:「在耶和華造化的起頭,在太初創造萬物之先,就有了我.從亙古,從太初,未有世界以前,我已被立.」這兩節經文,中文聖經比英文聖經譯得更清楚.在約翰福音一章一節說:「太初有道,道與神同在,道就是神.」太初比起初更早一點,這太初就是在世界起頭以先就有了,也就是先存的基督.
\newline
\newline
有一個名叫邁爾頓的詩人,他作了很多詩.他中年的時候雙目失明,但心靈的眼睛卻十分明亮.他在「失樂園」一書中描寫失去的樂園是什麼一回事,他想像父神與聖子商量,人犯了罪怎麼辦?父神認為既犯了罪,就要贖罪,聖子耶穌甘願受死代替世人贖罪.
\newline
\newline
在太初的時候,主耶穌做了創造世界的工程師,而且父神與聖子有密切的關係.約翰十七章廿四節說:「父阿,我在那裡,願你所賜給我的人,也同我在那裡,叫他們看見你所賜給我的榮耀；因為創立世界以前,你已經愛我了.這是主耶穌的禱告.祂與父神有深切的來往,父與子之間有愛的關係.愛一定要有對象,不會是單獨的；愛如果沒有適合的對象,必然是痛苦的.父神與聖子一直彼此相愛,並且一同享受極大的榮耀,這榮耀是未有世界以先,就已經有了!
\newline
\newline
「上帝從創立世界以前,在基督裡揀選了我們,使我們在他面前成為聖潔,無有瑕疵.」(弗一4)在創世以前聖父愛基督,也愛我們,那時,祂已想到我們,在基督裡揀選了我們,又為我們預備各樣屬靈的福氣.神有一種追隨的愛賜給我們,為的是要將我們帶到祂的愛中.
\newline
\newline
在彼前一章說到創世以前,父與子已計劃你我的救恩；計劃差遣基督為我們贖罪,可見耶穌基督之死,是父神的計劃與旨意.
\newline
\newline
提多書一章二節說到,神在萬古之先就應許將永生賜給我們,預備要在我們身上顯出祂的恩惠.父與子彼此之間有位格,有交通；同樣,我們都是有位格,有交通的人,三位一體的神互相了解,他們有位格,思想,愛,交通,我們信主的人也因此而有這一切.這是「進化論」中所沒有的寶貴信息.許多科學家試圖研究猴子與人類有何相同之處,但結果是浪費金錢,時間,徒勞無功.
\newline
\newline
我們所信的是昨日,今日,直到永遠的基督,祂是始終不改變的主,祂是阿拉法,也是俄梅戛；祂帶著先前的榮耀,這是無罪的榮耀,並願將這榮耀賜給所信的人.主耶穌多次提到祂是奉了差遣來到世界；祂是第一個奉父神差遣的傳教士,祂奉命來到世界傳救恩的信息.
\newline
\newline
現今,我們看見中國教會已經興起來了,許多青年奉獻自己出去傳福音.神已將福音的火把從西方傳到東方來,這是非常好的現象,因為福音要傳遍天下,主就快來了.主耶穌曾說過,有許多事情從前是隱藏起來的,現在都要顯明出來了!整部聖經都是神的信息,我們若要得著特別的亮光,就要自己在神的面前領受,因為我們所信的都包括在聖經裡面,求神讓我們明白祂的啟示和奧祕.
\newline
\newline
詩篇一百三十五篇五至九節提到上帝創造的奇妙,以後又說到上帝如何將以色列人從埃及救出來；十至廿一節說到上帝如何顧念他們,救他們脫離敵人的手.創造世界並非理論與學說,乃是實實在在的歷史事實,從以色列人的歷史中可以證明.我們認識耶穌基督,認罪悔改,知道自己得著復活的生命,因為我們與復活的主認識.並連結在一起.我們得到主耶穌基督的恩典,祂在永久的愛中為我們預備了拯救,無論貪富,貴賤；只要認識祂,相信祂,祂的恩召便臨到我們,我們藉著福音的真理進入神的愛中,享受祂為我們預備的永遠福份.這是何等寶貴啊!
\newline
\newline


\newpage

\section{第47屆~港九培靈會~包忠傑~第二講}
\label{sec:1134}
\textbf{舊約中的基督}
\newline
\newline
連結: \href{https://www.hkbibleconference.org/session-message/view/1134}{\texttt{https://www.hkbibleconference.org/session-message/view/1134}}
\newline
\newline
\hyperref[sec:1131]{< < < PREV SERMON < < <}
~
\hyperref[sec:toc]{[返目錄]}
~
\hyperref[sec:1136]{> > > NEXT SERMON > > >}
\newline
\newline
彌迦書 5:2-5:27
\newline
\begin{longtable}{cl}
\hline
\hline
章節 & 經文 (和合本修訂版)\\
\hline
5:2 & \begin{tabularx}{0.7\textwidth}{X} 伯利恆的以法他啊，你在猶大諸城中雖小，將來必有一位從你那裡出來，在以色列中為我作掌權者；他的根源自亙古，從太初就有。 \end{tabularx} \\ \\ \relax
5:3 & \begin{tabularx}{0.7\textwidth}{X} 因此，耶和華要將以色列人交給敵人，直到臨產的婦人生下孩子；那時，他其餘的弟兄必回到以色列人那裡。 \end{tabularx} \\ \\ \relax
5:4 & \begin{tabularx}{0.7\textwidth}{X} 他必倚靠耶和華的大能，倚靠耶和華－他神之名的威嚴，站立並牧養，使他們安然居住；因為現在他必尊大，直到地極。 \end{tabularx} \\ \\ \relax
5:5 & \begin{tabularx}{0.7\textwidth}{X} 這位就是和平。當亞述侵入我們領土，踐踏我們宮殿時，我們就立七個牧者，八個領袖攻擊它。 \end{tabularx} \\ \\ \relax
5:6 & \begin{tabularx}{0.7\textwidth}{X} 他們要用刀劍毀壞亞述地和寧錄地的關口。當亞述侵入我們領土，踐踏我們邊境時，他必拯救我們。 \end{tabularx} \\ \\ \relax
5:7 & \begin{tabularx}{0.7\textwidth}{X} 雅各的餘民必在許多民族中，如從耶和華降下的露水，又如甘霖降在草上；他們不倚靠人，也不仰賴世人。 \end{tabularx} \\ \\ \relax
5:8 & \begin{tabularx}{0.7\textwidth}{X} 雅各的餘民必在列國中，在許多民族中，如林間百獸中的獅子，又如少壯獅子在羊群中；他若經過就必踐踏撕裂，無人搭救。 \end{tabularx} \\ \\ \relax
5:9 & \begin{tabularx}{0.7\textwidth}{X} 願你的手舉起，高過敵人！願你的仇敵都被剪除！ \end{tabularx} \\ \\ \relax
5:10 & \begin{tabularx}{0.7\textwidth}{X} 耶和華說：到那日，我必從你中間剪除馬匹，毀壞戰車； \end{tabularx} \\ \\ \relax
5:11 & \begin{tabularx}{0.7\textwidth}{X} 除滅你國中的城鎮，拆毀你一切的堡壘； \end{tabularx} \\ \\ \relax
5:12 & \begin{tabularx}{0.7\textwidth}{X} 除掉你手中的邪術，你那裡也不再有占卜的人。 \end{tabularx} \\ \\ \relax
5:13 & \begin{tabularx}{0.7\textwidth}{X} 我必從你中間除滅雕刻的偶像和柱像，你就不再跪拜自己手所造的； \end{tabularx} \\ \\ \relax
5:14 & \begin{tabularx}{0.7\textwidth}{X} 我必從你中間拔除亞舍拉，毀滅你的城鎮； \end{tabularx} \\ \\ \relax
5:15 & \begin{tabularx}{0.7\textwidth}{X} 我必在怒氣和憤怒中報應那不聽從我的列國。 \end{tabularx} \\ \\ \relax
6:1 & \begin{tabularx}{0.7\textwidth}{X} 當聽耶和華說的話：起來，向山嶺爭辯，使岡陵聽見你的聲音。 \end{tabularx} \\ \\ \relax
6:2 & \begin{tabularx}{0.7\textwidth}{X} 山嶺啊，要聽耶和華的指控！大地永久的根基啊，要聽！因耶和華控告他的百姓，與以色列爭辯。 \end{tabularx} \\ \\ \relax
6:3 & \begin{tabularx}{0.7\textwidth}{X} 「我的百姓啊，我向你做了甚麼呢？我在甚麼事上使你厭煩？你回答我吧！ \end{tabularx} \\ \\ \relax
6:4 & \begin{tabularx}{0.7\textwidth}{X} 我曾將你從埃及地領出來，從為奴之家救贖你，我差遣摩西、亞倫和米利暗在你前面帶領。 \end{tabularx} \\ \\ \relax
6:5 & \begin{tabularx}{0.7\textwidth}{X} 我的百姓啊，當記念從前摩押王巴勒如何籌算，比珥的兒子巴蘭如何回應他，當記念從什亭到吉甲所發生的事，好使你們明白耶和華公義的作為。」 \end{tabularx} \\ \\ \relax
6:6 & \begin{tabularx}{0.7\textwidth}{X} 「我朝見耶和華，在至高神面前跪拜，當獻上甚麼呢？難道獻一歲的牛犢為燔祭來朝見他嗎？ \end{tabularx} \\ \\ \relax
6:7 & \begin{tabularx}{0.7\textwidth}{X} 耶和華豈喜悅千千的公羊，或是萬萬的油河嗎？我豈可為自己的過犯獻我的長子，為自己的罪惡獻我所親生的嗎？」 \end{tabularx} \\ \\ \relax
6:8 & \begin{tabularx}{0.7\textwidth}{X} 世人哪，耶和華已指示你何為善。他向你所要的是甚麼呢？只要你行公義，好憐憫，存謙卑的心與你的神同行。 \end{tabularx} \\ \\ \relax
6:9 & \begin{tabularx}{0.7\textwidth}{X} 耶和華向這城呼叫---看重你的名是真智慧---你們當聽懲罰和派定懲罰的人。 \end{tabularx} \\ \\ \relax
6:10 & \begin{tabularx}{0.7\textwidth}{X} 惡人家中不是仍有不義之財和惹人生氣的變小了的伊法嗎？ \end{tabularx} \\ \\ \relax
6:11 & \begin{tabularx}{0.7\textwidth}{X} 我若用不公道的天平和袋中詭詐的法碼，豈可算為清白呢？ \end{tabularx} \\ \\ \relax
6:12 & \begin{tabularx}{0.7\textwidth}{X} 城裡的有錢人遍行殘暴，其中的居民說謊話，口中的舌頭盡是詭詐。 \end{tabularx} \\ \\ \relax
6:13 & \begin{tabularx}{0.7\textwidth}{X} 因此，我也擊打你，使你受傷，因你的罪惡使你受驚駭。 \end{tabularx} \\ \\ \relax
6:14 & \begin{tabularx}{0.7\textwidth}{X} 你要吃，卻吃不飽，你的肚子仍是空空。你必被挪去，不得逃脫；如有逃脫的，我必交給刀劍。 \end{tabularx} \\ \\ \relax
6:15 & \begin{tabularx}{0.7\textwidth}{X} 你撒種，卻不得收割；踹橄欖，卻不得油抹身；有新酒，卻不得酒喝。 \end{tabularx} \\ \\ \relax
6:16 & \begin{tabularx}{0.7\textwidth}{X} 因為你遵守暗利的規條，行亞哈家一切所行的，順從他們的計謀；因此，我必使你荒涼，使你的居民遭人嗤笑，你們也必擔當我百姓的羞辱。 \end{tabularx} \\ \\ \relax
7:1 & \begin{tabularx}{0.7\textwidth}{X} 我有禍了！我好像夏日收割後的果子，又如收成之後剩餘的葡萄，沒有一掛可吃的，也沒有我心所渴想初熟的無花果。 \end{tabularx} \\ \\ \relax
7:2 & \begin{tabularx}{0.7\textwidth}{X} 地上的虔誠人滅盡了，人世間已無正直的人；他們都埋伏，為要流人的血，用羅網獵取自己的弟兄。 \end{tabularx} \\ \\ \relax
7:3 & \begin{tabularx}{0.7\textwidth}{X} 他們雙手善於作惡，君王和審判官都索取賄賂；位高的人吐出心中的慾望，彼此勾結。 \end{tabularx} \\ \\ \relax
7:4 & \begin{tabularx}{0.7\textwidth}{X} 他們當中最好的，不過像蒺藜；最正直的，不過如荊棘籬笆。你守候的日子，懲罰已經來到，他們必擾亂不安。 \end{tabularx} \\ \\ \relax
7:5 & \begin{tabularx}{0.7\textwidth}{X} 不可倚賴鄰舍，不可信靠密友；甚至對躺在你懷中的妻子也要守住你的口。 \end{tabularx} \\ \\ \relax
7:6 & \begin{tabularx}{0.7\textwidth}{X} 因為兒子藐視父親，女兒抵擋母親，媳婦抗拒婆婆，人的仇敵就是自己家裡的人。 \end{tabularx} \\ \\ \relax
7:7 & \begin{tabularx}{0.7\textwidth}{X} 至於我，我要仰望耶和華，等候那救我的神；我的神必應允我。 \end{tabularx} \\ \\ \relax
7:8 & \begin{tabularx}{0.7\textwidth}{X} 我的仇敵啊，不要向我誇耀。我雖跌倒，仍要起來；雖坐在黑暗裡，耶和華卻作我的光。 \end{tabularx} \\ \\ \relax
7:9 & \begin{tabularx}{0.7\textwidth}{X} 我要承受耶和華的惱怒，直到他為我辯護，為我伸冤，因我得罪了他；他要領我進入光明，我必得見他的公義。 \end{tabularx} \\ \\ \relax
7:10 & \begin{tabularx}{0.7\textwidth}{X} 那時我的仇敵看見這事就羞愧，他曾對我說：「耶和華－你的神在哪裡？」我必親眼見他遭報，現在，他必被踐踏，如同街上的泥土。 \end{tabularx} \\ \\ \relax
7:11 & \begin{tabularx}{0.7\textwidth}{X} 你的城牆重修的日子到了！到那日，邊界必擴展。 \end{tabularx} \\ \\ \relax
7:12 & \begin{tabularx}{0.7\textwidth}{X} 到那日，人必從亞述，從埃及的城鎮，從埃及到大河，從這海到那海，從這山到那山，都歸到你這裡。 \end{tabularx} \\ \\ \relax
7:13 & \begin{tabularx}{0.7\textwidth}{X} 然而，因居民的緣故，為了他們行事的結果，這地必然荒涼。 \end{tabularx} \\ \\ \relax
7:14 & \begin{tabularx}{0.7\textwidth}{X} 求你在迦密的樹林中，以你的杖牧放你獨居的民，你產業中的羊群；願他們像古時一樣，牧放在巴珊和基列。 \end{tabularx} \\ \\ \relax
7:15 & \begin{tabularx}{0.7\textwidth}{X} 我要顯奇事給他們看，好像出埃及地的時候一樣。 \end{tabularx} \\ \\ \relax
7:16 & \begin{tabularx}{0.7\textwidth}{X} 列國看見，雖大有勢力仍覺慚愧；他們必用手摀口，掩耳不聽。 \end{tabularx} \\ \\ \relax
7:17 & \begin{tabularx}{0.7\textwidth}{X} 他們要舔土如蛇，又如地上爬行的動物，戰戰兢兢離開他們的營寨；他們必畏懼耶和華---我們的神，也必因你而害怕。 \end{tabularx} \\ \\ \relax
7:18 & \begin{tabularx}{0.7\textwidth}{X} 有哪一個神明像你，赦免罪孽，饒恕他產業中餘民的罪過？他不永遠懷怒，喜愛施恩。 \end{tabularx} \\ \\ \relax
7:19 & \begin{tabularx}{0.7\textwidth}{X} 他必轉回憐憫我們，把我們的罪孽踏在腳下。你必將他們一切的罪投於深海。 \end{tabularx} \\ \\ \relax
7:20 & \begin{tabularx}{0.7\textwidth}{X} 你必按古時向我們列祖起誓的話，以信實待雅各，向亞伯拉罕施慈愛。 \end{tabularx} \\ \\
[1ex]
\hline
\hline
\end{longtable}
「伯利恆以法他阿,你在猶大諸城中為小；將來必有一位從你那裡出來,在以色列中為我作掌權的；他的根源從亙古,從太初就有.」(彌五2)這一節經文,是很重要的信息,論到主基督從太初就有.「根源」二字,可以說是祂的出現；英文聖經譯為「出發」,即是說從永遠,太初,亙古就有之意.
\newline
\newline
一,基督的影子與形象
\newline
\newline
在舊約中一次又一次地論到主耶穌的事,好像十字架旁邊的影子,這影子在祂未來之先已經有了.主基督生在世上特別是為我們流血捨命.我們要明白聖經,就要找一條適合的鑰匙,這鑰匙就是主耶穌；祂能為我們打開聖經的真理,祂用聖靈光照我們,引領我們.
\newline
\newline
「耶穌對他們說,無知的人哪!先知所說的一切話,你們的心,信得太遲鈍了.基督這樣受害,又進入他的榮耀,豈不是應當的麼.於是從摩西和眾先知起,凡經上所指著自己的話,都給他們講解明白了.」(路廿四25-27)那兩個門徒不明白復活的道理,主就耐心地向他們講清楚.關於祂自己的事,門徒便心裡火熱起來.
\newline
\newline
創世記三章十五節頭一次預言到基督將來要來.從該隱與亞伯獻祭的事上；神悅納亞伯所獻的祭,因為亞伯拉罕是預表主基督的.後來人犯了罪,神用洪水懲罰世界,但又叫挪亞預備方舟；挪亞一家進入方舟,就平安得救,免去洪水的災難.
\newline
\newline
當以色列人出埃及的時候,神用十大災難打擊埃及人.最後一災是擊殺埃及人所有的長子,那時,以色列人的長子得以存活,是因為他們按照神的吩咐,殺了逾越節的羔羊,用羊血塗在門楣門框上.
\newline
\newline
啟示錄十三章八節說:「凡住在地上,名字從創世以來,沒有記在被殺之羔羊生命冊上的人,.......」這「記在生命冊上」是在創世之前,羔羊就是從那時開始的.可以說,在神的心目中主基督代替人捨命,是太初就有的計劃,以色列人殺羔羊獻祭,是指望將有一位要來,代替他們贖罪.以色列人在會幕敬拜上帝的一切儀式當中,有許多象徵主基督的事情；大祭司亞倫一年一次進入至聖所,帶著血為百姓贖罪,亞倫也是預表主耶穌.我們看聖經的時候,要明白神的意思,從開始到末了都是預表基督的救恩.
\newline
\newline
在舊約聖經中,有幾個人的生平,可以預表主基督的:
\newline
\newline
(1)約瑟是基督的預表──(不斷線,不同)「以色列原來愛約瑟過於愛他的眾子,因為約瑟是他年老所生的,他給約瑟作了一件彩衣.」(創卅七3)
\newline
\newline
在我們的想像中,天父如何愛主耶穌呢?祂捨不得愛子受死,就用一些人預表祂,象徵祂.約瑟因受父親特別喜愛而招致哥哥的妒恨.主耶穌也是這樣,祂到自己的地方來,自己人卻不接待祂；當祂在猶太地傳道的時候,祭司長和文士都設法害祂.約瑟被賣至埃及,他的哥哥們脫去他的彩衣,塗上羊羔的血,拿回去給他父親看,說,約瑟被惡獸吃了!但是約瑟彩衣上的血不是約瑟的血,乃是羊的血.當主耶穌掛在十字架上的時候,害祂的人非常高興,以為把祂解決了,那裡曉得這是救恩成就的一天.約瑟的哥哥們賣了約瑟之後,以為約瑟解決了,誰知約瑟後來獲得升高,做了埃及的宰相.我們的主耶穌被釘十架,受死後,第三天復活,升上高天,坐在父神的右邊.約瑟做了埃及全地的宰相,沒有記恨他的哥哥們,反而用愛心接他們到埃及去同住,度過飢荒.主耶穌被釘在十字架的時候,還為釘祂的人祈禱,求父神赦免他們,這些事情都說明約瑟與主耶穌很相似,足以預表基督.
\newline
\newline
(2)摩西是基督的預表──「耶和華你的上帝要從你們弟兄中間,給你興起一位先知像我,你們要聽從他.......我心在他們弟兄中間,給他們興起一位先知像你,我要將當說的話傳給他,他要將我一切所吩咐的,都傳給他們.」(申十八18)
\newline
\newline
神自己說,要興起一位先知是誰呢?就是基督.摩西是帶領以色列人出埃及的元帥；主耶穌是救我們脫離罪惡統治的元帥.摩西在曠野四十年；主耶穌也曾在曠野四十天受魔鬼的試探.摩西在埃及被百姓棄絕,故逃到曠野；後來又回到以色列人中,與他們一同受若,當主耶穌被以色列人拒絕後,神又推開外邦人救恩之門；使我們可以進去蒙恩得救.摩西是神所揀選的領袖,他與神常有交通,他忠心事奉神；希伯來書三章二節兩次提到「盡忠」,說明耶穌在神的家中比摩西更盡忠,比摩西更配多得榮耀.
\newline
\newline
(3)大衛是基督的預表──大衛是以色列的王中最尊貴的一位,他與主耶穌同樣生在猶大伯利恆.他少年時哥哥們都看不起他,不明白為何要膏立他為王,可見神的恩召是何等的奇妙,大衛是神所喜愛的,稱為合神心意的人；主耶穌也是合神心意的人,神親自說:「這是我的愛子,是我所喜悅的.」大衛少年時,會用機弦甩石的方法,打死兇惡的巨人歌利亞,有人用這比喻來預表耶穌對付魔鬼的試探.神是一位滿有耐心的神,無論在舊約或新約的時代,祂對人仍然忍耐等候.大衛雖然被掃羅追迫,好幾次有機會可以殺掃羅,但大衛要等候神的時間；他不敢伸手傷害耶和華的受膏者,他對掃羅可謂百般忍耐,正係主耶穌忍耐等候我們一樣.大衛的朋友約拿單,本來是承繼掃羅王位的人；約拿單愛大衛,是超然的愛,非屬血氣的愛,這是出於神的愛,所以將王位讓給他.我們若愛主耶穌,就要將心中的王位讓給祂.大衛肯認罪悔改,肯流淚禱告；主耶穌也為耶路撒冷的人哀哭,祂說:「哀慟的人有福了；因為他們必得安慰.」(太五4)我們應該為自己的罪哀哭,更要以基督的心為心,為香港千千萬萬未悔改得救的人流淚.
\newline
\newline
二,詩篇中的基督
\newline
\newline
在詩篇中有十六篇稱為「彌賽亞之詩」:
\newline
\newline
第二篇論到主基督為王,將來祂要再來作王,執掌王權,直到永遠.
\newline
\newline
第八篇論到主基督創造宇宙萬物,一切被造者均彰顯主的榮美.
\newline
\newline
第十六篇論到主基督的尊貴,天父讓祂的兒子耶穌基督坐在父神的右邊.
\newline
\newline
廿二篇論到主受苦,受死.
\newline
\newline
廿三篇論到主基督是我們的牧者.
\newline
\newline
廿四篇論到榮耀的王要來,現在我們就要將心門打開迎接祂進來.
\newline
\newline
四十篇論到主基督樂意遵行父神的旨意.
\newline
\newline
四十一篇是論到主自己最親愛的朋友出賣祂.
\newline
\newline
四十五篇論到基督在榮耀中再來,祂有威嚴,尊貴.
\newline
\newline
六十八篇論到基督是有力的元帥,將魔鬼擄掠了.
\newline
\newline
六十九篇是述說義者孤單中的痛苦.
\newline
\newline
七十二篇論到公義的國度來到.
\newline
\newline
八十九篇是論到耶和華與大衛所立的約,直到萬代.
\newline
\newline
一○二篇論到主基督的痛苦.
\newline
\newline
一一○篇論到主基督為王,為祭司.
\newline
\newline
一一八篇論到主基督被賣之夜與門徒所唱的詩.
\newline
\newline
此外,在先知書中也有多處論到基督的事,因時間關係不能一一詳述,茲將經文臚列於後:
\newline
\newline
創三15；賽七14,十一1；五十三章全,六十一1-2；耶廿三6,三十三6；瑪四2.
\newline
\newline


\newpage

\section{第47屆~港九培靈會~包忠傑~第三講}
\label{sec:1136}
\textbf{道成肉身的基督}
\newline
\newline
連結: \href{https://www.hkbibleconference.org/session-message/view/1136}{\texttt{https://www.hkbibleconference.org/session-message/view/1136}}
\newline
\newline
\hyperref[sec:1134]{< < < PREV SERMON < < <}
~
\hyperref[sec:toc]{[返目錄]}
~
\hyperref[sec:1137]{> > > NEXT SERMON > > >}
\newline
\newline
腓立比書 2:5-2:8
\newline
\begin{longtable}{cl}
\hline
\hline
章節 & 經文 (和合本修訂版)\\
\hline
2:5 & \begin{tabularx}{0.7\textwidth}{X} 你們當以基督耶穌的心為心： \end{tabularx} \\ \\ \relax
2:6 & \begin{tabularx}{0.7\textwidth}{X} 他本有神的形像，卻不堅持自己與神同等； \end{tabularx} \\ \\ \relax
2:7 & \begin{tabularx}{0.7\textwidth}{X} 反倒虛己，取了奴僕的形像，成為人的樣式；既有人的樣子， \end{tabularx} \\ \\ \relax
2:8 & \begin{tabularx}{0.7\textwidth}{X} 就謙卑自己，存心順服，以至於死，且死在十字架上。 \end{tabularx} \\ \\
[1ex]
\hline
\hline
\end{longtable}
先知以賽亞說:「......主自己要給你們一個兆頭；必有童女懷孕生子,給他起名叫以馬內利.」(賽七14)以賽亞先知畫了一幅主基督受苦的圖畫,他看見主將一切的罪孽擔在自己身上.耶利米書十三章中,有兩次提到耶和華成為我們的義；父神要派一個代表來替我們受罪.瑪拉基書四章也提到主的日子必要來到.
\newline
\newline
今天所講的題目,有關主基督的降生,祂道成肉身,降生在猶太的伯利恆；祂降生時,受到時間的限制.靈是超過肉體和自然界的,神的事情遠超過人的思想；所以我們可以相信耶穌,得著救恩.如果用我們簡單的頭腦去思想基督是不夠的,因此,人要與神接觸,有交通聯繫；神自然有奇妙的辦法.歌羅西書一章十五節說:「愛子是那不能看見之上帝的像,是首生的,在一切被造的以先.」祂是肉眼不能看見之神的像,不是人手所造成的偶像,可以拿出來給人看.主基督既是不能看見之神的像,神要讓人看見,故此叫主成為人的身體,他也是教會全體之首,祂是元始,是從死裡首先復生的.主基督從死裡復活,這是神的大能,敗壞了掌死權的魔鬼.基督與父神在一切萬物之先,並預先定下效法祂兒子的人,都是主的弟兄,都得著兒子名分,但主基督是在眾弟兄之先,是眾多弟兄中的長子.
\newline
\newline
何謂道成肉身?基督要到世上來,是按著神計劃中的時間,祂在神預定的時間準時出現,絕不誤點.祂生在世上取了人的樣式,由童貞女馬利亞所生,其中包含寶貴的意思,創世記三章十五節說:「我又要叫你和女人彼此為仇,你的後裔和女人的後裔,也彼此為仇；女人的後裔要傷你的頭,你要傷他的腳跟.」這節經文清楚地說到救主是女人的後裔.現在有許多基督徒不承認這個道理,以為基督是沒有父親的,其實,基督的父親就是神,由聖靈感孕所生.基督到世上來,具有人的形像,又有神的性格,兩性在一個人的身體活動,正如提前三章十六節說:「大哉!敬虔的奧祕,無人不以為然；就是神在肉體顯現,被聖靈稱義,被天使看見,被傳外邦,被世人信服,被接在榮裡.」主基督是神,住在人中間,這就是神在肉身的顯現,神本性一切的豐盛,都有形有體的居住在基督裡面.
\newline
\newline
道成肉身的目的,是捨身流血,受死,救贖失喪的人.提多書二章十三節說:「等候所盼望的福,並等候至大的神,和我們救主耶穌基督的榮耀顯現.」省略略了「和」字,就是說,等候我們的救主耶穌基督.千萬不要小看祂,要清楚認識祂的地位是何等神聖,祂降世為人是要拯救我們,只有這個辦法,故此祂親自來了!這樣我們可以看見主基督有著神的位格.當以色列人看見主耶穌之時,人們只知道祂是木匠約瑟的兒子,所以我們需要神光照我們靈性的眼睛.西面在聖殿中,聖靈光照他,使他知道主耶穌是神的兒子,是基督.所以稱頌上帝說:「主啊,如今可以照你的話,釋放僕人安然去世；因為我的眼睛已經看見你的救恩.」(路二28-30)使徒約翰說:「道成了肉身,住在我們中間,充充滿滿的有恩典,有真理；我們也見過他的榮光,正是父獨生子的榮光.......從來沒有人看見神,只有父懷裡的獨生子將他表明出來.」(約一14,18)如果沒有基督降世為人,我們根本就不知道神是怎樣的,因為基督來了,把聖父彰顯給人看.
\newline
\newline
當神造人的時候,是照著自己的形像造的,人因為犯了罪才失去了神的形像,神又藉著基督再一次叫我們看見一個圓滿的神的樣式.主基督來到世界,為人捨命,復活,升天,坐在父神的右邊.祂來是帶著人的身體來,故許多人看見祂是神的兒子,因為他像普通的人一樣,聖靈光照誰,誰就認識祂.彼得作見證說:「我們從前,將我們主耶穌基督的大能,和他降臨的事,告訴你們,並不是隨從乖巧捏造的虛言,乃是親眼見過他的威榮.他從父神得尊貴榮耀的時候,從極大榮光之中,有聲音出來向他說,這是我的愛子,我所喜悅的.」(彼後一16-17)主登山變像的時候,讓三個門徒看見祂真正的榮耀.我們如果還沒有蒙恩,決不能看見祂的大榮光；但感謝主,祂讓我們瞻仰祂的榮美.
\newline
\newline
耶穌基督為馬利亞所生,到了時候祂就降世,正如加拉太書四章四節所說:「及至時候滿足,神就差遣他的兒子,為女子所生,且生在律法以下.」神曾應許大衛,堅立他的王位直到永遠.按照以色列人的家譜,就知道主基督是大衛的後裔,大衛的子孫；祂比所羅門的榮耀大得多.但以理書中所論到的巨像之十個腳趾,被大石砸碎就是論到主的國要滅絕一切的國,存到永遠.主基督既是大衛的後裔,祂就要作王,但祂在世的時候,有許多事情證明他和普通人一樣.祂有愛心,同情,喜樂,祂也曾軟弱,疲乏,憂傷；祂曾受試探,被迫害,並經歷過死亡,祂所感受的和我們一樣.不過祂是無罪的,無罪的基督才能擔當我們的罪,成為慈悲忠信的大祭司,為百姓的罪獻上挽回祭.基督自己又是被獻的羔羊,為世人贖罪.這就足以證明,基督確具有神性與人性,兩性彼此相交,是十分奇妙的.
\newline
\newline
有一次,主耶穌在路上行走,有一個患病的女人挨近祂,摸祂的衣裳的繸子,祂就知道了,立刻顯出大能,醫治那婦人.
\newline
\newline
馬大和馬利亞的弟弟拉撒路死了,主耶穌問她們將拉撒路放在甚麼地方?好像祂不知道似的,他們帶主到拉撒路的墓前；主叫拉撒路復活,復活的大能足以表明基督的神性.也許有人不明白,為何在一個身體當中同時具有兩性,這是因為主基督是完全的人,同時也是完全的神,一切事情,祂都知道,祂是無所不知的神.主耶穌受死後,三日復活,向門徒顯現,曾三次問彼得說:「你愛我麼?」彼得對耶穌說:「主阿,你是無所不知的,你知道我愛你.」耶穌呼召拿但業跟從祂的時候,指著拿但業說:「看哪,這是個真以色列人,他心裡是沒有詭詐的.」拿但業對耶穌說:「你從那裡知道我呢?」(參約一47-50)主基督是人也是神,祂取了人的身體；具有神的一切屬性,祂有權柄赦免人的罪,祂有能力醫治人的病；祂有與人一樣的心思,有時也會受到人所受的限制.祂的人性和神性是平衡的,也是合一的.
\newline
\newline
基督是神的兒子,有許多見證人,施洗約翰,天使加百列,聖父......也為祂作見證.基督從死裡復活,四十日後升天,坐在父神的右邊,祂帶著人的身體.將來我們復活之後所有的身體,也和主基督一樣.那時,我們所看見的主基督,也是一個人的樣子.將來大家在父神面前的時候,都要看見祂.
\newline
\newline
在天主教的教義中,他們對於基督為人,具有人性這一點模糊不清；他們向主祈禱要由主的母親馬利亞代求,由馬利亞作中保將祈禱轉達,這就是對基督的人性不理解之故.主耶穌曾說:「凡勞苦擔重擔的人,可以到我這裡來,.......」所以,我們若要親近主,就當毫無隔膜.我們應該知道,神是全能的,一切榮耀,豐盛都在主裡面.我們所相信的是三位一體的神,主基督將聖父彰顯出來,祂是我們的親密朋友,又是我們的長兄.祂道成肉身,將父神一切計劃,恩典顯明出來.因此,我們要好好清楚認識基督,常常與祂交通,並盼望祂快來!在祂再來的時候,能與祂一同享受永遠的福樂.
\newline
\newline


\newpage

\section{第47屆~港九培靈會~包忠傑~第四講}
\label{sec:1137}
\textbf{基督的恩典}
\newline
\newline
連結: \href{https://www.hkbibleconference.org/session-message/view/1137}{\texttt{https://www.hkbibleconference.org/session-message/view/1137}}
\newline
\newline
\hyperref[sec:1136]{< < < PREV SERMON < < <}
~
\hyperref[sec:toc]{[返目錄]}
~
\hyperref[sec:1140]{> > > NEXT SERMON > > >}
\newline
\newline
提多書 2:11-2:15
\newline
\begin{longtable}{cl}
\hline
\hline
章節 & 經文 (和合本修訂版)\\
\hline
2:11 & \begin{tabularx}{0.7\textwidth}{X} 因為，神救眾人的恩典已經顯明出來， \end{tabularx} \\ \\ \relax
2:12 & \begin{tabularx}{0.7\textwidth}{X} 訓練我們除去不敬虔的心和世俗的情慾，在今世過克己、正直、敬虔的生活， \end{tabularx} \\ \\ \relax
2:13 & \begin{tabularx}{0.7\textwidth}{X} 等候福樂的盼望，並等候至大的神和我們的救主耶穌基督的榮耀顯現。 \end{tabularx} \\ \\ \relax
2:14 & \begin{tabularx}{0.7\textwidth}{X} 他為我們的緣故捨己，為了要贖我們脫離一切罪惡，又潔淨我們作他自己的子民，熱心為善。 \end{tabularx} \\ \\ \relax
2:15 & \begin{tabularx}{0.7\textwidth}{X} 這些事你要講明，要充分運用你的職權勸勉人，責備人。不要讓任何人輕看你。 \end{tabularx} \\ \\
[1ex]
\hline
\hline
\end{longtable}
今天,我要講耶穌基督的恩典.在提多書二章十一節提到,上帝救眾人的恩典,已經顯明出來；提多書三章四至八節也提到主基督的恩典,和祂向人所施的慈愛.第七節說:「......我們因他的恩得稱為義,可以憑著永生的盼望成為後嗣.」各位!我相信我們作基督徒的,有許多人不明白恩典是甚麼意思.人類本來死在罪惡過犯當中,自己毫無辦法,如果不接受救主基督,我們自己絕無指望,所以我們要認清楚,基督的恩典到底是甚麼?
\newline
\newline
恩典──就是我們要得著一些不配得的東西；或無緣無故不必付代價而得著的東西.耶穌基督被釘死在十字架上,代替我們贖罪,完全是出於恩典.如此奇妙的恩典,本來是我們不配得的；這是白白得的恩典,也不必我們付代價.
\newline
\newline
上帝是威嚴的神,祂先要我們怕祂,後來因主基督的受死,使我們的罪得赦,就把懼怕除去了!主耶穌被釘在十字架上的時候,有人譏誚祂說:「他救了別人,不能救自己.」(太廿七42)其實,祂是有足夠的能力救自己的,但因為要救我們,所以祂不能夠救自己.
\newline
\newline
有一個小孩,失足跌落河中,他哭哭啼啼,沒有辦法；有一個人跳落水中將他救起來.這不是偶然,乃是出於恩典.
\newline
\newline
若有一個人患了重病,藥石無靈,醫生束手無策,但藉著禱告,求主醫治得以痊癒.這也是恩典,絕非偶然.
\newline
\newline
若有人掉在很深的井裡,自己無法爬上來；另外有一個人放一條繩子在井中,把他位上來；這個人就是他的救命恩人了!
\newline
\newline
若有一個人在房子裡睡覺,房子著火,他自己無法撲滅；後來消防人員將他救出,這是恩典.
\newline
\newline
在面臨飢荒的地方,有許多人餓得快要死了；你送給他們水和糧食,救子他們的命,這是恩典.
\newline
\newline
若有人在船上遇到風險沉船,救生員將他在海中救起,這是救命之恩.
\newline
\newline
若有人犯了大罪,判了極刑,在牢中等死；忽然有一個人肯代替他受刑,讓他活著；救了他一命,這就是救恩.為何主基督救恩如此寶貴,因為在天下人間,沒有第二個救法,只有祂能救我們.基督的恩典十分重要,是關乎生死關頭的事.林後八章九節說:「你們知道我們主耶穌基督的恩典；祂本來富足,卻為你們成了貧窮,叫你們因祂的貧窮,可以成為富足.」基督是帶著恩典與真理來的,從祂豐富的恩典裡,我們都領受了,而且恩上加恩.
\newline
\newline
今天是恩典的時代,正如主耶穌在會堂中所念的一段以賽亞書:「主的靈在我身上,因為祂用膏膏我,叫我傳福音給貧窮的人；差遣我報告被擄的得釋放,瞎眼的得看見,叫那受壓制的得自由,報告神悅納人的禧年.」(路四18-19).現在正是悅納的時候,拯救的日子；主耶穌大赦的時候,叫凡相信祂的人,罪得赦免.
\newline
\newline
天上的神是公義的神,祂赦免人一定要按公義付代價.「現在我們既靠著祂的血稱義,就更要藉著祂免去上帝的忿怒.」(羅五9)怎樣才能免去上帝的忿怒?必須倚靠基督的寶血.現在,主藉著聖靈設立一個恩典的國度,這恩典的國度設立在你我的心中,好像一個祕密的國家,藏在我們心裡一樣.我們如何才能進入這國度?我們生在那一個國家,就是那一個國家的人；我們要進入恩典的國度,就要藉著重生.耶穌說:「我實實在在的告訴你,人若不重生,就不能見上帝的國.」(約三3)不能進神的國就不能見神的國,所以必須重生才進入恩典的國度.我們蒙了神的恩典,罪得赦免,就被遷到愛子的國度裡.這個國度很普遍的,我們每逢遷居一個地方,必定在那裡找教會聚會；與當地的信徒交通,這是非常寶貴的.因為我們同一個國度,同一個家中；這國度沒有種族,國界,正如主耶穌所說,凡跟從主的,就是天父所賜的人.因此,凡是信主之人,都是同屬一個國度,同作天國的子民,這國度正在發展中.
\newline
\newline
我由美國來香港的時候,路經韓國,在那裡參觀一個世界最大的長老會；這個禮拜堂每個主日有五次崇拜,聚會的人數很多.我看見他們就高興快樂,因為大家都是屬主的人.主基督的國度是不住發展的,有些地方是有形的發展；有些地方是無形地發展.當主耶穌再來的時候,祂的國度都要實現.
\newline
\newline
我們相信基督,是因為主基督用恩典召我們來,若無恩典,我們不敢來就近祂,這恩典就是赦免的恩典；也是使我們稱義的恩典.我們現在做基督徒,成為神的兒女,都是靠這稱義的恩典；我們過去無論犯了甚麼罪,只要誠懇地接受耶穌基督的寶血,一切的罪都可赦免!
\newline
\newline
基督的恩典與慈愛雖然是豐富的,但不能隨便犯罪,若隨意犯罪,就是藐視基督的恩典了.基督的恩典不但能赦罪,而且能幫助我們過得勝的生活,因為我們在世界上生活,常常會遇到苦難和試探,但主的恩典夠我們用,我們只管向祂支取；祂的恩典豐豐富富,能叫我們在真道上站立得穩.有些人一生作搖搖動動的基督徒,好像電梯上上落落一樣,基督徒的靈性生活若是這樣,就站立不穩了!我們應作長進的基督徒,因主的恩典夠用,能使人長進.林前十五章十節說:「然而我今日成了何等人,是蒙神的恩才成；並且祂所賜我的恩,不是徒然的.......」保羅在哥林多前書多次提到神的恩典.我們若領取主的恩惠,就感到自己不配,不敢驕傲.許多人驕傲自滿,是因為不懂得自己靠神的恩典之故,神願意將恩典賜給謙卑的人,越謙卑就越蒙神恩.我們領受了基督的恩典,得了新的力量,新的幫助,才有能力出去為主工作.
\newline
\newline
保羅說:「若因一人的過犯,死就因這一人作了王,何況那些受洪恩又蒙所賜之義的,豈不更要因耶穌基督一人在生命中作王麼.」(羅五17)這是保羅個人的見證.今天,我們作基督徒的,是否有一位生命的王?是否與基督同掌王權,與主同享這福份?神為我們所預備的恩典,是叫我們作得勝的基督徒,在人生的過程中,常常作有見證,有恩惠的人.
\newline
\newline
保羅又說；「因主愛我們,就按著自己的旨意所喜悅的,豫定我們,藉著耶穌基督得兒子的名份；使他榮耀的恩典得著稱讚,這恩典是祂在愛子裡所賜給我們的.」(弗一5-6)我們現在正享受這恩典,將來要與主一同作王,要與眾天使一同稱讚祂的榮耀.
\newline
\newline


\newpage

\section{第47屆~港九培靈會~包忠傑~第五講}
\label{sec:1140}
\textbf{基督的寶血}
\newline
\newline
連結: \href{https://www.hkbibleconference.org/session-message/view/1140}{\texttt{https://www.hkbibleconference.org/session-message/view/1140}}
\newline
\newline
\hyperref[sec:1137]{< < < PREV SERMON < < <}
~
\hyperref[sec:toc]{[返目錄]}
~
\hyperref[sec:1142]{> > > NEXT SERMON > > >}
\newline
\newline
希伯來書 9:11-9:28
\newline
\begin{longtable}{cl}
\hline
\hline
章節 & 經文 (和合本修訂版)\\
\hline
9:11 & \begin{tabularx}{0.7\textwidth}{X} 但現在基督已經來到，作了已實現的美事的大祭司，經過那更大更全備的帳幕，不是人手所造，也不是屬於這世界的； \end{tabularx} \\ \\ \relax
9:12 & \begin{tabularx}{0.7\textwidth}{X} 他不用山羊和牛犢的血，而是用自己的血，只一次進入至聖所就獲得了永遠的贖罪。 \end{tabularx} \\ \\ \relax
9:13 & \begin{tabularx}{0.7\textwidth}{X} 若山羊和公牛的血，以及母牛犢的灰，灑在不潔的人身上，尚且使人成聖，身體潔淨， \end{tabularx} \\ \\ \relax
9:14 & \begin{tabularx}{0.7\textwidth}{X} 何況基督的血，他藉著永遠的靈把自己無瑕疵地獻給神，更能洗淨我們的良心，除去致死的行為，好事奉那位永生的神。 \end{tabularx} \\ \\ \relax
9:15 & \begin{tabularx}{0.7\textwidth}{X} 為此，基督作了新約的中保；因為他的死，贖了人在第一個約之時所犯的罪過，使蒙召的人能得著所應許永遠的產業。 \end{tabularx} \\ \\ \relax
9:16 & \begin{tabularx}{0.7\textwidth}{X} 凡有遺囑，必須證實立遺囑的人已經死了。 \end{tabularx} \\ \\ \relax
9:17 & \begin{tabularx}{0.7\textwidth}{X} 因為人死了，遺囑才有效力；立遺囑的人尚在，遺囑就不能生效。 \end{tabularx} \\ \\ \relax
9:18 & \begin{tabularx}{0.7\textwidth}{X} 所以，第一個約也是用血立的。 \end{tabularx} \\ \\ \relax
9:19 & \begin{tabularx}{0.7\textwidth}{X} 因為摩西當日照著律法將各樣誡命傳給眾百姓，就拿朱紅色絨和牛膝草，把牛犢、山羊的血和水灑在書上，又灑在眾百姓身上， \end{tabularx} \\ \\ \relax
9:20 & \begin{tabularx}{0.7\textwidth}{X} 說：「這血就是神與你們立約的憑據。」 \end{tabularx} \\ \\ \relax
9:21 & \begin{tabularx}{0.7\textwidth}{X} 他又照樣把血灑在帳幕和敬拜用的各樣器皿上。 \end{tabularx} \\ \\ \relax
9:22 & \begin{tabularx}{0.7\textwidth}{X} 按著律法，幾乎每樣東西都是用血潔淨的；沒有流血，就沒有赦罪。 \end{tabularx} \\ \\ \relax
9:23 & \begin{tabularx}{0.7\textwidth}{X} 這樣，照著天上樣式做的物件必須用這些禮儀去潔淨，但那天上的一切，自然當用更美的祭物去潔淨。 \end{tabularx} \\ \\ \relax
9:24 & \begin{tabularx}{0.7\textwidth}{X} 因為基督並沒有進了人手所造的聖所—這不過是真聖所的影像—而是進到天上，如今為我們出現在神面前。 \end{tabularx} \\ \\ \relax
9:25 & \begin{tabularx}{0.7\textwidth}{X} 他也無須多次將自己獻上，像大祭司每年帶著牛羊的血進入至聖所。 \end{tabularx} \\ \\ \relax
9:26 & \begin{tabularx}{0.7\textwidth}{X} 如果這樣，他從創世以來就必須多次受苦了。但如今，他在今世的末期顯現，僅一次把自己獻為祭，好除掉罪。 \end{tabularx} \\ \\ \relax
9:27 & \begin{tabularx}{0.7\textwidth}{X} 按著命定，人人都有一死，死後且有審判。 \end{tabularx} \\ \\ \relax
9:28 & \begin{tabularx}{0.7\textwidth}{X} 同樣，基督既然一次獻上，擔當了許多人的罪，將來要第二次顯現，與罪無關，而是為了拯救熱切等候他的人。 \end{tabularx} \\ \\
[1ex]
\hline
\hline
\end{longtable}
每一個基督徒,在自己的靈性上或讀經上,都應該認清楚神以甚麼為寶貴.許多人用自己的學說來講解聖經,並非按照真理來解說(如守望台,摩門教).我們要注意的是聖經以甚麼為貴重?在天父的心中,眼中,看祂兒子耶穌基督的寶血為最寶貴!
\newline
\newline
在舊約中,利未記與民數記都詳細記載流血獻祭的事,這些事都是預表耶穌基督,所以我們應耐心查看.神叫摩西到西乃,吩咐他造會幕,這會幕是天國真正的表徵；天上事情的影像,一切的建造按照天上的樣式,好比地上建造大廈的小模型,看模型就知道將來的情形一樣.
\newline
\newline
血是預表基督的寶血,舊約的祭司每一次為人獻祭,靠牛羊牲畜的血,這是指望將來有一位可以代替他們流血,使他們的罪可得赦免.
\newline
\newline
我們今天作基督徒的很有福氣,昔日的以色列人沒有聖經；現在神的啟示已經完全,主耶穌已經降生,受苦,被釘死在十字架上,第三日復活,四十日後升天,祂帶著自己的血到了天上的會幕去了!我們所相信的已圓滿了.
\newline
\newline
在舊約中,關於寶血的事是漸進的,頭一個用血獻祭的人是亞伯.希伯來書十一章四節說:「亞伯因著信獻祭與神,此該隱所獻的更美,因此便得了稱義的見證.......」亞伯藉著信心,用初生羔羊的血獻祭,就蒙神悅納.早期人類的歷史告訴我們,沒有血不能與神交往,不能到達神面前,享受神的恩惠.我們若要與神聯合,罪得赦免,就必須倚靠基督寶血代贖的功勞,寶血的真理就在此顯明了!
\newline
\newline
第二個是挪亞.洪水消落以後,人類已受了神的懲罰.挪亞從方舟出來,第一件事就是獻祭.
\newline
\newline
神試驗亞伯拉罕,要他帶著獨生子以撒,往摩利亞山上去獻為燔祭.以撒問亞伯拉罕:「燔祭的羊羔在那裡呢?」亞伯拉罕回答:「我兒,上帝必自己豫備作燔祭的羊羔.」這話十分寶貴,只要你肯獻上,上帝就有預備.因亞伯拉罕的心中,已經將以撒獻上了.
\newline
\newline
以色列人在埃及為奴受苦,神憐憫他們,要將他們救出來；神施行十大災難擊殺埃及人的時候,吩咐以色列人宰羊羔,把血塗在房屋的門楣和門框上,以色列人的長子得免於死.今天,我們因信主基督寶血的能力,心中已塗抹了基督的寶血,一切災禍就越過我們而去.
\newline
\newline
昔日以色列人在西乃,神叫摩西獻祭,要與摩西立約,是用祭牲的血.神與你我立約,是用基督的血,祂被釘在十字架上,所流出的寶血,就是立約的憑據.神為何要立新約呢?是因為舊約已經被人作廢了!現在世界上的國家,今天立約,明天作廢；但父神用自己獨生子基督的血所立的約,永遠不能作廢的.
\newline
\newline
亞伯拉罕在摩利亞山上獻了祭,代替了以撒,血流了之後,就有功效.西乃的血不但灑在以色列人的門楣與門框上,而是要灑在眾人的身上,表明與眾人有關係.今天,基督的寶血已灑在各人的心中,我們守聖餐就是要記念主所流的血.在舊約中看得很清楚,犯罪的人將手放在祭牲上,祭司殺這祭牲,代替犯罪的人受罰.今天,我們要在父神的面前認罪悔改,基督的寶血才有功效.假如我們得罪了神或得罪了人,就應該向神認罪,向人道歉,這是基督徒最大的見證.因為認了罪,心中才有平安,惟有悔改後,才能過得勝的生活.我們若要與魔鬼完全脫離關係,就一定要靠基督寶血得潔淨.
\newline
\newline
在舊約時代,以色列人的歷史中都可以看見,他們從埃及至西乃；從示羅至聖殿,都要用血,但聖殿拆毀以後,以色列人不知到那裡去獻祭了!現今猶太人守逾越節的時候,用羊腿放在桌上獻祭,已沒有祭司,也沒有用血了!所以猶太人離開了神的恩典了!但如今卻有許多青年猶太人信了基督,接受耶穌做救主,連他們的拉比都感到希奇.
\newline
\newline
耶穌基督如何論到血的問題呢?祂介紹自己是神羔羊.祂降生是為受死而來,惟有藉著祂的死,人類才有新生命.祂來,為要成就舊約一切關於流血的預表,祂在約翰福音六章五十至五十四節中,說到很希奇的話.祂說:「我是從天上降下來的糧；人若吃這糧,就必永遠活著；......我的肉真是可吃的,我的血真是可喝的.」基督的血已為我們流出,使我們信祂的人,罪得赦免!
\newline
\newline
神用一種血造全世界的人,但我們的生命已被魔鬼破壞,死在罪惡過犯之中.神的兒子來,給我們新的血,只要我們藉著信心,靠著祂流血之功,進入祂的生命裡面.
\newline
\newline
主基督很注重自己所流的血,在希伯來書中,有多次清楚地說到關乎基督寶血莫大的能力.基督進入了天堂；進入了真的會幕,祂現在做了我們的大祭司,藉著祂所獻上的寶血,我們得稱義.正如以弗所書一章七節說:「我們藉這愛子的血,得蒙救贖,過犯得以赦免,乃是照他豐富的恩典.」當魔鬼攻擊我們,叫我們靈性軟弱,信心動搖的時候,就用這節聖經對付它.讚美主,我們的神是信實的,公義的；主的寶血潔淨我們,祂能洗淨我們一切的不義.
\newline
\newline
基督的寶血將來有何價值呢?
\newline
\newline
基督復活後,祂的身體當然沒有血了,那麼,祂的血到那裡去呢?啟示錄曾多次提到羔羊寶座:「我又看見寶座與四活物並長老之中,有羔羊站立......四活物和四位長老,就俯伏在羔羊面前......大聲說,曾被殺的羔羊,是配得權柄,豐富,智慧,能力,尊貴,榮耀,頌讚的.......都歸給坐寶座的羔羊,直到永永遠遠.(五6,8,12)因著基督的血為我們打開天上的門,祂自己進了去,也把我們帶到天父面前.祂無罪之軀,卻因擔當世人的罪而上十字架.祂若不被釘流血,贖我們的罪,就連祂自己也不能進天堂去.祂曾在十字架上說:「我的神,我的神,為甚麼離棄我!」我們作基督徒的,要對付魔鬼,就要靠基督的寶血,因為祂的寶血有能力.
\newline
\newline
基督的寶血為我們成就了甚麼事?當然是洗淨我們的罪,又為我們打開墳墓的門.基督復活了,祂的墳墓打開,讓門徒和婦女可以親自看見復活的主.同時,基督的寶血將陰間的門打開,也把天上的門打開了!
\newline
\newline
基督的寶血在我們心中,既赦免我們的罪,又洗淨我們良心上的虧欠,叫我們心中爽快.而且藉著基督的寶血,敗壞了掌死權的魔鬼,魔鬼若攻擊我們,就用基督的寶血對付牠.當我們的「老我」出現的時候,也可以用基督的寶血去勝過!
\newline
\newline
撒迦利亞書十三章一節說:「那日必給大衛和耶路撒冷的居民,開一個泉源,洗除罪惡與污穢.」這是一句預言的話,開了一個泉源,川流不息.今天,無論那國那民,都可以享受基督血泉的功效,直到永遠.
\newline
\newline
基督所流的血,在父神眼中看為最寶貴,祂的寶血為我們打開天上的恩門；我們要倚靠主寶血的功勞,追求聖潔；也可以靠主的寶血使我們的祈禱,達到父神的面前.但願我們心中渴慕像祂,更認識基督寶血何等的寶貴,從而更愛這位為我們流血的基督.
\newline
\newline


\newpage

\section{第47屆~港九培靈會~包忠傑~第六講}
\label{sec:1142}
\textbf{基督的十字架}
\newline
\newline
連結: \href{https://www.hkbibleconference.org/session-message/view/1142}{\texttt{https://www.hkbibleconference.org/session-message/view/1142}}
\newline
\newline
\hyperref[sec:1140]{< < < PREV SERMON < < <}
~
\hyperref[sec:toc]{[返目錄]}
~
\hyperref[sec:1143]{> > > NEXT SERMON > > >}
\newline
\newline
路加福音 9:22-9:26
\newline
\begin{longtable}{cl}
\hline
\hline
章節 & 經文 (和合本修訂版)\\
\hline
9:22 & \begin{tabularx}{0.7\textwidth}{X} 又說：「人子必須受許多的苦，被長老、祭司長和文士棄絕，並且被殺，第三天復活。」 \end{tabularx} \\ \\ \relax
9:23 & \begin{tabularx}{0.7\textwidth}{X} 耶穌又對眾人說：「若有人要跟從我，就當捨己，天天背起自己的十字架來跟從我。 \end{tabularx} \\ \\ \relax
9:24 & \begin{tabularx}{0.7\textwidth}{X} 因為凡要救自己生命的，必喪失生命；凡為我喪失生命的，他必救自己的生命。 \end{tabularx} \\ \\ \relax
9:25 & \begin{tabularx}{0.7\textwidth}{X} 人就是賺得全世界，卻喪失了自己，或賠上自己，有甚麼益處呢？ \end{tabularx} \\ \\ \relax
9:26 & \begin{tabularx}{0.7\textwidth}{X} 凡把我和我的道當作可恥的，人子在自己的榮耀裡，和天父與聖天使的榮耀裡來臨的時候，也要把那人當作可恥的。 \end{tabularx} \\ \\
[1ex]
\hline
\hline
\end{longtable}
現代一般人的思想,都是人文主義的思想.他們以為神不過是人的思想造出來的,基督並無實際的存在.只要你思想祂,在你心中便有效驗；這是完全錯誤的,是中了人文思想的毒素.
\newline
\newline
許多人都不清楚十字架的重要和功效,主基督的目標是上各各他,祂就是為此來到世界.有些人以為,信耶穌只是叫我們做一個成功的人,使我們中心得快樂,這是錯的.當然,信耶穌可以心中得快樂,但神救贖我們的目的,是要在我們身上得榮耀.主耶穌上十字架的目的,也是榮耀天上的父.
\newline
\newline
十字架的真義是甚麼?十字架是最痛苦,最羞辱的刑罰.我們的主基督是神的羔羊,祂為擔當我們的罪上了十字架.「天路歷程」這本書甚能幫助我們走天路,基督徒走天路的時候背了大重擔,在窄路中進入了門,但到了十字架之時,只要仰望在十字架上的主,他身上的重擔就脫落,心中暢快,他深深地覺得因基督的愛得享安息,因祂的死而得到新生命.
\newline
\newline
基督釘十字架並非偶然而有,聖經說,他在創世之時已成為神的羔羊.所以,十字架有一個很長的陰影,這陰影從創世以前就有.在神與主之間的計劃中,第一步就是這樣成就救贖大功,所以基督定意成就天父所託付的聖工,上各各他釘十字架去!
\newline
\newline
苦行僧在受苦節的時候背十字架遊行,為的是要體驗主在十字架上的痛苦.
\newline
\newline
在聖經中多提到十字架的功效,還有一方面,在耶穌基督本身有很大的功效,就是基督提到自己心中所受的影響和關係.祂開始傳道的時候,魔鬼三次試探祂,目的是要基督不肯上十字架.魔鬼叫耶穌用石頭變餅,在肉體上試探祂；又帶主上聖殿的頂上,叫祂跳下,有天使用手托著......叫主在人面前得榮耀；魔鬼又說,你若拜我,我就把萬國榮華都賜給你.魔鬼已霸了世界,他是世界之王.末後主在客西馬尼園禱告,求父將這杯撤去.按照祂的肉體,祂是不想飲這杯的；因為祂是人,有人性的軟弱.但後來祂說:「然而不要照我的意思,只要照你的意思.」(參太廿六章)希伯來書五章八節說:「祂雖然為兒子,還是因所受的苦難學了順從.」主基督本是神的兒子,但祂成為人的樣式,在人的性情中受了試探,因受苦學了順從.主基督背負十字架,這十字架在祂心中成就功效.
\newline
\newline
我們作基督徒或聖工人員,都或多或少有自私的心；做傳道人要出名,要影響別人,不知不覺地便驕傲起來,敗壞了聖工.我們作基督徒或傳道人,都應該明白十字架的意義,就是不顧自己,對付老我.
\newline
\newline
主耶穌把上耶路撒冷受苦,上十字架的事預先告訴門徒,彼得說:「主阿,這是必不臨到你.」這話是出於魔鬼,故此主說:「撒旦,退到後邊去吧!」祂又說:「以後我不再和你們多說話,因為這世界的王將到,他在我裡面是毫無所有.」(約十四30)這是很寶貴的話,基督讓十字架的功效在祂裡面成功了.
\newline
\newline
我認識一位牧師,他很本事,又會寫作,但講道沒有亮光,所以教會不用他.為甚麼會這樣?就是驕傲,自私,貪財......這是隱而未現的罪,但十字架一來就要取銷這些.主耶穌說,世界之王在我裡面毫無所有.假如我們讓老我作王,靈性一定失敗.一個屬主的人,要認識清楚魔鬼的詭計.主對我們說,你們要跟從我,就當捨己,背起自己的十字架來跟從我.十字架在神的計劃中有很重要的價值,我們若要在靈性上得勝；就必定要治死老我,因為情慾的事是被咒詛的.
\newline
\newline
基督在十字架上所流的寶血甚有功效,它能對付一切世界的影響及邪情私慾.如果我們的老我尚未治死,既是未與主同釘十字架,就不能有得勝的生活.這些道理對未信主的人,或是屬肉體的基督徒,是認為愚拙的；在我們得救的人,卻為上帝的大能.走天路是越行越窄的,路是小的.我們若天天與主同釘十字架,就在生活上表現更加愛主；因為世界上的邪情私慾都都已釘死在十字架上了,自然就有美好的見證活出來.
\newline
\newline
這樣看來,基督的十字架有兩種功效:(一)基督被掛在十字架上替我們贖罪,這功效已做得十分完全.但在我們裡面,十字架的功效,能使我們逐漸成全.各人靈性的進步,有快有慢,如果我們愛主過於一切,十字架的功效就顯明了.我們若愛自己,靈性便是退步了!靈性的事要用屬靈的眼光去看,有些人外表雖然熱心；但內心沒有真正的用十字架去對付,也是沒有用的.因為聖靈不作工,十字架的功效,就不起作用!可見聖靈與十字架的功效,是十分重要的.
\newline
\newline
主耶穌雖然是無罪的,但祂尚且要與我們一樣受試探,受試煉,學習順從.他說,他自己分別為聖.祂本來就是聖潔的,祂在自己的肉體之中,還要對付魔鬼的試探和引誘.各位弟兄姊妹!聖靈作工在我們心中運行,叫我們看見自己的虧欠；好好對付,學習順從,將自己分別為聖.主基督不高舉自己,祂虛己,謙卑,把一切榮耀尊貴都倒空了.我們走天路的基督徒,豈能不虛己謙卑麼?
\newline
\newline
醫生若要知道一個人是否已經死亡,就用針刺他一下,看看他有無反應.我們作基督徒的,若要試試自己的老我有沒有治死；就要看看別人用話語刺激你的時候,心中反應如何?就可以知道了!假如別人用說話來破壞你的名譽,或是辱罵你,你是否要還口,報仇?主耶穌被打不還手,被辱罵不還口,這是十字架的功效,我們應該學習祂.神的兒子耶穌基督的寶血,十字架的功效要在我們身上產生能力,像主基督一樣.過去我們犯罪,叛逆,是受咒詛的,但基督已在十字架上流出了寶血,天父看我們的時候,因主基督的血遮蓋,已蒙潔淨了.
\newline
\newline
在舊約時代,有聖所及至聖所,中間有幔子.這幔子代表主的身體,為你,為我破碎；從上而下裂為兩半,使我們可以進到父神那裡去,這是何等寶貴的信息.聖經說到神的羔羊,在創世之前被殺,啟示錄十三章八節說:「凡住在地上,名字從創世以來,沒有記在被殺之羔羊生命冊上的人......」這節經文可譯為:「住在地上的名字,是沒有在被殺之羔羊生命冊上,......」就是神看來,因人犯罪而為人傷心難過,不但以前如此,以後也是如此；今天,父神也為世人之罪傷心.神為人類計劃了救恩,祂要付上痛苦的代價.主基督捨身流血這救恩已完全成功了；但祂仍嘆息,因為人仍然犯罪,仍然遭遇患難,痛苦.......彼前四章十二至十三節說:「親愛的弟兄阿,有火煉的試驗臨到你們,不要以為奇怪,倒要歡喜；因為你們是與基督一同受苦,使你們在他榮耀顯現的時候,也可以歡喜快樂.」我們是與基督一同受苦,我們要補滿基督患難的缺欠.神與祂的兒子一直嘆息,為祂的教會缺欠之緣故.所以我們要為教會興旺禱告,為教會中有不少情慾,失敗的事禱告.我們要與主一同受苦,與主同背十字架,讓十字架的功效成就在我們身上,使我們逐漸變成基督的樣式.
\newline
\newline
許多時候,人喜歡走捷徑,不喜歡背十字架,不願意受苦.怎樣才能愛這個十字架呢?首先就要有主的靈在我們心中,才能愛十字架,才能學習主的順從.我們若能為主的十字架受苦,在神面前才是金,銀,寶石的工作,十字架便將我們高舉,使我們更蒙神喜悅!「但願賜平安的神,就是那憑永約之血使群羊的大牧人我主耶穌,從死裡復活的神,在各樣善事上成全你們,叫你們遵行祂的旨意,又藉著耶穌基督在你們心裡行他所喜悅的事；願榮耀歸給他,直到永永遠遠.阿們.」(來十三20-21)
\newline
\newline


\newpage

\section{第47屆~港九培靈會~包忠傑~第七講}
\label{sec:1143}
\textbf{基督的復活與升天}
\newline
\newline
連結: \href{https://www.hkbibleconference.org/session-message/view/1143}{\texttt{https://www.hkbibleconference.org/session-message/view/1143}}
\newline
\newline
\hyperref[sec:1142]{< < < PREV SERMON < < <}
~
\hyperref[sec:toc]{[返目錄]}
~
\hyperref[sec:1144]{> > > NEXT SERMON > > >}
\newline
\newline
路加福音 24:13-24:46
\newline
\begin{longtable}{cl}
\hline
\hline
章節 & 經文 (和合本修訂版)\\
\hline
24:13 & \begin{tabularx}{0.7\textwidth}{X} 同一天，門徒中有兩個人往一個村子去；這村子名叫以馬忤斯，離耶路撒冷約有二十五里。 \end{tabularx} \\ \\ \relax
24:14 & \begin{tabularx}{0.7\textwidth}{X} 他們彼此談論所發生的這一切事。 \end{tabularx} \\ \\ \relax
24:15 & \begin{tabularx}{0.7\textwidth}{X} 正交談議論的時候，耶穌親自走近他們，和他們同行， \end{tabularx} \\ \\ \relax
24:16 & \begin{tabularx}{0.7\textwidth}{X} 可是他們的眼睛模糊了，沒認出他。 \end{tabularx} \\ \\ \relax
24:17 & \begin{tabularx}{0.7\textwidth}{X} 耶穌對他們說：「你們一邊走一邊談，彼此談論的是甚麼事呢？」他們就站住，臉上帶著愁容。 \end{tabularx} \\ \\ \relax
24:18 & \begin{tabularx}{0.7\textwidth}{X} 兩人中有一個名叫革流巴的回答：「你是在耶路撒冷的旅客中，惟一還不知道這幾天在那裡發生了甚麼事的人嗎？」 \end{tabularx} \\ \\ \relax
24:19 & \begin{tabularx}{0.7\textwidth}{X} 耶穌對他們說：「甚麼事呢？」他們對他說：「就是拿撒勒人耶穌的事。他是個先知，在神和眾百姓面前，說話行事都大有能力。 \end{tabularx} \\ \\ \relax
24:20 & \begin{tabularx}{0.7\textwidth}{X} 祭司長們和我們的官長竟把他解去，定了死罪，釘在十字架上。 \end{tabularx} \\ \\ \relax
24:21 & \begin{tabularx}{0.7\textwidth}{X} 但我們素來所盼望要救贖以色列民的就是他。不但如此，這些事發生到現在已經三天了。 \end{tabularx} \\ \\ \relax
24:22 & \begin{tabularx}{0.7\textwidth}{X} 還有，我們中間的幾個婦女使我們驚奇：她們清早去了墳墓， \end{tabularx} \\ \\ \relax
24:23 & \begin{tabularx}{0.7\textwidth}{X} 不見他的身體，就回來告訴我們，說她們看見了天使顯現，說他活了。 \end{tabularx} \\ \\ \relax
24:24 & \begin{tabularx}{0.7\textwidth}{X} 又有我們的幾個人往墳墓那裡去，所發現的正如婦女們所說的，只是沒有看見他。」 \end{tabularx} \\ \\ \relax
24:25 & \begin{tabularx}{0.7\textwidth}{X} 耶穌對他們說：「無知的人哪，先知所說的一切話，你們的心信得太遲鈍了。 \end{tabularx} \\ \\ \relax
24:26 & \begin{tabularx}{0.7\textwidth}{X} 基督不是必須受這些苦難，然後進入他的榮耀嗎？」 \end{tabularx} \\ \\ \relax
24:27 & \begin{tabularx}{0.7\textwidth}{X} 於是，他從摩西和眾先知起，凡經上所指著自己的話都給他們作了解釋。 \end{tabularx} \\ \\ \relax
24:28 & \begin{tabularx}{0.7\textwidth}{X} 他們走近所要去的村子，耶穌好像還要往前走， \end{tabularx} \\ \\ \relax
24:29 & \begin{tabularx}{0.7\textwidth}{X} 他們卻強留他說：「時候晚了，天快黑了，請你同我們住下吧。」耶穌就進去，要同他們住下。 \end{tabularx} \\ \\ \relax
24:30 & \begin{tabularx}{0.7\textwidth}{X} 坐下來和他們用餐的時候，耶穌拿起餅來，祝福了，擘開，遞給他們。 \end{tabularx} \\ \\ \relax
24:31 & \begin{tabularx}{0.7\textwidth}{X} 他們的眼睛開了，這才認出他來。耶穌卻從他們眼前消失了。 \end{tabularx} \\ \\ \relax
24:32 & \begin{tabularx}{0.7\textwidth}{X} 他們彼此說：「在路上他和我們說話，給我們講解聖經的時候，我們的心在我們裡面豈不是火熱的嗎？」 \end{tabularx} \\ \\ \relax
24:33 & \begin{tabularx}{0.7\textwidth}{X} 於是他們立刻起身，回耶路撒冷去，看見十一個使徒和與他們正在一起的人聚集在一處， \end{tabularx} \\ \\ \relax
24:34 & \begin{tabularx}{0.7\textwidth}{X} 說：「主果然復活了，已經顯現給西門看了。」 \end{tabularx} \\ \\ \relax
24:35 & \begin{tabularx}{0.7\textwidth}{X} 於是，兩個人把路上所遇到，和耶穌擘餅的時候怎麼被他們認出來的事，都述說了一遍。 \end{tabularx} \\ \\ \relax
24:36 & \begin{tabularx}{0.7\textwidth}{X} 正說這些話的時候，耶穌親自站在他們當中，說：「願你們平安！」 \end{tabularx} \\ \\ \relax
24:37 & \begin{tabularx}{0.7\textwidth}{X} 他們卻驚慌害怕，以為所看見的是魂。 \end{tabularx} \\ \\ \relax
24:38 & \begin{tabularx}{0.7\textwidth}{X} 耶穌對他們說：「你們為甚麼驚恐不安？為甚麼心裡起疑惑呢？ \end{tabularx} \\ \\ \relax
24:39 & \begin{tabularx}{0.7\textwidth}{X} 你們看我的手和我的腳，就知道實在是我了。摸摸我，看，因為魂無骨無肉，你們看，我是有的。」 \end{tabularx} \\ \\ \relax
24:40 & \begin{tabularx}{0.7\textwidth}{X} 說了這話，他就把手和腳給他們看。 \end{tabularx} \\ \\ \relax
24:41 & \begin{tabularx}{0.7\textwidth}{X} 他們還在又驚又喜、不敢相信的時候，耶穌對他們說：「你們這裡有甚麼吃的沒有？」 \end{tabularx} \\ \\ \relax
24:42 & \begin{tabularx}{0.7\textwidth}{X} 他們給了他一片烤魚， \end{tabularx} \\ \\ \relax
24:43 & \begin{tabularx}{0.7\textwidth}{X} 他接過來，在他們面前吃了。 \end{tabularx} \\ \\ \relax
24:44 & \begin{tabularx}{0.7\textwidth}{X} 耶穌對他們說：「這就是我從前和你們同在時所告訴你們的話：摩西的律法、先知的書，和《詩篇》上所記一切指著我的話都必須應驗。」 \end{tabularx} \\ \\ \relax
24:45 & \begin{tabularx}{0.7\textwidth}{X} 於是耶穌開他們的心竅，使他們能明白聖經， \end{tabularx} \\ \\ \relax
24:46 & \begin{tabularx}{0.7\textwidth}{X} 又對他們說：「照經上所寫的，基督必受害，第三天從死人中復活， \end{tabularx} \\ \\
[1ex]
\hline
\hline
\end{longtable}
關於基督復活,升天的事,在剛才所讀的經文中,門徒與婦女們都看不清楚,關於基督復活的事,主基督雖然曾對他們說過,但他們仍不明白.我們有時候也會被魔鬼迷惑,不明白聖經真理.人對於屬靈的事物,生來就瞎眼與愚昧的,因為撒旦弄瞎人的心眼,但等到聖靈光照我們的時候,才明白過來,好像在以馬忤斯路上與主同行的門徒一樣.所以每一個基督徒,都應該明白基督復活的事.
\newline
\newline
路加廿四章說到當時在以馬忤斯路上,基督與門徒講論自己的事情,他們被油蒙了心眼.以色列人不信主基督,就是他我們盼望的彌賽亞,主就向他們顯現.
\newline
\newline
第一個看見復活之基督的人,是抹大拉的馬利亞；主曾為她趕出七個鬼,她有一個報恩的心,常與主親近.主釘死在十字架後,她曾為主安葬的事盡了本份,她到墳墓中尋找主.她看見墳墓的石頭挪開了,屍體卻不在其中,以為主被人挪去了!主說:「為何在死人中找活人呢?」主基督所擔心的就是這一點,所以主一次又一次地向他們顯現.有一次,主顯現時對他們說:「你們要接受聖靈.」又有一次,主向十個門徒顯現,多馬不在其中；以後再一次特別為多馬顯現,為要多給他一次機會.
\newline
\newline
主向門徒解釋,祂如何受害,復活,都應驗了聖經的預言,門徒這才明白,直到他們進入了一座樓房,在擘餅的時候,才認出是主.
\newline
\newline
各位,我們在守聖餐的時候,若存著渴慕主的心；主必開我們的眼睛,讓我們看見祂的同在.
\newline
\newline
主特別關心的,就是三次不認祂的彼得.當我們受試探的時候,也許會像彼得一樣不認主,但主仍與我們接近.有人曾在試煉受苦時不認主,後來他流淚痛哭,自認是猶大,其實他也是彼得.主特別尋找跌倒的彼得,問彼得是否愛祂.祂特別關心彼得,忍耐他三次的不承認；但也給他三次機會,讓他再承認祂.當我們軟弱的時候,魔鬼便叫我們不認主,要我們灰心失望；但主愛我們,祂不願意見我們失落,祂來尋找我們.
\newline
\newline
主復活之後,仍關心門徒,只是門徒以為主已釘死,灰心欲絕,便打魚去了!誰知整夜勞力卻一無所得.所以主在海邊為他們預備了早餐,叫他們來同吃.照樣,主也為我們預備日用的飲食,祂雖然在天上,卻極關心我們.但願聖靈開我們的心眼,叫我們明白主的心.當我們失敗跌倒,或遇危險時,主就來幫助我們,並且眷顧我們,為我們預備所需要的一切.
\newline
\newline
我這次旅行,旅費不足,只好在禱告中交託主,直至在上機前一天,才籌足款購機票.走天路的人,常常都會遇到試煉,主要藉著各樣試煉成全我們的信心,好盡力為祂工作.
\newline
\newline
主耶穌復活後,雖然門徒有時見到主,有時又看不見主.主卻關心他們,愛他們,不但向他們顯現,也要他們認識祂,知道祂,並倚靠祂.
\newline
\newline
主升天之後,坐在父神的右邊,以弗所書一章廿二節說:「......萬有伏在祂的腳下,使祂為教會作萬有之首.」祂是教會的頭,指揮教會,對教會負責,這是祂在父神右邊主要的責任.希伯來書一次又一次地,提到主是大祭司,祂為我們禱告；因為我們的禱告不夠完善,所以需要祂的代求.啟示錄八章三至四節說:「有一位天使拿著金香爐,來站在祭壇旁邊；有許香賜給他,要和眾聖徒的祈禱,一同獻在寶座前的金壇上.那香的煙,和眾聖徒的祈禱,從天使的手中一同升到神面前.」這天使就是指主基督.天使有許多香,他將香放在禱告上一同獻上給神,天父聞到兒子獻上的香氣,就必垂聽禱告!因為主已將我們的祈禱,變成父神喜悅之香氣.
\newline
\newline
提前二章五節說:「因為只有一位神,在神和人中間只有一位中保,乃是降世為人的基督耶穌.」現在主基督正做中保的工作.我們若要上天堂,就需要一位中保擔保我們,這位中保,就是神的兒子耶穌基督.
\newline
\newline
主復活之後,為我們成就了甚麼事?羅馬書四章廿五節說:「耶穌被交給人,是為我們的過犯；復活,是為我們稱義.」祂為我們成就了寶貴的事,祂的血能為我們贖罪,祂復活是叫我們稱義.如果主沒有復活,又升到父神那裡去,祂的救贖就不完全.因為祂復活戰勝了死亡,才到父神那裡去.基督是完全無罪的,祂擔當了我們一切的罪,我們靠著祂；祂復活,在父神面前代求,所以救恩便成功了!
\newline
\newline
主復活,證明祂是神的兒子,同時證明我們也要復活,林前十五章十四至廿三節的經文足以證明.基督復活成了頭一個復活者的先鋒.祂作了模樣,打開了復活的門,帶我們經過死亡然後復活.
\newline
\newline
基督復活之後,門徒心中有了很大的改變,過去他們看見主死了,而且埋葬了,他們成了絕望的人；悶悶不樂,使隨彼得打魚去了!但他們看見復活的主向他們顯現,在以馬忤斯路上,和擘餅的時候,......便心裡火熱,不能不到耶路撒冷報信去；從此,憂愁變成喜樂,失望變成盼望,這是復活的主在門徒身上的果效.
\newline
\newline
主在臨終前,曾對門徒說過,祂去,原是為門徒預備地方去,並且還要差遣聖靈來!主在升天時吩咐門徒在耶路撒冷,等候聖靈降臨,就必得著能力.當時門徒還不明白是甚麼一回事,直到至五旬節到了,聖靈就降臨下來,充滿他們,叫他們有能力,有膽量為主作見證,傳揚福音,甚至肯為主殉命!
\newline
\newline
復活的主也要向你,向我顯現,叫我們認識祂,聽祂的話,相信祂!保羅作見證說:「我們傳揚祂,是用諸般的智慧,勸戒各人,教導各人；要把各人在基督裡完完全全的引到神面前.我也為此勞苦,照著祂在我裡面運行的大能,盡心竭力.」(西一28)各位,一個人蒙恩得救,是神用大能從魔鬼手中奪過來的；如果基督沒有復活,就沒有人能從魔鬼手中將我們奪回來.祂復活的能力要住在我們裡面,使用我們,藉著我們的身體為祂活在這世界.我們不要疏忽蒙恩的事情,如果我們真正為罪痛悔,主復活的能力,就會把我們改變過來.主的靈在我們裡面,使我們活著的就不再是自己,乃是基督在我們裡面活著了.
\newline
\newline
以色列人出埃及是一件大事,是神用大能大力救他們脫離埃及王的權勢；我們蒙恩的經驗也是一件大事,是基督用復活的大能,將我們拯救出來.我們有了這經驗,才有勇氣奔走天路.
\newline
\newline
主復活後便升上高天,祂與門徒離別之時,最後囑咐說:「......天上地下所有的權柄,都賜給我了.所以你們要去,使萬民作我的門徒,奉父子聖靈的名給他們施洗.凡我所吩咐你們的,都教訓他們遵守；我就常與你們同在,直到世界的末了.」(太廿八18-20)
\newline
\newline


\newpage

\section{第47屆~港九培靈會~包忠傑~第八講}
\label{sec:1144}
\textbf{再來之王}
\newline
\newline
連結: \href{https://www.hkbibleconference.org/session-message/view/1144}{\texttt{https://www.hkbibleconference.org/session-message/view/1144}}
\newline
\newline
\hyperref[sec:1143]{< < < PREV SERMON < < <}
~
\hyperref[sec:toc]{[返目錄]}
~
\hyperref[sec:1161]{> > > NEXT SERMON > > >}
\newline
\newline
馬可福音 11:1-11:10
\newline
\begin{longtable}{cl}
\hline
\hline
章節 & 經文 (和合本修訂版)\\
\hline
11:1 & \begin{tabularx}{0.7\textwidth}{X} 耶穌和門徒快到耶路撒冷，來到伯法其和伯大尼，在橄欖山那裡。耶穌打發兩個門徒， \end{tabularx} \\ \\ \relax
11:2 & \begin{tabularx}{0.7\textwidth}{X} 對他們說：「你們往對面村子裡去，一進去的時候會看見一匹驢駒拴在那裡，是從來沒有人騎過的，把牠解開，牽來。 \end{tabularx} \\ \\ \relax
11:3 & \begin{tabularx}{0.7\textwidth}{X} 若有人對你們說：『為甚麼做這事？』你們就說：『主要用牠，但會立刻把牠牽回到這裡來。』」 \end{tabularx} \\ \\ \relax
11:4 & \begin{tabularx}{0.7\textwidth}{X} 他們去了，看見一匹驢駒拴在門外街道上，就把牠解開。 \end{tabularx} \\ \\ \relax
11:5 & \begin{tabularx}{0.7\textwidth}{X} 在那裡站著的人，有幾個說：「你們解開驢駒做甚麼？」 \end{tabularx} \\ \\ \relax
11:6 & \begin{tabularx}{0.7\textwidth}{X} 門徒照著耶穌的話說，那些人就任憑他們牽去了。 \end{tabularx} \\ \\ \relax
11:7 & \begin{tabularx}{0.7\textwidth}{X} 他們把驢駒牽到耶穌那裡，把自己的衣服搭在上面，耶穌就騎上。 \end{tabularx} \\ \\ \relax
11:8 & \begin{tabularx}{0.7\textwidth}{X} 有許多人把衣服鋪在路上，還有人把田間的樹枝砍下來鋪上。 \end{tabularx} \\ \\ \relax
11:9 & \begin{tabularx}{0.7\textwidth}{X} 前呼後擁的人都喊著說：「和散那！奉主名來的是應當稱頌的！ \end{tabularx} \\ \\ \relax
11:10 & \begin{tabularx}{0.7\textwidth}{X} 那將要來的我祖大衛之國是應當稱頌的！至高無上的，和散那！」 \end{tabularx} \\ \\
[1ex]
\hline
\hline
\end{longtable}
經文:可十一1-10　啟十九11-16
\newline
\newline
在世界的國度中,必有人作主掌權.但真正配作王掌權的,只有一位,就是我們所盼望的主耶穌基督,祂是萬主之主,萬王之王.
\newline
\newline
神作事常常是雙重的.主耶穌第一次騎驢進入耶路撒冷,沒有威風,也沒有權柄；只是謙卑地進入耶路撒冷,因為那時候祂的使命是作神的羔羊.在啟示錄十九章十一至十六節的經文中,說到主基督騎著一匹白馬,眼睛如火燄,他頭上戴著許多冠冕；......在祂衣服的大腿上,有名寫著萬王之王,萬主之主.一切仇敵祂已勝過,祂要進到榮耀裡.基督復活之後,在以馬忤斯路上向兩個門徒顯現；祂告訴門徒,祂必須先受苦,將來要進入榮耀裡.
\newline
\newline
一,祂要震動世上的國
\newline
\newline
哈該書二章六節說:「萬軍之耶和華如此說,過不多時我必再一次震動天地,滄海,與旱地；我必震動萬國；萬國所羡慕的必來到,我就使這殿滿了榮耀.......」希伯來書十二章廿五至廿八節又說:「你們總要謹慎,不可棄絕那向你們說話的.因為那些棄絕在地上警戒他們的,尚且不能逃罪,何況我們違背那從天上警戒我們的呢?當時祂的聲音震動了地；但如今祂應許說,再一次我不單要震動地,還要震動天.這再一次的話,是指明被震動的,就是受造之物,都要挪去,使那不被震動的常存.所以我們既得了不能震動的國,就當感恩；照上帝所喜悅的,用虔誠敬畏的心事奉神.因為我們的神乃是烈火.」
\newline
\newline
過去,神震動了地,在聖經中曾有多次提到地震.將來神要震動天.為甚麼祂要震動這個宇宙呢?那不被震動的,常存的又是甚麼呢?在主再來之前,世界上有個兆頭,就是有許多國家要動搖,有些國家興起；有些國家敗亡,他們實在痛苦.神不但震動地,也震動世界的政治.在舊約記載,撒母耳領導以色列人打仗的時候,他禱告,地就震了!詩篇十八篇六至七節也說:「我在急難中求求耶和華,向我的上帝呼求；祂從殿中聽了我的聲音,我在祂面前入了祂的耳中.那時因祂發怒,地就搖撼戰抖,山的根基也震動搖撼.」不只大衛的禱告能震動地,許多聖徒的禱告也能震動地.故此聖經中說到聖徒禱告,就是有大事出現或地震.主耶穌被釘在十字架上的時候,地也震動,神正搖撼這個世界.根據科學家的統計,顯示出地震的數目正不住增加.
\newline
\newline
第一世紀──十五次.
\newline
\newline
第六世紀──三十次.
\newline
\newline
十四世紀──一百三十七次.
\newline
\newline
十五世紀──一百七十四次.
\newline
\newline
十六世紀──二百五十三次.
\newline
\newline
十七世紀──三百七十八次.
\newline
\newline
十八世紀──六百四十次.
\newline
\newline
十九世紀──二千一百一十六次.
\newline
\newline
有誰知道二十世紀會有多少次地震呢?
\newline
\newline
二,祂對祂的「小羊」說:「不要懼怕!」
\newline
\newline
從以上的數字看來,世界越來越搖撼了!聖經也告訴我們,因人犯罪,地就震動.神說,我現在搖動這個世界,將來天地都要搖動.這是神發怒的時候.地震實在使人害怕,有一次,我在美國加州講道,正放映幻燈片,忽然地震,眾人都很害怕.
\newline
\newline
在末日的時候,神要震動整個宇宙,目的是要顯露那不可動搖的,因為祂的國家不能搖動.那時,列國都驚慌.但我們信主的人,就不必懼怕.因為當主離開世界的時候,對我們說:「你們這小群,不要懼怕,因為你們的父,樂意把國賜給你們.」(路十二32)這是主永遠的國度.在聖經中多次提到不要怕,舊約中有三十多次；新約中超過十次,主耶穌也曾多次叫門徒不要怕.
\newline
\newline
記得有一位中國同工,講過一個真實的故事:一天,一位牧師坐船到別處佈道,船篷被風吹動；牧師的傭人很害怕,他想跳上岸去,但跳不到岸,返而跌落水中浸死了!今天,許多人在教會的船中,教會遇到風浪時,便想跳上去,他恐怕這隻船會沉.老牧師加上了一句話:「你跳上去恐怕有危險,只管安靜在船中好了.」現今世界罪惡多,風浪大,你只要在主的船中,就不必懼怕,因為祂對我們負責.
\newline
\newline
三,祂說:「天國的福音要傳遍天下.」
\newline
\newline
耶穌曾對門徒說:「這天國的福音要傳遍天下,對萬民見證,然後末期才來到.」(太廿四14)在末日的時候,神的計劃中,要將天國福音傳遍天下,祂要我們去傳,因為祂快再來了!時候不多了,我們要竭力工作；趕緊工作,恩惠的日子快要過去了!
\newline
\newline
現在正是榮耀的宣教時期,馬禮遜在一八○七年到廣州傳道,至今已經有一百七十多年了!神給中國人許多機會歸信祂.
\newline
\newline
我曾與一個中國醫生談話,他說,我用了十多年的時間學醫科,花了半生心血,才能成為一個醫生；但是,主若要用我的話,我願意放下醫務去傳道.真是感謝讚美主.現今神正在青年人當中作工,近年來許多青年人被神呼召,甘心樂意奉獻自己；他們以得著耶穌基督為至寶,因他們知道主再來的日子近了!
\newline
\newline
四,祂說:「還有一點點時候.」
\newline
\newline
主對門徒說:「還有一點點時候.」正如詩篇三十七篇十節說:「還有片時,惡人要歸於無有.......」我們舉目觀看,現在到處都有惡人；惡人實在太多了,信主的人太少了!我們要知道,我們有責任將這些人從罪惡中搶出來,因為還有片刻,惡人要歸於無有.
\newline
\newline
主耶穌說:「等不多時,你們就不得見我；再等不多時,你們還要見我.」(約十六16)等不多時,即一點點時候.一生如作夢,忽然過去,但我們的主再來時,是永遠的時刻；世上至暫至輕的苦楚總要過去的,黑夜已深,白晝將近,我們的主快再來了!
\newline
\newline
希伯來書十章卅七節又說:「還有一點點時候,那要來的就來,並不遲延.」這句話非常寶貴,或在今日,或在明日,主必快來!新約中有個最後的禱告:「證明這事的說,是了；我必快來.阿們!主耶穌阿,我願你來.」(啟廿二20)雖然世界上有許多苦難,動盪不安,但我們可以放心,因為主快再來,把我們帶到榮耀裡,將來我們要與祂同享福氣.但是今天,我們要儆醒等候,作好準備,常常禱告,多作主工,保守自己不沾染世界的污穢......迎接主再來.但願我們都從心裡說:「主耶穌阿,我願你來!阿們.」
\newline
\newline


\newpage

\section{第47屆~港九培靈會~滕近輝~第一講}
\label{sec:1161}
\textbf{聖靈好像火一樣}
\newline
\newline
連結: \href{https://www.hkbibleconference.org/session-message/view/1161}{\texttt{https://www.hkbibleconference.org/session-message/view/1161}}
\newline
\newline
\hyperref[sec:1144]{< < < PREV SERMON < < <}
~
\hyperref[sec:toc]{[返目錄]}
~
\hyperref[sec:1163]{> > > NEXT SERMON > > >}
\newline
\newline
感謝主的恩典!使我們每年在培靈會中,有十天的時間,彼此分享主的語語.兄弟戰戰兢兢來傳達主的信息,非常需要眾弟兄姊妹的代禱.
\newline
\newline
回想兄弟第一次在培靈會講道,遠在廿四年前,時光易逝,今昔相比,深感如今更蒙主的恩典.據去年與前年晚堂赴會人數,較廿四年前增加數倍,今年如何,尚未知道,還是要仰望主.
\newline
\newline
與港九培靈研經會同樣性質的聚會,在其他城市中,近年來也逐漸舉行了.兩年前在台北的培靈研經會,舉行了八天,最後一晚赴會人數有三千之眾.在星加坡,培靈會已舉行了八年,(分中文和英文的).中文的每年四月舉行,赴會人數達千人；英文的每年六月舉行,兩年來聚會人數高達三千三百.因為星加坡國家戲劇院只能坐滿三千三百人,其餘在外面不得進入的達數百人.此外吉隆坡,檳城,也都已開始了每年一度的培靈會.可以說香港的培靈研經大會對很多地方都有影響.英國的凱斯可大會百周年紀念於今年七月舉行,赴會人數達萬人,今年的人數特別多,不過以前也有七千人.這種培靈大會運動已遍及全世界.我們要求神,在前面的日子,使更多地區,興起同樣的聚會來.
\newline
\newline
神工作的方程式
\newline
\newline
我深深相信今天基督眾教會所需要的是什麼,就是聖靈的作為.神的工作,有個方程式:「聖靈的能力,加美好的靈性,再加上方法.」單有方法是不夠的,必須有美好的靈性；但還不夠,必須加上聖靈的能力.有了這三樣,神所要的條件就具備了,神就能完成祂最好的旨意.
\newline
\newline
但願我們在這十天晚上,用謙卑安靜的心在主的話語中查考領受,接受這個與我們每個人有密切關係的真理.
\newline
\newline
我們的主對於聖靈的注重,是無以復加的.在使徒行傳中,主耶穌對門徒說:「聖靈降臨在你們身上,你們就必得著能力.在耶路撒冷,猶太全地,撒瑪利亞,直到地極,作我的見證.」主耶穌對軟弱的門徒感到憂慮,因為他們缺乏信心.最大的使徒也三次不認主,祂對門徒實在很不放心,然而對聖靈的能力卻滿有信心,所以祂說:「聖靈降臨在你們身上,你們就必得著能力,從耶路撒冷開始,直到地極作祂的見證.」主耶穌在(約十四12)說了一句話,使我們乍聽時幾乎難以相信和接受.祂說:「我往父那裡去.」意思就是說祂必須往父那裡去,然後聖靈保惠師才來到世上.使門徒改變,得著能力,成了截然不同的人物.此後,他們就能作,比耶穌在世時,所作更大的事.因耶穌在世時,受了肉身限制,聖靈卻不受時間與空間限制的.當千千萬萬門徒被聖靈充滿時,他們可把福音,傳遍全世界,直到地極主的話是真實的,我們可以得著它,可以經歷它.
\newline
\newline
在這十個晚上,我要從兩個角度來講:
\newline
\newline
第一,在聖經中用十五個比喻來形容聖靈的工作.在這十五個比喻中,讓我們看見聖靈的能力,作為,和充滿.
\newline
\newline
第二,最後的晚上,我要特別講聖靈的充滿.
\newline
\newline
在聖經中多處可以看見,第一用以形容聖靈的比喻就是火.「聖火」,就在聖靈的火.聖火的異象都與聖靈的工作有關.
\newline
\newline
六個靈火的異象
\newline
\newline
第一個靈火的異象,這個異象其實是一對,(結一27-28)先知以西結看見耶和華榮耀的異象,這異象分為二部分,上面的部分中間是發光的真金,周圍是火.下面的部分,中間是燃燒的火,周圍是彩虹.這個異象,無論上下部分都有火.火是象徵耶和華的聖潔和公義,是公義聖潔的火；這火非常重要,這火也應燃燒在每個屬於耶和華的人之心中.這使我們連想到(賽四4)的話說:「主以公義的靈和焚燒的靈將錫安女子的污穢洗去.」
\newline
\newline
焚燒的靈是指聖靈的火,這火有潔淨的作用,神的聖潔像火般的燃燒.這火在神的本性裡燃燒,也在每個屬神的人重新的性格裡燃燒.
\newline
\newline
有一次我讀啟示錄,得了很大的幫助,發現四章三節用璧玉形容耶和華.二十一章11節又說新耶路撒冷城的榮耀好像璧玉.由此可知璧玉是極重要的象徵,璧玉象徵神的聖潔.神本性聖潔,凡與神有關的一切都是聖潔的,事奉神的人也必須要聖潔.
\newline
\newline
在座的眾弟兄姊妹們!今世代撒旦在那裡燃燒牠的火.這是戰火,牠到處引起戰火,魔鬼所點的是慾火,是情慾之火；慾火正在全世界燃燒.我們生存在這時代中,我們裡面必須有聖靈公義的燃燒,才能勝過外面的火.求主把聖火在我們裡面燃點,好叫我們在這充滿情慾的世代,過著聖潔的生活.可惜不少基督徒,被時代的慾火點燒了；犯了情慾的罪,失去了愛主的心,繼而失去了熱心；漸漸冷淡,退後,跌倒,這是何等可惜!當我們看見第一種靈火時,我們是覺得有莫大的需要.巴不得勝靈的火──焚燒的靈,在教會燃燒起來,把我們生活中一切的污穢焚燒淨盡.
\newline
\newline
這一對異象的第二部分在以賽亞書六章6節.先知以賽亞在聖殿,看見耶和華在寶座上；有一撒拉弗正在祭壇上燃燒的炭火,沾他的口.說:「看哪!這炭沾了他的嘴,你的罪孽便除掉,你的罪惡就赦免了.」那塊炭火是從祭壇上取來的,祭壇代表兩意思:一是贖罪的祭壇,一是奉獻的祭壇.贖罪的祭壇有潔淨的作用,耶穌基督救贖的寶血把我們的罪潔淨了.奉獻的祭壇也有潔淨作用,有的人不能勝過某些罪,但當他奉獻以後；他就得到一種新的力量,來勝過他原來所勝不過的罪.
\newline
\newline
莫dik 先生常帶著兩面放大鏡在衣袋裡,一面是完整的,另一邊的口袋,是放了一面破開成兩半的.他和人談話或講道時,都用這兩面鏡子.他說:「當陽光透過放大鏡,立刻出現一個焦點,焦點熱度很高可引起燃燒；但破碎的放大鏡,再無法造成焦點,因它的作用和熱力都失去了.」他用這比喻來說明奉獻的真理.
\newline
\newline
神透過我們的奉獻,造成一種能力,這能力含著熱度焦點.由此我們聯想耶穌在(約十五2)所說的比喻:若要枝子結果更多,必須修理乾淨.意思是把雜枝剪去,免得消耗葡萄樹上生命的汁漿而影響結果子,主耶穌所題的剪去就是潔淨的意思.奉獻能夠造成這種潔淨.有些平常的事情,本身不算為罪；但如果這些事攔阻我們愛神的心,與我們為主而活的基本原則不合；攔阻我們完全為主而活,消耗我們生命的力量,以致少結果子；凡此都屬不潔淨的,都是一個奉獻的基督徒人生的污點.如果我們奉獻了自己,就必須把這些都除去,讓我們得到潔淨.
\newline
\newline
第二靈火的異象(出三2-3)摩西在何烈山上看見這異象.荊棘被火燒著卻沒有燒毀,這是什麼意思?荊棘代表試煉,以色列人在埃及受到如火的試煉；但神在如火的試煉中保守他們,使他們安然渡過一切.這意義有三方面的應驗.(1)應驗在當時以色列人的處境之中,他們在埃及受法老的壓制,受很多的痛苦,好像在火中經過一樣,但耶和華保守他們.(2)應驗在摩西本人身,摩西被神呼召後,負著很重的擔子,要帶三百萬人經曠野往應許之地去.在環境上有許多困難,況且這些人是悖逆的以色列民；摩西肩上的擔非常沉重,甚至有一次他求死；(民十一)因他覺得重擔難挑,但神保守他到底,使他作成神所要他作的工.
\newline
\newline
我順便一題:聖經中有兩人求死:一是以利亞在蘿藤樹下求死.他說:「耶和華阿!罷了,求你取我的性命,因我不勝於我的列祖.」還有一個是摩西；這二人原是最堅強的人,但是他們卻求死.神給他們的答案和指示是相同的,神對摩西說:「你要擇選七十人,訓練他做以色列人的領袖.」神對以利亞說:「你回去,找以利沙,膏立他,訓練他,作你的繼承人.」後來以利沙開了先知學校,訓練了許多的門徒.今天是個訓練的時代,神告訴我們,教會發展增長的秘訣在於訓練人才.(3)應驗在以色列人將來的歷史上.兩千年之中,以色列人經過無數的試煉,神保守他們,能死裡逃生；最後使他們復國,讓神的美意成就.
\newline
\newline
火,另方面的意義是代表能力,神把這能力給摩西,摩西才能持續長久為神燃燒.如有神的能力在我們裡面,我們就能持久熱心,因來源是屬於神.願我們靠著聖靈的能力,再接再厲,為祂燃燒!.
\newline
\newline
在聖經中有四位老人,真是寶貝.
\newline
\newline
(一)亞伯拉罕　他愛主的高峰約在他一一七歲時達到的,(聖經沒有記載他一百十七歲時獻以撒)他把獨生愛子獻給神,那時他已年老了,但他愛神的心達到高峰,他裡面愛神的力量堅持不移.
\newline
\newline
(二)迦勒　(書十四6-11)告訴我們,迦勒自己作見證說:神當日打發我去窺探迦南地時是四十歲,現在我已八十五歲了；那時我的力量如何,今日我的力量仍然如何.
\newline
\newline
(三)但以理　他九十歲時,波斯王下禁令,不准向神或任何人求拜.但以理卻敬畏耶和華,向神忠心,有膽量,充滿信心,每天仍然打開窗戶朝向耶路撒冷,三次跪下禱告,與素常一樣.
\newline
\newline
(四)使徒約翰　他約九十歲時,被充軍到拔摩海島,為主受很多苦.後來被釋放回到以弗所,他寫了三封書信；字裡行間,充滿著愛......喜樂......的詞司.試默想,這位九十歲的老人家,曾為主受大苦,剛被釋放回來,卻滿口是愛和喜樂......真是語言難以形容的!因為神的愛在他裡面繼續燃燒,今天我們需要這力量來充滿我們!時間不能磨滅我們愛主的心,試煉也不能減少我們的愛主的心.
\newline
\newline
前幾年我經過羅省,有一位青年把一封信交給我看.原來是他父親由大陸,經香港轉到美國給他的,內容是勉勵他要愛主.試想,一個在宗教壓制下的人,非但能保守自己愛主的心；而且還寫信給他在美國的兒子,勉勵他要愛主.請問:在有壓力之下或是在自由的環境下愛主,何者較為容易呢?正確的答案當然是:「愛主不在乎環境.」有人在好的環境中冷淡了,有人在神的愛中跌倒了,但卻有人在大壓力之下能熱心愛主.愛主,是我們心裡的事,不是環境的事；我們可在任何環境中愛主.求主給我們這樣的火!叫我們愛神愛人.
\newline
\newline
第三個靈火的異象(結一13)這異象是事奉的異象,有四活物和基路伯,並有火上去下來,如燒著火炭的形狀,又如火把的形狀.四活物代表理想的事奉,異象中有火,乃是理想的事奉,所以我們要心裡火熱常常事奉主.
\newline
\newline
聖經裡有兩種特別的天使,即撒拉弗,與基路伯.撒拉弗的字根是火的意思.基路伯是知識之意.這兩種事奉的天使,神既給他們命名「火」「知識」,可見在神心目中,理想的事奉,火和知識,二者必兼而有之,無一可以缺少.有火而無知識,易變為狂熱；無知識的狂熱,極為危險!反言之,有知識而沒有火,知識也起不了作用.我曾讀(啟四6-8),發現基路伯和撒拉弗聚在一起,六節講到四活物的形狀,也就是基路伯的形狀.七節論基路伯的四個面.八節論撒拉弗的形狀,(以賽亞六章的異象裡也有這形像)有六個翅膀,他們晝夜不住的說「聖哉聖哉聖哉」,這也是撒拉弗的呼聲；所以我們看見撒拉弗基路伯和撒拉弗連在一起了.如果我問你這是基路伯或是撒拉弗?你一定回答不出,因為二者有同樣的特點.這裡給我們看見,對神的事奉,應該有知識也更應該有火,配合在一起,就是神的心意.
\newline
\newline
第四靈火的異象(王下六17)這是火車火馬的異象,這火是得勝的火,是勝過仇敵能力的火.以利沙和他年輕的僕人包圍,僕人很懼怕,聖經用哀哉二字來形容他的懼怕,(列王記下每逢有哀哉二字,表示陷於困境絕境)以利沙為僕人禱告,求神開他的眼睛.當他裡面的眼睛開了,就有新奇的發現,本來所見的,是滿山遍野的敵人,現在所見的卻是火車火馬與他們同在,保護他們.肉眼和信心之眼所見絕然不同.以利亞和以利沙同有一外號,當以利亞升天之際,以利亞說:「我父啊,以色列的戰車馬兵啊!」「以色列的戰車馬兵」這是以利亞的外號.以利沙本人臨死之際,有個王也說「以色列的戰車馬兵啊!」他們二人有同樣的名字.以利亞是個被聖靈感動的人,以利沙是個被聖靈加倍感動的人.他們二人被聖靈感動,被聖靈充滿,所以都有得勝仇敵的能力.這二人原是平凡之輩,一經聖靈充滿且加倍充滿時,他們就變成非常的人；因神的能力透過他的信心彰顯出來.我們看見有許多平常的人,因著信心被聖靈充滿,就大有能力作成了原來絕作不成的事.感謝主!火車火馬圍繞以利沙和他的僕人,與他們同在的,就比與敵人同在的更多,這是我們信心的發現.
\newline
\newline
有一天,我讀了一篇日本人寫的文章,我很高興.這日本人在香港相當出名,他是非基督徒,他曾到蘇聯旅行；回港後寫了一篇報道,由中文報翻刊出.該文載他在蘇聯發現當地有許多禮拜堂,在一個壓制數字的國家裡,竟然有許多青年人去赴會,這使他百思莫解──可見真正的信仰,是不畏任何力量攻擊的.兩週前,香港英文培靈會,有位講員,他最近到蘇聯去,進了一間禮拜堂,內坐滿一千六百人,很多都是青年；他也覺得很希奇.
\newline
\newline
感謝主!聖靈能保守我們的信心,使我們裡面的火勝過環境.當司打反被聖靈充滿時,他定睛望天,看見主耶穌站在神的右邊；雖然當時猶太人正用石頭打死他；他仍能夠下跪,注目天上,基督的愛充滿他的心.他禱告說:「主啊!不要將這罪歸於他們.」聖靈充滿的時候,他就有力量勝過殉道的痛苦,勝過被石頭打的痛苦.
\newline
\newline
聖靈要在全世界各處,照樣地工作.去年夏天,有一位弟兄從法國經香港回大陸.他對我說:他此行是為他的兒子主持婚禮.他的女兒是個很熱心的基督徒,曾參加浙江省一個大聚會；在黎明之前許多基督徒聚集在山上舉行退修會.感謝主!在任何環境中,神保守屬祂的人,聖靈的火花在屬祂的人心裡繼續燃燒.
\newline
\newline
第五靈火的異象(啟八3-6)金壇代表禱告,金壇裡有火,這一定是禱告的火.金壇的火倒在地上,表示神的計劃進行；可見禱告的火,和神計劃推行,有密切的關係.弟兄姊妹!千萬別輕看我們的禱告；叫我們輕看禱告是撒旦的破壞工作之一；我們若不禱告或少禱告,就是中了撒旦的詭計.神的話告訴我們,禱告與神的計劃進行有密切關係,神重視你我的禱告.我自從去年參加韓國福音大爆炸聚會以後,對於禱告的信心,真是大有幫助,每逢我禱告的信心較低落時,我馬上回想韓國兄弟姐妹的禱告,他們大部分教會,每天早晨四點鐘開始禱告會,廿多年如一日,這樣的禱告大有作用.請注意!全世界唯韓國基督徒增加的速度是人口增加的十倍.回想他們的禱告,我的心得到鼓勵.記得有一天下大雨,我去參加晚上聚會,我坐在車裡看見兩旁,一隊隊的青年人,用齊整的步伐往會場跑；口用韓語不斷喊著「耶穌革命」「聖靈能力爆炸」.他們在傾盆大雨中邊跑邊呼喊,我心深受感動,青年們不怕困難,長途跋踄參加聚會,雖然滂沱大雨,仍有幾十萬人聚會.真是表現聖靈爆炸的能力!弟兄姊妹!當我們禱告之火燃燒時,神的計劃就進行了.
\newline
\newline
第六靈火的異象(啟四5)火燈是代表神的七靈,七靈是代表聖靈完全的工作,在聖靈完全的工作裡面有火的作用.讓我們很自然地想到五旬節的情形,聖靈降臨時,有火焰顯現好像舌頭,落在門徒的頭上,他們被聖靈充滿,說起方言來,稱讚神的偉大.從前我讀到這段經文時,有個問題在心中「聖靈充滿和說方言有何關係呢?」似乎是兩不相干的事,為何在神安排之中,聖靈充滿說了方言呢?後來我明白了,那原因之一就是表示聖靈降臨時,門徒得著火熱的見證,舌頭在聖靈能力運行之下,把福音傳到地極,那裡有說各種方言的人,藉著聖靈的能力,把福音傳到講各種方言的人那裡去.所以講方言是個預表,預先表明在聖靈能力之下,福音傳遍全世界,講各種方言的人都要信主了.所以當聖靈降臨,我們就有了保證；從開始時就有最後效果的保證.哦!我們的主真是最好的主,凡事給我們留下見證,凡事給我們留下證據.一開始的時候,已經讓我們憑著信心,看見將來最後的結束,看見因聖靈的能力而有的大成功.今天這預表已經應驗了,福音已經傳遍全世界,這是兩千年前主講的話；這天國的福音要傳遍天下,然後末期才來到.
\newline
\newline
弟兄姊妹們!聖靈是保證耶穌基督是我們信心創始成終者；聖靈是大使命的創始成終者.
\newline
\newline


\newpage

\section{第47屆~港九培靈會~滕近輝~第二講}
\label{sec:1163}
\textbf{靈火的能力在以利亞身上的具體表現}
\newline
\newline
連結: \href{https://www.hkbibleconference.org/session-message/view/1163}{\texttt{https://www.hkbibleconference.org/session-message/view/1163}}
\newline
\newline
\hyperref[sec:1161]{< < < PREV SERMON < < <}
~
\hyperref[sec:toc]{[返目錄]}
~
\hyperref[sec:1164]{> > > NEXT SERMON > > >}
\newline
\newline
聖經用十五個比喻,形容聖靈的能力和工作,第一個比喻就是火,聖靈好像火一樣.
\newline
\newline
昨天晚上,我們從聖經看到六個靈火的異象；今天晚上,我們要從以利亞身上看見靈火的具體表現.
\newline
\newline
提起以利亞,我們就想到「火」.在迦密山上,他把祭物擺在祭壇上；神就從天降下火來,燒盡他的祭物,表示神的悅納.在(王下一章)我們看見神審判的火,藉著以利亞降在敵人身上.第二章我們看見以利亞被接上升時,有火車火馬.我再說,當我想起以利亞,我就想到火,以利亞信心的火,「禱告的火」,「奉獻之精神的火」,「能力的火」；以利亞這些表現都是因有聖靈的充滿.
\newline
\newline
以下有兩個理由可以證明:
\newline
\newline
(一)當以利亞快要升天之時,他的門徒以利沙緊緊跟著他；一步不離開,不放鬆.以利亞問他有何請求,以利沙說:「願感動以利亞的靈,加倍的感動我.」為何以利沙這樣請求呢?因他清楚看見以利亞的能力,是由聖靈而來.以利沙知道這根本的原因和秘訣後；他就求根本的事,他的要求非常好,非常寶貴.他知道以利亞所以有火的能力,是因為被聖靈充滿.
\newline
\newline
(二)(路一15):「施洗約翰從母腹裡就被聖靈充滿了.」17節:「他必有以利亞的心志能力.」從上面的經文可知以利亞的心志能力.是由被聖靈充滿而來.以利亞是個被聖充滿的人,是個有靈火在他裡面燃燒的人；從他的表現,可以看見聖靈的能力與工作.
\newline
\newline
去年七月,在瑞士舉行世界宣教會議.會中有位神的僕人,以影音節目向大家報導,全世界各地聖靈的作為.該節目名「聖靈行傳」,他發出一個使我留下深刻印象的問題:「今天以利亞的神在那裡?」接著,他又反覆地問:「今天神的以利阿們在那裡?」這個問題感動我的心,成了我從大會最大的領受.聽後,我心裡禱告說:「主阿!求你憐憫我,幫助我,使我多有以利亞的心志能力.」這個信息一直留在我心裡.
\newline
\newline
我讀啟示錄,發現(十一3)這段聖經非常重要.這裡告訴我們,在末世的時候,有兩個見證人要來,一般人都認為以利亞,是這兩個見証人之一.關於這一點,有兩種解釋:一是以利亞本人要來；一是有一位像以利亞的人要來,並非以利亞本人要來.就如耶穌說施洗約翰就是以利亞,他有以利亞的心志能力；在末世時也有位有以利亞的心志能力,他要到世上來.這兩種解釋,今晚我們暫不題孰是孰非.我們所著重的是一件不變的事──「以利亞心志的復興」.以利亞要來,表示以利亞心志能力的復興.既然神預定在末世以利亞要來,就是說,神的旨意,要末世教會有以利亞心志能力的復興.這是神對末世教會的計劃,神要在今天的教會中,有以利亞心志能力的復興.我們願意接受這個信息嗎?今晚我們要一同禱告說:「主啊!但願你這心意,在我身上不致落空；求主使我蒙憐憫,使我更多像以利亞.讓以利亞的心志能力,在這個軟弱的人身上,有更多的彰顯.」
\newline
\newline
以利亞的心志有六方面的表現:1,奉獻的心志.2,信靠神的心志.3,禱告的心志.4,順服的心志.5,愛神愛人的心志.6,受苦的心志.
\newline
\newline
以利亞的心志,是六重的心志.他心志的表現和末世的情形,兩方面的特點比較；我們可發現二者強烈的對比,而領受主給我們的信息.今晚我們謙卑俯伏在主面前,但願主的心和作為向我們顯明.願聖靈在我們心裡作工,藉著以利亞的表現,感動並提醒我們；讓我們有追求的心志,靈性向上進.
\newline
\newline
第一個對比,末世教會不冷不熱,和以利亞奉獻的心志.老底嘉時代,就是末世時代,有物質的享受.整個世界物質享受都提高了,但物質享受提高時,也就是靈性降落的時候.以利亞的心志,在迦密山上表現得最清楚,他獨自站在耶和華的一邊,面對多少敵人呢?我們常說四百五十個巴力先知.事實上,在(王上十八19)給我們看見,還有四百個耶洗別供養的先知.他們事奉另一神,名叫阿舍拉,所以合共是八百五十個先知.以利亞面對眾多敵人,他抱著信心,同時也抱著必死的決心；因他完全的奉獻,所以毫不懼怕.這種獻身的精神,流露出來,叫我們每逢讀這段聖經時,重新受感動.以利亞把祭物放在祭壇上,神降火下來表示神的悅納.我們看到祭物時,就連想到以利亞也把自己獻於祭壇上.以利亞真是毫無條件,毫無保留地擺在祭壇上,這樣的奉獻在神眼中極其寶貴.
\newline
\newline
感謝主!在今天這不冷不熱時代,聖靈讓我們有奉獻的心志!我曾在曼谷參加一個會議,有個代表從越南來的,因他的身量矮小,一直不被人注意；後來他的同工題到他的見證,於是我的感覺完全改變過來.這位矮小的傳道人,原是從越南西貢到山區去傳福音的,他有七次被越共捉去；他七次向他們傳講耶穌,最奇妙的,是七次都獲釋放.這是何等獻身的精神!他奉獻了自己,在壓力之下,危險之際,毫不退縮.我又看見過教會的一段記錄片,在泰國北部,有些人在痲瘋病人中工作；一方面為他們治病,一方面向他們傳福音.其中有位年輕的姐妹,是由荷蘭去的醫生；她離開祖國,離開溫暖的家,在遙遠陌生之地,為痲瘋病人服務.我看見有一鏡頭,她正為病人洗腐爛的腳.那腐爛不堪的腳,在影片裡我真是不想觀看.但她卻親自俯身為他們洗濯,我心很受感動!這位姐妹真正奉獻自己,不憎不惡,甘心情願地工作.
\newline
\newline
第二個對比(路十八8)主耶穌說:在末世的時候祂再來,遇得見世上有信德的嗎?照這段聖經上文,信德是指著禱告的信心而言.上文主設比喻說要常常禱告,不可灰心；跟著就說末世的時候祂再來.「遇得見有信德的嗎?」祂發出這問題,一方面是表示如此的信心很少,另方面是向我們挑戰.當祂再來時,能發現在祂的眾教會裡,有多少人有禱告的信心；你我有禱告的信心嗎?你我禱告的表現如何?
\newline
\newline
以利亞是個有禱告心志的人,我一想到他,不但想到「火」.也想到他臉藏在兩膝之間懇切禱告.他禱告了六次,仍然沒有動靜,毫無反應；到了第七次,他看見天空有如手掌大小的雲彩,他知道神聽了他的禱告.他看見的,只是手掌大小的雲彩,他的信心即非常的堅定,清楚,他立刻起來,有信心的行動.以利亞真是個有禱告信心的人,他的禱告大有能力；神特別聽他的禱告,因他有信心,又因他是個義人.他具備了這兩個條件,以信心和義行配搭起來,所以他的禱告特別蒙神應允.在末世時有這樣的信心嗎?感謝主!我見過一個最奇妙的例子.數年前我在東京領會,散會時有個中國人到台前找我,原來他是我在青島讀初中時,高我一班的同學,那時我們倆均非基督徒.我父親是中華基督教會的牧師,我雖生在牧師家中；身為牧師兒子,但我到了十八歲才重生得救,七七事變,中日戰爭,各人紛向外地逃難.後來我到西北讀書,他南向廣州,在中山大學讀醫科.到了第三年,他患上很嚴重的肺病；醫生認為沒有希望,因戰事延續,他被送往湖南一間出名的醫院.經詳細檢查,認為絕無希望了,就把他送去一個等死的地方.那時,有位醫生每天去看他,為他禱告,向他傳福音.有一天,醫生又來了,跪在他床邊為他禱告.醫生走後,他發現床邊被眼淚濕透了.他覺得很奇怪,素不相識,彼此無關,為何他這樣流淚代禱?於是他心深感動,向主敞開,接受了主.他開始禱告,奇蹟遂即出現,病完全痊癒；五天之後,他起身到各病房作見證.這就是在東京相遇時,寇弟兄對我講的見證.我們彼此緊握著手,快樂難以言喻.主救了我作祂的僕人,主也救他和全家得救.他的兒女都很熱心,他兒子是東京華僑教會的青年領袖,他女兒現在美國恩典神學院讀神學.這個活生生的見證,當時有兩個神蹟發生；一是他的身體得醫治,在很短期間,由無盼望轉為完全痊癒.二是愛心的神蹟,那位醫生對一個陌生人,竟然有如此的愛心；這個神蹟,真是比第一個神蹟更加奇妙.後來我見到于力工牧師,我問他:抗戰時期在湖南有間出名的醫院,其中有位醫生如此熱心,你有印象嗎?他說,記得有位王醫生,這件事好像是他作的.由此證實了寇弟兄的見證.感謝主!在這缺少禱告信心的時代中,仍看見屬主的人,有禱告的信心；有禱告的擔負,有禱告的愛心,有神蹟隨著發生出來.
\newline
\newline
第三個對比　我們看見末世時代,信仰幾乎破產,搖動；主耶穌在(太廿四)說:「且有好多假先知起來迷惑多人.」這些假先知很有辦法,很有口才,很有學問.因此,迷惑很多人信仰搖動.另方面我們看見以利亞,信靠神的心志,他的信心真是寶貝；因他信神又信神的話,所以神藉他作了大事.神藉他的信心,供應他的需要三年半；藉他的信心,使死人復活；藉他的信心,讓他勝過八百五十個先知；因著他的信心,神藉他行奇妙的大事.感謝主!我再說,一個沒信心者是個平常人,一個有信心者是個非常的人；如果一個有信心者後來失去信心,他就變成一個平常人；一個人是平常或非常,在乎他有沒有信心.求神使我們做個不平常的人,像我這平常的人,無可能成為不平常的人.感謝主!神有個辦法,叫我們每個人可以去做,神給我們一個秘訣,每個人都可實行.這秘訣和方法就是信心,有了信心；我們就變成一個不平常的人.何時我們失去信心,我們就立刻成一個最平常的人.以利亞在許多方面他是平常人,因他大有信心,就成了不平常的人.他的禱告,大大蒙神垂聽,神藉他的信心,作了非常奇妙的事.
\newline
\newline
有個日本牧師告訴我,當基督教傳入日本的早期,某處發生逼迫教會的事.在城外有大堆石頭,用來打基督徒的.那班受大逼迫的基督徒,不但有信心,也有愛心；他們聚在一起為逼迫他們的人禱告,神聽了他們禱告.那些逼迫他們的人中,後來有很多悔改信主；他們就商量利用那些石頭就地建造一間禮拜堂.結果,他們自己動手,建成了禮拜堂的後牆；牆前有個十字架,所用的石頭,都是那堆原來拿來打基督徒的.日本牧師說,那是一間特別的禮拜堂,在十字架的後面,有著活生生的見證.打基督徒的石頭,變成建造聖殿的材料,這是信心造成的變化.這變化見證牆前十字架,改變人生的大能大力.以利亞就有這樣堅強信靠神的心志.
\newline
\newline
第四個對比　(提後三1)保羅在此處題及末世的危險,道德的危險,屬靈的危險.這些危險中,我們特別注意「謗讟」(二節)這是對有權威的人的反抗,「不能自約」.(三節)意則不受約束,不受管制,反權威,反傳統.保羅說,這些都是末世反叛的精神.這真是今天社會的情形!許多人無所不及,好壞都反；只要是傳統就反,是權威就反,這是末世流行的精神和態度.另一方面,我們看見以利亞順服的心志；完全順服的神的命令,他所反叛的是罪惡.神吩咐他去見亞哈王,並對王說:「狗要餂你的血」.誰敢說這話呢?以利亞明知亞哈王,是暴君,但這句話未免太難出口；可是以利亞終於順服神的吩咐,去見亞哈王.將神的話照詞直言,亞哈王聽了,謙卑下來；撕裂衣服,慢慢的走,他認罪；雖是暫時的悔改,但我們看見神的話的權柄.在這反叛的時代,我們願意有以利亞那樣順服的精神.我們當順服的,不是屬於世界的一切；乃是當順服神的旨意,尊重神的吩咐,保守絕對順服的態度.有這樣的態度,就有屬靈的權威和能力；這種權威和能力,能使罪惡受轄制.感謝讚美主!我們在以利亞身上看見屬靈的權柄,因他順服耶和華.在印度有個墓碑寫著:「他什麼都不怕,因為他怕神.」這是個好基督徒的的墓碑.一個真敬畏神的人,對撒旦對罪惡是勇敢的.我們的腿只能向一邊彎,如果向神屈膝；向撒旦的一邊必是剛強的,如果向撒旦屈膝就向神剛硬.我們需要剛強的基督徒,剛強的傳道人；在這時代裡,求神給我們有以利亞心志的復興.
\newline
\newline
第五個對比　(太廿四12)末世時代是愛心冷淡的時代；因為不法的事增多,在無愛心的環境中,基督徒也失去愛心.以利亞卻有愛神愛人的心志,何以見得?大家都記得,當他所住地方那寡婦的孩子死了,他心裡非常難過,把死孩子從他母親懷中抱過來,放在自己的床上.這表示他要負責任,他有很沉重的負擔,他付上負責任的禱告；伏在孩子身上,面對面,口對口,手對手為孩子禱告.這樣表示他和孩子認同了,合而為一,成為一人了.在神面前,把孩子的生命,當作自己的生命；為所愛的人禱告,心是何等懇切!大衛為他的孩子禱告說:「押沙龍啊!我巴不得替你死.」有個母親禱告說:「主阿!求你醫治這孩子,把我的壽命加給他吧!」這是認同的禱告,以利亞就是這樣的禱告；以利沙也是這樣的禱告,我想他是從以利亞學來的.他非但學以利亞的外表,他更學習以利亞的心志；當以利沙伏在死去的孩子身上禱告時,不但外面姿勢和他一樣,其內心的負擔也是一樣.他實在是個好學生!從前宋尚節博士大大被主使用,有很多的小宋尚節出來了.他們用同樣的姿勢講道,我親眼見過一個傳道人,由口袋抽出手帕搖擺唱歌；另有一個傳道人把頭髮垂下來,學了宋尚節式的髮型.如此,不過單學了外貌,但裡面是否學到還成問題；他是否有宋尚節的愛心；信心,禱告；和他的犧牲?以利沙有個學生,就是單學他的表面.(王下五)載以利沙治愈乃縵將軍的痲瘋病,乃縵送來禮物,以利沙不肯接受；因他要乃縵知道,神的恩典是非錢可買.不管地位多高,金錢多富；他需要耶和華白白的恩惠,神的恩典是權勢金錢地位所不能買的.他唯一的目的,是要榮耀耶和華.以利沙的學生基哈西就不同了,他看見乃縵送來的禮物；主人三次不接受,他也三次不接受.但是到了第四次,基哈西接過禮物來了.他沒有真正學到以利沙的心志.以利沙是個好學生,他從以利亞有直正的學習和得著.
\newline
\newline
以利亞是個有愛心的人,他的一切表現都是因為被聖靈充滿.
\newline
\newline


\newpage

\section{第47屆~港九培靈會~滕近輝~第三講}
\label{sec:1164}
\textbf{聖靈好像耶和華的手}
\newline
\newline
連結: \href{https://www.hkbibleconference.org/session-message/view/1164}{\texttt{https://www.hkbibleconference.org/session-message/view/1164}}
\newline
\newline
\hyperref[sec:1163]{< < < PREV SERMON < < <}
~
\hyperref[sec:toc]{[返目錄]}
~
\hyperref[sec:1165]{> > > NEXT SERMON > > >}
\newline
\newline
「於是靈將我舉起帶我而去,我心中甚苦,靈性忿激；並且耶和華的靈,(原文作手)在我身上大能能力.」(結三14)這裡我覺得中文聖經譯得很好,譯者看得很清楚；耶和華的手,原意就是耶和華的靈,用「手」來形容「靈」,耶和華的靈如手要抓住我們.如果人被這手抓住,這手要在他的生命中施展工作.祂抓住誰,就在誰身上工作；然後又用這個人向別人工作.
\newline
\newline
有個傳道人說,神的手在水平線上來回的動,當人在水平線上時；這手立即抓住他,成了神手中有用的器皿.當一個基督徒或是傳道人,在他裡面有真正屬靈的渴慕和追求時,他就升上水平線；神的手就抓住他,他就變成不同的人,成為神手中合用的器皿.
\newline
\newline
當聖靈的手抓住人的時候,祂在這人身上,第一件事就是把他舉起；舉到屬靈的境界裡面.本來我們在自然界中生活,但聖靈的手抓住我們,就把我們舉起來,超過原來的水平線.
\newline
\newline
保羅在(林前二14)題到兩種人:一種是屬血氣的,一種是屬靈的.屬血氣的人就是屬自然的人,在自然的水平線生活.屬靈的人,在屬靈的高原上生活；屬靈的境界是超越的.
\newline
\newline
屬靈的境界,怎樣會高過自然的境界呢?現在要用兩個人來作例證.
\newline
\newline
一個是以利亞,他是代表聖靈感動的人.一個是以利沙,他是被聖靈加倍感動的人.我們原來都是屬血氣的,但聖靈在我們裡面,開始工作之後,就把我們改變了.祂把我們從屬血氣,改變成為屬靈；要我們順從聖靈的引導,體貼聖靈.在聖靈裡行事為人的,就是屬靈的人.這些都是聖靈的工作,把我們舉起來,放在屬天的高原上.
\newline
\newline
自然境界和屬靈境界的比較:在自然境界,有基本的規律,限制和管轄所屬的人.屬靈的人,卻超越自然的規律.
\newline
\newline
第一個比較,就是社會公平律.這個規律很寶貴,有用,而且重要.一般來說,公平律是以牙還牙,以眼還眼.在自然界的人,覺得公平是最崇高的；以眼還眼,以牙還牙,是天公地道的.但屬靈的境界,遠超公平的律,到達愛的水準.以利沙是個很好的例子:(王下六21-22)當時亞蘭軍隊落在以色列人手中,以色列王想殺他們,這是很自然的事,因為亞蘭人來侵略以色列.現在他們落在以色列人手中,以色列人大可以名正言順的殺掉他們.但以利沙說:「不可殺害他們,且為他們擺設筵席；讓他們吃飽喝足,然後送他們回到主人那裡去.」這件事似乎太反常,但他是站在屬靈水準上而作的.結果,亞蘭人從此不再侵犯以色列人了.假如我是個畫家,我必描繪亞蘭軍隊在筵席上面部的表情；他們既驚訝又受到感動,因為神的愛,澆灌在他們心裡.如果我會繪畫的話,我就會繪四幅畫,第一幅要畫的是:主耶穌釘在十字架上,禱告說:「父啊!赦免他們,因為他們所作的,他們不曉得.」第二幅是司提反即將被猶太人用石頭打死時,他舉目望天,被聖靈充滿,看見耶穌站在神的右邊.他跪著禱告說:「主啊!不要將這罪歸於他們.」第三幅是舊約的約瑟哥哥們把他賣到埃及.他離開父母家鄉,悲慘痛苦.當他做宰相見到哥哥們時,他抱著他們的頸項哭；這並非恨的哭,乃被神的愛所感動,流出饒恕的眼淚而哭.第四幅:以利沙招待被擊敗的亞蘭軍,他擺設筵席款待敵人.我們可以用一句聖經的話來形容,作為這幾幅畫的總描寫:「聖靈將神的愛,澆灌在我們心裡.」人纔有這四種的表現.
\newline
\newline
有一位宣教士在華北工作.在庚子年間,義和團事變,他自知處境危殆,急速寫信給他在美國讀書的兒子.內容說:「親愛的孩子!當你收到這信；可能我已為主被殺了.你聽到這消息後,卻不要仇恨殺我的人.我現在所禱告祈求的,是求主感動你的心,將來也到中國傳福音.我倒下去之後,你要跟著來,踏上我傳福音的腳蹤.」這封信非常的寶貝,據說現仍存在耶魯大學的圖書館裡；因為這位宣教士,是該校的校友.這種愛實非人原有的愛,乃是聖靈將神的愛,澆灌在人的心裡.
\newline
\newline
第二個比較.在自然境界,有一條生物的定律──「保全自己」.一切生物都盡力保全自己,有的生物學家把這條定律,稱為自然界的第一定律.但是在屬靈境界,第一條定律是「捨己」.
\newline
\newline
當我讀初中一時,開始學英文字母.我發現有一字,只有單獨的字母,而且是大寫的.I(我)真有出人頭地,唯我獨尊的涵義.我得救後,才發現God(神)也是大寫的.於是「我」和「神」在我思想上開始爭戰.「我」想得到統治權,「神」也要在我的生命中得到統治權.當神要得到我時,我卻緊緊抓住不肯放.我對神說:我絕對不能交出來給你,是我的主權；我絕不降服,絕不交出.感謝主!有一天,我才發現基督比我自己好得多.當我把自的一切,交在祂手中時；就是我最蒙恩,最幸福,最快樂.我發現耶穌基督,是最貴,最美麗,最崇高最完全的.祂的愛,祂的公義,真理,信實是最完美的.如果我把自己的生命,完全交在祂手裡,這才是我最有智慧的抉擇.所以當我把自己獻給主的時候,心裡滿有快樂；自己放下,得著基督.我已經與基督同釘十字架,現在活著的,不再是我；乃是基督在我裡面活著.不再是我乃是基督,這是最合算的交易；用我來交換基督,這是好得無比的.願在座每位弟兄姊妹!都看見基督的美好,甘願的,喜樂的無條件的把自己的生命,交在主手中.這時,聖靈的愛就要充滿你心,你就有完全不同的人格.
\newline
\newline
第三比較　在自然界中有一條數學的定律(2 + 2 = 4).但在屬靈的境界,超越了這個數字.列王記上十七章載:先知以利亞照神的吩咐去做一件難作的事.神要他向那寡婦啟齒.把僅有一把麵,一點油分享.寡婦毫不猶豫,照著先知以利亞所說的去作,用那最後所有的一把麵和一點油,做了餅先供給以利亞.這樣油與麵,並無短少,直到耶和華使雨降在地上的日子.(王上十七14)由此,這是一個屬靈境界的方程式:1 $\div$3 = 3 $\times$1260 這是屬靈境界的數學.從一點點變成豐富,我們的耶和華是使無變有,叫死人復活的神.「在耶和華豈有難成的事嗎?「我豈沒有對你們說過,你們必因信看見神的榮耀嗎?」只要我們有信心,就能看見神奇妙的大作為.屬靈境界的數學,是藉除得乘.除就是分,藉著分而增加,這是基督國度的定律.倫敦美術館內,有一幅畫,描繪一個人剛死了,躺在床上.那畫的重點,是兩隻抓緊的手；下面題著──「我所抓住的,我失去了；我所用去的,我抓住了.」這畫含意深長!許多人拼命的抓,不肯放手,結果兩手成空.好像空氣又像水；水流去了,空氣散掉了.生命結束時,何等虛空!主告訴我們:當我們與別人分享時,我們就得著了.保羅說:你們知道主耶穌基督,祂本來是富足,卻為你們成了貧窮；叫你們因祂的貧窮,可以成為富足.(林後八9)
\newline
\newline
第四個比較　自然境界,有一條生理定律.「勞動必然產生疲乏」這定律,相信沒有人否定吧!但是在屬靈的境界,卻超過了這定律:「奔跑卻不疲倦」(王上十九7-8)以利亞灰心了,他在蘿藤樹下求死；一個原來最堅強的人,現在竟然求死,後來他睡著了.睡覺也是主的恩典,人在灰心時,能甜睡一覺實在是很好的；如果你精神不佳,心靈愁煩時,別忘記好好睡一覺,這是以利亞的方法.人的身體和靈性有密切關係,身體差了；靈性也受影響,除非聖靈給予特別的恩典.我曾見過一個身體很軟弱的人,但聖靈給他特別恩典,於是軟弱變為剛強.一件自然的事,因聖靈的工作而變為超自然.當以利亞求死,非常灰心疲倦時；耶和華卻為他預備食物,天使來拍醒他,對他說,起來吃吧!他本來不吃喝；現在起來吃喝；這些飲食給他力量行走了四十晝夜,一直走到何烈山.超自然的食物,才能使他行走四十晝夜.「就是少年人也要疲乏困倦,強壯的也必全然跌倒；但那等候耶和華的,必從新得力,他們必如鷹展翅上騰,他們奔跑卻不困倦,行走卻不疲乏.」(賽四十30-31)超自然的食物,增加他超自然的力量,把他托起舉起,進入超自然的境界裡.有一次我到兵頭花園,見到一根堅固的大石柱,蔓藤環繞而上,蔓藤是一種軟弱的植物,無法自立；但依靠堅立的石柱,卻由最軟弱變成最強壯.雖經大風吹襲,仍然依附石柱屹立不動.當我們藉著禱告,抱住我們的神,我們就成了最強壯的；可以如鷹展翅,展開禱告的翅膀,在空中飛翔.鷹在空中飛翔而不疲乏,秘訣何在呢?因為鷹在空中展翅滑翔,牠知道利用風和氣流的作用,使牠平衡,安穩在空中,這就是鷹從新得力的秘訣.弟兄姊妹!我們若不學會禱告等候耶和華,有一天,我們也會疲乏,停止,退後,跌倒.等候耶和華,是我們必須學會的秘訣.鷹的翅膀是身體面積的廿五倍,牠翅膀伸開時；就可以把身體帶上天空.照樣,禱告的翅膀,也能把我們帶上去.以賽亞被聖靈感動說:「凡等候耶和華的,必從新得力,如鷹展翅膀上騰.」
\newline
\newline
第五個比較　自然界有一個物理的定律:「冷則收縮,熱則膨脹.」但在屬靈的境界,超過這個例子；約伯說:「耶和華使人在夜間歌唱.」夜間較冷而且黑暗,易使人懼怕灰心頹喪；感謝神!耶和華卻使人在夜間挺身昂首而歌唱.
\newline
\newline
在困難,試煉,寒冷之時,真能有信心擴展的歌唱嗎?以利沙身上有真實的表現(王下三末段),有三個王被困在曠野,無水喝,而且敵人迫近,他們來見以利沙.以利沙吩咐他們到處掘溝.他們不明白,以為水份已奇缺,若還掘溝,必使身體水份耗盡而死.但以利沙要他們憑信心挖溝時,神的作為立即彰顯；次日清早,所挖的溝,都滿了水.還有一次,以利沙的門徒死了.那寡婦來求他說:我丈夫死了,家裡很窮；欠下的債無力償還,現在債主要來捉我的兒子抵償.以利沙吩咐她回家去,向鄰舍借空器皿,(聖經載不要少借)這是信心的擴展.以利沙臨終時,他吩咐以色列王把窗門打開,向外射箭；接著又吩咐王多拿些箭,向地射去.王只射了三次,以為放無的之矢未免太糊塗.後來以利沙對王說,這是耶和華的得勝箭,多射多得勝,現在你射了三次；你只能得勝三次.信心的箭就是耶和華的得勝箭,我們若射信心的箭,就能看見耶和華為我們爭戰.
\newline
\newline
以利沙和以利亞同有「以色列戰車馬兵」的稱號.他們一人的力量,相等於以色列的全部戰車馬兵,以色列得勝的力量乃是以利亞,以利沙.這兩人同是被聖靈充滿的人,所以有力量使他們的國家得勝.
\newline
\newline
求主幫助,使我們追求,讓聖靈在我們心靈裡工作,把我們舉起來,超越屬血氣的境界──自然的境界；把我們舉起到屬靈的境界.在那裡,我們體貼聖靈,順從聖靈行事,受聖靈的引導.
\newline
\newline


\newpage

\section{第47屆~港九培靈會~滕近輝~第四講}
\label{sec:1165}
\textbf{聖靈的手降在以西結身上}
\newline
\newline
連結: \href{https://www.hkbibleconference.org/session-message/view/1165}{\texttt{https://www.hkbibleconference.org/session-message/view/1165}}
\newline
\newline
\hyperref[sec:1164]{< < < PREV SERMON < < <}
~
\hyperref[sec:toc]{[返目錄]}
~
\hyperref[sec:1166]{> > > NEXT SERMON > > >}
\newline
\newline
先知以西結七次說:「耶和華的靈(原文作手)降在我身上.」(結三32).當耶和華的手,降在以西結身上時,他看見了耶和華的榮耀,就俯伏於地.(23節)這就是自然的反應,任何人看見耶和華的榮耀,必然俯伏於地.保羅在大馬色的路上,主的榮光向他顯現,他就立刻仆倒在地.當約翰見主榮耀之異像時,也立刻仆倒.(啟一章)無論保羅,約翰,以西結,都有同樣的反應,沒有一個人能在耶和華的榮耀前站立得住.人的老我若是不倒下去,還未受對付,表示他沒有看見耶和華的榮耀.如果真的看見了,就必有反應,立刻仆倒在地.他必發現自己的不配,污穢,虧欠,抬不起頭來.「禍哉!我滅亡了,因為我是嘴唇不潔的人；又住在嘴唇不潔的民中,又因我眼見大君王萬軍之耶和華.」(賽六5)
\newline
\newline
感謝主!我們屬靈的經歷,不是停止在此.(24節):「靈就進入我裡面,使我站起來.」聖靈第一步工作,叫我們在神的榮耀前俯伏.第二步工作,俯伏以後,老我受了對付,自己被破碎；然後聖靈就扶持,叫我們站起來.其中先後次序很重要.自己必須先倒下去,而後站起來；如果未倒先站,就會發生很多問題.教會發生的問題是因為人在未被聖靈打倒,未肯俯伏在神面前；只想出頭站起來,高舉自己.這等人或早或遲,他必要倒下去.俯伏和起來,先後次序非常重要；先俯伏,被破碎,然後在聖靈扶持之下站起來.不是自己站起來,乃是聖靈在他裡面,使他必然站起來；掃羅倒下去了,保羅就要起來.
\newline
\newline
我們應該把「俯伏」,和「站起來」,兩部份的次序弄得清清楚楚.但是,還有第三部份也是同樣重要的.耶和華對以西結說:「你進房屋去,將門關上.」這很希奇!聖靈既使他站起來,卻沒有對他說:你可以大踏步出去工作,有聖靈在你裡面,有聖靈的能力充滿你,你可放心勇往直前去工作.卻說:「你進房去,將門關上.」將門關上,就是進入內室.(27節)告訴我們,耶和華接著又對他說:「但我對你說話的時候,必使你開口,你就要對他們說,主耶和華如此說......」在內室嚴密關閉之中,他聽見神的聲音說:「以西結啊!是你個啞吧,如果我還未對你說話你就開口,是沒用的；你先要安靜地聽,接受我的話.我的話就成了你的話,我的話透過你的口講出來,你便可以說:「耶和華如此說.」這話帶著耶和華所有的權威.由此,我們看到很重要的第三部:當人得著聖靈能力充滿時,必須記得,他受了極重要的限制.不是去傳自己被聖靈充滿,如何如何的好,用聖靈的恩賜來炫耀自己.若為了自己的好處而作,就辜負了聖靈的恩賜.耶和華對以西結說:聖靈使你起來,你要聽我的話進房屋去；將門關上,等候我的話.以西結遵行了,他成為神所重用的先知.
\newline
\newline
當耶和華的靈,耶和華的手降在以西結身上時,他學習了功課:(一)在榮耀的耶和華面前俯伏於地.(二)然後被聖靈舉起來,剛強站立.(三)接受屬靈原則的限制,等候耶和華的話,耶和華說講就講,說去就去,說隱藏就隱藏.這三部份都很重要.
\newline
\newline
「耶和華的靈(原文作手)在我身上大有能力.」(結三14)「在他們中間憂憂悶悶的坐了七日.」(15節末句),耶和華的手抓住了以西結,把他舉起來；那時他是否感覺飄飄然呢?不是,聖經說他「心中甚苦,靈性忿激,在河邊以色列人中間,憂憂悶悶的坐了七日.」此話何解?難道一個人被聖靈抓住就憂憂悶悶嗎?心裡感覺很苦,靈性忿激嗎?果真如此,誰還敢求聖靈充滿?感謝主!這裡所指的意思,乃是當聖靈的手,抓住一個人的時候；在他心靈中,必有重要的負擔.以西結為他的同胞就有負擔,他心裡的憂悶非為自己；乃是目睹他的同胞以列人離棄神,墮落,悖逆神,反抗神.他心裡痛苦,有很重的負擔,同時也有個願望；怎能將耶和華榮耀的異象──就是他親眼所見的實現在他同胞以色列人中間.這是他新的願望,新的負擔,新的感動；這負擔在他心裡一直蘊藏,直到了他成為耶和華的出口.
\newline
\newline
當尼希米聽見耶路撒冷城被毀壞,城門被焚燒,同胞受大災難的壞消息時,他就哭泣,一連哭了幾天,撕裂衣服,認罪禱告.(尼一1-4)聖靈把負擔給他,他想到了祖國,同胞,聖城,聖殿受凌辱；被破壞,受敵人壓制,他心極其痛苦!他愛耶和華,耶和華的百姓受到羞辱；他就撕裂衣服,在耶和華面前一連痛苦了幾天,禁食禱告.
\newline
\newline
保羅說:「我在基督裡說真話,並不謊言,有我良心被聖靈感動,給我作見證；我是大有憂愁,心裡時常傷痛.為我弟兄,我骨肉之親,就是被咒詛,與基督分離,我也願意.」(羅九1-3)我相信他的話不是虛假的,他說有他的良心被聖靈感動作見證.這真實的感情,沉重的負擔,叫我們看見保羅的心腸.(中英文皆譯為心腸)心腸是人身最內部的東西,在保羅心靈最深處受了攪動.他看見同胞以色列人剛硬不肯悔改,在那情形之下,他寧願自己被咒詛；這份感情真是叫我們受感動.我們為了主,為了同胞,為了教會,為自己的靈性,為主的榮耀,為靈魂的得救,我們的感情如何呢?
\newline
\newline
摩西為那些犯罪的以色列人禱告；因以色列人背棄神敬拜金牛犢.當摩西發現時,心裡極其難過,他向神禱告說:「這百姓犯了大罪,求你赦免他們,倘若不然.......」(出卅二32)這一些小點,我想是聖經中僅有的點滴.我每讀至此,彷彿看見摩西在流淚,這真是他的淚點.他禱告說:「倘或不然,求你從你所寫的冊上,塗抹我的名.」摩西的感情和保羅的感情相近,這是真正的負擔.當聖靈的手抓我們的時候,就把負擔托住我們.
\newline
\newline
我讀教會歷史,看見有一人名叫約拿單愛德華,他在美洲教會復興史中佔了很重要的地位.這位主的僕人,有一天他讀聖經,讀到保羅三天三夜不吃不喝,然後被聖靈充滿；出去為主作見證,成為福音運動的開始.他讀至此,心中大受感動,他也三晝夜不吃,不喝,不睡,不停地禱告；將他的心靈完全傾倒出來,這就成了美洲英格蘭區教會復興運動的開始.一位神的僕人,三晝夜不吃不喝,不睡覺,在施恩台前禱告；結果天門敞開,神豐富的恩典裂天而降.
\newline
\newline
另外一次,我讀教會名人傳記,看見摩爾根博士被稱為歐美傳道人之王子.他得到這樣崇高的榮譽,是因他將主的話吸收在自己心中.他日夜將他的心他的靈,滋浸在神的話中,然後奉主的名出去講道；他在倫敦西敏寺禮拜堂事奉.(不是國家那間大禮堂名字相似而已)他在那裡被主大大使用.他的傳記名為「聖經人」(直譯)舊約有「神人」,新約有「基督人」,他的傳記卻叫「聖經人」,是一本非常寶貴的傳記.書中也題到在他之先,有很多牧師的事蹟.其中有一位,每週由一至六,上午坐在每弟兄姊妹們座位上,(那時各弟兄姊妹,有固定的座位,且需繳座位費.)逐一的為他們禱告.他設身處地,形同一人,來到施恩台前懇切禱告.這是負擔,把弟兄姊妹的需要,放在自己肩頭上；站在他的地位上禱告.這樣的代禱多麼懇切!多麼有力量!我們真是不能不說:「神很難拒絕這樣的禱告.」
\newline
\newline
耶和華吩咐以西結,向左側臥三百九十天,然後向右四十天.(結四4)這是何意?乃是要用他的身體來代表以色列人的苦難；由此可見,先知以西結,和他的同胞的苦難打成一片.(我再用認同二字)以西結和他的同胞完全認同了.耶和華所吩咐他作的,如果他不了解神的心意,就會覺得很辛苦,甚至要後悔當了神的僕人.向左側臥三百九十天,向右側臥四十天,共一年多的時間,他若不了解神的心意,此事實難作下去.我相信以西結被聖靈的手抓住時,他的心和神的心在一起了；他體會父神的心,知道耶和華對祂的百姓是如何的感覺.他體會神對以色列人的愛和關心,他是出於了解而作.我們事奉主,有時覺得很苦,但我們若了解主的心,我們就能勝過一切的難處.如保羅說:「我體會耶穌基督的心腸.」這樣,無論主差遣我們作什麼,我們就從了解的心靈裡,發出所作的事,那時感覺不同了,那時我們就有真正屬靈的負擔.
\newline
\newline
今天我們所需要的,就是聖靈抓住我們.為了神的榮耀,為了人的得救,為了教會的復興,為了自己靈性的長進；要有這樣的負擔,來事奉主.若然,我們馬上就會看見聖靈的果效,看見聖靈的工作.這樣,我們就可以領受(詩五五22)的話:「你要把你的重擔,卸給耶和華；祂撫養你,祂永不叫義人動搖.」有一部英文的修正譯本:「你把耶和華所加給你的重擔,卸給耶和華.」這話真寶貝!我們的負擔是耶和華加給的,可能我們在這負擔下覺得沉重；但聖經教我們要把耶和華加給我們的負擔卸還給祂.這樣,負擔的感動有了,但負擔的壓力卻沒有了.這是秘訣,如果以西結沒學這秘訣,他必憂憂悶悶死了.他的健康是受不了.他的心靈受不了,擔子實在太沉重.我相信以西結,保羅和大衛,都學會了秘訣；從神領受了負擔之後將負擔的重擔交回給神.神藉這句話對我說:「我的孩子,我的僕人,我將重擔給你,不會過於你所能受的.我的孩子,我的僕人,把這重擔交回給我吧!」叫我們所有服事神的人,都學會這功課.我們被聖靈的手抓住,得到了負擔；然後把擔子交還給主.藉著禱告說:「聖靈啊,我自己做不成這事,這是你的重擔,你藉著我擔起這重擔.」那時,我們可以和主一樣說:「我的擔子是輕省的.」主在馬太十一章末兩節所講的,就是這個意思.主說:「你們要負我的軛,我的軛是容易的,我的擔子是輕省的.」如果有聖靈給我們的感動,給我們的負擔；我們應為此感謝主,接了過來,表示我們的順服；表示我們的信心,然後再交回給主.所有在教會得神使用的人,都應學會這個秘訣:先接受,後交回去；神有祂的美意要使用我們,叫我們分享祂的榮耀.
\newline
\newline
以西結書(八1-4)載有兩件事並排在一起:(一)他又看見神榮耀的異象.(二)擺在神榮耀異象的旁邊,就是在耶和華殿裡有偶像.這是惹動神震怒的事,當這兩件事並排在一起時,以西結的感覺和平常完全不同.他在神的榮耀中,看罪惡,在耶和華的榮耀中看偶像；他深深感到偶像奪去了耶和華的榮耀,這罪何等大而深重!先知以西結在耶和華的榮光中看見罪,看見自己不潔的嘴唇；以前他對自己的嘴唇沒什麼感覺,自以為區區之病,微不足道,何必小題大做,大驚小怪呢!但是現在他看見了耶和華的榮耀,他就很緊張激動地說:「禍哉!我有禍了!我是不潔的人；因為我看見了萬軍耶和華的榮耀.」這時,他的環境和感覺完全不同了.我們繼續下去,以西結看見在聖殿的外院門口；聖殿的內內外外,有爬蟲,走獸,塔模斯(他們愛情之神)日頭四樣的偶像.而且在院門口的牆上,有個很小的窟窿；神叫他打開,發現裡頭有門,推開門,看見牆上滿了彫刻圖畫的偶像.然後他又發現有七十長老,就是以色列的領袖,在那裡拜那些偶像.神對他說:「以色列家的長老各人在暗中所行的,你看見嗎?」(12節)繼之,他又看見耶和華殿外院朝北的門口有婦女坐著,為塔模斯哭泣(14節)塔模斯是他們愛情之神──淫亂的愛情；他們在那裡大動感情,深受感動為淫亂愛情之神哭泣.(16節)又載,在祭壇中間,(祭壇是聖殿最神聖之處)有廿五人,背向耶和華的殿,面向東方拜日頭；他們雖身在聖殿,心卻向著偶像,何等可悲!以西結被聖靈的手抓住,看見神的榮耀；同時在榮耀的旁邊,看見這些偶像和所發生可悲的事.他的心極其沉重,他的心破碎,他也體會神的心是何等的沉重破碎.
\newline
\newline
假如我寫一篇文章,題目:「一板之隔」這樣的題目你們作何感想?屬靈的文章怎會一板之隔?在(結四三8中句)「他們與我中間僅隔一牆,並且行可憎的事,」一邊是罪惡,一邊是祭壇,僅一牆之隔,多麼可憐!多少時候,在教會中,在基督徒心中,在我們心中；一邊「主」一邊「罪惡」僅一牆之隔!耶和華為祂的百姓傷心；以西結體會神的感覺,也為他們傷心.
\newline
\newline
另方面我們看見還有一種人,完全不同.他們為耶路撒冷全城之中所行可憎的事,嘆息哀哭；耶和華吩咐以西結尋找這班人,在他們額上畫記號.(九4)這些憎恨罪惡的人,是耶和華所要使用的,他們恨神所恨,愛神所愛,乃是耶和華所愛的人.我們聯想在(啟七),神的使者在屬祂的人額上打記號.(啟十四1)說:在錫安山上,有十四萬四千人,有耶和華的名,有耶路撒冷的名寫在他們額上.這班人是誰呢?(啟三12)說:「我要叫他在我神的殿中作柱子,我又要將我神的名,和我神城的名,並我的新名,都寫在他上面.」這是指得勝的基督徒說的,在他們額上有記號,所以我們得了一個結論,(14節)所指在錫安山上十四萬四千人是那些得勝的基督徒,他們額上有得勝的記號.感謝神!神所認識的人,在額上都有記號,使他們永遠與祂同在.得勝的基督徒,神特別賜他的榮耀；凡以西結作記號的人,都是神特別寶貝的人.在充滿罪惡的城裡,在充滿罪惡的潮流中；他們為罪憂傷,沒失去他們的義心,沒失去屬靈的敏感.他們有眼淚可流,他們的心,仍然向著耶和華.感謝主!在今天罪惡潮流中,神仍然保守許多屬祂的人,沒受罪惡潮流的影響.額上有記號的.是勝過罪惡的人,撒旦不能殺害他們；所以啟示錄九章載:無底坑開的時候,有黑煙冒出,黑煙中有蝗虫,蝗虫刺傷人；但蝗虫的刺不能傷害額上有神印記的人.(蝗虫和牠的剌代表撒旦用罪惡來傷害人),在罪惡潮流中勝過,逆流而上.惟有額上有印記的人,基督在他們裡面；基督的生命在他們裡面,所以能得勝.向撒旦誇勝.刺雖厲害,可怕,惡毒,但額上有印記的人可以勝過牠；撒旦無法可害那些有基督生命在他裡面的人.
\newline
\newline
哈利路亞!基督的生命,就是我們勝過罪惡的力量.願我們打開我們的心,以單純的信心,接受耶穌基督的生命進入我們裡面.我們就有力量,在罪惡的潮流中過得勝的生活.我們額上有神的印記.
\newline
\newline


\newpage

\section{第47屆~港九培靈會~滕近輝~第五講}
\label{sec:1166}
\textbf{聖靈的作為}
\newline
\newline
連結: \href{https://www.hkbibleconference.org/session-message/view/1166}{\texttt{https://www.hkbibleconference.org/session-message/view/1166}}
\newline
\newline
\hyperref[sec:1165]{< < < PREV SERMON < < <}
~
\hyperref[sec:toc]{[返目錄]}
~
\hyperref[sec:1167]{> > > NEXT SERMON > > >}
\newline
\newline
先知以西結七次說:「耶和華的靈(原文作手)降在我身上.」(結三32).當耶和華的手,降在以西結身上時,他看見了耶和華的榮耀,就俯伏於地.(23節)這就是自然的反應,任何人看見耶和華的榮耀,必然俯伏於地.保羅在大馬色的路上,主的榮光向他顯現,他就立刻仆倒在地.當約翰見主榮耀之異像時,也立刻仆倒.(啟一章)無論保羅,約翰,以西結,都有同樣的反應,沒有一個人能在耶和華的榮耀前站立得住.人的老我若是不倒下去,還未受對付,表示他沒有看見耶和華的榮耀.如果真的看見了,就必有反應,立刻仆倒在地.他必發現自己的不配,污穢,虧欠,抬不起頭來.「禍哉!我滅亡了,因為我是嘴唇不潔的人；又住在嘴唇不潔的民中,又因我眼見大君王萬軍之耶和華.」(賽六5)
\newline
\newline
感謝主!我們屬靈的經歷,不是停止在此.(24節):「靈就進入我裡面,使我站起來.」聖靈第一步工作,叫我們在神的榮耀前俯伏.第二步工作,俯伏以後,老我受了對付,自己被破碎；然後聖靈就扶持,叫我們站起來.其中先後次序很重要.自己必須先倒下去,而後站起來；如果未倒先站,就會發生很多問題.教會發生的問題是因為人在未被聖靈打倒,未肯俯伏在神面前；只想出頭站起來,高舉自己.這等人或早或遲,他必要倒下去.俯伏和起來,先後次序非常重要；先俯伏,被破碎,然後在聖靈扶持之下站起來.不是自己站起來,乃是聖靈在他裡面,使他必然站起來；掃羅倒下去了,保羅就要起來.
\newline
\newline
我們應該把「俯伏」,和「站起來」,兩部份的次序弄得清清楚楚.但是,還有第三部份也是同樣重要的.耶和華對以西結說:「你進房屋去,將門關上.」這很希奇!聖靈既使他站起來,卻沒有對他說:你可以大踏步出去工作,有聖靈在你裡面,有聖靈的能力充滿你,你可放心勇往直前去工作.卻說:「你進房去,將門關上.」將門關上,就是進入內室.(27節)告訴我們,耶和華接著又對他說:「但我對你說話的時候,必使你開口,你就要對他們說,主耶和華如此說......」在內室嚴密關閉之中,他聽見神的聲音說:「以西結啊!是你個啞吧,如果我還未對你說話你就開口,是沒用的；你先要安靜地聽,接受我的話.我的話就成了你的話,我的話透過你的口講出來,你便可以說:「耶和華如此說.」這話帶著耶和華所有的權威.由此,我們看到很重要的第三部:當人得著聖靈能力充滿時,必須記得,他受了極重要的限制.不是去傳自己被聖靈充滿,如何如何的好,用聖靈的恩賜來炫耀自己.若為了自己的好處而作,就辜負了聖靈的恩賜.耶和華對以西結說:聖靈使你起來,你要聽我的話進房屋去；將門關上,等候我的話.以西結遵行了,他成為神所重用的先知.
\newline
\newline
當耶和華的靈,耶和華的手降在以西結身上時,他學習了功課:(一)在榮耀的耶和華面前俯伏於地.(二)然後被聖靈舉起來,剛強站立.(三)接受屬靈原則的限制,等候耶和華的話,耶和華說講就講,說去就去,說隱藏就隱藏.這三部份都很重要.
\newline
\newline
「耶和華的靈(原文作手)在我身上大有能力.」(結三14)「在他們中間憂憂悶悶的坐了七日.」(15節末句),耶和華的手抓住了以西結,把他舉起來；那時他是否感覺飄飄然呢?不是,聖經說他「心中甚苦,靈性忿激,在河邊以色列人中間,憂憂悶悶的坐了七日.」此話何解?難道一個人被聖靈抓住就憂憂悶悶嗎?心裡感覺很苦,靈性忿激嗎?果真如此,誰還敢求聖靈充滿?感謝主!這裡所指的意思,乃是當聖靈的手,抓住一個人的時候；在他心靈中,必有重要的負擔.以西結為他的同胞就有負擔,他心裡的憂悶非為自己；乃是目睹他的同胞以列人離棄神,墮落,悖逆神,反抗神.他心裡痛苦,有很重的負擔,同時也有個願望；怎能將耶和華榮耀的異象──就是他親眼所見的實現在他同胞以色列人中間.這是他新的願望,新的負擔,新的感動；這負擔在他心裡一直蘊藏,直到了他成為耶和華的出口.
\newline
\newline
當尼希米聽見耶路撒冷城被毀壞,城門被焚燒,同胞受大災難的壞消息時,他就哭泣,一連哭了幾天,撕裂衣服,認罪禱告.(尼一1-4)聖靈把負擔給他,他想到了祖國,同胞,聖城,聖殿受凌辱；被破壞,受敵人壓制,他心極其痛苦!他愛耶和華,耶和華的百姓受到羞辱；他就撕裂衣服,在耶和華面前一連痛苦了幾天,禁食禱告.
\newline
\newline
保羅說:「我在基督裡說真話,並不謊言,有我良心被聖靈感動,給我作見證；我是大有憂愁,心裡時常傷痛.為我弟兄,我骨肉之親,就是被咒詛,與基督分離,我也願意.」(羅九1-3)我相信他的話不是虛假的,他說有他的良心被聖靈感動作見證.這真實的感情,沉重的負擔,叫我們看見保羅的心腸.(中英文皆譯為心腸)心腸是人身最內部的東西,在保羅心靈最深處受了攪動.他看見同胞以色列人剛硬不肯悔改,在那情形之下,他寧願自己被咒詛；這份感情真是叫我們受感動.我們為了主,為了同胞,為了教會,為自己的靈性,為主的榮耀,為靈魂的得救,我們的感情如何呢?
\newline
\newline
摩西為那些犯罪的以色列人禱告；因以色列人背棄神敬拜金牛犢.當摩西發現時,心裡極其難過,他向神禱告說:「這百姓犯了大罪,求你赦免他們,倘若不然.......」(出卅二32)這一些小點,我想是聖經中僅有的點滴.我每讀至此,彷彿看見摩西在流淚,這真是他的淚點.他禱告說:「倘或不然,求你從你所寫的冊上,塗抹我的名.」摩西的感情和保羅的感情相近,這是真正的負擔.當聖靈的手抓我們的時候,就把負擔托住我們.
\newline
\newline
我讀教會歷史,看見有一人名叫約拿單愛德華,他在美洲教會復興史中佔了很重要的地位.這位主的僕人,有一天他讀聖經,讀到保羅三天三夜不吃不喝,然後被聖靈充滿；出去為主作見證,成為福音運動的開始.他讀至此,心中大受感動,他也三晝夜不吃,不喝,不睡,不停地禱告；將他的心靈完全傾倒出來,這就成了美洲英格蘭區教會復興運動的開始.一位神的僕人,三晝夜不吃不喝,不睡覺,在施恩台前禱告；結果天門敞開,神豐富的恩典裂天而降.
\newline
\newline
另外一次,我讀教會名人傳記,看見摩爾根博士被稱為歐美傳道人之王子.他得到這樣崇高的榮譽,是因他將主的話吸收在自己心中.他日夜將他的心他的靈,滋浸在神的話中,然後奉主的名出去講道；他在倫敦西敏寺禮拜堂事奉.(不是國家那間大禮堂名字相似而已)他在那裡被主大大使用.他的傳記名為「聖經人」(直譯)舊約有「神人」,新約有「基督人」,他的傳記卻叫「聖經人」,是一本非常寶貴的傳記.書中也題到在他之先,有很多牧師的事蹟.其中有一位,每週由一至六,上午坐在每弟兄姊妹們座位上,(那時各弟兄姊妹,有固定的座位,且需繳座位費.)逐一的為他們禱告.他設身處地,形同一人,來到施恩台前懇切禱告.這是負擔,把弟兄姊妹的需要,放在自己肩頭上；站在他的地位上禱告.這樣的代禱多麼懇切!多麼有力量!我們真是不能不說:「神很難拒絕這樣的禱告.」
\newline
\newline
耶和華吩咐以西結,向左側臥三百九十天,然後向右四十天.(結四4)這是何意?乃是要用他的身體來代表以色列人的苦難；由此可見,先知以西結,和他的同胞的苦難打成一片.(我再用認同二字)以西結和他的同胞完全認同了.耶和華所吩咐他作的,如果他不了解神的心意,就會覺得很辛苦,甚至要後悔當了神的僕人.向左側臥三百九十天,向右側臥四十天,共一年多的時間,他若不了解神的心意,此事實難作下去.我相信以西結被聖靈的手抓住時,他的心和神的心在一起了；他體會父神的心,知道耶和華對祂的百姓是如何的感覺.他體會神對以色列人的愛和關心,他是出於了解而作.我們事奉主,有時覺得很苦,但我們若了解主的心,我們就能勝過一切的難處.如保羅說:「我體會耶穌基督的心腸.」這樣,無論主差遣我們作什麼,我們就從了解的心靈裡,發出所作的事,那時感覺不同了,那時我們就有真正屬靈的負擔.
\newline
\newline
今天我們所需要的,就是聖靈抓住我們.為了神的榮耀,為了人的得救,為了教會的復興,為了自己靈性的長進；要有這樣的負擔,來事奉主.若然,我們馬上就會看見聖靈的果效,看見聖靈的工作.這樣,我們就可以領受(詩五五22)的話:「你要把你的重擔,卸給耶和華；祂撫養你,祂永不叫義人動搖.」有一部英文的修正譯本:「你把耶和華所加給你的重擔,卸給耶和華.」這話真寶貝!我們的負擔是耶和華加給的,可能我們在這負擔下覺得沉重；但聖經教我們要把耶和華加給我們的負擔卸還給祂.這樣,負擔的感動有了,但負擔的壓力卻沒有了.這是秘訣,如果以西結沒學這秘訣,他必憂憂悶悶死了.他的健康是受不了.他的心靈受不了,擔子實在太沉重.我相信以西結,保羅和大衛,都學會了秘訣；從神領受了負擔之後將負擔的重擔交回給神.神藉這句話對我說:「我的孩子,我的僕人,我將重擔給你,不會過於你所能受的.我的孩子,我的僕人,把這重擔交回給我吧!」叫我們所有服事神的人,都學會這功課.我們被聖靈的手抓住,得到了負擔；然後把擔子交還給主.藉著禱告說:「聖靈啊,我自己做不成這事,這是你的重擔,你藉著我擔起這重擔.」那時,我們可以和主一樣說:「我的擔子是輕省的.」主在馬太十一章末兩節所講的,就是這個意思.主說:「你們要負我的軛,我的軛是容易的,我的擔子是輕省的.」如果有聖靈給我們的感動,給我們的負擔；我們應為此感謝主,接了過來,表示我們的順服；表示我們的信心,然後再交回給主.所有在教會得神使用的人,都應學會這個秘訣:先接受,後交回去；神有祂的美意要使用我們,叫我們分享祂的榮耀.
\newline
\newline
以西結書(八1-4)載有兩件事並排在一起:(一)他又看見神榮耀的異象.(二)擺在神榮耀異象的旁邊,就是在耶和華殿裡有偶像.這是惹動神震怒的事,當這兩件事並排在一起時,以西結的感覺和平常完全不同.他在神的榮耀中,看罪惡,在耶和華的榮耀中看偶像；他深深感到偶像奪去了耶和華的榮耀,這罪何等大而深重!先知以西結在耶和華的榮光中看見罪,看見自己不潔的嘴唇；以前他對自己的嘴唇沒什麼感覺,自以為區區之病,微不足道,何必小題大做,大驚小怪呢!但是現在他看見了耶和華的榮耀,他就很緊張激動地說:「禍哉!我有禍了!我是不潔的人；因為我看見了萬軍耶和華的榮耀.」這時,他的環境和感覺完全不同了.我們繼續下去,以西結看見在聖殿的外院門口；聖殿的內內外外,有爬蟲,走獸,塔模斯(他們愛情之神)日頭四樣的偶像.而且在院門口的牆上,有個很小的窟窿；神叫他打開,發現裡頭有門,推開門,看見牆上滿了彫刻圖畫的偶像.然後他又發現有七十長老,就是以色列的領袖,在那裡拜那些偶像.神對他說:「以色列家的長老各人在暗中所行的,你看見嗎?」(12節)繼之,他又看見耶和華殿外院朝北的門口有婦女坐著,為塔模斯哭泣(14節)塔模斯是他們愛情之神──淫亂的愛情；他們在那裡大動感情,深受感動為淫亂愛情之神哭泣.(16節)又載,在祭壇中間,(祭壇是聖殿最神聖之處)有廿五人,背向耶和華的殿,面向東方拜日頭；他們雖身在聖殿,心卻向著偶像,何等可悲!以西結被聖靈的手抓住,看見神的榮耀；同時在榮耀的旁邊,看見這些偶像和所發生可悲的事.他的心極其沉重,他的心破碎,他也體會神的心是何等的沉重破碎.
\newline
\newline
假如我寫一篇文章,題目:「一板之隔」這樣的題目你們作何感想?屬靈的文章怎會一板之隔?在(結四三8中句)「他們與我中間僅隔一牆,並且行可憎的事,」一邊是罪惡,一邊是祭壇,僅一牆之隔,多麼可憐!多少時候,在教會中,在基督徒心中,在我們心中；一邊「主」一邊「罪惡」僅一牆之隔!耶和華為祂的百姓傷心；以西結體會神的感覺,也為他們傷心.
\newline
\newline
另方面我們看見還有一種人,完全不同.他們為耶路撒冷全城之中所行可憎的事,嘆息哀哭；耶和華吩咐以西結尋找這班人,在他們額上畫記號.(九4)這些憎恨罪惡的人,是耶和華所要使用的,他們恨神所恨,愛神所愛,乃是耶和華所愛的人.我們聯想在(啟七),神的使者在屬祂的人額上打記號.(啟十四1)說:在錫安山上,有十四萬四千人,有耶和華的名,有耶路撒冷的名寫在他們額上.這班人是誰呢?(啟三12)說:「我要叫他在我神的殿中作柱子,我又要將我神的名,和我神城的名,並我的新名,都寫在他上面.」這是指得勝的基督徒說的,在他們額上有記號,所以我們得了一個結論,(14節)所指在錫安山上十四萬四千人是那些得勝的基督徒,他們額上有得勝的記號.感謝神!神所認識的人,在額上都有記號,使他們永遠與祂同在.得勝的基督徒,神特別賜他的榮耀；凡以西結作記號的人,都是神特別寶貝的人.在充滿罪惡的城裡,在充滿罪惡的潮流中；他們為罪憂傷,沒失去他們的義心,沒失去屬靈的敏感.他們有眼淚可流,他們的心,仍然向著耶和華.感謝主!在今天罪惡潮流中,神仍然保守許多屬祂的人,沒受罪惡潮流的影響.額上有記號的.是勝過罪惡的人,撒旦不能殺害他們；所以啟示錄九章載:無底坑開的時候,有黑煙冒出,黑煙中有蝗虫,蝗虫刺傷人；但蝗虫的刺不能傷害額上有神印記的人.(蝗虫和牠的剌代表撒旦用罪惡來傷害人),在罪惡潮流中勝過,逆流而上.惟有額上有印記的人,基督在他們裡面；基督的生命在他們裡面,所以能得勝.向撒旦誇勝.刺雖厲害,可怕,惡毒,但額上有印記的人可以勝過牠；撒旦無法可害那些有基督生命在他裡面的人.
\newline
\newline
哈利路亞!基督的生命,就是我們勝過罪惡的力量.願我們打開我們的心,以單純的信心,接受耶穌基督的生命進入我們裡面.我們就有力量,在罪惡的潮流中過得勝的生活.我們額上有神的印記.
\newline
\newline


\newpage

\section{第47屆~港九培靈會~滕近輝~第六講}
\label{sec:1167}
\textbf{新酵與舊酵}
\newline
\newline
連結: \href{https://www.hkbibleconference.org/session-message/view/1167}{\texttt{https://www.hkbibleconference.org/session-message/view/1167}}
\newline
\newline
\hyperref[sec:1166]{< < < PREV SERMON < < <}
~
\hyperref[sec:toc]{[返目錄]}
~
\hyperref[sec:1168]{> > > NEXT SERMON > > >}
\newline
\newline
馬太福音 13:31-13:33
\newline
\begin{longtable}{cl}
\hline
\hline
章節 & 經文 (和合本修訂版)\\
\hline
13:31 & \begin{tabularx}{0.7\textwidth}{X} 他又設個比喻對他們說：「天國好比一粒芥菜種，有人拿去種在田裡。 \end{tabularx} \\ \\ \relax
13:32 & \begin{tabularx}{0.7\textwidth}{X} 它原比所有的種子都小，等到長起來，卻比各樣的菜都大，且成了樹，以致天上的飛鳥來在它的枝上築巢。」 \end{tabularx} \\ \\ \relax
13:33 & \begin{tabularx}{0.7\textwidth}{X} 他又對他們講另一個比喻：「天國好比麵酵，有婦人拿來放進三斗麵裡，直到全團都發起來。」 \end{tabularx} \\ \\
[1ex]
\hline
\hline
\end{longtable}
(太十三31-33)這段經文,主向我們講兩個關於天國的比喻,告訴我們:天國發展的兩個基本動力.(一)是信心,主用芥菜種的比喻說明.(二)是聖靈的能力,主用麵酵的比喻說明.信心和聖靈的能力,這兩個動力,絕對不可缺少；何處有這兩個能力,教會就得以發展；基督的教會,是在這兩種基本動力之下發展的.
\newline
\newline
有人認為芥菜種的比喻,表示教會的變質和腐化.這裡說芥菜長成樹,豈不是反常嗎?還說,飛鳥宿在它的枝上；上文說:「飛鳥表示撒旦,撒種的人撒種子,有的撒在路邊,飛鳥來把它吃掉.」既然飛鳥代表撒旦,為甚麼芥菜種的比喻,又象徵信心呢?以下有三個理由可以證明.
\newline
\newline
(1)芥菜長成樹,在聖地不是反常的事.數年前,兄弟到耶路撒冷聖城去,在城門口,看見一棵芥菜樹；滕師母站在樹下,我為她拍了一張照片留念；芥菜樹比她高過三倍,真是不愧為樹,在聖地通常都可以看到.
\newline
\newline
(2)比喻中的飛鳥不一定是指撒旦,因為在聖經中亦同樣的比喻.有時是好的方面,有時是壞的方面.例如獅子,有時是象徵主耶穌；但有時是表示撒旦,不能一概而論.
\newline
\newline
(3)大多數的解經家,都認為是指好的方面.
\newline
\newline
所以兄弟把芥菜種的比喻,肯定的解作信心.
\newline
\newline
基督的國度,是按著信心的原則來推動發展的.我再說:「那裡有信心,那裡的教會就有發展.」這是不變的原則.
\newline
\newline
信心造成擴展,在舊約,我們從亞伯拉罕身上看見這個真理.亞伯拉罕曾建四個壇,這四個壇,都和信心的擴展有關.
\newline
\newline
第一個壇(創十二7)當亞伯拉罕到了迦南,發現那一帶地方都被別人佔據了；但他抓住神的應許,憑著信心去領受,就建立祭壇獻上感謝.神的應許尚未兌現,尚未成全,他事先已憑信心來接受,他的信心已經擴展了.
\newline
\newline
第二個壇建於伯特利和艾城之間,那是個高原地帶,居高臨下,俯覽約但河肥美的平原,他看見耶和華對他的應許,是何等寶貝.這是信心的展望,因信心舉目四望,看見耶和華寶貴的應許.
\newline
\newline
第三個壇(創十三末兩節)耶和華對亞伯蘭說,你起來,縱橫走遍這地；凡你腳所踏之地,我都賜給你為業.亞伯蘭就起來,舉起信心的腳步縱橫走遍那地.走完之後,他就建壇在希伯崙.這是信心腳步的擴展.
\newline
\newline
第四個壇(傳廿二9)亞伯拉罕帶他獨生愛子以撒到摩利亞山上.築了壇,親手將他的兒子捆綁放在壇上,要把他殺了當作燔祭獻給耶和華.這是信心的奉獻,當他將奉獻時,聽見有聲音說:「亞伯拉罕,千萬不可伸手傷害你的兒子,現在我知道你是敬畏耶和華的了.」接著就應許說:「你的後裔,要像天上的星,海邊的沙那樣多.」從信心的奉獻,有了將來的擴展.那裡有信心,那裡就有擴展；那裡沒有信心,那裡就有萎縮.
\newline
\newline
聖靈的能力也是天國進展的動力　主耶穌在(太十三32-33)告訴我們「天國好像麵酵」麵酵有「發」的作用,主耶穌用麵酵象徵聖靈.從前我覺得這樣講法似乎通不過.聖經明明說麵酵象徵罪惡,為何又說象徵聖靈呢,豈不是聖經自相矛盾嗎?不是的,現在讓我們把兩處聖經作個對比:
\newline
\newline
利未記廿三章,題到以色列人七個大節.七個大節期,代表耶和華救恩計劃的七個步驟.
\newline
\newline
逾越節　預表耶穌基督,作了逾越節的羔羊,完成了救贖之功.
\newline
\newline
除酵節　預表耶穌基督所流寶血,除淨我們一切的罪.這裡酵代表罪,除酵則把罪除去；耶穌基督的十字架,除去我們一切的罪過了.
\newline
\newline
初熟節　在逾越節的第三天,預表耶穌基督從死裡復活,成了睡了之人初熟的果子.我們的主復活了,有一天,所有屬於祂的人,也要從死裡復活.
\newline
\newline
五旬節　(這個節期我們要特別注意)在初熟節後第五十日,(16-17節)特別清楚告訴我們,在這天,以色列人要把新素祭獻給耶和華.新素祭是用細麵做兩個餅,要把酵加在裡面；然後獻給耶和華,這裡特別說要把酵加進去.令我們覺得稀奇的,是除酵節除了酵,為何過了五十天又要加酵?神的話並不矛盾,除酵節的酵是舊酵,象徵罪惡；五旬節的酵是新酵,乃是代表聖靈；五旬節是聖靈降臨的時候,聖靈一降下來,馬上起了「發」的作用；原來懼怕,軟弱,退後的門徒；現在有完全的改變,他們的信心,勇敢,熱心都發起來.因為從天上來的新酵,已臨到他們,所以教會也就不斷發展了.當聖靈工作的時候,就有三千人信主,五千人信主；使徒行傳最後的記載,在耶路撒冷有十幾萬的基督徒.在很短的時間,雖受極大的壓迫,政治上的攔阻和宗教上的攔阻,雙管齊下；但是在聖靈帶領之下,基督的教會反而發展,得救的人數增加.耶穌的十字架被高舉,耶穌基督復活的大能到處顯明；真是新酵發起的作用!毫無疑問,從聖經的預表,看見新酵就是代表聖靈的能力.
\newline
\newline
在教會進展過程中,我們必須追求信心,追求聖靈的能力；如果缺少這兩樣,教會不能興旺,不能發展,不能進步.我們的靈性不能登高進深,讓我們從主的話語,抓住這兩個原則.主在這兩個比喻中,清楚告訴所有的基督徒,和所有的教會領袖:若要教會興旺,必須有信心,也必須有聖靈的能力,二者均不可缺.　當新酵來時,能把舊酵作用除去；如果沒有新酵的作用,我們絕無可能把舊酵的作用除去.很多人幹了一件錯誤愚蠢的事,只想把舊酵(罪惡)的作用除去；但沒得到新酵的作用,因先後的次序錯了,所以他們追求的,永遠不能成功.我們必須先得到新酵的作用,那就是聖靈的能力,開始在我們身上工作.這時,舊酵的作用就要被消除了,作用就停止了.第一步造成第二步,而第二步是跟著第一步來的,不可先後顛倒.
\newline
\newline
新酵和舊酵的對比:
\newline
\newline
(一)舊酵造造成惡化(弗四22)「就要脫去你們從前的舊人,這舊人是因私慾的迷惑,漸漸變壞的.」食物腐化是因細菌侵入.聖經說人漸漸變壞,是因私慾(罪惡)的迷惑.19節說:「放縱私慾,貪行種種的污穢.」每個時代都是如此,但在末世時代更甚.啟示錄十七章載大淫婦的異象:大淫婦坐在眾水之上,眾水代表各國,各民.如果要描繪大淫婦坐在眾水之上,實無從下筆；我想只好畫個地球,大淫婦坐在地球上.大淫婦坐在整個世界上,她的作用,她的權勢,她的影響力普遍全世界.她手裡高舉一個金杯,裡面充滿牠行淫的污穢,她把那污穢可憎卑賤之物放在金杯裡.不是把污穢藏起來,而是高高舉起,以羞辱為榮耀,以淫亂為榮耀,這真是如今末世的情形!有人說,淫亂的事每個時代都有,不足為奇；但神的話告訴我們,末世時代,人在大淫婦影響之下；把淫亂當作榮耀,和過去的時代大有分別.過去時代是偷偷摸摸地作,不敢見天日；現在竟然公開明目張膽犯淫亂的罪,而且這種淫風,漸趨影響基督的教會.影響教會的青年,甚至影響教會的領袖.主的話告訴我們,大淫婦的權勢很強很盛.在這樣的世代,我們需要新酵──聖靈能力的工作.
\newline
\newline
新酵造成美化(弗一13)這裡特別題到「印記」古代的印記,是立體的形像或圖案.聖靈的印記,是耶穌基督的形像.當聖靈的印記,印在我們身上時,我們身上就開始顯出,耶穌基督的模樣.聖靈的印記,有寶貴的內容；聖靈的工作,要在我們身上顯出耶穌基督的樣式,這是何其美好!真是美化!香港政府計劃要美化香港,耶穌基督藉著聖靈,要美化我們的人生.讓我們的精神,心靈及整個人生；有基督的樣式,最美麗的情況,基督化的人生.(林後三18),我們眾人既然敞著臉得以看見主的榮光,好像從鏡子裡返照,就變成主的形狀,榮上加榮,如同主的靈變成的.這是聖靈的工作,把我們自己原來的醜惡改變,漸漸有更多基督的樣式.這個過程,可說是屬靈的拍照,菲林就是我們從基督得的新生命,菲林有感應,才照出像來.如果我們沒有基督的新生命,聖靈沒法在我們裡面工作,我們對基督的形像,也就沒有感應.我們應該把相機鏡頭的蓋子除去,正如(林後三16)所說:「把我們臉上的帕子除去.」與基督之間沒有任何隔膜,與主面對面,目不轉睛注視仰望主的榮面.這時聖靈的光在基督面上,使基督的面,向我們心靈彰顯出來.菲林還須經過黑房,就是我們的密室,內室,安靜與主面會之處；透過密室,顯影工作就完成了.感謝主!聖靈的光照在基督面上,而我們有更多基督的樣式.
\newline
\newline
(二)舊酵造成硬化　「心裡剛硬」(弗四18)就與神的生命隔絕.對神剛硬,對聖靈和屬靈的事也必剛硬.19節首句:「良心既然喪盡」原文是:「良心超過受感動的限度」中文譯得很好,意即不可能受感動.我親眼看見一個青年咒罵他的母親,正罵時,他聽見一個同學在樓下喊他；他立即下樓,和同學有講有笑地走了.咒罵母親,馬上又講又笑,竟然無動於中!人的心究竟會落到這樣剛硬的地步嗎?一個被稱為純潔的青年,他的良心至於這樣麻木嗎?這是不可否認的事實；聖經的話並不言過其實,也不誇張；人類正是墮落至如此地步.
\newline
\newline
新酵造成溶化　「因為我們兩下藉著祂被一個聖靈所感,得以進到父面前.」(弗二18)從前遠離神,背向神,向神剛硬；現在受聖靈感動,心開始軟化,改變他的方向,面朝向天父了.「不要叫神的聖靈擔憂.」(弗四30)這是完全的改變,聖靈感動他,他接受聖靈的感動；他怕聖靈擔憂,他的心柔和了.靈恩的雨,澆灌在他心田.「以恩慈相待,有憐憫的心彼此饒恕,正如神在基督裡饒恕了我們一樣.」(弗四32)往日剛硬的心,被聖靈溶化了,他想到基督的饒恕,就接受這饒恕用在別人身上.他清楚看見了屬靈的邏輯,他用基督對他的饒恕的去饒恕別人,他整個的心靈都軟化了,何等寶貴!你曾看見基督徒因得罪他人而流淚,因得罪他人而向人認罪；這是聖靈在他心裡作工,軟化了他的心,叫他柔和下來.
\newline
\newline
(三)舊酵使人分化　1. 神與人之間的分化(弗二13)人遠離神,是舊酵造成的分化作用,叫人與神越分越遠.　2. 人與人之間的分化(弗二14)題及一個中間隔斷的牆,牆意義之一是罪惡,罪惡成牆隔開於人與人之間.罪惡成為一道橫面的牆,使神和人隔開；罪惡成了正面的牆,使人與人隔開.那裡有罪,那裡人和神就隔開；那裡人與人的關係,就被破壞.為何今天國際間充滿了仇恨?為何人類的歷史是戰爭的歷史?為何家人自相爭鬥?為何主耶穌在世時,有個人求祂幫助,分家產?為何最親愛的人,後來變成最大的仇人?這都是罪惡的舊酵,造成的結果.聖經給我們看到,當初罪惡進到人類中間,馬上亞當逃避神,他藏起來.耶和華問說:「亞當你在那裡?」罪到人間,夫婦之間馬上生了隔膜,亞當和夏娃互相推諉責任；互相埋怨,彼此對立.繼而該隱和亞伯,兩弟兄衝突起來；該隱殺了亞伯,到了創世記第六章有句話說:「世界滿了強暴.」整個世界,因為罪已進入人類中間造成了分化.
\newline
\newline
新酵造成同化作用　「聖靈所賜,合而為一的心.」(弗四3)那裡有聖靈,大家合而為一,非常甜美的同化!保羅的意思說:「弟兄們!你們嚐過這甜美的滋味以後,你們就會竭力保守合而為一的心.」那一個教會有聖靈的運行,那一個教會就有合而為一的心.教會就會在地若天,在那裡有基督的榮耀,有靈性的長進,有工作的開展.第二節題到五種聯繫的力量:　1. 謙虛　2. 溫柔　3. 忍耐　4. 互相寬容　5. 和平,這五種聯合的力量,把基督徒緊緊地聯繫在一起.(4-6)又告訴我們:「聖靈叫我們看見,身體只有一個,就是只有一個教會；聖靈只有一個,同有一個指望；一主,一信,一洗,一神；這是所有基督徒合而為一的七個基礎.當我們看清楚認定以後,我們真是不能不合一了；當我們看見有不合一的情形時,我們的心真要破碎,這是聖靈的工作.
\newline
\newline
(四)舊酵造成奴化作用　在罪惡中的人,他們順服空中掌權者的首領.(魔鬼)牠有權柄,牠下令,他們都順服.罪惡叫這些人作撒旦的奴隸.有的國家推行奴化教育,撒旦正是如此；牠用各樣方法,各種工具在人類中推行奴化教育.
\newline
\newline
那時,你們在其中行事為人,隨從今世的風俗.」風俗是撒旦的王牌,是撒旦最善於利用的工具；撒旦最善於造成一種潮流,使人隨從.造成一種風氣,把人帶走；人人不甘落伍,極力追上潮流.潮流的力量實在很大,這股潮流把人類衝向牠的方向；這是撒旦拿手好戲,很少人可勝過潮流和風氣.這是撒旦的第一工具.第二個工具是蜜糖,引誘人,給人很多的承諾.第一工具叫人「怕」怕不合潮流,怕受人批評,怕和人不同.第二工具是蜜糖的引誘,這兩種工具是撒旦拿在手中,用以制服人的.這兩種工具,也曾用於主耶穌身上,撒旦對耶穌說,你餓了,你需要食物,可吩咐石頭變成食物.牠讓耶穌怕,怕沒有食物,怕挨餓.接著牠讓耶穌看萬國的榮華,如果耶穌肯拜牠,一切的好處都可以給.撒旦兩個工具,一個是「怕」一個是「引誘」.換言之,一個是「鞭子」一個是「蜜糖」.魔鬼手拿著這兩工具,觀看你是何等人,然後使用牠的工具.　感謝主!有種人不怕魔鬼的那些工具,正因為聖靈曾在他心裡工作.
\newline
\newline
新酵聖靈作成強化的工作(弗三16)神藉著祂的靈,就是聖靈,叫我們心裡的力量剛強起來.不怕魔鬼的鞭子,不怕為主有所損失,不怕不合潮流,不怕和人不同,不怕別人的歧視.有了這些力量,就不怕魔鬼,用牠的試探甜言蜜語來引誘我們；但是請注意!17節說:這力量是有基礎的,「使基督因你們的信,住在你們的心裡,叫你們的愛心有根有基.」基督的愛,是有根有基的愛.是長,闊,高,深的愛.有了這愛的基礎,我們就能剛強.
\newline
\newline
(五)舊酵造成毒化(弗四31)此節兩次提到「毒」,一切的苦毒,惡毒都是靈性的毒,罪惡使我們的靈性中毒.有個毒販說:「我們的雄心大志,就是要毒化整個世界.」出此語而不以為恥!這班人正是撒旦的傳聲者,牠要使人類的心靈充滿苦毒和惡毒.先毒別人,然後毒了自己；有時先毒自己,然後毒了別人,整個人類都受毒害.
\newline
\newline
新酵聖靈使教會聖化(弗二22)神藉聖靈居住在教會裡,教會是基督的身體.基督的身體是眾肢體合成的,每個屬基督的人,都是基督的肢體.許多肢體合在一起,就是基督的身體,這個基督的身體就是教會.神藉著聖靈住在教會中,換言之,神藉聖靈使教會聖化.
\newline
\newline
(六)舊酵造成暗昧化(弗四18首句)罪惡叫人的心越來越黑暗,且叫人把心門關閉；不讓聖靈的光明進去.撒旦最怕聖靈的光照,進人的靈裡.
\newline
\newline
新酵聖靈造成明朗化(弗一17-19)保羅說:「願神將祂啟示的靈賜給你們.」當神啟示的靈,賞給我們時,就叫我們看見三件事,讓這三件事明朗化.聖靈光照我們的心眼,叫我們清楚看見:　(1)讓我們真知道天父,真正認識祂,是何等寶貴.(2)叫我們清楚看見神,在我們身上的基業是何等的豐富.(3)叫我們看見神,向信祂的人所顯的能力,是何等浩大.當這三件事明朗化時,一個基督徒,就有明顯的改變.求主讓我們看見祂的可愛,看見在永世裡我們基業的豐富,值得我們全心全意去追求；求主讓我們看見祂能力何等浩大,使我們更加追求祂的能力.當這三件事明朗化時,我們整個的眼光改變,整個的反應改變,整個的追求改變.　願今天晚上,眾弟兄姊妹有聖靈在我們心中,讓這三件事明朗化.叫我們看見主的可愛,看見將來產業的豐盛,看見主向信祂的人所顯的能力,何等浩大,然後我們全心全意向著祂.
\newline
\newline


\newpage

\section{第47屆~港九培靈會~滕近輝~第七講}
\label{sec:1168}
\textbf{聖靈的油也是能力的油}
\newline
\newline
連結: \href{https://www.hkbibleconference.org/session-message/view/1168}{\texttt{https://www.hkbibleconference.org/session-message/view/1168}}
\newline
\newline
\hyperref[sec:1167]{< < < PREV SERMON < < <}
~
\hyperref[sec:toc]{[返目錄]}
~
\hyperref[sec:1169]{> > > NEXT SERMON > > >}
\newline
\newline
今晚我們要看兩個聖靈的比喻.
\newline
\newline
第一個請看(創一2)末句「神的靈運行在水面上.」運行二字,呂振仲博士新譯本是「覆煦」母雞孵在雞蛋上,那就是覆煦；聖靈的工作也是一樣.
\newline
\newline
覆煦有三種作用,也可說覆煦的工作包括三個因素:第一,覆煦裡面有熱力,母雞覆在雞卵上,這熱力幫助雞卵逐漸變化.如果沒有熱力,小雞無法形成.第二,覆煦裡面有愛力,母雞孵在卵上,有愛在其中.第三,覆煦裡面有生命力,雞蛋是從母雞生命產生出來的.母雞的熱力和愛力,都助雞蛋的生命孵以生長.那裡有聖靈的工作,那裡就有聖靈的熱力；就有聖靈的愛力,就有生命力.教會中如果有聖靈運行,那間教會就有福了；聖靈運行時,那時教會的弟兄姊妹靈性裡,就嚐到熱力和愛力.生命的力量,把他們托起來,而造就栽培他們.
\newline
\newline
聖靈工作時有熱力(徒四31)使徒們禱告,上文說:使徒們遭遇的逼迫,官府和宗教領袖用極大的壓力加在他們身上,禁止他們奉主耶穌的名講論.官府和宗教領袖的壓力,雙管齊下,擬將他們打倒!門徒聽到這消息,就聚在一起,同心合意禱告,他們的禱告滿了熱力；滿了信心,屬靈的能力,流露於字裡行間.「禱告完了,聚會才震動；他們就都被聖靈充滿,放膽講論神的道.」我們可以想像到,他們當時禱告的情形；相信他們不是默默無聲地禱告,必是大聲地禱告.並非故意大聲,乃是很自然地,從裡面湧出同心合意懇切的禱告；聖靈在他們中間,他們心靈裡,有著聖靈的熱力.
\newline
\newline
去年十月間,在馬來西亞的吉隆坡,舉行各教會領袖福音會議,各宗派的領袖都參加.有聖公會,浸信會,信義會,衛理公會,長老會,還有其他獨立的教會.大家聚在一起,好像大家庭一樣.我心裡有個感動,就向他們講聖靈充滿的信息,我所以講這題目,是因為在會前一個月中；我自己心靈經歷了聖靈的能力,經歷了聖靈的充滿.未講之前,我向主禱告說:「主啊!我要特別輕鬆地講,若有人受感動,那證明完全是聖靈的工作,不是任何的感情作用；我要很輕鬆地講,求聖靈來作工.」我開始講了一刻鐘,我感覺聖靈運行在會場裡面,聖靈也充滿在我裡面,我覺得不能再講了,我說:「我不需要講了,我們一齊禱告.」禱告的時候,我意想不到的情形發生了；那些大公會的傳道人,都以迫切宏亮的聲調禱告,(請各位聽這段錄音)照我想像他們在平時,沒有大聲禱告的習慣；我覺得很希奇,深知聖靈作工.那次的經驗,讓我對聖靈的工作,有進一步的體會；在聖靈工作時,禱告就有熱力,一變常態.
\newline
\newline
聖靈工作時有愛力(徒四7)司提反被聖靈充滿；當時猶太人用石頭要打死他,他跪下注目望天.看見主耶穌站在神右邊,他心靈流露了愛的禱告:「主啊!求你不要將這罪,歸於他們身上.」這樣的禱告,何等有熱,何等有愛力!這禱告給掃羅聽見了.(徒七末段)告訴我們,掃羅當時也在場,他看見司提反死的情形.他親耳聽見了司提反臨終的禱告,聖經沒記載當時他有任何改變；但我深信這件事,在他心靈深處像個炸彈一樣,當時他想要把這感動壓下去.從那時開始,他變本加厲逼迫教會；我相信這裡面,必有一個心理因素.他想壓抑內裡大力的感動,就加強外面逼迫的行動；但是有一天,復活的主,在大馬色的路上向他顯現.他見到復活之主的榮耀,他仆倒在地；他裡面的感動一觸就爆出來,三天三夜不吃不喝的禱告,認罪痛哭.他回想舊約的話,一面安靜,一面心如波浪翻騰大受感動.司提反的愛力,乃因他被聖靈充滿而產生；這愛力不落空,雖然司提反死了,但這愛力透過聖經,直至今日對千千萬萬人的心靈講話.兩千年來,每個人讀經,讀到這裡時,無不受感動,我相信許多基督徒,因著這段經文的記載,因這愛力復興起來!多少原來逼迫教會的人,讀到這愛力,他們就被感動,被溶化了.感謝主!司提反之死,不是徒然的,他在聖靈充滿時死.他顯出耶穌基督十字架的愛,也可在人的身上流露出來.如果這愛力單在耶穌基督身上給人看見,人就要說基督是神,不是人；祂能這樣,我做不到.當這愛也在司提反,這人身上顯出來時,我們再也無話可說了.我們就看見一件事實,當聖靈充滿時,神的愛就澆灌在人心裡.讓人的心也像基督的心一樣,基督神聖的愛,也能從人的心裡流露出來.
\newline
\newline
聖靈的第五個比喻　聖靈工作好像油一樣　(賽六一──第一句)「主耶和華的靈在我身上,因為耶和華用膏膏我.」這裡把靈和膏連在一起,聖靈的工作,好像膏油工作一樣；聖靈的膏油,非常寶貴.我們不但要得到聖靈的火,不但要被聖靈的手抓住,不但要被聖靈的酵發起來,不但要經歷聖靈的熱力,愛力；我們更要經歷聖靈膏油的工作.聖靈膏油的信息最少有五個:
\newline
\newline
(一)喜樂油(賽六一3)聖靈的油是喜樂油,誰心裡有了這油,就充滿喜樂了.聖靈充滿表現之一,就是屬靈的喜樂；如果一個人說他被聖靈充滿,卻沒有喜樂,我們很難相信.聖靈所充滿的人,有喜樂油膏抹在其中.(路十21)記載:「正當那時,耶穌被聖靈感動就歡樂,說:「父啊!天地的主,我感謝你,......因為你的美意本是如此.」祂被聖靈感動,歡樂起來；對父神滿了讚美,讚美之中,祂說:「父啊!你的美意本是如此,你的真理向聰明通達人就藏起來,向嬰孩就顯出來.」當時主耶穌受人反對,受人冷落；但是聖靈感動他的心,就歡樂起來.聖靈充滿時,無論環境如何,都有喜樂充滿在心中.記得我讀戴德生傳記:他早期在中國江蘇省鎮江縣工作,他住的房子被燒掉,在他的日記中寫了一句很寶貴的話:「我的環境就是神,神是我的環境.」他的環境,不是那被燒掉的房子,不是那不順利的環境；外面的一切,不是他真正的環境.心靈裡面另有的環境,就是他在神裡面生活,在神裡面工作.感謝主!保羅書信充滿了四個最寶貴的字:「在基督裡」這四字,有兩個基本的意思:(1)「在基督裡,」是我們到父神那裡去的通道,藉著耶穌基督可到天父面前.(2)「在基督裡,」表示基督是我們屬靈的環境,我們的言語,思想,行動,態度,反應,一切的一切,都在基督的範圍裡面,何等寶貴!一個信主的人,就是從世界遷入基督裡面；基督是我們的生命,我們的聖潔,我們的得勝,我們的力量.我們一切的一切,就是基督,在基督裡,我們可以享受取用祂的一切；不是說基督為我作甚麼,乃是說基督成了我的甚麼.「基督是我的一切,我的一切在基督裡.」這是保羅的訊息,是我們所信的福音；是天父在基督裡,為我們成全的救恩.我們又想到聖靈所結九樣的果子,第一是愛,愛是最高,因為神就是愛；所以愛排行第一,我們都很贊成.第二是喜樂,我想你不會不贊成吧!但是(加五22)明明告訴我們,聖靈所結的果子,就是仁愛,喜樂......」可見喜樂是何其重要!我們的天父是好天父,祂願意我們都有喜樂,神叫我們看見,有了祂的愛,我們就有喜樂；我們在祂愛中時,我們就被喜樂充滿.我再說,若一個人自認為被聖靈充滿,卻沒有喜樂的表現；我們就難以證明,他是真正被聖靈充滿.保羅說:「因為神的國不在乎喫喝,只在乎公義和平,並聖靈中的喜樂.」這是神國的三個基本要素,神的國,以公義,和平,並聖靈中的喜樂為原則.這是神國的基本性質,神的國若只有公義和平,恐怕太單調啦!當我們知道神國基本原則之一是喜樂時,我們便歡歡喜喜進神的國.感謝主將喜樂賜給我們!祂創造人的面,如果人的面不會笑,那成甚麼樣子呢?如果全世界的人都不會笑；那有甚麼意思,還是及早離開這個世界.中國人形容笑,不是口笑,而是眉開眼笑；人面上有笑容,整個面都活起來了.人生有笑,整個人生都活起來.基督的愛將屬靈的喜樂,充滿我們心中；使我們整個的人,活潑起來.充滿基督同在的喜樂,那時有基督同在,那裡就有喜樂.如果一個人愛主愛得很苦,這人有問題了；如果一個熱心禱告得很苦,那他的禱告,還沒有達到主旨意的境界.如果一個人熱心得很苦,這人很可憐.如果一個人覺得多多奉獻很苦,他的奉獻必有問題.如果聖靈充滿我們,基督的喜樂,基督的愛必充滿我們.那麼為主背十字架,也覺得這十字架是恩典；我們為主付代價,那代價成為詩歌.「因耶和華而得的喜樂,是你們的力量.」(尼八10)我查考聖經時,發現希西家王年間有一個大復興.(王下廿九至卅一)這三章聖經,記載這個大復興,有七次說以色列人喜樂.其中兩次說:「歡歡喜喜.」一次說「儘都歡喜.」一次「極大的喜樂.」一次「從所羅門以來,沒有這樣的喜樂過.」每次題到喜樂,前面加有形容詞,這是復興的表現.讀到此,我們不能不相信,他們有了真正的復興.照這三章的記載,當時以色列人守了七天的逾越節,他們意猶未盡,決定再守七天,這是超過律法.七天的逾越節,是律法的規定；每個以色列人無論歡喜,不歡喜都必遵守；但是現在復興來了,他們大家商議,一致決定再守七天.這第二個七天,是恩典的充滿,超過律法；不在律法之下而是進入恩典的範圍裡面,直昇到恩典的境界.感謝主!我們今天不再在律法之下,乃是在基督的恩典裡頭,這恩典是白白賜給我們,讓我們裡面有自由,無拘無束；讓我們以神為樂,以主的話為樂；以親近主為樂,這是真正的復興.你看見主的百姓,歡歡喜喜地到主面前,可知這是真正的大復興了.記得我讀詩篇百廿二篇2節:「耶路撒冷啊,我的腳站在你的門內.」當時我不甚了解,以為這話是多餘的；腳站在門內或外有甚麼好說的呢?為何大衛寫詩要插進這句話呢?直到有一天我到聖地去,下了飛機；我看見有個乘客,整個人俯伏在地,用嘴親地,心情激動.我一看,便知他是以色列人,他第一次腳踏祖國之地.如果他說:「耶路撒冷啊,我的腳踏在你的地上.」這話豈不充滿豐富深刻的感情!有的人寫詩,從描寫辭典找到美麗的詞句,但令人讀了不受感動；有的輕描淡寫,卻令人讀了很受感動,因他是從心裡真實的感情發出來.如果大衛不愛耶和華,他就寫不出那麼感動人的詩,感謝主!聖靈是喜樂油,叫我們充滿喜樂.
\newline
\newline
(二)醫治的油(雅五14-16)雅各說弟兄姊妹之中,有人病了,該請教會的長老來；奉主名用油抹他,為他禱告.這話何意?通常有三種解釋:第一種解釋,油代表醫藥,在那時代,油是很重要的藥品.抹油意思是一面禱告,一面服藥；醫藥是神的恩賜,禱告也神的恩賜.二者並用.第二解釋,油象徵聖靈的醫治,禱告時抹不抹油均可；抹油只是表面的事,注重於禱告求聖靈醫治.第三解釋,油是象徵聖靈的醫治,既有這個象徵,就用這個象徵.人需要象徵抽象屬靈的道理,如有記號表達,更容易得幫助.他們同意油代表聖靈的醫治,但認為應該用這象徵的辦法,所以用橄欖油抹了禱告.我並非指出那個解釋最好,我意思乃是說油代表醫治；聖靈能醫治我們軟弱的身體,也有能力醫治我們屬靈的疾病.雅各很均衡,他講完了身體的醫治之後,他說:「你們要彼此認罪,互相代求,使你們可以得醫治.」(16節)這醫治包括身體靈性都在內.弟兄姊妹!我們需要健康,身體健康是神的恩典,靈性健康更是神的恩典.我們的天父,在靈性上為我們預備了醫治,讓我們的身體得醫治,靈性也得醫治.
\newline
\newline
(三)分別為聖的油　會幕中有聖膏油,所有物件凡是為耶和華而用的,都要抹上聖膏油,分別為聖歸給耶和華.膏油有香氣,會幕中有三種香味,(一)是香壇燒香發出的,這香是用五種香料製成的,香代表禱告馨香的味.(二)是獻素祭時用的乳香,這種特別的香料,燃燒就發出香味,充滿全會幕,素祭預表主耶穌用完美的人生,奉獻給父神.素祭乃用最細的麵製作並抹上油,裡外都有油,表示主耶穌在世時內外都充滿了聖靈.乳香燃燒的香氣,充滿神的家,代表耶穌基督馨香的祭物.同時,素祭也代表我們把美好屬靈的品格,奉獻給神；基督徒若有美好的靈性,好的品格,就如馨香的素祭奉獻給神,令神滿意.(三)是聖別的香味,也是奉獻的香味,奉獻給耶和華；單單歸祂使用,歸耶和華為聖.禱告的香味,好品種的香味,分別為聖奉獻的香味.這三種香味,最令天父的心得到滿足；這三樣都與聖靈有關.我們的禱告與聖靈有關,聖靈感動我們禱告,也幫助我們禱告.我們分別為聖,與聖靈有關；聖靈潔淨我們,從罪惡中救出我們；把我們罪惡的思想習慣,言語行動都除去.我們好的品格,也是聖靈的工作；聖靈所結的果子,乃是仁愛喜樂和平......等.感謝主!聖靈的油,是分別為聖的聖膏油.
\newline
\newline
(四)教導的油　(約壹二27)約翰說:「你們從主所受的恩膏,常存在你們心裡；並不用人教訓你們,自有主的恩膏,在凡事上教訓你們；......」恩膏是主的話,加上聖靈的工作；聖靈用主的話,感動我們,造就我們,復興我們,得著我們,建立我們,那就是恩膏的作為.哦!我們需要恩膏；今日的教會需要恩膏,我們需要聽主的話,聖靈用神的話,在我們裡面發聲.我們聽見主的聲音,我們就不再聽而不聞視而無睹.我們需要恩膏,有了恩膏,一句最簡單的話,就變成最有能力的話.一句最普通的話,能進到人心裡,刺誘人心,改變人心；使人回轉,讓人復興.
\newline
\newline
(五)能力的油　聖靈的油,也是能力的油.先知撒迦利亞看見一個異象.異象中有燈台,燈台很特別,有七,七,四十九盞燈；另有四十九條的管子,連在燈台上；管子的另一端連在兩橄欖樹上.燈台表示以色列人,耶和華的旨意,要以色列人發光；發光需要能力,能力從橄欖樹的油供應而來,這油就是聖靈的油.這油有發光的能力,有榮耀神的能力；這油代表能力的油.
\newline
\newline


\newpage

\section{第47屆~港九培靈會~滕近輝~第八講}
\label{sec:1169}
\textbf{聖靈完全的能力}
\newline
\newline
連結: \href{https://www.hkbibleconference.org/session-message/view/1169}{\texttt{https://www.hkbibleconference.org/session-message/view/1169}}
\newline
\newline
\hyperref[sec:1168]{< < < PREV SERMON < < <}
~
\hyperref[sec:toc]{[返目錄]}
~
\hyperref[sec:1170]{> > > NEXT SERMON > > >}
\newline
\newline
在以色列國復興的過程中,有一件不可缺少的事,就是必須重建聖殿.他們建造聖殿過程中,曾遭遇很大的困難；波斯王下令他們停工,他們被迫於壓力之下停工了廿多年.這段悠長的年日,他們的心很悲傷失望痛苦.眼見聖殿的根基已安好,卻不得進行建造,他們的心極其傷痛,以致意志逐漸消沉.正當此時,突來一場大改變；他們振奮起來,上下同心；由消沉變為奮發.究竟有何事發生?何事使他們改變?
\newline
\newline
先知撒迦利亞看見了異象.(亞四)他明白了,領會了,接受了,這異象讓他自己先復興起來.然後把自己所看見的,所領會的,傳達給眾以色列人；於是上下同心,事情完全改觀.這個異象非常重要,可說是以色列人中興的異象,不但對當日以色列人,有重要的意義；今日對我們也有同樣的意義.如果我們領會這個異象的意義,也能叫我們從消沉中興起過來.
\newline
\newline
當撒迦利亞先知看見異象時,他像在沉睡中的人被喚醒一樣.(四1)他原來很消沉,現在卻從沉睡中醒來了；透過他,也叫以色列人從沉睡中起來.這異象的真正意義向我們顯明如下:
\newline
\newline
撒迦利亞看見這異象時,他聽見了四個呼聲；這四個呼聲非常重要而且寶貴.
\newline
\newline
第一個呼聲,(亞四6)「萬軍的耶和華說:不是倚靠勢力,不是倚靠才能；乃是倚靠我的靈,方能成事.」
\newline
\newline
第二個呼聲,(四7)「大山啊!你算甚麼呢?在所羅巴伯面前,你必成為平地.......」大山是不能越過的攔阻,橫阻在他們建殿的過程中,他們不得不停工.他們原以為無法越過的大山,在見了異象,從耶和華得了復興的呼聲之後；就知道大山在耶和華僕人面前,必成為平地.而且他必搬出一塊石頭,安在殿頂上,這塊石頭成為建造聖殿的材料.從攔阻變成建殿材料,真是奇妙的變化!這塊石頭第一個意思,是預表將來的基督.第二個意思,這塊石頭原取之於山,大山本是攔阻,現成建殿材料；這是耶和華奇妙的作為.當聖靈作工時,能把才能,智慧,權勢,都放在一邊,單單倚靠耶和華的靈；當聖靈作工時,大山就改變,不但被挪去,且變成建造神的工作的材料.這話是否會應驗?(拉六6-10)告訴我們:在波斯帝國河西地方,該地總督本來極反對以色列人建造聖殿；但當神向撒迦利亞啟示異象時,特別的事就發生了.大利烏王在位第二年,王下令河西總督,不可攔阻以色列人建殿的工作；並且要送上材料,和金錢幫助他們.河西總督原像大山,攔阻以色列人建聖殿的工作；但如今,聖靈工作,情勢全改；本為仇敵,為攔阻者,現在倒為他們預備材料和經費.
\newline
\newline
感謝主!我們需要聖靈這樣的工作.如果人心裡有聖靈的工作,他就有很明顯的表現.當他遭遇困難,打擊,意外,痛苦時；他的反應不是灰心發怨言,乃是禱告說:「主啊!求你使我看見這打擊,這攔阻,這痛苦的事變成你的祝福.」這是他的反應,聖靈在他裡面教導他,感動他,指示他在人生所遭遇的一切事上,看見神的賜福.　保羅的那根刺,讓他痛苦失望；但他在聖靈光照之下,曉得這根刺是神恩典的出口,所以他為這根刺感謝主.　我們也要為我們的刺感謝主,求主幫助我們,求聖靈在我們每個人心裡作工.當我們遭遇困難時,不為困難打倒,乃是以困難變成我們的挑戰；然後在挑戰之中領受主所賜的福.這真是奇妙寶貴的呼聲.
\newline
\newline
第三個呼聲,(10節)耶和華說:「誰藐神這日的事為小呢?」當時以色列人最多不過看見那塊石頭,但耶和華說:「不要輕看這塊石頭.」這是耶和華神施恩的開始.當先知以利亞在迦密山上,看見如手掌大小的一片雲彩,他沒有輕看,沒有藐視,他曉得神開始聽他的禱告,他馬上有了信心的動作.　神給我們的開始,我們要憑著信心去認識；神的工作常由小處開始.如果我們有聖靈所賜的信心,我們就不輕看藐視小的開始,乃是憑信心仰望主.當我們仰望,倚靠,往前走時,我們就看見小的開始,變為大的成功.如果戴德生看他當時銀行那十幾個先令的存款,他就會灰心,絕對不可開始他內地會的工作.他曉得銀行裡的十幾個先令,加上耶和華,就可以成功.他所看見的是神的手,神全能的手能夠抓住微小的開始；能夠把十幾個先令放在祂手中,使它增加千萬倍.戴德生有聖靈所賜信心,他沒有藐視聖靈最初給他的感動,這麼大的成功是從十幾個先令開始的.十幾個先令加上無限大的神,就有無限的成功.讓我們接受這個呼聲:「誰藐視今日的事為小呢?」
\newline
\newline
我願意在這裡作見證,目的為鼓勵青年們.我小時候不敢講話,有客來,我總是躲開,我的母親對我認識最清楚.為何我怕講話呢?因為我講話總是講不清楚,所以不如不講.直至我十八歲那年信了主,十九歲奉獻；當我母親知我奉獻時,她發出第一句話說:「你奉獻嗎?」她沒有繼續講,我知他言外之意是:「你不善辭令,怎能事奉主當牧師呢?」記得我讀小學時,全班同學要輪流上講台演說,輪到我時；我站在講台上,眼前一片模糊,一言發不出就下台了,直到我奉獻以後,主助我克服講話的恐懼.先父是個牧師,很有口才,前年我到台北聚會；張金蘭女大法官也來參加,她是先父為她施洗的,她來目的要聽我父親的兒子講道；聽後,她坦率對我說:「你的口才正像你父親,」我聽了很高興.先父於我十四歲時去世,我沒有留意他的口才很好；因為每當禮拜天先父講道,我坐在台下最後排的座位,和另一個少玩聖經後頁的地圖.那時我還未信主,雖然我有位好父親,但那時我心門未開,道聽不進去.直到我十八歲時才信主,十九歲奉獻,主一步步施恩,把我懼怕講話的心理改變.我母親自從我服事主以後,十八年之久為我禱告,直到她去世為止,這十八年的代禱,相信主已垂聽了.我寫第一部書「路標」第一頁寫「獻給為我禱告十八年的母親.」我奉獻以後,大多時間與我母親分開,抗戰期間分開,後來到國外讀神學,一直與她分開；她持續為我禱告十八年直到她離世為止.我講此事為要鼓勵在座的青年人,如果主呼召你奉獻給祂,不要說你沒有恩賜,如果是主呼召你,主必給你恩賜；祂能幫助你,祂能扶持你.廿五年前我開始在香港事奉主,廿四年前第一次在培靈會早堂講道,和弟兄姊妹一齊查考詩篇；到了最後兩天,我無法預備出講章.(這個秘密唯我師母知道)我絞盡腦汁,都想不到可講的,當晚我預備至凌晨二時,無論如何沒信息可傳.次早,我打電話給柯里培博士,他在當時剛開辦的浸信會神學院教書,我請求他說:「救我吧!我實在無辦法再講下去了,可否最後兩天你替我講?」他很有愛心,最後兩天柯牧師講道.當時我很難過,在主面前等候禱告,總是沒有可講的,沒有可造就弟兄姊妹的信息.但是主的恩典很奇妙,主能供應我們的需要,從那時至今,我已講了二千多篇的講章.主安排我在一個教會事奉,已經十八年了,對我很有幫助；每次都要講新的,不能重複.主的話是何等豐富,取之不盡,用之不竭.聖靈能將神的話,豐富地給我們發掘出來.「誰能藐視今日的事為小呢?」只要我們有信心,聖靈能帶領我們,進入耶和華的豐富裡面.
\newline
\newline
第四個呼聲,歡呼(7節下半)「人且大聲歡呼說,願恩惠恩惠歸與這殿.(殿或作石)」許多人為這塊石頭發出歡呼!願耶和華的恩惠,歸於這塊石頭,願耶和華的恩惠,在這塊石頭彰顯出來.一塊小小的石頭,後來成為聖殿；偉大的聖殿是從這塊石頭開始的.偉大的工作,是從聖靈最初微小的感動開始的.
\newline
\newline
(12節):在兩橄欖樹和燈台之間,有兩根枝子,橄欖樹的油,透過這兩根枝子,進到管子裡.(13節說),這兩根枝子,就是耶和華所揀選的受膏者；在當日而言,是指約書亞和所羅巴伯.第三章說,當時約書亞是大祭司,所羅巴伯是政治領袖；用現在的話來說,一位是聖職人員,一位是平信徒領袖.(有人不歡喜平信徒這名詞,不過我臨時借用)這兩根枝子,成為聖靈的通道；聖靈的油,藉著這兩種人,進入管子裡；使以色列人能發光.這裡特別題到這兩根枝子,表示這二人非常重要.今天在基督的教會中,傳道人和各部門的領袖,他們對教會有極重大的意義；如果這兩種人認識聖靈,對聖靈有信心,追求聖靈的能力；得到聖靈的潔淨和充滿,聖靈的油,就通過他們流通整個教會.這件事值得我們每個人注意,如果在座中任何一位,你是教會某部門的領袖或職員,或是傳道人.你在主手中有極重大的作用,聖靈的作為.要透過你臨到教會眾弟兄姊妹身上.我們不能輕看自己,要先在主面前俯伏,求主潔淨,復興,充滿,保守我有追求的心.然後在主手中,作主合用的器皿；如果這樣,全港九的教會有福了.我相信在座中,有很多弟兄姊妹在教會任職,求主先得著我們；我們就是燈台中的兩根枝子,主要用你用我.透過我們,把屬靈的恩典,加給眾弟兄姊妹.　第10節末句說:「見所羅巴伯手拿線鉈就歡喜.」耶和華為何歡喜?因祂看見這領袖手拿線鉈,準備動工了.他了解異象的意義,接受先知撒迦利亞的信息；伸出信心的手拿起線鉈,重新動工,耶和華看見就歡喜.弟兄姊妹!當我們受了聖靈的感動,當我們開始有新的追求,當我們伸出信心的手,拿起主所要我們做的工作,擔起主給我的托付時,我們的主就歡喜.祂為我們預備了聖靈的能力,在後面支持我們.祂全能的膀臂,從下面托住我們.到了有一天,我們要親眼看見神的大能彰顯.主耶穌說:「我豈沒有對你們說,你們必因信看見神的榮耀嗎?」
\newline
\newline
聖靈的第六個比喻
\newline
\newline
聖靈好像羔羊的七角.(啟五6)在聖經上,角是代表能力和勝利.有七種角,代表七種最大的能力.
\newline
\newline
1. 祭壇的四角　代表耶穌基督寶血的能力.耶穌基督的寶血,能洗淨我們一切的罪.基督寶血的能力,永不失效,這贖罪的能力,是寶貝最基本.
\newline
\newline
2. 香壇的角　對著約櫃的金香壇也有四角,香壇代表禱告；代表基督的代禱,也代表信徒的禱告.四角代表禱告的能力,禱告的能力最寶貴.
\newline
\newline
3. 聖殿中陳設餅台的四角　陳設餅的桌子,預表耶穌基督是生命的糧.祂的道,就是我們生命的糧.所以這些角,是代表神的道的能力.也可說就是福音的能力.「這福音本是神的大能,要救一切相信的人.」我們可以加上說,這也是主話的能力.耶和華對以賽亞說:「我的話必不徒然返回,他必然成就,我差他去所要成就的.」陳設餅台的四角,一方面表示福音的能力,同時也表示主話的能力.
\newline
\newline
4. 祭壇的角　(第四種角我再回到祭壇來講)祭壇不但是耶穌基督流寶血之處,也表明基督徒奉獻自己的愛與主之處.所以祭壇的四角,也代表基督徒奉獻自己所產生的能力.一個奉獻了的人,他的能力,是一個不肯奉獻自己的基督徒所沒有的.一個基督徒在奉獻之前,可能很軟弱；撒旦一口氣便可把他吹倒,罪惡的拳頭一打,就可把他打垮.但奉獻以後,他裡面有新的目標,新的動力,新的道路,新的方向,新的追求；這時他整個人都變了,使他從軟弱地步,進到剛強.感謝主!在教會,一個奉獻的人,能抵十個不肯奉獻的基督徒.不肯奉獻的人,起不了作用；懶洋洋地跟隨主,勉強勉強事奉主.但奉獻之後,有剛強的心志,要討主喜悅；心裡被主的愛感動,不再為自己活,乃為那死而復活的主活.這時,他表現完全不同了.
\newline
\newline
5. 第五種角(哈三4)「祂的輝煌如同日光,祂手裡射出光線；在其中藏著祂的能力.」能力原文是角,在耶和華手裡有祂的角,這是耶和華的能力,整個宇宙是耶和華創造,祂說有就有,命立就立；整個宇宙按祂旨意被創造而有；祂的計劃,最後都要完成.
\newline
\newline
6. 耶和華的寶座　啟示錄首三章是主耶穌給當時教會的書信,第四章是預言部份的開始.這部份一開始,就出現寶座與在寶座上的耶和華,何等寶貴!到了啟示錄最後,又看見寶座；以寶座開始以寶座結束.啟示錄此書,是坐在寶座上的神,用祂的能力權柄,成全祂的旨意和計劃.啟示錄是最後勝利的書,我們不看世界局勢如何,在這一切的背後,有耶和華的寶座.當先知以賽亞看見寶座時,正是王死的時候；王原來是他的後台,是他的朋友；王死了,他看見耶和華的寶座,耶和華高高坐在祂的寶座上.創造是寶座的彰顯,救贖的完成是寶座的建立；耶和華的寶座,透過了救贖之功,最後完全建立起來.到了最後,「成了」這二字是由寶座發出來的,從寶座有聲音發出來說都成了.(啟廿一3-5)自各各他山上開始,主在那裡說,「成了」,那「成了」是隱藏的.到了最後,神的寶座在整個宇宙顯明出來,那時候宇宙響應「都成了」的呼聲.
\newline
\newline
7. 羔羊頭上的七角,就是聖靈的能力.這個聖靈的能力,和上述六個聖靈的能力,都有密切關係.(因時間關係不能詳述)我讀使徒行傳,特別注意聖靈的工作.發現聖靈向初期教會有七種的得勝,弟兄姊妹因被聖靈充滿,就有這七種得勝,這七種得勝如同羔羊七角的勝利.
\newline
\newline
一,勝過懼怕(徒四8)彼得被聖靈充滿,當他被壓迫時,他發出聲音說:「我所聽見的所看見的不能不說,」聖靈充滿他,使他勝過懼怕.
\newline
\newline
二,勝過自私(徒四31-32)「禱告完了,聚會地方震動；他們就都被聖靈充滿,放膽講論神的道.那許多信的人,都是一心一意的,沒有一人說,他的東西,有一樣是自己的,都是大家公用.」聖靈的充滿,勝過他們自私的天性,都拿出為主而用.
\newline
\newline
三,勝過仇恨(徒七60)當司提反被聖靈充滿時,他跪下,為那班要用石頭打死他的人禱告.說:「主啊!不要將這罪歸於他們.」聖靈的能力,勝過仇恨,這真是大得勝!
\newline
\newline
四,勝過偏見(徒九18)保羅初次被聖靈充滿時,他的眼睛,好像有鱗掉下來.原來他的眼睛被矇蔽,聖靈光照他,如鱗脫落,眼睛明亮了.他立即有全面的改變,本來逼迫基督,現在見證基督；本來逼迫教會,現在熱心事奉耶和華.
\newline
\newline
五,勝過民族之間的隔膜(徒十44)彼得被聖靈帶領,到哥尼流家裡講道；他衝破民族界限,到外邦人家裡傳福音.正講道時,聖靈澆灌下來,以色列人和外邦人一齊蒙恩.
\newline
\newline
六,勝過嫉妒(徒十一24)安提阿教會裡面,主僕人巴拿巴被聖靈充滿,勝過了愛自己,愛作領袖的心.本來他自己作領袖,但經聖靈充滿,他想到保羅比自己更有恩賜,更有聖靈的能力,更得神的使用.所以他把自己撇在一邊,費時去找保羅來作教會領袖.他忘記自己,勝過嫉妒.
\newline
\newline
七,勝過撒旦(徒十三9)保羅初交出外佈道,到了居比路島上,方伯很歡喜聽道,但被行法術的以呂馬攔阻.這時保羅被聖靈充滿,就斥責以呂馬；使他眼睛暫時失明.說:「你混亂主的正道還不止住嗎?」他勝過了撒旦.
\newline
\newline
以上七次的得勝,都是聖靈充滿的得勝.聖靈的角,羔羊頭上的七角,聖靈完全的能力,被差遣到普天下去.
\newline
\newline
感謝主!今天為我們預備能力,如果你認識這些能力,追求能力,你也可以得著.求主給我們有追求的心.讓我們仰望主,憑著信心領受聖靈的能力,憑著信心領受聖靈的充滿；神能夠大大的恩待我們,神也能藉著我們大大恩待祂的眾教會.
\newline
\newline


\newpage

\section{第47屆~港九培靈會~滕近輝~第九講}
\label{sec:1170}
\textbf{三個關於聖靈的比喻}
\newline
\newline
連結: \href{https://www.hkbibleconference.org/session-message/view/1170}{\texttt{https://www.hkbibleconference.org/session-message/view/1170}}
\newline
\newline
\hyperref[sec:1169]{< < < PREV SERMON < < <}
~
\hyperref[sec:toc]{[返目錄]}
~
\hyperref[sec:1171]{> > > NEXT SERMON > > >}
\newline
\newline
靈往那裡去,四活物也往那裡去.」這裡「靈」是反映甚麼呢?兄弟曾參考過多處名註釋家的註釋,綜合而說,這裡所說的「靈」是指著神的靈；神的靈往那裡去,四活物也往那裡去.
\newline
\newline
四活物如何行動呢?(19節)「活物行走,輪也在旁邊行走,活物從地上升,輪也都上升.」聖靈好像輪推動四活物,聖靈和輪連在一起.「四活物的靈在輪中.」(20節)聖靈,輪和四活物聖靈三者密切連在一起；由此我們看見了很重要的真理.四活物的靈和聖靈合而在一起,不再分開了,二者有一致的行動,行動的方式就是那輪；聖靈是動力,如果教會要往前走,需要聖靈的輪推動.　從這段聖經,我們看見與聖靈有關的五個特點:
\newline
\newline
一,「活物上升,輪也在活物旁邊上升.」(20節)輪有上升的行動,從(20-21)我們看見輪有平面和上下兩種活動.聖靈的輪帶動我們,也有這兩種活動,先帶我們上去,讓我們到寶座那裡支取神的恩典；然後降下,有平面的行動,把神的恩典表顯在人與人之中的工作.平面的事奉需要上下的動作來支持；沒有上下的動作,平面的行動就缺少力量.
\newline
\newline
記得有一次我到葛培理的總部參觀,一進大門,轉彎上樓,梯口面對一幅葛培理低頭禱告的大畫像.他每次講完佈道會,必邀請決志信主的人走到前面,他就低頭禱告.當他禱告時,就是聖靈的輪推動的時候；把許多人推到救恩台前,叫他們決志接受主耶穌作他們的救主.葛培理工作了廿五年,至今還有莫大的吸引力；他每到一地佈道,都有數萬人聽道.有兩次機會,我參加他主領的佈道會,我仔細觀察,得到三個結論:(一)許多人赴會,目的是帶領未信主的親友聽福音.(二)為著看見眾多人們,歸向主而快樂,因為平時很少看見人肯決志信主.(三)是聖靈的工作.葛培理在廿五年之中,經常出現於美國各州,每週也在電視播映,美國居民,聽慣了他講道,也看慣了他的容貌.按理對他已沒多大興趣,但是廿五年來,他的聲望毫不減少.相反地,更看見他的工作,不斷在擴展中.這是聖靈的工作!主保守他謙卑,儆醒,向上看；所以經長久時間,仍然滿有能力.由此我們看見聖靈奇妙的作為,他禱告時,叫我們想到輪上下的活動；輪上下活動,平面的功效就顯出.
\newline
\newline
二,「輪中套輪」(16節)原文輪是單數,如果按單數解釋,即內外各有一輪.我相信這裡的意義不受文法的限制.這輪可以有很多輪子,輪中套輪,有各種樣式；各種大小,各種方向的輪子配合在一起,然後發生作用.讓外面的大輪子,有奇妙的成就.聖靈的工作實在如此,祂將各種恩賜,賜給基督的眾教會；叫我們藉著這些恩賜,好像輪子一樣,在聖靈大力推動下顯出我們的作用來.在這些輪子中,有兩種輪子是最易灰心的.小輪子不易被人發現,也不引人注意；以至自憐,灰心.覺得別輪龐大,自己渺小,因而灰心退後.還有一種隱藏的輪,從外面看不見；自覺無人稱讚,無人注意,功效也無人欣賞.如果你是這兩種輪之一,就願你從主的話語上,得著安慰吧!當這許多種輪子配合在一起時,就要讓大的,總的輪子發生大的作用,其中沒有任何一部份,是不需用的；也沒有一個小輪子,是不重要的,各有其作用.如果神安排我們在一處,無論你的工作是被人重視或被人輕看；你的崗位是明顯的或隱藏的,只要是主的安排,是聖靈的引導,我們就要存滿足喜樂的心事奉,肯這樣行的人有福了,這樣的工作,有永遠的果效.求主幫助在座每位弟兄姊妹,如果主安排你在某教會某部門工作；隱藏也好,明顯也好,人看為大或小的工作都好.只要為主而作,是作給主看的,目的要滿足主的心就好了.當你達到這目標時,你就可以心滿意足,主也心滿意足；稱讚你是又忠心又良善的僕人,你就可以得著神為你預備的福樂了.
\newline
\newline
三,「行走的時候,向四方都能直行,並不掉轉」(17節)四活物有四個面,每面旁邊有一輪.四個面朝四個方向直行,那就是各朝東南西北而行.照人的邏輯而言,是不可能的；四活物的面各向一方,四活物豈不是散開了嗎?但是在這異象,四個面都直往前行；每個面旁邊有一輪子,每個輪子直往前行,到達神要他去的地方.按常理來說是不通的,但這是屬靈的通；意思乃說,接受恩賜,我們就照主所賜的恩賜去用.聖靈暗中為我們配合恰當,只要照主的恩賜勇往前行；這樣去行,教會就有進步了.各肢體都有很恰當的配合,是聖靈智慧的配合；不必我們操心,我們的責任,就是善於運用主所給的恩賜.只要我們忠於恩賜,一直往前,神就成全祂整個的計劃.
\newline
\newline
四,「這活物往來奔走,好像電光一閃.」(14節)活物行走很快,因輪子走得很往,好像閃電,很快就到了目的地.當聖靈的輪推動,我們要乘輪而去,好像乘風駛帆；聖靈的輪轉動,快如閃電,眨眼就過去.我們必須及時把握機會,讓聖靈的輪把我們帶去,機會一失,我們將要後悔.　教會歷史給我們看見,聖靈的風吹動,聖靈的輪轉動的時候,有些先知先覺者,他們知道聖靈的引導,抓住機會.結果,他們成為那時代的先知,他們都是各時代被神所用的人.我們要乘著聖靈工作的機會而工作,這樣就可以事半功倍.
\newline
\newline
五,「輪的形狀和顏色,好像水蒼玉.」(16節)水蒼玉和壁玉的顏色差不多,都是屬天的顏色.意思即聖靈之輪的行動,是按屬天的原則,聖靈永不背乎屬天的原則.聖靈的法則,非屬血氣的法則,非人為的法則；聖靈按著屬天的法則行事,所以必有屬靈的成功.讓我鄭重說:「表面的成功,未必是成功；人眼中的成功,未必是主眼中的成功.」只要我們認定睛確的目標,將會得主的稱讚,會得到屬靈的成功,得到天上榮耀中的冠冕.這樣的人,有時會被人誤會,受人批評,但不必怕；只要認定目標,只求滿足主的心,所作的按照聖靈的法則.如此惟願足矣,可以歡喜快樂.有一天在榮耀中要聽見主的聲音說:「你這又良善又忠心的僕人,可以來享受主人為你預備的快樂.」
\newline
\newline
聖靈的工作好像眼睛一樣(啟五6)
\newline
\newline
聖靈不但是角也是眼,聖靈如眼.剛才我們看過聖靈的輪,「輪輞周圍滿有眼睛」(結一18)有了輪,加上眼睛,二者在一起更加寶貝.如果輪子沒有眼睛,可能帶你到錯誤的地方；輪有了眼睛,看得很清楚,能帶你行走在正確的道路上.眼代表智慧,聖靈的智慧地祂的輪子裡,在祂的引導裡,在祂的工作裡.聖靈從不作盲目的工作,祂所作的一切,都有最高的智慧.所以我們可以放心,只要是出於聖靈的感動,我們可以接受,可以順服.當聖靈攔阻保羅在小亞細亞工作時,保羅不明白；聖靈兩次攔阻保羅的工作(徒十六6-8),保羅卻認為那兩個地方可去工作.當時,他不明白為何聖靈兩次攔阻他,後來聖靈給他看見馬其頓的異象；他順服這個異象,轉了工作的方向到歐洲去.今日我們回顧之下,這轉變關係多麼重大.聖靈禁止他並沒有錯,聖靈為他關了這門,卻為他開了那門.如果我們進入聖靈所開的門,前途必是光明的；必定走在一條正確而成功的道路上
\newline
\newline
相信在今天,聖靈有眼睛的輪子,帶領中國教會走在這條差傳的道路上.差傳的意義,不在乎遠方或近處；重要的意義,是每間教會把心胸擴展,把負擔範圍擴展,屬靈眼光擴展.由本堂擴展到以外的地區,努力工作的範圍；由本堂小範圍伸展進入更大的範圍,伸展,擴展是教會增長的一件最重要的事.如果一個堂開始作差傳的工作,立即會發現弟兄姊妹,為傳福音的奉獻增加.繼之,發現弟兄姊妹的奉獻,所支持的傳道人；從本堂進行,到其他堂的傳道人.原來僅支持一個傳道人,當發動差傳的工作,就支持兩個傳道人.原來一個堂,逐漸成為兩個堂三個堂.聖靈已在推動,聖靈的風正吹起；聖靈的輪開始行動.讓我們抓住機會,接受感動,接受挑戰,全世界的華人都能醒起來,把傳福音的工作由小範圍擴展出去.每個堂都有發展力,聖靈能幫助我們,既然這事出於聖靈；我們順服,就會發現聖靈的同工.有了聖靈的同工,我們樣樣都有,沒有缺少.
\newline
\newline
(亞三9)給我們看見有塊特別的石頭.耶和華說:「我要親自彫刻這石頭,並要在一日之間,除掉這地的罪孽.」可見這石頭,是神所用來除掉罪孽的,這石頭很明顯是預表基督.這石頭上有七眼當然是預表基督的靈,基督的靈就是神的靈.四章7節:「這七眼乃是耶和華的眼睛遍察全地；見所羅巴伯手拿線鉈,就歡喜.」聖靈的眼遍察全地,發現有一人(就是所羅巴伯)憑信心,把工作的線鉈拿起來,祂就歡喜.哦!如果你我被聖靈的眼睛發現,有聖靈所賜的信心,勇氣,膽量,就可以拿起我們認為拿不起來的東西；作我們認為不可能的事.聖靈的眼睛看見因此歡喜,所羅巴伯作成了祂的工作.　感謝主!只要我們知道這件事是出於聖靈,讓聖靈的眼睛看見,就沒有不能成就的事了!
\newline
\newline
感謝主!聖靈充滿了以利沙,使沉下的斧頭也可以漂起來.斧頭下沉本是自然現象,但是感謝主,以利沙有屬靈的眼光,深信耶和華能叫下沉之斧,也可以浮上來.因他有這樣的信心,加上聖靈的工作,就算下沉了的斧頭,也會漂浮上來!照樣,聖靈也會把新的信心賜給我們.萬軍之耶和華,透過祂的七眼觀看,要看看有沒有信心的人,肯起來復興祂的工作,求聖靈給我們一個新的挑戰.讓我們呼求說:「主啊!我信你能叫下沉的斧頭浮起,也能叫我復興過來；叫下沉的傳道人興起,讓我這冷淡的心興起來,得勝過來.」
\newline
\newline
石頭的七眼另有一個意義,石頭有眼是活的,如果沒有眼就沒有生命.我們的主是活石,這活石就是基督教會的第一塊活石.基督的教會乃是建立在這活石上,所以是活的教會.一塊塊有生命的石頭,有七眼的石頭,都建立在基督的頭塊石頭上.這樣的教會滿了聖靈的眼睛,滿了聖靈的智慧和生命；這樣的教會是活的,動的,不是死的.聖靈工作時,教會變成活石,聖靈在裡面蓬勃生長,這樣的教會必然發展；死的教會不能生長,死的骸骨必先得生命,然後才成為強大的軍隊.
\newline
\newline
聖靈好像鴿子(太三16)
\newline
\newline
鴿子是和平的象徵,聖靈這位鴿子造成和平.
\newline
\newline
一,使人與神和好(弗二18)「藉著祂被一個聖靈所感,得以進到父面前.」我們被聖靈感動時,就進到父面前,聖靈使我們一天天靠近父神,把神人之間的隔膜除去.聖靈用說不出來的嘆息,為我們禱告；要把這些攔阻隔膜都除去,讓我們能直接與主面對面.祂不勉強我們,祂用祂的愛等候:「說不出來的嘆息」這是聖靈的聲韻,表示祂沉重的負擔.因為我們與主之間有隔膜,但祂不勉強我們；祂感動我們,不與我們繼續相爭.主的愛也是如此,主耶穌站在門外叩門,祂只說:「看哪!聽我聲音又開門的我就進去.」意思即是說有人聽見了就開門,有些人聽見卻不開門的.我們的主在門外忍耐等候,等我們聽見祂的聲音.啟示錄三章說:「聖靈向教會所說的話,凡有耳的,就應當聽.」主說:「看哪!我站在門外叩門,凡聽見我的聲音又開門的；我就進到他裡面；我與他,他與我一同坐席.」　有一次我參加一個婚禮,我早到十分鐘,坐在禮堂.我發現有幅主釘十字架的圖畫,畫法很特別；畫家從畫的前上方角度來畫,顯得耶穌的身體在十字架上懸掛.全身的重量掛在雙手上,這是我初次見到如此特別的筆法；引起我的注意.正看時,忽然我心深處有所感覺,我體會到主耶穌在十字架上經歷的痛苦,我的心受了感動.弟兄姊妹們!主的愛是何等長闊高深!難道這愛不能叩開你的心門嗎?
\newline
\newline
二,使基督藉著聖靈,將自己無瑕無疵獻給神.
\newline
\newline
「耶穌基督藉著永遠的靈,將自己無瑕無疵獻給神.」(來九14)這是聖靈在耶穌基督身上的工作,使祂無瑕無疵獻給美好的祭物.
\newline
\newline
三,使我們成為基督美麗的新婦.(歌一15)鴿子是美麗的,聖靈如同鴿子極其美麗.祂的工作,要使屬主的人變成美麗的,使我們成為無瑕無疵的祭物,把我們自己獻給主.聖靈改變我們屬靈的相貌,變成主美麗的樣式,把我們變為美麗可愛的新婦,獻給基督.
\newline
\newline
「你們想經上所說是徒然的嗎?神所賜住在我們裡面的靈,是變愛至於嫉妒嗎?」(雅四5)這經節另有一翻譯,為許多解經家所採納的:「神所賜住在我們裡面的靈是愛我們一直到嫉妒的地步嗎?」聖靈為了基督愛我們,為了基督而在我們身上有嫉妒的愛.意思即聖靈要得著我們如新婦一樣獻給基督,這是聖靈的願望.但當我們離開主,冷淡倒退時,愛世界多於愛主時；聖靈就有嫉妒的感覺.這嫉妒是代表最切的痛苦,弟兄姊妹!聖靈的愛裡有痛苦,因為這愛是最深的愛；聖靈愛基督也愛我們,祂願意我們成為基督的新婦奉獻給基督.如果我們了解聖靈的心意,我們就必體貼聖靈,將我們的愛奉獻給主；不讓任何的事任何的人,減少對主的愛.
\newline
\newline
香港有一間學院的團契,有一個不冷不熱的基督徒.他雖有恩賜,但他還未重生得救,他表面事奉,其實心不在焉.有一年受難節,團契演話劇,有一幕演主耶穌釘十字架；他們做了一個十字架,推舉這職員扮演主耶穌.在預演時他們把燈光放暗,用細鐵絲把他的手足綁在十字架上,下邊做了兩塊凸起的木頭墊起支持他的腳跟.他覺得很辛苦.回家之後,他想到今天被掛在十字架,不過是預演而已,尚且非常痛苦,何況當日主耶穌真的被釘十字架,是何等痛苦啊!他心裡受感動,他跪下,接受主,認罪,流淚,把他的心獻給主.從那時開始,他完全改變,做個很愛主的基督徒,他的事奉不同了；基督的愛在他裡面,使他愛主.
\newline
\newline
弟兄姊妹!聖靈好像鴿子,美麗的鴿子.把我們作成基督美麗的新婦,無瑕無疵的奉獻給基督；有一天,我們要與所愛的主永遠同在.
\newline
\newline


\newpage

\section{第47屆~港九培靈會~滕近輝~第十講}
\label{sec:1171}
\textbf{聖靈的充滿}
\newline
\newline
連結: \href{https://www.hkbibleconference.org/session-message/view/1171}{\texttt{https://www.hkbibleconference.org/session-message/view/1171}}
\newline
\newline
\hyperref[sec:1170]{< < < PREV SERMON < < <}
~
\hyperref[sec:toc]{[返目錄]}
~
\hyperref[sec:1148]{> > > NEXT SERMON > > >}
\newline
\newline
以西結書 1:12-1:21
\newline
\begin{longtable}{cl}
\hline
\hline
章節 & 經文 (和合本修訂版)\\
\hline
1:12 & \begin{tabularx}{0.7\textwidth}{X} 他們各自往前直行。靈往哪裡去，他們就往哪裡去，行走時並不轉彎。 \end{tabularx} \\ \\ \relax
1:13 & \begin{tabularx}{0.7\textwidth}{X} 至於四活物的形像，就如燒著火炭的形狀，又如火把的形狀。有火在四活物中間來回移動，這火有光輝，從火中發出閃電。 \end{tabularx} \\ \\ \relax
1:14 & \begin{tabularx}{0.7\textwidth}{X} 這些活物往來奔走，好像電光一閃。 \end{tabularx} \\ \\ \relax
1:15 & \begin{tabularx}{0.7\textwidth}{X} 我觀看活物，看哪，有四張臉的活物旁邊各有一個輪子在地上。 \end{tabularx} \\ \\ \relax
1:16 & \begin{tabularx}{0.7\textwidth}{X} 輪子的形狀結構好像耀眼的水蒼玉。四輪都是一個樣式，形狀結構好像輪中套輪。 \end{tabularx} \\ \\ \relax
1:17 & \begin{tabularx}{0.7\textwidth}{X} 輪子行走的時候，向四方直行，行走時並不轉彎。 \end{tabularx} \\ \\ \relax
1:18 & \begin{tabularx}{0.7\textwidth}{X} 至於輪圈，高而可畏；四個輪圈周圍佈滿眼睛。 \end{tabularx} \\ \\ \relax
1:19 & \begin{tabularx}{0.7\textwidth}{X} 活物行走，輪子也在旁邊行走；活物離地上升，輪子也上升。 \end{tabularx} \\ \\ \relax
1:20 & \begin{tabularx}{0.7\textwidth}{X} 靈往哪裡去，活物就往哪裡去；輪子在活物旁邊上升，因為活物的靈在輪中。 \end{tabularx} \\ \\ \relax
1:21 & \begin{tabularx}{0.7\textwidth}{X} 活物行走，輪子也行走；活物站住，輪子也站住；活物離地上升，輪子也在旁邊上升，因為活物的靈在輪中。 \end{tabularx} \\ \\
[1ex]
\hline
\hline
\end{longtable}
今天晚上的渴望
\newline
\newline
我看見在座中有很多青年人,我禱告的焦點,就是求主得著我們中間,每一個青年.在主的心目中,你們何等寶貝的；主要得著你們,而後要使用你們.祂要潔淨你們,祂潔淨的,祂要充滿；祂所充滿的,祂要使用.如果弟兄姊妹都有一顆心,讓主來使用；真實把自己放在祭壇上,成為奉獻的人,那能力和作用是偉大的!聖靈的能力,可以透過我們的心,我們的口,我們的腳,發生極大的作用；我相信主的心,是要得著今天在座的每一個人.你是主重價買來的,主在你身上有雙重的主權；祂創造你,是你的主,祂救贖你,更是你的主.但願祂的主權,在我們身上不致落空,不致歸於徒然.
\newline
\newline
讓我們把權歸給神的,歸給神；讓聖靈的手,抓住我們每一個人,在我們中間沒一人漏掉.主的手是慈愛的手,當我們的手,被祂的手抓住時；天上的恩典,透過祂的手,臨到我們身上；被聖靈的手抓住的人有福了.有首詩歌:「雖有別人已蒙應允,使我亦沾恩.」英文的詩歌:「不要越過我去.」既有許多人蒙主揀選,既有許多人被主的手抓住.我們就該禱告說:「主啊!既然這恩典臨到別人；主啊!不要越過我,願你施恩的手,也能抓住我.」今天晚上我們安息在主愛中,讓主愛充滿我們的心；不需要人講甚麼,願主的愛,藉聖靈澆灌在我們心裡.感謝讚美主!聖靈將神的愛,澆灌在我們心裡；那喜樂和滿足,是無法可以形容的.今晚願聖靈的能力,托住我們,主能力的手在我們以下；全能的膀臂,把我們托住,我們就心滿意足了.我們把自己的生命,和一切所有的,都擺在主恩手中.請記得!祭壇在主手中,當我把自己擺在祭壇上,也就是擺在主的恩手中.
\newline
\newline
我原來希望能講完聖經上關於聖靈的十五個比喻,到昨晚為止我講了九個；今晚我要講聖靈的充滿,再要提到三個聖靈的比喻.這樣,一共講了十二個比喻.
\newline
\newline
聖靈的充滿,是神給祂的教會最大的恩典.聖靈降臨,就是教會成立的時候,也是教會擴展的時候.感謝天父!把最大的恩賜給祂的教會；如果沒有聖靈降臨,就沒有教會；如果沒有聖靈充滿,福音就不能廣傳.若單靠當時的門徒,要把福音傳遍世界；那是絕對不可能.因為門徒都是軟弱的人,他們是漁夫無學問的小民；他們沒有地位,沒有權勢,怎能成功呢?感謝主!當他們軟弱,逃走,四散,絕望的時候；神的恩賜,就從天澆灌下來.按照舊約的應許,在神所預定的日期；聖靈就在五旬節降臨充滿門徒.把他們完全改變,我們確實知道,他們的改變不是感情作用；更沒有人為因素,完全是藉著從上而來能力.主對門徒沒有信心,但對聖靈的能力,卻有無限的信心.主說:「但聖靈降臨在你們身上,你們就必得著能力；並要在耶路撒冷,猶太全地,和撒瑪利亞,直到地極,作我的見證.」聖靈降臨,福音傳開,從最近處直到最遠處.都可以看見有福音使者,被主愛激勵去傳福音.主耶穌說:「這天國的福音,要傳遍天下.」主的話必要應驗,今天我們親眼看見主預言的應驗了.感謝讚美主!不是人的作為,乃是聖靈的能力!
\newline
\newline
(徒四31)說當時門徒──「禱告完了,聚會的地方震動,他們就都被聖靈充滿.」門徒被聖靈充滿,連地也震動起來,表示驚天動地的事發生了.當主耶穌被釘十字架時,地大震動,表示驚天動地的事發生了.主復活時,也有地震,也表示驚天動地的事發生了.保羅到了腓立比,那是向全世界福音廣傳的第一站,是全世界普傳福音的開始；當時有地震發生在腓立比的監獄裡.這就是說,馬其頓的異象正施行的時候,就有地震發生.神作事有祂的道理,有祂的原因；聖靈充滿時,就有力量把福音傳遍天下,所以有地震發生.可見主耶穌何等注重聖靈的能力.
\newline
\newline
我查考使徒行傳,發現初期教會在耶路撒冷有十五件大事發生,都與聖靈的能力有關.這十五件大事,造成初期教會的歷史；在這歷史中,我們看見聖靈的運行.有人說得好說得對,「使徒行傳,就是聖靈行傳.」一點也不錯.我們可以看見使徒行傳中,聖靈的能力,如何藉著那班有信心的人,如何藉著那些奉獻的人,作成他們原來所做不到的大事.
\newline
\newline
以下我要提到四種不同的聖靈充滿:
\newline
\newline
(1)在工作上的聖靈充滿　當主要把一樣工作托負給人的時候,主就把聖靈的能力充滿他,使他能完成工作；工作完了,聖靈就離開他.
\newline
\newline
(2)生活上的聖靈充滿　(加二22-23)當一個基督徒被聖靈充滿時,他就能結出九樣聖靈的果子.這種充滿不憑感覺,是四種聖靈充滿中之最高級的.這叫我們看見神工作的原則,最高級的不憑感覺,單憑信心；憑感覺越多,他就越低.當然低也是寶貴,但最高的是生活上的聖靈充滿,故結出九樣靈果.我再說,不憑感覺,乃憑信心接受；讓聖靈在一個已經潔淨了,同時也有信心的心裡,結出聖靈九樣寶貴的果子.
\newline
\newline
(3)感動上的聖靈充滿　(徒十三52)保羅和巴拿巴在一個地方受逼迫,被趕出城.當時,他們被聖靈充滿,滿心喜樂；這是感情上的聖靈充滿.有時我們禱告,覺得有特別甜蜜的經歷,心裡特別快樂,甚至覺得主特別同在,這些都是感情上的聖靈充滿.
\newline
\newline
(4)聖靈的澆灌　使徒行傳中,有四次題到聖靈澆灌,四次都有特別的事發生；可能會有說方言的,可能沒有說方言的.
\newline
\newline
有很多基督徒,對聖靈充滿沒有認識,所以就發生困難.如果你知道聖靈充滿,至少有四種,我相信你的問題就迎刃而解了.
\newline
\newline
怎樣能得勝靈充滿?
\newline
\newline
第一條件　相信.相信主已為我預備了,聖靈充滿的恩典；相信聖經上已有聖靈充滿的應許,我就可以得著.
\newline
\newline
第二條件　潔淨,讓聖靈潔淨你,把你所記得清,想得起的,所有罪過都認清楚,一概傾倒盡淨.
\newline
\newline
第三條件　祈求,神把最好的給求祂的人.求是表示你要,如果你不要,那就沒辦法勉強給你.禱告是心裡渴慕發生的聲音:「主啊!我渴慕你的恩典,渴慕你自己；渴慕聖靈的充滿,你所應許的,我看重,我認為寶貴.」　當我們求的時候,主就悅納我們這樣渴慕的態度,主不把恩典給不要的人.有的人說要懇切禱告.有人用理智思想,認為是感情主義,未求以前神已經知道我們要求甚麼,禱告不過是多餘的.我們禱告若是默默的,冷冷淡淡的,那就是代表無所謂的態度；得到也好,得不到也罷.懇切禱告是信心堅定的,必定要得；認為神的應許可靠,把神的恩典看為最寶貴.懇切的禱告,是由心裡發出來的；所以主悅納,主垂聽,主特別應允這樣的禱告.有時撒旦用一些理論,把聖經直接的吩咐抹煞遮蓋,使人覺得應用自己的頭腦理智,豈知理智多,轉幾轉就迷糊了.世上的哲學,興起一個便打倒老的一個.現在的否定從前的,沒有一個哲學能長久持續；這證明我們的理智不可靠.神安定在天的話,有永遠的智慧；讓我們憑著單純直接的信心來接受.到了有一天,我們會發現祂真正寶貴的理由.每一個憑著信心,接受神話語的人,遲早會看出神話語的真理；千千萬萬人,都有這樣經驗的見證.弟兄姊妹們!我們最大的智慧,就是憑著單純的信心,接受神的話.「深哉!神豐富的智慧和知識,祂的判斷,何其難測,祂的蹤跡,何其難尋,誰知道主的心,誰作過祂的謀士呢?」(羅十33-34)千萬不要自以為聰明,要憑信心俯伏在主面前,我們求就得著.
\newline
\newline
第四條件又回到信心,從信心開始,潔淨,祈求,然後憑信心來保持；當我們求的時候,神一定要把生活上的聖靈充滿賜給我們.至於感情上的充滿,澆灌的充滿,都應該交給主；主若喜悅,就給我們充滿；若不是主的旨意,我們就得不著.交在神權柄的手中,若神給你有感情的激動你就有；若神不給,你也可以安安靜靜,有聖靈充滿的能力.憑著信心接受主的應許,憑著信心來領取；保持在誠實善良,不但從信心得著,繼續又繼續的藉著信心來保守.
\newline
\newline
按兄弟的經驗,主曾給我兩次澆灌的充滿.第一次在多年前,那時我住在九龍衙前圍道.我正在繙譯「禱告出來的能力」此書時,心裡有很奇妙的充滿.我一面繙譯,眼淚一面滴下來,擬稿上有很多我的眼淚,我跪下禱告,充滿從裡面出來.一禱告便幾個鐘頭,裡面的喜樂和眼淚沖了出來.這樣一連保持了兩三個月,不是我要努力保持,我不曉得是怎樣一回事；我從未有此經驗,我惟有感謝讚美主.第二次是在去年九月,當我在日本領夏令會時；我向主求充滿,主給了我,心裡也有一種澆灌的感覺.使我感到很希奇的,同時候的聚會中,主給我在工作上從未曾有的力量.有一次我講了很平常的道之後,有一個廿多歲的弟兄站起來認罪.他講了幾句話,就仆倒在地,繼續哀哭認罪.我覺得那次講的,並非特別的信息,但聖靈使他知罪；使他認罪,使他復興.
\newline
\newline
聖靈充滿是非常真實的事,非人造的,是主為我們預備的.我們憑著信心,潔淨自己,祈求,得著以後,憑著信心繼續接受.
\newline
\newline
聖靈充滿有四個角度,藉此四角度,我講關於聖靈的三個比喻.
\newline
\newline
充滿的第一角度──「強度」當聖靈降臨時,好像一陣強風吹過.我曾題過,舊約原文,氣息,風,靈,原出一字.聖靈給我們得到生命的氣息,加強時就成了風.風是空氣流動而形成的,空氣靜止時,沒有作用,沒有力量.流動時即成了風,強風颶風力量很大,大風表示生命氣息的強化.我們不單有生命的氣息,不單有聖靈所給我們的生命,而且生命強化.主說:「風隨著意思吹......卻不曉得從那裡來,往那裡去；凡從聖靈生的,也是如此.」(約三8)這比喻表示聖靈重生工作的奧秘.我們不知聖靈怎樣使我們重生,但我們知道一個事實；重生的人自己知道,別人也知道他已重生.從他身上看見重生的記號.
\newline
\newline
聖靈降臨,充滿門徒,如同一陣強風吹過,門徒不但得到聖靈所賜的生命,且得了強有力的生命,超過他們原有的任何力量.　以西結卅七章記載,遍滿平原的骷骨,以西結奉耶和華之命發聲傳信息.有風吹來,進到骸骨,他們就起來,變成強大的軍隊.這是復興,是穌靈充滿；從死人骸骨變成強大的軍隊,這是得生命,也是被聖靈充滿.　我們不但得到生命,求主使我們的生命能起作用；有作用的生命,是能結果子的生命.　在醫院裡有病人,也有護士；在教會裡,有屬靈的病人,也有屬靈的護士.所有力量都用在照顧病人身上,成了醫院式的教會.軍隊式的教會,能為主工作,為主爭戰,為主服務,為主結果子；這樣的教會是聖靈充滿的教會,這叫做強化.
\newline
\newline
第二角度──「廣度」　聖靈充滿,表示聖靈進到人生的每個部份；被聖靈管治,被聖靈變化.發生作用,使我們的人生,完全被聖靈摸到.
\newline
\newline
聖靈如水　以西結卅七章記有一條生命的河,那就是聖靈充滿的生命河.聖靈好像充滿在生命河裡的水,這水從聖殿門檻流出；越流越廣越遠,達到四圍的地方.約但河很窄小,當我親臨約但河,簡直不相信是聖經上的約但河；但不要輕看它,每年發水時,河水漲過兩岸.寬度增加,河水瀰漫,澆灌滋潤；樹木結果子,各類的魚,很豐富的生產.這是廣度,聖靈的作用進到人生的每一部份,充滿我們的容量,把重擔攔阻除去.有的人自認為被聖靈充滿,那是某一程度的充滿.聖靈繼續工作時,會發現一些攔阻的傢俬,把那些傢俬除去,便得到進一步的充滿.再得潔淨,就有更進一步的充滿.聖滿的充滿,在祂的幅度上,廣度上,有各樣的增加.
\newline
\newline
第三角度──「深度」　這條河越流越廣,越深.以西結四七章量河的尺度,不是用尺用竿而是用人來量.初時水到他的踝子骨,往前走一千肘水漸深,到了膝.又往前走一千肘,到他的腰部.再往前走一千肘,水浸過他的頭；再也不能往前走,水把他漂起來.豐盛的生命,聖靈充滿,用人作為量度.最後整個人被漂起,腳不能站定,不能定自己的方向.河流的方向,就是他的方向,這時他不由自主；原來的自己被勝過了,被舉起來了,這就是聖靈充滿的表現.當人充滿老自己時,完全屬於血氣的味道,這證明沒有被聖靈充滿.一個人被聖靈充滿時,也可能倒退,回到屬血氣的地步；但還需要繼續充滿,他被充滿時,老自己動作不自由了.水到踝子骨時,行動還可以自由,水到了膝蓋,行動開始不方便.水到腰部,要緩慢地才能決定自己的腳步.但是越來越困難,老自己被攔阻,受對付.直到最後,自己失掉,被舉起來,被漂起,被沖去.聖靈的方向,才是他的方向,這就到了順服的階段.被充滿的人,不再說:「我要如何如何,我的意思如何,」乃是說:「不要照我的意思,願主旨成.」這是聖靈充滿的記號.以西結四七章:河的兩邊滿了樹,樹上每月結新果子,甚至乾旱季節也不例外；這是聖靈充滿的表現,不在乎環境,乃在乎裡面的力量.在南美洲有一種水蜘蛛,自己做個胞,裡面滿了空氣,然後下到水裡,浮沉自如.靠著胞內的空氣生活,身在水中卻呼吸陸上的空氣,在水中牠另有天地.我們也是如此,聖靈的充滿,使我們裡面另有天地,不受外界的影響,能在罪惡的環境中,過聖潔的生活.　聖經記載,約瑟是個有聖靈在裡面的人,也是個聖靈充滿的人；能勝過主母的試探,因他裡面另有力量,每月都結新果子.
\newline
\newline
第四角度──「濃度」　主說祂是葡萄樹,我們是枝子,枝子連在樹上,汁漿就進到枝子裡面.記得兄弟有一次參加香港英文的培靈會,有一位講員,他說葡萄樹的汁漿是代表聖靈.他用很多理由證明,他的解釋我接受了.父神是栽培的人,主耶穌是葡萄樹,我們是枝子.聖靈是葡萄樹的汁漿,聖靈把葡萄樹的生命運輸到枝子裡面；所以葡萄汁所到之處,就有葡萄樹的生命運行在那裡.
\newline
\newline
聖靈像風,聖靈像水,聖靈是葡萄樹的汁,當我們與主接連時,聖靈流通在我們裡面；這時,我們就有力量.主耶穌說:「你們離開我就不能作甚麼.」我們連在樹上時,我們在主裡面,主在我們裡面,這時就有能力的交流,就能結果子.基督的生命,進到我們裡面,我們就有結果子的能力.主說:「你們若常在我裡面,我的話也常在你們裡面.」(約十五7)主的話是真理,這就是加上真理的交流.主說:「你們要常在我的愛裡,正如我在父的愛裡一樣.」(9節)這是愛的交流,主的愛在我裡面,就引起我愛主的心.主說:「叫我的喜樂存在你們心裡,並叫你們的喜樂可以滿足.」這是喜樂的交流.
\newline
\newline
能力的交流,真理的交流,愛的交流,喜樂的交流；聖靈在裡面,把這些交流成為可能.
\newline
\newline
聖靈充滿重要的條件　(結四七)這裡我們看見,風成為生命之河.這裡特別題到,這條河繞經祭壇,從祭壇南邊往下流.聖靈充滿必須經過祭壇.祭壇的意義:(1)祭壇是主耶穌為我們死的地方.(2)是主的愛激勵我們奉獻自己的地方.經過祭壇,就是經過奉獻；如果我們要得勝靈充滿,如果得到之後,要想保持聖靈的充滿,我們必須奉獻.把自己無條件無保留的放在祭壇上.　兩年前有一位青年來電話告訴我,他不久將得到香港大學理科博士學位；但主感動他奉獻,他徵求我的意見,他應進那間神學院.一聞之下,我快樂極了.直到今天,他在一個神學院出版了一部書名「劍橋七傑」,我相信這書的出現,是聖靈的帶領.他向全世界的青年挑戰；七個劍橋大學的畢業生,個個都是出人頭地的,聖靈感動他們向遠方傳福音.
\newline
\newline
我們肯付代價嗎?中國人能不能付真的代價?如果你想得特別的權利,如果主給你特別權利,特別智慧,才幹,你認為那是屬於你的嗎?你認為奉獻給主可惜嗎?「焉知你得了王后的位分,不是為現今的機會嗎?」讓我奉主名向你們挑戰.主的呼召「今日誰樂意把自己獻給耶和華呢?」當聖靈的眼睛看我們的時候,祂了解我們的一切；祂明白我們的一切,祂愛我們.祂願意呼召我們,得著我們,讓你奉獻自己一生傳福音,一生事奉主；全時間的事奉,你願意嗎?主若呼召你,我希望你說:「我在這裡請差遣我.」
\newline
\newline


\newpage

\section{第47屆~港九培靈會~鮑會園~第一講}
\label{sec:1148}
\textbf{苦難與安慰}
\newline
\newline
連結: \href{https://www.hkbibleconference.org/session-message/view/1148}{\texttt{https://www.hkbibleconference.org/session-message/view/1148}}
\newline
\newline
\hyperref[sec:1171]{< < < PREV SERMON < < <}
~
\hyperref[sec:toc]{[返目錄]}
~
\hyperref[sec:1150]{> > > NEXT SERMON > > >}
\newline
\newline
哥林多後書 1:1-1:11
\newline
\begin{longtable}{cl}
\hline
\hline
章節 & 經文 (和合本修訂版)\\
\hline
1:1 & \begin{tabularx}{0.7\textwidth}{X} 奉神旨意作基督耶穌使徒的保羅和弟兄提摩太，寫信給在哥林多神的教會和全亞該亞的眾聖徒。 \end{tabularx} \\ \\ \relax
1:2 & \begin{tabularx}{0.7\textwidth}{X} 願恩惠、平安從我們的父神和主耶穌基督歸給你們！ \end{tabularx} \\ \\ \relax
1:3 & \begin{tabularx}{0.7\textwidth}{X} 願頌讚歸於神—我們主耶穌基督的父；他是發慈悲的父，賜各樣安慰的神。 \end{tabularx} \\ \\ \relax
1:4 & \begin{tabularx}{0.7\textwidth}{X} 我們在一切患難中，他安慰我們，使我們能用神所賜的安慰去安慰那些遭各樣患難的人。 \end{tabularx} \\ \\ \relax
1:5 & \begin{tabularx}{0.7\textwidth}{X} 正如我們跟基督同受許多苦楚，我們也靠基督得許多安慰。 \end{tabularx} \\ \\ \relax
1:6 & \begin{tabularx}{0.7\textwidth}{X} 如果我們受患難，那是為使你們得安慰，得拯救；如果我們得安慰，那也是為使你們得安慰，這安慰能使你們忍受我們所受同樣的苦楚。 \end{tabularx} \\ \\ \relax
1:7 & \begin{tabularx}{0.7\textwidth}{X} 我們為你們所存的盼望是確定的，因為知道你們分擔了我們的痛苦，也要分享我們的安慰。 \end{tabularx} \\ \\ \relax
1:8 & \begin{tabularx}{0.7\textwidth}{X} 弟兄們，我們不要你們不知道，我們從前在亞細亞遭遇苦難，因受到無法忍受的壓力，甚至連活命的指望都沒有了。 \end{tabularx} \\ \\ \relax
1:9 & \begin{tabularx}{0.7\textwidth}{X} 自己心裡也斷定是必死無疑，這是要使我們不依靠自己，只依靠使死人復活的神。 \end{tabularx} \\ \\ \relax
1:10 & \begin{tabularx}{0.7\textwidth}{X} 他曾救我們脫離那極大的死亡，他要繼續救我們，而且我們指望他將來還要救我們。 \end{tabularx} \\ \\ \relax
1:11 & \begin{tabularx}{0.7\textwidth}{X} 你們也要一同用祈禱來幫助我們，好使許多人為我們感恩，因著他們許多的禱告，我們獲得了恩賜。 \end{tabularx} \\ \\
[1ex]
\hline
\hline
\end{longtable}


\newpage

\section{第47屆~港九培靈會~鮑會園~第二講}
\label{sec:1150}
\textbf{生命的記號}
\newline
\newline
連結: \href{https://www.hkbibleconference.org/session-message/view/1150}{\texttt{https://www.hkbibleconference.org/session-message/view/1150}}
\newline
\newline
\hyperref[sec:1148]{< < < PREV SERMON < < <}
~
\hyperref[sec:toc]{[返目錄]}
~
\hyperref[sec:1152]{> > > NEXT SERMON > > >}
\newline
\newline
哥林多後書 1:12-1:12
\newline
\begin{longtable}{cl}
\hline
\hline
章節 & 經文 (和合本修訂版)\\
\hline
1:12 & \begin{tabularx}{0.7\textwidth}{X} 我們所誇的是：我們在世為人，特別是跟你們的關係，是憑著神所賜的坦率和真誠，不是靠人的聰明，而是靠神的恩惠；這是我們的良心可以作證的。 \end{tabularx} \\ \\
[1ex]
\hline
\hline
\end{longtable}
上述經文,保羅講到他生命中幾方面的態度,也是屬靈生命的記號.有這申命記號的,才能被神使用.因此,我們需要這記號來衡量我們的生活.
\newline
\newline
一,他的生活堅定
\newline
\newline
他的生活,並非反覆不定,似是而非的.他在神面前,知道應該怎樣行,不隨便改變.一個生活無目標的人,是被環境催逼來走；但一個生活有目標的人,不隨環境改變,他會堅定自己所走的方向.保羅應許哥林多信徒說,他會先到他們那裡去,但事實上他是先到馬其頓去,以後才到哥林多去.因此就有人批評他,說他反覆不定,似是而非.保羅解釋他行動的原因,他沒有為自己辯護,只是解釋.他的理論提到極高的標準,即是基督的標準.以神為出發點.意思是說我講福音,是按神的標準講；不但是福音,連我的生活也是按神的標準來過.
\newline
\newline
「豈是反覆不定麼」這句話在中文很難翻譯,故看不出它的真義.這是一句很重要的話,保羅說明他以前所做的一切,都是是的；神的應許,不論有多少,在基督都是是的；所以藉著祂,也都是實在的.(林後一20)原文是說:「神是應許不論有多少,在基督裡都是是的,「實在」小字原文作「阿們」；此字難在構造上,原文之意是「所以藉著祂(主)我們都向神說阿們」；前面的「是」是神的,神的話藉主是「是」的,一句不落空.我們要怎樣做呢?我們藉著耶穌,每一個命令都向神說阿們.意即:神給我們什麼吩咐,我們都完全順服.
\newline
\newline
保羅去何處傳道,阿們；要做什麼,阿們；走什麼路,阿們；自己沒有影響神的旨意,完全順服.神說什麼,阿們!不因環境改變神的旨意,不因自己的喜好改變神的旨意,不因別人的意見改變對神的態度.這件事保羅實在不容易做,特別是哥林多人對他有成見.神的話是阿們.完全順服神是屬靈生命的記號,非常寶貴.這個生命的記號是我們所不可缺少的.今天,你的生命面臨一個選擇,做什麼事,神要你追隨保羅有相同的態度,完全順服神的旨意,對神說阿們.
\newline
\newline
二,被神堅定的生命
\newline
\newline
「堅固」一詞保羅用得很美,如海上的浮標,在水面上受水流影響,但它有力量拉住根基,很牢固.我們的生命也是如此,我們生活在這個潮流中,受沖激,在教會中做事,受四方的力量沖激；沒有根基,便隨潮流走.保羅說不可以那樣,屬神的人當有堅固根基.
\newline
\newline
神怎樣堅固我們呢?
\newline
\newline
A,神在基督裡用膏來膏我們
\newline
\newline
「膏」指聖靈的膏抹,每一個屬基督的人,不論他知道真理有多少,生命長進有多大,都沒有關係；只要是重生得救的人,就是被聖靈所膏的.聖靈入住在他心中,「人若沒有基督的靈,就不是屬基督的.」(羅八9)
\newline
\newline
B,神用印印了我們
\newline
\newline
現代的人很少用印,古代的人很常用；遇有重要公文,信封上打上火漆印；叫人無法開信,這火漆印是人的保證.保羅說神用印印了我們,所以每一個基督徒有主的印印了,這記號是屬於主的.別人的印不行.今日,我們被神的印印過的人,能否在我們的生活中,把基督在我們身上的印彰顯出來呢?
\newline
\newline
C,聖靈在我們心裡作憑據
\newline
\newline
聖靈在我們心裡作憑據,在以弗所書,羅馬書中都有提及.在以弗所書中「憑據」譯成「質」,意即財物去當作抵押品.聖靈在我們裡面是質,等於神給我們所許保證,將來必要完全得著.
\newline
\newline
三,生命中應有靈敏的感覺
\newline
\newline
在我們生命中應有靈敏的感覺,特別是愛的真正同情心的感覺.(林後二1-4)保羅一再的對哥林多教會的信徒說,他第一封嚴厲的信責備犯罪的人,叫他心裡很難過.正如他說的:「我先前心裡難過痛苦,多多的流淚,寫信給你們；不是叫你們憂愁,乃是叫你們知道我格外的疼愛你們.」(林後二4)
\newline
\newline
一個在屬靈上有領受的人,有屬靈亮光的人；有時應該要去指責一些人的錯,這個責任不能逃避.但沒有人願意這樣做,這實在是不容易做的事.我們看見弟兄姊妹跌倒,聖靈感動我們去提醒他.我們的心有沒有為這個罪,如刀砍一樣痛苦呢?如果沒有為人的罪憂傷痛苦,那麼你就沒有資格去責備人的罪.心裡沒有同情的感覺,則同樣沒有資格去責備人的罪.心裡沒有同情的感覺,則同樣沒有資格去責備人.「同情」是很美的字,意即兩人一同感覺痛苦.心裡沒有與對方同感痛苦,則沒有資格去責備人的罪.
\newline
\newline
四,有真正饒恕人的心
\newline
\newline
保羅寫給哥林多教會,吩咐他們要把犯罪的人趕出教會去,他們受了刑罰後又如何呢?保羅說:「這樣的人,受了眾人的責罰,也就夠了.倒不如赦免他,安慰他,免得他憂愁太過,甚至沉淪了.」(林後二5-6)一個人犯了罪,若是他真正悔改了,就要真正的饒恕他.若是只指責他,沒給他悔改的機會,可能他因此沉淪了；這樣做在屬靈上沒有一點價值.可是在今日的教會中,卻不易饒恕人；因此教會中就有機會給撒旦侵入來,好像一個瘡口,容易給細菌侵入一樣.所以我們對於一個肯真正悔改的人,我們有責任去饒恕他,安慰他.
\newline
\newline
五,在生命中發出基督馨香之氣
\newline
\newline
發出基督的馨香之氣,有兩方面的表現:
\newline
\newline
A,生命的見證
\newline
\newline
一個真正順服基督的人,有生活的見證.要付代價才能發出有力的見證,保羅說:「感謝神!常帥領我們在基督裡誇勝,並藉著我們在各處顯揚那因認識基督而有的香氣.」(林後二14)「誇勝」是很美的字,也是個很生動的字,中文不易表達,像在羅馬的社會,將軍打仗勝利歸來,政府為他誇勝的功勞,列隊行車；後面的車中載著,被鎖鍊綁著的俘虜和戰利品.將軍非常榮耀,威武；被鎖鍊綁著的人滿了羞辱.保羅將自己比作是被鎖鍊綁著的俘虜,基督是將軍.他自己甘心樂意跟從主,顯出基督的勝利,自己的卑微.主的恩典越多,我們就更多發出基督的香氣來.
\newline
\newline
B,對神的話的態度
\newline
\newline
十七節說:「我們不像那許多人,為利混亂神的道；乃是由於誠實,由於神；在神的面前,憑著基督講道.」「混亂」也是很美的字,像是賣酒的人,不誠實,為了想多賺一些錢；便在酒中加入了醋,變成了假酒.這些是混亂的意思.中文譯得很好,「為利混亂神的道」.意即為「自己的好處」,故意錯解神的道,錯用神的道.你們有見過這樣的傳道人嗎?或者他們為了討人的喜歡,便將神的道理故意解釋得好聽一點,叫聽的人歡心.
\newline
\newline
神的道是真理,所以我們要忠實的解釋神的道,按著聖經的方法和啟示來解釋,才能發出基督的馨香之氣；叫四週的人能夠看見基督的榮美,與及屬靈的申命記號.
\newline
\newline


\newpage

\section{第47屆~港九培靈會~鮑會園~第三講}
\label{sec:1152}
\textbf{新約的執事}
\newline
\newline
連結: \href{https://www.hkbibleconference.org/session-message/view/1152}{\texttt{https://www.hkbibleconference.org/session-message/view/1152}}
\newline
\newline
\hyperref[sec:1150]{< < < PREV SERMON < < <}
~
\hyperref[sec:toc]{[返目錄]}
~
\hyperref[sec:1154]{> > > NEXT SERMON > > >}
\newline
\newline
哥林多後書 3:1-4:18
\newline
\begin{longtable}{cl}
\hline
\hline
章節 & 經文 (和合本修訂版)\\
\hline
3:1 & \begin{tabularx}{0.7\textwidth}{X} 難道我們又開始推薦自己嗎？難道我們像某些人那樣要用人的推薦信介紹給你們，或用你們的推薦信給人嗎？ \end{tabularx} \\ \\ \relax
3:2 & \begin{tabularx}{0.7\textwidth}{X} 你們就是我們的推薦信，寫在我們心裡，被眾人所知道、所誦讀的， \end{tabularx} \\ \\ \relax
3:3 & \begin{tabularx}{0.7\textwidth}{X} 而你們顯明自己是基督的書信，藉著我們寫成的。不是用墨寫的，而是用永生神的靈寫的；不是寫在石版上，而是寫在心版上的。 \end{tabularx} \\ \\ \relax
3:4 & \begin{tabularx}{0.7\textwidth}{X} 我們藉著基督才對神有這樣的信心。 \end{tabularx} \\ \\ \relax
3:5 & \begin{tabularx}{0.7\textwidth}{X} 並不是我們憑自己配做甚麼事，我們之所以配做是出於神； \end{tabularx} \\ \\ \relax
3:6 & \begin{tabularx}{0.7\textwidth}{X} 他使我們能配作新約的執事，不是文字上的約，而是聖靈的約；因為文字使人死，聖靈能使人活。 \end{tabularx} \\ \\ \relax
3:7 & \begin{tabularx}{0.7\textwidth}{X} 那用字刻在石頭上屬死的事奉尚且有榮光，以致以色列人因摩西臉上那逐漸褪色的榮光不能定睛看他的臉， \end{tabularx} \\ \\ \relax
3:8 & \begin{tabularx}{0.7\textwidth}{X} 那屬聖靈的事奉不是更有榮光嗎？ \end{tabularx} \\ \\ \relax
3:9 & \begin{tabularx}{0.7\textwidth}{X} 若是那使人定罪的事奉有榮光，那使人稱義的事奉的榮光就越發大了。 \end{tabularx} \\ \\ \relax
3:10 & \begin{tabularx}{0.7\textwidth}{X} 那從前有榮光的，因這更大的榮光，就算不得有榮光了； \end{tabularx} \\ \\ \relax
3:11 & \begin{tabularx}{0.7\textwidth}{X} 若是那逐漸褪色的有榮光，這長存的就更有榮光了。 \end{tabularx} \\ \\ \relax
3:12 & \begin{tabularx}{0.7\textwidth}{X} 既然我們有這樣的盼望，就大有膽量， \end{tabularx} \\ \\ \relax
3:13 & \begin{tabularx}{0.7\textwidth}{X} 不像摩西將面紗蒙在臉上，使以色列人不能定睛看到那逐漸褪色的榮光的結局。 \end{tabularx} \\ \\ \relax
3:14 & \begin{tabularx}{0.7\textwidth}{X} 但他們的心地剛硬，直到今日誦讀舊約的時候，這同樣的面紗還沒有揭去；因為這面紗在基督裡才被廢去。 \end{tabularx} \\ \\ \relax
3:15 & \begin{tabularx}{0.7\textwidth}{X} 然而直到今日，每逢誦讀摩西書的時候，面紗還在他們心上。 \end{tabularx} \\ \\ \relax
3:16 & \begin{tabularx}{0.7\textwidth}{X} 但他們的心何時歸向主，面紗就何時除去。 \end{tabularx} \\ \\ \relax
3:17 & \begin{tabularx}{0.7\textwidth}{X} 主就是那靈；主的靈在哪裡，哪裡就有自由。 \end{tabularx} \\ \\ \relax
3:18 & \begin{tabularx}{0.7\textwidth}{X} 既然我們眾人以揭去面紗的臉得以看見主的榮光，好像從鏡子裡返照，就變成了與主有同樣的形像，榮上加榮，如同從主的靈變成的。 \end{tabularx} \\ \\ \relax
4:1 & \begin{tabularx}{0.7\textwidth}{X} 所以，既然我們蒙憐憫受了這事奉的責任，就不喪膽， \end{tabularx} \\ \\ \relax
4:2 & \begin{tabularx}{0.7\textwidth}{X} 反而把那些暗昧可恥的事棄絕了，不行詭詐，不曲解神的道，只將真理顯揚出來，好在神面前把自己推薦給各人的良心。 \end{tabularx} \\ \\ \relax
4:3 & \begin{tabularx}{0.7\textwidth}{X} 即使我們的福音被遮蔽，那只是對滅亡的人遮蔽。 \end{tabularx} \\ \\ \relax
4:4 & \begin{tabularx}{0.7\textwidth}{X} 這些不信的人被這世界的神明弄瞎了心眼，使他們看不見基督榮耀的福音。基督本是神的像。 \end{tabularx} \\ \\ \relax
4:5 & \begin{tabularx}{0.7\textwidth}{X} 我們不是傳自己，而是傳耶穌基督為主，並且自己因耶穌作你們的僕人。 \end{tabularx} \\ \\ \relax
4:6 & \begin{tabularx}{0.7\textwidth}{X} 那吩咐光從黑暗裡照出來的神已經照在我們心裡，使我們知道神榮耀的光顯在耶穌基督的臉上。 \end{tabularx} \\ \\ \relax
4:7 & \begin{tabularx}{0.7\textwidth}{X} 我們有這寶貝放在瓦器裡，為要顯明這莫大的能力是出於神，不是出於我們。 \end{tabularx} \\ \\ \relax
4:8 & \begin{tabularx}{0.7\textwidth}{X} 我們處處受困，卻不被捆住；內心困擾，卻沒有絕望； \end{tabularx} \\ \\ \relax
4:9 & \begin{tabularx}{0.7\textwidth}{X} 遭受迫害，卻不被撇棄；擊倒在地，卻不致滅亡。 \end{tabularx} \\ \\ \relax
4:10 & \begin{tabularx}{0.7\textwidth}{X} 我們身上常帶著耶穌的死，使耶穌的生也在我們身上顯明。 \end{tabularx} \\ \\ \relax
4:11 & \begin{tabularx}{0.7\textwidth}{X} 因為我們這活著的人常為耶穌被置於死地，使耶穌的生命在我們這必死的人身上顯明出來。 \end{tabularx} \\ \\ \relax
4:12 & \begin{tabularx}{0.7\textwidth}{X} 這樣看來，死是在我們身上運作，生卻在你們身上運作。 \end{tabularx} \\ \\ \relax
4:13 & \begin{tabularx}{0.7\textwidth}{X} 但我們既然有從同一位靈而來的信心，正如經上記著：「我信，故我說話」，我們也信，所以也說話； \end{tabularx} \\ \\ \relax
4:14 & \begin{tabularx}{0.7\textwidth}{X} 因為知道，那使主耶穌復活的也必使我們與耶穌一同復活，並且使我們與你們一起站在他面前。 \end{tabularx} \\ \\ \relax
4:15 & \begin{tabularx}{0.7\textwidth}{X} 凡事都是為了你們，好使恩惠既藉著更多的人而加增，感恩也格外顯多，好歸榮耀給神。 \end{tabularx} \\ \\ \relax
4:16 & \begin{tabularx}{0.7\textwidth}{X} 所以，我們不喪膽。雖然我們外在的人日漸朽壞，內在的人卻日日更新。 \end{tabularx} \\ \\ \relax
4:17 & \begin{tabularx}{0.7\textwidth}{X} 我們這短暫而輕微的苦楚要為我們成就極重、無比、永遠的榮耀。 \end{tabularx} \\ \\ \relax
4:18 & \begin{tabularx}{0.7\textwidth}{X} 因為我們不是顧念看得見的，而是顧念看不見的；原來看得見的是暫時的，看不見的才是永遠的。 \end{tabularx} \\ \\
[1ex]
\hline
\hline
\end{longtable}
保羅說:「神叫我們承擔新約的執事,不是按字句,條文,乃是憑著精意,因為那字句是叫人死,精意是叫人活.」(參林後三6)這句話常被人錯用,沒有屬靈眼光的人,只看字句；但在靈裡有深刻領受的人,就能明白屬靈的意義.所以我們要追求,在靈裡長進.字面之意與精意是不同的,精意叫人活,這是一般人解經的標準.但保羅之意則與之相反.他是指新約的教訓和舊約的教訓來說的；舊約是憑字句,新約是憑精意.保羅說:「你們明顯是基督的信,藉著我們修成的；不是用墨寫的,乃是用永生的靈寫的；不是寫在石版上,乃是寫在心版上.」(林後三3)
\newline
\newline
起初,神賜下律法寫在石版上,交給摩西.但當摩西下山之時,看見百姓拜金牛,摩西便怒擲法版,因而破碎了.之後,摩西在神面前為百姓禱告,再回到神面前,於是神第二次再賜下法版.這事詳記在出埃及記,卅一章和卅二章裡頭.所以第一次和第二法的律法內容是不相同的,第一次以正義的標準,百姓就接受不了.按著這個審判,沒有一人能夠存在.但是神另賜的律法,顯明祂的公義和慈愛,正如聖經的記載；「耶和華說,日子將到,我要與以色列家和猶大家,另立新約.」(耶卅一31)按舊約的標準審判,則沒有一人能存活.
\newline
\newline
新約的執事,並非遵守神的命令,律法,乃憑神的恩典.新約時代的基督徒,要照神的指示,作新約的執事.
\newline
\newline
保羅如何論到作新約的執事呢?
\newline
\newline
(一)作新約執事的資格
\newline
\newline
我們憑己意不能遵守神的律法.保羅所說的執事,不是單講作使徒的執事,他乃是說新約的每一位基督徒都是執事；換言之,每一個屬主的人都是新約的執事.林後三章十七至十八節論到猶太人何時歸向主,便何時除去帕子.今日我們也是如此,好像從鏡子裡返照,就變成主的形狀,榮上加榮,如同從主的靈變成的.
\newline
\newline
「帕子」在何處?有人說,聖經被帕子遮住了,保羅說不對；帕子在我們心靈中,問題在我們身子,帕子除去便可看見主的榮光.但要有幾件事發生在我們身上,我們才有資格作新約的執事.
\newline
\newline
一,主的靈在那裡,那裡就有自由
\newline
\newline
在主的靈裡,才能得真正的自由.耶穌曾說:「你們若常常遵守我的道,就真是我的門徒.你們必曉得真理,真理必叫你們得以自由.」(約八31-32)有真理,就有自由,真理是聖靈的啟示.如何知道得自由?聖經譯作「知道」,「認識」,意思是指體驗過和經驗過的知道.譬如說,你若沒有嘗過番石榴的味道,無論怎樣說,你也不會明白的.
\newline
\newline
如你知道在神的真理中得自由?「自由」在聖經中用得很多,所以常被人誤會.保羅說:在肉體情慾之下,真自由並非想做什麼就做什麼.聖經說有幾件事神不能做的,就是背乎自己,失信和犯罪.但在神裡面有真自由.按著本性,我們沒有自由；立志為善由得我,行出來卻由不得我.正如保羅所說的:「我真是苦阿!」是因為肉體受罪捆綁,屬靈生活沒有自由.若是真在主裡,體驗到真理,便是真自由了.
\newline
\newline
二,我們眾人好像從鏡子裡返照
\newline
\newline
「從鏡子裡返照主的榮光」,此話很難翻譯.可作:不是主返照祂的榮光,我們屬主的人；似一面鏡子,返照主的榮光.主的榮光照在我們身上,我們在生活上,就把主的榮光返射出來.基督徒重要的責任是,叫世人在我們身上看見主的榮光,我們的生活是一面鏡子.但我們的生活如何呢?是否像哈哈鏡那樣,返照出來的生活,是變形的呢?生活完全走了樣,形狀一點都不像主.這樣的人沒有資格作新約的執事.
\newline
\newline
三,生命上不斷長進
\newline
\newline
生命上不斷長進,才能變成主的形狀,榮上加榮,如同從主的靈變成的.從一形狀變成另一形狀,是榮耀的.基督徒屬靈生命要有長進.羅馬書第八章有論及生命長進的問題,在神面前領受了真理而去遵行.有了進步,神便會將更深的真理啟示給我們；一步一步遵行,才能有長進.不遵行,領受了也沒有長進.
\newline
\newline
(二)新約執事的工作
\newline
\newline
林後四章五節說:「我們原不是傳自己,乃是傳基督耶穌為主,並且自己因耶穌作你們的僕人.」新約執事最主要的工作,有以下兩方面:
\newline
\newline
一,為要傳揚耶穌基督為主
\newline
\newline
作新約的執事,要傳揚基督耶穌為主.為何要傳揚基督耶穌為主呢?因為世人的眼睛被世界的神弄瞎了,而看不見靈裡的真正需要.世界上有享受,名譽,地位,成就等,但我們要叫人看見,主比這一切更好.將福音傳開,要用最好的方法去傳.憑人的方法不能勝任,要用屬靈的方法去作見證.不是靠口才,學問,經驗,乃是倚靠聖靈工作方能成事.
\newline
\newline
二,對工作的態度
\newline
\newline
A,有寶貝放在瓦器裡顯出神的大能
\newline
\newline
保羅說,我盡力為主工作,最好的也不過像瓦器一樣不值錢.好酒放在瓦缸裡,人買的是酒,很昂貴；不是買瓦缸,瓦缸也不值錢.傳福音不是靠自己的恩賜,能力,乃是把自己放在主的手中.寶貝貴重,瓦器不重要；酒雖然貴,但若沒有缸來盛是不行的.雖然瓦缸沒價值,但有寶貝裝在裡頭,便顯得重要了.
\newline
\newline
每位基督徒為主作工,是瓦器放在主手中.神使用的使徒保羅,馬丁路德,是他們肯完完全全順服在神前.在教會歷史中,有很多這樣的人,他們很有本事,但是神把他們放在一邊,沒有使用他們,因為他們不肯順服神.神沒有揀選耶路撒冷,有高深學問的大祭司,乃揀選一班打魚的人,因為他們肯順服神.
\newline
\newline
B,做神的工作要付代價
\newline
\newline
我們讀哥林多後書,講到很多苦難的事(林後四8-10),四面受敵,遭逼迫,帶著主的死.保羅為何會這樣做呢?因為他樂意做新約的執事.作主的工作,必須付上真正的代價,不付代價,生命就沒有意義.
\newline
\newline
作新約執事所走的道路,須付代價,有痛苦,犧牲,也許還有死亡在前頭,在這個挑戰之下,得著力量去作新約的執事,其力量之來源是「不喪膽；外體雖毀壞,內心卻一天新似一天.我們這至暫至輕的苦楚,要為我們成就極重無比的榮耀.」保羅忍受至暫至輕的苦楚,他沒有把眼睛放在至暫至輕的事上,乃是放在永遠的事上,看重極重無比的榮耀,勝過眼前至暫至輕的苦楚.他以永遠的榮耀作背景,所以他看現在的苦楚；是至暫至輕的,毫不足介意的.背景很重要,沒背景本身受苦難,也就沒有意義了.有一本雜誌,刊載捕鯨魚的事,所招請的明星極矮,只有三尺高；拍片之時,將人放大一倍,即六尺高,那麼鯨魚便顯得很巨大了.我們看永遠的事也是如此.我們有永遠的目標,譬如說,我們今日的損失是永遠的榮耀之一部份.我們若是看見永遠的榮耀,這至暫至輕的苦楚,就不當作是什麼一回事了.
\newline
\newline
我們要盡本份去作新約的執事,讓主的旨意成就.
\newline
\newline


\newpage

\section{第47屆~港九培靈會~鮑會園~第四講}
\label{sec:1154}
\textbf{基督的使者}
\newline
\newline
連結: \href{https://www.hkbibleconference.org/session-message/view/1154}{\texttt{https://www.hkbibleconference.org/session-message/view/1154}}
\newline
\newline
\hyperref[sec:1152]{< < < PREV SERMON < < <}
~
\hyperref[sec:toc]{[返目錄]}
~
\hyperref[sec:1156]{> > > NEXT SERMON > > >}
\newline
\newline
哥林多後書 5:1-5:21
\newline
\begin{longtable}{cl}
\hline
\hline
章節 & 經文 (和合本修訂版)\\
\hline
5:1 & \begin{tabularx}{0.7\textwidth}{X} 因為我們知道，我們這地上的帳篷若拆毀了，我們將有神所造的居所，不是人手所造的，而是在天上永存的。 \end{tabularx} \\ \\ \relax
5:2 & \begin{tabularx}{0.7\textwidth}{X} 我們在這帳篷裡嘆息，渴望得到那從天上來的居所，好像穿上衣服； \end{tabularx} \\ \\ \relax
5:3 & \begin{tabularx}{0.7\textwidth}{X} 倘若脫下也不至於赤身了。 \end{tabularx} \\ \\ \relax
5:4 & \begin{tabularx}{0.7\textwidth}{X} 其實，我們在這帳篷裡的人勞苦嘆息，並不是願意脫下地上的帳篷，而是願意穿上天上的居所，好使這必死的被生命吞滅了。 \end{tabularx} \\ \\ \relax
5:5 & \begin{tabularx}{0.7\textwidth}{X} 那為我們安排這事的是神，他賜給我們聖靈作憑據。 \end{tabularx} \\ \\ \relax
5:6 & \begin{tabularx}{0.7\textwidth}{X} 所以，我們總是勇敢的，並且知道，只要我們住在這身體內就是離開了主。 \end{tabularx} \\ \\ \relax
5:7 & \begin{tabularx}{0.7\textwidth}{X} 因為我們行事為人是憑著信心，不是憑著眼見。 \end{tabularx} \\ \\ \relax
5:8 & \begin{tabularx}{0.7\textwidth}{X} 我們勇敢，更情願離開身體，與主同住。 \end{tabularx} \\ \\ \relax
5:9 & \begin{tabularx}{0.7\textwidth}{X} 所以，無論是住在身內或住在身外，我們都立了志向要得主的喜悅。 \end{tabularx} \\ \\ \relax
5:10 & \begin{tabularx}{0.7\textwidth}{X} 因為我們眾人必須站在基督審判臺前受審，為使各人按著本身所行的，或善或惡受報。 \end{tabularx} \\ \\ \relax
5:11 & \begin{tabularx}{0.7\textwidth}{X} 既然我們知道主是可畏的，就勸導人；但是神是認識我們的，我盼望你們的良心也認識我們。 \end{tabularx} \\ \\ \relax
5:12 & \begin{tabularx}{0.7\textwidth}{X} 我們不是向你們再推薦自己，而是要讓你們有誇耀我們的機會，使你們好面對那憑外貌、不憑內心誇耀的人。 \end{tabularx} \\ \\ \relax
5:13 & \begin{tabularx}{0.7\textwidth}{X} 如果我們癲狂，是為神；如果我們清醒，是為你們。 \end{tabularx} \\ \\ \relax
5:14 & \begin{tabularx}{0.7\textwidth}{X} 原來基督的愛激勵我們；因我們這樣斷定，一人既替眾人死，眾人就都死了； \end{tabularx} \\ \\ \relax
5:15 & \begin{tabularx}{0.7\textwidth}{X} 並且他替眾人死，是叫那些活著的人不再為自己活，乃為替他們死而復活的主活。 \end{tabularx} \\ \\ \relax
5:16 & \begin{tabularx}{0.7\textwidth}{X} 所以，從今以後，我們不再按照人的看法來認識人，縱使我們曾經按照人的看法認識基督，如今卻不再這樣認識他了。 \end{tabularx} \\ \\ \relax
5:17 & \begin{tabularx}{0.7\textwidth}{X} 所以，若有人在基督裡，他就是新造的人，舊事已過，都變成新的了。 \end{tabularx} \\ \\ \relax
5:18 & \begin{tabularx}{0.7\textwidth}{X} 一切都是出於神；他藉著基督使我們與他和好，又將勸人與他和好的使命賜給我們。 \end{tabularx} \\ \\ \relax
5:19 & \begin{tabularx}{0.7\textwidth}{X} 這就是：神在基督裡使世人與自己和好，不將他們的過犯歸到他們身上，並且將這和好的信息託付了我們。 \end{tabularx} \\ \\ \relax
5:20 & \begin{tabularx}{0.7\textwidth}{X} 所以，我們作基督的特使，就好像神藉我們勸你們一般。我們替基督求你們，與神和好吧！ \end{tabularx} \\ \\ \relax
5:21 & \begin{tabularx}{0.7\textwidth}{X} 神使那無罪的，替我們成為罪，好使我們在他裡面成為神的義。 \end{tabularx} \\ \\
[1ex]
\hline
\hline
\end{longtable}
人的生命被神創造,基本目的非為在地上生活,人的生活非屬地的.地上的生活,只是一個訓練.神把人安放在地上,目的是為鍛鍊他們；使他們到有一天,能夠適合神所預備的居住之地.地上的生活是暫短的,神把人放在地上有一個目的,有一個使命,就是作基督的使者.使者的責任,去到另一國代表自己的國家工作.初到一國,把本國的意見,好處,表現在外國人面前；他講的話,並非自己的意思,乃代表國家講話；他講話的內容,自己不能負責.使者的立場,自己不能負責；因為國家已制定好,他的責任就是把國家的立場,表現出來!如何表現?什麼態度?叫人同情抑或令人反感?叫人尊敬或不尊敬?大使要負責任.與外國的關係友好或疏遠,就要看大使的工作好與壞了.大使的工作在該國,只是暫時性的,不是永住在該國的,待工作做完後,再派駐往別國工作.
\newline
\newline
保羅說,屬主的人是基督的使者,在地上代表我們的主.我們不屬這個世界,我們在地上重要的使命,責任,至終是使者.有一天我們要離開世界,回到我們的國家去.一個真正的愛國者,渴望回到自己的國家去.「使者」的原意正是如此.保羅說:我們在地上的帳棚若拆毀了,在天上有永存的房屋,為我們永久的居所.肉體完了,並非生命完了,乃是大使的工作完成了.結束那裡的工作,回到自己的國家去.保羅在此所說的的「房屋」與約翰福音十四章中所說的「地方」是同一意義的.主說:「我去原是為你們預備地方去.我若去為你們預備了地方,就必再來接你們到我那裡去.」
\newline
\newline
我們等待永久的居所,方向認清了,目標認準了,因此我們對地上的生活就有不同的觀念,地上的生活是暫時的,是一種客旅的生活.我們要如何為天家預備安排呢?我們要計算一下每天用多少時間為天家的生活安排.心在何處?看重什麼?保羅提醒我們要看重屬天的事,知道它是絕對真實的,支取力量,作使者的工作.基督的使者得力量的兩個因素:
\newline
\newline
一,眾人要在基督台前受審判
\newline
\newline
今天我們若不努力作使者的工作,將來有一天要受審判,站在基督台前每一件事要交待出來,我們的工作如何,私生活如何,辦事如何,都要在基督台前交帳.「在基督台前」一語,並非指將來得救或滅亡受審判；非為不信者預備的,乃為基督徒預備的.保羅在哥林多前書三章和十五章裡頭,講到基督徒在地上的工作和生活要受審判.工作和生活的好壞,是決定我們在主前領受賞賜的標準.
\newline
\newline
保羅在哥多前書三章裡,講到基督建造的根基是金,銀,寶石的；還有草,木,禾楷的呢?我們受審判時,生活和工作要經火的試驗.金,銀,寶石經火試驗後就顯得更美,可以保存到永遠；草,木,禾楷一經火便被燒毀了；只剩下生命得救,生活和工作燒毀了,不能永存.我們所行的,或善或惡要受報應,這善惡不是按人的標準定的.不信主的人,他的愛在神前沒有價值；違背神旨意的愛,是沒有價值的.神用的標準,是用永世的標準,祂是用這個標準來衡量我們的.我們的生命是在基督裡建造根基嗎?今天我們是用甚麼材料在生命的根基上建造呢?是金,銀,寶石,抑草,木,禾楷呢?所以,我們要在使者的工作上忠心.
\newline
\newline
二,基督的愛激勵我們
\newline
\newline
基督的愛激勵我們要對祂忠心.激勵推動我們,掌管我們.「激勵」一詞很美麗,有幾個意思:
\newline
\newline
A,勝過之意──一個被基督的愛激勵的人,他被基督的愛勝過.比方有一大石頭,人拿不起,洪水來把它沖走,洪水的力勝過石頭的重量.同樣道理,基督的力量是我們生活的力量,祂的力量勝過我們生活的力量.我們完全被基督的力量勝過,被祂的愛拖著而走；不知去向,生活被主的愛勝過.
\newline
\newline
B,掌管──被主的愛掌管,一點都不照自己的意思去生活.這個愛在我裡面成為力量,能夠做平時不願做或不能做的事.
\newline
\newline
C,勉強──意思與掌管相近,好像有一條大河,水流出去沒有被利用,因此給它開一條渠,攔住它.在適當地方建一個發電站,利用水沖的力量來推動發電；水被控制了,掌管了.
\newline
\newline
未被主的愛激勵的人,他的生命中就沒有力量.在基督的愛激勵,生活才有新的方向,也是基督徒應有的表現.
\newline
\newline
被主愛激勵的人,知道主已替眾人死了,我們活著不再為自己而活,乃是為主而活.生活要有這個記號,目標是為主而活.為主而活的基本意義,是指整個生活以主為中心.用主的標準來衡量.試問:我們今日是為誰而活?為自己?為自己的喜好?為自己的得失?還是為主而活?
\newline
\newline
我們生活的標準,並非憑外貌認人.(林後五16)「雖然憑著外貌認過基督,如今卻不再這樣認他了.」這句話是什麼意思呢?根據哥林多信徒對人外貌的認識判斷,亞波羅口才好,跟從他；彼得愛律法,跟從他；因此就產生亞波羅和彼得派.
\newline
\newline
今日我們也有這樣的看法:看人做事方法如何,穿衣服怎樣?家庭背景如何?講話態度如何?因而就憑外貌判斷人家是屬這種人,是屬那種人,產生批評人的態度.很多時候,教會中弟兄姊妹發生問題也不是這個原因.憑著外貌認識人是不深刻的,所以我們要小心,最好不要批評人.做基督的使者,是不憑外貌認人的.
\newline
\newline
基督使者在主前有一個託付,勸人與神和好,勸那些不信的人與神和好.
\newline
\newline
人不能與神和好有四個原因.
\newline
\newline
A,人與神隔離,人對神的事漠不關心,分配不承認神,在生活中從來沒有經驗過神,生活麻木.這種不會反對神,因為他不知道有神,是他的心眼被撒旦弄瞎的結果.
\newline
\newline
B,反抗神的人,他們並不是不承認神,乃是因錯誤的觀念所致,使他大大反抗神；樣樣事都覺得神不好,不公平；因此,使他的生活和經濟發生困難等等,都是神不好.他們對神採取沒有理由的反抗.
\newline
\newline
C,這種人明知道神,當神臨到他時,因為他有罪,便藉詞掩飾.他因為生活不好,不順服而引起自己的困難；反而把責任推卸給神.今日有很多人都是如此,不敢面對現實,對自己的犯罪,推卸責任,甚至責備神.
\newline
\newline
D,這種人用自己犯罪而失敗,軟弱,知道自己不好到了一個程度；認為神不會愛我這種人了.覺得自己太壞了,神豈能愛我呢?
\newline
\newline
基督的使者,帶人到神前與祂和好.什麼原因叫人與神和好呢?哥林多後書五章廿一節那裡說:「神使那無罪的,替我們成為有罪的；好叫我們在他裡面成為神的義.」祂在十字架上擔當了我們的罪.今日,我們因基督的緣故,本來要被定罪的,但被神稱為義了.
\newline
\newline
前面提到的四種人是與神隔離的,所以神在基督裡,勸你把他們帶到祂面前,與祂和好.作使者的人,極重要的使命,就是勸人和好.神已把責職託付了我們,並不是非做不可的,乃是喜歡做或不喜歡做的問題.所以每一個屬主的人都是基督的使者；既是使者,就要盡責忠心去做.靠神的恩典,作基督的使者,忠於託付,勸人與神和好!
\newline
\newline


\newpage

\section{第47屆~港九培靈會~鮑會園~第五講}
\label{sec:1156}
\textbf{天父的兒女}
\newline
\newline
連結: \href{https://www.hkbibleconference.org/session-message/view/1156}{\texttt{https://www.hkbibleconference.org/session-message/view/1156}}
\newline
\newline
\hyperref[sec:1154]{< < < PREV SERMON < < <}
~
\hyperref[sec:toc]{[返目錄]}
~
\hyperref[sec:1157]{> > > NEXT SERMON > > >}
\newline
\newline
哥林多後書 6:1-7:16
\newline
\begin{longtable}{cl}
\hline
\hline
章節 & 經文 (和合本修訂版)\\
\hline
6:1 & \begin{tabularx}{0.7\textwidth}{X} 我們與神同工的也勸你們，不可白受他的恩典； \end{tabularx} \\ \\ \relax
6:2 & \begin{tabularx}{0.7\textwidth}{X} 因為他說：「在悅納的時候，我應允了你；在拯救的日子，我幫助了你。」看哪，現在正是悅納的時候！看哪，現在正是拯救的日子！ \end{tabularx} \\ \\ \relax
6:3 & \begin{tabularx}{0.7\textwidth}{X} 我們不在任何事上妨礙任何人，免得這使命被人毀謗； \end{tabularx} \\ \\ \relax
6:4 & \begin{tabularx}{0.7\textwidth}{X} 反倒在各樣的事上表明自己是神的用人：就如在持久的忍耐、患難、困苦、災難、 \end{tabularx} \\ \\ \relax
6:5 & \begin{tabularx}{0.7\textwidth}{X} 鞭打、監禁、動亂、勞碌、失眠、飢餓、 \end{tabularx} \\ \\ \relax
6:6 & \begin{tabularx}{0.7\textwidth}{X} 廉潔、知識、堅忍、恩慈、聖靈的感化、無偽的愛心、 \end{tabularx} \\ \\ \relax
6:7 & \begin{tabularx}{0.7\textwidth}{X} 真實的言語、神的大能、藉著仁義的兵器在左在右、 \end{tabularx} \\ \\ \relax
6:8 & \begin{tabularx}{0.7\textwidth}{X} 榮譽或羞辱、惡名或美名。我們似乎是誘惑人的，卻是誠實的； \end{tabularx} \\ \\ \relax
6:9 & \begin{tabularx}{0.7\textwidth}{X} 似乎不為人所知，卻是人所共知；似乎是死了，卻是活著；似乎受懲罰，卻沒有被處死； \end{tabularx} \\ \\ \relax
6:10 & \begin{tabularx}{0.7\textwidth}{X} 似乎憂愁，卻常有喜樂；似乎貧窮，卻使許多人富足；似乎一無所有，卻樣樣都有。 \end{tabularx} \\ \\ \relax
6:11 & \begin{tabularx}{0.7\textwidth}{X} 哥林多人哪，我們對你們，口是誠實的，心是寬宏的。 \end{tabularx} \\ \\ \relax
6:12 & \begin{tabularx}{0.7\textwidth}{X} 你們的狹窄不是由於我們，而是由於你們自己的心腸狹窄。 \end{tabularx} \\ \\ \relax
6:13 & \begin{tabularx}{0.7\textwidth}{X} 你們也要照樣用寬宏的心報答我；我這話正像對自己的孩子說的。 \end{tabularx} \\ \\ \relax
6:14 & \begin{tabularx}{0.7\textwidth}{X} 你們不要和不信的人同負一軛。義和不義有甚麼相關？光明和黑暗有甚麼相連？ \end{tabularx} \\ \\ \relax
6:15 & \begin{tabularx}{0.7\textwidth}{X} 基督和彼列有甚麼相和？信主的和不信主的有甚麼相干？ \end{tabularx} \\ \\ \relax
6:16 & \begin{tabularx}{0.7\textwidth}{X} 神的殿和偶像有甚麼相同？因為我們是永生神的殿，就如神曾說：「我要在他們中間居住來往；我要作他們的神，他們要作我的子民。」 \end{tabularx} \\ \\ \relax
6:17 & \begin{tabularx}{0.7\textwidth}{X} 所以主說：「你們務要從他們中間出來，跟他們分別；不要沾不潔淨的東西，我就收納你們。 \end{tabularx} \\ \\ \relax
6:18 & \begin{tabularx}{0.7\textwidth}{X} 我要作你們的父，你們要作我的兒女。這是全能的主說的。」 \end{tabularx} \\ \\ \relax
7:1 & \begin{tabularx}{0.7\textwidth}{X} 所以，親愛的，既然我們有這樣的應許，就當潔淨自己，除去身體和靈魂一切的污穢，藉著敬畏神，得以成聖。 \end{tabularx} \\ \\ \relax
7:2 & \begin{tabularx}{0.7\textwidth}{X} 寬宏大量地接納我們吧！我們未曾虧負誰，未曾敗壞誰，未曾佔誰的便宜。 \end{tabularx} \\ \\ \relax
7:3 & \begin{tabularx}{0.7\textwidth}{X} 我說這話，不是要定你們的罪，我已經說過，你們常在我們心裡，我們情願與你們同生共死。 \end{tabularx} \\ \\ \relax
7:4 & \begin{tabularx}{0.7\textwidth}{X} 我對你們很是放心，多多誇耀你們；我滿有安慰，在我們一切患難中格外喜樂。 \end{tabularx} \\ \\ \relax
7:5 & \begin{tabularx}{0.7\textwidth}{X} 我們從前到了馬其頓的時候，身體沒有絲毫安寧，反而到處遭患難，外有紛爭，內有懼怕。 \end{tabularx} \\ \\ \relax
7:6 & \begin{tabularx}{0.7\textwidth}{X} 但那安慰灰心之人的神藉著提多來安慰了我們； \end{tabularx} \\ \\ \relax
7:7 & \begin{tabularx}{0.7\textwidth}{X} 不但藉著他來，也藉著他從你們所得的安慰，安慰了我們，因為他把你們的思念，你們的哀慟，你們對我的熱忱，都告訴了我，使我更加歡喜。 \end{tabularx} \\ \\ \relax
7:8 & \begin{tabularx}{0.7\textwidth}{X} 即使我先前那封信使你們憂愁，後來我曾懊悔，如今卻不懊悔；因為我知道，那封信使你們憂愁，不過是暫時的。 \end{tabularx} \\ \\ \relax
7:9 & \begin{tabularx}{0.7\textwidth}{X} 如今我歡喜，不是因你們曾憂愁，而是因憂愁導致你們的悔改。你們依著神的意思憂愁，凡事就不至於因我們受虧損了。 \end{tabularx} \\ \\ \relax
7:10 & \begin{tabularx}{0.7\textwidth}{X} 因為依著神的意思而憂愁，就生出沒有懊悔的悔改來，以致得救；但世俗的憂愁叫人死。 \end{tabularx} \\ \\ \relax
7:11 & \begin{tabularx}{0.7\textwidth}{X} 你看，你們依著神的意思而憂愁，這在你們當中產生了何等的殷勤、甚至辯白、甚至憤慨、甚至恐懼、甚至渴望、甚至熱忱、甚至責罰。在這一切事上，你們都表明自己是無可指責的。 \end{tabularx} \\ \\ \relax
7:12 & \begin{tabularx}{0.7\textwidth}{X} 所以，雖然我從前寫信給你們，卻不是為那虧負人的，也不是為那受人虧負的，而是要在神面前把你們顧念我們的熱忱表現出來。 \end{tabularx} \\ \\ \relax
7:13 & \begin{tabularx}{0.7\textwidth}{X} 因此，我們得了安慰。在我們所得的安慰之外，又因你們眾人使提多心裡暢快喜樂，我們就更加歡喜了。 \end{tabularx} \\ \\ \relax
7:14 & \begin{tabularx}{0.7\textwidth}{X} 我若對提多誇獎過你們甚麼，也不覺得慚愧，因為我對提多誇獎你們的話是真的，正如我對你們所說的話也向來都是真的。 \end{tabularx} \\ \\ \relax
7:15 & \begin{tabularx}{0.7\textwidth}{X} 提多一想起你們眾人的順服，怎樣恐懼戰兢地接待他，他愛你們的心就越發熱切了。 \end{tabularx} \\ \\ \relax
7:16 & \begin{tabularx}{0.7\textwidth}{X} 我如今歡喜，因為我在一切事上對你們有信心。 \end{tabularx} \\ \\
[1ex]
\hline
\hline
\end{longtable}
兒女與父母有密切的關係,生命從父母而來,所以生出來像父母.保羅說,我們今日作天父的兒女,我們的生活應在那一方面像天父呢?作兒女的,在父母前有極大的權利或享受大的福氣,是外人不能得著的.
\newline
\newline
保羅勸我們不可徒受神的恩典,我們既蒙了厚恩,我們的生活就應該與所蒙的恩相稱,才不會虧負神.保羅在經文中,提到若有以下幾種的生活表現,就等於徒受神的恩典了.
\newline
\newline
一,知道神一切恩典和救恩的意義
\newline
\newline
聖經教訓作基督徒的,要清楚明白神的救恩是怎麼一回事.若這些只是頭腦的知識,便與生活沒有關係；知道神的恩典,而沒有去體驗它,就是徒受神的恩典了.
\newline
\newline
二,相信基督並知道拯救的恩典
\newline
\newline
在生活中若不能把神豐富的恩典表現出來,他的生活是沒有影響力的.人的生活分兩部份:一是屬靈的,靈魂不再被定罪；離世之後,不經地獄而上天堂.另一是屬世界屬肉體的家庭生活,他所做的一切,與屬靈的事沒有關係,兩者互不相干.如果一個人信了耶穌,得救了；他的生活卻與未得救的人完全一樣!那麼,這個人便是徒受神的恩典了.
\newline
\newline
三,生命沒有長進
\newline
\newline
哥林多教會中有很多人,徒受神的恩典!他們知道得救的整個意義,和脫離死亡得永生；生命要長進,要過聖潔生活,結果子榮耀神.他們都明白了,卻不願去行；生活沒有長進的表現,而且時常過著失敗的生活.甚至一次又一次地,來到主面前流淚悔過祈禱,得著主的安慰和主賜力量；但下次又失敗了,原因是屬靈的原則在他生活中不能起作用.這種人也是徒受神的恩典.
\newline
\newline
哥林多教會有三種人徒受神的恩典,那麼要怎樣做,才能勝過這種徒受神恩典的生活呢?保羅在此提到幾件事,我們必須去做的.
\newline
\newline
一,經常體驗神的恩典
\newline
\newline
耶利米先知說到神的恩典,每一日都是新的.保羅說:「在悅納的時候,我應允了你；在拯救的日子,我搭救了你.看哪,現在正是悅納的時候,現在正是拯救的日子.」這句經文在佈道會時,常被人引用,鼓勵人抓住機會得拯救.但保羅解釋這經文,卻有更深的意義.哥林多信徒都是已得救的人,「拯救的日子」應譯作「救恩顯現的日子」.指神拯救的一切好處,現在在你的生活中發生作用,並在生活中表現出來.現在經驗救恩的工作,而表現出能力來.救恩現在在你生活中發生作用,救恩並非過去的事實.有這經驗的人,便有大喜樂,大福氣.
\newline
\newline
基督徒的生活,切勿靠得救那時的經驗過生活,但很多基督徒卻如此.今日沒有經驗神的恩典,以致無法過得勝的生活.遇上的試探便無法招架了,而完全失敗了.神的恩典每一天都是新的,所以我們要親自去,天天經驗神的恩典；這樣,我們的生活便會被神悅納.
\newline
\newline
二,過屬靈的生活永遠要用屬靈的方法
\newline
\newline
對於屬靈的事,永遠要用屬靈的原則,解決屬靈的問題.從(林後六4)開始,保羅說到生活中經受許多苦難和困苦,或受攻擊和欺騙.這種困苦和壓力,在人看來,是又大又難的擔子.保羅在這種情形下,仍然用基督徒的標準和方法去幫助人；他所用的兵器是仁義,在左在右.這句話「仁義的兵器在左在右」很難翻譯,保羅說,對方用窮困,逼迫,毀謗為武器攻擊我；我卻用神的仁義,握在左手和右手中與仇敵作戰.我們在生活中若用這原則和標準；過天父兒女的生活,就不會失敗.
\newline
\newline
有人為了達到目的,而不擇手段,可惜今日有許多基督徒也採用這種方法；為了要達到目的,用什麼方法也不在乎了.好像請人來禮拜堂聽道,竟用欺騙對方的方法,這是很錯誤的.基督徒的目標很高,所以用的方法要與目標相稱.即使人家受欺騙,人家也會發覺的,這樣我們如何在基督台前交帳呢?我們作神工作的人；一定要用誠實的方法去完成,否則便不能蒙神賜福.
\newline
\newline
主受撒旦試探時,撒旦帶祂上了一座最高的山,將世上的萬國,與萬國的榮華,都指給他看.對祂說:「你若俯伏拜我,我就把這一切都賜給你.」(太四8-9)這句話用我們的話是這麼說:「你到世上來做什麼,豈不是要將世界奪回交給神麼?」主說:「不錯!」撒旦問:「用什麼方法?」主答牠:「十字架.」於是撒旦對祂說:「釘十字架不好受,我給你的方法,要比上十字架容易多了；只要你拜我,我便把世界賜給你.」主對他說:「撒旦退去吧!因為經上記著說,當拜主你的神,單單要事奉他.」做屬靈的工作,如目的不堅定是不行的.
\newline
\newline
三,從憂愁中生出懊悔來
\newline
\newline
保羅說:「依著神的意思憂愁,就產生沒有後悔的懊悔來,以致得救.但世俗的憂愁,是叫人死.」(林後七10)保羅說你們要按著神的意思來悔改,叫你們聖潔,便在試煉中得救,脫離捆綁.唯一的方法,是向神悔改.
\newline
\newline
哥林多的信徒,由於悔改了,心裡便充滿喜樂,同時神的心也得著喜樂.因有悔改的心,整個屬靈的情況都改變了!在試煉中不會失敗,因而得救,脫離捆綁.
\newline
\newline
在神面前悔改,每一個基督徒都應該有這個經驗.但很多人解釋悔改只是一次的經驗,指在信耶穌時一次得救的悔改；信主後不悔改的基督徒是失敗的基督徒.所以基督徒對於在神面前的悔改,有幾件事是必須注意的.
\newline
\newline
一,基督徒犯錯後在神前應有悔改之心
\newline
\newline
悔改並非是一個動作,也不是在人面前和神面前痛苦一番；乃是在神面前的一種態度和存心,在靈裡重新經驗神的恩典.在神面前,人對罪沒有悔改的心,靈命就無法長進,生命便沒有轉機.所以失敗後要向神認罪悔改.
\newline
\newline
二,天父的兒女要過分別為聖的生活
\newline
\newline
天父的兒女,要遠離罪惡,過一個分別為聖的生活.(參看林後六12-18)我們與不信的人不相同,經上記著說:「你們務要從他們中間出來,與他們分別,......我就收納你們.」今天,聖靈在我們裡面,如果我們不聖潔,沾染污穢,那麼聖潔的靈,怎能在污穢的人身上彰顯出來呢?
\newline
\newline
我們要怎樣過分別為聖的生活呢?
\newline
\newline
首先,在我們的目標和思想中分別出來,心中的目標和意念要聖潔.
\newline
\newline
其次,在我們的生活行為上要分別為聖,要明白信主和不信主的,有什麼相干呢?特別在婚姻上,信與不信原不相配；不信的人,其生活沒有得救.所以在這事上要很小心.
\newline
\newline
在工作上,我們要與罪惡分離；不聖潔的工作方法,不能在基督徒的工作上有份.
\newline
\newline
在我們的生活中,不可有不聖潔的娛樂.
\newline
\newline
我們務要從他們中間出來,被神收納,作天父的兒女.求神保守我們,使我們的目標與崇高的生活相稱.
\newline
\newline


\newpage

\section{第47屆~港九培靈會~鮑會園~第六講}
\label{sec:1157}
\textbf{樂意的奉獻}
\newline
\newline
連結: \href{https://www.hkbibleconference.org/session-message/view/1157}{\texttt{https://www.hkbibleconference.org/session-message/view/1157}}
\newline
\newline
\hyperref[sec:1156]{< < < PREV SERMON < < <}
~
\hyperref[sec:toc]{[返目錄]}
~
\hyperref[sec:1158]{> > > NEXT SERMON > > >}
\newline
\newline
哥林多後書 8:1-9:15
\newline
\begin{longtable}{cl}
\hline
\hline
章節 & 經文 (和合本修訂版)\\
\hline
8:1 & \begin{tabularx}{0.7\textwidth}{X} 弟兄們，我們要把神賜給馬其頓眾教會的恩惠告訴你們： \end{tabularx} \\ \\ \relax
8:2 & \begin{tabularx}{0.7\textwidth}{X} 他們在患難中受大考驗的時候，仍然滿有喜樂，在極度貧窮中還格外顯出他們樂捐的慷慨。 \end{tabularx} \\ \\ \relax
8:3 & \begin{tabularx}{0.7\textwidth}{X} 我可以證明，他們是按著能力，而且超過了能力來捐助，主動 \end{tabularx} \\ \\ \relax
8:4 & \begin{tabularx}{0.7\textwidth}{X} 再三懇求我們，准他們在這供給聖徒的善事上有份； \end{tabularx} \\ \\ \relax
8:5 & \begin{tabularx}{0.7\textwidth}{X} 並且他們所做的，不但照我們所期望的，更照神的旨意先把自己獻給主，又給了我們。 \end{tabularx} \\ \\ \relax
8:6 & \begin{tabularx}{0.7\textwidth}{X} 因此，我們勸提多，既然在你們中間開始這慈善的事，就當把它辦成。 \end{tabularx} \\ \\ \relax
8:7 & \begin{tabularx}{0.7\textwidth}{X} 既然你們在信心、口才、知識、萬分的熱忱，以及我們對你們的愛心上，都勝人一等，那麼，當在這慈善的事上也要勝人一等。 \end{tabularx} \\ \\ \relax
8:8 & \begin{tabularx}{0.7\textwidth}{X} 我說這話，並不是命令你們，而是藉著別人的熱忱來考驗你們愛心的真誠。 \end{tabularx} \\ \\ \relax
8:9 & \begin{tabularx}{0.7\textwidth}{X} 你們知道我們主耶穌基督的恩典：他本是富足，卻為你們成了貧窮，好使你們因他的貧窮而成為富足。 \end{tabularx} \\ \\ \relax
8:10 & \begin{tabularx}{0.7\textwidth}{X} 我在這事上把我的意見告訴你們，是對你們有益，因為你們開始辦這事，而且起此心意已經有一年了。 \end{tabularx} \\ \\ \relax
8:11 & \begin{tabularx}{0.7\textwidth}{X} 如今就當辦成這事，既然有願做的心，也當照你們所有的去辦成。 \end{tabularx} \\ \\ \relax
8:12 & \begin{tabularx}{0.7\textwidth}{X} 因為人只要有願做的心，必照他所有的蒙悅納，並不是照他所沒有的。 \end{tabularx} \\ \\ \relax
8:13 & \begin{tabularx}{0.7\textwidth}{X} 我不是要別人輕鬆，你們受累，而是要均勻： \end{tabularx} \\ \\ \relax
8:14 & \begin{tabularx}{0.7\textwidth}{X} 就是要你們現在的富餘補他們的不足，使他們的富餘將來也可以補你們的不足，這就均勻了。 \end{tabularx} \\ \\ \relax
8:15 & \begin{tabularx}{0.7\textwidth}{X} 如經上所記：多收的沒有餘，少收的也沒有缺。 \end{tabularx} \\ \\ \relax
8:16 & \begin{tabularx}{0.7\textwidth}{X} 感謝神，把我對你們的熱忱同樣放在提多心裡。 \end{tabularx} \\ \\ \relax
8:17 & \begin{tabularx}{0.7\textwidth}{X} 他固然聽了我的勸告，但自己更加熱心，自願往你們那裡去。 \end{tabularx} \\ \\ \relax
8:18 & \begin{tabularx}{0.7\textwidth}{X} 我們還差遣一位弟兄和他同去，這人在傳福音的事上得了眾教會的稱讚； \end{tabularx} \\ \\ \relax
8:19 & \begin{tabularx}{0.7\textwidth}{X} 不但這樣，他也被眾教會選派跟我們同行，把所交託我們的這捐款送到了，為的是榮耀主，也表明我們的好意。 \end{tabularx} \\ \\ \relax
8:20 & \begin{tabularx}{0.7\textwidth}{X} 我們這樣做，免得有人因我們收的捐款多而挑剔我們。 \end{tabularx} \\ \\ \relax
8:21 & \begin{tabularx}{0.7\textwidth}{X} 我們留心做好事，不但在主面前，就是在人面前也是這樣。 \end{tabularx} \\ \\ \relax
8:22 & \begin{tabularx}{0.7\textwidth}{X} 我們又差遣一位弟兄同去。這人的熱忱，我們在許多事上屢次考驗過，現在他因為深深信任你們，就更加熱心了。 \end{tabularx} \\ \\ \relax
8:23 & \begin{tabularx}{0.7\textwidth}{X} 至於提多，他是我的夥伴，為服事你們作我的同工。至於那兩位弟兄，他們是眾教會的使者，是基督的榮耀。 \end{tabularx} \\ \\ \relax
8:24 & \begin{tabularx}{0.7\textwidth}{X} 所以，你們務要在眾教會面前向他們顯明你們的愛心和我所誇獎你們的憑據。 \end{tabularx} \\ \\ \relax
9:1 & \begin{tabularx}{0.7\textwidth}{X} 關於供給聖徒的事，我本來不必寫信給你們； \end{tabularx} \\ \\ \relax
9:2 & \begin{tabularx}{0.7\textwidth}{X} 因為我知道你們的好意，常對馬其頓人誇獎你們，說亞該亞人預備好已經有一年了。你們的熱心感動了許多人。 \end{tabularx} \\ \\ \relax
9:3 & \begin{tabularx}{0.7\textwidth}{X} 但我差遣那幾位弟兄去，要使你們照我的話預備妥當，免得我們在這事上誇獎你們的話落了空。 \end{tabularx} \\ \\ \relax
9:4 & \begin{tabularx}{0.7\textwidth}{X} 萬一有馬其頓人與我同去，見你們沒有預備好，就使我們所確信的反成了羞愧；你們的羞愧更不用說了。 \end{tabularx} \\ \\ \relax
9:5 & \begin{tabularx}{0.7\textwidth}{X} 因此，我想必須鼓勵那幾位弟兄先到你們那裡去，把從前所應許的捐款預備妥當，好顯出你們所捐的是出於樂意，不是出於勉強。 \end{tabularx} \\ \\ \relax
9:6 & \begin{tabularx}{0.7\textwidth}{X} 還有一點：「少種的少收；多種的多收。」 \end{tabularx} \\ \\ \relax
9:7 & \begin{tabularx}{0.7\textwidth}{X} 各人要隨心所願，不要為難，不要勉強，因為神愛樂捐的人。 \end{tabularx} \\ \\ \relax
9:8 & \begin{tabularx}{0.7\textwidth}{X} 神能將各樣的恩惠多多加給你們，使你們凡事常常充足，能多做各樣善事。 \end{tabularx} \\ \\ \relax
9:9 & \begin{tabularx}{0.7\textwidth}{X} 如經上所記：「他施捨，賙濟貧窮；他的義行存到永遠。」 \end{tabularx} \\ \\ \relax
9:10 & \begin{tabularx}{0.7\textwidth}{X} 那賜種子給撒種的，賜糧食給人吃的，必多多加給你們種地的種子，又增添你們仁義的果子。 \end{tabularx} \\ \\ \relax
9:11 & \begin{tabularx}{0.7\textwidth}{X} 你們必凡事富足，能多多施捨，使人藉著我們而生感謝神的心。 \end{tabularx} \\ \\ \relax
9:12 & \begin{tabularx}{0.7\textwidth}{X} 因為辦這供給的事，不但補聖徒的缺乏，而且使許多人對神充滿更多的感謝。 \end{tabularx} \\ \\ \relax
9:13 & \begin{tabularx}{0.7\textwidth}{X} 他們從這供給的事上得了憑據，知道你們宣認基督，順服他的福音，慷慨捐助給他們和眾人，把榮耀歸給神。 \end{tabularx} \\ \\ \relax
9:14 & \begin{tabularx}{0.7\textwidth}{X} 他們也因神極大的恩賜顯在你們身上而切切想念你們，為你們祈禱。 \end{tabularx} \\ \\ \relax
9:15 & \begin{tabularx}{0.7\textwidth}{X} 感謝神，因他有說不盡的恩賜！ \end{tabularx} \\ \\
[1ex]
\hline
\hline
\end{longtable}
這段經文在整個哥林多後書的結構上,可能是保羅有意的插入,但它卻是重要的一段.在生活中,金錢佔很重要的地位,所以保羅就論到金錢的奉獻問題.
\newline
\newline
金錢的奉獻,在教會中很很難講的問題；特別自己是傳道主任,或牧師的身份.保羅講到金錢奉獻時,提到極高的水準,有屬靈的原則和標準.我們簡單來思想一下,關於奉獻金錢的教訓.
\newline
\newline
一,為何要奉獻金錢
\newline
\newline
保羅講到奉獻金錢的原因時,常常這樣說:「我把神賜給馬其頓教會的恩告訴你們.」「格外顯出他們樂捐的厚恩.」「主耶穌基督的恩典.」他用「恩典」或「恩」的字.在生活上我們經驗神的恩典,被主恩激勵,而顯出蒙恩的生活.「恩典」一字,其意義非常廣泛,它不但給人好處；乃是指──知道自己的好處,也知道別人的需要.要把自己的好處,給有需要的人,與人分享神的恩典.自己心裡感到別人的需要,願意與他人分享自己的好處,並非佔有他人的好處.自己感到應當去行,按感覺須要去行.
\newline
\newline
二,奉獻的動作本身是恩典
\newline
\newline
馬其頓教會在極窮困中仍顯出厚恩,並非自己豐富,但因有神的恩典,奉獻給神.奉獻是恩典.一個人有力量奉獻,證明有神的恩典.有奉獻的力量,神便有恩典給他.奉獻是神賜的恩典.奉獻目的,叫屬靈的生命得著訓練.我們今生的生活,是為將來在天堂生活作好準備.其次,奉獻的目的,就是與人在工作上分享所得的益處.奉獻金錢應看作是基本的責任.
\newline
\newline
三,如何奉獻
\newline
\newline
林後八3說:「我可以證明他們是按著力量,而且也過了力量,自己甘心樂意的捐助.」這句話講出了奉獻的重要原則,按馬其頓信徒的力量,他們的奉獻超過了自己的力量.神看人的奉獻,不是看數目字,數目字不重要；乃是看他奉獻還剩下多少,祂要看你是否盡了力量或過了力量.今日講什一奉獻,按這標準做是不錯的.我覺得不單要按這標準來奉獻,而且更要按恩典的原則來奉獻.舉例來說,香港一家四口的家庭,每月入息一千元；奉獻十分一即一百元,我認為他的奉獻不但按著力量；而且是過了力量,九百元的開支實在是很吃力的.若有別一家庭,也是四口人,月入萬元,奉獻一千；我覺得他很吝嗇,沒有盡力量奉獻,更談不上過了力量奉獻了.
\newline
\newline
我們若在神前蒙恩,還要向神計較嗎?神要我們怎樣奉獻呢?
\newline
\newline
A,不是跟神計較數目,要按力量和過了力量來奉獻.馬其頓的教會在極窮困之時,還是超過了力量來奉獻.我們越奉獻,就會更豐富；奉獻之時,也會顯出樂捐的心.
\newline
\newline
B,要先預備妥當,所捐獻的是出於樂意,並非勉強.
\newline
\newline
「預備妥當」有幾方面的意思:
\newline
\newline
a,我們經過思想,用我們的智慧來分析和好好的計劃.想到神工作的需要,神有供給；便要盡力量去做要做的,忠實的準備好,作經常的奉獻給神.這是神感動我們去做,不是隨隨便便馬馬虎虎的.
\newline
\newline
b,表示我們的奉獻,是有系統的奉獻；有特別工作的需要,便作特別的奉獻.如建堂奉獻即是.但神的工作經常的進行,並非靠特別的奉獻,乃靠所有的弟兄姊妹,有系統的經常的奉獻.特別的奉獻──神會感動有力量的人來奉獻.每一個屬主的人,在奉獻的事奉上要有分；也有責任這樣做,所以要盡力量去做.好像聖經上記載的那個窮寡婦,奉獻兩個小錢的事；她已盡了全力去奉獻,因為她把養生的都獻上了,經常不斷的奉獻,是神前忠心的奉獻.
\newline
\newline
c,有一定的時間來奉獻,保羅說七日的頭一日便要拿出來奉獻.他不是說得著錢之後,等把開銷分配好,剩下的才拿來奉獻.乃是要先抽出來放在一邊,留著給神；然後其餘的錢,才是屬自己的,可分配使用.絕不是說,用完了剩下的,才拿來奉獻給神.
\newline
\newline
d,我們的奉獻,應出於樂意,並非出於勉強.換言之,奉獻金錢給神,並非是不得已的事；也不是別人看見而不好意思,才拿出一點來奉獻,或是怕口頭的壓力,或面子的壓力,或其他的壓力,這樣才來奉獻.神要我們的奉獻,是出於樂意的.
\newline
\newline
馬其頓教會的信徒,是極貧窮的,保羅不會這樣說:不要奉獻那麼多,教會有需要神會負責的.但保羅卻說:「他們求我將他們的錢拿去,他們再三的求我這樣做；他們的奉獻不是出於勉強,乃是出於樂意」.
\newline
\newline
四,奉獻的方法
\newline
\newline
照著神的旨意奉獻主,這種奉獻乃是我們真知道主的恩典,被主的恩激勵和感動.馬其頓教會是一個很好的例子,他們在極窮困之時,仍然把一切奉獻給主.
\newline
\newline
我們的主,本來是富足的,但因著我們成了貧窮.他掌管宇宙萬物,但祂甘心樂意來到世界上,將享受的榮耀放下.變成人的樣式,在人間過著卑微的生活.祂生在馬槽中,祂無佳形美容.祂被釘在十字架上,為什麼呢?為了我們的罪,被釘死在十字架上.祂叫我們因著祂的貧窮成了我們的富足.馬其頓的信徒不能與主相提並論.他們雖為了奉獻給主,願意放棄一切；但因為這是恩典,從主那裡來的恩典.
\newline
\newline
保羅說:奉獻之時.並非為數目的多少,乃是要先將自己奉獻給主.若沒有這個心意奉獻,那麼奉獻也是勉強的.奉獻不是犧牲,乃是主的恩典.有一宣教士到韓國,被牧師帶到田間去,看見兩人在田裡耕田.前頭的年輕人拉犁耙,後面的老年人推著,宣教士問:「他們沒有牛耕田,不是很苦嗎?」
\newline
\newline
那牧師回答他:「他們是我們的信徒,最近他們為了建堂奉獻,把牛賣了.」
\newline
\newline
「那不是很大的犧牲嗎?」
\newline
\newline
牧師答:「他們不覺得是大犧牲,乃是大福氣,牛賣去了,仍可耕田.」
\newline
\newline
那位宣教士聽了之後,大受感動.
\newline
\newline
五,奉獻的結果
\newline
\newline
奉獻之後,有四方面的結果:
\newline
\newline
A,屬靈的生命得造就,屬靈生命蒙福；不奉獻給神的人,生命不蒙福.
\newline
\newline
B,成為別人的見證,神的名得著榮耀,人的生命得到祝福.
\newline
\newline
C,神的工作能蒙福,分享的真正意義,是叫別人的工作蒙福.
\newline
\newline


\newpage

\section{第47屆~港九培靈會~鮑會園~第七講}
\label{sec:1158}
\textbf{屬靈的權柄}
\newline
\newline
連結: \href{https://www.hkbibleconference.org/session-message/view/1158}{\texttt{https://www.hkbibleconference.org/session-message/view/1158}}
\newline
\newline
\hyperref[sec:1157]{< < < PREV SERMON < < <}
~
\hyperref[sec:toc]{[返目錄]}
~
\hyperref[sec:1160]{> > > NEXT SERMON > > >}
\newline
\newline
哥林多後書 10:1-10:18
\newline
\begin{longtable}{cl}
\hline
\hline
章節 & 經文 (和合本修訂版)\\
\hline
10:1 & \begin{tabularx}{0.7\textwidth}{X} 我－保羅與你們見面的時候是溫和的，不在你們那裡的時候向你們是勇敢的，如今親自藉著基督的溫柔和慈祥勸你們。 \end{tabularx} \\ \\ \relax
10:2 & \begin{tabularx}{0.7\textwidth}{X} 有人認為我們是憑著血氣行事，我認為必須敢於對付這等人；我但求在那裡的時候，不必這樣勇敢。 \end{tabularx} \\ \\ \relax
10:3 & \begin{tabularx}{0.7\textwidth}{X} 我們雖然在血氣中行事，卻不憑著血氣爭戰。 \end{tabularx} \\ \\ \relax
10:4 & \begin{tabularx}{0.7\textwidth}{X} 因為我們爭戰的兵器本不是屬血氣的，而是憑著神的能力，能夠攻破堅固的營壘。我們攻破各樣的計謀， \end{tabularx} \\ \\ \relax
10:5 & \begin{tabularx}{0.7\textwidth}{X} 和各樣攔阻人認識神的高壘，又奪回人心來順服基督。 \end{tabularx} \\ \\ \relax
10:6 & \begin{tabularx}{0.7\textwidth}{X} 我已經預備好了，等你們完全順服的時候來懲罰所有不順服的人。 \end{tabularx} \\ \\ \relax
10:7 & \begin{tabularx}{0.7\textwidth}{X} 你們只看事情的外表。倘若有人自信是屬基督的，他要再想想，他屬基督，我們也屬基督。 \end{tabularx} \\ \\ \relax
10:8 & \begin{tabularx}{0.7\textwidth}{X} 主賜給我們權柄，是要造就你們，並不是要拆毀你們；我就是為這權柄稍微誇口也不覺得慚愧。 \end{tabularx} \\ \\ \relax
10:9 & \begin{tabularx}{0.7\textwidth}{X} 我說這話，免得你們以為我寫信是要恐嚇你們。 \end{tabularx} \\ \\ \relax
10:10 & \begin{tabularx}{0.7\textwidth}{X} 因為有人說：「他信上的語氣既嚴厲又強硬，他本人卻軟弱無能，言語粗俗。」 \end{tabularx} \\ \\ \relax
10:11 & \begin{tabularx}{0.7\textwidth}{X} 這等人當明白，我們不在那裡時信上怎麼說，見面時也必怎麼做。 \end{tabularx} \\ \\ \relax
10:12 & \begin{tabularx}{0.7\textwidth}{X} 因為我們不敢將自己和某些自我推薦的人並列相比；他們用自己度量自己，用自己比較自己，是不明智的。 \end{tabularx} \\ \\ \relax
10:13 & \begin{tabularx}{0.7\textwidth}{X} 我們不願意過分誇口，但是我們只在神劃定的界限內誇口。這界限甚至擴展到你們那裡。 \end{tabularx} \\ \\ \relax
10:14 & \begin{tabularx}{0.7\textwidth}{X} 我們擴展到你們那裡時並沒有越過了自己的界限，其實我們是首先到你們那裡傳基督福音的。 \end{tabularx} \\ \\ \relax
10:15 & \begin{tabularx}{0.7\textwidth}{X} 我們不靠別人所勞碌的過分誇口；我們只希望你們信心增長的時候，所劃定給我們的範圍也能夠因著你們更加擴展， \end{tabularx} \\ \\ \relax
10:16 & \begin{tabularx}{0.7\textwidth}{X} 使福音得以傳到你們以外的地方，而不在別人的範圍之內，以別人所成就的事誇口。 \end{tabularx} \\ \\ \relax
10:17 & \begin{tabularx}{0.7\textwidth}{X} 但「要誇耀的，該誇耀主」。 \end{tabularx} \\ \\ \relax
10:18 & \begin{tabularx}{0.7\textwidth}{X} 因為蒙悅納的，不是自我稱許的，而是主所稱許的。 \end{tabularx} \\ \\
[1ex]
\hline
\hline
\end{longtable}
哥林多教會有這麼多困難,是因有假師傅假先知在他們當中.保羅知道哥林多信徒悔改了,便願意回到他們當中去,要對付假師傅.用什麼方法對付他們?是憑主賜的權柄.
\newline
\newline
保羅的權柄根基是基督,所以他大有能力.我們生活在今日社會中,面臨許多的挑戰.若是我們不清楚立場和根基,就無法應付困難,因此我們的立場應建立在:
\newline
\newline
一,權柄非建立在人的智慧上
\newline
\newline
我們作戰的兵器非憑血氣,乃靠神的能力,才能攻破堅固的營壘；奪回人的心意,順服基督,攻下攔阻神的事.
\newline
\newline
「堅固營壘」,指人的心意,攔阻人認識神至高之事,保羅說,要將它攻破；將人的心意挽回,而跟從主.
\newline
\newline
人的心意是什麼呢?即自己的思想方法.人憑著自己的看法,決定自己的生活方向.神賜人智慧,今日的科學家百分之九十,并非基督徒,感謝神!也賜給他們智慧.哲學決定人生活方向的事,科學家處理物質的東西.人的哲學不可靠,不是建立在基督的根基上,乃建立在小範圍上,他們評論的標準也不準確.許多是似是而非的事,即使有一點點的道理,也只不過是一部份,而且是似是而非的.由於無神的哲學思想,就把自己特別限制在小圈子中,找理論支持或證實自己的道理可靠.他們是站在一個立場,要信什麼,不信什麼；接受這個,否認那個,是要根據自己的立場來作標準.可是,人的思想,看法,至終是攔阻人認識神.這是堅固的營壘,所以保羅說要攻破堅固的營壘,然後用真理決定生活的方向.
\newline
\newline
二,權柄的根基非建立在外表形狀上
\newline
\newline
有人批評保羅,說他講話又沉重又厲害,在很遠的地方也會講話；但,他真人的生活並非如此,他氣貌不揚,言語粗俗.他們這樣批評保羅.保羅反駁說:你以什麼來批評我?講的是什麼?教訓的是什麼?傳的是什麼?是我的外貌呢?還是真理呢?或者是我的教育?我不是要你們看我好看不好看,也許我的言語粗俗,我的目的是要你們接受神的真理,並非接受我的學問和知識.我是要你們接受主和祂釘十字架,並非看我外表的表現.
\newline
\newline
今日的社會,滿了試探,正如提摩太後書四章三至四節所說的:「因為時候要到,人必厭煩純正的道理；耳朵發癢,就隨從自己的情慾,增添好些師傅；並且掩耳不聽真道,偏向荒渺的言語.」以致講的雖是聖經的話,有什麼好聽呢?新的學問才好聽.但保羅說:我們要有神的亮光,不要把信仰根基建立在可看見的外表上面.哥林多信徒最會犯這錯誤,亞波羅有口才,會講聖經；他從亞力山大來,有好背景,結果就產生亞波羅派.同樣,今日教會也有類似的派別,求神保守我們信仰的根基,是接受在福音上.
\newline
\newline
三,權柄的根基非建立在人的評論上
\newline
\newline
保羅說這種假使徒,行事險詐,假作基督的使徒,裝著仁義,裝出外表.他們的權柄建立在人的評判上.但保羅說,權柄,不是建立在人的智慧上和外表上.
\newline
\newline
從積極上講,講柄應建立在幾方面:
\newline
\newline
A,建立在基督身上,建立在神的話語上.保羅說,這種權柄是主所賜的,他很注重這點,在他所寫的多本書信中都有提及.他的權柄是神所賜的,非教會所賜的,也非使徒所賜的.他憑什麼資格說這話呢?他說基督給他權柄,因在他的生命中,他已先經驗到基督自己.他說到在大馬色路上的經歷,主有聲音對他說:「掃羅,掃羅,你為什麼逼迫我?」他回答說:「主啊!你是誰?」「主」不可輕用,這字顯明一事,指神蹟發生在他身上的事.他完完全全順服主,承認主,在他身上有絕對的權柄.
\newline
\newline
一個人順服之後,才能知道主是誰,才可說主賜他權柄；生活中若沒有經驗過主,生命中就不能真正的順服主.即使你口才很好,能講很好聽的話,倘若沒有主的生命,那麼你就沒有主的權柄.隨便使用權柄是危險的.你如果沒有真心的奉獻,那麼你做主的工作也是危險的.你沒有能力,作奉獻的工人是危險的.例如在以弗所有幾個信徒,沒有順服主；他們看是耶穌趕鬼,結果他們也去趕鬼說:「我奉保羅所信的耶穌趕你出去.」鬼回答說:「你是誰?保羅我知道,耶穌我認得.」於是鬼跳到他們身上,附在他們身上；濫用主的權柄是危險的事.
\newline
\newline
B,因對工作的態度,顯出他的權柄.保羅的工作態度,是要早到他們那裡去,傳耶穌基督的福音.不但是勞苦,還指望他們的信心更加開展.因他這幾句簡單的話,顯出他工作的態度.他不是等他們傳了福音,自己去收果子,得別人的功勞.保羅說:「我把福音建立在主身上.」
\newline
\newline
今日很多工作的態度,與保羅相反,他們享受了別人的功勞.甚至不是自己的功勞,也歸功於自己；到處顯露自己,想得人的讚賞.
\newline
\newline
C,福音要傳到更遠的地方去.這是保羅的工作方法.保羅給我們一個原則,要將音傳到遙遠的地方.保羅建立他們,成為福音的根據地；是要把福音要傳到更遠的地方.並非要建立他們,成為他的地盤.這是初期教會傳道的重要原則.我們必須將這原則應用到我們的生活中.保羅傳福音受逼迫的目的,是叫福音傳到更遠的地方.我們今日的態度也是如此,不論作什麼,為要叫福音更加興旺.
\newline
\newline
四,權柄應建立在對真理的分辨上面
\newline
\newline
撒旦裝著光明的天使,仁義的使者.因為有真的,才會有假的出來跟隨著.例如有真鈔票,才有假鈔票出來；這假鈔票沒有用處,沒有價值的.真的玉石有價值,人造的假玉石；再好看也看沒有價值.
\newline
\newline
由於福音有價值,威脅著撒旦；叫人得生命,因此撒旦便推出假冒的福音.今日在教會中,有假冒的使者；淆亂福音,似是而非,甚至事奉也有假的.什麼是真的,什麼是假的呢?這須要從上頭而來的智慧,方能分辨出真假來.
\newline
\newline
五,權柄的根基,在身上帶有受苦記號
\newline
\newline
保羅在生活中受了很多苦難,經過苦難,得著神的安慰.茲將他受的苦難歸納如下:
\newline
\newline
他知道不能倚靠自己是福氣(林後一9).
\newline
\newline
受苦難使基督的生命顯明在他身上(林後四10-12).
\newline
\newline
受苦難要承受將來極重無比的榮耀(林後四17-18).
\newline
\newline
受苦難常有喜樂,樣樣都有,有屬靈的恩典.(林後六10)
\newline
\newline
因著這一切苦難,可得從神而來的安慰.(林後七5-6)
\newline
\newline
攻破堅固的營壘.(林後十4)
\newline
\newline
更多苦難沒人能了解.(林後十一章)
\newline
\newline
保羅描寫「苦難壓在他身上」,這句話很難翻譯.但可形容作,在火燭之時用棉被壓住火；又如老虎與小綿羊搏鬥之時,老虎壓在小羊身上一樣.保羅使用極有力的字.靈裡有重負,他甘心樂意去承受.這是主裡的權柄的記號.有此權柄,才能攻破堅固的營壘；攻擊裝作光明天使的假使者.
\newline
\newline


\newpage

\section{第47屆~港九培靈會~鮑會園~第八講}
\label{sec:1160}
\textbf{足夠的恩典}
\newline
\newline
連結: \href{https://www.hkbibleconference.org/session-message/view/1160}{\texttt{https://www.hkbibleconference.org/session-message/view/1160}}
\newline
\newline
\hyperref[sec:1158]{< < < PREV SERMON < < <}
~
\hyperref[sec:toc]{[返目錄]}
~
\hyperref[sec:1059]{> > > NEXT SERMON > > >}
\newline
\newline
哥林多後書 12:1-13:14
\newline
\begin{longtable}{cl}
\hline
\hline
章節 & 經文 (和合本修訂版)\\
\hline
12:1 & \begin{tabularx}{0.7\textwidth}{X} 雖然自誇無益，我還是不得不誇。我現在要提到主的異象和啟示。 \end{tabularx} \\ \\ \relax
12:2 & \begin{tabularx}{0.7\textwidth}{X} 我認識一個在基督裡的人，他在十四年前被提到第三層天上去；或在身內，我不知道，或在身外，我也不知道，只有神知道。 \end{tabularx} \\ \\ \relax
12:3 & \begin{tabularx}{0.7\textwidth}{X} 我認識的這樣的一個人—或在身內，或在身外，我都不知道，只有神知道— \end{tabularx} \\ \\ \relax
12:4 & \begin{tabularx}{0.7\textwidth}{X} 他被提到樂園裡，聽見隱祕的言語，是人不可說的。 \end{tabularx} \\ \\ \relax
12:5 & \begin{tabularx}{0.7\textwidth}{X} 為這人，我要誇口；但是為我自己，除了我的軟弱以外，我並不誇口。 \end{tabularx} \\ \\ \relax
12:6 & \begin{tabularx}{0.7\textwidth}{X} 就是我願意誇口也不算狂，因為我會說實話；只是我絕口不談，恐怕有人把我看得太高了，過於他在我身上所看見所聽見的； \end{tabularx} \\ \\ \relax
12:7 & \begin{tabularx}{0.7\textwidth}{X} 又恐怕我因所得的啟示太高深，就過於高抬自己，所以有一根刺加在我身上，就是撒但的差役來折磨我，免得我過於高抬自己。 \end{tabularx} \\ \\ \relax
12:8 & \begin{tabularx}{0.7\textwidth}{X} 為了這事，我曾三次求主使這根刺離開我。 \end{tabularx} \\ \\ \relax
12:9 & \begin{tabularx}{0.7\textwidth}{X} 他對我說：「我的恩典是夠你用的，因為我的能力是在人的軟弱上顯得完全。」所以，我更喜歡誇耀自己的軟弱，好使基督的能力覆庇我。 \end{tabularx} \\ \\ \relax
12:10 & \begin{tabularx}{0.7\textwidth}{X} 為基督的緣故，我以軟弱、凌辱、艱難、迫害、困苦為可喜樂的事；因為我甚麼時候軟弱，甚麼時候就剛強了。 \end{tabularx} \\ \\ \relax
12:11 & \begin{tabularx}{0.7\textwidth}{X} 我成了愚蠢人，是被你們逼出來的，因為我本該被你們讚許才是。雖然我算不了甚麼，卻沒有一件事在那些超級使徒以下。 \end{tabularx} \\ \\ \relax
12:12 & \begin{tabularx}{0.7\textwidth}{X} 我在你們中間，用百般的忍耐，藉著神蹟、奇事、異能顯出使徒的憑據來。 \end{tabularx} \\ \\ \relax
12:13 & \begin{tabularx}{0.7\textwidth}{X} 除了我不曾連累你們這一件事，你們還有甚麼事不及別的教會呢？這不公平之處，請你們饒恕我吧。 \end{tabularx} \\ \\ \relax
12:14 & \begin{tabularx}{0.7\textwidth}{X} 如今，我準備第三次到你們那裡去。我仍不會連累你們，因為我所求的是你們，不是你們的財物。兒女不該為父母積財，父母該為兒女積財。 \end{tabularx} \\ \\ \relax
12:15 & \begin{tabularx}{0.7\textwidth}{X} 我也甘心樂意為你們的靈魂費財費力。難道我越愛你們，就越少得你們的愛嗎？ \end{tabularx} \\ \\ \relax
12:16 & \begin{tabularx}{0.7\textwidth}{X} 罷了，我自己並沒有連累你們，你們卻有人說，我施詭詐，用心計牢籠你們。 \end{tabularx} \\ \\ \relax
12:17 & \begin{tabularx}{0.7\textwidth}{X} 我所差遣到你們那裡去的人，我何曾藉著他們中的任何人佔過你們的便宜呢？ \end{tabularx} \\ \\ \relax
12:18 & \begin{tabularx}{0.7\textwidth}{X} 我勸提多到你們那裡去，又差遣那位弟兄與他同去，提多佔過你們的便宜嗎？我們的行事為人不是同一心靈嗎？不是同一步伐嗎？ \end{tabularx} \\ \\ \relax
12:19 & \begin{tabularx}{0.7\textwidth}{X} 你們一直認為我們是在你們面前為自己辯護嗎？其實，我們本是在基督裡當著神面前說話。親愛的，一切的事都是為了造就你們。 \end{tabularx} \\ \\ \relax
12:20 & \begin{tabularx}{0.7\textwidth}{X} 我怕我再來的時候，見你們不合我所期望的，而你們見我也不合你們所期望的。我怕有紛爭、嫉妒、憤怒、自私、毀謗、讒言、狂傲、動亂的事。 \end{tabularx} \\ \\ \relax
12:21 & \begin{tabularx}{0.7\textwidth}{X} 我怕我再來的時候，我的神使我在你們面前蒙羞，並且又因許多人從前犯罪，行污穢、淫亂、放蕩的事，不肯悔改而悲傷。 \end{tabularx} \\ \\ \relax
13:1 & \begin{tabularx}{0.7\textwidth}{X} 這是我第三次要到你們那裡去。「任何指控都要憑兩個或三個證人的口述才能成立。」 \end{tabularx} \\ \\ \relax
13:2 & \begin{tabularx}{0.7\textwidth}{X} 對那些犯了罪的人和其餘所有的人，正如我第二次見你們的時候曾說過，現在不在你們那裡再次說：「我若再來，必不寬容。」 \end{tabularx} \\ \\ \relax
13:3 & \begin{tabularx}{0.7\textwidth}{X} 因為你們想求證基督是否藉著我說話。基督對你們並不是軟弱的，而是在你們裡面大有能力的。 \end{tabularx} \\ \\ \relax
13:4 & \begin{tabularx}{0.7\textwidth}{X} 他因軟弱被釘在十字架上，卻因神的大能仍然活著。我們在他裡面也成為軟弱的，但對你們，我們將因神的大能而與他一同活著。 \end{tabularx} \\ \\ \relax
13:5 & \begin{tabularx}{0.7\textwidth}{X} 你們總要省察自己是否在信仰中生活；你們要考驗自己。除非你們經不起考驗，你們自己豈不應該知道有耶穌基督在你們裡面嗎？ \end{tabularx} \\ \\ \relax
13:6 & \begin{tabularx}{0.7\textwidth}{X} 我希望你們知道，我們並不是經不起考驗的人。 \end{tabularx} \\ \\ \relax
13:7 & \begin{tabularx}{0.7\textwidth}{X} 我們祈求神使你們不做任何惡事；這不是要顯明我們是經得起考驗的，而是要你們行事端正，即使我們似乎經不起考驗也沒有關係。 \end{tabularx} \\ \\ \relax
13:8 & \begin{tabularx}{0.7\textwidth}{X} 我們不能做任何對抗真理的事，只能維護真理。 \end{tabularx} \\ \\ \relax
13:9 & \begin{tabularx}{0.7\textwidth}{X} 當我們軟弱而你們剛強時，我們也歡喜。我們所祈求的是：你們能成為完全人。 \end{tabularx} \\ \\ \relax
13:10 & \begin{tabularx}{0.7\textwidth}{X} 所以，我不在你們那裡的時候，把這些話寫給你們，好使我見你們的時候不用照主所給我的權柄嚴厲地待你們；這權柄原是為造就人，而不是為摧毀人。 \end{tabularx} \\ \\ \relax
13:11 & \begin{tabularx}{0.7\textwidth}{X} 末了，弟兄們，願你們喜樂。要追求完全；要接受鼓勵；要同心合意；要彼此和睦。如此，慈愛和平的神必與你們同在。 \end{tabularx} \\ \\ \relax
13:12 & \begin{tabularx}{0.7\textwidth}{X} 你們要用聖潔的吻彼此問安。 \end{tabularx} \\ \\ \relax
13:13 & \begin{tabularx}{0.7\textwidth}{X} 眾聖徒都向你們問安。 \end{tabularx} \\ \\ \relax
13:14 & \begin{tabularx}{0.7\textwidth}{X} 願主耶穌基督的恩惠、神的慈愛、聖靈的感動常與你們眾人同在！ \end{tabularx} \\ \\
[1ex]
\hline
\hline
\end{longtable}
保羅在哥林多後書提到許多的經驗,目的並非顯示自己的好處；乃是提到他極深的屬靈生活,顯示出神的大恩典,極為豐盛和寶貴.這個經驗也是我們生活所需要的.但是,神給我們的經驗,是沒有可誇的；我們之所以有這經驗,因為神要使用我們.或者從另一眼光來看:因為我比你好,所以神才給我們這個經驗；這樣就錯了!對領受神的恩典的認識就不夠,實在太膚淺了.
\newline
\newline
保羅在聖經中作了美好的見證,他講到幾個靈裡的經驗:
\newline
\newline
一,靈裡見到異象
\newline
\newline
保羅提到他十四年前被提到第三層天.他說:「或在身內,我不知道；或在身外,我也不知道,只有神知道.」這句話顯出事情不是出於他自己；但領受這經驗的,卻是他自己.十四年前得這經驗,是什麼時候呢?哥林多後書是主後五十六年冬天寫的,從使徒行傳的年代推算；十四年前就是他開始作第一次國外佈道的時候.他帶著巴拿巴,馬可同工同去；至中途,馬可離去.保羅去到路司得,有特別的經驗；行了神蹟,當地的人要向他們獻祭.保羅便向他們解釋說:「諸位,為什麼作這事呢?我們也是人,性情和你們一樣；我們傳福音給你們,是叫你們離棄這些虛妄,歸向那創造天,地,海,和其中萬物的永生神.」
\newline
\newline
那時,猶太人從各地來到,挑唆眾人用石頭打保羅,正以為他死了,把他拖到城外.神保守他,第二天又回到路司得城去.他被眾人打昏後,這個經驗是否就是他昏了之時,靈裡的經驗呢?神知道祂僕人的需要,所以給他深深的屬靈經驗,他不敢說是否那時,但總之是在他受深深擊打時,神給了他這個經驗.
\newline
\newline
保羅在靈裡得著極寶貴的經驗,他見到異象,是其他人不易得著的.不像今日的人,講靈裡的經驗時,大談召魂回來見面的事.然而保羅的經驗不可支持這行動,他所做的與世人所做的有很大的分別:
\newline
\newline
A,保羅清楚自己的經驗,這是基督徒的經驗.
\newline
\newline
B,他的經驗非在黑暗中進行.
\newline
\newline
C,沒有交鬼的人來替他進行.
\newline
\newline
D,不是他尋求這種經驗,乃是他不知不覺之時,神向他格外施恩.
\newline
\newline
E,他沒有勸人去尋求他同樣的經驗.
\newline
\newline
F,他的經驗完全在靈裡的,顯然不是肉體被提到第三層天去.
\newline
\newline
屬靈的經驗是神賜給祂的兒女的,但每一個人的經驗都不同,所以我們絕不可,以別人的經驗作為我們的標準.我們要在主裡追求,用屬靈的生命去得著屬靈的經驗.
\newline
\newline
二,保羅的肉體上有一根刺
\newline
\newline
保羅說:「恐怕我因所得的啟示甚大,就過於自高；所以有一根刺,加在我肉體上.就是撒旦的差役,要攻擊我,免得我過於自高.」為了這根刺,他曾三次求主叫這刺離開他.「刺」字很特別,並非是栽花時被花刺了一下；乃是一根很厲害的刺,是傷人很厲害的東西.為此,才能叫保羅整個身心不安.這根刺絕非指林後十一章末所記載的各種痛苦,這根剌比這些更厲害,令他沒有平安,他迫切求主拿開它,因為他忍受不了.有刺在身體上,以致不能在主面前好好工作.我們不知道他的那根刺是什麼,但解經家推測這根刺,是指他生活中的困難；身體有軟弱,眼睛差不多失明了.
\newline
\newline
保羅曾三次求主叫刺離開他.聖經譯作「三次」,原文非三次.「三」代表完全的數目,指他一而再,再而三的向主禱告,求主拿開它.他一定是禱告很多年了.主什麼時候給他這根刺呢?是否他在靈裡被提到第三層天時,主給他的呢?這樣說來,那就有十四年之久了.雖然他逼切的向主求,但主始終沒有將他身上的那根刺挪去.好像主根本沒有聽他禱告一樣.在我們生活中,有沒有這個經驗?有時我們為一件事求神,條件都合了,但就是神不做工作,主不垂聽,因此我們就灰心.為何主不聽呢?我們看保羅的經驗就明白.
\newline
\newline
三,主的恩典夠用
\newline
\newline
我們想像當時的情景.保羅說:「主阿!求你挪開我身上這根刺.」但主回答他:「我的恩典夠你用.」他多次這樣的求,然而主還是這樣回答他:「我的恩典夠你用.」我們有這個經驗嗎?有時候我們向神禱告一件事,求主成就.但主說:「我要給你另一條道路.」主有祂的計劃,智慧；主的意念高過我們的意念,祂的智慧高過我們的智慧.我們要讓主有機會給我們指示出當行的道路.
\newline
\newline
第九節說:「我的能力在軟弱的人身上顯得完全.」若果你是剛強的,為神工作有能力,恩賜,學問.這樣做工作至終可能是自己做的；若我們真承認自己一無所有,能作成工作；那麼就非你做的,乃是主恩典藉著你顯出來的作為.因為主的能力,在人的軟弱上顯得完全.
\newline
\newline
主為何加一根刺給保羅呢?免得他自高,看自己太高,越過神的地位,奪去神的榮耀和功勞.保羅說:「有一根刺加在我肉體上.」中文「加」字很特別,我們有留意到它嗎?呂振中翻譯作:「有一根刺給了我.」意即這根刺非我本性而有的,是另外有一位給了我的.下文講到撒旦的差役攻擊他,這是神所准許的.實際上這根刺是主所給的,目的是免得他自高.若是神給你一根刺,你會樂意接受嗎?但我不知祂給你的刺是一根什麼刺.
\newline
\newline
四,受凌辱和苦難仍有喜樂
\newline
\newline
主對保羅說:「我的恩典夠你用的,因為我的能力,是在人的軟弱上顯得完全.」所以他更喜歡誇自己的軟弱,好叫基督的能力護庇他.他為基督的緣故,就以軟弱,凌辱,急難,逼迫,困苦,為可喜樂的；因他什麼時候軟弱,甚麼時候就剛強了.保羅說:「我沒有口才,叫我說話能打動人心的,不是我自己的話；乃是聖靈的話,也是神的恩典.所以我更喜歡誇自己的軟弱,因為神的恩典夠我用的.」
\newline
\newline
五,他行事端正
\newline
\newline
保羅說:「我們求神,叫你們一件惡事都不作；這不是要顯明我們是蒙悅納的,是要你們行事端正,任憑人看出我們是被棄絕的吧!」遵主的話去行,顯出主的恩典.我自己怎樣不要緊,甚至被人批評,棄絕,最重要的是要行主的道.我們要按正義而行.即使如此行,受更多的痛苦也不要緊；我們若有這心志,主就會格外施恩典給我們.
\newline
\newline
保羅求主挪去他身上的一根刺,正如主耶穌在客西馬尼園禱告:「父阿,求你將這杯撤去,但不要照我的意思,而要照你的旨意成全.」主的禱告雖說叫杯撤去,但祂對父說:「按你的旨意而行.」與保羅之求主叫這根刺離開他,兩者有很大分別.保羅經過這次經驗後,他的禱告便如主耶穌那樣說:「父阿,照你的旨意而行吧.」
\newline
\newline


\newpage

\section{第48屆~港九培靈會~史祈生~第一講}
\label{sec:1059}
\textbf{生命的印證(一)}
\newline
\newline
連結: \href{https://www.hkbibleconference.org/session-message/view/1059}{\texttt{https://www.hkbibleconference.org/session-message/view/1059}}
\newline
\newline
\hyperref[sec:1160]{< < < PREV SERMON < < <}
~
\hyperref[sec:toc]{[返目錄]}
~
\hyperref[sec:1060]{> > > NEXT SERMON > > >}
\newline
\newline
經文:「耶穌看見這許多的人,就上了山,既已坐下,門徒到祂跟前來.祂就開口教訓他們說:虛心的人有福了,因為天國是他們的.哀慟的人有福了,因為他們必得安慰.溫柔的人有福了,因為他們必承受地土.飢渴慕義的人有福了,因為他們必得飽足.憐恤人的人有福了,因為他們必蒙憐恤.清心的人有福了,因為他們必得見神.使人和睦的人有福了,因為他們稱為神的兒子.為義受逼迫的人有福了,因為天國是他們的.人若因我辱罵你們,逼迫你們,捏造各樣壞話毀謗你們,你們就有福了.應當歡喜快槳,因為你們在天上的賞賜是大的.在你們以前的先知,人也是這樣逼迫他們.」(五1-12)
\newline
\newline
四福音中有三段經文──(約十三－十七,太五－七,太廿一－廿五)是三篇主耶穌有系統的講道.今所讀太五－七章登山寶訓的開端,是很有意思的.大文學家譚馬士積遜稱之為名著中最優美完善的佳作.韋氏大詞典作者描寫登山寶訓,為全世界律法的綱領；非常之週詳,深刻,高超.總而言之,神話是靈,是生命.飽學之士不會覺其膚淺,婦孺亦不覺其艱深.是天國之王宣講憲章.開頭講九福,因祂的國是有福之國.這段經文重在遵行.約翰說:小子們哪,不要被人誘惑,行義的才是義人.」「行事為人就當與蒙召的恩相稱.」(`弗四1)「總要在言語,行為,愛心,清潔上,都作信徒的榜樣.」(提前四12)
\newline
\newline
本段經文有多種說法:普通稱為「八福」,可與廿三章的「八禍」,遙相比對.有人認為10至12節乃是複述上文,應祗「七福」.也有人主張「九福」,因連11節所述,也是一福.更有人加上12節之歡喜快樂也是福份.故稱「十福」,符合十全十美之意.今我採取「九福」的說法.與加拉太書提到聖靈所結的「九果」相合.屬靈生命的歷程共分三段,1.生命的印證.2.生命的成長.3.生命之操練.每段包括三「福」,共為「九福」.如寶塔之九層,登殿臺階之九級.其上為耶穌基督.巴不得此經文像一幅圖畫,如希臘寓言；把銀燈放進漁夫之屋內,其中的器具都變成煉淨的銀子,我們都因主的活得改變.
\newline
\newline
一,倒空罪惡虛心受益
\newline
\newline
訊息之始,首重心靈.虛心的人有福了,為天國是他們的.佛說重戒口,回教尚潔心,但無非是枝葉.基督之國從心靈開始.虛心是指倒空罪惡.經上說:保守你的心勝過保守一切.因為心如平原跑馬,易放難收.箴言說:我兒,你要將心歸我,可見這是神對人首要的要求.虛心,是要將內裡的污穢倒盡,然後能接受福祉.人要前往美國,必先知道入境的手續.入天國何嘗不然,也應知道神所要求的.倒空罪惡是蒙福之第一步.王下四章以利沙先知門徒之妻往鄰舍取空瓶,然後將家內一瓶油將空瓶一一倒滿.主穌在迦拿婚筵行第一個神蹟──以水變酒,也使用了擺在門口的六口空石缸.基甸帶著神所揀選的三百人,對數不清的米甸人爭戰.要他們預備三百個空瓶子和三百個火把,藉此勝過了數不過來的強敵.這都證明神能使用倒空的器皿.經上有許多例證:如法老,巴蘭,亞干,猶大,掃羅等,未嘗不認識神,甚至為主作過工.但是悔而不改,不肯倒空罪惡,就落在悲慘的結局.
\newline
\newline
從前有個地主向僕人問得救的問題,僕人告訴他要進豬欄裡住幾天,在其中默想禱告,就能得救.主人以為太難.後來主人病重,重新考慮,需要得救.對僕人說:寧願進入了豬欄.僕人說:如今倒不需要這樣行,因為主人已經能謙卑下來接受救恩.「耶和華說:你們來,我們彼此辯論,你們的罪雖像硃紅,必變成雪白,雖紅如丹顏,必白如羊毛.」「我塗抹了你的過犯,像厚雲消散,我塗抹了你的罪惡,如薄雲滅沒.」(賽一18,四四22)「耶穌說:有錢財的人進神的國,是何等的難哪.駱駝穿過針的眼,比財主進神的國還容易呢.聽見的人說:這樣誰能得救呢?耶穌說:在人所不能的事,在神卻能.」(路十八24-27)駱駝能夠穿過針眼嗎?假如把駱駝變成一根線就通行無阻了.
\newline
\newline
二,靈裡貧窮以信為糧
\newline
\newline
(賽六六2):「但我所看顧的,就是虛心痛悔,因我話而戰兢的人.」聖經小字虛心原文作貧窮.(路六20):「你們貧窮的人有福了.」原文有貧乏至討飯度日之意.中世紀對此有錯誤的解釋,提倡修士貧窮,實行乞食苦行.今紐約街頭尚有該等修女.神叫我們悟知靈裡貧窮,需對祂絕對信靠.不應像老底嘉教會,以為樣樣都有；誰知是瞎眼可憐,赤身露體的.萬不可自以為屬靈的太多.進神的國,應該有根本改變,明瞭自己一無所有.
\newline
\newline
人向神禱告,不可像法利賽人,以為自己比別人好,把自己的好處盡量在神面前誇耀.有人也有類似這樣的毛病.自以為屬靈,所得亮光比別人多,得的啟示比別人大.就輕看他人,唯我獨尊；祗有他們一撮人是教會,成為無宗派的宗派.阿拉伯人有句諺語:觀看麥田的穗子,就知道神賜福給誰.越豐滿的頭越低垂,惟有虛空的,站得最直.我們在神面前,看清楚自己的敗壞.當效法施洗約翰的榜樣──祂必興旺,我必衰微.我們對神的態度,要如路加十八章那寡婦對不義的官要求.又好像雅各在毗努伊勒神的使者摔跤.你若不給我祝福,我就不讓你過去.人活著不是單靠食物,乃靠神口裡所出的一切話.該知道我們若沒有神的恩典就活不成.要好像乞丐般天天仰望祂.列王記下末後記載,約雅斤終身常在巴比倫王面前喫飯,王賜他所需用的食物,日日賜他一分.這就是我們寫照,祂要我們天天見祂的面,與祂親近,以祂信實為糧.
\newline
\newline
三,學生榜樣承受天國
\newline
\newline
主進入耶路撒冷,經上記著說:「要對錫安的居民說,看哪!你的王來到你這裡,是溫柔的,又騎著驢,就是騎著驢駒子.」榮耀的主,又是何等謙卑.主也曾對人發出呼召:「凡勞苦擔重擔的人,可以到我這裡來,我就使你們得安息.我心裡柔和謙卑,你們當負我的軛,學我的樣式,這樣,你們心裡就必得享安息.」(太十一28-29)此外,又看保羅怎樣為祂作見證:「祂本有神的形像,不以自己與神同為強奪的.反倒虛己,取了奴僕的形象,成為人的樣式.既有人的樣子,就自己卑微,存心順服,以至於死,且死在十字架上.」(腓二6-8)我們怎樣成為虛心的人,得以承受天國,應效法主虛己謙卑.
\newline
\newline
主說:「你們要進窄門,因為引到滅亡,那門是寬的,路是大的,進去的人也多；引到永生,那門是窄的,路是小的,找著的人也少.」大路多人行,小路卻要細意尋找.找到後樂意謙虛跟主腳步.凡聽見主話就去行的,好比聰明人將房子蓋在磐石上.雨淋,水沖,風吹,撞著那房子,房子總不倒塌.
\newline
\newline
主又對門徒說:「你們這小群,不要懼怕,因為你們的父樂意把國賜給你們.」(路十二32)這是何等安慰的應許.
\newline
\newline
某財主到其家中喫飯.把一切指給牧師看.問牧師說:「像我這樣,還需向神求什麼?」牧師說:「應當求神賜給你謙卑.」
\newline
\newline
屬靈道路的開始,首需倒空罪惡.在主前以一無所有的態度仰望主,跟隨主的腳步承受天國.
\newline
\newline


\newpage

\section{第48屆~港九培靈會~史祈生~第二講}
\label{sec:1060}
\textbf{生命的印證(二)}
\newline
\newline
連結: \href{https://www.hkbibleconference.org/session-message/view/1060}{\texttt{https://www.hkbibleconference.org/session-message/view/1060}}
\newline
\newline
\hyperref[sec:1059]{< < < PREV SERMON < < <}
~
\hyperref[sec:toc]{[返目錄]}
~
\hyperref[sec:1061]{> > > NEXT SERMON > > >}
\newline
\newline
馬太福音 5:4-5:4
\newline
\begin{longtable}{cl}
\hline
\hline
章節 & 經文 (和合本修訂版)\\
\hline
5:4 & \begin{tabularx}{0.7\textwidth}{X} 哀慟的人有福了！因為他們必得安慰。 \end{tabularx} \\ \\
[1ex]
\hline
\hline
\end{longtable}
經文:哀慟的人有福了,因為他們必得安慰.(太五4)
\newline
\newline
今日繼續講生命印證之第二步──哀慟的人有福.人是有感情的,發乎中,而形於外；悲則哀哭,喜則歡笑.這感情是超越種族語言的界限,互相傳達彼此心意的.哭不分中外,笑亦不分東西.一哭一笑,都能使人了解其人當時心境,這是與生俱來的本能.神予人屬靈生命,也必有憑據.哀慟乃是生命印證之二,乃傷心至極的境地.其內容可分作三方面:
\newline
\newline
一,世上的哀慟
\newline
\newline
神人摩西說:「我們經過的日子,都在你震怒之下,我們度盡的年歲,好像一聲嘆息.」(詩九十9)摩西描寫人生宛似一首哀砍,飽含憂患.(伯五7)說:「人生在世必遇患難,如同火星飛騰.」主耶穌也曾說:「你們在世上有苦難.」人好像被圈在其中,不能超越似的；主為人子,也有同樣的經歷.有人說:笑有廿八種,哭有卅六種.果然有各樣不同的形式,留心細察,自然知所判別.世間愁苦事,何止千萬般,故俗世稱為苦海.所謂生離死別,怨恨嗔痴；樂極尚且生悲,其他更勿論矣.創世記四十三至四十五章,寫約瑟與他兄弟異地相逢,情不自禁,覓地痛哭.及至見他胞弟便雅憫,更相擁抱頭,盡情大慟.約瑟貴為宰相,但當至情流露,亦難免痛哭流涕.又見約拿單與大衛,情逾手足.當掃羅逼迫大衛,使他不得不亡命他鄉時,兩人為分離傷痛,也是動人肺腑心腸的.
\newline
\newline
世人悲哀憂愁事,有可免,也有不能避免的.後者如天災人禍,使人措不及防；如近日各地發生之大地震,死難者百數十萬人；變生肘腋,無可奔逃.又記得第二次大戰時,香港淪陷.那時我適在本港,相信在座許多人也記憶猶新.那種光景真是不堪回首,難以避免的；但也有屬於杞人憂天,庸人自擾的.昨見報載一則趣聞:巴西里約熱內盧,有婦人之夫離家日久不歸,後來得接其夫死訊,以為是真,急購新衣鞋等前往,政府為其殮葬.事畢,婦人尚悲傷不已,後其夫卻無恙歸來,原來是馮京馬涼之誤.不特賠上許多眼淚,而且新衣新鞋被別人穿了,親夫卻穿舊衣敝屣回來.豈非枉哭嗎?又有一個八歲女孩,在園中玩耍正高興時,突然放聲大哭.母親來哄她,她也不肯聽,問她哭什麼?原來她行近井旁,見井深,嚇得發抖.媽媽安慰她,她說:我想到我幸而沒有掉下去.但我將來要出嫁,生了孩子,長到我這樣大；又到這裡玩耍,掉下去了怎辦?這女孩雖可笑,不過正有不少人也像她呢?其餘也有些人看了連環圖,聽了鬼故事,而為小說中的人物流淚的.這都屬於自尋煩惱之類,愚痴可哂!
\newline
\newline
但也有些悲哀是罪有應得的,(哀五16):「我們犯罪了,我們有禍了.」罪的工價乃是死,罪像繩索般捆人,如影子跟隨不捨.年前香港某工程公司,剉閉清盤,被揭發做假帳作弊；牽好幾位社會名流,一旦現身法庭,真是羞得無地可容.又如美國水門醜聞,被拆穿後,也把醜惡的面目揭露出來,他們都得到了應得的報應.此外如賭狗,賭馬,股票投機；漲則喜極而狂,敗則悲哀欲死,都屬自作自受,絕非是福祉的.
\newline
\newline
二,靈內的哀慟
\newline
\newline
(結九4),提到因城中行可憎之事歎息哀哭的人,在額上畫記號為記.(詩五一7):「神阿,憂傷痛悔的心,你必不輕看.」憂傷痛悔的原文是現在式,表示一直繼續之意.這樣的哀慟與世上哀慟不同,主說:哀慟的人有福了,他們必得著安慰.作天國的子民,對罪惡應當敏感,因他有生命的印證.把死人丟下坑中,不會感覺難受,他也不會掙扎；若把活人掉下就不同了,他必大呼救命!孩子一生下來就啼哭,哭聲洪亮,正表示這生命是強壯的.
\newline
\newline
以賽亞在聖殿中,見到主坐在高高的寶座上.就說:「禍哉!我滅亡了,因我是嘴唇不潔的,又住在唇嘴不潔的民中.」奧古士丁悔改之後,以詩篇五篇十七節作為座右銘,幫助自己時常儆醒.許多人卻不如此,反將犯罪作為歡笑,以羞恥為誇耀.傳道人與信徒,亦有具雙重人格的.外面一套屬靈的口吻,裡面卻別有居心,有婦人在街市上買菜,一面與攤販講價爭吵,但一面卻又滿口感謝神,讚美主,真是不知所謂.有人說猶大卅兩銀子賣主那麼便宜,如果我賣就起碼三百兩,其實許多人卻為了一針一線而出賣主.常為罪憂傷痛悔的人有福了.心裡有罪的人必不亨通,承認己罪的得蒙赦免.
\newline
\newline
其次是感激的眼淚.因體念到主的大愛,祂離開天上,降世人間,成為人子,取了奴僕的樣式；存心順服以至於死,且死在十字架上.當我們擘餅記念主的時候,想到主的身體為我擘開,祂的寶血為我流出,就想到我的罪和虧欠.路加福音第七章:耶穌進到一個法利賽人的家喫飯,有婦人將玉瓶打破把香膏抹耶穌的腳；把眼淚濕了祂的腳,又用頭髮來擦乾.法利賽人心中盤算:「如果祂是個先知,必知道挨近祂的,是個罪人.」耶穌就說:「她許多罪都赦免了,因為她的愛多.」我們是否也有這因感激而流的眼淚呢?
\newline
\newline
更有是按神的旨意憂愁,為別人憂傷.保羅說:「我是大有憂愁,心裡時常傷痛.為我骨肉之親,就是自己被咒詛,與基督分離,我也願意.」當保羅與以弗所長老惜別時說:「你們應當儆醒,記念我三年之久,畫夜不住的流淚,勸戒你們.」有像保羅這樣肯流淚付代價的人教會才有進步.今日何等需要肯這樣流淚的人,為眾教會,為主的工作,為許多失喪的靈魂,心裡哀慟懇切代求.從前有個地方遭遇大火,發覺尚有個孩子被困火場中.眾人正焦急,忽見一青年奮不身衝入救人,各人為之激動不已.及至青年從火中出來時,卻又使人從興奮變成失望；因他搶救的,不是孩子的生命,卻是盛著財寶的箱子.
\newline
\newline
我們的靈魂和肉體常常爭戰.羅馬書第七章很清楚描述這情形.甚至說:「我真是苦阿,誰能救我脫離這取死的身體呢?」巴不得我們也能為自己的軟弱失敗,像與主在客西馬尼園一同儆醒.
\newline
\newline
三,完全的安慰
\newline
\newline
信徒與世人在世上有同樣經歷.同會遭受到各種艱難痛苦.但一個醫生進到病人的家中,信主的和不信的就有顯著的分別.縱使所遭遇的是絕症,仍能顯示出有安慰和盼望.跳出現實的範圍,得著主自己的安慰.正如詩廿三篇說:「我雖然行過死蔭的幽谷,也不怕遭害,因為你與我同在,你的杖,你的竿,都安慰我.」
\newline
\newline
我們事奉主的人,可能碰到種種難處,筋疲力竭,灰心失望.但主給我們見到勞苦的果效,就可以滿意足.不但如此,神為我們預備了極重無比的榮耀,十字架可換取公義的冠冕.到那日,主擦乾我們一切的眼淚；也使在今日得嘗天恩的滋味.耶穌說有一個人,從耶路撒冷下耶利哥,路上遇強盜被打得半死.利未人和祭司經過都不理會,後來一個好撒瑪利亞人來了；用油與酒為他纏裹又扶上驢背,送到旅舍,估羅週全.這就是好鄰舍的寫照.
\newline
\newline
有一位傳道人從外國休假回紐約.抵達見群眾搖旗爽岸歡迎,心中詫異.後來才知道是羅斯福總統與他同時抵達,那些人不過是來歡迎總統的,原來接他的祗有一個人.我們還未到天家,將來我們都要到天家,一定有許多人這樣歡迎我們.但願我們見主面時,得祂稱讚是一個「又良善又忠心的僕人就夠了.」
\newline
\newline


\newpage

\section{第48屆~港九培靈會~史祈生~第三講}
\label{sec:1061}
\textbf{生命的印證(三)}
\newline
\newline
連結: \href{https://www.hkbibleconference.org/session-message/view/1061}{\texttt{https://www.hkbibleconference.org/session-message/view/1061}}
\newline
\newline
\hyperref[sec:1060]{< < < PREV SERMON < < <}
~
\hyperref[sec:toc]{[返目錄]}
~
\hyperref[sec:1062]{> > > NEXT SERMON > > >}
\newline
\newline
馬太福音 5:5-5:5
\newline
\begin{longtable}{cl}
\hline
\hline
章節 & 經文 (和合本修訂版)\\
\hline
5:5 & \begin{tabularx}{0.7\textwidth}{X} 謙和的人有福了！因為他們必承受土地。 \end{tabularx} \\ \\
[1ex]
\hline
\hline
\end{longtable}
經文:溫柔的人有福了,因為他們必承受地土.(太五5)
\newline
\newline
今日講到生命印證的第三步──溫柔的人有福.連著過去兩天訊息來講,好像雨過天晴,有豁然開朗的景象.有人說:今日世界科學雖然已到登峰造極,但人性卻淪於原始的野蠻.弱肉強食,巧取豪奪.溫柔即是軟弱,如何能生存於今世呢?有人以為當日的猶太人,受轄制於羅馬帝國,只好忍氣吞聲.所以耶穌才用這話安慰他們,他們若能委曲求全,就可免罹大禍.耶穌說:「溫柔的人有福了.」雖然現今的世界充滿了邪惡殘暴,正像挪亞,羅得的日子一樣.我們正需要有溫柔的心,面對現實.請留意以下去對溫柔的正確答案.
\newline
\newline
一,作新造的人
\newline
\newline
耶穌說:「溫柔的人有福.」是指屬靈方面來說的.「若有人在基督裡,他就是新造的人,舊事已過,都變成新的了.」(林後五17)有了主的生命,就是溫柔的人.有人生下來性情粗暴,易動肝火；逢人要讓他三分,動輒爭至面紅耳赤.有些人在家中做父親恍如法官,威嚴得令人發抖；一進門便鴉鵲無聲,一切以自我為中心.若非神的大能,誰能改變他呢?也有些人是好好先生,一味謙和退讓；與人無干,與世無爭,甚至比牧師的脾氣更好.他還需要重生嗎?但主耶穌對尼哥底母說:「你們必須重生.」像尼哥底母生於聖地,身為猶太人之高尚階級.主耶穌說:「我說你們必須重生,你不要以為希奇.」還有一種人涵養功夫做得到家,真是爐火純青,出神入化；正如漢末劉琨,使女捧湯翻污他的朝服,他不但泰然自若,反問女手有無燙傷.宋韓琦夜草文書,使女秉燭侍立.竟燃及公髯,驚懼而走.公止之日:「如今爾當學識如何持燈燭矣,何須走避.」人修養至如此地步,是否可與得救有關?
\newline
\newline
作神的子民,須有新生命.(結卅六26):「我也要賜給你一個新心,將新靈放在你們裡面；又從你們的肉體中除掉石心,賜給你肉心.」羅馬書二章末又說:外面作猶太人的不是猶太人,外面肉身的割禮也不是割禮.許多人在教會中,不過是糊塗人；正如某人到佛教辦的學校報名,在表上信仰項內,填上「佛教」；在另一欄問「何時開始?」上「祖傳」.試問信基督可否祖傳?他信父親信的,父親信牧師所信的.進一步問牧師信甚麼?祗有問牧師自己.我們不是這樣,乃是在生命中有得救經歷.不是按著儀文的舊樣,乃照著心靈的新樣.不是藉著傳統,習慣,和宗派主張；乃是在心靈之中有再造的生命.
\newline
\newline
二,結出屬靈果子
\newline
\newline
加拉太書提到聖靈所結的九樣果子,第八是「溫柔」,既是聖靈所結的果子,就不在乎人自己的努力.馬可四章28節:「地生五穀,是出於自然的.先發苗,後長穗,再後穗上結成飽滿的子粒.」這是植物正常生長的程序.孟子載:「揠苗助長.」非徒無益反而害,故需往下深深地扎根.詩一3:「他要像一棵樹栽在溪水旁,按時候結果子.」基督徒往不明白這道理,倚靠自己的力量.例如發脾氣,平常小心謹慎,處處留神,倒還應付得來.但遇到突如其來的刺激,就無法控制,以至原形畢露.靠聖靈而結出的果子,不是與自己完全無關.應當順服聖靈的帶領幫助,與聖靈合作.「保守自己常在神的愛中.」(猶21)約書亞領百姓過約但河,吩咐十二支派將十二塊石頭放在河底.表明我們的老我要被埋葬.石頭雖仍存在,但已被深埋於河底,我們也要藏在聖靈裡面,與聖靈合作,就能多結果子.摩西在埃及,見自己弟兄被欺壓,就動手將埃及人打死了以為憑自己的力量可以拯救百姓.卻不料被血氣的行動耽誤了四十年.亞當被造,神吩咐在園中看守修理,故必須與神合作,才能多結聖靈的果子.
\newline
\newline
在客西馬尼園,主勸勉門徒要儆醒.有人被過犯所勝,應把他挽回過來.保羅說:要彼此勸勉,不可停止聚會.這些都要與神合作,彼此幫助.彼得在五旬節之前,也是憑血氣行事；拔出刀來將大祭司僕人的耳朵削下,但為主所禁止.五旬節聖靈丕降,彼得憑聖靈的大能講道,有三千,五千人歸主；刺透人的心,出於聖靈的工作.神的靈像寶貝放在瓦器裡,就顯出莫大的能力,是出於神.雷子約翰,被主改變成為愛的使徒；晚年傳講訊息:「小子們,你們要彼此相愛.」這都是聖靈所作的工,故不可消滅聖靈的感動.
\newline
\newline
「摩西為人謙和,勝過世上眾人.」這是神為他作的見證.因為他將二百萬百姓帶領出來,四十年之久,藉聖靈學習了功課.「不輕易發怒的,勝過勇士；治服己心的,強如取城.」(箴十六32)拿破崙說:「世上有兩種力量,一是刀,一是愛.但愛勝於刀.」
\newline
\newline
天國之王謙謙和和地騎驢駒子,進入耶路撒冷.主從未動用一一卒,來擴張著祂的國度；但憑著聖靈的大能；仁愛為兵器.祂的境界遍於全地；祂的國度,永無窮盡.
\newline
\newline
香港與紐約同是地價高貴的地方,因為地少人稠,能擁有一片地皮,就價值不菲了.講到得地之法有二:一是高價收購,二是用武力佔領.耶穌不用這些,給我們另一個得到地之法:「溫柔的人有福了,因為他們必承受地土.」用溫柔之法可得土地,豈不奇哉?
\newline
\newline
古籍流傳；老子將死,門人圍侍床前,請留最後訓言.老子張口示眾人,問何所有?徒曰唯舌存焉.老子曰:「齒剛易折,舌柔長存,」印度聖雄甘地,採不抵抗政策,以柔克剛,卒贏得印度獨立自主.參創世記廿六章:以撒住在南地基拉耳,有一口他父親亞伯拉罕所挖的水井,被非利士人塞住了.他又另挖水井.如是者再三忍耐.到一地方又挖一井,敵人他不再為井爭競了.他給那井起名叫利河伯,意思說:「耶和華現在給我們寬闊之地,我們必在這裡昌盛.」這都證明主的說話是確實的,溫柔的人才能領受主的賞賜.
\newline
\newline
聖經又記載:亞伯拉罕帶同侄兒羅得,離開迦勒底吾珥,進入迦南.後來二人的牛羊財物漸多,亞伯蘭的牧人和羅得的牧人相爭.亞伯蘭就對羅得說:「你我不可相爭,你的牧人和我的牧人也不可相爭,因為我們是骨肉.......你向左,我就向右；你向右,我就向左.」羅得看中了約但河平原,在那裡定居.漸漸挪移帳棚直到所多瑪.所多瑪,蛾摩拉二城,卒為神降火所燬.亞伯拉罕卻得著神應許:「你舉目向東西南北觀看,凡你所看見的一切地,我都賜給你和你的後裔,直至永遠.」(創世記十三章)豈不又再證明,承受地土的,是溫柔的人.
\newline
\newline
文豪托爾斯泰曾講一寓言,他說:人生得地幾何?譬為一個大地主宣告將地廉值售出,祗需付一定額之價值,可得到從日出開始起步,日落跑回原處,所圍住的地.有一個人果然拼命奔跑,擴闊圈子；直至日落時將達終點,但已力盡不支,為免前功盡廢,只有仆身以手觸地終點.自云目的已達,誰料從此長眠不起.海明威以老人與海一書,蜚聲於世.寫一老人在深海捕得大魚,洋洋得意,繫於船邊,力划回程.及至抵步時,大魚祗剩一副骨骼而已.
\newline
\newline
這兩個故事都叫我們思想,我們在主裡面,得地有多少?你是否追求世界的一切?日光之下轉瞬虛空.約書亞記云:還有多地未得,我們如何?
\newline
\newline


\newpage

\section{第48屆~港九培靈會~史祈生~第四講}
\label{sec:1062}
\textbf{生命的成長(一)}
\newline
\newline
連結: \href{https://www.hkbibleconference.org/session-message/view/1062}{\texttt{https://www.hkbibleconference.org/session-message/view/1062}}
\newline
\newline
\hyperref[sec:1061]{< < < PREV SERMON < < <}
~
\hyperref[sec:toc]{[返目錄]}
~
\hyperref[sec:1095]{> > > NEXT SERMON > > >}
\newline
\newline
馬太福音 2:1-6:34
\newline
\begin{longtable}{cl}
\hline
\hline
章節 & 經文 (和合本修訂版)\\
\hline
2:1 & \begin{tabularx}{0.7\textwidth}{X} 在希律作王的時候，耶穌生在猶太的伯利恆。有幾個博學之士從東方來到耶路撒冷，說： \end{tabularx} \\ \\ \relax
2:2 & \begin{tabularx}{0.7\textwidth}{X} 「那生下來作猶太人之王的在哪裡？我們在東方看見他的星，特來拜他。」 \end{tabularx} \\ \\ \relax
2:3 & \begin{tabularx}{0.7\textwidth}{X} 希律王聽見了，就心裡不安；耶路撒冷全城的人也都不安。 \end{tabularx} \\ \\ \relax
2:4 & \begin{tabularx}{0.7\textwidth}{X} 他就召集了祭司長和民間的文士，問他們：「基督該生在哪裡？」 \end{tabularx} \\ \\ \relax
2:5 & \begin{tabularx}{0.7\textwidth}{X} 他們說：「在猶太的伯利恆。因為有先知記著： \end{tabularx} \\ \\ \relax
2:6 & \begin{tabularx}{0.7\textwidth}{X} 『猶大地的伯利恆啊，你在猶大諸城中並不是最小的；因為將來有一位統治者要從你那裡出來，牧養我以色列民。』」 \end{tabularx} \\ \\ \relax
2:7 & \begin{tabularx}{0.7\textwidth}{X} 於是，希律暗地裡召了博學之士來，查問那星是甚麼時候出現的， \end{tabularx} \\ \\ \relax
2:8 & \begin{tabularx}{0.7\textwidth}{X} 就派他們往伯利恆去，說：「你們去仔細尋訪那小孩子，找到了就來報信，我也好去拜他。」 \end{tabularx} \\ \\ \relax
2:9 & \begin{tabularx}{0.7\textwidth}{X} 他們聽了王的話就去了。忽然，在東方所看到的那顆星在前面引領他們，一直行到小孩子所在地方的上方就停住了。 \end{tabularx} \\ \\ \relax
2:10 & \begin{tabularx}{0.7\textwidth}{X} 他們看見那星，就非常歡喜； \end{tabularx} \\ \\ \relax
2:11 & \begin{tabularx}{0.7\textwidth}{X} 進了房子，看見小孩子和他母親馬利亞，就俯伏拜那小孩子，揭開寶盒，拿出黃金、乳香、沒藥，作為禮物獻給他。 \end{tabularx} \\ \\ \relax
2:12 & \begin{tabularx}{0.7\textwidth}{X} 因為在夢中得到主的指示，不要回去見希律，他們就從別的路回自己的家鄉去了。 \end{tabularx} \\ \\ \relax
2:13 & \begin{tabularx}{0.7\textwidth}{X} 他們走後，忽然主的使者在約瑟夢中向他顯現，說：「起來！帶著小孩子和他母親逃往埃及，住在那裡，等我的指示；因為希律要搜尋那小孩子來殺害他。」 \end{tabularx} \\ \\ \relax
2:14 & \begin{tabularx}{0.7\textwidth}{X} 約瑟就起來，連夜帶著小孩子和他母親往埃及去， \end{tabularx} \\ \\ \relax
2:15 & \begin{tabularx}{0.7\textwidth}{X} 住在那裡，直到希律死了。這是要應驗主藉先知所說的話：「我從埃及召我的兒子出來。」 \end{tabularx} \\ \\ \relax
2:16 & \begin{tabularx}{0.7\textwidth}{X} 希律見自己被博學之士愚弄，極其憤怒，差人將伯利恆城裡和四境所有的男孩，根據他向博學之士仔細查問到的時間，凡兩歲以內的，都殺盡了。 \end{tabularx} \\ \\ \relax
2:17 & \begin{tabularx}{0.7\textwidth}{X} 這就應驗了耶利米先知所說的話： \end{tabularx} \\ \\ \relax
2:18 & \begin{tabularx}{0.7\textwidth}{X} 「在拉瑪聽見號咷大哭的聲音，是拉結哭她兒女；她不肯受安慰，因為他們都不在了。」 \end{tabularx} \\ \\ \relax
2:19 & \begin{tabularx}{0.7\textwidth}{X} 希律死了以後，在埃及，忽然主的使者在約瑟夢中向他顯現， \end{tabularx} \\ \\ \relax
2:20 & \begin{tabularx}{0.7\textwidth}{X} 說：「起來，帶著小孩子和他母親回以色列地去！因為要殺害這小孩子的人已經死了。」 \end{tabularx} \\ \\ \relax
2:21 & \begin{tabularx}{0.7\textwidth}{X} 約瑟就起來，帶著小孩子和他母親進入以色列地去。 \end{tabularx} \\ \\ \relax
2:22 & \begin{tabularx}{0.7\textwidth}{X} 但是他因聽見亞基老繼承他父親希律作了猶太王，怕到那裡去；又在夢中得到主的指示，就往加利利境內去了。 \end{tabularx} \\ \\ \relax
2:23 & \begin{tabularx}{0.7\textwidth}{X} 他們到了一座城，名叫拿撒勒，就住在那裡。這是要應驗先知所說的話：「他將稱為拿撒勒人。」 \end{tabularx} \\ \\ \relax
3:1 & \begin{tabularx}{0.7\textwidth}{X} 在那些日子，施洗的約翰出來，在猶太的曠野宣講： \end{tabularx} \\ \\ \relax
3:2 & \begin{tabularx}{0.7\textwidth}{X} 「你們要悔改！因為天國近了。」 \end{tabularx} \\ \\ \relax
3:3 & \begin{tabularx}{0.7\textwidth}{X} 這人就是以賽亞先知所說的：「在曠野有聲音呼喊著：預備主的道，修直他的路。」 \end{tabularx} \\ \\ \relax
3:4 & \begin{tabularx}{0.7\textwidth}{X} 這約翰身穿駱駝毛的衣服，腰束皮帶，吃的是蝗蟲和野蜜。 \end{tabularx} \\ \\ \relax
3:5 & \begin{tabularx}{0.7\textwidth}{X} 那時，耶路撒冷、全猶太和全約旦河地區的人，都到約翰那裡去， \end{tabularx} \\ \\ \relax
3:6 & \begin{tabularx}{0.7\textwidth}{X} 承認他們的罪，在約旦河裡受他的洗。 \end{tabularx} \\ \\ \relax
3:7 & \begin{tabularx}{0.7\textwidth}{X} 約翰看見許多法利賽人和撒都該人也來受洗，就對他們說：「毒蛇的孽種啊，誰指示你們逃避那將要來的憤怒呢？ \end{tabularx} \\ \\ \relax
3:8 & \begin{tabularx}{0.7\textwidth}{X} 你們要結出果子來，和悔改的心相稱。 \end{tabularx} \\ \\ \relax
3:9 & \begin{tabularx}{0.7\textwidth}{X} 不要自己心裡說：『我們有亞伯拉罕為祖宗。』我告訴你們，神能從這些石頭中給亞伯拉罕興起子孫來。 \end{tabularx} \\ \\ \relax
3:10 & \begin{tabularx}{0.7\textwidth}{X} 現在斧子已經放在樹根上，凡不結好果子的樹就砍下來，丟在火裡。 \end{tabularx} \\ \\ \relax
3:11 & \begin{tabularx}{0.7\textwidth}{X} 我是用水給你們施洗，叫你們悔改；但那在我以後來的，能力比我更大，我就是給他提鞋子也不配，他要用聖靈與火給你們施洗。 \end{tabularx} \\ \\ \relax
3:12 & \begin{tabularx}{0.7\textwidth}{X} 他手裡拿著簸箕，要揚淨他的穀物，把麥子收在倉裡，把糠用不滅的火燒盡。」 \end{tabularx} \\ \\ \relax
3:13 & \begin{tabularx}{0.7\textwidth}{X} 當時，耶穌從加利利來到約旦河，到了約翰那裡，請約翰為他施洗。 \end{tabularx} \\ \\ \relax
3:14 & \begin{tabularx}{0.7\textwidth}{X} 約翰想要阻止他，說：「我應該受你的洗，你怎麼到我這裡來呢？」 \end{tabularx} \\ \\ \relax
3:15 & \begin{tabularx}{0.7\textwidth}{X} 耶穌回答他：「暫且這樣做吧，因為我們理當這樣履行全部的義。」於是約翰就依了他。 \end{tabularx} \\ \\ \relax
3:16 & \begin{tabularx}{0.7\textwidth}{X} 耶穌受了洗，隨即從水裡上來。天忽然為他開了，他看見神的靈降下，彷彿鴿子落在他身上。 \end{tabularx} \\ \\ \relax
3:17 & \begin{tabularx}{0.7\textwidth}{X} 這時，天上有聲音說：「這是我的愛子，我所喜愛的。」 \end{tabularx} \\ \\ \relax
4:1 & \begin{tabularx}{0.7\textwidth}{X} 當時，耶穌被聖靈引到曠野，受魔鬼的試探。 \end{tabularx} \\ \\ \relax
4:2 & \begin{tabularx}{0.7\textwidth}{X} 他禁食四十晝夜，後來就餓了。 \end{tabularx} \\ \\ \relax
4:3 & \begin{tabularx}{0.7\textwidth}{X} 那試探者進前來對他說：「你若是神的兒子，叫這些石頭變成食物吧。」 \end{tabularx} \\ \\ \relax
4:4 & \begin{tabularx}{0.7\textwidth}{X} 耶穌卻回答說：「經上記著：『人活著，不是單靠食物，乃是靠神口裡所出的一切話。』」 \end{tabularx} \\ \\ \relax
4:5 & \begin{tabularx}{0.7\textwidth}{X} 魔鬼就帶他進了聖城，叫他站在聖殿頂上， \end{tabularx} \\ \\ \relax
4:6 & \begin{tabularx}{0.7\textwidth}{X} 對他說：「你若是神的兒子，就跳下去！因為經上記著：『主要為你命令他的使者，用手托住你，免得你的腳碰在石頭上。』」 \end{tabularx} \\ \\ \relax
4:7 & \begin{tabularx}{0.7\textwidth}{X} 耶穌對他說：「經上又記著：『不可試探主—你的神。』」 \end{tabularx} \\ \\ \relax
4:8 & \begin{tabularx}{0.7\textwidth}{X} 魔鬼又帶他上了一座很高的山，將世上的萬國和萬國的榮華都指給他看， \end{tabularx} \\ \\ \relax
4:9 & \begin{tabularx}{0.7\textwidth}{X} 對他說：「你若俯伏拜我，我就把這一切賜給你。」 \end{tabularx} \\ \\ \relax
4:10 & \begin{tabularx}{0.7\textwidth}{X} 耶穌說：「撒但，退去！因為經上記著：『要拜主—你的神，惟獨事奉他。』」 \end{tabularx} \\ \\ \relax
4:11 & \begin{tabularx}{0.7\textwidth}{X} 於是，魔鬼離開了耶穌，立刻有天使來伺候他。 \end{tabularx} \\ \\ \relax
4:12 & \begin{tabularx}{0.7\textwidth}{X} 耶穌聽見約翰下了監，就退到加利利去； \end{tabularx} \\ \\ \relax
4:13 & \begin{tabularx}{0.7\textwidth}{X} 後來離開拿撒勒，往迦百農去，住在那裡。那地方靠海，在西布倫和拿弗他利地區。 \end{tabularx} \\ \\ \relax
4:14 & \begin{tabularx}{0.7\textwidth}{X} 這是要應驗以賽亞先知所說的話： \end{tabularx} \\ \\ \relax
4:15 & \begin{tabularx}{0.7\textwidth}{X} 「西布倫，拿弗他利，沿海的路，約旦河的東邊，外邦人的加利利— \end{tabularx} \\ \\ \relax
4:16 & \begin{tabularx}{0.7\textwidth}{X} 那坐在黑暗裡的百姓看見了大光；坐在死蔭之地的人有光照耀他們。」 \end{tabularx} \\ \\ \relax
4:17 & \begin{tabularx}{0.7\textwidth}{X} 從那時候，耶穌開始宣講，說：「你們要悔改！因為天國近了。」 \end{tabularx} \\ \\ \relax
4:18 & \begin{tabularx}{0.7\textwidth}{X} 耶穌沿著加利利海邊行走，看見兩兄弟，就是那叫彼得的西門和他弟弟安得烈，正往海裡撒網；他們本是打魚的。 \end{tabularx} \\ \\ \relax
4:19 & \begin{tabularx}{0.7\textwidth}{X} 耶穌對他們說：「來跟從我，我要叫你們得人如得魚一樣。」 \end{tabularx} \\ \\ \relax
4:20 & \begin{tabularx}{0.7\textwidth}{X} 他們立刻捨了網，跟從他。 \end{tabularx} \\ \\ \relax
4:21 & \begin{tabularx}{0.7\textwidth}{X} 耶穌從那裡往前走，看見另外兩兄弟，就是西庇太的兒子雅各和他弟弟約翰，同他們的父親西庇太在船上補網，耶穌就呼召他們。 \end{tabularx} \\ \\ \relax
4:22 & \begin{tabularx}{0.7\textwidth}{X} 他們立刻捨了船，辭別父親，跟從了耶穌。 \end{tabularx} \\ \\ \relax
4:23 & \begin{tabularx}{0.7\textwidth}{X} 耶穌走遍加利利，在各會堂裡教導人，宣講天國的福音，醫治百姓各樣的疾病。 \end{tabularx} \\ \\ \relax
4:24 & \begin{tabularx}{0.7\textwidth}{X} 他的名聲傳遍了敘利亞。那裡的人把一切病人，就是有各樣疾病和疼痛的、被鬼附的、癲癇的、癱瘓的，都帶了來，耶穌就治好了他們。 \end{tabularx} \\ \\ \relax
4:25 & \begin{tabularx}{0.7\textwidth}{X} 當時，有一大群人從加利利、低加坡里、耶路撒冷、猶太、約旦河的東邊，來跟從他。 \end{tabularx} \\ \\ \relax
5:1 & \begin{tabularx}{0.7\textwidth}{X} 耶穌看見這一群人，就上了山，坐下後，門徒到他跟前來， \end{tabularx} \\ \\ \relax
5:2 & \begin{tabularx}{0.7\textwidth}{X} 他開口教導他們說： \end{tabularx} \\ \\ \relax
5:3 & \begin{tabularx}{0.7\textwidth}{X} 「心靈貧窮的人有福了！因為天國是他們的。 \end{tabularx} \\ \\ \relax
5:4 & \begin{tabularx}{0.7\textwidth}{X} 哀慟的人有福了！因為他們必得安慰。 \end{tabularx} \\ \\ \relax
5:5 & \begin{tabularx}{0.7\textwidth}{X} 謙和的人有福了！因為他們必承受土地。 \end{tabularx} \\ \\ \relax
5:6 & \begin{tabularx}{0.7\textwidth}{X} 飢渴慕義的人有福了！因為他們必得飽足。 \end{tabularx} \\ \\ \relax
5:7 & \begin{tabularx}{0.7\textwidth}{X} 憐憫人的人有福了！因為他們必蒙憐憫。 \end{tabularx} \\ \\ \relax
5:8 & \begin{tabularx}{0.7\textwidth}{X} 清心的人有福了！因為他們必得見神。 \end{tabularx} \\ \\ \relax
5:9 & \begin{tabularx}{0.7\textwidth}{X} 締造和平的人有福了！因為他們必稱為神的兒子。 \end{tabularx} \\ \\ \relax
5:10 & \begin{tabularx}{0.7\textwidth}{X} 為義受迫害的人有福了！因為天國是他們的。 \end{tabularx} \\ \\ \relax
5:11 & \begin{tabularx}{0.7\textwidth}{X} 「人若因我辱罵你們，迫害你們，捏造各樣壞話毀謗你們，你們就有福了！ \end{tabularx} \\ \\ \relax
5:12 & \begin{tabularx}{0.7\textwidth}{X} 要歡喜快樂，因為你們在天上的賞賜是很多的。在你們以前的先知，人也是這樣迫害他們。」 \end{tabularx} \\ \\ \relax
5:13 & \begin{tabularx}{0.7\textwidth}{X} 「你們是地上的鹽。鹽若失了味，怎能叫它再鹹呢？它不再有用，只好被丟在外面，任人踐踏。 \end{tabularx} \\ \\ \relax
5:14 & \begin{tabularx}{0.7\textwidth}{X} 你們是世上的光。城造在山上是不能隱藏的。 \end{tabularx} \\ \\ \relax
5:15 & \begin{tabularx}{0.7\textwidth}{X} 人點燈，不放在斗底下，而是放在燈臺上，就照亮一家的人。 \end{tabularx} \\ \\ \relax
5:16 & \begin{tabularx}{0.7\textwidth}{X} 你們的光也要這樣照在人前，叫他們看見你們的好行為，把榮耀歸給你們在天上的父。」 \end{tabularx} \\ \\ \relax
5:17 & \begin{tabularx}{0.7\textwidth}{X} 「不要以為我來是要廢掉律法和先知。我來不是要廢掉，而是要成全。 \end{tabularx} \\ \\ \relax
5:18 & \begin{tabularx}{0.7\textwidth}{X} 我實在告訴你們，就是到天地都廢去，律法的一點一畫也不能廢去，直到一切都實現。 \end{tabularx} \\ \\ \relax
5:19 & \begin{tabularx}{0.7\textwidth}{X} 所以，無論誰廢掉這誡命中最小的一條，又教導人也這樣做，他在天國裡要稱為最小的。但無論誰遵行並如此教導人的，他在天國裡要稱為大。 \end{tabularx} \\ \\ \relax
5:20 & \begin{tabularx}{0.7\textwidth}{X} 我告訴你們，你們的義若不勝過文士和法利賽人的義，絕不能進天國。」 \end{tabularx} \\ \\ \relax
5:21 & \begin{tabularx}{0.7\textwidth}{X} 「你們聽過有對古人說：『不可殺人』；『凡殺人的，必須受審判。』 \end{tabularx} \\ \\ \relax
5:22 & \begin{tabularx}{0.7\textwidth}{X} 但是我告訴你們：凡向弟兄動怒的，必須受審判；凡罵弟兄是廢物的，必須受議會的審判；凡罵弟兄是白痴的，必須遭受地獄的火。 \end{tabularx} \\ \\ \relax
5:23 & \begin{tabularx}{0.7\textwidth}{X} 所以，你在祭壇上獻祭物的時候，若想起有弟兄對你懷恨， \end{tabularx} \\ \\ \relax
5:24 & \begin{tabularx}{0.7\textwidth}{X} 就要把祭物留在壇前，先去跟弟兄和好，然後來獻祭物。 \end{tabularx} \\ \\ \relax
5:25 & \begin{tabularx}{0.7\textwidth}{X} 你同告你的冤家還在路上，就要趕快與他講和，免得他把你送交給法官，法官交給警衛，你就下在監裡了。 \end{tabularx} \\ \\ \relax
5:26 & \begin{tabularx}{0.7\textwidth}{X} 我實在告訴你，就是有一個大文錢還沒有還清，你也絕不能從那裡出來。」 \end{tabularx} \\ \\ \relax
5:27 & \begin{tabularx}{0.7\textwidth}{X} 「你們聽過有話說：『不可姦淫。』 \end{tabularx} \\ \\ \relax
5:28 & \begin{tabularx}{0.7\textwidth}{X} 但是我告訴你們：凡看見婦女就動淫念的，這人心裡已經與她犯姦淫了。 \end{tabularx} \\ \\ \relax
5:29 & \begin{tabularx}{0.7\textwidth}{X} 若是你的右眼使你跌倒，就把它挖出來，丟掉。寧可失去身體中的一部分，也不讓整個身體被扔進地獄。 \end{tabularx} \\ \\ \relax
5:30 & \begin{tabularx}{0.7\textwidth}{X} 若是你的右手使你跌倒，就把它砍下來，丟掉。寧可失去身體中的一部分，也不讓整個身體下地獄。」 \end{tabularx} \\ \\ \relax
5:31 & \begin{tabularx}{0.7\textwidth}{X} 「又有話說：『無論誰休妻，都要給她休書。』 \end{tabularx} \\ \\ \relax
5:32 & \begin{tabularx}{0.7\textwidth}{X} 但是我告訴你們：凡休妻的，除非是因不貞的緣故，否則就是使她犯姦淫了；人若娶被休的婦人，也是犯姦淫了。」 \end{tabularx} \\ \\ \relax
5:33 & \begin{tabularx}{0.7\textwidth}{X} 「你們又聽過有對古人說：『不可背誓，所起的誓總要向主謹守。』 \end{tabularx} \\ \\ \relax
5:34 & \begin{tabularx}{0.7\textwidth}{X} 但是我告訴你們：甚麼誓都不可起。不可指著天起誓，因為天是神的寶座。 \end{tabularx} \\ \\ \relax
5:35 & \begin{tabularx}{0.7\textwidth}{X} 不可指著地起誓，因為地是他的腳凳；也不可指著耶路撒冷起誓，因為耶路撒冷是大君王的京城。 \end{tabularx} \\ \\ \relax
5:36 & \begin{tabularx}{0.7\textwidth}{X} 又不可指著你的頭起誓，因為你不能使一根頭髮變黑變白。 \end{tabularx} \\ \\ \relax
5:37 & \begin{tabularx}{0.7\textwidth}{X} 你們的話，是，就說是；不是，就說不是。若再多說，就是出於那惡者。」 \end{tabularx} \\ \\ \relax
5:38 & \begin{tabularx}{0.7\textwidth}{X} 「你們聽過有話說：『以眼還眼，以牙還牙。』 \end{tabularx} \\ \\ \relax
5:39 & \begin{tabularx}{0.7\textwidth}{X} 但是我告訴你們：不要與惡人作對。有人打你的右臉，連另一邊也轉過去由他打。 \end{tabularx} \\ \\ \relax
5:40 & \begin{tabularx}{0.7\textwidth}{X} 有人想要告你，要拿你的裡衣，連外衣也由他拿去。 \end{tabularx} \\ \\ \relax
5:41 & \begin{tabularx}{0.7\textwidth}{X} 有人強迫你走一里路，你就跟他走二里。 \end{tabularx} \\ \\ \relax
5:42 & \begin{tabularx}{0.7\textwidth}{X} 有求你的，就給他；有向你借貸的，不可推辭。」 \end{tabularx} \\ \\ \relax
5:43 & \begin{tabularx}{0.7\textwidth}{X} 「你們聽過有話說：『要愛你的鄰舍，恨你的仇敵。』 \end{tabularx} \\ \\ \relax
5:44 & \begin{tabularx}{0.7\textwidth}{X} 但是我告訴你們：要愛你們的仇敵，為那迫害你們的禱告。 \end{tabularx} \\ \\ \relax
5:45 & \begin{tabularx}{0.7\textwidth}{X} 這樣，你們就可以作天父的兒女了。因為他叫太陽照好人，也照壞人；降雨給義人，也給不義的人。 \end{tabularx} \\ \\ \relax
5:46 & \begin{tabularx}{0.7\textwidth}{X} 你們若只愛那愛你們的人，有甚麼賞賜呢？就是稅吏不也是這樣做嗎？ \end{tabularx} \\ \\ \relax
5:47 & \begin{tabularx}{0.7\textwidth}{X} 你們若只請你弟兄的安，有甚麼比別人強呢？就是外邦人不也是這樣做嗎？ \end{tabularx} \\ \\ \relax
5:48 & \begin{tabularx}{0.7\textwidth}{X} 所以，你們要完全，如同你們的天父是完全的。」 \end{tabularx} \\ \\ \relax
6:1 & \begin{tabularx}{0.7\textwidth}{X} 「你們要謹慎，不可故意在人面前表現虔誠，叫他們看見，若是這樣，就不能得你們天父的賞賜了。 \end{tabularx} \\ \\ \relax
6:2 & \begin{tabularx}{0.7\textwidth}{X} 「所以，你施捨的時候，不可叫人在你前面吹號，像那假冒為善的人在會堂裡和街道上所做的，故意要得人的稱讚。我實在告訴你們，他們已經得了他們的賞賜。 \end{tabularx} \\ \\ \relax
6:3 & \begin{tabularx}{0.7\textwidth}{X} 你施捨的時候，不要讓左手知道右手所做的， \end{tabularx} \\ \\ \relax
6:4 & \begin{tabularx}{0.7\textwidth}{X} 好使你隱祕地施捨；你父在隱祕中察看，必然賞賜你。」 \end{tabularx} \\ \\ \relax
6:5 & \begin{tabularx}{0.7\textwidth}{X} 「你們禱告的時候，不可像那假冒為善的人，愛站在會堂裡和十字路口禱告，故意讓人看見。我實在告訴你們，他們已經得了他們的賞賜。 \end{tabularx} \\ \\ \relax
6:6 & \begin{tabularx}{0.7\textwidth}{X} 你禱告的時候，要進入內室，關上門，向那在隱祕中的父禱告；你父在隱祕中察看，必將賞賜你。 \end{tabularx} \\ \\ \relax
6:7 & \begin{tabularx}{0.7\textwidth}{X} 你們禱告，不可像外邦人那樣重複一些空話，他們以為話多了必蒙垂聽。 \end{tabularx} \\ \\ \relax
6:8 & \begin{tabularx}{0.7\textwidth}{X} 你們不可效法他們。因為在你們祈求以前，你們所需要的，你們的父早已知道了。」 \end{tabularx} \\ \\ \relax
6:9 & \begin{tabularx}{0.7\textwidth}{X} 「所以，你們要這樣禱告：『我們在天上的父：願人都尊你的名為聖。 \end{tabularx} \\ \\ \relax
6:10 & \begin{tabularx}{0.7\textwidth}{X} 願你的國降臨；願你的旨意行在地上，如同行在天上。 \end{tabularx} \\ \\ \relax
6:11 & \begin{tabularx}{0.7\textwidth}{X} 我們日用的飲食，今日賜給我們。 \end{tabularx} \\ \\ \relax
6:12 & \begin{tabularx}{0.7\textwidth}{X} 免我們的債，如同我們免了人的債。 \end{tabularx} \\ \\ \relax
6:13 & \begin{tabularx}{0.7\textwidth}{X} 不叫我們陷入試探；救我們脫離那惡者。因為國度、權柄、榮耀，全是你的，直到永遠。阿們！』 \end{tabularx} \\ \\ \relax
6:14 & \begin{tabularx}{0.7\textwidth}{X} 你們若饒恕人的過犯，你們的天父也必饒恕你們； \end{tabularx} \\ \\ \relax
6:15 & \begin{tabularx}{0.7\textwidth}{X} 你們若不饒恕人，你們的天父也必不饒恕你們的過犯。」 \end{tabularx} \\ \\ \relax
6:16 & \begin{tabularx}{0.7\textwidth}{X} 「你們禁食的時候，不可像那假冒為善的人，臉上帶著愁容；因為他們蓬頭垢面，故意讓人看出他們在禁食。我實在告訴你們，他們已經得了他們的賞賜。 \end{tabularx} \\ \\ \relax
6:17 & \begin{tabularx}{0.7\textwidth}{X} 你禁食的時候，要梳頭洗臉， \end{tabularx} \\ \\ \relax
6:18 & \begin{tabularx}{0.7\textwidth}{X} 不要讓人看出你在禁食，只讓你隱祕中的父看見；你父在隱祕中察看，必然賞賜你。」 \end{tabularx} \\ \\ \relax
6:19 & \begin{tabularx}{0.7\textwidth}{X} 「不要為自己在地上積蓄財寶；地上有蟲子咬，能鏽壞，也有賊挖洞來偷。 \end{tabularx} \\ \\ \relax
6:20 & \begin{tabularx}{0.7\textwidth}{X} 要在天上積蓄財寶；天上沒有蟲子咬，不會鏽壞，也沒有賊挖洞來偷。 \end{tabularx} \\ \\ \relax
6:21 & \begin{tabularx}{0.7\textwidth}{X} 因為你的財寶在哪裡，你的心也在哪裡。」 \end{tabularx} \\ \\ \relax
6:22 & \begin{tabularx}{0.7\textwidth}{X} 「眼睛是身體的燈。你的眼睛若明亮，全身就光明； \end{tabularx} \\ \\ \relax
6:23 & \begin{tabularx}{0.7\textwidth}{X} 你的眼睛若昏花，全身就黑暗。你裡面的光若黑暗了，那黑暗是何等大呢！」 \end{tabularx} \\ \\ \relax
6:24 & \begin{tabularx}{0.7\textwidth}{X} 「一個人不能服侍兩個主；他不是恨這個愛那個，就是重這個輕那個。你們不能又服侍神，又服侍瑪門。」 \end{tabularx} \\ \\ \relax
6:25 & \begin{tabularx}{0.7\textwidth}{X} 「所以，我告訴你們，不要為你們的生命憂慮吃甚麼喝甚麼，或為你們的身體憂慮穿甚麼。生命不勝於飲食嗎？身體不勝於衣裳嗎？ \end{tabularx} \\ \\ \relax
6:26 & \begin{tabularx}{0.7\textwidth}{X} 你們看一看那天上的飛鳥，也不種也不收，也不在倉裡存糧，你們的天父尚且養活牠們。你們不比飛鳥貴重得多嗎？ \end{tabularx} \\ \\ \relax
6:27 & \begin{tabularx}{0.7\textwidth}{X} 你們哪一個能藉著憂慮使壽數多加一刻呢？ \end{tabularx} \\ \\ \relax
6:28 & \begin{tabularx}{0.7\textwidth}{X} 何必為衣裳憂慮呢？你們想一想野地裡的百合花是怎麼長起來的：它也不勞動也不紡線。 \end{tabularx} \\ \\ \relax
6:29 & \begin{tabularx}{0.7\textwidth}{X} 然而我告訴你們，就是所羅門極榮華的時候，他所穿戴的還不如這些花的一朵呢！ \end{tabularx} \\ \\ \relax
6:30 & \begin{tabularx}{0.7\textwidth}{X} 你們這小信的人哪！野地裡的草今天還在，明天就丟在爐裡，神還給它這樣的妝飾，何況你們呢？ \end{tabularx} \\ \\ \relax
6:31 & \begin{tabularx}{0.7\textwidth}{X} 所以，不要憂慮，說：『我們吃甚麼？喝甚麼？穿甚麼？』 \end{tabularx} \\ \\ \relax
6:32 & \begin{tabularx}{0.7\textwidth}{X} 這都是外邦人所求的。你們需要這一切東西，你們的天父都知道。 \end{tabularx} \\ \\ \relax
6:33 & \begin{tabularx}{0.7\textwidth}{X} 你們要先求神的國和他的義，這些東西都要加給你們了。 \end{tabularx} \\ \\ \relax
6:34 & \begin{tabularx}{0.7\textwidth}{X} 所以，不要為明天憂慮，因為明天自有明天的憂慮；一天的難處一天當就夠了。」 \end{tabularx} \\ \\
[1ex]
\hline
\hline
\end{longtable}
經文:饑渴慕義的人有福了,因為他們必得飽足.(太2-6)
\newline
\newline
過去三次聚會,思想過屬靈生命歷程的第一層──生命的印證.今進入第二步,生命的成長.基督徒不單得著生命,還需繼續長大,進到完全的地步.初為父母,緊張快樂難以形容.有個準備作爸爸的,在家中等待醫院來的電話；鈴一響,急不及待抓起聽筒,聲音發顫:「孩子生下來呢!我是否做爸爸?還是做媽媽?」對方回答:「平安生下來了,但無論是男是女,你總歸是爸爸,不會是媽媽的.」小孩子生下來之後,一月,二月,半載,一年......卻不長大,那是個嚴重的問題.必須請醫生研究,察看究竟.神的兒女也有經過三年,五年,十年八年,依然故我的.怎能不叫父神擔憂煩惱?因此,屬靈生命固需印證,也要進步成長.
\newline
\newline
一,饑渴
\newline
\newline
(路六21)「你們饑餓的人有福了.」比不上馬太此處經文:「饑渴慕義人有福了.」包括了飢饑和乾渴兩方面的難處.巴勒斯坦每年雨量稀少,舊約常提到有旱災發生.所以聽眾很容易了解主耶穌的話.在這裡自來水方便,很難想像缺少雨水,那種難耐的情景.世界各地,類此情況的很多.廿多年前在福州,洗澡水是黃的.前些時我在漢城,發現大多數水廁是缺水供應的.我曾到南美被擄,某地是終年不雨的,不知雨傘為何物.饑渴的人,才會領會到饑渴的苦況.有過饑腸轆轆經驗的人,就知道飢餓是怎樣一回事.有一位教師出題目給學生作文,題為「窮苦的家庭」.他的學生作答──「我的家庭真是窮苦阿,爸爸窮,媽媽也窮；廚師和老媽子也窮；花匠和司機更窮到不得了.......「小孩子不知所作的是何種妙文.我們要屬靈生命的長大,首需知道饑渴之苦,生命有重大的需要；那才能大大張口,吸取營養.今天許多基督徒,在屬靈上是有胃病的,所以不能吸收,也不長進.
\newline
\newline
詩篇講論渴慕神,形容為「如鹿渴慕溪水一樣.」世人在許多事上,表現如饑如渴.人為逐利奔忙,形形色色,錙銖必較.追求異性,日以繼夜,魂牽夢縈.競逐名位,不惜挖空心思,不擇手段.神的兒女在屬靈事上,更應有渴慕的心,詩人說:「我向你舉手,我的心渴想你,如乾旱之地盼雨一樣.」(詩一四三6)得生命之後,自然要藉著禱告與神交往.起初不肯開口禱告,過後就成為啞吧不會對神開聲.一位牧師在會中請人領禱,久久不見動靜,終於一人起來說:「會祈禱的今天未來,所以無人禱告.」有人禱告唸廿六個字母,據云,上帝已包含之意了；又有說:「我的祈禱跟日牧師所祈的一樣,阿們.」總之,形形式式,無奇不有.不禱告的基督徒,生命何來長進.應有如饑渴如的心,禱告親近神.魔鬼是不怕人作工,也不怕人聚會,卻頂怕人禱告.因為禱告是武器,所以仇敵最怕人祈禱.我常在半夜起來禱告.某次被鄰房的人聽到,次日,他問我為何這樣長篇.我說:「譬如一對男女在熱愛中,他們兩情相悅；雖然講來講去不外這些,但總是談不厭,說不完.」我們往往對人說話,絮絮滔滔,口若懸河；但對人作見證,向神祈禱,卻訥訥不能出口.
\newline
\newline
我們從前不懂愛神,反而貪愛世界.保羅說:「只是我先前以為與我有益的,我現在因基督都當作有損的......我為祂已經丟棄萬事,看作糞上,為要得著基督.」(腓三7-9)「我作孩子的時候,話語像孩子,心思像孩子意念像孩子,既成了人,就把孩子的事丟棄了.」(林前十三11)人漸漸長大起來,心思興趣也自然改換.有人問可不可以吸煙,飲酒,跳舞,等問題.認為聖經上又沒有明文禁止.基督徒何必這樣拘束?但是聖經上說:「豈不知你們是神的殿,神的靈住在你們裡頭麼?」(林前三16)可以或不可以,不需多講,也可知道的了.孩童時放風箏,騎竹馬是頂自然的事；但當長大成人還作這些,豈不成了怪事麼?這不是犯罪的問題,乃是與身份不相稱.有人說:基督徒生活太枯燥,不抽煙,打牌,喝酒等；簡直毫無人生樂趣,如何過活.人生需及時行樂呀!等到將來老大才信還未遲,這是完全錯誤的觀念.人有生命,必有憑據,必會表現出如饑如渴的對神傾慕.
\newline
\newline
二,慕義
\newline
\newline
我們所慕的義,並非如中國古籍所載之仁義道德,乃指主耶穌基督；祂是我們尋求的目標.但我們往往所求的,是主的工作,主的恩賜,而非主自己.這是不夠的,必須要如饑如渴地尋求祂,祂就是我們的義,是我們的滿足.義又指神的道,彼得稱「傳義道的挪亞」,就是指此.「看你們學習的工夫,本該作師傅,誰知還得有人將神聖言小學的開端,另教導你們,並且成了那必須喫奶,不能喫乾糧的人.」(來五12)故當離開基督道理的開端,竭力進到完全地步.不要儘在曠野繞圈子,應當進入迦南享受應許的美地.豈可以永遠祗食嗎哪,也要喫迦南地的土產.平均來說,一百個基督徒僅有兩個能夠自己讀經,從神的話吸取養料.這個比率無疑太過可憐.大部份都是祗能喫奶的,就是別人消化過的.從前餵孩子,是由大人先把飯放在口中嚼爛,然後哺給嬰孩.許多營養都失去了.但大多數人就是懶於自己咀嚼,祗喫別人加工的東西.有人聽道也揀好聽悅耳的,稍為難接受的就打瞌睡,這都是不長進的表明.喫乾糧就是能下功夫,細加咀嚼,嚐出味道來,有些人對我說,我從少就讀聖經,會考又有好成績.這也未必能作準.可能他也一竅不通.其實一望而知；有時台上講舊約,他卻盡在新約那部份翻.宣佈讀猶大書,又問幾章幾節.勤翻的聖經,必然外面好裡面爛,有完整聖經的,必是懶讀聖經的人.牧師來探訪時,才把聖經撲淨塵埃,甚至打開一看,原來失蹤許久的眼鏡夾在那裡.有些所謂教會學校,其實包括先生,學生,及校董在內,對聖經仍是一無所知,也一無所愛.那得有如鹿渴慕溪水的態度,追求接受主的話呢?
\newline
\newline
三,得飽足
\newline
\newline
「得飽足」這字是進行式的,是指一直進行下去,不是一陣子的.天天享受神,得著滿足.「因祂使心裡渴慕的人,得以知足；使心裡饑餓的人,得飽美物.」(詩一○七9)「我在床上記念你,在夜更的時候思想你,我的心就像飽足了骨髓肥油.」(詩六十三6)渴慕神,思想神的話,就能心靈得到飽足,中國人講究飲食,舉世知名.所謂世界廚師,僅有一個半.半個是法國人,一個屬於中國.外國人得利用蒸汽發動機械；我們早已懂得用它來蒸叉燒飽.廣東人更厲害吃的分有天上飛的,除飛機外都可以喫.地上有四隻腳的；除桌椅之外,也都可喫.菜式花樣翻新,層出不窮.人們稱「食得是福」.但基督徒的福份不在乎肚腹.是天天享受主,真是有主萬事足.有人問我作牧師到底如何?我說:「還用問嗎?幾十年來,我從不懊悔.愈做愈開心,越過越滿足.真是世上難找,難以形容.」主說:「人若渴了,可以到我這裡來喝.信我的人,就如經上所說,從他腹中流出活水的江河.」(約七37-38)主所賜的平安喜樂,沒有人能奪去；我們住在主裡面,何等奇妙,何等平安.祂用各樣美物使我們所願得以知足.縱使在世上也有不如意的事,但可在凡事上得到幫助.詩人說:「除你以外,在天上我還有誰呢?除你以外,在地上我也沒有所愛慕的.」有了主耶穌,何等滿足.
\newline
\newline
得救之後,當求長進.愛慕神的話,祈禱親近主.生命自然會改變.將來在神的國中,不再饑,不再渴；天上莫大的榮耀,會得預嚐矣.
\newline
\newline


\newpage

\section{第48屆~港九培靈會~史祈生~第五講}
\label{sec:1095}
\textbf{生命的成長(二)}
\newline
\newline
連結: \href{https://www.hkbibleconference.org/session-message/view/1095}{\texttt{https://www.hkbibleconference.org/session-message/view/1095}}
\newline
\newline
\hyperref[sec:1062]{< < < PREV SERMON < < <}
~
\hyperref[sec:toc]{[返目錄]}
~
\hyperref[sec:1097]{> > > NEXT SERMON > > >}
\newline
\newline
馬太福音 5:7-5:7
\newline
\begin{longtable}{cl}
\hline
\hline
章節 & 經文 (和合本修訂版)\\
\hline
5:7 & \begin{tabularx}{0.7\textwidth}{X} 憐憫人的人有福了！因為他們必蒙憐憫。 \end{tabularx} \\ \\
[1ex]
\hline
\hline
\end{longtable}
經文:憐恤的人有福了,因為他們必蒙憐恤.(太五7)
\newline
\newline
昨天講過饑渴慕義的人,必得飽足.今天講憐恤的人有福,他們必得憐恤.前者是內心的經驗,後者是外面的行動.誠於中則形於外；來源滔滔汨汨,輸送自必廣闊綿長.
\newline
\newline
一,蒙憐恤
\newline
\newline
憐恤乃是神的本性,保羅最能體驗出神這份心腸.「我今日成了何等人,乃是神的恩才成的.......在罪人中我是個罪魁,然而我蒙了憐憫.」大衛說:「父親怎樣憐恤他的兒女,耶和華也怎樣憐恤敬畏祂的人.」(詩一○三13)「我喜愛憐恤,不喜愛祭祀.」(何六6)主耶穌降生前,聖靈感動,祭司撒迦利亞豫言說:「因我們神憐憫的心腸,叫清晨的日光從高天臨到我們.」又回想創世記所載:當神用洪水毀壞全地之後,就以虹為記號,與人類立約,不再被洪水滅絕,表明神的顧念和憐恤.
\newline
\newline
我們要知道神的心,可以從兩個窗戶望見──有個孩子問媽媽,怎知道媽的心?媽媽叫他看自己的眼睛.孩子就看見媽媽眼中祗有小弟弟.神的心意也可從祂的看顧體驗出來.一個孩子看母親從市上回來,拿著各物又抱著西瓜.孩子哭嚷:「媽媽愛西瓜不愛孩子.」媽媽哄他,「你愛媽媽?」「孩子說愛.」「愛多少呢?」孩子不能形容,用手比一比說,「愛這麼多.」媽再問:「祗這麼多麼?」孩子將兩手大張:「我愛媽媽這樣多.」媽媽感到滿足了,放下所有的抱起孩子,主在十字架上,就表明神大張兩手環抱我們的大愛.
\newline
\newline
「我必憐恤他們,如同人憐恤服事自己的兒子.」(瑪三17)這又寫出了神憐恤的另一面.世上祗有兒子服事父母,焉有父母服事兒子的呢?神的憐恤就是如此.有些姐妹把心思意念都集中在兒子身上,彷彿一切都為了他.雅各回來與以掃重逢之後,以掃要他同往.雅各說:「求我主在僕人前頭走,我要量著在我面前群畜和孩子的力量慢慢的前行.」神恩待我們也是這樣,我們本是軟弱,容易跌倒.我們的父恩待憐恤我們,量我們的力量慢慢前行.
\newline
\newline
翰福音十章描寫主是好牧人,把羊領出來,既放出自己的羊,就在前頭走.人生走屬靈道路,也須經過迷茫曠野,黑暗森林；但有主作我們的牧者就不必懼怕.香港近年改變得面目全非,若非友好領路,我將無所適從.但主滿有慈悲憐憫.壓傷蘆葦祂不折斷,將殘火祂不吹滅.我們沒有誰能避免軟弱失敗,就是少年人也疲乏困倦,強壯的也全然跌倒,但等候耶和華的必重新得力.所以無論誰若失敗跌倒了,當立刻爬起來.因主使我們坦然無懼到祂施恩座前,蒙憐恤,得恩惠,作隨時幫助.
\newline
\newline
二,憐恤人
\newline
\newline
孟子說:「惻隱之心人皆有之.」是神把祂的本性造在人裡面.據蘇東坡曾費五百文購得一所房子,繼而夜間聞有人在戶外哭泣.是原住者售屋後無處棲身故然.東坡見而憐之,歸其故居而解其困.林肯總統一日策馬途中,見道旁有一豬陷身泥澤掙扎不已,不忍.下騎設法拯起之,雖衣服玷污,在所不計.可知人皆有憐恤人的心,但被罪過所包圍,漸漸便麻木了.別人有了過犯錯失,需要體恤赦免,用溫柔慈愛的心把他挽回過來.主教導我們的禱詞:「免我們的債,如同我們免了人的債.」若不懂得赦免人,就不配事奉我們的主.美國文學家史蒂文生每晚有家庭禮拜,又有背誦主禱文的習慣.一次,史蒂文生悄悄地離開了聚會,他妻子以為因身體不適；原來他覺得心中仍對弟兄懷怨,故自感不配唸主禱文.
\newline
\newline
馬太福音十八章記載彼得曾問耶穌說:「弟兄得罪我,饒怨他到七次夠了麼?」主耶穌說:「不是到七次,乃是七十個七次.」隨後又講一個比喻:有僕人欠主一千萬銀子,不能償還.俯伏哀求寬免,主人答應了他.出門後遇見一個欠他十兩銀子的人,卻捏住他的喉嚨逼他還債.主人知道後,勃然大怒,將他交給掌刑的治罪.登山寶訓中主一再提到兩種義──你們的義,若不勝於文士和法利賽人的義,斷不能進天國.我們既然蒙過憐恤,也應憐恤他人.見弟兄無衣無食,不可塞住憐憫人的心.彼此相愛,不是祗在言語舌頭上；更要在行為誠實上.對一個窮人說:「你平平安安的回去罷,願你喫得飽,穿得暖.」這是無濟於事的.若不愛看得見的弟兄,怎能夠愛那看不見的神呢?有時我們對弟兄生了怨嫌,而牢不肯放,甚至避面不在同一地方聚會.這又何必呢?當知道主怎樣愛我們,寬容我們的過犯.有人存有偏見,祗愛那強壯,富足,尊貴的；不愛那些軟弱,貧窮,可憐的.那就是雅各書二章所說的:按外貌待人,偏心,用惡意斷定人了.我們斷不可如此,反倒應該看顧那些窮苦,疾病,居住本屋的人.這才配得作為天父的兒女.
\newline
\newline
在某次聚會中,一人戴帽坐於前列,旁人為之側目.牧師詫問其故.他說:「我來此聚會多時,無人理會我；今日若非作此打扮,牧師亦不會問及我吧?.」事奉主的應該俯就卑微的人,要體恤扶助,有病的為他們祈禱,經濟困難的幫助他們解決.香港物質繁榮,有錢財的人很多,對世上窮乏的應當樂於施與,而且作的時候,左手不要讓右手知道.許多為主作工的工人,好像牛馬良辛勞,卻得不到充足的草料.一位牧師出門代步的馬很肥,會友問他的馬為何這樣肥壯?牧師說:「這馬是我養的,所以很肥；我是你們養的,不過很瘦.」馬肥人瘦,每見中外皆然.也有人雖不富足,但喜歡打扮得體面些.因不是要得人的可憐,若出自愛心的扶助,方為可貴.
\newline
\newline
三,有福了
\newline
\newline
經上說:「憐恤人的人有福了,因為他們必蒙憐恤.」在曠野的人群中,一個小童帶著五餅二魚.因見到眾人有需要,就把來交在主手中.擘開餵飽了五千人.自己喫飽,也滿足了眾人,這是雙重得福.以利亞先知在遍地大旱饑荒之時,來到撒勒法寡婦家,母子二人僅有一把麵,和一點油.但遵著先知的話作成了餅給先知,此後瓶裡的油,罈中的麵毫不短缺.所以箴言十九章17節:「憐憫貧窮的,就是借給耶和華.」借給神,比存入世上銀行的利息還優厚.當遭遇患難的時候,神必定施恩搭救.能夠憐恤人的,是踏入生命長進的第二步,他有了神的性情,也是神所賜福的人.
\newline
\newline
(雅二13)又說:「因為那不憐憫人的,也要受無憐憫的審判,憐憫原是向審判誇勝.」當審判的時候,憐憫人的可坦然無懼站在神面前.主也是何等的憐恤人,祂提過:當人子在榮耀中降臨的時候,祂要將萬民分別出來,如同牧人把綿羊山羊分成左右兩邊.祂要對右邊的說:「你們這蒙我父賜福的,可以承受那創世以來為你們預備的國.因為我餓了,你們給我喫；渴了,你們給我喝；我作客旅,你們留我住；.......」義人回答說:「主阿,我們甚麼時候作過這些事呢?」主說:「我實在告訴你們,這些事你們既作在我弟兄中一個小的身上,就是作為我身上了.」(參太廿五31-46)
\newline
\newline
(西三12-15)「所以你們既是神的選民,聖潔蒙愛的人,就要存憐憫,恩慈,謙虛,溫柔,忍耐的心......在這一切之外,要存著愛心,愛心是聯絡全德的.」
\newline
\newline
今天我們到神面前蒙了憐憫.就當存憐憫人的心；用行為,見證,把人帶到主面前.在靈程上,主也必一生一世以恩惠慈愛相加.不在乎定意,也不在乎奔跑,祗在發憐憫的神.
\newline
\newline


\newpage

\section{第48屆~港九培靈會~史祈生~第六講}
\label{sec:1097}
\textbf{生命的成長(三)}
\newline
\newline
連結: \href{https://www.hkbibleconference.org/session-message/view/1097}{\texttt{https://www.hkbibleconference.org/session-message/view/1097}}
\newline
\newline
\hyperref[sec:1095]{< < < PREV SERMON < < <}
~
\hyperref[sec:toc]{[返目錄]}
~
\hyperref[sec:1099]{> > > NEXT SERMON > > >}
\newline
\newline
馬太福音 5:8-5:8
\newline
\begin{longtable}{cl}
\hline
\hline
章節 & 經文 (和合本修訂版)\\
\hline
5:8 & \begin{tabularx}{0.7\textwidth}{X} 清心的人有福了！因為他們必得見神。 \end{tabularx} \\ \\
[1ex]
\hline
\hline
\end{longtable}
經文:清心的人有福了,因為他們必得見神.(太五8)
\newline
\newline
今天講到屬靈生命成長的第三步,是清心的人有福.是內心裡面的扎根.許多人有馬大的事奉,卻忽略了馬利亞的安靜.所以成長的再進一步是往下扎根.基督徒有扎根和不扎根的；不扎根的,經不起雨淋,水沖,風吹,試煉一來就跌到了.唯有扎根的,好像利巴嫩的香柏樹,雖屢經狂風暴雨,仍穩站在高處.末世聖徒也當如此.倫敦有一次遭遇大旱,許多植物都漸漸枯毀.唯有一棵葡萄樹仍然繼續生長結果；原來它的根深深地扎下,直透到泰晤士河.我們傳道人有個毛病,就是祗顧忙於工作,忽略了有規律的靈修讀經,這是一層往裡面深入的工作.我們的心好像儲水池子,必須引水道收集,然後能流通供應各方.儲水池又需保持清潔,否則所流到的地方,都不能蒙福,反受其害.所以須注重內心的建設,今就此思想三件事:
\newline
\newline
一,聖潔的國度
\newline
\newline
「惟有你們是被揀選的族類,是有君尊的祭司,是聖潔的國度,是屬神的子民」(彼前二9).「聖潔」一辭,新舊庫譯本及現代福音譯本繙作「心中純潔」,比較更為恰當.不但內心要潔淨,更要純真.摩西律法中有非常清楚的潔淨條例,(利十10-11):「耶和華曉諭亞倫說,你和你兒子進會幕的時候；清酒濃酒都不可喝,免得你們死亡.這要作為你們世世代代的定例,使你們可以將聖的,俗的,潔淨的,不潔淨的,分別出來.」(又參申廿二9-11):便吩咐不可把兩樣種子種在葡萄園；不可並用牛驢耕地；不可穿羊毛細麻兩樣攙雜料作的衣服.這都表明神對百姓清潔純真的要求.利未記對聖潔最為著重,書中共提到一三一次.
\newline
\newline
到現在,紐約猶太人店號所發售的牛羊肉,都蓋有特別的記號.表示是經由拉比所認可的,所以他們對於食物的清潔特別嚴格.但這卻不過是外表的,照新約的亮光,神要我們追求內裡的聖潔.正如保羅說:「這些規條使人徒有智慧之名......其實在克制肉體的情慾上,是毫無功效.」請參(提前一5；提後二22；彼前一22).又在(提前三9；提後一3),都提到「清潔的良心」是非常重要.主在登山寶訓上:「有吩咐古人的話說:『不可殺人.』只是我告訴你們,凡向弟兄動怒的難免受審判.又有話說:『不可姦淫.』只是我告訴你們,看見婦女就動淫念的,心裡已經與她犯姦淫了.」「光明與黑暗有甚麼相通,信的和不信的有甚麼相干.你們務要從他們中間出來；與他們分別,不要沾不潔淨的物,我就收納你們.我要作你們的父,你們要作我的兒女.」(參林後六14-18),這都指神的子民分別聖俗,勿沾不潔.
\newline
\newline
大衛禱告說:「神阿,求你鑑察我,知道我的心思；試煉我,知道我的意念.看在我裡面有甚麼惡行沒有,引導我走永生的道路.」又在大衛懺悔詩有這樣的話說:「神阿,求你為我造清潔的心,使我裡面得重新有正直的靈.」這都是長成的生命,內心的追求.
\newline
\newline
保羅說:「你要逃避少年的私慾,同那清心禱告主的人追求公義,信德,仁愛,和平.」約瑟逃避主母的試誘,寧可丟失自己的外衣.但以理不以王膳玷污自己,亦不委曲求全.西諺云:「如不喝酒,就勿繫馬於酒店門前.」都是叫我們避之若浼.經上說:主來的日子,好像挪亞和羅得的時代.那時,人照常喫喝嫁娶,耕種買賣,不知不覺人子就來了.照人看來,這些豈不都是正常的事.但神看內心,祂洞悉人心裡所存的惡念.希西家王將內庫展示於巴比倫使者,大衛數點他的百姓,都如出一轍.唯有神能鑑察人心.所以我們行事為人,有聖靈的恩膏凡事教導我們.使我們能認識好歹,分別聖俗,過成聖的生活.
\newline
\newline
二,有福的確據
\newline
\newline
彼得前書說:「如今雖不得看見,卻因信祂就有說不出來,滿有榮光的大喜樂.」(彼前一8):「信就是所望之事的實底,是未見之事的確據.」(來十一1)清心的人能見到上帝.因信者的心中滿有榮光的大喜樂.賈玉銘老牧師常對人說:「奇妙,奇妙,真是莫明其妙.」對他來說,神是可以摸著那樣地接近.
\newline
\newline
我們怎樣見神,不是用肉眼看祂,乃是透過信心的眼睛才能見到.我們從幾方面,都可以看見神.
\newline
\newline
1.聽覺　童子撒母耳在殿中,聽見神呼叫他的名字,初以為是祭司以喚他.後來明白就說:「請說,僕人敬聽.」摩西在何烈山放羊,也曾聽見神對他說話:「不要近前來,當將你腳上的鞋脫下來,因為你所站地是聖地.」
\newline
\newline
2.觸覺　有時人被神摸著,但有時人又可摸到神.以利亞被耶洗別追殺,逃到南地在羅騰樹下求死.神的使者兩次拍醒他:「起來喫吧,你當跑的路尚遠.」被希律捉拿下在監牢的彼得,睡夢中有天使拍他的肋膀,帶他出監.
\newline
\newline
3.味覺　神也是可用味覺嚐到的.(詩卅四8):「你要嘗嘗主的滋味,便知道祂是美善.」彼前二章也作這見證,舊約雅歌描述:「我的良人在男子中,如同蘋果樹林中.我歡歡喜喜坐在祂的蔭下,嘗祂果子的滋味覺得甘甜.」
\newline
\newline
4.嗅覺　林後二章寫出:藉我們在各處顯揚那因認識基督而有的香氣.雅歌稱頌良人；你的膏油馨香,你的名如同倒出來的香膏.
\newline
\newline
5.視覺　約伯說:我從前風聞有你,現在親眼看見你.保羅為以弗所教會禱告:求神將那賜人智慧和啟示的靈,賞給他們,使他們真知道祂.(弗一17)
\newline
\newline
當我們生命越發長進的時候,屬靈的感應,應當更操練得敏銳.從各方面接觸神,摸著神與祂同行.愛因斯坦有一次在一火車站餐廳預備喫晚餐.黑人侍者將餐單放在他面前時,才發覺忘掉了眼鏡,自然看不出所以來.侍者嘆說:「原來你與我同是目不識丁的.」弄得愛因斯坦啼笑皆非.少了一副眼鏡,就使鼎鼎大名的科學家變成與文盲無異.我們是否具備了屬靈的眼鏡使我們能看見祂?「清心的人有福了,你們必得見神.」應抓住這應許,常親近主,在苦難中經歷主的同在.
\newline
\newline
三,見祂的真體
\newline
\newline
參(約壹三2,彼前一8,林前十三12),同樣講到將必看見祂的真體.約伯說:「我知道我的救贖主活著,末了必站立在地上.我這皮肉滅絕之後,我必在肉體之外得見神.」(伯十九25)想到這裡,我們應當如何快樂.有一天我們要與主面對面,看見祂的真體.黑夜已深白晝將近,晨星出現,主必快來；我們與祂在空中相會,同在至永遠.
\newline
\newline
丹麥有一座耶穌的雕像,從別的角度看不出其特別之處.唯有跪下抬頭仰望,則見其慈容無不受感動.巴西里約熱內盧,夜景可算世界最美的,有一處名叫耶穌山,有耶穌像張開雙手在暗處發光.使人想到主再來時的景象.彼得後書一章,教我們在先知的預言上留心.二次大戰時,珍珠港突襲掀起了序幕,日本席捲了東南亞.美軍統帥麥克阿瑟將軍撤離菲島時曾許下諾言:「我必再來.」之豪語,後來果然兌現.一國之名將,許下諾言尚且如此.我們怎可對聖經上「主再來」的應許失去信心嗎?
\newline
\newline
方言現世的局面,尼布甲尼撒所夢之大像,有半鐵半坭的十趾；民主與極權的兩種制度豈不正並存於世.(亞十四4)預言主來時欖橄山要裂兩成半,今這預言又告實現.以色列復國,六日戰爭,都證明不可能的事成為可能.正是無花果樹發芽長葉的時候.夏天近了,人子已在門口了,祂必快來.
\newline
\newline


\newpage

\section{第48屆~港九培靈會~史祈生~第七講}
\label{sec:1099}
\textbf{生命的操練(一)}
\newline
\newline
連結: \href{https://www.hkbibleconference.org/session-message/view/1099}{\texttt{https://www.hkbibleconference.org/session-message/view/1099}}
\newline
\newline
\hyperref[sec:1097]{< < < PREV SERMON < < <}
~
\hyperref[sec:toc]{[返目錄]}
~
\hyperref[sec:1100]{> > > NEXT SERMON > > >}
\newline
\newline
馬太福音 5:9-5:9
\newline
\begin{longtable}{cl}
\hline
\hline
章節 & 經文 (和合本修訂版)\\
\hline
5:9 & \begin{tabularx}{0.7\textwidth}{X} 締造和平的人有福了！因為他們必稱為神的兒子。 \end{tabularx} \\ \\
[1ex]
\hline
\hline
\end{longtable}
經文:使人和睦的人有福了,因為他們必稱為神的兒子(太五9)
\newline
\newline
屬靈生命的歷程,分為三個階段:1.生命的印證.2.生命的成長.3.生命的操練.得著了生命；繼續長大,還要不斷向前邁進.既如是,則不能缺少操練,所謂「操練身體,益處還少,惟獨敬虔,凡事都有益處.」(提前四8)今講操練的第一層─使人和睦,便為有福.
\newline
\newline
一,和睦的基礎
\newline
\newline
「和睦」一辭,原文有和好,和諧,平安,和平等意思；新舊庫譯本作「締造和平」.主說:「他們必稱為神的兒子.」世人講武力強權,愛好和平反涉空論.野獸尚少相食,人卻自相殘殺.人於孩童之時,已喜以槍炮為現具,可見人性本善之非.據說滿清皇帝頭戴之紅頂,是用人血染成.中日交戰之時,鈴木大將也曾作和平之囈語,說必須征服莫斯科,打到華盛頓,就能達到和平,簡直是痴人說夢.國際間也不乏唱此高調,締結了數不清的盟約.但今天簽訂,難保他日不撕毀.聯合國之建立,安全理事會之籌設,初為維護世界之和平.當聯合國大廈奠基之時,將一本聖經放置其下；意謂以聖經之原則為其基礎.至今視之,無疑已將聖經埋葬.其後又將壁上石刻之禮運大同篇拆除,遷於中國城.孔孟之道,大同之理想,人以為不合時宜.反將安理會之議席,暫作小戰場.唇槍舌劍永無寧日.如果說能表現一致的話,那麼祗有飄揚於廣場上之百多面旗幟而耳.至於能和平相處,則惟有集郵者搜購各國的郵票,將它貼在同一冊上,就可以相安無事耳.
\newline
\newline
聯合國廣場上豎立了一塊石,上面刻著以賽亞書二章4節:「他們要將刀打成犁頭,把槍打成鐮刀,這國不舉刀攻擊那國,他們也不再學習戰事.」也不過是具文而已.真正的和平需要有基礎.回溯人類的第二代,已發生骨肉相殘；原因是遠離神,與生命之源頭隔斷.故基礎在基督身上,在於再得的生命.耶路撒冷之聖殿有牆,將猶太人與外邦人分開.以弗所書二章:「因他使我們和睦,將兩下合為一,拆毀了中間隔斷的牆.」主耶穌說:「我來並不是叫地上太平,乃是叫地上動刀兵.......人的仇敵,就是自己家裡的人」(太十34-36)這裡指出必經過相爭,才能達致真正和平.因為信與不信不能共負一軛,真和平需有生命的基礎.一九一四年第一次世界大戰,英德雙方互相對壘.聖誕前夕,卻唱起平安夜來；不期使雙方軍士,在陣前握手高歌,共慶主臨.可見耶穌基督才是和平之基.
\newline
\newline
二,和諧的生活
\newline
\newline
我們得稱為神的兒子,有神的生命,神的性情是和平的,試看宇宙的運轉何等和諧；四時運行,無言無語.地球每小時約移動一千哩,其上的人渾無所覺.時序更移,滋生萬物.何等寧靜安詳.祂又叫我們得安息,並且進一步享安息.世人有如鴨子泛流,外表安靜,裡面卻忙個不了.「我們的心若不責備我們,就可以向神坦然無懼了.」(約壹三21)良心好像燈光,照亮我們離開罪惡,享受安息.詩一三一篇:「耶和華阿,我的心不狂傲,我的眼不高大.......我的心平穩安靜,好像斷過奶的孩子,在他母親的懷中.」這是一幅何等美麗寧謐的畫面.若我們在生命操練進步,也可享受這如詩如畫的人生.
\newline
\newline
進一步要推廣這愛,締造和睦.俗語說:「家和萬事興.」夫婦之間尤需和睦相處.若雙方意見相左,勢成水火,則無論錢財,地位,學識,都變成毫無意義了.有對夫婦也感到這方面的重要,但苦不能得.終商得一計,若丈夫在外遭遇不快,則戴帽欹側；反之,妻在家如感心情抑鬱,則倒繫圍裙.彼此相約,互相遷就,果能相安一時.一次丈夫回家,啟門後彼此瞠目相對.原來夫帽欹側,妻之圍裙倒繫,仍是勢成冰炭.信徒家中,常見掛上一銅牌,上面寫著:「基督是我家之主.」這才是至美至善的道路.如果生活上尊主為大,讓基督在家中作主,自然產生平安.舊約歷史,大衛及所羅門作王時,國家興盛富強.後來卻由盛而衰,原因分裂為南國猶大,及北國以色列.從此國勢漸衰.教會亦然,若同心合意,自然蒙神福佑；如果分爭結黨,則災禍立致.有一道歌詞說:世間好似兩軍對敵,即是善與惡.注意要我們與惡相爭,不是在教會中肢體彼此相爭.貽外人笑柄,羞辱主名.從前我曾在鄉下佈道,有位老者初因好奇上前來聽,後來卻捏著子輕蔑地離開.我覺得奇怪而問他.原來他的鄰舍住著一位稱為基督徒的,但生活為人卻不堪承教,這就成了別人的絆腳石.
\newline
\newline
主說:若有弟兄向你懷怨,應先去與他和好,然後在壇上獻祭.人或想,既然是他人恨我,何以卻要我向對方和解賠不是.孰知若要效主榜樣,就當照著行.因不問是否合理,乃要救這軟弱弟兄.和平的平字是把兩顆心,放在同一水平之上.不可驕傲,自以為是,輕看別人.看見弟兄眼中有剌,卻看不見自己有梁木；原來對方的剌是由本身反映過去的.「你的仇敵若餓了,就給他喫.若渴了就給他喝.因為你這樣行,就是把炭火堆在他的頭上.」(羅十二20)就是說要愛你的仇敵.一同作工的也要彼此相愛,三股合成的繩子不能折斷.經上又說:「你們五個人要追趕一百人,一百人要追趕一萬人.」(利廿六8)這是同心合力所產生的屬靈功效.
\newline
\newline
三,和好的使命
\newline
\newline
諾貝爾發明炸藥之後,因良心內疚,設立和平獎金,每年發頒給對和平有貢獻的人.神的兒女,有責任傳福音使人與神和好,因祂藉著基督已將使人與神和好的職份託付我們了.教會要行動,同心合意傳神的福音,不要祗在紙上談兵,因傳福音是整個教會的行動.有一個笑話說:有教會開佈道大會,第一晚講員登上講台問大家,知否我今晚講甚麼?眾人說:「不知.」牧師說既不知我要講甚麼,那我就不必講了,就此離開.第二晚開會時,眾人鑑於昨夜所發生的事,就準備著應付.牧師果然照問如故,眾人答曰:「知道.」既然知道,又不需講了.第三晚,會眾怕牧師故技重施,共謀對策；果然又發問,這回一半答「知」,一半答「不知」.牧師說:那就好了,可叫知道的告訴不知道的.從這話也講出了一番道理,我們知道的,去告訴不知的,福音豈不很易遍傳嗎?每天每時都有不少人離開世界,送殯行列鬧哄哄的,但有些人關心又一個靈魂淪亡呢?
\newline
\newline
英名歷史家湯恩伯,八十生辰時答紐約時報記者訪問,曾表示廿一世紀麼中國人的時代.將佔世界人口五分之三,中國人向同胞傳福音最為方便,言語的隔膜容易消除.中國人做生意無孔不入,我曾到過南美洲某一個小島,該處人跡少到,但卻見店舖內有中國貨的產品,福音未傳到,糖果罐頭卻先運到了.我們為主得人應有比商人牟利更熱切的心情.名佈道叨雷曾說,每年一個人歸主,循此以往,卅三年整個世界都傳遍福音,都歸信主了.一年三百六十五天,領一人歸主也不能作到,可算閒懶不結果子.
\newline
\newline
南北戰爭之時,壯丁入伍,甚至未經訓練的也要即時驅使出戰.一次軍隊吹號撤退,內一新兵不諳號令懵然不知.既退,眾以為必死.誰知不久睹其無恙歸來,並兩手抓到兩個敵人.各人奇怪問之.他鬥然回答,樹林之中有許多,祗要去抓就必得著.
\newline
\newline
主將傳福音的使命,不交予天使,卻交給我們.這是何等奇妙的恩典.若不傳福音就有禍了.應體貼主的心,多作主工,候主再臨.
\newline
\newline


\newpage

\section{第48屆~港九培靈會~史祈生~第八講}
\label{sec:1100}
\textbf{生命的操練(二,三)}
\newline
\newline
連結: \href{https://www.hkbibleconference.org/session-message/view/1100}{\texttt{https://www.hkbibleconference.org/session-message/view/1100}}
\newline
\newline
\hyperref[sec:1099]{< < < PREV SERMON < < <}
~
\hyperref[sec:toc]{[返目錄]}
~
\hyperref[sec:1101]{> > > NEXT SERMON > > >}
\newline
\newline
馬太福音 5:10-5:12
\newline
\begin{longtable}{cl}
\hline
\hline
章節 & 經文 (和合本修訂版)\\
\hline
5:10 & \begin{tabularx}{0.7\textwidth}{X} 為義受迫害的人有福了！因為天國是他們的。 \end{tabularx} \\ \\ \relax
5:11 & \begin{tabularx}{0.7\textwidth}{X} 「人若因我辱罵你們，迫害你們，捏造各樣壞話毀謗你們，你們就有福了！ \end{tabularx} \\ \\ \relax
5:12 & \begin{tabularx}{0.7\textwidth}{X} 要歡喜快樂，因為你們在天上的賞賜是很多的。在你們以前的先知，人也是這樣迫害他們。」 \end{tabularx} \\ \\
[1ex]
\hline
\hline
\end{longtable}
經文:為義受逼迫的人有福了,因為天國是他們的.人若因我辱罵你們,逼迫你們,捏造各樣壞話毀謗你們,你們就有福了.應當歡喜快樂,因為你們在天上的賞賜是大的.在你們以前的先知,人也是這樣逼迫他們.(太五10-12)
\newline
\newline
時間過得很快,今天是最後一天了.起初聽人傳說香港很可怕,但我現在倒覺得它很可愛.願你們對紐約也是這樣.
\newline
\newline
今天繼續講屬靈生命的歷程,第三段──生命的操練,最後兩個福份.
\newline
\newline
操練的功課,要常常學習,一直到見主的面.主滿有溫柔,祂告訴我們要學祂的樣式.但溫柔並非千依百順.請看主在傳道之時,進入聖殿,用繩子作的鞭子,趕出販賣牛羊鴿子的；又推翻兌換銀錢的桌子.顯明主的態度,是有原則和立場的.主愛罪人卻恨惡罪惡；愛弟兄而拒絕魔鬼.主說:「我來不是叫地上太平,乃是叫地上動兵刀.因為我來,是要叫人與父親生疏,女兒與母親生疏,媳婦與婆婆生疏.」注意都是從下而上的.豈不是與中國素常所注重的孝道相互衝突?絕不,保羅教訓我們要在主裡聽從父母.(弗六2)就是有原則不盲從之意.
\newline
\newline
基督徒在世,為主的緣故,過公義的生活,不免遭受難處逼迫.但這是生命的操練上必需的.摩西帶領百姓出埃及,起行時數近二百萬.但得進入迦南的,當代的人僅剩迦勒和約書亞二人.他們敬畏神,堅持所信的.當我們在天路上追求,到達屬靈的高峰時,保羅說:「你們中間或有人似乎是趕不上了.」當我們平安無事時,勇猛好像雄獅；事情臨到,卻膽怯如.緬想許多基督徒,今天還處身在患難桎梏之中,我們應如何禱告記念他們呢!以下要將為義受逼迫,和為主受苦二題,合在一次講完.
\newline
\newline
一,為義受逼迫
\newline
\newline
主呼召門徒來跟從祂,並不用花言巧語.從來沒有應許他們,道路通達不遇艱難.反倒說:「狐狸有洞,天空飛鳥有窩,只是人子沒有枕頭的地方.」祂也沒有把我們帶到曠野或另一個星球去,使我們過安逸無憂的生活.反倒要背著十字架跟隨祂；捨棄生命才能得著生命.主明言在世上有苦難.為屬祂的人禱告:「我不求你叫他們離開世界,只求你保守他們脫離那惡者.」祂又對彼得說:「撒旦想要得著你們,好篩你們像篩麥子一樣,但我已經為你祈求,叫你不至於失了信心,你回頭以後,要堅固你的弟兄.」這都給我們見到屬靈生命的操練；總要登峰造極,苦難是免不了的,主耶穌說:「被召的人多,選上的人少.」巴不得主來的時候,我們屬於被選的少數.
\newline
\newline
論到所得的福份,經上說:「天國是他們的.」應當注意和頭一個福是相同的.是為得勝的人所預備的.就是為義受逼迫的人,能夠承受的福祉.但也當清楚為義受逼迫的真意.在保羅的經歷中,他寫出:我身上常帶著耶穌的死；他冒死是屢次有的；我們這活著的人,是常為耶穌被交於死地.(參林後四10,十一23,四11)我曾在國內當軍中牧師,也常關心那些軍人內心的情況.他們都預備好了,面對死亡,一無恐懼.作主的精兵,就是如此.
\newline
\newline
但是有些人受苦,卻不是這樣.(彼前二20):「你們若因犯罪受責打能忍耐,有什麼可誇的呢?」(彼前四15)「你們中間卻不可有人,因為殺人,偷竊,作惡,好管閒事而受苦.」許多時候人會混淆,以為所遭遇的,是為義所帶來的難處.作神子民選擇義是應當的.主說:「先求神的國和祂的義,這一切東西都要加給你們了.」「我們既從仇敵手中被救出來,就可......用聖潔公義事奉祂.」所以天國的子民,不可與世同流合污,隨夥行惡.寧可喫虧,遭人歧視,蒙受損失,卻是討主喜悅的.若有人因好管閒事,或因種種惡行而受苦,那是罪有應得,理當如此,絕不會在神裡面得福.嘗見有人在祈禱會中,指摘別人,甚至控告自己的丈夫.及至讓丈夫知道,自然挨打了.有些人傳福音卻沒有智慧,冒失地在言語上得罪人,挨耳光也絕非奇事.又有些人在教會之中,自恃財勢,擅作威福.遇到有人不賣他的賬,就滿腹牢騷,以為是為義受苦.其實,凡屬這些都是咎由自取.絕非神的旨意,與賜福的應許無份無關.我曾在某雜誌見到一段消息:美國有某大學生團契佈道團,在一地舉行公開佈道會.在美國人民是享有宗教自由的,祗要循正當手續申請便行了.但那班學生,自命為神的使者,罔顧政府法紀.受到干涉,捉將官去.還藉此大肆宣傳,以羞辱為榮耀.有些富足的人,在教會也玩弄權勢,稍不如意,又以為為義受苦.其實,正合於亞拿尼亞,撒非喇之流,愛體面,欺哄神.何義之有?
\newline
\newline
我們屬主的人,在今世過公義的生活是應當的.雖然暫受虧損,被人指摘,目為迂腐.但卻是蒙神所賜福的.美國有不少規模宏大,生意興盛的公司,他們都是守住公義待人的宗旨；反之就不能長久留存.(約壹一17)說:「這世界和其上的情慾,都要過去,惟獨遵行神旨意的,是永遠長存.」有個初嫁的媳婦,歸寧之日,對母親流淚訴苦.不是因丈夫不良,乃因家姑難於服侍.經過母親多方勸解還是不能釋然.後來母問女兒:「婆婆幾歲?」答說:「八十八歲.」母親說:「既然八十八歲了,還有多少日子受苦呢?」這才使女兒轉悲為喜.安徒生童話也寫過一隻醜小鴨,在眾子中常被欺侮.直至一天,一隊天鵝飛過,醜鴨子就逐隊而飛,原來這被藐視的是隻天鵝.我們要忍受這至暫至輕的苦楚,想到將來的榮耀,就不足介意了.
\newline
\newline
二,為主受苦
\newline
\newline
在印度許多教徒,躺在地上讓載神像的車子軋過去.他們的信仰雖錯誤,但忠誠卻可嘉.宋尚節博士講道中曾用兩圓球做比喻.一球是用坭土作的,一摔便碎；另一球是真的皮球,用力摔下去,高高地彈起來.兩球表面一樣,表明真假兩種信徒,平常不易分辨；直到苦難試煉來臨,真假就判然若揭.美國總統競選當中那些候選人往各處,巡迴演講,但常受到人反對,甚至跑到面前唾罵他們.好像華萊士曾競選失敗,被行刺成為殘廢.當其東山復起再度競選時,人對之譏笑辱罵,造成醜陋的人像而窘辱之.但那些競選人都能為黨國的緣故,夷然處之.基督徒對此作何感想.豈不也應為主的名,忍受各樣的辱罵,逼迫,凌辱.因為這是魔鬼常用的伎倆.中國的諺語說:「宰相腹中可撐船.」基督徒也應有容人之量.
\newline
\newline
經文說:「因為在你們以前的先知,人也是這樣逼迫他們.」耶穌說:「從義人亞伯的血起,直至你們在殿和壇中間所殺的巴拉加的兒子撒迦利亞的血為止.」(太廿三35)有的殺害,釘十字架在會堂裡鞭打,從這城追逼到那城.自古以來為主殉道的先知義人.不勝枚舉,主後三百年間,羅馬帝國屢次對教會興起大逼迫.萬千聖徒捐軀流血,寫成光榮的史蹟.門徒中司提反最先為主殉道,引致掃羅歸正.其後眾使徒均為主捨命,作出美好的見證.據教會遺傳:彼得安得烈被釘十字架,雅各馬太達太被殺,約翰充軍,腓力路加被絞死,巴拿巴馬提亞被石頭打死,巴多羅買被剝皮,多馬用槍貫胸.......保羅與提摩太都在羅馬殉道.若沒有初代聖徒慷慨赴義,何能有今日教會.西班牙考古學家普頓參觀羅馬殉道聖徒的墓園,致信與友人書其感想:因信而死的人數何多,其碑文記事,斑斑可考,使人緬想前賢,油然有踵有步武之感.希伯來書十一章提到許多因信受苦的見證人,忍受各樣苦難,不肯苟且得釋放,為要得著更美的復活.保羅又說:「末了所毀滅的仇敵就是死.」(林前十五26)能不伯死就是得勝者,能夠與主同行.耶穌說:「殺身體不能殺靈魂的,不要怕他們,惟有能把身體和靈魂都滅在地獄裡的,正要怕他.......就是你們的頭髮都被數過.」
\newline
\newline
末世主來更近,應當預備迎見主.以受苦的心志為兵器,十架換取榮耀的冠冕.在主前配被稱為又良善,又忠的僕人.主說:「樂意把國賜給你們.」
\newline
\newline


\newpage

\section{第48屆~港九培靈會~孫德生~第一講}
\label{sec:1101}
\textbf{聖靈的火}
\newline
\newline
連結: \href{https://www.hkbibleconference.org/session-message/view/1101}{\texttt{https://www.hkbibleconference.org/session-message/view/1101}}
\newline
\newline
\hyperref[sec:1100]{< < < PREV SERMON < < <}
~
\hyperref[sec:toc]{[返目錄]}
~
\hyperref[sec:1103]{> > > NEXT SERMON > > >}
\newline
\newline
列王記上 18:12-18:40
\newline
\begin{longtable}{cl}
\hline
\hline
章節 & 經文 (和合本修訂版)\\
\hline
18:12 & \begin{tabularx}{0.7\textwidth}{X} 恐怕我一離開你，耶和華的靈就把你提到我所不知道的地方去。這樣，我去告訴亞哈，他若找不到你，就必殺我。僕人是自幼敬畏耶和華的。 \end{tabularx} \\ \\ \relax
18:13 & \begin{tabularx}{0.7\textwidth}{X} 耶洗別殺耶和華先知的時候，我把耶和華的一百個先知藏了，每五十人藏在一個洞裡，拿餅和水供養他們，難道沒有人把我做的這事告訴我主嗎？ \end{tabularx} \\ \\ \relax
18:14 & \begin{tabularx}{0.7\textwidth}{X} 現在你說：『你去告訴你主人說，看哪，以利亞在這裡』，他一定會殺我。」 \end{tabularx} \\ \\ \relax
18:15 & \begin{tabularx}{0.7\textwidth}{X} 以利亞說：「我指著所事奉永生的萬軍之耶和華起誓，我今日要讓亞哈見到我。」 \end{tabularx} \\ \\ \relax
18:16 & \begin{tabularx}{0.7\textwidth}{X} 於是俄巴底去迎見亞哈，告訴他這事。亞哈就去見以利亞。 \end{tabularx} \\ \\ \relax
18:17 & \begin{tabularx}{0.7\textwidth}{X} 亞哈見了以利亞，就說：「真的是你嗎？你這使以色列遭殃的人！」 \end{tabularx} \\ \\ \relax
18:18 & \begin{tabularx}{0.7\textwidth}{X} 以利亞說：「使以色列遭殃的不是我，而是你和你的父家，因為你們離棄耶和華的誡命，去隨從巴力。 \end{tabularx} \\ \\ \relax
18:19 & \begin{tabularx}{0.7\textwidth}{X} 現在你要派人去召集以色列眾人，以及耶洗別所供養的四百五十個巴力的先知和四百個亞舍拉的先知，叫他們都上迦密山到我這裡來。」 \end{tabularx} \\ \\ \relax
18:20 & \begin{tabularx}{0.7\textwidth}{X} 亞哈就派人到以色列眾人那裡，召集先知上迦密山。 \end{tabularx} \\ \\ \relax
18:21 & \begin{tabularx}{0.7\textwidth}{X} 以利亞近前來對眾百姓說：「你們心持二意要到幾時呢？如果耶和華是神，就當順從耶和華；如果是巴力，就當順從巴力。」百姓一言不答。 \end{tabularx} \\ \\ \relax
18:22 & \begin{tabularx}{0.7\textwidth}{X} 以利亞對百姓說：「作耶和華先知的只剩下我一個；巴力的先知卻有四百五十人。 \end{tabularx} \\ \\ \relax
18:23 & \begin{tabularx}{0.7\textwidth}{X} 請給我們兩頭牛犢，巴力的先知可以為自己挑選一頭牛犢，切成小塊，放在柴上，不要點火；我也預備一頭牛犢放在柴上，也不點火。 \end{tabularx} \\ \\ \relax
18:24 & \begin{tabularx}{0.7\textwidth}{X} 你們求告你們神明的名，我也求告耶和華的名。那應允禱告降火的就是神。」眾百姓回答說：「好主意。」 \end{tabularx} \\ \\ \relax
18:25 & \begin{tabularx}{0.7\textwidth}{X} 以利亞對巴力的先知說：「因為你們人多，先挑選一頭牛犢，預備好了，求告你們神明的名，卻不要點火。」 \end{tabularx} \\ \\ \relax
18:26 & \begin{tabularx}{0.7\textwidth}{X} 他們把所給他們的牛犢預備好了，從早晨到中午，求告巴力的名說：「巴力啊，求你應允我們！」卻沒有聲音，也沒有回應。他們就在所築的壇四圍蹦跳。 \end{tabularx} \\ \\ \relax
18:27 & \begin{tabularx}{0.7\textwidth}{X} 到了正午，以利亞嘲笑他們，說：「大聲求告吧！因為它是神明，它或許在默想，或許正忙著，或許在路上，或許在睡覺，它該醒過來了。」 \end{tabularx} \\ \\ \relax
18:28 & \begin{tabularx}{0.7\textwidth}{X} 他們大聲求告，按著他們的儀式，用刀槍刺割自己，直到渾身流血。 \end{tabularx} \\ \\ \relax
18:29 & \begin{tabularx}{0.7\textwidth}{X} 中午過去了，他們狂呼亂叫，直到獻晚祭的時候，卻沒有聲音，沒有回應的，也沒有理睬的。 \end{tabularx} \\ \\ \relax
18:30 & \begin{tabularx}{0.7\textwidth}{X} 以利亞對眾百姓說：「你們到我這裡來。」眾百姓就到他那裡，他把那已經毀壞了的耶和華的壇修好。 \end{tabularx} \\ \\ \relax
18:31 & \begin{tabularx}{0.7\textwidth}{X} 以利亞按照雅各子孫支派的數目，取了十二塊石頭；耶和華的話曾臨到雅各，說：「你的名要叫以色列。」 \end{tabularx} \\ \\ \relax
18:32 & \begin{tabularx}{0.7\textwidth}{X} 以利亞用這些石頭為耶和華的名築一座壇，在壇的四圍挖溝，可容納二細亞穀種。 \end{tabularx} \\ \\ \relax
18:33 & \begin{tabularx}{0.7\textwidth}{X} 他又在壇上擺好了柴，把牛犢切成小塊放在柴上，說：「你們用四個桶盛滿水，倒在燔祭和柴上。」 \end{tabularx} \\ \\ \relax
18:34 & \begin{tabularx}{0.7\textwidth}{X} 他又說：「倒第二次。」他們就倒第二次。他又說：「倒第三次。」他們就倒第三次。 \end{tabularx} \\ \\ \relax
18:35 & \begin{tabularx}{0.7\textwidth}{X} 水流到壇的四圍，溝裡也滿了水。 \end{tabularx} \\ \\ \relax
18:36 & \begin{tabularx}{0.7\textwidth}{X} 到了獻晚祭的時候，先知以利亞近前來，說：「耶和華－亞伯拉罕、以撒、以色列的神啊，求你今日使人知道你是以色列的神，我是你的僕人，我遵照你的話做這一切事。 \end{tabularx} \\ \\ \relax
18:37 & \begin{tabularx}{0.7\textwidth}{X} 求你應允我，耶和華啊，應允我，使這百姓知道你－耶和華是神，是你叫他們回心轉意的。」 \end{tabularx} \\ \\ \relax
18:38 & \begin{tabularx}{0.7\textwidth}{X} 於是，耶和華降下火來，燒盡燔祭、木柴、石頭、塵土，又燒乾了溝裡的水。 \end{tabularx} \\ \\ \relax
18:39 & \begin{tabularx}{0.7\textwidth}{X} 眾百姓看見了，就臉伏於地，說：「耶和華是神！耶和華是神！」 \end{tabularx} \\ \\ \relax
18:40 & \begin{tabularx}{0.7\textwidth}{X} 以利亞對他們說：「拿住巴力的先知，不讓任何人逃走！」眾人就拿住他們。以利亞帶他們到基順河邊，在那裡殺了他們。 \end{tabularx} \\ \\
[1ex]
\hline
\hline
\end{longtable}
經文:為義受逼迫的人有福了,因為天國是他們的.人若因我辱罵你們,逼迫你們,捏造各樣壞話毀謗你們,你們就有福了.應當歡喜快樂,因為你們在天上的賞賜是大的.在你們以前的先知,人也是這樣逼迫他們.(太五10-12)
\newline
\newline
時間過得很快,今天是最後一天了.起初聽人傳說香港很可怕,但我現在倒覺得它很可愛.願你們對紐約也是這樣.
\newline
\newline
今天繼續講屬靈生命的歷程,第三段──生命的操練,最後兩個福份.
\newline
\newline
操練的功課,要常常學習,一直到見主的面.主滿有溫柔,祂告訴我們要學祂的樣式.但溫柔並非千依百順.請看主在傳道之時,進入聖殿,用繩子作的鞭子,趕出販賣牛羊鴿子的；又推翻兌換銀錢的桌子.顯明主的態度,是有原則和立場的.主愛罪人卻恨惡罪惡；愛弟兄而拒絕魔鬼.主說:「我來不是叫地上太平,乃是叫地上動兵刀.因為我來,是要叫人與父親生疏,女兒與母親生疏,媳婦與婆婆生疏.」注意都是從下而上的.豈不是與中國素常所注重的孝道相互衝突?絕不,保羅教訓我們要在主裡聽從父母.(弗六2)就是有原則不盲從之意.
\newline
\newline
基督徒在世,為主的緣故,過公義的生活,不免遭受難處逼迫.但這是生命的操練上必需的.摩西帶領百姓出埃及,起行時數近二百萬.但得進入迦南的,當代的人僅剩迦勒和約書亞二人.他們敬畏神,堅持所信的.當我們在天路上追求,到達屬靈的高峰時,保羅說:「你們中間或有人似乎是趕不上了.」當我們平安無事時,勇猛好像雄獅；事情臨到,卻膽怯如.緬想許多基督徒,今天還處身在患難桎梏之中,我們應如何禱告記念他們呢!以下要將為義受逼迫,和為主受苦二題,合在一次講完.
\newline
\newline
一,為義受逼迫
\newline
\newline
主呼召門徒來跟從祂,並不用花言巧語.從來沒有應許他們,道路通達不遇艱難.反倒說:「狐狸有洞,天空飛鳥有窩,只是人子沒有枕頭的地方.」祂也沒有把我們帶到曠野或另一個星球去,使我們過安逸無憂的生活.反倒要背著十字架跟隨祂；捨棄生命才能得著生命.主明言在世上有苦難.為屬祂的人禱告:「我不求你叫他們離開世界,只求你保守他們脫離那惡者.」祂又對彼得說:「撒旦想要得著你們,好篩你們像篩麥子一樣,但我已經為你祈求,叫你不至於失了信心,你回頭以後,要堅固你的弟兄.」這都給我們見到屬靈生命的操練；總要登峰造極,苦難是免不了的,主耶穌說:「被召的人多,選上的人少.」巴不得主來的時候,我們屬於被選的少數.
\newline
\newline
論到所得的福份,經上說:「天國是他們的.」應當注意和頭一個福是相同的.是為得勝的人所預備的.就是為義受逼迫的人,能夠承受的福祉.但也當清楚為義受逼迫的真意.在保羅的經歷中,他寫出:我身上常帶著耶穌的死；他冒死是屢次有的；我們這活著的人,是常為耶穌被交於死地.(參林後四10,十一23,四11)我曾在國內當軍中牧師,也常關心那些軍人內心的情況.他們都預備好了,面對死亡,一無恐懼.作主的精兵,就是如此.
\newline
\newline
但是有些人受苦,卻不是這樣.(彼前二20):「你們若因犯罪受責打能忍耐,有什麼可誇的呢?」(彼前四15)「你們中間卻不可有人,因為殺人,偷竊,作惡,好管閒事而受苦.」許多時候人會混淆,以為所遭遇的,是為義所帶來的難處.作神子民選擇義是應當的.主說:「先求神的國和祂的義,這一切東西都要加給你們了.」「我們既從仇敵手中被救出來,就可......用聖潔公義事奉祂.」所以天國的子民,不可與世同流合污,隨夥行惡.寧可喫虧,遭人歧視,蒙受損失,卻是討主喜悅的.若有人因好管閒事,或因種種惡行而受苦,那是罪有應得,理當如此,絕不會在神裡面得福.嘗見有人在祈禱會中,指摘別人,甚至控告自己的丈夫.及至讓丈夫知道,自然挨打了.有些人傳福音卻沒有智慧,冒失地在言語上得罪人,挨耳光也絕非奇事.又有些人在教會之中,自恃財勢,擅作威福.遇到有人不賣他的賬,就滿腹牢騷,以為是為義受苦.其實,凡屬這些都是咎由自取.絕非神的旨意,與賜福的應許無份無關.我曾在某雜誌見到一段消息:美國有某大學生團契佈道團,在一地舉行公開佈道會.在美國人民是享有宗教自由的,祗要循正當手續申請便行了.但那班學生,自命為神的使者,罔顧政府法紀.受到干涉,捉將官去.還藉此大肆宣傳,以羞辱為榮耀.有些富足的人,在教會也玩弄權勢,稍不如意,又以為為義受苦.其實,正合於亞拿尼亞,撒非喇之流,愛體面,欺哄神.何義之有?
\newline
\newline
我們屬主的人,在今世過公義的生活是應當的.雖然暫受虧損,被人指摘,目為迂腐.但卻是蒙神所賜福的.美國有不少規模宏大,生意興盛的公司,他們都是守住公義待人的宗旨；反之就不能長久留存.(約壹一17)說:「這世界和其上的情慾,都要過去,惟獨遵行神旨意的,是永遠長存.」有個初嫁的媳婦,歸寧之日,對母親流淚訴苦.不是因丈夫不良,乃因家姑難於服侍.經過母親多方勸解還是不能釋然.後來母問女兒:「婆婆幾歲?」答說:「八十八歲.」母親說:「既然八十八歲了,還有多少日子受苦呢?」這才使女兒轉悲為喜.安徒生童話也寫過一隻醜小鴨,在眾子中常被欺侮.直至一天,一隊天鵝飛過,醜鴨子就逐隊而飛,原來這被藐視的是隻天鵝.我們要忍受這至暫至輕的苦楚,想到將來的榮耀,就不足介意了.
\newline
\newline
二,為主受苦
\newline
\newline
在印度許多教徒,躺在地上讓載神像的車子軋過去.他們的信仰雖錯誤,但忠誠卻可嘉.宋尚節博士講道中曾用兩圓球做比喻.一球是用坭土作的,一摔便碎；另一球是真的皮球,用力摔下去,高高地彈起來.兩球表面一樣,表明真假兩種信徒,平常不易分辨；直到苦難試煉來臨,真假就判然若揭.美國總統競選當中那些候選人往各處,巡迴演講,但常受到人反對,甚至跑到面前唾罵他們.好像華萊士曾競選失敗,被行刺成為殘廢.當其東山復起再度競選時,人對之譏笑辱罵,造成醜陋的人像而窘辱之.但那些競選人都能為黨國的緣故,夷然處之.基督徒對此作何感想.豈不也應為主的名,忍受各樣的辱罵,逼迫,凌辱.因為這是魔鬼常用的伎倆.中國的諺語說:「宰相腹中可撐船.」基督徒也應有容人之量.
\newline
\newline
經文說:「因為在你們以前的先知,人也是這樣逼迫他們.」耶穌說:「從義人亞伯的血起,直至你們在殿和壇中間所殺的巴拉加的兒子撒迦利亞的血為止.」(太廿三35)有的殺害,釘十字架在會堂裡鞭打,從這城追逼到那城.自古以來為主殉道的先知義人.不勝枚舉,主後三百年間,羅馬帝國屢次對教會興起大逼迫.萬千聖徒捐軀流血,寫成光榮的史蹟.門徒中司提反最先為主殉道,引致掃羅歸正.其後眾使徒均為主捨命,作出美好的見證.據教會遺傳:彼得安得烈被釘十字架,雅各馬太達太被殺,約翰充軍,腓力路加被絞死,巴拿巴馬提亞被石頭打死,巴多羅買被剝皮,多馬用槍貫胸.......保羅與提摩太都在羅馬殉道.若沒有初代聖徒慷慨赴義,何能有今日教會.西班牙考古學家普頓參觀羅馬殉道聖徒的墓園,致信與友人書其感想:因信而死的人數何多,其碑文記事,斑斑可考,使人緬想前賢,油然有踵有步武之感.希伯來書十一章提到許多因信受苦的見證人,忍受各樣苦難,不肯苟且得釋放,為要得著更美的復活.保羅又說:「末了所毀滅的仇敵就是死.」(林前十五26)能不伯死就是得勝者,能夠與主同行.耶穌說:「殺身體不能殺靈魂的,不要怕他們,惟有能把身體和靈魂都滅在地獄裡的,正要怕他.......就是你們的頭髮都被數過.」
\newline
\newline
末世主來更近,應當預備迎見主.以受苦的心志為兵器,十架換取榮耀的冠冕.在主前配被稱為又良善,又忠的僕人.主說:「樂意把國賜給你們.」
\newline
\newline


\newpage

\section{第48屆~港九培靈會~孫德生~第二講}
\label{sec:1103}
\textbf{聖靈是誰}
\newline
\newline
連結: \href{https://www.hkbibleconference.org/session-message/view/1103}{\texttt{https://www.hkbibleconference.org/session-message/view/1103}}
\newline
\newline
\hyperref[sec:1101]{< < < PREV SERMON < < <}
~
\hyperref[sec:toc]{[返目錄]}
~
\hyperref[sec:1104]{> > > NEXT SERMON > > >}
\newline
\newline
約翰福音 14:15-14:18
\newline
\begin{longtable}{cl}
\hline
\hline
章節 & 經文 (和合本修訂版)\\
\hline
14:15 & \begin{tabularx}{0.7\textwidth}{X} 「你們若愛我，就會遵守我的命令。 \end{tabularx} \\ \\ \relax
14:16 & \begin{tabularx}{0.7\textwidth}{X} 我要求父，父就賜給你們另外一位保惠師，使他永遠與你們同在。 \end{tabularx} \\ \\ \relax
14:17 & \begin{tabularx}{0.7\textwidth}{X} 他就是真理的靈，是世人不能接受的。因為他們既看不見他，也不認識他；你們卻認識他，因他常與你們同在，也要在你們裡面。 \end{tabularx} \\ \\ \relax
14:18 & \begin{tabularx}{0.7\textwidth}{X} 我不會撇下你們為孤兒，我必到你們這裡來。 \end{tabularx} \\ \\
[1ex]
\hline
\hline
\end{longtable}
我十分高興見到在座有許多青年人.我的年紀雖大,但我自覺不老,我還能打網球,還能回味十八歲的感覺,也能明白青年人心裡的感受.各位年青人!你們所有的經驗我已試過了,我曾和你們一樣的淘氣.感謝主!祂遇見我,改變了我的生命.
\newline
\newline
我年輕時代,不多聽見人題到聖靈,我常到禮拜堂,但從未聽到題及聖靈.聖靈是三位一體之中最易被人遺忘的；但感謝神!如今趨勢不同了,常有人題及聖靈；可是許多人對於聖靈是誰,聖靈作何工作,仍然模糊不清；甚至基督徒們因對聖靈見解不同,而發生激烈的辯論.有時所辯論的是關於形容聖靈的詞句,及人和聖靈所發生的關係那種經歷而辯論.這樣使魔鬼最開心,牠最喜歡那些膽怯,不敢題聖靈和聖靈工作的人；牠驅使一些人,對聖靈有過份的反應.我曾經在中國看見類此事件發生.在一九四七年我有機會探訪中國,到過很多地方,也到新疆和西藏；看見聖靈所作的工作,和有些人對聖靈的誤解；致使教會起了紛爭,見證受了損壞.有的人聽到題說聖靈,就退避三舍了.然而我誠心相信聖靈是奇妙的一位,祂很願意在我們的生命中賜福給我們.
\newline
\newline
今晚我願意把五十年來從聖靈所學習的心得,與大家分享,我不想把任何的看法加以批評,也不將任何人所說的話加以反對；惟將聖經記載關於聖靈的事與大家分享.
\newline
\newline
聖靈是一種力量,有影響或有位格的一位呢?聖靈是否如電力一樣,你可以開掣也可以關掣呢?聖靈是否真確活著的一位?聖靈是可使我們利用以達目的,抑或祂願意使用你我來就祂旨意的呢?祂是很智慧值得我們愛祂的,或是毫無人格,毫無感情的等待我們去利用的呢?祂是否一位可使用人去作工的?這些都是重要的問題,你不能憑著某人的經驗來取得答案,唯一的答案乃記載在聖經.
\newline
\newline
聖靈是一位智慧的靈,值得我們愛上祂.很多年前我參加一個晚餐會,我準備了適合於該會的訊息；但臨時心裡很深地感動,不想講所預備的訊息,那除此我講什麼呢?進餐時我一直向主說:「主啊!你要我講些什麼呢?」飯後請我講道；我起身趨前站在那裡無道可講,我就打開聖經讀約翰福音十四章,希望主從那裡給我訊息.當我讀至第七節:「你們若認識我,也就認識我的父.」我恍然大悟!我確是認識主,祂是我的拯救者,我的主宰,是我人生的掌管者；我也認識我的父,祂是我天父.我禱告說:「我們在天上的父......」祂是慈愛的,是關心我的那位父親.這時我開始覺得此節聖靈的重要；但還未得到訊息；我繼續讀到17節:「就是真理的聖靈,乃世人不能接受的,因為不見祂,也不認識祂；你們卻認識祂；因為祂常與你們同在,也要在你們裡面.」聖靈就將這節聖經向我顯明,我認識天父,也認識耶穌基督；同樣,我認識聖靈乃真正活著的一位.我愛主耶穌,愛天父,也愛聖靈；我敬拜天父,敬拜主耶穌,也敬拜聖靈.我講道時常題到聖靈,我對聖經學院的學生也常題聖靈,我知道聖靈神學的信條；但已往聖靈對於我似非活生生的,現在對我卻成為真實的,活著的了.
\newline
\newline
使徒保羅見到十二位以弗所的長老,就問他們說,你們信的時候受了聖靈沒有?他們回答說,沒有,也未曾聽見有聖靈賜下來.」這班人甚至連聖靈都不認識,可能這乃教會軟弱的原因,他們的教會只有十二個人,教會沒有增長；保羅認為這是很重要的事,當他發現他們回答「沒有」,立即為他們祈禱,聖靈就降在他們身上；因此以弗所教會,成為最強大有力的教會.我們確知道關於聖靈很多的事,但你是否直接認識祂呢?認識並非單是在頭腦裡知道,乃是個別的經驗中認識.比如我自己可以說我認識福特總統,也知道很多關於他的事；但我和他之間並無個別的認識,實際上連見他面的機會都沒有.許多基督徒與聖靈的關係也是如此,雖知道關於聖靈的事,卻沒有切實認識聖靈.譬如我們覽閱食譜,其中有各種美味的菜式和營養豐富的食品,但是看完之後,仍然不能果腹；若是你按著食譜烹製食物,而後吃得津津有味,那才是絕不相同的一回事!憑著頭腦知識的食物,不能給你得到飽足.論到頭腦裡食物的知識和食物在肚子裡的飽足作用,我寧願選擇飽足實際的經驗,所以,我寧願得著聖靈在我生命裡面,所起的作用那種經驗,勝過徒然有頭腦的知識.
\newline
\newline
聖靈是真實的一位,祂的可愛,正如耶穌基督值得我們愛祂一樣；我們認識耶穌基督,照樣我們也認識聖靈.我們都有人間的父親,當我們思想天上的父親時,我們也就容易吸收這個觀念了.題及父時,我們的腦海裡就有一幅圖畫,我們想到主耶穌,曾經成為肉身來到世上,和我們有同樣的身體.至於聖靈,我們似乎很難想像.
\newline
\newline
聖靈是怎樣的靈?聖靈沒有身體,卻是有位格；聖經說神是靈,所以天父是靈.很多時候我們把身體和位格搞亂了.當我死了以後,我的身體變成塵土,但是我的個性人格是否停在那裡不存在呢?基督徒死後失去身體,直到復活時得回；但是他的個性仍然活在那裡,不一定要有身體才有個性和人格.聖靈雖無身體而有位格.聖經告訴我們,聖靈是三位一體之神的第三位,當你受洗禮或浸禮,乃是奉聖父聖子聖靈的名.我們所敬拜的神惟獨一位,這位神將父,子,聖靈都顯明出來.
\newline
\newline
三位一體的神,三位並非三個不同的存在；父,子,聖靈完全和諧地合作,父與子的工作從不對立,父與聖靈也從不對抗.主耶穌說:「我對你們所說的話,不是憑著自己說的:乃是父賜給我的話,我所作的乃是父差遣我作的工.」雖是三位,但工作卻是和諧的,祂們分享同一生命,雖然三位一體的神各有獨特的工作,但他們的和諧一致在聖經中清楚地明顯出來.耶穌說:「我與父原為一,人看見了我,就是看見了父.」你若想知神的樣式,請你熟讀聖經上所顯明的耶穌,從聖經你看見耶穌基督所行的事,就是看見父所行的.主耶穌對人的愛心,正是天父所作的；主耶穌所說的話,正是天父要說的話,父與子原為一.
\newline
\newline
耶穌基督所顯明的不但與天父連為一,也與聖靈連為一.「我要求父,父就就另外賜給你們一位保惠師.」(約十四16)原文聖經題到保惠師時,有個很特別的意義,希臘文題到「另外」有兩字:「不同類」,「另一個」.「另一個」的意思是同一類,但另外多了一個.本節聖經耶穌所說的「另外」就是這個意思.耶穌乃是說:「我要賜一個幫助者給你們,是「同一類」之「另一個」,同我一樣之「另外一個.」耶穌的意思是說,聖靈和祂一樣地可愛,溫柔,明事理,肯犧牲.聖靈不是耶穌,但像耶穌,祂來世上代替耶穌.主耶穌對門徒說:「我要離開你們,離開這世界,我不撇下你們為孤兒,我要差一位安慰者來代替我.」所以,聖靈來到世上,要維護耶穌基督的名,為耶穌基督辯護,來推進神的國度.聖靈帶著祂的位格到世上,代替耶穌基督在世上時所作的工.
\newline
\newline
當耶穌對門徒說,祂要關離開他們時,門徒很是戰兢；他們跟隨耶穌已經三年,耶穌成了他們生命的中心；一旦耶穌離開,怎能生存?但耶穌進一步對他們說:「我去是與你們有益的.」假如一個母親對她愛女說:「我將死了,但我死了對你有很大益處.她女兒一定要說:「媽媽你是待我最好的人,為何說這樣最不好的事呢?」當時門徒不明白為何聖靈屬靈的同在,比耶穌肉體的同在更好呢?為何要耶穌離開了反而對他們有益呢?耶穌在世時受了身體的限制,同一時間不能身處兩地.那麼,拿撒勒的人有耶穌同在,而耶路撒冷的人,就不能享有耶穌同在了!我想,如果彼得和耶穌說話,約翰在旁也要與耶穌說話,一定要等得急不及待的.耶穌離開地上後,門徒所得的耶穌不再是受身體限制的耶穌,乃是無所不在的耶穌了.耶穌對門徒說:「我要差聖靈來,永遠與你們同在.」這豈不更好嗎?
\newline
\newline
聖靈要永永遠遠住在我們裡面,祂不影響我們的外面,也將要常在我們裡面,感動我們作正確的思想.祂在我們的心靈裡,有所啟示,叫我們明白真理；祂將我裡面潔淨,使我專一愛主.祂加強我的意志,叫我遵行神的旨意.這就是耶穌所說的,祂去與門徒有益,否則聖靈就沒有降下來.藉著聖靈,耶穌就進到裡面.聖靈是耶穌基督的另一面,或可說是耶穌基督的替身；祂時常與我們同在,堅強我們,指引我們.
\newline
\newline
聖靈與耶穌基督的關係；聖靈主要的工作不是要使人成為偉人,不是要叫我們有非常奇異興奮的經驗,也不是要將屬於靈界的奇特活動忽然給我們.這樣的事祂曾作過,祂使人靈裡滿有能力,使人有聖靈充滿奇妙的經歷；(我也經歷過)祂確曾將屬靈的恩賜,給祂的教會；但這些不過是聖靈次要的工作而已.
\newline
\newline
聖靈主要的工作──耶穌說:「真理的聖靈,祂來了,就要為我作見證.」(約十五26)「祂要榮耀我,因為祂要將受於我的,告訴你們.」(約十六14)聖靈主要的工作,乃要將耶穌基督顯為大,並非要把人顯為大.親愛的弟兄姊妹!很多時候人要將自己顯為大,沒有將主顯為大.耶穌說祂要藉著聖靈,把真理向我們講明.如何能作成這工作?何處可找到關於基督的事呢?當然在聖經裡,而聖靈將聖經神的話發出亮光,叫我們能夠明白.
\newline
\newline
遠在一九二一年,我第一次參加培靈會,會場在紐西蘭南端一個帳幕裡,(時我己信耶穌)神在那處摸著我,使我得到聖靈實際的經驗而大有改變.原本我最愛閱讀小說,家父雖有許多屬靈書籍,但對我不感興趣；我每晚臨睡前也讀聖經,讀到章節越短的我越高興.自從主摸著我之後,有一次在我度假之地,發現書架上有很多小說；但那時對我毫無吸引力.而吸引我的卻是書架上,僅有的一本屬靈書藉「愚人的亮光和生命」.這書是萬宇博士解釋約翰福音的著作,這書實在作了我心靈的糧.從那時開始至今,我真正喜愛閱讀屬靈的書.並非我自己想這樣做,乃是聖靈在一夜之間把我從前的愛好改變；使我對耶穌基督的事,有了新的愛慕.這事連我自己也意想不到,聖靈帶領我對於屬靈的事有興趣,並且使我認識主耶穌,是我生命中最真實的一位；使我渴望像耶穌基督,渴望將祂的福音告訴別人.
\newline
\newline
各位!如果你遇到有些基督徒說了許多自己的經歷；但很少講到耶穌基督的話,無疑他的重點錯了.聖靈的工作乃為基督作見證,耶穌說:「聖靈要榮耀我,要將我的事向你們顯明出來.」這是真誠屬靈經驗所帶來的結果.
\newline
\newline
聖靈是真正有位格的一位.耶穌說:「我要差遣聖靈到你那裡,祂要指引你,把一切事教導你.」耶穌說這位聖靈是位指引者,是一位教師.使徒行傳一章十六節說:聖靈藉大衛的口在聖經上預言.第十三章記載安提阿教會領袖聚集禱告.毫無疑問,他們是為向普天下傳福音禱告.第二節:「聖靈說,要為我分派巴拿巴和保羅,去作我召他們所作的工.」神的聖靈呼召保羅和巴拿巴,出去傳福音給萬民聽.
\newline
\newline
當聖靈發聲向安提阿教會領袖說話時,他們正以禱告和禁食服事主.神的聖靈有沒有向你的教會領袖說話呢?有沒有吩咐他們把某弟兄姊妹分別出來,奉祂的呼召去作工呢?站在教會的地位,他們是否願意聖靈向他們發聲呢?「聖靈向眾教會所說的話,凡有耳的都應聽.」
\newline
\newline
聖靈能愛也能擔憂.神的大愛真是奇妙無窮!耶穌基督的大愛,驅使祂走上十字架；聖靈的愛也是同樣偉大.(羅十五30)保羅題到聖靈的愛,(弗四30)保羅說:「不要叫聖靈擔憂.」惟愛你的人才為你擔憂；正如不順眼的兒子令父母擔憂.照樣,聖靈十分愛我們,祂為不順眼的基督徒擔憂.聖靈既是有位格的,基督徒應把人格交在聖靈手中；讓聖靈透過了我,更大的得著我.使我順眼,好叫耶穌基督,得著更大的榮耀.
\newline
\newline
願每一個神的兒女,都能經歷這位聖靈是三位一體之神的第三位.
\newline
\newline


\newpage

\section{第48屆~港九培靈會~孫德生~第三講}
\label{sec:1104}
\textbf{聖靈降臨}
\newline
\newline
連結: \href{https://www.hkbibleconference.org/session-message/view/1104}{\texttt{https://www.hkbibleconference.org/session-message/view/1104}}
\newline
\newline
\hyperref[sec:1103]{< < < PREV SERMON < < <}
~
\hyperref[sec:toc]{[返目錄]}
~
\hyperref[sec:1106]{> > > NEXT SERMON > > >}
\newline
\newline
使徒行傳 2:1-2:8
\newline
\begin{longtable}{cl}
\hline
\hline
章節 & 經文 (和合本修訂版)\\
\hline
2:1 & \begin{tabularx}{0.7\textwidth}{X} 五旬節那日到了，他們全都聚集在一起。 \end{tabularx} \\ \\ \relax
2:2 & \begin{tabularx}{0.7\textwidth}{X} 忽然，有響聲從天上下來，好像一陣大風吹過，充滿了他們所坐的整座屋子； \end{tabularx} \\ \\ \relax
2:3 & \begin{tabularx}{0.7\textwidth}{X} 又有舌頭如火焰向他們顯現，分開落在他們每個人身上。 \end{tabularx} \\ \\ \relax
2:4 & \begin{tabularx}{0.7\textwidth}{X} 他們都被聖靈充滿，就按著聖靈所賜的口才說起別國的話來。 \end{tabularx} \\ \\ \relax
2:5 & \begin{tabularx}{0.7\textwidth}{X} 那時，有從天下各國來的虔誠的猶太人，住在耶路撒冷。 \end{tabularx} \\ \\ \relax
2:6 & \begin{tabularx}{0.7\textwidth}{X} 這聲音一響，許多人都來聚集，各人因為聽見門徒用他們各自的鄉談說話，就甚納悶， \end{tabularx} \\ \\ \relax
2:7 & \begin{tabularx}{0.7\textwidth}{X} 都詫異驚奇說：「看哪，這些說話的不都是加利利人嗎？ \end{tabularx} \\ \\ \relax
2:8 & \begin{tabularx}{0.7\textwidth}{X} 我們每個人怎麼聽見他們說我們生來所用的鄉談呢？ \end{tabularx} \\ \\
[1ex]
\hline
\hline
\end{longtable}
有一位偉大的聖經教師史博士曾講一個題目:「停留在復活節和五旬節之間.」他的意思是說,很多基督徒雖已信主重生得救,也相信耶穌曾從死裡復活；但他們屬靈的生命沒有從此起點進發,以致不知五旬對屬靈生命的意義；屬靈經驗靜止於此,沒有增長.
\newline
\newline
我剛在英文培靈會講道,題及我的妻子患上癌症,在離世與主同在之前十天；我盡量服事她,希望減輕她的痛苦.她對我說:「不必這樣週到服事,我病至如今,不必太講究啦!我仍要在靈裡長進呀.」我說:「你在世日子僅有幾天,還談得上靈命長進嗎?」
\newline
\newline
我的弟兄姊妹!你對屬靈的生命有這樣的態度嗎?如果你能認識五旬節的意義,那對你靈命有莫大的幫助.
\newline
\newline
在教會歷史中,除了加略山之外,五旬節就是最重要的經歷.初期教會較之現今教會,真有天淵之別!神的聖靈大大彰顯於初期教會中間,當時得救人數日日加增.初期教會的特徵是能力；今日教會的特徵卻是軟弱；原因是日我們沒有機會,讓聖靈在我們中間彰顯大能.
\newline
\newline
主預備人迎接五旬節,耶穌對門徒說:「你們要在耶路撒冷等候,直到你們領受從上頭上來的能力.」(路廿四49)我們的主知道門徒所需要的,乃是從天上來的能力,叫他們在耶路撒冷等候直到這能力顯現.四十天之後,耶穌回到天父那裡去；那時門徒在一間樓房裡,同心合意恆切禱告十天之久,參加的共有一百廿名敬虔的男女信徒.這是多好的禱告聚會!神早就預備聖靈降臨之日.
\newline
\newline
五旬就是五十天,耶穌被釘十字架後的五十天聖靈要降臨.五旬節也稱為初熟節(利廿三17)那初熟的果子,要經過五十天的時間；在利未記之前,神早已預定五旬節賜下聖靈.
\newline
\newline
我想門徒們十天禱告的過程中,他們必用很多時間查察自己的心求主鑑察；深信在那段時間,彼得一定情不自禁地說:「哦!主啊!我怎能三次起咒發誓否認你?「多馬一定譴責自己說:「主啊!我竟然對你起了懷疑.」雅各和約翰一定禱告說:「主啊!我倆怎能如此自私,想佔據天國的高位；為何那時忍心求神降火燒撒瑪利亞人?」我們可以想像十二門徒中的十一位,責怪自己竟然在主最需要的時候,忍心離主逃跑.當時這些人都謙卑在主的面前,我親愛的弟兄姊妹!你若肯謙卑在主面前,主要向你顯明你心靈的錯誤.
\newline
\newline
主曾對門徒說:「你們要為我作見證,使萬民作我的門徒；你們要往普天下去,傳福音給萬民聽.」門徒想到主吩咐,決心要負起這個責任；然而他們深恐軟弱無能勝任.但他們想起主的話說:「聖靈降臨在你們身上,你們就得著能力.」他們在那裡等候,呼求將聖靈的能力賜下,他們的心渴望耶穌基督成就祂的應許.後來,神的答覆臨到了,忽然從天上有響聲下來,好像一陣大風吹過.(多年前我在台灣遇到一次強烈的颱風,風速每小時二百哩.那次風災有二百人喪生,我永不忘記那令人戰兢的強風,風速無法抗拒的強風.)此乃聖靈能力的象徵,不但可聽得見,且可以看得見的象徵.繼而,舌頭如火焰,分開落在各人頭上；這並非真火,這火是由神那裡降下來作為聖靈的象徵.火能把任何可燃之物燒盡.聖靈如火,要尋找把凡屬於你的人身上的雜質,全然燒滅.神用風和火,來表明聖靈,不可抗拒的能力,門徒就都被聖靈充滿了.
\newline
\newline
被聖靈充滿,乃信徒生活中最重要的功課.當時,那一百廿人中並非部分充滿,乃是全被聖靈充滿,聖經說:「他們就都被聖靈充滿,按著聖靈所賜的口才,說起別國的話來.」為何神給他們這個記號,使他們說起別國的話?我想這與傳福音有關；當時的五旬節日,是公眾的假期,有十六個國家的人來赴節.他們各說不同的語言,門徒就用別國的語言傳福音；這些來自異他地的人,竟能聽到自己的鄉語,所以能夠向他們傳福音.論到說方言,有很多不同看法,雖然許多人不同意我的意見,但我都尊重他們.行傳二章和林前十二和十四章,所題的「說方言」並非指同一件事,兩處所題都是聖靈超自然的工作；但兩者之間確有分別.林前十二和十四章保羅說的是,如果有人說方言必須有人繙譯；但行傳二章五旬那天門徒說方言卻不必繙譯,聽者都從自的鄉談,明白使徒所說的.到底使徒那時說話的內容是什麼?聖經並未題及.只是說這班門徒用鄉談講說神的大作為.當時有沒有人聽了相信耶穌呢?到底有何效果?第二章6-7節:「各人聽見門徒用眾人的鄉談說話,就甚納悶......都驚訝希奇.」「眾人就都驚訝猜疑,」12節甚至「還有人譏誚,」13節聖經有題及有人信主；但他們聽了使徒所說的,並看見當時的情形,他們了解神果然作了大事.這時候,彼得走來,用自己的話傳福音,彼得的舌頭如火焰,他帶著正燃燒的倫理邏輯,對他們講.你們知道嗎?神的真理,人的聲音,聖靈的能力這三樣同在一起,其力量是無法抗拒的.當彼得帶著聖靈的大能說話,那些人就無法抗拒了.當司提反滿有聖靈能力作見證時,聖經說,他們就抗受不住司提反的能力和智慧.彼得宣揚福音時,那些人才信主.那班在耶路撒冷的人,確曾見過耶穌和使徒同行；他們「眾人就都驚訝猜疑,彼此說,這是什麼意思呢?還有人譏誚說,他們無非是新酒灌滿了.」(12-13)彼得說:這些人其實不是醉了,乃是被神的聖靈灌滿了.
\newline
\newline
為何有許多人要嗜酒呢?有是因人生受到很大的壓力,藉酒麻醉自己；有的人憂傷借酒消愁；有的人借酒以排除生活的單調；有的人利用酒來鼓起自己的勇氣.彼得說,你們若以為酒能使使你們成全這些事；神的聖靈也都能成就,且千倍地成就.神賜下聖靈,使你得到從神而來的刺激,勝過人生的種種壓力；使你憂愁的人生轉為喜樂,使你軟弱之中得著鼓勵.聖靈是一切問題的答案.
\newline
\newline
保羅說:「不要醉酒,酒能使人放蕩,乃要被聖靈充滿.」(弗五18)保羅將酒與聖靈作比較.大家都知道酒對人身體的影響,當醉酒時,人一反常態,懦夫卻勇如獅子；沉默寡言,滔滔不絕；吝嗇者,揮霍其畢生積蓄.你千萬別以為酒的用途很廣,能使人大衛改變,以為酒能使人得到一時之釋放；當心!隨之而來的是陷阱呀!其惡果將較前痛苦十倍.
\newline
\newline
親愛的朋友!當我們被聖靈充滿時,一切的改變都在我們身上成就了.祂使懦弱的人得到勇氣,使拙口笨舌的人,有口才為基督作見證；使吝嗇者,為神慷慨解囊奉獻；使憂傷者喜樂,聖靈乃是神為人的需要而賜下的答案.
\newline
\newline
耶穌基督成就了祂的應許,且另外賜下一位保惠師聖靈,祂是幫助我們的.請各位注意五旬節聖靈降臨,在使徒身上所起的改變.昨晚講過,當耶穌對門徒說要離開他們；門徒的心非常憂愁,行傳二至四章關於此事的記載,是何等生動,真實,明顯,榮耀!門徒深切體會耶穌基督的同在,當怹們出去傳耶穌基督的名,主顯然在他們身旁.現在耶穌雖在天上,但聖靈將祂在我們的生命中,顯得非常真實.
\newline
\newline
門徒對聖經神的話有了新的看見,彼得在行傳二章的講章,只不過引述舊約聖經而已；不像現今的傳道人,用很多時間預備講章.當時彼得沒有時間預備講章,乃是聖靈幫助彼得.他怎可能隨口引述諸多舊約聖經呢?因為他早已讀熟了.各位是否熟念聖經呢?巴不得我也如此,可惜我現已無法做到.要趁著年青念熟聖經,你若盼望在神手中有用；須讓你的頭腦充滿神的話.感謝神!在我年幼時,我的父母及老師教導我,要熟念聖靈；當我現在需要時,聖靈就將聖經神的話提醒我.彼得引述經文,其中有件趣味的事,是彼得前所未想用來配搭使用的經文.彼得說今日的事「正是先知約珥所說.神說:在末後的日子,我要將我的靈灌溉我的僕人和使女,他們就要說預言.」(16)聖靈光照彼得的思想,使彼得善為引用聖經證明神的信實.這是聖靈的作為,使我們從聖靈有新的看見和了解.聖靈賜與門徒滿有權柄的話,在聽眾心中,發出深的感動.
\newline
\newline
彼得講道的果效(二27)是否聽眾在那裡打盹呢?是否彼得在勸導他們呢?彼得乃是控告這班人,把神的兒子釘十字架.聖靈賜給彼得滿有能力的說話,使聽者覺得扎心.
\newline
\newline
弟兄姊妹!當你為主作見證時,是否帶著權柄和能力?當主耶穌講道時,人都說祂的教訓,不像那班法利賽人和文士,祂是帶著權力的語氣來教訓人.神的心意要為祂傳話的人,帶著屬靈的權柄和能力.保羅說:「我說話帶著屬靈的權柄,彰顯神的大能.」
\newline
\newline
聖靈降臨以後,門徒們大有膽量,毫無畏懼；離此數日前,他們因懼怕猶太人,關上門躲在屋裡；現在竟勇敢如獅子.「他們就都被聖靈充滿,放膽講論神的道.」(四31)他們現在面對逼迫他們的人.雅各為主殉道捨命,但是不能嚇倒他們；並非他們生來有此膽量,乃是聖靈使然.你是否膽怯不敢為基督作見證?是否因懼怕而閉口不言?我不敢說現在我毫無畏懼；但我相信聖靈給我勇氣,放膽為基督作見證.當這班使徒講道時,有一種敬畏的心臨到眾人；這是今日我們難以見到的,我們也沒有時候見到在神聖潔之下震抖.多年前,有一位在中國的著名韓國宣教師,他心裡渴慕聖靈的能力,悉心研究聖經中,關於聖靈的教訓.全卷聖經中所有聖靈的真理,他都一一查考；每有亮光,就和同工們分享.有一次他和十位中國同工出去佈道,會場已坐滿了,十位中國同工坐在講員後面.他一開始講道,就有一種聖靈同在的不尋常感覺,他說:「看哪!神的羔羊,除去世人罪孽的.」那時他發現聽眾之中,起了極大的敬畏.講完道,全體會眾都站起來；他覺得很驚異,不知如何是好.他一轉身,看見那十位同工,有著同樣向神敬畏的表情；其中一位對他說:「親愛的弟兄啊!我們所尋求的那位,今晚在聖靈同在之下,真正臨到我們了.」全體會眾都歸向神,神的話語帶著聖靈的能力,在那晚彰顯出來.這就是使徒面對五旬節,五旬節對他們的意義.其結果不但是撒種,同時也有收割.我們將神的道來撒種,固然重要；因沒有撒種,便無可收割.為何今日見不到收割,為何見不到更多人,得進入神的國度?
\newline
\newline
請留心五旬節後有何事發生?五旬節那天有三千人信主,另一次有五千人信主,主將得救的人天天加給他們.
\newline
\newline
五旬節乃是神在教會作工,其中的例證榜樣之一.五旬節的事實不會重新出現,正如加略山的十字架,不需要再出現一樣.我們的耶穌基督只需一次受死,神賜下聖靈只有一次.但五旬節可以延續,當日發生的事,今日也可以再發生.因有同樣一位聖靈,條件仍然是一樣.神並無改變,耶穌基督沒有改變,福音也無改變.當有人起來,願意符合聖靈能力的條件時,我們就會發現當日五旬節的情況延續下去,正如當日五旬節所發生的一樣.
\newline
\newline


\newpage

\section{第48屆~港九培靈會~孫德生~第四講}
\label{sec:1106}
\textbf{罪}
\newline
\newline
連結: \href{https://www.hkbibleconference.org/session-message/view/1106}{\texttt{https://www.hkbibleconference.org/session-message/view/1106}}
\newline
\newline
\hyperref[sec:1104]{< < < PREV SERMON < < <}
~
\hyperref[sec:toc]{[返目錄]}
~
\hyperref[sec:1108]{> > > NEXT SERMON > > >}
\newline
\newline
以弗所書 4:25-4:32
\newline
\begin{longtable}{cl}
\hline
\hline
章節 & 經文 (和合本修訂版)\\
\hline
4:25 & \begin{tabularx}{0.7\textwidth}{X} 所以，你們要棄絕謊言，每個人要與鄰舍說誠實話，因為我們是互為肢體。 \end{tabularx} \\ \\ \relax
4:26 & \begin{tabularx}{0.7\textwidth}{X} 即使生氣也不要犯罪；不可含怒到日落， \end{tabularx} \\ \\ \relax
4:27 & \begin{tabularx}{0.7\textwidth}{X} 不可給魔鬼留地步。 \end{tabularx} \\ \\ \relax
4:28 & \begin{tabularx}{0.7\textwidth}{X} 偷竊的，不要再偷；總要勤勞，親手做正當的事，這樣才可以把自己有的，分給有缺乏的人。 \end{tabularx} \\ \\ \relax
4:29 & \begin{tabularx}{0.7\textwidth}{X} 一句壞話也不可出口，只要隨著需要說造就人的好話，讓聽見的人得益處。 \end{tabularx} \\ \\ \relax
4:30 & \begin{tabularx}{0.7\textwidth}{X} 不要使神的聖靈擔憂，你們原是受了他的印記，等候得救贖的日子來到。 \end{tabularx} \\ \\ \relax
4:31 & \begin{tabularx}{0.7\textwidth}{X} 一切苦毒、憤怒、惱恨、嚷鬧、毀謗，和一切的惡毒都要從你們中間除掉。 \end{tabularx} \\ \\ \relax
4:32 & \begin{tabularx}{0.7\textwidth}{X} 要仁慈相待，存憐憫的心，彼此饒恕，正如神在基督裡饒恕了你們一樣。 \end{tabularx} \\ \\
[1ex]
\hline
\hline
\end{longtable}
聖靈的工作,其中最重要的乃與罪有關；因祂是聖潔,所以任何不聖潔的,和罪惡有關係的都令祂擔憂傷痛.
\newline
\newline
人都不願意談及罪,然而罪和現世代關係重大.罪是聖經記載的重要中心之一,也是人生過程重要遭遇.
\newline
\newline
人生最大的問題是罪嗎?聖經由創世記第三章開始至啟示錄廿二章,一直都題及罪.全部聖經也都題及,神如何為我們預備救恩,叫人從罪中得釋放.
\newline
\newline
罪是一切傷心痛苦的來源,罪把耶穌基督釘在十字架上；罪使基督徒失敗跌倒；罪使人與神失去交通；罪把我們的喜樂搶奪了；罪把神的律法毀壞,令神的心破碎.
\newline
\newline
聖靈感動未信者知罪(約十六8-9)「祂既來了,就要叫罪人知道自己的罪......因他們不信我.」這裡所說的罪,不是雙數式的；並不是說聖靈來了,要叫人為許多罪責備自己；而是單數式的,乃是說有一件大罪幪罩著人類.並非說關於殺人姦淫,嫉妒的思想上等等的罪,乃是不信主耶穌的罪,他們自己的良心,起來定自己的罪；聖靈所作的,乃是要叫世人有新的看見,覺得不信主罪的嚴重.
\newline
\newline
對世人而言,罪是外在的行動,只要能掩人耳目就算幸運；這不是耶穌基督的教訓.主耶穌說罪非你手所犯的動作,而是頭腦心思意念所犯的；心的嫉妒,憎恨,和手謀殺是同等的罪.不潔淨的心思,和犯姦淫的行動,是同樣的罪；世人就不相信這一套,惟聖靈感動人時；人才開始校正自己的看法,才和神的看法一樣.
\newline
\newline
罪是神最憎恨的,也是件最可怕的事.自從耶穌降世,罪惡轉向一個不同的中心.耶穌來到世上之後,罪的基本要素,是人不相信耶穌；所以聖靈要感動叫人知罪,但世人不明白罪和信耶穌有何相干.神說不信耶穌就是罪,這罪可使你下到地獄去.
\newline
\newline
聖靈表明耶穌基督是救主,你可相信祂也可拒絕祂；接受耶穌的,就是接受救恩和永生.拒絕耶穌,不相信祂的,就永遠失喪.故此,罪是極其嚴重的；但世人不曉得罪的嚴重性.五旬節之日,猶太人沒有想到把耶穌釘十字架有何嚴重；當聖靈降臨,事情就完全改觀了.彼得向他們說:「這位耶穌,你們就藉著無法之人的手,把祂釘在十字架上殺了.」(徒二23)「你們釘在十字架上的這位耶穌,神已經立祂為主為基督了.」(36節)甚至彼得呼籲:「你們各人要悔改,奉耶穌基督的名受洗,就必領受所賜的聖靈；眾人聽見這話,覺得扎心!就對彼得和其餘使徒說,弟兄們,我們當樣行?」(37節)情形完全改變了.耶穌說:「祂既來了,就要叫世人為罪,......為罪,是因他們不信我.」感謝神!在那天有三千人歸信主.
\newline
\newline
可能在座中有多人未信主,已信者,你是否接受耶穌作你的救主?你能說:「我認識祂,信祂,愛祂；祂拯救我,我已從祂得到永遠的生命嗎?」如果你不敢這樣說,今晚聖靈要感動你,使你知道不相信耶穌的罪.最奇妙的,是你可在你的座位上,甚至分寸不移動,只要對主說:「主啊!我打開心門求你進來,我相信你.」
\newline
\newline
多年前有一次我在澳洲長老會講道,我講福音,講得到救恩的方法.我說:「你們中間有誰願意心裡相信,口裡承認耶穌?」這句話不過是我講章的內容,想不到詩班座中,有一牧師的兒子回答說:「孫牧師!我願意.」繼之,前後左右又有八至九次「我願意」的聲音呼出；他們願意心裡相信,口裡承認耶穌；聖靈使這班青年覺罪,願意得到救主.後來那第一位舉手的青年,當了醫生傳教士事主多年.
\newline
\newline
聖靈感動叫人信主使人知罪.未知各位的經驗如何?以我說來,信主之後,我為罪所流的眼淚比前更多；主使我更加知道自己的罪；因為聖靈愛我們,要叫我們的生命免受罪的摧殘,叫我們知道罪的可怕；聖靈復興我們的時候,叫我們覺得自己有罪.卅年前我在中國,我的朋友擔任英文培靈會的主席,他和韋牧師及我三人往各地旅行訪問到了重慶.韋牧師和我有機會向重慶神學院學生講道,約有百人赴會,他們恆切禱告,求神賜下恩雨；但他們卻沒有認識是怎樣的恩雨.聖靈感動他們認自己的罪,從第一次聚會開始,我們便看見聖靈掌管這聚會；神帶領韋牧師和我互相配搭,聖靈真是對這班學生說話；他們的良心受激動,一切隱藏和被遺忘的罪,當時都顯露出來；(該院院長陳崇桂牧師曾寫書「火焰的燈」記載此次復興的經驗)他們流淚,傷痛,悔改,他們接二連三地將心傾倒向神認罪.每次會後禱告認罪達兩小時,如校方允許,他們實在不肯停止,當時學校停課一週,供大家有足夠時間禱告.禱告的靈臨到各人身上,個個切實認罪；聖靈的工作使他們對罪有憎恨感,他們不單認罪,同時也償還,錢債一一償還,且紛紛寫信道歉或見證,一下子學校信箋和郵票全部用完,後來有人奉獻,繼續供應.當時學生們三五成組澈夜禱告.當聖靈降臨,他們的心都受了感動,對付罪惡,和神有正當關係,內裡有了奇異的改變,對神有了新的奉獻,重新把自己交在神的手中.後來時局影響,神學院被迫停辦,那是最後一次的聚會使他們復興,神的恩典為他們預備,教他們如何渡過前面艱難的日子.我不知那些神學生現在那裡；但我相信他們確曾經歷聖靈同在的喜樂；神的靈降在他們身上,叫他們知罪悔改,與神建立正當的關係.
\newline
\newline
罪惡叫聖靈擔憂,「不要叫神的聖靈擔憂」(弗四30).擔憂的意思,就是把傷痛加在別人身上.我們怎能把痛加在聖靈身上呢?擔憂包含著愛的成份,惟愛我的人才為我擔憂；敵我的人只有激憤仇恨,絕無擔憂可言.聖經說我們的身體是聖靈的殿；換言之,聖靈住在每個基督徒心靈裡,如果我容許罪惡進入,聖靈便為我傷痛擔憂.不要叫神的聖靈擔憂,意思是指現在要停止犯罪,不叫神的聖靈擔憂.耶穌說:「聖靈來了,要住在你們裡面直到了永遠.」但我們常忘記有位聖靈住在我們裡面,以致犯罪叫祂傷痛.聖靈住在我們裡面,為我們有生之年工作.
\newline
\newline
不順服聖靈的命令,叫聖靈擔憂.惟獨耶穌從未叫聖靈擔憂.耶穌說:「我常常行天父所喜悅的事.」我們當學習主耶穌的榜樣,不叫聖靈擔憂.讓聖靈完全掌管我們的生命.
\newline
\newline
以弗所書中,保羅題及足以叫聖靈擔憂的:
\newline
\newline
1.不合宜的話　聖靈是真理的靈,所以任何不合真理,不誠實的事,都會叫聖靈擔憂.「所以你們要棄絕謊言,各人與鄰舍說實話；因為我們是互相為肢體.」(四25)有人以為撒一點小謊言沒有關係.各位做父母的!當你的兒女初次撒謊,那是不是小朋友的謊言呢?總之,一切不誠實的話,甚至出自小孩子口中的,都是謊言.虛偽或誇大其詞的言語,都屬謊言之列；欺騙奸詐污穢的言語,(29節)都會叫聖靈擔憂.「淫詞,妄語,和戲笑的話,都不相宜.」(五4)保羅並非要我們把人生的樂趣完全排除,連幽默的話都不可講；他的意思是不說無助於人,與人無益,不能造就人的話.
\newline
\newline
2.不受約束的感情叫聖靈擔憂　「生氣卻不要犯罪；不可含怒到日落；......」(26節)「一切苦毒,惱恨,忿怒,嚷鬧,譭謗,並一切的惡毒,都當從你們中間除掉.」(31節)論到發怒,聖經有數次記載:當耶穌看見那班所謂虔誠的法利賽人,搾取寡婦的東西時,祂發怒了.當主耶穌要醫治手枯乾的病人,法利賽人加以攔阻；那時,主發怒了.當主看見聖殿所發生的事,祂用繩當鞭子,把牛羊和買賣的人趕出去；鴿子放出去,兌換銀錢的桌子翻倒.耶穌說:「我父的殿,必作禱告的殿,你們倒使它成為賊窩了.」耶穌基督沒有自私,沒有犯罪,祂的發怒是義怒.當耶穌基督被人戲弄,受人凌辱；人揪祂的鬍鬚,給祂戴棘荊冕冠.甚至釘祂在十字架上,祂都沒有發怒；但當祂看不公義的事,發生在人身上,天父的名被污辱的時候,祂發怒了.弟兄姊妹!惟不自私之怒,才是無罪之怒.
\newline
\newline
不誠實的行為叫聖靈擔憂,我也曾經有過偷竊的行為,但聖靈感動我,使我以同等價值的金錢去償還.偷竊的方式很多,如遲到早退,事倍功半,偷竊僱主的時間,無異竊取財物.
\newline
\newline
不受控制的情緒,也令聖靈擔憂.(五3)(至於淫亂,並一切污穢,或是貪婪,在你們中間連題都不可,方合聖徒的體統.)不道德與兩性之間不清白的事,現社會充滿了這種污穢.保羅說這樣的事,在你們中間連題都不可.貪婪是種奢慾,這是現社會的趨勢.作為神的兒女,現應渴望神工作擴展和增多；但魔鬼卻叫我們對財物貪得無厭.保羅說:「貪心的,就和拜偶像一樣.」(5節)「我無論在什麼景況,都可以知足......我知道怎樣處卑賤,也知道怎樣處豐富；或飽足,或飢餓,或有餘,或缺乏；隨事隨在,我都得了秘訣.」(腓四11-12)
\newline
\newline
當聖經神的話語對你失去吸引力時,聖靈為你擔憂你禱告讀經較前減少,你從前曾克服罪惡；但現在罪惡管理你,從前你深感有主同在,如今已漸減退了.聖靈所結喜樂的果子已不在你心中了；事奉和見證都失去能力了,聖靈為你擔憂.在你生命之中是否有確據叫你知道聖靈確曾為你擔憂?聖靈為我們擔憂?聖靈為我擔憂,是否離棄我們呢?大衛對神說:「求你不要從我收回你的聖靈.」(詩五十一11)這樣的情形,在五旬節之前或者可能；但主耶穌:「我要求父就另外賜給你們一位保惠師,叫祂永遠與你們同在.」這是何等蒙福的應許!聖靈固然為我們擔憂,但祂卻永遠不離開我們.
\newline
\newline
聖靈叫我們知道罪的嚴重性　聖靈好比室內的陽光,如果我們拉下窗簾,看不見.是你自甘與陽光隔開.聖靈樂意隨隨時幫助我們,但因罪惡攔阻,聖靈也就愛莫能助.聖靈現在不再向我們彰顯基督了,聖靈的果子也不再在我們身上結出來；我們將因此一敗塗地了,直到聖靈不斷感動你在神面前悔改；你將窗簾揭開,太陽仍然供應熱力和亮光.
\newline
\newline
我們若要得回聖靈的同在和喜樂,第一,當潔淨自己(林後七1)「親愛的弟兄阿!我們既有這等應許,就當潔淨自己,除去身體靈魂一切的污穢,敬畏神,得以成聖.」有種潔淨是神賜給的,還有一種潔淨必須自己去做,唯一方法就是對聖靈向我們顯明令祂擔憂的罪.下定決心,起來對付,要承認自己的罪；接受耶穌基督的寶血來洗脫潔淨.罪是神最憎惡的,所以我們當與神同樣地憎恨罪惡.認罪是自己的本份,洗罪在於神,兩者必相輔而行,罪方得潔淨.第二,相信並感謝神.第三,再次把你的生命順服在聖靈管理之下,一生讓聖靈管理.
\newline
\newline
耶穌基督從未令聖靈擔憂,因為祂讓聖靈完全管理.
\newline
\newline


\newpage

\section{第48屆~港九培靈會~孫德生~第五講}
\label{sec:1108}
\textbf{聖靈充滿}
\newline
\newline
連結: \href{https://www.hkbibleconference.org/session-message/view/1108}{\texttt{https://www.hkbibleconference.org/session-message/view/1108}}
\newline
\newline
\hyperref[sec:1106]{< < < PREV SERMON < < <}
~
\hyperref[sec:toc]{[返目錄]}
~
\hyperref[sec:1110]{> > > NEXT SERMON > > >}
\newline
\newline
以弗所書 5:15-6:9
\newline
\begin{longtable}{cl}
\hline
\hline
章節 & 經文 (和合本修訂版)\\
\hline
5:15 & \begin{tabularx}{0.7\textwidth}{X} 你們要謹慎行事，不要像無知的人，要像智慧的人。 \end{tabularx} \\ \\ \relax
5:16 & \begin{tabularx}{0.7\textwidth}{X} 要把握時機，因為現今的世代邪惡。 \end{tabularx} \\ \\ \relax
5:17 & \begin{tabularx}{0.7\textwidth}{X} 不要作糊塗人，要明白主的旨意如何。 \end{tabularx} \\ \\ \relax
5:18 & \begin{tabularx}{0.7\textwidth}{X} 不要醉酒，酒能使人放蕩；要被聖靈充滿。 \end{tabularx} \\ \\ \relax
5:19 & \begin{tabularx}{0.7\textwidth}{X} 要用詩篇、讚美詩、靈歌彼此對說，口唱心和地讚美主。 \end{tabularx} \\ \\ \relax
5:20 & \begin{tabularx}{0.7\textwidth}{X} 凡事要奉我們主耶穌基督的名常常感謝父神。 \end{tabularx} \\ \\ \relax
5:21 & \begin{tabularx}{0.7\textwidth}{X} 要存敬畏基督的心彼此順服。 \end{tabularx} \\ \\ \relax
5:22 & \begin{tabularx}{0.7\textwidth}{X} 作妻子的，你們要順服自己的丈夫，如同順服主。 \end{tabularx} \\ \\ \relax
5:23 & \begin{tabularx}{0.7\textwidth}{X} 因為丈夫是妻子的頭，如同基督是教會的頭；他又是這身體的救主。 \end{tabularx} \\ \\ \relax
5:24 & \begin{tabularx}{0.7\textwidth}{X} 教會怎樣順服基督，妻子也要怎樣凡事順服丈夫。 \end{tabularx} \\ \\ \relax
5:25 & \begin{tabularx}{0.7\textwidth}{X} 作丈夫的，你們要愛自己的妻子，正如基督愛教會，為教會捨己， \end{tabularx} \\ \\ \relax
5:26 & \begin{tabularx}{0.7\textwidth}{X} 以水藉著道把教會洗淨，使她成為聖潔， \end{tabularx} \\ \\ \relax
5:27 & \begin{tabularx}{0.7\textwidth}{X} 好獻給自己，作榮耀的教會，毫無玷污、皺紋等類的缺陷，而是聖潔沒有瑕疵的。 \end{tabularx} \\ \\ \relax
5:28 & \begin{tabularx}{0.7\textwidth}{X} 丈夫也應當照樣愛妻子，如同愛自己的身體；愛妻子就是愛自己了。 \end{tabularx} \\ \\ \relax
5:29 & \begin{tabularx}{0.7\textwidth}{X} 從來沒有人恨惡自己的身體，總是保養愛惜，正像基督待教會一樣， \end{tabularx} \\ \\ \relax
5:30 & \begin{tabularx}{0.7\textwidth}{X} 因我們是他身體的肢體。 \end{tabularx} \\ \\ \relax
5:31 & \begin{tabularx}{0.7\textwidth}{X} 「為這個緣故，人要離開父母，與妻子結合，二人成為一體。」 \end{tabularx} \\ \\ \relax
5:32 & \begin{tabularx}{0.7\textwidth}{X} 這是極大的奧祕，而我是指基督和教會說的。 \end{tabularx} \\ \\ \relax
5:33 & \begin{tabularx}{0.7\textwidth}{X} 然而，你們每個人都要愛妻子，如同愛自己一樣；妻子也要敬重她的丈夫。 \end{tabularx} \\ \\ \relax
6:1 & \begin{tabularx}{0.7\textwidth}{X} 作兒女的，你們要在主裡聽從父母，這是理所當然的。 \end{tabularx} \\ \\ \relax
6:2 & \begin{tabularx}{0.7\textwidth}{X} 當孝敬父母，使你得福，在世長壽。這是第一條帶應許的誡命。 \end{tabularx} \\ \\ \relax
6:3 & \begin{tabularx}{0.7\textwidth}{X} 見上節 \end{tabularx} \\ \\ \relax
6:4 & \begin{tabularx}{0.7\textwidth}{X} 作父親的，你們不要激怒兒女，但要照著主的教導和勸戒養育他們。 \end{tabularx} \\ \\ \relax
6:5 & \begin{tabularx}{0.7\textwidth}{X} 作僕人的，你們要懼怕戰兢，用誠實的心聽從你們肉身的主人，好像聽從基督一般； \end{tabularx} \\ \\ \relax
6:6 & \begin{tabularx}{0.7\textwidth}{X} 不要只在人的眼前這樣做，像僅是討人的喜歡，而是作基督的僕人，從心裡遵行神的旨意， \end{tabularx} \\ \\ \relax
6:7 & \begin{tabularx}{0.7\textwidth}{X} 甘心服侍，好像服侍主，不像服侍人， \end{tabularx} \\ \\ \relax
6:8 & \begin{tabularx}{0.7\textwidth}{X} 因為知道每個人所做的善事，不論是為奴的或是自主的，都必按所做的從主得到賞賜。 \end{tabularx} \\ \\ \relax
6:9 & \begin{tabularx}{0.7\textwidth}{X} 作主人的，你們待僕人也是一樣，不要威嚇他們，因為知道他們和你們在天上同有一位主，他並不偏待人。 \end{tabularx} \\ \\
[1ex]
\hline
\hline
\end{longtable}
未知我們中間有多少人,對自己屬靈的生命現已感到滿足?我想多數人都覺得還有長遠的路程,需繼續奔跑.聽說有兩位牧師,彼此是朋友；有一早晨,他們從不同的方向走來,其中一位走路快捷,另位問他,「你急於往那裡去?」答:「我要趕向一個十全十美之地,請勿攔我的去路,因前面當跑的路遙遠.」
\newline
\newline
使徒行傳有九次記載被聖靈充滿,可見其重要性了.
\newline
\newline
主說:「不多幾日,你們要受聖靈的洗(浸).」(徒一5)主告訴門徒,他們即將要受聖靈的洗(浸),到了五旬節之日,他們新的經歷果然成就了,聖經說:「他們就都被聖靈充滿了.」自此,聖經不再用「你們要受聖靈的洗」來記載了；這兩句話的意思雖然不是絕對相同,但都同在五旬節成就了.關於「聖靈的洗」有很多不同的看法,與我看法不同者,我都敬重他們.(林前十二13):「都從一位聖靈受洗,」在新約書信中題及聖靈的洗僅此一次,其他在四福音和行傳裡,論及此事均謂「聖靈充滿」,到了(林前十二13),作者寫「都從一位聖靈受洗,成了一個身體.」這是每個信徒必需的經歷,保羅回顧這個已經成就的事實,他說此乃聖靈的工作；藉著聖靈來臨,把信徒連結起來成為基督的身體,這事在五旬節之日成就了.當人信了耶穌,他就受聖靈的洗,成為基督的身體,成為基督身體的肢體.全卷新約的聖經無一處命令我們一定要受聖靈的水禮.照我的看法,每一信徒在相信耶穌的剎那間,已受了聖靈的洗禮；但信徒都在這命令之下要被聖靈充滿.保羅在以弗所五章18節說「不要醉酒」,不要使酒來壟斷你的生命.「乃要被聖靈充滿」讓聖靈來佔有你,管理你；這是神所命令的事,並非模稜兩可,亦非可有可無的事.五旬節之日,所有在樓房上聚集的門徒都被聖靈充滿,事後他們都出去為主作見證,這些見證都反應出被聖靈充滿的結果.保羅寫信給以弗所教會,將此命令給他們.以弗所教會是較諸其他教會堅強有力的,信徒們都從那一次得著新的生命,聖靈進入信徒生命中居住.他們受了聖靈的印記,得著聖靈為基業；但保羅對這成熟的教會說:「乃被聖靈充滿」這表明他們雖是好的信徒,但對於聖靈充滿罔然不知.
\newline
\newline
為何我們要被聖靈充滿?
\newline
\newline
1.需要聖靈的能力過聖潔的生活.任何人都無法靠自己的力量,過聖潔生活；肉體的引誘使我們犯罪,所以需要比自己更強的力量.
\newline
\newline
2.需要聖靈的力量,來履行基督徒的本分為耶穌基督作見證.主耶穌說:「但聖靈降臨在你們身上,你們就必得著能力,並要......作我的見證.」(徒一8)主耶穌吩咐門徒等候在耶路撒冷,直到得著從上面來的能力.他們需要聖靈充滿,神確將此經驗賜給他們:他們帶此能力去為主作見證.所以我們需要聖靈的充滿,過聖潔的生活.在這滿有需要的社會,為耶穌基督作有效見證.
\newline
\newline
被聖靈充滿並非神秘的經驗,以致有無被聖靈充滿卻不自知.當神將這命令給我們時,我們自己可以知道,究竟有沒有遵守這個命令.
\newline
\newline
有一次我聽到一位講員說,一隻玻璃杯,雖是空的,但裡面卻佈滿空氣,如果把水倒滿杯中,就空杯中原有的空氣,被逐出而充滿著水.他用此例證解釋人被聖靈充滿,我認為此例證不甚完滿.人並非一個死沉沉的水杯,人是活的,有選擇的權利；聖靈也是有位格活著的,我的人性現在已被聖靈的格控制和管理.祂進入我的生命裡,管理我的生命；但除非我邀請聖靈進來,否則不肯這樣做；因為聖靈從不勉強人之所為.耶穌說:「看哪!我站在門外扣門,聽到我的聲音就開門的,我要進到他裡面去.」祂在門外等待.照樣,聖靈也必等我們願意時才進來.
\newline
\newline
我年青時住在紐西蘭的最南端,該市鎮沒有什麼可娛樂的,甚至連電影都沒有.有一次,外地來了個催眠術的,大家都紛紛去看表演,那人請了數位當地商界名人坐於台上,其中有位是我認識的；他同意讓那人向他施展催眠術,果然被催入睡眠狀態.那人對他發問說:「假如你是一國之王,你要作什麼?」答:「假如我是國王,我要推開木門,塞入一塊肥肉.」所答的話,到底是被催眠者的意見,或是催眠者的意見?當然是催眠者把意見放入被催眠者身上；所以若非被催眠者願意接受催眠者施展工夫,催眠者則無能為力.當你我願意邀請聖靈進入我們的生命裡,聖靈才能夠管理我們.聖靈十分願意管理我們的生命,好叫我們能像耶穌基督.
\newline
\newline
當聖靈管理我們的個性時,祂作些什麼?我們邀請聖靈進來,祂就在我們心靈的深處照亮我們的思想,好叫我們明白屬靈的真理；聖經說,屬血氣的人,不可能明白屬靈的真理.然後聖靈潔淨我們污穢的情緒,有時我們准許情緒進入罪中,讓情慾脾氣來管理自己；但聖靈潔淨使我們的愛情,轉移到耶穌基督身上,叫我們惟獨愛主.
\newline
\newline
聖靈堅強我們的意志.以我來說,感興趣的事我就專心去做,不感興趣的事就無心去作.有時強烈的引誘臨到,我們的內心,一部分要抗拒,一部分不抗拒；在因此心思起伏,意志薄弱,甚至跌倒罪中.當聖靈管理我們,在最危急之時,祂幫助堅強我們的意志.
\newline
\newline
怎能知道自己被聖靈充滿.不錯,聖靈確曾把屬靈恩賜賜給人；但有恩賜的人不一定是被聖靈充滿.我認識一位滿有屬靈恩賜的人,但是他的脾氣很可怕,使人難以和他相處,被聖靈充滿,必定有明顯的象徵,首先在個性和生命上.保羅說:「......乃要被聖靈充滿.當用詩章,頌詞,靈歌,彼此對說,口唱和的讚美主；凡事要奉我們耶穌基督的名,常常感謝父神.」(五18-20)保羅認為聖靈必然成就這些事,在一個被聖靈充滿之人的生命裡.用詩章,頌詞,靈歌,彼此對說,就是彼此談話,著重屬靈的成份.各位!你每日所談論的,是屬靈或屬世的事?你是個讚美或埋怨的基督徒呢?
\newline
\newline
先父個子矮小,他是個具有愛心的基督徒；我結婚後他和我們同住多年.每晨他一起床便唱設,他的心充滿了讚美；我不敢說他曾將這恩賜傳給他的兒子；但這確是神的心意要我們以讚美來開始一天的生活.「口唱心和」或者你難以做到,但你若將生命交給聖靈管理,祂必為你成就.
\newline
\newline
凡事獻上感謝,你是否每時每事都獻上感謝?若靠自己的確做不到；但聖靈能幫助我們,保羅熟揀這個技巧,他說:「無論任何景況,我都學會了知足,我都要獻上感謝.」我很欣賞他說「學會了.」保羅在感謝的恩典上有長進；甚至在患難中他都獻上感謝.這樣高超的境界,實非一時可達；但你有否向這個境界長進?
\newline
\newline
彼此順服(21節)就是凡事勿一意孤行；聖經說「當有謙讓順服的心.」
\newline
\newline
保羅題及被聖靈充滿後之家庭生活(22節)「你們作妻子的,當順服自己的丈夫,如同順服主.」順服丈夫,並非說婦女比男人低了一級,任何聰明人都知道妻子比本身更好,(此非本節聖經所指)保羅說在家庭中丈夫是妻子的頭,有總決事情的權力,如果抗拒神的程序,就將失去神的賜福.固然保羅說丈夫是妻子的頭,要順服自己的丈夫；但21節說「要彼此順服」這是雙方面的.保羅又說:「你們作丈夫的,要愛你們的愛」實在高超!世俗的夫婦看見將感到基督徒的夫婦生活,是何等幸福；實在與他們有所分別.可惜基督徒常常距離神的目標太遠!「丈夫也當照樣愛妻子,如同愛自己的身子；」(28節)保羅覺得作丈夫的常常沒有盡到本分；所以他鄭重再說:「丈夫也當照樣愛妻子,如同愛自己一樣.」(33節)
\newline
\newline
被聖靈充滿的家庭　夫婦要彼此順服,作丈夫的要愛妻子,作妻子的要順服丈夫.至於作兒女的,要在主裡聽從父母......」(六1-2)青年弟兄姊妹!你是否聽從父母孝敬父母?你需要聖靈幫助你.「你們作父親的,不要惹身兒女的氣,只要照著主的教訓和警戒,養育他們.」(3節)這裡沒有題及作母親的應如何待兒女,可能作母親的較為瞭解兒女.保羅繼續又題及主僕之間應如何對待,說:「你們作僕人的,要懼怕戰兢,用誠實的心聽從你們肉身的主人,好像聽從基督一般.」(5節)而作主人的不要壓搾僕人的血汗,保羅說:「你們作主人的,待僕人,也是一理,......因為知道他們和你們,同有一位主在天上,祂並不偏待人.」(9節)
\newline
\newline
一個被聖靈充滿的人,他的個性,渴望進展越發像耶穌；日漸看見自己醜陋的面目,變成像主一般可愛的形像.
\newline
\newline
怎能知道確實得到聖靈充滿的經驗,當怎樣做才能被聖靈充滿?
\newline
\newline
1.如果你查察自己未被聖靈充滿,你要坦白承認這件事實.說:「主啊!我未被聖靈充滿,我渴望被聖靈所充滿.」
\newline
\newline
2.如果神給你看見自己的罪惡,你要認罪悔改,撇棄罪惡.
\newline
\newline
3.要認定耶穌基督作你生命的主宰.
\newline
\newline
4.邀請聖靈進來聖靈的目的,乃要耶穌基督成為你生命之主；我們要邀請聖靈來管理我們的生命.
\newline
\newline
當我十九歲時,我確曾邀請聖靈,我毫無保留地,聖靈在我生命裡管理我的一生；從此情形完全不同已往了.生命得到神賜給我的豐盛,情感的激動不再時有時無了.
\newline
\newline
當你邀請聖靈進來,你就確實相信聖靈為你成就一切:你讀經時,這位聖靈要帶領你明白真理,祂要向你發聲,使你在禱告時深感與神親近,使你的禱告蒙神垂聽；當你為主作見證時,聖靈要加給你口才和能力.
\newline
\newline


\newpage

\section{第48屆~港九培靈會~孫德生~第六講}
\label{sec:1110}
\textbf{聖靈所結的果子}
\newline
\newline
連結: \href{https://www.hkbibleconference.org/session-message/view/1110}{\texttt{https://www.hkbibleconference.org/session-message/view/1110}}
\newline
\newline
\hyperref[sec:1108]{< < < PREV SERMON < < <}
~
\hyperref[sec:toc]{[返目錄]}
~
\hyperref[sec:1111]{> > > NEXT SERMON > > >}
\newline
\newline
加拉太書 5:16-5:23
\newline
\begin{longtable}{cl}
\hline
\hline
章節 & 經文 (和合本修訂版)\\
\hline
5:16 & \begin{tabularx}{0.7\textwidth}{X} 我說，你們要順著聖靈而行，絕不可滿足肉體的情慾。 \end{tabularx} \\ \\ \relax
5:17 & \begin{tabularx}{0.7\textwidth}{X} 因為肉體的情慾和聖靈相爭，聖靈和肉體相爭，這兩個彼此敵對，使你們不能做所願意做的。 \end{tabularx} \\ \\ \relax
5:18 & \begin{tabularx}{0.7\textwidth}{X} 但你們若被聖靈引導，就不在律法之下。 \end{tabularx} \\ \\ \relax
5:19 & \begin{tabularx}{0.7\textwidth}{X} 情慾的事都是顯而易見的；就如淫亂、污穢、放蕩、 \end{tabularx} \\ \\ \relax
5:20 & \begin{tabularx}{0.7\textwidth}{X} 拜偶像、行邪術、仇恨、紛爭、忌恨、憤怒、自私、分派、結黨、 \end{tabularx} \\ \\ \relax
5:21 & \begin{tabularx}{0.7\textwidth}{X} 嫉妒、醉酒、荒宴等類。我從前告訴過你們，現在又告訴你們，做這樣事的人必不能承受神的國。 \end{tabularx} \\ \\ \relax
5:22 & \begin{tabularx}{0.7\textwidth}{X} 聖靈的果子就是仁愛、喜樂、和平、忍耐、恩慈、良善、信實、 \end{tabularx} \\ \\ \relax
5:23 & \begin{tabularx}{0.7\textwidth}{X} 溫柔、節制。這樣的事沒有律法禁止。 \end{tabularx} \\ \\
[1ex]
\hline
\hline
\end{longtable}


\newpage

\section{第48屆~港九培靈會~孫德生~第七講}
\label{sec:1111}
\textbf{聖靈與禱告}
\newline
\newline
連結: \href{https://www.hkbibleconference.org/session-message/view/1111}{\texttt{https://www.hkbibleconference.org/session-message/view/1111}}
\newline
\newline
\hyperref[sec:1110]{< < < PREV SERMON < < <}
~
\hyperref[sec:toc]{[返目錄]}
~
\hyperref[sec:1126]{> > > NEXT SERMON > > >}
\newline
\newline
路加福音 11:1-11:13
\newline
\begin{longtable}{cl}
\hline
\hline
章節 & 經文 (和合本修訂版)\\
\hline
11:1 & \begin{tabularx}{0.7\textwidth}{X} 耶穌在一個地方禱告。禱告完了，有個門徒對他說：「主啊，求你教導我們禱告，像約翰教導他的門徒一樣。」 \end{tabularx} \\ \\ \relax
11:2 & \begin{tabularx}{0.7\textwidth}{X} 耶穌對他們說：「你們禱告的時候，要說：『父啊，願人都尊你的名為聖；願你的國降臨； \end{tabularx} \\ \\ \relax
11:3 & \begin{tabularx}{0.7\textwidth}{X} 我們日用的飲食，天天賜給我們。 \end{tabularx} \\ \\ \relax
11:4 & \begin{tabularx}{0.7\textwidth}{X} 赦免我們的罪，因為我們也赦免凡虧欠我們的人。不叫我們陷入試探。』」 \end{tabularx} \\ \\ \relax
11:5 & \begin{tabularx}{0.7\textwidth}{X} 耶穌又對他們說：「你們中間誰有一個朋友半夜到他那裡去，對他說：『朋友！請借給我三個餅； \end{tabularx} \\ \\ \relax
11:6 & \begin{tabularx}{0.7\textwidth}{X} 因為我有一個朋友旅途中來到我這裡，我沒有東西招待他。』 \end{tabularx} \\ \\ \relax
11:7 & \begin{tabularx}{0.7\textwidth}{X} 那人在裡面回答：『不要打擾我，門已經關了，孩子們也同我在床上了，我不能起來給你。』 \end{tabularx} \\ \\ \relax
11:8 & \begin{tabularx}{0.7\textwidth}{X} 我告訴你們，雖不因他是朋友起來給他，也會因他不顧面子地直求，起來照他所需要的給他。 \end{tabularx} \\ \\ \relax
11:9 & \begin{tabularx}{0.7\textwidth}{X} 我又告訴你們，祈求，就給你們；尋找，就找到；叩門，就給你們開門。 \end{tabularx} \\ \\ \relax
11:10 & \begin{tabularx}{0.7\textwidth}{X} 因為凡祈求的，就得著；尋找的，就找到；叩門的，就給他開門。 \end{tabularx} \\ \\ \relax
11:11 & \begin{tabularx}{0.7\textwidth}{X} 你們中間作父親的，誰有兒子求魚，反拿蛇當魚給他呢？ \end{tabularx} \\ \\ \relax
11:12 & \begin{tabularx}{0.7\textwidth}{X} 求雞蛋，反給他蠍子呢？ \end{tabularx} \\ \\ \relax
11:13 & \begin{tabularx}{0.7\textwidth}{X} 你們雖然不好，尚且知道拿好東西給兒女，何況天父，他豈不更要把聖靈賜給求他的人嗎？」 \end{tabularx} \\ \\
[1ex]
\hline
\hline
\end{longtable}
禱告是基督徒生命中,最重要的一件事,禱告的生活非常強盛,才能成為一個健壯的基督徒.無論是傳道人,牧師,青年人以及任何基督徒,在禱告的生活中,也常會發生問題.為何我們的禱告沒有進步；因為我們不肯讓聖靈帶領我們,按照祂的方法去禱告.
\newline
\newline
主耶穌經常禱告,門徒都和祂在一起,當耶穌禱告與天父交通時,門徒心裡深切地渴慕.門徒常常聆聽耶穌講道,但他們從未要求耶穌教導他們講道；他們常見耶穌行神蹟奇事,他們也沒有懇求教導他們行神蹟奇事；惟他們看見耶穌禱告與天父交通,他們也求主教導他們禱告,他們並非求主教導他們怎樣禱告,乃是直接求主教導他們去禱告；教導他們與父有交通,怎樣在父前向祂傾心吐意,俯伏敬拜像耶穌一樣；教導他為別人代求,像耶穌為別人禱告一樣.但願今晚我們都向主呼求,求主教導我們禱告!
\newline
\newline
主耶穌差遣聖靈,來教導我們禱告　先知撒迦利亞稱聖靈為「那施恩叫人懇求的靈.」(亞十二10)聖靈是叫人懇求叫人禱告的聖靈.保羅說:「靠著聖靈隨時多方禱告祈求,並要在此儆醒不倦,為眾聖徒祈求.(弗六18)他題到聖靈如何幫助我們禱告:「況且我們的軟弱有聖靈幫助,我們本不曉得當怎樣禱告,只是聖靈親自用說不出來的歎息,替我們禱告.鑑察人心的,曉得勝靈的意思；因為聖靈照著神的旨意,替聖徒祈求.」(羅八26-27)你是否覺得禱告的生命非常軟弱,覺得禱告是件難事呢?這裡說神差遣聖靈,在我們軟弱時,幫助教導我們禱告,我們不曉得有何祈禱的內容；但神賜下聖靈來,顯明我們應當為何事祈禱,故此,我們每逢禱告必須倚賴聖靈的幫助.好比我們要和朋友通電話,必須雙方都有電話設備,更需要有電流流通,否則不發生效力；聖靈就如同電力,在我與之間作為聯絡.真奇怪!有時我們很不想禱告,這便魔鬼最快樂；有首詩歌:「軟弱信徒跪下禱告,撒見絕望震抖.」魔鬼千方百計攔阻我們禱告；但聖靈要幫助我們,叫我們渴慕禱告.我初但看信主時常為著自己向神禱告,這種自私的禱告實在無益.有一次我聽講道的人說:平衡之禱告必須有敬拜,感謝,認罪,祈求,(為自己禱告)代求,(為別人禱告)如果要得到完全的禱告生活,必須具有這五個部分在禱告之中.
\newline
\newline
「敬拜」　當我們敬拜時,我們就不能想到自己和自己的需要,也不必想到別人的需要；應當想到神,想到神的兒子耶穌基督,也想到聖靈；想到神的偉大和榮耀,想到主耶穌為我們在十字架上受死；我們的心向祂敬拜.門徒求主教導他們禱告,主教他們一個模範的禱告說:「你們禱告的時候,要說,我們在天上的父,願人都尊你的名為聖.願你的國降臨.願你的旨意行在地上,如同行在天上.」這是主禱文的上半段,充滿了神的榮耀和神所關心的事,沒有題到關於人本身的事.主耶穌以神的榮耀和神所關心的事,在禱告生活中居首位.我承認在禱告生活中,我們沒有讓神居首位,常以自我居首位,關心自己的需要.但是主耶穌教導我們要讓神居首位,禱告時首先要敬拜.
\newline
\newline
「感謝神」　新約的聖經中,禱告和感謝常是相題並論.要為我們所得蒙的福向神感謝.神經常不斷地,賜福給我們；但是我們卻習以為常,毫不在乎.保羅說:「當存感謝的心.」主耶穌醫治了十個長大痲瘋的人,只有一個回來感謝；耶穌很傷心地說:「你回來了,其他的九位呢?」我們常是如此,接受了祝福,卻忘記了感謝.
\newline
\newline
「認罪」　基督徒無時不需耶穌基督寶血的潔淨；可能不知不覺我們犯了罪,我們對罪的知覺遲鈍.我們若是親密與神同行,就愈感到自己虧欠了神的榮耀,未達到神的標準.有時我們上了魔鬼的勾當致落在罪惡中；所以在我們禱告生活中必須認罪,當我們發覺自己犯罪,應立即認罪,求主潔淨,切勿拖延!勿找尋理由為自己的罪辯護,托詞掩蓋；要將罪敞開在主面前.與主站在同一陣線,憎恨罪惡；主要幫助我們,使我們痛恨罪,與罪隔絕,不再沾染.如果得罪了神,只需向神認罪；若是得罪人,應該先向神認罪,然後再向人認罪,若得罪了教會,也應先向神認罪,再向教會認罪.
\newline
\newline
「祈求」　(為個人的需要)主禱文的下半段就是祈求,求神供給我目前之需要；日用的飲食,求神天天賜給.也為屬靈的需要向神祈求,赦免我的罪,保守我脫離試探,引誘,兇惡.
\newline
\newline
「代求」　為別人禱告應告應佔有禱告的重要地位,不要想到自己,只要顧及別人.有首詩歌說:「求主幫助我每天的生活,讓我忘記自己；禱告的時候,只為別人祈求.」
\newline
\newline
聖靈帶領我們到天父面前(弗二18)「因為我們上下藉著祂被一個聖靈所感,以進到天父面前.」我們禱告時,當仰望這位聖靈,祂能幫助我們頭腦的無知,幫助我們的軟弱,這是主所應許的.「因為聖靈照著神的旨意,替聖徒祈求.」除了在人裡頭的靈,誰知道人的事；像這樣,除了神的靈,也沒有人知道神的事.(林前二11)聖靈十分清楚知道神的旨意.「我們若照祂的旨意求什麼,祂就聽我們；這是我們向祂所存坦然無懼的心.既然知道祂聽我們一切所求的,就知道我們所求於祂的無不得著.(約壹五14-15).
\newline
\newline
聖靈用聖經教導我們,知道神的意旨,照神的意旨禱告.在聖經中一切關於道德的屬靈原則,都賜下給我們.關乎個性的生活的原則,都在聖經；如果合乎聖經,我們可以確知必然接受.可是在生活中,有很多實例聖經沒有題及；聖經中沒有結婚對象的名字,沒有你讀大學應修的科系,也沒有說明你的職業；但聖靈給你內在的確據和感動,聖靈向你發聲音.啟示錄一再的說「聖靈向教會所說的話,凡有耳的,就應聽.」耶穌說:「我的羊聽我的聲音.」我們若誠懇尋求神的旨意,聖靈必給我們確據,叫我們心裡明白神的旨意.在行傳十六章,保羅和西拉想將福音帶到亞細亞,但「聖靈禁止,」他們想往庇推尼,「耶穌的靈不許,」在這班神僕心裡有聖靈的聲音,他們就停下來；大家禱告時,聖靈對他們說話,眾人就認為這就是神的旨意.神的旨意向我們啟示的方法,是藉聖靈對我們說話,我們的心靈就充滿平安.「神所賜出人意外的平安,必在基督耶穌裡,保守你們的心懷意念.」(腓四7)你心裡若有平安,就是屬神的旨意,若沒有平安就必須等候.
\newline
\newline
聖靈另有一方法,教導我們明白神的旨意,祂替別人代禱的負擔,重壓在我們肩上；這是聖靈彰顯神旨意的方法.多年前我到中國大陸旅行,牛們騎馬到了貴州的中心；有人告訴我們該地常有強盜出沒,話還未說完,看見有一人死在那裡；我們非常驚恐,頓時草本皆兵；但是結果總算平安渡過.後來內子來信說,一天我覺得須為你祈禱；我就用很長的時間祈禱,直到我覺得禱告蒙神垂聽,你的危險已經渡過.」原來她為我禱告時,正是找們經過賊窩最危險之際.神的聖靈,將禱告的負擔,放在處紐西蘭的人；為著千萬里外,處身危險之中的人禱告；遙遠的路程並不影響禱告.
\newline
\newline
聖靈要幫助我們繼續禱告　我們有時頹喪不想禱告,聖經說:「要不住禱告.」保羅對帖撒羅尼迦的人說:「我不住為你們禱告.」所謂不住禱告,並非停止其他一切的工作；好像患咳嗽.雖不斷地咳嗽,間中都有停止的時候；我想保羅說的不住禱告,就是這個意思,當他頭腦不思想別事時,就為帖撒羅尼迦教會眾人禱告.各位!這一點大家都可做到,不住向神獻上禱告,正如聖靈不住在我們裡面工作一樣.保羅的生命中,自發自動地關心弟兄姊妹,不住地為他們禱告.
\newline
\newline
聖靈教導我們奉耶穌基督的名禱告　在約翰福音(十四13-14,十五16,十六23,24,26)耶穌六次題及奉祂自己的禱告.耶穌的名,即代表耶穌本身,奉祂的名即如耶穌一樣；何等偉大!比如我的銀行戶口,我將我的名給內子,她可以簽我名取款；我就是她,她就是我.我們和主耶穌的關係亦如此,我們禱告非奉自己的功德,是要奉主的聖名；意即我禱告的內容必須是主耶穌禱告的內容.故此,我們禱告時,務必認清是否屬於主耶穌的禱告.我的禱告,主能否同意簽名呢?這樣可以免去許多自私的禱告.我們奉主名禱告「叫父因兒子得榮耀.」(約十四23)
\newline
\newline
為何我們的禱告不蒙父垂聽?　我想我們的禱告常常不誠實,我們應該省察一下.聖經記載,一次有人帶了個被鬼附的孩子求門徒醫治,但門徒束手無策,他們問耶穌,「為什我們不能將那個鬼趕出去?」他們將自己的失敗坦白在主面前；主耶穌清楚告訴門徒失敗的原因是「你們不信我.」如果我們的禱告不蒙主垂聽,應當到主面前省察自己,聖靈要幫助我們記起聖經的話:「我若心裡注重罪孽,主必不聽.」(詩六十六18)這聖經並非說我已跌倒陷落罪中,不肯放棄罪惡的時候,我的禱告主必不聽；所以我們必須放棄罪惡,撫心自問何罪攔阻,以致禱告不蒙主應允.或者聖靈要帶領你思想,雅各書四章三節:「你們求也得不著,是因為你們妄求,要浪費在你們的宴樂中.」禱告的動機非常重要,必須正確,自問所求,是否為自私的享受呢?或是讓神的名得著榮耀,神的子民得福,神的教會被建立呢?或者聖靈要帶領你思想,希伯來書十一章6節:「人非有信,就不能得神的喜悅；因為到神面前來的人,必須信有神,且信祂賞賜那尋求祂的人.」聖靈向你指明,祈禱當存著信心.
\newline
\newline


\newpage

\section{第48屆~港九培靈會~孫德生~第八講}
\label{sec:1126}
\textbf{傳福音給萬民}
\newline
\newline
連結: \href{https://www.hkbibleconference.org/session-message/view/1126}{\texttt{https://www.hkbibleconference.org/session-message/view/1126}}
\newline
\newline
\hyperref[sec:1111]{< < < PREV SERMON < < <}
~
\hyperref[sec:toc]{[返目錄]}
~
\hyperref[sec:1128]{> > > NEXT SERMON > > >}
\newline
\newline
使徒行傳 11:19-11:30
\newline
\begin{longtable}{cl}
\hline
\hline
章節 & 經文 (和合本修訂版)\\
\hline
11:19 & \begin{tabularx}{0.7\textwidth}{X} 那些因司提反的事遭患難而四處分散的門徒，直走到腓尼基、塞浦路斯和安提阿。他們不向別人講道，只向猶太人講。 \end{tabularx} \\ \\ \relax
11:20 & \begin{tabularx}{0.7\textwidth}{X} 但內中有塞浦路斯和古利奈人，他們到了安提阿也向希臘人傳講主耶穌的福音。 \end{tabularx} \\ \\ \relax
11:21 & \begin{tabularx}{0.7\textwidth}{X} 主的手與他們同在，信而歸主的人數很多。 \end{tabularx} \\ \\ \relax
11:22 & \begin{tabularx}{0.7\textwidth}{X} 這風聲傳到耶路撒冷教會的人耳中，他們就打發巴拿巴到安提阿去。 \end{tabularx} \\ \\ \relax
11:23 & \begin{tabularx}{0.7\textwidth}{X} 他到了那裡，看見神所賜的恩就歡喜，勸勉眾人要立定心志，恆久靠主。 \end{tabularx} \\ \\ \relax
11:24 & \begin{tabularx}{0.7\textwidth}{X} 這巴拿巴原是個好人，滿有聖靈和信心，於是有許多人歸服了主。 \end{tabularx} \\ \\ \relax
11:25 & \begin{tabularx}{0.7\textwidth}{X} 他又往大數去找掃羅， \end{tabularx} \\ \\ \relax
11:26 & \begin{tabularx}{0.7\textwidth}{X} 找著了，就帶他到安提阿去。他們足有一年和教會一同聚集，教導了許多人。門徒稱為「基督徒」是從安提阿開始的。 \end{tabularx} \\ \\ \relax
11:27 & \begin{tabularx}{0.7\textwidth}{X} 當那些日子，有幾位先知從耶路撒冷下到安提阿。 \end{tabularx} \\ \\ \relax
11:28 & \begin{tabularx}{0.7\textwidth}{X} 內中有一位，名叫亞迦布，站起來，藉著聖靈指示普天下將有大饑荒；這事在克勞第年間果然實現了。 \end{tabularx} \\ \\ \relax
11:29 & \begin{tabularx}{0.7\textwidth}{X} 於是門徒決定，照各人的力量捐錢，送去供給住在猶太的弟兄。 \end{tabularx} \\ \\ \relax
11:30 & \begin{tabularx}{0.7\textwidth}{X} 他們就這樣做了，託巴拿巴和掃羅的手送到眾長老那裡。 \end{tabularx} \\ \\
[1ex]
\hline
\hline
\end{longtable}
如果神要賜福給人,也需要有人肯接受祂的福.要在行動上有所反應,這才是真正的福.你若承認已得了極大的屬靈福氣,那麼,我要問你:你將這福氣如何表達,才能把福氣輸送給別人?主耶穌說:「信我的人,從他腹中要流出活水的江河來.」聖經告訴我們,耶穌這話是指賜下聖靈而言.聖靈賜福給我們,並要我們把福氣輸送給人.耶穌基督是到神那裡的唯一道路,除此沒有人能到神那裡去.作基督徒的,應將福音告訴世人,否則他們的「血」就在我們的手上.
\newline
\newline
聖經由始至終是一部傳福音的書.新約聖經所有的作者都是傳教師,每一封書信都是作者用傳福音的方法,寫給他所建立的教會.使徒行傳是傳福音歷史的記載,記載了一個時代卅年久傳福音的歷史.神賜下使徒行傳作為例證,表明神藉使徒們所成就的大事.
\newline
\newline
初期教會所作的,被人誤為教會使天下變為繚亂的局面；事實上,初期教會卻是把繚亂的世局扭轉進入完整的地步.
\newline
\newline
耶穌基督復活之後深切關懷的,是福音能夠傳遍普天下.馬太福音記載耶穌差遣門徒住普天下去,使萬民作祂的門徒.馬可福音載:要把福音帶給每一個人.路加福音載:這福音必須傳遍萬民.約翰福音耶穌說:「父怎樣差遣我,我也照樣差遣你們.」使徒行傳一章八節是耶穌在世上最後囑咐的話:「要在耶路撒冷,猶太全地,和撒瑪利亞,直到地極,作我的見證.」耶穌從死復活,傳福音是祂心裡最大的負擔,祂最後說:「直到地極作我的見證.」耶穌基督十分願意每個人聽到福音,祂應許要賜給門徒能力.說:「天下地上的權柄都賜給我了」「看哪!我要與你們同在,直到世界的末了.」
\newline
\newline
使徒行傳十三章一節揭開了傳福音事工的第一頁,這節聖經首要的字句是「聖靈說,」耶穌基督把傳福音事工的責任交付聖靈呼召人作宣教士.在第一世紀有兩個深有影響力的教會,第一個當然是五旬節次日,在耶路撒冷所設立的教會；有段時期這教會經常得著復興.主將得救人數天天加給他們,但可惜復興也停頓在耶路撒冷；直到耶路撒冷受逼迫,信徒四散,他們才重新傳揚福音.
\newline
\newline
主耶穌說要向萬國萬民傳福音:「那些因司提反的事遭患難四散的門徒......他們不向別人講道,只向猶太人講.」耶路撒冷教會只向猶太人傳福音,這教會沒有與主同心,主要撇開這樣的教會.親愛的弟兄妹姊!無論國家,民族,教會,家庭,個人都具有愛自己過於愛別人的特徵,這完全與神的心意相背；神愛世人,神的愛,是普遍的愛；我們當學基督的樣式,甚至愛最卑微的人,關心普世的人.
\newline
\newline
多年前我曾收到一位,從事於福音廣播工作的日本牧師,巴斯哈托利來信:「我國基督徒太自私,只顧自己；耶穌基督要我們往普天下傳福音給萬民聽,盼望孫牧師來此講道,激起弟兄姊妹往普天下傳福音的熱誠.」我很願意有此機會事奉,就到日本數天領會,赴會者除基督徒之外,約有四百位傳道人.當大會結束那晚,主席邀請青年們凡願意往普世傳福音的有所表示,(據說往外傳福音在日本從未有人提及)當時約有百位青年弟兄姊妹(多數是大學生)願意當海外宣教士,直至如今,他們還在海外工作.
\newline
\newline
當耶路撒冷教會得悉安提阿教會需要幫助時,就差遣巴拿巴去幫助他們.巴拿巴是個好人,一位勸慰者.他鼓勵安提阿教會,於是工作日日增長,超過他所能勝任.他思索誰可幫助他,禱告之中,聖靈指示邀請大數的掃羅.(即後來的保羅)他們足有一年之久同工,工作蓬勃增長,成為有史以來最偉大的教會.安提阿教會使全世界教會聯結起來,工作開展初期,它就以全世界作為焦點,非單向猶太人傳福音；在他們領袖之中,有從居比路來的,有從大數,以東,古利奈(非洲)來的,教會中完全沒有種族歧視,不分國籍,不問身份,真是合一的教會!達到主耶穌的命令.安提阿教會工作方面是多姿多采的；教會中有先知和教師,有同工的隊伍；神所賜各有不同的恩賜,在教會中一齊發揮,被主使用.一般人以為事奉是作在會友身上,但安提阿這班領袖,他們所認定的,是事奉主,為首要,然後才服事弟兄姊妹；他們認真工作大家一齊禁食,當時聖靈對他們說:「要為我分派巴拿巴和掃羅,去作我召他們所作的工.」我們中間有教會領袖嗎?你是否讓聖靈向你發聲?你對聖靈的聲音有動於中嗎?
\newline
\newline
安提阿教勇往直前為主作見證,「門徒稱為基督徒,是從安提阿起首.」(十一26)他們的見證使人知道他們是屬於基督的.這教會不住增長,「主與他們同在,信而歸主的人就很多了.」(21節)「......教訓了許多人；」(26節)「他們既奉了差遣,就下安提阿去,聚集眾人,......」(十五30)據記載,初期教會在那城市有十萬基督徒,佔全城之半數.這是何等有效的見證!
\newline
\newline
安提阿教會在受託的事工上非常慷慨(十一28-30)當時有一位先知亞迦布的,指明天下將有饑荒；於是弟兄姊妹都很激動,收集獻金送去供給有需要的肢體,這是聖經首次題及,教會弟兄姊妹集捐項送給外地教會.安提阿的弟兄姊妹滿有神的愛心,他們各盡所能奉獻.
\newline
\newline
我們的主是位好的猶太人,祂遵守法律,按當時規定,要獻上十分之一.我們的主非但獻上十分之一,祂的一切都獻上給天父了.我們當學主的好榜樣,請問你有沒有把收入十分之一奉獻給主?如果你沒有在奉獻上忠心,你將失去蒙福的機會.信徒忠心的奉獻,不但教會和海外傳道的需要豐足,而本身生命豐滿是意想不到的.數年前我到馬尼拉,有一個教會發動海外傳福音工作；菲律濱本是傳福音的工場,他們一向沒想到向外傳福音,在此之兩年,他們積蓄基金作為海外傳道之用.兩年中差派了六個青年出去做宣教士,積蓄了二萬六千菲幣供給他們.過了一週,我到泰國去,說出這事,會後即有一對年青夫婦來和我握手,說他倆就是由菲律濱差派來的.在曼谷有間小小的華人教會,建立僅兩年；為海外傳奉獻了七萬三千美元,這教會在管家的職份上慷慨.
\newline
\newline
安提阿教會存著禱告的心來計劃同工們禁食禱告,奉聖靈差派到各地去傳福音.安提阿教會成為全世界傳福音的中心,香港的差傳工作根源,也是出於當時安提阿教會的禱告會.最令人鼓舞的,所謂第三世界的國家──亞洲,非洲,拉丁美洲；聖靈給予傳福音的負擔,從這些國家差派出去的傳教士共有三千四百人；在西方國家徵募宣教士日走下坡之際,其他地方徵募之舉卻漸抬頭.數年來韓國教會,以二千宣教士為目標將差遣出去；日本也正計劃要加派一千位,還有許多地方現在都有差傳的組織:香港也有這樣的工作,但不敷應用.
\newline
\newline
數年前,新加坡的大學生,請我主領退修會研經會.參加人數約五六十人,每晚領他們研經.到了最後一晚,他們要求以傳福音為主題.據他們說,從未有人勉勵他們向海外傳福音,一向都被認為他們所處的是福音的工場.他們很想知道在世界各地的工作,當晚我就將這幅傳福音的圖畫指示他們,會後有十五位青年來對我說,他們渴望當海外傳教師,應如何進行?我就給他們一些建議,之後一週內,他們又來找我；因他們想更多知道如何預備作個教士,我囑他們安排另一退修會,每日有十堂準宣教士的聚會.結束之日,他們都有一種新的感受,他們從未想到將作西方人的宣教士.
\newline
\newline
神不是叫我們個個都往海外傳道,但神要我們在自己的崗位上,為主作見證.無論年輕或年長的,都當將自己交在主手中,甘心樂意地事奉祂.我在二十歲時奉獻作海外傳教士,到了五十二歲才進入海外福音工場.七十歲那年主又領我轉入新畿內亞,在那裡度過兩年美好的時光.我不知道前面的路途如何,但我隨時作好準備候主差遣.數週前,我們在紐西蘭歡送兩位年約五十歲的；其中一位是某間相當大的中學校長,一位是市政局文員,他們蒙主呼召離開了世俗的工作,現在一位在印度,一位在尼泊爾傳福音.
\newline
\newline
作為基督徒都同樣有責任向人傳福音,耶穌說,你們要往普天下傳福音給萬民聽.這命令是向所有的信徒發出的,主並非要我們個個遠涉重洋,乃是期望我們肩負傳福音的使命,叫人認識主耶穌.
\newline
\newline
耶穌說:「我是道路,若不藉著我,沒有人能到父那裡去.」人若沒有聽過耶穌的名,不認識耶穌,怎能到父那裡去呢?耶穌說,人若不重生,就不能見神的國.彼得說除了耶穌以外,沒有別的名,我們可以靠著得救.我並非說未聽過福音者就必失喪,犯罪不信耶穌者才是失喪的人.但你我因聽到福音有了耶穌,祂將我們的生命從失敗死亡的邊緣扭轉過來.我們既得了這樣的經驗,因得救恩而喜樂,我們對那沒有指望的人是否關心呢?
\newline
\newline
我們的差會在菲律濱海島工作時,第一個歸信主的是個年高體弱的祖母,她一聽到福音即接受耶穌作救主.她非常愛主,當她受浸時,牧師問她,「你是否相信耶穌作你救主了」她回答說:「當然相信,如果你較早來傳福音,我早就信了.」很多人在沒有指望之中離世了,乃因基督徒不關心他們,沒有將救恩的福音告訴他們.你是否經常為海外傳道事工代禱?你所關心的是否只限於本國或本地的傳福音工作?
\newline
\newline
有一位偉大的皇室外科醫生,有一天他乘火車出外旅行,途中火車肇事,許多人受了傷；那醫生卻沒有受傷.他跳出火車,看見司機躺在那裡,他連忙進行急救,發現他傷勢嚴重,旁的人聽見醫生用低沉的語調說:「只要我有外科醫療器材,才能挽回他的生命.」我們的主從天上看見,世界三分之二的人口,沒有機會聽見得救的福音；我們的主救恩經已齊備,只期待器材為祂效忠.我的弟兄姊妹!你願意作主的器皿嗎?你願意毫無保留將生命交在主的手中嗎?或者你認為自己現在沒有任何好處可以獻給主；但是可將你整個生命誠心交託主,安靜等候,任主使用!同時,要在自己的崗位上為主作美好的見證.
\newline
\newline


\newpage

\section{第48屆~港九培靈會~孫德生~第九講}
\label{sec:1128}
\textbf{迦勒專心跟從耶和華}
\newline
\newline
連結: \href{https://www.hkbibleconference.org/session-message/view/1128}{\texttt{https://www.hkbibleconference.org/session-message/view/1128}}
\newline
\newline
\hyperref[sec:1126]{< < < PREV SERMON < < <}
~
\hyperref[sec:toc]{[返目錄]}
~
\hyperref[sec:1129]{> > > NEXT SERMON > > >}
\newline
\newline
約書亞記 14:6-14:14
\newline
\begin{longtable}{cl}
\hline
\hline
章節 & 經文 (和合本修訂版)\\
\hline
14:6 & \begin{tabularx}{0.7\textwidth}{X} 猶大人來到吉甲，約書亞那裡，基尼洗族耶孚尼的兒子迦勒對約書亞說：「耶和華在加低斯‧巴尼亞指著我和你對神人摩西所說的話，你都知道。 \end{tabularx} \\ \\ \relax
14:7 & \begin{tabularx}{0.7\textwidth}{X} 耶和華的僕人摩西從加低斯‧巴尼亞差派我窺探這地的時候，我剛四十歲。我把心裡的話向他報告。 \end{tabularx} \\ \\ \relax
14:8 & \begin{tabularx}{0.7\textwidth}{X} 雖然同我上去的眾弟兄使百姓膽戰心驚，我仍然專心跟從耶和華－我的神。 \end{tabularx} \\ \\ \relax
14:9 & \begin{tabularx}{0.7\textwidth}{X} 那日，摩西起誓說：『你腳所踏之地必要歸你和你的子孫永遠為業，因為你專心跟從耶和華－我的神。』 \end{tabularx} \\ \\ \relax
14:10 & \begin{tabularx}{0.7\textwidth}{X} 現在，看哪，耶和華照他所說的使我活了這四十五年。當以色列人在曠野飄流的時候，耶和華曾對摩西說了這話。現在，看哪，我已經八十五歲了。 \end{tabularx} \\ \\ \relax
14:11 & \begin{tabularx}{0.7\textwidth}{X} 現今我還很健壯，像摩西差派我去的那天一樣；無論是戰爭，是出入，我現在的力量和那時的力量一樣。 \end{tabularx} \\ \\ \relax
14:12 & \begin{tabularx}{0.7\textwidth}{X} 請你將耶和華那日所說的這山區給我。那日你也曾聽說，這裡有亞衲族人，以及寬大堅固的城，或許耶和華會照他所說的與我同在，我就把他們趕出去。」 \end{tabularx} \\ \\ \relax
14:13 & \begin{tabularx}{0.7\textwidth}{X} 於是約書亞為耶孚尼的兒子迦勒祝福，把希伯崙給他為業。 \end{tabularx} \\ \\ \relax
14:14 & \begin{tabularx}{0.7\textwidth}{X} 所以希伯崙成了基尼洗族耶孚尼的兒子迦勒的產業，直到今日，因為他專心跟從耶和華－以色列的神。 \end{tabularx} \\ \\
[1ex]
\hline
\hline
\end{longtable}
人生在世,都是隨著年齡逐漸衰退；但迦勒在世的日子從不停止長進,他卻日益強壯,年日消逝並不使他屬靈的生命衰退.迦勒此名是誠信忠實的意思,他確名符其實；他的一生足以為青年人,中年人及老年人的榜樣.
\newline
\newline
迦勒的生平可分為三部分,聖經對於他首四十年的記載不多,第一次的記載是他被差遣窺探迦南地.第二段的四十年,他在曠裡飄蕩；因以色列人不肯進入應許之地.最後的階段,我們不知有多久；但我們知道他得到生平最偉大的勝利.從他出生至四十歲,我們稱他為青年期.他被揀選作為窺探迦南地的十二個探子之一.至於他早年生活,聖經并沒有題及；正如耶穌的早年,聖經也很少題到一樣.耶穌十二歲時所作的事,聖經僅用兩節記載；但是我們從祂最後的三年半,可以見到祂前卅年的生活.
\newline
\newline
聖經告訴我們,被揀選窺探迦南的探子,都是從領袖中選拔出來的.通常作領袖的,都須經嚴密的紀律訓練；由此可知迦勒在年青時,就是個守紀律的人.
\newline
\newline
一,迦勒大有勇敢　當迦勒和約書亞窺探迦南地回來,他們報告該地實況；而其他十個探子報告壞的消息,反要用石頭將他們二人打死.迦勒和約書亞卻大有勇敢,不但有匹夫之勇,且有道德上的勇氣；道德之勇較匹夫之勇更難能可貴,他們站穩立場,在群眾反對的民族面前,果敢力排眾議,卓然獨立,真難得!當人極力反對基督,而你卻堅立不移,實在很難!神的眼中看此為寶貴,必得神的獎賞.
\newline
\newline
二,迦勒大有信心　信心非常重要,聖經說:「人非有信,就不能得神的喜悅,」(來十一6)迦勒和約書亞是大有信心的人,當時所面臨的境遇,更顯出他們信心的偉大和可貴.那十個探子窺探之後的報告,是可怕並令人悲觀的；他們的心充滿不信,說「那地的民強壯,域邑也堅固寬大.」(民十三28)「我們不能上去攻擊那民,因為他們比我們強壯......我們所窺探之地,是吞喫居民之地......人民都身量高大......亞衲族人,就是偉人；我們看自己就如蚱蜢一樣,據他們看我們也是如此.」(31-33)他們的報告,簡直使人心寒膽顫；但「迦勒在摩西前安撫百姓,說我們立刻上去得那地罷,我們足能得勝.」(30節)「約書亞和迦勒撕裂衣服,對以色列全會眾說,我們所窺探經過之地,是極美之地.耶和華若喜悅我們,就必將我們領進那地；把地賜給我們,那地原是流奶與蜜之地.但你們不可背叛耶和華,也不要怕那地的居民,因為他們是我們的食物......有耶和華與我們同在,不要怕他們.」(民十四6-9)這是何其大的對比!同一個迦南地的人,事,物,報告卻大不相同.十個探子看亞衲族人又高又大,把自己比作蚱蜢；迦勒卻把亞衲族人,和全能的神對比.十個探子所看見的是巨大的人,迦勒所看見的是偉大的神,他將那地居民當作食物.眾人聽信十個探子的報告,膽怯埋怨,彼此說:「我們不如回埃及去罷.」他們忘記在埃及為奴被鞭之苦,但迦勒永不回頭；他認定有耶和華同在,不要怕,要勇住直前.青年弟兄姊妹!你是否有迦勒般的信心和勇氣,為基督力排眾議,面對巨人,深信神能領你勝過呢?
\newline
\newline
迦勒年青時代奠定了好的根基,進入中年時代的工作.中年人有其獨特的考驗,或者不如青年的考驗猛烈；但詭詐的程度更高,牽涉的範圍也廣.成家立室之後,許多人常因事業和家庭的牽累,以及魔鬼的引誘；對主的熱誠幾乎都失去了.靈修沒有經常化了,事奉也敷衍形式化了；不像從前是因著被主愛激勵而事奉.到了中年有個自然的傾向,就是「放鬆一點」似乎有種特權可以稍微鬆懈,不再勇往直前事奉主；卻坐在安樂椅上電視機旁怡然自樂.有時候,人到了中年便開始失去自己的理想,緬懷青年時理想沒有實現；今日世事無可避免得過且過罷了.何西阿書先知的預言,論到以法連頭髮灰白而不自知.頭髮灰白表明人已衰老,一般人初發現自己有白髮,便嘆老之將至.有人還白髮染黑,不願意露出白髮於人；但如果屬靈方面出現白髮,卻無法可染,惟澈底對付連根拔掉.各位弟兄姊妹!你現在還如年青時一樣地愛主,或已失去了起初的愛心呢?你必須拔除白髮；向主認罪,求主賜給你起初愛祂的心.
\newline
\newline
迦勒滿有榮耀經過年青的考驗.迦勒和約書亞對主絕對真誠；主叫他們在曠野上來回,度過他們的中年.那四十年如同一隊很長的送殯行列,他們眼見同年代者一個又一個死了,直到最後剩下他們二人.雖然如此,迦勒並不灰心,而且這時正是也精力最充沛的時候,他精猛如昔,為主奔路；到了他人生第二個四十年將要結束時,他的信心如同年青時一樣的強壯.精神同樣的活潑,他預備接受神給他下一個階段的考驗.他年青時正如鷹展翅上騰了到了中年,他奔跑不知疲倦；他年青時熱心,中年時仍貫澈如一；到了老年他還有冒險勇敢的精神.聖經裡題到的老年人,惟迦勒給我們最大的鼓勵.
\newline
\newline
迦勒是個精明果敢,並滿有成就的老年人；他在八十五歲時,面臨最大的挑戰.一般而論,最遲六十五歲已告退休；但迦勒到了八十五歲,聖經並沒有記載他退休,或稍微放鬆工作；他對於老年與眾有不同的看法.他認為老年人並非日走下坡,反而向著更高之處進發.他從未失去屬靈的雄心大志,年老的階段,是他人生最好的年日；也是他對以色列人,作出最大的貢獻之時.
\newline
\newline
迦勒整個人生的每一階段,都是美好的榜樣；他年青時單獨站,中年時單獨行走,老年時單獨往上爬.神對迦勒應許說:「惟獨我的僕人迦勒,因他另有一個心志,專一跟從我；我就把他領進他所去過的那地,他的後裔也得那地為業.」這個應許,在迦勒生平佔有極重要的地位.經過了四十五年,迦勒才提起此事；他的記憶力沒有衰退,他曾五次提起神的應許.四十年來他仰望神的應許,這應許幫助他經過曠野的日子；雖然那日子是可怕的,但他深信神的應許必然成就.
\newline
\newline
迦勒八十五歲時,他說:「我還是強壯,像摩西打發我去的那天一樣；無論是爭戰,是出入,我的力量那時如何,現在還是如何.」(書十四11)八十五高齡的迦勒,不是穿著拖鞋坐在安樂椅上的老太爺；卻是穿著鐵鞋上山爭戰的勇士.迦勒的勝利不是身體的勝利,乃是靈裡的勝利；身體遲早會衰退軟弱,最重要是內裡的靈.許多身體軟弱的人更有屬靈的強壯,比身體強大的神更多成就.
\newline
\newline
迦勒所關心的惟神的勝利　當時約書亞將神應許之地,分給以色列各支派.迦勒領取應得之分時,他非以其年老為理由,要求平原美地,他只要求高山.　你面前有沒有高山?巴不得你今天求神賜給你高山,給你能夠克服如高山般的罪惡；是否有些尚未蒙神允的禱告,成為你的高山；多年來你向神所求的,至今尚未得著,你須求神把高山賜給你.你家庭中有無高山?讓你今日在神前支取高山,在你的事奉上,工作中,是否有高山使你不能爬越過去?哦!神啊!求你把這座高山賜給我.迦勒所要求的高山是希伯崙,最敵人最堅固之城邑,土地肥沃,風景絕佳；他除了神之最好的,不肯以次好的代替.很多時候,我們只求好的而已,沒有求最好的；如果你將最好的獻給主,主也將以最好的賜給你.迦勒年老時出去爭戰竟然得勝,他征服了應許之地,將亞衲族三個族長都驅逐出去；聖經說其他的人不能將敵人全部逐去,惟獨迦勒是能得的.這是給年長的弟兄姊妹很大的鼓勵.
\newline
\newline
我們差會有位西國傳教士威爾理,他在中國四川工作至七十歲,退休後到加拿大.他覺得無所事事,就開始學習希臘文；因他年青時沒有機會讀希臘文的新約聖經,後來他果然學成了.到他九十歲時,他在多倫多浸信會神學院和神學生們在一起,他繼續學習；到了他百歲時,我在多倫多講道,他遠路來赴會.禱告時,他非常踴躍首先站起來領禱,會後他來和我傾談,提起一些從前中國信徒的名字.我發覺他衣袋裡有本破舊的希臘文字典,原來是他用來在乘火車途中參閱的,到了百歲,仍然孜孜不倦.
\newline
\newline
親愛的朋友!你正讀些什麼?研究些什麼?你是否將你的頭腦經常保持活潑新鮮?是否將你思想最好的部分獻給主,或是讓你的頭腦擱置生草呢?你若能讀些查經課程,你屬靈的生命就能繼續增長,能有更好的東西獻給主.
\newline
\newline
有一次我在多倫多講道,大約過了四年,我再到該地,有個老姊妹來對我說:「四年前聽到你講關於迦勒的那篇道,當晚神對我說話,我想我雖已七十歲,但身體康健,生活沒有負擔,經濟也能獨立；難道神沒有工作叫我做嗎?我從迦勒身上學了功課,所以我到巴西作宣教工作,三年來實在是我美好的日子.」這事離今已有十年了.去年我到新畿內亞講道,我又講到迦勒,並提起這位老姊妹的經驗.原來他們也認識她,因她曾到此當中學教師直到她八十歲；她是教師中最受歡迎的一位,她現已八十二歲了,身體還很健康.年長的弟兄姊妹!何必思慮你的人生將渡盡了,可能神要將更好的年日賜給你.迦勒豈不是八十五歲之後才勝利完成他生命中最大工作嗎?在澳洲我有位朋友,她一百零九歲才離世,當她百歲時,她和兩個女兒,(一個七十歲,一個八十歲.)她們到另一個沒有教會的住宅工作,開始家庭禱告會,請鄰舍參加.兩年前我經新加坡時,那八十歲的女兒告訴我；她們的家庭禱告會已經成立一間教會；現有會友二百人聚會,且教會差派了廿位青年往海外傳福音.還有一位美國哥倫比亞聖經大學的學生,她畢業後從事福音錄音工作,用多種語言的錄音帶傳福音.
\newline
\newline
我的弟兄姊妹!只要你把生命奉獻給主,任讓主來使用你吧!
\newline
\newline
迦勒人生的秘訣,叫他從青年至老年一直平穩長進給主使用；他的秘訣並非深奧,而是很簡單的,只要肯付代價.他的秘訣聖經題到三次,第一次迦勒自己說:「我專心跟從耶和華我的神.」(書十四8)第二次是摩西說的:「因為你專心跟從耶和華我的神.」(9節)除此還有更好的見證,神說:「惟獨我的僕人迦勒,因他另有一個心志,專一跟從我.」(民十四24)這是很簡單的秘訣,因為迦勒專心跟從主,所以全面撲滅仇敵,榮耀的得勝.聖經記載,自此那地戰爭平息,大享安寧.
\newline
\newline
當你專心跟從主的時候,你心裡的內戰即將停止,得著神賜給你全然的勝利.
\newline
\newline
巴不得我們都向神禱告說:「神啊!求你將這座山賜給我,將我的生命拿去,幫助我像迦勒專心跟從你.」
\newline
\newline


\newpage

\section{第48屆~港九培靈會~孫德生~第十講}
\label{sec:1129}
\textbf{高山與深淵}
\newline
\newline
連結: \href{https://www.hkbibleconference.org/session-message/view/1129}{\texttt{https://www.hkbibleconference.org/session-message/view/1129}}
\newline
\newline
\hyperref[sec:1128]{< < < PREV SERMON < < <}
~
\hyperref[sec:toc]{[返目錄]}
~
\hyperref[sec:1051]{> > > NEXT SERMON > > >}
\newline
\newline
創世記 45:1-45:15
\newline
\begin{longtable}{cl}
\hline
\hline
章節 & 經文 (和合本修訂版)\\
\hline
45:1 & \begin{tabularx}{0.7\textwidth}{X} 約瑟在所有侍立在他旁邊的人面前情不自禁，就喊叫說：「每一個人都離開我，出去吧！」約瑟和兄弟相認的時候沒有一人站在他那裡。 \end{tabularx} \\ \\ \relax
45:2 & \begin{tabularx}{0.7\textwidth}{X} 他放聲大哭，埃及人聽見了，法老家中的人也聽見了。 \end{tabularx} \\ \\ \relax
45:3 & \begin{tabularx}{0.7\textwidth}{X} 約瑟對他兄弟們說：「我就是約瑟。我的父親還在嗎？」他兄弟們不敢回答他，因為他們在他面前都很驚惶。 \end{tabularx} \\ \\ \relax
45:4 & \begin{tabularx}{0.7\textwidth}{X} 約瑟又對他兄弟們說：「靠近我一點。」他們就近前來。他說：「我是被你們賣到埃及的兄弟約瑟。 \end{tabularx} \\ \\ \relax
45:5 & \begin{tabularx}{0.7\textwidth}{X} 現在，不要因為把我賣到這裡而憂傷，對自己生氣，因為神差我在你們以先來，為要保全性命。 \end{tabularx} \\ \\ \relax
45:6 & \begin{tabularx}{0.7\textwidth}{X} 現在這地的饑荒已經二年了，還有五年不能耕種，沒有收成。 \end{tabularx} \\ \\ \relax
45:7 & \begin{tabularx}{0.7\textwidth}{X} 神差我在你們以先來，為要給你們在世上存留餘種，大施拯救，保全你們的性命。 \end{tabularx} \\ \\ \relax
45:8 & \begin{tabularx}{0.7\textwidth}{X} 這樣看來，差我到這裡來的不是你們，而是神。他又使我如同法老之父，作他全家之主，和埃及全地掌權的人。 \end{tabularx} \\ \\ \relax
45:9 & \begin{tabularx}{0.7\textwidth}{X} 你們要趕緊上到我父親那裡，對他說：『你兒子約瑟這樣說：神已立我作全埃及之主，請你下到我這裡來，不要耽擱。 \end{tabularx} \\ \\ \relax
45:10 & \begin{tabularx}{0.7\textwidth}{X} 你和你的兒子孫子，羊群牛群，以及一切所有的，都可以住在歌珊地，與我相近。 \end{tabularx} \\ \\ \relax
45:11 & \begin{tabularx}{0.7\textwidth}{X} 我要在那裡奉養你，因為還有五年的饑荒，免得你和你的家屬，以及一切所有的，都陷入窮困中。』 \end{tabularx} \\ \\ \relax
45:12 & \begin{tabularx}{0.7\textwidth}{X} 看哪，你們的眼睛和我弟弟便雅憫的眼睛都看見，是我親口對你們說話。 \end{tabularx} \\ \\ \relax
45:13 & \begin{tabularx}{0.7\textwidth}{X} 你們要把我在埃及一切的尊榮和你們所有看見的事情都告訴我父親，也要趕緊請我父親下到這裡來。」 \end{tabularx} \\ \\ \relax
45:14 & \begin{tabularx}{0.7\textwidth}{X} 於是約瑟伏在他弟弟便雅憫的頸項上哭，便雅憫也在他的頸項上哭。 \end{tabularx} \\ \\ \relax
45:15 & \begin{tabularx}{0.7\textwidth}{X} 他又親眾兄弟，伏著他們哭。過後，他的兄弟就和他說話。 \end{tabularx} \\ \\
[1ex]
\hline
\hline
\end{longtable}
當你讀經覺得乏味時,我提議你查考聖經人物的事蹟,把有關的經節抄下來；你可以從中學習功課,得著鼓勵.
\newline
\newline
約瑟年少俊美；聖形容他秀雅俊美.他得父親偏愛,引致他哥哥們對他因妒成恨.
\newline
\newline
約瑟一生,「夢」佔了很重要的部分.約瑟所發的兩個夢,大不受他哥哥歡迎.約瑟對他弟兄們說:我夢見:「在田間捆禾稼,我的捆起來站著,你們的捆起來圍著我的捆下拜.(創三七7)「後來他又作了一夢,也告訴他的哥哥們說:「看哪!我又作一夢,夢見太陽,月亮,與十一個星,向我下拜.」(9節)這次約瑟的父親聽了,「就責備他說:你作的這是什麼夢,難道我和你母親,你弟兄果然要來俯伏在地向你下拜嗎?」約瑟的夢使他哥哥們更加憎恨他.有一天他奉父親之命,窺探哥哥們的行為；可能他們作賊心虛,恐約瑟報告實情,所以不許他回家,把他擲下野地的坑裡.
\newline
\newline
約瑟一生可分為三個階段:
\newline
\newline
一,約瑟受逆境的考驗　他早年過著波浪式的生活,起伏不平,忽低忽高；時而抓起希望,忽又陷入失望之中；甚而淪落監牢,人竟將他遺忘了.他從小受父親寵愛,想不到被擲下枯坑,頓時孤單又飢渴,似乎毫無指望；後來哥哥們把他拉上來,他以為可以回家,那知卻把他賣給米甸的以實瑪利人.因此,素來驕生慣養的約瑟開始過為奴生活了.他又被轉賣到法老的護衛長波提乏家中.波提乏高居官職,權勢浩大,波提乏的妻,可算是埃及國的第二夫人了.主人覺得他是個非凡青年,不但秀雅俊美,聰明過人,且忠心耿耿；於是約瑟的新希望又萌芽了,主人使他專權管理家政,深得信任；似乎樣樣盡都順利,他想釋放有期了.誰知重大的考驗突然臨到,他抗拒波提乏妻子的引誘,致被憎恨拘捕入獄.雖是如此,約瑟在獄中仍然對神盡忠,也忠於職守；不久,他成為監犯中最令人喜愛者.忠於神者,始終要居優越位置,現在約瑟再次升高了；他管理所有獄中囚犯,獄卒非常信任他.約瑟在獄中,先後為法老的酒政和膳長解夢,都一一應驗；膳長被殺,酒政得釋放.約瑟對酒政說:「但你得好處的時候,求你記念我,施恩與我,在法老面前題說我,救我出這監牢(四十14)酒政滿口答應,但卻忘記了.約瑟日夜等候,一直過了兩年仍在獄中.
\newline
\newline
約瑟忠於也忠於職守,並抗拒引誘；但他的遭遇,卻如此惡劣.一般人,定必唉聲嘆氣,灰心失志,以為善無善報,不如回到罪惡中去.但約瑟並不失去對神的信靠,他沒有埋怨；而是埋頭苦幹,力爭上游.雖一再失望,他卻始終忍耐；在他四圍與他相處,都屬不信神的人,他卻對神堅信不移.
\newline
\newline
約瑟少年時代,在逆境中所學習的功課,也是我們應學習的.有時神要我們升高之前降卑；我們的主登上寶座之前,曾經七次往下降.祂「有人的樣子,就自己卑微,存心順服,以至死,且死在十字架上.所以神將祂升為至高,賜給祂那超乎萬名之上的名.叫一切在天上的,地上的,和地底下的,因耶穌的名,無不屈膝.(腓二8-10)你現在是否正遭遇艱難逆境呢?是否正灰心絕望?請記得!當約瑟被人忘記,困在獄中,似乎他一生的命運已成定局；但神卻由此情此境中給他高升.
\newline
\newline
二,約瑟被引誘試探的考驗　創世記卅九章記載波提乏的妻子引誘約瑟的經過；這是我們勝過試探的典型榜樣.「性」的方面是人最大的引誘.其實,性慾本身並無錯誤,是神給人的禮物；若能正當使用,就成為莫大的福氣.若在神的旨意之外濫用,便成為極大的咒詛.今日的世界充滿情慾,從四面八方來襲擊神的兒女；所以我們要特加注意,約瑟如何勝過這樣的試探.
\newline
\newline
「引誘」在舊約原文有兩不同的字,意思亦異；一為考驗,舊約時代神試驗亞伯拉罕；但聖經記載神從不引誘使人犯罪.從字眼看,靈經說神引誘亞伯拉罕,其實則神考驗亞伯拉罕.金屬溶於爐火提煉,除去雜質,得到純淨；這就是第一個字的意思,神將考驗放進我們的生命中.考驗的第二個意思是誘惑人犯罪；魔鬼有六千年引誘人的經驗,牠先尋找弱點然後進攻,只要有隙可乘,便進來引誘我們犯罪.有時對你成為引誘,對我卻不成引誘,而是考驗耳.撒旦引誘我們犯罪企圖摧毀我們；神卻考驗我們,為使我們更加強壯更加潔淨.在約瑟的一生,魔鬼的引誘和神的考驗同時並行,魔鬼千方百計要摧毀他,神卻考驗他使他更堅強為神站立.當魔鬼要摧毀你時,主在你身旁,祂要叫你經過考驗,得到更好的結果.
\newline
\newline
約瑟經過試探時,他是個廿七歲的青年,正是血氣方剛,熱血奔騰之時,突如其來的試探,確難保全他的貞潔.相反地,很容易引他陷入罪中；當時他遠離家庭,如果犯罪,家人無從知悉.靈性方面也沒有親人,或長者可幫助他；同時波提乏妻子,又盡一切引誘的能事；試!約瑟此時所處是何等有利於誘惑的環境!那一天早晨,約瑟如常進行管家的工作,那想到一生的成敗繫此一霎,誰能逆料一國之第二夫人,竟誘惑個為奴隸的青年.試探簡直不容人防範,撒旦的突擊,常使人不及穿戴屬靈的軍裝.所以我們應當時加儆惕,隨時穿上屬靈軍裝,與主同行,不斷與主有交通,儆醒禱告；否則,試探來到,便無法抵擋.另一面,約瑟面對的試探,不僅一次而己,波提乏的妻子日日不斷來引誘他,不正當的思想很容易在他腦中扎更新長；頭腦如同戰場,思想實際的爭戰在頭腦裡.約瑟始終不容這種思想在自己腦中有餘地,他說:「我怎能作這大惡,得罪神呢?」(卅九9)他對罪有與神同感的態度.各位!我們遇見罪惡就戰兢,這對我們有很大的保險.約瑟認犯罪不是得罪人,乃是得罪神；因為怕得罪神,保守他離開罪惡.聖經說:當波提乏家中沒有一個人在那屋裡,假如約瑟犯罪,也無第三者知道(除神之外)而且對他很有利.可是這位青年人過慣規律的生活,思想受到約束；不容許頭腦考慮邪惡,隨時作好防範準備.各位!你是否作好準備抗拒罪惡?約瑟的經歷有許多寶貴的教訓,尤其是青年弟兄姊妹.
\newline
\newline
「情慾」的本身並非罪,神這樣創造我們,這事可以成為純淨美好,試探的本身也非罪,主耶穌「祂也曾凡事受過試探,與我們一樣,只是祂沒有犯罪.」(來四15)行神律法之外的,就是犯罪.神預備了婚姻,作為人滿足性慾的方法,若在神命定的方法之外,去尋求滿足就是罪.甚至在婚姻之內,過分使用也是罪.
\newline
\newline
約瑟起來逃命,戰勝這場罪惡；他急於奔命以致衣裳丟在婦人手裡.抗拒撒旦最好的方法,就是毅然決然背向而馳.
\newline
\newline
約瑟與大衛有著強烈的對比約瑟經考驗得到榮耀的勝利,經此考驗更加堅強,更加潔淨；而犯罪的經驗使大衛成一個軟弱的人.聖經說,作為以色列王不可多妻；而大衛有許多妻妾,當他犯罪時,他非常懿惰.在他應該領軍作戰的那段時間,他卻睡到日頭平西才起床.他放縱自己,忽略職責,致使魔鬼有機可乘來摧殘他.各位弟兄姊妹!我們不可離自己的崗位,要儆醒戒備以防引誘.
\newline
\newline
約瑟得到勝利有何賞賜?是否得到榮耀的職位?他所得到的,乃是下到監牢；於是將約瑟帶進第三個考驗.
\newline
\newline
三,富貴環境的考驗　法老作夢,全國的術士和博士,沒有人能圓解；酒政忽想起約瑟來,將他帶進皇宮.法老對約瑟說:「我聽見人說,你聽了夢就能解.」約瑟回答法老說,「這不在乎我,神必將平安的話回答法老.」(四一15-16)約瑟與神同行,他認識神,信靠神；他為法老解夢,並指出埃及遍地,將來必有可怕的荒年.神將計劃指示約瑟,讓他向法老建議,如何挽救將來的荒年；建議法老當揀選一個,聰明有智慧的人管理埃及地.「法老對臣僕說,像約瑟這樣的人,有神的靈在他裡頭；我們豈能找得著呢?法老對約瑟說,神既將這事都指示你,可見沒有人像這樣有聰明有智慧.你可以管我我家,我民都必聽必聽從你的話,惟獨在寶座上我比你大......我派你治理埃及全地.法老就摘下手上打印的戒指,戴約瑟手上；給他穿上細麻衣,把金鍊戴在他的頸項上.」(四一38-42)約瑟由奴隸一躍而為埃及國一人之下,萬人之上的宰相,榮華富貴皆備.他原有奴隸的衣裳,留在波提乏妻子手上；現在換上了華麗高貴的細麻衣.
\newline
\newline
約瑟好的品德,是否得到善報?他十三年的逆境考驗值得嗎?我們常不明白神的作為,正如約瑟不明白他的一生.他似乎覺得所發生的事,對他的人生毫無意義；一生動向無定,未能按計劃進發.但他現今已達寶座,回想十三年來遭遇的逆境一切都是造成今日寶座的條件.他的兄弟因嫉妒把他下在曠坑中,又把他賣給米甸奴販；然後轉賣給波提乏.波提乏的妻子引誘他不遂,下到監獄,又被酒政遺忘,最後為法老解夢,每次的遭遇,都是達到寶座的步驟.弟兄姊妹!或者你不明白神在你生命所作的,或者你以為所遇的苦難和憂傷太神秘；但請你記得!耶穌說:「我現在所作的事你們不明白,將來必能曉得.」
\newline
\newline
為何神帶領約瑟經長期的考驗呢?正如建築物越高越需要深固的地基；神給約瑟重大的使命,足以影響整個世界.神在約瑟身上十三年的預備,為要叫他挽救埃及全地；以色列眾民和全世界民族的生命.
\newline
\newline
親愛的弟兄姊妹；可能神有極大的使命要交託給你,要把為主作見證,傳福音直到地極的使命交託給你.約瑟對他弟兄說:「這是神差我在你們以先來,為要保全生命.」「這樣看來,差我到這裡來的不是你們,乃是神「神的手一直在約瑟身上,約瑟被神的手摸著皆有所反應,神在他身上的對付,他並不反感埋怨；他從逆境中出來,成為堅強滿有信心為主工作的人.在那段日子,神給他鋼鐵般的個性,好叫他能治理埃及全地.
\newline
\newline
你願否重新將你的生命,交在主的手中?或你還未明白主在你生命中所作的；當看見加略山的十字架時,就知道愛我們.雖然我們不全明白,但我們信靠主,現在不明白,將來必定明白.你若交託主,主必不令你失望,要記得!主帶你上升以前,必帶你下降,要有高山也必須有深淵.
\newline
\newline


\newpage

\section{第48屆~港九培靈會~張子華~第一講}
\label{sec:1051}
\textbf{神的兩個方法}
\newline
\newline
連結: \href{https://www.hkbibleconference.org/session-message/view/1051}{\texttt{https://www.hkbibleconference.org/session-message/view/1051}}
\newline
\newline
\hyperref[sec:1129]{< < < PREV SERMON < < <}
~
\hyperref[sec:toc]{[返目錄]}
~
\hyperref[sec:1052]{> > > NEXT SERMON > > >}
\newline
\newline
在這段經文中,叫我們明白神的道路,神用兩個方法叫我們明白神的道路.神對付世界和對付每一位基督徒經歷,都是常用這兩個方法.在我們屬靈的經歷中,神也是用同樣的方法,所以我們並不感到驚奇,因為神的道路本是如此.
\newline
\newline
第一個方法:神默然不語
\newline
\newline
在第一至第二節的經文中,先知哈巴谷看見國中有可怕的大事:罪惡充盈,有強暴,有奸惡；違背神的律法典章和祂的公義,因此得不到耶和華的喜悅.先知呼求神,但祂久久不答覆.神似乎不答覆他,神在那個環境中不作什麼事.同樣,這種情形也會發生在我們的生活中,因此在這種情況之下,就有人問神:為何准許這件事發生?為何准許新派神學興起?但是神卻不回答忠實子民的呼求,依然默默不語.神為何不作聲呢?我不明白.
\newline
\newline
這是我們基督徒面臨到的思想上的問題.很久之前,有一位很愛主的弟兄,他在經歷中有了困擾,他對我說:「牧師,到目前為止我仍然相信耶穌,我知道有永生,我相信聖經是神的話,我接受聖經的教訓和教義.不過在我生命中的經歷,使我不明白,因為神從來未向我笑過.」
\newline
\newline
我問他:「你怎麼知道神沒有向你笑過?」
\newline
\newline
他答:「有一件事我不明白,我信主後至今已經十多年了,雖然我屢次求神將一切困難從我身上取去；但是十多年來,神一直安靜,沒有把它拿去.因此,我一直在困難中過生活.我常在信仰上對神有了懷疑.」
\newline
\newline
弟兄姊妹,神不作聲,神安靜,這是祂所用的第一個方法.
\newline
\newline
第二個方法:神的答覆與我們所求的相反.
\newline
\newline
我們必須要學習這個方法,當我們向神祈求時,也許經過長長久的時間祂還不答覆,祂若答覆,卻與我們所期望的相反.哈巴谷先知求神復興他的國家,將仇敵,不公義和強暴的事除去,使人民有安舒的日子.
\newline
\newline
似乎神是答覆了問題,但神卻興起以色列人的仇敵,迦勒底人來對付他們,這是先知所難以明白的一件事.有一位弟兄,有一日他忽然覺得神離他很遠；因此,他跪在神前切禱,求神復興他,將愛人之心放在他裡面.但是過了一段很長的時間,不蒙神答覆,似乎離神越來越遠.後來神答覆他的祈求,但出乎他意料之外,與他所要的剛好相反.神將他放在大試煉中,使他受剌激,受打擊,遭毀壞；之後神將他抬起,給他大使命,叫他去幫助要些陷在罪中之人.可是當時他不明白神如此的答覆.
\newline
\newline
我們必須接受一個事實,神的答覆可能與我們所期待的相反,或出乎我們的意料之外,但我們也須接受.我們是否求健康,神卻將疾病加在我們身上?是否求生活上的需要,神卻使我們更貧窮?是否求喜樂,享受如人生的豐裕,神卻把荊棘痛苦加在我們身上......你是否嘗試過呢?
\newline
\newline
我們的祈求,可能神長期不答覆；我們的祈求,可能神的答覆卻與我們所求的剛好相反.不久之前,在我牧養的教會發生一件事,一個剛信主不久的弟兄,他太太和兒女都信了主.一日,我與一位長老去探望他.他對我說:「牧師,我是一個初出世的嬰孩,我求神帶領我在屬靈上長進,有一個愉快美滿的基督化家庭.」當時,他懷著一顆熱心.接著我為他禱告,求神答覆他的心願.
\newline
\newline
這位弟兄一直等待神的答覆,試想神如何答覆他呢?他的一個廿四歲的兒子,大學剛畢業,在紐約工作,買了車子,租了地方住.在七五年感恩節時,他駕車回家探望父母,卻不幸因汽車失事而喪命.我十五年牧會,從來沒有見一次喪事禮拜比這一次更傷心,更悲慘的.雖然如此,這對夫婦經過這件慘遇後,更加在靈裡扎根,信心更堅固.
\newline
\newline
神的道路奇妙,祂給我們的答覆也是奇妙的,神的道路比人的道路更奇妙.敬虔的基督徒,有時也會懷疑神的道路,好像先知哈巴谷一樣,不認識神的作為.當然,迦勒底人是神所興起對付以色列人的,正如興起未信主的人對付我們一樣,這是叫人難以明白的事.
\newline
\newline
弟兄姊妹,愛神的人一定要接受這個事實.在座的弟兄姊妹,你是否多年來生病或生活困苦；求神除去,神卻沒有答覆.在座的青年,你是否求神開路,神卻向你關門呢?你向神求健康,神卻將疾病加給你?做父母的祈望子女愛神,而他們卻離神更遠!神沒有答覆,你是否就對神懷疑,信心搖動,以至於迷惘?
\newline
\newline
我們到神面前親近祂,要明白祂道路是奇妙的,我想採用著名傳道人慕利鐘的四個事實,作為我的結束.世界上所發生的事實,無論是個人的或家庭的,我們不明白.如果所發生的事,是不如意的,我們當如何接受呢?
\newline
\newline
一,要知道一切事情皆由神掌管.
\newline
\newline
神為何興起迦勒底人呢?每一件事情的發生,皆由神掌管.
\newline
\newline
二,神要施行祂的神聖計劃.
\newline
\newline
神有祂美好的計劃在個人,在教會,在家庭中；無論是美滿的,或痛苦的,都是神所計劃進行的.
\newline
\newline
三,所發生的一切的事,都是神所定下的時間表.
\newline
\newline
神有一個美好的時間表安排,這個時間表是絕對的,準確的正如你所求的.
\newline
\newline
四,神要成就神的國度和神的旨意在我們身上.
\newline
\newline
一切事情的發生,最後的結果,無非叫神的旨意在我們身上成全.叫我們更加愛神,明白祂的旨意的事,更加榮耀祂.並且,叫我們更相信祂；祂的答覆,叫我們更加看見祂的美意本是如此,叫我們用信心去倚靠祂.
\newline
\newline


\newpage

\section{第48屆~港九培靈會~張子華~第二講}
\label{sec:1052}
\textbf{義人必因信得生}
\newline
\newline
連結: \href{https://www.hkbibleconference.org/session-message/view/1052}{\texttt{https://www.hkbibleconference.org/session-message/view/1052}}
\newline
\newline
\hyperref[sec:1051]{< < < PREV SERMON < < <}
~
\hyperref[sec:toc]{[返目錄]}
~
\hyperref[sec:1053]{> > > NEXT SERMON > > >}
\newline
\newline
哈巴谷書 1:12-2:10
\newline
\begin{longtable}{cl}
\hline
\hline
章節 & 經文 (和合本修訂版)\\
\hline
1:12 & \begin{tabularx}{0.7\textwidth}{X} 耶和華－我的神，我的聖者啊，你不是從亙古就有嗎？我們必不致死。耶和華啊，你派他為要行審判；磐石啊，你立他為要懲治人。 \end{tabularx} \\ \\ \relax
1:13 & \begin{tabularx}{0.7\textwidth}{X} 你的眼目清潔，不看邪惡，也不看奸惡，為何你卻看著人行詭詐呢？惡人吞滅比自己公義的人，為何你保持沉默呢？ \end{tabularx} \\ \\ \relax
1:14 & \begin{tabularx}{0.7\textwidth}{X} 你為何使人如海中的魚，又如無人管轄的爬行動物呢？ \end{tabularx} \\ \\ \relax
1:15 & \begin{tabularx}{0.7\textwidth}{X} 他用鉤子把他們全拉上來，用羅網捕獲他們，拉漁網聚集他們。因此，他歡喜快樂， \end{tabularx} \\ \\ \relax
1:16 & \begin{tabularx}{0.7\textwidth}{X} 向羅網獻祭，向漁網燒香；因為他藉此得豐盛的收穫與肥美的食物。 \end{tabularx} \\ \\ \relax
1:17 & \begin{tabularx}{0.7\textwidth}{X} 但他豈可因此屢屢倒空羅網，時常殺戮列國的人，毫不顧惜呢？ \end{tabularx} \\ \\ \relax
2:1 & \begin{tabularx}{0.7\textwidth}{X} 我要站在我的瞭望臺，立在城樓上觀看，看耶和華要對我說甚麼，我可用甚麼話向他訴冤。 \end{tabularx} \\ \\ \relax
2:2 & \begin{tabularx}{0.7\textwidth}{X} 耶和華回答我，說：將這默示清楚地寫在看板上，使人容易朗讀。 \end{tabularx} \\ \\ \relax
2:3 & \begin{tabularx}{0.7\textwidth}{X} 因為這默示有一定的日期，論及終局，絕不落空。它雖然耽延，你要等候；因為它必臨到，不再遲延。 \end{tabularx} \\ \\ \relax
2:4 & \begin{tabularx}{0.7\textwidth}{X} 看哪，惡人自高自大，心不正直；惟義人必因他的信得生。 \end{tabularx} \\ \\ \relax
2:5 & \begin{tabularx}{0.7\textwidth}{X} 他因酒詭詐、狂傲、不安於位；他張開喉嚨，好像陰間，如死亡不能知足，他聚集萬國，招聚萬民全歸自己。 \end{tabularx} \\ \\ \relax
2:6 & \begin{tabularx}{0.7\textwidth}{X} 這些人豈不都要提起詩歌和俗語，嘲諷他說：禍哉！你增添不屬自己的財物，靠押金發財，要到幾時呢？ \end{tabularx} \\ \\ \relax
2:7 & \begin{tabularx}{0.7\textwidth}{X} 咬傷你的豈不忽然興起，擾害你的豈不突然崛起，你就成為他們的擄物嗎？ \end{tabularx} \\ \\ \relax
2:8 & \begin{tabularx}{0.7\textwidth}{X} 因你搶奪許多國家，流人的血，向土地、城鎮和全城的居民施行殘暴，各國殘存之民都必搶奪你。 \end{tabularx} \\ \\ \relax
2:9 & \begin{tabularx}{0.7\textwidth}{X} 禍哉！那為本家積蓄不義之財、在高處搭窩、指望得免災禍的人！ \end{tabularx} \\ \\ \relax
2:10 & \begin{tabularx}{0.7\textwidth}{X} 你圖謀剪除許多民族，犯了罪，使自己的家蒙羞，自害己命。 \end{tabularx} \\ \\
[1ex]
\hline
\hline
\end{longtable}
在人生的道路上,發生的事情並非是我們所想像的,若遇到不如意的事發生時,對神的信實,愛,和公義就有了懷疑.
\newline
\newline
在我們日常生活中,所發生的事,有些是我們難以接受的.就像神若愛我,為甚麼使我患癌症呢?為甚麼我祈禱了許久,神仍未答覆呢?為甚麼神讓我家中發生不如意的事呢?因此,你就開始懷疑神的愛了!哈巴谷書的開端,很誠實的告訴我們,有很多的遭遇,我們不明白神的用意.
\newline
\newline
每一個屬靈的問題,都有它解決的辦法.每一個試煉,神必定開出路；我們受的試煉,必不至擔當不起的,神不會把難擔的擔子加在我們身上的.從哈巴谷書中,我們看到解決思想上的方法,這是最重要的信息.我們不是講病,乃是講醫病的方法.
\newline
\newline
我們受困擾,越受困擾就越難過；是因為我們沒有用神的正當方法去解答.很多時候,我們的解答,非從神而來,而是從自己而來,哈巴谷先知如何從生活中,信仰上,惆悵裡,恐懼中,再一次得到復興呢?他用甚麼方法呢?
\newline
\newline
當我們受到困擾時,應當怎樣去應付呢?
\newline
\newline
(一)不要把困擾向人多講,否則困擾就越陷越深.
\newline
\newline
(二)不要急促下結論.在我所事奉的教會,有一位姐妹是從香港來的,她到美國,很不如意.丈夫給她許多痛苦,使她很失望難過!她沒有在安靜中思想這問題,指望得人的同情,就隨便向人講；那知越講心就越痛苦,越難過,一年多沒有喜樂的日子.只有埋怨說:「神不看顧我」.
\newline
\newline
對於這樣的事,我們必須安靜等候神.
\newline
\newline
一,要思想神的公義,信實和應許
\newline
\newline
當我們來到神前,安靜思想所發生的問題之時,不要思想問題的本身.一章一至十一節講到問題的本身.但十二節忽然插入先知的一句話,這話有非常重要的屬靈教訓；就是想神信仰中不變的原則.神的永恆,公義,和慈愛；祂有豐的憐憫和無盡的恩典.要思想神的信仰,神不變的原則；這是解決問題的最好方法.
\newline
\newline
哈巴谷雖求解答問題,神卻緘默.神為何興起迦勒底人摧毀以色列國呢?先知不明白,但卻相信,神是聖潔的,不變的,公義的,祂對作惡的人終底必施以報應.倘若我們願意效法先知,問題就迎刃而解了.我昨天講到一對愛主的夫婦,他們所愛的兒子,因車禍而喪生.這位姐妹幾乎失去理智,口裡喃喃的抱怨說:「神為何如此待我?」我就勸她到休息室,為她禱告.我真不知道該用甚麼話來安慰她?使她信心不至失落.若說:「事情已經過去了,何必再哭呢!神會幫助你的.」這樣的話似乎不合情理.於是對她說:「我們所信的神,從不會做錯事的.我們縱然失信,祂仍是可信的,因祂不能背乎自己.祂是公義的,也是永恆的!」
\newline
\newline
當我講到信仰的原則時,奇妙的事就發生了,她的眼淚忽然停止了.我對她講及哈巴谷的原則,她就受感動.照樣當困擾發生在我們身上的時候,最好不要思想問題的本身,要相信神不變的原則.這樣,就能幫助我們了解神,對神所作的事,不再起疑惑.
\newline
\newline
二,要禱告將問題交託給神
\newline
\newline
先知告訴我們,我們對基本問題有懷疑時,要將問題交託給神,把問題帶到守望台.守望台是觀察軍情的,守望的兵跑到上面,用遠大的眼光警衛城池,今日基督徒也應當這樣看守神的教會.
\newline
\newline
三,要等候從神那裡尋求答案
\newline
\newline
由(二4至20)的經文中,看見神給先知哈巴谷的答案,基督徒祈禱未蒙應允,對神起懷疑是因為沒有到守望台等待神的答覆.先知把問題交託給神,卻在守望台等候神的答覆.神答覆是:
\newline
\newline
第一,不必怕,迦勒底人不過是神興起的杖,以後神會消滅他們.
\newline
\newline
第二,義人必因信得生.換言之,迦勒底人自高自大,心不正直,但義人必因信得生.對基督徒來說,我們在世上過生活,有兩條道路,須要選揀其一的:
\newline
\newline
A,信心的道路.
\newline
\newline
B,你自己理智的道路.
\newline
\newline
迦勒底人自高自大,以自己為中心,什麼事情要用自己的方法來解決,自己的方法就是自高自大,心不正直.一個用信心信靠神的人,必得著神的公義；反過來說,一個不信神的人,必得不著神的公義.弟兄姊妹,你是站在那一方面?
\newline
\newline
什麼是信心生活的定義呢?就是不以自己為中心,而過信靠神的生活.信心最高的表現,就是相信神的話.神的話如何講,我就如何去行,這才是絕對信靠神的生活.
\newline
\newline
兩條道路中,基督徒只可走一條,可能接受神的路,百分之一百信神；或走自以為義的路,中間的路線絕對沒有的.以自我為中心,自高自大的生活,有一日必要滅亡；義人必有一日勝利.
\newline
\newline
弟兄姊妹,你站在那一條路上?神對你我的要求,絕對站在信心的道路上,在一生中,過信心的生活.
\newline
\newline
我們基督徒只有一條路可走,以自我為中心的人一定滅亡,義人必有一日勝利.
\newline
\newline


\newpage

\section{第48屆~港九培靈會~張子華~第三講}
\label{sec:1053}
\textbf{苦難中的喜樂}
\newline
\newline
連結: \href{https://www.hkbibleconference.org/session-message/view/1053}{\texttt{https://www.hkbibleconference.org/session-message/view/1053}}
\newline
\newline
\hyperref[sec:1052]{< < < PREV SERMON < < <}
~
\hyperref[sec:toc]{[返目錄]}
~
\hyperref[sec:1054]{> > > NEXT SERMON > > >}
\newline
\newline
哈巴谷書 3:1-3:19
\newline
\begin{longtable}{cl}
\hline
\hline
章節 & 經文 (和合本修訂版)\\
\hline
3:1 & \begin{tabularx}{0.7\textwidth}{X} 哈巴谷先知的禱告，調用流離歌。 \end{tabularx} \\ \\ \relax
3:2 & \begin{tabularx}{0.7\textwidth}{X} 耶和華啊，我聽見你的名聲；耶和華啊，我懼怕你的作為。求你在這些年間復興你的作為，在這些年間將它顯明出來；在發怒的時候以憐憫為念。 \end{tabularx} \\ \\ \relax
3:3 & \begin{tabularx}{0.7\textwidth}{X} 神從提幔而來，聖者從巴蘭山臨到；（細拉）他的榮光遮蔽諸天，頌讚遍滿全地。 \end{tabularx} \\ \\ \relax
3:4 & \begin{tabularx}{0.7\textwidth}{X} 他的輝煌如同日光，從他手裡發出光芒，那裡隱藏他的能力。 \end{tabularx} \\ \\ \relax
3:5 & \begin{tabularx}{0.7\textwidth}{X} 在他前面有瘟疫流行，在他腳下有熱症發出。 \end{tabularx} \\ \\ \relax
3:6 & \begin{tabularx}{0.7\textwidth}{X} 他站立，震動大地，他觀看，震動列國。永久的山崩裂，長存的嶺塌陷，他的作為與古時一樣。 \end{tabularx} \\ \\ \relax
3:7 & \begin{tabularx}{0.7\textwidth}{X} 我見古珊的帳棚遭難，米甸地的幔子動搖。 \end{tabularx} \\ \\ \relax
3:8 & \begin{tabularx}{0.7\textwidth}{X} 耶和華啊，你豈是向江河發怒，向江河生氣，向海洋發烈怒嗎？你騎在馬上，坐在得勝的戰車上， \end{tabularx} \\ \\ \relax
3:9 & \begin{tabularx}{0.7\textwidth}{X} 你的弓全然顯露，箭是發誓的言語；（細拉）你以江河分開大地。 \end{tabularx} \\ \\ \relax
3:10 & \begin{tabularx}{0.7\textwidth}{X} 山嶺見你，無不戰抖；大水氾濫而過，深淵發聲，洶湧翻騰。 \end{tabularx} \\ \\ \relax
3:11 & \begin{tabularx}{0.7\textwidth}{X} 因你的箭射出光芒，你的槍閃出光耀，日月都停在原處。 \end{tabularx} \\ \\ \relax
3:12 & \begin{tabularx}{0.7\textwidth}{X} 你發怒遍行大地，以怒氣責打列國，如打穀一般。 \end{tabularx} \\ \\ \relax
3:13 & \begin{tabularx}{0.7\textwidth}{X} 你出來拯救你的百姓，拯救你的受膏者；你打破惡人之家的頭，暴露其根基，直到頸項。 （細拉） \end{tabularx} \\ \\ \relax
3:14 & \begin{tabularx}{0.7\textwidth}{X} 你以其戈矛刺透他戰士的頭；他們如旋風將我颳散，他們喜愛暗中吞吃困苦的人。 \end{tabularx} \\ \\ \relax
3:15 & \begin{tabularx}{0.7\textwidth}{X} 你騎馬踐踏海，踐踏洶湧的大水。 \end{tabularx} \\ \\ \relax
3:16 & \begin{tabularx}{0.7\textwidth}{X} 我聽見這聲音，身體戰兢，嘴唇發顫，骨中朽爛，在所立之處戰兢；但我安靜等候災難之日臨到那上來侵犯我們的民。 \end{tabularx} \\ \\ \relax
3:17 & \begin{tabularx}{0.7\textwidth}{X} 雖然無花果樹不發旺，葡萄樹不結果，橄欖樹也不收成，田地不出糧食，圈中絕了羊，棚內也沒有牛； \end{tabularx} \\ \\ \relax
3:18 & \begin{tabularx}{0.7\textwidth}{X} 然而，我要因耶和華歡欣，因救我的神喜樂。 \end{tabularx} \\ \\ \relax
3:19 & \begin{tabularx}{0.7\textwidth}{X} 主耶和華是我的力量，他使我的腳快如母鹿，又使我穩行在高處。這歌交給聖詠團長，用絲弦的樂器。 \end{tabularx} \\ \\
[1ex]
\hline
\hline
\end{longtable}
當你讀到本書第三章的時候,有沒有一個感覺,發現先知對神的態度,與第一章的時候完全不同了.進入第三章,是喜樂的一章；先知從發問題到順服神,他沒有問神為什麼.他在痛苦中有歡欣,有快樂.
\newline
\newline
在第三章裡面,先知禱告的態度,有幾點值得我們注意的:
\newline
\newline
一,謙卑的態度
\newline
\newline
他在第二節說:「耶和華阿!我聽見你的名聲,就懼怕.耶和華阿,求你在這些年間復興你的作為,在這些年間顯明出來；在發怒的時候,以憐憫為念.」這話正可表明他的謙卑態度,他沒有靠自己,沒有為自己求什麼.他的心思,集中在神的聖潔和偉大方面.
\newline
\newline
基督徒禱告,常缺乏謙卑的態度,甚至祈禱時向神討價；注力集中在自己的要求上,作有條件的禱告.但,先知來到神面前,伏服在祂腳前,沒有為自己求什麼,也沒有求神把苦難拿去；他聽見神的名字,就懼怕.今日在教會中的信徒,也應該有這種態度.
\newline
\newline
二,敬畏神的心
\newline
\newline
先知聽見神的名字就懼怕,但這種懼怕不是為了拜偶像獻長子的害怕神,乃是事奉祂,敬畏祂而畏懼祂.
\newline
\newline
三,先知的祈求
\newline
\newline
先知祈求什麼呢?他沒有求神將苦難挪去.他反而祈求神說:求你在這些年間復興你的作為.換言之,神既准許仇敵來攻打我的國家,乃是神的美意；願神懲治我們之後,再復興我們.
\newline
\newline
「這些年間」是指受痛苦,受逼迫的時候.他求在自己的國家中,成就神的旨意,他關心神的計劃和神的工作,卻沒有關心到自己的利益和工作,但求神藉苦難成就祂的計劃.
\newline
\newline
我們的祈求卻與先知的祈求相反.當苦難臨到時,我們會求神將苦難挪去.這種祈禱是不成熟的祈禱,也是不信靠神的祈禱.基督徒要關心神的計劃和工作,當苦難打擊發生時,應要祈求神「在這些年間復興祂的作為.」求神藉著苦難,「復興祂的作為在我身上」,除此以外,我們不應有別的關懷了.
\newline
\newline
我住的地方,離開蔡蘇娟姐妹寓所不遠,開車半小時可達.所以我常有機會去探她.當我進入她的暗室之後,感受到的是神,是光,是心靈的滿足.她時常這樣說:「我躺在床上四十多年,從無問神為什麼?我只有一個問題問神:神阿,你要在我身上作什麼?」
\newline
\newline
我不知在座弟兄姊妹,有沒有生活上的難處?疾病的纏繞?所愛的人患了不治之症?在家中沒有喜樂,夫婦是否不和?但我們來到神面前,要像先知的態度一樣,求神藉這些苦難成就祂的計劃.
\newline
\newline
這一次我回到香港和東南亞事奉神,曾經受過神的對付.我答應過神回來遠東事奉,事後我覺得似乎對神的答覆是錯了,我想到環境,兒女,工作,居住的地方,師母的那一份工作,心裡起了惆悵.本來在去年就要回來的,但心中一直掙扎,掙扎不捨得離開自己的崗位.然而神有方法對付我,祂用苦難加在我們身上,起初,我像先知哈巴谷一樣,求神拿去我身上的苦難,我不明白神的旨意.神藉著苦難,趕我回到遠東.
\newline
\newline
哈巴谷書第三章,在信仰上能夠幫助我們.這一章講到神曾在以色列民族中的所作所為,神從提幔而來.先知在祈禱中,忽然想起神非常愛祂的子民,帶領他們過紅海,四十年在曠野,過約但河,進入應許的流奶與蜜之地.神在以色列人中施行恩典,叫先知大得安慰.
\newline
\newline
哈巴谷先知的信心,有事實作為根據,他看到神在他國家中所做的工作,給他很大的信心.因此,基督徒也應該用這個原則,想到神在我們身上所做的工作,就能克服對神的懷疑了.
\newline
\newline
有一位弟兄,忽然遭受很大的打擊,精神上和物質上受極大的痛苦.因此,對神起了懷疑,他也埋怨神,說神不愛他.我對他題及神不變的原則,勸他祈禱,等候神.並問他說:「弟兄!你幾時得救的?」他答:「十年前.」我再問他:「你能否將你得救的經過講給我聽?」他答:「我是個大罪人,我的得救實在是一個大神蹟.」於是他把得救的大神蹟告訴我,並且他一連講了四,五個信心的經歷,忽然之間就醒悟過來；睜大了眼睛,因為他已經驗到神的慈愛和信實.最後我對他說:「你這次雖遭遇苦難,神沒有理由不愛你的.」他答:「是啊!」
\newline
\newline
哈巴谷先知從過去的事,再從自己實在的經驗,恢復了對神的信心.
\newline
\newline
一個信靠神的人,雖落在百般試煉中仍要以神樂.喜樂是神所賜給我們的,並非是我用勞力賺回來的；喜樂是神給每一個基督徒的禮物.無論落在任何環境中,都應當以神為樂.
\newline
\newline
在任何環境中,都不要失去喜樂,神是信實的.在絕對痛苦時,請採用哈巴谷先知所持守的態度,仍要因耶和華歡欣,他說:「雖然無花果樹不發旺；葡匋樹不結果,橄欖樹也不效力,田地不出產糧食,圈中絕了羊,棚內也沒有牛；然而我要因耶和華歡欣,因救我的神喜樂.」以色列人多數務農,以務農為生.雖然地裡不出產,甚至牛羊也絕種了,落在這個環境中,先知對神說:「我要因耶和華歡欣,因救我們的神喜樂.」這就是哈巴谷書的信息.
\newline
\newline
瑪拉基書:一個離開神的心所發的怨言
\newline
\newline


\newpage

\section{第48屆~港九培靈會~張子華~第四講}
\label{sec:1054}
\textbf{懷疑神的愛}
\newline
\newline
連結: \href{https://www.hkbibleconference.org/session-message/view/1054}{\texttt{https://www.hkbibleconference.org/session-message/view/1054}}
\newline
\newline
\hyperref[sec:1053]{< < < PREV SERMON < < <}
~
\hyperref[sec:toc]{[返目錄]}
~
\hyperref[sec:1055]{> > > NEXT SERMON > > >}
\newline
\newline
瑪拉基書 1:1-1:5
\newline
\begin{longtable}{cl}
\hline
\hline
章節 & 經文 (和合本修訂版)\\
\hline
1:1 & \begin{tabularx}{0.7\textwidth}{X} 耶和華的話，藉瑪拉基傳給以色列的默示。 \end{tabularx} \\ \\ \relax
1:2 & \begin{tabularx}{0.7\textwidth}{X} 耶和華說：「我曾愛你們。」你們卻說：「你在何事上愛我們呢？」耶和華說：「以掃不是雅各的哥哥嗎？我卻愛雅各， \end{tabularx} \\ \\ \relax
1:3 & \begin{tabularx}{0.7\textwidth}{X} 惡以掃，使他的山嶺荒涼，把他的地業交給曠野的野狗。」 \end{tabularx} \\ \\ \relax
1:4 & \begin{tabularx}{0.7\textwidth}{X} 以東若說：「我們雖被毀壞，卻要重建荒廢之處。」萬軍之耶和華如此說：「任他們建造，我必拆毀；人必稱他們為『邪惡之境』，為『耶和華永遠惱怒之民』。」 \end{tabularx} \\ \\ \relax
1:5 & \begin{tabularx}{0.7\textwidth}{X} 你們必親眼看見，你們要說：「耶和華在以色列疆界之外必尊為大！」 \end{tabularx} \\ \\
[1ex]
\hline
\hline
\end{longtable}
神十分看重這卷先知書,神藉先知傳給以色列人,這是一個有負擔重著負擔的默示.(原文直譯)神叫我們看重本書,且要應用在生活中.
\newline
\newline
耶和華說:「我曾愛你們.」我讀到這裡的時候,就心中發生一個疑問,是否神現在不愛我們呢?但是,我給我一個亮光.耶和華說:「我曾愛你們.」並非說神現在不愛我們,乃是說現今「你們不接受我的愛.」我們必須接受這個觀念:神對我們的愛永不改變,當我們與神的關係發生障礙之時,我們可能懷疑神的愛,不接受神給予我們永恆的愛.
\newline
\newline
我們所遭遇的,以我們環境,都可能會埋怨神,對神說:「神啊!你在那一件事上,叫我知道你愛我呢?」一個基督徒若落到這樣的地步,覺得神不愛他.懷疑神對他的愛,實在太可憐了,特別當我們的環境像一片烏雲來到時,會令我們見不到神的愛.
\newline
\newline
有一個姊妹,家庭本來很幸福快樂,但苦難卻接踵臨到.她有三個兒女,大女兒因發高熱,理智失常,痴呆了!第一次打擊臨到這個家庭.
\newline
\newline
過了三年,十四歲大的兒子因呼吸器官與肺之間患瘤腫,又因注射麻醉藥,呼吸困難而死亡,第二次打擊再臨到這個家庭.當我得知這消息便去探望這個家庭.她的丈夫第一句話問我說:「神的愛在那裡呢?」
\newline
\newline
弟兄姊妺!是否在我們的家庭中,也有不如意的事發生呢?疾病苦難,死亡......許多打擊接踵而至,使你對神的愛發生了懷疑?因此,你會對神說:「神啊!你太不愛我,你在那件事上叫我看見你的愛呢?」也許你會說:「信主的人是應該喜愛的,因為神是我們患難中隨時的幫助；但我信主後,至今還沒有經歷過.不但如此,在我信主之後卻比未信主前更壞更差,早知如此,我寧願不信了.神啊!你在那一件事上,叫我覺得你愛我呢?」
\newline
\newline
我們是否已開始懷疑神的愛呢?
\newline
\newline
也許大家認識艾德理先生,他多年在中國傳道.她的女兒是學護士的,預備和醫科畢業的夫婿到東南亞事奉神的.但這個青年醫生在那年接受恩典神學院,暑期聖經訓練時,不幸在游泳時溺斃了.請問該如何解釋呢?一對愛主的青年夫婦,肯犧牲名利而奉獻到東南亞工作,竟然有此遭遇.神若愛她,證據在那裡?
\newline
\newline
若懷疑神的愛,就會使我們與神的關係受到致命傷.因為我們與神的關係,乃是建立在愛的根基上.如果我們對神有懷疑,就破壞了這個關係.正如夫妻,乃是建立在無條件的愛上,互相尊重；若懷疑對方,家庭就分裂.我牧會十五年,看見許多離開神的人,開始之時只不過懷疑神的愛.當生活受到打擊時,要求神給予愛的證據；如果得不著證據,就離開神越來越遠了.這是非常嚴重的事,懷疑神的愛,是與神的關係破裂的致命傷!
\newline
\newline
你若懷疑神的愛不妨看以下的兩個人,就是雅各和以掃.神叫我們所看見的重點不是雅各,而是以掃.從以掃的遭遇,你會知道神明明是看顧雅各的.很多人都以痛苦來解釋痛苦問題,那就是說:如果你覺得痛苦,你就看一看比你更痛苦的人.
\newline
\newline
今日世界有一個哲學:若要解決你的痛苦,就要看一看比你更痛苦的人.神叫我們看見一個事實,神對一個人的處罰,像雅各的哥哥以掃一樣.神並非叫我們為了解決痛苦,就看那些比我們更痛苦的人.這不是神的方法.也不是叫我們看到一個更慘人,而我們比他稍好,我們就會得到安慰,這都不是神的方法.這一段經文給我們一個屬靈的寶貴含義在裡面,乃是叫我們在受苦的時候,神說:你要仰望那一位受神鞭打的哥哥.我自己得到一個結論,我們基督徒只有一位長兄,聖經說我們的長兄就是耶穌基督.仰望那愛我們的主耶穌基督,讓我們再回來十字架上面,這是解決問題最重要的關鍵.
\newline
\newline
若心中懷疑神不愛我們,而要求神給證據我們看的時候,就要回到主為我們釘十字架的事上,這是神愛我們的最大證據了.以掃和主耶穌惟一不同之處,他只背負自己的罪,而主卻背負我們的罪.我們仰望主的十字架,便可消除我們懷疑神的愛.神的愛長,闊,高,深,都在十字架上為我們顯明了.
\newline
\newline
(一)十字架是我個人生活的力量,也是每一個人生活的力量.我記得我有一個朋友,有一年他在夏令會中滿有懷疑,無法聽道,他逃避聚會.正在困倦沒有辦法之時,忽然求神帶領他回到各各他的山上.當他思想完十字架的時候,他一切的懷疑便消失了.
\newline
\newline
我相信在座的每一個人,一定會同意我說的,我們事奉神有佷多的痛苦,是無法講給人聽的.當我們受人的批評和攻擊時,我真覺得走事奉神的道路不容易,每當試探或攻擊來到,唯一的方法叫我得釋放的,就是仰望主的十字架.主為我被神鞭打傷透了,這樣,我們還需要第二個證據嗎?不需要了.因為耶穌基督的愛是神愛我們的最大中心,也是神的啟示.
\newline
\newline
(二)我們從神的誡命看見神的愛,所以我們要好好思想神一切的誡命.神的出發點是為了愛我們.例如六日工作,第七日休息；保羅說:我們應當一無掛慮.神說:要把一切重擔卸給我......神看顧我們,一切好處都是為我們.
\newline
\newline
(三)從神的鞭打看見神的愛,做父母的就看見這個真理,鞭打是愛的最大真理.父母不鞭打兒女就不愛兒女,鞭打是證明父母愛兒女最高的表現.兒女犯罪要鞭打,這個功課是學不完的.神愛我們,才鞭打我們.我們在神的鞭打中,應該看見神的愛.
\newline
\newline
今日,神對我們的信息很簡單,但對我們很重要很需要.我不是先知,但我知道在座的弟兄姊妹,可能在你的生活中痛苦,你正在懷疑神.在此,我們一同抬頭仰望我們的長兄耶穌基督,這是神給我們最大的愛的證據,這個愛永遠不改變.
\newline
\newline


\newpage

\section{第48屆~港九培靈會~張子華~第五講}
\label{sec:1055}
\textbf{藐視神的名}
\newline
\newline
連結: \href{https://www.hkbibleconference.org/session-message/view/1055}{\texttt{https://www.hkbibleconference.org/session-message/view/1055}}
\newline
\newline
\hyperref[sec:1054]{< < < PREV SERMON < < <}
~
\hyperref[sec:toc]{[返目錄]}
~
\hyperref[sec:1056]{> > > NEXT SERMON > > >}
\newline
\newline
瑪拉基書 1:6-1:6
\newline
\begin{longtable}{cl}
\hline
\hline
章節 & 經文 (和合本修訂版)\\
\hline
1:6 & \begin{tabularx}{0.7\textwidth}{X} 萬軍之耶和華對你們說：「兒子孝敬父親，僕人敬畏主人；我既為父親，孝敬我的在哪裡呢？我既為主人，敬畏我的在哪裡呢？你們這些藐視我名的祭司啊！」你們卻說：「我們在何事上藐視你的名呢？」 \end{tabularx} \\ \\
[1ex]
\hline
\hline
\end{longtable}
事奉神的人,常被人藐視.我們被人藐視,有時並非是別人的錯,乃是我們有錯.神責備當時的祭司,藐視祂的名,不尊敬祂的名.
\newline
\newline
祭司問:「我在何事上藐視你?不尊重你呢?」神回答說:「你們將污穢的食物獻在我的壇上,且說,我們在何事上污穢你呢?因你們說,耶和華的桌子,是可藐視的.你們將瞎眼的獻為祭物,這不為惡麼?將瘸腿的,有病的,獻上,這不為惡麼?」
\newline
\newline
為何祭司會這樣做呢?他們工作一日復一日,年復一年,每日對著工作,已到厭煩的地步.開始對工作覺得討厭時,對神的事奉就會馬虎起來.祭司本來很仔細檢查獻祭的牛羊,但工作一兩年後,對於獻祭人拖來的牛羊,就以最快速手法完成；他們既討厭神的工作,態度就馬虎起來了.他們在事奉上,再不會體會神的心意而行.
\newline
\newline
在座的同工們,我不知道你在教會中事奉的光景如何?如果你從來沒有討厭過神的工作,我要為你感謝神.在我工作的經驗中,我曾對神的工作馬虎,甚至討厭神的工作,我討厭教友找我,我討厭有電話來.回想我剛奉獻讀神學時,無論如何痛苦,艱難,我很堅決要事奉主.但在事奉上,年復一年就好容易變成討厭工作了,因而對工作馬虎了.今日教會中的傳道人,有多少個能付代價去預備講章的?有多少個能為弟兄姊妹解答問題的?是否覺得討厭,沒法叫他們快速離去呢?
\newline
\newline
因為工作的忙碌,使我們從早到晚要坐在辦公室裡,要處理會務策劃工作,這樣久而久之,我們就對工作開始覺得厭煩了.既厭煩,就會馬虎起來!甚至把神的桌子當作是污穢的,在工作上不尊敬他的名.不但傳道人有這個危險,即使平信徒也是如此,好像教主日學,你是否重視?你的工作如何?教導真理的工作是榮耀的,上算的,但你準備功課是否忠心?有沒有用時間去探訪學生?你是否在巴士上預備教主日學的功課?你有沒有關心到主日學生的靈性呢?
\newline
\newline
長老,執事第一次被神選上時,是很快樂的；因為可以為神的家,為兄弟姊妹而事奉,以盡做長執的本分.但經過年日的磨煉後,工作的熱忱就沒有了,只會在教會中命令別人,做這樣,做那樣,變成了馬虎的執事長老了.
\newline
\newline
今天,讓我們看見耶和華如何對付以色列的祭司,如何責備他們,等如責備我們一樣.你應該省察,你在事奉上,到達了什麼地步?你工作的態度如何?與神的關係如何?你在工作上忠心否?是否已盡你的職份.
\newline
\newline
祭司是應該敬畏耶和華的,可惜他們反而在事奉上,藐視耶和華的名.如困你真已落在這樣的光景裡,你就要懇求神,求祂好施恩給你(一9).神不恂情面的對祭司責罵說?「甚願你們中間有一人關上殿門,免得你們徒然在我壇燒火.萬軍之耶和華說,我不喜悅你們,也不從你們手中收納供物.」如果一個人到了這光景,就要求神給足夠的恩典,在事奉上謙卑悔改.這是解決問題的第一個方法.
\newline
\newline
解決問題的第二個方法,就是我們在事奉中要榮耀神.以榮耀神為我們最高的目標.若果我們在事奉上,時時看重神的榮耀,神就會幫助我們,使我們不再馬虎地去工作.難道傳道人不盡心盡意講道就不能榮耀神麼?試問傳道人在每一週的生活中,用多少時間在看視,多少時間上茶樓,是否這些時間已多過事奉神的時間,這樣的生活難道會榮耀神麼?傳道人以自己的事為念多於事奉神的事為念,是否可以榮耀神嗎?真正榮耀神才是我們事奉神的目標.
\newline
\newline
至於解決問題的第三個方法看重神的話.這是每一個事奉神的人,所必需要的.事奉神的人,要與神有密切的交通,建立關係.這關係可以從日常靈修生活中表現出來；除祈禱以外,要將神話常放在心內,不斷的追求思想,對神的話,每日要有新鮮的感覺,這正是每一個事奉神的人,最大的需要.
\newline
\newline
如果我們對神工作的感厭煩,正因為我們破壞了對神的關係.如果對神的話語,得不到新的供應；那麼,我們的工作就會感到沉悶,成了重擔,感到討厭.所以我們事奉神,每日要與神交往,與神說話,要在神的話語上好好地扎根,才能得到新的供應；才有活力,向前奔走.
\newline
\newline
倘若在神的話語上,失去了興奮性；那麼,走神的路,做神的工作,就會感到厭煩.我常有這個痛苦的經驗.為了工作,為了家庭,為了兒女,終日忙碌,而忽略了祈禱讀聖經,就無法獲得新鮮的供應,既無新鮮供應,神也不再說話,我們在工作上自然就會生厭了.所以我們不但要悔改,而且要看重神的供應.
\newline
\newline
當我們在苦悶中之時,我們要想起神以前我們所立的約(二4)祂說:「我曾與他立生命和平安的約,我將這兩樣賜給他,使他存敬畏的心,他就敬畏我,懼伯我的名.」這裡提醒我們想著一件事,就是以前對主所講的話和所立的約.要知道我們事奉神是極榮耀的,每一個做牧師的有上好的福分,是世人所沒有的,是神所選召的.在我十五年來的牧會生命中,很容易灰心失意.我曾對神說:「我不幹了.」所以要常常思想與神所立的約,提醒自己!
\newline
\newline
在座中的同工,主日學教員,長老,執事,要回想你與神所立的約；在你灰心之時,可求祂賜你力量,使你從新得力,繼續事奉祂.
\newline
\newline
弟兄姊妹,先知書寫在幾千年前,神知道我們的軟弱,祂藉著先知提醒我們每一個事奉祂的人,不可厭煩和馬虎地為神工作,我們要在這事上悔改,從新得力,一生一世事奉祂,跟隨祂.
\newline
\newline


\newpage

\section{第48屆~港九培靈會~張子華~第六講}
\label{sec:1056}
\textbf{懷疑神的公義}
\newline
\newline
連結: \href{https://www.hkbibleconference.org/session-message/view/1056}{\texttt{https://www.hkbibleconference.org/session-message/view/1056}}
\newline
\newline
\hyperref[sec:1055]{< < < PREV SERMON < < <}
~
\hyperref[sec:toc]{[返目錄]}
~
\hyperref[sec:1057]{> > > NEXT SERMON > > >}
\newline
\newline
瑪拉基書 2:17-2:17
\newline
\begin{longtable}{cl}
\hline
\hline
章節 & 經文 (和合本修訂版)\\
\hline
2:17 & \begin{tabularx}{0.7\textwidth}{X} 你們用言語使耶和華厭煩，卻說：「我們在何事上使他厭煩呢？」因為你們說：「凡行惡的，耶和華看為善，並且喜愛他們；」又說：「公平的神在哪裡呢？」 \end{tabularx} \\ \\
[1ex]
\hline
\hline
\end{longtable}
今日要講的問題就是很多基督徒在生活上掙札,因看不到公義的事,所以懷疑.(三1-7)講到主再來的事,到那日作惡的人得不到神的喜悅,但敬虔受苦的人就得著神的喜悅.詩篇七十三篇特別講到懷疑神供應的問題.
\newline
\newline
弟兄姊妹!我在過去牧養教會的經驗中,看見很多基督徒受困擾時,看不見神的慈愛；在困擾的環境中,使傳福音的工作受到攔阻.有人覺得這個世界真不公平,我不偷竊,不搶劫；但那些惡人什麼都做,搶劫,殺人,偷竊......他們卻樣樣都好,真是不公平了,這世界那裡有神呢?若果有神存在,有神管理,世界上就不致有那麼多不公平的事了.這是思想上的束縛,這篇詩可能幫助我們解除.寫詩的人有這個經驗,本篇詩的結論在第一節:「神實在恩待以色列那些清心的人.」
\newline
\newline
二至十四節,作者寫出自己的經驗,心裡有過這些掙扎,就如我們看見惡人,狂妄人順利,就差不多跌倒.他們作惡褻慢神,反而住好屋,坐汽車,他們無須受苦,為人卻非常驕傲.他們豐富到一個地步,連沒有想到的都有了.他們欺壓窮人不但作惡,用口褻瀆神；還用各種方法咒詛福音.第十節詩人轉一個角度來講,我們相信神,卻喝盡了滿杯的苦水.第十一至十二節形容惡人的情形.第十三節是描寫我與惡人的比較,覺得自己太冤枉了,我愛神,遵神的律例；卻每天受苦,受神的懲治,這樣我如何能看見神的公義呢?神阿!你不公平.
\newline
\newline
弟兄姊妹,不知道你的生活,與不信神的親戚比較,在思想中是否有掙扎?有很多基督徒在思想上確有大掙扎,這是非常嚴重的事.我們懷疑神的愛和公義,若我們不得勝這個掙扎,會叫我們離開神,甚至做出不應當做的事來.
\newline
\newline
作者在困擾中如何得勝呢?
\newline
\newline
第一,禁止我們的舌頭
\newline
\newline
不講不應講的,如果對人講心懷不平的話,就是得罪了人,所以要禁止舌頭.
\newline
\newline
第二,要站在神的立場上看問題
\newline
\newline
站在人的立場上,用理智,思想,哲學去看問題,一定是沒有辦法明白的.第十七節是本篇詩的轉捩點:「等我進了神的聖所,思想他們的結局.」作者轉移方向,站在神的立場上,用神的眼光,透過神來看問題,就明白了一切.這是非常重要的事,也是基督徒處世的最重要原則.
\newline
\newline
第十七節「等我進了聖所」這句話,有三種不同的看法和解釋:
\newline
\newline
A,基督徒對財寶的看法與世人不同.世人以金錢多少來衡量人生的地位,用來量度人生的意義和關係.基督徒若站在神的立場,進入聖所來看這問題；就不以金錢衡量人生的失敗與成功,快樂或憂愁,生活是否有意義.我們看十八至十九節,就會知道,原來金錢是靠不住的.弟兄姊妹,不知你對金錢的看法如何?若想以金錢來量度成就,用它來買快樂,這就是一種錯誤的看法.
\newline
\newline
金錢多少不能代表祝福,富有的人并不能代表全是神賜福,神把富有的人在滑跌和沉淪之中,使他們驚慌不安.我見過很多有錢人,心靈不安,有些人正往地獄裡走.
\newline
\newline
B,對災禍的看法與世人不同.世人看災難是不幸的,痛苦的,或者以為自己前生作過什麼惡事,但第廿一至廿二節告訴我們,原來災難臨到我們,有兩個可能性:(a)出於神的管教.(b)出於神的試煉.因神要把更大的福賜給我們.
\newline
\newline
弟兄姊妹,在靈裡未受過災難痛苦的基督徒,不是成熟的基督徒.人在順景之時,就忘記神,所以神要使災難臨到我們,是要管教我們,使我們不要離開祂的恩典.我曾探望過蔡蘇娟姊妹,她說:「我雖不明白神為何准許苦難臨到我,但神的美意本是如此!」你去看她時,似乎覺得她很可憐；不過每日卻有很多人去看她,許多人聽她作見證後,也受感動而信了主.她還說:「神藉著我寫的書,我的禱告,都幫了別人.倘若我健康時,反而不能幫助人哩!也許我還留在國內不能出來.你看,苦難反而成了我的祝福.」
\newline
\newline
弟兄姊妹；不知你對苦難的看法如何?從神的立場看,原來苦難是對我們,有那麼大的好處.
\newline
\newline
C,基督徒對人生最高目標和美滿享受與世人不同,有人說,有錢有學問就是享受,這是世人的看法.基督徒卻要愛慕神的事情,在神裡面得著滿足,以神為樂.覺得有神的同在是一種最大的福份.神的話是最美好的糧食,正如第廿五節詩人所說的:「除你以外,在天上我有誰呢?除你以外,在地上我也沒有所愛慕的.」
\newline
\newline
弟兄姊妹,我們為甚麼會覺得心靈空虛,患得患失呢?
\newline
\newline
倘若我們學好了這功課,以神為樂,以神為我們的一切；在神裡面得著一切的滿足.那麼「除了神以外,在地上什麼也沒有所愛慕了!」
\newline
\newline
弟兄姊妹,你的人生觀念是什麼呢?你以為什麼是人生最高的目標?人生最高的享受是什麼?詩人說:除你以外,在天上我有誰呢?除你以外,在地上我也沒有所愛慕的!」願神幫助我們,一起到神面前,學習這三個寶貴的功課吧.
\newline
\newline


\newpage

\section{第48屆~港九培靈會~張子華~第七講}
\label{sec:1057}
\textbf{奪取神的供物}
\newline
\newline
連結: \href{https://www.hkbibleconference.org/session-message/view/1057}{\texttt{https://www.hkbibleconference.org/session-message/view/1057}}
\newline
\newline
\hyperref[sec:1056]{< < < PREV SERMON < < <}
~
\hyperref[sec:toc]{[返目錄]}
~
\hyperref[sec:1058]{> > > NEXT SERMON > > >}
\newline
\newline
瑪拉基書 3:8-3:12
\newline
\begin{longtable}{cl}
\hline
\hline
章節 & 經文 (和合本修訂版)\\
\hline
3:8 & \begin{tabularx}{0.7\textwidth}{X} 人豈可搶奪神呢？你們竟搶奪我！你們卻說：『我們在何事上搶奪你呢？』其實就是在你們當納的十分之一奉獻和當獻的供物上。 \end{tabularx} \\ \\ \relax
3:9 & \begin{tabularx}{0.7\textwidth}{X} 因你們全國上下都搶奪我的供物，詛咒就臨到你們身上。 \end{tabularx} \\ \\ \relax
3:10 & \begin{tabularx}{0.7\textwidth}{X} 你們要將當納的十分之一全然送入倉庫，使我家有糧，以此試試我，是否為你們敞開天上的窗戶，傾福與你們，甚至無處可容。這是萬軍之耶和華說的。 \end{tabularx} \\ \\ \relax
3:11 & \begin{tabularx}{0.7\textwidth}{X} 我必為你們斥責蝗蟲，不容牠毀壞你們的土產。你們田間的葡萄樹，果實未熟以先也不會掉落。這是萬軍之耶和華說的。 \end{tabularx} \\ \\ \relax
3:12 & \begin{tabularx}{0.7\textwidth}{X} 萬國必稱你們為有福的，因你們必成為喜樂之地。這是萬軍之耶和華說的。」 \end{tabularx} \\ \\
[1ex]
\hline
\hline
\end{longtable}
我從美國回來,就知道在香港搶扒手,非常之厲害；這些可憎惡可怕的罪行,每天不斷地發生著.
\newline
\newline
弟兄姊妹!我們在神的面前,在奉獻的事上,用不同的理由和藉口,來搶劫神的供物,若非神的話,我實在沒有膽量講這題目了.
\newline
\newline
今天的經文內,我可用四方面來講解與各位分享.
\newline
\newline
一,奉獻十分之一給神
\newline
\newline
不納當納的十分之一給神,是一種嚴重的罪行.我相信有一日在天上,我們站在神的面前；會看見很多基督徒,在這方面失敗.很可惜!我們對神的話,沒有用嚴謹的態度去接受；因此就輕看了十分之一的奉獻.又多少時候,我們對於神的誡命,覺得有些是重要的,有些卻不是重要的.在座的每一位,當然我們沒有膽量去殺人,拿槍去打劫；但我們卻會用自己的聰明和藉口,和不擇手段的去得著我們想要得的,這在神看來是極嚴重的罪行啊!
\newline
\newline
我們有沒有想過,不獻十分之一是犯罪呢?在神眼中看我們所行是賊呢?這事在人看來,卻不覺得重要,當作很隨便.所以,對於我們屬靈的長進,靈裡的活潑,對神的信心,對愛神的心,就成了致命的打擊.盼望我們能夠接受神的話在心中,在奉獻的態度上,確信若不奉獻十分之一,就是非常嚴重的罪行.
\newline
\newline
二,十分之一是當納的
\newline
\newline
我們講慣了十分之一奉獻,其實,看清楚一點,就知道十分之一不是奉獻,那是當納的數目.那部份不是屬於我們的,是屬於神的.在十分之一之外才是奉獻,我們有自由要奉獻多少,但應隨聖靈的帶領而奉獻.
\newline
\newline
神恩待我們,眷顧我們,祂知道我們的需要.真正愛神的人,和靈命長進成熟的基督徒,他必會因愛神而奉獻,神斷不會強迫他的.這十分之一不是我們的,是神的；所以這是當納的,應該屬於祂.若將十分之一拿來私用,就等於搶劫給神的供物.
\newline
\newline
三,基督徒要學習倚靠神並澈底愛祂
\newline
\newline
在這段聖經中,神給我們一個屬靈的挑戰,要學習全然的信靠祂,並切實的愛祂.
\newline
\newline
許多教會中的信徒,認為牧師講奉獻的道理是不屬靈的,有問題的；其實,基督徒是欺侮傳道人,虧欠傳道人.做傳道人的,認為奉獻,是很重要的真理.其重要性,并非金錢的本身；在十分之一的後面,乃是基督徒對神的信心和愛心的表現.若在十分之一的奉獻上虧欠神,你就無法告訴人知道你是愛神的,是信靠神的.若在生活上,沒有愛神之心的表現,這個愛心就是虛假的.神是可靠,可信的.不肯奉獻的人,總有很多的藉口,窮有窮的藉口,富有富的藉口,總之,有他們的藉口.窮的說,我所賺的錢,就算不奉獻還不夠開支；富的說,我雖然有錢,但手上沒有現金,因為已放在房屋,股票上,我拿不出來.即使過得去的也說:我要儲蓄多一些,留待發展生意,等到有錢賺時才奉獻吧!
\newline
\newline
我們若肯奉獻十分之一,神就必傾福與我們,甚至無處可容.并且還要把更多屬靈的福氣賜給我們.叫我們有平安有喜樂,能看見祂是全然可信的.我們若肯在奉獻上有信靠神的表現,神必然大大的賜福給我們.
\newline
\newline
神要求我們要用最大的信心去信靠祂,愛祂；這樣,當我們有缺乏時,神必敞開天上的窗戶,傾福與我們,甚至無處可容.
\newline
\newline
四,十分之一是神給我們的恩典
\newline
\newline
第十一至十二節的經文,我們要了解到以色列是一個農業的國家.蝗蟲對農作物,是一個嚴重的大損害.但神說:「我必須為你們斥責蝗蟲,不容牠毀壞你們的土產；你們田間的葡萄樹在未熟之先,也不掉果子.」換言之,我們能將十分之一獻給神,這已經是神的恩典了,如果不是神所賜,恐怕連應該奉獻的也不夠呢!
\newline
\newline
有人以為基督徒是瞎眼的,因為他想錢是我自己賺回來的.那裡是神所賜的呢?不錯,你每天作工,賺錢,但也須有健康的身體.那麼,健康的身體是誰所賜的?原來,我們成了何等樣的人,完全是出於神的恩典.誰給你恩典?是神.
\newline
\newline
願神光照我們,在祂面前謙卑悔改,不要強詞奪意說,我們生活在新的時代,毋須遵守舊約十分之一奉獻的規條.這樣說法,雖有部份理由,但主耶穌對新約的基督徒要求是更大的,因為要捨己,全所有奉獻.
\newline
\newline
舊約不過十分之一的奉獻,但在新約時代,神賜給我們要的恩典既然更多；那麼,就算是十分之二的奉獻也十分應該的.所以,我們要在神的光照之下,徹底悔改!以前偷竊,現在不偷竊,在奉獻上,做一個神所喜愛的人.
\newline
\newline


\newpage

\section{第48屆~港九培靈會~張子華~第八講}
\label{sec:1058}
\textbf{兩個問題}
\newline
\newline
連結: \href{https://www.hkbibleconference.org/session-message/view/1058}{\texttt{https://www.hkbibleconference.org/session-message/view/1058}}
\newline
\newline
\hyperref[sec:1057]{< < < PREV SERMON < < <}
~
\hyperref[sec:toc]{[返目錄]}
~
\hyperref[sec:777]{> > > NEXT SERMON > > >}
\newline
\newline
瑪拉基書 2:10-2:16
\newline
\begin{longtable}{cl}
\hline
\hline
章節 & 經文 (和合本修訂版)\\
\hline
2:10 & \begin{tabularx}{0.7\textwidth}{X} 我們豈不都有一位父嗎？豈不是一位神創造了我們嗎？為何互相行詭詐，褻瀆了神與我們列祖所立的約呢？ \end{tabularx} \\ \\ \relax
2:11 & \begin{tabularx}{0.7\textwidth}{X} 猶大行事詭詐，在以色列和耶路撒冷中行了可憎的事；因為猶大褻瀆耶和華所喜愛的聖殿，娶外邦神明的女子為妻。 \end{tabularx} \\ \\ \relax
2:12 & \begin{tabularx}{0.7\textwidth}{X} 凡做這事的，無論是清醒的或回應的，即使獻供物給萬軍之耶和華，耶和華也要將他從雅各的帳棚中剪除。 \end{tabularx} \\ \\ \relax
2:13 & \begin{tabularx}{0.7\textwidth}{X} 你們又再做這樣的事，使哭泣和嘆息的眼淚遮蓋耶和華的祭壇，以致耶和華不再理會那供物，也不喜歡從你們的手中收納。 \end{tabularx} \\ \\ \relax
2:14 & \begin{tabularx}{0.7\textwidth}{X} 你們還說：「這是為甚麼呢？」因為耶和華在你和你年輕時所娶的妻之間作證。她雖是你的配偶，你誓約的妻，你卻背棄她。 \end{tabularx} \\ \\ \relax
2:15 & \begin{tabularx}{0.7\textwidth}{X} 一個人如果還剩下一點靈性，他不會這麼做。這人在尋找甚麼呢？神的後裔！當謹守你們的靈性，誰也不可背棄年輕時所娶的妻。 \end{tabularx} \\ \\ \relax
2:16 & \begin{tabularx}{0.7\textwidth}{X} 耶和華－以色列的神說：「我恨惡休妻的事和衣服外面披上暴力的人。所以當謹守你們的心，不可行詭詐。這是萬軍之耶和華說的。」 \end{tabularx} \\ \\
[1ex]
\hline
\hline
\end{longtable}
第一個問題:與弟兄姊妹的關係?
\newline
\newline
經文:瑪二10-16
\newline
\newline
在教會裡面有罪惡,就可使教會在社會的地位上,完全失去見證.弟兄之間,若發生了問題,就干涉到神的聖潔,因此就影響我們與神之間的關係,祂就不再接納我們的供物和禱告.舊約如此講,新約也如此講.我們獻祭時,若與弟兄之間有問題,就先要放下祭物,先與弟兄和好,然後再回去獻祭.
\newline
\newline
在教會中肢體的生活,與世人不同的,這要叫世人知道,我們是主的門徒,證明我們是神的兒女.如果失去這見證,我們對社會的影響力就完全失去了.今日在教會中有一件痛心的事,就是用詭詐彼此相待.「詭詐」在英文或希臘文,意思是包括彼此之間,互不信任,狡猾,內裡沒有仁慈寬大的心.常見做生意的人,彼此為著利害關係,見面時雖面露笑容,而心裡卻為了利潤.巴不得我成就,他失敗.用詭詐的方法,叫人在背後失敗.商場常有這種事,稀奇的在教會中也有這種事發生.主說:你們要彼此相愛.可惜!在今日的教會中,奸臉惡毒的事太多了,因此,便叫弟兄姊妹受到虧損!
\newline
\newline
弟兄姊妹中,有猜疑,用政治手段來待人,用猜疑的心和奸詐的手法彼此相待,都是不好的.可憐!教會竟陷在這樣的光景中,在不信的人面前告狀.因此,神就嚴嚴的警告我們,祂用的是人與人之間的關係；因為人與人之間的關係,是最親密的關係,這個關係就是夫妻的關係,竟然發生問題,這是一個等待解決的問題,這也是末世的通病.我們好多時候不明白,問題是娶了外邦女子；作者瑪拉基用這個人與人之間,最親密的關係來形容這件事的可怕,竟然存在神的選民當中,應如何解決呢?
\newline
\newline
十五至十六節講到一切的假面孔和狡猾,卻是從人的心發出,所以我們應當謹守自己的心,應該有有新的心志.有很多的事,是從我們罪惡的心發出；求神幫助我們,能約束和謹守自己的心,在弟兄姊妹之間,彼此的關係,應如何相待呢?我有三點心志提供給大家:
\newline
\newline
(一)將誠實心志放在我們當中
\newline
\newline
一個最好的方法,就是用「誠實」彼此相待.有了誠實的心,我們就不會虛假說大話.主說:是就說是,不是就不是,不然就陷在不能自拔的道路上.只要有一次的虛假,就會用各種的虛假來掩蓋本來一個很小的虛假了.說了謊話,就要用多次的謊話來遮掩!所以我們要在教會中,用誠實的心彼此相待.
\newline
\newline
(二)求神叫我們有主耶穌的愛心
\newline
\newline
世人講愛,講得很虛浮,將建立在條件,利益,互相利用的關係上.耶穌基督的愛,只有一個特徵,祂通過自己的犧牲,將快樂賜給人.祂的一生,終日所想的所行的,從不為自己的利益打算.祂在一生中,沒有一件事,或一句話,是以自己為中心的.主透過自己的苦難,犧牲將永久的平安,喜樂,好處給人.我們要效法主的愛心,怎樣捨己去為人.不要個人的利益,而虧損別人.
\newline
\newline
(三)求主賜下寬大的心
\newline
\newline
保羅很讚賞哥林多的教會.因為他們以完全的心境待保羅.他們有愛心和誠實的心,而且還有寬大的心；因此在這方面,我們要有兩個表現.
\newline
\newline
A,無條件的彼此接納.在教會中可能人與人之間的關係有問題,刺眼和不開心的性格,叫你難以忍受的.其實許多衝突,是可以避免的.但我們都是神所造的,應該無條件的接納別人,因看神的恩,祂已無條件接納了我們.試問你肯付上這個代價嗎?
\newline
\newline
B,不記念別人的錯過.要在主裡面彼此饒恕,彼此赦免.無論事情發生在何時何處,都要靠著神的恩典來饒恕.無論事情發生在何時何處,都要靠著神的恩典來饒恕.如果神要以公義來對付我們的罪惡,我們還能站立得住嗎?
\newline
\newline
我們與別人的關係是很重要的,是能影響到神與我們的關係.所以我要求神憐憫我們,幫助我們攪好這方的關係.
\newline
\newline
第二個問題:事奉神是徒然的嗎?
\newline
\newline
經文:瑪三13-四6
\newline
\newline
這個問題與「懷疑神的公義」有關係.一個基督徒常因工作多,又得不到別人的稱讚,賞識而發怨言；以為事奉神是徒然的,這是基督徒事奉上常犯的毛病.
\newline
\newline
我們基督徒有什麼情形下,會陷入這個困境裡呢?有以下四個原因:
\newline
\newline
一,太看重人的稱讚
\newline
\newline
在得不到別人稱讚時,就覺得事奉神是徒然的.
\newline
\newline
二,太過看重工作的果效
\newline
\newline
我們作工,若看不到果效時,就以為事奉神是徒然的.保羅說:這人撒種,那人收割.我們的勞苦不會徒然的.什麼時候收割,我們不知道.我們對主忠心,主一定會記念我們勞苦.
\newline
\newline
三,失去愛神的心
\newline
\newline
我們工作忙碌,慢慢的對神的奉獻,對神的愛,就會隨著時間的磨煉而失去!若非神的愛,我們的事奉是徒然的.
\newline
\newline
四,對於別人的批評太過敏感
\newline
\newline
在工作中,我們要靠神的恩典,不能倚靠自己的聰明.有些時,做傳道人的也很困難；因為努力去做,受人的批評就越大.你做多,別人說你愛出風頭；你做少,別人說你懶惰.你關心人,別人說你要討好人；如果你少關心人,別人會說你沒有愛心.因此,到了這個地步,就灰心起來,覺得對神的事奉是徒然的.
\newline
\newline
我給大家的,是關乎生活,事奉,信仰上爭戰的原則.透過主再來的光景,得勝在信仰上一切的困擾；當主再來時,為祂受苦的人,一定要與祂一同享受榮耀.
\newline
\newline


\newpage

\section{第49屆~港九培靈會~張子華~第一講}
\label{sec:777}
\textbf{順服神權柄的重要性}
\newline
\newline
連結: \href{https://www.hkbibleconference.org/session-message/view/777}{\texttt{https://www.hkbibleconference.org/session-message/view/777}}
\newline
\newline
\hyperref[sec:1058]{< < < PREV SERMON < < <}
~
\hyperref[sec:toc]{[返目錄]}
~
\hyperref[sec:778]{> > > NEXT SERMON > > >}
\newline
\newline
本人帶著感謝與戰兢之心情,接受港九培靈會委辦會之邀請擔任講道多謝神給我再有此機會;但我也感到負擔重大,敢請兄姊代禱,願神使用我站立在弟兄姊妹面前,能夠分享所領受的,把信息講得清楚.當我答應了大會講道時,心中經過長久的掙扎.還未知應該講些什麼?是否應該查考某卷書?或特別專題.感謝神,祂對專心尋求祂旨意的人,必顯明祂的引導,使我明白祂的引導,心裡得平穩安靜.也使我明白祂放在我心中的意念.剛才艾牧師的講道,在我心中發生共鳴,基督教並非信條理論,乃是與神交通團契的生活;它的特徵是和平或和諧是從神開始,又將和諧生活,帶進基督徒生活的每個小圈子裡以弗所書記載,那些本來與神為敵的人,神藉著祂兒子為我們釘在十字架上;就將那隔絕的牆拆毀,使人與神和好並在這世上,成為傳平信息的人我打算把十天講道作如下分配,前三天講基督徒不能過和諧生活的三方面原因:
\newline
\newline
(1)不承認神的主權.
\newline
\newline
(2)不在神的旨意下過生活 - 不知應用聖經,在每一生活細則上.
\newline
\newline
(3)失掉對神起初的愛心.
\newline
\newline
這是頭三天的講論然後進入第二部份,如何對人過和諧的生活:.消極方面,怎樣對付成見,仇恨;如何去接納每一弟兄姊妹積極方面,如何用基督的愛去愛人,這是每一教會最強有力的見證主耶穌在(約十三34-35)說出一條新的命令乃是彼此相愛祂說:.「你們若有彼此相愛的心,眾人因此就認出你們是我的門徒了」我很奇怪,為什麼不說當讀經,禱告,或作其他敬虔的事;而單說我們有彼此相愛的心,眾人就認出你們是我們門徒了可見這.是最強有力的見證,足以叫人認識主.第三方面將講到與自己和諧.你知否你對自己經歷了幾許的爭戰.如何克服自卑,自憐,接納自己為神所創造的工作.如何借聖經勝過種種弱點,而重過和諧生活這部份也包括三點信息:.家庭之和諧生活,從男女互相認識開始,夫婦,婆媳之間願神憐憫,在祂語之亮光中,帶進基督徒生活的每個小圈子,享受十字架賜下之平安福樂.
\newline
\newline
壹,與神和好生活之障礙 - 不接受神在生命中絕對的主權(詩一一○1-3)
\newline
\newline
我作基督徒生命的經歷,相信也可代表大部份基督徒.這障礙就是不瞭解到每位重生基督徒,都有個不能代替的經驗──享受「主」在他身上絕對的權柄.我們許多時候會懷疑:人說信主之後,必得著平安,得享安息；但我卻沒有平安,且充滿了掙扎,請問弟兄姊妹!你的生活是否靜止,而毫無掙扎呢?我們應先瞭解與神和諧關係的絕對性.我們追求靈命長進和教會復興；曾嘗試用種種方法和舉行不同的聚會,卻得不到預期效果.詩一一○篇前三節,啟示真復興與真和諧的途徑,就是接受主在生命中,掌管絕對的主權.如此才可享受生命成熟帶來的福份.
\newline
\newline
第一節,大衛受膏為王時,寫成此詩:父神對聖子說:「你坐在我右邊,等我使禰仇敵作禰腳凳.」人對這詩可能有種種解釋,但明顯的意思,應論到主以大能勝過仇敵.對於成熟的基督徒,耶穌「主」的信息,從創世記到啟示錄都明白提出.除非由主掌權,否則談不上屬靈的生活.這節又論到,耶穌是世界的統治者.有一首聖詩內容:「日月所照,萬國萬邦,耶穌必為,統治君王；南北西東,主治普及,千秋萬古,悠久無疆.不絕禱聲,向祂奉獻,無量頌歌,成祂冠冕；到處聖名,好比馨香,隨同神祭,上升於天.」無論人承認與否,第一節經文所啟示的,耶穌基督乃是君王.
\newline
\newline
「耶和華必使禰從錫安伸出能力的杖來,禰要在禰仇敵中掌權.」這節經文顯示祂與信徒有個人的關係.第一節啟示祂是世界的統治者,第二節卻顯出祂是一位掌權者.「能力的杖」代表權柄.掌權的意思是用能力管理我的生命.每個基督徒不會否定耶穌是統治者,正如萬物是藉著祂所造.但要承認耶穌是掌權者,卻會有一番掙扎.因接受主在生命中掌權,就需要放下許多喜愛的,而拿起許多不願意的.許多基督徒在這取捨之間,發生許多掙扎.雖知道神有權柄,但不願使用在祂兒女身上.弟兄姊妹,若我們不接受讓神的權柄施用在身上,則可斷然地說,必不能過和諧的生活.反抗神的主權者,永不能享受和諧.最近我有機會往英國領會,有機會往白金漢宮觀光,但見守衛軍兵制服莊嚴,儼有封建景象.但熟悉英國政制的人,莫不知英國是采虛君制度；政權操於首相和國會之手.所有大小政事都由首相和國會決定,然後呈到女皇面前請其簽署.幾多基督徒也是如此,將主耶穌當英女皇看待,我們已將大小事情決定好了,然後交給祂,由祂簽字認可.誰操你生命的權柄呢?是你自己；還是耶穌；要記得一個重要的原則!神要在我們生命的每一部份掌權.(不是某一部份.)我們當承認祂在我們身上主權的絕對性.許多人對神的話,並非絕對地順服,或憑己意妄加衡量.讓我重複地申明一次:當我們將生命讓耶穌當權時,祂的權柄要在我生命的每一部份彰顯,是絕對的不容其他解釋.對於權柄祇有兩樣結果:一是反抗,一是順服.對於主的權柄,聖經的話語,許多時我們是按己意選擇,失去了絕對的價值.
\newline
\newline
舉例說明:我的大兒喜食朱古力糖,若一盒之中含許多種朱古力.他唯獨愛吃內有果仁的一種,不作他選.某次,我和妻子出街,朋友送一盒朱古力糖來,他就打開享受,吃了一些,我回來一看,就知他揀過所有的糖,何以見得?因為每一粒都有被捏過的指印,可見他捏到軟的便丟下,捏到硬的就吃了,我倆不禁相對而笑.多少時候,我們對神的權柄透過祂的話語,何嘗也不採這態度.對待神的誡命,好像孩子揀朱古力糖.若對我自由有影響,與我思想有衝突,叫我吃虧,我就把它捏過後,又復放下.如捏到一處誡命,是帶有神的豐盛,赦罪之恩,就歡喜接受.否則,就放下.幾時,我們在生活上,若以此眼光,態度對神,我們必然失敗!因為神要在你我生命中掌權,每一部份,無論大事小事,有絕對的有權柄.
\newline
\newline
講到神的絕對性,使我想起一位姊妹過去的事.她已屆三十年華,有結婚的需要.在某一機會之下,認識一位擁有博士學位而樣貌英俊的男子,但可惜他對基督教抱激烈反對態度.他們發展至相愛階段.我不知用了多少時間,規勸,討論,反覆申論到神對於信徒婚姻的定意.但她卻認為神的話,對她無絕對性.例如經上說:「信與不信不能同負一軛」.英文用 May Not 而非 Must,她據此認為不肯定的憑證,若這理由成立,則整本聖經的絕對性可被推翻.為著遷就她本身的意願,不削減神話語的絕對權柄.弟兄姊妹!今日信息中不能忘記的要點,是神要在你生命中掌權,普及每一項細則,具有絕對性.請問神是否在你身上如此掌權；在家庭,婚姻,生意,教會的事奉,靈修生活,奉獻的態度,都要清楚交代.基督徒不能用「照我的看法」,或「也許另有解釋」,這些模稜兩可,取巧的話,在神權柄中都用不著,若然神在你身上,生活細則之中無絕對權柄,就不能享有生活上一個大特徵──和諧.「大能的杖」,「仇敵中掌權」；在生活上與神的旨意,聖經記載相違背的,就是神的仇敵,都要打碎.這是重要的態度和方法.
\newline
\newline
第三節,當神掌權,就有三種表現:
\newline
\newline
1.「禰的民要以聖潔的妝飾為衣」.
\newline
\newline
2.「甘心犧牲自己」.
\newline
\newline
3.「禰的民多如清晨的甘露」.
\newline
\newline
願以這些作為一面鏡子,或用為試題,細察在我們生活的細則中,是否讓神掌權?
\newline
\newline
1.以聖經為妝飾,是以神的旨意,誡命,教導,帶領,作為生活最高標準.神不允許縱容任何罪惡.在你生活中有何不應有的事物,明知而不曾對付?在生意場中,在學生生活的範疇裡,或作家庭主婦的,有否隱而未現的罪,或不法之行為.愛慕聖潔是追求最高生活標準,主掌權的第一種流露.
\newline
\newline
2.甘心犧牲自己.神很明白要在信徒身上掌權,但有極大的敵人,就是自己.所以祂呼召:「人若跟從我,就當捨己,天天背著十字架來跟從我.」在生活上對聖經的態度,若你不接受神的最高標準,就是以己意來代替神.在生活的原則上,一方是神,一方是己.若讓神掌權的話,就祇有神是對的,而我是錯的.我很喜歡約翰福音四章一段話:「你們豈不是說:到收割的時候,還有四個月麼?我告訴你們,舉目向田觀看,莊稼已經熟了,可以收割了.」清楚有兩種說法,一是我們說的,一是主說的.顯出一個非常重要的生活原則,讓我下結論,在每個生活圈子裡,對事情的看法,有你的意思和神的意思相對立.但神的意思永遠是對,而你的意思永遠是錯.
\newline
\newline
3.你的民多如清晨的甘露,或作你少年時光耀如清晨的甘露.基督徒若接受神在其生命的每一部份掌權,定必感到人生無窮高度的享受.清晨的甘露是指在靈裡更新,在主裡活潑,生活和諧.好比清晨的甘露,每早晨滋潤我們的心靈.若你覺得基督徒的生命是古板的,掙扎的,失去興趣,缺乏興奮；滿是痛苦勞碌重擔.很可能,你未享有神的權柄.因神最瞭解你我.神實施祂的權柄,將誡命賜給我們,非為把擔子加在我們身上,為使我們的生活新鮮,愉快,和諧.人必須由此立場去瞭解神的愛和計劃.例如:有人對參加崇拜覺不耐煩,聽講道如背重擔.但神知道人六日作工,一天安息守主日是最合宜的.若用神的愛,存順服的心守主日,自然覺得甘甜.心被恩感,靈內復興.
\newline
\newline
弟兄姊妹!神在你身上有無絕對權柄?過去生活,由你作主,還是由神控制.巴不得神憐憫,進入到完全的地步,祗要你作順命的兒女.
\newline
\newline


\newpage

\section{第49屆~港九培靈會~張子華~第二講}
\label{sec:778}
\textbf{活在神話語旨意中的和諧生活}
\newline
\newline
連結: \href{https://www.hkbibleconference.org/session-message/view/778}{\texttt{https://www.hkbibleconference.org/session-message/view/778}}
\newline
\newline
\hyperref[sec:777]{< < < PREV SERMON < < <}
~
\hyperref[sec:toc]{[返目錄]}
~
\hyperref[sec:779]{> > > NEXT SERMON > > >}
\newline
\newline
也許我們中間有些弟兄姊妹昨天沒有來,所以我要將全部講道簡略題說.我經過幾個月來的祈禱等候,從神那裡領受信息,也認明是教會所需要的.在研經會中,艾牧師多次提到,我們所信的是一個生命;也要在生活中見到它的真面目基督徒生活有一個大特徵,就是和平(或和諧)既藉著耶穌並祂釘十字架得與神和好,和諧就進到人生的每一關係上可惜在神家裡,少看到和平的事實弟兄姊妹之間,家人之間,彼此不和睦;!又往往不接受神創迼的事實,產生自卑,自憐,這都是不和諧實況所以我預備與眾人分享的信息:.前三日講基督徒怎樣保守同神有和諧的關係昨日講過順服神權柄的絕對性,它觸摸到生活的每一部份.基督徒由於順服權柄,就產生和諧.基督徒的自由,也由此產生.我前在菲律賓講道,有一位姊妹帶著眼淚對我傾訴家:.原來她父親很愛主,也盼望兒子能愛主過敬虔生活但兒子卻不願相從,演變至勢成水火,家無寧日;甚至惡劣情況到持刀相向及至後來接到她的來信,經過祈禱和輔導,兒子瞭解一個重要真理. - 就是要服從父親的權柄信上提到自此家庭一改舊觀,充滿了和諧其實和諧的生活,何等需要應用到基督徒的生活上.;你若反抗神,你就無法享受到神所賜的和諧與甘甜唯有停止對抗,接受神掌權,才會帶來生命的意義與甘美今天所講:應該活在神話語,旨意中;使你過著成聖,和平,喜樂的生活明日將提到與神和諧的一大攔阻 - 「失去了起初的愛心」講完了對神方面三講之後,進而談論兩方面:一方面與人和好,怎樣靠神恩典來對付仇恨與成見,用神的愛去愛別人;因為凡信主的人,都同在一團契中然後再講到與自己和諧,這聽來似乎頗新鮮;事實上,我們對自己有許多不滿;從容貌,健康,學問,......種種欠缺上,產生自卑,自憐,我們必須知道如何去克服這種種困難,進而與自己和好.最後講家庭的和諧,夫婦,父母,子女之間;種種磨擦糾紛,如何避免,進而建立美滿和諧的家庭.
\newline
\newline
活在神話語,旨意中,過和諧生活(詩十九7-14;賽五十五8-9,太二13;腓二5.)
\newline
\newline
我曾見過一段記載:關乎一位美國太空人抒寫他的感受.當太空船將進入太空之後,發覺有一事實,天空所有之寧靜和諧,非地上所能覓得.有位牧師對此作出評論,何故太空會呈現寂靜,和諧,有條不紊的系統；原因是在神所指定之軌道運行.反觀世界,充滿了戰爭,仇恨,不平,不公之事；主要因神所造之人,離開神所安排的軌道；過著以己為中心的生活,走自己的道路.當基督徒的生活,越出神定下的範疇時,我們就不能享受神的平安!也不能發揮神在我們身上的功能.發明火車的人,使車行在軌道上,就能行駛順利；也能達到其本身的作用,若離了軌道就不能動彈,甚至帶來嚴重災害.
\newline
\newline
各位!可以清楚看見,神拯救了我們,也有一定的生活規則給祂的每一兒女.祂所給的規則,是為了他們的幸福,而非重擔；與未信前的生活方式,傾向,絕不相同.以色列民本來在埃及為奴,神用大能將他們救出.未進入迦南地,先來到西乃山下,神將另一個生活的模式賜給他們；要他們忘記之前為奴時的生活方式.照神所定之方法,進入迦南享受安息,神把甘甜流奶與蜜之地賜給他們.所以第一,每人都要知神透過祂的話,給我們與世人不同的生活方式,來作信徒生命的表現.我常感到教會太多理論家,他們知道許多教條,或背熟許多聖經.但早上講員提醒這並非真基督教,我們要從生命中表現生活的新方式.
\newline
\newline
(一)每一信徒,神有特別的方式,在他的生命之中.
\newline
\newline
(二)神話語(旨意)為給人生有幸福的享受.讀(詩十九):「耶和華的律法全備,能甦醒人心.」甦醒是復甦之意.「耶和華的法度確定,能使愚人有智慧.......」我非常喜愛此詩,我要向大家作一見證:神的話語是最高度的喜樂和享受.很可惜!弟兄姊妹常以神的旨意難守,視為重擔.為何我不能惱恨別人,要赦免仇敵?為何我不能在夫婦之外,另有男女的關係?若你還有這疑問,我真要為你擔心.因你仍未能體會到,神所賜的法度,並非為叫我們背重擔,乃是叫我們得到幸福和享受.我的大兒子十四歲,今年讀高二.有一次他對我和師母說:他是如何愛我們.我就問他:父母曾打罵過你,為甚麼還愛父母呢?他答道:我知道這是為我好的緣故.他的話激發我對神感謝的心.讓我想到,神所要我們作的,或禁止我們的；是為了我們的好處,我們是否會改變對神旨意的態度?
\newline
\newline
(三)神的意念高過我們的意念,神的道路高過我們的道路.(賽五五8-9)所載,是每位基督徒追求生命長進的動機.我們是罪人蒙恩得救,在人生路上瞭解一項事實,人的意念常與神的不同.但結論是:神的永遠比我的好,因此,我要以此作為我人生的目標:以基督的心,為我心；祂的意念,為我的意念.有生之日,不斷追求,叫我意念更似主耶穌.這樣的追求長進,是基督徒最興奮的事,因為到了一地步,又有另一境界展開.叫我們的信心永不停止.
\newline
\newline
不久之前,我們一家人回港度假,想到我在美牧會時,曾領一位從香港赴美的姊妹歸主.趁我回港時,請求父親接待我；因此,我就被邀請往參觀那富豪的別墅.照主人的安排引導,我看過那裡的五個客廳.每間客廳的陳設按各地而風格不同.由美國式,進而北歐式,再進中國的古典方式；一廳比一廳更為華麗,確是美不勝收.一路行來,我心中默默地感謝神.這是基督徒追求完美的一幅寫照,到達一個美的地步,再追尋另一更美的,永不停止.不要自我陶醉,與那些冷淡比較,自己以為不錯.應該追求更似耶穌基督,更合乎神的旨意.
\newline
\newline
(四)神的旨意祗有一個途徑,可以明白在青年夏令會中,或特別座談會中,我們常有難題解答,多數的人問我「怎樣能知道神的旨意?」像這一類問題.容我引用一位教授的話:「若你每日過著與神同行的生活,你永不會不明白神的旨意.」若你平常少有讀經,祈禱,與神相交.到有事時,怎能明白神的旨意.若你祈禱,讀經,過著健全正常的屬靈生活；神應許祂的旨意必向祂兒女顯明.因此,我要與各位分享一些個人經驗.過去我曾在港某一神學院攻讀.這學院受巴特神學之影響極大,那時,我因成績優良,院方擬派我往德國深造,所以我要在德文方面用功,預備前往.這時我與神有深切的相交,每早五時起床個人靈修,神就帶領我的道路,祂帶領我認識一位姊妹,繼而彼此相愛.就使我原來想往德國讀神學的準備,改往美國哥倫比亞聖經學院就讀.由於認識我今日的妻子,就改變了我的道路,神學的立場和我的事奉.因此結論是:你若真與神同行,就永不會不明白神的旨意.故此問題不在於去著意追求神的旨意,而要建立你與神的健全關係,神的恩手必攜帶你的的路程.
\newline
\newline
(五)神話語(旨意),要成為基督徒生活最高原則的標準.當人過基督徒生活時,神的話語是最高標準,也是最高享受.分兩方面來說:其一,神許多說話,是直接教訓人如何生活.例如:不可姦淫,不可偷盜,要愛人如己......等.清楚顯示出,神的要求為何,不能有其他解釋.另一方面,有關生活的興趣和習慣,聖經上並無明白指示可否?就要藉著聖靈活潑的帶領,將原則應用於生活的每一部份.舉例說明:當我與女友相愛時,常喜歡駕車往人煙罕見之處,互訴衷曲.但我兩人都是信主的.當讀到主禱文:「不叫我們遇見試探.」這經文就引我們去思想,是否我們把自己放在試探之中.因此神的話語又成為我們生活上的最高原則,保守自己不陷在試探之中.又例如:可否吸煙也是人常發的問題.我並非在此講律法,不可吸煙.祇是聖經有話說:「豈不知你們的身子,就是聖靈的殿麼?」若有人毀壞神的殿,神就要毀壞那人.照今天科學的證明,吸煙對身體有害是毋庸置疑的.所以我可以直接了當的說:不可吸煙,因為吸煙是殘害自己的身體,吸煙就是毀壞神的殿.其他如在飲食上,睡眠上,沒有好好照顧自己,影響健康精神,同樣在這原則上,不是神所喜悅的.又例如說:「凡事謝恩.」聖靈引導我活潑地應用這原則.凡不能向神感恩的事,我不去作.我在華盛頓教會工作時,知道幾位會友喜歡賭錢.常聚在某個家庭打麻雀,我知道這事,就找機會,突然到他家探訪.當開門時,他們見是牧師到訪就大衛緊張.但來不及擋駕,我已進到他們當中,裝作若無其事,並謂想學打麻雀.他們正不知是真是假.我說:讓我們先禱告,然後請你們教我打牌.其中有人說:不能祈禱.我於是為他們讀出「凡事謝恩」這段經文.感謝神使他們從此脫離這惡習.
\newline
\newline
住在聖靈裡面,與神同行,聖靈能透過神的話語產生屬靈的原則.(羅十四23):「凡不出於信心的都是罪.」無論作一事或不作某事,如果不是坦然無懼,而是猶疑不決,那就應該小心.要在每件所行的事上,憑著信心而行,就是討神喜悅的,若然,我們基督徒的生活.就會像在愉快的搖籃中.我們的生活充滿意義,甘甜,幸福和安息.你願意活在神話語的亮光中嗎?祂的意念永遠高過自己的意念.若與神同行就永不致迷失,而能夠過成聖的生活.
\newline
\newline


\newpage

\section{第49屆~港九培靈會~張子華~第三講}
\label{sec:779}
\textbf{與神和諧關係破壞的原因－ 失去起初的愛心}
\newline
\newline
連結: \href{https://www.hkbibleconference.org/session-message/view/779}{\texttt{https://www.hkbibleconference.org/session-message/view/779}}
\newline
\newline
\hyperref[sec:778]{< < < PREV SERMON < < <}
~
\hyperref[sec:toc]{[返目錄]}
~
\hyperref[sec:780]{> > > NEXT SERMON > > >}
\newline
\newline
啟示錄 2:1-2:4
\newline
\begin{longtable}{cl}
\hline
\hline
章節 & 經文 (和合本修訂版)\\
\hline
2:1 & \begin{tabularx}{0.7\textwidth}{X} 「你要寫信給以弗所教會的使者，說：『那右手拿著七顆星，在七個金燈臺中間行走的這樣說： \end{tabularx} \\ \\ \relax
2:2 & \begin{tabularx}{0.7\textwidth}{X} 我知道你的行為、勞碌、忍耐，也知道你不容忍惡人。你也曾察驗那自稱為使徒卻不是使徒的，看出他們是假的。 \end{tabularx} \\ \\ \relax
2:3 & \begin{tabularx}{0.7\textwidth}{X} 你能忍耐，曾為我的名勞苦而不困倦。 \end{tabularx} \\ \\ \relax
2:4 & \begin{tabularx}{0.7\textwidth}{X} 然而，有一件事我要責備你，就是你把起初的愛心拋棄了。 \end{tabularx} \\ \\
[1ex]
\hline
\hline
\end{longtable}
過去兩日提醒自己與各位,雖然借十字架功勞,從作神的仇敵而成神的朋友,已經與神和好了.但我們在靈命中多方顯出,與神關系仍缺乏和諧.我要提出三方面,來說明如何保守與神和諧之關系:
\newline
\newline
1.順服神絕對的主權,在生活的每一細則上;生活若出現反抗神的權柄,與神和諧就受到損傷2.借著聖靈活潑的能力,將神話語應用在每個生活圈子里,活在神完美旨意中建立與神同行的健全關系;愛神的人永不會迷失神的旨意的今日講第三方面:失去和諧的另一原因,對神失去起初的愛心;明日進入對人的關系,如何克服論斷,憎恨等困難;後日講積極地用神的愛,來建立與人和平的關系;禮拜六講接受自己,脫離自卑,自憐;接受神的創造和神允許我們的遭遇,享受內心的平安;禮拜一至禮拜三講家庭,夫婦之間,父子之間的和諧關系.
\newline
\newline
我們與神的關係,建立在愛的基礎上.不是我們先愛神,是神先愛我們,我們的心受到激勵,才能愛神.兩者之間,需如此保持健全的關係.但事實上我們不能不承認,在經歷上會遇到各種原因,而失去對神的愛.由此,同神的和諧關係就出問題,今日一同思想,為何失去了對神起初的愛.
\newline
\newline
(一)為何失去起初的愛
\newline
\newline
是否因為不順服?事實上並不盡然；許多人繼續勉強順服,甚或真順服,也會失掉初愛所帶來之興奮和活力.因不讀聖經嗎?可能,但法利賽人何嘗不多讀聖經.缺少禱告嗎?一個會禱告的人,亦會在愛中失落,靈性冷淡.不為主作見證嗎?啟示錄記載以弗所教會曾勞碌作工,並不怠倦；但耶穌基督透過聖靈,指出它是失去了起初的愛心.弟兄姊妹,失去對神的初愛,其中最大,主要的原因,而且是唯一的,乃是失去了需要神的感覺.請不要忘記一件事實,當聖靈在人心中工作,人受到光照時；人就會見到自己光景可憐,滿了罪污,沉淪絕望,痛苦孤單；何等需要神的愛,需要拯救,因而接受救恩.在一剎那間,在聖靈光照之下,我們感到極度需要神,而接受救恩.這需要的感覺,是一個原動力.但慢慢過去,在物質上,靈命上,因蒙神賜福,對神的需要,反而下降了.這就是對神失掉起初的愛心的緣由.
\newline
\newline
當我們讀到以色列民族歷史的時候,就發覺到,神用奇妙的方法教導祂的百姓.他們脫離在埃及為奴的地位,過曠野生活.每一步路程,神都為他們預先準備,供給他們食物和水.神不用他們耕種土地,神卻賜他們嗎哪為食物.但嗎哪是不能收藏起來的,隔天,就會變壞不能食；神這樣做,為了使百姓每步需要神.
\newline
\newline
以色列民的旅程,到達迦南地邊境.神開始預言祂的百姓會忘記祂,失去倚靠祂的心；初愛的重點──需要神的感覺.這是(申八11-14)所記載的.後來神的預言果然應驗,到他們進到流奶與蜜之地,蒙神賜福,財富增加,就逐漸忘記神.今日新約的信徒,同樣有這危險.當神賜福滿足了我們基督徒的生活時,我們就不會感到沒有神不能活下去的需要了.這感覺消失了,那就是失了起初的愛了.如果你覺得錢財的能力,足以代替神；也就等如失去了對神的初愛.
\newline
\newline
當我讀神學時,曾經參加過一項試驗,如今我把這試驗提出來作為弟兄姊妹的借鏡.
\newline
\newline
1.當你愛慕神的心,比不上任何人或事物時,可能你已失去對神的初愛.你生活中有沒有愛任何事物,人,比愛神更多?看某一件事比神更重要；為它勞碌,佔有的思想更多?若然,則你已失去對神起初的愛心了.
\newline
\newline
2.在自己的心靈深處,沒有極度渴望藉著禱告與神交往,那就是失去對神的初愛了.(換言之,你的喜樂與甘甜,不再建立在神的話語中.)
\newline
\newline
3.在娛樂或空閒時間,不想念神,就是失去祂起初的愛.
\newline
\newline
4.若常藉口；神知道我的軟弱,了解我的難處.我不過是人,用許多推諉的話,來放縱自己；故意作神所不喜悅的事,就是失去對祂起初的愛.
\newline
\newline
5.若我在金錢奉獻的態度上,不慷慨,不甘心,不盡力奉獻,卻沒有學習超過能力的獻上,就失去了初愛.許多基督徒自以為實行十一奉獻就夠.然而,神卻要看我們奉獻的態度如何?神若對你有所要求,你是否甘心,樂意,滿足神的心?若不到此地步,就是失去了初愛.
\newline
\newline
6.若你停止對弟兄的關懷,沒有真誠愛弟兄；看見弟兄姊妹窮乏,受苦,卻無動於中.別人的重擔,苦難已不能觸動你的內心；那就是失了起初的愛.
\newline
\newline
7.如果我對遵守神的誡命,不視為對神愛的反應；而看作重擔,並非是幸福,快樂,那就是失去了初愛.
\newline
\newline
8.若我將自己放在愛世界事上,正如(約壹二15):「人若愛世界,愛父的心就不在他裡面了.」就證明失去起初的愛.
\newline
\newline
9.當我堅決拒絕放棄某些絆跌人的喜好和活動,就是失去了對神的初愛.
\newline
\newline
10.若心中不能絕對澈底去饒恕某一個人,這都使我們失去起初的愛.以上十項試驗,當各位回家後,可以再詳細拿來對照自己的屬靈生活；就可以找出你對主初愛的光景如何.在每一項試驗,都可以從聖經中得到神話語為根據的.
\newline
\newline
(二)如何堅守起初的愛
\newline
\newline
首先我提出一個最常見的錯誤觀念.許多時候信徒認為,靈命越長進時,就越屬靈,就越接近於完全.但在屬靈的實際經驗上,卻完全兩樣.原來當信徒越追求長進時,就越缺乏,越需要主.保羅在(林前十五:9)曾自稱:「我原是使徒中最小的,不配稱為使徒.」當他靈命最剛強老練時,越見到自己的缺乏和需要主,這就是保守對神初愛的秘訣.
\newline
\newline
耶穌說:「靈裡貧窮的人有福了,因為天國是他們的.」現今我提供幾個步驟,或者說:新的觀念來得回,保守,這起初的愛.
\newline
\newline
1.從新的禱告經驗裡保守起初的愛　我要在禱告上介紹給你們一個新的感受.當我發現這新經驗時,就感覺到以前我禱告,何等缺乏.在禱告中有讚美,代求,祈禱,每部份都重新有需要主的感覺.真誠地向主說:「我需要你!」靈裡有這渴慕,你就會得回起初的愛.當我向神讚美,因祂所作的,是我所不能作的.當將一個「需要」擺在神面前時,我必需了解,我何等地需要主.若我沒有這觀念,就是失掉起初的愛.法利賽人和稅吏兩種禱告的分別,就是有沒有需要主的感覺.法利賽人自誇:我如何如何,我不像別人.......許多基督徒也是如此,沒有需要主的感覺；甚至你為別人禱告時,也不覺得有這需要.舉一眼前之例:當我們每餐謝飯時,雖然口裡說:「主阿,感謝你賜我飲食.」但何嘗在心中有這意念?(主若不賜給,我就不能有此飲食.)因為在現實上,我們都不愁衣食.但無論我們有幾許貲財,倘若我們知道需要主,才能作這感謝──我們需要主每天賜我飲食,這是一個非常重要的祈禱新觀念.
\newline
\newline
2.新的讀經經驗保守初愛　一般人以為讀經為使生命微長進,是不錯的.但我要介紹一新觀念,讀經時,要用神的話,面對我屬靈的貧窮光景,滿足我屬靈需要.因此,真正讀經要有此感覺.多少時,我們讀經只成了例定的,責任的,讀過便了事.不會面對神的話語,如飢如渴地,讓祂光照見你的需要,缺陷,敗壞.讓祂教導,責備,達到祂的標準.應該在每日讀經中了解,需要主的話來提醒,責備,滿足一切屬靈的需要.
\newline
\newline
3.從新奉獻觀念得回初愛　現在不是講十分之一或十分之五的問題.乃是提到一個基本的觀念.當你將錢投入奉獻袋時,應有一觀念,你一切所有物質財富都是屬神的.神若要你把全部財產為祂所用,你也能順服不予保留.在每次奉獻上你要有此態度,否則,你便不是個需要神的人.特別對一些在錢財上無欠缺的人,你要知道需要神.如果你認為有權支配一切,這是錯誤的.當你感覺錢財是你用本領賺來的,你有權來支配使用它,那就錯了.如果神要對付時,恐怕你擔當不起.要開始你對金錢,對奉獻的新觀念,使你能保守那起初的愛.
\newline
\newline
4.從新的禁食觀念保守初愛　禁食是最有效的提醒,神是我們每日所需之糧.從禁食經驗中,你會深切地認識飢餓之苦,從而體會到神對我需要的供給.
\newline
\newline
5.從新事奉觀念保守初愛　許多時我們以為事奉,是神要我勞碌.如今見到,藉著供給別人屬靈需要,體會到自己靈裡的缺乏,一個真事奉的人,一定時常跪在神前說:「神阿,你若不滿足我的需要,我就不能供應別人,不過是一廢壞的器皿而已.」若你自以為能夠,不需要神；那種事奉,不能想像了!但我有機會事奉時,必須向主呼求,若祂不滿足我,就不能供應別人所需.靠著祂,在講台上釋放祂的信息,來幫助眾弟兄姊妹!就是從這事奉觀念上,體會對神的需要更為迫切.
\newline
\newline
以上最後所講五方面,本是很顯淺的道理.但可惜許多人,在這幾方面都失去需要神的感覺,漸成麻木.我們要保守這需要神的感覺.無論在任何環境之中；不論是富有,貧窮,有學問與否?求神教我們用這種種方法,保守我們對祂有強烈的渴想和需要.常常如此,就能得回,保守起初愛神的心.
\newline
\newline


\newpage

\section{第49屆~港九培靈會~張子華~第四講}
\label{sec:780}
\textbf{與神和諧的關係－ 順服神}
\newline
\newline
連結: \href{https://www.hkbibleconference.org/session-message/view/780}{\texttt{https://www.hkbibleconference.org/session-message/view/780}}
\newline
\newline
\hyperref[sec:779]{< < < PREV SERMON < < <}
~
\hyperref[sec:toc]{[返目錄]}
~
\hyperref[sec:781]{> > > NEXT SERMON > > >}
\newline
\newline
撒母耳記上 15:1-15:23
\newline
\begin{longtable}{cl}
\hline
\hline
章節 & 經文 (和合本修訂版)\\
\hline
15:1 & \begin{tabularx}{0.7\textwidth}{X} 撒母耳對掃羅說：「耶和華差遣我膏你為王，治理他的百姓以色列，現在你要聽從耶和華的話。 \end{tabularx} \\ \\ \relax
15:2 & \begin{tabularx}{0.7\textwidth}{X} 萬軍之耶和華如此說：『以色列人從埃及上來的時候，在路上亞瑪力人怎樣待他們，怎樣抵擋他們，我都要懲罰。 \end{tabularx} \\ \\ \relax
15:3 & \begin{tabularx}{0.7\textwidth}{X} 現在你要去攻打亞瑪力人，滅盡他們所有的，不可憐惜他們，將男女、孩童、吃奶的，以及牛、羊、駱駝和驢全都殺死。』」 \end{tabularx} \\ \\ \relax
15:4 & \begin{tabularx}{0.7\textwidth}{X} 於是掃羅在提拉因召集百姓，數點他們，共有二十萬步兵和一萬猶大人。 \end{tabularx} \\ \\ \relax
15:5 & \begin{tabularx}{0.7\textwidth}{X} 掃羅到了亞瑪力的京城，在谷中設下埋伏。 \end{tabularx} \\ \\ \relax
15:6 & \begin{tabularx}{0.7\textwidth}{X} 掃羅對基尼人說：「你們離開亞瑪力人下去吧，免得我把你們和亞瑪力人一同殺滅，因為以色列眾人從埃及上來的時候，你們曾恩待他們。」於是基尼人離開了亞瑪力人。 \end{tabularx} \\ \\ \relax
15:7 & \begin{tabularx}{0.7\textwidth}{X} 掃羅攻打亞瑪力人，從哈腓拉直到埃及東邊的書珥， \end{tabularx} \\ \\ \relax
15:8 & \begin{tabularx}{0.7\textwidth}{X} 生擒了亞瑪力王亞甲，用刀殺盡亞瑪力的眾百姓。 \end{tabularx} \\ \\ \relax
15:9 & \begin{tabularx}{0.7\textwidth}{X} 掃羅和百姓卻憐惜亞甲，愛惜上好的牛、羊、牛犢、羔羊，以及一切美物，不肯滅絕。但是凡看不上眼和沒有價值的，他們盡都殺了。 \end{tabularx} \\ \\ \relax
15:10 & \begin{tabularx}{0.7\textwidth}{X} 耶和華的話臨到撒母耳說： \end{tabularx} \\ \\ \relax
15:11 & \begin{tabularx}{0.7\textwidth}{X} 「我立掃羅為王，我感到遺憾，因為他轉去不跟從我，不遵守我的命令。」撒母耳就很生氣，終夜哀求耶和華。 \end{tabularx} \\ \\ \relax
15:12 & \begin{tabularx}{0.7\textwidth}{X} 撒母耳清早起來，去見掃羅。有人告訴撒母耳說：「掃羅到了迦密，看哪，他在那裡為自己立了紀念碑，又轉身下到吉甲。」 \end{tabularx} \\ \\ \relax
15:13 & \begin{tabularx}{0.7\textwidth}{X} 撒母耳到了掃羅那裡，掃羅對他說：「願耶和華賜福給你，耶和華的命令我已遵守了。」 \end{tabularx} \\ \\ \relax
15:14 & \begin{tabularx}{0.7\textwidth}{X} 撒母耳說：「我耳中聽見有羊叫、牛鳴的聲音，又是甚麼呢？」 \end{tabularx} \\ \\ \relax
15:15 & \begin{tabularx}{0.7\textwidth}{X} 掃羅說：「這是百姓從亞瑪力人那裡帶來的，因為他們愛惜上好的牛羊，要獻給耶和華－你的神。其餘的，我們都滅盡了。」 \end{tabularx} \\ \\ \relax
15:16 & \begin{tabularx}{0.7\textwidth}{X} 撒母耳對掃羅說：「住口吧！等我把耶和華昨夜向我所說的話告訴你。」掃羅說：「請說。」 \end{tabularx} \\ \\ \relax
15:17 & \begin{tabularx}{0.7\textwidth}{X} 撒母耳說：「你雖然看自己為小，你豈不是作了以色列諸支派的元首嗎？耶和華膏你作了以色列的王。 \end{tabularx} \\ \\ \relax
15:18 & \begin{tabularx}{0.7\textwidth}{X} 耶和華差遣你，吩咐你說：『你去除滅那些犯罪的亞瑪力人，攻打他們，直到把他們完全滅盡。』 \end{tabularx} \\ \\ \relax
15:19 & \begin{tabularx}{0.7\textwidth}{X} 你為何沒有聽從耶和華的話呢？你為何急著撲向掠物，行耶和華眼中看為惡的事呢？」 \end{tabularx} \\ \\ \relax
15:20 & \begin{tabularx}{0.7\textwidth}{X} 掃羅對撒母耳說：「我聽從了耶和華的話，行了耶和華派我行的路，擒了亞瑪力王亞甲來，滅盡了亞瑪力人。 \end{tabularx} \\ \\ \relax
15:21 & \begin{tabularx}{0.7\textwidth}{X} 百姓卻從掠物中取了牛羊，是當滅之物中最好的，要在吉甲獻給耶和華－你的神。」 \end{tabularx} \\ \\ \relax
15:22 & \begin{tabularx}{0.7\textwidth}{X} 撒母耳說：「耶和華喜愛燔祭和祭物，豈如喜愛人聽從他的話呢？看哪，聽命勝於獻祭，順從勝於公羊的脂肪。 \end{tabularx} \\ \\ \relax
15:23 & \begin{tabularx}{0.7\textwidth}{X} 悖逆與占卜的罪相等，頑梗與拜偶像的罪孽相同。因為你厭棄耶和華的命令，耶和華也厭棄你作王。」 \end{tabularx} \\ \\
[1ex]
\hline
\hline
\end{longtable}
這段經文大家非常熟識,以色列的第一位君王 - 掃羅,如何在神面前失落了今日許多教會,弟兄姊妹,外表上生命活潑,如同獻各樣的祭,行神所喜悅的事.但,這段經文裡的重要事實,恐怕會成為今日教會不蒙福,和生命成長的攔阻;就是某些小節上,未能完全順服神這事明顯是極其嚴重的,若想知道這事的歷史背景,(可以翻閱撒上十章)神從眾百姓,以及一小支派中揀選掃羅作王;先喜悅掃羅後轉為憎惡,並對立他為王后悔;最後將他權柄收回掃羅.被揀選,膏立,而至被丟棄﹀失信,究竟為了什麼請勿忘記裡面一句說話:?掃羅於被神厭棄的境地,因為沒有完全遵守神的話部份順服神,是非常可怕的; 「它」欺騙過千萬基督徒,使他們失去屬靈能力!陷入嚴重罪行,是信徒靈命致傷.容我發問,「掃羅有沒有順服?」誰敢說掃羅?有順服事實上他作了大部份神所吩咐的從第20節:「我實在聽從了耶和華的命令,行了耶和華所差遣我行的路,擒了亞瑪力王亞甲來,滅盡了亞瑪力人」他在大部份的事上順服神,但神說:.「你沒有順服我」耶和華叫他殺盡所有的人,他殺了百份之九九‧九,祗留下一人神要他滅盡所有牲畜,但他留了肥壯的據說:.作獻祭之用他確曾順服,甚至比你們的順服更高我每逢讀到這經文時,內心大有感受;神的要求如此之高,我滿足神的光景又如何難怪許多神所應許的福份,我不能得到按人看來,掃羅已順服神;?但神看為不順服,因為神所要求的,是百份之百的順服.願神憐憫我們,叫我們看到自己的光景,又看到神的心意如何.叫我們知道部份順服的可怕,就使我們朝著神要求我們完全的方向走,神是會悅納我們有這心意的.
\newline
\newline
為何部份順服神,後果如此嚴重可怕歸納為下列數點理由?
\newline
\newline
1.因部份順服違背了神完美旨意,神叫掃羅攻打亞瑪力人,吩咐他將所有男女老幼,牛羊牲畜,盡行殺滅.掃羅殺了所有的人,祇留下他們的王亞甲.消滅了大部份牲畜,祇留下少數肥壯的.讓我們不要忘記,他誠然作了許多,但他在神要完全消滅這命令上,是違背了.(雅二10):「因為凡遵守全律法的,只在一條上跌倒,他就是犯了眾條.」若再舉例說明,在神的十誡,來代表祂完全的旨意；若我犯了其中的一條──偷盜,就等於違背神的全律法.必須十條都守住,才算是遵守神的全律法.所以上面說過,部份順服是嚴重的錯誤.因為它常欺騙信徒,以為神知道我的軟弱,容許我馬虎點過去.要知我們在一件事上不順服神,就是過一個違背神的生活.
\newline
\newline
有些兄姊作基督徒許久,各方面也見到他循規蹈距.靈修,崇拜,事奉,奉獻,主日學等都有他份；但總覺得他生命裡缺少甚麼?沒有活力,喜樂,生命裡缺少興奮.主要原因,我想,他必定是在某事上違背神.那麼他就不能過一個與神和諧的生活了!無論那攔阻是大事小事,容我分享一點個人見證:當我在某時期,追求靈命長進,一直過程順利.但到一個地步,突然停頓下來,不曉得甚麼緣故.雖有讀經,但不興奮；雖祈禱,不過虛掩其事,缺乏甘甜；雖有事奉,也不能滿足.我一直對神尋找原因,神就給我看到我生命中的不順服,在神完美的旨意中,我作了悖逆的事.原來在教會中,有弟兄得罪我,我心中仍舊懷恨他；神在我心中工作,要把仇恨除去.可惜我總不能用積極的愛去愛他,和他恢復交通.神在這事上嚴厲對付我,使我在與祂交通上,失去甘甜活潑；直至我起來重建我與弟兄間應有的正常關係.這在我是非常困難,因為提起他,就再次觸發過往的舊痕,但神終於勝過我.一天,我靠著祂的恩典,主動去按他的門鈴,向他道歉.用一個敞開的心,用積極的愛,來與他重新建立完美的關係.藉著今日的經文,使我看到一件事是肯定的；若你生命中有一件錯事,有一件喜愛的,不願放下；就足以影響神與你的關係,如同神對掃羅的關係；至終,在神的恩典中失落了.
\newline
\newline
2.部份順服毀壞(不完成)神的工作,神要掃羅把亞瑪力人口,牲畜盡行殺絕.到天上時,就必知道神何以定要作這工.但神的工作若半途而廢,是因掃羅不澈底地遵行神的旨意.弟兄姊妹曾否想過一件美麗的事?我每逢念及,就感謝讚美神.因神在每個弟兄姊妹身上,有祂完美的計劃,要叫我們完美像耶穌.另一方面,我們的前途,神預備了非常美好的道路,使我們完成祂所託付之工.但幾時,我們若在某事上不順服神,在某一項決定上違背了神的旨意；就足以影響神在我們身上所定的計劃.
\newline
\newline
有一位姊妹,未信前,是一位令人退避三舍之人,貌醜,性惡,好罵人.總之,不堪承教.信了主後,起初還沒有多大改變.後來跟一位姊妹交通,這位姊妹告訴她,神可以藉她彰顯祂的榮美.姊妹聽後受感動,澈底地對付自己,直至到把原來個性,完全改換過來.不但性情,連樣貌也變為可人；散發信徒的馨香,掛著可親的微笑.神實在能在人的個性或其他方面工作達到一地步,成就祂極美麗的工作.但如果你在興趣上,個性上,有些神要對付的地方,你不願放下,那就攔阻神完美的工程.我們何等盼望,奇妙的工作成就在屬祂的人身上,反照祂極大的榮光!
\newline
\newline
從另一角度來看,神在基督徒身上有工作,託付和奇妙安排.若作下一個錯誤決定,違背神的旨意,可能使神原來在你身上的計劃受到攔阻.我曾經有過這樣經驗,清楚了神的呼召,但看看現實環境；傳道人生活的種種艱難,父母的影響,親戚的勸阻.......我就沒有入神學院,反而進了某大學讀書.那一年生活,總而言之,用「迷惘」二字可以形容了.直至我再次對神順服,回想神在我身上的恩典和工作；實在要感謝祂!沒有因這一次背叛,而全部摧毀.我要對青年們講出另一個人的見證.一位朋友在十八歲時接受了神的呼召,但一直背叛直到如今.據說四十年來生涯,不獨未曾體會到神在他身上的完美計劃；更沒有享受甘甜和諧的交通,中國人諺語說:「一失足成千古恨.」一次的背扳,足以把人帶進痛苦的深淵.
\newline
\newline
3.部份順服,是以人意代替神的榮耀.掃羅對撒母耳說:「我有罪了,我因懼怕百姓,聽從他們的話,就違背了耶和華的命令.」(24節)每次不順服神,都會告訴自己,有一個意思比神的意思更好.這是很自然的表現,比如我不順服神的時候,理由是,我還有幾個兄弟需要照顧,怎能去做傳道?我應當跟父親去作生意,將來賺錢的話,可以支持別的傳道人來補回對神的損失.你想我的想法何等幼稚可笑.每次不聽從神的命令,就是人的意思代替神.神不喜悅這種態度；神會聽你的解釋,因你的解釋必定達不到神的旨意.
\newline
\newline
4.部份順服,帶來虛偽的面孔.徒有美麗的外表,而毫無屬靈的實際,成為不折不扣的假冒為善的人.13節:「撒母耳到了掃羅那裡,掃羅對他說:願耶和華賜福與你,耶和華的命令我已經遵守了.」類似惡人先告狀,若從心理立場看這幅圖畫,何等滑稽.好像我的小兒子,當他偷糖吃被我發覺時,就主動地對我表示親熱.掃羅心知自己有錯,也知道自己沒有澈底順服神.於是擺出一付虛假的偽裝,來遮蓋他內心的悖逆.不久前,我教會來了一位姊妹.在交通之間,她常透露她用多少時間讀經祈禱,她如何愛主,然而我心中暗忖,這恐怕將是教會中之問題人物.果然不出所料.這都是屬靈外表來遮蓋本身的虧欠,和靈裡的空虛.部份順服帶來的可怕結果.
\newline
\newline
5.部份順服,可能帶來說謊的機會.掃羅說:留下上好的牛羊,用以奉獻給神.我覺他可能是個藉口.一個違背神的人,往往能用屬靈的話作為藉口；而又往往講得美麗和動聽.有位朋友生意蒙神賜福.六年前因教會有需要,而向他提及奉獻的事.當時他說:目前他不能奉獻一文,因他有大計劃可賺大錢,並預備兩年後可將一百萬獻給神.不必說都屬子虛烏有.由於心的不順服,故編出許多美麗的謊言.
\newline
\newline
6.部份順服,一日自動爆發,會影響對神的見證.掃羅講完:耶和華的命令我已遵守了.但卻聽到羊叫牛鳴,牛羊不同他合作,他悖逆的罪行,便爆發出來.隱藏的罪,儘管我們竭力遮掩,到一日,必暴露出來,就影響神在我們身上的見證.經上說:掩蓋的事,沒有不露出來的；隱藏的事,沒有不被人知道的.唯有愚蠢的基督徒,才去作隱藏罪惡的事.
\newline
\newline
7.部份順服可挽救否?答案:可!雖然神無恢復掃羅的王權.有種情形很可怕,背叛神的人偷生於福樂中之損失,卻失卻了神的赦免.舉例說明,有一塊木髹得很漂亮,一天你在上面釘一口釘；後來覺得不滿意,從新把釘拔下來.釘雖拔下,但木上的釘痕不能消滅的!弟兄姊妹,生活應極小心.我們經過反叛神之後,雖然神仍赦免我們,但存留下不能磨滅的痕跡,和不能追回的損失.
\newline
\newline
神雖沒有永保掃羅的王權,但掃羅認罪,我相信他必蒙神赦免.但跟著要實行一件事,殺掉亞甲王和所有上好牛羊.弟兄姊妹!在你生命中,甚麼是你的亞甲王和上好牛羊呢?神要你除盡,而不要部份的順服.
\newline
\newline


\newpage

\section{第49屆~港九培靈會~張子華~第五講}
\label{sec:781}
\textbf{與人和諧關係－ 如何解憎恨,不和}
\newline
\newline
連結: \href{https://www.hkbibleconference.org/session-message/view/781}{\texttt{https://www.hkbibleconference.org/session-message/view/781}}
\newline
\newline
\hyperref[sec:780]{< < < PREV SERMON < < <}
~
\hyperref[sec:toc]{[返目錄]}
~
\hyperref[sec:782]{> > > NEXT SERMON > > >}
\newline
\newline
馬太福音 5:21-5:25
\newline
\begin{longtable}{cl}
\hline
\hline
章節 & 經文 (和合本修訂版)\\
\hline
5:21 & \begin{tabularx}{0.7\textwidth}{X} 「你們聽過有對古人說：『不可殺人』；『凡殺人的，必須受審判。』 \end{tabularx} \\ \\ \relax
5:22 & \begin{tabularx}{0.7\textwidth}{X} 但是我告訴你們：凡向弟兄動怒的，必須受審判；凡罵弟兄是廢物的，必須受議會的審判；凡罵弟兄是白痴的，必須遭受地獄的火。 \end{tabularx} \\ \\ \relax
5:23 & \begin{tabularx}{0.7\textwidth}{X} 所以，你在祭壇上獻祭物的時候，若想起有弟兄對你懷恨， \end{tabularx} \\ \\ \relax
5:24 & \begin{tabularx}{0.7\textwidth}{X} 就要把祭物留在壇前，先去跟弟兄和好，然後來獻祭物。 \end{tabularx} \\ \\ \relax
5:25 & \begin{tabularx}{0.7\textwidth}{X} 你同告你的冤家還在路上，就要趕快與他講和，免得他把你送交給法官，法官交給警衛，你就下在監裡了。 \end{tabularx} \\ \\
[1ex]
\hline
\hline
\end{longtable}
今日開始時,神將重要的問題問我們:????!你有沒有憎恨你的弟兄?有沒有與人不和睦?有沒有對人存成見?在肢體生活上有沒有攔阻?若有,求聖靈能力工作,使我們作得勝的基督徒弟兄姊妹,今日世人用甚麼眼光看教會,教會應怎樣向他們發出呼召:???你們要看真誠的愛嗎?要看真正的和平嗎?要看人與人之間的和諧相處嗎?可以到教會來看,任何作神子民的,應該能夠發出這樣邀請.但在這方面,我們卻對神多有虧欠.與人和好的關係,是由於我們既與神和好;這和好就進入到每個人中間,這人與人之間的關係,乃人神關係的反照(約壹四20):「人若說我愛神,卻恨他的弟兄,就是說謊話的.不愛他所看見的弟兄,就不能愛沒有看見的神.」保守兄弟姊妹之關係,其中一項重要的功課,要學習去原諒得罪我的人.今日要將這功課,分三方面思想:
\newline
\newline
(一)從聖經的教訓,找出原諒的真義
\newline
\newline
從聖經教訓,和神饒恕人的經驗,看對此名詞的解釋:
\newline
\newline
1.不再記念別人所作的過錯.經上一再說過,神若記念我們的過犯,誰能站立得住?神的饒恕,就像把它塗抹,不再能辨認.許多時我們說饒恕人,但見面時,舊事又再呈現眼前.甚至幾十年長,到死不能忘記.真的原諒是不再記念.
\newline
\newline
2.靠著主的恩典,完全接納這人進入你生命之中.
\newline
\newline
3.不單不記念人的惡,又接納他進入你生命之中.要用神的愛來積極地愛他.有時別人得罪我,我不對他懷恨報復便夠了.其實還不夠,更須用積極的愛,去愛那得罪我的人,才是聖經上「原諒」的真義.
\newline
\newline
(二)為何要原諒人,不原諒有何害處?
\newline
\newline
可分四方面來看:
\newline
\newline
1.饒恕是神吩咐的命令.主耶穌說:「你們饒恕人的過犯,你們的天父也必饒恕你們的過犯；你們不饒恕人的過犯,你的天父,也必不饒恕你們的過犯.」愛神的人,就必遵守神的命令.但有個情形很奇怪,人對誡命看得很嚴重.例如說:不可姦淫,不可偷盜.人若觸犯,就會被人不齒.對於不履行十一奉獻,人會敏感地覺察到.但對背後批評人,卻熟視無睹.我們要看重饒恕人的命令,正如看重不犯姦淫,不偷盜一般,這是基督徒應有的態度.
\newline
\newline
2.不原諒人,會直接影響個人的健康.心靈苦悶,長久不息地對人懷恨.據醫生的警告,許多疾病由此而生.它會引導精神,心臟,血壓等,產生惡果.所以恨人是直接地毀壞神的殿.
\newline
\newline
3.不原諒人,會影響靈命.傷害人與神之間關係,至少三方面:
\newline
\newline
a.不原諒人的就不能愛神.(約壹四20-21)「人若說我愛神,卻恨他們的弟兄,就是說謊話的.不愛他所看見的弟兄,就不能愛沒有看見的神.愛神的也當愛弟兄,這是我們從神所受的命令.」特別注意第20節的「用」字:不愛弟兄,就不能愛神.若對弟兄懷怨,不用積極的愛來愛弟兄,對神的愛終受影響,成為靈命的致命傷.
\newline
\newline
b.不饒恕人,就是拒絕神的饒恕.(參太六14-15,12)倘若我不願意僥恕人,就是我不需神饒恕我.不知是否你到一個程度,不需要神的饒恕；我卻每日需要神的憐憫,恩典,饒恕,否則我不能生活.耶穌教門徒祈禱:「免我們的債,如同我們免了人的債.」是非常重要的屬靈原則.
\newline
\newline
c.影響祈禱生活,與神的交往.參(太五21-26)用今天的文字形容:當你跪在神面前祈禱時,若想起與某人有問題,就要先與人和好.否則,神不會聽你祈禱,你與神交往就受攔阻.以上的問題是否嚴重,一個愛神的人,只能抉擇取捨.
\newline
\newline
4.不原諒人,會影響人的情緒.今日社會流行兩種精神病:
\newline
\newline
一,是情緒低沉,
\newline
\newline
二,是焦慮.
\newline
\newline
這都與情緒有關.若不注意這方面的健康,也會陷入同樣困境.上述兩種病癥的有無,與我們心境是否舒暢,與人關係是否和諧；有無對人懷恨,與人交惡.都會影響個人情緒生活,構成致命傷.以上四方面,但願我們考慮,是否願意過原諒人的生活.
\newline
\newline
(三)如何處理對我有過犯的人:
\newline
\newline
1.要開始學習思想,神用何等大愛對待我們.其實不原諒人的試探是最大的.做一個教會的牧師,他會受無理的攻擊,批評,背後講閒話的機會最多.傳道人往往成為議論的中心,他的生活,行為,講道,態度......都可成為人的話柄.當我遭受他人的毀謗時,就跪在神面前,思想神的愛.我原本如何不配,祂卻饒恕我；神的愛成為一股力量,使我有寬大的心；使我進入神的愛中,被祂溶化.
\newline
\newline
2.當人得罪我時,應有正確的觀念.剛才所讀(撒下十六11),大衛受便雅憫人示每辱罵時說:「由他咒罵罷,因為這是耶和華吩咐他的.」他把得罪他的人,看為神使用的工具.用來教導他,對付他,以達到神的心意.大衛的一生,對神有湛深的經歷.如今他在犯罪失敗當中,家庭巨變,兒子押沙龍正尋索他的性命；又被此莽漢咒罵.本來他可略一吩咐,就可把他解決了.但他說:不可,這是神對付我的工具,使我學成未學之功課.受到人的攻擊,可能是神所使用的方法.為要磨煉我們有更大的愛心和信心,更澈底的忍耐.從約瑟的生平看,當與諸兄相見時,兄長見面前的大臣,就是過去他們害過的約瑟,他們驚震不已!但約瑟說:你們雖然害我,但是出於神,因要差我在你們之先來到這裡,要保存許多人的性命.我最近也在這方面,學到些難學的功課,也感謝神帥領我得勝!我在神學院畢業後,一直在華盛頓一間教會作牧師.數年來教會非常蒙恩,發展很快；各方面都好,一帆風順.家庭,子女,工作,各方面都非常完滿.心中不禁泛起終老是鄉之感.但另一方面,神又似有引導,在東南亞為祂作點工作.於是斷續地收到幾處神學院的邀請,初不介意,後來接二連三地接到數地神學院的來函,才內心有點被打動.最後接到本港中國神學研究院的邀請,還在猶疑不決.某日,教會中一位很屬肉體的人,到來問我.他說:像他這樣的人,為何我不給他作教會領袖.後來這人糾合一些向來疏遠教會的人,到來抓起風波,惹事生非.當時我非常氣忿,但當我跪在神面前禱告,就知道原來這是神手中所使用的工具.當我遲延不肯照神旨意回東南亞工作時,神有祂的辦法.興起這人來,作祂管教的杖和竿.當我明白過來,就開始看這與我為敵的人,看是神手中的工具,為要使我這悖逆神旨意的僕人,肯改變來遵行祂的道路.因此對於得罪傷害我們的人,應有正確的態度,看出可能是神管教我的工具.
\newline
\newline
3.須採取積極的態度,到對方面前,與他有敞開的交通.這誠然不易學到的功課,每個弟兄姊妹,當受到人毀謗中傷時.錯在對方,必須他首先來向我認罪,我才能不予計較.鮮有能主動地去與對方和好.如果你愛主,正確的態度去對待那些惡意中傷你的人.求神賜給力量,主動地到他面前,敞開心門,問他究竟我作了甚麼錯事,得罪你?這樣敞開自己與對方交通,往往得到想不到的果效.神所用的有效武器,使我們能勝過敵對我們的人.
\newline
\newline
4.必須相信主曾說:伸冤在我,我必報應.基督徒在任何情形之下,不能自行報復.照主的話:寧可讓步.照你我個性,在這等教訓面前,該受神何等磨煉.無論如何,不能以牙還牙,以眼還眼.神是公義的,當信祂的話:伸冤在我,我必報應.不必自行解決,否則越弄越壞.舉例說明:教會有 AB 二人,二人因故失和.立刻可以產生兩個反應:一是立刻予以還擊,半斤八襾,旗鼓相當.於是雙方的挈友定必連成陣線,彼此敵對.教會就產生結黨紛爭,終而導致分裂.如果另一個反應是,弟兄得罪我,我還報之以愛心和饒恕,就不會產生對立.因為我相信主的話:伸冤在我,我必報應.
\newline
\newline
最後,或者神使用你去幫助那敵對者,要在主裡面得著長成.意思是,某人中傷你,一路造謠,肆行毀謗.但你一味以善意回報,好像毫不介懷.就會令他生下思想,何以令你如此?你好像有某些東西,是他所缺少的,會興起一個尋求的心,去追求你所具有的謙卑,溫柔,忍耐,恩慈.其次,對外來的逼迫,祇加強你為那人禱告的心志.可能藉著你的代禱,改變了這弟兄,教會中處理人與人之間的問題,用的方法若對的話；可避免產生爭論,更可成為那些攻擊者──軟弱弟兄,有成長的機會.
\newline
\newline
你要用積極的愛,去對待恨你的人.如你有任何一個兄弟姊妹,是你眼中釘的話,這問題你要立刻解決清楚.
\newline
\newline


\newpage

\section{第49屆~港九培靈會~張子華~第六講}
\label{sec:782}
\textbf{接受你自己}
\newline
\newline
連結: \href{https://www.hkbibleconference.org/session-message/view/782}{\texttt{https://www.hkbibleconference.org/session-message/view/782}}
\newline
\newline
\hyperref[sec:781]{< < < PREV SERMON < < <}
~
\hyperref[sec:toc]{[返目錄]}
~
\hyperref[sec:783]{> > > NEXT SERMON > > >}
\newline
\newline
以弗所書 2:10-2:10
\newline
\begin{longtable}{cl}
\hline
\hline
章節 & 經文 (和合本修訂版)\\
\hline
2:10 & \begin{tabularx}{0.7\textwidth}{X} 我們是他所造之物，在基督耶穌裡創造的，為要使我們行善，就是神早已預備好要我們做的。 \end{tabularx} \\ \\
[1ex]
\hline
\hline
\end{longtable}
今日所講的是陌生的題目,可能對一部份人來說,對內容會暫時不明白.但聖靈會繼續祂感動啟示的工作,對各人靈命有益.許久之前有一姊妹,致電給我夫婦;透露準備用她自己的方法,來了結殘生當我們告訴她未採取行動前,可否來我們處,彼此交通;她勉強答應下來,我們就彼此會面.姊妹看來十分憂鬱,既自卑又自憐,對人生充滿失望;.過往的完全失敗,未來又毫無把握,覺得活下去毫無意義,後來,我們請她把心裡所藏的盡情吐露,她終於詳細地述說她的心事(在此我要請作父母的特別留意這位姊妹的話,知道怎樣對待你的兒女,也請教會的弟兄姊妹,切勿隨便對別人妄下評價)以下是​​她的自白:.我兩歲時,母親常說我生得最醜,最初我不太明白,後來我年齡漸長,較為懂事時;到有一日,著實解答我心中疑團,我對著鏡子仔細檢視自己,我果然覺得自己很難看,眼耳口鼻都不相稱,越看越醜,以後逐漸形成一個觀念.我是一個不幸的,醜陋的人.及至後來與神有了關係,還是時常對神提出疑問;為何祂造別人美麗,而對我卻如此不公平,我不能滿意神的安排,一直內心裡對神反抗,由於堅持著這種態度,使我心境一直趨於下沉,對生命一無留戀.所以,今日我要和各位討論一個問題:對自己有健全的認識,接受自己,是和諧生活的重要步驟和秘訣.
\newline
\newline
(一)正確的認識自己
\newline
\newline
最佳的定義,是明白而且接受,神在我身上所創造的;包括容貌,性情,才幹......和一生的遭遇「我們原來是祂的工作,在基督耶穌裡造成的」(弗二10)「我未成形的體質,你的眼早己看見了.......」(詩一三九16)倘若我們只在鏡子裡,看到容貌的種種缺點,而看不見神在你身上創造的計劃和大愛,則你是一個有問題的人.當接受和明白神在我身上的作為.就是承認一項事實,要在神所賦與,安排給我的一切條件裡,存感謝的心;盡最大努力,靠著神的恩典,過神所使用我之生活.
\newline
\newline
(二)假如不接受這對自己正確的認識
\newline
\newline
你的人生會帶來空虛,惆悵,憂慮,痛苦.從下面幾方面,表現出這等人所隱藏的難處:
\newline
\newline
1.成為一個自我批評者.你不能開朗地與人相交,趨向孤單獨處,情緒低沉,不願完全信賴神,自憐自怨,進而成為一埋怨神的人,我們要自問有無這等現象；自憐自卑,要求苛刻,離群獨居,不敢信賴神.另一個表現是；注重外表,引人注意.這都是不接受神所創造的事實.
\newline
\newline
2.情緒不穩引致埋怨神.一個對自己諸多不滿的人,很自然地發展成為對神不滿.這樣的人,必然內心充滿了問號?對神總是:「為何?......」為甚麼我容貌比不上人?為甚麼別人富裕我貧窮?為甚麼我不像別人那麼有才幹?......歸根結蒂,就是不接受神的創造和安排.
\newline
\newline
3.反抗權柄.接著上述的充滿疑問之後,進一步為反抗權柄.對父母,教會,帶領者,政府,都表現不合作,對之反抗.因為覺得神對待他不公平,內心的怨懟形成對權柄反抗.作父母的應該明白此點,好好的幫助兒女.教會青年如出現反抗行為,領袖也應知其成因.
\newline
\newline
4.成為過度注重物質者.舉例來說:人如果不接受神的創造,對自己不作正確的認識；就會作成錯誤的認識,發生於別人所給與之評價.受他人意見所影響,因為自覺才幹,樣貌,......種種不如人,又怕別人看為不中用的.於是就用外表事物來遮蓋,用金錢,物質來轉移.用各樣享受來彌補他人的評價.弟兄姊妹!若你的心不在神裡面得享受,卻一生追求物質享受,表明你對神不滿.
\newline
\newline
5.不滿現實,對任何事物都不滿意.
\newline
\newline
(三)應如何接受自己　靠著神恩典,產生健康的自我認識:
\newline
\newline
1.靠著神的恩典,每個人都當讚美神.因每個人都為神所貴重的.無論美與不美,富與貧,有學問與否,甚至瞎眼瘸腿；記得我們都被耶穌看為寶貴,都當讚美祂.因為至少有三件事是每個基督徒所共有的:
\newline
\newline
a.是按神完美計劃所造的(創二7)
\newline
\newline
b.是用主耶穌基督的寶血所買贖的(羅五8)
\newline
\newline
c.是聖靈所居住的(林前六9).建立健康的自我認識,首先是明白神如此寶貴我們.每當我們失敗跌倒時,每想到神的愛和創造,耶穌基督的拯救,聖靈的提醒；就不禁用詩人的話說:「人算甚麼?你竟顧念他.」千萬不要忘記寶貴的事實.神可能造你不太美麗,無多大才幹......但經由神所創造的,祂的智慧,祂的工作,豈是我等有限的人,配作論斷.
\newline
\newline
2.不能讓人的價值觀,改變對自己的認識.記得使徒保羅的名言:「連我自己也不論斷自己.」別人說他氣貌不揚,他滿不在乎說:「連我也不論斷自己.」唯有主能判斷,不受人對他的評價,來改變神創造的事實,所產生對自己的認定.弟兄姊妹!如今我要提出一個很重要的真理:基督徒的團契生活,彼此相愛,要在主裡面,無條件不帶任何評價地,彼此接納.若你按喜好,有條件地來與弟兄姊妹相交,則你會造成弟兄姊妹不正確地去認識自己.若我們因人的容貌奇怪,或缺乏某些才能,帶有某些缺憾,而加以岐視的話,就是得罪神.因為他也是照神形像造的,是神所愛的,我們哪有資格評判他.但記得,你不要讓人對你的評價影響你對自己的認識.
\newline
\newline
3.相信神所創造,都是好的.讀創一,二章,神看一切所創造的,都是好的,不知你作何感想.我在哥倫比亞聖經學院讀書時,有一位姊妹生來殘廢,後來更不幸患病使她雙腳也被鋸掉.但我每想念到這姊妹時,心中就充滿了感謝神.因為這一切事實,並不影響她發出基督的香氣.她是同學中最快樂的一位姊妹.當與她交通時提及到她的身體,會否影響她的事奉道路.她說:我不知道這會有甚麼壞影響,但我知既出於神所創造和安排,就必是最好的.弟兄姊妹,是否你稍不如意就埋怨神,例如牙痛或其他小毛病.讓我們讀一遍(賽四十五9-10).
\newline
\newline
4.神在我們身上,改造的計劃還未完成.「我們原是神的工作,在耶穌基督裡造成的,為要叫我們行善,就是神所預備叫我們行的.」「我深信那在你們心裡動了善工的,必成全這工,直到耶穌基督的日子.」(弗二10,腓一6)舉例來說:你知否有些外表並不美麗的人,卻可以叫人察覺到她具有另一方面更美的地方.我在教會事奉時,有位姊妹來與我交通,她說:「牧師,你有沒有發覺,教會中許多夫婦,丈夫很漂亮,太太卻不大相稱.」意思是:我生得比她們好看,為何男士們都找錯對像?當時我祇能回答她:「女性吸引人的地方,有外在美和內心美兩方面,而內在美更為重要.基督要在人的個性上薰陶,磨煉.鑄造出基督徒獨有的溫柔,忍耐,愛神愛人等特性.」我一路思想教會中的年輕太太,果如她所講,都不大美麗.但大都具有神所薰陶,塑造,那種不同凡俗的形質.所以我們要建立正確的認識,神在你身上改造的工作,還在繼續中.祂要在你身上造成,帶有祂的榮美,智慧,力量.
\newline
\newline
5.要相信神會使用不幸遭遇,來達成祂完美計劃.有一次讀一本書,使我深受感動；在澳洲有一位善寫詩的女子,年幼時患上風濕症,以致手足殘廢不能運用,終生要靠輪椅幫助.究竟這樣的人,在神計劃中有何用處?如果不用正確的觀念,來面對這事實人,就祇會等死.但她說:神叫我手足不能自由行動,她必定有一樣工作給我做,是我如果健康正常時永遠不會作的.於是用通訊的辦法,幫助了成千上萬心靈痛苦憂傷的人,遠過於任何教會牧師所作的.每次我去見蔡蘇娟姊妹時,都得到幫助.覺得她具有最健全的自我認識.一個本來殘廢,不能再見日光的人,她所作的比我們更多更偉大,我們豈不應認定,縱使不幸臨到我身,但這遭遇是達到神完美計劃所必需的.沒有它,就作不成.
\newline
\newline
6.必須將自己看得合乎中道.這是在教會事奉的重要原則和正確方法.有兩種人是有問題的:一種是把自己看得太高.過於高抬自己,好出風頭人物；否則又看成懷才不遇,滿腹牢騷.另一種,卻把自己看得太低.樣樣都以為作不來,祇有成為別人的負累.這兩種態度都不可有.應該把自己看得合乎中道,在自己的崗位上,盡其所有來事奉神.
\newline
\newline
最後,用四個步驟來作為結束:
\newline
\newline
1.隨時注意在你生命之中,有無不能感謝神的地方?如果有的話,可能是一個警號.
\newline
\newline
2.要在禱告中,或對著鏡子的時候,要從心裡發出讚美的話:神阿!你的創造實在奇妙!
\newline
\newline
3.求神繼續改造你,對付你的軟弱,所有不能夠祝福人的地方.使你能更大,澈底地,發出基督的香氣,展露你內在的美麗.
\newline
\newline
4.奉獻你自己給神,求神用你,使別人因你產生神工作的果實.
\newline
\newline
今日我選講這題目,因為教會和家庭的難處,多因此有虧欠.願聖靈幫助我們,使多思想自己,多看神.建立起對自己正確的認識.
\newline
\newline


\newpage

\section{第49屆~港九培靈會~張子華~第七講}
\label{sec:783}
\textbf{家庭的和諧－ 向父母服權}
\newline
\newline
連結: \href{https://www.hkbibleconference.org/session-message/view/783}{\texttt{https://www.hkbibleconference.org/session-message/view/783}}
\newline
\newline
\hyperref[sec:782]{< < < PREV SERMON < < <}
~
\hyperref[sec:toc]{[返目錄]}
~
\hyperref[sec:784]{> > > NEXT SERMON > > >}
\newline
\newline
以弗所書 6:1-6:3
\newline
\begin{longtable}{cl}
\hline
\hline
章節 & 經文 (和合本修訂版)\\
\hline
6:1 & \begin{tabularx}{0.7\textwidth}{X} 作兒女的，你們要在主裡聽從父母，這是理所當然的。 \end{tabularx} \\ \\ \relax
6:2 & \begin{tabularx}{0.7\textwidth}{X} 當孝敬父母，使你得福，在世長壽。這是第一條帶應許的誡命。 \end{tabularx} \\ \\ \relax
6:3 & \begin{tabularx}{0.7\textwidth}{X} 見上節 \end{tabularx} \\ \\
[1ex]
\hline
\hline
\end{longtable}
由今日開始,一連三天講家庭的和諧.這是人際關係主要的一環.我深心地體會,家庭若不和諧,教會就受虧損.盼望每一位牧師同工,需要注意對家庭的輔導.各個家庭所潛伏之問題,若不得到解決,則不獨影響各人的靈性,也會將問題帶進教會.這三天所傳講的訊息,我自覺很困難.今天對青年講順服權柄,明天講為父之道,後天講關乎兩方面的訊息.記得十多年前我曾作導師,一位十四歲青年人問我:父親不准我參加夏令會,我應怎辦?當時我對他說,要背起你的十字架來跟從主.回想當年,這是何等不合經訓的輔導.教會中也可常聽到如此的輔導工作.應用背十字架跟隨主,就破壞兒子與父母的關係,這是非常嚴重的錯誤.今天講兒女對父母權柄應有的態度.觀念若正確,就可減少許多無謂的磨擦.
\newline
\newline
上面幾處經文,講到兒女對父母,態度上的幾個原則.然後我會引用幾個實例:例如父母反對你的婚姻,反對你全時間事奉的決定,應該採怎樣的態度.或者在座聽眾們,你心中有特別的難題,在今日講道中未獲得解決；你應該請教教會的長者,替你作輔導的工作.
\newline
\newline
原則一,神藉著父母的權柄和管教,成就祂在兒女身上的計劃.無論他們是否基督徒,他們代表神在家中,管教兒女.我們必需要有這觀念,而且思想正確；因此對兒女有個要──服從權柄.在神的國度,每個組織之中,要有權柄的次序.任何反抗權柄的行為,是神所不喜悅的.國家之掌權者,是總統和執政者.為著和平與秩序,人民必須遵守國家之法律.如開車超過速度就需罰款,破壞法紀就須坐牢.聖經說:反對在上掌權的,就是反抗上帝.(參羅十3)沒有權柄不是出於神.神在教會中,也有祂權柄的次序,神設立牧者帶領,因為他們是施教導之責,然後長者,執事,協助共同擔負教會聖工.香港教會應該對此有認識,神的僕人應在他地位上,得到他當有的權柄.會眾應該要順服權柄!倘若你今天聽道心中作難的話,你就有反抗權柄的傾向.家庭方面,假定兒女是信徒,而父母縱使是非基督徒的話；父母也代表神的權柄來治理這家,反抗父母就是反抗神的權柄.這是第一個非常重要的原則.
\newline
\newline
原則二,由第一原則,聽眾會產生下面疑問:若我父母無理,不信,不了解我,行為敗壞,那我是否仍要順服?如何順服?由此引進聖經予我們的第二項原則:我們順服,是因他們的地位,而非因他們的德性.我們應在這順服的要點上分別清楚,不可混淆.舉例來說:感謝神,我父親在晚年戒除吸煙的嗜好,但向來他有這習慣.記得十四,十五歲時,我和弟弟偷偷躲在廁所裡學吸煙；被爸爸發覺後,拖出打一大場.這真是一個很深刻的經驗.當被打完一頓之後,我心中很不平；想他也吸煙,有何理由責罰我.但我繼又思想,他在父親的地位上,有權禁止我吸煙.至於父親有無資格的問題,留待明天再講.但願青年人明白,父親的權柄,是由神而來的.順服是因他的地位,而非個性.經上說:在上有權柄的,人人當順服他；因為沒有權柄不是出於神的......因為他不是空空的佩劍,他是神的用人.(參羅十三1-4)
\newline
\newline
原則三,神在許多時,用看為不合理的事,去成就祂的計劃.祂的計劃就是造就你或你父母的屬靈方面.由於你的順服,甚至叫不信的父母歸信基督.試舉一例:教會有個十五,六歲青年來問導師:我的媽媽不許我參加夏令會,我應怎辦?或者嚴重點,父母不許參加主日學,祗准我間中參加崇拜,應怎樣應付?導師可能有兩個回答:一是教他不必聽父母的話,因為父母不認識神,不明白屬靈的事.於是青年回家,可能與父母的關係惡化,這是最常見的錯誤方法.試問你能夠為這個家庭作些甚麼,不能叫神的計劃在這個家庭實行.另一方法,用婉言勸導,也許父母不讓你參加夏令會是因為怕你吃不好,睡不好,不安全.可能父母的動機是出於愛護他的子女.應該在主裡順服父母,為他們祈禱.相信兒女願意如此聽從父母,就因這美好的見證建立良好的關係；神的計劃也可能實現在這家庭之中,這是反抗權柄者所永遠達不到的.反抗權柄是神所不喜悅的.基督徒要注意,不要錯誤地使用.天天背著十字架來跟從主,聽從神不聽從人是應當的；不要魯莽地用在對父母的關係上,以致家庭關係破碎,失去神所賜的和諧.
\newline
\newline
原則四,順從父母,由於年齡關係.至此,稍為提及父母方面,作為明日講道的預備.管教兒女分作三個時期:
\newline
\newline
1.完全控制──由出生至十四歲.
\newline
\newline
2.控制與自由各半──中學至大學階段.父母勿對子女要求過高,致弄出難題.
\newline
\newline
3.兒女對上面權柄之改易,改為對神第一,對父母第二.
\newline
\newline
因年齡及家庭背景,子女若仍在其父母的養育管教之下,則父母之權柄更具絕對性.繼續由半輔導,進至能自立時,明白首先順服神,然後父母居輔導地位.由年齡之漸長,而對權柄的要求有這三方面.在座青年,對此應知有所抉擇,在你求學的年代,在父母養育之下,應以父母為第一順服對象；在你自立之後,你要首先順服神,父母則站於輔導者的地位.
\newline
\newline
原則五,須學習,明白如何處理,上面權柄帶來之吩咐,明顯與聖經原則違背者.這原則有四點:
\newline
\newline
1.應在開始時有好的根基,在許多事情上,你讓在上有權者,知道你的順服.在你裡面,真正有順服的靈.
\newline
\newline
2.當命令來到時,雖然與聖經教訓衝突,但他的出發點乃由於愛.父母或反對你奉獻讀神學,或反對你跟某人結婚；雖不合理,但他的動機由於愛.心中有這順服的靈非常重要,是神的教會,人際關係的要點.
\newline
\newline
3.以完全順服的態度,將此不合神旨意的命令,向神祈求.
\newline
\newline
4.為此事恆切祈禱.既有順服態度的好根基,存心信賴父母的愛；然後對神有完全順服的態度,不求自己的益處.還要加上恆切的祈禱.你可相信,神的話在聖經中:「王的心在耶和華手中,好像隴溝的水,隨意流轉.」(箴廿一1)舉例說明這點:青年但以理,被擄到巴比倫王宮.他在許多事上順服王命,學巴比倫的學問,言語.因此建立美好的關係,予人有順服的印象.等到王頒下,要吃王所賜的飲食──祭偶像之物.他一面的求神,他知道作神子民是不應吃的.但他無反抗王的命令,用討價還價的方式,向管轄的人商量.然後設想王的用意要使我健康.他不懷疑神的權柄,裡面有順服的靈.然後神給他勇氣求王,給他祇食菜蔬,得王准許,他也果然健康更佳.這是聖經明顯的例證.
\newline
\newline
又再講另一例子,某青年要奉獻傳道,但父母反對.這人素常孝順,一向建立好的根基.今他清楚神的呼召,但父母不同意.他祇將心意清楚地解釋給父母知道,然後溫柔忍耐地過了一,二年.就見到父母從起初反對的態度,漸轉趨溫和；最後雖然他們仍不情願,但至終說:我兒,生命是你的,可以自作決定.我是對一些較長大的弟兄姊妹講,要用這原則,來應付那些看來不合理,攔阻神旨意的命令.希望你裡面滿有順服的靈,仰望神來改變這事實.
\newline
\newline
原則六,可能神用這權柄,攔阻你走上錯路.有個青年,大學畢業後,愛上一女子；父母反對他倆的婚事.反對的理由,絕無人的觀點成份,而是用理性,禱告,作出決定,當時我也曾勸這青年,應該考慮父母的意見,可能神藉他們攔阻他走錯路.二人終於結婚,短期內發現女方許多問題不能忍受；又短短二年間分開.然後他回來說:到此我才知道父母的權柄不是無理的,神曾藉著他攔阻我走錯路.神能用基督徒或非基督徒父母,作為祂的工作.
\newline
\newline
原則七,若我有順服的靈對父母,表明是一個神所喜悅的人；耶和華能成就我所不能成就的事.(箴十六7):「人所行的若蒙耶和華喜悅,耶和華也使他的仇敵與他和好.」今日之訊息似乎是一面倒,但明日會將它平衡；首先要兒女作得對,要聽從父母,作神所喜悅的人.然後神能作出我們看為難的事.
\newline
\newline
最後我將三樣顯著的難處,處理的辦法,給每位青年.
\newline
\newline
1.在父母養育之下,還未曾自立的.盼望在任何情形之下,縱使為屬靈的益處,但有當為你家人屬靈益處,讓「當信主耶穌,你和你一家都必得救.」的應許成就在你家中.
\newline
\newline
2.若父母反對你婚姻,應如何!
\newline
\newline
a.一定要延遲訂婚或結婚的時間.
\newline
\newline
b.重新在神面前禱告,求神給你確據,是否神用你父母提醒你.
\newline
\newline
c.平心靜氣地對父母提出對方的優點,即你為何選擇他的理由.
\newline
\newline
d.讓對方得與父母彼此談論認識.
\newline
\newline
e.虛心請教父母,對對方的意見和看法.
\newline
\newline
3.若父母反對你奉獻作傳道:
\newline
\newline
a.等候一段時間,在神面前祈禱等候.但非無限期延遲,祇讓事態不致惡化.
\newline
\newline
b.接受父母出於善良的動機.
\newline
\newline
c.先選擇入另一間大學,修讀有助於傳道的學識,以為他日讀神學先作準備.
\newline
\newline
d.用順服的靈,求神感動改變父母的心.
\newline
\newline
當你作到以上步驟之後,將會見到神在父母心中工作.若不然,則將順服之對象歸於神.每位青年當省察,是否裡面有順服的靈,你在家以何態度對父母.神要求你作順命的兒女.
\newline
\newline


\newpage

\section{第49屆~港九培靈會~張子華~第八講}
\label{sec:784}
\textbf{家庭的和諧－ 父母的責任}
\newline
\newline
連結: \href{https://www.hkbibleconference.org/session-message/view/784}{\texttt{https://www.hkbibleconference.org/session-message/view/784}}
\newline
\newline
\hyperref[sec:783]{< < < PREV SERMON < < <}
~
\hyperref[sec:toc]{[返目錄]}
~
\hyperref[sec:785]{> > > NEXT SERMON > > >}
\newline
\newline
以弗所書 6:4-6:4
\newline
\begin{longtable}{cl}
\hline
\hline
章節 & 經文 (和合本修訂版)\\
\hline
6:4 & \begin{tabularx}{0.7\textwidth}{X} 作父親的，你們不要激怒兒女，但要照著主的教導和勸戒養育他們。 \end{tabularx} \\ \\
[1ex]
\hline
\hline
\end{longtable}
今天的講道著重對父母講,首先要明白一項真理,神如何將權柄賦與.在任何屬靈的組織中,論到權柄的地方,都不是唯我獨尊,暴君式專權的.神賜下權柄,也帶來責任,這是我們首先要明白的.例如聖經說:「你們作妻子的,當順服自己的丈夫,如同順服主.因為丈夫是妻子的頭,如同基督是教會的頭,他又是教會全體的救主.教會怎樣順服基督,妻子也要怎樣凡事順服丈夫.」很顯然夫妻之間,由丈夫作頭.但作頭者,並不等於可以任意發號施令；而是領受當盡的責任.換言之,權柄所帶來之責任感,應遠超高高在上之意念.這也是教會負責階層所應具備之正確觀念.
\newline
\newline
其次講父母對兒女的期望,我國傳統觀念上,養育兒女,大多望將來可藉兒女享福.為人子女者,要記念父母,報答親恩,是理所當然的.但作父母的,若本此動機來養育兒女,則與聖經的原則相背.神給你做父母的責任,就當盡責把他們撫養成人；使他敬神,愛神,能夠自立,為父母的責任才算完成.必須帶有這動機來完成責任,不是為了報酬.但也盼望青年們注意,父母若用正確觀念對待,則兒女也應出乎自然地,盡各樣方法去報答父母養育之恩.如此兩方面平衡,自然減少灰心,失望,種種痛苦.今日訊息,特別注意到兩方面:
\newline
\newline
(一)不要惹兒女的氣
\newline
\newline
「惹兒女的氣」這話的意思,有許多不同的見解.我自己的領受:在家庭環境生活中,父母不作任何事,或疏忽其責任,使兒女身心導致不健全.惹兒女的氣,會產生後果,到兒女長大後,出現情緒,心靈不健全的病徵.按照心理學,輔導學之研究結果:許多情緒不穩定的病症,例如:情緒低沉,焦慮,恐懼及許多其他毛病,絕大部份與他的家庭背景有關.因此輔導者要穩定青年們的情緒,常見的輔導工作,祇是研究對方家庭背景；希望從其中可找到影響青年情緒不穩定之成因.惹兒女的氣,就是不明白兒女的需要,怎樣使他們滿足.如此長大後就會形成兒女身心的不健全.這樣的人,可能對己,對人都有害.情緒上略有三方面,為人父牧者需留意:
\newline
\newline
1.每個家庭,兒女要得到父母的愛才能滿足.缺乏父母的愛與關懷的,長大後會成為一個問題人物,對社會和教會都有害處.形成孤僻,不會認識神的愛,也不會愛人.請求每一位父親,母親注意:要在愛上,滿足你兒女的需要.我要提醒:許多基督徒家庭,雖然有愛,但卻是偏愛.你知道偏愛對兒女的害處嗎?因為它會引起惹兒女的氣.試看聖經中以掃,雅各的生平,若從心理學上看,是非常有趣的事.這家庭充滿著問題,撇開神預定揀選來說:父愛以掃,母愛雅各.偏愛使兒女的心靈情緒損傷,被溺愛者固然不健全,被忽略者更不健全.這是神所不喜悅的,因神絕不偏待人.容我問一問題:你的兒女是否感受到愛?愛,不是在乎有吃,有穿.愛是一件微妙抽象的事,很難解釋出來,但必能被感受出來.正如我知道基督愛我,因我感受得到；但要我具體描畫,則很難解釋.每一位父母,都應在愛情上使兒女得到滿足；如果偏愛,就是虧欠神.有個基督徒家庭,母親對待兩個女兒,就犯了這個錯誤；對一個溺愛,對一個憎嫌.我一直為她擔憂,後來兩位女兒都長大出嫁；都可見到不健全的性格,情緒,兩者都受到損害.所以偏愛的父母,就是惹兒女的氣.
\newline
\newline
2.兒女不單需要父母的愛,還需要知道他在家中受到重視.要明白自己的身份,要知道他是誰.他在父母眼中是重要的.因此,作父母的就應避免用世人的話咒罵兒女.例如:你是一個最無用的東西......；全世界最......是你,不知為甚麼有你這樣的兒子.如果你重複用這種說話來對待你的兒女,他們會想,我究竟是個怎麼樣的人?兒女需要知道自己的身份.若他在家庭中,得不著這方面滿足；長大後,他會用自己的方法,叫人知道他是誰.報紙上常見到在美國許多這樣莫明其妙的人.例如美國總統候選人華萊士,在馬利蘭州一家百貨公司附近；競選演講時,被一個人用鎗狙擊,以致半身殘廢.兇手被捉拿後,經心理學家調查他的身世和他的家庭背景.原來是出自一個破碎家庭,被父母輕看；他為要使人知道有他這個人存在,就不惜採用這愚蠢的方法,為要使無數的人注意他,於是另一個就被犧牲了!他對心理學家說:我對華萊士並非憎恨,與他素無仇怨,他不過是我揀選的對像而已.我祇要用這簡單方法──開一鎗；就使全世界的人知道我這人存在.又以前,在德薩斯州立大學,有一個人帶著十幾枝鎗,數萬發子彈,跑上鐘樓,向下面掃射,凡會動的都被殺害.他講出唯一的理由,就是他要全世界人都知道,有他這人存在.家中的兒女如會感覺失去身份,無重要性,是個多餘的,可有可無.在積極方面,他可能成為轟天動地的人；但消極方面,也會成為孤獨,消沉,無信心,在任何事上都感到自己無能的人.我更要提醒為父牧者,對子女不要要求過高.要明白子女才智,當他們在學業上,已竭盡所能時；就應該對他表示嘉獎,使他知自己在父母眼中被重視.在香港教育制度競爭劇烈之下,將會產生極多情緒不穩的青年.父母若處理不當,則影響更甚.考80分則要求85,得85分要求90.升中試,大學入學試,都給青年們過重的擔子.如此過重的壓力,會造成青年們精神的損害.對兒女才智無正確認識,憑一己之慾望；無饜的要求,會造成青年對他身份存在的懷疑.他對自己失望,懷疑,不感到所盡力爭取的,是一次成就.應當留意在這方面使他們滿足.應當問問自己是否在這方面有惹兒女的氣?
\newline
\newline
3.兒女需要安全感,父母要供給他們這方面的需要.當兒女幼少時,他飢餓,懼怕而啼哭；祇要抱在懷中,他就會感到安全,但當兒女年歲漸長,進入大學階段時,就更需要有安全感；使他們覺得家庭對他們是安全可以接納的.一個缺乏安全感的人,會帶來情緒上的不安,失去自信心,懼怕,猶疑不決......等心理缺憾.安全感之產生,與夫妻之感情關係甚大.一個破碎的家庭,離婚的夫婦,終日詬誶不寧的家庭；都不能給予子女有安全感.必須夫妻相愛,齊心努力,家庭和諧,才能產生安全感.我們應該問問自己是否在此惹兒女的氣?不要讓任何事物,環境使兒女在不正常的氣氛下長大,以致情緒,心靈受傷!
\newline
\newline
(二)在主裡教養兒童
\newline
\newline
父母對子女有屬靈責任.「教養孩童,使他走當行的道,就是到老他也不偏離.」(箴廿二6)這話照我自己的領受是──當兒女年少時父母要盡屬靈的責任來教養他；使他們信主,愛主,這樣就算到老,他也不會離開主道.按照教會禮儀來說:長老會與浸信會對兒童有不同作法.長老會對兒童施行灑水禮,浸信會行奉獻禮.我就長老會的立場來說:灑水禮並不代表給予救恩.當一對夫婦將他們的子女,抱來到神面前,為他行奉獻禮或灑水禮,都是現出一幅極美的圖畫；並不是在孩子身上,乃在那對夫婦身上.他們一定不會是對非基督徒,今天藉著奉獻或水禮,與神立約；應承在家中盡上屬靈的責任,來教養他們的子女,相信主必定成就祂自己那部份.神從來是不會失誤的,父母若將子女帶到神面前,將來他若不信,反叛神；我們應當自己判斷,神不會失信,可能由於我們失責而已.我很喜歡約翰衛斯理母親所講的話:神賜給我兒女,好像將我帶入宣道工場,我變成一個宣教士.面對著工場──家庭,神給我大責任,帶領兒女信主,愛主,在屬靈上從各方面,引導他們走主之道路,使了解神的計劃.你在屬靈的責任上有沒有盡上教養之責呢?讓我向你發問幾個問題:
\newline
\newline
1.有沒有神的祭壇?有沒有因教導兒女讀經祈禱,留給他們優良的印象；使他們知道讀經,祈禱是非常甘甜的福份.要知道家庭祭壇,是和諧福樂的基礎.
\newline
\newline
2.自己屬靈生活,有沒有為兒女留下好榜樣?父母應首先自己以身作則,才能用真理正道來匡正青年子弟.否則徒然憑身份地位,使子女服從,到底失去見證.
\newline
\newline
3.有沒有盡力帶領兒女對教會愛慕,及有正確的認識?在我的意念中,有一幅最美好的圖畫.就是主日早晨,一家人在吃完早餐之後；一齊去同一間禮拜堂參加崇拜,這是好得無比的.但可見的是,夫妻,兒女,分別在不同的教會聚會,這實在不是良好的現象.希望弟兄姊妹留意,改正這病態.家人一同禮拜的好處,父母有機會,幫助兒女明白牧師所傳訊息.使一家有共同的目標,可齊心努力.我初在美國華盛頓教會事奉,感覺聚會中多數是成年人,不見他們的兒女.原來他們都到用英語的禮拜堂聚會.我立即建議改為繙譯英語,把青年人爭取回來,使他們有機會全家一致地參與禮拜.
\newline
\newline
4.有沒有用各樣辦法,來影響你的兒女尋求神的旨意,走神的道路?作父母的何曾有坐下來和你的兒女有屬靈的交通,利用為父母的影響,輔導的能力,使他遵行神的旨意.耶穌曾說:舉目向田觀看,莊稼已經熟了,可以收割了.要收的莊稼多,作工的人少,你知否這需要越來越嚴重.如今教會缺乏工人,原因很多,但其中主要原因之一:父母不影響他們的兒女關心神的需要,鼓勵他們欣然答應主的呼召.反多有一面倒的,勉強兒女在世界道路上.甚至有見兒女走奉獻道路的,為他憂傷,認為失望!
\newline
\newline
神賜給我兩個聰明的兒子,按人的思想大兒至二十歲即可大學畢業.但我與妻同心向神多方禱告,與神立約:
\newline
\newline
a.神使他們怎樣都可以,但不能不信.
\newline
\newline
b.我願意盡量影響他們,使他們知道神國的需要；作傳道的喜樂,使他們能走上奉獻的道路.
\newline
\newline


\newpage

\section{第49屆~港九培靈會~張子華~第九講}
\label{sec:785}
\textbf{建立夫妻和諧生活}
\newline
\newline
連結: \href{https://www.hkbibleconference.org/session-message/view/785}{\texttt{https://www.hkbibleconference.org/session-message/view/785}}
\newline
\newline
\hyperref[sec:784]{< < < PREV SERMON < < <}
~
\hyperref[sec:toc]{[返目錄]}
~
\hyperref[sec:786]{> > > NEXT SERMON > > >}
\newline
\newline
創世記 1:26-1:27
\newline
\begin{longtable}{cl}
\hline
\hline
章節 & 經文 (和合本修訂版)\\
\hline
1:26 & \begin{tabularx}{0.7\textwidth}{X} 神說：「我們要照著我們的形像，按著我們的樣式造人，使他們管理海裡的魚、天空的鳥、地上的牲畜和全地，以及地上爬的一切爬行動物。」 \end{tabularx} \\ \\ \relax
1:27 & \begin{tabularx}{0.7\textwidth}{X} 神就照著他的形像創造人，照著神的形像創造他們；他創造了他們，有男有女。 \end{tabularx} \\ \\
[1ex]
\hline
\hline
\end{longtable}
在我過去十五年的牧會經驗中,對和諧生活,發覺問題最多的,是出自家庭.家庭的問題,又多因夫妻關係不正常而引起.昨天講家庭父母的責任,不要惹兒女的氣,不要讓家中任何事影響到兒女的身心的不健全.我提出三方面要滿足兒女:愛,身份,安全感,在這三方面都要使兒女滿足.今日請稍為深思,既然兒女有這三方面需要,則夫妻的本身,又何等需要有良好的關係,若夫妻的關係不好,就不能給兒女愛的滿足,也使兒女失去其身份,更缺乏安全感.昨天提到父母另一重大責任,以屬靈的真理教養兒童；使他們認識主,愛主,帶領他們有屬靈的鍛鍊,築家庭的祭壇；更要影響他們尋求神的旨意,走主的道路.如果夫妻不一致,關係破裂,就無法在屬靈上得增長.今日願意用餘下時間,思想一重要題目.如何建立夫妻和諧生活.請注意我所用字眼,「建立」而非「創造」.創造是使無變有.建立則不然,是先具備了材料和模樣,以之來建立夫妻和諧的生活.(來三3-4)提到神是建造萬物的建築者；夫妻關係,也必需承認神是和諧生活的建造者.(出廿五8)神吩咐製作一切事物,都要照天上的模型.現稍為提醒,神要人照著天上的樣式來製作.我們試看:父,子,聖靈,三位一體,和諧合一；這就是夫妻和諧生活的原則,要照從上頭啟示的模式.再次,了解建築和諧關係,必需明白材料.如果明白是那一種材料,就求神的憐憫.許多夫婦不和,因不明白材料的緣故?
\newline
\newline
1.每個人有他的特性:有男女之別,心,生理,思想,情緒,每個人對環境的感受,都有顯著的分別,這都要認識和接受.
\newline
\newline
2.若以此認識個人,每個人有他的性格.性格的定義不是作一事,是如何作一事,每人會用不同方法,去作同一件事.去年我住在美孚新村,也有同事住在那裡.雖然起點和終點都同,但採用的方法卻常不一樣；這就是個性的解釋.明白夫妻是兩個不同的性格,可用不同的方法,來作同一的事.應知任何一方,無資格要求對方同一樣,這是夫婦爭吵的主因.
\newline
\newline
3.我們都是失落的人:都如羊走迷,各人偏行己路.有這認識,就知道建立和諧生活之先決條件了:a.自己努力, b.神的憐憫.綜合上述,我們有建築師,有天上的模型,又知道了材料,就可循以下步驟.以下先與青年略談如何選擇伴侶,不能按下列條件作決定:
\newline
\newline
1.樣貌的吸引力.若單憑美麗,英俊作為結婚條件,就帶進危險的地步.我不能否認荷里活有許多美女明星,若外型美麗能夠作為婚姻幸福的保證的話；則娶了伊利沙白泰萊的男士,又怎能與她離婚.青年們應特別注意,內心的美善比外表的美麗更為重要.
\newline
\newline
2.不能用你們的愛情作決定,要明白你是何等材料──不過是失落的人,是偏行己路之人.若用愛來決定,一旦愛情改變了又如何?應該知道我們的愛是會改變的,甚至愛主的心也常動搖.這並非說作為夫妻,毌須有愛；乃是說,單憑彼此相愛就決定成婚是不可靠的.
\newline
\newline
3.不能以彼此的認識,作為先決的條件.有一項非常要的原則請勿忘記.若神要我結婚,則我在母腹之中,神已為我揀選伴侶.要對神信靠祂的安排,在創世之前,祂已為我揀選我的伴侶.我們要有堅強的信心站在其中.
\newline
\newline
下列有四個條件,作為揀選的原則:
\newline
\newline
1.必須是基督徒.神在新舊約聖經中,清楚啟示,基督徒選擇婚配對象,對方必須信仰相同.若果為結婚,讓我借用一位輔導者的話,不幸福的婚姻,比地獄還可怕!所以我們需要信靠神的帶領.特別有些已過了結婚年齡的人,神對他的配偶還未有安排,以致動搖信心；與不信者結婚,就與神的旨意違背了；你不能期望得神賜福.又用某人作生例子；來反對聖經上重要原則的.我認識一位姊妹,約卅歲時嫁與一未信者.可惜當時沒有人給適當的輔導,以致婚後十多年來,無日不爭吵度日.她後悔當時抱著「年晚煎堆」,人有我有的心態而結婚,以致夫婦生活苦不堪言.
\newline
\newline
2.能夠在屬靈事上同負一軛.一對真基督徒的戀愛生活,在一起時不應單是談情說愛；應在屬靈事上彼此分享,看對方的負擔和異象,彼此能否一致.「信與不信不能同負一軛.」可擴大來看；不應祗限於對方已經是信主的人,就不理一切與他結婚.要在基督徒的人生,前面的道路,工作負擔,領受的異象等著眼.從各方面來看,對方是否能夠與我同走天路.夫婦關係的重要原則:二人開始同走天路,必須抓緊.我曾為許多預備結婚的人作輔導,我會問他們有關十一奉獻,對教會,對事奉等見解.若發覺他們意見太分岐的話,會對他們提出警告,供他們重新考慮.
\newline
\newline
3.彼此在肉體上與愛情上,能互相吸引.有了以上兩者,然後加上雙方愛情和彼此的吸引.
\newline
\newline
4.聖靈之印證.每一位基督徒,與神有美好的關係,有聖靈居住在心中,對每件事加以印證；對與不對,這是個人的經驗,唯有你與神知道.當你充份地知道神的指示,就能堅定你的信心.創廿四章是非常美麗的愛情故事,重要的是聖靈的帶領和印證.老僕人帶著利百加回來,遇見以撒在田間默想；經上寫出二人的心境,有非常美麗的描寫.利百加非常欣悅地急忙下了駱駝,以撒領她進了母親的帳棚；得到了自母親不在後,從未得的安慰.讀這故事,了解到神的印證何等奇妙!若我們與神的關係完好,就能享用這原則,在每件事上了解符合神的旨意.
\newline
\newline
若我們妥善地運用上面原則,就進入夫婦關係之幸福中.但踏上這步之後,還需繼續努力.到此讓我們略談到教會一種錯誤觀念:認為每個人都需要結婚.有些姊妹很喜歡替人作媒,她會告訴人說:某某姊妹,又美麗,又有學問,人好,又有生命的流露......祇可惜!她未有男朋友.似乎不找個伴侶結婚便是大錯.你要看聖經是否這樣告訴我們?這一種觀念形成了無形的壓力,推使一些弟兄姊妹,隨便找伴侶結婚,這種壓力不應在神家中存在.我看有些獨身弟兄姊妹,神不在這方面預備；祂要好好用這器皿,是一般結了婚的人不能達成的.所以不需要太緊張,太熱衷替人撮合.每個青年應安靜等候,神必帶領你進到為你預備的境地,不要讓人加以干涉.
\newline
\newline
以上講過結婚,下面一段講婚後.是否一結婚,就事事如意,一切平坦呢?無人體會到,夫婦關係需要共同的努力,神特別的憐憫.既知道我們不過是失落的材料,就當靠神的恩典,努力追求,使幸福和諧生活每日增添.達致夫婦和諧生活的重要原則,乃二人成為一體.但這是一個奧秘.我必須引用天上的模型,三位一體之真神作為解釋.無人能識透,解釋,如何三位真神,又合而為一成為獨一的真神.他們的和諧,合作,分工,可以從聖經稍為得到亮光.夫婦二人成為一體,也可從這三一神的奧秘來比方,叫人可以領悟.現提出幾件事需努力:
\newline
\newline
1.按以弗所書,二人一體的夫妻生活,需要按照聖經,在家庭,在權柄的次序.丈夫乃是妻子的頭,這是神的模型.然而丈夫是妻子的頭是在一定條件之下的.不要誤解,妻子需順服丈夫如同順服基督.要解作:妻子順服丈夫如同丈夫順服基督.弟兄若要作家庭之頭,就要接受基督為你之頭.應知由這權柄所帶來之責任何等的大.他要帶領整個家庭,遵行神的旨意,同走天路；他要作一家屬靈的領袖,領導家人在靈裡長進,所以要謹慎對家庭所花時間,所下的工夫.男人需出外工作,或不免有應酬,也在教會負較重的責任；但也要知道神所給你對家庭的責任,你必須花時間在家人的身上,才能夠成就神的旨意.一個作牧師的人,會受到很大的試探,盡上自己力量,日夜苦幹.因為教會內有作不完的工夫,但我經禱告後,與教會有了協議.我必需有時間對我的家,盡帶頭的本份.做頭的意思,並非單指意氣用事,乃是需要盡上屬靈的責任.我對作妻子的說:你的丈夫,若是愛主遵行神旨意的人；你當無條件地將自己交託在他手中,作順命之妻子.若丈夫未符合上述的標準,作妻子的要努力幫助他；使他成為一家之頭.若感覺到丈夫沒有這資格,而自己作頭的話,這是大錯!一家有兩個頭,這是不和諧的主要原因.
\newline
\newline
2.需要有單純,潔淨的愛情,在兩人之間.這愛包括捨己,為對方而活；有虧欠時原諒對方,包容對方的軟弱等.接納對方好像基督接納我.讀弗五章感到一件奇事,祗寫丈夫應當愛妻子,如同基督愛教會.未見寫妻子要愛丈夫,由此感到神的智慧實在奇妙!聖經這樣寫有兩個理由:a. 神知道丈夫愛妻子比妻子愛丈夫艱難.男子在外工作,接觸人多,試探大,愛情較易變遷,在性方面跌倒的機會更大.b. 在創造計劃上,女性屬於被動,若丈夫能對妻子真誠相愛,則不難得到愛的回應.愛情之表現在家庭中極其重要,叫子女也能得到愛情,是很需要有外表行動來表明.約略地誘露點我家庭的事,使大家得到幫助.我的小兒子六歲時,一天當我們在客廳休息的時候,他忽然對我們說:「爸爸,為什麼從不見你吻媽媽?」又問母親同樣的問題.雖然他在美國長大,略與中國保守思想有點出入.但夫妻和諧生活,需有愛情外表的流露.不獨夫妻雙方需要,而兒女也需有此愛來滿足.縱使年華老大,也需此愛的表現,產生和諧幸福.
\newline
\newline
3.夫婦需有敞開的心,無攔阻的交通.沒有一秘密不能彼此分享；沒有一重擔不能彼此分擔.夫婦難處的形成,每每由於交通中斷.為何在婚前兩人的交通甜蜜而投契,而婚後這情形漸漸消失?在神計劃中,原應隨著年日增加越發增長.夫妻豈不應多找機會,彼此敞開心懷,傾訴內心的負擔；然後藉著禱告,帶到神面前.家庭要建立祭壇,夫妻每日需有共同禱告.如果一晚你覺得不能共禱時,勸你勿去睡覺；寧可敞開你的心,去解決你夫妻間之難處.一晚我和妻子討論一些問題,有點爭持不下,就各自睡去,但聖靈在我們心中工作；不約而同地轉過來,相視而笑,起來祈禱然後再睡.
\newline
\newline
4.要知撒旦在心中有工作的可能.撒旦破壞之工作,可以將小問題,擴大至不可收拾.所以非常需要設法約束節制.若一次夫妻反目,就讓撒旦在你夫妻關係上,留下了工作的地步；故此應該竭力避免,要約束自己.
\newline
\newline
5.學習喜歡對方的興趣.各人有不同的愛惡.但當進到二人一體的真理時,就需努力學習適應對方,培養愛對方所愛的事物.若果結婚數十年,夫妻依然故我,分道揚鑣實非好事.我夫妻二人個性大不相同；妻愛音樂,較趨文靜；我是老粗,愛著打球.起初較難遷就,但逐漸大家努力發展共同興趣,就建立起家庭的共同生活.想不到有一日,妻主動要我帶她看球賽,我又要她和我去聽音樂.我輔導學識也因妻子修這科而得來,共同生活是夫婦生活的秘訣.
\newline
\newline
6.從聖經觀點對性的認識.若祗求滿足自己,是錯誤的看法.若從聖經的角度,就知神對此之安排,藉此將整個人奉獻給予對方,為對方而活；如此就能保持肉體關係的新鮮.你夫妻關係是否和諧?是否知道不和諧的夫妻,會為社會造不正常的青年.神看家庭比教會更重要,有健全的家庭才有建全的教會.
\newline
\newline
要齊心合力建立健全的夫妻關係,裡面要有尊重,服事,感謝,安靜之靈.
\newline
\newline


\newpage

\section{第49屆~港九培靈會~戴紹曾~第一講　初期教會的耶穌}
\label{sec:786}
\textbf{第一講　初期教會的耶穌}
\newline
\newline
連結: \href{https://www.hkbibleconference.org/session-message/view/786}{\texttt{https://www.hkbibleconference.org/session-message/view/786}}
\newline
\newline
\hyperref[sec:785]{< < < PREV SERMON < < <}
~
\hyperref[sec:toc]{[返目錄]}
~
\hyperref[sec:787]{> > > NEXT SERMON > > >}
\newline
\newline
當我們讀到新約聖經第五卷──使徒行傳時,我們就以為這卷書的主角是一些使徒.因為本書上半卷是記載彼得的事奉,後半卷是論及保羅的事奉；主藉著他們的見證使福音廣傳.有人說,本書應命名為聖靈行傳,這樣的建議不錯；因為從第一章至廿八章,使我們清楚看見聖靈的工作,完全照著耶穌基督的應許,聖靈降在祂的門徒身上,在聖靈引導之下,門徒得著能力去傳福音.
\newline
\newline
當我們仔細查考使徒行傳時,我們就看見重點不在於聖靈而是在於主耶穌身上.路加在使徒行傳一章一節,就把它和路加福音連接一起,說:「提阿非羅阿,我已經作了前書,論到耶穌開頭一切所行所教訓的,」前書乃指路加福音,是關於耶穌降生,耶穌在世上的事工,耶穌的死和復活,都一一記載；可見後書是繼續講到耶穌.耶穌自己說:「只等真理的聖靈來了,祂要引導你們明白一切的真理......祂要榮耀我,因為祂要將受於我的,告訴你們.」使徒行傳第一章,是福音和教會的橋樑,也是路加福音和教會歷史的橋樑,也可說是耶穌在世之事工的結論.
\newline
\newline
耶穌復活向祂的門徒顯現,是使徒行傳非常要的信息.路加告我們,在這時間,四十天之久耶穌活活向祂的門徒顯現.復活的真理,是使徒行傳開始的重點.
\newline
\newline
耶穌復活四個證據:
\newline
\newline
(一)空的墳墓
\newline
\newline
當我們看到那空的墳墓,我們將作如何解釋?按邏輯只有三個不同的解釋:第一是耶穌的門徒把耶穌的屍體帶走.這可能嗎?事實上這班門徒對耶穌的復活毫無盼望；雖主曾對他們說:「我要復活.」但這信息沒有進入他們心裡,他們根本不明白,甚至到了最後,他們還不信主已復活了.第二,當時彼拉多安排許多兵丁嚴密守衛；倘若門徒要偷竊主的屍體,必須大力對抗衛卒.第三,此點更重要,如果耶穌沒有復活,門徒將屍體帶走,然後秘密收藏,那麼,寫下的歷史還可能嗎?請大家思想,若門徒明明知道耶穌沒有復活,是他們作假,這樣,門徒肯不肯一直受苦下去呢?為了一句謊言,為了一件虛假的事,許多門徒不可能一直甘心受苦.由此三點,可見耶穌復活並非門徒所造成的事.有人說是敵人把耶穌的屍體帶走；但我們要問:
\newline
\newline
1.敵人為何要將耶穌屍體帶走?事實上,當時是他們把穌的屍體帶走,並且請求彼拉多派遣士兵嚴加守衛的.
\newline
\newline
2.如果是那些敵人所為,當門徒開始講耶穌復活之道時,敵人儘可把耶穌的屍體擺出,使人知道耶穌沒有復活,來證明門徒的見證是假的；由此可見不是敵人所為.根據我的了解只有一個可能:耶穌復活,既不可能是耶穌的門徒所作,也非敵人所為；聖經記載耶穌從死裡復活了,耶穌復活的第一個明證是空的墳墓.
\newline
\newline
(二)耶穌向門徒顯現
\newline
\newline
聖經至少有十次題到耶穌活活地向祂的門徒顯現；向馬利亞顯現,向以馬忤斯路上的倆門徒顯現；當十一個門徒聚集的時候,復活的主在他們中間；還有一次共有五百男女同時看見耶穌.從這許多次耶穌顯現,我們知道耶穌復活了.
\newline
\newline
(三)門徒的改變
\newline
\newline
門徒情緒的改變,態度的改變.當主耶穌被釘在十字架上後,門徒們就完全失去盼望；甚至在以馬忤斯路上走的倆門徒,他們以為遇見的是位陌生人.他們對陌生人說:「我們素來所盼望要贖以色列民的就是祂.......」現在我們卻看見這班門徒絕望的態度轉變過來,滿有盼望；甚至為復活的主而死也願意.他們竟然敢把猶太人的安息日改為主日；而且在主日聚會,記念主復活.猶太人是不容易改變固有傳統的；但那些基督徒,卻敢主日聚集在一起記念主的復活.
\newline
\newline
(四)門徒傳講福音
\newline
\newline
主耶穌復活不但叫我們看見空的墳墓,看見主多次向門徒顯現；看見門徒的改變；更使我們看見門徒傳講福音.耶穌復活成為他們傳福音最重要的一環,他們引證所傳的是真實可信的,耶穌復活就成為他們信息中最重要的部份,第一章是重要的橋樑,因為復活的主與祂的門徒說話.祂把賜下聖靈的應許賜給他們,並吩咐他們在耶路撒冷等候,直到聖靈降臨下來.最後還吩咐他們:「但聖靈降臨在你們身上,你們就必得著能力,並要在耶路撒冷,猶太全地,和撒瑪利亞,直到地極作我的見證.」(徒一8)路加寫行傳一章,就把耶穌的復活和顯現,並應許賜下聖靈,這三件事擺在我們面前.耶穌吩咐門徒傳福音作祂的見證,然後升天去.當耶穌升天時,這些門徒還不甚明白,他們呆呆駐足望天,突然兩個天使站在旁邊提醒他們:「你們見祂怎樣往天上去,祂還要怎樣再來.」
\newline
\newline
耶穌升天之後,聖靈未降臨以先,第一個應許就是「主必再來.」保羅在他的書信中,特別提到主再來.
\newline
\newline
從門徒的見證──他們所傳的道,我們可發現有幾個重點,特別是指著主耶穌的:
\newline
\newline
一,耶穌是主
\newline
\newline
如今福音廣傳,世界每角落都有人傳福音；不過第一篇信息卻是彼得傳的.五旬節到了,彼得站起來講道,眾人以為這些人喝醉了；彼得說我們不是喝醉了,乃要應驗聖經的預言,在末後的日子,神的靈要充滿凡有血氣的.他的信息的最高峰,乃是傳揚耶穌是主.到了第十章,彼得在哥尼流家裡領會,這是使徒在外邦人之家傳福音的開始；彼得仍然傳揚耶穌是主.無論在第二章36節及第十章36節,都是強調耶穌基督是主.　我們讀使徒行傳,發覺「主」字出現最多次數.有時我們不大清楚,到底「主」是指聖父還是指聖子呢?若細讀,就知各有所指；有時指聖父,但多次是指著耶穌.在舊約,「主」是指耶和華.當時彼得傳信息,聽眾多數是猶太人,他們熟識舊約聖經的背景；彼得講到主,他們就會聯想到舊約的信息.彼得所強調的是主耶穌,同時彼得更強調「主」是有權柄的.在第一章,記載門徒問耶穌說:「主阿!你復興以色列國就在這時候嗎?」是對主客氣的稱呼.但是到了第二章,彼得宣佈耶穌是主,這個稱呼就帶著新的意義了.耶穌基督在彼得和其他門徒身上有祂的主權,門徒講道時喜歡作比較,一方面講耶穌,一方面讓眾人對耶穌的看法和態度提出來,另方面講到神對耶穌的看法迥然不同.五旬節的信息彼得對眾人說:「你們釘在十字架上的耶穌,神已經立祂為主為基督了.」彼得又站在那裡說:「你們把祂釘在木頭上,」熟悉舊約聖經的人,都曉得把一個人掛在木頭上是一種極羞辱和被咒詛的記號,彼得說了這話,強烈的反應立刻出現了；神將耶穌高舉起來,不但叫祂從死復活,叫祂升天,並叫祂坐在自己的右邊；可以說:神叫祂進到極大的榮耀裡.弟兄姊妹!請你看看這個對比；敵人用最卑鄙的方法來羞辱耶穌；但父神卻賜給祂最高的榮耀.
\newline
\newline
二,耶穌是神的僕人
\newline
\newline
門徒的信息中,不但強調耶穌是主,彼得特別強調耶穌是神的僕人.他們在一處禱告時,更把「聖潔」兩字加在僕人上面:「聖僕耶穌」這字在舊約聖經有很深的背景和意義；以賽亞五十三章,先知看見這位受苦的僕人,也用預言講到耶穌是無罪的,為了眾人的罪接受痛苦.　埃提阿伯的太監讀至這段聖經時,他不明白,後來腓利給他講解.　耶穌是受苦的僕人,彼得是要聽眾知道耶穌來世,不但是帶著使命,而且還有要完成神的旨意.馬可福音特別講到耶穌是僕人,祂到世上是要作眾人的僕人.彼得也強調此點,他題及兩件事,一是摩西的預言,摩西對以色列人說:「神要從你們弟兄中間興起一位先知像我,你們要聽他.」(申十八5)意思是神興起一位像我的先知時,祂將帶來新的啟示,讓你們對神有新的認識.耶穌已經完成了這個使命.一是亞伯拉罕的預言；他得到神的應許,神對他說:「我必賜福給你......地上的萬族,都要因你得福.」(創十二2-3)彼得告訴我們,這應許不但應驗在亞伯拉罕身上,同時他的後裔「耶穌」也把這福帶給全世界.(加三16)耶穌是聖潔的僕人,是帶著啟示的僕人,更是帶著祝福的僕人.
\newline
\newline
三,耶穌是救主
\newline
\newline
彼得和其餘的門徒不但強調耶穌是主,祂是救主.彼得和那些宗教領袖最初的衝突,(四至七章)是因彼得和約翰醫治了一個瘸腿的人,當他們二人進聖殿禱告時,看見這可憐的人坐在那裡,就奉耶穌的名醫好他；這人便起來行走,跳躍,和他們一齊進入聖殿敬拜.因此,這人從殘廢,不自由,而轉變成個自由,釋放,得拯救的人.於是引起了眾人的注意,彼得藉此機會講論耶穌是救主,可是這班領袖絕不以此為然,他們認為在聖殿禱告是可以的,但是奉耶穌的名講道決不容許,他們就開始禁止,捉拿彼得和約翰下到監獄.當時他們講到救主來的時候,基本使命是要拯救祂的百姓脫離罪惡.正如聖經記載:天使對約瑟說:「馬利亞將要生一個兒子,你要給他起名叫耶穌,因他要將自己的百姓,從罪惡裡救出來.」(太一21)彼得清楚地講,耶穌是作世人的救主,是唯一的救主,除祂以外,沒有別的名我們可以靠著得救.(徒四12)從第三章至第五章,十二次題到耶穌的名,耶穌的名就大大地傳開了,宗教領袖想要鎮壓,禁止傳揚耶穌的名；但彼得認為非傳耶穌的名不可；因天下人間沒有賜下別名我們可以靠著得救.門徒早已向全世界宣佈,唯一的救法就是耶穌基督,祂是世人的救主.如果今晚還有一位未信耶穌的,巴不得你從今能認識這位從死裡復活的救主.
\newline
\newline
在台灣有個名為五教廟的,有一次我路經該地,很想知道這廟的究竟,我就進去參觀.原來他們拜佛,拜孔子；也有回教,道教；我發現櫃內有個像,手持聖經的.我問:「這是誰?」答:「是耶穌.」五教廟竟然有耶穌在內,據他們的廟書載,耶穌是偉大的老師.我覺得很矛盾!耶穌說祂是道路；彼得說祂是唯一的救主；怎能與其他宗教相比,相連呢?　當時彼得勇敢站起來對那班宗教領袖說:「我們所看見的,不能不說.」(徒三20)耶穌是救主,祂不但是我們的救主,是猶太人的救主,也是全世界人類的救主.
\newline
\newline
四,耶穌復活要審判活人死人
\newline
\newline
使徒行傳有兩次題到未來的審判.一次是在哥尼流的家,一次是在雅典地方.在雅典地方是應該提的,因雅典是個敗壞,罪惡滿城之地；但哥尼流卻是個虔誠的人,禱告的人.不過一個虔誠禱告的人,也是同樣需要救主,不能靠自己的善行得救；並且哥尼流也需要知道,到了世界末日,耶穌基督要再來.彼得,保羅兩次講到耶穌的審判,特別說明耶穌復活,證明祂是審判活人,死人的主.這是復活的真理,因耶穌復活了,死人也要復活,然後有審判.
\newline
\newline
耶穌對門徒說:「你們要去使萬民作我的門徒......我就常與你們同在,直到世界的末了.(太廿八19-20)在使徒行傳,我們看見復活的主成就了祂的應許,不斷與門徒同在.當司提反面對眾多敵人時,他們用石頭要打死他,他看見了耶穌,他所看見的是站起來的耶穌；彼得告訴我們,耶穌升天之後是坐在父神的右邊.意思是工作完成了,帶著榮耀高升了.司提反看見復活的耶穌,是站在父神右邊,帶給他無限的安慰.還有保羅在哥林多時,傳福音的工作極其艱難,全城的人起來攻擊他,攔阻他傳福音.當天,主耶穌向保羅顯現,對他說:「不要怕,只管講,不要閉口；有我與你同在.」(徒十八9)還有一次保羅回到耶路撒冷,他的同胞捉拿他,要打死他,後來從監獄解去羅馬,在大海中遇風浪,正危急時；復活的主再次向他顯現.安慰他說:「保羅!不要害怕︳你必定站在該撒面前,並且與你同船的人,神都賜給你了.」(徒廿七24)
\newline
\newline
親愛的弟兄姊妹!今天我們在福音的工作上,復活的主與我們同在；無論我們是遭受許多患難時,被逼迫時,主都與我們同在.
\newline
\newline
另一方面,我們看見復活的主初次向保羅顯現,那時,他是逼迫教會的,是反對主的；但主耶穌向他顯現,開他的心眼,叫他悔改,成為新造的人.福音初傳入歐洲時,保羅在腓立比的河邊講道,呂底亞聽見了,聖經說:「主開導她的心.」她就信了,感謝主!祂把人的成見除掉,把人剛硬的心改變,賜給人新的生命,在祂裡面成為新造的人.
\newline
\newline
復活主在門徒的工作上,不斷地幫助他們,主藉他們的手,施行神跡奇事,證明祂的恩道.(徒十四3)可見他們的工作,有主的能力和主的同在.復活的主,祂的應許永不落空.
\newline
\newline
多年前,有位傳道人騎馬到某村莊,下馬到一房屋叩門；有婦人應門,傳道人問她:「耶穌是否住在這裡?」婦人莫明其妙,不知如何作答；傳道人再三問她:「耶穌是否住在這裡?」婦人始終沒有回答,於是傳道人騎馬離去.後來婦人的丈夫回家,晚膳時,婦人向丈夫提起剛才的事,她的丈夫問她如何回答?婦人說:我不知所答.夫說,你為何不告訴那傳道人,我們每主日到附近禮拜堂禮拜,我們已在那禮拜堂受洗,而且我是該教會的執事.婦人說:他都沒問這些問題,他只問:耶穌是否住在這裡?
\newline
\newline
弟兄姊妹!初期教會,有主同在,在艱難的情形下,在福音的工作上,主與他們同在.今天復活的主也要成為我們見證的信息.因耶穌復活了,我們才有新的生命；因耶穌復活了,我們的同胞才有盼望.
\newline
\newline
求主幫助我們!高舉耶穌,好像初期的教會一樣.
\newline
\newline


\newpage

\section{第49屆~港九培靈會~戴紹曾~第二講　聖靈和祂的工作}
\label{sec:787}
\textbf{第二講　聖靈和祂的工作}
\newline
\newline
連結: \href{https://www.hkbibleconference.org/session-message/view/787}{\texttt{https://www.hkbibleconference.org/session-message/view/787}}
\newline
\newline
\hyperref[sec:786]{< < < PREV SERMON < < <}
~
\hyperref[sec:toc]{[返目錄]}
~
\hyperref[sec:788]{> > > NEXT SERMON > > >}
\newline
\newline
昨晚我已提過,有人認為使徒行傳可以稱為聖靈行傳；因為從第一至廿八章都是論及聖靈的工作.第一章二節首句「藉著聖靈,」可說是使徒行傳的鑰節；甚至有解經家說是本卷書神學方面的中心.
\newline
\newline
聖靈是誰?是我們今晚須要認識的.
\newline
\newline
(一)聖靈是有位格的
\newline
\newline
聖靈非抽象的勢力,也非一種主觀的感覺；雖然本書常提到聖靈的能力,但祂卻不是勢力；雖然本書也常說到聖靈的恩賜,但祂卻不是物件.祂是具有位格的,絕非可怕的鬼怪現象.
\newline
\newline
在第五章,教會發生了一件醜事,就是亞拿尼亞和他的妻子撒非喇說謊.彼得對他說:「撒旦充滿了你的心,叫你欺哄聖靈.」(3節)「你們為甚麼同心試探主的聖靈呢?」(9節)從這兩句話,可知聖靈是有位格的,不容人向祂撒謊,試探.彼得向那班宗教領袖講道說:「我們為這事作見證,神賜給順從之人的聖靈也為這事作見證.」(徒五32)可知聖靈是有位格的.保羅在(弗四30)說:「不要叫神的聖靈擔憂.」由此更可知聖靈是有位格的神.
\newline
\newline
(二)聖靈是神是有神性的
\newline
\newline
五章3節彼得說:「......你欺哄聖靈,......」接著,第4節末句「是欺哄神了」足證聖靈是神,是有神性的.
\newline
\newline
有人以為在四福音書中,耶穌曾應許要賜下聖靈；但在使徒行傳之前是沒有聖靈的；這樣的想法是錯誤的.我們不僅在使徒行傳中看見聖靈的工作；在四福音中,更看見耶穌受洗時,聖靈降臨在祂身上.彼得和保羅在使徒行傳中,也講到聖靈在舊約時代的工作.聖靈感動大衛；也感動以賽亞說話.可見舊約時代聖靈也作工；創世記一章2節也說:「神的靈運行在水面上.」神創造天地之初就有聖靈的工作.聖靈有位格,是神,有神的神性,也是永遠的；不但在現時,在使徒時代,在耶穌時代,以至在舊約時代,創世以來,我們都看見聖靈的工作.
\newline
\newline
聖靈充滿居住在聖徒心裡(林前六19)保羅寫信給當時的基督徒,叫他們過聖潔的生活,盼望他們對自己有更清楚的認識.他說:「豈不知你們的身體是聖靈的殿麼?」舊約以色列人有會幕,有聖殿；但是現今聖靈都居住在我們身體裡面.巴不得在座弟兄姊妹,清楚認識這個真理.
\newline
\newline
路加醫生是敘述歷史的作者,他用許多不同的動詞；來形容聖靈在我們裡面,在聖徒裡面.現在讓我簡單地提出七個關於聖靈住在我們裡面的不同動詞:
\newline
\newline
1.聖靈的洗
\newline
\newline
約翰曾說:「我不是基督,是奉差遣在祂前面的,」(約三28)耶穌說:「約翰是用水施洗；但不多日,你們要受聖靈的洗.」(徒一5)
\newline
\newline
2.聖靈的降臨
\newline
\newline
「但聖靈降臨在你們身上,你們就必得著能力.」(一8)這裡特別給我們看見,聖靈從天而降,是神所賜降在我們身上的.
\newline
\newline
3.聖靈充滿
\newline
\newline
「他們就都被聖靈充滿.」(二4)這是講到五旬節那一百廿人聚在一起,同心合意禱告,聖靈降臨充滿他們的心.
\newline
\newline
4.聖靈的澆灌
\newline
\newline
(二17)彼得講道特別引用舊約,約珥的經文論到聖靈的澆灌.這裡有很深的背景,聖靈的澆灌,也是指著聖靈降臨在他身上的事.
\newline
\newline
5.聖靈的恩賜
\newline
\newline
恩賜本非動詞而是名詞；但我必須解釋,我發現中文不容易解釋這個字,且較容易引起誤解；中文和英文譯本有許多不同,和原文也有許多不同之處.中文很微妙,單數與複數不分,必須注意上下文才能明白；英文只要加上S便是複數.聖靈的恩賜,有兩個清楚不同的解釋；如果是單數乃指一件事,如果是複數,意思就完全不同了.照第十章45節:「那些......信徒,見聖靈的恩賜也澆在外邦人身上,就都希奇.」這裡聖靈的恩賜,到底是單數或複數呢?或者弟兄姊妹以為是聖靈所賜各樣不同的恩賜,有各種不同的表現；其實不然,這裡所提是單數的.第二章38節,譯者幫助我們,沒有用「聖靈的恩賜」諸字,而用「領受所賜的聖靈」,按照原文是「聖靈的恩賜」.總之,當這句話是單數時,乃指聖靈自己,是神的應許,成為一件禮物賜給人；是聖靈的本身,不是別的表現；神將聖靈賜給信祂的人.
\newline
\newline
6.受聖靈
\newline
\newline
(八15)彼得發現撒瑪利亞的人,還未受聖靈.保羅到了以弗所,問那些門徒:「你們信的時候,受了聖靈沒有?」(十九2)聖靈的洗,聖靈的降臨,聖靈的充滿,聖靈的澆灌,聖靈的恩賜和受聖靈,這許多不同的字眼,基本上是同指一件事──聖靈居住在信徒裡面.耶穌對門徒說:「聖靈要降臨在你們身上.」到了第二章,這件事就應驗了,百廿人等候時；他們受聖靈,被聖靈充滿.在行傳中,路加有五次講到聖靈,充滿一些人；其實聖靈充滿不僅五次.在許多次聖靈充滿之中,有五次是一些信的人,他們從未受過聖靈,從未曾被聖靈充滿的,當時他們就被聖靈充滿了.第一次門徒共百廿人被聖靈充滿.(二1)耶穌的應許應驗在他們身上.第二次被聖靈充滿的,不是猶太人,是撒瑪利亞人.他們聽了腓利講道,受感動相信了耶穌；當彼得為他們禱告,他們就受聖靈,被聖靈充滿.聖靈的工作擴充了,首先受聖靈的是猶太人,繼而伸延到撒瑪利亞人.照著一章8節耶穌所說的:「要在耶路撒冷,猶太全地,和撒瑪利亞......作我的見證.」當他們作見證時,聖靈隨著他們；在撒瑪利亞有信徒被聖靈充滿,是從來未有的經歷.第三次記於九章17節,這裡所提極為簡單；掃羅(後來的保羅)在大馬色的路上,去的時候,是個逼迫主的基督徒,當他遇見耶穌,他的生命全改變了!他忘記了原來的目的,他進到大馬色竟然在那裡禱告；後來神對在大馬色信徒亞拿尼亞說:「你去按手在掃羅身上；對他說:在你來的路上,向你顯現的主就是耶穌,打發我來,叫你能看見,又被聖靈充滿.」第四次載於第十章44節,這裡我們又看見福音廣傳了,從猶太人,撒瑪利亞人,延伸到外邦人,神差遣彼得到哥尼流的家,向他傳講耶穌基督的福音；當彼得講道時,哥尼流和他的家人都被聖靈充滿.第五次,我們又看見福音傳到以弗所,(十九6)保羅到那裡要開始傳福音,遇見十二門徒,他們已受了約翰的洗,他們雖信,但未曾聽見聖靈；保羅很關懷地問:「你們信的時候受了聖靈沒有?」回答說:「沒有,也未曾聽見有聖靈賜下來.」「保羅按手在他們頭上,聖靈便降臨在他們身上.」(6)弟兄姊妹!在使徒行傳廿八章中,這五次是一些人空前的被聖靈充滿,當他們聽了福音相信耶穌時；聖靈就降臨在他們身上,他們被聖靈充滿了.
\newline
\newline
聖靈的工作,在這些不同章節的記載中,顯有不同之處:
\newline
\newline
1.地點不一
\newline
\newline
我們不清楚五旬節聖靈在何處降臨,當然是在耶路撒冷,但解經家們於地點各有不同的解釋,有的說:「在那樓房上他們聚集禱告的地方.」有的說:「在聖殿.」在五旬節聖靈充滿那一百廿人,他們是在耶路撒冷；到了第八章,不在耶路撒冷,也不在聖殿,是在大馬色,該撒利亞,以弗所,這些不同的地方,有時在信徒聚會之處,有時在家中.
\newline
\newline
2.人數不一
\newline
\newline
有時一百廿人,在以弗所是十二人；在哥尼流的家,是哥尼流和他的家人；在大馬色僅掃羅一人.
\newline
\newline
3.情形不同
\newline
\newline
有時門徒按手,(大概有三次)在哥尼流的家,彼得根本不及按手；他正講道之中,聖靈就降臨在哥尼流的家；所以聖靈降臨,並非一定按手,也不一定要說方言.總之,我們不能限定聖靈的工作,不能說祂必有一定的表現.從上面五處經文,我們可以見到一個真理,就是我們不能限制聖靈的工作.耶穌說:「風隨著意思吹,你聽見風的響聲,卻不曉得從哪裡來,那裡去；凡從聖靈生的,也是如此.」(約三8)可惜今天教會,要想盡辦法來限定聖靈的工作.這是錯誤的.從使徒行傳,我們看見有很多不同的現象,不同的情形；可是有兩件相同的事,我們切勿忘記,且須緊緊抓住.第一,聖靈自己充滿在那些聖徒裡面,居住在他們裡面；不論在五旬節,在以弗所,在哥尼流的家,作工充滿他們的同是聖靈.第二,如果你仔細考查聖經,可發現每次被聖靈充滿之後的人,必開口為耶穌作見證,傳講福音的.「都聽見他們......講說神的大作為.」(二11-14)提及彼得開口傳福音,到了第九章,當保羅被聖靈充滿時,掃羅被改變成為保羅.「就在各會堂裡宣傳耶穌,說祂是神的兒子.」(20)第十章46節說:「因聽見他們說方言,稱讚神為大.」請特別注意末句「稱讚神為大」.他們說方言,為的是要稱讚神為大；在哥尼流家裡,福音廣傳了；他們被聖靈充滿之後,都開口傳講耶穌基督,高舉耶穌,讓眾人對耶穌有新的認識.
\newline
\newline
聖靈充滿的意義
\newline
\newline
(一)聖靈的工作和「得救」有關
\newline
\newline
在行傳中,路加用許多章節加以說明.第一件事是得救的問題.第十五章記載,使徒和長老在耶路撒冷開個很重要的會議,討論人到底怎樣得救?要不要靠摩西的律法?引起這問題的是那班猶太人.彼得說:「我們得救,乃是因主耶穌的恩,和他們一樣,這是我們所信的.」(11)「知道人心的神,也為他們作了見證；賜聖靈給他們,正如給我們一樣.」聖靈沒有分別,我們得救是因信,他們得救也是因信.
\newline
\newline
(二)聖靈的工作和「潔淨」有關
\newline
\newline
在第十五章特別提到,聖靈潔淨他們的心.弟兄姊妹還記得嗎?第二章曾提到五旬節聖靈降臨,有舌頭如火焰顯現出來,分開落在一百廿人個人的頭上.(1-3)在聖經中,火是預表潔淨的工作.約翰說,「耶穌來了,祂要用聖靈與火給你們施洗.」(路三16)彼得往哥尼流家之前,在約帕地方他需要啟示；神向他顯現,給他看見異象,要除掉他心裡的成見；因彼得是猶太人,他有民族優越感；看不起外邦人,所以神用異象幫助彼得.當彼得「看見天開了,有一物降下......裡面有地上各樣四足的走獸和昆蟲,並天上的飛鳥；又有聲音向他說,彼得,起來,宰了吃.彼得卻說,主啊!這是不可的,凡俗物,和不潔淨的物,我從來沒有吃過.第一次有聲音對他說,神所潔淨的,你不可當作俗物.」(十11-15)這種異象幫助彼得明白聖靈的工作；不僅傳福音給猶太人,也應該給撒瑪利亞人分享.
\newline
\newline
聖靈的工作帶來新的生命
\newline
\newline
彼得從哥尼流的家回耶路撒冷,有些基督徒追問他為何進入外邦人的家,彼得見證說:「我是見了異象遵命而去的,是主打碎我個人成見,在聖靈帶領之下去.當我講道時,聖靈降臨在那些外邦撒瑪利亞人身上,神叫他們悔改得生命了.」(十一18)這件事和五旬節的預表是相同的；剛才所說的火焰,就是在五旬節的時候,不但看到火焰,而且聽見天上有聲音下來,好像一陣大風吹過；「風」有很深的預表,是代表聖靈.記得舊約以西結卅七章,神的先知看見一片荒涼,遍地骸骨,神問先知「這些骸骨能復活嗎?」按常理是不可能的；但後來一陣大風吹過,帶著新的生命,死人竟然都活了,成為極大的軍隊.風也預表新的生命.主耶穌對尼哥底母說:「你必須重生.」他不明白,耶穌向他解釋時也提到「風隨著意思吹......不曉得從哪裡來,往哪裡去；凡從聖靈生的,也是如此.」　聖靈的風帶來新的生命,聖靈就居住在其中,因為這身體要成為神的殿.當聖靈居住在我們裡面時,自然而然有聖靈所結的果子,有聖靈的恩賜.弟兄姊妹!我們從這幾處經文,發現每事發生非單在門徒身上,在教會領袖身上；而且凡信的人都有這樣的經歷,得救時,聖靈作潔淨的工作,賜給新的生命,居住在其中充滿他們,結出靈果來.
\newline
\newline
或者有弟兄姊妹要說,在使徒行傳中,不僅此五次聖靈充滿他們；不錯,當門徒遭遇困難時,有特別需要時,聖靈再次充滿,但和以前的情形不同,這是一種特別的充滿,是為一時當前的需要.例如彼得站在那些宗教領袖面前,他們控告彼得；那時他需要膽量和智慧來與他們對答,遂被聖靈充滿了.這完全符合耶穌復活前所應許的,耶穌對門徒說:「我離開你們,將來你們到各處為我作見證,有時坐監,有時被帶到公會,你們受審時,不要怕,聖靈將賜給你們當說的話.」我相信弟兄姊妹,可站起作此見證,在任何時候；可能在家中,在學校,在極其困難時,忽然有聖靈的感動,聖靈將當說的話賜給我們.耶穌說的話在行傳中應驗了.
\newline
\newline
聖靈賜安慰
\newline
\newline
有時門徒們需要安慰,他們所受的壓力太大.感謝主!在基督徒的路程中,有時遭遇困難,四面楚歌；我們需要幫助,聖靈居住在我們裡面,隨時隨地賜給我們安慰.
\newline
\newline
聖靈賜喜樂
\newline
\newline
在各種情形下,信徒被聖靈充滿時大有喜樂.保羅告訴我們,聖靈所結的果子,第一是仁愛,第二是喜樂.當保羅和巴拿巴在安提阿被趕出去時；聖經說:「他們被聖靈充滿,」就是被喜樂充滿.感謝主!在困難時,不但有聖靈的安慰,還有喜樂.
\newline
\newline
聖靈不斷地引導
\newline
\newline
聖靈不但賜給門徒安慰和喜樂,還不斷引領他們的道路,無論何時何地,聖靈一直都在帶領他們.
\newline
\newline
今晚我們特別看到行傳中,耶穌基督的應許應驗在早期教會；在不同的地方,他們相信了耶穌,被聖靈充滿,聖靈居住在他們裡面.
\newline
\newline
弟兄姊妹!讓我們思想,你的身體是聖靈的殿；言語,行為,心的深處,思想,態度,都要被聖靈管理；因為聖靈要得著我們,潔淨我們,居住在我們裡面；舊約的聖殿已經沒有了,現在聖靈要住在你裡面；雖然我們不配,但祂仍然要住在我們裡面.請問!聖靈是否完全得著你?你的身體──是否合乎聖靈居住?
\newline
\newline
從神的話語中,從初期教會的見證,看見聖靈住在他們裡面；有新的生命,有潔淨的生活,有能力的見證.
\newline
\newline
但願主的靈今天晚上在你我心裡,在我們身上運行,完成神的美意.
\newline
\newline


\newpage

\section{第49屆~港九培靈會~戴紹曾~第三講　被聖靈充滿的司提反}
\label{sec:788}
\textbf{第三講　被聖靈充滿的司提反}
\newline
\newline
連結: \href{https://www.hkbibleconference.org/session-message/view/788}{\texttt{https://www.hkbibleconference.org/session-message/view/788}}
\newline
\newline
\hyperref[sec:787]{< < < PREV SERMON < < <}
~
\hyperref[sec:toc]{[返目錄]}
~
\hyperref[sec:789]{> > > NEXT SERMON > > >}
\newline
\newline
在使徒行傳中,影響保羅最深的,有兩位信徒.他們不是門徒中的一個,也不是教會中的柱石或任何人.保羅在加拉太書特別聲明他所領受的不是從人來的,而是直接從主領受的.但是他的見證,並非說沒有人影響他.影響他的在行傳中至少有兩位信徒.
\newline
\newline
我曾提過,「聖靈充滿」是個動詞,是一種動作.在五處地方,聖靈降臨,充滿了從未被聖靈充滿的人,後來聖靈也再次感動,充滿祂的門徒.聖靈在人身上大大地作工,可是在行傳中有幾處經文,提到聖靈充滿卻不是一些動作,而是指著一種情況；就是說這個信徒或門徒,他一般的情形,他日常的情形是滿有聖靈的.
\newline
\newline
第六章有兩次提到聖靈充滿,那兩次不是在他們將揀選司提反的時候；不是他被題名時,忽然被聖靈充滿；或者他站在公會時需要膽量,他突然被聖靈充滿,並不是這個意思.我們若看原文,就曉得這次聖靈充滿,是指此人的本質,經常的情形；今天充滿,明天也一樣充滿.乃是說:司提反是常有聖靈充滿的.路加寫這本書告訴我們,司提反滿有信心,滿得恩惠,滿有能力；神的恩典在他身上非常豐富,無論作何事,為耶穌作見證,都充滿了神的能力.路加一而再,再而三地告訴我們:司提反是一位被聖靈充滿的人.
\newline
\newline
司提反的生平非常短暫,現在我們要從三個不同角度來看:
\newline
\newline
(一)作執事的司提反(徒六章)到司提反是執事
\newline
\newline
我自己曾經過一番的掙扎；因為聖經並沒有提到執事兩字,或者我們就用事奉主的司提反吧.司提反事奉主的背景很特別；路加告訴我們,當時教會出了一個問題.(我想青年人對路加的為人,一定很欣賞的.)路加從來不隱瞞事實,是就說是,不好的事他也坦然地說了.那時教會裡有人發怨言；因為有在處理膳食的事上不公平,路加就將這事誠誠實實告訴我們.早期教會竟有如此不如意的事情發生,但是感謝主,卻因此而帶給我們安慰.我們滿以為初期教會是完全的教會,是毫無問題的教會；其實不然,當教會得勝,福音廣傳,人數增多的時候,也難免帶來許多新的問題.我順便講個小故事:在美國佛羅里達州有個小教會,來了位青年牧師,他對教會有很大的抱負,首先他注意主日學和青年工作.過了不久工作很蒙主賜福,人數加多；在每主日的上午,主日學大有容納不下之勢；有位年老的基督徒,見狀搖頭說:「巴不得教會像以住一樣.」他好像覺得人數增多麻煩也多,座位不夠啦,教師也不夠啦,真是問題多多!初期教會當時有些說希利尼話的猶太人發怨言,這些人是誰呢?我們先要了解一下他們的背景.以色列人有些地方和中國人很相像的,他們人口分散各處；正如華僑散居在南洋以及世界各角落一樣.舊約聖經告訴我們,亞述國把北部的以色列十個支派俘擄去,隔了一段時期,巴比倫又把南部兩支派擄去；所以很多人都在巴比倫生根傳代.但是也有一部份的人,歸回聖地的.其餘則散居於北非,土耳其,希臘,羅馬及地中海四圍.這些猶太人都操希臘話,所以他們回來的時候,不一定講希伯來文或亞蘭文,他們都講希臘語；這班人住在耶路撒冷,也在聖殿聚會；當福音傳開時,這些人也成為福音的對象.在第二章,五旬節彼得講道,有從各地來的猶太人,他們聽福音,聽到了自己的鄉談.扎心受感動而信主.教會很關心那些貧窮的人和孤兒寡婦；但是在公賬的事上會有所虧欠,所以問題就呈到使徒面前；使徒看重這個問題,立刻召集眾教會.各位同工!當教會發生問題,求主賜給我們像初期教會使徒的膽量,面對事實,務求解決.使徒把重要的原則佈告眾人:
\newline
\newline
關於自己工作的原則(六2)「十二使徒叫眾門徒來,對他們說,我們撇下神的道,去管理飯食,原是不合宜的.」他們認清楚最重要的責任是傳講耶穌的道.他們多方考慮,計劃「選出七佰有好名聲,被聖靈充滿,智慧充足的人管理這事.」(六3)可見使徒對於人選的條件非常審慎.今天的教會選舉長老,執事或同工,有沒有考慮到這些問題?教會選舉的條件應當平衡,不但要有好的靈性,還須兼有才幹和本身的工作能力.眾人都很贊成使徒所提的方案；於是選出七位.解經家告訴我們,這七位都是說希利尼話的人,這些名字不是希伯來文,乃是希臘文.可能教會選舉時,他們有寬大的心腸；因為十二門徒都是猶太人,在加利利或猶太地生長的；現在耶路撒冷有從各地來的人,他們就從說希利尼話的人中選了七位,帶到使徒面前；他們接納了,這表明他們在工作上同心興旺福音,分工合作；使徒的責任是傳神的道,多多祈禱；被選的七位負責處理教會總務方面的事.結果如何呢?
\newline
\newline
1.「神的道興旺起來,」(7節)這原是聖徒所冀望的.2.「門徒數目加增的甚多.」由此可見,教會和睦,循屬靈原則去解決問題,對福音廣傳大有幫助；相反地,如果驕傲,各執己見,不照屬靈原則去做,傳福音的事工必受攔阻.3.「也有許多祭司信從了這道.」(7節末句)路加提及此點,包含很多的衝突；因為這些祭司原是在聖殿事奉的.
\newline
\newline
司提反被揀選後,他很謙卑處理一切事務,並不以為管理飯食是大材小用.除了忠心負起自己的責任,還用很多時間事奉主；並竭力為耶穌的真道爭辯,為主的真理作見證.在福音的工作上,一個管理膳食的人,常會存自卑感；司提反卻不然,雖然他「滿得恩惠能力,在民間行了大奇事和神蹟.」(六8)但他不以此自傲,他把精力專注於這一方面的工作,他明白神的道；在當時教會中,他是最重要的發言人,他把彼得在五旬節所講的道,作進一步的解釋.司提反在當時是否一位執事,我們暫不分辯,但至少我們曉得門徒分派他管理膳食；他關心基要真理,有條不紊地一一解釋；甚至使會堂眾人敵擋不住.哦!巴不得今天在中國教會,有更多的執事,更多的信徒關心基要真理.
\newline
\newline
多年來中國教會忽略了神學教育,甚至有加以輕視的.感謝主!近十年來中國教會有個新的動向；在香港,台灣,新加坡,菲律賓的華人教會,有了新的負擔；許多青年人願意一生事奉主,很多帶職的青年關心基要真理.
\newline
\newline
(二)為主作見證的司提反　我們看司提反的見證,應特別注意他所講的.
\newline
\newline
使徒行傳第七章是最長的一章,共有60節,從中我們發現司提反清楚了解基要真理.第六章是行傳中最短的一章.第二章論到有一班猶太基督徒和其他一般猶太人一模一樣,上聖殿禱告,參加聖殿的聚會,在聖殿獻祭.到了第三章,彼得和約翰在禱告的時間,上聖殿去禱告.一般虔誠的猶太教徒,上聖殿去的時候,這些基督徒也去；唯一的分別,就是他們認耶穌為彌賽亞,且盡力傳講耶穌是彌賽亞.彼得要證明耶穌是彌賽亞,特別提及舊約,也特別強調耶穌的復活.他說:耶穌的復活證明祂是彌賽亞.這件事是那些撒都該人所難容忍的,他們不信耶穌復活的事；所以他們是第一起來反對基督徒的.到了第七章,我們發現那些基督徒又增加了敵人；過去法利賽人對這事漠不關心；加瑪列就是法利賽人,他是位偉大的老師,他說:「現在我勸你們不要管這些人,任憑他們罷；他們所謀的,所行的,若是出於人,必要敗壞；若是出於神,你們就不能敗壞他們；恐怕你們倒是攻擊神了.」(五38-39)意思是任憑他們自生自滅,這是已往一般法利賽人的政策.可是現在問題來了,因司提反所講的,證明不是這麼簡單,耶穌基督的福音,和猶太教恰巧相反.衝突是必然有的,雖然早期的基督徒循規蹈矩到聖殿去,但司提反這時也把衝突講出來.彼得講道重點,是引用舊約先知所預言的彌賽亞,而司提反卻像個歷史家,講道時卻把以色列整個歷史都搬出來.從他們的國父亞伯拉罕直到耶穌基督,二千年的歷史都搬出來解釋；其重點有二:
\newline
\newline
第一,說明敬拜神不能受時間和空間的限制,是普世的,是屬靈的；神呼召亞伯拉罕非在聖地,神向他顯現是在吾珥,後來到了哈蘭,神再向他顯現；到了聖地時,這地不屬於亞伯拉罕,但神仍然與他同在.雅各把以色列人帶到埃及,埃及也非聖地；但神仍然與以色列人同在.摩西領以色列人出埃及,經過曠野地方,那裡也非聖地；但神仍然與他們同在.感謝主!這個信息應帶給我們中國教會很大的安慰,我們不需去參拜聖地,我們在這地方可以敬拜主.耶穌在世上時,有一天祂走路,到了一個井邊,有一撒瑪利亞婦人和耶穌談話；談了很久,她發現耶穌是先知,她對耶穌說:「我們的祖宗在這山上禮拜,你們倒說,應當禮拜的地方在耶路撒冷.」(約四20)耶穌說:「神是個靈；所以拜祂的,必須用心靈和誠實拜祂.」(約四23)司提反明白這個真理,當時他清楚提出:「主說,天是我的座位,地是我的腳凳；你們要為我造何等的殿宇,哪裡是我安息的地方呢?這一切不都是我手所造的麼?」(徒七49-50)他引用(賽六六1-2)的話來證明敬拜神不一定要在巴勒斯坦,也不一定要在聖殿.
\newline
\newline
第二,司提反提及在以色列的歷史上,一向有人拒絕神所差派來的；當神興起約瑟時,他自己的弟兄拒絕他.當神用摩西的時候,他的同胞拒絕他,反對他.當神興起眾先知,以色列人一一拒絕.司提反說:「你們這硬著頸項,心與耳未受割禮的人,常對抗聖靈,你們的祖宗怎樣,你們也怎樣.......」這些話真是叫眾人抵擋不住.司提反的意思是要眾人明白,敬拜神不是單靠聖殿；神要的不是表面的,而是內心的；如果外表虔誠敬拜神,內心卻遠離神,這絕非神所喜悅的.
\newline
\newline
今天在教會裡,我們也應該注意,雖然我們已經受洗,已經加入教會,殷勤事奉主；為主作見證,在教會擔任很多職務.但是不能單靠這些,神所要的,不單是外表的形式；耶穌說:「神是個靈,所以拜祂的,必須用心靈和誠實拜祂.」先知彌迦告訴我們,神所要的,不是獻祭的事奉,」他說:「世人哪,耶和華已指示你何為善；祂向你所要的是甚麼呢?只要你行公義,好憐憫,存謙卑的心,與你的神同行.」(彌六8)請問,你心靈的深處夠不夠這個條件?夠不夠這個標準呢?求主憐憫!
\newline
\newline
(三)為主殉道的司提反
\newline
\newline
司提反幫助他的同胞了解,耶穌基督的福音,不能受時間和空間的限制；不單是外表的問題,也是內心的事實.這位為主殉道的司提反,是位被聖靈充滿的弟兄.聖經說:「眾人聽見這話,就極其惱怒,向司提反咬牙切齒,但司提反被聖靈充滿,定睛望天,看見神的榮耀,又看見耶穌站在神的右邊.」(徒七54-56)當時他面對眾人的「極其惱怒」「咬牙切齒」,但他仍安然站立.天上來的平安,是世界所不能奪去的.在第七章開頭司提反說:「當日我們的祖宗亞伯拉罕......榮耀的神向他顯現.」而今司提反親眼「看見神的榮耀,又看見耶穌站在神的右邊.」聖經曾告訴我們,耶穌升天,座在父神右邊；但這裡為何說耶穌站在神的右邊呢?至少有兩個解釋:第一,耶穌說:「你們若在眾人面前不承認我,我在父面前也不認你們.」也就是說,若在人面前承認耶穌,耶穌在父面前也認我們.是否那時耶穌正站在父的右邊,見證這位殉道者司提反呢?第二,耶穌站在父的右邊要接納祂忠心的僕人,因那時司提反在世的時間不多了.司提反說:「我看見天開了,人子站在神的右邊.」這裡值得我們注意,在全卷新約聖經,「人子」都是耶穌自稱的,惟此次司提反稱呼耶穌「人子」.
\newline
\newline
司提反的禱告包括兩件事:
\newline
\newline
第一,他的禱告是向著主耶穌呼籲說:「求主耶穌接收我的靈魂.」當耶穌被釘在十字架上,最後也向父說:「父啊,我將我的靈魂交在你的手中.」司提反和主耶穌說了同樣的話.
\newline
\newline
第二,司提反最後的一句話,他跪下向主禱告說:「主啊!不要將這罪歸於他們.」主耶穌在十字架上也曾向父神禱告說:「不要將這罪歸到他們身上.」一個被聖靈充滿的人,那樣安詳,交託,饒恕敵人.
\newline
\newline
有人說,使徒行傳六至八章,有很多資料是腓利告訴路加的,此說很可能；因為提到執事的事以及腓利的工作；可是我相信還有一位把自己的見證告訴路加；因為當司提反站在公會時,「在公會裡坐著的人,都定睛看他,見他的面貌,好像天使的面貌.」當司提反跪下禱告,安然把自己的靈魂交託主,以愛心為敵人禱告,這些情形由始至終,掃羅在現場,親眼看見司提反的面貌好像天使,親耳聽見他的見證,心裡了解他的意思,就是耶穌基督的福音必須接替猶太教,且要超過猶太教,不能受猶太教的限制；但當時的掃羅卻是站在反對的立場.
\newline
\newline
弟兄姊妹!這個見證向我們說明,有一位被聖靈充滿的人,他謙卑雖然在教會管理膳食,卻能忠心作教會所分派他作的工.因此福音被廣傳,他也能行奇事,卻不因此而驕傲,更不以此為重,他反而全力注重在耶穌基督的真道.他忠心,貞節,勇敢,毫不保留的,為一次託付聖徒的真道竭力爭辯,他因此而殉道.
\newline
\newline
弟兄姊妹!我覺得更重要的,不是你我是否願意為主殉道的問題,而是在日常生活中,你是否願意為主而活呢?你是否願意將你的身體當作活祭獻上,好像司提反,在日常生活中忠心為主而活.今天中國教會,主耶穌基督的教會,需要這樣的青年,需要這樣的基督徒.
\newline
\newline


\newpage

\section{第49屆~港九培靈會~戴紹曾~第四講　巴拿巴完全的奉獻}
\label{sec:789}
\textbf{第四講　巴拿巴完全的奉獻}
\newline
\newline
連結: \href{https://www.hkbibleconference.org/session-message/view/789}{\texttt{https://www.hkbibleconference.org/session-message/view/789}}
\newline
\newline
\hyperref[sec:788]{< < < PREV SERMON < < <}
~
\hyperref[sec:toc]{[返目錄]}
~
\hyperref[sec:790]{> > > NEXT SERMON > > >}
\newline
\newline
當聖靈充滿一個聖徒時,就會賜給他新的生命；要作潔淨的工作,就賜給他能力.
\newline
\newline
保羅開始傳福音時,對他有重大影響的兩個人；一個是司提反,一個是巴拿巴.司提反的生活與見證,和他所講神的道大大地感動了保羅.司提反被聖靈充滿；巴拿巴也被聖靈充滿.路加告訴我們「巴拿巴是個好人,被聖靈充滿,大有信心.」也許我們覺得這三句話太空洞,不太具體了；但當我們讀使徒行傳時,從路加敘述巴拿巴的見證中,就有清楚的交代了.
\newline
\newline
按聖經的原則,當兩件事擺在一起時；就會有強烈的對比,以引起我們特別的注意.路加寫這段歷史,有意把(四章32-37至五章1-11)作強烈的對比；他受聖靈感動寫下這史實,留作我們的儆誡.
\newline
\newline
耶路撒冷教會,是當時發展極速的教會,他們都同心合意興旺福音.從(四章32-35),我們看見教會的同心,他們一心一意事奉主；這是多麼美麗的見證!教會若能同心,福音就易於傳開；反之,不同心就難以引人歸主.他們的同心有相當具體的表現；因為沒有一人說,東西是自己的,都是大家公用,這是很崇高的理想.使徒在這階段採取了這種方式,是依照耶穌帶領十二門徒同樣的原則；在一個較小的基督徒團契中,這原則是行得通的；但是當團契的人數多起來的時候,就非常困難了,所以日後的教會就沒有採納它.巴拿巴在安提阿時,教會就再沒有採納這原則,保羅出外佈道,設立教會時,他也沒有採納這原則.當初的教會實行凡物公用,完全是出於甘心自願,並非出於規定,勉強,迫令；而是出於他們甘心樂意的表現.他們「分物」乃按各人的需要,不是平均分配,既不規定也不限制.我們不但看見他們同心的表現,也看見神的能力在使徒身上,傳福音滿有能力,眾人大大蒙神的恩典.同心,有能力,神的恩典就臨到教會.「沒有一個人缺乏」不單指物質方面的需要,使徒講道,眾人靈裡也得到豐富的領受,沒有一人缺少；神的愛激勵他們,同心合意興旺福音.
\newline
\newline
在第四章,他們正好有個禱告聚會,彼得和約翰剛出監牢,和門徒在一起禱告.說:「主啊!他們恐嚇我們,現在求主鑑察,一面叫你僕人大放膽量,講你的道:一面伸出你的手來,醫治疾病,並且使神蹟奇事,因著你聖僕耶穌的名行出來.」禱告完了,聚會的地方震動,他們就都被聖靈充滿,放膽講論神的道.他們沒有要求主保護,只要求主賜膽量；好放膽講明主的道；他們的禱告果然迅速應驗,有能力傳講主的道,為主作見證.
\newline
\newline
我們對教會有個錯誤的觀念,只顧屬靈的需要,卻忽略了其他方面,這不是初期教會的典範.
\newline
\newline
使徒認為不能撇下神的道,他們有責任按真理傳講神的道；但同時他們也感覺教會裡許多在物質方面缺乏的人,教會應該關心他們,需要積極來解決這方面的問題.請問:在中國教會,在你聚會之處,有沒有這樣的平衡,除了顧及靈性方面的供應,也要兼顧物質方面的需要呢；這並非社會福音,而的確是教會應該要注意到的；如果忽略了,就離開了使徒時代教會的榜樣.
\newline
\newline
路加特別說明愛心,和奉獻的表現是怎麼一回事,路加提到巴拿巴的見證,先將他的背景告訴我們:「巴拿巴是個利未人,」意思是說他在服侍神的事上原已有份.巴拿巴曾否在聖殿事奉主,聖經沒有明載.巴拿巴生長於居比路,即今之塞普魯斯島,這島相當熱鬧,在亞歷山大帝時,已有很多猶太人僑居該地,那裡有會堂；巴拿巴和保羅到居比路,他們先在會堂向猶太人傳福音.在居比路有個假先知名巴耶穌,他和保羅之間有很大的衝突.至於巴拿巴何時離開居比路到耶路撒冷,我們不得而知.
\newline
\newline
巴拿巴本名是約瑟,這名甚好,熟悉舊約聖經的人,很容易聯想到舊約的約瑟；那樣的聖潔,那樣的受苦,後來成為他父家的救命恩人,他實在是位很好的青年.當巴拿巴生下來時,可能他的父母想到舊約的約瑟；盼望孩子長大成人之後,在主的工作上亦有貢獻,所以以此命名.但使徒卻給他一個新名,我想並非他的原名不好；在新約,名字很重要,是具有意義的.在外國,人名大都沒有意思；但是中國人的名字多數有含意,有的含意很寶貴,有的含意很偉大；我有一朋友名超群,此名真是出類拔萃!亞拿尼亞的名意思是「神的厚恩」,撒非喇意思是「美麗.」使徒帶領巴拿巴,發現他滿有恩賜,所以給他新名叫巴拿巴；可見初期教會很重視恩賜.巴拿巴的名是勸勉的意思,他很會勸勉人,教導人；另方面也是安慰者的意思.許多解經家都不懂得如何把這兩方面的意思連在一起,可惜他們不識中文；否則二者連在一起,就易如反掌了.中文譯為「勸慰子」,「子」意思就是他的本質,真是妙極了!
\newline
\newline
聖經說巴拿巴有很多田地,都變賣了把價銀奉獻,擺在使徒腳前.巴拿巴的見證顯明了他的心:
\newline
\newline
(一)巴拿巴有愛心
\newline
\newline
他們奉獻完全出於自動,毫不勉強,可以說是主的愛激勵他,他關心別人的需要,關心主的教會；就甘心樂意的奉獻,把價銀送到教會；這是他愛心的表現,願意在別人的需要上有所供獻.弟兄姊妹!今天我們在教會的奉獻上是否也出於愛心呢?
\newline
\newline
(二)巴拿巴有信心
\newline
\newline
他變賣了田地,全部價銀都奉獻,現在是一無所有了.一般來說,奉獻的人可能有兩方面的顧慮,第一,全都奉獻了,將來靠甚麼來養生過活呢?第二,教會到底有沒有前途呢?事實上這個見證已擺在兩次的大逼迫中,撒都該人和猶太人已經起來逼迫基督徒,彼得和約翰已經被下在監獄,不久司提反也殉道,這一切情形巴拿巴都看見了；如果羅馬人再起來反對,那麼,教會前途堪虞.你我今日的奉獻,恐怕難以經得起如此重大的考驗!教會的歷史迄今將近兩千年了,全世界都有信徒；但當時的教會,不過是個極小的團體,巴拿巴的信心卻勝過了一切的顧慮.
\newline
\newline
戴德生未到中國之前,神大大考驗他的信心,要他學習如何憑信心過生活.有一天他在倫敦傳福音,一個窮人來見他,請他去診治妻子的病.戴德生原是學醫的,於是就去為她診病,診罷留下藥方,但心裡即刻起了感動,有聲音對他說:「把你袋中的錢拿出來給這人,因他貧窮無錢買藥.」不過戴德生的生活也很艱苦,銀行裡沒有存款,家中也沒有隔宿之糧,僅有的是袋中的一個銅板.如果給他,明天吃飯也成問題；於是心中產生極大的掙扎,就對那人說:「我們跪下一起禱告吧!」事後,戴德生自己說:「我覺得當時的禱告是沒有意思的,甚至天花板好像是鐵的打不穿的.我們的禱告,主不垂聽.」禱告之後,他內心還是不斷掙扎,直到後來他順服了主,把錢交給那人.說:「朋友!這是我唯一的錢,你拿去買藥吧!」當他離開那家後,按他的日記簿所記:「我的心和我的口袋一樣的輕.」意思是他的心輕鬆沒有負擔了,他的口袋沒有錢了.回家之後天色已晚,準備休息；因為整個下午在外傳福音,頗覺疲乏；忽然聽見有敲門聲,原來是郵差送信,有個素不相識的人,寄了幾塊錢給他.關於此事,戴德生在日記上記著:「把錢借給主用,利息是最多的,不到一天,就連本帶利還給我了.」
\newline
\newline
七年前,神在台灣興起中華福音神學院,這學院完全沒有教會和差會的支持.(上述者是一個英國人的見證,現在要提到一個中國同工的見證)有個同工,他原有份很好的工作,在飛利浦公司有高職高薪.當神學院需物色一位總務,人選方面考慮了很久,後來想到這位弟兄,就和他商談；他一口答應,便向飛利浦公司辭職.但老板極力挽留,若他的新職有較高的薪金,老板也願意照給,只希望他留任.但這位弟兄告訴老板說:「將就任之職薪金僅及現在的一半.」老板為之愕然.這位弟兄存著奉獻的心,憑信心走上主帶領的道路.感謝主!不但在過去的教會歷史,看見西方青年有奉獻的心；就算在今日的中國教會,也有許多同樣的事情發生.
\newline
\newline
(三)巴拿巴有專心
\newline
\newline
他決意專心事奉主,他不願在主的工作上有所分心,他破釜沈舟地撇下一切.
\newline
\newline
回想宋尚節從美國返中國的途中,船行在太平洋上,他把文憑撕碎丟棄海洋中.這事引起了很多錯解,有的人批評他藐視教育；其實不是,他在船上時須作最後的決定.當然在他心中,難免有很大的掙扎；神在他身上,已經有使命要他傳福音,可是他的父親是教會的老牧師,盼望兒子學成歸來可以為他養老.而且國內多間大學已聘定他當教授,因此,父親的盼望,大學的聘書,產生了他心裡很大的衝突；所以毅然決定把文憑撕掉,要一生專心事奉主.
\newline
\newline
(四)巴拿巴同心的表現
\newline
\newline
他無意設立自己的基金會,無意注重個人的工作,無意創造自己的天下；他是屬於主耶穌的身體,有份於這身體；耶穌受苦他也該受苦,肢體有需要他也該有份,他真是願意把自己的一分力量擺在教會.
\newline
\newline
弟兄姊妹!今日主的教會太過分散,個人主義,分散了教會的力量；宗派主義,也分散了教會的力量；許多工作有待教會來作,卻因彼此不願意同心,沒有進展.論到大眾傳播的工作,如果中國要好好地利用大眾傳播的媒介,不是任何教會團體能作到的；如果中國教會要施行神學教育,也非任何團體能夠做到的；我們必須集中神給我們的力量,教會的力量,同心合意興旺福音.現在中國教會面對差傳工作,這工作也會把我們的力量分散,求主幫助!使我們能同心在主的工作上,使我們同胞看見這樣的見證,也使我們能作到非同心不能作到的事.
\newline
\newline
(五)巴拿巴有謙虛的心,他是利未人,是受過教育的人,他願意把奉獻送到使徒腳前,這表示他的謙卑
\newline
\newline
弟兄姊妹!我們在教會,是否願意順服權柄?多少時候,我們看不起牧師,看不起神的僕人,以為自己比別人強,不願意像巴拿巴謙虛的把奉獻,把力量,擺在教會中.為甚麼「使徒大有能力見證主耶穌復活?」(四33)因為弟兄姊妹同心,接著說:「眾人也都蒙大恩.」為甚麼?這是巴拿巴和其他聖徒在神領導之下帶來的大恩典.
\newline
\newline
強烈的對比(五1-11)昨晚我曾提過路加敘述歷史非常真實,毫不隱瞞；這裡有件嚴肅,可怕,失敗的事情,路加清楚提出.
\newline
\newline
教會未曾有殉道者之先,神的憤怒已經臨到教會,神的審判臨到教會.「因為時候到了,審判是從神的家起首；若是先從我們起首,那不信從神福音的人,將有何等的結局呢?」(彼前四17)亞拿尼亞虛假之奉獻的事,和舊約約書亞記有件事很相似.舊約是講到神的選民,新約是講到神的教會；舊約以色列人進迦南地,攻打耶利哥城.神曾吩咐凡耶利哥城裡,任何東西都不可取；但是亞干私自留下當滅之物.路加在這裡所用的動詞,和舊約時所用的動詞一樣；而且這時所發生的事也很像舊約亞干的事一樣.亞拿尼亞和撒非喇賣了田地,即把價銀私自留下一部份.這裡我們要認識清楚,田地是他們自己的；賣與不賣,奉獻或不奉獻；奉獻全部或一部份,都有他們的主權.他們的罪有三:第一,撒謊 :當時,在猶太人中,商人作假是普通的事；可是在教會中,道德的標準,不是按照社會的標準.彼得對他們說:你欺哄聖靈.」親愛的弟兄姊妹!基督徒的言語非常重要,是就說是,不是就說不是.聖經說:「若多說就是出於那惡者.」第二,奉獻的動機不正 : 根本非出於愛心,他們以自我為中心,私自留下一部份價銀.可能他們看見巴拿巴全部奉獻,使許多人大受感動；因此在教會中聲望日高,眾人都尊敬他；所以亞拿尼亞夫婦目的是為騙取名聲和地位.
\newline
\newline
耶穌在登山寶訓中,特別在馬太六章提及這問題,當你禱告,奉獻,禁食時,你的動機何在?是否要叫人看,叫人稱讚呢?耶穌說:「他們已得了他們的賞賜.」耶穌幫助我們看見作事的動機是很重要的.保羅在林前十三章3節也提到此事:「若將所有的賙濟窮人,又捨己身叫人焚燒,卻沒有愛,仍然與我無益.」保羅也幫助我們看見作事動機的重要.第三,奉獻的心不完全 : 聲言全部奉獻給主,實際上卻沒有完全奉獻.各位!我們偶不小心,不但欺騙自己,也欺騙別人.有的青年在夏令會時奉獻自己,當時實在受了神的感動；但後來又想辦法妥協,以為帶職事奉也是一樣.有的青年心裡想,多得個學位,也許更被主重用；不知不覺,因他的妥協,就失去奉獻的心.並非說帶職事奉是不對的,但若神的旨意要你一生全心事奉,你帶職事奉就不對了.你若再得一個學位而不完全擺上,任你如何解釋都是欺騙神；神要的是你自己,是你的心.未知在座中有沒有全時間事奉主的弟兄姊妹,開始時真是像巴拿巴全心的事奉主；但不知不覺成了亞拿尼亞；可能有許多不同的因素,如今讓你變成分心了.
\newline
\newline
去年聽說有位宣教士,本來蒙召要全時間事奉主；但到了工場,他分心了,因該地有利於經商,他忙於投資,無暇事奉了.
\newline
\newline
各位同工!我們的奉獻如何?
\newline
\newline
亞拿尼亞完全沒有像巴拿巴的愛心,信心,專心,同心,虛心；彼得對他說:「撒旦充滿你的心.」巴拿巴卻是位被聖靈充滿的人.
\newline
\newline
弟兄姊妹!有一天約翰衛斯理走路時遇見一酒鬼,形狀可憐,他說:「若非神的恩典,那個就是衛斯理.」行傳第五章亞拿尼亞和撒非喇失敗的見證,我們只有低頭說:「若非神的恩典,那不是我嗎?」主對我們說:「凡稱呼祂名的人,總要離開不義.」
\newline
\newline
但願聖靈在你我心裡說話!
\newline
\newline


\newpage

\section{第49屆~港九培靈會~戴紹曾~第五講　巴拿巴衷心饒恕他的弟兄}
\label{sec:790}
\textbf{第五講　巴拿巴衷心饒恕他的弟兄}
\newline
\newline
連結: \href{https://www.hkbibleconference.org/session-message/view/790}{\texttt{https://www.hkbibleconference.org/session-message/view/790}}
\newline
\newline
\hyperref[sec:789]{< < < PREV SERMON < < <}
~
\hyperref[sec:toc]{[返目錄]}
~
\hyperref[sec:791]{> > > NEXT SERMON > > >}
\newline
\newline
主給我今年有這樣好的機會,來一同查考聖經；願主幫助我們叫我們知道,人活著不是單靠食物,乃是靠神口裡所出的一切話.願主使我們在查考聖經時,明白怎樣領受主所賜的亮光.
\newline
\newline
前數晚提到初期教會的形態,我們看見教會在使徒時代如何傳講主耶穌；我們也看見聖靈在教會作工,聖靈如何居住在神兒女裡面.我們特別提到司提反,和巴拿巴對保羅很大的影響.司提反是教會第一個殉道者.他是教會的執事,是個護教者,有條有理去對付敵人.昨晚我們看見巴拿巴的見證,聖經說他也是被聖靈充滿的,路加告訴我們,巴拿巴時時刻刻被聖靈充滿；我們特別看到他的愛心,信心,專心,他和教會同心,並他謙卑的心.這二人影響了保羅,這並非說保羅所領受的福音由他們而來,乃是他們生活的見證影響了保羅,保羅當時名叫掃羅；現在為了方便起見,我們一概稱他為保羅.
\newline
\newline
當保羅離開耶路撒冷往大馬色的路上,他手持宗教領袖發給他的證書,許可他捉拿基督徒；想不到這位逼迫基督徒的,在那裡遇見了耶穌,他一生掀起了一百八十度的轉變；神抓住了保羅,他認識耶穌了；之後他進到大馬色,大發熱心為耶穌作見證.許多猶太人都很希奇,因為這人原是逼迫基督徒的,為何現在為耶穌作起見證呢?經過這一段時間,保羅不能繼續留在那裡；猶太人起來了,大馬色當局要捉拿保羅；有位弟兄幫助他,用筐子把他由城牆上縋下去,讓他逃跑；所以事情不能一概而論.保羅在大馬色逃跑,神救了他的生命；在腓立比監獄他沒有逃跑,結果,主救了看守監牢之人的命.
\newline
\newline
保羅要到耶路撒冷見彼得；在加拉太書一章,說他願意與教會的柱石有交通,願意在福音的工作上和其他使徒取得連繫；所以他從大馬色回耶路撒冷時,沒有去找昔日老友,這是他的見證.一個基督徒改變了,屬於教會,有新的生命,教會是他新的家,有新的朋友；他回到耶路撒冷願意和素不來往的基督徒們有交通.路加說:「掃羅到了耶路撒冷,想與門徒結交；他們都怕他,不信他是門徒.」(徒九26)因為他們想不到保羅如此大改變.那保羅信主大概將近三年,在大馬色,在阿拉伯曠野,現在回耶路撒冷,首次和基督徒來往；但是他們卻怕他,不信他是門徒.保羅當時絕無可能重逢他的故知,同時,他也不作此打算,且老友們也不會予以接納；這時保羅的處境,真是進退維谷!際此之時,被聖靈充滿的巴拿巴出現了.
\newline
\newline
「掃羅到了耶路撒冷想與門徒結交,他們卻都怕他,不信他是門徒,惟有巴拿巴接待他,領去見使徒.」「接待」中文的意思是接待客人,但是這裡接待兩字有更深的意思,可能在中文聖經,這兩字有多次出現；但我很想知道原文除上此處用「接待」之外,還有在哪裡用過；於是我就查考原文,當我找到時,給我很大的鼓勵.馬太第十四章記載:耶穌行神蹟,以五餅二魚給五千人吃飽.事後,耶穌發現門徒和眾人有意擁祂為王；所以耶穌就催促門徒上船,渡過加利利海,解散了眾人,就獨自上山禱告.到了夜裡四更天,耶穌下山,在水面上走,徒門徒那裡去.彼得和其他門徒看見了,非常害怕,以為是鬼怪；耶穌說:「你們放心,是我,不要怕.」彼得說:「主啊,如果是你,請叫我從水面上走到你那邊.」主說:「你來吧!」彼得就下船,在水面上走,要到耶穌那裡去；只因見風甚大,就害怕!將要沉下去.便喊說:「主啊!救我.」耶穌趕緊伸手拉住他.「拉住」這動詞在原文和行傳第九章「接待」同出一字.當彼得沒有信心將要下沉時,耶穌「拉住」他的手,救他上來.當保羅到耶路撒冷,弟兄姊妹沒有信心更沒愛心,不肯接待他；但聖經說,有一位巴拿巴被聖靈充滿,伸出手來接待他.這裡有很重要的意思,就是巴拿巴肯饒恕人.　根據亞歷山大革里免告訴我們,(他是早期教會的領袖),當主耶穌打發七十個門徒兩個兩個出去時,其中有巴拿巴.(此非四福音的記載)如果這位早期的教父革免的見證是對的；當保羅逼迫基督徒時,巴拿巴也是被逼迫對象之一了.第八章3節:「掃羅卻殘害教會,進各人的家,拉著男女下在監裡.」第九章一節:「掃羅仍然向主的門徒,口吐威嚇兇殺的話.」過去保羅如此逼害教會和基督徒；巴拿巴接待保羅,必須先由心中饒恕他.司提反殉道時,跪地求主不要將罪歸於他們.巴拿巴有同樣的靈來饒恕保羅.
\newline
\newline
弟兄姊妹!我們肯衷心饒恕得罪我們的人嗎?他曾毀謗你,害你,在傳福音的事上和你衝突；實在太對你不起,你肯不肯饒恕他呢?
\newline
\newline
主耶穌非常注重饒恕的真理.在登山寶訓中,耶穌講了很多寶貴的真理,其中包括有主禱文.路加告訴我們；耶穌的門徒問祂說:「主啊!你能教導我們禱告,像約翰教導他的門徒禱告一樣嗎?」(太六9-13)耶穌以主禱文教導門徒禱告,其中說:「你們不饒恕人的過犯,你們的天父也必不饒恕你們的過犯.」到了十八章彼得問耶穌:「主啊!我的弟兄得罪我,我當饒恕他幾次呢?到七次可以麼?」耶穌說:「我對你說,不是到七次,乃是到七十個七次.」繼之,耶穌講了一個天國的比喻:有個主人,他的僕人欠他一千萬銀子,主人明知他無法償還,吩咐把他和他的妻兒以及一切所有的變賣還債.僕人跪地求饒,主人便動了慈心釋放他,並免了他的債.當僕人出來,遇見一個欠他十襾銀子的同伴,便以兇狠的態度對待他,把他下在監裡,迫他還債.眾同伴看見了,抱打不平,把這事告訴主人,於是主人叫了他來,怒責他說:「你這惡奴才,你央求我,我就免了你所欠的債；你豈不應當憐恤你的同伴,如同我憐侐你嗎?」主人大怒,把他交給掌刑的,等他償還了所欠的債.
\newline
\newline
親愛的弟兄姊妹!我們的主是無所不知的,當審判來到,何等可怕!也許你覺得這是神學上難以解釋的,在心理方面是不能解釋的；可是在福音的真理上,如果你不肯饒恕人,被捆綁的,不是那人,而是你自己.是你把自己捆綁了,而不是主人再次把你捆綁.今晚在我們中間,不知有多少弟兄姊妹自己受捆綁呢?這是因為你不肯從心裡饒恕別人.我們常背主禱文說:「免我們的債,如同我們免了別人的債.」如果我沒有赦免人,這個標準怎樣執行呢?如果我不肯赦免的的弟兄,不肯赦免你的姊妹；主怎能赦免我們呢?若照我們所定的標準,主就不肯赦免我們了.親愛的弟兄姊妹!因巴拿巴從心裡饒恕了保羅,不但釋放了保羅；巴拿巴自己也得釋放,福音也得到釋放了.「那時猶太,加利利,撒瑪利亞,各處的教會都得平安,被建立；凡事敬畏主,蒙聖靈的安慰,人數就增多了.」(徒九31)讓聖靈在我們心裡作工.司提反說:「主啊!不要將這罪歸於他們.」巴不得感動司提反的靈,今晚感動你的心；感動巴拿巴的靈,也感動我們.
\newline
\newline
巴拿巴使人和睦　巴拿巴接待了保羅,覺得僅他一人饒恕他是不夠的,還把他帶到使徒中間.加拉太書告訴我們,保羅到耶路撒冷,見了彼得和雅各,當時沒有看見其他的使徒；所以現在巴拿巴把他帶到使徒中間.　使人和睦的工作實在不易,常會吃力不討好,會引致雙方對自己的誤會.巴拿巴卻除了一切的顧慮,有信心,有膽量,有愛心；有智慧的把保羅帶到使徒中間,叫他們知道保羅是個屬靈的人.
\newline
\newline
在使徒行傳中,我們看見聖靈的工作,常是透過使徒,開他們的心竅,然後展開工作(彼得到哥尼流的家是個顯然的例子).
\newline
\newline
巴拿巴領保羅到使徒中間,說了三次「怎麼」:
\newline
\newline
1.他在大馬色的路上「怎麼看見主.」主的光照使他有悔改的心,巴拿巴說出保羅得救的經歷.
\newline
\newline
2.「主怎麼向他說話.」不但保羅看見異象,同時主還向他說話.
\newline
\newline
3.他在大馬色,「怎麼奉耶穌的名放膽傳道,」保羅為主作見證.巴拿巴滿有智慧,將保羅的見證一一向使徒敘述；於是耶路撒冷的門徒便接待保羅了.巴拿巴是個被聖靈充滿的人,他使人和睦.
\newline
\newline
今天中國的教會,香港的教會,和你聚會的教會,都需要使人和睦的人；到處都有很深的鴻溝,需要你作橋樑；需要你付上代價,為主背負十字架.在教會中,有時年長的看年輕的不順眼；但是我此次來香港,看見會中許多青年人都很同心.我承認,在我們的教會中,年長的一輩都堅持著傳統是絕對不可更改的；年青的一代,卻認為年長者因循保守,難以吸引人信主,應該採取新的方法去傳福音；如此教會就產生了代溝.誰願作橋樑?好叫教會弟兄姊妹能彼此相愛.年青的尊敬長者,年長的也重視青年.有時教會中因言語不同,產生了隔膜.各地教會,幾乎有此情形；有粵語,有潮語,有國語；在福音工作上造成阻力.有時中西不能同心,西國人看輕中國人；中國人卻看西國人不順眼,因而工作上發生衝突.弟兄姊妹!誰願作橋樑?使兩下和好,一齊在福音事工上,同心合意事奉主.有時受過教育,有錢財,有地位的基督徒,輕看沒有受教育,沒錢財地位的基督徒,求主憐憫!這些深溝在教會是不該存在的.有時受過神學教育,和沒有受過神學教育的同工,彼此不和,互相歧視,求主憐憫!有時在教會選舉方面,因為自己落選就與人衝突；請看!巴拿巴在選舉執事名單上未被提名,在十二使徒行列中也沒有他的名；他卻饒恕人,幫助別人,使人和睦.我們屬於一個身體的,保羅在林前十二章講得非常清楚,眼睛不能對手說,我用不著你；巴拿巴沒有對保羅說:「弟兄!我用不著你,你已往的失敗太大了,我們不能用你.」保羅說:「頭不能對腳說,我用不著你.」可是在今天教會裡,我們常常講這樣的話；這不是聖靈的工作,聖靈的工作乃使人和睦.
\newline
\newline
今晚我所講的,不是鼓勵弟兄姊妹回到教會去管閒事,而是我們經過禱告有了聖靈的感動,衷心去饒恕人；我們才有資格幫助別人,彼此和睦.我們常諸多藉口,說:「這不關我的事.」「我不能負這樣的責任」.「有很多的困難.......」主耶穌說:「要天天背起你的十字架.」各位同工!可能使人和睦將成為你所背的十字架,但耶穌說:「使人和睦的人有福了!因為他們必稱為神的兒子.」求主幫助我們,背起十字架來跟隨祂.主耶穌作了我們的榜樣,成全了人與神和好的福音；祂要我們講這福音.
\newline
\newline
巴拿巴衷心饒恕人,使保羅與教會和好.保羅在耶路撒冷只有十五天的時間,在這短短的日子,巴拿巴卻完成了如此重任,帶給了耶路撒冷教會一次大復興.弟兄姊妹!如果我們願意行在光明中,衷心饒恕別人,在教會中使人和睦；我們將看見神大大震動香港教會,復興教會,幫助我們向同胞傳福音；四百多萬的同胞需要耶穌,他們看見我們彼此和睦,這是最大的吸引力.因為最大的吸引力不是因我們的神學,我們的熱心,我們任何方面的工作,我們的組織,我們的教堂；乃是因看見我們彼此相愛,就被主的靈所吸引.我相信你的教會；如果弟兄姊妹彼此相愛,明年此時一定座位容納不下.為何福音沒有廣傳?是否因為我們在這兩方面虧欠了主?
\newline
\newline
巴拿巴是一個好人,被聖靈充滿,大有信心.路加在此提出一個相當具體的見證.聖經並無提到巴拿巴有過一次說方言；但聖經說他饒恕一個弟兄,他使別人和睦.
\newline
\newline
我們願意學習巴拿巴的榜樣嗎?
\newline
\newline


\newpage

\section{第49屆~港九培靈會~戴紹曾~第六講　巴拿巴的事奉}
\label{sec:791}
\textbf{第六講　巴拿巴的事奉}
\newline
\newline
連結: \href{https://www.hkbibleconference.org/session-message/view/791}{\texttt{https://www.hkbibleconference.org/session-message/view/791}}
\newline
\newline
\hyperref[sec:790]{< < < PREV SERMON < < <}
~
\hyperref[sec:toc]{[返目錄]}
~
\hyperref[sec:792]{> > > NEXT SERMON > > >}
\newline
\newline
願主的恩典在我們中間,據我所知道有的弟兄姊妹是從香港來的,也有的是從新界來的,感謝主!給我們有這樣的自由的機會,聚在一起查考聖經.
\newline
\newline
這次我所講的主題是初期教會的型態.我們看到在使徒時代的教會,主耶穌所佔有的地位；也看到聖靈在教會中的工作,特別看到了具體的表現,在司提反和巴拿巴的身上；這兩位主的僕人大大影響了保羅.在使徒行傳中,當我們讀到「被聖靈充滿」此句話時,我們的感想如何呢?可能我們立刻想到這是一種的表現；可是路加不斷讓我們看見,這些被聖靈充滿的人,在日常生活中,彰顯了主的榮耀.前晚我們看見巴拿巴奉獻的心,是一種完全的奉獻；但,亞拿尼亞和撒非喇,他們的奉獻,是保留的奉獻,是不誠實的奉獻,是博取人的榮耀的奉獻.巴拿巴的奉獻卻是愛的奉獻,是信心,專心,同心,謙卑的奉獻.昨晚我們看見巴拿巴衷心饒恕一個弟兄,這是一件不容易做到的事.主的愛感動他的心,他不但饒恕了弟兄,並且使別人和好；進而把保羅帶入教會.使保羅能與教會的弟兄姊妹,有美好的交通,享受教會裡面的愛,也幫助保羅在當時,當地,為主作見證.巴拿巴作為橋樑,他饒恕別人,使人和好,作主的橋樑.教會實在很需要,基督徒作為橋樑.
\newline
\newline
過去的兩晚,我們看到巴拿巴生活上的表現；今天晚上,我們要看他在事奉上的表現.他如何在事奉主的道路上,表現他是個被聖靈充滿的弟兄.首先我們要注意此事的背景,這件事和司提反的死有十分密切的關係.
\newline
\newline
(十一19-21)這段經文,我為它命名「無名英雄.」路加沒有寫明這些基督徒的名字,他們確是偉大的弟兄!關於這班無名英雄,有五件事值得我們特別注意:
\newline
\newline
(一)他們是信徒
\newline
\newline
當然我們也可以稱他們為門徒；因為在使徒行傳中,所有信主的,都稱為門徒.但是我發現其中並沒有一個人是使徒,也沒有一個人是耶路撒冷教會的領袖.他們都是平信徒.腓利到撒瑪利亞為主作見證,他是個執事；彼得到哥尼流的家為主作見證,他是神所重用的使徒.現在我們看見這些平信徒,照著剛才所提到的兩位工人的原則,把福音傳給外邦人,完全像(徒一8)所載,福音「在耶路撒冷,猶太全地,和撒瑪利亞,遍傳直到地極.」在我們的世代已到了地極的階段.我們在這重要的階段裡,所發現的竟然是一班信徒帶路,這是何等大的突破!他們向外邦希利尼人傳福音,他們並不因居於信徒身份輕視自己,以為無分於福音的工作；他們邊走邊傳,到腓尼基,到安提阿；無論往哪裡去,他們都為主作見證.現今在教會裡,很多基督徒常有一種錯誤的觀念；以為只有牧師,傳道人才是負起傳福音責任的人.其實主的福音不是單靠少數全時間事奉的傳道人,可以傳到地極的.有許多地方是傳道人無法到的,好像你辦公的地方,你的工廠,學校；許多可以到的地方,傳道人卻不能到；主耶穌安排你在那處做光,做鹽,你就要在那裡為主作見證.
\newline
\newline
(二)他們是受苦受逼迫的信徒
\newline
\newline
因著司提反殉道,他們被迫四散.表面看來,他們好像無家可歸的人,不幸和不如意的事,雖然不斷地臨到他們,他們並沒有懷疑到神的存在.假如稍有一些痛苦的事臨到我們,魔鬼又乘機來試探我們；我們對神的信心,就很容易起了懷疑了.請問你經得起這的考驗嗎?當時那些信徒,沒有懷疑神的存在,更沒有懷疑神的愛；不但沒有發怨言,卻到處為主傳福音.這是何等寶貴的見證!人要跟從主,就當捨己,背起他的十字架.可惜許多人在教會中,所要的只是精神的寄託；當他們得不著而困難又來到時,立刻就動搖了.這些無名英雄卻不然,相反地,逼迫成為一股推動的力量,使他們順服主的旨意,為主而活,為主作見證.
\newline
\newline
(三)他們不受傳統的限制
\newline
\newline
一般猶太人信了耶穌,為主作見證,只向自己的同胞傳福音；但他們卻不然.他們開始向別人傳福音,無論是自己的同胞,是猶太人；只要有一個人蒙恩,他們就高興.主的福音臨到外邦人,無論是哪一個人蒙福,他們也高興.當他們向外邦人傳福音時,他們不提彌賽亞；因恐聽者不易明白,他們傳揚耶穌是主,眾人一聽瞭然；神給他們智慧,讓他們知道如何向新的對象傳福音.
\newline
\newline
(四)他們不受環境的影響
\newline
\newline
他們到了安提阿,這是一個重要商埠,福音傳到這裡,是極重要的；安提阿是當時羅馬帝國的第三大城,(第一是羅馬,第二是北非的亞歷山大.)那城是中西文化交流地,也是重要的商埠.那裡有許多傳統的宗教,像以弗所居民一樣,拜亞底米女神；是個敗壞,淫亂,腐化的都市.羅馬有位詩人寫詩諷刺安提阿:「安提阿的垃圾,飄流在海上,一直飄到羅馬.」當那些基督徒進到安提阿,他們非但不受該地的影響,反而影響了那地,改造了腐敗的社會.感謝主!這是主給我們的挑戰!這班基督徒作光明之子,安提阿有各種宗教信仰,他們不因此低頭,他們絕不說:「條條大道通羅馬.」他們立刻擺出見證,他們的見證是──耶穌是主,」正像彼得所說:「天下人間,沒有賜下別名,我們可以靠著得救.」(徒四12)可惜今天許多基督徒,在這方面沒有明確的立場,以為所有的宗教,不過是大同小異,都是勸人為善；不敢放膽宣揚耶穌基督的名.當時,這班無名英雄,大發熱心為主作見證,在那污穢,淫亂的城市,傳講耶穌基督聖潔的福音.我們在香港也有同樣的使命；可能在多方面,今日的香港好像昔日的安提阿.我們實在需要效法這些無名英雄的榜樣.
\newline
\newline
(五)主與他們同在
\newline
\newline
感謝主!他們雖然受逼迫,甚至無家可歸,淪落到了一個風化敗壞的地方.但主與他們同在,正如耶穌所說:「天上地下的權柄都賜給我了；所以你們要去,使萬民作我的門徒......我就必與你們同在,直到世界的末了.」哦!親愛的弟兄姊妹!難怪這班無名英雄有能力為主作見證.這裡更清楚地說:「主的手在他身上.」主的手是甚麼?就是在福音的工作上有積極的表現,「許多人都歸服主了.」我最喜歡這句話,歸服主,就是人悔改,生命改變,有新的生命,生活絕對的改變了方向.已往是為己而活,如今卻為主而活.他們在安提阿地方,立下教會的根基,為耶穌作見證.求主幫助我們!讓我們看見安提阿成為第一個海外傳福音的中心,是因一班信徒的見證而設立.這實在帶給我們香港的弟兄姊妹很大的鼓勵.求主幫助我們了解自己的責任,從今以後,你不再以為傳福音是牧師的責任；在你的四周圍,你有責任傳講基督的福音.
\newline
\newline
巴拿巴的工作(十一22-24)我們應注意三件事:
\newline
\newline
(一)巴拿巴的順服
\newline
\newline
耶路撒冷的弟兄姊妹,得知福音已經廣傳的好消息,也知道在安提阿有很多人信了主；教會就打發巴拿巴北上.這件事顯出他們心裡的矛盾,本來他們聽到撒瑪利亞有復興的火,就立刻打發彼得和約翰去,為何現在沒有打發使徒去呢,聖經沒有明說,是否他們對於新的工作,有所不放心?可是至少他們選上巴拿巴確是明智之舉；因巴拿巴自小生長在居比路,該地和安提阿距離很近；且巴拿巴又是個有智慧,被聖靈充滿的弟兄.我曾經提過巴拿巴把變賣田地的錢,帶到使徒腳前；這是表現他的謙卑,就在此時,他被派北上；他不作考慮,毅然遵命.哦!多少時候我們在教會中,不願像巴拿巴一樣地順服；對於教會所安排的工作,無意去做,卻諸多藉口,更不願謙卑接受,或說:「等聖靈感動吧!」自己不願意去做,反而推諉責任.
\newline
\newline
(二)巴拿巴的看見
\newline
\newline
他看見神所賜的恩就歡喜,這是巴拿巴到安提阿觀察之後的感受；在福音的工作上這是個極重要的關鍵.我非常歡喜這句話,你呢?你有沒有看見神的恩；你若看見了神的恩,你將如何描述呢?
\newline
\newline
巴拿巴所看見的有三:
\newline
\newline
第一,他看見一些信徒為主作見證,「逼迫」不能攔阻他們的見證,新的環境,也不能攔阻他們為主作見證.
\newline
\newline
第二,他看見因著這些信徒作見證,有許多外邦人信了耶穌；希利尼人信了耶穌,這是新的突破.
\newline
\newline
第三,他看見猶太人和希利尼人在一起,同心興旺福音；這是空前的突破.巴拿巴看見了神的恩,就歡喜；雖是那些無名英雄作工的果效,但巴拿巴心胸寬大,看見主賜福別人的工作就歡喜.多少時候,我們的心胸窄小,只要主賜福自己和自己的家,自己的教會；不然,就紅了眼.巴拿巴是利未人,他看見一些信徒,一些沒有上過神學的人,蒙主使用,他一樣地歡喜.
\newline
\newline
(三)巴拿巴完全照著神的恩賜作勸勉的工
\newline
\newline
真是名符其實!他名巴拿巴意則勸慰子；他勉勵信徒要恆久靠主,立定志向；因他知道有許多波折,許多困難會臨到基督徒,他曉得猶太人和外邦人不易融合,不彼此相愛；所以鼓勵他們繼續以愛心相處,熱心為主作見證.
\newline
\newline
巴拿巴被聖靈充滿,正是他剛從耶路撒冷北上；如果是我,剛從神學院畢業被派到安提阿,恐怕我的作風和巴拿巴完全不同；可能我會對他們說:「已往的工作都是一班平信徒做的,現在我來了,我是受過神學教育的牧師,以後的工作從頭做起.」巴拿巴雖然知道原有的基礎膚淺不穩；但他在這個基礎上,已經看見了主的恩典,所以他願意在主的恩典上繼續發展,在原有的小小基礎上,建主安提阿教會.他並沒有把教會建立在自己的身上,這是巴拿巴的見證,因此,許多人歸服了主.教會的基礎,非在傳道人身上,非在執事身上,非在長老身上；不管你是被聖靈充滿,恩賜多大；但是教會的根基是耶穌基督；人,事,物都要過去,惟主的教會仍然存在.巴會巴以智慧將教會建立在磐石上,因此教會增長.
\newline
\newline
巴拿巴和同工　凡事要進行都不能一手包辦,不論你是牧師,傳道,長老,在教會事工上負重任的；如果你認為可一手包辦所擔任的工作,那就太可憐了!一個事奉主的人,很需要同工,巴拿巴是個聖靈充滿,是主重用的僕人,也不例外；無論哪一方面的事工,我們需要更多的同工.巴拿巴感覺他需要同工,他先有這樣的感動.「選擇」實在很重要,巴拿巴十分重視「選擇」.在(十一26)有兩個動詞「找」和「帶」,第一,巴拿巴需要找同工；第二,他需要帶同工.換言之,首先他需要發掘人才,進一步他需要培養人才.為了重視這個問題,他不是寫信,也不是差遣別人去做,而是親自到大數去.
\newline
\newline
今天中國教會,面對海外差傳工作,應當學習巴拿巴對於人選問題的重視.
\newline
\newline
巴拿巴往大數去找保羅,「找」這個動詞,表示是費了很多的工夫才尋到保羅.保羅在(腓三8)說:「我也將萬事當作有損的,因我以認識我主基督耶穌為至寶；我為祂已經丟棄萬事,看作糞土,為要得著基督.」解經家告訴我們,當時保羅回家,他的父母因為他變成了一個基督徒,非常不滿,公然與他脫離了父子關係.巴拿巴找到了保羅,就帶他到安提阿去.在第九章題過巴拿巴接待保羅,現在帶他和自己在一起.帶一個人實在很不容易,尤其像保羅,他曾受過高等教育,個性很強,又是一個有多方恩賜的人；巴拿巴能夠帶他真是不簡單!
\newline
\newline
一九七○年,台北舉行福音會議,為我所安排的題目是「中國教會與西差會的關係」.我認為這是一個難題,因為其中有很多的「文章」,很多的困難,和很多失敗的地方；為此我費了個多月的時間,去拜訪很多年老的傳道人和西國教士.使我畢生難忘的,就是有一天我到長老會拜訪周牧師,他在高雄任職多年.我問他:「你工作多年覺得西差會最大的失敗何在?請別客氣,開門見山的講吧!」他說:「西差會最大的失敗,就是在中國教會,沒有找出良好的人選,加以培養.」我繼續問他:「是否他們不善於物色人選和培養人才?」答:「最大的問題,就是他們害怕培養了恩賜多的人才,將來要代替他們的工作.」我無言對之,惟低頭求主憐憫.宣教士在中國作甚麼?是建立自己的地盤嗎?只顧自己的工作嗎?他曾否想到有一天他不在這裡?哦!主內的中國弟兄姊妹!這問題非但出現在外國教士身上,不但西教士犯有此種錯誤；中國牧師也未嘗不然,看見教會中有恩賜,有抱負,有學問,有心事奉主的青年就害怕,不敢提拔,藉詞「恐年輕得志自高自大,」其實怕的是代替自己的地位.教會中每個部門的負責人,幾乎都有此現象；為了鞏固本身的地位,不敢培養人才,不願發掘人才,不願鼓勵愛主的人事奉主.若從此以往,實是教會的大損失,教會的悲劇!但巴拿巴卻把最難帶領的一個青年發掘出來,帶到福音的工場上.加爾文說:「這是巴拿巴的單純,」我很歡喜這句話,所謂「單純」則不顧自己,完全為主,活著專心為耶穌.
\newline
\newline
有一天,施洗約翰的門徒給他一生最大的考驗.說:「拉比!從前同你在約但河外,你所見證的那位,現在施洗,眾人都往祂那裡去了.」約翰卻毫不保留地回答:「我曾說:我不是基督.是奉差遣在祂面前的,你們自己可以給我作見證.」求主給我們有這樣的看見:不是我,乃是主；是主的榮耀,是主的教會.施洗約翰說:「祂必興旺,我必衰微.」
\newline
\newline


\newpage

\section{第49屆~港九培靈會~戴紹曾~第七講　巴拿巴的見證}
\label{sec:792}
\textbf{第七講　巴拿巴的見證}
\newline
\newline
連結: \href{https://www.hkbibleconference.org/session-message/view/792}{\texttt{https://www.hkbibleconference.org/session-message/view/792}}
\newline
\newline
\hyperref[sec:791]{< < < PREV SERMON < < <}
~
\hyperref[sec:toc]{[返目錄]}
~
\hyperref[sec:793]{> > > NEXT SERMON > > >}
\newline
\newline
前三個晚上,我們特別提到一位被聖靈充滿的門徒巴拿巴,他有美麗而具體的見證.今晚我們要繼續來看他的見證.
\newline
\newline
「巴拿巴原是個好人,被聖靈充滿,大有信心.」(徒十一24)甚麼叫做好人?(絕對不是好好先生).甚麼叫做被聖靈充滿,大有信心?
\newline
\newline
首先,我們看見巴拿巴奉獻的心,是一種完全的奉獻繼之我們看見巴拿巴從心裡饒恕一個弟兄,這弟兄原是逼迫基督徒的;巴拿巴竟然饒恕他,且把這位弟兄帶進教會團契中.巴拿巴作為橋樑,使人和好,讓教會能夠同心合意興旺福音.昨晚我們看見巴拿巴順服的心.有一些受逼迫的基督徒,到處為主作見證,到了安提阿;他們以為福音對猶太人那么好,對外邦人也一樣好;所以就向外邦人傳福音.當耶路撒冷教會聽到這消息,就打發巴拿巴北上,巴拿巴遵命而去;這是今天基督徒應當學習的功課,教會安排的工作,我們應當謙卑接受.巴拿巴到了安提阿,我們看見他單純的心,他看見神賜福別人所作的工就歡喜,並且親自投入,來堅固這個工作當工作擴展,巴拿巴需要同工,我們再次看見他單純的心;他要找個同工,他所帶的同工是一位受過高等教育,有恩賜,個性很強的人,真是不容易.巴拿巴須付上代價!他不顧自己,只為了主名的緣故,巴拿巴將保羅帶到安提阿,他們足有一年的工夫,在教會一同聚集,教訓了許多人.本來,巴拿巴是勸勉安慰人的,我們又看見巴拿巴很關心教會的教育,有系統地將神的話教導神的兒女;他們需要明白真理,明白舊約和新約的基要真理.巴拿巴有條有理教導信徒們,讓他們知道怎樣把真理,在日常生活中活出來,使他們知道安提阿是個淫亂的社會,在此黑暗的地方,怎樣過聖潔的生活.商人作假不誠實,巴拿巴幫助弟兄姊妹作光明之子,在虛假的社會作真人;巴拿巴也幫助基督徒,知道在被逼迫之中怎樣在真道上站穩.
\newline
\newline
在已往的廿卅年,許多弟兄姊妹經過了相當大的考驗;這個考驗,也可說是福音工作上百年來的考驗.如果弟兄姊妹站不起來,是否証明我們已往在福音的工作上有所虧欠,沒有幫助基督徒知道,在逼迫之下應如何在真道上站穩,巴拿巴不但把基要真理教導信徒,幫助他們如何活出真理來;更重要的,他成全信徒各盡其責,讓弟兄姊妹都起來在福音的工作上有份,不單作個聽眾,更要作基督的精兵;不但聽道,還要傳道,若自己不明白真理,怎能事奉主!
\newline
\newline
在新約,特別是在四福音,主耶穌強調訓練信徒的重要耶穌說:「天上地下所有的權柄,都賜給我了,所以你們要去,使萬民作我的門徒」(太廿18 -19)耶穌給門徒的大使命,是使萬民作祂的門徒;接著有兩個步驟,耶穌說:「奉父子聖靈的名,給他們施洗」意思是悔改得新生命,第一步要有新的生命,才有資格作主的門徒,第二,耶穌說:「凡我所吩咐你們的,都教訓他們遵守.」請弟兄姊妹注意:使萬民作耶穌的門徒,包括兩方面,第一,傳福音,領人歸主,有新的生命;這是洗禮所帶來的意思；第二,教導他們,如保羅在以弗所書說:「要成全聖徒」在安提阿教會,他們被稱為基督徒是有原因的,因為他們活出基督的樣式.有一天,彼得和約翰站在公會中,反對他們的人,認出他們曾經跟隨過耶穌的;因為在他們的身上體驗到主的聖潔,主的能力,主的膽量.當時安提阿的人,看見信徒,就稱他們為基督徒;因為他們確實活出耶穌基督的樣式,並且他們在福音的工作上,日夜傳講耶穌基督,勸人信主,帶領人歸主.
\newline
\newline
我們在教會,很會「生孩子」;我不是說家庭計劃不好,而是教會的工作不夠,我們常開佈道會,吸引人來聽道,帶領人來信耶穌;信了之後,我們沒有計劃怎樣讓他們在真道上扎根.讓我冒昧地說,中國教會的主日學,(指台灣方面)只有兒童主日學,少有成年主日學,沒有好的計劃讓整個教會,對神的話有全盤的了解.求主幫助我們!
\newline
\newline
巴拿巴和保羅竭盡全力,幫助弟兄姊妹對神的真理有完全的了解,這件事對以後的工作關係非常密切.巴拿巴和保羅在安提阿只是暫時的一段期間；如果他們沒有把工作做好,就無法離開,但是由於他們對同工和弟兄姊妹訓練有方,所以當他們離開之後,教會的工作毫不受到影響；仍然有人負責傳道,有人負責教導弟兄姊妹.這是很好的榜樣!你自己所擔任的工作,你是否考慮:怎樣培養人才；到了有一天,你不在那裡,仍然有人代替你的工作.
\newline
\newline
關心同工的需要　第十一章末段告訴我們,當時的教會在差傳工作開始之前,先關心人在生活方面的需要.多少時候,今天的教會走上極端,有的只注意社會問題；在講台上講的是社會問題,工作上也是解決社會問題；有的卻從不考慮到社會問題,不顧到人的需要.
\newline
\newline
初期教會是個平衡的教會；不但注意靈命的培養,同時也關心社會的需要.當時有位先知指明耶路撒冷將有大饑荒,教會很關心這事,立即展開籌募；計劃先讓弟兄姊妹克盡己力,在這工作上有份,這是美好的見證.他們收集了奉獻,由巴拿巴和保羅帶到耶路撒冷.安提阿教會的弟兄姊妹,從未到過耶路撒冷,與彼得及眾門徒也從未謀面；但他們被主的愛所激勵；慎重其事,有計劃,負責任來進行,讓弟兄姊妹各盡所能.今天中國教會,也有很多福音工作；如大眾傳播工作,神學工作,差傳工作等,都進行緩慢,正因為弟兄姊妹沒有盡其力量.求主幫助我們!巴拿巴和保羅親自把奉獻送到耶路撒冷,然後回到安提阿.他們忠於職守,有始有終,清清白白,行在光明之中；他們活出的見證,足以作我們的好榜樣.
\newline
\newline
有個歷史家告訴我們,十七世紀時,英國幾乎走上法國內戰的路線；後來神興起約翰衛斯理,他和同工們到處傳福音.因此教會得到復興,扭轉了當時社會的危機.那時英國社會非常黑暗,在上有權有勢的貴冑,漠不關心窮人的困苦,沒有用公義管理國家.當福音傳開,這福音並非政治方面革命的福音；而是道德方面革命的福音,藉著福音改變了人心.約翰衛斯理每天早晨五點鐘,站在礦工調班之處,向上班及下班的工人傳福音,勸他們信耶穌,他們的生命改變了.同時,他用書信及各種方法,使掌權執政者大大改變；因此,當時英國就避免了一場流血的內戰.
\newline
\newline
香港的教會,有沒有關懷社會的事?青少年問題,吸毒問題,墮胎問題,窮人問題等等.......今天當我從新界回來,看見那些高樓大廈,又看見山邊窮人破爛不堪的房舍,我就想到一個問題:今天在香港的教會,如果來了一個沒有受過教育而且很貧窮的人,是否受歡迎呢?主耶穌說:「你們是世界的光,你們是世界的鹽.」基督徒應當關心這些問題,社會才不致繼續腐化.
\newline
\newline
有異象的教會　第十三章我們看見一個差傳的教會,在真道上有根基,有同工的帶領,有對社會的關懷.「在安提阿的教會中,有幾位先知和教師,就是巴拿巴,和稱呼尼結的西面,古利奈人路求,與分封之王希律同養的馬念,並掃羅.」(1節)當你讀至這些名字,是否覺得淡而無味,要打瞌睡呢?可是當我看到這些名字,我幾乎要跳躍起來；這裡證明教會的合一.合一的見證是差傳工作必要的條件.一般而論,先知的個性和作風古怪,脾氣急躁.教師則不然,慢慢地,逐句有條有理地講解教導弟兄姊妹；所以先知和教師在一起,很容易衝突的.至於巴拿巴是利未人,保羅是法利賽人,他們能在一起,實在是美好的見證!還有稱呼尼結的西面,「尼結」按英文是黑人,是從北非來的黑人.路求,這裡說他是古利奈人；但這名卻是羅馬人的名字.還有與分封之王同養的馬念,意思是馬念和希律王在皇宮一齊長大的,是相當有地位的人；還有青年人馬可從耶路撒冷來的.這許多人,其中有年青的,有年長的；有猶太人,有非洲人,羅馬人；有地位的,無地位的；保羅是受過高等教育的,他們共聚一堂,同心合意興旺福音,這是得勝的見證,是教會增長的表現；我們不能說他們之間沒有界線,但這界線並沒有攔阻福音的工作；黑人仍然是黑人,有地位和無地位的背景仍舊沒有改變；但是,不同的背景,不同的種族,並沒有影響福音的工作.
\newline
\newline
「聖靈說,要為我派巴拿巴和掃羅,去作我召他們所作的工.」(徒十三2)他們被聖靈差遣,開始在居比路傳福音.居比路是巴拿巴的老家,傳福音的工作由此為起點.佈道團出發時,由巴拿巴率領；但是當他們離開居比路往他處去；路加說:「掃羅和他的同人......」這時候,顯然保羅成為領隊者.由此我們再次看見巴拿巴的單純.保羅原是巴拿巴一手提拔的,現在竟代替了他的地位；然而巴拿巴卻不以為意.還有馬可就在此時離開.有的解經家說他是不同意保羅領隊,有的說他是思家；有的說他害怕,聖經卻沒有明說.這是一件嚴重的問題,多少時候,我們只顧自己,沒想到這是主的工作.當保羅前往許多新的地方傳福音,他需要同工；如果此時巴拿巴有所不滿而離去,福音的工作豈不半途而廢嗎?保羅需要同工,雖然巴拿巴現在不居於領導地位,但他願意支持神所重用的僕人.求主給我們有此寬大的心胸.聖經沒記載巴拿巴有沒有說方言,有沒有趕鬼,有沒有叫死人復活；但聖經中的記載,叫我們看見這位被聖靈充滿的僕人,有單純的心,有愛主的心；當主提拔他的同工代替了領隊的地位時,他欣然支持,同心合力,到處宣傳福音,設立教會.聖經記載,巴拿巴真是不顧性命安危,一次被驅逐,又一次被聖靈充滿；有喜樂,勇往前進.到了路司得,那地的人要拜他們,要宰牛獻祭；巴拿巴禁止他們,說:「不可以,我們是人.」巴拿巴幫助他們認識獨一無二的真神.在路司得,他們遭受極大的逼迫,那些反對的人,用石頭把保羅打得半死,拖到城門外去；但門徒為他禱告,他又起來了.這兩個同工離開路司得往哪裡去呢?弟兄姊妹!你認為查考聖經是件很困難的事嗎?有一次,我讀使徒行傳,邊讀邊看地圖；也許你覺得我不夠屬靈為何要參考地圖?但我從中受了很大的感動；他們在路司得遭遇最厲害的逼迫；其實越過一山便可到大數去；如果我是巴拿巴,我一定對保羅說:「回去罷,這裡離你老家很近,工作艱難,休息半年再說吧!」巴拿巴並無此心,他想到最初建立的教會；從路司得回到特庇,到以哥念,安提阿,一站又一站；他們堅固信徒,選舉長老,將整個工作交托主.
\newline
\newline
他們工作的過程:
\newline
\newline
1.在各地傳福音,他們在會堂,在家庭,在街道各處傳福音.
\newline
\newline
2.使人作門徒.
\newline
\newline
3.堅固信徒,造就信徒,教導信徒.
\newline
\newline
4.實行選舉.
\newline
\newline
5.禱告,把他們交託主.
\newline
\newline
多少時候,我們沒有按步就班,第一步,就跨向選舉.如果今晚是向西國同工講道,講至此,最少我還要講半小時；因為這是我們最失敗的地方,不懂得如何撒手.巴拿巴和保羅,他們懂得交託,讓主帶領祂的教會,讓當地的人起來作主的工.這是初期教會的見證.
\newline
\newline
巴拿巴和保羅為真理站穩　他們回安提阿向教會負責報告他們在外一段時間的工作情形.這時候,安提阿教會發生一個真理問題,他們倆位為真道竭力爭辯「因信稱義」的問題.當時有一些信徒從耶路撒冷來,告訴弟兄姊妹「如果不守割禮,不遵守摩西的律法,不先作猶太人,就無法得救.」我們早在第七章已經看見司提反的見證,他說明福音不能受地方的限制,這福音是為萬民的.猶太教不過是個過渡時期,福音要從猶太教延伸出去,不能重返猶太教.保羅和巴拿巴為了因信稱義這重要的真理站起來.中世紀時,神也興起馬丁路得為著同樣的真理；人得救不是靠行為,乃是因信稱義.我們需要明白真理,為神一次托付聖徒的真道竭力爭辯；不是作些吹毛求疵的事,乃是為基要真理站穩.可惜人常為了一些皮毛的字句爭辯；而自以為是為真理爭辯,以致教會分裂.保羅在加拉太二章告訴我們,巴拿巴幾乎有一度搖動,因從耶路撒冷來的壓力太大.彼得也幾乎搖動；保羅責他的見證不正.到了耶路撒冷,幸而巴拿巴和保羅為真理並肩站穩,在同一陣線,才不致造成教會分裂.感謝主!
\newline
\newline
可能你認為巴拿巴和保羅後來的見證不好；因他們不歡而散,當他們同心準備出發時,巴拿巴對保羅說:「我的親戚馬可,我們帶他一同去；雖然過去他虎頭蛇尾,我們再次提拔他吧!」但是保羅不同意,因他認為馬可是個失敗而沒有見證的弟兄；此去乃為堅固教會,非為提拔弟兄,應把他留在安提阿,讓教會幫助他.不錯,他們確因此不歡而散.
\newline
\newline
保羅在羅馬的監牢裡寫歌羅西書和腓利門書時,他沒有提到巴拿巴,卻提到馬可,那時馬可在保羅身邊.保羅給提摩太最後的信說:「我離世的時候到了,那美好的仗我已經打過了;當跑的路我已經跑盡了;所信的道我已經守住了; ......你來的時候要把馬可帶來,因為他在傳道的事上於我有益處」(提後四7-11)巴拿巴成功了,他把失敗的弟兄帶領起來;當保羅在監獄時,馬可來服侍他.新約聖經,沒有保羅福音,沒有巴拿巴福音;可是有馬可福音,因為巴拿巴關心一個失敗的弟兄.
\newline
\newline
可能在教會中有位弟兄姊妹失敗了,需要你作他的巴拿巴,可能在座中也許有青年弟兄姊妹,好像馬可一樣.你有很好的開始,後來卻軟弱了,失敗了!但你已經回頭,你就需要起來,跟隨巴拿巴,專心事奉主.
\newline
\newline
巴拿巴是個好人,他被聖靈充滿,大有信心;他饒恕別人,使人和睦;肯謙卑順服,有單純的事奉,他不但饒恕保羅,他也饒恕馬可;因此福音的工作就踏進一步.
\newline
\newline
但願感動巴拿巴的靈,感動我們每一個人!
\newline
\newline


\newpage

\section{第49屆~港九培靈會~戴紹曾~第八講　尋求神的旨意}
\label{sec:793}
\textbf{第八講　尋求神的旨意}
\newline
\newline
連結: \href{https://www.hkbibleconference.org/session-message/view/793}{\texttt{https://www.hkbibleconference.org/session-message/view/793}}
\newline
\newline
\hyperref[sec:792]{< < < PREV SERMON < < <}
~
\hyperref[sec:toc]{[返目錄]}
~
\hyperref[sec:794]{> > > NEXT SERMON > > >}
\newline
\newline
初期教會,是神所帶領的教會.路加敘述初期教會歷史,著重教會和信徒蒙神的帶領.
\newline
\newline
順從神是應當的　彼得向反對的人說:「順從神,不順從人,是應當的.」(徒五29)我們服從神過於服從人是應當的.
\newline
\newline
願主的旨意成就　當時教會中有班弟兄姊妹很作難,他們一直勸保羅別上耶路撒冷去；保羅不聽,最後他們就說:「願主的旨意成就.」
\newline
\newline
今年暑假,可能成為你最重要的一個暑假；因為你必須選擇,特別是你剛從學校畢業,你更需要選擇.選擇也是基督徒必要的事.我很欣賞古時孟子對他自己的選擇所說的話,他先講到他所吃的,然後講到他的人生哲學.他說:「魚我所欲也,熊掌亦我所欲也；二者不可得兼,捨魚而取熊掌者也.」下面的話,真是給基督徒很大的鼓勵.他說:「生我所欲也,義亦我所欲也；二者不可得兼,捨生而取義者也.」吃的東西要選擇,人生也必須有選擇.
\newline
\newline
基督徒對於選擇,最重要的是神的旨意如何?主要我選擇甚麼?我想在座中的青年們,必沒一人像我那樣糊塗；當我大學畢業前夕,我必須選擇.我就向主禱告:「主啊!我有順服你的心；可是我不知你的旨意如何?你可否寫封信,擺在我的書檯上；我必遵照你的信,絕對地順服你.」翌晨,我的書檯並沒有收到主給我的信.
\newline
\newline
神帶領祂的兒女　你關心知道神的旨意；但我敢說,神更關心讓你知道祂的旨意.早期的教會,他們在生活,真理,事奉三方面,卻蒙神的啟示,蒙神的帶領.
\newline
\newline
神特別的啟示　就是神非凡的帶領.在我們生活中,若不小心,容易落在「主觀」裡面,那是很危險的.我如此說,並非輕視這方面的帶領；而是提醒弟兄姊妹,不必規定求主怎樣帶領；如果你沒有得到所求的帶領,絕對別想:「難道我不是基督徒嗎?我不夠屬靈嗎?」這是魔鬼的手段,常迫使我們存著主觀的思想,抓住我們,控制我們偏左偏右.
\newline
\newline
在使徒行傳中,至少有五點論及神特別的啟示:
\newline
\newline
(一)天使
\newline
\newline
行傳中數次提到天使.當撒都該人反對基督徒,禁止他們奉耶穌的名講道,把使徒下在監牢；天使在半夜開了監門,使徒出來,到聖殿繼續為主作見證.彼得也有同樣的經歷(十二章)希律王殺了雅各,正準備殺彼得,當晚,主差遣祂的使者,為彼得打開鎖鍊,越過重重的監門,帶領彼得離開.第八章提及腓利被一天使帶領到曠野,向埃提阿伯太監傳道.第廿七章保羅在船上遇風浪,當船將被破壞時,主的使者向他顯現,對他說:「不要怕!你必站在該撒面前,並且與你同船的人,我都賜給你了.」
\newline
\newline
(二)異象
\newline
\newline
保羅看見從天上來的異象,耶穌向他顯現,這異象是保羅畢生難忘的；因為這異象使他的人生有一百八十度的轉變；原是逼迫基督徒的掃羅,變成了為主受苦的保羅.在大馬色的門徒亞拿尼亞看見異象,在異象中主吩咐他去為掃羅禱告.哥尼流是個禱告,渴慕,虔誠的人,但他不認識耶穌；有一天,他在禱告中看見異象；神對他說:「你要打發人到約帕去,帶彼得回來；因為他要將我的話告訴你.」我們最熟悉的馬其頓異象,也稱之為馬其頓呼聲.許多人以為這是保羅工作的開始；其實,此時保羅已經工作多年了.當福音須由亞洲傳入歐洲時,保羅在半夜看見異象；有一個人呼召他向歐洲人傳福音.此外,還有一次彼得「魂遊象外」據路加的解釋,他把這事和異象相題並論.當彼得等候午膳,當哥尼流的使者將到他家的時候；聖經說:彼得魂遊象外,看見天開了,有一塊大布,裡面有地上各樣四足的走獸和昆蟲,並天上的飛鳥,又有聲音向他說:「彼得!起來,宰了吃.」彼得卻說:「主啊!這是不可的,凡俗物,和不潔淨的物,我從來沒有吃過.」於是有聲音向他說:「神所潔淨的,你不可當作俗物.」主叫他到外邦人,哥尼流的家去傳福音.
\newline
\newline
(三)預言
\newline
\newline
在使徒行傳中,有先知多次說預言.先知亞迦布到安提阿,聖靈感動他說預言,天下將有大饑荒,特別耶路撒冷一帶更甚.因此,弟兄姊妹把奉獻,從安提阿送去耶路撒冷.亞迦布第二次說預言,是當保羅要上耶路撒冷時；他到了該撒利亞,面向耶路撒冷,亞迦布阻止他,說:「因為反對的人很多.」但保羅堅決要去.後來亞迦布,好像舊約的先知,把保羅的腰帶拿來捆綁自己的手腳,說:「猶太人在耶路撒冷,要如此捆綁這腰帶的主人.」亞迦布雖被聖靈感動說預言,且事情也應驗了；但保羅卻不聽從,因為他另有一種帶領.我們應當明白,亞迦布被聖靈感動說預言,沒有限制保羅的行動；保羅仍在神帶領之下行動,自己向神負責.保羅到了耶路撒冷,而且更免費到了羅馬.(他乘船由羅馬政府負責)
\newline
\newline
(四)聖靈的感動
\newline
\newline
腓利被聖靈感動,聖靈對他說:「你去貼近埃提阿伯的太監.」聖靈帶領腓利.有一次聖靈對彼得說:「有三人等候你.」這是聖靈的感動.當時安提阿教會也蒙聖靈的帶領,聖靈對安提阿教會說:「要為我分派巴拿巴和掃羅；去作我召他們所作的工.」他們被聖靈差遣出去作工.
\newline
\newline
(五)聖靈的攔阻
\newline
\newline
聖靈不但感動,差遣,帶領作祂的工,聖靈有時也「禁止」「攔阻」.當保羅想向西行,耶穌的靈不許；要往北行,聖靈又不准.在十六章我們又看見聖靈對保羅的攔阻.
\newline
\newline
弟兄姊妹!我們應當學習順服聖靈的攔阻；我們也別忘記約翰說:「親愛的弟兄啊!一切的靈你們不可都信,總要試驗那些靈是出於神的不是；因為世上有許多假先知已經出來了.」(約壹四1)我們應當小心分辨,加以考驗；因為有的靈不是從主來的.　舊約聖經告訴我們怎樣考驗:
\newline
\newline
1.若有先知說預言沒有應驗,可知那是假先知,非聖靈的感動.
\newline
\newline
2.如果先知所說的預言應驗了,但藉此引領神的百姓遠離主,這必不是真正的先知；因為神絕不叫先知說預言,令祂的百姓離開主.
\newline
\newline
記得我讀大學時,校長對我說:當他年青時正談戀愛,(那時未當校長)他和一個美貌的女子彼此相愛,尤其女的更加愛他；可是後來他深覺這事不是神的旨意,就對女方表示從此要分手,女子卻說,她有主的感動,需要持續友情；他就對女子說:「你有聖靈的感動,但是我為甚麼沒有呢?」許多時候,基督徒以自己認可的事,強調是聖靈的感動；若有人反對,豈不就等於反對聖靈否?這是很危險的!在主的工作上,有意見提出是很好的,可能別人也有和你同樣的意見,大家應當同心,在神的話語帶領之下以作決定；如果輕易引用「聖靈的感動」,就是得罪聖靈了.
\newline
\newline
聖靈的感動是有原則的,保羅告訴我們:第一,榮耀耶穌,高舉基督.第二,符合聖經,聖靈的感動,絕對不與聖經的記載有所衝突；因為聖經是出於聖靈啟示,作者把神的話記下來留給我們的,今日聖靈的感動,和昔日聖靈的感動作者,絕不矛盾.
\newline
\newline
我再說,我絕對不輕視上述五方面的帶領；因為聖經明文記載,但是,親愛的弟兄姊妹!特別是青年們!不一定有這種經歷才算是神的帶領；神要用更清楚,更易於令我們了解的方法來帶領我們.　以下我要提到較為平常的帶領；較為普通的引導:
\newline
\newline
聖經是我們生活和信仰最高的準則,無論舊約,新約都是這樣.改教運動馬丁路得帶領教會回到這個真理.
\newline
\newline
神的話是我們信仰和生活最高的準則.我們都有聖經,到底我們怎樣用聖經?耶穌受試探時,用聖經的話抵擋魔鬼；但魔鬼也引用聖經的話再次引誘耶穌.所以我們要小心,用聖經要有正當的方法,也會有錯誤的方法.有的基督徒用「魔術式」讀經,把聖經打開,手指一點,就當作是神的啟示.聽說有次,有個人採取此法,頭一次打開聖經,看到「猶大到外面吊死了」,接著又打開聖經另處:「你要如此行.」我想這非神的帶領,也非使用聖經的正當方法.你每早靈修,可能神藉著一節聖經抓住你的心；但是奉勸弟兄姊妹,千萬勿採取亂點方法.查考聖經,主要藉著聖經的話指引你的路.
\newline
\newline
有的人用「寓意式」的方法來解釋聖經,把耶穌所講的比喻「好撒瑪利亞人」其中的地名,人物和事情的經過,都一一作出代表；這樣解釋聖經實太過份!耶穌所講的話,難道我們不看上下文嗎?我們常是斷章取義,抽出幾個字或是一節經文來作解釋.
\newline
\newline
正當讀經法是「歸納式」的方法,要以全卷來看,要注意每節上下文的銜接.題目是甚麼,這樣才可領受到聖經的光,不至於斷章取義.也要注意歷史的背景,和講者的意思；不一定要懂原文,當然懂得原文要感謝主；但是不懂原文並非說你不懂聖經,中文聖經同樣能帶給我們主的亮光.大衛說:「你的話是我腳前的燈,是我路上的光.」弟兄姊妹!巴不得神的話是你所愛慕的,常把神的話藏在心裡.
\newline
\newline
我們從使徒行傳中,看見早期教會很重視主的話.在第一章,當他們要補選使徒時,(那時猶大已經死了)彼得站起來說:「必須從那常與我們作伴的人中,立一位與我們同作耶穌復活的見證.」這是根據聖經預言而行的.第二章,彼得要幫助聽眾知道耶穌是彌賽亞,他並非用崇高的理論；因為聽眾是猶太人,所以他由始至終都引用舊約聖經,來證明耶穌是基督,是彌賽亞.到了十五章,教會開重要會議討論真理問題,討論非猶太人加入教會的規矩問題；彼得,雅各,保羅,巴拿巴都一一發言.彼得講他傳福音的見證,保羅和巴拿巴也起來見證；稍後,雅各也站起來.(他是耶穌的兄弟)雅各在當時是耶路撒冷教會的領導人；但他並不以身份自居,來解決教會的問題.他只是說:「我們已經聽了這許多工作和生活的見證,我們當以這些見證,結合聖經所說的話來作決定.」至於重要的選擇與決定,聖經能給我們明確的答案.
\newline
\newline
從行傳中,我們看見使徒們到處傳道,他們教導工作的內容,不外乎新約聖經中的真理；因為他們所教導的許多功課,成為後來的福音和書信；雖然當時沒有記載下來,但根本上那就是新約的啟示,當時他們以舊約教導信徒,於是才有後來新約的啟示.
\newline
\newline
求主幫助中國教會,把根基扎穩在神的話語裡面.基督徒當在聖經裡扎根,如大衛所說:「他要像一棵樹栽在溪水旁,按時候結困子,葉子也不枯乾；」晝夜思想主話的是有福的人,第一,要神的帶領,回到聖經裡面；第二,要禱告.我們看見渴慕主話的哥尼流,他禱告,尋求,就得著了.聖經說:「祈求就給你們,尋找就尋見,扣門就給你們開門.」哥尼流是個見證,他正禱告時,主啟示他,藉著禱告,他得救了.
\newline
\newline
多年前,我的家庭面臨一個很大的決定,那時我父母在河南傳福音.我和我的姊,弟,妹四人在山東的煙台就讀；因為戰爭爆發,日軍佔領山東,又進入河南；我父母無法留在那裡傳福音,所以到山東和我們在一起.父親告訴我:「我們要回美國去,船票已購妥；因為日本人不讓我們在河南傳福音；且你們在山東求學也有危險.」我知道父母一直為這事禱告,我卻沒有禱告,我聽說要回美國,非常高興,因我認為美國是天堂.(我想許多青年人都和我有同感.)過了一段很長的時間,父親對我說:「我們要退票,因我們不斷為此事禱告,深覺此時離開中國是不合適的.你們仍留在這裡求學,至於我們,既不能在河南工作；我們就要到西北去,在陝西傳福音.」當時面臨極大的選擇,禱告實在非常重要.
\newline
\newline
父親曾經給我講個小故事:有一天他造訪一位傳道人的家,孩子正吵鬧要求一件事；他母親不答應,孩子繼續吵鬧.後來他的母親叫小孩把事告訴戴牧師；聽取戴牧師的意見.我父親先問小孩:「你有沒有禱告?」孩子說:「我沒有禱告,我不禱告；因我知道我所要的主不給我.」我想很多基督徒也是如此,明明知道不是神的旨意,所以不肯禱告尋求神的旨意.
\newline
\newline
主耶穌面向各各他,在客西馬尼園向父神禱告說:「不要照我的意思,只要照你的意思.」讀經,禱告,主要光照你的心.你現在是否面對極大的選擇,重要的決定?我勸你回家之後,跪在主面前禱告.
\newline
\newline
主藉著教會顯明祂的旨意　我們常忘記教會的權柄,忘記教會是主所用的；主要透過祂的教會來顯明祂的旨意.巴拿巴本來在耶路撒冷作工,可是教會經過禱告,分派他到安提阿去,他終於順服去了.彼得和其他使徒決定要彌補猶大的門徒空缺,條件和教會選執事時一樣,要有好名聲,被聖靈充滿,有智慧；這些是使徒所定的標準和資格；神透過祂的教會讓他們知道祂的旨意.耶路撒冷會議,初時教會未清楚神的旨意；他們聚集,禱告,討論,作見證,查考聖經；直到決定了,寫信通知眾教會,說:「我們以為好,聖靈也以為好.」意即不單是他們的意思,也是聖靈的感動和引導.
\newline
\newline
各位青年!各位同工!我們作事,應根據聖經的原則,也應順服教會的權柄.
\newline
\newline


\newpage

\section{第49屆~港九培靈會~戴紹曾~第九講　初期教會增長的情形}
\label{sec:794}
\textbf{第九講　初期教會增長的情形}
\newline
\newline
連結: \href{https://www.hkbibleconference.org/session-message/view/794}{\texttt{https://www.hkbibleconference.org/session-message/view/794}}
\newline
\newline
\hyperref[sec:793]{< < < PREV SERMON < < <}
~
\hyperref[sec:toc]{[返目錄]}
~
\hyperref[sec:795]{> > > NEXT SERMON > > >}
\newline
\newline
基督教傳入中國,至今已有百七十年了.當馬禮遜牧師來澳門,正是最艱難的時期.雖然福音傳入中國,歷時已久,但是信主的比例仍是相當少.今天晚上,我們並非要分析基督教在中國的歷史,也非檢討得失；要重要的,我們回溯使徒時代的教會,可以得到一些原則.那時使徒門徒傳福音,蒙神使用,在行傳中有很多例子.
\newline
\newline
路加是位醫生,是科學家,他很注意歷史；當他敘述這段歷史時,在路加一章一節和行傳一章一節都告訴我們.路加敘述這段歷史的原則,他自己考察並聽許多人的見證,收集資料才按次序記錄下來.路加很注意事實,也很注重統計；他敘述歷史有事實的根據,他把事實逐一分析:一章8節敘述福音在地理上的擴展；這是耶穌的預言,也是耶穌的吩咐,耶穌對使徒說:「你們要在耶路撒冷,」(一至七章是在耶路撒冷的時期)「猶大全地,和撒瑪利亞；」(由八至十二章記載這段歷史)「直到地極,作我的見證.」(由十三章安提阿教會的見證,直到廿八章保羅在羅馬監牢開始記載這段歷史)(我說「開始」是因尚在廣傳中)
\newline
\newline
在使徒行傳中,路加曾五次作出統計,(六7,九31,十二24,十六5,十九20)當福音歷史每告一段落時,路加就作出一個結論；讓我們從中看見福音的能力和教會的擴展.
\newline
\newline
從上面五處經文,可以看到教會增長的原則及其結果:
\newline
\newline
(一)「神的道興旺起來,在耶路撒冷門徒數目加增的甚多.」(六7)
\newline
\newline
這是結果.從一至五章記載福音廣傳有一定的原則,然後產生人數增多的結果.到底福音廣傳的原則是甚麼?第一件事完全照著主說的:「你們要作我的見證,現在作沒有能力,要等候聖靈的充滿.」所以第一可說是這段歷史的出發點.是她能力的源頭.四福音記載彼得失敗的見證,一個使女問他是否耶穌的門徒；他竟不敢誠實回答,他三次不認耶穌；不但彼得如此,其他門徒也是遠遠地跟隨耶穌；這些門徒怎能為主作見證呢?但是,到了第二章,他們得了新的能力,自然而然他們的見證感動人心,吸引人歸主；這是聖靈的能力.弟兄姊妹!今天中國教會,以及你所屬的教會要有力傳福音,實在太需要聖靈的能力,你需要聖靈的充滿!
\newline
\newline
教會需要解決內在和對付外在的問題,外在的問題是敵人的攔阻和逼迫.敵人把使徒下監,不准他們奉耶穌的名講道；但使徒雖然挨打,也不願屈服.說:「我們不能不說,我們順從神,不順從人是應當的.」親愛的弟兄姊妹!不知你們有沒有這樣的感覺?主給我們的見證,不是單講些理論.當時使徒們雖受極大壓力,他們仍然在聖殿裡,在家中不斷傳道,教導人,為主作見證.教會遭逼迫,務要站穩；不能閉口不言,見證我們所見,所經驗的；主在我們心裡,在我們生活中,何等寶貴!我們不能不說.
\newline
\newline
教會內在的問題　當時教會因為人數增多,處事未免漸趨不協調,不平均；有的人受虧損,沒有得到公平的分配,以致有人發怨言；外在問題既已對付,內在問題也急需解決；否則福音不能廣傳.使徒面對此些問題,認清自己的責任,遵照耶穌時常的教導；以祈禱傳道為念,照常講解真道,供應弟兄姊妹靈裡的需要；使徒們站穩立場,不讓問題影響工作；所以就選舉七位執事,分工處理內在問題；而且他們也開始傳福音.這是多美的見證!使徒們懂得屬靈的原則,守住自己的崗位；讓新選的執事分擔解決內在問題；因此,初期教會大衛增長,由三千人增至五千人,且不斷增長.
\newline
\newline
(二)「那時猶太,加利利,撒瑪利亞,各處的教會都得平安,被建立；凡事敬畏主,蒙聖靈的安慰,人數就增多了.」(九31)
\newline
\newline
這是一幅美麗的圖畫!屬靈方面顯然增長,基督徒長大成人,遂即人數增多.輔成教會增長的因素有二:
\newline
\newline
1.逼迫
\newline
\newline
逼迫使他們分散各處；前曾提及那些無名英雄受逼迫,到了撒瑪利亞,到了安提阿；他們不因逼迫而閉口,相反的,逼迫把他們帶到新的地方為主作見證.腓利到撒瑪利亞可能也是因受逼迫,不能留在耶路撒冷.撒瑪利亞人和猶太人,原非好友,他們結怨根深；當以色列人被擄之後,亞述王把外人帶入,撒瑪利亞人的背景始於那時；他們多方面遵守舊約.當耶穌到撒瑪利亞,和該地婦人談道,那婦人頗覺希奇,猶太人怎肯和撒瑪利亞人講話呢?平時,猶太人視撒瑪利亞人如犬類；撒瑪利亞人視猶太人也是如此.有一天耶穌帶門徒路經撒瑪利亞,被他們發現是上耶路撒冷的,因而遭拒絕；當時門徒之中雅各和約翰兩人,膽量大,脾氣躁；他們問耶穌:「你要我們禱告求父降火燒滅他們,像昔日先知以利亞所作的嗎?」耶穌責備他們說:「你們的心如何?你們並不知道.」跟隨耶穌的門徒,心裡的成見,無異一般的猶太人.到了第八章腓利下撒瑪利亞去傳福音,很多人信了耶穌；耶路撒冷教會得知這個消息,立即打發約翰和彼得前往.原來失敗的門徒約翰和他的兄弟雅各,脾氣暴躁,被人稱為雷子的；但是這次約翰來的時候,聖靈已經除去他心裡的成見；他們看見腓利的工作,看見腓利為撒瑪利亞人禱告；聖靈降臨在他們身上.這禱告和前雅各和約翰,所要禱告的完全不同；這次的禱告是聖靈的火,聖靈的能力使這班人有新的生命,有潔淨的心來為主作見證.你看!成見除掉,福音廣傳了.親愛的弟兄姊妹!在教會,我們心中存有成見,為此,我們不肯對外傳福音.求主除去我們的成見!
\newline
\newline
2.保羅得救
\newline
\newline
這件事成為福音廣傳的第二因素.司提反講道時的智慧,感動保羅的靈；他的理論,引用舊約聖經的話語,使反對他的人抵擋不住!包括保羅在內,好像有根刺刺入保羅的心；他看見司提反的死,是那樣平靜,那樣饒恕人,那樣交託主；司提反的死感動了保羅.殉道者的血,成了教會的種子!有人以為基督徒殉道,會影響福音廣傳,事實上適得其反.中國歷史上,庚子年間基督徒殉道,藉此福音廣傳,眾多同胞歸向了主.求主幫助我們!為著主,我們應當站穩!司提反的死,改變了掃羅,興起了保羅；同時,神的工作在保羅身上,大馬色路上逼迫教會的保羅,成為耶穌基督的忠僕,為主作見證.經過一段時間,他回到耶路撒冷,與教會和好.福音廣傳需要福音的使者；神的原則,不是靠天使,不是用別的方法；乃是藉著人使福音廣傳.今天在中國教會,神正尋找合乎祂用的人.可惜,像保羅一樣的知識份子,很少肯謙卑接受耶穌,很少肯謙卑被主所用.如果沒有保羅,哪裡有使徒行傳?在教會歷史,我們看見神興起保羅,興起馬丁路得,在英國神興起衛斯理,在瑞士,法國,神興起加爾文；都是大有知識,受過高等教育的人；他們的文字工作,影響了在上有權柄的人；他們對人的愛心和帶領,也使在上的人得救.宋尚節是個例子,他受過高等教育,帶了博士學位從美國回中國傳福音；雖然時間短短,但他在中國及東南亞一帶忠心傳福音；神大大的使用他.今天在中國教會,受過高等教育而有抱負有恩賜的青年在哪裡?誰肯擺上呢?這段歷史,甚至整個教會歷史,若沒有保羅,我們怎麼看,怎麼解釋呢?當然,在保羅之上有聖靈的能力,神的旨意,神的主權,使保羅歸主.基督徒被逼迫分散帶來福音.保羅歸主成為教會得力的傳道人.
\newline
\newline
(三)「神的道日見興旺,越發廣傳.」(十二24)
\newline
\newline
這個時期可說是安提阿時期,教會得到向外邦人傳福音的異象；本來猶太人以為福音僅是為了猶太人；如今他們曉得這福音,也是為全世界的人.彼得講道,自己不完全明白,他在五旬節之日講道,結論是──「凡求告祂名的人,必然得救.」他是引用舊約的話；到了第十章,主感動彼得向哥尼流傳福音,起初彼得不肯；但哥尼流求告主名,彼得低頭謙卑順服,結果福音廣傳.
\newline
\newline
抗戰時期,有許多西教士還在中國工作；主開始把海外傳福音的負擔,放在一個中國弟兄心裡,那時我父母在西北陝西,西安,保定一帶地方傳福音.在西北聖經學院,有位中國同工馬可牧師,有一天,他應邀到另一間學校講道；當日破曉,他起床出外,安靜預備心；主開始感動他的心,對他說:「你看教會歷史,福音傳開時,從耶路撒冷到安提阿,到希臘,羅馬,歐洲,英國,美國,中國,一直往西,現在福音已傳到中國.可是從中國到耶路撒冷,還有一段距離；這最後一環的工作,是留給中國教會的,應派傳教士出去傳福音.」馬牧師衷心感謝主!接著主又對他說:「我要你去傳福音,」馬牧師回答主說:「我不能去!因我家有七個孩子,我還要教書,要訓練工人.」主說:「你要去!」當時他心裡起了很大的掙扎；感謝主!他終於順服了.遂即組織「邊傳福音團」由中國西部開始,要把福音傳回耶路撒冷.這一帶地方,都是信奉佛教,更多信奉回教；傳福音實在不易!這是卅年前中國教會小小差傳工作的開始,至今情形如何,我們不得而知；但今天主的靈在中國教會,把差傳的負擔和使命賜給中國教會.求主幫助我們放棄各樣的成見,勝過各樣困難,把福音傳出去.以往外國人在中國失敗的地方,作為我們前車之鑑,勿重蹈覆轍!求主給我們聖靈的能力和智慧,把福音傳開.
\newline
\newline
當日在安提阿,使徒向外邦人傳福音,諸多逼迫,他們都一一勝過.希律王派人殺害雅各,為了討猶太人歡喜；希律王又捉拿彼得下監；我們看見神的權柄,差遣天使救出彼得.到了這件事的結束,路加有趣地記載驕傲的希律王,有一天他演講,聽眾說是神講話,是神的聲音,非人的聲音.希律王因不將榮耀歸給神致死.接著,聖經說:「但神的道日見興旺,越發廣傳.」希律王的話,人說是神的話；但神把手收回,希律王立刻氣絕身亡,他的話立刻沒有了；但神的話興旺起來.這是一個強烈的對比,人的話有限,神的話興旺.
\newline
\newline
(四)「眾教會信心越發堅固,人數天天加增.」(十六5)
\newline
\newline
這可稱為加拉太階段,教會被堅固,信心被堅固,人數天天加增.為甚麼?因為教會解決了一個重大的問題──得救問題,誰可得救?(WHO?)怎樣得救?(HOW?)我們已經看見；神幫助祂的兒女,知道福音是為萬民的.當福音傳開,傳到外邦人中,在安提阿,在加拉太和許多地方.已往的基督徒都是猶太人；但現在不是,這些非猶太人到底怎樣成為基督徒,其中實在很複雜.如果在今日,或可說是護照問題.因為這班人要先入猶太籍,信耶穌須有猶太護照,才能做基督徒；要受摩西的律法,男丁要受割禮,要遵守很多規條,作個標準的猶太人；具備許多條件,才可以受洗,才可被聖靈充滿.可是保羅站起來說:「不是這樣,」彼得也站起來說:「錯了!猶太派的人都錯了!我到哥尼流的家講道之時,他們已經被聖靈充滿；保羅和巴拿巴去的時候,那些人還沒有受猶太的規矩；但他們已經領受聖靈得救了.」保羅說我們得救是因著信,不在乎這些規條.保羅此時為因信稱義的真理竭力爭辯.感謝主!耶路撒冷的教會聚集在一起,商討這個重要的問題,得到聖經的答案.因此,教會有合一的靈,有同一的立場,在真理上站穩；沒受猶太派及外面宗教所影響,站在純正的信仰立場；所以福音廣傳.「耶穌會」在中國傳福音時,和其他的「修道會」大衛發生衝突,以致攔阻工作.今天教會與教會,門徒與門徒之間,是否各自為政?若然,必影響福音進展.我並非說不能有不同的宗旨,神藉著各教會把福音傳開,但問題在於態度方面.求主賜給中國教會,有合一的靈,在福音真道上齊心努力,使我們的同胞認識主.事實上,在福音的教會中,相同之處多於不同之處；但是,多少時候,我們強調與眾不同,忘了彼此共同之點的基要真理；其實,彼此不同之處,僅屬於皮毛.
\newline
\newline
(五)「主的道大大興旺,而且得勝,」(十九20)
\newline
\newline
這階段是希臘時期.福音傳到希臘國的第一個地方是腓立比；保羅向呂底亞傳福音.呂底亞是個虔誠人,但她不認識耶穌.保羅對她說:「行為不能救你,你必須信耶穌!」主開導她的心,呂底亞和她全家都得救了.在腓立比有個女孩子被鬼附著,有一天,女孩子跟著保羅；保羅轉過身來,奉耶穌基督的名釋放她；為此保羅被下監牢,在監裡傳福音,看守的兵得救了.這是保羅在腓立比的工作.後來保羅不能久留腓立比.我發現保羅每到一個地方,都是聖靈的帶領；每離開一地,都是敵人趕走他們；這也是聖靈的帶領,在敵人逼迫,追趕,殺聲緊緊的情況下,保羅仍繼續傳福音的前程.保羅到了帖撒羅尼迦三週,此段記載我很喜歡,尤其是英文聖經的記載更為生動.中譯:「......那攪亂天下的,也到這裡來了.」(十七6)英譯:「These Men Who have turned the World upsidedown have come here also,」意即「這些把世界顛倒的人也到我們這裡來.」這話可以代表當時福音的工作,真是把世界顛倒,把許多人的生命都顛倒過來.保羅在此傳福音,許多人歸主；雖然受逼迫,但見證出來了!接著,保羅在庇哩亞,在雅典,在會堂,在街上講道.在雅典大廟之前講道；他看見祭壇上寫著「未識之神,」就以此為題,向他們傳福音,向他們介紹耶穌基督是創造天地之主.保羅說:「我不以福音為恥,這福音是神的大能,要拯救一切信祂的人.」繼之,保羅到哥林多,這是一個腐敗淫亂之地；一般的希臘人,常以「哥林多化」為詞,來諷刺人之敗壞淫亂的惡行.保羅在哥林多,定意要講耶穌基督的十字架,要把人帶到十字架前.讓十字架來對付罪惡,讓耶穌的寶血洗淨罪污；保羅在此腐敗淫亂之地高舉了十字架.在哥林多書上,清晰可見福音的大能.接著,保羅到以弗所,他三次出外佈道,這裡的工作,深蒙神賜福,他先十二個人禱告,他們被聖靈充滿.保羅在會堂講道,不受歡迎；後來到一間學校,據英文翻本的註釋,說保羅在那學校,從上午十一時至下午四時傳福音；在那熱帶地方,在那段酷熱的時間,福音竟有極大的吸引力,使聽者對付罪惡.本來是信邪術,拜假神的都受感動,紛紛將那些敗壞的邪書帶來焚燒.感謝主!保羅把福音傳開,信的人得了耶穌基督新的生命,靠著聖靈的能力,過著得勝,聖潔的新生活.福音完全改變那地.
\newline
\newline
今晚,我們看見福音廣傳的原則,聖靈的能力；被逼迫時站穩,他們解決了內在的問題,除去成見,為真理站穩.他們所傳的是對付罪惡的福音.
\newline
\newline
今天在香港,若要教會增長,福音廣傳；必須注意這些原則.求主光照我們的心,使我看見福音廣傳的原則.願聖靈在我們中間作工!
\newline
\newline


\newpage

\section{第49屆~港九培靈會~戴紹曾~第十講　為主受苦}
\label{sec:795}
\textbf{第十講　為主受苦}
\newline
\newline
連結: \href{https://www.hkbibleconference.org/session-message/view/795}{\texttt{https://www.hkbibleconference.org/session-message/view/795}}
\newline
\newline
\hyperref[sec:794]{< < < PREV SERMON < < <}
~
\hyperref[sec:toc]{[返目錄]}
~
\hyperref[sec:1385]{> > > NEXT SERMON > > >}
\newline
\newline
感謝主!這幾天晚上有機會和弟兄姊妹分享主的話.雖然我們工作地方不同,面貌不一樣,但我們好像一家人.最近在美國有一些新名稱,如ABC(American Born Chinese)那麼,我是CBA了.還有,僑居在美國多年華人的後代,他們是在美國出生的,被稱為「香蕉」,因為外黃裡白.我聽了哈哈大笑,他們說:「你別笑!你是煮熟的雞蛋,外白裡黃.」可是我兩個孩子都不以為然,他們在中國學校讀書已經五年了；有一天我送他們上學,他們不要我送到校門口；我覺得很奇怪,原來他們不願意同學們,看見他們的父親是外國人；因為他們和同學已經打成一片了.這幾天我感到和你們也打成一片了.
\newline
\newline
在未講道之前,我願意和青年人說幾句話:這樣的聚會是很難得的,不分宗派,聚在一起；好像主的靈特別在我們中間作工,許多人有特別的感動.但危險也就在此,如果我們只憑感覺,聚會進行期間,作個好的基督徒；有追求,有事奉,有讀經,很愛主；但聚會過了,感動就有所不同.我們是否就這樣冷淡下來?求主幫助我們!不僅作聽道的基督徒,也願意出去行道.求主幫助每一位,特別是數晚以來舉手的弟兄姊妹.主要幫助你遵行祂的旨意,你已經在眾人面前站起來了,求主使你站穩!
\newline
\newline
今天晚上,我們要一同思想的,仍然是初期教會的見證,在新約,這方面的信息很多.
\newline
\newline
耶穌帶領門徒,不斷對他們說:「你們在世上要我受苦.」耶穌未離世之前,門徒已開始為主受苦.耶穌離開他們之後,聖靈降臨在他們身上,他們到處為主作見證,他們真是明白主的心意!在使徒行傳中,彼得和保羅(尤其是保羅)在各處帶領教會,提醒信徒要預備心為主受苦.
\newline
\newline
當保羅的基督徒路程開始的時候,主藉祂的僕人亞拿尼亞告訴他:「主對亞尼拿亞說,你只管去,他是我所揀選的器皿,要在外邦人和君王並以色列人面前,宣揚我的名.我也要指示他,為我的名必須受許多的苦難.」(九5-16)亞拿尼亞本不願意去見保羅,因保羅原是個逼迫基督徒的人；但主叫他「只管去,」別追究已往的事；因為「他是我所揀選的器皿」這句話我很喜歡.我們都願意作主重用的器皿.「要外邦人和君王並以色列人面前,宣揚我的名.」如果我們有機會在眾人面前見證神的恩典,傳福音,高舉耶穌,這是多美的事!「我要指示他,為我的名必須受許多的苦難.」作主重用的器皿誰都願意,站在君王面前,許多基督徒也願意；至於為主受苦,基督徒卻另眼相看,必須加以考慮了.關於保羅為主受苦,相信我們大家都很熟知；不過,多看別人為主之故受苦,對我們很有幫助；可能我們有著同樣的經歷,那就更能得到鼓舞.
\newline
\newline
首先,我們來看保羅在身體,物質,和環境方面所遭遇的.
\newline
\newline
聖經告訴我們,保羅身上有根刺,到底這根刺是甚麼,保羅自己沒有說明.有人說他的眼睛不好,需要別人替他寫信；當他簽名之後,他說:「你看!我簽的字多麼大!」他也對加拉太的信徒說:「你們愛我,甚至願意把自己的眼睛剜出來給我.」可能他的眼睛非常不好,這帶給他極大的痛苦；但主對他說:「我的恩賜夠你用的.」
\newline
\newline
在教會歷史上,主所重用的信徒,有的身體殘缺不全；但主藉著他身體上的痛苦來彰顯祂自己的榮耀.有個失明的姊妹,她從主所體驗的,作成詩歌帶給許許多多的人；蒙主賜福從幾百年前直到現在.我們真是看見了主賜福.
\newline
\newline
保羅在物質方面時常遭遇困難,有時吃不得飽,穿不得暖；當他乘船要到羅馬去,海中船被破壞；保羅告訴我們,曾一連三次他有這樣的經歷,甚至在海上晝夜飄泊.為主的緣故,保羅常受苦難.現今我們在物質方面的享受較好,但可能未來的年日,主要我們為祂而受更大的考驗；一旦苦難臨到,我們能否站穩呢?廿一年前,我曾在香港居留半年；有位從前內地會的教士,他曾經多年在山東事奉主,他也是位醫生,後來他在香港北角行醫；我初到香港時,他請我住在他家裡,特別為我在他的候診室之一角落,裝上簾子,擺好床鋪；我就住在那裡約有一兩個月.之後我到台灣的高山傳福音,在一個小村莊作開荒工作.那裡未有基督徒,也沒有旅店,我家住在山地同胞的石屋,那裡也是警方的派出所.過了一段時間我重臨該地,那時已有基督徒了；我就住在他們的家裡.在中華福音神學院,有一天,同學來報告說:「另一位同學不睡自己的床上,卻睡別處的地上.」經我了解之後,才知原來他是學習過簡陋克苦的生活,準備將來出去作開荒的工作,向自己的同胞傳福音.
\newline
\newline
保羅在各方面為主受苦　他在大馬色初信主時,當地的官長要捉拿他,看守城門防他逃脫；他到腓立比,挨打,被下監牢,受虐待；到耶路撒冷,百夫長和君王,在保羅身上施加壓力,處處都是不如意的事；但是保羅為主的緣故,謙卑忍受.保羅到了雅典,受到另一種壓力,乃精神的痛苦；他是個學者,在那些哲學家面前講耶穌復活,他被譏笑,受蔑視,然而他仍不以福音為恥.他在哥林多,那時眼睛已經失靈,他立定意志,不作別的,不講別的,不講任何高等學問；只講耶穌基督和祂的十字架.保羅受到自己的同胞攻擊,這是極其難堪的；他在耶路撒冷,起來反對他的,不只是羅馬人,而多數是自己的同胞；他在海外傳福音,反對最激烈的是猶太人.保羅說:「我是猶太人,我自己的同胞為何如此對待我?」他說:「為我弟兄,我骨肉之親,就是自己被咒詛,與基督分離我也願意.」(羅九3)保羅這樣的愛他的同胞,但他們卻誤會他.
\newline
\newline
戴德生在世上傳福音,很願意和中國人接近；我想這是許多人都記念他的原因.當他在上海開始工作,有次他到鄉村,穿著一般英國人所穿的衣服；他盡力講道,發現座中有位非常專心的聽者.他感到很欣慰,越講越起勁；講完道,看見此人舉手.說:「我有問題要問,」戴德生以為他大受感動,十分高興；想不到他問:「講道先生,你衣服上的釦子是甚麼意思?」從那時開始,戴德生不再穿西裝,他的穿著,盡量適應聽眾:不過,經此服裝改革,除了本身覺得不習慣之外,他本國英人對他的服裝大不滿意,認為這個英國青年,有失國家尊嚴,應穿本國的服裝才對.他的同胞不了解他,對他施加壓力；然而為了愛主,為了福音的緣故,他不顧一切.
\newline
\newline
我曾提過,巴拿巴需要同工,到大數去找保羅.解經家根據「找」這個字,證明保羅因信耶穌之故,回家以後與家人分手.如果我沒有記錯,席勝魔傳記,他信了主,家人極力反對,後來他帶領一位范先生認識耶穌；他信了主,回家時也受到家人厲害的反對.數年來,我發現神學院每年進院的新生,其中必有一兩個受到家庭極大的壓力.在四年前有個新生,他是已經讀完大學,服滿兵役,且是已結婚的青年；當時他的妻子將要分娩,進入醫院,寫了一封信給我.內容是說:「昨天他的父親到醫院來看我,對我說:「我兒子信了耶穌,住在神學院；如果他再不回家,我將與他脫離父子關係.」院長:請你鼓勵他繼續走奉獻的道路.」在兩年前一個晚上,有位女同學來告訴我,次日早晨她的父親要來領她回家.當晚,我和她為此事同跪在辦公室禱告；次日早禱聚會之後,她父親果然來了；我避免和他見面,因恐事態演變複雜,所以請同工輔導主任代為接見.他們又談了個多鐘頭,聲音越談越大,毫無結果；後來他說要見院長.我只好接見.我們大談了個多小時,最後他說:「戴先生!問題很簡單,今天我女兒若不跟我回家,就循法律途徑,脫離父女關係.」但是,感謝主!直到現在,那位姊妹仍在神學院,還有一年將畢業了,她的父親並沒有與她脫離關係.(時間關係不能詳述)因保羅一人站穩想給他全家蒙福.因上述那位弟兄和那位姊妹的站穩,想給家庭極大的蒙福；雖然在當時幾乎構成家庭分裂,但我多次看見因著一個青年基督徒肯站穩,不搖動,結果全家蒙福.正如保羅說:「當信主耶穌,你和你的一家都必得救.」約書亞說:「至於我和我家,我們必定事奉耶和華.」
\newline
\newline
保羅信主不久回到耶路撒冷,基督徒都不敢相信他,這實在是很大的考驗!過了一段時間,一班猶太派的信徒常不了解保羅；以為他不愛國,不愛摩西,破壞舊約；所以他們反對保羅.有更多的例子記載在哥林多後書,保羅詳細說明他為主的緣故受苦.
\newline
\newline
保羅看「為主受苦」的角度和感受　保羅認清為主受苦是必然的,他對提摩太說:「你要和我同受苦難,好像基督的精兵.」(提後二3)這句話我很喜歡,保羅並非自己高高在上,過著舒適的生活,而給前線上的小兵寫信；保羅卻是在監牢中寫信給提摩太.感謝主!提摩太知道!他家住在路司得,親眼看見保羅第一次到路司得,幾乎被敵人用石頭打死；他知道作為基督徒要為主作見證必須受苦.保羅也對提摩太說:「凡立志在基督耶穌裡敬虔度日的,也都要受逼迫.」(提後三12)他們看為主受苦是一種榮耀.有的基督徒受到苦難時,怨天尤人,以為主不愛他.第五章41節記載,那些門徒挨打,遍體鱗傷；但「他們離開公會,心裡歡喜；因被算是配為這名受辱.」求主給我們有這樣的看見!
\newline
\newline
各位青年!各位弟兄姊妹!在學校,在家中,在工作處為主的緣故；受欺侮,遭逼迫,有壓力臨到,我們能否欣然接受?因被算是配為主受苦；但是我們對此應加小心,因為在中世紀,有種錯誤的觀念進入教會,有些基督徒故意尋找麻煩；教會歷史記載,一直到第四世紀,有許多基督徒願意殉道,此後羅馬帝國漸受基督徒影響,甚至包括皇帝在內；所以基督徒再不受苦了.
\newline
\newline
保羅所受的苦很多,他把自己所受的苦難一一顯示出來,來證明他是神的僕人.求主給我們有這樣的心,當苦難來到,主看得起你,才讓苦難臨到你；因為你是主的僕人.
\newline
\newline
保羅清楚看見比身體的舒服,安全,甚至比他生命更重要的,就是當跑的路要跑盡；所以有一天,保羅到了該撒利亞腓利的家,門徒攔阻他上耶路撒冷去.保羅說:「你們為甚麼這樣痛哭,使我心碎呢?我為主耶穌的名,不但被人捆綁,就是死在耶路撒冷,也是願意的.」(徒廿一13).
\newline
\newline
保羅關心弟兄姊妹,他要堅固教會,他關懷弟兄姊妹愛主的心,由多方面可以見得,保羅在路司得被敵人用石頭打得半死,他起來之後沒有回家；雖然距他老家「大數」僅一山之隔,但他繼續前往特庇又轉回路司得,就是不久之前敵人要把他打死的地方.然後他到以哥念,到安提阿,為要堅固各教會的弟兄姊妹,因恐他們所受的逼迫過於嚴重,為了弟兄姊姝的緣故,他冒著生命的危險!他輕視自己的身體,輕看自己所受的痛苦.其實,他最後一次回耶路撒冷,並沒有重要的信息要傳,他可以不去.聖經告訴我們,他是把哥林多,帖撒羅尼迦,腓立比的弟兄姊妹愛心之奉獻帶去；讓耶路撒冷的弟兄姊妹看見,這些外邦教會願意站在同一陣線；看見希臘的教會,和耶路撒冷的教會是合一教會.保羅為了合一的見證,費盡心血,遭受苦難,冒大危險；他更關心的是主耶穌的榮耀,他看見將來的榮耀,輕看目前的苦難；兩者相比,苦難不過是至輕至暫的,榮耀才是永遠的.
\newline
\newline
保羅具有上述種種的認識,才能夠在苦難中站穩,他決意要順服主.當他在亞基帕王面的時候說:「我沒有違背從天上來的異象.」主耶穌特別強調基督徒要學習順服.主說:「遵行天父旨意的人才能進天國.」重點在於順服.主耶穌所講蓋房子的比喻,兩個蓋房子的人的分別要緊在於基礎.保羅說:「那已經立好的根基是主耶穌,此外沒有人能立別的根基.」(林前三11)耶穌說:「凡聽見我的話不去行的,好比一個無知的人,把房子蓋在沙土上.」「凡聽我的話就去行的,好比一個聰明人,把房子蓋在盤石上.」(太七24-27)我們不但聽道,也要行道,奠定牢不可破的基礎.我們應當順服!
\newline
\newline
保羅常與基督徒有交通,他初得救時,在大馬色常與門徒在一起,繼之；他到耶路撒冷,也是和門徒在一起；後來到了羅馬,從弟兄們得著鼓勵,得了安慰,重新得力.各位基督徒!個人孤立,容易跌倒,和基督徒在一起,有一種能力.(來十25)說:「你們不可停止聚會,好像那些停止慣了的人,倒要彼此勉勵；既知道那日子臨近,就更當如此.」感謝主!香港的信徒,許多工作都以教會為中心,這是應該的；因為在神的家中,彼此有交通,軟弱時,互相幫助,得著安慰.當一個基督徒覺得無需別人時,他必站立不住.求主幫助我們!但以理有三個朋友,當他有困難時,立刻把問題帶去,請他們代禱.我們需要禱告的同伴,基督徒的交通實在蒙福的,可帶來力量.在主裡的交通,主裡的關懷,實在太重要!兩個人在一起,主就在他們中間.
\newline
\newline
主應許祂的同在,不是一種感覺；有時保羅的感覺很不舒服；但主在他身邊,他有喜樂.他是第一個到歐洲的傳道者,不受歡迎,敵人用棍打他,被下監牢,全身痛楚；在潮濕的腓立比監獄,他不能睡；此情此景,他卻唱歌禱告；因主在他身旁.感謝主!主與我們同在.
\newline
\newline
我們若在苦難中,忠心順服主,與弟兄姊妹保持交通；有主的應許,有主的同在,我們能勝過一切的考驗.前面的預言,有許多考驗,可能有許多苦難；求主幫助我們,定睛在祂的身上,祂作為我們的榜樣.我們應當忠心順服祂!
\newline
\newline
當你看完使徒行傳,你是否有奇異的感覺?使徒行傳斷了嗎?結論在哪裡?我的回答是:結論在你我的手中!因為行傳還未寫完,歷史尚在延續.未來的一年,在港九培靈研經會未臨之前；你的生活,你的見證,要成為這歷史的重要部分.但願感動使徒的靈,感動你,帶領你,賜能力給你,充滿你!阿們.
\newline
\newline


\newpage

\section{第49屆~港九培靈會~楊濬哲~序}
\label{sec:1385}
\textbf{序}
\newline
\newline
連結: \href{https://www.hkbibleconference.org/session-message/view/1385}{\texttt{https://www.hkbibleconference.org/session-message/view/1385}}
\newline
\newline
\hyperref[sec:795]{< < < PREV SERMON < < <}
~
\hyperref[sec:toc]{[返目錄]}
~
\hyperref[sec:768]{> > > NEXT SERMON > > >}
\newline
\newline
本講道集與各位見面時,適逢是港九培靈研經大會,第五十屆金禧紀念之期；但願我們帶著萬分歡悅的心情,來迎接這盛大的紀念會.
\newline
\newline
由一九二八年一九七八年,這漫長的歲月中；神的眼,神的心,常看顧賜福與培靈會.五十年來播下的種子,經過多少眼淚和禱告的孕育,終於成長了.教會由極荒涼極冷淡的景況下,已甦醒過來,而且不斷地增長,更新,成熟；培靈會不敢居功,但對港九教會帶來的復興,也盡上了綿力!
\newline
\newline
我們更要感謝神的,近十年來,各神學院的質素,不斷提高,奉獻入學的青年,也陸續增加,而且多間神學院,甚至供不應求,有遺憾之感!培靈會在這方面,也有直接或間接的影響.
\newline
\newline
回憶初期的講道集,因限於人力,物力,都非常簡陋.有兩年合刊的,也有數屆停了出版的,總題亦欠完整的.直至六零年以後,方漸循完善.
\newline
\newline
本講道集是第四十九屆的講道記錄.三位講員的信息,都值得向各位推薦介紹.研經會由艾理德牧師主講.艾牧師是建道神學院的副院長,原文精湛,詳於解經,對約翰壹書有獨特的見解,很能幫助我們.講道會由張子華牧師主講,張牧師是中國神學院研究院的副院長,他對輔導工作,有特長之處；這次講解和諧生活,對基督徒個人與神,與人,與家庭的關係,都有深切的交代.奮興會是由戴紹曾牧師主講,提起戴牧師,我們會聯想到他的曾祖父戴德生牧師,中國內地會的創辦人.我們很幸福,得見到戴德生牧師的曾孫,而且他能繼承曾祖父的心志事奉主.戴紹曾牧師是中華福音神學院的院長.他對使徒行傳與今日教會的需要──聖靈的工作,分析得十分透切,誠屬不可多得的時代信息.戴牧師的中文修養是湛深的,可謂名實相符的中國通,可見他是如何的愛中國教會.深望每位讀者都能重溫本書的信息,使你靈性增益.
\newline
\newline
第四十九屆大會的人數,可說是空前的,不但正副堂滿座,連浸信會的主日學課室,詩班室,會客室,都擠滿了人；我們親眼看見神豐富厚恩,實在不能不發出感謝頌讚,歸榮耀給神!
\newline
\newline
蒙九龍城浸信會多年來接待,同工們的合作和招待,實在萬二分的感激不盡!
\newline
\newline
本講道集蒙鄧作良牧師,朱雷麗端姊妹記錄,朱建磯牧師編輯,朱楊麗華姊妹校對,得如期出版,他們忠心勞苦,願主記念!
\newline
\newline


\newpage

\section{第49屆~港九培靈會~艾理德~第一講}
\label{sec:768}
\textbf{當心你的生活}
\newline
\newline
連結: \href{https://www.hkbibleconference.org/session-message/view/768}{\texttt{https://www.hkbibleconference.org/session-message/view/768}}
\newline
\newline
\hyperref[sec:1385]{< < < PREV SERMON < < <}
~
\hyperref[sec:toc]{[返目錄]}
~
\hyperref[sec:769]{> > > NEXT SERMON > > >}
\newline
\newline
本講道集與各位見面時,適逢是港九培靈研經大會,第五十屆金禧紀念之期；但願我們帶著萬分歡悅的心情,來迎接這盛大的紀念會.
\newline
\newline
由一九二八年一九七八年,這漫長的歲月中；神的眼,神的心,常看顧賜福與培靈會.五十年來播下的種子,經過多少眼淚和禱告的孕育,終於成長了.教會由極荒涼極冷淡的景況下,已甦醒過來,而且不斷地增長,更新,成熟；培靈會不敢居功,但對港九教會帶來的復興,也盡上了綿力!
\newline
\newline
我們更要感謝神的,近十年來,各神學院的質素,不斷提高,奉獻入學的青年,也陸續增加,而且多間神學院,甚至供不應求,有遺憾之感!培靈會在這方面,也有直接或間接的影響.
\newline
\newline
回憶初期的講道集,因限於人力,物力,都非常簡陋.有兩年合刊的,也有數屆停了出版的,總題亦欠完整的.直至六零年以後,方漸循完善.
\newline
\newline
本講道集是第四十九屆的講道記錄.三位講員的信息,都值得向各位推薦介紹.研經會由艾理德牧師主講.艾牧師是建道神學院的副院長,原文精湛,詳於解經,對約翰壹書有獨特的見解,很能幫助我們.講道會由張子華牧師主講,張牧師是中國神學院研究院的副院長,他對輔導工作,有特長之處；這次講解和諧生活,對基督徒個人與神,與人,與家庭的關係,都有深切的交代.奮興會是由戴紹曾牧師主講,提起戴牧師,我們會聯想到他的曾祖父戴德生牧師,中國內地會的創辦人.我們很幸福,得見到戴德生牧師的曾孫,而且他能繼承曾祖父的心志事奉主.戴紹曾牧師是中華福音神學院的院長.他對使徒行傳與今日教會的需要──聖靈的工作,分析得十分透切,誠屬不可多得的時代信息.戴牧師的中文修養是湛深的,可謂名實相符的中國通,可見他是如何的愛中國教會.深望每位讀者都能重溫本書的信息,使你靈性增益.
\newline
\newline
第四十九屆大會的人數,可說是空前的,不但正副堂滿座,連浸信會的主日學課室,詩班室,會客室,都擠滿了人；我們親眼看見神豐富厚恩,實在不能不發出感謝頌讚,歸榮耀給神!
\newline
\newline
蒙九龍城浸信會多年來接待,同工們的合作和招待,實在萬二分的感激不盡!
\newline
\newline
本講道集蒙鄧作良牧師,朱雷麗端姊妹記錄,朱建磯牧師編輯,朱楊麗華姊妹校對,得如期出版,他們忠心勞苦,願主記念!
\newline
\newline
第一講　「當心你的生活」
\newline
\newline
首先多謝港九培靈研經大會委辦會,邀請我擔任此次研經會講員.香港教會具有優良的傳統,其中心乃在一班教會領袖身上.我到香港工作,我的工作對象是中國的弟兄,而不是西方的弟兄；如今有機會與大家一同研經交通,實在是本人莫大的權利.這次選讀約翰壹書,每早晨從大綱中抽出一段；盼望盡我所能的,為大家講解得完備,使大家能明白經上教訓,並應用在我們的生活上.
\newline
\newline
引言
\newline
\newline
(一)歷史,宗教相關的問題
\newline
\newline
甲,當時背景教會成立初期,就有一新宗教運動發生.它有一基本假設──各宗教從不同的途徑達到同一目標.這說法和今日許多人主張的也相像.他們所提出的特別教訓大致有四點:
\newline
\newline
(1)物質世界是邪惡的.
\newline
\newline
(2)人的靈在肉身內被囚束縛,必須將受束縛之靈釋放出來.
\newline
\newline
(3)分別有黑暗世界和光明世界,光明世界是神的居所,而黑暗世界乃是世人所處之地.
\newline
\newline
(4)從黑暗世界渡到光明世界,必須憑著特別的知識；所以本書中常見「知識」,「光」,「黑暗」,等字眼.
\newline
\newline
在現今世代也常遇到這些.不久之前我在建道神學院辦公室內,秘書告訴我有兩位西方人來訪,接見後原來是兩個摩門教的傳教士.開始談論時,他們鄭重地向我介紹,他們具有的所謂特別知識和啟示；如果我肯接受就能得著真理云云.許多人著重特別知識,又有不少人主張各種宗教殊途同歸；以致當日信徒中間產生混亂.約翰就當時的需要,寫本書信,使教會信徒能辨別是非.
\newline
\newline
乙,基督徒與教會面臨的問題　當這新運動與教會接觸時,就產生以下的兩種難題:
\newline
\newline
(1)那些人能接納基督教,而基督教能否接納他們?我從前在加利福尼亞州攻讀博士學位,研究宗教哲學,教授為世界有名哲學家之一；我功課中之一項,要選擇一哲學家,分析他的思想並加以批判.我也曾在另外之哲學研究部門寫過類似的論文,並有很高的成績.於是就把那些資料交給我的教授.後來他約我到辦公室交談,對我的論文所作研究和批判表示讚賞.但他對我說:「你也應讓這哲學家改變你基督徒之信念.」基督徒批評他的思想,但也應接受他的思想；兩者調和才能有新認識,許多人對基督徒如此說,約翰當時也有這問題.(2)有許多人說:「教會須改變他的信仰,用現代的言辭來表達.」約翰稱這些人為假先知.他們的錯誤教訓有兩方面:
\newline
\newline
(1)否認「道成肉身」,以為神不能有身體.
\newline
\newline
(2)自以為具有特別知識,犯罪也不算為罪.故此,教會大受影響!教會合一大受破壞,正確教訓已經受混亂.使徒約翰關心教會面臨種種問題,故寫成本書,將正確的教訓述明,使信徒得到喜樂.
\newline
\newline
(二)解釋的原則:「相信與實行兩者的和諧.」
\newline
\newline
約翰作出一個假設,是他一切教訓中心──信仰與實行必須一致.人若認識神,要從生活行為上表明出來,不可背道而馳.為要更明白這中心原則,我們試從約翰福音和壹書中,歸納合併來看,以四句話來分別說明.
\newline
\newline
甲,基督教並非宗教,乃是生活方法.許多青年人說:我若能得多些關乎基督教的知識,就能夠成為一好信徒；也能與人交談辯論,將基督教的優越性向人表明.但基督教的真義決非這樣,乃從基督接受生命並與人分享.記得許多年前我在加拿大大學裡唸書,某個早上乘公共汽車回校上課.途中,想向鄰座男子分享我的信仰.一方面受到聖靈催迫,另一方面又怕對方提出難答的問題.終於打開話匣,並將話題轉到信仰上去.對方突然打斷話柄,反問我一問題:「神是否全能?」我答:「當然.」他說:「若然,則神會不會造一塊巨石,龐大到自己不能舉起的呢?」當時我不識怎樣回答,覺得非常沮喪.多年後我才找到問題的癥結:我們若把自己的信仰,作成一個思想體系好像別的宗教一樣.則這立論一被攻破,整個信仰就站不住腳.因此我重申:基督教是一個生活方法,我們從主接受生命,與神有交通,然後將生命與人分享.
\newline
\newline
許多人作基督徒時,自己定下許多行事的標準.例如每週崇拜聚會,每日如何讀經祈禱.然後循序漸進.但往往過去相當時日,基督徒的生活會從高峰掉到深淵,由於屢起屢蹶,就落在非常沮喪的光景,認為基督徒的理想生活是不能達到的.我們從教會歷史上,也可認明過去中世紀時,信徒也曾在這錯誤觀念上度過不少年日.今讓我們將這昇降梯子側放,倒能令我們找著正確的生長過程.
\newline
\newline
其次又有人以為,神掌管著一本其大無朋的書,書上載有各別信徒的名字.然後我們的一言一動,無不載在書卷內.但這也不過對神法則的誤解.
\newline
\newline
讓我們來作個比喻:我有三個子女,我的幼女現正開始學行路.我對他說:「行路有兩原則,一是要讓身體平衡；其次要將一足移到另足之前.這就是行路的方法.」我對她解釋清楚,然後放她自己學步.當然她不久傾跌於地,於是我又扶她起來,對她再解釋；又復叫她行走,結果不例外地又再跌倒,再三嘗試都是一樣,然後我又打她一頓,再將規則重頭說一番.許多時我們對神的設想也如此,我們的神有許多規則,希望有人將神所有規則教訓我們,我們就能作好基督徒.其實,我們教兒童學行,首先是放開她,讓她試著行路；又站在她旁邊注視著扶她.當她跌下時,我們就扶起她,哄她,並且抹乾她的眼淚.這樣,不久小女孩就能行,又跑,又跳.神在信祂人身上的工作也一樣.神不總是在計算我們的失敗錯誤.祂愛我們,樂意像慈愛的父親一樣,幫助我們生活得更完滿.當我們走錯路時,祂管教我們.好像孩子跌倒時祂又扶起.至終使我們能行,能跑,能跳.使我們知道如何敬拜,認罪,悔改,信靠,順服,愛神,而生活在基督裡.
\newline
\newline
乙,基督教並非信條,乃是團契交通.約翰福音曾用一特別名詞,英文繙為「交通」,「團契」.說明屬神之事能以交通分享.參林前十一章25-26節:「......若一個肢體受苦,所有肢體就一同受苦；若一個肢體得榮耀,所有的肢體就得榮耀.」對於這經文我們曾否實行?一位肢體受苦,其他與他一同受苦較易；若一位肢體得榮耀,各人與他一同快樂則較難.應知我們眾信徒在基督裡是合一的團契,彼此相愛應該表明在信徒之間.數年前,我曾乘渡海輪由九龍往香港,擬赴某處講道.正沉思講道內容,忽見坐在前列位有一女士,從座中起立高歌.揮動手臂,狀若曾服食迷幻藥.坐在她旁邊的人紛紛離座他往.我注視她的樣貌,恍忽見到她曾在這世上歷盡滄桑.我忽然想到是否在世上,有一人肯愛這樣的人呢?是否她還有一些地方,值得被人接納的呢?我便開始為這人祈禱,並且覺得有一個地方她應該去的,就是信徒的團契.在那裡她將得到接納和愛,她不需要隱藏她的罪過,她可以得著無條件的愛.在教會中,有些人將自己的難題深深地掩藏起來,裝作若無其事,因為恐怕別人知道後會有不良影響,其實這對教會屬靈生活有莫大傷害.過去曾有美國同工,為他女兒有些煩惱事,請我去與他女兒作個人交通及祈禱.事後,我問他為何將這些事告訴我?他說:「我知你知道這一切事,也不會改變你對我的看法.」如果有弟兄能與你分享一些內心的秘密時,你知道後是否會改變對他們的看法,抑或繼續不變地愛他?
\newline
\newline
丙,基督教並非倫理,乃作主門徒的回應.
\newline
\newline
丁,基督教並非定義,乃是宣告.基督教並非著重在許多名詞上加以解說.正如中國的儒家學說,論語對「仁」一字,孔子並不下定義而清楚解說.孔子明白到,重要的不在這字本身的意義,乃重在實行出來.哲學上多用定義來說明真理,但基督徒都知真理是甚麼.基督教不是定義,乃是宣告.從前我學中文,首先學寫自己的名.我照著教師所寫的來模倣.寫艾理德的「德」字時,先寫右邊,後寫左旁.教師糾正我,要從左邊到右邊.我初覺教師太武斷,何以不能兩者均可.但我知道不是教師的主觀,乃是中國文字的通則.這也可藉之說明基督教,並非了解許多知識,然後與人作學理的辯論；乃是基督所教導的,祂的門徒必須實行出來.人若認識神,必須與所行的一致.故此,基督教的本身,就是一項宣告.
\newline
\newline
今天講完本書簡單之背景,然後往後數天繼續查考內文.此書著重信仰與生活的實行.望各位能從本書得到幫助.
\newline
\newline


\newpage

\section{第49屆~港九培靈會~艾理德~第二講}
\label{sec:769}
\textbf{與神與弟兄相交}
\newline
\newline
連結: \href{https://www.hkbibleconference.org/session-message/view/769}{\texttt{https://www.hkbibleconference.org/session-message/view/769}}
\newline
\newline
\hyperref[sec:768]{< < < PREV SERMON < < <}
~
\hyperref[sec:toc]{[返目錄]}
~
\hyperref[sec:770]{> > > NEXT SERMON > > >}
\newline
\newline
約翰一書 1:1-1:4
\newline
\begin{longtable}{cl}
\hline
\hline
章節 & 經文 (和合本修訂版)\\
\hline
1:1 & \begin{tabularx}{0.7\textwidth}{X} 論到從起初原有的生命之道，就是我們所聽見、所看見、親眼看過、親手摸過的－ \end{tabularx} \\ \\ \relax
1:2 & \begin{tabularx}{0.7\textwidth}{X} 這生命已經顯現出來，我們看見了，現在又作見證，把原與父同在，並且向我們顯現過的那永遠的生命傳揚給你們－ \end{tabularx} \\ \\ \relax
1:3 & \begin{tabularx}{0.7\textwidth}{X} 我們把所看見、所聽見的傳揚給你們，為要使你們也與我們有團契，而我們的團契是與父和他兒子耶穌基督所共有的。 \end{tabularx} \\ \\ \relax
1:4 & \begin{tabularx}{0.7\textwidth}{X} 我們把這些事寫給你們，使我們的喜樂得以滿足。 \end{tabularx} \\ \\
[1ex]
\hline
\hline
\end{longtable}
我今早搭長洲渡輪來港,船上有一英國人在讀報紙遙見報紙上標題:「世界正步向毀滅的邊緣」第三次世界大戰的危機正日增,如果沒有主,世人活下去還有什麼盼望今早讓我們來思想約翰壹書的序言,講出一重要主題.因為這一個主題,能使人活在世上有盼望.
\newline
\newline
這主題就是「生命之道」(1)學者對它有幾種不同見解;.第一,認為指耶穌基督自己第二,以為是主耶穌所帶來的信息但我的意見是屬第三派綜合上述二者,不能將耶穌與祂的信息分開講到主耶穌,也就是祂宣告的信息在此,約翰對「生命之道」寫出五個特徵;.並分別用五個字來描寫:
\newline
\newline
甲,起初──道先存的事實「論到從起初原有生命之道.」(1)「......原與父同在......」(2)參看約翰福音一章一節:「太初有道,道與神同在.」歸納兩家經文,有幾點意思:生命之道非人所創立,乃神所賜而永恆的.祂不改變,是可信靠的；帶給人類永恆的信息.這就是約翰對生命之道的宣告.但馬上有人質疑:既說從起初就有,何以我卻一無所知?試舉一比方來說明.假設整個世界祇有一個顏色──紫色──無論人物,樓房,街道,通是紫色.後來有個人作見證說,他見過紅色,起初當然沒有人信有這個可能.後來見證的人多起來,人就確知有紅色.當人未知有紅色時,其實它早已存在.因紫色原是紅,藍二色合成的.約翰說:這生命之道,是起初原有的,非人所創；當人將生命打開,就能夠經歷到祂,是無比奇妙的經驗.
\newline
\newline
乙,顯現──道顯現於歷史中讀第1-2節,約翰用了三個動詞:「聽見」,「看過」,「摸過」.用以答覆他們的敵對者,「神不能有身體」的理論.約翰聲言,他曾經聽過,看過,摸過.就徹底否定對方的論調.特別「摸」一詞,乃包含詳細考察,驗證之意.這「生命之道」不單太初存在,還進入歷史人世之中.這是一個頗重要的問題,神的兒子為何需要道成肉身,來到這世界?我想有兩個原因:第一,人犯罪是在這世界之中,並不在世界以外.創世記記載:撒旦如何向夏娃進行試探,他提出一個宗教性問題:「神真是好嗎?」夏娃答:「祂當然是好的.」正在她維護神的時候,就不暇細察撒旦的詭計,他繼續找到弱點進攻:「神既是好的,就不會因你食這樹上果子叫你死亡.祂一定不會為這小事責備你.」夏娃就陷在惡者的網羅.神要求我們在一切事上順服祂.如果我們有一件事不順服,主就是要為你在這一點上,為你來到世界,解決你的罪.第二個原因,神如果祇是一個抽像觀念,就不能成為敬拜的對象和信仰的基礎.舉例來說:我想你認識我兩女兒中之一,於是我對你縷述她種種特點:她是女性,姓我的姓氏,住在我家；縱使我不厭其詳地再加描述:她的重量,體型,高度,髮色.你也無法構成我女兒的具體形像.照樣我們對神的特徵,譬如說:「神是愛.」我們何由證實.如果我們當中有人需要愛,關懷,他的問題可在基督裡得著解答.祂在十字架上為我們死.也就是祂親身進入世界的理由.
\newline
\newline
丙,傳給──信息的見證你見證有兩個條件:約翰在此提到:「所看見,所聽見.」
\newline
\newline
(3)是指宣告的信息,經由個人的經歷.若有人被召到法庭作證人,他的證供必需是他親眼所見,親耳所聞.第二個條件,約翰表明,乃由神所差派來傳這信息的.希臘文「傳給」一詞有重要的意思,即「確實而原原本本」傳給的意思.
\newline
\newline
丁,相交──屬靈的團契　「使你們與我們相交,我們乃是與父並祂兒子耶穌基督相交的.」(3)「相交」一詞非常重要.人為福音作見證仍然未夠,必需帶進團契,使眾人與神相交.成為基督身體上的肢體,是何等榮耀的屬靈團契.今早在座的眾弟兄姊妹,許多人互不認識.但為著同一目的要領受神的話,就何等和諧地共聚一堂.基督召我們,使我們都屬於祂,因此我們彼此相屬,將我們聯在一起.有一天,神從各方,各國,把稱呼主名的人帶到祂寶座之前,那時他們唱詩讚美俯伏敬拜神.如今,一個肢體受苦,眾肢體也一同受苦；一個肢體得榮耀,眾肢體也隨之得榮耀.有一位朋友,是為總統駕駛飛機的:一次,到達羅馬尼亞國內,他辦完手續後,在街頭漫步.突然聽見有歌聲送來,雖然不明白歌詞的說話,但音調是熟識的.他趕忙上前就近那人,但苦於言語不通；但當他用英文唱出同一聖詩時,他們就完全明白,兩人擁抱起來.我們屬於一個家庭,彼此相交團契,也是與神相交.
\newline
\newline
但我也要鄭重地說,這相交團契也很容易被摧毀,罪會使之被破壞.我們不能這樣說:「作某些錯事與別人無關.」須知同屬一個身體,不能個別地從其中抽出來.因一人一次的不順服,雖說他不願意傷害這身體,但身體的相交團契到底被損害了.我們不單需有好的聚會,有好的講道和講經；我們更需要敞開己心,求主鑑察使彼此交通,不受攔阻而得復原.我們得學習流淚痛哭,分擔痛苦；學習彼此相愛,這在信徒中間,乃是十分重要.
\newline
\newline
戊,喜樂──福音的果效「我們將這些話寫給你們,使你們的喜樂充足.」(4)約翰所用「喜樂」一詞,有喜樂繼續不斷充足的意思,也有恩典之意.尼希米八章10說:「因靠耶和華而得的喜樂是你們的力量.」福音帶給相信的人有喜樂.也是基督徒的特徵,也給信徒有生活的能力.世人也知歡喜,但絕不是基督徒所得著的喜樂.詩篇廿三篇說:「使我的福杯滿溢.」當我們在聚會之中,聆聽了主的話,真使我們的福杯滿溢.有何等美好的經驗,證明主裡面的喜樂,是何等地甘美.當我們運用信心讚美神,神賜給我們問題的答案,使我們明白,神賜給的是何等的幫助,何等的愛.
\newline
\newline
哈巴谷三章16-17節是非常美的兩節聖經.先知為他的國,他的同胞向神呼求,當神在一章6節說:「我必興起迦勒底人......」來答覆他的禱告,先知不禁對神發出怨言來.當我為教會為信徒求神幫助時,神說:要興起仇敵來與我們爭戰,也是沒有人肯接受的.先知開頭也如此.但後來,他思想到神的作為,祂是一切難題的解答,神必明白祂自己的作為.所以先知說:「雖然無花果樹不發旺,葡萄樹不結果,橄欖樹也不效力；田地不出糧食,圈中絕了羊,棚內也沒有牛；然而,我要因耶和華歡欣,因救我的神喜樂.」親愛的朋友,你能否如此讚美神?我們不是在世上作孤兒,在天上的父關心地愛我們.除非經祂許可,否則事情不能臨到我們身上.無論在生存死亡中,神都看顧施恩.當我們如此感謝讚美神,喜樂與能力就從我們身上得釋放.
\newline
\newline
有一位事奉神的人,多年來為他兄弟,不斷忠心地代禱.直至有一日神在對他心中說話,使他停止祈求,而對神感謝讚美.繼而,開始見到神在動工,他兄弟突從東邊原來工作處調往西邊.他又遇到一位基督徒同事,與他分享自己的信仰,兩月後成了一位真正的基督徒.我們要學習喜樂,祂本身有大能力,神是知道如何成就祂自己工作的.
\newline
\newline
約翰說:「我們將這些話寫給你們,使你們的喜樂充足.」學習在主裡喜樂,這是基督徒的特徵.怎能喜樂?因神愛我們,關心我們；在困難中,在夜幕深沉之中,祂與我們同在.
\newline
\newline
蘇格蘭有位有名傳道人,年輕時在一家店舖內作工.一日,工作至很晚,需走五哩路回家.半路有森林,常有強盜出沒.當這少年行到一段黝黑地帶時,因心中恐懼,不期哼著詩歌來壯膽.突然聽見遠處有呼他名字的聲音,便愕然卻步；接著明白不遠處站的是他父親,就不期將恐懼完全驅除.父親與他同行,聽見父親的腳步聲,就恍如有回到家裡的寧靜感覺.
\newline
\newline
生命的信息,永恆的神為你我的罪,來到世間.我們應該停下來聽,我們天上的父在這裡；無論夜如何黑暗,在失望中,祂與你同在,等待著要與你同行.屬祂的兒女們,神要照顧你,叫你們的喜樂充足.
\newline
\newline


\newpage

\section{第49屆~港九培靈會~艾理德~第三講}
\label{sec:770}
\textbf{與基督教不能分開的倫理宣告}
\newline
\newline
連結: \href{https://www.hkbibleconference.org/session-message/view/770}{\texttt{https://www.hkbibleconference.org/session-message/view/770}}
\newline
\newline
\hyperref[sec:769]{< < < PREV SERMON < < <}
~
\hyperref[sec:toc]{[返目錄]}
~
\hyperref[sec:771]{> > > NEXT SERMON > > >}
\newline
\newline
約翰一書 1:5-2:29
\newline
\begin{longtable}{cl}
\hline
\hline
章節 & 經文 (和合本修訂版)\\
\hline
1:5 & \begin{tabularx}{0.7\textwidth}{X} 神就是光，在他毫無黑暗；這是我們從主所聽見，又報給你們的信息。 \end{tabularx} \\ \\ \relax
1:6 & \begin{tabularx}{0.7\textwidth}{X} 我們若說，我們與神有團契，卻仍在黑暗裡行走，就是說謊話，不實行真理了。 \end{tabularx} \\ \\ \relax
1:7 & \begin{tabularx}{0.7\textwidth}{X} 我們若在光明中行走，如同神在光明中，就彼此有團契，他兒子耶穌的血就洗淨我們一切的罪。 \end{tabularx} \\ \\ \relax
1:8 & \begin{tabularx}{0.7\textwidth}{X} 我們若說自己沒有罪，就是欺騙自己，真理就不在我們裡面了。 \end{tabularx} \\ \\ \relax
1:9 & \begin{tabularx}{0.7\textwidth}{X} 我們若認自己的罪，神是信實的，是公義的，必要赦免我們的罪，洗淨我們一切的不義。 \end{tabularx} \\ \\ \relax
1:10 & \begin{tabularx}{0.7\textwidth}{X} 我們若說自己沒有犯過罪，就是把神當作說謊的，他的道就不在我們裡面了。 \end{tabularx} \\ \\ \relax
2:1 & \begin{tabularx}{0.7\textwidth}{X} 我的孩子們哪，我把這些話寫給你們，是要你們不犯罪。若有人犯罪，在父那裡我們有一位中保，就是那義者耶穌基督。 \end{tabularx} \\ \\ \relax
2:2 & \begin{tabularx}{0.7\textwidth}{X} 他為我們的罪作了贖罪祭，不單是為我們的罪，也是為普天下人的罪。 \end{tabularx} \\ \\ \relax
2:3 & \begin{tabularx}{0.7\textwidth}{X} 我們若遵守神的命令，就知道我們確實認識他。 \end{tabularx} \\ \\ \relax
2:4 & \begin{tabularx}{0.7\textwidth}{X} 人若說「我認識他」，卻不遵守他的命令，就是說謊話的，真理就不在他裡面了。 \end{tabularx} \\ \\ \relax
2:5 & \begin{tabularx}{0.7\textwidth}{X} 凡遵守他的道的，愛神的心確實地在他裡面達到完全了。由此我們知道我們是在他裡面。 \end{tabularx} \\ \\ \relax
2:6 & \begin{tabularx}{0.7\textwidth}{X} 凡說自己住在他裡面的，就該照著他所行的去行。 \end{tabularx} \\ \\ \relax
2:7 & \begin{tabularx}{0.7\textwidth}{X} 親愛的，我寫給你們的不是一條新命令，而是你們從起初所受的舊命令；這舊命令就是你們所聽過的道。 \end{tabularx} \\ \\ \relax
2:8 & \begin{tabularx}{0.7\textwidth}{X} 然而，我寫給你們的是一條新命令，在基督裡是真實的，在你們也是真實的，因為黑暗漸漸消逝，真光已經在照耀。 \end{tabularx} \\ \\ \relax
2:9 & \begin{tabularx}{0.7\textwidth}{X} 人若說自己在光明中，卻恨他的弟兄，他到如今還是在黑暗裡。 \end{tabularx} \\ \\ \relax
2:10 & \begin{tabularx}{0.7\textwidth}{X} 那愛弟兄的，就是住在光明中，他不會使人失足犯罪。 \end{tabularx} \\ \\ \relax
2:11 & \begin{tabularx}{0.7\textwidth}{X} 惟獨那恨弟兄的，是在黑暗裡，也在黑暗裡行走，不知道往哪裡去，因為黑暗使他的眼睛瞎了。 \end{tabularx} \\ \\ \relax
2:12 & \begin{tabularx}{0.7\textwidth}{X} 孩子們哪，我寫信給你們，因為你們的罪藉著基督的名得了赦免。 \end{tabularx} \\ \\ \relax
2:13 & \begin{tabularx}{0.7\textwidth}{X} 父老們啊，我寫信給你們，因為你們認識從起初就有的那一位。青年們哪，我寫信給你們，因為你們勝過了那惡者。 \end{tabularx} \\ \\ \relax
2:14 & \begin{tabularx}{0.7\textwidth}{X} 孩子們哪，我曾寫信給你們，因為你們認識父。父老們啊，我曾寫信給你們，因為你們認識從起初就有的那一位。青年們哪，我曾寫信給你們，因為你們剛強，神的道常存在你們心裡，你們也勝過了那惡者。 \end{tabularx} \\ \\ \relax
2:15 & \begin{tabularx}{0.7\textwidth}{X} 不要愛世界和世界上的東西，若有人愛世界，愛父的心就不在他裡面了。 \end{tabularx} \\ \\ \relax
2:16 & \begin{tabularx}{0.7\textwidth}{X} 因為凡世界上的東西，好比肉體的情慾、眼目的情慾和今生的驕傲，都不是從父來的，而是從世界來的。 \end{tabularx} \\ \\ \relax
2:17 & \begin{tabularx}{0.7\textwidth}{X} 這世界和世上的情慾都要消逝，惟獨那遵行神旨意的人永遠常存。 \end{tabularx} \\ \\ \relax
2:18 & \begin{tabularx}{0.7\textwidth}{X} 孩子們哪，如今是末世的時光了。你們曾聽過那敵基督者要來，現在有好些敵基督者已經出來了；由此我們就知道，如今是末世的時光了。 \end{tabularx} \\ \\ \relax
2:19 & \begin{tabularx}{0.7\textwidth}{X} 他們從我們中間出去，卻不是屬我們的，若是屬我們的，就必仍舊與我們同在。他們出去，這就顯明他們都不是屬我們的。 \end{tabularx} \\ \\ \relax
2:20 & \begin{tabularx}{0.7\textwidth}{X} 你們從那聖者受了恩膏，並且你們大家都知道。 \end{tabularx} \\ \\ \relax
2:21 & \begin{tabularx}{0.7\textwidth}{X} 我寫信給你們，不是因你們不認識真理，而是因你們認識，並且知道一切虛謊都不是從真理出來的。 \end{tabularx} \\ \\ \relax
2:22 & \begin{tabularx}{0.7\textwidth}{X} 誰是說謊話的呢？不就是那不認耶穌為基督的嗎？那不認父與子的，這個人就是敵基督的。 \end{tabularx} \\ \\ \relax
2:23 & \begin{tabularx}{0.7\textwidth}{X} 凡不認子的，就沒有父；宣認子的，連父也有了。 \end{tabularx} \\ \\ \relax
2:24 & \begin{tabularx}{0.7\textwidth}{X} 論到你們，務要將那從起初所聽見的常存在心裡；若將從起初所聽見的存在心裡，你們就會住在子裡面，也會住在父裡面。 \end{tabularx} \\ \\ \relax
2:25 & \begin{tabularx}{0.7\textwidth}{X} 基督所應許我們的就是永生。 \end{tabularx} \\ \\ \relax
2:26 & \begin{tabularx}{0.7\textwidth}{X} 我將這些話寫給你們，是論到那些迷惑你們的人說的。 \end{tabularx} \\ \\ \relax
2:27 & \begin{tabularx}{0.7\textwidth}{X} 至於你們，你們從基督所受的恩膏常存在你們心裡，並不用人教導你們，自有他的恩膏在凡事上教導你們。這恩膏是真的，不是假的，你們要按這恩膏的教導住在他裡面。 \end{tabularx} \\ \\ \relax
2:28 & \begin{tabularx}{0.7\textwidth}{X} 孩子們哪，你們要住在基督裡面。這樣，他若顯現，我們就可以坦然無懼；當他來臨的時候，在他面前不至於慚愧。 \end{tabularx} \\ \\ \relax
2:29 & \begin{tabularx}{0.7\textwidth}{X} 你們若知道他是公義的，就知道凡行公義的人都是他所生的。 \end{tabularx} \\ \\
[1ex]
\hline
\hline
\end{longtable}
今早先和各位思想一個問題,倘若我們各人都向神求一件事,那麼你求什麼呢?首先講我自己吧!是否我要求神,使我作個好講員,或者作個好教師;能將神的話語教導我的學生,幫助他們明白真理,但我也不會這樣求又或我時常求神經常與我同在,與我相交;當我遇見任何問題,我就求祂與我同在,為我解決呢我也不會這樣求因為,神的同在是祂應許我已經講過一章頭四節,永恆之主所給我們的信息:?.就是我們與神相交,並與弟兄彼此相交,祂就必與我們同在0.4節講相交的結果,使我們得到喜樂.
\newline
\newline
繼續,約翰開始講論本書兩個主題 - 神是光,神是愛關於論到神的教義,乃是基督教的中心約翰說「神是光」,因為祂是光;所以要求祂的兒女行在光明中.也就是叫祂的兒女行事為人,若憑著愛,就是行在神的光明中了.神是光與神是愛,兩者是互相和諧吻合的.我們若要與神相交,必需知道祂是那一位.明白神的位格,才能過真正的基督徒生活,約翰首先闡明「神就是光」的意義.
\newline
\newline
(一)要鼓吹的生活倫理(一5)
\newline
\newline
「神就是光,在祂毫無黑暗,這是我們主所聽見,又報給你們的信息.」神就是光,這句簡單的話是很難解釋的.特別那時敵對者也套用這句話,但他們並不真正地明白約翰提出三句話來說明真義:
\newline
\newline
一,這話是象徵,因為文字的有限性,不足以表達抽像的觀念.在中國文學和哲學上,對此更易接受.例如莊子云:「筌者所以得魚,得魚而忘筌；蹄者所以在兔,得而忘蹄；言者所以在意,得意而忘言.」筌者,捕魚之餌.蹄者,誘兔之物也.魚既得,則筌蹄可不再理會.意義既得,則言詞可忘.這段文字結論說:「吾安得忘言之人而與之言哉!」這是一句非常令人嚮往的名言.約翰借用神就是光一語,是要將神的品格解釋出來.試問更有何字,能夠更澈底,更豐富地把神描寫出來.
\newline
\newline
二,是神本質的啟示.參(約一6-9；詩廿七1,卅七6；羅十三12-14.)合併來看,就將光的意思解釋清楚.把光的種種特徵,如公義,純潔,知識......等,表示出來.
\newline
\newline
三,神本性的道德基礎.神如何,祂所造的也要如何,兩者是和諧的.若我們認識神,則必要在生活中表明出來,不可背道而馳.行事為人必須公義,敬虔,聖潔.(彼前一15):「那召你們的既是聖潔,你們在一切所行的事上也要聖潔.」最後,用(林後四6)作本段的結論:「那吩咐光從黑暗裡照出來的神,已經照在我們心裡,叫我們得知神榮耀的光,顯在耶穌基督的面上.」我曾聽過一故事,一人現已作神學院院長,他父親在蘇聯作傳道.一次帶他離家探望朋友.當地的門窗很粗笨,費卻一番工夫,才把所有門窗弄妥.那兒子向來頑劣,此次訪問,父親得到許多關於他劣跡的投訴.歸家時,父親把兒子召到跟前教訓他.室內本懸一十字架,是能在黑暗裡發光的.父親吩咐把燈熄滅,十字架卻不發光.原來因為多關閉門窗,無陽光照射入來這十字架就不能再在黑暗裡發光.許多時候我們也把生命關閉著,不使耶穌福音真光照進心門.
\newline
\newline
(二)要避開的生活錯誤(一6-10)
\newline
\newline
在第六,八,十節,開首同樣用「我們若說」來帶起.約翰提出三個當時教會面對的問題.首先要解釋「我們若......」,就希臘文用法來講,有三種用法.
\newline
\newline
一,是假定那事真發生.
\newline
\newline
二,是純屬假設.
\newline
\newline
三,可真可假,不予確定.
\newline
\newline
舉例來說:譬如座中有人曾打劫銀行,以為我未知,而我想鼓勵他自首,扮作不知.就用第三種語氣:「如果那一位曾打劫銀行,他就錯了.」這是一種很客氣不固定是否的語態.
\newline
\newline
甲,否認罪間斷我們與神相交(6-7)　約翰說:「我們若說是與神相交,卻仍在黑暗中行,就是說謊話的.......」因有人說:縱使我犯罪,但我與神的交通並無間斷.不要受欺!不能過犯罪生活而仍然與神有交通.得罪神,未認罪得赦免,決不能過喜樂平安的生活.不能說,不過犯輕微的罪無關重要,如此就不行在光明中.基督徒必需每一點順服在神面前.
\newline
\newline
我以前曾在加拿大一大學教書,又在一些教會講道.又組成男聲合唱團同往教會唱詩,有一隨軍牧師邀我們往他的教會講道.會後駕車返校已到翌晨.一位同行的男同學希望與我交談,我們就一起回到宿舍.我看出他內心有極大波動,他流著淚對我傾訴:在神學第一年當中,他曾在許多次考試中作弊.但沒有給人發覺,卻給人看為屬靈長進的人；被選為學生會會長.但他卻失去與神交通,沒有內心喜樂.後來,我陪他到校長面前,對校長說明真相.神的賜福再臨到這人,神也特別地使用他.罪必間斷人與神的交通.若繼續地不順服神,則他的生命與生活必然是一片荒漠.當我們不順服神,卻裝作與神有相交的樣子,那是太容易了.他可以仍然祈禱,參加聚會,讀經.但卻不順服神,讓罪停留在他的身上.
\newline
\newline
「我們若在光明中行,如同神在光明中,就彼此相交；祂兒子耶穌的血也洗淨我們一切的罪.」(7)經上的話給我們大勉勵.我們要從內心渴想對神有完全的順服,就可得有和諧的生活.請非常小心,撒旦非常狡猾時常抓住機會試探我們.當我們向神承認自己的罪,就當相信神兒子耶穌的寶血,洗淨我們一切的罪.並不需要作別的事情,好像天主教辦告解立善功,這都不合乎聖經真理.
\newline
\newline
乙,否認罪存在我們心裡(8-9)　「我們若說自己無罪就是自欺,真理不在我們心裡了.」約翰提出第二個當時面臨的錯誤.是以為罪不在人本性中.所以說:「若說自己無罪就是自欺!」因為神說:「世人都犯了罪,虧欠了神的榮耀.」人內心中都有反叛神的傾向.約翰指明,如果不承認有這傾向,就是自欺了.如果叫人誠實地作答:「你愛神嗎?願意順服神嗎?」答案當然是否定的.因為人是自私的,愛好自作主張,這就是人的本相.讓我們用點幻想作個比方,來說明此真理.我們得到一隻好看的小豬,把牠洗濯清潔,又加上打扮裝飾,並且教曉了種種規矩儀注.甚至能與家人一同起居,儼然是我家一個新成員.小豬也安之若素.有一日,門戶無意敞開,小豬瞥見泥潭污水,頓時興高彩烈,衝出門外投身其中,如享至樂.人的內心具有對神反叛的本質,與渴望與神相交,同樣存在.所以,人若說自己無罪,就是自欺,真理不在我們心裡了.第9節:「我們若認自己的罪,神是信實的,是公義的,必要赦免我們的罪,洗淨我們一切的不義.」這裡:「信實」,「公義」,是神的兩個特徵,給我們大安慰.因祂的信實,必照應許實行.因祂是公義的,祂有權赦免我們.我們也可以此擊退撒旦的控告,得赦免乃是藉著神的信實和公義.有一年我在美國大學教書,同時擔任做一個教會的牧師.我在教學時說:基督教非宗教體系,乃是因信與神相交.當時教會有位姊妹與一未信青年相愛,後來這男子與我相約見面,在幾次晤談中,我對他解釋怎樣作主門徒.與他一同禱告,帶領他認罪並接受主耶穌.這男子離去之後,我沉思這是否做得對,因我未將許多基督教信仰教曉他,例如如何加入教會,和作信徒的本份等.但這樣我又把基督教當成行為體繫了.後來女子父親──執事會主席問我,這男子是否真作了基督徒.我知他回到家中時,作了見證,對家人承認悔改接受基督.
\newline
\newline
丙,否認罪在我們行為上表現出來(10)　這是約翰指出的第三個問題:「我們若說自己沒有犯過罪,便是以神為說謊的,祂的道也不在我們心裡了.」當時的人,自誇有特別知識,看所犯的不算為罪.現今世代同樣有這傾向,對罪用種種藉口來掩飾.說是身體的軟弱,心理的問題.並不逕看為罪.神稱為罪便是罪,得罪神自然是罪.伊甸園中始祖墮入試探背逆神,當神來尋找他們時,他們互相推諉.相信神如責備蛇時,蛇亦會委過於神.
\newline
\newline
應該在神面前承認:「我有罪了!」「我唯獨得罪了你!」「求你為我造清潔的心,使我裡面重新有正直的靈.」然後我們才能得蒙神的賜福.
\newline
\newline
你們要在光明中行,如同神在光明中.又需小心避免以上三種錯誤.
\newline
\newline


\newpage

\section{第49屆~港九培靈會~艾理德~第四講}
\label{sec:771}
\textbf{與基督教不能分開的倫理宣告(續上)}
\newline
\newline
連結: \href{https://www.hkbibleconference.org/session-message/view/771}{\texttt{https://www.hkbibleconference.org/session-message/view/771}}
\newline
\newline
\hyperref[sec:770]{< < < PREV SERMON < < <}
~
\hyperref[sec:toc]{[返目錄]}
~
\hyperref[sec:772]{> > > NEXT SERMON > > >}
\newline
\newline
約翰一書 1:5-2:6
\newline
\begin{longtable}{cl}
\hline
\hline
章節 & 經文 (和合本修訂版)\\
\hline
1:5 & \begin{tabularx}{0.7\textwidth}{X} 神就是光，在他毫無黑暗；這是我們從主所聽見，又報給你們的信息。 \end{tabularx} \\ \\ \relax
1:6 & \begin{tabularx}{0.7\textwidth}{X} 我們若說，我們與神有團契，卻仍在黑暗裡行走，就是說謊話，不實行真理了。 \end{tabularx} \\ \\ \relax
1:7 & \begin{tabularx}{0.7\textwidth}{X} 我們若在光明中行走，如同神在光明中，就彼此有團契，他兒子耶穌的血就洗淨我們一切的罪。 \end{tabularx} \\ \\ \relax
1:8 & \begin{tabularx}{0.7\textwidth}{X} 我們若說自己沒有罪，就是欺騙自己，真理就不在我們裡面了。 \end{tabularx} \\ \\ \relax
1:9 & \begin{tabularx}{0.7\textwidth}{X} 我們若認自己的罪，神是信實的，是公義的，必要赦免我們的罪，洗淨我們一切的不義。 \end{tabularx} \\ \\ \relax
1:10 & \begin{tabularx}{0.7\textwidth}{X} 我們若說自己沒有犯過罪，就是把神當作說謊的，他的道就不在我們裡面了。 \end{tabularx} \\ \\ \relax
2:1 & \begin{tabularx}{0.7\textwidth}{X} 我的孩子們哪，我把這些話寫給你們，是要你們不犯罪。若有人犯罪，在父那裡我們有一位中保，就是那義者耶穌基督。 \end{tabularx} \\ \\ \relax
2:2 & \begin{tabularx}{0.7\textwidth}{X} 他為我們的罪作了贖罪祭，不單是為我們的罪，也是為普天下人的罪。 \end{tabularx} \\ \\ \relax
2:3 & \begin{tabularx}{0.7\textwidth}{X} 我們若遵守神的命令，就知道我們確實認識他。 \end{tabularx} \\ \\ \relax
2:4 & \begin{tabularx}{0.7\textwidth}{X} 人若說「我認識他」，卻不遵守他的命令，就是說謊話的，真理就不在他裡面了。 \end{tabularx} \\ \\ \relax
2:5 & \begin{tabularx}{0.7\textwidth}{X} 凡遵守他的道的，愛神的心確實地在他裡面達到完全了。由此我們知道我們是在他裡面。 \end{tabularx} \\ \\ \relax
2:6 & \begin{tabularx}{0.7\textwidth}{X} 凡說自己住在他裡面的，就該照著他所行的去行。 \end{tabularx} \\ \\
[1ex]
\hline
\hline
\end{longtable}
約翰壹書重要訊息 - 與神相交,是每一位信徒所渴望的;約翰也在此盡所能地教導我們如何與神相交第一,二章主要是論到這主題與神相交,必需認識神是怎樣的故此,約翰寫出:「神就是光.」從這真理引申出五點重要教訓上文已經講過兩點:
\newline
\newline
(一)要鼓吹的生活倫理(一5)神是光,是基督徒生活的準則神是聖潔,公義,純潔的,所以祂的兒女也要聖潔;如同在光明中,與神同行.
\newline
\newline
(二)要避開的生活錯誤(一6-10)
\newline
\newline
甲,否認罪間斷我們與神相交當時有假傅教導人說:.「罪不能間斷人與神相交」弟兄姊妹們應記得,不能一面得罪神,一面又與神相交罪必摧毀人與神的交通.
\newline
\newline
乙,否認罪存我們心裡假教師說:「人性中無罪.」(約壹一8)聲明:「我們若說自己無罪,便是自欺,真理不在我們心裡了」我們試從(太四章)記載耶穌受試探的經文很有意思:「當時耶穌被聖靈引到曠野,受魔鬼的試探」我們有沒有對這節經文仔細地想過為什麼魔鬼試探耶穌不在別的地方,卻在曠野經上說:.聖靈將耶穌引到曠野受魔鬼的試探,因為神的靈知道,試探非從身外開始,乃由內心發生耶穌是要解決內心的問題試進一步來問說:耶穌能否犯罪答案是:??「能」何能如此,一位完美的神人也可以犯罪要點是,耶穌所面臨的,非倫理,道德的試探,狡猾的仇敵,乃照耶穌本身的情況來試探祂:「?你是神的兒子」「.你怎樣可以證明是神的兒子,除非把石頭變成餅我才相信」但耶穌並不照他的話來做祂是神的兒子,與父完全合一,祂遵行天父旨意;我們就將注意力集中此點原來我們的內心是悖逆的,不肯順服神丙,否認罪在我們行為上表現出來第10節.:「我們若說自己沒有犯過罪,便是以神為說謊的,祂的道也不在我們心裡了.」不視乎我們自己斷定是否犯罪,神看為「罪」就是犯罪了.
\newline
\newline
(三)要期望的生活鼓勵(二1-2)
\newline
\newline
鼓勵分為兩要點:
\newline
\newline
甲,目的:至此,約翰將提出題目的方法改變.不再用「你們若......」,而改作「我小子們哪!我將這些話寫給你們,是要叫你們犯罪.」朋友阿!神的旨意是要我們不犯罪.我們若犯罪得罪神,生活就變成沮喪.我們雖常有軟弱,但神的目的是要我們不犯罪.如果我們軟弱犯罪,在神那裡為我們有所預備.見於1,2兩節.首先第1節:「若有人犯罪,在父那裡我們有一位中保,就是那義者耶穌基督.」注意「中保」一詞,乃是指一位來鼓勵你的人.多見於法律上,律師與此同一字根.我有過一次上法庭的經驗,當時我正攻讀神學,經歷一次交通意外.美國的習慣,對方往往就著其保險額要求賠償.那次對方的律師竭力想找證據來控告我,但抓不著我的錯失.他給我一封信要求賠償十萬美元.以一個神學生身份,何來這筆鉅款,心裡非常焦急.幸得保險公司派律師來替我辯護,為我解決了問題.照樣,我們若犯了罪,誰能替我們向神申訴呢?有耶穌基督,祂是那義者；祂替我們講話,來鼓勵,安慰我們.
\newline
\newline
乙,不但有中保,更有「挽回祭」(見第2節),就是主耶穌.聖經學者譯挽回祭的意思,現試舉一事來形容:這裡有幾個小孩子是兄弟姊妹,其中一人正在吃糖果；弟妹來到時,他把那些糖果分給他們.其中有嫌少的,不滿就吵鬧起來.爸媽來要他多予些,哥哥再給時,那孩子更故扭不吃,爸媽說把糖果全數予他罷.那後來為挽回他不吵鬧的糖果,就是回祭.這詞被許多世俗宗教套用,所以學者認為如此譯法不宜,但在此無法詳細討論.約翰講出,在耶穌基督裡,使神的怒氣得以止息；我們違反神的命令,但義者耶穌基督,卻為我的罪而把自己獻上.
\newline
\newline
當我們面對自己的罪時,如何使罪得著潔除?只有兩個答案:耶穌為你辯護,為你作挽回祭.基督徒們應得到鼓勵,若你失卻與神交通團契,不要失望,神要把你帶到祂面前,耶穌為你代求,將身體獻上,為你的罪作挽回祭.
\newline
\newline
(四)要遵行的生活榜樣(二3-6)
\newline
\newline
論到基督徒生活榜樣有二要點:都用重要的字來引出,3節:「我們若遵守祂的誡命,就曉得是認識祂.」5節末句:「......從此我們知道我們是在主裡面.」在生活上作榜樣的第一點:遵守祂的誡命.「我們若遵守祂的誡命,就曉得是認識祂.人若說我認識祂,卻不遵守祂的誡命,便是說謊話的,真理也不在他心裡了.」(二3-4)約翰在此最關心的重點是──人若說認識神,但生活卻不符合,就自相矛盾.在座每位相信都已經認識神,但你是否在生活中遵守神的命令,作順命的兒女.「遵守」一詞,在希臘文有儆醒順服之意.
\newline
\newline
我有一個朋友在越南作宣教士,遭遇許多問題.一日,他約同幾位宣教同工懇切祈禱,神也為祂兒女帶來復興.當時越南正在戰火之中,信徒生存岌岌可危.當神的靈對人說話時,就曉得當日信徒家中,竟有偷偷供奉偶像.約翰提出警告,這人口裡說認識神,在暗中卻不順服；就是說謊話的,真理也不在他心裡了.
\newline
\newline
作榜樣的第二點:「凡遵守主道的,愛神的心在他裡面實在是完全的.」(5)在生活上讓神的愛彰顯出來.約翰講出藉著遵守祂的命令,是向神表示愛祂的最好方法.那位越南宣教士,當他們同心向神呼求時,神也對他們講話:一夜,神的愛浸透他的心,他就流淚俯伏在神前.他的妻子也被驚醒,知道是聖靈工作時；就一同流淚認罪,讓聖靈充滿一個預備好的器皿.我過去曾讀偉大的宣教士戴德生的傳記:記載他在一個下午,用整個下午禱告；直到主的愛滿盈在他心中,他甚至躺在地板上,對主說:求你賜給我工作,使用我,不論多少,輕重.神也使用了他傳福音給中國.弟兄姊妹們!要遵守主的命令,來對神表明你愛祂.
\newline
\newline
「人若說他住在主裡面,就該自己照主所行的去行.」(5)是這一段的總結──按照耶穌基督本身的榜樣去行.保羅幫助我們更加明白,(參腓二5)「你們當以基督耶穌的心為心.」保羅在此表現出耶穌的三種態度.(1)祇要是父神的旨意,祂甘於捨棄所有權利.祂本來與父同尊同榮,但卻甘心捨棄神兒子的權利.我們試用想像力來說明此點:父神與聖子同在天上,下望塵世.父對子說,我要差禰到世間去,他們將不喜歡禰,厭棄禰,並且把禰釘在十字架上.但禰要去為他們死.耶穌本可以推辭,留在天上不來為人釘十字架.但耶穌為了遵行父旨,甘心捨棄一切.我們是否願意為神撇捨,雖然是可享的權利,為了神將它捨棄.
\newline
\newline
試再作個比方.許多年前,我在加拿大教書.我是宣道會的會友.本會的制度,宣教士需聽從差會調動,不能自行選擇工作崗位.那時我在鑽研哲學,也很喜愛教師工作.但一天,神的靈感動我,要我作宣教士,於是我掙扎很厲害.我向神辯論:你不是呼召我作教師嗎?學生們豈不是也很歡迎我?但神祇問我是否願意順服?兩個星期以來,不斷地掙扎,心裡覺得沉重.我再次在辦公室內獨自對神,神向有一個方法對付我,祂說:你可按照你的心意而行,但我不會再用你.我終於降服在神面前,答應不教書.
\newline
\newline
約翰所講照著主的榜樣去行.我們不能選擇自己所喜好的,對主說:別的我都依從了,但這是我所應有的,我不能放棄這權利.
\newline
\newline
主耶穌作榜樣第(2)為遵行父的旨意,願將一切阻礙除去.(參腓二7)「反倒虛己,取了奴僕的形像,成為人的樣式.」祂不但來到世間,甚至成為奴僕,我們是否願意付出代價,為著討神的喜悅.(3)「既有人的樣子,就自己卑微,存心順服,以至於死,且死在十字架上」(腓二8)耶穌基督的態度,是不計代價,順服到底,甚至十字架的死亡.
\newline
\newline
(約壹二6)「就該自己照主所行的去行.」讓我作個見證:當我和妻子結婚時,曾在神前立約,如果神對我或她講話時,對方需為他禱告,彼此鼓勵對神順服.我曾在一間大教會中作牧師,一次講道很蒙神的使用.我正鼓勵會眾如有覺得需要,可來到講台前,我要為他們禱告,或對個別問題進行輔導.誰知行出來的有我妻子在內,我感到很尷尬.但我又想到順服神,當神向我們發出祂的旨意,無論怎樣為難,都要對神順服.於是請另一位牧師來結束聚會,我跪在妻旁,同禱告.所以無論我們覺得怎樣困難,總要照著主所行的去行.
\newline
\newline


\newpage

\section{第49屆~港九培靈會~艾理德~第五講}
\label{sec:772}
\textbf{生活的證據}
\newline
\newline
連結: \href{https://www.hkbibleconference.org/session-message/view/772}{\texttt{https://www.hkbibleconference.org/session-message/view/772}}
\newline
\newline
\hyperref[sec:771]{< < < PREV SERMON < < <}
~
\hyperref[sec:toc]{[返目錄]}
~
\hyperref[sec:773]{> > > NEXT SERMON > > >}
\newline
\newline
約翰一書 2:7-2:27
\newline
\begin{longtable}{cl}
\hline
\hline
章節 & 經文 (和合本修訂版)\\
\hline
2:7 & \begin{tabularx}{0.7\textwidth}{X} 親愛的，我寫給你們的不是一條新命令，而是你們從起初所受的舊命令；這舊命令就是你們所聽過的道。 \end{tabularx} \\ \\ \relax
2:8 & \begin{tabularx}{0.7\textwidth}{X} 然而，我寫給你們的是一條新命令，在基督裡是真實的，在你們也是真實的，因為黑暗漸漸消逝，真光已經在照耀。 \end{tabularx} \\ \\ \relax
2:9 & \begin{tabularx}{0.7\textwidth}{X} 人若說自己在光明中，卻恨他的弟兄，他到如今還是在黑暗裡。 \end{tabularx} \\ \\ \relax
2:10 & \begin{tabularx}{0.7\textwidth}{X} 那愛弟兄的，就是住在光明中，他不會使人失足犯罪。 \end{tabularx} \\ \\ \relax
2:11 & \begin{tabularx}{0.7\textwidth}{X} 惟獨那恨弟兄的，是在黑暗裡，也在黑暗裡行走，不知道往哪裡去，因為黑暗使他的眼睛瞎了。 \end{tabularx} \\ \\ \relax
2:12 & \begin{tabularx}{0.7\textwidth}{X} 孩子們哪，我寫信給你們，因為你們的罪藉著基督的名得了赦免。 \end{tabularx} \\ \\ \relax
2:13 & \begin{tabularx}{0.7\textwidth}{X} 父老們啊，我寫信給你們，因為你們認識從起初就有的那一位。青年們哪，我寫信給你們，因為你們勝過了那惡者。 \end{tabularx} \\ \\ \relax
2:14 & \begin{tabularx}{0.7\textwidth}{X} 孩子們哪，我曾寫信給你們，因為你們認識父。父老們啊，我曾寫信給你們，因為你們認識從起初就有的那一位。青年們哪，我曾寫信給你們，因為你們剛強，神的道常存在你們心裡，你們也勝過了那惡者。 \end{tabularx} \\ \\ \relax
2:15 & \begin{tabularx}{0.7\textwidth}{X} 不要愛世界和世界上的東西，若有人愛世界，愛父的心就不在他裡面了。 \end{tabularx} \\ \\ \relax
2:16 & \begin{tabularx}{0.7\textwidth}{X} 因為凡世界上的東西，好比肉體的情慾、眼目的情慾和今生的驕傲，都不是從父來的，而是從世界來的。 \end{tabularx} \\ \\ \relax
2:17 & \begin{tabularx}{0.7\textwidth}{X} 這世界和世上的情慾都要消逝，惟獨那遵行神旨意的人永遠常存。 \end{tabularx} \\ \\ \relax
2:18 & \begin{tabularx}{0.7\textwidth}{X} 孩子們哪，如今是末世的時光了。你們曾聽過那敵基督者要來，現在有好些敵基督者已經出來了；由此我們就知道，如今是末世的時光了。 \end{tabularx} \\ \\ \relax
2:19 & \begin{tabularx}{0.7\textwidth}{X} 他們從我們中間出去，卻不是屬我們的，若是屬我們的，就必仍舊與我們同在。他們出去，這就顯明他們都不是屬我們的。 \end{tabularx} \\ \\ \relax
2:20 & \begin{tabularx}{0.7\textwidth}{X} 你們從那聖者受了恩膏，並且你們大家都知道。 \end{tabularx} \\ \\ \relax
2:21 & \begin{tabularx}{0.7\textwidth}{X} 我寫信給你們，不是因你們不認識真理，而是因你們認識，並且知道一切虛謊都不是從真理出來的。 \end{tabularx} \\ \\ \relax
2:22 & \begin{tabularx}{0.7\textwidth}{X} 誰是說謊話的呢？不就是那不認耶穌為基督的嗎？那不認父與子的，這個人就是敵基督的。 \end{tabularx} \\ \\ \relax
2:23 & \begin{tabularx}{0.7\textwidth}{X} 凡不認子的，就沒有父；宣認子的，連父也有了。 \end{tabularx} \\ \\ \relax
2:24 & \begin{tabularx}{0.7\textwidth}{X} 論到你們，務要將那從起初所聽見的常存在心裡；若將從起初所聽見的存在心裡，你們就會住在子裡面，也會住在父裡面。 \end{tabularx} \\ \\ \relax
2:25 & \begin{tabularx}{0.7\textwidth}{X} 基督所應許我們的就是永生。 \end{tabularx} \\ \\ \relax
2:26 & \begin{tabularx}{0.7\textwidth}{X} 我將這些話寫給你們，是論到那些迷惑你們的人說的。 \end{tabularx} \\ \\ \relax
2:27 & \begin{tabularx}{0.7\textwidth}{X} 至於你們，你們從基督所受的恩膏常存在你們心裡，並不用人教導你們，自有他的恩膏在凡事上教導你們。這恩膏是真的，不是假的，你們要按這恩膏的教導住在他裡面。 \end{tabularx} \\ \\
[1ex]
\hline
\hline
\end{longtable}
(一)要提出的生活證據(二7-27)
\newline
\newline
論到與神相交,上文講到就該自己照主所行的去行按著大綱,在「神就是光」這大題之下,續講出:要提出的生活證據,證明是行在光中本節又分為三點:
\newline
\newline
甲,倫理的問題或試驗(二7-14)
\newline
\newline
(1)新時期的命令 (i)命令的來源:「親愛的弟兄阿!我寫給你們的,不是一條新命令,乃是你們從起初所受的舊命令,這舊命令就是你們所聽見的道.再者,我寫給你們的是一條新命令,在主是真的,在你們也是真的,因為黑暗漸漸過去,真光已經照耀.」(7-8)這命令當人一信主時,就開始領受的.是信仰的中心.這命令曾在舊約時頒賜給人,但當耶穌基督來到,更以新形式表明出來.真光為何?耶穌是真光(參約一9)命令是甚麼?「我賜給你們一條新命令,乃是叫你們彼此相愛,我怎樣愛你們,你們也要怎樣相愛.你們若有彼此相愛的心,眾人因此就認出你們是我的門徒了.」(約十三34-35)這舊命令與新命令究竟指甚麼?這是一個證據證明生活在光明中,如同神在光明中.要彼此相愛如同基督愛我們.愛必需要表達出來的,要愛弟兄像對自己一樣.在美國有一位神的僕人,某次,帶同他的幼子參加聚會,孩子逐漸顯得不耐煩,父親為哄他安靜,就除下手上的錶給他玩,那是一隻帶有記念性的錶；過了一會,孩子突然不慎讓錶墮在三合土地上,他慌張地拾起來,錶被跌停了,他不斷地搖著它使它繼續行下去,因為他怕因此受責備；但父親反抱著他說:「爸爸愛你比愛這錶更多.」孩子覺得希奇,回家對哥哥說:「爸爸愛我比愛錶更多.」
\newline
\newline
弟兄姊妹!神要求你愛弟兄比其他一切更多.尋求每一個機會去愛弟兄姊妹,要遵守這命令使別人認出是主的門徒.也許人對彼此相愛此命令難於了解,讓我們試用另一字──「寬恕」來說明.當你祈禱的時候,若想起弟兄得罪你,應該先與他和好.饒恕人好像我們被基督饒恕.若不寬恕人,也不能被主寬恕.這是非常重要的.不愛弟兄,就不能表達神的愛；就不能證明一個與神相交的生活.大學同學中,每有彼此捉弄的事.我在加拿大讀書時,冬季嚴寒,各層宿舍都有煖氣裝置.我住四樓,三樓一位同學愛好作弄人,從氣喉的孔中,把人引到洞下,又潑水下去,作為戲耍.一次,新來的同學,卻是不易與者,被潑水後大肄咆哮,又跑上來跟他理論.後雖被訓導員勸止,友亦忿忿不平；到夜間在神面前讀經,祈禱時,就失卻與神交通的滋味.再進而成績下降,他知道原因為何.但總難於對他的同學認罪而寬恕別人.經上說若不饒恕人,天父亦不饒恕我們.這就是主的命令──你們要彼此相愛.
\newline
\newline
(ii)命令的範圍:「人若說自己在光明中,卻恨他的弟兄,他到如今還是在黑暗裡.愛弟兄的就是住在光明中,在他並沒有絆跌的緣由.惟獨恨弟兄的是在黑暗裡；且在黑暗裡行,也不知往那裡去,因為黑暗叫他眼睛瞎了.」(9-11)人有許多途徑恨弟兄,經文把這些表達出來.對弟兄發怒,反感,嫉妒.......心中充滿這些,就不被神的愛充滿.與神相交也從而中斷了.
\newline
\newline
記得多年前一次講道完畢,有些弟兄姊妹前來接受輔導.有一位姊妹心中有極大的需要,得到神的赦免,經歷神愛的充滿；她願意作任何神所要她作的.但奇怪神的福卻不臨到她.原來她與另一姊妹不和,心中懷有苦毒.聖靈是十分信實的,若我們行在黑暗裡而非行在光明中,就不能經歷聖靈的工作.若對弟兄懷怒就摧殘了身體的見證.
\newline
\newline
(2)新時期的各階層(二12-14)在此約翰的講論稍作分岔.在講完彼此相愛之後,讓信徒們知道他們並非不信任.他又將基督徒生活分成三個階段:(i)小子們,(ii)少年們,(iii)父老.這分法並不在於年齡的差別,其特徵乃在屬靈方面,分成這三種人.
\newline
\newline
(i)小子們──指屬靈的初階,開始對神委身,得著赦免.「小子們哪!我寫信給你們,因為你們的罪藉著主名得了赦免......小子們哪!我曾寫信給你們,因為你們認識父.」(12,13)讓我再作一次比方.在美國有一位海軍朋友,在某次聚會中,接受了耶穌基督作他個人救主.當晚他就寢時,想到他應該早起讀經祈禱.繼而想他需要一個鬧鐘,幫助他的靈修生活.他就起來跪下把他的需要告訴神,不久,聽到有人叩門；另一位朋友來說,他有多一個鬧鐘,願意送給他.你想我這朋友的驚奇和喜樂是何等的大.但我們都有這經驗,當你將需要告訴神,必會得著神的看顧.
\newline
\newline
(ii)少年人.「少年人哪!我寫信給你們,因為你們勝了那惡者.......少年人哪!我曾寫信給你們,因為你們剛強.神的道常存在你們心裡,你們也勝了那惡者.」(13,14)這是一個爭戰與受苦的階段.除非忍受苦難,否則靈命不能進深.回到剛才那海軍朋友的經過,當神答應他禱告,賜給他一個鬧鐘時,他感到何等興奮!作基督徒真是件奇妙的事；於是他告訴神,承認祂是至高主宰；無論祂吩咐任何事,他都願意去作.但朋友,請小心.神果然聽那朋友的禱告,原來他是在海軍醫藥部門任職的,常常悄悄帶些藥品回家,時日既久積聚甚多.如今神對付這件事,要他把藥品送回,這叫他十分為難；經過許多掙紮,於願意順服神的命令,將兩大袋藥物帶回海軍基地.我們可以想像,經過閘口守衛時,對這些東西如何解釋；正不知如何是好,突然一位重要人物行來,朋友對他講出神在他身上所行的事,終於隨著這重要人物順利通過,把藥品全部歸還.這事很難作到,但我們必需經過這順服的階段,否則不能得到神的賜福.必須甘心下到山谷的深處,才能攀登屬靈的高峰；若非甘心順服,徹底底地降卑,就不能獲得權柄與能力.
\newline
\newline
(iii)父老──經歷過與神深交的階層.(13,14)節重複說:「父老阿!我寫信給你們,因為你們認識那起初原有的.」一個與神有密切交往的人,他一定學會信靠神,與惡者爭戰；過順服主的生活,實行對弟兄彼此相愛.若非如此,就無法對外人證明他是個屬主的人.
\newline
\newline
乙,世界的問題或試驗(二15-17)
\newline
\newline
這段包含三層重要教訓.15節:「不要愛世界.」約翰用「世界」一詞,不是指穿甚麼衣服,或文化的模式,甚至頭髮的長短；約翰是指內心的態度,指一個人不容神作其生命的主宰.約翰用命令的方式來吩咐:「不要愛世界」,希臘文在此用一特別結構,意思是「一路在愛世界」.約翰要他們停止,是這命令的意思.第二點描寫世界的特徵:就像「肉體的情慾」,也是個特別的字眼,意為某種態度或行為方式,注重身體情慾的需要,一切都為滿足其私慾而已.我們的身體是屬於神的.但卻從未將身體交託給神.我們每被試探勝過,就是為此.並不是說:身體,胃口,愛好是邪惡的.但「世界」一詞,是指單為身體,胃口,愛好而生存的人.其次提到「眼目的情慾」.意即只看外表而忘卻內在的價值.記得少時與妹妹跟同媽媽往探望姨媽.大人們談興正濃,我們卻感到饑腸轆轆,忽然瞥見房間架上有一盆生果.於是與妹妹相約行事,她在門口把風,我去偷取一個.出到門口,我正要大嚼一口,卻不料是蠟造的.而且旁邊也有了咬過的痕跡.這正是約翰所要表達的,祗看外表,而忽略內在價值的意思.
\newline
\newline
第三個特徵,「今生的驕傲」.是一種虛假誇張的態度.人每有這種錯誤,虛張聲勢使人誤認其價值更形重要.譬如某人去到加拿大,朋友遇見彼此寒暄.他說:感謝神賜福,他在香港過得很好.有兩間房子,又開餐館.其實真相並非如此,不過僅有一間房子而已.這就是跟從世界之人的另一特徵.
\newline
\newline
以上一段的結論:「這世界,和其上的情慾,都要過去,惟獨遵行神旨意的,是永遠長存」(17)結束上文提到的兩個試驗,測量出是否與神相交.(1)是否彼此相愛.(2)不要愛世界.惟獨遵行神旨意,能得永恆的賜福.
\newline
\newline


\newpage

\section{第49屆~港九培靈會~艾理德~第六講}
\label{sec:773}
\textbf{生活的試驗}
\newline
\newline
連結: \href{https://www.hkbibleconference.org/session-message/view/773}{\texttt{https://www.hkbibleconference.org/session-message/view/773}}
\newline
\newline
\hyperref[sec:772]{< < < PREV SERMON < < <}
~
\hyperref[sec:toc]{[返目錄]}
~
\hyperref[sec:774]{> > > NEXT SERMON > > >}
\newline
\newline
約翰一書 2:18-2:27
\newline
\begin{longtable}{cl}
\hline
\hline
章節 & 經文 (和合本修訂版)\\
\hline
2:18 & \begin{tabularx}{0.7\textwidth}{X} 孩子們哪，如今是末世的時光了。你們曾聽過那敵基督者要來，現在有好些敵基督者已經出來了；由此我們就知道，如今是末世的時光了。 \end{tabularx} \\ \\ \relax
2:19 & \begin{tabularx}{0.7\textwidth}{X} 他們從我們中間出去，卻不是屬我們的，若是屬我們的，就必仍舊與我們同在。他們出去，這就顯明他們都不是屬我們的。 \end{tabularx} \\ \\ \relax
2:20 & \begin{tabularx}{0.7\textwidth}{X} 你們從那聖者受了恩膏，並且你們大家都知道。 \end{tabularx} \\ \\ \relax
2:21 & \begin{tabularx}{0.7\textwidth}{X} 我寫信給你們，不是因你們不認識真理，而是因你們認識，並且知道一切虛謊都不是從真理出來的。 \end{tabularx} \\ \\ \relax
2:22 & \begin{tabularx}{0.7\textwidth}{X} 誰是說謊話的呢？不就是那不認耶穌為基督的嗎？那不認父與子的，這個人就是敵基督的。 \end{tabularx} \\ \\ \relax
2:23 & \begin{tabularx}{0.7\textwidth}{X} 凡不認子的，就沒有父；宣認子的，連父也有了。 \end{tabularx} \\ \\ \relax
2:24 & \begin{tabularx}{0.7\textwidth}{X} 論到你們，務要將那從起初所聽見的常存在心裡；若將從起初所聽見的存在心裡，你們就會住在子裡面，也會住在父裡面。 \end{tabularx} \\ \\ \relax
2:25 & \begin{tabularx}{0.7\textwidth}{X} 基督所應許我們的就是永生。 \end{tabularx} \\ \\ \relax
2:26 & \begin{tabularx}{0.7\textwidth}{X} 我將這些話寫給你們，是論到那些迷惑你們的人說的。 \end{tabularx} \\ \\ \relax
2:27 & \begin{tabularx}{0.7\textwidth}{X} 至於你們，你們從基督所受的恩膏常存在你們心裡，並不用人教導你們，自有他的恩膏在凡事上教導你們。這恩膏是真的，不是假的，你們要按這恩膏的教導住在他裡面。 \end{tabularx} \\ \\
[1ex]
\hline
\hline
\end{longtable}
培靈會是寶貴的日子,神著祂的僕人帶來信息,這在我們是重大的權利；但也向我們發出重要的問題.神既將亮光賜下,我們是否願意按著而行,是否預備好一個敞開的心,對神說:求你鑑察我的內心,看在我裡面有甚麼是你不喜悅的.願我行在光明中,如同神在光明中；神就會在每個生命中,釋放祂的大能力.
\newline
\newline
現在繼續查考約翰壹書,他曾提出警告:要行在光明中,如同神在光明中.為的,他提過五點必需履行的教訓.上回講到第五點:要提出的生活證據.在此條之下,分作三次試驗:甲,倫理的問題或試驗──是否彼此相愛.乙,世界的問題或試驗──跟隨世界的方式還是跟從神的旨意.今日繼續講第三項.
\newline
\newline
丙,神學的問題或試驗(二18-27)
\newline
\newline
(1)要宣告的末世教訓(18-21)「小子們哪!如今是末時了,你們曾聽見說:那敵基督的要來,現在有好些敵基督的出來了.」為使這教訓能被明白,願提供三點簡單事實,好叫我們領會:
\newline
\newline
(i)新約之執筆人無足夠字眼來表明末期.
\newline
\newline
(ii)新約執筆人並無將時間表列明出來.
\newline
\newline
(iii)查考新約,會發現作者著重三個時期: a.現今, b.要來的, c.末後日子.現今是當前生活的那段時間.將要來的是指主耶穌再來.祂第一次帶來新生命；新約作者仰望末世來臨,而對末世也生問題,在末期之中更有末期.讀者需要留意.
\newline
\newline
18節講及末世,約翰提到敵基督的要來作為證據.「敵基督」一詞,新約作者雖有涉及這意思,但直接使用的,祇有約翰.約翰提出耶穌基督要再來,現在就是末時了.這是否錯誤呢?對此我願意提三項原則來說明:
\newline
\newline
(i)約翰在此乃表達出神學教訓,不是寫一時間表.
\newline
\newline
(ii)關於主再來的時間日期,主曾說:我們不能知道.「所以你們要儆醒,因為那日子,那時辰,你們不知道.」(太廿五13)約翰在此所鄭重的,是一神學上的真理.讓我作個比方,當飛機作遠程飛行時,駕駛員必需預備一份航程表.又在航線上定一點稱為 PNR (意即永不回轉之點).超過此點就無論如何直往目的地,不能回頭.約翰的觀點彷彿越過此點,進至末後之時代.因無可能知道準確的時間.
\newline
\newline
(iii)必需預備迎接主耶穌基督再來.(可十三32-37):「但那日子,那時辰,沒有人知道,連天上的使者也不知道,子也不知道,惟有父知道.你們要謹慎,儆醒祈禱,因為你們不曉得那日期幾時來到.這事正如一個人離開本家,寄居外邦；把權柄交給僕人,分派各人當作的工,又吩咐看門的儆醒,所以你們要儆醒,因為你們不知道家主甚麼時候來；或晚上,或半夜,或雞叫,或早晨.恐怕他忽然來到,看見你們睡著了.我對你們說的話,也是對眾人說,要儆醒.」(參羅十三11-14；彼後三8-13)等都有相同的意思.
\newline
\newline
因此,與耶穌基督再來教訓最有關係的是要儆醒.有一日,當主回來時,我們要站立祂審判台前.主要察看我們內心存有的.我們能否對主說:我不知道你這時回來,所以沒有準備,不殷勤工作.主會望著我們憂傷,祂在十字架上替我們死；把工作交給我們,許多世事都已應驗指向耶穌基督的再來.雖然過去了許多年,主再來的應許仍未實行,不可忘記經上的話!「有人以為祂是耽延,其實不是耽延,乃是寬容你們,不願有一人沉淪,乃願人人都悔改.」我們是否把握著機會.小子們!現在是末時了.每時期都有敵基督者出現,應當儆醒準備.多年前,我有一位朋友,善於游背泳,但在受訓練時,卻不好好地照所有規則來做,我曾對他說:你遲早要後悔的.但他不以為意,在一次重要比賽中,全加州許多好手雲集；他參加的是一百公尺背泳比賽,當游畢全程時,他的成績第二.正在為這榮譽高興時；裁判卻宣佈他因轉身犯規而被取銷資格,何等可惜!
\newline
\newline
有一天,我們都要站在主面前.到時,主是否稱讚你說:又忠心,又良善的僕人,可以進來享受我的快樂呢?還是你令主心中充滿憂傷呢?照著神的教訓,活出基督生命並非容易.(路九62):「耶穌說:手扶著犁向後看的,不配進神的國.」在希臘文的語氣中能幫助我們明白.「在這裡有一人,他面對基督徒對神信託委身並不清楚,於是耶穌說:你不配進神的國.」請聽,這是非常認真的事,當你對主說:「我要跟從你.」耶穌說:「不能往後看,若不面對基督徒的委身認真跟從,不配進神的國.」
\newline
\newline
(2)要告誡離道者的教訓(二22-23)「誰是說謊話的呢?不是那不認耶穌為基督的麼?不認父與子的,這就是敵基督的.凡不認子的就沒有父,認子的連父也有了.」對這段不想用太多時間講論,但要提出重要警告；現今世代的人許多喜歡將各種宗教,拉雜混合在一齊.神的兒女也應知道這等教訓.他們將耶穌與基督分開,基督是神的啟示,耶穌卻不同；基督比耶穌更大,耶穌不過是基督之一而已.基督的啟示分別在各宗教中.約翰說:誰是說謊的?就是那否認耶穌是基督的,要把握這教訓的重要性.因為不認子的就沒有父.失去了天父就失去生命.這是被約翰定罪的假教訓.
\newline
\newline
(3)要護衛的使徒教訓(二24-27)此段要分三點講出:
\newline
\newline
(i)教會的重要:「他們從我們中間出來,卻不是屬我們的.若是屬我們的,就必仍與我們同在.他們出去,顯明都不是屬我們的.」(19)約翰說:假教師離開信徒團體,這基督的身體乃真理的聖靈所居.假教師不屬真理,因此離開信徒的團體──基督的身體.當留心這重要的應用,現今有許多人,對教會感到失望；他們發覺教會滿了問題,信徒言行不一,所以離開教會.但我們應當知道教會是聖靈的居所,是神的權柄能力彰顯的所在.教會雖非十全十美,但神仍在其中.耶穌說:陰間的權勢不能勝過教會,有耶穌基督福音的權柄.約翰對這屬靈團體有絕大信心.今日不見神的榮耀彰顯,是因不信任神,輕看祂的權柄.英文有三個字:a.相信,b.不信,c.非信.假設我藏有一萬,對眾人宣告誰是第一位相信的,願意前來接受這一萬元；這個宣告就可能有三種反應: a.不信者,自忖艾牧師是一位宣教士,不能有一萬元,所以他不會來. b.相信者,他信我是一位基督徒,必不會撒謊,所以他前來表示相信我的話. c.非信者,他相信我有一萬元,但不信我會履行諾言,將錢給他,這第三種人就叫做非信者.他知道神的應許,知道神凡事都能,但卻不信祂能為我做甚麼.讀舊約聖經,神一再提醒祂的百姓,祂過去在他們中間作過甚麼.當百姓又遇到新問題時,卻自言自語說:「過去神曾幫助我們,如今祂不幫助了.」因此,神就大大震怒.弟兄姊妹!今天你是否活在失望之中?你以為神向別人有許多偉大應許,但神就不理會你所需要的?這就是非信者,是不能在生活中經驗到神能力的人.
\newline
\newline
(ii)必須住在所領受的主的教訓中(24-26),不可忽略福音的大能力.它足以改變你的生命.所聽見的真理,足能供給你所需要的一切.不要受人的誘惑,要將從起初聽見的道,存在心裡.
\newline
\newline
(iii)聖靈的恩膏(20,27)「你們從那聖者受了恩膏.......」「你們從主所受的恩膏,常存在你們心裡；並不用人教訓你們,自有主的恩膏在凡事上教訓你們.......」單有上文的教訓還不夠,並且需要經驗聖靈的能力.這裡並非指膏立,乃是以抹油來形容聖靈.小心這裡有重要的地方不能誤解的.約翰在此作個對比.他說:「聖靈已經教導你們,就不需要那些假教師的教訓.祇要繼續遵守真理的聖靈的教訓,在生活中行出來.」
\newline
\newline
在結束時,可以將上述教訓應用在我們生活上.我們要再經歷聖靈的恩膏,需將生命交給聖靈管理.須要聖靈幫助我們傳福音,管理我們；不可叫聖靈擔憂,使聖靈的能力受攔阻,不能流露出來.今日神已對我們發聲,聖靈能力已顯明出來,要使信徒得更新,得能力!乾旱的生命,可以成為泉源.你把生命敞開,對主說,不要使我攔阻聖靈在我身上工作,作牧師,教師,宣教士的,都可能攔阻神施恩之工.生命有隱藏的罪,不順服神.對弟兄懷恨.都足以攔阻神聖靈的工作.我常警告自己,也向各位同工勉勵；若我自己不愛人,就不能要求弟兄姊妹愛人.我不謙卑順服,就不能要求別人謙卑順服.作神僕人的,我們都需要神所賜聖靈的恩膏；如果我們疲乏困倦,我們要求神使我們復得喜樂平安.
\newline
\newline
我少時在加拿大鄉間,看見有棵大橡樹,我從其中學到了功課嚴冬風雪後,樹上祇餘零星地幾片葉子；但卻連得很牢,甚至孩子可以拉著打千秋.但到春來生命更新,舊葉子就自然脫落.在生命中是不該有的,神的生命來到就自然推出,(路四18)提到世上那些受苦的,有各樣需要的,被罪惡所擄的.巴不得神興起順服祂的人,滿有聖靈能力,完成祂的託負.
\newline
\newline


\newpage

\section{第49屆~港九培靈會~艾理德~第七講}
\label{sec:774}
\textbf{殷勤去作神的工}
\newline
\newline
連結: \href{https://www.hkbibleconference.org/session-message/view/774}{\texttt{https://www.hkbibleconference.org/session-message/view/774}}
\newline
\newline
\hyperref[sec:773]{< < < PREV SERMON < < <}
~
\hyperref[sec:toc]{[返目錄]}
~
\hyperref[sec:775]{> > > NEXT SERMON > > >}
\newline
\newline
約翰一書 2:28-2:29
\newline
\begin{longtable}{cl}
\hline
\hline
章節 & 經文 (和合本修訂版)\\
\hline
2:28 & \begin{tabularx}{0.7\textwidth}{X} 孩子們哪，你們要住在基督裡面。這樣，他若顯現，我們就可以坦然無懼；當他來臨的時候，在他面前不至於慚愧。 \end{tabularx} \\ \\ \relax
2:29 & \begin{tabularx}{0.7\textwidth}{X} 你們若知道他是公義的，就知道凡行公義的人都是他所生的。 \end{tabularx} \\ \\
[1ex]
\hline
\hline
\end{longtable}
來到(約壹二28-29),約翰將本書第一個主題結束,而引入第二個主題第一個主題,是向神兒女提出一項挑戰:「神就是光,要行在光明中,如同神在光明中」你若說認識神,就必須完成這知識所帶來的責任遵守神的命令,彼此相愛,不貪愛世界;並須記住已到末時了,必須慇勤去作神的工這就是一,二章的教訓到此約翰著重兩點.:
\newline
\newline
(i)在28節:「小子們哪!你們要住在主裡面,這樣,祂若顯現,我們就可以坦然無懼；當祂來的時候,在祂面前也不至於羞愧.」如何在耶穌來的時候,可以坦然無懼,不至羞愧呢?約翰給我們答案:是要住在主裡面.但他沒有進一步解釋這話的意思,不過(約十五1-10)有較詳盡的說明.由1至5節,留意兩個重要的字眼:「枝子若不常在葡萄樹上,自己就不能結果子」(4)「因為離了我,你們就不能作甚麼.」(5)當約翰提出:「神就是光,要行在光明中,如同神在光明中.」在這主題他下了一個結論──要住在主裡面.意即:基督徒生活離了主,就不能作甚麼.原來這功課也不容易學會.基督徒往往嘗試靠自己來活出基督徒生活,結果是引致一而再的失敗.讓我再作個比方:基督徒好像一部車子,需要藉賴裡面的汽油才能發動；當他發覺汽油快將用完時,就要到附近油站添油.但他祇有五元,所以祇能加一點點,再不多時就乾竭了.許多基督徒,每年一次來參加培靈會；以為大大地充滿一次,來支持一年的需要,是不可能的.或到有特別需要時,就去尋求特別聚會來滿足他的需要.或臨到有問題時,才到神面前求祂恩典.這都是未能學到常住在主裡面的秘訣.約翰在總結上文主題,提出第一個要點:除了依靠聖靈能力之外,不能過基督徒生活.弟兄姊妺,你曾否向神如此禱告,求聖靈管理我的生活,思想和私慾.我願將一切交託主,願意實行主的旨意.當你如此祈禱,就是進入主裡面,依靠祂而生活.美國加州出產葡萄很有名,我家鄉附近是釀酒的地方.少時常到葡萄園摘葡萄吃.那些農夫常用利剪修理葡萄,當一棵樹若不結出佳美的果子,那些農夫就曉得從何處下功夫,他會小心不太深入地割下去,使樹之漿汁流出但不會讓樹受創,而使樹的生命循一新方向運行,以達到結果子的要求.親愛的朋友!聖靈要在你身上作這修剪之工.若未曾經過對付,就不能領會這出於神的恩典.我們要得神喜悅,似乎神卻把更多的煩惱帶給我們,為要教我們住在祂裡面,單倚靠主.
\newline
\newline
(ii)「你們若知道祂是公義的,就知道凡行公義之人,都是祂所生的.」(29)另一個重點,不單要住在主裡面,還需要實行.曾見青年人參加夏令會,受到感動將生命獻上委身基督,然後得到喜樂平安；帶著美好的見證,屬天的喜樂,回到家中.過三幾個星期後,所得的漸漸過去,隨之而來的是失望.於是對自己說,神的應許不知能否實現在我身上.以為一次就足以維持長遠的需要.應當將生命交託神,並專心遵行神所吩咐的.關於有所實行,保羅在(加五22-23)對我們有幫助:「聖靈所結的果子,就是仁愛,喜樂,和平,忍耐,恩慈,良善,信實,溫柔,節制,這樣的事,沒有律法禁止.」約翰對本書第一主題作出結論──住在主裡面.靠著實行神的事情住在主裡,總括有三意:第一,認識自己不足,要甘心地對主說:主阿!我若靠自己,就不能過基督徒生活.我也曾在此學過功課.當我在神學院攻讀第三,四年時,常有機會在教會講道見證.神也很賜福,到我回校讀最後一年時,我家鄉的教會請我再講一次道.我就選了一篇自己認為很不錯的講章,誰知講道不久,聽眾懨懨欲睡.我講完坐下,心中十分難過,於是神對我說:離了我,你就不能作甚麼.其次,要求我要與神有好的交通.然後神隱藏的生命帶來力量.再次,要預備好被神差遣.三,表彰基督教的倫理──「神就是愛」三1至五12至此轉入另一主題:「神是愛」.怎樣形容與神相交?表明是由愛的聯繫.
\newline
\newline
(一)神兒女生活中行事的特徵(三1-24)神的兒女,為家庭的成員,每一份子都有他的責任和權利.甲,指望──要珍貴超越我們一切期望的遠景(1-3)特別提出「指望」一詞,乃根據三個事實:(1)「你看父賜給我們是何等的慈愛,使我們得稱為神的兒女.世人所以不認識我們,是因為未曾認識祂.」(1)這指望根深於神.當思想,創造萬有之神,竟揀選受造之人作祂的兒女,經歷祂的大愛,與祂有交通.當想到這一層的時候,得著何等鼓勵.去年暑假某日,我放工回家時,發覺女兒很不開心；我試著用種種方法哄她,也無法使她開腔.歇一會再問她時,她就哭起來,原來她把我放書桌上的錄音機弄壞了,心中懼怕.於是我對女兒說:我是你爸爸,愛你比世上一切物件更重,你有難處時應該儘快對我說出；無論大小事情,我都要為你解決.然後我抱著女兒吻她說:你是家庭一份子,我要從家中把失望懼怕除去.神今日對我們這樣說:「你看父賜給我們是何等的慈愛,使我們得稱為神的兒女.」如果你失望犯罪,盡所能地把它藏過來,因為不順服而失去喜樂.神對你說:你是我的兒女,鼓勵你來到祂面前,用祂的愛來解決問題；帶給你喜樂,用寶血洗淨你一切罪,看你為家裡的一員,這是何等美善.
\newline
\newline
世人不能認識,因為他們未認識神.你是神的兒女,願你接受這莫大的權利.關於這點,還有些人難以接受神愛他這事實.因為他不愛自己,不接受自己.以自己種種缺點為藉口,以為神也像他一般看法.聖經說你要盡心,盡性,盡意,盡力,愛主你的神,並要愛鄰舍如同自己.如果不愛自己,怎能愛鄰舍.為此,我要提六點來解決這難題.
\newline
\newline
(i)將生命之否定特徵加以分析.在自己生命中甚麼是不善的,時常引起我作自我批判的.
\newline
\newline
(ii)消除對失敗之恐懼.今日世人總要求別人成功,其實是最危險思想.憂慮考試失敗,生意不成就......等,都造成很大壓力.其實我從錯處得著的,比成功得著的還多.
\newline
\newline
(iii)將不愉快之回憶消除.每人都有過痛苦經歷,或犯過重大的過錯；未信主前敗壞行為.作基督徒又有失敗跌倒,於是自以為神不再愛我這樣的人了.快快地忘掉神應許:東離西有多遠,祂叫我們的過犯離我們也有多遠.
\newline
\newline
(iv)學習恕己.撒旦每每利用此錯誤思想,使我們不能寬恕自己,而導致失敗.
\newline
\newline
(v)深信神的愛,靠著那加給我力量的主,凡事都能作.
\newline
\newline
(vi)每人在神計劃中有獨特之地位,神要使用我們.我們對自己的本性也不能盡知；但神在我們身上,有獨特的做法.記得我初入神學院時,我個性很畏羞.首需作一次三分鐘講道,是我一次難忘的經驗.用心寫好了講章,站起來費力地講完它,好像經過一段漫長的時間.但誰知教授對我說:你講得太短了,不過卅秒鐘,這樣怎能作一個傳道人.我們不能知自己隱藏的生命,但把自己放神手中；因祂獨特美好之計劃,必需確信這事實.被稱為神的兒女,是何等榮耀的權利.現從第2節講作神兒女的遠景.此處較為複雜,分兩點思想:(1)神自己命定的.(2)神旨意所定.「親愛的弟兄阿,我們現在是神的兒女,將來如何,還未顯明；但我們知道若顯現,我們必要像祂,因為必得見祂的真體.」(2)神的兒女,要像基督.神最終目的為何?也未能盡知成為神的形像的意義.但約翰說:「主若顯現,我們卻要像祂.」以弗所書一章3-4節,論到天上各樣屬靈的福氣.就如神從創立世界以前,使我們在祂面前成為聖潔,無有瑕疵.神呼召神的兒女們,要到達一目的地──(5)「又因愛我們,就按著自己旨意所喜悅的；預定我們,藉著耶穌基督得兒子的名分.」論到神揀選的教義,涉及到神的預定和預知.但今早無充分時間來分析這教義.茲舉出兩點略為說明:(1)教義非靠本身成立,乃有一目的在鼓勵我們.(2)這教義與神為其兒女所定的旨意有關係.神在創立世界之前,為其兒女設立了目的地.又在基督裡揀選我們,定下了我們與神之間的美好關係.最後,第3節:「凡向祂有這指望的,就潔淨自己,像祂潔淨一樣.」凡有權利作兒女的人,也帶來本節的責任──要潔淨自己,像祂潔淨一樣.凡不敬虔生活,都摧毀在神計劃中的榮耀與盼望.若不順服神,甘心讓聖靈管理,就使盼望摧殘.當向神承認自己的不順服,使回復那完美的關係.
\newline
\newline


\newpage

\section{第49屆~港九培靈會~艾理德~第八講}
\label{sec:775}
\textbf{神兒女的權利和義務}
\newline
\newline
連結: \href{https://www.hkbibleconference.org/session-message/view/775}{\texttt{https://www.hkbibleconference.org/session-message/view/775}}
\newline
\newline
\hyperref[sec:774]{< < < PREV SERMON < < <}
~
\hyperref[sec:toc]{[返目錄]}
~
\hyperref[sec:776]{> > > NEXT SERMON > > >}
\newline
\newline
約翰一書 3:1-4:21
\newline
\begin{longtable}{cl}
\hline
\hline
章節 & 經文 (和合本修訂版)\\
\hline
3:1 & \begin{tabularx}{0.7\textwidth}{X} 你們看父賜給我們的是何等的慈愛，讓我們得以稱為神的兒女；我們也真是他的兒女。世人不認識我們的理由，是因他們未曾認識父。 \end{tabularx} \\ \\ \relax
3:2 & \begin{tabularx}{0.7\textwidth}{X} 親愛的，我們現在是神的兒女，將來如何還未顯明。我們所知道的是：基督顯現的時候，我們會像他，因為我們將見到他的本相。 \end{tabularx} \\ \\ \relax
3:3 & \begin{tabularx}{0.7\textwidth}{X} 凡對他有這指望的，就潔淨自己，像他是潔淨的一樣。 \end{tabularx} \\ \\ \relax
3:4 & \begin{tabularx}{0.7\textwidth}{X} 凡犯罪的，就是做違背律法的事；違背律法就是罪。 \end{tabularx} \\ \\ \relax
3:5 & \begin{tabularx}{0.7\textwidth}{X} 你們知道，基督曾顯現是要除掉罪；在他並沒有罪。 \end{tabularx} \\ \\ \relax
3:6 & \begin{tabularx}{0.7\textwidth}{X} 凡住在他裡面的，不犯罪；凡犯罪的，未曾看見他，也未曾認識他。 \end{tabularx} \\ \\ \relax
3:7 & \begin{tabularx}{0.7\textwidth}{X} 孩子們哪，不要讓人迷惑了你們；行義的才是義人，正如基督是義的。 \end{tabularx} \\ \\ \relax
3:8 & \begin{tabularx}{0.7\textwidth}{X} 犯罪的是出於魔鬼，因為魔鬼從起初就犯罪。神的兒子顯現出來，是為了要毀滅魔鬼的作為。 \end{tabularx} \\ \\ \relax
3:9 & \begin{tabularx}{0.7\textwidth}{X} 凡從神生的，不犯罪，因神的道存在他裡面，他也不能犯罪，因為他是由神所生的。 \end{tabularx} \\ \\ \relax
3:10 & \begin{tabularx}{0.7\textwidth}{X} 這就顯明誰是神的兒女，誰是魔鬼的兒女了。凡不行義的，不是出於神，不愛他弟兄的，也是如此。 \end{tabularx} \\ \\ \relax
3:11 & \begin{tabularx}{0.7\textwidth}{X} 我們要彼此相愛。這就是你們從起初所聽到的信息。 \end{tabularx} \\ \\ \relax
3:12 & \begin{tabularx}{0.7\textwidth}{X} 不要像該隱；他是屬那邪惡者，殺了自己的弟弟。為甚麼殺了他呢？因為自己的行為是邪惡的，而弟弟的行為是正直的。 \end{tabularx} \\ \\ \relax
3:13 & \begin{tabularx}{0.7\textwidth}{X} 弟兄們，世人若恨你們，不要驚訝。 \end{tabularx} \\ \\ \relax
3:14 & \begin{tabularx}{0.7\textwidth}{X} 我們知道，我們已經出死入生了，因為我們愛弟兄。沒有愛心的，仍住在死中。 \end{tabularx} \\ \\ \relax
3:15 & \begin{tabularx}{0.7\textwidth}{X} 凡恨自己弟兄的，就是殺人的；你們知道，凡殺人的，沒有永生住在他裡面。 \end{tabularx} \\ \\ \relax
3:16 & \begin{tabularx}{0.7\textwidth}{X} 基督為我們捨命，我們從此就知道何為愛；我們也當為弟兄捨命。 \end{tabularx} \\ \\ \relax
3:17 & \begin{tabularx}{0.7\textwidth}{X} 凡有世上財物的，看見弟兄缺乏，卻關閉了惻隱的心，神的愛怎能住在他裡面呢？ \end{tabularx} \\ \\ \relax
3:18 & \begin{tabularx}{0.7\textwidth}{X} 孩子們哪，我們相愛，不要只在言語或舌頭上，總要以行為和真誠表現出來。 \end{tabularx} \\ \\ \relax
3:19 & \begin{tabularx}{0.7\textwidth}{X} 從這一點，我們會知道，我們是出於真理的，並且我們在神面前可以安心， \end{tabularx} \\ \\ \relax
3:20 & \begin{tabularx}{0.7\textwidth}{X} 即使我們的心責備自己，神比我們的心大，他知道一切。 \end{tabularx} \\ \\ \relax
3:21 & \begin{tabularx}{0.7\textwidth}{X} 親愛的，我們的心若不責備我們，在神面前就可以坦然無懼了。 \end{tabularx} \\ \\ \relax
3:22 & \begin{tabularx}{0.7\textwidth}{X} 我們一切所求的，就從他得著，因為我們遵守他的命令，行他所喜悅的事。 \end{tabularx} \\ \\ \relax
3:23 & \begin{tabularx}{0.7\textwidth}{X} 神的命令就是：我們要信他兒子耶穌基督的名，並且照他所賜給我們的命令彼此相愛。 \end{tabularx} \\ \\ \relax
3:24 & \begin{tabularx}{0.7\textwidth}{X} 遵守神命令的，住在神裡面，而神也住在他裡面。從這一點，我們知道神住在我們裡面，這是由於他所賜給我們的聖靈。 \end{tabularx} \\ \\ \relax
4:1 & \begin{tabularx}{0.7\textwidth}{X} 親愛的，一切的靈不可都信，總要察驗那些靈是否出於神，因為有許多假先知已經來到世上。 \end{tabularx} \\ \\ \relax
4:2 & \begin{tabularx}{0.7\textwidth}{X} 凡宣認耶穌基督是成了肉身而來的靈就是出於神的，由此你們可以認出神的靈來； \end{tabularx} \\ \\ \relax
4:3 & \begin{tabularx}{0.7\textwidth}{X} 凡不宣認耶穌的靈，不是出於神。這是那敵基督者的靈；你們從前聽見他要來，現在他已經在世上了。 \end{tabularx} \\ \\ \relax
4:4 & \begin{tabularx}{0.7\textwidth}{X} 孩子們哪，你們是屬神的，並且勝過了假先知，因為那在你們裡面的比那在世界上的更大。 \end{tabularx} \\ \\ \relax
4:5 & \begin{tabularx}{0.7\textwidth}{X} 他們是屬世界的，所以講論世界的事，而世人也聽從他們。 \end{tabularx} \\ \\ \relax
4:6 & \begin{tabularx}{0.7\textwidth}{X} 我們是屬神的，認識神的就聽從我們；不屬神的就不聽從我們。從此我們可以認出真理的靈和錯謬的靈來。 \end{tabularx} \\ \\ \relax
4:7 & \begin{tabularx}{0.7\textwidth}{X} 親愛的，我們要彼此相愛，因為愛是從神來的。凡有愛的都是由神而生，並且認識神。 \end{tabularx} \\ \\ \relax
4:8 & \begin{tabularx}{0.7\textwidth}{X} 沒有愛的就不認識神，因為神就是愛。 \end{tabularx} \\ \\ \relax
4:9 & \begin{tabularx}{0.7\textwidth}{X} 神差他獨一的兒子到世上來，使我們藉著他得生命；由此，神對我們的愛就顯明了。 \end{tabularx} \\ \\ \relax
4:10 & \begin{tabularx}{0.7\textwidth}{X} 不是我們愛神，而是神愛我們，差他的兒子為我們的罪作了贖罪祭；這就是愛。 \end{tabularx} \\ \\ \relax
4:11 & \begin{tabularx}{0.7\textwidth}{X} 親愛的，既然神這樣愛我們，我們也要彼此相愛。 \end{tabularx} \\ \\ \relax
4:12 & \begin{tabularx}{0.7\textwidth}{X} 從來沒有人見過神，我們若彼此相愛，神就住在我們裡面，他的愛在我們裡面得以完滿了。 \end{tabularx} \\ \\ \relax
4:13 & \begin{tabularx}{0.7\textwidth}{X} 因為神將他的靈賜給我們，由此我們知道我們是住在他裡面，而他也住在我們裡面。 \end{tabularx} \\ \\ \relax
4:14 & \begin{tabularx}{0.7\textwidth}{X} 父差子作世人的救主，這是我們所看見並且作見證的。 \end{tabularx} \\ \\ \relax
4:15 & \begin{tabularx}{0.7\textwidth}{X} 凡宣認耶穌為神兒子的，神就住在他裡面，而他也住在神裡面。 \end{tabularx} \\ \\ \relax
4:16 & \begin{tabularx}{0.7\textwidth}{X} 我們知道並且深信神是愛我們的。神就是愛，住在愛裡面的就是住在神裡面；神也住在他裡面。 \end{tabularx} \\ \\ \relax
4:17 & \begin{tabularx}{0.7\textwidth}{X} 由此，愛在我們裡面得以完滿：我們可以在審判的日子坦然無懼，因為基督如何，我們在這世上也如何。 \end{tabularx} \\ \\ \relax
4:18 & \begin{tabularx}{0.7\textwidth}{X} 在愛裡沒有懼怕；完滿的愛把懼怕驅逐出去，因為懼怕裡含著懲罰，懼怕的人在愛裡尚未得到完滿。 \end{tabularx} \\ \\ \relax
4:19 & \begin{tabularx}{0.7\textwidth}{X} 我們愛，因為神先愛我們。 \end{tabularx} \\ \\ \relax
4:20 & \begin{tabularx}{0.7\textwidth}{X} 人若說「我愛神」，卻恨他的弟兄，就是說謊了；不愛他看得見的弟兄，就不能愛看不見的神。 \end{tabularx} \\ \\ \relax
4:21 & \begin{tabularx}{0.7\textwidth}{X} 愛神的，也要愛弟兄；這是我們從神所受的命令。 \end{tabularx} \\ \\
[1ex]
\hline
\hline
\end{longtable}
今天再繼續與大家分享約翰壹書三,四章,本書第二個主題. - 「神就是愛」這個主題是講作兒女的權利和義務神就是愛,乃是從神是光而來的基督徒要行在光明中,如同神在光明中,是指基督徒的生活行為與神的本性相符合約翰將應有的表現作成摘要. - 「神就是愛」今日要盡可能.將第三章教訓鄭重講出,過去幾天著重實用方面,今天要較深入地轉看神學教訓方面,昨日講第三大段第一點:神兒女生活中行事的特徵分成四項真理說明:
\newline
\newline
(1)神采取主動,呼召人稱他們為兒女(1節).作神兒女的意思,表示得著奇異的權柄.(2節)透露神目的之遠景,將來如何雖還未顯明,但必要像基督想到神兒的將來,就連想到常接觸到三個使人混淆的字眼 - .揀選,預定,預知這三個名詞實際內容常被人誤解按新約用這三個名詞,其上下文有目的在鼓勵信徒並不是神定下了某人得救,某人滅亡,這些字的使用與神兒女的命運有關繼續分析三個名詞的意義.:
\newline
\newline
(1)揀選,按希臘文原意:乃有特別的職任,工作,就選某人去作的意思新約的用法相同,為著某特別的責任,神就揀選.
\newline
\newline
(2)預定,表示神已經決定好,神兒女最後的命運;但並不表示,神已定用何種方法決定其命運.
\newline
\newline
(3)預知,三者以此為最難為了幫助聽眾較易明白,我要把三點結論講出:
\newline
\newline
一,神預早知道,不表示揀選.此字在舊約出現七百次,都沒有和揀選拉上關係.到新約翻作預知.
\newline
\newline
灣預知與特別關係有關聯.譬如一對青年男女準備結婚,對婚姻有一概念,此概念是與日後生活無關,但這特別關係,就是將此概念之知識給予他們.參(弗一3- 5)
\newline
\newline
保羅提到在創立世界以先,在基督裡我們被揀選；按著祂旨意所喜悅的,預定我們藉著耶穌基督得兒子的名分.並且預備好至終的結果,使我們在祂面前成為聖潔,無有瑕疵.神預備好揀選的計劃,叫我們像主基督.又從(羅八15):「你們所受的不是奴僕的心,仍舊害怕；所受的乃是兒子的心,因此我們呼叫阿爸父.」我們若將此節與(弗一4)合起來看.(16節)聖靈與我們的心,同證我們是神的兒女.這是關心神兒女將來的命運.(17節)神兒女承受為後裔之權利何等偉大,但至今未能明白作後裔之意.但當我們面對面見基督時,就能明瞭.作者又說:一切受造之物,一同歎息勞苦,不知結果為何.但不須掛慮,萬事都互相效力,叫愛神的人得益處(28節),就是按祂旨意被召的人(28節)神為祂的兒女定有奇妙的計劃(29節)因為祂預先知道的人,祂預先知道是那些人會成為基督徒.我再講預知,並不根據祂的揀選；而揀選反根據預知.基於神的預知,定規那些人與基督有關係.因此就決定他們的命運如何.(29-30)預先定下傚法祂兒子的模樣......所召來的人,又稱他們為義,又叫他們得榮耀.比方我要買一幢房屋,我決定它必須在加州,我就不買別的地方的屋宇.然後我預早決定在裡面住的人.除了我們的兒女親戚之外,他人不能住我的房屋.故我所選的不是張三或李四等,某些特別的人,乃是決定以其種類來作為條件.(約壹三章)說:我們得稱為神兒女,凡願與耶穌發生關係的,就能進入這奇異的命運中.(彼前一2):「就是照父神的先見被揀選.......」表達出揀選的條件,是根據於父神的先見.但切勿將次序掉錯,若掉轉過來,則誤解為神預知,揀這些不揀那些,預知那一種人會作基督徒.據此,神揀選的條件:
\newline
\newline
(1)藉著聖靈成為聖潔.
\newline
\newline
(2)順服耶穌基督,蒙祂寶血所灑.不過這真理容易被人混亂,所以要用幾句說話,將它概括說出.
\newline
\newline
(i)神的掌管乃按其性格之善良而決定,非按其愛惡而取捨.
\newline
\newline
(ii)神的全權,祂揀選某一班人而非個人.約翰說:我們在基督裡蒙揀選.
\newline
\newline
(iii)神的全權,按其旨意下決定.例如我決定買屋,和定好入住其中的人.神也自行決定祂的條件,若滿足這條件就可入住.
\newline
\newline
(iv)神預先知道,但神不預先決定,人對這些作出的反應.因為神預早知道的緣故,就為祂的兒女預先分別出.如果這方法有益,神就用這方法；神為祂的兒女的安排實在奇妙!
\newline
\newline
乙,問題──要潔淨我們一切的過犯(4-10)
\newline
\newline
提到在許多罪中,求神潔淨我們一切過犯.在這裡須要明白經文的中心思想,首先約翰並非解釋罪惡的性質:如羅馬書五章所闡釋罪的性質.這段乃講出道德行為的來源,犯罪的原因是靠魔鬼生活；若靠神過公義的生活,才能蒙神的喜悅.「凡不行義的,就不屬神.」(10節)小子們哪!不要被人誘惑,行義的才是義人,正如主是義的一樣.」我們的生活行為反映出生命之來源,現在又來到重要的問題,如何解釋6,9,二節?「凡住在祂裡面的,就不犯罪,凡犯罪的,是未曾看見祂,也未曾認識祂.(6)「凡從神生的就不犯罪,因神的道存在祂心裡,他也不能犯罪,因為他是由神生的.」(9節)這裡是否說:基督徒不能犯罪,若犯罪就反映非從神而來?這與(二1-2):「若有人犯罪,在父那裡我們有一位中保,就是那義者耶穌基督.祂為我們的罪作了挽回祭.......」兩者是否矛盾?讓我根據希臘文最簡單地解釋這節聖經,意思是:沒有過基督徒生活,而對罪漠不關心.若人犯罪得罪神,卻漠不關心,顯示非從神而生,如基督徒生活有麻煩,失去喜樂,因他的生命沒有順服.從神生的,一旦犯罪,心裡就應該關切.約翰不是說:信徒不會犯罪；乃是說:信徒必須關心會犯的罪.
\newline
\newline
丙,行動──要實行一切待人接物的最佳境界(11-17)
\newline
\newline
「我們應當彼此相愛,這就是你們從起初聽見的命令.」(11節)寫出一大前提──若彼此相愛,就最像神.(參四7-8):「親愛的弟兄啊!我們應當彼此相愛,因為愛是從神來的.凡有愛心的,都是由神而生,並且認識神.沒有愛心的,就不認識神,因為神就是愛.」再看(12-15):論到弟兄相愛的問題,我們很有可能怨恨弟兄.神十分想根深蒂固在我們心靈中.神的話說:若對弟兄生氣,批評弟兄,不按祂的愛來愛,那就與殺人者無異.這教訓何等嚴厲.愛能使人認是主的門徒,就能得到新人悔改.不愛弟兄的,就不能成為新造的人；如果憎恨,就是對人摧殘了.
\newline
\newline
你有沒有見過那些失去志趣的老年人呢?我曾在老年人中工作也曾見過一些難以相處的人.如果你仁慈心,就能愛別人所不能愛的,就能使接受的人,有反應表現.彼此相愛能在真道上互相建立；不然的話,就變成摧殘基督身上的工作了!如果論斷別人的是非,聖靈的能力就無從彰顯；因這問題的緣故,我們要在生活上採取紀律,立下標準.而對這些標準,使我們不致得罪神,為生活立下標準,免違反神的話,因為我們太知道人的本性.
\newline
\newline
讓我作個比方,假如我有一新屋,屋後有曠地長約百呎,地的盡處乃是懸崖,我覺大有危險,便在距離廿五呎處建一道籬笆,當我唸大學的兒子回家度假時,見這籬笆認為不合,何需耗費多地.於是將它後撤.兒子年齡越大,智識越高,回家後數次遷改,將籬笆推近懸崖,後卒將之拆掉.不久,果有人由此處墜落.那些籬笆圍牆,是我們訂立之標準.雖然因人而異,不盡相同,因為太知道我們的本性.立這些為限止我們到達一個地步.免得漫無終止,就會在生活行為上得罪神.
\newline
\newline
(16-17節):「主為我們捨命,我們從此就知道何為愛,我們也當為弟兄捨命.凡有世上財物,看見弟兄窮乏,卻塞住憐恤的心,愛神的心怎能存在他裡面呢?」神從上頭看下來,見到我們,愛我們；為預備美好的將來.但見到弟兄有需要,卻塞住憐憫的心,不予援手.神的愛就不能藉我們表明出來.這是一個行動需要遵行.
\newline
\newline
丁,應許──要支取應許,以致我們除去一切懼怕(18-24)
\newline
\newline
提到應許時,約翰首先定下一個原則:「小子們哪!我們相愛,不要只在言語和舌頭上,總要在行為和誠實上.」(18節)我建議你安靜下來,列一張表,把所有你所不愛的人,或事一一寫明.然後請你跪下,為這些禱告,我原本不喜歡這樣人和事,但求主將祂的愛賜給我.你會感到驚奇,神能將祂的愛賜下,改變了你.這是非常重要的真理.
\newline
\newline
(18-21)節論到另一問題:「從此就知道我們是屬真理的,並且我們的心在神面前可以安穩.我們的心若責備我們,神比我們的心大,一切事沒有不知道的.」有時我們被良心責備,就會思想到神到底是否與我同在?怎樣來堅固我們的信心?對這話題有兩個解答:(1)靠著彼此相愛來確知神與我們同在,但這也不是真實答案,單憑愛,也不足以確知神的同在.(2)在(20節)中；神比我們的心大,一切事沒有不知道的.我們心中有疑惑,如何能勝過?靠著對神本性的認識,就能得勝我們內心的猶疑.
\newline
\newline
神對你所施的慈愛,大於人對自己的看法.雖然心會內疚自責,但神所看到的更為深入.祂明瞭我們的個性與動機.遠過於自己所能看到的.要倚靠神!讓祂豐富的憐憫,深入地判斷你,過於自己輕率所作下的判斷.
\newline
\newline


\newpage

\section{第49屆~港九培靈會~艾理德~第九講}
\label{sec:776}
\textbf{愛裡沒有懼怕}
\newline
\newline
連結: \href{https://www.hkbibleconference.org/session-message/view/776}{\texttt{https://www.hkbibleconference.org/session-message/view/776}}
\newline
\newline
\hyperref[sec:775]{< < < PREV SERMON < < <}
~
\hyperref[sec:toc]{[返目錄]}
~
\hyperref[sec:736]{> > > NEXT SERMON > > >}
\newline
\newline
約翰一書 3:18-3:24
\newline
\begin{longtable}{cl}
\hline
\hline
章節 & 經文 (和合本修訂版)\\
\hline
3:18 & \begin{tabularx}{0.7\textwidth}{X} 孩子們哪，我們相愛，不要只在言語或舌頭上，總要以行為和真誠表現出來。 \end{tabularx} \\ \\ \relax
3:19 & \begin{tabularx}{0.7\textwidth}{X} 從這一點，我們會知道，我們是出於真理的，並且我們在神面前可以安心， \end{tabularx} \\ \\ \relax
3:20 & \begin{tabularx}{0.7\textwidth}{X} 即使我們的心責備自己，神比我們的心大，他知道一切。 \end{tabularx} \\ \\ \relax
3:21 & \begin{tabularx}{0.7\textwidth}{X} 親愛的，我們的心若不責備我們，在神面前就可以坦然無懼了。 \end{tabularx} \\ \\ \relax
3:22 & \begin{tabularx}{0.7\textwidth}{X} 我們一切所求的，就從他得著，因為我們遵守他的命令，行他所喜悅的事。 \end{tabularx} \\ \\ \relax
3:23 & \begin{tabularx}{0.7\textwidth}{X} 神的命令就是：我們要信他兒子耶穌基督的名，並且照他所賜給我們的命令彼此相愛。 \end{tabularx} \\ \\ \relax
3:24 & \begin{tabularx}{0.7\textwidth}{X} 遵守神命令的，住在神裡面，而神也住在他裡面。從這一點，我們知道神住在我們裡面，這是由於他所賜給我們的聖靈。 \end{tabularx} \\ \\
[1ex]
\hline
\hline
\end{longtable}
在九天培靈會過程中,神使眾人蒙福.神的話語使我們得鼓勵和安慰.由於聽從祂的話,經驗到大喜樂.今天最後一次藉著講道勉勵神的兒女們.再返到(三章18-24節),約翰使讀者從懼怕中支取應許.我們要明白,接受這應許,就可以從中享受自由.
\newline
\newline
經文中藏有三點要訓:
\newline
\newline
(一)一個原則
\newline
\newline
「小子們哪!我們相愛,不要只在言語和舌頭上,總要在行為和誠實上.」(18節)原則就是基督徒需要愛.在這幾天聚會中,也曾重複地講論這原則.而且愛在新約中,也是最重要的主題.如果懂得愛就最像神.(參看四11,12,16節):「親愛的弟兄啊!神既是這樣愛我們,我們也當彼此相愛.從來沒有人見過神,我們彼此相愛,神就住在我們裡面,愛祂的心在我們裡面得以完全了.神愛我們的心,我們也知道也信.神就是愛,住在愛裡面的,就是住在神裡面,神也住在他裡面.」神先將愛藉著約翰指示我們,愛並非外表和言語上,乃誠實地相愛.討論和講道都容易,但實行這愛卻不容易.若實行愛則似神愛人,瞭解人,像神對人.彼此相愛要在行為和誠實上.書信用「誠實」一詞是表屬乎真理.屢見於本書(參一8,二4,三18等節).研究哲學的人,許多時候講到真理.基督徒也應該對此瞭解,簡單來說真理是和諧的理論.讓我們稍為思想,何謂真理?或說:真理就是對的事.但也未能充份回答問題.甚麼是對呢?平常人就將對的與真理作為同義使用.聖經中對此有特別的用法,約翰提到在真理中,表示口講和心中是和諧一致的.假如我說:昨天早上十時我在三藩市.那大家都知道與事實不相符.約翰就在此教導基督徒彼此相愛,而且要在言語和實行上彼此符合.
\newline
\newline
對於三章18節經文,我要用四句話,將實用的意義講出:
\newline
\newline
(1)愛乃基督徒和諧生活的原則.比方我在大學讀書時,為對生活上幫補得找些工作.做點替貨車卸下貨物的粗工,就需要與一些粗魯的工人有接觸.他們滿口粗言穢語,甚至咒詛我們的主；我為此感覺非常煩惱.有一日在祈禱中,神對我講話,要學習認識這樣的人,他們可說是不虔不義之典型.當我能用主的愛來愛他們,實在見到神的恩典,喜樂與和平,一切聖靈的果子運行,實現基督徒生活和諧的原則.
\newline
\newline
(2)愛最像神.愛能知道別人,如同主知道我們一樣.當我們讀(林前十三章)時,我們是否有更深的領會呢?這篇論到愛的經文,也講到知識.「愛」與「知識」有複雜的關係.簡單地說:「愛是永不止息,先知講道之能,終必歸於無有；說方言之能,終必停止；知識也必歸於無有.我們現在所知道的有限,先知所講的也有限,等那完全的來到,這有限的必歸於無有了.我作孩子的時候,話語像孩子,心思像孩子,意念像孩子,既成了人,就把孩子的事丟棄了.我們如今彷彿對著鏡子觀看,模糊不清,到那時就要面對面了.我如今所知道的有限,到那時就全知道,如同主知道我一樣.」(林前十三8-12)保羅所講現今世界的知識不能全知各事,不過局部而已；好像一面不好的鏡子,不能將形像清晰地反照出來.此處用字很有意思,指一特別的鏡──哈哈鏡,當人站在鏡前時,反照出高,瘦,矮,胖等歪曲的形像.保羅用這字來形容世界的知識便是這樣.但他又顯示有方法,得神的知識,如同神知道我們一樣.就是面對面見主時所見的.這乃是從愛而來的知識.去愛別人,就有新鮮的方法明白對方,如主明白他一樣.
\newline
\newline
(3)這原則發出問題,弟兄得罪我,當饒恕他幾次?經常對弟兄展現愛心,愛別人和寬恕別人同樣的難.為此「愛」和「寬恕」應連在一起.彼得問主說:「弟兄得罪我,要饒恕他幾次?」相信彼得自忖,等主將次數說出之後,到此限度就不再寬恕了.但主對他說:「不是七次,而是七十個七次.」就是無限量地饒恕人,寬恕不是等對方作甚麼?而是自動的.不是等待弟兄改善了,才去寬恕他.責任在我們身上,神要我們現在就開始愛對方.所以約翰說:「小子們哪!我們相愛,不要只在言語和舌頭上,總要在行為和誠實上.」
\newline
\newline
(4)第18節之第四點形容──愛的果子乃對神坦然無懼.
\newline
\newline
(二)一個問題
\newline
\newline
「從此,就知道我們是屬真理的,並且我們的心在神面前可以安穩.我們的心若責備我們,神比我們的心大,一切事沒有不知道的.」(19-20)首先需要明白上下文.約翰講及一種經驗,自己的心責備自己.經驗到好像神對自己發怒,不知如何去作.約翰說:「從此知道我們是屬真理,並在神面前可得安穩.」但盡己所能,仍覺有罪.我們都曾有過這樣經驗；自覺得罪神,心中的喜樂被奪去.撒旦乘機來試探,你心都對你說有罪,就必然是有罪了.我記得年青時,常赴培靈會,當講員指責罪時,心中常有罪惡感,好像失掉了心中的安穩.為此,讓我們來看三件事
\newline
\newline
(1)看是否有明知而未認的罪,對過往已經承認的罪勿再理會,勿妄製造事端.
\newline
\newline
(2)是否愛你所有弟兄姊妹,不向你弟兄懷怨.
\newline
\newline
(3)勿再分析你每一感覺,而應用20節的話:「神比你的心大,一切事沒有不知道的.」
\newline
\newline
有一位學者對這節經文另解為:神比我們的心大,神的審判也比我們的心大.這解釋是不對的.這是將神的廣大,作為安慰的調節,若我們感到失卻安穩,但我們又經省察沒有犯任何罪和得罪弟兄；就祇管簡單地交託主,忘記此事.約翰說神滿有憐憫心腸,神愛你比一切更甚；祂也知道你內心有責備.所以應當簡單地忘掉它,把它交在主的足前.這是何等的鼓勵.
\newline
\newline
(三)一個應許
\newline
\newline
當我們在神面前坦然無懼的時候,祂的應許就臨到:「並且我們一切所求的,就從祂得著,因為我們遵守祂的命令,行祂所喜悅的事.」(22)約翰又在第五章再重複這事:「我們若照祂的旨意求甚麼,祂就聽我們,這是我們向祂所存坦然無懼的心.既然知道祂聽我們一切所求的,就知道我們所求於祂無不得著.」(14-15)在此我要將一些重要教訓與各位分享,其次要解答一個問題.首先,我們向神所求的是否都得著?按照應許是應該得著的.如何達到?約翰提到一個方法──我們的意念與神的旨意一致,就必按我們所想所求賜給我們.就如三章22節記載的,即使神的答案是否定的,也是我心所願意接受的答案.因為我要得著神心中的事情.第二,祈禱所帶來之問題,乃是神學上難題之一,神既是掌管一切,何必祈禱?神既完全知道,在我們未祈禱以先,祂已經知道；又何必祈禱呢?經上說:「並且我們一切所求的,就從祂得著.」面對為何要向神祈求的問題,我願以三點與各位分享:
\newline
\newline
(1)神以全能掌管一切,並非在永恆時間上之一剎那,譬如我劃一條線,從一點到另一點,代表世界的開始和結束.然後又在線上放一點,讓它與時間的終始拉上關係.在每一時間上所發生之事,原來在神之點上已經完成了.將世間上所發生之事,總結在神那一點之上.用這方法來明白神,就是錯誤,以致發生許多難題.例如我今早起床,會選擇配哪一條領帶.能否說我用那一條領帶,神也在那一點上完成了呢?其實神關心的非我用哪條領帶的問題.其次,為那些未相信的人祈禱,神若是將一切在時間和空間所發生的事都定好,人的得救沉淪也經有所規定,則又何需祈禱?
\newline
\newline
由此帶出教訓之(2)神的全能乃按神的旨意目的而決定.不能將神限制於其本身定義中,如此則無異將神摧殘了.
\newline
\newline
(3)神之全能乃自由和有動力的.
\newline
\newline
我們若緊握以上三點,就能鼓勵我們對祈禱的意義.讓我們再進一;步對祈禱加以分析:
\newline
\newline
(1)祈禱不應祇作神學上的知識；除非實行,否則不能明白.有些哲學家研究神,看祂是生命和經歷以外的一位；可以用研究數學的方法來認識祂,這是同樣的錯誤.我們不能說認識神,而不向祂祈禱,敬拜.祈禱也一樣,除非實行,否則不能認識.
\newline
\newline
(2)祈禱主要是與神的關係,這關係可在父子的關係中表達出來.我對天上的父祈禱,也將一切需要和關心向祂表達.在禱告之中,可以對天父有所學習；而父的旨意也藉此表明.
\newline
\newline
(3)祈禱在基本上,並不是與神聯繫,彼此互通消息.而基本上帶到另一要點,乃是信心的行動；又是分享神生命的方法.當我藉著祈禱與天父發生關係時,必需有絕對的信心.有時到祂面前本不曉怎樣禱告,就將這心意表達出來；求祂將這禱告帶去,到那永恆完美的旨意之中.我祈禱時,甚至不知父的答案是甚麼?但我會將感覺最好的帶到父面前,對父絕對地相信,絕對地交託.
\newline
\newline
我們又要十分小心,許多時我們祈禱中所講,所作是錯的.當我們祈禱完畢,願禰旨意成就,阿們!是否內心真是如此?抑或另有想法,要我們一切所想望的都行得通,要神照我們一切的思想辦妥.如此,我們不是仰望一位在天上的父,是向著命運,是犯罪的所為.我們的天父關心我們,祂知道如何將最好的賜給我們.
\newline
\newline
祈禱會將一個很重要的分別帶出來.就是將祈求與願望分開.有件事西方孩子很容易學會,小孩子喜歡吃糖和雪糕,我最少的女兒每喜歡坐我膝上,抱我吻我表示親熱.我當然知道她下一步是甚麼?她接著會索食雪糕.我每每藉此教導屬靈真理,父親對你的愛非如你所想的.無論我給予或不給予；我的愛都大過這些.有時我們對主說:我有這大的需要,神不給予我,表示神不愛我.這也是神對祂子民的功課.神賜給食物,雲火柱引導,但這些還未使他們學會秘訣.他們必須知道,人活著不是單靠食物,乃靠神的話語.願神振奮我們,使我們知道神在我們身上有美好計劃,無論祂對你的祈求,賜給或者不賜給.
\newline
\newline


\newpage

\section{第50屆~港九培靈會~楊濬哲~第一講}
\label{sec:736}
\textbf{為甚麼我們要被聖靈充滿}
\newline
\newline
連結: \href{https://www.hkbibleconference.org/session-message/view/736}{\texttt{https://www.hkbibleconference.org/session-message/view/736}}
\newline
\newline
\hyperref[sec:776]{< < < PREV SERMON < < <}
~
\hyperref[sec:toc]{[返目錄]}
~
\hyperref[sec:738]{> > > NEXT SERMON > > >}
\newline
\newline
此次港九境靈研經大會講員所傳的信息,我深深覺得有神的引導；研經會是講怎樣研讀聖經,講道會是論到救贖,和怎樣被聖靈充滿,奮興會是講優美的人生.這個次序很奇妙!當我們明白了聖經,得了救贖之後,就應該追求聖靈的充滿；得著聖靈充滿的人,然後才會有優美的人生.
\newline
\newline
論到被聖靈充滿的信息,是聖經中最重要的.信徒若不被聖靈充滿,靈命則無法長進,教會也無法興旺.我們看教會情形和信徒的靈性,實在需要聖靈來充滿我們.但被聖靈充滿是非常複雜的題目,歷來解經的人,追求的人,看法不同,見解也不同；以至莫衷一是,叫人無所適從.我們真需要聖靈親自來帶領我們.對於這真理,有客觀的認識,又有主觀的分別.客觀方面常偏重被聖靈充滿的道理,多用心思來研究；但忽略了靈性追求的經歷.而主觀方面,剛巧相反；只注意被聖靈充滿的經歷,而不用心思去想.因此客觀方面就變成空洞,在主觀方面則變成不穩和錯誤!求主憐憫我們,使我們能夠雙方並重.
\newline
\newline
為甚麼我們要被聖靈充滿?查考聖經我們看見有三個最大的原因.
\newline
\newline
(一)因為有父神的應許(路二十四49；徒一4-5,二33)
\newline
\newline
主復活後將升天的時候,祂對門徒說:「我將父所應許的降在你們身上.」這幾處聖經都清楚地說明,我們之所以要被聖靈充滿,乃是根據父神的應許.為甚麼呢?因為我們始祖在未犯罪前,與神的交通是毫無阻隔的；但人犯罪後,就斷絕了關係.而且死亡就臨到眾人,因人與神隔絕了!雖然舊約有記載神的靈充滿人,使人可以為神工作.但按原則來說,在舊約時候聖靈還未進入到人裡頭.正如(創六3)說「人既屬乎血氣,我的靈就不永遠住在他們裡面了.」(詩五十一11)大衛知道神會在身上收回祂的聖靈.我們又看見主來世之後,約翰清楚說:「那時還沒有賜下聖靈,因為耶穌還未得著榮耀.」(約七39)舊約時代的人,還未得著聖靈進到他們裡面,永遠居住；聖靈雖臨到人,使用人,使人有超人力量,能爭戰.按原則說舊約時代的人還未得著聖靈進到他們的生命裡面.參孫是個好例子,他雖然常被聖靈感動行事,聖靈卻不在他生命裡面.但神的誠意,乃要人得著聖靈,祂藉以西結說:(結三十六26)要賜給祂的選民新的靈.又說要從他們的肉體中,除去石心,賜他們肉心.先知約珥亦預言說:「以後我要將我的靈澆灌凡有血氣的,你們的兒女要說預言,你們的老年人要作異夢,少年人要見異象.在那些日子,我要將我的靈澆灌我的僕人和使女.」(珥二28-29)可見聖靈不但臨到選民,也臨到每一個信主的人,所以主耶穌復活後,祂就一而再題到父神的應許,彼得在五旬節時,也題到聖靈充滿的奇事.請注意誰能成就神所應許的呢?只有主耶穌,因為祂來世時重大的使命,就是要恢復神所賜下的聖靈；就是人犯罪後所失去的.主耶穌在受難的一週,一直提到聖靈的降臨,究竟主怎樣完成神的應許呢?祂首先要為人受死,才能成就神的應許.神處置罪惡的方法有二:
\newline
\newline
(1)在罪人身上定了罪案,
\newline
\newline
(2)或藉主代替人定了罪案(羅八3)
\newline
\newline
神卻採了第二個方法.照聖經預言應驗(賽五十三6)神使我們眾人的罪都歸到祂身上.保羅也說:「神使那無罪的,替我們成為罪,好叫我們在祂裡面,成為神的義.」(林後五21)主必須受死復活升天,然後才降下聖靈來；才完成神應許的.施洗約翰見證說:主是神的羔羊,又說祂要用聖靈施洗.
\newline
\newline
(二)因為是主耶穌的命令(弗五18)
\newline
\newline
「不要醉酒,酒能使人放蕩,乃要被聖靈充滿.」這節聖經乃是主耶穌藉著保羅所說的,這是主的命令.當時以弗所教會有高貴的長老,保羅並沒有單對他們說,乃是對每一個信徒說.比如聖潔,神吩咐有的以色列人都要聖潔一樣.這節聖經所說:不要醉酒,乃是消極的；要被聖靈充滿,是積極的.這是雙重的命令,同具重要性.醉酒固然是罪,不被聖靈充滿,同樣是罪.在某城裡,有某教會的執事,因為醉酒；被教會開除執事職.葛培理博士說,不被聖靈充滿的長執們,為何你們不開除呢?言下之意,是叫我們知道被聖靈充滿的重要性了.既是主的命令,我們就要遵守,聖經說:「聽命勝於獻祭,順從勝於公羊的脂油.」(撒上十五 22)所以我們知道,不順從,不遵守是很大的罪.主耶穌說:「你們若愛我,就必遵守我的命令.」祂在約翰福音裡面,一再題到愛祂的人,必遵守祂的命令；因此,我們很需要主憐憫我們,叫我們知道被聖靈充滿是主的命令,我們必須遵守.
\newline
\newline
(三)因為這是信徒的權利,我們再留心看(弗五18)
\newline
\newline
就知道被聖靈充滿是我們的權利.這節聖經指「你們」不是指一人,乃是指多數,是聖徒,是蒙恩得救的人.以弗所信徒因為信了耶穌,就受了聖靈為印記,聖靈居住在你心裡,你要被聖靈充滿.初期教會,連管理飯食事務的執事也要被聖靈充滿.五旬節時他們有一百廿人,在馬可樓,他們不分階級,不分性別的聚集祈禱.當五旬節到了的時候,他們就都被聖靈充滿了.所以這是信徒所享有的權利.可惜教會中人,不知這方面的重要性,結果只做成生活上的失敗,貪圖世福的人多,追求聖靈充滿的人少；所以沒有得勝罪惡的能力,更沒有忍受苦難的美德；愛人靈魂的心太缺少,對於失喪的靈魂毫不在乎；教會內紛爭結黨,好像哥林多教會一樣.求主憐憫我們,也叫聖靈充滿我們!
\newline
\newline
一九三二年,全世界都盛行著牛津大學團契運動；許多人被聖靈充滿,生活上有奇妙的改變!其中有一個中年商人,發了大財,他不注重屬靈的事,只顧享受.有一天,他聽見牛津團契的人,到他的城裡面.他所遇見的朋友,都對他談論這件事.他聽見後十分厭煩,後來到一位不信的人家中吃飯.但那家的人也和他題及到牛津團契的事.他就離開那裡,跑到遠方去,與人談生意,那人對他說,我們已加入牛津團契.我要向你認罪,因為我曾經欺騙過你兩次.他聽見後,心中更加難過.後來他又另到一城裡尋求快樂,在那城裡,他遇見了多年相交的女朋友,生活與前大不相同,他問女朋友何故?她答:乃是受了牛津團契的影響,就受感動悔改歸信了主,而且接受了聖靈充滿.這位女朋友知道這中年的商人是向來不誠實的,常常假冒欺騙,所以特別約他晚上去赴一個聚會.到了會場,許多人赴會,很難才找一個座位來坐.當晚有一個少年人站起來作見證說:從前我是如何的敗壞,如今在基督裡面蒙主的拯救,領受了聖靈,得勝靈的充滿,成了主的僕人.這些話就在中年商人的心中,使他感到十分的難堪!心如刀割,聖靈就感動他了使他知道以前的罪惡,假冒為善,要他徹底悔改；成為主的兒女,而得勝靈充滿.我們又再看到保羅在(弗五18)說乃要被聖靈充滿.這話在原文之意,乃是在繼續式,表示在進行中,不是過去已經成全的事.因此這句話可翻作──乃要不斷的在聖靈裡得著的充滿.即主在(約七37)所指的渴,不是一次就夠；電池充電一次就夠,電車是一直不住的充電.使徒們也是不住的被聖靈充滿,好像主以水變酒一樣.我們要一直不住的被充滿,而被主使用.夏甲的一皮袋水是不足用的,她與兒子彼此大哭,幸而天使指示她,看見一口水井.李叔青博士說:在大飢荒的時候,許多人飢餓而死,我們不覺得希奇；但如果有人對著飯,眼睛看不見飯而餓死,那就奇怪了,今日信徒眼睛未開,而得不著聖靈充滿就太奇怪了!(弗一17-18)求主憐憫我們,保羅祈禱的話,也是今日我們所需要祈求的話.「求我們主耶穌基督的神,榮耀的父,將那賜人智慧和啟示的靈,賞給你們,使你們真知道祂!」
\newline
\newline


\newpage

\section{第50屆~港九培靈會~楊濬哲~第二講}
\label{sec:738}
\textbf{甚麼是被聖靈充滿}
\newline
\newline
連結: \href{https://www.hkbibleconference.org/session-message/view/738}{\texttt{https://www.hkbibleconference.org/session-message/view/738}}
\newline
\newline
\hyperref[sec:736]{< < < PREV SERMON < < <}
~
\hyperref[sec:toc]{[返目錄]}
~
\hyperref[sec:739]{> > > NEXT SERMON > > >}
\newline
\newline
昨天散會後,有一對信主的夫婦約我交通,我們花了許多時間來交通聖靈充滿的問題.我實在感覺到聖靈充滿是一個非常重要的問題.這是現在各教會所需要的,尤其是港九教會信徒的需要；所以盼望各位在這數天內,多花點時間來聽這個信息.
\newline
\newline
甚麼是被聖靈充滿呢?從反面來看,充滿聖靈,並非領受聖靈,亦非重生得救,更不是領受更多的聖靈.如果我們不明白就不知道甚麼才是被聖靈充滿.
\newline
\newline
(1)充滿聖靈,不是領受聖靈；當然是領受聖靈之後,然後充滿聖靈.我們對這兩件事不可混亂,首先要知道,誰能接受聖靈呢?常有弟兄姊妹要追求聖靈進到他的心,有些甚至禱告懇切,極其悲傷掙扎,非求得勝靈居在心中不可.有一位姊妹說,她的丈夫已經得著聖靈了；其實她的丈夫老早就信了耶穌,但她以為信了耶穌還不夠,還要追求領受聖靈.這觀點是錯誤的,因為除了不信的世人不肯接受之外；所有信主的人都是已經接受聖靈的了.請看(約十四17)「就是真理的聖靈,乃世人不能接受的；因為不見祂,也不認識祂；你們卻認識祂,因祂常與你們同在,也要在你們裡面.」請留意頭兩句話,世界的人未信主的人,是不能接受聖靈的.(羅八9)「如果神的靈住在你們心裡,你們就不屬肉體,乃屬聖靈了；人若沒有基督的靈,就不是屬基督的.」這節聖經已很清楚的說明了信主時,就可得著聖靈.又(弗二22)「你們也靠祂同被建造,成為神藉著聖靈居住的所在.」保羅在(林前三16,六19)一連兩次講到,我們的身子是聖靈的殿,聖靈要居住在其中.不過領受聖靈是一件事,被聖靈充滿,又是另外一件事,我們要分別清楚.以弗所的信徒,他們因聽見真理的道,就接受聖靈,聖靈也在他們身上印了印記,但他們還要更進一步,被聖靈充滿.(弗五18)以弗所教會的情形,就給我們有更清楚的認識.領受聖靈之後,又要被聖靈充滿.領受聖靈是在得救的時候,充滿聖靈是在奉獻的時候,是有所不同的.領受聖靈是聖靈進入來,充滿聖靈就是我出去；領受聖靈,聖靈就恆久永遠居住我裡面.(約十四17)我小時看見過佈道用的掛畫明心圖.有一張印有人的心,心裡有鴿子；另一幅是心中有蛇,有狐狸和其他的東西；因此,心內的鴿子就飛去了.這一幅圖畫表明聖靈離開了,我覺得這不合真理；因為聖靈既然來到我們心中,就永不離開.但充滿聖靈呢?可能是有變遷的,當我們犯罪,得罪了主或得罪了人,聖靈就不能充滿我們.所以在以弗所書,和帖撒羅尼迦書,題醒我們不要叫聖靈擔憂,不可消滅聖靈的感動!這是我們追求聖靈充滿的人要注意的.
\newline
\newline
領受聖靈的人,仍舊有肉體的；正如哥林多的教會,他們雖然領受了聖靈,他們仍舊有嫉妒紛爭等等屬肉體的事情.但充滿聖靈,乃是叫我們成為屬靈的基督徒.總括來說,領受聖靈就是把我們永遠的靈魂,交托給神；充滿聖靈是把我們短暫的生命,和自己交托給神.親愛的弟兄姊妹!我們既然已把永遠的靈魂交托給神,但有沒有把自己來交托神呢?假若有人說,我有一萬元存在銀行裡,而另外他有一百元,你叫他把一百元存放銀行；但他說銀行靠不住；這豈非笑話嗎?那麼為甚麼我們不把自己交托給神呢?
\newline
\newline
(2)充滿聖靈不是重生得救.換言之,重生得救的人,未必是被聖靈充滿.人必須先有重生得救,然後才能被聖靈充滿.一個罪人必須受聖靈的責備,又被聖靈感動接受了主耶穌,他就可以重生得救.有神的生命,成為神的兒女了；於是聖靈就進入他的心裡面,使他能夠呼叫阿爸父.這人雖然重生得救,但未必是被聖靈充滿,未能真正活出基督的生命.當主復活後,來到門徒那裡,主向他們吹了一口氣,他們都受了聖靈.(這裡須加解釋,並非說主復活後,使徒才重生得救)他們重生得救之時,應該是在馬太十六章,彼得代表眾人承認耶穌是神的兒子之時.但他們仍未有聖靈,因為那時還沒有賜下聖靈來；直等到主復活後向對他們吹氣,再等到五旬節到了,聖靈就降下來了.(徒八章)講撒瑪利亞的信徒,已聽見腓利傳講基督就信了,他們必定受了聖靈；但必須等到使徒按手在他們頭上,聖靈才降在他們身上；他們就經歷到聖靈的充滿了.再看保羅在(徒九章)他在大馬色的路上,天上忽然發光,照耀著他.他就仆倒在地上,聽見主的聲音,又與主講話.他說:「主呀,你是誰?」主答:「我就是你所逼迫的耶穌.」從保羅與主對答的話,證明保羅已經接受了主.(林前十二2)清楚說,如果沒有聖靈感動,就沒有人能稱耶穌為主的.主又說:「我就是你所逼迫的耶穌.」「耶穌」之名,照天使的解釋,乃要將自己的百姓從罪惡裡救出來之意.這可以證明保羅已經重生得救了.但保羅還未被聖靈充滿.直等到(徒九17)亞拿尼亞為他按手,他才被聖靈充滿.按神的心意,實在願意每個基督徒都被聖靈充滿.但我們看今日的教會的信徒,正如羅馬書七章保羅所經歷的一樣.肢體中另有個律,與心中的律交戰；所過的生活,是波浪式的.保羅說:「我真是苦啊!誰能救我脫離這取死的身體呢?」我們若在重生後,不追求聖靈的充滿,就難免過著流浪飄蕩的生活.如果遇著試探,就會跌倒屈服了.以色列人本來以十一天的路程可以進入迦南地,但他們竟然在曠野飄流了四十年.今日教會的信徒,也有許多正落在以色列人的光景裡面的.求主真要憐憫我們!
\newline
\newline
聖靈充滿,不是得更多的聖靈,教會有不少的人有這樣錯誤的觀念.有位牧師見證說,他初信主時,渴想聖靈充滿；好像駱駝入帳幕,是逐漸佔據的.有個亞拉伯人有一匹駱駝,他帶了小帳幕旅行.晚上開了帳幕歇宿,駱駝捆在外面.夜裡天氣寒冷,駱駝就請求主人,准許牠的頭進去.主人可憐牠就准了,誰知駱駝的頭入了之後,牠就連身軀也擠了進去.牧師說被聖靈充滿也正如駱駝入帳幕一樣,這是錯誤的!三位一體的神,父子聖靈,我們是不能分開的,更不能分割來領受.聖靈充滿好像一座機器,是由大電力廠發出,電力已充沛在機器裡面,只要一按鈕,整部機器就會發動.我們整個人也像海綿一樣,入水就脹；所以我們知道聖靈充滿是要整個得著我們.
\newline
\newline
從正面看聖靈充滿是甚麼意思?
\newline
\newline
(一)是聖靈在我們身上得著完全的地位,被聖靈充滿,原文是「普利羅」,用一件東西,來充滿另一件東西.請參閱兩處經文(徒十三52)「門徒滿心喜樂,又被聖靈充滿.」門徒在彼西底雖受逼迫,但內心充滿喜樂,因被聖靈充滿.(弗五18)「不要醉酒,乃要被聖靈充滿.」兩者意思是相同的.按以弗所書五章所載,一個被聖靈充滿的人,他的生活是何等的愉快,何等的有意義.被聖靈充滿二字,按聖經原文是用「普利麗斯」,是形容滿有聖靈充滿之意.(路四1)原文是耶穌滿有聖靈,因此祂能抵禦試探.
\newline
\newline
(二)聖靈充滿在我們的身上,來彰顯祂完全的能力,祂不但充滿我們裡面,更充滿到我們外面.(徒二4)門徒被聖靈充滿,裡外都有力量；內裡是使我們過屬靈的生活,外面是使我們工作有能力.(徒一8)充滿聖靈,我們就可得能力為主工作,更有膽量為主作見證.使徒不怕權勢,逼迫,大有膽量；今日教會正需要這外面的力量.
\newline
\newline
有一位弟兄在一個海島上面作工,原來是沒有膽量的人.那海島有些很有權勢的人,他們不要神,拒絕福音.有一日他們拉了這弟兄去,加一罪名,把他囚禁處死.但他毫不畏懼,而且有膽量反駁他們.後來他們釋放了他,使他感到莫明其妙!這實在是聖靈充滿的力量.
\newline
\newline
(三)聖靈在我們身上實施完全的統治.(弗五18-六9)該段經文說明一個被聖靈充滿的人,聖靈在他身上,要實施完全的統治.
\newline
\newline
(1)能讚美主.
\newline
\newline
(2)能感謝父神.
\newline
\newline
(3)存敬畏基督的心,彼此順服.
\newline
\newline
(4)過和諧生活,妻子順服丈夫,丈夫愛妻子.
\newline
\newline
(5)兒女孝順父母,父母不惹兒女的氣.
\newline
\newline
(6)僕人能用誠實的心服事主人,主人能用公平的心待僕人.
\newline
\newline
在這一段聖經可見被聖靈充滿的人,才可過和諧的生活.夫妻間的和諧生活是最難的.一對夫婦信主後,常不和諧,一次相吵,夫以聖經教訓妻,要妻順服；妻即教訓丈夫要愛妻子.他們生活甚不和諧,除非被聖靈充滿,被祂管治,生活才可和諧.
\newline
\newline


\newpage

\section{第50屆~港九培靈會~楊濬哲~第三講}
\label{sec:739}
\textbf{被聖靈充滿的途徑}
\newline
\newline
連結: \href{https://www.hkbibleconference.org/session-message/view/739}{\texttt{https://www.hkbibleconference.org/session-message/view/739}}
\newline
\newline
\hyperref[sec:738]{< < < PREV SERMON < < <}
~
\hyperref[sec:toc]{[返目錄]}
~
\hyperref[sec:740]{> > > NEXT SERMON > > >}
\newline
\newline
我們知道每一個屬主的人,都必須要被聖靈充滿,但怎樣纔能被聖靈充滿呢?我們應該研究思想這問題.兄弟今日就根據查考聖經所得著的,來和大家分享.
\newline
\newline
(一)必須將器皿倒空(王下四1-7)這段經文論到有一個先知門徒的妻子,因為她丈夫死了,欠下了許多債,現在債主要來討債.她無力償還,債主就要取她的兒子做奴僕；於是她來哀求先知以利沙.以利沙對她說:你家裡有甚麼?我能夠為你作甚麼呢?她說婢女家中除了一瓶油之外,沒有甚麼了.以利沙就叫她向鄰舍借空器皿,她與兒子在家中倒油,一直至沒有器皿,油就止住了.她就藉著賣油還了債,而且靠剩餘的可以度日.油是預表聖靈,女人借器皿是最緊要的.器皿不空,油就無法倒進去.(弗五18)保羅說:「不要醉酒,酒能使人放蕩,乃要被聖靈充滿.」他講論到被聖靈充滿,先在(弗四1-17-五17)是論到我們信徒的心思意念,言語行為,應該有新的對付,然後就題到要被聖靈充滿.主耶穌在世時,常教訓我們,必須除去一切的攔阻,被聖靈所充滿.多年前,有一位著名的牧師,到一個高級的官員家中講道,那官員本已信主.牧師的信息是論到聖靈充滿的問題.他說,我們信必須要被聖靈充滿,但我們先要丟掉一切的攔阻；我們不但要向神謙卑,也要向人謙卑.講道後,這位高級的官員就與牧師交通,講出他的虧欠來.他說,我肯在神的面前謙卑；卻不肯在人的面前謙卑,因我得罪了家中的工人.我已向神認罪,但不肯向工人認罪.如今請牧師同我一齊見我的工人,我要向他認罪,求主赦免.這是我們每一個追求聖靈充滿的人,必須要對付的.
\newline
\newline
(二)必須心裡渴慕(約七37)「人若渴了,可以到我這裡來喝,信我的人,就如經上所說,從他腹中,要流出活水的河來.」從主所講的這番話,就知道心中渴慕要得勝靈充滿的人,然後可以得著聖靈的充滿.不渴慕,就斷不能得著；在八福中,要飢渴慕義,纔能得飽足.(詩一○七9)「因祂使心裡渴慕的人,得以知足；使心裡飢餓的人,得飽美物.」(賽四十四3)「因為我要將水澆灌口渴的人,將河澆灌乾旱之地；我要將我的靈,澆灌你的後裔；將我的福,澆灌你的園子.」口渴表示有渴慕和需要,如鹿渴慕溪水一樣.舉世聞名的慕迪先生,自小家貧,沒有受過高深的教育,長大後教主日學；以後撇下了一切跟隨主,奉獻作傳道,成為偉大的佈道家.當他初次出來傳道的時候,講完道,許多人與他握手稱讚他講得好.但有兩位姊妹卻對他說,我們為你祈禱.慕迪問,你們為我祈禱,要為我求甚麼?她們說,要求主給你能力,慕迪不高興,懷疑自己講道是否無力量.但後來他詳細思想,這兩位姊妹所講的,其實是指要被聖靈充滿,然後纔有力量事奉主!於是他從那時起,他就有渴慕追求的心.過了四十天,他經過紐約往歐洲佈道時；就大大的被聖靈充滿,不但在歐洲工作大有果效,結了許多果子；後來回到紐約,和美國,他的工作就大不相同了.因此我們每一個人,必須有飢渴的心,來追求聖靈的充滿.請注意主所講的話.「人若渴了,可以到我這裡來喝；信我的人,就如經上所說,從他腹中,要流出活水的江河來.」人若渴了,可以到我這裡來喝；信我的人,就如經上所說,從他腹中,要流出活水的江河來.」這節經文有三個字要注意的.「渴」,「喝」,「流」,先要飢渴,然後飲下,纔可流出活水的江河來.甚麼江河呢?就是祈禱的江河,讀經的江河,聚會的江河,見證的江河,信心的江河,愛心的江河,熱心的江河,以及一切的江河.這樣不但可以叫自己蒙福,更是可以叫人蒙福.
\newline
\newline
(三)必須有懇切祈禱的心(詩八十一10)「我是耶和華你的神,曾把你從埃及地領上來,你要大大的張口,我就給你充滿.」請留意下面兩句話,你要大大張口,我就給你充滿.「張口」何意?當然是想吃東西,或是想講說話.照樣我們在屬靈方面亦然,我們若大大張口,主就將祂天上的嗎哪,賜給我們吃.我們大大張口,是要讚美感謝神,也是向祂祈禱,來支取聖靈的充滿,主就藉著我們的禱告,要充滿我們.這是主藉詩人的話來激勵我們.在使徒行傳記載,門徒一百二十人在馬可樓同心禱告,聖靈就充滿了他們,使他們滿有能力的,為主作工.到了第四,使徒為了主的道被拉去,被禁止傳福音；他們卻毫不畏懼,後來回到耶路撒冷,眾人同心禱告,迫切到地大震動,然後他們被聖靈充滿.如果不祈禱,就不能被聖靈充滿.甚至連主耶穌,我恭敬說,主被聖靈充滿,亦因為祂迫切的禱告.(路三21-22)「眾百姓都受了洗,耶穌也受了洗,正禱告的時候,天就開了；聖靈降臨在祂身上,形狀彷彿鴿子.......」主禱告後,聖靈就降在祂身上,可見禱告是何等重要.
\newline
\newline
我們讀到查理士芬尼的傳記,他蒙主大用,成了大奮興家,這不是偶然的,他有一次感到靈性上枯乾可憐,他就離開城市,跑到郊外森林裡面,在神面前迫切祈禱,一共廿四小時,因著他的禱告,他就得著聖靈的充滿,大大復興了教會.
\newline
\newline
我們又聽見美國叨雷博士,他有一次感到自己過份誇張,成了無能力的熱心,他就停止所有的工作；他在神面前說,神阿,求你可憐我!因著他迫切祈禱,神的靈就臨到他身上,使他得著力量；成為美國著名的佈道家,他與慕dik 同期.祈禱和等候是有密切關係.主在復活後升天前,一連兩次囑咐門徒,要在城裡等候,領受從上頭來的能力.主的門徒祈禱等了十日,聖靈就降下來充滿他們.他們在等候時,必定祈禱,省察自己,尤其彼得三次否認主.雅各和約翰也省察自己,何等自私,愛慕世界和高位名譽；他們必定想到他們的衝勁,求火燒消撒瑪利亞人,多馬也省察自己,想到自己的虧欠,信心的不足；他必須在主面前認罪.十一個門徒,他們都在主面前省察和等候；求主赦免,然後聖靈纔可以充滿他們.
\newline
\newline
(四)必須殷勤讀經,因為聖經是神的話語,是聖靈流通的管子.聖靈的恩賜和能力,要藉著神的話湧進信徒的心裡面,叫我們變成活水的江河.(弗五 18-20；西三16-17)這兩處經文有類同的地方,歌羅西書說,要把基督的道理,豐豐富富的藏在心裡；而以弗所書就說,要被聖靈充滿.可見聖靈充滿與基督的話語,有密切關係；而且要基督的道理存在心裡,然後可以得著聖靈充滿.慕建一生被主使用,因為他充滿了神的話語.倘若我們不先充滿神的話,是沒有可能會被主使用的.我們必須先在神的話語上好好的扎根.
\newline
\newline
(五)必須憑著信心,我們所走的路程,由得救,教義,以至成聖,都是一條信心的道路,在聖經中可清楚看見.照樣,被聖靈充滿,也必須憑著信心；正如主說,信我的人,要從祂腹中流出活水的江河來.所以不必掙扎,奮鬥,只要憑著信心來求就夠了.(來十一1)「信就是所望之事的實底,是未見之事的確據.」信心就是憑據,我們若肯信而接受主的聖靈；主的聖靈,就充滿我們.英國一位主所重用的僕人梅爾博士,他有這樣的經歷,他覺得沒有被聖靈充滿,心中很枯乾,他就到一個安靜的地方,向主迫切祈禱,後來他想到,主在聖經內曾應許說,信就可以得著(加三14)「使我們因信得著所應許的聖靈.」這節聖經忽然臨到他,他就歡喜到跳起來!憑信心接受了,梅爾博士後來講道,就大有能力,使許多人得著復興.
\newline
\newline
(六)必須謙卑順從(徒五32)「我們為這事作見證,神賜給順從之人的聖靈,也為這事作見證.」彼得得著聖靈充滿之後,他就有這經歷,他想到自己從前的背逆,不聽主的話,在五旬節之前,他就順從主,這樣就得著聖靈的充滿了.請留意,有一件很奇妙的事,就是在舊約聖經裡面,從來沒有題到我們要順從聖靈,但在新約聖經,卻多次題到要我們順從聖靈.特別在加拉太書五章六章有題到,而(林後三17)說:「主就是那靈,主的靈在那裡,那裡就得以自由.」最令我們感到重要的,就是我們必須順從聖靈；即在聖靈管制之下,然後可得勝靈的充滿.在聖經裡面有兩個人是值得我們效法的；一位是主耶穌,主是何等謙卑順服,然後被聖靈充滿.馬太三章載主未被聖靈充滿之先,祂已謙卑順服；尊貴的神子,竟要約翰為祂施洗.祂要盡諸般之義,然後祂就被聖靈充滿.第二位是保羅(徒九 17)他謙卑接受了一個無名小卒,亞拿尼亞為洗禮和按手.他雖是迦瑪列門下的弟子,尚且謙卑順服,所以被聖靈充滿.我們中國一位得主重用的僕人,他講述見證,他讀了梅爾博士寫的書,心中很受感動.一次他見到圖書館中有馬可福音的講義,他私下拿了來要預備講道,聖靈光照他要向校長認錯,但書已遺失了,只好賠償；肯謙卑順服聖靈,聖靈就充滿他.
\newline
\newline
(七)必須完全奉獻,(路一15)「他在主面前將要為大,淡酒濃酒都不喝,從母腹中就被聖靈充滿了.」這節聖經是指著施洗約翰說的,他分別為聖,好像撒母耳一樣,一生作拿細耳人,完全奉獻,跟從主事奉主.他少年時就離家事奉神.他原可以承繼父親的職份作祭司.但他犧牲了前途,到曠野去,身穿駱駝毛的衣服,腰束皮帶,吃的是蝗蟲野蜜；到處宣傳天國的福音,勸人悔改.有一本書描寫他少年時的光景,他父母已離世,他在這情形下,就奉獻一生事奉主.他親友雖勸他,他說我已奉獻不能收回,一直忠心到底,為主殉道.我們為甚麼不能被聖靈充滿呢?恐怕是不肯完全奉獻吧!
\newline
\newline
葛培理博士講到一件很有趣的事,他說有一個青年人,他想把他的房子出租,後來有一個人要租他的房子.青年人說我要留一個房間自用,原來他要留這個房間,是用來養他的大寵物──「老虎」；那人然不租了.我們今日也許像這青年一樣,不是把房子完全歸主用呢!
\newline
\newline


\newpage

\section{第50屆~港九培靈會~楊濬哲~第四講}
\label{sec:740}
\textbf{被聖靈充滿的果效}
\newline
\newline
連結: \href{https://www.hkbibleconference.org/session-message/view/740}{\texttt{https://www.hkbibleconference.org/session-message/view/740}}
\newline
\newline
\hyperref[sec:739]{< < < PREV SERMON < < <}
~
\hyperref[sec:toc]{[返目錄]}
~
\hyperref[sec:731]{> > > NEXT SERMON > > >}
\newline
\newline
兄弟這幾天在神面前迫切的禱告,今日是培靈最後一天,也是講道會最後一堂了.求主藉著這次的信息,打動我們的心,叫我們知道被聖靈充滿,這是非常重要的事.現在我要繼續講到聖靈充滿的果效是甚麼?
\newline
\newline
(一)驅除罪惡,一個人被聖靈充滿之後,就能夠驅除罪惡.在(徒二章)記載五旬節到了,門徒都聚集在一處；以下是講到聖靈降臨,如同火焰,分開落在他們各人頭上；他們就被聖靈充滿.請留意聖靈降臨下來,好像舌頭如同火焰.火能燒盡一切污穢,門徒從前的驕傲,嫉妒,自誇,以及其他等等的罪惡,在那天聖靈降臨時,他們被聖靈清潔,聖靈就充滿他們.各位!五旬節聖靈降臨的能力,並未過去,祂仍要在我們中間運行,要潔淨我們一切的污穢.在十七世紀,英國倫敦發生黑死病,叫許多人死亡；當局雖盡量設法,仍不能防止此病蔓延.直到有一天,倫敦忽然發生大火,將全城燒盡,此後黑死病就停止了.由此可見火之潔淨消毒力,是何等的大.聖靈也要如火焚燒我們.(羅八2)「因為賜生命聖靈的律,在基督耶穌釋放了我,使我脫離罪和死的律了.」請留意,人被聖靈充滿之後,就不再受罪捆綁,因而進入到聖潔的生命中.罪就是鎖鏈,就是捆綁；而聖靈的工作,就是解脫,使我們得自由,得釋放!這不是靠著自己,乃是倚靠聖靈.正如(亞四 6)「不是倚靠勢力,不是倚靠才能,乃是倚靠耶和華的靈,方能成事.」今日很需要聖靈來充滿我們,使我們得釋放,得自由.我們常會發脾氣,尋找人的錯處,恨人,論斷人,辱罵人,攻擊人,以及驕傲,嫉妒,貪圖虛榮,自大,我們有種種的罪惡,自己也覺得像是處在(羅馬七章)的情形下,我們又像保羅所說:「我真是苦呀!誰能救我脫離取死的身體呢?」但感謝神!當聖靈充滿我們之後,我們就有(羅八2)那種自由了.在第七章裡面,我們看不見有「聖靈」二字,而在第八章卻題了廿多次.讚美主!多年前一間聖經學院,請了一位講員,他講聖靈充滿的問題.雖然聚會拖得很長,但其中有一學生,原來恨惡一個人,非常傲慢,撕裂了別人的內衣；聖靈就感動他認罪,全院因此而大得復興.
\newline
\newline
(二)有力量作工(徒一8)主耶穌說:「但聖靈降臨在你們身上,你們就必得著能力；並要在猶太全地,撒瑪利亞,直到地極,作我的見證.」這是主耶穌臨升天前,應許門徒的.這叫我們知道被聖靈充滿的人,才有力量來作見證為主作工.原來我們信主後,很想為主作見證,但作見證卻無力量.彼得本來很軟弱,在一個使女面前,尚且不敢認主!但五旬節到了,他就被聖靈充滿,他的見證大有力量,叫人札心,就有三千五千人歸信主.保羅得著聖靈充滿之後,不只一次,而且是多次,他的講道,大有力量.「講的道,不是用智慧委婉的言語,乃是用聖靈和大能的明證.」(林前二4)這節聖經可以說明保羅講道得力的秘訣何在?他完全是靠著聖靈的大能來為主作見證.美國的叨雷博士,被聖靈充滿他的時候,他在各地講道,大有能力,多人札心歸信主.一次他到澳洲講道,有一萬六千人聽道,其實那會場只可容納八千人.他在那裡數週,已帶領了八千人信主.他另外到一地方講道,又有八千人歸主.總計他在英國,美國,澳洲,印度等地講道,有十萬人歸信主.叨雷博主靠聖靈大能,才有這樣的力量.
\newline
\newline
(三)充滿喜樂(約十五11)「這些事我已經對你說了,是要叫我的喜樂,存在你們心裡,並叫你們的喜樂可以滿足.」請留意在約翰十五章之前主耶穌應許降下聖靈來,到了十五章就說,因祂的喜樂,可以滿足.因此,一個人被聖靈充滿之後,就充滿喜樂,這種喜樂不是人可以給的,乃是主所賜的.充滿喜樂的人,就把一切的疑惑,懼怕,痛苦,都擠出去；又用知足的平安喜樂,來充滿我們的心.舊式用的風琴,為甚麼能發出聲音呢?不是靠著琴鍵,乃是由於風箱裡面的氣,發出了力量.如果風箱沒有氣,在鍵盤上彈琴的琴師,就無法施展其技了.神的靈若不充滿我們的心,我們唱的詩也會冰冷得如死一般.但聖靈充滿我們的心,我們的喜樂就盈溢,唱出來的詩也不同.多年前挪威有一個琴師名叫南遜,他離家往北極探險,已有兩年,他知道愛妻掛心他就寫封平安信,綁在鴿子的腳上,鴿子經過一千里的海面,回到南遜的家,他妻子自然歡喜若狂了.鴿子是預表聖靈,聖靈充滿我們之後,我們就有平安和喜樂.今天的基督徒為何沒有平安和喜樂,自然是因為未被聖靈充滿.
\newline
\newline
(四)結出靈果(加五22-23)「聖靈所結的果子,就是仁愛,喜樂,和平,忍耐,恩慈,良善,信實,溫柔,節制,這樣的事,沒有律法禁止.」這九個果子,是一個被聖靈充滿的人,從他內心裡流露出來的.請問各位,你到底有沒有結出這九樣果子來呢?」一個姊妹她信主之後,很盼望丈夫也一同信,她勸勉丈夫信,丈夫不肯.但有一天,丈夫卻自動與他妻子,同到禮拜堂赴祈禱會；向牧師表示,我要決心信耶穌.在歸途中,妻問說,你為何忽然決心信主呢?夫答我本來心裡很剛硬不信,但與你結婚後,因為見你有仁愛,喜樂,和平......等美德,所以我決心要信耶穌.這九個果子的次序,第一是仁愛,為甚麼把仁愛排在第一呢?這是聖靈奇妙的安排,叫我們知道愛心是何等重要.保羅在(林前十三13)清楚說:「其中最大的是愛.」本章聖經前面說,即使能說萬人的方言,天使的話語,有信能移山倒海......卻沒有愛,仍然與我無益.可見聖靈所結的果子頭一件就是愛的重要.
\newline
\newline
(五)要遵從引導,在使徒行傳記載,有幾位被聖靈充滿的人,他們都是順從聖靈引導的,好像腓利,他原來在撒瑪利亞,工作非常有效果的,他不是要靠自己,看見環境好而留戀.聖靈叫他去曠野,他就立刻去.埃提亞伯的太監,在車上讀聖經.聖靈對腓利說,你去靠近那車；於是腓利就去靠近那車,向太監說,你明白聖經嗎?於是根據聖經來帶領他信主.
\newline
\newline
保羅彼得,都不是靠自己,完全是倚靠聖靈行事.但願主幫助我們,每個人被聖靈充滿之後,能夠遵從聖靈的引導帶領.
\newline
\newline
(六)榮耀父神(弗三19)保羅禱告求神叫一切所充滿的,充滿他們,然後願神在教會中得著榮耀,直到世世代代,永永遠遠!可見一個被聖靈充滿的人,才能夠榮耀神.主耶穌的臨別贈言,所講的話,也是如此.(約十六14)「祂要榮耀我,因為祂要將受於我的,告訴你們.」聖靈不但叫人為罪為義,為審判責備自己；祂更要引導人明白真理,榮耀基督.所以我們知道一個被聖靈充滿的人,不是榮耀自己,乃是榮耀神.被聖靈充滿的人,常過著平淡的生活,而且受逼迫.雅各約翰,被聖靈充滿之後,也過著貧窮金銀都沒有的生活,雅各更為主而殉道.司提反被聖靈充滿也為主殉道榮耀神.我們被聖靈充滿,就有一個高尚的動機.若詳細看(徒八14-24)彼得去撒瑪利亞為門徒按手,叫他們受聖靈.行邪術的西門,拿錢給使徒,要買這權柄,為人按手得勝靈.他的動機不對,彼得就責備他:「你在這道上無分無關,因為你在神面前,你的心不正.」西門不純正的動機,應作為我們的鑑戒!因此,我們追求聖靈的充滿,必須動機純正.有一次美國叨雷博士到一個地方講道,講的題目都是聖靈充滿.會後一個牧師與他同行,牧師說我很想得著聖靈充滿,叨雷博士問何故,他答,我被聖靈充滿後,教會就必器重我,信徒也尊重我；因為我做牧師,受人藐視.叨雷博士叫他不要追求了,因他動機不對!
\newline
\newline
(七)能夠進入國度(太二十五1-3)這段經文所講是十童女的比喻,五個被撇下是代表沒有被聖靈充滿的；而五個聰明的,卻代表得著了聖靈的充滿的.因為愚拙的拿著燈,卻不預備油；聰明的拿著燈,又預備油在器皿裡.歷來解釋聖經的人都說,油是預表聖靈.愚拙的說請分點油給我們,因為我們的燈要滅了.這五個聰明的童女,得著聖靈的充滿,因此可以在主國度裡,同赴羔羊的筵席(啟十九9).那五個愚拙的童女回來時,國度的筵席已經過了,她們雖然苦求,主卻回答說:我不認識你們.(原文是我不嘉許你們)可惜今日教會不少的人,有重生的生命,卻沒有聖潔的生活；只有聖靈,卻不被聖靈充滿；只是清楚得救,卻沒有追求國度的賞錫.我們不但得救就算,還要得著賞賜!
\newline
\newline


\newpage

\section{第50屆~港九培靈會~鄭德音~第一講}
\label{sec:731}
\textbf{神是誰?}
\newline
\newline
連結: \href{https://www.hkbibleconference.org/session-message/view/731}{\texttt{https://www.hkbibleconference.org/session-message/view/731}}
\newline
\newline
\hyperref[sec:740]{< < < PREV SERMON < < <}
~
\hyperref[sec:toc]{[返目錄]}
~
\hyperref[sec:732]{> > > NEXT SERMON > > >}
\newline
\newline
使徒行傳 17:24-17:25
\newline
\begin{longtable}{cl}
\hline
\hline
章節 & 經文 (和合本修訂版)\\
\hline
17:24 & \begin{tabularx}{0.7\textwidth}{X} 他是創造宇宙和其中萬物的神；他既是天地的主，就不住在人手所造的殿宇裡， \end{tabularx} \\ \\ \relax
17:25 & \begin{tabularx}{0.7\textwidth}{X} 也不用人手去服侍，好像缺少甚麼似的；自己倒將生命、氣息、萬物賜給萬人。 \end{tabularx} \\ \\
[1ex]
\hline
\hline
\end{longtable}
神是與人有關係的主
\newline
\newline
「並且被造的,沒有一樣在祂面前不顯然的；原來萬物,在那與我們有關係的主眼前,都是赤露敞開的.」(來四13)
\newline
\newline
──神與人的關係,正如樹與枝子的關係.
\newline
\newline
「若有幾根枝子被折下來,你這野橄欖得接在其中,一同得著橄欖根的肥汁；你就不可向舊枝子誇口,若是誇口,當知道不是你托著根,乃是根托著你.」(羅十一17-18)
\newline
\newline
「你們要常在我裡面,我也常在你們裡面.枝子若不常在葡萄樹上,自己就不能結果子；你們若不常在我裡面,也是這樣.我是葡萄樹,你們是枝子；常在我裡面的,我也常在他裡面,這人就多結果子；因為離了我,你們就不能作甚麼.人若不常在我裡面,就像枝子丟在外面枯乾,人拾起來,扔在火裡燒了.」(約十五 4-6)
\newline
\newline
(壹)永生神乃是創造主
\newline
\newline
「起初神創造天地」.(創一1)
\newline
\newline
「創造宇宙和其中萬物的神,既是天地的主......」(徒十七24)
\newline
\newline
聖經宣告創造主的存在
\newline
\newline
「你且問走獸,走獸必指教你；又問空中的飛鳥,飛鳥必告訴你；或與地說話,地必指教你；海中的魚,也必向你說明.看這一切,誰不知道是耶和華的手作成的呢?凡活物的生命和人類的氣息,都在祂手中.」(伯十二7-10)
\newline
\newline
「耶和華我們的主阿,你的名在全地何其美；你將你的榮耀彰顯於天.你因敵人的緣故,從嬰孩的口中,建立了能力,使仇敵和報仇的,閉口無言.我觀著你指頭所造的天,並你所陳設的月亮星宿,便說,人算甚麼,你竟顧念他；世人算甚麼,你竟眷顧他.」(詩八1-4)
\newline
\newline
「諸天述說神的榮耀；穹蒼傳揚祂的手段.這日到那日發出言語；這夜到那夜傳出知識.」(詩十九1-2)
\newline
\newline
「......我們傳福音給你們,是叫你們離棄這些虛妄,歸向那創造天,地,海,和其中萬物的神.祂在從前的世代,任憑萬國各行其道；然而為自己未嘗不顯出證據來,就如常施恩惠,從天降雨,賞賜豐年,叫你們飲食飽足,滿心喜樂.」(徒十四15,17)
\newline
\newline
「神的事情,人所能知道的,原顯明在人心裡；因為神已經給他們顯明.自從造天地以來,神的永能和神性是明明可知的,雖是眼不能見,但藉著所造之物,就可以曉得,叫人無可推諉......」(羅一19-20)
\newline
\newline
「因為房屋都必有人建造；但建造萬物的就是神.」(來三4)
\newline
\newline
創造主對人的創造
\newline
\newline
「神說,我們要照著我們的形像,按著我們的樣式造人,使他們管理海裡的魚,空中的鳥,地上的牲畜,和全地,並地上所爬的一切昆蟲.神就照著自己的形像造人,乃是照著祂的形像造男造女.」(創一26-27)
\newline
\newline
「創造宇宙和其中萬物的神......」(徒十七24)
\newline
\newline
創造主是世人生命的根本
\newline
\newline
「因為萬有都是本於祂,倚靠祂,歸於祂；願榮耀歸給祂,直到永遠.阿們.」(羅十一36)
\newline
\newline
「因為在你那裡,有生命的源頭；在你的光中,我們必得見光.」(詩三十六9)
\newline
\newline
人的生存在乎神
\newline
\newline
「神從一本造出萬族的人,住在全地上,並且預先定準他們的年限,和所住的疆界；要叫他們尋求神,或者可以揣摩而得,其實祂離我們不遠；我們生活,動作,存留,都在乎祂......」(徒十七26-28)
\newline
\newline
人有神作生命的根本,才得有福生存,否則有禍.
\newline
\newline
「蒲草沒有泥,豈能發長；蘆荻沒有水,豈能生發.尚青的時候,還沒有割下,比百樣的草先枯槁.凡忘記神的人,景況也是這樣；不虔敬人的指望要滅沒；他所仰賴的必折斷,他所倚靠的是蜘蛛網.」(伯八11-14)
\newline
\newline
──這正如枝子有樹作生命的根本,則生存在有福中,否則有禍.
\newline
\newline
「你們要常在我裡面,我也常在你們裡面.枝子若不常在葡萄樹上,自己就不能結果子；你們若不常在我裡面,也是這樣.我是葡萄樹,你們是枝子；常在我裡面的,我也常在他裡面,這人就多結果子；因為離了我,你們就不能作甚麼.人若不常在我裡面,就像枝子丟在外面枯乾,人拾起來,扔在火裡燒了.」(約十五 4-6)
\newline
\newline
(貳)永生神乃是統治主
\newline
\newline
「......是天地的主......自己倒將生命,氣息,萬物,賜給萬人.」(徒十七24-25)
\newline
\newline
「耶和華在天上立定寶座；祂的權柄統管萬有.」(詩一○三19)
\newline
\newline
統治主乃是世人生命的倚靠
\newline
\newline
「萬有都是......倚靠祂......」(羅十一36)
\newline
\newline
「我們生活,動作,存留,都在乎祂......」(徒十七28)
\newline
\newline
「你且問走獸,走獸必指教你......看這一切,誰不知是耶和華的手作成的呢.凡活物的生命,和人類的氣息,都在祂手中.」(伯十二7-10)
\newline
\newline
倚靠統治主的人則生命在有福中,否則有禍.
\newline
\newline
「倚靠耶和華,以耶和華為可靠的,那人有福了.他必像樹栽於水旁,在河邊扎根,炎熱來到,並不懼怕,葉子仍必青翠,在乾旱之年毫無掛慮,而且結果不止.」(耶十七7-8)
\newline
\newline
「耶和華如此說,倚靠人血肉的膀臂,必中離棄耶和華的,那人有禍了；因為他必像沙漠的杜松,不見福樂來到,卻要住曠野乾旱之處,無人居住的鹼地.」(耶十七5-6)
\newline
\newline
──正如枝子倚靠樹則有福,離棄樹則有禍.(約十五4-6)
\newline
\newline
(參)永生神乃是救贖主
\newline
\newline
「......你的後裔和女人的後裔,也彼此為仇；女人的後裔要傷你的頭,你要傷他的腳跟.」(創三15)
\newline
\newline
「雅各阿,創造你的耶和華,以色列阿,造成你的那位,現在如此說,你不要害怕,因為我救贖了你.」
\newline
\newline
「世人蒙眛無知的時候,神並不鑑察,如今卻吩咐各處的人都要悔改.」(徒十七30)
\newline
\newline
(一)救贖主乃是得救贖之人屬靈生命的「根本」
\newline
\newline
「萬有都是本於祂......」(羅十一36)
\newline
\newline
「因為在你那裡,有生命的源頭；在你的光中,我們必得見光.」(詩三十六9)
\newline
\newline
世人因悔改相信則得蒙救贖
\newline
\newline
「我差你到他們那裡去,要叫他們的眼睛得開,從黑暗中歸向光明,從撒旦權下歸向神；又因信我,得蒙赦罪,和一切成聖的人同得基業.」(徒二十六18)
\newline
\newline
「因為世人都犯了罪,虧缺了神的榮耀.如今卻蒙神的恩典,因基督耶穌的救贖,就白白的稱義.神設立耶穌作挽回祭,是憑著耶穌的血,藉著人的信,要顯明神的義......也稱信耶穌的人為義.」(羅三23-26)
\newline
\newline
得救贖的人,則有救贖主作屬靈生命的根本
\newline
\newline
「盼望那無謊言的神,在萬古之先所應許的永生.」(多一2)
\newline
\newline
「願頌讚歸與我們主耶穌基督的父神,祂曾照自己的大憐憫,藉耶穌基督從死裡復活,重生了我們,叫我們有活潑的盼望,可以得著不能朽壞,不能玷污,不能衰殘,為你們存留在天上的基業.」(彼前一3-4)
\newline
\newline
「認識你獨一的真神,並且認識你所差來的耶穌基督,這就是永生.」(約十七3)
\newline
\newline
「這見證,就是神賜給我們永生,這永生也是在祂兒子裡面.人有了神的兒子就有生命；沒有神的兒子就沒有生命.」(約壹五11-12)
\newline
\newline
有救贖主作生命的根本則有各樣屬靈福氣
\newline
\newline
「......神在基督裡曾賜給我們天上各樣屬靈福氣.」(弗一3)
\newline
\newline
──正如枝子有樹作生命的根本,則有生存的福:(約十五4-5)否則有禍.(約十五5-6)
\newline
\newline
(二)救贖主是得救贖之人,屬靈生命的「倚靠」
\newline
\newline
「萬有都是本於祂,靠於祂......」(羅十一36)
\newline
\newline
「......耶和華要將救恩定為城牆......使守信的義民得以進入.堅心倚賴你的,你必保守他十分平安,因為他倚靠你.你們當倚靠耶和華直到永遠；因為耶和華是永久的磐石.」(賽二十六1-4)
\newline
\newline
離棄救贖主的人則有禍,倚靠人的則有福
\newline
\newline
「耶和華如此說,倚靠人血肉的膀臂,心中離棄耶和華的,那人有禍了；因他必像沙漠的杜松,不見福樂來到,卻要住曠野乾旱之處,無人居住的鹼地.倚靠耶和華,以耶和華為可靠的,那人有福了.他必像樹栽於水旁,在河邊扎根,炎熱來到,並不懼怕,葉子仍必青翠,在乾旱之年毫無掛慮,而且結果不止.」(耶十七5-8)
\newline
\newline
──正如枝子離樹則有禍；倚靠樹則有福.(約十五4-6)
\newline
\newline
有救贖主作生命的倚靠,才能有豐盛生命和生活
\newline
\newline
「神的神能已將一切關乎生命和虔敬的事賜給我們......叫我們脫離世上從情慾來的敗壞,就得與神的性情有分.正因這緣故,你們要份外的殷勤；有了信心,又要加上德行；有了德行,又要加上知識；有了知識,又要加上節制；有了節制,又要加上忍耐；有了忍耐,又要加上虔敬；有了虔敬,又要加上愛弟兄的心；有了愛弟兄的心,又要加上愛眾人的心.你們若充充足足的有這幾樣......必叫你們豐豐富富的,得以進入我們主救主耶穌基督永遠的國.」(彼後一3-11)
\newline
\newline
有豐盛生命和生活的人,才能將榮耀歸神
\newline
\newline
「因為你們是用重價買來的；所以要在你們的身子上榮耀神.」(林前六20)
\newline
\newline
「所以你們或喫或喝,無論作甚麼,都要為榮耀神而行.」(林前十31)
\newline
\newline
「並靠著耶穌基督結滿了仁義的果子,叫榮耀稱讚歸與神.」(腓一11)
\newline
\newline
「你們的光也當這樣照在人前,叫他們看見你們的好行為,便將榮耀歸給你們在天上的父.」(太五16)
\newline
\newline
(三)救贖主乃是得救贖之人,屬靈生命的歸宿
\newline
\newline
「萬有都是本於祂,靠於祂,歸於祂......」(羅十一36)
\newline
\newline
「主阿,你世世代代作我們的居所.」(詩九十1)
\newline
\newline
「永生的神是你的居所......」(申三十三27)
\newline
\newline
有救贖主作屬靈生命的歸宿,則有永福
\newline
\newline
「就說,耶穌阿,你得國降臨的時候,求你記念我.耶穌對他說,我實在告訴你,今日你要同我在樂園裡了.」(路二十三42-43)
\newline
\newline
「於是王要向那右邊的說,你們這蒙我父賜福的,可來承受那創世以來為你們所預備的國.」(太二十五34)
\newline
\newline
「......看哪,神的帳幕在人間；祂要與人同住；他們要作祂的子民,神要親自與他們同在,作他們的神.」(啟二十一3-4)
\newline
\newline
「因為你必不將我的靈魂撇在陰間；也不叫你的聖者見朽壞.你必將生命的道路指示我；在你的面前有滿足的喜樂；在你右手中有永遠的福樂.」(詩十六10-11)
\newline
\newline
──正如枝子長在樹裡則與樹共存(約十五4-5)
\newline
\newline
(四)救贖主乃是得救贖之人,屬靈生命的榮耀
\newline
\newline
「......願榮耀歸給祂,直到永遠.」(羅十一36)
\newline
\newline
「我的拯救,我的榮耀,都在乎神......」(詩六十二7)
\newline
\newline
「那城內又不用日月光照；因有神的榮耀光照；又有羔羊為城的燈.」(啟二十一23)
\newline
\newline
「不再有黑夜；他們也不用燈光日光；因為主神要光照他們；他們要作王,直到永永遠遠.」(啟二十二5)
\newline
\newline
有救贖主作屬靈生命的榮耀,則有永榮
\newline
\newline
「死人復活也是這樣,......復活的是榮耀的......」(林前十五42-43)
\newline
\newline
「我們卻是天上的國民；並且等候救主,就是主耶穌基督,從天上降臨.祂要按著那能叫萬有歸服自己的大能,將我們這卑賤的身體改變形狀,和祂自己榮耀的身體相似.」(腓三20-21)
\newline
\newline
「睡在塵埃中的,必有多人復醒；其中有得永生的,有受羞辱永遠被憎惡的.智慧人必發光,如同天上的光；那使多人歸義的,必發光如星,直到永永遠遠.」(但十二2-3)
\newline
\newline
「那時義人在他們父的國裡,要發出光來,像太陽一樣.......」(太十三43)
\newline
\newline
「......你們的生命與基督一同藏在神裡面.基督是我們的生命,祂顯現的時候,你們也要與祂一同顯現在榮耀裡.」(西三3-4)
\newline
\newline
──正如使樹得榮耀的枝子則與樹共榮(約十五8).
\newline
\newline
「叫我們主耶穌的名,在你們身上的榮耀,你們也在祂身上得榮耀,都照著我們的神並主耶穌基督的恩.」(帖後一12)
\newline
\newline
(肆)永生神乃是審判主
\newline
\newline
「因為祂已經定了日子,要藉著祂所設立的人,按公義審判天下,並且叫祂從死裡復活,給萬人作可信的憑據.」(徒十七31)
\newline
\newline
「你這人哪......你以為能逃脫神的審判麼.還是你藐視祂豐富的恩慈,寬容,忍耐,不曉得祂的恩慈是領你悔改呢；你竟任著你剛硬不悔改的心,為自己積蓄忿怒,以致神震怒,顯祂公義審判日子來到.祂必照各人的行為報應各人......」(羅二3-6)
\newline
\newline
「父不審判甚麼人,乃將審判的事全交與子.」(約五22)
\newline
\newline
審判主使接受救贖的人,得永生永福的報應
\newline
\newline
「凡恆心行善,尋求榮耀尊貴；和不能朽壞之福的,就以永生報應他們.」(羅二8)
\newline
\newline
「這正是主降臨要在祂聖徒身上得榮耀,又在一切信的人身上顯為希奇的那日子......叫我們主耶穌的名,在你們身上得榮耀,你們也在祂身上得榮耀,都照著我們的神並主耶穌基督的恩.」(帖後一10-12)
\newline
\newline
「當人子在祂榮耀裡,同著眾天使降臨的時候,要坐在祂的榮耀寶座上；萬民都要聚集在祂面前,祂要把他們分別出來......於是王要向右邊的說,你們這蒙我父賜福的,可來承受那創世以來為你們所預備的國......」(太二十五31-40)
\newline
\newline
「論到睡了的人,我們不願意弟兄們不知道,恐怕你們憂傷,像那些沒有指望的人一樣.我們若信耶穌死而復活了,那已經在耶穌裡睡了的人,神也必將他與耶穌一同帶來.我現在照主的話告訴你們一件事；我們這活著還存留到主降臨的人,斷不能在那已經睡了的人之先；因為主必親自從天降臨,有呼叫的聲音,和天使長的聲音,又有神的號吹響；那在基督裡死了的人必先復活.以後我們這活著還存留的人,必和他們一同被提到雲裡,在空中與主相遇；這樣,我們就要和主永遠同在.所以你們當用這些話彼此勸慰.」(帖前四13-18)
\newline
\newline
「我們卻是天上的國民；並且等候救主,就是主耶穌基督,從天降臨.他要按著那能叫萬有歸服自己的大能,將我們這卑賤的身體改變形狀,和祂自己榮耀的身體相似.」(腓三20-21)
\newline
\newline
「在頭一次復活有分的,有福了,聖潔了；第二次的死在他們身上沒有權柄；他們必作神和基督的祭司,並要與基督作王一千年.」(啟二十6)
\newline
\newline
「我聽見有大聲音從寶座出來說,看哪,神的帳幕在人間；祂要與人同住,他們要作祂的子民,神要親自與他們同在,作他們的神；神要擦去他們一切的眼淚；不再有死亡,也不再有悲哀,哭號,疼痛,因為以前的事都過去了.坐寶座的說,看哪,我將一切都更新了.......」(啟二十一3-5)
\newline
\newline
審判主使棄絕救贖的人,得永死永禍的報應
\newline
\newline
「耶穌在諸城中行了許多異能,那些城的人終不悔改,就在那時候責備他們說,哥拉汎哪,你有禍了!伯賽大阿,你有禍了!因為在你們中間所行的異能,若行在推羅西頓,他們早已披麻蒙灰悔改了.但我告訴你們,當審判的日子,推羅西頓所受的,比你們還容易受呢.迦百農阿,你已經升到天上；將來必墜落陰間；因為在你那裡所行的異能,若行在所多瑪,他還可以存到今日；但我告訴你們,當審判的日子,所多瑪所受的,比你還容易受呢.」(太十一20-24)
\newline
\newline
「惟有結黨不順從真理,反順從不義的,就以忿怒惱恨報應他們.」(羅二8)
\newline
\newline
「神既是公義的,祂必將患難報應那加患難給你們的人；也必使你們這受患難的人,與我們同得平安；那時,主耶穌同祂有能力的天使,從天上在火焰中顯現,要報應那不認識神,和那不聽從我主耶穌福音的人.他們要受刑罰,就是永遠沉淪,離開主的面和祂權能的榮光.」(帖後一6-9)
\newline
\newline
「王又要向那左邊的說,你們這被咒詛的人,離開我,進入那為魔鬼和他的使者所預備的永火著去.」(太二十五41)
\newline
\newline
「我又看見一個白色的大寶座,與坐在上面的；從祂面前天地都逃避,再無可見之處了.我又看見死了的人,無論大小,都站在寶座前；案卷展開了；並且另有一卷展開,就是生命冊；死了的人都憑著這些案卷所記載的,照他們所行的受審判.於是海交出其中的死人,死亡和陰間也交出其中的死人；他們都照各人所行的受審判.死亡和陰間也被扔在火湖裡；這火湖就是第二次的死.若有人名字沒記在生命冊上,他就被扔在火湖裡.」(啟二十11-15)
\newline
\newline
「認識你獨得的真神,並且認識你所差來的耶穌基督,這就是永生.」(約十七3)
\newline
\newline
──神與救贖──
\newline
\newline
1．經題:神是誰?
\newline
\newline
經文:徒十七24,25,30,31；羅十一36
\newline
\newline
──神是與人有關係的主.(來四13)這正如樹與枝子一樣.(羅十一17,約十五4-6)
\newline
\newline
(一)永生神乃是創造主(創一；徒十七24)
\newline
\newline
創造主是世人生命的「根本」.(羅十一36)
\newline
\newline
世人的「生存」在乎神.(徒十七28；耶十七7-8)
\newline
\newline
(二)永生神乃是統治主(創二；徒十七25)
\newline
\newline
統治主是世人生命的「倚靠」.(羅十一36)
\newline
\newline
世人的「生活」在乎神.(徒十七28；耶十七7-8)
\newline
\newline
(三)永生神乃是救贖主(創三；賽四十三1；徒十七30)
\newline
\newline
教贖世人的主,祂是屬靈生命的「根本」.(羅十一36)所以得救贖的人,因有主作屬靈生命的根本,便有屬靈的生命和生存.(羅十一17)
\newline
\newline
救贖世人的主,祂又是屬靈生命的「倚靠」.(羅十一36)所以得救贖的人,因有主作屬靈生命的倚靠,便有屬靈的豐盛生活(約十五4-5)
\newline
\newline
救贖世人的主,又是屬靈生命的「歸宿」.(羅十一36)所以得救贖的人,因有主作屬靈生命的歸宿,便有屬靈的永遠安息和福樂.(路二十三43；太二十五34；啟二十6,二十一3-4；詩十六11)
\newline
\newline
救贖世人的主,終是屬靈生命的「榮耀」.(羅十一36；啟二十二5,二十一23)所以得救贖的人,因有主作屬靈生命的榮耀,便享屬靈的永遠榮耀.(但十二3；太十三43)
\newline
\newline
(四)永生神乃是審判主(創三；徒十七31；羅二5-6)
\newline
\newline
審判主使接受救贖的人,得著永福的報應.(約三18,36；羅二7；太二十五34；啟二十6,二十一3-4)
\newline
\newline
審判主使棄絕救贖的人,受永禍的報應.(約三18,36；羅二8；太二十五41；啟二十15；十四9-11)
\newline
\newline


\newpage

\section{第50屆~港九培靈會~鄭德音~第二講}
\label{sec:732}
\textbf{「人的罪惡」與「神的救贖」}
\newline
\newline
連結: \href{https://www.hkbibleconference.org/session-message/view/732}{\texttt{https://www.hkbibleconference.org/session-message/view/732}}
\newline
\newline
\hyperref[sec:731]{< < < PREV SERMON < < <}
~
\hyperref[sec:toc]{[返目錄]}
~
\hyperref[sec:734]{> > > NEXT SERMON > > >}
\newline
\newline
創世記 3:1-3:24
\newline
\begin{longtable}{cl}
\hline
\hline
章節 & 經文 (和合本修訂版)\\
\hline
3:1 & \begin{tabularx}{0.7\textwidth}{X} 耶和華神所造的，惟有蛇比田野一切的走獸更狡猾。蛇對女人說：「神豈是真說，你們不可吃園中任何樹上所出的嗎？」 \end{tabularx} \\ \\ \relax
3:2 & \begin{tabularx}{0.7\textwidth}{X} 女人對蛇說：「園中樹上的果子，我們都可以吃； \end{tabularx} \\ \\ \relax
3:3 & \begin{tabularx}{0.7\textwidth}{X} 只是園子中間那棵樹的果子，神曾說：『你們不可吃，也不可摸，免得你們死。』」 \end{tabularx} \\ \\ \relax
3:4 & \begin{tabularx}{0.7\textwidth}{X} 蛇對女人說：「你們不一定死； \end{tabularx} \\ \\ \relax
3:5 & \begin{tabularx}{0.7\textwidth}{X} 因為神知道，你們吃的日子眼睛就開了，你們就像神一樣知道善惡。」 \end{tabularx} \\ \\ \relax
3:6 & \begin{tabularx}{0.7\textwidth}{X} 於是女人見那棵樹好作食物，又悅人的眼目，那樹令人喜愛，能使人有智慧，她就摘下果子吃了，又給了與她一起的丈夫，他也吃了。 \end{tabularx} \\ \\ \relax
3:7 & \begin{tabularx}{0.7\textwidth}{X} 他們二人的眼睛就開了，知道自己赤身露體，就編織無花果樹的葉子，為自己做成裙子。 \end{tabularx} \\ \\ \relax
3:8 & \begin{tabularx}{0.7\textwidth}{X} 天起了涼風，那人和他妻子聽見耶和華神在園中來回行走的聲音，就藏在園裡的樹木中，躲避耶和華神的面。 \end{tabularx} \\ \\ \relax
3:9 & \begin{tabularx}{0.7\textwidth}{X} 耶和華神呼喚那人，對他說：「你在哪裡？」 \end{tabularx} \\ \\ \relax
3:10 & \begin{tabularx}{0.7\textwidth}{X} 他說：「我在園中聽見你的聲音，我就害怕；因為我赤身露體，我就藏了起來。」 \end{tabularx} \\ \\ \relax
3:11 & \begin{tabularx}{0.7\textwidth}{X} 耶和華神說：「誰告訴你，你是赤身露體呢？莫非你吃了那樹上所出的，就是我吩咐你不可吃的嗎？」 \end{tabularx} \\ \\ \relax
3:12 & \begin{tabularx}{0.7\textwidth}{X} 那人說：「你賜給我、與我一起的女人，是她把那樹上所出的給我，我就吃了。」 \end{tabularx} \\ \\ \relax
3:13 & \begin{tabularx}{0.7\textwidth}{X} 耶和華神對女人說：「你怎麼會做這種事呢？」女人說：「那蛇引誘我，我就吃了。」 \end{tabularx} \\ \\ \relax
3:14 & \begin{tabularx}{0.7\textwidth}{X} 耶和華神對蛇說：「你既做了這事，就必受詛咒，比一切的牲畜和野獸更重。你必用肚子行走，終生吃土。 \end{tabularx} \\ \\ \relax
3:15 & \begin{tabularx}{0.7\textwidth}{X} 我要使你和女人彼此為仇，你的後裔和女人的後裔也彼此為仇。他要傷你的頭，你要傷他的腳跟。 」 \end{tabularx} \\ \\ \relax
3:16 & \begin{tabularx}{0.7\textwidth}{X} 又對女人說：「我必多多加增你懷胎的痛苦，你生兒女時必多受痛苦。你必戀慕你丈夫，他必管轄你。」 \end{tabularx} \\ \\ \relax
3:17 & \begin{tabularx}{0.7\textwidth}{X} 又對亞當說：「你既聽從你妻子的話，吃了那樹上所出的，就是我吩咐你不可吃的，土地必因你的緣故受詛咒；你必終生勞苦才能從土地得吃的。 \end{tabularx} \\ \\ \relax
3:18 & \begin{tabularx}{0.7\textwidth}{X} 土地必給你長出荊棘和蒺藜來；你也要吃田間的五穀菜蔬。 \end{tabularx} \\ \\ \relax
3:19 & \begin{tabularx}{0.7\textwidth}{X} 你必汗流滿面才有食物可吃，直到你歸了土地，因為你是從土地而出的。你本是塵土，仍要歸回塵土。」 \end{tabularx} \\ \\ \relax
3:20 & \begin{tabularx}{0.7\textwidth}{X} 那人給他妻子起名叫夏娃，因為她是眾生之母。 \end{tabularx} \\ \\ \relax
3:21 & \begin{tabularx}{0.7\textwidth}{X} 耶和華神用獸皮做衣服給亞當和他的妻子穿。 \end{tabularx} \\ \\ \relax
3:22 & \begin{tabularx}{0.7\textwidth}{X} 耶和華神說：「看哪，那人已經像我們中間的一個，知道善惡，現在恐怕他又伸手摘生命樹所出的來吃，就永遠活著。」 \end{tabularx} \\ \\ \relax
3:23 & \begin{tabularx}{0.7\textwidth}{X} 耶和華神就驅逐他出伊甸園，使他耕種土地，他原是從土地裡被取出來的。 \end{tabularx} \\ \\ \relax
3:24 & \begin{tabularx}{0.7\textwidth}{X} 耶和華神把那人趕出去，就在伊甸園東邊安設基路伯和發出火焰轉動的劍，把守生命樹的道路。 \end{tabularx} \\ \\
[1ex]
\hline
\hline
\end{longtable}
論到「人的罪惡」與「神的救贖」的經文
\newline
\newline
「雅各阿,創造你的耶和華,以色列阿,造成你的那位,現在如此說,你不要害怕,因為我救贖了你；我曾題你的名召你,你是屬我的.」(賽四十三1)
\newline
\newline
「以色列阿,你當仰望耶和華；因祂有慈愛,有豐盛的救恩.祂必救贖以色列脫離一切的罪孽.」(詩一三○7-8)
\newline
\newline
「神愛世人,甚至將祂的獨生子賜給他們,叫一切信祂的,不至滅亡,反得永生.因為神差祂的兒子降世,不是要定世人的罪,乃是要叫世人因祂得救.」(約三16-17)
\newline
\newline
「因為世人都犯了罪,虧缺了神的榮耀.如今卻蒙神的恩典,因基督耶穌的救贖,就白白的稱義.」(羅三23-25)
\newline
\newline
創世記三章論人的罪惡,與神的審判和救贖
\newline
\newline
(一)神給亞當的抉擇(創二16-17)
\newline
\newline
「耶和華神吩咐他說,園中各樣樹上的果子,你可以隨意喫；只是分別善惡樹上的果子,你不可喫,因為你喫的日子必定死.」(創二16-17)
\newline
\newline
「看哪,我今日將生與福,死與禍,陳明在你面前.」(申三十15)
\newline
\newline
「我今日呼天喚地向你作見證,我將生死,禍福,陳明在你面前,所以你要揀選生命,使你和你的後裔都得存活.且愛耶和華你的神,聽從祂的話,專靠祂,因為祂是你的生命；你的日子長久,也在乎祂......」(申三十19-20)
\newline
\newline
(二)亞當的悖逆(創三6)
\newline
\newline
「於是女人見那棵樹的果子好作食物,也悅人眼目,且是可喜愛的,能使人有智慧,就摘下果子來喫了；又給她丈夫,她丈夫也喫了.」(創三6)
\newline
\newline
「他們卻如亞當背約,在境內向我行事詭詐.」(何六7)
\newline
\newline
「因一人悖逆,眾人成為罪人......」(羅五19)
\newline
\newline
亞當之於罪惡──隱藏,推諉
\newline
\newline
「......便拿無花果的葉子,為自己編作裙子......那人和他的妻子聽見神的聲音,就藏在園裡的樹木中,躲避耶和華神的面.」(創三7-8)
\newline
\newline
「我若像亞當遮掩我的過犯,將罪孽藏在懷中.」(伯三十一33)
\newline
\newline
「遮掩自己罪過的,必不亨通；承認離棄罪過的,必蒙憐恤.常存敬畏的,便為有福；心存剛硬的,必陷在禍患裡.」(箴二十八13-14)
\newline
\newline
「那人說,你所賜給我,與我同居的女人,她把那樹上的果子給我,我就喫了.」(創三12)
\newline
\newline
(一)亞當犯了悖逆之罪的後果(創三7-24)
\newline
\newline
一,在罪惡之下(創三6)
\newline
\newline
「因為世人都犯了罪,虧缺了神的榮耀.」(羅三23)
\newline
\newline
「......因我們已經證明,猶太人和希利尼人都在罪惡之下......」(羅三9)
\newline
\newline
「惡人必被自己的罪孽捉住；他必被自己的罪惡如繩索纏繞.」(箴五22)
\newline
\newline
「我看出你正在苦膽之中被罪惡捆綁.」(徒八23)
\newline
\newline
二,在神定罪之下(創三11-18)
\newline
\newline
「耶和華是忌邪施報的神......大有能力,萬不以有罪的為無罪......」(鴻一2-3)
\newline
\newline
「......普世的人都伏在神審判之下......」(羅三19)
\newline
\newline
「如此說來,因一次的過犯,罪人都被定罪......」(羅五18)
\newline
\newline
三,在罪刑之下(創三19)
\newline
\newline
「我(神)必因邪惡刑罰世界,因罪孽刑罰惡人......」(賽十三11)
\newline
\newline
「神判定行這樣事的人是當死的......」(羅一32)
\newline
\newline
「你們現今所看為羞恥的事,當日有甚麼果子呢；那些事的結局就是死.」(羅六21)
\newline
\newline
「因為斷定罪名,不立刻施刑,所以世人滿心作惡.」(傳八11)
\newline
\newline
四,在罪權之下(創三8)
\newline
\newline
「......我是屬乎肉體的,是已經賣給罪了.因為我所作的,我自己不明白；我所願意的,我並不作；我所恨惡的,我倒去作.......既是這樣,就不是我作的,乃是住在我裡頭的罪作的.......我真是苦阿,誰能救我脫離這取死的身體呢?」(羅七14-24)
\newline
\newline
五,在罪性之下(創三8)
\newline
\newline
「這就如罪是從一人入了世界,死又是從罪來的,於是死就臨到眾人,因為眾人都犯了罪.」(羅五12)
\newline
\newline
「......罪惡是人民的羞辱.」(箴十四34)
\newline
\newline
「睡在塵埃中的,必有多人復醒；其中有得永生的,有受羞辱永遠被憎惡的.」(但十二2)
\newline
\newline
六,亞當不能自救(創三7)
\newline
\newline
「他們二人的眼睛就明亮了,才知道自己是赤身露體,便拿無花果樹的葉子,為自己編作裙子.」
\newline
\newline
「......我們仍犯罪,這景況已久；我們還能得救麼；我們都像不潔淨的人,所有的義都像污穢的衣服；我們都像葉子漸漸枯乾.我們的罪孽好像風把我們吹去......」(賽六十四6)
\newline
\newline
「......這樣誰能得救呢?......在人這是不能的；在神凡事都能.」(太十九25-26)
\newline
\newline
七,神公義的審判(創三13-19)
\newline
\newline
「耶和華是忌邪施報的神,......大有能力,萬不以有罪的為無罪......」(鴻一2-3)
\newline
\newline
「主耶和華阿,你若究察罪孽,誰能站得住呢；但在你有赦免之恩,要叫人敬畏你.」(詩一三○3-4)
\newline
\newline
八,神慈愛的救贖(創三15,21)
\newline
\newline
「他們在一切苦難中,祂也同受苦難......祂以慈愛和憐憫救贖他們......」(賽六十三9)
\newline
\newline
「以色列阿,你當仰望耶和華；因祂有慈愛,有豐盛的救恩.祂必救贖以色列脫離一切的罪孽.」(詩一三○7-8)
\newline
\newline
「我離棄你不過片時,卻要施大恩將你收回.我的怒氣漲溢,頃刻之間向你掩面,卻要以永遠的慈愛憐恤你；這是耶和華你的救贖主說的.」(賽五十四7-8)
\newline
\newline
九,亞當因順從而得救贖(創三21)
\newline
\newline
「主耶和華說:我指著我的永生起誓,我斷不喜悅惡人死亡；惟喜悅惡人轉離所行的道而活；以色列阿,你們轉回,轉回罷,離開惡道；何必死亡呢?」(結三十三11)
\newline
\newline
「當趁耶和華可尋找的時候尋找祂,相近的時候求告祂；惡人當離棄自己的道路；不義的人當除掉自己的意念,歸向耶和華,耶和華就必憐恤他；當歸向我們的神,因為神必廣行赦免.」(賽五十五6-7)
\newline
\newline
「......除了我以外,再沒有神；我是公義的神,又是救主,......地極的人都當仰望我,就必得救......」(賽四十五21-22)
\newline
\newline
「神是為我們施行諸般救恩的神；人能脫離死亡,是在乎主耶和華.」(詩六十八20)
\newline
\newline
「但如今神的義在律法以外已經顯明出來,有律法和先知為證；就是神的義,因信耶穌基督,加給一切相信的人,並沒有分別......」(羅三21-22)
\newline
\newline
「我們不至消滅,是出於耶和華的諸般慈愛,是因祂的憐憫,不至斷絕.......凡等候耶和華,心裡尋救祂的,耶和華必施恩給他.」(哀三22,25)
\newline
\newline
2．經題:「人的罪惡」與「神的救贖」
\newline
\newline
經文:創三；羅三23-25
\newline
\newline
──創世記三章,論到「人的罪惡」與「神的救贖」.聖經的主題,也是論到「人的罪惡」與「神的救贖」.
\newline
\newline
(一)神給亞當的抉擇(申三十15,19-20；創二16-17)
\newline
\newline
(二)亞當的悖逆(創三6；羅五19；何六7)
\newline
\newline
(三)亞當悖逆的後果:
\newline
\newline
1.在罪惡之下(創三6；羅三23,9；箴五22)
\newline
\newline
2.在定罪之下(創三11-18；羅三19,五18)
\newline
\newline
3.在罪刑之下(創三19；羅一32,六23)
\newline
\newline
4.在罪權之下(創三8；羅七14-24)
\newline
\newline
5.在罪性之下(創三10；羅七17；詩五十一5)
\newline
\newline
(四)亞當不能自救(創三7；賽六十四6；太十九25-26)
\newline
\newline
(五)神公義的審判(創三13-19；鴻一2；詩一三○3-4)
\newline
\newline
(六)神慈愛的救贖(創三15,21；詩一三○7-8；賽五十四7-8)
\newline
\newline
(七)亞當因順從而得蒙救贖(創三21；羅三22；腓三9)
\newline
\newline


\newpage

\section{第50屆~港九培靈會~鄭德音~第三講}
\label{sec:734}
\textbf{神的救贖工作}
\newline
\newline
連結: \href{https://www.hkbibleconference.org/session-message/view/734}{\texttt{https://www.hkbibleconference.org/session-message/view/734}}
\newline
\newline
\hyperref[sec:732]{< < < PREV SERMON < < <}
~
\hyperref[sec:toc]{[返目錄]}
~
\hyperref[sec:714]{> > > NEXT SERMON > > >}
\newline
\newline
以賽亞書 43:1-43:8
\newline
\begin{longtable}{cl}
\hline
\hline
章節 & 經文 (和合本修訂版)\\
\hline
43:1 & \begin{tabularx}{0.7\textwidth}{X} 雅各啊，創造你的耶和華，以色列啊，造成你的那位，現在如此說：「你不要害怕，因為我救贖了你；我曾提你的名召你，你是屬我的。 \end{tabularx} \\ \\ \relax
43:2 & \begin{tabularx}{0.7\textwidth}{X} 你從水中經過，我必與你同在，你渡過江河，水必不漫過你；你在火中行走，也不被燒傷，火焰必不燒著你身。 \end{tabularx} \\ \\ \relax
43:3 & \begin{tabularx}{0.7\textwidth}{X} 因為我是耶和華－你的神，是以色列的聖者－你的救主；我使埃及作你的贖價，使古實和西巴代替你。 \end{tabularx} \\ \\ \relax
43:4 & \begin{tabularx}{0.7\textwidth}{X} 因我看你為寶貝為尊貴；又因我愛你，所以使人代替你，使萬民替換你的生命。 \end{tabularx} \\ \\ \relax
43:5 & \begin{tabularx}{0.7\textwidth}{X} 你不要害怕，因我與你同在；我必領你的後裔從東方來，又從西方召集你。 \end{tabularx} \\ \\ \relax
43:6 & \begin{tabularx}{0.7\textwidth}{X} 我要對北方說，交出來！對南方說，不可扣留！要將我的兒子從遠方帶來，將我的女兒從地極領回， \end{tabularx} \\ \\ \relax
43:7 & \begin{tabularx}{0.7\textwidth}{X} 就是凡稱為我名下的人，是我為自己的榮耀創造的，是我所塑造，所做成的。」 \end{tabularx} \\ \\ \relax
43:8 & \begin{tabularx}{0.7\textwidth}{X} 你要將有眼卻瞎、有耳卻聾的民都帶出來！ \end{tabularx} \\ \\
[1ex]
\hline
\hline
\end{longtable}
神對世人的救贖
\newline
\newline
「雅各阿,創造你的耶和華,以色列阿,造成你的那位,現在如此說,你不要害怕,因為我救贖了你......」(賽四十三1)
\newline
\newline
「......凡有血氣的,必都知道我耶和華是你的救主,是你的救贖主,是雅各的大能者.」(賽四十九26)
\newline
\newline
「以色列阿,你當仰望耶和華；因祂有慈愛,有豐盛的救恩.祂必救贖以色列脫離一切的罪孽.」(詩一三○7-8)
\newline
\newline
「神愛世人,甚至將祂的獨生子賜給他們,叫一切信祂的,不至滅亡,反得永生.因為神差祂的兒子降世,不是要定世人的罪,乃是要叫世人因祂得救.」(約三16-17)
\newline
\newline
「耶穌就對他們說,我父作事直到如今,我也作事.」(約五17)
\newline
\newline
「我在父裡面,父在我裡面,你不信麼.我對你們所說的話,不是憑著自己說的,乃是住在我裡面的父作祂自己的事.」(約十四10)
\newline
\newline
「......證明他們沉淪,你們得救,都是出於神.」(腓一28)
\newline
\newline
「但你們得在基督耶穌裡,是本乎神,神又使祂成為我們的智慧,公義,聖潔,救贖；如經上所記,「誇口的當指著主誇口.」(林前一30)
\newline
\newline
神的救贖工作,正如栽培人的工作
\newline
\newline
「我是真葡萄樹,我父是栽培的人.」(約十五1)
\newline
\newline
「不錯,他們因為不信,所以被折下來；你因為信,所以立得住；......可見神的恩慈和嚴厲......而且他們若不是長久不信,仍要被接上；因為神能夠把他們從新接上.」(羅十一20-23)
\newline
\newline
「......你這野橄欖得接在其中,一同得著橄欖根的肥汁......」(羅十一17)
\newline
\newline
「凡屬我不結果子的枝子,祂就剪去；凡結果子的,祂就修理乾凈,使枝子結果子更多.」(約十五2)
\newline
\newline
「......惟有萬靈的父管教我們,是要我們得益處,使我們在祂的聖潔上有分.凡管教的事.當時不覺得快樂,反覺得愁苦；後來卻為那經練過的人,結出平安的果子,就是義......」(來十二10-11)
\newline
\newline
(一)神在基督裡向在罪惡之下的人,施行救贖工作
\newline
\newline
──神要救贖人脫離罪惡
\newline
\newline
「因為世人都犯了罪,虧缺了神的榮耀.」(羅三23)
\newline
\newline
「......因我們已經證明,猶太人和希利尼人都在罪惡之下......」(羅三9)
\newline
\newline
「惡人必被自己的罪孽捉住；他必被自己的罪惡如繩索纏繞.」(箴五22)
\newline
\newline
「耶穌對他說,......我在父裡面,父在我裡面,你不信麼.我對你們所說的話,不是憑自己說的,乃是住在我裡面的父作祂自己的事.」(約十四9-10)
\newline
\newline
「耶穌就對他們說,我父作事直到如今,我也作事.」(約五17)
\newline
\newline
「耶穌對他們說,我實實在在的告訴你們,子憑著自己不能作甚麼,惟有看見父所作的,子才能作,父所作的事,子也照樣作.」(約五19)
\newline
\newline
(二)神使悔改相信的人,得以歸在基督裡脫離罪惡,成為得救的人
\newline
\newline
「神愛世人,甚至將祂的獨生子賜給他們,叫一切信祂的,不至滅亡,反得永生.」(約三16)
\newline
\newline
「因為神差祂的兒子降世,不是要定世人的罪,乃是要叫世人因祂得救.信祂的人,不被定罪；不信的人,罪已經定了,因為他不信神獨生子的名.光來到世間,世人因自己的行為是惡的,不愛光倒愛黑暗,定他們的罪就是在此.」(約三17-19)
\newline
\newline
「......誰從古時指明,誰從上古述說,不是我耶和華麼?除了我以外,再沒有神；我是公義的神,又是救主,除了我以外,再沒有別神.地極的人都當仰望我,就必得救；因為我是神,再沒有別神.」(賽四十五21-22)
\newline
\newline
「......這是證明他們沉淪,你們得救,都是出於神.」(腓一28)
\newline
\newline
「但你們得在基督耶穌裡,是本乎神,神又使祂成為我們的智慧,公義,聖潔,救贖；如經上所記,誇口的當指著主誇口.」(林前一30)
\newline
\newline
「神救了我們,以聖召召我們,不是按我們的行為,乃是按祂的旨意和恩典；這恩典是萬古之先,在基督耶穌裡賜給我們的；但如今藉著我們救主基督耶穌的顯現,才表明出來了.」(提後一9-10)
\newline
\newline
「基督耶穌降世,為要拯救罪人；這話是可信的,是十分可佩服的；在罪人中我是個罪魁.然而我蒙了憐憫,是因耶穌基督要在我這罪魁身上,顯明他一切的忍耐,給後來信祂得永生的人作榜樣.」(提前一15-16)
\newline
\newline
「......你要給祂起名叫耶穌；因祂要將自己的百姓從罪惡裡拯救出來.」(太一21)
\newline
\newline
「基督釋放了我們,叫我們得以自由,所以要站立得穩,不要再被奴僕的軛挾制.」(加五1)
\newline
\newline
「耶穌回答說,我實實在在的告訴你們；所有犯罪的,就是罪惡的奴僕.......所以天父的兒子若叫你們自由,你們就真自由了.」(約八34-36)
\newline
\newline
(三)神在基督裡向在定罪之下的人,施行救贖的工作
\newline
\newline
──神要救贖人脫離定罪
\newline
\newline
「我們曉得律法上的話,都是對律法之下的人說的,好塞住各人的口,叫普世的人都伏在神審判之下......」(羅三19)
\newline
\newline
「如此說來,因一次的過犯,眾人都被定罪......」(羅五18)
\newline
\newline
「他們既然故意不認識神,神就任憑他們存邪僻的心,行那些不合理的事......他們雖知道神判定,行這樣事的人是當死的,然而他們不但自己去行,還喜歡別人去行.」(羅一28-32)
\newline
\newline
「耶穌大聲說,信我的,不是信我,乃是信那差我來的.人看見我,就是看見那差我來的.我到世上來,乃是光,叫凡信我的,不住在黑暗裡.若有人聽見我的話不遵守,我不審判他；我來本不是要審判世界,乃是要拯救世界.棄絕我不領受我的話的人,有審判他的；就是所講的道,在末日要審判他.因為我沒有憑著自己講；惟有差我來的父,已經給我命令,叫我說甚麼,講甚麼.我也知道他的命令就是永生；故此我所講的話,正是照著父對我所說的.」(約十二44-50)
\newline
\newline
神使悔改相信的人,得以歸在基督裡脫離定罪,成為得救稱義的人
\newline
\newline
「耶和華救贖祂僕人的靈魂；凡投靠祂的,必不至定罪.」(詩三十四22)
\newline
\newline
「誰能控告神所揀選的人呢；有神稱他們為義了.」(羅八33)
\newline
\newline
「我們藉愛子的血,得蒙救贖,過犯得以赦免,乃是照祂豐富的恩典.」(弗一7)
\newline
\newline
「我們既因信稱義,就藉著我們的主耶穌基督,得與神相和.我們又藉著祂,因信得進入現在所站的這恩典中,並且歡歡喜喜盼望神的榮耀.」(羅五1-2)
\newline
\newline
(四)神在基督裡向在罪刑之下的人,施行救贖工作
\newline
\newline
──神要救贖人脫離罪刑
\newline
\newline
「罪的工價乃是死；惟有神的恩賜,在我們的主基督耶穌裡,乃是永生.」(羅六23)
\newline
\newline
「看哪......犯罪的他必死亡.」(結十八4)
\newline
\newline
「你們現今所看為羞恥的事,當日有甚麼果子呢；那些事的結局就是死.」(羅六21)
\newline
\newline
「主所應許的尚未成就,有人以為祂是耽延；其實不是耽延,乃是寬容你們,不願有一人沉淪,乃願人人都悔改.」(彼後三9)
\newline
\newline
「神......叫他作君王作救主,將悔改的心,和赦罪的恩賜給以色列人.」(徒五31)
\newline
\newline
「......這樣看來,神也賜恩給外邦人,叫他們悔改得生命了.」(徒十一18)
\newline
\newline
神使悔改相信的人,得以歸在基督裡脫離罪刑,成為得永生的人
\newline
\newline
「神愛世人,甚至將祂的獨生子賜給他們,叫一切信祂的,不至滅亡,反得永生.因為神差祂的兒子降世,不是要定世人的罪,乃是要叫世人因祂得救.」(約三16-17)
\newline
\newline
「信祂的人,不被定罪；不信的人,罪已經定了,因為他不信神獨生子的名.」(約三18)
\newline
\newline
「信子的人有永生,不信子的人得不著永生,神的震怒常在他身上.」(約三36)
\newline
\newline
「我實實在在的告訴你們,那聽我話,又信差我來者的,就有永生,不至於定罪,是已經出死入生了.」(約五24)
\newline
\newline
「我的羊聽我的聲音,我也認識他們,他們也跟著我.我又賜給他們永生；他們永不滅亡,誰也不能從我手裡把他們奪去.」(約十27-28)
\newline
\newline
(五)神在基督裡向在罪權之下的人,施行救贖工作
\newline
\newline
──神要救贖人脫離罪權
\newline
\newline
「人子阿,你要對以色列家說,你們常說,我們的過犯罪惡在我們身上,我們必因此消滅,怎能存活呢.你對他們說,主耶和華說,我指著我的永生起誓,我斷不喜悅罪人死亡；惟喜悅惡人轉離所行的道而活；以色列家阿,你們轉回,轉回罷,離開惡道,何必死亡呢.」(結三十三10-11)
\newline
\newline
「人仰望耶和華,靜默等候祂的救恩,這原是好的.......因為主必不永遠丟棄人.主雖使人憂愁,還要照祂諸般的慈愛發憐憫.因祂並不甘心使人受苦,使人憂愁.」(哀三26-33)
\newline
\newline
「以色列阿,你當仰望耶和華；因祂有慈愛,有豐盛的救恩.祂必救贖以色列脫離一切的罪孽.」(詩一三○7-8)
\newline
\newline
「主以色列的神,是應當稱頌的；因祂眷顧祂的百姓,為他們施行救贖；在祂僕人家中,為我們興起了拯救的角,(正如主藉著從創世以來,聖先知的口中所說的話.)拯救我們脫離仇敵......」(路一67-71)
\newline
\newline
「祂(救主)為我們捨了自己,要贖我們脫離一切罪惡,又潔淨我們,特作自己的子民,熱心為善.」(多二14)
\newline
\newline
「......我也知道,在我裡頭,就是我肉體之中沒有良善；因為立志為善由得我,只是行出來由不得我.......但我覺得肢體中另有一個律,和我心中的律交戰,把我擄去叫我附從那肢體中犯罪的律.我真是苦阿!誰能救我脫離這取死的身體呢?感謝神,靠著我們的主耶穌基督就能脫離了......」(羅七14-25)
\newline
\newline
神使那常在基督著的人,靠著祂的不斷拯救,得以脫離罪權,成為生活成聖的人
\newline
\newline
「那召你們的既是聖潔,你們在一切所行的事上也要聖潔；因為經上記著說,你們要聖潔,因為我是聖潔的.」(彼前一15-16)
\newline
\newline
「親愛的弟兄阿,我們既有這等應許,就當潔淨自己,除去身體靈魂的污穢,敬畏神得以成聖.」(林後七1)
\newline
\newline
「......惟有萬靈的父管教我們,是要我們得益處,使我們在祂的聖潔上有分.凡管教的事,當時不覺得快樂,反覺得愁苦；後來卻為那經練過的人,結出平安的果子,就是義......」(來十二10-11)
\newline
\newline
「但你們得在基督耶穌裡,是本乎神,神又使祂成為我們的智慧,公義,聖潔,救贖.」(林前一30)
\newline
\newline
「並靠著耶穌基督結滿了仁義的果子,叫榮耀稱讚歸與神.」(腓一11)
\newline
\newline
「......感謝神,因為祂從起初揀選了你們,叫你們因信真道,又被聖靈感動成為聖潔,能以得救.」(帖後二13)
\newline
\newline
「但現今你們既從罪裡得了釋放,作了神的奴僕,就有成聖的果子,那結局就是永生.」(羅六22)
\newline
\newline
(六)神在基督裡向在罪性之下的人,施行最後的救贖工作
\newline
\newline
──神要救贖人脫離罪性
\newline
\newline
「我是在罪孽裡生的；在我母親懷胎的時候,就有了罪.」(詩五十一5)
\newline
\newline
「惡人一出母胎,就與神疏遠；一離母腹,便走錯路,說謊話.」(詩五十八3)
\newline
\newline
「若我所作的,是我所不願意的,......乃是住在我裡頭的罪作的.」(羅七16-17)
\newline
\newline
「以色列阿,你當仰望耶和華......祂必救贖以色列脫離一切的罪孽.」(詩一三○7-8)
\newline
\newline
「願賜平安的神,親自使你們全然成聖；又願你們的靈,與魂,與身子,得蒙保守,在我主耶穌基督降臨的時候,完全無可指摘.那召你們的本是信實的,祂必成就這事.」(帖前五23-24)
\newline
\newline
神要那等候基督再來的人,靠著祂的最後拯救,得以脫離罪性,成為生命得榮耀的人
\newline
\newline
──等候基督再來的人
\newline
\newline
「你且去等候結局,因為你必安竭；到了末期,你必起來,享受你的福份.」(但十二13)
\newline
\newline
「義人死亡,無人放在心上；虔誠人被收去,無人思念；這義人被收去是免了將來的禍患.他們得享平安,素行正直的,各人在墳裡安歇.」(賽五十七1-2)
\newline
\newline
「論到睡了的人,我們不願意弟兄們不知道,恐怕你們憂傷,像那些沒有指望的人一樣.我們若信耶穌死而復活了,那已經在耶穌裡睡了的人,神也必將他與耶穌一同帶來.我們現在照主的話告訴你們一件事；我們這活著還存留到主降臨的人,斷不能在那已經睡了的人之先；因為主必親自從天降臨,有呼叫的聲音,和天使長的聲音,又有神的號吹響,那在基督裡死了的人必先復活.以後我們這活著還有留的人,必和他們一同被提到雲裡,在空中與主相遇；這樣,我們就要和主永遠同在.所以你們當用這些話彼此勸慰.」(帖前四14-18)
\newline
\newline
「我們卻是天上的國民；並且等候救主,就是主耶穌基督,從天上降臨.祂要按著那能叫萬有歸服自己的大能,將我們這卑賤的身體改變形狀,和祂自己榮耀的身體相似.」(腓三20-21)
\newline
\newline
「神救眾人的恩典,已經驗明出來,教訓我們除去不敬虔的心,和世俗的情慾,在今世自守,公義,敬虔度日；等候所盼望的福,並等候至大的神,和我們救主耶穌基督的榮耀顯現.」(多二11-13)
\newline
\newline
──在基督裡得榮耀的人
\newline
\newline
「我們講的,乃是從前所隱藏,神奧秘的智慧,就是神在萬世以前,預定使我們得榮耀的......」(林前二7)
\newline
\newline
「預先所定下的人又召他們來；所召來的人,又稱他們為義；所稱為義的人,又叫他們得榮耀.」(羅八30)
\newline
\newline
「我為這福音受苦難,甚至被捆綁,像犯人一樣；然而神的道不被捆綁.所以我為選民凡事忍耐,叫他們也可以得著那在基督耶穌裡的救恩和永遠的榮耀.」(提後二9-10)
\newline
\newline
「我想現在的苦楚,若比起將來要顯於我們的榮耀,就不足介意了.受造之物,切望等候神的眾子顯出來.因為受造之物服在虛空之下,不是自己願意,乃是因那叫他如此的.但受造之物仍然指望脫離敗壞的轄制,得享神兒女自由的榮耀.」(羅八18-21)
\newline
\newline
「神既是公義的,就必將患難報應那加患難給你們的人；也必使你們這受患難的人,與我們同得平安；那時主耶穌同祂有能力的天使,從天上在火燄中顯現,要報應那不認識神,和那不聽從我主耶穌福音的人.他們要受刑罰,就是永遠沉淪,離開主的面和祂權能的榮光；這正是主降臨要在他聖徒的身上得榮耀,又在一切信祂的人身上顯為希奇的那日子......叫我們主耶穌的名,在你們身上得榮耀,你們也在祂身上得榮耀,都照著我們的神並主耶穌基督的恩.」(帖後一6-12)
\newline
\newline
「......祂公義的審判日子來到,祂必照各人的行為報應各人；凡恆心行善,尋求榮耀尊貴,和不能朽壞之福的,就以永生報應他們；惟有結黨不順從真理,反順從不義的,就以忿怒惱恨報應他們.」(羅二5-8)
\newline
\newline
「小子們哪,你們要住在主裡面；這樣,祂若顯現,我們就可以坦然無懼；當祂來的時候,在祂面前也不至於慚愧.」(約壹二28)
\newline
\newline
「親愛的弟兄阿,我們現在是神的兒女,將來如何,還未顯明；但我們知道主若顯現,我們必要像祂；因為必得見祂的真體.凡向祂有這樣指望的,就潔淨自己,像祂潔淨一樣.」(約壹三2-3)
\newline
\newline
「基督是我們的生命,祂顯現的時候,你們也要與祂一同顯現在榮耀裡.」(西三4)
\newline
\newline
「原來那為萬物所屬,為萬物所本的,要領許多的兒子進榮耀裡去......」(來二10)
\newline
\newline
「在頭一次復活有分的,有福了,聖潔了；第二次的死在他們身上沒有權柄；他們必作神和基督的祭司,並要與基督一同作王一千年.」(啟二十6)
\newline
\newline
「我聽見有大聲音從寶座出來說,看哪,神的帳幕在人間；祂要與人同住,他們要作祂的子民,神要親自與他們同在,作他們的神；神要擦去他們一切的眼淚；不再有死亡,也不再有悲哀,哭號,疼痛,因為以前的事都過去了.」(啟二十一3-4)
\newline
\newline
神是為我們施行諸般救恩的神；人能脫離死亡,是在乎主耶和華.(詩六十八20)
\newline
\newline
3．經題:神的救贖工作
\newline
\newline
經文:賽四十三1；詩一三○3-4,7-8；約五17,十四10；林前一30；腓一28
\newline
\newline
──神的救贖工作,正如栽培人的工作.(約十五1；羅十一23,17；林後五5)
\newline
\newline
(一)神在基督裡向在罪惡之下的人,(羅三23,9；徒八23)施行救贖的工作.(約十四10)叫悔改相信的人,(徒二十21,二十六18)得以歸在基督裡,(林前一30)成為得救的人.(腓一28；提前一15-16；太一21；加五1)
\newline
\newline
(二)神在基督裡向在定罪之下的人,(羅三19,五18)施行救贖的工作.叫悔改相信的人,得以歸在基督裡,成為得稱義的人.(羅三24-26)
\newline
\newline
(三)神在基督裡向在罪刑之下的人,(羅一32,六23)施行救贖的工作.叫悔改相信的人,得以歸在基督裡,成為得永生的人.(羅六23；約十七3；約壹五11-12)
\newline
\newline
(四)神在基督裡向在罪權之下的人,(羅七14-24)施行救贖工作.叫悔改相信的人,得以歸在基督裡,成為生活成聖的人(帖前五23；約十七17)
\newline
\newline
(五)神在基督裡向在罪性之下的人,(羅七17；詩五十一5)施行救贖工作.叫悔改相信的人,得以歸在基督裡,成為得榮耀的人.(羅八30；但十二3；太十三43)
\newline
\newline


\newpage

\section{第50屆~港九培靈會~陳終道~第一講}
\label{sec:714}
\textbf{聖經奇妙的影響力}
\newline
\newline
連結: \href{https://www.hkbibleconference.org/session-message/view/714}{\texttt{https://www.hkbibleconference.org/session-message/view/714}}
\newline
\newline
\hyperref[sec:734]{< < < PREV SERMON < < <}
~
\hyperref[sec:toc]{[返目錄]}
~
\hyperref[sec:715]{> > > NEXT SERMON > > >}
\newline
\newline
感謝主!兄弟今天有機會與各位一同研究神的話語.我要講的題目是「怎樣研讀聖經」.這信息將由淺入深,並先由基本的問題,漸次深入.
\newline
\newline
很多基督徒不愛讀經,連禱告也只在用膳時形式化的念幾句.有些基督徒只不過在主日聽聽牧師講道而已,不讀經,這是個嚴重的問題!可能我們不知道讀經的好處,也未曾嘗到主話的甘甜；也可能我們想讀經,卻不知如何入門.常有基督徒問我:「應當怎樣讀聖經?」有的人對聖經發生懷疑,所以我必須題到聖經的可靠.有的人對解經方面發生爭辯；有的反對靈意解經；有的濫用靈意解經；到底怎樣才是正確的解經,我也將會題到這方面的問題.很多基督徒對舊約的預表發生困難,預表是否可以隨意解釋呢?沒有一定的原則?諸如此類我將一一題及.有些基督徒對聖經上的比喻感到困難,因為教義上的爭論也是由比喻開始的；我將會題到比喻的讀法.如果時間許可,也將會題到怎樣讀預言.許多基督徒讀經,覺得前後矛盾,就是因對聖經的寫法沒基本的認識.有人說聖經這麼多卷,幾時才讀得完；現在有很多書籍,是教人如何讀聖經,到底以何種方法讀經為最好呢?我希望在這次聚會中,能夠提到這些問題.
\newline
\newline
聖經是一部奇妙的書,影響世界最大；聖經的偉大是沒有任何書籍可比擬的.
\newline
\newline
聖經對個人的影響力(提後三15-17；彼前二1-2；弗六17；來十7)這幾處經文都題到聖經對個人,有奇妙的影響.聖經使人有得救的智慧.它與世上的智慧是不同的.因為聖經的智慧,能把人引到得救的那邊；世界的智慧,卻引人到滅亡那邊.人無論赴奮興會得救,或赴佈道會得救,或看屬靈的書受感動,或看到福音單張得救,他們都要根據聖經的話語.
\newline
\newline
基督徒憑著聖經的話,確知自己已經得救.有人說得救之後心裡喜樂,那麼,心裡不快樂的人是否得救?聖經說:「信子的人有永生,不至被定罪,是已經出死入生了.」我們若認自己的罪,神就根據祂的公義和信實,赦免我們一切的罪,洗淨我們的不義；所以我們得救是因著聖經,我們知道自己得救也是由於聖經.但是世界的智慧叫人驕傲,世人稍有才幹知識,就自高自大.
\newline
\newline
奧古斯丁是位聞名的教士,他未信主之前,行為極度放蕩.有一天,他聽到聖經的話:「黑夜已深,白晝將近；我們就當脫去暗昧的行為,帶上光明的兵器.行事為人要端正,好像行在白晝；不可荒宴醉酒；不可好色邪蕩；不可爭競嫉妒.」(羅十三12-13)因此,他受感動就得救了.
\newline
\newline
有個煙草公司的職員到非洲去,看見當地土人讀聖經,就笑他,說:「這已是落後的書,我們那邊沒有讀了.」土人定睛看他,說:「幸虧此書對我們並不落後,否則,今天你已被吞下我的肚腹裡了.」可見聖經實在有奇妙的影響力量,能改變人.
\newline
\newline
聖經使基督徒的靈命長大.彼得說:「要愛慕那純淨的靈奶,像纔生的嬰孩愛慕奶一樣.」奶怎樣叫嬰孩長大,神的話也叫我們的靈命長大.我們經常吃食物才能維持生命；我們天天讀聖經,靈命才能長大.摩西上山與神面對面,他的臉上發光而不自知；他下山時,以色列人卻發現他臉上發光.我們天天讀經,就像被神的榮光照耀,不知不覺生命改變並長大.聖經能改變更新人的心思意念；靈命長大了,就丟棄了孩童時代的愛好.
\newline
\newline
教會是個生命的團體,要生命長大成熟,才能生養兒女；教會要增長,就必須生命長大,要追求靈性長大,主才把得救人數加給教會.今天的教會,好像個未長成的女孩,而兒女已眾多；雖然勞碌辛苦,對兒女仍感照顧不週.
\newline
\newline
聖經是對付仇敵的屬靈兵器,基督徒若沒有神的話,就沒有屬靈兵器抵擋試探.有些基督徒羨慕做個得人的漁夫,卻自嘆不如,這是因為沒有屬靈的兵器,不讀聖經,神的話不在你裡面的緣故.(弗六17)說:「要拿著聖靈的寶劍,就是神的道.」「拿著」,就可以隨時準備,有寶劍還要持之在手,不要收藏起來；如果你一週或偶然讀次聖經,當然沒力量勝過試探,也沒力量作得人的漁夫.掃羅時代有場可笑的戰爭,當時掃羅吹角,吹號準備和非利士人打仗；但全國僅掃羅和他兒子約拿單各拿一刀,而非利士的戰車卻有六千.基督徒要勝過試探,必須裝備自己,接受訓練.聖經稱基督徒是基督的精兵,精兵應有良好的裝備和訓練,且能善用裝備,具有堅強的心志,然後可以出去打仗.今日的基督徒,雖然有了必勝的武器,只可惜還沒有接受訓練,未能善用武器,因懶惰成聖,便輕忽了寶貝的武器.
\newline
\newline
聖經是叫我們明白神旨意的最好根據,希伯來書告訴我們:「祂的事在經卷上已經記載了.」聖經能助我們更多認識主耶穌基督.神賜聖經給我們,叫我們透過聖經認識祂.明白祂的旨意,了解祂的作為,愛好和恨惡；知道祂的本性是何等善良聖潔,知道祂行事的法則.我們從聖經認識,了解祂,然後才能在凡事上接受恩膏的教訓,使聖靈藉著聖經來教導我們.
\newline
\newline
聖經中記載兩個絕對不同的使徒──彼得和猶大,猶大賣耶穌,後來似乎悔悟把錢拿回來.彼得三次不認耶穌,甚至發誓起咒；這樣的罪可說不輕,但是後來悔改,被主使用,在使徒中起了帶頭作用.彼得和猶大到底有何不同呢?當彼得三次不認耶穌時,主轉過身來看彼得,他便想起主對他所說的話......就出去痛哭.彼得想起主的話,猶大卻沒想起主的話；猶大沒有信心,彼得卻有信心,因主的話在他心裡.我們要將主的話藏在心中；聖經能叫我們想起主的話,幫助我們明白神的旨意,接受聖靈的引導；聖經的話如火車軌道,聖靈如火車頭,必先有軌道而後才能行車.聖靈藉著神的話在我們心中作工；基督徒不讀經是莫大的損失!
\newline
\newline
聖經能叫你得救,叫你生命長大,並叫你勝過試探,叫你接受聖靈的旨意,明白神的旨意；若沒有聖經的話語在你心中,你就無能為力,得不到神的幫助.聖經對個人具有奇妙的力量,對整個世界亦然.
\newline
\newline
今日世界普遍趨向自由平等的思想,顯然是由於聖經的影響；聖經給我們看見,我們在神面前都是罪人的地位,都是憑主耶穌的恩典得救；不論貧富貴賤,同樣站在罪人的地位；但蒙恩以後都是站在得救的地位；不分種族,國家,為主,為奴性別,階級,在耶穌基督裡都成為一了.最早反對剝削勞工的是聖經,雅各書五章公開責備富人剝削工人,在當時還未有人敢反對.美國解放黑奴名聞的總統林肯說:「聖經是部偉大的書,是神賜給人最好的禮物.」可見他的思想深受聖經影響.歷史上許多不平等及封建的制度,都是聖經所反對的.出名的宣教士威廉克里在印度傳道時,竭力反對當時一種風俗,要與亡夫陪葬.中國古時女人的纏足,受到傳教士提摩太劇烈的反對,中國最初的女子學校也是教會創辦的；中國女子能得自由平等解放,都是聖經之賜；若非靠聖經傳入中國,現在女子走起路來,恐怕還是三寸金蓮哩!
\newline
\newline
今日有人高唱:「基督教神學本色化.」意思是要遷就本地的風俗文化.這樣,還有福音可傳嗎?那些為傳福音而來中國,且曾作出許多貢獻的西教士；如果入鄉隨俗,入港隨灣,那還有什麼意思呢?聖經不是文化侵略,把西方的文化帶來東方勉強人接受；無論西方或東方,其文化有違反聖經者,均受到聖經的影響而改變.聖經非但改變東方的文化,也要改變西方的文化；因人的文化很多時候是與神敵對的.神的僕人,曉得神的道高過人的思想；因為神的道不會受人的思想而有影響的.從前多妻制度普遍,但聖經早已主張一夫一妻；現在世界各國已公認納妾為不合法了,這就是聖經思想的影響力了.一九一九年中國發起五四學生運動.開始時,只不過一班大學生因為不滿政府的外交政策,而示威遊行.後來演變為新文化運動,這運動的重點,是主張使用語體文,反對原有的「之乎者也」的文言.不過語體文和合本聖經,是在同年出版的.早在廿七年前,已著手進行繙譯文理為語體文了；當時除聖經外,別無其他語體文的巨著流傳於民間.所以當時連非教會學校,也採用聖經的四福音為課本；鄉間的文盲,讀白話聖經兩年就認識了.從前傳福音開始時,多辦識字班,連老農婦讀了兩年白話聖經,便識字了.聖經實在影響中國的文化.可是人並不歸功於聖經；其實和合本聖經對白話文貢獻殊大.
\newline
\newline
世界很多有名望的人,也受聖經的影響,有人問英女皇維多利亞說:「有什麼秘訣可以使英國成為強盛大國?」女王舉起聖經說:「聖經」.可惜現代許多基督教國家已沒有基督教了,以致國家日漸衰弱.
\newline
\newline
偉大科學家牛頓是個熱心的基督徒,他寫了啟示錄和但以理兩卷書的註解.這兩卷書原是最令人難以接受的,因其中神奇的事太多.但許多偉大的科學家,都不反對聖經.而且愛慕聖經.
\newline
\newline
你和聖經的關係如何?有沒有屬於你自己的聖經?有沒有好好讀聖經?抑或把你的聖經置於書架上呢?因你沒有讀聖經,所以聖經的福份沒有臨到你.
\newline
\newline
有些地方的基督徒非常渴慕聖經,據說他們寫信給遠東福音廣播電台,要求廣播聖經節時逐字慢讀,以便他們抄錄,每人記一節然後合攏成章.我們在這裡,讀聖經十分方便,為何不好好的讀聖經呢?有一日,我們將與他們相會在天,那時我們就會覺得何等慚愧萬分!聖經是否成為我們的裝飾品呢?抑或能放在心中成為生命的力量和滋潤呢?
\newline
\newline
求主使我們愛慕祂的話,來赴會,不但聽人所說的話；同時也愛慕神的話,把神的話藏在心中,就不至於得罪主了.
\newline
\newline


\newpage

\section{第50屆~港九培靈會~陳終道~第二講}
\label{sec:715}
\textbf{聖經的可信}
\newline
\newline
連結: \href{https://www.hkbibleconference.org/session-message/view/715}{\texttt{https://www.hkbibleconference.org/session-message/view/715}}
\newline
\newline
\hyperref[sec:714]{< < < PREV SERMON < < <}
~
\hyperref[sec:toc]{[返目錄]}
~
\hyperref[sec:716]{> > > NEXT SERMON > > >}
\newline
\newline
在這末世的時代,魔鬼要極力誘人不信聖經的話,甚至在教會的圈子裡,也有人懷疑聖經的話.有人問:我們是否能從科學立場證明聖經是神的話?又能否從聖經之外來證明聖經是神的默示?這些都是不合理的要求,正如顯微鏡是研究細菌不可缺的儀器；怎能要求研究細菌的人,不用顯微鏡呢?聖經是唯一記載神默示的書,若否認聖經的價值,又想曉得神的默示,那是不可能的.我們豈能限制醫生看病,不根據醫學原理,而採用會計學呢?各種科學必有研究它的工具和原理.照樣,我們也必須根據聖經,才能認識神的話.
\newline
\newline
聖經的權威和可信,最叫人發生動搖的,就是有人認為今日是科學進步的時代,人不易受騙；也就是說從前科學落後,人易受騙,所以才相信聖經.如果聖經是本騙人的書,根本就經不起時間的考驗,惟有真實才能經得起時間考驗.從前科學落後,是現代人的想法；當那時代的人照樣認為自己是進步的.現在的人可與從前的人比較,但也無法和將來的人比較.聖經完成至今約兩千年,當然那時的人無法與現代的人比較；如果以兩千年前的眼光來看現在,簡直把人看成神仙了.即使科學進步可以搖動聖經的權威,也絕不可能歷時至今才搖動的.
\newline
\newline
根據大英百科全書記載,埃及的金字塔建築於主前二至四千年之間,庫弗金字塔是用二百三十萬塊,每塊兩噸半重的石頭建造的；當時未有起重機也沒有直昇機,用何種方法堆成金字塔呢?直到如今科學家尚且覺得這是值得研究的問題.古時也可用藥物保存屍體幾千年不腐,這是他們的科學；不過古時科學專供帝王和有權勢的人享受,而今科學進步多為普遍大眾享受而已.我在四川讀書時,曾參觀博物院,其中有漢朝皇帝用的銅製洗臉盆,盆底雕刻三條魚,盆兩邊各有扶手；摩擦之,水中發出悠美的響聲.繼續摩擦,水呈縐紋,聲響悅耳,並由三魚口噴水；這是關於物理學震動音波的原理,雕刻的準確恰使三條魚口中能噴出水來,真是高度深奧的技巧!所以不可藐視古代的科學,漠視古人較今人愚拙.聖經不但沒有被科學進步所搖動,反而在科學不斷進步中,神藉祂的僕人把祂的默示寫下.雕
\newline
\newline
摩西是一位充滿埃及學問的學者,是聖經首五卷五經的作者.所羅門王是一位著名的詩人,也是生物學家,他能講論花草樹木,鳥獸蟲魚.由此可知聖經的作者,有當代的科學家,有淵博的學者.科學進步並不能搖動聖經,因這是兩種不同的範圍；科學研究物質方面；聖經卻告訴我們神的啟示和神的話語.
\newline
\newline
基督徒對神的話語,應有堅固的信心,才能從屬靈的經歷,體驗神話語的能力；如果沒有信心,聖經的話對你就不產生力量.因為人非有信,不能得神的喜悅；到神面前的人,必須信有神,而且信祂必賞賜尋求祂的人.如果存著不信的心來讀聖經,就得不著聖靈的教導.
\newline
\newline
雖然聖經曾被科學家懷疑,特別是關於宇宙的來源.聖經說天地是神創造的,科學家則謂天地是自然演進而成的.聖經說天地是神創造的,早有記載,且是最後的定論.至於科學家所說的,是否可作最後定論呢?若科學家尚未有最後定論,便沒有資格和聖經挑戰.研究科學者有人說:「非研究科學者勿批評科學.」然而他們是否能以此語自作反省:「非研究聖經者,勿批評聖經.」凡批評聖經不可信的人,多是不勤讀聖經的人.我曾從事文字工作,發現那些喜歡批評書的人,都不是存心閱書的人.他們隨便翻閱即作出批評,而真正讀書的人,如有所批評亦屬善意.如果科學家認為宇宙的創造問題已有答案,膽敢宣佈進化論是最後的結論；那麼則儘可向聖經來挑戰!正如分兩組比賽回答問題,一組已有答案；另一組不但未有答案,且譏笑別組的答案錯誤,豈不可笑嗎?所以那些根據進化論批評聖經的人應質問科學家:「進化論是否宇宙的答案?」
\newline
\newline
有一位生物學家愛司丁康克里說:「如果生命的來源憑偶然的可能性,這種惑然的情形；則有如印刷所經一場爆炸之後,可以印出了一本完整的字典來.」你信嗎?
\newline
\newline
有人說,是否聖經因流傳時間過久,傳播範圍太廣而發生錯誤呢?這是個很好的問題,但我們不要忽略一點:神是否早已估計聖經在傳播中會發生錯誤?我們不應把傳播上的錯誤歸咎於聖經,傳播之錯並非聖經之錯；聖經告訴我們,神所默示的這部聖經,非屬偶然湊集的書,而是神計劃中要賜給人的書.亞伯拉罕時還沒有聖經,神對他說:「萬國都必因你的後裔得福.」保羅說:「聖經既然預先看明,神要叫外邦人因信稱義,就早已傳福音給亞伯拉罕.」(加三8)神對亞伯拉罕所應許的話,正是神啟示聖經要給我們的一部份,那時聖經的腹稿已在神那裡,所以神對亞伯拉罕所說的話,就是後來聖經的話.這節經文說明舊約聖經的啟示,早在神的計劃中.(羅十五14)說明在舊約神啟示之目的,是要教訓新約的信徒:「從前所寫的聖經,是為教訓我們寫的,叫我們因聖經所生的忍耐和安慰,可以得著盼望.」(彼後三15-16)彼得證明保羅所寫的新約書信,是出於神的靈感.
\newline
\newline
聖經是神預早計劃要完成的一部書,在神啟示給人的時候,是有目的,有主題的.因此,聖經是一貫性的,不因人在傳播上的錯誤而受損害.舊約許多預表和先知的預言,都有它的中心點是指向耶穌；新約四福音記載耶穌的生平,耶穌如何降生,受死,復活,正是舊約先知所預言的,正是舊約預表所顯明的.然後在新約中,解釋這位耶穌基督,如何照著神的旨意而來,如何死而復活,將再來接我們；解釋新約照著舊約預言,完成神的旨意；所以萬一有傳播上的錯誤,聖經也不致受影響.
\newline
\newline
至於抄寫上的錯誤,不如我們想像的那麼嚴重；因為從前抄寫聖經者,都是很虔誠的人,他們逐字朗誦抄寫；再作校對.一頁如有四個錯字,全頁作廢重抄,抄完一卷如發現錯誤,全卷作廢重抄；他們工作虔誠慎重,錯誤的可能性絕少.一旦因人的軟弱,抄卷有錯；但古時文字傳播工作進行緩慢,可以及時追回修改.聖經傳播開始時,神曾告誡以色列人要謹慎教訓他們的兒女；所以作君王的,也有聖經抄本置於身旁.
\newline
\newline
有人以為聖經是由教會審定,這樣豈不是教會比聖經更有權威?我們應注意,聖經從神默示至完成約經一千五百年之久；神容許祂的話語經過漫長的時間才完成,是要讓它接受時間的考驗；聖經中很多歷史都是具有目的之挑選；因為「從前所寫的聖經,都是為教訓我們寫的.......」神將那些可表明祂旨意的歷史,挑選載於聖經,是經過長時間,從人的歷史中選出的.聖經有超時代的性質,所含原則適用於任何時代.神並沒有把祂的啟示在一個時代中,全給一個人；而是在每個不同的時代,按著祂的旨意啟示給祂所挑選的人.把歷代集合起來的,就成了神完全的啟示.天地萬物是神創造的,但並非交給一人管理；而是交給亞當和他的後裔,全世界的人類管理.耶穌基督的救贖,是他一人成功的；而傳揚救贖的恩典,乃所有信徒之責；凡蒙恩得救的人,就要共同傳揚救恩.聖經是聖靈所啟示,非由一人領受,乃給神所揀選的眾僕人一同領受.三位一體的神所創造,所完成的救贖,所啟示我們的聖經,是要眾信徒一同完成共同領受的.神容許祂的話語經長期考驗,是要顯明那些假冒神默示的經書都經不起考驗.保羅在(帖後二1-2)提醒當時聖徒:「弟兄們,論到我主耶穌基督降臨,和我們到他那裡聚集,我勸你們,無論有靈,有言語,有冒我名的書信,說主的日子現在到了,不要輕易動心,也不要驚慌.」使徒提醒當時的信徒,要分辨清楚使徒書信的真假.彼得後書三章也題到當時有假冒的經書出現,提醒信徒要分辨何為神的啟示.約翰也在他的書信勸勉信徒說:「一切的靈不可都信.」又提醒信徒勿輕易相信靈感.初期教會使徒,提醒信徒小心認識神的啟示.教會中有屬靈的敬虔,又對聖經專門研究有經驗的人；他們經長期分辨經書,且細聽敬虔者的朗誦,然後開大會把真假經書分別出來；所以如果說加太基會議是審定聖經,那就錯了!那時教會不過是經過長時間,來辨認經書的真假而已.既經辨認清楚,就把神的話彙集起來；教會只是辨認神的話,教會無權審定神的話.聖經的權威僅屬聖經本身,絕非人所能加給的；因為神的話安定在天,直到永遠.
\newline
\newline
我們當愛慕神的話,知道神的話是可信的,將神的話藏在心裡,就不至於得罪神,就可勝過試探,可作得人的漁夫.神的話是我們腳前的燈,指引我們前面的路.最可怕的就是心中沒有真理的基督徒,作事憑自己的思想意念,憑世界的智慧.
\newline
\newline
求主提醒我們,讓我們寶貴神的話語!
\newline
\newline


\newpage

\section{第50屆~港九培靈會~陳終道~第三講讀聖經應有的態度}
\label{sec:716}
\textbf{第三講讀聖經應有的態度}
\newline
\newline
連結: \href{https://www.hkbibleconference.org/session-message/view/716}{\texttt{https://www.hkbibleconference.org/session-message/view/716}}
\newline
\newline
\hyperref[sec:715]{< < < PREV SERMON < < <}
~
\hyperref[sec:toc]{[返目錄]}
~
\hyperref[sec:717]{> > > NEXT SERMON > > >}
\newline
\newline
以賽亞書 50:4-50:4
\newline
\begin{longtable}{cl}
\hline
\hline
章節 & 經文 (和合本修訂版)\\
\hline
50:4 & \begin{tabularx}{0.7\textwidth}{X} 主耶和華賜我受教者的舌頭，使我知道怎樣用言語扶助疲乏的人。主每天早晨喚醒，喚醒我的耳朵，使我能聽，像受教者一樣。 \end{tabularx} \\ \\
[1ex]
\hline
\hline
\end{longtable}
在未講怎樣讀聖經之前,我們應該先知道讀聖經應有的態度,態度不對,可能使方法失效.有的人學問很好,但因處事態度傲慢,以致工作效果很低；有的人學問稍差,但處事態度謙虛謹慎,所以效果不錯.比如,駕駛的技術固然要緊,修養也很重要；有人開車失事,並非技術不高,乃因缺乏修養.如果我們在神面前的事奉態度錯誤,一切事奉就打了折扣.
\newline
\newline
讀聖經應有的態度:
\newline
\newline
(一)當謙卑受教
\newline
\newline
先知以賽亞說:「主耶和華賜我受教者的舌頭,使我知道怎樣用言語扶助疲乏的人；主每早晨提醒,提醒我的耳朵,使我能聽,像受教者一樣.」以賽亞是個先知,他是為神說話的；他每早晨必先在神面前接受神提醒的聲音,安靜聽神的教訓.他先聽,像個受教者；然後他出去講受教者所說的話；因他深恐說話像個沒受過屬靈教育的人,所以他禱告求神賜他受教者的舌頭.基督徒若沒有存心接受神的教訓,就在讀聖經時,甚麼也得不到.
\newline
\newline
怎樣算是受教?就是接受聖經話語的光照之後；改變自己,有錯認錯,有罪認罪,該被神修剪就因神的話接受修剪.有時聖經的話指正人的錯處,人非但不受教；反而抗拒聖經,挑剔及批評,為了保護自己.在教會中挑撥是非,其實這人正落在軟弱中.我們讀聖經,當在神面前赤露敞開,不遮掩自己的錯,不保護自己的軟弱；讓整個人照在神的光中,指出錯誤,讓神修理.
\newline
\newline
(二)要有飢渴慕義的心
\newline
\newline
先知耶利米說:「耶和華萬軍之神啊!我得著你的言語,就當食物喫了；你的言語,是我心中的歡喜快樂；因為我是稱為你名下的人.」(耶十五16)古時神所使用的先知,雖然奉神的名說話,多年傳道；對神的話何等的愛慕,因此,神話語的信息臨到他們.我們當反省一下,對神的話是否存不在乎的態度；果真如此,神的話向你就不發出亮光.食物是人所必需的,耶利米把神的話當作必需的食物；但有的基督徒讀聖經猶如進補,只是時而有之.耶利米成為偉大的先知,因他對神的話飢渴,這樣的人才能支取(太五6)耶穌的應許:「飢渴慕義的人有福了,因為他們必得飽足.」真正愛慕神話語的,必能從神的話語中得到亮光.很多方法都從「愛」來的,有「愛」就有方法.求主給我們愛慕祂的話語!
\newline
\newline
(三)當依靠聖靈
\newline
\newline
「只等真理的聖靈來了,祂要引導你們明白一切的真理.......」
\newline
\newline
除了神的靈,沒有人知道神的事,所以讀聖經一定要依靠聖靈.神既給我們智慧聰明,難道不用?比方我們有文字的學識,難道不用嗎?聖經乃是說不要依靠屬世的學問和聰明；因為聖經基本上非人間的作品,所以不能憑人間的學識聰明智慧去研究,必須在人的智慧之外,依靠聖靈的能力,聖靈的教導,才能明白.其實這些話不需我們來強調；可是我發現許多人,讀聖經倚靠自己的頭腦,我們所贊成的一個理論和我們實行的有時竟然脫節,只在道理上接受.我們讀聖經不能倚靠自己,要倚靠聖靈的教導；存著仰望倚賴聖靈的態度,才能夠從神的話語摸到門路.凡是技術都有它的原理,說起來極為簡單；原理只有一個,但實地製作則各人的成績不一定相同.因為有的人已經摸到門路,有的沒有摸到.例如同一菜譜,製出的菜卻因人而異；因為有人能體會其中技巧.讀經的體會必需聖靈來教導,好像學彈鋼琴,老師對每個學生教法一樣,但各人對琴鍵的感觸,這方面就無法傳授；必要憑自己的體會,而這體會就是決定成功與否的因素.
\newline
\newline
(四)要願意為主使用
\newline
\newline
讀聖經要有飢渴慕義的心,要存受教者的態度；這些固然是為自己,然而也是為了主工作的需要.有的人讀經專為別人讀,聽道專為別人聽；我們讀聖經要先為自己讀,然後領受.不但自己受教導,也要為神用.主耶穌說:「凡我所吩咐你們的,都教訓他們遵守.」我們必須先接受主的吩咐,接受神話語的修理教導,然後能教導別人,叫別人得益.有些人在屬靈的追求上,專為自己的享受,只為自己領受而不與別人分享,這樣的基督徒很難熟讀聖經.讀和用要同時並進,才能熟練；背熟的經節如果不常用,也會忘記的.
\newline
\newline
保羅說他自己是執事,也是神奧秘事的管家.管家就是管理家務,如膳食和打掃屋子等.保羅所說:「奧秘事」是指福音真理的奧秘；他是奧秘真理的管家.基督徒受了委托,神把救贖的真道,聖經真理中重要的信息,教導我們,交托給我們,叫我們成為真理的管家；所以我們要先明白聖經真埋,善於運用聖經真理.主耶穌在三個僕人的譬喻中,稱讚那領五千和領兩千的僕人,為良善忠心的僕人；而責備那領一千的僕人,說他是又懶又惡的僕人；因他沒有運用從主所領來的一千兩銀,他浪費了機會.今天福音的真理是委托給我們的,這部聖經是托負給我們的,不但要我們讀,還要我們應用,這是我們的責任,是主給我們的托負；可是我們連這個都不知道,既不讀又不用；正是又惡又懶.主提醒我們,叫我們要甘心樂意為祂；讀聖經並應用所讀的聖經,然後神的話在我們心中才活潑有力量.
\newline
\newline
(五)讀聖經要有恆心
\newline
\newline
徒十七11說:「天天考查聖經.」恆心乃是續而不斷；恆心不是熱心,有的人每天讀經廿章,但只讀兩天；恆心也非雄心,恆心乃是天天一樣,日日如此.正如母親餵養孩子,天天不斷地餵養,在不知不覺中,自然長大.庇哩亞人天天考查聖經,結果他們靈性就有長進.
\newline
\newline
(六)要有信心
\newline
\newline
「人非有信,就不能得神的喜悅,因為到神面前來的人,必須信有神,且信祂賞賜那尋求祂的人.」(來十一6)我們讀經必須有信賴神的心,信賴神的話語,信賴神的話語必不會錯；如果沒有完全信任神話語的心,就得不到聖靈的教導.有的人問了問題,為其解答後還是不相信,仍然我行我素；其實他是存心試探,人若存不信的心讀經,必得不到聖靈的教導.疑問分信與不信兩種,耶穌在世上時,法利賽人,撒都該人,律法師,常問耶穌問題；耶穌對付這等人的方法,就是反問他們；因為他們並非存心明白問題.耶穌的門徒也常問耶穌問題；但他們的疑問非出於不信,而是因他們不明白.學生問老師是「請問」,老師問學生是考問；敵人為抓把柄所發問題是「試探的問」.我們讀經如有不明之處.是期待明白之問；可在神前禱告,聖靈必教導我們.
\newline
\newline
有人說,科學要存疑問的態度,例如牛頓看見蘋果從樹上掉下,因心裡有疑問,後來才發現地心吸力；但我們可說牛頓的發明,不是因疑問而是信心.他相信蘋果憑空而墜,其中必有原因,結果找出道理,他起初的疑問是在信心圈子裡.
\newline
\newline
但以理讀聖經的態度(但九2-3)但以理從書上得知耶和華的話,臨到先知耶利米,論耶路撒冷荒涼的年數,七十年期滿.他完全相信這話,且因這話大受感動,便禁食,披麻蒙灰,認罪,為他的國家禱告,他並非赴了奮興會,他是讀聖經,讀後有立刻遵行的心.他讀聖經的態度,是我們應當學習的；聖經的話影響但以理的內心,他立即採取行動；這是個蒙福基督徒讀經應有的方法.有的人讀聖經,猶如無讀一樣,毫無影響.
\newline
\newline
讀聖經最重要的,還不在於讀後所得的教訓,乃在讀了聖經的話,行事為人有所改變；我們信了主,心思意念應有更新,靈性需求繼續長大；好像新生的嬰孩,雖有生命,但距離父母的心意尚遠；所以我們要心思繼續更新,靈命繼續長大,讓我們的心思與天父更加接近,使我們的思想被聖經感化,心意與神更靠近.這樣,我們就容易明白聖經,體會神的心意,在世上真正為神而活.
\newline
\newline
從(但九2-3)這段經文,雖是一鱗半爪；但亦可見神的話,在但以理心裡有權威,所以他說的話也有權威.
\newline
\newline
求主恩待!叫我們非把聖經當作裝飾品,熟讀聖經非為在人面前誇耀,博取人的尊敬；乃是願意來到神話語面前,好像站在鏡子前.在光中,被照耀,被修理,在神榮光反照之下,漸漸有神的形像；新的生命漸漸長大,心意更新而變化,察驗何為神的善良,純全可喜悅的旨意.
\newline
\newline
求主保守我們的心,讓我們存著受教者的心來讀祂的話,來用祂的話；有飢渴慕義的心,依靠聖靈；願意恆心信賴神的話,決心遵行神的話.
\newline
\newline


\newpage

\section{第50屆~港九培靈會~陳終道~第四講}
\label{sec:717}
\textbf{靈修讀經和熟讀聖經}
\newline
\newline
連結: \href{https://www.hkbibleconference.org/session-message/view/717}{\texttt{https://www.hkbibleconference.org/session-message/view/717}}
\newline
\newline
\hyperref[sec:716]{< < < PREV SERMON < < <}
~
\hyperref[sec:toc]{[返目錄]}
~
\hyperref[sec:718]{> > > NEXT SERMON > > >}
\newline
\newline
申命記 6:4-6:9
\newline
\begin{longtable}{cl}
\hline
\hline
章節 & 經文 (和合本修訂版)\\
\hline
6:4 & \begin{tabularx}{0.7\textwidth}{X} 「以色列啊，你要聽！耶和華－我們的神是獨一的主。 \end{tabularx} \\ \\ \relax
6:5 & \begin{tabularx}{0.7\textwidth}{X} 你要盡心、盡性、盡力愛耶和華－你的神。 \end{tabularx} \\ \\ \relax
6:6 & \begin{tabularx}{0.7\textwidth}{X} 我今日吩咐你的這些話都要記在心上， \end{tabularx} \\ \\ \relax
6:7 & \begin{tabularx}{0.7\textwidth}{X} 也要殷勤教導你的兒女。無論你坐在家裡，走在路上，躺下，起來，都要吟誦。 \end{tabularx} \\ \\ \relax
6:8 & \begin{tabularx}{0.7\textwidth}{X} 要繫在手上作記號，戴在額上作經匣； \end{tabularx} \\ \\ \relax
6:9 & \begin{tabularx}{0.7\textwidth}{X} 又要寫在你房屋的門框上和你的城門上。 \end{tabularx} \\ \\
[1ex]
\hline
\hline
\end{longtable}
這段經文題醒當時的以色列民,要畢生遵行神的誡命,要盡心,盡意,盡力愛神.摩西教導他們無論坐下,躺臥,走路,在家或在外在城門上,都要把神的話語明顯地擺在眼前；要在各方面接觸神,要記住神的話.這段經文也教導我們要愛神,要將神的話牢記在心.
\newline
\newline
我常說基督徒有兩件終身大事:(1)是找配偶.(2)是讀聖經.讀聖經是我們開始作基督徒,直到見主面一輩子的事.
\newline
\newline
怎能使讀經有效果呢?當然方法是要講究的；但基本上,首先我們必須有愛神的心,有愛就有方法,不愛即使有方法也不會用.一個愛慕神話語的基督徒,總能找出讀經的方法；相反地,不愛慕神話語者必拒方法於千里.我曾發現有些基督徒專找讀經的方法；但卻不讀經,實際上,這一類的人是沒有方法的.
\newline
\newline
基督徒讀經可分三大類,或者說是有三個不同之目的:第一,為靈修:基督徒每天應有靈修,為靈修而讀經之目的與普通目的不同；因為靈修目的是要與神親近.第二,為熟讀聖經:不同目的當然要用不同的方法.第三,要求明白:不但讀聖經,還要明白聖經,所以不能快讀,要謹慎地讀,且須閱參考書,不能限定每天讀多少章；如果這樣,不免匆忙,影響與神親近.我們活在世上,有很多叫我們與神遠離的事物.神是靈；而我們還活在肉身裡面,肉身纏繞,整天都是接觸物質與肉身有關的事.基督徒接受了救恩,是神的兒女,與天上的父應發生感情；我們的神非同世人所造的假神；真神愛我們,我們愛祂,這就是感情.靈修的作用,就是使我們與肉眼不能見的神增加真正的感情.在靈修之中,禱告,唱詩,默想和讀經完全配合在一起.在神面前挨近祂,反覆思想聖經神的話語；傾心吐意向神禱告,唱詩來發出讚美的聲音；這樣便達到靈修真正之目的.欲達此目的須有幾個條件:
\newline
\newline
要有安靜的地方
\newline
\newline
在香港,居住環境實屬困難,為適應環境,只好當作旁若無人.戰爭時期,我們在內地讀書,全宿舍住了數百個學生,一排排的床,各人的床就是自己的天下,我們都在床上靈修,這方面須要特別學習才能適應.靈修是親近神,要有合適的時間和地點,虔誠者甚至可以跪下靈修.至於靈修時所讀經文,可選擇自己喜歡的經節,也可按次序讀經,真正目的乃與神親近.與神有切身關係,這並非為別人,所以要把自己置於神話語的光中,被光照,看見自己有錯就認罪,看見自己不夠就禱告；藉著讀經和禱告打成一片,好像哈拿在神面前傾心吐意.但靈修有一先決條件,必須先與人和睦,沒有罪的隔阻才能和神交通；完全除去與人之間的仇恨,最好能達忘我的境界；並且忘記時間和地點,但這些很難做到,因各有所忙.關於時間方面,可用鬧鐘報時,(用烘蛋糕用的鬧鐘較準)這樣可以忘記時間,專心靈修；若時間充裕的人,就不必受時間的限制了.
\newline
\newline
靈修時與神親近是否感到甘甜?讀聖經要有恆心,要有信心,我們必須深信來到神的面前,不論情緒,感覺如何,已經達到親近神之目的了.深信與神同在,乃憑信心,非憑感覺；因為人的情緒──喜,怒,哀,樂時有變化,沒有人整天笑,也沒有人整天哭,(神經病者除外)不能憑著情緒和感覺,而評估靈修的成功與失敗.我們是罪人,我們的良心要憑真理的審察；一旦發覺得罪人,虧欠人,就要向人認錯,與人和好,這樣才可以維持經常的靈修.有時因生活問題致影響靈修中斷,萬勿灰心,要立即恢復靈修,以免受影響；若不繼續下去,那才是真正的失敗.
\newline
\newline
熟讀聖經和靈修讀經的方法完全不同,若用靈修的方式,欲達熟讀聖經之目的,我想:讀到上天堂時還不能熟讀,因為在靈修中沒有充分的時間讀聖經,所以必須在靈修之外.
\newline
\newline
讀經是基督徒應有的本分；真理是委托給教會,委托給所有的蒙恩者.聖經是基督徒的兵器和糧食,也是神奧秘事的委托；我們是奧秘事的管家,接受福音真理管家的責任.教會是真理的柱石,教會要維持真理；真理如同柱石,棟樑.世界沒有真理,因為沒有神的話,神的話是賜給祂的兒女們；神的兒女有責任明白神說的話,有責任維持真理的純全.這是基督徒和教會共同的責任,當然傳道人有更大的責任.
\newline
\newline
靈修讀經可能受時間和地點的限制,如要熟讀聖經可用快讀的方法,免受時間,地點的限制；隨身攜帶聖經,隨時隨地都可讀.青年人閱長篇小說,十萬字一兩天便可看完；如果要快讀聖經,可利用零碎時間.你試過用五分鐘讀幾章聖經嗎:青年人五分鐘看三章,最少看兩章；我奉勸青年人要多讀聖經,因為年紀越大就閱讀越慢.如果一下要抽出半小時或較為困難,五分鐘較為容易,一天之間,集數個五分鐘便是半小時了,最少可讀十二章,這樣一個月便讀完新約並有餘；一年可讀十二遍,那麼,用三年或五年便可把新約聖經熟讀；且快讀的方法因時間短,注意力很集中；有一次,我在聚會時試驗,很多人在五分鐘內讀了三章,可是十分鐘卻讀不了六章；因為注意力不能維持較長的時間.最要緊禱告求神給你有愛慕神話語的心.
\newline
\newline
我讀聖經最多是在神學院求學時,那時甚至進洗手間都帶聖經；同學批評我不聖潔,我說:「我帶了一節背熟的經節出來,你帶了什麼出來呢?」我並非叫你們效法我,我乃是說明,如果你愛慕聖經,你就想盡辦法來讀聖經.求神給你有愛慕神話語的心!今日有許多人研究聖經,他們躍過這一階段,沒有在聖經的本文下過工夫；卻閱讀許多參考書和註解書,以為這樣就可明白聖經.其實,基本上先熟讀聖經,然後再參考其他的書,那情形就絕不相同了；能夠融匯貫通,能夠瞭解巧妙之處.兄弟有段時間任教神學,發覺中學畢業程度而熟悉聖經的神學生,對課程的領會,勝過大學程度而聖經不熟的神學生.比方當我講述批評小說紅樓夢,如果你曾閱此小說,熟悉其內容；當然領會力勝過那班一無所知的人.所以我們必須先熟讀聖經神的話語,然後進行研究,領會的程度就絕然不同了.不過,快讀當然沒有深刻的印象,但因重溫機會較多；一遍又一遍,卻給你聖經全貌重複的印象.這樣的概念實在是今天基督徒非常需要的!很多基督徒對於聖經全面的概念極其缺少,僅些片段；好像兒童玩拼圖,憑零星小塊無法明白整個圖樣；拼成之後,便可一目了然.許多人認為讀聖經有不明之處,暫時擱下,後來就會明白,也就是多讀就熟悉了.
\newline
\newline
此外,要熟讀聖經需靠記憶.記憶分兩方面；一是記憶概念,一是記憶經句.特別的經句需特別背熟,記憶概念最重要是每章的題目必須牢記,此法非常有用.每章選定你所印象最深刻的事作為題目,例如:馬太福音一章記基督降生.二章記博士朝拜.三章記約翰施洗.四章記耶穌受試探.五章記登山寶訓有八福.六章記登山寶訓有禱告.七章記登山寶訓有永生的門和滅亡的門.八章記五個神蹟其中一個大痲瘋.九章記五個神蹟其中一個是癱子.十章記耶穌選召門徒.十一章記施洗約翰問耶穌.十二章記辯論安息日的問題.十三章記天國的比喻.十四章記施洗約翰被殺.十五章記辯論洗手污穢的問題.十六章記彼得認耶穌是基督.十七章記山上變像.十八章記謙卑與饒恕.十九章記辯論休妻.廿章記兇惡的園戶.廿一章記耶穌騎驢入京.廿二章記娶妻的筵席.廿三章記法利賽人的七禍.廿四章記預言.廿五章記三個比喻.廿六章記設立新約.廿七章記耶穌受難.廿八章記耶穌復活升天.
\newline
\newline
如果牢記每章的題目,不必翻開聖經,可在腦子裡溫習,既省時且方便.第一章耶穌降生之前記耶穌的家譜；第二章博士朝拜之後,約瑟帶耶穌和馬利亞逃往埃及；每章以題目為中心,前後聯想；這樣可隨時隨地概念溫習聖經.若能時常這樣溫習在心,一方面你心中常有神的話,且不知不覺之中,聖經漸熟.所以記憶聖經每章的概念,對熟悉聖經大有作用.讀經時,無妨熟讀標題,但內容也要讀；若只讀標題,不知內容,就作用不大.你可試把馬可福音共十六章,自己定下題目並牢記,然後思想其前後所記載,這樣,你對聖經就能漸熟；可能你以為這樣太麻煩,但不勞而獲是沒有可能的.求主給我們有願付代價的心!
\newline
\newline
熟讀聖經的第二個方法是背熟金句.許多人都以為自己記性差,其實大家都是一樣.背金句必須經過忘記之後再次背熟；申命記早已暗示我們,無論在家,無論行路,躺下,起來,要把神的話繫在手上為記號,戴在額上為經文.意思就是要用種種方法多接觸經文.所以遇有金句,必定要背,一遍一遍；無論何時何地,重複背誦,且要多找機會應用.可把金句卡,袖珍聖經畫上顏色,隨身攜帶,隨時閱讀.我還用過一法,是把聖經金句錄音,在家中整天收聽,這樣也頗有幫助.無論用何方法,都不可灰心,趁著年輕,記憶力強,熟讀聖經,多記神的話語在心中,神要使用我們.
\newline
\newline
求主使我們將祂的話藏在心中,常常操練自己,務求熟悉和明白神的話語.
\newline
\newline


\newpage

\section{第50屆~港九培靈會~陳終道~第五講}
\label{sec:718}
\textbf{精細讀經}
\newline
\newline
連結: \href{https://www.hkbibleconference.org/session-message/view/718}{\texttt{https://www.hkbibleconference.org/session-message/view/718}}
\newline
\newline
\hyperref[sec:717]{< < < PREV SERMON < < <}
~
\hyperref[sec:toc]{[返目錄]}
~
\hyperref[sec:720]{> > > NEXT SERMON > > >}
\newline
\newline
詩篇 1:1-1:6
\newline
\begin{longtable}{cl}
\hline
\hline
章節 & 經文 (和合本修訂版)\\
\hline
1:1 & \begin{tabularx}{0.7\textwidth}{X} 不從惡人的計謀，不站罪人的道路，不坐傲慢人的座位， \end{tabularx} \\ \\ \relax
1:2 & \begin{tabularx}{0.7\textwidth}{X} 惟喜愛耶和華的律法，晝夜思想他的律法；這人便為有福！ \end{tabularx} \\ \\ \relax
1:3 & \begin{tabularx}{0.7\textwidth}{X} 他要像一棵樹栽在溪水旁，按時候結果子，葉子也不枯乾。凡他所做的盡都順利。 \end{tabularx} \\ \\ \relax
1:4 & \begin{tabularx}{0.7\textwidth}{X} 惡人並不是這樣，卻像糠秕被風吹散。 \end{tabularx} \\ \\ \relax
1:5 & \begin{tabularx}{0.7\textwidth}{X} 因此，當審判的時候惡人必站立不住，罪人在義人的會眾中也是如此。 \end{tabularx} \\ \\ \relax
1:6 & \begin{tabularx}{0.7\textwidth}{X} 因為耶和華知道義人的道路，惡人的道路卻必滅亡。 \end{tabularx} \\ \\
[1ex]
\hline
\hline
\end{longtable}
詩篇第一篇告訴我們,怎樣才能夠作個按時結果子的基督徒,怎能勝過各種不同的環境；那就是要接上源頭耶穌基督,接上神話語的源頭.神常有信息臨到你,若你對神的話經常有領受；就像樹栽在溪水旁,經常有直接的領受,所以能按時結果子.
\newline
\newline
昨天題到怎樣讀聖經,背經句和記住每章的大意,就可以幫助我們熟悉聖經.我也曾請求弟兄姊妹,自定馬可福音各章的題目並且牢記,這樣有助於記憶全卷內容.此外還要清楚全部聖經的概念.有些經文的記載,可能產生混亂的感覺,此類經文特別需要清理,才不致在觀念上發生混亂.然後,一遍又一遍閱讀,使印象加深；否則多遍閱讀也不會加深印象.例如出埃及記載會幕,許多人讀來讀去對會幕還是弄不清楚；這就是因為對會幕構造的佈置未先作了解.其實,只要明白構造的佈置,把會幕的佈置擺清楚就行了,至聖所有約櫃,聖所北邊有陳設餅的桌子,靠近幔子有香壇,南邊有燈台；聖所外面有洗濯盆,靠近大門口有燔祭壇.隔著幔子就是外院,這是主要的佈置.在腦子裡先有簡單的印象,讀時更加深印象,就容易明白了.讀舊約有許多地方必須弄清楚的,會幕是其中之一.
\newline
\newline
創世記較容易讀,因為所載都是人物,且按次連接,但是到出埃及記就難了.摩西領以色列民出了埃及,十九章他們到了西乃山的曠野；摩西和以色列人立舊約,在西乃曠野停留約十一個月,直到民數記十章十一節他們才離開西乃曠野.換言之,他們自出埃及記十九章以後,包括利未記,直到民數記十章十一節,他們才離開西乃山曠野.如果你先把握清楚歷史的時間和地點,讀時就不致混亂.讀舊約還有部份易引起我們發生混亂的,就是歷史的一些名字；以色列到了所羅門的兒子羅波安時,分開南北兩國.北邊是以色列國,有十九個王；南邊是猶大國,有廿個王.南北兩國的王,名字有相同的,甚至同名的王也有,同一時代在位的也有；所以易於混亂,必須整理清楚.北國從耶羅波安開始,拿答,巴沙,以拉,心利,暗利,亞哈；如先整埋清楚再讀,非但不致混亂,還能加深印象.
\newline
\newline
先知書也是使我們讀時頗覺莫明其妙的,因為它的信息次序,不按章數次序的先後；且先知的預言,時而時白的,時而象徵的,時而比喻,時而私意誇張的講法.時而指當時的人,時而在責備當時的人中就說出豫言,以致歷史和預言參雜.但先知和歷史的君王,不計算在一起,那些先知是列王在位時興起為先知的；以色列人沒有君王以前,他們沒有人統帥掌權.摩西之後是約書亞,繼而士師,那時非混亂；到了列王時,才立了君王執政,重建祭司制度,祭司管理宗教方面的事.至於先知的職份,包括寫歷史的事；神在此之外還興起了一些先知.他們不一定有國家的職份,但神供給他們信息；當君王,祭司或百姓陷入罪惡之中,神藉先知責備警告他們.先知從神領受信息,預先警告百姓前面的危險,勸勉他們回頭；並預言將來可能發生的事,先知以超人姿態出現,也包括在歷史之內.所以,要明白先知書,同時也應明白歷史書；因歷史書是先知書的背景,要以歷史對照才能明白先知書.
\newline
\newline
以色列前面幾個王:掃羅,大衛,所羅門；那時有撒母耳,拿單,迦特為先知.到了耶羅波安分國時,有亦多,亞希雅為先知.到亞哈王時,有以利亞,以利沙,有個叫米該亞的先知.到了以色列將亡國時,北國有何西阿,南國有以賽亞,彌迦為先知；被擄到巴比倫時,有耶利米,以西結,但以理為先知.被擄回國,重建聖殿時,有哈該,撒迦利亞,最後有瑪拉基為先知.歷史君王的程序認清楚,先知配合在內,這樣舊約就有較完整的概念.能找出這些概念,先在腦裡作鋪平的工作,讀時:程序,時間,地點就清楚；人物就不至混亂.
\newline
\newline
至於新約耶穌基督的生平記於四福音,雖然耶穌生平合參已有幾種出版,但我們也無法記得整個的合參.較簡單記載基督生平的聖經,祂童年以前的事,路加二章載耶穌十二歲上聖殿,包括祂降生和孩童時候.到了耶穌出來傳道以後,約翰福音記載耶穌出來傳道經過四個逾越節,每年一個逾越節,耶穌傳道三年半.第一個逾越節記在約翰二章,第二個記在六章,第三個在十一章,第四個在十九章；記憶中有了這些概念,就幫助我們記住耶穌的生平.使徒行傳,前十二章記載使徒彼得的事；十三章開始記保羅的事.保羅工作主要有四行程,十三和十四章他遊行佈道,十五章上段遊行佈道之中到耶路撒冷開會；末段記他開始第二次遊行佈道.直到十八章下段又開始第三次佈道行程,廿章他回耶路撒冷被捉拿,後來到該撒利亞,上告於該撒.第四個行程是以囚犯身份到羅馬；行傳記載保羅到羅馬是他第一次下監獄.
\newline
\newline
我們一生讀聖經,必須記得四點:
\newline
\newline
(1)讀,如果有了許多方法而不肯讀,則等於零.
\newline
\newline
(2)熟讀能生巧.
\newline
\newline
(3)求明白,熟而不明,恐有弊病.
\newline
\newline
(4)用,既熟又明,要將所明的運用.以上四點,必須按著次序.
\newline
\newline
怎樣精細讀聖經?第一步,先反覆地讀,知其概念,把每章所記的事列出.第二步,把每件事當作一小段,其中有何人物,情節,道理.有何關於神的事,人的事,撒旦的事.有何關於神的應許,神的命令.有何關於人的失敗,其失敗原因和結果.有何人的成功,其原因及結果.何為我們當立即遵行的,何為當時的人所作,何為神所憎惡和喜悅的事.這些問題可以從多方面,提問此段經文；當然不是每段經文所合用.比方此段經文沒有任何人物,也沒有講論任何道理,或有講論責備的話.總之,要靈活應用,非公式化,非電腦化,要與我們的頭腦發生共鳴,從許多角度去思想.第三步,把疑難之點列出,心得列出,作個總結；而後閱參考書,這樣收益必大.如要精細讀經,應先熟讀,明白主題,把主要和次要信息分開；這樣可操練讀經的能力,然後再閱參考書.參考書不可不看,因為神藉祂許多僕人所領會的寫下,目的為留給我們；如果不看,對自己的損失很大.然而,如不熟讀聖經先看參考書,就未免有不消化,難吸收之嫌；我們要學習如何從神領受,也參考別人如何從神領受；要隨時願意改正自己的錯誤,不存專依賴別人的心.日久你就漸有從神話語領受信息的能力,心竅習練通達,就有從神話語領受信息的力量.看書目的非單從書中領受教訓,同時也領受作者從神所領受的話語,這才是重要之目的.
\newline
\newline
例如創世記第三章的記載,有神,有人,有鬼；有罪惡,救贖,審判,應許；有人的失敗,有神的恩典；有分別善惡樹,有生命樹；有許多神學上的難題,可供神學家辯論半天,此章齊備.在你未看參考書之前,你一遍遍地讀,首先自己決定可分成幾段,一段是亞當夏娃犯罪的經過,一段是神審判他們的經過,一段是神審判的結果.至於人物:有神,亞當,夏娃,魔鬼,一至七節是亞當,夏娃犯罪的經過,首先看到魔鬼怎樣試探人,蛇對女人說:「神豈是真說,不許你們喫園中所有樹上的果子嗎?」其實；神並非這樣說,魔鬼把神的話改變；一方面,神吩咐說:「園中各樣樹上的果子,你可以隨意喫；只是分別善惡樹上的果子,你不可吃.......」神並沒說:不許吃園中所有樹上的果子,魔鬼竟把兩個問題混雜來說,目的是要引夏娃交談,從中插入試探.魔鬼引誘人犯罪常用試探方法,把真假混淆,使人思想紛亂.為何魔鬼先試探夏娃而不先試探亞當?為何夏娃不讓亞當回答魔鬼的問題呢?神給此命令時到底夏娃是否在場?其實,神的吩咐在先,夏娃被造在後；亞當直接從神領受命令,清楚這個命令；夏娃則不然,以致回答魔鬼的話出了問題；有處增加,有處減少.神在創二章說不可吃分別善惡樹上的果子,並沒有說:「不可摸」,夏娃卻加上「不可摸」,且偏去摸那棵樹；神說:「喫的日子必定死.」夏娃卻說:「免得你們死,」夏娃對於神的命令模糊不清,所以上了魔鬼的當.夏娃吃了禁果犯了大罪,但罪名卻由亞當擔當；聖經說由亞當一人犯罪,罪就入了世界.當時亞當乃奉神之命修理看守伊甸園.為何要修理看守?那時全世界惟亞當一人,難道還要看守嗎?此乃暗示要防備敵人──魔鬼.亞當的失敗是魔鬼侵入他還不知道,以致夏娃被引誘,她不加考慮吃了分別善惡樹的果子,而且給亞當吃；他們立即眼睛明亮,才知道自己赤身露體.到了最後,神來到他們中間,亞當聽見神的聲音就藏起來.犯罪的結果就是受審判.如果我們從許多不同的角度,思想這段經文,其中必會得著許多的教訓.
\newline
\newline
聖經是供每個人讀的,然而許多基督徒卻把讀經當作是傳道人的本分,這不合神的旨意!基督徒若不普遍明白聖經,易於發生異端,神的僕人就難推動工作.我們需要回到神的話語上!
\newline
\newline
新約馬太福音一章分兩段,一段記耶穌的家譜,一段載耶穌降生,天使向約瑟報告.新約一開頭就記耶穌的家譜,介紹祂是亞伯拉罕的後裔,大衛的子孫；因為舊約重視祭司和君王的家譜；若家譜不清楚,身份有問題就不能任職.神應許亞伯拉罕的後裔得福,後裔乃指耶穌基督而言；神也應許大衛寶座存到永遠,實際上大衛的寶座由其子所羅門繼承；到了所羅門之子羅波安即告中斷.寶座存到永遠,是指耶穌基督說的:猶太人的觀念,大衛的子孫將坐他寶座的是彌賽亞.家譜的開始就題到大衛的子孫耶穌基督,然後才題到約瑟；到底哪個重要,應當分辨清楚.按歷史是約瑟為主角,但照這段信息是論耶穌基督為主.試把18至20節列出論到耶穌之點；論到約瑟之點有何可作我們的榜樣,他如何處事,他怎樣受感動得天使的指示,娶過馬利亞,沒有與她同房；凡此均足作為我們屬靈的教訓和知識.至於不明白之處,非屬懷疑不信,而是自認為不明白的,記錄下來,然後找機會解決問題,或閱參考書,或請教神的僕人,或以後查經時得以明白.
\newline
\newline
每個基督徒可找出時間,仔細讀一段聖經；或仔細讀福音書中的神蹟,操練精細深入聖經中.這樣,無論你到那裡,都像一棵樹栽在溪水旁,直接從神領受；如果你永遠接受從人來的,那麼人給你好的,你就得好的；人給你壞的,你就得了壞的,而且還會受了壞的損害.
\newline
\newline
從神領受是每個基督徒的本分!
\newline
\newline


\newpage

\section{第50屆~港九培靈會~陳終道~第六講}
\label{sec:720}
\textbf{如何避免誤解聖經}
\newline
\newline
連結: \href{https://www.hkbibleconference.org/session-message/view/720}{\texttt{https://www.hkbibleconference.org/session-message/view/720}}
\newline
\newline
\hyperref[sec:718]{< < < PREV SERMON < < <}
~
\hyperref[sec:toc]{[返目錄]}
~
\hyperref[sec:721]{> > > NEXT SERMON > > >}
\newline
\newline
提摩太後書 2:15-2:15
\newline
\begin{longtable}{cl}
\hline
\hline
章節 & 經文 (和合本修訂版)\\
\hline
2:15 & \begin{tabularx}{0.7\textwidth}{X} 你當竭力在神面前作一個經得起考驗、無愧的工人，按著正意講解真理的話。 \end{tabularx} \\ \\
[1ex]
\hline
\hline
\end{longtable}
保羅勸勉提摩太竭力在神面前,作無愧的工人,按著正意分解真理的道.傳道人按著正意分解神的道,亦即信徒按著正意領受神的道.我們讀聖經,須按聖經的正意來領會；但很多時候,我們讀經卻存著自己的偏見；這樣,會影響我們不正確地領會聖經.一個熱心愛主的基督徒,必愛讀聖經；但若錯誤領會聖經,就可能引致相反的效果,因為聖經的根據錯誤了.
\newline
\newline
使我們誤解聖經的因素:
\newline
\newline
(一)個人主觀的偏見,原由多是先入為主
\newline
\newline
初信主時,聽到關於某節聖經錯誤的講解,那時尚未有分辨的能力就接受了；於是根深蒂固,影響後來那段聖經正確的了解.有的人對某節聖經有特別的愛好,所以更歡喜那節聖經.但領受神話語的方式,是很寶貴的；我們必需正確地領會聖經.
\newline
\newline
例如在抗日戰爭時期,有個基督徒逃避日機轟炸,他和朋友們躲在溝中；當時有句話一直臨到他心中「起來向南走!」他便向南走.走了數百碼,又躲在溝裡；後來日機用機關鎗掃射,那溝裡的人就全死了；他卻死裡逃生.他非常感謝主!因主實在對他說話.「起來!向南走.」這句話是聖靈對腓利說的,這句話確實救了那人一命.當時他確有被聖靈感動的經驗,可惜他後來竟然構成個人的觀念,他以為凡聖經裡提到向南邊,都是好的；向北邊就是壞的,因他得的經驗,印象深刻.
\newline
\newline
存著偏見讀聖經是不對的.我們相信神有時會讓類似的經驗臨到我們,我們因神的話而得到深刻的帶領,甚至恐裡逃生；但那並非一成不變的,不過僅是臨時的應付辦法而已.聖經有正確的辦法,神是要我們明白整部聖經的意義,如同明白知己的一樣.我們應當藉著聖經而認識神,這才是我們讀經之目的.
\newline
\newline
(二)因個人工作立場,為逃避困難而產生的偏見
\newline
\newline
洗禮宗的教會,是主張受洗禮的,讀聖經讀到的都是洗禮；但浸禮宗的教會就必讀受浸禮,既有浸禮的宗派,所以讀聖經讀到的都是浸禮.我並非說不該有立場,也不是說那立場才對,而是讀經必須放下工作立場；不容任何必需的理由存在,才能從聖經神的話,找出正確的意思.
\newline
\newline
我在東南亞工作多年,有許多青年信了主不願受洗；他們以(路二十三43)耶穌釘十字架時對旁邊的強盜說:「我實在告訴你,今天你要同我在樂園裡了.」引證那個十字架旁邊的強盜,也未曾受洗；但耶穌卻明確應許他可以得救.這些青年不受洗的原因,並非他們真正被掛在十字架旁邊,而是要逃避家庭的反對.許多不信主的家庭,只願孩子到教會,卻不准孩子受洗；他們認為受洗,就是信了耶穌.故此,教會就形成了一個問題；這班確已蒙恩的青年,他們思想中存在不必受洗可以領受聖餐的觀念.他們認為既已蒙恩,便可記念施恩拯救的主；問題是要逃避面對的困難.其實,在這種情形之下,更應該受洗,以表明自己確已信了耶穌；我們必須面對現實的挑戰.使徒們在初期教會時就是這樣,使徒行傳中可見到一些例子；受洗與得救是互相關連的,受洗與赦罪是相提並論的.在當時信而受洗,兩件事是分不開的；他們本是猶太人,但怎能顯出如今是個基督徒?他們本來就信有神,也是信這位獨一的神；受洗就是個明顯的證據來公開他們的信仰,既然信,就必須受洗.看聖經必須放開面對的困難,才能找到正確的意思.
\newline
\newline
例如關於十一奉獻的問題有兩種主張:牧養教會的,與在教會機構團體工作的,見解就不同.有人主張十分之一要完全歸給所屬教會,這是按照(瑪三10)說:「萬軍的耶和華說,你們要將當納的十分之一,全然送入倉庫,使我家有糧.......」家是神的教會,但在教會機構工作的則認為十分之一,不一定獻於本教會,只要獻於教會範圍內便可.為何有兩種不同的看法?我不說誰是誰非,但我認為以工作立場來看聖經,就看得不正確了.我們若想把聖經看得正確,不至誤解,並能深切的領會；就必須放開個人一切的利益關係.
\newline
\newline
還有一些人,看聖經另有偏誤的原因,例如:研究科學的認為聖經不合科學,因此以聖經的解釋來遷就科學；用科學的理論,解釋聖經,就會極易發生錯誤.我曾經提過,科學對宇宙的來源尚未有最後的答案,而且論調時常改變；所以若以科學解釋聖經,聖經就要跟著科學改變了.其實聖經所載與科學並不衝突,但我們解經毋須依賴科學才算權威；因聖經本身是絕對的,不在乎人之承認與否.
\newline
\newline
例如創造天地萬物六天的過程,有許多不同的解釋.因為創世記和現代進化論有所衝突,假使為了遷就進化論,而把六天的創造,當作六個長時期來解釋；就算可以這樣解釋；不過是為著遷就進化論,在原則上已是錯誤了.把六天創造當作六個長時期,雖然可解決創世記第一章的一些困難；但是第二章和第五章卻有更多的難題.因第一章多次說有晚上,有早晨.如果以一天當作一個長時期,那麼,晚上和早晨就變作各半個長時期了.試問是否一個長時期的黑夜,然後又一個長時期的白天呢?那就很困難了!因為創世記第一章,六天的創造多次提到有晚上有早晨；每天提一次,這樣就把那一天的時間管住了.第二個困難,第七天是安息日,和前六天完全一樣；第七日就是後來神和以色列人立約的憑據.出埃及記卅一章末段,神吩咐以色列人,要世世代代守安息日.而且特別註明,神六天創造天地,第七天安息舒暢；明顯的指創世記二章的話.以色列人守安息日,是根據創世記二章開頭兩節的話；這安息日不是個長時期,乃是指著第七日而言.安息是指長達十萬年的長時期,那麼,我們就永遠安息,不必作工了.第三個問題,創世記二章,提到亞當被造較為詳細的描寫,神用塵土造人,然後從鼻孔吹口氣；如果說亞當被造是經過進化而來,則創世記二章變成一種寓意的描寫,非正確的史實.那麼,吹一口氣經過很長時間才造一個人,後來又造夏娃,造多少年呢?第二章實在很多問題!到了第五章說亞當活到九百歲就死了；如果神創造一個人需經萬年,僅活了九百多歲,未免太差勁!
\newline
\newline
我們不可為符合科學用來解釋聖經,科學只不過是個見證人,不是審判官.我們絕不能以科學來審定聖經,科學有時不過是證明聖經的話是對的而已.我們必須對聖經有完全絕對的信心,才不致誤解聖經.
\newline
\newline
要避免誤解聖經,也要了解聖經的字意和經意；如果我們不了解經文的正確意思,單咬文嚼字,可能誤解其意.
\newline
\newline
我在新加坡時,曾參加過青年人的禱告會,傳了信息後,他們準備禱告,就關門熄燈.(那地天氣炎熱無冷氣設備)禱告完畢,我問他們為何要關門熄燈?答:「聖經記載,耶穌說:禱告時要進入內屋,關上門,禱告暗中的父.」有時人只遵照聖經字面的意思,其實字面的意思,不一定是準確的意思,我們必須從其上下文來領會聖經的意思.
\newline
\newline
許多人極端反對靈意解經,因為不單字意誤解,靈意也常有誤解.耶穌說:「若一隻手叫你犯罪,就把這手砍掉；你缺少一手進入永生強如兩手丟在地獄裡；若一眼叫你跌倒,寧可挖掉那眼,缺了隻眼進入永生強如雙眼掉在永火裡.」若按字面的意思,手跌倒砍手,眼叫你跌倒挖眼；如果肚子整天想吃東西,把胃挖掉,腳犯罪砍腳；假如這樣,等你進天堂時甚麼都沒有了.我們不應該單從字面來領會經意；耶穌說這話不過是比喻,意思教我們要決心對付罪惡,要願意付上代價；這樣的人才能勝過罪.
\newline
\newline
曾經有個基督徒對我說,耶穌說:「禱告時不要像法利賽人用重複的話.」而我禱告時常用重複的話,所以我禱告之後,深恐耶穌不垂聽.雖然這不是基督徒的一般情形；但我想這樣的基督徒也不少.關於重複禱告,有兩種原因:法利賽人是故意延長禱告的話；像寫文章的人,沒有資料可寫,故意裝腔作調,無病呻吟.另一種原因是懇切之故,心裡有迫切的需要；緊急,惦記,可能也會重複禱詞.耶穌並無定下禱告的規則,叫我們講過的話不可再講.主耶穌在客西馬尼園禱告,三次的話都是一樣(太二十六42,44)；但耶穌禱告的重複和馬太六章禱告的重複完全不同.我們領會聖經的話,要按照上下文的意思；約翰十七章提到「世界」,同章同字,意思卻不同；因為上下文的用法不同.所以不能單憑字面的原意,來解釋那段聖經,必須看該字在全句上的用法.約翰十七章五節,耶穌求神叫自己與神同享榮耀；這裡「世界」和十一節「世上」在原文同出一字；但五節「世界」乃指整個宇宙,未有整個宇宙以前,主耶穌與神同樣榮耀.十一節「世上」乃指地球,那時耶穌已來到世上.14節「世界」指世界上的人；因為無可能世界會恨人,屬世界的人才恨不屬世界的人.16節耶穌說:「他們不屬世界,正如我不屬世界一樣.」這裡「世界」包括一切邪惡的勢力.耶穌和門徒都生活在世界,但他們不屬於世界的邪惡勢力範圍.魔鬼是罪惡勢力,影響全世界；但基督徒不屬此世界.屬物質的世界於我們是暫時的,我們不在罪惡勢力底下,我們是在基督裡面.
\newline
\newline
讀聖經必須根據上下文來領會,才是真正的領會.求主給我們智慧!
\newline
\newline
耶穌在世上時,曾告訴門徒不要請有錢人吃飯；要請那些貧窮的,在路上求乞的,他們無法報答的.這話並非叫我們不請有錢的朋友吃飯.有個弟兄對我說:「只能稱你弟兄,不稱呼你牧師；因耶穌曾說:「你們中間都是弟兄,只有一位是夫子.」(太二十三8-10)其實這段聖經是耶穌教訓門徒,不可自以為比別人強,要謙卑彼此服事之意.如果這段聖經是以教導我們關於對人方面的稱呼,那我們豈不是連肉身的父親也不稱為父嗎?難道要稱呼父親為老弟兄嗎?
\newline
\newline
領會聖經常會有自己的偏見,叫我們加上自己的意思,以致失去原來的經意.讀經必須按正意領會真理的道,讀經真正的秘訣必須留意上下文:熟了然後求精,才不致犯咬文嚼字的毛病.
\newline
\newline


\newpage

\section{第50屆~港九培靈會~陳終道~第七講}
\label{sec:721}
\textbf{查經可能遇到的難題}
\newline
\newline
連結: \href{https://www.hkbibleconference.org/session-message/view/721}{\texttt{https://www.hkbibleconference.org/session-message/view/721}}
\newline
\newline
\hyperref[sec:720]{< < < PREV SERMON < < <}
~
\hyperref[sec:toc]{[返目錄]}
~
\hyperref[sec:723]{> > > NEXT SERMON > > >}
\newline
\newline
彼得後書 1:20-1:21
\newline
\begin{longtable}{cl}
\hline
\hline
章節 & 經文 (和合本修訂版)\\
\hline
1:20 & \begin{tabularx}{0.7\textwidth}{X} 第一要緊的，你們要知道，經上所有的預言是不可隨私意解釋的， \end{tabularx} \\ \\ \relax
1:21 & \begin{tabularx}{0.7\textwidth}{X} 因為預言從來沒有出於人意的，而是人被聖靈感動說出神的話來。 \end{tabularx} \\ \\
[1ex]
\hline
\hline
\end{longtable}
關於避免誤解聖經的一些補充,我們別忽略避免誤解聖經；要放下個人的偏見,要存心正確,置自己於神面前來讀神的話語.不但要避免字意的誤解,並須從經文瞭解其中的經意；更重要的,讀經時萬勿斷章取義；這點我們早已知道；但或不甚明白.很多基督徒所發的聖經問題；如果留意上下文,亦自可回答.要避免斷章取義,就要留心上下文,因為這是讀經方法中之最重要的.所有解釋聖經的重要原則,均須與上下文配合.有人說:「神是愛,是光.」看神為抽象的,而非實際存在的；光是摸不到沒有位格的,愛也是沒位格的；這樣就把神解釋為人最高的理想,不是實際存在的一位神.完全誤解聖經了.約翰壹書提到神是光,神是愛；當我們讀了全卷之後,就可以知道神絕非沒有位格的；這卷書非常強調我們怎樣與神交通,所以當然神是有位格的,我們才可與祂交通.假如神不存在,那麼,你與神交通等於自己與自己說話；但是約翰一書,開頭就證明:「論到從起初原有的生命之道,就是我們所聽見,親眼看過,親手摸過的.」他證明神的存在.
\newline
\newline
聖經有一個字曾引起很大的爭論,(賽七14)研究希伯來文的說「童女」也可作「少婦」解；所以有人就把童女譯成少婦.如果注意上下文,則應翻成童女才對；若把新約馬太和路加福音連起來看,更可肯定此字應翻作童女.若先知說:「我給你個兆頭,必有少婦懷孕生子,這根本算不得兆頭,少婦生子有甚麼希奇,怎能算是特別的道理?何必勞動先知宣告.但童女懷孕生子,才是特別的兆頭,才值得先知宣告.假如這字是「少婦」的話,馬太福音一章的記載就毫無意義了.天使向約瑟報夢,向他解釋,馬利亞所懷的孕,是從聖靈來的.第一章約瑟想暗暗地把馬利亞休了,為何無緣無故要把她休了呢?那麼,馬太一章全部記載屬於多餘了.當天使向馬利亞報喜訊時,(路加一章)預先告訴她要從聖靈懷孕.馬利亞對天使說:「我還沒出嫁,怎有這事呢?」天使答說:「聖靈要臨到你身上,至聖者的能力要蔭庇你.」在人不能,在神凡事都能.若馬利亞是個少婦,這些記載均屬多餘.如果連貫上下文讀,就不致曲解聖經,也不致強解經意.
\newline
\newline
林後六章頭兩節,保羅說他是與神同工的,勸勉信徒不可徒受神的恩典.「因為祂說,在悅納的時候,我應允了你；在拯救的日子,我搭救了你.」很多基督徒用這節聖經傳福音,「看哪,現在正是悅納的時候,現在正是拯救的日子.」其實,這兩節聖經並不是對非基督徒說,現在要信耶穌；細看上下文即知道這是對基督徒說的；保羅說:「我們與神同工的也勸你們,不可徒受神的恩典；因為神曾經在今天悅納的時候應允了我們,在搭救的日子我們得了拯救.看哪現在正是悅納的時候,現在正是拯救的日子.」這是對基督徒說,要趁著神悅納的時候,和神搭救的日子領人歸主,才不徒受神的恩典；信徒要趁今天神悅納的時候去拯救人；因為我們正是神悅納人時蒙拯救的.此段經文,連貫與不連貫上下文看,意思就不同了.
\newline
\newline
讀聖經不要斷章取義,謹慎留心上下文；就算領會錯了,也不致於錯到嚴重的地步.若不留心上下文,就算以經解經；其實只不過以經串經,不一定可以解釋清楚.只不過僅勉強串連相同的經句而已.若留心上下文,加上別處經文的補充；領會聖經就不至於有大差錯.
\newline
\newline
查經可能遇到的難題可分兩大類:
\newline
\newline
第一,因為聖經不熟悉,或因知前不知後；知此不知彼,認識聖經不夠完全；所以自以為是難題.
\newline
\newline
第二,因聖經沒有給我們足夠的啟示,作為根據以解決該問題.那就是本來無法解決的問題,這類的難題,我們應加以留心.
\newline
\newline
現在先提到些較為容易的問題.
\newline
\newline
有人常會發覺聖經似有矛盾之處,例如(太十二30)耶穌和門徒反駁法利賽人的話:「不與我相合的,就是敵我的；不同我收聚的,就是分散的.」界限非常清楚.(路九50)耶穌說:「不要禁止他,因為不敵擋你們的,就是幫助你們的.」這經文似乎和上面經文相反,但如果留心它的上下文；耶穌說不與我相合的,就是敵我的.......這話是對不信祂的人而言；那班法利賽人和文士們,知耶穌靠神的能力趕鬼；卻說耶穌靠鬼王別西卜趕鬼.耶穌在下文說他們是毒蛇的種類.路加九章是指一班信耶穌的人,因不跟隨門徒們去傳道,他們奉耶穌的名趕鬼,所以門徒禁止他們；耶穌對門徒說:「不要禁止他.......」可見耶穌對那信奉祂的人,並非和耶穌在一起,才算是信奉祂；有的人耶穌要他跟從,有的耶穌不要他跟從.有個原來被鬼附而得醫治的人,要求跟從耶穌；耶穌不許,要他回自己的家鄉,見證耶穌為他作了何等大的事.這兩句話看來有所衝突,其實是互相補充的；聖經多處的真理,都是互相補充的.我們應合併起來看,別把互相補充的孤立；當作互相衝突看待.保羅對傳基督的人有很寬大的心,甚至嫉妒紛爭,他也以為對他無所妨礙；只要是傳基督,所以基督就被傳開了.不過保羅對於那班傳異端的人,卻絕不容忍；他稱他們假弟兄,就是一刻工夫,也沒有容讓順服他們.(加二4-5)真理的兩方面是互相補充,而非互相衝突的.
\newline
\newline
聖經的記載,未必是聖經所贊同的；有時我們對此有所誤會,以為聖經所記的事才可作,沒記載的事不可作.茲列舉數例:
\newline
\newline
亞伯拉罕娶夏甲.聖經記載撒拉提議以使女夏甲給亞伯拉罕作妾；後來夏甲懷孕,就小看她的主母,於是發生爭吵.撒拉苦待夏甲,夏甲就逃走；但神要她回家待產.這是創世記十六章的簡單記載.聖經並無提及,亞伯拉罕娶夏甲是對與不對,只把事情的經過記載；不加置許,也不表示贊成與否.聖經從未提及夏甲是亞伯拉罕的妾,只提到撒拉建議以夏甲給亞伯拉罕為妾；事後聖經記夏甲均以使女稱之.新約加拉太書引用創世記十六章的記載,也是以使女的地位稱夏甲.由創世記十六章,我們看見亞伯拉罕的家庭,因多了個妾而發生糾紛.當時的猶太婦人若不生育,視為莫大的羞恥；所以撒拉把自己的使女給丈夫為她生子,所生孩子歸於自己名下,以除去羞恥.夏甲懷孕在身時,生男或生女尚未可知,可惜主婢時常爭吵,因家庭糾紛造成痛苦.後來夏甲所生的子以實瑪利和以撒爭吵,撒拉就把夏甲趕逐出去.神同意撒拉這樣做,因為得應許生的才可承受產業.今日的阿拉伯人就是以實瑪利後代,和以色列為世仇,全世界都因此受影響.我們若細讀聖經,可知這事非屬神的意思.
\newline
\newline
士師記十一章記載士師耶弗他,是神興起他要救以色列人,脫離亞捫人的手.他未出去打亞捫人之前在神前許願,(30-31)從來他打勝仗回來,第一個出來迎接的竟是他的獨生女.很多人發出問題說,到底這是神的旨意嗎?其實他許願和打勝仗毫無關係,並非因他許願才打勝仗；神的靈早已降在耶弗他身上.上文記載,神的靈降在耶弗他身上是在先,他許願是在後；神早已應許要把亞捫人交他手中,他卻自己畫蛇添足,結果悔之晚矣.
\newline
\newline
諸如此類記載,聖經並無置評是非,只是記錄當時史實而已；目的在於顯明士師時代以色列人在靈性上和真理上何等落後.甚至被神興起的士師在屬靈真理上尚且糊塗,在那時代的士師,半熱心半糊塗者大有其人,他們被神使用並非因他們的完全.
\newline
\newline
聖經記載奧秘事的部份,也給我們感到許多困難；奧秘事是屬於神的,祂沒有完全向我們顯明.許多人以為聖經是無所不載的書,其實不然,聖經只是記載神所要我們知道的事；而非記載我們想欲知道的事.可能我們想知道的事很多,但不一定是神要我們知道的；神所要我們知道的,一切都記在聖經裡了.有些屬於靈界的事,聖經沒記載.
\newline
\newline
但以理書第十章,我們看見因但以理禁食禱告；天使長和波斯魔君,與希臘魔君發生爭戰.靈界的戰爭是怎回事,我們不得而知；但我們知道,禱告會影響靈界發生爭戰.看來那些魔君似有地區性,因魔鬼非全能非無所不在的；所以有波斯魔君,希臘魔君.
\newline
\newline
猶大書九節提到天使長米迦勒,曾為摩西的屍首與魔鬼爭辯.聖經別處沒詳載,屬靈境界裡,有時天使與魔鬼爭辯.關於摩西的屍首,我們只能從申命記卅四章推想,摩西死時神親自為他安葬；這是摩西一生最大的尊榮,但人都不知摩西葬在哪裡.為何神不讓人知道摩西安葬之處?我們只能推想,可能因為魔鬼要利用摩西的屍首引誘以色列人,而神就不容許；所以引致天使長和魔鬼爭辯.摩西在世時雖遭以色列人反對,但他死後,以色列人很敬重他.我們若留心看福音書,發現那些控告耶穌的人,都是以摩西律法為題；可見摩西的律法,在猶太人心中的地位.摩西是頗受以色列人崇拜的一位.關於此類史實,我們只能憑聖經所記略知一二,詳情無法查考,惟有說:「奧秘的事是屬於耶和華；顯明的事才是屬於我們和我們的子孫.」
\newline
\newline
另有一點順便在此略題,聖經有些字句,在聖經中有特別專門的用法.各地有不同的方言和流行語,各處有它當地特別的說話；至於聖經,不但有希伯來人的說法,並有聖經原本的用語.例如耶穌在路加十六章論不義的管家說「你們要藉那不善的錢財,結交朋友.......倘若你們在不義的錢財上不忠心,誰還把那真實的錢財託付你們呢?」很多人對這話都不明白,根本基督徒不該要那不義之財,還說藉不義之財結交朋友.......其實,耶穌所說不義之財,乃對錢財的一種特別名稱,所有錢財都是屬於不義的世代,只在不義的世代有用；在永義裡卻無用,不義的錢財乃指今生有用的錢財.
\newline
\newline
還有一例:哥林多人喜歡彼此告狀,保羅責備他們說:「你們竟敢在不義的人面前求審.」保羅的意思,乃指凡不信耶穌的,都是不義的人:都屬於不義的世代,這話也是聖經獨有的語法.比方說:「神的靈離開掃羅.」聖經說:「從耶和華那裡來的惡魔攪擾掃羅!」難道神那裡有惡魔嗎?此類均屬聖經獨有的語法,因為惡魔攪擾掃羅是經神許可的；好像魔鬼在神前控告約伯,後來神許可魔鬼攻擊他；聖經說:「撒旦從神面前退去.」有時聖經為要表明神的權柄,一切主權都屬於神；神是掌管萬有的,所以有此類記載.例如大衛數點百姓,這是大衛的錯.聖經說:「耶和華激動大衛數點百姓.」但是又說:「大衛數點百姓犯了罪,大衛也自己認錯.」不過,如果我們看歷代志上廿一章便知真正激動大衛,數點百姓的是撒旦；因為經過神的許可,所以聖經如此記載.此類記載較為困難,所以我們應存謙卑的心,並且存著信任聖經的態度來讀.有些不是真正的難題,不過是前後補充；只要留意上下文即可明白.但是我們應當知道,在屬靈的境界,有很多事,神今天不能告訴我們；要等到我們與祂面對面時才能知道.在這有限的肉身中,我們當常存謙虛的心.
\newline
\newline
求神教導我們,讓我們做個明白神話語的人,讓我們不強解聖經.
\newline
\newline


\newpage

\section{第50屆~港九培靈會~陳終道~第八講}
\label{sec:723}
\textbf{聖經的預表和靈意}
\newline
\newline
連結: \href{https://www.hkbibleconference.org/session-message/view/723}{\texttt{https://www.hkbibleconference.org/session-message/view/723}}
\newline
\newline
\hyperref[sec:721]{< < < PREV SERMON < < <}
~
\hyperref[sec:toc]{[返目錄]}
~
\hyperref[sec:724]{> > > NEXT SERMON > > >}
\newline
\newline
歌羅西書 2:16-2:17
\newline
\begin{longtable}{cl}
\hline
\hline
章節 & 經文 (和合本修訂版)\\
\hline
2:16 & \begin{tabularx}{0.7\textwidth}{X} 所以，不要讓任何人在飲食上，或節期、初一、安息日等事上評斷你們。 \end{tabularx} \\ \\ \relax
2:17 & \begin{tabularx}{0.7\textwidth}{X} 這些原是未來的事的影子，真體卻是屬基督的。 \end{tabularx} \\ \\
[1ex]
\hline
\hline
\end{longtable}
許多基督徒以為聖經的預表,可以隨便解釋,所以有很多人反對預表的解釋.
\newline
\newline
我們講預表時,不能不講靈意,因這會牽涉靈意解經的問題很大；所以有很多人反對靈意解經,按我個人的領會,那些反對靈意解經的人實際上他們是反對私意解經.我們不挑剔靈意的不是,但應該反對私意解經的錯誤.靈意的用法確是聖經原本的用法,且非偶然用之,而是多次用之.私意解經不單只憑靈意,甚至連字意也用私意解釋.
\newline
\newline
茲略舉數例說明:
\newline
\newline
(一)聖經的確有靈意預表的用法......
\newline
\newline
「不是憑著字句,乃是憑著精意；因為那字句是叫人死,精意是叫人活.」(林後三6)「意」在中文或為「義」之錯字,精義的意思即最合理之道.「......預言中的靈意,乃是為耶穌作見證.」(啟十九10)上文「精意」與這裡「靈意」,在原文同是一字,原文 Pneuma 的意思是「靈」有時指神的靈；英文聖經靈字用大寫來表明聖靈的靈(Spirit)；其實這字不一定是神的靈；鬼的靈,人的靈都同是此字,所指何靈要視其上下文而定.聖經中有兩次此字不指神,不指人,也不指鬼；乃指神話語的靈,神律法的靈,神預言中的靈意,和(林後三16)「精意」同是一字指靈,聖經惟此兩次指神話語的靈.話語和文章哪裡有靈?所以中文翻成「精義」這個意思非常正確,文章的靈魂就是其精華.(啟十九10)提到預言；用「靈意」即與「精義」相同的用法.
\newline
\newline
在聖經中有些處的確是「靈意」或「精義」的用法.本來翻成精義已是夠好的,為何聖經和合本神的僕人要用「靈意」和「精義」兩不相同的字呢?因須視該處適用何字；比方耶穌說:「我的肉是可以吃的,我的血是可以喝的；」如果憑字意可能發生問題,(天主教的變體,則用字意來解釋這話；他們說聖餐的餅經過禱告之後,就變成耶穌的身體.)耶穌說這話不是比喻,而是實際上我們可領受祂的生命；我們可得祂寶血的潔淨.這是屬靈的事實,非屬物質的事實；所以只能按靈意論.
\newline
\newline
保羅對哥林多人說:「我們逾越節的羔羊基督已經被殺獻祭了；所以守這節不用舊酵.」他乃是勸勉哥林多信徒要聖潔,要除去犯罪的舊人.如細讀(林前五 7),「羔羊」字旁有小點,意即原文無此字.為何譯者要加上「羔羊」?因舊約逾越節確是用羔羊獻祭；若此處經文不加上「羔羊」,則「我們逾越節的基督已經被殺獻祭了」,這樣就不通了；因為逾越節的時候還沒有基督,那時基督還未降生成為肉身；保羅顯然以基督取代羔羊.無可否認保羅用靈意,把舊約的羔羊用在基督身上.
\newline
\newline
耶穌和法利賽人辯論安息日的問題.法利賽人批評耶穌的門徒,在安息日掐麥穗吃是不對的.於是耶穌引證兩件事.
\newline
\newline
(1)(太十二1-8)大衛和跟從他的人逃避掃羅時,曾經到聖殿吃了祭司的陳設餅,那餅惟祭司可吃,大衛卻吃了.
\newline
\newline
(2)律法規定祭司,在聖殿裡犯安息日也沒有罪.如果耶穌非按靈意解釋,那麼,祂所引證這兩件事,與門徒在安息日掐麥穗吃有何關連?耶穌接著說「我喜愛憐恤,不喜愛祭祀.你們若明白這話的意思,就不將無罪的,當作有罪的了.」耶穌把三件事連貫而言.在以色列諸王中,大衛是合神心意的王,大衛曾受膏於神；他為了逃避掃羅的追殺,逃到祭司處吃陳設餅.按著神的憐恤來說是可以吃的；若按律法和條文則不可.為何祭司在聖殿裡犯安息日無罪?因為祭司的工作是在聖殿裡,而安息日的真義,是叫人放下一切的工作,專心事奉敬拜神.(中文安息日稱為禮拜日最為合適)祭司每天都在聖殿中事奉神,犯與不犯安息日,都是一樣的.所以這兩件事的精義,都是教導人敬畏神；神喜愛憐恤,根據祂的憐恤來體諒人,了解人,這是神的宗旨.門徒在安息日,因事奉忙碌,經過麥地就掐麥穗吃.神不喜愛我們以條文儀表事奉,而是喜愛我們以心靈誠實事奉；然後祂就按祂的憐恤施恩給我們.
\newline
\newline
(二)預表的靈意有聖經本身的解釋
\newline
\newline
預表的靈意不可隨便解,要完全照聖經本身的解釋來解釋.(約壹二27)論信徒的生活,恩膏乃指聖靈；但是約翰不解釋恩膏是預表聖靈,他直接以恩膏代表聖靈.實際上,恩膏在舊約是指聖膏油；會幕一切東西都要用聖膏油塗抹.惟獨聖所可以配料製造聖膏油,不准以色列人複製或仿製；經聖膏膏過的,便分別為聖,無論人或物都屬於神了.在新約,神不是用香膏膏基督徒,乃用聖靈；神賜聖靈給我們,「祂用膏膏我們,」意思即賜下聖靈給我們作憑據.(林後一 21-22)如按字意,「恩膏」和「聖靈」似乎互不相關；但聖經裡的預表,非我們可隨意解釋的,因預表是按靈意的.有時我們聽人說到聖經的預表時,我們不知有它一定的根據,自以為可隨便解釋；其實,靈意絕不可隨便解釋.比方說,亞當是預表耶穌基督,為什麼?(羅五14)「亞當乃是那以後要來之人的預像,」則指基督.(林前十五45)稱亞當為「首先的人亞當,成了有靈的活人.末後的亞當,成了叫人活的靈」.首先的亞當是被造的,末後的亞當是耶穌基督,是給人生命成了叫人活的靈；末後的亞當顯然是根據首先被造的亞當之預言說的.亞當和夏娃,預表基督和教會,亞當與夏娃二人成為一體,乃預表基督與教會合一的奧秘.(弗五30)保羅提到夫婦二人的合一,說:「夫婦二人成為一體.」他加上解釋:「這是極大的奧秘,但我是指著基督和教會說的.」可見亞當和夏娃的關係,乃是神設立婚姻制度的開頭；要用神聖的婚姻制度,來表明基督與教會合一的奧秘.因男女結合,有表明基督與教會合一的奧秘作用,所以婚姻制度必須看重為神聖；不合法的婚姻關係,會叫人對基督與教會合一的領會,發生錯誤.此類預表完全根據聖經本身的解釋.
\newline
\newline
以撒也是預表耶穌基督,是應許中要賜下的兒子,以撒是憑神應許而生的,神曾多次應許亞伯拉罕,「萬國要因你得福」惟有一次起誓向他應許,是在他獻以撒之後,(創二十二18)神起誓應許他「萬國都必因你的後裔得福.」後裔乃指耶穌基督；(加三16)指明:「神並非說眾子孫......乃是說那一個子孫......就是基督.」以撒就是神所應許那位要來的耶穌基督之預表.神要亞伯拉罕把以撒獻上,當他舉刀時,神差天使攔阻他,說:「不可在這童子下手......現在我知道你是敬畏神的了.」後來他看見一隻羊羔,就用以代替以撒.神藉亞伯拉罕的獨生子,而賜下自己的獨生子；因亞伯拉罕獻上獨生子,而得著神賜下獨生子的應許.
\newline
\newline
摩西有很多事可以預表耶穌基督,但聖經所題不多,惟提到摩西預表那將要來的先知.(申十八15-18):耶和華你的神要從你們弟兄中間,給你興起一位先知像摩西.後來在使徒行傳中,使徒們也見證稱耶穌是「那先知」.以色列人心中所盼望的那先知,是指應許中,要來傳神旨意的那位先知,就是猶太人的彌賽亞觀念.摩西預表將要來的那位先知,是根據聖經本身的用法.有一次,耶穌和尼哥底母談重生問題,祂說:「摩西怎樣在曠野舉銅蛇；人子也要照樣被舉起來.」可知摩西在曠野舉銅蛇,乃預表耶穌被釘在十架上；這也是根據聖經的用法,非隨私意表達,而是具有聖經的根據和用法.以色列人在曠野行程中,有雲柱火柱引導；過了紅海,因無水喝,摩西擊打磐石流出水來.他們絕糧時,神降嗎哪；凡此經歷均有預表的意思.根據什麼?(林前十1-4)「弟兄們!我不願意你們不曉得,我們的祖宗從前都在雲下,都從海中經過；都在雲裡海裡受洗歸了摩西.並且都吃了一樣的靈食,也都喝了一樣的靈水；所喝的是出於隨著他們的靈磐石,那磐石就是基督.」按史實,以色列人曾被雲柱火柱引導,紅海水分開,然後經過.保羅引用這段史實,說以色列人都在雲下海中受洗歸了摩西.因為他們過紅海之後,便由摩西帶領,完全脫離了埃及.這經歷預表基督徒受洗歸入耶穌基督,就向世界死了；和耶穌基督同死同埋葬.繼之,保羅又說,他們都吃一樣的靈食,他們事實上吃的是嗎哪.「靈食」是靈意的用法,表明基督徒領受靈命的糧食.「都喝了一樣的靈水......是出於......靈磐石,那磐石就是基督.」磐石流出活水,是預表基督為我們的滿足.磐石預表耶穌基督,是根據保羅自己的用法；我們無法解釋.保羅用以色列人的經歷為根據,應用在基督身上.所以聖經的預表,都有聖經本身的用法.關於此類例子很多,因時間關係,僅能提出其中我們較為熟悉之一二.
\newline
\newline
(三)讀預表應注意幾個要點:
\newline
\newline
1.預表只能按聖經的用法來領會.新約引用舊約所記,那才是預表的真正根據；若新約沒引用就不算預表.許多解經家都認為舊約的約瑟可預表基督生平；因他的遭遇和耶穌頗多相同.約瑟是他父親所愛,耶穌是天父所愛；他被兄弟賣了,耶穌被門徒猶大賣了.他被冤枉下監,出獄當埃及宰相；耶穌被人棄絕,死而復活被升為至高.類似的經歷很多,但新約沒引用作耶穌基督的預表；所以我們只能說約瑟的事很多像耶穌.不能說他是基督的預表,因新約沒引證.凡聖經有根據,我們才根據,才肯定；否則,我們不能根據也不能肯定.
\newline
\newline
2.預表實際含有預言的性質.預表實際上是另一種方式的預言,非用言語,乃用人物作象徵來說神將來必成之事.所有預言應驗之後,就不能再作別的解釋.我們常說舊約時代,是神給我們觀看的一個模型；各不同的預表則各不同的模型.說明共同的題目,並神所要為我們成就的救恩；救恩的主要內容,是耶穌基督和祂所救贖的教會.
\newline
\newline
3.預表的主要功用.乃說明與救贖有關的真理,亦可證明新舊約的連貫性.
\newline
\newline
4.預表必定不如所表明的.如能完全表明就不算為預表,正如模型的房屋,不等於實際的房屋；模型房屋只供人參觀,實際房屋才供人居住.所以舊約用以預表耶穌的,毋須每點都像耶穌基督；以撒預表耶穌,其要點乃應許中要來的那一位,餘者都不能作預表.例如以撒貪吃以掃狩獵的野味,那當然不能用來預表耶穌.摩西乃預表基督的那先知；因他是為神全家盡忠的僕人；但我們都知道摩西因發怒,神罰他不得進迦南.神在歷史上收集許多不同的人物特出之點像耶穌,可用作表明耶穌基督的工作；連在一起才足證明,並非某個人的生平便可；因為耶穌是集所有人之最好的都包括在內.實際上,人所有最優之處,都是從祂至善至美的事上分出來的一小點.
\newline
\newline
預表不能單獨作教義上,或信仰上的根據,只可與聖經明白的教訓,配合作信仰的根據.
\newline
\newline
我們信仰的根據,是根據使徒明白的教訓,主耶穌明白的命令.
\newline
\newline


\newpage

\section{第50屆~港九培靈會~陳終道~第九講}
\label{sec:724}
\textbf{聖經的一貫性}
\newline
\newline
連結: \href{https://www.hkbibleconference.org/session-message/view/724}{\texttt{https://www.hkbibleconference.org/session-message/view/724}}
\newline
\newline
\hyperref[sec:723]{< < < PREV SERMON < < <}
~
\hyperref[sec:toc]{[返目錄]}
~
\hyperref[sec:750]{> > > NEXT SERMON > > >}
\newline
\newline
路加福音 24:27-24:27
\newline
\begin{longtable}{cl}
\hline
\hline
章節 & 經文 (和合本修訂版)\\
\hline
24:27 & \begin{tabularx}{0.7\textwidth}{X} 於是，他從摩西和眾先知起，凡經上所指著自己的話都給他們作了解釋。 \end{tabularx} \\ \\
[1ex]
\hline
\hline
\end{longtable}
在我心靈裡,總覺得今天的基督徒,缺乏一種認識,沒有體會到聖經是神完整的啟示.
\newline
\newline
神的啟示雖不是一時間給我們,整部聖經的完成,幾經一千五百年之久；可是裡面的信息是一貫的,是有主題的一部書.聖經非零星湊合的經卷,實在是神有計劃寫的書；有些地方經文也可見於此.(羅馬十五4):「從前所寫的聖經」(乃指舊約)「都是為教訓我們的.」「我們」乃指新約時代的信徒；以外邦為主的教會.神有計劃,有目的完成這部聖經,為要教訓我們.(林前十11):「他們遭遇這些事,都要作為鑑戒；並且寫在經上,正是警戒我們這末世的人.」「他們」乃指以色列人,他們的歷史,有蒙恩的經過,也有失敗的歷史；保羅說,都是為了警戒我們需寫在經上的.換言之,以色列人的歷史,他們的遭遇,經過,是有意挑選記載在聖經上的；為著可作我們的鑑戒.這是說明寫聖經是神原有的計劃,從人類歷史中,選出可用的材料記於聖經上.
\newline
\newline
另一方面給我們看見,舊約聖經和新約信徒關係的密切.舊約以色列人的經驗,選記在聖經作我們的鑑戒；作為新約使徒講教訓時,可引用的例子.
\newline
\newline
新舊約含有合一性,含有一貫性.「聖經既然預先看明,神要叫外邦人因信稱義；就早已傳福音給亞伯拉罕說,萬國都必因你得福.」(加三8)亞伯拉罕的時候,尚未有聖經,可以說是神傳福音給他,對他說話.保羅所以這樣說,可見神對亞伯拉罕所說的話,也就是神要啟示給人聖經內容的一部份.當神未啟示之前已有腹稿,神對亞伯拉罕所說的話,也就是後來所寫的聖經.神把聖經給我們,不是一下子全部給一人,而是經過約一千五百年的時間,經許多朝代,在不同的時代中選召祂合用的器皿,把啟示給予一部份,逐漸啟示給人.
\newline
\newline
聖經確是神要給我們的書,不過因歷時長久,且作者眾多；所以完成之後,人常誤為聖經是許多不同的經卷湊合而成.其實,由始至終乃神原有的計劃要給予我們的完整啟示.很少基督徒能在短時間內,一遍又一遍地讀聖經；所以很少基督徒對於聖經有完整的概念,只是抓住許多的片段；因此對聖經的了解難有完整的觀念.
\newline
\newline
聖經確是一部極奇妙的書,並無組織編輯委員會,而且作者多不同時代,甚至有相距數百年者；但聖經的內容卻有一貫性,可見有位莫測的神管理整個聖經的寫作.
\newline
\newline
我們可從三方面來看聖經的一貫性:
\newline
\newline
(一)從聖經的結構上看
\newline
\newline
聖經之第一卷是創世記,歷史的人物按次連串而下；亞當生該隱和亞伯.聖經記載他的後代直到挪亞,他的後裔代代而下,直至亞伯拉罕,以撒,雅各.出埃及記第一章銜接創世記末章,創世記末章記載約瑟在埃及當宰相多年直到壽終.出埃及記一章記約瑟死後,以色列人在埃及遭遇逼迫.利未記,民數記,所載歷史都是緊接的.申命記溫習以色列民出埃及的經過；摩西五經在結構上非常緊湊.約書亞記記載約書亞被神揀選,第一章八節:「這律法書」指前摩西之五經,「不可離開你的口」,這說明約書亞記緊接律法書；事實上,約書亞繼承了摩西的地位,從約書亞開始的歷史,都有記載以色列人是否遵行律法書；遵守都蒙福,違背者遭罰.到了士師時代,非常混亂,因他們沒聽從神藉祂僕人的吩咐,消滅迦南異族.列國時代,列王立國；凡君王遵守神給摩西的律法書,則國家蒙福；不遵守者,國家衰微.整個歷史書所載都是根據摩西的律法書；描述以色列百姓對律法書是否遵行,結構緊湊,絲毫沒中斷.假如沒有摩西五經,和約書亞記；士師記的事就不會發生,而所記的事就沒有背景.以色列人在舊約時代,祭司和君王均為世襲制度；以色列中利未人必須屬亞倫的子孫方可當祭司；這兩個世襲的系統,祭司管理宗教方面；君王管理政治方面.所以它的歷史貫串而下,使整個舊約的啟示有連貫性.
\newline
\newline
新舊約的關係非常密切.新約的馬太福音一章一節論耶穌基督的家譜說:「亞伯拉罕的後裔,大衛的子孫.」可見這兩句話,把新約和舊約連繫起來.日前我曾提過,亞伯拉罕的獨生子以撒是應許中要來的兒子,乃預表耶穌基督；而耶穌基督,就是成全大衛的寶座,直到永遠的那一位.四福音論耶穌基督的生平,多處記載耶穌基督的事；都題「經上說」或「應驗經上所說.」可見新約所記乃舊約的預言,是連貫的.使徒行傳也是一樣,第一章乃是接續四福音的記載.書信是解釋福音,有人以為兩者有衝突之處；因為四福音注重行為,書信注重信心.其實,我們忽略了一件事,寫四福音最初的兩人,都和書信有密切關係.路加福音的作者路加,是保羅最得力的助手；保羅傳福音,他始終跟著.馬可福音的作者馬可,是彼得最得力的同工,彼得稱他為福音的兒子；保羅在殉道之前,非常欣賞馬可.咐吩提摩太要將馬可帶來,可見馬可,是保羅和彼得所共同器重的.
\newline
\newline
馬可所寫的福音和保羅及彼得所寫的書信,為何不覺有衝突呢?因為福音和書信,乃從不同的角度來講明救恩.福音只注重記錄事實.書信則講明救恩的原理.書信更講明這位降世的耶穌基督的生,死,復活；把基督所成功的救恩之道解釋清楚.
\newline
\newline
有人以為啟示錄是新約特別的預言書.但是從開始所記的和新約舊約都有關係；論到七燈台就把舊約的燈台,和新約的教會連起來.第二至第三章提到七個教會時,引證舊約的巴蘭,耶洗別,提到生命樹.這些都是新約的背景,而啟示錄的災禍,昔時在埃及經已行過.整卷聖經在結構上是一致,連貫不可分開的.
\newline
\newline
(二)從聖經的立場上看
\newline
\newline
凡「書」都有它的立場,有的「書」是政治立場,有的「書」是戀愛觀的立場；而聖經在立場上是一貫性的.
\newline
\newline
1.對不信者都看為同站一邊的.無論在歷史上或使徒時代,凡不信者都看為同類的人；信者也都看為同類的人.比方敵擋神的,主耶穌說:「從義人亞伯的血起,直到你們在殿和壇中間所殺的巴拉加的兒子撒迦利亞的血為止.我實在告訴你們,這一切的罪,都要歸到這世代了.」(太二十三35-36)聖經豈非告訴我們,父不擔當子的罪,子不擔當父的罪嗎?自己的罪應自己擔當,為何耶穌說:「......這一切的罪,都要歸到這世代.」亞伯是亞當的兒子,撒迦利亞是列王時代稍為後起的祭司；相距時間很遠,他們兩人被殺的血要歸於耶穌的世代.其實,耶穌這話是因祂的觀念.凡敵擋神,犯罪的人,從亞伯起至撒迦利亞,至耶穌一直至現在,都是一個世代；人類犯罪背逆神,根本沒有任何新的世代；所以耶穌說:「都要歸到這世代.」意即歸併一起受刑罰；因為人類一向都是敵擋神的.至於使徒保羅的觀念,他在(提後三5)提到:「有敬虔的外貌,卻背了敬虔的實意,這等人你要躲開.」他舉例出埃及記七至八章說明,有兩個行法術的如何抵擋摩西,這等人也照樣抵擋真道；保羅提醒要慎防教會中敵擋真道的人,乃指當時以弗所的假使徒.啟示錄二章,主稱讚以弗所教會試驗認清假的使徒.我們應注意,假使徒就是教會圈內的假弟兄和異端；而抵擋摩西的二人是行法術的,在時間上相隔很遠,且一為異教之流,一為教會圈中之輩；但保羅以二者同日而語.他們所處時代雖不同,表現也不同,背景卻是一樣；聖經把此類人物,看為同站一邊的.彼得也說,從前在百姓中間有假先知起來,將來在你們那裡,也有假師傅私自引人陷入異端；他把將來在信徒中產生的假師傅和從前的假先知同等看待；從前的假先知是以色列中的人；將來要出現的假師傅,是教會圈中的人.
\newline
\newline
2.對一切信主的人都看為屬神的.一次,主耶穌稱讚一個百夫長的信心,其中有句責備猶太人的話說:「從東從西,將有許多人來,在天國裡與亞伯拉罕,以撒,雅各,一同坐席；惟有本國的子民,竟被趕到外邊黑暗裡去.」從各地各族來的有信心的人,都要與亞伯拉罕他們一同坐席.反言之,以色列人按肉身雖是神的百姓；但沒信心的,照樣被趕到外邊黑暗裡去.主耶穌把有信心者看為同類的人,保羅也是如此；他提到以色列人餘數時,把以利亞時代七千不向巴力屈膝的人,引用到當時信徒身上.可見保羅的觀念,也是把有信心者歸於一邊.雅各也是如此,他說:以利亞是和我們一樣有信心的人,他所信的神也是我們所信的神,他的禱告和我們一樣有功效；他憑信心蒙神悅納,我們也是這樣.
\newline
\newline
(三)從聖經主題信息上看
\newline
\newline
由創世記至啟示錄唯一題目是:「耶穌基督和祂的救恩.」亞當預表耶穌基督.亞伯用羊羔獻祭蒙神悅納,羊羔預表耶穌基督.施洗約翰出來傳道時,說:「看哪!神的羔羊,除去世人罪孽的.」羔羊實際上乃指基督.希伯來書稱讚亞伯說:「亞伯因著信獻祭與神,比該隱所獻的更美.......」信是有對象的,信那後者要來的羔羊.耶穌說:「亞伯拉罕仰望我的日子,既看見了就歡喜.」猶太人聽見這話,大不以為然,用石頭要打耶穌.亞伯拉罕所信的,是後來的耶穌基督；他的信心從等候所應許的以撒表現出來.摩西逃到曠野牧羊,希伯來書說的摩西──「他看為基督受的凌辱,比埃及的財物更寶貴.」摩西所信的是基督.以色列人出埃及,得神給他們屬靈的基業迦南地；他們有個盼望,有個國度,有座城；不是他們自己建立的,他們等候那有根基的城.迦南預表在耶穌基督裡所要得的福份；希伯來書四章直接以「安息」取代「迦南」.他說約書亞沒有真正叫以色列人得到安息；因約書亞領以色列人得迦南地,沒有把當地敵人消滅；所以他說以色列人要進入迦南,實則指我們在基督裡要得到屬靈的福份,那福份是耶穌基督,成功了救恩所帶來給我們的.
\newline
\newline
整部聖經無論新約舊約,都是以耶穌基督為中心.行傳是見證耶穌基督.書信是講解耶穌基督.啟示錄是預言基督再來.聖經按結構,立場,信息,都是一貫的.
\newline
\newline
基督徒必須注意讀全部聖經.不少基督徒以專題方式查經,這樣可以節省時間.雖然如此,但必須逐卷而查,因專題查經等於讀速成班.速成班就非正規的教導,思想易受講者牽制而不自知.如果逐卷讀,受著上下文管理,不易患斷章取義；且具有全面性的概念,這是今日基督徒非常需要的,讀法當然須費工夫,斷非不勞而獲.
\newline
\newline
讀聖經是基督徒的終身大事,應列為畢生不可缺少而且是經常性的大事,我們當定下讀經計劃,才能在神手中把神的話融匯貫通,作為神手中的器皿.
\newline
\newline
求神幫助我們,讓我們愛慕主的話,得到主的話當食物吃了；我們行事為人,知道是不是神的旨意.
\newline
\newline


\newpage

\section{第50屆~港九培靈會~高文~第一講　承認罪過}
\label{sec:750}
\textbf{第一講　承認罪過}
\newline
\newline
連結: \href{https://www.hkbibleconference.org/session-message/view/750}{\texttt{https://www.hkbibleconference.org/session-message/view/750}}
\newline
\newline
\hyperref[sec:724]{< < < PREV SERMON < < <}
~
\hyperref[sec:toc]{[返目錄]}
~
\hyperref[sec:751]{> > > NEXT SERMON > > >}
\newline
\newline
歷代志下 20:1-20:12
\newline
\begin{longtable}{cl}
\hline
\hline
章節 & 經文 (和合本修訂版)\\
\hline
20:1 & \begin{tabularx}{0.7\textwidth}{X} 此後，摩押人和亞捫人，連同一些米烏尼人來攻擊約沙法。 \end{tabularx} \\ \\ \relax
20:2 & \begin{tabularx}{0.7\textwidth}{X} 有人來報告約沙法說：「從海的那邊，以東有大軍來攻擊你，看哪，他們在哈洗遜‧他瑪，就是隱‧基底。」 \end{tabularx} \\ \\ \relax
20:3 & \begin{tabularx}{0.7\textwidth}{X} 約沙法懼怕，就定意尋求耶和華，在全猶大宣告禁食。 \end{tabularx} \\ \\ \relax
20:4 & \begin{tabularx}{0.7\textwidth}{X} 於是猶大人聚集，求耶和華幫助，甚至他們從猶大各城前來尋求耶和華。 \end{tabularx} \\ \\ \relax
20:5 & \begin{tabularx}{0.7\textwidth}{X} 約沙法站在猶大和耶路撒冷的會眾中，在耶和華殿新的院子前， \end{tabularx} \\ \\ \relax
20:6 & \begin{tabularx}{0.7\textwidth}{X} 說：「耶和華－我們列祖的神啊，你不是天上的神嗎？你不是萬邦萬國的主宰嗎？在你手中有大能大力，無人能抵擋你。 \end{tabularx} \\ \\ \relax
20:7 & \begin{tabularx}{0.7\textwidth}{X} 我們的神啊，你不是曾在你百姓以色列面前驅逐這地的居民，將這地賜給你朋友亞伯拉罕的後裔永遠為業嗎？ \end{tabularx} \\ \\ \relax
20:8 & \begin{tabularx}{0.7\textwidth}{X} 他們住在這地，又為你的名建造聖所，說： \end{tabularx} \\ \\ \relax
20:9 & \begin{tabularx}{0.7\textwidth}{X} 『若有禍患臨到我們，或刀兵的懲罰，或瘟疫饑荒，我們在急難的時候，站在這殿前向你呼求，你必垂聽並且拯救，因為你的名在這殿裡。』 \end{tabularx} \\ \\ \relax
20:10 & \begin{tabularx}{0.7\textwidth}{X} 現在，看哪，以色列人出埃及地的時候，你不容許以色列人侵犯亞捫人、摩押人和西珥山人，以色列人就離開他們，不滅絕他們。 \end{tabularx} \\ \\ \relax
20:11 & \begin{tabularx}{0.7\textwidth}{X} 看哪，他們這樣回報我們，要來驅逐我們離開你賜給我們為業之地。 \end{tabularx} \\ \\ \relax
20:12 & \begin{tabularx}{0.7\textwidth}{X} 我們的神啊，你不懲罰他們嗎？因為我們無力抵擋這來攻擊我們的大軍。我們不知道該怎麼做，我們的眼目單仰望你。」 \end{tabularx} \\ \\
[1ex]
\hline
\hline
\end{longtable}
兄弟能夠參加培靈會的五十週年金禧紀念,對我來說實在是莫大的福份；因為能與其他講員,一同分享為基督所過的生活.今天晚上,當詩班歌頌神的時候,我的心已高高的向著我的神.當我想到天堂裡的音樂,到那時候我們惟一所知道的音樂,都是讚美我們的主.各位在唱詩時已學了天國的言語,這是何等大的喜樂.
\newline
\newline
這十天晚上,我願意將耶穌基督賜給我們的生命－優美的人生,與大家講解.優美的人生乃是神願意每一個屬乎祂的兒女所過的生活.若今天我們沒有支取在基督裡的豐盛,就是有了權利而沒有好好的使用.我們何等需要再次奮興過來!
\newline
\newline
「奮興」的意思就是在神面前再次活過來.惟有在基督裡腳踏實地的過生活,才能應付這世代的危機.中文的「危機」是用兩個中國字拼成；危－危險,機－機會；這詞已含了許多智慧.確實地,一個危機擺在面前,也是將機會帶到面前；一個危機叫我們知道必須倚靠神的幫助.只有在黑夜裡,我們才會舉頭望見星光.今天這時代,乃是整個世界面對危機的時代.如果今天給每人有新的機會,就是你我所能佔有的時刻.當見到這世界的危險時,讓我們仰望倚靠神所給予的機會!
\newline
\newline
歷代志下廿章,就將這個真理再次刻劃出來.一個細小的王國──猶大,正被三個異教之國所侵略,他們的處境看來沒有盼望,連他們的王都懼怕起來!請留心第三節:當他在懼怕中時,他就定意尋求耶和華；在整個猶大國裡宣告禁食.這個危機將他帶到一個情景,他知道一切的盼望惟有在乎神.於是召聚百姓,叫他們預備心向著神；然後站在他們當中,藉著一個禱告,將他們的心都提起來.第六節開始將當時的情形描寫至第十二節,在第十二節結束時,他說:「我們的神阿,你不懲罰他們麼?因為我們無力抵擋這來攻擊我們的大軍.我們也不知道怎樣行,我們的眼目,單仰望你.」這裡有兩個復興的原則；承認自己的軟弱,無力量；又承認堅信神的能力.除非認清自己的軟弱無助,否則不能支取神恩典的豐富.
\newline
\newline
人歸向神之第一次記載,在(創四26)；塞特生了一兒子,起名叫以挪士.「以挪士」是軟弱無助之意.這事以前,人類剛才經歷到第一次死亡,因為亞伯之死不是死於自然,是被他的兄弟所殺.如果這情形繼續演變下去,整個人類就會被刀劍暴力摧殘.在這時,人纔醒悟回轉歸向神.塞特將他的兒子起名叫以挪士,因為他認識自己的軟弱,無能為力；需要神,也影響了那一世代的人,求告耶和華神.
\newline
\newline
當主要將奮興帶到祂子民的時候,這些事情都要發生；就是我們必需承認自己是軟弱無助的,認定我們的幫助是從神而來.要誠實對己是十分困難的,今日在教會裡,多少人就是這樣過一個虛假,偽裝的生活；只有經歷宗教形式,而沒有經驗裡面的能力；好像一架車,剛出廠,但三個車呔已洩氣,沒有電油一樣.又像啟示錄中老底嘉教會:他們是富足,發了財,以為自己一樣都不缺.但神見到他們心靈的空虛,充滿污穢；只有外表宗教的儀式,卻失去了作工的能力.主勸勉他們買眼藥擦眼,使他們能看見,誠實看清楚自己的境況.主呼召他們悔改,除非他們悔改,否則主要將他們從口中吐出去.因為主深愛他們,不忍他們離開祂.當我們察覺自己有罪,原來是主的愛,想我們能歸向祂,求祂的幫助.
\newline
\newline
在我任教的大學裡,我們也經歷過這種情形,我們誠實地承認自己的軟弱,以求神的復興臨到.一九七○年二月,在早會中,有同學見證神如何在他們生活中作工.一位高班同學坦白承認自己是個偽裝,假冒為善的人；因他曾用嘴唇承認耶穌,但在心裡卻仍未承認主的主權；他的眼淚開始流下在他的臉上,他請求同學原諒.他見證在早會前一晚,他已在神面前對付過,向主承認他的罪過；主就賜他清潔的心,使他經歷到一個真正基督徒所有的喜樂.他的誠實就鼓勵了其他同學都誠實起來,隨後其他同學都紛紛站起來承認自己的罪.非常明顯,神的靈就降到這群同學全體的身上.有人甚至要求在講台上用擴音器作見證,承認自己的罪；許多同學都走到台前等候作見證.一時間,整個房子內都充滿了喊叫聲和歡呼聲,彼此的仇恨就在神聖的愛裡溶化了,他們心靈裡得著釋放和喜樂.曾有一位教授站起來說；六年來,我自己靈性失敗,他請求同學寬恕他.又有一位師母在聚會中說她對她的教會充滿憤恨,求她的會友寬恕他.隨即她的丈夫,那作牧師的也站起來,個別提名請會友原諒他所做錯的事；然後他見證神如何將愛心充滿他的心.當日大概有一百多人站在講台上求神的寬恕.
\newline
\newline
在校園裡這樣的奮興,從早到晚,每日如是；一百八十五小時的復興並沒受攔阻.我們只好暫停上課,再沒有人打籃球,或是約異性朋友.惟一所關心的,就是必須與神恢復優美的關係；在禱告中等候神,又出外宣告神美好的工作.復興的火燄一直伸展到其他校園裡,他們都感覺到這所學院發生的復興在他們身上的影響力.隨之同學們分批到千間以上教會作見證,以至那一帶教會都有了復興.當聖靈管理一個聚會時,我們會見到神是大而可畏的.只有經驗如此復興的人,才會體驗神如何藉著破碎的心靈來彰顯祂的作為.我們要將假面具除去,真誠的相待.有人說:復興－是開始不計較別人的罪,乃是承認自己的罪.
\newline
\newline
今天晚上,你有罪要向神承認嗎?你心裡有憤恨,嫉妒嗎?你的生活是否仍帶著假面具?你雖是教會會友,但在你心靈裡所有的都是失敗而已!在你裡面失去喜樂,沒有愛心的洋溢.你所過的生活,實在不像基督要你所過的豐盛生活.你有沒有向神承認自己罪惡的需要呢?承認自己是軟弱無助的呢?約沙法向神說:「神阿,我們不曉得怎樣行,現在這敵軍來了,我是軟弱無助的!」然而約沙法還向神這樣說:「耶和華阿,我們的眼目,單單仰望你.」我們必須承認自己的軟弱無助；但更重要的是讓我們堅定相信神的信實,將注意力放在偉大的神,祂的榮耀,權柄上.
\newline
\newline
列王記下第六章另外記載一件事情,十五至十六節:某天,以利沙見到在他城外有敵軍圍繞他們,目的是要捉拿他.當那少年僕人見這情況時,就跑到主人以利沙面前,說:「主人阿,我們要怎麼辦呢?」你要明白為何這少年人懼怕起來?因為他們沒有防禦敵軍的進襲.但以利沙知道他的幫助是從高處,是從神而來.於是他轉過身來對少年人說:「不要懼怕,與我們同在的,比與他們同在的更多.」之後,他開始禱告,是一個非常突出的祈禱.他說:「主阿,求你開這少年人的眼,讓他能看見!」你有沒有聽過有人這樣祈禱呢?神卻應允了禱告.聖經告訴我們:正當少年人看的時候,見到四周圍有耶和華所賜的火車火馬圍繞他們的城.這少年人現在所見到的,正是先知早已看見的異象.神早已知道這個形勢,神十分願意幫助祂的兒女；神深愛祂的子民,願意將拯救給予屬祂的人.但願我們張開眼睛,看見神的能力,榮耀和權柄!
\newline
\newline
多少時候我們只是四顧環境,要從人得著力量；只看見自己的難題,自己的麻煩.我們何等需要將注意力放在耶穌基督的榮耀上!那位永生的耶和華,是我們的主,是整個宇宙的創造者.不要懼怕,也不要氣餒.我們的神明白我們的情況,是人不能測度的.祂將能力賜給軟弱的人,看哪!這是神的話語,在祂伸展的手臂裡,祂能創造天和地.在耶和華豈有難成的事呢?祂是我們的神,我們的保衛者.祂會將恩典,榮耀賜下,對那行動正直的人,祂未嘗將一樣的好處留下不賜給屬祂的人.在神的話裡還有成千上萬的應許.
\newline
\newline
我們既知道神如此深愛我們,就讓我們的注意力放在祂永恆的話語上.若耶和華幫助我們,誰能敵擋我們呢!祂甚至將自己獨生子捨棄,拯救我們,難道祂不會將一切白白賜給我們嗎?
\newline
\newline
各位,是否你有試探難掙脫?是否因世界的難題,嚇怕你不能為主作見證呢?留心歷代志下廿章:當約沙法王禱告完,猶大眾人與王都一同站在耶和華面前；聖靈感動其中一青年說出預言.(二十15),他說:「不要因這大軍恐懼驚惶,因為勝敗不在乎你們,乃在乎神!」17節:「你們不要爭戰,要擺陣站著,看耶和華為你們施行拯救.」十八節:「約沙法就面伏於地,俯伏在耶和華面前；他深信神的先知所講的話.猶大全軍和所有耶路撒冷的人,都俯伏敬拜耶和華神.」留心廿節:「現在敵人進來了!翌日清早,這作王的站在眾百姓面前說:你們要信耶和華你們的神,就必立穩；相信祂的先知,就必亨通.」勸告完後,王就設立歌唱的人,吩咐這詩班唱詩稱頌神的慈愛永遠長存.在敵人面前能夠歌唱.各位,你能否想想這美麗的情景!亞捫人,米烏利人和摩押人一同向這猶大王國進攻,在軍事上來說,他們的情況是沒有指望的,但現在他們不會因四周圍的敵人懼怕,原因是他們聽見神的話,又深信祂的話.王說:現在我所關心的,是這詩班；因為現在是機會,讓他們的生命發出詩歌.我的軍隊殘缺,人數又少,這詩班要作他們的前鋒,看清楚戰場的情形.但他們要唱出安息日所唱的詩歌,就是要讚美耶和華,因祂的慈愛永遠長存.請注意廿二節:當眾人剛唱歌讚美的時候,耶和華就派伏兵擊殺敵人,所有敵軍都戰敗了；他們見到神的子民唱詩讚美神,就竟然驚慌至互相殘殺.
\newline
\newline
當世人看見基督徒竟然提起來歌頌這位全能的神,他們都要驚惶起來.在戰爭的歷史上,寫下了難忘驚奇的一頁!各位,這是歷史的事實,這是一個明證,要將我們今日的屬靈生命境況來作一個儆惕!在各種不同的環境下,你我都要面對仇敵.今日多少人受到異教的影響；我們的生命要受威嚇；容易叫我們灰心,而憐愛自己離開主.讓我們認定,靠自己是不能戰勝敵人的,沒有人能靠己活出基督徒的生命來；教會必須將她的眼目單單仰望神；讓主充滿我們,使用我們,好活出優美的人生來!
\newline
\newline
衛斯理約翰垂死時,一直在唱歌.他說:「當我呼吸最後一口氣時,我仍要歌頌我的救贖主.」當我們心靈面對死亡時,惟有頌讚才能喚醒我們.他見到周圍的人為他灰心難過,他就對他們說:「一切最好的,乃是神仍與我們同在.」他告訴那恐懼的人一同跪下與他禱告,向神發出讚美.這位垂死的人伸開他的手,望著說:「我仍要讚美祂.」病房裡的人,一而再,再而三聽見衛斯理垂死的呼叫:「我要讚美祂,.......」讓我們不論活著或在死亡的邊緣,都要讚美祂!
\newline
\newline
各位,今日你的生活是否如此呢?如果你還沒有這樣經驗,你需要復興!你需要進入神深願你過的生活景況裡──一個得勝的生活!
\newline
\newline


\newpage

\section{第50屆~港九培靈會~高文~第二講　活在和諧禱告生活中}
\label{sec:751}
\textbf{第二講　活在和諧禱告生活中}
\newline
\newline
連結: \href{https://www.hkbibleconference.org/session-message/view/751}{\texttt{https://www.hkbibleconference.org/session-message/view/751}}
\newline
\newline
\hyperref[sec:750]{< < < PREV SERMON < < <}
~
\hyperref[sec:toc]{[返目錄]}
~
\hyperref[sec:752]{> > > NEXT SERMON > > >}
\newline
\newline
馬太福音 18:15-18:20
\newline
\begin{longtable}{cl}
\hline
\hline
章節 & 經文 (和合本修訂版)\\
\hline
18:15 & \begin{tabularx}{0.7\textwidth}{X} 「若是你的弟兄得罪你，你要去，趁著只有他和你在一起的時候，指出他的錯來。他若聽你，你就贏得了你的弟兄； \end{tabularx} \\ \\ \relax
18:16 & \begin{tabularx}{0.7\textwidth}{X} 他若不聽，你就另外帶一個或兩個人同去，因為『任何指控都要憑兩個或三個證人的口述才能成立』。 \end{tabularx} \\ \\ \relax
18:17 & \begin{tabularx}{0.7\textwidth}{X} 他若是不聽他們，就去告訴教會；若是不聽教會，就把他看作外邦人和稅吏。 \end{tabularx} \\ \\ \relax
18:18 & \begin{tabularx}{0.7\textwidth}{X} 「我實在告訴你們，凡你們在地上所捆綁的，在天上也要捆綁；凡你們在地上所釋放的，在天上也要釋放。 \end{tabularx} \\ \\ \relax
18:19 & \begin{tabularx}{0.7\textwidth}{X} 我又實在告訴你們，若是你們中間有兩個人在地上同心合意地求甚麼事，我在天上的父必為他們成全。 \end{tabularx} \\ \\ \relax
18:20 & \begin{tabularx}{0.7\textwidth}{X} 因為，哪裡有兩三個人奉我的名聚會，哪裡就有我在他們中間。」 \end{tabularx} \\ \\
[1ex]
\hline
\hline
\end{longtable}
某些年前,英國威爾斯地方,神復興的火正在燃燒；千萬人歸向神,有人要找出復興之秘訣,於是參加了其中一個聚會；詢問那復興之原因.其中一人站立起來,舉起雙手說:根本沒有秘訣可言,「你們尋求就必得著!」這是神賜下復興之方法.請留心(太十八19-20),聖經特別應許:當有兩三個人,同心合意禱告的時候,神就在他們中間.此外,耶穌在這裡還刻劃三件重要的事:
\newline
\newline
(一)必須有同一個心意的合一.原文意思是用在交響樂曲,大家在音響上是合一,和諧的.例如兩件樂器同奏,豈不是比一件樂器更加悅耳嗎?但不容易,必須付上代價才能得著.若有不和諧的表現,就應避免.在一個交響樂裡,若有某樂章或樂器出現困難,不協調,整個樂曲就受影響.同樣,在我們的禱告生活裡也是如此.整體禱告比個人更重要,而同一時候,個人本身的價值並沒有減少.事實上,可以藉此學習與整體融洽.在教會裡也是如此,每一個人在基督的身體上有一重要位置,但必須要有和諧.在婚姻上也是一樣,神造男造女,叫他們彼此補足；作丈夫和妻子分享肉身生命時,在他們的禱告上也要分享.
\newline
\newline
禱告合一廣義來說:除夫妻外,整個家庭也當合一.若家裡能在愛心禱告上,彼此合一,是何等美好的事,他們都要成為神的殿.這樣,在教會也可以找到；若整個會眾彼此同心聚集,聖靈就降臨在其中.若不能全體會眾一同禱告,可否先與兩三人一同開始.
\newline
\newline
合一禱告必先有靈裡的合一.這段經文之上下文,提到教會裡糾正弟兄的錯處.(太十八15-18),提及對犯罪弟兄的紀律,要親自到他面前幫助他；若不聽從,再次帶一些人去見他；經多人勸告後仍不聽從時,才帶到會眾前.這是全體教會應有相愛的心,也是經常關心的事.
\newline
\newline
耶穌在禱告的事上還有其他應許:凡你們在地上所捆綁的,在天上也要捆綁；凡你們在地上所釋放的,在天上也要釋放.(太十六19)又說:「若你們中間有兩個人在地上,同心合意的求甚麼事,我在天上的父,必為他們成全.」(太十八19)
\newline
\newline
(二)必須在身份上與基督合一.要奉耶穌的名聚集.(約十四13-14,十六24)這禱告的內容正反映主的個性,和祂的身份,異象,降世救人目的.奉主名禱告意思是正如主在天上的禱告一樣.這樣比起單講幾句話,更能討主歡喜；因為我自己的意思與神的旨意是合一的.為此,禱告乃是神採取主動,我們只不過是個器皿,透過聖靈將禱告獻給神.換言之,在禱告裡,先察驗自己,我的意思是否與神的旨意一致.首先,認識我與神是否合一,現在願意照祂的意旨；這不是頭腦上的同意,而是一個要離棄罪惡的決心.
\newline
\newline
一九四九年,在快樂島上,人們日以繼夜,禱告求賜下復興.某夜,一弟兄讀出神的話:誰能登耶和華的山?誰能在耶和華面前站立?就是那些手潔心清的人.那時,他們就彼此說:我們這樣聚集,若心裡對神不正直；倒不如先察驗自己,心和手是否聖潔?於是,他們就發現自己有某些罪要承認,他們開始與神取得和諧.他們的生命先復興,然後復興的火就在島上燃燒起來.成千上萬的人受影響,他們都歸進神的國度裡.
\newline
\newline
真正的禱告必須付代價.當主禱告說:「父阿!不要照我的意思,乃是照你的意思.」主所付的代價就是自己的生命.同樣,我們也要付出代價,向著十字架,將「老我」除去.我們的心必須校正,與主的旨意配合,這才能經驗與主同在的喜樂.
\newline
\newline
(三)必須得著神賜下答案的確據.我們祈求就必得著,主在新約聖經中有多處的應許:(約十五7,十六23；太二十一22).這樣的禱告在我們屬靈的練習裡,實在太膚淺了.多少時候,我們向神討價還價,沒有信心.難道神不知道祂兒女們的需要嗎?禱告不單只有交通,還有是我們在祂面前的蒙福.當然,並不是說凡奉主名祈求,就立刻得著答應.神可以延遲答覆,藉此教導我們學忍耐的功課.只要活在聖靈裡,無論求什麼,祂必會垂聽.因為我們所作的,正是祂要求我們去做的.為此,聖靈充滿的生命,乃是以禱告為主的生命.在新約初期教會的禱告,沒有一個是神沒有答覆的；乃因他們肯同心合意的求.
\newline
\newline
在我任教的校園,發生復興時,他們終日所關心的是禱告.校園內每一角落,都有人組織禱告聚會.在地庫裡,學院行政大樓內的辦公室,宿舍,甚至在洗手間也有人禱告.在這奇異的和諧裡,神的同在是明顯的.
\newline
\newline
我最初認識的「小組禱告」,是當我正讀大學時.某天晚上,一位由貝立大學來的弟兄向我說:我們要復興.當我們分手告別時,他提議一起禱告,對我來說,很不自然和不尋常；因為這樣開聲禱告,又是在宿舍內,恐怕別人都聽見了.所以第二回:我提議在別處禱告,於是我們在課室裡跪下禱告.將我們心裡所盼望校園復興的事向神陳明.不久,我發覺那弟兄將話題轉向自己,不再是單為別人求；他為自己的罪,在神面前對付.特別是他的驕傲.當時,我聽見有叩鐵的聲音,我不奇然張大眼睛,見他在自己手錶上取下一塊金屬片,然後把它拋出窗外.原來這是他在運動上贏得來的紀念品,因為他這紀念品,他非常驕傲.如此真誠懇切的態度,我從未見過的.當時我亦因此而在神面前對付自己,求主粉碎那假冒為善的我.我發現我們是融合在與神和諧合一的景況裡.這次就成為我生命的轉戾點.自此之後,我在宿舍,校園裡找人和我一同禱告,直到如今.
\newline
\newline
各位,你今天晚上當開聲禱告時,你相信神會為你成就什麼呢?在你家庭,校園,你的禱告生活如何呢?你有小組與你一同禱告嗎?你的心是否在神面前對付清楚?有否讓主成為交響樂曲的指揮?惟有真正的禱告,才能經驗主的同在,你的人生才得到優美!
\newline
\newline


\newpage

\section{第50屆~港九培靈會~高文~第三講　得主寶血潔淨}
\label{sec:752}
\textbf{第三講　得主寶血潔淨}
\newline
\newline
連結: \href{https://www.hkbibleconference.org/session-message/view/752}{\texttt{https://www.hkbibleconference.org/session-message/view/752}}
\newline
\newline
\hyperref[sec:751]{< < < PREV SERMON < < <}
~
\hyperref[sec:toc]{[返目錄]}
~
\hyperref[sec:754]{> > > NEXT SERMON > > >}
\newline
\newline
希伯來書 9:11-9:12
\newline
\begin{longtable}{cl}
\hline
\hline
章節 & 經文 (和合本修訂版)\\
\hline
9:11 & \begin{tabularx}{0.7\textwidth}{X} 但現在基督已經來到，作了已實現的美事的大祭司，經過那更大更全備的帳幕，不是人手所造，也不是屬於這世界的； \end{tabularx} \\ \\ \relax
9:12 & \begin{tabularx}{0.7\textwidth}{X} 他不用山羊和牛犢的血，而是用自己的血，只一次進入至聖所就獲得了永遠的贖罪。 \end{tabularx} \\ \\
[1ex]
\hline
\hline
\end{longtable}
在非洲肯雅,一日有一少年到醫務所求醫,因為他在森林剪草受了傷.每一步他腳踏之處都留下了血跡.數小時後,那少年的母親在醫務所找到他.人就希奇問她如何到這地方,原因是那位母親從未到過這裡.她就回答說:「這是很容易的,我祇要跟著那些血跡便可以找到了.」同樣,我們今日找著耶穌基督,也是跟著祂的血跡而來!
\newline
\newline
聖經中有四百六十次提及「血」,每次講到「血」就會思想三個非常美好的字:生命,犧牲和立約.
\newline
\newline
(一)生命──神賜下生命的原則.我們屬靈生命的性質也正反映出血.每一個活物都有紅色的血在他體內,血將食物和氧氣放在體內,來支持我們的生命；血也戰勝病菌,又將無用的東西排出體外.沒有血就不能有生命.在人體上平均每一分鐘血會作兩次循環,這樣每個細胞就得潔淨.每一個成年人平均有五至七品脫血,每一立方毫厘血大概有五百五十萬個細胞；這些細胞有一百一十至廿日之生命.為了達到新陳代謝的作用,每秒鐘身體內便會製造新的細胞；實在使我們太驚奇了!在血裡確實見到生命.所以聖經提及生命的血,我們就更易明白了.心臟是整個血所經過循環的地方,因此,心臟也使我們想起生命.(這次培靈會的紀念章似乎作了這個例證.)所以聖經又說:在人心裡出了各樣邪惡的思想,罪人必須心得潔淨；神要將一個新的心賜給那些從心裡相信祂的人.血既然代表生命,生命是從神來的,因此提到「血」就不能輕易處理.
\newline
\newline
早在舊約時代,以色列人已懂得先取出血然後才吃肉,因為在律法上說,血是在生命裡.神將血賜下,是救贖的預表,是神聖的,是神所創造的；血也是指為我們捨去的那一位.為這緣故,當思想血時,我們要存一個尊敬的心.二千年前這血在加略山的十字架上為我們流出,藉著血耶穌與我們合而為一.所以耶穌說:「凡喝我血,吃我肉的就得永生.」耶穌不單只將人生哲學或倫理準則給我們,乃是將祂自己賜給我們.基督徒的生命不單單只是一個信條,乃是加入那位深愛我們,將祂生命賜給我們的主.
\newline
\newline
有人將血儲藏在血庫裡,當有需要血時,便從血庫提取出來.許多時候一個垂死的人,因為有人肯將血輸給他,就得回生命.加略山的血庫豈不是更加美好嗎?因為在那處有無限的供給.現在這血與初期的一樣有功效,能夠幫助每一種類型,每一個的需要；對那些肯存著信心接受的人,這血庫是白白施出的.
\newline
\newline
(二)犧牲──是無罪的代替有罪的犧牲,將生命傾倒出來!因此犧牲的愛都在這血裡面.在舊約的祭物裡,就將這真理表明出來.神藉著人心願獻出的祭物,教導人明白敬拜的意義.生命是白白傾倒出來以至於死,人們所能明白及付出的,以這個為最終極.但是神卻不要低於這些獻祭,當時的敬拜是要帶一隻無瑕疵的動物作祭牲,代替人的罪.祭司要按手在祭牲的頭上,表明牠是被殺的,是代替那有罪的人,然後那人提出獻祭的原因,祭司就用鋒利的刀從祭牲的喉嚨割開,血就傾流出來；祭司用器皿盛上這些血,然後灑在祭壇上及周圍之處.這方法代表我們如何將生命呈獻給神,也代表了神給人的審判就是死亡,同時亦表明人何等渴望與神和好,是人在神的憐憫上,信心的表現.整個獻祭是否有效,視乎那人是否真將自己與血應同.恐怕宗教最大的危險,就是以宗教的儀式,代替真誠的意義；以儀式成為全套的信仰和敬拜.外在的儀式其用意是幫助人明白屬靈的真理.若人忘記這用意,那些宗教儀式可變成拜偶像.在初期的人出毛病就在此,他們經過許多宗教儀式,但只有儀式卻沒有真誠的心.神多次警告他們,叫他們不要誤解儀式的用意,否則會招引神的審判而已.神不會容忍假冒為善,也不會容許褻瀆祂的事情.耶穌基督第一件工作就是潔淨聖殿,將在殿裡作買賣的人趕走,因為他們將獻祭的真正意義溶毀了.聖經提醒我們不要將關乎敬拜神的聖事當作外表,而內心毫無誠意.在這一切裡神所教導我們的,就是愛.靠神的話語所過的生活,正好反映出我們是否有愛.順服是勝於獻祭.在最高意義方面,獻祭是表示愛；表明出心裡渴想發表,但口裡不能表達出來.那位敬拜者真的在神面前說:我是真正的倒空自己,希望靠這獻祭我罪得洗除.也表明他心裡想真正願意盡心,盡性,盡意愛神.同樣寶貴的,就是在祭壇上的血,也表明神愛那敬拜者.聖經說:神不計代價,十分願意與那敬拜者恢復交通,彼此合一.祂是一位掌權的神,是十全十美的神；祂深願盡祂所能挽回失喪的人類.在祭壇上的血就表明神的方法,好叫我們按神的公義和聖潔,得以挽回.
\newline
\newline
當然在舊約的獻祭還未完全,只代表以後那真正的獻祭.新約聖經記載當耶穌被帶到城外,一切牛羊都要被帶出來,這點是非常有意義的.當耶穌在十架上受死時,他們在殿裡也獻上祭物.三小時之久,我們的主懸掛在十架上,祂被人藐視,棄絕,被列在罪犯中,在祂兩旁都有罪犯與祂同釘.祂只說了幾句話,耶穌在那群譏笑祂,使祂受苦的人面前,幾乎不成聲；當祂難於呼吸之時,祂的身體因痛苦而抽筋,祂提高聲音大聲喊叫:「成了!」當時在聖殿中的幔子從上至下裂開,從大理石之牆上掉下來.現在人可以直接看見殿裡的至聖所,那代表舊約儀文律例的幕現在已完了,誰有心就可以直接到神面前,在舊約所代表獻祭的一切現在實現了.神用自己獨生兒子獻上,作無瑕無疵的羔羊,現在人可以存著信心到神面前,見到神的榮耀,神的聖潔.你們看主的大愛,當我們還作罪人的時候,基督就為我們受死,將神的愛向我們顯明.
\newline
\newline
(三)立約──稱為「那應許的約.」這約是忠心至死的標誌,因此在許多國旗會見到紅色,意思在這標誌下生活的人；甚至願意為這標誌所代表的,忠心至死防衛自己的國家,甚至犧牲自己的性命.在原始民族是用真正流血作出這種表示.例如有兩個人想彼此立約,他們會割開自己的手掌,然後彼此緊握,兩隻有創痕的手的血就彼此混合.有人說這是今日握手的起源.將你的忠貞向對方承諾.神說:「看哪!我將你們銘刻在我的掌上.」(賽四十九16)就是這意思.當一個人要在眾人面前表明他赤忠的心,也是用這方法,他將右手前臂割開,使眾人見到他流出的血,直流到他腳前成小血灘,就宣告他所說的是認真又誠實的.當以賽亞說:耶和華藉祂右手發誓時,我相信這種情景在他心意裡.這古舊的習俗在今日法庭裡仍有它的遺痕,無論如何表達,血總是與立約有關係的.
\newline
\newline
當人反叛神之後,在伊甸園裡就有這樣的先聲.神聲明:有一天,女人的後裔要傷蛇的頭.亞當並沒有做任何事情,足以承受神這應許,神只是將這應許白白賜下.然後神確保這約,就在當地立一個祭；神親自宰殺一隻羊,用牠的皮作為始祖的皮衣,用來遮蓋他們.這表明了神的愛.
\newline
\newline
挪亞從方舟出來時立刻就獻祭.聖經說:「當耶和華聞那馨香氣,就心裡說,我不再因人的緣故咒詛地,也不再按著我才行的,滅各樣的活物了.」(創八21)神應許人要生養眾多,遍滿全地.在天空上作見證的是天虹,現在在祭壇上作見證的是血.
\newline
\newline
神應許要藉亞伯拉罕,興起祂要揀選的民族.創世記所記的歷史:亞伯拉罕得神的命令要將自己的兒子獻上,在上山的路上,兒子問父親說:我見到柴,見到火,但在何處找到祭牲呢?父親回答他說:神會供給祭牲.在山上築了壇,亞伯拉罕把以撒捆綁好,放在柴上,伸手拿刀要殺他的兒子,突然有神的聲音在天上發出說:亞伯拉罕,你不可在童子的身上下手；我現在知道你是敬畏神的了,因為你沒有將你的獨生子留下不給我.於是亞伯拉罕看見一隻公羊,兩角扣在稠密的小樹中；那羊是沒有瑕疵的,可作為祭牲；亞伯拉罕就取下代替他兒子作祭物,神再從天上對他說:我要賜福你和你的後裔,多如天上之星,海邊的沙；藉著你世上萬族要得福.這個見證乃是用血所建立的.
\newline
\newline
當以色列民在埃及為奴時,神應許拯救他們.要每家獻上羊羔為祭物,為代替那家裡頭生的而死；又用牛膝草將血塗在門框和門楣上,誰也不可離開自己的房子之外,直到早晨.當夜神的使者見到凡門框和門楣上有血的,就越過那家,不擊殺那家裡頭生的.這門上的血,就證明這家深信永生神的話,在家裡堅心信靠等候神的拯救.
\newline
\newline
(來九11-12)記載:這血就是神與你們立約的憑據.先洒在律法書上,然後應用在萬民身上.眾民所仰望的祭牲,在加略山上耶穌被獻上,這救恩就實現了；神的兒子要這樣傾倒流下自己的寶血.當主與門徒吃逾越節最後晚餐時,他拿起無酵餅,禱告,祝謝後,就擘開,分給門徒說:你們拿著吃,這是我的身體.然後拿起杯來說:這是我與你們所立的新約,是我的血,為多人流出.吃完了晚餐,耶穌又領他們同唱逾越節的詩歌,可能是詩篇一一五篇至一一八篇之精華.你能否想像耶穌基督在那小群的門徒面前,當祂唱完第一節,門徒就在副歌加入......匠人所棄的石頭,已成了房角的頭塊石頭,這是耶和華所作的,在我們眼中看為希奇.這是耶和華所定的日子.我們在其中要高興歡喜.......各位,你能否想像耶穌的眼淚不住流落在祂的臉上,當祂想到用繩索把祭牲拴住,牽到壇角那裡時,仍唱出:你是我的神,我要尊崇你的慈愛,直到永遠.他們唱完詩後,進了園子；不久祂就被掛在十字架上,鮮血從祂身上流出,成為十架腳前的血庫時,這就是神最終所立的約了.這約直到留遺命的人死了才生效,由耶穌舉起雙手,整個受造之物都看見了,天軍都作見證,神在祂自己所說過的事情,現在神都成就了.正如神從開始就說將救恩賜給人類,現在神實現了祂曾說過的.這是天上的見證,新的約現在被血所確立了.因為耶穌基督靠這血進入不是人手所作的聖所,現在主在那聖所,在神面前代表我們站立.這是見證神是十分愛我們,將自己的生命賜給我們,是口舌述說不盡的恩賜,是不能改變的聖言.
\newline
\newline
許多年前,在美國西部,某些律法受到批評,及有不少波折.強盜搶劫那些全無保衛的村落,人們也無法幫助那些鄉民.後來村民一定要捕捉那些強盜,如果找到任何犯罪證據,必將違法者處決.不久,他們在肯薩斯州發現這群強盜,就決定判決死刑,有人在群眾中來作見證,其中一個作見證的人看見在死囚中,有一個是他親密的朋友；他就走到官長面前請求釋放那死囚.但他們解釋這些人都有犯罪證據,要一同處死.但那見證人說出那死囚有家庭,有兒女,他的家人十分需要他；他請求釋放那罪犯,但官長不肯這樣做.最後,在無可辦法之下,他要去代替那朋友受死刑.這實在是太驚人的請求,他們只有容讓他去取代那死囚的位置,與其他死囚一同被處死.最後,那罪犯回到這可怕的刑場,請求官長讓他埋葬他的好友；在墓碑上他這樣刻上:「我最好的朋友,他取代我的位置,他為我死.」這紀念碑至今還在肯薩斯州附近的小村落裡.
\newline
\newline
各位弟兄姊妹,耶穌基督祂豈不是為了我們,取替我們的位置,為我們而死嗎?不是耶穌犯了罪而走上十字架,乃是祂深愛我們這無助軟弱的罪人.因為祂與我們的生命相同,所以祂在神面前能真正支持我們的生命.祂代替我們,將自己的生命傾倒出來；祂為我們受死,這罪債現在已一筆勾銷；祂完全為我們付上,將白白的救恩賜給我們.惟有接受祂的寶血,進到祂裡面,才能接受救恩.祂十分深願我們每一位能過著優美的人生,無論何人願意都可以來,在祂裡面永遠活著.
\newline
\newline


\newpage

\section{第50屆~港九培靈會~高文~第四講　聖靈充滿的生活}
\label{sec:754}
\textbf{第四講　聖靈充滿的生活}
\newline
\newline
連結: \href{https://www.hkbibleconference.org/session-message/view/754}{\texttt{https://www.hkbibleconference.org/session-message/view/754}}
\newline
\newline
\hyperref[sec:752]{< < < PREV SERMON < < <}
~
\hyperref[sec:toc]{[返目錄]}
~
\hyperref[sec:755]{> > > NEXT SERMON > > >}
\newline
\newline
約翰福音 14:12-14:19
\newline
\begin{longtable}{cl}
\hline
\hline
章節 & 經文 (和合本修訂版)\\
\hline
14:12 & \begin{tabularx}{0.7\textwidth}{X} 我實實在在地告訴你們，我所做的工作，信我的人也要做，並且要做得比這些更大，因為我到父那裡去。 \end{tabularx} \\ \\ \relax
14:13 & \begin{tabularx}{0.7\textwidth}{X} 你們奉我的名無論求甚麼，我必成全，為了使父因兒子得榮耀。 \end{tabularx} \\ \\ \relax
14:14 & \begin{tabularx}{0.7\textwidth}{X} 你們若奉我的名向我求甚麼，我必成全。」 \end{tabularx} \\ \\ \relax
14:15 & \begin{tabularx}{0.7\textwidth}{X} 「你們若愛我，就會遵守我的命令。 \end{tabularx} \\ \\ \relax
14:16 & \begin{tabularx}{0.7\textwidth}{X} 我要求父，父就賜給你們另外一位保惠師，使他永遠與你們同在。 \end{tabularx} \\ \\ \relax
14:17 & \begin{tabularx}{0.7\textwidth}{X} 他就是真理的靈，是世人不能接受的。因為他們既看不見他，也不認識他；你們卻認識他，因他常與你們同在，也要在你們裡面。 \end{tabularx} \\ \\ \relax
14:18 & \begin{tabularx}{0.7\textwidth}{X} 我不會撇下你們為孤兒，我必到你們這裡來。 \end{tabularx} \\ \\ \relax
14:19 & \begin{tabularx}{0.7\textwidth}{X} 再過不久，世人不再看見我，你們卻會看見我，因為我活著，你們也要活著。 \end{tabularx} \\ \\
[1ex]
\hline
\hline
\end{longtable}
今天晚上講到耶穌基督升天之後要降下的聖靈,經文內容是耶穌親自對我們應許；祂吩咐門徒要等候,直到聖靈降臨.除非耶穌基督在世上完成了祂的工作,否則聖靈不會降臨,當耶穌基督在世上完成了祂的工作,回到天父的右邊時,聖靈就帶著全然的能力降臨.培靈會的徽章就含有這意思:「心形」表明我們在基督裡所認識和接受的生命,「鴿子」代表聖靈,祂來將耶穌基督救贖大功散播在我們身上.基督教不是一些關於耶穌基督的教條,也不是叫我們盡量學習耶穌的生活方式；乃是基督在我裡面活著.從三方面可以看見這真理:
\newline
\newline
(一)聖靈將耶穌基督的性情輸入我們裡面,在(約十四18-19)中主耶穌解明了這真理.在提到聖靈時,耶穌說:「我不撇下你們為孤兒,我必到你這裡來；還有不多時候,世人不再看見我,你們卻看見我,因我活著,你也要活著.」父神計劃了,耶穌基督執行了,現在聖靈要將這一切顯明出來.換句話說,聖父是行政者,聖子是一位顯明者,聖靈是一位運行者.創世記第一章介紹了聖靈,神藉著祂的靈創造了宇宙；該處說,神的靈運行在水面上.約伯說:靠著神的聖靈,天地就因此建立起來了.同樣神創造人時,神將氣吹進人的鼻孔裡；「吹氣」與聖靈的「靈」字同一字根.意思是神將所創造的人藉著所吹入的氣,使他成為靈,好叫人在神裡可以有動作,有生活.罪惡的悲劇就是人的生活敗壞了.人仍然可能有神的生命,但無法過著屬神的生活.神的靈現在從人身上撤退了,因此人的性情被罪惡侵蝕就敗壞了.這是每一個在世上的人的光景.一個敗壞了的人怎可能被救贖呢?這就需要一位十全十美的,肯為我們被獻上的耶穌,祂肯來到世上才能為我們成全.究竟耶穌如何與人認同呢?當然是靠著神的聖靈,聖靈將生命的種子放進童貞女馬利亞的胎裡,馬利亞受了孕；就將神的兒子生產出來.耶穌基督為女子所生,祂的生命,從開始就是一個獨特的生命.在肉身上祂與我們一樣,但祂沒有任何一樣從罪惡來的敗壞；所以祂是一位十全十美的人.在這種能力下,耶穌向我們呼召,我們也要得著這種改變.那將耶穌成為肉身的能力,同樣臨到我們身上,使我們有屬靈的生命；耶穌稱這過程為「重生」.耶穌說:「從神的聖靈而生.」祂說:叫人活著的乃是靈,肉體是無益的；我對你們所說的話,就是靈,就是生命.(約六63)當然我們還有一個敗壞的肉體,在我們裡面還帶著罪惡的遺痕.但我們裡面的心靈,開始享受,分享耶穌的靈命,在學習像主的事情上不斷有長進.這是靠著那位聖潔的聖靈而作的,將本屬於耶穌基督的性情輸入我們的心靈裡,這就是被聖靈所結的果子.好像:仁愛,喜樂,和平等等；這一切都是耶穌基督生命的特徵,無人能靠自己過著像耶穌基督獨特的生命,神也不會祈望我們如此；惟有靠住在我們心裡的聖靈的幫助才能過屬主的生活.當我們在主的形像上不斷長進,以致榮上加榮.有一天這古舊的身體都會改變.當離世回天家時,聖靈會同樣改變我們的身體.
\newline
\newline
(二)聖靈賜能力給我們作神的工.(約十四12)主說:「我實實在在的告訴你們,我所作的事,信我的人也要作；並且要作比這更大的事,因我往父那裡去.」
\newline
\newline
當我們親近神的時候,祂的聖靈就會降臨在我們身上.當然聖靈早已在世上工作,對某些人在某些時間裡,聖靈已工作；例如:當約瑟在埃及時,神的聖靈就與他同在.又神的聖靈將智慧和指引賜給摩西,因此他才能作成神所交付他的工；當工作增加時,神教導他找七十位長老,聖靈又將智慧賜給那七十人.當他們建造會幕時,神又賜下聖靈使他們當中有人,有才幹去作成這工.(出三十一3)記載比撒列建造會幕時,被神的靈充滿,因為那會幕是預表耶穌基督在二千年後要完成的工作.(出二十八3)提到婦女為祭司做袍,她們同樣要被聖靈充滿,因為祭司的袍是預表耶穌基督大祭司的責任；袍是潔白的,代表主的公義.聖經記載只有一位祭司沒有穿上祭司袍的,那就是在加略山上的主.祂赤裸著上身掛在十架上,這本是褻瀆神的；但因祂是獨特無瑕疵的,是照麥基洗德的等次為我們獻上,這獻上也是靠神的靈而成的.在士師時代也是這樣,神的聖靈充滿基甸,和參孫的身上.聖經裡最大一個悲劇,就是有關參孫的一句話:神的靈離開了他,他也不知道；還希望再作神的工；這實在是令人最感到痛苦的悲劇.先知靠著聖靈,才能說同一的話；他們帶著能力才可以寫出神的話,為這緣故,聖經就與世上其他的書籍截然不同了.因為這書是神的靈所默示的,聖靈確保寫書的人不會將錯誤加上；因為他們是按聖靈的感動而寫,所以我們相信並接受這書是神的話.耶穌說:「你們查考聖經,因你們以為內中有永生,給我作見證的,就是這經.」(約五39)這是認識耶穌的最好方法,因為祂是聖靈最終又最完整的見證.當主在世上工作時,聖靈帶領著祂.有一次,主在自己家鄉講道,祂翻開舊約以賽亞書說:「耶和華的靈在我身上,祂膏我.......」(路四18-19)主向會堂的人宣告說:今天這經應驗在你們耳中了.主在世上的工作是靠聖靈而作的,祂的講道,教訓,醫病和趕鬼,藉這些工作神呼召人信靠救主,當人在聖靈的感動推迫下還是不信,就是褻瀆神.因為惟有靠著主,人才能到父那裡.耶穌說:聖靈來了是要叫人知罪,人靠自己是不能得救,一切都是靠恩典.因這世界的王已經受了審判.世人將神的兒子釘在十架上,以為從此可將祂除去,但聖靈使耶穌基督從死裡復活過來.聖靈今晚叫我們知道,耶穌基督是獨一戰勝死亡的門,他就是道路,真理,生命.我雖然講道,但惟有聖靈使你能夠順服真理；我能夠指導你歸向耶穌基督,但惟有聖靈使你真正知道自己的需要.感謝神,每當我們高舉耶穌時,就能吸引萬人來歸向祂；這不是人的工作,乃是聖靈使人知罪,吸引人歸向主.
\newline
\newline
(三)聖靈將神的同在顯明出來每一件事神所造的都在耶穌身上顯明出來.祂是道,是神自我的顯露,是我們經常可以看見的.「從來沒有人看見過神,只有在父懷裡的獨生子,將祂表明出來.」(約一18).約十四16提及保惠師,就是安慰者之意,祂站在你旁邊,加你力量,另外一個翻譯是稱為辯護者或律師.當你在法庭上,祂是一位站在你旁邊,明白你的處境,能同情你,有影響力的,祂是能幫助你的一位.三年來耶穌在門徒中間就好像門徒的律師一樣,一同奔走人生的道路.現在耶穌要回到天父那裡,在肉身上不再與他們同在；或許你想知道門徒為何憂愁起來,是因為他們已經完全靠賴了耶穌.所以耶穌對他們說:「你們心裡不要憂愁,你們信神,也當信我.」(約十四1)我不是這樣就離開你們,我明白你們的處境；我回到父那裡,我會求父另外賜下保惠師.耶穌豈不是關心我們嗎?我們是永不會孤單的.「另外」就使整件事成為美好甘甜.原文「另外」有兩個意思:
\newline
\newline
第一,兩者截然不同.沒有共通地方；例如:人與檯.
\newline
\newline
第二,種類相同,但不同個體；例如:講員與翻譯員,同是人,講道一樣,動作一樣,思想也是一樣,甚至可以成為同胞兄弟.
\newline
\newline
這樣,當耶穌說我另外賜下保惠師,是指兩個不同個體,但在質素上完全相同,神是靈同是耶穌基督的靈.是一位有位格的,祂與耶穌相似,祂來要引領我們進入真理；要把將來的事向我們表明,也要教導我們如何禱告,祂要榮耀主耶穌.世人不認識主就不能接受這個真理.這不是理論,乃是實實在在的,今日聖靈的同在正如昔日主與門徒同在一樣.事實上,聖靈永久的同在,比當時主與門徒更親密；在聖靈裡,肉體的限制被除去,祂永不會丟棄我們.這應許要臨到每一位願意接受聖靈的人.所以當時主對門徒說:要等候聖靈的降臨,惟有得到從上而來的能力才有力量出去作工,廣傳基督.因此在不久,五旬節的時候,聖靈就降臨.有風和有火焰如同舌頭一樣,這一切叫我們想到神偉大的洗禮.這事以後,就不再需要重演.因五旬節持續的果效,就是主充滿在信徒心中的方法.聖靈不但在五旬節時充滿在門徒心中,在使徒行傳其他時間,聖靈同時充滿他們,使他們有膽量和喜樂.叫世人都認出他們是跟隨過主耶穌的,他們的生活方式,正可以表明耶穌是活在他們生命裡面.
\newline
\newline
各位今天晚上,你是否經驗過活在你生命中?你有沒有這應許?只要你完全將自己交給神,一切罪惡向神承認,憑信心接受主.惟有一個倒空了又向聖靈張開口的,聖靈才充滿他.一個茶杯不能將海水充滿,但將杯倒轉放在水平下,它就能把海水盛滿.
\newline
\newline
在某年的冬天,一日我與紐約一位會督同在街上走；忽然一個小孩把一個雪球拋落在會督的帽上.小孩心裡十分害怕,他以為這人必會帶他到警察局.但當會督見到他,他身上穿的衣服鞋襪都十分破舊；會督心裡充滿憐憫就帶他到對面的商店,他要送給小孩一些生果,雖然生果十分昂貴,但他盡量將生果放進小孩的衣袋裡,並對小孩說:回家告訴你母親,神愛你們.當時會督心裡充滿喜樂.說:如果我們將地上屬世的東西送與別人,神豈不是更將天上屬靈的厚恩賜給我們嗎?
\newline
\newline
各位,在你心裡是否有些東西攔阻聖靈將主的恩賜下,你是否願意將它除去,讓聖靈充滿你.只要你的心繼續向神敞開,神必向你施恩,賜下應許.
\newline
\newline


\newpage

\section{第50屆~港九培靈會~高文~第五講　事奉永活的神}
\label{sec:755}
\textbf{第五講　事奉永活的神}
\newline
\newline
連結: \href{https://www.hkbibleconference.org/session-message/view/755}{\texttt{https://www.hkbibleconference.org/session-message/view/755}}
\newline
\newline
\hyperref[sec:754]{< < < PREV SERMON < < <}
~
\hyperref[sec:toc]{[返目錄]}
~
\hyperref[sec:757]{> > > NEXT SERMON > > >}
\newline
\newline
約翰福音 17:17-17:26
\newline
\begin{longtable}{cl}
\hline
\hline
章節 & 經文 (和合本修訂版)\\
\hline
17:17 & \begin{tabularx}{0.7\textwidth}{X} 求你用真理使他們成聖；你的道就是真理。 \end{tabularx} \\ \\ \relax
17:18 & \begin{tabularx}{0.7\textwidth}{X} 你怎樣差我到世上，我也照樣差他們到世上。 \end{tabularx} \\ \\ \relax
17:19 & \begin{tabularx}{0.7\textwidth}{X} 我為他們的緣故使自己分別為聖，為要使他們也因真理成聖。 \end{tabularx} \\ \\ \relax
17:20 & \begin{tabularx}{0.7\textwidth}{X} 「我不但為這些人祈求，也為那些藉著他們的話信我的人祈求， \end{tabularx} \\ \\ \relax
17:21 & \begin{tabularx}{0.7\textwidth}{X} 使他們都合而為一。正如父你在我裡面，我在你裡面，使他們也在我們裡面，好讓世人信是你差我來的。 \end{tabularx} \\ \\ \relax
17:22 & \begin{tabularx}{0.7\textwidth}{X} 你所賜給我的榮耀，我已賜給他們，使他們合而為一，像我們合而為一。 \end{tabularx} \\ \\ \relax
17:23 & \begin{tabularx}{0.7\textwidth}{X} 我在他們裡面，你在我裡面，使他們完完全全合而為一，讓世人知道是你差我來的，也知道你愛他們，如同愛我一樣。 \end{tabularx} \\ \\ \relax
17:24 & \begin{tabularx}{0.7\textwidth}{X} 父啊，我在哪裡，願你所賜給我的人也同我在哪裡，使他們看見你所賜給我的榮耀，因為創世以前，你已經愛我了。 \end{tabularx} \\ \\ \relax
17:25 & \begin{tabularx}{0.7\textwidth}{X} 公義的父啊，世人未曾認識你，我卻認識你，這些人也知道是你差我來的。 \end{tabularx} \\ \\ \relax
17:26 & \begin{tabularx}{0.7\textwidth}{X} 我已讓他們認識你的名，還要讓他們認識，好讓你愛我的愛在他們裡面，我也在他們裡面。」 \end{tabularx} \\ \\
[1ex]
\hline
\hline
\end{longtable}
(今天晚上看見許多年青人,因為人多的緣故,委屈了你們,有人要坐在地板上.我知道各位從港九各教會而來,盼望你們明晚會帶另一位未信主的親友來參加,叫他們也得著這寶貴的福份.)
\newline
\newline
剛才讀出的經文,是主耶穌為年青的門徒的祈禱,盼望他們的前途,能過著優美人生的生活.請各位注意耶穌禱告的五件事項:
\newline
\newline
(一)主祈禱求神賜這班年青人對失喪靈魂有異象.18節:「你怎樣差我到世上,我也照樣差他們到世上.」21節:祂祈求世人能相信.23節:相信這位天父差來的耶穌.作為耶穌的門徒必須明白我們在世上的責任,有人在肉身上有痛苦,我們必須關懷；希望他們得著公平,良善的待遇.最重要是叫人認識主,和主所賜下的生命.當耶穌看見耶路撒冷,就哭了；因他們不明白祂來世的目的,他們好像失喪的羊,無牧羊人一般.今日我們想到有三十億人口,還未認識耶穌,你的心是否傷痛?每廿四小時就有十四萬以上的人,沒有盼望的了結他們的生命.如果將失喪的人排成隊伍,可以圍繞地球卅個圈.基督耶穌降生為要拯救這些人.
\newline
\newline
在美國有海防軍,負責海岸巡邏,搶救被溺者.一次暴風侵襲,要出海找一隻沉船；隊長發號施令要各人準備好搶救遇難的人.其中一個隊員害怕說:隊長,恐怕我們出去,一去不復回!隊長大聲說:我們不一定要回來,但一定要出去,海岸巡邏隊目的不是救自己,乃是出去搶救在苦難中的人.讓我們負起搶救靈魂的責任:在我們心靈裡必須對失喪的人有異象!
\newline
\newline
(二)主祈禱求神使每一個基督徒,都有主所託負的使命,被差遣到世界去.20節:「我不但為這些人祈求,也為那些因他們的話信我的人祈求.」(約十四12)耶穌說:「我所作的事,信我的人也要作.」這是神的呼召,各教會裡的信徒,都有不同的恩賜,不同的職位；但每一肢體必須有事奉.今日你的教會有多少位牧師?可能是一位,兩位或三位,他們是全時間事奉的.但正確的答案應是每人都是牧師.禮拜日大家聚集在一起服事神,禮拜一開始,各人就在不同的崗位服事了.牧師在教會輔導人,商人在商店,學生在學校,主婦在家庭,每一位信徒都是在服事神.牧師最主要的工作是裝備信徒,講解神的話,為弟兄姊妹禱告,與弟兄姊妹同工.今日何等需要全時間受訓練的工人,或作女傳道,或作宣教士,或作護士.......若你將自己生命投資在服事主的事上,是何等美的事.
\newline
\newline
(三)主祈禱求神使門徒在服事的工作上,能夠效法主的榜樣.主的生命是一個明證,讓我們去效法.主週遊四方,繼續行善事,幫助有病的,貧窮的,受壓逼的.......一有機會主就為福音作見證,有時向成千上萬的人講道,但主卻用大部份時間,在訓練他揀選的十二個門徒上.他們十分渴想學習神的真理.他們不是猶太的宗教領袖,只是平信徒而已.除了那賣主的猶大之外,其他的門徒都注目在主身上.耶穌三年多工作裡,大部份與門徒在一起,無論食,住,行,在會堂教訓人,讀經祈禱,幫助人等等,一直以來,他們看見主如何作工,主也讓他們有不同學習的機會,使他們找到自己的恩賜,傳福音服事他人.(太二十八19-20)的大使命意義深長,其中最重要的一詞:「作門徒」,其他的都是文法上的分詞,全歸附在動詞之下,所以大使命最重要是使人作主的門徒.主要求他們作平常主所作的工作,這是主的命令,人人都要作.有些人有向群眾佈道的恩賜,但使人作門徒卻是單對單,一領一的事.每人都有機會使人作主的門徒,學生與主婦,宣教士與牧師,同樣有機會領人歸主,這實在是使人興奮的工作!在我們的崗位上,需要找人作門徒,這必須要有忍耐,為他們禱告,栽培他們直到這人也能使別人作門徒.
\newline
\newline
(四)主祈禱求神使我們願意奉獻自己.19節:成聖是分別出來之意.主自己分別為聖,專心執行父要祂所行的使命,直到上十字架.主也祈求祂的門徒同樣分別為聖,使他們能經驗主,而投身在主的使命上；不再有保留,是絕對的順服,以至於死.在十架旁的人都嘲笑祂說:祂能救別人,卻不能救自己!但這正說明出真理來:主來並不是要救自己,乃是為救你和我!所以祂犧牲自己,祂來要服事人,代價極高；我們都要如此效法主.當耶穌預言要上耶路撒冷受死時,彼得卻阻止祂；但主責備彼得說:撒旦退我後邊去罷,你是絆我腳的；因為你不體貼神的意思,只體貼人的意思.這表明彼得還未明白主上十字架的目的.彼得會這樣想是因為從未想到十字架,他只希望坐在寶座右邊,成為重要人物,他想得到快樂,過奢侈的生活；他不想有苦難,更不想他的主在人手下受死.今日在教會裡許多人都像彼得,從十字架退後.主責備彼得被撒旦迷惑,然後主轉身對門徒說,若有人要跟從我,就當捨己,背起他的十字架來跟從我.因為凡要救自己生命的,必喪掉生命；凡為我喪掉生命的,必得著生命.(太十六21-25)這是得著優美人生的秘訣.一粒麥子若不是掉落在地上死了,就不能結出子粒來.除非你我為主捨己,否則不能有什麼成就.若你們關心的是世界的享受,你就不能經歷豐盛的生命.
\newline
\newline
(五)主祈禱求神使我們被神的愛感動,樂意的奉獻,更使我們能存著喜樂的心接受十字架.愛是神的本性,十字架只不過將神的愛刻畫出來.(約十七23-24):主祈求這愛在門徒身上,使他們知道天父如何愛,他們也要如何的相愛.是愛激勵我們進入到世界,正如主一樣.當我們真正愛主,就要跟從祂直走完這路程.
\newline
\newline
我曾在一間大學執教,見許多青年喜歡三五成群的；但忽然我發現某一男生不再出現在他們的小群中,我覺得很奇怪.他是否病了?到底為什麼?一天晚上,我終於找到他,在情人徑上,他與她一同行走,從他倆目投視,我就想到他們已墮入情網了.當你墮入情網時,總希望常見到對方,討他的歡喜；好像對世界上其他東西都失去興趣,因為對方實在太美好了,我們對主也該如此.盼望多與主親近,做討祂喜悅的事,為了愛主,願意捨棄一切,不計較任何代價.
\newline
\newline
我聽聞有一青年男子要輸血給他的姊妹,由於輸血者必須曾有過病人的病源,因此只有這青年才是最合適的.經醫生一番解釋後,最後他流著淚,嘴唇戰抖著說:醫生,我願意.當輸血過程將近完成時,那青年又問醫生說:醫生,我何時會死,這時醫生才明白,為什麼他遲延才決定輸血給他的姊妹,又為什麼他如此發抖流淚.因為那在他所能理解之中,他以為輸血給他姊妹,就是連生命也給她.在那一剎那,他作了偉大的決定.我想我們站在十字架前,也作了決定,想將生命交付給主,願意盡心,盡性,盡意,盡力愛祂.神不是逼迫,鞭打我們,乃是用祂的愛激勵我們.所以基督徒的生命能夠如此美好,甘願犧牲,就算有試煉和困難,也能歌唱.因為主愛你,你也愛主.主吩咐什麼,你也願意去作.各位青年人,今天晚上你是否甘願奉獻自己,完全獻上為主而活?
\newline
\newline


\newpage

\section{第50屆~港九培靈會~高文~第六講　得著新生命}
\label{sec:757}
\textbf{第六講　得著新生命}
\newline
\newline
連結: \href{https://www.hkbibleconference.org/session-message/view/757}{\texttt{https://www.hkbibleconference.org/session-message/view/757}}
\newline
\newline
\hyperref[sec:755]{< < < PREV SERMON < < <}
~
\hyperref[sec:toc]{[返目錄]}
~
\hyperref[sec:758]{> > > NEXT SERMON > > >}
\newline
\newline
路加福音 19:1-19:10
\newline
\begin{longtable}{cl}
\hline
\hline
章節 & 經文 (和合本修訂版)\\
\hline
19:1 & \begin{tabularx}{0.7\textwidth}{X} 耶穌進了耶利哥，要從那裡經過。 \end{tabularx} \\ \\ \relax
19:2 & \begin{tabularx}{0.7\textwidth}{X} 有一個人名叫撒該，作稅吏長，是個財主。 \end{tabularx} \\ \\ \relax
19:3 & \begin{tabularx}{0.7\textwidth}{X} 他要看看耶穌是怎樣的人，只因人多，他的身材又矮，所以看不見。 \end{tabularx} \\ \\ \relax
19:4 & \begin{tabularx}{0.7\textwidth}{X} 於是他跑到前頭，爬上桑樹，要看耶穌，因為耶穌要從那裡經過。 \end{tabularx} \\ \\ \relax
19:5 & \begin{tabularx}{0.7\textwidth}{X} 耶穌到了那裡，抬頭一看，對他說：「撒該，快下來！今天我必須住在你家裡。」 \end{tabularx} \\ \\ \relax
19:6 & \begin{tabularx}{0.7\textwidth}{X} 他就急忙下來，歡歡喜喜地接待耶穌。 \end{tabularx} \\ \\ \relax
19:7 & \begin{tabularx}{0.7\textwidth}{X} 眾人看見，都私下議論說：「他竟然到罪人家裡去住宿。」 \end{tabularx} \\ \\ \relax
19:8 & \begin{tabularx}{0.7\textwidth}{X} 撒該站著對主說：「主啊，我把所有的一半給窮人；我若勒索了誰，就還他四倍。」 \end{tabularx} \\ \\ \relax
19:9 & \begin{tabularx}{0.7\textwidth}{X} 耶穌對他說：「今天救恩到了這家，因為他也是亞伯拉罕的子孫。 \end{tabularx} \\ \\ \relax
19:10 & \begin{tabularx}{0.7\textwidth}{X} 人子來是要尋找和拯救失喪的人。」 \end{tabularx} \\ \\
[1ex]
\hline
\hline
\end{longtable}
這段經文是耶穌多次呼召人跟從祂之一祂呼召人用各種不同的方法;但祂降臨在人心裡之真確性卻常是一樣在座各位朋友,或許你還未經驗到耶穌就在這裡,但願你今天晚上就要經驗到祂.雖然祂出現的地方和時候不同,但仍是真實無改變的.
\newline
\newline
撒該是個稅吏,常坐在稅關收稅;有些時候他會把一部份稅銀放進自己的錢袋裡,他擁有世上一切的享受,華麗的房子,一切金錢能買到的東西!但他心裡沒有平安當然他可以在忙碌的工作中,把緊張孤單及需要的心情暫時麻木的放下不理.有時他或許會想到究竟還有沒有人關心他的心靈.
\newline
\newline
這一天,他聽見耶穌要來.他也曾聽過別人講及耶穌,這人能夠使瞎子看見,生病的得醫治,甚至死人也得復生.但是他是做生意的,哪裡有空閒去聽耶穌的傳聞呢!不過,令他最嚮往的,就是耶穌竟然成為罪人和稅吏的朋友；甚至在他的門徒中,有一位也是稅吏,名叫馬太.有人告訴撒該耶穌可以憑著祂的權柄赦免人的罪,對那些願意跟從祂的人,耶穌應許賜他永遠的生命.撒該至少不是一個對宗教熱心的人,但至少他懂得想,我們就應該稱讚他.因為撒該心裡想:若不是耶穌是欺騙人的,就是耶穌的神經不正常,要不然祂必定是神,必定是神的兒子!於是撒該對耶穌產生了好奇,就想更多認識耶穌.
\newline
\newline
他在辦公地方聽見有人大聲喊叫:耶穌來了!拿撒勒人耶穌來了!內裡驅使的心越來越重,現在是撒該的機會,可以親眼看見耶穌.於是決定暫時離開稅關口,走到街上去,在人群中要看耶穌.但當他們見到他,這個矮子,是罪人中的惡人,是個賣國賊；他們就更緊緊的彼此築成人牆,把他推到外面,不讓他進入人群中.他只有更覺得孤單,無人關心.你有沒有嘗過這經驗,人是很多,但自己卻有孤單之感呢!撒該仍不灰心,當他見到那邊的桑樹,很快他就爬上樹上,他看見耶穌來了.第一次他看見耶穌,耶穌走在門徒前面.大概撒該這回覺得很出奇,因為在人看來耶穌並沒有美貌可言!聖經說:耶穌基督常經憂患,無佳形美容(賽五十三2).但撒該見到耶穌流露的面容,發現在世間從未見過的.似乎很神秘,在耶穌身上有一神奇的力量,祂對四周圍的人都表現出極大的愛心.現在撒該心中何等渴想認識耶穌,更多的注視著祂.
\newline
\newline
突然,耶穌停住腳步；抬頭一看,看見撒該；就對他說；撒該,快些下來,今晚我要到你家裡住.你能否想像撒該當時的感受是如何呢!設身處地,你望見耶穌的目光,你有何反應呢?在人群中耶穌向你講話.撒該當然會想:我想耶穌都認識我!這是何等優美的事,耶穌叫你的名字.多少時候,我也會輕忽這回事,但忽然間我明白了,因為耶穌明白我一切,別人不知我的事,但祂知道.祂知道我常到教會裡,但我心靈裡還未滿足,仍然渴望常認識神；祂也知道我心裡的罪,只有靠神的恩典,我能成為新造的人,祂已為我計劃一個美好的生命.如果我知道耶穌知透一切是何等快慰的事!朋友,你心靈還未滿足的,還未有平安,你有渴望；多少時候你在想,究竟人生有何意義?有人關心我嗎?耶穌不單知道你的思想而且能改變你,一切的罪祂能赦免,可以使你在神面前,你一切人生的夢想得以成就,賜你優美的人生!
\newline
\newline
撒該想:耶穌不單叫我的名字,而且還要住在我家裡.他知道耶穌要他愛祂,是何等奇妙的事阿!不錯,世上沒有人愛你,連朋友都會令你失望,但耶穌愛你；為這緣故祂來這世界.小時候,我有一隻小狗.一次我叫不到牠,我就吹口哨,是牠所認識的聲音,牠就走回來.神也是如此,神知道我,認識我；無論我在何處,祂叫我的名字,我就回到祂身旁.
\newline
\newline
保羅對主的愛無法解釋,實在不能盡言語表達；他說:「基督的愛是何等長闊高深,......這愛是過於人所能測度的.」(弗三18-19)
\newline
\newline
某講員曾作以下的例子:一小男孩認為家裡的人沒有關心他和愛他.他就計算每日為母親工作的代價,掃地廿五仙,買東西五十仙等等,大概該日的工資是一元五角.吃晚餐時他就把工資單放在母親的碟裡.次日早餐,他果然從自己的碟中發現母親的回條:我帶你到這世上來生產之苦,不計代價；多年來辛勞教養你,不計代價；未出世就為你祈禱,將你奉獻給神,直到如今,不計代價；一切都不計算代價,若要計算起來,就等如「我愛我的兒子」!各位,這豈不是作母親對兒女的大愛嗎?這例證也不能盡量表達神的大愛,這愛不是父母,夫婦,兄弟姊妹之愛；乃是全能神的大愛.祂是一位靠自己能力創造穹蒼的神,是整個受造之物要向祂敬拜的神,是萬王之王,萬主之主；除了神的憐憫,我們的氣息也不是本能所得的；祂是一位無限量的神,祂愛你,也愛我.我們曾作何事,配得上祂如此愛我們.我們卻是離棄祂,走自己的道路,為自己而活,活在罪惡中.但祂仍然愛我們,祂的愛沒有離開我們.
\newline
\newline
假若你與母親在家裡,忽然有人叩門,站在門外是一位士兵,他手拿死囚罪狀,是控告你曾犯過的罪；你以為沒有人知道這事的,但現在被人揭發了；就在你沒有任何反應之際,母親走到士兵前說:「帶我回去罷,我承認這一切的罪,我接受這判刑.」這時候你還會假裝不理會嗎?你會在法庭上,忍心讓法官定你母親有罪嗎?你是否還敢說:「我要隨我意思行,我是我命運中的主人.」耶穌也是如此代替你,我的罪,祂的愛甚至遠超過這位母親的愛.聖經說:當我們還作罪人的時候,耶穌基督就為我們死.耶穌在十架上背負你,我的罪,你,我憂傷難過的事.當你見到耶穌在十架上,你不能再說:無人關心我!因為神關心你,愛你；你不能再推卸責任了!
\newline
\newline
現在撒該要作決定了.撒該歡喜因為他並沒有遲延,他趕快下來,在街上與神和好.可能四週圍的人會笑他,不信的人可能譏笑撒該發宗教狂.各位!你以為撒該會理會他們嗎?絕對不!因為這是他在人生中第一次經驗真正的快樂,心裡得釋放,心靈充滿喜樂；經驗主的手引領他,他知道回家後救恩就臨到他一家.
\newline
\newline
耶穌來了,祂要作你個人的救主.祂是否已到你的心?你是否經驗到祂?你的罪得赦免沒有?你是否知道你已向著天國?你可願意接受祂?現在就決定罷!
\newline
\newline


\newpage

\section{第50屆~港九培靈會~高文~第七講　活在神全美的旨意裡}
\label{sec:758}
\textbf{第七講　活在神全美的旨意裡}
\newline
\newline
連結: \href{https://www.hkbibleconference.org/session-message/view/758}{\texttt{https://www.hkbibleconference.org/session-message/view/758}}
\newline
\newline
\hyperref[sec:757]{< < < PREV SERMON < < <}
~
\hyperref[sec:toc]{[返目錄]}
~
\hyperref[sec:759]{> > > NEXT SERMON > > >}
\newline
\newline
羅馬書 8:14-8:18
\newline
\begin{longtable}{cl}
\hline
\hline
章節 & 經文 (和合本修訂版)\\
\hline
8:14 & \begin{tabularx}{0.7\textwidth}{X} 因為凡被神的靈引導的都是神的兒子。 \end{tabularx} \\ \\ \relax
8:15 & \begin{tabularx}{0.7\textwidth}{X} 你們所領受的不是奴僕的靈，仍舊害怕；所領受的是兒子名分的靈，因此我們呼叫：「阿爸，父！」 \end{tabularx} \\ \\ \relax
8:16 & \begin{tabularx}{0.7\textwidth}{X} 聖靈自己與我們的靈一同見證我們是神的兒女。 \end{tabularx} \\ \\ \relax
8:17 & \begin{tabularx}{0.7\textwidth}{X} 若是兒女，就是後嗣，是神的後嗣，和基督同作後嗣。如果我們和他一同受苦，是要我們和他一同得榮耀。 \end{tabularx} \\ \\ \relax
8:18 & \begin{tabularx}{0.7\textwidth}{X} 我認為，現在的苦楚，若比起將來要顯示給我們的榮耀，是不足介意的。 \end{tabularx} \\ \\
[1ex]
\hline
\hline
\end{longtable}
在我故鄉德薩斯州,曾有這樣的故事;某日下午,有兩個人在草原上行走,忽然走來一隻野牛,他們嚇得魂不附體;一個快爬上小樹,另一個立刻鑽進附近的地洞裡;野牛見追不到他們就走開去了忽然在洞裡的夥伴猛然跳出來,野牛見狀就走回追逐他,嚇得他連忙鑽回地洞裡​​;於是野牛又走開了.突然在樹上的掉了下來,野牛看見了,又來追逐他,他也害怕得急忙爬上小樹上.如此情況出現多次,不久,小樹上的夥伴就對地洞裡的說:為什麼你不在洞裡逗留多一會兒在洞裡的夥伴回答說:?!誰也不想呀但洞裡有一隻野蚊,好不易受啊如果你在這情形下,會不會進退維谷呢!感謝神,因為人生的一切決定,不是像那故事的複雜,世上許多人卻確實生活在混亂當中主曾說:.許多人都如羊失迷,不知道人生的方向但基督徒卻不是這樣(羅八1 4-18)就是關於生活方向的信息請留意14節,其中四點真理:
\newline
\newline
(一)神為祂的兒女預備了計劃,而且是十全十美的,可從耶穌基督身上看得出來主是道路,是教會的元首,又是首生的;要叫一切信祂的,得以長大成人,滿有基督長成的身量(弗一4-5)「神從創立世界以前,在基督裡揀選了我們,使我們在祂面前成為聖潔,無有瑕疵;又因愛我們,就按著自己旨意所喜悅的,預定我們,藉著耶穌基督得兒子的名分」11節說:「我們也在祂裡面得了基業,這原是那位隨己意行作萬事的,照著祂旨意所預定的.」在創世以前,神早已計劃好,這豐盛的生命是神所預定要賜給祂的子民.神沒有預定誰人失喪,乃願意萬人得救,人人得著真理而悔改.當人要拒絕真理,就是失喪的由來,這是何等大的悲劇.留心(羅八29)「祂預先所知道的人,就預先定下傚法祂兒子的模樣......」神為我們所定的計劃,有主作我們的榜樣,祂的榮耀成為我們生命的終結.神深願我們生活的每一部份都符合神那十全十美的計劃,包括我們的家庭,職業,友誼,休息,活動等等.地上的百合花,天上的飛鳥,神都看顧;!何況我們是被祂兒子寶血所買贖回來的呢祂豈不更關心我們麼!
\newline
\newline
(二)神的聖靈不斷引領我們進入真理,將屬基督的事向我們顯明瞭.其中的方法是用聖經.聖靈默示寫聖經的人,為這緣故,我們要十分確信聖經的記載全都是真實的.當我們用正確的方法查考聖經,就能認識基督,這就是永生.主說:「這聖經是為我而寫的.」既是如此,我們要恆心研讀,以神的話作我們生活的準則.另一方法,聖靈幫助我們是用祈禱.26-27節:聖靈親自用說不出來的歎息,替我們禱告.因為聖靈照神的旨意替聖徒祈求.你有沒有每日忠心於禱告生活,在神面前安靜靈修,等候祂呢?若你將這特權放棄不理,難怪你的生活混亂了.
\newline
\newline
(三)神的兒女必須跟隨神的道路.耶穌基督要求我們跟隨祂,放棄舊的路,走上新的道路.信心不是頭腦的知識和信條,乃是要我們必須將生命信託在神的話語上.這也是愛心的表示,耶穌說:有了我命令,又遵守的人,就是愛我的.信心還需要有長進,按照聖靈的帶領,就不斷跟從,向主學習.每日有新挑戰,是何等美好.今日我們需要順服,按我們所知的順服神；也是跟隨主應有的存心.
\newline
\newline
在南美厄瓜多爾,有五位宣教士被殺為主殉道.當記者到現場訪問時,很驚訝的問道；神為什麼讓他們被殺?其中一位宣教士的妻子回答說:神可能將我丈夫從不順服中拯救出來.「不順服」正是基督徒所應該懼怕的,死亡還是其次,因基督已從死裡復活,除去死的毒咒.美國某雜誌認為他們的殉道是愚蠢的行為,是生命的浪費.這些人如同世人一樣,忘記其中的關鍵；保存生命是其次,最重要是持守對神順服的心志.這些殉道者生前在大學唸書時,在日記上寫著:一個人將自己不可能持守的東西放棄,而換取不可失去的東西,並不是愚蠢的行為.在他出發前曾對他妻子說:若是不能回來,我甘願為歐家族人丟棄自己的生命.正因為他深知主要他將福音傳給印第安人,他要順從主在他身上發的號令.因此保羅寫信給羅馬教會說:我們沒有一人為自己活,也沒有一人為自己死.我們若活著,是為主而活；若死了,是為主而死；所以我們或活或死,總是主的人.(羅十四7-8)正因這種順服,纔能使教會成為榮耀,凱旋的教會.
\newline
\newline
神不單在我們身上有十全十美的旨意,不單有聖靈引領我們遵行神的旨意,神的兒女也必須順,服跟隨祂的道路.
\newline
\newline
(四)當我們順服之後,就有完全的平安.在先父離世前對我說:「兒阿,當一個人與神同行,死後去哪裡呢?」他見我無言回答,就微笑著說:「去哪裡不要緊,最重要的是與神同在.」這叫我想起亞伯拉罕出吾珥時,他不知道往哪裡去,但他甘願作世上的客旅,作尋找者；因為神與他同行,他仰望那有根基的城.今日我們何等需要在人生旅途上與主同行,這樣才得完全的平安.
\newline
\newline
(羅八15)「你們所受的不是奴僕的心,仍舊害怕；所受的乃是兒子的心,因此我們呼叫阿爸父.」「聖靈也與我們的心同作見證.」(羅八16)這是另外得到的滿足,有聖靈甜美的作見證:我們是神的兒女.17節:「既是兒女,便是後嗣,就是神的後嗣,和基督同作後嗣；如果我們和祂一同受苦,也必和祂一同得榮耀.」「作後嗣」是在法律上所用的詞句,意義是雖被收養,但在家裡同樣得到承繼權.所以凡屬耶穌的,都可與耶穌一同分享.不錯,我們可在天上與主同在,同享榮耀；但在地上仍有苦難的.有人因信主而受逼迫,苦難,我們不明白,為何在世上多有苦難?但請留心28節:神為此叫萬事互相效力,目的是叫愛神的人得著益處.
\newline
\newline
數年前,我在葛培理佈道團工作,正在酒店房間內等候外出,忽然電話鈴響了,我從對方得悉佈道團中負責總務的弟兄突然去世,這真使人驚愕；對方繼續告訴我有關詳情,同時希望我作傳棒者,將這消息轉達給葛牧師.事發是這樣的:兩天前,弟兄與一群青年人在湖上划艇；忽然一少女不慎掉入湖裡,弟兄見狀立刻跳入水裡將她救起；當船上的人把少女抱上來後,很快發覺不見了那弟兄.直到翌日潛水人才在湖邊找到他的浮屍.他的妻子兩天來,都站在湖邊,親眼看見潛水人員打撈她丈夫的屍體；她仍如此鎮定,且吩咐人將這消息轉達時,要加上一節聖經──(詩一一五3).他們都相信可能這弟兄因心臟病突發而死.但使人感到驚奇的,就是這寡婦的信心與平安,雖然不明白神為何要如此對她,但仍深信神的話和順服祂的安排,還用聖經安慰別人:「然而我們的神在天上,都隨自己的意旨行事.」(詩一一五3)這使我不禁流出淚來!
\newline
\newline
正如保羅繼續在羅馬書上這樣寫道:「既是這樣,還有什麼說的呢?神若幫助我們,誰能敵擋我們呢?......誰能使我們與基督的愛隔絕呢?難道是患難麼?是困苦麼?是逼迫麼?是飢餓麼?是赤身露體麼?是危險麼?是刀劍麼?......然而靠著愛我們的主,在這一切的事上,已經得勝有餘了.因為我深信無論是死,是生,......都不能叫我們與神的愛隔絕；這愛是在我們的主基督耶穌裡.」(羅八31-39)這是基督徒的信靠.這世界很快便要消逝,但順服神旨意的人卻要永遠長存.這是每一個順服神的人的安全保證；雖然四周有紛亂,但在你心裡有完全的平安.各位,你是否有不可移動的信靠呢?你是否活在主十全十美旨意中呢?
\newline
\newline


\newpage

\section{第50屆~港九培靈會~高文~第八講　活在聖潔生活中}
\label{sec:759}
\textbf{第八講　活在聖潔生活中}
\newline
\newline
連結: \href{https://www.hkbibleconference.org/session-message/view/759}{\texttt{https://www.hkbibleconference.org/session-message/view/759}}
\newline
\newline
\hyperref[sec:758]{< < < PREV SERMON < < <}
~
\hyperref[sec:toc]{[返目錄]}
~
\hyperref[sec:761]{> > > NEXT SERMON > > >}
\newline
\newline
以弗所書 5:25-5:33
\newline
\begin{longtable}{cl}
\hline
\hline
章節 & 經文 (和合本修訂版)\\
\hline
5:25 & \begin{tabularx}{0.7\textwidth}{X} 作丈夫的，你們要愛自己的妻子，正如基督愛教會，為教會捨己， \end{tabularx} \\ \\ \relax
5:26 & \begin{tabularx}{0.7\textwidth}{X} 以水藉著道把教會洗淨，使她成為聖潔， \end{tabularx} \\ \\ \relax
5:27 & \begin{tabularx}{0.7\textwidth}{X} 好獻給自己，作榮耀的教會，毫無玷污、皺紋等類的缺陷，而是聖潔沒有瑕疵的。 \end{tabularx} \\ \\ \relax
5:28 & \begin{tabularx}{0.7\textwidth}{X} 丈夫也應當照樣愛妻子，如同愛自己的身體；愛妻子就是愛自己了。 \end{tabularx} \\ \\ \relax
5:29 & \begin{tabularx}{0.7\textwidth}{X} 從來沒有人恨惡自己的身體，總是保養愛惜，正像基督待教會一樣， \end{tabularx} \\ \\ \relax
5:30 & \begin{tabularx}{0.7\textwidth}{X} 因我們是他身體的肢體。 \end{tabularx} \\ \\ \relax
5:31 & \begin{tabularx}{0.7\textwidth}{X} 「為這個緣故，人要離開父母，與妻子結合，二人成為一體。」 \end{tabularx} \\ \\ \relax
5:32 & \begin{tabularx}{0.7\textwidth}{X} 這是極大的奧祕，而我是指基督和教會說的。 \end{tabularx} \\ \\ \relax
5:33 & \begin{tabularx}{0.7\textwidth}{X} 然而，你們每個人都要愛妻子，如同愛自己一樣；妻子也要敬重她的丈夫。 \end{tabularx} \\ \\
[1ex]
\hline
\hline
\end{longtable}
神深願我們過著美麗,聖潔的生活,整本聖經都是這主題.剛才所讀出的經文,就將聖潔帶到這焦點裡.保羅關心我們是否過著優美的人生.在此有五點真理:
\newline
\newline
(一)教會是要聖潔,毫無瑕疵(弗五26-27)聖潔,成聖或潔淨同一字根,基本意思是分別出來;我們要分別出來歸於神,是屬神的人若有人.要見神的面貌,就可以從聖徒身上看見證如羅馬人所敬拜的神是木星,非常兇惡;這是從羅馬人兇惡的性情而知的敬拜神也會影響我們道德生活水準.如同以弗所城的人,他們敬拜的女神,使他們過著不道德的生活.我們所過的生活不能高於神.今日基督徒所敬拜的是一位聖潔的神,祂的眼睛不能看見邪惡;聖潔是神的本性,所以祂願意屬祂的人也聖潔為此,聖經說:神在創世前在基督裡揀選我們,叫我們有祂的形像;但人犯了罪,這形像破壞了,人要面臨死亡的審判;所以耶穌要來這世界,目的是要將我們從罪中挽回,祂的死除去罪惡,借祂的寶血我們得以成聖;有一天,在神的聖潔中,教會要神面前成為聖潔,榮耀的教會.
\newline
\newline
(二)聖潔在基督的愛中彰顯.(25節)我們要以基督對教會的愛,為教會捨己；用此來衡度我們的愛.這種愛全無自私,不尋求自己的益處.因此聖潔並不是宗教外在的表現,乃是內心的虔誠；亦不是法律的條文,乃是靈裡的自由；同時,真正的聖潔不會產生驕傲,一個聖徒所誇口的,不過是在耶穌基督的恩典裡誇口而已.這種成聖的生活在禱告,敬拜中表現出來；也不應將聖潔與聖靈的果子相題並論.聖潔不是恩賜,乃是將聖靈的本性彰顯出來.分別為聖的生活,不是說不再遇見試探；我們仍會犯罪,但靠著神的恩典可以不犯罪,勝過試探.一次當我在花園工作時,我三歲的兒子見我汗流肩背,他想爸爸一定很口渴了；於是他走入廚房,拿起一隻杯,盛滿了水；當我走入屋,見他笑咪咪對我說:「爸爸,我想你必定很渴了.」雖然杯很骯髒,水又熱；但見他臉上天真活潑的笑容,這真是純潔的愛的表現；他盡他所能的討父親的歡喜.當我們見主面時也是一樣,雖然做錯了,但是我們的心願意討祂的喜悅,盡心,盡性,盡意的愛祂.這愛是一切先知和律法的總歸,也是聯絡全德的.
\newline
\newline
(三)在人際關係裡,愛要湧流.整段經文都是提及與人相處的信息,特別是在家裡,每日實踐過聖潔的生活.從22至23,28至33節提到作丈夫與妻子的責任.第六章作兒女對父母的責任.再提到作主人的要公平對待僕人,僕人也要誠實行事服侍主人.每一個人都有本份,以愛心待人,以尊貴,公義的態度待人；將基督的愛向他們彰顯.為此,聖潔成為傳福音,社會工作的基礎.原來基督的愛激勵我們關心別人,雖然別人不關心我們,也不要忽略他們的需要.主說:你要人怎樣待你,你也要怎樣待人.屬神的人必須經常傾流愛心,為人服務.
\newline
\newline
(四)除去生命中不聖潔的東西.罪與神的聖潔相違反,有衝突.如果我們經常在罪中生活,表示我們還被魔鬼捆鎖.基督徒要憎恨罪惡,討救主的喜悅.我們的生活需要校正,以符合基督的形像.可能今日你忽略了,或放棄了讀經,禱告的權利,不關心別人,這些事都不能討神的喜悅.自我的表現,如妒忌,驕傲,這一切都反映出人類墮落的本性；軟弱,疾病等都會一直使我們煩惱；我們必須每日對付清楚,向神承認.(約壹一9):「我們若認自己的罪,神是信實的,是公義的,必要赦免我們的罪,洗淨我們一切的不義.」基督徒不是不會犯罪,但犯罪後,他有一位代求者,赦罪的主.今晚在你心中是否有一些反叛神的事,如憎恨,妒忌.......在第二次世界大戰時,一機師要駕機飛過英倫海峽.突然他在機倉裡聽見一些聲音,是老鼠咬東西所發出的；他發覺原來那隻老鼠正在咬那些氣喉,可能會咬斷,使飛機出毛病,全機的人都會掉入大海洋裡.聰明的機師就將飛機繼續往上升,升至一個程度,使老鼠得不到氧氣而死亡,基督徒也是如此,在我們的生命中,可能有雜聲,有抵擋神的東西；不要仍活在這狀況裡,應當去親近神.既然基督的血能洗淨我們的罪,除去我們的不潔；就預備我們過成聖的生活.
\newline
\newline
(五)惟有信靠主才能過聖潔的生活,不是任何東西可以換取,也不是我們能賺回來的；乃完全是恩典.祂愛教會,為教會捨己.我們所能做的,只是接受祂的恩典.我記得這段經文在我婚禮時很有用.六月三日下午,我倆站在小禮拜堂內的壇前,接受這訓誨,彼此立下誓言.從新娘在婚禮前用心準備的潔白的嫁衣裳和她臉上的表情,表示她完全的愛我,願意終生委託於我.過去廿七年來,我越發明白婚禮誓詞的意義,而這些應用在當時婚禮中是沒有想到,也是無法領會的.我亦越來越體驗愛情是什麼,年月來,我對她的愛是與日俱增的.但相反,若新娘進入禮堂時穿上的,是骯髒的嫁衣,從她的面貌表現出一些不信賴,她還有一些自己的計劃,不想完全的信託於我；那時我的感受會如何呢!今日主見到教會的生活,祂的感受又如何呢?我們是祂的新婦,祂要看看我們是否有優美的表現.基督為教會捨去自己,為要叫教會成為聖潔,有一日無瑕疵的獻給神.各位,今晚你的心在神面前是否如此無瑕無疵?你的生活是否聖潔榮耀神,過討祂喜悅的生活?
\newline
\newline


\newpage

\section{第50屆~港九培靈會~高文~第九講　活在主再來的盼望中}
\label{sec:761}
\textbf{第九講　活在主再來的盼望中}
\newline
\newline
連結: \href{https://www.hkbibleconference.org/session-message/view/761}{\texttt{https://www.hkbibleconference.org/session-message/view/761}}
\newline
\newline
\hyperref[sec:759]{< < < PREV SERMON < < <}
~
\hyperref[sec:toc]{[返目錄]}
~
\hyperref[sec:762]{> > > NEXT SERMON > > >}
\newline
\newline
約翰一書 3:1-3:3
\newline
\begin{longtable}{cl}
\hline
\hline
章節 & 經文 (和合本修訂版)\\
\hline
3:1 & \begin{tabularx}{0.7\textwidth}{X} 你們看父賜給我們的是何等的慈愛，讓我們得以稱為神的兒女；我們也真是他的兒女。世人不認識我們的理由，是因他們未曾認識父。 \end{tabularx} \\ \\ \relax
3:2 & \begin{tabularx}{0.7\textwidth}{X} 親愛的，我們現在是神的兒女，將來如何還未顯明。我們所知道的是：基督顯現的時候，我們會像他，因為我們將見到他的本相。 \end{tabularx} \\ \\ \relax
3:3 & \begin{tabularx}{0.7\textwidth}{X} 凡對他有這指望的，就潔淨自己，像他是潔淨的一樣。 \end{tabularx} \\ \\
[1ex]
\hline
\hline
\end{longtable}
過去幾天晚上,我們看見神的心意,祂對我們在這世上生活的關懷.今天晚上讓我們注意將來的事,乃是進入永世的時代,得永遠的生命的日子.
\newline
\newline
剛才所讀的經文,提及神對祂兒女的大愛.為了證明這真理,約翰說:「世人不認識我們,因他們不認識基督.我們將來如何,還未顯明.」若要知道人生之終結如何,可以看聖經如何教導我們.歷史並不是互不相關的事實,乃是向著創造者的計劃一同演變.為此作為神的兒女,我們可以盼望什麼呢?
\newline
\newline
首先,耶穌基督要再來.我們的生命與耶穌基督的顯現,緊密的連結在一起.
\newline
\newline
在德國,一間由女執事及女傳道開辦的醫院,在第二次世界大戰時多次受空襲,屢次重建；後來一群護士在牆壁上寫著:「我們不知道將來如何,但知道誰要再來!」
\newline
\newline
今天我們亦同樣有此心願,不知前面年日有何事情會發生,或許在社會裡有大動亂,世界的大勢力會將我們的生活方式隨便轉移,患難或許會勝過我們.但無論任何事發生,我們總知道有一天,我們的主要再來!
\newline
\newline
當耶穌升天後,門徒聽見天使對他們說:「你們為什麼望著天呢?這離開你們被接升天的耶穌,你們見祂怎樣往天上去,祂還要怎樣來.」耶穌多次對門徒保證祂要再來.祂說如同閃電,從這方到那方；照樣人子要再來.祂要駕著雲,帶著榮耀和權柄再來.祂要差遣天使走在前頭,吹響號筒；從地球之東,西,南,北結集祂的子民,他們要在神榮耀寶座裡與基督同坐.每人按他們為主所作的,接受賞賜.那天,所有仇敵都要被踐踏在祂腳下.基督要在祂的國度裡作王,從永遠到永遠!
\newline
\newline
耶穌常用「人子」這稱號,早在但以理書已見到了.(但七13-14)記載:「人子駕著天雲而來,被領到亙古常在者面前,得了權柄,榮耀,國度,使各方,各國,各族的人,都事奉祂；祂的權柄是永遠的,不能廢去,祂的國必不敗壞.」在耶穌時代被稱的「人子」及「人子將來的事情」都有相同的意思.新約提及那先存在的天上者,現在來到世上.祂會將一切不敬虔者滅絕,拯救一切公義的人.祂要與祂的子民一同統治那聖潔,永不衰殘,永不過去的國度.耶穌也稱自己為那將要來的人子.祂一共用了八十二次關於「人子」的事情.雖然祂過去的生活卑微,但祂絕不低過「人子」的地位.耶穌第一次來是要做贖罪的羔羊,第二次來要作榮耀的君王,這兩次並沒有分開.祂常想到自己就是那將要來,用公義統治的那一位.在祂心中常提到神的國已經來臨了,到人子再來時,神的國就要彰顯,使每人的眼目都得以看見.難怪耶穌的人生充滿快樂.對祂來說天國是真實的.在祂腦海中常有天使組成的詩班,天使們吹起銀色號筒的聲音.祂知道神要興起屬神子民向祂歌頌.所以祂對這世界並沒有幻想,祂明知在這世界上會遭受人的厭棄,這世界並非祂的家鄉,祂的家鄉是在天上.耶穌教導你我用同樣的眼目去看這一切的事情,有長遠的眼光,用永恆的角度去觀看現今這世界.每一個消逝的時刻,將你我越髮帶近主的再來!在那時候,我們信心所仰望的都成為事實.或許這世上的事物很不稱心,但神卻要用這些事來完成祂的目的.一個不能動搖的確信,就是耶穌是萬王之王,萬主之主,祂要再來!這是教會榮耀的盼望.讓諸天歡樂,讓屬神的子民今天來歡呼:耶穌基督要再來!我們仍在地上的,要存快樂的心等候祂來!
\newline
\newline
「將來如何,還未顯明；但我們知道主若顯現,我們必要像祂.」(約壹三2).在此約翰將第二個偉大的真理向我們講明:神的兒女們會成為耶穌基督身體的形像.若主未回來,我已離世,身體便要歸塵土,但我的靈要與主同在.直到主回來時,祂會帶我一同回來,與自己身體相連.保羅在(帖前四14-18)已講明這真理.他又說:「祂要按著那能叫萬有歸服自己的大能,將我們這卑賤的身體改變形狀,和祂自己榮耀的身體相似.」(腓三21)
\newline
\newline
信徒復活後之身體是怎樣呢?請看耶穌復活的身體:門關了,耶穌能在房間出現,又能升上天去.這世上的權力不能控制束縛我們的主.主叫門徒用手摸祂手裡的釘痕,肋旁的傷痕；他們很清楚知道仍然是耶穌的身體,不過其中的質素不同了.保羅在寫給哥林多教會的信說:所種的是羞辱的,復活的是榮耀的；所種的是軟弱的,復活的是強壯的；所種的是血氣的身體,復活的是靈性的身體.(林前十五43-44)將來復活的身體是最美,最強壯的,無年齡的限制,是有史以來「靈」真正的出現,彼此能辨認,互有交通.在這情形下與耶穌的形狀相連,相同.
\newline
\newline
約翰再加上第三點寶貴的真理:我們能面對面看見祂的身體,乃是我們的靈與神的靈直接交通.因我們被造時是按祂的形像而造,同樣我們的靈也像基督的靈.
\newline
\newline
當約翰被放逐在拔摩海島上,他看見神的異象:「聖城耶路撒冷由神那裡從天而降,預備好了,就如新婦妝飾整齊,等候丈夫.」他又聽見有大聲音從寶座出來說:「看哪!神的帳幕在人間,他要與人同住；他們要作祂的子民,神要親自與他們同在,作他們的神.......(啟二十一2-5)這表明基督要與祂子民親密的關係,連整個聖城內的設計都反映出神的榮耀:如同極名貴的寶石,好像碧玉,明如水晶.(啟二十一11)城是四方的,長闊一樣(啟二十一16).這和諧的設計反映出我們與神緊密的關係.城是精金的,如同明淨的玻璃.城牆的根基是用寶石修飾.(啟二十一18-19)城門是一顆珍珠,城內的街道是精金.(啟二十一21)但約翰舉目卻不見聖殿,然後他才知道神的城不再用人手所作的殿來表明神的同在,因主全能者和羔羊為城的殿.(啟二十一22)列國要在城的光裡行走,地上的君王必將自己的榮耀歸與那城.(啟二十一24)城當中有一寶座,一道生命水的河,明亮如水晶,從寶座流出.河邊種有生命樹,果子用來醫治萬民.(啟二十二1-2)祂的民都要服事祂,有名字刻在他們的額上.不再有黑夜,也不用燈光和日光.因主要光照他們,他們要作王,直到永永遠遠!(啟二十二3-5)這是我作子民的盼望和心裡的確信.
\newline
\newline
我所任教的神學院──亞斯伯理神學院之院長,喜歡提及他一次的經驗.當他從非洲回國時,正與羅斯福同船；這著名的政治家剛狩獵回來.船到達紐約港口,整個紐約市的居民都出來迎接羅斯福,岸上有樂隊奏樂,海上的小船也響起笛來歡迎,連紐約市長也出來歡迎他.但與羅斯福同船的院長卻頓時感到最孤單,被人奚落.他想起自己在非洲,為宣教師如何辛勞,犧牲.現在他回到自己的故鄉,卻沒有一人迎接他.但忽然間,紐約市在他眼前慢慢消失；他在想,在他眼前出現的是一座美麗的城,用精金作的城,主耶穌坐在城的寶座中；他聽見親愛主的聲音；莫利遜呀!歡迎你回來你真正的家鄉,我等候你已久!雖然將來如何還未知,但榮耀的盼望必再來!
\newline
\newline
最後,請留心約翰在(約壹三3):「凡向祂有這指望的,就潔淨自己,像祂潔淨一樣.」約翰又說:「既是神的兒女,就不要再故意犯罪.因為一切污穢,行可憎惡的事,不能見神的國.」我們不能帶著不潔的手,不純的心到神的面前.所以聖經說:新婦必須妝飾整齊,等候丈夫.衣裳必須潔白.
\newline
\newline
作為教會一份子的你,你是否預備好迎接主呢?你的生活是否如新婦等候新郎呢?但願我們都預備好見我們的主!
\newline
\newline


\newpage

\section{第50屆~港九培靈會~高文~第十講　得勝的生活}
\label{sec:762}
\textbf{第十講　得勝的生活}
\newline
\newline
連結: \href{https://www.hkbibleconference.org/session-message/view/762}{\texttt{https://www.hkbibleconference.org/session-message/view/762}}
\newline
\newline
\hyperref[sec:761]{< < < PREV SERMON < < <}
~
\hyperref[sec:toc]{[返目錄]}
~
\hyperref[sec:663]{> > > NEXT SERMON > > >}
\newline
\newline
啟示錄 12:11-12:11
\newline
\begin{longtable}{cl}
\hline
\hline
章節 & 經文 (和合本修訂版)\\
\hline
12:11 & \begin{tabularx}{0.7\textwidth}{X} 弟兄勝過那條龍是因羔羊的血，和因自己所見證的道。雖然至於死，他們也不惜自己的性命。 \end{tabularx} \\ \\
[1ex]
\hline
\hline
\end{longtable}
我們思想過優美的人生,優美的人生會從讚美湧溢出來,為這原因,基督徒應經常唱詩.在啟示錄中記載天上所唱的詩歌,第四章至第九章不斷的出現天上唱詩的記載.大部份記載主要在地上施行審判,但偶然地上的情景轉移到天上,在短暫時間我們見到神的寶座,在寶座周圍見到一大群各方,各族的人正在歌唱,他們向神唱讚美詩,其中一首在(啟十二10)「我聽到天上有大的聲音說,我神的救恩,能力,國度,並祂基督的權柄,現在都來到了,因為那在我們神面前晝夜控告我們弟兄的,已經被摔下去了」在神的寶座周圍聽見這種詩歌,真是美好了但我不想大家太過天真,聖經記載我們還有一位仇敵,就是魔鬼;整個地獄都會起來威嚇教會,這些邪惡的勢力要盡它所能向善良的摧殘,只要在地上生活,你我都要投身在一場神聖的戰役裡;聖說在末日臨近時,這爭戰的程度就越增加.耶穌基督說在末日會有更厲害的患難,但不論發生何事,教會必須成為得勝的教會.撒旦是戰敗的仇敵,在神永恆的計劃中,撒旦吃敗仗是必然的事實.各位是否注意詩歌的內容呢?救恩現在來臨了,神的權柄,國度現在來臨了,而那位控告弟兄的已經被摔下了.跟著第十一節提到勝利的原因.
\newline
\newline
(一)弟兄勝過它是用羔羊的血,這是地方教會得勝信息的總結,對一切使徒是好消息,事實上我們的盼望惟獨在耶穌基督的寶血和祂的身上；聖經提到血,就提到神捨棄祂兒子的生命,為我們的罪而流出寶血.人藉寶血就得到各樣的福氣──稱義,成聖,得救贖等等.聖經用這個「血」字一共四百六十次,血講到救贖,挽回,成為祭司等等的意思.既然血用了那麼多次,究竟在聖經中有沒有一頁是不提到血的呢?其實整本聖經就如被一條染滿血的繩穿過一樣,來證明神的恩典；當然在十字架就將寶血帶到焦點.一些地方將羔羊與血相提並論,在舊約時候,用到血是指明神與人和好的關係,人得到救贖；在祭壇上,寶血表示神對人的審判；亦表明神接納那些來獻祭的人.當那些獻祭的人殺牛羊的時候,乃是承認他們的罪；讓牛羊擔當他們的罪,而他們在神面前看自己如死去一般.因此就得以與神重新建立一個全新的關係.舊約的獻祭乃是應許將來要來的那位耶穌基督,他將自己的血獻上,好將我們帶到神的面前.當耶穌基督在十字架快流盡寶血的時候,祂大聲呼叫:「成了!」祂乃是宣告舊約一切獻祭制度上所預表的,現在已經完成.我記得一位浸信會牧師在參觀聖地後,講出他的感受,當他眼見到十字架山時,他從旅行團中首先的走上去,站在山腳下,他的頭低下,最後,旅行團嚮導打破緘默.他問:「先生,你以前曾否到此?」靜默了一刻,那位牧師就說:「大約二千年前,我已經來過」.親愛的弟兄姊妹,二千年前你我都曾到過這裡；當耶穌受死時,祂取替了我們的位置,我們本來都在定死罪的情況下,但祂的血使我們再一次得著清潔,祂的血今天仍然生效,這是教會得勝的信息.
\newline
\newline
(二)藉著神自己的道,耶穌基督在十字架上的勝利必須被傳揚,叫所有的人都要來到主的面前.如果我們不傳揚這個信息,就是攔阻了神的工作.所以,我們每一個基督徒都是耶穌基督的見證人.我記得有一次我學校的一位學生要坐巴士去另一處地方,他一路坐車一路讀聖經,而旁邊的搭客說他從未見過有人一路坐車一路讀聖經的.這成為那位同學的機會向旁邊的搭客講見證.他說出耶穌基督在他的生命中是真正活著的一位,當他說完後,那位搭客就說:「我從來沒有聽過這樣的事,其實,在巴士中,還有一位需要這樣的信息.」那位同學就說:巴不得我有機會向他講:所以,他們安排調換座位,現在那位同學再有機會講見證.當他快講完時,前面的男搭客回頭向他發出問題:「剛才最後的一句,你可否再重複一次,因為我都想聽.」而旁邊一行的女士也非常好奇,最後有人要求他大聲講,他就走到巴士的前面,盡量的大聲向他們述說耶穌愛他們.耶穌基督死在十字架上贖我們的罪,然而祂從死裡復活,藉著祂的能力,這位同學的生命得到改變,最後他加上一句話:「因為耶穌基督的緣故,我能愛巴士上的每一位.」在這時候.巴士已到了終站,巴士司機站起來對他說:「你還有什麼要講的呢?多謝神他很有心,要講的差不多都講完了.」那位同學就低聲說:「哈利路亞!」各位,當你面上帶著得救的笑容,和有喜樂的共鳴時,世人就從你身上認出救主.各位!請問上一次有人聽見你講有關耶穌基督的話,是什麼時候呢?
\newline
\newline
教會的得勝和耶穌基督的道,被見證在彼得的身上表明出來,耶穌曾問門徒到底祂是誰?彼得就說出:「你是基督,是永生神的兒子.」這句話得主的喜悅.耶穌基督指出這句話不是人指示彼得講的,乃是從神那裡所領受的；耶穌跟著說:「我要將我的教會建立在這磐石上面,陰間的權勢都不能勝過他.」有時我們想到磐石時,就想到耶穌基督,或者彼得所講出的是正統的答案,成為教會的代表；在此我們認識彼得正在宣告他的救主,惟有彼得作此見證後,耶穌宣告要有教會的建立.單單知道耶穌是那一位應許而來的救主,是不足夠的,這好消息必須被宣告出去,不單成為一項信條,乃要成為基督徒個人經歷的活生生的見證.如果耶穌基督未曾被人宣告出去,世人如何能相信呢?人惟有聽道,才能信道.除非我們今日出去宣講這個信息,否則下一個世代不會有教會.我們相信每逢耶穌基督被高舉的時候,就能吸引萬人來歸祂.
\newline
\newline
(三)作見證者的生命必須與他的見證符合「他們雖至於死,也不愛惜自己的生命.」在此使我們見到有效的傳揚福音的方法就是見證.「見證」這二字的原意是殉道,如果我們接受主的話,今日我們為主作見證,就是要為主殉道；我們存著愛主的心,順服主的心,對自己否定.不久前,我在美國的印第安納州開汽車,見到一間教會的門前,寫著:「當稱謝進入祂的門,當讚美進入祂的院.」(詩一○○4)我想教會的門有這句話是非常恰當的.各位,當你們進入神的門時,有這樣的字,你們的心一定很歡喜的.那些字是從詩篇一百篇所引述出來,是否記得在這些話之前,提到「你們是祂草場上的羊」.其實提到草場上的羊,我只想到一件事:就是能夠入到耶路撒冷,真正入到神的門.進入神的門作甚麼呢?豈不是要成為祭壇上的祭牲嗎?其實這裡所講的是指著我們,我們是祂的羊.然而我們想到我們高舉耶穌基督時,我們願意委身給祂,就算死也在所不惜.親愛的弟兄啊!我以神的慈悲勸你們,將身體獻上,當作活祭.當我們完全放在祭壇上時,我們就完全屬於祂的了.這是成聖的原則.「聖徒」的意思乃是為神分別出來的人,一個屬於神的人；在一個基督徒的人生開始時,成聖就是這個意思,我們必須委身在神的旨意中.當我們在耶穌基督的恩典和經驗中增長時,我們的委身也就跟著增長.我們向神委身這件事,從頭到尾都是如此認真的.聖經說:「凡屬基督耶穌的人,是已經把他的肉體,連肉體的邪情私慾,同釘在十字架上了.」(加五24)我們對罪看自己是死了,向神卻當看自己是活著的,所以我們或活或死,總是主的人.(羅十四8)
\newline
\newline
不久前,在非洲的一個地方有一次叛亂,一天晚上,那些叛亂份子坐吉普車來到一個福音站的前面,當時他們喝令福音站內的校長和傳道,都要上吉普車.他們就將吉普車駛到最近的一條河邊,停在那裡.他們為了要對付這兩個基督徒,而起了爭論；那位牧師就知道前面將遭遇什麼事情.他遂問校長說:「你在耶穌基督裡是否有了確據?」那位校長說:「我已經信了主」.牧師就記下那日的時間,然後加上一句；他是在那日死的?答:就在此時.那些叛軍喝令牧師從橋上走過去,當他走在橋上時,他就唱起歌來:「我憑信離遠就可以看見,我天父在那住處常等待著我,天父在那處為我們預備美好的地方,在那甜蜜的地方,隨著日子的消逝,我們要彼此相見,在那美麗的海岸相見,我們要在那榮美的地方再見.」那日,他根本沒有機會唱完那首詩歌,當他走到橋中時,他已被鎗殺了.那些叛軍走過去,將屍體拋落河中.那位校長想到第二個就是輪到他了,但他想不到那些叛軍又開始爭吵起來,然後他們走回汽車,開車去了.後來他述說這件事,提到那些叛軍看見牧師殉道前,仍然唱著詩歌,覺得非常詫異；從來沒有人見過唱詩來殉難的.這種向主的委身,一直都是令世人覺得希奇的,這種榮耀的盼望能夠帶著歌唱到主的面前,無論代價是甚麼,總要跟從基督到底.我們知道在座當中可能有親戚在國內為主的緣故這樣殉道,或許神今日沒有呼召我們作這種肉身上的受難殉道；但在我們的心靈中,向主獻上這個心願,已經成為事實；我們應該認定自己成為基督的殉難者.所以不論外在的環境如何惡劣,我們靠著耶穌基督已經得勝有餘了.我們將來如何還未顯明,但我們知道主若顯現,我們必要像祂(約壹三2)萬膝都要在主的面前下拜,萬口都要稱耶穌基督乃是主,使榮耀歸與父神.(腓二10-11)就不住的為這緣故,歡呼歌唱.因為基督使我們從罪惡,從死亡裡得著釋放.祂已經得勝,魔鬼完全的被打敗.各位!難道你現在聽不見天上凱旋的歌聲嗎?
\newline
\newline
我們得勝了,是靠著羔羊的寶血,靠著我們見證的道,也因為我們不愛惜自己的生命,忠心至死.各位!你現在是否有此經驗?這是得勝教會的經驗.在天上已經成為的事實,在今日也應該成為我們心中的經驗,耶穌基督昨日,今日,直到永遠,都是一樣的.各位!假若今日你的生活不在這種得勝之中,你就是活在一種可憐無力的生活當中.讓我們一齊低頭祈禱省察吧!
\newline
\newline


\newpage

\section{第51屆~港九培靈會~唐佑之~第一講　四匹馬的異象}
\label{sec:663}
\textbf{第一講　四匹馬的異象}
\newline
\newline
連結: \href{https://www.hkbibleconference.org/session-message/view/663}{\texttt{https://www.hkbibleconference.org/session-message/view/663}}
\newline
\newline
\hyperref[sec:762]{< < < PREV SERMON < < <}
~
\hyperref[sec:toc]{[返目錄]}
~
\hyperref[sec:674]{> > > NEXT SERMON > > >}
\newline
\newline
撒迦利亞書 1:1-1:17
\newline
\begin{longtable}{cl}
\hline
\hline
章節 & 經文 (和合本修訂版)\\
\hline
1:1 & \begin{tabularx}{0.7\textwidth}{X} 大流士王第二年八月，耶和華的話臨到易多的孫子，比利家的兒子撒迦利亞先知，說： \end{tabularx} \\ \\ \relax
1:2 & \begin{tabularx}{0.7\textwidth}{X} 「耶和華曾向你們祖先大發烈怒。 \end{tabularx} \\ \\ \relax
1:3 & \begin{tabularx}{0.7\textwidth}{X} 你要對以色列人說，萬軍之耶和華如此說：你們要轉向我，這是萬軍之耶和華說的，我就轉向你們，這是萬軍之耶和華說的。 \end{tabularx} \\ \\ \relax
1:4 & \begin{tabularx}{0.7\textwidth}{X} 不要效法你們的祖先。從前的先知呼叫他們說：『萬軍之耶和華如此說，當回轉離開你們的惡道惡行。』他們卻不聽，也不順從我。這是耶和華說的。 \end{tabularx} \\ \\ \relax
1:5 & \begin{tabularx}{0.7\textwidth}{X} 你們的祖先在哪裡呢？那些先知能永遠存活嗎？ \end{tabularx} \\ \\ \relax
1:6 & \begin{tabularx}{0.7\textwidth}{X} 然而我的言語和律例，就是我所吩咐我僕人眾先知的，豈不臨到你們的祖先嗎？他們就回轉，說：萬軍之耶和華定意按我們的所作所為對待我們，他也已經照樣行了。」 \end{tabularx} \\ \\ \relax
1:7 & \begin{tabularx}{0.7\textwidth}{X} 大流士第二年十一月，就是細罷特月二十四日，耶和華的話臨到易多的孫子，比利家的兒子撒迦利亞先知，說： \end{tabularx} \\ \\ \relax
1:8 & \begin{tabularx}{0.7\textwidth}{X} 「我夜間觀看，看哪，有一人騎著紅馬，站在窪地的番石榴樹中間。在他身後有紅色、褐色和白色的馬。」 \end{tabularx} \\ \\ \relax
1:9 & \begin{tabularx}{0.7\textwidth}{X} 我說：「主啊，這是甚麼意思？」與我說話的天使說：「我要指示你這是甚麼意思。」 \end{tabularx} \\ \\ \relax
1:10 & \begin{tabularx}{0.7\textwidth}{X} 那站在番石榴樹中間的人回答說：「這是奉耶和華差遣，在遍地巡邏的。」 \end{tabularx} \\ \\ \relax
1:11 & \begin{tabularx}{0.7\textwidth}{X} 他們對站在番石榴樹中間耶和華的使者說：「我們在遍地巡邏，看哪，全地都安息平靜。」 \end{tabularx} \\ \\ \relax
1:12 & \begin{tabularx}{0.7\textwidth}{X} 於是，耶和華的使者說：「萬軍之耶和華啊，你惱恨耶路撒冷和猶大的城鎮已經七十年了，你不施憐憫要到幾時呢？」 \end{tabularx} \\ \\ \relax
1:13 & \begin{tabularx}{0.7\textwidth}{X} 耶和華就用美善的話和安慰的話回答那與我說話的天使。 \end{tabularx} \\ \\ \relax
1:14 & \begin{tabularx}{0.7\textwidth}{X} 與我說話的天使對我說：「你要宣告，萬軍之耶和華如此說：我為耶路撒冷而妒忌，為錫安大大妒忌。 \end{tabularx} \\ \\ \relax
1:15 & \begin{tabularx}{0.7\textwidth}{X} 我非常惱怒那享安逸的列國，因我從前稍微惱怒，他們就越發加害。 \end{tabularx} \\ \\ \relax
1:16 & \begin{tabularx}{0.7\textwidth}{X} 所以耶和華如此說：現在我回到耶路撒冷，仍要施憐憫，我的殿要重建在其中，準繩必拉在耶路撒冷之上。這是萬軍之耶和華說的。 \end{tabularx} \\ \\ \relax
1:17 & \begin{tabularx}{0.7\textwidth}{X} 你要再宣告，萬軍之耶和華如此說：我的城鎮要再度繁榮發達。耶和華必再安慰錫安，揀選耶路撒冷。」 \end{tabularx} \\ \\
[1ex]
\hline
\hline
\end{longtable}
「神的榮耀」在以色列民族,和教會歷史中的事實,都成為我們屬靈的感受.撒迦利亞書通常被稱為小先知書,這非因先知有大小之分,只是指書的篇幅長短而言,對信息及內容的重要性並無影響.撒迦利亞書在小先知書中最長,共十四章,內容主要分作兩大段:即一至八章及九至十四章.前者屬先知文學,亦即預言文學；後者則屬啟示文學,其內容與舊約但以理書,和新約的啟示錄有點相似.在這寶貴的書卷中,以不同的重點和方法表達出神奇妙的啟示.
\newline
\newline
現在我們先看一至六章的八個異象,這些異象可作為教會更新的指導.在培靈會期間,我們會注意這八大異象,同時更特別注重啟示文學中末世的預言.
\newline
\newline
每當研究啟示文學的預言時,有兩件必須注意的事:即預言的歷史背景,和先知信息的特點.這樣可看出神在以色列民族歷史信仰的經驗,然後我們便可應用於今日教會的需要,並知道自己該作些甚麼.
\newline
\newline
在舊約的先知書中,除了幾部以外,大都有詳細的歷史背景.我們不要以為那些信息只屬於神對以色列人古老的啟示,事實上這些啟示是永遠常新的.因為神不但在以色列民族,世界歷史,每一個民族及歷代教會之中；更在個人生中讓我們能認識祂,相信祂並親自聽見及知道祂對我們的旨意.因此讀舊約的歷史書,先知書時,基督徒必須要有歷史的認識,不但對過去的歷史有所認識,更應注意在今天的時代環境中,神要我們作些甚麼.這帶給我們基督徒一種緊急感,知道還有不多的時間,那要來的就來了；若有人不愛主,這人是可咒可詛的,因為主再來的日子近了.面對這時代我們要有緊急感,快點起來為主作工；因為神給我們時代的使命感,好像先知撒迦利亞一樣.
\newline
\newline
撒迦利亞是以色列人被擄歸回時代的一位先知,雖然以色列人是神的選民,當他們不愛神時,神的公義還是向他們彰顯,要臨到他們,可是以色列民族敗亡,甚至被擄到外邦,而整個民族陷於瓦解.人失敗了,但神永不失敗!祂是信實的,正如經上所記載的.以色列敗亡七十年後,民族還要復興,那些歸回的人,「崇拜」是他們第一件注意的事,聖殿必須重建.耶路撒冷和民族的重建中,崇拜必須先恢復,見證的事重新建立起來.他們重新要得著神的恩惠,注意神託付他們的使命.撒迦利亞與哈該是復興的先知(還有較後期的瑪拉基),因此我們通常都把哈該與撒迦利亞相提並論.他們特別注重聖殿重建.這是波斯王大利烏在位的時期,即主前五百二十年.哈該回到耶路撒冷時,年紀老邁,他曾見過從前的聖殿,因此極了解當時群眾的感受.想到神的管教,歷史的苦難和民族的憂患；他便安慰當時的百姓,說將來的殿的榮耀要遠勝過舊日的.雖然沒有所羅門建的那麼榮華,但屬神非物質而屬靈的榮耀是永遠的,內在的；所以哈該與撒迦利亞都把握著這重要的意義.哈該這位年老的先知,講的信息似乎沒有多大收穫,卻是能直接,簡單而有力的說出.年青的撒迦利亞可能是生長於外邦,從沒有見過舊日的聖城和聖殿,因此撒迦利亞書的信息與哈該書不同.它充滿屬靈的魄力,想像力,更有屬靈的理想,把神的信息,更有力和美麗地表達出來.神分別用哈該,撒迦利亞,我們若把自己完全交託給神,那麼我們必能得著神的信息,作這時代的見證人.
\newline
\newline
先知信息的特點是非常注意歷史的意識,神藉著歷史把祂的心意表達出來.沒有一件歷史事實是未經神許可的,因祂更在歷史之上統管萬有,把祂的公義慈愛表達出來.歷史具有承續性,非偶然產生的；是一種歷史的因果,正如將來要承受過去和現在的後果,因此你我都要負起歷史的責任.若今天我們在主面前盡心竭力,為明日教會的屬靈光景,負完全的責任,這就是先知所有的感受了.
\newline
\newline
在先知書中很注重兩件事情.信息一詞在翻譯上用不同的方法,有不同意義.它可以作「話語」解,本來是負擔之意,有時是默示,神把負擔放在先知身上,先知把這託負帶出來,神託負我們,一如先知屬靈的感受.第二方面,有時是「默示」,是指看見的異象之意.先知不但有特別的見解,且看見和體會別人所不能的.
\newline
\newline
異象是異乎尋常的事,神以超自然的方法使人看見異象,但有時並非如此,甚至只是以生活日常的經驗使人有屬靈的感受.箴二十九18說:「沒有異象,民便放肆.」其實異象在這裡是指先知的話語.如果我們沒有神的話語,將必變得墮落,馬虎.神有時要我們看見特別的景象,但有時給的只是屬靈的感受,使我們在神的話語上下功夫,體會祂的心意.神使先知看見,明白,且去傳揚見證,更要順服；你我也需如此.
\newline
\newline
經文中有兩處不同的表達,一句是「萬軍之耶和華如此說」,另一句是「這是萬軍之耶和華說的.」這兩句話看來差別不大,實則大有不同；前句是表明神是非常嚴正,嚴肅,嚴厲,且帶著權威,有絕對全能的力量和命令,人必須要謹慎且敬畏地領受這一項正式公義的宣告.「這是耶和華說的」一句即不如前者,「說的」在翻譯上是一種嘆息,極輕微的聲調,充滿了溫柔憐憫,有時甚至帶著安慰,也像悲愴的歎息,表示神非常失望,難過,因為祂的選民失敗了!這屬神的人卻沒有注意神的話語,沒有追求祂的道,不能符合祂對我們的要求,這便是耶和華的嘆息.當我們讀完先知書的時候,我們也該有同樣的感受；求神憐憫並感動我們的心,叫我們與主有同樣的心,而且通過信心付諸行動.
\newline
\newline
「耶和華如此說,你們轉向我,我就轉向你們」.神不需要轉動,只有人需要.正如那浪子與父親背道而馳,直至他醒悟後,才返回父親的懷裡.何西阿先知帶著淚和歎息大聲疾呼的說:「來吧!我們歸向耶和華.」(歸向就是悔改之意).當以色列人歸向神時,神就如晨光出現,神不改變也沒有轉動的影兒,如果信徒感覺神遙遠,抽象的話；問題不在神,而是我們需要轉向祂而已.
\newline
\newline
「你們要回頭離開,」我們要離開,回頭和轉回.這是第一篇信息或者是全書的序言.他們的列祖和先知在哪裡呢?人會過去,神的僕人亦然,但神的工作卻永遠長存.培靈會已是五十一屆,五十年間有多少神的僕人在傳神的信息；有的已成過去,但神的話語,工作不會過去.今天你我亦將過去,但凡遵行神旨意的仍永遠長存.讓我重複這點,我們必須要對歷史負責,向神及以後的人交代.這就是神所告訴撒迦利亞的.
\newline
\newline
(一)四匹馬的異象(1-8)
\newline
\newline
現在我們正式來看第一個異象(第八節).「我夜間觀看」,這是一個晚上的景象,若注意先知書的內容,在被擄以前,異象不在晚上而在白天.所以連先知對晚上的異象也感到疑問.自被擄以後,撒迦利亞或但以理卻特別注意晚間的異象.因為神要有黑夜的信息,要提醒先知們在歷史的黑夜中,在神的光中,看見亮光.同樣我們也要注意黑夜的信息,處境,以致在夜間所該看見且要作的事.自亞當夏娃犯罪後,那一夜的黃昏,人類歷史便陷入罪惡,無知痛苦的漫漫長夜中,而這正是整本聖經要告訴我們的重要信息.當以色列人在埃及地時,因埃及王的阻攔,便有了黑暗之災；但在以色列人居住之地,卻有亮光,這表明世界都在黑暗中,以賽亞書說:「在阿拉伯的曠野,有人聲說,守望的人夜裡如何?」黑夜是否會過去呢?記得主耶穌來到這世界時,並不如現在聖誕節那麼富美麗的傳奇性.事實上,神的兒子降生在伯利恆的晚上,那明星──人類歷史希望之星,升起了；光來到世界,但世界的人因犯罪不愛光,倒愛黑暗.主耶穌曾說:祂是世界的光,跟從祂的就不在黑暗裡走.但以後的人卻把祂釘在十字架,那時從午正至申初,遍地都黑暗了.人類只有在主耶穌從死裡復活的那天,才有了早晨.我們現正等待另一個早晨,主耶穌從天降臨的那個榮耀的早晨,但現在還未到.在啟示錄中,約翰看見榮耀的人子,祂的面貌如烈日放光；當這榮耀的人子,行走在金燈臺之間,表明了教會仍在歷史的黑夜之中.我們一直在守主的晚餐,是黑夜的守望者；正如先知撒迦利亞,又如耶路撒冷的祭司們,按著更次守夜,靜待天亮.忍耐,忠心,信心,恆心,直至東方發白,號角傳來之時,耶路撒冷之聖殿就要獻上早祭.這是當時歷史的背景,不單是猶太教會的習俗,也給我們一幅活潑的圖畫.我們都是黑暗中的守望者,都不能昏昏欲睡,要像滿了油的聰明童女,需要儆醒.在歷史長且深的黑夜中,我們怎能睡覺呢?因此彼得曾說:「萬物的結局近了,所以我們必須謹慎自守,儆醒禱告.」你我都是黑夜中的守望者,正如撒迦利亞,他看見一人騎著紅馬,站在窪地番石榴樹中間.在低地,山谷的地方,或許有象徵的意義,今天我們都活在山谷的地方,是潮濕且陰暗,我們都在這環境中,世人都沒有希望,但我們仍有.讓我們抬起頭來,仰望高山︴白雲和陽光,我們的幫助從何而來?我們要舉目往上看,因為我們的幫助從耶和華而來.
\newline
\newline
再者,番石榴是極矮小的樹木,其味卻幽香.我們在這世界,環境也該如此,不能隨波逐流,要成為中流砥柱；見證主,發出祂的香氣.雖然不顯著,威武,但卻有地位,價值.基督徒就是如此卑微,平凡,但卻又是偉大；因為神的生命,在我們身上照常顯大.
\newline
\newline
這人就站在那裡,在他背後,還有數匹馬,有紅,黃和白色的.有注重末世預言的人解釋牠們顏色所帶有的屬靈意義.其中易產生過份的危險,有幾點是需要注意的:「紅」是代表鮮血,爭戰 ;「白」是純潔,得勝 ;「黃」在中文該譒為混雜,即斑點,不紅不白.這些馬是表明世界在烽火中,沒有太平,安全,戰爭是末世的現象.
\newline
\newline
於是我對與我說話的天使說:「主阿!這是甚麼意思呢?」「主阿」在這裡不一定是道成肉身以前的基督,但最少也是尊貴的稱呼.他說:「我要指示你這是什麼意思」.十一節:「那站在番石榴樹中間的人說,這是奉耶和華差遣,在遍地走來走去,見全地都安息平靜.」世界在烽火之中,矛盾地是安息,平靜.外表看來,世界是繁榮平靜；卻又是混亂,充滿爭戰.人在外面粉飾太平,實在卻沒有一點平安.在耶利米書和以西結書中,都提及假先知,他們都說平安的話,其實不然.愛聽那些輕描淡寫,旁敲側擊,輕鬆幽默的話,這並非聽道的態度,神的信息是嚴厲及有份量的.求神憐憫,賜福教會的講台,不要逃避現實,不能以信心的眼光,來否認混亂的事實.先知的信息常常是現實的,也是理想的,兩者皆需注意.因此神極其惱恨.「萬軍之耶和華阿!你惱恨耶路撒冷要到幾時呢?已經七十年,你不施憐憫要到幾時呢?」那些作惡的外邦人,他們還在享受,以為是平安的生活.若研究歷史,可知當時波斯在政治上不安,人卻在那裡粉飾太平.我們不能如此,要有先知的眼光,作為當代的見證人,傳道者,見證神的公義,傳審判末世的信息.
\newline
\newline
耶和華就是美善的安慰者.(十三節)安慰與憐憫同義,是神內心的歎息和焦急,而且說到耶和華神「心裡火熱」,又說神「不單安慰錫安,也揀選了耶路撒冷.」我們可見神是憐憫的,是心裡火熱的,火熱是忌邪之意,祂火熱即擊打那些作惡者和仇敵們.在以賽亞書中說你們要安慰我的百姓,對錫安說安慰話.安慰有兩個字,一是指神有迫切說不盡憐憫的心賜；另一個是摸到他們的心意.我們的主就是這樣,作為教會工人的也該如此.我們是守望者,不單看見,還要見證我們的神──是滿有憐憫,安慰的神,更要摸到人的心.今天有許多人在痛苦之中,神要我們用溫柔的手來包裹破碎的心靈.我們有安慰的功能,更可治邪,我們不能與世界同流合污,不能妥協.神的心如何,我們也要如何.我們必須像神一樣恨惡罪惡,總要心裡清清潔潔的,合乎主用,預備行各樣的善事.
\newline
\newline


\newpage

\section{第51屆~港九培靈會~唐佑之~第二講}
\label{sec:674}
\textbf{四角與準繩的異象}
\newline
\newline
連結: \href{https://www.hkbibleconference.org/session-message/view/674}{\texttt{https://www.hkbibleconference.org/session-message/view/674}}
\newline
\newline
\hyperref[sec:663]{< < < PREV SERMON < < <}
~
\hyperref[sec:toc]{[返目錄]}
~
\hyperref[sec:675]{> > > NEXT SERMON > > >}
\newline
\newline
撒迦利亞書 1:18-2:13
\newline
\begin{longtable}{cl}
\hline
\hline
章節 & 經文 (和合本修訂版)\\
\hline
1:18 & \begin{tabularx}{0.7\textwidth}{X} 我舉目觀看，看哪，有四隻角。 \end{tabularx} \\ \\ \relax
1:19 & \begin{tabularx}{0.7\textwidth}{X} 我問那與我說話的天使：「這是甚麼意思？」他對我說：「這是擊散猶大、以色列和耶路撒冷的角。」 \end{tabularx} \\ \\ \relax
1:20 & \begin{tabularx}{0.7\textwidth}{X} 耶和華又把四個匠人指給我看。 \end{tabularx} \\ \\ \relax
1:21 & \begin{tabularx}{0.7\textwidth}{X} 我問：「這些人來做甚麼呢？」他說：「那是擊散猶大的角，使人不敢抬頭；但這些匠人前來威嚇他們，打掉列國的角，因為他們舉起角來擊散猶大地。」 \end{tabularx} \\ \\ \relax
2:1 & \begin{tabularx}{0.7\textwidth}{X} 我舉目觀看，看哪，有一人手拿丈量的繩。 \end{tabularx} \\ \\ \relax
2:2 & \begin{tabularx}{0.7\textwidth}{X} 我問：「你到哪裡去？」他對我說：「要去丈量耶路撒冷，看有多寬多長。」 \end{tabularx} \\ \\ \relax
2:3 & \begin{tabularx}{0.7\textwidth}{X} 看哪，與我說話的天使出去，另有一位天使迎著他來， \end{tabularx} \\ \\ \relax
2:4 & \begin{tabularx}{0.7\textwidth}{X} 對他說：「你跑去告訴這個年輕人說，耶路撒冷必有人居住，如同無城牆的鄉村，因為其中的人和牲畜很多。 \end{tabularx} \\ \\ \relax
2:5 & \begin{tabularx}{0.7\textwidth}{X} 耶和華說：『我要作耶路撒冷四圍火的城牆，並要作城中的榮耀。』」 \end{tabularx} \\ \\ \relax
2:6 & \begin{tabularx}{0.7\textwidth}{X} 耶和華說：「來，來！你們要從北方之地逃回；因我曾把你們分散到天的四方。這是耶和華說的。」 \end{tabularx} \\ \\ \relax
2:7 & \begin{tabularx}{0.7\textwidth}{X} 來！住巴比倫的錫安百姓啊，逃吧！ \end{tabularx} \\ \\ \relax
2:8 & \begin{tabularx}{0.7\textwidth}{X} 萬軍之耶和華在顯出榮耀之後，差遣我到擄掠你們的列國那裡，他如此說：「碰你們的就是碰他自己眼中的瞳人。 \end{tabularx} \\ \\ \relax
2:9 & \begin{tabularx}{0.7\textwidth}{X} 看哪，我要揮手攻擊他們，他們就必作自己奴僕的擄物。」你們就知道萬軍之耶和華差遣了我。 \end{tabularx} \\ \\ \relax
2:10 & \begin{tabularx}{0.7\textwidth}{X} 耶和華說：「錫安哪，應當歡樂歌唱，因為，看哪，我要來，要住在你中間。 \end{tabularx} \\ \\ \relax
2:11 & \begin{tabularx}{0.7\textwidth}{X} 在那日，必有許多國家歸附耶和華，作我的子民。我要住在你中間。」你就知道萬軍之耶和華差遣我到你那裡去。 \end{tabularx} \\ \\ \relax
2:12 & \begin{tabularx}{0.7\textwidth}{X} 耶和華必收回猶大，作為他聖地的產業，他必再度揀選耶路撒冷。 \end{tabularx} \\ \\ \relax
2:13 & \begin{tabularx}{0.7\textwidth}{X} 凡血肉之軀都當在耶和華面前靜默無聲，因為他從他的聖所奮起了。 \end{tabularx} \\ \\
[1ex]
\hline
\hline
\end{longtable}
今次研經的總題是神的榮耀與教會更新,最重要者是在個人的讀經生活.「成聖須用工夫」,多讀這卷書,藉聖靈的光照可領受更多.
\newline
\newline
神的榮耀與教會更新,是神的心意和目的,這也該是我們的心志.神的榮耀顯現在教會中,如同在以色列一樣.我們因此而復興,更新,並能見證,使多人歸主.這該是我們的異象和願望,而個人也因此得著更新.記得神的榮耀曾顯現在曠野的會幕,所羅門的聖殿,敬拜者的場合等,我們若信靠順服,祂的榮耀就必彰顯.主耶穌曾說:「你們若信,就必能看見神的榮耀」,求主幫助我們不足的信心,從祂的話語中,看見祂的榮耀並教會的更新.
\newline
\newline
(二)四角的異象(一18-21)
\newline
\newline
這段經文在古卷中該屬第二章,是第二個異象──神的審判.「我舉目觀看,見有四角.」角在聖經中常象徵力量,特別是指軍事力量.古時迫害以色列的外邦人,都是窮兵黷武之輩,戰爭繁複,迫害著神的百姓.先知哈巴谷懷疑神的公義是否彰顯,實有疑問.當時以色列面臨被擄,他眼見外邦人,在迫害神的選民,而感到不明白.但我們從撒迦利亞書中便可得著答案.
\newline
\newline
有人以為四角所指的,是東之亞摩利人,南之以東人,西之非利士人,及北之撒瑪利亞人.也有按但以理書解作巴比倫,瑪代波斯,希臘和羅馬四帝國.然而聖經既沒有明說,就要避免人意的解釋；尤以預言,有人常勉強逐一解釋它如何應驗,就會產生不少錯誤.我們必須注意經文的歷史背景和歷史意義,後者更是信仰的內容,在研究舊約時,尤為重要.這裡的四角是指異邦罪惡權勢的總合,即整個世界的屬世權勢.
\newline
\newline
「於是我就問與我說話的天使說:這是甚麼意思?他回答說:這是打散猶大,以色列和耶路撒冷的角.」(十八十九節)這些都是要迫害神選民的人,屬魔鬼的力量.
\newline
\newline
「耶和華又指四個匠人給我看.」(二十節)解經家中有以為四匠人,即四個彌賽亞──屬神的人戰勝屬世者.但若以經解經來說,按結二十一31,「善於殺滅的畜類人手中.」這一句話明顯是指那些外邦人,罪惡的人.其中「善於」一詞與工匠同義.因此匠人並非神的代表,乃是指外邦的勢力.換言之,罪惡互相爭競︴殘殺,或許神降罰以致罪人,惡人自相毀滅；這就是神對以色列人要說的話.那些陷害他們的,看來像很厲害,恐嚇使他們抬不起頭來；然而神是歷史的主宰,祂要藉歷史表達祂公義的審判.罪惡一定會敗壞,祂的兒女不應失望,心懷不平；因歷史本身就是審判.現代流行述語:「歷史的危機.」其中危機一詞在英語(crisis)原文上意即「審判」.歷史每件事都在說明神的審判,我們所見的是局部和片斷的,因此感到不清楚.
\newline
\newline
神允許外邦人陷害祂的兒女,是因為他們失敗,該受管教和公正的責罰.其中有一項原則──在神家中就該接受祂的審判:「審判是從神的家起首,跟著是外邦人.」按這原則就可解決先知哈巴谷的疑問.此點在阿摩司書中也有說明.今日的教會也該認識這真理,正如在啟示錄中,神審判世界以先,也先審判教會(啟二至三章).七個教會中,非拉鐵非和士每拿分別是傳福音和受苦的教會；他們不需要受審判,神且要特別看顧.信徒常以為不用受審判,因早已蒙恩得救,不在神的震怒之下.但聖經明說信徒要受審判,因為神的兒女在恩典中墮落,沒有長進,不受審判將必繼續墮落.為此信徒在世仍遭受患難,連罪惡的勢力都能打倒他們.
\newline
\newline
同時,我們也把傳福音看得太容易,以為主耶穌既已成就十字架的功勞,單是信便可以了.好像是給他面子,所以隨便的相信和隨便過生活.教會之所以要受審判,是因神的恩典太寶貴,而又不需代價的白白得著,我們並沒有珍惜它.以色列也是如此,至今他們經過數次歷史浩劫──耶路撒冷陷落,聖殿拆毀,民族敗亡.因著神的審判,教會要復興,除非澈底悔改,否則必受審判.
\newline
\newline
在這異象中,神特別教導撒迦利亞,也叫以色列人知道,神是信實的,外邦人也終必受審判.這是以色列信仰經歷所學習的功課──先受審判,然後要審判人.聖經很著重這個意思.我們將必與主作王,一同審判世界,這固然是將來的事,現在我們也要做審判的工作,教會傳福音,就是審判世界.傳福音不可隨便,不該注重人的興趣,以詼諧的方法作為吸引,而根本沒有使人為罪,為義,為審判自己責備自己.傳福音當以聖經為原則,要傳審判的信息.在教會歷史中的大復興,所傳信息的內容都是帶著審判的.一次愛德華先生講道時說,若罪人落在永生神的手中是極可怕的,很多人緊抓所坐的椅子,迫切的請他幫助,好像害怕立即就要下到地獄去似的.今日傳福音最大的錯誤是 : 只注重方法,不重視信息.然而要傳審判的信息實在不容易,因為我們要先受審判──受罪的責備和悔改；然後才懂得帶著眼淚,給別人傳審判的信息.這是第二個異象.
\newline
\newline
(三)準繩的異象
\newline
\newline
第三個異象是解釋第一章十六節,「所以耶和華如此說:現今我回到耶路撒冷,仍施憐憫!我的殿必重建其中,準繩必拉在耶路撒冷之上,這是萬軍之耶和華說的.」神的榮耀與教會更新,是神的應許和願望.我們戰兢的仰望祂,就必看見祂要在我們身上所作的工.
\newline
\newline
神要先知量度耶路撒冷,然後跑去告訴那少年人──少年是指屬靈感受的程度,意思是要先知看見聖殿復興的現象,讓他曉得自己當作的事.聖殿重建時是這樣,在主前的五八一年,耶路撒冷城陷落,聖殿被毀,以色列人在被擄的人群當中.雖然他們是神的選民,因罪也要受審判,這是神的公義；並且在神的計劃中,審判是從神的家開始.被擄至異邦的以色列人,深切痛悔,感到過去先知的話都是真實的.耶利米哀歌三章中,詳細描述以色列人遭受酷劫的悲慘境況,但又說相信他們並不永遠毀滅,人縱然失信,但神卻是信實的,神仍按祂救恩的計劃,祂不會後悔祂的揀選,以色列不會滅亡,因神恩典廣大.果然在五三八年,巴比倫敗亡,波斯王古列允許以色列人歸回,且在故土上重建.當他們回到耶路撒冷城,第一件事要在聖殿舊址築起祭壇；但十八年後聖殿還沒有被建成,先知哈該與撒迦利亞便起來,鼓勵他們.主前五二零年,哈該說:「你們要省察你們的行為,你們自己住在天花板的房屋,神的殿卻仍荒涼.」他的信息得著反應,他因此不再責備且鼓勵他們剛強,充滿希望,因為聖殿的榮耀將遠勝舊者.聖殿果然大約在主前五一五年建立起來,這殿通稱第二聖殿或所羅巴伯聖殿.這並非表示他們可以一勞永逸,在尼希米的時代(即主前四四四年),因為城牆還沒有修好,尼希米奉神的命令,經過百般的艱難,建造了耶路撒冷城的城牆.聖城的城牆是必須有的,我們也須注意這真理,我們要建立個人,家庭以及教會的城牆.要不然,教會將成為社會社團,不但失卻影響社會的力量,反受世俗勢力的入侵,教會便失去屬靈的功能和性質.家庭與個人也是如此,要讓神管理我們,就不要讓罪惡輕易進入我們的城牆.
\newline
\newline
很奇怪,在第二章卻說不用城牆──「像設有城牆的鄉村一樣」(亞二4).這似乎與以上所說的顯得有點矛盾,城牆是需要的,但並非最完美和理想,城牆實在不是神的理想.讓神使我們屬靈的眼睛擴大,逐漸長進,由有城牆至沒有,當時以色列人勢力脆弱,沒有能力抗拒外邦人和撒瑪利亞人,尼希米曾竭力地勸以色列人,要分別為聖及守安息日,不要與外邦人混在一起；然而在每次尼希米離開後,他們便失敗.竟然沒有人喜歡居住在耶路撒冷,在被擄以前,這京城是宗教,政治及經濟的中心；被擄以後,聖殿與城牆都重建了,卻沒有人居住.結果尼希米命令,每族十丁抽一,要他們居住其中.從這異象中可看出教會的更新,一如聖殿和城牆的重建,都是要注重質與量.撒迦利亞人看見耶路撒冷城滿了人居住,就像沒有城牆一樣.
\newline
\newline
「耶和華說,我要作耶路撒冷四圍的火城,並要作其中的榮耀.」神的榮耀,祂自己的同在,作為我們個人,家庭及教會最完美和理想的城牆.這是神的計劃.就歷史研究,尼希米記指出耶路撒冷城牆的構造是不錯的,但仍有過份狹窄的民族主義缺點.比方說,他們不願與有混雜血統的撒瑪利亞人來往,種族的歧見一直沒有打破,在主耶穌的行程中,祂必須經過撒瑪利亞而且在那裡傳道.教會建立以後,腓利也下到撒瑪利亞去,把傳統狹隘的觀念完全打破.主耶穌來,除掉這種敗壞的態度,他會不遺餘力地責備那些猶太教的人.今日我們雖然應當分別為聖,對於教會之間,教會與機構及宗派之間的城牆,我們需要排除這種不正常的現象.建造城牆,目的只在抵禦罪惡的勢力,非為自我孤立,明哲保身.但這不表示中國教會不需要宗派的存在；標榜屬靈而排斥宗派者會極錯誤.我們所不需要的,只是宗派主義那種自以為義的城牆.過份著重城牆就容易漠視他人的需要,特別今天在大城市的環境下,常有不願與人接觸,過於謹慎小心的危險,福音的見證,因此不能順利帶出來；真正的城牆是非人為的,有神的榮耀在裡面.神的保護才是最好的,祂不會任由外邦人欺侮祂眼中的瞳人.(申三十二10)我們的神乃是烈火,祂要成為我們的火牆；祂是公義的,也表明了祂的熱心；必定成就這事,祂看顧和恩待我們.火也是忌邪之意,又代表祂的愛,亮光和溫暖；這才是我們真正需要的城牆,同時,又有神的榮耀居住在其中.
\newline
\newline
我們讀到神的榮耀曾在以色列人的會幕中,在聖殿上,但在以西結書中,有一幅極其悲慘的圖畫(結九,十章),神的榮耀離開以色列.以迦博(參撒上四 21)因為拒絕神,(結九,十章)說神的榮耀,捨不得離開,從基路伯上升,慢慢移至聖所的門口,殿的東門,繼而停留在城東的山上,再三等候,直至以色列人不肯悔改時,才真的走了.今日神對我們也是一樣,我們軟弱了,神的靈會捨不得走,流著眼淚帶著歎息,一直等待著我們悔改.在結四十三章,我們看見神的榮耀再次回來,神不會永遠丟棄我們,我們悔改之時,是在歡樂歌唱以前,又當在耶和華面前靜默無聲,因為祂升起,從聖所出來了(亞二13),正如(代上六7)所說:「我的名永遠在那裡,我的眼,我的心在那裡.」神的名在我們身上,好像倒出來的香膏,沒有其他的名我們可以靠著得救.祂的名就是奇妙,祂與我們同在,看顧的眼和開懷的心,常在我們身上.我們應當靜默無聲等候神.
\newline
\newline


\newpage

\section{第51屆~港九培靈會~唐佑之~第三講}
\label{sec:675}
\textbf{約書亞的異象}
\newline
\newline
連結: \href{https://www.hkbibleconference.org/session-message/view/675}{\texttt{https://www.hkbibleconference.org/session-message/view/675}}
\newline
\newline
\hyperref[sec:674]{< < < PREV SERMON < < <}
~
\hyperref[sec:toc]{[返目錄]}
~
\hyperref[sec:676]{> > > NEXT SERMON > > >}
\newline
\newline
撒迦利亞書 3:1-3:10
\newline
\begin{longtable}{cl}
\hline
\hline
章節 & 經文 (和合本修訂版)\\
\hline
3:1 & \begin{tabularx}{0.7\textwidth}{X} 天使指給我看：約書亞大祭司站在耶和華的使者面前，撒但站在約書亞的右邊控告他。 \end{tabularx} \\ \\ \relax
3:2 & \begin{tabularx}{0.7\textwidth}{X} 耶和華向撒但說：「撒但哪，耶和華責備你！揀選耶路撒冷的耶和華責備你！這不是從火中抽出來的一根柴嗎？」 \end{tabularx} \\ \\ \relax
3:3 & \begin{tabularx}{0.7\textwidth}{X} 約書亞穿著污穢的衣服，站在那使者面前。 \end{tabularx} \\ \\ \relax
3:4 & \begin{tabularx}{0.7\textwidth}{X} 使者吩咐那些侍立在他面前的說：「脫去他污穢的衣服。」又對約書亞說：「你看，我使你的罪孽離開你，要給你穿上華美的衣服。」 \end{tabularx} \\ \\ \relax
3:5 & \begin{tabularx}{0.7\textwidth}{X} 我說：「要將潔淨的冠冕戴在他頭上。」他們就把潔淨的冠冕戴在他頭上，給他穿上華美的衣服，耶和華的使者在旁邊站立。 \end{tabularx} \\ \\ \relax
3:6 & \begin{tabularx}{0.7\textwidth}{X} 耶和華的使者告誡約書亞說， \end{tabularx} \\ \\ \relax
3:7 & \begin{tabularx}{0.7\textwidth}{X} 萬軍之耶和華如此說：「你若遵行我的道，謹守我的命令，就可以管理我的家，看守我的院宇；我也要使你在這些侍立的人中間來往。 \end{tabularx} \\ \\ \relax
3:8 & \begin{tabularx}{0.7\textwidth}{X} 約書亞大祭司啊，你和坐在你面前的同伴都當聽，因為他們是作預兆的：看哪，我必使我僕人大衛的苗裔長出。 \end{tabularx} \\ \\ \relax
3:9 & \begin{tabularx}{0.7\textwidth}{X} 看哪，這是我在約書亞面前所立的石頭，這一塊石頭上有七眼。看哪，我要親自雕刻這石頭，並在一日之間除掉這地的罪孽。這是萬軍之耶和華說的。 \end{tabularx} \\ \\ \relax
3:10 & \begin{tabularx}{0.7\textwidth}{X} 在那日，你們各人要請鄰舍坐在葡萄樹和無花果樹下。這是萬軍之耶和華說的。」 \end{tabularx} \\ \\
[1ex]
\hline
\hline
\end{longtable}
(亞一14-17)這一段經文中,第二,三個異象分別出現於第十五,十六節.今天我們要看的第四個異象,則在第十七節──「你要宣告說,萬軍之耶和華如此說,我的城邑必再豐盛發達,耶和華必再安慰錫安,揀選耶路撒冷.」
\newline
\newline
在前三個異象中,是要以色列人特別注意歷史的苦難,他們從苦難中看見神的恩惠.這是兩方面的──是公義與慈愛,也是懲罰與揀選.第四個異象,是要他們更看出自己,同樣教會也可以在神面前看見自己的光景──因此,審判必須從神的家起首.神是公義的,所以當教會失敗,軟弱和退後時,神一定要管教,但祂並不丟棄屬祂的人,祂要保護他們如同眼中的瞳人.罪惡必須消除,而最重要者,還是我們必須有屬靈的眼光,得見神的作為.不論在任何環境,不可失望,灰心,要知道神是奇妙的主；即使處身在低窪的山谷中,我們還須抬頭仰望神,因為祂必向我們施恩.在第三個異象的末了,先知要以色列人低頭靜默敬拜.敬拜是首要的,不然我們就沒有好的領袖和事奉.因此,在第四章中,特別看到祭司的工作,自被擄歸回,以色列人中有兩個領袖.所羅巴伯是政治的領袖,大祭司約書亞是宗教的領袖；此外還有先知哈該和撒迦利亞同期工作.
\newline
\newline
第三章記載的第四個異象中,主要的人物是大祭司約書亞.這異象好像是一個法庭,神是審判官,歷史每一件事,都與神的審判有直接的關係.歷史中的危機,都在表明神公義的審判.審判者是神,撒旦是原告,大祭司約書亞卻是被告.當約書亞被起訴時,神責備撒旦,這並非說撒旦的控告沒有憑據,身為大祭司的約書亞該是潔淨的,不但是內在,更需要外表的潔淨.不過經文告訴我們,他穿著污穢的衣服,他就沒有資格事奉神.這樣看來撒旦的控告是有憑據的,但神責備的卻是撒旦.對於認識撒旦,新約比舊約解釋較為清楚,舊約只有約伯記曾提到撒旦,其次本書也有提到；撒旦受責備是因為牠不尊重神的心意.牠的控告雖然是對的,約書亞實在是犯了不潔的罪,但神的心意,卻要安慰錫安,揀選耶路撒冷,所以神要責備撒旦.神有揀選的恩典,這非撒旦和一般人所能明白的.神的工作,要藉祂工人表達出來；所以尊重神的,就必須尊重神的心意和神的工作,特別是祂的工人,因為他們是神的代表.在這裡,我們要注意,以色列本是神所揀選的；在創世記中,特別強調神的揀選,而出埃及記對神揀選的恩典,有更透澈的說明(出十九5).神說:「你們若實在聽從我的話,遵守我的約,就要在萬民中作屬我的子民.」神揀選以色列的目的,並非因他們的民族優越,有輝煌的歷史和文化的遺產.他們本是一個弱小落後的民族,居於中東的一隅,過著半游牧的生活,實在沒有可值得揀選的地方,唯一的原因,是神的恩典.我們得蒙揀選,亦並非因我們有甚麼好處,完全是出於神的恩典.今天我們在神面前成為何等樣人,都是神恩典造成的.所以我們必須尊重神,並祂的工作,心意和祂的工人.若不然,我們必像撒旦一樣,受到責備.
\newline
\newline
大祭司約書亞獨自面對審判,並非因他罪大惡極,乃因他是全以色列會眾的代表；這是希伯來人的思想,稱為積聚人格.換言之,約書亞是神的僕人,他失敗,就要受到責罰和審判；因為審判是從神的家起首,他要代表全以色列受審判.神在公義中還充滿無限慈愛,神說約書亞是從火中抽出來的一根柴(也包括全以色列).神是烈火,故公義的審判,要臨到以色列.神是信實的,祂的恩典不改變,祂有無限的憐憫,祂是守約施慈愛的主,以色列像是「從火中抽出來的一根柴.」阿摩司書四章十一節也有同樣的說法.阿摩司是主前第八世紀以色列第一位寫作的先知,是四位重要先知之一(即阿摩司,何西阿,以賽亞及彌迦),神分別使用他們.身為南方人的阿摩司,是向北傳道,何西阿亦向本身的北方人傳道,而生於皇宮貴族的以賽亞,則向自己南方人傳道.相反的,彌迦是在鄉村長大的,他仍用不同的方法,把神的真理表達出來.阿摩司講說公義責備人,極其嚴厲.何西阿身為北方人,是慈愛的先知,心存愛同胞之心,流淚,嘆息地傳神的愛.以賽亞出身貴族,有特別的風度；彌迦卻是鄉下人,說話極率直.以賽亞重聖潔,彌迦總匯三者之內容.他說神要求我們行公義,好憐憫,存謙卑的心與神同行.我們讀聖經,須下功夫.神的啟示是一貫的,我們若能了解這四位先知的信息,就能明白其他先知所傳的,能以經解經了.「從火中抽出來的一根柴」這句話,可使我們聯想到傳公義的阿摩司.不過神在公義中仍有憐憫,如果人願意悔改,就可以獲得神的憐憫.審判中有一個極重要的信息是「餘數」或稱「餘民」.神實在捨不得滅絕所有的人,甚至是外邦人.神盼望人悔改,所以在先知書中,說出「餘數」的真理.這些人都是敬虔的,而且是切實悔改的以色列人,他們有真信心仰望神,其中更有外邦人的餘數.聖經中明說,若外邦人願意悔改,神也照樣讓他們蒙恩得救.因此,我們不該有狹窄的思想,以為以色列人必得救,倘若他們犯罪,照樣是會滅亡的.當他們不接受神揀選的愛時,就會在恩典中墮落,神只得把他們放棄,這是極重要的真理.
\newline
\newline
我常感覺中國教會,受許多不大正確的觀念影響.比方我們非常強調一種意義,舊約是律法時代,新約是恩典時代；時代的真理是正確的.但我不同意以上所說,因為舊約也有恩典,恩典與律法是相輔相承的.若我們相信神的啟示是合一的話,就不能以此劃分新舊約,當我們讀新約時,必會想起律法,就如約翰一章17 節所說,律法是藉摩西傳的,恩典和真理都是從耶穌基督來的,這是非常寶貴的真理,先律法,後恩典,沒有人能靠律法稱義.但舊約,我看恩典還是在律法之前,神先有揀選,是以色列的恩典,然後才叫他們愛神.如在出二十章,傳十誡以前,神說,我是耶和華你們的神,曾把你們從埃及為奴之地領出來.出埃及完全是神的恩典,以色列人並沒有作甚麼,是白白得著的.新約保羅在以弗所書前三章說恩典,神的拯救,基督把其中的牆拆毀,使我們能到父神的面前.然而在四章一節說:「你們既然蒙恩,行事為人就當與蒙召的恩相稱.」這是律法,律法中有恩典,恩典中亦有律法.保羅好像不著重律法,因為是針對猶太人以為行律法可以得救.對於西方的思想,就受到希臘的影響,而中國教會的傳統,則受宣教士的影響.猶太人具有希臘影響的西方思想,所以對新約比較舊約明白,東方人較容易了解舊約,但可惜跟隨了西方人的路線,好像要跟猶太人一樣；先進入了律法,然後又從律法中出來,實在沒有這個必要.在舊約中其實有無限的恩典,神白白把救恩給以色列人,從埃及拯救他們出來,完全是無條件和無須付代價的.我們若不重視恩典,恩典就變得太隨便了.我們雖然不好,神卻仍舊愛我們.因為各人都有軟弱,會隨便,馬虎,以為自己已蒙恩得救,就放縱自己,這是非常危險的.神的恩典每每賜給我們,也附帶有條件的；條件是我們要信靠順服,並且接受祂,不然,祂的恩典不會臨到我們.
\newline
\newline
中國教會也很喜歡講預定論.我們得救都是神所預定的,不能得救的人也是神預定的,這種思想完全錯誤.神預定所有人得救,不得救的,是他個人不願意得救.聖經早已說明,我們個人,家庭以及教會,在神面前都要負責任,這就是「餘數」的道理.神願意恩待每一個人,神愛世人,所以揀選以色列,給他們差傳的工作,作全世界的宣教士.在賽四十二章說以色列是神揀選的僕人,他們要成為眾民的中保,作外邦人的光.所有的人都在神的愛中.神有自己的方法,祂先揀選以色列,然後由他們把恩典帶出去.這也是教會的真理,神揀選我們,要我們把恩典帶出去；若我們失敗,神的審判就從祂的家起首.但不是每個人都失敗,我們會先受警戒,好叫我們悔改,作敬虔的人,我們好像從火中抽出來的一根柴一樣,悔改後,我們這些「餘數」,將成為復興的核心；教會的復興,就是從這些願意切實悔改的人起首.
\newline
\newline
當時,以色列要有大祭司,我們的大祭司是主耶穌.祂沒有犯罪,祂豈不是不會有約書亞的不潔問題嗎?事實上以色列全族,都是祭司的國度,這思想在新約也有,彼得前書二章九節說每個信徒都是君尊的祭司.在教會歷史上,曾把人分作聖品與俗品,這並非聖經的真理；明白真理的人,都知道每個信徒都是君尊的祭司.馬丁路德改教時,也重視這真理.但近年來,全世界教會都有一種趨向,在信徒運動聲中,以為平信徒可作任何職事,因此就漠視專職的傳道人,這一點實在是錯誤的.不懂敬拜,是今天教會最大的問題；缺少等候及個人靈修,自然不可能在團體中有好的敬拜.今日的牧師,傳道人及幹事因事務太忙,以致缺少了禱告等候的時間,甚至長執們也是如此.主日崇拜只是象徵式,在教會舉行過,來參加崇拜的人,也得不著甚麼.須知道祭司的工作是獻祭與奉獻,熱心的信徒只注重事奉,而缺乏真正的敬拜,詩班更鮮有做到帶領會眾讚美神,更甚者,是常聽道而不行道的,甚至有人人連聽道也不會.二百年前福音傳到中國,神興起不少奮興佈道家,中國教會雖未因重敬拜儀式,而忽略聖經話語,然而卻常停留在道理的開端.佈道家的信息雖好,他們卻沒有栽培的恩賜,以至不能幫助中國教會在神的話語上建立.一般人聽道的程度,是易的,喜歡故事,見證,笑話一類輕鬆愉快的；卻沒有在神話語上下功夫,這是教會的重大損失.到教會聽道的人太多,卻缺少行道的人；聽道的人,只喜歡聽外來的講者,而忽略了自己的牧者.反觀舊約祭司是教導神話語的人,他們多不解釋,只直接地告訴人,今天中國教會,極缺乏神的話語；講故事在講道中雖然有此需要,但也有取代神話語的危險.好的敬拜,要安靜等候,不單聽人講,也要自己去領受.祭司不隨便解釋神的話語,但常見人在佈道會中,說太多而不正確的話,忽略了神話語本身的功效.祭司且有醫治的責任,除了敬拜,奉獻和帶領別人,還要安慰幫助有需要的人.此外,祭司要把人帶到神面前,祭司必須清潔,正如約書亞要換上華美的衣服,戴上潔淨的冠冕,準備作神的工.今天我們有很多重要的責任,使命等待著,但必須先潔淨,而且要有祭司型的屬靈經驗.「萬君之耶和華如此說,你若遵行我的道,謹守我的命令,你就可以管理我的家,看守我的院宇,我也要使你在這些站立的人中間來往.」要管理神的家,必須先管理自己的家,熱心的信徒,卻沒有使自己的家榮耀神.我們需要先自潔,然後才能管理神的家,且在神的使者面前來往.
\newline
\newline
有三點我們要注意的:
\newline
\newline
第一,「我必使我僕人大衛的苗裔發出」,苗裔即樹枝,不像柴被砍下,大衛的苗裔是指彌賽亞,但教會是彌賽亞的團體,這說明我們有君尊的身份.
\newline
\newline
第二,「看哪,我在約書亞面前所立的石頭,在一塊石頭上有七眼.」又說要雕刻這石頭,石頭像摩西說的,耶和華是磐石,或彼得所說的,我們都是活石.七眼,表明神的看顧,祂是無所不知,奇妙的神,祂把智慧賜給我們.
\newline
\newline
最後,「你們要請鄰舍坐在葡萄樹和無花果樹下.」意味著不再有戰爭苦難,我們都在永久的平安中.現在還有許多人沒有在平安中,這是我們的失敗和忽略,因為仍穿著污穢祭司的衣袍.神的工作已成就,把罪人提升至義人,祂使我們因信稱義,但請注意,千萬不可在恩典中墮落,求神興起我們,管理神的家,在神使者中蒙恩更深.
\newline
\newline


\newpage

\section{第51屆~港九培靈會~唐佑之~第四講}
\label{sec:676}
\textbf{教會更新的果效}
\newline
\newline
連結: \href{https://www.hkbibleconference.org/session-message/view/676}{\texttt{https://www.hkbibleconference.org/session-message/view/676}}
\newline
\newline
\hyperref[sec:675]{< < < PREV SERMON < < <}
~
\hyperref[sec:toc]{[返目錄]}
~
\hyperref[sec:677]{> > > NEXT SERMON > > >}
\newline
\newline
撒迦利亞書 4:1-4:2
\newline
\begin{longtable}{cl}
\hline
\hline
章節 & 經文 (和合本修訂版)\\
\hline
4:1 & \begin{tabularx}{0.7\textwidth}{X} 那與我說話的天使又來叫醒我，好像人睡覺時被喚醒一樣。 \end{tabularx} \\ \\ \relax
4:2 & \begin{tabularx}{0.7\textwidth}{X} 他問我：「你看見甚麼？」我說：「我看見了，看哪，有一個純金的燈臺，頂上有燈座，其上有七盞燈，每盞燈的上頭有七根管子； \end{tabularx} \\ \\
[1ex]
\hline
\hline
\end{longtable}
今天我們要思想撒迦利亞所看見的第五個異象.第四章第一節說:「那與我說話的天使又來叫醒我,好像人睡覺被喚醒一樣.」看完第四個異象之後,撒迦利亞處在一個非常深切的屬靈感受中,他的生命溶化在異象裡面,可是神還有進一步的啟示,還有更重要的真理.於是便叫醒他,問他說:「你看見了甚麼?」記得耶穌曾經有一次問一個瞎眼的人說:「你看見甚麼?」那瞎子因為得到耶穌的醫治,非常興奮,忙說:「我看見了!」只是他所看見的仍不清楚,他看見人好像樹木一樣,而且行走,他雖然開了眼,但所看見的並不正確.今天在神的教會內,你和我也好像瞎眼的一樣,看得不清楚,不準確,不澈底.我們真要求問主說:「這到底是甚麼意思呢?」我們不要自以為是,自圓其說,乃要像撒迦利亞一樣,求主給我們解釋明白.
\newline
\newline
撒迦利亞所看見的這異象,和前面的異象有非常密切的關係,這章聖經的主要人物是所羅巴伯,前一個異象即第三章內,主要人物則為約書亞,這二人同是被擄歸回的主要人物.大祭司約書亞是當時屬靈的領袖,而所羅巴伯則是行政的領袖,是個省長.在第四章末了的地方,天使問撒迦利亞說:「你知不知道這兩根橄欖枝在兩個流出金色油的金嘴旁邊,是甚麼意思?」撒迦利亞回答說不知道,天使便告訴他,這是兩個受膏者,站在普天下主的旁邊.這兩位受膏者,就是大祭司約書亞和行政領袖所羅巴伯,為甚麼兩人這麼重要呢?原來約書亞在第四個異象中,是代表以色列整個的民族,因為以色列是祭司的國度,神要藉他們把福音的計劃帶給普世,但以色列人失敗了,他們祭司的外袍污穢了,所以神責備他們.雖然如此,但神並沒有丟棄他們,神要潔淨大祭司約書亞,要他誠心實意的做大祭司的工作,在潔淨之後再蒙差遣.但這個異象還未講完,第五個異象便是更進一步的說明,其中出現的人物,雖然只是行政首長所羅巴伯,但他卻是非常重要.
\newline
\newline
所羅巴伯的重要,是接續第四章所講到大衛的苗裔,大衛的子孫,這和彌賽亞的預言有非常密切的關係.神曾和大衛立約,應許他的家有永遠的王位,可是自所羅門以後,國家便分裂為二；十個支派在北國,由耶羅波安統治.他們和大衛的系統完全沒有關係,所以他們國內政治混亂,而且速速敗亡.但南國的兩個支派,就是猶大和便雅憫的支派就不同了,神始終保存猶大支派大衛家裡面的人；甚至到被擄時期,還是大衛家的人.後來到波斯王讓以色列歸回的時候,設巴薩仍然是大衛家的人,至於所羅巴伯,他是撒拉鐵的兒子,撒拉鐵是約雅敬的孫,換句話說,所羅巴伯就是約雅敬的曾孫,還是大衛家的人,一直到彌賽亞,仍然是屬大衛家的.由此讓我們看見神的信實,和祂奇妙的保守.彌賽亞必須是君王,這是先知以賽亞的預言中所非常著重的.彌賽亞既然是君王,是祭司,有榮耀的身份,我們接受了祂的生命,相信祂的人,也同樣有君尊的身份,也同樣是君王,是祭司,這是第四和第五個異象的主要意思.
\newline
\newline
跟著第四章第二節:「他問我說:你看見了甚麼?我說:我看見了一個純金的燈台,頂上有燈盞,燈台上有七盞燈,每盞有七個管子.」燈台是在會幕和聖殿內的,一個純金的燈台,是純粹而且尊貴的,教會是燈台,它應該有聖潔和尊貴的地位；如果我們查考出埃及記,便知道金燈台是用一塊金子造成的,這表示基督的身體只有一個,聖靈只有一位,直正的教會是合一的；有一樣的信仰,一樣的見證.當然,金子要變成純金,必須經過錘打,那就是說,教會也必須經歷許多的操練.
\newline
\newline
基督徒的生活,是一種操練的生活.我們要在敬拜上操練,敬拜的真義,就是體會神的同在,不是單在外表儀式的聚集,乃要體會神與我們同在,從而起敬畏的心,這經驗是我們一生要學習的功課,我們要在主面前預備,等候,讓神向我們顯現,引導我們,這是一種不斷的操練.另外,我們要在禱告上操練,我們是祭司,是負責清理燈台的,燈台清理乾淨,加上油之後,才能發光,照亮他人.如果我們不好好學習禱告,操練禱告,教會怎麼會有生氣呢?因此我們務必在敬拜和靈修的生活上操練自己.還有,我們的事奉也必須經過操練,這事奉並不單指工作上的事奉,還有心靈的事奉,敬拜的事奉,生活的事奉等等,這一切都需要我們不斷去操練.很多時候我們說事奉神,其實只是為了滿足自己的虛榮,表現自己的才幹而已,並不是真的服事神.教會裡面如果產生問題,往往是那些最熱心的弟兄姊妹,有事奉的弟兄姊妹產生的,原因是各人以自我為中心,缺少禱告,缺少靈修,缺少弟兄姊妹的相交而形成的.也有不少信徒,在主日沒有敬拜,他們有太多的事奉,以致沒有在神面前安靜等候.美國一位很出名的哲學家曾說,今日教會最大的問題,就是後排的基督徒太多了.甚麼是後排的基督徒呢?就是那些袖手旁觀的信徒,是一位旁觀者,不參與任何工作,甚至可以遲到早退,這是後排的基督徒,不過,我也發現很多有事奉的人坐在後面,因為他們要做招待或別樣的工作,所以坐在後排,但他們沒有多的領受,沒有好好的敬拜,奉獻,結果完了聚會依然故我,在屬靈方面沒有長進,最熱心的人也不長進,教會的毛病便是在此.我們要好好做祭司的工作,整理燈台,燈台本身並不能發光,但這裡說有七盞燈,七個管子,加起來便有四十九個,可以讓油能正常不斷的供應.我們需要在神面前領受恩典,能力,才可以發光.
\newline
\newline
撒迦利亞並不明白這異象的意思,萬軍之耶和華便說:「不是倚靠勢力,不是倚靠才能,乃是倚靠我的靈,方能成事.」根據撒迦利亞書本來的字眼,勢力是指軍事和政治的力量,我們可以簡稱它為「人力」,即人的力量；才能原來是指財力,物力,那就是說,不是倚靠人力,不是倚靠財力,乃是倚靠靈力,這是個很重要的真理.保羅也曾說,所有的事都本於神,依靠祂而歸於祂,如果我們的工作不是出於神,神不會負責,即使我們很熱心,很努力的做,我們所建造的不是靈宮,只是巴別塔而已.
\newline
\newline
今天我們看見差傳的工作很興旺,但我們不要犯上西國教會宣教工作的錯誤,他們把本末倒置了.我們做神的工作,必須有三方面的力量,就是靈力,人力和財力,但他們卻把財力放在第一,開宣教會議的時候,最主要是大家捐錢,然後便找人,找到之後便發展工作,到宣教士或工作人員寫信回來作報告時,才說要禱告,這是本末倒置了.錢不是不重要,但不是最重要的,我們有了錢,但不一定有合適的宣教士,有合適的宣教士,也不一定能有效的做宣教工作；我們應把聖靈放在第一,惟有耶和華的靈,方能成事；有了靈,便一定有人,因為聖靈會差遣合適的工人,至於金錢,在神眼中實在是微乎其微的.同樣,我們若單在宣教工作上捐錢,這也不表明我們是真正關心宣教的工作,我們乃要在神面前等候,讓神來差遣我們.很多時候我們犯了一個很嚴重的錯誤,好像亞伯拉罕的兒子以撒一樣,不曉得自己就是那燔祭的羊羔.我們常禱告求神差派工人,但我們自己卻不準備回應神的呼召,這樣的禱告是空洞的.我們是祭司,應該隨時蒙潔淨,隨時接受操練,又隨時蒙差遣.
\newline
\newline
跟著的聖經說,撒迦利亞看見有一塊石頭在殿頂,這和第四個異象中冠冕栽在人的頭上的象徵很相似,這石頭不單像冠冕栽在殿的頂上,而且要成為地的根基,讓神的恩典藉著這些工作得以完成.誰能藐視這件事為小呢?雖然開始時不為人注意,但聖靈的工作一直在發展,而且非常偉大.只有當我們真正做主的工作時,神的恩典就帶出了.今天,教會是金燈台,我們需要神的靈不斷供應我們,使我們能在夜深黑暗的時代中,為神大發光明.還有,為甚麼說約書亞和所羅巴伯二人是受膏者呢?我們說受膏者的時候,不是指彌賽亞嗎?不錯,彌賽亞意思即受膏者,但我們知道大祭司是要經過按立的,就好像有首詩篇形容把油倒在亞倫頭上,一直流到他的鬍鬚,衣襟,油是表示神的能力,我們也必須這樣.另外,君王也是要受膏的,所以這裡提到受膏者,是表明神把權授與他們,要他們做神的工作.今天,我們是君尊的祭司,也是君王的身份,神的靈澆灌我們身上,所以教會應該注重聖靈的工作,讓神的靈在我們身上發生功效.
\newline
\newline


\newpage

\section{第51屆~港九培靈會~唐佑之~第五講}
\label{sec:677}
\textbf{教會應走的路線}
\newline
\newline
連結: \href{https://www.hkbibleconference.org/session-message/view/677}{\texttt{https://www.hkbibleconference.org/session-message/view/677}}
\newline
\newline
\hyperref[sec:676]{< < < PREV SERMON < < <}
~
\hyperref[sec:toc]{[返目錄]}
~
\hyperref[sec:678]{> > > NEXT SERMON > > >}
\newline
\newline
撒迦利亞書 6:9-6:13
\newline
\begin{longtable}{cl}
\hline
\hline
章節 & 經文 (和合本修訂版)\\
\hline
6:9 & \begin{tabularx}{0.7\textwidth}{X} 耶和華的話臨到我，說： \end{tabularx} \\ \\ \relax
6:10 & \begin{tabularx}{0.7\textwidth}{X} 「你要拿從巴比倫歸來的被擄之人黑玳、多比雅、耶大雅所獻的，當日就要進到西番雅的兒子約西亞的家裡， \end{tabularx} \\ \\ \relax
6:11 & \begin{tabularx}{0.7\textwidth}{X} 拿這金銀做冠冕，戴在約撒答的兒子約書亞大祭司的頭上； \end{tabularx} \\ \\ \relax
6:12 & \begin{tabularx}{0.7\textwidth}{X} 對他說，萬軍之耶和華如此說：『看哪，那名稱為大衛苗裔的，要在本處生長，並要建造耶和華的殿。 \end{tabularx} \\ \\ \relax
6:13 & \begin{tabularx}{0.7\textwidth}{X} 就是他，要建造耶和華的殿，他要承受尊榮，坐在位上掌王權；又有一位祭司坐在自己的位上，兩職之間籌劃和平。 \end{tabularx} \\ \\
[1ex]
\hline
\hline
\end{longtable}
撒迦利亞書六章九至十三節:「耶和華的話臨到我說,你要從被擄之人中取黑玳︴多比雅,耶大雅的金銀,這三人是從巴比倫來到西番雅的兒子約西亞的家裡.當日你要進他的家,取這金銀作冠冕,戴在約撒答的兒子大祭司約書亞的頭上,對他說,萬軍之耶和華如此說:看哪,那名稱為大衛苗裔的,他要在本處長起來,並要建造耶和華的殿.他要建造耶和華的殿,並擔負尊榮,坐在位上,掌王權,又必在位上作祭司,使兩職之間籌定和平.」
\newline
\newline
這裡提到在被擄的人中找三個人,他們把金銀帶來,帶到約書亞的家裡,約書亞大概是在巴勒斯坦的地方.當猶大國敗亡的時候,整個以色列國的民族,包括北國的以色列和南國的猶大,都分散在各地,有的是在被擄的地方,有的則流落在巴勒斯坦,神要叫他們兩方面的人,都聯合起來合而為一.怎樣合一呢?就是叫被擄的人回來,但不可空手回來,他們在被擄的地方,雖然身份是奴隸,但是在巴比倫做生意,經濟環境也不錯,所以這三個代表被擄的人,將奉獻的金銀帶到耶路撒冷來,交給約書亞；然後用這些金銀作冠冕,不是作一個,而是作好幾個冠冕.關於冠冕的解釋,請各位再留意一件事,就是神的榮耀與教會更新,我們在神面前一定要有好的奉獻,但奉獻的目的是為甚麼呢?是不是為教會的經費,傳道人的生活,傳福音的事工呢?這當然是,但最主要的,是為彰顯神的榮耀.在這裡,冠冕這字是複數的,而不是單數的.怎樣可以做多個冠冕呢?而這些冠冕又戴在誰的頭上呢?原來不是戴在幾個人的頭上,乃是戴在一個人的頭上,他就是大祭司約書亞,因為約書亞代表整個以色列民族.那麼為甚麼是多數呢?原來這是希伯來語文文法的方式,多數是表示華麗,偉大的意思.有時聖經記載神說:「我們」,這就是三位一體的道理,大部份是指三位一體,但有時不一定是.我們必須真正明白聖經用字的意思.
\newline
\newline
約書亞為甚麼要戴冠冕呢?這是表明他祭司的工作是榮耀的工作.聖經很清楚的告訴我們:「他要建造耶和華的殿,並擔負尊榮,坐在位上,掌王權,又必在位上作祭司,使兩職之間籌定和平.」甚麼是兩職呢?就是君王和祭司的職份.以色列民族是君尊的祭司,我們信徒也有君尊的職份,我們要尊重祭司的職份,好好的敬拜神,並要教導律法,做醫治的工作,把人帶到神的面前.傳福音是祭司的工作,差遣也是祭司的工作.只可惜我們常對祭司有錯誤的觀念,以為單有外表的敬拜便足夠了,舊約時代也的的確確是這樣.我們知道,舊約時代有些祭司,單單注重外表的形式,沒有注重屬靈的光景,所以神便興起先知,特別是作早期寫作的先知,如阿摩司,何西亞,以賽亞,彌迦等,這四個先知是主前第八世紀的先知,他們非常反對祭司型的宗教,他們說如果單有外表形式的宗教是沒有用的.我們可以注意這些先知的信息,阿摩司說,我連先知都不算,因為我本來就是耕田的農人；最要緊是得到神真正的差遣,得到神的命令,將神的話傳出去.在先知裡面,大概祭司型的宗教都帶著反對的態度,甚至在耶利米的時候,當猶大王約西亞改變宗教,修理聖殿之時,很多人去敬拜神,外表看來好像很好,但耶利米覺得,如果他們在內心沒有悔改,外表的敬拜還是沒有用的.耶利米書第七章開始,便記載耶利米站在聖殿的門口,傳講耶和華的話,他說:「你們不要倚靠虛謊的話,這是耶和華的殿,是耶和華的殿,是耶和華的殿.」他他提醒百姓,不要以為來到耶和華的殿便可以蒙福,表面看來好像是這樣,但如果他們內心沒有悔改的話,照樣會滅亡,果然他們是滅亡了,這是多悲慘的收場.可是,當我們讀以西結書的時候,便發現他重新注意祭司型的宗教.為甚麼呢?因為當時他們沒有聖殿,沒有敬拜.是不是他們不應該敬拜,不需要聖殿,不需要儀式呢?不是,他們實在需要,但是他們不能單單注重外表,還要注重內心,這才是最重要的.所以,被擄歸回之後,祭司的宗教,有先知信息,而先知的信息也有祭司的重點,這是今天我們華人教會應走的路線.
\newline
\newline
有時我們太注重先知的信息,我們喜歡聽道,希望讀經之時有亮光,而不大注重形式,其實形式是應該注重的.我們知道,教會是無形的教會,是普世的教會,我們不要分宗派,中外,地點,時間,凡是蒙恩得救者,無論是在哪裡,是屬於哪一個團體,在哪一間教會,都是一樣的,因為身體只有一個,教會只有一個,聖靈只有一位；不過,在另一方面,我們也要注重地方教會的真理,教會一定是有地方性的,是有形有體的,因為教會是基督的身體.為甚麼基督要來到世間,道成肉身呢?因為要讓人具體的,真實的看見神,正如約翰一書所說,這生命之道是我們親眼看過,親手摸過,是親身經歷過的.現在主耶穌已經復活升天,他是不是就沒有身體呢?不是,他還有身體,就是祂的教會,教會是繼續的道成肉身,有形有體,很真實的把神的道表現出來.因此,我們今天便要注重形式,組織,系統,好像舊約時代祭司所做的一樣,我們也要照著規矩一樣一樣的去做.
\newline
\newline
有些人說,我們甚麼都不要,禮拜堂不是教會,我們不需要組織,而宗派更是罪惡.其實我覺得,今天我們中國的教會需要宗派,沒有宗派是注重人,不是注重神的工作.教會如果有屬靈的人在帶領,那就很好；但如果這人一旦過去了,或者他所作的,令人失望的話,教會便分裂了；但宗派的情形就不同,神的工人會過去,但是神的工作不會過去.我們做神的工作,不單要想到今天,還要想到明天,如果單有幾個熱心愛主的信徒聚集在一起,那是不夠的,地上的教會,是繼續把神的道彰顯,成為肉身的,宗派曾經成為歷史上的必須,有它的價值.可是,近幾十年來,中國教會受了一種影響,認為宗派不屬靈,組織不屬靈,形式不屬靈,到底是否真的不屬靈呢?可能是不屬靈,但如果沒有這些,是否便屬靈呢?也不一定.我們聖餐擘餅的時候,為甚麼預備有餅和杯呢?這不是形式嗎?我們為甚麼不單閉上眼睛,有人領導我們吃餅和喝杯,而不預備餅和杯呢?我們需要形式來幫助我們屬靈的感受.當然,如果我們單有外面形式和象徵是絕對不夠的,如果單注意外表的祭司組織是絕對錯誤的,但如果我們只講屬靈而不把神的工作做好,同樣也是錯誤.華人教會沒有充份的長進,我看這也是原因之一.我們需要先知的話,也要有祭司的重點,我們的教會要平衡,不要偏離左右.教會的路是甚麼呢?是十字架的路,有人說,我們從使徒行傳去看教會的路,又有人說,我們應該回到耶路撒冷的路,也有人主張教會走回安提阿的路,但這不是回復,因為神在歷史中一直往前走,直到主耶穌基督再來的.
\newline
\newline


\newpage

\section{第51屆~港九培靈會~唐佑之~第六講}
\label{sec:678}
\textbf{教會復興的真義}
\newline
\newline
連結: \href{https://www.hkbibleconference.org/session-message/view/678}{\texttt{https://www.hkbibleconference.org/session-message/view/678}}
\newline
\newline
\hyperref[sec:677]{< < < PREV SERMON < < <}
~
\hyperref[sec:toc]{[返目錄]}
~
\hyperref[sec:679]{> > > NEXT SERMON > > >}
\newline
\newline
撒迦利亞書 7:1-7:7
\newline
\begin{longtable}{cl}
\hline
\hline
章節 & 經文 (和合本修訂版)\\
\hline
7:1 & \begin{tabularx}{0.7\textwidth}{X} 大流士王第四年九月，就是基斯流月初四，耶和華的話臨到撒迦利亞。 \end{tabularx} \\ \\ \relax
7:2 & \begin{tabularx}{0.7\textwidth}{X} 那時伯特利人已經差遣沙利色和利堅‧米勒，並他們的人，去懇求耶和華的恩， \end{tabularx} \\ \\ \relax
7:3 & \begin{tabularx}{0.7\textwidth}{X} 問萬軍之耶和華殿中的祭司，又問先知：「我當如歷年以來所行，在五月哭泣齋戒嗎？」 \end{tabularx} \\ \\ \relax
7:4 & \begin{tabularx}{0.7\textwidth}{X} 萬軍之耶和華的話臨到我，說： \end{tabularx} \\ \\ \relax
7:5 & \begin{tabularx}{0.7\textwidth}{X} 「你要向這地全體百姓和祭司說：『你們這七十年來，在五月、七月禁食悲哀，豈是真的向我禁食嗎？ \end{tabularx} \\ \\ \relax
7:6 & \begin{tabularx}{0.7\textwidth}{X} 你們吃喝，不是為自己吃，為自己喝嗎？ \end{tabularx} \\ \\ \relax
7:7 & \begin{tabularx}{0.7\textwidth}{X} 當耶路撒冷和四圍的城鎮有人居住，享繁榮，尼革夫和謝非拉也有人居住的時候，耶和華藉從前的先知所宣告的，你們不當聽嗎？』」 \end{tabularx} \\ \\
[1ex]
\hline
\hline
\end{longtable}
撒迦利亞書七章一到七節和八章十八到廿三節中,記載從伯特利有人來問,怎樣才是正常的屬靈生活,是不是需要禁食,禱告呢?我們需不需要在五月和七月禁食呢?在摩西和律法書裡面,七月是個重要的月份,七月初一是吹角節.吹角節以後,各人要集中耶路撒冷守節；七月初十是贖罪日,各人要悔改,禁食禱告.大祭司要進到至聖所,為自己和眾百姓贖罪,之後就是過住棚節；特別紀念他們在曠野寄居的日子,怎樣得神眷祐.只可惜當時聖殿已被毀,他們被擄至外邦地方,想到自己民族遭遇空前的浩劫,便更覺需要禁食禱告,哭泣痛悔.除七月之外,五月也要禁食,因為五月是記念聖殿被毀的日子.另外,四月是耶路撒冷被巴比倫軍隊攻破；十月是整個耶路撒冷被巴比倫圍困,所以他們在四,五,七,十月都要禁食.但現在他們卻問,到底需不需要禁食,因為禁食好像是禮儀,是祭司型的屬靈生活,到底現在還有沒有這個需要呢?
\newline
\newline
撒迦利亞便回答他們說 : 禮儀固然需要,但更重要是先知的話,除了儀式外,我們還要有真正悔改的心.這樣,禁食的日子就必成為猶大家歡樂的日子,因為經歷痛悔之後,神就賜我們歡樂.就如歷代志記載希西家王在獻燔祭以後,就有非常歡樂的現象.如果我們生活中沒有歡樂,那是因為我們沒有獻上燔祭.舊約內有三個祭是我們每天都必須獻的,就是燔祭,把自己完全的擺上；平安祭,和人,和神有平安的關係；素祭,就是感恩的祭.神不單要求我們有形式上的禁食,獻祭,更要緊是內心真誠的悔改,這樣才會帶來歡樂的光景.同樣,個人的復興,教會的復興,都是悔改而來的；當人願意切實悔改之時,神的福氣便馬上臨到他們.威爾斯在大復興之前,社會情況非常低落,酒吧門庭若市,教會則冷冷清清,但是自從有幾個人在神面前願意切實認罪,復興就來了,教會突然多了許多人聚會,甚至酒吧也要關門,因為所有人都到禮拜堂；法庭,警察,監獄也不需要,整個社會好像脫胎換骨一樣翻新過來.人主動來尋找基督,就像當日人主動去找猶太人,說:「我聽見神與你們同在了.」他們要和猶太人一起去敬拜神.今天,香港教會有沒有這樣復興的光景呢?
\newline
\newline
跟著我們來看第八章,特別注意七,八兩節,在這裡,我們看見神招聚祂的兒女,祂要拯救各處的人,成就祂恩惠的約,這約就是神要作他們的神,他們要作神的子民.這是教會更新的異象.十四,十五節說:「萬軍之耶和華如此說:你們列祖惹我發怒的時候,我怎樣定意降禍,並不後悔.現在我定意施恩與耶路撒冷,和猶大家,你們不要懼怕.」耶和華是一位定意的神,祂有一定的旨意,絕不會出爾反爾；祂是信實的,祂說施恩給人,便施恩給人,因為祂是慈愛的神,但祂也是公義的神.祂定意刑罰犯罪的人,決不後悔.可是,有的時候,我們也看見神的後悔,就如神在挪亞的世代中後悔造人,當尼尼微城的人悔改時,神便後悔不降災禍與他們.到底神會不會後悔呢?祂又為甚麼要後悔呢?其實這只是用字方面的問題,神的後悔不是意志的,乃是情緒的,那就是說,神絕不會改變祂的意思,也不會懊悔祂所做的事,只是在情緒方面,神真的覺得後悔了,因為祂內心難過憂傷,如同父母在對子女失望時講出後悔的話.
\newline
\newline
不過,神的恩典也是帶有條件的,從一方面看,神的恩典是白白的,無條件的賜予,不須我們付任何代價；但另一方面,神的恩典也是有條件的,如果我們不悔改,神的福氣便不能臨到我們.神定意降禍,定意施恩,一切的權都在祂手中,絕不能更改.我們常說,禱告能改變一切.不錯,禱告是有改變的力量,但禱告不是改變事情,只是改變人,神的旨意定了,不能更改,禱告只能改變人；就像主耶穌在客西馬尼園的禱告一樣,不要照我的意思,只要照神的意思.
\newline
\newline
上文第六章講到,有三個人帶著金銀從巴比倫來,他們要把這些金銀,就是很多人的奉獻,交給西番雅的兒子賢,或稱為約西亞.這三人是被擄去巴比倫的人,而約西亞則在本地,沒有被擄,他們都一同齊心合意的興旺神的工作,要把金銀做成冠冕,戴在大祭司約書亞的頭上.約書亞是以色列眾民的代表,所有的以色列人都是祭司.之後,這冠冕要歸約西亞和帶來奉獻的那三個人,意思是將尊榮歸給奉獻者.並要把這冠冕放在聖殿裡作紀念,讓大家在敬拜時常看見,好提醒他們一個以色列的敬拜者有怎樣的身份,使他們不單想到自己的尊貴,更想到彌賽亞君王,因為只有彌賽亞才配稱為尊貴榮耀的王.同樣,基督是教會的頭,尊榮都是屬於祂的,當我們把尊榮歸給祂時,我們也得著尊榮.
\newline
\newline
跟著我們看八章十二節,那裡說:「我要使這剩餘的民,享受這一切的福.」這個剩餘的道理非常重要,前面第三章提過,約書亞好像是從火中抽出的一根柴,餘數就是指少數敬畏神,願意切實悔改的人,這少數人,或稱為餘數,便成為以色列民復興的核心.這好比一個倒置的三角形,上面是很寬的,但慢慢愈來愈窄,最後只剩下一點.神本來是要拯救世界上所有人類的,但因為人墮落失敗,所以神就無法廣泛施恩給普世的人,祂只揀選了以色列民族；但以色列民族中,也不是所有都是虔誠的,結果就只剩得一小撮的餘數.而最後唯一的餘數,就是主耶穌基督；這冠冕就是要留給這餘數的餘數,因為只有祂是萬王之王,萬主之主.第九章九至十二節便形容這位王怎樣騎驢進入耶路撒冷,從四福音記載中,我們知道祂此行是要為人類的緣故,上十架犧牲,擔當世人的罪.祂的死,就像一粒麥子埋在地裡死了,然後結出許多子粒來,這個本來是三角形的頂點,現在又逐漸擴展,而且愈來愈寬,神把得救的人數,天天加給教會.
\newline
\newline
耶穌是尊貴的王,祂有尊貴的冠冕,權柄,可是祂到世間來的時候,卻成為一個最卑微,最貧窮的人.祂自己也明說,祂不是要接受人的服事,乃是要服事人.祂在上十架的前一夜還替門徒洗腳,服事他們,把愛的榜樣留給他們.耶穌既然這樣,我們豈不也要這樣嗎?我們是神的兒女,應該有尊貴感,可是我們在神面前也要存心卑微,學效耶穌一樣,作個服事人的僕人.耶穌連門徒的腳都洗了,我們豈不應當照樣行嗎?今天,只有基督是教會的主,我們都是僕人,牧師傳道也不是主人,只是神的僕人,換句話說,教會是一個僕人的團體,我們要在神面前盡心事奉.以賽亞先知講及彌賽亞預言的時候,前半部介紹祂是位很尊貴的君王,從大衛家出來的；但後半部卻介紹祂是個受苦的僕人,有幾首著名的詩都是從這個角度來描寫的.四十二章描寫這僕人有神的靈與祂同在,祂要做傳道的工作,壓傷的蘆葦,祂不折斷；將殘的燈火,祂不吹滅,祂要將神的公理傳於外邦.祂的職分是先知的職分.四十九章二節形容這僕人好像箭藏在箭袋裡,要隱藏起來,但當他出來的時候卻很鋒利,為甚麼呢?因為祂要成為耶和華的口,成為神的發言人,祂的工作是先知的工作.第三首僕人之詩,形容祂受逼迫,遭鞭打,但是祂毫不在乎,因為祂要執行神的使命,所以祂求神給祂一個受教者的舌頭,讓祂用言語扶助軟弱的人.在這裡,祂的工作是先知和祭司合併的工作.至於第四首僕人之詩,是從五十二章末了開始,直至五十三章,先知嘆息說,我們所傳的,有誰信呢?耶和華的膀臂,向誰顯露呢?這位僕人看見很多人像羊走迷了路,心裡難過非常.祂是一位常經憂患,多受痛苦的主,但是祂受的鞭傷,使我們得醫治,祂受的痛苦,使我們得平安.祂是祭司,又是代罪的羔羊,祂像羊羔被牽到宰殺之地.換句話說,祂不但是祭司,也是祭品,把自己完完全全,毫無保留的獻給神,於是成就了得救的根源.所以耶穌不單說要服事人,他還說要捨命作多人的贖價.祂在設立聖餐的時候,拿起杯來說:「這杯是用我血所立的新約.」約是用血來立的,耶穌為我們成就了恩惠的約,我們若不吃祂的肉,喝祂的血,就與祂的生命無份.我們的王已經來過,這是歷史的事實,但是祂回去了,不過祂還要回來；歷史上每一件事實都在告訴我們,祂快要再來了,讓我們預備好自己,迎接這榮耀的王.
\newline
\newline


\newpage

\section{第51屆~港九培靈會~唐佑之~第七講}
\label{sec:679}
\textbf{好牧人}
\newline
\newline
連結: \href{https://www.hkbibleconference.org/session-message/view/679}{\texttt{https://www.hkbibleconference.org/session-message/view/679}}
\newline
\newline
\hyperref[sec:678]{< < < PREV SERMON < < <}
~
\hyperref[sec:toc]{[返目錄]}
~
\hyperref[sec:680]{> > > NEXT SERMON > > >}
\newline
\newline
撒迦利亞書 10:1-10:7
\newline
\begin{longtable}{cl}
\hline
\hline
章節 & 經文 (和合本修訂版)\\
\hline
10:1 & \begin{tabularx}{0.7\textwidth}{X} 春雨的季節，你們要向耶和華求雨。耶和華發出雷電，為眾人降下大雨，把田園的菜蔬賜給人。 \end{tabularx} \\ \\ \relax
10:2 & \begin{tabularx}{0.7\textwidth}{X} 因為家中神像所言的是虛空，占卜者所見的是虛假，他們講說假夢，徒然安慰人。所以眾人如羊流離，因無牧人就受欺壓。 \end{tabularx} \\ \\ \relax
10:3 & \begin{tabularx}{0.7\textwidth}{X} 我的怒氣向牧人發作，我必懲罰那為首的；萬軍之耶和華眷顧他的羊群，就是猶大家，必使他們如戰場上的駿馬。 \end{tabularx} \\ \\ \relax
10:4 & \begin{tabularx}{0.7\textwidth}{X} 房角石從他而出，橛子從他而出，戰爭的弓也從他而出，每一個掌權的都從他而出。 \end{tabularx} \\ \\ \relax
10:5 & \begin{tabularx}{0.7\textwidth}{X} 他們必如戰場上的勇士，踐踏仇敵如街上的泥土。他們必爭戰，因為耶和華與他們同在，他們必使騎馬的羞愧。 \end{tabularx} \\ \\ \relax
10:6 & \begin{tabularx}{0.7\textwidth}{X} 我要堅固猶大家，拯救約瑟家，我要領他們歸回，因我憐憫他們，他們必像我未曾棄絕他們一樣；都因我是耶和華－他們的神，我必應允他們。 \end{tabularx} \\ \\ \relax
10:7 & \begin{tabularx}{0.7\textwidth}{X} 以法蓮人必如勇士，他們心中暢快如同喝酒；他們的兒女看見就歡喜，他們的心必因耶和華喜樂。 \end{tabularx} \\ \\
[1ex]
\hline
\hline
\end{longtable}
我們從第十章的信息中,可以看見神對祂子民以色列人,和對全人類的心意；讓我們看見彌賽亞是拯救者,要拯救以色列民,更要拯救全世界的人.當以色列人在苦難中的時候,他們需要神施恩,賜下甘霖,解決他們物質的生活；但更緊要的,是屬靈生活的需要,當神看見祂的羊流離失所,沒有牧人的時候,便心裡哀痛,因此,彌賽亞要親自成為他們的牧人,主耶穌要成為人類的善牧,牧養他們.
\newline
\newline
首先我們注意第四節,這裡有幾個特別描寫彌賽亞的象徵的詞,有房角石,釘子和爭戰的弓.舊約非常著重建造的事,他們建造會幕,聖殿來敬拜神.聖經告訴我們,人類的一切活動不外乎有兩種可能:一是創世記十一章所記載的建造巴別塔,它是人類社會的活動,是人文主義的表現,人不需要神,他們自己興建巴別塔,結果他們失敗了,巴別塔被拆毀,無法建成.雖然這是古代所發生的事,但時至今日,人類仍是在建造自己的巴別塔.另一個可能,就是彼前二章所說的建造靈宮,凡蒙恩得救的人,都像活石,被建造成為靈宮,作聖潔的祭司.所以我們如果不是建造靈宮,便是建造巴別塔.別以為教會的事工當然是建造靈宮,如果我們倚靠自己,可能在教會中也建造巴別塔.真正的靈宮不是用泥土或磚塊建造的,而是用活石,就是用我們的生命來建成的,這靈宮就是教會.教會的功能是要建造在使徒和先知的根基上,而二者的重點又各有不同,使徒主要是作廣的工作,傳揚福音,使人得救；先知主要是作深的工作,造就信徒,培育信徒.教會需要一面傳福音,增加得救人數,同時也要有深度方面的長進,不能老是停留在道理的開端,我們要在神的話語上扎根,下功夫.今天中國教會對神的話仍然認識不夠,一九四九年大陸解放,很多佈道家和有名的傳道人都到了香港,所以港九便起了一個空前的教會復興；現在,三十年已經過去了,我們可以很冷靜很仔細的回顧一下教會的進展,不錯,這三十年來,人數是加增了,教會的活動也多了,但可惜信徒仍很少注重神的話語.一直以來,香港教會都有一種現象,就是信徒只喜歡聽道；聽不同的人,用不同的方法,講不同的道.甚至在同一教會內,每主日都換不同的講員,這本來是好的；因為我們可以從不同的情況下,接受不同的神的話語；可是,這並不是正常的現象,信徒真正得到神的恩典,是在教會的牧者身上,因為他知道信徒的需要,能按時分糧.有的信徒說,我們的牧者沒有供應,為甚麼呢?因為我們聽道有問題,我們愛聽一些淺白的,有趣的,還要詼諧一點的道；如果講得分量太重,我們便聽不下去,覺得太枯燥,太艱澀.結果我們所喝的不是有營養的奶,而是放很多糖的水；喝起來雖甜,但沒有多少營養,我們真需要在神的話語上下功夫,有了好的根基,然後再在上面建造.
\newline
\newline
第二樣提到的是釘子,釘子就是樁.搭帳棚的時候,帳棚的繩子要綁在木樁上,這樣帳棚的根基才得穩固.以賽亞書五十四章給我們看見一幅非常活潑的圖畫,我們要把帳棚擴大,將繩子拉長,將釘子,即樁打得穩固,這可以說是傳福音的異象,差傳的異象,或教會增長的異象；只有我們的生命在彌賽亞裡面,才得穩固,才能真正建立起我們的帳棚.
\newline
\newline
第三是爭戰的弓,英國一位文學家曾說:弓在神手裡就是恩惠的記號,在人手裡便成為爭戰的武器.聖經第一次出現弓形的樣式,是在創世記第九章內.當挪亞從方舟出來,抬頭看的時候,天空出現了一道彩虹,這是神恩惠的記號,和平的記號,表明神不再用洪水毀滅世界.虹是美麗的,多姿多采的,基督徒的人生也是一樣,能帶給世界希望和快樂,這是教會應有的見證.只可惜它落在人手中的時候,便成了爭戰的武器；舊約聖經多次提到神為爭戰的弓,祂稱為萬軍之耶和華,使我們想到舊約非常著重聖戰,意思是我們要為耶和華爭戰；要靠耶和華而爭戰,我們的耶和華是萬軍之耶和華.今天,教會是個爭戰的教會,我們要與罪惡宣戰,而這場爭戰是沒有終結的；教會如果太安穩,信徒如果太安逸,便是有問題,我們必須爭戰.啟示錄記載最後有哈米吉多頓的大戰,這是一場包括地上,天空,以至整個宇宙的激烈戰爭；但彌賽亞是得勝的君王,祂是全能的神,永在的父,和平的君.
\newline
\newline
跟著我們看第十章的下半段,這裡我們看到出埃及的經驗,我們知道以色列人曾從埃及為奴之地出來,這是第一次的出埃及,而這裡所記載的,是第二次的出埃及.以色列人被擄至外邦,現在從外邦之地出來,這是第二次的出埃及.至於第三次的出埃及,是耶穌登山變像時,和摩西以利亞談論去世的事,中文聖經的去世,原來就是出埃及,意思就是耶穌為我們釘在十字架上,成就救恩,這是個出埃及的經驗.還有一次出埃及,就是當主再來拯救我們之時,我們再也不屬於地上,這一切都要成為過去,這也是出埃及.在這裡我們看見神的作為和計劃,神只有一個心意,就是與人來往,亞當夏娃雖然失敗,但神沒有失敗；神叫亞當夏娃教他們的兒子敬拜,後來,希伯來的族長也敬拜；會幕,聖殿就是人敬拜神的地方,而敬拜的意思,就是體會神的同在.耶穌基督到世間來,聖經稱為以馬內利,就是神與我們同在.
\newline
\newline
接著我們來看第十一章,這章給我們看見的,都是些灰暗的,悲慘的圖畫；前面三節是牧長的哀歌,跟著講到羊沒有好牧人.這裡提到三個牧人,到底他們是誰呢?我們可以在歷史中找到這三個人；所謂牧人,就是領導羊群的人.在以色列的國家內,有三種領導的人,就是君王,祭司,先知.君王是治理人民的,但人作王常常會有問題,我們看見到列王時代的混亂,他們怎樣不遵從神,不效法他們的列祖,所以君王失敗了.祭司只注重外表的形式,沒有實際屬靈的生命,所以也失敗了.先知本來是傳神的信息,是以色列的良心,但他們卻講些不負責任的說話,他們說平安了,平安了,其實沒有平安,所以先知也失敗了.這是不是今天教會的問題呢?教會的失敗主要就是傳道人的失敗,是神的工人有問題,傳道人太注重名譽,錢財,沒有在屬靈方面追求.今天教會若要進步,若要更新,傳道人首先要在神面前悔改,而信徒,長執也要為傳道人禱告.很多時候,信徒袖手旁觀,帶著觀望,批評,論斷的態度；結果教會裡就有了惡性循環,信徒不尊重神的工人,長執們要控制傳道人,甚至是講台.這樣,傳道人便覺得連一點尊貴的身份也沒有,於是便懶散下來,或者是注重世俗了.所以教會若要復興,我們所有的人都要悔改.因為教會的問題不是別人的問題,而是我們自己的問題.很可惜今天教會與教會之間,弟兄姊妹之間,傳道人與信徒之間起了爭競；其實我們是一家的人,怎麼變成仇敵呢?求主憐憫我們.
\newline
\newline
神的心非常難過,因為他的百姓好像沒有牧人的羊群一樣.牧羊人有兩根杖,一根叫做榮美,一根叫做聯索,譯得好一些,就是一根稱為杖,一根稱為竿.杖是責打的,也是驅散野獸的；而竿是引導走錯路的羊群的.這裡提到兩根杖,榮美是表明神的恩惠,是神恩惠的約,但我們既然背約失信,所以神只好放棄.另一根聯索,是聯合的意思,以色列人南北二國要聯合起來,可是失敗了.今天教會應當聯合起來,但我們卻失敗了,神怎能施恩給我們呢?所以神說:我要興起一個壞的牧人,一個無知的牧人,不會牧羊,不會引導,只會陷害,使羊遭受禍害；結果刀必臨到他的膀臂和右眼上,膀臂表示有力量,而右眼則表明看顧,也是智慧,結果,他甚麼都沒有了.今天,教會需要牧人,需要有恩賜又有忠心的牧人,真正作牧養的工作.
\newline
\newline
另一方面,百姓本來就沒有仰望神,也不尊重神,甚至不尊重神的牧人,他們給他們的牧人三十塊錢,以為這是美好的價值,其實這只是一個奴僕的價錢,所以神吩咐撒迦利亞要將這三十塊錢去給yiu戶.我們注意第十二章十節:「我必將那施恩叫人懇求的靈,澆灌大衛家,和耶路撒冷的居民,他們必仰望我,就是他們所扎的.」以往我們太蔑視神的教會,神的工人,神的工作,如今我們必須仰望祂,因為有一天我們要面對面的見祂,到時我們抱怎樣的態度呢?所以趁著還有今天,讓我們靠著主的恩典,切實悔改.
\newline
\newline


\newpage

\section{第51屆~港九培靈會~唐佑之~第八講}
\label{sec:680}
\textbf{泉源和活水}
\newline
\newline
連結: \href{https://www.hkbibleconference.org/session-message/view/680}{\texttt{https://www.hkbibleconference.org/session-message/view/680}}
\newline
\newline
\hyperref[sec:679]{< < < PREV SERMON < < <}
~
\hyperref[sec:toc]{[返目錄]}
~
\hyperref[sec:682]{> > > NEXT SERMON > > >}
\newline
\newline
撒迦利亞書 13:1-13:1
\newline
\begin{longtable}{cl}
\hline
\hline
章節 & 經文 (和合本修訂版)\\
\hline
13:1 & \begin{tabularx}{0.7\textwidth}{X} 「在那日，因罪惡與污穢的緣故，必有一泉源為大衛家和耶路撒冷的居民而開。」 \end{tabularx} \\ \\
[1ex]
\hline
\hline
\end{longtable}
教會復興是我們每一個人在神面前的心願,就是讓神的榮耀彰顯出來,不是自己得滿足,而是讓主的心得到滿足.撒迦利亞書的精要,就是神的榮耀與教會更新；前半部我們看到八個異象,是論及教會更新的異象,是從以色列民族歷史信仰的經驗來看教會的真理,也應用到我們身上；後半部我們看見撒迦利亞將這些真理概括起來,不但想到過去,想到現在,也想到將來.基督徒也要學習這個功課,我們想到過去就感恩,想到現在要審查,想到將來要盼望,因為神有說不盡的恩賜.
\newline
\newline
今天我們看撒迦利亞書十三章,第一節說:「那日必給大衛家和耶路撒冷的居民,開一個泉源,洗除罪惡與污穢.」十四章八節又說「那日必有活水從耶路撒冷出來,一半往東海流,一半往西海流,冬夏都是如此.」泉源和活水象徵神的生命,屬靈的恩典,聖靈的工作,這是我們必須體會的.神是我們一切的泉源,是我們取之不盡,用之不竭的.先知書特別把外邦人的泉源和以色列人的泉源作個強烈的對比.何西阿書十三章十五節講到外邦人的泉源是會乾的,因為他們沒有繼續的供應；耶利米書五十一章說外邦人的泉源是會枯竭的,但以色人的泉源卻不會,因為耶和華是他們泉源.但可惜以色列人沒有完全倚靠神,所以在耶利米書二章那裡,先知特別責備以色列人做了兩件錯事,就是離棄了活水的泉源,為自己鑿出一個破裂不能存水的池子,這是個非常嚴重的警告.我們常為自己建築水池,想到經濟的力量,物質的需要,知識和技能的培養,但真正的泉源卻是在神那裡.我們常說:「人的盡頭是神的開始.」這話對一個沒有信主的人來說,是非常恰當的,當他們到了盡頭之時,如果願意接受福音,神的工作便開始,但這句話只能用在福音事工上,卻不能應用在信徒的生活中.我們不能靠自己的方法,努力去做事,到沒有辦法時才求神憐憫,我們應該一開始便倚靠神.
\newline
\newline
如果我們要求長進,對神便要有深度的認識,以西結書四十七章描寫在聖殿祭壇裡面,流出一股活水來,這水愈流愈長,愈遠,愈廣,愈深,到底有甚麼意思呢?從新約的觀點來看,它代表我們教會的復興,教會更新的異象.耶穌基督在十架上所成就的救恩,為人流出一股生命的活水來,叫人可以靠祂而得著豐盛的生命.這活水是從十字架的祭壇流出來的,起初很淺,只能到踝子骨；彼得和約翰曾在聖殿的美門口那裡,用耶穌基督的名,醫好一個瘸腿的人,他的腳和踝子骨便立刻健壯,蹦蹦跳跳的讚美神.然後水往上漲,直流到膝蓋那裡,表示我們與主同行,過禱告的生活；之後水流到腰那裡,腰表明事奉,是束腰和爭戰,我們不但禱告,更有力量與主同工；最後這水把整個人蓋過,不能走動了,只能游過去,這表示我們完全藏在神的愛和聖靈的大能中,完全被聖靈充滿,與神同在,與神合一.撒迦利亞所講的泉源和活水,也有同樣的意思.
\newline
\newline
跟著我們看十三章二節至九節,這裡又重複提到牧人的事,就是指到君王,祭司,先知那些領導人物,他們都失敗了.大衛雖是個好君王,但在道德方面也有污點,所羅門在晚年也失敗,分裂以後,北國幾乎沒有一個好王,南國的君王大部份也是不好,有的雖然好,但也只是短暫的,結果還是失敗.祭司們也失敗,他們把偶像與真神的敬拜混在一起,沾染污穢,沒有聖潔.先知們也是失敗,他們輕忽地醫治百姓的損傷,虛報平安的信息,即使有些好先知,忠心傳神的活,但是百姓卻掩耳不聽.總括來說,舊約沒有一個先知,祭司或君王是成功的,他們統統失敗了.所以神要為他們開一個泉源,洗除罪惡與污穢,神要從地上除滅偶像的名,再沒有假先知與污穢的靈.今天,教會最迫切需要的,便是牧師.無可否認,有些傳道人牧師,不一定是蒙主選召的,所以他們傳道的職分被人誹謗.求神憐憫我們,我們所需要的,是那些真正蒙召的牧者.可是,教會的弟兄姊妹也不能專靠牧師一人工作,牧師的職責,是要專心以祈禱,傳道為事,至於其他的事務,便應該由長老,執事或弟兄姊妹們來分擔.不過,初期教會還給我們看見一個很好的榜樣,司提反和腓利,本來是管理飯食的執事,但結果他們都做了傳道工作,司提反因傳福音而殉道,腓利則到撒瑪利亞傳道,後來更帶領埃提阿伯的太監信主.所以我們看見,牧師固然要關心信徒的需要,作牧養的工作；而信徒也應該參與,如果在屬靈方面有長進的,也應當擔負牧養的工作,換句話說,我們都要一同起來,關心教會的工作,特別是關心弟兄姊妹屬靈的需要,真正的關懷每一個人.聖經說,當牧人被擊打,羊群就分散了,主耶穌被擊打時,羊群不都是分散了嗎?如果教會的牧師被擊打,那麼信徒不都是成為可憐的分散的羊嗎?今日,魔鬼常利用他的使者來擊打教會的牧人,為的是拆散教會,我們要小心,不要成為撒旦利用的工具.第七節說,我必反手加在微小者的身上,如果我們不去照顧那些軟弱的小羊,神便要親自的照顧他們.
\newline
\newline
我們曾經講過餘數的道理,這裡又再次提到餘數,第八節說:「耶和華說,這全地的人,三分之二必剪除而死,三分之一仍必存留.我要使這三分之一經火,熬煉他們,如熬煉銀子,試煉他們,如試煉金子.」今天教會也需要熬煉,如果我們真正敬畏神,神必把工作託付我們；如果神使用我們,祂一定要熬煉我們,因為我們的渣滓太多,不夠純淨,所以要經過神的熬煉,這就是為甚麼我們在日常生活中會有疾病,憂傷,困苦等事.神熬煉我們,為的是潔淨我們,叫我們在潔淨之後可以為神所用.瑪拉基書三章形容神像一位熬煉金子的工匠,坐在爐子旁邊,有時火比較慢,但有時又會很猛,甚至是我們難以忍受的痛楚,神雖然捨不得,但是火若不猛,渣滓便無法浮起來,渣滓不除淨,我們便仍不純淨；所以神在我們旁邊,非常忍耐的熬煉我們,直到把渣滓除掉.金子和銀子的表面純淨到好像一面鏡子的時候,祂才住手,那就是說,直到我們能反映祂榮耀的形象為止.今天,我們若要為神所用,也必受祂嚴厲的熬煉.
\newline
\newline
十四章一到三節講及最後的爭戰,我們記得有歌革和瑪各的爭戰,有哈米吉多頓的爭戰,世界所有的罪惡都要消滅,耶和華要為我們爭戰,我們的主是一位得勝者.第四節說:「那日,他的腳必站在耶路撒冷前面朝東的橄欖山上,這山必從中間分裂,自東至西,成為極大的谷.」這橄欖山是在耶路撒冷前面朝東的地方,是日出之地,黑夜快將過去,公義的日頭快要上升,耶和華的日子要來到,這不是光明的日子,對罪惡的人來說,這是黑暗的日子,因為這日必沒有光,三光都要退縮；但是對我們信靠神的人來說,到了晚上才有光明,因為黑夜快要成為過去,早晨就要來到.主來的時候,地大震動,山谷要分開,主耶穌要和一切的聖者同來,這是祂的恩典.
\newline
\newline
最後,我要提到兩件事情,第一是住棚節,十六節說我們最要緊是守住棚節,「所有來攻擊耶路撒冷國中剩下的人,必年年上來,敬拜大君王萬軍之耶和華,並守住棚節.」他們本來是攻擊耶和華的外邦人,可是外邦人中也有餘數,凡是信靠神的人,都要敬拜耶和華.這裡特別提到守住棚節,是因為住棚節尚未完全應驗,我們知道逾越節已經應驗了,耶穌就是逾越節的羊羔；五旬節也應驗了,因為五旬節時教會被建立起來；但是住棚還未應驗,今天我們還居住在地上,要等候有一天神的帳幕永遠住在人間,我們需要住棚節的實現.而今天我們最要緊的,是學習敬拜祂,在教會,在家庭,在個人的生活中好好的敬拜祂.這裡有很清楚的敬拜儀式,中文聖經用大字把它標明出來,就是「歸耶和華為聖」,這裡說連馬也歸耶和華為聖,耶和華殿內的鍋,都要像祭爐前的碗一樣.馬本來是打扙的,但神也可以叫牠成為聖潔,鍋是耶和華殿中最污穢的器具,但現在要像祭爐前的碗一樣的聖潔.在耶和華的殿中,必不再有迦南人,就是不再有犯罪和污穢的人.啟示錄最後一章說,主再來的時候,污穢的仍舊叫他污穢,聖潔的仍舊聖潔,所以我們現在便要聖潔,否則到主再來時就太遲了.我們要過聖潔的生活,要歸耶和華為聖,好像新婦裝備整齊,歡歡喜喜的迎接祂一樣.
\newline
\newline


\newpage

\section{第51屆~港九培靈會~張慕皚~第一講}
\label{sec:682}
\textbf{認識我們的時代}
\newline
\newline
連結: \href{https://www.hkbibleconference.org/session-message/view/682}{\texttt{https://www.hkbibleconference.org/session-message/view/682}}
\newline
\newline
\hyperref[sec:680]{< < < PREV SERMON < < <}
~
\hyperref[sec:toc]{[返目錄]}
~
\hyperref[sec:683]{> > > NEXT SERMON > > >}
\newline
\newline
感謝神!給我們有第五十一屆的港九培靈研經大會.過去神實在賜福給我們,很多教會,信徒都因培靈大會而得到造就,實在感謝神.
\newline
\newline
在普世華人教會中,香港佔有一個重要的地位,可以說是一個屬靈重要的基地,產生出不少得力的信徒,和神的僕人.神實在復興香港的教會,以致普世的華人教會,因而受影響,也得到復興.
\newline
\newline
我今次在培靈會的講題,是在末世時代,怎樣操練自己,成為神手中一個成熟合用的器皿.我們會從多方面的跡象去看這問題.因這末世充滿危機,挑戰；我深信神在等候一些人,到一個成熟合用的時候,給予神使用.因此我以這個題目與大家分享.今天的題目是「認識我們的時代」.以後的八天再繼續思想我們所需要的,進而成為神合用的器皿,請大家多多禱告,記念這個聚會.
\newline
\newline
請打開馬太二十四章,在這章聖經中,耶穌將這末世時代的輪廓清楚地描述出來.因時間關係,我們只看馬太二十四章三至十四節.
\newline
\newline
為何要認識時代?
\newline
\newline
聖經給我們看見,神在每一個世代中,在教會的每一個時代中,都要找一些「先知型」的人替祂工作.先知的特點,是針對那個時代的需要,向那時代的人說話；先知乃是神所興起的,所以先知要具備三個條件:
\newline
\newline
1.認識神:包括知識和頭腦上的認識,並在生命中確實經歷神的存在.
\newline
\newline
2.認識神的話:深刻地認識神的話,才可以得到造就.
\newline
\newline
3.認識時代:我們要用聖經的話語,向這時代說話.例如:軍隊進攻,要知己知彼,才能百戰百勝；醫生要對症才能下藥.求主幫助我們有一個時代感,使我們認識自己所處的時代.
\newline
\newline
耶穌很留意在自己的時代中有甚麼事情發生,他說:這是一個麻木不仁的時代,又是一個邪惡淫亂的世代.所以主嚴厲的責備猶太人,不曉得分辨所處的時代.耶穌對那時代的人說:「晚上天發紅,你們就說天必要晴；早晨天發紅又發黑,你們就說今日必有風雨.你們知道分辨天上的氣色,倒不能分辨這時候的神蹟.」(太十六2-3)聖經中很清楚地將這時代描繪出來,而這末世時代是指耶穌第一次與第二次回來的中間的一段時間.(彼前一20)記載基督在創世以前,是豫先被神知道的,卻在這末世纔為你們顯現.(希一1-2)也說「神既在古時藉著眾先知,多次多方的曉諭列祖,就在這末世藉著他兒子曉諭我們；又早已立祂為承受萬有的,也曾藉著祂創造諸世界.」這說明耶穌道成肉身來到世界,這也是末世的開始.
\newline
\newline
看聖經末世預兆,我們要避免兩方面的極端.第一,是過份看重聖經中末世的預兆.第二,是不留意聖經所說末世的預兆,不留意世界的潮流.人更不能將聖經與世界合起來,而儆醒等候被神使用.舊約有一半以上的經文,新約也有二分一的經文,是論到耶穌再來的.這些經文所說的話,都是要催促我們儆醒,成為神合用的器皿.
\newline
\newline
如何知道耶穌快再來?
\newline
\newline
聖經給我們的提示,就是耶穌回來時,預兆發生得愈緊密,愈嚴重.太二十四12提到只因不法的事增多,許多的人愛心纔漸漸冷淡了.提後三13也記載:「只是作惡的和迷惑人的必越久越惡,他欺哄人也被人欺哄.」人類的罪惡越來越可怕,這是聖經給我們的預兆；當我們看見這些預兆時,我們就要儆醒,求主幫助.
\newline
\newline
馬太廿四章很清楚地說明,我們的時代有三個預兆.
\newline
\newline
(一)失去控制的時代.「耶穌在橄欖山上坐著,門徒暗暗的來說,請告訴我們,甚麼時候有這些事?你降臨和世界末了,有甚麼預兆呢?耶穌回答說,你們要謹慎,免得有人迷惑你們；因為將來有好些人冒我的名來說,我是基督,並且要迷惑許多人.......且有很多假先知起來,迷惑多人.」(太二十四3-5,11)耶穌說必有很多假先知,假彌賽亞起來.因為人心飢渴,希望有超自然的力量,來幫助他們；可惜他們卻不知道有一位真的彌賽亞,正在等候要幫助他們.保羅在提後三1說:「你該知道,末世必有危險的日子來到.」「危險」原文的解釋,是失去控制之意.耶穌在這裡向我們說到,有假彌賽亞的出現.彌賽亞是神所膏立的,是要拯救世界的那一位.人類看重自己的科學知識,不錯,這些帶來很多方便；卻也帶來更多的危險,使我們走進迷失的時代,失去控制的時代.倘若教會向人傳講真的彌賽亞,那麼假先知,假彌賽亞便不會應運而生.
\newline
\newline
現今的人喜歡看邪靈,邪神,占卜,星相的電影,可見他們心靈的空虛,他們處在一個迷失的時代中；更反映出他們渴望有超自然力量,來幫助他們走出這時代.
\newline
\newline
美國醫學界提出全人醫治,不單用藥物,更借用神秘主義來醫病.因他們發覺人們的疾病,多由心靈引起.在這裡給我們知道,只有真的彌賽亞才可以醫治他們.
\newline
\newline
這是一個失去控制的時代,環境,生態各方面都失去控制.在六十代年初,已有很嚴重的環境污染,很多科學家預言,人類將會自己消滅自己.
\newline
\newline
經濟失去控制──通貨膨脹很厲害,在南美洲達到百分之四百.
\newline
\newline
政治失去控制──全世界最強的國家-美國,總統卡達雖然是基督徒,但仍不能把美國治理好.
\newline
\newline
教育失去控制──各地的教育制度都不成功,不能引導學生向善.知識灌輸也不足,學生畢業後隨即失業,造成社會上嚴重的青少年罪犯問題.
\newline
\newline
道德失去控制──沒有絕對標準.
\newline
\newline
宗教失去控制──異端,邪派勃起,不道德的事,甚至發生在教會內.
\newline
\newline
在這些情形之下,我們看見人類正墮落在恐懼迷失中,和無所適從的情況之下.人也不能再說「人定勝天」的話,人是衡量一切的標準.但到這末世時代只是瞎子領瞎子,所以教會應該起來扮演先知角色,信徒也應該起來作先知的事奉.我們都應該起來向世人發出責備和安慰的說話；因為他們渴望真彌賽亞的出現,求主幫助我們!
\newline
\newline
(二)人際關係大崩潰的時代(太二十四6,10,12)在末世時候,有一個很特別的現象,而且也是個嚴重的現象,就是國家,民族,家庭中的分裂,和教會內愛心的失落.自從人類墮落後,罪便產生自我的生活,自私自利,而且拒絕聖經的教訓.求主教導我們,曉得將愛心推廣出去,因為只有愛才能將人類連結起來.
\newline
\newline
人的自私與自利,帶來多方面的分裂.
\newline
\newline
1.人與神的分裂:人以自己的能力,做自己不能做的事.在香港有三分一人口,就患有輕微的精神病；這是因為他們的罪,使他們與神分離.
\newline
\newline
2.人與人之間的分離:聖經記載,民要攻打民,國要攻打國,彼此陷害,恨惡.因此有一位哲學家說:「別人就是我們的地獄」我們不能接受別人.
\newline
\newline
3.人與人自然的分離:有飢荒,地震.
\newline
\newline
4.家庭分裂:十九世紀的詩人說:「家庭是甜蜜的.」但現在美國的家庭已經變質了!人要離家出走,因為失去家庭的愛.我們看見家庭的分裂,實在是人際關係大崩潰的時代.求主加給我們聰明與智慧,使我們在這時代裡做時代的先知.教會應起來,將耶穌基督的愛伸展出去.因為在人際關係大崩潰的時代中,我們首先要由家庭開始,這才能發展到世界,和社會.家庭是社裡最基本的單位,盼望我們每一個信徒,能首先穩定自己的家庭,然後發揚出去,用愛心作為包裹,願主幫助我們.
\newline
\newline
(三)福音傳遍的時代太二十四14記載:「這天國的福音要傳遍天下,對萬民作見證,然後末期纔來到.」在末世時代是福音傳遍的時代,神要復興祂的子民,到世界各地傳福音.
\newline
\newline
神已經作好多方面準備,祂已經預備好人心,藉假先知與迷失的世代,使人渴望真的彌賽亞到來.求主幫助我們.神已經預備好莊稼,莊稼亦已經成熟,我們必須起來工作了.
\newline
\newline
神又預備了快速的工具,將世界縮小.大眾傳播媒介,也成我們廣傳福音的好工具.神又加給教會新的醒覺,使我們在這時代,能作全面性的傳福音.祂興起差傳與總動員工作.弗四11說:神給我們有先知,祭司,牧師.為叫他們在教會內,全時間事奉神,和造就裝備每一信徒,各人在自己的崗位上,負起自己的責任.
\newline
\newline
神學教育與研經課程,都是神所興起的,為要幫助更多的信徒,使肢體彼此配合為神工作.祂又興起很多屬靈的機構,使他們願意成為神的僕人；將福音廣傳,使福音傳得更速,更快.
\newline
\newline
一切神已經預備好了,只等待你和我起來,你願否成為神合用的器皿嗎?倘若願意,我們必須先有時代感,才能成為時代的先知.求主施恩,賜福,憐憫我們；使我們成為合用的器皿,在這末世時代中榮耀神.
\newline
\newline


\newpage

\section{第51屆~港九培靈會~張慕皚~第二講}
\label{sec:683}
\textbf{屬靈生命的更新}
\newline
\newline
連結: \href{https://www.hkbibleconference.org/session-message/view/683}{\texttt{https://www.hkbibleconference.org/session-message/view/683}}
\newline
\newline
\hyperref[sec:682]{< < < PREV SERMON < < <}
~
\hyperref[sec:toc]{[返目錄]}
~
\hyperref[sec:685]{> > > NEXT SERMON > > >}
\newline
\newline
感謝主!帶領培靈研經進入第三天.昨天香港遭遇大風暴,正如主耶穌基督再來之前有大災難一樣.在末世的時代,有大風暴,而又有片刻寧靜,祂要來之前,我們要儆醒,不可沉睡,要儆醒自守.神要裝備我們,好讓大風暴來時,可以給神用,作神合用的器皿.
\newline
\newline
今日要思想的題目:「屬靈生命的更新.」我們要認識這個末世的危機,我們要求神恩待我們；因為只有祂才能使我們更新,使我們成為神手中合用的器皿.
\newline
\newline
神今天要用怎麼樣的人呢?
\newline
\newline
有人以為神要用的是一班神學生,牧師,傳道人.其實當我們讀聖經的時候,發覺在末世時代,神要選用的,乃是一些屬祂的人.(徒二16-18)記載先知約珥所說的話:「神說,末後的日子,我要將我的靈澆灌凡有血氣的；你們的兒女要說豫言,你們的少年人要見異象,老年人要作異夢.在那些日子,我要將我的靈澆灌我的僕人和使女,他們要說豫言.」所以不論男,女,老,少,都要裝備自己,迎接挑戰,因為有大風暴要臨到這世代.
\newline
\newline
聖經告訴我們,一個人的屬靈生命可能要經過三個階段:
\newline
\newline
1.未信主時,我們的生命是死在罪惡過犯中,與神隔離,與神為敵,好像以色列人,在埃及一樣,過著為奴痛苦的生活,被失望侵蝕,等待死亡.
\newline
\newline
2.主耶穌為我們的罪,已死在十字架上,當我們仰望祂的時候,死的生命就有生氣；當我們死在過犯中的時候,神便叫我們與基督一同復活,一同坐在天上.(弗二5)
\newline
\newline
3.我們初信主時,罪雖然得了赦免,有了復活的生命,卻仍未有基督徒的特質,單有基督徒的名份；因此我們要不斷更新,經歷,才能作一個實實在在的基督徒,否則我們必會生活在自己肉體(罪性)中,脫不了從亞當所承受而來的罪性.所以我們要在生命中不斷更新,長進,攀登高峰.倘若不這樣做,便好像以色列人出埃及後在曠野的生活一樣；充滿埋怨,黑暗,枯燥無味,充滿神的管教與懲罰.聖經說這是一個嬰孩的情況,求主今日更新我們,使我們好像以色列人經過四十年曠野生活,進入迦南的情形一樣,是進到流奶與蜜之地,充滿神的恩典,光明與喜樂,今天我們願不願意進入流奶與蜜之地呢?
\newline
\newline
現在我要以聖經中兩個幔子與大家思想生命更新的問題.
\newline
\newline
(一)第一個幔子記載在(可十五33-39):「從午正到申初,遍地都黑暗了,申初的時候耶穌大聲喊著說,以羅伊,以羅伊,拉馬撒巴各大尼；翻出來就是我的神!我的神!為甚麼離棄我?旁邊站著的人,有的聽見就說,看哪!他叫以利亞呢.有一個人跑去把海絨蘸滿了醋,綁在葦子上,送給祂喝.說,且等著看以利亞來不來把他取下.耶穌大聲喊叫,氣就斷了,殿裡的幔子,從上到下裂為兩半；對面站著的百夫長,看見耶穌這樣喊叫斷氣,就說,「這人真是神的兒子.」這段聖經很明顯地,形容耶穌為我們的罪死在十字架上的情形.在救恩歷史中有兩個重要的黑暗的時刻.第一個黑暗的時刻記載於人類的始祖,受魔鬼引誘,墮落犯罪,這是人類歷史黑暗的一面.第二個黑暗的時刻是神的兒子,為人類的罪死在十字架上.為甚麼是最黑暗呢?因祂擔當我們的罪,這是多麼的可怕.主耶穌在十字架上也痛苦地呼喊說:「我的神,我的神,為甚麼離棄我?」(可十五34).罪叫我們與神分離,耶穌基督擔當我們的罪,因此神不得不與祂的兒子分離.當耶穌死在十字架上時,有甚麼事情發生呢?「殿裡的幔子由上至下裂為兩半.」(可十五38).幔子是進入至聖所的一個阻隔,因主耶穌的死,打破了世人與神之間的阻隔.(來十19-20)記載說:「弟兄們,我們既因耶穌的血,得以坦然進入至聖所,是藉著祂給我們開了一條又新又活的路,從幔子經過,這幔子就是他的身體.」主耶穌犧牲了自己的身體,承擔了我們的罪,我們才能與神面對面的交通,毫無阻隔.
\newline
\newline
不錯,聖經的中心思想是救贖,但也是以馬內利,就是神與我們同在；叫我們再一次與神過一種甜蜜的生活.
\newline
\newline
在救恩歷史當中,神藉著三個階段與人同在:
\newline
\newline
1.舊約時代是藉獻祭在會幕中與神同在,這是間接的交通,要通過大祭司.
\newline
\newline
2.耶穌在十字架上死了,擔當我們的罪,幔子裂開,聖靈居住在人的中間,但仍有罪和物質在我們當中.
\newline
\newline
3.人最後的盼望記載在(啟二十一1-5)「我又看見一個新天新地,因為先前的天地已經過去了,海也不再有了,我又看見聖城新耶路撒冷由神那裡從天而降,豫備好了,就如新婦妝飾整齊等候丈夫.我聽見有大聲音從寶座出來說,看哪,神的帳幕在人間,祂要與人同住,他們要作祂的子民,神要親自與他們同在,作他們的神.神要擦去他們一切的眼淚,不再有死亡,也不再有悲哀,哭號,疼痛,因為以前的事過去了.坐寶座的說,看哪!我將一切都更新了,又說你要寫上,因這些話是真實的.」新天地的情形是人進一步與神親近,天,地和海都要過去,是指物質,罪惡之惡,這一切都要被除去,人與神完全合而為一.
\newline
\newline
求神幫助我們,今天我們還未到第三個階段,仍處於第二個階段,我們當怎樣尋求更新呢!除去第二個幔子?第一個幔子已被耶穌除去,但人仍未脫離第二個幔子的綑縛.
\newline
\newline
第二個幔子記載在(林後三14-17):「但他們的心剛硬,直到今日誦讀舊約的時候,這帕子還沒有揭去,這帕子在基督裡已經廢去了；然而直到今日每逢誦讀摩西書的時候,帕子還在他們心上,但他們的心幾時歸向主,帕子就幾時除去了.主就是那靈,主的靈在那裡,那裡就得以自由.」帕子等如幔子,帕子就是自我為中心的思想,在我們內心交織而形成的.人處處只為自己著想,像撒旦一樣,(賽十四13-14)記載說:「你的心裡曾說,我要昇到天上,我要高舉我的寶座在神眾星以上,我要坐在聚會的山上,在北方的極處,我要昇到高雲之上,我要與至上者同等.」只顧「我」,「我要......我要.......」因罪的根源,而產生了驕傲,我們不再以神為自己生活的中心,只能揀選撒旦.耶穌在(路十二16-21)裡說了一個比喻:一個財主自私自利,以自己為中心,他收藏糧食,叫自己的靈魂安安樂樂的度日.但神說:我今天要收回你的靈魂,你所有的要歸誰呢?聖經稱那財主為愚笨的人,因為他有了錢之後,就以自我為中心.
\newline
\newline
我們可細看財主與和拉撒路的比喻.財主身穿紫色袍,而拉撒路全身長滿瘡,但財主不顧念拉撒路的需要.當財主死後,下到陰間,因他只顧自己；拉撒路死後卻在亞伯拉罕的懷抱裡,因為他一生倚靠神.耶穌也叫我們要先求神的國和祂的義,這些東西都要加給我們(太六33),這是神的應許.
\newline
\newline
一個自私的人去到教會,只是想著自己會得著甚麼,我們到教會崇拜,是先將神的榮耀歸給祂,其次才為自己求福.只想自己的得著,是幼稚不成熟的表現.
\newline
\newline
帕子有不同的質和模式,但都顯明我們有罪,各人的罪不同,但都出於自我.
\newline
\newline
亞伯拉罕的罪,是看重兒子過於看重神,當他將帕子除去後,神就大大賜福給他；這事以後,神稱亞伯拉罕為祂的朋友(雅二23).今天我們的生命應進到與神為朋友的階段.聖經上說:我們的神是忌邪的,嫉妒的,我們不能讓人奪去我們愛神的心.
\newline
\newline
今天在你手中是否有帕子阻隔你與神的關係呢?但願我們能像亞伯拉罕一樣,除去心中的帕子,成為神手中一個合用的器皿,所以我們必須更新.
\newline
\newline
以掃的帕子是一碗紅豆湯,為著這碗湯,他放棄了長子名份,長子的名份,是屬天的祝福,但以掃卻甘願放棄了,所以他是神所恨惡的.
\newline
\newline
今天的人只追求物質的享受,貪戀世俗的事物；所以不能將帕子除去,他們就與神的愛隔絕了.
\newline
\newline
摩西的帕子是他生長於皇宮中,受埃及的學問,於是以為自己有能力拯救同胞；他驕傲,自大,他的帕子是自高,所以神管教他,要祂成為完全謙卑的人.四十年的曠野訓練,神除去他驕傲的帕子(申三十四10)他死了,聖經說以後以色列中再沒有興起先知像摩西的,他是耶和華面對面所認識的.因為他放棄了自高,自大的帕子.
\newline
\newline
約伯的帕子是他對神只有頭腦上的認識,因此神大大的對付他,教訓他,使他厭惡自己,最後他說:「從前我風聞有你,現在我親眼看見你.」(伯四十二5)
\newline
\newline
今天人只希望追求認識神的知識,頭腦上的知識,但神要我們通過這些知識而經歷神的實在,否則我們對神的認識只不過好像約伯一樣.因此在聖經(林前八 1)警告我們說:「知識是叫人自高自大,惟有愛心能造就人.」求主幫助我們.神今日給我們很多研經課程,求神使我們能藉著這一切經歷,支取祂那豐盛的慈愛.
\newline
\newline
耶穌基督給我們立了一個好榜樣,就是犧牲自己釘在十字架上,祂說:「若有人要跟從我,就當捨己,背起他的十字架來跟從我.」(太十六24)
\newline
\newline
耶穌來到世上不單成全救恩,而且做了一個好榜樣,祂過一個無罪的生活.所以我們靈命的更新就是要不斷向耶穌學習,完全捨己.求主幫助!
\newline
\newline
人靈命長進有三個階段:
\newline
\newline
1.單方面的同在,這是軟弱的生命.
\newline
\newline
2.主的同在是偶然的,在困難得解決時才知道有主的同在,而在平常的生活中則只有些少領略到.
\newline
\newline
3.面對面的與神同在,像摩西一樣.這是我們可以達到的,但必定要付上代價與追求；聖經告訴我們要敬虔,要成為一個成熟的器皿,我們必要付出很大的代價,才能進到更善的地步.求主幫助我們除去心中的帕子.
\newline
\newline
(弗四11)說:「祂所賜的有使徒,有先知,有傳福音的,有牧師和教師,為要成全聖徒,各盡其職,建立基督的身體.」他們在教會的要合時間事奉神和造就裝備每一信徒,各人在自己的崗位上負起責任.
\newline
\newline
神學教育與延伸都是神所興起的,為要幫助更多的信徒,使肢體彼此配搭為神工作.祂又興起很多福音的機構,使他們作教會的僕人,將福音廣傳,傳得更快,更速.
\newline
\newline
神已經預備好一切,只等待你,我起來!你願意成為神合用的器皿嗎?我們必須先有時代感,才能成為時代的先知.求主祂恩,賜福,憐憫我們,使我們成為合用的器皿,在這末世時代中榮耀神.
\newline
\newline


\newpage

\section{第51屆~港九培靈會~張慕皚~第三講}
\label{sec:685}
\textbf{禱告生活的更新}
\newline
\newline
連結: \href{https://www.hkbibleconference.org/session-message/view/685}{\texttt{https://www.hkbibleconference.org/session-message/view/685}}
\newline
\newline
\hyperref[sec:683]{< < < PREV SERMON < < <}
~
\hyperref[sec:toc]{[返目錄]}
~
\hyperref[sec:686]{> > > NEXT SERMON > > >}
\newline
\newline
昨天我們講到屬靈生活的更新,更講及一個合乎神使用的人,必須是一個成熟的器皿.要成為一個成熟的器皿,就必須要不斷的更新,與主不斷接近.怎樣才能與主不斷接近呢?祈禱便是一個重要的環節,因為祈禱可以使我們與神更加親近,使我們更加成熟；所以我們今天要思想的題目是──「祈禱告活的更新.」
\newline
\newline
在可十四32-42中,記載耶穌在客西馬尼園的禱告.
\newline
\newline
耶穌在世三十三年,最後被釘在十字架上,為我們的罪而死.耶穌在世末後的三年裡,用來傳道和替人趕鬼,而且在這三年中,祂曾接受了兩次重大的試探.第一次的試探是主在曠野禁食禱告四十天,但魔鬼卻藉此來試探祂.第二次是主在被釘上十字架之前的試探,這是屬於肉體上的試探,因為主也是人,祂亦有軟弱,會疲倦與飢渴,他也會驚恐,魔鬼要乘虛而入.在這兩次的試探中,主都得勝了,祂就成了我們得勝的根源.(來二18)很清楚地告訴我們說:「祂自己既然被試探而受苦,就能搭救被試探的人.」(來四14-16)說:「我們既然有一位已經升入高天尊榮的大祭司,就是神的兒子耶穌,便當持定所承認的道,因我們的大祭司並非不能體恤我們的軟弱,祂也曾凡事受過試探與我們一樣,只是祂沒有犯罪.所以我們只管坦然無懼的來到施恩的寶座前,為要得憐恤,蒙恩惠作隨時的幫助.」又在(來五7-9)說:「基督在肉體的時候既大聲哀哭,流淚禱告懇求那能救祂免死的主,就因祂的虔誠,蒙了應允.祂雖然為兒子,還是因所受的苦難學了順從.祂既得以完全,就為凡順從祂的人,成了永遠得救的根源.」由以上三段經文中,我們可以體會到主既身為神的兒子,尚且要學習和忍受這苦難與順服的功課,何況我們呢?
\newline
\newline
現在我們要從客西馬尼園中,看一些屬靈的教訓.
\newline
\newline
(一)客西馬尼園是一個危機的時刻.在聖經的救恩內,也有兩個危機時刻.第一是伊甸園內的試探,亞當受魔鬼的引誘,終於失敗在魔鬼的手下.第二個是客西馬尼園內的試探,末後的亞當,終於克服了一切,勝過魔鬼,通過十字架,為我們成全了救恩,將祂的子民帶到永生的福份中.
\newline
\newline
末世時代是一個危機時代
\newline
\newline
第一天我們亦有講及末世之危機,不論在經濟,教育或道德上都是充滿了危機.
\newline
\newline
今天我們要從教會外和教會內兩方面去看末世之危機.我們不單關心教會內的危機,而又要關心教會外的危機,雖然我們成為神的子民,但主耶穌卻對我們說:「他們不屬世界,正如我不屬世界一樣.」(約十七16)所以內外之危機,我們都要關心.
\newline
\newline
末世帶來三個危機:
\newline
\newline
(一)知識增長帶來之危機──今天知識增加和增長得很快,因為科學昌明,帶給人類很多的貢獻；人類在每一分,每一秒之中,都有新知識的產生,但因為增長的快速,令人類無從適應,以至產生衝激.這些衝激帶給人類很多的煩惱.
\newline
\newline
(二)戰爭危機──世界上最強的兩個國家,他們在軍備上都互相競爭,務求勝過對方.在一九四五年墨西哥州內,美國發明了原子彈,科學家看見原子彈發彈之時,他說:「我的神啊!我們已為自己創造了地獄.」不錯,現今的人生活在一片戰爭的恐懼之中,軍備上的競爭,隨時隨地會帶給人類滅亡之途.
\newline
\newline
(三)人口膨脹的危機──人口不斷的膨脹,導致一個嚴重的危機,每一天有一萬人死於營養不良,飢餓之下.我們生活於舒適之香港,所以不曉得飢餓是什麼的一回事.在世界上有百分之六十的人死於營養不良,百分之三十的人口是接近死亡的邊緣上,為的是什麼?是因為糧食的缺乏,終有一天,糧食缺乏會帶來一次嚴重的糧食戰爭.
\newline
\newline
除了以上三點外,科學家還告訴我們有人口,糧食,能源,戰爭,環境污染等等五大危機,這些東西會隨時隨地將人類全部消滅.
\newline
\newline
教會內亦發生危機,啟示錄所提到的七個教會,只有兩個是不受主責備的,其餘的五個都要受到責備.
\newline
\newline
五個受神責備而又充滿危機的教會
\newline
\newline
第一個是以弗所教會──它是一個失去心的教會,我們看見這教會十分勞碌地工作,事奉主.但這一切只是出於責任感,因此得不到神的喜悅,只有事工變成制度化.在(林前十三1-3)中保羅告訴我們說:「我若能說萬人的方言,並天使的話語,卻沒有愛,我就成了鳴的鑼,響的鈸一般.我若有先知講道之能,也明白各樣的奧秘,各樣的知識,而且有全備的信,叫我能夠移山,卻沒有愛,我就算不得甚麼!我若將所有的賙濟窮人,又捨己身叫人焚燒,卻沒有愛,仍然與我無益.」由此我們可以知道並且明白愛心是何等的重要,當教會失去愛心時,便會受到主的懲罰.
\newline
\newline
第二個是別迦摩的教會──它是一個世俗化的教會,所以神說要用口中的利劍攻擊他們,意思是將教會與世俗分開來.今天很多神學家以為教會應該世俗化；甚至聚會也不需要,只要走到社會裡,服務與貢獻社會便可以了.而且基督教與非基督教應該合一,但他們卻不知道今天的教會大都受社會影響,而不是教會去影響社會.所以求主儆惕我們,使我們不要受社會的影響而走入世俗化之途.
\newline
\newline
第三個是推雅推喇的教會──它充滿了異端的危機.在教會歷史裡,是極其嚴重的對付異端；但可惜在今天人們不再持守純正的信仰了.由異端可能產生了兩個大危機:
\newline
\newline
1.關乎耶穌基督的,他們認為耶穌只是一個人,不過人將祂神格化了.但祂始終是人,所以不能救人脫離罪惡.
\newline
\newline
2.不再承認聖經是完全的,正確的.他們對這真理的持守方面已逐漸的放鬆.求主使我們儆醒,使我們不致陷落異端之中.
\newline
\newline
第四個是撒狄教會──它是一個失去聖靈工作的教會.在末世時代也是聖靈工作的時代,教會如果失去聖靈的工作,便等於身體失去靈魂一樣,不能生存,所以主責備說:「按名你是活的,其實你是死的.」(啟三1)在今天的教會中,很多人害怕提及聖靈的工作,其實我們要進深認識聖經教導我們的真理.盼望今天的教會能更多注重聖靈的工作,在聖靈的真理上不致偏差.
\newline
\newline
第五個是老底嘉教會──它是一個失去主治的教會.耶穌原本是教會的頭,但在教會內他不發生作用,他要在外面叩門.失去主治的教會是一個屬世的民主教會,而我們所需要的是一個屬靈的民主教會,是由聖靈帶領一班屬乎神的人一同事奉神.盼望我們能更加更新,不要效法老底嘉教會,使我們能過著一個讓基督管治的教會.
\newline
\newline
(二)在客西馬尼園內,門徒呼呼沉睡,他們並沒有理會耶穌當時那懇切的禱告,他們更加不知道那時刻的嚴重.在這之前,他們以為夫子很快會得國了,所以在那裡等候.
\newline
\newline
末世時代也是教會黑暗的時代,因為充滿了罪,在(帖前五1-6)說:「弟兄們,論到時候日期,不用寫信給你們,因為你們自己明明曉得主的日子來到,好像夜間的賊一樣.人正說平安穩妥的時候,災禍忽然臨到他們,如同產難臨到懷胎的婦人一樣,他們絕不能逃脫.弟兄們,你們卻不在黑暗裡,叫那日子臨到你們像賊一樣.你們都是光明之子,都是白晝之子,我們不是屬黑夜的,也不是屬幽暗的,所以我們不要睡覺,像別人一樣,總要儆醒謹守.」這段經文說明人會在不知不覺中沉睡了,亦會在不知不覺中跌入罪中,所以我們要儆醒.
\newline
\newline
為甚麼會隨流失去?
\newline
\newline
隨流失去就是自己倚靠自己,不會醒覺,不向神支取力量,於是因自恃而隨流失去.
\newline
\newline
現今教會設有許多訓練班來訓練人,但人卻倚賴自己以為完整的計劃,對自己的俱備充滿信心,於是慢慢沉睡,不向神祈求力量,所以陷入魔鬼的手中.求主儆惕我們.
\newline
\newline
(三)主耶穌單獨在那危機來時作禱告.當主耶穌面臨最大危機的時候,他並沒有運用自己的聰明與恩賜,到處去奔跑,尋求解決的方法.祂只是安靜地尋找一處可以禱告的地方,靜靜的與神交通,支取力量.讀馬可一章我們可以知道耶穌是一個重視禱告的人,在他揀選門徒之前的一夜,他也是整夜的禱告.主耶穌是如何的重視禱告的生活.
\newline
\newline
禱告是甚麼?
\newline
\newline
禱告就是倚靠神,一個願意倚靠神的人,他必定時常禱告,禱告乃是一條通天的橋樑.
\newline
\newline
禱告可以改變教會歷史,又可以改變末世教會.求主幫助我們.
\newline
\newline
中國教會需要復興,首先得復興的,乃是一群願意禱告的人,因為他們已在禱告的功課上下了功夫.
\newline
\newline
有很多人追求不同的恩賜,但我們要求神先給我們禱告恩賜的人,曉得恆切禱告,儆醒禱告,付上代價的禱告,有異象的禱告.
\newline
\newline
馬丁路德改教時,他的工作十分忙碌,但每一天仍抽出二,三小時來禱告,否則那一天便過得毫無意義.很多時候我們愈忙碌愈忽略禱告生活,但禱告是我們支取力量的根源.求主幫助我們,使我們有更多人願意禱告.有一位牧師說:「教會內若有十分一的信徒,參加每星期一次的祈禱會,那麼這個教會便十分成功了.」求主真的幫助我們,使我們願意在禱告上有負擔,願意下功夫,免得我們的主在寶座上孤軍作戰,祈求天父真的幫助保守我們.
\newline
\newline


\newpage

\section{第51屆~港九培靈會~張慕皚~第四講}
\label{sec:686}
\textbf{奉獻生活的更新}
\newline
\newline
連結: \href{https://www.hkbibleconference.org/session-message/view/686}{\texttt{https://www.hkbibleconference.org/session-message/view/686}}
\newline
\newline
\hyperref[sec:685]{< < < PREV SERMON < < <}
~
\hyperref[sec:toc]{[返目錄]}
~
\hyperref[sec:688]{> > > NEXT SERMON > > >}
\newline
\newline
彼得前書 4:7-4:8
\newline
\begin{longtable}{cl}
\hline
\hline
章節 & 經文 (和合本修訂版)\\
\hline
4:7 & \begin{tabularx}{0.7\textwidth}{X} 萬物的結局近了。所以你們要謹慎自守，要警醒禱告。 \end{tabularx} \\ \\ \relax
4:8 & \begin{tabularx}{0.7\textwidth}{X} 最要緊的是彼此切實相愛，因為愛能遮掩許多的罪。 \end{tabularx} \\ \\
[1ex]
\hline
\hline
\end{longtable}
末世教會有一個非常嚴重的危機,很多信徒對救恩的真理都有錯誤的觀念；他們對耶穌基督的救恩只接受一半,就是只讓耶穌成為他們的救贖主,卻沒有成為生命的主.他們只領受主的救恩,但不願意讓主在他們的生命中居首,所以雖然信了主,生活仍和以前一樣,這是非常錯誤的.今天我們要從約翰福音廿一章內看看有關奉獻生活的真理.
\newline
\newline
很多解經家都認為約翰廿一章是附加的經文,並不重要,二十章31節已經是整本書的總結,證明了耶穌基督是神的兒子.但我並不同意這個看法,約翰廿一章非但不是附加的經文,更是不可缺少的.沒有約翰廿一章,就好像羅馬書只有前面的十一章,而沒有十二章和以後的記載一樣.前十一章很清楚的講到,我們接受耶穌做救贖主,同時也要接受祂做生命的主；第六章把這真理講得特別清楚,我們不能再過犯罪的生活,乃要將肢體獻給神作義的器皿.但到十二章的時候,保羅再特別強調奉獻的生活,耶穌基督是生命的主的真理,並勸勉信徒要把身體獻上,當作活祭,意思是把整個人奉獻給神,毫無保留的獻在壇上；為主燒盡,還要心意更新而變化,過一個新的,屬主的生活.同樣,使徒約翰在說過耶穌基督是神的兒子,是救主這些真理之後,便再次將祂要求我們完全跟隨祂,完全被祂得著的真理闡釋清楚,所以第廿一章的記載是同樣重要的.
\newline
\newline
首三節記載彼得帶領其他六個門徒回去打魚,恢復他們以前所過的生活,可是打了一夜,竟連一點收穫也沒有,在這裡,耶穌再次提醒他們,跟隨基督是要專心一致,不能半途而廢的.他們跟隨耶穌,就要將身心完全奉獻與主,生活完全由主管理；如果他們要回頭重過以往的生活,便一無所獲了.門徒原來是主所呼召的,他們和主的關係,是一個主僕的關係,說得正確些,應該是奴僕或奴隸.奴僕在主人面前沒有自己的意思,主人的意思就是他的意思；主人要他做什麼,他便要依足去做,如果做不出來,那只好接受主人的懲罰.我們和主的關係也是這樣,既然跟隨主,應了祂的召,就要聽祂的指揮,主要我們學甚麼,做甚麼；我們便要跟著學,跟著做.
\newline
\newline
耶穌曾經在路加九章六十一,六十二節教訓過他們這個真理,有人說:「主,我要跟從你,但容我先去辭別家裡的人.」耶穌說:「手扶著犁向後看的,不配進神的國.」耶穌並非不近人情,連辭別家人都不允准,祂乃是看到這個人的心,知道他不能完全脫離以前的生活,不能放下俗世的牽掛,所以祂才這樣說的.一個留戀過往,追昔過去的人,是無法完完全全跟從主；就像羅得的妻子一樣,雖然離開了所多瑪,但仍回頭觀望,捨不得丟在那裡的家產財物,結果變了鹽柱.
\newline
\newline
門徒為甚麼起先跟隨主,到現在才變節呢?因為他們經歷不到主的同在.耶穌未死之前,他們天天和主一起生活,現在耶穌雖然在復活後也向他們顯現,但他們總覺得主離他們很遠,因此他們以為主再不會在他們中間,於是便重操故業,打漁去了.耶穌提醒他們,奉獻的生活絕不能這樣,他們既然一次將自己交託給主,便要一直向前奔跑,不能再回顧的了.
\newline
\newline
要保持起初的熱誠,我們便須不斷尋求主在我們生命中同在的真實經歷,那就是說,我們信了主之後,便須把自己的生命交託給主,讓主來管理,切不可被撒旦利用.保羅在羅馬書第六章很嚴重的警告我們,我們信主之後,便當把肢體獻給主,再不能讓罪在我們身上作主,服從身子的私慾,受罪來支配.一個不讓主管理的生命,後果是不堪設想的,撒旦必然利用我們的軟弱,再一次得著我們.因此,我們要時刻警惕自己,看看今日是甚麼來支配我們的生命?是不是金錢,事業,還是我們一些特別的嗜好?我覺得今日有很多人,包括信徒,都為電視機所支配,他們一天到晚的生活,都是環繞電視機的.我不是反對看電視,電視是一個良好的傳播工具,能夠被神使用；其中也有很好的節目值得我們看,但問題是如果我們讓一部電視機來支配控制我們的生活,那就不合聖經的教訓了.我在芝加哥讀神學時,看見一間教會的晚堂崇拜,早在七點鐘便開始,平常教會的晚堂聚會,都是八時開始的,我便問問他們其中的原因,原來禮拜天晚上大家都喜歡早點崇拜,崇拜完可以回家看九時的電視偵探片.到底我們的生活是由甚麼來支配呢?
\newline
\newline
一個只接受耶穌為救贖主,而未過奉獻生活的人,就好像一個女子,結婚不久便離開丈夫,過回獨身的生活一樣.保羅說他是把哥林多的信徒好像貞潔的童女,許配給耶穌基督；我們應該和主過一個非常緊密的生活,怎能信主以後,又回去過以前的生活呢?
\newline
\newline
跟著在第四到第七節中,耶穌吩咐門徒把網撒在船的右邊,就必得著,他們把網撒下去竟拉不上來.這裡提醒我們,一個完全奉獻的生活,最重要的就是事奉.路加五章一到十一節的經文也有類似的記載.門徒整夜勞碌都打不到魚,但耶穌告訴他們撒網在那裡,他們就得著魚了.這不是希奇的事麼?門徒中多是漁夫,是專業打漁的人,不但有自己多年的經驗,也有上一代所留下寶貴的經驗,他們對這海哪裡有魚,和甚麼時候有魚,都瞭如指掌,他們打了整夜都沒有所獲；現在由一個外行人吩咐他們在那裡撒網,這豈不是笑話?但門徒都很順服,因為耶穌是他們的主,當他們聽從主的份咐,不憑自己的經驗時,收穫就來了.
\newline
\newline
事奉的意義也是這樣,屬世的工作和屬靈的事奉完全不同,屬世的工作要講經驗,但屬靈的工作,主要是順服神的引導,經驗是其次的.我們所做的每一個事奉,都要先帶到主面前,求主教導我們怎樣計劃,怎樣去做,我們才去做,這樣的事奉,才是奉獻生活的事奉.單是人的經驗,只會產生傳統,組織,制度,成見,偏見和自高,這一切都會令我們產生一種屬靈的硬化症,就是對神的帶領和聖靈的感動缺乏敏銳,變得反應遲鈍,甚至麻木不仁,這是奉獻生活中一個極大的危機.
\newline
\newline
在耶穌揀選十二個門徒的事件中,我們發現一個很奇怪的事實,耶穌並沒有揀選當時一般受過宗教訓練的熱心人士,祂沒有揀選法利賽人,撒都該人,文士或祭司成為祂的門徒,其實這些人有經歷,而且受過訓練,又熟聖經,如果耶穌揀選了他們,便不須給他們太多的訓練；但耶穌的看法不是這樣,他們雖然好像具備很多優越條件,但他們卻患上屬靈的硬化症,神不能在他們身上下功夫,所以耶穌只得把他們放在一邊.這情形就好像耶穌所說,新布不能補在舊衣上,新酒不能放進舊皮袋裡,因為舊布和舊皮袋都已失去彈性,若把新的加上去,結果只有破爛.不過,新約的保羅卻是主所重用的另一個法利賽人,他雖然受過高深的猶太教訓練,但他的生命仍有彈性,仍有屬靈的伸張力,可以接受主的陶造.從馬其頓異象的這段經文中,我們很清楚看見保羅這種屬靈的敏感,他雖然已經計劃好佈道的行程,但聖靈既然禁止他們在亞西亞講道,他們便立即停止,到了每西亞的境界,想往庇推尼去,耶穌的靈又不許.後來,保羅在異象中看見一個馬其頓人,聖經說他們便「隨即」往馬其頓去.保羅有經驗,也有計劃,但是卻願意順服在神的帶領下來事奉.今天我們在教會的事奉,是先有自己的計劃,組織,然後禱告求神來成全我們的意思；還是我們先禱告,求主帶領我們,讓我們能按祂的旨意來事奉呢?耶穌在世的時候,尚且順服父神的意旨,何況我們呢?
\newline
\newline
最後,耶穌提醒門徒,奉獻的生活是專一愛主的,十五節至十九節記載耶穌三次問西門彼得,你愛我比這些更深麼.耶穌所要求的,不是表面的事奉,好像以弗所的教會,雖然勞碌事奉,卻失去起初的愛心,主不會欣賞.耶穌要彼得在愛祂和愛俗世打漁的事上作個選擇,祂所要的是我們的心,和對祂專一的愛；這專一的愛是出於自願的,不是強迫的,雖然聖經也有這樣的命令,要我們盡心,盡性,盡意,盡力的愛神；可是若非出於自願的愛,是非常痛苦的.有一個故事講到一個婦人嫁了丈夫後,發覺她丈夫是個暴君,對她諸多苛求,管束極嚴,每天上工前寫下一切要她做的事,晚上回來逐樣核對,看她做妥沒有；妻子雖然樣樣照做,但內心極為痛苦.後來丈夫死了,她又重嫁另一位丈夫,這位丈夫非常愛她,體貼她,她發覺現在她為這位丈夫和家庭所做的,比以前為多,但她一點不覺辛苦,反而內心充滿愛和快樂.
\newline
\newline
山森多夫是德國的一名伯爵,有錢有勢,只可惜玩世不恭,一天他參觀一個畫展時,被一幅名畫深深的吸引住了,這是幅耶穌被釘十架的畫；他看見耶穌頭載荊棘冠冕,血從手掌穴中流下來,畫的下面寫著一行字,意思是我已經為你做了這件事,你為我做了甚麼呢?山森多夫深受感動,便在畫前低下頭來,將自己完全獻給主,說:「你這樣愛我,為我死在十字架上,我怎能為自己而活呢?我的財產,家庭以及一切都是你的,求你使用.」這個禱告幫助了近代宣教事業的發展.
\newline
\newline
只有當我們不斷體會主在十架的犧牲和大愛時,我們才會被祂的愛吸引,願意專一的愛祂.只有當我們願意把自己完全擺上,真正過奉獻的生活時,中國教會,以至普世的教會才得以復興.
\newline
\newline


\newpage

\section{第51屆~港九培靈會~張慕皚~第五講}
\label{sec:688}
\textbf{彼此相愛的更新}
\newline
\newline
連結: \href{https://www.hkbibleconference.org/session-message/view/688}{\texttt{https://www.hkbibleconference.org/session-message/view/688}}
\newline
\newline
\hyperref[sec:686]{< < < PREV SERMON < < <}
~
\hyperref[sec:toc]{[返目錄]}
~
\hyperref[sec:690]{> > > NEXT SERMON > > >}
\newline
\newline
未世時代是個靈戰非常激烈的時代,撒旦要傾巢而出作最後的戰爭,因此,我們的生命也需要作全面的更新,好成為神手中合用的器皿.今天我們要思想的,是彼此相愛這方面的更新.我們先來看看幾處有關的經文:(彼前四7-8)「萬物的結局近了,所以你們要謹慎自守,儆醒禱告,最要緊的是彼此切實相愛,因為愛能遮掩許多的罪.」(太二十四10-12)「那時必有許多人跌倒,也要彼此陷害,彼此恨惡,且有好些假先知起來,迷惑多人,只因不法的事增多,許多人的愛心才漸漸冷淡了.」(約十三34)「我賜給你們一條新命令,乃是叫你們彼此相愛,我怎樣愛你們,你們也要怎樣相愛,你們若有彼此相愛的心,眾人因此就認出你們是我的門徒了.」
\newline
\newline
末世時代最需要的,就是在主面前切實的彼此相愛,在彼此相愛的生活中尋求聖靈的更新；因為末世時代,人要彼此陷害,愛心冷淡,且有許多人跌倒,所以我們要特別實行彼此相愛的道理.為甚麼末世時代人的愛心會逐漸冷淡呢?如果我們研究歷史,便會發現工業革命帶來人類生活一個很大的轉變,人類慢慢從農村生活進入大都市的生活中；也慢慢從彼此關懷,守望相助的愛心生活進入彼此傾軋,互相競爭的自私生活去.人與人之間只有利害關係的存在,各人都想凌駕別人之上,教會處在這樣的背景和風氣下,自然也大受影響.除了受到一般性的社會因素影響外,還有一個重要的原因,就是教會過份注重事工和計劃,以致忽略愛心生活的重要.
\newline
\newline
當我們過份注重事工和計劃的時候,便很容易失去與弟兄姊妹之間的關懷和聯繫,彼此沒有接觸,也不知道別人的需要,如何能愛呢?還有,當我們過份注重工作時,也會過份高舉恩賜,把某種恩賜列為是最好的；於是人人都想得到該種恩賜,而忘記了恩賜是聖靈為了教會的需要而個別賜予的；哥林多教會分裂的原因之一,就是過份高舉恩賜.我知道美國有些教會,他們選同工或聘請傳道人之時,要規定他們有某方面的恩賜,訓練,學識,他們看重恩賜過於人屬靈生命的成長,這是本末倒置的.聖經對神工人的要求,主要是他們屬靈的程度和愛心的表現,恩賜是其次的.保羅在哥林多前書十三章也把各樣的恩賜和才幹列於愛之後,即使我們有超然的恩賜,才能,甚至能說萬人的方言,天使的話語,也有先知講道之能,如果沒有愛,一切都歸於徒然.
\newline
\newline
約翰福音十三章告訴我們,教會的特質就是弟兄姊妹能彼此相愛,如果一個教會裡看不見弟兄姊妹彼此相愛,那麼它已經失去教會的本質,不能見證主了；這就如同鹽失了味一樣,沒有甚麼能使它變成鹹的,只有丟在外面,被人踐踏.為甚麼鹽會失去味呢?今天我們用的鹽都是經過煉淨的,但耶穌時代的鹽則在海邊或山地掘出來的,含有很多的渣滓,這些渣滓很多時候和鹽很相似,質地也差不多一樣,不過如果放在廚房太久,受著熱力的蒸發,慢慢便失去鹽的本質,只剩下一些渣滓,一點味道也沒有,根本不能用來調味,只得把它丟掉.教會和信徒也是一樣,如果我們不求主不斷的更新我們,叫我們有彼此相愛的心,恐怕我們也像耶穌所說失去味的鹽一樣,再沒有功用了.
\newline
\newline
現在我們來思想一下聖經所講的愛的特質.聖經中最少有四種不同的愛,第一是男女的愛,是以情感為出發點,所以有時也是盲目的.第二是朋友的愛,這愛是以理智為出發點,別人愛我,關心我,他值得做我的朋友,但如果有一天我發覺他背後罵我,又不守約,那麼我便和他一刀兩斷.當我們計算過他值得我愛時,我才愛他.第三是骨肉的愛,這愛以天性為出發點,神造人時,已經把這愛放在人的天性內,不管他值不值得我愛,只因彼此是骨肉,我們之間便有愛的存在.最後一種是從上而來的愛,是超越人間一切的愛.這愛是以我們的意志為出發點,由意志開始而達致我們的全人.
\newline
\newline
這種從上而來,以意志為出發點的愛,正是主耶穌所要求的.耶穌教訓我們要愛仇敵,為逼迫我們的禱告,這不是情感所能達成的,只有堅強的意志才能幫助我們去愛.雖然某弟兄或姊妹不可愛,但既然主耶穌所賜的命令是彼此相愛,我們便得用意志為出發點,立志將好處帶給他,當我們遵照主的命令,實行愛弟兄姊妹時,我們便發現弟兄姊妹可愛的地方了.在好撒瑪利亞人的比喻中,他所幫助的是他們民族的世仇猶太人,在情感來說,他絕對不會愛他,但他所表現的愛,乃是由意志出發的,因此他能關心他,把好處帶給他,這是耶穌命令我們要追求的那種愛心.
\newline
\newline
今日的世代,最缺乏的是愛.美國哈佛大學一位社會學家所羅基,致力研究社會的需要,他發現在人的本性裡面,每個人都需要愛,沒有愛人不能過活,沒有愛人所過的是悲慘的生活,因為愛是神放在人裡面一種屬靈的維生素.在他的著作中提到一件這樣的事實,有些科學家用人來做試驗,他們把卅多個被父母遺棄的嬰兒送進收容所去,給他們最好的食物,最好的照料,和最好的生長環境,然後觀察他們的生活,這些嬰兒因為失去父母的愛,都變得食慾不振,慢慢萎縮,最後死亡；有些雖然長大,但都變成神經有問題或心理不健全的人,這是缺乏父母的愛的緣故.這位社會學家研究的結果,就是我們需要產生一些愛,把愛好像原子能一樣產生出來,可惜他只是個研究社會問題的人,他不認識神,不曉得神是愛的源頭,所以至終無法找到一樣能產生愛的動力,結果他失敗了.
\newline
\newline
世人既然這樣迫切的需要愛,為甚麼在神的家,在教會裡面沒有愛呢?我們為甚麼表現不到彼此相愛的生活呢?我想有很多原因,我只提出幾個:
\newline
\newline
第一,我們在教會中找不到愛,原因是我們只講愛神,但沒有愛人.聖經說,如果我們不愛看得見的弟兄,就不能愛看不見的神；只講愛神而不愛人,是不合神的教訓的.
\newline
\newline
第二,我們的愛沒有實際的行動,只在我們的心裡,在我們的禱告裡,但絲毫沒有行動,這也是不對的.耶穌教導我們要愛仇敵,為他們禱告,這就是行動；愛心若沒有具體的表現,就如雅各所說信心沒有行為一樣,是死的,沒有用的.
\newline
\newline
第三,我們往往只愛遠方的人,而忽略近處的人.愈遠的人愈容易愛,我們可以為非洲,南美洲,澳洲的人迫切禱告,甚至流淚的禱告,但卻不知道怎樣愛自己教會內的弟兄姊妹.有一個故事,可能也是事實,說美國有一所教會,提倡白種人和黑種人各有自己的教會,一個禮拜天早上,一位黑人青年穿著筆挺的西裝,到白人的教會聚會；在座的白人看見,個個怒目而視,終於找兩個兇猛而有力的人把他押出去.那個黑人青年莫明其妙,便問他們為何這樣待他；接著他道出自己來這教會聚會的原因,原來有白人宣教士到非洲傳道,使他得聞福音,所以他心中一直希奇白人為甚麼會有這麼大的愛心,願意到非洲大陸去；如今因為有機會到美國唸書,所以特地到白人教會看看.這真是個多麼大的諷刺.
\newline
\newline
到底聖經教導我們彼此相愛的原則是怎樣呢?我提出四個可以實行的原則.第一,聖經說我們愛,因為神先愛我們；又說愛神的,也當愛弟兄,這表示彼此相愛和我們愛神有很大的關係.人與人之間有問題,人際關係惡劣,主要是我們與神之間有問題,有間隔,沒有愛的聯繫,因此便不能愛弟兄姊妹.這好比一個三角形 ABC,頂點 A 代表神,而 B 和 C 則代表人,如果 B 和 C 愈靠近 A 時,BC 之間的距離便愈近；照樣,如果我們和神的關係愈密,和人的關係也自然愈好.第二,耶穌教導我們要愛鄰舍如同自己,將自己放在別人的立場上,為別人著想,愛別人好像愛自己一樣,我喜歡別人怎樣待我,我便照樣去待別人.另一個比較進一步的原則,就是腓立比書二章三節所說:「凡事不可結黨,不可貪圖虛浮的榮耀,只要存心謙卑,各人看別人比自己強.」人都有一個弱點,就是喜歡看別人的短處,弱點,所以不能愛人.我們若要愛人,就得欣賞別人的長處,看別人比自己強.亞伯拉罕和他姪兒羅得分家的時候,他並沒有因為自己是長輩,便先行選擇土地；他乃是看羅得比自己強,當他是長兄或者前輩般看待,請他先選擇肥美的地,結果神大大的賜福他.最後一個原則,也是更進一步的原則,就是馬太福音廿五章四十節王所說的話:「我實在告訴你們,這些事你們既作在我這弟兄中一個最小的身上,就是作在我身上了.」這原則就是教導我們對待別人,好像對待耶穌一樣,如果我們真能這樣,教會的問題,糾紛便自然大大減少了.保羅在以弗所書教導一班作奴僕的,要甘心服事主人,好像服事基督一樣.今天,如果我們接待弟兄姊妹好像接待耶穌基督一樣,我們便能真正實行所命令的彼此相愛的生活了.
\newline
\newline


\newpage

\section{第51屆~港九培靈會~張慕皚~第六講}
\label{sec:690}
\textbf{崇拜生活的更新}
\newline
\newline
連結: \href{https://www.hkbibleconference.org/session-message/view/690}{\texttt{https://www.hkbibleconference.org/session-message/view/690}}
\newline
\newline
\hyperref[sec:688]{< < < PREV SERMON < < <}
~
\hyperref[sec:toc]{[返目錄]}
~
\hyperref[sec:691]{> > > NEXT SERMON > > >}
\newline
\newline
今天我們要思想有關崇拜生活的更新,先請看兩處的經文:(詩九十六9)「當以聖潔的妝飾敬拜耶和華,全地要在祂面前戰抖.」(約四24)「神是個靈,所以拜祂的,必須用心靈和誠實拜祂.」
\newline
\newline
敬拜神是人的本能,是神在造人的時候放在人本性裡面的,這本能就是當人接觸到神的永能和神性時,一種自然的流露,因此每個人都知道應該怎樣敬拜神.只可惜當人犯罪墮落後,人的本性便受到虧損,他雖然仍有敬拜神的本能,但卻失去敬拜真神的能力；所以便將這種本能轉移去敬拜偶像,將昆蟲,禽獸等當作神一樣的來敬拜；或是敬拜自己,以自己為神,這是人類今日極可憐的光景.
\newline
\newline
但當我們蒙神的恩典,接觸過神,在生命中遇見過神之後,我們必然知道怎樣敬拜祂.就如亞伯拉罕在幔利橡樹遇見神之時,立即俯伏向神下拜；又當以色列民在西乃山看見神顯現之時,都非常震驚,從震驚中產生一種對神的敬畏,在神面前戰抖害怕,因為看見了神的威榮.新約時代也是一樣,當西門彼得遵照耶穌的吩咐,在水深處撒網打魚,又大有所獲時,便立即跪在耶穌面前敬拜祂,因為他看到了祂所顯的神蹟,知道祂不是個普通的人,乃是神.
\newline
\newline
既然人遇見了神便知道怎樣敬拜神,為甚麼今天教會崇拜的光景這樣冷落,弟兄姊妹不願意到教會敬拜神,又不曉得怎樣敬拜呢?原因是和我們屬靈的生命有關,或許是我們未曾真正遇見過神,又或者我們的生命太過軟弱,以致我們不能按神的心意來敬拜祂.我提出幾個教會常見的情形,有些可能是西方教會比較顯著的.
\newline
\newline
第一種現象,就是很多人將崇拜變成團契；我在美國曾參加過一個崇拜聚會,沒有任何秩序,儀式,大家只是聊聊天,唱一些愛唱的詩,有些弟兄姊妹分享一些見證,這就叫做講道,他們把崇拜和團契混亂了.
\newline
\newline
第二種現象,是把講道變為一種談話.很多西方的年青人落在一個反權威的狀況中,他們說人是平等的,所以不應該讓牧師在台上教訓他們；因此傳道人要迎合會眾的需要,不能大聲宣講神的話,只能細聲地和他們談話.
\newline
\newline
第三,很多教會將講道變成對話.這是因為會友希望先是牧師在台上說,他們也要有反應,有參與,所以牧師講道時,會眾可以隨時發問,或是在講道後有一段發問的時間,讓會眾也可以表達他們的意見.
\newline
\newline
第四,在崇拜聚會中加插娛樂性節目,以此吸引會眾；他們等約一些明星,球星來表演某些節目,以迎合會眾的需要.
\newline
\newline
第五,也是中國教會中常見的現象,就是以崇拜為我們得著屬靈福氣的途徑；所以赴會的唯一目的便是要享受崇拜,聽道必須有得著,詩歌必須有感力,否則他們便不到這處聚會,改到別處能令他們有「得著」的地方去.以上所提到的五種情況都是不正確的,主要的原因是我們將神的位置和人的位置調換了.
\newline
\newline
聖經告訴我們,我們所敬拜的對象是神,我們崇拜的目的,是讓神有所得著,把祂應得的頌讚,尊榮歸給祂.當然,我們也希望在聚會中有所領受,但我們主要的目的,乃是讓神有所得著.其實,教會有各種不同類型的聚會,比方團契聚會是幫助弟兄姊妹互相溝通︴聯繫,彼此在主內建立；主日學可以幫助信徒認識真理；一些教導性的聚會或講座也是叫我們在真理方面得著造就；但崇拜聚會卻是為神的.聖經雖然沒有清楚講明崇拜應有的形式,但我們從初期教會的崇拜中,看見兩個主要的項目,就是講道和守主的聖餐,兩樣均以主為中心,講道是宣講主的話,而守聖餐則為了記念主.
\newline
\newline
在改教時期,馬丁路德和達爾文都非常注重講道,他們根據聖經逐章逐節的慢慢解釋,把神的話完整而有系統的表達出來.至於守聖餐,初期教會則非常著重,在聖餐之前,還先享受愛筵,就是各人帶點食物來一同享用,愛筵之後便舉行聖餐.如有吃不完的食物,便在聖餐後分給貧苦的弟兄姊妹,或是主外的朋友,使他們也能共享,這是非常有意義的.
\newline
\newline
今天,我們真需要更新我們的崇拜生活,剛才看的兩節經文便提醒我們幾個重要的原則,
\newline
\newline
第一,敬拜時,我們的心靈和儀表應該是合一的.耶穌教訓我們要以心靈和誠實去拜父,神所著重的,是我們心靈的敬拜,但並非說儀表一點也不緊要.耶穌特別強調心靈的敬拜,主要是和舊約在聖殿裡單重儀表的敬拜作一對照,舊約時代他們獻上一些動物的祭牲作為敬拜,這只是影兒,到新約時代耶穌完成救恩後,我們的敬拜便進入敬拜的實體裡去,摩西所訂立種種敬拜的律例和制度,已不再需要,今天我們是以心靈來直接敬拜神.法利賽人只有敬虔的外貌,而無敬虔的實意,他們的禱告,禁食,都是做在人前,為要得人的稱讚,耶穌非常嚴厲的斥責他們.今天,很多人雖然在崇拜聚會中唱詩,禱告,但他們的心卻牽掛俗世的事,或是看見某個和他不合的弟兄或姊妹,於是整堂聚會都想著別人的不是,這種只有外表而沒有內心的敬拜,是神所憎惡的.所謂有諸內必形於外,我們外表的舉動,儀表往往就是反映我們內心的態度.不過,有些教會卻走了另一個極端,就是完全不著重聚會的形式；曾經有一位神學生作見證說,他是睡在床上做靈修的,因為他嘗試過很多靈修的方式,都覺得不好,最後他發覺躺在床上讀聖經祈禱是最好的.當然,敬拜神是用我們的心靈,但正確的儀表也是需要的,比方我們在崇拜聚會時的姿勢和服飾,也應配合我們內心的敬虔.
\newline
\newline
第二,敬拜的行動和聖潔的生活應該合一.神不悅納該隱的供獻,不是因為神不喜悅他所獻上的田裡的出產,而是因為他的生活不夠聖潔,不能討神的喜悅.神在以賽亞書一章十三節那裡講了一番非常嚴厲的話,神不喜悅他們的祭物和朝拜,因為他們犯罪作惡,除非他們洗濯,自潔,除掉惡行,過聖潔的生活,這樣,他們在神殿裡的敬拜才得神的喜悅.很多信徒都有一種將功贖罪的心理,以為禮拜天到教會崇拜,便可補回六日生活的不聖潔；然而神所看重的,乃是我們每日都過聖潔的生活,而不是單在主日的崇拜.假如我們每日不過聖潔的生活,神也不會悅納我們在主日的敬拜.
\newline
\newline
第三,崇拜的心靈和真理的認識應該合一.耶穌說:「神是個靈,所以拜祂的,必須用心靈和誠實拜祂.」中文聖經所翻譯的「誠實」有三種不同的解釋,第一,是指一種真誠懇切的態度.第二,是指我們進入敬拜的實際裡,不像舊約時代的敬拜,只是一些預表,影兒.第三,也是含意最廣的,就是指真理.那就是說,我們敬拜神,要按照真理的教導,用心靈來敬拜.如果我們對神真理的認識有缺乏或錯誤時,我們的敬拜便不能討神的喜悅.
\newline
\newline
舉例來說,有些異端邪派不信三位一體的神,他們說神只有一位,耶穌不過是一位聖人先知,而聖靈是從神發出來的一種能力,是沒有位格,沒有性情的.他們對真理的看法有這樣大的偏差,因此,即使他們有很美麗的廟宇,教堂,又真心誠意的敬拜神,神也不能悅納.另外有些人不信聖經是神的啟示,是神真確無誤的聖言,他們對聖經的看法是這樣,如果他們讀到某些經節是對他們有幫助的,那些經節便變成神的話,其他沒有幫助的經節,便僅是以前的先知和使徒所寫的話語,持這樣看法的人,他們的敬拜能不能討主喜悅呢?故此我們今日敬拜神,必須按照我們心靈的虔敬,並加上我們對真理的認識來敬拜.
\newline
\newline
我們對神的認識怎樣呢?這和我們的敬拜很有關係,你心目中的神是怎樣的一位神呢?有人或者以為祂是一位慈祥的老公公,有求必應,從不會生氣的；另外有人或許認為祂像一位嚴厲的警察,時刻偵查我們,一有過失便會懲罰我們,所以我們在神面前充滿懼怕；其實兩種看法都有偏差,如果我們抱著這樣的觀念去敬拜神,又怎能討祂的喜悅呢?因此我們必須以聖經作為我們信仰的基石,使我們對神有正確的認識,但這些知識不能單停留在我們腦海中,更要進入我們的生命裡.只有這樣,我們對神的敬拜才能蒙祂的悅納.
\newline
\newline


\newpage

\section{第51屆~港九培靈會~張慕皚~第七講}
\label{sec:691}
\textbf{人價值觀的更新}
\newline
\newline
連結: \href{https://www.hkbibleconference.org/session-message/view/691}{\texttt{https://www.hkbibleconference.org/session-message/view/691}}
\newline
\newline
\hyperref[sec:690]{< < < PREV SERMON < < <}
~
\hyperref[sec:toc]{[返目錄]}
~
\hyperref[sec:692]{> > > NEXT SERMON > > >}
\newline
\newline
保羅在羅馬書十二章三節說:「我憑著所賜我的恩,對你們各人說,不要看自己過於所當看的,要照著神所分給各人信心的大小,看得合乎中道.」到底我們用甚麼來衝量自己,衡量別人,才合乎神的標準呢?撒母耳是個神所喜悅的人,又是個非常屬靈的先知,但當他要從耶西的眾子中挑選其中一個為王時,他的眼光也有偏差,所以神提醒撒母耳說:「不要看他的外貌,和他身材高大,我不揀選他；因為耶和華不像人看人,人是看外貌,耶和華是看內心.」這是神對人的價值觀.
\newline
\newline
自有人類以來,人就不斷探索「我究竟是誰?我從哪裡來?我要往何處去?」這些問題的答案,直至今日,一般的哲學家和科學家對人主要有三種看法:
\newline
\newline
第一,人不過是一件物件,和石頭,椅子,鋼琴等物件無異；所分別的就是人是一個活動的物件,因此人類的價值被貶得非常低.
\newline
\newline
第二,人似一部機器,特別是現代新發明的電腦人,更酷似人,人會說話,電腦也會說話,且有超人的記憶力,人會精神崩潰,電腦也會出錯,因此人只不過是一部機器而已.
\newline
\newline
第三,人是動物,自從達爾文開始提倡進化論之後,人類便喜歡以進化論來解釋自己的來源；可惜時至今日,講進化論的人,一點也不知進化的過程和方法,便相信人是由猴子而變成的,最近澳洲的一位科學家在其研究結果中,卻發現人不是從猴子而變,乃是猴子從人而變,那就是說猴子比人更進一步.當人遠離神,不認識神的時候,人便無法認識自己；只有當人認識真神時,才能認識自己,知道自己的價值.所以不認識神的人,往往對人的價值看得極低,他們以屬世的物質,屬地的眼光來衡量人,撒母耳雖然是神人,也犯了同樣的毛病,只有神是從屬天的角度來衡量人,認識人真正的價值.
\newline
\newline
在傳統方面,我們對人的衡量方法最少有四個.
\newline
\newline
第一是按他的外貌,好像撒母耳立王的時候一樣,他看見大衛的長兄身材魁梧,孔武有力,便憑他的外貌以為他可以為王.聖經裡面還有其他的例子,就如保羅,因著樣貌不好而遭受很多委屈,哥林多的信徒看不起他,因為他樣子醜陋.根據歷史的記載和描述,保羅的個子非常矮小,兼且禿頭,駝背,彎鼻,眉頭深鎖,雙腳也是彎曲的；只是他的身上卻充溢著神的恩典.然而哥林多教會卻有一位儀表出眾,風度翩翩的人物,名叫亞波羅,若論外表儀容,保羅是望塵莫及的；因此亞波羅不但惹人好感,很多人甚至崇拜他說:「我是屬亞波羅的.」保羅為此非常憤怒,因為竟然有人覺得他樣貌不好,不值得稱為使徒,因此他在信中極力為自己使徒的職分來辯護.又有一次,當保羅和巴拿巴在路司得地方傳福音時,因為能夠醫病,所以有人便以為他們是神,把祭物獻給他們,他們稱巴拿巴為丟斯,就是主要的神；稱保羅為希耳米,就是次一等的神,這是他們按保羅和巴拿巴的儀容來分的.由此可見,人對儀表看得多麼重要,甚至有人因為自己生得醜陋而產生自卑,或者是千方百計的去整容,改型,務使自己能有好的外貌.可惜這並不是神衡量人的方法.
\newline
\newline
第二個衡量人的方法是按他的職業,以前的封建時代,人的身價在乎他是否出自名門望族,如果是,即使生得蠢鈍,樣子醜陋,無才無德,但身價依然極高.相反來說,如果出身寒微,就算是一表人才,聰明絕世,也未必受人重視.現今雖然不是封建時代,但人的看法,想法依然一樣,我們慣用職業來衡量人,認為從事某種行業的人,身份會較高,也被人器重珍惜；因此,不少人不惜一切要爬上高職位,來表示自己的重要.這種對人的價值觀也是錯誤的.
\newline
\newline
第三,人往往用財富來衡量別人,一個人收入的多少就決定他身份的高低,薪水少的人會覺得自己的價值很低,比不上那些薪水高的人,這是錯誤的.雅各書第二章講到基督徒的圈子內也出現這樣的情形,信徒以外貌待人,那些雍容華貴,珠光寶氣的便坐在好位上；窮困潦倒,骯髒破爛的便站著,或坐在別人腳凳下,這是神所不喜歡的.
\newline
\newline
第四種衡量方法,是以人的才幹為標準,才幹大的等於價值高,才幹小的價值低,沒有才幹的,便是沒有用了.因此一些比較平庸,資質又差的人,和那些年紀老邁,不事生產的人,便被人看輕,甚至厭棄,這對他們是非常不公平的.
\newline
\newline
神衡量人的標準絕不是這樣,耶穌基督就是最明顯的例子.如果我們用上面四種尺度來衡量耶穌,耶穌是絕無價值可言的.論外表,以賽亞先知形容耶穌是無佳形美容的,我們看見祂的時候,也無美貌使我們羨慕祂；論職業,耶穌出身寒微,只是木匠的兒子；論財富,耶穌異常貧苦,經常不名一文,納稅的時候,還須靠從魚口中取銀,他連枕首棲身的地方都沒有,在十架上,別人把祂僅有的衣服都拿走；至於才幹,耶穌是神的兒子,大有才幹,能力,但祂並沒有使用這些來拯救世人.祂雖然有權差十二營的天使來救祂脫離十架之苦,但祂沒有這樣做,祂乃是用愛來拯救世人.
\newline
\newline
究竟神用什麼尺度來衡量人呢?我們先從創造的角度來看,神造每一個人都是平等的,祂是照自己的形象來造人,雖然人犯罪,把神的形象損毀,但人在神面前仍有價值,因為人是祂親手造的.特別是神的兒女,是蒙耶穌基督的寶血買贖回來的,就如保羅所說:「你們受洗歸入基督的,都是披戴基督；並不分猶太人,希利尼人,自主的,為奴的,或男或女,因為你們在基督裡都成為一了.」另外,神是以我們生命成熟的程度來衡量我們.神的心意是要我們追求認識祂,我們與主的關係,就好像枝子和葡萄樹一樣,是非常緊密的.神賜我們有祂自己的生命,這生命是能與神交往,與神合一,是豐盛的生命.保羅在信主之後追求的目標是什麼呢?他說:「只是我先前以為與我有益的,我現在因基督都當作有損的；不但如此,我也將萬事當作有損的,因我以認識我主基督耶穌為至寶,我為祂已經丟棄萬事,看作糞土,為要得著基督.」今日,神要用我們與祂的關係來衡量我們的價值,到底我們追求認識神到怎樣的地步呢?神所看重的,是那些追求完全被祂得著的人.神所揀選的器皿也是這樣,並不在乎他的恩賜,才幹,學識,乃在乎他屬靈的生命,與主的關係.
\newline
\newline
另一個神衡量人的標準,就是在神的價值觀裡面,我們屬靈生命的長進比恩賜更重要.哥林多教會是個非常追求屬靈恩賜的教會；各樣恩賜都有,但是他們卻因為過份高舉恩賜的重要,以致教會發生分裂.保羅在林前十三章便再次強調神對人的價值觀,說明愛比其他一切更重要,因為愛是聯絡全德的；一個有愛心的人,就願意背起十架跟從主,以愛心捨己為人,像耶穌基督一樣.今日,我們要追求屬靈生命的長進,能夠結滿屬靈的果子,能夠有愛的表現.神所能使用的,是一班有美好的靈性,又加上恩賜的人.聖經描寫大衛是個有美好靈性合神心意的人,而且他也有好的才幹,善戰,能彈琴,作詩,具備多項恩賜.不過,神揀選他,不是因為他有才幹,恩賜,而是因為他有美好的靈性.所羅門,但以理也是一樣,教會歷史中的奧古斯丁,馬丁路德,加爾文,衛斯理,芬尼等,都是有美好的靈性又加上恩賜的人,他們在神手中被神大大的使用.在某些情況下,如果神找不到兼具有美好的靈性和高超的恩賜的人；神寧可以人的靈性為先,使用那些有美好的靈性,而恩賜較差的人,而不使用那些有好的恩賜而靈性低的人,這是神揀選和衡量人的標準.今日,教會的領袖,團契的導師,都應以靈性為追求的首要目標.我曾問一所神學院的教務主任,神學教授應具備些甚麼條件,這神學院很有名氣,所請的教授都是很有學識的,但這位教務主任回答我說,學校請教授,不是先看他有多少個學位,而是先看他有多愛主,跟主是否有好的關係,這也是神的標準.
\newline
\newline
最後,神看人的忠心,勝過他工作的果效.當然,每一個屬主的人都要結果子,努力工作,但神看重的,不是教會的人數,不是我們工作的量,而是看我們是否有忠心；你以神所給你的一千銀子,努力去發展,強如那領五千銀子但懶惰不作工的人；因為神所看重的是我們的內心,是我們的生命,是我們的忠心.
\newline
\newline


\newpage

\section{第51屆~港九培靈會~張慕皚~第八講}
\label{sec:692}
\textbf{時代使命感的更新}
\newline
\newline
連結: \href{https://www.hkbibleconference.org/session-message/view/692}{\texttt{https://www.hkbibleconference.org/session-message/view/692}}
\newline
\newline
\hyperref[sec:691]{< < < PREV SERMON < < <}
~
\hyperref[sec:toc]{[返目錄]}
~
\hyperref[sec:1383]{> > > NEXT SERMON > > >}
\newline
\newline
今天我們要藉著以斯帖和她所處的時代,並她所作的事情,來思想今日中國信徒所應具備的時代感.猶大人被擄之後七十年,神恩待他們,讓他們有機會回國,重建聖殿,並耶路撒冷的城牆；只是能夠回國重建家園的猶大人數目極少,大部份仍留在波斯的地方.以斯帖記便是記載這班留在波斯地的人和他們的遭遇.今天我們集中思想該書卷的第四章7至17節的這段經文.
\newline
\newline
首先我們看見以斯帖是處身於一個絕境的時代裡面.以斯帖是住在王宮裡的,她不知道外間所發生的事,也不曉得同胞正是面臨絕境；因為奸臣哈曼要將末底改並所有波斯地的猶大人滅絕.今天,教會以外的人,都像當時的猶大人一樣面臨絕境,只是我們躲在教會內,不知道外面的人的需要.
\newline
\newline
到底這個時代的人,對自己的看法是怎樣呢?是否像以斯帖時代的猶大人,感到自己毫無希望,要渴慕得著救恩呢?如果我們看人類近代的歷史,便知道人類一直到中世紀的時代,他們的生活,世界觀,都是以神為中心；但其後,人文主義抬頭,人開始以自己為中心,以自己來衡量宇宙萬物一切；後來又有科學主義,將人慢慢帶到無神的生活去.二十世紀又有一種叫做實存主義的思想,專門研究和關心人生的問題.他們要找出人生有甚麼意義,人為甚麼要生存等問題的答案,這思想對現代人的影響極大.
\newline
\newline
我們舉些例子來看看教會外的一班人,怎樣落在一個絕望的光景中.實存主義的鼻祖尼采,提倡人類要進化到超人的地步,但他對人生有極大的失望,覺得人生全無意義；就如一個人行鋼錢,行到中央時,被人推了一把,倒在地上而跌死,人生也是這樣充滿危險,了無意義.另一位實存主義作家瓊烏,形容人生就像處身在一個流行瘟疫的城市內,雖然為著蔓延的疫症焦急發愁,但是卻無法倖免,只有死路一條.又有一位作家形容人生像一堆垃圾,骯髒,無意義.小說家海明威在一本著作中說,人生根本無可救藥,死是解決一切不幸事情的最有效方法.另一位劇作家說,人生唯一的意義就是死,它在寒冷的冬夜中一張溫暖的毛氈一樣,所以讓我們來愛它,不要憎恨它.以上就是現代人對人生的看法和評價,他們發出絕望的呼聲,非常可憐.
\newline
\newline
昔日的猶大人也是一樣,面對絕種的困難,救助無門,王后以斯帖雖是猶大人,卻沉醉在王宮享樂的生活中,對外間同胞的痛苦毫不知情,有人到王宮將這消息告訴她,但她覺得自己沒有辦法,無能為力.今天,教會有聽見世人絕望的呼求嗎?我們在教會的溫室內,享受一切屬靈的福樂,卻忘記外面的人,沒有一點時代感.因為我們忘記神怎樣帶領我們.神的帶領有三方面:第一是藉聖經的話來告訴我們處身在一個怎樣的時代裡面.第二是用教會的歷史,藉前人的經驗來讓我們知道時代的需要.第三是藉聖靈在我們周圍的環境工作,給我們看見時代的需要,讓我們知道自己的使命.很多時候,我們忽略聖靈在我們四圍環境所做的工作,以致看不清神所交託我們的使命,就如以馬忤斯路上的門徒一樣,因著主的死而憂愁悲傷,眼睛模糊,看不見復活的主就在他們身邊.又如以利沙的門徒,只見敵軍重重包圍,而看不見神的火車火馬,正四面護衛他們.今日,我們屬靈的眼睛也太過模糊,看不見神的工作,以致我們過一個毫無意義,因循苟且的生活,在屬靈方面失去使命感.美國人很喜歡搭順風車,他們把想去的地方寫在一個牌子上,然後掛在身上,在公路旁等順路的人接載他們.但有一次,我發現有一個牌寫得非常特別,他的目的地是「任何地方」(anywhere).很多時候,我們的生活也像這個人所掛的牌一樣,沒有目的,隨波逐流,失去神給我們的時代使命感.
\newline
\newline
王后以斯帖就是這樣,自己可以在宮廷內享福便算,毫不顧同胞們的苦情,後來末底改便警惕以斯帖,告訴她今日她有皇后的身份,完全是神特別的恩賜,神既然給她這樣特殊的地位,就是為現今的機會.今日,神也是將我們放在特殊的地位上,為要叫我們能幫助在絕境裡的人,但很多時候,我們不單沒有以自己的身份為榮,反以作基督徒為羞恥；有人怕帶聖經到禮拜堂去,因為不想人知道他是基督徒.求主恩待我們,我們要珍惜我們基督徒的身份,聖經說我們是大祭司,是在人和神之間做溝通的工作,我們的職責是傳福音,勸人與神和好.特別在這末世的時代,神把傳福音的責任放在中國人身上；今日最多失喪的靈魂是中國人,而最能把福音傳給他們的,就是中國的信徒,是華人的教會,是我們這班人.在普世的華人教會中,香港教會的地位又顯得特殊,我們看見,香港有很多神學院,屬靈機構,有神很重用的僕人信徒；他們可以把教會做得好,把福音傳得完備.而且信徒中,很多又有特別的恩賜,青年的弟兄姊妹又多,其中不少更有奉獻的心志,願意獻身事主,這是香港教會特殊的條件,我們有否為此感謝神?並更覺自己責任的重大?
\newline
\newline
其實,神在這時代已經藉著各方面預備人心,但問題卻在於我們,今日傳福音最大的攔阻不是這班不信的人,而是我們.就如尼尼微傳福音最大的攔阻,不是尼尼微城的人,而是神的僕人約拿；我們有太多的藉口推搪,不願到外邊去,只沉醉在教會屬靈的享受中,像以斯帖沉醉在王宮的享受一樣.研究教會增長的人說,今日教會增長緩慢的一個現象,就是信徒和未信主的人的接觸愈來愈少,於是慢慢聽不到外間的呼聲,不知道外間的需要,這是我們需要警惕的.
\newline
\newline
今天,神已經在中國教會中興起教會增長的異象,差傳的異象,但這只是個開始,還須我們積極去實行.很多人問,福音還未傳遍香港,我們就往外面傳嗎?初期教會給我們一個很好的榜樣,安提阿教會雖然是個新成立的教會,但他們沒有等福音傳遍安提阿,便立即差派保羅和巴拿巴出去,做宣教的工作.由此可見,差傳的工作是極需要的.另一方面,當我們做差傳工作的時候,最要緊的不是金錢,而是禱告和工人.禱告是最要緊的,禱告能帶來復興,能產生屬靈的能力.耶穌對門徒說:「要收的莊稼多,作工的人少,所以你們當求莊稼的主差派工人出去,收祂的莊稼.」禱告是第一重要的,只有禱告才能開始差傳的工作,其次是工人的差派.神要在這時代興起許多合祂使用的人,作傳福音的工作.過去這幾年內,神在北美的青年中間做很厲害的工作；很多從香港或台灣去的大學生,在那裡信主,蒙召,將自己奉獻給神,每年都有幾百青年人獻身,樂意事奉主的.我們應該珍惜他們的奉獻,好好帶領他們走上事奉的道路.香港教會也有青年獻身的好現象,我們應該對這些弟兄姊妹多加鼓勵,教導,讓他們能成為神合用的工人.
\newline
\newline
耶穌設過一個比喻,說一個被鬼附著的人,鬼雖然離開了他,但找不到歇腳之處,看見原來那人的地方打掃得非常乾淨,便去帶了七個比自己更惡的鬼來,進到那人裡面.今日神藉著種種災難,已經把人的心打掃乾淨,這些人心靈空虛,正等待神的福音,但如果我們沒有這樣的使命感,撒旦便很容易乘虛而入,擄掠人的心靈了.不久以前,我看見一本美國雜誌中有一篇文章,論到現代人渴慕救恩；這雜誌不是屬靈刊物,但卻指出人心何等渴慕神的救恩,他舉了一個例子,說洛杉磯有一個人,自小有嚮往宗教的心,但可惜找了多年,仍找不到一個能滿足他心靈的宗教；後來他有一個醒悟,覺得不需要繼續尋找,因為他自己就是神,於是便創立一種宗教,勸人來報名.今日人類心靈苦悶空虛,只有耶穌基督的福音才能滿足人心的需要.
\newline
\newline
當以斯帖聽完末底改所講的一番警惕的話時,便有所醒覺了,她請猶大人為她禁食三日三夜,然後她違例去見王,倘若死就死吧.今日,我們聽見世人絕望的呼求時,我們能否對主說:「主阿,我在這裡,請差遣我.」我們不要像約拿那樣,被神用苦難管教之後才願順服；乃要像以賽亞一樣,對神的呼召有即時的回應.但願我們各人都在主面前,盡上傳福音的責任.
\newline
\newline


\newpage

\section{第51屆~港九培靈會~楊濬哲~序}
\label{sec:1383}
\textbf{序}
\newline
\newline
連結: \href{https://www.hkbibleconference.org/session-message/view/1383}{\texttt{https://www.hkbibleconference.org/session-message/view/1383}}
\newline
\newline
\hyperref[sec:692]{< < < PREV SERMON < < <}
~
\hyperref[sec:toc]{[返目錄]}
~
\hyperref[sec:693]{> > > NEXT SERMON > > >}
\newline
\newline
蒙　神恩領,港九培靈研經會自一九二八年創立至今,經已過了半世紀.在這漫長的歲月中,神藉著祂所重用的僕人使女,每屆傳達祂的信息,叫罪人得拯救,叫信徒得造就,叫教會得復興.這是聖靈奇妙的作為,我們實在不能不異口同聲讚美神,榮耀歸給神.除了感謝神外,也要感謝海內外教會的關懷與代禱,尤其是港九教會的關懷與代禱；他們於每年七月間會期前,都分別在各教會同心合意的,在神面前舉手禱告.我深深覺得,這是本會蒙神賜福最大的因素.當然在精神上 的支持,與金錢上的奉獻,也是佔著重要的一環.
\newline
\newline
一九七九年第五十一屆大會經已如期於八月一日舉行.八月初本港天氣異常炎熱,且值颶風襲港,大雨滂沱,以致大會的翌日無法舉行.然而在那一天,仍有不少兄姊頻來電話,詢問是否如常舉行；也有些兄姊,冒險前來大會場地同心禱告.大會的第三天繼續舉行,直至十日畢會.每日赴會人數極其擁擠,尤其是晚上堂內,堂外,副堂,會客室,詩班室,均告滿座.於是又再開放二樓,三樓的主日學課室,還是滿座；而向隅的兄姊,只有在禮堂中席地而坐.大家「聚首一堂」,諦聽真道.這種飢渴要主的心,使人受感甚深!
\newline
\newline
本講道集是第五十一屆的講道記錄.三位講員的信息,都值得推薦.研經會由唐佑之牧師主講.唐牧師是美國加州金門神學院教授,對於預言甚有心得,本屆主講撒迦利亞書,使赴會兄姊清楚知道,主再來及其所成的事；因而作聰明的童女,儆醒預備等候主來.講道會由張慕皚教授主講.張教授是加拿大雷城聖經神學院教授,年青而蒙主重用.此次總題為「末世時代成熟合用的追求」,使聽眾獲益良多.奮興會由滕近輝牧師主講.滕牧師是香港北角宣道會主任牧師,兼中國神學研究院院長,多年為培靈會講員.此次總題為「靈程高處的經歷」,發揮淋漓盡致.他列舉聖經中靈程高處的人物,使每一個人聽了,深覺自己的靈程,應如何往高處行.
\newline
\newline
荷蒙九龍城浸信會繼續借用堂址為會場,同工,執事,會友等協助大會籌備,佈置和招待等工作,在人事上,得到他們很大的幫助,謹綴數言,以誌不忘.
\newline
\newline
本講道集蒙鍾少英姊妹,梁壽華弟兄,朱雷麗端姊妹記錄,朱建磯牧師編輯,梁淑貞姊妹校對,林立基弟兄封面設計,乃得如期出版.深信他們的勞苦,在主裡面不是徒然的.歷年以來,本講道集深受讀者歡迎,現各期講道集均已售罄.足見人心渴慕主道,不祇在耳朵方面,也在眼睛方面.年來荷蒙福音傳播中心,代錄製錄音帶發售,使本港信徒,尤其是遠方海外信徒得以分享.願神繼續使用本講道集和錄音帶,使其成為沙漠活泉,曠野道路.
\newline
\newline
去秋大會結束後,余即離港返美,並遠征加拿大.詎料不久,即聞委辦黃仲凱牧師蒙主寵召,息勞歸天的消息.黃牧師一生熱心事主,對於本會的事工,也不遺餘力.如今息了自己的勞苦,作工的果效也隨著他.甚願我們在這短暫的人生中,盡忠於主的託付,至終得以領受主的獎賞,阿們.
\newline
\newline


\newpage

\section{第51屆~港九培靈會~滕近輝~第一講}
\label{sec:693}
\textbf{屬靈的境界}
\newline
\newline
連結: \href{https://www.hkbibleconference.org/session-message/view/693}{\texttt{https://www.hkbibleconference.org/session-message/view/693}}
\newline
\newline
\hyperref[sec:1383]{< < < PREV SERMON < < <}
~
\hyperref[sec:toc]{[返目錄]}
~
\hyperref[sec:696]{> > > NEXT SERMON > > >}
\newline
\newline
出埃及記 24:9-24:10
\newline
\begin{longtable}{cl}
\hline
\hline
章節 & 經文 (和合本修訂版)\\
\hline
24:9 & \begin{tabularx}{0.7\textwidth}{X} 摩西、亞倫、拿答、亞比戶，以及以色列長老中的七十人都上去， \end{tabularx} \\ \\ \relax
24:10 & \begin{tabularx}{0.7\textwidth}{X} 看見了以色列的神。在他的腳下，彷彿有藍寶石鋪道，明淨如天。 \end{tabularx} \\ \\
[1ex]
\hline
\hline
\end{longtable}
各位弟兄姊妹,今天開始了港九培靈會下半世紀的歷史.五十年是個很重要的分界線,照一般情形來講,五十歲以上的人就開始衰退,力量減少,生命萎縮；但對於港九培靈研經會,我們仰望施恩的神,使下半世紀蒙更大的恩典.我要將一節寶貴的聖經送給港九培靈研經會:「他們年老的時候,仍要結果子；要滿了汁漿而常發青.」(詩九十二14)這裡形容一個蒙恩的人,雖年紀老了；但仍然結果子,滿了生命的汁漿常發青；仍然有豐盛的生命.我們也想到詩篇一○三篇五節「如鷹反老還童」的話.在摩西和大衛的時代,他們相信鷹到老時,會重新出一次羽毛；所以年老的鷹,羽毛特別豐滿.求神照樣恩待港九培靈研經大會,求神保守我們屬靈的追求,使我們越久蒙恩越多.
\newline
\newline
當兄弟為了講題而禱告仰望主的時候,我心裡有個感動,我想到基督徒最寶貴的一件事,就是靈性逐漸提高,繼續有追求的心志,向著恩典的高處追求.聖經記載,但以理大概在他將近九十歲時,仍然打開住處的窗子；面向聖城耶路撒冷,一日三次禱告,像素常一樣.他的年紀雖大,但心靈追求恆切負擔的程度仍然不減.求神賜給我們這樣追求的心志!追求之時,由逐漸到靈程的高處.
\newline
\newline
聖經中提到很多的山,共有十四座是特別的；因為在這些山頭上曾經發生過很特別的事,這些事像一幅幅的圖畫,把靈程高處的景況向我們表達.我們將在這十天晚上領受其中靈訓.
\newline
\newline
西乃山,我們不是要看西乃山下百姓的經歷,乃是要看西乃山上摩西的經歷.當摩西上山之時,他進入一種新的且很特別的境界；我們可稱之屬天的境界,它的生活有幾個特點:
\newline
\newline
(一)天品──屬天的品性,其中最要緊的就是聖潔.出四9-10提到七十四位以色列的領袖都上西乃山,他們在山上看見以色列的神腳下彷彿有平鋪的藍寶石,如同天色明淨.神的所在好像藍寶石,就是天的顏色,是聖潔的意義；神在那裡,那裡就有聖潔.如果利未記是首詩歌的話,其中的副歌,必多次重複提到,每提一次神的律例之後,接著就說:「你們要聖潔,因為我是聖潔的.」這話一出現,神是聖潔的神,所以與神有關的一切,也該是聖潔的；神也要祂的百姓,成為聖潔的百姓.在聖經中,藍寶石和壁玉,都是象徵屬天的顏色,也都是象徵聖潔；這是個很重要的信息.在今日黃色時代中,所有屬神的人,必須學這功課,必須追求聖潔；所有屬主的人,要從這淫亂世代中分別出來,讓父神的聖潔充滿在我們心中,讓我們常唱這屬靈的副歌:「神是聖潔的,所以我們也要聖潔.」聖潔的基督徒,聖潔基督徒的家庭,聖潔基督的教會；會對世界產生成聖的力量.
\newline
\newline
(二)天道──屬天的真道,摩西上了西乃山,耶和華向他說話；神的話就是天道,是從上面來的聲音,很清楚,很確定的,帶著權柄,帶著能力.進到摩西的耳中,也進到他的心中,又透過他傳達到所有的以色列人；西乃山上充滿了天道的境界.天道有四個特點:
\newline
\newline
1.知識和行動配合在一起,不是單有知識,也有行動；不是單有行動,也有真理的知識.當神向摩西講話時,聖經用兩種方式來表達,「耶和華曉諭摩西說」,「耶和華吩咐摩西說」,這兩種說法意義不同.「曉諭,」要他知道是神心意的傳達,讓他透過曉諭的話,認識神的心意.「吩咐」是命令,目的在於產生行動；神吩咐摩西是要他順服,順服的行動是必須的.茲舉兩節聖經為例:「耶和華曉諭摩西說.」(出卅一1)「耶和華吩咐摩西說,......」(三十34)這裡我們看見二者的基本意義是一樣,但其間有很重要的分別；神的話語,要兩樣合起來,我們讀聖經就漸漸得到更多屬靈的知識；我們絕對不要忽略,神要我們遵行祂的道.在神學中有個理論是特別關於主的道,這理論頗有意義.有個特別名詞叫「相與論」,意思是說,當神的話臨到一個人的時候,那人就透過神的話與神相面對,立刻在他生命裡產生反應,產生效果,產生改變,產生行動；從此以後,這人就和已往不同了.神的話有特別的作用,賽五十五10-11告訴我們,耶和華的話絕不徒然返回,好像雨和雪降落；使地面有生命生發,種子發出來,植物長出來,生命的活動彰顯.雨的作用較快,雪的作用較慢；但或快或慢,或明顯,或隱藏,都有功效.
\newline
\newline
2.天道和人道互相配合,在神的話裡面,一方面講到天道,同時也講到人道,天道和人道配合起來.例如十條誡命,前四誡是關於神的,後六誡是關於人的.關於神的在先,是基礎；關於人的在後,建立在神道的基礎上.神道,天道是基礎,人道是天道基礎上的建造物；兩者合在一起,就是神的話.感謝主!神的話何等全備!教導我們如何對神,也教導我們如何對人；如何作天父的好兒女,如何做人.真正明白聖經真道的人,一定作個很好的人,待人接物,一定滿有好的表現；他是個義人,也是個虔誠人；虔誠是對神,義人是對人；對神對人都有好表現.
\newline
\newline
3.神的公義和慈愛配合在一起,在神的律法裡面,表現了清楚的公義的律例.神審判的法則確定,從不含糊,耶和華的律例確定,按照祂的話執行祂公義的審判；但是同時,我們在出埃及記所記載的律法裡面,看見了神的慈愛,令人景仰的慈愛!(出二十三5)告訴我們,如果人有仇人,當他看見仇人的驢被重馱壓倒時,不可裝作視而不見遠避,乃要協力相助,要向仇敵顯出愛心.11節又告訴我們,所有以色列人的田地,年年耕種,到了第七年,不可耕種,要把這一年的出產都給窮人.(利十九18):「不可報仇,也不可埋怨你本國的子民,卻要愛人如己.......」這也是在神的律法中.「愛人如己」這四個字,似乎在新約上才有；但事實上在西乃山耶和華給摩西的律法已經有了.34節:「和你們同居的外人......並要愛他如己.......」剛才那一節說的是對本國人的態度,這一節是對外國人的態度,都說「愛人如己.」心胸何等廣闊!超越國際界限.神的話有公義也有慈愛,兩樣都彰顯出來,配合在一起,就是有情有理的律法,既合情又合理,何等完美!
\newline
\newline
4.以感恩的心真誠順從,耶和華的命令,要以色列人遵守,要他們以感恩的心為動機.出埃及記第廿章提到十條誡命.有人說,十條誡命命是從第三節開始；因為第一誡記在第三節「除了我以外,你不可有別的神.」有一次,我發現九龍有間禮拜堂講台後面寫著十條誡命,是由第二節開始,「我是耶和華你的神,曾將你從埃及地為奴之家領出來.」然後才開始第一誡；我忽然領悟到一件以往從未注意到的事,正確地說,十條誡命是由第二節開始的,先講到神的救恩,神的恩典慈愛臨到以色列人,所以神的百姓應該遵守神的話,應以感恩的心為出發點,感恩是動機.各位弟兄姊妹!我們這班蒙恩的人,對慈愛的天父,對恩惠的主,對保惠師聖靈,應由感恩的心發出真正的順從.
\newline
\newline
聖經告訴我們,迦南地是流奶與蜜之地,預表屬天的境界；希伯來書第四章說流奶與蜜之地是預表在基督裡面的生活,在基督裡的境界.在基督裡的境界就是屬天的境界,因為(弗二6)說:「神使我們與基督耶穌一同復活,一同坐在天上.」「天上」也可譯作屬天的境界.我們與基督一同復活,然後與祂一同在屬天的境界生活.迦南地有個特點,就是到處充滿了神的話.當時摩西向以色列人述說又述說神的話,「申命」就是重述命令；好叫以色列人清楚知道,確實記住.摩西說:「以色列人啊!將來你們進入迦南地,千萬要記著申命記第六章的話.」這也是我們每個人所牢記的,因為這是屬天生活的重點.(6-9)節提到神的話要記在心上,戴在額上,繫在手上,寫在房門上和城門上；這些是代表五個生活的中心.我們的心思想,一生的果效都由心中發出；心乃中心之處,我們要讓耶和華的話充滿心中.手上代表行為,就是人生動作的中心,這中心要充滿神的話.額上代表人最明顯的表現,神的話是人生表現的中點.家門上是家庭生活的中心,整個家庭要以神的話為重心,家庭中充滿了神的話.城門上代表社會生活的中心,在這中心也充滿了神的話.這是屬天境界的情況.當以色列人進入迦南之後,約書亞作他們的領袖,在他帶領之下,做了幾件很重要的事；他們過約但河時,將廿四塊大石分為兩處擺上,每處十二塊,上面刻了「神的作為,也就是神的話；在以巴路山上,(參書八30-35)約書亞把神的律法用大字刻在那裡,給以色列人常看見這話；那山是令人注目之地,人經過必看見神的話.當約書亞年老臨終時,他到示劍地方,在一棵大橡樹下立了塊大石頭,把耶和華的話刻在其上,並對以色列人說,這是神和你們之間的見證.(參書二十四26)在整個迦南地充滿了神的話,在屬天的境界裡充滿了神的話.神的話是我們一切的準則.
\newline
\newline
(三)天樂──屬天的快樂(出二十四11)七十四位以色列的領袖上了西乃山,他們到了神的面前,「又喫又喝」,這話使我們頗覺希奇,我讀經至此,真是覺得奇怪,我心裡想,他們到西乃山上,面對耶和華,應當敬拜神才對,應當俯伏在地,戰戰兢兢才對；但是這裡神的話說,這些人在耶和華面前又喫又喝,蒙神悅納.喫喝代表快樂；他們享受神的同在,享受神的面光快樂.真是屬天的喜樂!保羅說:「神的國......只在乎公義,和平,並聖靈中的喜樂.」(羅十四17)公義和平喜樂,這三樣是神國的要素.感謝天父!在神的國度充滿了喜樂,這是屬天的喜樂；如果天國有公義,有聖潔,卻沒有喜樂,恐怕我心不嚮往,雖然想去,卻是馬虎而已；心靈枯乾緊張之地,誰能受得起?終日緊張,嚴肅,沒有歌唱,沒有喜樂,沒有笑容,心靈沒有釋放,相信誰都受不了.感謝天父!祂是慈愛的父,在祂面前有聖潔公義,又有喜樂和平；我們到祂那裡,將享受永遠的福樂.感謝天父!祂為我們預備了永遠的家鄉──甜美的家.今晨聽張慕皚博士講道,提到從前的「甜蜜的家」這首歌,現今在歐美已被破壞,已經不真實了,到處是破碎的家.原是Homesweethome如今變成Brokenhome,這是實在的情形!我們只有在神家中,才得到祂永遠的恩典和福樂.天父是位喜樂的神；詩人說:「我要去,到我最喜樂的神那裡去.」(詩四十三4)有一位基督徒說神是幽默的神,我初聽見時覺得他不大敬虔,後經慢慢體會,頗有同感.神給人幽默感是很大的恩賜.純潔高尚的幽默是喜樂的泉源之一.幽默感的人一來,令人高興,大家露出笑容,大家欣賞,表示這共鳴是創造者給於我們的.有一次,我看見長頸鹿那長長的頸,使我感到創造的主真是有點幽默.我看見企鵝站著多麼神氣,每想起來我就發笑；鸚鵡學人語更加有趣.又有一次,我造訪朋友,一進門就聽見哈哈大笑,我以為朋友那麼高興,因為笑聲是我朋友的聲音；後來才知道是鸚鵡學他笑.這是創造主給牠的本能.我們的神是喜樂的神,我們裡面的喜樂一經主話過濾,經過聖靈消毒；就成為高尚純潔幽默喜樂的歌唱,因耶和華而得的喜樂,是我們的力量.讓我們都作個喜樂的人,這是符合神心意的.
\newline
\newline
(四)天榮──屬天的榮耀(出二十四15-17)耶和華的榮耀停於西乃山,雲彩遮蓋山六天；第七天耶和華呼召摩西到更高的山峰上.在舊約,雲彩代表耶和華的榮耀；而且是耶和華與祂百姓同在的榮耀.在屬天境界中有耶和華的榮耀彰顯,有神的同在.某次我聽一位傳道人講道,他說,全部聖經,他認為最寶貴的就是神的同在.他說,到了新耶路撒冷之時,神的帳幕在人間.看哪!神與人同在,那就是最後的結果,最後的光景.那時候,神的榮耀親自充滿在祂的百姓中間,在屬靈的高處多享受神的同在.
\newline
\newline
(五)天式──屬天的樣式(出二十五40)耶和華對摩西說,要謹慎小心,絕不可忽略,不可隨便；無論作何物件,都要照我在山上指示你的樣式.屬天的樣式,無論作什麼都要按照這法則；神所指定的樣式都有預表的意義,後來應驗在耶穌身上,耶穌基督就是屬天的樣式.在約翰福音裡,主介紹自己是──「在地又在天的人子」在地仍然在天；在祂裡面天地合一.感謝主!天上的樣式在地上彰顯.耶穌基督到世上來,將父神的榮耀給我們看；沒有人見過神,只有在父懷裡的獨生子將祂表明出來.我們在耶穌基督裡面進入屬天的境界.
\newline
\newline
保羅的神學有四個字「在基督裡.」我們與基督一同復活,又在基督裡承受天上一切的福氣,又與基督一同在屬天的境界裡.我們在基督裡別有天地,這世界充滿罪惡痛苦,但我們在基督裡另有一個世界,別有一個天地,這天地是個新的境界；充滿了屬天的和平,屬天的愛,屬天的榮耀,屬天的話,在地猶天.如果我們中間有人到過荷蘭,一定到阿姆斯特丹參觀,此處離市區不遠,有個大花園,其中有全世界最美麗的花；特別是鬱金香花的種類繁多,樣式和顏色千變萬化.二月時節,外面下雪；但花棚中充滿各樣美麗的花朵,外面冰天雪地,裡面卻花天草地,真是兩個世界,兩個境界,冬天和春天的境界.我們在基督裡別有天地,在地如在天,我們在屬天的境界裡生活.
\newline
\newline
(六)天才──屬天的才能或謂屬天的恩賜(出三十一1-6)這段經文中,有幾個難念難記的名字,比撒列是個很重要的名字,是個很有意義的人；亞何利亞伯,這名字更加難讀,也是個很重要的名字,他們二人是同工.耶和華在西乃山上對摩西說,你要作此作彼,幾十百樣,當時摩西聽了,一定心裡想:耶和華啊!你講這麼多有何用,誰能作呢?尺寸材料都講這麼清楚,何處可得此材料呢?我作不了,你別講了.很可能摩西有此感覺；但神對他說,不要擔心,我已預備了兩個領袖互相幫助,也預備一些與他們同心的人,給他們恩賜才能,有工作的心志,他們可以作成我的工作；不要擔心,下山後按照我的話宣告.我們再看(卅五25-26,34),這裡提及有些婦女拿了材料來,她們有智慧,也就是有恩賜,同時她們心受感動要為耶和華作工.還有些有恩賜的人,自己作工,也教導別人.恩賜是可以學的,也就是表示有人的恩賜是隱藏的,可以引發出來；而有的人可以引發別人的恩賜.這些人不但有聰明,有恩賜,並且很細心；不是照自己的意思作工,乃是按照耶和華所吩咐的去作.摩西自己很謹慎,把謹慎也帶給他的同工.耶和華興起這麼多人,有恩賜,有心志,心裡有感動,又能遵照神的吩咐去作；所以會幕浩大的工作易於作成,成了古代奇蹟之一.感謝天父!在屬天的境界裡面,我們更學會了如何有效的運用神給我們恩賜.
\newline
\newline
(七)天契──屬天的團契,摩西到了西乃山頂上,與神同在四十晝夜,領受十誡,兩次共八十晝夜.第一次記於出埃及記廿四章18節.第二次記於卅四章28節.摩西八十晝夜與耶和華同在,耶和華為他作見證.神不是藉異象,不是藉間接的方式；乃是直接明說,神直接向摩西明說,好像人與朋友面對面談話.(出三十三11；民十二8)神對摩西說,你到山上來住在這裡.(出二十四12)「住」意思是長期和神在一起.摩西兩次上西乃山,兩次之間,有很重要的分別；我相信有預表重大的意義.第一次上去,神將十條誡命給他,他帶了法板從西乃山下來,看見以色列人犯罪,心裡憂傷痛苦,發了義怒；他把法板摔碎,法板沒有了!第一次的法板失敗了,神再讓他上山,又給他法板；他再次帶了下山,卅四章說他下山的時候臉面發光!因他與耶和華談話,與神面對面；所以面上發光,以色列眾人都都害怕!摩西要用帕子遮面.第二次把法板帶下來,帶著光輝；雖然當時面上有帕子遮住,但百姓已看見光輝.神的律法第一次給人類,是在西乃山上；律法失敗,因為好的律法對沒有生命的人沒好作用,所以律法失敗了.律法不能救人,律法只能審判,只能定人的罪；但是神第二次賜下祂的真理,在基督耶穌裡賜下來的,在祂裡面充充滿滿的有恩典,有真理.從前以色列人在五旬節時慶祝神賜下律法,後來神就揀選那一天賜下聖靈.當時所有的猶太人都慶祝記念,當日神將律法賜給他們；正在那時聖靈就降臨,這是神揀選的最好時刻,要表示一個重要的真理.聖靈來了,攜帶生命,神的百姓,不再在律法底下；乃在賜生命聖靈的律之下,恩典,生命,都是藉著主耶穌來的.感謝讚美主!摩西兩次得到法板,第一次摔破,失敗了；第二次他的臉發光.雖然當時是隱藏的,但後來耶和華發個預言說:將來要在以色列人中間興起一位先知像摩西,這就是耶穌基督；摩西是耶穌基督的預表.摩西是先知,耶穌基督是更大的先知.那預言說,神的百姓個個都要聽那位先知；所以摩西開始的工作,成全在耶穌身上.摩西把神的話帶給人,這是先知的工作；但是先知第一段的工作失敗,直到後來,那更大的先知來到時；在祂裡面充充滿滿的有恩典,有真理.
\newline
\newline
摩西在西乃山上,與耶和華面對面；當我們到靈程高處之時,我們更會與神面對面.今晚所講的,是屬天境界的七種表現；求主幫助我們領受,使我們的追求朝著標竿直跑.
\newline
\newline


\newpage

\section{第51屆~港九培靈會~滕近輝~第二講}
\label{sec:696}
\textbf{愛慕屬靈的福份 ― 信心的探險}
\newline
\newline
連結: \href{https://www.hkbibleconference.org/session-message/view/696}{\texttt{https://www.hkbibleconference.org/session-message/view/696}}
\newline
\newline
\hyperref[sec:693]{< < < PREV SERMON < < <}
~
\hyperref[sec:toc]{[返目錄]}
~
\hyperref[sec:697]{> > > NEXT SERMON > > >}
\newline
\newline
申命記 34:1-34:3
\newline
\begin{longtable}{cl}
\hline
\hline
章節 & 經文 (和合本修訂版)\\
\hline
34:1 & \begin{tabularx}{0.7\textwidth}{X} 摩西從摩押平原登上尼波山，到了耶利哥對面的毗斯迦山頂。耶和華把全地指給他看：從基列到但， \end{tabularx} \\ \\ \relax
34:2 & \begin{tabularx}{0.7\textwidth}{X} 拿弗他利全地，以法蓮、瑪拿西的地，猶大全地直到西邊的海， \end{tabularx} \\ \\ \relax
34:3 & \begin{tabularx}{0.7\textwidth}{X} 尼革夫，從棕樹城耶利哥的平原到瑣珥。 \end{tabularx} \\ \\
[1ex]
\hline
\hline
\end{longtable}
感謝主!今天晚上雖是天氣惡劣,傾盆大雨；想不到仍然座滿,連詩班都坐滿了.
\newline
\newline
今晚,我們要看第二和第三座山.
\newline
\newline
毘斯迦山(申三十四1-3)摩西在這山頂上的經歷有五個重點:
\newline
\newline
(一)愛慕和追求,摩西上了毘斯迦山,看見耶和華所應許之地何等美好!迦南全地,清楚展現在他眼前,他就生出了愛慕的心.
\newline
\newline
在屬靈靈程的高處,我們看見神屬靈福份的美好,就產生愛慕的心,由愛慕產生追求；從此不再輕視屬靈的福份,相反的,把屬靈福份置於最高地位,全心全意追求.這是何等美好!每個人到了靈程高處時,都有這樣的感受.我們聯想到保羅在(弗一18)說的話；讓我們從中看見一件事,這件事有兩方面,一方面看見神給我們的恩召,有何等指望；另一方面又看見神在聖徒身上所得的基業,有何等豐盛的榮耀.第一方面是我們所得的,神給我們恩召.這恩召的內容,有何等寶貴的指望,包括天上各樣屬靈的福份,這一切都是神賜給我們的.保羅緊緊跟著說第二方面,就是神在我們身上所得到的是何等豐盛的榮耀.我們從神得到,神又從我們得到,合在一起；保羅說,但願聖靈開我們的心眼,讓我們清楚看見這兩方面.在以弗所第一章有十個「得」字,其中六個是講我們從神「得」,四個是講神從我們「得」.兩方面的「得」合在一起,就是神寶貴的永遠計劃.當我們的心眼被光照而看見這榮耀,豐富的時候；立刻在我們心靈的深處,產生愛慕和追求；沒有靈程高處的人,在低處所看見的是這個世界.世界的一切充滿他心,榮華富貴物質的享受充滿他肉體的眼睛；因此,他追求世界.
\newline
\newline
懇求主藉著啟示的靈,照明我們心中的眼睛,把我們的心靈帶到高處；好像摩西到了毘斯迦山頂上.按著神的指示舉目觀看,看見迦南全地,何等美麗!何等豐盛!清清楚楚呈現在他眼前.
\newline
\newline
兩年前兄弟被邀赴中東領會,有一天,數位同工帶我到毘斯迦山上,居高臨下,整個約但河平原一覽無遺,遠眺基列山地,何等壯巍美麗!回溯當日摩西看見神的應許美地,他的心是何等快樂!
\newline
\newline
(二)只憑信心卻有眼見耶和華對摩西說:「這就是我應許之地.」過去摩西是在心中思想,想像；現在耶和華的應許,很具體的,清楚的,實在的,呈現在他眼前.神的應許何等真實!但是,雖然他看見了,而以色列人得到迦南地還是在將來；所以,眼見還要加上信心的展望.感謝天父!祂賜給我們信心,透過信心給我們眼見；藉著眼見我們的信心更加強；當信心加強時,我們就看見更大更奇妙的事.因而信心更加強,然後更強的信心,又有新的,大的看見；這樣,從信心到眼見,從眼見到信心；一直升上去,一直開擴︴開展.保羅說:「只憑信心,不憑眼見.」這話非常正確,但是請注意!只憑信心,雖然不憑眼見,卻有眼見；我再說:雖然不憑眼見,卻有眼見.神把眼見給了摩西,當他看見時,他的信心就得了證實.得到證實之後,他又有信心的展望；將來在神帶領之下,以色列人要進到迦南美地.感謝天父!祂給我們信心,從信心產生眼見,然後從眼見又產生更大的信心,這樣螺旋上去.有人說在將來,信心都要變成眼見；所以將來在天上沒有信心只有眼見.這話也對也不對；因為將來在天上仍然有信心,有聖經的話作根據.保羅說:「如今常存的有信,有望,有愛.」保羅並非說,信,望到了某個階段就不存在,乃是說「常存」,信和望一直存在下去；我們到了天家之時,仍然向前看；因為神是無限的,祂的慈愛無限,祂的偉大無限,祂的奇妙無限,祂的創造無限.我們一直在永世裡面探險,時有新的發現.這是無限的神,與我們有個美好的關係.我們是有限的,如果我們完全知道神的一切,我們就變成無限的了,這是不可能的；所以我們在永世裡的歷程,是個探險的歷程.如今常存的有信,有望,有愛,其中最大的是愛；但是信和望和愛打成一片.神藉著我們看見的,來相信我們未看見的；神給我們很多屬靈的經驗,好像雲彩一樣環繞我們作證據.讓我們根據所看見的和經驗的,來相信那未曾看見的.我們看見了十字架的愛,我們根據這個實際,來相信凡神的作為都是愛；雖然有時我們不明白為何某些事情發生,但是神要求我們,根據十字架的歷史的愛的事實,來接受天父的全愛.神就是愛,已經在歷史中向我們證明了；神對我們的大愛,就是十字架的愛,當我們心裡有任何問題的時候,聖靈的感動,聖靈的引導,就要我們仰望十字架；一望之下,我們心裡一切不明白的都解決了,並不是我們明白了,乃是解決了.根據看見的來相信未曾看見的,是聖經基本的原則,是由信以致有信的意義之一.
\newline
\newline
(三)神的公義向摩西顯明(4末句)「現在我使你眼睛看見了,你卻不得過到那裡去.」這話令摩西何等難過!這句話給我們何等重要的信息,給我們很重要的提醒,也給我們很重要的警告.為何神不讓摩西進入應許美地呢?「耶和華對摩西,亞倫說,因為你們不信我,不在以色列人眼前尊我為聖,所以你們必不得領這會眾進我所賜給他們的地去.」(民二十12)大家都記得,當以色列人在曠野時,耶和華兩次為他們預備水；第一次記載於出埃及記十七章,第二次記於民數記廿章.第一次耶和華吩咐摩西要用杖擊打磐石,就必有水流出來.摩西照樣做,水就流出來了；讓以色列人得到滿足的供應.第二次耶和華對摩西說,吩咐磐石,就有水流出來.但是摩西卻用杖擊打磐石兩次.摩西沒有接受神的話,可能有兩個原因；第一,他根據經驗作事,上一次耶和華叫他擊打磐石,現在他照做,他依靠過去的經驗.過去的經驗是死的,因為每一件事實,每一個經驗,有其個性；我們不能用過去的方法來應付前面活潑的處境,我們必須隨時隨地在每件事上,重新,活潑的來尋求神的引導.這樣,我們就常在聖靈活潑,正確的引導之下.第二,摩西不放心神對他說的話,用自己的方法來做這件事,以為第一次用「打」,水就流出來；這次單「吩咐」恐怕不生效吧!他信不過神.摩西過去曾經受過順服的訓練,在西乃山上,耶和華曾將各樣物件的造法告訴他,他已經學習了完全的,絕對的,百分之百的,按照神的指示來「造」的功課,故此,現在他不應該這樣犯錯!但是他有軟弱,此時他缺少信心,信不過神的方法,要用自己的方法.構成這個錯誤的因素有三:
\newline
\newline
1.奪取了神的榮耀(民二十20下半句)「我為你們使水從這磐石中流出來.」「我」他奪取神的榮耀,沒有尊神為聖.
\newline
\newline
2.破壞了預表的意義,摩西沒有順服神,神之所以兩次的吩咐不同,其中有很重要預表的原因.擊打磐石流出水來,是預表耶穌基督受苦,由祂裡面流出生命的活水；主受苦只有一次；所以摩西擊打磐石,只可有一次.主一次受苦就永遠繼續供應我們一切的需要.摩西自作聰明,不順服神的吩咐,破壞了重要的預表.缺少敬畏神的態度,摩西第二次擊打磐石,一連擊打兩下；第一下是根據過去的經驗,但是第二下顯然是出於極不耐煩的心情之下擊打.他在這嚴肅的事上發怒,沒有尊耶和華為聖.
\newline
\newline
以上三個因素,構成摩西第二次的錯誤.
\newline
\newline
摩西犯了這兩個錯誤,以致不能進入迦南應許美地.由此我們想到保羅所說的話:「可見神的恩慈和嚴厲......,」神有慈愛,極大的慈愛；但是同時,在另方面,神有祂公義的嚴厲,祂公義的嚴厲向祂所愛所重用的摩西顯明出來.神是輕慢不得的!
\newline
\newline
(四)日子如何力量也必如何(申三十四7)摩西死的時候,年一百二十歲,眼目沒有昏花,精神沒有衰敗.摩西至死仍是健康的,他的力量一直持續到最後.證實了神的話.「你的日子如何,你的力量也必如何.」當神把責任給人的時候,神也同時供給完成責任所需的力量；所以我們接受責任,也必接受力量；除非那責任不是從神來的.有人自己去抓責任,因為其中有地位,結果抓到時不能夠作成,因為不是神給他的責任.凡是神給的責任,如果人以謙卑和信心接受,神就同時給於足夠的恩典,有力量來完成責任.當我們求恩典和力量之時,我們要先清楚責任和地位是否神給的；如非神給的,恐神不給你恩賜和力量.如果你的接受是謙卑的,是為了主,是為了基督的國度,為了神的榮耀,為了人的得救；你可謙卑而具信心地接受神的託付,神必然成全祂所給你的託付.感謝讚美主!讓我們想到摩西蒙召時所看見的異象,他在何烈山上,看見荊棘被火燃燒而沒燒毀.火有兩個意思,一是試煉,以色列人接受如火的試煉；但在火煉中得到神的保守.第二個意思,火燒荊棘而不毀,代表神是無窮盡的能源.這個意義在摩西身上顯明出來,他四十年之久作神忠心的管家,在神的全家盡忠,四十年曠野的勞苦,身心靈都在很大重擔之下；但是神的恩典夠他用.當他離世之時,身體仍然保持健康；眼目未曾昏花,精神未曾衰敗.耶和華是他的恩典,是他的力量.讓我們又想到迦勒的話:「從耶和華差派我的日子開始,至今已經四十五年,現在我八十五歲了；那時如何,現在仍然如何.」神的恩典夠他用.
\newline
\newline
摩西在神的家中有六個特點,是聖經記載摩西的事最突出的六個特點,其中有三個特點載於申命記第卅四章.
\newline
\newline
第一個特點(10節)摩西是個偉大的先知,而且預表將來那更偉大的先知──耶穌基督；他偉大到可預表耶穌.
\newline
\newline
第二個特點(10節下句)他是耶和華的朋友,與神面對面,認識耶和華特別深；因此得到神的心,也體會神的心；並有神的負擔.他在西乃山上四十晝夜下來之時,看見以色列人犯罪,他禱告說:「耶和華啊!求你赦免這些犯罪的百姓；倘或你肯赦免他們的罪......不然,」其中有些小點,顯出他的話中斷,表示他在哭；接著說:「求你從你所寫的冊上塗抹我的名.」他願意用自己的生命來代替以色列人；這是神的心,他認識體會神的心,與神有同感.
\newline
\newline
第三個特點(12節)「又在以色列人眼前顯大能的手,行一切大而可畏的事.」神藉摩西所行的神蹟比任何人更多.
\newline
\newline
第四個特點(民十二7)耶和華親自為摩西作見證,說:「摩西在我全家盡忠.」這是很高的稱讚!
\newline
\newline
第五個特點(來十一25-26)這裡提到摩西的犧牲,他甘願放棄埃及皇宮一切的福氣,為要與神的百姓同受苦,為神大有犧牲,他為神放下的實在很多.何等難得!如果我們為神放下摩西所放下的十分之一已經很好了.
\newline
\newline
第六個特點(來十一27)摩西看見那不能看見的主.這是他的信心.
\newline
\newline
以上六個特點構成了摩西這個特殊的人物.
\newline
\newline
(五)對末世教會的旨意(申三十四6)說:神將摩西埋葬,但沒人知道他的墳墓在哪裡.相信其中必有很特別的原因,到底是怎麼我們不清楚；不過這事卻與末世有關係.(啟十一3-6)說,在大災難之時有兩個見證人要到世上來.聖經沒有提他們的名字,但按照他們的行動看來,一個好像是以利亞,一個好像是摩西.以利亞沒見死,摩西被神自己埋葬.所以很多解經家都認為,那兩個見證人,是指以利亞和摩西而言.關於末世之時有兩個見證人要來,這話有三種解釋:
\newline
\newline
(一)是他倆本人要來.
\newline
\newline
(二)是有像他倆的心志之人要來.(如主耶穌說施洗約翰就是以利亞.因為約翰與以利亞有同樣的心志.)
\newline
\newline
(三)是在末世之時,有摩西和以利亞的能力和心志的復興,並非他倆要來,而是他們的能力要復興,他們的心志精神要復興.無論採取哪種解釋,有件事是不可避免的,就是以利亞和摩西心志的復興,他們精神的復興,這是神對末世教會的旨意.在末世之時,我們需要以利亞一樣的人,像摩西一樣的人.但願神在我們的時代,多興起這樣的人.像摩西作先知,傳達神的心意,有神的負擔；像摩西那樣與神面對面,與神有美好的關係；像摩西那樣顯出神的能力；也像摩西那樣盡忠；像摩西那樣有犧牲的心志；像摩西那樣大有信心.求神,懇求神!在我們自己的生命裡面也作這樣的工作.
\newline
\newline
信心的探險
\newline
\newline
第三座山是山地.摩西派十二個探子到迦南地窺探,也可說試探查；他們到了希伯崙山地,在那高高的地勢上,鳥瞰周圍,看見神的應許美地；其中有約書亞和迦勒二人回來報訊,那地真是流奶與蜜之地,他倆抬了一大串葡萄回來,證明迦南地實如神所說那樣的豐富美好.(民十三17-22)
\newline
\newline
希伯崙山地代表靈程高處信心的探險.這山地的意義和毘斯迦山的意義相近.在這個經驗裡面也有幾個要點；讓我們很簡單來領受.
\newline
\newline
(一)要放開膽量(20節)摩西對十二探子說,你們去探查時要放開膽量.這是信心,如果要探險,進入新的境界,有新的發現,有新的得著,就必須有信心,這是絕對少不得的.
\newline
\newline
(二)要帶果子回來,摩西吩咐探子們不要空手回來,要帶果子回來,向以色列人證明迦南地實在美好.他們從以實各谷帶回來大穗葡萄,兩人抬一穗葡萄；有人說太言過其實,太誇張,哪裡有這麼大的葡萄?其實,這裡面有個意思,他們不但要把葡萄帶回來,還要保持葡萄的完整.兩人抬葡萄有兩個意思,一是表示很大,一是不致損壞,保持完整；並不言過其實,也並不誇張.聖經的話都是真實的.他們小心勞苦的帶回來,但是很快樂,他們的探險不落空,有果實,不徒然；他們把證據帶回來,向整個以色列人證明神是可信的.
\newline
\newline
感謝神!在聖經中,在教會歷史中,歷世歷代以來,很多有信心的人,他們作屬靈的探險,進入新的境界,到了新的高原,有新的深度,有重要偉大的發現.他們把果子留下,有的留在聖經中,有的留在教會歷史裡,有的留在書籍裡,有的留在個人見證裡,有的留在講道裡面；多少的果子,好像見證人周圍環繞著我們；這些人都見證了神的真實.
\newline
\newline
(三)要作有真信心的領袖,第二節告訴我們,每個探子都是個領袖,都是族長.為何摩西要派族長出去探查?其中必有很重要的意思,因為族長具有影響力,如果他們經歷了神的信實,他們就會影響很多人,也使全族的人都得到莫大的幫助.一個有信心的人他會作屬靈的領袖.一個作屬靈領袖的人,必須是個有信心的人,互為因果.一個人有信心,很自然就作屬靈的領袖；有了信心,很自然會影響別人.換言之,每個作領袖的人,必須有真的信心.
\newline
\newline
求神在今天的教會中,多多興起真有信心的領袖,神可透過他們,鼓勵很多的人,興起很多的人；使基督的教會,大大的發展.
\newline
\newline
(四)要用信心面對難處,當十二探子到迦南地,據說他們發現強大的敵人；敵人的身體高大,城也很堅固.(28節)這些都是可怕的攔阻,都是難以解決的問題.在信心的探險中,必然有難處,否則就不必探險.面對難處,我們就能得到屬靈的挑戰,這時我們要運用信心,藉著信心,得到神大能的彰顯.神的信實和人的信心合在一起,探險就成功了.感謝天父!祂是信實的,祂的話一點一劃都不能廢去.
\newline
\newline
(五)充分的證明「過了四十天,他們窺探那地才回來.」(25節)在聖經中,四十這個數目有很寶貴的意思,就是充分的證明.這裡四十天表示完全證明神的應許是真的.在聖經歷史中,四十的數字常常出現,都是歷史的事實,後代的人不能假造的.這些史實所表現的意義都相同.例如以色列人在曠野四十年,充分證明人的敗壞和神的保守.申命記中神說:「四十年你們的鞋沒有穿破.」神完全的供應到完全的證明.先知以利亞靠著天使給他的食物,走路四十天；完全證明神的力量夠他用.主耶穌復活之後,四十天在地上,向門徒顯現,完全的證明祂真的復活了.四十,四十,四十,在歷史事實上有同樣的意義,歷史的主有奇妙的安排.
\newline
\newline
感謝主!我們所信的是可靠的.
\newline
\newline
(六)信心的果效,十二個探子之中,迦勒和約書亞二人大有信心,他們報喜信；他們講話對以色列人起了大作用.他們的信心可用四個字來代表:「能」(30節末句)「我們足能得勝.」「美」(十四7)「我們經過之地是極美之地.」「必」(8節)「耶和華若喜悅我們,就必將我們領進那地.」「有」(9節末句)「有耶和華與我們同在.」有了這四個字,最後的結果就是「進」(24節)「惟獨我的僕人,因他另有一個心志,專一跟從我,我就把他領進他所去過的那地；他的後裔也必得那地為業.」迦勒和約書亞有信心察看那地,在信心的探險以後,向以色列人報喜信；因此,耶和華領他們進入美地.
\newline
\newline
各位弟兄姊妹!希伯崙高原是屬靈的高原,在那裡我們可以有信心的探險.求神加強我們的信心,帶領我們進入新的境界.有新的經驗,有新的領受,有新的能力,與神有新的關係；經歷神的信實,看見神奇妙的偉大.沒有信心的人,永遠不能經歷；求神帶領我們到靈程的高處.
\newline
\newline


\newpage

\section{第51屆~港九培靈會~滕近輝~第三講}
\label{sec:697}
\textbf{舉手得勝 ― 何烈山上的復興}
\newline
\newline
連結: \href{https://www.hkbibleconference.org/session-message/view/697}{\texttt{https://www.hkbibleconference.org/session-message/view/697}}
\newline
\newline
\hyperref[sec:696]{< < < PREV SERMON < < <}
~
\hyperref[sec:toc]{[返目錄]}
~
\hyperref[sec:698]{> > > NEXT SERMON > > >}
\newline
\newline
如果今天晚上,在這裡所有的弟兄姊妹都被主得著,那是多麼的寶貴!如果所有的弟兄姊妹們都有愛主的心,都有事奉主的心志,都有跟從主的決心；那真是一支強而有力的屬靈軍隊.我們心裡的禱告,就是求神在所有弟兄姊妹們生命中深深地作工.
\newline
\newline
第四座山──利非訂山(出十七8)告訴我們,以色列人來到利非訂,這地方有座山；沒有提到山名,我們就叫它利非訂山.在這座山上所發生的事,給我們留下很重要屬靈的信息.在這山上的經歷有七個要點,讓我們很簡單來領受.
\newline
\newline
(一)要點第9節告訴我們,在山下有爭戰,山上也有爭戰.這兩種爭戰性質不同,山下的爭戰是用刀槍；山上的爭戰,是耶和華的杖,它是大大的兵器.一般人感到山下的爭戰才是真正的戰爭,常認為山上的爭戰不是爭戰,或不當作爭戰來看.今天晚上的信息,有個最重要之點,讓我們領受神給我們的一個觀念；山上所進行的實在是爭戰,兩樣的戰爭都是重要,合在一起,才能得到真正的勝利.摩西對約書亞說,你我分工合作,你在山下作戰；我們三人(摩西,亞倫,戶珥)在山上作戰.今天我們需要山上作戰的軍隊,我們須認識手中的兵器.耶和華的杖常在摩西手中,有時亞倫用這杖；但這杖是屬摩西的,是耶和華用的杖.摩西在山上,他舉手時就把杖舉起來.我們常說,摩西在山上舉手；更正確地講,不是舉手.摩西的手算不得甚麼,任何人的手都算不得甚麼.摩西真正舉起來的,是耶和華的杖,代表耶和華的權柄,代表耶和華的能力,也代表耶和華的同在.當杖舉起時,撒旦就懼怕；神的軍隊就得勝,這是何等重要的一件事!所有事奉神的人,每個教會的領袖,每個團契,團體,基督教機構的領袖,都應該清楚明白這件事.在神的話裡面,如此清楚記載這件事,目的是叫我們從中學習屬靈重要的功課.在聖經中,我們看見至少七次藉著禱告得勝,當時的戰爭很劇烈,構成很大的危機；但在這七次戰爭的危機中,那些信神的人,都藉著禱告得到勝利.茲略一提:
\newline
\newline
使徒行傳二章記載,當時教會和希律王之間有個爭戰；雅各被殺,彼得被下監獄；那時耶路撒冷全教會天天禱告.結果整個局勢改變,耶路撒冷教會大大興旺.
\newline
\newline
創世記卅二章記載,當時以掃帶了四百精兵要對付雅各；那時雅各真是把它當作一場戰爭,他原自知必敗；但他在耶和華面前禱告,好像摔跤時那樣禱告.他說,如果你不給我祝福我就不容你走；結果,他得到勝利.以掃的態度完全改變,他們和好了；一場戰爭變成快樂的相會,化敵為友.
\newline
\newline
以賽亞書卅七章記載,當亞述大軍來到耶路撒冷,那時以色列的領袖都進聖殿禱告.希西家王俯伏在神面前,他自知毫無辦法；惟懇切禱告依靠耶和華,結果亞述十八萬大軍被打敗.
\newline
\newline
歷代志下廿章記載,當約沙法作王時,三國聯軍進攻,約沙法禁食禱告,全國人也禁食禱告.那時,王作了一件很希奇的事,讓詩班走在軍隊前面,結果打了勝仗.用詩班打仗,用耶和華的讚美歌唱勝敵,使三個聯軍退敗.
\newline
\newline
以斯帖記四章載,當時以色列人面臨整個民族被屠殺的危險,波斯國王可能下令殺死全猶太人.在緊急關頭,以斯帖對所有猶大人說,你們要三晝夜禁食禱告,向耶和華懇求；結果,大的危機變成了他們大的祝福.
\newline
\newline
在客西馬尼園裡,我們看見歷史中一場最大的戰爭,是神和撒旦之間的爭戰,是各各他山頭的爭戰.這場戰爭,是萬世的戰爭；戰爭前夕,我們的主耶穌俯伏在客西馬尼園,祂禱告之中,流汗如血滴在地上；如果我們了解這次爭戰的嚴重性,我們就不會希奇主的汗為何那樣滴下.藉著那個禱告,主得勝了,各各他山上的勝利,奠定於得勝禱告之上.
\newline
\newline
在此我們看見利非訂的山頭上,神的僕人摩西舉起耶和華的杖,以色列人就得勝.這是一個重要的歷史給我們功課.
\newline
\newline
(二)要點　以色列的敵人亞瑪力軍很狡猾.(申二十五18)清楚告訴我們,當時以色列人非常疲乏困倦,亞瑪力人從後面悄悄地來,用偷襲的方式進攻.他們從最後面入手,在最軟弱的人身上下手；極其詭詐!面對這樣詭詐的敵人,以色列人需要耶和華的幫助.
\newline
\newline
今天我們所面對的敵人是古蛇,這名字代表魔鬼的特點,極其狡猾；我們非牠的對手,我們很容易像亞當,夏娃上牠的當.神的百姓易於受撒旦迷惑,所以你我必須主的幫助.當我們疏忽大意之時,當我們依靠自己之時,就是我們最愚蠢,最糊塗之時；也是撒旦最歡喜之時,牠可趁機四面八方,施其詭詐狡猾的計謀.各位弟兄姊妹!小心我們的弱點!撒旦常注意我們的弱點,牠認定我們的弱點,針對我們的弱點；趁著我們不防備,不儆醒,不禱告之時,立刻偷偷狠毒下手.當日亞瑪力人如何的偷襲；今日魔鬼也照樣的偷襲.
\newline
\newline
感謝天父!摩西作了一件最聰明的事,他到了山上,舉起耶和華的杖,他何時舉起手來,山下的以色列人何時得勝.
\newline
\newline
(三)要點(出十七12)　「但摩西的手發沉.」開始時他有力量,他的手伸的很直,他覺得很輕鬆,覺得沒問題；但越久越覺累,後來他的手發沉,舉不起來.各位弟兄姊妹!我相信在座每一位,都有禱告的手發沉的經驗；摩西如此,我們也是如此；所以我們需要摩西的經驗作為提醒.
\newline
\newline
(四)要點(12節)　當摩西手發沉之時,他有兩個朋友,同伴,同工；在他兩邊幫助,扶住他禱告的手.這是神所用的方法,也是我們必須用的方法.主藉著禱告的同志,托住我們禱告的手.每個主的僕人,每個基督徒,應有同心禱告的人,這是不可缺少的.一個人禱告很容易消沉,我敢說,任何人若沒禱告的同志,若無禱告的同工；他在禱告的生活上,在禱告的工作上,必定時常軟弱,時常消沉,禱告的時候,手常會垂放下來.在此我們看見個很實際的方法,我們每人都可運用這個方法；我們找禱告的同伴,找禱告的同工,找禱告的同志,(我覺得這兩字用在這裡很好)同樣心志的人,同樣看見的人,同樣認識自己的軟弱,同樣知道自己需要的人,有約定的時間一齊禱告.如果你沒有別人,可以和你太太；如果沒有別人和你一齊禱告,可以約定你的先生一同禱告.有了同心禱告的人,我們的禱告工作,禱告生活就興起穩定.
\newline
\newline
(五)要點　「摩西的手發沉.」這手是多數的,「摩西舉手」這手是單數的；先講單數,後講雙數.由此我們看出一件事,摩西舉手只舉一隻手,拿著耶和華的杖.但下面說他兩手都發沉,言外之意,就是他一隻手舉上去,累時放下；另舉一隻手,不斷輪流更替；到了後來,兩隻手都累得發沉.從單數至複數,兩手都發沉,沒有辦法了!就在這時,他的倆同工來扶助,搬了塊石頭讓摩西坐上.這事做得很智慧,一方面對摩西有好處；因站著舉手倍加的累,坐下可以保持身體的力量.一方面對兩位同工也有好處,因摩西坐下,降低高度,兩同工扶助他；手就不必舉得太高,也可節省氣力.我們在禱告上需要用智慧,有的人作事很聰明,有各樣好的方法；但在禱告上卻很糊塗,有時消耗所有的力量.我們要用智慧來保持我們的禱告,保持我們禱告的心志,保持我們禱告的熱誠,保持我們禱告的負擔.智慧的方法對每個人不同,但我們都應當運用神給我們的智慧；找最合適的地方一同禱告.
\newline
\newline
(六)要點(12末句)　「直到日落的時候.」他們舉手到底,沒有半途而廢,直到勝利完成.他們堅持的禱告非常寶貴.
\newline
\newline
(七)要點(15節)　「摩西築了一座壇,起名叫耶和華尼西.」摩西到了勝利之時,他築了座壇為紀念,這是耶和華吩咐他的.神要摩西記下來以免遺忘,又要摩西告訴約書亞,讀給他聽；意思要他常常牢記,要把這事代代傳下；直到今日傳到我們的手中,傳到我們的耳中,傳到我們的心中.「耶和華尼西」的意思,即耶和華是「我們的勝利,是我們得勝的旗幟,我們的勝利從耶和華而來.這座壇是永遠的紀念,永遠的提醒.各位弟兄姊妹!今晚,讓我們接受在心中,靠著主,在禱告的工作上下工夫,願所有的弟兄姊妹都成為在山上舉手的人；如果我們這樣做,立刻港九教會加添了一股屬靈的新力量.
\newline
\newline
舉手是禱告一幅美麗的圖畫.在聖經中,有七幅禱告的圖畫,第一幅是舉手.第二幅是摔跤.第三幅是叩門；叩門要給他開門.第四幅,一個人站在缺口上,在那裡堵住破口,第五幅燒香,香氣升到神的面前,得神的悅納.第六幅守望樓,我們禱告的時候就到了守望樓上去.第七幅如小鳥大大張口,主就給他充滿.一隻小鳥看見母鳥來時就張口,母鳥可以填滿他的嘴,口張的越大,得到的越多；小鳥很聰明,主也要我們大大張口,主就給我們充滿.這七幅圖畫,每幅有其重點,請記在心!舉手的重點是耶和華的能力.摔跤是不放棄永恆的禱告.叩門是很確定的禱告；如果我叩門,我知道裡面有我要找的人.主說,祈求就給你,尋找就尋見.這兩樣較為一般性,尋找,祈求,還不十分確定,但是主說叩門,那就很確定；你定不會冒然叩門,你必知道所找的人就在裡面,然後叩門,雖聽不到回聲,但你知道人必在裡面,繼續確實地叩.站在缺口的重點,就是防止撒旦的攻擊.燒香是神寶貴禱告.守望是儆醒.張口是信心.
\newline
\newline
這七幅圖畫給我們七個禱告的重點,我們從這山上的經歷得到這個靈程高處的經歷──禱告,舉手,得勝.
\newline
\newline
何烈山上的復興
\newline
\newline
何烈山上有復興的經歷,有兩個人在何烈山上復興.
\newline
\newline
摩西在何烈山上的復興.(出三1-10)摩西淪落放羊的地步,這是他從未想到的；我們可以想像他悲傷的心情,我們可體會到他的軟弱,他的失望.他原有雄心大志,他曾為耶和華犧牲,他曾經厭棄皇宮,他並非不得已才離開皇宮；希伯來書十一章告訴我們,他離開皇宮是他自動的選擇,他為了和神的百姓同受苦難,甘願放棄埃及皇宮的一切享受,同時也放棄了他一切的前途.他原有繼承法老的希望,他原是個有才智的人,曾接受埃及皇室最高的教育.(徒七22)說,摩西說話行事都有才能.他是個前途遠大的青年,但他為了主的緣故,放棄一切的一切；似乎他一切的犧牲都落空了,最後的結局逃難到曠野,在那裡放羊,日夜和羊群為伍.一個有雄心大志的人潦倒至此地步,他的心情何等痛苦失望；但是,感謝主,在這樣情況之下,神彰顯了祂的恩典.神的恩典是在一切的失敗變成了恩典,這是神奇妙的作為.有兩個基本的恩典是摩西所絕對需要的,第一,他學會了不再依靠自己,不再憑血氣；這功課是他非學會不可的,因為是後來真正成功的條件,他曾經伸出拳頭,因為他是皇宮出來的人,他是太子,居高地位,用拳頭解決問題；但是後來他學會,伸拳頭的結果是曠野,失敗,所以他不再靠自己.當他不敢依靠自己時,他開始依靠神；一個人越有才能,越難不依靠才能,越有本事越難依靠耶和華.一個沒本事,沒恩賜的人,無所依靠,當然依靠神；依靠神是很自然的,比較容易.但如果一個人大有才能學問,有特別的背景,有很高的條件；當神要用這人之時,神不得已,必須讓他不再依靠自己.這談何容易!怎能叫個大有本事的人不依靠自己呢?所以神用的方法必須厲害一點加在摩西身上,四十年曠野的生活,把他自己打碎,把這樣強的人,有這麼高條件的人打倒在地,叫他認識自己,謙卑下來,完全依靠耶和華.這個功課實非易學,或者我們較容易些,但對摩西是最難最難的功課.神讓他失敗,教導他學習這功課；他的失敗,在神手中變成恩典.
\newline
\newline
感謝主!神就是這樣的一位神,我們有一切的軟弱,失敗；但當我們回頭歸向祂之時,學會了功課之時,就都變成了神的恩典.
\newline
\newline
在聖經中充滿了這樣的例子:
\newline
\newline
彼得三次不認主,主對他說:「你回頭以後,要堅固你的弟兄.」彼得三次的失敗,成了以後堅固別人的經驗根據.他何時想到三次不認主,立刻想盡辦法,用盡力量去幫助,堅固別人.失敗成了恩典的出口.
\newline
\newline
約翰原是雷子,大概這代表他的性格,他發怒時大發雷霆,相當厲害!他作了主的門徒,主改變他,後來成為一個最了解「愛」的使徒.約翰福音書中,記載「愛」的信息特別多,是其他福音書所無僅有的.這表示他在愛主,愛人上,受了特別的造就；因為他自知缺點,他原是沒有愛的人；當他知道自己軟弱失敗之時,在主帶領之下,特別在這方面追求；結果他的缺點成了他的優點,成了他的強點.
\newline
\newline
在主手中,人的軟弱失敗成了恩典的出口.
\newline
\newline
加在保羅身上的那根刺,是因他所得的啟示太大,主怕他有驕傲的可能或者有驕傲的傾向,所以讓這根刺加在他的身上.結果,這根刺就成為他蒙恩的緣由.基督藉著這根刺大大恩待他,更看見基督的能力,基督的恩典在他身上彰顯.
\newline
\newline
摩西也是如此,他的失敗成為神的恩典,他失敗後跑到曠野.曠野四十年,成了皇宮和事奉之間必須經過的階段；如果摩西從皇宮一躍而為以色列的領袖,帶領以色列人過紅海,過曠野的生活.四十年的時間,他絕對受不了；習慣皇宮生活的人,怎能適應曠野生活呢?一進曠野,必成為第一個受不了而發怨言的人,他的體格,精神,心志,一切生活習慣都顯露他不能適應曠野的生活.神知道他的需要,藉著他的失敗,給他操練,讓他的心志,身體,都能適合神用.摩西在何烈山上看見異象,復興起來,在異象中,他認識了神.我要特別提到摩西認識神的幾點:
\newline
\newline
神向摩西呼召.(出三4)說:神在火中,神在荊棘裡,神在被火燒的荊棘中.請注意我講兩句話:「神在火上,神也在火中」,神在火上支配這火；當以色列民被火試煉之時,神操縱在火之上,保守在火中受試煉的以色列民.神超越這火,祂是火的主,主宰掌管一切；同時,也告訴我們神在火中.有人說,神在火中發生,不過是說,該事件中有神；不錯,這異象神在其中,是出於神.神要用這異象,但還有個意思我們不可忽略,神在火裡與以色列民同在,和他們一同經過苦難.下文有個很重要的意思,神對摩西說:我看見以色列民的痛苦,他們的痛苦我知道,我體會,我了解,所以我要拯救他們.有一處聖經給我很大的幫助,特別介紹給大家:「他們在一切苦難中,祂也同受苦難......」(賽六十三9)「他們」是以色列民,「祂」是耶和華；以色列民在一切苦難中,耶和華與他們同受苦難.這話讓我們了解神的心意,這話解決了我們很多的問題；苦難是個很大的問題,為何人受苦難?為何義人受苦難?為甚麼神創造人的世界上有各樣苦難?為甚麼?為甚麼?記得有一次我到耶路撒冷城外,參觀六百萬以色列人被殺的地方,在紀念處門口；有個約伯的銅像,抬頭向天發問為甚麼?這是以色列人的問題.六百萬人被殺,以色列人極其痛苦,他們想到約伯是他們的代表人物.約伯在那裡問耶和華:「為甚麼?為甚麼?」照樣,我們很多時候抬頭向著神,我們內心的深處,從痛苦中發出呼聲,「神啊!為甚麼?為甚麼有此事發生,似乎與你的慈愛違背,似乎與你的公義違背；我不明白,神啊!巴不得我會明白.」聖經給我們最好的答案,解決了我們的問題.耶和華說:「我的百姓經過這一切苦難,我也與他們一同經過.」人為的苦難,耶和華與他們一同經過；十字架是個證據,神不是潔身自愛,神非漠不關心.祂自己與人類同受苦難,十字架上的苦難,比我們每個人的苦難都大.十字架的事實,給我們得到答案；神就是愛,祂所創造的人類經歷苦難,是必須的,是不可避免的.為甚麼不可避免?我們不得而知,但是我們知道不可避免；如果可避免,神就不這樣作,神與我們同受患難,祂的愛對我們有了保證,最後,神要解決這一切的問題,祂的成功,真正的成功,創造的成功,救恩的成功,祂為我們付了最高的代價,比我們付了更高的代價；任何人的痛苦怎能比各各他山上的痛苦呢?神在火中,以色列民經過埃及的苦難,神與他們一同經過.現在這位火中的神呼召摩西,叫他奉獻,叫摩西也經過火；所以摩西的奉獻,是與神的百姓同受苦難的奉獻.十字架上的主,所呼召的人是背十字架的人.經歷上十字架苦難的主,祂所呼召的同工,是願意背十字架,願意經歷十字架的同工.想到此,我們所付的代價,算不得什麼,還有很多當付的代價我們不願意付.願意付代價的主,在哪裡呼召,祂今日的呼召,如當日神呼召摩西一樣；在火中的神,呼召摩西經歷火；經歷十字架的主,呼召我們和祂一同經歷十字架.
\newline
\newline
以利亞在何烈山上的復興(王上十九)這章記載以利亞的灰心和復興.以利亞的灰心是最令人驚奇的,他由最高峰跌下深淵,但是,感謝主!主使他復興起來.神怎樣對待他,用甚麼方法?第一個方法,請注意5-6節,這裡使我聯想復活以後的主,一天早上,為門徒預備了早餐,主說:「你們來吃!」我相信希伯來人的習慣,老師不做早餐給門徒吃的；這次主卻為門徒預備早餐.主的預備,主的體恤,主的同情和了解是最完全的.以利亞跑到羅騰樹下求死,說,「罷了」意思是說,算了,我放棄了,一切工作都放下了.這樣的情況,為甚麼主知道?
\newline
\newline
第一個原因,當時他的身,心,靈都極其疲乏；因他長期面對緊張關頭,情勢非常危急；到了迦密山高峰,雖事已過,但他整個人垮下來,身,心,靈到了最疲乏的地步.主知道他需要好好睡覺,主給他預備身體所需要的,讓他好睡.聖經沒提他睡了多久,我相信一定等他睡足了,天使才來拍醒他吃早餐；神為他預備的早餐,使他的心靈,身體,精神得到恢復.許多時候,人在靈性上不健康,極需好好睡覺,這是實在的需要.人的身體欠健康,靈性就受了虧損!一個靈性最好的人,當他疲乏時聽道也聽不進去,雖然很想聽,但邊聽邊打瞌睡.從前,我講道時看見有人瞌睡,我很難過；現在卻不然,我知道有幾個最愛主的人,聽道時睡覺,因他身體太疲乏.當我看見他眼睛瞌上又勉強要睜開,我很受感動；他努力掙扎,因他有個願意的心.但是,有的人聽道時,一坐下就睡覺,這是身體的缺點；健康和靈性有密切的關係,主了解我們的需要.
\newline
\newline
第二件事,神在山洞裡問以利亞說:「你在這裡作甚麼?」神要他有安靜思想的機會.
\newline
\newline
第三件事,神使烈風大作,又有地震又有火.這三樣在舊約時代,代表神在那裡.這三樣正是迦密山上的現象,神曾經在迦密山上降下火,使狂風帶來大雨,滿足以色列全國的需要.當耶和華得勝之時,在迦密山頂上真有震動,全國震動,偶像震動,亞哈王震動,耶洗別震動；敬畏耶和華的人也都快樂震動.這迦密山上的經驗,耶和華讓它重現；提醒以利亞,你的神是火的神,是風的神,是震動的神；過去如此,今天仍然如此.後來有個微小的聲音,(12節)為何神要用微小的聲音?我想其中有個意思,神溫柔,微小的聲音,使以利亞立刻有了特別的反應.他從山洞出來,用外衣蒙上臉,出來站在洞口.神第一次吩咐他出來,他沒有出來；他聽到神微小的聲音即有回音.神微小的聲音,有力量地進到他心靈的深處；神慈愛,溫柔,了解的微小聲音,對他講三件事；
\newline
\newline
第一提醒,神說:「以利亞啊!你在這裡作甚麼?」神又用微小的聲音問以利亞:「你為甚麼跑到這裡?」偉大的以利亞,曾經得勝的以利亞!為何在這裡?為何至此光景?慈愛,同情,體貼的提醒!
\newline
\newline
第二命令,「耶和華對他說,你回去......」(15節)神對以利亞說了兩句話:「出來」「去」不要再動怒,「出來」「回去」回到工作的崗位.這是神微小的聲音講出來的命令.
\newline
\newline
第三鼓勵(18節)神安慰他說:「我在以色列中為自己留下七千人,是未曾向巴力屈膝,未曾與巴力親嘴的.」以利亞!你不至孤單,還有七千的同志和你在一起,為耶和華站得住,為信仰站得住；有耶和華為你作見證,面對敵人,仍然堅強立定.感謝主!祂了解我們一切的軟弱.主還作一件很寶貝的事,為以利亞預備個好同工.神說:你去!膏以利沙和你同工,以後要接續你的工作.從此以後,一個人變兩個人,一個先知等於兩個先知；孤單的力量變成加倍的力量,聖靈充滿以利亞,加倍充滿以利沙.神的恩典何等豐富!以利亞求死的意念被神復興過來.
\newline
\newline
今天,我們接受神對以利亞說的話:「出來!」從灰心的洞出來,然後,「去!」去事奉,去工作.
\newline
\newline
迦密山的神是你的神,昨天的神是你今天的神；神要為你預備同工,大同工七千人,小同工以利沙.你去!你去!
\newline
\newline


\newpage

\section{第51屆~港九培靈會~滕近輝~第四講}
\label{sec:698}
\textbf{完全的奉獻}
\newline
\newline
連結: \href{https://www.hkbibleconference.org/session-message/view/698}{\texttt{https://www.hkbibleconference.org/session-message/view/698}}
\newline
\newline
\hyperref[sec:697]{< < < PREV SERMON < < <}
~
\hyperref[sec:toc]{[返目錄]}
~
\hyperref[sec:700]{> > > NEXT SERMON > > >}
\newline
\newline
今天晚上,我心裡有三種的負擔:第一,信息的負擔,懇求聖靈用主的話,讓每位弟兄姊妹心靈裡面得到造就.第二,剛才看見很多站立的弟兄姊妹,我負擔你們太累.後來看見逐漸坐下,我的負擔也逐漸的減輕了；如果你們站著,我講道的長度就要減短了.第一種負擔是要長,第二種負擔是要短,似乎有點矛盾；我們看主怎樣引導吧.第三種負擔,為了還沒吃晚飯的弟兄姊妹,我實在不忍心!求主加給你們力量,但願主的話成為你們的滿足；好像主當日對撒瑪利亞婦人講道,那婦人悔改,主的心滿了喜樂安慰.門徒買食物回來請主吃,主說:「我有食物,我的食物就是遵行那差我來者的旨意.」今晚你們得飽足,因你們的食物就是主的話.
\newline
\newline
剛才所唱的詩都是關係奉獻的詩,領詩的人不曉得我今晚要講甚麼；我所要講的正是奉獻.很明顯在主帶領之下,我們的唱詩,我們所聽的信息,主帶領我們走上奉獻的路.培靈詩選第五十七首副歌:「將你最好的獻於主,獻你年青的力量.」獻你壯年的力量,獻你老年的力量.老年還有力量嗎?有,有的弟兄姊妹老當益壯,在靈程上,在事奉上,還是有力量,「滿了汁漿而常發青.」
\newline
\newline
在我講道之前,我發出奉獻的呼召；不必你們在外表上有任何表示；在你內心把自己奉獻給主.平時,呼召在講道之後,但是今天我要在講道以前,還沒任何人講道,沒任何人的話語,完全是主自己作工；完全是聖靈的工作,完全是主愛的激勵.如果在座的弟兄姊妹,都把自己奉獻,就會成教會一股大的力量.求神在每一個弟兄姊妹心中工作!如果你還沒有奉獻過,今晚你奉獻的時候；響應主的愛,把你的心,你的生命,你的恩賜,你的生活,你一切的一切都奉獻給愛你的主.
\newline
\newline
第六和第七座山,山上所發生的事情,幫助我們體會靈程高處的經歷.
\newline
\newline
第六座山──摩利亞山,就是亞伯拉罕獻以撒之地,他把最好的獻給耶和華；他完全奉獻的愛,在創世記廿二章中有很寶貴的要點.現在我們按著時間所許可來領受.
\newline
\newline
(一)經得起考驗的愛(1節)
\newline
\newline
亞伯拉罕經過了最厲害的試驗,證明他對神真實的愛.
\newline
\newline
每個人都感覺自己愛神,但有的人的愛不過是感覺而已,一經試驗就垮了,就消失了.我們的愛,要受時間和環境的考驗；有時我們周圍看看,五年前熱心愛主的弟兄姊妹,現在不見了.五年試驗的結果,好像那首詩所說的,他的愛消失了,他的愛變成抱怨了.
\newline
\newline
兩週前,我遇見一位巴黎的弟兄,他不久以前回故鄉溫州；他曾在那裡參加個四百人的聚會,在會中,弟兄姊妹熱切禱告.有人問他:「是否用很小的聲音禱告?」他說:「正相反,他們用很大的聲音禱告,一個緊接著一個爭先禱告,因恐錯過機會.」那些弟兄姊妹,至今仍然站得住,仍然滿有愛主的心,他們的信心站得住,愛主的心仍然堅強.
\newline
\newline
(二)計算代價而又不計算代價的愛(2節)
\newline
\newline
這裡有幾句重複的話,神對亞伯拉罕說:「你帶著你的兒子,」「就是你獨生的兒子」,「你所愛的以撒」,為什麼神重複又重複的說呢?難道神覺得亞伯拉罕的心還不夠痛嗎?每加一句如加一刀!為何神要一句一刀的加上去呢?神要他清楚自己的代價,清楚計算奉獻給耶華要付的代價；然後不惜代價奉獻給神.計算代價而又不計算代價,才是真正的愛,是真正的奉獻.亞伯拉罕的愛就是這樣的愛.
\newline
\newline
近來有四位醫生完全奉獻自己來事奉主,一位是新加坡大學醫學院的講師,到非洲西部當宣教士.第二位是本港的蔡元雲醫生,放下醫務投身在突破運動中傳福音,完全擺上.第三位謝興光醫生,放下醫務也放下印尼一間醫學院講師的工作,現在西加利曼丹傳福音.第四位龍維耐醫生,現在申請到西加利曼丹去.四位醫生,四位弟兄；我相信他們事先有很多夜不能入寐.計算又計算,計算清楚了；知道他們所要付的是甚麼代價,然後被主愛激勵感動,完全獻上.
\newline
\newline
(三)順服的愛(1末句)
\newline
\newline
「他說,我在這裡.」當神呼召亞伯拉罕時,他說:「神啊!我在這裡.」這四個字在聖經中有特別的意義.以賽亞說:「我在這裡,請差遣我.」摩西看見何烈山上的異象時,他說:「我在這裡.」他答應神的呼召.撒母耳半夜聽見耶和華的聲音,他起來,說:「我在這裡.」他要聽神的話,遵行神的話.「我在這裡」表示順服,表示預備好了,表示決心要按照所聽見神的話去做；亞伯拉罕一貫的態度如此.主將祂的心意告訴願意接受的人.亞伯拉有了順服的態度,神就告訴他要作甚麼.
\newline
\newline
亞伯拉的順服有三個很寶貴的表現,代表他的順服是清清楚楚,確確實實的順服.
\newline
\newline
(一)「亞伯拉罕清早起來......」(3節)樣樣預備齊要動身,他不拖延,不耽擱,他依照神的話,他積極行動；清早起來,代表他的決心,他澈底的順服.
\newline
\newline
(二)「劈好了燔祭的柴,就起身往神所指示他的地方去了.」(3節)其實他可以不必這樣做,但他心想,將去之處情形如何未知.可能那地無樹,無柴,祭就獻不成；就不能順服耶和華的話,那是萬萬不可的!因此,他在原地預備了柴,帶著柴走遠路實在麻煩,可能人會笑他熱心得糊塗；但是亞伯拉罕的心意,就是準備好要順服.何等寶貴真實的順服!
\newline
\newline
(三)「亞伯拉罕對他的僕人說,你們和驢在此等候......」(5節)亞伯拉罕把僕人打發開,免得獻祭時受僕人攔阻；他預備的多麼週到,他所想所安排的,都是為了實行順服,何等寶貴!
\newline
\newline
(四)無以復加的愛(7節)
\newline
\newline
「以撒對他父親亞伯拉罕說,父親哪!亞伯拉罕說,我兒,我在這裡.」他第二次說,「我在這裡,」這次不是向神說的,是向以撒說的,以撒天真爛漫向父親講話；他父親說:「我兒」這兩個字裡面充滿了愛,說不出來痛苦的愛.「我兒,我在這裡」意思是說:「孩子啊!我預備好了,一切都是為你.」這是他平常對孩子講的話,他對他的獨生兒子就是這樣一切都預備好,他一向是這樣的態度,他的孩子一說甚麼,他立刻就預備好.他整個的心,整個的愛,整個的時間財富,都給他的孩子.他的孩子是他的生命,甚至比自己的生命更寶貴.兩個「我在這裡」似乎是衝突的；這就是他痛苦之處.他愛神,要順服神；他說:「神啊!我在這裡」.他愛他的孩子,他對孩子永遠是這個態度,他說:「孩子啊!我在這裡,」我的生命是你的,你就是我的生命.兩個「我在這裡」衝突,他要在這衝突中作選擇,作決定；兩個愛在一起之時,就顯出比較.這樣,我們立刻想到主對彼得講的話「約翰的兒子西門,你愛我比這些更深嗎?」真的愛,是經得起比較的；在衝突之時作選擇,實不容易!亞伯拉罕在不容易決定之中,表現了他的愛,他把對孩子的愛奉獻給神,他愛神的心超過一切.我相信我們中間沒有人能了解一個到了百歲才得到兒子的父親,是多麼愛他的孩子；不曉得我們中間有沒有到五十才得唯一的兒子,能否了解亞伯拉罕呢?他承受神的應許,經過廿五年,到了他百歲才得到以撒.等他得到時,不但他兩隻手,他整個心,連他整個的生命,把這個孩子抱住；他對孩子的愛,實無以復加!而今他把這愛奉獻給神.我講這件事容易,你們聽這件事也容易；但亞伯拉罕作這件事不容易.
\newline
\newline
在聖經中,有兩個人的經歷稍為靠近亞伯拉罕.大衛作王之前,他知道耶和華已經應許他作王；但是當時掃羅很強盛,有兩次機會,大衛可以輕易殺死掃羅.為了尊重耶和華,大衛不伸手傷害耶和華的受膏者.當時擺在大衛面前的,是尊重耶和華呢?是自己王位榮耀的寶座呢?機會就在眼前,如何選擇呢?如果是你如何?感謝神!大衛尊重耶和華,愛耶和華；放下自己的寶座.我再說:我說很容易,你聽也容易；但當日大衛做這件事不容易,大不容易!
\newline
\newline
摩西從小在法老皇宮長大,他應有盡有,很可能繼承法老的王位；那時埃及是榮耀的大帝國,是世界最文明之地,擺在摩西面前的,一邊是榮華富貴,作最大帝國最榮耀的王.另一邊是與神的百姓,同受苦難.如果是你如何選擇?摩西有信心,有愛心,有忠心,他信耶和華,愛他的同胞；他作了個非常寶貴選擇的決定；他寧願與神的百姓同受苦難.我再說,我說容易,但當日摩西作此選擇實不容易.
\newline
\newline
當日大衛和摩西兩人,還不過是選擇身外之物；但是,亞伯拉罕所面對的選擇,是本身的生命.如果他的孩子死了,按常情說,他活不成；他的孩子比他自己的生命更寶貴.這是切身的選擇,他把這個愛奉獻給神.何等寶貴!
\newline
\newline
(五)無保留的愛(12節)
\newline
\newline
「天使說,你不可在這童子身上下手,一點不可害他!現在我知道你是敬畏神的了,因為你沒有將你的兒子；就是你獨生的兒子,留下不給我.」
\newline
\newline
(六)完全的愛
\newline
\newline
「耶和華的使者從天上呼叫他說,亞伯拉罕,亞伯拉罕!他說:我在這裡.」(11節)亞伯拉罕第三次說:「我在這裡.」這是在他已經舉起了刀,以撒已經被捆綁了,刀正要砍下時,神的聲音來了；亞伯拉罕再次說:「我在這裡.」何等寶貴!他第一次說:「我在這裡,」還是個未知數,能否順服,尚未可知,那是事先的態度.但是現在他說:「我在這裡.」正是適合的話,他已經獻上以撒；已經把最心愛,最寶貴,比生命更寶貴的以撒獻上.他已經在行動中,已經表現他的決心,已經成為事實；這時他帶著事實說:「我在這裡.」
\newline
\newline
感謝主!亞伯拉罕不但奉獻了完全的奉獻,他也奉獻了完全的愛.這是他在摩利亞山上所作的.
\newline
\newline
亞伯拉罕獻以撒預表的意義,亞伯拉罕蒙了大恩典,因為他被神認為值得作神愛的預表；這是何等大的榮譽,沒有任何其他的榮譽比這榮譽更大,他竟然配得預表神把祂的獨生愛子賜給世人的愛.當我們看見這預表的意義時,我們對亞伯拉罕就有一點了解,也可以幫助我們對神的愛有多一點了解.有一次我問一個人:「誰最了解神的愛?」回答說:「百歲有獨生子而獻上的那個人,是最能了解神將獨生子,賜給我們的愛.」我說:「那麼我沒有希望了,大家都沒有希望.」亞伯拉罕很特別,在他那個時代神特別的作為行神蹟；讓他在百歲之時；如聖經所說,好像一個死了的人,生出子孫來.好像天上的星,地上的沙那樣多.我們實在不容易完全了解神的愛,我們只能懇求多一點的了解.
\newline
\newline
神將祂的獨生子,祂所愛的耶穌賜給我們；我們是誰?我們還作罪人的時候,耶穌就為我們死.......這是我十八歲時得救的經句,那天晚上,我睡不著,別人都睡著了；我在牆角,點著一盞油燈,心裡痛苦,我知道我是個罪人；我打開聖經來讀,一直讀到羅馬書五章八節:「惟有基督在我們還作罪人的時候為我們死,神的愛就在此向我們顯明了.」讀到這裡,我明白了,我接受了,這是為我這個罪人的愛,是為我這個罪人的救恩；天父把祂的兒子給不配接受的人.
\newline
\newline
「亞伯拉罕把獻祭的柴放在他兒子以撒身上......」(6)以撒背那燒他的柴,正如耶穌基督背那釘祂的十字架.多麼清楚的預表!
\newline
\newline
神對亞伯拉罕說,你往摩利亞地去,在我所指示你的山上,把以撒獻為燔祭.為何耶和華要他走三天這麼遠的路?為何要到這麼遠處去?這是神的計劃的預表.摩利亞山就是後來所羅門建聖殿之處,聖殿裡面有個祭壇,那祭壇所在.就是昔日亞伯拉罕獻以撒的所在,那祭壇是預表耶穌基督為我們捨身流血.這是何等辛苦預表的連繫!神對亞伯拉罕說:你獻以撒作燔祭,不是隨便任何地方可以獻,一定要往我所指示的地方去.神這樣作,有祂重要的理由；亞伯拉罕當時不明白,但他順服,不明白的順服,是他的智慧.有的人很驕傲,他不明白的,認為一定不對,不管神的話怎樣說；他所明白的才對,多麼狂妄!多麼驕傲!多麼愚蠢!以為自己比上帝更有智慧.昨天我們提到摩西擊打磐石,可能他想,打石頭流出水來,豈非糊塗之事嗎?如果打了流不出水來,豈被人取笑嗎?摩西就是用杖打石頭流出水來的.摩西,他順服,他雖不明白；但神既這樣說,神有道理,有理由.我的本分,我的智慧乃惟神之命是從.弟兄姊妹!我們最高的智慧,就是順服神的話.撒旦叫我們驕傲,叫我們高舉理性,一定要我們明白,用理性審判神,用理性審判神的話,用理性審判神神蹟的作為,用被造的理性來審判創造理性的神.何等愚昧!這是智慧人的愚昧,智慧人中了自己的詭計.亞伯拉罕順服,神要他往那裡,他就往那裡；因此成了預表的奇妙.今天,我們回顧亞伯拉罕,我們為亞伯拉罕的順服而快樂,我們為亞伯拉罕的順服而感謝神.
\newline
\newline
「他們到了神所指示的地方,亞伯拉罕在那裡築壇,把柴擺好,......」(9)亞伯拉罕築壇用石塊,建造起來,一塊一塊的疊上去.這位老先生,那時大概一百十七歲了,至少也有一百十五歲；那時以撒約十五至十七歲之間.老先生搬石建壇,也許以撒在旁,說:「爸爸,我來幫你好不好!」我相信亞伯拉罕不要以撒幫忙,乃因不忍心看見以撒建築獻他自己的壇.亞伯拉罕用年老的身體來築壇,他築這祭壇；預表神在歷史中預備祭壇,耶穌基督在祭壇上獻上自己.舊約的祭壇是預表,是出於神的安排,一直指向各各他山,指向十字架.舊約的祭壇到處皆是,特別是會幕和聖殿裡面.歷史中所有的祭壇,是出於耶和華,神在歷史中,一步步地預備安排這個預表；最後完成十字架,不是偶然的,是神預備的中心.神動用了整個舊約時代,動用了預表,動用了先知的預言,動用了歷史的事實,動用了以色列人的制度,動用了一切的一切；為要作一個個的箭頭,指向各各他,指向十字架.
\newline
\newline
「亞伯拉罕在那裡築壇,把柴擺好,捆綁他的兒子以撒,放在壇的柴上.」(9下)這時是亞伯拉罕最痛苦的時候；舉刀的時候不是他最痛苦的時候,那時已經過了痛苦的高峰,心已麻木,痛苦的感覺已經麻木了.一個十五至十七歲的猶太孩子,身材都已相當高大強壯；如果以撒不肯被捆綁,百多歲的老亞伯拉罕也真是無可奈何.這裡我們看見一個最寶貴的預表,耶穌基督說,「沒有人奪我的命去,是我自己捨的.」主耶穌自己上了十字架,祂自己的腳走上各各他的路；祂伸出自己的手被人釘上,自動的,甘願的.一個自願跟蹤主的人,他自己的走上十字架的路,自己的手伸出來與主同釘十字架.「寫給無名的傳道人」這首詩是中國基督教的傑作,是聖靈的出口,是對青年人的挑戰；也是對所有年青基督徒的挑戰,對你對我的挑戰.我們自己的腳來跟從主,自手來事奉主,作主要我們做的.主手上有釘痕,我們手上沒有釘痕；但是,一個與主同死的人,他的手上有無形的釘痕；當然釘痕與主不同,但是在那裡有個願意被釘死的心志.
\newline
\newline
在最緊急之時,亞伯拉罕發現一隻羊的兩角扣在樹上.豈不是一隻掛在樹上的羊嗎?豈不是十字架嗎?十字架是甚麼?看哪!神的羔羊除去世人罪孽的,掛在十字架上,流出寶血擔當我們的罪,洗淨我們一切的不義.一隻羊掛在樹上,還有甚麼其他更可能清楚的預表?不可能,沒有.神的預表安排奇妙至極,我們不能再要求更多,我們看的這麼清楚啦.這時候,有聲音發出「耶和華以勒」耶和華必預備,在哪裡預備?在耶和華的山上有預備；甚麼山?摩利亞山；甚麼山?各各他山；甚麼預備?被殺的羔羊,贖罪的救主.
\newline
\newline
「並且地上萬國都必因你的後裔得福,因為你聽從了我的話.」(十八)地上萬國都必因亞伯拉罕的後裔得福,這後裔是誰?從上文連起來讀,只有一個很清楚的結論,就是「耶穌基督.」當然還有廣義,以撒是亞伯拉罕的後裔,從以撒下來的以色列人,猶太人都是亞伯拉罕的後裔.今天我們信主的人,都是亞伯拉罕信心屬靈的子孫；但是最重要的就是主耶穌,地上萬國要因主耶穌得福,因祂得救；但是神的旨意,還不只停於此；「後裔」有廣泛的意義,後裔包括你,我在內.神的旨意要所有亞伯拉罕屬靈的子孫,好像主耶穌一樣,讓別人得福,作賜福的器皿.
\newline
\newline


\newpage

\section{第51屆~港九培靈會~滕近輝~第五講}
\label{sec:700}
\textbf{各各他山上的經歷}
\newline
\newline
連結: \href{https://www.hkbibleconference.org/session-message/view/700}{\texttt{https://www.hkbibleconference.org/session-message/view/700}}
\newline
\newline
\hyperref[sec:698]{< < < PREV SERMON < < <}
~
\hyperref[sec:toc]{[返目錄]}
~
\hyperref[sec:701]{> > > NEXT SERMON > > >}
\newline
\newline
第七座山──各各他山　在這座山上的經歷是最寶貴的,這山上所發生的事,改變了七個人的生命.他們在各各他山上,有很寶貴的經歷；願我們也得著他們的經驗.
\newline
\newline
(一)古利奈人西門(可十五21)
\newline
\newline
古利奈人西門有個很特別的經歷,他替耶穌背十架上各各他山；當時他一定覺得很難過,心裡一定充滿抱怨,很不甘心被羅馬兵丁拉夫作此苦差事.他原不認識主耶穌,剛從鄉下進城,意料不到碰上這件不愉快的事.但是,後來到了各各他山上,他把十字架放下以後,看見耶穌被釘在十字架上；又看見十字架上一幕幕的情景,看見主耶穌的表現,看見山頂上所發生的各樣事情.他的內心受感動,他深深知道,釘十架的耶穌不是個平常人；他從所見所感,後來成為基督徒.在(羅十六13)保羅告訴我們,後來西門和他一家都成了基督徒.十字架上的感動,完全改變他的人生.(徒十三1-2)提到安提阿教會一些領袖的名字,其中有個稱呼尼結的西面；有解經家說,那人很可能就是古利奈人西門.果真如此,古利奈人西門不但信了主,且後來成了教會的領袖事奉主,被主大大的使用.古利奈人西門的經歷很特別；其中有神豐富奇妙的恩典,他成為一個最瞭解背十字架意義的基督徒.他會常常回想,並常常告訴別人,可能他常站起來作見證.說:「當時我不明白,一面背十字架,一面抱怨；一面背十字架,一面咒詛.但是現在我明白了,這是天父給我的大恩典,使我有這樣的權利；有這樣的福氣來為我的主耶穌背十字架.當主倒在地上,不能起身背十字架之時,我能得著權利,為主付上一點代價.古利奈人西門永遠永遠為此經驗感謝神.」
\newline
\newline
各位弟兄姊妹!在各各他山上,有個很高靈程的經歷,就是為主背十字架；存著感恩的心來背,經歷了認識上,感覺上的大改變.從勉強到甘心,從抱怨到讚美；歡歡喜喜為主背十字架.
\newline
\newline
(二)十字架上的一個強盜(路二十三42)
\newline
\newline
有兩個強盜與主同釘十字架,一個在那裡諷刺耶穌,至死不悟,心剛硬到底.另外一個強盜,看見主的表現,看見各各他山上各種特別的情形,他的心大受感動；有剛強的表現,坦白說出心裡的話.他想到過去曾聽主耶穌講天國的道理,現在看見主的十字架上有王的牌子；他把過去所聽的,和如今所見的連繫一起.他知道那牌子是諷刺的意思,但是那牌子使他回想過去聽的,天國和天國之王的道理；他深信耶穌就是天國國度的王.所以這時他就講出了他的信心:「耶穌啊!你得國降臨的時候,求你記念我.」耶穌掛在十架上最不像個王,但是祂屬靈的表現正像個王.從表面看十字架上的耶穌,是完全失敗；但在屬靈的表現上,祂是勝利的王.這個強盜在他心靈的深處,感受主耶穌基督屬靈勝利的表現,他接受耶穌基督作王.這是非常寶貴的經歷,他不但接受耶穌基督作他的救主,也接受了耶穌基督作他的王.他在同一時間得到兩樣.在我們屬靈的經歷中,求主幫助我們,接受耶穌基督作救主；也接受祂作王.讓主坐在我們心中的寶座上,管理我們的一生,讓主指導我們.
\newline
\newline
耶穌的十字架上有個牌子,用羅馬,希臘,希伯來,三種文字寫著,「猶太人之王.」這三種文字代表當時所有的語言,羅馬文,代表當時政治的力量；希臘文,代表當時文化和商業的力量；希伯來文,代表神的選民,代表猶太教的力量.羅馬帝國,希臘文化,以色列的宗教,在這三種文字中表達出來.當他們寫「猶太人之王」這牌子時,他們夢想不到他們寫了預言；這位被釘十字架的耶穌就是王.這三種力量都在祂之下,在這三種力量範圍之內祂都是王.後來基督教成為羅馬帝國的國教,基督教大大的影響希臘文化；後來猶太教直接變成基督教.耶穌基督超越一切,祂是王；但是當時那牌子的意義是諷刺的,我們看見兩種情形,「諷刺的王」和「真正的王.」今天耶穌基督在我們心中是哪一種王?許多基督徒,事實上掛了一個諷刺的王的牌子,在他的人生中,他隨便說,耶穌是我的主我的王.實際上他自己坐在他人生的寶座上,主耶穌卻被置於角落,甚而至主耶穌被重釘十字架.「耶穌啊!你是我的主我的王」這樣的講這樣的唱是諷刺；好像當時羅馬兵,吐唾沫在主臉上,打祂的臉,用荊棘作冠冕戴在祂的頭上.圍著在祂的面前戲弄,說,恭喜猶太人的王.諷刺王的,是不認識主的人所作；但是不少屬主的人,蒙主寶血救贖的人,事實上也在那裡讓耶穌作被諷刺的王.當時的人用荊棘作冠冕戴在主的頭上,是戲弄的意思,他們所為是出於侮辱；但是在神奇妙安排之下,裡面有個恩典的記號.荊棘代表咒詛,耶穌基督替我們擔當了罪惡的咒詛,祂忍受了代罪的刑罰.咒詛變成了恩典,從另一角度來看,荊棘的冠冕是假的,是羞辱的,是侮辱的冠冕.許多基督徒把荊棘的冠冕放在主的頭上,讓主的心充滿痛苦；表面上是讓主戴冠冕,事實上是羞辱主.耶穌基督在你我人生中到底有甚麼地位?好不好讓我們一同到各各他山上,設身處地如兩千年前我們在那裡；用心靈的眼睛觀看,用我們的心靈來感受.讓我們看見十字架上的主之時,心裡有清楚的感動.主啊!我要你作我的救主,也作我的君王,我要將我人生最高的地位給你,願你在我心中被加冕,我的主!我的神!
\newline
\newline
(三)羅馬的百夫長(路二十三47；可十五39)
\newline
\newline
羅馬的百夫長是個有實權的軍官,軍階相當高,權柄相當大；這個百夫長,就是指揮羅馬兵丁釘耶穌在十架的,這人竟然說耶穌真是個義人.按理他是不可能講這話的,他這樣說,就等於定自己的罪,打自己的嘴吧,審判自己.因為當時是他指揮兵丁釘死耶穌的,如果耶穌是義人,他是個害死義人的人；當他說耶穌真是義人時,言外之意,難道他自己不了解嗎?這話的嚴重性,難道他不知道嗎?當然知道,不過這時,他心中的感動太強烈,很自然地湧出情不自禁強烈的話.人深受感動時,常不能控制自己的言語；不當講的也講了,不顧自己如何,只感到事實如此.羅馬軍以殘暴著名,他們殺人如麻,置人於死地,司空見慣,名符其實的鐵石心腸,這樣的人,竟然受感動而呼喊說:「這真是一個義人.」表示他感動之深,同時等於承認自己罪惡之深.他不但說耶穌是義人,也說:「這真是神的兒子,」他看見發生在十架上的情形,看見當時大自然的情形,各種的情況合在一起.他深感有種超自然的因素在一切之中,這超自然的神蹟,屬靈的因素把他征服了!這時他謙卑下來,我相信他本非信神的人,更不是敬畏神的人；但是現在他內心的感受,他坦白地講「這真是神的兒子.」
\newline
\newline
各位弟兄姊妹!在各各他山上充滿了證據,從不同的角度給我們看見,耶穌是神為世人預備的救主,至少在十字架上,或說,在各各他山上,有七個證據.記載於約翰福音第十九章有七個證據,證明耶穌基督真是神為人類所預備的救主.
\newline
\newline
第一個證據(18二句)　有兩個強盜與耶穌同釘十字架,應驗了以賽亞書五十三章12節的話:「......祂也被列在罪犯之中；......」事前當猶太人讀到這話之時,他們以為其中不過是大概的意思,但是現在按字面應驗.聖經的預言非常奇妙,有時是象徵性的預言,有時很奇妙的按著字面應驗.既有很多實例擺在我們面前,所以我們解釋將來的預言時要小心,要常保持字句應驗的可能性；雖然有時看起來不可能按字面應驗,但事實上過去如何這樣應驗,將來還是照樣應驗.神是奇妙的神,我們萬勿限制神的作為,在我們的觀念中,在我們的信心中,要為神行神蹟留餘地.神是行神蹟的神,神使我們驚奇；有人說,神是驚奇的神,這話非常寶貴!我們讀神的話,要保持高度的信心.
\newline
\newline
第二個證據18　「他們就在那裡釘祂在十字架上,......」詩篇廿二篇預言說,他們刺穿祂的手,刺穿祂的腳.這是聖靈藉著大衛講的預言,彌賽亞的手腳被扎.被扎就是被刺穿,被刺穿就是被釘透；連主耶穌死的方式都預先講明.
\newline
\newline
第三個證據23-24　兵丁既然將耶穌釘在十字架上,就拿他的衣服分為四分,每兵一分；又拿他的裡衣；這件裡衣,原來沒有縫兒,是上下一片織成的.他們就彼此說,我們不要撕開,只要拈鬮,看誰得著；這要應驗經上的話說,他們分了我的外衣,為我的裡衣拈鬮.兵丁果然作了這事.多麼奇妙的應驗!大衛說這預言,他自己一定不明白,他受聖靈的感動,說了他當時不明白的話；到了時候,應驗在耶穌身上.
\newline
\newline
第四個證據28-29　「這事以後,耶穌知道各樣的事已經成了,為要使經上的話應驗,就說,我渴了.有一個器皿盛滿了醋,放在那裡；他們就拿海絨蘸滿了醋,綁在牛膝草上,送到祂口.」詩篇六九篇大衛的預言,在此又奇妙的應驗了；大衛說「我渴了,他們拿醋給我喝.」(21)想不到這話按字面應驗了.大衛此語有如詩意的描寫；但在聖靈感動之下,在神的安排之中,竟然按字面應驗.
\newline
\newline
第五個證據30　「猶太人因這日是預備日,又因那安息日是個大日,就求彼拉多叫人打斷他們的腿,把他們拿去,免得屍首當安息日留在十字架上.」主耶穌在逾越節的預備日被釘十字架,真是舊約的預表,耶穌就是逾越節的羔羊,主耶穌被釘死之時,正是逾越節羔羊被殺之時.何等奇妙的應驗!沒有人能決定何時生,也不能決定何時死,除非自殺；生和死的時間,不是自己所能控制的.主耶穌的生和死,都是按照預言,在神清楚確定的安排之中.多麼清楚的證據!
\newline
\newline
第六個證據33　「只是來到耶穌那裡,見他已經死了,就不打斷他的腿.」在逾越節羊羔身上,耶和華吩咐摩西告訴所有以色列人,當他們喫逾越節羔羊時,絕對不可有折斷任何一根骨頭.當時摩西不明白為何耶和華要定下這樣的條例；他不明白,但他照樣吩咐如此做.以色列人也不明白,但他們順服神的話去做.我們事後來看,看見神作每件事有祂的理由；神的命令神的吩咐,有祂的道理.讓我再說,我們最高的智慧,乃是用信心來順服神的話；當我們用自己的經驗,和理智來否定神的話時,是我們最愚蠢之時.
\newline
\newline
第七個證據34　「惟有一個兵拿槍扎他的肋旁,隨即有血和水流出來.」從前我不明白這話的意思,我心思想,當主被刺之時,槍刺入心房,當然有血流出來,這點我明白；但為何有水流出來,我不明白；直到有一天,我讀一位醫生所寫的書,當人最痛苦之時,心臟外面的心包(亦可謂心囊)充滿水分.當時主掛在十字架,就是如此.兵丁用槍刺祂時,有血和水流出來.約翰不是現代的醫生,他不可能曉得醫學上的事；但他把親眼所見的事實寫下來.這叫做見證,把所看見的證明出來.聖經是可信的,是親眼看見的人寫的；這樣的證明是側面的證明,比正面的證明更有力量.為何這事記載下來呢?因為裡面有屬靈象徵的意思,應驗了撒迦利亞書十三章1節的話,在彌賽亞的日子,耶和華要為以色列人開贖罪的泉源.到了那日,耶和華在聖山上,在各各他山上；十字架上所懸掛的主耶穌的肋旁,有活水的泉源流出來,是洗罪的泉源,洗淨我們一切的罪過.主的寶血,主生命的水,使我們罪得赦免,也使我們脫離罪惡的權勢.雙重的救恩,都是從十字架上的主留下來的,完全的救恩!奇妙的救恩!我們的天父為你為我預備好的,人若忽略這麼大的救恩,怎能逃罪呢?證據,證據,證據,各樣的證據,完全的證據,擺在我們面前的,是不容否認的事實,是歷史的見證；我們的信仰,不是建立在理論上,乃建立在歷史的事實上,聖經的信仰,或說基督教的信仰,是建立在歷史事實上的信仰；我們若忽略這麼大的救恩,怎能逃罪呢?今晚,不曉得各位當中有哪位還沒信主的,也許還在慕道期間；看哪!看哪!神的話,神的救恩,神的啟示,神的證據,擺在你的面前.今晚,請你真誠的,確實的,來接受主作你的救主.
\newline
\newline
(四)約瑟
\newline
\newline
(五)尼哥底母(38-39)
\newline
\newline
約瑟和尼哥底母兩人有個共同點,都是因怕猶太人,偷偷作門徒,暗中跟從主.但是,經過了各各他山以後,他們看見了十字架上的主以後,他們改變了,兩個改變:
\newline
\newline
(1)現在公開表白他們的信仰,約瑟公開請求彼拉多將耶穌的身體交給他,帶去葬在在他自己新的墳墓裡.這是最公開的公開,這事一定傳遍眾人中間,成了重要新聞之一,成了人人講論的題材.他是個有地位有錢的人,現在竟然公開表白他的信仰立場.尼哥底母是猶太人的老師,用今天的話來講,是猶太人的教授,猶太人的導師.大家都高看他,他有領導力,有學問,有影響力；是個領袖,他曾偷偷來見耶穌,一直偷偷的作門徒.主所講重生的道理感動了他,他雖很有學問,但從未想到重生的道理；他說,年老了豈能再生出來?他這話有很深的哲學,不像表面上那麼蠢.意思是說,人的本性怎可能改變；但是,他聽了主的話之後,有新的理解.他看見主的作為,看見主的表現,聽主的道,看主的品格,看主的神蹟,他暗中一一觀察.他是個有思想,有頭腦的人；不是一時衝動的人,是個思考週到的人.不然他不會夜裡去見耶穌,他詳細觀察思考,最後,他的觀察和思考,在十字架上達到高峰.十字架給他作了結論,這位耶穌是彌賽亞,是我所等候的；祂是耶和華為我們預備的救主.我要接受,我要相信.從此以後,他表白立場,表白信徒的身份.當主死之時,他帶了一百斤香料,和約瑟一齊在眾人面前作了這件事.再不懼怕,再不猶豫,再不倔強；毅然決定表明立場,承認自己是耶穌的門徒.表現膽量,表現勇敢.
\newline
\newline
(2)這兩個人不但如此,他們還有最好的奉獻,不但有信仰,他們的信仰作了他們最好的奉獻,把自己的墳墓獻出來.是家庭的墳墓,把主搬入他的家庭關係裡面.聖經考據家告訴我們,這是猶太人很難作的一件事.他儘可預備個不相干的墳墓葬耶穌,他可以出高價買塊新墳地；但讓出自己家族的墳地給主用,這是非常難做到的,是難得的奉獻.尼哥底母奉獻百斤香料,非常的貴重!他們兩人都表明信仰的立場,又有難得的物質奉獻.
\newline
\newline
(六)巴拉巴(路二十三18)
\newline
\newline
「眾人卻一齊喊著說,除掉這個人,釋放巴拉巴給我們.」巴拉巴作夢也想不到,他本預備等死,等最後宣告死期來臨；不料忽有「釋放巴拉巴,釘死耶穌」的聲音發出.後來耶穌被釘死,巴拉巴真的被釋放.教會歷史相傳,後來巴拉巴信了耶穌.我相信巴拉巴被釋放後,跟耶穌到各各他山上,他看見耶穌被釘十字架.在人群當中,他還有別人所沒有的感受,別人所無的經歷；他好像古利奈人西門的經歷,他個人與耶穌有了關係.耶穌代他釘十字架,世上最了解這件事實的就是巴拉巴.按字面說,耶穌替他死,代他受刑,讓他得釋放.他作基督徒以後,我相信他的見證是最寶貴的；他最會講主代死的道理,好像古利奈人西門,最會講為主背十字架的道理.我猜想一下,會不會在使徒行傳的教會舉行佈道會,有兩個講員,一個是巴拉巴,一個是古利奈人西門.巴拉巴先講「耶穌替我死」；用他自己的經驗來講；叫人受感動,叫人明白.事實擺在眾人的面前,叫人明白了,相信了.古利奈人西門站起來講,你們信主以後,要為主背十字架,要與主同背十字架.他把自己的經驗講出來,真是現身說法!最真實,最透澈,叫聽見的人明白.這兩人的經歷,讓我們想到,我們必須是個感恩的人.你我面對十字架的救恩,我們了解什麼叫代死,什麼叫贖罪；我們體會到這大恩大愛,我們也應該和他們那樣的作見證.
\newline
\newline
(七)約翰
\newline
\newline
約翰的經歷很寶貴.約翰福音第十九章27節記載,當耶穌釘十字架時,很多女門徒還在十字架旁；但男門徒除了約翰之外都跑掉了.所以人都說姊妹好得多,因為姊妹知道沒有危險,所以放心來看耶穌釘十字架；男門徒知道自己危險,所以都逃跑.請問:難道約翰不知身歷險境嗎?他不顧危險,一直跟從主到十字架旁邊,他是個特別的門徒,他有從主來的特別託負；耶穌說:「約翰!看你的母親.」主將祂最愛的人,祂的母親馬利亞交給約翰照顧.約翰接受這個託負,從此以後,他每逢回想釘十字架的主最後對他講的話,他就有新的決心,要好好完成這個託負.他每逢閉上眼睛,十字架的情況重現,立刻他內心就得到新的完成託負的力量.
\newline
\newline
以上七人在各各他山上,有不同的感受,不同的經歷；配合在一起,幾乎是全部靈程的縮影.讓我作個先後次序的編排:
\newline
\newline
第一部　知罪像百夫長
\newline
\newline
第二部　了解主和經歷主代死的愛像巴拉巴
\newline
\newline
第三部　接受主作王好像十字架上的強盜
\newline
\newline
第四部　奉獻給主好像約瑟和尼哥底母
\newline
\newline
第五部　接受主的託負好像約翰
\newline
\newline
第六部　為主背十字架付代價好像古利奈人西門
\newline
\newline
完全認識主,透澈了解主的愛,接受主作王,奉獻自己,接受託負,為主背十字架.這是靈程高處的經歷.
\newline
\newline


\newpage

\section{第51屆~港九培靈會~滕近輝~第六講}
\label{sec:701}
\textbf{全心全意歸向神}
\newline
\newline
連結: \href{https://www.hkbibleconference.org/session-message/view/701}{\texttt{https://www.hkbibleconference.org/session-message/view/701}}
\newline
\newline
\hyperref[sec:700]{< < < PREV SERMON < < <}
~
\hyperref[sec:toc]{[返目錄]}
~
\hyperref[sec:703]{> > > NEXT SERMON > > >}
\newline
\newline
雖然今天又有風球也有雨,但是各位弟兄姊妹仍然冒雨來赴會.一星期來,這是第三次打風了,好像風很歡喜這個禮拜；或者說,風很注意這個禮拜.在風雨之中,我們親近主,領受主的話,歡喜,順服,恭敬的聚集在主的面前.求主賜給我們真正需要的信息.培靈會沒有用.我再說,培靈會沒有用；除非我們用信心和順服來接受主的話,我們願意完全謙卑在主面前；應該破碎的求主破碎,主破碎以後就能夠建立.
\newline
\newline
今晚我們來到第八座山,就是迦密山上的經歷.全心全意歸向神,是靈程高處的表現.
\newline
\newline
以利亞在迦密山上,為耶和華作見證；有許多以色列人,也有巴力的先知和亞舍拉的先知,和他一同在那裡.我們從列王記上十八章18節開始至末了,這段經文中,看見兩種的情形,成為強烈的對比；一方面我們看見以色列人一般的情形,他們心持二意,亦即三心兩意,他們半信半疑向著耶和華.另一方面,先知以利亞全心全意向著耶和華.
\newline
\newline
首先,我們詳細看以色列人當時的一般情形:以利亞對以色列人說,你們為何心持兩意呢?「心持兩意」這四個字,很正確形容當日以色列人的情況；他們的信仰是一半的,他們對神的態度是一半的,他們在神面前是破碎的器皿.他們三心兩意,其因素何在?
\newline
\newline
(一)隨從潮流
\newline
\newline
以色列民受當地迦南各民族風俗的影響,用現今的話來形容,他們隨從潮流.18節有「隨從」兩字,隨從巴力.迦南人敬拜巴力,以色列民住在他們之中,受了一般風氣的影響,和他們一齊去拜巴力,他們在信仰上沒有站穩,信仰沒有根基,何等可惜!
\newline
\newline
今天屬神的基督徒們,也面對潮流的威脅,面對風俗的挑戰.
\newline
\newline
聖經中記載兩種風,一是風俗,以弗所書二章2節用這兩字.很多人的生活隨從今世的風俗.第二個風,就是聖靈的風.一個真正屬神的人,應該是在聖靈的風帶領的方向.世風日下,如果我們隨世俗的風,我們的道路就日趨下坡；如果我們跟隨聖靈的風,靈風日上.但願我們人生的船帆掛起,讓聖靈的風吹動我們人生之船,讓我們靈性的方向朝上,不是隨流而下.求神給我們這班屬祂的人,有屬靈的福份,有堅定的信仰.對基督的信仰,有根有基；讓我們在今世潮流中,為主站立得住!作中流砥柱,是今日基督教會極大的需要!
\newline
\newline
(二)有名無實
\newline
\newline
當時的以色列民,表面上他們敬拜耶和華,事實上他們敬拜巴力.巴力的名字是「主」的意思,以色列人另有其主,表面上耶和華是他們的主,他們口說「我們的主我們的神」；事實上他們擁抱巴力,讓巴力作他們實際的主人.我們這班屬主的人,要記得主的話；主說,你們為甚麼稱呼我「主啊!主啊」卻不按照我的話行呢?近三年來,美國有部份的弟兄姊妹,青年的美國人,他們見面時有個新的招呼方式,握手時彼此發問；「耶穌基督是你人生之主嗎?」這是很好問安的話,彼此提醒.在使徒時代的教會,信徒見面時也有很特別問候的話,他們彼此說:「基督復活了.」這代表他們的信心,從信心裡流露出來的喜樂.在四十多年前,中國大陸有不少的信徒,見面時彼此問:「你得救了嗎?」這是很好彼此提醒的問題,有很多人得到這問題的幫助.我是否個掛名基督徒?是否真的得救?真的重生了嗎?我與主的關係如何?得到重要的提醒.這三種問安的話合在一起,非常寶貴!主復活了,我真是屬主嗎?我重生得救嗎?然後問:「耶穌基督是我人生之主嗎?」讓我們接受耶穌基督的主權,讓主管理我們,讓主作我們的王.
\newline
\newline
有人問我:「耶穌基督是哪種王?」君主制度有兩種,一是立憲的君主制度,一是絕對的君主制度.這青年問:「耶穌基督是立憲的王,或是絕對的王?」我說:「二者皆是」主是立憲的王,祂的憲法是祂的本性──聖潔,仁愛,公義.祂永遠不背乎自己,祂不會違反本性,祂所作一切都符合慈愛,都合乎公義,都合乎聖潔；祂按照祂的本性治理我們.何等好的一位王!同時,主耶穌也是位絕對的君王,祂要求我們絕對的順服.如果有位這樣的王來治理我們,引導我們；當我們順從祂接受祂的引導,我們有福了.若在暴君之下,很可憐,受壓制,沒自由,作奴隸；但是,人若在滿有公義,慈愛,聖潔的王治理之下就有福了.當日彼拉多問耶穌:「你是王嗎?」基督說:「你說的是,我為此而生,我到世上來為真理作見證；凡屬真理的人就聽我的話.」我們的王是真理的王,祂按照真理帶領我們,管理我們.接受基督王權的人有福了.
\newline
\newline
(三)騎牆情況
\newline
\newline
以利亞問以色列眾人,你們為何心持兩意呢?你們這樣做不對,你們要改變,要悔改.他們聽了這話,有何反應?他們「一言不答,」因為他們覺得以利亞說的有理,但是,他們覺得耶和華好,巴力也好；所以一言不答；兩個都要,不作決定,他們騎在牆上.事實上,他們的不決定,成了他們的決定；他們所以走這條可憐的道路,就是因為他們沒有決定.他們以為可以不作決定,或者延遲決定；他們想避開正面的問題,他們想模稜兩可.他們這樣做,事實上,他們就是已經作了決定,接受巴力,離棄耶和華,是不決定的決定.我們在神和巴力之間,必須作積極的選擇；在主和世界之間,必須作積極決定的選擇.
\newline
\newline
今天晚上,讓我們在這裡好像在迦密山上一樣,面對以利亞的挑戰作個決定.是事奉耶和華或是事奉世界?我們不要推辭我們的選擇,如果我們不決定,事實上就是作了相反的決定.若繼此以往,我們最後的道路,必定是愛世界的道路.當日以色列人所以拜巴力,就是這個因素逐漸形成的；逐漸的過程很厲害,很可怕；許多基督徒乃因此而跌倒.讓我們看清楚這件事!今晚,我們面對基督徒生活的重要決定,讓我們一同站在以利亞這方面,全心全意向著耶和華.感謝天父!有七千人站在耶和華這邊,除了以利亞之外,還有許多同道,很多的同志.
\newline
\newline
巴力的意義是什麼?當時在迦南地,有巴力和他的王后成雙成對；王后亞舍拉(19節所見之名稱)另名亞絲巴力.在迦南地,巴力和亞舍拉的偶像雙雙排列；尤以亞舍拉的偶像更令人矚目.她騎在一隻獅子上,胸部敞露,一手拿百合花,另隻手拿著一條蛇,騎在獅子上,代表她是戰爭的神,勇猛威武!胸部敞開代表農業發達,大地豐收,她是農神；當時農業發達代表興旺,成功,豐富,享受.一手拿著百合花,代表性愛.另一手拿著蛇,代表神權──宗教的權柄.她是農神,性神,戰神,宗教神.這考古學家,聖經的學者告訴我們的.這樣的神在迦南地作王,是迦南人所拜的對象；神的百姓受了迦南人的影響,也拜這個偶像；這偶像對我們是何等真實!這四種偶像構成了今日的潮流.基督徒生存在這潮流中,正面對這偶像的引誘,很多基督徒已被這浪潮衝走了.在這偶像面前下拜,接受了這四種意義,成了失敗之基督徒.非常可惜!多少人在那裡敬拜,成功發達,物質享受.多少人以性為神,揀選性,不要主.為了不正當性的滿足,而出賣主,混入淫亂之中.受一般觀念的影響,接受自由性愛的觀念；落在不聖潔的淫亂生活中.如果在港九的各教會,把眾基督徒的生活調查一下,我相信我們會發現不少的基督徒,在這件事上已經跌倒,已經失去聖潔的立場；離開聖經真理的道路,已經把聖潔的神棄於一旁,失了聖徒的立場.受了潮流的衝擊,站立不住.沒有聖經真理的根基,沒有真正認識聖潔的主；沒有讓基督的寶血,洗淨自己的生命；沒有對聖經真理的認真,沒有決心要跟從基督的腳步.撒旦到處誇勝,甚至如啟示錄二章所說:在基督寶血所買贖的教會,有撒旦的座位.今天有多少人在那裡敬拜暴力,戰爭；世界充滿戰爭和暴力,基督徒不知不覺受了暴力觀念的影響；已經離棄基督溫柔的標準,已經忘記基督愛仇敵的真理.何等可惜!有多少人正受蛇的毒害,這條蛇代表撒旦的權勢,透過假宗教,進入人生命血液中.世上有許多含毒素的理論,害人的學說,叫許多人中毒,離開真道.例如有人提倡「無原則絕對的自由」為所欲為.這種自由隨著哲學思想,泛濫在全世界人的心中,特別在青年人的思想；從撒旦口中滴出的毒,一滴滴實存主義的哲學,滴入人的心靈中.物質的享受,性的放任,假宗教的信仰,有毒的理論,戰爭暴力的觀念；這四種的偶像,正是今日事實情形的寫照.當日神的百姓以色列人,在這些偶像面前跌倒失敗；今日基督的教會如何?我們站得住嗎?能勝過這些偶像嗎?這四重偶像正向我們施展壓力.求主使我們像以利亞一樣,來到迦密山頂上,全心全意歸向神.
\newline
\newline
有節聖經對我自己有很大的提醒:「他又領我到耶和華殿的內院,誰知,在耶和華的殿門口,廊子和祭壇中間,約有二十五人,背向耶和華的殿,面向東方拜日頭.」(結八16)這裡提到有廿五個人在祭壇處拜偶像.這是何等重要的提醒!我們奉獻自己以後,並非一勞永逸；我們走在奉獻的路上要小心.我們中間有許多已經奉獻的弟兄姊妹,你已經決志了,已經向主向人表示,你已經是個奉獻的人.但是,請你記得!撒旦能使一個奉獻的人,在祭壇中間拜偶像.奉獻不是一勞永逸的,不是一次過的；奉獻是一條路,我們當一步一步行走.讓我們像以利亞一樣,奉獻以後,一直儆醒謹慎走在奉獻的路上；不容偶像在不知不覺之中進到我們心裡,以致我們成為名義上是奉獻,而實際上是拜偶像的人.
\newline
\newline
以利亞在迦密山上,表現全心全意向著耶和華.茲將他的表現仔細分析(王上十八22)「以利亞對眾民說,作耶和華先知的,只剩下我一人；......」其他的人是被殺呢?還是有別的原因而不作先知呢?兩個可能都有,他們可能被殺,也可能在壓力之下退卻；改變了作先知事奉耶和華的志向,放棄了先知的責任,受不起壓力和引誘.雖然還有七千人沒向巴力屈膝,但他們不是先知,就先知講只剩以利亞一人,只有他一個人站得住,在潮流中得勝.「只有」兩字實在太寶貴!在聖經中有非常寶貴的見證,聖經中曾用於四個人身上.第一個人挪亞,「惟有挪亞在耶和華眼前蒙恩.」(創六8)「耶和華見人在地上罪惡很大,終日所思想的盡都是惡.」(5)惟有挪亞和別人不同.分別出來,為神站得住,並且有見證.何等寶貴!第二個人和第三個人是約書亞和迦勒(民四十24)耶和華親自為他們作見證說:「這許多以色列人十次試探我,惟獨我的僕人約書亞和迦勒,他們另有一個心志,」一個不同的心志,不同的態度,走另一條路,完全不同的表現.第四個人是以利亞,雖然是他自己作見證；但這是事實,他能夠在那樣不敬畏神的潮流中站立得住.他有膽量,一人面對八百五十個假先知──巴力四百五十個先知,加上亞舍拉四百個先知.(19)他有膽量,他肯站起來,挺身而出為耶和華站起來作見證.何等寶貴!在聖經中,我們看見有幾個人也表現這樣的膽量,我們首先想到約拿單.當非利士大軍進攻以色列國,有一晚,約拿單帶著拿兵器的少年人,夜進敵人陣地,攪亂敵人陣腳,大獲全勝.我們又想到基甸的三百勇士,面對敵人大軍,他們敢吹號角,打破瓶子,點起火把,大聲呼喊「耶和華和基甸的刀,」故意引敵前來.他們有信心,有膽量,有敢死和必勝的決心,結果他們大獲全勝.我們為了這樣的決心感謝主,末世時代,我們的主需要有這樣心志的基督徒.英文把這樣的心志叫背脊骨,基督徒需要有背脊骨.不是被撒旦一口氣可吹倒的人,不是為一點利益就出賣基督的人,不是為一點滿足就犯罪的人,不是為名利就把主丟在一邊的人,不是為看重卅兩銀就讓主受羞辱的人；乃是打破玉瓶,願意把最好的奉獻給主的,又好像亞伯拉罕的奉獻,這樣的人.今天教會的背脊骨是屬靈的力量.求主恩待我們!求主激勵我們!求主堅固我們!也求主建立我們!
\newline
\newline
第三點,請看31-32,30節,以利亞重修,重建耶和華的壇,在修建之前,耶和華的壇毀壞不堪.迦密山頂上,原有耶和華的壇,而且是個很重要的壇；過去以色列人曾在那裡向耶和華獻祭,但現在已經荒廢了,因為沒有人用,沒人向耶和華獻祭.多麼荒涼!多麼可憐!這時以利亞起來重建耶和華的壇；並且呼召以色列民一齊修建耶和華的壇.他們修建時用十二塊石頭,代表以色列十二支派,每塊石頭代表一支派,十二塊石頭代表所有神的百姓；用這十二塊石頭來建立神的祭壇,這不是石頭的壇,是以色列人──神的百姓──上帝的兒女的壇.以利亞領先把自己獻上,他有奉獻的心志,他做這事正是代表他的內心,先把自己奉獻,然後和以色列民重建神的壇；這是不可缺少的次序.以利亞實在是個奉獻的人,他有絕對透徹獻身的精神；在他號召之下,眾以色列人和他一齊起來重建神的壇.講到這裡,我們很自然想到基甸建造耶和華的壇；有三個步驟和以利亞所作的三個步驟,給我們很大的提示.在士師記第六章,有兩件似乎是重複的事,其實不是.基甸為耶和華築壇,名叫耶和華沙龍.(24)在25-26節又提到兩件事:
\newline
\newline
其一,基甸拆毀巴力的壇,
\newline
\newline
其二,基甸整整齊齊再為耶和華築一座壇.
\newline
\newline
現在我們把先後過程再看看,第一步基甸建造耶和華的壇.第二步拆毀巴力的壇.第三步整整齊齊為耶和華築壇.第一件事和第三件事是否重複?不是重複的,是兩件不同的事；第一件事是初步建壇,第三件事是完完整整為耶和華築壇.何時建壇完成呢?就是第二步做到的時候,第二步就是拆毀巴力的壇,把壇上巴力和亞舍拉的偶像砍毀.用偶像的木頭作柴燒向耶和華獻祭；偶像的木頭成了向耶和華獻祭的柴,有了這柴,耶和華的祭壇就完成了.在這祭壇上,有真正的祭,祭被火燒,升騰上去；柴是被燒毀巴力的偶像.何等寶貴的祭!何等完整的奉獻!
\newline
\newline
弟兄姊妹!如果我們要把自己奉獻,要重建神的壇,有了第一步的奉獻還不夠；那不過是開始,不過是態度,是起步.還要拆毀偶像,世界上一切的一切,不再是我們追求的目標；要把這一切奉獻給神,包括我們的偶像以撒.當我們燒毀了偶像之後,就完成了我們向耶和華的奉獻,這時的祭壇是完整的祭壇.那麼,我們要問；建壇在前或拆壇在前?感謝神!兩樣都在前,看你指哪一步而言.合在一起,就是神帶領我們走的路,也是神帶領以利亞走的路,更是神藉以利亞,帶領以色列人走的路.第四點,以利亞築好壇之後,擺上祭物,「於是耶和華降下火來,燒盡燔祭,」(38)焚燒燔祭,表示神的悅納；神不能不悅納這樣的祭物.因為這祭物是真的祭物,何等寶貴!何等真誠!芬芳的香味上騰,滿足主的心.當祭物被燃燒之時,我們心眼中所看見的,正是以利亞獻上自己為祭在壇上.天上的火,焚燒以利亞所奉獻的祭,以利亞的人生,成了實際上的祭.以利亞的人生是火的人生,以利亞的特點就是火,讀列王記上,一想到以利亞時,很容易聯想到火.以利亞和火真有密切之關係,我們想到以利亞信心的火,聯想到以利亞禱告的火,想到以利亞奉獻的火,也想到以利亞身上聖靈的火.四樣的火有寶貴的表現.
\newline
\newline
信心的火　「以利亞對亞哈說,你現在可以上去喫喝,因為有多雨的響聲了.」(41)惟以利亞信心的耳朵聽見多雨的響聲,別人沒有聽見.「對僕人說,你上去,向海觀看.僕人就上去觀看,說,沒有甚麼.」(43)僕人無所聞亦無所見,惟以利亞聽見了神所應許的雨,現在要來了.本章的第一段,神已經告訴以利亞,要將雨賜給以色列人；但在恩雨未降之前,必須先有迦密山上的經歷.迦密山上的勝利,迦密山上的回轉,迦密山上的奉獻,迦密山上的全心全意；一個人領導,許多人跟從,造成一股信仰的新洪流.以利亞的信心,使他有先知先覺之明.
\newline
\newline
禱告的火　「亞哈上去喫喝.以利亞上了迦密山頂,屈身在地,將臉伏在兩膝之中.」(42)兩人的表現完全不同,亞哈上去喫喝.以利亞上迦密山頂禱告.他禱告的姿勢很特別,屈身在地,將臉伏在兩膝之中.這代表他的懇切,特別的姿勢代表他特別的懇切.他禱告七次,每次禱告之後,問僕人有沒有雲彩?有風雨的兆頭嗎?僕人每次回來都一樣回答:「沒有甚麼.」他一次又一次地禱告,直到第七次,有一小片如手掌大的雲彩從海裡上來,他知道神的時候到了,神的作為來臨,神的榮耀彰顯,神的能力榮耀就在眼前.這時,他不等雨下來,就差僕人去告訴亞哈當套車下去,免得受雨阻擋.突然間,天陰雲黑,降下大雨.何等的禱告!這禱告裡面真是有火,和他的信心打成一片.當他聽見雨的響聲,亞哈王上去喫喝,他去禱告.請問,他既然聽見雨的響聲,何必再禱告,可坐等雨來.若是你我,我們的信心聽見雨聲來了,還禱告嗎?以利亞知道神的應許要加上人信心的懇求,配合在一起,這是神的旨意.施恩的神定了原則,要人在恩典上與神合作,這是神的恩典.原來神不需要我們禱告,我們的禱告算不得甚麼?我們不禱告,神也可以工作.但神在祂的慈愛中定了個旨意,要我們在祂恩典中與祂合作,祂要施恩,在我們禱告之時,祂就應允我們的禱告.神讓我們同工,這是全部聖經給我們看見的寶貴道理.全能的神,不需要我們別的；只要我們與祂合作,為要把恩典給我們,使我們成為祂的同工.因工作將來我們得賞賜,神目的乃叫我們得賞賜.祂的慈愛何等豐富而奇妙!我們的本分就是禱告,神的原則,如果沒有與祂同工,祂不工作,神等你我在禱告上與祂同工.以利亞在這方面的表現,成為永遠的紀念.
\newline
\newline
聖靈的火　「耶和華的靈降在以利亞身上,他就束上腰,奔在亞哈前頭,直到耶斯列的城門.」(46)耶和華的靈降在以利亞身上,給他超自然的能力,跑在亞哈王馬車前面.我想亞哈的馬車一定套馬多匹,但是以利亞有超自然的能力.神在前面,神永遠是在前的；祂的靈降臨之時,我們打頭陣,我們能看見神的新作為.以利亞在他本應疲倦之時變為剛強,跑在亞哈前面,圍繞城門,宣告,宣告耶和華的作為；宣告耶和華的恩典,彰顯見證耶和華的能力.
\newline
\newline
只有聖靈能幫助我們,我們何等需要聖靈的能力!
\newline
\newline


\newpage

\section{第51屆~港九培靈會~滕近輝~第七講}
\label{sec:703}
\textbf{看見基督的榮耀}
\newline
\newline
連結: \href{https://www.hkbibleconference.org/session-message/view/703}{\texttt{https://www.hkbibleconference.org/session-message/view/703}}
\newline
\newline
\hyperref[sec:701]{< < < PREV SERMON < < <}
~
\hyperref[sec:toc]{[返目錄]}
~
\hyperref[sec:704]{> > > NEXT SERMON > > >}
\newline
\newline
如果有人要對付我,我一定心裡很不快樂；但是,如果主的話要對付我,我就很快樂.願我們每個人讓主的話來對付,存著順服謙卑的心來領受主的話.有人說,在過去,港九大多數的基督徒,只歡喜聽奮興會的講道；因為有很多故事和不少的譬喻；故此,有人就擔心他們對主的話本身的容量不夠,亦缺少深度；但是近年來我相信這種情形已改變了很多,此次三位講員,都是差不多用查經的方式來講.以我自己說,過去的幾個晚上,我沒有講過一個故事,也幾乎沒用過譬喻；只乾巴巴講聖經,但是各位仍然每晚來聽道；這就證明大家愛慕主的道,也表示各位有愛主的心,有愛主的心來聽主的話,這樣,無論用何方式講都是有味道.
\newline
\newline
今晚我們要講第九和第十兩座山.我大膽宣佈,無論時間夠或不夠,一定要做到,我預備用七十分鐘來講,請大家心理上有所準備.如果住在新界或遠處,必須早退者；我不會怪你,相信大家也都會原諒你；不過,可能的話,仍請你們聽完.
\newline
\newline
路加福音第九章,我們看見主耶穌基督在高山上改變形像.聖經沒有題出山的名字,據教會的歷史傳統來講；這座山被認為是他泊山.但不少的聖經學者卻認為不是他泊山,而是黑門山,因為黑門山比他泊山更高,還有其他的理由使他們這樣相信.對我們來說,他泊山也好,黑門山也好,總之是在一座山上.我們稱此山為變像山,那就永遠不會錯,也永遠不會有爭論.耶穌基督在變像山上,顯出祂神性的榮耀,顯出祂天國之王的榮耀.讓我們用對比方式來領受這段聖經的信息.我要將對比研究聖經的方式介紹給大家,當然研究聖經無論用何方式,旨在得著聖經的真義,在於得著而不在於方法.對比研經的方法,一方面能幫助我們把些隱藏的信息發掘,同時也較易引入注意.今晚,我們就用此方式來查考這段聖經,經節除此段之外,還有馬太福音十七章,現在我們把這兩處經文一起來領受.
\newline
\newline
第一對比──暗暗和顯現(路九28-31)　暗暗和顯現在意義上好像相反,但是合在一起給我們明白一個完全的真理.耶穌基督帶著三個門徒暗暗的上山,顯然主的旨意是不讓別人知道.在四福音中,至少有五次提到,主耶穌行了神蹟以後很嚴厲的吩咐人不要宣揚.為何主行了神蹟不讓人注意,要暗暗去做呢?我想最少有三個原因,這三個原因能叫我們更清楚的認識主.
\newline
\newline
第一,主耶穌不要人把信心建立在神蹟上.神蹟好是不錯的,如果人把信心建立在神蹟上,這個根基就不大穩固；很多基督徒跌倒原因在此.
\newline
\newline
第二,馬可一章45節告訴我們,有一次耶穌醫好一個長大痲瘋的人,他得醫治後,耶穌吩咐切勿傳揚這事；但是他到處講,結果,耶穌就不能明明進城；被人群重重包圍,不能走動,攔阻了主的工作.
\newline
\newline
第三,主耶穌不願意把工作的重點放在不正確的地方.主來是要救人,主來要講天國的道理,主來是要人認識神,真正悔改歸向天父,不是來行奇事叫人驚訝羨慕,主知道這重點非常重要.例如主使五千人得飽,他們看見神蹟,當時好像相信,但後來喊叫釘祂十字架的就是這班人.
\newline
\newline
主在復活之前,常將祂神蹟奇事的作為隱藏；在這相信也是一樣,主吩咐他們不要告訴人,門徒也聽見這話,馬太和馬可都是這樣講.主說:「在我復活以前不要講,等我復活以後才講.」意思很明顯不要祂的神蹟和復活連在一起,要人按照祂復活的事來相信祂.當這兩方面合在一起之時,我們的信心就有堅固的歷史基礎.既是暗暗的,為何31節又說,耶穌,摩西,以利亞在三門徒面前顯現呢?主目的是要他們明白,要他們認識主,這是主對門徒的造就.好像主講天國度理之時都是用譬喻,門徒問主為甚麼用譬喻不明說呢?主說:「天國的奧秘只叫你們知道,不叫他們知道.」主這話是什麼意思呢?全面來看,得到很清楚的結論.主的道理──天國的奧秘,只叫有心追求真理的人知道,如果人無心追求,永遠不可能明白；就好像把珍珠丟在豬的面前,牠們踐踏珍珠還來咬你.沒有心的人,永遠不明白基督的真理.惟有愛慕的人,需要的人,謙卑的人,追求的人；神的道,天國的奧秘,是為這樣的人.只要我們有一顆愛慕和追求的心,立刻我們有屬靈的容量,立即有屬靈的眼光,立刻神的恩典豐富的賜給我們,這個條件非常重要.主把最好的給有愛慕之心的人.暗暗上了高山,還有另一方面的意義；我們想想,如果當時十二門徒都和主在一起,主不可能只帶三個人上山.由此可見,這三人特別常和主在一起；而且主帶他們上山禱告,可見這三個門徒是常和主一齊禱告的人.他們是最內圈的門徒,和主最接近；別個門徒有時離開主,所以主帶這三個門徒上高山去,他們都不知道.這三個門徒緊緊跟從主,因此,他們享受了這個特權.一方面這經驗是恩典,同時,神的恩典有祂的條件；神的恩典特別賜給這三個人,因為他們有特別屬靈的條件.他們與主密切接近,他們特有禱告的恩典.主的恩典有時是白白的,毫無條件的；但很多時候,主的恩典也需要條件.求主給我們蒙恩的條件!追求恩典就等於追求蒙恩的條件；如果追求恩典而不追求蒙恩的條件,那就是緣木求魚.
\newline
\newline
第二對比──「禱告」「打盹」　28節提到上山禱告,29節說正禱告的時候.請問!誰正在禱告?答「耶穌.」那麼還有三個門徒呢?答:「睡覺.」對嗎?我說也對也不對；因為開始時的確是耶穌和三個門徒一起禱告,他們不會一上山就睡覺的.他們上了山,找個合適地方禱告；但是到了後來,只剩下一人禱告.你們說對嗎?好像也不十分對；因為他們也不會三個人同時睡著的.開始時四個人,一個繼一個的睡著了；最後,剩下耶穌禱告.「彼得和他的同伴都打盹,」(32)我們所看見的對比是「主在禱告」「門徒在打盹」這次很特別嗎?不見得是例外吧!大概主已試過多次了.我們再想到更嚴重的一次,我相信我們挑選不出更嚴重的時刻,就是在客西馬尼園中.當主正禱告,汗流如血點滴下的時候,主所揀選的三個門徒在那裡睡覺；他們在變像山上睡覺,在客西馬尼園睡覺.這是門徒實際的情形.多少時候我們也是這樣,主是禱告的主,我們是睡覺的門徒.
\newline
\newline
我要和大家一同記得,有三個時候千萬不可睡覺的,如果睡覺就不得了.第一,是戰爭的時候.我們曾經看見利非訂山上,摩西舉起耶和華的杖禱告,他舉手之時,山下的以色列人就打勝仗；山下戰爭,山上也在戰爭,禱告爭戰.山下用刀,山上用舉手,舉起耶和華的杖；如果下面爭戰,上面睡覺,以色列人必敗北,責任在於山上睡覺的人身上.打仗之時,屬靈的戰爭,千萬不能睡覺!很多失敗,都是因睡覺.第二,收割時不能睡覺.箴言書說,收割時睡覺是貽羞之子.莊稼熟了,大家都在收莊稼,都努力工作,都在勞碌；睡覺者是貽羞的人,要受眾人給他的羞辱.今日是末世的時候,是收莊稼之時；主快來以前,莊稼已經熟了,可以收割了,在我們周圍有很多機會,很大的需要,很大的呼聲；當有人在搶救靈魂之時,有人在流淚工作之時,有人流汗工作之時,有人流血工作之時；卻有人睡覺,他是貽羞之子.時候到了,基督的教會,眾教會,所有屬主的人,所有已經接受福音的人；讓我們傳福音,讓我們收莊稼,因為莊稼已經熟了.第三,結婚的時候不能睡覺,你見過有人結婚的時候睡覺嗎?若有,真是頭條新聞,人人必笑他；我想連他自己永遠不會忘,別人也永不忘記,必成為笑料.當主再來之時,主要迎接祂的新婦,我們是基督屬靈的新婦；這時我們不可睡覺,燈要預備油,要顯亮,儆醒等候主再來.如果有人在這三個時候睡覺,那是難以原諒的過失.
\newline
\newline
「正禱告的時候,祂的面貌就改變了,」(29)這話很寶貴,充滿了重要的信息；禱告之時,面貌改變.主是如此,我們也是如此.屬靈的改變和我們的禱告有密切的關係；一個不禱告的人,很難改變他屬靈的相貌.路加福音還有一節很寶貴的經文:「耶穌正禱告的時候,天就開了.」(三21)有神的聲音下來,我們禱告之時,神的聲音和指示和引導就下來.(徒十二5)記載,當彼得被下在監裡,耶路撒冷教會切切禱告；到了第7節說,彼得手鐐開了.他們正禱告的時候,鎖鍊就開了.正禱告之時,面貌改變了.正禱告之時,天開了.正禱告之時,神的話下來.正禱告之時,捆綁鬆開了.正禱告之時,神蹟奇事發生了.正禱告之時,主的拯救出現.弟兄姊妹們!如果我們反過來看,相信要講的很長；教會不禱告就如何如何?你不禱告就如何如何?我不禱告就如何如何?如果基督的百姓不禱告,就會跌倒,失敗,軟弱,羞辱主名；離開主就不能作甚麼.
\newline
\newline
第三對比──「榮光天國的君王」「受苦的僕人」　(31)「他們在榮光裡顯現......」請注意這句話!下文說「談論耶穌去世的事.」這裡提到「榮光」和「去世」.主顯現,天國的君王顯現；主的榮光是天國君王的榮光.「去世」是受苦的僕人.舊約先知預言當關於這兩方面,證明彌賽亞是君王,是勝利者；無限的榮耀,大能大力.第二方面,受苦的僕人,以賽亞書五十三章,以色列人只接受第一方面,他們不了解第二方面,不接受第二方面.他們對彌賽亞的認識是片面的,因此,當主來時,以色列人不認識祂,他們只想到榮耀的君王彌賽亞.他們不曉得神的計劃,是要他們的彌賽亞經過苦難,進入榮耀；因此,他們拒絕主耶穌.如果一個人不認識神的話,或片面了解神的話,或斷章取義,多麼危險!如果按照自己的意思,單要一部份,不要另部份,如果自己審核聖經此對彼錯,我們也落在同樣錯誤之中,也落在同樣危險之中,也會落在同樣的結果裡面.我們要小心!要接受神全部的話!耶穌基督是君王,祂也是受苦的僕人,祂被殺,然後經過苦難,進入永遠的榮耀,這是神的計劃；我們也要與主同作僕人,然後與主同作王.以色列人單看主的王權,而今天卻有很多人單看見耶穌的卑微,看見祂的人性,看見祂是僕人；忘記了祂是榮耀的得勝者.這都是極端.若單看見主的王權,或單看見主卑微的人性,這都是殘缺不全的.進一步說,我們單想與主同作王,不肯學習和主一同作僕人,那麼,將來這王作不成.有人夢想作王,忘記很大的步驟,忘記很大的冠冕；不肯作僕人,不走主的路；他的王座和冠冕永遠是夢,簡直白日發夢!
\newline
\newline
第四對比──「律法的預表」「先知的預言」　(30)摩西和以利亞顯現,摩西代表律法,以利亞代表先知.律法,先知,都是為耶穌作見證,且特別為耶穌去世的事作見證.耶穌基督的死,律法為主的死作見證.先知也為主的死作見證；但是見證的角度不同；律法的見證是預表的見證,先知的見證是預言的見證.預表乃藉著人或事或制度,來表示耶穌將來要作的；預言是用話語或異象,說明將來耶穌要作的；雖方式不同,但目的一樣.用兩不相同的箭嘴,指向主耶穌.在高山上,主顯出祂的榮耀,顯明祂是從神來的,是神為人類所預備的彌賽亞,祂是先知,祂是祭司,祂是君王,三個職份在這山頂上都顯出來.下文說:「你們要聽祂」這是主先知的職份.主要去世,他們在山上談論這事,此乃祭司的職份.祂顯出祂的榮耀,這是君王的職份.先知,祭司,君王,三重的身份,完完全全,在此顯明出來,有見證,有過去的見證；還有,三個門徒是將來的見證.過去,將來,在當時出現,這異象是超越時間的異象,是超越空間的異象.我不應該稱此為異象,因為這是三個門徒實際的經驗,是事實；但是屬靈來講是異象,有心靈的看見,有心靈的體會,有心靈的接納,產生了效果,這個經驗,超越「時」「空」,把「過去」「將來」「現在」連於一起.耶穌基督有證明,還有各種的證明,有全備的證據,有各種的見證人.整個的歷史,全部的聖經,整個的教會歷史,整個的聖經歷史,甚至聖經以外的歷史,有很多很多都是為基督作見證,基督是人類歷史的中心；因為祂是神救贖世人計劃的中心,耶穌基督應驗了摩西和以利亞的見證.
\newline
\newline
第五對比──「喜樂」「懼怕」(33-34)　彼得說了兩句話,一句對,一句不對,他說,主啊!我們在這裡真好!這話非常正確；把他的感覺形容出來,流露他的心情,這個經驗對彼得太寶貴!感謝主!彼得睡覺沒睡到底,他打盹之後醒了過來；如果主改變形像過了他才醒過來,我相信彼得在他整個生命歷史中將有最大的後悔.他將終身抱怨,一輩子不原諒自己.如果我們是個旁觀者,看見彼得在那重大關頭之下睡覺.巴不得叫醒他.叫他及時醒過來,看見主的榮耀,快點說,主啊!我們在這裡真好!我們立刻生出渴慕與主同在的心.每個看見主榮耀的人,都會生出新的渴慕和追求,說,主啊!榮耀的主啊!與你同在真好!我們渴慕與榮耀的主同在.同時,彼得看見神的威嚴就懼怕,他的「喜樂」「懼怕」是相對的,合在一起,是每個屬神的人在神面前應有的態度.(詩二11):「當存畏懼事奉耶和華,又當存戰兢而快樂.」戰兢是敬畏的心,快樂是從神的愛產生的,為了神的公義而敬畏祂；為了神的愛而喜樂,敬畏就是愛裡的怕,又愛又怕.我們對神的態度應當如此,如果單怕,我們是奴隸；如果單有愛,可能放肆.敬畏和愛要配合在一起.
\newline
\newline
第六對比──「看見」「只見」(32-36)　「......既清醒了,就看見耶穌的榮光,並同祂站著的那兩個人.」這裡說彼得看見三個人,耶穌,摩西和以利亞.「聲音住了,只見耶穌一人在那裡.」只見耶穌,看見他們,這兩樣要相配合.摩西和以利亞是事奉主的,他們在那裡見證主,有他們的崗位,有他們的職份,有他們的工作；他們事奉主之時,別人看見,這是好的；他們和主在一起,別人看見主,也看見他們；但是到了最後只有主；主超過一切,無人可以佔據祂的地位,無人可以奪取主的榮耀；只有耶穌,惟有耶穌,絕無一人一事可奪取主的地位和榮耀.
\newline
\newline
第七對比──(太十七5)　耶穌基督是父神所喜悅的,因為主聽父神的話,蒙父喜悅；因祂順服父神的旨意,蒙父的喜悅.何等寶貴!讓我們想到約翰福音十四章有兩節經文,也是同樣的連在一起；主說:「有了我的命令又遵守的,這人就是愛我的；愛我的,必蒙我父愛他,我也要愛他,並且要向他顯現.」(21)主耶穌接著又說:「但要叫世人知道我愛父,並且父怎樣吩咐我,我就怎樣行,起來,我們走罷!」這是主的榜樣,主的見證,他愛父,作父所吩咐祂作的.主吩咐我們的事祂自己作了,祂要我們遵行祂的命令,就證明我們是愛祂,正如祂遵行父的命令,證明祂愛父；我們聽主的話,因為主聽父的話.這裡,我們看見真正屬靈的權柄是從那裡來的；如果人遵行主的道,他有順服,他所講的道,人就願意聽；人聽他的話,因他聽主的話,他有權柄,屬靈的權柄從主而來.
\newline
\newline
最後,我們又回到路加第九章,28節說:他們上山.到了36節提及他們下山.37節他們下了山.他們上了山,得到這個經驗,蒙了這個福氣,有新的看見,新的感受,新的得著,然後他們下山,把所得的和別人分享.神給我們恩典,為要叫我們和別人分享；我們得著,為的要給；聖靈進來,為要充滿我們,要從我們湧流出來.我們所有的一切,神有計劃,神要用我們；我們上山,還要下山,不能停在山上.彼得在山上,有一半的話說錯了,他說主啊!我們在這裡搭三座棚.意思要永遠住在山上,他真是不知道他說的是甚麼,他說錯了,他怎可永遠停留在山上?當時山上的經歷,是為了山下的需要.彼得年老快將殉道之前,他寫彼得後書,第一章下半,他提到這個經驗,說,弟兄們,我盡心竭力的提醒你們,我要鼓勵你們,要幫助你們；因為我在山上有這個經驗,我知道我所信的是誰,所以我要鼓勵你們,要叫你們在已有的恩典上長進.
\newline
\newline
現在我們來看第十座山.這座山也沒有名字；但是主耶穌上了這山就講登山寶訓,所以我們稱之為寶訓山.這山大概在加利利海邊,到過聖地的人,大都上山參觀建立在這山頭上的小禮拜堂.
\newline
\newline
一言而論,在基督國度真善的境界,是整個登山寶訓的中心信息.馬太福音第五,六,七這三章聖經裡,主告訴我們何謂真善.我們要進入真善的境界.現在我們用很快速的時間來看.
\newline
\newline
主在登山寶訓告訴我們:
\newline
\newline
(一)真的善是內在的善　主講了一個例子,當人看見婦女就動淫念時,他在心裡已經犯了姦淫.(太五27)淫行不但是外面的行為,也是內心的,所以善是內在的善.一個聖潔的人,不但行為聖潔,內心也聖潔.
\newline
\newline
(二)真的善是律法精義的善　超越了律法的字句,進入律法的精義裡面.例如；人認為殺人是犯罪,但是主告訴我們,凡動怒的,意思即謂,恨人的就是殺人.(太五21)殺是從恨產生,殺是外面的行動,恨是裡面的行動.律法告訴我們不可殺人,這字句裡有不可恨人的精義；所以真的善不是字句的善,而是精義的善.再舉一例:按律法說,不可起誓.主告訴我們:「你們的話,是就說是,不是就說不是；若再多說,就是出於那惡者.」(太五37)主意思即說,律法的精義不在於起誓,律法的精義不過是字句,律法裡面的精義是甚麼?一個屬基督的人,應有好品格,話就是他品格的代表,是就說是,他的品格擔保是；如果他說不是,他的品格擔保不是；別人可以接受他品格的擔保,不需要說別的.如果別人不信你,恐怕是你無理,因你不值得信任；你的品格難以令人相信,如果你發誓發咒才叫人相信；那是出於惡者,是撒旦從中破壞基督徒的品格.真正的意義就是說,基督徒的品格應是最好的擔保.
\newline
\newline
(三)真善是積極的善(太七12)　你要人怎樣待你,你就先怎樣待人；不是讓人先待你,而是你先待人.何等積極!和孔子所說「己所不欲勿施於人」的話,有很大的分別,孔子說的「不,不,不」是消極的,主說「要」是積極的.積極的善比消極的善更高.
\newline
\newline
(四)真善是勝過惡的善　以善勝惡是真善.主說:「有人打你的右臉,連左臉也轉過來由他打.」(太五39)有人說這是奴隸道德,其實,這完全是門外話,他們不了解主的心意,沒摸到門路.主說這話是最高的標準,需要最大的力量；以愛勝惡是最剛強的；以善勝惡是最大的勝利,是最難得的勇敢.平常人被打還打是最簡單的,任何人都會；但是要學習以愛勝恨,那是最高級的,最有思想,最有深度的人,最有力量,最有堅強意志的人,最有智慧的人,瞭解了最重要的真理.仇恨不能解決人類重要的問題,仇恨只能增加問題,報仇只能增廣仇恨的圈子,只能使更多無辜者受累,只能增加世界的痛苦和黑暗,只能使人類愈變愈成野獸,愈失去人性.惟有愛──基督的愛是最高的道德,最有理由的道德,最有功效的道理,是整個人類社會所最需要的.
\newline
\newline
(五)真善是主動的善(太五41)　「有人強逼你走一里路,你就同他走二里.」第一里是不得已被強迫的,第二里是主動的.要我走一里,我就走二里,這背景是甚麼?當時羅馬兵丁拉夫背行李,一站又一站,(大概一站即一里路)給強迫服務.如果有人勉強你做件事,你就做兩件事；第一件事是不得已,第二件是主動的善.平常的反應,都是被動的.真正的說,人都是反動人物,對事物的反應都是反動的；我們受了刺激必有反應.主在這裡說,我們要在道德上採取主動；我們要靠著主,靠著主的愛,在這世界上創造真的善,創造真的愛,這是主所作的.
\newline
\newline
(六)真的善是有永恆價值的善(太六20)　主說:「你們要積儹財寶在天上.」有永遠的價值,不是單求暫時的福份；追求真善的人有長遠的眼光,能夠分辨暫時和永恆的價值.
\newline
\newline
(七)真善是生命本質的善(太七21)　主說:凡好樹都結好果子,因為它有好的生命.如果我們有好的生命在裡面,就從這好生命的本質發出真的善果.
\newline
\newline
(八)真善不但是內在的也是行動的(太七21)　「凡稱呼主啊主啊的人,不能都進天國；惟獨遵行我天父旨意的人,才能進去.」
\newline
\newline
(九)真善產生真福(太五1-11)八福是真福.真福從哪裡來?是從真善而來.
\newline
\newline
登山寶訓,主伸出祂的指頭,指點我們走上真善的道路.
\newline
\newline


\newpage

\section{第51屆~港九培靈會~滕近輝~第八講}
\label{sec:704}
\textbf{合神心意的教會接受大使命}
\newline
\newline
連結: \href{https://www.hkbibleconference.org/session-message/view/704}{\texttt{https://www.hkbibleconference.org/session-message/view/704}}
\newline
\newline
\hyperref[sec:703]{< < < PREV SERMON < < <}
~
\hyperref[sec:toc]{[返目錄]}
~
\hyperref[sec:706]{> > > NEXT SERMON > > >}
\newline
\newline
今晚,我們來到第十一和第十二座山,這兩座山在聖經中都沒有名字,我們就按著這兩座山上發生的事給它們起名.
\newline
\newline
第十一座山──聖殿山(結四十1-2)耶和華的靈將先知以西結帶到至高的山上,這山是異象中的一座山.山上有聖殿的異象,在這山上,我們看見了合乎神心意的聖殿；是按照神的啟示而建立的聖殿.
\newline
\newline
舊約的聖殿和新約的教會,有彼此遙應之處.從以西結書四十至四十七章裡面,我們來思想合神心意的教會.
\newline
\newline
當我們靈程到了高處之時,我們要愛基督的教會；同時,我們也看見了神心意中教會應該如何.我們對教會的標準提高了,我們大家一齊在教會中追求的目標也提高了；對教會的期望提高了,因此,教會中所有基督徒的標準也提高了.一個人的靈性越高,他越愛他的教會,他不但愛整個教會──基督的身體,他也愛所屬個別的教會.在教會中就有更好的事奉,更符合神的心意.如果一個人說自己靈性高,但在自己的教會中沒有貢獻,沒盡忠事奉,沒努力幫助弟兄姊妹；我就很難相信他真是有高的靈程.
\newline
\newline
神所啟示的聖殿「殿的法則,乃是如此......」(四十三12)這句話是承接上文的,意思是說,神關於聖殿的啟示就如上面所述.接著又作了下面的結論:「殿在山頂上,四圍的全界,要稱為至聖；這就是殿的法則.」這裡又看見相同的話「這就是殿的法則.」由此可見,當中的兩句話,就是神對聖殿法則的總綱:第一,聖殿是在山頂上.第二,聖殿要稱為至聖.聖殿是高的,聖殿是聖的；兩樣合在一起,就代表神對祂教會心意的總綱.今晚我們要看的,就在這兩要點之下.
\newline
\newline
第一個要點,聖殿在山頂上,包含許多很寶貴的意思:
\newline
\newline
(一)教會是屬天的　聖殿在高山上,表示離天很近.教會中應有屬天的原則,來帶領和管制；不是按照人的意思,不是按你我的意思來作教會的事；乃是按照神的啟示,按照從天上來的信息.略舉幾節經文作代表.第四十章3-4節告訴我們,有神的使者,顏色如銅,手拿麻繩和量度的竿.銅代表審判,麻繩和竿代表量度的標準；按照神的標準來審判的教會；這是何等嚴肅的事件!基督的教會,有一天要按照神的標準受審判；我們在基督的教會事奉祂,有一天我們的事奉,要按神的標準受審判.神的使者說,人子啊!以西結啊!我特來指示你,你要特別的領受；用眼看,用耳聽,用心記,用口傳給以色列人聽.今天我們要從這四方面來接受,用我們的靈耳清楚的聽；用我們的眼來看主的話,清楚的看；用我們的心來接受,清楚記住；然後用我們的口去傳達；不但自己得著,也和別人分享.在眾教會裡造成一個普遍的態度觀念,就是我們對於屬靈的原則,有個嚴肅的態度,和大家一齊追求.造成追求的風氣,造成認真的潮流,接受主的標準.第四十章5節很值得我們注意.神的使者手中有竿,這竿有一定的長度.此處提到兩種標準,第一個標準,是神的使者手中之竿的標準.第二個標準,是當時流行的標準.現在我們將這兩個標準作比較,神使者手中的標準和流行的標準不同；第一個標準是所羅門建造聖殿時的長度.第二個數字是流行的標準.各位弟兄姊妹!我們何等受流行標準的影響!因為一般如此,因為別人都這樣作,所以我們也這樣作.不認識神的標準的人是飄流的人,是被世風吹來吹去的人,是無定向的人.當時神要以西結清楚領受,建造神的殿必須按照神的標準；這道理很簡單,但我們很容易忽略.讓我們有個認真的態度,以主的心意代替一般的作法.第四十一章6節告訴我們一件很值得注意的,殿牆的石頭不可用鑿來鑿.(王上六7)說,建造聖殿所用的石頭,都是預先在山上按著神的尺寸作好的.樣樣都備齊,建造時只把石頭擺上就成,不需臨時由任何人開鑿修正.這是很重要的,一切都按照神的方式預備好.我們作的,就是把神已作好的擺上去,按照神的方式擺上去,一切都按照神的指示；這樣的聖殿,才能合乎神心意.
\newline
\newline
(二)神的能力　聖殿在高山頂上代表力量.詩人說「我要向山舉目.......我的幫助從造天地的耶和華而來.」(詩一二一1-2)山代表創造天地萬物神的能力.聖殿山上,耶和華的能力同在.還有一個意思,(但二34-35)告訴我們,但以理所見的異象中,有塊非人手鑿出來的石頭由天而降,打在大像身上.把它打倒,打碎,然後這塊非人手所鑿的石頭；長大發展成為一座大山,充滿天下.這座山代表神的權能,神的國度,充滿天下,這國度非人所作成,是從上面來的.神的國度乃由神的能力完成的,開始時很小,後來成為充滿天下的大山；聖殿就在這座山的山頭上.基督的教會,在神的國度能力的山頂上.這座山也讓我們想到(太十六18)的話；主耶穌說:「我要把我的教會建造在這磐石上,陰間的權柄不能勝過她.」山頂上的聖殿,也讓我們想到不能搖動的磐石.力量就是得勝,有力量就有得勝；神的能力,是我們得勝的力量.在山頂上的聖殿裡,我們看見有種特別寶貴的雕刻,聖殿的柱子上,牆上,周圍的牆上,都雕刻棕樹.(四十31；四十一18-20)聖殿的門上也有棕樹.(25)走廊上也有棕樹的雕刻.(26)聖殿的柱上,牆上,門上,走廊上,四處都有棕樹的雕刻；都是在最顯眼之處,給人容易看見,容易引人注目.棕樹代表勝利,在聖殿中,充滿了神的能力所造成的勝利.感謝主!我們可以放心,可以用信心順服,支取神高山的能力.沒有學會支取神能力的人,永遠是軟弱的.神已經為我們預備了能力的資源,讓我們支取,領受,運用,成全教會的勝利.在聖地有棕樹,香柏樹,橄欖樹,這三種樹都是常青的,代表勝過環境的力量.好像中國的梅,竹,松,在冬季仍是青翠的,梅花在雪天開放,勝過環境,向嚴冬誇勝.照樣,在聖殿裡面有三種寶貴的樹,香柏樹在高山上,迎著強風直立,勝過風的威力,有誇勝的力量.橄欖樹生長在石頭中間,凡到巴勒斯坦觀光的人,最先注意的是到處皆石；可能石頭之多,冠於全球.該地有笑話流傳,當神創造世界時,吩咐一位天使,背著石頭到全世界去散佈；到了巴勒斯坦上空,袋子破了,所有的石頭都落在那裡.這笑話實在很幽默!最奇妙的是在遍地石頭之處,竟有青翠茂盛而根深蒂固,結果子的橄欖樹.棕樹生長在沙漠地；耶利哥有棕樹林,所以外號棕樹城.棕樹在沙漠乾旱的環境中誇勝.感謝神!神為祂的教會預備力量,在這充滿罪惡的世界,充滿痛苦的世界裡,充滿撒旦作為的世界裡；在這彎曲的世代,能夠常見神的恩典和能力.
\newline
\newline
(三)靈程向上　高山當然是高的,在神的旨意中,教會是屬靈的高原.聖殿的構造有個很顯著的特點,(四十22)告訴我們,聖殿外門口的前面有七層台階；如果要進外院,必須登七層台階,級級登高.31節又告訴我們,聖殿內院的門口,有八層台階,如要進這門,須上八層台階.到了49節,在聖殿本身門口,又有台階,沒提多少層；但在聖經七十士譯本中,就是舊約聖經的希臘文譯本中,提到這裡有十層台階；由七至八至十層台階,越登越高.這是聖殿裡面的情形,登高努力一直向上,越靠近至聖所越高.我們的靈程應該是向上的,教會中人應當有向上的路.你我的靈性正在向上或向下呢?你走屬靈的上坡或下坡呢?求主今晚提醒我們每個人,我們願意靠著主,體會主這樣的心意,有追求的態度,讓我們的靈程一直向上去.求主保守我們追求的心永遠不失落.「這圍殿的旁屋,越高越寬；......」(四十一7)我們的靈程越高似乎路越窄,所謂「曲高和寡」,唱低音時很多人都一齊唱；唱高調時,能跟得上者少.當最好女高音的歌聲聳入雲霄時,大家望了興嘆；心裡很佩服,但不能與她同唱.許多人以為靈程越高路越窄,其實正是相反,那是撒旦的話,是人的幻想；靈程越高越有喜樂,越有主的同在,越有得勝；越經歷神的豐富,愈能超越環境,越能預嚐天味.
\newline
\newline
(四)神的同在　聖殿在高山上,越高越親近神.(四十三5-7)說,聖靈把以西結舉起來,進入聖殿的內院；發現耶和華的榮光充滿在那裡,同時聽見有聲音,宣告三件事:第一,神說:「這是我寶座之地.」第二「是我腳掌所踏之地.」這話代表此地是屬神的.第三,耶和華說:「我要在這裡住.不是作客,不是路經此地；而是神要在這裡常住,這裡有耶和華同在.聖殿高處是耶和華同在的地方.感謝主!聖殿是禱告的殿,今日的教會,應是禱告的教會.在舊約中,禱告的殿是個特別的形容詞,是代表聖殿特點的特點.在新約時代,基督的教會應該是個禱告的教會；在禱告中,我們與神面對面.靈程高處是與神親近的地方.
\newline
\newline
(五)教會是神的見證者　聖殿在山的高處,代表教會有見證.主耶穌說:城造在山上,是不能隱藏的.山城一目了然.香港是山城,一進港口,向岸上一望,從最前面看到最頂上；雖然前面房子很高,但擋不住後面的房子,因為房子建造在山坡上.教會建立在山頭上,教會是神的見證者,這是多麼重要!(四十一18-19)說:聖殿牆上雕刻基路伯和棕樹,棕樹的兩邊有兩個臉,一邊是獅子的臉,另一邊是人的臉,是四活物的兩個臉.四活物原有四個臉,現在只見兩臉,另兩臉看不見；那就是說,兩臉是顯露的,兩臉是隱藏的.為何這樣分開呢?有人說,牆上的雕刻是平面的,當然只可看見兩面；如果是立體的就可看見四面,這是勢所必然,沒什麼特別意思.但問題就是為甚麼要揀選人的臉和獅子的臉呢?為什麼不是牛的臉或鷹的臉?這就是表示裡面有選擇,既有選擇,就有意義作為選擇的標準.兩個顯露的臉,獅子的臉代表得勝；人的臉,代表神的形像.這是教會應該顯露的見證.這兩個臉,在聖殿的牆上,廊子上,柱子上,門上,到處彰顯,被人看見,被人注意.神的形像,就是基督的樣子；基督彰顯神的樣子,而教會應當彰顯基督的樣子.教會裡面的弟兄姊妹,應該彰顯神的樣式,彰顯基督的樣式.基督是道成肉身,基督彰顯道；基督徒在次要的意義上,也是道成肉身,意思就是說彰顯道,這是重要的見證.是品格的見證,這品格的見證就是我們生命的事奉.單有工作的事奉,而沒有生命的事奉是不完全的；工作的事奉要加上生命的事奉,再加上本位的事奉.本位的事奉就是敬拜.我們站在被造者的本位上敬拜創道者,這是我們的事奉.本位的事奉──敬拜,加上工作的事奉；又加上生命的事奉,這三個因素造成完全的事奉.獅子的臉代表得勝,我們得勝之時就是見證；我們失敗之時就失去見證.人的臉,獅子的臉,應該到處可見,這是神對祂教會的心意.
\newline
\newline
(六)教會是充滿神榮耀的地方　聖殿在高山上,還有個很寶貴的意思.(四十2)提到聖殿在高山上的南邊.在中國北方,建屋最好的方向是朝南；因為冬暖夏涼,最好的陽光在南邊.(四十三4)說,聖殿的正門朝東.東面是日出之地,聖殿在光裡面；不但得到最早的陽光,也得到耶和華的榮光.無論是從自然方面說,或是從屬靈方面說；聖殿就是充滿光明的地方,有陽光光照,有神的榮光充滿.神的榮光是什麼?第一章末段告訴我們,神的榮光是普天上的榮光,是虹的榮光,是精金所發的光；是火的光,是藍寶石的光.精金代表神的神格.藍寶石代表神的聖潔,神聖潔的光,公義的光,神聖榮耀的光,神的愛光,合在一起；燦爛光輝,充滿在聖殿中.神的旨意,教會是充滿神榮耀的地方；非充滿人榮耀的地方,非充滿世界,充滿榮華富貴的地方；乃充滿神榮光的地方.
\newline
\newline
(七)教會最崇高　聖殿在高山上,高於一切,超越一切；神的心目中,教會是最崇高的.神的計劃就是教會,所以我們必須把教會放在第一位.我們的人生,我們的追求,我們的志向,我們的一切；應該重視神所重視的教會.感謝主!有很多人愛教會,真是愛教會,對教會有很多貢獻；把教會看作最重要,為主在教會裡面活著.在教會中很活躍,很活潑,很火熱,忠心事奉主.摩西在神的全家盡忠,各位弟兄姊妹!你用何態度看教會?重視教會呢?或是在教會的事奉可有可無呢?教會的興旺或冷淡你關心沒有?當教會冷淡之時,你心裡難過嗎?你曾為教會的興旺禱告嗎?你關心基督的國度嗎?你對教會的態度如何?神的心意乃要教會中的每位弟兄姊妹都愛教會,重視教會.主講天國的譬喻裡面,有一人羨慕珠子,變賣一切來買真珠.又有一人發現地裡有珠子,就變賣一切來買這塊地.人都笑他糊塗,別人不知地裡有寶貝；但他知道,他付了最高的代價,完全的代價,要得這個寶貝；因他知道價值.當我們知道教會價值的時候,我們願意為教會付一切代價.天國和教會的意義是相近的,我們為教會,應該有基督的心；基督愛教會,為教會捨已；我們愛教會,為教會作了什麼?
\newline
\newline
(八)從教會流出豐盛生命的江河(四十七1)　這裡我們看見有一條河,從聖殿裡流下；從殿門口流出,從高山上流下來；直流下四海,改變了四海的水.這是教會的一幅圖畫,在神的家中,應該有生命水流出來；恩典的江河,生命作為的江河,豐盛生命的江河,從教會裡流出來.越流越深越廣,滋潤全地；河的兩邊有樹木,每月結果子,年給十二次果子.不受環境限制,不受氣候季節影響；因為她的根吸收生命河的水.這是神對教會的心意,神要用教會作為賜福的器皿,讓祂的恩典,祂的救恩,祂的能力；從教會流出來,使世界苦水變甜.死海是死亡的地方,但生命的河水流進去,改變了枯萎；使死變生,生產魚類,各樣的魚生長,兩岸滿了漁夫打魚；成為豐富出產之地,有豐富的收成,結美麗寶貴超環境的果子.多麼美好!
\newline
\newline
(九)耶和華是唯一的主唯一的王　聖殿裡面有王,真正的王耶和華.(四十四1-2)說,聖殿最重要的門關閉不開,永不能開；因為耶和華已從這裡進去了.祂是至高者,祂進來與祂的百姓同在；作他們的主,作他們的王.以後,再沒其他的主,其他的王進來；祂進來以後,門封住,意思是說,耶和華是唯一的主,是唯一的王.城裡有王是人,不過是耶和華的代表,真正的主,真正的王是耶和華.教會是神的國度,神在這裡掌權；神的國不在乎言語,乃在乎權能,神的權能在那裡彰顯,那裡是神的國.神的權能如何彰顯?只有一個方法,神的權能藉著我們的順服彰顯出來.教會裡面,必須有順服,順服我們的主,我們的王.
\newline
\newline
現在我要用一刻鐘來講第十二座山.這座山也沒有名字,我們從(太二十八16)給它個名,叫加利利山.在這座山上,復活的主把大使命交託給祂的門徒.第十一座山和這座山的意義連在一起,一個合乎神心意的教會,就是個擔起大使命的教會；一定接受大使命,也成全大使命.
\newline
\newline
(一)大使命的重要　第16節說「耶穌約定的山上,」這件事是耶穌約定的事.上文告訴我們,至少有三次,主要門徒到加利利去；第一次在主復活之前.(二十六32)第二次是天使向門徒講話.(二十八1)第三次是主耶穌復活之後親自向門徒講的.三次重複的,鄭重的對門徒說,你們要到加利利去.這件事有如此鄭重的預備,表示這是主所重視的.「祂不在這裡,照祂所說的,已經復活了,......」(6-7)這裡有三件事連在一起.第一件事提到主已經復活了.「......你們來看安放主的地方,」第二件事提到主的死,主的埋葬.「快去告訴祂的門徒說,祂從死裡復活了；並且在你們以先往加利利去,在那裡你們要見他；.......」主的死,復活,大使命,三件事在同一句話中提到,這三件事都是最根本的事.
\newline
\newline
基督的教會,如果要符合神的心意,必須要接受主的死；必須相信主的復活,必須擔起大使命.何等重要!每個信主的人必須傳福音,必須擔起大使命的責任；這是我們永遠不能夠饒過的一件事.這話多麼清楚!由此可知,一個不關心大使命的教會,還沒有懂得主的心意,還有很大的虧欠.今天主等候所有華人教會,擔起大使命,就是差傳工作.
\newline
\newline
(二)大使命的根據18　「耶穌進前來,對他們說,天上地下所有的權柄,都賜給我了.」天上地下一切權柄都在主手中.接著,主說:「所以你們要去!......」(19)「所以」就是根據主的權柄,我們去,是根據主的權柄；如果主沒有權柄,我們就不要去,去了就會失敗,就會徒然.因為人能作什麼呢?人哪裡是撒旦的對手呢?但是這裡主說:「天上地下一切的權柄,都在我的手中；因此,所以你們要去,擔起大使命.」不但有基督的權柄,也有聖靈的能力；在基督的權柄裡面,包括著聖靈的能力.聖靈降臨時,要彰顯三種的能力:
\newline
\newline
(1)火像舌頭一樣降在門徒的頭上,那是言語見證的能力.
\newline
\newline
(2)一陣大風吹過.我們知道,風是代表強的生命氣息.這個解釋是從以西結書卅七章來的,在那風吹來之時,氣息就進到枯乾的骸骨裡面.風就是氣息,也就是聖靈；所以,靈,風,氣息,都是一樣；希伯來文也是同一個字,故意把這三樣連起來說的.聖靈降臨之時,有一陣強風吹過；代表強的生命氣息.這是生命見證的能力.他們說起方言來,這是聖靈特別的恩賜,這是聖靈恩賜見證的能力,言語見證的能力,生命見證的能力,恩賜見證的能力,這三樣都需要.感謝天父!祂藉著聖靈已將這三樣的能力都賜給我們；權柄又有能力,所以我們要去!
\newline
\newline
(三)大使命的方向　主耶穌說:「你們要去.」大使命的方向就是「去」,是外向的,神的心意不是要祂的教會內向.很多教會都屬內向,就像本堂,一切的努力,追求,奉獻,禱告,關心,目標,都是本堂.向裡面生,越生越縮；內生的教會是個縮小的教會.神的旨意,要我們去,向外看,舉目向天觀看.當我們舉目之時,我們的心胸擴大,眼光擴大,關心擴大,代禱的圈子擴大,愛心的圈子擴大,工作擴大.我們就為神作更大的工作,成為神手中更有用的器皿.
\newline
\newline
(四)大使命的範圍　「使萬民作主的門徒.」直到天涯地角,整個世界是我們的教區；是我們的天地.華人教會都應有此眼光.
\newline
\newline
(五)大使命的內容　一方面帶人信主得救,奉父子,聖靈的名施洗,這是認信；接受三位一體的神,和三位一體之神的工作.三位一體神的恩典,我們用信心來接受聖父,聖子,聖靈.第二方面,使萬民作主的門徒,不但叫人相信,第二步的工作叫人作門徒,從信徒到門徒.「門徒」是使徒時代最著重的名詞.在行傳中,「門徒」兩字用了四十次.「弟兄」用了廿五次.「聖徒」用了四次.「信徒」用了一次.「基督徒」用了一次.大家一聽便知,使徒行傳的重點在那裡.「門徒」用了四十次,充滿在使徒行傳中,這是在聖靈引導之下,在使徒的領導之下.這個選擇,這個名稱,代表重點.所有的信徒,相信以後作門徒,門徒是跟從主的人；每個信了主的人,要學習,要受訓練來跟從主.使徒時代的教會,是興旺發展的教會；不錯,這六個名字都有其重要性,有其意義,都不可忽略.弟兄,很寶貴啊!代表親愛.聖徒很寶貴,代表他們的聖潔.信徒代表他們的信心.會友代表他們團契的朋友.基督徒代表他們和基督有密切的關係.六個名字各有好處,各有不同重點；所以聖經在聖靈引導之下都用了；但是很明顯,重點是在「門徒.」我相信今天很多教會所以軟弱,乃因有很多信徒而缺少門徒.
\newline
\newline
(六)大使命的應許　「我就常與你們同在,直到世界的末了.」誰敢這樣講?惟有主的能力,主的全知全能,加上祂的全愛；給我們這個應許.我們執行大使命的時候,大使命的主,要與我們同在；大使命的主,與執行大使命的教會同在.
\newline
\newline
感謝主!祂的應許是給那些傳福音的教會.求主幫助我們!
\newline
\newline


\newpage

\section{第51屆~港九培靈會~滕近輝~第九講}
\label{sec:706}
\textbf{儆醒預備等候主再來}
\newline
\newline
連結: \href{https://www.hkbibleconference.org/session-message/view/706}{\texttt{https://www.hkbibleconference.org/session-message/view/706}}
\newline
\newline
\hyperref[sec:704]{< < < PREV SERMON < < <}
~
\hyperref[sec:toc]{[返目錄]}
~
\hyperref[sec:647]{> > > NEXT SERMON > > >}
\newline
\newline
在培靈會的十天當中,有九天是風風雨雨的；但是各位弟兄姊妹仍然來赴會,更加顯出各位愛主的心,和愛慕主話的心.有人為天氣禱告,有兩種禱告的方式:一種禱告是求主讓雨不要下；第二種禱告求主按你的旨意成就.如果是你的旨意,求你叫雨停止；如果你的旨意要下雨,求你給我們能勝過雨.第二種禱告是更成熟的禱告.感謝主!在任何環境中,祂也賜給我們得勝的力量.港九培靈會有個特點,就是沒有任何節目,只有聽神的話；所以凡是到主面前來的人,都是單純為了主的話而來.這是我們要感謝主的.求主使港九各教會都有主話能力的吸引.
\newline
\newline
第十三座山──橄欖山(太二十四2-3)主耶穌在橄欖山上,和門徒在一起,門徒提出兩個問題.第一個問題是關於聖殿被拆毀的事,他們問主甚麼時候有這事發生呢?「耶穌對他們說,你們不是看見這殿宇嗎?我實在告訴你們,將來在這裡,沒有一塊石頭留在石頭上不被拆毀了.」第二個問題:「你降臨和世界的末了,有什麼預兆呢?」是關於主再來的事.在主的答案裡面,這兩件事有時參在一起；所以我們讀時要留心,要把這兩件事的話分開；不然,就恐有解釋上的困難.耶路撒冷被毀,聖殿被拆除；和主第二次來,這兩件事合成了橄欖山上的信息.我們到了橄欖山上,對於主的再來有了愛慕,我們就儆醒,預備,等候主再來.我們的靈程到了高處之時,就更加愛慕主的再來,就更加儆醒等候主再來.
\newline
\newline
關於聖城和聖殿被毀,主耶穌這預言果然在主後七十年應驗了.當主講這預言時,任何人都覺得沒有可能會發生這事；當時的門徒是何等的欣賞殿宇的偉大.路加福音所載,更加清楚表達門徒羨慕稱讚的心情.當時的人感覺,這麼堅固的聖殿,這樣雄偉的建築物,絕對不會被拆毀；但是主的話終於應驗了.這件事情的應驗是個保證,保證主所講的話.
\newline
\newline
關於主再來的預言也必應驗　主講到聖殿被毀和祂的再來,兩方面的預言；人都看為不可能應驗；但是到了時候,主後七十年,聖殿果然被拆毀,按照主的話應驗了,這是個保證,保證主第二次來的預言也必應驗.今天亦有很多人不相信主將再來,覺得沒有這可能；好像當日門徒覺得勝殿不可能被拆毀一樣.但是,神的話必然應驗.馬太福音廿四至廿五章,是主在橄欖山上所說的話,其中絕大部份都是關於主再來的信息.可分為三方面:
\newline
\newline
(一)主再來之前的情形,廿四章告訴我們八種情形.茲簡單來領受:
\newline
\newline
第一種情形(5)　「因為將來有好些人冒我的名來,說,我是基督,並且要迷惑許多人.」主這個預言是雙關的；先指著耶路撒冷被毀時,有此情形發生；另外指著祂再來時也有這樣情形發生.教會歷史告訴我們,聖城被毀之時,有好些假彌賽亞,他們起來迷惑許多人,正應驗了主的話.在末世主快再來之前,也會有這情形,現在已經開始有了.在我書櫃有一本書,是位在美國工作的韓國人名叫文先明所寫,書內說他是彌賽亞；在美國有很多人擁護他,他已經在一百二十個大城市建立他的教會.美國的新聞週刊,時代週刊,數次提到他的事.他說他是彌賽亞,真是應驗了主預言假基督的話.我相信類似的事以後還會發生.
\newline
\newline
第二種情形是關於自然界(7下半)　「多處必有飢荒,地震.」這裡特別提到地震.各時代都有地震；但現在有個明顯的現象,地震情形越來越嚴重.我們舉個最出名的例子,中國大陸唐山的大地震,是有史以來傷亡最多的一次大地震,死者達五十至百萬之間.這裡又提到飢荒,目前好像不太顯明；但是在聯合國,有個特別部門,是要解決人類所面對的飢荒問題.有不少專家,認為這是個很嚴重的實際大問題；所以特別組織部門來解決.在啟示錄第六章提到四匹馬,第三匹是黑馬,代表飢荒.有兩種飢荒,一是糧食,一是燃油；等別提到油.我不曉得這預言和今日的油荒有否關係,我不曉得那預言是否指目前油荒而言,我們不敢確定,但是有可能；各種飢荒和缺少,都會在末世發生.關於大自的情況,還有一處經文值得我們注意.主在橄欖山上說的話,有部份載於路加福音卅一章,在25節特別提到末世之時,人因海中波浪的響聲就慌慌不定.海中波浪的響聲有何可怕呢?這讓我們想到在新聞週刊,過去曾經刊載一些科學家的報導；到了一九八二年,太陽系的各行星將排成一條直線,在太陽的同一邊,九個行星排成一條直線.科學家們很擔心,在那特殊情況之下,地球表面的水有很大變動的可能,將而造成空前的災害.這是科學家的預測.天文學家對星球位置的計算很準確,他們有此報導.
\newline
\newline
第三,治軍事方面的現象(6-7)　「你們也要聽見打仗和打仗的風聲......民要攻打民,國要攻打國......」打仗的風聲是冷戰,打仗是熱戰,冷戰熱戰都有.「民要攻打民,國要攻打國.」民和國有何分別?有人認為兩者之間沒有分別,這是希臘文中有道理的一種看法.但是,為何這裡要重複的講,這重複很可能有其意義,把國和民分開來看；不但有國際間的問題,也有種族間的問題.這正是今日的情形,今天種族之間的仇恨比國際間的問題,有時更複雜更棘手.政治和軍事方面,世界局勢的重點是在中東.這也是聖經預先告訴我們的,在啟示錄九章告訴我們,末世之時,軍事行動的中心是在中東.(14-16節)提到,爭戰的地區在伯拉大河,就是現今的約法拉底斯河,位在中東.這裡又提到,動員的軍隊,馬軍有二萬萬.馬軍在當時的意義,是機械化的軍隊；當時全世界的人口,不過二萬萬五千萬.但是這裡說,末世爭戰之時,動員的人數達二萬萬,這個數字正是如今的情況.若果十年前有人說世界軍事政治的中心是在中東,沒有人相信,而今人人都清楚看見了.我去年在印尼的加里曼丹,那地沒有中文和英文報紙,只有的印尼文報我不會看,簡直和世界隔絕；唯一聯繫的線是收音機.我每天收聽遠東福音廣播電台和英文新聞報告,整整一年,無日新聞不提中東.這正是聖經所載末世的情形.
\newline
\newline
第四,教會的情形(9-11節)　提到兩種情形,一是基督徒要受逼迫,許多屬主的人,要受很嚴重的逼迫.這正是現今的情況,今天在世界上很多地區,屬基督的人正受逼迫,在重壓之下.但感謝主!主保守他們,堅固他們；其中有許多人仍能站得住.叫我深深感覺保羅的話何等真實:「這福音本是神的大能,要拯救一切相信的人.」福音的能力,在最強壓力之下顯明出來.有一次,我讀華盛頓郵報,有一整頁是加拿大著名記者的報導.他用三個月時間,到蘇聯實地觀察教會情形,後來他寫了這篇詳細的報導.根據他的觀察和實際接觸,所得的結論說,現蘇聯基督徒人數在三千萬和六千萬之間.真是令人難以置信!實在出人意料之外!但這是事實的情形.若三千萬就是蘇聯全人口六分之一,六千萬是其全人口三分之一強,在壓力之下,宗教不自由情況之下,還有這麼多敬拜基督的人嗎?有,仍然有.這是見證.雖然這些基督徒中間,有的信仰有偏差；但基本上講,他們是敬拜耶穌基督的.這福音本是神的大能,在任何環境下,任何局勢中,任何處境裡,福音的能力仍然彰顯.求神給我們堅固的信仰!教會第二種的情形「且有好些假先知起來,迷惑多人.」(11)假先知就是冒神的名講自己的話之人；他們有地位,有學問,有口才；所以迷惑許多人.今日有多少理論在教會,其作用在破壞人對聖經的信仰,叫人離開福音的道路.何等可惜!這些都是主預先告訴我們的情形；所以我們在心靈中已有準備,這種情形必然發生,而正發生在我們的時代中.我們務要持守聖經的真道,勿受迷惑,在此潮流中,要站立得住,為真道作見證.
\newline
\newline
第五方面社會情況(12)　「只因不法的事增多......」這是今天的情況,全世界每個國家,罪案都不斷增加.記得數年前,在一個青年夏令會中,約有百人聚會,我在講道時問:「你們中間有多少人,或你的家人,同學,親友被搶過?」結果,舉手人數超出三分之一.這樣的情況香港人聽了都見怪不怪；不只香港如此,在印尼椰城,有位姊妹來聚會；在汽車途中,旁座者拿刀指她,命令把錢交出,到手後下車走了.你若看見這樣的情形,你不知道是椰城,還以為是香港呢!到處一樣,全世界沒一個在犯罪上落後的國家；有的在教育上落後,但在犯罪上卻不落後.恐怖主義,是今時代特別的產品,精通科學的人,正在幫助他們.主耶穌說,祂快再來之時的社會情形,和昔日挪亞時代洪水快來之時一樣,(37-38)這裡用喫,喝,嫁,娶,四個字來代表；換今天的話來說,就是物質主義.貪求享受是當時代人生的目標,也是他們人生的意義,是他們人生中最高追求.主說在祂快來之日,社會情形正是如此.物質主義充滿世界,連基督徒也受了影響.在東南亞某國,其中有不少華人教會；有間神學院,共有卅位同學,廿七位女生,三位男生.為何獻身事主的僅十分之一弟兄呢?因為該處華人所有基督徒作父母的,都用盡力量攔阻他們的兒子,奉獻給主作傳道,都要他們的兒子做生意.這是實際的情形,基督徒攔阻自己的子女奉獻,為了受這世界潮流的影響.末世之時正像老底嘉教會一樣,那時代,他們在物質上發達,說:「我們發了財,我們樣樣都有.」但主說,在事實上,他們不知道他們是貧窮的,可憐的,瞎眼的.主用這話很重.這是當時實際的情形.末世教會很多的基督徒,是基督寶血所救贖的人,是天父的兒女；但他們之人生目的乃物質享受.把十字架棄於一旁,把世界擺在前面,他們追求的對象就是愛世界.多麼可惜!
\newline
\newline
各位弟兄姊妹!在這樣的時代,這樣的潮流中,我們應當如何?求主愛激勵我們,讓弟兄姊妹都好像約書亞和迦勒一樣,另有一個心志要跟從主；因此,進入屬靈的應許美地.我們不欠世界的債去愛世界；我們欠主愛的債來愛主.
\newline
\newline
第六種情形,末世之時福音要傳遍天下,然後主再來.(14)　這是主說的預言.現福音已經傳遍天下,每個國家都有福音,所有基督的教會都有十字架的記號,都有基督的見證,證明主兩千年前講這預言.我再次說,當時聽起來,這預言不可能應驗；但是今天我們親眼看見主的預言已經應驗.因此,我們得著一個更強的理由,知道講這話的主,祂的話是可信的；祂已經給了我們許多的證據.讓我們可相信「主再來」的預言,這話不但是預言,也是挑戰!主再來以前福音要傳遍天下；這是教會的責任,是每個基督徒的責任.
\newline
\newline
昨天晚上,我們看見教會的大使命.感謝主!有百多位弟兄姊妹到台前來,願意全時間事奉主,要負起大使命的責任；在近處,在遠處,為主作見證,傳揚福音.如果主呼召你全時間事奉祂,你不要以代職事奉,來替代主的呼召；主對每個人有祂的旨意,你不能用其他方式來代替神的旨意.你順服之時,神的恩典,能力都要臨到你；神的恩典是藏在我們的順服裡面.求神給我們另有一個心志,是約書亞和迦勒的心志,他們跟從神,進入應許美地.
\newline
\newline
第七,敵基督的情形(15)　「你們看見先知但以理所說的,那行毀壞可憎的,站在聖地.」這裡提到先知但以理所說的預言,這預言說,敵基督要站在聖地,要出現在聖殿裡面.這敵基督是誰?我們從保羅的預言清楚看見,在末世要有敵基督出現.敵基督出現和十個國度有密切的關係.許多人想到這十個國度,都是偏向歐洲方面；但是我以為,我們應向中東方面著想；因為聖經的指示很清楚.啟示錄所說的敵基督是指著過去七個與猶太人為敵的大帝國,都是與選民為敵的.所以在末世之時,那個敵基督勢力的集團,一定仍然與神的選民為敵；那是中東地區有十個國.在他們中間,後來有領袖出來,站在以色列人所重建的聖殿中,自稱是神.最可能這件事的發生,就是一個回教的領袖佔領了聖殿,對以色列人說,你們拜神吧!我就是你們的神.啟示錄十一章提到重建聖殿,聖殿在末世出現；而且聖殿出現之時,以色列人要落在敵人的手中,遭他踐踏.所以中東的局勢令人擔心.今天以色列很強,我們不知道啟示錄廿一章的話,何時要應驗；那時以色列人要經過大災難.當然,對啟示錄的解釋,無一人是完全的,我們留下餘地；神的話是奇妙的.但是,這些可能性,相當清楚的顯出來.啟示錄的預言好像一塊鑽石,有名貴的面,每面有它的意義,有它的正確性；合在一起,乃是神的話全部意義.有人看見這面,有人看見那面；也許有一天,我們發現很多看的那一面都是在神的話意義裡面.我們保留可能性在心中,讓我們儆醒預備!
\newline
\newline
第八種情形,以色列的情形(32-33)　無花果樹在當時是象徵以色列人.主說:「無花果樹發嫩長葉之時,就知道夏天近了.」接著說,人子正在門口.意思很明顯,當末世以色列復國之時,即近主再來之時.兩千年前的話,在一九四八年應驗；以色列亡國將兩千年後復國,按照主的話應驗.這就給我們一個更強的證據,講這預言的主,祂的話是可信的.我們應該儆醒,預備,等候主再來.
\newline
\newline
主要我們儆醒預備,怎樣儆醒預備呢?在廿五章主題了三個譬喻,都是針對末世的情形講的.「那時」就是主再來之時,那時如何?主要我們按照這三個譬如來預備.
\newline
\newline
第一,童女的譬喻.
\newline
\newline
第二,管家的譬喻.
\newline
\newline
第三,羊的譬喻.
\newline
\newline
這三個譬喻給我們三個重點.
\newline
\newline
第一個重點,要保持我們對主的愛,我們以童女的身份來預備；童女意義的重點是愛.保羅說,我將你們當作貞潔的童女獻給基督,所以你們不要失去對基督純一的心.這是保羅所關心的事.主說,祂快要來之時,我們好像童女.我們要作聰明的童女,要預備好,要保持我們對主的愛,向主有純一的心.
\newline
\newline
弟兄姊妹,讓我們接受主這個重點.主再來之時,天國好像十個童女；你我都在其中,我們都預備好了；有愛主的心,純一的心向主.
\newline
\newline
第二個譬喻的重點,忠心事奉.每個蒙恩的人都是管家,主把祂的事工託付我們,把祂所愛的教會託付我們,把救恩的傳揚託付給我們；同時主也給我們恩賜,給我們力量,來完成祂的託付.但是,主要求我們一個條件,要忠心.忠心運用主所給我們的恩賜,在神的家中,努力事奉；發展你的恩賜,發揮你的恩賜,不要埋沒你的恩賜.這是我們儆醒預備等候的一個方式.
\newline
\newline
第三個譬喻帶來第三個重點,(二十五31至末了)羊的譬喻,有山羊,有綿羊；綿羊,是那些愛主而愛人的人.山羊,是那些口說愛主而沒愛人的.末世之時,因為愛主而向人表現我們的愛,這是我們的準備.
\newline
\newline
橄欖山上的三個譬喻給我們三個重點,清楚告訴我們應該如何準備.我們要準備一顆愛主的心,永遠不失去愛主的心,向主有純一的心,全心全意向著主.我們要忠於主的託付,做個好管家,為主而活.充分又充分的運用神給我們的恩賜,將來見主之時,就是主再來之時；我們可以將祂所託付我們的,另有所賺而獻上.同時,我們也不要忘記,主要我們因愛祂而愛人.末世之時也是教會彰顯愛心之時,在服務中彰顯,用許多方式彰顯.這三個重點合在一起,就是基督心中要我們作的準備；我們透過主的話,摸到了主的心.我們要有準備的路線,有準備的重點.
\newline
\newline
第十四座山是兩座山,一座是預表,一座是成全.
\newline
\newline
亞拉臘山(創八4)當日挪亞方舟在洪水之中,最後洪水停止,方舟停在亞拉臘山,這是預表救恩.俗曆七月十七日,方舟落在亞拉臘山上.俗歷七月十七日,等於以色列人聖歷一月十七日；這日正是初熟節,亦即主耶穌復活之日.請大家記住!俗歷七月十七日等於聖曆一月十七日,而一月十四日是逾越節,神的羔羊被殺,預表主被殺.過了三天,十七日初熟節,預表主復活；主在逾越節被釘十字架,第三天以後復活,正是初熟節之日復活.利未記廿三章提及以色列人的七個節期,表示神救恩計劃中的七部,預表意義非常顯明.任何人不可能推翻這章預表的意義,因為太清楚!太確定!逾越節就是主被殺之日,三天以後初熟節,正是主復活之日；主成了復活初熟的果子.過了十天是五旬節,聖靈降臨之時；七月一日吹角節是主再來之時,那就是利未記廿三,廿四章所載.那麼,我們要問:七月一日在洪水和方舟記載之中有否提及?請大家注意13節,正月初一日,就是七月一日；因為俗曆和聖曆相差六個月,七月十七日就是一月十七日,是初熟節,是主復活之日.一月一日是等於七月一日,是吹角節,那天是預表主再來之日；在新約裡面說,號筒末次吹響之時,主就再來.這是很清楚確定的一句話.吹角節預表主的再來,是一月一日即七月一日發生的事.13節說:「地上的水都乾了,」救恩完成了.救恩完成的第一步,方舟落在山頂上,脫離了水；主的復活是救恩的第一步；因主復活,我們得救了,我們真的得救了!有了救恩的證據,有了救恩的把握,主再來之時,救恩完成了.這裡清清楚楚告訴我們,一月十七日方舟停在山上；七月一日地上的水都乾了.奇妙嗎?豈是偶然的嗎?我們相信絕非偶然.神在祂的話語中,在在給我們留下很多證據；讓我們知道,祂的話是可信的,是超越我們所想的.如果人憑自己的理智來否定神的話；我再說,是很愚蠢的.
\newline
\newline
亞拉臘山的意思,就是完全的救恩,又是完成了的救恩.神給我們的救恩,是全備的救恩,完全的救恩；有第一步的完成,有最後的完成.主後死裡復活,祂的大能顯明祂是神的兒子,祂的復活,是救恩的保證,是歷史的事實.神在歷史的事實中,給我們確定的保證；主的復活,表示祂勝過死亡,勝過罪惡,勝過撒旦.祂是得勝的主,這是第一步的保證；還有,主再來是最後的完成,這話是保證的保證.
\newline
\newline
感謝主!我們所信的主是可信的,祂是創始成終的主.我們用信心來接受祂,我們信祂之時,就逐漸發現更多可信的憑據.我們以信心開始,然後我們的信心變眼見；我們經歷的越多,我們的眼見越多,信心越有把握,就進到更高層的信心裡面.神給我們眼見就是給我們經歷信心,證實信心.然後又帶我們進入更高的,不是憑眼見的信心.我們因有信心,就得到經歷；因有經歷,得到那更大的信心.一直增高上去,這是神在我們身上的作為.
\newline
\newline
感謝主!給我們留下許多的見證.
\newline
\newline
一對山的第二個(啟二十一10)「我被聖靈感動,天使就帶我到一座高大的山,將那由神那裡從天而降的聖城耶路撒冷指示我.」這裡提到一座高大的山,使徒約翰被聖靈感動,在這座山上看見新耶路撒冷,這是救恩最後的完成.這座山和亞拉臘山互相遙應,亞拉臘山是預表,是清楚的預表,確定的預表,奇妙的預表；後來啟示錄廿一章提到高大的山是最後的成全.
\newline
\newline
感謝主!我們將來所等候的是新耶路撒冷.廿一章二節說,新耶路撒冷從天而降,預備好了,就如新婦妝飾整齊,等候丈夫.新耶路撒冷的特點,就如新婦等候丈夫──就是基督.這是愛中的結合,新耶路撒冷和耶穌基督永遠在愛中合而為一.所有在耶路撒冷的人,都有這個特點,和基督在愛中結合.每個在耶路撒冷的人,你我都有這個特點；接受了主的愛,被主的愛得著,有個愛主的心.不是一時的愛,乃是永遠的愛.神以永遠的愛愛我們,十字架是最高的證明；我們被主愛激動,我們也以永遠的愛愛基督.我們和基督在愛中,永遠結合.這是我們的特點,是新耶路撒冷的特點；這個特點是全部聖經的特點,是永遠至永遠的特點.全部聖經最後的一句話:「主啊!我願你來.」請大家一同來讀這節聖經作為結束,作今年培靈會的結束.啟示錄廿二章20節是最後祝福前的一句話,21節是祝福,20節是最後的話.請我們一同來讀:「證明這事的說,是了；我必快來.阿們.主耶穌啊!我願你來.」我們再次讀下半節:「阿們,主耶穌啊!我願你來.」「阿們.主耶穌啊!我願你來.」
\newline
\newline


\newpage

\section{第52屆~港九培靈會~史保羅~第一講　愛是恆久忍耐又有恩慈}
\label{sec:647}
\textbf{第一講　愛是恆久忍耐又有恩慈}
\newline
\newline
連結: \href{https://www.hkbibleconference.org/session-message/view/647}{\texttt{https://www.hkbibleconference.org/session-message/view/647}}
\newline
\newline
\hyperref[sec:706]{< < < PREV SERMON < < <}
~
\hyperref[sec:toc]{[返目錄]}
~
\hyperref[sec:649]{> > > NEXT SERMON > > >}
\newline
\newline
多謝楊牧師和委辦們,讓我有機會在這歷史悠久的港九培靈研經大會中有份事奉,我真是感到萬分榮幸!聽說這大會是在廣州開始的,後來香港相繼效尤舉行,歷年來均蒙神賜福.本屆我能與焦源濂牧師,薛玉光先生配搭事奉這是何等大的福份.只可惜我的困難是聽不懂他們所說的話,假如我懂中文我就可以用中文和他們交通了.
\newline
\newline
請打開羅馬書五章,一至五節這裡說:「我們既因信稱義,就藉著我們的主耶穌基督,得與神相和.我們又藉著祂,因信得進入現在所站的這恩典中,並且歡歡喜喜盼望神的榮耀.不但如此,就是在患難中,也是歡歡喜喜的,因為知道患難生忍耐.忍耐生老練,老練生盼望,盼望不至於羞恥.因為所賜給我們的聖靈,將神的愛澆灌在我們心裡.」常有人問我:「基督教是甚麼?」我相信最好的解答是引用剛才所讀的經文其中一節半.
\newline
\newline
基督教的信仰可以分為兩部分.(羅五1)是基督教信仰的第一部份.這是指著因信稱義,藉著耶穌基督得與神相和好.換言之就是藉著耶穌得與我們永活的神接觸.如果不是藉著耶穌基督,就沒有下文了.基督教並非單指我們與朋友和家人的關係,乃是要我們先與神有正確良好的關係,然後才能與其他的人有正確的關係,基督教乃是要我們憑信,藉著耶穌基督與永活的神接觸作為開始的.當我們與永活的神接觸了之後,纔能將神的能力向其他人彰顯.一個基督徒要學習與家人,朋友,同學和同事有正當的相處；這就是基督徒信仰的第二部份.亦即(羅五5)下半節所說的:「所賜給我們的聖靈將神的愛澆灌在我們心裡.」這裡說當我們接受了耶穌為我們的救主之後,神便將祂的愛藉著聖靈放在我們的心裡面.各位當我們成為了基督徒之後,我們就與神的愛有接觸；在心裡面有平安,我的罪得蒙赦免,又得被稱為義了.其次是在我裡面有了神的愛,當我們讀到這一段經文時,便覺得它的確非常美好；在講員傳達這篇信息時,我們會說這確是個偉大的真理.「愛」是最偉大的神學,「愛」是一最好的教義.可惜大多數人,聽了道卻不肯去行；在何處聽道,仍把道留在何處.似乎在禮拜堂所聽的道,與我們日常生活毫無關係.我們在教會,詩班,主日學的時候是一種人；在學校,辦公室和路上又是另一種人.今晚我要向各位傳遞的信息,乃是神的愛,這愛澆灌在我們心裡,是鐵一般的事實.但如果神的愛,只在教會裡發生功效的話；那麼神的愛就簡直是未曾工作過.假若我們要活出一個真正基督徒樣式的話,我們就必須將神的愛在學校,家庭,辦公室,甚至在街道上也要彰顯出來.保羅明白到這一點,因他知道很多人雖在信仰和教義上很純正；但卻在日常生活裡,不能實行出來.許多人雖明白到羅馬書內所說的教義,可是在他們的日常生活裡,卻不能表現出他們是明白這些道理的.因此他寫了一封信給哥林多教會說；不能把他們當作屬靈的,只得把他們當作屬肉體,在基督裡為嬰孩的.(林前三1)故此他在林前十三章說明神的愛是怎樣發生功效的.由四至七節保羅對神的愛下了一個定義.指出我們怎樣可以在日常的生活裡,把神的愛活出來.如果你真的因信稱義,你就必須有以下經節所講的行為；如果神的愛果真澆灌在你的心裡,你便應該有以下經節所講的表現；如果我是屬神的子民,這段聖經便是為我而寫的了.
\newline
\newline
保羅開始就說愛是恆久忍耐,又有恩慈.這句話裡面有三點是我們要注意的.
\newline
\newline
(一)受苦　因為全世界的人都在受苦.多年前我在自己的教會,也以「忍耐」受苦為講道題目.有位姊妹收聽電台的廣播,聽到我教會的崇拜.便寫信給我說:「史保羅牧師,我上主日聽到你講論受苦的信息,我覺得你並不知道自己在說甚麼.因為你一生從未曾嘗過甚麼苦.你未真正捱過餓,你常有足夠的衣服穿,也住在舒適的屋子裡.可是我少時卻很窮,有時屋內從早到晚沒有一點食物,我的肚子時常是空空的.我家沒有足夠金錢來買衣服,所穿的都是由別人施捨來的,既不合潮流也不稱身；穿在身上使我們感到羞愧!為了避免給別人看見了而譏笑,我每天都要從後巷小路上學.當父親沒有能力付租時,我們便要整夜露宿街頭.史牧師啊!你知道甚麼叫做受苦嗎?只有我才知道肚子空空,穿破舊衣服和露宿街頭的個中滋味和感受.」不錯,一直以來我都吃得飽,穿得暖,住得舒適,其實這位姊妹不知道每個人都要受苦,只是情形不而已.有些人要受肚子空空的苦,這很慘!有些人要受心靈空虛的苦,這更慘!有些人受苦是因為背痛,每移一步都會楚楚作痛；多年來也不能挺起腰骨做人,這是很悲慘的受苦.另外有些人受苦是因為心痛,心痛比背痛更悲慘了.偶然有些青年男女來向我說:「牧師啊!我多年來都在受苦.我半生都在尋找一位女子與我結婚,但從未找到,我真是苦.」另外有些青年人卻這樣對我說:「我這一生都在受苦,因為我找到一位惡女子與我結婚.」有些人在受苦,因為他們找不到一塊居住的地方；另外有些人受苦,是因為他們不懂得在居住地方與人相處.今天晚上我盼望大家看出這點,我們總要受苦,只是各人的方式不同而已,成千上萬的人,無論是否基督徒,或其他宗教的人,懂不懂得心靈的平安或曾否嘗過主恩的滋味都在受苦.受苦這件事,本身不會使我們成為基督徒.
\newline
\newline
(二)忍耐,保羅說:「愛是恆久忍耐.」愛是讓你恆久地無限地「忍」.原文聖經「忍耐」這詞之後是沒有逗號的,可是有些人以為愛是恆久忍耐便夠了.我所能忍的已經忍夠了,甚至要握著拳頭在神面前說:「神啊!好了!到此為止了!我不可能再忍了!」
\newline
\newline
多年前我教會的一位長老,到我辦公室與我傾談.看他的樣子似乎正想發怒.我問:「占美!你的麻煩在哪裡?」他舉起拳頭用力槌在我台面上說:「牧師!我忍夠了!我對這兒子已沒有辦法了!我叫他做的事,他偏不去做.他晚上去街總是深夜才回來.我不知他往哪裡去,也不知他幹甚麼事.而且是一晚比一晚遲回來.昨天我終於對他說:『兒啊!若你今夜出外,我就要你在某一段時間內必須回來.』我向他說明某一段時間他若果不回來,我便會把門鎖上,而他就要另外找地方住了!牧師,我兒子昨夜仍然出了外,我不知道他往哪裡,只是到了我所指定的時間他仍未回家,我守候了多半小時後,便決定起來把門鎖上了.他果然是一去不返,我也不知他去了哪裡.」各位這長老的神學究竟出了甚麼毛病呢?他是一位聰明的好人,他明白聖經,他覺得對自己兒子的忍耐可以到某階段便停止.他當然知道一個十來歲的孩子會有很多問題,十七歲的兒子與四十歲的父親對事情的看法會有不同.年輕人有時會有反叛和不服從的心理.不過他竟以為自己有權說:「我對這兒子的反叛只能忍到此為止!」我有權對神說我容忍到某程度,便停止不再忍了嗎?大家可留意到這位長老忘記了神的愛,並不是這樣的,神的愛子能說我只能到此為止,乃是要我們按神的旨意做多少便多少.我無權來決定自己究竟要容忍到哪裡.也許是一星期,一個月,數年甚至一生之久.
\newline
\newline
(三)恩慈　直到現在我尚未能描繪出,一個基督徒的真正形態.保羅最後加上一點說:「愛是又有恩慈」.請留心「恩慈」這字在聖經上不是單獨存在的.「恩慈」與「恆久忍耐」是有一個連接詞,連在一起的.似乎恩慈是指向那受苦忍耐的來源.只說神的愛是有恩慈是不夠的.我們是要看到神的愛,指望著受苦的源頭,而常存著恩慈,神的愛是使我們面對那些令我們受苦的人,而存著恩慈.你可曾記得勝經上說若你的仇敵饑餓你會怎樣做呢?是否趁這時踢他一腳呢?聖經說:「不!如果你的仇敵餓了,你就給他吃.若渴了,就給喝.」(羅十二20)你若真是主的門徒,便要向你的仇敵存恩慈,好好地對待令你受苦的人.」在耶穌時代,巴勒斯坦地方,是由羅馬帝國所統治的.那時有一條法律可稱為「一哩路法律」就是羅馬士兵有權要求,在他管轄範圍以內的男子為他做事.這士兵有權命令某一男子立即放下工作,為他擔起擔子走路.不過這士兵只能夠要求那人陪他走一哩路,在這一哩路終結後,就必須讓這人返回原地工作.但你可以想像到,當時的猶太人已對這條法律何等的憎恨,何況耶穌還教訓他們說:如果下一次有人要求你們行一哩路,你們不要向那士兵發怒,反而要帶著和氣微笑地望著他說:先生我非常喜歡背負你的重擔啊!我今天橫直也不太忙.隨即便歡喜快樂地挑起這些重擔而行.當你走完一哩路後,放下你的重擔然後休息一會,再對那羅馬兵笑笑:『我與你同行剛才的一哩路感到非常好,你反不反對我再為你背負第二哩路嗎?』假如初期的基督徒,確實是這樣履行的話,那麼當他行到第二哩路的終點時,必定是很容易向那羅馬兵,打開話題說及耶穌基督了.為甚麼在我們當中,有些基督徒沒有機會為主作見證呢?恐怕是因為我們只顧履行別人對我們最低限度的要求,我們只走我們被要求走的路程,卻不願走到第二哩路的終點.但神的愛是要我們走到第二哩路的終點,這也就是向普世佈道的起點.亦即為主作見證的開始.
\newline
\newline
晚上開始聚會的時候,我提到神的愛若不能在家中實行出來的話,恐怕你家裡就有問題發生了.但每個人總會有一些煩惱,因為兩人同住在一起便會有難處.我妻子生於北愛爾蘭,雖然她在七歲時便全家遷往加拿大,可是婚姻卻仍沿用北愛爾蘭方式.她一向確信一位好丈夫要負責倒垃圾.我認為這簡直是太荒唐,太侮辱了!我不喜歡倒垃圾,而她偏要我倒,你說我該怎樣做呢?我是一個真實的基督徒,而她卻要我倒垃圾?在我所住的地方,倒垃圾的方法是把垃圾盛入垃圾膠袋的.只要把垃圾袋放入屋外收集垃圾的大鐵桶內,市政事務處自然會按時來收集垃圾了.我的責任是把這膠袋從屋內,拿到屋外的大鐵桶內.每次倒垃圾時,我會從她手中接轉包拉圾說:「算了!我到此為止夠了!」然後用沉重大力的步伐,走往大鐵桶那裡,用力一擲,隨即返回教會工作.如果我是基督徒我便不能如此.當她給我一大包垃圾的時候,我反要微笑,我要一面拿著垃圾,一面唱著福音短歌,好好的把垃圾放在大鐵桶內.這就是第一哩路.我若要做第二哩路的基督徒時,我就應該回家對妻子說:「太太,我已倒了那袋垃圾,你可否讓我再為你清潔房子呢?」在我們當中已結了婚的弟兄姊妹們,請問你們在婚姻的路上行了多少的第二哩路呢?做丈夫的請問你走第二哩路之後,到現在又隔了多久呢?作妻子的請問你上次走了第二哩路,到現在又相隔了多久呢?作青年人的,你走過第二哩路到現在有多久沒有再走過呢?這就是基督徒與非基督徒的分別!基督徒是憑信和耶穌基督與永活的神有接觸,基督徒具有聖靈所澆灌的愛在心裡；基督徒是要在家裡彰顯出神的愛來.愛是恆久忍耐,又有恩慈.讓我們低頭祈禱.
\newline
\newline


\newpage

\section{第52屆~港九培靈會~史保羅~第二講　愛是不嫉妒不自誇不張狂}
\label{sec:649}
\textbf{第二講　愛是不嫉妒不自誇不張狂}
\newline
\newline
連結: \href{https://www.hkbibleconference.org/session-message/view/649}{\texttt{https://www.hkbibleconference.org/session-message/view/649}}
\newline
\newline
\hyperref[sec:647]{< < < PREV SERMON < < <}
~
\hyperref[sec:toc]{[返目錄]}
~
\hyperref[sec:651]{> > > NEXT SERMON > > >}
\newline
\newline
哥林多前書 13:4-13:5
\newline
\begin{longtable}{cl}
\hline
\hline
章節 & 經文 (和合本修訂版)\\
\hline
13:4 & \begin{tabularx}{0.7\textwidth}{X} 愛是恆久忍耐；又有恩慈；愛是不嫉妒；愛是不自誇，不張狂， \end{tabularx} \\ \\ \relax
13:5 & \begin{tabularx}{0.7\textwidth}{X} 不做害羞的事，不求自己的益處，不輕易發怒，不計算人的惡， \end{tabularx} \\ \\
[1ex]
\hline
\hline
\end{longtable}
請翻開羅馬書五章一至五節:「我們既因信稱義,就藉著我們的主耶穌基督,得與神相和.我們又藉著他,因信得進入現在所站的這恩典中,並且歡歡喜喜盼望神的榮耀.不但如此,就是在患難中,也是歡歡喜喜的,因為知道患難生忍耐,忍耐生老練,老練生盼望,盼望不至於羞恥；因為所賜給我們的聖靈,將神的愛澆灌在我們心裡.」我提議將這五節經文濃縮為一節半,就是將第一節與第五節連在一起:「我們既因信稱義,就藉著我們的主耶穌基督得與上帝和好.所賜給我們的聖靈,將神的愛澆灌在我們心裡.
\newline
\newline
羅馬書第五章,說明當我們成為基督徒時所會發生的事:我們因信稱義,得與神和好,聖靈就將神的愛,澆灌在我們心裡.
\newline
\newline
林前十三章告訴我們,這等人的生活是怎樣的呢?羅馬書第五章是論到我們與神的關係,林前十三章是論到我們與人的關係,可見從這兩節聖經的話,我們必預先與神有好的接觸,才能與人有良好的相處.
\newline
\newline
有時我想到基督徒也有兩種,第一種可稱為羅馬書第五章的基督徒,可能他們在神學上是正確的,在教義上是十全十美的；一切關於聖經的知識他們都有了,只不過他們不能與人相處,在生活上完全不像基督徒.
\newline
\newline
第二種是林前十三章的基督徒.他們的生活正如林前十三章所描繪的,他們已熟讀一切關於愛的美好事情,和聖經的教訓,怎樣彼此相愛.心裡說:「這是對的,我要愛世上所有的人!」
\newline
\newline
你們當記起在六十年代,大部份青年人都在說同樣的話,在加拿大青年人唱這首歌:「這世界需要的是更多的愛!」(What the world needs now is more love)在美國大巡遊時,也唱過這首歌,並在顯眼的地方,亦貼上這些徽號.
\newline
\newline
這些青年人說這世界委實太多仇恨,戰爭和怒恨了,現在我們正需要的,是彼此相愛!他們所說的是絕對正確的,但以一個普通人來說,世上所有的人,是絕不可能的；因為世上每一個人都不是好相處的.在你所居住的地方,工作的場所,就讀的學校,甚至所參加的教會,你可能會遇到很多難以相處的人.有人以為作為一個人我實在沒有能力,去愛那些那不可愛的人.也許我閱讀完林前十三章後,會說我一定要這樣做:明早起床,我要開始愛所有的人.當第二天早上起床,我第一個見到的是我的妻子,愛我的妻子,並不困難；因她是一個可愛的女士.之後,我見到我的孩子,我有一個二十九歲的兒子；記得當他讀中學時,正是「披頭四」流行的年代,他的頭髮比較長一些.所有的中學生都認為這樣才對,頭髮長才是醒目；我不把他的髮掉,因為他比我重五十磅,肌肉比我強壯.況且這與聖經無關,我愛他並不困難.我有一個二十六歲的女兒,當她讀中學時,女孩子都穿著短裙；在加拿大有這樣的一句說話:「若你准你女兒只穿一吋,她也有膽量只穿一吋.」我還有一個女兒,她差不多十全十美,她外表非常美麗,又聰明,是一個可愛的女孩,因為她像父親嘛.所以我愛她並不困難,我差不多忘記了家庭裡還有一個成員,名叫「柏」,牠是一頭(愛爾蘭)狗.在三月十七日出生,牠經常對人非常友善.有時當「柏」不聽話,我就向牠發怒,打牠；牠就把尾巴放在兩腳之間,走到牆角去.我放工回家時,她又非常高興地,向我微笑舔我的臉；我家裡沒有其他人會這樣做,所以愛「柏」這頭狗,對我來說並不困難.
\newline
\newline
當我離家上班工作時,我便遇到一些世界上最難相處的人,我不能愛他們如同愛我的家人一樣.我作為一個人,實在沒有能力,去愛這些人.
\newline
\newline
我雖然要了羅馬書第五章的經歷,但我不能夠做到,林前十三章所說的話；因為我已沒有那種愛.我若要過那種愛的生活,就必須要與神接觸,所以要在羅馬書五章一至五節註明:「參閱林前十三章四至七節」又在林前十三章四至七節旁寫上:要參閱羅馬書五章一至五節」.我們會察覺到,一個真正的基督徒,不單是羅馬書第五章或林前十三章的基督徒；而是要必須將這兩段經文放在一起的.昨天晚上,我們談了第四節的上半節說,當神愛澆灌在我們心裡,我們會有甚麼行為表現呢?
\newline
\newline
(一)愛是不嫉妒　今天晚上我們會查下半部,這部份說到神的愛能令我們做到生活行為的三方面:首先,愛是不嫉妒.嫉妒就是因為不滿別人的擁有和成就而引起的,惱怒,是因為別人得到了,而自己沒有,所以心中起惱怒.在舊約希伯來文,「嫉妒」的字根是焚燒的意思.很多時不同國家的文化,也用同一的意思,來表達某些事情.譬如:我們向別人發怒,就會說「怒從火中燒」,或被人激怒了會說「火滾」；「或是他惹起我把火」.妒忌的意思,在拉丁文聖經妒忌的音是「現斯驕」意思是緊望著一些東西,而巴不得那東西受到損害.這個字在掃羅與大衛之間的事曾經用過.在當時掃羅王是一個戰場上的偉大領袖,年青的大衛,曾跟隨掃羅王出外爭戰；戰爭完畢後,人群為著這些戰士的勝利凱旋,而夾道大聲歡呼:「掃羅殺死千千!」掃羅非常歡喜,非常陶醉這些說話,心裡非常高興.然後當大衛經過,那些人竟然大聲歡呼說:「大衛殺死萬萬!」聖經跟著說:掃羅回轉身來,目露兇光在大委身上,想要把他除掉.各位見到嗎?掃羅嫉妒的原因,是因為大衛有超越他的才能,這就是使徒保羅在這段聖經用妒忌的意思了,只看到一些東西,又希望得到並不是妒忌.有時別人有些東西,我們沒有便說:「巴不得我也可以得到.」這思想不算是錯的,但譬如一天早上,我在家裡看見鄰舍買了一輛簇新的汽車,我便對太太說:愛妻啊!請看鍾先生的新車吧!我也希望有一部如鍾先生一樣的新車啊.這不是妒忌,只不過是表示你有些志氣；因為很多人也像你,希望能買到一部新車.又假如有一天我看見鍾先生的新車,而指給太太看並說:「你看鍾先生買了一部新車,我的妒火從中燒了!」我遂發怒,因為鍾先生擁有一部新車.這就叫做妒忌,聖經說當我與神有良好的關係,神的愛就在我心中,就不會因別人所擁有的東西或好處而生妒忌了.在多倫多民眾教會有很多客座講員,我是負責邀請講員的,有時,我所邀請的講員,講得並不好,會令到我感到尷尬,有些講員當他開始講時,就知道他是有恩賜的,他對神的話語有極透徹了解；在傳達真理上極有恩賜.我看見會眾聽道的神情,就知道神的靈,已在會場當中運行了.如果這時我心裡羨慕說:「我假如也能講得像他這麼好便好了.」這就不是嫉妒,只不過有志氣而已.效法一個人,羨慕一個人,以他為榜樣,這並非錯的.可惜今天的大眾傳播媒介,已把所有人拉到不分彼此的水平了.我們找到一個人雖然在某方面有很大的成就,但我們同樣地會發現他在某方面的缺點；結果青年人就找不到一個好模範.年青人呀!讓我今天晚上告訴你,若你在生命上找不到模範的人,恐怕你的人生難以成功.如果你想做一個好醫生,或者是個好律師,你就必須有一個好模式,知所效法.讓我們效法神的僕人,我認為這就是神在聖經上,給我們這麼多偉人的其中一個原因.我們讀了摩西的事蹟,好叫我們能夠像摩西；讀到以利亞的生平,好叫我們像以利亞；我們讀到馬大,馬利亞的故事,好叫我們可以學習她們.聖經裡面充滿了許多蒙神賜福者的事蹟,可以使我們效法.當我有志氣要效法某一位講員時,這事的本意並非錯；但如果我聽一位偉大的講員講道時,心裡感到不滿,我發現他講道遠比我好,他傳揚神的訊息的能力是我遠遠不及的；這樣我便向他發怒.他講道比我更好,是我不能接受的事實,我真巴不得在講道時,有些事情騷擾他就好了!譬如有一個嬰孩大聲哭起來,他的聲沙啞了!他口疲乏起來；他說一個有趣的故事卻沒有人笑,或是其中一位會眾突然心臟病復發,需要即刻抬出會堂......誰知他一路傳得非常順利,而神又賜福他,存這種心理對他的惱恨便是嫉妒.使徒保羅說如果上帝的愛澆灌在我心裡,我們就永不會對別人有嫉妒的感覺.
\newline
\newline
(二)愛是不自誇　一個屬神的人不會四處招搖,給別人看出他是一個何等偉大的人.他不會在別人面前使人覺得他的重要性,他不會將神所賜給他的恩賜來抬高自己；他不是一個在別人面前,自我吹噓的人.不過我們不要對謙卑有所誤解.倘若以為謙卑是四處說自己甚麼也不會,這不過是虛偽的謙卑；因為神給每一個人總有些恩賜,只是我們所有的恩賜都不儘同.我們要憑著信心,使用這恩賜,讓別人因此蒙福；而不是使用這些恩賜,自高自傲,產生自我重要的感覺!記得上兩年葛培理博士,在我的城市主持一次佈道大會,當他在多倫多時,神給我有機會和福份與葛牧師相敘數小時,我提及神給他的恩賜作為佈道家,我說:「葛博士,神給你一個偉大的恩賜,你死怕是世上最偉大的佈道家了!」然後,譬如葛培理博士這樣的回答我:「其實我並不怎懂得講道,說到佈道家我亦不曉得佈道!」若葛牧師這樣說就完全不是謙虛,這是一個謊話,因為各位都同意葛博士有佈道的恩賜,他真是有講道的恩賜,神不會期望葛培理牧師說:「我沒有啊!」神都期望葛牧師能夠用這些恩賜將榮耀歸神,使全世界的人蒙福.保羅所說的正是這道理.我們應當憑著極大的信心,使用這些恩賜,為要叫神得著榮耀.不要自己去奪取利益來利己,而應當使用所得的恩賜來幫助他人.我要一生使用這些恩賜上,好叫許多人因此蒙福.
\newline
\newline
(三)愛是不張狂　有些人歡喜誇張自己,以至成為張狂.他們誇張自我重要性,似乎許多人都有個大狂,在我們裡面有個脹大了的感覺.這裡有兩件事是非常相似,「自誇」是別人可從我們外表看出來的:「自我張狂」是我們對自己內部的看法,我們試圖去感動別人,讓人知道我們自己的重要.我相信這是最好的比喻.今晚我手上有一個大紅汽球,它的直徑有二十吋,當我舉高它時,每個人都會看到；甚至現時開閉路電視傳播的也看到了,將一個大的紅色汽球高舉時就可以看出它的重要性；這個汽球已吹大了,當舉起它時就說:「看我有多大,多麼重要,我是鮮艷的紅色,我又大又美,人人都應該看我.」我們有汽球就是這個原因,但是用一口針來刺這汽球時,人人都會知道它的命運了,漏了空氣的氣球,剩下的會是甚麼呢?是殘餘,破碎,沒有用的膠.一個汽球被吹大了,就給人一個自我重要感；在它裡面有一個自我重要的張狂,現在吹大了,要讓人看見它的重要性.張狂的事,就是四處招搖,要給人知道；因為他對自我重要,有一個大的感覺.張狂是在外面要儘量給人看見,因為他們裡面只有一個自我重大的感覺.如果我們裡面充滿了神的愛,就不會這樣；我會用神給我的恩賜,讓它賜福與他人.因此,就不會對自己有一個重要感了.如果知道我的一切都是由神而來的,無論我有的是甚麼,無論我能做甚麼,都是神給我的；如此,我就應該使用神給我的恩賜,來造福別人,好叫神得榮耀.讓我們思想今晚研讀的經文,我希望知道,我講道的時候,各位正在想甚麼?聖經說,愛是不嫉妒,舉凡因別人擁有而發怒的人；因別人有所成就而眼紅的人,那就不是愛了.我們聽了,覺得真對,這段聖經真是對!因它說中了我很多朋友正是這樣.聖經又說,愛是不自誇,不張狂；愛不是在別人面前炫耀自己,誇張自己所做的.但世界上的確有很多這樣行事自誇,張狂的人.若整晚聚會你只是想到別人的過失,就浪費了整晚寶貴的時間.
\newline
\newline
這段聖經所指的對象,不是在會場以外的,也不是指今晚沒有來的人；你知道神正在向誰說的嗎?乃是向你和我說的.反正是基督徒在人生上的反映,也是基督徒常遇到的問題,隨時都可以發生在你我身上.在我們當中沒有一個人是經常嫉妒的,但我們每一個人在某些時間上,會有嫉妒的傾向.沒有一個人會給別人看到自己的重要性,但我們每個人間中都有這個傾向.論到張狂,可能沒有一個人會經常四處去表現自己的重要性,也沒有一個人經常自我吹噓,以至別人說我們張狂的.有些時候我們雖有這樣的傾向,這些事可能我史保羅也會做的.如果自誇和張狂,成為一個試探時,我應當怎樣做呢?若別人有,而我卻沒有的時候,因此我就向他發怒,我當怎辦呢?當我將神給我的恩賜,用來表現自己的重要性時；我當怎樣辦呢?正在這一剎那,我就會離開了林前十三章而走回羅馬書五章了.所以我必須與永活的神重新建立關係!我今日成為何等樣的人,完全是靠神的恩典.我們伏在神面前,求神寬恕我,我對神說:神啊!我因信稱義,雖然我是罪人,祂卻使我稱義；祂赦免了我所有的罪；神將平安放在我心,祂將祂的愛藉著聖靈澆灌在我心裡.無論我有甚麼,都是靠賴神的恩典；我成為何等的人,亦是靠賴神的恩典.主呀!求你赦免我對別人的妒忌,求你赦免我錯用你所賜給我的恩典；我當再與神重新建立關係,我必須活在林前十三章裡面.兄弟姊妹!你們必須先有羅馬書第五章的經驗,然後你才能把林前十三活出來.你要記著要藉主耶穌與永活的神有接觸,你要確知因為神的愛澆灌在你心裡,你才能與別人有正確的關係.
\newline
\newline


\newpage

\section{第52屆~港九培靈會~史保羅~第三講　愛是不作害羞的事不求自己的益處不輕易發怒}
\label{sec:651}
\textbf{第三講　愛是不作害羞的事不求自己的益處不輕易發怒}
\newline
\newline
連結: \href{https://www.hkbibleconference.org/session-message/view/651}{\texttt{https://www.hkbibleconference.org/session-message/view/651}}
\newline
\newline
\hyperref[sec:649]{< < < PREV SERMON < < <}
~
\hyperref[sec:toc]{[返目錄]}
~
\hyperref[sec:652]{> > > NEXT SERMON > > >}
\newline
\newline
哥林多前書 13:5-13:5
\newline
\begin{longtable}{cl}
\hline
\hline
章節 & 經文 (和合本修訂版)\\
\hline
13:5 & \begin{tabularx}{0.7\textwidth}{X} 不做害羞的事，不求自己的益處，不輕易發怒，不計算人的惡， \end{tabularx} \\ \\
[1ex]
\hline
\hline
\end{longtable}
今晚讓我們注意哥林多前書十三章第五節.在這節聖經的開始,保羅說到三點真理,第一,愛是不作害羞的事.第二,愛是不求自己的益處.第三愛是不輕易發怒.如果我們將這節聖經的原意完全的保留,但加上一個連接詞在這幾句說話之間,就會變成因為愛是不求自己的益處,所以就不作害羞的事,故此不輕易發怒.」或者我用一句淺白的說話來說,愛是經常使你有好的禮貌,從來不會粗魯,不會失禮人的.也就是說,愛是令你的行動,不會有傷害別人的感受.我想最容易明白這道理的,便是讓我們用吃飯時的禮貌,來作個比喻；我發覺世界各處的地方,各有不同的禮貌風俗.在某些地方餐館台上你可以這樣做；在另一些地方,你不可以這樣做.我們可以用某種方式,來吃某種的食物,但在某處卻要用另一種方式來進食其他,那要視乎你在那一個國家而定.
\newline
\newline
(一)愛是不作害羞的事　每一個誕生到世界上的嬰孩,都是用他們的手來拿食物的；他們需要一段長時間適應；漸長,才能學習用刀,叉進食；長大了才能學習我們進食的習慣.年輕人!如果明天你和我共進午膳,我們吃的是生炒牛肉飯；我們對坐,我很餓,於是不用筷子或刀叉,便用手來吃,吃得雖飽,味道無窮.你說:「史博士!」你可否讓我吃一些?我便立刻用手抹了才上的東西給你吃,你見我這樣,便會感到惡心作嘔,立即離開餐桌.各位看見嗎!我的好禮貌對我本身本無關係,我之所以需要有好禮貌,為要與別人和睦相處；使我不會影響別人對我的感受.這就是保羅所教訓的,我們基督徒的要在世界上作基督的樣式,和好的見證.如果我們的行為使人反感,我們便虧欠主的名,不能為主作美好的見證.這樣我就不能與我同住的人,接受我所傳的福音了.不是他們反對主耶穌基督,或反對聖經所說的,只是因為他們不能接受我的行為.
\newline
\newline
當宣教士到另一個不同文化的地方,他們便要學習在不同地方,怎樣和別人相處；如果這基本的功課也學不會,他便沒有機會傳揚福音.無論他對聖經如何熟習,就讀神學院有多少時間,甚至得過很多神學士的學位,也是徒然的.曾經有兩位美國人寫了一本書,名叫「醜惡的美國人」其實也可稱「醜惡的加拿大人」「醜惡的俄國人」,「蘇聯人」或「中國人」.世界每一處地方,也有這些醜惡的人.這本書的中心,說到一群外交家,他們去到一個細小國家；在那處代表自己的國家,他們建築一個美國式的小圈子；他們不肯花時間去學習當地的言語,文化,風土人情和當地的習慣.這群失敗的外交家,他們所做的,沒有帶給美國有甚麼好處,反而羞辱了自己的國家.在這段聖經裡面,保羅說,你不要做一個醜惡的基督徒,不要活在世界上,而使人對基督有反感!你必須顧及別人的感受,你必須對他們有禮貌,能與他們和睦相處.我們有最偉大的信息要宣告,如果你不學習與別人相處,便永遠不能宣告這個信息.
\newline
\newline
(二)愛是不求自己的益處　在我們的世界上,有兩種哲學思想,是世俗和耶穌基督的思想.世人的哲學思想如下:我活在世上,便盡在這世界上搾取,讓自己得益處.在選擇一行職業時,我是盼望得到一份高薪的職業；希望我的工作,能給我很多假期和更多的利益,這就是我所盼望的.最初門徒遇見耶穌的時候,門徒對成功,繁榮,權力,受人擁戴,十分有興趣；他們盼望主耶穌建立一個王國,每人都問耶穌,在你的國度裡,我們中間誰是最重要的人物呢?誰是坐在你的右邊和左邊的呢?他們說:主啊!我們自己能夠得到甚麼呢?當你建立起王國時,我們會得到甚麼呢?耶穌對他們說:你們想錯了!你們忘記我來世上的目的是要服事別人,不是要去去搾取益處,乃是要將自己獻出.耶穌基督又說:你若要作我的跟從者,就必須有這樣的思想；我作為你們的主,是要服事他人,那你也要服事他人.在我們的一生的日子裡面,我必須要問,我能夠給甚麼世人呢?今天我與世人相處,我可以貢獻多少給他們呢?我能為世人留下多少益汗呢?這便是保羅要傳的信息.
\newline
\newline
你們曾否記得前兩晚,我所說的信息,是保羅描述一個得救的人,一些因信稱義與神和好的人；在生命得到聖靈,聖靈將神的愛,充滿了他心.保羅說:我們若因信稱義,與神和好,而聖靈將神的愛充滿我內心時,主就會祈望我們過這樣的生活；我們的人生,是在乎我們能夠為世人留下多少.
\newline
\newline
各位今天晚上,當你鑑察你的生命時,我想請問你,你為甚麼而生活呢?在你的事業上,你所尋求的是甚麼?你基本人生所關心的是甚麼?你是跟隨別人這種思想嗎?或是你採用耶穌基督的人生哲學,你今天從事的職業是否對他人有貢獻?如果聖靈將神的愛充滿你內心,我們的人生就不會尋求自己的益處了.可能你說:「史牧師呀!其實在這世界每人都去搾取呀!」我知道世人真是這樣做,他們因為是屬世的人.神的子民,並不盼望過這種的人生,我們自問是否真的屬神的人,我們是否有正確基督徒的人生.
\newline
\newline
(三)愛是不輕易發怒　淺白的解釋,愛從來不會激氣的.你是否發現在你教會中,多少人給人氣得發抖呢?在民眾教會,就是我所牧養的教會裡,我們每一行通道,都有一些招待員.有時候,我覺得那些招待員,以為通道是他們的,如他們遲到,發覺別人站了他的崗位,便會生反感.某些詩班員,長久坐同一位置,漸漸地自以為這張椅是他們的；如果有一天他們發覺別人佔了他們的位置,便會不高興.或者教會有一位獨唱者,看見牧師忘記宣佈獨唱的程序,她回家後便很生氣.
\newline
\newline
保羅說:如果神的愛是充滿了我的內心,無論發生甚麼事,我總不生氣.我相信受激怒有兩種醫治的方法:一是保羅在聖經上說明,我不是在世上要搾取益處歸自己,我就永遠不會生氣.當我飛來香港主領這個培靈研經會時,可以存著兩種心情,我想究竟這個培靈會怎樣接待我呢?是否有足夠時間讓我講道?當他們介紹我的時候,會否說很多話來讚賞我呢?我究竟在今次培靈會之後,會得到甚麼益處呢?如果我存著這樣心理來香港,當我抵步後的數小時,我便很失望!因為各位見到嗎?我所盼望的得不到,我便失望,起反感,我會哭和激怒.我來香港第二種心情,我要去那個偉大的培靈大會,親愛的神呀!求你給我完成我的工作,我能夠將甚麼福氣,留給香港的弟兄姊妹?求你給我知道,怎能使弟兄姊妹在生活上得到神的福份；給我知道怎樣將我自己留下給香港人,求你給我知道怎樣將一些貢獻留給這個城市.各位!假若我存這個心情來香港,我便不會激怒,我根本不為求自己的益處而擔憂.要是真的得不著甚麼,也不會反感.有第二個方法使我們避免這反感心情,就是我們要投身在使命上；將這個使命和我自己來比較,便會覺得這使命偉大得多了.我們曾否想到世界上的政治家,怎樣可以忍受每天這麼多的批評.英國的首相常受到批評,無論他做甚麼事情,最少有百分之五十的人民都說他做錯；美國總統也同樣在美國各地受到大肆批評,報紙刊登的,都是在批評他.一個政治家怎樣可以面對這麼多的批評,為甚麼他們又不會哭起來,離開政壇呢?因為他們熱心於一個使命,這使命和他們本身來比較,要偉大得多了.有時候,對他最重要是他的政治黨派.他說:「我不理會對我本身的說法,只要我的政黨繼續前進,我就快樂了.」有時候,那政治家很盼望某個法案,能夠在國會上通過；或者他們整個人生,是為要通過法案而努力.對他來說,這是全世界最重要的事,談論我本人,我根本不理.我不理會報章怎樣批評我,我不會反感；我所關心的是法案是否能通過.我們若用基督徒與政治家來比較,基督徒便好像沒有這樣容忍,願意接受各種批評,這因為太多基督徒,還沒有真正投身,委身於我們那偉大的使命上.
\newline
\newline
談到世界上各種的使命,你能否說出哪一個使命,偉大得過一個神兒女所投身的使命呢?我們生活在一個很繁複的世界,充滿困難,似乎任何國家,都沒有辦法解法這些困難.例如人口膨脹,疾病,戰爭,份子活動等等.這一切問題都無法解決,在世上每一處地方的人,都為著自己的罪而自疚.各人的心靈裡面是虛空的,各自尋找人需要的緣由.各人都盼望自己的罪得到赦免,找到人生生存的目標.可以有一天,能夠有勇敢面對死亡.在剛才所提到的問題,只有基督徒才有答案.我們的教會建築物外面有人等待我們,對他們的難題提出解答,但可是我們大部份的人,掛著自己的感受,而忘記了神給我們重大的使命.
\newline
\newline
有一次,耶穌對門徒說:「你們要像世上的鹽.」在家裡,可以放在兩個合宣的地方,可以放在廚房的樽裡,或湯裡.有時,我幻想兩粒細細的鹽,彼此對話:「在樽裡面,非常安全!」第一粒鹽對第二粒鹽說:「我們在這裡真好呀!這裡又好又乾；出面有陽光,有雨水,我們在樽裡都很安全,這裡又暖又好,彼此讚賞.各人都是這樣細粒,同樣顏色,我們真是在享受呀?」另一鹽答說:「是呀!我也喜歡這處,彼此交通,和同類共聚相處；我們是一樣的思想,我們在這地方一直享受下去多好；但我們有一天應該,要被倒進熱湯裡.」第一粒鹽立刻抗議說:「你有否想像過,被倒進熱湯裡,裡面的滋味如何?我們很快就給溶化了,整粒也失掉了,這對我們根本沒有甚麼益處?熱湯是這麼大,我鹽是這麼小的；我們留在樽裡,豈不是很快活的麼?我們何不留在原地,樽上有個圓蓋,很值得我們留在這裡.」另一粒鹽答說:「沒錯我們喜愛這處,但我們的使命是入到熱湯裡的呀!」
\newline
\newline
各位弟兄姊妹!今晚我們想到在自己的禮拜堂裡,和弟兄姊妹之間的交通,是很好的享受；一起唱詩,覺得靈裡有溫暖,當你聽牧師,傳道人講道時,你能享受神話語的教導；基督徒彼此交通,得到福氣,你們有好的事奉,教會是所好地方.但教會像是一個樽,禮拜堂外面,正有臣型的湯壼和熱湯,人們正在自己的問題裡打滾,需要尋求到答案.我聽神聲音向我呼召,時候到了!要從樽裡出來,要從廚房餐桌出來,要倒入世界的滾湯裡.你大半基督徒的人生,是否已花在教會的鹽樽裡,並沒有倒入滾湯裡呢?我可以聽到你心底的說話；史牧師呀!這個世界很大,香港也有很多人,在香港以外的世界有四十億人口；我們只是一小群人,就是召集全世界的基督徒,都是不足夠用?我們可發出甚麼作用呢?將我倒入湯裡,我又怎能起作用呢?我沒有能力領世界的人歸信主.但今天晚上,我要你們緊記,當你接受主耶穌為你個人的救主,你可能是全家第一位信主的,但當你以神兒女的身份回家裡的時候,你便要在家裡成為鹽.
\newline
\newline
在你未成為基督徒之前,你家裡是沒有鹽的.但因你信主的緣故,而有你在裡面,全家便與以前的味道有分別.如果你全家都是基督徒,當你遷入到一個新社區；雖然你還未領社區的人歸信主,也許最初你對一個人傳福音有點成功,或者有些鄰居對你的見證有些反應,但大部份鄰居卻未曾理會到.或者你以為基督教家庭,和鄰居之間起不到甚麼作用.但我告訴你,當你這基督教家庭來遷入這個社區之前,整個社區是沒有鹽的,那些人在世上沒有機會聽到福音.但當你搬入這個社區那一天,就是鹽入到社區裡,你便成為他們的第一個見證.
\newline
\newline
在我結束之前,我要向你發出幾個問題,你今天是否還在樽裡?或已倒進熱湯裡?愛是不作害羞的事,因愛是不求自己的益處,愛是不輕易發怒.生氣的醫治方法,就是忘記自己,投身為耶穌基督作見證,和遵行大使命.
\newline
\newline


\newpage

\section{第52屆~港九培靈會~史保羅~第四講　愛是不計算人的惡不喜歡不義只喜歡真理}
\label{sec:652}
\textbf{第四講　愛是不計算人的惡不喜歡不義只喜歡真理}
\newline
\newline
連結: \href{https://www.hkbibleconference.org/session-message/view/652}{\texttt{https://www.hkbibleconference.org/session-message/view/652}}
\newline
\newline
\hyperref[sec:651]{< < < PREV SERMON < < <}
~
\hyperref[sec:toc]{[返目錄]}
~
\hyperref[sec:654]{> > > NEXT SERMON > > >}
\newline
\newline
哥林多前書 13:5-13:6
\newline
\begin{longtable}{cl}
\hline
\hline
章節 & 經文 (和合本修訂版)\\
\hline
13:5 & \begin{tabularx}{0.7\textwidth}{X} 不做害羞的事，不求自己的益處，不輕易發怒，不計算人的惡， \end{tabularx} \\ \\ \relax
13:6 & \begin{tabularx}{0.7\textwidth}{X} 不喜歡不義，只喜歡真理； \end{tabularx} \\ \\
[1ex]
\hline
\hline
\end{longtable}
今晚請大家翻開(腓三17-22):「弟兄們!你們要一同效法我,也當留意著那些照我們榜樣行的人.因為有許多人行事,是基督十字架的仇敵,我屢次告訴你們,他們的結局,就是沉淪,他們的神就是自己的肚腹；他們以自己的羞辱為榮耀,專以地上的事為念.我們卻是天上的國民,並且等候救主,就是主耶穌基督,從天上降臨.他要按著那能叫萬有歸服自己的大能,將我們這卑賤的身體,改變形狀和他自己榮耀的身體相似.」這可以說是哥林多前書第十三章的註釋.
\newline
\newline
(一)不計算人的惡　在古舊的希伯來文和希臘文聖經,是沒有分章分節的,我相信神後來感動一些人做了這分章節的功夫；但我認為有些地方應該分多一節的,但他們沒有分.例如林前十三章四節的末句,不計算人的惡,應該是第六節「不計算人的惡.」意思是有了愛,我便不會對於人的錯誤,好像作運動比賽的一樣計分.保羅說:基督徒的生活,是不計算別人在我們身上所做成的傷害.
\newline
\newline
聽聞在北美洲有百分之五十的婚姻,是以離婚作收場,這因為佔大部份的夫婦,都計較他們伴侶的過失.我記得有一對已經結婚了二十年的夫婦,來找我說他們要離婚,他們所提出的理由,是因為一個小問題.其實這小問題不過是計分表上的一個最後小數；因此就成了他們婚姻破裂的致命傷,這好像繩折斷了駱駝頸一樣.大家都知道一條繫在駱駝頸上的繩,本身不能令駱駝的頸折斷,但如果在那繩子上負力,駱駝頸便會折斷.世界上任何一種關係就是這樣,假如我們只顧計算著別人的錯誤,這關係總有一天會破裂.世界上沒有一位太太或丈夫,能忍受對方計算他(她)的錯誤；這最終只會導致婚姻失敗.也沒有父母願意子女,經常數算著他們做父母的過失,這只會終止父母子女間之愛.亦沒有子女接受父母,經常記著他們的過錯,只會令到他們不能忍受,離開家庭.在友誼上如此,在教會裡也是這樣.牧師計算會友的過錯,會友數算牧師的過失,教會就會因此不和諧.踢足球,打網球,打水球等運動,我們都需要計分.但保羅說:如果神的愛澆灌在你們心裡的話,你們便不要對別人的錯誤,好像作運動比賽一樣的看.
\newline
\newline
(二)不喜歡不義　在林前十三章第六節:「愛是不喜歡不義,只喜歡真理.」這節經文可說是,最古舊的方法,來描述我們這世代的道德觀.作為基督徒的我們可以跟從兩種生活方式.第一種是接納聖經的絕對.譬如聖經說不可殺人,不可說謊,不可偷盜.我們便絕不殺人,說謊,或是偷盜.第二種是我們在社會上,可改變的法則.例如在加拿大有一半會眾,都會在崇拜後往餐廳喝茶,可是英國的基督徒,卻認為在崇拜後不該往公眾場所去.又在我的教會,只有很少姊妹,會在禮拜堂裡戴帽子；但有些國家卻認為每一次在禮拜堂裡都要戴帽子.記得我曾經到過一處地方的教會,那裡的弟兄都坐在右面；所有的姊妹都坐在左面；但在我加拿大教會,弟兄姊妹都是混合而坐.這都是聖經沒有說對,或不對的事.可是如果我們將這可變的法則,變成不可變的法則；成為神絕對的吩咐,便會有麻煩.或者我們將神絕對的標準,變成可變的法則,那就會有更大麻煩.保羅在這節經文說到,如果神的愛充滿我心,我便不喜歡不義.凡是神絕對的吩咐,而我不能遵守,我便會覺得很悲傷.我不會對謀殺,偷盜,不道德事情和行為感到快樂.但有另外一種人會對自己所做的錯事感到快樂.罪入到世界,便將我們的思想曲解過來,對於正確與錯誤,我們會感到迷糊；我們會對於自己所做的錯事,用諸般理由來解釋.聖經說不可說謊,不可偷盜,我們會加上去說,在某些情形下,是可以說謊和偷盜的.在英國小說裡面,羅賓漢是人所共知的英雄,幾乎每個人都稱他為劫富濟貧的好人.而事實上羅賓漢做的是錯事,他首先做了錯事,才使他有機會做好事.有些人為求達到目的而不擇手段.就譬如上週在義大利有一連串恐怖事情發生,恐佈份子自以為榮耀,他殺了一百人,傷害了三四百人.誰用任何標準來看他們所做的這些事都是錯的,不過他們卻以自己有重大的使命,腦海裡有一個好目標,要是殺害百人來達到這目的,又有何不可.這正如腓立比書第三章十九節說到,有等人以自己的羞辱為榮耀,他們已犯了罪,但卻炫耀自己的罪；他們不順服神,但卻誇耀自己的不順服；他們做錯了事,但卻炫耀自己的錯事.
\newline
\newline
(三)只喜歡真理　有一種人為著罪惡而流淚,以自己做錯事為羞辱.保羅在腓立比書第三章說,這種人要以保羅為榜樣.因為保羅已經認清楚到,在我們的日常生活裡面,並沒有甚麼創新事情.我們很希望自己能夠有創新的舉動,說話,偉論,意見和產品等等.然而最精美的產品,也只不過是許多人所做的,產品的最後一個成果而已.在舊約聖經有一個聰明人,就已經說了他環顧四週世界,日光之下並無新事.現今所做的事,只是模仿前人所做過的.一個嬰孩由初生到五歲這段時間,是從觀察四週而學會東西.所以加拿大孩子學會英文,中國小孩學會中文.每一個嬰孩在本能上都是會爬,但這些會爬的嬰孩,為何還要學站立行走呢?嬰孩頭重腳輕,站立行走對他們來說是何等麻煩,不過他們並沒有抱怨甚麼.每次當他們跌下,看見四週的人,都是用雙腳走路的,他們便從跌倒中,再試站起來用雙腳走路.我想像得到,如果嬰孩沒有看見人是用雙腳走路,他們斷不會用雙腳走路,這因為嬰兒,是在模仿四週人的行為而走路.
\newline
\newline
做基督徒也是這樣,我們總要學習某些人.也許有很多人說,我們只要仰望耶穌便行了.不錯,這是一個很真實意見,只不過我們在家裡,詩班席上,講壇和街上從不曾看見過耶穌.但不要忘記保羅說過,我們的身體是聖靈的殿,正因神的靈住在我們的身上；我們活在這世代,所以神便住在這世代.我們想見耶穌,便要看看我們四週生存著的聖靈的殿.
\newline
\newline
保羅在這節經文裡說,我們可以效法兩種人.第一種是不敬虔的人,他們以自己的羞辱為榮耀,跟隨著他們只會墮落.另一種是敬虔的人,是一個活得像耶穌而又給聖靈充滿的人,在這裡你必須作出抉擇.你或許是剛剛信主的,是神的兒女,在基督裡是個嬰孩.但你要學習一個基督徒應有的語言,思想和行事為人.如果你要找一個人來模仿和效法的話,保羅說你可以找他.今天究竟有多少基督徒,可以成為別人效法的對象?身為人母的,有沒有想到妳的女兒認為妳是全世界最好的人,最好的廚子,最聰明的女士,最美麗的人?這女兒自出世至懂事,都希望能夠像母親,但妳能否對你的女兒說,我雖然不是如妳所說是全世界之最,不過我已經盡力地為主而活；盡心竭力地走向神的方向,我的女兒,跟從我罷!或者妳對妳的女兒說,我犯了很多錯,一次又一次的失敗,我的一生是一團糟,你還是找別人來學習罷!記得我兒子四歲的時候,每天早上都看著我刮鬍鬚,我便給他一柄沒有刀片的剃刀和剃鬚膏.當我在剃鬚時,我發現他每一個動作都在跟著我,我仔細的思想,今天我生活的模式,人生的經歷,所走的方向,是否可以讓我的兒子,學習和跟從呢?你知否大部份的小男孩,都認為自己的父親是全世界最強壯,工作最好的人.你能否對你的兒子說,我並不是全世界最強壯,做事最好和最英俊的人；但我是為耶穌而活,給聖靈充滿,每一天都是朝向神的方向,你學我罷!也許你說我時常走錯路,犯了很多錯誤,時常違背神,兒子你找另一個人跟從罷!做牧師和傳道人的是要帶領群羊,作一個敬虔屬神的好模範,給聖靈充滿,活在崇高的境界,是一個基督徒應有生活的好榜樣.我不知你能否站在講台對會眾說,我不是全世界最好教育的人,也不是全世界好的講員或最聰明的牧者；但我求聖靈充滿我生命,盡我所能為主而活,我已竭力服事神；你們可以跟從我,因為我的人生路程是正確的.今晚坐著的弟兄姊妹當中,有些是作主日學教師的,是否有留意到每個主日學學生,都認為他們的主日學老師,是一個好基督徒的模範?雖然他們並不以為你是全世界最好的教師,有時候你所說的話他們不願意聽,但他們希望知道一個基督徒應該怎樣做的時候,他們會找你作榜樣.在下一個主日,你上主日學課時,有甚麼向那些學生說呢?你是否說我口裡所說的話你們可以依從,只是我的行為卻千萬不要以我作榜樣.或者你說我每一天求神的聖靈充滿我,我竭力去跟從主耶穌；按我所能而活出基督徒生活的樣式,你們如果跟從我便是走向正路.
\newline
\newline
保羅說,我們每一個人都會跟從某一個人,他可以是一個敬虔或不敬虔的人,我們亦有另外一些人來跟從我們.
\newline
\newline
我們要思想一下,在這世界上,有誰跟隨著我們,以我們作他生活的模式.這人或許是一個小孩,一位主日學生,或一批會眾.另外我們要想想,看我能否像保羅一樣的說,你們可以跟從我,以我作為你們的模範.
\newline
\newline


\newpage

\section{第52屆~港九培靈會~史保羅~第五講　凡事包容,相信,盼望,忍耐}
\label{sec:654}
\textbf{第五講　凡事包容,相信,盼望,忍耐}
\newline
\newline
連結: \href{https://www.hkbibleconference.org/session-message/view/654}{\texttt{https://www.hkbibleconference.org/session-message/view/654}}
\newline
\newline
\hyperref[sec:652]{< < < PREV SERMON < < <}
~
\hyperref[sec:toc]{[返目錄]}
~
\hyperref[sec:656]{> > > NEXT SERMON > > >}
\newline
\newline
哥林多前書 13:1-13:13
\newline
\begin{longtable}{cl}
\hline
\hline
章節 & 經文 (和合本修訂版)\\
\hline
13:1 & \begin{tabularx}{0.7\textwidth}{X} 我若能說人間的方言，甚至天使的語言，卻沒有愛，我就成為鳴的鑼、響的鈸一般。 \end{tabularx} \\ \\ \relax
13:2 & \begin{tabularx}{0.7\textwidth}{X} 我若有先知講道的能力，也明白各樣的奧祕，各樣的知識，而且有齊備的信心，使我能夠移山，卻沒有愛，我就算不了甚麼。 \end{tabularx} \\ \\ \relax
13:3 & \begin{tabularx}{0.7\textwidth}{X} 我若將所有的財產救濟窮人，又犧牲自己的身體讓人誇讚，卻沒有愛，仍然對我無益。 \end{tabularx} \\ \\ \relax
13:4 & \begin{tabularx}{0.7\textwidth}{X} 愛是恆久忍耐；又有恩慈；愛是不嫉妒；愛是不自誇，不張狂， \end{tabularx} \\ \\ \relax
13:5 & \begin{tabularx}{0.7\textwidth}{X} 不做害羞的事，不求自己的益處，不輕易發怒，不計算人的惡， \end{tabularx} \\ \\ \relax
13:6 & \begin{tabularx}{0.7\textwidth}{X} 不喜歡不義，只喜歡真理； \end{tabularx} \\ \\ \relax
13:7 & \begin{tabularx}{0.7\textwidth}{X} 凡事包容，凡事相信，凡事盼望，凡事忍耐。 \end{tabularx} \\ \\ \relax
13:8 & \begin{tabularx}{0.7\textwidth}{X} 愛是永不止息。先知講道之能終必歸於無有；說方言之能終必停止；知識也終必歸於無有。 \end{tabularx} \\ \\ \relax
13:9 & \begin{tabularx}{0.7\textwidth}{X} 我們現在所知道的有限，先知所講的也有限， \end{tabularx} \\ \\ \relax
13:10 & \begin{tabularx}{0.7\textwidth}{X} 等那完全的來到，這有限的必消逝。 \end{tabularx} \\ \\ \relax
13:11 & \begin{tabularx}{0.7\textwidth}{X} 我作孩子的時候，說話像孩子，心思像孩子，意念像孩子；既長大成人，就把孩子的事丟棄了。 \end{tabularx} \\ \\ \relax
13:12 & \begin{tabularx}{0.7\textwidth}{X} 我們現在是對著鏡子觀看，模糊不清；到那時，就要面對面了。我如今所認識的有限，到那時就全認識，如同主認識我一樣。 \end{tabularx} \\ \\ \relax
13:13 & \begin{tabularx}{0.7\textwidth}{X} 如今常存的有信，有望，有愛這三樣，其中最大的是愛。 \end{tabularx} \\ \\
[1ex]
\hline
\hline
\end{longtable}
在哥林多前書十三章第七節,保羅說到神的愛有四點.第一,愛是凡事包容.第二,愛是凡事相信.第三,愛是凡事盼望；第四,愛是凡事忍耐.
\newline
\newline
我希望各位從兩個角度來看這四點真理.首先我要說的；假如全世界祇剩下我一個人,這節聖經對我來說,是甚麼意思?各位可曾想過,在我們一生中最大的難題,就是怎樣與自己相處.當你到一個精神學或精神病專家那裡時會說:「我與太太,或者與丈夫,又或者與僱主很難相處.」其實這些事專家已經知道,問題是在你自己身上.我們常常對於自己總有點不歡喜的地方,例如種族,國籍,膚色,經濟狀況,甚至自己的樣子等等.就以我自己來說,我不喜歡自己的頭髮越來越少,我希望有很濃密的頭髮；但事實上我是一個頭髮稀少的人,如果我這樣也不能接受自己,又怎能與家人相處呢?然後請大家再從我與別人的關係,的角度來看這四點.
\newline
\newline
(一)愛是凡事包容　這裡的字眼翻譯出來的意思是有一個屋頂,包著它裡面的東西.屋頂的作用,是可以保溫,使室內空氣調節,防止雨水,保障安全.換句話說；屋頂是防止裡面的東西走出外,又包裹著那些在裡面的東西.
\newline
\newline
記得第一天晚上,我們思想哥林多前書第十三章的時候,保羅一開始就給我們下了「愛」的定義是恆久忍耐,而在恆久忍耐這辭上下文裡面,保羅指出愛是將一切都包含在內.愛是有個容量的,面對一切東西都好像一個屋頂,這樣假若我是屬神的人,我便有容量來包容一切臨到我身上的擔子,我不會把我的重擔放到別人身上.保羅的意思不是叫我們面對人生的大重擔時,也要獨力去支撐,因為有些重大的屬靈重擔,是你自己擔當不了的,那時你應該去牧師或傳道人那裡去告知他.或者有一天我會遇到一個醫學上的大重擔,我便應該去找一個醫生將我的重擔告訴他.若面對小毛病細小擔子,保羅說如果聖靈將神的愛澆灌在我身上,我便應該有容量,去包容這些每日的小擔子.
\newline
\newline
我第一次離家是往美國去升學,在差不多每一所美國大學裡面,都有一個男學生被稱為「模範美國男學生」或「男子漢」,我很快便知道原來在我間大學裡面,有一位青年人,正是這樣的美國男子漢,他差不多樣樣事情都可以做到,他是一個偉大的公開演說家,是一位優秀演員,全校最佳的運動員,最英俊的男孩子；此外他還身軀健碩結實得令我羨慕不已,他名叫大衛.他是我的朋友.有一天,我在校園裡碰見他,就問候他說:「大衛啊,今天好嗎?」這位個子高大英俊的青年竟花了半小時向我傾訴他一切難題傷痛的事.
\newline
\newline
也許你想這樣的人,還有甚麼不妥嗎?我自己在那時已經夠多煩惱了,那時還要擔起他那一份嗎?但若在另一方面想,我們會發覺我們是會將重擔,也分給那些我們最愛的人.每天早上,當我還是睡眼惺忪的時候,內子已經在廚房哼著福音詩歌,她會轉過身來對我微微笑著說:「今天你好嗎?」我立刻便將我的擔子放在她身上.有甚麼不好呢?當然沒有甚麼不好,不過她已有自己的煩惱了,她未必需要在這個時候來擔我的擔子啊!
\newline
\newline
各位,如果神的愛充滿我,我便能夠包容我的重擔,可以不讓我的重擔分加在別人身上；我學習包容,在我煩惱上面蓋上屋頂；但當我看見別人有煩惱的時候,也能夠將自己的屋頂放在他的煩惱上,幫助他去包容.那就是說誰有難題的時候,他可以絕對放心地來告訴我；因為無論他的煩惱難題是多麼尷尬或麻煩,他知道祇要他向我說的話,我便不將他的事情告知其他人.各位有沒有留意世界上,有多少人確能在你的煩惱上蓋上屋頂呢?有沒有遇見過一個人可以使你將你的大能處告訴他呢?就是在基督徒中亦很難找到,因為我們很少人學習怎樣去分擔別人的重擔.我們祇是聽見別人有甚麼毛病,便很想將這毛病分播出去,讓教會其他人也知道.
\newline
\newline
(二)愛是凡事相信　意思就是「愛」要對每件事情作一個最好的看法.這段聖經如何應用到我們身上呢?我無論在任何壞的環境,也思想到好的方面和作最好的想法.
\newline
\newline
譬如有人對我說:「史牧師,我現在病得凄慘啊!」「沒錯!你很慘,但是你尚未破產啊,這還好!」另一位走來對我說:「史牧師我快破產了!」這件事真的很慘呀!可是,雖然你快要破產,但你沒有病倒,這還好.又有另一位走來對我說:「我又病又破產,我這次是雙重慘慘!」「但這個世界上還有人愛你,關心你啊!在世界上你並不孤單,這件是好事啊!」跟著有人走來向我說:「史牧師啊!外面下著傾盆大雨啊!」「這件事可能很差,但我有傘子啊!」現在你明白沒有,保羅的意思是「愛」的注意力是在這把傘子,不是在那場豪雨上.
\newline
\newline
「愛」集中在人生的福氣上,不是著重在人生的重擔上.
\newline
\newline
在部份的詩歌集裡面,有以下一首的詩歌.我不知道這首詩,寫作的背景是怎樣,但我可以想像這首詩歌的起源.有一個人,他與一位經常數著自己重擔的女子結婚.每一天早上,這位做妻子的,總用重擔和早餐來服侍丈夫,這一件事經過了很多年.有一天早晨,這男人在他進食早餐的時候,在紙上寫上以下的說話:「你要數算主的恩典,主的恩典樣樣要記清楚,你便會立驚訝而歡呼!」原來這首詩是在吃早餐的時候,由愛你的丈夫獻給你這位妻子的.然我不清楚,這首詩歌寫作的情形是否這樣,但這首詩寫的事情正是要指出這樣.正如保羅開始時所說的,「愛」本身是省察事情發生的好一方面,並不是在壞方面.應用在對人方面,也應該是這樣.我對別人有個最好,最樂觀的想法；我要拒絕對別人有更壞的想法.
\newline
\newline
在文明的社會,大部份的法庭都有以下一個原則；被告人在未定罪以前,法庭看他是清白的,直到法庭證明那人是有罪為止.想信各位很重視和讚同這法律程序,但各位有沒有將法庭的審判程序,和我們對待朋友的方法來作比較呢?我們會發覺自己竟比法律程序更苛刻.只要我們聽到有關朋友或最愛的人的謠言,我們便有一個傾向去相信那謠言,立刻假定了他們是有罪,然後觀看他們能否證明他們是無罪的.
\newline
\newline
昨天晚上,我說到我有三個兒女,當他們十多歲的時候,有一次他們要外出,我在他們出門前規定他們在那時間前返家.他們出去很久,直至我所限定的時間到了還未回來.我便決定多等二十分鐘.二十分鐘過了,但尚未見他們,我再等多二十分鐘,到第二段的二十分鐘他們還沒有返來.我開始坐下準備我的講辭.最後門被輕輕地敲著,我便去開門,在他們未曾發出一句說話前,我便大罵他們一頓.
\newline
\newline
各位,這樣做是甚麼意思呢?我一開始已經假定他們有罪了!其實我們有沒有想到給機會他們解釋清楚,為何未能準時返家呢?我的兒女說他們忘記看鐘點,因為那裡沒有鐘可以看時間.我很難證實他們的說話,但我希望各位明白一點,假如我的兒女當時被押上法庭,法庭在開始時是當他們是清白的,直到法庭有足夠證據證明他們是有罪.可是我們對待家裡的兒女時,竟然連法庭所給予他們的權利也否認了.
\newline
\newline
保羅說,在我們與別人的相處關係上面,我們要以最好的想法來對待他們.在我們的教會中,我們要看每一個都是清白的,不要立刻論斷他們是有罪的.這樣的話,教會便有更多的人,與我們一同事奉主了.但可惜我們一聽到謠言便信了,以致那些弟兄姊妹要黯然地離開教會.
\newline
\newline
(三)愛是凡事盼望　意思是無論環境怎樣,仍存著盼望.我永遠不會自己放棄或說:「不可能了」.就算那件事情在我人生裡,已壞到無可再壞的地步,祇要我是神的兒女,我從不會投降和失望；祇要我是神的兒女,我從不想去自殺；我要仍存著愛的盼望和眼光來看我人生.在我的生活裡,我經常生凍瘡.以我累積了這麼多年生凍瘡的經歷,我相信自己是全世界凍瘡權威之一了.在凍瘡裡,我學會了兩點是還未在醫學雜誌上被引述過的.第一點,就是凍瘡是無藥可治的.曾經有很多會眾給我有關治凍瘡的方法,但依然是沒有效.有一次差不多連我的舌也燒傷.好多年前我在路易士安羅州舉行一個聚會後,有一個人拿一樽油走來對我說:「這樽油油可醫治你的凍瘡,但是你要在未生凍瘡之前預先塗上去.」這更使我相信凍瘡是無藥可醫的.有關生凍瘡的第二點是,我只要等十日,這些凍瘡必要死謝了.假若明天早上,我望望鏡子,又發現自己生了一粒凍瘡,我要挺起二十七又四分三吋的胸,我要定睛望著這粒凍瘡說:「凍瘡先生啊!我知道要醫你,你是無藥可治的,但我說給你聽,只要我在這十多天繼續生存下去,到那時你已經謝死,我仍活!」
\newline
\newline
各位兄弟姊妹!恐怕我們在屬靈生命裡面,都有不同的挫折,這可能是你的挫折,你完全沒有辦法解決；在你四週圍都好像有牆逼著你,使你動彈不得.有些弟兄姊妹也許在商業上遇到挫折,你在銀行的一切現款都用光了,不知道解決問題的方法.有些弟兄姊妹在道德生活上有頓挫,你繼續犯那些罪,你曾經參加很多的類此這樣的培靈研經會；也讀過許多種屬靈的書籍,但卻未曾在道德上所犯過錯上得勝.其他的弟兄姊妹,在你們的生命裡面可能有其他的頓挫,但你們要挺直你們的身子,定睛望著你心靈上的頓挫說:「凍瘡先生啊!我是神的兒女.我不知道這個難題的答案在那裡,也不知這煩惱會何時過去；但因我仍是神兒女,有一天總會有一種辦法,煩惱會被消滅,全能的神與我仍然要在這處呢!」在羅馬書第八章裡面,保羅豈不是要講論到這一點的真理嗎?他說給羅馬人聽:「萬事都互相效力,叫愛神的人得益處.」他正指明我們人生的一切凍瘡.他說神知道答案,結果,和怎樣去解決這個問題.我們只需學習把這些頓挫,放在全能的神的手裡面.又當我看別人的時候,我仍要存著盼望的眼光.但許多時我們對人卻有失望的傾向.我這個丈夫沒有希望了,我這位太太沒有盼望了,我的子女永遠不會成才的!我們把自己當作是一個細小的神,自以為我們有權來判定四周,誰有希望,誰沒有希望.這段聖經要我們向著每一個人,都帶著盼望的眼光.你可否記得主耶穌曾經講及農夫與麥田的事呢!一次田裡面有麥子和種子,有人看見便對農夫說:「假若往我們田那處,把所有的稗子都拔出來丟去,豈不是好嗎?在田當中所剩下的便全都是麥子了!」但耶穌說:「不要這樣,讓麥子稗子一同生長,直到收割的時候,我便會做分辨的工作.」各位可知道耶穌親自解釋這節經文的意思.祂說這塊田地就是世界,那些麥子就代表神的兒女們；那些稗子代表了魔鬼的兒女.但因為從來沒有一個人,可以每一次都準確地分別出誰是基督徒,誰不是基督徒,不論是牧師,傳道人,都會執事都有可能會有別的錯誤,把麥子當野草拔掉.惟有一絕個對準確的方法,是將來由神自己分別.令人興奮的是,神正在從事將稗子,也改變成為麥子呢?這就是整個福音的作用,福音乃是神將一個失喪的人拯救過來；乃是神的大能把盼望給予一個沒有盼望的人.神將人從大路上搶救出來放回窄路上.各位可以見到整個耶穌基督的福音,神的拯救和改變人的大能力.
\newline
\newline
或者今天我看見一個人,會說他沒有盼望了,但也許明天神已改變了這個人的生命.今天我看見一個人,我認為他是稗子,便把他拔出來逐出教會門外,但在下禮拜神要拯救這個人.耶穌基督說讓麥子與稗子一同生長吧!好叫到收割的時候,我主耶穌親自做分辨的工作,我將這些麥子放進我的會庫裡面；又將這些稗子拔出來拿去焚燒.
\newline
\newline
所以在今生的時候,我們對別人要常存盼望的眼光,因為這個原因,使徒保羅才能說愛是凡事忍耐,因為祇有對面有盼望的人,纔能忍受現在.
\newline
\newline


\newpage

\section{第52屆~港九培靈會~史保羅~第六講　無比的愛}
\label{sec:656}
\textbf{第六講　無比的愛}
\newline
\newline
連結: \href{https://www.hkbibleconference.org/session-message/view/656}{\texttt{https://www.hkbibleconference.org/session-message/view/656}}
\newline
\newline
\hyperref[sec:654]{< < < PREV SERMON < < <}
~
\hyperref[sec:toc]{[返目錄]}
~
\hyperref[sec:657]{> > > NEXT SERMON > > >}
\newline
\newline
哥林多前書 13:1-13:3
\newline
\begin{longtable}{cl}
\hline
\hline
章節 & 經文 (和合本修訂版)\\
\hline
13:1 & \begin{tabularx}{0.7\textwidth}{X} 我若能說人間的方言，甚至天使的語言，卻沒有愛，我就成為鳴的鑼、響的鈸一般。 \end{tabularx} \\ \\ \relax
13:2 & \begin{tabularx}{0.7\textwidth}{X} 我若有先知講道的能力，也明白各樣的奧祕，各樣的知識，而且有齊備的信心，使我能夠移山，卻沒有愛，我就算不了甚麼。 \end{tabularx} \\ \\ \relax
13:3 & \begin{tabularx}{0.7\textwidth}{X} 我若將所有的財產救濟窮人，又犧牲自己的身體讓人誇讚，卻沒有愛，仍然對我無益。 \end{tabularx} \\ \\
[1ex]
\hline
\hline
\end{longtable}
「我若能說萬人的方言,並天使的話語,卻沒有愛,我就成了鳴的鑼響的鈸一般.我若有先知講道之能,也明白各樣的奧秘,各樣的知識,而且有全備的信,叫我能夠移山,卻沒有愛,我就算不得甚麼.我若將所有的賙濟窮人,又捨己身叫人焚燒,卻沒有愛,仍然與我無益.」(林前十三1-3)
\newline
\newline
在這三節經文裡,保羅說到有些事情不及愛那樣有能力.他提到四件事情:
\newline
\newline
一,愛比講道的能力更大.
\newline
\newline
二,愛比行神蹟的力量更大.
\newline
\newline
三,愛比用金錢來賙濟窮人更有力.
\newline
\newline
四,最後,他說愛甚至比為一個偉大的使命而捨身更有力,更偉大!
\newline
\newline
這四件事情,每一件都是一股無可匹比的力量.
\newline
\newline
在我們身上的愛,其實比剛才所說的任何一件事更有能力.
\newline
\newline
首先,保羅用了四種方法來描述講道的能力.指出講道者善於運用舌頭的力量.而這不是可思議的力量,可以完成很多偉大的事情.
\newline
\newline
在第二次世界大戰的時候,有三位偉大的演說家.第一個是英國的邱吉爾爵士.當時英國的軍隊剛撤出鄧扣克,把所有的槍械和武器都留在歐洲的海岸上.他們已沒能力去防衛英倫三島和作戰.邱吉爾先生當然比其他人更明瞭這劣勢.但當他第一次以首相的身份向國會講話時,他手無寸鐵,沒有裝甲車或是足夠的飛機.他祇發揮他聲音的力量向國會說,我們要在海灘街頭和家中作戰.他說服國家的人民說:「英國還還有足夠的力量來抵擋希特拉.」他甚至向希特拉說明,英國一定會成功,這一切的成就,都是憑著人類的聲音.在邱吉爾先生的地方隔了一個海峽,有另外一個人也用聲音的力量.他就是希特拉.當時德國與任何國家來比較,是在最低潮和不景氣的時候,數以萬計的德國人,沒有工作可做,德國的錢幣下降到不值錢.這時希特拉開始他的演說,他說服德國人民；他可以使德國在不景氣當中復甦過來,甚至可以征服歐洲；成為世界上最大的民族,這一切都是憑藉著他的聲音而達到.德國的將軍,有很多都有作戰經驗和打仗的方法,但都因為希特拉的聲音,大有能力的原因,以致他們都要遵從.
\newline
\newline
在這同時,在美國有另一位演說家,他就是美國歷史上,任期最長的羅斯福總統.當時是在一九三○年代初期,美國經歷到一連串的不景氣,羅斯福總統時常透過電台,向全國人民廣播,有些人稱他的演說,為火爐邊的談話；因他沒有邱吉爾先生的巧言妙語般的口才,也沒有希特拉大聲般的聲調疾呼.他只用那些去家裡與人談話的聲調和語氣,去說服人,安慰人.說事情會一步步前進,他的民族要繼續抬起頭來,美國人便相信他.羅斯福總統藉著人的聲音,而獲致一部份的成就.
\newline
\newline
使徒保羅在經文所說的先知講道能力,便是這種能力,他說每一個男人,女人,在他們用聲音的能力上,都可以有很大的能力表現.但保羅接著又說；假如你是一個有好口才的人,同時說話亦好像天使一樣,相信每一個人都會問天使怎樣說話呢?這裡所說天使言語,與說方言是兩件不同的事.我找不到任何理由,來支持天使言語,就等於說方言.聖經給我們知道,天使的話語,通常是採用說話對象,所採用的語言.例如天使向以利亞顯現,及說馬利亞所採用的言語.當天使和施洗約翰的父親說話的時候,也用他能說的言語.在聖經裡沒有記載,天使採用陌生的言語,天使所採用的言語,都是所說的對象能明白的,意思是天使都能用言語,及聲音向人傳達.當然大家都知道,天使經常都是神的使者,代表神說話,我們能說天使的言語,就是我們有能力,將神的訊息傳給人知道.保羅又說到第三點,如果這人有講說話的能力,正如其他人都有講說話的能力一樣.他同時還有多一點其它的恩賜,就是他能將神的信息傳揚出去.保羅說他不但可將神的訊息傳達,他還能作先知講道,意思是他講話傳道的時候,能帶著神先知的聲音去講道.接著保羅舉例說:一個人有先知講道之能以外,也可以明白各樣的奧秘.奧秘在新約聖經通常是指那些真理,在舊約的時候是被隱藏起來的.但在新約時代,就向使徒顯現出來.譬如保羅提及耶穌基督復活及再來的奧秘.但在舊約聖經中,那些人對於耶穌基督的復活及升天都不明白,這些真理在當時都是被隱藏起來.保羅在提及耶穌基督與教會之間的關係,稱為奧秘.在舊約聖經中,神與人有關係,祇限於猶太人而已.他們以為神除了猶太人以外,對其他的人,神都不關心,但事實上,神竟然對所有的人都關心,這件真理在舊約聖經中都被隱藏,他們不曉得,但保羅都明白這個奧秘.道成肉身這個真理,亦是神向使徒保羅所要顯明的,但在舊約時代,當時的人,根本無辦法明白神怎可能成為人；他們很困難地明白到「彌賽亞」,一方面是神,但另一方面同時又是人呢!保羅說這點真理在舊約是隱藏起來的.現在他在描寫一個有先知講道能力的人.他能夠像天使般向人傳揚,同時又帶著先知講道的權柄；而且他又明白到神話語的奧秘,最後他又有著全備的知識.知識在新約聖經,常是指有關屬靈事情的知識；在哥林多前書第二章中,保羅用了很多篇幅來提及屬靈的知識.他說一個普通屬血氣的人,沒有辦法明白屬靈的事情,因為他們會認為屬靈的事情是愚拙的.保羅說一個普通的人,沒有辦法知道屬靈的事情.當保羅說到知識恩賜的時候,他指有關於認識屬靈的事.在加拿大,有些人認為知識的恩賜,一定是有些像幻術性的恩賜,使一個人面對一大群人能瞬即知曉誰患了病,說出坐著的某處的人是心臟衰弱,在樓臺上某某人患有癌症,但保羅所寫的書信裡,知識的恩賜永遠不是如以上所述說的.知識的恩賜,永遠是說明白有關於屬靈的真理.所以一個人能夠好像普通人一般講說話,同一時候,他有著天所有的權柄；他也有著先知講道的權能,這個人竟然完全明白神的話語的奧秘.最後,他對於屬靈的事情,都有完全的知識.這就是保羅在此,正描述一個傳揚福音的人是怎樣的.保羅說到一個人如果有一切的才幹,但假如他的推動力,他的全心,不是出於愛的推動；他所做的一切都沒用處的.他的講道好像是一個走了音的樂器.我們可以成為一個偉大的佈道家,但可能與神失去了聯繫；惟有那些與神有生命接觸的人,才可以真知道神的大愛.保羅說一位弟兄或一位姊妹,可以有一切的恩賜,但仍然可以與神失去了團契及交通.愛比會講道傳福音的能力更偉大.保羅講過一件令我感嘆的事實,竟然比可以行神蹟的更有力量.正如在(林前十三23)說就算我們有了全備的信心能夠移山,這裡所說的信,不是指到我們信耶穌而得救,那種初步信心的信,信耶穌得救這種初步的信心,是世界上任何人都可以得到的.世界上每一個人都可以靠著信靠耶穌基督為救主,而有初步的信心.當我們相信耶穌為救主的時候,我們便成為神的兒女；這種信,我們稱為因信而得救的信心.但在這經文中所說的信,是能行神蹟的信,有心能夠醫治人脫離某種病；或者是有能力,能在水面上行走；或是有能力將水變為酒,亦能夠有耶穌基督般的能力去行神蹟.保羅說如果你亦擁有這種能力,但在內裡卻沒有神的愛來推動,這種情形是可以發生的.我們也要十分小心,在現今世代裡有很多人,自稱有能力能行神蹟；每一個能行神蹟的人,未必真是憑著神的能力而行,有些人甚至所謂行神蹟,是靠著魔鬼的力量.相信大家都知道有一天,在這世界裡會出現敵基督.這敵基督要統治這個世界,這個與敵基督一同工作的人,聖經稱之為「假先知」,假先知能做的工作,包括行神蹟,能做耶穌基督所能做的事,但他所得的能力,乃是從魔鬼而來.能夠行神蹟的人,並不充份表示他是一個屬神的人；他既然可以成為魔鬼的使者,而又有權柄去做所謂神蹟；保羅說這些事若不是出於愛心的話,便是沒有價值.
\newline
\newline
接著,保羅說到愛比賙濟窮人更有能力,所謂賙濟,包括站在人道上去幫助他人.拯救在餓死邊緣的人；照顧在醫藥上有需要的人和沒有地方居住的人等等.這些都是好事.但保羅說,這種捐贈包括對教會各項聚會,和一切宗教團體的奉獻,必須是被神的愛所推動,否則這種捐輸,便是毫無價值.作為一個基督徒,你曾否做過不正當動機的奉獻嗎?在民眾教會,我們收獻金是把獻金袋,在會眾中傳遞的.有時當獻金袋傳遞到某些人的時候,他們看見周圍的人都有奉獻,恐怕自己不奉獻會引起尷尬.為著不要失面子的緣故,他們總要放些少獻金在獻金袋裡.這種奉獻的態度是錯誤的.但在某些時候,我們都會嘗試過這種的錯過.很多年前,我在南非約翰尼斯堡舉行一連串佈道大會,這是神賜給我的福份；在每天晚上都有成千上萬的人參加聚會,每日中午聚會亦有好幾百人.有一天中午,在聚會完畢後,我返回酒店的時候,看見路旁坐著一個衣服破舊,拿著小罐子的乞丐；這使我想起加拿大,有些乞丐是以此為職業.他們所賺的錢比做生意的還多；所以我們小心,免得我們拿錢給這類乞丐.我心想或許這個乞丐亦是職業性,他會將所得的金錢拿去買酒.根據我過往的經驗及習慣,我最後是不會拿錢給他.但當我差不多走入酒店的時候；但我會想起,在這個地方有幾百人聽我講道,而在這星期中,我都在講述神的愛；假如有些聽眾看見我經過,連一個錢不施捨給乞丐的話,他們可能會認為我講是一套,做出來的卻是另一套.為了避免有人看見我不施捨而覺得尷尬；我便施捨了給他.但當我回到房子的時候,我開始思想我剛才的行為.我對自己說:「史保羅,你的施捨態度是錯誤的,因你很懼怕在人面前失去了面子,你不是因為關心他,愛他,憐憫他而施捨給他.」神說這種捐輸對奉獻者,及受惠者,是毫無價值的!你施捨的錢,不但不能對需要者有多大的幫助；而且在神的眼中,也是毫無益處.有些人認為他奉獻給神,神就會以大數目來回報給他,對這一類人而言,他是把神看為銀行.我若在星期日奉獻了十元,在本月結束的時候,神會回報給我一百元.當然神確實應許,當我們奉獻的寺候,神就會賜福給我們,但祂卻從來未曾說明這種賜福是甚麼福氣.這些福氣也許是在金錢上,家人上,教會裡,健康方面,壽命上或其他許多許多的方式；總言之,祂會把最好的東西來賜給你.可是如果我奉獻給神是存著一個目的,好叫神在金錢上賜福給我,我便陷於一個錯誤的奉獻目的.最後,保羅說愛是比為一個偉大的使命,而捨生命更有能力.
\newline
\newline
換句話說:即使捨己身叫火焚燒,但除非這種行為,是出於神的愛所推動,否則這種殉道,仍對我們毫無益處.在加拿大,有人認為如果一個軍兵為國家戰死在戰場上,便可以得救和立即到達天堂.這當然全沒有聖經根據.在哥林多前十三章一至四節,保羅將這些誤解糾正.他說:除非是由於神的所推動,否則這種殉道對我們是毫無益處.如果我們講道,行神蹟,捐獻金錢,捨身等要帶有能力,便必須要出於神的愛所推動.但是我們要從哪裡才可以得到這種愛呢?在這個世界裡,哪一類人才可以有這種愛呢?在羅馬書第五章一至五節,保羅在這裡說,我們得到這種愛,是因相信耶穌得稱為義人,每一個人如果接受了耶穌基督,作為我們個人的救主；在我們的心裡有神所賜的平安,聖靈便將上述神的愛,澆灌在我們的心裡.假如我們沒有與這位永遠活著的神接觸,不接受主耶穌作我的救主；我們便沒有辦法得到這種愛.不管我們做了許多好事,甚至又得到了別人聘請為傳道人；或可以行神蹟,我顧意為一個偉大的使命而捨身殉道.若果我沒有與神的愛接觸,我們所做的一切仍與我無益.今天晚上我們應問問自己,是否那些與神有接觸的人?已接受了耶穌基督,作為我個人的救主,在我心靈裡得到神所賜給我的平安?是在神家庭裡的一份子?我的罪已得到主的赦免?我已確知自己是走向天堂的道路?除非我們對以上的問題有了正確的答案,否則我們便無從開始為神而活,所做的一切好事都看為不重要.除非我們有了正確的答案,我們所做的一切,才會有價值及寶貴.這樣,我為神所做的工,都變為重要.
\newline
\newline


\newpage

\section{第52屆~港九培靈會~史保羅~第七講　愛的增長}
\label{sec:657}
\textbf{第七講　愛的增長}
\newline
\newline
連結: \href{https://www.hkbibleconference.org/session-message/view/657}{\texttt{https://www.hkbibleconference.org/session-message/view/657}}
\newline
\newline
\hyperref[sec:656]{< < < PREV SERMON < < <}
~
\hyperref[sec:toc]{[返目錄]}
~
\hyperref[sec:659]{> > > NEXT SERMON > > >}
\newline
\newline
哥林多前書 13:8-13:13
\newline
\begin{longtable}{cl}
\hline
\hline
章節 & 經文 (和合本修訂版)\\
\hline
13:8 & \begin{tabularx}{0.7\textwidth}{X} 愛是永不止息。先知講道之能終必歸於無有；說方言之能終必停止；知識也終必歸於無有。 \end{tabularx} \\ \\ \relax
13:9 & \begin{tabularx}{0.7\textwidth}{X} 我們現在所知道的有限，先知所講的也有限， \end{tabularx} \\ \\ \relax
13:10 & \begin{tabularx}{0.7\textwidth}{X} 等那完全的來到，這有限的必消逝。 \end{tabularx} \\ \\ \relax
13:11 & \begin{tabularx}{0.7\textwidth}{X} 我作孩子的時候，說話像孩子，心思像孩子，意念像孩子；既長大成人，就把孩子的事丟棄了。 \end{tabularx} \\ \\ \relax
13:12 & \begin{tabularx}{0.7\textwidth}{X} 我們現在是對著鏡子觀看，模糊不清；到那時，就要面對面了。我如今所認識的有限，到那時就全認識，如同主認識我一樣。 \end{tabularx} \\ \\ \relax
13:13 & \begin{tabularx}{0.7\textwidth}{X} 如今常存的有信，有望，有愛這三樣，其中最大的是愛。 \end{tabularx} \\ \\
[1ex]
\hline
\hline
\end{longtable}
「愛是永不止息,先知講道之能,終必歸於無有,說方言之能,終必停止,知識也終必歸於無有.我們現在所知道的有限,先知所講的也有限,等那完全的來到,這有限的必歸於無有了.我作孩子的時候,話語像孩子,心思像孩子,意念像孩子；既成了人,就把孩子的事丟棄了.我們如今彷彿對著鏡子觀看,模糊不清.到那時就要面對面了.我如今知道的有限,到那時就全知道,如同主知道我一樣.如今常存的有信,有望,有愛,這三樣,其中最大的是愛.」(林前十三8-13)我們昨晚說到愛比講道,行神蹟,為一個使命而捨己身,一件重要事而捨生命等事,更有力量和更偉大.此外保羅從第四節至第七節說明,當我們心裡被愛充滿的時候,應該怎樣與人相處,最後他說到愛究竟可以留存多久.
\newline
\newline
有人對我說:當戲劇被發明的初期,觀眾如果對某些戲子不歡迎的時候,便會在台下喧嘩,直至到那戲子離場為止.當保羅說到愛是永不止息的時候,也便引用這個比喻說到在我們人生的舞台上,彷如一位好戲子,永遠不會被觀眾討厭,從台上被趕走.但是先知講道,知識,講方言等事都是有限的；在最後終結,那「完全」來到的時候,這些「不完全」的事都會成為過去.
\newline
\newline
就單獨以先知講道之能來說,我們發現舊約裡面的先知,總會發生一些錯誤.摩西的錯是曾經不信服神；以利亞的錯是滿以為只有自己,是全世界唯一信神的人；所以我們將來在天上的時候,不再需要先知.
\newline
\newline
知識不論我們學會了多少,我們仍有許多許多的不曉得.每一個偉大的科學家都會對你說,對於知識這一方面我們只知道皮毛,因為知識是永遠學不完的.
\newline
\newline
在十二和十四章保羅提及講方言這事,他說講方言是一件好事,不過有一個問題是很容易引起混亂.有些時候基督徒也不明白,基督徒所講的方言,而當一些未信的人,聽到一些信徒在講方言的時候,他們會以為這群信徒在發瘋.這些偉大的恩賜,都會因為某些原因而消失,而保羅所要告訴我們的,是基督徒必須要長進.
\newline
\newline
第一方面我們在知識上要增長.我們做基督徒越久,便要對屬靈的事認識更多.
\newline
\newline
第二方面我們要在品格在行動上有增長.保羅在此用一個小孩,從生長到成人的過程來作比喻.他說當我們是一個小孩的時候,在行事上便好像是小孩.一個小孩發怒時,他會不理會甚麼便胡亂打人；當他難過的時候,他不會理會有人正在睡覺,便放聲大哭.保羅指出一個基督徒不能永遠在屬靈上作小孩.他(她)應該長成為一個男人或女人.若果是信了主十年的,便應該在行事為人上,比十年前更加成熟,倘若我們的品格,在每一年都沒有長進的話,這很明顯是我們裡面,有些錯誤和問題了!
\newline
\newline
第三方面保羅提到我們在眼光方面要有增長.他用鏡子作了一個比喻.因為在保羅時代,鏡子並不能像我們如今一樣,反映出這樣完美清楚的影像來,他們可以看到自己,但不會很清楚,好像是在薄霧裡看東西一樣.保羅說我的眼光在這世界上,受到一些損害和限制.可是當我成為神的兒女之後,每一天我的眼光都應該在神面前有增長,我對所見的事物,以前所認識的人,教會裡的工作和神；都應該越來越清楚,就正如我的身體,需要成長一樣.我們都應該問自己到底要做些甚麼,才可以使我們成長呢?
\newline
\newline
大家都知道當我們悔改信主的時候,便是重生.不論我們的年齡大小,我們在開始信主的時候,從屬靈的角度來看是一個嬰孩.要使一個嬰孩成長,在基本上是有四點.
\newline
\newline
(一)嬰孩必須有食物　今天世界有很多地方的難民嬰孩正在死亡,因為他們沒有食物.去年我探訪泰國的難民營,他們已沒有國家,嬰孩只剩下皮膚包著骨頭,這因為沒有食物,他們都會死亡.「像纔生的嬰孩愛慕奶一樣,叫你們因此漸長.」彼得前書第二章第二節,彼得說到嬰孩是必須要得到食物.我想起當我女兒還是嬰孩的時候,我們每隔三,四小時便要餵她飲奶,就算是半夜三更,我和內子總要互相輪晚起來餵她.有時,當我讓她飲完奶後,我上床還未立刻睡著,心想我們每一晚都要起來餵這嬰孩,確是麻煩,要是我有辦法要她一直飲,直至到她飲完了整個星期的奶,然後把她放在牆角讓她睡足一星期便好了.大家都會笑,因這是沒有可能的事.這個嬰孩只有一個很細小的胃,每一天只能讓她飲很小份量的奶,這小小的生命若要繼續生存；這小小的胃便要常常的被餵飽.彼得前書正是說明這一樣餵養的方法.當你重生的時候你是一個屬靈的嬰孩,但是一個嬰孩需要長大,方法便是在神的話語上得到餵養,所以如果不讀聖經,便完全沒有希望能夠成長.有很多人知道這一回事,但仍然不肯像餵嬰孩一樣,來讀神的話語.有些人只是在禮拜日讀聖經,他以為在禮拜日飽讀,便能飽到下禮拜日.換言之,他們相信餵一個肉身嬰孩,雖然不能一禮拜餵一次,但在屬靈上卻可以.其實一天讀經的時間,決不能足夠七天用.我們若要按照神的心意來成長的話,就必須要每天讀經,我是一個屬靈的嬰孩,只有一個很細小的胃,我必須時常讓這個容量充滿,我知道惟一的方法,是我每一天選定了一個固定時間來讀聖經.記得當主耶穌說到祈禱的時候,他這樣說:「當你進到你的密室裡祈禱的時候.......」你記不記得祂說「當」.主耶穌的意思是,祂知道我們每個人,都有一段時間需要祈禱.未知在座當中是否每一位,都為自己定下一段時間祈禱呢?你們是否知道自己明天,在哪一段時間祈禱呢?或許有些人說:我每天都有讀經,但我是隨著自己在每一天喜歡哪一段時間讀經才讀經.如果你是這樣的話,我保證你不會是每天都在讀經.
\newline
\newline
當我們在讀經的時候,如果遇到一些疑難的章節,我們應該怎樣辦呢?我們是不是說:「牧師呀!我讀不明白,所以我何需讀經呢?」其實讀聖經應該好像吃魚一樣.魚有魚骨,當你吃魚發現魚骨時,你會否說:「啊!沒希望了!我只能消化魚肉,不能消化魚骨,我以後不吃魚了!」這樣你便把魚拋掉.沒有人這樣吃魚的,當你吃魚發現魚骨的時候,你會將那條魚骨放到另外一邊,然後繼續吃魚.再發現第二條魚骨的時候,也是這樣.在讀經時,當你讀到有不明白的地方,你要說:「呀!這是魚骨,我要將它拿到一邊,繼續讀下去.」有一點是只有讀經會發生,而吃魚不會發生的；當你重複又重複閱讀一段經文時,你會發現裡面的魚骨越來越少.以前明明是魚骨,現今突然變為魚肉.你知不知道當你再次重讀的時候；你已經長成了一些,你的胃口已變得強壯一點,可以消化以前不能消化的東西.你讀經的時候,你的靈命便在成長中；你越加成長時,你便更加明白有關聖經的事.
\newline
\newline
(二)嬰孩必須有運動　當一個嬰孩成長的時候,他一定要有運動.想想看把一個嬰孩放進了一個沒有半點可供移動的箱裡會有甚麼後果.他一會成為一個毫無用處的人,這不是因他沒有食物,只是他沒有機會走動而已.有些基督徒就是這樣永遠不成長,因為他沒有接受一些基督徒應有的運動.保羅在腓立比書第三章說:「我不是以為自己已經得著了,我只有一件事,就是忘記背後,努力面前的,向著標竿直跑.」他說他現今正在做屬靈運動,並不是單單坐著領受禮的說話；而是要在這個賽跑場上,吃從天上所得的食物,然後纔奔跑.這是否已參與這一個屬靈的運動呢?在座中大多數的弟兄姊妹,都屬於某一間教會,你在教會裡有沒有教主日學,參加詩班,出外探訪或其他事奉呢?你有沒有向牧師說你希望可以服事神呢?你知否在今日的教會裡面,只有百分之十的基督徒是活躍的.其餘百分之九十,相信是在每主日,都只是坐著吃大餐.他們領受了很多力量,卻沒有使用出來.你們知不知道吃得太多,而不運動是會引起消化不良；會覺得不舒服,開始與人產生磨擦；你是一個難相處的人,你必須要有這些運動.
\newline
\newline
(三)嬰孩必須有交通和溝通　才可以成長為一個正常人,在監獄,最嚴厲的刑罰是將一個人隔離,獨自關在監牢裡面；每次都從門上的洞口供給他食物,但不准他看到其他人或與別人談話.在這情形下,他就無法成為一個正常的人.可惜有些基督徒卻不知道要與其他基督徒交通的重要,以為自己可以單獨成長,不要要參與教會事奉；他們接受了主耶穌之後,便好像成了一個被隔離的囚犯一樣.怎會成長呢?保羅在哥林多前書第十二章,說到基督徒好像一個屬靈的身體；身體上有不同的肢體,有手,有腳,有眼睛各式各樣；這些肢體是互相倚賴.手不能沒有腳,腳不能沒有眼；眼不能沒有耳朵,肢體要彼此相顧,纔可以儘量發揮潛能.每個基督徒都需要在同一個身體內,與其他肢體,其他肢體也需要你.神將一些特別的東西賜給你,好叫你對這身體有所貢獻.你是否每次都有機會參與教會的聚會?或者當教會有聚會時,你說:「我應該要去!」若我要令到其他的基督徒蒙福的話,我便要去到他們當中,我若要增長,我便要與兄弟姊妹有交通.
\newline
\newline
(四)嬰孩必有良好的環境　試想若把一個嬰孩放入雪櫃裡面會有甚麼後果?大家都會微笑,但我卻肯定地說,在我們當中,有些人會在下週往雪櫃裡參加聚會,因為你會往一所不明白福音:不相信聖經的地方聚會.又試想想將一個嬰孩,放入焗爐裡又會如何?這嬰孩必死無疑.這樣我們知道要在正常的環境下,我們屬靈上才可以成熟.
\newline
\newline
很多時在我們的生活裡,會與很多不信的人接觸.譬如和不信的人工作,上課,同住等等.這些事情都很重要,因為你要向這些人作見證.只是有一點更重要的,是你要有一段時間和神的子民在一起.保羅說:「你要從他們中間出來,與他們分別.」我們在行動上需要與世俗的人有分別,我們生命的目標是要效法耶穌基督的樣式.這也就是保羅在(羅八29)所說的.而這一節經文的上一節,便是:「我們曉得萬事都互相效力,叫愛神的人得益處.」這也就是說無論哪一件事發生在我們身上,都是好的,是為了榮耀神,雖然某些事情發生在我的身上,是不如意的；但這是神准許發生的,祂會使事情變為如意；使我們越來越像耶穌基督.所以今晚你坐在這裡聚會,明天有困難臨到你身上,或是家裡有人患上疾病等等,都是要幫助你效法主耶穌的模樣.而這過程是從我們重生開始,直到一生的終結.
\newline
\newline
我們做基督徒有三個階段.第一個階段是好像嬰孩一樣生下來.第二個階段是在生命裡不斷的成長.第三個階段是與主面對面.在這時候主耶穌要介紹你和我,給他的父神認識.那時我們會好像耶穌基督一樣,沒有任何瑕疵皺紋和污點.總會有一天,每一個基督徒都要像耶穌；這點當你和不信的人在一起,或在散會裡遇到你討厭的人時要記著.我未曾認識一個人是沒有缺點的,沒有屬靈皺紋和沒有錯處的.我只看自己比任何人多缺點,多毛病.但我不可灰心,因保羅說:你知道神仍然在我身上工作.使我每一天越加像耶穌.於是我便回顧過往,為著有進步而感謝神,沒有進步的便要尋求解決,這可能是與讀經,屬靈運動,與肢體交通和屬靈環境有關係.
\newline
\newline


\newpage

\section{第52屆~港九培靈會~史保羅~第八講　當代神的聲音}
\label{sec:659}
\textbf{第八講　當代神的聲音}
\newline
\newline
連結: \href{https://www.hkbibleconference.org/session-message/view/659}{\texttt{https://www.hkbibleconference.org/session-message/view/659}}
\newline
\newline
\hyperref[sec:657]{< < < PREV SERMON < < <}
~
\hyperref[sec:toc]{[返目錄]}
~
\hyperref[sec:661]{> > > NEXT SERMON > > >}
\newline
\newline
當我年青的時候,我覺得我有些不妥,因為我從來沒有聽到神的聲音.偶然聽見別人在講台上,分享一些話說:「神呼召我,又差遣我去；神吩咐我作這事又作那事.」我認為我有些不妥,因為神從來沒有在我身上做甚麼；我周圍的人都好像聽到神的聲音,而我就偏偏聽不到.自從某次,我學會了聽神的聲音後,我就知道原來是有五個很基本的方法,能夠聽到神的聲音.當然我相信除了這五樣之外,還有其他的方法；而且可能你聽見神聲音的方式,與我所聽見的方式不同.
\newline
\newline
(一)神的聲音是人可以聽見的　聖經記載亞當,夏娃在伊甸園的故事.那時,天起了涼風,他們在園中行走的時候,就聽見耶和華的聲音.有些人就說:哦!他們只是在做夢吧,他們是否在做夢呢?不,這不是聖經所說的.聖經很清楚的說,他們聽見的是一個聲音.有些人又說:他們可能是見到異象吧!但聖經明說他們聽到的是聲音.另外我們看見聖經記載,關於以利亞,這個偉大先知的事情；以利亞在當時,是最偉大的宗教領袖,以現在來說,以利亞就好像現代的葛培理一樣.人人都曉得以利亞這個名字,以利亞是一個祈禱的偉人；你記不記得以利亞祈禱,求天上不下雨,就結果三年半沒有降下來.這是一種非常有力的禱告,後來以利亞又祈禱,傾盆的大雨就降下來,我們又記得以利亞在迦密山上,和巴力的先知來比較真假,看看真神還是巴力會從天上降下火來,將他們的祭物焚燒.你記得那些巴力先知嘗試求火降下來,他們雖然嘗試,但卻失敗了!因他們的神是假的,所以不能從天上降下火來.然後以利亞便向耶和華禱告,聖經就說:火從天上降下來.跟著發生了一件很奇怪的事,在這偉大的屬靈勝利之後,以利亞的情緒,卻變得非常之低沉!以利亞雖然是一個偉大的先知,他和我們卻是一樣性情的人.你有沒有發現到一件事,可能在你覺得是最勝利的時候,也就是你最失敗的時候；在在你的人生裡你以為到達高峰的時候,可能突然跟著來是個低潮.每一個做牧師的,都有這種經驗.他們都喜歡禮拜日,而且非常享受這一日,因為禮拜日他們有機會向人宣講神的話,而且可以感受到會友溫暖回應.因為他在會友當中,有這個屬靈的交通.但是禮拜一卻不同了.他可能告訴會友,禮拜一永遠不要打電話來,你若要死便找禮拜二來死好了!以利亞就好像這樣,他低潮到一個地步,結果就去到一個山洞裡躲藏.以利亞還有第二個理由他要逃跑.因為有一個女人追趕他,這個女人就是王后耶洗別要出來殺他的頭；所以他一有機會就要逃走.聖經說以利亞跑到一個山洞裡躲藏,他很希望也很焦急能在這時候聽到神的聲音,希望得到神的指示,使他知道下一步應該怎樣做；他很希望知道神要他往那裡去.他心裡很焦急,但是當他坐在山洞裡的時候,覺得到的只是悽慘孤單,他為神大發熱心說:「以色列人背棄了你的約,毀壞了你的壇,用力殺了你的先知尋索我的命.」你曾否想過,你在這世上是唯一剩下來的基督徒呢?或者說；整個城市只剩下你一個人傳福音.事實上神永不會因沒有你便能行事的,有些青年人,他可能在一個大工廠或大公司裡任職的.他說:「牧師哪!我在這大工廠裡,只剩下我一個人是愛主的.」其實還有其他的人也同樣的愛主,不過你還沒有發然而已.你所發出的光,已微弱得到看不見的地步了,你必須將那些火挑旺起來,別人就能見到光了!神對以利亞說我在以色列人中,為自己留下七千人是未曾向巴力屈膝的.以利亞需要聽到神的聲音,他聽見在山洞外有一陣大風吹過的聲音,但耶和華卻不在風中,風後地震,耶和華卻不在其中,地震後有火,耶和華也不在火中,火後有一微小聲音發出來,這纔是神的聲音.這聲音直到啟示錄仍有發生,約翰也就見過.聖經記載,從創世記至啟示錄,有很多人曾聽到神的聲音,今晚在我們當中也許有些人也曾聽過神的聲音,知道神要他做些甚麼事.不過也有許多人說,我從未曾聽過神聲音的.我們要彼此尊重,不可批評別人因為這是你與神之間的事；反過來說如果我說沒有聽到神的聲音,你決不可把我看作是低一等的基督徒呀!
\newline
\newline
(二)神的聲音是可以透過徵兆就是的　在士師時代,因國中沒有王,各人任意而行,他的並沒有法律的裁判；各人都是按自己的尺度來作審判.神呼召基甸作士師,要讓他成為國家的領袖,可是基甸卻未有清楚得到神的指示,恐怕自己弄錯了.這時基甸就祈禱求神給他一個徵兆.他把一團羊毛放在禾場上,若羊毛有霧水,別的地方是乾的,就知道神要藉他的手來拯救以色列人.到了第二天基甸發現他所擺置在地面的羊毛,果然可以搾出露水來；但地面卻仍然是乾的,基甸認為這確是神的聲音.可是他卻又猶疑這也許是偶然巧合.所以再次求神如果祂的旨意要他去拯救以色列人的話,便讓他這次所放的羊毛是乾的,而地面是濕的,翌日他又發現所放的羊毛果然是乾的,而地面卻是濕透.這次他就清楚確定是神的聲音,是神呼召他!他就在那困難的環境裡,作了以色列人的師士拯救他們.
\newline
\newline
(三)神的聲音是可以透過人與我們說話的　譬如神呼召大衛時,並沒有直接向大衛說話；反而叫撒母耳往耶西那裡去,在他的兒子中,要揀選一人作以色列的王.我不知道神為甚麼不命撒母耳直接膏大衛,而要看過大衛的弟兄後才揀選他出來.我只知道神有權用任何方法做任何事情,我卻沒有權命令神做任何事情.許多時我們認為神不照我們的方法做事這想法便是錯了.不要忘記神是全權的,祂可以在任何時間.用任何方法,做任何事情.以大衛一生來說,他曾親自聽見神的聲音,但當他被召作王的時候；他並沒有直接聽到神的聲音,是神藉先知撒母耳向他說話的.
\newline
\newline
在加拿大多倫多一百哩以外,有一條鄉村名叫愛堡.七十五年前,那鄉村只有一小所屬長老會的教會.每月牧師會到這教會,主持一次主日崇拜.最近這地方有另一所教會,距離十哩路遠；在當時還沒有汽車行駛,而馬匹也不多,所以在這社區裡的人,很少機會可以聽到神的話語.有一農夫的女兒名叫「Grace」有見及此,便向他父親說；她希望用一學校,在每星期日下午開辦主日學；好讓那些鐵路工人,和農夫們的小孩子,有機會聽到聖經,這樣她便出了告示,之後每一個星期,有些小孩子來聽她教導神的話語.有一主日,只有十一個小男孩來上主日學,她在開始上課之前,向這些小孩提及有機會,全時間事奉的事.她說在你們當中,每一位都有可能成傳道人.這時坐在左邊有一個衣著破舊,身體軟弱的赤足小童；他心裡對這翻話起了反應,他心想如果神願意的話,我要成為這傳道人.後來神就讓這小童長大後,環球十五次,在七十五個國家裡面傳道,著作了三十五本書；並且被翻譯成一百二十八種語言文字.神還讓他作了不下一千二百首聖詩和福音歌辭,每主日晚上在自己教會,向二千個會眾講道,又用這鄉村小童牧養這教會,用一千六百萬金錢來支持宣教工作,這小童便是我父親.他在今年十一月便屆九十一歲了,如果你有機會與他交談的話；你可能會問,他在年少時,只不過是一個很窮和很軟弱的鄉村小童；沒有甚麼機會往教會裡,但現已活了九十年,而且還作了全球佈道.這究竟是如何會發生的呢?我相信我的父親會回答說他是個鄉村小童,他聽到神的聲音呼召要他傳道.其實,在當時我父親並沒有直接聽到神的聲音,神就是用一個寂寂無名的女子向我父親說話.也許在她有生之日,她不能看到我父親日後所發生的事；她所知道要作的,是忠心於神所交託給她的.有許多時神是透過別人,向我們傳達祂的聲音.
\newline
\newline
(四)神的聲音　是透過環境來向我們說明,在使徒行傳二十一章的中間,我們會看到保羅怎樣被人囚禁著；他當時正在耶路撒冷然卻被羅馬政府捉拿,從那時開始保羅便一直被囚直至到死.在使徒行傳裡面,有六章聖經,講述保羅還有許多地方要去,和有許多事要做,神就在第廿一章開始用環境來指引保羅.在此之前,保羅從不同的方式聽到神的聲音；譬如在大馬色路上保羅就聽到神的聲音；在他立刻悔改之後他又透過先知亞拿尼亞而聽到神的聲音；有一次神透過他所見到的異象,向他說明馬其頓的呼聲.但在廿一章中神用環境來帶領保羅,他當時被羅馬政府,從耶路撒冷轉押上一艘船上.他並沒有選擇在哪一時間,上哪一艘船；但在這艘船上正是一個適合的人.在正適當的時間,向正需要的人,傳講最適合的信息.他做神要他所做的事,神又把他從這艘船,移到一小島上面,而這小島是保羅從沒有計劃要往佈道的.神就以保羅正最適合的人,在適當的時候向正需要的人,傳講最適合的信息.此後神又將保羅從這小島,移到另一艘往羅馬的船上.到羅馬後神又將保羅,放在一個很特別的監房裡；守這監房的人也就是保護皇族的人,這正是一個在地上最具戰略性的講台.因為保羅所說的一切,都會被守衛聽到,而向皇族報告.也就是說從這時開始神的福音,已達到羅馬帝國的中心,在世界上沒有任何一所教會或差會,能夠這樣做到.當我們回顧自己的時候,我們會發覺到神許多時是藉著環境的方法,來帶領我們.我們最著緊的是,神把我們放在那環境,要我們做甚麼我們便要忠心.或許在我們當中有些人還不明白,有些人也許會說要是我親身聽到神的聲音,我也不會相信的.徵兆也不能用科學來解釋,人所說的話我不會信任；我一生生活是胡混潦倒,我決不可倚賴我的環境,好像我這一類人可否聽到神的音呢?
\newline
\newline
(五)神的聲音是藉著聖經來傳達　這聖經就是我個人的小書叢,每次我讀經時,我會聽到神的聲音；我相信在座當中或很多人,希望知道神要他做甚麼東西,要往哪裡去!過些甚麼生活?只要他聽到神的聲音,他便會立刻有反應.在座當中有絕大部份人,我是不認識的,但我能夠準確地說給你們聽,神要你們做些甚麼事情.有三件事神說是要我們做的,神說我要你們做一個祈禱的人；聖經說我們要常常禱告,不可灰心.但有些人說只要我看見很奇妙的徵兆,和神蹟我便會祈禱；但神已經清楚告訴我們要祈禱,祂知道若果我們不願這樣做,就算是有神蹟或奇妙的徵兆,我們也不會祈禱.神又說你要成為我話語的讀書與學生,一個初生嬰孩,要渴慕奶才能夠成長.使你成為無愧的工人,按正義分解神的道；但你也許說如果我能夠看到神,或遇見有人從天上回來,我便會讀聖經.大家可曾記起拉撒路的故事,拉撒路死了往樂園,財主死了往地獄；那財主看到拉撒路在樂園與亞伯拉罕在一起；財主要求拉撒路從樂園來安慰他.但亞伯拉罕說這是不可能的,因為沒有人可以經過那深淵.跟著財主要求亞伯拉罕為他做一件事,那就是差人到人間,向他尚活的五位兄弟說明這地獄的慘況；因為如果有人從死裡復活向那些兄弟說明這些事；他們一定會相信.但亞伯拉罕說,就算有人從死裡復活,他們也不會相信的.因為他們已相信摩西與先知,可是他他們仍不服從所聽到神的話.因此你不要以為如果自己不讀經,神便會行一些神蹟要我們讀聖經,最後神說我要你做見證,「但當聖靈降臨在你們身上,你們就必要在耶路撒冷,猶太全地,和撒瑪利亞直到地極,作我的見證.」正也就說神將我們留在這世界要為神作見證.你有沒有向鄰舍,街坊作見證.有誰為世界未聽過福音的人祈禱和奉獻呢?這些就是今天晚上神要向我們發出的聲音.
\newline
\newline


\newpage

\section{第52屆~港九培靈會~史保羅~第九講　神的呼召}
\label{sec:661}
\textbf{第九講　神的呼召}
\newline
\newline
連結: \href{https://www.hkbibleconference.org/session-message/view/661}{\texttt{https://www.hkbibleconference.org/session-message/view/661}}
\newline
\newline
\hyperref[sec:659]{< < < PREV SERMON < < <}
~
\hyperref[sec:toc]{[返目錄]}
~
\hyperref[sec:662]{> > > NEXT SERMON > > >}
\newline
\newline
請翻閱使徒行傳第十六章,我把這章聖經稱為救恩章.因為這裡記載三個在腓立比城的人,如何接受救恩.第一個是職業婦女呂底亞；第二個是被污鬼附身的使女,第三個是一名普通禁卒.這三人在未接受福音之前,神向保羅已作了一個特別的呼召,往一個特別的地方,做一件特別的工作.這也就是馬其頓的呼聲.對保羅來說,這馬其頓的呼召是要他往那地方作宣教工作.可能今天晚上,神會向我們在坐當中某些人,作同樣的呼召；對其他一些人卻要他們在職位上為主作證.可能今天晚上,神會向我們在坐當中某些人,作同樣的呼召；對其他一些人卻要他們在職位上為主作證.只是我們很容易會混淆不清,不知道我們現在要做何事,是否神呼召我們做的.其實神向我們作出特殊的呼召,必定有三個基本呼召之後.
\newline
\newline
(一)是救恩的呼召　耶穌說:「凡勞苦擔重擔的人,可以到我這裡來,我就使你們得安息.」這是一個對全世界的男女老少的呼召,是一個很個人性的呼召.很多時,我們將這個人性的救恩呼召,與國家宗教混淆一起.在世界上,有許多國家都有國教,提起日本,我們會想起佛教；提到北非,我們會想起回教；提到印度,會想起印度教,當提到加拿大和某些國家時,我們會想起基督教.因為大部份加拿大人,都稱自己為基督徒；而這個情形卻帶來很大的混亂,有些人以為自己在一個基督教國家出生,便順理成章的成為基督徒.這當然是錯的,因為我可以在日本出生,在印度等地方出生；卻不可能是佛教徒,或印度教徒.如果我們將國家宗教與我們的個人救恩混淆,是一件很危險的事.又假如我是在一個基督他家庭裡長大,我的父母是基督徒,他們從少便帶我來到教會,但這卻不會自動使我成為基督徒的.到父親做了一個很敬虔的基督徒,有九十年以上；但他的敬虔,決不可以自然地使我敬虔.我母親是一位新祈禱偉人,但她的祈禱生活亦不會自動地使我成為基督徒的.除非我在一生裡的某一個時刻,回應神的呼召；接受主耶穌為我個人的救主,否則我永遠不能成為一個基督徒.
\newline
\newline
我十八歲進入大學,在大學裡規定每一個學生,必要選修演講這一個科目.教我的老師說,演講最艱難的地方,是要站在群眾面前說話.解決的方法,是找三位我所認識的同學,然後有力地指著他們說,現在要與他們三位說話,我要他們坐好留心聽我說.我永遠不會忘記我當時的第一次嘗試,我望著屋頂,窗門和皮鞋說:「我現在對你說話.」那一次,沒有人留心聽我說話.老師要我們一直的練習,直至我們可以辦得到在人前演講.事隔很久,我站在這一類聚會的講台已經有很多次,但每一次,當我知道再過數分鐘,我要站在講台上向大家講解聖經的時候,我心裡總是這樣禱告:「親愛的主!請讓大家忘記在這聚會裡面,有其他人存在；求神使他們亡記鄰座有人；求你幫助他們忘記是我站在講台上,讓他們知道他們現在是與你面對面.求你讓他們知道,現在是你指著他們說話,是主要他們挺直身兒聽主的說話.」
\newline
\newline
(二)成聖的呼召　當我們對神救恩的呼召,作了回應之後,神會向我們發出成聖的呼召.這呼召只會臨到基督徒的身上.關於這一點,在聖經裡有很多地方也有記載.譬如在羅馬書第十二章第一節:「所以弟兄們,我以神的慈悲勸你們,將身體獻上,當作活祭；是聖潔的,是神所喜悅的,你們如此事奉,乃是理所當然的.」神說如果我們答應神救恩之呼召的話,我們的新生命便要成聖,生命的每一部份都要交給神.當我們接受救恩呼召,我主耶穌作我們個人的救主之後；我們便好像把一所房屋的大門鑰匙交給神,其實這還是不夠的,我們應要將這屋內的其他門匙也交給神.對於做父母的,成聖便是把孩子的睡房門匙交給耶穌基督.你們是否這樣禱告:「神呀!我謝謝你,將兒女賜給我,我要盡我所能的看顧他們,讓他們不受傷害,訓練他們成為有用的人,我會盡心的愛他們.但親愛的神呀!我現在將他們奉獻給你,因為他們是屬於你的人,我就讓你塑造他們.你若是要他們長大後,差他們到別的地方作宣教工作,我們此生再沒有機會相見,我是願意的.又若你只把他們讓我看管五,六年,然後把他們送回到你那裡去,我也是情願的.」各位做父母的,你們有沒有將這一條鑰匙交給神呢?你要明白,這才是真正奉獻的意義.
\newline
\newline
很多時候,我們會在類此這種聚會裡面,將自己奉獻給神,這是一件好事,但只是開始而已.我們要將一生的每一部份獻給神.以經商的基督徒來說,成聖的意思是我們把生意的鑰匙交給神.不過在我們交匙的時候,首先要注意我們所經營的,這類生意是否可以給神；因為有些邪惡的生意,是不能榮耀神的.我們是要把這類生意放棄,這可能是一個很掙扎性的抉擇；但你既然接受了主耶穌基督,你的生活方式便要完全的改變.其次我們要注意的是,我們生意的經營方式,是否可以榮耀神,是否可以讓神看我們的經營手法和賬簿.我是否誠實和仁慈呢?
\newline
\newline
對學生來說,成聖意思是我們要將教育生活的匙獻給神.許多時我們不理會神,便任意選擇學校和要讀的科目,溫習功課時,也把神當作不存在.只要你是基督徒,神對你教育生活的每一部份都有興趣.你要求問神,你該就讀哪一所學校,選修哪一個科目,該花多少時間在功課溫習上,這樣才是真正的奉獻.
\newline
\newline
(三)準備的呼召　在提摩太後書第二章十五節說:「你當竭力,在神面前得蒙喜悅,作無愧的工人；按著正意,分解真理的道.」神要我們在祂的特別呼召來臨前,作好一切的準備.聖經裡面有許多這樣的例子.譬如摩西,他有很好的埃及和希伯教育背景,神還要他作四十年的準備,才讓他把以色列人帶出為奴之地.新約裡面的保羅也是一個例子.試想想當他在大馬色路上悔改後,我們會怎樣.我們也許立刻邀請他,在下一個主日作見證,然後在九龍城浸信會講道,以後便以此為終生工作.若果保羅是這樣的生活,那麼他只是花了三十年在世界講台上,作一個屬靈的嬰孩.每次他可以分享的,只是一個他自己的簡單得救故事.但神決不容許保羅這樣做.保羅信主之後,便花了三年在阿拉伯的曠野,作準備的工作.此後,他才聽見神特別的呼召,作他第一次的傳道和旅行宣教工作.對某些人來說,這準備工作也許是要你入某所神學院得造就,準備做一位全時間的牧師或宣教士.或者對有些人來說,自己太老了,不能再入學讀書,也沒有時間作宣教士的話,你仍可以參加類此這樣的培靈會裡面餵養和教導我們,每天我們這個培靈會都有三小時以上,講說神的話語；但在香港仍有數以千計的基督徒,沒有把握這個機會,很明顯地,他們不著意自己究竟有沒有準備好事奉主.這一類基督徒,永不會聽到神這特別呼召的.另外一個準備自己的方法,是利用屬靈書籍.在現今世代所出版的屬靈書籍甚多,這是教會歷史裡從來未見過的.在我的教會裡面,有些人在一個月裡面,連一本好書也沒有閱讀過,便呆坐著問神為甚麼沒有差使他們工作.譬如教一大班主日學,或參加詩班,作教會執事,甚至差派他到外面作一些特別見證.這只有一個答案,乃是你未曾準備好事奉神!你沒有好好地接受各類屬靈教育的訓練.你可知道在世界上有很多不同行業的人,都要再受教育呢?他們雖然已經成長了,但仍須再受訓練,好讓他們的專業知識可以追得上時代.試想假如我們到一位醫生那裡求診.他說:「我很願意為你診症,但有一點我要向你說明,我是在三十年前的醫學畢業的,從那天起,我一直沒有看過任何的醫學書籍.」聽到這裡,我一定會盡快離開診所.今天我們是否好像那醫生一樣,已經多年沒有在聖經上進修,很久沒有讀過屬靈書籍,我們卻只會問神為甚麼不用我們呢?
\newline
\newline
當我們接受了神救恩,成聖和預備等呼召之後,我們會聽到神特別的呼召,他要我們往馬其頓那裡去.我要你作一個律師,一個家庭裡的母親,一個田裡的農夫,一個醫生；神就這樣給我們作一個見證的崗位.在一九四○年的早期,德國擁有全世界最強大的空軍.在一個晚上,德國大舉轟炸英國,英人聽到警報聲,都奔入防空洞裡面.就在倫敦的近郊,有一個很細小的空軍基地,指揮官連忙命令發高射砲的作好準備；但只有數名士兵在那裡,他又轉向探射燈的部份,只見到寥寥數人；他又命飛機出發與敵人交戰,但卻見機師不足,於是在指揮官旁的有人說:「在我們不遠的地方,有一個藏了百萬人以上的防空洞,我們可以在那裡叫一萬人來看管探射燈,把天空照得如白晝,又叫二萬人來發吧,便天空佈滿砲彈.敵人的飛機無從經過,又命二萬人駕駛飛機迎戰不可嗎?」我相信那指揮官會很憂傷的答道:「我也希望可以這樣做,只是在三年前,國王呼籲國人加入軍隊,大多數的人都不關心這事,他們只管做自己的事情,最後只得數百人來加入軍隊,被訓練如何發砲,用探射燈和駕駛飛機,只有這些已受訓練的人,才能夠擔任這工作,防空洞裡面的人對我來說是沒有用的.」我們可以幻想到主耶穌從天上望下,看見地上有超過四十億的人,其中有三分之一人未曾聽過福音,在日本,印度,中國只有少數的基督徒在工作.有人站在耶穌身旁,看見這情景,便打開各教會的會友名冊,看見有許多名字,便問主耶穌為甚麼不用那些記錄在名冊上的,數以百計基督徒呢?為甚麼不差遣一萬人往印度,一萬人往日本,一萬人往南美呢?我相信主耶穌會向這人露愁容說:「我盼望可以這樣做,但有許多在名冊上有名的人從不接受救恩的呼召,有些已接受救恩呼召,但從不接受成聖的呼召；有些已接受這兩種的呼召；但卻從不接受準備的呼召.而我所可以用的,只是那些接受了三種呼召的人,我已為每個人準備了一個特別工作,為每一個基督徒預備了一個特別的範圍,是他們可以作見證的,只是他們沒有準備好,所以沒有人可以到這些範圍裡面.
\newline
\newline
今天晚上,我們聽了這四種呼召,不知在座各位是屬於哪一類人呢?
\newline
\newline


\newpage

\section{第52屆~港九培靈會~史保羅~第十講　神的旨意}
\label{sec:662}
\textbf{第十講　神的旨意}
\newline
\newline
連結: \href{https://www.hkbibleconference.org/session-message/view/662}{\texttt{https://www.hkbibleconference.org/session-message/view/662}}
\newline
\newline
\hyperref[sec:661]{< < < PREV SERMON < < <}
~
\hyperref[sec:toc]{[返目錄]}
~
\hyperref[sec:1384]{> > > NEXT SERMON > > >}
\newline
\newline
我曾經答應各位今晚要講到神的旨意.關於神的旨意我們起碼要知道五件事實.
\newline
\newline
(一)神旨意是永恆不改變的　首先我要講的,我們有可能疏忽了神的旨意,神從來沒有強迫任何人來遵行祂的旨意和祂要人所做的事；我確信,神掌萬有的一條信條,就是神掌管萬有,簡單來說,神在任何時間,地方喜歡做甚麼,便做甚麼.我相信預定論,就是神能夠由起初到末了,已經預定了一切；這也許很難明白,但這是聖經所記載的.我相信「預知論」,那就是說在事情未曾發生之前,神早已知道這事情的發生.但神仍容許我在某一個程度上有自由；好叫我對人生作出決定.我不應在讀了聖經之後,再存有一種悲觀的人生.很多時我們的錯誤是既然神已經計劃了一切,我就沒有選擇的餘地.因為任何事情若已預定發生,就總會發生,我是毫無能力去改變它.這樣的消極想法,並不是聖經的教訓.讓我們看看聖經所記載的事實吧!舊約基本上乃是敘述以色列人的歷史,記載神怎樣呼召他們,成為神的選民；又怎樣從埃及出來經過曠野,然後到那應許之地.最後以色列人又怎樣在應許之地,建立起王國.在這整個以色列民族歷史裡面,神一直向以色列人說明祂的心意；但以色列人卻不是經常遵守神的旨意,有時候,他們順服神,但很多時候他們卻違背神.這樣看來,他們是有權選擇遵行神或違背神的旨意.此外,約拿的事蹟又是其中的一個例證,記得神呼召約拿往尼尼微城傳道,當然約拿可以選擇去或不去.最後,約拿決定不去,便坐船往相反的方向走.他違背了神的旨意.神在約拿時代是統管萬有,祂十分清楚約拿的違背會有甚麼後果,但約拿仍然固執去違背神的旨意.在新約時代,神設立了救恩計劃,這計劃原來是為那些違背神旨意的人而設的；因為人向神轉臉,叛逆神,以至人墮落沉淪而需要救恩,如果神的旨意是自由施行而人肯遵行的話,我們根本上便毋須轉臉向神拯救了.
\newline
\newline
(二)人怎樣違背神的旨意　任何人,如果聽過福音,卻不肯接受耶穌基督做他個人的救主,他便是離開了神的旨意.如果我從來不肯為我的罪悔改,我就已經離開了神的旨意.彼得說:「主所所許的尚未成就,有人以為他是耽延；其實不是耽延,乃是寬容你們,不願有一人沉淪,乃願人人都悔改.」(彼後3)在這段聖裡面,彼得說明神的旨意是叫我們人人都悔改,歸向耶穌,得著拯救.但聖經其餘的部份卻說,在世界上並不是人人都肯悔改的,也就是說有些人是永遠得不著救恩的.耶穌基督曾經說過有一條寬闊的門路是引向滅亡的,有一條窄小的門路是帶你進入永生的.但不是每個人都走向那引進永生的窄路,亦即是說,不是每個人都肯定向那天堂的道路；耶穌說走在那寬路上面的,有許多人,他們是正走向地獄.如果我是向著那寬大的路走,我便是向著滅亡,已經離開了神的旨意.我雖然是一個好人,衣著光鮮,受很好的教育,有很多錢；但假如我從來未曾為自己的罪悔改的話,我仍然是在神的旨意之外.所以神的旨意是可以被忽略的.
\newline
\newline
(三)神的旨意,著重現在,而並非將來　對你與我來說,神的旨意通常所關心的,都是現在與過去,而較少關心到將來.當我年青的時候,我很喜歡參加培靈會,這類型聚會中,有時我知道講員會提到神的旨意,我就十分盼望講員會給我知道神要我在前面的日子裡,做些甚麼事情.我真的希望能夠知道,神為我準備餘下人生的計劃,最低限度,希望知道神在未來三年中,為我擬定計劃.請各位留心,所謂認識神的旨意,與對前途的預測是無關的.如果我們一向只想著「將來」,「將來」,這是徒然的.神賜我有機會在各地方旅行,講道的時間,很多時會遇到一些有深度的弟兄姊妹,和為主而被拘禁在集中營多年的信徒交通.有時我也有機會能夠與一些世界上最知名的宣教士談話.我認識有些人他將整個生命奉獻給主.我也有機會與一些偉大的祈禱勇士分享,在這些人物當中,我從找不到有一位可以說得出神明天要要他做甚麼事情,在那處有甚麼事情會發生,他們是很有智慧不為前面將來作打算的.我們今天應該盡一切能力對前面稍為作出判斷,應該為著神也許要我們去完成的事情作出計劃.不過,我們必須緊記神的計劃,可能與我們的計劃有所不同.有時,我們可能聽到講員說神為我們的人生,預備了一個美好的計劃；我當然全心的相信這句說話,但是要我將神的計劃影印本拿在手,我便沒有法子了.有時候我們說神已為我的人生劃出了一份藍圖,但我真的是從未看過,誰人拿著這人生計劃的藍圖.當我離開香港之後,我會搭飛機回到多倫多.我已經購買了機票,預訂了座位,知道是那一間航空公司；經過大約二十小時的飛行,我就會在多倫多著陸.我離開機場之後,我便會搭的士回家；回到家後,我就有機會見到內子,小兒,女兒.跟著的禮拜,我會回到民眾教會的講壇,作崇拜的講道,和晚堂講道.這一切只是我的計劃而已.但究竟我能否由香港居留的青年會宿舍,平安去到機場,尚是未知之數；因為香港的交通梗塞,能不能夠幫助我依時到達機場,是很難說的.或者在路上,我遇到交通意外便喪命了也說不定.其次是我的航程安全與否,我也不知道,可能飛機會在海洋上失事也不定.就算我安全到了多倫多,我乘坐的士也會有問題.其實當我回到家後,內子是否在家,我也不能確定；他們可能離家去了旅行也不定.就算家人果真安坐在家,但在下個禮拜我能不能夠回去講道,也是未知.或者,當我今次離加來港的時候,教會已決定將我辭退了,他們不再用我作他們的牧師.各位,計劃我要作出,我是按照我所知道的計劃去做,但是唯有神才知道,我的計劃能否完成.甚麼事情我現在不能確知,但是有一件事,是我知道的,就是現今我是否在神的旨意中,我今天是否正在做著神要我做的事.記著神的旨意著重我的現在,而並非我的將來.
\newline
\newline
(四)神的旨意永遠不會與神的話語互相矛盾　(詩四十8):「我的神阿!我樂意照你的旨意行,你的律法在我心裡.」詩篇是按照猶太人的詩歌體裁而寫成的.猶太人的詩歌表達方式,是在同一節用兩種方式來表達同一件事和意思.這稱為猶太人詩歌的「平衡句」.在聖經的「詩篇」和「箴言」,有很多這體裁的例子.這節聖經正是這樣,他將神的旨意與神的律法,神的話語連接在一起.在詩人的腦海裡面神的旨意,與神的律法是同一件事情,詩人認為假如我要遵守神的旨意的話,我惟一的方法是遵行神的律法.在舊約的時代,神的律法乃是神寫下的說話,要祂子民所做的事情；而這律法是具體性的,適合任何時間時代閱讀.除了說假如我要遵行神的旨意的話,我便必須要熟知神的律法之外,耶穌基督也曾如此說過.(太七21):「凡稱呼我主阿,主阿的人,不能都進天國；惟獨遵行我天父旨意的人,纔能進去.」我相信當耶穌基督在教導這一群聽眾的寺候,那些聽眾並不會陌生,他們大家都知道只有那些遵行神旨意的人,纔能夠上天堂.但我相信在他們心裡面,有第二個問題是:「主阿,我們知道如果要往天堂的話,我們要遵行你的旨意,但我們怎樣可以知道你的旨意呢?」於是耶穌回答說:「有兩個人同建造房屋,一個是聰明人,另一個是無知的人.聰明人所建的房屋,是立在一個穩固的根基上；無知的人所建的房屋,是立在沙土上.此後,有風雨來打擊這兩所房屋；無知的人所建的房屋,在風雨中倒塌了；而聰明人所建造房屋卻依然吃立.主耶穌隨著說,聰明人就是聽了我話,而又肯遵行的人.而無知的人,就是那些聽了我話,但卻不去行的人.這也是說,要明白神的旨意,就必須要知道神所說的話.而全世界只有聖經,和引述聖經的書籍,是記載神的話.因此聖經在基督徒的生命上,是十分重要的.如果你稍為失去了對神話語的絕對權威,和正確的信心,你便永遠找不到神的旨意,我現在已深信這本聖經是絕對的真確,裡面絕無錯誤,是我人生的惟一權威.如果我所做事情,是與聖經和神的話語有違背的話,我就知道這件事並不是神的旨意了.如果我以為神要我去某處地方,但這要去的地方,是與神的話語有相違背的話,這便不是神的旨意了.神的旨意,永遠受著神的話語所控制；我們必須熟習這本聖經,才可以知道神的旨意.
\newline
\newline
(五)神的旨意是決定我們在神面前的境況　神的旨意更多與我們的屬靈境況有關,多過於我們在地理上的衡量.很多時候,我們想到神的旨意,我們就顧慮到地理上的問題.譬如我們考慮到結婚的對象時,假如我與本處的人結婚,我便會在本處居住；如果我與別處的人結婚,我便要跟隨著那人居住,如果我在這所學校讀書,我便向另外一個方向；如果我往另外一所學校的話,我便向另外一個方向.以職業來說,如果我做醫生的話,我便在那處居住；如果我成為宣教士,我便往海外的國家居住；如果我做教師,我便會整生人在一所學校裡居住,如果我是護士,我便會將時間花在醫院裡.但如果我們讀聖經的時候,你會發神的旨意原來與地點無關.神的旨意亦與我們在哪處無關,而是決定於我們在神面前的境況.關於神的旨意,這裡有三節的三節經文.第一節聖經是(羅十二1)「所以弟兄以神的慈悲勸你們,將身體獻上!當作活祭,是聖潔,是神所喜悅的；你們如此事奉,乃是理所當然的.」昨天晚上我們知道這節聖經要求我們將生命奉獻給神,這與今日我們與誰結婚無關；與我們選擇哪一行業無關,與我們選擇哪一間學校無關.這節聖經是說到我個人向全能的神所奉獻的心,但在(羅十二2):「不要效法這個世界,只要心意更新而變化,叫你們察驗何為神的善良,純全可喜的旨意.」保羅將神的旨意與勸我們將自己奉獻給神連結在一起,他說神的旨意,乃是我們奉獻神有關.另一節聖經送給各位是在(帖後四2),讓我首先先讀一次錯誤的:「神的旨意就是要你與正確的人結婚.」你也許會說,聖經根本沒有指名道姓說與誰結婚.各位,你可發現到原來與誰結婚,在這節聖經與神的旨意無關.你會說:「牧師呀?我與一位正確的人結婚,豈是錯呢?」你當然是對的,與正確的人結婚固然是重要的.只是有這可能是結婚對象正確,但與神的關係卻不正確.神真的為我的人生選擇了一位男子或女子,而我果真與那位結了婚；但如果我個人與神的關係不正確的話,我仍然離開神的旨意.各位見到嗎?假如在地點上,我確是在神要我在的地點上,但是我一樣可以在神旨意之外.這節聖經也從來沒有提到選擇適當的學校.你可曾發現到原來你選擇那所學校與神的旨意有關.你也許說我往一所正確的學校,豈不應該的嗎?不錯,你需要往神要去的學校；但是假如所選擇的學校是對,而你個人與神的關係不對,你仍然在神旨意之外.在地點上你可以往正確的地點,但個人關係與神不正確,你仍然是在神旨意之外.讓我們再看最後一節經文,(西四12)下半節,保羅說:「願你們在神一切的旨意上,得以完全.」請問各位「完全」是在哪裡才可以找到呢?是在哪個地方,城市呢?你說這與地點無關,所謂「完全」是不是指世界上任何的地點；乃是一個人在神的面前的光景,個人與神的關係.這是神的旨意,要你與神有正確的關係.各位可曾記得我前數晚提及如果有三,四扇敞開的門,我應該選擇哪一扇呢?或者你說有三位女子(或男子)要與我結婚.首先提到結婚條件,對方一定是基督徒,你知道你能夠與他們(她們)相處得很好,你知道你們兩者之間,有很多相同的地方；但你仍然不知道神為你預備了哪一位,弟兄或姊妹做你終身的伴侶.你面對一個這樣重要的問題時,讓我給你一個建議；首先用些時間忘記那三,四扇敞開的門.你要進入祈禱的內室關上門,臉伏於地,在神的面前省察你的內心；你一定要弄清楚,將一切的罪在神面前承認.你要確知你在神面前,要與神有正確的關係,你要確知你與神有完全的交通.那程度猶如一位黑人的詩:「完全無隔膜,我與救主深交.在神與我之間,全無阻隔.讓神的福氣臨到我,在我心靈之間,不再有罪惡攔阻我與神之間的交往.在我心裡面,不再有內疚,使我與神的恩賜的福隔離.」現在我可以起來深知道我與神之間全無隔膜.這樣你起來,從這扇門走出去,第一個經過的男子(或女子),你便與他(她)結婚吧!當你肯這樣做的時候,神就會負責任,要那一位男子或女子,首先經過你的門前.各位你可見到很多人,把神的旨意說得太複雜,但聖經卻將這件事說得很簡單.聖經說你要在神的旨意當中,你必須與神有正確的關係.聖經說當我救主與我之間,不再有攔阻的時候,我便是在神的旨意裡面.各位主會負我的責任,祂會經常查察我屬靈的光景,至於在地點上,和時間上,這是神的責任,就讓神照顧這些問題好了.如果我作為父親的,我要問一問神,神啊!是否一個你所喜歡的父親?如果我為人母親的,我就要確知道我究竟是不是一位神要我這樣做的母親.如果我準備結婚的,我必須在神的面前查察自己的光景,我是不是神要我做這樣的人,如果我要選擇職業的話,同樣,我要查察我究竟在神的面前,是不是神所要這樣子的人,當今天晚上我們聚會結束的時候,我盼望每一位向自己發問這個問題.我是不是神要我成為那個樣子的青年男子呢?我作為僱主的,我是不是真的按照神的旨意作為這樣的僱主呢?我若作為受僱者的,我是不是按照神要我這樣子去做呢?今天晚上我們在這聚會裡面,我們能否說:「在我與神之間再無罪惡和任何的隔膜?」如果我真的能夠說這句話,我知道現在,我是在神的旨意當中；否則,我便是在神的意之外了!
\newline
\newline


\newpage

\section{第52屆~港九培靈會~楊濬哲~序}
\label{sec:1384}
\textbf{序}
\newline
\newline
連結: \href{https://www.hkbibleconference.org/session-message/view/1384}{\texttt{https://www.hkbibleconference.org/session-message/view/1384}}
\newline
\newline
\hyperref[sec:662]{< < < PREV SERMON < < <}
~
\hyperref[sec:toc]{[返目錄]}
~
\hyperref[sec:614]{> > > NEXT SERMON > > >}
\newline
\newline
籌備一次大型的聚會,真是不簡單,以我們有限的人力,物力,若非有神大能的幫助,和主內同工,同道們的合作支持,實在是無法成事的.每年一度的港九培靈研經大會,情形越來越覺得增長快速,已不是我們力量所能及的.當聚會結束時,我們纔能鬆一口氣,戰戰兢兢不住的感謝神,將榮耀歸給祂!
\newline
\newline
回憶第五十二屆大會,(八零年)末後的幾個晚上；赴會的人潮已在六時湧到,未到七時全堂滿座；甚至副堂,詩班室,會議室,主日學課室,全數爆滿.另一方面,站在行人路上,禮拜堂前走廊的,還不計其數；驚動了警署要派出六位交通警察來維持秩序.
\newline
\newline
主的道興旺,是無可諱言的:若不是耶和華建造房屋,建造的人,就枉然勞力!
\newline
\newline
按照過去的經驗,八一年我們大膽嘗試,借用宣道中學禮堂為副會場,並裝置閉路電視.求主恩待,賜我們好的成果.也期望著會眾不住的祈禱.
\newline
\newline
本講道集是五十二屆的講道紀錄,值得向讀者推荐和介紹.焦源濂牧師主講的約伯記,實在是受苦信徒的慰心劑,使受苦中的信徒充滿盼望.薛玉光先生主講的──教會的合一真理與見證,乃是針對末世教會所需要的信息,也是牧會的同工,不可不讀的.史保羅牧師主講的林前十三章──愛,分七講:另三晚的專題,都是極適合書中信徒的信息.希望讀者好好珍惜本書,活學活用.
\newline
\newline
本講道集蒙朱雷麗端姊妹,梁壽華弟兄,陳士琦弟兄記錄,又蒙陳士琦弟兄之以琳設計公司,設計封面及海報,使本書增光不淺!本書由朱建磯牧師編輯校對,乃得如期出版與讀者見面.以上各位之辛勞,願主記念和悅納.
\newline
\newline
蒙九龍城浸信會每年來惠借場地,及同工們的合作,義工們的盡力勷助,每年大會乃得如期舉行.統此致謝.
\newline
\newline
請多多記念兩位元老委辦,黃原素牧師和鄭德音牧師,身體的軟弱,年老多病；願主憐憫看看他們,大家也切切為他們祈禱.
\newline
\newline
願主藉著本書,造福多人.
\newline
\newline


\newpage

\section{第52屆~港九培靈會~焦源濂~第一講}
\label{sec:614}
\textbf{約伯記簡介}
\newline
\newline
連結: \href{https://www.hkbibleconference.org/session-message/view/614}{\texttt{https://www.hkbibleconference.org/session-message/view/614}}
\newline
\newline
\hyperref[sec:1384]{< < < PREV SERMON < < <}
~
\hyperref[sec:toc]{[返目錄]}
~
\hyperref[sec:618]{> > > NEXT SERMON > > >}
\newline
\newline
感謝神的恩典,給我這次有機會,和大家分享約伯記的真理和信息.在這八天中,我們要講全部約伯記共四十二章:因時間急促,希望各位預先閱讀這卷書,以免臨時難以瞭解.
\newline
\newline
首先我們須明白約伯記在聖經中的地位.全部舊約的聖經,可分為四大部份:一,律法書,有五本.二,歷史書,有十二本,論以色列的立國,興旺,衰微,滅亡.三,詩歌書,共有五本,文字體裁以詩歌形式寫出,表達神人之中的情感和愛慕.四,先知書,共十七本.全部舊約論神選民的歷史.以色列立國,根基非建立於人的智慧和才能上,乃建立在神的律法上；惟神的話是磐石,是隱固的根基.歷史證明,何時離開神的律法,國家便衰微；遵守神的律法,國家就穩固,強壯,興旺.多數人主張詩歌書是寫於選民國家強大興旺時,我們看見神人之間正常的關係,是多有讚美.人屬靈生命也是一樣,在神面前有詩歌讚美,他屬靈的光景必然美好；如果滿懷憂愁,長聲嘆息,他所唱的是哀歌,靈性處低潮.多數先知書,是寫於選民國度衰微之際；因神興起時代的工人,故先知的出現,是為警誡教導和勸勉.
\newline
\newline
約伯記作者,至今仍無統一的說法,有說是摩西,有說是約伯,有說是以利戶,各說法都有所根據,但卻無絕對之把握.根據本書內容,可以肯定作者必具下列資格和條件.本書事實的根源,第一,作者是個屬靈人,有屬靈的啟示.首兩章就記載天上所發生奧秘的事,將屬靈所見到的奧秘,有所啟示,他必然是個屬靈的人,才能參透屬靈的事.第二,作者是個很有智慧,很有學問的人.本書所有探討的問題,非一般人所能寫出,言論大有智慧,深討範圍廣泛,包括天文,地理,物理,動物,植物,人生哲學,論理等,絕非凡夫俗子所能作.第三,作者是個猶太人.此書原文是用最美麗的希伯來文寫的.保羅曾說,神的聖言是交託猶太人的.(羅三2)意即神將真理的啟示交給猶太人.舊約聖經共卅九卷,都是經過猶太人傳達給全世界的人.我們如果瞭解以上三方面,對作者就略有概念.總之,我們不可忘記,真正聖經的作者是神的聖靈:因為聖經都是神所默示的.所有聖經學者告訴我們,此書是全部聖經六十六卷中最古老的.根據以下數點可資證明:
\newline
\newline
(一)無律法的痕跡,聖經是猶太人得到摩西傳神律法之後所寫,其中難免有律法的規條,禮儀,祭物,會幕,祭司等等之事:凡此在約伯將全無記載.
\newline
\newline
(二)本書記載人的財產,非用金銀,乃用實際的物質.約伯是大財主,他的財富用牛羊,僕婢表明,和亞伯拉罕時代相同.
\newline
\newline
(三)約伯受苦時,年七十歲,受苦後又活了一百四十歲.神加倍賜他年日,可能活了二百一十歲.他壽命長度,也和亞伯拉罕時代差不多:亞伯拉罕活了百多歲,以撒活到一百八十五歲才死.
\newline
\newline
(四)約伯是個好家長,為兒女獻祭.那時的人還沒有利未人當祭司,他們是以家長代祭司職份.創世記八章,說挪亞出了方舟為家庭獻祭,廿二章記亞伯拉罕為家庭獻祭.約伯實際情形也和他們相似.
\newline
\newline
以上四點我們可以斷定本書是聖經中最古老的書,甚至早於摩西五經.
\newline
\newline
常有人問我:「義人為何受苦?」這是約伯記的主題.相信在座中,沒有人有生以來從未受過苦的,或多或少必會有.以利法說:「人生在世,必遇患難,如同火星飛騰.」(五7)苦難不但使人肉體受苦,也使人心靈問題重重,一般的苦難都令人困擾.我們難以明白義人為何受苦,這問題是基督徒信仰方面最大的挑戰,雖難尋答案,但作為基督徒的,不能置此問題不理；因這問題接觸我們信仰的根基,我們如無清楚的看法,不能瞭解,一旦問題臨到,我們若非閉眼不顧消極逃避,便是軟弱跌倒.
\newline
\newline
大家都知道,我們的神是全能的神,在祂無所不能,是滿有慈悲,憐憫,全愛的神；當人受苦時,祂要救人脫離苦難；但為何神容許義人受苦?為何見義人受苦而袖手旁觀?是否神愛莫能助呢?義人受苦,使我們對神的全能和全愛發生懷疑；要解決這個問題,是我們心所渴望和要求,也是神願意給我們清楚的答案；所以神藉著約伯記,揭露一個真正的義人,約伯受苦的事實和經歷,來打開我們的眼睛.
\newline
\newline
第一章一至五節,是義人約伯,在受苦之先的情形,給我們對約伯有相當全面的了解.以歷史背景來看,約伯的名字是受苦和遭難的意思.根據(創四13),約伯是以薩迦的兒子,也就是雅各的孫子；除此節經文之外,無法找到更確鑿的證據.約伯住在烏斯地,該地是聖地,位在巴勒斯坦東南,阿拉伯西南的一帶地方.
\newline
\newline
本書再三強調,介紹約伯個人的特點,第一節論及他為人完全正直,敬畏神,遠離惡事.這是千真萬確的事實,有三次同樣的介紹,一次聖靈藉作者介紹,兩次神親自作見證.
\newline
\newline
一,約伯是個完全人,本文所指「完全」的意思,不是說他從未犯罪,非罪人也,乃是說他生活行為經常完全,難找出缺點.約伯是個好父親,好丈夫,好的朋友,好主人,好學者,好領袖,好財主；大有才能,性格完善,難找瑕疵；世人之中是少有的.
\newline
\newline
二,正直行事為人的存心和態度,對神信心純正,對人愛心無虛偽,不自私；如果雅各果真是他的祖父,那麼,孫子遠勝於祖父.創世記廿八章載,雅各蒙神啟示和揀選,神在伯特利向他顯現；應許賜福給他,讓萬國因他得福,國度向四方擴展.後來雅各向神有五項要求；願神在他所行的路上與他同在,一路保佑他,給他食物,給他衣服,平安回到祖家.然後把他答應神作三件事:以神作他的神,用石頭建成神的殿,獻上十分之一；這是約伯祖父對神的態度.至於約伯對神信心純正,神看不出他任何虛偽,神說他完全正直裡外如一.當約伯受苦時,他沒有內疚.一個不正直的人,外強中乾,外表耀武揚威,內心卻軟弱；如伯沙撒王在神天秤上顯出虧欠,裡面的彎曲敗壞顯明出來.約伯在百般痛苦中,自知在神的天秤上,神知道他的純正,所以裡面有大平安.(三十一5-6)
\newline
\newline
我們遇苦難時可怕的不是外面的風波,而是裡面魔鬼的控告和罪惡.
\newline
\newline
三,敬畏神,遠離惡事這是約伯信仰的特色,屬靈的實際並非他天然的性格,非後天的修養,也非人生學問的培養,完全正直的根基乃是他與神之間,有正常,良好的關係,是屬靈關係的產物.他敬畏神,遠離惡事,是他人生智慧的原因.聖經說:「敬畏耶和華是智慧的開端.」他為人的智慧是屬天的,神是他美好的見證,能力的根源；約伯從神支取,取之不盡,用之不竭.
\newline
\newline
約伯生活中,家道富足.有的人對富足有所誤會,說「為富不仁」,有的人以為富足都是投機取巧而得,眼見許多富有者,犯罪墮落；其實,富足本身並非壞事.伊甸園是富足之地,乃神恩典之一,神願祂的兒女身體健康,生活富足.約伯家產宏厚,估價現值一百廿五萬美元,當時列為首富,約伯富足而受苦,這見證是個深明的見證,最打動我們的心靈.他的財富毀於一夕之間,絲毫不留,但他仍然讚美神.
\newline
\newline
約伯兒女眾多,有七男三女.七和三均屬完全數,相加得十也是完全數.他的兒女都長大成人,興家立業,循規蹈矩,彼此相愛,可以說,約伯是個成功的父親.
\newline
\newline
約伯有豐富的智慧和知識.知識是對客觀事物有正確清楚的認識.智慧是正確使用知識的能力.有知識者未必有智慧.
\newline
\newline
在台灣有個資本家開設大工廠,他屬下的工程師多是博士,而他本身卻目不識丁.
\newline
\newline
我曾參加美國MIT麻省理工大學一位名教授的喪事聚會,此人有學問,有成就和有財富,可惜他人生的失敗；他的家庭分裂,妻兒神經失常；他身為教授,大有學問,知識,然而卻沒有智慧做人,結果憂鬱可憐而死.
\newline
\newline
約伯知識,智慧兼有,他以智慧治理家庭,(這是現今信徒很缺少的)他暗中按次為兒子祈禱,恐他們犯罪,心中離棄神.可見他有屬靈的智慧.
\newline
\newline
約伯這樣的人也有受苦的時候；因為有神的美意在其中.
\newline
\newline


\newpage

\section{第52屆~港九培靈會~焦源濂~第二講　撒旦對約伯的第一次攻擊}
\label{sec:618}
\textbf{第二講　撒旦對約伯的第一次攻擊}
\newline
\newline
連結: \href{https://www.hkbibleconference.org/session-message/view/618}{\texttt{https://www.hkbibleconference.org/session-message/view/618}}
\newline
\newline
\hyperref[sec:614]{< < < PREV SERMON < < <}
~
\hyperref[sec:toc]{[返目錄]}
~
\hyperref[sec:620]{> > > NEXT SERMON > > >}
\newline
\newline
約伯記 1:6-1:22
\newline
\begin{longtable}{cl}
\hline
\hline
章節 & 經文 (和合本修訂版)\\
\hline
1:6 & \begin{tabularx}{0.7\textwidth}{X} 有一天，神的眾使者來侍立在耶和華面前，撒但也來在其中。 \end{tabularx} \\ \\ \relax
1:7 & \begin{tabularx}{0.7\textwidth}{X} 耶和華對撒但說：「你從哪裡來？」撒但回答耶和華說：「我從地上走來走去，在那裡往返。」 \end{tabularx} \\ \\ \relax
1:8 & \begin{tabularx}{0.7\textwidth}{X} 耶和華對撒但說：「你曾用心察看我的僕人約伯沒有？地上再沒有人像他那樣完全、正直、敬畏神、遠離惡事。」 \end{tabularx} \\ \\ \relax
1:9 & \begin{tabularx}{0.7\textwidth}{X} 撒但回答耶和華說：「約伯敬畏神，豈是無故呢？ \end{tabularx} \\ \\ \relax
1:10 & \begin{tabularx}{0.7\textwidth}{X} 你豈不是四面圈上籬笆圍護他和他的家，以及他一切所有的嗎？他手所做的都蒙你賜福，他的家產也在地上增多。 \end{tabularx} \\ \\ \relax
1:11 & \begin{tabularx}{0.7\textwidth}{X} 但你若伸手毀他一切所有的，他必當面背棄你。」 \end{tabularx} \\ \\ \relax
1:12 & \begin{tabularx}{0.7\textwidth}{X} 耶和華對撒但說：「看哪，凡他所有的都在你手中；只是不可伸手加害於他。」於是撒但從耶和華面前退出去。 \end{tabularx} \\ \\ \relax
1:13 & \begin{tabularx}{0.7\textwidth}{X} 有一天，約伯的兒女正在他們長兄的家裡吃飯喝酒， \end{tabularx} \\ \\ \relax
1:14 & \begin{tabularx}{0.7\textwidth}{X} 有報信的來見約伯，說：「牛正耕地，母驢在旁邊吃草， \end{tabularx} \\ \\ \relax
1:15 & \begin{tabularx}{0.7\textwidth}{X} 示巴人忽然闖來，把牲畜擄去，並用刀殺了僕人；惟有我一人逃脫，來報信給你。」 \end{tabularx} \\ \\ \relax
1:16 & \begin{tabularx}{0.7\textwidth}{X} 他還說話的時候，又有人來說：「神從天上降下火來，把羊群和僕人都吞滅了；惟有我一人逃脫，來報信給你。」 \end{tabularx} \\ \\ \relax
1:17 & \begin{tabularx}{0.7\textwidth}{X} 他還說話的時候，又有人來說：「迦勒底人分成三隊忽然闖來，把駱駝擄去，並用刀殺了僕人；惟有我一人逃脫，來報信給你。」 \end{tabularx} \\ \\ \relax
1:18 & \begin{tabularx}{0.7\textwidth}{X} 他還說話的時候，又有人來說：「你的兒女正在他們長兄的家裡吃飯喝酒， \end{tabularx} \\ \\ \relax
1:19 & \begin{tabularx}{0.7\textwidth}{X} 看哪，有狂風從曠野颳來，襲擊房屋的四角，房屋倒塌在年輕人身上，他們就都死了；惟有我一人逃脫，來報信給你。」 \end{tabularx} \\ \\ \relax
1:20 & \begin{tabularx}{0.7\textwidth}{X} 約伯就起來，撕裂外袍，剃了頭，俯伏在地敬拜， \end{tabularx} \\ \\ \relax
1:21 & \begin{tabularx}{0.7\textwidth}{X} 說：「我赤身出於母胎，也必赤身歸回；賞賜的是耶和華，收取的也是耶和華。耶和華的名是應當稱頌的。」 \end{tabularx} \\ \\ \relax
1:22 & \begin{tabularx}{0.7\textwidth}{X} 在這一切的事上，約伯並沒有犯罪，也不以神為狂妄。 \end{tabularx} \\ \\
[1ex]
\hline
\hline
\end{longtable}
昨天我們講約伯記大綱,本書討論一個最古老的問題,是人人所必發生的；它使基督徒困擾,它對於信仰是個極具挑戰的問題.神為了解答義人為何受苦的問題,為使基督徒能在義人受苦的事實,找到清楚屬天的啟示；教我們受苦時,知道如何面對事實,也知道如何勝過；神就揀選了義人約伯,作為真理的懿範.
\newline
\newline
從第一章一至五節,我們看見受苦之前的約伯.昨天曾經講過,第一,他有最幸福的家庭.第二,有最美好的靈性.
\newline
\newline
今天我們來看約伯另兩個為人的特點:
\newline
\newline
第三,他有難能可貴的朋友,一般財主的朋友,不外牌友(雀友)馬友,酒肉朋友,玩友,假友,勢利朋友非屬此輩.我們從聖經看見,當他失去一切之時,他的三個有道德,有修養,有學問,靈性高的朋友；不遠千里而來,為他所遭遇的痛苦而安慰他.他們看見他所受的苦,非憑言語可以安慰；所以想尋找他受苦的根源,然後給予澈底幫助.由四至卅一章,都是他們長篇談話的記錄,他們的話講了又講,或者會使讀者覺得厭煩；但這些冗長的談話,反使我們覺得這樣的友情,何等可貴!他們都是民族領袖,是學者,是道德家,是約伯的知友.
\newline
\newline
第四特點,約伯有尊榮的地位.「在東方人中,就為至大.」在當時的社會裡,他是一個在各個不同階層人士,共同敬愛的領袖,連王子對他也肅然起敬.他能得尊榮,並非因財富多,勢力大所形成的；乃因他是個天然的領袖,是受人尊敬的對象.無論學問,修養,性格,靈性,個人生活,家庭生活各方面都完全,而眾望所歸,成為一個公正的好領袖.以上各點,我們可下個結論；約伯乃古今罕有的完全人,世上少有的幸福人.照人的眼光看,約伯是沒有理由要受苦的.
\newline
\newline
按一般情形而論,人受苦是必有原因的:例如因貧窮,無知,兒女不孝,社會歧視,身體軟弱,結交損友,性格古怪等等都是造成痛苦的因素；但這些因素都與約伯無關.另一方面,因犯罪和犯錯誤也能使人受苦；大衛犯罪受苦,摩西犯了錯誤受苦；但是約伯卻完全正直,敬畏神,遠離惡事,為何他也受苦呢?他在未受苦之前,自己十分有把握地說:「我必平安離開我的安樂窩.」誰料到他不但受苦,而且還苦上加苦.
\newline
\newline
約伯記告訴我們:義人為何受苦?首先要知道,義人受苦是個屬靈的問題,並非一般人所能了解明白的.因它的根源是屬靈的,是隱藏的.魔鬼卻在背地工作.叫義人受苦.這一點,就是令人費解之處.按我們肉眼所見的都是人間的活動,但本書在開始時就把我們帶到天上,讓我們明白苦難發生之前,天上曾經發生過一件事.魔鬼製造苦難來對待約伯,預謀破壞他的信仰；牠要藉此以不費吹灰之力,也破壞其他信徒的信仰.牠若打倒了屬靈的領袖,其餘的就可以不攻自破了.這就是撒旦使約伯受苦的主要原因.
\newline
\newline
主的門徒彼得,是撒旦在主未上十架之前,所要打倒的對象.主為他禱告說:「西門,西門,撒旦要篩你們,如同篩麥子一樣.你回頭以後,要堅固你的弟兄.」西門是領袖,是撒旦所注目的；所以主為他禱告,因為他跌倒會引致其他弟兄跌倒；他剛強,可使其他的弟兄堅固剛強.
\newline
\newline
神容許義人受苦,為要高舉義人,叫撒旦蒙羞；這是屬靈的爭戰,有爭戰就必有苦難.創傷纍纍的兵士,乃是經歷苦戰的記號,義人就是基督的精兵.保羅對提摩太說:「你要和我同受苦難,好像基督耶穌的精兵.」(提後二3)義人受苦,為要把義人的尊貴顯明出來.經歷苦難的義人,對神纔有更深的認識；神藉苦難把義人的生活,帶到更深的境界.
\newline
\newline
本書末了約伯說:「我從前風聞有你,現在親眼看見你.因此我厭惡自己,在塵土和爐灰中懊悔.」(四十二5-6)苦難使義人踏入屬靈新境界,打開義人屬靈的眼睛,看見屬靈的新事.
\newline
\newline
從約伯和三個朋友指責他的談話中,我們看見約伯對自己的認識膚淺,對神的認識也膚淺.保羅雖有三層天的經歷,但他說:「我不是以為我已得著了,我乃是忘記背後,努力面前的,向著標竿直跑.」約伯在屬靈的事上停頓,不知追求長進,所以神容許撒旦攻擊他；但是臨到約伯的一切苦難,都會經過神的衝量,沒過於他所能忍受的.神限制撒旦,不准牠越過範圍來傷害義人；因為義人在主的眼中是極為寶貴的.我們在主裡,主藉著這樣一個義人苦難的奧秘,和實際的經歷,來表達在我們面前；這個教訓,不是道理,乃是事實,是見證.義人應從本書得著極大的激勵.
\newline
\newline
撒旦對約伯的第一次攻擊(二1-3)
\newline
\newline
(一)天上的辯論(一6-12)義人在地上發生的事,根源在於天上.約伯是神,撒旦和天使共同注意的中心.基督徒應時常提醒自己,注意事情發生的根源在哪裡,不要小看自己,以為自己算不得甚麼.約伯成為神和撒旦爭辯的對象,天使是見證者.「神的眾子」在詩篇中是指眾天使而言(參詩二十九1)本書第一,第二章都有說:「神的眾子來待立在耶和華面前.」從聖經看來,眾天使在天上經常有聚集:但在這段時間,「撒旦也來在其中.」撒旦突然出現,一定有牠的動機和目的.撒旦是墮落的天使長,牠的名是敵擋的意思；當神問撒旦「從哪裡來?」回答說:「我從地上走來走去,往返而來.」地上是撒旦活動的範圍,由牠回答神的話,約略可見牠的性格和態度是何等的狂傲；但也可見到牠失喪和不安的情形；好像動物園的狂獸不斷走動侷促不安,意味著內裡受限制沒自由.另一方面亦見到牠的狡猾和詭詐,神出鬼沒令人冷不及防,防不勝防.所以保羅教我們要小心,免得撒旦趁機勝過我們；人以為平安穩妥之時,災禍忽然臨到.
\newline
\newline
「從地上走來走去,往返而來,」這話也表明撒旦終日忙碌作工.「魔鬼如同吼叫的獅子,遍地遊行,尋找可吞吃的人.」(彼前五8)所以遊手好閒的人最易被魔鬼得著.當大衛睡到日頭平西時,他就犯了罪.懶惰的人,沒有目標的人,最容易被撒旦所吞吃.
\newline
\newline
神對撒旦提到約伯之後,約伯就開始受苦.神為約伯辯明:「地上再沒有人像他完全正直,敬畏神,遠離惡事.」(一8)撒旦要苦害約伯,早有預謀；可是牠走來走去,往返而來不得其門而入因為神「四面圈上籬笆,維護約伯和他的家,並他一切所有的；他手所作的,都蒙神賜福,......」(一10)基督徒應該曉得,你和你的家和一切有平安,並非我們有甚麼本事,也非偶然；而是神用籬笆維護,我們應當感謝神.撒旦在神前控告約伯,目的是要剝奪他一切所有而打倒他；但神預先知道的,約伯能夠站立得住.
\newline
\newline
(二)辯論的發生　在宇宙間神所創造的萬物之中,人是神最獨特的創造.按軀體來說,人是軟弱的.詩篇第八篇說:「人算甚麼?」但在靈界說,人較天使微小一點；比天使軟弱的人和神之間可有純潔,親密,自發,鞏固的關係,這是天使所比不上的.這關係在穹蒼和主宰之中非常甘美,是宇宙中最奇妙的事.約伯和常人到底有何不同,他的美麗,尊貴何在?不是單看他的外表,性格和美好的生活,有好的家庭,好的成就；神使他有純潔美好的信,非任何人事物可以將他毀壞的.神對撒旦說:「你曾用心察看我的僕人約伯沒有?」但是撒旦對約伯另有看法,牠認為約伯和神之間的關係,不十分純潔,也不尊貴；約伯無非是是出於自私自利而信靠神.神稱讚約伯,撒旦大不贊同,向神抗議說:「你且伸手毀他一切所有的,他必當面棄掉了.」但神卻不願意伸手使約伯受苦,神十分清楚,相信約伯的為人,必能勝過撒旦所加害給他的一切試探,所以才將約伯交給撒旦；因為神的手是施恩的手,不是降災的手.由此可見人的尊貴和受苦,有很大的關係.經過撒旦的試探,約伯真正的尊貴就顯明出來,他和神的維繫不是財物,而是屬靈純正的信心.經過試驗,可見財主約伯,遠超乎一般財主之上；他有尊貴的信心,純潔的信仰.證明詩篇八2-8節這段經文,我們看見人是軟弱,算不得甚麼；然而神卻高舉人,給人管理萬有,戴尊貴榮耀的冠冕.神讓約伯受苦是要他向全世界證明,苦難是可以得勝的；撒旦雖然猖狂,但神若不許可,牠不能在神兒女身上任意妄為.人勝過苦難,撒旦便蒙羞逃遁,人雖軟弱,但靠主就可向撒旦誇勝,因我們爭戰乃是與天空屬靈氣的魔鬼爭戰.
\newline
\newline
約伯受苦彰顯了神的榮耀.一般宗教,人和所信仰的對象,都是互利的,惟有基督教能顯明我們與神的關係是特別的.神無條件愛世人,神也盼望祂所愛的人能真誠愛祂.不是因為任何好處,乃因神是我們的主；我們的好處不在祂以外,這樣才是好的見證.
\newline
\newline
但願我們學習大衛說:「你是我的主,我的好處不在你以外.除你以外,在天上我還有誰呢?除你以外,在地上我也沒有所愛慕的.」
\newline
\newline
但願我們在主面前,從約伯身上學習功課.
\newline
\newline


\newpage

\section{第52屆~港九培靈會~焦源濂~第三講　撒旦對約伯的第二次攻擊}
\label{sec:620}
\textbf{第三講　撒旦對約伯的第二次攻擊}
\newline
\newline
連結: \href{https://www.hkbibleconference.org/session-message/view/620}{\texttt{https://www.hkbibleconference.org/session-message/view/620}}
\newline
\newline
\hyperref[sec:618]{< < < PREV SERMON < < <}
~
\hyperref[sec:toc]{[返目錄]}
~
\hyperref[sec:621]{> > > NEXT SERMON > > >}
\newline
\newline
約伯記 2:1-3:26
\newline
\begin{longtable}{cl}
\hline
\hline
章節 & 經文 (和合本修訂版)\\
\hline
2:1 & \begin{tabularx}{0.7\textwidth}{X} 又有一天，神的眾使者來侍立在耶和華面前，撒但也來在其中。 \end{tabularx} \\ \\ \relax
2:2 & \begin{tabularx}{0.7\textwidth}{X} 耶和華問撒但說：「你從哪裡來？」撒但回答說：「我從地上走來走去，在那裡往返。」 \end{tabularx} \\ \\ \relax
2:3 & \begin{tabularx}{0.7\textwidth}{X} 耶和華對撒但說：「你曾用心察看我的僕人約伯沒有？地上再沒有人像他那樣完全、正直、敬畏神、遠離惡事。你雖激起我攻擊他，無故吞滅他，他仍然持守他的純正。」 \end{tabularx} \\ \\ \relax
2:4 & \begin{tabularx}{0.7\textwidth}{X} 撒但回答耶和華說：「人以皮代皮，情願捨去一切所有的，來保全性命。 \end{tabularx} \\ \\ \relax
2:5 & \begin{tabularx}{0.7\textwidth}{X} 但你若伸手傷他的骨頭和他的肉，他必當面背棄你。」 \end{tabularx} \\ \\ \relax
2:6 & \begin{tabularx}{0.7\textwidth}{X} 耶和華對撒但說：「看哪，他在你手中，只要留下他的性命。」 \end{tabularx} \\ \\ \relax
2:7 & \begin{tabularx}{0.7\textwidth}{X} 於是撒但從耶和華面前退出去，擊打約伯，使他從腳掌到頭頂長毒瘡。 \end{tabularx} \\ \\ \relax
2:8 & \begin{tabularx}{0.7\textwidth}{X} 約伯就坐在灰燼中，拿瓦片刮身體。 \end{tabularx} \\ \\ \relax
2:9 & \begin{tabularx}{0.7\textwidth}{X} 他的妻子對他說：「你仍然持守你的純正嗎？你背棄神，死了吧！」 \end{tabularx} \\ \\ \relax
2:10 & \begin{tabularx}{0.7\textwidth}{X} 約伯卻對她說：「你說話，正如愚頑的婦人。唉！難道我們從神手裡得福，不也受禍嗎？」在這一切的事上，約伯並沒有以口犯罪。 \end{tabularx} \\ \\ \relax
2:11 & \begin{tabularx}{0.7\textwidth}{X} 約伯的三個朋友，提幔人以利法、書亞人比勒達、拿瑪人瑣法，聽說這一切的災禍臨到他身上，各人就從自己的地方相約同來，為他悲傷，安慰他。 \end{tabularx} \\ \\ \relax
2:12 & \begin{tabularx}{0.7\textwidth}{X} 他們遠遠地舉目觀看，認不出他來，就放聲大哭。各人撕裂外袍，向空中撒塵土，落在自己的頭上。 \end{tabularx} \\ \\ \relax
2:13 & \begin{tabularx}{0.7\textwidth}{X} 他們同他七天七夜坐在地上，一句話也不對他說，因為他們見到了極大的痛苦。 \end{tabularx} \\ \\ \relax
3:1 & \begin{tabularx}{0.7\textwidth}{X} 此後，約伯開口詛咒自己的生日。 \end{tabularx} \\ \\ \relax
3:2 & \begin{tabularx}{0.7\textwidth}{X} 約伯說： \end{tabularx} \\ \\ \relax
3:3 & \begin{tabularx}{0.7\textwidth}{X} 「願我生的那日滅沒，說『懷了男胎』的那夜也滅沒。 \end{tabularx} \\ \\ \relax
3:4 & \begin{tabularx}{0.7\textwidth}{X} 願那日變為黑暗，願神不從上面尋找它，願亮光不照於其上。 \end{tabularx} \\ \\ \relax
3:5 & \begin{tabularx}{0.7\textwidth}{X} 願黑暗和死蔭索取那日，願密雲停在其上，願白天的昏暗恐嚇它。 \end{tabularx} \\ \\ \relax
3:6 & \begin{tabularx}{0.7\textwidth}{X} 願那夜被幽暗奪取，不在一年的日子中喜樂，也不列入月中的數目。 \end{tabularx} \\ \\ \relax
3:7 & \begin{tabularx}{0.7\textwidth}{X} 看哪，願那夜沒有生育，其間也沒有歡樂的聲音。 \end{tabularx} \\ \\ \relax
3:8 & \begin{tabularx}{0.7\textwidth}{X} 願那些詛咒日子且能惹動力威亞探的，詛咒那夜。 \end{tabularx} \\ \\ \relax
3:9 & \begin{tabularx}{0.7\textwidth}{X} 願那夜黎明的星宿變為黑暗，盼亮卻不亮，也不見晨曦破曉； \end{tabularx} \\ \\ \relax
3:10 & \begin{tabularx}{0.7\textwidth}{X} 因它沒有把懷我胎的門關閉，也沒有從我的眼中隱藏患難。 \end{tabularx} \\ \\ \relax
3:11 & \begin{tabularx}{0.7\textwidth}{X} 「我為何不出母胎而死？為何不出母腹就氣絕呢？ \end{tabularx} \\ \\ \relax
3:12 & \begin{tabularx}{0.7\textwidth}{X} 為何有膝蓋接收我？為何有奶哺養我呢？ \end{tabularx} \\ \\ \relax
3:13 & \begin{tabularx}{0.7\textwidth}{X} 不然，我現在已躺臥安睡，而且，早已長眠安息； \end{tabularx} \\ \\ \relax
3:14 & \begin{tabularx}{0.7\textwidth}{X} 與那些為自己重建荒涼之處，地上的君王和謀士在一起； \end{tabularx} \\ \\ \relax
3:15 & \begin{tabularx}{0.7\textwidth}{X} 或與把銀子裝滿房屋，擁有金子的王子在一起； \end{tabularx} \\ \\ \relax
3:16 & \begin{tabularx}{0.7\textwidth}{X} 我為何不像流產的胎兒被埋藏，如同未見光的嬰孩？ \end{tabularx} \\ \\ \relax
3:17 & \begin{tabularx}{0.7\textwidth}{X} 在那裡惡人止息攪擾，在那裡困乏人得享安息， \end{tabularx} \\ \\ \relax
3:18 & \begin{tabularx}{0.7\textwidth}{X} 被囚的人同得安逸，不再聽見監工的聲音。 \end{tabularx} \\ \\ \relax
3:19 & \begin{tabularx}{0.7\textwidth}{X} 大的小的都在那裡，奴僕脫離主人得自由。 \end{tabularx} \\ \\ \relax
3:20 & \begin{tabularx}{0.7\textwidth}{X} 「遭受患難的人為何有光賜給他呢？心中愁苦的人為何有生命賜給他呢？ \end{tabularx} \\ \\ \relax
3:21 & \begin{tabularx}{0.7\textwidth}{X} 他們等死，卻不得死；求死，勝於求隱藏的珍寶。 \end{tabularx} \\ \\ \relax
3:22 & \begin{tabularx}{0.7\textwidth}{X} 他們尋見墳墓，就歡喜快樂，極其高興。 \end{tabularx} \\ \\ \relax
3:23 & \begin{tabularx}{0.7\textwidth}{X} 這人的道路遮隱，神又四面圍困他。 \end{tabularx} \\ \\ \relax
3:24 & \begin{tabularx}{0.7\textwidth}{X} 我吃飯前就發出嘆息，我的唉哼湧出如水。 \end{tabularx} \\ \\ \relax
3:25 & \begin{tabularx}{0.7\textwidth}{X} 因我所恐懼的臨到我，我所懼怕的迎向我； \end{tabularx} \\ \\ \relax
3:26 & \begin{tabularx}{0.7\textwidth}{X} 我不得安逸，不得平靜，也不得安息，卻有患難來到。」 \end{tabularx} \\ \\
[1ex]
\hline
\hline
\end{longtable}
主耶穌說:「在世上你們有苦難.」(約十六33)因為世界有罪的存在,罪在那裡,苦難也必跟著在那裡.魔鬼是世界的王,牠名為世界的神,牠在天上雖沒有地位,但在地上,牠狂妄地說:「我從地上走來走去,往返而來.」世界掌握在惡者手中.一家之長不良,全家受苦；社會領袖不善,社會就有災難；一國的執政者有問題,國民就在災難之中.魔鬼是神的仇敵,是撒謊的,殺人的,挑撥離間的,是製造痛苦的；所以受牠轄制的,怎能脫離苦難?至於神的兒女,為何也受苦呢?
\newline
\newline
我們已經講過,魔鬼在地上雖然自由,走來走去,往返而來；但是屬神的義人約伯,神圈上籬笆維護他和他的一切,使魔鬼無可奈何；所以牠想盡辦法要打擊約伯,破壞約伯.主耶穌為門徒祈禱說:「我不求你叫他們離開世界,只求你保守他們脫離那惡者.他們不屬世界,正如我不屬世界一樣.」(約十七15-16)
\newline
\newline
義人受苦是因撒旦的攻擊,是爭戰的徽號.神容許撒旦攻擊約伯,拆去維護約伯的籬笆；乃因祂要將衪所創造萬物之中的人類,最奇妙的尊貴顯明出來.外表看來,人算不得甚麼,比天使微小一點,能力軟弱,軀體受限制；但卻與造物之主有最甘美,自然,純潔,堅強的關係.這種關係,是任何力量所不能破壞的；人的尊貴,就在義人受苦的事上顯明出來.約伯真可以說為人爭了一口氣.他的尊貴,不在於錢財,聰明和智慧上；乃在於他內心,屬神純潔的信仰上.
\newline
\newline
(一)為何神容許撒旦存在並活動?
\newline
\newline
(1)為造就信徒的靈命　我們雖有生命,但生命不夠豐盛；雖然得救了,仍然是嬰孩式幼稚的生命.我們需要生長,成熟,剛強,需要受試探以達到這目的.(雅一12)說:「忍受試探的人是有福的,因為他經過試驗以後,必得生命的冠冕.......」所以撒旦的試探,能使信徒的生命成熟和剛強.
\newline
\newline
(2)要信徒學習屬靈的爭戰　將來在永世裡面,可得到神永遠的獎賞.保羅在(提後四7-8)離世前說:「那美好的仗我已經打過了.」屬靈的仗不是在和平之國打的,乃是在世上有問題之處打的.地上屬靈的仗已經勝利了,耶穌基督已經藉著死,敗壞那掌死權的魔鬼；但是如今還有一些時間給基督徒在這爭戰上有份.保羅說:「那美好的仗我已經打過了.......從此以後,有公義的冠冕為我存留.」得救是無條件的,但得勝必須經過爭戰.
\newline
\newline
(3)要彰顯神的大能　為何神不親自對付魔鬼呢?如果這樣,神未免太高舉牠了,神只用軟弱的人,就足以對付魔鬼.主耶穌在地上時,是以人子的地位,並非以神子的權能,來對付魔鬼.主全心信服,以至於死,且死在十字架上；神卻將主耶穌升為至高,賜給祂那超乎萬名之上的名.神用軟弱的,勝過強壯的.
\newline
\newline
我是最厭惡螞蟻的,但有一次,我在廚房發現了成群的螞蟻圍食我剩下的食物；我就一一把牠處死,心中還引以自豪.忽然我又感到,就算我殺死千萬螞蟻,也不足為奇；但如果有一天,一隻螞蟻咬我至死的話,就將成為報章上的頭條新聞了.
\newline
\newline
(詩八2)所指嬰孩和吃奶的,乃是軟弱的形象.第四節又說:「人算甚麼,你竟願念他；世人算甚麼,你竟眷顧他.你叫他比天使微小一點,並賜他榮耀尊貴為冠冕.你派他管理你手所造的,使萬物......都服在他的腳下.」在萬物之中,人是神最獨特的創造,人是軟弱的；但人也是強壯的,人和神能有奇妙的愛情,有極大的能力從他的瓦器裡彰顯出來；人是神的榮耀,安慰,滿足和享受.
\newline
\newline
(4)神的心意盼望罪人悔改　人在平安無事時,很難尋找神,苦難卻叫人親近神.聖經所記載的浪子,直到他一無所有,挨餓受苦之時；就使他想起父家來,終於回到他父親的懷抱,保羅在林前五章說,教會中有人犯了淫亂；要把這樣的人交給撒旦,敗壞他的肉體,使他的靈魂可以得救,正是此意.
\newline
\newline
(5)因為得救的人數尚未滿足　(羅十一25)所以神暫且容許撒旦有機會活動.當神毀滅魔鬼時,就是審判之日來到之時；魔鬼被扔到火湖裡,那時,未蒙恩得救的人再也沒有悔改的機會了.時間的結局一到,得救的恩門就關閉.神若今日把撒旦捆綁丟在火湖裡,恐怕我們許多親朋將會失去蒙恩的機會.神寬容忍耐,乃是盼望得救人數滿足.
\newline
\newline
(6)基督徒至今尚未預備好　主來的日子耽延,使基督徒可以作好準備；不至於羞愧,可以無瑕無疵的見主面,得見神的榮耀.(啟十九7)說:「當主再來時,新婦也自己預備好了.」可惜今日許多基督徒還在醉生夢死,在黑暗中浪費神所賜的光陰!
\newline
\newline
(7)神知道基督徒的需要　藉著撒旦可以使基督徒學習謙卑,因為有恩賜的人常常自高自大.保羅也有過這樣的經驗,神容許一根刺在他身上,因為恐他所得的啟示太大,以至驕傲自大.神藉著撒旦的攻擊,使他知道了自己的軟弱,因而他說:「我甚麼時候軟弱,甚麼時候就剛強了.」他能保持謙卑依靠神來工作,纔不至於在不知不覺之中偏離神.
\newline
\newline
撒旦仍有機會在地上活動,雖非神所喜悅,但對於基督徒正是合時的需要.
\newline
\newline
當撒旦得神許可之後,「於是撒旦耶和華面前退去.有一天......」( 一12-13 )由此可知撒旦並非立刻行動,牠等到「有一天」認為那天是最好的一天,最可沉重打擊約伯的一天.約伯在撒旦手中,好像小小的羔羊,處於千萬狼群之中受到殘害,毫無反抗的力量；可見人落在撒旦手中是何等軟弱!
\newline
\newline
(二)撒旦的手段有四個性質.
\newline
\newline
(1)撒旦的攻擊出其不意(一15)「迦勒底人分作三隊,忽然闖來......」(17)「不料有狂風從曠野颳來......」(19)真是閉門家中坐,禍從天上來.許多時候人受苦,不是因犯罪；是由於撒旦的攻擊,是背後有個蠻不講理的撒旦,興波作浪.
\newline
\newline
(2)連連打擊　四次的痛苦連續而來(16-19)瞬息間,接連有人來報惡訊,真是令人莫測,措手不及的攻擊.
\newline
\newline
(3)澈底的毀壞　第一個災難,牛,驢全被示巴人擄去,僕人也被殺.第二個災難,群羊和僕人被火燒滅.第三個災難駱駝被迦勒底人擄去,僕人也被殺.第四個災難,七兒三女全被房屋倒塌壓死.約伯本來財豐力厚,頃刻竟化為烏有,蕩然無存.這樣的打擊,何等厲害殘忍!
\newline
\newline
(4)遮蓋真相　利用示巴人,迦勒底人,火和狂風,以為天災人禍；沒想到是撒旦的作為,叫人怨天尤人,或埋怨自己不幸；使基督徒因此失敗,達到撒旦攻擊之目的.全部約伯記的人物,都不知道是撒旦隱藏的詭計.
\newline
\newline
(三)約伯如何擔當苦難
\newline
\newline
(1)約伯有正確的財產觀念(一20-22)這是約伯得勝的見證,他是人,有人性軟弱的表現,他有人之常情,遭打擊,受災難,當然有所反應.財物盡失,兒女都死亡,他痛苦傷心；這是人性方面的軟弱.但另一方面,他靈性裡面卻非常剛強,他對喪失一切,有正確的態度.他說:「賞賜的是耶和華,收取的也是耶和華.」一切都是神所賞賜,所交託給我們的,主權在於神.但人常忘記主權屬於神,卻把主權據為己有.有些人信心不正確,只建立在神的恩典上,而不是建立於神的本身.越有錢的人,信心越軟弱,愛錢財勝過愛主；越容易被撒旦乘機把信心拿去,以致被牠打倒.你是否這樣?遭遇財物損失時,你埋怨,不滿；信心幾乎被撒旦奪去,那麼,你就應當以約伯作為學習的榜樣了.
\newline
\newline
(2)約伯有正確的人生觀　當他富有時,他未曾以錢財作他的主；耶和華是他的神,他真是堪配說:「我的好處不在神以外.」當他失去一切時,他仍以耶和華為他的神,仍然榮耀神,敬拜神,稱頌神,他說:「耶和華的名,是應當稱頌的.」他以榮耀神為他至高之目的,這是最美麗的人生,是最剛強的人.
\newline
\newline
(3)約伯有敬虔的操練　他平時常禱告,和神之間不斷連絡.一個戰士平日的訓練很重要,不能臨上陣才練鎗.一個基督徒經常讀經,禱告,事奉,有愛心,似乎不覺有奇妙的經歷；但到了遭遇難處時,就感到平時在敬虔上操練大有益外；不致驚惶失措,自然而然穩步邁進正確方向.
\newline
\newline
(4)約伯有康強的體魄　約伯之得勝,我認為身體康強是原因之一.人常常心靈願意,肉體卻軟弱了!當人失去身外之物,必大受打擊,神經就失去常態了.身體健康和心理狀態有關；不過,基督徒要緊記切勿斷章取義來解釋聖經.比如「操練身體益處還少.」倒不如不操練身體了.我們應當像保羅所說:「又願你們的靈,與魂,與身子,得蒙保守,在我主耶穌基督降臨的時候,完全無可指摘.」(帖前五23)
\newline
\newline
撒旦對約伯第二次的攻擊,也是從天上發動的.牠第一次攻擊失敗後仍不甘心,再次出現在天上,要求再次打擊約伯.可見撒旦的性格死不認錯,完全推翻曾經說過的話.人也常有撒旦的情.基督徒得救之後,作風才漸改變；因為「聖靈來了,就要叫世人為罪,為義,為審判,自己責備自己.」突然打擊最能表達內心的真相；但是撒旦對約伯突然施以打擊,取去他的牛羊,駱駝,兒女,並一切所有,約伯卻無動於中；所以撒旦向神要求第二次攻擊約伯.神是無錯誤的神,祂瞭解人勝過任何人對人的瞭解.神對約伯的認識毫不含糊,我們的神是無所不能的神；約伯的純正,何等安慰神的心.
\newline
\newline
巴不得今日的基督徒,真有尊貴的靈性,站立起來；在這苦難的世界,為主作美好的見證.
\newline
\newline


\newpage

\section{第52屆~港九培靈會~焦源濂~第四講　失敗的安慰者-以利法(一)}
\label{sec:621}
\textbf{第四講　失敗的安慰者-以利法(一)}
\newline
\newline
連結: \href{https://www.hkbibleconference.org/session-message/view/621}{\texttt{https://www.hkbibleconference.org/session-message/view/621}}
\newline
\newline
\hyperref[sec:620]{< < < PREV SERMON < < <}
~
\hyperref[sec:toc]{[返目錄]}
~
\hyperref[sec:623]{> > > NEXT SERMON > > >}
\newline
\newline
約伯記 4:1-4:21
\newline
\begin{longtable}{cl}
\hline
\hline
章節 & 經文 (和合本修訂版)\\
\hline
4:1 & \begin{tabularx}{0.7\textwidth}{X} 提幔人以利法回答說： \end{tabularx} \\ \\ \relax
4:2 & \begin{tabularx}{0.7\textwidth}{X} 「人想與你說話，你就厭煩嗎？但誰能忍住不發言呢？ \end{tabularx} \\ \\ \relax
4:3 & \begin{tabularx}{0.7\textwidth}{X} 看哪，你素來教導許多人，又堅固軟弱的手。 \end{tabularx} \\ \\ \relax
4:4 & \begin{tabularx}{0.7\textwidth}{X} 你的言語曾扶助跌倒的人；你使軟弱的膝蓋穩固。 \end{tabularx} \\ \\ \relax
4:5 & \begin{tabularx}{0.7\textwidth}{X} 但現在禍患臨到你，你就煩躁了；它挨近你，你就驚惶。 \end{tabularx} \\ \\ \relax
4:6 & \begin{tabularx}{0.7\textwidth}{X} 你的倚靠不是在於你敬畏神嗎？你的盼望不是在於你行事純正嗎？ \end{tabularx} \\ \\ \relax
4:7 & \begin{tabularx}{0.7\textwidth}{X} 「請你追想：無辜的人有誰滅亡？正直的人何處被剪除？ \end{tabularx} \\ \\ \relax
4:8 & \begin{tabularx}{0.7\textwidth}{X} 按我所見，耕罪孽的，種毒害的，照樣收割。 \end{tabularx} \\ \\ \relax
4:9 & \begin{tabularx}{0.7\textwidth}{X} 神一噓氣，他們就滅亡；神一發怒，他們就消失。 \end{tabularx} \\ \\ \relax
4:10 & \begin{tabularx}{0.7\textwidth}{X} 獅子吼叫，猛獅咆哮，少壯獅子的牙齒被敲斷。 \end{tabularx} \\ \\ \relax
4:11 & \begin{tabularx}{0.7\textwidth}{X} 公獅因缺獵物而死，母獅的幼獅都離散。 \end{tabularx} \\ \\ \relax
4:12 & \begin{tabularx}{0.7\textwidth}{X} 「有話暗中傳遞給我，耳朵聽其微小的聲音。 \end{tabularx} \\ \\ \relax
4:13 & \begin{tabularx}{0.7\textwidth}{X} 世人沉睡的時候，從夜間異象的雜念中， \end{tabularx} \\ \\ \relax
4:14 & \begin{tabularx}{0.7\textwidth}{X} 恐懼戰兢臨到我身，使我百骨戰抖。 \end{tabularx} \\ \\ \relax
4:15 & \begin{tabularx}{0.7\textwidth}{X} 有靈從我面前經過，我身上的毫毛豎立。 \end{tabularx} \\ \\ \relax
4:16 & \begin{tabularx}{0.7\textwidth}{X} 那靈停住，我卻不能辨其形狀；有形像在我眼前。我在靜默中聽見有聲音： \end{tabularx} \\ \\ \relax
4:17 & \begin{tabularx}{0.7\textwidth}{X} 『必死的人能比神公義嗎？壯士能比造他的主純潔嗎？ \end{tabularx} \\ \\ \relax
4:18 & \begin{tabularx}{0.7\textwidth}{X} 看哪，主不信靠他的僕人，尚且指他的使者為愚昧， \end{tabularx} \\ \\ \relax
4:19 & \begin{tabularx}{0.7\textwidth}{X} 何況那些住在泥屋、根基在塵土裡、被蛀蟲所毀壞的人呢？ \end{tabularx} \\ \\ \relax
4:20 & \begin{tabularx}{0.7\textwidth}{X} 早晚之間，他們就被毀滅，永歸無有，無人理會。 \end{tabularx} \\ \\ \relax
4:21 & \begin{tabularx}{0.7\textwidth}{X} 他們帳棚的繩索豈不從中拔出來呢？他們死，且是無智慧而死。』」 \end{tabularx} \\ \\
[1ex]
\hline
\hline
\end{longtable}
約伯受撒旦第一次的攻擊,他的意念堅強是何等的完美；撒旦因為未能打倒約伯,仍不甘心；所以牠要求第二次攻擊約伯.神對於義人約伯的純正,對於他屬靈生命的堅固,從開始就有十分把握；所以毫不猶豫地答應了撒旦,一次,再次的要求；神何等尊重約伯!約伯的價值何等寶貴!
\newline
\newline
撒旦兩次打擊約伯,手法完全不同；首次撒旦選了最適當的時候下手,須要等待時機.第二次撒旦從神拿到權柄,便立刻行動.雖手法不同,表明惡毒的心卻是相同；使約伯痛苦之中,加上難以忍受身體的痛楚.
\newline
\newline
平常人在苦難中,灰心失志,只要略加壓力即可完全崩潰.約伯受極大的痛苦,處身一無所有,心靈極度憂傷,充滿疑問之時；撒旦乘機用苛疾加在約伯身上.撒旦利用人禍,天災；動用自然界的現象,以達其目的來破壞人的信心；用疾病對付世上的人,特別是神的兒女.
\newline
\newline
約伯患了極其痛苦古怪的病,從腳掌到頭頂長毒瘡,體無完膚.(二7)本來他身體健康,人都喜歡親近他,現在容貌全非,「坐在爐灰中拿瓦片括身體,」(8)因為毒瘡不斷流膿水,極度痕癢.「肉體以蟲子和塵土為衣......皮膚破裂.」(七5)因病使他呼吸困難,(九18)日不能寐,夜發惡夢；(七4-14)以致口臭牙爛,形成皮包骨.(十九17-20)除了皮外疼痛,甚至痛入骨頭和神經；日夜不安寧,因此精神大受打擊.實在是一般人所難忍受!約伯第一次受苦是頃刻間,頓成一無所有的打擊；第二次是長期不治痛楚難當之症,健康喪失.真是貧病交加!
\newline
\newline
約伯未受苦之先,所受的恩,所蒙的福,對他都是一種考驗.因他有正確的人生觀,富足和平安未嘗奪去他純正的心；當他喪失一切,他的心也毫不動搖.他勝過了考驗,且成了他日後遭受苦難考驗的巨大力量；他受得起富足的考驗,才能抵擋貧窮的考驗.
\newline
\newline
基督徒在地上,時常會遇到考驗；撒旦可能用各種方法來試探,使我們在平安豐富之時,失落純正的信心.不知不覺中了撒旦詭計,如在電子瓦煲裡一樣；溫暖舒適,豈料歷時既久,熟爛溶化.內地出來的弟兄姊妹們!不要忘記你們從前在苦難中有美好的見證,不要以為現在到了自由舒適之地就可以放鬆了；應當小心謹慎,提防撒旦,在你不知不覺中破壞你的見證.
\newline
\newline
約伯第二次被撒旦攻擊,除身體受苦外,加上心靈的痛苦；聖經告訴我們,他因貧病,被迫住在城外.遠離人群,獨坐爐灰中,平時的威風盡失；他的怪病被視為大痲瘋患者,與親朋隔絕；下流之輩來戲弄他.約伯本是個有信心祈禱的人,素常為兒女為家庭禱告,蒙神垂聽；如今他的禱告似付諸東流.
\newline
\newline
有的人禱告無效便灰心,有的人雖不灰心,但卻抱著「無所謂」的態度；神垂聽也好,不垂聽也罷.
\newline
\newline
約伯是個篤信獨一真神的人,他信神是慈愛,公義,全能,真理的神,但現在發生在他身上的,卻是矛盾的事實.從他話中透露他心裡疑雲密佈,黑影重重,這是他心靈最大的痛苦.在此打擊之下,他的妻子對他說:「你仍然持守你的純正麼?」
\newline
\newline
我以為約伯之妻,是個愛主而有屬靈深度的姊妹,她是約伯的好妻子；當約伯受第一次打擊,財富和兒女全失,可以說她受的打擊過於約伯所受的.一般而論,為牧者疼惜兒女勝過為父者；女人對財物關係之深,也過於男人；約伯受了雙重的打擊而站穩,他妻子亦依然站穩,這是一般姊妹望塵莫及的.他們十個兒女都循規蹈矩,雖然歸功於約伯天天為他們祈禱；但約伯妻子之教導有方也功不可沒,他們夫婦同心合意,按照屬靈原則教導培養兒女；這位姊妹真是位賢良配偶；所以約伯平時享有溫暖蒙恩的家庭,使約伯第一次受苦難時,得到極大的安慰和扶持.
\newline
\newline
彼得前書說姊妹是軟弱的器皿.當第二次苦難發生於約伯本身時,他的妻子卻軟弱了；因她愛丈夫勝過所有的財物,她覺得神不合理；她見丈夫患不治頑疾,難渡餘年；她因愛心過切,致顯出屬靈的軟弱,對約伯說:「你棄掉神,死了罷.」這是撒旦的本意,牠第一次對神說,你伸手剝奪約伯所有的,他就當面棄掉你.第二次又說,你伸手加在他身上,他就當面棄掉你.撒旦最大目的,要屬神的人棄掉神.撒旦於約伯未能達到目的,而冀利用約伯妻子可達到目的.「純正」是神誇獎約伯屬靈的美德,是約伯尊貴之處,是他有生以來的見證；現在妻子要他變節不再持守純正,要他放棄他的見證,並且叫他「死了罷」!因她不忍見丈夫受苦.
\newline
\newline
親愛的弟兄姊妹!約伯夫妻實在恩愛!夫妻相愛本是好的；但是愛心若不建立在真理上,便成為撒旦最厲害攻擊的工具.基督化家庭乃是夫婦共處,尊主為大,受主管理.
\newline
\newline
約伯在危險關頭反責備他的妻子,出言愚頑的婦人；由此使我們想起主耶穌責備彼得的事.當主對彼得說,祂將要上耶路撒冷,受人凌辱,被釘死在十字架上.彼得說:「主啊!千萬不可如此.」主責備他說:「撒旦退我後面去罷.」約伯正是這樣責備他的妻子,他對神的純正仍然不變.他說:「難道我們從神手裡得福,不也受禍麼.」這話雖不甚正確,神的手是施恩的手,非降禍的；不過他說話的精神是對的.他認為禍也是神降下的,他的觀念乃是一切的事都出於神,都有神的主權在其中；所以他對神的心,永遠謙卑順服.
\newline
\newline
我們在主裡,很多時候失敗,是因對神的尊敬不夠；常用環境來解釋神,很難以神來解釋環境.用環境來解釋神,就是不信,把神的榮耀,尊貴降低於地.以神來解釋環境,就是信心的表現；雖不瞭解環境,但深信神必有美意,信心提高至超越地步.
\newline
\newline
約伯有美好的信心,他不憑眼見,乃憑信心；因此他站立得住,信心毫不動搖.
\newline
\newline
約伯何其認識神,愛慕神.(羅八35-39)論及兩方面的愛,先是神對人的愛,後是人對神的愛.神對人的愛是超越一切的愛,是任何力量所不能隔絕的；「無論是死,是生,是天使,是掌權的；是有能的,是現在的事,是將來的事,是高處的,是低處的,是別的受造之物,都不能叫我們與神的愛隔絕,這愛是在我們的主基督耶穌裡的.」因這毫無攔阻的愛臨到人,所以我們可以誇勝地說:「誰能使我們與基督的愛隔絕呢?難道是患難麼,是困苦麼,是逼迫麼,是飢餓麼,是赤身露體麼,是刀劍麼......然而靠著愛我們的主,在這一切的事上,已經得勝有餘了.」我們所以得勝,不是因為我們愛主,乃是靠著主的愛.凡此都可應用在約伯身上,他終於靠主得勝了,他的得勝使其真相大白.神容許撒旦攻擊約伯,有祂的美意,有多方面之目的.
\newline
\newline
當約伯喪失一切時,他的三個朋友從遠方而來.約伯三個朋友都是地方領袖,有學問,有道德,可以說是約伯患難之交,良朋密友.他們認為約伯此次痛苦太大,個人難以相助,所以約同而來；他們憑真誠的愛心,目的為他分憂,與哀哭的人同哭.安慰他,不但對他深表同情,並且盼望他受傷的心靈得著鼓勵；他們一見約伯就放聲大哭,各人撕裂外袍,把塵土向天揚起來,落在自己的頭上.他們就同他七天七夜,坐在地上,一言不發.他們目的要解除他的痛苦,希望能把他痛苦的根源除掉；他們好像醫生,抱著治病救人的宗旨輪流會診；雖醫術欠高明,但熱誠精神可嘉.由第三章至卅一章都是他們長篇談話記錄.以利法說了三次話,比勒達也說了三次話,瑣法說了兩次話,約伯說了十次話.約伯長篇的談話,內容大都和他三個朋友不同；雖然討論的題目都是「約伯為何受苦」可是他們的方法和態度不同.約伯之所以言語多,是因切身的痛苦多.約伯自從第二次受苦,對他妻子責備一番之後,一直沒發言,現在和他三朋友長篇大論,可見他的心靈逐漸甦醒；他希望從痛苦中求解脫,希望神給予啟示得以解釋,約伯談話之中,開始向神說話.第七章1-6,約伯說到人生的悲慘,用各樣事物來衡量.由第七節開始至末了,他向神說話,藉著和三朋友談話,他向神談論,希望從中得到至清楚的最後答案.約伯十次的說話,和他三朋友共八次的說話,讀之甚難明白,好像前後脫節,各人的談話都主觀表達個人的學識和見解,各人有其不同的根據,詞句優美.
\newline
\newline
第三章可以說是約伯的哀歌,是受苦人心情的流露.
\newline
\newline


\newpage

\section{第52屆~港九培靈會~焦源濂~第五講　失敗的安慰者-以利法(二)}
\label{sec:623}
\textbf{第五講　失敗的安慰者-以利法(二)}
\newline
\newline
連結: \href{https://www.hkbibleconference.org/session-message/view/623}{\texttt{https://www.hkbibleconference.org/session-message/view/623}}
\newline
\newline
\hyperref[sec:621]{< < < PREV SERMON < < <}
~
\hyperref[sec:toc]{[返目錄]}
~
\hyperref[sec:624]{> > > NEXT SERMON > > >}
\newline
\newline
約伯記 8:1-8:22
\newline
\begin{longtable}{cl}
\hline
\hline
章節 & 經文 (和合本修訂版)\\
\hline
8:1 & \begin{tabularx}{0.7\textwidth}{X} 書亞人比勒達回答說： \end{tabularx} \\ \\ \relax
8:2 & \begin{tabularx}{0.7\textwidth}{X} 「這些話你要說到幾時？你口中的言語如狂風要到幾時呢？ \end{tabularx} \\ \\ \relax
8:3 & \begin{tabularx}{0.7\textwidth}{X} 神豈能偏離公平？全能者豈能偏離公義？ \end{tabularx} \\ \\ \relax
8:4 & \begin{tabularx}{0.7\textwidth}{X} 或者你的兒女得罪了他，他就把他們交在過犯的掌控中。 \end{tabularx} \\ \\ \relax
8:5 & \begin{tabularx}{0.7\textwidth}{X} 你若切切尋求神，向全能者懇求； \end{tabularx} \\ \\ \relax
8:6 & \begin{tabularx}{0.7\textwidth}{X} 你若純潔正直，他必定為你興起，使你公義的居所興旺。 \end{tabularx} \\ \\ \relax
8:7 & \begin{tabularx}{0.7\textwidth}{X} 你起初雖然微小，日後必非常強盛。 \end{tabularx} \\ \\ \relax
8:8 & \begin{tabularx}{0.7\textwidth}{X} 「請你詢問上代，思念他們祖先所查究的。 \end{tabularx} \\ \\ \relax
8:9 & \begin{tabularx}{0.7\textwidth}{X} 我們不過從昨日才有，一無所知，因我們在世的日子好像影子。 \end{tabularx} \\ \\ \relax
8:10 & \begin{tabularx}{0.7\textwidth}{X} 他們豈不指教你，告訴你，說出發自內心的言語呢？ \end{tabularx} \\ \\ \relax
8:11 & \begin{tabularx}{0.7\textwidth}{X} 「蒲草沒有泥豈能生長？蘆荻沒有水豈能長大？ \end{tabularx} \\ \\ \relax
8:12 & \begin{tabularx}{0.7\textwidth}{X} 它還青翠，沒有割下的時候，比百樣的草先枯槁。 \end{tabularx} \\ \\ \relax
8:13 & \begin{tabularx}{0.7\textwidth}{X} 凡忘記神的人，路途也是這樣；不虔敬人的指望要滅沒。 \end{tabularx} \\ \\ \relax
8:14 & \begin{tabularx}{0.7\textwidth}{X} 他所仰賴的必折斷，他所倚靠的是蜘蛛網。 \end{tabularx} \\ \\ \relax
8:15 & \begin{tabularx}{0.7\textwidth}{X} 他要倚靠房屋，房屋卻站立不住；他要抓住房屋，房屋卻不能存留。 \end{tabularx} \\ \\ \relax
8:16 & \begin{tabularx}{0.7\textwidth}{X} 他在日光之下茂盛，嫩枝在園中蔓延； \end{tabularx} \\ \\ \relax
8:17 & \begin{tabularx}{0.7\textwidth}{X} 他的根盤繞石堆，鑽入石縫。 \end{tabularx} \\ \\ \relax
8:18 & \begin{tabularx}{0.7\textwidth}{X} 他若從本地被拔出，那地就不認識他，說：『我沒有見過你。』 \end{tabularx} \\ \\ \relax
8:19 & \begin{tabularx}{0.7\textwidth}{X} 看哪，這就是他道路中的喜樂，以後必另有人從塵土而生。 \end{tabularx} \\ \\ \relax
8:20 & \begin{tabularx}{0.7\textwidth}{X} 看哪，神必不丟棄完全人，也不扶助邪惡人的手。 \end{tabularx} \\ \\ \relax
8:21 & \begin{tabularx}{0.7\textwidth}{X} 他還要以喜笑充滿你的口，以歡呼充滿你的嘴唇。 \end{tabularx} \\ \\ \relax
8:22 & \begin{tabularx}{0.7\textwidth}{X} 恨惡你的要披戴羞愧，惡人的帳棚必歸於無有。」 \end{tabularx} \\ \\
[1ex]
\hline
\hline
\end{longtable}
約伯記清楚告訴我們,義人在世上是要顯明神的榮耀,義人是神所創造萬物之中最奇妙的創造,雖然軟弱,卻能抵擋撒旦的攻擊；在苦難中,對神仍然有無偽的信心,有純潔的愛心.義人受苦是要顯明人尊貴的價值,是人超過其他萬物的明證.神信任約伯,所以當撒旦千方百計要攻擊約伯時,神容許撒旦所為；因為神知道約伯必然會站穩.
\newline
\newline
弟兄姊妹!我們若不是因犯罪而受苦,就應當感謝神；因為神看重你,視你如珠如寶為萬物中最奇妙的創造.雖然世人鄙視我們,神卻把我們高舉起來,使祂的大仇敵撒旦閉口無言.人算甚麼,神竟如此高舉我們!
\newline
\newline
約伯第二次所受的,確是極大的痛苦,他幾乎忍無可忍；但神知道他必然站住.當他一切都失落時；神差遣三個屬神的人來安慰他,幫助他,陪伴他.三人與約伯無言同坐了七晝夜之後；約伯開始說話,他怨嘆自己,痛不欲生.似乎他三個朋友的安慰,所收是適得相反的效果；其實卻是表明他心靈逐漸甦醒的好現象；他開始感嘆人生遭遇,尋找痛苦的原因和答案.
\newline
\newline
約伯和三朋友的談話,共有卅多章之長,未免令人讀之乏味；有人如走馬看花看了就算了；但是在神眼中,這些談話是神看為寶貴的,所以記在聖經上.瑪拉基書三章16節說:「那時敬畏耶和華的彼此談論,耶和華側耳而聽,且有記念冊在祂面前,記錄那敬畏耶和華思念祂名的人.」
\newline
\newline
約伯三個朋友,安慰和幫助他們受苦的弟兄,探討屬靈的奧秘；雖然他們有錯,但存心良善,談論內容十分重要.神看重基督徒在主內的交通.談論開始,約伯唱了哀歌.
\newline
\newline
第三章是約伯的哀歌,可分為兩段:
\newline
\newline
(一)約伯咒詛自己的生日(1-10)此段經文共有十個願字,咒詛的願望共有十次.他所以咒詛的原因,他說:「因沒有把懷我胎的門關閉,也沒有將患難對我的眼隱藏.」
\newline
\newline
在舊約聖經中,除了約伯之外,先知耶利米,也因受苦而發出荒謬的反應,他甚至連給他父親報信的人,他都咒詛.
\newline
\newline
耶利米是先知,約伯是個屬靈的領袖,他們並非普通信徒；但是因受苦竟然發生奇怪的思想.這要提醒我們,當人受苦時,思想非常荒謬；只想到自己,不想到別人.撒旦叫人在受苦時專想自己,思想狹窄,想來想去想不開,越想自己,越沒出路；怨天尤人,滿懷憎恨見別人平安無事,巴不得人人與自己共受苦難；這與神的心意完全相反,神要藉著苦難寬廣人的心懷意念.大衛說:「我未曾受苦之先走迷了路,現在轉回跟隨你的律法.我受苦與我有益；為要使我學習你的律例.」人若在受苦之中,放下自己,到神前多多尋求神；開自己的眼睛,看見別人的需要；這樣可以減少自己的痛苦,而且可以蒙神賜福.詩篇第四篇一節是大衛獻給受苦人的經歷,人受困於苦難中；如果到神面前,神可帶你到寬廣之地,使你心胸擴大.
\newline
\newline
(二)約伯咒詛自己的生命(11-24)那時約伯的語氣轉變,連聲說:「為何,」「為何不出母胎而死；為何不出母腹絕氣；為何有膝接收我；為何有奶哺養我.」他認為生不如死.他說:「受患難的人,為何有光賜給他呢?心中愁苦的人,為何有生命賜給他呢?」他以為生而受苦,何必生於人世.生命應是享受,希望,喜樂,幸福,是光輝燦爛的；而他卻是受苦.弟兄姊妹!你有這樣的思想嗎?你會不會想到約伯在順境時何等美滿,至此地步何等凄涼.其實,這時才是約伯一生最有意義的日子.
\newline
\newline
常人在受苦時,就顯出思想膚淺,在苦難中好像被捆綁；完全失去自由,懷才不遇,自覺人生毫無意義.
\newline
\newline
雅各十二個兒子中最出色的是約瑟,他有崇高的理想,有美好的靈性,大有才能；他得了長子名份,即雙倍的祝福；可是他曾被賣為奴,被下監牢,全身被鎖鍊捆綁在黑暗之中,寸步難行.一個有為青年經歷諸般試煉,終於曉得忠心事奉神.學習如何管理周圍的環境,先受環境的影響,然後藉著經歷去影響環境；他到了埃及以後和在自己家裡,判若二人.他到波提乏的家,耶和華與他同在,於是凡事順利；波提乏把一切交給他管理.他被下監,雖全失自由；但獄中氣氛因而改變,他和酒政,膳長並其他的政要人物同囚,從中學習有關政治方面的事；十三年歲月,看來似乎浪費光陰,事實上大有收穫；一切的學習和操練,奠定了將來成就的基礎,他當了宰相管理全埃及.
\newline
\newline
種子埋在地下的時期,泥土表面一無所有,看起來令人乏味；但到了發芽,生長,開花,結果子,漸漸給人覺得有意思.其實埋在地下那段期間,才是最有意義的；雖不見天日,但往下扎根,才能向上結果.
\newline
\newline
求神讓我們曉得,若不因犯罪而受苦難,是人生最重要的關頭,是最有意義的,是更大榮耀,更尊貴的生命,將要顯明出來的開始.
\newline
\newline
約伯咒詛生命,他情願死.關於死亡,是約伯記重要討論的題目.約伯對死亡的觀念如何?請看(三13-19)約伯不認為死亡結束人生一切,如同消滅一樣；人有靈魂,肉身死亡後另有一個世界；但是他對死亡也有很多不正確的觀念.他沒有天堂和地獄的觀念,他以為人死後,無論君主,謀士,財主,王子,懷孕未到期而落的胎,嬰孩,惡人,困乏人,被囚的人,督工,大的,小的,奴僕,主人都在一起,都是死人階級,和平共處,永遠安息.那時約伯不得安息,日夜唉哼.一個多苦者的思想錯誤,片面,主觀,幻想；根據自身想像別人的問題,當某種觀念出現時,自以為完全正確,任何人無法加以攔阻.
\newline
\newline
人的性命在神手中,如果神給一個愛主者或犯罪者的時間未完而想離開世界,不一定是件幸福的事.人遇苦難時常羨慕死亡,但內心不一定願意死,這是人之常情.約伯形容求死勝於求隱藏的珍寶,但他卻不自殺；那不過是他心情的流露,痛苦的表達.
\newline
\newline
約伯開口說話,他三個朋友中最年長的.猶太人習俗與我國略同,尊重長者.以利法說:「年老的當先說話,受膏的當以智慧教訓人.」(三十二7)「我們這裡有白髮的,和年紀老邁的,比你父親還老.」(十五10)這話十分可能表明他自己；他勸導約伯:「若不聽老人言,吃苦在眼前.」
\newline
\newline
以利法是約伯三個朋友中經歷最多的.從他話中發現他多用自己的經歷,作勸導約伯的根據；他有屬靈的經歷,帶有啟示；也有人生的經歷,閱歷豐富,對人生百態有仔細研究的態度；見識多,瞭解人生.解經者都認為他是個經驗派大師.
\newline
\newline
以利法也富有同情心,言語溫柔.(四3-5)他開始說話時略帶責備的語氣；但他不過是提醒約伯,不可忘記原是屬靈的人,以往曾經幫助過許多受苦者；現在應當引用那些勸人之道自勉.
\newline
\newline
彼勒達語氣完全不同(八3-4)他一開口,就大膽假設約伯有罪,說神是公義的；或者是約伯的兒女得罪神,以致受報應；使約伯聽之,有口難言.
\newline
\newline
瑣法的話更加厲害(十一3-6)他說神追討約伯的,比他罪孽該得的還少.
\newline
\newline
以利法話語的大前提是種瓜得瓜,種豆得豆,犯罪受罪.他找不到約伯犯何罪；但認定他受苦乃因犯罪.他說出個人閱歷(四1-8)說出自己屬靈的經歷和啟示.(五3)第四章(12-21)所說的,有人認為以利法當時所見的可能是魔鬼,所以使他恐懼,戰兢；但我認為他非屬普通信徒,且舊約有載,許多先知見異象都有懼怕的反應.當摩西見神在荊棘火焰中向他顯現,他不敢向前.大大蒙神眷愛的但以理,當神的使者向他顯現,他仆倒在地,三個七日似乎身體軟弱.以利法見到神的尊榮聖潔,深覺自己有罪；好像以賽亞覺得自己卑賤污穢而戰兢恐懼.我個人對以利法所說的經歷覺得無所懷疑；不過該異象是否適用於約伯情形之下呢?異象是以利法的,他濫用屬靈經歷,且擴大其解釋,片面看神是孤獨乖僻,人難以相親.(18)不錯,神是至尊無比,是獨一真神；但是神愛世人,與心靈痛悔的人同居,祂名為以馬內利.神信任約伯,所以讓魔鬼攻擊他.
\newline
\newline
弟兄姊妹!我們在約伯記,看見個屬靈長者,他有寶貴的經歷和啟示；但要小心!恐怕他看自己的啟示,為獨一的啟示,可以解決一切的問題.年長的雖經歷多,卻不一定是全面的經歷,不能以老賣老.教會中傳道人切勿以為自己所說的話清楚包括一切.多謝主!祂在地上醫治人,不是單用一個方法,乃是用各種各樣的方法.講道和經應相輔而行,道理有助知識；經歷使人真正得幫助；但萬勿強調經歷過於主自己,不願謙卑順服聖靈的引導,若然,其經歷非但無濟於事,反易於破壞人.
\newline
\newline
約伯最佩服三朋友中之以利法,所以默然受他責備；但若是當日約伯佩服,聽從這位有經歷的人而妄自認罪在神前,那就極為危險!今日許多所謂有經歷者,心有餘而屬靈深度不足；教會諸多紛亂都出於此.有的教會死氣沉沉,有的亂七八糟,有經歷的人太多,各執自己的經歷是唯一的,是至高最完全的,蔑視他人.
\newline
\newline
求神教我們,越屬靈,越謙卑,越小心,越客觀.
\newline
\newline


\newpage

\section{第52屆~港九培靈會~焦源濂~第六講　失敗的安慰者-比勒達}
\label{sec:624}
\textbf{第六講　失敗的安慰者-比勒達}
\newline
\newline
連結: \href{https://www.hkbibleconference.org/session-message/view/624}{\texttt{https://www.hkbibleconference.org/session-message/view/624}}
\newline
\newline
\hyperref[sec:623]{< < < PREV SERMON < < <}
~
\hyperref[sec:toc]{[返目錄]}
~
\hyperref[sec:626]{> > > NEXT SERMON > > >}
\newline
\newline
約伯記 11:1-31:20
\newline
\begin{longtable}{cl}
\hline
\hline
章節 & 經文 (和合本修訂版)\\
\hline
11:1 & \begin{tabularx}{0.7\textwidth}{X} 拿瑪人瑣法回答說： \end{tabularx} \\ \\ \relax
11:2 & \begin{tabularx}{0.7\textwidth}{X} 「這許多的話豈不該回答嗎？多嘴多舌的人豈可成為義呢？ \end{tabularx} \\ \\ \relax
11:3 & \begin{tabularx}{0.7\textwidth}{X} 你誇大的話豈能使人不作聲嗎？你戲笑的時候豈沒有人使你受辱嗎？ \end{tabularx} \\ \\ \relax
11:4 & \begin{tabularx}{0.7\textwidth}{X} 你說：『我的教導純全，我在你眼前是清潔的。』 \end{tabularx} \\ \\ \relax
11:5 & \begin{tabularx}{0.7\textwidth}{X} 但是，惟願神說話，願他張開嘴唇攻擊你。 \end{tabularx} \\ \\ \relax
11:6 & \begin{tabularx}{0.7\textwidth}{X} 願他將智慧的奧祕指示你，因為健全的知識是兩面的。你當知道，神使你忘記你的一些罪孽。 \end{tabularx} \\ \\ \relax
11:7 & \begin{tabularx}{0.7\textwidth}{X} 你能尋見神的奧祕嗎？你能尋見全能者的極限嗎？ \end{tabularx} \\ \\ \relax
11:8 & \begin{tabularx}{0.7\textwidth}{X} 高如諸天，你能做甚麼？比陰間深，你能知道甚麼？ \end{tabularx} \\ \\ \relax
11:9 & \begin{tabularx}{0.7\textwidth}{X} 其量度比地長，比海更寬。 \end{tabularx} \\ \\ \relax
11:10 & \begin{tabularx}{0.7\textwidth}{X} 他若經過，把人拘禁，召集會眾，誰能阻擋他呢？ \end{tabularx} \\ \\ \relax
11:11 & \begin{tabularx}{0.7\textwidth}{X} 因為他知道虛妄的人；當他看見罪惡，豈不留意嗎？ \end{tabularx} \\ \\ \relax
11:12 & \begin{tabularx}{0.7\textwidth}{X} 空虛的人若獲得知識，野驢生下的駒子也成了人。 \end{tabularx} \\ \\ \relax
11:13 & \begin{tabularx}{0.7\textwidth}{X} 「至於你，若堅固己心，又向主舉手； \end{tabularx} \\ \\ \relax
11:14 & \begin{tabularx}{0.7\textwidth}{X} 你若遠遠脫離你手中的罪孽，不容許不義住在你帳棚之中； \end{tabularx} \\ \\ \relax
11:15 & \begin{tabularx}{0.7\textwidth}{X} 這樣，你必仰起臉來，毫無瑕疵；你也必安穩，無所懼怕。 \end{tabularx} \\ \\ \relax
11:16 & \begin{tabularx}{0.7\textwidth}{X} 你必忘記你的苦楚，就是想起來，也如流過的水。 \end{tabularx} \\ \\ \relax
11:17 & \begin{tabularx}{0.7\textwidth}{X} 你在世要升高，比正午更明，雖有黑暗，仍像早晨。 \end{tabularx} \\ \\ \relax
11:18 & \begin{tabularx}{0.7\textwidth}{X} 你因有指望就必穩固，也必四圍察看，安然躺下。 \end{tabularx} \\ \\ \relax
11:19 & \begin{tabularx}{0.7\textwidth}{X} 你躺臥，無人驚嚇，並有許多人向你求恩。 \end{tabularx} \\ \\ \relax
11:20 & \begin{tabularx}{0.7\textwidth}{X} 但惡人的眼睛要失明；他們無路可逃，他們的指望就是氣絕身亡。」 \end{tabularx} \\ \\ \relax
12:1 & \begin{tabularx}{0.7\textwidth}{X} 約伯回答說： \end{tabularx} \\ \\ \relax
12:2 & \begin{tabularx}{0.7\textwidth}{X} 「你們果真是人物啊！智慧要與你們一同去死。 \end{tabularx} \\ \\ \relax
12:3 & \begin{tabularx}{0.7\textwidth}{X} 但我也有聰明，跟你們一樣，並非不及你們。這些事，誰不知道呢？ \end{tabularx} \\ \\ \relax
12:4 & \begin{tabularx}{0.7\textwidth}{X} 我這求告神、蒙他應允的人竟成了朋友所譏笑的；又公義又完全的人竟遭受譏笑。 \end{tabularx} \\ \\ \relax
12:5 & \begin{tabularx}{0.7\textwidth}{X} 安逸的人心裡藐視災禍，這災禍在等待失足滑跌的人。 \end{tabularx} \\ \\ \relax
12:6 & \begin{tabularx}{0.7\textwidth}{X} 強盜的帳棚安寧，惹神發怒的人穩固，他們把神握在自己手中。 \end{tabularx} \\ \\ \relax
12:7 & \begin{tabularx}{0.7\textwidth}{X} 「你問走獸，走獸必指教你；你問空中的飛鳥，飛鳥必告訴你； \end{tabularx} \\ \\ \relax
12:8 & \begin{tabularx}{0.7\textwidth}{X} 或者你與地說話，地必指教你；海中的魚也必向你說明。 \end{tabularx} \\ \\ \relax
12:9 & \begin{tabularx}{0.7\textwidth}{X} 在這一切當中，有誰不知道這是耶和華的手做成的呢？ \end{tabularx} \\ \\ \relax
12:10 & \begin{tabularx}{0.7\textwidth}{X} 凡動物的生命和人類的氣息都在他手中。 \end{tabularx} \\ \\ \relax
12:11 & \begin{tabularx}{0.7\textwidth}{X} 耳朵豈不辨別言語，正如上膛品嘗食物嗎？ \end{tabularx} \\ \\ \relax
12:12 & \begin{tabularx}{0.7\textwidth}{X} 年老的有智慧，壽高的有知識。 \end{tabularx} \\ \\ \relax
12:13 & \begin{tabularx}{0.7\textwidth}{X} 「在神有智慧和能力，他有謀略和知識。 \end{tabularx} \\ \\ \relax
12:14 & \begin{tabularx}{0.7\textwidth}{X} 看哪，他拆毀，就不能重建；他拘禁人，人就不得釋放。 \end{tabularx} \\ \\ \relax
12:15 & \begin{tabularx}{0.7\textwidth}{X} 看哪，他使水止住，水就乾了；他把水放出，水就淹沒大地。 \end{tabularx} \\ \\ \relax
12:16 & \begin{tabularx}{0.7\textwidth}{X} 在他有能力和智慧，走迷的和使人迷路的都屬他。 \end{tabularx} \\ \\ \relax
12:17 & \begin{tabularx}{0.7\textwidth}{X} 他把謀士剝衣擄去，使審判官變為愚妄。 \end{tabularx} \\ \\ \relax
12:18 & \begin{tabularx}{0.7\textwidth}{X} 他解除君王的權勢，用帶子捆住他們的腰。 \end{tabularx} \\ \\ \relax
12:19 & \begin{tabularx}{0.7\textwidth}{X} 他把祭司剝衣擄去，使有權能的人傾覆。 \end{tabularx} \\ \\ \relax
12:20 & \begin{tabularx}{0.7\textwidth}{X} 他廢去忠信者的言論，奪去長者的見識。 \end{tabularx} \\ \\ \relax
12:21 & \begin{tabularx}{0.7\textwidth}{X} 他使貴族蒙羞受辱，放鬆勇士的腰帶。 \end{tabularx} \\ \\ \relax
12:22 & \begin{tabularx}{0.7\textwidth}{X} 他從黑暗中彰顯深奧的事，使死蔭顯出光明。 \end{tabularx} \\ \\ \relax
12:23 & \begin{tabularx}{0.7\textwidth}{X} 他使邦國興旺而又毀滅，使邦國擴展又被掠奪。 \end{tabularx} \\ \\ \relax
12:24 & \begin{tabularx}{0.7\textwidth}{X} 他將地上百姓中領袖的聰明奪去，使他們迷失在荒涼無路之地。 \end{tabularx} \\ \\ \relax
12:25 & \begin{tabularx}{0.7\textwidth}{X} 他們在無光的黑暗中摸索；他使他們搖晃像醉酒的人一樣。」 \end{tabularx} \\ \\ \relax
13:1 & \begin{tabularx}{0.7\textwidth}{X} 「看哪，這一切，我眼都見過；我耳都聽過，而且明白。 \end{tabularx} \\ \\ \relax
13:2 & \begin{tabularx}{0.7\textwidth}{X} 你們所知道的，我也知道，並非不及你們。 \end{tabularx} \\ \\ \relax
13:3 & \begin{tabularx}{0.7\textwidth}{X} 然而我要對全能者說話，我願與神理論。 \end{tabularx} \\ \\ \relax
13:4 & \begin{tabularx}{0.7\textwidth}{X} 但你們是編造謊言的，全都是無用的醫生。 \end{tabularx} \\ \\ \relax
13:5 & \begin{tabularx}{0.7\textwidth}{X} 惟願你們全然不作聲，這就是你們的智慧！ \end{tabularx} \\ \\ \relax
13:6 & \begin{tabularx}{0.7\textwidth}{X} 請你們聽我的答辯，留心聽我嘴唇的訴求。 \end{tabularx} \\ \\ \relax
13:7 & \begin{tabularx}{0.7\textwidth}{X} 你們要為神說不義的話嗎？要為他說詭詐的言語嗎？ \end{tabularx} \\ \\ \relax
13:8 & \begin{tabularx}{0.7\textwidth}{X} 你們要看神的情面嗎？要為他爭辯嗎？ \end{tabularx} \\ \\ \relax
13:9 & \begin{tabularx}{0.7\textwidth}{X} 他查究你們，這豈是好事嗎？人欺騙人，你們也要照樣欺騙他嗎？ \end{tabularx} \\ \\ \relax
13:10 & \begin{tabularx}{0.7\textwidth}{X} 你們若暗中看人的情面，他必定要責備你們。 \end{tabularx} \\ \\ \relax
13:11 & \begin{tabularx}{0.7\textwidth}{X} 他的尊榮豈不叫你們懼怕嗎？他豈不使驚嚇臨到你們嗎？ \end{tabularx} \\ \\ \relax
13:12 & \begin{tabularx}{0.7\textwidth}{X} 你們可記念的諺語是灰燼的箴言；你們的後盾是泥土的後盾。 \end{tabularx} \\ \\ \relax
13:13 & \begin{tabularx}{0.7\textwidth}{X} 「你們不要向我作聲，讓我說話，無論如何我都承當。 \end{tabularx} \\ \\ \relax
13:14 & \begin{tabularx}{0.7\textwidth}{X} 我為何把我的肉掛在我的牙上，將我的命放在我的手掌中呢？ \end{tabularx} \\ \\ \relax
13:15 & \begin{tabularx}{0.7\textwidth}{X} 看哪，他要殺我，我毫無指望，然而我還要在他面前辯明我所行的。 \end{tabularx} \\ \\ \relax
13:16 & \begin{tabularx}{0.7\textwidth}{X} 這要成為我的拯救，因為不虔誠的人不可到他面前。 \end{tabularx} \\ \\ \relax
13:17 & \begin{tabularx}{0.7\textwidth}{X} 你們要細聽我的言語，讓我的申辯入你們耳中。 \end{tabularx} \\ \\ \relax
13:18 & \begin{tabularx}{0.7\textwidth}{X} 看哪，我已陳明我的案，知道自己有義。 \end{tabularx} \\ \\ \relax
13:19 & \begin{tabularx}{0.7\textwidth}{X} 還有誰要和我爭辯，我現在就緘默不言，氣絕而死。 \end{tabularx} \\ \\ \relax
13:20 & \begin{tabularx}{0.7\textwidth}{X} 惟有兩件事不要向我施行，我就不躲開你的面： \end{tabularx} \\ \\ \relax
13:21 & \begin{tabularx}{0.7\textwidth}{X} 就是把你的手縮回，遠離我身；又不使你的威嚴恐嚇我。 \end{tabularx} \\ \\ \relax
13:22 & \begin{tabularx}{0.7\textwidth}{X} 這樣，你呼叫，我就回答；或是讓我說話，你回答我。 \end{tabularx} \\ \\ \relax
13:23 & \begin{tabularx}{0.7\textwidth}{X} 我的罪孽和我的罪有多少呢？求你叫我知道我的過犯與我的罪。 \end{tabularx} \\ \\ \relax
13:24 & \begin{tabularx}{0.7\textwidth}{X} 你為何轉臉，拿我當仇敵呢？ \end{tabularx} \\ \\ \relax
13:25 & \begin{tabularx}{0.7\textwidth}{X} 你要驚動被風吹的葉子嗎？要追趕枯乾的碎秸嗎？ \end{tabularx} \\ \\ \relax
13:26 & \begin{tabularx}{0.7\textwidth}{X} 你寫下苦楚對付我，又使我擔當幼年的罪孽。 \end{tabularx} \\ \\ \relax
13:27 & \begin{tabularx}{0.7\textwidth}{X} 你把我的腳鎖上木枷，察看我一切的道路，為我的腳掌劃定界限。 \end{tabularx} \\ \\ \relax
13:28 & \begin{tabularx}{0.7\textwidth}{X} 人像滅絕的爛物，像蟲蛀的衣裳。」 \end{tabularx} \\ \\ \relax
14:1 & \begin{tabularx}{0.7\textwidth}{X} 「人為婦人所生，日子短少，多有患難。 \end{tabularx} \\ \\ \relax
14:2 & \begin{tabularx}{0.7\textwidth}{X} 他出來如花，凋謝而去；他飛逝如影，不能存留。 \end{tabularx} \\ \\ \relax
14:3 & \begin{tabularx}{0.7\textwidth}{X} 這樣的人你豈會睜眼看他，又叫我來，在你那裡受審嗎？ \end{tabularx} \\ \\ \relax
14:4 & \begin{tabularx}{0.7\textwidth}{X} 誰能使潔淨出於污穢呢？誰也不能！ \end{tabularx} \\ \\ \relax
14:5 & \begin{tabularx}{0.7\textwidth}{X} 既然人的日子限定，他的月數在於你，你劃定他的界限，他不能越過； \end{tabularx} \\ \\ \relax
14:6 & \begin{tabularx}{0.7\textwidth}{X} 求你轉眼不看他，使他得歇息，直到他像雇工享受他的一天。 \end{tabularx} \\ \\ \relax
14:7 & \begin{tabularx}{0.7\textwidth}{X} 「因樹有指望，若被砍下，還可發芽，嫩枝生長不息。 \end{tabularx} \\ \\ \relax
14:8 & \begin{tabularx}{0.7\textwidth}{X} 樹根若衰老在地裡，樹幹也死在土中， \end{tabularx} \\ \\ \relax
14:9 & \begin{tabularx}{0.7\textwidth}{X} 及至得了水氣，還會發芽，長出枝條，像新栽的樹一樣。 \end{tabularx} \\ \\ \relax
14:10 & \begin{tabularx}{0.7\textwidth}{X} 但壯士一死就消逝了；人一氣絕，他在何處呢？ \end{tabularx} \\ \\ \relax
14:11 & \begin{tabularx}{0.7\textwidth}{X} 海中的水枯竭，江河消散乾涸。 \end{tabularx} \\ \\ \relax
14:12 & \begin{tabularx}{0.7\textwidth}{X} 人一躺下就不再起來，等到諸天沒有了，仍不復醒，也不能從睡中喚醒。 \end{tabularx} \\ \\ \relax
14:13 & \begin{tabularx}{0.7\textwidth}{X} 惟願你把我藏在陰間，把我隱藏，直到你的憤怒過去；願你為我定下期限，並記得我。 \end{tabularx} \\ \\ \relax
14:14 & \begin{tabularx}{0.7\textwidth}{X} 壯士若死了能再活嗎？我在一切服役的日子中等待，直到我退伍的時候來到。 \end{tabularx} \\ \\ \relax
14:15 & \begin{tabularx}{0.7\textwidth}{X} 你呼叫，我就回答你；你手所做的，你必期待。 \end{tabularx} \\ \\ \relax
14:16 & \begin{tabularx}{0.7\textwidth}{X} 但如今你數點我的腳步，不察看我的罪。 \end{tabularx} \\ \\ \relax
14:17 & \begin{tabularx}{0.7\textwidth}{X} 我的過犯被你密封在囊中，你遮掩了我的罪孽。 \end{tabularx} \\ \\ \relax
14:18 & \begin{tabularx}{0.7\textwidth}{X} 「然而，山崩變為無有，磐石從原處挪移。 \end{tabularx} \\ \\ \relax
14:19 & \begin{tabularx}{0.7\textwidth}{X} 流水沖蝕石頭，急流洗去地上的塵土；你也照樣滅絕人的指望。 \end{tabularx} \\ \\ \relax
14:20 & \begin{tabularx}{0.7\textwidth}{X} 你終必勝過人，使他消逝；你改變他的容貌，把他送走。 \end{tabularx} \\ \\ \relax
14:21 & \begin{tabularx}{0.7\textwidth}{X} 他的兒子得尊榮，他不知道；他們降為卑，他也不曉得。 \end{tabularx} \\ \\ \relax
14:22 & \begin{tabularx}{0.7\textwidth}{X} 他只覺得身上疼痛，心中為自己悲哀。」 \end{tabularx} \\ \\ \relax
15:1 & \begin{tabularx}{0.7\textwidth}{X} 提幔人以利法回答說： \end{tabularx} \\ \\ \relax
15:2 & \begin{tabularx}{0.7\textwidth}{X} 「智慧人豈可用虛空的知識回答，用東風充滿自己的肚腹呢？ \end{tabularx} \\ \\ \relax
15:3 & \begin{tabularx}{0.7\textwidth}{X} 他豈可用無益的話，用無濟於事的言語理論呢？ \end{tabularx} \\ \\ \relax
15:4 & \begin{tabularx}{0.7\textwidth}{X} 你誠然廢棄敬畏，不在神面前默想。 \end{tabularx} \\ \\ \relax
15:5 & \begin{tabularx}{0.7\textwidth}{X} 你的罪孽指教你的口；你選用詭詐人的舌頭。 \end{tabularx} \\ \\ \relax
15:6 & \begin{tabularx}{0.7\textwidth}{X} 你自己的口定你有罪，並非是我；你自己的嘴唇見證你的不是。 \end{tabularx} \\ \\ \relax
15:7 & \begin{tabularx}{0.7\textwidth}{X} 「你是頭一個生下來的人嗎？你受造在諸山之先嗎？ \end{tabularx} \\ \\ \relax
15:8 & \begin{tabularx}{0.7\textwidth}{X} 你曾聽見神的密旨嗎？你要獨自得盡智慧嗎？ \end{tabularx} \\ \\ \relax
15:9 & \begin{tabularx}{0.7\textwidth}{X} 甚麼是你知道，我們不知道的呢？甚麼是你明白，我們不明白的呢？ \end{tabularx} \\ \\ \relax
15:10 & \begin{tabularx}{0.7\textwidth}{X} 我們這裡有白髮的和年老的，比你父親還年長。 \end{tabularx} \\ \\ \relax
15:11 & \begin{tabularx}{0.7\textwidth}{X} 神的安慰和對你溫和的話，你以為太小嗎？ \end{tabularx} \\ \\ \relax
15:12 & \begin{tabularx}{0.7\textwidth}{X} 你的心為何失控，你的眼為何冒火， \end{tabularx} \\ \\ \relax
15:13 & \begin{tabularx}{0.7\textwidth}{X} 以致你的靈反對神，你的口說出這樣的言語呢？ \end{tabularx} \\ \\ \relax
15:14 & \begin{tabularx}{0.7\textwidth}{X} 人是甚麼，竟算為潔淨呢？婦人所生的是甚麼，竟算為義呢？ \end{tabularx} \\ \\ \relax
15:15 & \begin{tabularx}{0.7\textwidth}{X} 看哪，神不信任他的眾聖者；在他眼前，天也不潔淨， \end{tabularx} \\ \\ \relax
15:16 & \begin{tabularx}{0.7\textwidth}{X} 何況那污穢可憎，喝罪孽如水的世人呢！ \end{tabularx} \\ \\ \relax
15:17 & \begin{tabularx}{0.7\textwidth}{X} 「我指示你，你要聽我；我要陳述我所看見的， \end{tabularx} \\ \\ \relax
15:18 & \begin{tabularx}{0.7\textwidth}{X} 就是智慧人從列祖所受，傳講而不隱瞞的事。 \end{tabularx} \\ \\ \relax
15:19 & \begin{tabularx}{0.7\textwidth}{X} 這地惟獨賜給他們，並沒有外人從他們中間經過。 \end{tabularx} \\ \\ \relax
15:20 & \begin{tabularx}{0.7\textwidth}{X} 惡人一生的日子絞痛難熬，殘暴人存留的年數也是如此。 \end{tabularx} \\ \\ \relax
15:21 & \begin{tabularx}{0.7\textwidth}{X} 驚嚇的聲音常在他耳中；在平安時，毀滅者必臨到他。 \end{tabularx} \\ \\ \relax
15:22 & \begin{tabularx}{0.7\textwidth}{X} 他不信自己能從黑暗中轉回；他被刀劍看守。 \end{tabularx} \\ \\ \relax
15:23 & \begin{tabularx}{0.7\textwidth}{X} 他飄流在外求食：『哪裡有食物呢？』他知道黑暗的日子在他手邊預備好了。 \end{tabularx} \\ \\ \relax
15:24 & \begin{tabularx}{0.7\textwidth}{X} 急難困苦叫他害怕，而且勝過他，好像君王預備上陣。 \end{tabularx} \\ \\ \relax
15:25 & \begin{tabularx}{0.7\textwidth}{X} 因他伸手攻擊神，逞強對抗全能者， \end{tabularx} \\ \\ \relax
15:26 & \begin{tabularx}{0.7\textwidth}{X} 挺著頸項，用盾牌堅厚的凸面向全能者直闖； \end{tabularx} \\ \\ \relax
15:27 & \begin{tabularx}{0.7\textwidth}{X} 又因他的臉蒙上油脂，腰上積滿肥肉。 \end{tabularx} \\ \\ \relax
15:28 & \begin{tabularx}{0.7\textwidth}{X} 他住在荒涼的城鎮，房屋無人居住，成為廢墟。 \end{tabularx} \\ \\ \relax
15:29 & \begin{tabularx}{0.7\textwidth}{X} 他不得富足，財物不得常存，產業在地上也不加增。 \end{tabularx} \\ \\ \relax
15:30 & \begin{tabularx}{0.7\textwidth}{X} 他不得脫離黑暗，火焰要把他的嫩枝燒乾；因神口中的氣，他要離去。 \end{tabularx} \\ \\ \relax
15:31 & \begin{tabularx}{0.7\textwidth}{X} 不要讓他倚靠虛假，欺騙自己，因虛假必成為他的報應。 \end{tabularx} \\ \\ \relax
15:32 & \begin{tabularx}{0.7\textwidth}{X} 他的日期未到之先，這事必實現；他的枝子不得青綠。 \end{tabularx} \\ \\ \relax
15:33 & \begin{tabularx}{0.7\textwidth}{X} 他必像葡萄樹，葡萄未熟就掉落；又像橄欖樹，一開花就凋謝。 \end{tabularx} \\ \\ \relax
15:34 & \begin{tabularx}{0.7\textwidth}{X} 因不敬虔之輩必不能生育，受賄賂之人的帳棚必被火吞滅。 \end{tabularx} \\ \\ \relax
15:35 & \begin{tabularx}{0.7\textwidth}{X} 他們所懷的是毒害，所生的是罪孽，肚腹裡所預備的是詭詐。」 \end{tabularx} \\ \\ \relax
16:1 & \begin{tabularx}{0.7\textwidth}{X} 約伯回答說： \end{tabularx} \\ \\ \relax
16:2 & \begin{tabularx}{0.7\textwidth}{X} 「這樣的話我聽了許多；你們全都是使人愁煩的安慰者。 \end{tabularx} \\ \\ \relax
16:3 & \begin{tabularx}{0.7\textwidth}{X} 如風的言語有窮盡嗎？或者甚麼惹動你回答呢？ \end{tabularx} \\ \\ \relax
16:4 & \begin{tabularx}{0.7\textwidth}{X} 我也能說你們那樣的話，你們若處在我的景況，我也可以堆砌言詞攻擊你們，又可以向你們搖頭。 \end{tabularx} \\ \\ \relax
16:5 & \begin{tabularx}{0.7\textwidth}{X} 但我必用口堅固你們，顫動的嘴唇帶來舒解。 \end{tabularx} \\ \\ \relax
16:6 & \begin{tabularx}{0.7\textwidth}{X} 「我若說話，痛苦仍不得緩解；我若停止，痛苦就離開我嗎？ \end{tabularx} \\ \\ \relax
16:7 & \begin{tabularx}{0.7\textwidth}{X} 但現在神使我困倦，你使所有的親友遠離我， \end{tabularx} \\ \\ \relax
16:8 & \begin{tabularx}{0.7\textwidth}{X} 你抓住我，成為見證起來攻擊我；我的枯瘦也當著我的面作證。 \end{tabularx} \\ \\ \relax
16:9 & \begin{tabularx}{0.7\textwidth}{X} 神發怒撕裂我，逼迫我，向我咬牙切齒；我的敵人怒目瞪我。 \end{tabularx} \\ \\ \relax
16:10 & \begin{tabularx}{0.7\textwidth}{X} 他們向我大大張口，打我的耳光羞辱我，聚在一起攻擊我。 \end{tabularx} \\ \\ \relax
16:11 & \begin{tabularx}{0.7\textwidth}{X} 神把我交給不敬虔的人，把我扔到惡人的手中。 \end{tabularx} \\ \\ \relax
16:12 & \begin{tabularx}{0.7\textwidth}{X} 我本是安逸，他折斷我，掐住我的頸項，把我摔碎，又立我作他的箭靶。 \end{tabularx} \\ \\ \relax
16:13 & \begin{tabularx}{0.7\textwidth}{X} 他的弓箭手圍繞我。他刺破我的腎臟，並不留情，把我的膽汁傾倒在地上。 \end{tabularx} \\ \\ \relax
16:14 & \begin{tabularx}{0.7\textwidth}{X} 他使我破裂，破裂又破裂，如同勇士向我直闖。 \end{tabularx} \\ \\ \relax
16:15 & \begin{tabularx}{0.7\textwidth}{X} 「我把麻布縫在我的皮膚上，把我的角放在塵土中。 \end{tabularx} \\ \\ \relax
16:16 & \begin{tabularx}{0.7\textwidth}{X} 我的臉因哭泣變紅，我的眼皮上有死蔭。 \end{tabularx} \\ \\ \relax
16:17 & \begin{tabularx}{0.7\textwidth}{X} 我的手中卻沒有暴力，我的祈禱也是純潔的。 \end{tabularx} \\ \\ \relax
16:18 & \begin{tabularx}{0.7\textwidth}{X} 「地啊，不要遮蓋我的血！不要讓我的哀求有藏匿之處！ \end{tabularx} \\ \\ \relax
16:19 & \begin{tabularx}{0.7\textwidth}{X} 現今，看哪，在天有我的見證，在上有我的保人。 \end{tabularx} \\ \\ \relax
16:20 & \begin{tabularx}{0.7\textwidth}{X} 我的朋友譏誚我，我卻向神眼淚汪汪。 \end{tabularx} \\ \\ \relax
16:21 & \begin{tabularx}{0.7\textwidth}{X} 願人可與神理論，如同人與朋友一樣； \end{tabularx} \\ \\ \relax
16:22 & \begin{tabularx}{0.7\textwidth}{X} 因為再過幾年，我必走那往而不返之路。」 \end{tabularx} \\ \\ \relax
17:1 & \begin{tabularx}{0.7\textwidth}{X} 「我的靈耗盡，我的日子消逝；墳墓為我預備好了。 \end{tabularx} \\ \\ \relax
17:2 & \begin{tabularx}{0.7\textwidth}{X} 戲笑的人果真陪伴著我，我的眼睛盯住他們的悖逆。 \end{tabularx} \\ \\ \relax
17:3 & \begin{tabularx}{0.7\textwidth}{X} 「願你親自為我付押擔保。誰還會與我擊掌呢？ \end{tabularx} \\ \\ \relax
17:4 & \begin{tabularx}{0.7\textwidth}{X} 因你蒙蔽他們的心，使不明理，所以你必不高舉他們。 \end{tabularx} \\ \\ \relax
17:5 & \begin{tabularx}{0.7\textwidth}{X} 控告朋友為了分享產業的，他兒女的眼睛要失明。 \end{tabularx} \\ \\ \relax
17:6 & \begin{tabularx}{0.7\textwidth}{X} 「神使我成為人群中的笑談，他們吐唾沫在我臉上。 \end{tabularx} \\ \\ \relax
17:7 & \begin{tabularx}{0.7\textwidth}{X} 我的眼睛因憂愁昏花，我的肢體全像影兒。 \end{tabularx} \\ \\ \relax
17:8 & \begin{tabularx}{0.7\textwidth}{X} 正直人因此必驚奇；無辜的人要興起攻擊不敬虔之輩。 \end{tabularx} \\ \\ \relax
17:9 & \begin{tabularx}{0.7\textwidth}{X} 然而，義人要持守所行的道，手潔的人要力上加力。 \end{tabularx} \\ \\ \relax
17:10 & \begin{tabularx}{0.7\textwidth}{X} 至於你們眾人，再回來吧！你們中間，我找不到一個智慧人。 \end{tabularx} \\ \\ \relax
17:11 & \begin{tabularx}{0.7\textwidth}{X} 我的日子已經過去了，我的謀算、我心的願望已經斷絕了。 \end{tabularx} \\ \\ \relax
17:12 & \begin{tabularx}{0.7\textwidth}{X} 他們以黑夜為白晝，即使面臨黑暗，以為亮光已近。 \end{tabularx} \\ \\ \relax
17:13 & \begin{tabularx}{0.7\textwidth}{X} 我若盼望陰間為我的家，若下榻在黑暗中， \end{tabularx} \\ \\ \relax
17:14 & \begin{tabularx}{0.7\textwidth}{X} 若對地府呼叫：『你是我的父親』，若對蟲呼叫：『你是我的母親、姊妹』， \end{tabularx} \\ \\ \relax
17:15 & \begin{tabularx}{0.7\textwidth}{X} 這樣，我的盼望在哪裡呢？我所盼望的，誰能看見呢？ \end{tabularx} \\ \\ \relax
17:16 & \begin{tabularx}{0.7\textwidth}{X} 這盼望要下到陰間的門閂嗎？要一起在塵土中安息嗎？」 \end{tabularx} \\ \\ \relax
18:1 & \begin{tabularx}{0.7\textwidth}{X} 書亞人比勒達回答說： \end{tabularx} \\ \\ \relax
18:2 & \begin{tabularx}{0.7\textwidth}{X} 「你們尋索言語要到幾時呢？你們要明白，然後我們才說話。 \end{tabularx} \\ \\ \relax
18:3 & \begin{tabularx}{0.7\textwidth}{X} 我們為何被視為畜牲，在你們眼中看為愚笨呢？ \end{tabularx} \\ \\ \relax
18:4 & \begin{tabularx}{0.7\textwidth}{X} 在怒氣中將自己撕裂的人哪，難道大地要因你見棄、磐石要挪開原處嗎？ \end{tabularx} \\ \\ \relax
18:5 & \begin{tabularx}{0.7\textwidth}{X} 「惡人的亮光必要熄滅，他的火焰必不照耀。 \end{tabularx} \\ \\ \relax
18:6 & \begin{tabularx}{0.7\textwidth}{X} 他帳棚中的亮光要變黑暗，他上面的燈也必熄滅。 \end{tabularx} \\ \\ \relax
18:7 & \begin{tabularx}{0.7\textwidth}{X} 他強橫的腳步必遭阻礙，他的計謀必將自己絆倒。 \end{tabularx} \\ \\ \relax
18:8 & \begin{tabularx}{0.7\textwidth}{X} 他因自己的腳陷入網中，走在纏人的網子上。 \end{tabularx} \\ \\ \relax
18:9 & \begin{tabularx}{0.7\textwidth}{X} 羅網必抓住他的腳跟，陷阱必擒獲他。 \end{tabularx} \\ \\ \relax
18:10 & \begin{tabularx}{0.7\textwidth}{X} 繩索為他藏在土裡，羈絆為他藏在路上。 \end{tabularx} \\ \\ \relax
18:11 & \begin{tabularx}{0.7\textwidth}{X} 四面的驚嚇使他害怕，在他腳跟後面追趕他。 \end{tabularx} \\ \\ \relax
18:12 & \begin{tabularx}{0.7\textwidth}{X} 他的力量必因飢餓衰敗，禍患要在他的旁邊等候， \end{tabularx} \\ \\ \relax
18:13 & \begin{tabularx}{0.7\textwidth}{X} 侵蝕他肢體的皮膚；死亡的長子吞吃他的肢體。 \end{tabularx} \\ \\ \relax
18:14 & \begin{tabularx}{0.7\textwidth}{X} 他要從所倚靠的帳棚被拔出來，帶到使人驚恐的王那裡。 \end{tabularx} \\ \\ \relax
18:15 & \begin{tabularx}{0.7\textwidth}{X} 不屬他的必住在他的帳棚裡，硫磺必撒在他所住之處。 \end{tabularx} \\ \\ \relax
18:16 & \begin{tabularx}{0.7\textwidth}{X} 下邊，他的根要枯乾；上邊，他的枝子要剪除。 \end{tabularx} \\ \\ \relax
18:17 & \begin{tabularx}{0.7\textwidth}{X} 他的稱號從地上消失，他的名字不在街上存留。 \end{tabularx} \\ \\ \relax
18:18 & \begin{tabularx}{0.7\textwidth}{X} 他必從光明中被驅逐到黑暗裡，他必被趕出世界。 \end{tabularx} \\ \\ \relax
18:19 & \begin{tabularx}{0.7\textwidth}{X} 他在自己百姓中必無子無孫，在寄居之地也沒有倖存者。 \end{tabularx} \\ \\ \relax
18:20 & \begin{tabularx}{0.7\textwidth}{X} 以後的人要因他的日子驚訝，以前的人也被驚駭抓住。 \end{tabularx} \\ \\ \relax
18:21 & \begin{tabularx}{0.7\textwidth}{X} 不義之人的住處總是這樣，這就是不認識神之人的下場。」 \end{tabularx} \\ \\ \relax
19:1 & \begin{tabularx}{0.7\textwidth}{X} 約伯回答說： \end{tabularx} \\ \\ \relax
19:2 & \begin{tabularx}{0.7\textwidth}{X} 「你們攪擾我的心，用言語壓碎我要到幾時呢？ \end{tabularx} \\ \\ \relax
19:3 & \begin{tabularx}{0.7\textwidth}{X} 你們這十次羞辱我，苦待我也不以為恥。 \end{tabularx} \\ \\ \relax
19:4 & \begin{tabularx}{0.7\textwidth}{X} 果真我有錯，這錯是在於我。 \end{tabularx} \\ \\ \relax
19:5 & \begin{tabularx}{0.7\textwidth}{X} 若你們真要向我誇大，以我的羞辱來責備我， \end{tabularx} \\ \\ \relax
19:6 & \begin{tabularx}{0.7\textwidth}{X} 就該知道是神傾覆我，用羅網圍繞我。 \end{tabularx} \\ \\ \relax
19:7 & \begin{tabularx}{0.7\textwidth}{X} 看哪，我喊冤叫屈，卻不蒙應允；我呼求，卻沒有公正。 \end{tabularx} \\ \\ \relax
19:8 & \begin{tabularx}{0.7\textwidth}{X} 神攔住我的道路，使我不得經過；他使黑暗籠罩我的路徑。 \end{tabularx} \\ \\ \relax
19:9 & \begin{tabularx}{0.7\textwidth}{X} 他剝去我的榮光，摘去我頭上的冠冕。 \end{tabularx} \\ \\ \relax
19:10 & \begin{tabularx}{0.7\textwidth}{X} 他在四圍攻擊我，我就走了；他將我的指望如樹拔出。 \end{tabularx} \\ \\ \relax
19:11 & \begin{tabularx}{0.7\textwidth}{X} 他向我發烈怒，以我為他的敵人。 \end{tabularx} \\ \\ \relax
19:12 & \begin{tabularx}{0.7\textwidth}{X} 他的軍隊一齊上來，修築道路攻擊我，在我帳棚的四圍安營。 \end{tabularx} \\ \\ \relax
19:13 & \begin{tabularx}{0.7\textwidth}{X} 「他把我的兄弟隔在遠處，使我認識的人全然與我生疏。 \end{tabularx} \\ \\ \relax
19:14 & \begin{tabularx}{0.7\textwidth}{X} 我的親戚都離開了我；我的密友都忘記了我。 \end{tabularx} \\ \\ \relax
19:15 & \begin{tabularx}{0.7\textwidth}{X} 在我家寄居的和我的使女，都當我是陌生人；我在他們眼中被視為外邦人。 \end{tabularx} \\ \\ \relax
19:16 & \begin{tabularx}{0.7\textwidth}{X} 我呼喚僕人，他卻不回答；我必須親口求他。 \end{tabularx} \\ \\ \relax
19:17 & \begin{tabularx}{0.7\textwidth}{X} 我口的氣味令我妻子厭惡，我的同胞都憎惡我。 \end{tabularx} \\ \\ \relax
19:18 & \begin{tabularx}{0.7\textwidth}{X} 連小男孩也藐視我；我起來，他們都嘲笑我。 \end{tabularx} \\ \\ \relax
19:19 & \begin{tabularx}{0.7\textwidth}{X} 我的知心朋友都憎惡我；我平日所愛的人向我翻臉。 \end{tabularx} \\ \\ \relax
19:20 & \begin{tabularx}{0.7\textwidth}{X} 我的皮和肉緊貼骨頭，我得以逃脫，僅剩牙齒。 \end{tabularx} \\ \\ \relax
19:21 & \begin{tabularx}{0.7\textwidth}{X} 我的朋友啊，可憐我！可憐我！因為神的手攻擊我。 \end{tabularx} \\ \\ \relax
19:22 & \begin{tabularx}{0.7\textwidth}{X} 你們為甚麼彷彿神逼迫我，吃我的肉還不滿足呢？ \end{tabularx} \\ \\ \relax
19:23 & \begin{tabularx}{0.7\textwidth}{X} 「惟願我的言語現在就寫上，都記錄在書上； \end{tabularx} \\ \\ \relax
19:24 & \begin{tabularx}{0.7\textwidth}{X} 用鐵筆和鉛，刻在磐石上，存到永遠。 \end{tabularx} \\ \\ \relax
19:25 & \begin{tabularx}{0.7\textwidth}{X} 我知道我的救贖主活著，末後他必站在塵土上。 \end{tabularx} \\ \\ \relax
19:26 & \begin{tabularx}{0.7\textwidth}{X} 我這皮肉滅絕之後，我必在肉體之外得見神。 \end{tabularx} \\ \\ \relax
19:27 & \begin{tabularx}{0.7\textwidth}{X} 我自己要見他，親眼要看他，並不像陌生人。我的心腸在我裡面耗盡了！ \end{tabularx} \\ \\ \relax
19:28 & \begin{tabularx}{0.7\textwidth}{X} 你們若說：『我們怎麼逼迫他呢？事情的根源是在於他』， \end{tabularx} \\ \\ \relax
19:29 & \begin{tabularx}{0.7\textwidth}{X} 你們就當懼怕刀劍，因為憤怒帶來刀劍的刑罰。這樣，你們就知道有審判。」 \end{tabularx} \\ \\ \relax
20:1 & \begin{tabularx}{0.7\textwidth}{X} 拿瑪人瑣法回答說： \end{tabularx} \\ \\ \relax
20:2 & \begin{tabularx}{0.7\textwidth}{X} 「這樣，我的思念叫我回答，因為我心中急躁。 \end{tabularx} \\ \\ \relax
20:3 & \begin{tabularx}{0.7\textwidth}{X} 我聽見那羞辱我的責備；我悟性的靈回答我。 \end{tabularx} \\ \\ \relax
20:4 & \begin{tabularx}{0.7\textwidth}{X} 你豈不知道嗎？亙古以來，自從人被安置在地， \end{tabularx} \\ \\ \relax
20:5 & \begin{tabularx}{0.7\textwidth}{X} 惡人歡樂的聲音是暫時的，不敬虔人的喜樂不過是轉眼之間。 \end{tabularx} \\ \\ \relax
20:6 & \begin{tabularx}{0.7\textwidth}{X} 他的尊榮雖達到天上，頭雖頂到雲中， \end{tabularx} \\ \\ \relax
20:7 & \begin{tabularx}{0.7\textwidth}{X} 他必永遠滅亡，像自己的糞一樣。看見他的人要說：『他在哪裡呢？』 \end{tabularx} \\ \\ \relax
20:8 & \begin{tabularx}{0.7\textwidth}{X} 他必如夢飛去，不再尋見；他被趕走，如夜間的異象。 \end{tabularx} \\ \\ \relax
20:9 & \begin{tabularx}{0.7\textwidth}{X} 親眼見過他的，必不再見他；他自己的地方也不再見到他。 \end{tabularx} \\ \\ \relax
20:10 & \begin{tabularx}{0.7\textwidth}{X} 他的兒女要向窮人求恩；他的手要賠還錢財。 \end{tabularx} \\ \\ \relax
20:11 & \begin{tabularx}{0.7\textwidth}{X} 他的骨頭雖然滿有年輕的活力，卻要和他一同躺臥在塵土之中。 \end{tabularx} \\ \\ \relax
20:12 & \begin{tabularx}{0.7\textwidth}{X} 「他口中以惡為甘甜，把惡藏在舌頭底下， \end{tabularx} \\ \\ \relax
20:13 & \begin{tabularx}{0.7\textwidth}{X} 愛戀不捨，含在口中。 \end{tabularx} \\ \\ \relax
20:14 & \begin{tabularx}{0.7\textwidth}{X} 他的食物在肚裡卻要翻轉，在他裡面成為虺蛇的毒液。 \end{tabularx} \\ \\ \relax
20:15 & \begin{tabularx}{0.7\textwidth}{X} 他吞了財寶，還要吐出；神要從他腹中掏出來。 \end{tabularx} \\ \\ \relax
20:16 & \begin{tabularx}{0.7\textwidth}{X} 他必吸飲虺蛇的毒汁，毒蛇的舌頭必殺他。 \end{tabularx} \\ \\ \relax
20:17 & \begin{tabularx}{0.7\textwidth}{X} 他不再看見溪流，流奶與蜜之河。 \end{tabularx} \\ \\ \relax
20:18 & \begin{tabularx}{0.7\textwidth}{X} 他勞碌得來的要賠還，不得吞下；賺取了財貨，也不得歡樂。 \end{tabularx} \\ \\ \relax
20:19 & \begin{tabularx}{0.7\textwidth}{X} 他欺壓窮人，棄之不顧，強取非自己所蓋的房屋。 \end{tabularx} \\ \\ \relax
20:20 & \begin{tabularx}{0.7\textwidth}{X} 「他的肚腹不知安逸，所貪戀的連一樣也不放過， \end{tabularx} \\ \\ \relax
20:21 & \begin{tabularx}{0.7\textwidth}{X} 剩餘的沒有一樣他不吞吃，所以他的福樂不能長久。 \end{tabularx} \\ \\ \relax
20:22 & \begin{tabularx}{0.7\textwidth}{X} 他在滿足有餘的時候，必有困苦臨到；凡受苦楚之人的手必加在他身上。 \end{tabularx} \\ \\ \relax
20:23 & \begin{tabularx}{0.7\textwidth}{X} 他的肚腹正要滿足的時候，神必將猛烈的憤怒降在他身上；他正在吃飯的時候，神要將這憤怒如雨降在他身上。 \end{tabularx} \\ \\ \relax
20:24 & \begin{tabularx}{0.7\textwidth}{X} 他要躲避鐵的武器，銅弓要將他射透。 \end{tabularx} \\ \\ \relax
20:25 & \begin{tabularx}{0.7\textwidth}{X} 箭一抽，就從他背上出來，發亮的箭頭從他膽中出來；有驚惶臨到他身上。 \end{tabularx} \\ \\ \relax
20:26 & \begin{tabularx}{0.7\textwidth}{X} 他的財寶隱藏在深沉的黑暗裡；有非人吹起的火要把他吞滅，把他帳棚中所剩下的燒燬。 \end{tabularx} \\ \\ \relax
20:27 & \begin{tabularx}{0.7\textwidth}{X} 天要顯明他的罪孽，地要興起去攻擊他。 \end{tabularx} \\ \\ \relax
20:28 & \begin{tabularx}{0.7\textwidth}{X} 他家裡出產的必消失，在神憤怒的日子被沖走。 \end{tabularx} \\ \\ \relax
20:29 & \begin{tabularx}{0.7\textwidth}{X} 這是惡人從神所得的份，是神為他所定的產業。」 \end{tabularx} \\ \\ \relax
21:1 & \begin{tabularx}{0.7\textwidth}{X} 約伯回答說： \end{tabularx} \\ \\ \relax
21:2 & \begin{tabularx}{0.7\textwidth}{X} 「你們要細心聽我的言語，這就算是你們的安慰。 \end{tabularx} \\ \\ \relax
21:3 & \begin{tabularx}{0.7\textwidth}{X} 請寬容我，我又要說話；說了以後，任憑你嗤笑吧！ \end{tabularx} \\ \\ \relax
21:4 & \begin{tabularx}{0.7\textwidth}{X} 我豈是向人訴苦？我為何不是沒有耐心呢？ \end{tabularx} \\ \\ \relax
21:5 & \begin{tabularx}{0.7\textwidth}{X} 你們要轉向我而驚奇，要用手摀口。 \end{tabularx} \\ \\ \relax
21:6 & \begin{tabularx}{0.7\textwidth}{X} 我每逢思想，心就驚惶，戰兢抓住我身。 \end{tabularx} \\ \\ \relax
21:7 & \begin{tabularx}{0.7\textwidth}{X} 惡人為何存活，得享高壽，勢力強盛呢？ \end{tabularx} \\ \\ \relax
21:8 & \begin{tabularx}{0.7\textwidth}{X} 他們的後裔與他們一起，堅立在他們面前，他們得以眼見自己的子孫。 \end{tabularx} \\ \\ \relax
21:9 & \begin{tabularx}{0.7\textwidth}{X} 他們的家宅平安無懼，神的杖不加在他們身上。 \end{tabularx} \\ \\ \relax
21:10 & \begin{tabularx}{0.7\textwidth}{X} 他們的公牛傳種而不斷絕，母牛生牛犢而不掉胎。 \end{tabularx} \\ \\ \relax
21:11 & \begin{tabularx}{0.7\textwidth}{X} 他們打發小男孩出去，多如羊群，他們的孩子踴躍跳舞。 \end{tabularx} \\ \\ \relax
21:12 & \begin{tabularx}{0.7\textwidth}{X} 他們隨著琴鼓歌唱，因簫聲歡喜。 \end{tabularx} \\ \\ \relax
21:13 & \begin{tabularx}{0.7\textwidth}{X} 他們度日諸事亨通，在平安中下到陰間。 \end{tabularx} \\ \\ \relax
21:14 & \begin{tabularx}{0.7\textwidth}{X} 他們對神說：『離開我們吧！我們不想知道你的道路。 \end{tabularx} \\ \\ \relax
21:15 & \begin{tabularx}{0.7\textwidth}{X} 全能者是誰，我們何必事奉他呢？求告他有甚麼益處呢？』 \end{tabularx} \\ \\ \relax
21:16 & \begin{tabularx}{0.7\textwidth}{X} 看哪，他們亨通不是靠自己的手；惡人的計謀離我好遠。 \end{tabularx} \\ \\ \relax
21:17 & \begin{tabularx}{0.7\textwidth}{X} 「惡人的燈何嘗熄滅？患難何嘗臨到他們呢？神何嘗發怒，把災禍分給他們呢？ \end{tabularx} \\ \\ \relax
21:18 & \begin{tabularx}{0.7\textwidth}{X} 他們何嘗像風前的碎秸，如暴風颳去的糠秕呢？ \end{tabularx} \\ \\ \relax
21:19 & \begin{tabularx}{0.7\textwidth}{X} 神為惡人的兒女積蓄罪孽，不如本人遭報，好使他親自知道。 \end{tabularx} \\ \\ \relax
21:20 & \begin{tabularx}{0.7\textwidth}{X} 願他親眼看見自己敗亡，親自飲全能者的憤怒。 \end{tabularx} \\ \\ \relax
21:21 & \begin{tabularx}{0.7\textwidth}{X} 他的歲月既盡，他身後還顧他的家嗎？ \end{tabularx} \\ \\ \relax
21:22 & \begin{tabularx}{0.7\textwidth}{X} 誰能將知識教導神呢？是他審判那些居高位的。 \end{tabularx} \\ \\ \relax
21:23 & \begin{tabularx}{0.7\textwidth}{X} 有人至死身體強壯，盡得平順安逸； \end{tabularx} \\ \\ \relax
21:24 & \begin{tabularx}{0.7\textwidth}{X} 他的肚腹充滿奶汁，他的骨髓滋潤。 \end{tabularx} \\ \\ \relax
21:25 & \begin{tabularx}{0.7\textwidth}{X} 有人至死心中痛苦，從未嘗過福樂的滋味； \end{tabularx} \\ \\ \relax
21:26 & \begin{tabularx}{0.7\textwidth}{X} 他們同樣躺臥於塵土，蟲子覆蓋他們。 \end{tabularx} \\ \\ \relax
21:27 & \begin{tabularx}{0.7\textwidth}{X} 「看哪，我知道你們的意念，並殘害我的計謀。 \end{tabularx} \\ \\ \relax
21:28 & \begin{tabularx}{0.7\textwidth}{X} 你們說：『權貴的房屋在哪裡？惡人住過的帳棚在哪裡？』 \end{tabularx} \\ \\ \relax
21:29 & \begin{tabularx}{0.7\textwidth}{X} 你們沒有詢問那些過路的人嗎？你們不承認他們的證據嗎？ \end{tabularx} \\ \\ \relax
21:30 & \begin{tabularx}{0.7\textwidth}{X} 就是惡人在患難的日子得存留，在憤怒的日子得逃脫。 \end{tabularx} \\ \\ \relax
21:31 & \begin{tabularx}{0.7\textwidth}{X} 他所行的，有誰當面給他說明？他所做的，有誰報應他呢？ \end{tabularx} \\ \\ \relax
21:32 & \begin{tabularx}{0.7\textwidth}{X} 然而他要被抬到墳地，並有人看守墓穴。 \end{tabularx} \\ \\ \relax
21:33 & \begin{tabularx}{0.7\textwidth}{X} 他要以谷中的土塊為甘甜；人人要跟在他後面，在他前面去的無數。 \end{tabularx} \\ \\ \relax
21:34 & \begin{tabularx}{0.7\textwidth}{X} 你們怎能以空話安慰我呢？你們的對答全都錯謬！」 \end{tabularx} \\ \\ \relax
22:1 & \begin{tabularx}{0.7\textwidth}{X} 提幔人以利法回答說： \end{tabularx} \\ \\ \relax
22:2 & \begin{tabularx}{0.7\textwidth}{X} 「人能使神有益嗎？智慧人能使他有益嗎？ \end{tabularx} \\ \\ \relax
22:3 & \begin{tabularx}{0.7\textwidth}{X} 你為人公義，豈能叫全能者喜悅呢？你行為完全，豈能使他得利呢？ \end{tabularx} \\ \\ \relax
22:4 & \begin{tabularx}{0.7\textwidth}{X} 他豈是因你敬畏的心就責備你，審判你嗎？ \end{tabularx} \\ \\ \relax
22:5 & \begin{tabularx}{0.7\textwidth}{X} 你的罪惡豈不是大嗎？你的罪孽不是沒有窮盡嗎？ \end{tabularx} \\ \\ \relax
22:6 & \begin{tabularx}{0.7\textwidth}{X} 因你無故強取弟兄的抵押，剝去赤身者的衣服。 \end{tabularx} \\ \\ \relax
22:7 & \begin{tabularx}{0.7\textwidth}{X} 疲乏的人，你沒有給他水喝；飢餓的人，你沒有給他食物。 \end{tabularx} \\ \\ \relax
22:8 & \begin{tabularx}{0.7\textwidth}{X} 有能力的人得土地；尊貴的人住在其中。 \end{tabularx} \\ \\ \relax
22:9 & \begin{tabularx}{0.7\textwidth}{X} 你打發寡婦空手回去，你折斷孤兒的膀臂。 \end{tabularx} \\ \\ \relax
22:10 & \begin{tabularx}{0.7\textwidth}{X} 因此，有羅網環繞你，有恐懼忽然使你驚惶； \end{tabularx} \\ \\ \relax
22:11 & \begin{tabularx}{0.7\textwidth}{X} 或有黑暗使你看不見，有洪水淹沒你。 \end{tabularx} \\ \\ \relax
22:12 & \begin{tabularx}{0.7\textwidth}{X} 「神豈不是在高天嗎？你看星宿的頂點何其高呢！ \end{tabularx} \\ \\ \relax
22:13 & \begin{tabularx}{0.7\textwidth}{X} 你說：『神知道甚麼？他豈能透過幽暗施行審判呢？ \end{tabularx} \\ \\ \relax
22:14 & \begin{tabularx}{0.7\textwidth}{X} 密雲將他遮蓋，使他不能看見；他周遊穹蒼。』 \end{tabularx} \\ \\ \relax
22:15 & \begin{tabularx}{0.7\textwidth}{X} 你要依從上古的道嗎？這道是惡人行過的。 \end{tabularx} \\ \\ \relax
22:16 & \begin{tabularx}{0.7\textwidth}{X} 他們未到時候就被抓去；他們的根基被江河沖去。 \end{tabularx} \\ \\ \relax
22:17 & \begin{tabularx}{0.7\textwidth}{X} 他們向神說：『離開我們吧！』全能者能把他們怎麼樣呢？ \end{tabularx} \\ \\ \relax
22:18 & \begin{tabularx}{0.7\textwidth}{X} 然而，是神以美物充滿他們的房屋；惡人的計謀離我好遠！ \end{tabularx} \\ \\ \relax
22:19 & \begin{tabularx}{0.7\textwidth}{X} 義人看見他們的結局就歡喜；無辜的人嗤笑他們： \end{tabularx} \\ \\ \relax
22:20 & \begin{tabularx}{0.7\textwidth}{X} 『攻擊我們的果然被剪除，剩餘的都被火吞滅。』 \end{tabularx} \\ \\ \relax
22:21 & \begin{tabularx}{0.7\textwidth}{X} 「你要與神和好，要和平，這樣，福氣必臨到你。 \end{tabularx} \\ \\ \relax
22:22 & \begin{tabularx}{0.7\textwidth}{X} 你當領受他口中的教導，將他的言語存在心裡。 \end{tabularx} \\ \\ \relax
22:23 & \begin{tabularx}{0.7\textwidth}{X} 你若歸向全能者，就必得建立。你要從你帳棚中遠離不義， \end{tabularx} \\ \\ \relax
22:24 & \begin{tabularx}{0.7\textwidth}{X} 你要將黃金丟到塵土裡，將俄斐的金子丟在溪河石頭之間； \end{tabularx} \\ \\ \relax
22:25 & \begin{tabularx}{0.7\textwidth}{X} 全能者就必作你的黃金，作你成堆的銀子。 \end{tabularx} \\ \\ \relax
22:26 & \begin{tabularx}{0.7\textwidth}{X} 那時，你要以全能者為喜樂，向神仰臉。 \end{tabularx} \\ \\ \relax
22:27 & \begin{tabularx}{0.7\textwidth}{X} 你要向他禱告，他就聽你；你也要還你的願。 \end{tabularx} \\ \\ \relax
22:28 & \begin{tabularx}{0.7\textwidth}{X} 你定意要做何事，必然為你成就；亮光也必照耀你的路。 \end{tabularx} \\ \\ \relax
22:29 & \begin{tabularx}{0.7\textwidth}{X} 當人降卑，你說：是因驕傲；眼目謙卑的人，神必然拯救。 \end{tabularx} \\ \\ \relax
22:30 & \begin{tabularx}{0.7\textwidth}{X} 不是無辜的人，神尚且要搭救他；他必因你手中的清潔得蒙拯救。」 \end{tabularx} \\ \\ \relax
23:1 & \begin{tabularx}{0.7\textwidth}{X} 約伯回答說： \end{tabularx} \\ \\ \relax
23:2 & \begin{tabularx}{0.7\textwidth}{X} 「如今我的哀告還算為悖逆；我雖唉哼，他的手仍然重重責罰我。 \end{tabularx} \\ \\ \relax
23:3 & \begin{tabularx}{0.7\textwidth}{X} 惟願我知道哪裡可以尋見神，能到他的臺前， \end{tabularx} \\ \\ \relax
23:4 & \begin{tabularx}{0.7\textwidth}{X} 我就在他面前陳明我的案件，滿口辯訴。 \end{tabularx} \\ \\ \relax
23:5 & \begin{tabularx}{0.7\textwidth}{X} 我必知道他回答我的言語，明白他向我所要說的。 \end{tabularx} \\ \\ \relax
23:6 & \begin{tabularx}{0.7\textwidth}{X} 他豈用大能與我爭辯呢？不！他必理會我。 \end{tabularx} \\ \\ \relax
23:7 & \begin{tabularx}{0.7\textwidth}{X} 在那裡正直人可以與他辯論，我就必永遠脫離那審判我的。 \end{tabularx} \\ \\ \relax
23:8 & \begin{tabularx}{0.7\textwidth}{X} 「看哪，我往前走，他不在那裡；往後退，也沒有察覺他。 \end{tabularx} \\ \\ \relax
23:9 & \begin{tabularx}{0.7\textwidth}{X} 他在左邊行事，我卻看不見他；他轉向右邊，我也見不到他。 \end{tabularx} \\ \\ \relax
23:10 & \begin{tabularx}{0.7\textwidth}{X} 然而他知道我所走的路；他試煉我，我就如純金。 \end{tabularx} \\ \\ \relax
23:11 & \begin{tabularx}{0.7\textwidth}{X} 我的腳緊跟他的步伐；我謹守他的道，並不偏離。 \end{tabularx} \\ \\ \relax
23:12 & \begin{tabularx}{0.7\textwidth}{X} 他嘴唇的命令，我未曾背棄；我看重他口中的言語，過於我需用的飲食。 \end{tabularx} \\ \\ \relax
23:13 & \begin{tabularx}{0.7\textwidth}{X} 只是他心志已定，誰能使他轉意呢？他心裡所願的，就行出來。 \end{tabularx} \\ \\ \relax
23:14 & \begin{tabularx}{0.7\textwidth}{X} 因此，為我所定的，他必做成，這類的事他還有許多。 \end{tabularx} \\ \\ \relax
23:15 & \begin{tabularx}{0.7\textwidth}{X} 所以我在他面前驚惶；我思想就懼怕他。 \end{tabularx} \\ \\ \relax
23:16 & \begin{tabularx}{0.7\textwidth}{X} 神使我喪膽，全能者使我驚惶。 \end{tabularx} \\ \\ \relax
23:17 & \begin{tabularx}{0.7\textwidth}{X} 但我並非被黑暗剪除，只是幽暗遮蓋了我的臉。 \end{tabularx} \\ \\ \relax
24:1 & \begin{tabularx}{0.7\textwidth}{X} 「為何全能者不定下期限？為何認識他的人看不到那些日子呢？ \end{tabularx} \\ \\ \relax
24:2 & \begin{tabularx}{0.7\textwidth}{X} 有人挪移地界，搶奪群畜去放牧。 \end{tabularx} \\ \\ \relax
24:3 & \begin{tabularx}{0.7\textwidth}{X} 他們拉走孤兒的驢，強取寡婦的牛作抵押。 \end{tabularx} \\ \\ \relax
24:4 & \begin{tabularx}{0.7\textwidth}{X} 他們使貧窮人離開正道；世上的困苦人盡都隱藏。 \end{tabularx} \\ \\ \relax
24:5 & \begin{tabularx}{0.7\textwidth}{X} 看哪，他們如同野驢出到曠野，殷勤尋找食物，在野地給孩童餬口。 \end{tabularx} \\ \\ \relax
24:6 & \begin{tabularx}{0.7\textwidth}{X} 他們收割別人田間的莊稼，摘取惡人剩餘的葡萄。 \end{tabularx} \\ \\ \relax
24:7 & \begin{tabularx}{0.7\textwidth}{X} 他們終夜赤身無衣，在寒冷中毫無遮蓋。 \end{tabularx} \\ \\ \relax
24:8 & \begin{tabularx}{0.7\textwidth}{X} 他們在山上被大雨淋濕，因沒有避身之處就擁抱磐石。 \end{tabularx} \\ \\ \relax
24:9 & \begin{tabularx}{0.7\textwidth}{X} 又有人從母懷中搶走孤兒，在困苦人身上強取抵押品。 \end{tabularx} \\ \\ \relax
24:10 & \begin{tabularx}{0.7\textwidth}{X} 困苦人赤身無衣，到處流浪，餓著肚子扛抬禾捆， \end{tabularx} \\ \\ \relax
24:11 & \begin{tabularx}{0.7\textwidth}{X} 他們在圍牆內榨油，踹壓酒池，自己卻口渴。 \end{tabularx} \\ \\ \relax
24:12 & \begin{tabularx}{0.7\textwidth}{X} 在城內垂死的人呻吟，受傷的人哀號；神卻不理會狂妄的事。 \end{tabularx} \\ \\ \relax
24:13 & \begin{tabularx}{0.7\textwidth}{X} 「又有人背棄光明，不認識光明的道，不留在光明的路上。 \end{tabularx} \\ \\ \relax
24:14 & \begin{tabularx}{0.7\textwidth}{X} 殺人者黎明起來，殺害困苦人和貧窮人，夜間又作盜賊。 \end{tabularx} \\ \\ \relax
24:15 & \begin{tabularx}{0.7\textwidth}{X} 姦夫的眼等候黃昏，說：『沒有眼睛能見我』，就把臉蒙住。 \end{tabularx} \\ \\ \relax
24:16 & \begin{tabularx}{0.7\textwidth}{X} 盜賊黑夜挖洞；他們白日躲藏，並不認識光明。 \end{tabularx} \\ \\ \relax
24:17 & \begin{tabularx}{0.7\textwidth}{X} 他們全都看早晨如死蔭，因為他們熟悉死蔭的驚駭。 \end{tabularx} \\ \\ \relax
24:18 & \begin{tabularx}{0.7\textwidth}{X} 「惡人在水面上快速飄盪，他們在地上所得的產業被詛咒；無人再回到他們的葡萄園。 \end{tabularx} \\ \\ \relax
24:19 & \begin{tabularx}{0.7\textwidth}{X} 乾旱炎熱融化雪水；陰間也如此吞沒犯罪的人。 \end{tabularx} \\ \\ \relax
24:20 & \begin{tabularx}{0.7\textwidth}{X} 懷他的母胎忘記他；蟲子要吃他，覺得甘甜；他不再被人記念；不義的人必如樹折斷。 \end{tabularx} \\ \\ \relax
24:21 & \begin{tabularx}{0.7\textwidth}{X} 「他與不懷孕不生育的婦人交往，卻不善待寡婦。 \end{tabularx} \\ \\ \relax
24:22 & \begin{tabularx}{0.7\textwidth}{X} 然而神用能力保全有勢力的人；那性命難保的人仍然興起。 \end{tabularx} \\ \\ \relax
24:23 & \begin{tabularx}{0.7\textwidth}{X} 神使他安穩，他就有所倚靠；神的眼目看顧他們的道路。 \end{tabularx} \\ \\ \relax
24:24 & \begin{tabularx}{0.7\textwidth}{X} 他們高升，不過片刻就沒有了；他們降為卑，被除滅，與眾人一樣，又如穀的穗子被割下。 \end{tabularx} \\ \\ \relax
24:25 & \begin{tabularx}{0.7\textwidth}{X} 若不是這樣，誰能指證我是說謊的，以我的言語為毫無根據呢？」 \end{tabularx} \\ \\ \relax
25:1 & \begin{tabularx}{0.7\textwidth}{X} 書亞人比勒達回答說： \end{tabularx} \\ \\ \relax
25:2 & \begin{tabularx}{0.7\textwidth}{X} 「神有統治之權，威嚴可畏；他在高處施行和平。 \end{tabularx} \\ \\ \relax
25:3 & \begin{tabularx}{0.7\textwidth}{X} 他的軍隊豈能數算？他的光向誰不會升起呢？ \end{tabularx} \\ \\ \relax
25:4 & \begin{tabularx}{0.7\textwidth}{X} 這樣，在神面前人怎能稱義？婦人所生的怎能潔淨？ \end{tabularx} \\ \\ \relax
25:5 & \begin{tabularx}{0.7\textwidth}{X} 看哪，在神眼前，月亮無光，星宿也不皎潔， \end{tabularx} \\ \\ \relax
25:6 & \begin{tabularx}{0.7\textwidth}{X} 更何況是如蟲的人，如蛆的世人呢！ \end{tabularx} \\ \\ \relax
26:1 & \begin{tabularx}{0.7\textwidth}{X} 約伯回答說： \end{tabularx} \\ \\ \relax
26:2 & \begin{tabularx}{0.7\textwidth}{X} 「無能的人蒙你何等的幫助！膀臂無力的人蒙你何等的拯救！ \end{tabularx} \\ \\ \relax
26:3 & \begin{tabularx}{0.7\textwidth}{X} 無智慧的人蒙你何等的指教！你向他顯出豐富的知識。 \end{tabularx} \\ \\ \relax
26:4 & \begin{tabularx}{0.7\textwidth}{X} 你向誰發出言語？誰的靈從你而出？ \end{tabularx} \\ \\ \relax
26:5 & \begin{tabularx}{0.7\textwidth}{X} 在大水和水族以下，陰魂戰兢。 \end{tabularx} \\ \\ \relax
26:6 & \begin{tabularx}{0.7\textwidth}{X} 在神面前，陰間顯露；冥府也不得遮掩。 \end{tabularx} \\ \\ \relax
26:7 & \begin{tabularx}{0.7\textwidth}{X} 神將北極鋪在空中，將大地懸在虛空。 \end{tabularx} \\ \\ \relax
26:8 & \begin{tabularx}{0.7\textwidth}{X} 他將水包在密雲中，盛水的雲卻不破裂。 \end{tabularx} \\ \\ \relax
26:9 & \begin{tabularx}{0.7\textwidth}{X} 他遮蔽寶座的正面，把他的雲彩鋪在其上。 \end{tabularx} \\ \\ \relax
26:10 & \begin{tabularx}{0.7\textwidth}{X} 他在水面上劃一圓圈，直到光明與黑暗的交界。 \end{tabularx} \\ \\ \relax
26:11 & \begin{tabularx}{0.7\textwidth}{X} 天的柱子震動，因他的斥責驚奇。 \end{tabularx} \\ \\ \relax
26:12 & \begin{tabularx}{0.7\textwidth}{X} 他以能力攪動大海，藉知識打傷拉哈伯。 \end{tabularx} \\ \\ \relax
26:13 & \begin{tabularx}{0.7\textwidth}{X} 他藉自己的靈使天空晴朗；他的手刺殺爬得快的蛇。 \end{tabularx} \\ \\ \relax
26:14 & \begin{tabularx}{0.7\textwidth}{X} 看哪，這不過是神工作的些微；我們聽見他的話，是何等細微的聲音！他大能的雷聲誰能明白呢？」 \end{tabularx} \\ \\ \relax
27:1 & \begin{tabularx}{0.7\textwidth}{X} 約伯繼續發表他的言論說： \end{tabularx} \\ \\ \relax
27:2 & \begin{tabularx}{0.7\textwidth}{X} 「我指著奪去我公道的永生神，並使我心中愁苦的全能者起誓： \end{tabularx} \\ \\ \relax
27:3 & \begin{tabularx}{0.7\textwidth}{X} 只要我的生命尚在我裡面，神所賜的氣息仍在我鼻孔內， \end{tabularx} \\ \\ \relax
27:4 & \begin{tabularx}{0.7\textwidth}{X} 我的唇絕不說不義，我的舌也不說詭詐。 \end{tabularx} \\ \\ \relax
27:5 & \begin{tabularx}{0.7\textwidth}{X} 我斷不以你們為義；我至死不放棄自己的純正！ \end{tabularx} \\ \\ \relax
27:6 & \begin{tabularx}{0.7\textwidth}{X} 我持定我的義，並不放鬆；在世的日子，我的心不責備我。 \end{tabularx} \\ \\ \relax
27:7 & \begin{tabularx}{0.7\textwidth}{X} 「願我的仇敵如惡人一樣；願那起來攻擊我的，如不義之人一般。 \end{tabularx} \\ \\ \relax
27:8 & \begin{tabularx}{0.7\textwidth}{X} 不敬虔的人有甚麼指望呢？神要剪除他，取他的性命。 \end{tabularx} \\ \\ \relax
27:9 & \begin{tabularx}{0.7\textwidth}{X} 患難臨到他，神豈聽他的呼求？ \end{tabularx} \\ \\ \relax
27:10 & \begin{tabularx}{0.7\textwidth}{X} 他豈以全能者為樂，隨時求告神呢？ \end{tabularx} \\ \\ \relax
27:11 & \begin{tabularx}{0.7\textwidth}{X} 神手所做的，我要指教你們；全能者所行的，我也不會隱瞞。 \end{tabularx} \\ \\ \relax
27:12 & \begin{tabularx}{0.7\textwidth}{X} 看哪，你們自己也都見過，為何全變為這樣虛妄呢？ \end{tabularx} \\ \\ \relax
27:13 & \begin{tabularx}{0.7\textwidth}{X} 「這是神為惡人所定的份，殘暴人從全能者所得的產業： \end{tabularx} \\ \\ \relax
27:14 & \begin{tabularx}{0.7\textwidth}{X} 倘若他的兒女增多，仍被刀所殺；他的子孫必不得飽食。 \end{tabularx} \\ \\ \relax
27:15 & \begin{tabularx}{0.7\textwidth}{X} 他遺留的人必死而埋葬，他的寡婦也不哀哭。 \end{tabularx} \\ \\ \relax
27:16 & \begin{tabularx}{0.7\textwidth}{X} 他雖積蓄銀子如塵沙，堆積衣服如泥土， \end{tabularx} \\ \\ \relax
27:17 & \begin{tabularx}{0.7\textwidth}{X} 他儘管堆積，義人卻要穿上，無辜的人卻要分取銀子。 \end{tabularx} \\ \\ \relax
27:18 & \begin{tabularx}{0.7\textwidth}{X} 他建造房屋如蟲做窩，又如守望者所搭的棚。 \end{tabularx} \\ \\ \relax
27:19 & \begin{tabularx}{0.7\textwidth}{X} 他雖富足躺臥，卻不得收殮，他張開眼睛，就不在了。 \end{tabularx} \\ \\ \relax
27:20 & \begin{tabularx}{0.7\textwidth}{X} 驚恐如洪水將他追上，暴風在夜間將他颳去。 \end{tabularx} \\ \\ \relax
27:21 & \begin{tabularx}{0.7\textwidth}{X} 東風把他吹去，他就走了；風將他颳離原地。 \end{tabularx} \\ \\ \relax
27:22 & \begin{tabularx}{0.7\textwidth}{X} 風無情地擊打他，他試圖逃脫風的手。 \end{tabularx} \\ \\ \relax
27:23 & \begin{tabularx}{0.7\textwidth}{X} 風要因他拍掌，並要發叱聲，使他離開原地。」 \end{tabularx} \\ \\ \relax
28:1 & \begin{tabularx}{0.7\textwidth}{X} 「銀子有礦；煉金有場。 \end{tabularx} \\ \\ \relax
28:2 & \begin{tabularx}{0.7\textwidth}{X} 鐵從土裡開採，銅從礦石鎔出。 \end{tabularx} \\ \\ \relax
28:3 & \begin{tabularx}{0.7\textwidth}{X} 人探索黑暗的盡頭，查究礦石直到極處，那是幽暗和死蔭； \end{tabularx} \\ \\ \relax
28:4 & \begin{tabularx}{0.7\textwidth}{X} 他在無人居住之處開鑿礦穴，在無足跡之地被遺忘，與人遠離，懸空搖擺。 \end{tabularx} \\ \\ \relax
28:5 & \begin{tabularx}{0.7\textwidth}{X} 地出產糧食，地底翻騰如火。 \end{tabularx} \\ \\ \relax
28:6 & \begin{tabularx}{0.7\textwidth}{X} 地的石頭是藍寶石之處，那裡還有金沙。 \end{tabularx} \\ \\ \relax
28:7 & \begin{tabularx}{0.7\textwidth}{X} 鷙鳥不知那條路，鷹眼也未曾見過。 \end{tabularx} \\ \\ \relax
28:8 & \begin{tabularx}{0.7\textwidth}{X} 狂傲的野獸未曾踩踏，猛烈的獅子也未曾經過。 \end{tabularx} \\ \\ \relax
28:9 & \begin{tabularx}{0.7\textwidth}{X} 「人動手鑿開堅石，翻倒山的根基， \end{tabularx} \\ \\ \relax
28:10 & \begin{tabularx}{0.7\textwidth}{X} 在磐石中鑿出水道，親眼看見各樣寶物。 \end{tabularx} \\ \\ \relax
28:11 & \begin{tabularx}{0.7\textwidth}{X} 他封閉河川不得涓滴，使隱藏之物顯露出來。 \end{tabularx} \\ \\ \relax
28:12 & \begin{tabularx}{0.7\textwidth}{X} 「然而，智慧何處可尋？聰明之地在哪裡？ \end{tabularx} \\ \\ \relax
28:13 & \begin{tabularx}{0.7\textwidth}{X} 智慧的價值無人能知，活人之地也無處可尋。 \end{tabularx} \\ \\ \relax
28:14 & \begin{tabularx}{0.7\textwidth}{X} 深淵說：『不在我裡面。』滄海說：『不在我這裡。』 \end{tabularx} \\ \\ \relax
28:15 & \begin{tabularx}{0.7\textwidth}{X} 智慧不可用黃金換取，也不能用白銀秤她的價值。 \end{tabularx} \\ \\ \relax
28:16 & \begin{tabularx}{0.7\textwidth}{X} 俄斐的金子和貴重的紅瑪瑙，以及藍寶石，不足與她比擬； \end{tabularx} \\ \\ \relax
28:17 & \begin{tabularx}{0.7\textwidth}{X} 黃金和玻璃不足與她比較；純金的器皿不足兌換她。 \end{tabularx} \\ \\ \relax
28:18 & \begin{tabularx}{0.7\textwidth}{X} 珊瑚、水晶都不值得提；智慧的價值勝過寶石。 \end{tabularx} \\ \\ \relax
28:19 & \begin{tabularx}{0.7\textwidth}{X} 古實的紅璧璽不足與她比較；純金也不足與她比擬。 \end{tabularx} \\ \\ \relax
28:20 & \begin{tabularx}{0.7\textwidth}{X} 「智慧從何處來呢？聰明之地在哪裡？ \end{tabularx} \\ \\ \relax
28:21 & \begin{tabularx}{0.7\textwidth}{X} 她隱藏，遠離眾生的眼目，她掩蔽，遠離空中的飛鳥。 \end{tabularx} \\ \\ \relax
28:22 & \begin{tabularx}{0.7\textwidth}{X} 毀滅和死亡說：『我們風聞其名。』 \end{tabularx} \\ \\ \relax
28:23 & \begin{tabularx}{0.7\textwidth}{X} 「神明白智慧的道路，知道智慧的所在。 \end{tabularx} \\ \\ \relax
28:24 & \begin{tabularx}{0.7\textwidth}{X} 因為他鑒察直到地極，遍觀普天之下， \end{tabularx} \\ \\ \relax
28:25 & \begin{tabularx}{0.7\textwidth}{X} 要為風定輕重，又度量諸水， \end{tabularx} \\ \\ \relax
28:26 & \begin{tabularx}{0.7\textwidth}{X} 為雨定律例，為雷電定道路。 \end{tabularx} \\ \\ \relax
28:27 & \begin{tabularx}{0.7\textwidth}{X} 那時他看見智慧，就談論她，堅定她，並且查究她。 \end{tabularx} \\ \\ \relax
28:28 & \begin{tabularx}{0.7\textwidth}{X} 他對人說：『看哪，敬畏主就是智慧；遠離惡事就是聰明。』」 \end{tabularx} \\ \\ \relax
29:1 & \begin{tabularx}{0.7\textwidth}{X} 約伯繼續發表他的言論說： \end{tabularx} \\ \\ \relax
29:2 & \begin{tabularx}{0.7\textwidth}{X} 「惟願我如從前的歲月，如神保護我的日子。 \end{tabularx} \\ \\ \relax
29:3 & \begin{tabularx}{0.7\textwidth}{X} 那時他的燈照在我頭上，我藉他的光行過黑暗。 \end{tabularx} \\ \\ \relax
29:4 & \begin{tabularx}{0.7\textwidth}{X} 在我壯年的時候，神親密的情誼臨到我的帳棚中。 \end{tabularx} \\ \\ \relax
29:5 & \begin{tabularx}{0.7\textwidth}{X} 全能者仍與我同在，我的兒女都環繞我。 \end{tabularx} \\ \\ \relax
29:6 & \begin{tabularx}{0.7\textwidth}{X} 我的腳洗在乳酪當中；磐石為我流出油河。 \end{tabularx} \\ \\ \relax
29:7 & \begin{tabularx}{0.7\textwidth}{X} 我出到城門，在廣場安排座位， \end{tabularx} \\ \\ \relax
29:8 & \begin{tabularx}{0.7\textwidth}{X} 年輕人見我而迴避，老年人起身站立。 \end{tabularx} \\ \\ \relax
29:9 & \begin{tabularx}{0.7\textwidth}{X} 王子都停止說話，用手摀口； \end{tabularx} \\ \\ \relax
29:10 & \begin{tabularx}{0.7\textwidth}{X} 領袖靜默無聲，舌頭貼住上膛。 \end{tabularx} \\ \\ \relax
29:11 & \begin{tabularx}{0.7\textwidth}{X} 耳朵聽見了，稱我有福；眼睛看見了，就稱讚我。 \end{tabularx} \\ \\ \relax
29:12 & \begin{tabularx}{0.7\textwidth}{X} 因我拯救了哀求的困苦人和無人幫助的孤兒。 \end{tabularx} \\ \\ \relax
29:13 & \begin{tabularx}{0.7\textwidth}{X} 將要滅亡的為我祝福，我使寡婦心中歡呼。 \end{tabularx} \\ \\ \relax
29:14 & \begin{tabularx}{0.7\textwidth}{X} 我穿上公義，它遮蔽我；我的公平如外袍和冠冕。 \end{tabularx} \\ \\ \relax
29:15 & \begin{tabularx}{0.7\textwidth}{X} 我作瞎子的眼，瘸子的腳。 \end{tabularx} \\ \\ \relax
29:16 & \begin{tabularx}{0.7\textwidth}{X} 我作貧窮人的父；我不認識之人的案件，我也去查明。 \end{tabularx} \\ \\ \relax
29:17 & \begin{tabularx}{0.7\textwidth}{X} 我打破不義之人的大牙，從他牙齒中奪走他所搶的。 \end{tabularx} \\ \\ \relax
29:18 & \begin{tabularx}{0.7\textwidth}{X} 我說：『我要增添我的日子如塵沙，我必死在自己家中。 \end{tabularx} \\ \\ \relax
29:19 & \begin{tabularx}{0.7\textwidth}{X} 我的根伸展到水邊，露水夜宿我的枝上。 \end{tabularx} \\ \\ \relax
29:20 & \begin{tabularx}{0.7\textwidth}{X} 我的榮耀在我身上更新，我的弓在我手中日新。』 \end{tabularx} \\ \\ \relax
29:21 & \begin{tabularx}{0.7\textwidth}{X} 「人聽我說話而等候，為我的教導而靜默。 \end{tabularx} \\ \\ \relax
29:22 & \begin{tabularx}{0.7\textwidth}{X} 我說話之後，他們就不再說；我的言語滴在他們身上。 \end{tabularx} \\ \\ \relax
29:23 & \begin{tabularx}{0.7\textwidth}{X} 他們等候我如等雨水，又張口如切慕春雨。 \end{tabularx} \\ \\ \relax
29:24 & \begin{tabularx}{0.7\textwidth}{X} 我向他們微笑，他們不敢相信；他們不使我臉上的光失色。 \end{tabularx} \\ \\ \relax
29:25 & \begin{tabularx}{0.7\textwidth}{X} 我為他們選擇道路，又坐首位；我如君王在軍隊中居住，又如人安慰哀傷的人。」 \end{tabularx} \\ \\ \relax
30:1 & \begin{tabularx}{0.7\textwidth}{X} 「但如今，比我年輕的人譏笑我；我曾藐視他們的父親，不放在我的牧羊犬中。 \end{tabularx} \\ \\ \relax
30:2 & \begin{tabularx}{0.7\textwidth}{X} 他們的精力既已衰敗，手中的氣力於我何益？ \end{tabularx} \\ \\ \relax
30:3 & \begin{tabularx}{0.7\textwidth}{X} 他們因窮乏飢餓，沒有生氣，在荒廢淒涼的幽暗中啃乾燥之地。 \end{tabularx} \\ \\ \relax
30:4 & \begin{tabularx}{0.7\textwidth}{X} 他們在草叢之中採鹹草，羅騰樹的根成為他們的食物。 \end{tabularx} \\ \\ \relax
30:5 & \begin{tabularx}{0.7\textwidth}{X} 他們從人群中被趕出，人追喊他們如賊一般， \end{tabularx} \\ \\ \relax
30:6 & \begin{tabularx}{0.7\textwidth}{X} 以致他們住在荒谷，住在地洞和巖穴中。 \end{tabularx} \\ \\ \relax
30:7 & \begin{tabularx}{0.7\textwidth}{X} 他們在草叢中叫喚，在荊棘下擠成一團。 \end{tabularx} \\ \\ \relax
30:8 & \begin{tabularx}{0.7\textwidth}{X} 這都是愚頑卑微人的兒女；他們被鞭打，趕出境外。 \end{tabularx} \\ \\ \relax
30:9 & \begin{tabularx}{0.7\textwidth}{X} 「現在這些人以我為歌曲，以我為笑談。 \end{tabularx} \\ \\ \relax
30:10 & \begin{tabularx}{0.7\textwidth}{X} 他們厭惡我，躲避我，不住地吐唾沫在我臉上。 \end{tabularx} \\ \\ \relax
30:11 & \begin{tabularx}{0.7\textwidth}{X} 神鬆開我的弓弦使我受苦，他們就在我面前脫去轡頭。 \end{tabularx} \\ \\ \relax
30:12 & \begin{tabularx}{0.7\textwidth}{X} 這夥人在我右邊起來，他們推開我的腳，築災難之路攻擊我。 \end{tabularx} \\ \\ \relax
30:13 & \begin{tabularx}{0.7\textwidth}{X} 他們毀壞我的道，加增我的災害；他們毋須人幫助。 \end{tabularx} \\ \\ \relax
30:14 & \begin{tabularx}{0.7\textwidth}{X} 他們來，如同闖進大缺口，在暴風間滾動。 \end{tabularx} \\ \\ \relax
30:15 & \begin{tabularx}{0.7\textwidth}{X} 驚恐傾倒在我身上，我的尊榮被逐如風；我的福祿如雲飄去。 \end{tabularx} \\ \\ \relax
30:16 & \begin{tabularx}{0.7\textwidth}{X} 「現在我的心極其悲傷，困苦的日子將我抓住。 \end{tabularx} \\ \\ \relax
30:17 & \begin{tabularx}{0.7\textwidth}{X} 夜間，我裡面的骨頭刺痛，啃著我的沒有止息。 \end{tabularx} \\ \\ \relax
30:18 & \begin{tabularx}{0.7\textwidth}{X} 我的外衣因大力扭皺 ，內衣的領子把我勒住。 \end{tabularx} \\ \\ \relax
30:19 & \begin{tabularx}{0.7\textwidth}{X} 神把我扔在淤泥之中，我就像塵土和灰燼一樣。 \end{tabularx} \\ \\ \relax
30:20 & \begin{tabularx}{0.7\textwidth}{X} 我呼求你，你不應允我；我站起來，你只是望著我。 \end{tabularx} \\ \\ \relax
30:21 & \begin{tabularx}{0.7\textwidth}{X} 你對我變得殘忍，大能的手追逼我。 \end{tabularx} \\ \\ \relax
30:22 & \begin{tabularx}{0.7\textwidth}{X} 你把我提到風中，使我乘風而去，使我消失在烈風之中。 \end{tabularx} \\ \\ \relax
30:23 & \begin{tabularx}{0.7\textwidth}{X} 我知道你要使我歸於死亡，到那為眾生所定的陰宅。 \end{tabularx} \\ \\ \relax
30:24 & \begin{tabularx}{0.7\textwidth}{X} 「然而，人在廢墟豈不伸手？遇災難時一定呼救。 \end{tabularx} \\ \\ \relax
30:25 & \begin{tabularx}{0.7\textwidth}{X} 人遭難的日子，我豈不為他哭泣呢？人貧窮的時候，我豈不為他憂愁呢？ \end{tabularx} \\ \\ \relax
30:26 & \begin{tabularx}{0.7\textwidth}{X} 我仰望福氣，災禍就來到；我等待光明，黑暗便來臨。 \end{tabularx} \\ \\ \relax
30:27 & \begin{tabularx}{0.7\textwidth}{X} 我內心煩擾不安，困苦的日子臨到我身。 \end{tabularx} \\ \\ \relax
30:28 & \begin{tabularx}{0.7\textwidth}{X} 我在陰暗中行走，沒有日光，我在會眾中站立求救。 \end{tabularx} \\ \\ \relax
30:29 & \begin{tabularx}{0.7\textwidth}{X} 我與野狗為弟兄，我跟鴕鳥為同伴。 \end{tabularx} \\ \\ \relax
30:30 & \begin{tabularx}{0.7\textwidth}{X} 我的皮膚變黑脫落，我的骨頭因熱燒焦。 \end{tabularx} \\ \\ \relax
30:31 & \begin{tabularx}{0.7\textwidth}{X} 我的琴音變為哀泣；我的簫聲變為哭聲。」 \end{tabularx} \\ \\ \relax
31:1 & \begin{tabularx}{0.7\textwidth}{X} 「我與眼睛立約，怎能凝望少女呢？ \end{tabularx} \\ \\ \relax
31:2 & \begin{tabularx}{0.7\textwidth}{X} 從至上的神所得之分，從至高全能者所得之業是甚麼呢？ \end{tabularx} \\ \\ \relax
31:3 & \begin{tabularx}{0.7\textwidth}{X} 豈不是禍患臨到不義的，災害臨到作惡的嗎？ \end{tabularx} \\ \\ \relax
31:4 & \begin{tabularx}{0.7\textwidth}{X} 神豈不察看我的道路，數點我所有的腳步嗎？ \end{tabularx} \\ \\ \relax
31:5 & \begin{tabularx}{0.7\textwidth}{X} 「我若與虛謊同行，我腳若緊跟詭詐， \end{tabularx} \\ \\ \relax
31:6 & \begin{tabularx}{0.7\textwidth}{X} 願神用公道的天平秤我，願他知道我的純正。 \end{tabularx} \\ \\ \relax
31:7 & \begin{tabularx}{0.7\textwidth}{X} 我的腳步若偏離正路，我的心若隨從我眼目，我的手掌若黏有污穢； \end{tabularx} \\ \\ \relax
31:8 & \begin{tabularx}{0.7\textwidth}{X} 願我栽種，別人來吃，我的農作物連根拔出。 \end{tabularx} \\ \\ \relax
31:9 & \begin{tabularx}{0.7\textwidth}{X} 「我心若因婦人受迷惑，在鄰舍的門外等候， \end{tabularx} \\ \\ \relax
31:10 & \begin{tabularx}{0.7\textwidth}{X} 就願我妻子給別人推磨，別人與她同寢。 \end{tabularx} \\ \\ \relax
31:11 & \begin{tabularx}{0.7\textwidth}{X} 因為這是邪惡的事，審判官裁定的罪孽。 \end{tabularx} \\ \\ \relax
31:12 & \begin{tabularx}{0.7\textwidth}{X} 這是一場火，直燒到毀滅，必拔除我一切的家產。 \end{tabularx} \\ \\ \relax
31:13 & \begin{tabularx}{0.7\textwidth}{X} 「我的僕婢與我爭辯，我若藐視不聽他們的冤情， \end{tabularx} \\ \\ \relax
31:14 & \begin{tabularx}{0.7\textwidth}{X} 神興起的時候，我怎樣行呢？他察問的時候，我怎樣回答他呢？ \end{tabularx} \\ \\ \relax
31:15 & \begin{tabularx}{0.7\textwidth}{X} 造我在母腹中的，不也是造了他嗎？在母胎中使我們成形的，豈不是同一位嗎？ \end{tabularx} \\ \\ \relax
31:16 & \begin{tabularx}{0.7\textwidth}{X} 「我若不讓貧寒人遂其所願，或是叫寡婦眼中失望， \end{tabularx} \\ \\ \relax
31:17 & \begin{tabularx}{0.7\textwidth}{X} 或獨自吃自己的食物，孤兒沒有吃其中些許； \end{tabularx} \\ \\ \relax
31:18 & \begin{tabularx}{0.7\textwidth}{X} 從我年輕時，孤兒就與我一同長大，我好像他的父親，我從出母腹就扶助寡婦； \end{tabularx} \\ \\ \relax
31:19 & \begin{tabularx}{0.7\textwidth}{X} 我若見人因無衣死亡，或見貧窮人毫無遮蓋； \end{tabularx} \\ \\ \relax
31:20 & \begin{tabularx}{0.7\textwidth}{X} 我若不使他真心為我祝福，不使他因我羊的毛得暖； \end{tabularx} \\ \\
[1ex]
\hline
\hline
\end{longtable}
我昨天已經介紹約伯三位朋友之中以利法,他是其中最年長,富有同情心,最有屬靈經驗和知識的.他發言開始時是提醒和勸勉,繼之以自己屬靈經歷和知識幫助；他自以為對約伯的勸勉是毫無疑問的,是大有權威的,是約伯不應抗拒的.在(五27)所說的意思,命令約伯須聽他的話,才能得益,不然將受虧損；他自認為真理之言,殊知約伯卻不接受.因此,當他第二次發言時,同情心就減少並轉為嚴厲；並且他重申自己屬靈的經歷.(十五14)是他在第四章所見異象的解釋,再次提醒約伯.
\newline
\newline
有些人屬靈寶貴的經歷和見證多年如一.其實見證不應停留,經歷應該常新.
\newline
\newline
以利法第二次說話時,對約伯有更厲害的責備.他根據他在社會的經歷,(十五17)曾兩次說話,亦不外乎本身經歷和社會見聞,尤以第二次更加嚴重.他說犯罪者的一切,必轉眼成空.(20節)人生路上充滿不安和懼怕,罪惡影響人的心靈和對神的態度.從上述的話語中,足見以利法是個靈程很高深的人.罪惡帶來的痛苦不單是物質方面,也關乎到心靈和態度方面.以利法有相當深入的看見!22節說有罪的人,雖然祈禱,聽道都不適合自己,不信自己能從苦境轉回,罪影響人心靈的力量.以利法看見約伯不斷自艾自怨,訴說痛苦；認為是有罪的現象.他也認為罪影響人對神的態度頑強背逆；痛苦的遭遇對自己軟弱站不起來.我們常見罪大惡極的人最反對神,不怕神,不願悔改.(24-27)以利法把約伯看作沒有希望的人,不聽勸導,反對他就是反對神.這裡我們看到,以利法的判斷雖然錯誤,但道理卻是對的；這是他較比勒達和瑣法高明之處；他們二人雖也講及痛苦與罪有關；但觀念庸俗,只見肉身的痛苦,錢財的損失和家庭的災難,均屬顯而易見的.他們未嘗強調犯罪所受敗壞的影響.
\newline
\newline
以利法第二次勸勉約伯的話,更加嚴厲和深入；希望用濃藥治重病,就知約伯仍無動於中；且越聽越怒.第三次的談話記在廿二章,他至此大怒,溫柔態度盡失,話中開始暗示約伯有罪,說約伯有罪是件事實,現在更進一步明指約伯犯罪.(6-11)以利法作假見證,還以為自己是為真理辯護；說約伯犯貪得無厭的罪,剝奪窮人；(6-7)包庇有權勢的人(8)苦待孤兒寡婦.(9)假冒偽善,遮蓋自己的敗壞.(13-14)以上諸罪可說是富人常犯的,不擇手段,惟利是圖.為爭取地位,專與社會名流過從,冀達名利,地位,權勢全勝的慾望；運用權力壓迫無辜,剝削,虛偽.以利法想像中的約伯正是如此,成見使他把約伯看作個罪大惡極的人；所以他對約伯說:「因此有網羅環繞你,有恐懼忽然使你驚惶；或有黑暗蒙蔽你,並有洪水淹沒你.」(10-11)那時以利法對約伯的態度和初來時完全不同了.我們讀至此,可能有人想,以利法的言談錯誤.其實他的道理是對的,只是應用不當；所以神仍然把它記在聖經.倘若有人在你受苦時,用以利法的話來幫助你,也許可以適用；因為像約伯純正的人極少.以利法的話記在聖經,一方面叫我們勸慰人之時,務必小心,雖是真理良言,也必須適合時宜.另一方面,用屬靈真理的話來光照自己,引導我們避免因犯罪而受苦,得真理之前,從罪的痛苦中解脫.
\newline
\newline
讀約伯記務求細心研讀,不可跳讀；必須連貫細讀才不致錯解其意義.
\newline
\newline
以利法在廿二章說了許多又嚴厲又惡的錯話之後,覺得對待約伯未免太過份；所以最後一段語氣急轉溫柔,述說識主之道.(21-30)言語美麗,極可安慰人；論及認識神的福樂,(21)大有平安.(23)與神建立穩固的人生.(27)禱告蒙垂聽,得神成就凡他所作盡都順利.(28)路上有光,神要把他升高.(29)蒙神拯救.(30)第二方面論及認識神的方法:(1)必須領受神言.(22)從神的話開始,且將神的話藏在心裡.(2)認罪悔改.(23)(3)渴慕追求神的話,有飢渴慕義的心.(24-26)(4)在神前禱告尋求.
\newline
\newline
聖經多次論認識神的途徑,這位屬靈的長者以利法,也提出認識神的方法.第二個安慰使者比勒達也同樣地失敗,他的同情心不及以利法；他一發言就說約伯的家屬,兒女有罪,所以致死.(八1-7)他形容約伯是個忘記神的人,是不虔敬(應譯為假冒偽善)；這樣的人要滅亡.他也明指約伯有罪,與神之間關係有問題,所以受苦.比勒達所根據的是古人的學說,可稱他為學院派的大師.第八章(8-10節)顯示他在真理方面,比古代的人研究得清楚,他是個崇古保守的人.保守不一定是純正,保守是接受苦代一切的傳統,過於保守崇古的人常鄙視今世的人膚淺無知.他說:「我們不過從昨天才有,一無所知；我們在世的日子好像影兒.」(9)看今世人生如風,沒有實際的學問在其中；他強調自己的學問最高,是傳統學者中研究最深的人.有的人說話,常引用孔子的話,表示自己學問淵博.比勒達述說古人之話的教訓,說人和神的關係正如蒲草,蘆荻和泥水之關係(11)照古人教訓,人一切禍福都是和神關係的結果,人和神的關係是決定一切事物的因素.以此推論,他認定約伯的根基有問題.比勒達這段話給我們很大的教訓,神是我們一切的根基；無神的人生就沒有根基,行為也沒目標,在永世裡一切皆空.
\newline
\newline
為甚麼所羅門王也說:「虛空的虛空」呢?有人認為他是個花花公子；其實他卻是個很有上進心的人.他大有學問,對國家多有建設；治國有條有理,除了妃嬪太多之外,他樣樣都好；為何他還說虛空的虛空?因為他說:當他年青時忘記了造他的主.沒有根基,所以沒有實際永恆的價值.為何許多人的努力成空?因他不知道真正好的和壞的是甚麼.人都看自己最好,各持己見,以為自己的定義才是對的.其實,凡世上人的定義都不對；惟神的定義才是實在的根基；向著這根基去努力才好.神是我們的好根基；向著這根基也是我們最大的倚靠.在家倚靠父母,在社會上依靠朋友,作工依靠自己；但是這些有時也靠不住.(14-15)唯獨信神,我們找到了奇妙的依靠；因神是全能的神.大衛說:「我依靠神必不懼怕,依靠耶和華強似依賴王子.」倚靠耶和華的人一無所缺,軟弱成為剛強.神是我們最強有力的倚靠,祂愛我們；所以真正信主的人不易灰心,不易絕望；祂成就超乎我們力所能及的事.有神才有方向,才有真正長久的結局；依靠神,則人生好像黎明的光,愈照愈明直至日午.大衛在詩篇廿三篇最後說:「我一生一世,必有恩惠慈愛隨著我；」他回憶已往神的恩惠和慈愛,成為他的依靠和幫助.指望將來,他說:「我且要住在耶和華的殿中,直到永遠.」他以耶和華家中為他永遠的歸宿.「永遠」對他並非模糊可怕的.
\newline
\newline
比勒達對神有熟悉的知識.他說,凡忘記神的人,好像一陣風吹過,不留痕蹟(18)(羅十一36)也論到我們和神的關係,我們本於祂,祂是我們的根基；依靠祂,祂是我們人生的保障；歸於祂,祂是我們人生的結局；所以,願榮耀歸給祂.
\newline
\newline
比勒達用古人研究屬靈的知識幫助約伯,這段話實在非常美麗；可惜他先決的判斷錯誤,致不受歡迎.他第二次說話,記在十八章,較第八章為短,同時也沒新亮光；因古人的道理已於第八章說盡了,他無以幫助約伯；所以再次引用舊的亮光,把人犯罪的苦果說得更加可怕以恐嚇約伯.他說惡人的苦難難以想像,惡人無論在家,在外；到處有災難,日夜不平安,得怪病,絕子絕孫.(絕後是猶太人最大的咒詛)比勒達此行目的,到底是同情或是無情對待呢?他初來時所懷感情何在?我們從中學習屬靈境界的功課.真理可使人同情,但若錯用真理,就可能成為最無情的人.約伯處於嚴重的創傷中,痛苦難當；所謂屬靈人比勒達如此無情對待,指責約伯罪有應得.
\newline
\newline
求主教我們,曉得自己對真理的領受有限,對人的事,必須祈禱求神給我們客觀的態度去觀察,求神賜給適當的真理去幫助人.
\newline
\newline
人是運用表面的真理,卻反使人遭受其害.
\newline
\newline
當你有此經歷時,勿以為沒有甚麼意義；藉此,或者從神得到真理的新亮光,使你有更好的用處.
\newline
\newline
相信經此一番對約伯的幫助,以利法,比勒達和瑣法三人,在屬靈的見識和助人態度的方法上,得到更大的祝福.我們也學習了許多的功課.
\newline
\newline


\newpage

\section{第52屆~港九培靈會~焦源濂~第七講　失敗的安慰者-瑣法}
\label{sec:626}
\textbf{第七講　失敗的安慰者-瑣法}
\newline
\newline
連結: \href{https://www.hkbibleconference.org/session-message/view/626}{\texttt{https://www.hkbibleconference.org/session-message/view/626}}
\newline
\newline
\hyperref[sec:624]{< < < PREV SERMON < < <}
~
\hyperref[sec:toc]{[返目錄]}
~
\hyperref[sec:628]{> > > NEXT SERMON > > >}
\newline
\newline
約伯記 11:1-11:12
\newline
\begin{longtable}{cl}
\hline
\hline
章節 & 經文 (和合本修訂版)\\
\hline
11:1 & \begin{tabularx}{0.7\textwidth}{X} 拿瑪人瑣法回答說： \end{tabularx} \\ \\ \relax
11:2 & \begin{tabularx}{0.7\textwidth}{X} 「這許多的話豈不該回答嗎？多嘴多舌的人豈可成為義呢？ \end{tabularx} \\ \\ \relax
11:3 & \begin{tabularx}{0.7\textwidth}{X} 你誇大的話豈能使人不作聲嗎？你戲笑的時候豈沒有人使你受辱嗎？ \end{tabularx} \\ \\ \relax
11:4 & \begin{tabularx}{0.7\textwidth}{X} 你說：『我的教導純全，我在你眼前是清潔的。』 \end{tabularx} \\ \\ \relax
11:5 & \begin{tabularx}{0.7\textwidth}{X} 但是，惟願神說話，願他張開嘴唇攻擊你。 \end{tabularx} \\ \\ \relax
11:6 & \begin{tabularx}{0.7\textwidth}{X} 願他將智慧的奧祕指示你，因為健全的知識是兩面的。你當知道，神使你忘記你的一些罪孽。 \end{tabularx} \\ \\ \relax
11:7 & \begin{tabularx}{0.7\textwidth}{X} 你能尋見神的奧祕嗎？你能尋見全能者的極限嗎？ \end{tabularx} \\ \\ \relax
11:8 & \begin{tabularx}{0.7\textwidth}{X} 高如諸天，你能做甚麼？比陰間深，你能知道甚麼？ \end{tabularx} \\ \\ \relax
11:9 & \begin{tabularx}{0.7\textwidth}{X} 其量度比地長，比海更寬。 \end{tabularx} \\ \\ \relax
11:10 & \begin{tabularx}{0.7\textwidth}{X} 他若經過，把人拘禁，召集會眾，誰能阻擋他呢？ \end{tabularx} \\ \\ \relax
11:11 & \begin{tabularx}{0.7\textwidth}{X} 因為他知道虛妄的人；當他看見罪惡，豈不留意嗎？ \end{tabularx} \\ \\ \relax
11:12 & \begin{tabularx}{0.7\textwidth}{X} 空虛的人若獲得知識，野驢生下的駒子也成了人。 \end{tabularx} \\ \\
[1ex]
\hline
\hline
\end{longtable}
兩個失敗的安慰者,以利法注重經驗和見識；比勒達根據古人的學說和真理勸導約伯,說人若離棄神必發生問題.這話是不錯的,但卻不能說,沒發生問題的人,就與神有良好的關係.比勒達堅定相信自己的話是絕對的,所以刻薄殘忍的責備約伯.真理若應用錯誤,不但無助於人,反叫人心靈受傷和痛苦.比勒達第二次說的話仍使約伯無法接受；因而他第三次發言很短(載於二十五章).一方面顯出他的話已經說完,另一方面顯示他不甘承認失敗,繼續爭取最後發言權；他強調神有治理的權柄,人無法持強取勝.
\newline
\newline
第三個失敗的安慰使者瑣法　第十一章1-12節顯示瑣法的性格粗鹵爽直,他以衛道英雄的姿態出現,以家長的態度教導約伯；他輕視約伯,指他是個虛妄,蠢笨如野驢駒子的人；他不能忍受約伯之狂言亂語,巴不得神罰他更加痛苦.瑣法是個自命不凡的人,以為自己所說都是真理奧秘之言,視約伯為無知,不可教之材.
\newline
\newline
約伯對以利法的話雖然不滿,表示失望；但他能忍受.約伯對比勒達則有諷刺.對瑣法更毫不客氣地說:「你們所說的,誰不知道呢?」(1-3節)說瑣法所知有限,豈非藉著地和走獸,空中飛鳥,海中的魚所指教,告訴,說明,才能明白嗎?(7-8)約伯把瑣法所說的道理,作進一步的批判:「你們是編造謊言的,都是無用的醫生.惟願你們全然不作聲,這就算為你們的智慧.」(十三4-5)「你們以為可記念的箴言,是爐灰......可靠的堅壘,是淤泥.......」(12)約伯對瑣法的態度非常頑強.瑣法是約伯三朋友中言論實質最短少的,只說了兩次話便無言可說.瑣法見識多,是古時有學問屬靈的長者,他首次發言記於第十一章,可分為兩段,首段論神奇妙的作為.(7-12)第二段論悔改之道,悔改可蒙福(13-20)人悔改必須先預備自己的心,切實向神求告；禱告之後,必須離棄所知的罪,然後在人生中改變行為.這篇悔改之道,因時間關係,未能詳說.總之,和新約論悔改是完全一樣的.
\newline
\newline
瑣法第二篇的說話,記於第二十章,也可分為兩段,論惡人的報應.約伯和他辯論,約伯說義人受苦,惡人有時也享福.瑣法說惡人享福是非常短暫的.(二十1-19)繼續,他又說惡人必定有報應,而且來得突然.(20-29)按字面說,惡有惡報,善有善報；這話是不錯的,在人看來,這是世界普通的現象.因果律和屬靈的律有關,但世上有很多事實,是這律所不能解釋的；這律並非絕對的律,乃神所立道德之律；凡律是永不改變的.但在現今邪惡的世代,關於這律何時絕不改變,卻得不到最完滿的答案.不過等主再來時在永世裡,才可以見到這律乃是真正的律,惡人必定受苦,善人必定蒙福.現在是恩典時代,神使日頭照好人,也使日頭照歹人.罪人仍有自由的時間；如果神今日執行這律,那麼,一切罪人都要滅亡.瑣法誤解了這律,以為這樣在今世可解答人一切的問題.
\newline
\newline
以利法認為神對罪人報應單在物質方面.瑣法屬靈的深度不及以利法.其實,最大的咒詛不止於物質方面；而罪人因罪致影響性格,受到損害,才是最大的刑罰.
\newline
\newline
我們從瑣法得到教訓,千萬不要憑眼見判斷人,憑成敗論英雄.在屬靈方面,家庭興旺,便以為靈性好,神施恩；世事和屬靈事上很多時候可說糢糊不清.希伯來書第十一章,論古時信心的英雄；我們讀後覺信心真是了不得!「因信制服了敵國,堵住獅子的口,滅了烈火的猛勢,脫了刀劍的鋒刃.」(來十一33-34)但是；不可忘記,他們因著信,「有人忍受嚴刑,不願苟且得釋放......,又有人忍受戲弄,鞭打,捆鎖,監禁,各等的磨練,被石頭打死,被鋸鋸死,受試探,被刀殺......受窮乏,患難,苦害 ; 在曠野,山洞,地穴飄流無定.」(來十一35-38)在人看來,或以為他們的信心不足；所以受苦致死；但是,在神眼中,他們確是信心英雄.
\newline
\newline
以世事來說,夫婦感情破裂,離婚,他們彼此都覺痛苦；他們受苦乃因有錯,犯罪；但是他們的兒女卻無辜受苦.常有人問我,世上因戰爭造成許多孤兒寡婦,許多人飄流無定.為何神要叫這樣的事發生呢?今天的世界常因人犯罪而帶來苦難；人常以此歸咎於神.如今還是神容忍的時候,必然還有罪惡,所以必須禱告,求主快來,求主多加能力,叫我們引人歸主；讓無辜受苦的人得以減少,脫離罪惡；信主之後,自有平安.家庭方面也是如此.
\newline
\newline
關於約伯三個朋友,因時間關係,只能如上作簡介.
\newline
\newline
現在我們來看約伯,對於三個朋友勸勉的反應.歸納而言,從約伯十次說話,可知他對三朋友的反駁和辯論.請弟兄姊妹注意!約伯的話,可表明無辜受苦人的心態,他裡面的思想,對於因受苦在他身上所產生,靈裡的光表達出來；所以許多人在受苦時愛讀約伯記,覺得約伯是最同情自己的人；自己有苦難言的話,約伯都代為表達了.茲分數方面來介紹:
\newline
\newline
(一)約伯受苦時對於人生的觀感　人受苦時,常想到人生問題,平安幸福之時則不然.受苦時感嘆人生如夢,灰暗悲觀,毫無意義.(請參看七1-10；九23-37；10:18-22；9全；十四全)約伯那時說,他的人生如同戰場,你爭我奪,不得安息；人生似牛馬,充滿勞苦血汗；人生如織布,日月如梭,生命不過一口氣；人生如煙雲,轉眼消散.他感慨生命之弱,人生之凄涼；他又形容人生充滿災難,如空虛的夢,深感人生沒意義.
\newline
\newline
人若不怕死,可算為勇敢；不怕受苦而活,更為勇敢.當我們受苦時,就逐漸忘記人生另有一精彩的形態；基督徒的悲觀,是個屬靈的問題,因為他榮耀神的人生觀未曾建立得完全.人多受苦時對社會的觀察,非常敏感,不滿現實,當人在苦難中,平時毫無問題的事,也成為問題重重.
\newline
\newline
(二)約伯受苦時,對社會充滿嚕囌疑問　他說,人類歷史充滿疑問,很多問題都難尋答案；人類歷史混亂,沒有是非道德的原則.(二十一1-16):社會是惡人享福,義人受苦的社會.正與瑣法之「惡人必然受苦」的論調相反.第二十四章約伯論混亂的世界在(十二13-25)首先論神在大自然的能力,無人能阻；但他看不出全能者無所不能的神,對人類社會作事清楚的原則.在這段經文,他論到各種不同的人,也論到國家興衰所有的變化.他看不出有何道德原則,又說到神對待人有各種不同方式.有的聰明人忽變為愚蠢,有能力的突然失敗；並非他們犯了罪而遭變化,也非因愛神而突然興旺.約伯對於這些問題,始終不明白,都是在他未受苦前從沒想到的.
\newline
\newline
在二十四章約伯看到社會黑暗的一面,(5-8)說被壓迫者飄流無定,受苦可憐的情形.(9-17)說壓迫者強暴橫行,所以社會充滿悲哀.約伯埋怨為何袖手旁觀,置之不理.(18-25)
\newline
\newline
(三)約伯對神的反應　約伯受苦時沒有失去信心,仍以神為獨一無二,奇妙大能的真神.他說神創造天地萬物,行大事不可測度,行奇事不可勝數.(九8-12)這話以利法在五章9節曾經說過.但是約伯不明白,神為何要對付他至此地步.(七17-20)說,偉大行奇事的神,為何在苦難上看中了他:他覺得神對他真是小題大作,絲毫不願放過他.
\newline
\newline
記得我小的時候,有一次家母問我們幾個兄妹,長大了作甚麼?她對大哥二哥和妹妹各人的志向都感欣慰；惟有我說長大要打魚,令她非常失望難過.這是出於愛.
\newline
\newline
人在神的心目中何等寶貴!「人算得甚麼,你竟顧念.」奇妙的愛,未有十字架之先,神己經愛我們；十字架是愛最高最明顯的標誌.
\newline
\newline
另一方面,約伯對神的性格,覺得很難理解,(九13-24)約伯大膽訴說:神對他發怒,蠻不講理,不垂聽他的禱告,又殘忍無情,持勢欺人,獨斷獨行,善惡不分,幸災樂禍,放縱罪人.(十六6-14)約伯更把自身痛苦的原因歸咎於神,埋怨神對他不公平；使親友都遠離他；興起惡人聚會打擊他,且把他當作箭靶.三個朋友如箭手四面圍攻他,神親自攻擊他.怒目相看,向他咬牙切齒,逼迫他,搯住他的頸項,撕裂他的肺腑,傾倒他的膽,將他破裂又破裂.雖然約伯對神諸多指摘,但神卻不伸手打他:可見神之寬大,民主.無所不知的神深明約伯的內心,約伯雖滿口埋怨神,但卻緊纏著神.
\newline
\newline
人生最大的痛苦莫如被棄絕.舊約該隱殺了兄弟；神責備,他仍不怕；後來神把他逐出人群,不能見神面.於是他懼怕地說:「神啊!你使我不能見你的面.」主耶穌在十字架上,忍受一切的羞辱,門徒的背叛；祂都默然無聲；但祂卻呼喊:「我的神,我的神,為甚麼離棄我?」
\newline
\newline
約伯不單肉身的痛苦,他認為那些算不得甚麼,最大的痛苦乃神向他隱藏,棄絕他；這叫他擔當不住；他極力要尋找原因和答案,他發出被棄絕的哀聲；他的話無形中表達了神的愛．(七7-21)他兩次向神撒嬌,情感自然流露.真是個愛神瞭解神心意的人!他懷著極大的盼望,在神面前多次要求,一再辯論.
\newline
\newline


\newpage

\section{第52屆~港九培靈會~焦源濂~第八講　疑團盡消}
\label{sec:628}
\textbf{皆大歡喜}
\newline
\newline
連結: \href{https://www.hkbibleconference.org/session-message/view/628}{\texttt{https://www.hkbibleconference.org/session-message/view/628}}
\newline
\newline
\hyperref[sec:626]{< < < PREV SERMON < < <}
~
\hyperref[sec:toc]{[返目錄]}
~
\hyperref[sec:632]{> > > NEXT SERMON > > >}
\newline
\newline
約伯記 38:1-42:17
\newline
\begin{longtable}{cl}
\hline
\hline
章節 & 經文 (和合本修訂版)\\
\hline
38:1 & \begin{tabularx}{0.7\textwidth}{X} 那時，耶和華從旋風中回答約伯說： \end{tabularx} \\ \\ \relax
38:2 & \begin{tabularx}{0.7\textwidth}{X} 「誰用無知的言語使我的旨意暗昧不明？ \end{tabularx} \\ \\ \relax
38:3 & \begin{tabularx}{0.7\textwidth}{X} 你要如勇士束腰；我問你，你可以讓我知道。 \end{tabularx} \\ \\ \relax
38:4 & \begin{tabularx}{0.7\textwidth}{X} 「我立大地根基的時候，你在哪裡？你若明白事理，只管說吧！ \end{tabularx} \\ \\ \relax
38:5 & \begin{tabularx}{0.7\textwidth}{X} 你知道是誰定地的尺度，是誰把準繩拉在其上嗎？ \end{tabularx} \\ \\ \relax
38:6 & \begin{tabularx}{0.7\textwidth}{X} 地的根基安置在何處？地的角石是誰安放的？ \end{tabularx} \\ \\ \relax
38:7 & \begin{tabularx}{0.7\textwidth}{X} 那時，晨星一同歌唱；神的眾使者也都歡呼。 \end{tabularx} \\ \\ \relax
38:8 & \begin{tabularx}{0.7\textwidth}{X} 「當海水衝出，如出母胎，誰用門將它關閉呢？ \end{tabularx} \\ \\ \relax
38:9 & \begin{tabularx}{0.7\textwidth}{X} 是我用雲彩當海的衣服，用幽暗當包裹它的布， \end{tabularx} \\ \\ \relax
38:10 & \begin{tabularx}{0.7\textwidth}{X} 為它定界限，又安門和閂， \end{tabularx} \\ \\ \relax
38:11 & \begin{tabularx}{0.7\textwidth}{X} 說：『你只可到這裡，不可越過；你狂傲的浪要到此止住。』 \end{tabularx} \\ \\ \relax
38:12 & \begin{tabularx}{0.7\textwidth}{X} 「你有生以來，曾命定晨光，曾使黎明知道自己的地位， \end{tabularx} \\ \\ \relax
38:13 & \begin{tabularx}{0.7\textwidth}{X} 抓住地的四極，把惡人從其中驅逐出來嗎？ \end{tabularx} \\ \\ \relax
38:14 & \begin{tabularx}{0.7\textwidth}{X} 地改變如泥上蓋印，萬物出現如衣服一樣。 \end{tabularx} \\ \\ \relax
38:15 & \begin{tabularx}{0.7\textwidth}{X} 亮光不照惡人，高舉的膀臂也必折斷。 \end{tabularx} \\ \\ \relax
38:16 & \begin{tabularx}{0.7\textwidth}{X} 「你曾進到海之源，或在深淵的隱密處行走嗎？ \end{tabularx} \\ \\ \relax
38:17 & \begin{tabularx}{0.7\textwidth}{X} 死亡的門曾向你顯露嗎？死蔭的門你曾見過嗎？ \end{tabularx} \\ \\ \relax
38:18 & \begin{tabularx}{0.7\textwidth}{X} 地的廣大，你能測透嗎？你若全知道，只管說吧！ \end{tabularx} \\ \\ \relax
38:19 & \begin{tabularx}{0.7\textwidth}{X} 「往光明居所的路在哪裡？黑暗的地方在何處？ \end{tabularx} \\ \\ \relax
38:20 & \begin{tabularx}{0.7\textwidth}{X} 你能將它帶到其領域，能辨明其居所之路嗎？ \end{tabularx} \\ \\ \relax
38:21 & \begin{tabularx}{0.7\textwidth}{X} 你知道的，因為那時你已出生，你活的日子數目也多。 \end{tabularx} \\ \\ \relax
38:22 & \begin{tabularx}{0.7\textwidth}{X} 「你曾進入雪之庫，或見過雹的倉嗎？ \end{tabularx} \\ \\ \relax
38:23 & \begin{tabularx}{0.7\textwidth}{X} 雪雹是我為災難的時候，為打仗和戰爭的日子所預備。 \end{tabularx} \\ \\ \relax
38:24 & \begin{tabularx}{0.7\textwidth}{X} 光亮從何路分開？東風從何路分散遍地？ \end{tabularx} \\ \\ \relax
38:25 & \begin{tabularx}{0.7\textwidth}{X} 「誰為大雨分道，誰為雷電開路， \end{tabularx} \\ \\ \relax
38:26 & \begin{tabularx}{0.7\textwidth}{X} 使雨降在無人之地，在無人居住的曠野， \end{tabularx} \\ \\ \relax
38:27 & \begin{tabularx}{0.7\textwidth}{X} 使荒廢淒涼之地得以豐足，青草得以生長？ \end{tabularx} \\ \\ \relax
38:28 & \begin{tabularx}{0.7\textwidth}{X} 「雨有父親嗎？露珠是誰生的呢？ \end{tabularx} \\ \\ \relax
38:29 & \begin{tabularx}{0.7\textwidth}{X} 冰出於誰的胎？天上的霜是誰生的呢？ \end{tabularx} \\ \\ \relax
38:30 & \begin{tabularx}{0.7\textwidth}{X} 諸水堅硬如石頭，深淵之面凝結成冰。 \end{tabularx} \\ \\ \relax
38:31 & \begin{tabularx}{0.7\textwidth}{X} 「你能為昴星繫結嗎？你能為參星解帶嗎？ \end{tabularx} \\ \\ \relax
38:32 & \begin{tabularx}{0.7\textwidth}{X} 你能按時領出星宿嗎？能引導北斗與其眾星嗎？ \end{tabularx} \\ \\ \relax
38:33 & \begin{tabularx}{0.7\textwidth}{X} 你知道天的定律嗎？你能使地歸其權下嗎？ \end{tabularx} \\ \\ \relax
38:34 & \begin{tabularx}{0.7\textwidth}{X} 「你能向密雲揚起聲來，使傾盆的雨遮蓋你嗎？ \end{tabularx} \\ \\ \relax
38:35 & \begin{tabularx}{0.7\textwidth}{X} 你能發出閃電，使它們行走，並對你說『我們在這裡』嗎？ \end{tabularx} \\ \\ \relax
38:36 & \begin{tabularx}{0.7\textwidth}{X} 誰將智慧放在朱鷺中？誰將聰明賜給雄雞？ \end{tabularx} \\ \\ \relax
38:37 & \begin{tabularx}{0.7\textwidth}{X} 誰能用智慧數算雲彩？誰能傾倒天上的瓶呢？ \end{tabularx} \\ \\ \relax
38:38 & \begin{tabularx}{0.7\textwidth}{X} 那時，塵土聚集成團，土塊緊緊結連。 \end{tabularx} \\ \\ \relax
38:39 & \begin{tabularx}{0.7\textwidth}{X} 「你能為母獅抓取獵物，使少壯的獅子飽足嗎？ \end{tabularx} \\ \\ \relax
38:40 & \begin{tabularx}{0.7\textwidth}{X} 那時，牠們在洞中蹲伏，在隱密處埋伏。 \end{tabularx} \\ \\ \relax
38:41 & \begin{tabularx}{0.7\textwidth}{X} 誰能為烏鴉預備食物呢？那時，烏鴉之雛哀求神，因無食物飛來飛去。」 \end{tabularx} \\ \\ \relax
39:1 & \begin{tabularx}{0.7\textwidth}{X} 「你知道巖石間的野山羊幾時生產嗎？你能觀察母鹿下小鹿嗎？ \end{tabularx} \\ \\ \relax
39:2 & \begin{tabularx}{0.7\textwidth}{X} 你能數算牠們懷胎的月數嗎？你知道牠們幾時生產嗎？ \end{tabularx} \\ \\ \relax
39:3 & \begin{tabularx}{0.7\textwidth}{X} 牠們屈身，生下幼兒，就解除了陣痛。 \end{tabularx} \\ \\ \relax
39:4 & \begin{tabularx}{0.7\textwidth}{X} 其子漸漸肥壯，在荒野長大；牠們出去，不再歸回。 \end{tabularx} \\ \\ \relax
39:5 & \begin{tabularx}{0.7\textwidth}{X} 「誰放野驢自由？誰解開快驢的繩索？ \end{tabularx} \\ \\ \relax
39:6 & \begin{tabularx}{0.7\textwidth}{X} 我使曠野作牠的住處，使鹽地當牠的居所。 \end{tabularx} \\ \\ \relax
39:7 & \begin{tabularx}{0.7\textwidth}{X} 牠嘲笑城內的喧嚷，不聽趕牲口的喝聲。 \end{tabularx} \\ \\ \relax
39:8 & \begin{tabularx}{0.7\textwidth}{X} 諸山是牠漫遊的草場，牠尋找各樣青綠之物。 \end{tabularx} \\ \\ \relax
39:9 & \begin{tabularx}{0.7\textwidth}{X} 「野牛豈肯服事你？豈肯在你的槽旁過夜？ \end{tabularx} \\ \\ \relax
39:10 & \begin{tabularx}{0.7\textwidth}{X} 你豈能用套繩將野牛繫於犁溝？牠豈肯隨你耙鬆山谷之地？ \end{tabularx} \\ \\ \relax
39:11 & \begin{tabularx}{0.7\textwidth}{X} 你豈可因牠力大就倚靠牠？豈可把你的工交給牠做呢？ \end{tabularx} \\ \\ \relax
39:12 & \begin{tabularx}{0.7\textwidth}{X} 你豈能靠牠把你的穀物運回，又收聚在你的禾場上嗎？ \end{tabularx} \\ \\ \relax
39:13 & \begin{tabularx}{0.7\textwidth}{X} 「鴕鳥的翅膀歡然拍動，但豈是鸛的翎毛和羽毛嗎？ \end{tabularx} \\ \\ \relax
39:14 & \begin{tabularx}{0.7\textwidth}{X} 因牠把蛋留在地上，使蛋在塵土中得溫暖， \end{tabularx} \\ \\ \relax
39:15 & \begin{tabularx}{0.7\textwidth}{X} 卻忘記腳會把蛋踹碎，野獸會踐踏它。 \end{tabularx} \\ \\ \relax
39:16 & \begin{tabularx}{0.7\textwidth}{X} 牠粗暴待雛，似乎不是自己生的；雖徒然勞苦，也不懼怕。 \end{tabularx} \\ \\ \relax
39:17 & \begin{tabularx}{0.7\textwidth}{X} 因為神使牠忘記智慧，也未將悟性分給牠。 \end{tabularx} \\ \\ \relax
39:18 & \begin{tabularx}{0.7\textwidth}{X} 牠幾時挺身展開翅膀，就嘲笑馬和騎馬的人。 \end{tabularx} \\ \\ \relax
39:19 & \begin{tabularx}{0.7\textwidth}{X} 「馬的力量是你所賜的嗎？牠頸項上的鬃是你披上的嗎？ \end{tabularx} \\ \\ \relax
39:20 & \begin{tabularx}{0.7\textwidth}{X} 是你叫牠跳躍像蝗蟲嗎？牠噴氣之威嚴使人驚惶。 \end{tabularx} \\ \\ \relax
39:21 & \begin{tabularx}{0.7\textwidth}{X} 牠用蹄在谷中挖地，以能力歡躍；牠出去迎擊仇敵。 \end{tabularx} \\ \\ \relax
39:22 & \begin{tabularx}{0.7\textwidth}{X} 牠嘲笑懼怕，並不驚惶，也不因刀劍退卻。 \end{tabularx} \\ \\ \relax
39:23 & \begin{tabularx}{0.7\textwidth}{X} 箭袋在牠身上錚錚有聲，槍和短槍閃閃發亮。 \end{tabularx} \\ \\ \relax
39:24 & \begin{tabularx}{0.7\textwidth}{X} 牠震顫激動，將地吞下 ；一聽角聲就站不住。 \end{tabularx} \\ \\ \relax
39:25 & \begin{tabularx}{0.7\textwidth}{X} 每逢角聲一響，牠說：『啊哈！』牠從遠處聞到戰爭的氣息，聽見軍官如雷的吼聲和吶喊。 \end{tabularx} \\ \\ \relax
39:26 & \begin{tabularx}{0.7\textwidth}{X} 「鷹展開翅膀向南飛翔，豈是藉著你的智慧嗎？ \end{tabularx} \\ \\ \relax
39:27 & \begin{tabularx}{0.7\textwidth}{X} 大鷹上騰在高處搭窩，豈是聽你的指示嗎？ \end{tabularx} \\ \\ \relax
39:28 & \begin{tabularx}{0.7\textwidth}{X} 牠住在山巖，以山峰和堅固之所為家， \end{tabularx} \\ \\ \relax
39:29 & \begin{tabularx}{0.7\textwidth}{X} 從那裡窺察食物，眼睛自遠方瞭望。 \end{tabularx} \\ \\ \relax
39:30 & \begin{tabularx}{0.7\textwidth}{X} 牠的雛吸血；被殺的人在哪裡，牠也在哪裡。」 \end{tabularx} \\ \\ \relax
40:1 & \begin{tabularx}{0.7\textwidth}{X} 耶和華繼續對約伯說： \end{tabularx} \\ \\ \relax
40:2 & \begin{tabularx}{0.7\textwidth}{X} 「強辯的豈可與全能者爭論？與神辯駁的可以回答吧！」 \end{tabularx} \\ \\ \relax
40:3 & \begin{tabularx}{0.7\textwidth}{X} 於是，約伯回答耶和華說： \end{tabularx} \\ \\ \relax
40:4 & \begin{tabularx}{0.7\textwidth}{X} 「看哪，我是卑賤的！我用甚麼回答你呢？我只好用手摀住我的口。 \end{tabularx} \\ \\ \relax
40:5 & \begin{tabularx}{0.7\textwidth}{X} 我說了一次，就不回答；說了兩次，不再說了。」 \end{tabularx} \\ \\ \relax
40:6 & \begin{tabularx}{0.7\textwidth}{X} 於是，耶和華從旋風中回答約伯說： \end{tabularx} \\ \\ \relax
40:7 & \begin{tabularx}{0.7\textwidth}{X} 「你要如勇士束腰；我問你，你可以讓我知道。 \end{tabularx} \\ \\ \relax
40:8 & \begin{tabularx}{0.7\textwidth}{X} 你豈可廢棄我的判斷？豈可定我有罪，好顯自己為義嗎？ \end{tabularx} \\ \\ \relax
40:9 & \begin{tabularx}{0.7\textwidth}{X} 你有神那樣的膀臂嗎？你能像他那樣發雷聲嗎？ \end{tabularx} \\ \\ \relax
40:10 & \begin{tabularx}{0.7\textwidth}{X} 「你要以榮耀莊嚴為妝飾，以尊榮威嚴為衣服。 \end{tabularx} \\ \\ \relax
40:11 & \begin{tabularx}{0.7\textwidth}{X} 你要發出你滿溢的怒氣，見一切驕傲的人，使他降卑； \end{tabularx} \\ \\ \relax
40:12 & \begin{tabularx}{0.7\textwidth}{X} 你見一切驕傲的人，將他制伏，把惡人踐踏在原來地方。 \end{tabularx} \\ \\ \relax
40:13 & \begin{tabularx}{0.7\textwidth}{X} 你將他們一同埋藏在塵土中，把他們的臉遮蔽在隱密處。 \end{tabularx} \\ \\ \relax
40:14 & \begin{tabularx}{0.7\textwidth}{X} 這樣，我也向你承認，你的右手能救你自己。 \end{tabularx} \\ \\ \relax
40:15 & \begin{tabularx}{0.7\textwidth}{X} 「看哪，我造河馬，也造了你；牠吃草像牛一樣。 \end{tabularx} \\ \\ \relax
40:16 & \begin{tabularx}{0.7\textwidth}{X} 看哪，牠的力氣在腰間，能力在肚腹的肌肉上。 \end{tabularx} \\ \\ \relax
40:17 & \begin{tabularx}{0.7\textwidth}{X} 牠挺直尾巴如香柏樹，牠大腿的筋緊密結合。 \end{tabularx} \\ \\ \relax
40:18 & \begin{tabularx}{0.7\textwidth}{X} 牠的骨頭好像銅管；牠的肢體彷彿鐵棍。 \end{tabularx} \\ \\ \relax
40:19 & \begin{tabularx}{0.7\textwidth}{X} 「牠在神所造之物中為首，只有創造牠的能攜刀臨近牠。 \end{tabularx} \\ \\ \relax
40:20 & \begin{tabularx}{0.7\textwidth}{X} 諸山為牠產出食物，百獸也在那裡遊玩。 \end{tabularx} \\ \\ \relax
40:21 & \begin{tabularx}{0.7\textwidth}{X} 牠伏在蓮葉之下，在蘆葦和沼澤的隱密處。 \end{tabularx} \\ \\ \relax
40:22 & \begin{tabularx}{0.7\textwidth}{X} 蓮葉的陰影遮蔽牠，溪旁的柳樹環繞牠。 \end{tabularx} \\ \\ \relax
40:23 & \begin{tabularx}{0.7\textwidth}{X} 看哪，河水氾濫，牠不慌張；連約旦河漲到牠口邊，牠也安然自若。 \end{tabularx} \\ \\ \relax
40:24 & \begin{tabularx}{0.7\textwidth}{X} 誰能在牠眼前捉拿牠呢？誰能以圈套穿牠鼻子呢？」 \end{tabularx} \\ \\ \relax
41:1 & \begin{tabularx}{0.7\textwidth}{X} 「你能用魚鉤釣上力威亞探嗎？能用繩子壓下牠的舌頭嗎？ \end{tabularx} \\ \\ \relax
41:2 & \begin{tabularx}{0.7\textwidth}{X} 你能用繩索穿牠的鼻子嗎？能用鉤子穿牠的腮骨嗎？ \end{tabularx} \\ \\ \relax
41:3 & \begin{tabularx}{0.7\textwidth}{X} 牠豈向你連連懇求，向你說溫柔的話嗎？ \end{tabularx} \\ \\ \relax
41:4 & \begin{tabularx}{0.7\textwidth}{X} 牠豈肯與你立約，讓你拿牠永遠作奴僕嗎？ \end{tabularx} \\ \\ \relax
41:5 & \begin{tabularx}{0.7\textwidth}{X} 你豈可拿牠當雀鳥玩耍？豈可將牠繫來給你幼女？ \end{tabularx} \\ \\ \relax
41:6 & \begin{tabularx}{0.7\textwidth}{X} 合夥的魚販豈可拿牠當貨物？他們豈可把牠分給商人呢？ \end{tabularx} \\ \\ \relax
41:7 & \begin{tabularx}{0.7\textwidth}{X} 你能用倒鉤扎滿牠的皮，能用魚叉叉滿牠的頭嗎？ \end{tabularx} \\ \\ \relax
41:8 & \begin{tabularx}{0.7\textwidth}{X} 把你的手掌按在牠身上吧！想一想與牠搏鬥，你就不再這樣做了！ \end{tabularx} \\ \\ \relax
41:9 & \begin{tabularx}{0.7\textwidth}{X} 看哪，對牠有指望是徒然的；一見牠，豈不也喪膽嗎？ \end{tabularx} \\ \\ \relax
41:10 & \begin{tabularx}{0.7\textwidth}{X} 沒有那麼兇猛的人敢惹牠。這樣，誰能在我面前站立得住呢？ \end{tabularx} \\ \\ \relax
41:11 & \begin{tabularx}{0.7\textwidth}{X} 誰能與我對質，使我償還呢？天下萬物都是我的。 \end{tabularx} \\ \\ \relax
41:12 & \begin{tabularx}{0.7\textwidth}{X} 「我不能緘默不提牠的肢體和力量，以及健美的骨骼。 \end{tabularx} \\ \\ \relax
41:13 & \begin{tabularx}{0.7\textwidth}{X} 誰能剝牠的外皮？誰能進牠的鎧甲之間呢？ \end{tabularx} \\ \\ \relax
41:14 & \begin{tabularx}{0.7\textwidth}{X} 誰能開牠的腮頰？牠牙齒的四圍是可畏的。 \end{tabularx} \\ \\ \relax
41:15 & \begin{tabularx}{0.7\textwidth}{X} 牠的背上有一排排的鱗甲，緊緊閉合，封得嚴密。 \end{tabularx} \\ \\ \relax
41:16 & \begin{tabularx}{0.7\textwidth}{X} 這鱗甲一一相連，氣不得透入其間， \end{tabularx} \\ \\ \relax
41:17 & \begin{tabularx}{0.7\textwidth}{X} 互相連接，膠結一起，不能分開。 \end{tabularx} \\ \\ \relax
41:18 & \begin{tabularx}{0.7\textwidth}{X} 牠打噴嚏就發出光來，牠的眼睛好像晨曦。 \end{tabularx} \\ \\ \relax
41:19 & \begin{tabularx}{0.7\textwidth}{X} 從牠口中發出燒著的火把，有火星飛迸出來； \end{tabularx} \\ \\ \relax
41:20 & \begin{tabularx}{0.7\textwidth}{X} 從牠鼻孔冒出煙來，如燒開的鍋在沸騰。 \end{tabularx} \\ \\ \relax
41:21 & \begin{tabularx}{0.7\textwidth}{X} 牠的氣點著煤炭，有火焰從牠口中發出。 \end{tabularx} \\ \\ \relax
41:22 & \begin{tabularx}{0.7\textwidth}{X} 牠頸項中存著勁力，恐懼在牠面前蹦跳。 \end{tabularx} \\ \\ \relax
41:23 & \begin{tabularx}{0.7\textwidth}{X} 牠的肉塊緊緊結連，緊貼其身，不能搖動。 \end{tabularx} \\ \\ \relax
41:24 & \begin{tabularx}{0.7\textwidth}{X} 牠的心結實如石頭，如下面的磨石那樣結實。 \end{tabularx} \\ \\ \relax
41:25 & \begin{tabularx}{0.7\textwidth}{X} 牠一起來，神明都恐懼，因崩潰而驚慌失措。 \end{tabularx} \\ \\ \relax
41:26 & \begin{tabularx}{0.7\textwidth}{X} 人用刀劍扎牠，是無用的，槍、標槍、尖槍也一樣。 \end{tabularx} \\ \\ \relax
41:27 & \begin{tabularx}{0.7\textwidth}{X} 牠以鐵為乾草，以銅為爛木。 \end{tabularx} \\ \\ \relax
41:28 & \begin{tabularx}{0.7\textwidth}{X} 箭不能使牠逃走，牠看彈石如碎秸。 \end{tabularx} \\ \\ \relax
41:29 & \begin{tabularx}{0.7\textwidth}{X} 牠當棍棒作碎秸，牠嘲笑短槍的颼颼聲。 \end{tabularx} \\ \\ \relax
41:30 & \begin{tabularx}{0.7\textwidth}{X} 牠肚腹下面是尖瓦片；牠如釘耙刮過淤泥。 \end{tabularx} \\ \\ \relax
41:31 & \begin{tabularx}{0.7\textwidth}{X} 牠使深淵滾沸如鍋，使海洋如鍋中膏油。 \end{tabularx} \\ \\ \relax
41:32 & \begin{tabularx}{0.7\textwidth}{X} 牠使走過以後的路發光，令人覺得深淵如同白髮。 \end{tabularx} \\ \\ \relax
41:33 & \begin{tabularx}{0.7\textwidth}{X} 塵世上沒有像牠那樣的受造物，一無所懼。 \end{tabularx} \\ \\ \relax
41:34 & \begin{tabularx}{0.7\textwidth}{X} 凡高大的，牠盯著看；牠在一切狂傲的野獸中作王。」 \end{tabularx} \\ \\ \relax
42:1 & \begin{tabularx}{0.7\textwidth}{X} 約伯回答耶和華說： \end{tabularx} \\ \\ \relax
42:2 & \begin{tabularx}{0.7\textwidth}{X} 「我知道，你萬事都能做；你的計劃不能攔阻。 \end{tabularx} \\ \\ \relax
42:3 & \begin{tabularx}{0.7\textwidth}{X} 誰無知使你的旨意隱藏呢？因此我說的，我不明白；這些事太奇妙，是我不知道的。 \end{tabularx} \\ \\ \relax
42:4 & \begin{tabularx}{0.7\textwidth}{X} 求你聽我，我要說話；我問你，求你讓我知道。 \end{tabularx} \\ \\ \relax
42:5 & \begin{tabularx}{0.7\textwidth}{X} 我從前風聞有你，現在親眼看見你。 \end{tabularx} \\ \\ \relax
42:6 & \begin{tabularx}{0.7\textwidth}{X} 因此我撤回，在塵土和爐灰中懊悔。」 \end{tabularx} \\ \\ \relax
42:7 & \begin{tabularx}{0.7\textwidth}{X} 耶和華對約伯說話以後，耶和華就對提幔人以利法說：「我的怒氣向你和你兩個朋友發作，因為你們議論我，不如我的僕人約伯說的正確。 \end{tabularx} \\ \\ \relax
42:8 & \begin{tabularx}{0.7\textwidth}{X} 現在你們要為自己取七頭公牛、七隻公羊，到我的僕人約伯那裡去，為自己獻上燔祭，我的僕人約伯就為你們祈禱。我必悅納他，不按你們的愚妄處置你們。你們議論我，不如我的僕人約伯說的正確。」 \end{tabularx} \\ \\ \relax
42:9 & \begin{tabularx}{0.7\textwidth}{X} 於是提幔人以利法、書亞人比勒達、拿瑪人瑣法遵照耶和華所吩咐的去做，耶和華就悅納約伯。 \end{tabularx} \\ \\ \relax
42:10 & \begin{tabularx}{0.7\textwidth}{X} 約伯為他的朋友祈禱。耶和華就使約伯從苦境中轉回，並且耶和華賜給他的比他從前所有的加倍。 \end{tabularx} \\ \\ \relax
42:11 & \begin{tabularx}{0.7\textwidth}{X} 約伯的兄弟、姊妹，和以前所認識的人都來到他那裡，在他家裡跟他一同吃飯。他們因耶和華所降於他的一切災禍，都為他悲傷，安慰他。每人送他一塊可錫塔和一個金環。 \end{tabularx} \\ \\ \relax
42:12 & \begin{tabularx}{0.7\textwidth}{X} 這樣，耶和華後來賜福給約伯比先前更多。他有一萬四千隻羊，六千匹駱駝，一千對牛，一千匹母驢。 \end{tabularx} \\ \\ \relax
42:13 & \begin{tabularx}{0.7\textwidth}{X} 他也有七個兒子，三個女兒。 \end{tabularx} \\ \\ \relax
42:14 & \begin{tabularx}{0.7\textwidth}{X} 他給長女起名叫耶米瑪，次女叫基洗亞，三女叫基連哈樸。 \end{tabularx} \\ \\ \relax
42:15 & \begin{tabularx}{0.7\textwidth}{X} 在全地的婦女中找不著像約伯的女兒那樣美貌的。她們的父親使她們在兄弟中得產業。 \end{tabularx} \\ \\ \relax
42:16 & \begin{tabularx}{0.7\textwidth}{X} 此後，約伯又活了一百四十年，得見他的四代兒孫。 \end{tabularx} \\ \\ \relax
42:17 & \begin{tabularx}{0.7\textwidth}{X} 這樣，約伯年紀老邁，日子滿足而死。 \end{tabularx} \\ \\
[1ex]
\hline
\hline
\end{longtable}
約伯對於三個朋友勸勉的反應,可以反映出一個受苦的人,心靈的思想和狀態.這些話對於我們,無論是用在勉勵別人,或是瞭解自己方面,都有很大的幫助；可以作為我們受苦時候的借鏡,藉以得著力量勝過苦難.
\newline
\newline
約伯對神,內心有迫切的要求；他有時覺得神是一位難以應付的神.另一方面他覺得三個朋友越說越不對；因為他們沒有可能得到別人的瞭解.雖然他有時寄望於神,但又感到神好像沒有給他機會.約伯切慕和神面對面講理.從九章1-4；二十三章3-10；九章32-十章2的話中,清楚表明他自信無辜之心,在痛苦中始終沒動搖,他在神前敬虔的操練,為人的屬靈原則永不動搖.他雖多次埋怨神,甚至控告神；但他為人的方式仍不改變,從他的話中可見他堅強的信心,永不動搖.十三章節15另一翻譯是:「神雖殺我,我雖無指望,我仍要信靠祂.」約伯話中也常題神創造的大功,和他三個朋友對於神的智慧,創造宇宙的大能,有同樣的認識.在他所有的辯論,神暗中把屬靈真理賜給他；所以約伯言談之中,常說含有屬靈亮光帶有啟示的道理,連他自己都不瞭解.雖然約伯覺得痛苦使他近於死亡,他仍是清潔,他希望有一日可表白他的清潔和無辜.(十六15-21)他說:「他要為他作見證.」我相信他一定是想到舊約,亞當的兒子義人亞伯,無辜被他兄弟該隱殺害；雖無人為他伸冤；但聖經說,地為他哀號.約伯說:「現今在天有我的見證,在上有我的中保.」(19)「中保」有譯為「記錄」這是當朋友誤責約伯時,他內心的感受:是聖靈給他的安慰,也是事實.他的無辜,不但天地為他作證,他和朋友的談話在天上有記錄.他說:「惟願我的言語.....都記錄在書上；用鐵筆鐫刻,用鉛灌在磐石上,直存到永遠.」(1九23)約伯的話真是超乎他所求所想的；神讓它全寫在聖經上.約伯又說:「我知道我的救贖主活著,末了必站立在地上.」(25)這是何等的啟示!
\newline
\newline
「救贖主」在舊約是個隱藏而重要的啟示.救贖主與救贖者有血統之關係；是大有才能,大有愛心,為受冤,窮苦,無指望者的利益勞苦奔跑.按舊約律法,窮苦人家變賣產業,其本族富有之最近親,有責任贖回；若有同族死了無後留下者的寡婦,應由其最近親屬娶之,為死者名下留後代.舊約路得記,清楚論及此事；其屬靈意義,則將我們的救贖主,和罪人的關係表明.當我們墮落,犯罪,貧窮,變賣一切毫無所有時,主耶穌是我們救贖主；我們是人,祂要和我們建立血統關係,道成肉身,住在我們中間.稱我們為弟兄,不以為恥；憑祂的愛心和大能救贖我們.「救贖主」這啟示是最奇妙的；約伯看見了這奇妙的啟示.
\newline
\newline
約伯對死亡有個更加新的看法,他本以為死亡是個和平共處的地方,繼之他對死亡感到恐懼；現在他卻存著復活的盼望,肉身死了,靈魂永存,且與主同在.這一點在舊約中,很少有這樣的看見；經歷大苦難的約伯,竟然看見了這奇妙的啟示.此外,神的聖靈暗中啟示約伯,義人受苦,非因罪,非因神不喜歡他；乃是神對他的考驗和試煉.(二十三10)神知道他,而且使他經試煉之後更尊貴,更純全.
\newline
\newline
約伯和三個朋友彼此長談,三個朋友越講心靈越空洞,終成無稽之談.至於約伯言談漸見新亮光,步進真理.三個朋友堅持成見,停留於和人談論神的止境中.至於約伯,除了和人談論神之外,進而和神談話；他有渴慕求明白真理的心.真是「飢渴慕義的人有福了,他們必得飽足.」
\newline
\newline
當約伯以為神隱藏之時,事實上神未曾離棄他；神如慈母,暗中觀察剛學走路的小孩,急需時予一臂之助.大衛說:「因為我遭遇患難,祂必暗暗的保守我.」(詩二十七5)
\newline
\newline
在約伯和三個朋友會談中,我們看見神對他的恩典和同在.可惜約伯第三次說的話,因三個朋友對他嚴厲指責,定他的罪；所以極力為自己的無辜辯護,越辯越痛苦,靈越趨黑暗；雖聖靈給他亮光,但對他發不出大的力量和安慰,其痛苦無法接受人的安慰.
\newline
\newline
約怕在廿九至卅一章的說話,可說是最厲害,最癲狂.廿九章論及「從前的我」何等快樂幸福,最使他羨慕的,是與神親密的關係,(2)他認為屬靈的福才是最大的福.他論及他早年時的豐富,受人尊敬,是困苦人的及時雨,需要者的亮光,無法無天者的救星,是無冠冕的君王.卅章論「現在的我」是可憐的,身,心,靈都極其痛苦；物質損失,兒女死亡,身患惡疾；被譏笑,親朋遠離；最使他難當的痛苦,是感到神向他變心.卅一章論「美麗的我」他對鏡自我欣賞.許多釋經家著重這三章的「我」字,廿九章共42次,卅章54次,卅一章62次.越講越多「我」,屢講美的幸福人生,屢講身體,心靈的痛苦；屢講受冤枉,受委屈,都以「我」為大前題.到了最後,他很不甘心,奈何為仁,為義的「我」反而受苦!「我」要暴發忿怒,「我」要抗議,「我」要伸冤.約伯不懼怕所有的控告,他說:「惟願有一位肯聽我,......,願全能者回答我.」他希望神出來,他希望他的辯駁成為他榮耀的明證,裝飾在肩上非常美麗.這時約伯認為人無法和他辯論,連神也藏起來了.其實,神在等候.
\newline
\newline
約伯最大的痛苦,不是錢財物質的損失,其痛苦最大的根源,亦非神隱藏；而是他自我唉嘆.他是個純正且敬畏神的人,他的生活公義,聖潔；但他並非公義,聖潔的化身,誰能比耶穌基督更公義,聖潔?祂被罵不還口,受害不說威嚇的話；只將自己交託公義審判的主,祂曉得神了解一切.主耶穌超過萬人之上,凡事為我們留下懿範.約伯雖然公義,哪能與主比擬,雖然有信心,哪像耶穌基督完全美好.耶穌基督是我們信心創始成終的主,我們必須仰望祂.
\newline
\newline
當約伯逞強之後,三個朋友非常害怕,大家不歡將散,忽然出現個青年以利戶.
\newline
\newline
青年學者以利戶(卅二-卅七)此人不鳴則已,一鳴驚人；他連續講了六章長篇的話.他是個屬靈的青年學者,說話條條有理,大有聰明；話中談論天文地理,學問淵博；他的話有說服力,有亮光,顯有屬靈的啟示.他在三十二8:三十三4說,神使年長者有屬世的經驗,屬靈的智慧；然而這一切的根源都出自神的靈.此點是他較其他三人更可幫助約伯的.這少年人言談穩健,為人謙卑.他的來龍去脈,聖經隻字不提.無論如何,所有約伯和三個朋友的辯論過程,他身在其中；他聽見約伯三個朋友亂定約伯的罪,也聽見約伯越說話越癲狂；他有意尋求真理的答案,他雖言語滿懷,心靈激動；(三十二 18-19)但他仍耐心靜聽,直到各年長者無言,他才開始發言.他也是個熱愛真理的人；感到能參加其中和各屬靈長者一同探討「義人受苦」的問題,真是千載難逢的機會.他尊重長者,說:「年老的當先說話,受膏的當以智慧教訓人.」(7)可是後來聽完了他們的談話,他說:「尊貴的不都冇智慧,受膏的不都能明白公平.因此我說,你們要聽我言,我也要陳說我的意見.」(9-10)青年人難免有天然未成熟的表現.
\newline
\newline
以利戶一方面反駁批判約伯的言談,一方面發表他對人生苦難的認識,對此真理之瞭解.他把約伯的話歸納,然後逐條回答.
\newline
\newline
第一,反駁約伯所說:「神攻擊無辜的人」.(三十二9-11)以利戶說:「你這話無理,因神比世人更大.你為何與祂爭論呢?因祂的事都不對人解說.神說一次,兩次,世人卻不理會.」
\newline
\newline
有人說:「如果神將祂要作的事預先告訴你,你的快樂將減少一半；痛苦將增加一倍.」比方說,神告訴你卅歲將結束你的性命；你一定終日為此恐懼憂鬱,悶悶不樂,未到時候便死了.如果神預告你卅歲將發財；你就坐待那日,作事失去信心學習；到時也不覺快樂了.
\newline
\newline
神比世人更大,因祂的愛心和智慧,所以祂的事不對人解說；但是,神有時對人說話,人置之不理；有時神也藉著異象和夢或疾病對人說話；但人心剛硬,不肯靜聽.
\newline
\newline
第二,約伯說神專權奪理,神的性格很難測度.卅四章以利戶解釋這個問題.
\newline
\newline
第三,約伯說神賞罰不明.以利戶在三十五9-11說:義人在受苦之中,尊貴超過其他萬物；空中,地上的飛禽走獸,受痛苦的反應非哀求,則掙扎逃避；惟人痛苦時,可在夜間歌唱；因人是有靈性的動物.可惜人常把自己的地位降低.同時他也覺得人在痛苦中有神的大能,有神的作為.
\newline
\newline
以利戶把約伯的話歸納三點反駁,由卅六至卅八章,論及他本身對各樣事物真理的看法.論人類的苦難問題.(三十六1-23),(三十六24-三十七),論神的大作為.
\newline
\newline
以利戶對真理研究的心得,請看三十六8-10這裡說為何有時神使人受苦；因為有的人必須受苦,不受苦就易於放蕩.苦難如鎖鍊繩索捆綁,使人不自由,也可使人坐下靜聽神的話.指出人的過犯,罪惡,內心的驕傲；人必須安靜等候神的亮光和聖靈的提醒,才能明白裡面的惡行.大衛在(詩篇一三九23-24)所說的,不是單外面的問題:他求神鑑察他裡面的惡行.一個基督徒──所謂屬靈的人,裡面有問題；只要安靜下來,神就開他耳目,叫他明白心裡的光景.這是以利戶深入的見解.約伯三個朋友,好像審判官地對待約伯；以利戶認為神是教師,苦難是他給人學習的功課；他以客觀,心平氣和的態度勸慰約伯.
\newline
\newline
(三十六15-16)論及苦難的作用,有時苦難救我們脫離患難；好像注射預防針,小痛可免大痛；今生的小苦可免將來永世的大苦,這是神奇妙的作為.接著,以利戶又說及神,在宇宙萬象之中的大智慧和大權能.當以利戶說了六章長篇的話之後；神開始說話.為何以利戶的話,能引出神的話語呢!
\newline
\newline
時間非常重要(三十八1)「那時耶和華從旋風中回答約伯.」「那時」,就是人的話都講完之時:人的盡頭,就是神的開頭.當人嘈吵之後,心靈安靜下來,才能聽見神的聲音；神的話使人清醒.至此,三個年長者謙卑下來,他們具有學者風度.
\newline
\newline
神說了兩次話,茲略述綱要:
\newline
\newline
(一)神創造萬有的奇妙作為,論及大地和山嶺.(三十八1-18)
\newline
\newline
(二)神創造穹蒼萬象的奇妙作為(三十八19-38)論光明及黑暗,人的由來；氣候變化的原則,天體的運行,大自然的定律.
\newline
\newline
(三)地上生物的奇妙(三十八39-三十九30)神說祂創造各樣活物,牠們家庭龐大；野獸的性情,活物動向各異.真是奇妙的世界!
\newline
\newline
(四)論神另一方面的權能(四十-四十一)這裡所說不同,這裡多次論及驕傲.四十一章末節另一繙譯:「他是一切驕傲之子的王,一切狂傲之物的王子.」
\newline
\newline
綜合而言,論及道德的世界,是非,善惡,驕傲,謙卑；隱藏的,抽象的道德世界,一切由神管理.這裡論及看不見的物質世界,卅九章舉出河馬和鱷魚,河馬不是現今我們所見的；那時的河馬,尾巴如香柏樹,可能指古時一種巨大恐怖的活物.鱷魚也非今日所能見的.驕傲之意和三十三17同出一字根,指一切驕傲之子的王,是狂傲之物的王子──背後就是撒旦.
\newline
\newline
神對約伯共提出184個問題:「你知道嗎?曉得嗎?對我有何意見?」神問約伯諸多問題,並非顯赫自己無所不能,而是藉此開約伯的眼目.讓他看見有形的宇宙,天地萬物,都由神妥善管理,供給一切所需.讓他曉得驕傲的人,陷落自己的詭計中.神開約伯信心之眼,使他明白神創造管理宇宙,和萬物毫無錯誤,難道對待約伯有錯嗎?
\newline
\newline
約伯開始明白過來,心有安息.經歷試煉必須保持永恆的信心,才能蒙神賜福；因為神的美意永不改變.約伯仆倒在神面前,恍然大悟；信心把他的靈眼打開.約伯從苦境轉回.(四十二2-5)他現在明白神在他身上是無所不能的神.約伯素有堅強的信心,有屬靈的生活,無形中他是個完全人,因而自以為義.神所愛者常自鳴得意.
\newline
\newline
聖經舊約的雅各,一向不受人的管教,自從在母胎裡,就為難他的母親和哥哥,都是他對付人,人無法對付他；他用一碗紅豆湯從他哥哥換取了長子名份.真是一本萬利!唯獨神能對付他；他到晚年,在神面前何等溫柔；臨終扶著杖頭敬拜神.
\newline
\newline
當約伯明白神的對付時,他屬靈眼界大開,說:「從前風聞有你,現在親眼看見你.」他進入屬靈實際中,看見奇妙的事實,到達至上無比的境界；如以弗所書一章所說:「遠超過一切所有的.」約伯此時完全謙卑下來,為他三個朋友禱告,他們的友誼比前更加鞏固.
\newline
\newline
神賜福約伯比從前倍增,特別他三個女兒也名列聖經之中；可能是當時最美麗,最屬靈的女子.長女名耶米瑪,中譯美晨.次女基洗亞,意思是肉桂,桂花,桂香.三女基蓮哈樸,意思是一種畫眼的顏料,可稱為慧玉.
\newline
\newline
神的美意在約伯身上都顯明出來.
\newline
\newline
感謝神!藉此給我們屬靈的榜樣,作我們處此苦難世界的借鏡,叫我們在苦難中得以剛強.
\newline
\newline


\newpage

\section{第52屆~港九培靈會~薛玉光~第一講　認識獨一的真神}
\label{sec:632}
\textbf{第一講　認識獨一的真神}
\newline
\newline
連結: \href{https://www.hkbibleconference.org/session-message/view/632}{\texttt{https://www.hkbibleconference.org/session-message/view/632}}
\newline
\newline
\hyperref[sec:628]{< < < PREV SERMON < < <}
~
\hyperref[sec:toc]{[返目錄]}
~
\hyperref[sec:633]{> > > NEXT SERMON > > >}
\newline
\newline
今次我所講的總題是:「合一的真理和見證」,前四次是講論教會合一真理的基礎,後四次則講合一的見證,共八次講完.我覺得今年培靈會的幾個講題都很不平常,我們沒有彼此商量,只在主面前等候；焦牧師早堂所講的約伯記是關乎基督徒的受苦,今天教會對受苦的教導太缺少,很多基督徒不明白受苦的真理,所以稍遇試煉便跌倒退後,故此神要藉著焦牧師來教導我們.晚堂的史牧師會講愛心,今天信徒非常缺乏愛心的見證,史牧師會用十晚的時間和我們分享林前十三章的真理.而我則講到合一的問題,因為我覺得這也是今天教會的需要.我相信如果香港的信徒和同工都注意這真理,並照著去行,則香港教會便會有大大的不同.
\newline
\newline
我們先來看幾處的經文:(申六4)「以色列阿,你要聽,耶和華我們神是獨一的主.」(可十二29)「耶穌回答說:第一要緊的,就是說,以色列阿,你要聽,主我們神是獨一的主.」(約五44)「你們互相受榮耀,卻不求從獨一之神來的榮耀,怎能信我呢?」(羅十六27)「願榮耀因耶穌基督歸與獨一全智的神,直到永遠,阿們.」
\newline
\newline
今次我為甚麼決定講教會合一的道理呢?因為教會在今天一般情形來說,都失去合一的見證,基督徒在主裡應該合一,這是我們的主的心願.耶穌在約翰福音十七章和門徒分離之前的禱告中,多次求神教導門徒合而為一,但是今天我們做不到,以致教會不能榮耀神,不能勝過仇敵,個人靈命也不能好好的長進.講合一要用很多的時間,我只把我最近領受的和大家分享.在未講這題目以前,感謝主讓我遇到不少位主重用的僕人,他們也有同樣的看法,真可說是不謀而合.
\newline
\newline
前四次我要講合一真理的基礎,我們必須先認識真理,然後照著去行,這樣的合一才是真實的,有根基的.今天早上先講頭一個關於真理的基礎,就是認識獨一的真神.教會今天有許多難題,主要的原因是我們還沒有真正的認識神.耶利米書九3提到耶和華說:「我的民並不認識我.」耶和華的百姓就是猶太人,猶太人為甚麼不認識神呢?他們從祖宗亞伯拉罕就認識神,但是神卻對耶利米說他的百姓並不認識神.這是甚麼意思呢?這是指他們的事奉與他們行事為人和他們所認識的神不相符合.今天我們也是神的百姓,在屬靈的地位和身份上說和猶太人相似,我們在唱詩,禱告,讀經和聽道等事情上對神的認識都沒錯,如果來個考試,準可以拿得九十五到一百分,特別是福音派的教會；可是我們的生活,工作,事奉和我們的唱詩,禱告卻差了一大截.如果也來考試一下,可能我們只得六十分,甚至是三十分,二十分.如果我們在神面前謙卑和誠實一點,便可以看見自己實際的光景,很多教會都唸使徒信經,那是信的綱要,我們說相信獨一的真神,創造天地的主,唸了又唸,但我們的行事為人,工作,和我們所唸的是否相符呢?我們真應該好好的再思想一下,認識我們這位真神.我們所信的神,是以色列人列祖的神,是亞伯拉罕,以撒,雅各的神,是舊約時代先知的神,是新約時代使徒的神,不是今天西方一些神學家所講的神,也不是教會中某些長執心目中的神.所以我們要澈底真正的認識祂.
\newline
\newline
第一,神是獨一的真神.「真」就是「原來」的意思.(創一1)說「起初神創造天地」,(約一1)說「太初有道」,祂是太初的那一位.凡是真的東西都是起初的東西,假冒的東西是後來才有的.中國人很注重名家的字畫,其中有一些我們稱為真跡,就是當時原作者寫的；後來則有些冒充的,仿效而寫的,所以是假的,真與假要由藝術界的一些鑑定家才能判別.我們所敬拜的真神耶和華,是原先的那位,不是後來人假冒的,甚至魔鬼也會冒充神,我們必須當心.
\newline
\newline
神也是獨一的,是僅此一位的,這是個事實,事實也就是最真實的,就是真理,英文的真理一字也就是真實.有人批評基督教太不民主,不能容納別的宗教,別的神.美國教育哲學家杜威,他的民主理論曾經哄動一時,但現在美國卻身受其害,杜威說:「在這個民主的時代,我們怎可以信一位專制的神呢?」有些信徒也有這個疑問,以為各宗教都差不多,只是基督教比較強一點而已.我們絕不能這樣信,因為獨一的真神是全然的事實,惟有祂是真實的.這不是基督徒自己理論或主張,乃是事實存在於先,是不能改變的.神在聖經中啟示摩西說,神是獨一的神；耶穌來到世上,也是照樣的說；聖靈也是這麼見證,保羅在羅馬書的末了也說,神是獨一的神.由此可見,這是三位一體神自己的宣告,是個絕對不能改變的事實.如果我們說神不是獨一的,各宗教都是一樣,那我們就是把神當作說謊了.
\newline
\newline
第二,這位獨一的真神是有啟示的真神.這啟示藉著不同的方法,一個是一般性普遍的啟示,乃是藉著祂所造的自然界,天地萬物,來顯明神的永能和神性,這是每一個人都可以看見,可以明白的.請問有人能造高山,大海,太陽,星球等等嗎?這些都是明顯的證據,人如果不承認,就是故意的不信,就是犯了是無可推諉的不信的大罪.
\newline
\newline
另一個是明白的啟示,就是聖經.對一切願意尋求祂的人,神為他們預備一本書,而這個啟示,非常寶貴和重要的.基督教和其他宗教有根本不同的地方,其他宗教是人尋找神,但耶穌的福音是神尋找人,獨一的真神向人顯示,這是神的恩典,祂開恩把屬靈的事向我們顯明,我們才知道.我們藉著聖經的啟示,便知道神是怎樣的一位神.
\newline
\newline
第三,獨一的真神也是立約的神,這是很奇妙的.人與人之間有立約,做生意的人要簽合同,這種立約是人和人所立的.但人和貓,和狗能立約嗎?用甚麼文字來立呢?因為彼此此不同類便不能立約.神是至高的神,是創造的主宰,而人是被造的,本來彼此不能立約,但神降卑祂自己,降到站在與人立約的地位上,這真是奇妙.神要向人證祂是信實的神,所以祂與人立約.舊約時代,神與信心的祖宗亞伯拉罕立約,還有以撒,雅各,那時立約是用口頭來立的,神對亞伯拉罕說要賜他有後裔,這後裔要叫萬國得福,神這樣說,事情後來就這樣成就,神守了祂的約,因為是信實的神.到新約時代,神又藉祂的兒子耶穌基督與人立約,這約是新約,是恩典之約,是用耶穌的寶血來立的.這約的內容說人可以藉著信耶穌,罪便得赦免,並且得永生.這約是恩典之約,是永遠有效,不會廢棄的.神立了約,便守到底,祂是獨一的真神.
\newline
\newline
第四,神是施恩的神,這叫我們得著極大的安慰.很多人作了多年的基督徒,仍不認識神的恩典,很多人以為神就像警察一樣,專門找人的錯.所以做基督徒很苦,沒有自由,終日愁眉苦臉.但神並不是這樣的,祂是滿有恩典,滿有慈愛的,祂所賜人的恩典不只一樣,而是恩上加恩,是數算不盡的.最大而又最基本的恩典,就是神把祂的兒子耶穌賜給人,「神既不愛惜自己的兒子為我們眾人捨了,豈不也把萬物和祂一同白白的賜給我們麼?」我們要認識這位獨一的真神,是滿有恩典的,我們在神面前是蒙恩,蒙愛的人,我們要注意與祂的關係,我們既是屬於祂的,我們應該事奉祂.(徒二十七22)保羅的手帶著鎖鍊,向與他同船的人作見證說:「我所屬,所事奉的神」.我們作基督徒的,如果認識這真理,我們是屬於獨一的真神,我們的合一便不會有問題.可惜今天很多人不認識這個真理,不知道自己屬於神,以為只是屬於某一教會宗派,有些傳道人高舉宗派,說要忠於自己的宗派教會,這話表面上可說沒有大錯,但實際很有問題.如果單講屬於宗派,忠於宗派,而不講屬於神,忠於神,是有很大毛病的,會攔阻彼此相通與合一,甚至破壞了合一.我們若認識到自己乃屬於神,是把神放在第一,我們與神的關係是永遠的,但與宗派的關係只是暫時的,天上是沒有分宗派的.我並非攻擊宗派,我自己也是屬於宗派教會,宗派本身的存在有它的背景與原因,但若把宗派放在神之上,那便大有問題,因為這不是聖經的教訓.
\newline
\newline
今天的教會,好像士師時代的光景.(士五8):「以色列人選擇新神,爭戰的事就臨到」.以色列人是神的百姓,從亞伯拉罕開始就相信獨一的真神耶和華,怎麼這裡說他們選擇新神呢?是指甚麼神呢?士師時代的特點,就是各人任意而行,自己各有主張,不遵守神的話,那就是他們心中的新神.(耶二28)話:「猶大啊,你的神的數目,與你城的數目相等.」猶大國有幾十個大小的城市,在猶大人心中也竟有幾十個神,這句話好像說中了我們福音派的教會,我們像猶大人一樣,是正統屬大衛的子孫.但是他們竟厭棄真神,拜異邦人的神,今天的教會,在唱詩,禱告的時候,承認獨一的真神；但行事為人卻另有新神,另有偶像,或是把人當作神,或是把金錢,把知識當作神.有很多受差會管轄的教會,把外國的監督,主席,差會的負責人當作神,聽從他們的話過於聽神的話.有的比較富有的教會,出錢多的人說話便大聲,沒出錢的人連說話也無資格,錢便成了那教會的神.今天一般的神學,看重知識,忽略靈命,特別是西方的一般神學院,更是強調知識,但有知識的不一定有智慧,若把知識放在第一,把它當作神,那就危險了.更可怕的,是人竟把色情看作神,西方有些名為基督教的國家,竟有同性戀的團契,同性戀的牧師,真是背道到了極點.他們選擇了新神,他們教會中的金牛犢就是色情.神是聖潔的,他恨惡罪惡,但今天竟有教會墮落到這樣的地步.
\newline
\newline
弟兄姊妹,你所認識的是怎樣的一位神呢?你所事奉的到底是哪一位神呢?聖經說:「以別神代替耶和華的,他們的愁苦必加增.」認識獨一的真神,是建立合一見證頭一塊的基石,是最重要的基礎.
\newline
\newline


\newpage

\section{第52屆~港九培靈會~薛玉光~第二講　認識唯一的救主耶穌}
\label{sec:633}
\textbf{第二講　認識唯一的救主耶穌}
\newline
\newline
連結: \href{https://www.hkbibleconference.org/session-message/view/633}{\texttt{https://www.hkbibleconference.org/session-message/view/633}}
\newline
\newline
\hyperref[sec:632]{< < < PREV SERMON < < <}
~
\hyperref[sec:toc]{[返目錄]}
~
\hyperref[sec:634]{> > > NEXT SERMON > > >}
\newline
\newline
約翰福音 17:1-17:3
\newline
\begin{longtable}{cl}
\hline
\hline
章節 & 經文 (和合本修訂版)\\
\hline
17:1 & \begin{tabularx}{0.7\textwidth}{X} 耶穌說了這些話，就舉目望天，說：「父啊，時候到了，願你榮耀你的兒子，使兒子也榮耀你； \end{tabularx} \\ \\ \relax
17:2 & \begin{tabularx}{0.7\textwidth}{X} 因為你曾賜給他權柄掌管凡血肉之軀的，使他把永生賜給你所賜給他的人。 \end{tabularx} \\ \\ \relax
17:3 & \begin{tabularx}{0.7\textwidth}{X} 認識你—獨一的真神，並且認識你所差來的耶穌基督，這就是永生。 \end{tabularx} \\ \\
[1ex]
\hline
\hline
\end{longtable}
我們先來看幾處的經文:(約十七1-3)「耶穌說了這話,就舉目望天說:父阿,時候到了,願你榮耀你的兒子,使兒子也榮耀你.正如你曾賜給他權柄,管理凡有血氣的,叫他將永生賜給你所賜給他的人.認識你獨一的真神,並且認識你所差來的耶穌基督,這就是永生.」十一節「從今以後,我不在世上,他們卻在世上,我往你那裡去.聖父啊,求你因你所賜給我的名保守他們,叫他們合而為一,像我們一樣.」22-23節「你所賜給我的榮耀,我已賜給他們,使他們合而為一,像我們合而為一.我在他們裡面,你在我裡面,使他們完完全全的合而為一,叫世人知道你差了我來,也知道你愛他們如同愛我一樣.」
\newline
\newline
(林前一10-13)「弟兄們,我藉我們主耶穌基督的名,勸你們都說一樣的話.你們中間也不可分黨,只要一心一意彼此相合.因為革來氏家裡的人,曾對我提起弟兄們來,說你們中間有分爭.我的意思就是你們各人說,我是屬保羅的,我是屬亞波羅的,我是屬磯法的,我是屬基督的.基督是分開的麼?保羅為你們釘了十字架麼?你們是奉保羅的名受了洗麼?」
\newline
\newline
世人不認識耶穌,這還沒有多大的關係,但最可惜的是教會的人不認識耶穌,這問題就大了.因為基督徒的責任,是把耶穌基督告訴世人,如果基督徒本身對耶穌認識不清楚,或認識錯誤,或根本與耶穌沒有關係,他怎能向別人傳耶穌呢?即使他願意傳,也是傳錯的道理吧了.
\newline
\newline
今天很多信徒對耶穌的認識大有問題,例如有些人說耶穌主張博愛的,另有些人信耶穌,是由於祖傳的信仰,也有人仍認為耶穌與其他宗教都是一樣的.這些人對耶穌的認識都不夠清楚,他們聚在一起,並不能使教會合一,因為他們根本沒有合一的生命.
\newline
\newline
究竟我們對耶穌應該有怎樣的認識呢?在此我只講最重要的幾點.
\newline
\newline
(一)要認識耶穌是永生神的兒子,並不是普通的人.何以證明呢?因為是耶穌的自證,是他自己說的,他在禱告中稱天上獨一的真神為父,稱自己為兒子.有人說各宗教都一樣,但實際上耶穌的福音和別的宗教根本不一樣.別宗教的教主,如佛教的釋迦牟尼,他是印度的一個王子,他從來沒有說自己是神的兒子,而佛教對創造天地的真神根本沒有答案；回教的穆罕默德,也沒有稱自己是神的兒子.只有耶穌這樣說,所以基督教根本和別的宗教不一樣.另一個證明,耶穌所說的話,所作的事,所顯明的權柄能力,都證明他的確是神的兒子.人如果不信耶穌是神的兒子,他就不是真實的基督徒,只是掛名的；他也不是真神的兒子,只是教會的會友,因為他的信沒有內容.很多人以為耶穌是個好人,是個聖賢,是個教主,但耶穌自己說他是神的兒子,這是非常重要的.
\newline
\newline
(二)要認識耶穌是個無罪的人.從古到今,沒有一人無罪,除了耶穌以外.祂是個完全的人,又是完全的神.祂在世上三十三年,住在罪惡的世界,受過罪惡的試探,和我們一樣,可是祂卻沒有犯過罪.有人敵視祂,要找祂的把柄,甚至猶太教的領袖,常派偵探去偵查祂,要找祂的錯處,但都找不著.最後,耶穌被人釘死,罪名是別人加上的,是不能成立的.彼拉多是羅馬的巡撫,但他按羅馬的法律,那是比較公道的法律來審查祂時,卻查不出祂有甚麼罪.用猶太教的舊約律法來查,也得不出甚麼結果.以耶穌是個完全無罪的人,當祂死的時候,天地昏黑,因天地主宰獨一的神顯出神蹟,證明祂是不該死的；這是人類的滔天大罪,把神的兒子,這位義者耶穌釘死了.
\newline
\newline
(三)要認識耶穌從死裡復活.死亡是人類的大仇敵,從來沒有人能打敗它,但耶穌卻勝過死亡,祂死了,埋葬了,第三天從死裡復活.我們怎麼知道呢?那些不信耶穌的人,說這是基督徒自己講的,是心理作用,是不能證明的.但我覺得此事可以證明,耶穌的復活,千真萬確,我且舉出兩個證明.第一,祂今天仍然掌權.世界上沒有人死了仍可以掌權的,但耶穌死了又活了,並且在教會中掌權.祂曾說祂要建立教會,陰間的權柄不能勝過祂.自從教會成立在地上,不知有多少的政治首領和君王,要用他們的力量摧毀教會,但都不能成功,因為有一位永活的耶穌在掌權.這就證明一件事,任何的政權,從羅馬的該撒開始,直到今天的無神政權,地上所有人的力量聯合起來,都沒法摧毀教會,因為教會有一位復活的主,這是個有力的證明.第二,耶穌是活的,祂改變了我,救了我.今天站在各位面前的我,不是以前的我,如果沒有耶穌的復活,我今天根本不會站在台上,但可能已經不活在世上.我的生命曾經歷過危險,是復活的主救了我的.我早年在大學是專修教育的,而且已堅決立志,終身從事教育的工作,可是大學畢業後,我沒有教書,竟作了傳道的工作,一直傳道到今天.是誰改變了我呢?是復活的主,只有祂能改變個性強烈的我,這就是最明顯的證明.
\newline
\newline
還有,耶穌基督是唯一的救主.很多人說基督教太狹窄,只承認耶穌是救主,真是豈有此理.我覺得如果承認各宗教都是一樣,才是豈有此理,因為各宗教根本不一樣.耶穌有三個資格,祂是神的兒子,祂完全沒有犯罪,祂死了又復活,祂救了世上的罪人,也救了我,所以祂是救主.這是事實,是真實的,也是真理.別的宗教有嗎?佛教,回教,印度教,儒教被稱為世界主要的大宗教,這些教主有沒有自稱是神的兒子呢?沒有.有沒有哪位教主從未犯過罪呢?沒有.孔子也承認人有罪,包括他在內.有沒有那一位從死裡復活呢?更是沒有.信的人只有去看教主的墓,回教徒每年一次到麥加朝聖,就是朝拜他們教主的墓.但耶穌死了,有基督徒到耶路撒冷去拜祂的墓嗎?沒有.耶穌的墳墓是空的.可見根本不一樣,怎能說各宗教都是一樣呢?耶穌是唯一的救主,這是基督教最大的特點,別的宗教只有教主,只是會教導人而已,但耶穌是救主,把我們從罪中救出來.如果我們否認這事,就等於說耶穌是個大騙子,我們敢如此說嗎?
\newline
\newline
可是教會今天的情況很令人奇怪,竟然有不信耶穌是神的兒子的言論,不信祂復活.在西方的大神學院中,常有這樣的理論.如果我們要選擇進神學院,就必須注意這些,任何神學院若不信耶穌是神的兒子,不信復活,我們就不應當去.這些神學的信仰,敗壞到極點.今天外國有很多大禮拜堂,因為沒有會友,所以把禮拜堂賣給回教徒作寺院,這真是基督教極大的羞恥.為甚麼教會會墮落到這樣的地步?這是因為神學信仰敗壞,所以訓練出來的傳道人,不信耶穌是神的兒子,不信耶穌復活,以致教會成為死的教會,會友都先後離去,終致教會關門.
\newline
\newline
我們要認識聖經的真理,耶穌是神的兒子,祂是無罪的,是死而復活的,是唯一的教主.我們需要祂救我們,可是,我們對得救的問題清楚嗎?很多人在教會中有很多事奉,受洗已經多年,卻不知自己得救了沒有,真是可憐.我自己也是這樣,我母親是衛理公會的會友,所以我在嬰孩時便受了洗,自幼也學會念聖經,也懂得一點禱告,可是卻沒有主的生命.直到我在大學參加基督徒團契時,才真正的清楚得救,以後才獻身服事主.沒有得救的基督徒不是真正的基督徒,一個真正的基督徒必然有永生的把握,知道不論何時離世,都可以到神那裡去.如果你沒有把握上天,那就是很有把握下地獄,因為我們只有兩個去處.所以我們必須很清楚自己的得救,相信耶穌是神的兒子,相信祂為我死,相信祂復活了.而我是一個罪人,按照神的公義,我應當沉淪,我不能自救,我需要救主.我對主說:「主阿,我相信你替我的罪而死,我接受你做我個人的救主,請你進到我中心,與我同在,保守我,直到永遠.」
\newline
\newline
已經得救的人所要作的,是傳揚耶穌,把這位救主告訴別人,這叫做服事.在服事主的時候,我們要注意一件事,就是體會主的心.耶穌在快要離世之時,對那班受過祂訓練三年的門徒有個期望,就是要他們同心合意,在基督耶穌裡合而為一.我們服事主的,也要做合神心意的事.甚麼是合神心意的事呢?是不是蓋一間世界最大的禮拜堂呢?還是在敬拜的時候,一同守某一種的儀式呢?這都不是.約翰福音十七章告訴我們,合而為一才是合神的心意.可是今天教會卻忽略了,我們只注意一些次要的事,而且為著那些事彼此爭吵,分門別類,這是合乎神的心意嗎?
\newline
\newline
或許你會問,合而為一的範圍有多大呢?和甚麼人合一呢?聖經給了我們答案.約翰福音十七章有句話說,除了那滅亡之子以外,耶穌都保守了他們.耶穌這篇禱告,是為他的十一個門徒的,被稱為滅亡之子的猶大則是除外.猶大是個未得救而滅亡的人,他代表那些雖然做禮拜,雖然作主的工,卻仍未得救的人.這些人不在合而為一的範圍內,只有真正得救的人才能合一,才應該合一.
\newline
\newline
可是,今天很多自稱為信徒,信仰純正的人卻分門別類,像哥林多的教會一樣,雖然他們信仰是對了,但裡面卻分了四派,有保羅派,亞波羅派,磯法派和基督派.保羅責備他們分門別類,頭一派都是高舉人,所以不對,但第四派是基督派,應該是好的,可是保羅也責備他們?為甚麼呢?因為他們犯了一個錯誤,以為只有他們才是屬基督的,別人都不是,這便是宗派主義,只有他們對,而別人不對,高舉某種看法過於主自己,或者高舉某個屬靈的人,而成為一個派別,這便破壞了教會的合一.求主憐憫我們,認識耶穌基督,祂是神的兒子,是唯一的救主,是教會的元首,惟有高舉祂,傳揚祂,才能合一.
\newline
\newline
保羅是主忠心的僕人,他被主拯救,蒙主呼召,傳揚這位救主,但當時有些人也傳耶穌,動機卻不大對,腓立比書說有些人傳基督並不是出於真心,要造成紛爭嫉妒,但保羅卻能容忍他們,而且感到快樂高興,因為基督的名被傳開了,這是何等大的度量.這樣的傳道人,才是合神心意的忠僕.所以我們講合一的時候,應特別清楚這個真理,要高舉神,不要高舉人,也不要接受人的高舉與特別擁護,惟有這樣,教會才能真正的合一.
\newline
\newline


\newpage

\section{第52屆~港九培靈會~薛玉光~第三講　認識聖靈}
\label{sec:634}
\textbf{第三講　認識聖靈}
\newline
\newline
連結: \href{https://www.hkbibleconference.org/session-message/view/634}{\texttt{https://www.hkbibleconference.org/session-message/view/634}}
\newline
\newline
\hyperref[sec:633]{< < < PREV SERMON < < <}
~
\hyperref[sec:toc]{[返目錄]}
~
\hyperref[sec:636]{> > > NEXT SERMON > > >}
\newline
\newline
約翰福音 16:7-16:14
\newline
\begin{longtable}{cl}
\hline
\hline
章節 & 經文 (和合本修訂版)\\
\hline
16:7 & \begin{tabularx}{0.7\textwidth}{X} 然而，我把真情告訴你們，我去對你們是有益的。我若不去，保惠師就不會到你們這裡來；我若去，就差他到你們這裡來。 \end{tabularx} \\ \\ \relax
16:8 & \begin{tabularx}{0.7\textwidth}{X} 他來的時候，要為罪、為義，為審判，指證世人； \end{tabularx} \\ \\ \relax
16:9 & \begin{tabularx}{0.7\textwidth}{X} 為罪，是因他們不信我； \end{tabularx} \\ \\ \relax
16:10 & \begin{tabularx}{0.7\textwidth}{X} 為義，是因我到父那裡去，你們將不再見到我； \end{tabularx} \\ \\ \relax
16:11 & \begin{tabularx}{0.7\textwidth}{X} 為審判，是因這世界的統治者已受了審判。 \end{tabularx} \\ \\ \relax
16:12 & \begin{tabularx}{0.7\textwidth}{X} 「我還有好些事要告訴你們，但你們現在擔當不了。 \end{tabularx} \\ \\ \relax
16:13 & \begin{tabularx}{0.7\textwidth}{X} 但真理的靈來的時候，他要引導你們進入一切真理。因為他不是憑著自己說的，而是把他所聽見的都說出來，並且要把將要來的事向你們傳達。 \end{tabularx} \\ \\ \relax
16:14 & \begin{tabularx}{0.7\textwidth}{X} 他要榮耀我，因為他要把從我領受的向你們傳達。 \end{tabularx} \\ \\
[1ex]
\hline
\hline
\end{longtable}
請先看約翰福音十六7-14,這裡請特別注意耶穌的幾句說話,第八節,「祂既來了」,這個「祂」就是聖靈.「叫世人為罪,為義,為審判自己責備自己」,「世人」就是不信的人.第十三節,「只等真理的聖靈來了,要引導你們明白一切的真理」.「你們」就是門徒,是已經信主的人.還有第十四節,「祂要榮耀我」；「祂」就是聖靈,「我」就是耶穌自己.請再看使徒行傳一8,「但聖靈降落在你們身上,你們就必從著能力,並要在耶路撒冷,猶太全地,和撒瑪利亞,直到地極作我的見證.」還有約翰壹書四1-6.
\newline
\newline
我曾說過為甚麼這次要講合一的真理和見證,因為神叫我看見教會有此需要,教會卻沒有做到耶穌所禱告的 - 叫教會合而為一.如果教會不能合一,便不能榮耀神；在工作上便不能打勝仗,基督徒的靈性長進也不能正常；所以為教會,為個人,為主的工作,都要在此事上有正確認識.過去兩個早上,我講到「認識真神,認識耶穌」.今天講「認識聖靈」,就是說我們要認識三位一體的神.這是我們合一的一個很重要的基礎,認識父神是這基礎第一塊基石,認識主耶穌是第二塊基石；今天講的認識聖靈,是第三塊基石.我應該先講甚麼叫做認識,這是以往曾講過的.
\newline
\newline
耶穌在約翰十七章所講的話,認識獨一真神,並認識你差來的基督,這就是永生,這「認識」二字有很深的意思,這是指很密切的關係；不是泛泛之交,耶穌所講的認識真神,認識基督指屬靈經歷,不是指屬靈的知識,不過兩者也有關係,由知識可進到經歷.但知識可能只停在知識那裡,所以知識並不等於經歷.屬靈可叫我們心裡得飽足,沒有經歷,我們與神就沒有關係,就不能做到合而為一了.
\newline
\newline
現在我們講「認識聖靈」這個題目.耶穌到世上來,為我們做成救贖工作；祂離世之前,告訴門徒說,祂去,是要差保惠師聖靈來.到五旬節的時候,聖靈降臨,開始了教會時代,也是聖靈工作的時代.今天教會的工作,是聖靈的工作,但教會對認識聖靈的教導卻非常缺乏.教會一年的講道,講聖靈有多少次呢?恐怕五次也不夠,甚至兩次也不夠；講耶穌卻很多,講天父,信心,也很多,但講聖靈就很少了.記念耶穌的降生很隆重,但聖靈降臨節期我們就沒有記念了,甚至連講道時也沒有提.因為教會多講耶穌,少講聖靈,就有很多信徒甚至傳道人不甚認識聖靈.
\newline
\newline
今天教會對聖靈的看法,有兩大錯誤.第一個錯誤就是以為聖靈就是能力,傳道人沒有能力就求聖靈充滿他,使他大有能力.聖靈是不是能力?剛才讀的使徒行傳一8,是否說聖靈就是能力?不是!那裡只說,因為聖靈降臨在你們身上,你們就必得著能力；而不是說聖靈就是能力.聖靈是有位格的,祂能給人能力,祂本身並非能力；所以人不能控制祂,但祂能在人身上作工.能力是沒有位格的.有一首詩歌說,「多禱告,多有能力；少禱告,少有能力」,這是很有問題的,看今天教會的情況就不是這樣.這個問題我們可以從約書亞記得著答案.約書亞攻打艾城失敗了,約書亞跪在耶和華面前禱告,從早上到晚上；多多禱告,但仍沒有能力.神沒有對他說,他須更多禱告就給他能力；而是說,不要禱告,起來.神要他對付罪惡,把亞干除掉,然後就有能力打勝仗.我們禱告時若不注意自己與神的關係,雖然多多禱告,仍然沒有能力的.不要用禱告代替悔改,先把自己的心弄乾淨,才禱告.
\newline
\newline
人把聖靈當作能力是錯誤的.因為聖經沒有這樣的說法.你心中若有罪惡,聖靈就不能與你同在,你雖然多禱告,求聖靈充滿你,聖靈也不能充滿你.
\newline
\newline
第二個錯誤,人以為有神蹟奇事就是聖靈的工作.人對「神蹟」一詞常有錯誤觀念.聖經記載先知,聖徒和耶穌本人,行了許多神蹟,這些神蹟是神的靈所做的.如今我們若見有人用神蹟醫病,就以為有聖靈；沒有醫病,就認為沒有聖靈,這是錯誤的.甚麼是神蹟?神的工作就叫神蹟.男女結婚生子,也是神蹟,不能因為看慣了,就不以為是神蹟.我們若以啞巴能講話就是神蹟,那麼就會產生很大的錯誤.例如我們說神的靈可以醫病,但轉過來說凡是醫病的都是聖靈,這就大有問題了.聖靈在聖經中做了某些事情,我們今天凡看見這些事,就認為是聖靈的工作,這是很危險的.今天很多基督徒把行一點超自然的事的都當作是聖靈,不加以分別,真是糊塗.
\newline
\newline
我們要靠著主恩,根據聖經認識聖靈.耶穌在聖經上說,祂去是差遣聖靈來；耶穌更把聖靈是甚麼?做甚麼工作介紹給我們知道.所以從耶穌的說話認識聖靈,那一定不會錯的.
\newline
\newline
第一,耶穌說,聖靈來了,叫世人為罪,為義,為審判,自己責備自己.聖靈來了,是叫罪人悔改,聖靈在人心工作,叫人看見自己的罪.約翰福音第三章記載耶穌和尼哥底母談到聖靈叫人重生,這是聖靈很重要的工作,可以說是頭一樣的工作.聖靈降臨的時候,使徒行傳第二章記載彼得講道有三千人悔改,是聖靈叫不信的人悔改,歸在耶穌名下.
\newline
\newline
第二,耶穌說,聖靈來了,教導門徒,使他們進入真理,叫門徒明白神的道,神的道就是記載在聖經中的真理.我們讀時會不大明白,因為只憑著知識,除非是得到聖靈的幫助.
\newline
\newline
如果仔細看使徒行傳第二章,就可以見到聖靈降臨時做了這兩樣工作.請看四十一,四十二節說:「於是領受他話的人,就受了洗,那天門徒約添了三千人.」這不是彼得的工作,而是百分之一百是聖靈的工作；還有,四十二節說,都恆心遵守使徒的教訓,彼此交接,擘餅,祈禱.這是對內的栽培,在主話裡勸勉.他們之所以能遵守主道,是有使徒教導,另有聖靈的幫助,他們才能明白.
\newline
\newline
聖靈在教會的工作,是叫許多未信主的人相信主,又叫信主的人更明白真理.
\newline
\newline
第三,聖靈來了,要榮耀耶穌.這是耶穌說的.因為耶穌自己成就了救贖工作,榮耀該歸給這被殺的羔羊.今天有些運動,掛上基督教招牌,其實是異端；他們有一個特點,就是高舉人,高舉創辦人,甚至有時把他寫的書,放在與聖經同等的地位上.例如安息日會就是這樣,該會創始人所寫的書被認為與聖經有同等價值,其他異端如基督科學教會等,也是一樣高舉人.又有一些極端的教派,也是高舉人,甚麼姊妹甚麼弟兄,有甚麼恩賜,就被人高舉起來.聖靈的工作是不高舉人的,只高舉主耶穌.
\newline
\newline
第四,聖靈來了,將神的愛澆灌在我們心中.羅馬書五章五節說,聖靈把神的愛澆灌在我們心中.人墮落以後,愛心成為大問題,人變成自私自利,只愛自己愛的人,不會愛仇敵；但神的愛很奇妙,神叫你愛你所不愛的人.
\newline
\newline
第五,叫我們更恨惡罪惡.使徒行傳五章五節記載,教會中有隱藏的罪惡,聖靈把它顯露出來,而且施行審判.亞拿尼亞,撒非喇夫婦就因此而死了.很多人以為新約時代的神與舊約的不同,新約講恩典,神很寬容,所以馬虎一點也沒不要緊,不像舊約時神那麼嚴厲.利未記十章亞倫兩個兒子作祭司,在聖殿中犯了錯誤,神立刻燒死他們；新約的時代,亞拿尼亞在教會中,在錢財上做了欺騙的事,欺哄了聖靈,結果聖靈叫他們死去,可見神是不變的,祂仍然是一位嚴厲的神.神恨惡罪過,從舊約到新約的都是一樣的.聖靈要審判罪惡,聖靈不是能力,能力不會審判人.
\newline
\newline
今天是末世時代,聖經清楚的警告我們,一切的靈不都可信,總要分辨.有些靈在今天做了許多奇怪的事情,我們應加以分辨.靈所做的工作,所說的話,都要照聖經的記載來分辨.耶穌已告訴我們,聖靈來了,要做甚麼工作,乃是叫罪人悔改.若有甚麼運動稱是聖靈工作.像宋尚節博士,慕迪先生,芬尼先生,現在的葛培理,他們講道,很多人悔改,這是聖靈的工作.因為這些工作,都與耶穌的話相符.又有些運動,有很多人領受神的話,而且按著正意解釋聖經,很多人得造就,我才會承認那是聖靈的工作.又若有一些新興的工作,我就會注意負責的人生活如何?是否比以前更聖潔更誠實,若然,我就會肯定這是聖靈的工作.因為聖靈叫人聖潔誠實.若有教會突然增加了許多人數,但這些人是從別的教會拉過來的；這不能說是教會的增長,乃是偷羊；敢信聖靈決不會做這樣的工作.彼得在五旬節佈道沒有偷羊,保羅也沒有偷羊.所以當有人說某教會大有增長時,你暫且不要開心,可能是假的,因為魔鬼最會賣假貨.當信徒不會分辨靈的時候,牠就會用假的東西給我們.
\newline
\newline
求主幫助我們認識主耶穌所講的那一位聖靈,祂是保惠師,我們想信這位聖靈,祂就會保守我們,在基督耶穌裡合而為一的心；這種合而為一才是有真理基礎的,不會虛假.
\newline
\newline


\newpage

\section{第52屆~港九培靈會~薛玉光~第四講　信服永不改變的聖經}
\label{sec:636}
\textbf{第四講　信服永不改變的聖經}
\newline
\newline
連結: \href{https://www.hkbibleconference.org/session-message/view/636}{\texttt{https://www.hkbibleconference.org/session-message/view/636}}
\newline
\newline
\hyperref[sec:634]{< < < PREV SERMON < < <}
~
\hyperref[sec:toc]{[返目錄]}
~
\hyperref[sec:638]{> > > NEXT SERMON > > >}
\newline
\newline
耶利米書 32:36-32:39
\newline
\begin{longtable}{cl}
\hline
\hline
章節 & 經文 (和合本修訂版)\\
\hline
32:36 & \begin{tabularx}{0.7\textwidth}{X} 現在論到這城，就是你們所說，已經因刀劍、饑荒、瘟疫被交在巴比倫王手中的，耶和華－以色列的神如此說： \end{tabularx} \\ \\ \relax
32:37 & \begin{tabularx}{0.7\textwidth}{X} 「看哪，我曾在怒氣、憤怒和大惱怒中，將以色列人趕到各國；我必從那裡將他們召集出來，領他們回到此地，使他們安然居住。 \end{tabularx} \\ \\ \relax
32:38 & \begin{tabularx}{0.7\textwidth}{X} 他們要作我的子民，我要作他們的神。 \end{tabularx} \\ \\ \relax
32:39 & \begin{tabularx}{0.7\textwidth}{X} 我要使他們彼此同心同道，好叫他們永遠敬畏我，使他們和他們後世的子孫得享福樂。 \end{tabularx} \\ \\
[1ex]
\hline
\hline
\end{longtable}
我們先看幾處經文:(耶三十二36-39；結十一17-21；約十七20-23)「我不但為這些人祈求,也為那些因他們的話信我的人祈求,使他們都合而為一；正如你父在我裡面,我在你裡面,使他們也在我們裡面,叫世人可以信你差了我來.你所賜給我的榮耀,我已賜給他們,使他們合而為一,像我們合而為一.我在他們裡面,你在我裡面,使他們完完全全的合而為一,叫世人知道你差了我來,也知道你愛他們如同愛我一樣.」(弗四1-3)「我為主被囚的勸你們,既然蒙召,行事為人就與蒙召的恩相稱.凡事謙虛,溫柔,忍耐,用愛心互相寬容,用和平彼此聯絡,竭力保守聖靈所賜合而為一的心.」
\newline
\newline
前三天我們已經講過合一基礎的三塊基石,就是認識獨一的真神,認識惟一的救主耶穌和認識聖靈,今天我們要看第四塊的根基石頭,就是信服一本永不改變的聖經.
\newline
\newline
現今,在基督教裡面有很多派別,令人感到困惑.其實只有兩派,一派是信仰純正的,雖然他們有不同的會別名稱；這一派我們稱為福音派,他們的信仰是根據一本可靠的聖經.另外信仰不純正的教派則不只一本聖經,天主教除了信新舊約聖經之外,還有次經和歷代教皇的諭令.摩門教也講一點聖經,但更重要的,是他們自己的摩門經.所有的異端教派,都是不只根據一本聖經的.聖經只有一本,共六十六卷,分舊約和新約兩部份；這是我們信仰的根基,是信徒生活最高的準則.如果我們不信這個,便不能算是純正的信仰.為甚麼聖經只有一本呢?因為一本已經夠了,聖經一開始就說:「起初神創造天地.」講及人類歷史的開始,而最後一卷則論到人類歷史的終結；即耶穌再來,有新天地,直到永遠,聖經是為人類的需要而有的,使我們知道自己的由來,命運,以及今日怎樣生活,將來的結局如何等等,這一切都在聖經內.任何教會所傳講的道理,如果是根據這本聖經的,便是對的；凡越過這本聖經的,都是錯的.如果有人說這本聖經不夠,還有新的啟示,是不在聖經裡的；這樣的講法就是錯誤的,是從另一個靈來的.
\newline
\newline
那麼,聖經是甚麼呢?聖經就是神的話,是絕對的,是永不改變的真理.世人的哲學,思想都會改變,因為要迎合人的需要；惟有聖經的真理,就是神的話,是永不改變的.我們有神的話,何等有福,在動亂的世代中,有一塊穩當,堅固磐石,我們站在這磐石上,永不動搖.現今人類種種邪惡的行為,都是由邪惡的思想而來的；近代西方所稱最進步的思想,就是說世上沒有絕對的事,都是相對的.我們若接受這種理論,便會影響我們對聖經的看法；把聖經的話也看為是相對的,那就錯了.其實,聖經的真理是絕對的,世上有些事物也是絕對的,比方有人打你一下,你必然會覺得痛,而且很想還手,這是個事實,是絕對的.神的話是真理,安定在天,直到永遠.而且事實可以證明,神的話確是真理.神的話經過時間的考驗,幾千年來沒有改變；很多人要否認它,推翻它,但是它仍屹立不動.另外,神的話也經得起考驗.雖然各宗教都有它的經典,佛教有佛經,回教有可蘭經；但這些經都只在某一些地方流行,而且有的還需要藉武力來保護；惟有聖經是傳到全世界的,任何國家,任何地方都有聖經的流傳,而且不須用武力保護.所有接受聖經的人,都說這是真理,有許多人為了傳講聖經而奉獻自己,甚至犧牲自己的性命.這就證明聖經是真理,因為只有真理比生命更寶貴.
\newline
\newline
既是這樣,我們應當用甚麼態度來對待聖經呢?第一,要完全相信.不信聖經,或只信一部份聖經的人,都不是基督徒.我們只能完全相信聖經,或是完全拒絕,斷不能信一半或是一部份.有人說,我不明白聖經,怎能相信呢?這不要緊,你可以先相信,然後才慢慢明白.因為這是神的話,不是人的話；人的話是先明白而後相信；但神的話是先相信而後明白的,這和學識的深淺並沒有關係,問題只在乎我們是否信神的話.人除非先相信,否則永遠沒法明白.另外有人說,我信了,但仍是不明白；這是可能的,因為沒有人能明白全本聖經.雖然這樣,我們仍是相信,因為所不明白的與我們所相信的並無阻礙；而且正因為我們不明白,才證明聖經是可信的,是神的話.這正如一個傑出的太空科學家,很有學問；他有個三歲的兒子,雖然對他父親所學的一竅不通,但是他卻完全相信他的爸爸.我們必須相信聖經都是神的話,好像這小孩子相信他爸爸一樣.我們都是神的兒女,我們相信神的話,並且慢慢明白,慢慢追求.由於不斷追求,我們才感覺到快樂.
\newline
\newline
今天有很多不相信聖經的人,卻領導教會的工作.那些新派的傳道人,從新派的神學院出來,領導信徒.西方的神學界犯了一個很大的錯誤,以為那些會寫大本書,會講一大套理論的人,都是神學家,而不管他是否相信聖經；所以神學家的定義,變得非常混亂.英國著名的牧師史托得,是當代一位很傑出的屬靈的人.他說:「神學就是對於全能的神的知識.」也就是相信獨一真神的知識,才算是神學.這和西方神學界一般人所講的神學不同.我第一次到美國講道的時候,他們正流行一種叫做「神死了」的神學理論,教會有很多神學家便研究這種最流行,最進步的神學思想,真是荒謬之極!一間所謂神學院,要訓練神的工人,卻相信神死了,這種講法怎能稱為神學呢?那不過是掛羊頭賣狗肉的神學,但竟然有人去相信,去接受.這些神學家對神為甚麼有許多不同的講法呢?為甚麼會這樣混亂呢?最大的原因,是他們並不信聖經是神的話.很多神學家根本沒有好好去讀神的話,有的神學家主張神是慈愛的,所以沒有審判地獄,他們說神的責任就是赦免人的罪,所以殺人放火的都可以上天堂.但是聖經卻明明的說,神也是公義的,祂絕對不把有罪的當作無罪,祂是聖潔的神,是烈火.如果我們有讀聖經的話,就不會想出這一套的怪道理來.所以我們今天作基督徒,作傳道人的,要完全相信神的話,否則我們不配稱為基督徒.
\newline
\newline
聖經是神對人的默示,一方面叫人知道神是一位怎樣的神,另一方面神又藉著聖經教導我們,督責我們.神有資格批評我們,但我們沒有資格批評神以及批評聖經；如果我們批評聖經,就是批評神本身了.譬如要批評一個人,那麼自己必須要比人更有學問,更有智慧,至少也和他同等,才有資格批評他.誰能批評真理呢?誰比真理更高明呢?但是今天竟有許多人,自以為聰明,胡亂批評神的話.從歷史上看,所有批評聖經,攻擊聖經的人,都失敗了,沒有一個成功的.這幾年就有很多這樣的人,他們以為自己的言語比聖經更有價值,自己就是神,但結局是失敗的,羞恥的.我們要完全相信聖經.我曾經兩次進神學院,我自己持守一個原則,就是絕對不進一間不信聖經是神的話的神學院.今天西方的神學界非常可憐,有的不信聖經,有的則相信聖經中可能有神的話,這是個非常詭詐的說法,破壞人對聖經的信仰.英國有位出名的牧師司布真說:「我的天父看為愚拙的人,我不敢稱他是聰明的.」神所說的愚妄人是誰呢?就是那些說沒有神,不信神,不信聖經的人.雖然有人稱他們為博士,為神學家,但他們在神眼中是愚妄人,我們怎能稱他們為聰明的呢?
\newline
\newline
除了相信聖經之外,我們還要遵行聖經.聖經說:「你信有神,你信的不錯,但鬼魔也信.」魔鬼也是信有神的,但是他不順服神,不遵從神的話.所以我們除了相信聖經,還要遵從聖經,而且要遵從聖經中一切的話.雖然在今天邪惡的世代中,要遵守聖經的話很難,但卻不是不可能的.神的話要完全遵行,不能單選擇某些方便的便遵行.摩西臨終時,把領袖的責任交給約書亞,對他說,你要謹記神的話,這本律法書,要常常誦讀,書上一切的話都要遵行.我們遵守神的話,不可單遵守一部份,正如相信祂的話,也不能單信一部份,乃要全部遵行.舉例說,十條誡命我都遵行,獨有孝敬父母這條例外,這樣便是違背神命令的人.我們如果只遵守一部份,這不能叫主的心喜悅.而且,我們也不能用這一件來代替另一件.如果我們需要悔改就要悔改,絕對不能逃避,或拚命以禱告來代替悔改.雖然認罪悔改並不好受,會丟我們的面子,但我們所需要的卻正是這個,不能以別的來代替.很多人做不法的生意,以為奉獻金錢便可以彌補罪過,可以博取神的通容,但神是不受騙的.我們必須切實遵行神的話,才能得神的喜悅,即使是付代價,也必須這樣行.
\newline
\newline
聖經很多地方都吩咐我們合一,舊約的先知耶利米和以西結都有提到.以色列的十二個支派,因為得罪神,所以被神分散到全世界,但神並不長久發怒,祂要把他們從各處招聚回來,而且仍然是一個民族.神歡喜祂的子民合而為一,士師時代他們確互相殘殺,但是神卻要他們重新成為合一的民族.新約時代,耶穌最後一次為門徒禱告時,從半段是為今天的教會禱告,而不單為當時十一個門徒禱告.(約十七20)「我不但為這些人祈求,也為那些因他們的話信我的人祈求.」「這些人」是指十一個門徒,這些門徒在五旬節以後便被差出去,建立教會,一代接一代的直到今天.所以耶穌的禱告是包括我們的,祂要我們在祂名下的人,都合而為一,像父子聖靈合一一樣.以弗所書所指的,特別是外邦的信徒,也要合而為一,這是神的心意,也是聖經的真理.
\newline
\newline
有人說,今天很多人都說信聖經,也說按正意解釋聖經,但彼此仍然有差別,應當怎樣去解決呢?如果是按正意解釋聖經的,當留意幾個要點.第一,是榮耀神的,若是這樣的解釋不能榮耀神,就不是按正意.第二,叫人得益處,我們遵行神的話,如果沒有偏差,一定叫人得益處.第三,共同的承認.凡是有生命的基督徒,有聖靈同在的人,內心都能印證這是對的,是出於神的.根據這三個要點,我們可以分辨到底是否真的按正意來解釋.合而為一是信徒應該實行的,這是聖經的教訓,是照正意的解釋,並且可以叫神得榮耀,叫遵守的人得益處.
\newline
\newline


\newpage

\section{第52屆~港九培靈會~薛玉光~第五講　神學信仰的合一}
\label{sec:638}
\textbf{第五講　神學信仰的合一}
\newline
\newline
連結: \href{https://www.hkbibleconference.org/session-message/view/638}{\texttt{https://www.hkbibleconference.org/session-message/view/638}}
\newline
\newline
\hyperref[sec:636]{< < < PREV SERMON < < <}
~
\hyperref[sec:toc]{[返目錄]}
~
\hyperref[sec:639]{> > > NEXT SERMON > > >}
\newline
\newline
耶利米書 2:36-2:36
\newline
\begin{longtable}{cl}
\hline
\hline
章節 & 經文 (和合本修訂版)\\
\hline
2:36 & \begin{tabularx}{0.7\textwidth}{X} 你為何東奔西跑改變你的道路呢？你必因埃及蒙羞，像從前因亞述蒙羞一樣。 \end{tabularx} \\ \\
[1ex]
\hline
\hline
\end{longtable}
請先看兩處的經文:(耶二36)「你為何東跑西奔要更換你的路呢?你必因埃及蒙羞,像從前因亞述蒙羞一樣.」(耶六16)「耶和華如此說:你們當站在路上察看,訪問古道,那是善道,便行在其間；這樣,你們心裡必得安息.他們卻說,我們不行在其間.」
\newline
\newline
過去四天我們已講過教會合一的真理,就是關於我們對三位一體神正確的認識,對聖經全部的相信和遵守,這些話都告訴我們,應該在基督耶穌裡合而為一.既然我們認識了真理,便應當遵行.我們每次聽道時應該注意,我們聽道是有責任的,耶穌說,多給誰便向誰多要；今天教會有一些聽道專家,也有些遊行聽道家,所有好的道理都聽,但靈性卻沒有長進,因為他聽的道理多,愈聽便愈驕傲,而且還會批評別人講道.聽道是要行道的,我們明白真理,真理要求我們去遵行,不然明白道理也是無用的.既然聖經教導我們要在基督耶穌裡合一,我們便該去實行.但怎樣實行呢?從誰開始呢?這是一個問題.而且,要達到合一的地步,不是人的力量,乃要靠聖經的能力,聖經將合一的心放在我們裡面,又幫助我們克服一切攔阻合而為一的東西.所以從今天開始,我們要怎樣達到合而為一,和誰應該做這事.如果我們用自己的聰明,而達到合一,必然是做不到的.
\newline
\newline
過去多年來,教會都看到合一的需要,很多的文章和講道,都提到我們應該怎樣合而為一,更有些教會領袖,開始進行促進教會的合一.有的用基督教的名義,組織一個很大的會,是世界性的,來協進教會的合一.初開始時這會的信仰還不錯,但後來卻變了質,新派的人也加入了,直到今天,這個會並沒有實際協助教會達到合一,只是把一些基督教的教派拉在一起做會員而已.這不是神所要求的合而為一.另一個組織,是立一些基要信仰的信條,凡是贊成的,都在信條下簽名,加入為會員.但是,他們所注重的是信條,不是耶穌基督的生命,結果也達不到合一的目的.有時我們想藉一個屬靈領袖,能在眾教會中最有號召力的,組織一個會,促進教會的合一,但這些仍是人的方法,不是神的方法.這樣的合一,沒有穩固的基礎,不符合神的旨意.真正的合一,是應該先有真理的基礎,我們從神的話中明白甚麼是合一,有甚麼方法能達到合一,然後跟著去做,便可以成功.
\newline
\newline
當我們實行合而為一的時候,誰應該開始呢?我認為應該由神學開始,因為神學對教會的關係非常重大,神學院是栽培傳道人的,神學院怎樣,將來的教會也是怎樣.西國有句話說:「甚麼樣的傳道人,就帶出甚麼樣的教會.」我們也可以說,今天中國的神學院如何,未來的中國教會也是如何.所以從事神學教育的人,有很重大的責任.神學是決定教會的；神學的信仰,影響到教會的信仰；神學生的生活,影響到將來教會的生活,神學教育的成功或失敗,也直接影響到將來教會聖工的成功或失敗.我深深的覺得,神學院應該作眾教會的榜樣,而這榜樣應該是好的榜樣.所以從事神學教育的人,責任非常大,要辦神學院.作神學的董事或信託委員,責任也非常重大,如果不好好的做,在神面前要負很大的責任.
\newline
\newline
神學教育有甚麼責任呢?神學院的職責是甚麼呢?我覺得這時代的神學院,就像神的先知一樣,在所處的時代中,代替神向當代的人說話.神學是領導教會,教導真理的.神學所教導的真理若有偏差,便會大大的影響教會.還有,辦神學是應該為神而辦,因為教會是神的,不是人的組織,所以應該單純是為神而辦,求神的榮耀,而不是為宗派或為著人.神學也應該是屬於神,高舉神的,要絕對順從真理,不然便失去它的意義,不能對神的國度有貢獻,甚至會對教會帶來害處.神學院最重要的是信仰,信仰應該放在第一,彼得後書第一章說先有信心,然後有德行,知識.信心就是人與神的關係,也就是信仰,這是最重要的.如果信仰不清楚而追求知識,那知識便有問題.今天特別是西方的神學院,比較著重知識的追求.因為西方的文明,是注重科學技術的,這是世上的東西.當教會效法世界時,世界的東西便跑進入教會,使教會信仰墮落,生活不聖潔.神學院的信仰偏差,便會效法世界,注重知識,甚至於改變信仰,不講信仰.所以信仰在神學中是最重要的.
\newline
\newline
隨信仰而來的,便是生活.先有信心,後有德行,德行就是生活.若單講信仰而無生活,這信仰便是空的,是死的.我們有稱為死的正統派,有些教會信仰沒有錯,但裡面卻死氣沉沉.神學院若單講信仰是不夠的,他的生活,行事為人,是否照著他的信仰呢?有了信仰和生活,然後才加上知識便很有用處.
\newline
\newline
今天中國人的神學院,一般都是跟著西方的神學院跑,因為最早的神學院是西差會所辦的,跟著才有憑信心辦的神學,而這種神學院是佔少數的.現今在神學教育中負主要責任的人,多半是在西方神學院受過訓練的；那些留美,留英的博士,站在神學院最高層的領導位置上.試問西方神學院的光景怎樣呢?就像耶利米書第二章所講的,他們是東奔西跑的信仰方向不一定.先是西方哲學影響了神學,有些所謂神學家,只不過是哲學家,他們把聖經中的一些意思斷章取義,滲入政治的思想.講出一套理論,更成為當代的前進神學,這是跟西方跑.但今天的西方人,因為心靈空虛,又厭煩了西方的哲學,便轉向東方.美國有很多人追求東方的哲學,接受東方的宗教信仰,印度教在美國非常吃香,這就是東奔西跑的做法.我不是說所有的神學都如此,而且一般的神學,能持守純正的,少而又少!
\newline
\newline
現今西方的神學,就像聖經的士師時代,以色列中沒有王,各人任意而行.德國的神學家有一套,英國的神學家有一套,美國的神學家又有一套,非常混亂.前幾年菲律賓有位青年,在中國人的聖經學院中讀了兩年,對神學院不滿意,覺得他們的方法太古老,神學院內的生活訓練太束縛,所以便跑到美國去,兩年中轉換三間神學院,後來便不讀了,因為看見西方神學信仰太混亂.這是個事實,如果你經歷多幾間西方的神學便會看見這樣的情形.因為信仰上混亂,所以根本無法合而為一.今天的宗派主義,高舉宗派過於耶穌,也是從西方來的,是西教士傳入來的.西教士來中國傳福音,我們欽佩他們的犧牲,感激他們的愛心,但同時他們也帶來了一些錯誤,這並不是我一個人這麼說,很多西教士自己也承認,宗派主義的思想,是從西方來的.所以很多神學院,都有宗派主義的觀念,還有各種不同的神學思想所產生的工作,方法也常改變.中國領導神學教育的人,從西方訓練回來,把西方那一套的東西全搬回來,而沒有自己的東西,他們所教的一些新的東西,都是從外國來的.中國教會提倡的一些運動,都是從西方學來的,沒有一種運動,是中國傳道人自己看聖經,從聖經發現的.
\newline
\newline
神學院和神學教育的地位,就像先知與使徒一樣.舊約的先知,新約的使徒,都是像今天辦神學的人一樣.先知以利亞,以利沙都訓練過學生,先知撒母耳也辦過先知學校；耶穌在世上訓練過十二個使徒.保羅是外邦使徒,他訓了提摩太和提多,這是聖經中神學的樣式,是聖經為我們留下的模樣,都是注重信仰,和與信仰符合的生活,注意聖潔和敬畏神.因為這緣故,這些辦神學的人,能影響全國的屬靈空氣,保羅影響當時的教會,撒母耳領導當時的以色列復興；但今天的神學院,因為信仰變了質,所以便失去了該有的地位.神學院應該領導教會,指責世界的罪惡,但西方的神學院不能領導教會,而且破壞了教會.
\newline
\newline
另一方面,因為神學院沒有注重信仰,所以不能訓練出合神使用的工人.在整個屬靈的工作上,舊約時代是看先知,新約時代是看使徒.今天神的教會需要屬靈的領袖,好像先知和使徒的一樣.但今天很多神學院所訓練出來的人,不是先知,也不是使徒,而是文士.他們懂得一點聖經的知識,但沒有屬靈的經歷,沒有敬畏神的心；和神的關係不清楚,所以不能使教會走上合一的道路.
\newline
\newline
聖經告訴我們,要尋找古道,就是古老的一條道路.教會若要復興,便要走古老的路,不要走新派的路.科學的技術,愈新愈高明,但信仰則愈古老愈高明.新派就是不信派,聖經中所講的新神就是假神.今天我們要達到合而為一,首先要注意神學的信仰,是基要的信仰,是根據聖經所講的基要信仰,而且把它放在最重要地位.其他如儀式,制度,或解釋預言的亮光等,都是次要的,不應該成為合一的攔阻.今天很多神學院說明是為宗派辦的,他們非常看重該宗派一些特別的儀式和講法；但也有一些沒有高舉他們的宗派特點,而他們所設立的神學院,是為眾教會辦的,這樣就能促進教會的合一.不過,另外有些不是由宗派辦的神學院,也可能太著重自己的特點,自以為最屬靈,所以不能與別人合作,其實他們在無形中已成為一個的宗派.
\newline
\newline
所謂在基督耶穌裡的合一,並不是指所有的事情都要相同,乃是指生命的相同.正如一個家庭有五個兄弟,他們都是同一對父母所生的,雖然他們的容貌,高矮,肥瘦各有不同,但他們是自然合一的,並不需要組織五兄弟合一的促進會,外面的不同並不影響兄弟間的合一.教會的合一也是這樣,是生命的合一,而神學信仰直接關係到屬靈生命.神學院需要信仰相同,才能合一.今天一個很可惜的現象,就是應該合一的不願合一,不應合一的卻合在一起.福音派的神學院,都是信基要真理的,但常常不能合作；可是另一方面,有些福音派的神學院竟和新派的神學院合作,接受福音派的信仰,但卻參加不信派的活動,這就是不應該合一的,卻合一起來.
\newline
\newline
今天,神學院應該注意基要的信仰,至於次要的東西,不要高舉.福音派的神學院應該合一,彼此密切合作,因為他們都是為神而辦,為訓練神的工人而辦神學.求主幫助我們,在實行合一之時,神學院首先走在前面,作教會的好榜樣.
\newline
\newline


\newpage

\section{第52屆~港九培靈會~薛玉光~第六講　信徒生命的合一}
\label{sec:639}
\textbf{第六講　信徒生命的合一}
\newline
\newline
連結: \href{https://www.hkbibleconference.org/session-message/view/639}{\texttt{https://www.hkbibleconference.org/session-message/view/639}}
\newline
\newline
\hyperref[sec:638]{< < < PREV SERMON < < <}
~
\hyperref[sec:toc]{[返目錄]}
~
\hyperref[sec:641]{> > > NEXT SERMON > > >}
\newline
\newline
約翰一書 1:1-1:4
\newline
\begin{longtable}{cl}
\hline
\hline
章節 & 經文 (和合本修訂版)\\
\hline
1:1 & \begin{tabularx}{0.7\textwidth}{X} 論到從起初原有的生命之道，就是我們所聽見、所看見、親眼看過、親手摸過的－ \end{tabularx} \\ \\ \relax
1:2 & \begin{tabularx}{0.7\textwidth}{X} 這生命已經顯現出來，我們看見了，現在又作見證，把原與父同在，並且向我們顯現過的那永遠的生命傳揚給你們－ \end{tabularx} \\ \\ \relax
1:3 & \begin{tabularx}{0.7\textwidth}{X} 我們把所看見、所聽見的傳揚給你們，為要使你們也與我們有團契，而我們的團契是與父和他兒子耶穌基督所共有的。 \end{tabularx} \\ \\ \relax
1:4 & \begin{tabularx}{0.7\textwidth}{X} 我們把這些事寫給你們，使我們的喜樂得以滿足。 \end{tabularx} \\ \\
[1ex]
\hline
\hline
\end{longtable}
請看(約壹一1-4)「論到起初原有的生命之道,就是我們所聽見,所看見,親眼看過,親手摸過的,(這生命已經顯現出來,我們也看見過,現在又作見證,將原與父同在,且顯現與我們那永遠的生命,傳給你們).我們所看見,所聽見的,傳給你們,使你們與我們相交,我們乃是與父並祂兒子耶穌基督相交的.我們將這些話寫給你們,使你們的喜樂充足.」(約壹四20-五3)「人若說我愛神,卻恨他的弟兄,就是說謊話的.不愛他所看見的弟兄,就不能愛沒有看見的神.愛神的,也當愛弟兄,這是我們從神所受的命令.凡信耶穌是基督的,都是從神而生,凡愛生他之神的,也必愛從神生的.我們若愛神,又守他的誡命,從此就知道我們愛神的兒女.我們遵守神的誡命,這就是愛他了,並且他的誡命不是難守的.」(腓二1-2)「所以在基督裡若有甚麼勸勉,愛心有甚麼安慰,聖靈有甚麼交通,心中有甚麼慈悲憐憫,你們就要意念相同,愛心相同；有一樣的心思,有一樣的意念,使我的喜樂可以滿足.」
\newline
\newline
第一,我要講的是信徒生命的合一,是和我們是有切身關係的.教會的合一,不是靠開一兩次的大會,或是成立一個長久性組織,乃是從生命方面入手.基督徒是有生命的,一切的表現都從生命而來,而合一也是生命必有的表現之一.首先我們要注意的,是聖徒在基督裡的合一,是應當的,是必須的,因為這是主的命令.主的命令是要我們彼此相愛,這是合一的表現.基督徒在主裡面彼此相愛,就能達到合一；能夠以基督的愛對待基督裡的肢體,合而為一便不成問題.還有,這也是神的心意,耶穌在約翰福音十七章的禱告中,再三的求父保守門徒合而為一；如果我們不合一,便是叫主傷心.
\newline
\newline
第二,信徒在基督裡是能夠合一的,因為我們有神的生命.神是生命之源,這生命是愛的生命,而我們是一同領受這生命的.這生命在本質上是合一的,所以合一是很自然的事,不須自己追求,只須保守合而為一的心.當人初信主之時,生命非常幼稚,需要不斷長大；耶穌說:「我來了,是要叫人得生命,並且得的更豐盛.」但教會今天的困難,是很多人做掛名的基督徒,沒有生命,這成為教會合一的最大的攔阻.還有一些基督徒,以得救重生為滿足,不再追求,這是很大的錯誤.我們需要不斷長進,達到豐盛的生命,讓基督的生命在我們裡面掌權.一個有豐盛生命的人,與人合一是沒有問題的,沒有豐盛生命的人,與人合一才會有問題.
\newline
\newline
在一些宗派教會作領袖的人,有的是有美好靈性,但也有的卻沒有.他們只注重高舉自己的宗派,藉著某些條件,或是辦事能幹,或是有外國得來的學位,可以作高層的領導者,這些人常常是靈性枯乾的.但另一方面,有些作信心差會的領導人,他們有豐盛的生命,有很好屬靈經歷,他們能以愛心對待從各宗派而來的宣教士,雖然各人背景不同,卻能為同一的目標爭戰.這並不是容易的事,必須各人有合一的態度,特別是領導的人,所以我們做基督徒的,做教會執事長老的,做神學院院長的,要注意我們自己的靈性,是否有豐盛的生命.最好的測驗便是看我們的愛心,我們能否和許多背景不同的人相處和睦,是否有謙卑,愛人的態度.
\newline
\newline
第三,我們要思想的,信徒生命合一的見證是甚麼,我們要在那些事情上表現我們是真正合一的.約翰一書多次提到相交,相交就是團契,在主裡面彼此有交通.合一的見證,就是與別人有交通.在甚麼事上有交通呢?
\newline
\newline
(一)在禱告上.(羅十五30)「弟兄們!我藉著我們主耶穌基督,又藉著聖靈的愛,勸你們與我一同竭力,為我祈求神.」今天教會的禱告能否表現合一的見證呢?我覺得仍需要追求.今天的教會在禱告上有缺點.(A),禱告會只是一種例常會,不是真正的禱告.一個鐘頭的禱告會,有時講道已佔了一半時間,會眾參加禱告會是為了聽道,而不是禱告.這是錯的.禱告會應注重禱告,讓每一位參加的都有機會會禱告.教會今天對信徒的禱告不夠操練,所以我們禱告很弱,很多人沒有信心禱告的經歷,這是個缺點.(B),我們的禱告常以自我為中心,只為自己禱告.其實這不是禱告,而是作宣傳,叫別人都聽他奇異的經歷.還有,教會的禱告會都是為自己本身的需要禱告,為自己的家庭,為自己的教會,這不是合一；乃是自私自利,令神的豐富恩典不能彰顯.很少人為別的宗派代禱,各人都活在自己的圍牆裡面.耶穌教導門徒禱告時說,我們在天上的父,神是放在第一；願人都尊你的名為聖,都是稱父的名,是父的事情.我們應該為神的眾教會禱告,因為教會是屬神的國度.如果我們要合一,我們禱告的範圍先要擴大,而且和眾聖徒有禱告中的團契.(C),我們通常的禱告,都是要傳道人為信徒禱告,但有多少信徒為傳道人禱告呢?我們提到牧師,常常是說批評的話.使徒行傳第四章記載教會的公禱,那是全教會的信徒為傳道人禱告.我們今天為牧師禱告多少呢?今天眾教會不能彼此合一,原因是我們沒有彼此代禱.要達到合一,必須先由禱告開始,禱告應該彼此相通,而不應分門別類,更不應自己只顧自己的事,也要顧別人的事.
\newline
\newline
(二)信徒的相交,應在奉獻錢財上,這是個好的見證.(腓四16)「就是我在帖撒羅尼迦,你們也一次,兩次的,打發人供給我的需用.」今天教會很多的工作需要錢,我們的錢應該為神而用,為教會而用.可是為哪一間教會呢?很多大宗派要把信徒的錢集中起來歸總會支配.有一次我在一本英文的基督教雜誌上看到一篇消息,標題是「打死宣教士的子彈是用信徒的錢.」有一個講合一的世界性的大協進會,很多大宗派都是他的會員,他的錢非常多,可是掌握的卻是一些新派的人.今天西方的神學院,走上離經背道的路.他們說不需要派宣教士傳福音叫人悔改,因為耶穌來乃是救我們脫離政治的捆鎖,所以教會要幫助人得解放.非洲有些游擊隊,都是由當地的革命份子領導的,而這個世界性的協進會,竟一次又一次撥大筆的款項支持他們,供他們買武器,去奪取政權.有一天他們取得了政權,這些教會的牧師,宣教師都被他們抓入監牢,加上反革命的罪名,有的甚至被槍斃.這些錢是從那裡來呢?是教會供給他們的.教會今天因為信仰的錯誤,竟到了這樣可憐的地步.因為這緣故,最近有幾間教會已經退出了這世界性的組織.教會的錢要用得適當,要用在神的聖工上.很多大宗派的教會,因為錢多,便花很多在一些表面好看的工作上,但真正傳福音的工作,卻少人支持.有些傳道人告訴信徒,他們的錢只可奉獻給自己的教會,不能送到外面去.但今天有很多戰略性的福音工作,就如福音廣播,學生工作,神學教育,文字工作等,都不是一間教會的力量所能辦到的；論人力,財力,沒有一間教會能做得好,要使工作有效,必須靠所有信徒的支持.我們支持了多少呢?今天我們要支持這些工作很容易,因為這些工作大部份都在香港,我們可以在人力和財力上支持他們,以此達到合一的美好見證.
\newline
\newline
(三)我們應該在靈裡有交通.(羅十五29)「我也曉得去的時候,必帶著基督豐盛的恩典而去.」保羅寫羅馬書的時候,還未到過羅馬,他和當地的基督徒並不相識,可是他有個心願,就是有一天要到羅馬去,和那裡的信徒有靈裡的交通.保羅本來是開荒佈道的人,他不願在基督的名被傳過的地方傳福音,但他想念羅馬的基督徒,所以他寫這封羅馬書；並透露他極願到羅馬去,和他們交通的心願.初期的教會,明顯有合一的見證,新約的書信,有很多是眾教會輪流宣讀的.保羅在歌羅西書四章十六節那裡,叫他們把這信也給老底嘉的教會讀,而他們也要讀從老底嘉來的書信.可見當時眾教會在屬靈供應上是相通的,我們也應活出合一的見證,因為神把他屬靈的恩典賜給眾教會.一個教會若要長進得好,需要和別的教會有交通；除非信仰不同,否則必然可以交通；而且別的教會也必然有一些東西是你所沒有的.雖然你或許比他長進,比他屬靈,但是你並不完全.我國古語說:「智者千慮,必有一失；愚者千慮,必有一得.」在屬靈的追求上,也是一樣,所以聖經教導我們,不要看自己比別人強,因為有很多東西是別人有而你沒有的.我們應該找一些屬靈的朋友,可以彼此交通,彼此代禱,彼此扶持.
\newline
\newline
作為一個基督徒,除了讀聖經之外,也要選擇訂一些屬靈的刊物,得到屬靈的供應,還應該有一些刊物,是能夠告訴你一些遠地的傳福音情況.如香港學生工作,香港的福音廣播錄音室,他們經常都有出版的刊物；我們應該看這些刊物,從其中得幫助,也為他們的需要禱告,為他們的工作奉獻.基督徒需要和別人交通,傳道人也同樣需要.保羅是大使徒,但他經常要求信徒在禱告上支持他.很多人以為傳道人禱告比較有效,會督禱告更有力,其實不是的.我們的禱告,在神面前是平等,是一樣的.如果人心中隱藏罪惡,雖然他是牧師,是會督,但禱告依然無效；如果一個目不識丁的平信徒,但他內心聖潔,他的禱告便大有功效.所以我們若要合而為一,便應當在屬靈的恩典上,和別人有交通.特別是傳道人,應該在合一的見證上,作為信徒的榜樣,因為傳道人與傳道人之間,往往更難合一,更難相愛.但我們必須靠著神的恩典,去遵守這合一的命令.
\newline
\newline


\newpage

\section{第52屆~港九培靈會~薛玉光~第七講　教會事奉的合一}
\label{sec:641}
\textbf{第七講　教會事奉的合一}
\newline
\newline
連結: \href{https://www.hkbibleconference.org/session-message/view/641}{\texttt{https://www.hkbibleconference.org/session-message/view/641}}
\newline
\newline
\hyperref[sec:639]{< < < PREV SERMON < < <}
~
\hyperref[sec:toc]{[返目錄]}
~
\hyperref[sec:642]{> > > NEXT SERMON > > >}
\newline
\newline
帖撒羅尼迦前書 1:9-1:10
\newline
\begin{longtable}{cl}
\hline
\hline
章節 & 經文 (和合本修訂版)\\
\hline
1:9 & \begin{tabularx}{0.7\textwidth}{X} 因為他們自己已經傳講我們是怎樣進到你們那裡，你們是怎樣離棄偶像，歸向神來服侍那又真又活的神， \end{tabularx} \\ \\ \relax
1:10 & \begin{tabularx}{0.7\textwidth}{X} 等候他兒子從天降臨，就是神使他從死人中復活的那位救我們脫離將來憤怒的耶穌。 \end{tabularx} \\ \\
[1ex]
\hline
\hline
\end{longtable}
請先看幾處經文:(帖前一9-10)「因為他們自己已經報明我們是怎樣進到你們那裡,你們是怎樣離棄偶像歸向神,要事奉那又真又活的神.等候他兒子從天降臨,就是他從死裡復活的,那位救我們脫離將來忿怒的耶穌.」(林前四6-7)「弟兄們,我為你們的緣故,拿這些事轉比自己和亞波羅,叫你們效法我們不可過於聖經所記,免得你們自高自大,貴重這個,輕看那個.使你與人不同的是誰呢?你有甚麼不是領受的呢?若是領受的,為何自誇,彷彿不是領受的呢?」還有(林前十二12-20).
\newline
\newline
今天我們要思想一個非常重要的題目,就是教會事奉的合一.我們每一位都是屬乎主的教會,雖然所屬教會的名稱不同,卻同是一個屬靈大教會,就是主的身體.教會就是神的家,是屬於我們的；而教會的工作,包括對內的栽培和牧養,對外的傳福音,我們應當有份.第一,神的要求就是要我們事奉祂.神救了我們,把我們仍放在世上,就是為要事奉祂.帖撒羅尼迦前書那裡講得很清楚,我們以前是拜偶像,屬於偶像的；但當我們離棄偶像,歸向真神之後,便要事奉那位又真又活的神；等候祂兒子從天降臨,救我們脫離那極大的忿怒.神為了要用我們,才救我們,所以基督徒要作一個有用的人,要作一個服事神的人.服事神並不是重擔,也不是責任,乃全然是恩典,是特殊權利,絕對不是被神勉強的,這個我們必須記住.不信神的人,即使有再好的才幹,也沒權利在神的教會中工作,這是神兒女的特權.
\newline
\newline
還有,不服事神的人,是很大的虧欠.(耶四十八10),「懶惰為耶和華行事的,必受咒詛,禁止刀劍不經血的,必受咒詛.」一個不為神工作的人,就是懶惰的人,神的百姓不應懶惰,因為神的百姓就是神的軍隊.以色列百姓出埃及的時候,聖經說；那一天,耶和華把他的軍隊從埃及帶出來.以色列每一個男丁都是軍隊裡的軍兵,教會也是神屬靈的軍隊.新約很多地方說我們要打屬靈的仗,要穿戴神所賜的全副軍裝,這是對每一個基督徒說的,不單是對傳道人說.當兵的人在打仗時若不打仗,便要受軍法的審判,非常嚴重；同樣,一個不為神爭戰,懶惰不事奉的人,聖經說他必受咒詛.耶穌在比喻中提到的那個懶僕人,他領了錢之後,連想也不想,就把錢埋在地裡,結果他被主人丟在黑暗裡.可見事奉是每一個基督徒都必須的,我們有神所賜的恩賜,便應該用來事奉神,而且要把事奉看作是恩典,是特別的福氣.
\newline
\newline
第二,我們的事奉應該有合一的見證.教會合一才是正常的,若分裂則反常.我們若在事奉上不能合一,便不能榮耀神.今天很多神的工作,不但不能合一,而且彼此明爭暗鬥,這是非常羞辱神的.教會往往有過這樣的事,基督徒為了爭禮拜堂打架而被刊登在報紙上,成為大新聞的也有.社會上打架的人到處都有,而且打到頭破血流,但報紙不登出來；而教會只不過打一次架,報紙便把它列為大新聞,因為這實在是太反常,羞辱神的事.我們是神的兒女,我們在教會中事奉,應該合一,合一才能榮耀神,才有力量.如果我們相咬相吞,即使沒有敵人來攻打,也會自己消滅自己.所以合一是必須的.
\newline
\newline
第三,我們要思想怎樣來達到合一.有三樣事情要注意的.
\newline
\newline
(一)靈性和工作的問題,二者應該有平衡.但有些教會只講靈性的追求,要多讀經,禱告,要作屬靈的人,但不注意工作,不向外發展.他們還解釋說,這是重質不重量.耶路撒冷的教會,開始的時候,一次便有三千人加入教會,另一次則有五千人加入,很快便由一百多人變成一萬多人.這不是彼得的工作,乃是聖靈的工作.聖靈在教會中作栽培的工作,使人靈性長進,又使得救的人數增加.另有一些教會則拚命工作,用各樣新的方法來工作,開研討會,但卻忽略了與神的關係,忽略了往下扎根,這也是錯的.二者必須有很清楚的關係.當保羅在沉船遇險的時候,他沒有害怕,還安慰同船的人.他說:「我所屬,所事奉的神.」他把屬神放在第一,事奉神放在第二,二者都需要.教會若要合一,便要注意靈性的長進,和工作的發展；把與神的關係放在第一,然後根據神的心意去工作,便很自然達到合一.
\newline
\newline
(二)我們所擔任的工作,是神工作的一部份,而不是全部.我們要明白屬靈工作的整體性,才能達到合一.神把各樣不同的恩賜給不同的人,但這些人都屬於一個身體,就好像身體上的肢體,不是完全獨立的,乃是彼此需要,又彼此合作的.這樣合作不是一時的,而是從有生命就開始,便要合作,直到生命的結束.譬如說,一個人眼睛有毛病,需要戴眼鏡,這本來是眼的事；但是眼鏡卻偏偏不是掛在眼睛上；而是掛在耳朵上,和被托在鼻子上；這就是肢體彼此之間的關顧,是需要密切合作的.很多人為主工作,以為自己做的最重要,而別人的都不重要,這就成為不能合一的原因.如果我們知道自己所負擔的工作,不過是主工作的一部份,需要別人的支持,幫助,合作,便很容易達到合一.
\newline
\newline
有些不是教會單獨能做卻又是很重要的工作,像福音廣播,學生工作,文字工作,神學教育等,都需要眾教會的支持.有人有廣播的技術,有預備節目的恩賜,但卻缺乏錢,需要別人奉獻.他們不能單靠自己的才幹,更需要別人的禱告；其他的工作也是一樣,我們是彼此都有需要的.教會應當看這些是他們的工作,因為這是主的工作,而教會是主的教會,所以關係是很密切的.從事福音工作的人,也應該知道他們需要教會的支持,他們不能代替全教會,他們也不是教會,他們自己也應該有所屬的教會.他們的工作只是幫助教會的,他們需要教會的支持,而教會也應該積極的參與和支持.
\newline
\newline
(三)甚麼事情最會妨礙合一呢?當我們不高舉主,只高舉人的時候,合一便受破壞了.保羅責備哥林多教會分黨分派,破壞合一,把人高舉起來；有的擁護保羅,便成為保羅派,有的擁護亞波羅,便成了亞波羅派；一共分有四派,彼此不合作,常有爭吵,保羅和亞波羅對此有甚麼態度呢?保羅說他把這些事轉比他的亞波羅,意思就是說,他教訓信徒的話,他們自己先實行,先應用在他們身上.怎樣應用呢?保羅不要他們擁護,不要他們高舉人.教會今天的毛病也是在此,看見某人被神重用,便高舉他,以他為誇口,於是便破壞了合一.他們不明白某人之所以被神重用,是因為神用他,而不是他有甚麼了不起.宋博士被神重用,不是他本身的了不起,而是有千萬人為他禱告.葛培理被神重用,也不是他有甚麼特別,乃是全世界有數不盡的人為他禱告,這些禱告的人,和葛培理的工作是合而為一的.所以信徒和做教會工作的人,不要高舉任何一個人,而那些屬靈和特別有恩賜的人,也不要接受別人的高舉,不高舉人已經不容易,不接受人的高舉更不容易.耶穌變餅給五千人吃飽之後,許多人擁護他,甚至強迫他作王；但耶穌不接受,他們退到山上去.
\newline
\newline
一個真正屬靈的人,是不接受人擁護的,因為當人擁護他的時候,宗派便開始了,合一就會受到破壞.高山是怎樣形成的呢?是泥土裡重疊泥土而上的,到最高的時候,只有一點,所以愈高愈孤單.而且愈是高山的泥土,愈沒有用,不能種東西.但大河流的河底是非常低的,有很多水,大大小小的支流都匯到河裡,成為大江河,江河所到之處,澆灌土地,使樹木生長.我們不要像山一樣,愈爬愈高,沒有真正的用處；乃要像江河一樣,能容納許多的水；當有神豐盛的生命,能滋潤許多的人,而不是高高在上,受人擁護,得自己的榮耀.
\newline
\newline
我們為主工作,在教會事奉,有一個很重要的原則,就是先有生命,然後有生活的表現,跟著才是工作.有屬靈的生命,有美好的生活見證,然後才工作,這工作必然是金銀寶石,而不是草木禾楷的.那些為主所蒙福的工作,能夠經過一百年,二百年,甚至是更長的時間,我們從其中也可以找到這原則.今天,很多人的工作剛好相反,他們先注意工作,但生活沒有見證,生命也大有問題,這樣的工作不能蒙主賜福.
\newline
\newline
英國有間神學院,以前是沒有的,是由三間的福音派神學院合併而成的.這三間神學院雖是福音派,但工作發展很困難,老師不足,學生又少,經費總是不夠.後來,其中一間學院的董事和另一間學院的董事商量,看是否可以合併.經過長時間的禱告,他們便決定合併,最後第三間也加入合併；三間神學院的名稱都沒有了,他們用一個新的名字代替.這新學院開始以後,非常蒙神賜福,教師多起來,除了三間原有的教師之外,還有新的教師加入,學生也多,經費也夠了.後來他們訓練出來的學生,遠超過三間神學院加起來的,他們的見證,也好過三間神學院的見證.這說明了甚麼是合一.為了主的榮耀,我們自己的工作應當隨時可以放下.香港的神學院很多,是否也可以合作,或者實行合併呢?神學應該領導教會,如果是神的意思,如果對神的工作更有幫助,應該願意如此合作.神學先有合一的榜樣,才能帶領教會合一.我們做合一運動,不能建立自己的巴別塔,我們只能順服真理,遵從真理.
\newline
\newline
今天中國的教會,人數少,力量弱,就是因為分門別類,不能合一,所以力量抵銷.我們要緊記,只有耶穌基督,才是惟一領袖.各宗派教會的不同,只是名稱的不同,但都是屬於基督的肢體,都有同一的生命,並且都服事那位復活的主,將來都要受主的審判.求主憐憫我們,在教會的事奉上能有合一的見證.
\newline
\newline


\newpage

\section{第52屆~港九培靈會~薛玉光~第八講　福音工場的合一}
\label{sec:642}
\textbf{第八講　福音工場的合一}
\newline
\newline
連結: \href{https://www.hkbibleconference.org/session-message/view/642}{\texttt{https://www.hkbibleconference.org/session-message/view/642}}
\newline
\newline
\hyperref[sec:641]{< < < PREV SERMON < < <}
~
\hyperref[sec:toc]{[返目錄]}
~
\hyperref[sec:591]{> > > NEXT SERMON > > >}
\newline
\newline
馬太福音 28:16-28:20
\newline
\begin{longtable}{cl}
\hline
\hline
章節 & 經文 (和合本修訂版)\\
\hline
28:16 & \begin{tabularx}{0.7\textwidth}{X} 十一個門徒往加利利去，到了耶穌指定他們去的山上。 \end{tabularx} \\ \\ \relax
28:17 & \begin{tabularx}{0.7\textwidth}{X} 他們見了耶穌就拜他，然而還有人疑惑。 \end{tabularx} \\ \\ \relax
28:18 & \begin{tabularx}{0.7\textwidth}{X} 耶穌進前來，對他們說：「天上地下所有的權柄都賜給我了。 \end{tabularx} \\ \\ \relax
28:19 & \begin{tabularx}{0.7\textwidth}{X} 所以，你們要去，使萬民作我的門徒，奉父、子、聖靈的名給他們施洗， \end{tabularx} \\ \\ \relax
28:20 & \begin{tabularx}{0.7\textwidth}{X} 凡我所吩咐你們的，都教導他們遵守。看哪，我天天與你們同在，直到世代的終結。」 \end{tabularx} \\ \\
[1ex]
\hline
\hline
\end{longtable}
教會唯一的使命,就是向世人傳福音.耶穌在升天之前,把這使命交給門徒,祂說:「天上地下所有的權柄都賜給我了.」那就是說,祂是得勝的大君主,所以大使命就是大君王的命令.在這命令的後面,耶穌給了一個非常寶貴的應許,「我就常與你們同在,直到世界的末了.」教會和其他團體的不同,就是教會有主的同在,所以教會不可能被消滅,被打倒.
\newline
\newline
教會是以傳福音開始的,聖靈降臨建立了教會後,立時就有傳福音的工作,有三千人信主,加入教會；而到使徒行傳最後一章,仍是傳福音,雖然保羅坐監牢,但沒有停止傳福音的工作.所以教會在世上,無論得時不得時,都是傳福音.教會在耶路撒冷剛開始的時候,困難比較少,所以他們傳福音較易,但後來逼迫困難來了,在使徒行傳的末了,困難很大；最重要的使徒,手上帶著鎖鍊,身體失去自由,但沒有停止傳福音.所以今天的教會,一定要傳福音,若不傳福音,便是放棄了主對教會所交託的使命.
\newline
\newline
今天西方有些神學家有這樣的謬論,說傳福音的時代已經過去,我們不須再傳福音了,因為各國都有自己的宗教,而且各宗教都是一樣；如果硬把福音傳給他們,就是干涉他們的自由.所以他們認為宣教士應該回家,宣教士的時代已經過去,這種說法是荒謬的,他們不是神的兒女,是假先知.有一個憑信心設立,專門向非洲人工作的差會,今天便缺少二百五十位宣教士.包括醫生,工程師,聖經教師等等；他們找不到這許多人,他們盼望各國的教會能差派這樣的人才,肯進到非洲國家去.耶穌對門徒說,要到普天下去傳福音,直等到他來.今天耶穌還未來,所以我們仍要傳福音,而且要加緊傳福音,因為耶穌快來了.
\newline
\newline
另外我們要思想的,就是中國的教會對傳福音負了甚麼責任?我們和別的國家一樣,不過是基督身體內的一個肢體,我們要和其他的合在一起,負起傳福音的責任.有人認為傳福音的工作就好像接力賽跑一樣,猶太人跑了第一棒,跑完之後,便把棒交給歐洲教會,歐洲人跑一段；又把差傳的棒交給美洲的教會,到現在,輪到中國教會接棒了.我並不贊成這樣的講法,因為是不合聖經的,聖經說教會要傳福音,教會存在多久,傳福音就要多久,差傳工作是每個教會都要作的,棒子不能放下,更不能交給別人來代替.我可以舉出事實來證明.很多人以為猶太人沒有傳福音了,但是今天的宣教團契中,仍有猶太人的福音會存在,而且在歐美等地發展.而在歐洲國家,如英國,德國,猶太人還有數千人在宣教的工場上繼續工作.至於美國人的猶太宣教士,則佔數量最多.今天全世界的宣教士,大約共有五萬多人,其中美國人佔三萬多.如果這宣教的棒,全落在中國人身上,中國人能接得起嗎?中國能派出五萬多名跨越文化的宣教士,去取代現有的宣教士嗎?感謝主,近二十年來,中國人在這方面已有醒悟,今天參與差傳工作的華人教會,大概已有一,二百間.
\newline
\newline
可是,如果我們用聖經來對照一下,便會發現中國的差傳,有一些缺點,是應該改正的.
\newline
\newline
第一,今天的差傳,多半是差錢,而不是差人.差傳應該是差人,聖經第一次的差傳,是安提阿教會把保羅和巴拿巴二人差出去,後來他們把福音從亞洲帶到歐洲,果效非常好.可是那裡並沒有提到錢,但今天中國人的差傳卻是注重錢,注重差傳基金,卻沒有多人作宣教士,只是把那些錢用來支持已有工作的傳道人,這樣的差傳是有缺點的.內地會的創辦人戴德生牧師,他作工有一個原則,值得我們學習.他說:「神的工作,如果合神的旨意,用神的方法,神的供應是不成問題的.」
\newline
\newline
第二,聖經所講的差傳,是先有教導,後有行動的.安提阿是建立在外邦的教會；建立教會以後接下去即有教導,由保羅和其他一些人作教導的工作,共有一年多的時間.後來他們靈命長進了,生活也有良好表現.(徒十一26)說:「門徒稱為基督徒,是從安提阿起首.」這是安提阿教會最美的見證,他們的生活和不信的人完全不同,他們聖潔,謙卑,彼此相愛,使外邦人非常希奇,覺得他們就像基督,所以稱他們為基督徒.可見他們先有信仰的根基,生命長進,生活有表現,然後才有差傳的行動的.
\newline
\newline
中國教會為甚麼沒有宣教士呢?因為我們對信徒沒有教導,特別沒有關於為主受苦的教導.耶穌說:「若有人要跟從我,就當捨已,天天背起他的十字架來跟從我.」作為基督的門徒,是要受苦的.我們不是自己找苦來吃,乃是要預備受苦的心,這就是我們的兵器.有的教會很有錢,但卻從來無人願意獻身傳道,他們出錢很容易,而且出的很多；但是傳道就不同了,傳道是要受苦的.我們不敢作宣教士,是因為怕受苦,我們太缺少這方面的教導.
\newline
\newline
中國教會今天差傳的工場有多大呢?到目前為止,我們已經做了二十年的差傳,但大多數都是差錢,也有的是差人的.可惜是借別教會的人,大教會出錢,小教會出人；是做那些富足教會的宣教士,而且他們多數在中國人中間工作.我們支持一個人在泰國,馬來西亞,沙巴等地工作,但他向誰工作呢?向中國人,廣東人,客家人,福建人工作,所以只是中國人向中國人工作.但(徒十一19-21)告訴我們:「那些因司提反的事遭患難四散的門徒.直走到腓尼基和居比路並安提阿,他們不向別人講道,只向猶太人講.但內中有居比路和古利奈人,他們到了安提阿,也向希利尼人傳講主耶穌.主與他們同在,信而歸主的人就很多了.」起初教會的差傳,就是那些因逼迫四散的猶太人,他們只在各處找猶太人傳福音,但後來有人也向希利尼人傳福音,聖經特別說:「主與他們同在,信而歸主的人就很多了.」當他們傳福音,跨越自己的民族,文化的時候,神便格外的賜福他們.因為耶穌的使命,是要他們向萬民傳福音,包括各國各族人.中國人作差傳,只注意自己同族的人,不注意外族的人,這是我們的弱點.如果我們只向本國的人傳福音,便等於把耶穌的命令打了折扣.
\newline
\newline
有些外國專家說,中國人若要傳福音,不必到別的地方去,就留在亞洲好了,因為亞洲信徒佔全人口的比例最低.這是確實的,一般亞洲的國家,在全國總人口內,基督徒所佔的比例少而又少,有的不到百分之一,有的是百分之二或三；最高是韓國,據說已經達到百分之二十,但是再沒有更高的.所以西方專家便勸中國人只留在亞洲,不要到別處去.這樣的說法對嗎?今天世界很明顯的分成兩個壁壘,就是信神的和不信神的；照經濟和文化方面的分法,則非洲和亞洲同屬第三世界.到底耶穌說向萬民傳福音,是否也包括非洲在內呢?全非洲現在有四萬萬的人,他們當中仍有許多人沒有聽過福音.初期教會,福音曾經達到非洲的北部,使徒行傳中有亞力山大城的名字.今日仍叫亞力山大,是埃及一個很出名的城.但回教興起以後,便把這些教會通通摧毀,非洲人陷在黑暗之中,有一千二百年之久,埃及和埃塞俄比亞雖然有教會,但卻是死的教會,因為教會不傳福音,所以沒有發展.
\newline
\newline
直到三百年前,才有李文斯敦到非洲去,藉著醫療服務,用他的一生在非洲傳道,最後死在那裡.當時非洲被稱為黑暗大陸.後來到一八九三年,北美洲有三個青年,他們要到非洲內地傳道.那時的西非洲,還沒有獨立的國家.當他們初到海岸的時候,遇到一些年長的白人宣教士,他們告訴這三個青年人不能進去；因為當時的非洲內地,根本沒有路可通,各樣的疾病,野獸,非常可怕.但這三個青年人仍遵照主的命令,繼續前行.不到一年,三人之中已經死了兩個,剩下的一個又得了重傷,只好返回北美洲.他回國後,報告非洲的需要,有些青年人響應,於是他們第二次再起程,可是卻無法進去,各人都得了重病.後來第三批再去,經過很久很久,才在內地建立一個小小的茅舍,這就是外國宣教士在非洲內地傳道的開始.他們在那裡辛苦的工作.起初整整十七年的時間,但只得十三個黑人信主,而十七年中,宣教士死在非洲的,比黑人信主的數目還要多.他們付了這麼大的代價.到如今這三個青年所建立的工作,已經發展成一個很大的差會,他們憑信心工作,有一千多個宣教士,從世界各國加入這個差會,分佈在東非和西非工作,所建立的教會,共有四千間之多.
\newline
\newline
在東北非洲地方有一個國家即叫依索比亞,是蘇丹內地差會的重要福音工場,曾經有三百多位宣教士在那裡工作,建立了兩千多間教會.到了一九七四年,這國發生政變,當地的政黨掌權.當時蘇丹內地差會曾來信給兩個很有名的西差會,向他們請教,可是沒有回答.就在那時,神便帶領我到非洲工作,先到非洲工作,先到依索比亞,共八天的時間幫助信徒機會去工作.我發現非人十分歡迎且渴望有中國人去幫助他們,傳福音給他們聽.
\newline
\newline
當然,我並非說每一個人都要到非洲去,如果神差遣你,那麼我們都應該同心說阿們.另外有些人,雖然自己不能去,但他們知道這是神的工作,所以願意盡上本份,用金錢去支持差傳的工作,用禱告去支持差傳的工作.一個宣教士,背後需要很多人的禱告,一個差會也需要很多人的經濟支持.所以,傳福音給萬民聽是教會當行的事,是合乎聖經教訓,我們應該切實照著去行,不要把神的命令任意打折扣,中國教會應當快快差派更多宣教士,去到給萬民傳福音,遠如非洲,南美洲的工場也是我們的傳福音工場,我們應當關心,應當有分.這樣才能夠見證我們教會傳福音的工場是合而為一的工場.
\newline
\newline


\newpage

\section{第53屆~港九培靈會~戴紹曾~第一講}
\label{sec:591}
\textbf{聖靈與差傳事工}
\newline
\newline
連結: \href{https://www.hkbibleconference.org/session-message/view/591}{\texttt{https://www.hkbibleconference.org/session-message/view/591}}
\newline
\newline
\hyperref[sec:642]{< < < PREV SERMON < < <}
~
\hyperref[sec:toc]{[返目錄]}
~
\hyperref[sec:593]{> > > NEXT SERMON > > >}
\newline
\newline
使徒行傳 1:1-1:11
\newline
\begin{longtable}{cl}
\hline
\hline
章節 & 經文 (和合本修訂版)\\
\hline
1:1 & \begin{tabularx}{0.7\textwidth}{X} 提阿非羅啊，我在第一本書中已論到耶穌從開頭所做和所教導的一切事， \end{tabularx} \\ \\ \relax
1:2 & \begin{tabularx}{0.7\textwidth}{X} 直到他藉著聖靈吩咐所揀選的使徒後，被接上升的日子為止。 \end{tabularx} \\ \\ \relax
1:3 & \begin{tabularx}{0.7\textwidth}{X} 他受害以後，用許多確據向使徒顯明自己是活著的，在四十天之中向他們顯現，並講說神國的事。 \end{tabularx} \\ \\ \relax
1:4 & \begin{tabularx}{0.7\textwidth}{X} 耶穌和他們聚集的時候，囑咐他們說：「不要離開耶路撒冷，但要等候父的應許，就是你們聽見我說過的。 \end{tabularx} \\ \\ \relax
1:5 & \begin{tabularx}{0.7\textwidth}{X} 約翰是用水施洗，但過了不多幾天，你們要在聖靈裡受洗。」 \end{tabularx} \\ \\ \relax
1:6 & \begin{tabularx}{0.7\textwidth}{X} 他們聚集的時候，問耶穌：「主啊，你就要在這時候復興以色列國嗎？」 \end{tabularx} \\ \\ \relax
1:7 & \begin{tabularx}{0.7\textwidth}{X} 耶穌對他們說：「父憑著自己的權柄所定的時候和日期，不是你們可以知道的。 \end{tabularx} \\ \\ \relax
1:8 & \begin{tabularx}{0.7\textwidth}{X} 但聖靈降臨在你們身上，你們就必得著能力，並要在耶路撒冷、猶太全地和撒瑪利亞，直到地極，作我的見證。」 \end{tabularx} \\ \\ \relax
1:9 & \begin{tabularx}{0.7\textwidth}{X} 說了這些話，他們正看的時候，他被接上升，有一朵雲彩從他們眼前把他接去。 \end{tabularx} \\ \\ \relax
1:10 & \begin{tabularx}{0.7\textwidth}{X} 他升上去的時候，他們定睛望天，看哪，有兩個人身穿白衣站在他們旁邊， \end{tabularx} \\ \\ \relax
1:11 & \begin{tabularx}{0.7\textwidth}{X} 說：「加利利人哪，你們為甚麼站著望天呢？這離開你們被接升天的耶穌，你們見他怎樣升上天去，他也要怎樣來臨。」 \end{tabularx} \\ \\
[1ex]
\hline
\hline
\end{longtable}
這幾個晚上,我們要一同思想的題目是初期教會的型態,和四年前的題目一樣；但以前所講的,我們特別注意耶穌在初期教會的眼光,當時教會如何看耶穌,如何傳耶穌,耶穌在教會的地位如何.我們也講到聖靈的工作和聖靈的位格；並列舉些被聖靈充滿的實例:如司提反的見證,巴拿巴的見證,尤其注意他被聖靈充滿時日常生活的表現.凡此,在使徒行傳中有相當具體的記載.我們也一同思想如何尋求神的旨意.這是許多青年求知的問題.我們也看見初期教會增長的情形.最後一晚,我們也看見初期教會為主受苦.
\newline
\newline
這十個晚上,我們又回到使徒行傳,再看初期教會的型態,這次我們著重於差傳事工,回顧初期教會如何關心福音工作.
\newline
\newline
差傳與神的國有密切關係,我們常禱告說:「願你的國降臨.」行傳始末都論到神的國.路加告訴我們,耶穌受害之後,以許多憑據將自己活活地顯給使徒看,四十天之久向他們顯現,講說神國的事.(一3)到了廿八章,那時福音已傳開,保羅在羅馬沒有自由,視為囚犯,他在自己租的房子裡住足了兩年,凡來見他的人,他全都接待,放膽傳講神國的道,將主耶穌基督的事教導人,並沒有禁止.(30-31)行傳開始,主耶穌向門徒傳講神國的事,過了卅年,保羅在羅馬生活在不自由情形之下,照樣傳講神國的事.
\newline
\newline
各位青年弟兄姊妹!我們的心地有時很狹窄.主耶穌對門徒講神國的事,但他們心裡另有所想,「問耶穌說,主阿!你復興以色列國,就在這時候嗎?」(6)門徒想到自己的國家,是神應許給亞伯拉罕後裔之地,但是他們的國家卻被羅馬國佔領了.其實,主所指非地上的國而是天國,是關乎萬民的福音,這福音將影響世上列國.我們的心是否像耶穌門徒那樣狹窄,以自己,自己的家,自己的同胞為中心?復活的主所關心的是天下四海的人,祂說:「這福音要傳到地極.」
\newline
\newline
求主幫助我們,以祂的心為心,以祂的事為念,開闊我們的心胸,擴充我們的異象；不只在近處,還有遠方的需要.
\newline
\newline
使徒行傳共廿八章,也可稱為聖靈行傳.全卷至少有55至60次題到聖靈:聖靈在舊約,在新約,在教會,在個人身上,在各方面工作.路加告訴我們:耶穌吩咐所揀選的使徒是藉著聖靈.(一2)
\newline
\newline
有位聖經學者布路斯說:「『藉著聖靈』是使徒行傳在神學首要的論調.」耶穌不只藉著聖靈吩咐門徒,且福音也是藉著聖靈傳開的,門徒也是藉著聖靈把福音傳到地極.
\newline
\newline
思想差傳事工,當然也以聖靈的工作為首,我們要特別注意聖靈與差傳的工作,應注意者有四方面:－
\newline
\newline
(一)聖靈潔淨人心的工作　主耶穌吩咐門徒不可離開耶路撒冷,要等候父所應許的.(一4)為甚麼接著說明,約翰是用水施洗,但不多日,你們要受聖靈的洗.(5)潔淨工作,是差傳事工的預備工作,如無潔淨工作,就沒有事奉,也沒有差傳的事工了.保羅在耶路撒冷向教會領袖解釋為何進哥尼流的家.在當時這是件難事,因猶太人不肯和外邦人來往.福音要進到外邦人中,須預先解決猶太人心中的優越感,劈開他們心中的成見,和老傳統觀念.有人反對彼得到外邦人的家；他解釋當他進哥尼流的家.聖靈降臨在哥尼流家人身上.那時聖靈的降臨,和五旬節聖靈降臨一模一樣；因信而心得潔淨,不分你們他們.因此,我們看見聖靈潔淨的工,不只在哥尼流的家,因信潔淨了他們的心；五旬節門徒在耶路撒冷,也是因信得潔淨.是的,器皿需要聖靈預先潔淨,才能合乎主用；否則談不上為主作工.
\newline
\newline
以賽亞書第六章記載青年以賽亞蒙召的事實.當治國多年的烏西雅王駕崩之年,以賽亞帶著沉重的心情進入聖殿.他恍然大悟.雖國之寶座空置,沒有君王；但是他看見神坐在祂自己的寶座上,他看見權能的主,更看見聖潔的主；於是,這位驕傲的青年低下頭來,本以為自己樣樣都行,今覺自己不配,大聲呼喊:「禍哉!我滅亡了.」他看見聖潔的神,也看見自己；強烈對比之下,給他對神對自己有種新的認識；以往覺得自己不錯,比別人強,多有事奉,做神的出口；但看見異象以後,完全不同了,說；「禍哉!我滅亡了；因為我是嘴唇不潔的人,又住在嘴唇不潔的民中.」各位青年!我相信這話也是我們常想的,此念頭也是常壓在我們心上的.我們處此淫亂的社會,怎麼辦呢?尤其要蒙主使用怎可能呢?以賽亞聽見神說:「我可以差遣誰呢?誰肯為我去呢?」他說:「我在這裡,請差遣我.」他看見聖潔的神,認識自己的需要,謙卑承認自己的污穢,然後完全在神恩典之下,蒙神潔淨,並蒙神使用.「我們若認自己的罪,神是信實的,是公義的,必要赦免我們的罪,洗淨我們一切的不義.」(約壹一9)感謝主!
\newline
\newline
(二)聖靈賜能力的工作　當然潔淨也是能力的一種表現,因為永遠無法靠自己潔淨自己,聖靈要幫助我們過得勝的生活；在我們尚未談及工作能力之先,應注意生活的能力.雖然我們今天談的是差傳事工,我們必須先想到生活的問題；因為我是何等樣的人,會影響我所做的；實際上,我是何等樣的人,比我所作的更為重要.多少時候,我們想為主做這做那,可是祂要先在我們身上作工.一兩年前,我曾聽過使徒德牧師演講基督徒蒙召的問題,開始講時,一般人的觀念,蒙召一定是該作這個該作那個,他並非如此.他說:「神召你們,叫你們聖潔.」今天人在世上,價值觀念顛倒,打開報紙,每天都有訃文,寫的大都是專家,也有蝴蝶專家；其實更為重要的,是這人是何等樣的人,他內裡的美德,他的身份及其生活表現.各位青年!當我們想到差傳事工,讓我們先想到生活上的問題,聖靈的能力要幫助我們在日常生活中過得勝生活.保羅給帖撒羅尼迦教會寫信說:「願賜平安的神親自使你們全然成聖,又願你們的靈,與魂,與身子,得蒙保守,在我主耶穌基督降臨的時候,完全無可指摘；那召你們的,本是信實的,祂必成就這事.」(帖前五23-24)保羅指出信徒的生活,要聖潔,過得勝的生活；靠自己是辦不到的,惟召我們的主是信實的,祂必成就這事.巴不得今晚主抓住我們的心,看見衪是信實的,衪必成就這事．神要賜我們能力,我們要順服祂,要愛我們的弟兄；這是很重要的；無論大小事上要順服神的旨意；因為神把聖靈賜給順服祂的人.我們若渴慕聖靈,卻在日常生活中不肯順服,就永遠得不到聖靈的能力.求主幫助我們,無論在教會,家庭,學校,在工作上,在社會上,有一致的標準,在日常生活中靠聖靈得勝,這種能力只彰顯在我們生活中,這是基礎,在工作上更顯得清楚；所以主耶穌告訴門徒「聖靈降臨在你們身上,你們就必得著能力.」意思是說,福音傳到地極,乃靠聖靈的能力.聖靈的能力運用在我們身上,若宣揚自己,高抬自己,彰顯自己的恩賜,那就錯了；我們應像施洗約翰說:「祂必興旺,我必衰微.」不是高抬自己,乃要高舉耶穌.聖靈要給我們能力作見證.(二11)說他們用鄉談講說神的大作為.他們開始傳福音,福音就傳開了.到了21節,福音傳開的範圍更大「到那時候,凡求告主名的,就必得救.」在耶路撒冷,猶太全地,和撒瑪利亞,直到地極,福音傳開了.一種抵擋不住的能力,六章說,彼得站在那些官長面前,他們說所發生的事情我們不能說沒有.司提反是以聖靈和智慧說話,這種能力使眾人抵擋不住.巴不得今天我們傳福音時有一樣的能力,使聽眾無不受感動.如果我們細看聖靈每降臨在一些聖徒身上,都有同一的表現；不是醫病,也非其他表現；乃是歌頌神的聖名把福音傳開.但願主幫助我們,在聖靈能力之下,更有效地把福音傳開.聖靈的能力給我們膽量,勇氣.彼得站在那些官長面前,他被聖靈充滿.(四8)他們看見聖靈能力在他身上的表現,乃是膽量的表現.彼得不久之前曾三次否認耶穌,現在完全不同了,有聖靈的能力在他身上,他說:「你們這些官長文士禁止我們奉耶穌的名講道,傳耶穌復活的訊息；可是我告訴你們,「天下人間沒有賜下別的名,我們可以靠著得救.」我們都需要聖靈的能力作為見證主的力量.
\newline
\newline
求主除去我們所有的懼怕,賜給我們聖靈的膽量聖靈的能力.
\newline
\newline
(三)聖靈的引導　聖靈在使徒行傳許多方面,許多不同地方作引導之工,舊約聖靈的啟示也是一種引導.在行傳中,我們看見聖靈差傳的工作,不斷帶領指引使徒的路.腓利傳福音,已經到了撒瑪利亞,聖靈卻吩咐他到曠野去；聖靈感動他離開耶路撒冷,離開撒瑪利亞,到偏僻地方.因著腓利的順服,一個非洲人埃堤阿伯的太監聽了福音,得蒙拯救,這豈不是差傳工作嗎?各位!神帶領我們不一定到大的地方,過舒服的生活；祂可能要我們進到曠野,到被人忽略之地傳福音,隱藏起來,帶著聖靈的能力傳福音.今天在東南亞,許多國家很難進去傳福音,中東一帶很多回教徒,沒有人去傳福音.聖靈是否對你說:「你要去!」香港一些青年,使我欽佩；最近有位伍弟兄在新加坡,準備到泰國中部做超越文化的工作,泰國中部有種病使人死亡；這位弟兄雖然身體不很健康,但他願意到那被人遺忘,人所不肯去的佛教國家；去向泰國人傳福音,帶領他們認識耶穌.
\newline
\newline
當彼得禱告時,聖靈對他說:樓下有三人等你一同到哥尼流的家；聖靈帶領彼得,是彼得空前的舉動,聖靈要他作本來不可能作的事；後來彼得站在耶路撒冷教會時,他說,他去時是聖靈吩咐他去的.(十一12)
\newline
\newline
當安提阿教會的弟兄姊妹同心合意傳福音,禁食禱告時,聖靈說要分派巴拿巴和掃羅去作工.聖靈又向教會顯現,帶領教會；我們看見當時教會的順服,揀選最好最會講解聖經的人才,有佈道恩賜的人才出去作差傳工作.
\newline
\newline
聖靈不只作差派的工作,有時也作攔阻的工作.當保羅第二次出外時,他想到以弗所,但聖靈不許,耶穌的靈攔阻保羅.
\newline
\newline
感謝主!初期教會對聖靈的帶領有敏銳的心.相信初期教會不隨便說是聖靈的感動.我們應當小心!在教會中常有輕易地說:「聖靈的感動」其實卻是自己的意思；因為自己的意思,恐怕別人反對；聖靈的感動,別人就不敢反對；若反對,便是反對神.這樣的思想真可怕!
\newline
\newline
各位青年!我們需要聖靈的帶領,你前面的道路,你的學問,你的婚姻,你一生的工作,你將來的事奉,你自己打算嗎?你是否放在禱告中,讓聖靈帶領?各位青年!我們要等候聖靈的啟示,明白神的旨意.我們若在凡事上認定祂,祂必指引你的路.
\newline
\newline
(四)聖靈的恩典　聖靈不但作潔淨工作,聖靈也賜能力,聖靈帶領我們在福音事工上,聖靈也賜恩典.茲簡舉三例:
\newline
\newline
司提反在極惡劣環境之下,敵人切齒痛恨他,要用石頭打死他,他抬頭看見天開,內心何等平靜.路加告訴我們,司提反被聖靈充滿,在敵人面前危險之中,他平靜的表現,實在是聖靈所賜的恩典.
\newline
\newline
舊約的以利沙,四圍受敵；他的僕人驚惶失措,報告主人；以利沙禱告,求神開他僕人的眼睛,讓他看見幫助我們的,比敵擋我們的還多.在最危險時,聖靈要幫助我們看見神的同在,耶穌站在神的右邊.
\newline
\newline
耶路撒冷教會經過一番的逼迫；聖經說,他們得到聖靈的安慰,人數也增多了.聖靈的安慰,聖靈的喜樂,都在使徒行傳中,在傳福音的過程,保羅和巴拿巴到一個地方傳福音,被同胞拒絕,後被趕走.路加說:保羅和巴拿巴兩個神的僕人都被聖靈充滿了,滿有喜樂.
\newline
\newline
感謝主!平靜,安穩,喜樂都是聖靈所賜的恩典.聖靈在差傳的工作上,要帶領我們,不只在初期教會,今天更是如此.
\newline
\newline
願我們今晚省察自己,聖靈是否在心裡作工?
\newline
\newline


\newpage

\section{第53屆~港九培靈會~戴紹曾~第二講}
\label{sec:593}
\textbf{禱告與差傳}
\newline
\newline
連結: \href{https://www.hkbibleconference.org/session-message/view/593}{\texttt{https://www.hkbibleconference.org/session-message/view/593}}
\newline
\newline
\hyperref[sec:591]{< < < PREV SERMON < < <}
~
\hyperref[sec:toc]{[返目錄]}
~
\hyperref[sec:595]{> > > NEXT SERMON > > >}
\newline
\newline
使徒行傳 4:23-4:31
\newline
\begin{longtable}{cl}
\hline
\hline
章節 & 經文 (和合本修訂版)\\
\hline
4:23 & \begin{tabularx}{0.7\textwidth}{X} 二人既被釋放，就到自己的人那裡去，把祭司長和長老所說的話都告訴他們。 \end{tabularx} \\ \\ \relax
4:24 & \begin{tabularx}{0.7\textwidth}{X} 他們聽見了，就同心合意地高聲向神說：「主宰啊！你是那創造天、地、海和其中萬物的； \end{tabularx} \\ \\ \relax
4:25 & \begin{tabularx}{0.7\textwidth}{X} 你曾藉著聖靈託你僕人—我們祖宗大衛的口說：『外邦為甚麼擾動？萬民為甚麼謀算虛妄的事？ \end{tabularx} \\ \\ \relax
4:26 & \begin{tabularx}{0.7\textwidth}{X} 地上的君王都站穩，臣宰也聚集一處，要對抗主，對抗主的受膏者。』 \end{tabularx} \\ \\ \relax
4:27 & \begin{tabularx}{0.7\textwidth}{X} 希律和本丟‧彼拉多，同外邦人和以色列民，果然在這城裡聚集，要攻打你所膏的聖僕耶穌， \end{tabularx} \\ \\ \relax
4:28 & \begin{tabularx}{0.7\textwidth}{X} 做了你手和你旨意所預定必成就的事。 \end{tabularx} \\ \\ \relax
4:29 & \begin{tabularx}{0.7\textwidth}{X} 主啊，現在求你鑒察，他們的威嚇，使你僕人放膽講你的道， \end{tabularx} \\ \\ \relax
4:30 & \begin{tabularx}{0.7\textwidth}{X} 伸出你的手來，讓醫治、神蹟、奇事藉著你聖僕耶穌的名行出來。」 \end{tabularx} \\ \\ \relax
4:31 & \begin{tabularx}{0.7\textwidth}{X} 他們禱告完了，聚會的地方震動；他們都被聖靈充滿，放膽傳講神的道。 \end{tabularx} \\ \\
[1ex]
\hline
\hline
\end{longtable}
全卷使徒行傳強調聖靈的工作,教會藉聖靈的能力把福音傳開.
\newline
\newline
禱告與差傳是分不開的.耶穌說:「要收的莊稼多,作工的人少；所以你們當求莊稼的主,打發工人出去,收他的莊稼.」一個不禱告的教會,不可能是個差傳的教會,一個懂得禱告的教會,才知道怎樣作差傳的教會；不但可以作,會作,也應該作.
\newline
\newline
安提阿教會設立不久,弟兄姊妹同心合意一齊事奉主,禁食禱告；聽見神的聲音,莊稼的主對教會說,要分派保羅和巴拿巴出去,教會要他們作差傳的工.初期教會差傳工作就由此而開始.
\newline
\newline
福音傳到中國時,特別在近代教會史；百五年前,無論馬禮遜,戴德生,及其他早期的宣教士,福音運動是產生於禱告；當時西方教會開始關心其他地方的許多人.以內地會來講,戴德生和他的同人,恆切禱告,懇切祈求主,要打發24位宣教士到中國內地,每省兩位,邊疆地方也有兩位；雖然人數很少,但內地會工作就由禱告開始了.
\newline
\newline
今天華人教會很關心差傳工作,香港,菲律賓,台灣,新加坡,東南亞一帶甚至北美洲.
\newline
\newline
當我們談及差傳,常把注意力放在錢財上,教會經濟條件許可,就開始差傳工作,沒有著重於禱告；其實,教會有禱告,錢自然不成問題；事實上,工人比錢更為重要.所以主耶穌說,要求莊稼的主打發工人.求主幫助我們,叫我們的注意力放在主所注意的；主告訴我們「要禱告」!
\newline
\newline
初期教會的榜樣,實在帶給我們很大的幫助,從初期教會看見禱告的重要性.
\newline
\newline
使徒行傳全卷,作者路加至少有卅多次題及教會禱告,不斷題到教會禱告生活,由此可見初期教會重視禱告.使徒時代的教會是個差傳的教會；我們看見禱告與差傳,差傳與禱告.
\newline
\newline
禱告的重要:當時,初期教會面臨分裂,有些寡婦發怨言說:「我們被忽略了,我們這些外地來的猶太僑民,初回到耶路撒冷時,教會關心那些會講本地話的寡婦,講希利尼語的寡婦卻被忽略了.」使徒們發現教會潛伏分裂的危機,所以處事審慎,招聚全會眾,說:「我們撇下神的道,去管理膳食是不可以的.」這句話是今日許多傳道人所不敢講的；他們忙於東奔西跑,做些不關重要的事；但初期教會要專心以祈禱傳神的道為事.當時使徒不肯撇下禱告,不肯撇下神的道,絕非他們驕傲,沒有憐憫的心,不關心這些被忽略的寡婦,他們感到自己的工作須優先考慮,以專心祈禱傳道為事,其餘的事,可交給忠心事奉的人去作.
\newline
\newline
求主幫助我們看見當時使徒如何重視禱告.
\newline
\newline
當主耶穌升天以後,主和門徒回耶路撒冷,在房裡同心合意恆切禱告.五旬節,彼得講道；許多人覺得扎心,彼得要他們向神悔改認罪,至少有三千人歸主.這些歸主的人,「都恆心遵守使徒的教訓,彼此交接,擘餅,祈禱.」(二42)不只使徒重視禱告,連剛歸主的三千人也重視禱告.
\newline
\newline
一年來,我在外看見一些新的信息；其實,我相信在座的弟兄姊妹早就看見了.當我思想使徒時代教會的禱告,我讀使徒行傳,發現基督徒有許多名稱,有聖徒,信徒.......初期教會的信徒,也被稱為「求告主名的人.」由此可知他們多麼重視禱告.安提阿地方,一般人稱這些聖徒為基督徒；因他們終日講基督,過基督的生活,四圍的人稱他們基督徒.為何還有人稱他們為求告主名的人呢?豈不是他們常常禱告嗎?不知今天香港一般人如何看我們呢?巴不得戴紹曾和在座各位兄姊都有同樣的稱呼,我們同是一群求告主名的人.
\newline
\newline
初期教會有很多重要的功課,我們需要學習:
\newline
\newline
禱告不受時間限制.他們在任何時間禱告.彼得和約翰,在下午三點鐘上聖殿禱告.哥尼流也是同樣時間禱告；也許這是猶太人的規矩.有一次彼得在約帕地方,肚子餓等吃午餐時,趁機禱告,禱告之中他得到啟示；藉著禱告,福音傳開了.我們又看見教會在夜間禱告.保羅和西拉也在半夜禱告.詩篇112篇作者說:「祂必不叫你的腳搖動,保護你的必不打盹；保護以色列的,也不打盹,也不睡覺.」早,午,晚,隨時隨地可以禱告；神聽禱告,絕不受時間限制.至於禱告的地點,初期教會在很多不同的地方禱告,在耶路撒冷,在撒瑪利亞,大馬色,該撒利亞,約帕,安提阿,腓立比,各地都是他們禱告之處；有時也在聖殿禱告,但不只限於聖殿.司提反在路上禱告,彼得在屋頂禱告,他們也在河邊,保羅在沙灘,在海島禱告,地方不一；這是當時許多猶太人不明的真理.我們再看司提反在第七章所講的信息,其重要目的,是要他的同胞對真神有正確認識.有人反對司提反說:「你反對我們的聖殿.」司提反開始講的一篇信息釋放了教會,使教會能夠離開耶路撒冷,將福音傳到地極.當時猶太人的宗教以聖殿為中心,認為有聖殿才能敬拜神,才能獻祭禱告；然而司提反對他們說:「你們錯了,神是無所不在的神.」第七章開頭說:榮耀的神,在吾珥地方向亞伯拉罕顯現.不是耶路撒冷,不是在聖殿,乃在吾珥.過段時間,亞伯拉罕到了哈蘭,神再次向他顯現,叫他到迦南,在迦南地,神又向他和他的後裔顯現.在埃及,神向約瑟,向摩西顯現在曠野.(英國人腓立斯寫了本書「你的神太微小了」此書頗有意思)司提反所講的信息,要聽眾領悟他們所知道的神太微小了.我們若明白禱告的真理,就不至有此錯誤觀念.禱告不限定在禮拜堂.家裡,學校,工作場所,渡船中,無論何時何地都可以禱告；因我們的神是無所不在的.
\newline
\newline
禱告也不受人數限制,或多或少均可.彼得個人禱告,神向他顯現；哥尼流個人禱告,神垂聽,救恩臨到他的家；彼得約翰二人禱告,非常投契；我們看四福音,這二人平時保持競爭的關係,當耶穌復活那天,他倆向著墳墓賽跑,彼得年紀大,跑步慢；但他作事積極,起步卻很快；約翰雖年青卻膽怯,起步蹣跚,但是跑的很快,早抵墓前；不過他膽小不敢擅入墓穴,所以站住,結果還是彼得先進去的.到了使徒行傳,彼得和約翰的關係,不再是賽跑競爭式的了,因有神的愛在他們心裡,聖靈潔淨了他們的心,心地不再狹窄,兩人同心合意上聖殿禱告.(第三章)感謝主!主自己告訴我們；若有兩三個人奉祂的名禱告,祂在我們中間,無論求甚麼,祂都給我們.感謝主!禱告不在於人數多少,一兩個人的禱告神垂聽,福音也被傳開了.
\newline
\newline
有一次我參加九龍一個禱告會,海外基督使團同工費述凱牧師講道,他講個有關禱告的例子；他剛得救那時,無論大小事他都禱告,出門之前,先求問主要他行左或行右,以便向人傳福音.別以為小事主不理會,不,大小事情都可以帶到主面前.
\newline
\newline
在行傳中,一般的事情都有,我們可見到一些代禱的項目.彼得和約翰到了撒瑪利亞,信徒還未受聖靈,需要見証,讓猶太人和全世界的人知道救恩沒有分別；所以彼得禱告,求神賜下聖靈,撒瑪利亞人得到了聖靈.第九章,彼得為病人祈禱.保羅到羅馬去時,船遇破壞,有段時間在海島,他也為病人禱告.再看教會,要分派一些人的工作,為那七人禱告.(六章)他們的選擇出於禱告,這是很重要的.教會選舉,若不出於禱告,常會出問題.安提阿教會要分派人出去,他們的選擇,也產生於禱告.許多的失敗,乃因沒有禱告,後果實不堪設想!保羅設立的教會選立長老,將這些長老,在禱告中交托主,他自己就離開.這是西教士們需要學習的功課:撒手放下,有信心,讓當地的信徒負責自己的工作.保羅這樣做並非靠自己的智慧和自己的斷定；乃是出於禱告,把教會長老交托所信的主.
\newline
\newline
教會的祈求,分很多方面,有時需要神的帶領往何處去.第四章,他們面對眾多仇敵；但他們並非禱告求主使他們脫離仇敵,而是求主賜給膽量,放膽向仇敵傳福音.
\newline
\newline
禱告的內容廣泛,請問:各位弟兄姊妹的禱告,到底多廣?範圍多大?我發現很多人的禱告,首先為自己的需要,路途通順時不禱告,有困難時才禱告.初期教會的信徒,被稱為求告主名的人；無論得意與否,工作得力與否,他們都禱告,且禱告的內容很廣.
\newline
\newline
求主教導我們如何禱告!
\newline
\newline
請看行傳中三個實例:
\newline
\newline
第一是彼得,約翰和他們同人的禱告,第四章是行傳記載中最長的禱告；其實,我們覺得並不長；可是神聽禱告不在乎長短.我相信行傳中他們恆切的禱告必定很長的,不過禱詞無記載.他們禱告開頭就說:「主啊!」英文翻本為「掌全權的主」這樣禱告的開始是我最喜歡聽的.他們面對希律王和許多官長和凱撒,都是大有權威者；但是他們禱告的對象是坐在寶座上的主,有權柄有主權的主,他們不注目希律和凱撒,他們只看見創造天地偉大的主,全世界全宇宙是祂創造的.他們講及主永恆的計劃,早在大衛時主曾說過,會有人起來反對受膏者－我們的救贖主,現在事情竟然發生了,他們受了有地位的人恐嚇,不可奉耶穌的名傳福音.面臨這樣的危險,他們禱告求主賜予膽量並求神蹟奇事隨著他們的工作；因為猶太人被說服需要看見神蹟奇事.感謝主!巴不得今天我們也有這樣的禱告,華人教會有這樣的禱告做差傳的工作；震動香港,新加坡,台灣和東南亞.現在已經有很多弟兄姊妹在中國禱告,已有神蹟奇事隨著他們.禱告得到膽量,放膽傳福音,得到聖靈的充滿；我們該注意聖靈充滿的表現,藉著聖靈的能力,使福音更廣傳,不是為顯耀自己,讓人看見自己所得的恩賜；乃要高舉耶穌,放膽傳福音.
\newline
\newline
因時間關係,我們簡率來講.第十和十一章是很長的記載,請各位注意彼得,暫時不談哥尼流,雖然他禱告,他被稱為福音的對象；但福音使者是彼得,他這次的禱告產生了差傳的異象；為了順服,他違背他的傳統,違背了他的背景,違背了他從小所接受的風俗習慣；所以當差傳異象臨到彼得,第十章說:彼得「心裡正在猜疑.」彼得有所顧慮.有時受了神的啟示,不一定即時瞭解,可能神要我們作從未作過的事,是自己覺得無辦法作到的事,可能是引人誤會的事,是別人所不諒解的事.彼得繼續思想,繼續禱告,神的命令臨到他身上,說:「有三個人來找你,起來!下去!和他們同往,不要疑惑,因為是我差他們來的.」彼得順服,到了哥尼流的家,他得到新啟示,他明白了凡求告主名的人都可以得救.」神愛猶太人,也愛其他的人.彼得本不明白這真理；但是他禱告,差傳的異象就臨到他身上.
\newline
\newline
對華人教會超越文化差傳的工作,可能我們不容易了解,也不容易參與；以為向自己同胞傳福音已經夠了,全世界華人佔1/4,自顧不暇,還要向非華人傳福音,作超越文化的工作.可是,如果神揀選我們,如同揀選彼得,到哥尼流的家向非猶太人傳福音,我們肯去嗎?昨晚曾題起伍弟兄,他從香港到泰國,並非向華僑傳福音；因為神感動他的心,他想到香港有許多教會,聽福音的機會很多；而泰國許多鄉村,連一個教會都沒有,他說:「神要我去,我必須去.」
\newline
\newline
第十一章,我們看見彼得傳遞異象,非常寶貴!當時了不起的彼得被耶路撒冷領袖召來,查問為何進入外邦人的家?若是五旬節前的老彼得,一定拍案厲聲地說:「你們不知我是教會的柱石嗎?難道還懷疑我嗎?我要做就做,要去就去,不必問為甚麼.」但是彼得那時卻謙卑回答,按著神的帶領把事情次序講解,這是彼得平時難以做到的；他平日做事沒條理.他們起初不贊成彼得所為,並加以批評.彼得有條有理講解神的帶領,說:「神的帶領產生於禱告,哥尼流也禱告；我到那裡一開講,聖靈便降在他們身上,正如降在我們身上一樣；可見我的禱告,不是亂想,我的感動,是神的帶領.」彼得解釋之後,他們閉口不言.感謝主!他們不再批評,問題解決了,而且他們歸榮耀與神.
\newline
\newline
巴不得現今的教會,學習初期教會,可以減少許多的批評.
\newline
\newline
福音要傳進歐洲(十六章)保羅到了歐洲,到一個禱告的地方.從差傳策略看來,保羅差傳的對象,是些懂得禱告的人,他們未信耶穌之前,就懂得向真神禱告,保羅覺得這些人,心已有預備,所以就帶領他們認識耶穌.
\newline
\newline
請各位注意!保羅的觀念和今天許多基督徒的觀念大不相同,許多人以為道道通羅馬,只要信教,不論信何教,誠實你所信的教便可.
\newline
\newline
保羅到歐洲放膽向那些已懂禱告的人,傳耶穌基督的福音.聖經說:「主開導呂底亞的心,使他放下成見,接受耶穌,他和他全家都歸主.」我覺得保羅很有智慧,當時如果他沒有離開,很可能教會分裂；保羅在禱告的地方也沒另設團體,他繼續在禱告的地方,那時正有傳福音的機會,也有幫助他人的機會.第三章記載約翰和彼得進聖殿禱告時,幫助一個瘸腿的；同時保羅到禱告的地方途中,幫助一個被鬼附的女孩子；她很可憐,不自由；人利用她賺錢,她每天跟在保羅後面叫喊.有一天,保羅轉身,奉耶穌的名吩咐鬼離開,使她得釋放,得自由.
\newline
\newline
第十六章,還有一次的禱告,保羅和西拉被下監,他們挨打,全身痛苦,在黑暗潮濕的羅馬監牢中,半夜高歌禱告.
\newline
\newline
弟兄姊妹!求主給我們有新的看見,教會的差傳產生於禱告.初期教會如此,今日教會也要如此.
\newline
\newline


\newpage

\section{第53屆~港九培靈會~戴紹曾~第三講}
\label{sec:595}
\textbf{神的話與差傳}
\newline
\newline
連結: \href{https://www.hkbibleconference.org/session-message/view/595}{\texttt{https://www.hkbibleconference.org/session-message/view/595}}
\newline
\newline
\hyperref[sec:593]{< < < PREV SERMON < < <}
~
\hyperref[sec:toc]{[返目錄]}
~
\hyperref[sec:596]{> > > NEXT SERMON > > >}
\newline
\newline
當福音傳入中國時,在一八○七年,馬禮遜離開英國到中國；可以說,他的工作是以聖經為基礎.他自己沒有設立很多的教會,也沒有帶領很多人信耶穌；可是他一到廣州,即開始把聖經翻成中文,經許多困難.十二年之久才把新舊約全書譯成中文.欲知詳情請參閱「聖經與中華」該書寫馬禮遜如何學習華語.當時,外人學華語實在不容易；因清廷不准中國人教授外國人.據說馬禮遜的老師不是基督徒,他身藏毒藥,一旦被清廷發覺,擬服毒自盡.學習語言,如此困難!那時香港一帶地方,天主教極力反對馬禮遜,不容他留在澳門.印刷方面,從前均用木刻字排版,在華南都有白蟻為患,刻成的字版常被白蟻毀壞.東印度公司也非常反對馬禮遜的工作,他要離開英國時,該公司不讓他搭他們的船；所以他只好搭比荷輪船,先到美國而後轉到中國.有人問馬禮遜:「你到中國作甚麼?」他答:「傳福音」.那人說:「你以為像你這麼年青的人,能夠影響古老的國家嗎?」他說:「我不能,神能.」神藉著自己的話,影響了那時的中國人.今晚我們在此舉行培靈會是怎樣開始的呢?乃在百七十多年前,有個英國青年馬禮遜到中國,學習中文,把聖經譯成我們的語言.感謝主!是由聖經起首的.
\newline
\newline
如今在許多宣教工作上,仍然是如此；數月前我在泰國北部,住在同工家裡幾天,在少數民族中間,看見不少同工正作翻譯工作.如今海外使團一些同工也進行翻譯聖經的工作；因為還有些地方的人,沒有自己的聖經.在中國也有許多弟兄姊妹渴慕神的話,去年我在中國,曾看見他們的手抄本；因為在某時期,聖經都被焚毀了.有一個愛主的老弟兄,他家曾遭三次搜索,為的是找聖經焚燒.到了一九七八年,聖經在中國極其缺少,倖存一本,大家便借用抄錄,有一位弟兄用了五年時間,把全部聖經抄了兩次.未知在座弟兄姊妹,愛慕主的話的程度如何?也許家中有聖經多本,卻很少打開,真是可憐!飽經風雨的教會,弟兄姊妹懂得勝經的寶貴；在自己生活中,福音工作上,神的話多麼重要.
\newline
\newline
聖經是神的話和差傳工作是不能分開的.讓我們看初期教會,神的話扮演甚麼角色?當時聖經,不比現在方便,大部份用口傳；在一些會堂中的領會者可能有古卷,每逢安息日應用；一般家庭極少有,幾乎沒有.保羅給提摩太寫信,吩咐他把那些古卷帶來；是否聖經,不得而知,但很可能是聖經.
\newline
\newline
我們讀行傳,全書不斷題到神的話,神的道,聖經.逐字引用舊約經文,至少有廿五次以上題舊約所發生的事.不計其數次題神的話.當時使徒們對於神的話,其態度值得今日傳道人,宣教士,事奉主的弟兄姊妹所注意的.
\newline
\newline
第六章載,教會面臨分裂,有些寡婦發怨言；使徒們極關心此事,立刻召集全會眾,說:「我們撇下神的道,去管理飯食,原是不合宜的.」這裡潛伏一種危險,教會有屬靈的供應和物質的供應,彼此攻擊,兩方面都要做；寡婦被忽略,眾弟兄姊妹也需要神的話,二者不可兼,到底熟輕熟重?使徒站穩立場,說:「我們不能撇下神的道,我們要專心以祈禱傳道為事.」這話極其寶貴!可見他們對神的話何等看重.結果「神的道興旺起來.」(7)因為使徒們不肯撇下神的道,不肯輕易將神的話交給別人,他們守著本位,遵行耶穌的吩咐.
\newline
\newline
有時我想,今天神的話不能興旺起來,是否在香港在台灣在華人教會中,專心祈禱以神的道為事的人沒有這樣做呢?也許有人要說,保羅也是帶職事奉的；不過,那是不得已的.第十八章載,保羅在哥林多織造帳棚；因為當時哥林多教會不明奉獻的真理；加以有些所謂「走江湖」之輩,滲雜教會中.保羅要弟兄姊妹認清他傳福音非為錢,乃為搶救失落的靈魂.第十八章題保羅在哥林多時織造帳棚,當西拉和提摩太來的時候,他立刻放下工作,專心傳道.
\newline
\newline
感謝主!我們看見使徒們的見證,他們重視神的話；偉大的宣教師保羅,也同樣重視神的話,初期教會對神的話何等重視!
\newline
\newline
初期教會有一個稱呼「求告主名的人」,由此可知使徒時代的教會,是個禱告的教會.另有一個稱呼「信奉這道的人,」「這道」就是神的道.馬丁路德時代,高舉神的話,以神的話為信仰生活之準則.各位青年基督徒!不應有雙重標準；在教會做個規矩的基督徒；進入學校,踏入社會,接受社會道德觀,接受社會論理觀；成了個兩面的人.使徒時代的基督徒被稱為「信奉這道的人」,他們在日常生活中,進行神的道,以身作則,到處將神的話活出來.在行傳中,每題大衛的話必題明乃由聖靈感動而說.題摩西說的話,乃由神的靈感動而說.他們認識神的話,不是傳統,不是哲學,乃是神的道聖靈感動人,真理來自神.
\newline
\newline
初期教會引用神的話有多方面；不過,我們著重差傳方面,餘者略為一題.
\newline
\newline
初期教會的禱告,他們常引用聖經的話.在第四章26節,他們提醒神,早日大衛說:「世上的君王一齊起來......要敵擋主,並主的受膏者.」禱告的人對主說:「現在這事已經應驗,且發生在我們中間.」摩西禱告向主講話之時,他提醒主說「主啊!你是信實的,你以往怎樣說,現在要怎樣作.」我們敢不敢向主這樣禱告呢?我們應當抓住主的應許向主禱告.
\newline
\newline
第二方面,我們特別注意他們傳福音,講道時怎樣應用神的話.行傳中至少有廿四篇講章,九章是保羅講的,十章是彼得講的；還有司堤反講的,也有其他的人講的.他們講時善用神的話,結束時總有個目標；不單講些舊約聖經故事,不只敘述舊約歷史；無論保羅或彼得,都有個共同目的,吩咐聽眾要悔改接受耶穌為救主,罪才能得赦免.傳道並非講完故事就算了,他們講道過程中,有時題神的話,講一些預言；且說:「看哪!舊約的預言現已應驗了,從這樣的神跡可以看見神的作為,看見神的偉大.」
\newline
\newline
王明道先生曾說:「以色列復國帶給我極大的鼓勵；因為神早在祂的話語中,應許了這事.亡國千多年的國家竟在一九四八年復國,完全照著神的應許.在我受逼迫時,經重大考驗時；神應許的應驗,堅固了我的信心.」
\newline
\newline
聖靈降臨,當時許多人都不明白.彼得告訴他們,這是約珥所應許的,末世時有此現象.主耶穌的死,許多人不明白他們的彌賽亞要受苦,受死,復活.彼得和保羅把聖經的預言,讓聽眾明白,相信這位死而復活的耶穌.他們常用聖經預言證明神的大作為.彼得說:主耶穌升天坐在神的右邊,大衛所講的話現在應驗了.耶穌基督是主乃根據大衛講的.神對亞伯拉罕應許要賜福給他,萬人也要因他蒙福.(創十一,十二)神召亞伯拉罕出來,應許要賜福他.也許我們想到流奶與蜜之地,就是那時神所應許給他的福.我們也想到猶太人很會賺錢,是否這福呢?彼得解釋,說明都不是這兩方面,乃是罪得赦免的,這福先臨到猶太人；其含意很寶貴,不只臨到猶太人,而且先臨到猶太人──當時代住在耶路撒冷的猶太人.另有一個含意,這福份也要臨到凡求告主名的人必然得救；我們也可享受神應許亞伯拉罕的福.
\newline
\newline
司提反引用聖經時,不大講到預言,他好像個歷史家,他講的是一篇歷史,目的要聽眾明白,真正敬拜神的意義,明白聖殿是甚麼；所以他講解一段很長的歷史,以證明兩件事:
\newline
\newline
第一,神是無所不在的神.
\newline
\newline
第二:從舊約歷史看見,總有人硬心抗拒神,神興起的先知,許多人不接受,拒絕他們,殺害他們；到了今天,仍有人拒絕耶穌.司提反講歷史,目的乃高舉耶穌,勸聽眾接受耶穌,不要硬著心拒絕耶穌,免得神的刑罰臨到.
\newline
\newline
腓利受感動,走過去和埃提阿伯太監談話,他正讀以賽亞書不懂得解,就問腓利.腓利藉著該段經文向他傳講耶穌,從神的話介紹耶穌.結果,這北非人聽了福音,接受了耶穌.類似之例實在很多.
\newline
\newline
弟兄姊妹!用神的話要有智慧,尤其在差傳工作上,濫用神的話,不能感動人；因所用的方法,不符合聽眾需要.
\newline
\newline
上列數例是使徒向猶太人傳道,他們講舊約,講經文使聽眾明白.當使徒向外邦人傳福音時,他們引用神的話就不同了,可能不是逐字引用,保羅工作上有實例:他在安提阿會堂講道(十三章)開始時對猶太人說:「以色列人和一切敬畏神的人.」由此可見聽眾有兩種人.行傳中曾數次介紹「敬畏神的人」,好像哥尼流,呂底亞,他們原非猶太人,亦非以色列人,他們開始接受猶太教的影響,叫做敬畏神的人,他們的背景不同,可是稍為接觸舊約經文.十三章保羅講道,顯然是對一些懂得勝經的人,因他不斷引用聖經.如果我們再往下看,十六,十七章保羅在歐洲,腓立比,帖撒羅尼迦,庇利亞；特別在雅典,他講道用詞不同.向猶太人他不斷引用聖經；對那些不明聖經,從未接觸聖經的人,他很少引用聖經的話.
\newline
\newline
我們應該學習,有時不是不用聖經,而是用的方法要特別小心,要按照聽眾的背景和理解力引用.
\newline
\newline
今日華人教會對福音與文化一事,有很大興趣；不久之前,在星加坡召開第二屆華福會題及此事,有些人要進行研究福音與文化的關係,這是值得做的；可能一百五十年前許多地方因西教士對中國文化不夠了解,福音傳開,要求加上和減少一些,必要和不必的東西；所以研究這問題,使我們在差傳工作上,更有效地傳福音；不過,有句話要題醒大家:當進行時,對神的話夠不夠熟悉?台灣有位很愛主的弟兄寫了本書,他的背景相當不錯,讀大學時研究文學；但是對神的話不夠清楚,對解釋孔子思想也許他是權威；可是解釋聖經,或兩者參在一起時,卻有點危險.我常勸這位弟兄早日進神學院,然後寫這本書,可以寫得更好,可能對教會貢獻更大.
\newline
\newline
教導工作和差傳工作是分不開的,使徒時代的教會非常重視教導神的話；剛得救者,不論在耶路撒冷,安提阿,哥林多,如果不教導,他們怎能站穩?怎能以神的話為生活基礎呢?五旬節三千人歸主,五章42節有同樣記載他們如何專心遵守使徒的教訓.巴拿巴和保羅在安提阿,重視差傳工作的教會,在教會未開始差傳工作,先有系統的教導,使信徒和聖徒們站穩在神的話語上.這教會也成為差派工人的教會.
\newline
\newline
華人教會需要注意訓練教導的工作,特別是聖經的教導.我們的主日講壇,是有計劃的或是臨時的感動而講呢?我們可以在講台上供應弟兄姊妹各方面的需要.基督徒有種觀念,覺得臨時的感動較為準確,長期的計劃乃出於人.我覺得很奇怪,為何長期計劃不能出於神呢?為何臨時感動才算是神的感動呢?摩西在山上四十晝夜,神把會幕藍圖告訴他,神把十條誡命的律法告訴他；也是經過了很長的時間.神是有計劃的神.教會當有計劃教導聖徒.我覺得主日學是個很好下手的地方,藉此可對全教會進行教導工作.以往我們對主日學抱著狹窄的觀念；很多教會以為主日學是兒童主日學,所以成年人和老年人都沒參加.近數年來,漸漸有些教會改變了,許多教會的主日學是為全教會預備的,各等年齡都有自己的班,按照程度,按照生活需要,有系統查考聖經.
\newline
\newline
數年前,張明哲博士寫「傷亡率」一書,題及青年,特別是學生工作.有的青年讀大學時信耶穌,離了大學進入社會謀事,結婚,難免遇見新考驗.張博士說:按統計,過了十年,本來愛主的青年,有百分九十與教會脫節.真是可怕的現象!我相信有很多原因,暫不分析；但是我覺得如果華人教會,給廿五至卅五歲的青年,在教會有系統查考聖經的機會,大可減少可怕的傷亡率；因為這段時間,是青年遇考驗最厲害的時候,最迫切需要神道的時候；如果教會有他們的查經班,亦即他們的主日學班.年齡差不多的青年剛踏入社會,他們可以一同根據聖經,討論家庭問題,社會問題,道德問題,和個人工作上遭遇的困難,遵照聖經神的話的教導.
\newline
\newline
如果教會的主日學,不單為小朋友,也是為全教會；我相信神要大大賜福與教會,神的道要興旺起來,福音要傳開,不只在香港；我們將看見差傳的工作進行.
\newline
\newline
願主賜福祂自己的話語.
\newline
\newline


\newpage

\section{第53屆~港九培靈會~戴紹曾~第四講}
\label{sec:596}
\textbf{基督徒與差傳}
\newline
\newline
連結: \href{https://www.hkbibleconference.org/session-message/view/596}{\texttt{https://www.hkbibleconference.org/session-message/view/596}}
\newline
\newline
\hyperref[sec:595]{< < < PREV SERMON < < <}
~
\hyperref[sec:toc]{[返目錄]}
~
\hyperref[sec:603]{> > > NEXT SERMON > > >}
\newline
\newline
使徒行傳 8:1-8:8
\newline
\begin{longtable}{cl}
\hline
\hline
章節 & 經文 (和合本修訂版)\\
\hline
8:1 & \begin{tabularx}{0.7\textwidth}{X} 掃羅也贊同處死他。從那一天開始，耶路撒冷的教會遭受到大迫害，除了使徒以外，眾門徒都分散在猶太和撒瑪利亞各處。 \end{tabularx} \\ \\ \relax
8:2 & \begin{tabularx}{0.7\textwidth}{X} 有些虔誠的人把司提反埋葬了，為他大大哀哭。 \end{tabularx} \\ \\ \relax
8:3 & \begin{tabularx}{0.7\textwidth}{X} 掃羅卻殘害教會，挨家挨戶地進去，拉著男女關在監裡。 \end{tabularx} \\ \\ \relax
8:4 & \begin{tabularx}{0.7\textwidth}{X} 那些分散的人往各地去傳福音的道。 \end{tabularx} \\ \\ \relax
8:5 & \begin{tabularx}{0.7\textwidth}{X} 腓利下撒瑪利亞城去，向當地人宣講基督。 \end{tabularx} \\ \\ \relax
8:6 & \begin{tabularx}{0.7\textwidth}{X} 眾人都聚精會神，同心合意地聽腓利所說的話，一邊聽他的話，一邊看他所行的神蹟。 \end{tabularx} \\ \\ \relax
8:7 & \begin{tabularx}{0.7\textwidth}{X} 因為有許多人被污靈附著，那些污靈大聲呼叫，從他們身上出來；還有許多癱瘓的、瘸腿的都得了醫治。 \end{tabularx} \\ \\ \relax
8:8 & \begin{tabularx}{0.7\textwidth}{X} 那城裡，有極大的喜樂。 \end{tabularx} \\ \\
[1ex]
\hline
\hline
\end{longtable}
宣教士包括哪些人?誰是被差派的人?參與差傳工作的是哪些基督徒?可能有人馬上回答:初期教會所差派的是使徒.耶穌揀選了十二個門徒和接續他們工作的人.可是,我們看使徒行傳,發現參與差傳工作的,不只那十二個門徒；司提反,腓利,巴拿巴,保羅,西拉,提摩太,他們都參與差傳工作.
\newline
\newline
差傳人選不能限定狹窄.基督耶穌打發門徒時,馬太廿八章最後幾節經文,耶穌對門徒說:「所以你們要去,使萬民作我的門徒.」這是工作的目標.接著,主耶穌又對門徒說:「凡我所吩咐你們的,(包括「你們要去」這個命令)你們都要教導他們遵守.」所以,這些要聽福音,要接受耶穌基督的人,從使徒領受耶穌基督給他們的使命─差派的使命,使萬民作門徒的使命.這使命不只是給那十二位門徒；因著他們的見證而成為門徒的基督徒,也要遵守這吩咐；可以說十二位門徒是代表所有的門徒,他們代表全教會；主怎樣吩咐他們,也就是怎樣吩咐我們.使萬民作耶穌基督的門徒,是所有門徒的責任,是全教會的責任,每個基督徒都有責任.這是我們必須明白的事實.
\newline
\newline
第二問題是工作地區的問題.到底工場多廣?第一章8節,耶穌對門徒說:「並要在耶路撒冷,猶太全地,和撒瑪利亞,直到地極,作我的見證.」由此可知,工作範圍何等廣泛.
\newline
\newline
這幾個晚上,我們已漸漸看見,神如何開導門徒的心；特別是彼得,他本不明白這真理,他還不知道福音要傳到地極；雖然耶穌在四福音已有這樣的命令,在行傳一章也有同樣的命令；可是彼得不完全明白.到了第二章,聖靈已降臨在教會,彼得向他們解釋,這是約珥所應許的,是舊約的應許.舊約時代,先知已有這樣的看見.早在約珥先知時,神藉祂的僕人,已經把工作範圍題出「我要將我的靈澆灌凡有血氣的.」不只猶太人,也包括希利尼人,並包括全世界各國的人.到了第二章,五旬節彼得講道時說:「凡求告主名的人,都必得救.」當時彼得用「凡」和「都」字,他自己不知道範圍到底多廣；他以為凡猶太人,凡是自己的同胞都包括於這範圍內.到了第十章,彼得才醒悟「說......我真看出......原來各國......都為主所悅納.祂是萬有的主......祂是一切的主宰,是世上所有的人的主.」(34-36)「凡信祂的人......聖靈降在一切聽道的人身上.」(43-44)由此,我們看到,神如何帶領彼得一步步明白這個真理.我想彼得的反應和我們的反應是一樣的,首先關心自己的人.
\newline
\newline
舊約聖經有個不遵命的宣教士約拿,神吩咐他往西到尼尼微城,他偏偏往東到他施；因為他不喜歡尼尼微城的人脫離神的忿怒,不願他們得救.親愛的弟兄姊妹!今天我們基督徒,我們這些作門徒的,已經領受主賜的福份；主的恩典在我們身上是那麼豐富,我們的責任如何?
\newline
\newline
在華福會議,兄弟代表基督使團,有機會發言,我把心裡負擔告訴大家,也向華人教會提出挑戰.我特別提及日本；其實,華人感到向日本傳福音是件難事,因半世紀之前,日本在中國及東南亞一帶的行為,於日本手下受害的華人實在很多；且人們都認為日本是個工業發達,教育程度很高的國家；各方面及生活條件都很不錯,毫無想到他們是傳福音的對象.可是,以日本人口而言,是世界第六大國,亟需福音.基督福音使團,單在日本,就需要廿五位同工.我們向主禱告,是否在廿五位新同工之中,差派14位男同工.感謝主!香港差派的宣教士已有蘇恩覺姊妹代表,她在日本北海道忠心傳主福音.但是,那邊的需要很大,那裡的人極其需要耶穌基督的福音.誰肯關心呢?
\newline
\newline
我們又想到菲律賓,也許認為這是個天主教國家；其實,在南島很多回教徒；許多部落根本沒有福音.海外基督使團只有一位女同工在那邊.弟兄在哪裡?還有哪些人可去協助呢?須要把聖經翻譯成當地的話,在少數民族之中,須要設立教會,需要新的同工.
\newline
\newline
至於泰國,人口有四千六百多萬,福音在泰國已有百五年；然而信主的人,不到十萬.感謝主!從香港已有些代表到泰國,有的向華人傳福音,有的向難民傳福音,有的開始做超越文化的差傳工作.
\newline
\newline
上述三個國家,不過是代表東亞的需要.到底福音範圍多廣呢?使徒行傳告訴我們:「要往普天下去.」巴不得主開導我們的心胸,把世界放在我們心裡；使我們以基督的心為心,以耶穌基督的事為念.
\newline
\newline
差傳使命的範圍(八1-8,十一19-21,十三1-3)人才方面包括全教會.工作地區方面包括全世界.
\newline
\newline
有三種不同的宣教士,(關於這方面,我得到菲律賓青年工作副總幹事瑪嘉烈之助).
\newline
\newline
第一,守本位者.當逼迫臨到教會,基督徒四散,大部份不能留在耶路撒冷；惟有使徒們留下來,我想他們絕非避免逼迫.路加敘述這段歷史說:「掃羅卻殘害教會,進各人的家,拉著男女下在監裡.」(八3)這是當時耶路撒冷的情形.(二十六10-11)告訴我們,留在耶路撒冷是不能避免逼迫的；但是這些使徒,因帶著使命,所以留下來,他們要專心以祈禱傳道為事,他們覺得留在耶路撒冷是他們的本分；雖然面臨危險,仍持守本位.
\newline
\newline
抗戰時期,神帶領家父留在中國；我相信那時他作此決定是很難的.本來他決定把全家從日本佔領區帶去美國；因他的工作既不能進行,孩子們也有危險,為父牧者自然會這樣做；可是家父禱告,我相信是聖靈感動他,他改變原意；家父母到了西北去,在那裡繼續傳福音,設立聖經學校,訓練同工.孩子們都留在山東.感謝主!去年我回去,有機會看見我父母的許多同工,他們有不少的學生就在那時得救.神的美意,當時讓我父母決定不離開本位,願意冒著危險繼續工作.到了去年,我才明白為何那時父母撇下孩子們,五年多不相見,因為神要拯救一些人.我也看見我父母的同工們,卅多年來忠心為主而活,忠心為主作見證；是否當時受訓練,更清楚認識主的話.宣教士持守本位是不易的,相信在一九四九年,有千千萬萬基督徒,尤其是傳道人,都考慮離開中國,等平靜下來,再回到自己的地方；可是經過禱告,決定不能離開,就忠心留下來帶領神的教會.我聽過一個見證:當一九七五年越南將變色時,有位牧師,該宗派的教會派他去美進修；那時留在越南相當危險.有人問他「你是否離開?」他說:「牧者離開,神的群羊怎麼辦呢?」就在那年四月,西貢淪陷了.他留在本位,持守本位.
\newline
\newline
我們可以說,那些使徒就是宣教士,在聖靈帶領之下,忠心守住神給他們的使命；面對危險,靠著聖靈,不懼怕,他們成了中流砥柱.這種宣教士所接受的使命最難.
\newline
\newline
第二,隨走隨傳者.這是主自己吩咐門徒的.當逼迫臨到,這班人覺得不能留在耶路撒冷,不得不離鄉別井；逼迫並沒攔阻他們傳福音,雖無家可歸,但耶穌基督的福音一定要傳.所以他們隨走隨傳.路加敘述一些被聖靈充滿的人,他常說:「他們放膽傳講神的話.」
\newline
\newline
例如腓利,是個隨走隨傳者,他本非十二門徒之一；當教會面臨分裂之時,教會選出七位有好名聲,被聖靈充滿,有智慧的人,其中有腓利,他有好見證,有能力地為主作見證.到了廿一章8節,路加稱他為「傳福音的腓利」,(英文稱為佈道家)他有力地傳福音,他的見證從他的家開始.聖經記載,腓利有四個女兒,都有好的見證,都說預言,都能傳講神的話.我相信她們的父母影響力很大.這位傳福音的腓利,絕非逃脫不負家庭責任；他傳福音乃基於他的家庭先有好見證,然後他出去傳福音.本來教會選他管理飯食,因為使徒們曾說,「撇下神的道去管理飯食,是不應該的.」腓利是位有好名聲,被聖靈充滿的,讓他管理飯食,未免大材小用；但是腓利在耶路撒冷時,持守本位,忠心管理飯食.我們在小事上忠心事奉主,在大事上主也要用我們.
\newline
\newline
各位青年,萬勿想在教會中,一躍登天,應一步步事奉主；在微不顯眼的小事上忠心,必有機會在大事上事奉主.如果在小事上不肯忍耐忠心事奉,恐怕有一天在大事上會挫敗.你看!忠心管理飯食的腓利,當逼迫臨到時,他離開耶路撒冷到撒瑪利亞.該地的人一向不歡迎猶太人,照樣,猶太人也不喜歡撒瑪利亞人,他們彼此鄙視如犬,他們之間的仇恨很深.門徒曾到過撒瑪利亞,看見該地的人硬著心,不接受耶穌；就對耶穌說:「你要我們禱告,求父降下火來,燒滅他們嗎?」彼得這樣的態度當然不對,耶穌嚴嚴責備他.腓利雖然無家可歸,他來到撒瑪利亞；放膽為耶穌作見證,帶領多人歸主.後來他從北方到南方,從都市到僻壤,在曠野向那位太監埃提阿伯傳福音.他對於主的帶領,有敏銳的反應；主叫他北上,叫他南下,他都聽從.雖然歷史對這事記載不詳；但今日在埃及阿比西尼亞的科普特教,是否當時因埃提阿伯太監一人接受耶穌的緣故,有人說很有密切關係.是否腓利隨走隨傳,在撒瑪利亞,在曠野地方影響了北非.
\newline
\newline
此次在華福會,有幾位講員特別強調一件事,我也起了共鳴.鮑會園院長題及基督徒為主作見證.唐崇榮弟兄題到今日許多基督徒無動於衷,在日常生活中為主作見證.正在那段時間,我有很特別的經歷.華福會之前,有一晚我剛從海外回來,約有六週時間在外奔跑,到了新加坡已相當累,理應多事休息；可是過了兩天,我必須到西馬；因海外基督使團的學校在那裡,該校是從前山東煙台學校的後代.名叫曲阜學校也叫煙台學校,是宣教士的子弟學校,也是我的母校,時逢百週年紀念,於是勉強搭夜車,心裡向主埋怨.主責備我,提醒我此行有工作.車開了,我發現有個馬來亞青年坐在旁邊,他是新加坡控制吊動車的工人.我開始對他傳福音,發覺他的心完全預備好了；正等人傳福音給他,當晚他表示願意接受耶穌.我低頭向主認罪,因我發怨言上車,沒想到主把工作放在車上；要我在車上傳福音,帶領這青年認識主.
\newline
\newline
隨走隨傳的宣教士,其中有些無名英雄,聖經沒有題名.第十一章19節:「那些因司提反的事遭患難四散的門徒......」四散的門徒也是隨走隨傳的宣教士,他們北上,一面離開耶路撒冷,一面傳福音；有的向自己的同胞猶太人傳福音,部份有新突破的也向希利尼人傳福音.主與他們同在,信而歸主的人很多.這是安提阿教會的開始,是教會做差派工作的開始.安提阿教會的基礎,是隨走隨傳的宣教士；神使用他們,把福音傳開了.
\newline
\newline


\newpage

\section{第53屆~港九培靈會~戴紹曾~第五講}
\label{sec:603}
\textbf{教會與差傳}
\newline
\newline
連結: \href{https://www.hkbibleconference.org/session-message/view/603}{\texttt{https://www.hkbibleconference.org/session-message/view/603}}
\newline
\newline
\hyperref[sec:596]{< < < PREV SERMON < < <}
~
\hyperref[sec:toc]{[返目錄]}
~
\hyperref[sec:604]{> > > NEXT SERMON > > >}
\newline
\newline
使徒行傳 13:1-13:3
\newline
\begin{longtable}{cl}
\hline
\hline
章節 & 經文 (和合本修訂版)\\
\hline
13:1 & \begin{tabularx}{0.7\textwidth}{X} 在安提阿的教會中，有幾位先知和教師，就是巴拿巴和稱為尼結的西面、古利奈人路求，與希律分封王一起長大的馬念，和掃羅。 \end{tabularx} \\ \\ \relax
13:2 & \begin{tabularx}{0.7\textwidth}{X} 他們在事奉主和禁食的時候，聖靈說：「要為我分派巴拿巴和掃羅去做我召他們做的工作。」 \end{tabularx} \\ \\ \relax
13:3 & \begin{tabularx}{0.7\textwidth}{X} 於是他們禁食禱告後，給巴拿巴和掃羅按手，然後派遣他們走了。 \end{tabularx} \\ \\
[1ex]
\hline
\hline
\end{longtable}
差傳的同工是由教會打發出去的,他們工作完畢回到教會,報告傳道工作,神在他們身上的恩典.
\newline
\newline
昨晚題到三種不同的宣教士,一種是持守本位的同工,在極大危險之中,使徒仍然留在耶路撒冷；主安排他們在那裡工作,他們持守本位.第二種是隨走隨傳的同工,他們因著逼迫不能留在耶路撒冷.解經家告訴我們,特別是生長在猶太以外的猶太人,他們遭受更厲害的逼迫,以致無家可歸；可是他們沒忘記自己的主,他們離開家庭卻不離開主；主給他們使命為祂作見證,所以他們邊走邊傳.他們北上,神使用他們.路加說:「主與他們同在,信而歸主的人很多.」可能有人說,今之教會不需要這種同工；有持守本位的同工,有隨走隨傳的同工夠了.
\newline
\newline
近數年來在各地舉行大規模的福音會議,都特別強調這件事,現世界人口已增至40億,其中基督徒最多僅佔十億.(此統計乃包括天主教和一般掛名基督徒在內)換言之,世上還有四分之三的人不認識耶穌；且卅億世人中,僅十億人住的較接近基督徒,持守本位的同工可以向他們傳福音,不必被差派出去.但是,世界人口還有半數以上(廿億)不靠近基督徒,也沒有教會,這樣的人怎得聽福音蒙拯救呢?或有些隨走隨傳的人可以把福音帶給他們,可是相當有限.針對這廿億人口,差派同工出去的問題,是值得考慮的.由此可見差派的重要性,因為全世界半數以上的人口需要福音,若沒有人把福音帶給他們,就永遠聽不見耶穌基督的名字,永遠聽不到福音.保羅說:「凡求告主名的,就必得救.然而人未曾信祂,怎能求祂呢?未曾聽見祂,怎能信祂呢?沒有傳道的,怎能聽見呢?若沒有奉差遣,怎能傳道呢?」(羅十13-15)廿億人口,等待基督徒關心,他們的靈魂,需要被主差派的人；才能得到福音,得蒙拯救,和我們同享救恩之樂.
\newline
\newline
昔日安提阿教會作差派事工,有三方面的特點:
\newline
\newline
第一個特點　安提阿教會開始有基督徒的名稱.一班使徒到安提阿傳福音.教會便成立,基督徒的名稱是從安提阿開始的；顧名思義,一定是該教會的弟兄姊妹,有特出的表現,給人感覺他們有種力量,基督的能力在他們身上；他們不是為自己活,乃是為基督而活；他們生活的表現,實在令人想到耶穌基督.相信另有個原因,他們整天嘴唇講基督；勸人信基督,告訴人信耶穌基督的好處,且告訴人如何信法,基督徒應如何生活.
\newline
\newline
各位兄姊!青年們!學生們!別人看見我們的生活方式,言行舉止；能否令人想到耶穌基督那樣的愛心,溫柔,關懷別人?我們生活表現能否與基督徒的名相稱呢?
\newline
\newline
第二特點　開始社會關懷,社會關懷亦即教會關懷.安提阿教會收奉獻送交其他教會,因聽說那處有飢荒,有孤兒寡婦.安提阿教會不只自己生活有見證,個人生活有見證,大家集中力量；基督的愛在他們心裡,他們收奉獻送到耶穌撒冷；幫助難民,幫助那些飢荒中缺乏日用飲食的人.
\newline
\newline
第三特點　安提阿教會作超越文化的差傳工作,打發巴拿巴和掃羅出去；工作以教會為中心,差傳工作與教會是不能脫節的.持守本位的同工以教會為中心,隨走隨傳的同工也出發於教會.在聖靈帶領之下,安提阿教會打發兩位同工出去,出發點仍然是教會；他們出去有很清楚的目標,清楚的使命,要設立教會.安提阿教會的同工團在一年之中,至少設立了四個教會,而去處不止四個；教會增長,福音傳開.
\newline
\newline
安提阿教會是合一的教會　弟兄姊妹若不同心,若分門結黨,實在談不上差傳工作.弟兄姊妹同心,不一定是背景相同.第十三章告訴我們,那班人是來自不同的背景.那些到居比路去的,和古利奈人,都是無名信徒,直到今日,我們都不知道他們的名字.從聖經記載的人名,我們可知這些基督徒的背景.路加說:「在安提阿教會有幾位先知和教師.」這是教會的職份,是教會的恩賜,作先知的不一定說預言,可能述說神的救恩,包括兩方面.聖經上所記的先知,有的很特別,在舊約和新約,在行傳中,至少有兩次題及一位先知名叫亞迦布,他預言耶路撒冷將遭飢荒.在廿一章,亞迦布要攔阻保羅南下耶路撒冷,「就拿保羅的腰帶,捆上自己的手腳.說:聖靈說,猶太人在耶路撒冷要如此捆綁這腰帶的主人,把他交在外邦人手裡.」(11)但是保羅不顧而去,果然事情發生在他身上.先知的個性很特別,很強,作事劇烈,動作戲劇性；無論舊約新約都是這樣.我有時想,幾十年前中國的先知必是宋尚節,他在中國,東南亞作工,他的個性很強；有時候他講道,傳譯的人如果不夠快不夠順,他就把那人推下.
\newline
\newline
至於教師的個性是慢吞吞的,循規蹈矩,正與先知的個性相反；兩者若在一起,真是不得了!
\newline
\newline
安提阿教會,先知和教師一起事奉,一起禱告,一起禁食.
\newline
\newline
求主幫助我們,別堅持個性,不與別人合作；當然你不能,但是,神能幫助我們同心,祂已將各種不同的恩賜給祂的兒女.我們的身體,不只有眼睛,還有耳朵；不只有手,還有腳；如果身體只有一個肢體,那是多怪多可怕呀!身體的美,有不同的肢體；身體的利害,是不同肢體的運用.照樣,教會的美,在於不同的樣子,不同的表現,不同的恩賜.如果每個人都像你或像我,那麼,教會的力量就大大減少了.
\newline
\newline
利未人巴拿巴,生長於居比路.西面被稱呼尼結:(「尼結」英文即黑人)他是北非人.路求是羅馬人的名字.猶太人,非洲人,羅馬人都在安提阿教會,不分種族.馬念與分封希律王是在皇宮一齊長大的,出身於貴族.在安提阿教會,有名的,無名的；有地位的,無地位的都在一起.保羅也在其中,他是受過高等教育大有學問的人,他也謙卑和那些沒受過教育的人在一起.安提阿教會是個參與差傳工作的教會,是同心合意的教會,傾聽主的話.保羅在(以弗所書四章11-16節)描寫教會,他所描寫的,可以說是安提阿教會.安提阿是同心的教會,是禱告的教會；差傳的負擔乃產生於禱告,他們迫切禱告,禁食禱告.
\newline
\newline
今日的教會常忘記禁食禱告.在差傳事工上,我們不要忘記禁食禱告.多年來,神在內地會,在海外基督使團,藉著同工禁食禱告,幫助我們勝過很多難處.沒有禱告,沒有迫切禁食禱告；有時我們實在缺少帶領,缺少聖靈的能力.安提阿教會是個順服的教會,順服神的帶領.差傳的教會,必須順服神.
\newline
\newline
但是,請各位注意!神的帶領有很多不同的方法；單就這段教會歷史,可以看見各樣不同的方法.例如那些來自安提阿的弟兄姊妹,隨走隨傳的同工,他們的帶領,不是馬其頓異象,也沒有聖靈的感動；神用逼迫的環境來帶領他們,在逼迫的環境中完成神的美意.過了一段時間,耶路撒冷教會聞安提阿方面有事發生；這裡給我們看到第二種帶領,教會打發巴拿巴北上,不一定是個人的感動.親愛的青年們!我們要很小心,多少時候,基督徒有個人主義,自己沒有感動,不去,不幹.這裡我們看見教會的分派.昨晚我們也題及撒瑪利亞人聽見福音,乃由教會派彼得和約翰去的；都非出於個人的感動.教會打發巴拿巴出去,他即時遵命.求主給我們有謙卑的心；教會所安排的工作,我們須謙卑接受,或者你以為自己的力量,恩賜,經驗不夠；但是,教會的安排,是因牧師信任我們,覺得我們有這樣的潛力,靠神能力得以勝任.巴拿巴工作很忙,覺得需要同工,就帶保羅同去.保羅在加拉太書信特別強調聖靈在他身上的帶領；可是路加告訴我們,另一面是巴拿巴帶領他去的.是的,有時神藉著一個屬靈的人的勸勉,把祂的旨意表明出來.所以,有時我們的思想過於單純,滿以為自己沒有感動,也沒有馬其頓異象.以保羅來說,巴拿巴帶他去之前,也沒有馬其頓異象,只是一位屬靈的弟兄,到他的家把他帶出來的.神用各樣方法帶領祂的兒女.第四個方法是,聖靈在教會中作感動之工.有時藉著先知說預言,聖靈感動安提阿教會,帶領教會打發同工出去,初期教會是順服的教會,是遵命的教會；雖然神的帶領不一,方法不同,無論如何一概遵命.
\newline
\newline
巴拿巴的勸勉,堅固教會的信心；保羅有恩賜,講神學,幫助信徒們有堅強的信仰.還有一些無名的基督徒,來到安提阿,神使用他們；其實,那時教會可以從無名基督徒中選出數人派出去；但是,教會卻把最好的同工釋放出去.
\newline
\newline
今天華人教會,肯否把最好的打發出去?華人教會中,有許多基督徒父母,可能想把最聰明的孩子栽培做大生意,多多賺錢,或當醫生,或在政府機構做事；而把最沒出色,沒恩賜的孩子,讓他如有感動,就去傳福音,或是女兒,神有揀選,就讓她作主工.這種情形,真是可憐!
\newline
\newline
請看保羅,他是個受了高等教育的人,當他離開耶路撒冷到大馬色時,他未信耶穌,他手持一些證書.數月前我讀經時得到新的亮光,聖經說,那時保羅的證書是從祭司長得來的.是否年青的保羅,連以色列人的大祭司,祭司長,他都認識,與他們常過從.當保羅出發至大馬色,他的生命全面改變；因耶穌基督得到了他,他就稱耶穌為主.進到大馬色,那些證書無用了；他去找基督徒,在該地為主大發熱心.過了一段時間,他回到耶路撒冷,他沒有去找大祭司,也沒有找以往的朋友；他所找的是基督徒,這個小群,被人鄙視,被人厭棄,保羅願和他們在一起.保羅說:「我們不以福音為恥,這福音本是神的大能.」保羅所看見的,是今日許多基督徒所沒有看見的；他出去的時候,是天國的大使,是耶穌基督所特派的,要在世界講耶穌基督的福音；他昂首前往,因他知道所事奉的是誰.
\newline
\newline
想到宣教的歷史,在美國有位耶魯大學的畢業生名威廉波頓,他是當時百萬富商波頓牛奶公司的少東.他聽見主的聲音,要他去中國傳福音.有一天他想到差傳工作,當他看見工人搬木頭,一端有數人,另端僅一人；他心裡思想,如果自己去幫忙,當然幫助一人的那邊；因此他決定到中國傳福音,他撇下一切,離開家鄉,起程前往.到了埃及,他害病身亡；可是他的臉仍然朝向中國,要傳福音.
\newline
\newline
英國劍橋七位青年,拋下一切,獻身到中國傳福音.
\newline
\newline
教會歷史上,我們看見這樣的見證,使徒時代,近百年來,都有同樣的情形.
\newline
\newline
親愛的弟兄姊妹!華人教會面對差傳工作,我們的反應如何?求主幫助我們!神的靈在我們中間,祂在你心裡說話.我可以不必講完今晚的信息,明天可以續講.我深深感到神在我們的中間,祂要揀選我們.
\newline
\newline
我們看中國教會史,有許多見證.宋尚節離開福建到美國攻讀博士學位,得到化學博士銜；當時他心裡大受感動,感到神要用他.那時美國大學聘請他,他認為美國不需他的學位,中國更需要.他堅決回中國,那時中國南京北京數間大學都請他,他都拒絕；其實,教書也是正當的事業,照樣可以榮神益人；然而他決定將一生年日完全用於神的工作上.於是他帶著完全奉獻的心回國,在中國許多地方及東南亞一帶,忠心熱心為主作見證.今日許多地方的教會傳福音工作,乃因宋尚節順服神的帶領,願意把最好的奉獻給主.
\newline
\newline
主對我們說:「要先求神的國和祂的義.」我們願意順服否?神對安提阿教會說:「要把最好的打發出去.」我們願意嗎?
\newline
\newline
基督徒的父母!如果神要用你的兒女,你願意嗎?像哈拿一樣；像教會歷史上許多人一樣,用禱告,用金錢,用眼淚,用書信,來支持你的兒女,順服神,傳講主的福音.
\newline
\newline


\newpage

\section{第53屆~港九培靈會~戴紹曾~第六講}
\label{sec:604}
\textbf{差傳與事奉}
\newline
\newline
連結: \href{https://www.hkbibleconference.org/session-message/view/604}{\texttt{https://www.hkbibleconference.org/session-message/view/604}}
\newline
\newline
\hyperref[sec:603]{< < < PREV SERMON < < <}
~
\hyperref[sec:toc]{[返目錄]}
~
\hyperref[sec:605]{> > > NEXT SERMON > > >}
\newline
\newline
使徒行傳 14:19-14:28
\newline
\begin{longtable}{cl}
\hline
\hline
章節 & 經文 (和合本修訂版)\\
\hline
14:19 & \begin{tabularx}{0.7\textwidth}{X} 但有些猶太人，從安提阿和以哥念來，挑唆眾人，並且用石頭打保羅，以為他死了，就把他拖到城外。 \end{tabularx} \\ \\ \relax
14:20 & \begin{tabularx}{0.7\textwidth}{X} 當門徒圍著他的時候，他站了起來，走進城去。第二天，保羅同巴拿巴往特庇去。 \end{tabularx} \\ \\ \relax
14:21 & \begin{tabularx}{0.7\textwidth}{X} 保羅和巴拿巴對那城裡的人傳了福音，使好些人成為門徒後，又回路司得、以哥念、安提阿去， \end{tabularx} \\ \\ \relax
14:22 & \begin{tabularx}{0.7\textwidth}{X} 堅固門徒的心，勸他們持守他們的信仰，說：「我們進入神的國，必須經歷許多艱難。」 \end{tabularx} \\ \\ \relax
14:23 & \begin{tabularx}{0.7\textwidth}{X} 二人在各教會中選立了長老，禁食禱告後，把他們交託給他們所信的主。 \end{tabularx} \\ \\ \relax
14:24 & \begin{tabularx}{0.7\textwidth}{X} 二人經過彼西底來到旁非利亞， \end{tabularx} \\ \\ \relax
14:25 & \begin{tabularx}{0.7\textwidth}{X} 在別加講了道，就下亞大利去， \end{tabularx} \\ \\ \relax
14:26 & \begin{tabularx}{0.7\textwidth}{X} 從那裡坐船回安提阿去。當初，眾人就在這地方，把他們交託在神的恩典中，要完成現在所做的工。 \end{tabularx} \\ \\ \relax
14:27 & \begin{tabularx}{0.7\textwidth}{X} 他們一到那裡，就聚集了會眾，述說神藉他們所行的一切事，並且神怎樣為外邦人開了信道的門。 \end{tabularx} \\ \\ \relax
14:28 & \begin{tabularx}{0.7\textwidth}{X} 二人在那裡同門徒住了一段日子。 \end{tabularx} \\ \\
[1ex]
\hline
\hline
\end{longtable}
差傳工作是有目標有策略的.
\newline
\newline
有些基督徒以為工作不需要有計劃有策略,因這樣太不屬靈；只要靠聖靈,聖靈作工,工就完成.這是很大的錯誤,這種思想不準確.
\newline
\newline
從多方面可以看見智慧的神,工作是有系統,有次序,有計劃的.神創造天地,每天做每日的工.神吩咐以色列人建造會幕,有清楚的藍圖告訴摩西；如果不屬靈,神為何要把這麼清楚的計劃告訴祂的僕人呢?
\newline
\newline
保羅說:「時候滿足,神差遣祂的兒子到世上來.」連時間都是神安排的,「在神永恆的計劃中,神所預定的時間,耶穌基督降生在世上.」(加四4).
\newline
\newline
主耶穌做自己的工作,選了十二門徒,主有目的,門徒和主在一起,主打發他們出去；主自己要去的村莊,祂先打發門徒去,作預備的工作.五千人得餅吃飽,耶穌叫眾人坐下,有五十人一排,有百人一排的；耶穌供應眾人的需要,也是有計劃的.
\newline
\newline
從舊約到新約,都看見神是有計劃的神,祂的工作有計劃.當然我們工作之前要先禱告,讓一切計劃產生於禱告；當工作進行時,要用禱告托住工作；當工作完成,要以禱告獻上感謝.華人教會工作應有目標,有策略,工作應產生於禱告.
\newline
\newline
上月在新加坡召開第二屆華福會,會中提出未來十年的計劃,華人教會應走的路線,福音工作如何傳開,文字,佈道,差傳,神學工作等等問題.有人以為不必思想太多,誰知道十年後教會情形如何；不肯思想的人,如同沒有計劃的農夫.
\newline
\newline
求主幫助我們!牧師要有計劃,主日講台要有計劃；如無計劃,就會重複某方面的信息而忽略另方面重要的信息.傳道人也應考慮如何訓練新同工,在教會中運用同工的恩賜.多方面的工作,牧師,傳道人要有計劃.長執會也應有計劃,策略各方面的工作；主日學師資訓練,班級制度,教材是否合適,整個主日學教育要有計劃.政府教育有計劃,為何教會教育沒有計劃呢?青年工作,婦女工作,弟兄工作都應有計劃；否則,就沒有目標,每週節目索然無味.神學院應有計劃,栽培師資,充實圖書,可能整個課程也要修訂.今年你的教會,你的主日學,你的青年團契有沒有計劃?去年度有沒有計劃?我們對主的盼望是甚麼?我們盼望主在我們身上,在教會中,在香港作甚麼?
\newline
\newline
求主叫我們看見,工作要有計劃的重要!
\newline
\newline
從行傳中明顯看見使徒們的工作有計劃,他們先禱告,第十三章告訴我們,大家在一起敬拜主；禁食禱告,他們的計劃和工作,產生於禱告.「聖靈說,要為我分派巴拿巴和掃羅,去作我召他們所作的工.」(2)我特別注意末四字「所作的工」再看十四章26節巴拿巴和掃羅回到安提阿,「當初他們被眾人所託蒙神之恩,要辦現在所作之工......」這句又在此出現.出發之前,有一定的工；回來之後,向教會報告所作之工,他們差傳工作是有計劃的.
\newline
\newline
第十四章記載差傳工作有五方面:
\newline
\newline
(一)傳福音(21)保羅和巴拿巴二位同工出去作差傳事工首先傳福音.保羅說:「福音是神的大能,要拯救一切相信的人,先是猶太人,後是希利尼人.」(他們的聽眾常包括這兩種人)他們告訴聽眾怎樣得蒙拯救,罪得赦,與神和好,在耶穌基督裡過新生的生活.由此可見他們有計劃,傳福音的內容,地點,他們預先有計劃,選擇那些戰略性的都市.安提阿是個中心地帶,在當時羅馬帝國三大城；除羅馬和北非的亞歷山大外,它位居第三.他們出去傳福音,建立教會,並非無定向.保羅和巴拿巴每出外佈道,都是找中心區.第一次先到他們較熟悉之處.巴拿巴是居比路人,(即今之賽普魯斯)他領隊在兩個大都市傳福音.後來又到今之土耳其,那地保羅較為熟悉,他們先到伯賽大的安提阿,(非他們之出發地安提阿)然後到以哥念,路司得,特庇.這些都是重要地方.保羅本有計劃,他想第二個要地是以弗所；但聖靈攔阻,耶穌的靈不許可,因時機未到.於是保羅改變計劃,先到希臘國,腓立比,帖撒羅尼迦,比利亞,雅典,哥林多傳福音,設立教會,先環繞個大圈.到了第三次,神的時候到了,保羅就直接到以弗所去了,在那裡傳福音；以弗所成了早期教會的要地.還有許多地方,保羅無法去,他衡量優先的價值,注意很多不同的因素；從人口稠密,有會堂,語言相通,工作易於被人接納之處進行.他們在中心地區完工之後,其他周圍地方,我們看見有很美的見證；在大都市裡的弟兄姊妹,從中心區把福音傳開了.保羅寫信給帖撒羅尼迦教會,說:「不用我們說甚麼話,主的道從你們那裡已經傳揚出來.」(帖前一8)在希臘北部,馬其頓地方；南部,亞該亞,福音傳開了,真是奇妙!保羅選擇了有戰略性的都市,設立教會,讓弟兄姊妹把福音傳開.
\newline
\newline
我們在工作上,應該注意哪裡的工作影響較大,收穫較快,門已經開了,我們可以進去.保羅就是這樣計劃,神這樣的帶領他.
\newline
\newline
另方面,傳福音的方法不一,對象不同,聽眾也不一樣.並非福音改變了,保羅傳的福音並沒改變,只是傳達的方法不同；有時用的語言不同,平時他對著懂希臘語的人講道,所以希臘文成了他們所用的語言.可是,有一次在耶路撒冷,保羅用希伯來語講道,大家都安靜地聽.保羅明白聽眾的心理；如果當時他用希臘文或希利尼語,以色列人絕對不聽,因為當時他們最反對外國人.那時傳福音,最少有三種不同的對象,保羅常到猶太人的會堂,還有些進猶太教的希利尼人,路加稱他們為敬畏神的人；這兩種人較熟悉舊約摩西五經和舊約的先知書；保羅從舊約舉例,不斷引用舊約聖經的話,向他們傳福音.另有一種聽眾,在路司得那些迷信拜偶像的人,從未接觸舊約人物；可能他們對神的觀念都不正確,他們拜多神,拜人手所造的神.保羅向他們傳福音,要獻祭時,他們也把保羅當作神；保羅立即指出他們的錯誤,說,我們和你們一樣,都是人,沒有甚麼了不起,不要拜我們.然後保羅就從另方面把福音介紹給他們,從大自然神的創造,有規律,有次序,證明有神.最明顯是保羅在雅典講道,根本未提舊約,他先講創造天地的神；繼而講神的恩典；神是慈愛的,常降下雨來,太陽也晒在地上.這些都證明神的恩典,由此可知必有一位真神.保羅有時也講到審判,善有善報,惡有惡報,不是不報,日子未到.保羅向這些拜偶像迷信的外邦人傳福音,方法不同,我們應當學習.
\newline
\newline
傳福音是必然的,但是我們如何去傳呢?教會常用自己的方言,所謂方言即術語.不能否認,教會中的基督徒,大多受過教育,在社會上是有地位的.可能我們的思想,生活方式,和未受教育的人大不相同,我們向作工的人傳福音,如果堅持要他們照著我們的樣子；相信一輩子見不到他們歸向主,雖同是中國人,但思想,生活方式各有不同.我們是否願意遷就他們,越過文化的溝?馬其頓人對保羅說:「過到我們這裡來.」保羅便過海去了.我們傳福音,肯否過去那邊,學習他們的語言及生活方式?
\newline
\newline
在明朝,耶穌會傳福音入中國,利馬竇和其他同工,特別注意服裝,以為服飾應像中國人；他們看見中國的和尚,自認是西方來的和尚,也穿上和尚的衣服.有人警告說:「錯了!因為一般人對和尚的印象不好,看他們個個懶惰.」因此,耶穌會的人便改穿儒教的老師裝.戴德生來中國,正值清朝,他頭梳辮子,身穿中國服裝；因他認為應該投入接近聽眾.如果我們向工人傳福音,我們肯這樣做嗎?在台灣,工人們上主日禮拜,服裝隨便,腳穿木屐.牧師對此很不滿.若迫他們在語言,服裝,禮拜內容,各方面像我們,未免使他們覺得尷尬；所以我們當改變方法,符合他們的思想,適應他們的生活方式,惟有福音不改變.我相信這樣的傳福音,才能叫他們歸向耶穌.保羅說:「我向甚麼樣的人,就變成甚麼樣的人.」福音絕不妥協,十架真理並不減少,只是其他方面符合聽眾需要.
\newline
\newline
願主給我們這樣的智慧.
\newline
\newline
還有,保羅傳福音,無論得時不得時都要傳.傳福音不限於會堂,禮拜堂；保羅在河邊,在獄中,在亞基帕王面前,耶路撒冷,羅馬,他都傳福音,他對提摩太說:「得時不得時,要傳福音.」求主給我們有此迫切的心!有一天,保羅在腓立比獄中,看守的人問他:「我該怎樣得救呢?」相信定是保羅在獄中向他傳過福音的；雖然他身繞鎖鍊,頭痛流血,他照樣傳福音.
\newline
\newline
(二)使人作門徒,是保羅傳福音的目標,並非傳了福音就算盡了責任,他巴不得帶領每個人都作門徒.耶穌給門徒的使命:「所以你們去,使萬民作我的門徒.」怎麼作門徒呢?差傳的大使命有兩個動詞,「給他們施洗,施浸.」和「教導」凡主所吩咐的,要教導他們遵守.這兩個動詞則說,怎樣領人作門徒.「洗禮」有很深的預表,保羅在羅馬書說:「我們浸在水裡,好像埋葬死人,以往不認識神,因有罪與神隔絕；現在向罪死,從水中出來,如同復活,一舉一動有新生的樣式；現在活著的不再是我,是基督活在我裡面.」第一個動詞告訴我們,在基督耶穌裡得新生命.第二個動詞「教導」即作門徒要明白神的真理,不但明白,還要遵守.耶穌吩咐門徒教導他們遵守.要尊主為主.主教導門徒說:「不是每個喊主啊主啊的人都能進天國,唯獨遵行天父旨意的人,才能進去.」求主幫助我們!做耶穌基督的門徒,不但要信耶穌,接受福音,還要尊主為主.保羅之目的,不但讓人聽福音,還要帶領人作門徒.
\newline
\newline
(三)堅固門徒(22)這是初期教會沒有忽略的工作.保羅的榜樣是很寶貴的例子,他在安提阿,再到以哥念,後到路司得,在那裡被人用石頭打得半死；敵人以為他死了,把他拖到城外丟掉.感謝神!保羅又起來了,不但起來,次日還越山往特庇,他的家在大藪,離特庇不遠；他若南下,可抵家稍事休息；若北上,卻朝著石頭和敵人；可是北上也是教會,也是門徒,他們需要堅固.結果,他起身北上路司得,去堅固那些門徒,讓他們明白,作主門徒的須要經歷逼迫.保羅對提摩太說:「要預備心,像基督的精兵,要有一顆為主受苦的心.」他告訴聖徒們:「作門徒,作基督徒必須經過許多苦難,我也同樣受許多苦難.」
\newline
\newline
各位兄姊!我覺得保羅傳福音和我們不大一樣,我們常講福音的好處；保羅卻開門見山說:「信耶穌固然有好處,但必定有許多苦難會臨到你們.」第一,保羅幫助他們預備心.第二,面對著世界,世界的引誘要勝過；因為那些都市是強亂的都市,保羅深恐他們同流合污,勸他們站穩.因為社會污穢不堪,道德淪亡,必須站穩才能作中流砥柱.這也是今日基督徒必須學習的,在社會上工作,黑幕重重,常為錢財違背良心.保羅告訴他們,耶穌基督的信仰,是基督徒生活的準則；我們雖然生活在世上,卻不要效法世界；當靠耶穌基督,不但要明白祂的旨意,更要遵行祂的旨意.靠著聖靈,我們遵行主的旨意,過得勝的生活.當時保羅恐四周的異端邪教影響他們,故以逼迫,世界,信仰,三方面堅固那些門徒,勸他們恆守所信的道.
\newline
\newline
(四)在各教會中選立長老(25)教會的組織很要緊.保羅和巴拿巴在各教會組織起來,選派教會領袖.福音在那教會,只有短短的時間,他們選立了長老.多麼偉大!
\newline
\newline
求主給我們像保羅那樣的信心和智慧.
\newline
\newline
(五)藉著禱告將教會並弟兄姊妹交托主.主是莊稼的主,是他們信仰的對象.主自己說:「我要將我的教會建立在磐石上,陰間的權柄不能勝過他.」保羅把教會交托在主的手中,他給羅馬教會寫信說:「甚麼都不能使我們與神的愛隔絕.」他的結論是「甚麼都不能」.
\newline
\newline
求主幫助我們,像使徒時代的教會,學習差傳,學習教會的事奉.
\newline
\newline


\newpage

\section{第53屆~港九培靈會~戴紹曾~第七講}
\label{sec:605}
\textbf{差傳工作與家庭}
\newline
\newline
連結: \href{https://www.hkbibleconference.org/session-message/view/605}{\texttt{https://www.hkbibleconference.org/session-message/view/605}}
\newline
\newline
\hyperref[sec:604]{< < < PREV SERMON < < <}
~
\hyperref[sec:toc]{[返目錄]}
~
\hyperref[sec:608]{> > > NEXT SERMON > > >}
\newline
\newline
使徒行傳 18:1-18:4
\newline
\begin{longtable}{cl}
\hline
\hline
章節 & 經文 (和合本修訂版)\\
\hline
18:1 & \begin{tabularx}{0.7\textwidth}{X} 這些事以後，保羅離開雅典，來到哥林多。 \end{tabularx} \\ \\ \relax
18:2 & \begin{tabularx}{0.7\textwidth}{X} 他遇見一個生在本都的猶太人，名叫亞居拉。不久前，他帶著妻子百基拉從意大利來，因為克勞第命令所有的猶太人都離開羅馬。保羅去投靠他們。 \end{tabularx} \\ \\ \relax
18:3 & \begin{tabularx}{0.7\textwidth}{X} 他們本是製造帳棚為業。保羅因與他們同業，就和他們同住，一同做工。 \end{tabularx} \\ \\ \relax
18:4 & \begin{tabularx}{0.7\textwidth}{X} 每逢安息日，保羅在會堂裡辯論，勸導猶太人和希臘人。 \end{tabularx} \\ \\
[1ex]
\hline
\hline
\end{longtable}
聖經記載亞居拉和妻百基拉,遭遇了不幸的事,他們本在羅馬行商；因凱撒的命令,革老丟要驅逐猶太人離開羅馬；他們就來到哥林多.據歷史家說,當時猶太人中因有一人的名字(Christer)和基督的名相似,引致猶太人起了爭論,所以該撒不讓猶太人留在羅馬.尼羅王一向憎恨猶太人,尤其是基督徒.不過,所發生的事是在尼羅作王以先.那時基督的福音已傳入到羅馬,可能這對夫婦那時就信了.我也能接受這個說法,因為保羅每當題及他們時,始終沒題到是他帶他們信主的.至於其他由保羅帶領信主的人,保羅稱他們為兒子,或者題及怎樣帶他們信主；所以很可能這對夫婦未離開羅馬以先,已經接觸了福音,信了耶穌.他們被逐來到哥林多,工作須由頭做起,重建家園,困難重重,那段日子實在難渡!可是遭遇並無影響他們的信仰,他們沒有埋怨神.
\newline
\newline
各位青年!我們的神並沒有應許我們在一生的道路上,因信祂就都順利亨通.有時,我們難免有非凡的痛苦遭遇.保羅覺得身上有根刺,巴不得把它拔掉；可是欲脫不能,他呼求主,求主拔掉那刺.主對他說:「我的恩典夠你用的.」願我們抓住主的話,我們當學習這重要的功課.前面道路,無論你的家,學業,職業,可能遇風波,路途崎嶇險阻；看見別人亨通,以為神愛別人不愛你,以致放棄信仰,那你就錯了.請看這對夫婦的榜樣,遭遇逆境,仍然持守所信的道；在不違背信仰之下,他們聽從凱撒命令離開羅馬；然而對耶穌基督的愛,堅貞不渝.他們原來經商,遷徙異地,要重新建立市場,處處費盡心思精神,但是他們不讓工作影響信仰,他們明白「人活著不是單靠食物」的真理.固然要謀生,更要親近神.
\newline
\newline
現時代的青年基督徒,出了學校,踏入社會,立家立業,壓力重重,常常連禮拜的時間都沒有,沒有時間讀經禱告,也沒時間事奉主.這樣是錯誤的；應該讓信仰來影響我們的家庭和工作才對.
\newline
\newline
當時哥林多是個繁榮的商業地區,東西各有港口,有如香港.那對夫婦來到哥林多,非但不讓工作影響信仰,反而因信仰改變他們的工作.他們尊主為主.求主幫助青年們,特別是近年進入社會的青年,你掙扎,立業,思想,千萬別把時間用盡；信仰,教會,事奉,甚至你的主,你是否都不要了.這是多麼可憐!保羅有位同工底馬,愛現今的世界,離開主,離開保羅去了.
\newline
\newline
哥林多是個腐敗之地,大致上各處港口都是如此.數年前我到荷蘭阿姆斯特丹,那地實在很美!但另一面是個很骯髒腐化之地.哥林多就是這樣.古時如果有人說,某人是哥林多化,即表示他犯了淫亂的罪；若說,那人是個哥林多朋友,乃指那人是妓女.地名和淫亂的罪有關連,可見腐敗地步至此.
\newline
\newline
保羅在(林前六9-11):「......不義的人不能承受神的國......無論是淫亂的,拜偶像......都不能承受神的國.你們中間也有人從前是這樣；但如今你們奉主耶穌基督的名,並藉著我們神的靈,已經洗淨,成聖稱義了.」
\newline
\newline
感謝主!亞居拉夫婦和保羅,把耶穌基督的真理帶到腐敗的哥林多,那地的人就改變了.亞居拉夫婦,靠著主,不同流合污,顯出聖潔的,新的見證.哥林多教會,後來有些人失敗了；(五1-2,六15有記載)可是這個家庭在腐敗黑暗環境中,靠主聖潔.
\newline
\newline
當保羅從雅典下到哥林多,有很大屬靈的爭戰.向那些哲學家,猶太人,傳講福音不是容易的,林前二章3節,保羅描寫自己:「我在你們那裡,又懼怕,又甚戰競.」可能那時保羅的病常發作,他到了哥林多；神預備了亞居拉夫婦的家,他們把門打開,保羅便進到他們的家.
\newline
\newline
多年前,我在台灣作山地工作,那地從無人傳福音；初抵該地,不知住何處；根本沒有基督徒接待,也沒有旅店.當晚是週末,我住在派出所,很多小朋友把我腰帶周圍咬的遍滿紅斑,非常痕癢.當我第二次重返山地傳福音,情形大不相同；基督徒把門打開,接待我,表示他們信了主,很是高興.
\newline
\newline
保羅到了哥林多,有個敬畏神的家庭,歡迎他,同情他,和他同工,並帶給他很大的幫助.「如有人接待你,就住在他的家.」我想保羅根據主教導門徒的真理,就在亞居拉夫婦的家住下了.他在哥林多十八個月,享受他們的接待,實屬難得!據路加說,保羅投奔他們,還另有原因,他們之間本系同行,保羅小時曾學過織造帳棚的技術；亞居拉夫婦正是作此行業.
\newline
\newline
話至此,我想到神在每個人身上有祂永恆的計劃.各位青年!當你在學時學習的功課,你不曉得,將來長大踏入社會,如何應用.
\newline
\newline
保羅在神永恆的計劃中,小時所學的正與傳福音事工互相配搭,他一邊親手作工織造帳棚,同時也傳福音,神的計劃奇妙!
\newline
\newline
(第十八章18-21節)記載亞居拉夫婦因逼迫離開羅馬,來到哥林多,此次離開哥林多他往,是為了福音工作與保羅配搭；保羅要離開哥林多到以弗所,他們認為須與保羅同去.保羅卻把他們留在那裡,他們幫助保羅立將來工作的基礎.以弗所是保羅早就要去的地方,他第二次出外佈道,心向以弗所.神的靈不許,(我曾題過保羅工作的計劃是有戰略性的都市)因神的時候未到,要他直往特羅亞；後來他看見馬其頓異象,先進入歐洲!在腓立比一帶傳福音；但是保羅心裡嚮往以弗所.現在帶了亞居拉夫婦到以弗所,留下他們先作初步工作.他們不讓工作影響信仰,也不讓工作影響差傳的事奉.相反地,他們將工作和差傳事奉配搭,他們手作的工,符合差傳工作.這是多美的見證!早期教會,他們的心在福音工作上,在差傳工作上逼切地步至此!為了福音的緣故,離開哥林多,就到以弗所.
\newline
\newline
以弗所的工作展開了(十八24-28)保羅早就願意進到這個大都市；可是他來了遂即離開,他有一種敏銳的心,他明白神的旨意.我們看見一個對比,在腓立比時,保羅被下監,半夜地大震動,門打開了,保羅可以離開；但是他不離開；結果,帶領了一些人歸主.此次來到以弗所,門開了,猶太人歡迎他；他把同工留下來,自己卻回耶路撒冷.
\newline
\newline
求主給我們,個人,家庭,教會,有敏銳的心,無論去留,遵主命令.
\newline
\newline
有個猶太人名叫亞波羅,他生長在北非的亞歷山大,來到以弗所；他是有口才有學問的人,最能講解聖經.這人已經在主的道上受了教訓,心裡火熱；可是在真道上還不完全,所明白的不夠；他只曉得約翰的洗禮,對基督的救贖,復活的主,聖靈的工作,都不明白.他講解只是一部份,帶領人只能帶到一個地步；這樣有恩賜的人,當時在以弗所工作.亞居拉夫婦到會堂聽亞波羅講道,他們極其留心聽道；不像一般人,只注意講員的口才,內容的興趣；他們感到亞波羅所講的道美中不足；若能講到耶穌的復活,耶穌所應許的聖靈,那就更寶貴；更加造就多人,使人得到更大的祝福和幫助.教會的工作,差傳的工作,福音任何方面的工作都很重要.亞居拉夫婦在聽道不滿意之情形下,並非起而公開指摘,若然,不但分裂會堂,且不能幫助保羅立下好基礎,反而造成破壞.他們有智慧,不單懂得勝經的真理,更懂得勝靈的工作,聖靈在他們裡面.有些人為了一次托負信徒的真道,竭力爭辯；到了一個地步,他們的態度,作法,違背了所爭辯的真理.亞波羅不明聖靈的工作,有聖靈在他們身上的夫婦,他們肯接待亞波羅到他們的家,這是基督徒的交通,基督徒的愛；他們將神的道給他講解更加詳細.說耶穌接受約翰的洗禮,約翰為耶穌作先鋒；耶穌死了復活了,向祂的門徒顯現；對他們說:「父所應許的,我都要給你們,我升到父那裡,要將聖靈賜給你們,在聖靈能力之下,把福音帶到地極.」亞居拉夫婦如此教導這位有才學的亞波羅,不以自己才疏學淺以織造帳棚為生的小民而膽怯,不敢啟齒.這是聖靈的工作,他們以謙虛的心幫助亞波羅.感謝主!
\newline
\newline
親愛的青年!我們聽道,當有辨別的靈,不要吹毛求疵,任意批評.
\newline
\newline
亞居拉和百基拉的心,實在很美!
\newline
\newline
聖靈說:「要用愛心說誠實話.」主告訴我們,如果你知道有錯.人家得罪你,你該去找那人.在教會中,我們常違背耶穌的命令.聖經明文,你若不去做,就變成雅各書三章所說的「有毒的舌頭」.舌頭是全身最可怕的肢體,應當榮神益人才對；可是基督徒常用舌頭敗壞人.哥林多教會也是這樣,他們很會說方言,但是他們敗壞人.
\newline
\newline
保羅的同工不是這樣,該在人面前講的就講,該在家講的就在家講；結果,整個教會得幫助.以弗所沒有分別.到了哥林多,亞波羅幫助這許多那蒙恩信主的人,在眾人面前極有能力,駁倒猶太人；無論主內的人,或是猶太人,亞波羅都有貢獻,有幫助.這是多美的見證!過了些時,保羅到以弗所開始工作,就在亞居拉和百基拉的小小基礎上傳福音.保羅再次和他們一起配搭.這幾個晚上,我們將看見以弗所如何蒙神拯救；這對夫婦在差傳工作上擔當了重要的角色,他們離開家,到了哥林多.他們得勝,他們的家成了基督化的家庭.保羅題及他們,不但說是他的同工,還常說,問候在他們家的教會.可見在哥林多,在以弗所,在羅馬,他們的家的就是家庭的教會,他們的家是傳福音的家,是差傳的家.我們看見他們的同心,愛心,熱心.求主幫助我們,在福音工作上,能有更多青年,更多年青夫婦和中年夫婦,效法亞居拉夫婦的榜樣.相信福音在香港,在東南亞,在全世界將越發傳開.
\newline
\newline


\newpage

\section{第53屆~港九培靈會~戴紹曾~第八講}
\label{sec:608}
\textbf{差傳工作的原動力}
\newline
\newline
連結: \href{https://www.hkbibleconference.org/session-message/view/608}{\texttt{https://www.hkbibleconference.org/session-message/view/608}}
\newline
\newline
\hyperref[sec:605]{< < < PREV SERMON < < <}
~
\hyperref[sec:toc]{[返目錄]}
~
\hyperref[sec:610]{> > > NEXT SERMON > > >}
\newline
\newline
腓立比書 3:7-3:16
\newline
\begin{longtable}{cl}
\hline
\hline
章節 & 經文 (和合本修訂版)\\
\hline
3:7 & \begin{tabularx}{0.7\textwidth}{X} 只是我先前以為對我是有益的，我現在因基督的緣故而當作是有損的。 \end{tabularx} \\ \\ \relax
3:8 & \begin{tabularx}{0.7\textwidth}{X} 不但如此，我已把萬事當作是有損的，因我以認識我主基督耶穌為至寶。我為他已經丟棄萬事，看作糞土，為要贏得基督， \end{tabularx} \\ \\ \relax
3:9 & \begin{tabularx}{0.7\textwidth}{X} 並且得以在他裡面，不是有自己因律法而得的義，而是有信基督的義，就是基於信，從神而來的義， \end{tabularx} \\ \\ \relax
3:10 & \begin{tabularx}{0.7\textwidth}{X} 使我認識基督，知道他復活的大能，並且知道和他一同受苦，效法他的死， \end{tabularx} \\ \\ \relax
3:11 & \begin{tabularx}{0.7\textwidth}{X} 或許我也得以從死人中復活。 \end{tabularx} \\ \\ \relax
3:12 & \begin{tabularx}{0.7\textwidth}{X} 這不是說我已經得著了，已經完全了；而是竭力追求，或許可以得著基督耶穌所要我得著的。 \end{tabularx} \\ \\ \relax
3:13 & \begin{tabularx}{0.7\textwidth}{X} 弟兄們，我不是以為自己已經得著了；我只有一件事，就是忘記背後，努力面前的， \end{tabularx} \\ \\ \relax
3:14 & \begin{tabularx}{0.7\textwidth}{X} 向著標竿直跑，要得神在基督耶穌裡從上面召我來得的獎賞。 \end{tabularx} \\ \\ \relax
3:15 & \begin{tabularx}{0.7\textwidth}{X} 所以，我們中間凡是成熟的人，總要存這樣的心；若在甚麼事上存別樣的心，神也會把這些事指示你們。 \end{tabularx} \\ \\ \relax
3:16 & \begin{tabularx}{0.7\textwidth}{X} 然而，我們達到甚麼地步，就當照這個地步行。 \end{tabularx} \\ \\
[1ex]
\hline
\hline
\end{longtable}
耶穌基督復活,帶給初期教會的聖徒,在福音工作上有莫大的力量.
\newline
\newline
保羅說:「因信神而來的義,使我認識基督,曉得祂復活的大能；」(腓三9-10)還有行傳一章3節「祂受害之後,用許多憑據,將自己活活的顯給使徒看,四十天之久向他們顯現,講說神國的事.」主耶穌復活之後向祂的門徒顯現.(一21-22)在一章8節,主已經說:「你們要作我的見證.」使徒們聚在一起,要選一位和他們一同作主復活的見證.
\newline
\newline
主耶穌復活,成為初期教會很大的力量.(三14-15,四33)首先,我們看,在耶路撒冷有個空的墳墓；這是主的門徒,和主的敵人都承認的.關於這個空置的墳墓,解釋各有不同.按邏輯說,若非人把主耶穌的屍體帶走,就是照著主所應許的,祂從死裡復活了.若是人把主的屍體帶走,有兩種可能性；一是敵人所為,一是門徒所為.敵人能不能,會不會這樣做呢?彼拉多派兵丁嚴密防守,敵人根本無法下手；另方面,當時敵人告訴彼拉多:「這迷惑人的,在世時曾說死了要復活；為避免門徒把屍體帶走,必須嚴密看守.」當門徒見證主復活時,敵人保持緘默,這是很重要的證據；如果敵人知道主的屍體在何處,他們必出而證明耶穌復活是假的.然而他們除了禁止耶穌的門徒,不可談論耶穌復活的事之外,他們無話可說；因此足以證明絕非敵人把主的屍體帶走.至於門徒,若要帶走主的屍體,也全無可能；因為兵丁守衛森嚴,門徒們也不甚相信主說:「必復活」的話,門徒幾乎無此盼望,更沒有可能突然改變常態,製作,捏造虛假的道理.四福音記載,耶穌死了,門徒充滿憂愁絕望.可是到了使徒行傳時,他們開始講主的復活,他們心態完全轉變；這種喜樂心情的改變,是不可能寄於「謊言」的.還有,數位門徒以往不大同心,彼得和約翰,也曾爭論誰為大,誰坐在耶穌的左右邊；如果他們頓時合作做假事,經年累月,總有一天會露出真相的；難道門徒願為句謊言而死亡嗎?相信絕不.那時他們傳道以耶穌復活的事實為信息的中心,甚至他們的生活,猶太人的風俗習慣也改變了.猶太人一向頑固,可是為了記念耶穌復活,他們不守舊時的安息日；而守主耶穌復活的聖日,實在是很大的改變!他們在主復活之日,聚集講道,收奉獻.這樣大的改變,不可能基於一句謊言一件假事；門徒們深信耶穌從死裡復活了.所以,空的墳墓,不是敵人把耶穌的屍體帶走,更不是耶穌的門徒所為；乃是照主自己所應許的應驗.
\newline
\newline
主復活的事實,帶給門徒新的喜樂,新的力量；成為傳福音的原動力.復活的主與他們同在,他們勇敢外出見證耶穌基督.復活的重要性,在行傳中多處出現.
\newline
\newline
「復活」或謂「主從死裡被興起,」在使徒行傳中至少題了廿八次,平均每章一次.第一章題及主耶穌升天,門徒聚集一起,按照主的命令,同心合意禱告.他們原有十二人,現僅十一人；他們感到十二的數字意義很深,舊約時代有十二個支派；耶穌立了十二個門徒.有人批評他們選了一位門徒代替猶大是不對的.無論如何,聖經並不題他們不對.雖然保羅後起作使徒；可是他和耶穌所揀選的使徒情形不同.保羅的確是位帶著使命的使徒.當時門徒們認為必須立一位門徒.條件「要有好名聲,被聖靈充滿,有智慧,有信心.」在第一章隻字不題,惟獨題「必須......立一位與我們同作耶穌復活的見證.」(22)當然對主的生平:從主受約翰的洗,直到祂升天,都必須瞭解.可是彼得題出一件非常重要的事情:要能見證主的復活；一定是主曾向他顯現的人.由此可見主復活的重要性.同時,我們也看見這些門徒的改變.
\newline
\newline
保羅在大馬色路上,因復活的主向他顯現,有了徹底的轉變；並非有人說服他,不是思想的改變；或是他自己突然作了新的決定,乃是復活的主向他顯現.
\newline
\newline
彼得,司提反,保羅,及其他使徒的講道,「主復活」成為他們信息的中心,無論向猶太人或非猶太人；無論在耶路撒冷,帖撒羅尼迦,或以前在雅典,都講復活的信息；講時引起人的反對,尤其是撒都該人不信復活,起來反對,他們討厭基督徒,不歡喜他們傳講主復活,蓄意加害,造成教會初次遭受逼迫；撒都該人禁止他們不得傳講主復活的道理.當保羅工作將結束時；有個官名非斯都,他清楚看見產生衝突的原因,他對亞基帕王說:「當時猶太人中有基督徒傳講耶穌復活的信息；可是其他的猶太人否認這事,於是同胞之間發生衝突.」至於外邦人,當保羅到了雅典,不以福音為恥,向知識份子傳福音,清楚對他們說:「耶穌基督死了復活.」但是,人譏笑他,不信此事,而保羅卻仍然傳講復活的信息.這個信息成為他們傳福音重要的一環.
\newline
\newline
復活的意義,對我們或是對當時的聽眾來講,其意義有四:
\newline
\newline
(一)主的復活是個證據　主耶穌復活,特別對猶太人成了重要的見證,這位復活的耶穌是舊約預言的彌賽亞.他們對聽眾說:「這事足以證明祂是我們所等候的彌賽亞；因為歷史上從未有這樣的事.主也曾叫此人復活；但是主的復活不一樣,祂乃是照著自己的應許「三天以後,我要復活」「主果然復活了.」
\newline
\newline
彼得講道時,曾經數次對聽眾說:「請注意!你們對主耶穌的看法,與神對耶穌的看法完全不一樣；你們棄絕祂,將祂處死；神卻叫祂從死裡復活了.」彼得用舊約的經文說:「眾人所棄的那塊石頭,成了大建築物重要的石頭.你們厭棄祂,拒絕祂；神卻叫祂死裡復活.」彼得根據主的復活,指他們看法錯了；強調此事足以證明耶穌就是我們所等待的彌賽亞.保羅在彼西底的安提阿也如此說,在歐洲帖撒羅尼迦的會堂也如此說.
\newline
\newline
(二)主復活使信者得救贖　使徒們講道時,特別強調,因著耶穌復活,勝過死亡,勝過魔鬼,更勝過罪惡；所以我們可以脫離罪惡.在三章26節彼得講,我們脫離罪,和主的復活連在一起,復活的主,藉祂復活的大能帶給我們救贖,遠超過舊約；正如希伯來書作者所題一樣.保羅說,本來律法不能使我們脫離的；靠著復活的耶穌我們得以釋放.保羅在羅馬書,加拉太書題的更清楚,靠自己遵守律法,在多方面我們實無能為力；他說:「我裡面有個律,本來立志行善由得我,行出來卻由不得我.」他說出自己體驗的痛苦.「過去我不認識耶穌,不知道有位復活的主,不知道主復活的大能.我真苦呀!我想行善；但行出來卻由不得我.」
\newline
\newline
有齣舞台劇名「第三害」敘述中國歷史周處的故事.當時,地方有三害,山上有老虎,海中的蛟龍；第三害是周處本身,他卻不自知.有一天,周處奮勇上山入虎穴,經搏鬥,滅老虎；得意洋洋回家,以為鄉民必稱讚他.過些日子,他又戰勝蛟龍,除去第二害；得意之餘,他自言自語:「既對付了老虎,又戰勝蛟龍,卻無改變自己.」他和保羅同感「老我最難對付.」
\newline
\newline
保羅因著復活的主,不但有力立志行善；且有力量行出來.他說:「若有人在基督裡,他就是新造的人；舊事已過,都變成新的了.」
\newline
\newline
親愛的兄姊!福音的本身,福音的中心在於耶穌基督的復活；主復活滿足了神的心意,使我們與神和好,帶給我們人所未有的能力,使我們得勝.「惟獨神所復活的,祂並未見朽壞.......赦罪的道是由這人傳給你們的.你們靠摩西的律法,在一切不得稱義的事上,信靠這人,就都得稱義了.」(十三37-39)
\newline
\newline
(三)主耶穌復活證明祂是主　彼得清楚對他的同胞說:「故此,以色列全家當確實的知道,你們釘在十字架上的這位耶穌,神已經立祂為主為基督了.」(二36)保羅在大馬色路上遇見耶穌,他說:「主啊!你是誰?」主說:「我是耶穌,是你所逼迫的耶穌.」這時,保羅明白耶穌復活了；他聽見主的聲音,看見主的光,得著主的使命；他稱祂為主,對主說:「主啊!要我為你作甚麼?」主對保羅說:「我要你作我的器皿,為我的名要受許多苦難；你要為我作見證,要站立在君王面前.」保羅認識復活的主,他接受主,且尊主為主.
\newline
\newline
我再次題醒弟兄姊妹!不是單口裡說主啊!主啊,這樣是不夠的,心裡要尊主為主.主自己告訴祂的門徒,許多人口裡說:「主啊!我們傳道,醫病,趕鬼,見證.」可是主要對他們說:「離開我去吧!」口裡說主,工作表面上似乎順服主；但事實上沒尊主為主.主的復活,證明祂是主.我們是否在日常生活中尊祂為主呢?
\newline
\newline
(四)主的復活要叫信祂的人復活　(十七30-31)耶穌從死裡復活給萬人作可信的憑據.「憑據」就是讓我們知道有一天這位復活主要審判活人死人.保羅給提摩太寫信說:「我在神面前,並在將來審判活人死人的基督耶穌面前；憑著祂的顯現和祂的國度囑咐你,務要傳道!」(提後四1-2)主必再來,那時這位復活的主要審判活人死人.保羅說:「因這緣故,我勉勵自己,提醒自己,警告自己,在日常生活中,應該討祂的喜悅,言行一致,榮耀神.」保羅知道因為有一天要站在神的審判台前.
\newline
\newline
耶穌基督復活的意義,證明祂是舊約聖經預言的彌賽亞,證明祂是救主；天下人間沒有賜下別的名,可以靠著得救.祂要在你我心裡作主作王.主的復活,告訴我們將來世界末日,主必再來；在審判天下的時候,我們要站在祂的面前.求主幫助我們,對主復活有更準確的認識.
\newline
\newline
我擔心我們沒有把復活的信息放在心裡,深恐我們沒有重視主的復活；不像使徒時代,他們重視主的復活,相信主的復活；因主復活,他們得著能力,見證主的救恩,他們體驗主復活的大能.請問!在我們生活中,有沒有這樣的經驗?復活的主有沒有赦免我們的罪?赦免你的罪?你有沒有將這信息傳給別人,讓別人享受這位復活之主的恩典呢?使徒時代的信徒已明白主復活的重要性了.
\newline
\newline
求主幫助華人教會,幫助我們在香港見證主從死裡復活,作世人的救主.
\newline
\newline


\newpage

\section{第53屆~港九培靈會~戴紹曾~第九講}
\label{sec:610}
\textbf{差傳工作的訓練和裝備}
\newline
\newline
連結: \href{https://www.hkbibleconference.org/session-message/view/610}{\texttt{https://www.hkbibleconference.org/session-message/view/610}}
\newline
\newline
\hyperref[sec:608]{< < < PREV SERMON < < <}
~
\hyperref[sec:toc]{[返目錄]}
~
\hyperref[sec:611]{> > > NEXT SERMON > > >}
\newline
\newline
使徒行傳 16:1-16:5
\newline
\begin{longtable}{cl}
\hline
\hline
章節 & 經文 (和合本修訂版)\\
\hline
16:1 & \begin{tabularx}{0.7\textwidth}{X} 後來，保羅來到特庇，又到路司得。在那裡有一個門徒，名叫提摩太，是信主的猶太婦人的兒子，他父親卻是希臘人。 \end{tabularx} \\ \\ \relax
16:2 & \begin{tabularx}{0.7\textwidth}{X} 路司得和以哥念的弟兄都稱讚他。 \end{tabularx} \\ \\ \relax
16:3 & \begin{tabularx}{0.7\textwidth}{X} 保羅要帶他同去，只因那些地方的猶太人都知道他父親是希臘人，就給他行了割禮。 \end{tabularx} \\ \\ \relax
16:4 & \begin{tabularx}{0.7\textwidth}{X} 他們經過各城，把耶路撒冷使徒和長老所決定的規條交給門徒遵守。 \end{tabularx} \\ \\ \relax
16:5 & \begin{tabularx}{0.7\textwidth}{X} 於是眾教會信心越發堅固，人數天天增加。 \end{tabularx} \\ \\
[1ex]
\hline
\hline
\end{longtable}
過去八晚,我們看見在座許多青年,都渴慕神的話.想到華人教會的未來,我們若真正肯把自己擺上,這是很美的現象.相信我們中間有不少的青年,有心事奉主；不但帶職事奉主,而且有心一生全時間事奉主.
\newline
\newline
要事奉主,必需要訓練和裝備.凡事奉主的人和作為一個屬靈的領袖,有很多因素,不是與生俱來的.在神永恆的計劃中,無論家庭,教會,神學院,在事奉主者身上,有很大的影響,也有很大的貢獻.昔日神所重用的僕人摩西,家對他有很大的影響.撒母耳母親的虔誠,使他從小懂得勝經.保羅過去的背景,他所讀的大學,或是在迦瑪列的手下,都對他以後的事奉有很大的影響.教會歷史所載的馬丁路德,加爾文,約翰衛斯理；這些人都不是一下子登峰造極；他們得學位,有的因家庭背景,有的因教會背景,有的是在神學院的背景.以中國教會的宋尚節來說,他生在愛主的家庭,他父親是位牧師,終身事奉主.王明道先生也是生於基督徒的家庭；雖然他的遭遇相當不幸,沒見過他自己的父親；但虔誠的母親,對他影響殊深.
\newline
\newline
使徒行傳中論及訓練和裝備,最明顯的是提摩太的見証.在四福音中,我們看見耶穌怎樣帶領門徒.
\newline
\newline
我們萬勿以為一個人,瞬息間便可成為神所重用的人；幾乎每個青年心裡有此矛盾.當我們矚目看田,看見教會的荒涼,同胞的需要；許多人失喪,許多人不認識耶穌；基督的愛激勵我們,我們得先會對此關懷.以兄弟的經驗,當我讀大學時,暑期就去工作,有時做青年學生工作；有時他們問我:「為甚麼你還要讀大學讀神學呢?現在的需要太大,你快出來傳福音吧!」我聽了有點動心,覺得他們言之有理.可是,感謝主!祂給我理智,願意多學習,多明白聖經；多明白教會歷史,多懂得神學,多瞭解如何事奉主.當時實在有種矛盾的現象「趕快進入工場作工.」感謝主!祂攔阻我,叫我應多下工夫,接受訓練與裝備.
\newline
\newline
各位青年!在福音工作上,如果主給你有機會念大學,讀神學,你要抓住機會,善用機會；把你最好的時間,全心擺上.
\newline
\newline
請看!提摩太的家成為他第一環訓練的地方.保羅寫信給提摩太提及說:「你無偽的信」這是提摩太的見証；他純正的信心,無偽的信是先在他的外祖母和母親心裡的；外祖母影響外孫,母親影響兒子；在家庭中生活與信仰沒有脫節.
\newline
\newline
在座為父母的,請注意!並非教會才是家,若是把責任推給教會,不可能幫助孩子.要使孩子生長在基督化的家庭,站穩在真道上,接受訓練和裝備.
\newline
\newline
提摩太在家中,不但有榜樣,還有教導；保羅對提摩太說:「並且知道你是從小明白聖經.」雖然他的父親是希利尼人；但這位敬畏主的母親,教導他從小學習聖經.
\newline
\newline
記得當我十幾歲時,在中國經過了五年之久與父母分開；到了抗戰結束,我們出了集中營,到陝西與父母相聚.父親說他要到河南,問我和他同去嗎?當時我感到十分榮幸；因為父親要帶我同去他的工作崗位.過了段時間,相信是我青年時期最大的收穫；父親教導我讀聖經.他本身沒機會上神學；可是他讀的是教會大學,有機會念聖經；加以自己肯在神的話語上用功.我們從陝西到河南,沿途中有時搭火車,有時搭卡車,有時坐馬車；途中父親教導我如何讀經,特別教我歸納式查經方法.
\newline
\newline
未知提摩太的母親如何教導兒子讀經,不過,我相信她對提摩太日後的事奉,他所明白的舊約,有很大貢獻.
\newline
\newline
我再次說,你們為父母的,負有重大責任,在家中要教導兒女明白聖經.
\newline
\newline
十七八世紀,衛斯理約翰的家庭,對英國有很大的貢獻.直到今天,我們還唱著衛斯理的詩歌,無論是中英文本都極受歡迎.歷史學家說:「英國不像法國有流血的內戰,是神使用衛斯理約翰；無論在國教,在社會中,帶來福音大復興.」衛斯理的母親有十九個孩子,這位偉大的母親,沒把孩子們送托兒所,自己卻開辦托兒所.記得有一次在美國夏天佈道大會,佈道家講道,有位母親大受感動,會畢趨前對講員說她要奉獻傳福音；傳道人看出她是位母親,問她有幾個孩子,她說有十幾個；傳道人說；「感謝主!主不但呼召你傳福音,並且連聽眾都給你了；你可以好好地在家裡傳福音.」衛斯理的母親大概就是如此,她把孩子個別分開；每週輪流,各個孩子都有她自己帶領的機會.據衛斯理的日記,他最喜歡週四；因為這天母親自己照顧他,個別談話,教導他明白聖經.
\newline
\newline
做母親的!你沒有那麼多孩子,較不吃力；可是你肯否在你三兩個孩子身上下工夫呢?在家的訓練和裝備非常重要.
\newline
\newline
第二個訓練的地方是在教會.當時提摩太生長之處還沒有教會,他母親帶他上會堂；某天,傳道人保羅和巴拿巴來傳福音,提摩太得救了.保羅稱提摩太為「我兒」因他帶領提摩太認識耶穌.雖然提摩太的母親虔誠,但他需要重生的經歷.可能在座中有不少青年,生長在基督化家庭；但是,你知道耶穌基督是你個人的救主嗎?父母愛主虔誠熱心事奉主；然而我們不能依靠父母的信仰和事奉,我們必須自己有重生的經驗.
\newline
\newline
當提摩太看見保羅,又看見敵人來了,用石頭幾乎把保羅打死；相信此事給提摩太很大的考驗,難道信耶穌要遭遇這樣的逼迫和考驗嗎?然而提摩太並沒有搖動,沒有離開主.他看見保羅把教會組織起來,選立了執事,然後保羅離開.提摩太仍然持守所信的道,因他所信的不是保羅.有時我們錯了,把注意力放在牧師或傳道人身上；可能我們很愛他,因他幫助我們；可是一旦他不在我們中間,我們的心便冷淡,離開了主.第十六章告訴我們,「保羅回到路司得,在那裡有一個門徒,名叫提摩太,」這証明保羅離開一兩年之久,提摩太仍然站穩,他是個門徒,第二節說,他是那帶地方的弟兄所稱讚的門徒；即是說他在福音工作上,在基督的信仰上,有好的見証.到了第六章,教會需要新血,使徒人手不足,他們要選七位有好名聲的；我想路加在這裡題了「提摩太是弟兄所稱讚的門徒」就是一樣的意思.有好名聲,生活有好見証；不但有純正信仰,且雖身處拜偶像淫亂之地；而生活則有聖潔的,得勝的見証,這行一致,信仰和生活沒有脫節.
\newline
\newline
或者有人要說,保羅豈不是寫信給加拉太教會說:「你們不是靠行律法,靠受割禮得救.」為何保羅為提摩太行割禮呢?提摩太不是已得救了是個門徒了嗎?我們若細讀保羅在別處所講的,並無衝突.這位青年提摩太的父親是希利尼人,母親是猶太人；這是當地一帶的猶太人所共知,為了福音的緣故,為了要在各地有影響,必須這樣做；不是為了得救,不是靠這成為基督徒；保羅這樣做,乃是為了除掉傳福音的攔阻.弟兄姊妹!我們再次看見保羅如何重視福音,他能適應各地不同的情形；為的是高舉基督,為了福音益發廣傳；不是妥協,乃明智之舉,叫福音不受莫須有的攔阻.
\newline
\newline
提摩太受了眾長老的按手,保羅也為他按手；分派他出去,成為早期的宣教士.保羅常把按手和屬靈恩賜,連在一起.
\newline
\newline
早期教會,非常重視聖靈所賜屬靈的恩賜；今日則較少人注意及此,不過常重視有沒有受神學訓練.初期教會不只注重說方言,而是注重講道；說預言,教導人,幫助人等等各樣不同的恩賜,根據這些恩賜以鼓勵青年.
\newline
\newline
各位兄姊!主給你的恩賜是甚麼?你有沒有發掘,鍜鍊,運用?我們應當教導的該教導,見証主,幫助人；發掘自己的恩賜,忠心運用,榮耀主.
\newline
\newline
提摩太在教會中,保羅題醒他要好好運用主所給的恩賜,挑旺起來如火一樣.
\newline
\newline
接受訓練和裝備,在使徒行傳中,不只給我們看見提摩太在家中,在教會裡有這樣的經歷,這樣的栽培；神也興起一位屬靈的領袖,將他帶起來,幫助他再受裝備.
\newline
\newline
今天教會需要新血,但是我們感到太缺少；是否因今天的教會,太多傳道人牧師沒有足夠的時間去關心青年,去發掘人才,培養人才?多少時候,我們忙忙碌碌,以致下一代的青年缺少這樣的訓練和裝備.
\newline
\newline
保羅對提摩太說:「你要注意教會那些忠心能教導別人的人.」保羅到了路司得,「想要帶提摩太同去」這句話很寶貴!他看中這位青年,看他可以造就；要帶領他.可是,這位青年未嘗沒有缺點,他的恩賜還沒有有好好接受鍜鍊,他也膽怯!一個屬靈的領袖保羅,發現提摩太有潛力,是可造之材；他深信靠著聖靈；這位青年雖然軟弱缺少信心,沒有經驗；可能在講道操練上不夠,可是保羅要帶領他,期待有日可以把工作責任交給他.
\newline
\newline
新愛的弟兄姊妹!特別是我們中間的傳道同工；求主給我們這樣的信心,這樣的看見,肯提拔青年,栽培造就人才.
\newline
\newline
保羅分為四方面進行訓練提摩太:
\newline
\newline
(一)以身作則　(提後三10)在原文,也許你要覺得更囉囌「你服從了我的教訓,我的品行,我的志向,我的寬容,我的愛心,我的忍耐.」保羅特別強調自己的榜樣,藉以影響提摩太.
\newline
\newline
一個屬靈的領袖,若要影響下一代的領袖；不是口傳而是身傳,必須以身作則.
\newline
\newline
(二)用功教導　(提後一13,二2,三14)論保羅有條有理地教導提摩太.下工夫個別教導,對以後福音的發展廣傳,教會的增長,有相當大的戰略性.雖然有一天保羅要離開世界；但他已訓練了接班的人了.
\newline
\newline
(三)保羅事奉主的榜樣　保羅講道,辯論,勸勉猶太人,親手作工,禱告,組織教會的時候；提摩太都在他身邊,有機會看見一切榜樣.
\newline
\newline
感謝主!廿五年前我開始在差傳工作事奉主；主使我很大的蒙福,讓我有機會跟父母學習.我敢說,從父母所學的和我在大學,神學一樣寶貴,並不少於神學院所學的.我有機會看見他們如何傳福音,個別傳道；如何在神學院教書,如何作文字工作.牧師們若肯把這方面教導神學生,他們一定會很大的蒙福,兄弟多年神學教育,對於神學院派學生到教會實習,深感難過；因為教會的人太忙,不能好好的教導神學生如何事奉.或許叫他掃地,使他學習謙卑本是不錯的；可是應當教他掃地如何掃法.此外,開會,也應該教導怎樣開會.傳道人有沒有時間教導呢?我個人感覺,神學教育很大的失敗,是在於時間上,或在實習工作上.我們在神學院,已經學了很多理論,到了週末,要到教會；可是教會和神學院配搭不夠好.巴不得各神學院派學生到教會,牧師能把他看作你的提摩太,個別帶領他；你雖很忙,週末工作很重；你若真像保羅一樣地帶領提摩太,神學院和教會就不會有脫節；這樣的合作,對華人教會是最大的貢獻.
\newline
\newline
(四)保羅肯把責任交托提摩太\b0 　有時保羅未免擔心這青年能否勝任；但終於交托給他,他也接受了.某次保羅身在雅典又在哥林多,很想知道帖撒羅尼迦信徒的信心情況；就打發提摩太去堅固那邊信徒的信心,他回來時,使保羅非常高興.有時保羅打發提摩太到哥林多,這教會是很難搞的教會；分門結黨,勾心鬥角.保羅告訴教會說,「年青的提摩太膽怯,不要驚嚇他；但是他愛主,你們要接受他.」當保羅下監時,打發提摩太到以弗所,交給他重任,去帶領那些長老守住信仰；提防異端,保護純正的信仰；教導聖徒,身負重任.這都是提摩太學習的機會.
\newline
\newline
經歷了這些困難,結果產生了一個接棒的青年,造就了一個能把福音帶給第二代的青年.保羅自己見証說:「提摩太待我像兒子待父親一樣」提摩太敬保羅,像兒子孝順父親一樣.這是多美的榜樣!
\newline
\newline
各位青年!我們當學習提摩太的順服.有時他不明白,然而他總是順服.
\newline
\newline
保羅非但稱提摩太「我兒」也稱他「同工,」兩代全無代溝；一起配搭,一起事奉主.保羅說:「沒有別人,心和我一樣像提摩太.」這是因為保羅肯下工夫在提摩太身上.保羅數次說:「提摩太勞力作工.」其實,保羅自己擺出見証,他費財費力要傳福音.提摩太的心和他老師一樣,至死忠心；當保羅初次下監,他不以福音為恥,不以保羅的鎖鍊為恥.保羅第二次在監裡,給提摩太寫信叫他速來,提摩太即來,他如此忠心!希伯來書十三章3節告訴我們,提摩太已經出監.他也下監,為主多方受苦.
\newline
\newline
經訓練和裝備的過程,在神恩典之下,造就忠心神的僕人.
\newline
\newline
家的訓練,教會的訓練,在屬靈領袖帶領之下,各方面的訓練和裝備；也是今時代工人所需要的,萬勿輕視!
\newline
\newline
求主幫助我們看見,要事奉主,必須經過相當的訓練；在神的話語上,在禱告方面,在生活方面；求主幫助,要忍耐接受訓練.
\newline
\newline
神要幫助我們,幫助我們中間許多青年,成為下一代忠心的僕人.
\newline
\newline


\newpage

\section{第53屆~港九培靈會~戴紹曾~第十講}
\label{sec:611}
\textbf{差傳與戴氏家族(略述)}
\newline
\newline
連結: \href{https://www.hkbibleconference.org/session-message/view/611}{\texttt{https://www.hkbibleconference.org/session-message/view/611}}
\newline
\newline
\hyperref[sec:610]{< < < PREV SERMON < < <}
~
\hyperref[sec:toc]{[返目錄]}
~
\hyperref[sec:574]{> > > NEXT SERMON > > >}
\newline
\newline
十八世紀的英國,(那時還沒有美國)有位青年名戴雅各,他就是戴德生的曾祖.在外國,同時代用同名名是可以的,也是常有的.JamesTarlor的名於八代之間已有六代出現.
\newline
\newline
戴德生的曾祖戴雅各JamesTarlor在他的時代,英國很亂.他雖然也到教會,但還未得救.約翰衛斯理和傳道人到他的地方,都不受他歡迎,有時他甚至反對得很厲害,竟然準備爛蕃茄和壞雞蛋,在開會時,目標向著傳道人搗亂.某次,正當他站著瞄準時,忽然聽見神的話,講員講約書亞記廿四章15節:「至於我和我家,我們必定事奉耶和華.」他聽了就擲出所攜帶的臭物,擾亂聚會；但神的話卻銘刻在他心中.過了一段時間,當他結婚之日清晨,他想；今日結婚是我終身的大事,必須作好準備,於是就到田間安靜思想.這節經文不斷在他腦中,雖然極力想擺脫；可是欲罷不能,聖靈在他心裡作工.最後,他跪在田間,接受耶穌基督作他個人救主.並對主說:「主啊!我現在真的說:「至於我和我家,我必定事奉耶和華.」婚禮的時間已到,他急忙回家更衣向著禮拜堂直跑；他的朋友見他來了,才鳴鐘進行婚禮.禮畢茶會,沒有講道；他卻站起來說:「今天婚禮我遲到,很對不起大家.」他繼續說,他遲到的原因；新娘一聽之下,氣昏了,很不高興地說:「他信了耶穌,難道我嫁的是個傳道人?」,可是也無可奈何!戴雅各的心裡一直禱告,求主讓他的新娘早日接受主,同心事奉主；豈料他越禱告,她的心越剛硬.他問主:「主啊!你不是要我和我家事奉你嗎?為何她的心如此剛硬?」有一天他回家,覺得有點受不了,忽然抱起新娘到樓上臥室；用手壓她,迫她跪下；自己也跪下,流淚禱告說:「主啊!這些日子我所求的是照著你的話「全家要事奉你,」如今全家還有一半不認識你,求你感動她,叫她早日認識你.」這時,她也哭了；他以為手壓她太重；原來她也在禱告,認罪禱告,接受主.
\newline
\newline
各位青年,今晚一開始,我就講了這個見証,並非告訴你們一個新的佈道方法.如果你也這樣做,說不定會給太太打你一拳；這要視乎是否出於神的感動才好做.
\newline
\newline
十八世紀的英國,神作工,戴雅各和妻子接受耶穌；他們一生事奉主,他們的子孫,以至他們的曾孫一直在事奉主.
\newline
\newline
戴雅各的曾孫戴德生是我的曾祖,他生長在基督化的家庭,他父親是個藥劑師,母親是虔誠的基督徒；在家常聽見他父母的見証,他父親常常在教會講道.戴德生出了學校,在銀行工作；那段日子,結交不良的少年,誤入歧途；他心裡痛苦,明知神要他走傳道的道路,可是他愛世界.有一個假日,他進父親書房看書,發現有些福音單張標題寫著:「基督完成的工作.」他原想閱讀的是故事.不料這個新鮮的題目吸引了他；他越讀就越明白得救的真理:只要我們信祂,接受祂,不須我們付任何代價.耶穌基督在十字架上,已經為我們預備救贖,使我們可與神和好,罪得赦免.原來就在這一天早晨,他在遠處的母親去探望親戚；有了感動,就跪下為兒子禁食禱告,她心裡得著平安,深信神必聽禱告.同日,戴德生對妹妹說:「有件事我只告訢你,就是我今天得救了.」他還囑咐妹妹別告訴任何人.過了幾天,母親回家,戴德生想告訴母親,母親卻說:「我已經知道了.」「是妹妹說的嗎?」「不是的,那天當我為你禱告時,深知神聽了我的禱告；今天你親自告訴我,我更加歡喜.」
\newline
\newline
過了不久,戴德生開始為中國有負擔；其實,他的父親早就有意到中國傳福音,所以凡有關於中國題材的刊物必買.因此,家裡有很多這樣的資料.他從中發現到,每月有百萬不認識耶穌的中國人離開世界,這個統計打動了他的心；他幾乎受不住,覺得需到中國傳福音.與此同時,他認識了一位很有音樂天份的女朋友；他想,若一同到中國去傳福音,互相配搭,是很不錯的,初談此事,頗為順利；有一天,女朋友對他說:「談到我們倆的事,我父親不反對,他很歡喜你；但是他覺得到中國去的觀念要早日取消；傳福音,當牧師,這很好；在英國也是一樣；如果你堅持到中國去,我倆的婚事就不必再談了.」
\newline
\newline
各位青年!面對這樣的選擇,怎麼辦呢?隨著感情,是最容易的；而且按理,留在英國也是一樣事奉主,也可一生事奉主.可是,神的呼召不是這樣,不遵主旨意,就是妥協.鮑院長在早堂的聚會,一再的問:「我們為主到底撇下甚麼?」你甘心放棄甚麼?我們肯不肯為主的緣故這樣作呢?相信當時戴德生心裡有很大的掙扎；到中國去路途遙遠崎嶇；不如留在英國吧!各位青年!可能你的父母,你的先祖有機會聽福音；乃因著一位英國的青年人,作了決定.女朋友可以不要,中國卻不能不去；因為這是神的旨意.巴不得我們中間許多青年有相同的心志.
\newline
\newline
保羅說:「我活著為基督,我願意為祂撇下萬有.」他在腓立比書三章說他把萬事看作糞土,他的地位,背景,家庭,一切都願放下,他愛主愛家.他說:「為了骨肉之親,與主隔絕；若能叫他們可以得救,我也願意.一個基督徒並非不愛家,但必須愛主過於一切.
\newline
\newline
戴德生開始訓練自己,他覺得需要學醫；就在一位醫生手下學習.他也開始訓練自己的信心,臨別時,父母對他說:「你若有任何需要,我們願意在經濟上支持你.」他回答說:「請放心吧!一切都預備好了.」當他到了醫生那裡,醫生也對他說:「你有任何需要,請勿客氣；告訴我,我一定供給你.」他說:「請放心吧!一切都預備好了.」於是,醫生以為有他的父母負責；而他父母以為一切有醫生負責.至於戴德生的觀念,卻是一切有神負責；祂是耶和華以勒,祂要我去中國,祂一定負責.戴德生懷著這樣的信心開始學醫,週末他在人口頻密地區傳福音.有一個主日,他整個下午傳福音,派單張完畢,已相當疲倦；在回家途中,覺得背後有人緊隨著他；回頭一看,有個衣服襤褸的貧民.對他說:「戴先生,我妻病重在家,你可以為她診病嗎?」他這時雖然需要休息,但他認為有更重要的事要做,遂轉身和那人同去,到了一個又黑又小的房間,看見婦人臥在地上；就急忙為她診病,然後吩咐買藥.但他心中覺得有聲音對他說:「你這假冒為善的人,明知他沒錢請醫生,那有錢買藥呢?為何不把你口袋裡五塊錢給她呢?」他說:「主啊!我只有五塊錢,明天還要吃飯哩!我不能給.」主一再摧促,他也再三推辭；這時,他還要顯出一副屬靈的面孔來,對那人說:「我為你的太太禱告.」(這都是一般基督徒的所為)於是他跪下來禱告.事後他自己作見証,說那次的禱告通不到天上.各位!請不要弄錯,奉獻有奉獻的時候,禱告有禱告的時候.當他離去,神的聲音就對他說:「戴德生!你要到中國去,你今天在倫敦敢不信我,你到中國怎能信我?到時你的朋友離你很遠,誰來幫助你,你怎麼辦呢?」他遂拿出那五塊錢,心裡想,「我早就對主說,如果五塊錢可以分得開,我是會分給他的.」各位青年!有時主在我們身上的要求是絕對的,主要我們完全地順服.戴德生終於拿出那五塊錢,吩咐那人趕快去買藥.關於這事,戴德生在他的日記寫著:「當我走出門口,我的心和我的口袋一樣輕.」口袋裡沒有錢,可是心裡有平安；他相信有一位信實可靠的主.他回到家裡,把門關好,上樓準備睡覺；忽聞有扣門聲,原來郵差來送信；打開信封,內有五拾多元,他說:「主啊!赦免我,我的信心太小.」後來他自己說:「我們的主是信實的,快速勝過銀行,連本帶利都還給我了；且主給的利息比銀行更高,我要把錢存在主的銀行裡.」他一生在福音工作上,從來不接受教會奉獻的錢.內地會工作要開辦,他不肯去捐錢.他說:如果弟兄姊妹有奉獻的心,神有感動,他可以寄來；但我們不捐錢,相信神自己有預備.一個同工,十個同工,百個同工,或一千個同工在中國；神是信實的,祂為一個工人預備,也必為百千個同工預備；因為我們的神沒有難成的事.
\newline
\newline
今天因時間關係,我只好帶著各位走馬看花.
\newline
\newline
戴德生到中國之時,正是太平天國之時代,洪秀全攻佔上海.戴德生到了上海,學習中國語言,過了一段時間,他開始下鄉傳福音.有一次他想去南京向太平天國的人傳福音,可是沒有法子進去.他覺得穿著洋服,會攔阻他傳福音；他認為必須與中國人更加接近,聽者才能接受福音.因此,他改穿中國服裝.那時,上海的英國人對他非常不滿,認為這英國青年有損國體.可是,他好像保羅,為了福音的緣故不顧一切,同胞不諒解仍屬次要；傳福音使人接受耶穌作個人救主,是最重要的,他的同工也是如此.
\newline
\newline
神帶給戴德生作配偶的女子,名叫以馬利亞是個孤女；她的父親也是宣教士,神奇妙地帶她到中國寧波.(今晚無時間講,請參閱戴德生以馬利亞)戴德生在一八五四年到上海,過了六年才離開；當時他身體甚弱,醫生告訴他,恐怕一輩子不能再回中國.在那段期間,他進行翻譯工作.四年之久他不在中國,他禱告神將新異象給他,啟示他:「中國沿海一帶,大港口已經開放,有很多宣教士在那裡傳福音；可是內地有千千萬萬的人沒有福音可聽,你要進內地去.」他以為沒有可能,因為他本來不屬任何宗派,可憑著甚麼去呢?而且還要帶領同工同去.他到各宗派詢問,到處被拒絕.到了一八六五年某主日,可說是他重要之日,那時他在英國南邊；是日清晨,他無心去禮拜,他想,在英國每主日有人傳福音,許多人都聽到福音；然而中國內地,成千上萬人從未聽到福音.他在沙灘徘徊；主對他說:「你肯不肯去?是否願意接受挑戰?照著我給你的異象,開始禱告,也為24位同工禱告.」很奇妙!神帶領24位同工,組成中國內地會,他們就開始在中國內地傳福音.
\newline
\newline
由一八六五年至一九○五年,四十年之久,神幫助祂的僕人傳福音至每個省,雖福音還沒有完全傳開,可是在中國內地各省,都有祂的使者傳福音.
\newline
\newline
我的祖父生在中國,後來又回中國傳福音,他名戴存仁.我的曾祖按著中國人的規矩,老大名存仁,老二存義,還有存禮,存智,存信.戴存義也在中國傳福音.
\newline
\newline
戴德生在山東煙台,為他同工的子女設立學校,兄弟也曾在該校就讀.
\newline
\newline
差傳工作應面對同工子女教育問題,教會打發工人出去,孩子們的教育需要解決；若無妥善計劃,可能被派出的宣教士,很難久留在原工作地.
\newline
\newline
當我的祖父回英國就讀大學時,有一天,他的父親來信囑咐他回煙台,剛創辦的學校任教.(時在一八八一年離今正百年)祖父放棄讀大學,順服父命,離英赴中國山東,幫助教育宣教士的子女,他這方面的奉獻,是為了福音的緣故.他年老退休,那時我被關在日本集中營,有機會服侍祖父.他恆心愛主,每清晨五時起床,唱歌禱告,終身事奉主.
\newline
\newline
家父名戴永年,生於蘇格蘭,未滿週歲到中國；他也是一生在中國傳福音,他在上海得救的；後來到開封內地會醫院當藥劑師,作福音醫療工作.他有渴慕的心,自知雖已得救,但未過得勝聖潔的生活,亟需聖靈潔淨和充滿,故竭力追求,神大大賜他,得到聖靈充滿聖靈潔淨.他想繼續進修,因而寫信和他的叔叔聯繫,遂即到美國讀書.這件大事一方面考驗他的信心；他離中國時,盤資只夠到美國西岸,距目的地尚遠,但他深信神必有預備.到了西雅圖,住在朋友的地方；雖然囊空如洗,卻仍收拾行裝往火車站,上了櫃台；忽有人近前對他說:「弟兄!我有感動,我覺得應將此奉獻給你.」後來家父讀完大學,回到中國,帶著家母同往傳福音.
\newline
\newline
最後,我要見証,一九三九年中國抗戰時期,家父無法留在開封工作,就到山東,擬帶全家赴美；因為戰爭爆發,日本佔領之地,無法進行福音工作.家父到了山東,訂了船票；可是經懇切禱告,終而放棄此行；於是把孩子們留在山東,家父母又進內地去了.先到開封,再到陝西,西安一帶.家父心有負擔創辦聖經學校；當時西北常遭日機轟炸,朋友們認為不可能招得學生.家父為此禱告,求神預備地方,第一要水源充足,第二要安全地帶可以避免日機轟炸.過些日子,西北內地會同工報告家父,在保基有適合地點；於是設立西北聖經學校.五年之久,家父母在那裡教書.那裡院子很大,有五口井,用水不成問題,也沒有日機騷擾.到了一九四八年十二月八日,報紙大事標題日本轟炸珍珠港,家母閱報大衛震驚,她想到四個孩子在日軍佔領之地；她無法繼續閱報,急忙回到房裡,跪下痛哭.她不會禱告,只想到日軍進入南京和上海時,那些殘酷行為的報導,她想到自己的女兒,她哭了很久,直到主的聲音在她心裡:「女兒,你忘記我的話嗎?你還在山東時,我曾對你說,你若順服我,我必看顧他們.」主曾給家母一節經文:「要先求神的國和神的義,這些東西都要加給你們了.」(太六33)有一位老牧師將此經文翻成自己的話:「如果你關心神的事情,神要關心你的事情.」請看保羅,他關心主的事,一切都為福音的緣故；所以神關心保羅並使用他.家母跪在主面前,主說:「你若先求我的國關心我的事,我必照顧你的孩子.」「神的事情,」對家母來說,是在西北地方,有許多沒有聽到福音的人；許多中國同胞需要耶穌基督的福音.這是神所關心的.家母對主說:「主啊!我願意.」
\newline
\newline
過了不久,我們全校師生被擄囚禁集中營,和父母隔絕了五年半.感謝主!一九四五年八月,我們出了集中營,搭飛機到西安.去年我重臨該機場,回想卅五年前我初次抵達這個機場,不禁淚下,說:「主啊!你是信實的.」
\newline
\newline
當時,家父為我們一切都安排妥當；但他沒有告知家母.我們到西安的翌日早晨,有中國同工來接機,帶我們乘搭火車到小鎮,轉坐馬車；約經十五英哩路,時屆黃昏,抵家已八點鐘了.沿途羊腸小徑,人地生疏,加以言語不通；但是,神終於帶領我們四兄弟姊妹安抵目的地.有個西北聖經學校的學生迎面而來,領我們回到家中；老師們正開會.我們一到,他們欣然散會,留下家父母和我們,全家團聚.我實無法形容當時的快樂!五年半我們見不到父母,但信實的主沒忘記我們；雖然日本管轄之下,我們並不在日人手裡,卻在主的手中.
\newline
\newline
感謝主!你若先求神的國和祂的義,先關心祂的事情,以主和福音為前提；神會帶領你一生,神會使用你.
\newline
\newline


\newpage

\section{第53屆~港九培靈會~林克己~第一講}
\label{sec:574}
\textbf{蒙神呼召的神蹟}
\newline
\newline
連結: \href{https://www.hkbibleconference.org/session-message/view/574}{\texttt{https://www.hkbibleconference.org/session-message/view/574}}
\newline
\newline
\hyperref[sec:611]{< < < PREV SERMON < < <}
~
\hyperref[sec:toc]{[返目錄]}
~
\hyperref[sec:576]{> > > NEXT SERMON > > >}
\newline
\newline
出埃及記 2:23-2:25
\newline
\begin{longtable}{cl}
\hline
\hline
章節 & 經文 (和合本修訂版)\\
\hline
2:23 & \begin{tabularx}{0.7\textwidth}{X} 過了許多年，埃及王死了。以色列人因做苦工，就嘆息哀求；他們因苦工所發出的哀聲達於神。 \end{tabularx} \\ \\ \relax
2:24 & \begin{tabularx}{0.7\textwidth}{X} 神聽見他們的哀聲，就記念他與亞伯拉罕、以撒、雅各所立的約。 \end{tabularx} \\ \\ \relax
2:25 & \begin{tabularx}{0.7\textwidth}{X} 神看顧以色列人，神是知道的。 \end{tabularx} \\ \\
[1ex]
\hline
\hline
\end{longtable}
在這次講道會中,我們要從聖經的神蹟中看八十年代的挑戰.研究這些神蹟之前,我們先要面對一個問題:這些神蹟是否可信?請看(詩五十21)神說:「你行了這些事,我還閉口不言,你想我恰和你一樣,其實我要責備你,將這些事擺在你眼前.」這節經文的中心意思有二:
\newline
\newline
(1)神的偉大和人的渺小.
\newline
\newline
(2)神在這裡責備人的驕傲,對人說,你想我恰和你一樣.
\newline
\newline
這是人最基本的錯.我們軟弱,所能做的有限,對神的大能力有錯誤的觀念,以為神與我們一樣,作事有問題；事實上神比我們大,祂配做人崇拜和信心的對象.聖經高舉祂作我們全能的天父,在祂無難成的事,祂有超自然的能力,因此聖經上的神蹟非常可信.如果聖經沒有神蹟的記載,我們才會懷疑祂的偉大.
\newline
\newline
神蹟的可信,除了在神學上看得出來之外,考古學也提供了寶貴的見證,尤其是主耶穌在世上所行的神蹟,足以見證聖經的可靠和可信性.英國著名的考古學家蘭塞博士年青的時候,因為德國新派神學的影響,便對聖經的真確性產生疑惑；但經多年在中東研究考古學,對路加所作的使徒行傳研究過後.便發表意見說:「路加醫生實在是一位可靠認真而有才幹的歷史家,我在他所作的使徒行傳中,找不出一處歷史上的錯誤.」蘭塞博士的見證很有意義,因為如果使徒行傳是這樣可靠,我們想信路加所作的福音書也一樣可靠.路加福音記載了很多神蹟,包括最偉大最重要的神蹟,就是耶穌的復活,這是我們信仰的基礎,福音的中心.因此今年的講道會,我們一共要看八件有興趣有意義的神蹟.
\newline
\newline
第一件是蒙神呼召的神蹟.出埃及記第二,第三章,這件蒙神呼召的神蹟有五方面:
\newline
\newline
(一)是神的憐憫.(二23-25)「過了多年,埃及王死了,以列人因作苦工,就歎息哀求,他們的哀聲達於神.神聽見他們的哀聲,就記念他與亞伯拉罕,以撒,雅各所立的約,神看顧以色列人,也知道他們的苦情.」(三7-9)「耶和華說,我的百姓在埃及所受的困苦,我實在看見了.他們因為督工的轄制所發的哀聲,我也聽見了,我原知道他們的痛苦.我下來是要救他們脫離埃及人的手,領他們出了那地,到美好寬闊流奶與蜜之地,就是到迦南人,赫人,亞摩利人,比利洗人,希未人,耶布斯人之地.現在以色列人的哀聲達到我耳中,我也看見埃及人怎樣欺壓他們.」
\newline
\newline
這兩處經文高舉了耶和華的慈愛和無所不知.信徒的苦楚,神能看見；信徒的禱告,神能聽見.信徒在苦難中學到了神要他們學的功課之後,神便必然下來釋放他們,使他們見證他的憐憫和恩典.難怪神的僕人王明道作了一本受苦有益的書,也有一位宣教士說:「神的兒女,雖然有時候吃苦,但是永不吃虧.」吃苦和吃虧有極大的分別,我們所面對的試煉和患難,雖然好像難經過的高山,可是經過以後,我們便能作別人的嚮導,挑戰和安慰；使他們在我們身上可以看到,作基督徒有甚麼價值和好處,在凡事上認定主有甚麼好處.
\newline
\newline
至於禱告方面,這兩段經文也有寶貴的教訓.信徒多次的禱告,恆切的禱告,雖然這樣,神卻好像聽不見,也不應允.實際上我們每逢禱告,天父都聽見,時候到了便應允.(羅八26)說:「況且我們的軟弱有聖靈幫助,我們本不曉得當怎樣禱告,只是聖靈親自用說不出來的嘆息,替我們禱告.」這節經文的主題是信徒的祈禱,我們的軟弱有聖靈幫助.我們的禱告,聲音雖然非常微小,但是聖靈是天父的愛心,他為我們預備擴音器,能將我們的禱告清楚的送到神的耳中.我們隨時隨地可以到神施恩的寶座前,為自己和別人求恩；深信神必聽見和應允我們的禱告.
\newline
\newline
(二)是神的顯現.(出三1,6)「摩西牧養他岳父米甸祭司葉忒羅的羊群,一日領羊群往野外去,到了神的山,就是何烈山.耶和華的使者從荊棘裡火燄中,向摩西顯現.摩西觀看,不料,荊棘被火燒著,卻沒有燒燬.摩西說:我要過去看這大異象,這荊棘為何沒有燒壞呢?耶和華神見他過去要看,就從荊棘裡呼叫說,摩西,摩西.他說:「我在這裡.」神說:「不要近前來,當把你腳上的鞋脫下來,因為你所站之地是聖地.」又說:「我是你父親的神,是亞伯拉罕的神,以撒的神,雅各的神.摩西蒙上臉,因為怕看神.」
\newline
\newline
神為甚麼要用燒著的荊棘來向摩西顯現呢?其中我們可以看見神的智慧,神呼召和帶領他兒女們的時候,必然用一些適合他們環境和背景的方法.祂知道怎樣向我們說話,引起我們的詫異.摩西在曠野牧羊,必定多次看過燒著的荊棘,但是摩西知道今次有神蹟,因為荊棘沒有被火燒壞,和他以前見的不同.所以我們可以確實知道,神要啟示他為我們生命所預備的計劃,所定的旨意,是必定按我們的環境,背景,性格,向我們說話.
\newline
\newline
另一方面,從全本聖經的教訓中,我們知道火是神與人同在的象徵,神曾經三次用火向他的僕人以利亞顯現,第一次在迦密山,第二次在曠野,第三次在以利亞乘旋風升天的時候.新約聖靈降臨的時候,有風和火；風代表神的大能,火代表神的同在.所以,神在火燒的荊棘中向摩西顯現,是提醒祂和我們,祂呼召我們為祂作工,必定每時刻與我們同在,加添我們的力量,引導我們的路程,供給我們的需要,使用我們的見證.還有,這顯示出信徒在受試煉和逼迫的時候,全能的神能保護他們.雖然火燒,卻沒有燒燬.以色列人雖然在埃及受壓迫,但在全能的神的領導下；他們還有光明的前程,平安經過四百年的試煉,最終到達應許之地.
\newline
\newline
(三)是神的揀選.(出三10)「故此我要打發你去見法老,使你可以將我的百姓以色列人從埃及領出來.」(四13)「摩西說:主阿!你願意打發誰,就打發誰去吧!」這兩節經文告訴我們,萬王之王耶和華,用自己絕對的主權,揀選摩西的時候,摩西覺得自己很不適合作主的僕人,以色列人的領袖和代表,去見法老,把以色列人從埃及地領出來.摩西以為神揀選錯了,便求神打發他人,因為摩西感到自己的軟弱和不配.事實上,這是被主使用的條件和秘訣,是靈性生活健康的表證.(林後十二9),祂對我說:「我的恩典夠你用的,因為我的能力,是在人的軟弱上顯得完全,所以我更喜歡誇自己的軟弱,好叫基督的能力覆庇我.」保羅見證神的恩典怎樣夠他用.只可惜多數信徒讀這節聖經的時候,只留意上半節而忽略下半節,其實下半節有信徒得著靈力的秘訣,被神使用的條件.大部份人都喜歡誇自己的好處,優點,只有真正屬靈的人,才會誇自己的軟弱,樂意把自己的弱點擺在神面前,這是神所喜悅的一種屬靈的誇口.這樣的信徒,雖然知道自己不配,但是並不推辭神的呼召,他只是緊緊的倚靠主,將人的軟弱換取神的大能.
\newline
\newline
(四)是神的安慰.(出三11-12)「摩西對神說:我是甚麼人,竟能去見法老,將以色列人從埃及領出來呢?神說:我必與你同在,你將百姓從埃及領出來之後,你們必在這山上事奉我,這就是我打發你去的證據.」(四10-12)「摩西對耶和華說:主啊!我素日不是能言的人,就是從你對僕人說話以後,也是這樣,我本是拙口笨舌的.耶和華對他說:誰造人的口呢?誰使人口啞,耳聾,目明,眼瞎呢?豈不是我耶和華麼?現在去吧!我必賜你口才,指教你所當說的話.」摩西這個反應,相信是很多現代基督徒蒙神呼召時內心的感想,我不配也不會.
\newline
\newline
二十多年前,兄弟蒙召作宣教師的時候,心裡的疑惑,口裡的推辭都很多.我記得我對主說:「主阿!我知道你揀選錯了,我不配作宣教師,我不夠聖潔,我也不會作這個工作；恩賜,口才,經驗都沒有.我要用六個禮拜時候才預備一次講道；還有見證方面,也不會大大的吸引人.」但感謝主!我雖然推辭,但是主很忍耐,繼續感動呼召我,為我解決一切問題.當摩西感到不配,說我是甚麼人的時候,主馬上說我必與你同在；感到自己不會的時候,主便說我必賜你口才,指教你所當說的話.主呼召人起來工作的時候,必給他恩典和力量.
\newline
\newline
(五)是神的計劃.(徒七22)「摩西學了埃及人一切的學問,說話行事都有才能.」神為每一個信徒都有完整的計劃.摩西蒙神呼召時,雖然深感不配,但事實上卻很有資格負起這個責任.他已經有足夠的學問,訓練,經驗；因為在王宮長大,就學習了管理百姓的功課；因為在曠野有牧羊的經驗,就知道怎樣作以色列人的牧者.行政牧養方面,都適合作以色列人的領袖.所以我們蒙召服事主的時候,可以坦然無懼.神是無所不知的,祂能用我們各樣的經歷,預備訓練我們,成為合祂心意的僕人.
\newline
\newline
有一個年青人,蒙召作宣教師,在海外的土人中作開荒的工作.很多認識他的人便勸他不要去,第一因為那些土人是吃人的.第二因為這位年青人是跛腳的,只有一隻腳,另一隻是木造的.雖然人再三的勸他不要去,但是這位弟兄卻清楚神的呼召,一定要去.一天,土人果然抓住這位弟兄說:「我們的火和鍋都預備好了,我們今天要吃你.」這位宣教師便馬上拿起一把刀,從他木造的卻上割下一塊來,遞給土人說:「你們在吃我以前,應該先嚐一嚐我的味道,到底好不好?你們先試試這個.」他們試過後,覺得味道不好,而且太硬,結果便放了他.宣教師後來回國休假的時候,便見證主的無所不知和奇妙的智慧,神為甚麼要揀選一個殘廢跛腳的人,傳福音給吃人肉的人聽呢?結果受感動的人很多,他們知道神揀選人之前,必定先預備訓練他.今天主也必照樣恩待我們,願主在他的工作上,大大的使用我們.
\newline
\newline


\newpage

\section{第53屆~港九培靈會~林克己~第二講}
\label{sec:576}
\textbf{苦水變甜的神蹟}
\newline
\newline
連結: \href{https://www.hkbibleconference.org/session-message/view/576}{\texttt{https://www.hkbibleconference.org/session-message/view/576}}
\newline
\newline
\hyperref[sec:574]{< < < PREV SERMON < < <}
~
\hyperref[sec:toc]{[返目錄]}
~
\hyperref[sec:578]{> > > NEXT SERMON > > >}
\newline
\newline
出埃及記 15:22-15:26
\newline
\begin{longtable}{cl}
\hline
\hline
章節 & 經文 (和合本修訂版)\\
\hline
15:22 & \begin{tabularx}{0.7\textwidth}{X} 摩西領以色列人從紅海起程，到了書珥的曠野，在曠野走了三天，找不到水。 \end{tabularx} \\ \\ \relax
15:23 & \begin{tabularx}{0.7\textwidth}{X} 到了瑪拉，他們不能喝瑪拉的水，因為水是苦的；所以那地名叫瑪拉。 \end{tabularx} \\ \\ \relax
15:24 & \begin{tabularx}{0.7\textwidth}{X} 百姓就向摩西發怨言，說：「我們喝甚麼呢？」 \end{tabularx} \\ \\ \relax
15:25 & \begin{tabularx}{0.7\textwidth}{X} 摩西呼求耶和華，耶和華指示他一棵樹。他把樹丟在水裡，水就變甜了。耶和華在那裡為他們定了律例、典章，在那裡考驗他們。 \end{tabularx} \\ \\ \relax
15:26 & \begin{tabularx}{0.7\textwidth}{X} 他說：「你若留心聽從耶和華－你神的話，行我眼中看為正的事，側耳聽我的誡令，遵守我一切的律例，我就不將所加於埃及人的疾病加在你身上，因為我是醫治你的耶和華。」 \end{tabularx} \\ \\
[1ex]
\hline
\hline
\end{longtable}
今年講道會的總題是主奇妙的作為,我們要看八件有興趣有挑戰的神蹟.學習怎樣跟從服事主,其中有苦水變甜的神蹟.這段聖經也是描寫以色列人的領袖摩西的經過,在(出十五22-26),看他怎樣作我們的榜樣.
\newline
\newline
這段聖經可以分五段來研究.
\newline
\newline
(一)摩西怎樣作領袖二十二節上「摩西領以色列人從紅海往前行」.前天我們已經看到,摩西因為心裡謙卑,沒有高舉自己作以色列人的領袖；他因為不想吩咐別人行這條路,做那件事,所以他雖然有才幹經訓練,但是他仍諸多推辭.神需要多多的感動鼓勵他,他才願意出來做領袖.可是,這一顆謙卑的心,是神所賜福喜悅並使用的,是基督徒得神使用的秘訣,和不能缺少的屬靈資格.(民十二3)稱讚摩西說:「摩西為人極其謙和,勝過世上的眾人.」這裡我們看見一個合神心意,被神使用的僕人,是在神和人面前都保持一顆謙卑的心,凡事都有適合的態度.同樣,我們做教會的工作,不是因為自己揀選這責任,乃是因為神呼召我們出來為他工作.我們需要有摩西這樣謙卑的態度,才能被神使用.
\newline
\newline
兄弟頭幾年學習講國語的時候,有差會的華文課程,要求我背熟一句成語,就是「人不知己過,牛不知力大.」這句話非常有意思.人因為內心的驕傲,犯罪的結果,看不出自己的過失毛病.聽說有一位神學生非常驕傲,常常誇口,同學忍受不住,便派代表團去對付他.他們到他房間說話,他一語不發,很安靜的聽,他們以為這位弟兄肯接受規勸.想不到這弟兄卻回答說:「謝謝你們的愛心,但你們責備我的時候,我連一句話也沒有說,相信再沒有別人比我這樣謙卑了.你們說我驕傲,但我以為全世界大概沒有人比我更謙卑的了.」看不出自己的軟弱,因自己的謙卑而驕傲起來,是這位弟兄的毛病.
\newline
\newline
(二)受試驗二十二節下「到了書珥的曠野,在曠野走了三天找不著水.」這是第一次想不到的困難.以色列人在摩西帶領之下,已經看見主奇妙的作為；雖然困難重重,但是神大能的手已經把他們領出來,使他們經過紅海而到達曠野.可是,摩西和以色列人都發現,紅海和曠野的地理環境是大不相同的.在書珥的曠野他們要受試驗,看他們的信心到底怎樣.口渴的時候,走了三天也找不到水喝,這時他們要灰心喪膽還是仰望神呢?他們要批評摩西還是信任他呢?
\newline
\newline
神為甚麼要用這樣的方法去試驗祂所愛的人呢?(代下三十二31)「惟有一件事,就巴比倫王差遣使者來見希西家,訪問國中所現的奇事.這件事神離開他,要試驗他,好知道他心內如何.」在試驗方面,摩西和希西家都有一樣的經歷；當他們事事如意的時候,耶和華好像離開他們,不再施恩,不再行神蹟.不單如此,神在試驗之上還加試驗,(出十五23)「到了瑪拉,不能喝那裡的水,因為水苦,所以那地名叫瑪拉.」摩西作領袖,遇到兩個想不到的難處.第一是走三天也找不著水；第二是找到的水竟是苦的.在人來說,很容易對神的大能和大愛發生疑問.還有第三方面的怨言,24節「百姓就向摩西發怨言說:我們喝甚麼呢?」我們可以想像到摩西所經過的,不但有想不到的難處,也有想不到的批評.但是摩西在這裡有美好的反應,他忍受著,並沒有辭職.
\newline
\newline
(三)禱告神二十五節上「摩西呼求耶和華.」有些基督徒問:今日是恩典時代,是聖靈降臨福音傳遍的時代,神為甚麼還給我們舊約聖經,這部份是描寫古代以色列人的風俗,記載他們的宗教歷史,對二十世紀的信徒有甚麼用處呢?(林前十10-11)「你們也不要發怨言,像他們有發怨言的,就被滅命的所滅.他們遭遇這些事,都要作為鑑戒,並且寫在經上,正是警戒我們這末世的人.」這是保羅所給我們的答案.保羅在這裡論到以色列人的失敗,並以此警戒我們.對於好的當然我們要效法,就像摩西一樣,在困難的時候,不但沒有辭職,反而用堅固的信心禱告神,用堅定不移的信仰仰望神,這是我們要學習的.
\newline
\newline
在我們服事主的生命中,神必定有時也領我們到書珥的曠野,是一個不能想辦法而只能求辦法的地方,為要粉碎我們的驕傲,使我們承認自己的有限,只有全能全智的主才有解決的方法.這些經歷非常寶貴,能夠訓練信徒,專心倚靠神而不靠人,求神賜智慧與能力,跟從摩西的腳步呼求耶和華.
\newline
\newline
有一位蘇格蘭的女宣教師,蒙召到非洲去開荒傳道.在這個環境底下,她要接受神許多試煉.一天,她病倒了,而且病得很重,她又沒有錢,不要說買藥,就連每天的食物也沒有.這位姊妹沒有辦法,只得跪下向神禱告,求神解決她的問題.結果,就在她禱告的時候,郵差拍門,給了她一個很大的包裹.她打開看看,裡面沒有別的,只有蘇格蘭人早餐所吃的燕麥粉.女宣教師因為沒有藥,又沒有別的食物,只有這個,於是她便不只早餐吃燕麥粉,連午餐和晚餐都要吃,這樣吃了六個星期,吃得她都膩了.可是,奇怪的是,當她愈吃燕麥粉的時候,她的身體便愈加健康,六個星期後,她就全康復了.後來她回國休假的時候,便作這個奇妙的見證.見證會完畢,一位基督徒的醫生便對她說:姊妹,神真是奇妙的.你在非洲所患的那種病,治療方法就是吃六個星期的燕麥粉,此外不吃別的.神不單給了你所需的食物,更給了你所需的藥.這一個見證,給我們確實的知道禱告的果效.
\newline
\newline
(四)蒙啟示二十五節「摩西呼求耶和華,耶和華指示他一棵樹.」請注意,摩西如果靠自己的聰明智慧,很難知道附近的這一棵樹就是神為他解決的方法.所以神要給他特別的啟示,他才明白.我們有困難的時候,也應該求主開我們的眼睛,使我們看清楚祂的能力,方法,作為是怎樣奇妙,還合我們的需要.
\newline
\newline
我們看另一個舊約的例子.(王下六15-18)「神的僕人清早起來出去,看見車馬軍兵圍困了城.僕人對神人說:「哀哉!主阿,我們怎樣行才好呢?神人說:不要懼怕,與我們同在的,比與他們同在的更多.以利沙禱告說:「耶和華阿,求你開這少年人的眼目,使他能看見.耶和華開他的眼目,他就看見滿山有火車火馬圍繞以利沙.」敵人下到以利沙那裡,以利沙禱告耶和華說:「求你使這些人的眼目昏迷.耶和華就照以利沙的話,使他們的眼目昏迷.」這段經文講到一個心裡害怕的年青人,他只看見環境的危險,所以便說哀哉.相信這是一個屬肉體的基督徒遭遇難處時的自然反應,因為有習慣倚靠自己而不仰望神,遇到危機不能想辦法,便說哀哉!但是這卻沒有用.摩西和以利沙遭遇危機時並不是這樣,摩西注重禱告,呼求耶和華,耶和華便指示他一棵樹；以利沙一樣注重禱告,他求神開少年人的眼睛,使他看見屬靈的真理和事實.少年人的眼睛開了,才看見屬靈隱藏的真理,滿山不單有敵軍,也有火車火馬,代表神的同在,大能和拯救.讓我們效法摩西的榜樣,受試驗就禱告神,禱告神就蒙啟示.
\newline
\newline
(五)看神蹟(二十五下)「他把樹丟在水裡,水就變甜了.」這不是主奇妙的作為嗎?這節經文關於信心的,有四個功課.第一,信心的仰望,摩西呼求耶和華.第二,信心的賞賜,耶和華指示他一棵樹.第三,信心的順服,他把樹丟在水裡.第四,信心的神蹟,水就變甜了.仰望,順服,然後便看神蹟,這不單是個屬靈的次序,也是屬靈的原則.
\newline
\newline
這個神蹟還有一個挑戰,安慰和勉勵,就是耶和華要摩西與他同工,神蹟才出現.耶和華雖然指示摩西一棵樹,但是摩西要出力把樹丟在水裡,水才能變甜.所以這神蹟包含了兩方面:神的啟示和人的出力.神有啟示的責任,而人有出力的責任.在我們屬靈的生活和工作上,神很願意負起啟示的責任,但我們肯一樣的殷勤嗎?這實在是我們的特權和福氣,雖然神無所不能,但是祂願意與我們同工.
\newline
\newline
有一位英國弟兄蒙召到日本作宣教師,他把這感動告訴牧師,牧師很高興,也找機會訓練他.有一天,牧師告訴他,教會今晚要在街上開露天佈道會,到時請他作見證.但是那位弟兄雖然清楚神的呼召,卻推卻牧師的邀請說:我從小有口吃的習慣,我不行.但是牧師說:如果你相信神已經呼召你作宣教師,他一定願意為你解決口吃的習慣.你起來作見證之前,要默禱,求主釋放你的舌頭,給你當說的話.你這樣用信心呼求神,神必叫你看見他奇妙的作為.果然當這弟兄作見證的時候,不單有從上頭來的智慧和能力,有吸引力和感動力；同時更有神蹟出現.他所講的話,就像流動的江河一樣,很自然的從他口中出來.以後他就再沒有口吃發生了.倚靠神的大能,問題就得到根本解決.後來這弟兄成為神很使用的旅行佈道家；在日本很多地方,傳揚福音,領人歸主.
\newline
\newline
摩西所經歷的這個神蹟,相信在很多方面都適合我們,但願我們從其中學習它的教訓.
\newline
\newline


\newpage

\section{第53屆~港九培靈會~林克己~第三講}
\label{sec:578}
\textbf{力上加力的神蹟}
\newline
\newline
連結: \href{https://www.hkbibleconference.org/session-message/view/578}{\texttt{https://www.hkbibleconference.org/session-message/view/578}}
\newline
\newline
\hyperref[sec:576]{< < < PREV SERMON < < <}
~
\hyperref[sec:toc]{[返目錄]}
~
\hyperref[sec:583]{> > > NEXT SERMON > > >}
\newline
\newline
創世記 11:31-11:32
\newline
\begin{longtable}{cl}
\hline
\hline
章節 & 經文 (和合本修訂版)\\
\hline
11:31 & \begin{tabularx}{0.7\textwidth}{X} 他拉帶著他兒子亞伯蘭和他孫子，哈蘭的兒子羅得，以及他的媳婦，亞伯蘭的妻子撒萊，一同出了迦勒底的吾珥，要往迦南地去；他們來到哈蘭，就住在那裡。 \end{tabularx} \\ \\ \relax
11:32 & \begin{tabularx}{0.7\textwidth}{X} 他拉共活了二百零五年，就死在哈蘭。 \end{tabularx} \\ \\
[1ex]
\hline
\hline
\end{longtable}
今天所講的神蹟,相信是每一位基督徒都必須經驗的,就是力上加力的神蹟.奔跑天路的時候,雖然遭遇難處,信心會軟弱,但是因為仰望主,就可以抬起頭來,重新得力.在研究這寶貴的真理之前,我們先要注意基督徒應該躲避的一種錯誤和危險,就是半途而廢.因為不進步就停步,停步之後就退步；因為不力上加力,就半途而廢,不再蒙神賜福並使用的熱心信徒.今日的教會,特別是年青的信徒,很多是半途而廢的.雖然屬靈的生活有良好的開始,讀經禱告,參加聚會,領人歸主,成為聖潔,在凡事上得主的喜悅；但是過了一段日子,便不再那樣熱心,不再有跟從主到底的心志了.
\newline
\newline
聖經起碼有三個人可以作我們的鑑戒.第一個是底馬.(提後四10)「因為底馬貪愛現今的世界,就離棄我往帖撒羅尼迦去了.」底馬是保羅的同工,他跟從主雖然有好的開始,但因為魔鬼的試探和世界的引誘,半途而廢,放下傳福音的工作,到帖撒羅尼迦去,尋找物質上的享受.第二個是新約聖經第二本福音書的作者馬可,他在事奉方面半途而廢.但是感謝主,他的跌倒和軟弱上是暫時的.(徒十五36-39)「過了些日子,保羅對巴拿巴說,我們可以回到從前宣傳主道的各城,看望弟兄們景況如何.巴拿巴有意,要帶稱呼馬可的約翰同去.但保羅因為馬可從前在旁非利亞離開他們同去作工,就以為不可帶他去.於是二人起了爭論,甚至彼此分開,巴拿巴帶著馬可,坐船往居比路去.」我們注意馬可離開保羅的地方,就是旁非利亞.按照聖經歷史家的發現,古代的旁非利亞,不單是一個天氣悶熱的地方,而且蚊子也很多；因為蚊子帶給人一種很厲害的熱病,所以經過旁非利亞而生病的旅客很多.結果有一班解經家發表意見說,馬可在這裡離開保羅回到耶路撒冷,是因為怕經過旁非利亞,他是這樣過份擔心自己的健康.這班解經家,也相信保羅在哥林多後書十二章所作的見證,關於他身體上的一根刺,就是撒旦的差役來攻擊他.論到他在旁非利亞所生的病.還有一班解經家,讀了加拉太書,就猜想保羅的刺是一種眼病.但我相信基督徒所生的病,若是神的旨意,他必能行神蹟醫治我們.不過,聖經沒有告訴我們保羅的那根刺是甚麼,是要叫我們為自己身體或生活上的每一根刺,用信心仰望他,從而經歷他的恩典夠我們用.
\newline
\newline
還有一處關於馬可的經文,(提後四11)「獨有路加在我這裡,你來的時候要把馬可帶來,因為他在傳道的事上於我有益處.」提摩太後書是保羅所寫最後的一封書信,在這裡年老的保羅告訴我們一個好消息,馬可在事奉方面的半途而廢不過是暫時的,他跌倒之後便起來.他為我們留下一個很好的榜樣.
\newline
\newline
第三個例子,就是今天我們要特別注意的,他就是亞伯拉罕的父親他拉.(創十一31-32)「他拉帶著他兒子亞伯蘭,和他孫子哈蘭的兒子羅得,並他兒婦亞伯蘭的妻子撒萊；出了迦勒底的吾珥,要往迦南地去；他們走到哈蘭,就住在那裡.他拉共活了二百零五歲,就死在哈蘭.」我們不曉得他拉有沒有歸向耶和華,但是這一個人,雖然不一定是個半途而廢的例子,但一定是個半途而廢的比方.我們在他身上可以看清楚,基督徒要躲避的錯誤和危險；就是雖然有好的開始,但沒有力量和志氣.你想做一個跟從主的屬靈旅客,相信他拉的經歷,值得我們多默想.因為教會中有很多熱心的信徒,特別是年青人,雖然開始很好,但幾年之後,情形又怎樣呢?
\newline
\newline
我們在剛才的兩節經文中,看看他拉所作五件有意義的事:
\newline
\newline
(一)離開吾珥他拉和他的家人,按著他的建議和計劃,出了迦勒底的吾珥.我們特別注意「出了」這兩個字.一九二九年,英國一位出名的考古學家伍雷博士,發掘到他拉所住的迦勒底的吾珥,便發現到在主前兩千年左右,就是亞伯拉罕的時代中,住在吾珥的人很有學問.這是一個文明的社會,他們不單會在泥版上寫字,也有歷史記錄和數學課本,可是在宗教信仰上很有問題.在吾珥城的中心,有一座很壯觀的神廟,廟分四層樓,分別塗上銀色,紅色,藍色和黑色.原來這是他們所拜的月亮的廟宇,每種顏色都代表不同的意思,但均與月亮有關.最高的一層是他們聚集敬拜的地方,銀色就像月亮一樣.紅色代表太陽下山月亮出來的時候,藍色代表天空普通的顏色,黑色則代表月亮所復活的範圍,就是黑夜.這間華美的廟宇很有影響力,引誘全城的人都來拜月亮.然而他拉離開吾珥,可以說是很好的決定,也可以說是基督徒信主之時所作的決定.今天決定接受耶穌的人也很多,他們願意在宗教信仰上,離開吾珥,不再拜偶像,或是站在無神論的立場上；乃是敬拜事奉創造之主.吾珥神廟的引誘力,遠比不上福音的大能.他拉所作第一件有意義的事,就是離開吾珥.
\newline
\newline
(二)要往迦南地去很多人以為亞伯拉罕是頭一個要往迦南地去的人,其實是他的父親他拉.當他拉離開吾珥之後,他心目中的去處是迦南地.迦南地是流奶與蜜之地,在舊約聖經上,在我們屬靈的生活上,它都有非常寶貴的意思,是代表事奉的生活,豐盛的生命.在迦南可以享受恩典,服事恩主,這樣過一個愉快美滿的人生.雖然先要經乾燥可怕的曠野,遭遇患難,接受挑戰,勝過仇敵,但是到了迦南,可以享受存到永遠的豐盛生命,今生來生都有屬靈的福氣.
\newline
\newline
(三)走到哈蘭哈蘭是離開吾珥前往迦南的半路,這城離吾珥和迦南均有四百多哩遠；無論地理方面,屬靈方面,哈蘭都很重要.它是半途,也是邊疆.這裡的人和吾珥的人一樣,都是拜月亮.以為月亮是他們應該拜的地方神.吾珥到哈蘭,是地方神月亮所統治的範圍.經過哈蘭往前走,便進入另一個地方神的領土.在這方面,古代人信仰,他們承認和仰望一些地方神.當雅各欺騙他哥哥以掃之後,他便逃命,離開迦南回到哈蘭.(創二十八10)「雅各出了別是巴,向哈蘭走去.」當時雅各對無所不知,無所不在的神認識不夠,他以為以色列人的神只是以色列人在迦南所拜的地方神.後來他做夢以後,對神才有了新的信仰,十六節「雅各睡醒了說,耶和華真在這裡,我竟不知道.」神的無所不在,是我們基督徒一個很寶貝的信仰.因此,哈蘭地不單是半途,屬靈方面也是邊界,他拉要突破,要經過,才能真正認識那位創造天地,無所不在,唯一的真活神.今天教會也有許多人,在屬靈上雖然有好的開始；但是沒有經過哈蘭,就是事奉的生命,豐盛的生命的邊界；因為沒有將自己完全交給神,所以不但不能永遠享福,也不能被神使用.
\newline
\newline
(四)住在那裡請問在你屬靈的生命中,你有沒有停步,不再往前行呢?在外表看來,你還有心追求主,但裡面是否像他拉一樣,住在哈蘭呢?聖經告訴我們,在屬靈的追求中；人不是進步,便是退步.信靠耶穌,傳揚福音,都好像踏單車一樣,雖然慢慢向前走,但是不成問題；不過如果停下來,便要跌倒了.願主幫助我們,當走到這邊界地哈蘭的時候,不要停下來；乃要經過,突破,進入豐富的迦南美地.這樣看來,吾珥雖然代表信主之前的犯罪生活,但哈蘭卻代表半途而廢,有所保留,軟弱的靈性生活.
\newline
\newline
(五)死在那裡這是他拉所作最可惜最不能挽救的事,他雖然很長命,但卻死在哈蘭,永遠達不到他離開吾珥之時所定的目的地.聖經為甚麼記載這些事呢?他拉的生平,和我們八十年代的基督徒有何關係呢?(林前十1-4)「弟兄們,我不願意你們不曉得,我們的祖宗從前都在雲下,都從海中經過,都在雲裡海裡受洗歸了摩西.並且都吃了一樣的靈食,也都喝了一樣的靈水,所喝的是出於隨著他們的靈磐石,那磐石就是基督.」請注意這幾節經文中的「都」字,他們都在雲下,都從紅海中經過,都跟從摩西,都有一樣認識主服事主的機會,但結果怎樣呢?第五節「但他們中間,多半是神不喜歡的人,所以在曠野倒斃.」聖經為甚麼要記載這些事呢?答案在第六,七兩節,這些事都是我們的鑑戒,叫我們不要貪戀惡事,像他們那樣貪戀的.也不要拜偶像,像他們有人拜的,如經上所記,百姓坐下吃喝,起來玩耍.」我們特別提「玩耍」二字.我們跟從主,服事主,並不是玩耍,乃是戰爭,在屬靈的戰場上與撒旦爭戰.做基督徒也不是精神寄託,逃避現實的事,乃是一個一生跟從主的決定.哈蘭是事奉的生活,豐盛的生命,要享受這些,就必須經過邊界.
\newline
\newline
聖經把他拉的事蹟寫下,是有意思的.保羅在林前十11有這樣的批評「他們遭遇這些事,都要作為鑑戒,並且寫在經上,正是警戒我們這末世的人.」
\newline
\newline
然而半途而廢只是消極的,我們所著重的,乃是力上加力.(詩八十四7)「他們行走,力上加力,各人到錫安朝見神.」神要我們行走,不要跌倒,更不要死在哈蘭；乃是要力上加力,這真是一個屬靈的神蹟.因為人行走,通常不是力上加力,而是累上加累,好像他拉一樣,漸漸走不動.但是聖經詩篇卻應許我們要力上加力,向前直跑,一直到迦南美地.願主加力量行神蹟,叫我們能過事奉的生活,有豐盛的生命.
\newline
\newline


\newpage

\section{第53屆~港九培靈會~林克己~第四講}
\label{sec:583}
\textbf{降火下雨的神蹟}
\newline
\newline
連結: \href{https://www.hkbibleconference.org/session-message/view/583}{\texttt{https://www.hkbibleconference.org/session-message/view/583}}
\newline
\newline
\hyperref[sec:578]{< < < PREV SERMON < < <}
~
\hyperref[sec:toc]{[返目錄]}
~
\hyperref[sec:585]{> > > NEXT SERMON > > >}
\newline
\newline
列王記上 18:20-18:46
\newline
\begin{longtable}{cl}
\hline
\hline
章節 & 經文 (和合本修訂版)\\
\hline
18:20 & \begin{tabularx}{0.7\textwidth}{X} 亞哈就派人到以色列眾人那裡，召集先知上迦密山。 \end{tabularx} \\ \\ \relax
18:21 & \begin{tabularx}{0.7\textwidth}{X} 以利亞近前來對眾百姓說：「你們心持二意要到幾時呢？如果耶和華是神，就當順從耶和華；如果是巴力，就當順從巴力。」百姓一言不答。 \end{tabularx} \\ \\ \relax
18:22 & \begin{tabularx}{0.7\textwidth}{X} 以利亞對百姓說：「作耶和華先知的只剩下我一個；巴力的先知卻有四百五十人。 \end{tabularx} \\ \\ \relax
18:23 & \begin{tabularx}{0.7\textwidth}{X} 請給我們兩頭牛犢，巴力的先知可以為自己挑選一頭牛犢，切成小塊，放在柴上，不要點火；我也預備一頭牛犢放在柴上，也不點火。 \end{tabularx} \\ \\ \relax
18:24 & \begin{tabularx}{0.7\textwidth}{X} 你們求告你們神明的名，我也求告耶和華的名。那應允禱告降火的就是神。」眾百姓回答說：「好主意。」 \end{tabularx} \\ \\ \relax
18:25 & \begin{tabularx}{0.7\textwidth}{X} 以利亞對巴力的先知說：「因為你們人多，先挑選一頭牛犢，預備好了，求告你們神明的名，卻不要點火。」 \end{tabularx} \\ \\ \relax
18:26 & \begin{tabularx}{0.7\textwidth}{X} 他們把所給他們的牛犢預備好了，從早晨到中午，求告巴力的名說：「巴力啊，求你應允我們！」卻沒有聲音，也沒有回應。他們就在所築的壇四圍蹦跳。 \end{tabularx} \\ \\ \relax
18:27 & \begin{tabularx}{0.7\textwidth}{X} 到了正午，以利亞嘲笑他們，說：「大聲求告吧！因為它是神明，它或許在默想，或許正忙著，或許在路上，或許在睡覺，它該醒過來了。」 \end{tabularx} \\ \\ \relax
18:28 & \begin{tabularx}{0.7\textwidth}{X} 他們大聲求告，按著他們的儀式，用刀槍刺割自己，直到渾身流血。 \end{tabularx} \\ \\ \relax
18:29 & \begin{tabularx}{0.7\textwidth}{X} 中午過去了，他們狂呼亂叫，直到獻晚祭的時候，卻沒有聲音，沒有回應的，也沒有理睬的。 \end{tabularx} \\ \\ \relax
18:30 & \begin{tabularx}{0.7\textwidth}{X} 以利亞對眾百姓說：「你們到我這裡來。」眾百姓就到他那裡，他把那已經毀壞了的耶和華的壇修好。 \end{tabularx} \\ \\ \relax
18:31 & \begin{tabularx}{0.7\textwidth}{X} 以利亞按照雅各子孫支派的數目，取了十二塊石頭；耶和華的話曾臨到雅各，說：「你的名要叫以色列。」 \end{tabularx} \\ \\ \relax
18:32 & \begin{tabularx}{0.7\textwidth}{X} 以利亞用這些石頭為耶和華的名築一座壇，在壇的四圍挖溝，可容納二細亞穀種。 \end{tabularx} \\ \\ \relax
18:33 & \begin{tabularx}{0.7\textwidth}{X} 他又在壇上擺好了柴，把牛犢切成小塊放在柴上，說：「你們用四個桶盛滿水，倒在燔祭和柴上。」 \end{tabularx} \\ \\ \relax
18:34 & \begin{tabularx}{0.7\textwidth}{X} 他又說：「倒第二次。」他們就倒第二次。他又說：「倒第三次。」他們就倒第三次。 \end{tabularx} \\ \\ \relax
18:35 & \begin{tabularx}{0.7\textwidth}{X} 水流到壇的四圍，溝裡也滿了水。 \end{tabularx} \\ \\ \relax
18:36 & \begin{tabularx}{0.7\textwidth}{X} 到了獻晚祭的時候，先知以利亞近前來，說：「耶和華－亞伯拉罕、以撒、以色列的神啊，求你今日使人知道你是以色列的神，我是你的僕人，我遵照你的話做這一切事。 \end{tabularx} \\ \\ \relax
18:37 & \begin{tabularx}{0.7\textwidth}{X} 求你應允我，耶和華啊，應允我，使這百姓知道你－耶和華是神，是你叫他們回心轉意的。」 \end{tabularx} \\ \\ \relax
18:38 & \begin{tabularx}{0.7\textwidth}{X} 於是，耶和華降下火來，燒盡燔祭、木柴、石頭、塵土，又燒乾了溝裡的水。 \end{tabularx} \\ \\ \relax
18:39 & \begin{tabularx}{0.7\textwidth}{X} 眾百姓看見了，就臉伏於地，說：「耶和華是神！耶和華是神！」 \end{tabularx} \\ \\ \relax
18:40 & \begin{tabularx}{0.7\textwidth}{X} 以利亞對他們說：「拿住巴力的先知，不讓任何人逃走！」眾人就拿住他們。以利亞帶他們到基順河邊，在那裡殺了他們。 \end{tabularx} \\ \\ \relax
18:41 & \begin{tabularx}{0.7\textwidth}{X} 以利亞對亞哈說：「你現在可以上去吃喝，因為有暴雨的響聲了。」 \end{tabularx} \\ \\ \relax
18:42 & \begin{tabularx}{0.7\textwidth}{X} 亞哈就上去吃喝。以利亞上了迦密山頂，屈身在地，把臉伏在兩膝之中。 \end{tabularx} \\ \\ \relax
18:43 & \begin{tabularx}{0.7\textwidth}{X} 他對僕人說：「你上去，向海觀看。」僕人就上去觀看，說：「沒有甚麼。」以利亞說：「你再去。」如此七次。 \end{tabularx} \\ \\ \relax
18:44 & \begin{tabularx}{0.7\textwidth}{X} 第七次，僕人說：「看哪，有一小片雲從海裡上來，好像人的手掌那麼大。」以利亞說：「你上去告訴亞哈，當套車下去，免得被雨阻擋。」 \end{tabularx} \\ \\ \relax
18:45 & \begin{tabularx}{0.7\textwidth}{X} 霎時間，天因風雲黑暗，降下大雨。亞哈就坐上車，往耶斯列去了。 \end{tabularx} \\ \\ \relax
18:46 & \begin{tabularx}{0.7\textwidth}{X} 耶和華的手按在以利亞身上，他就束上腰，奔在亞哈前頭，一路到耶斯列。 \end{tabularx} \\ \\
[1ex]
\hline
\hline
\end{longtable}
今天我們要從這三段經文中,思想降火下雨的神蹟.關於這個偉大的神蹟,我們要注意五個特點,五個問題:
\newline
\newline
(一)是心持兩意的問題(王上十八21)「以利亞前來,對眾民說,你們心持兩意要到幾時呢?若耶和華是神,就當順從耶和華；若巴力是神,就當順從巴力.眾民一言不答.」神的忠心僕人以利亞,向以色列人發出挑戰,在一個宗教爭鬥的世代中問他們:「你們心持兩意要到幾時呢?」以利亞這三個字很有意思,就是神是我的力量；這不單是他的名意,也是他的口號.他不怕王后耶洗別所注重的迦南人的神巴力；雖然迦南人強調說,人若不向他們的神禱告,就不會下雨.因為只有巴力才是那地的雨神,天氣好不好,莊稼多不多,都在巴力的主宰中.可是以利亞不承認這個巴力教,也不尊重迦南公主耶洗別與亞哈王結婚時所帶來的四百五十個巴力先知.在屬靈方面,耶洗別對當時的以色列人很危險,她引誘很多人去拜巴力.耶和華便興起以利亞作他的僕人先知,指出選民心持兩意的問題,用降火下雨的神蹟,來證明耶和華是真神.
\newline
\newline
在原文希伯來文中,心持兩意就是像跛子走路一樣.以利亞問以色列人,你們在宗教信仰和生活上,就好像跛子走路一樣辛苦,也沒有甚麼進步,要到幾時呢?這不單是當日的以色列人所遇見的問題,八十年代的信徒也要面對,在我們中間也有耶洗別的影響,以色列人怎樣尊重耶洗別,我們也尊重一些不信主的人.他們有學問,金錢,名譽,才幹,我們就大受他們影響.被他們的人生觀所吸引,結果在基督徒的人生中像跛子走路般；心持兩意,信心搖動,不曉得信耶穌好,還是不信耶穌好.你今天所面對的,是心持兩意的問題嗎?在你屬靈的生命中,有沒有耶洗別的影響,有沒有世界的引誘力?
\newline
\newline
(二)沒有水喝的問題(王下三9)「於是以色列王和猶大王,並以東王,都一同去繞行七日的路程,軍隊和所帶的牲畜沒有水喝.」沒有水喝怎麼辦呢?第十節「以色列王說,哀哉!」可是這樣說有用嗎?遇到不如意的事,不容易解決的問題,你會說哀哉這兩個字嗎?哀哉雖然是自然的反應,但不是屬靈的反應.我們不要效法以色列王,乃要注意先知以利沙怎樣仰望耶和華.(十七,十八節)「因為耶和華如此說,你們雖不見風,不見雨,這谷必滿了水,使你們和牲畜有水喝.在耶和華眼中這還算為小事,他也必將摩押人交在你們手中.」這裡教訓我們,在人沒有力量,沒有辦法的時候,神還能用超自然的方法供給我們的需要.我們看難解決的問題,但在神眼中還算小事.所以,哀哉並不是我們應有的反應,我們應該仰望神,因為他必能幫助我們.
\newline
\newline
從前在印度常有旱災,沒有水喝,所以當地的信徒便定了一個日子,開特別祈禱會,求神下大雨解決他們的問題.日子到了,有幾位教會長者一同來參加禱告會,因為烈日當空,陽光很猛,所以都戴了黑眼鏡.但也有一個小女孩來參加,她不但不戴黑眼鏡,手中還帶著雨傘.長老們笑她,輕看她的信心,問她為甚麼要帶雨傘.她便反問長老們開禱告會的目的,她深信神必聽她的禱告,到時天必要下雨.後來他們進去禱告,果然神改變了天氣,太陽不見了,外面有打雷下雨的聲音.散會時,小女孩拿著她的雨傘歡歡喜喜的回家去.至於那些長老們,黑眼鏡當然沒有用,衣服都濕透了,他們不得不承認自己缺少那小女孩活潑單純的信心,也沒有她屬靈的看法.我們個人方面,團體方面,在遇到困難之時,反應是怎樣呢?
\newline
\newline
(三)滿處挖溝的問題(王下三16)「他便說,耶和華如此說,你們要在這谷中滿處挖溝.」(詩八十四6)「他們經過流淚谷,叫這谷變為泉源之地,並有秋雨之福,蓋滿了全谷.」這兩處經文教訓我們,神雖然願意聽我們的禱告,在特殊的情況下,為仰望祂的人降火下雨,但我們要作適當的準備,與神同工,給他機會來幫助我們.以色列人要在谷中滿處挖溝,要流汗要出力,接受神的預備.同樣在詩篇那裡,他們經過流淚谷的時候,就是一片枯乾,很少下雨的曠野地方.有時會裂開,也會有旋風吹過,雖然神會預備秋雨之福；但在經文中有一個很有意義的字,就是「叫」字.他們經過流淚谷,叫這谷變為泉源之地.一個好的基督徒不是被動的,他不冷淡,也不懶惰,靠著主他願意出力,叫人生的流淚谷變為泉源之地,這樣應付人生的難處,人生的重擔和壓力.
\newline
\newline
到底(詩八十四6)的「叫」字和(王下三16)的「滿處挖溝」,有甚麼屬靈的教訓和應用呢?相信在讀經禱告方面,神很願意為我們下雨,把他的恩雨靈雨,豐豐富富的賜給我們,滋潤我們的生活,灌溉我們的靈性；使我們的靈性生活活潑有力,滿有信心.但我們要先作準備,藉著每天讀經禱告,滿處挖溝；這樣,叫人生的流淚谷變為泉源之地,享受秋雨之福.請問你有把讀經禱告養成習慣嗎?你有預備自己接受主恩嗎?
\newline
\newline
有一天,美國加州一位八十四歲的老人家,站起來作見證,他說我感謝神保守我直到今天,我更感謝神給我力量把全本聖經從頭到尾讀了一百六十一次.這五十五年來,我只能用一隻眼睛看聖經,因為另一隻瞎了；但是主的恩典是多方面的；我雖然讀了聖經這樣多,這樣久,可是每逢拿起聖經讀時,屬靈的生命重新得力.有新的功課,有新的靈糧.這位八十四歲的老弟兄的見證,不是很有挑戰性嗎?我們中間把聖經從頭到尾讀了一次的人,有多少呢?求主把每天讀經禱告的功課,銘刻在我們的心版上；使我們永不忽略,永不忘記.
\newline
\newline
(四)如此七次的問題(王上十八41-44)「以利亞對亞哈說,你現在可以上去吃喝,因為有多雨的響聲了.亞哈就上去吃喝.以利亞上了迦密山頂,屈身在地；將臉伏在兩膝之中,對僕人說,你上去,向海觀看.僕人就上去觀看說:沒有甚麼.他說你再去觀看,如此七次.第七次僕人說:我看見有一小片雲從海裡上來,不過如人手那樣大.以利亞說:「你上去,告訴亞哈,當套車下去,免得被雨阻擋.」這裡講到先知以利亞和亞哈王,一個很有興趣的對比.天空一片雲也沒有,但以利亞卻對亞哈說,有多雨的響聲了,以利亞怎麼能聽見雨的響聲呢?他屬靈的聽覺怎麼這樣靈敏呢?這裡有一個寶貴的屬靈功課,就是以利亞緊緊的跟隨主,已經訓練自己屬靈的耳朵,聽清楚主的聲音,明白神降火下雨的時間.求主幫助我們,訓練自己屬靈的耳朵,聽清楚主的聲音,與神同心同工.
\newline
\newline
可是還有如此七次的問題,是以利亞禱告的時候所受的試煉.第四十二節描寫以利亞禱告的姿態,他謙卑懇切,屈身在地,將臉伏在兩膝之中；四十三節便講到他禱告的結果.他僕人上去觀看的時候,看不見甚麼.許多時候,我們謙卑迫切的禱告,也得不著甚麼,也許這正是我們要學的如此七次的功課.可是以利亞並不灰心,也不停止禱告,他叫僕人再去觀看,如此七次.
\newline
\newline
弟兄二十多年前在劍橋大學唸書時,基督徒團契中有一位很愛主的顧問阿金遜博士,有一次他在聚會中作見證,為兩件事感恩.頭一件是他已經信主五十年,今天是他蒙恩的日子.難怪他對同學的幫助那麼大,那麼深.另一件是他今天收到一封信,是一位才信主的同事寄給他的；他為這同事已經禱告了二十五年,沒有一日間斷,天天為他祈求.頭二十年的時間中,他總是施反對和輕視福音的態度.這五年才漸漸對福音發生興趣,後來終於信了主.阿金遜博士的見證,正是如此七次的問題；禱告若不馬上見果效,千萬不要灰心,還要如此七次.
\newline
\newline
最後,降下大雨的問題.(王上十八45-46)「霎時間天因風雲黑暗,降下大雨.亞哈就坐車,往耶斯列去了.耶和華的靈降在以利亞身上,他就束上腰,奔在亞哈前頭,直到耶斯列的城門.」神任何屬靈的工作,雖然通常有微小的開始,就好像起初只有一小片雲,但繼續禱告的時候,小能變大.在這裡我們看見神工作的三個階段:第一,沒有甚麼.第二,一小片雲.第三,降下大雨.因此我們不要灰心,要一面禱告,一面作準備.後來神的靈降在以利亞身上,他就束上腰,重新得力,奔在亞哈的前頭.從迦密山的山頂,一千七百三十呎高；一直到耶斯列的城門,有二十哩之遠.他跑得比馬車更快,確實是個神蹟.願主在這個時代中,一樣感動充滿我們.使我們在教會的工作上,多準備束上腰,這樣在神降下大雨的時候,才能像以利亞那樣.快快的跑到耶斯列.在原文上,耶斯列有撒種子的意思.是個撒種子的地方.我們在福音的工作上,能來到這個地方,就能被聖靈充滿,能被天父使用.
\newline
\newline


\newpage

\section{第53屆~港九培靈會~林克己~第五講}
\label{sec:585}
\textbf{靠主喜樂的神蹟}
\newline
\newline
連結: \href{https://www.hkbibleconference.org/session-message/view/585}{\texttt{https://www.hkbibleconference.org/session-message/view/585}}
\newline
\newline
\hyperref[sec:583]{< < < PREV SERMON < < <}
~
\hyperref[sec:toc]{[返目錄]}
~
\hyperref[sec:586]{> > > NEXT SERMON > > >}
\newline
\newline
哈巴谷書 3:16-3:19
\newline
\begin{longtable}{cl}
\hline
\hline
章節 & 經文 (和合本修訂版)\\
\hline
3:16 & \begin{tabularx}{0.7\textwidth}{X} 我聽見這聲音，身體戰兢，嘴唇發顫，骨中朽爛，在所立之處戰兢；但我安靜等候災難之日臨到那上來侵犯我們的民。 \end{tabularx} \\ \\ \relax
3:17 & \begin{tabularx}{0.7\textwidth}{X} 雖然無花果樹不發旺，葡萄樹不結果，橄欖樹也不收成，田地不出糧食，圈中絕了羊，棚內也沒有牛； \end{tabularx} \\ \\ \relax
3:18 & \begin{tabularx}{0.7\textwidth}{X} 然而，我要因耶和華歡欣，因救我的神喜樂。 \end{tabularx} \\ \\ \relax
3:19 & \begin{tabularx}{0.7\textwidth}{X} 主耶和華是我的力量，他使我的腳快如母鹿，又使我穩行在高處。這歌交給聖詠團長，用絲弦的樂器。 \end{tabularx} \\ \\
[1ex]
\hline
\hline
\end{longtable}
我們先來看看靠主喜樂的神蹟,怎樣在使徒保羅的生活中顯現出來.(腓四4)「你們要靠主常常喜樂.我再說,你們要喜樂.」保羅寫這節經文時,從人看來,非常可憐.環境可怕,前途黑暗.他因為為主作見證,坐監牢,受捆鎖,等候在羅馬王前受審,審訊結果可能會被釘死,為主殉道.他能這樣勉勵別人,必定要靠主喜樂,這真是一件奇妙超自然的事.明顯地,保羅的喜樂,並不是靠環境,而是有一個屬靈天上的來源.保羅因為靈性健康,信心堅固,甚麼都不怕也不問；深知無論是生是死,信徒都有一個光明的前途；因為本身的罪已得赦免,必定可享受永遠的福氣.按照教會的傳說,保羅在監牢中喜樂的經歷,大有吸引力.看守他的兵,多有歸向主,他們的軍官不曉得怎辦,因為這些人不能讓他們長時間和保羅在一起,免得他們相信耶穌,像保羅一樣不停的傳揚福音.如果我們查考腓立比書,就會看見保羅常提到他在主裡面的喜樂,所以有基督徒提議這書信應該稱為喜樂的書信.環境無論如何惡劣,前途無論如何黑暗,仍然在主裡有喜樂.
\newline
\newline
現在看舊約(哈一1)「先知哈巴谷所得的默示.」哈巴谷是個靠主喜樂的人,他在這裡自稱先知.在舊約所有的先知中,只有哈巴谷一個這樣做；自稱先知,注重自己的身份,這並不是驕傲的表現.他在當時靈性退步的世代中,在主前六百年左右,但卻清楚神的呼召,又自稱先知.在(三18)作見證說:「然而我要因耶和華歡欣,因救我的神喜樂.」因此第一章一節提醒我們,信徒要清楚神在他身上的要求,謹寫神在他們身上所定的旨意,才能心滿意足,過一個靠主喜樂的生活.神要我作甚麼,我就作甚麼,主要是神的旨意,安排.如果我們清楚自己在社會中的身份,和在教會中的職份,都是神為我們揀選；是見證他的範圍和機會,我們便能心滿意足,靠主喜樂.
\newline
\newline
哈巴谷這名字有甚麼意思呢?在聖經上,人的名字,通常和他們的性情品格和見證有關.譬如說:司提反是冠靈的意思,為主殉道,就得到生命的冠冕.俄巴底亞是敬拜事奉耶和華.哈巴谷有兩個意思:一是雙手抱住,二是一根植物或一朵花.我們查考哈巴谷書之時就看見,這位偉大的先知對神的態度,真是雙手抱住.雖然在他的時代中,靈性退步的人很多,但哈巴谷仍是愛他的主,願意順從服事主.另一方面,他也與神摔交,好像雅各一樣,夜間與神摔交.哈巴谷在他所寫的書中,提到一些他心裡面的問題,一些他不明白的屬靈的事；他就願意在禱告和默想中與神摔交,得到問題的答案.這為我們留下一個美好的榜樣,在我們心中,若有問題的時候,不要怕在主面前提出來；同時間在主面前等候,藉著默想禱告與神摔交,為要得著答案.雙手抱住有這兩個意思:表示愛心和與神摔交.
\newline
\newline
哈巴谷的第二個意思,就是一根植物或是一朵花.相信在這裡有挑戰,教訓和安慰.有時基督徒卑賤,好像一根植物一樣,學問,錢財,名譽都沒有；可是在生活見證上,卻美如一朵花,用聖潔的行為,發出香氣,一種領人歸主的香氣,將榮耀歸給神.
\newline
\newline
現在我們看哈巴谷先知所默想的五個問題,雖然內心有這些問題,但仍靠主喜樂.哈巴谷關於神不明白的五件事.
\newline
\newline
(一)神的方法(一6)「我必興起迦勒底人,就是那殘忍暴躁之民,通行遍地,佔據那不屬自己的住處.」這裡提到神管教他的選民所用的方法；祂要興起迦勒底人,攻打他們,藉著戰爭使他們悔改.承認神的聖潔,和絕對的主權.可是這在哈巴谷心中成為一個問題；請注意他怎樣描寫迦勒底人.他說他們殘忍暴躁,這也是神所承認的.哈巴谷就問,神公平嗎?雖然神的子民應該受管教,受刑罰,但神使用這些比以色列民更壞的迦勒底人,做祂的工具,這個公平合理嗎?神所用的方法,不是不公道嗎?哈巴谷不明白神的方法,便在禱告中與神摔交,要得到神的啟示和答案.然而神的意念高過我們的意念,他所用的方法器皿,真是奇妙.在教會歷史上例子很多.十八世紀約翰衛斯理傳福音的時候,有一次,有人帶一個無神論者參加衛斯理所主持的一個露天佈道會；無神論者對衛斯理所講的不感興趣,便用雙手塞著耳朵.突然飛來了一隻蒼蠅,停留在無神論者的鼻子上；他雖然搖頭,但是蒼蠅卻不飛走.無奈他痕癢非常,只得用塞著耳朵的手去撥開蒼蠅.就在這時,衛斯理便引用啟示錄的一句話說:「凡有耳的,就應當聽.」無神論者就被這句說話吸引著,很留心的聽下去,結果受感動歸向主.如果神能使用一隻蒼蠅,在哈巴谷的時代中,使用殘忍暴躁的迦勒底人；那麼,神在我們的時代中,更願意使用愛祂的人了.只要我們對主的愛雙手抱住,雖然卑賤,但亦能成為神使用的方法和器皿；把福音傳到地極,享受服事主的喜樂.
\newline
\newline
(二)神的靜默(一13)「你眼目清潔不看邪僻,不看奸惡,行詭詐的,你為何看著不理呢?惡人吞滅比自己公義的,你為何靜默不語呢?」無可否認的,世界上不公平的事很多,不單是哈巴谷的時代,在我們的時代中,也有惡人吞滅義人的事；然而滿有慈愛能力的神,為甚麼容許這些事存在?為甚麼看著不理,不除掉這些事呢?這是哈巴谷內心的問題,也是今日很多人的問題.神的靜默怎樣解釋呢?哈巴谷在與神摔交之後,得到這問題一部份的答案；就是聖潔公道的神,必刑罰邪惡,賞賜良善.(二14)「認識耶和華榮耀的知識,要充滿遍地,好像水充滿海洋一般.」我們不要因為神的靜默而疑惑祂的關心和慈愛.信心之父亞伯拉罕禱告時提到一個審判的原則；「他對耶和華說,審判全地的主,豈不行公義嗎?」聖經告訴我們,神在耶穌再來之時,要設立祂審判萬民的大寶座,使萬民得到應得的報應；使罪人無可推諉,不得不承認祂的聖潔公義.神的靜默,必須從聖經所定的觀點來看,才能明白.
\newline
\newline
先知在(一13)提到神的靜默,不單在審判方面,也在聽禱告方面.哈巴谷雖然在代禱中,多為神的選民求恩；但是神好像睡覺,不動手幫助；也不禁止惡人吞滅他們,難怪哈巴谷兩次的問「為甚麼」.神的靜默實在是一個很大的問題.不過無論發生甚麼,我們都不要疑惑神的關心,慈愛和無所不知.幾十年以前,英國和德國打仗的時候,有兩位基督徒從倫敦上船要到南非去,他們多禱告,求主一面保佑他們,一面為他們預備適合的天氣,使他們一路順風.誰知他們愈多禱告,天氣便愈壞,一直到他們抵達南非的時候,風雨仍繼續不停.這兩位基督徒很不明白,上船的時候便對船長說出自己的感受；船長不是基督徒,但聽完之後卻說.不要讓你們的信心這樣快搖動.你們的耶穌沒有睡覺,這次的大風雨,正是你們神慈愛和能力的明證,從英國到這裡,差不多全路都有一隻德國戰船在跟著我們,但是神藉著暴風雨保守了我們.所以我們要記得,不要誤解神的靜默.
\newline
\newline
(三)神的時間(二1-3)「我要站在守望所,立在望樓上觀看,看耶和華對我說甚麼話,我可用甚麼話向他訴冤.他對我說,將這默示明明的寫在版上,使讀的人容易讀.因為這默示有一定的日期,快要應驗,並不虛謊,雖然遲延,還要等候.因為必然臨到,不再遲延.」這段經文有個屬靈的原則,神甚麼默示,甚麼計劃,都有一定的時間.我們需要有第三節的態度,雖然遲延,仍要等候,要忠心的事奉,等候神動工的時間.
\newline
\newline
兄弟在馬來西亞宣教的時候,也有機會學這個功課.我四年和一位年老的女傳道一同牧養一個青年團契,頭兩年的工作很難,沒有甚麼明顯的果效,我們只能抓住哈巴谷書二章三節的應許,將這原則應用在我們的工作上.雖然遲延,還要等候,因為必然臨到,不再遲延.兩年後,一位十六歲的女孩信主,之後她又帶領她的姐姐和弟弟信主.後來的兩年中,常有年青人悔改信主,並受洗禮.因此我們不必灰心,神的計劃,默示,引導,都有一定的時間,我們應該等候信任他.
\newline
\newline
(四)神的作為(三2)「耶和華阿,我聽見你的名聲就懼怕.耶和華阿!求你在這些年間復興你的作為；在這些年間顯明出來,在發怒的時候,以憐憫為懷.」這問題是現代很多信徒所面對的.如果神滿有慈愛,偉大無比,大有力量,他為甚麼不早日復興他的作為呢?如果他關心教會,關心社會,為甚麼教會工作這麼冷淡,不發生功效?祂為甚麼不提醒我,復興我呢?這是一個難以明白的問題.然而一個屬靈的原則就是,神工作的發展,是在乎祂絕對的主權.祂甚麼時候工作,用甚麼方法,用甚麼人,這都是衪的權力.神的時間和作為是有連帶關係的.我們的責任只是忠心事奉,忍耐等候,懇切禱告,雖然遲延,還要等候,學習哈巴谷一樣.
\newline
\newline
(五)神的大能(三17-19)「雖然無花果樹不發旺,葡萄樹不結果,橄欖樹也不效力,田地不出糧食,圈中絕了羊,棚內也沒有牛.然而我要因耶和華歡欣,因救我的神喜樂.主耶和華是我的力量,他使我的腳快如母鹿的蹄,又使我穩行在高處.」這段經文描寫一件神蹟,就是靠主喜樂的神蹟.哈巴谷在他各方面的工作上,雖然還未看見神的大能和榮耀,但是他在日常事奉的生活中,已經得了得勝的秘訣.雖然在他心內,有五個不明白的問題,但是他有一個態度,就是因耶和華歡欣,快樂.原文的歡欣和喜樂都很有意思,如果按照原文來翻譯,就是「然而我要因耶和華歡欣就跳起來；因救我的神賜給我的喜樂就旋轉.」這是屬靈的跳舞,因為心裡快樂,不得不跳起來旋轉.因此,我們喜樂,不是因為環境,乃是因為有主的力量在我們身上,有主的快樂在我們心裡.願主使這靠主喜樂的神蹟,成為我們每一個人屬靈的經歷.
\newline
\newline


\newpage

\section{第53屆~港九培靈會~林克己~第六講}
\label{sec:586}
\textbf{五餅二魚的神蹟}
\newline
\newline
連結: \href{https://www.hkbibleconference.org/session-message/view/586}{\texttt{https://www.hkbibleconference.org/session-message/view/586}}
\newline
\newline
\hyperref[sec:585]{< < < PREV SERMON < < <}
~
\hyperref[sec:toc]{[返目錄]}
~
\hyperref[sec:588]{> > > NEXT SERMON > > >}
\newline
\newline
馬可福音 6:30-6:44
\newline
\begin{longtable}{cl}
\hline
\hline
章節 & 經文 (和合本修訂版)\\
\hline
6:30 & \begin{tabularx}{0.7\textwidth}{X} 使徒們聚集到耶穌那裡，把一切所做的事、所傳的道全告訴他。 \end{tabularx} \\ \\ \relax
6:31 & \begin{tabularx}{0.7\textwidth}{X} 他就說：「你們來，同我私下到荒野的地方去歇一歇。」這是因為來往的人多，他們連吃飯的時間也沒有。 \end{tabularx} \\ \\ \relax
6:32 & \begin{tabularx}{0.7\textwidth}{X} 他們就坐船，私下往荒野的地方去。 \end{tabularx} \\ \\ \relax
6:33 & \begin{tabularx}{0.7\textwidth}{X} 眾人看見他們走了，有許多認識他們的，就從各城步行，一同跑到那裡，比他們先趕到了。 \end{tabularx} \\ \\ \relax
6:34 & \begin{tabularx}{0.7\textwidth}{X} 耶穌出來，見有一大群的人，就憐憫他們，因為他們如同羊沒有牧人一般，於是開始教導他們許多事。 \end{tabularx} \\ \\ \relax
6:35 & \begin{tabularx}{0.7\textwidth}{X} 天已經很晚，門徒進前來，說：「這地方偏僻，而且天已經很晚了， \end{tabularx} \\ \\ \relax
6:36 & \begin{tabularx}{0.7\textwidth}{X} 請叫眾人散去，他們好往四面的鄉鎮村莊去，自己買些東西吃。」 \end{tabularx} \\ \\ \relax
6:37 & \begin{tabularx}{0.7\textwidth}{X} 耶穌回答他們說：「你們給他們吃吧！」門徒對他說：「我們要拿兩百個銀幣去買餅給他們吃嗎？」 \end{tabularx} \\ \\ \relax
6:38 & \begin{tabularx}{0.7\textwidth}{X} 耶穌說：「你們有多少餅？去看看。」他們知道後就說：「有五個，還有兩條魚。」 \end{tabularx} \\ \\ \relax
6:39 & \begin{tabularx}{0.7\textwidth}{X} 耶穌吩咐他們，叫眾人一組一組地坐在青草地上。 \end{tabularx} \\ \\ \relax
6:40 & \begin{tabularx}{0.7\textwidth}{X} 眾人就一群一群地坐下，有一百的，有五十的。 \end{tabularx} \\ \\ \relax
6:41 & \begin{tabularx}{0.7\textwidth}{X} 耶穌拿著這五個餅和兩條魚，望著天祝福，擘開餅，遞給門徒，擺在眾人面前，也把那兩條魚分給眾人。 \end{tabularx} \\ \\ \relax
6:42 & \begin{tabularx}{0.7\textwidth}{X} 他們都吃，並且吃飽了。 \end{tabularx} \\ \\ \relax
6:43 & \begin{tabularx}{0.7\textwidth}{X} 門徒把餅和魚的碎屑收拾起來，裝滿了十二個籃子。 \end{tabularx} \\ \\ \relax
6:44 & \begin{tabularx}{0.7\textwidth}{X} 吃餅的男人共有五千。 \end{tabularx} \\ \\
[1ex]
\hline
\hline
\end{longtable}
今天我們要查考耶穌所行五餅二魚的神蹟,記載在新約翰福音第六章一到十五節.我們要從其中注意五件事,主在五方面怎樣作我們的榜樣.
\newline
\newline
(一)耶穌的關心(六1-5)「這事以後,耶穌渡過加利利海,就是提比哩亞海.有許多人,因為看見他在病人身上所行的神蹟,就跟隨祂.耶穌上了山,和門徒一同坐在那裡.那時猶太人的逾越節近了.耶穌舉目看見許多人來,就對腓力說:「我們從哪裡買餅,叫這些人吃呢?」這段經文兩次告訴我們,主怎樣關心別人.他們無論在哪一方面有需要；生病的要醫治,肚子餓的要食物,被罪捆綁的要得赦免,耶穌就用自己的大能大力,為他們行神蹟,想辦法.我們在今天的社會中,要多留意耶穌這美好的榜樣.耶穌和現代許多人是相反的,他雖然是主,但不求自己的利益,不想到自己的方便和舒服,這是基督化的品格應該有的特色.信主之前,我們很自然會自私自利,很少想到別人的益處,以自己的需要為重要.但信主之後,我們這種觀念便應改變,自私要變為關心,要學習看見別人的需要,要訓練自己求別人的益處.
\newline
\newline
我們在聖經中看看兩種人生觀的比較.(傳二4)這裡有所羅門王為自己利益所作的見證,「我為自己動大工程,建造房屋,栽種葡萄園.」第八節:「我又為自己積蓄金銀.」請注意「為自己」這幾個字,這是基督徒應該避免的自私的人生觀.而(腓一21-24)為見證主的緣故而坐監的保羅,描寫他的人生觀,他對神對人所存的態度,和他對死的觀念.他是不怕死的,因為離世與基督同在是好得無比的.二十一節他見證說,因我活著就是基督.在保羅的生活中,基督是居首位的,保羅要跟從祂,順服祂,傳揚祂.保羅得想早日離開世界上天堂,永遠與主在一起；但是他卻先求別人的益處.所以他說:「然而我在肉身活著,為你們更是要緊的.」傳道書是「為自己」,這裡是「為你們」.在保羅日常的生活中有一個次序,就是基督第一,別人第二,自己第三.這是真正的基督化人生觀,和舊約聖經所羅門王為自己打算的態度大有分別.願主幫助我們,把關心別人養成我們的習慣,好顯明主吸引人的大愛.
\newline
\newline
(二)耶穌的計劃(約六5-7)「耶穌舉目看見許多人來,就對腓力說,我們從哪裡買餅,叫這些人吃呢?他說這話,是要試驗腓力,他自己原知道要怎樣行.腓力回答說,就是二十兩銀子的餅,叫他們各人吃一點,也是不夠的.」耶穌供給人的需要,有一個完整的計劃.他雖然問腓力的意見,但是他原知道自己要怎樣行.我們在人生的道路上與耶穌同行,在事奉的生活中與耶穌同工,衪能夠預備一個美好的計劃,適合我們各方面的需要.我們在祂面前等候禱告,他就會使我們知道他的旨意和安排如何.同樣,這也教導我們,信徒服事主不要糊裡糊塗,要有計劃像主一樣.
\newline
\newline
跟著我們注意第六節的另一件事,就是耶穌對腓力的態度.他是耶穌的門徒,但是他不單要與主同工,在工作上被主使用,更要受試驗.這裡絕錄耶穌怎樣訓練腓力,上半節說:「他說這話,是要試驗腓力.」現代的基督徒服事主,遇見困難不要以為奇怪,這是主訓練和試驗我們的好方法,提醒我們要在凡事上仰望他,而且學習一個重要的功課,就是主使用誰便試驗誰,這樣使我們成為更合符祂使用的僕人.
\newline
\newline
我上一次回國休假時,因大女兒生病不能準時回港,便暫時作英國南部一個小教會的牧師.有一天,一對愛主的夫婦,邀請我到他家開佈道會,他們是別個教會的會友.當我到他們家之時,覺得很奇怪,一間大房子內共有九十七人來聚會,這對夫婦事先宣佈,講員講道之前,先給他們看一些香港的幻燈片,這樣吸引他們來.當晚我呼召時,有五個人決志.其中有個婦人,他的兩位朋友邀請她來聽道之前,已經多為他代禱.婦人歸向主之後,有人告訴我,一年以前,她很盼望和一位愛主的弟兄結婚；但那弟兄堅守聖經的原則,不肯答允.可是婦人不灰心,繼續等候,希望他改變主意.但有一天這位弟兄走路時,心臟病發便去世了,婦人非常痛苦,在這個時候,邀請她參加佈道會的兩位姊妹便和她做朋友,整天為她禱告,打電話給她,勉勵安慰她,傳福音給她；但是婦人不聽,不過兩位姊妹不灰心,仍一樣的探望照顧她,結果她們的恆心得主的賞賜,婦人願意接受她們的邀請,當晚流淚悔改接受主.因此我們不要怕主的試驗,因為試驗以後,主必叫我們看見果效.
\newline
\newline
(三)耶穌的同工主很歡迎人來與他同工,在這個神蹟中,不單有腓力和其他門徒,也有帶五餅二魚的小童；他雖然年小,但耶穌願意使用他.我們在這方面要效法主的榜樣,學習怎樣與別人同工.舊約時代的所羅門,也用智慧提到這一點,(傳四9)「兩個人總比一個人好,因為二人勞碌同得美好的果效.」信徒一同工作能被主使用,我們在教會中要彼此尊重,承認別人的恩賜經驗和貢獻,同心合意的興旺福音.主怎樣和信祂的人分擔責任,是值得我們默想的模範.
\newline
\newline
兄弟二十多年前,在英國東邊的一間小教會講道,會畢有個十歲的女孩子到台前來和我談談,她說她名叫瑪加列,她很想接受耶穌為救主,不知自己是否太小.我還記得她的頭髮很美,有兩條長長的辮子.我們用幾分鐘時間談道之後,她便低下頭來,誠心的接受主.兩個星期之後,她又來和我談話,她說她已經傳揚福音給三十位同學聽,其中兩人還肯接受救主.她們每天下課便到我家,我們一同唱詩讀經禱告,十分高興.信主只有兩個星期,已經作了這樣美好的工作.兩年之後,有一天我在附近另一所禮拜堂領會,我便邀請瑪加列和我一同去,請她獨唱,因為她唱得很好.可是她聽完之後,卻提出另一個條件；她說她除了獨唱之外,還希望有十分鐘時間講見證.結果,那天的聚會開得很成功,很多人被她的見證所感動,神大大的使用她.所以,「兩個人總比一個人好」,我們要學習與別人同工,和弟兄姊妹同心的服事主.雖然耶穌是全能的主,他原知道自己要怎樣行,然而他卻樂與腓力,門徒,並帶五餅二魚的小孩同工.
\newline
\newline
(四)耶穌的祝謝(約六8-14)「有一個門徒,就是西門彼得的兄弟安得烈,對耶穌說,在這裡有一個孩童,帶著五個大麥餅,兩條魚,只是分給這許多人,還算甚麼呢?耶穌說:你們叫眾人坐下.原來那地方的草多,眾人就坐下,數目約有五千.耶穌拿起餅來,祝謝了,就分給那坐著的人,分魚也是這樣,都隨著他們所要的.他們吃飽了,耶穌對門徒說,把剩下的零碎,收拾起來,免得有蹧蹋的.他們便將那五個大麥餅的零碎,收拾起來,裝滿了十二個籃子.眾人看見耶穌所行的神蹟,就說,這真是那要到世間來的先知.」在這段經文裡,耶穌的祝謝可以作我們的榜樣.他雖然承認這五餅二魚是人帶來的.但同時也是天父的預備,所以要獻上祂的感謝.保羅在他的書信中也以四個字來勉勵我們,就是「凡事謝恩」.五個餅,兩條魚,一個孩童,能算甚麼呢?可是在耶穌的祝謝後,他不單能吃飽五千人,更有剩下的.可惜有基督徒疑惑這神蹟,他們說在耶穌祝謝時,大家都拿出自己所帶來的食物吃.這雖然好像很有理由,但他們對聖經的記載和耶穌的神性都有錯誤的觀念.聖經說這五千人整天和耶穌在一起,因此到晚上他們所帶來的食物都已經吃光了；還有耶穌的神性,聖經高舉他是神全能的兒子,在他無難成的事.他有超自然的能力,聖經所載的任何神蹟都十分可信.
\newline
\newline
耶穌行這個神蹟還有一個目的,就是教訓我們,你願意奉獻甚麼,主願意使用甚麼.孩童雖然只得五個餅,兩條魚,算不得甚麼,但已經是他的一切所有；當他願意甘心獻上的時候,主就能使他餵飽五千人.所以,你願意奉獻甚麼,主便願意使用甚麼.請看一個屬靈的原則,寶貴的應許,(林後八12)「因為人若有願作的心,必蒙悅納,乃是照他所有的,並不是照他所無的.」主所要的,不是我們的恩賜,而是我們自己；他的要求是照我們所有的,而不是照我們所無的.我們獻上自己,他便使用我們.
\newline
\newline
(五)耶穌的謙卑(約六15)「耶穌既知道眾人要來強逼他作王,就獨自又退到山上去了.」耶穌雖然有機會加冕作王,但他知道要拯救世人,就必須掛在十字架上,再無別的方法.所以他就躲避那些想要他作王的人,在山上把自己隱藏起來.祂有行神蹟的大能,也有機會作主,但他卻願走十字架的道路.主的選擇和謙卑,值得我們默想.
\newline
\newline
關於救人的好消息,主要的是十字架,我們領人歸主,也要高舉十字架,耶穌的寶血和復活.難怪保羅說:「我曾定了主意,在你們中間不知道別的,只知道耶穌基督,並他釘十字架.」這是我們傳福音的時候所必須堅守的原則,不是傳自己,乃是傳基督耶穌和他的十字架.
\newline
\newline


\newpage

\section{第53屆~港九培靈會~林克己~第七講}
\label{sec:588}
\textbf{耶穌履海的神蹟}
\newline
\newline
連結: \href{https://www.hkbibleconference.org/session-message/view/588}{\texttt{https://www.hkbibleconference.org/session-message/view/588}}
\newline
\newline
\hyperref[sec:586]{< < < PREV SERMON < < <}
~
\hyperref[sec:toc]{[返目錄]}
~
\hyperref[sec:590]{> > > NEXT SERMON > > >}
\newline
\newline
馬太福音 8:23-8:27
\newline
\begin{longtable}{cl}
\hline
\hline
章節 & 經文 (和合本修訂版)\\
\hline
8:23 & \begin{tabularx}{0.7\textwidth}{X} 耶穌上了船，門徒跟著他。 \end{tabularx} \\ \\ \relax
8:24 & \begin{tabularx}{0.7\textwidth}{X} 海裡忽然起了猛烈的風暴，以致船幾乎被波浪淹沒，耶穌卻睡著了。 \end{tabularx} \\ \\ \relax
8:25 & \begin{tabularx}{0.7\textwidth}{X} 門徒去叫醒他，說：「主啊，救命啊，我們快沒命啦！」 \end{tabularx} \\ \\ \relax
8:26 & \begin{tabularx}{0.7\textwidth}{X} 耶穌說：「你們這些小信的人哪，為甚麼膽怯呢？」於是他起來，斥責風和海，風和海就大大平靜了。 \end{tabularx} \\ \\ \relax
8:27 & \begin{tabularx}{0.7\textwidth}{X} 眾人驚訝地說：「這是怎樣的一個人？連風和海都聽從他。」 \end{tabularx} \\ \\
[1ex]
\hline
\hline
\end{longtable}
今天我們要查考的經文,是馬太福音十四章二十二到三十三節,這是耶穌和彼得一同履海的神蹟.他們因著有從上而來的能力,所以能勝過地心吸力的力量；我們特別要從彼得身上,學習五樣寶貴的功課.
\newline
\newline
(一)信心的異象雖然環境危險可怕,有黑夜狂風巨浪,但是彼得倚靠主,甚麼都不怕,他對主說:「請叫我從水面上走到你那裡去.」這是信心的異象.走在水面上,雖然是人類從未做過的,雖然從科學的觀點來看是不可能；但是信心比任何定律更大,仰望耶穌的時候,甚麼都能作.(徒二16)講到彼得在五旬節聖靈降臨時怎樣提到舊約聖經先知約珥的預言.他說:「這正是先知約珥所說的,神說,在末後的日子,我要將我的靈澆灌凡有血氣的；你們的兒女要說預言,你們的少年人要見異象,老年人要作異夢.」被聖靈充滿,被基督得著,就能見異象,去作人類和教會從來沒有作過的事.願主幫助我們在這二十世紀的時代中,好像彼得一樣見異象,一種服事主傳福音的異象,勇敢的往前走,作見證時有信心所發的膽量.
\newline
\newline
第二次世界大戰的時候,英國有一位很有才幹的科學家,他想辦法發明一種特別的地震彈,這種極具爆炸力的地震彈,要打破德國水塘的牆壁,這樣使水流出來,好攔阻德國人.這位科學家腦中有這個異象,有這個計劃,便在自己的花園內,用一盆水和一個塑膠袋來試驗.因為他想像中的地震彈,要在水面上跳幾次,才可碰到水塘的牆壁.他在自己花園中試驗時,很有膽量,也不怕人取笑誤會他,因為有人經過的時候停下來看,而且大聲笑,看不起這位科學家,說他在玩耍.然而後來地震彈作成了,這些人只得承認他們的錯.同樣在我們的世代中,也有人誤會我們傳福音的熱誠,以為我們糊塗,不明白我們服事主的異象.然而到主再來之時,萬膝都要跪拜他.
\newline
\newline
(二)信心的對象(太十四28上)「彼得說,主,如果是你.」如果叫他在海面上行走的,是這位全能的主,他就甘心樂意,坦然無懼的去做.彼得已經見過耶穌平靜風浪,也聽過耶穌對風和海所說的話.(路八24)「門徒來叫醒了祂說:夫子,我們喪命喇.耶穌醒了,斥責那狂風大浪,風浪就止住,平靜了.」這裡說耶穌斥責風浪,馬太福音第八章也是這樣說,但在四本福音書中,只有馬可告訴我們耶穌對風和海說話時所用的字,原因是馬可的主題是耶穌作神的僕人,耶穌的說話行事都有能力.我們看(可四39)「耶穌醒了,斥責風,向海說,住了吧,靜了吧.風就止住,大大的平靜了.」請注意「住了吧,靜了吧」這六個字.原文這幾個字很有意思,耶穌所選用的希臘文動詞,是一個人對自己所養的狗所說的話:祂吩咐風浪,像人吩咐一隻狗一樣.按原文的意思,耶穌是說:「不要跳起來,要戴口套.」耶穌能吩咐風浪,改變天氣,好像人吩咐他所養的狗一樣容易.彼得看見了耶穌所作的事,信心得堅固；所以後來他想要在海面上行的時候,他只問一件事,「主,如果是你.」這是信心的對象,因為彼得在海面唯一所能倚靠的就是主.耶穌沒有叫船上的人預備繩子,好隨時把下沉的彼得拉上來.請問如果主今天呼召你做一件事,你看見了異象,你敢不敢去做?或者是作教會的義工,幫忙教會的青少年工作,或者是作全時間的傳道人,或是宣教師,在本地或海外工作,請問你願意嗎?巴不得我們都對主說:「主,如果是你,我便願意站在使徒保羅信心的立場上；相信靠著那加給我的力量的主,凡事都能作.」
\newline
\newline
(三)信心的搖動當彼得定睛注視耶穌,倚靠主一步一步往前走的時候,黑夜狂風巨浪都不成問題.但是第三十節告訴我們,只因見風甚大,就害怕,將要沉下去.這裡關於信心有一個基本的問題,就是基督徒都有兩隻眼睛,肉身的和屬靈的；我們用屬靈的眼睛注視仰望耶穌,就能經過任何黑夜,能抵擋任何大風,應付任何巨浪,過得勝的屬靈生活.但是我們一用肉身的眼睛看環境,看困難看人,像彼得一樣見風甚大,就必定沉下去,不再在海面上勇敢的行走.彼得為甚麼停步看環境呢?或者他不滿意自己的進步,太慢了.但是我們不要忘記龜免賽跑的故事,免子雖然起步快,但因為中途睡覺,所以終於比不上一步一步慢行的烏龜.同時,在屬靈的路途中,雖然有人有好的開始,起步很快,但過了一段日子卻冷淡下來,不再往前.我們要求主給我們有倚靠衪的心,一步一步的往前走,一生跟隨主.
\newline
\newline
在聖經裡面記載的,不單有失敗的使徒彼得,也有失敗的舊約的先知以利亞.(太十四30)說:「見風甚大」,(王上十九3)也有相同的幾個字,「以利亞見這光景,就起來逃命,到了猶大的別是巴,將僕人留在那裡.」以利亞「見這光景」,就好像彼得「見風甚大」一樣.兩個人都有一樣的錯；就是停止仰望主,開始看環境.以利亞所見的光景,就是一位拜偶像的王后怎樣死嚇他,他就起來逃命,忘記保佑他的全能的神,比任何王后任何危險還大.然而,彼得在海面上行走,雖然暫時遭遇失敗,但將要沉下去時,同禱告仰望耶穌,大聲呼喊:「主阿,救我!」結果耶穌趕緊伸手拉住他.
\newline
\newline
很多基督徒問:信徒向神禱告,要多長才能蒙神悅納呢?第三十節就給了我們答案.彼得信心搖動的時候,快要沉下,便大聲呼喊:「主阿救我!」他沒有時間把禱告拉長,只能喊出主阿救我這四個字；同樣,信徒禱告,必須注意一個屬靈的原則,主所要求的,不是禱告的長短,而是禱告是否懇切,誠實,彼得在這裡的禱告正是模範的禱告,是蒙神悅納的禱告.神必定立即幫助我們,扶持我們.
\newline
\newline
(四)信心的得勝這是第卅二節的主題「他們上了船,風就住了.」這節經文很短,但我以為最重要的兩個字是「他們」,原因是有兩個人同行,牽著手一同勝過狂風巨浪,不再有危險.這裡有得勝的基督徒生活的秘訣.彼得在這次的經歷中清楚明白,要勝過黑夜狂風,不是他這一個軟弱的人,而是他和耶穌.他和救主同行,就有力量有辦法,我們有抓住這個真理,了解得勝的靈性生活的秘訣未呢?我們是單獨過活,還是與主同行呢?信心的搖動是見風甚大,但信心的得勝是與主同行.
\newline
\newline
有一位愛主的弟兄,信主之前常到酒店去喝酒醉酒,被飲酒的壞習慣捆綁.酒店在鄉下地方,交通不便,不能駕車只能騎馬,到了目的地,就在酒店門前的杆子上綁馬.用繩綁好馬之後,便進去喝酒.酒店隔壁有間禮拜堂,門外也有綁馬的杆子.這位弟兄信主之後,第一個主日到禮拜堂做禮拜,在門外杆子上綁馬之時,酒店的老闆便出來對他說:今天你的眼睛有毛病嗎?你看錯了,這裡有我們綁馬的地方,你綁好馬進來喝杯酒吧!那位弟兄便說了一句很有意思的話,他說:「我們已經經過了.」老闆莫名奇妙,就問:你怎麼說我們?你是說你和你的馬嗎?弟兄便說:已經經過的,不是我的馬,而是我和我的主.我在屬靈方面是個新造的人,永遠與基督聯合的人.以後我不再進來喝酒,我有力量經過了.從今以後,我也不再用你的杆子綁馬,這間禮拜堂是我屬靈的家.所以,一個與主同行的生活,才是得勝的生活.
\newline
\newline
第五,信心的高峰.因為在試驗中,認識主更深；在以後的日子裡,便愛主更深.這是信心的高峰,第三十三節說:「在船上的人都拜他說,你真是神的兒子了.」很可惜有些基督徒,在牧師來探訪或勉勵他的時候,很會講一套很好聽的屬靈話；口裡雖然很愛主,開口便說感謝的話,但是內心還未到達信心的高峰.這些在愛主方面有問題的基督徒,就好像我所聽過的某一個小男孩,在自己的睡房裡打死一隻老鼠,便拉著老鼠尾到廚房給他媽媽看.他以為媽媽一定很喜歡.但當他正在拉的時候,牧師來了.孩子很害怕,他以為牧師一定不喜歡一隻死的老鼠,所以便想出一個很聰明的方法,就是講幾句很屬靈的話給牧師聽.他說:「這裡有一隻很危險的老鼠,是我剛才用信心打死的.牠在我的睡房裡出現,不把牠打死,牠便可能在我睡覺時咬我的腳.我拿棍子去打牠.真感謝主!耶穌加添我的力量,結果老鼠斷氣,耶穌也接了牠回天家.」一隻老鼠上天堂,各位有聽過嗎?口裡說的話雖然很好聽,但是不一定是信心的高峰.
\newline
\newline
(太八27)「眾人希奇說,這是怎樣的人,連風和海也聽從他了.」這是耶穌第一次顯出他的大能,改變天氣,平靜風浪；門徒就產生一個問題,這是一個怎樣的人呢?到了第十四章,他們第二次看見主顯出大能力的時候,門徒不疑惑,猜想,討論,問是怎樣的人；他們已經認識耶穌,於是全船的人都拜耶穌.說這真是神的兒子.這是信心的高峰,我們要注意,從馬太福音第八章到十四章,有一年的時間,其間門徒在認識主的事上很有長進,因為認識他不單是人,也是神,所以便愛他拜他.我們在屬靈上有進步嗎?第八章是猜想討論的時候；但第十四章他們已經認識他,而且拜他了.求主幫助我們,使我們不單經驗到信心的得勝,與主同行；並且來到信心的高峰,愛主更深.
\newline
\newline


\newpage

\section{第53屆~港九培靈會~林克己~第八講}
\label{sec:590}
\textbf{得勝有餘的神蹟}
\newline
\newline
連結: \href{https://www.hkbibleconference.org/session-message/view/590}{\texttt{https://www.hkbibleconference.org/session-message/view/590}}
\newline
\newline
\hyperref[sec:588]{< < < PREV SERMON < < <}
~
\hyperref[sec:toc]{[返目錄]}
~
\hyperref[sec:1386]{> > > NEXT SERMON > > >}
\newline
\newline
羅馬書 8:1-8:2
\newline
\begin{longtable}{cl}
\hline
\hline
章節 & 經文 (和合本修訂版)\\
\hline
8:1 & \begin{tabularx}{0.7\textwidth}{X} 如今，那些在基督耶穌裡的人就不被定罪了。 \end{tabularx} \\ \\ \relax
8:2 & \begin{tabularx}{0.7\textwidth}{X} 因為賜生命的聖靈的律，在基督耶穌裡從罪和死的律中把你釋放出來。 \end{tabularx} \\ \\
[1ex]
\hline
\hline
\end{longtable}
(羅八37)「然而靠著愛我們的主,在這一切的事上,已經得勝有餘了.」今天我們要從這節經文中,看得勝有餘的五方面真理.
\newline
\newline
(一)得勝的範圍「在這一切的事上」,是得勝的範圍.可是,在這一切的事上是指甚麼呢?要解答這個問題,便需要用一個解釋聖經很重要很基本的原則,就是看上下文前後的話.第三十五節是這句的背景,提到患難困苦逼迫的事,雖然大受試煉,但卻是得勝的範圍.在這些事上得勝有餘,在見證上打美好的勝仗.有些基督徒說,我雖然天天讀經禱告,但是讀聖經好像沒有味道,對我的幫助不很大.不過羅馬書第八章三十五和三十七節,可以幫助我們,用聖經培養自己靈性生活的方法；就是將這兩節的屬靈景況,應用在我們的生活中.在基督徒的生活中,誰沒有遇到患難困苦逼迫呢?這是每一個人必有的遭遇,但是我們怎樣應付呢?第三十七節就有很實際的教訓,也有我們得勝的範圍.無論是患難困苦逼迫,我們總要在這一切事上得勝有餘,好吸引人來歸向耶穌.
\newline
\newline
在印度一個開荒宣教的地方,有一個小教會,裡面有一位愛主的姊妹.她的丈夫很有學問,但是反對宗教,甚麼都不信.於是這位姊妹和教會的牧師便多為他禱告.結果有一個主日,這人來了,以後還常常參加教會的聚會.牧師很高興,每逢講道還儘量預備特別的信息,多用科學的比方,適合知識份子的例證.過了不久,姊妹的丈夫便信了耶穌,而且申請接受洗禮.牧師很快樂便問他說:「是甚麼道理使你相信耶穌.」這位才信主不久的弟兄便說:「我信耶穌,不是因為我所聽的道,而是因為我所看見的見證.任何道理我都可以抵擋,只是我妻子在遭受我反對嘲笑和逼迫時所表現的那份愛心和關懷；我無法抵擋,我願意也接受耶穌.」這位姊妹雖然受到逼迫,然而卻是得勝有餘.
\newline
\newline
(腓一27)「只要你們行事為人與基督的福音相稱」,這是保羅另一次的勸勉,我們特別留意「只要」這兩個字.保羅用只要這兩個字,是說明行為的重要,能吸引感動人來信主.或者你對真理的認識不很深,信主也沒有多久,但是只要你的行為聖潔光明,在信心在得勝的範圍內有長進,神就能使用你的見證.
\newline
\newline
(二)得勝的態度注意「靠著」這兩個字,「靠著愛我們的主」.(腓四13),「我靠著那加給我力量的,凡事都能作.」這是聖經另一次所提到的「靠著」,是我們所需要的態度.因為我們的智慧力量都有限,就要在凡事上倚靠那位全能全智的主.有一次,我有機會到死海去遊覽,看見一件我永遠不能忘記的事.在死海的水面,有一個人浮來浮去,手腳都不動,就像木頭一樣很容易的浮在水面,他頭上有一頂帽,手中有一份報紙,也有他的朋友在岸上替他拍照.死海的水很奇妙,鹽份很高,所以浮力很大,不動手不動腳也不會沉下去.耶穌的大能大愛就好像死海一樣,你靠著他,浮在祂的大愛之上,便永不會沉下去.耶穌在患難困苦逼迫中,能夠伸手扶持你.
\newline
\newline
現在我們看(羅八1-2),看看保羅怎樣用一個比方,來說明和發揮這個真理.「如今那些在基督耶穌裡的,就不定罪了.因為賜生命聖靈的律,在基督耶穌裡釋放了我,使我脫離罪和死的律了.」保羅用兩種律法的比方告訴我們,在基督徒的生活中,有兩個律相鬥相爭,有罪和死的律,好像地心吸力把我們拉下去,使我們跌下去.因為不仰望主,就受罪和壞習慣捆綁.另外是被聖靈充滿的人所享受的,是一種奇妙的釋放；是賜生命聖靈的律,在我們的生活所發生的效力.如果我這本聖經代表一個倚靠自己的基督徒,就必定跌倒,不能抵擋罪和死的律；但是現在我用手來代表聖靈的充滿,基督的福氣,我把手放在聖經之下,問題就解決了.因為手有能力有生命,這本聖經靠著我的手,就不會跌下去.因此,我們得勝的態度,是在乎靠著.
\newline
\newline
(三)得勝的秘訣請注意「愛我們的主」這五個字.基督徒遭遇患難困苦逼迫的時候,因為信心不堅固,信仰不堅定,常懷疑神對我們的慈愛,雖然本章二十八節提醒我們「萬事都互相效力」,無論享福或受苦,總是對我們有益處,但是我們在受試煉時,信心仍是搖動,所以我們需要保羅在這裡指出「愛我們的主」作為得勝的秘訣.我們雖然不明白神為甚麼要這樣對待我們,要暫時隱藏在患難後面,但是我們千萬不要停止信任他,倚靠他,認識他.
\newline
\newline
在希臘原文裡面,「愛」字很有意思,它不是現在時態,乃是過去時態.保羅在這裡講的,是愛了我們的主,然而靠著愛了我們的主.神為甚麼這樣默示保羅,要他用過去時態呢?這真是我們得勝的秘訣.我們在遭遇苦難之時,千萬不要灰心,主要的秘訣是多默想主的愛；同時提醒自己,主怎樣愛了我們.在歷史上為我們的罪死在十字架上,這是不能推翻的事實.主在十字架上已經愛了我們,他犧牲性命,流出寶血,這是歷史的證據.不但如此,主還用恩典多次聽了我們的禱告.我們要默想主怎樣愛了我們,解決我們的困難,供給我們的需要,加添我們的力量,引導我們的路程,使用我們的見證.只要我們回想主怎樣在歷史上,在恩典中愛了我們,想到這些過去的恩典和事實,我們就能夠得勝.主怎樣愛了我們,也必愛我們到永遠.因為他是永不改變的,昨日,今日到永遠都一樣.所以,我們要靠著愛我們的主,過得勝的生活.
\newline
\newline
(四)得勝的程度請注意「有餘」這兩個字.基督徒不單要得勝,而且要得勝有餘.可是這個標準和程度,有誰達得到呢?如果有人達到,這不是屬靈的神蹟嗎?兄弟在埃及當兵的時候,有一晚黃昏,忽然有軍官來說,曠野有仇敵來攻擊我們,要我們拿來福槍出去應戰.當時天色很暗,甚麼都看不清楚,我們前進的時候,軍官忽然命令我們開槍；我們因為看不見目標,只得朝前面亂開.後來甚麼聲音都沒有了,我們便以為敵人已經逃跑或是被打死,於是我們很高興的回軍營.第二天清早起來,令我們感到尷尬的是,我們作夜所射中的,不是仇敵,只是一個大紙皮箱,滿了我們所射的彈孔.恐怕我們那次假的勝利,就像很多基督徒的靈性生活一樣；沒有真的勝利,只是假的得勝.我們的仇敵魔鬼比紙皮箱還大；我們要認識他各樣的詭計,熟悉聖經,熟悉撒旦的作為,和知道自己要裝備的武器,才能夠得勝.
\newline
\newline
第二次世界大戰的時候,英國將軍蒙哥馬利,在他房間的牆壁上,掛了敵軍軍長的相.蒙哥馬利有個奇怪的習慣,就是每天預備打仗的計劃之前,必花幾分鐘站在像前,一直看敵軍軍長.他不單定睛看,還繼續不斷的說:「認識你的仇敵,認識你的仇敵.」因為他這樣做,就能明白對手的思想.歷史家知道,沒有一個德軍在戰場上能敵擋蒙哥馬利的.願主幫助我們,把得勝有餘的基礎立好,就是多查考聖經,認識仇敵.社會的潮流雖然常有改變,但撒旦所用的方法卻沒有多大改變；他在聖經中怎樣攻擊神的兒女,今日也同樣攻擊我們.
\newline
\newline
(五)也是最重要的,就是得勝的時間請注意「已經」兩個字.保羅說:「然而靠著愛我們的主,在這一切的事上,已經得勝有餘了.」還沒有打仗,已經得勝,是甚麼意思呢?有一位屬靈愛主的弟兄,預備有挑戰性的定義來解釋,他說一個得勝者,和一個得勝有餘者,有很大的分別.一個得勝者,與仇敵交戰時,還未知道那次打仗的結果,需要等到最後才知道自己是否得勝者.但是一個得勝有餘者,還未到戰場,已經把打仗的基礎立好,已經作美好的準備,遇到仇敵就不能不打勝仗.這個解釋非常好.我們的得勝,是在乎我們讀經禱告與主相交的時間,這能支持我們一天過得勝的生活.
\newline
\newline


\newpage

\section{第53屆~港九培靈會~楊濬哲~序}
\label{sec:1386}
\textbf{序}
\newline
\newline
連結: \href{https://www.hkbibleconference.org/session-message/view/1386}{\texttt{https://www.hkbibleconference.org/session-message/view/1386}}
\newline
\newline
\hyperref[sec:590]{< < < PREV SERMON < < <}
~
\hyperref[sec:toc]{[返目錄]}
~
\hyperref[sec:555]{> > > NEXT SERMON > > >}
\newline
\newline
為了港九培靈研經會的工作,我內心有一個很重的負擔.我雖然已在美國定居,為了這負擔,每年五月底,我就由美國回到香港參與籌備的工作,直到大會結束才離開.眼見每年大會赴會的人數,不斷增加,信徒愛慕主道的飢渴態度,真令到大會欲罷不能.我真要從內心裡發出感謝神的聲音,而願將一切榮耀歸與神.
\newline
\newline
八一年第五十三屆大會,有一個新的嘗試,開闢第二會場.借用宣道中學的宣中堂為副會場,裝置閉路電視和音響完善的設備,以便晚堂奮興會時,可以容納更多的人.在經費方面雖然開支增加了,感謝神!卻豐豐富富的供給,使我們一無所缺.
\newline
\newline
大會的三位講員:研經會鮑會園牧師主講哥林多前書,雖然只有八堂,卻能把書內的精要,難以講解的真理,都儘量表達出來了.對教會,對夫婦關係,對信徒的日常生活,有明確指導.林克己牧師主領的講道會,他是一個英國人,卻能講出流利的國語,令人羨慕不已.他所講的總題是主奇妙的作為,用八個神蹟去表達,使人耳目一新.奮興會由戴紹曾牧師主領,他是第二次來領會,講的仍然是使徒行傳,初期教會的形態；而這次卻注重差傳方面,在教會增長呼聲殷切的今日,這信息實在是需要而合適的.
\newline
\newline
歡迎葉紹蔭長老,加入委辦會與我們併肩事奉；朱小皿弟兄擔任執行幹事職,願主大大使用他們.
\newline
\newline
張有光牧師多年來任委辦會的副主席,對會場的安排,講員的招待,工作的分配,同工的合作,給我們很大的幫助.他雖已退休,在美國定居,但仍繼續在西雅圖有事奉.他勞苦的榜樣,仍在我們眼前,願主記念!
\newline
\newline
本講道集蒙黃民弟兄,梁壽華弟兄,雷麗端姊妹筆記,朱建磯牧師整理編輯,陳綺華姊妹協助校對,冼錦棠弟兄設計封面,乃得如期出版.願主使用本書,造福多人.
\newline
\newline
多謝九龍城浸信會執事會,宣道中學校董會,借用場地.同公義工合作,招待佈置；以及所有為大會工作出錢出力的主內兄姊同道們,委辦們,謹書數言,統此致謝!
\newline
\newline


\newpage

\section{第53屆~港九培靈會~鮑會園~第一講}
\label{sec:555}
\textbf{本書簡介及教會的分爭結黨}
\newline
\newline
連結: \href{https://www.hkbibleconference.org/session-message/view/555}{\texttt{https://www.hkbibleconference.org/session-message/view/555}}
\newline
\newline
\hyperref[sec:1386]{< < < PREV SERMON < < <}
~
\hyperref[sec:toc]{[返目錄]}
~
\hyperref[sec:558]{> > > NEXT SERMON > > >}
\newline
\newline
我非常感謝主,再有機會在培靈會與大家一起思想神的話.當我研讀哥林多前書的時候,心裡覺得非常沉重,因為當時哥林多教會竟發生意想不到的困難；但我也為此感謝神,因為這些問題,也看見神話的啟示,主的恩典勝過那些困難.哥林多前書是一本非常實際的書,雖然內容有深奧的教訓:如靈恩問題,主耶穌與信徒復活問題；但多論及教會實際的問題,本書幫助我們知道怎樣處理這些問題.所以未來幾堂講道中,會思想生活中實際的問題,而不著重在複雜的神學教義上.
\newline
\newline
保羅在第二次佈道旅程中,建立了哥林多教會.這城市是當時中西交匯之地,在文化經濟各方面很發達；因此,許多哥林多人心目中看重生意,金錢和物質；而沒有更高的生活目標,更可憐連教會也受了影響.教會裡的人也效法這些人的思想,追求物質,表面上愈發達愈好,卻不追求更高的目標.其實神已格外賜福哥林多教會,哥林多信徒追求恩賜,知識,口才,看他們這種情況,能否給香港基督徒一些警惕呢?香港社會非常發達,因此「發達」成為香港人追求的目的,而教會是否也受了影響呢?外面愈發達愈好,奉獻愈多愈好；聚會人數愈多愈好,禮拜堂夠大就好了.甚至這思想影響了一些傳道人,以為愈忙愈好.最近我與一位醫生太太談話,她說:「大概牧師與醫生差不多,醫生若沒有病人來找他；牧師若沒有人來請講道,他就不好了.」難怪有些傳道人盡量使自己忙碌,當他心裡漸漸感到枯乾,心靈得不到滿足時；於是再加一些工作,就用工作來補充心靈的虛空,這是多麼可憐的現象.
\newline
\newline
其實很多基督徒和教會都有這種危機,整個教會將注意力放在外表上,而忽略了靈裡實在的需要.哥林多教會正是這樣,保羅指出基督徒的問題不在外表,而乃是內在生命的問題.這本書結構簡單而實際,保羅論及的問題可分兩類:
\newline
\newline
(一)保羅從別人處聽到哥林多教會的事情.第一章至第六章,「因為革來氏家裡的人,曾對我題起弟兄們,說你們中間有分爭.」(一11)「風聞在你們中間有淫亂的事.」(五1)六章前半是保羅聽說他們中間有爭訟的事,六章後半是關乎吃祭偶像之物的問題.
\newline
\newline
(二)哥林多教會寫信問保羅的問題,第七章至第十五章.「論到你們信上所題的事」(七1),七章前半講及婚姻,七章廿五節起論到童身的問題.八章至十一章論吃祭偶像之物的問題.十二章是論到屬靈恩賜的問題.
\newline
\newline
在思想這些具體問題前,我們先思想哥林多教會,寫信問保羅的神學問題:婚姻,童身,靈恩,身體復活等,這都是很深奧的思想,似乎顯出他們的知識很深,可以問深奧的神學問題,但卻不明白生活中罪惡的問題；可以問及復活,但卻不知如何對付犯姦淫的弟兄.雖然有許多的知識,但卻與生活脫節,這就是危機.保羅時代如此,現今香港教會也是如此.近來許多神學院與基督教團體舉辦的講座課程,似乎以題目愈深,收費則愈高,參加人數就愈多.為此我感謝主,但我擔心弟兄姊妹注重知識,而不知道把知識應用在生活中,這是很重要的問題.故此哥林多教會很逼切的寫信問保羅.保羅卻並未立即回答,他乃先提及他們生活中的罪,然後再回答他們的神學思想問題.從這樣的次序,就可以看見保羅的看法重點,所以我們要按照保羅的次序來思想這些問題.今天我們思想(一10-18)
\newline
\newline
哥林多信徒因為知識與恩賜很多,因此各人的看法都不同；那一方面如擁有恩賜,知識,才幹的人,便有人去擁護他；另外的人就去擁護其他的人,所以教會便有結黨和分爭的事情發生.「你們各人說」(一12),他們每人都選擇自己的立場,參加一個黨派.哥林多教會分開四大黨派,四個黨派都用「人物」作代表,是極具代表性的,他們在教會的圈子有特別的意思:
\newline
\newline
(一)磯法黨　彼得是使徒,受猶太人所擁護.他有有時明知是真理,但因受到傳統的影響,像猶太人是不可與外邦人吃飯的；所以他為了害怕猶太人的批評,便不與外邦信徒吃飯.這種人是極端守猶太律法,只受某些人尊敬.我們守規則,若做得到就好,表示自己在某方面有成就；人時常有個缺點,是要靠自己的功勞,來得到神的喜悅.因自己多賺一點功勞,多做一些工作,多捐一點錢,便可以使自己在人面前得一點誇口.不認識神的人也甘心樂意作善事捐錢,表示自己有能力作善事；甚至有人藉著受苦來賺得功勞.馬丁路德未明白救恩以前,晚上睡在地板上；用鞭子打身體,令自己全身流血；用膝頭爬樓梯向神父告解.這不單祇是未信主的人所有的思想,甚至有些基督徒也有這種將功贖罪的觀念.其實最重要是向神認罪,靠律法立功,這不是方法；心靈的問題,若不在主面前對付,你作多少工作,捐多少錢,參加多少聚會也沒有用.靠律法立功,這是磯法黨的錯誤.
\newline
\newline
(二)保羅黨　這黨的人比磯法黨開通,保羅是外邦人的使徒,單傳恩典的福音；只憑信心,相信耶穌就夠了,不靠守律法得救.生活與救恩有甚麼關係呢?我們所尊重的是保羅,聖經沒有教導我們靠行為得救,所以我們作甚麼都可以.這些人很可能習慣用恩典作藉口,放縱自己的肉體生活；在教會中我們常看見這種情況,既已承認相信主耶穌,但生活卻不像基督徒的樣子.這種人當他犯罪後,你絕不能斥責他,因他會答你說:「你休想用律法捆綁我,我是在恩典之下,而不在律法之下.」多年前我到一個教會講道,講題是「律法的功用」.在講道前,一位長老先禱告,他好像講一篇道說:「神啊!我們得救了,不是靠律法,求你不要用律法捆綁我們.」後來我纔知道,原來這位長老在教會犯了非法的事；當他的傳道人與教會一位執事,與他交通時,要他接受教會的懲治,但他卻說:「此地不留人,自有留人處.」他在生活中明明犯了罪,但你卻不能斥責他.
\newline
\newline
(三)阿波羅黨　(徒十八,十九章)阿波羅很有口才,也很會講解聖經；如在今天來說,他確是一位好的學者.彼得算甚麼呢?他只不過是個漁夫；而保羅也其貌不揚,寫信時,又沉重,又厲害,言語又粗俗,說話又不文雅；但阿波羅卻不同了,他是位學者,這班人就以阿波羅為號召了.可能他們是高舉知識學問,其實追求知識學問並非錯誤,但我們的知識學問有沒有被主耶穌掌管呢?可能這些知識學問會絆倒很多人!今日教會祇注重追求某些有知識學問的學者的觀點,他們甚至不怕公開違背聖經的真理.有時更以追求知識為目標.故此,青年人要小心,作為鑑戒,而成年人也要小心,以免逼使子女犯上追求知識的危機.他們以為有了知識學問就可尋找到好的工作,獲得優厚薪酬；沒有知識的人在社會中便沒有地位.因此追求知識,便成了生命中最重要的目標,最好能先進大學,或者專上學院,其次是師範.等到落選時,纔考慮進入神學院.這樣就形成了世界知識,比起聖經知識更重要；等到世界知識得不到的時候,纔追求聖經的知識,這是很錯誤的觀念.
\newline
\newline
(四)基督黨　難道跟隨基督是錯誤嗎?跟隨基督不會錯,但這班人的態度錯了,這些人大致上可能有兩個錯誤:(1)在某些方面有屬靈的驕傲.他們以為別人追求磯法,保羅,阿波羅.其實,主的僕人只不過是人,是主的門徒而已；但他們卻以為他們所追求的,乃是以基督為目標,他們認為他們一黨乃是最好的.當然以追求基督目標是最好的,不過若用這些作屬靈藉口,卻是一種錯誤.我們不能認為別人與我不相同,就錯了,別人的真理看法,教會處事方式與我不同,他就錯了,全世界只有我纔對.(2)這些人可能在屬靈上有特別的經驗,可能有人在異象中看見主,會說方言,繙方言,被醫治等.有這些經驗是神的恩典,但恩典經歷,並不是用來高舉自己.我們不能認為別人沒有這種經驗就不屬靈,甚至別人沒有這種經歷,就不是基督徒.其實每個人蒙恩都不一樣的,神給我們每個人都有不同的恩典；但卻不能用這恩典,來批評他人和教會.最近新興的怪異教會,報紙稱他們為邪教.他們高舉「死亡」,說死了比活著還好,他們批評不與他們追求的,都不是教會.如果有這種態度,在基本上就誤解了神的啟示.如果神給我們有特別亮光,這亮光和恩賜,是合乎聖經的教訓,我們就應當感謝主；但我們不能說,別人與我們不同的,他是錯了.
\newline
\newline
以上是哥林多教會所犯的毛病,保羅說基督是不可以分開的,跟隨主的人都是基督身上的肢體.這個實際的道理,是神的大能；我們要同心合意建立基督的身體,而不是高舉某肢體或自己.但願哥林多教會的錯處問題,成為我們的鑑戒.
\newline
\newline


\newpage

\section{第53屆~港九培靈會~鮑會園~第二講}
\label{sec:558}
\textbf{三等人}
\newline
\newline
連結: \href{https://www.hkbibleconference.org/session-message/view/558}{\texttt{https://www.hkbibleconference.org/session-message/view/558}}
\newline
\newline
\hyperref[sec:555]{< < < PREV SERMON < < <}
~
\hyperref[sec:toc]{[返目錄]}
~
\hyperref[sec:559]{> > > NEXT SERMON > > >}
\newline
\newline
哥林多前書 2:14-3:7
\newline
\begin{longtable}{cl}
\hline
\hline
章節 & 經文 (和合本修訂版)\\
\hline
2:14 & \begin{tabularx}{0.7\textwidth}{X} 然而，屬血氣的人不接受神的靈的事，他反倒以這為愚拙，並且他不能了解，因為這些事惟有屬靈的人才能領悟。 \end{tabularx} \\ \\ \relax
2:15 & \begin{tabularx}{0.7\textwidth}{X} 屬靈的人能看透萬事，卻沒有一人能看透他。 \end{tabularx} \\ \\ \relax
2:16 & \begin{tabularx}{0.7\textwidth}{X} 「誰曾知道主的心？誰會教導他？」至於我們，我們有基督的心。 \end{tabularx} \\ \\ \relax
3:1 & \begin{tabularx}{0.7\textwidth}{X} 弟兄們，我從前對你們說話，還不能把你們當作屬靈的，只能把你們當作屬肉體的，你們在基督裡僅是嬰孩。 \end{tabularx} \\ \\ \relax
3:2 & \begin{tabularx}{0.7\textwidth}{X} 我用奶餵你們，沒有用飯餵你們，因為那時你們不能吃。就是如今還是不能， \end{tabularx} \\ \\ \relax
3:3 & \begin{tabularx}{0.7\textwidth}{X} 因為你們仍是屬肉體的。你們中間有嫉妒、紛爭，這豈不是屬乎肉體，照著世人的樣子生活嗎？ \end{tabularx} \\ \\ \relax
3:4 & \begin{tabularx}{0.7\textwidth}{X} 有人說：「我是屬保羅的」；有人說：「我是屬亞波羅的」；這樣你們豈不是和世人一樣嗎？ \end{tabularx} \\ \\ \relax
3:5 & \begin{tabularx}{0.7\textwidth}{X} 亞波羅算甚麼？保羅算甚麼？我們都是神的執事，藉著我們，你們信了；這不過是照著主給各人的恩賜去做罷了。 \end{tabularx} \\ \\ \relax
3:6 & \begin{tabularx}{0.7\textwidth}{X} 我栽種了，亞波羅澆灌了，惟有神使它生長。 \end{tabularx} \\ \\ \relax
3:7 & \begin{tabularx}{0.7\textwidth}{X} 可見，栽種的算不了甚麼，澆灌的也算不了甚麼；惟有神能使它生長。 \end{tabularx} \\ \\
[1ex]
\hline
\hline
\end{longtable}
基督教與其他宗教最主要的分別,乃是信的人,與不信的人,基本上生命的不同.哥林多教會分爭的最主要原因,是因為他們依照世人的生活方式,卻不像基督徒的樣式；所以保羅在本章要講及這個基本的分別.他提及兩大類三種不同的人.基本上世人可分為兩類:第一類是「屬血氣的人」(二14),這是未信主,沒有主生命的人.第二類是「弟兄們」(三1),顯然這是基督徒,但在弟兄們中,又可分為「屬肉體的」基督徒,他們信主耶穌後,仍過屬肉體的生活；另一種是「屬靈的」基督徒,信主後,他們的生活可以見證他們正是跟隨主的人.
\newline
\newline
(一)屬血氣
\newline
\newline
屬血氣可說是生來自然的人,看似與基督徒沒有多大分別,但他們整個生活目標是按自己所定；因為他們沒有新生命在內,故此沒可能有另外的標準.他們是屬於世界的,所以仍不是基督徒.這些人生活的表現可能很有愛心,謙卑,受人尊敬；在商場可能是個誠實的好人,生活不像世人的污穢罪惡.有時他們也到教堂聚會,覺得倚屬教會是一種身份和地位；甚至願意對教會負上一些責任,捐錢給教會,在教會圈子裡很有名望,受人尊敬.但是他們沒有被主改變,沒有主的生命,主耶穌對尼哥底母說:「人若不重生,就不能見神的國.」(約三3),這些人沒有重生的經驗,整個人的生活按自己本性來表現,這是屬血氣的人；他所有的生命,永不能帶他到神的面前.
\newline
\newline
(二)屬肉體
\newline
\newline
保羅稱他們為弟兄們,但他們是屬肉體的.這些人是重生得救的,有真正屬靈生命；但仍按照自己的肉體過生活,生命與生活沒有發生關係；表現像不信的人,不像真正的基督徒.保羅形容屬肉體的人在基督裡為嬰孩,他不能用飯喂他們.嬰孩有很多地方很可愛,家中有了新生命,全家都喜樂感謝神；但過了幾年,嬰孩仍未長大,家長會很憂愁擔心.照樣很多基督徒信主五年十年,表現仍是嬰孩一樣,令天父擔心.
\newline
\newline
嬰孩有以下的表現:
\newline
\newline
a.嬰孩不會餧養自己　嬰孩餓的時候會哭,但絕不會用奶餧養自己；你需要放奶入他口中,他才會吃.有很多基督徒就是這樣；不知道用屬靈的糧食來餧養自己,一星期做一次禮拜,聽到一篇信息很快樂,待下星期再來聽,他永不知道怎樣讀神的話,怎樣把神的話應用在自己生命裡,不會照顧自己的生命.以前我剛到一個教會事奉,有一位弟兄來找我,問我如何處理家中的問題,我便幫助他明白聖經的教訓.他便很高興地離開了,但下星期再來,又是問同樣的問題,我又幫他去解決.以後他每星期來一兩次,問同樣的問題,但他自己不知道如何應付；起初我很高興他來找我,但後來我就覺得有點煩了.
\newline
\newline
b.嬰孩不理會別人　家裡正忙碌的時候,嬰孩不會明白,到吃奶的時候,便會哭,絕不肯延遲兩個小時才吃.家裡經濟困難,他也不會替媽媽擔心,吃飽了便安安靜靜睡覺,別人捱餓,與他無關,他只想及自己,不會關心別人.有很多人屬靈生命正是如此,只會滿足自己的願望.教會中有不滿意的事,便會立刻發怨言,絕不理會這事的因由與後果,只想及自己的需要,得不到便不高興了.整個生命目標只為自己,從未想到弟兄姊妹的需要；更不會用禱告關心,或以物質幫助有需要的人.
\newline
\newline
c.嬰孩喜歡表現自己　嬰孩有機會表現自己便快樂,否則便不開心了.有一次我去探訪一對傳道人夫婦,他們有兩個四五歲的兒女,大的是哥哥,小的是妹妹.我便請這位哥哥唱歌,待他唱完後,那位妹妹躲進房去；原來她在房裡哭起來,我覺得很奇怪,她為甚麼會哭?做媽媽的就最知道兒女的心情,媽媽說:「因為你請哥哥唱,沒有叫她唱.」有很多基督徒,只要有機會表現自己,出多少錢,做多少；也願意拚命去作,整個工作的目的,要在人前表現自己.在我以前工作的教會,是把主日獻花者的芳名刊發在週刊上,有一次忘記刊登獻花者的名字,到月底結賬時,這人便不肯交錢.他說:「秩序單沒有刊登我的名字,人家怎麼知道是我奉獻的,我何必要給錢呢?」
\newline
\newline
d.嬰孩只顧外表喜好　保羅說:「你們是屬肉體的,因為在你們中間有嫉忌分爭.」(三3),嫉妒分爭是老我的表現,我喜歡就喜歡,我不喜歡我永遠也不會喜歡.如果你所得的稱讚不及我多,我便心裡快樂；樣樣按自己的喜好,這是嫉妒分爭的來源.事事以自我為中心,整天關心自己喜歡做的事,是否按照自己的意思,是否有應得的報酬與獎賞.
\newline
\newline
e.常在肉體中失敗　這些人並非喜歡罪惡,也並非不知道罪惡是不好的；只是有願意對付的心,卻不能勝過罪惡捆綁.譬如看見不如意之事便發脾氣,但事後心裡知道錯,很難過!在主面前痛哭悔改,不過下次仍依然故我.整個生活常在失敗中掙扎,有人說這種人是常常悔改,常常犯罪的基督徒.
\newline
\newline
f.在屬靈事上不結果子　屬肉體的基督徒不是不願意工作和見證,但卻永遠作不出有永遠價值的工作來.因為生命是按照老我過生活.主耶穌說:「你們若常在我裡面,我也常在你們裡面；我的話若常在你們裡面,如此你們就必結果子.」倘若我們沒有住在主裡面,主的話就不常在我們裡面,我們永遠不能結出有屬靈價值的果子了.
\newline
\newline
(三)屬靈的
\newline
\newline
屬靈的基督徒與屬肉體的基督徒,一樣同有主的生命,但分別是在生活上.他們整個生活是按屬靈的原則,屬靈生活是他最高的標準.但由於仍然要在肉體中活著,因此一個屬靈的基督徒也可能會犯罪；不過在犯罪之後卻會在掙扎中勝過罪惡的捆綁.經過試探和失敗,並在主面前受對付過後,下次就不會再被同樣罪惡所轄制了.一個屬靈的基督徒,並不是沒有失敗經驗的基督徒,因為我們在今生,不可能過一個完全聖潔的生活,但我們整個生活方式,應不斷向上有長進.我們可以用爬山的方法來比喻基督徒的生活.普通的山路不是一直向上的,一般都是爬一段路後,就有一段平路,甚至有時向下走一小山谷,然後再爬另一山峰；向上再一段路後,又有一段平路或向下走幾路,然後又向上走.一個基督徒屬靈生活的長進,有時候會有一段平淡的生活,甚至有失敗跌倒的經驗；但經過這個經驗後,我們可以爬更高的山峰.當爬到這個山峰後,我們又會走一段平淡生活或跌倒；但是整個生活方式,卻是向上有長進和得勝的生活.
\newline
\newline
第二方面,屬靈基督徒的生命,不是以自己為中心,他乃是以主耶穌和神的旨意來過生活的標準,雖然有時候在考驗中會跌倒,但他真正生活目標是按主耶穌的意思來過生活.不再以自己的心意,乃是以主的心意為生活的標準.主耶穌在客西馬尼園禱告說:「父啊!倘若可行,求你叫這杯離開我......我心裡甚是憂傷,幾乎要死.」但他最後說:「然而不要照我的意思,只要照你的意思.」
\newline
\newline
第三方面,屬靈基督徒有結屬靈果子的生命,能結出有價值常存的果子.他的生命不是為結果子而活,乃是為遵行神的旨意；故這種生命,自然能結出屬靈的果子.
\newline
\newline
我們已看過這三種人的生活,若我們發現自己仍在屬血氣的生活裡,我們該怎麼辦?我覺得唯一的事情,現在俯伏在主面前,對主說:「我知道我的需要,我需要你的救恩,求你救我.」如果我們發現自己在屬肉體基督徒的生活裡,我們需要坦白與主面對現實,將這需要帶到主面前,求主加添力量；靠著主的恩典,我以後要勝過這些屬肉體的表現,追求過屬靈的生活.
\newline
\newline


\newpage

\section{第53屆~港九培靈會~鮑會園~第三講}
\label{sec:559}
\textbf{屬靈生命的建造}
\newline
\newline
連結: \href{https://www.hkbibleconference.org/session-message/view/559}{\texttt{https://www.hkbibleconference.org/session-message/view/559}}
\newline
\newline
\hyperref[sec:558]{< < < PREV SERMON < < <}
~
\hyperref[sec:toc]{[返目錄]}
~
\hyperref[sec:565]{> > > NEXT SERMON > > >}
\newline
\newline
哥林多前書 3:10-3:17
\newline
\begin{longtable}{cl}
\hline
\hline
章節 & 經文 (和合本修訂版)\\
\hline
3:10 & \begin{tabularx}{0.7\textwidth}{X} 我照神所給我的恩典，好像一個聰明的工頭，立好了根基，別人在上面建造；只是各人要謹慎怎樣在上面建造。 \end{tabularx} \\ \\ \relax
3:11 & \begin{tabularx}{0.7\textwidth}{X} 因為，那已經立好的根基就是耶穌基督，此外沒有人能立別的根基。 \end{tabularx} \\ \\ \relax
3:12 & \begin{tabularx}{0.7\textwidth}{X} 若有人用金銀、寶石，草木、禾秸，在這根基上建造， \end{tabularx} \\ \\ \relax
3:13 & \begin{tabularx}{0.7\textwidth}{X} 各人的工程必將顯露，因為那日子要將它顯明，有火把它暴露出來，這火要試煉各人的工程怎樣。 \end{tabularx} \\ \\ \relax
3:14 & \begin{tabularx}{0.7\textwidth}{X} 人在那根基上所建造的工程若能保得住，他將要得賞賜。 \end{tabularx} \\ \\ \relax
3:15 & \begin{tabularx}{0.7\textwidth}{X} 人的工程若被燒了，他將損失，雖然他自己將得救，卻要像從火裡經過一樣。 \end{tabularx} \\ \\ \relax
3:16 & \begin{tabularx}{0.7\textwidth}{X} 難道不知你們是神的殿，神的靈住在你們裡面嗎？ \end{tabularx} \\ \\ \relax
3:17 & \begin{tabularx}{0.7\textwidth}{X} 若有人毀壞神的殿，神一定要毀滅那人；因為神的殿是神聖的，這殿就是你們。 \end{tabularx} \\ \\
[1ex]
\hline
\hline
\end{longtable}
保羅說:「因為我們是與神同工的,你們是神所耕種的田地,所建造的房屋.」(三9)基督徒是神所耕種的田地,神是耕種的主,好像主耶穌在福音書比喻一般；人子是撒種者,福音種子是由神播出來的.但這節聖經,在解釋上很困難,一方面是神所耕種的田地,而另一方面又是神所建造的房屋,這句話的主詞是誰呢?按思想是神,但在文法上不是最好的解釋.另外在神學思想也是錯誤的觀念,生命是神所播種,若神不放種子在我們身上,我們就沒有盼望.但當我們有這生命之後,保羅又在這生命上轉了另外的比喻方式來講.他好像看見一座房屋,而這座房屋的建造也是我們.神將生命放在我們裡面,生命的成長是我們每一個基督徒的責任；生命立在主耶穌基督根基之上,好像神把種子放在我們生命裡面一樣,我們不能作這個工夫,但當根基立好後,要在上面建造房屋；表示生命需要繼續長大,這是我們每一個基督徒應該負起的責任.聖經上從沒有說神要叫我們生命怎樣長大,為此原因保羅在書信中多用房屋來比喻基督徒的生命,在以弗所書也是這樣.生命是一樣,而根基也是一樣,但在根基上建造的房屋會有很大的分別,大小也不同.我們基督徒如何建造自己的生命,怎樣栽培自己的生命,我們不能埋怨神,乃是自己要負責任.
\newline
\newline
保羅說我好像一個聰明的工頭一樣,立好了根基,將生命種子放在哥林多教會裡,然後有別人在上面建造,「只是各人要謹慎怎樣在上面建造」(三10),這裡各人是指每一個人,每一個有主耶穌生命的人,我們要謹慎怎樣在生命裡建造.下面保羅提及我們的生命可以用有價值或無價值的材料來建造,如果有人用金銀寶石來建造,他的工程就能夠存留得住,如果用草木禾楷來建造,在審判時會就被火燒毀.
\newline
\newline
在未思想我們要用甚麼材料,在屬靈生命根基上來建造之前,我們先思想一個重要的真理,不論你用甚麼材料來建造你的生命,你的生命也會一天天的被建造出來.以前的日子時間已經過去了,你永遠無機會改善以往被建造的生命,你可以重修或修改一幢建築物；但我們的屬靈生命已被建造,你永不能改變它,這是多麼嚴肅的思想.你作錯的事情,永遠存在你的生命裡,每一個過去的時間,都無法改變自己的生命.我們要謹慎自己的生命,如果我們在上面建造沒有價值的材料,我們將來怎能向主交賬；所以我們每人在主面前檢討自己,思想到我們未來的生活方向.
\newline
\newline
另一方面,保羅說我們用金銀寶石,草木禾楷在這根基上建造,將來每一個人的工程在主面前被顯露,顯露出用甚麼材料來建造,金銀寶石指有永遠價值,可經得起考驗的.草木禾楷指沒有價值,經不起考驗,保羅不是說有些人可用金銀寶石,有些人可用草木禾楷,來作建造的材料；保羅指出基督徒都有主的生命,沒有人一生都是草木禾楷,一生中沒有作過一件有價值的工作；同樣地,沒有一個基督徒一生全用金銀寶石來建造,他一生沒有錯誤,沒有作過任何一件沒有價值的事,沒有基督徒能過一個百分之百完全的生活.我們的生命是用許多不同材料來建造,這是無可避免的,因我們還活在老我的生命裡,但我們有一個責任,盡量用金銀寶石來建造.
\newline
\newline
但是這種好壞的材料混雜在一起,很多時候是基督徒的悲劇.我們有時心裡受感動,很熱心作事,便用金銀寶石在生命上加上兩塊磚頭或一層材料；但有時候我們受環境影響,整個生命用草木禾楷建造了一段,這是多麼可憐的建造.造成這種情形的可能有幾方面的原因.
\newline
\newline
(一)有時候在生命中選擇錯誤,將注意力放在某些事上,而看輕了另外一些事情.譬喻有些人非常注重靈修的生活,每日逼切禱告,渴慕主的話,在主面前對付自己,他在這一方面是用金銀寶石建造的生命；但另外一方面,卻不注意對待人的態度,在家裡對父母,妻子,兒女非常不好,因此在這方面他用草木禾楷來建造生命.整個人好像活在兩個圈子裡,這是非常錯誤的觀念.
\newline
\newline
(二)有時候我們單單注重工作,對神的工作的託負,非常認真,但對人的責任,對自己的責任卻忽略了,其實這是同樣的重要.每個人都生活在不同的圈子裡,思想方式及生活都受自己喜好影響,很容易對一些事情劃分屬靈與不屬靈.很多年前,我剛從外國返港事奉,被請到一個青年聚會講道,這會場很特別,它的四週都有水龍頭,講道前主席出來報告,請弟兄姊妹把水龍頭關掉,以免浪費用水.當我講完道出外準備返回神學院,當時我對香港情況不大熟悉,所以便登上了一輛白牌車離開.第二天有一位弟兄來神學院找我,告訴我昨天坐白牌車是非常錯誤的事情；你作這種事,等於不是基督徒.這位弟兄看坐白牌車的事很重要,但對水龍頭有否關掉,是否浪費食水卻不要緊.一方面要守得很嚴緊,一方面看卻認為不重要.
\newline
\newline
有時我們會花許多的時間去作一些沒有價值的事,是因為我們將注意力放在某些錯誤的事上.有時候是受環境影響,好像彼得懼怕人的話,不敢承認主.雅各在伯特利與神面對面,心裡受感動,與神立約,願意為神而活；但看見拉班有許多牛羊,而自己還沒有甚麼財產,於是用詭詐方法奪取舅父的牛羊.不論為怕環境或為了自己的愛好,我們眼睛總是矇蔽了,把某些東西看作重要,但把某些事情卻看輕了.實際上這會影響我們屬靈生命的建造.保羅提及這種前後不相符的生活,在我們一生中作過多少草木禾楷的工作,我們用多少金銀寶石,同時又用多少草木禾楷來建造生命?
\newline
\newline
跟著保羅又提出一個警告:「各人的工程必然顯露,因為那日子,要將他顯明出來,有火發現,這火要試驗各人的工程怎樣.」(三13),在那日子乃指主耶穌再來的日子,我們要面對面見主,我們每一日所建造的工程要在主面前受考驗,今天可能無人考驗我們,我們的思想,工作,可以嚴密隱藏,但有一天,我們每一件事都要在主前受審判.到那時候草木禾楷的東西都要被燒毀,雖然建造在主的根基上,仍然可以得救；但生活中卻沒有一件事被主悅納稱讚的,這種生活多麼可憐呀!
\newline
\newline
保羅又提醒我們,每一個房屋建造的人是我們,我們的生命是神的殿(三16).神的靈住在我們裡面,我們不能說我沒有聖靈住在心裡,聖靈是有位格的神,不是物質.因此不能說我們的聖靈有多少?通常我們有一種錯誤思想,以為聖靈充滿就增加聖靈,好像倒水注滿一個杯子一樣.其實分別不是在你有聖靈多少,而乃是聖靈在你的生命掌握握多少主權；你順服聖靈有多少,我們認定聖靈是神.所以每一個的決定要追求聖靈的喜歡,如果我們尊敬聖靈在我們的生命裡,我們不會違背聖靈,作聖靈不喜歡的事情.「破壞」這個字在原文是看作平常的意思,有人把這殿看作平常.(徒十章)記載彼得看見異象,天上有布袋墜下來.裡面有很多昆蟲走獸,彼得說這是不潔淨「浴」物.神說我潔淨的,你不可當作俗物,這裡有兩個相對的意思:潔淨是歸主為聖,俗是平常.我們不要把自己的身體看作與世人一樣,因為神的靈住在裡面,我們不能把神的生命看作平常,與世人一樣；否則神會把我們看作該滅亡的世人一樣.
\newline
\newline
我們被主潔淨從世界分別出來的人,我們的動機,作事方式,生活方式,應該與世界不同,要過著分別為聖的生活,我們能否像世人一樣的生活嗎?我們能否再像世人作草木禾楷的生活,思想,態度呢?若我們看自己與世人一般,神會毀壞我們.生命是神的殿,而神的殿是聖的,所以讓我們時常謹慎自己的生活,用金銀寶石來建造自己的生命；經得起考驗,將來能在主面前存到永遠.
\newline
\newline


\newpage

\section{第53屆~港九培靈會~鮑會園~第四講}
\label{sec:565}
\textbf{錯誤的思想}
\newline
\newline
連結: \href{https://www.hkbibleconference.org/session-message/view/565}{\texttt{https://www.hkbibleconference.org/session-message/view/565}}
\newline
\newline
\hyperref[sec:559]{< < < PREV SERMON < < <}
~
\hyperref[sec:toc]{[返目錄]}
~
\hyperref[sec:567]{> > > NEXT SERMON > > >}
\newline
\newline
哥林多前書 4:1-4:13
\newline
\begin{longtable}{cl}
\hline
\hline
章節 & 經文 (和合本修訂版)\\
\hline
4:1 & \begin{tabularx}{0.7\textwidth}{X} 人應該把我們看為基督的執事，為神的奧祕的管家。 \end{tabularx} \\ \\ \relax
4:2 & \begin{tabularx}{0.7\textwidth}{X} 所求於管家的，是要他忠心。 \end{tabularx} \\ \\ \relax
4:3 & \begin{tabularx}{0.7\textwidth}{X} 我被你們評斷，或被別人評斷，我都以為是極小的事；連我自己也不評斷自己。 \end{tabularx} \\ \\ \relax
4:4 & \begin{tabularx}{0.7\textwidth}{X} 雖然我不覺得自己有錯，卻也不能因此判為無罪；審斷我的是主。 \end{tabularx} \\ \\ \relax
4:5 & \begin{tabularx}{0.7\textwidth}{X} 所以，時候未到，在主來以前甚麼都不要評斷，他要照出暗中的隱情，揭發人的動機。那時，各人要從神那裡得著稱讚。 \end{tabularx} \\ \\ \relax
4:6 & \begin{tabularx}{0.7\textwidth}{X} 弟兄們，為你們的緣故，我拿這些事應用到我自己和亞波羅身上，讓你們從我們學到「不可過於聖經所記」這話的意思，免得你們自高自大，看重這個，看輕那個。 \end{tabularx} \\ \\ \relax
4:7 & \begin{tabularx}{0.7\textwidth}{X} 使你與人不同的是誰呢？你所有的有哪一個不是領受的呢？若是領受的，為何自誇，彷彿不是領受的呢？ \end{tabularx} \\ \\ \relax
4:8 & \begin{tabularx}{0.7\textwidth}{X} 你們已經飽足了，已經富足了，用不著我們，自己就作王了。我願意你們果真作王，讓我們也可以與你們一同作王！ \end{tabularx} \\ \\ \relax
4:9 & \begin{tabularx}{0.7\textwidth}{X} 我想，神把我們作使徒的明顯地列在末後，好像定死罪的囚犯，因為我們成了一臺戲，給世界、天使和眾人觀看。 \end{tabularx} \\ \\ \relax
4:10 & \begin{tabularx}{0.7\textwidth}{X} 我們為基督的緣故成為愚拙的；你們在基督裡倒是聰明的。我們軟弱，你們倒強壯；你們有榮耀，我們倒被藐視。 \end{tabularx} \\ \\ \relax
4:11 & \begin{tabularx}{0.7\textwidth}{X} 直到如今，我們還是又飢又渴，又赤身露體，又挨打，又到處漂泊， \end{tabularx} \\ \\ \relax
4:12 & \begin{tabularx}{0.7\textwidth}{X} 並且勞碌，親手做工；被人咒罵，我們就祝福；被人迫害，我們就忍受； \end{tabularx} \\ \\ \relax
4:13 & \begin{tabularx}{0.7\textwidth}{X} 被人毀謗，我們就勸導。直到如今，人還把我們看作世上的污穢，萬物中的渣滓。 \end{tabularx} \\ \\
[1ex]
\hline
\hline
\end{longtable}
今天所讀的經文是整大段的結束,第一章講到哥林多教會的紛爭嫉妒,然後在第二章是保羅分析,因為他們很多人屬肉體,按照世人的方式行事；第三章指出他們不認識自己生命的目標.保羅在第四章解明他們所以有黨派之爭,是因為:
\newline
\newline
(一)他們對傳道人看法錯誤.
\newline
\newline
(二)他們對聖經看法錯誤.
\newline
\newline
(三)他們對自己看法錯誤.
\newline
\newline
(一)他們對傳道人看法錯誤
\newline
\newline
哥林多人把傳道人放在特別的地位,說我是屬磯法的,我是屬保羅的,我是屬阿波羅的.他們就以「他」為這派的領袖和榜樣,甚至與別人對抗,因「他」是他們最尊重的一個.保羅對他們說:你們根本錯了,因為不認識傳道人的工作和地位.「人應當以我為基督的執事,為神奧秘事的管家.」(四1)這節經文是特別指著全時間事奉的傳道人,神在眾人中特別呼召傳道人；要他整個生命放在全時間事奉上,特別在傳福音牧養教會方面事奉神.當然從另一個角度看,每一個基督徒都是事奉神的人,因為你不是事奉神就是事奉撒旦.只有這兩種勢力對立,沒有一個基督徒承認自己是事奉撒旦的人,所以無論作甚麼職業,我們都按照神的旨意生活,所以每一個基督徒都是事奉神的人.但是在所有事奉神的工作中,全時間事奉是最蒙福的工作；因他一生的工作,都直接與傳福音有關.「他所賜的有使徒,有先知,有傳福音的,有牧師和教師」(弗四11)所以我們不能忽略這些人的重要.
\newline
\newline
雖然我們看見有些教會沒有全時間做工的人,這個教會仍然很蒙福；但從整個聖經和教會歷史來看,這種情況是例外的.正如一個傳道人不受神學訓練,也可以傳福音,但這些人是例外.能夠沒有唸過大學,而可以教大學的人是很少的,王雲五只唸過小學,後來在大學文學院作院長,但全中國只有一個王雲五.傳道人把一生的生命放在事奉上,是不能兼職的工作,若是這樣,兩方面都不能完全.聖經上稱傳道人為神的僕役,乃是專責服事大君王的,所以必須專心工作,絕不能兼職.
\newline
\newline
但另一方面我們要看清楚傳道人的地位.哥林多人把傳道人看作黨派的領袖,這裡說傳道人是基督的執事.「執事」這個字比奴僕身份還要低,普通奴僕是用錢買來的.但在希臘人的社會裡,他們很喜歡坐船出海,搖船的奴僕是奴隸中最低賤的人.這些船的構造很特別,上面有甲板,下面有這些僕役坐在艙底兩邊；所以船上的人看不見他們,他們也看不見甲板上的人.但當掌舵的人打特別訊號給他們,他們便按命令來工作.這些搖船的人整個生命責任,就是聽從主人的訊號和命令.保羅說傳道人就是這種僕役,我們整個生活方式是按命令生活.不是要表現自己,建立甚麼黨派；我們最主要是聽從發命令的神,若果我失去主給我的身份,我便不配作基督的僕役.我們要留意基督的僕役,乃是要向基督負責任,而不是向任何人負責任.我不是要傳道人與教會不合作,因為我們要遵守組織的規定,工作才能順利進行；但我們基本工作方向乃是要聽從大君王耶和華的命令,向祂交賬.所以任何時候我在教會中看別人的臉色作事的話,我便不是基督的僕役.照樣,教會的弟兄姊妹亦應看重傳道人,傳道人的工作是按照基督的引導和命令；若教會要決定傳道人的工作方向,而不是讓他順服主的命令,這是錯了.哥林多教會把傳道人的身份看錯了,沒有看清楚傳道人是基督的執事,有了這個觀念,所以教會便受虧損.
\newline
\newline
(二)他們對聖經看法錯誤
\newline
\newline
「叫你們效法我們,不可過於聖經所記.」(四6),這句話很難解釋,是否聖經教導我們效法保羅,到一個地步就要停止?事實這裡的標點符號不清楚,其實保羅說:「你們效法我們,不可過於聖經所記.」這是解釋他們應該效法保羅的內容,在甚麼事情上效法我們?乃是不可過於聖經所記.保羅的生活單單按聖經的標準,所以在這一件事,你們該效法我.不是用哥林多人的智慧或知識來決定,這正是哥林多人的毛病,他們以自己的聰明及知識來領受神的話.在遵守神的話之外,再加上其他事情,他們把知識恩賜看著與主的話平等,甚至還高於主的話.人常常有這毛病,就是遵守神話語之外,再加上其他的事情；但聖經說不可在神的話語上增加或減少.夏娃就是犯了這個毛病,神吩咐她不可以吃分別善惡樹上的果子,但她說不可以吃也不可摸.今天有人批評基督教裡的異端邪教,這些異端都普遍有一個特徵；就是他們的信仰,在聖經以外,還加上其他東西.我們可以有把握地說,若有任何教會團體教導你遵守聖經；加上其他東西,這個團體便有問題.最近香港有一個邪教鼓勵人自殺,他們有很多教義,歷史,圖片,幻燈；在聖經以外,加上其他東西,這是錯誤.
\newline
\newline
有時候對聖經的解釋,超過聖經本身的意義,也是錯誤.把自己的解釋,加於聖經,也是超過聖經所記,聖經沒有這樣講,而我們堅持一定要作,這是錯誤.每個宗派及教會都有它自己的敬拜儀式,但若有人在敬拜上加上一些特別的儀式,看作基督徒必要的條件；你不按照我的方式,就不是教會或基督徒了,這就是超過聖經所記.我們需要認清楚我們基本生活標準,需要按照聖經的記載,也需要按照聖經的正確意義去解釋.哥林多教會的毛病就在這裡,他們把神的話看錯了.
\newline
\newline
(三)他們對自己看法錯誤
\newline
\newline
保羅問他們:「使你與人不同的是誰呢?你有甚麼不是領受的,為何自誇,彷彿不是領受的呢?」(四7),哥林多基督徒覺得自己比別人好,我是屬保羅黨的,我比阿波羅黨好；我是屬磯法黨的,我比保羅黨好.保羅說你們真正能得著的,是神的恩典,根本你們成為基督徒,都是領受的,但怎麼你們以自己所領受的自誇.恩典乃是我們不應得的好處,神給我們恩典,並不是我們比別人好；但我們不能以領受後誇口,這是多麼愚拙的事.四章七至十三節,保羅教導他們在主面前謙卑,他們高抬自己,其實根本他們沒有資格高抬自己,因為這是神的恩典.你們應該存感恩的心來領受恩典,認清自己與別人沒有不同；若有不同是神格外的恩典,所以要謙卑,承認自己毫無好處；便不會擁護哪一黨,哪一派.「謙卑」不是客氣,口頭上推讓的意思,真正的謙卑乃是認清自己的真正身份.我們原是蒙恩罪得赦免的罪人,若果我們有這種態度,便不會以任何事情來自誇.
\newline
\newline
保羅自己作榜樣,他說:「我們還是又飢又渴,又赤身露體,又挨打,又沒有一定的住處.並且勞苦,親手作工,被人咒罵,我們就祝福,被人逼迫,我們就忍受.被人毀謗,我們就善勸.直到如今,人還把我們看作世界上的污穢,萬物中的渣滓.」(四11-13)傳道人是耶穌的僕役,我們是人類中最低賤的僕役,若你們以我為黨派的領袖,你應該成為像我這樣謙卑的人.哥林多人若看清楚這點,他們原來所尊崇的人；乃是僕役中的僕役,他們便不再自誇了.如果認清這個真理,教會便不會有紛爭黨派了.
\newline
\newline
求神藉著保羅的勉勵,讓我們更清楚認定我們在神的家中應站的地位.
\newline
\newline


\newpage

\section{第53屆~港九培靈會~鮑會園~第五講}
\label{sec:567}
\textbf{淫亂的問題}
\newline
\newline
連結: \href{https://www.hkbibleconference.org/session-message/view/567}{\texttt{https://www.hkbibleconference.org/session-message/view/567}}
\newline
\newline
\hyperref[sec:565]{< < < PREV SERMON < < <}
~
\hyperref[sec:toc]{[返目錄]}
~
\hyperref[sec:569]{> > > NEXT SERMON > > >}
\newline
\newline
哥林多前書 5:1-5:13
\newline
\begin{longtable}{cl}
\hline
\hline
章節 & 經文 (和合本修訂版)\\
\hline
5:1 & \begin{tabularx}{0.7\textwidth}{X} 我確實聽說在你們中間有淫亂的事；這種淫亂連外邦人中也沒有，就是有人和他的繼母同居。 \end{tabularx} \\ \\ \relax
5:2 & \begin{tabularx}{0.7\textwidth}{X} 你們還自高自大！你們不是該覺得痛心，把做這事的人從你們中間趕出去嗎？ \end{tabularx} \\ \\ \relax
5:3 & \begin{tabularx}{0.7\textwidth}{X} 我人雖然不在你們那裡，心卻在你們那裡，好像親自與你們同在。我奉我們主耶穌的名，已經判斷了做這事的人。 \end{tabularx} \\ \\ \relax
5:4 & \begin{tabularx}{0.7\textwidth}{X} 你們聚會的時候，我的心和你們同在。你們藉著我們主耶穌的權能， \end{tabularx} \\ \\ \relax
5:5 & \begin{tabularx}{0.7\textwidth}{X} 要把這樣的人交給撒但，使他的肉體敗壞，好讓他的靈魂在主的日子可以得救。 \end{tabularx} \\ \\ \relax
5:6 & \begin{tabularx}{0.7\textwidth}{X} 你們這樣自誇是不好的。你們不知道一點麵酵能使全團發起來嗎？ \end{tabularx} \\ \\ \relax
5:7 & \begin{tabularx}{0.7\textwidth}{X} 既然你們是無酵的麵，要把舊酵除淨，好使你們成為新團；因為我們逾越節的羔羊—基督已經被殺獻為祭牲了。 \end{tabularx} \\ \\ \relax
5:8 & \begin{tabularx}{0.7\textwidth}{X} 所以，我們來守這節，不可用舊酵，就是不可用惡毒、邪惡的酵，只用純潔真實的無酵餅。 \end{tabularx} \\ \\ \relax
5:9 & \begin{tabularx}{0.7\textwidth}{X} 我先前寫信告訴過你們，不可與淫亂的人交往。 \end{tabularx} \\ \\ \relax
5:10 & \begin{tabularx}{0.7\textwidth}{X} 此話不是泛指這世上所有行淫亂的，或貪婪的，勒索的，或拜偶像的；若是這樣，你們非離開這世界不可。 \end{tabularx} \\ \\ \relax
5:11 & \begin{tabularx}{0.7\textwidth}{X} 但現在，我寫信告訴你們，若有稱為弟兄的人卻仍犯淫亂，或貪婪，或拜偶像，或辱罵，或醉酒，或勒索，這樣的人不可跟他交往，就是跟他吃飯都不可以。 \end{tabularx} \\ \\ \relax
5:12 & \begin{tabularx}{0.7\textwidth}{X} 因為審判教外的人與我何干？教內的人豈不是你們要審判嗎？ \end{tabularx} \\ \\ \relax
5:13 & \begin{tabularx}{0.7\textwidth}{X} 至於外人有神審判他們。如經上說：「要從你們中間把那邪惡的人趕出去。」 \end{tabularx} \\ \\
[1ex]
\hline
\hline
\end{longtable}
保羅在本章裡給我們非常嚴厲的教訓,如果我們明白他是對付罪時,便不會覺得他的教訓過份嚴厲了.有人批評聖經,因為聖經舊約記載了很多污穢犯罪的行為.不過聖經記載這些事與報紙雜誌記載的是有基本的分別.人看了報紙雜誌的淫亂故事,心裡便會引起污穢的思想；但讀了聖經所記載的......我們就有恨惡罪惡的感覺；永不會引起我們愛好這罪,也使我們對這罪覺得可怕可惡.保羅在這段經文裡,教導我們怎樣對付淫亂的事.
\newline
\newline
保羅所講的對象是哥林多教會,他吩咐哥林多教會不能容許這樣的罪存在在教會裡,因為這樣的事連外邦人也不會做.當然有很多的罪不分大小,在主面前同樣是罪.例如屬靈的冷淡,但冷淡的罪並沒有公開羞辱主的名；惟獨淫亂的罪,就公開羞辱了主的名.因為這些事傳到外邦人中,就羞辱了主的名,失掉了見證.所以教會不應該容許這樣的人存在.
\newline
\newline
另外方面,保羅提到哥林多教會的態度,是值得我們留意的.保羅說他們中間有淫亂的事,但他們沒有對付,而且還自高自大.自高自大可能指他們用別的事情來誇口,他們誇耀自己的恩賜才幹；我覺得保羅的意思,指出他們知道他們中間有人犯了罪,但仍然大量寬容,沒有對付這罪.可能他們說:「雖然這人犯了罪,但我們其實也犯過罪；只不過罪的性質不同而已,所以我們很大量寬容.」他們不知這羞辱主名的罪的嚴重,在他們中間不覺得罪的可怕,沒有為罪難過的心,這正是保羅指責他們的原因.這可能是今日教會最大的咒詛,把罪看作平常和無關重要.今日我們要在靈裡得著力量,在見證上發生作用,首先要對罪惡有靈敏的感覺.大衛是一個生活中很軟弱的人,他感情很衝動,他一生中犯了很多罪,甚至是很嚴重的罪,但聖經仍稱他為合神心意的人；因為他一時被肉體所勝,但他甚麼時候甦醒過來；就立即在神面前認罪痛哭流淚悔改.初期教會有許多特別經歷,使徒行傳第五章,記載亞拿尼亞和撒非喇的故事,兩人看見巴拿巴將田產賣了,把價錢拿到教會.也許他們覺得自己不好意思,別人奉獻了,而他們未有奉獻.既應奉獻,卻又捨不得,所以他們私自留下部份的金錢；結果神的靈啟示使徒知道,他們立即受神的刑罰,死在眾人面前.神的懲罰是非常嚴厲可怕的,因為教會必須聖潔,絕不能容許罪的存在.我們的教會雖然有同樣的罪存在,卻不覺得罪的可怕,不覺得在神面前需要悔改.所以保羅吩咐哥林多教會,不要自高自大.
\newline
\newline
我們看見這罪的性質和當時教會的需要,現在我們要看保羅吩咐教會,應如何對付這罪惡吧.「我身子雖不在你們那裡,心卻在你們那裡,好像我親自與你們同在,已經判斷行了這事的人.就是你們聚會的時候,我的心也同在,奉我們主耶穌的名,並用我們主耶穌的權能,要把這樣的人交給撒旦；敗壞他的肉體,使他的靈魂,在主耶穌的日子可以得救.」(3-5)首先保羅吩咐哥林多教會,他們聚會時,奉主耶穌的名應該怎樣作.教會中有人犯了公開羞辱主名的罪,這事不是一兩個人知道,乃是大家都知道,連外邦人也知道,我們不能在教會替他隱瞞,不能暗暗對付,因為這件事得罪了主,也是得罪了整個教會的弟兄姊妹.所以對付的時候要在整個教會公佈出來,保羅認為這不是愛心的問題；乃是整個教會的見證問題,應在教會中公開這罪.
\newline
\newline
「奉主耶穌的名」這句話還有另外意思,如果這件事在執事會,甚至全教會討論過,當決定的時候,需要整個教會來處罰他；而不是單單由傳道人一個作主.這樣有幾個好處,首先對付這罪的權柄在整個教會手上,讓整個教會知道這事是錯.其次讓犯事的人知道這是全教會所定的刑罰,而不是傳道人與他過不去.當然若弟兄姊妹不清楚真理,傳道人有責任教導,但應由執事會及全部弟兄姊妹作決定.第二點在作這事時候,是要奉主的名,表示這是主的意思.若不是奉主名而作,這就錯了!因此我們不能用別的理由把一個人從教會趕出去,乃是按照聖經的教訓,不能以自己的意思作標準.有時候在教會中,有人並不是違背聖經的教訓,只是因一兩個人的意見不合,便被趕出去,這是錯了.如果我不能說是奉主的名,而實際上不是主的意思,那我便是假傳聖旨了.
\newline
\newline
然後保羅說要把那人交給撒旦,敗壞他的肉體,這話意思很簡單清楚.要把這人交給撒旦,乃是放在教會以外,因為教會是屬神的.一個人信主後,是已經脫離了撒旦的權柄之下,進入神愛子的國度.但這人污穢了主的名,我們應該把他趕出去,使他不再有這機會及資格得到保護,他便落在撒旦的權勢之下.我們這樣作是很嚴厲,保羅說:「敗壞他的肉體」,表示讓他在罪惡的肉體去敗壞；因他在犯事前,他一定不是沒有聖靈的警告,聖靈會提醒他,只是他拒絕順從聖靈的感動.主耶穌在馬太福音第十八章的教訓也是差不多；若有人犯罪,先要暗暗向他儆告,如果他不悔改；就帶同教會一兩位領袖向他指責.他若悔改,你便得了一位弟兄；如果他不肯悔改,此後就要把他看作外邦人,不再看他是弟兄.但是另外一方面,他的肉體在罪惡中敗壞,使他的靈魂可以得救.這教訓與第三章一樣,若他的生命是草木禾楷建造,將來在審判的日子,他的工作會被燒毀,但他自己卻得救.雖然他的肉體敗壞,但他的生命仍然屬於主耶穌,我們把這個人從教會趕出去,但至終主仍拯救他的命,他會僅僅得救.
\newline
\newline
另外這樣作還有一個重大的目的(9-11),在這裡保羅似乎對基督徒比對非基督徒更加嚴厲,如果有非基督徒的,犯了污穢的事,我們仍須與他們接觸,例如你到米舖買米或到飯店吃飯,那位店主可能是個罪人,問題你是為交易與他來往,他是不是罪人也不要緊,你們來往不是要證明他是罪人.但是我們與基督徒來往就不同了；基督徒是在靈裡,交通來往的,你與他來往,就承認他是弟兄了.假如他現在犯了罪,你仍與他來往交通；可能給他有一種錯誤觀念,以為犯了罪也不要緊,因為主內弟兄仍接納我是基督徒,而教會仍然接納我是會友.保羅說這種人連跟他吃飯也不可.因為若與他吃飯,他便覺得自己雖然犯了罪仍然可被接納為弟兄,我們這樣做就等於鼓勵他仍在罪惡中生活了.但若果我與他吃飯,不是在屬靈的根基上；而是在自然和家庭關係上,這樣一起吃飯也沒問題.最近我聽見一位姊妹因她的丈夫在外面犯罪,她便不與丈夫吃飯要離開家庭,這種思想是錯誤的.
\newline
\newline
吃飯在原文意思,「不可與他一同吃」,可能這裡所指的吃,是指吃聖餐,指主的杯和餅.一個人若犯了罪,不聽勸告,不肯悔改；就可以停止他領聖餐.這一點有人不贊成,他們以為聖餐,乃是領受屬靈的生命,和蒙福的途徑；一個犯罪的弟兄,不可以領聖餐,豈不是失去了這種得力量悔改的機會嗎?我們要明白聖餐的意義,基本乃是與神在靈裡的交通,也是與同領聖餐的弟兄姊妹靈裡的交通,一個人明知故犯,怎可與神,與弟兄姊妹靈裡交通?若我們容納犯罪的人領聖餐,實際上是向他表示,你犯罪也不要緊,這等於鼓勵他仍在罪中生活,不需要對付罪了.另一方面,領受聖餐確是得福的來源,但聖經說要憑信心領受.聖餐不是特效藥,你有沒有信心不要緊,功效乃在特效藥內,並不是因有牧師祝福後,吃下了聖餐,屬靈的生命便健壯.其實最重要你在主面前對付自己的罪,憑信心與主有屬靈的交通,在這樣情況下,我們領受這杯和餅,纔能成為我們屬靈的福份.所以我們停止犯罪的人領聖餐,一方面使他知道自己犯罪,不配作一個基督徒,另外因他帶著罪來到聖餐桌前,他是不能夠蒙福的.
\newline
\newline
保羅不是沒有愛心,他在其他地方講到他內心非常的痛苦.照樣我們在教會中對付罪惡,為了愛心,需要在主面前付出代價；若我們責備他的罪,心靈沒有感覺痛苦,我就沒有資格來責備他的罪.保羅心靈痛苦,正是顯明了他的愛心,求主的話也激勵我們,叫我們成為教會的祝福.
\newline
\newline


\newpage

\section{第53屆~港九培靈會~鮑會園~第六講}
\label{sec:569}
\textbf{基督徒生活的原則}
\newline
\newline
連結: \href{https://www.hkbibleconference.org/session-message/view/569}{\texttt{https://www.hkbibleconference.org/session-message/view/569}}
\newline
\newline
\hyperref[sec:567]{< < < PREV SERMON < < <}
~
\hyperref[sec:toc]{[返目錄]}
~
\hyperref[sec:570]{> > > NEXT SERMON > > >}
\newline
\newline
哥林多前書 6:12-6:20
\newline
\begin{longtable}{cl}
\hline
\hline
章節 & 經文 (和合本修訂版)\\
\hline
6:12 & \begin{tabularx}{0.7\textwidth}{X} 「凡事我都可行」，但不是凡事都有益處。「凡事我都可行」，但無論哪一件，我都不受它的轄制。 \end{tabularx} \\ \\ \relax
6:13 & \begin{tabularx}{0.7\textwidth}{X} 「食物是為肚腹，肚腹是為食物」；但神要使這兩樣都毀壞。身體不是為淫亂，而是為主；主也是為身體。 \end{tabularx} \\ \\ \relax
6:14 & \begin{tabularx}{0.7\textwidth}{X} 神已經使主復活，也要用他自己的能力使我們復活。 \end{tabularx} \\ \\ \relax
6:15 & \begin{tabularx}{0.7\textwidth}{X} 你們豈不知道你們的身體是基督的肢體嗎？我可以把基督的肢體作為娼妓的肢體嗎？絕對不可！ \end{tabularx} \\ \\ \relax
6:16 & \begin{tabularx}{0.7\textwidth}{X} 你們豈不知道與娼妓苟合的，就是與她成為一體嗎？因為主說：「二人要成為一體。」 \end{tabularx} \\ \\ \relax
6:17 & \begin{tabularx}{0.7\textwidth}{X} 但與主聯合的，就是與主成為一靈。 \end{tabularx} \\ \\ \relax
6:18 & \begin{tabularx}{0.7\textwidth}{X} 你們要遠避淫行。人所犯的，無論甚麼罪，都在身體以外；惟有行淫的，是得罪自己的身體。 \end{tabularx} \\ \\ \relax
6:19 & \begin{tabularx}{0.7\textwidth}{X} 你們豈不知道你們的身體是聖靈的殿嗎？這聖靈是從神而來，住在你們裡面的。而且你們不是屬自己的人， \end{tabularx} \\ \\ \relax
6:20 & \begin{tabularx}{0.7\textwidth}{X} 因為你們是重價買來的。所以，要在你們的身體上榮耀神。 \end{tabularx} \\ \\
[1ex]
\hline
\hline
\end{longtable}
今天所讀的三段經文,是保羅所討論的一個題目,教我們如何生活,在吃,穿,娛樂上都有一些的標準.這些標準叫我們知道怎樣行,才可以配合基督徒的樣式.首先我們要思想哥林多城的背景,哥林多處在當時特別的環境,此城在道德上很敗壞,把某些事情看作很平常.當時有一句流行的俗語說:「一個人若生活敗壞,他就像哥林多人一樣.」當地的教會正將生活的注意力,去效法哥林多城作為標準.他們本來是屬乎主,應該有高尚的標準；可惜他們的生活被世人改變,很多事情與教外人沒有分別.這正是今天我們教會的危險和儆惕,忘記了我們基督徒的身份,去效法世人的生活標準和方式.
\newline
\newline
保羅提到吃祭偶像之物,和淫亂的事兩方面的問題.在一般哥林多人而言,淫亂和吃偶像之物,也是很普遍的事情.至於吃祭偶像之物,解經家有不同的解釋,有說他們宰殺的牛羊牲畜,都經過祭祀,才拿到市場去買.因此市場所賣的肉都是拜過偶像的.另有解經家說他們很熱心拜偶像,拜祭後,一部份的肉留在廟裡,另一部份的肉帶回家去.若他不喜歡吃,可以帶到市場去賣,因為拜祭過偶像,所以價錢較廉,很多人為了節省緣故,就買這些肉來吃.所以吃拜偶像食物是很平常的事,這也成為教會很嚴重的問題.基督徒可以吃這些祭偶像的肉嗎?非基督徒請基督徒吃這肉,基督徒怎麼辦?淫亂的事在基督徒立場上有甚麼看法呢?這三段經文都是講到應當怎樣行,簡單說可以有三方面的主題.
\newline
\newline
(一)為我們自己的緣故應該怎樣行(第六章)
\newline
\newline
「凡事我都可行,但不都有益處.凡事我都可行,但無論哪一件,我總不受他的轄制；食物是為肚腹,肚腹是為食物.」(六12-13)保羅論到一個要的原則,基督徒既已得救,在主耶穌裡有絕對的自由.食物不能救我,也不能叫我滅亡；萬物都是神的創造,若存感恩的心,甚麼都可以吃.首先保羅說「凡事」,這是不包括犯罪污穢的事情,乃是指本身沒有錯誤的事；它本身沒有犯罪,我們甚麼都可以作.早餐時你要吃粥,飯,中餐或西餐,你隨便吃甚麼都沒有關係,這不是錯誤的事.但你不可以偷竊或淫亂,保羅說這些事根本不應該行的.「身子不是為淫亂,乃是為主.」(六13).
\newline
\newline
凡事都可行,但卻有一些條件,這裡保羅說:「凡事我都可行,但不都有益處.凡事我都可行,但無論哪一件,我總不受他的轄制.」(六12)這兩個條件直接影響我們基督徒的生活.首先這裡所說的「益處」,即指對我們屬靈生命的益處；保羅說我們甚麼事都可以作,但卻不是每件事對我們屬靈生命有益處的,故此我們要避免那些對我們屬靈生命沒有益處的事情.看書是好的,看書本身不是罪；但不一定每一本書都對我們有幫助.實際上有些書籍對我們有害,所以我們要選擇.我們有自由可作甚麼事,但我們的時間很有限,故我們要選擇對我們屬靈有益處的事去作.
\newline
\newline
然後保羅說第二個原則「我總不受他的轄制,若我們作一些事,影響我們在主裡真正的自由；使我們屬靈生命受轄制,我便不作.這些事本身不是錯,但當我們覺得這些事是最重要,使我受捆綁,以致我不能作更有意義的工作,我便受這些事的轄制.譬如喝咖啡不是罪,但如果有某國家請我去領會；我要先問那裡有咖啡喝才去,那我便受了轄制.有一次我到一間教會培靈會講道,星期五晚那天聚會參加者少得可憐,後來該教會牧師說:「我知道今晚電視有很精彩的節目,所以今晚人少,情有可原.」凡事皆可行,但不可受它的轄制.
\newline
\newline
(二)為弟兄姊妹緣故我們應該怎樣行(第八章)
\newline
\newline
「論到祭偶像之物,我們曉得我們都有知識,但知識是叫人自高自大,惟有愛心能造就人.論到吃祭偶像之物,我們知道偶像在世上算不得甚麼.」(八1,4)在這裡「我們知道偶像算不得甚麼」,還有另外的翻譯「我們知道偶像根本不存在」,保羅並不是指那些石頭,木頭雕刻的偶像不存在；乃是說這些偶像所代表的神根本不存在.故此與那些拜祭過偶像的肉並沒有分別,也不會影響那些肉的質素和營養,因此我們可以吃.但保羅給我們一些原則,「但人不都有這種知識」.我們真正明白偶像是不存在的,但有些真正信主的人可能沒有這些知識；可能他們以前拜過這些偶像,仍覺得這些神是存在的.所以當他看見我吃拜過黃大仙的肉,他以為可以拜黃大仙；因為牧師可以作,我也可以作.我不能因我的自由,使軟弱的弟兄跌倒,因為這樣作,就是得罪弟兄亦即得罪基督.按知識我是知道,但為了愛心的緣故；甘心放棄自己的權利,所以我寧願不吃.很多時候我們太過抓緊自己的權利,所以失敗,我不因弟兄姊妹因我而跌倒.很多年前,我在芝加哥讀書的時候,該城有很多醉酒鬼；有些教會的弟兄便向這些人工作,其中有一間教會帶領了一位弟兄信主.這人因酗酒,弄致家破,失去職業,但信主後有很大改變；脫離了酒,有正當職業,與妻子和好,並加入教會福音隊向人傳福音.一天晚上,他們佈道後,一位教會領袖邀請他們到家中慶祝.因為這位領袖家境富裕,而在這種情形下,喝一點酒是沒有關係的,所以他便請這位弟兄喝.那位弟兄起初非常震驚和拒絕,但其他人也說,我們一起慶祝,喝一點酒是沒有關係的；這弟兄在不好意思之下,便喝了一小杯.很快這人又成為醉酒鬼,幾個月後,便失去職業和家庭.
\newline
\newline
第十章裡保羅說:若有非基督徒請我吃飯,我可以憑信心去,即使吃那些拜祭過偶像的食物,我覺得是沒有問題,可以吃.但倘若那非基督徒說這是拜過偶像的,我便不可吃.保羅解釋這不是基督徒裝假,乃是為那告訴你的人的良心緣故；因他說這是拜過偶像的,表示在他心目中,祭過偶像和未祭過偶像的食物是有分別的.所以為他良心的緣故,我便不吃.同樣原因,若碰巧買了祭過偶像的食物也不要緊.但基督徒不應參加廟裡的宴會,因為這種情況,不是單單吃那些肉,乃是直接參與祭偶像的活動.
\newline
\newline
(三)為神的緣故我們應該怎樣行(第十章)
\newline
\newline
「所以你們或吃或喝,無論作甚麼,都要為榮耀神而行.」(十31)不單注意我們靈性生命長進,不單注意對弟兄姊妹的影響；而且不論作甚麼事,都要為榮耀神.當我去看電影,打麻將,賽馬及任何事,應自問我所作的能否藉此榮耀神.「榮耀神」這三個字很難解釋,是否我們把更多的榮耀加給神呢?這根本是無可能,神有完全的榮耀,而我們是犯罪的人,怎能把榮耀加給神.聖經說榮耀神,乃是藉著我們的生活把神的榮耀反映出來,這是神造人的目的.一位畫家畫了一幅正在彈琴的美女圖畫,我們看見這幅畫,不是欣賞稱讚那位女子很美麗,她的姿態很好,乃是稱讚那位畫家的技術.神造我們,叫我們的生活在人面前把神的榮耀彰顯出來；若我們生活不似神,充滿了罪惡,不能反映神的榮美；我們便不能榮耀神,不配作神的兒女.「世人都犯了罪,虧缺了神的榮耀.」(羅三23)就是這個意思.很多時候我們忘記對神和自己的基本身份責任,作出與神形像不符合的事,每逢我們作甚麼事前,先自問作這事是否像神的形像；若我們不能夠,便不要作,所以我們要在生活中反映神的榮耀.
\newline
\newline


\newpage

\section{第53屆~港九培靈會~鮑會園~第七講}
\label{sec:570}
\textbf{婚姻的生活}
\newline
\newline
連結: \href{https://www.hkbibleconference.org/session-message/view/570}{\texttt{https://www.hkbibleconference.org/session-message/view/570}}
\newline
\newline
\hyperref[sec:569]{< < < PREV SERMON < < <}
~
\hyperref[sec:toc]{[返目錄]}
~
\hyperref[sec:571]{> > > NEXT SERMON > > >}
\newline
\newline
哥林多前書 7:1-7:17
\newline
\begin{longtable}{cl}
\hline
\hline
章節 & 經文 (和合本修訂版)\\
\hline
7:1 & \begin{tabularx}{0.7\textwidth}{X} 關於你們信上所提的事，男人不親近女人倒好。 \end{tabularx} \\ \\ \relax
7:2 & \begin{tabularx}{0.7\textwidth}{X} 但為了避免淫亂的事，男人當各有自己的妻子，女人也當各有自己的丈夫。 \end{tabularx} \\ \\ \relax
7:3 & \begin{tabularx}{0.7\textwidth}{X} 丈夫對妻子要盡本分；妻子對丈夫也要如此。 \end{tabularx} \\ \\ \relax
7:4 & \begin{tabularx}{0.7\textwidth}{X} 妻子對自己的身體沒有主張的權柄，權柄在丈夫；丈夫對自己的身體也沒有主張的權柄，權柄在妻子。 \end{tabularx} \\ \\ \relax
7:5 & \begin{tabularx}{0.7\textwidth}{X} 夫妻不可忽略對方的需求，除非為了要專心禱告，在兩相情願下暫時分房；以後仍要同房，免得撒但趁著你們情不自禁而引誘你們。 \end{tabularx} \\ \\ \relax
7:6 & \begin{tabularx}{0.7\textwidth}{X} 我說這話是出於容忍，不是命令。 \end{tabularx} \\ \\ \relax
7:7 & \begin{tabularx}{0.7\textwidth}{X} 我願眾人像我一樣；但是各人都有來自神的恩賜，一個是這樣，一個是那樣。 \end{tabularx} \\ \\ \relax
7:8 & \begin{tabularx}{0.7\textwidth}{X} 我對沒有嫁娶的和寡婦說，他們若能維持獨身像我一樣就好。 \end{tabularx} \\ \\ \relax
7:9 & \begin{tabularx}{0.7\textwidth}{X} 但他們若不能自制，就應該嫁娶，與其慾火攻心，倒不如結婚為妙。 \end{tabularx} \\ \\ \relax
7:10 & \begin{tabularx}{0.7\textwidth}{X} 至於那已經嫁娶的，我吩咐他們—其實不是我，而是主吩咐的：妻子不可離開丈夫， \end{tabularx} \\ \\ \relax
7:11 & \begin{tabularx}{0.7\textwidth}{X} 若是離開了，不可再嫁，不然要跟丈夫復和；丈夫也不可離棄妻子。 \end{tabularx} \\ \\ \relax
7:12 & \begin{tabularx}{0.7\textwidth}{X} 我對其餘的人說—是我，不是主說—倘若某弟兄有不信的妻子，妻子也情願和他一起生活，他就不可離棄妻子。 \end{tabularx} \\ \\ \relax
7:13 & \begin{tabularx}{0.7\textwidth}{X} 妻子有不信的丈夫，丈夫也情願和她一起生活，她就不可離棄丈夫。 \end{tabularx} \\ \\ \relax
7:14 & \begin{tabularx}{0.7\textwidth}{X} 因為不信的丈夫會因著妻子成了聖潔；不信的妻子也會因著丈夫成了聖潔。不然，你們的兒女就不潔淨了，但現在他們是聖潔的。 \end{tabularx} \\ \\ \relax
7:15 & \begin{tabularx}{0.7\textwidth}{X} 倘若那不信的人要離開，就由他離開吧！無論是弟兄是姊妹，遇著這樣的事都不必拘束。神召你們原是要你們和睦。 \end{tabularx} \\ \\ \relax
7:16 & \begin{tabularx}{0.7\textwidth}{X} 你這作妻子的怎麼知道不能救你的丈夫呢？你這作丈夫的怎麼知道不能救你的妻子呢？ \end{tabularx} \\ \\ \relax
7:17 & \begin{tabularx}{0.7\textwidth}{X} 無論如何，要照主所分給各人的恩賜和神所召各人的情況生活。我在各教會裡都是這樣規定的。 \end{tabularx} \\ \\
[1ex]
\hline
\hline
\end{longtable}
今天我們一同思想婚姻生活的事,「論到你們信上所題的事.」(七1),很明顯哥林多教會寫信給保羅,問及婚姻生活應採取甚麼方式,可能今天基督徒也面對同樣的問題.首先保羅說:「男不近女倒好」(七1),「我對著沒有嫁娶的寡婦,若他們常像我就好.」(七8),(七章25節)後來他論及童身,好像保羅對婚姻態度非常消極,婚姻是不得已的惡事；但卻可以肯定不是保羅在以弗所書,歌羅西書一貫的教訓.聖經將婚姻放在崇高的地位,保羅甚至在以弗所書把大丈夫與妻子的關係,比喻基督與教會的關係；多處地方指出婚姻是神的旨意.保羅在這裡對婚姻的態度很消極,原因在(七章26節),當時哥林多教會面對很多困難.其中一個是外教人的逼迫,很多基督徒要忍受苦難,甚至連生命也有危險.我想為此原因,在這種環境下,不結婚可以減少很多不必要的困難.但是這教訓並不適合任何的情況,所以我們應將整本聖經關於婚姻的教訓來看神的旨意.
\newline
\newline
婚姻是神聖的,婚姻是神所賜的福氣.保羅提及不結婚,童身,這是神給一些人特別的帶領和旨意.有些人神要他結婚,神要藉他們的家庭來見證神；另有些人神要他不結婚,藉獨身來見證事奉神.因此結婚或不結婚,都要追求神在我身上的旨意,順從神的帶領.結婚有婚姻生活中的試驗,需要神的恩典來得勝榮耀神；同樣不結婚也有他的試驗,神也有足夠的恩典,使他能勝得過,如果我們違背了神的恩典,走自己的道路,神不會賜給我們恩典和力量去勝過.為此,按照神的旨意,結婚或不結婚都是應當的.今天有些宗教團體高舉獨身,以為獨身更加聖潔,更專心為神工作,使徒保羅在(提前四章三節)說:禁止嫁娶的教訓是魔鬼的主意,不是聖經或神的意思.一個人的婚姻可以使事奉更榮耀神,若他們獨身,但思想受騷擾,反而不能專心事奉神.所以我們說在某種情況下,是最屬靈呢?順從神的旨意,是最聖潔和最屬靈的.
\newline
\newline
「男不近女倒好」(七1)這裡顯明指著婚姻關係裡的生活,而不是婚姻以外的男女關係,因為這根本是錯的.哥林多教會的情況,若不結婚會更好,「但要免淫亂的事,男子當各有自己的妻子,女子也當有自己的丈夫.」(七2),「倘若自己禁止不住,就可以嫁娶；與其慾火攻心,倒不如嫁娶為妙.」(七9),保羅講到今日生活中最基本的需要,性慾是神造人的一種本能；人身體成熟的時候,自然有性的慾望,所以性慾並不是罪.但若這慾望強到不能控制,用一種不合聖經的方法去滿足這種慾望,我們整個生命被這種慾望所控制和勝過,走上不應走的道路,這便是「慾火攻心」.神用婚姻來解決這種性慾,而且丈夫妻子的責任是去滿足對方,他們兩人成為一體,在肉體的結合是兩人成為一體的一部份.所以不可彼此虧欠(七4).丈夫和妻子之間應當互相體恤和適應,使兩人婚姻生活成為祝福而不是阻礙；若性生活成為夫妻兩人的負擔,兩人便錯了.
\newline
\newline
保羅又解釋在某種情況下可以分房,這三個條件,少一個也不可以.
\newline
\newline
第一,「兩相情願」,夫妻大家同意,不能單獨作決定；沒有一個人作主,必須兩相情願.
\newline
\newline
第二,「暫時的」.
\newline
\newline
第三,「專心禱告」,兩人心裡為某些事情有負擔禱告,才可以同意分房.
\newline
\newline
但到一個時候,禱告的負擔會過去,我們不能一生為某一件事專心禱告；當這負擔過去後,以後仍要同房,免得撒旦引誘你們.
\newline
\newline
婚姻生活是正常的生活,若有人以聖潔的藉口來阻礙這正常的生活,只是為了自己特別喜歡,這是違背神的旨意,便是不聖潔和錯誤的思想.今天有些人說馬利亞一生過童身的生活,高舉馬利亞到絕對聖潔的生活；她從來沒有與約瑟有關係,其實這是貶抑馬利亞的人格.馬利亞是聖潔的猶太女子,她最少與約瑟有十二年的婚姻生活,當時猶太女子盼望成為彌賽亞的母親,馬利亞最初不明白耶穌是誰,她已經與約瑟結婚了.所以她與約瑟有正常夫妻關係,高舉馬利亞的童身,這是玷污馬利亞的人格；因她失掉了妻子的身份,違背了神的旨意.
\newline
\newline
然後保羅提到婚姻中的特別困難,家庭中間若有不信的,怎麼辦?(七12-13)這裡指結婚後,夫妻中其中一個信了主；丈夫或妻子信主後,生活方式有改變,與配偶不同,可能引起夫妻不和.過往你的家庭喜歡打麻將,現在你信了主,便不打麻將,可能你的夫或妻子不高興.在這種情況,信主的應該:第一,不能在見證上妥協,不能因討好丈夫或妻子,而陪他打麻雀.如果主感動我知道這是錯的,我便不作了.第二,不能因他對我不滿意,為此我要與他分離.一方面我在信仰立場站穩,另一方面我要表現更大的愛心,溫柔,忍耐,謙卑,仍然盡我作妻子或丈夫的責任,帶領對方信主.我認識一位牧師,他們夫婦結婚前尚未信主,後來妻子信主了,丈夫便不高興,事事與她為難；但由於妻子生命有改變,對丈夫的溫柔忍耐,打動了他的心；後來他便信了主,而且奉獻作傳道,在一間教會事奉,有很多的成績.有時候妻子或丈夫對配偶謙卑溫柔,而尚未能結出果子來.十多年前一美國歌星信了主,他來過香港很多教會講見證,很多人得幫助.他們夫婦過去在娛樂圈生活,丈夫信主後,用他的恩賜來事奉神；但妻子仍然反對他,事事與他為難,丈夫上台見證,她在台下抽煙,直到如今,她仍未悔改信主,但是神用這位弟兄成為見證.若果你的配偶尚未信主,不要灰心難過,仍然保持你的見證.
\newline
\newline
但保羅說,若果不信的要求離開,你就讓他離開；他受不了這種生活,不是因你的見證生活不好,乃是你的聖潔成為他的壓力.他硬著心不肯悔改而要求離開,你應不必勉強.這種離開不是離婚,他故意離去,乃是反對你的屬靈見證；信主的仍無權利可以重婚,保羅說最好盼望他以後信了主,重新和好.但你現在不可因此作離婚,因為這婚姻關係仍然存在,在神面前所立的婚姻,雖然配偶仍未信主,但因婚約仍在神的面前,夫婦不論哪一方信主,家庭便蒙福.家庭生活建立在神面前,婚約成為聖潔,為此緣故,你們的兒女也成聖潔.即使你的配偶多年尚未信主,你們的兒女仍是聖潔的.我們應按照聖潔來對他們,因此我們不能忽略教導兒女真理的責任；夫婦兩人,只有一方是信主,信主的要在屬靈方面幫助兒女,我們不能勉強配偶負這個責任.我們要在主面前非常尊重我們的責任,怎樣教導我們的兒女,使他們在真理上有根有基.
\newline
\newline
今天的經文是非常難明,但我覺得保羅給我們重要的生活原則；願神的話幫助我們在家庭裡更加快樂去事奉神.
\newline
\newline


\newpage

\section{第53屆~港九培靈會~鮑會園~第八講}
\label{sec:571}
\textbf{屬靈的權柄}
\newline
\newline
連結: \href{https://www.hkbibleconference.org/session-message/view/571}{\texttt{https://www.hkbibleconference.org/session-message/view/571}}
\newline
\newline
\hyperref[sec:570]{< < < PREV SERMON < < <}
~
\hyperref[sec:toc]{[返目錄]}
~
\hyperref[sec:545]{> > > NEXT SERMON > > >}
\newline
\newline
哥林多前書 9:1-9:27
\newline
\begin{longtable}{cl}
\hline
\hline
章節 & 經文 (和合本修訂版)\\
\hline
9:1 & \begin{tabularx}{0.7\textwidth}{X} 我不是自由的嗎？我不是使徒嗎？我不是見過我們的主耶穌嗎？你們不是我在主裡面工作的成果嗎？ \end{tabularx} \\ \\ \relax
9:2 & \begin{tabularx}{0.7\textwidth}{X} 假若對別人來說，我不是使徒，對你們來說，我總是使徒；因為你們在主裡正是我作使徒的印證。 \end{tabularx} \\ \\ \relax
9:3 & \begin{tabularx}{0.7\textwidth}{X} 對那些質問我的人，這就是我的答辯。 \end{tabularx} \\ \\ \relax
9:4 & \begin{tabularx}{0.7\textwidth}{X} 難道我們沒有權利靠著傳福音吃喝嗎？ \end{tabularx} \\ \\ \relax
9:5 & \begin{tabularx}{0.7\textwidth}{X} 難道我們沒有權利帶著信主的妻子一起出入，如同其餘的使徒，和主的兄弟們，和磯法一樣嗎？ \end{tabularx} \\ \\ \relax
9:6 & \begin{tabularx}{0.7\textwidth}{X} 只有我和巴拿巴沒有權利不做工嗎？ \end{tabularx} \\ \\ \relax
9:7 & \begin{tabularx}{0.7\textwidth}{X} 有誰當兵而自備糧餉呢？有誰栽葡萄園而不吃園裡的果子呢？有誰牧養牛羊而不喝牛羊的奶呢？ \end{tabularx} \\ \\ \relax
9:8 & \begin{tabularx}{0.7\textwidth}{X} 我說這些話豈是照一般人的看法？律法不也是這樣說嗎？ \end{tabularx} \\ \\ \relax
9:9 & \begin{tabularx}{0.7\textwidth}{X} 就如摩西的律法記著：「牛在踹穀的時候，不可籠住牠的嘴。」難道神所掛念的是牛嗎？ \end{tabularx} \\ \\ \relax
9:10 & \begin{tabularx}{0.7\textwidth}{X} 他不全是為我們說的嗎？的確是為我們說的！因為耕種的要存著指望去耕種；收割的也要存著分享穀物的指望去收割。 \end{tabularx} \\ \\ \relax
9:11 & \begin{tabularx}{0.7\textwidth}{X} 我們既然把屬靈的種子撒在你們中間，若從你們收取養生之物，這還算大事嗎？ \end{tabularx} \\ \\ \relax
9:12 & \begin{tabularx}{0.7\textwidth}{X} 假如別人在你們身上享有這權利，何況我們呢？然而，我們並沒有用過這權利，倒是凡事忍受，免得基督的福音受到阻礙。 \end{tabularx} \\ \\ \relax
9:13 & \begin{tabularx}{0.7\textwidth}{X} 你們豈不知在聖殿供職的人吃聖殿中的食物嗎？在祭壇伺候的人分享壇上的供物嗎？ \end{tabularx} \\ \\ \relax
9:14 & \begin{tabularx}{0.7\textwidth}{X} 主也是這樣命令，要傳福音的人靠著福音養生。 \end{tabularx} \\ \\ \relax
9:15 & \begin{tabularx}{0.7\textwidth}{X} 但這權利我全然沒有用過。我寫這些話，並非要你們這樣待我，因為我寧可死也不讓人使我所誇的落了空。 \end{tabularx} \\ \\ \relax
9:16 & \begin{tabularx}{0.7\textwidth}{X} 我傳福音原沒有可誇耀的，因為我是不得已的，若不傳福音，我就有禍了。 \end{tabularx} \\ \\ \relax
9:17 & \begin{tabularx}{0.7\textwidth}{X} 我若甘心做這事，就有賞賜；若不甘心，責任卻已經託付給我了。 \end{tabularx} \\ \\ \relax
9:18 & \begin{tabularx}{0.7\textwidth}{X} 這樣，我的賞賜是甚麼呢？就是我傳福音的時候，使人不花錢得福音，免得我用盡了傳福音的權利。 \end{tabularx} \\ \\ \relax
9:19 & \begin{tabularx}{0.7\textwidth}{X} 我雖然是自由的，不受人管轄，但我甘心作了眾人的僕人，為贏得更多的人。 \end{tabularx} \\ \\ \relax
9:20 & \begin{tabularx}{0.7\textwidth}{X} 對猶太人，我就作猶太人，為要贏得猶太人；對律法以下的人，我雖不在律法以下，還是作律法以下的人，為要贏得律法以下的人。 \end{tabularx} \\ \\ \relax
9:21 & \begin{tabularx}{0.7\textwidth}{X} 對沒有律法的人，我就作沒有律法的人，為要贏得沒有律法的人；其實我在神面前，不是沒有律法，而是在基督的律法之下。 \end{tabularx} \\ \\ \relax
9:22 & \begin{tabularx}{0.7\textwidth}{X} 對軟弱的人，我就作軟弱的人，為要贏得軟弱的人。對甚麼樣的人，我就作甚麼樣的人。無論如何我總要救一些人。 \end{tabularx} \\ \\ \relax
9:23 & \begin{tabularx}{0.7\textwidth}{X} 凡我所做的，都是為福音的緣故，為要與人共享這福音的好處。 \end{tabularx} \\ \\ \relax
9:24 & \begin{tabularx}{0.7\textwidth}{X} 你們不知道在運動場上賽跑的，大家都跑，但得獎賞的只有一人？你們也要這樣跑，好使你們得著獎賞。 \end{tabularx} \\ \\ \relax
9:25 & \begin{tabularx}{0.7\textwidth}{X} 凡參加競賽的，在各方面都要有節制，他們不過是要得會朽壞的冠冕；我們卻是要得不會朽壞的冠冕。 \end{tabularx} \\ \\ \relax
9:26 & \begin{tabularx}{0.7\textwidth}{X} 所以，我奔跑，不像無目標的；我鬥拳，不像打空氣的。 \end{tabularx} \\ \\ \relax
9:27 & \begin{tabularx}{0.7\textwidth}{X} 我克制己身，使它完全順服，免得我傳福音給別人，自己反而被淘汰了。 \end{tabularx} \\ \\
[1ex]
\hline
\hline
\end{longtable}
在當時哥林多教會中,顯然有人詢問關於保羅的權柄,保羅對他們的嚴厲責備,是否單憑保羅的感覺,或憑他在他們中間的成就來說這些話.開始的時候,保羅似乎為作使徒的職份來辯護,但實際上他沒有這樣做.或者保羅憑他的工作成就來證明他是真正的使徒,保羅在哥林多教會有顯著的成就；因他一個人有很多人信了主,也因著他的工作哥林多教會被建立起來,異端也被攔阻辯駁；若要看工作果效,他的工作很顯著,但保羅也沒有用這個理由來證明他有這個權柄.這給我們有重要的原則,工作的果效仍不能證明保羅有使徒的資格,否則他們便不會疑惑保羅的權柄.
\newline
\newline
在今天的時代,只注重外表工作果效和表現；例如:教會發展,人數增加；自己多麼受人歡迎,但單憑這些仍不能證明我們有真正屬靈權柄.今天香港的異端,好像摩門教發展最快,所以單憑人數多不能證明有真正的能力和權柄.有時表面上工作沒有成就,仍不能否定沒有工作的能力在內.最近我讀到有關一對宣教士夫婦在中國西北的見證,他們晚年時在青海一帶工作；工作果子很少,信主的人一隻手可以數算,但因著他們的工作,當地福音之門被打開了.有些人一生在同一工作上,沒有真正的成就,但真正屬靈工作,不能單憑外表及果效判斷；很多時候我們太注重外表來證明工作的性質,或一個人內在屬靈力量,這是不正確的衡量方法.今天教會中最流行的學說理論,使我們走錯了方向,單單注重數量和數目.當保羅說及他真正的權柄來源後,哥林多人便不能再與他辯論,我想他們看完信後,會覺得羞愧!自己怎能用想法與保羅相比?後來在哥林多後書中,看到哥林多教會真正悔改,我想這地方的教訓會有一定的位置.
\newline
\newline
保羅在這裡責備非常嚴厲,但他沒有單單責備他們,而自己卻沒有這種生命過另外的生活方式.如果我們的見證要有效,我們的見證必須要出自裡面屬靈生命.主耶穌說學生不能大過先生,屬靈的事情也一樣,我們事奉神的人必要先達這地步,才能引導弟兄姊妹達到這地步；我們無可能帶領一班比自己生命更長進的弟兄姊妹.下面保羅分辯他屬靈權柄的來源是甚麼:
\newline
\newline
(一)傳福音的人有一定的權柄和應得的享受
\newline
\newline
「難道我們沒有權柄娶信主的姊妹為妻,帶著一同往來,彷彿其餘的使徒麼?」(九5)難道我傳福音是沒有權享受生活嗎?然後保羅說傳福音的,難道沒有權柄享受這些收穫嗎?他舉出很多理由(九7),證明這是人之常情,然後他舉出聖經的吩咐(九9),神也有清楚的指示和命令.「我們若把屬靈的種子撒在你們中間,就是從你們收割奉養肉身之物,這還算大事嗎?」(九11),保羅在這裡提出一個原則,事奉神的人應得他物質上的好處,傳福音的人應該靠福音養生.今天教會對這個問題很敏感,不容易講；我覺得今天教會對待傳道人應按照保羅所定的規矩來作.教會對傳道人的待遇能否足夠供養傳道人呢?如果把傳道人的待遇,與教會最高薪或最低薪的人是不當的,如果我們將傳道人薪酬與掃街倒垃圾的人相比,表明我們覺得傳道人工作與掃街的有同等價值,我們把傳道人工作看得不重要.
\newline
\newline
我們要記得保羅是向哥林多教會說的,但保羅從來沒有接受哥林多教會的供給.他卻願意接受馬其頓腓立比教會一次兩次的供給.保羅因哥林多教會整個思想仍在屬肉體的階段,恐怕他們誤解了供給傳道人的意思,所以保羅沒有使用這些權柄.保羅實在為了福音緣故,甘心放棄了這些權柄.如果我們要使屬靈生命有真正果效,我們必須要有保羅的思想,為福音緣故甘心放下我應有的好處.「甘心」(九17)這個字中文好像歡歡喜喜的意思,但其實原文是義務和自告奮勇.為了主緣故,我願否放棄一些具體的東西,這是很實際的問題,其實聖經的教訓是很實際的.(羅十二1)將身體獻上,這是很實際的,不但用心,也是用手,腳來獻給神.以前我們的身體作犯罪的行動,但現在獻給神來工作.具體地我可以為主做甚麼?一位牧師對一小孩說,你願意為主獻上甚麼,那小孩想一想後說:「如果我有一百萬,我願意獻給神.」牧師點點頭稱讚他,然後問他身上有多少,他說有十塊錢在身上.牧師要他奉獻這十塊錢,他不願意,牧師便不明白地問他:「既然你願意奉獻一百萬,為甚麼你不肯奉獻十塊錢?」那小孩說:「我奉獻一百萬,因為我沒有,但現在這十塊錢是我所有的.」
\newline
\newline
(二)保羅整個人生目標是盼望與人同得福音的好處(九19-23)
\newline
\newline
a.保羅為了與人同得福音的好處,他整個生活方式也就沒有固定,他沒有固定的愛好,沒有固定生活圈子,目的是叫人得到福音.香港教會與海外教會有所不同,海外的基督徒家庭生活,都環繞教會的工作活動而定；譬如我今天請客,我先思想教會中有沒有特別的活動如果有禱告會,我便不請客,不但一個家庭是這樣,而是很多家庭是這樣.但香港很少看見基督徒家庭按照教會的活動來編排生活；我們需要把教會看作我生命中的一部份,為主的緣故若有影響教會的事奉和活動,我們願否把一些應酬放棄?讓整個教會生活,成為我重要的部份.
\newline
\newline
b.保羅提及他沒有特別固定的生活圈子,向猶太人就作猶太人；向律法下的人就作律法下的人,為的要藉著我的生活去影響那些人的生命.我看見很多教會中有很多的小圈子,外人沒有辦法打進去；圈子以外的事情都與我無關,不關心其他人的需要.我們有責任在他們中間去過生活,叫他們與我同得福音的好處.今天還有很多基督徒,以為海外宣教是我基本責任以外的事情,這不是我基本,只是某些人才需要去作.其實宣教工作是我們每個人的責任,向所有人傳福音是基督徒的責任.
\newline
\newline
(三)保羅的生活目標有竭力追求的見證(九24-27)
\newline
\newline
保羅提及到一個人要得能壞的冠冕需要去角力爭戰,我們為得的是不能壞的冠冕,豈不要更加角力爭勝嗎?在這裡保羅用運動賽跑的比喻,哥林多人很熟識賽跑,但賽跑不是容易得獎的,賽跑需要留意幾件事情:
\newline
\newline
a.賽跑的人必須按照固定範圍和規則來跑,他不能亂跑,照樣基督徒屬靈生命也一樣,走天路也要認定天國的規則,才會有長進.你若生活自由隨便,你只能作一個馬馬虎虎的基督徒；但若你要作一個真正的基督徒,你就要按照聖經的規則而行.
\newline
\newline
b.賽跑的人需要事前準備,以前我讀大學時有一位同房同學,他是一個長跑好手.賽跑的人不作我們能夠作的事情,我們不會去作的事情他卻去作.在當地的比賽通常在春天舉行,但他整年不停地練習,冬天下雪的時候,我們要躲在房裡,但他卻每天在外練習.我留意他幾年來不喝汽水,因為汽水會影響他的氣力.他這樣作是希望得獎賞,同樣,我們若要得獎賞,也要訓練我們屬靈的生命.
\newline
\newline
保羅的整個生命和工作是有目標的,他不受肉體的限制.「所以我奔跑,不像無定向的；我鬥拳,不像打空氣的.我是攻克己身,叫身服我,恐怕我傳福音給別人,自己反被棄絕了.」(九26-27)保羅不願意傳福音給別人,自己反被棄絕了,我們也不要一生作個表面上的基督徒,我們也看過一些事奉神的人和基督徒；一生在教會裡工作,但結果晚年時被主丟棄,這是多麼可憐.
\newline
\newline
保羅對盤問的人的答辯,他屬靈生命有這樣的見證,他為主放棄自己的權利,他整個生活目標,是叫人與他同得福音的好處.他一生是過鍛鍊自己的生活,攻克己身,不讓肉體在生活中控制自己；將來見主面的時候,不被丟棄.願主使保羅的見證成為我們的激勵；使我們有這樣的心志事奉主.
\newline
\newline


\newpage

\section{第54屆~港九培靈會~唐佑之~第一講}
\label{sec:545}
\textbf{苦澀與甘甜}
\newline
\newline
連結: \href{https://www.hkbibleconference.org/session-message/view/545}{\texttt{https://www.hkbibleconference.org/session-message/view/545}}
\newline
\newline
\hyperref[sec:571]{< < < PREV SERMON < < <}
~
\hyperref[sec:toc]{[返目錄]}
~
\hyperref[sec:546]{> > > NEXT SERMON > > >}
\newline
\newline
以西結書 2:1-3:11
\newline
\begin{longtable}{cl}
\hline
\hline
章節 & 經文 (和合本修訂版)\\
\hline
2:1 & \begin{tabularx}{0.7\textwidth}{X} 他對我說：「人子啊，你站起來，我要和你說話。」 \end{tabularx} \\ \\ \relax
2:2 & \begin{tabularx}{0.7\textwidth}{X} 他對我說話的時候，靈進入我裡面，使我站起來，我就聽見他對我說話。 \end{tabularx} \\ \\ \relax
2:3 & \begin{tabularx}{0.7\textwidth}{X} 他對我說：「人子啊，我差你往悖逆我的國家，以色列人那裡去，他們是悖逆我的。他們和他們的祖先違背我，直到今日。 \end{tabularx} \\ \\ \relax
2:4 & \begin{tabularx}{0.7\textwidth}{X} 這些人厚著臉皮，心裡剛硬。我差你到他們那裡去，你要對他們說：『主耶和華如此說。』 \end{tabularx} \\ \\ \relax
2:5 & \begin{tabularx}{0.7\textwidth}{X} 他們是悖逆之家，他們或聽，或不聽，必知道在他們中間有了先知。 \end{tabularx} \\ \\ \relax
2:6 & \begin{tabularx}{0.7\textwidth}{X} 你，人子啊，雖有荊棘和蒺藜在你那裡，你又住在蠍子中間，總不要怕他們，也不要怕他們的話；他們雖是悖逆之家，但你不要怕他們的話，也不要因他們的臉色驚惶。 \end{tabularx} \\ \\ \relax
2:7 & \begin{tabularx}{0.7\textwidth}{X} 他們或聽，或不聽，你只管將我的話告訴他們；他們是極其悖逆的。 \end{tabularx} \\ \\ \relax
2:8 & \begin{tabularx}{0.7\textwidth}{X} 「但是你，人子啊，要聽我對你說的話，不要像那悖逆之家一樣悖逆，要開口吃我所賜給你的。」 \end{tabularx} \\ \\ \relax
2:9 & \begin{tabularx}{0.7\textwidth}{X} 我觀看，看哪，有一隻手向我伸來；看哪，手中有一書卷。 \end{tabularx} \\ \\ \relax
2:10 & \begin{tabularx}{0.7\textwidth}{X} 他在我面前展開書卷，它內外都寫著字，上面所寫的有哀號、嘆息、悲痛的話。 \end{tabularx} \\ \\ \relax
3:1 & \begin{tabularx}{0.7\textwidth}{X} 他對我說：「人子啊，要吃你所得到的，吃下這書卷；然後要去，對以色列家宣講。」 \end{tabularx} \\ \\ \relax
3:2 & \begin{tabularx}{0.7\textwidth}{X} 於是我張開了口，他就使我吃這書卷。 \end{tabularx} \\ \\ \relax
3:3 & \begin{tabularx}{0.7\textwidth}{X} 他對我說：「人子啊，要吃我所賜給你的這書卷，塞滿你的肚腹。」我就吃了，口中覺得其甜如蜜。 \end{tabularx} \\ \\ \relax
3:4 & \begin{tabularx}{0.7\textwidth}{X} 他對我說：「人子啊，你要到以色列家那裡去，對他們傳講我的話。 \end{tabularx} \\ \\ \relax
3:5 & \begin{tabularx}{0.7\textwidth}{X} 你奉差遣不是往那說話艱澀、言語難懂的民那裡，而是往以色列家去； \end{tabularx} \\ \\ \relax
3:6 & \begin{tabularx}{0.7\textwidth}{X} 你不是往那說話艱澀、言語難懂的許多民族那裡去，他們的話你不懂。然而，我若差你往他們那裡去，他們會聽從你。 \end{tabularx} \\ \\ \relax
3:7 & \begin{tabularx}{0.7\textwidth}{X} 以色列家卻不肯聽從你，因為他們不肯聽從我；原來以色列全家是額頭堅硬、心裡剛愎的人。 \end{tabularx} \\ \\ \relax
3:8 & \begin{tabularx}{0.7\textwidth}{X} 看哪，我使你的臉堅硬，對抗他們的臉；使你的額頭堅硬，對抗他們的額頭。 \end{tabularx} \\ \\ \relax
3:9 & \begin{tabularx}{0.7\textwidth}{X} 我使你的額頭像金剛石，比火石更堅硬。他們雖是悖逆之家，但你不要怕他們，也不要因他們的臉色而驚惶。」 \end{tabularx} \\ \\ \relax
3:10 & \begin{tabularx}{0.7\textwidth}{X} 他又對我說：「人子啊，我對你說的一切話，你心裡要領會，耳朵要聽。 \end{tabularx} \\ \\ \relax
3:11 & \begin{tabularx}{0.7\textwidth}{X} 要到被擄的人，到你本國百姓那裡去，他們或聽，或不聽，你要對他們宣講，告訴他們這是主耶和華說的。」 \end{tabularx} \\ \\
[1ex]
\hline
\hline
\end{longtable}
教會在這時代須接受時代的挑戰.教會在這時代到底應該怎樣明白神的心意呢?我們就應該讀以西結書.
\newline
\newline
以西結是神的工人,是先知,他站在歷史的苦難危機中,他在神前領受並在人前傳揚.以西結素稱為「被擄的先知」.
\newline
\newline
為何以色列人遭受擄掠到外邦?為何他們要面對歷史的苦難?是因他們犯了罪,不能滿足神的心意；雖然他們是神的選民,他們在神的恩惠範圍裡；但神公義的,不能以有罪為無罪.
\newline
\newline
我們是屬於神的,正如以色列人屬於神一樣；如果我不能滿足神的心意,我們也要面對神公義的審判.
\newline
\newline
聖經告訴我們,審判要從神的家起首.(彼前五17)神要審判以色列人,然後審判外邦人；神要審判教會,然後審判整個世界.
\newline
\newline
以西結約在主前五百多年時,出來傳神的信息.神的話不但向當時代發出,且在每個時代都有同樣的信息,甚至此時此地,神同樣用祂的話告訴我們.先知的聲音沒有淹沒在時代的廢墟中,反而經過時代的走廊發出巨響；經過這個時代,經過現代每個角落進入人心,發生回響.先知以西結當時對以色列人所說的話,也是對我們說的；雖然是對當時的以色列人所說的,但也是對我們今日的教會說的.神的真理永遠長存,神的話語永不過去,神的話帶著能力,至今仍發出功效.審判要從神的家起首,不但在主前五百多年之先,甚至在主後的今天以至將來,我相信神不斷要對祂的教會說話.
\newline
\newline
以西結在神面前受了感動,當他看見異象之時,有隻手伸出來,手上有書卷；神要他把書卷吃了,這就是今時代教會應該明白的信息.以色列人失敗了,所以神要審判他們；每個時代教會失敗了,所以審判要從神的家起首.教會在這時代,要面對這個審判；因為神審判世界之前,要審判教會.在啟示錄,神審判世界時,先要審判教會,審判亞細亞的七個教會.示每拿教會沒有受審判,因他們在苦難當中仰望神,神就憐憫他們,幫助他們.非拉鐵非教會沒有受審判,因他們傳福音,神覺得無須審判他們；但是,神卻要審判其他的教會.
\newline
\newline
今天世界很多地方的人,甚至我們的同胞,也有多苦難的.很多教會在苦難中,他們仰望神,他們不須受審判；但是你我,很多教會,在這時代,包括自由世界的教會須受審判；因為審判要從神的家起首.先知有此感受,他是神的僕人,他有時代的知識；他知道在那個時代的環境,知道歷史中有神的手,在歷史過程中沒有一件事是偶然發生的.歷史的每個危機,都是神的審判；所以他在神前體會到,他應作當代的見證人.他須在神前領略屬靈的感受,他必須把書卷吃下.由此可知神的話何等重要!多麼有能力!
\newline
\newline
在士師和撒母耳時代,當時耶和華的言語稀少,不常有默示.難道神的話不在當代傳出來嗎?難道神的心意不願表達出來嗎?不是的,乃是因祂找不到合乎祂心意的人.當撒母耳屬靈知識還不甚充分的時代；神選召他成為當代的見證人.神要在我們中間尋找為祂傳話的人.今天的撒母耳在哪裡?今天的以西結在哪裡?
\newline
\newline
一個真正願意為神發言的人,要具有時代的意識,屬靈的感受.神的手伸出來,我們感受得到嗎?我們看得見嗎?我們能看見別人所看不到的,感受別人所不能感受的嗎?以西結就有這樣的經驗；神要他把書卷吃下,他吃時感到書卷滿有甘甜；因為神的話比一切都有意味.
\newline
\newline
在今天的世界裡,有許多人得不到神的話；是因為缺少見證人,沒有為神傳話的人.有的人聽不見神的話；因為看不見真理在我們基督徒身上表達出來,這就是當時以色列人失敗的原因.
\newline
\newline
神揀選以色列人,是否他們比別的民族優秀?是否他們的文化道德超越於其他民族?絕對不是.以色列不過是中東地帶,一個弱小的民族,幼稚而沒有文化；埃及,巴比倫諸國的文化比以色列更優秀.神所以揀選以色列民,乃要他們事奉神,神恩待他們,乃要他們得到恩典之後,把恩典帶出去；因為神的救恩不是單為以色列人,神的救恩是為全世界的.
\newline
\newline
以色列人失敗了,他們在恩典中墮落了.你我也是如此,神揀選我們作祂的兒女,並非你我比別人強,唯一的說明,乃本於神的恩典.我們是否在恩典中長進呢?或是在恩典中墮落呢?我們想到神的恩典,神真是奇妙的神,祂的恩典白白賜給我們；不須我們付上任何代價,不須我們任何功德,儘可白白享受.神的恩典太多,我們卻忽略了神的恩典,沒有珍貴神的恩典.恩典雖可白白得來；但是貴重無比,我們並無法以任何代價換取；所以神把恩典白白地賜給我們.可惜!我們在恩典中失落了,沒認真在恩典中長進.
\newline
\newline
神的教會所需要的是公義,審判的信息；就是昔日神要以西結傳出的信息.神的公義裡面有無限的慈愛,慈愛裡面有許多公義,且兩者分不開.神吩咐以西結要向以色列人傳道；並非到言語難懂的外邦人中去.
\newline
\newline
教會在這時代須有屬靈的敏感,教會須作差傳,須傳福音；但我相信,若我們沒有在神面前切實悔改,切實復興；要做差傳工作,福音工作是不可能的.神告訴以西結,以色列人很硬心,不願意接受神的恩典；他們在屬靈方面麻木不仁.所以神的僕人先知以西結,必須具有屬靈的感受,屬靈的敏感力.教會也必須有屬靈的敏感,靈敏感受神的恩典；否則,神的審判就會臨到.若我們沒有受過神的管教,怎能澈底悔改呢?所以我們多麼需要在神前思想我們的光景.神的心意在這時代,特別要叫我們明白祂的公義,明白審判從神的家起首.須在神前切實悔改,然後我們纔能復興；承受神的使命,有力量把神的恩典帶出去.這樣,福音纔能發生功效,差傳工作纔有果效；不然,恐怕我們就不能切實達到神的心意.
\newline
\newline
所以以西結先知必須開口吃書卷,因書卷充滿了神的話；先知切實吃了下去,先有屬靈的感受,然後纔可以把神的話帶出來,這樣的話語纔有能力,纔有功效,纔能到達與很多需要的人.
\newline
\newline
教會在這時代,對神的話語實在太膚淺,太幼稚；沒有真正的得著,沒有真正的領受,沒有真正的享受,也沒有真正的承受.故此,神的話語就被我們綁著,沒有出去.
\newline
\newline
有一位文學家曾說,世上有兩種人,一種是無話可講說的,但卻滔滔不絕,頭頭是道.其實廢話連篇,毫無內容,全無真理；可是卻有很多人聽,很多人接受.另外一種人有話可講,話有內容,有真理；但他們卻保持緘默,令人不解其意.(這樣的人乃指基督徒而言)基督徒有話語,有信息,可是不講.這個世界充滿了虛謊,所見所聞都是世俗淫亂,虛謊,嘆息,哭泣悲哀的聲音；我們怎麼還保持緘默呢?教會應該成為先知的教會,成為先知的團體.我們每個人都應領受神的話語,將神的話語吃下去；然後發出來,這樣就必看見神話語的功效.每個信徒都有責任,將神的話傳揚出來.
\newline
\newline
先知以西結吃了書卷,他感到口中非常甘甜.各位!我們有這經驗嗎?這是今天教會應有的經歷.很多時候,因為我們沒有真正感到神話語的甘甜,以致我們在恩典中墮落；在真理中矛盾,在信仰中逃避,在現實中被腐蝕.
\newline
\newline
何以我們會在恩典中墮落,乃因沒有讓神的話進入我們裡面,不明白也不體會恩典的實際,對神的話一知半解,所了解的不過一鱗半爪；沒有真正在神的話語上下工夫,沒有真正將神的話如書卷吃下去.我們因為忽略了神的話語,所以在恩典中墮落.我們往往知道了一點真理,卻不甚清楚,以為聖經的話不過是理論,演講,說說聽聽而已.我們不能把真理應用在實際生活中,在教會聽道時.雖然覺得真理是對的；出了教會,因面對種種的挑戰,看見世俗的情況,於是心裡又發生矛盾.覺得神的話不夠實際,不能應用在我們實踐的生活中.所以就在恩典中墮落,在真理中矛盾,在信仰中逃避.另外有種信念,雖然知道神是慈愛公義信實的,但只是一種觀念,沒有在生活中去真正體驗；結果,信心就呈逃避的態度.而且我們在需要時,纔鍜煉一下信心；以為信心具魔術性,可改變事情.其實不然,信心乃是神的恩典,聖靈的感動,加以我們需要在神的話語中切實下工夫；然後我門在主的恩典和知識上有長進.如果我們在恩典中墮落,在真理中矛盾,在信仰中逃避；我們定會在現實中腐蝕；因我們的生命不堅強,太軟弱,經不起考驗,我們生活的力量就越加脆弱了!這是教會的現象嗎?但願不是吧!可惜事實卻如此!我們當在神前切實悔改,若不在神的真道上建立自己,怎能知道要在聖靈中禱告呢?又怎會知道仰望神的憐憫呢?
\newline
\newline
教會在這時代,需要把書卷吃下,纔能覺得甜蜜.我們也需要有這經驗,否則,神的審判和公義就會臨到我們身上.「耶和華啊!你若究察罪孽,誰能站得住呢?然而你有赦免之恩.」我們需要神的赦免；因我們在話語上失敗,輕忽,必須重新奮興；讓聖靈在我們心裡動工,祂的話無不帶著能力；一出來就有功效,一解開就發出亮光,一進入就覺甘甜.
\newline
\newline
另一方面,我們也看見先知以西結對神的話有苦澀的感受.神的話語,神的信息,一方面表達祂的慈愛,恩典和公義；一方面也表達祂的震怒和說不出來忿怒的態度.「......其上所寫的有哀號,嘆息,悲痛的話.」二-10)吃下時是甘甜的,傳出來時是苦澀的.這是先知的感受,你呢?當我們得到神的話語時,我們覺得甘甜；神有說不盡的恩慈,神有無限的愛；但另一方面,神的話不是福音,不是佳音,乃是禍音,叫人感到害怕,痛苦.是否我們傳福音太隨意,單講神的慈愛,沒有講神的公義.先講福音,再講禍音,後講佳音；然後我們纔領會神無限的慈愛.
\newline
\newline
怎樣吃書卷,茲題出三方面.我們讀神的話語,必定要注意三方面.
\newline
\newline
(一)必定要默念真理的奇妙,一遍又一遍反覆細讀
\newline
\newline
如詩篇第一篇說:「惟喜愛耶和華的律法,晝夜思想,這人便為有福.他要像一棵樹栽在溪水旁,按時結果子......」晝夜不住思想默念神恩典的奇妙,就是我們吃書卷的經驗.
\newline
\newline
(二)尋求話語的引導,求問神今天對我說甚麼
\newline
\newline
因為神的心要我們明白,祂所表達的話語要引導我們,叫我們從神的話語,來思想祂的心意.
\newline
\newline
有的人說,信心不該有理智,有理性的信心是不堅定的；但是信心若沒有理性就是迷信.有許多基督徒是迷信的,如果沒有把悟性放在神裡面,沒有把理性放在神手裡；基督徒的信心很可能會變成迷信.信仰不是叫我們逃避,而是叫我們面對現實,幫助我們勇往直前.
\newline
\newline
(三)思想神的作為,明白神的心意
\newline
\newline
神在這時代,在此時此地的環境,叫我們明白教會在這個時代要作甚麼,我們在這時代能作甚麼?
\newline
\newline
我們尋求神的引導,明白神的心意；然後,我們的禱告,如詩人在詩篇119篇32節的禱告:「你開廣我心的時候,我就往你命令的道上直奔.」
\newline
\newline


\newpage

\section{第54屆~港九培靈會~唐佑之~第二講}
\label{sec:546}
\textbf{離愁與復合}
\newline
\newline
連結: \href{https://www.hkbibleconference.org/session-message/view/546}{\texttt{https://www.hkbibleconference.org/session-message/view/546}}
\newline
\newline
\hyperref[sec:545]{< < < PREV SERMON < < <}
~
\hyperref[sec:toc]{[返目錄]}
~
\hyperref[sec:547]{> > > NEXT SERMON > > >}
\newline
\newline
以西結書 9:3-9:19
\newline
\begin{longtable}{cl}
\hline
\hline
章節 & 經文 (和合本修訂版)\\
\hline
9:3 & \begin{tabularx}{0.7\textwidth}{X} 在基路伯之上，以色列神的榮耀從那裡上升，到殿的入口處。神召那身穿細麻衣、腰間繫著墨盒的人前來。 \end{tabularx} \\ \\ \relax
9:4 & \begin{tabularx}{0.7\textwidth}{X} 耶和華對他說：「你去走遍耶路撒冷全城，那些為城中所做可憎之事嘆息哀哭的人，你要在他們額上做記號。」 \end{tabularx} \\ \\ \relax
9:5 & \begin{tabularx}{0.7\textwidth}{X} 我耳中聽見耶和華對其餘的人說：「要跟隨他走遍全城去擊殺。你們的眼不要顧惜，也不要可憐他們。 \end{tabularx} \\ \\ \relax
9:6 & \begin{tabularx}{0.7\textwidth}{X} 要將年老的、年輕的、少女、孩童和婦女，從我的聖所開始全都殺盡，只是不可挨近凡有記號的人。」於是他們從殿前的長老殺起。 \end{tabularx} \\ \\ \relax
9:7 & \begin{tabularx}{0.7\textwidth}{X} 他對他們說：「要使這殿污穢，使院中遍滿被殺的人。你們出去吧！」他們就出去，在城中擊殺。 \end{tabularx} \\ \\ \relax
9:8 & \begin{tabularx}{0.7\textwidth}{X} 他們擊殺的時候，只剩我一人，我就臉伏在地上，呼喊說：「唉！主耶和華啊，你將憤怒傾倒在耶路撒冷，豈要把以色列所剩餘的人都滅絕嗎？」 \end{tabularx} \\ \\ \relax
9:9 & \begin{tabularx}{0.7\textwidth}{X} 他對我說：「以色列家和猶大家的罪孽極其重大。遍地都有流血的事，滿城有冤屈，因為他們說：『耶和華已經離棄這地，他看不見我們。』 \end{tabularx} \\ \\ \relax
9:10 & \begin{tabularx}{0.7\textwidth}{X} 因此，我的眼必不顧惜，也不可憐他們，要照他們所做的報應在他們頭上。」 \end{tabularx} \\ \\ \relax
9:11 & \begin{tabularx}{0.7\textwidth}{X} 看哪，那身穿細麻衣、腰間繫著墨盒的人回覆這事說：「我已經照你所吩咐的做了。」 \end{tabularx} \\ \\ \relax
10:1 & \begin{tabularx}{0.7\textwidth}{X} 我觀看，看哪，在穹蒼之中，也就是基路伯的頭上，有藍寶石的形狀，彷彿寶座的形像顯在他們上面。 \end{tabularx} \\ \\ \relax
10:2 & \begin{tabularx}{0.7\textwidth}{X} 耶和華對那身穿細麻衣的人說：「你進到基路伯下面旋轉的輪子中，從基路伯之間取出火炭裝滿兩手掌，撒在城上。」我親眼看見他進去。 \end{tabularx} \\ \\ \relax
10:3 & \begin{tabularx}{0.7\textwidth}{X} 那人進去的時候，基路伯站在殿的南邊，雲彩充滿了內院。 \end{tabularx} \\ \\ \relax
10:4 & \begin{tabularx}{0.7\textwidth}{X} 耶和華的榮耀從基路伯那裡上升，到殿的入口處；殿內滿佈雲彩，院子也充滿了耶和華榮耀的光輝。 \end{tabularx} \\ \\ \relax
10:5 & \begin{tabularx}{0.7\textwidth}{X} 基路伯翅膀的響聲傳到外院，好像全能神說話的聲音。 \end{tabularx} \\ \\ \relax
10:6 & \begin{tabularx}{0.7\textwidth}{X} 耶和華吩咐那身穿細麻衣的人說：「要從基路伯之間旋轉的輪子中取火。」那人就進去站在一個輪子旁邊。 \end{tabularx} \\ \\ \relax
10:7 & \begin{tabularx}{0.7\textwidth}{X} 基路伯中的一個基路伯伸手到基路伯中間的火那裡，取一些放在那身穿細麻衣人的手掌中，那人拿了就出去。 \end{tabularx} \\ \\ \relax
10:8 & \begin{tabularx}{0.7\textwidth}{X} 在基路伯翅膀以下，顯出有人手的樣式。 \end{tabularx} \\ \\ \relax
10:9 & \begin{tabularx}{0.7\textwidth}{X} 我又觀看，看哪，這些基路伯的旁邊有四個輪子。一個基路伯旁有一個輪子，另一個基路伯旁也有一個輪子；輪子的形狀好像水蒼玉石。 \end{tabularx} \\ \\ \relax
10:10 & \begin{tabularx}{0.7\textwidth}{X} 至於四輪的形狀，都是一個樣式，好像輪中套輪。 \end{tabularx} \\ \\ \relax
10:11 & \begin{tabularx}{0.7\textwidth}{X} 輪子行走的時候，向四方都能直行，行走時並不轉彎。頭轉向何方，它們也隨著向何方行走，行走時並不轉彎。 \end{tabularx} \\ \\ \relax
10:12 & \begin{tabularx}{0.7\textwidth}{X} 基路伯的全身，連背帶手和翅膀，並輪子周圍都佈滿眼睛。他們四個的輪子都是如此。 \end{tabularx} \\ \\ \relax
10:13 & \begin{tabularx}{0.7\textwidth}{X} 我耳中聽見這些輪子稱為「旋轉的輪」。 \end{tabularx} \\ \\ \relax
10:14 & \begin{tabularx}{0.7\textwidth}{X} 基路伯各有四張臉：第一是基路伯的臉，第二是人的臉，第三是獅子的臉，第四是鷹的臉。 \end{tabularx} \\ \\ \relax
10:15 & \begin{tabularx}{0.7\textwidth}{X} 基路伯升上去了；這就是我在迦巴魯河邊所看見的活物。 \end{tabularx} \\ \\ \relax
10:16 & \begin{tabularx}{0.7\textwidth}{X} 基路伯行走，輪子也在旁邊行走。基路伯展開翅膀，離地上升，輪子也不轉離他們的旁邊。 \end{tabularx} \\ \\ \relax
10:17 & \begin{tabularx}{0.7\textwidth}{X} 基路伯站住，輪子也站住；基路伯上升，輪子也跟著上升，因為活物的靈在輪中。 \end{tabularx} \\ \\ \relax
10:18 & \begin{tabularx}{0.7\textwidth}{X} 耶和華的榮耀離開殿的入口處，停在基路伯之上。 \end{tabularx} \\ \\ \relax
10:19 & \begin{tabularx}{0.7\textwidth}{X} 基路伯展開翅膀，在我眼前離地上升；他們離去的時候，輪子在旁邊，都停在耶和華殿的東門口。在他們上面有以色列神的榮耀。 \end{tabularx} \\ \\ \relax
10:20 & \begin{tabularx}{0.7\textwidth}{X} 這是我在迦巴魯河邊所見的活物，他們在以色列神之下；因此我知道他們是基路伯。 \end{tabularx} \\ \\ \relax
10:21 & \begin{tabularx}{0.7\textwidth}{X} 他們各有四張臉、四個翅膀，翅膀以下有人手的樣式。 \end{tabularx} \\ \\ \relax
10:22 & \begin{tabularx}{0.7\textwidth}{X} 至於他們臉的模樣，以及身體的形像，正是我從前在迦巴魯河邊所看見的。他們各自往前直行。 \end{tabularx} \\ \\ \relax
11:1 & \begin{tabularx}{0.7\textwidth}{X} 靈將我舉起，帶我到耶和華聖殿面向東方的東門。看哪，門口有二十五個人。我見其中有百姓的領袖押朔的兒子雅撒尼亞和比拿雅的兒子毗拉提。 \end{tabularx} \\ \\ \relax
11:2 & \begin{tabularx}{0.7\textwidth}{X} 耶和華對我說：「人子啊，他們就是圖謀罪孽，在這城中設計惡謀的人。 \end{tabularx} \\ \\ \relax
11:3 & \begin{tabularx}{0.7\textwidth}{X} 他們說：『蓋房屋的時候尚未臨近；這城是鍋，我們是肉。』 \end{tabularx} \\ \\ \relax
11:4 & \begin{tabularx}{0.7\textwidth}{X} 人子啊，因此你當說預言，說預言攻擊他們。」 \end{tabularx} \\ \\ \relax
11:5 & \begin{tabularx}{0.7\textwidth}{X} 耶和華的靈降在我身上，對我說：「你當說，耶和華如此說：以色列家啊，你們所說的，你們心裡所想的，我都知道。 \end{tabularx} \\ \\ \relax
11:6 & \begin{tabularx}{0.7\textwidth}{X} 你們在這城裡大行屠殺，被殺的人遍滿街道。 \end{tabularx} \\ \\ \relax
11:7 & \begin{tabularx}{0.7\textwidth}{X} 所以主耶和華如此說：你們在城中殺的人是肉，這城是鍋；你們卻要從其中被帶出去。 \end{tabularx} \\ \\ \relax
11:8 & \begin{tabularx}{0.7\textwidth}{X} 你們怕刀劍，我卻要使刀劍臨到你們。這是主耶和華說的。 \end{tabularx} \\ \\ \relax
11:9 & \begin{tabularx}{0.7\textwidth}{X} 我要把你們從這城中帶出去，交在外邦人的手裡，且要在你們中間施行審判。 \end{tabularx} \\ \\ \relax
11:10 & \begin{tabularx}{0.7\textwidth}{X} 你們要仆倒在刀下；我必在以色列的邊界審判你們，你們就知道我是耶和華。 \end{tabularx} \\ \\ \relax
11:11 & \begin{tabularx}{0.7\textwidth}{X} 這城必不作你們的鍋，你們也不作鍋中的肉。我要在以色列的邊界審判你們， \end{tabularx} \\ \\ \relax
11:12 & \begin{tabularx}{0.7\textwidth}{X} 你們就知道我是耶和華；因為你們不遵行我的律例，也不順從我的典章，卻隨從你們四圍列國的規條。」 \end{tabularx} \\ \\ \relax
11:13 & \begin{tabularx}{0.7\textwidth}{X} 我正說預言的時候，比拿雅的兒子毗拉提死了。於是我臉伏在地，大聲呼叫說：「唉！主耶和華啊，你要把以色列剩餘的人都滅絕淨盡嗎？」 \end{tabularx} \\ \\ \relax
11:14 & \begin{tabularx}{0.7\textwidth}{X} 耶和華的話臨到我，說： \end{tabularx} \\ \\ \relax
11:15 & \begin{tabularx}{0.7\textwidth}{X} 「人子啊，耶路撒冷的居民對你的兄弟、你的本家、你的親屬、以色列全家所有的人說：『你們遠離耶和華吧！這地是賜給我們為業的。』 \end{tabularx} \\ \\ \relax
11:16 & \begin{tabularx}{0.7\textwidth}{X} 所以你當說：『主耶和華如此說：我雖將以色列全家遠遠流放到列國，使他們分散在列邦，我卻要在他們所到的列邦，暫時作他們的聖所。』 \end{tabularx} \\ \\ \relax
11:17 & \begin{tabularx}{0.7\textwidth}{X} 你當說：『主耶和華如此說：我必從萬民中召集你們，從分散的列邦中聚集你們，又將以色列地賜給你們。』 \end{tabularx} \\ \\ \relax
11:18 & \begin{tabularx}{0.7\textwidth}{X} 他們到了那裡，必從其中除掉一切可憎之物、可厭的事。 \end{tabularx} \\ \\ \relax
11:19 & \begin{tabularx}{0.7\textwidth}{X} 我要使他們有合一的心，也要將新靈放在你們裡面，又從他們的肉體中除掉石心，賜給他們肉心， \end{tabularx} \\ \\ \relax
11:20 & \begin{tabularx}{0.7\textwidth}{X} 使他們順從我的律例，謹守遵行我的典章。他們要作我的子民，我要作他們的神。 \end{tabularx} \\ \\ \relax
11:21 & \begin{tabularx}{0.7\textwidth}{X} 至於那些心中隨從可憎之物、可厭的事的人，我必照他們所做的報應在他們頭上。這是主耶和華說的。」 \end{tabularx} \\ \\ \relax
11:22 & \begin{tabularx}{0.7\textwidth}{X} 於是，基路伯展開翅膀，輪子都在他們旁邊；在他們上面有以色列神的榮耀。 \end{tabularx} \\ \\ \relax
11:23 & \begin{tabularx}{0.7\textwidth}{X} 耶和華的榮耀從城中上升，停在城東的那座山上。 \end{tabularx} \\ \\ \relax
11:24 & \begin{tabularx}{0.7\textwidth}{X} 靈將我舉起，在異象中神的靈將我帶回迦勒底地，到被擄的人那裡；之後我所見的異象就離我上升去了。 \end{tabularx} \\ \\ \relax
11:25 & \begin{tabularx}{0.7\textwidth}{X} 我就把耶和華指示我的一切事都說給被擄的人聽。 \end{tabularx} \\ \\ \relax
12:1 & \begin{tabularx}{0.7\textwidth}{X} 耶和華的話臨到我，說： \end{tabularx} \\ \\ \relax
12:2 & \begin{tabularx}{0.7\textwidth}{X} 「人子啊，你住在悖逆之家中；他們有眼可看卻看不見，有耳可聽卻聽不到，因為他們是悖逆之家。 \end{tabularx} \\ \\ \relax
12:3 & \begin{tabularx}{0.7\textwidth}{X} 所以人子啊，你要收拾被擄時需用的物件，白天在他們眼前離去，在他們眼前離開你所住的地方，移到別處去；他們雖是悖逆之家，或者可以領悟。 \end{tabularx} \\ \\ \relax
12:4 & \begin{tabularx}{0.7\textwidth}{X} 你要白天在他們眼前拿出你被擄時需用的物件。到了晚上，要在他們眼前離去，像被擄的人離去一樣。 \end{tabularx} \\ \\ \relax
12:5 & \begin{tabularx}{0.7\textwidth}{X} 你要在他們眼前挖通牆壁，從其中將物件帶出去。 \end{tabularx} \\ \\ \relax
12:6 & \begin{tabularx}{0.7\textwidth}{X} 到天黑時，在他們眼前背在肩上帶走，並要蒙住臉看不見地，因為我要使你成為以色列家的預兆。」 \end{tabularx} \\ \\ \relax
12:7 & \begin{tabularx}{0.7\textwidth}{X} 我就照著所吩咐的去做，白天拿出被擄時需用的物件。到了晚上，用手挖通牆壁；天黑的時候，在他們眼前背在肩上帶走。 \end{tabularx} \\ \\ \relax
12:8 & \begin{tabularx}{0.7\textwidth}{X} 次日早晨，耶和華的話臨到我，說： \end{tabularx} \\ \\ \relax
12:9 & \begin{tabularx}{0.7\textwidth}{X} 「人子啊，以色列家，就是那悖逆之家，豈不是問你說：『你在做甚麼呢？』 \end{tabularx} \\ \\ \relax
12:10 & \begin{tabularx}{0.7\textwidth}{X} 你要對他們說：『主耶和華如此說：這是關乎耶路撒冷君王和其中以色列全家的默示。』 \end{tabularx} \\ \\ \relax
12:11 & \begin{tabularx}{0.7\textwidth}{X} 你要說：『我是你們的預兆：我怎樣做，他們所遭遇的也必這樣，他們必被擄去，作俘虜。』 \end{tabularx} \\ \\ \relax
12:12 & \begin{tabularx}{0.7\textwidth}{X} 他們中間的君王也必在天黑時把物件背在肩上帶走。他們要挖通牆壁，從其中帶出去。他必蒙住臉，眼看不見地。 \end{tabularx} \\ \\ \relax
12:13 & \begin{tabularx}{0.7\textwidth}{X} 我要把我的網撒在他身上，他就被我的羅網纏住。我要帶他到迦勒底人之地的巴比倫；他沒有看見那地，就死在那裡。 \end{tabularx} \\ \\ \relax
12:14 & \begin{tabularx}{0.7\textwidth}{X} 我要把四圍幫助他的和他所有的軍隊分散到四方，也要拔刀追趕他們。 \end{tabularx} \\ \\ \relax
12:15 & \begin{tabularx}{0.7\textwidth}{X} 我把他們驅逐到列國，分散在列邦的時候，他們就知道我是耶和華。 \end{tabularx} \\ \\ \relax
12:16 & \begin{tabularx}{0.7\textwidth}{X} 我卻要留下他們當中幾個人得免刀劍、饑荒、瘟疫，使他們在所到的列國述說自己所做一切可憎的事；他們就知道我是耶和華。」 \end{tabularx} \\ \\ \relax
12:17 & \begin{tabularx}{0.7\textwidth}{X} 耶和華的話臨到我，說： \end{tabularx} \\ \\ \relax
12:18 & \begin{tabularx}{0.7\textwidth}{X} 「人子啊，你吃飯時必戰抖，喝水時必驚惶憂慮。 \end{tabularx} \\ \\ \relax
12:19 & \begin{tabularx}{0.7\textwidth}{X} 你要對這地的百姓說：主耶和華論以色列地的耶路撒冷居民如此說，他們吃飯時必憂慮，喝水時必驚惶，因其中居民所行殘暴的事，這地必然荒廢，一無所存。 \end{tabularx} \\ \\ \relax
12:20 & \begin{tabularx}{0.7\textwidth}{X} 有人居住的城鎮必變為廢墟，地必荒涼；你們就知道我是耶和華。」 \end{tabularx} \\ \\ \relax
12:21 & \begin{tabularx}{0.7\textwidth}{X} 耶和華的話臨到我，說： \end{tabularx} \\ \\ \relax
12:22 & \begin{tabularx}{0.7\textwidth}{X} 「人子啊，在以色列地你們怎麼有這俗語說：『日子延長，一切異象卻落了空』呢？ \end{tabularx} \\ \\ \relax
12:23 & \begin{tabularx}{0.7\textwidth}{X} 你要告訴他們說：『主耶和華如此說：我必令這俗語止息，以色列中不再有人引用這俗語。』你卻要對他們說：『日子臨近，一切的異象都必應驗。』 \end{tabularx} \\ \\ \relax
12:24 & \begin{tabularx}{0.7\textwidth}{X} 從此以後，以色列家不再有虛假的異象和奉承的占卜。 \end{tabularx} \\ \\ \relax
12:25 & \begin{tabularx}{0.7\textwidth}{X} 我---耶和華說話，所說的必定實現，不再耽延。你們這悖逆之家啊，你們在世的日子，我所說的話必定實現。這是主耶和華說的。」 \end{tabularx} \\ \\ \relax
12:26 & \begin{tabularx}{0.7\textwidth}{X} 耶和華的話臨到我，說： \end{tabularx} \\ \\ \relax
12:27 & \begin{tabularx}{0.7\textwidth}{X} 「人子，看哪，以色列家的人說：『他所見的異象是許多日子以後的事，所說的預言是指著遙遠的時候。』 \end{tabularx} \\ \\ \relax
12:28 & \begin{tabularx}{0.7\textwidth}{X} 所以你要對他們說：『主耶和華如此說：我的話不再有一句耽延，我所說的話必定實現。』這是主耶和華說的。」 \end{tabularx} \\ \\ \relax
13:1 & \begin{tabularx}{0.7\textwidth}{X} 耶和華的話臨到我，說： \end{tabularx} \\ \\ \relax
13:2 & \begin{tabularx}{0.7\textwidth}{X} 「人子啊，你要說預言，攻擊以色列中說預言的先知，對那些隨心說預言的人說：『你們當聽耶和華的話。』」 \end{tabularx} \\ \\ \relax
13:3 & \begin{tabularx}{0.7\textwidth}{X} 主耶和華如此說：「禍哉！那些愚頑的先知，隨從自己的心意，卻一無所見。 \end{tabularx} \\ \\ \relax
13:4 & \begin{tabularx}{0.7\textwidth}{X} 以色列啊，你的先知好像廢墟中的狐狸， \end{tabularx} \\ \\ \relax
13:5 & \begin{tabularx}{0.7\textwidth}{X} 沒有上去堵住缺口，也沒有為以色列家重修城牆，使它在耶和華的日子來臨時，可以在戰爭中站得住。 \end{tabularx} \\ \\ \relax
13:6 & \begin{tabularx}{0.7\textwidth}{X} 他們看見的是虛假，是謊詐的占卜，說是耶和華說的；其實耶和華並沒有差遣他們，他們卻指望那話必站立得住。 \end{tabularx} \\ \\ \relax
13:7 & \begin{tabularx}{0.7\textwidth}{X} 你們豈不是見了虛假的異象嗎？豈不是說了謊詐的占卜嗎？你們說，這是耶和華說的，其實我沒有說過。」 \end{tabularx} \\ \\ \relax
13:8 & \begin{tabularx}{0.7\textwidth}{X} 所以主耶和華如此說：「因你們說的是虛假，見的是謊詐，所以，看哪，我要敵對你們。這是主耶和華說的。 \end{tabularx} \\ \\ \relax
13:9 & \begin{tabularx}{0.7\textwidth}{X} 我的手必攻擊那見虛假異象、用謊詐占卜的先知，他們必不列在我百姓的會中，不錄在以色列家的名冊上，也不能進入以色列地；你們就知道我是主耶和華。 \end{tabularx} \\ \\ \relax
13:10 & \begin{tabularx}{0.7\textwidth}{X} 他們誘惑我的百姓，說：『平安！』其實沒有平安，就像有人築牆壁，看哪，他們倒去粉刷它。 \end{tabularx} \\ \\ \relax
13:11 & \begin{tabularx}{0.7\textwidth}{X} 所以你要對那些粉刷的人說：『牆要倒塌，暴雨漫過。你們大冰雹啊，要降下 ，狂風要吹裂這牆。』 \end{tabularx} \\ \\ \relax
13:12 & \begin{tabularx}{0.7\textwidth}{X} 看哪，這牆倒塌，人豈不是要問你們說：『你們所粉刷的在哪裡呢？』」 \end{tabularx} \\ \\ \relax
13:13 & \begin{tabularx}{0.7\textwidth}{X} 所以主耶和華如此說：「我要發怒，使狂風吹裂它，在怒中令暴雨漫過，又發怒降下大冰雹，毀壞它。 \end{tabularx} \\ \\ \relax
13:14 & \begin{tabularx}{0.7\textwidth}{X} 我要這樣拆毀你們那粉飾的牆，把它夷為平地，以致根基露出；牆一倒塌，你們也要在其中滅亡。你們就知道我是耶和華。 \end{tabularx} \\ \\ \relax
13:15 & \begin{tabularx}{0.7\textwidth}{X} 我要對牆和粉刷它的人發盡我的憤怒，我要對你們說：『牆沒有了！粉刷它的人也沒有了！』 \end{tabularx} \\ \\ \relax
13:16 & \begin{tabularx}{0.7\textwidth}{X} 這就是以色列的先知，他們指著耶路撒冷說預言，見到這城平安的異象，其實沒有平安。這是主耶和華說的。」 \end{tabularx} \\ \\ \relax
13:17 & \begin{tabularx}{0.7\textwidth}{X} 「你，人子啊，要面向你百姓中隨心說預言的婦女們，說預言攻擊她們， \end{tabularx} \\ \\ \relax
13:18 & \begin{tabularx}{0.7\textwidth}{X} 說，主耶和華如此說：『這些婦女有禍了！她們為眾人的手腕縫驅邪帶，替身材高矮不同的人做頭巾，為要獵取人的性命。難道你們要獵取我百姓的性命，使自己存活嗎？ \end{tabularx} \\ \\ \relax
13:19 & \begin{tabularx}{0.7\textwidth}{X} 你們為幾把大麥、幾塊餅，在我的百姓中褻瀆我，對那肯聽謊言的百姓說謊言，讓不該死的人死，讓不該活的人活。』」 \end{tabularx} \\ \\ \relax
13:20 & \begin{tabularx}{0.7\textwidth}{X} 所以主耶和華如此說：「看哪，我要對付你們那用以獵取人，如獵飛鳥般的驅邪帶。我要把驅邪帶從你們的手腕扯去，釋放那些如飛鳥被你們獵取的人。 \end{tabularx} \\ \\ \relax
13:21 & \begin{tabularx}{0.7\textwidth}{X} 我也必撕裂你們的頭巾，救我百姓脫離你們的手，使他們不再被獵取，落在你們手中；你們就知道我是耶和華。 \end{tabularx} \\ \\ \relax
13:22 & \begin{tabularx}{0.7\textwidth}{X} 我未曾使義人傷心，你們卻以謊話使他傷心，且又堅固惡人的手，不使他回轉離開惡道得以存活。 \end{tabularx} \\ \\ \relax
13:23 & \begin{tabularx}{0.7\textwidth}{X} 所以，你們必不再看見虛假的異象，也不再行占卜的事；我要救我的百姓脫離你們的手；你們就知道我是耶和華。」 \end{tabularx} \\ \\ \relax
14:1 & \begin{tabularx}{0.7\textwidth}{X} 有幾個以色列的長老到我這裡來，坐在我面前。 \end{tabularx} \\ \\ \relax
14:2 & \begin{tabularx}{0.7\textwidth}{X} 耶和華的話臨到我，說： \end{tabularx} \\ \\ \relax
14:3 & \begin{tabularx}{0.7\textwidth}{X} 「人子啊，這些人在心中設立偶像，把陷自己於罪的絆腳石放在面前，我真的能讓他們求問嗎？ \end{tabularx} \\ \\ \relax
14:4 & \begin{tabularx}{0.7\textwidth}{X} 所以你要告訴他們，對他們說：『主耶和華如此說：以色列家的人，凡在心中設立偶像，把陷自己於罪的絆腳石放在面前，卻來到先知那裡的，我---耶和華在他所求的事上，必因他拜許多偶像向他施行報應， \end{tabularx} \\ \\ \relax
14:5 & \begin{tabularx}{0.7\textwidth}{X} 為要奪回以色列家的心，他們全都拜偶像，與我疏遠了。』 \end{tabularx} \\ \\ \relax
14:6 & \begin{tabularx}{0.7\textwidth}{X} 「所以你要對以色列家說：『主耶和華如此說：回轉吧！回轉離開你們的偶像，轉臉離開一切可憎的事。』 \end{tabularx} \\ \\ \relax
14:7 & \begin{tabularx}{0.7\textwidth}{X} 因為以色列家的人，或在以色列中寄居的外人，凡與我隔絕，在心中設立偶像，把陷自己於罪的絆腳石放在面前，卻來到先知那裡，要為自己的事求問我的，我---耶和華必親自報應他。 \end{tabularx} \\ \\ \relax
14:8 & \begin{tabularx}{0.7\textwidth}{X} 我要向那人變臉，使他成為警戒和笑柄，並且我要把他從我民中剪除；你們就知道我是耶和華。 \end{tabularx} \\ \\ \relax
14:9 & \begin{tabularx}{0.7\textwidth}{X} 先知若被騙說了一句預言，是我---耶和華騙了那先知，我要伸手攻擊他，把他從我百姓以色列中除滅。 \end{tabularx} \\ \\ \relax
14:10 & \begin{tabularx}{0.7\textwidth}{X} 他們必擔當自己的罪孽。先知的罪孽和求問之人的罪孽都一樣， \end{tabularx} \\ \\ \relax
14:11 & \begin{tabularx}{0.7\textwidth}{X} 使以色列家不再走迷離開我，也不再因各樣的罪過玷污自己，卻要作我的子民，我也作他們的神。這是主耶和華說的。」 \end{tabularx} \\ \\ \relax
14:12 & \begin{tabularx}{0.7\textwidth}{X} 耶和華的話臨到我，說： \end{tabularx} \\ \\ \relax
14:13 & \begin{tabularx}{0.7\textwidth}{X} 「人子啊，若有一國犯罪干犯我，我也伸手攻擊它，斷絕他們糧食的供應，使饑荒臨到那地，將人與牲畜從其中剪除； \end{tabularx} \\ \\ \relax
14:14 & \begin{tabularx}{0.7\textwidth}{X} 雖有挪亞、但以理、約伯這三人在那裡，他們只能因自己的義救自己的命。這是主耶和華說的。 \end{tabularx} \\ \\ \relax
14:15 & \begin{tabularx}{0.7\textwidth}{X} 我若使惡獸經過那地，大肆蹂躪，使地荒涼，以致因這些獸，人都不得經過； \end{tabularx} \\ \\ \relax
14:16 & \begin{tabularx}{0.7\textwidth}{X} 雖有這三人在其中，主耶和華說：我指著我的永生起誓，他們不能救兒子女兒，只有他們自己可以得救，那地仍然荒涼。 \end{tabularx} \\ \\ \relax
14:17 & \begin{tabularx}{0.7\textwidth}{X} 或者我使刀劍臨到那地，說：『讓刀劍穿越那地』，以致我把人與牲畜從其中剪除； \end{tabularx} \\ \\ \relax
14:18 & \begin{tabularx}{0.7\textwidth}{X} 雖有這三人在其中，主耶和華說：我指著我的永生起誓，他們不能救兒子女兒，只有他們自己可以得救。 \end{tabularx} \\ \\ \relax
14:19 & \begin{tabularx}{0.7\textwidth}{X} 或者我叫瘟疫流行那地，把我的憤怒帶著血傾在其中，好使人與牲畜從其中剪除； \end{tabularx} \\ \\ \relax
14:20 & \begin{tabularx}{0.7\textwidth}{X} 雖有挪亞、但以理、約伯在那裡，主耶和華說：我指著我的永生起誓，他們不能救兒子女兒，只能因自己的義救自己的命。 \end{tabularx} \\ \\ \relax
14:21 & \begin{tabularx}{0.7\textwidth}{X} 「主耶和華如此說：我若將這四樣大災，就是刀劍、饑荒、惡獸、瘟疫降在耶路撒冷，將人與牲畜從其中剪除，豈不是更嚴重嗎？ \end{tabularx} \\ \\ \relax
14:22 & \begin{tabularx}{0.7\textwidth}{X} 看哪，在那裡必有倖免於難的人帶著兒子女兒；看哪，他們來到你們這裡；你們看見他們的所作所為，就會因我降給耶路撒冷的災禍，因我降給它的一切，得到安慰。 \end{tabularx} \\ \\ \relax
14:23 & \begin{tabularx}{0.7\textwidth}{X} 你們因看見他們的所作所為，得到安慰，就會知道我在耶路撒冷所做的並非毫無緣故。這是主耶和華說的。」 \end{tabularx} \\ \\ \relax
15:1 & \begin{tabularx}{0.7\textwidth}{X} 耶和華的話臨到我，說： \end{tabularx} \\ \\ \relax
15:2 & \begin{tabularx}{0.7\textwidth}{X} 「人子啊，葡萄樹比一切其他的樹，就是樹林裡眾樹木的樹枝，有甚麼長處呢？ \end{tabularx} \\ \\ \relax
15:3 & \begin{tabularx}{0.7\textwidth}{X} 可以從其中取木料來做工嗎？人可以拿來做釘子，掛東西在上面嗎？ \end{tabularx} \\ \\ \relax
15:4 & \begin{tabularx}{0.7\textwidth}{X} 看哪，它已經拋在火中當柴燒，火既燒了兩頭，中間也燒焦了，它還有甚麼用處呢？ \end{tabularx} \\ \\ \relax
15:5 & \begin{tabularx}{0.7\textwidth}{X} 看哪，它完整的時候尚且不能拿來做工，何況被火燒焦了，還能拿來做工嗎？ \end{tabularx} \\ \\ \relax
15:6 & \begin{tabularx}{0.7\textwidth}{X} 所以，主耶和華如此說：我怎樣使林中樹裡的葡萄樹在火中當柴燒，我也必照樣對待耶路撒冷的居民。 \end{tabularx} \\ \\ \relax
15:7 & \begin{tabularx}{0.7\textwidth}{X} 我必向他們變臉；他們雖從火中逃出來，火仍要燒滅他們。我向他們變臉的時候，你們就知道我是耶和華。 \end{tabularx} \\ \\ \relax
15:8 & \begin{tabularx}{0.7\textwidth}{X} 我必使這地荒涼，因為他們做了背叛的事。這是主耶和華說的。」 \end{tabularx} \\ \\ \relax
16:1 & \begin{tabularx}{0.7\textwidth}{X} 耶和華的話臨到我，說： \end{tabularx} \\ \\ \relax
16:2 & \begin{tabularx}{0.7\textwidth}{X} 「人子啊，你要使耶路撒冷知道它那些可憎的事。 \end{tabularx} \\ \\ \relax
16:3 & \begin{tabularx}{0.7\textwidth}{X} 你要說，主耶和華對耶路撒冷如此說：你的根源，你的出身，是在迦南地；你的父親是亞摩利人，母親是赫人。 \end{tabularx} \\ \\ \relax
16:4 & \begin{tabularx}{0.7\textwidth}{X} 論到你出世的景況，在你出生的日子沒有人為你斷臍帶，也沒有用水清洗，使你潔淨；沒有人撒鹽在你身上，也沒有人用布包你。 \end{tabularx} \\ \\ \relax
16:5 & \begin{tabularx}{0.7\textwidth}{X} 沒有人顧惜你，為你做一件這樣的事來可憐你。你卻被扔在田野上面，因你出生的日子就被厭惡。 \end{tabularx} \\ \\ \relax
16:6 & \begin{tabularx}{0.7\textwidth}{X} 「我從你旁邊經過，見你在血中打滾，就對你說：『你雖在血中，卻要活下去！』我又說：『你雖在血中，卻要活下去！』 \end{tabularx} \\ \\ \relax
16:7 & \begin{tabularx}{0.7\textwidth}{X} 我使你成長如田間所生長的；你就漸長，美而又美，兩乳成形，頭髮秀長，但你仍然赤身露體。 \end{tabularx} \\ \\ \relax
16:8 & \begin{tabularx}{0.7\textwidth}{X} 「我從你旁邊經過看見你，看哪，正是你渴慕愛情的時候，我就用我衣服的邊搭在你身上，遮蓋你的赤體；又向你起誓，與你立約，你就歸我。這是主耶和華說的。 \end{tabularx} \\ \\ \relax
16:9 & \begin{tabularx}{0.7\textwidth}{X} 那時我用水洗你，洗淨你身上的血，又用油抹你。 \end{tabularx} \\ \\ \relax
16:10 & \begin{tabularx}{0.7\textwidth}{X} 我使你身穿錦繡衣裳，腳穿海狗皮鞋，用細麻布裹著你，精緻衣料披在你身上。 \end{tabularx} \\ \\ \relax
16:11 & \begin{tabularx}{0.7\textwidth}{X} 我用首飾打扮你：我把手鐲戴在你手上，項鏈在你頸上， \end{tabularx} \\ \\ \relax
16:12 & \begin{tabularx}{0.7\textwidth}{X} 我也把環子戴在你鼻上，耳環在你耳上，華冠在你頭上。 \end{tabularx} \\ \\ \relax
16:13 & \begin{tabularx}{0.7\textwidth}{X} 這樣，你就有金銀的首飾，穿的是細麻衣和精緻衣料，以及錦繡衣裳；吃的是細麵、蜂蜜和油。你也極其美貌，配登王后之位。 \end{tabularx} \\ \\ \relax
16:14 & \begin{tabularx}{0.7\textwidth}{X} 你美貌的名聲傳到列國，因我加給你榮華，使你完美。這是主耶和華說的。 \end{tabularx} \\ \\ \relax
16:15 & \begin{tabularx}{0.7\textwidth}{X} 「只是你仗著自己美貌，又憑著你的名聲行淫。你向路人縱情淫亂，你的美貌就屬於他的了。 \end{tabularx} \\ \\ \relax
16:16 & \begin{tabularx}{0.7\textwidth}{X} 你拿你的衣服為自己做成彩色丘壇，在其上行淫。這樣的事本不該有，以後也不該發生。 \end{tabularx} \\ \\ \relax
16:17 & \begin{tabularx}{0.7\textwidth}{X} 你拿我所賜給你的那些美麗的金銀寶物，為自己製造男性的偶像，與它們行淫； \end{tabularx} \\ \\ \relax
16:18 & \begin{tabularx}{0.7\textwidth}{X} 你拿你的錦繡衣裳為它們披上，把我的膏油和香料擺在它們面前； \end{tabularx} \\ \\ \relax
16:19 & \begin{tabularx}{0.7\textwidth}{X} 你把我賜給你的食物，就是我賜給你享用的細麵、油和蜂蜜，都擺在它們面前作為馨香的供物。事情就是這樣。這是主耶和華說的。 \end{tabularx} \\ \\ \relax
16:20 & \begin{tabularx}{0.7\textwidth}{X} 你拿你為我所生的兒女獻給它們吞噬。你的淫亂豈是小事？ \end{tabularx} \\ \\ \relax
16:21 & \begin{tabularx}{0.7\textwidth}{X} 你竟把我的兒女殺了，使他們經火獻給它們！ \end{tabularx} \\ \\ \relax
16:22 & \begin{tabularx}{0.7\textwidth}{X} 你做這一切可憎和淫亂的事，並未追念你幼年的日子，那時你赤身露體，在血中打滾。」 \end{tabularx} \\ \\ \relax
16:23 & \begin{tabularx}{0.7\textwidth}{X} 「你有禍了！你有禍了！這是主耶和華說的。你做這一切惡事之後， \end{tabularx} \\ \\ \relax
16:24 & \begin{tabularx}{0.7\textwidth}{X} 又為自己建造土墩，在各廣場上築起高臺。 \end{tabularx} \\ \\ \relax
16:25 & \begin{tabularx}{0.7\textwidth}{X} 你在各個街頭建造高臺，使你的美貌變為可憎；又向所有過路的人招手，多行淫亂。 \end{tabularx} \\ \\ \relax
16:26 & \begin{tabularx}{0.7\textwidth}{X} 你也和你那放縱情慾的鄰邦埃及人行淫，增添你的淫亂，惹我發怒。 \end{tabularx} \\ \\ \relax
16:27 & \begin{tabularx}{0.7\textwidth}{X} 看哪，我伸手攻擊你，減少你的福分，卻將你交給恨惡你的非利士人，讓他們任意待你。他們為你的淫行也感到羞恥。 \end{tabularx} \\ \\ \relax
16:28 & \begin{tabularx}{0.7\textwidth}{X} 你尚且不滿意，又與亞述人行淫，但與他們行淫之後，仍不滿足； \end{tabularx} \\ \\ \relax
16:29 & \begin{tabularx}{0.7\textwidth}{X} 於是你與那稱為貿易之地的迦勒底多行淫亂，即使這樣，你仍不滿足。 \end{tabularx} \\ \\ \relax
16:30 & \begin{tabularx}{0.7\textwidth}{X} 「你的心何等脆弱！這是主耶和華說的。你做這一切事，都是不知羞恥的妓女所做的， \end{tabularx} \\ \\ \relax
16:31 & \begin{tabularx}{0.7\textwidth}{X} 在各個街頭建造土墩，在各廣場上築高臺；但你藐視行淫的賞金，又不像妓女。 \end{tabularx} \\ \\ \relax
16:32 & \begin{tabularx}{0.7\textwidth}{X} 你這行淫的妻子啊，竟然接外人，替代丈夫。 \end{tabularx} \\ \\ \relax
16:33 & \begin{tabularx}{0.7\textwidth}{X} 凡妓女都是得人贈禮，你反倒餽贈你所愛的人，倒貼他們，使他們從四圍來與你行淫。 \end{tabularx} \\ \\ \relax
16:34 & \begin{tabularx}{0.7\textwidth}{X} 你的淫行與其他婦女相反，不是人要求與你行淫；是你給人賞金，不是人給你賞金；你是相反的。」 \end{tabularx} \\ \\ \relax
16:35 & \begin{tabularx}{0.7\textwidth}{X} 「你這妓女啊，要聽耶和華的話。 \end{tabularx} \\ \\ \relax
16:36 & \begin{tabularx}{0.7\textwidth}{X} 主耶和華如此說：因你放縱情慾，露出下體，與你所愛的行淫，因你敬拜一切可憎的偶像，就像自己兒女的血獻給它們， \end{tabularx} \\ \\ \relax
16:37 & \begin{tabularx}{0.7\textwidth}{X} 所以，看哪，我要聚集所有與你交歡的情人，不論是你所愛的或你所恨的，聚集他們從四圍到你那裡來；我要在他們面前暴露你的下體，使他們看盡你的下體。 \end{tabularx} \\ \\ \relax
16:38 & \begin{tabularx}{0.7\textwidth}{X} 我也要審判你，如審判淫婦和流人血的婦女一樣。我要在憤怒和妒忌中使流血的罪歸到你身上。 \end{tabularx} \\ \\ \relax
16:39 & \begin{tabularx}{0.7\textwidth}{X} 我要把你交在他們手中；他們必拆毀你的土墩，毀壞你的高臺，剝去你的衣服，奪取你美麗的寶物，留下你赤身露體。 \end{tabularx} \\ \\ \relax
16:40 & \begin{tabularx}{0.7\textwidth}{X} 他們必聚集眾人攻擊你，用石頭打死你，用刀劍刺透你， \end{tabularx} \\ \\ \relax
16:41 & \begin{tabularx}{0.7\textwidth}{X} 用火焚燒你的房屋，在許多婦女眼前審判你。我必使你不再行淫，你也不再給賞金。 \end{tabularx} \\ \\ \relax
16:42 & \begin{tabularx}{0.7\textwidth}{X} 我止息了向你所發的憤怒，我的妒忌也離開了你；這樣，我就平靜，不再惱怒。 \end{tabularx} \\ \\ \relax
16:43 & \begin{tabularx}{0.7\textwidth}{X} 因你不追念幼年的日子，反而在這一切的事上惹我發烈怒，所以，看哪，我必照你所做的報應在你頭上。在你一切可憎的事上，你不是還行了淫亂嗎？這是主耶和華說的。」 \end{tabularx} \\ \\ \relax
16:44 & \begin{tabularx}{0.7\textwidth}{X} 「看哪，凡說俗語的必用這俗語攻擊你，說：『有其母必有其女。』 \end{tabularx} \\ \\ \relax
16:45 & \begin{tabularx}{0.7\textwidth}{X} 你實在是你母親的女兒，厭棄丈夫和兒女；你也是你姊妹的姊妹，厭棄丈夫和兒女。你的母親是赫人，父親是亞摩利人。 \end{tabularx} \\ \\ \relax
16:46 & \begin{tabularx}{0.7\textwidth}{X} 你的姊姊是撒瑪利亞，她和她的女兒們住在你北邊；你的妹妹是所多瑪，她和她的女兒們住在你南邊。 \end{tabularx} \\ \\ \relax
16:47 & \begin{tabularx}{0.7\textwidth}{X} 你不只效法她們的行為，照她們可憎的事去做，不消多時，你所做的一切就比她們更惡。 \end{tabularx} \\ \\ \relax
16:48 & \begin{tabularx}{0.7\textwidth}{X} 主耶和華說：我指著我的永生起誓，你的妹妹所多瑪與她的女兒們並未做你和你女兒們所做的事。 \end{tabularx} \\ \\ \relax
16:49 & \begin{tabularx}{0.7\textwidth}{X} 看哪，你的妹妹所多瑪的罪孽是這樣：她和她的女兒們都驕傲，糧源充足，大享安逸，卻不扶持困苦和貧窮人的手。 \end{tabularx} \\ \\ \relax
16:50 & \begin{tabularx}{0.7\textwidth}{X} 她們狂傲，在我面前做可憎的事，我看見了就把她們除掉。 \end{tabularx} \\ \\ \relax
16:51 & \begin{tabularx}{0.7\textwidth}{X} 撒瑪利亞所犯的罪不及你的一半，你所做可憎的事比她更多；比起你所做這一切可憎的事，你的姊妹倒顯為義。 \end{tabularx} \\ \\ \relax
16:52 & \begin{tabularx}{0.7\textwidth}{X} 你既為你的姊妹辯護，就要擔當自己的羞辱。因你所犯的罪比她們更可憎，她們比你倒顯為義；你既使你的姊妹顯為義，就要抱愧，擔當自己的羞辱。」 \end{tabularx} \\ \\ \relax
16:53 & \begin{tabularx}{0.7\textwidth}{X} 「我必使她們被擄的歸回，使所多瑪和她的女兒們、撒瑪利亞和她的女兒們，並與你一起被擄的都歸回； \end{tabularx} \\ \\ \relax
16:54 & \begin{tabularx}{0.7\textwidth}{X} 好使你擔當自己的羞辱，為所做的一切抱愧，讓她們得到安慰。 \end{tabularx} \\ \\ \relax
16:55 & \begin{tabularx}{0.7\textwidth}{X} 你的妹妹所多瑪和她的女兒們必回復原狀；撒瑪利亞和她的女兒們必回復原狀；你和你的女兒們也必回復原狀。 \end{tabularx} \\ \\ \relax
16:56 & \begin{tabularx}{0.7\textwidth}{X} 在你驕傲的日子，你的妹妹所多瑪豈不是你口中的笑柄嗎？ \end{tabularx} \\ \\ \relax
16:57 & \begin{tabularx}{0.7\textwidth}{X} 在你的惡行顯露以前，那受了凌辱的亞蘭女兒們和亞蘭四圍非利士的女兒們，都在四圍藐視你。 \end{tabularx} \\ \\ \relax
16:58 & \begin{tabularx}{0.7\textwidth}{X} 耶和華說：你的淫蕩和可憎之事，你自己要擔當。」 \end{tabularx} \\ \\ \relax
16:59 & \begin{tabularx}{0.7\textwidth}{X} 「主耶和華如此說：你這輕看誓言而背約的，我必照你所做的報應你。 \end{tabularx} \\ \\ \relax
16:60 & \begin{tabularx}{0.7\textwidth}{X} 然而我要追念在你幼年時我與你所立的約，也要與你立定永約。 \end{tabularx} \\ \\ \relax
16:61 & \begin{tabularx}{0.7\textwidth}{X} 當你接納你的姊姊和妹妹時，你要追念你所行的，自覺慚愧；並且我要將她們賞給你做女兒，卻不是按著我與你所立的約。 \end{tabularx} \\ \\ \relax
16:62 & \begin{tabularx}{0.7\textwidth}{X} 我要堅定與你所立的約，你就知道我是耶和華， \end{tabularx} \\ \\ \relax
16:63 & \begin{tabularx}{0.7\textwidth}{X} 使你在我赦免你一切惡行時，心中追念，自覺慚愧，又因羞辱就不再開口。這是主耶和華說的。」 \end{tabularx} \\ \\ \relax
17:1 & \begin{tabularx}{0.7\textwidth}{X} 耶和華的話臨到我，說： \end{tabularx} \\ \\ \relax
17:2 & \begin{tabularx}{0.7\textwidth}{X} 「人子啊，你要向以色列家出謎語，設比喻， \end{tabularx} \\ \\ \relax
17:3 & \begin{tabularx}{0.7\textwidth}{X} 說，主耶和華如此說：有一隻大鷹，翅膀大，翎毛長，羽毛豐滿，色彩繽紛；牠飛到黎巴嫩，啄去香柏樹梢， \end{tabularx} \\ \\ \relax
17:4 & \begin{tabularx}{0.7\textwidth}{X} 啄斷它頂端的嫩枝，叼到貿易之地，放在商業城中。 \end{tabularx} \\ \\ \relax
17:5 & \begin{tabularx}{0.7\textwidth}{X} 牠又從這地取了一些種子，種在肥沃的田裡，栽於豐沛的水源旁，如種植柳樹。 \end{tabularx} \\ \\ \relax
17:6 & \begin{tabularx}{0.7\textwidth}{X} 它漸漸生長，成為低矮蔓生的葡萄樹；樹枝伸向那鷹，根部在牠下面。這樣，它就長成了一棵葡萄樹，生出枝子，長出枝幹。 \end{tabularx} \\ \\ \relax
17:7 & \begin{tabularx}{0.7\textwidth}{X} 「有一隻大鷹，翅膀大，羽毛多。看哪，葡萄樹從栽種它的苗圃向這鷹伸出根來，長出枝子，期盼從牠得到澆灌。 \end{tabularx} \\ \\ \relax
17:8 & \begin{tabularx}{0.7\textwidth}{X} 這棵樹栽於肥田豐沛的水源旁，原是為了生枝、結果，成為佳美的葡萄樹。 \end{tabularx} \\ \\ \relax
17:9 & \begin{tabularx}{0.7\textwidth}{X} 你要說，主耶和華如此說：這棵葡萄樹豈能發旺呢？鷹豈不拔出它的根來，摘光它的果子，使它枯乾，連長出的嫩葉都枯萎了嗎？要把它連根拔除，並不需要費大力或動用許多人。 \end{tabularx} \\ \\ \relax
17:10 & \begin{tabularx}{0.7\textwidth}{X} 看哪，葡萄樹雖然栽種了，豈能發旺呢？一經東風擊打，豈不全然枯乾了嗎？它必在生長的苗圃中枯乾了。」 \end{tabularx} \\ \\ \relax
17:11 & \begin{tabularx}{0.7\textwidth}{X} 耶和華的話臨到我，說： \end{tabularx} \\ \\ \relax
17:12 & \begin{tabularx}{0.7\textwidth}{X} 「你要對那悖逆之家說：你們不知道這些事是甚麼意思嗎？你要這樣說，看哪，巴比倫王曾到耶路撒冷，把其中的君王和官長帶到巴比倫去， \end{tabularx} \\ \\ \relax
17:13 & \begin{tabularx}{0.7\textwidth}{X} 又從以色列王室後裔中選取一人，與他立約，令他發誓，又擄走國中有勢力的人， \end{tabularx} \\ \\ \relax
17:14 & \begin{tabularx}{0.7\textwidth}{X} 使王國衰弱，不再強盛，只能靠守盟約方得生存。 \end{tabularx} \\ \\ \relax
17:15 & \begin{tabularx}{0.7\textwidth}{X} 他卻背叛巴比倫王，差派使者前往埃及，要求埃及人給他馬匹和許多人。他豈能亨通呢？這樣做的人豈能逃脫呢？他背了約豈能逃脫呢？ \end{tabularx} \\ \\ \relax
17:16 & \begin{tabularx}{0.7\textwidth}{X} 主耶和華說：我指著我的永生起誓，他定要死在巴比倫，就是巴比倫王所在之處；因為巴比倫王立他為王，他竟輕看向王所起的誓，背棄王與他所立的約。 \end{tabularx} \\ \\ \relax
17:17 & \begin{tabularx}{0.7\textwidth}{X} 當敵人建土堆，築堡壘，要殲滅許多人時，法老雖有強大軍隊和大批人馬，在戰場上還是不能幫助他。 \end{tabularx} \\ \\ \relax
17:18 & \begin{tabularx}{0.7\textwidth}{X} 他輕看誓言，背棄盟約，看哪，雖已投降，卻又做這一切的事，他必不能逃脫。 \end{tabularx} \\ \\ \relax
17:19 & \begin{tabularx}{0.7\textwidth}{X} 所以主耶和華如此說：我指著我的永生起誓，他既輕看我的誓言，背棄我的約，我必使這罪歸到他頭上。 \end{tabularx} \\ \\ \relax
17:20 & \begin{tabularx}{0.7\textwidth}{X} 我要把我的網撒在他身上，他就被我的羅網纏住。我要帶他到巴比倫，在那裡因他背叛我的罪懲罰他。 \end{tabularx} \\ \\ \relax
17:21 & \begin{tabularx}{0.7\textwidth}{X} 所有逃跑的軍隊必倒在刀下；剩餘的也必分散四方。你們就知道說這話的是我---耶和華。」 \end{tabularx} \\ \\ \relax
17:22 & \begin{tabularx}{0.7\textwidth}{X} 主耶和華如此說：「我要從香柏樹高高的樹梢摘取並栽上，從頂端的嫩枝中折下一嫩枝，栽於極高的山上， \end{tabularx} \\ \\ \relax
17:23 & \begin{tabularx}{0.7\textwidth}{X} 栽在以色列高處的山上。它就生枝、結果，成為高大的香柏樹，各類飛禽中的鳥都來宿在其下，宿在枝子的蔭下。 \end{tabularx} \\ \\ \relax
17:24 & \begin{tabularx}{0.7\textwidth}{X} 田野的樹木因此就知道是我---耶和華使高樹矮小，使矮樹高大，使綠樹枯乾，使枯樹發旺。我---耶和華說了這話，就必成就。」 \end{tabularx} \\ \\ \relax
18:1 & \begin{tabularx}{0.7\textwidth}{X} 耶和華的話又臨到我，說： \end{tabularx} \\ \\ \relax
18:2 & \begin{tabularx}{0.7\textwidth}{X} 「你們在以色列地何以有這俗語，『父親吃了酸葡萄，兒子牙齒就酸倒』呢？ \end{tabularx} \\ \\ \relax
18:3 & \begin{tabularx}{0.7\textwidth}{X} 主耶和華說：我指著我的永生起誓，你們在以色列必不再引用這俗語。 \end{tabularx} \\ \\ \relax
18:4 & \begin{tabularx}{0.7\textwidth}{X} 看哪，所有的生命都是屬我的；父親的生命怎樣屬我，兒子的生命也照樣屬我；然而犯罪的，他必定死。 \end{tabularx} \\ \\ \relax
18:5 & \begin{tabularx}{0.7\textwidth}{X} 「人若是公義，行公平公義的事： \end{tabularx} \\ \\ \relax
18:6 & \begin{tabularx}{0.7\textwidth}{X} 未曾在山上吃祭物，未曾向以色列家的偶像舉目；未曾污辱鄰舍的妻，也未曾在婦人的經期間親近她； \end{tabularx} \\ \\ \relax
18:7 & \begin{tabularx}{0.7\textwidth}{X} 未曾虧負人，而是將欠債之人的抵押品還給他；未曾搶奪人的物件，卻把食物給飢餓的人吃，把衣服給赤身的人穿； \end{tabularx} \\ \\ \relax
18:8 & \begin{tabularx}{0.7\textwidth}{X} 未曾向人取利息，也未曾索取高利，反倒縮手不作惡，在人與人之間施行誠實的判斷； \end{tabularx} \\ \\ \relax
18:9 & \begin{tabularx}{0.7\textwidth}{X} 遵行我的律例，謹守我的典章，按誠實行事；這人是公義的，必要存活。這是主耶和華說的。 \end{tabularx} \\ \\ \relax
18:10 & \begin{tabularx}{0.7\textwidth}{X} 「他若生了兒子，兒子作強盜，流人的血，作父親的雖然未犯此過，兒子卻對弟兄行了以上所說的惡，在山上吃祭物，污辱鄰舍的妻； \end{tabularx} \\ \\ \relax
18:11 & \begin{tabularx}{0.7\textwidth}{X} 見上節 \end{tabularx} \\ \\ \relax
18:12 & \begin{tabularx}{0.7\textwidth}{X} 虧負困苦和貧窮的人，搶奪別人的物件，不歸還抵押品，卻向偶像舉目，做可憎的事； \end{tabularx} \\ \\ \relax
18:13 & \begin{tabularx}{0.7\textwidth}{X} 向人取利息，索取高利，這人豈能存活呢？他不能存活。他因做這一切可憎的事，必要死亡，他的血要歸到自己身上。 \end{tabularx} \\ \\ \relax
18:14 & \begin{tabularx}{0.7\textwidth}{X} 「看哪，他若生了兒子，兒子見父親所犯的一切罪，他見了，卻不照樣去做； \end{tabularx} \\ \\ \relax
18:15 & \begin{tabularx}{0.7\textwidth}{X} 他未曾在山上吃祭物，未曾向以色列家的偶像舉目，未曾污辱鄰舍的妻； \end{tabularx} \\ \\ \relax
18:16 & \begin{tabularx}{0.7\textwidth}{X} 也未曾虧負人，未曾取人的抵押品，未曾搶奪人的物件，卻把食物給飢餓的人吃，把衣服給赤身的人穿， \end{tabularx} \\ \\ \relax
18:17 & \begin{tabularx}{0.7\textwidth}{X} 縮手不害困苦人，未曾向人索取利息或高利；反倒順從我的典章，遵行我的律例；如此，他必不因父親的罪孽死亡，定要存活。 \end{tabularx} \\ \\ \relax
18:18 & \begin{tabularx}{0.7\textwidth}{X} 至於他父親，因為施行欺壓，搶奪弟兄，在百姓中行不善，看哪，他必因自己的罪孽死亡。 \end{tabularx} \\ \\ \relax
18:19 & \begin{tabularx}{0.7\textwidth}{X} 「你們還說：『兒子為甚麼不擔當父親的罪孽呢？』兒子若行公平公義的事，謹守遵行我一切的律例，他必要存活。 \end{tabularx} \\ \\ \relax
18:20 & \begin{tabularx}{0.7\textwidth}{X} 惟有犯罪的，卻必死亡。兒子不擔當父親的罪孽，父親也不擔當兒子的罪孽。義人的善果要歸自己，惡人的惡報也要歸自己。 \end{tabularx} \\ \\ \relax
18:21 & \begin{tabularx}{0.7\textwidth}{X} 「惡人若回轉離開所做的一切罪惡，謹守我的一切律例，行公平公義的事，他必要存活，不致死亡。 \end{tabularx} \\ \\ \relax
18:22 & \begin{tabularx}{0.7\textwidth}{X} 他所犯的一切罪過都不被記念；他因所行的義，必要存活。 \end{tabularx} \\ \\ \relax
18:23 & \begin{tabularx}{0.7\textwidth}{X} 惡人死亡，豈是我所喜悅的呢？我豈不是喜悅他回轉離開所行的道而存活嗎？這是主耶和華說的。 \end{tabularx} \\ \\ \relax
18:24 & \begin{tabularx}{0.7\textwidth}{X} 至於義人，他若轉離義行而作惡，照著惡人所做一切可憎的事去做，豈能存活呢？他所行的一切義都不被記念；反而因所行的惡、所犯的罪死亡。 \end{tabularx} \\ \\ \relax
18:25 & \begin{tabularx}{0.7\textwidth}{X} 「你們卻說：『主的道不公平！』以色列家啊，你們要聽，我的道不公平嗎？你們的道不是不公平嗎？ \end{tabularx} \\ \\ \relax
18:26 & \begin{tabularx}{0.7\textwidth}{X} 義人若轉離義行而作惡，他就因這些惡而死亡。他要死在他所作的惡中。 \end{tabularx} \\ \\ \relax
18:27 & \begin{tabularx}{0.7\textwidth}{X} 惡人若回轉離開所行的惡，行公平公義的事，他必救自己的命； \end{tabularx} \\ \\ \relax
18:28 & \begin{tabularx}{0.7\textwidth}{X} 因為他省察，回轉離開所犯的一切罪過，他必要存活，不致死亡。 \end{tabularx} \\ \\ \relax
18:29 & \begin{tabularx}{0.7\textwidth}{X} 以色列家還說：『主的道不公平！』以色列家啊，我的道不公平嗎？你們的道不是不公平嗎？ \end{tabularx} \\ \\ \relax
18:30 & \begin{tabularx}{0.7\textwidth}{X} 所以，以色列家啊，我必按你們各人所做的審判你們。當回轉，回轉離開你們一切的罪過，免得罪孽成為你們的絆腳石。這是主耶和華說的。 \end{tabularx} \\ \\ \relax
18:31 & \begin{tabularx}{0.7\textwidth}{X} 你們要把所犯的一切罪過盡行拋棄，為自己造一個新的心和新的靈。以色列家啊，你們為甚麼要死呢？ \end{tabularx} \\ \\ \relax
18:32 & \begin{tabularx}{0.7\textwidth}{X} 我不喜歡有任何人死亡，所以你們當回轉，要存活！這是主耶和華說的。」 \end{tabularx} \\ \\ \relax
19:1 & \begin{tabularx}{0.7\textwidth}{X} 你當為以色列的領袖們唱哀歌， \end{tabularx} \\ \\ \relax
19:2 & \begin{tabularx}{0.7\textwidth}{X} 說：你的母親在獅子中是怎樣的母獅呢？牠蹲伏在少壯獅子中，養育小獅子。 \end{tabularx} \\ \\ \relax
19:3 & \begin{tabularx}{0.7\textwidth}{X} 牠養大了其中一隻小獅子，成了少壯獅子，學會抓食，牠就吃人。 \end{tabularx} \\ \\ \relax
19:4 & \begin{tabularx}{0.7\textwidth}{X} 列國聽見了就把牠逮住在他們的坑裡，用鉤子拉牠到埃及地去。 \end{tabularx} \\ \\ \relax
19:5 & \begin{tabularx}{0.7\textwidth}{X} 母獅見自己等候，期望落空，就從小獅子中取一隻，養為少壯獅子； \end{tabularx} \\ \\ \relax
19:6 & \begin{tabularx}{0.7\textwidth}{X} 牠在眾獅子中徜徉，長大成為少壯獅子，學會抓食，牠就吃人。 \end{tabularx} \\ \\ \relax
19:7 & \begin{tabularx}{0.7\textwidth}{X} 牠拆毀他們的宮殿，使他們的城鎮變為廢墟；因牠咆哮的聲音，遍地和其中所充滿的都荒廢了。 \end{tabularx} \\ \\ \relax
19:8 & \begin{tabularx}{0.7\textwidth}{X} 於是四圍列國從各省前來攻擊牠，把網撒在牠身上，把牠逮住在他們的坑裡。 \end{tabularx} \\ \\ \relax
19:9 & \begin{tabularx}{0.7\textwidth}{X} 他們又用鉤子鉤住牠，把牠放入籠中，帶到巴比倫王那裡，把牠押進城堡，以色列山上就不再聽見牠的聲音。 \end{tabularx} \\ \\ \relax
19:10 & \begin{tabularx}{0.7\textwidth}{X} 你的母親如葡萄樹，在葡萄園中，栽於水邊，因為水多，就多結果子，多生枝子； \end{tabularx} \\ \\ \relax
19:11 & \begin{tabularx}{0.7\textwidth}{X} 它長出堅固的枝幹，可作統治者的權杖。這枝幹高舉在茂密的樹枝中，可見樹身高大，枝子繁多。 \end{tabularx} \\ \\ \relax
19:12 & \begin{tabularx}{0.7\textwidth}{X} 但在烈怒中它被拔出，摔在地上；東風吹乾其果子，那堅固的枝幹因折斷而枯乾，被火燒燬； \end{tabularx} \\ \\ \relax
19:13 & \begin{tabularx}{0.7\textwidth}{X} 如今這葡萄樹移植於曠野，在乾旱無水之地， \end{tabularx} \\ \\ \relax
19:14 & \begin{tabularx}{0.7\textwidth}{X} 火從枝幹中發出，燒滅它的枝條和它的果子，以致不再有堅固的枝幹，可作統治者的權杖。這是哀傷之歌，成為一首哀歌。 \end{tabularx} \\ \\ \relax
20:1 & \begin{tabularx}{0.7\textwidth}{X} 第七年五月初十，有以色列的幾個長老前來求問耶和華，坐在我面前。 \end{tabularx} \\ \\ \relax
20:2 & \begin{tabularx}{0.7\textwidth}{X} 耶和華的話臨到我，說： \end{tabularx} \\ \\ \relax
20:3 & \begin{tabularx}{0.7\textwidth}{X} 「人子啊，你要告訴以色列的長老，對他們說，主耶和華如此說：你們來是為求問我嗎？主耶和華說：我指著我的永生起誓，我必不讓你們求問。 \end{tabularx} \\ \\ \relax
20:4 & \begin{tabularx}{0.7\textwidth}{X} 人子啊，你要審問他們嗎？你要審問嗎？你當使他們知道他們祖先那些可憎的事； \end{tabularx} \\ \\ \relax
20:5 & \begin{tabularx}{0.7\textwidth}{X} 你要對他們說，主耶和華如此說：當日我揀選以色列，對雅各家的後裔起誓，在埃及地向他們顯現，起誓說：我是耶和華---你們的神； \end{tabularx} \\ \\ \relax
20:6 & \begin{tabularx}{0.7\textwidth}{X} 那日我向他們起誓，要領他們出埃及地，到我為他們所找到的流奶與蜜之地，就是全地中最美好之地。 \end{tabularx} \\ \\ \relax
20:7 & \begin{tabularx}{0.7\textwidth}{X} 我對他們說，你們各人要拋棄眼中所喜愛的可憎之物，不可用埃及的偶像玷污自己。我是耶和華---你們的神。 \end{tabularx} \\ \\ \relax
20:8 & \begin{tabularx}{0.7\textwidth}{X} 他們卻悖逆我，不肯聽從我，不拋棄他們眼中所喜愛的可憎之物，離棄埃及的偶像。「我就說，在埃及地，我要把我的憤怒傾倒在他們身上，向他們發盡我的怒氣。 \end{tabularx} \\ \\ \relax
20:9 & \begin{tabularx}{0.7\textwidth}{X} 我這麼做是為了我名的緣故，免得我的名在他們所居住之列國眼中被褻瀆；我曾在這些列國眼前向他們顯現，領他們出了埃及地。 \end{tabularx} \\ \\ \relax
20:10 & \begin{tabularx}{0.7\textwidth}{X} 我領他們出埃及地，帶他們到曠野。 \end{tabularx} \\ \\ \relax
20:11 & \begin{tabularx}{0.7\textwidth}{X} 我將我的律例賜給他們，將我的典章指示他們；人若遵行就必因此存活。 \end{tabularx} \\ \\ \relax
20:12 & \begin{tabularx}{0.7\textwidth}{X} 我將我的安息日賜給他們，在我與他們中間作記號，讓他們知道我---耶和華是使他們分別為聖的。 \end{tabularx} \\ \\ \relax
20:13 & \begin{tabularx}{0.7\textwidth}{X} 以色列家卻在曠野中悖逆我，不順從我的律例，厭棄我的典章；人若遵行就必因此存活。他們卻大大干犯我的安息日。「因此我說，我要在曠野把我的憤怒傾倒在他們身上，滅絕他們。 \end{tabularx} \\ \\ \relax
20:14 & \begin{tabularx}{0.7\textwidth}{X} 我這麼做是為了我名的緣故，免得我的名在列國眼中被褻瀆，因為在這些列國眼前我領了他們出來。 \end{tabularx} \\ \\ \relax
20:15 & \begin{tabularx}{0.7\textwidth}{X} 並且我在曠野向他們起誓，必不領他們進入我所賜的流奶與蜜之地，就是全地中最美好之地； \end{tabularx} \\ \\ \relax
20:16 & \begin{tabularx}{0.7\textwidth}{X} 因為他們厭棄我的典章，不順從我的律例，干犯我的安息日，他們的心隨從自己的偶像。 \end{tabularx} \\ \\ \relax
20:17 & \begin{tabularx}{0.7\textwidth}{X} 雖然如此，我的眼仍顧惜他們，不毀滅他們，不在曠野把他們滅絕淨盡。 \end{tabularx} \\ \\ \relax
20:18 & \begin{tabularx}{0.7\textwidth}{X} 「我在曠野對他們的兒女說：『不要遵行你們祖先的律例，不要謹守他們的規條，也不要用他們的偶像玷污自己。 \end{tabularx} \\ \\ \relax
20:19 & \begin{tabularx}{0.7\textwidth}{X} 我是耶和華---你們的神，你們要順從我的律例，謹守遵行我的典章， \end{tabularx} \\ \\ \relax
20:20 & \begin{tabularx}{0.7\textwidth}{X} 且以我的安息日為聖。這日必在我與你們中間作記號，使你們知道我是耶和華---你們的神。』 \end{tabularx} \\ \\ \relax
20:21 & \begin{tabularx}{0.7\textwidth}{X} 只是他們的兒女悖逆我，不順從我的律例，也不謹守遵行我的典章；人若遵行就必因此存活。他們卻干犯我的安息日。「因此我說，我要在曠野把我的憤怒傾倒在他們身上，向他們發盡我的怒氣。 \end{tabularx} \\ \\ \relax
20:22 & \begin{tabularx}{0.7\textwidth}{X} 但我卻縮手而未如此行；我這麼做是為了我名的緣故，免得我的名在列國眼中被褻瀆，因為在這些列國眼前我領了他們出來。 \end{tabularx} \\ \\ \relax
20:23 & \begin{tabularx}{0.7\textwidth}{X} 並且我在曠野向他們起誓，要把他們驅散到列國，分散在列邦； \end{tabularx} \\ \\ \relax
20:24 & \begin{tabularx}{0.7\textwidth}{X} 因為他們不遵行我的典章，厭棄我的律例，干犯我的安息日，眼目向著他們祖先的偶像。 \end{tabularx} \\ \\ \relax
20:25 & \begin{tabularx}{0.7\textwidth}{X} 我也任他們遵行那無益的律例，隨從那不能使人存活的規條。 \end{tabularx} \\ \\ \relax
20:26 & \begin{tabularx}{0.7\textwidth}{X} 他們使所有頭生的經火，我就任憑他們在這供物上玷污自己；我令他們驚恐，他們就知道我是耶和華。 \end{tabularx} \\ \\ \relax
20:27 & \begin{tabularx}{0.7\textwidth}{X} 「人子啊，你要告訴以色列家，對他們說，主耶和華如此說：你們的祖先在背叛我的事上再次褻瀆了我； \end{tabularx} \\ \\ \relax
20:28 & \begin{tabularx}{0.7\textwidth}{X} 我領他們到我起誓應許賜給他們的地，他們看見各高岡、各茂密的樹，就在那裡獻祭，獻上惹我發怒的供物，也在那裡焚燒馨香的祭，獻澆酒祭。 \end{tabularx} \\ \\ \relax
20:29 & \begin{tabularx}{0.7\textwidth}{X} 我就對他們說：你們去的那丘壇叫甚麼呢？它名叫巴麻，直到今日。 \end{tabularx} \\ \\ \relax
20:30 & \begin{tabularx}{0.7\textwidth}{X} 所以你要對以色列家說，主耶和華如此說：你們仍要照你們祖先所做的玷污自己嗎？還要照他們可憎的事行淫嗎？ \end{tabularx} \\ \\ \relax
20:31 & \begin{tabularx}{0.7\textwidth}{X} 當你們獻上供物，使你們兒子經火的時候，你們仍用各樣的偶像玷污自己，直到今日。以色列家啊，我豈能讓你們求問呢？主耶和華說：我指著我的永生起誓，我必不讓你們求問。 \end{tabularx} \\ \\ \relax
20:32 & \begin{tabularx}{0.7\textwidth}{X} 「你們說：『我們要像列國和列邦的宗族一樣，去事奉木頭與石頭。』你們所起的心意萬不能成就。」 \end{tabularx} \\ \\ \relax
20:33 & \begin{tabularx}{0.7\textwidth}{X} 「主耶和華說：我指著我的永生起誓，我要作王，用大能的手和伸出的膀臂，並傾倒出來的憤怒治理你們。 \end{tabularx} \\ \\ \relax
20:34 & \begin{tabularx}{0.7\textwidth}{X} 我必用大能的手和伸出的膀臂，並傾倒出來的憤怒，把你們從萬民中領出來，從被趕散到的列邦聚集你們。 \end{tabularx} \\ \\ \relax
20:35 & \begin{tabularx}{0.7\textwidth}{X} 我必帶你們到萬民的曠野，在那裡當面審判你們。 \end{tabularx} \\ \\ \relax
20:36 & \begin{tabularx}{0.7\textwidth}{X} 我怎樣在埃及地的曠野審判你們的祖先，也必照樣審判你們。這是主耶和華說的。 \end{tabularx} \\ \\ \relax
20:37 & \begin{tabularx}{0.7\textwidth}{X} 我要使你們從杖下經過，按著約的拘束帶領你們。 \end{tabularx} \\ \\ \relax
20:38 & \begin{tabularx}{0.7\textwidth}{X} 我必從你們中間除盡叛逆和得罪我的人；我將他們從所寄居的地方領出來，他們卻不得進入以色列地，你們就知道我是耶和華。 \end{tabularx} \\ \\ \relax
20:39 & \begin{tabularx}{0.7\textwidth}{X} 「你們，以色列家啊，主耶和華如此說：你們若不聽從我，從今以後就讓各人去事奉他的偶像吧，只是不可再以你們的供物和偶像褻瀆我的聖名。 \end{tabularx} \\ \\ \relax
20:40 & \begin{tabularx}{0.7\textwidth}{X} 「在我的聖山，就是以色列高處的山，以色列全家，那地所有的人，都要在那裡事奉我。在那裡我悅納他們，並要你們獻供物和初熟的土產，以及一切的聖物。這是主耶和華說的。 \end{tabularx} \\ \\ \relax
20:41 & \begin{tabularx}{0.7\textwidth}{X} 我把你們從萬民中領出來，從被趕散到的列邦聚集你們，那時我必悅納你們如同悅納馨香之祭，我要在列國眼前，在你們中間顯為聖。 \end{tabularx} \\ \\ \relax
20:42 & \begin{tabularx}{0.7\textwidth}{X} 我領你們進入以色列地，就是我起誓應許賜給你們列祖之地，那時你們就知道我是耶和華。 \end{tabularx} \\ \\ \relax
20:43 & \begin{tabularx}{0.7\textwidth}{X} 你們在那裡要追念那玷污自己的所作所為，又要因所行的一切惡事厭惡自己。 \end{tabularx} \\ \\ \relax
20:44 & \begin{tabularx}{0.7\textwidth}{X} 以色列家啊，我為我名的緣故，沒有照著你們的惡行和你們的敗壞對待你們；你們就知道我是耶和華。這是主耶和華說的。」 \end{tabularx} \\ \\ \relax
20:45 & \begin{tabularx}{0.7\textwidth}{X} 耶和華的話臨到我，說： \end{tabularx} \\ \\ \relax
20:46 & \begin{tabularx}{0.7\textwidth}{X} 「人子啊，你要面向南方，向南方傳講，向尼革夫田野的樹林說預言。 \end{tabularx} \\ \\ \relax
20:47 & \begin{tabularx}{0.7\textwidth}{X} 你要對尼革夫的樹林說，要聽耶和華的話。主耶和華如此說：看哪，我要在你那裡點火，燒滅你們中間所有的綠樹和枯樹，猛烈的火焰必不熄滅；從南到北，人的臉都被燒焦。 \end{tabularx} \\ \\ \relax
20:48 & \begin{tabularx}{0.7\textwidth}{X} 凡血肉之軀都知道是我---耶和華點了火，這火必不熄滅。」 \end{tabularx} \\ \\ \relax
20:49 & \begin{tabularx}{0.7\textwidth}{X} 於是我說：「唉！主耶和華啊，人都指著我說：他不是說比喻的人嗎？」 \end{tabularx} \\ \\ \relax
21:1 & \begin{tabularx}{0.7\textwidth}{X} 耶和華的話臨到我，說： \end{tabularx} \\ \\ \relax
21:2 & \begin{tabularx}{0.7\textwidth}{X} 「人子啊，把你的臉正對著耶路撒冷，對著聖所傳講 ，向以色列地說預言。 \end{tabularx} \\ \\ \relax
21:3 & \begin{tabularx}{0.7\textwidth}{X} 你要向以色列地說，耶和華如此說：看哪，我與你為敵，拔刀出鞘，把義人和惡人從你中間剪除。 \end{tabularx} \\ \\ \relax
21:4 & \begin{tabularx}{0.7\textwidth}{X} 因為我要剪除你當中的義人和惡人，所以我的刀要出鞘，從南到北攻擊所有的血肉之軀； \end{tabularx} \\ \\ \relax
21:5 & \begin{tabularx}{0.7\textwidth}{X} 凡血肉之軀都知道我---耶和華已拔刀出鞘，刀必不再入鞘。 \end{tabularx} \\ \\ \relax
21:6 & \begin{tabularx}{0.7\textwidth}{X} 你，人子啊，要嘆息，在他們眼前斷了腰，愁苦地嘆息。 \end{tabularx} \\ \\ \relax
21:7 & \begin{tabularx}{0.7\textwidth}{X} 若有人對你說：『你為甚麼嘆息呢？』你就說：『因為有風聲傳來，人心惶惶，雙手發軟，精神衰敗，膝弱如水。看哪，它臨近了，一定會發生。』這是主耶和華說的。」 \end{tabularx} \\ \\ \relax
21:8 & \begin{tabularx}{0.7\textwidth}{X} 耶和華的話臨到我，說： \end{tabularx} \\ \\ \relax
21:9 & \begin{tabularx}{0.7\textwidth}{X} 「人子啊，你要預言說，耶和華如此吩咐，你要說：有刀，刀已磨快，又擦亮了； \end{tabularx} \\ \\ \relax
21:10 & \begin{tabularx}{0.7\textwidth}{X} 磨快為要大大殺戮，擦亮為要像閃電。我們豈能快樂呢？它藐視我兒的權杖和一切的木頭。 \end{tabularx} \\ \\ \relax
21:11 & \begin{tabularx}{0.7\textwidth}{X} 它已經交給人擦亮，可以掌握使用；這刀已經磨快擦亮，好交在行殺戮的人手中。 \end{tabularx} \\ \\ \relax
21:12 & \begin{tabularx}{0.7\textwidth}{X} 人子啊，你要呼喊哀號，因為這刀將臨到我的百姓，臨到以色列所有的領袖身上。他們和我的百姓都要交在刀下，所以你要捶胸。 \end{tabularx} \\ \\ \relax
21:13 & \begin{tabularx}{0.7\textwidth}{X} 因為這是一個考驗，若它藐視權杖，也不算一回事，又怎麼樣呢？這是主耶和華說的。」 \end{tabularx} \\ \\ \relax
21:14 & \begin{tabularx}{0.7\textwidth}{X} 「人子啊，你要拍掌預言，使這刀三番兩次臨到；這是致人死傷的刀，就是包圍人，使人大受死傷的刀。 \end{tabularx} \\ \\ \relax
21:15 & \begin{tabularx}{0.7\textwidth}{X} 我設立這恐嚇的刀，攻擊他們一切的城門，為要使他們的心驚慌害怕，許多人因而跌倒。唉！它造得像閃電，磨得尖利，要行殺戮。 \end{tabularx} \\ \\ \relax
21:16 & \begin{tabularx}{0.7\textwidth}{X} 刀啊，要行動一致，向右邊，或指向左邊；面向哪方，就向哪方。 \end{tabularx} \\ \\ \relax
21:17 & \begin{tabularx}{0.7\textwidth}{X} 我也要拍掌，使我的憤怒平息。這是我---耶和華說的。」 \end{tabularx} \\ \\ \relax
21:18 & \begin{tabularx}{0.7\textwidth}{X} 耶和華的話臨到我，說： \end{tabularx} \\ \\ \relax
21:19 & \begin{tabularx}{0.7\textwidth}{X} 「人子啊，你要畫定兩條路線，使巴比倫王的刀過來，這兩條路必從同一地分出來；要在通往城裡的路口畫手作指標。 \end{tabularx} \\ \\ \relax
21:20 & \begin{tabularx}{0.7\textwidth}{X} 你要劃定一條路，使刀來到亞捫人的拉巴，來到猶大，在堅固城耶路撒冷。 \end{tabularx} \\ \\ \relax
21:21 & \begin{tabularx}{0.7\textwidth}{X} 因為巴比倫王站在岔路上，在兩條路口占卜。他搖籤求問神像，察看肝臟； \end{tabularx} \\ \\ \relax
21:22 & \begin{tabularx}{0.7\textwidth}{X} 右手是耶路撒冷的占卜，以便安設撞城槌，張口喊殺，揚聲呼叫，建土堆，築堡壘，以撞城槌攻打城門。 \end{tabularx} \\ \\ \relax
21:23 & \begin{tabularx}{0.7\textwidth}{X} 在那些曾鄭重起誓的猶大人眼中，這是虛假的占卜；但巴比倫王要使他們想起自己的罪孽，以便俘擄他們。」 \end{tabularx} \\ \\ \relax
21:24 & \begin{tabularx}{0.7\textwidth}{X} 於是，主耶和華如此說：「因你們的過犯顯露，你們的罪孽被記得，以致你們的罪惡在你們一切的行為上都彰顯出來；你們既被記得，就被擄在掌中。 \end{tabularx} \\ \\ \relax
21:25 & \begin{tabularx}{0.7\textwidth}{X} 你這褻瀆行惡的以色列王啊，你的日子，最後懲罰的時刻已來臨。 \end{tabularx} \\ \\ \relax
21:26 & \begin{tabularx}{0.7\textwidth}{X} 主耶和華如此說：當除掉榮冕，摘下華冠，景況已不復從前；要使卑者升為高，使高者降為卑。 \end{tabularx} \\ \\ \relax
21:27 & \begin{tabularx}{0.7\textwidth}{X} 我要將這國傾覆，傾覆，再傾覆；這國必不存在，直等到那應得的人來到，我就將國賜給他。」 \end{tabularx} \\ \\ \relax
21:28 & \begin{tabularx}{0.7\textwidth}{X} 「人子啊，你要說預言；你要說，論到亞捫人和他們的凌辱，主耶和華吩咐我如此說：有刀，拔出來的刀，已經擦亮，為了行殺戮；它亮如閃電以行吞滅。 \end{tabularx} \\ \\ \relax
21:29 & \begin{tabularx}{0.7\textwidth}{X} 他們為你見虛假的異象，行謊詐的占卜，使你倒在褻瀆之惡人的頸項上；他們的日子，最後懲罰的時刻已來臨。 \end{tabularx} \\ \\ \relax
21:30 & \begin{tabularx}{0.7\textwidth}{X} 你收刀入鞘吧！我要在你受造之處、生長之地懲罰你。 \end{tabularx} \\ \\ \relax
21:31 & \begin{tabularx}{0.7\textwidth}{X} 我要把我的憤怒傾倒在你身上，把我烈怒的火噴在你身上；又將你交在善於殺滅、畜牲一般的人手中。 \end{tabularx} \\ \\ \relax
21:32 & \begin{tabularx}{0.7\textwidth}{X} 你要成為火中之柴，你的血必在地裡；你必不再被記得，因為這是我---耶和華說的。」 \end{tabularx} \\ \\ \relax
22:1 & \begin{tabularx}{0.7\textwidth}{X} 耶和華的話臨到我，說： \end{tabularx} \\ \\ \relax
22:2 & \begin{tabularx}{0.7\textwidth}{X} 「你，人子啊，你要審問，審問這流人血的城嗎？要使它知道它一切可憎的事。 \end{tabularx} \\ \\ \relax
22:3 & \begin{tabularx}{0.7\textwidth}{X} 你要說，主耶和華如此說：那在其中流人血的城啊，它的時刻已到，它製造偶像玷污了自己。 \end{tabularx} \\ \\ \relax
22:4 & \begin{tabularx}{0.7\textwidth}{X} 你因流了人的血，算為有罪；因所製造的偶像，玷污自己；你使你的日子臨近，你的年數已來到。所以我使你承受列國的凌辱和列邦的譏誚。 \end{tabularx} \\ \\ \relax
22:5 & \begin{tabularx}{0.7\textwidth}{X} 你這惡名昭彰、混亂的城啊，離你或遠或近的國家都必譏誚你。 \end{tabularx} \\ \\ \relax
22:6 & \begin{tabularx}{0.7\textwidth}{X} 「看哪，以色列的領袖在你那裡，為了流人的血各逞其能。 \end{tabularx} \\ \\ \relax
22:7 & \begin{tabularx}{0.7\textwidth}{X} 你那裡有輕慢父母的，在你當中有欺壓寄居者的，你那裡也有虧負孤兒寡婦的。 \end{tabularx} \\ \\ \relax
22:8 & \begin{tabularx}{0.7\textwidth}{X} 你藐視我的聖物，干犯我的安息日。 \end{tabularx} \\ \\ \relax
22:9 & \begin{tabularx}{0.7\textwidth}{X} 你那裡有為流人血而毀謗人的，你那裡有在山上吃祭物的，在你當中也有行淫亂的， \end{tabularx} \\ \\ \relax
22:10 & \begin{tabularx}{0.7\textwidth}{X} 有露父親下體的，有玷辱經期中不潔淨之婦人的。 \end{tabularx} \\ \\ \relax
22:11 & \begin{tabularx}{0.7\textwidth}{X} 這人與鄰舍的妻子行可憎的事，那人行淫污辱媳婦，在你那裡還有人污辱他的姊妹，父親的女兒。 \end{tabularx} \\ \\ \relax
22:12 & \begin{tabularx}{0.7\textwidth}{X} 你那裡有收取報酬而流人血的。你取利息，又索取高利；欺壓鄰舍，奪取財物；你竟然忘了我。這是主耶和華說的。 \end{tabularx} \\ \\ \relax
22:13 & \begin{tabularx}{0.7\textwidth}{X} 「看哪，我因你所得不義之財和你們中間所流的血，就擊打手掌。 \end{tabularx} \\ \\ \relax
22:14 & \begin{tabularx}{0.7\textwidth}{X} 到了我對付你的日子，你的心豈能忍受呢？你的手還能有力嗎？我---耶和華說了這話，就必成就。 \end{tabularx} \\ \\ \relax
22:15 & \begin{tabularx}{0.7\textwidth}{X} 我要把你驅散到列國，分散在列邦。我也必除掉你們中間的污穢。 \end{tabularx} \\ \\ \relax
22:16 & \begin{tabularx}{0.7\textwidth}{X} 你在列國眼前因自己所做的被侮辱，你就知道我是耶和華。」 \end{tabularx} \\ \\ \relax
22:17 & \begin{tabularx}{0.7\textwidth}{X} 耶和華的話臨到我，說： \end{tabularx} \\ \\ \relax
22:18 & \begin{tabularx}{0.7\textwidth}{X} 「人子啊，我看以色列家為渣滓。他們是爐中的銅、錫、鐵、鉛，是煉銀的渣滓。 \end{tabularx} \\ \\ \relax
22:19 & \begin{tabularx}{0.7\textwidth}{X} 所以主耶和華如此說：因你們全都成為渣滓，所以，看哪，我必將你們聚集在耶路撒冷中。 \end{tabularx} \\ \\ \relax
22:20 & \begin{tabularx}{0.7\textwidth}{X} 人怎樣把銀、銅、鐵、鉛、錫聚在爐中，吹火使它鎔化；照樣，我也要在我的怒氣和憤怒中聚集你們，把你們安置在城中，使你們鎔化。 \end{tabularx} \\ \\ \relax
22:21 & \begin{tabularx}{0.7\textwidth}{X} 我必聚集你們，把我烈怒的火吹在你們身上，你們就在其中鎔化。 \end{tabularx} \\ \\ \relax
22:22 & \begin{tabularx}{0.7\textwidth}{X} 銀子怎樣在爐中鎔化，你們也必照樣在城中鎔化，因此就知道是我---耶和華把憤怒傾倒在你們身上。」 \end{tabularx} \\ \\ \relax
22:23 & \begin{tabularx}{0.7\textwidth}{X} 耶和華的話臨到我，說： \end{tabularx} \\ \\ \relax
22:24 & \begin{tabularx}{0.7\textwidth}{X} 「人子啊，你要向這地說：你是未被潔淨之地，在我盛怒的日子，沒有雨水在其上。 \end{tabularx} \\ \\ \relax
22:25 & \begin{tabularx}{0.7\textwidth}{X} 其中的先知同謀背叛，如咆哮的獅子抓撕掠物。他們吞滅人命，搶奪財寶，使這地寡婦增多。 \end{tabularx} \\ \\ \relax
22:26 & \begin{tabularx}{0.7\textwidth}{X} 其中的祭司曲解我的律法，褻瀆我的聖物，不分別聖與俗，也不使人分辨潔淨和不潔淨，又遮眼不顧我的安息日；在他們中間連我也被褻慢了。 \end{tabularx} \\ \\ \relax
22:27 & \begin{tabularx}{0.7\textwidth}{X} 其中的領袖彷彿野狼抓撕掠物，流人的血，傷害人命，為得不義之財。 \end{tabularx} \\ \\ \relax
22:28 & \begin{tabularx}{0.7\textwidth}{X} 其中的先知為他們粉刷，見虛假的異象，行謊詐的占卜，說：『主耶和華如此說』，其實耶和華並沒有說。 \end{tabularx} \\ \\ \relax
22:29 & \begin{tabularx}{0.7\textwidth}{X} 這地的百姓慣行欺壓搶奪之事，虧負困苦和貧窮的人，欺壓寄居者，沒有公平。 \end{tabularx} \\ \\ \relax
22:30 & \begin{tabularx}{0.7\textwidth}{X} 我在他們中間尋找一人重修城牆，在我面前為這地站在缺口上，使我不致滅絕它，卻連一個也找不著。 \end{tabularx} \\ \\ \relax
22:31 & \begin{tabularx}{0.7\textwidth}{X} 所以我把憤怒傾倒在他們身上，用烈怒之火消滅他們，照他們所做的報應在他們頭上。這是主耶和華說的。」 \end{tabularx} \\ \\ \relax
23:1 & \begin{tabularx}{0.7\textwidth}{X} 耶和華的話臨到我，說： \end{tabularx} \\ \\ \relax
23:2 & \begin{tabularx}{0.7\textwidth}{X} 「人子啊，有兩個女子，是一母所生， \end{tabularx} \\ \\ \relax
23:3 & \begin{tabularx}{0.7\textwidth}{X} 她們在埃及行淫，年少時就開始行淫；在那裡任人擁抱胸懷，撫弄她們少女的乳房。 \end{tabularx} \\ \\ \relax
23:4 & \begin{tabularx}{0.7\textwidth}{X} 她們的名字，大的叫阿荷拉，妹妹叫阿荷利巴。她們都歸於我，生了兒女。論到她們的名字，阿荷拉是撒瑪利亞，阿荷利巴是耶路撒冷。 \end{tabularx} \\ \\ \relax
23:5 & \begin{tabularx}{0.7\textwidth}{X} 「阿荷拉歸我之後卻仍行淫，戀慕所愛的人，就是亞述人，都是戰士， \end{tabularx} \\ \\ \relax
23:6 & \begin{tabularx}{0.7\textwidth}{X} 穿著藍衣，作省長、副省長，全都是俊美的年輕人，騎著馬的騎士。 \end{tabularx} \\ \\ \relax
23:7 & \begin{tabularx}{0.7\textwidth}{X} 阿荷拉與亞述人中所有的美男子放縱淫行，她因拜所戀慕之人的一切偶像，玷污了自己。 \end{tabularx} \\ \\ \relax
23:8 & \begin{tabularx}{0.7\textwidth}{X} 她從埃及的時候，就沒有離開過淫亂；因為她年輕時，有人與她同寢，撫弄她少女的乳房，和她縱慾行淫。 \end{tabularx} \\ \\ \relax
23:9 & \begin{tabularx}{0.7\textwidth}{X} 因此，我把她交在她所愛的人手中，就是她所戀慕的亞述人手中。 \end{tabularx} \\ \\ \relax
23:10 & \begin{tabularx}{0.7\textwidth}{X} 他們暴露她的下體，擄掠她的兒女，用刀殺了她；他們向她施行審判，使她在婦女中留下臭名。 \end{tabularx} \\ \\ \relax
23:11 & \begin{tabularx}{0.7\textwidth}{X} 「她妹妹阿荷利巴雖然看見了，卻還是縱慾，比姊姊更加腐敗，行淫亂比姊姊更甚。 \end{tabularx} \\ \\ \relax
23:12 & \begin{tabularx}{0.7\textwidth}{X} 她戀慕亞述人，就是省長和副省長，披掛整齊的戰士，騎著馬的騎士，全都是俊美的年輕人。 \end{tabularx} \\ \\ \relax
23:13 & \begin{tabularx}{0.7\textwidth}{X} 我看見她被污辱，姊妹二人同行一路。 \end{tabularx} \\ \\ \relax
23:14 & \begin{tabularx}{0.7\textwidth}{X} 阿荷利巴又加增淫行，她看見牆上刻有人像，就是鮮紅色的迦勒底人雕刻的像。 \end{tabularx} \\ \\ \relax
23:15 & \begin{tabularx}{0.7\textwidth}{X} 它們腰間繫著帶子，頭上有飄揚的裹頭巾，都是將軍的樣子，巴比倫人的形像；迦勒底是他們的出生地。 \end{tabularx} \\ \\ \relax
23:16 & \begin{tabularx}{0.7\textwidth}{X} 阿荷利巴一看見就戀慕他們，派遣使者往迦勒底他們那裡去。 \end{tabularx} \\ \\ \relax
23:17 & \begin{tabularx}{0.7\textwidth}{X} 巴比倫人來到她那裡，上了她愛情的床，與她行淫污辱她。她被污辱，隨後她的心卻與他們生疏。 \end{tabularx} \\ \\ \relax
23:18 & \begin{tabularx}{0.7\textwidth}{X} 這樣，她既暴露淫行，暴露下體；我的心就與她生疏，像先前與她的姊姊生疏一樣。 \end{tabularx} \\ \\ \relax
23:19 & \begin{tabularx}{0.7\textwidth}{X} 她仍繼續增添淫行，追念她年輕時在埃及地行淫的日子， \end{tabularx} \\ \\ \relax
23:20 & \begin{tabularx}{0.7\textwidth}{X} 戀慕情人的身壯精足，如驢似馬。 \end{tabularx} \\ \\ \relax
23:21 & \begin{tabularx}{0.7\textwidth}{X} 這樣，你就渴望年輕時的淫蕩；那時，埃及人因你年輕時的胸懷，撫弄你的乳房。」 \end{tabularx} \\ \\ \relax
23:22 & \begin{tabularx}{0.7\textwidth}{X} 阿荷利巴啊，主耶和華如此說：「看哪，我要激起先前你喜愛，而後生疏的人前來攻擊你。我必使他們前來，在你四圍攻擊你； \end{tabularx} \\ \\ \relax
23:23 & \begin{tabularx}{0.7\textwidth}{X} 有巴比倫人、迦勒底眾人、比割人、書亞人、哥亞人，還有亞述眾人與他們一起，都是俊美的年輕人。他們是省長、副省長、將軍、有名聲的，全都騎著馬。 \end{tabularx} \\ \\ \relax
23:24 & \begin{tabularx}{0.7\textwidth}{X} 他們用兵器、戰車、輜重車，率領大軍前來攻擊你。他們要拿大小盾牌，戴著頭盔，在你四圍擺陣攻擊你。我要把審判交給他們，他們必按著自己的規條審判你。 \end{tabularx} \\ \\ \relax
23:25 & \begin{tabularx}{0.7\textwidth}{X} 我要向你傾洩我的妒忌，使他們以憤怒對待你。他們必割去你的鼻子和耳朵，你剩餘的人必倒在刀下。他們必擄去你的兒女，你所剩餘的必被火焚燒。 \end{tabularx} \\ \\ \relax
23:26 & \begin{tabularx}{0.7\textwidth}{X} 他們必剝去你的衣服，奪取你美麗的寶物。 \end{tabularx} \\ \\ \relax
23:27 & \begin{tabularx}{0.7\textwidth}{X} 這樣，我必止息你的淫行和你從埃及地就開始犯的淫亂，使你不再仰望亞述，也不再追念埃及。 \end{tabularx} \\ \\ \relax
23:28 & \begin{tabularx}{0.7\textwidth}{X} 主耶和華如此說：看哪，我必把你交在你所恨惡的人手中，就是你心與他生疏的人手中。 \end{tabularx} \\ \\ \relax
23:29 & \begin{tabularx}{0.7\textwidth}{X} 他們要以恨惡對待你，奪取你勞碌得來的一切，留下你赤身露體。你淫亂的下體，連你的淫行和淫蕩，都必顯露。 \end{tabularx} \\ \\ \relax
23:30 & \begin{tabularx}{0.7\textwidth}{X} 人必向你行這些事；因為你隨從外邦人行淫，用他們的偶像玷污自己。 \end{tabularx} \\ \\ \relax
23:31 & \begin{tabularx}{0.7\textwidth}{X} 你走了你姊姊的路，所以我必把她的杯交在你的手中。」 \end{tabularx} \\ \\ \relax
23:32 & \begin{tabularx}{0.7\textwidth}{X} 主耶和華如此說：「你必喝你姊姊的杯，那杯又深又廣，盛得很多，使你遭受嗤笑譏刺。 \end{tabularx} \\ \\ \relax
23:33 & \begin{tabularx}{0.7\textwidth}{X} 你必酩酊大醉，滿有愁苦。你姊姊撒瑪利亞的杯，驚駭和淒涼的杯， \end{tabularx} \\ \\ \relax
23:34 & \begin{tabularx}{0.7\textwidth}{X} 你必喝它，並且喝乾。甚至咀嚼杯片，撕裂自己的胸脯；因為我曾說過。這是主耶和華說的。」 \end{tabularx} \\ \\ \relax
23:35 & \begin{tabularx}{0.7\textwidth}{X} 主耶和華如此說：「因你忘記我，將我丟在背後，所以你要擔當你的淫行和淫蕩。」 \end{tabularx} \\ \\ \relax
23:36 & \begin{tabularx}{0.7\textwidth}{X} 耶和華對我說：「人子啊，你要審問阿荷拉與阿荷利巴嗎？要指出她們所做可憎的事。 \end{tabularx} \\ \\ \relax
23:37 & \begin{tabularx}{0.7\textwidth}{X} 她們行姦淫，手中有血。她們與偶像行姦淫，使她們為我所生的兒女經火，給它們當食物。 \end{tabularx} \\ \\ \relax
23:38 & \begin{tabularx}{0.7\textwidth}{X} 此外，她們還向我這樣做：同一天又玷污我的聖所，干犯我的安息日。 \end{tabularx} \\ \\ \relax
23:39 & \begin{tabularx}{0.7\textwidth}{X} 她們殺了兒女獻給偶像，當天又進入我的聖所，褻瀆了它。看哪，這就是她們在我殿中所做的。 \end{tabularx} \\ \\ \relax
23:40 & \begin{tabularx}{0.7\textwidth}{X} 「況且你們兩姊妹派人從遠方召人來。使者到了他們那裡，看哪，他們就來了。為了他們，你們沐浴，畫眼影，佩戴首飾， \end{tabularx} \\ \\ \relax
23:41 & \begin{tabularx}{0.7\textwidth}{X} 坐在華美的床上，前面擺設桌子，把我的香料和膏油放在其上。 \end{tabularx} \\ \\ \relax
23:42 & \begin{tabularx}{0.7\textwidth}{X} 在那裡有一群人歡樂的聲音；有許多的平民，從曠野來的醉漢，把鐲子戴在她們手上，把華冠戴在她們頭上。 \end{tabularx} \\ \\ \relax
23:43 & \begin{tabularx}{0.7\textwidth}{X} 「我論到這久行姦淫而色衰的婦人說：現在人們還要與她行淫，她也要與人行淫。 \end{tabularx} \\ \\ \relax
23:44 & \begin{tabularx}{0.7\textwidth}{X} 人去到阿荷拉和阿荷利巴二淫婦那裡，好像與妓女行淫。 \end{tabularx} \\ \\ \relax
23:45 & \begin{tabularx}{0.7\textwidth}{X} 義人必按照審判淫婦和流人血之婦人的規條，審判她們；因為她們是淫婦，她們的手中有血。」 \end{tabularx} \\ \\ \relax
23:46 & \begin{tabularx}{0.7\textwidth}{X} 主耶和華如此說：「我要讓軍隊上來攻擊她們，使她們驚駭，成為擄物。 \end{tabularx} \\ \\ \relax
23:47 & \begin{tabularx}{0.7\textwidth}{X} 這軍隊必用石頭打死她們，用刀劍殺害她們，又殺戮她們的兒女，用火焚燒她們的房屋。 \end{tabularx} \\ \\ \relax
23:48 & \begin{tabularx}{0.7\textwidth}{X} 我必使這地不再有淫行，所有的婦女都受警戒，不再效法你們的淫行。 \end{tabularx} \\ \\ \relax
23:49 & \begin{tabularx}{0.7\textwidth}{X} 人必因你們的淫行報應你們；你們要擔當拜偶像的罪，因此你們就知道我是主耶和華。」 \end{tabularx} \\ \\ \relax
24:1 & \begin{tabularx}{0.7\textwidth}{X} 第九年十月初十，耶和華的話臨到我，說： \end{tabularx} \\ \\ \relax
24:2 & \begin{tabularx}{0.7\textwidth}{X} 「人子啊，你要記錄這一天的名稱，這特別的一天，巴比倫王圍困耶路撒冷，就在這特別的一天。 \end{tabularx} \\ \\ \relax
24:3 & \begin{tabularx}{0.7\textwidth}{X} 你要向這悖逆之家設比喻，對他們說，主耶和華如此說：把鍋放在火上，放好了，倒水在其中； \end{tabularx} \\ \\ \relax
24:4 & \begin{tabularx}{0.7\textwidth}{X} 要將肉塊，一切肥美的肉塊，腿和肩都放在鍋裡，要裝滿上等的骨頭； \end{tabularx} \\ \\ \relax
24:5 & \begin{tabularx}{0.7\textwidth}{X} 要取羊群中最好的，把柴堆在下面，把它煮開，骨頭煮在其中。 \end{tabularx} \\ \\ \relax
24:6 & \begin{tabularx}{0.7\textwidth}{X} 「主耶和華如此說：禍哉！這流人血的城，就是長鏽的鍋。它的鏽未曾除掉，要將肉塊從其中一一取出，不必抽籤。 \end{tabularx} \\ \\ \relax
24:7 & \begin{tabularx}{0.7\textwidth}{X} 這城所流的血還在城中，血倒在光滑的磐石上，沒有倒在地上，用土掩蓋； \end{tabularx} \\ \\ \relax
24:8 & \begin{tabularx}{0.7\textwidth}{X} 是我使這城所流的血倒在光滑的磐石上，不得掩蓋，為要惹動憤怒，施行報應。 \end{tabularx} \\ \\ \relax
24:9 & \begin{tabularx}{0.7\textwidth}{X} 所以主耶和華如此說：禍哉！這流人血的城，我必親自加大柴堆。 \end{tabularx} \\ \\ \relax
24:10 & \begin{tabularx}{0.7\textwidth}{X} 你要添上木柴，使火著旺，將肉煮爛，加上香料，烤焦骨頭； \end{tabularx} \\ \\ \relax
24:11 & \begin{tabularx}{0.7\textwidth}{X} 你要把空鍋放在炭火上，將鍋燒熱，把銅燒紅，鎔化其中的污穢，除淨其上的鏽。 \end{tabularx} \\ \\ \relax
24:12 & \begin{tabularx}{0.7\textwidth}{X} 然而這一切勞碌無效，它厚厚的鏽，即使用火也除不掉。 \end{tabularx} \\ \\ \relax
24:13 & \begin{tabularx}{0.7\textwidth}{X} 雖然我想潔淨你污穢的淫行，你卻不潔淨，你的污穢再也不能潔淨，直等我止息了向你發的憤怒。 \end{tabularx} \\ \\ \relax
24:14 & \begin{tabularx}{0.7\textwidth}{X} 我---耶和華說了這話，時候到了，就必成就；必不退縮，不顧惜，也不憐憫。人必照你的所作所為審判你。這是主耶和華說的。」 \end{tabularx} \\ \\ \relax
24:15 & \begin{tabularx}{0.7\textwidth}{X} 耶和華的話臨到我，說： \end{tabularx} \\ \\ \relax
24:16 & \begin{tabularx}{0.7\textwidth}{X} 「人子，看哪，我要以災病奪取你眼中所喜愛的，你卻不可悲哀哭泣，也不可流淚， \end{tabularx} \\ \\ \relax
24:17 & \begin{tabularx}{0.7\textwidth}{X} 只可嘆息，不可出聲，不可辦理喪事；裹上頭巾，腳上穿鞋，不可摀著鬍鬚，也不可吃一般人的食物。」 \end{tabularx} \\ \\ \relax
24:18 & \begin{tabularx}{0.7\textwidth}{X} 到了早晨我把這事告訴百姓，晚上我的妻子就死了。次日早晨我就遵命而行。 \end{tabularx} \\ \\ \relax
24:19 & \begin{tabularx}{0.7\textwidth}{X} 百姓對我說：「你這樣做跟我們有甚麼關係，你不告訴我們嗎？」 \end{tabularx} \\ \\ \relax
24:20 & \begin{tabularx}{0.7\textwidth}{X} 我對他們說：「耶和華的話臨到我，說： \end{tabularx} \\ \\ \relax
24:21 & \begin{tabularx}{0.7\textwidth}{X} 『你告訴以色列家，主耶和華如此說：我要使我的聖所被褻瀆，就是你們憑勢力所誇耀、眼裡所喜愛、心中所愛惜的；並且你們所遺留的兒女必倒在刀下。 \end{tabularx} \\ \\ \relax
24:22 & \begin{tabularx}{0.7\textwidth}{X} 那時，你們要照我所做的去做。你們不可摀著鬍鬚，也不可吃一般人的食物。 \end{tabularx} \\ \\ \relax
24:23 & \begin{tabularx}{0.7\textwidth}{X} 你們頭要裹上頭巾，腳要穿上鞋；不可悲哀哭泣。你們必因自己的罪孽衰殘，相對嘆息。 \end{tabularx} \\ \\ \relax
24:24 & \begin{tabularx}{0.7\textwidth}{X} 以西結必這樣成為你們的預兆；凡他所做的，你們也必照樣做。那事來到，你們就知道我是主耶和華。』」 \end{tabularx} \\ \\ \relax
24:25 & \begin{tabularx}{0.7\textwidth}{X} 「你，人子啊，那日當我除掉他們所倚靠的保障、所歡喜的榮耀，並眼中所喜愛的，心裡所重看的兒女時， \end{tabularx} \\ \\ \relax
24:26 & \begin{tabularx}{0.7\textwidth}{X} 逃脫的人豈不來到你這裡，使你耳聞這事嗎？ \end{tabularx} \\ \\ \relax
24:27 & \begin{tabularx}{0.7\textwidth}{X} 那日你要向逃脫的人開口說話，不再啞口無言。你必這樣成為他們的預兆，他們就知道我是耶和華。」 \end{tabularx} \\ \\ \relax
25:1 & \begin{tabularx}{0.7\textwidth}{X} 耶和華的話臨到我，說： \end{tabularx} \\ \\ \relax
25:2 & \begin{tabularx}{0.7\textwidth}{X} 「人子啊，你要面向亞捫人說預言，攻擊他們。 \end{tabularx} \\ \\ \relax
25:3 & \begin{tabularx}{0.7\textwidth}{X} 你要對亞捫人說，當聽主耶和華的話。主耶和華如此說：我的聖所遭褻瀆，以色列地變荒涼，猶大家被擄掠；那時，你因這些事說『啊哈』， \end{tabularx} \\ \\ \relax
25:4 & \begin{tabularx}{0.7\textwidth}{X} 所以，看哪，我要把你交給東方人為業；他們必在你中間安營居住，設立居所，吃你的果子，喝你的奶。 \end{tabularx} \\ \\ \relax
25:5 & \begin{tabularx}{0.7\textwidth}{X} 我必使拉巴成為牧放駱駝之地，使亞捫成為羊群躺臥之處，你們就知道我是耶和華。 \end{tabularx} \\ \\ \relax
25:6 & \begin{tabularx}{0.7\textwidth}{X} 主耶和華如此說：因你們拍手頓足，幸災樂禍，藐視以色列地， \end{tabularx} \\ \\ \relax
25:7 & \begin{tabularx}{0.7\textwidth}{X} 所以，看哪，我要伸手攻擊你，把你交給列國作為擄物。我必從萬民中剪除你，從列邦中消滅你。我必除滅你，你就知道我是耶和華。」 \end{tabularx} \\ \\ \relax
25:8 & \begin{tabularx}{0.7\textwidth}{X} 「主耶和華如此說：因摩押和西珥人說『看哪，猶大家與列國無異』， \end{tabularx} \\ \\ \relax
25:9 & \begin{tabularx}{0.7\textwidth}{X} 所以，看哪，我要破開摩押邊界的城鎮，就是摩押人所誇耀的城鎮，伯‧耶施末、巴力‧免、基列亭， \end{tabularx} \\ \\ \relax
25:10 & \begin{tabularx}{0.7\textwidth}{X} 令東方人前來攻擊亞捫人。我必將亞捫交給他們為業，使亞捫人在列國中不再被記念。 \end{tabularx} \\ \\ \relax
25:11 & \begin{tabularx}{0.7\textwidth}{X} 我也必向摩押施行審判，他們就知道我是耶和華。」 \end{tabularx} \\ \\ \relax
25:12 & \begin{tabularx}{0.7\textwidth}{X} 「主耶和華如此說：因為以東向猶大家報仇，因向他們報仇而大大顯為有罪， \end{tabularx} \\ \\ \relax
25:13 & \begin{tabularx}{0.7\textwidth}{X} 所以主耶和華如此說：我要伸手攻擊以東，將人與牲畜剪除，使以東從提幔起，直到底但，地變荒涼，人也都倒在刀下。 \end{tabularx} \\ \\ \relax
25:14 & \begin{tabularx}{0.7\textwidth}{X} 我要藉我子民以色列的手報復以東；他們必照我的怒氣，按我的憤怒對待以東，以東人就知道施報的是我。這是主耶和華說的。」 \end{tabularx} \\ \\ \relax
25:15 & \begin{tabularx}{0.7\textwidth}{X} 「主耶和華如此說：因非利士人報仇，就是心存輕蔑報仇；他們永懷仇恨，意圖毀滅， \end{tabularx} \\ \\ \relax
25:16 & \begin{tabularx}{0.7\textwidth}{X} 所以主耶和華如此說：看哪，我要伸手攻擊非利士人，剪除基利提人，滅絕沿海剩餘的居民。 \end{tabularx} \\ \\ \relax
25:17 & \begin{tabularx}{0.7\textwidth}{X} 我要大大報復他們，發怒斥責他們。我報復他們的時候，他們就知道我是耶和華。」 \end{tabularx} \\ \\ \relax
26:1 & \begin{tabularx}{0.7\textwidth}{X} 第十一年某月初一，耶和華的話臨到我，說： \end{tabularx} \\ \\ \relax
26:2 & \begin{tabularx}{0.7\textwidth}{X} 「人子啊，因推羅向耶路撒冷說：『啊哈！那眾民之門已經破壞，向我敞開；它既變為廢墟，我必豐盛。』 \end{tabularx} \\ \\ \relax
26:3 & \begin{tabularx}{0.7\textwidth}{X} 所以，主耶和華如此說：推羅，看哪，我與你為敵，使許多國家湧上攻擊你，如同海洋使波浪湧上一樣。 \end{tabularx} \\ \\ \relax
26:4 & \begin{tabularx}{0.7\textwidth}{X} 他們要破壞推羅的城牆，拆毀它的城樓。我也要刮淨它的塵土，使它成為光滑的磐石。 \end{tabularx} \\ \\ \relax
26:5 & \begin{tabularx}{0.7\textwidth}{X} 推羅必成為海中的曬網場，因為我曾說過， 這是主耶和華說的。它必成為列國的擄物， \end{tabularx} \\ \\ \relax
26:6 & \begin{tabularx}{0.7\textwidth}{X} 推羅鄉間鄰近的城鎮必遭刀劍滅絕，他們就知道我是耶和華。」 \end{tabularx} \\ \\ \relax
26:7 & \begin{tabularx}{0.7\textwidth}{X} 主耶和華如此說：「看哪，我必使諸王之王，就是巴比倫王尼布甲尼撒，率領馬匹、戰車、騎兵、軍隊和許多人從北方來攻擊推羅。 \end{tabularx} \\ \\ \relax
26:8 & \begin{tabularx}{0.7\textwidth}{X} 他必用刀劍殺滅你鄉間鄰近的城鎮，也必築堡壘，建土堆，舉盾牌攻擊你。 \end{tabularx} \\ \\ \relax
26:9 & \begin{tabularx}{0.7\textwidth}{X} 他要安設撞城槌攻破你的城牆，以刀劍拆毀你的城樓。 \end{tabularx} \\ \\ \relax
26:10 & \begin{tabularx}{0.7\textwidth}{X} 因他馬匹眾多，塵土必揚起遮蔽你。他進入你的城門，如同進入已有缺口之城。那時，你的城牆必因騎兵、車輪和戰車的響聲震動。 \end{tabularx} \\ \\ \relax
26:11 & \begin{tabularx}{0.7\textwidth}{X} 他的馬蹄必踐踏你所有的街道；他必用刀劍殺戮你的居民。你堅固的柱子必倒在地上。 \end{tabularx} \\ \\ \relax
26:12 & \begin{tabularx}{0.7\textwidth}{X} 人必擄獲你的財寶，掠奪你的貨財；他們要破壞你的城牆，拆毀你華美的房屋，將你的石頭、木頭、塵土都拋在水中。 \end{tabularx} \\ \\ \relax
26:13 & \begin{tabularx}{0.7\textwidth}{X} 我要使你唱歌的聲音止息；人不再聽見你彈琴的聲音。 \end{tabularx} \\ \\ \relax
26:14 & \begin{tabularx}{0.7\textwidth}{X} 我必使你成為光滑的磐石，作曬網的場所。你不得再被建造，因為我---耶和華已這樣說了。這是主耶和華說的。」 \end{tabularx} \\ \\ \relax
26:15 & \begin{tabularx}{0.7\textwidth}{X} 主耶和華對推羅如此說：「在你中間行殺戮，受傷的人唉哼時，海島豈不都因你傾倒的響聲震動嗎？ \end{tabularx} \\ \\ \relax
26:16 & \begin{tabularx}{0.7\textwidth}{X} 那時沿海的君王都要從寶座下來，除去朝服，脫下錦衣，披上戰兢，坐在地上，不停發抖，為你而驚駭。 \end{tabularx} \\ \\ \relax
26:17 & \begin{tabularx}{0.7\textwidth}{X} 他們必為你作哀歌，向你說：『你這聞名之城，航海之人居住，海上最為堅固的，你和居民使所有住在沿海的人無不驚恐，現在竟然毀滅了！ \end{tabularx} \\ \\ \relax
26:18 & \begin{tabularx}{0.7\textwidth}{X} 如今在你傾覆的日子，海島都要戰兢；海中的群島見你歸於無有就都驚惶。』」 \end{tabularx} \\ \\ \relax
26:19 & \begin{tabularx}{0.7\textwidth}{X} 主耶和華如此說：「推羅啊，我要使你變為荒涼，如無人居住的城鎮；又使深水漫過你，大水淹沒你。 \end{tabularx} \\ \\ \relax
26:20 & \begin{tabularx}{0.7\textwidth}{X} 那時，我要使你和下到地府的人同去，到古時候的人那裡；我要使你和下到地府的人一同住在地的深處，在久已荒廢的地方，使你那裡不再有人居住；我要在活人之地顯榮耀。 \end{tabularx} \\ \\ \relax
26:21 & \begin{tabularx}{0.7\textwidth}{X} 我必叫你令人驚恐，使你不再存留於世；人雖尋找你，卻永不尋見。這是主耶和華說的。」 \end{tabularx} \\ \\ \relax
27:1 & \begin{tabularx}{0.7\textwidth}{X} 耶和華的話臨到我，說： \end{tabularx} \\ \\ \relax
27:2 & \begin{tabularx}{0.7\textwidth}{X} 「人子啊，要為推羅作哀歌。 \end{tabularx} \\ \\ \relax
27:3 & \begin{tabularx}{0.7\textwidth}{X} 你要對位於海口，跟許多海島的百姓做生意的推羅說，主耶和華如此說：推羅啊，你曾說：『我全然美麗。』 \end{tabularx} \\ \\ \relax
27:4 & \begin{tabularx}{0.7\textwidth}{X} 你的疆界在海的中心，造你的使你全然美麗。 \end{tabularx} \\ \\ \relax
27:5 & \begin{tabularx}{0.7\textwidth}{X} 他們用示尼珥的松樹作你的甲板，用黎巴嫩的香柏樹作桅杆， \end{tabularx} \\ \\ \relax
27:6 & \begin{tabularx}{0.7\textwidth}{X} 用巴珊的橡樹作你的槳，用鑲嵌象牙的基提海島黃楊木為艙板。 \end{tabularx} \\ \\ \relax
27:7 & \begin{tabularx}{0.7\textwidth}{X} 你的帆是用埃及繡花細麻布做的，可作你的大旗；你的篷是用以利沙島的藍色和紫色布做的。 \end{tabularx} \\ \\ \relax
27:8 & \begin{tabularx}{0.7\textwidth}{X} 西頓和亞發的居民為你划槳；推羅啊，你們中間的智慧人為你掌舵。 \end{tabularx} \\ \\ \relax
27:9 & \begin{tabularx}{0.7\textwidth}{X} 迦巴勒的長者和智者在你中間修補裂縫；海上一切的船隻和水手都在你那裡進行貨物交易。 \end{tabularx} \\ \\ \relax
27:10 & \begin{tabularx}{0.7\textwidth}{X} 「波斯人、路德人、弗人在你的軍營中作戰士；他們在你們中間懸掛盾牌和頭盔，彰顯你的尊榮。 \end{tabularx} \\ \\ \relax
27:11 & \begin{tabularx}{0.7\textwidth}{X} 亞發人和你的軍隊都駐守在四圍的城牆上，你的城樓上也有勇士；他們懸掛盾牌，成全你的美麗。 \end{tabularx} \\ \\ \relax
27:12 & \begin{tabularx}{0.7\textwidth}{X} 「他施因你多有財物，就作你的客商，他們帶著銀、鐵、錫、鉛前來換你的商品。 \end{tabularx} \\ \\ \relax
27:13 & \begin{tabularx}{0.7\textwidth}{X} 雅完、土巴、米設都與你交易，以人口和銅器換你的貨物。 \end{tabularx} \\ \\ \relax
27:14 & \begin{tabularx}{0.7\textwidth}{X} 陀迦瑪族用馬匹、戰馬和騾子換你的商品。 \end{tabularx} \\ \\ \relax
27:15 & \begin{tabularx}{0.7\textwidth}{X} 底但人與你交易，許多海島成為你的碼頭；他們拿象牙、黑檀木與你交換。 \end{tabularx} \\ \\ \relax
27:16 & \begin{tabularx}{0.7\textwidth}{X} 亞蘭 因你貨品充裕，就作你的客商；他們用綠寶石、紫色布、刺繡、細麻布、珊瑚、紅寶石換你的商品。 \end{tabularx} \\ \\ \relax
27:17 & \begin{tabularx}{0.7\textwidth}{X} 猶大和以色列地都與你交易；他們用米匿的小麥、餅、蜜、油、乳香換你的貨物。 \end{tabularx} \\ \\ \relax
27:18 & \begin{tabularx}{0.7\textwidth}{X} 大馬士革也因你貨品充裕，多有各類財物，就帶來黑本酒和白羊毛與你交易。 \end{tabularx} \\ \\ \relax
27:19 & \begin{tabularx}{0.7\textwidth}{X} 威但和從烏薩來的雅完人為了你的貨物，以加工的鐵、桂皮、香菖蒲換你的商品。 \end{tabularx} \\ \\ \relax
27:20 & \begin{tabularx}{0.7\textwidth}{X} 底但以騎馬用的座墊毯子與你交換。 \end{tabularx} \\ \\ \relax
27:21 & \begin{tabularx}{0.7\textwidth}{X} 阿拉伯和基達所有的領袖都作你的客商，用羔羊、公綿羊、公山羊與你交換。 \end{tabularx} \\ \\ \relax
27:22 & \begin{tabularx}{0.7\textwidth}{X} 示巴和拉瑪的商人也來與你交易，他們用各類上好的香料、各類的寶石和黃金換你的商品。 \end{tabularx} \\ \\ \relax
27:23 & \begin{tabularx}{0.7\textwidth}{X} 哈蘭、干尼、伊甸、示巴商人、亞述和基抹都與你交易。 \end{tabularx} \\ \\ \relax
27:24 & \begin{tabularx}{0.7\textwidth}{X} 這些商人將美好的貨物包在藍色的繡花包袱內，又將華麗的衣服裝在香柏木的箱子裡，用繩索捆著，以此與你交易。 \end{tabularx} \\ \\ \relax
27:25 & \begin{tabularx}{0.7\textwidth}{X} 他施的船隻為你運貨，你在海中滿載貨物，極其沉重。 \end{tabularx} \\ \\ \relax
27:26 & \begin{tabularx}{0.7\textwidth}{X} 划槳的把你划到水深之處，東風在海中將你擊破。 \end{tabularx} \\ \\ \relax
27:27 & \begin{tabularx}{0.7\textwidth}{X} 你的財寶、商品、貨物、水手、掌舵的、修補船縫的、進行貨物交易的，並你那裡所有的戰士和你中間所有的軍隊，在你傾覆的日子都必沉在海底。 \end{tabularx} \\ \\ \relax
27:28 & \begin{tabularx}{0.7\textwidth}{X} 因掌舵者的呼聲，郊野就必震動。 \end{tabularx} \\ \\ \relax
27:29 & \begin{tabularx}{0.7\textwidth}{X} 所有划槳的都從他們的船下來；水手和所有在海上掌舵的，都要登岸。 \end{tabularx} \\ \\ \relax
27:30 & \begin{tabularx}{0.7\textwidth}{X} 他們必為你放聲痛哭，撒塵土於頭上，在灰中打滾； \end{tabularx} \\ \\ \relax
27:31 & \begin{tabularx}{0.7\textwidth}{X} 又為你使頭光禿，用麻布束腰，號咷痛哭，痛苦至極。 \end{tabularx} \\ \\ \relax
27:32 & \begin{tabularx}{0.7\textwidth}{X} 他們哀號的時候，為你作哀歌，為你痛哭：有何城如推羅，在海中沉寂呢？ \end{tabularx} \\ \\ \relax
27:33 & \begin{tabularx}{0.7\textwidth}{X} 你由海上運出商品，使許多民族充裕；你以許多財寶貨物令地上的君王豐富。 \end{tabularx} \\ \\ \relax
27:34 & \begin{tabularx}{0.7\textwidth}{X} 在深水中被海浪打破的時候，你的貨物和你中間所有的軍隊都下沉。 \end{tabularx} \\ \\ \relax
27:35 & \begin{tabularx}{0.7\textwidth}{X} 海島所有的居民為你驚奇，他們的君王都甚恐慌，面帶愁容。 \end{tabularx} \\ \\ \relax
27:36 & \begin{tabularx}{0.7\textwidth}{X} 萬民中的商人向你發噓聲；你令人驚恐，不再存留於世，直到永遠。」 \end{tabularx} \\ \\ \relax
28:1 & \begin{tabularx}{0.7\textwidth}{X} 耶和華的話臨到我，說： \end{tabularx} \\ \\ \relax
28:2 & \begin{tabularx}{0.7\textwidth}{X} 「人子啊，你要對推羅的君王說，主耶和華如此說：你心裡高傲，說：『我是神明；我在海中坐諸神之位。』雖然你把你的心比作神明的心，你卻不過是人，並不是神明！ \end{tabularx} \\ \\ \relax
28:3 & \begin{tabularx}{0.7\textwidth}{X} 看哪，你比但以理更有智慧，任何祕密都不能向你隱藏。 \end{tabularx} \\ \\ \relax
28:4 & \begin{tabularx}{0.7\textwidth}{X} 你靠自己的智慧聰明得了財寶，把金銀收入庫房； \end{tabularx} \\ \\ \relax
28:5 & \begin{tabularx}{0.7\textwidth}{X} 你靠自己的大智慧以貿易增添財寶，又因你的財寶心裡高傲； \end{tabularx} \\ \\ \relax
28:6 & \begin{tabularx}{0.7\textwidth}{X} 所以主耶和華如此說：因你把你的心比作神明的心， \end{tabularx} \\ \\ \relax
28:7 & \begin{tabularx}{0.7\textwidth}{X} 所以，看哪，我必使外國人，就是列國中兇暴的人臨到你這裡；他們要拔刀摧毀你用智慧得來的美物，污損你的榮光。 \end{tabularx} \\ \\ \relax
28:8 & \begin{tabularx}{0.7\textwidth}{X} 他們必使你墜入地府；你要像被刺殺之人的死，死在海中。 \end{tabularx} \\ \\ \relax
28:9 & \begin{tabularx}{0.7\textwidth}{X} 在殺你的人面前，你還能說『我是神明』嗎？在殺害你的人手中，你不過是人，並不是神明。 \end{tabularx} \\ \\ \relax
28:10 & \begin{tabularx}{0.7\textwidth}{X} 你要死在陌生人手中，像未受割禮之人的死，因為我曾說過，這是主耶和華說的。」 \end{tabularx} \\ \\ \relax
28:11 & \begin{tabularx}{0.7\textwidth}{X} 耶和華的話臨到我，說： \end{tabularx} \\ \\ \relax
28:12 & \begin{tabularx}{0.7\textwidth}{X} 「人子啊，要為推羅王作哀歌，對他說，主耶和華如此說：你曾是完美的典範，智慧充足，全然美麗。 \end{tabularx} \\ \\ \relax
28:13 & \begin{tabularx}{0.7\textwidth}{X} 你在伊甸---神的園中，佩戴各樣寶石，就是紅寶石、紅璧璽、金剛石、水蒼玉、紅瑪瑙、碧玉、藍寶石、綠寶石、紅玉；你的寶石有黃金的底座，手工精巧，都是在你受造之日預備的。 \end{tabularx} \\ \\ \relax
28:14 & \begin{tabularx}{0.7\textwidth}{X} 我指定你為受膏的基路伯，看守保護；你在神的聖山上；往來在如火的寶石中。 \end{tabularx} \\ \\ \relax
28:15 & \begin{tabularx}{0.7\textwidth}{X} 你從受造之日起行為正直，直到後來查出你的不義。 \end{tabularx} \\ \\ \relax
28:16 & \begin{tabularx}{0.7\textwidth}{X} 你因貿易發達，暴力充斥其中，以致犯罪，所以我污辱你，使你離開神的山。守護者基路伯啊，我已將你從如火的寶石中殲滅。 \end{tabularx} \\ \\ \relax
28:17 & \begin{tabularx}{0.7\textwidth}{X} 你因美麗心中高傲，因榮光而敗壞智慧，我已將你拋棄在地，把你擺在君王面前，好叫他們目睹眼見。 \end{tabularx} \\ \\ \relax
28:18 & \begin{tabularx}{0.7\textwidth}{X} 你因罪孽眾多，貿易不公，褻瀆了你的聖所；因此我使火從你中間發出，燒滅了你，使你在所有觀看的人眼前變為地上的灰燼。 \end{tabularx} \\ \\ \relax
28:19 & \begin{tabularx}{0.7\textwidth}{X} 萬民中凡認識你的都必為你驚奇。你令人驚恐，不再存留於世，直到永遠。」 \end{tabularx} \\ \\ \relax
28:20 & \begin{tabularx}{0.7\textwidth}{X} 耶和華的話臨到我，說： \end{tabularx} \\ \\ \relax
28:21 & \begin{tabularx}{0.7\textwidth}{X} 「人子啊，你要面向西頓，向它說預言。 \end{tabularx} \\ \\ \relax
28:22 & \begin{tabularx}{0.7\textwidth}{X} 你要說，主耶和華如此說：『西頓，看哪，我與你為敵，我要在你中間得榮耀。』我在它中間施行審判、顯為聖的時候，人就知道我是耶和華。 \end{tabularx} \\ \\ \relax
28:23 & \begin{tabularx}{0.7\textwidth}{X} 我必令瘟疫進入西頓，使血流在街上。刀劍從四圍臨到它，被殺的要仆倒在其中；人就知道我是耶和華。」 \end{tabularx} \\ \\ \relax
28:24 & \begin{tabularx}{0.7\textwidth}{X} 「四圍恨惡以色列家的人，對他們必不再如刺人的荊棘，傷人的蒺藜；他們就知道我是主耶和華。」 \end{tabularx} \\ \\ \relax
28:25 & \begin{tabularx}{0.7\textwidth}{X} 主耶和華如此說：「我將分散在萬民中的以色列家召集回來，在列國眼前向他們顯為聖的時候，他們仍可在我所賜給我僕人雅各之地居住。 \end{tabularx} \\ \\ \relax
28:26 & \begin{tabularx}{0.7\textwidth}{X} 他們要在這地上安然居住。我向四圍恨惡他們的眾人施行審判之後，他們要建造房屋，栽葡萄園，安然居住，他們就知道我是耶和華---他們的神。」 \end{tabularx} \\ \\ \relax
29:1 & \begin{tabularx}{0.7\textwidth}{X} 第十年十月十二日，耶和華的話臨到我，說： \end{tabularx} \\ \\ \relax
29:2 & \begin{tabularx}{0.7\textwidth}{X} 「人子啊，你要面向埃及王法老，向他和埃及全地說預言。 \end{tabularx} \\ \\ \relax
29:3 & \begin{tabularx}{0.7\textwidth}{X} 你要說，主耶和華如此說：埃及王法老，你這臥在自己江河中的海怪，看哪，我與你為敵。你曾說：『我的尼羅河是我的，是我為自己造的。』 \end{tabularx} \\ \\ \relax
29:4 & \begin{tabularx}{0.7\textwidth}{X} 我必用鉤子鉤住你的腮頰，令江河中的魚貼住你的鱗甲；我要把你和所有貼著鱗甲的魚從你的江河中拉上來。 \end{tabularx} \\ \\ \relax
29:5 & \begin{tabularx}{0.7\textwidth}{X} 我要把你和江河中的魚全都拋棄在曠野；你必仆倒在田間，無人收殮，無人掩埋。我已將你給了地上的走獸、空中的飛鳥作食物。 \end{tabularx} \\ \\ \relax
29:6 & \begin{tabularx}{0.7\textwidth}{X} 「埃及所有的居民必定知道我是耶和華。因為你已成為以色列家蘆葦的杖； \end{tabularx} \\ \\ \relax
29:7 & \begin{tabularx}{0.7\textwidth}{X} 他們用手掌一握，你就斷裂，傷了他們的肩；他們靠著你，你卻折斷，閃了他們的腰。 \end{tabularx} \\ \\ \relax
29:8 & \begin{tabularx}{0.7\textwidth}{X} 所以主耶和華如此說：我必使刀劍臨到你，把人與牲畜從你中間剪除。 \end{tabularx} \\ \\ \relax
29:9 & \begin{tabularx}{0.7\textwidth}{X} 埃及地必荒蕪廢棄，他們就知道我是耶和華。「因為法老說『尼羅河是我的，是我所造的』， \end{tabularx} \\ \\ \relax
29:10 & \begin{tabularx}{0.7\textwidth}{X} 所以，看哪，我必與你和你的江河為敵，使埃及地，從密奪到色弗尼，直到古實邊界，全然廢棄荒蕪。 \end{tabularx} \\ \\ \relax
29:11 & \begin{tabularx}{0.7\textwidth}{X} 人的腳不經過，獸的蹄也不經過，四十年之久無人居住。 \end{tabularx} \\ \\ \relax
29:12 & \begin{tabularx}{0.7\textwidth}{X} 我要使埃及地成為荒蕪中最荒蕪的地，使它的城鎮變為荒廢中最荒廢的城鎮，共四十年之久。我必將埃及人分散到列國，四散在列邦。 \end{tabularx} \\ \\ \relax
29:13 & \begin{tabularx}{0.7\textwidth}{X} 「主耶和華如此說：滿了四十年後，我必招聚分散在萬民中的埃及人。 \end{tabularx} \\ \\ \relax
29:14 & \begin{tabularx}{0.7\textwidth}{X} 我要令埃及被擄的人歸回，使他們回到本地巴特羅。在那裡，他們必成為弱小的國家， \end{tabularx} \\ \\ \relax
29:15 & \begin{tabularx}{0.7\textwidth}{X} 成為列國中最低微的，不再自高於列邦之上。我必使他們變為小國，不再轄制列邦。 \end{tabularx} \\ \\ \relax
29:16 & \begin{tabularx}{0.7\textwidth}{X} 埃及必不再作以色列家的倚靠，卻使以色列家想起他們仰賴埃及的罪。他們就知道我是主耶和華。」 \end{tabularx} \\ \\ \relax
29:17 & \begin{tabularx}{0.7\textwidth}{X} 第二十七年正月初一，耶和華的話臨到我，說： \end{tabularx} \\ \\ \relax
29:18 & \begin{tabularx}{0.7\textwidth}{X} 「人子啊，巴比倫王尼布甲尼撒令他的軍兵大力攻打推羅，以致頭都光禿，肩都磨破；然而他和軍兵雖然為攻打推羅花這麼多力氣，卻沒有從那裡得到甚麼犒賞。 \end{tabularx} \\ \\ \relax
29:19 & \begin{tabularx}{0.7\textwidth}{X} 所以主耶和華如此說：我要將埃及地賜給巴比倫王尼布甲尼撒；他必擄掠埃及的財富，搶奪它的擄物，擄掠它的掠物，用以犒賞他的軍兵。 \end{tabularx} \\ \\
[1ex]
\hline
\hline
\end{longtable}
耶和華的榮光充滿聖殿.在過去的日子,耶和華的榮耀曾與以色列人同在,但亦曾經離開過聖殿和以色列人；因為以色列人失敗犯罪,虧缺了神的榮耀,而且忽略了神的榮耀.神不得已離開他們,神實在不願意,不捨得!但神是聖潔公義的,所以不能不離開犯罪的以色列人,和祂的聖殿.
\newline
\newline
耶和華的榮光離開以色列民,這正是他們歷史信仰的經歷.教會在這時代有沒有這個危險,能否儆覺信仰的危機?希望悲劇不重演.求主的靈光照我們,感動我們；激動責備俄們,讓我們知道怎樣在神面前謹慎.
\newline
\newline
神的榮耀離開以色列人,有一個專有名詞叫「以迦博」(撒上四21)當以利的孫子生下來時,起名以迦博,是因耶和華的榮耀離開了他們.這不但是以利家族的悲劇,也是整個以色列民的悲劇.祭司以利是撒母耳從小和他學習事奉神的那位神的老僕人,他有屬靈的經歷,他教導撒母耳認識神呼召的聲音；他也有事奉的經驗；多年在聖殿事奉神.當他年老時,沒有在神前儆醒忠心事奉；而且他的家庭完全失敗,他兒子犯罪作惡,甚至把聖潔之地成為淫污之所.
\newline
\newline
各位!我們多麼需要在神面前儆醒!一個神的僕人,雖然有經歷,有經驗,曾經被神使用,尚且會被神丟棄；以利失敗,兒子也失敗,甚至他孫子的名也叫「以迦博」；以致以色列民也陷在非常可怕的前途中.
\newline
\newline
求主憐憫教會!教會在這時代必須在神面前好好省察,讓神的僕人在我們禱告中支持；而我們也要用禱告支持他們；以免因一人的失敗,成一個家庭的失敗,就影響到整個團體,整個教會.一個亞干足以影響以色列全會眾,叫他們在敵前失敗.一個以利不但使他的家庭失敗,也叫整個民族陷於險境.所以你我所處的是多麼重要的地位,所肩負的是多麼重要的責任.求神憐憫我們!叫我們知道必須為自己,為神的工人,為教會負責人禱告,免得神的榮耀離開了我們.
\newline
\newline
教會在這時代,要警覺信仰的危機；因為神是輕慢不得的.人種的是甚麼,收的也是甚麼；我們必須在神面前恭敬,聖潔,規矩,謹慎；不要倚靠以往的虔誠,屬靈的成功；因為我們今天的熱心不能確保昨天的成功,所以我們必須時時刻刻隨時隨地在主面前謹慎.
\newline
\newline
以西結所見的異象,有榮耀在聖殿裡.神的寶座本在基路伯護衛之下.當以色列人失敗時,神就必須離開他們.當神的榮耀離開時,寶座就上昇.但神捨不得離開祂的百姓,還在那裡等候,纏綿,留戀.神的愛如此!雖然神的公義,要責備以色列人,因為審判要從神的家起首.但神還寬容等待,希望以色列人悔改；神一再忍耐等待,先知流淚嘆息地等待以色列人回轉；可惜他們沒有回轉,所以神的榮耀再離開了.
\newline
\newline
有一位蘇格蘭詩人馬德遜,他是神重用的僕人；他在大學讀書時奉獻作傳道,突然他的眼瞎了,因此未婚妻就和他解除婚約.他清楚神呼召他,他就越加盡力親近神.在他寫作的詩中,有幾句很出名的話說:「神啊!我曾經為玫瑰感謝你千次,但從未一次為玫瑰的刺感謝你.」還有一次他禱告說:「神啊!感謝你!因為你愛我,感謝你!因為你的愛裡有智慧.」他真是個精神新鮮的人；始終如一忠心於神.還有一首詩歌是極美的:「噢!神的愛不會離開我,一定不會離開我,是要我仰望祂；讓我透過流淚的眼睛,能看見永恆的彩虹.」
\newline
\newline
是的,神實在捨不得離開我們,但我們若一直犯罪,不肯悔改；並且惱羞成怒,將錯就錯；甚至每況愈下,變本加厲,神就不得不離開了.神在殿門口等候,依依不捨,但終於不得不離開了.「以迦博」神的榮耀離開了,實在太悲慘!
\newline
\newline
歷史的事實記載在聖經中,雖然神的信實會再回來；可憐以色列人已遭受了歷史空前的酷劫了.當他們痛定思痛時,纔會想到先知的話；那時不能不悔改,不能不在神前低頭敬拜,說:「神的道聖哉!善哉!」這幅圖畫,希望不再出現.
\newline
\newline
我們讀教會歷史,看見神的榮耀已多次離開!回來又離開.教會在這時代,不能讓神的榮耀離開；我們也不願又不敢讓神的榮耀離開,我們當禱告,求神的榮耀不離開我們.沒有任何事情比較神離開我們更加悲慘的!
\newline
\newline
當神的榮耀離開了掃羅,他就成了悲劇的人物.士師參孫,多麼悲慘!神曾用他,可惜他在婚姻上失敗,在道德上失敗；能力失去了,竟不自知.神的榮耀離開是逐漸而穩定的.當神看見我們的軟弱,就一再警告,提醒,施恩,責備我們,一直等我們回轉；我們若不回轉,神的榮耀,神的恩典就離開了,而我們自己卻不知也不覺.我們一點點的軟弱,一點點的過失,和罪惡污穢,就足以逐漸讓神的恩典離開我們.別以為少一次的聚會,少讀一節聖經,少一點禱告,少一點奉獻,少一點事奉沒有甚麼關係；其實卻大有關係；因為逐漸地,不知不覺地；神的能力,恩典,福樂就離開了.
\newline
\newline
各位弟兄姊妹!我們需要神的憐憫!我們不要像參孫,能力失去也不自知,連眼都瞎了,結果如牲畜一樣在磨坊推磨,落在厭煩的生活中.
\newline
\newline
摩西上西乃山,在神的面前禱告,四十晝夜之後；下山時,臉皮就發光而不自知.屬靈的長進自己不知道,屬靈的破壞自己也不知道.一樹棵倒下,但未倒之先已出問題.火車關掣後,還可繼續行兩步；所以,別以為自己還可以把恃!保羅說:「自己以為站立得穩的,須要謹慎,免得跌倒.」(林前十13)求神光照我們!
\newline
\newline
神的榮耀怎樣進到聖殿裡面呢?當神創亞當和夏娃,每一天開始之時,神都和他們交往,他們可以瞻仰神的榮光,生活充滿光亮.但他們吃了禁果犯了罪,就逃避神的面.神來尋找亞當,祂說:「你在哪裡?」人躲避神就失去了神的同在,失去了神的榮耀；所以神就把他們從伊甸園趕出去.人失敗了,但神永不失敗,神就為他們預備敬拜贖罪的方法.他們敬拜,也教導兒子們敬拜.當人敬拜時,就體會到神的同在,體會到神的榮耀.以色列人在曠野飄流,神的僕人摩西照著神指示的樣式建造會幕；神不厭其願,仔細周密地把敬拜的儀節指示他們；於是摩西就忠心按著神的命令指示去做.當會幕建成獻祭之時,耶和華的榮光就充滿了會幕.所以何時有敬拜,何時就體會到神的同在,就能看見耶和華的榮光.由會幕我們聯想到所羅門建造的聖殿,當聖殿完成獻祭之時,耶和華的榮光就充滿整個聖殿.可惜以色列人沒有專心跟隨神,他們一面敬拜神,一面卻去拜偶像；因此,神甚不喜悅.他們雖然經過屢次宗教改革,但他們還沒澈底悔改.猶大國最大一次的復興,是在約西亞作王年間,他重修聖殿,復興敬拜的事.表面看來,宗教已復興；可惜,他們只有敬拜的外表,而忽略了內裡的敬虔；以為有聖殿就有安全,可以不受外邦侵略,耶路撒冷就沒有外患了.然而,耶和華並不居住人手所造的殿；若人離開祂,祂的同在也必離開他們,那麼,聖殿何用之有呢?耶利米書七章告訴我們,神叫先知站在聖殿門口,一而再,再而三地說:「這是耶和華的殿.」你們若改正行動作為,我就使你們在這地方仍然居住.倘若他們不悔改就一定會敗亡,聖殿一定被毀壞；因為沒有耶和華的榮耀,沒有耶和華榮耀的同在.我們清楚看到,這是歷史的事實,歷代神用祂的話警告教會,叫我們儆醒.教會在這時代須要儆惕,因已面臨信仰上的危機.
\newline
\newline
主耶穌臨世,道成肉身,充充滿滿有恩典有真理,住在我們中間.意即主是會幕,在我們中間支搭,帳棚是神與人相會之處.主受難前,與門徒在一起,有一個門徒對祂說:「這是何等的殿宇!」耶穌卻回答說:「將來在這裡,沒有一塊石頭留在石頭上,不被拆毀了.」祂曾說:「你們拆毀這殿,我三日內要建立起來.祂這話是以自己的身體為殿說的,因為祂乃是神榮耀的顯現.我們卻沒有看見過神和神的榮耀,只有父懷裡的獨生子可以表明出來.主是道路,真理,生命,若不藉著祂,就沒有人能到父那裡.主被釘十字架,祂的榮光好像熄滅了；但主從死裡復活,而且那天,祂的身體豈不是已經得了榮耀嗎?因祂的身子道成肉身,表彰神的榮耀.現在祂的榮耀祂的身體在哪裡?祂的身體就是教會,換言之,教會是神的道繼續地成為肉身,教會有神的榮耀.「教會是祂的身體,是那充滿萬有者所充滿的.」(弗一23)
\newline
\newline
我們的教會有沒有神的榮耀?神是輕慢不得的；如果我們忽略,犯罪,墮落,神的榮耀可能離開我們,所以我們要細心省察!
\newline
\newline
昨晚我們特別注意先知的經歷,因教會是先知的團體,今晚我們所注意的是祭司的團體；因我們都是君尊的祭司,我們敬拜神,向神獻祭；以神的話語教導人,醫治人,幫助人,體恤人,好像祭司一樣.如果我們失敗,神的榮耀就會離開我們.我們有否好好地敬拜神?
\newline
\newline
今天的教會在敬拜上是失敗了,我們個人缺少敬拜,沒注重靈修生活,沒有用心靈誠實事奉神；終日以工作為前提,以為參加教會許多活動就是事奉,其實不然.真正的事奉,要從心靈開始.保羅在(羅一9)說:「用心靈所事奉的神.」事奉的真正意思,就是敬拜.所以羅馬書十二章說:「要將身體獻上,當作活祭......如此事奉,乃是理所當然的.」約書亞記廿四15說:「至於我,和我家,我們必定事奉耶和華.」事奉不是指工作而是指敬拜,要誠心實意事奉神.我們個人和家庭,有沒有好好的敬拜神?很多基督徒家庭,沒有家庭禮拜.在家庭若沒有敬拜怎能在教會的大家庭有好的敬拜呢?我們在敬拜上失敗了!我們以為上禮拜堂聽道就是敬拜,這樣的敬拜神是不悅納的.真正的敬拜要從心靈開始,個人開始；從小家庭至大家庭,然後從大家庭又帶回小家庭.般基督徒雖然每主日到教會,不過聚會稍長便不勝其煩,缺乏真正敬拜的態度.敬拜是漫長的路,帶我們步步更親近神.
\newline
\newline
創世記載,亞伯和該隱,在父母教導之下曉得敬拜；他們獻祭,該隱敷衍塞責,亞伯卻認真奉獻；所以神看中了亞伯和他的供物,卻看不中該隱和他的供物.因此該隱變了臉色,他在敬拜上失敗了.
\newline
\newline
弟兄姊妹!如果你在生活中沒有快樂,有憂慮,嫉妒,煩惱；表明你在敬拜上失敗了.當我們在奉獻上失敗之時,我們在敬拜上就失敗；如果我們在敬拜上失敗,整個基督徒生活就完全失敗了.神怎能與我們同在呢?神實在不願離開我們.當時神不願離開該隱,還安慰他說:「你若行得好,豈不蒙神悅納!你若行得不好,罪就伏在門前.」後果是不堪設想!該隱沒有悔改,殺了他的兄弟亞伯；以致一生過著流離飄蕩的生活.
\newline
\newline
我們平日的敬拜是「零」的,主日的敬拜是「整」的；我們應該化零為整,進教會敬拜以後,到家庭還要再敬拜,化整為零.這是基督徒必須有的.
\newline
\newline
教會常講奉獻的道理,林後也告訴我們「損得樂意是神所喜悅的.」我們常把這節聖經應用錯了,其實,這裡所指是附帶的奉獻.舊約燔祭是主要的奉獻,平安祭素祭乃屬附帶的奉獻；不獻也不算犯罪,但若獻得樂意是神所喜悅的.至於燔祭乃屬主要的奉獻,不論家庭個人,樂意或勉強,一定要獻；不獻就是犯罪.教會對此常沒有弄清楚.求神憐憫我們!我們應該拿主要的奉獻教會,然後奉獻額外的附帶的,我們獻得樂意是神所喜悅的.
\newline
\newline
我們每個人都是君尊的祭司,我們每個人都要敬拜,每家庭都要敬拜,然後大家庭好好敬拜.我們個人獻上,家庭獻上,全教會獻上,這纔是教會的復興.
\newline
\newline
各位弟兄姊妹!教會在這時代,已面臨信仰上危機；只片面而不切實,不能真正蒙受神的恩典.我們是活石,被建造成為靈宮.我們是君尊的祭司.(彼前二9)由此可知,我們參加培靈會,我們也要參加教會；要把培靈會的熱心,帶到教會.我們每個人都要下決心,說:「主啊!我要愛教會.」因為教會是基督的身體,今日教會仍繼續「道成肉身.」教會是神榮耀的居所；培靈會會過去,等到明年纔再有；但教會是恆常的,讓我們衷心思想,我們對教會如何?
\newline
\newline
我們必須建造神的靈宮,創世記十一章載,有好多人在建造巴比塔；但我們是屬神的,我們當建造靈宮.很多時候,我們受了世俗的影響,受了時代逆流的衝擊；身在教會,不建造靈宮,卻建造巴比塔,如同當日的以色列人一樣.神不能住在巴比塔,只能住在靈宮；我們若建造巴比塔,神就離開我們了.我們必須建造靈宮,纔真正在神裡面,神也在我們裡面.祂的榮耀,祂榮耀的同在和我們一起,我們真正地蒙福且恆久蒙福.看哪!神的帳幕在人間.(啟廿一3)以馬內利,就是神與我們同在,這是神始終如一的心意,祂必要實現.讓我們恆心仰望祂吧!
\newline
\newline


\newpage

\section{第54屆~港九培靈會~唐佑之~第三講虛}
\label{sec:547}
\textbf{謊與真理}
\newline
\newline
連結: \href{https://www.hkbibleconference.org/session-message/view/547}{\texttt{https://www.hkbibleconference.org/session-message/view/547}}
\newline
\newline
\hyperref[sec:546]{< < < PREV SERMON < < <}
~
\hyperref[sec:toc]{[返目錄]}
~
\hyperref[sec:548]{> > > NEXT SERMON > > >}
\newline
\newline
神選召以西結說預言攻擊那些假先知,因為他們所說的話,不是耶和華的信息；不是真理,而是虛謊的話.
\newline
\newline
神的教會應有儆惕的心,教會在這時代,應分辨真偽的話.
\newline
\newline
以色列人走錯了路,因聽信那些假先知說不正確的話.以西結照著神的吩咐向他們說明,並指明假先知的錯謬.有的假先知完全用迷信方法說預言,他們有男的也有女的；可能他們所作的具有吸引力,容易使人受迷惑.
\newline
\newline
以西結照著神的吩咐說話,可能當時的人都不歡喜聽；但以西結照著神的吩咐向他們說明,並指明假先知的錯謬.有的假先知完全用迷信方法說預言,他們有男的也有女的；可能他們所作的具有吸引力,容易使人受迷惑.
\newline
\newline
以西結照著神的吩咐說話,可能時的人都不歡喜聽；但以西結是神忠心誠實的工人,必要傳講神的話.當時的人反應如何,不得而知；然而神的話有能力,一出來就有功效,一解開就有亮光.不過,神藉祂僕人所講的話,不一定是甘甜的,有時卻是苦澀的；但是神的話,人不能不尊重,不能不重視,不能不順從,便不能不遵行.神的話或者在當時沒有果效,可是以後仍有功效；在每個時代,今時今日都有功效.
\newline
\newline
教會在這時代,須明白神的真理,須遠離虛謊的話.今日教會有兩種危險,歷代教會也是這樣.不大熱心的基督徒,很容易受世俗影響；稍為熱心的,很容易受虛謊的影響.魔鬼用的似乎是真理,其實是謊言,來迷惑這樣的人.求神憐憫!我們既有追求的心,但決不可受迷惑；我要在主面前,要在主的真道上站穩.
\newline
\newline
有一種先知是完全迷信完全錯誤的,他們用的是拜偶像的方式.這一煩較易預防.一提到假先知,即想到他們是異端邪說傍門左道之輩；使我們隨時提防,遠離他們,不上勾當.但是神叫以西結特別攻擊的那些假先知,表面上卻看不出來；他們隨己意行事說話,沒有明白神的心意,沒有真正屬靈的感受.為了吸引人,滿足私慾,就用各種方法叫多人受迷惑.
\newline
\newline
教會在這時代也有此危險.有的人看來信仰很純正,他們標榜屬靈,言語屬靈；但所作所為,非屬神旨意.凡此都是假先知,我們要格外小心,因真假先知殊難分辨.我們多麼需要神的憐憫,讓真理的聖靈引導我們.
\newline
\newline
一個真先知,可能有時也說虛謊的話.在以西結書,我們發現,有時神要以西結講話,但有時卻不許他講話；如果不准他講話,他卻講話,那就可能他也會成為假先知了.有一些知名的傳道人,他們被神重用；但如果他們講了神沒有吩咐他們講的話,那就是虛謊的話.總而言之,人彼此都有這樣的危險,神的工人也有此危險,聽道的人也有危險.教會在這時代,須有分辨的力量.華人教會對此不甚注意,以至造成很大的虧損.
\newline
\newline
回憶一九四九年,很多人來到香港,那時人心渴慕真道,教會很多聚會,聽眾很多.當時在九龍和香港,甚至有連續百天的聚會；天天人群擁擠到教會聽道,願意在真道上領受,教會非常蓬勃.可惜有的人沒有過正常的教會生活；只追隨講員,這是極不正常的現象.聽道者因為態度有問題,結果產生了一些特殊講員,他們大有吸引力,也很會講道；然而他們講的,有放在神前從新衡量的必要；因為我們不敢隨便批評論斷.若出於神的話是會存在的,若不是出於神的話,就會使人受迷惑,反而得不到造就.那些熱心聽道的人,今天在哪裡呢?我們看見現在的培靈會,大部份赴會者都是年青的一輩.求神憐憫我們,讓我們不要重複卅多年前,那些熱心聽道的錯誤.聽道者和講道者,都該明白應有的態度；有時是出於講員的錯誤,領導聽眾產生不正常的心態.而另一方面亦因聽眾的錯誤,鼓勵講員往錯誤的方向走.講員為了迎合會眾的心理,講得動聽叫人注意,讓自己出名.歷代常有這樣的危險,當時以色列人就是這樣.其實先知本來不一定是壞的,不過有些為了聽眾的要求,而聽眾亦為了滿足講員的虛榮心而變了質.有人為名為利,竟把神的話扭轉來講,真是危險!求神憐憫!教會在這時代,不能再有此錯誤.
\newline
\newline
講員怎會領導聽眾走錯路呢?因為沒有清楚神給他的託付.講道者恩賜各有不同,有的有傳福音的恩賜,有的有教師的恩賜,有的有牧師的恩賜,而且傳講的性質也不同.華人教會聽眾所要求的,乃是屬奮興家姿態的講員.奮興家當然很可吸引人,在傳福音事工上很有功效；講得活潑生動,使未信者明白福音真理；但是,並非每個講員者都是奮興家.弟兄姊妹常覺自己教會的牧師,不及奮興家會講道；因此迫使牧師要以奮興家講道的方式來講道.在過往的數十年,神在中國教會曾興起幾個奮興家；他們走遍大江南北,到處傳福音,在國內外使許多人歸主,蒙受主的福份.但是這些人並非是神僅賜的傳道人,神還賜給教會,各種不同恩賜的傳道人.一般信徒以為奮興家,纔是神所重用的傳道人,他們樣樣都行；因此鼓勵奮興家,建立教會,牧養教會.在教會專事查經培靈工作,但神的託付卻非如此,以致多年來造成中國教會極大的損失.教會在這時代,不要再重複以往的失敗!我們應看清楚恩賜,尊重每個恩賜；那些特殊講員,對我們的幫助不是很大.對我們最大幫助的,乃是教會中神的工人,他們有牧養的恩賜,惟有他們纔真正知道我們的需要.有人傳福音,一定要講些故事,見證,有時滑稽幽默,這是無可厚非的；但弟兄姊妹常會弄錯,以為聽道講道都該如此.結果,中國教會只能站在道理的開端,卻難以進到完全的地步,因為沒有在主的恩典和知識上長進.請不要誤會,我不是說講故事見證不對；有時確需要這些來解釋聖經的真理,來見證神的能力,讓人看見福音的好處.有時實在需要幽默；因為人們覺得苦悶,需要清新的感覺.有佈道恩賜的人,一步步引人到主前信靠福音,蒙恩得救；他們實在功不可沒.但是若所有講道都用此方式,那麼,我們一直不會長進；我們只會吃奶,不會吃乾糧,只能做小孩；太幼稚,太膚淺,太表面；整日講恩典,卻在恩典中墮落.我們只聽道,不能行道；只覺得聽道是一種享受,聽了算了,這太危險了!因為大家有這要求,給了講員極大的試探,為名為利而採用這種的講道方式；可能所講非神所許,他們竟作了假先知,講了虛謊的話；弟兄姊妹也就走上虛謊的路,永不長進.故事雖動聽,但神的話比一切更重要,見證雖好,但有時未免太過渲染,要見證神的能力,卻不知不覺離開了神的話.許多人沒有屬靈的吸收力,只聽簡單有趣的；乍吃覺得甘甜,不過,甜也有程序上之分；如瘋狂下咽,甜一甜就完了.我們所需是純淨不加水,不混雜質,且富營養的食物,纔能保持身體健康.希望華人教會,在座每一位,我們須在聽道的程度上長進,吃得乾糧,真正接受神說不盡恩典的豐富.關於解經方面也有很大錯誤,華人教會有很多的解經是錯誤的.為了滿足聽眾好奇心,過份靈意解經,過份用寓意方法解經；講的頭頭是道,聽眾耳所未聞,覺得講者對神的話很有亮光.其實完全離開原意,聖經並無記載,乃出於人的推測；不是出於神的話,那是假先知,是虛謊.
\newline
\newline
有次我到某教會講道,主席禱告說:「求主幫助今天的講員,讓他所講的是我們從未聽過的.」這使我實在難以啟齒,因為我所要講的都是人所聽過的.我們聽道不應為了滿足好奇心和屬靈的虛榮心.講員用己意解釋預言,成為虛謊的話.為迎合聽眾興趣,不敢直接了當把罪講解；有時把神的話講得旁敲側擊,迂迴婉轉,其實不一定屬神的意思；說的平安,其實沒有平安.彌迦書三章說,有的先知偽報平安.如果大家輕輕鬆鬆地參加聚會,會畢平安無事回家；這樣到底有助於靈性嗎?絕對沒有.耶利米是位哀哭的先知,他為了以色列民的厄運；為了他們遭受神的審判,為了許多人走上滅亡的道路而哀哭.當時卻有一班先知,專講輕描淡寫的話「平安了,平安了,」其實沒有平安.「他們輕輕忽忽的醫治我百姓的損傷.」其實有時需要動手術,痛苦的經驗是必要的；因為神的話好像兩刃的利劍,把我們刺入,割開,叫我們能赤露敞開,澈底得醫治.以色列民所以敗亡,乃因沒有澈底醫治.
\newline
\newline
教會在這時代,需要澈底的醫治,不能輕忽,隨便,馬虎,間接的講解神的話.盼望弟兄姊妹多為講道的人禱告,且切實分辨聽的是真理抑或虛謊.這樣,華人教會纔能長進,健壯,真正地往前走.我們個人的生命也應如此.
\newline
\newline
第十節說,那些假先知好像未泡透的灰抹在牆上,表面而已,不能持久.神的教會,須在神的話語裡,澈底泡透,纔有功效；神的工人應在真道上,在神的話語裡下工夫.華人教會每個信徒當在神的話語裡下工夫,澈底得著神的話,浸透在神的真理中；有分辨的力量,纔會見證,遵從而且實行.
\newline
\newline
何西阿書說,以色列人好像個沒有翻過的餅.一面焦黑,一面生而不熟.今日華人教會就是這樣不平衡.我們太少讀經,對神的真理不瞭解,有的人只讀新約,不讀舊約；讀保羅書信,反而不多讀耶穌基督的生平；讀詩篇沒有讀啟示錄；讀箴言沒有讀先知書.我們分題查經只有一個題目,慕道班,主日講道,查經班,禱告會,主日學都是只有一個題目；並非說系統的分題的不好,不過是我們對神的真理沒有整個的瞭解；正像沒有翻過的餅一樣.求主憐憫!叫我們平衡地瞭解神的話語.神的教會要長進,弟兄姊妹要長進,惟有在神的真道上造就.
\newline
\newline
第五節說那些假先知沒有上去堵破口,沒有重修垣牆,也沒有在陣上作戰.以色列人認為垣牆很重要,可以防守敵人任意侵入,和隨意破壞；若有破口,城牆陷落,必定失敗.今日教會到處是破口,亟需修牆準備爭戰.有些基督徒以為所信的是位慈愛的神,神是我們的避難所,是我們的山寨,是我們的力量；信了主可以安穩上天堂.這是舊日一般的觀念；但我相信今之青年基督徒不是這樣,我們知道教會在這時代是爭戰的教會,我們須負起爭戰的任務；靠著神的力量出去爭戰,非但靠神得安而已.神是萬軍的耶和華,我要作祂的精兵,勇敢為祂爭戰.青年基督徒們!我們要起來作戰,好好修理教會的垣牆.
\newline
\newline
當尼希米發現耶路撒冷城牆破壞,他立刻禁食禱告,在神前澈底悔改.難道是他一人的罪過嗎?絕對不是.今天有很多青年基督徒很愛主,很熱心；看見教會不愉快的情形,就論斷批評,這是不健全的.尼希米得波斯王的許可回猶大國,他沒有馬上建造城牆,而是先行非常冷靜觀察破口的情形,然後著手吩咐各人對著自己的房屋修造；於是破口堵塞,城牆建立.這就是教會復興的途徑,是我們應當作的.假先知沒有堵塞破口,建設城牆,所以不能長大.我們不能隨意聽信一班人講道,我們當在自己的教會,謙卑在神工人面前學習,在神面前領受,在神話語裡下工夫,真正浸入真理裡面,讓真理的靈感動我們,切實在真理中長進.求神叫我們各人明白這個意思,切勿隨便聽道,要有分辨的力量；若有人講錯,沒有照神的意思講,我們萬勿鼓勵他,甚至要照著我們所瞭解的糾正,為他禱告；不是由我們來指摘而是讓神來指責.我們當在教會中安守本分,來領受神的真理.
\newline
\newline


\newpage

\section{第54屆~港九培靈會~唐佑之~第四講}
\label{sec:548}
\textbf{脆弱與堅韌}
\newline
\newline
連結: \href{https://www.hkbibleconference.org/session-message/view/548}{\texttt{https://www.hkbibleconference.org/session-message/view/548}}
\newline
\newline
\hyperref[sec:547]{< < < PREV SERMON < < <}
~
\hyperref[sec:toc]{[返目錄]}
~
\hyperref[sec:549]{> > > NEXT SERMON > > >}
\newline
\newline
以西結書 15:1-15:8
\newline
\begin{longtable}{cl}
\hline
\hline
章節 & 經文 (和合本修訂版)\\
\hline
15:1 & \begin{tabularx}{0.7\textwidth}{X} 耶和華的話臨到我，說： \end{tabularx} \\ \\ \relax
15:2 & \begin{tabularx}{0.7\textwidth}{X} 「人子啊，葡萄樹比一切其他的樹，就是樹林裡眾樹木的樹枝，有甚麼長處呢？ \end{tabularx} \\ \\ \relax
15:3 & \begin{tabularx}{0.7\textwidth}{X} 可以從其中取木料來做工嗎？人可以拿來做釘子，掛東西在上面嗎？ \end{tabularx} \\ \\ \relax
15:4 & \begin{tabularx}{0.7\textwidth}{X} 看哪，它已經拋在火中當柴燒，火既燒了兩頭，中間也燒焦了，它還有甚麼用處呢？ \end{tabularx} \\ \\ \relax
15:5 & \begin{tabularx}{0.7\textwidth}{X} 看哪，它完整的時候尚且不能拿來做工，何況被火燒焦了，還能拿來做工嗎？ \end{tabularx} \\ \\ \relax
15:6 & \begin{tabularx}{0.7\textwidth}{X} 所以，主耶和華如此說：我怎樣使林中樹裡的葡萄樹在火中當柴燒，我也必照樣對待耶路撒冷的居民。 \end{tabularx} \\ \\ \relax
15:7 & \begin{tabularx}{0.7\textwidth}{X} 我必向他們變臉；他們雖從火中逃出來，火仍要燒滅他們。我向他們變臉的時候，你們就知道我是耶和華。 \end{tabularx} \\ \\ \relax
15:8 & \begin{tabularx}{0.7\textwidth}{X} 我必使這地荒涼，因為他們做了背叛的事。這是主耶和華說的。」 \end{tabularx} \\ \\
[1ex]
\hline
\hline
\end{longtable}
當耶路撒冷未陷落,聖殿未被毀,以色列國還未敗亡之前；先知以西結已被擄到外邦.表面上,似乎他的工作受了限制,他不能向整個以色列會眾傳講神的話語；但另一方面,神保守他在外邦地,使他在巴比倫國的迦巴魯河邊,向被擄的以色列人作牧養的工作.他傳講神的信息；也讓那些回耶路撒冷的以色列人,帶回去神審判的信息,好叫整個以色列人知道.
\newline
\newline
以色列國好像葡萄樹,在聖經中不論舊約新約均屢次有題及.葡萄樹是一種奇特的植物,一方面看來很脆弱,樹枝柔軟,不能作有用的木材；它的真正價值,乃在於結果子.它的樹枝雖柔軟,但結果纍纍,枝子彎彎下垂卻不易折斷,所以,說它脆弱,或說它堅韌,兩方面說法都對.這表明神的選民,在性質方面既脆弱,而又是很堅韌的.照樣,神的教會,以及你我凡屬於祂的人也是如此.脆弱或堅韌,乃視乎樹枝本質的功用而定.神對以色列人所說的話,正是這時代教會應該明白的信息.
\newline
\newline
舊約中有數處關於葡萄樹的信息,新約也有幾章重要的聖經,是關於論到葡萄樹的.首先我們思想到以賽亞書第五章,稱為葡萄園之歌；它是一章極美麗的希伯來抒情詩,但另一方面也可說是非常悲慘的信息.
\newline
\newline
「我要為我所親愛的唱歌......論他葡萄園的事.」神如何賜福以色列人,為他們預備栽種美麗且豐盛的葡萄園.「我所親愛的有葡萄園,在肥美的山崗上.」意思是既有美好的環境,又朝著陽光,地質極肥美,而且栽植了上好品種的葡萄樹.這裡充分說明神對以色列人的恩典,神揀選以色列為選民,並非說他們是特別優秀的民族,有高尚的道德,悠久的文化；而是完全出於神的恩典,揀選這個弱小民族來彰顯祂的大能.叫他們成為恩典的導管,藉以把神救贖的福音傳出去.因為神的救恩,不單為以色列人,也是為整個世界.神為以色列預備,美好的環境和最有利的條件,讓他們成為真正上好的葡萄園；神指望他們結好葡萄,但是反而結出野葡萄；因以色列人在神的恩典中墮落了,使神的心傷痛.照樣,我們今天也是如此,我們所以稱為屬於神的人,唯一的解釋乃是神的恩典.我們怎可以在神的恩典中墮落呢?因此,教會在這時代,當認清自己的地位,功能和應有的本份.(耶二21)耶利米先知所傳的是同樣的信息.「然而我栽你是上等的葡萄樹,全然是真種子；你怎麼向我變為外邦葡萄樹的壞枝子呢?」在何西阿書第十章,神使以色列成為茂盛的葡萄樹結果繁多.這不是指屬靈方面,是指物質方面.神是豐富的,使他們一無所缺；但是相反的,他們卻日漸墮落,沒有將榮耀歸給神,甚至去拜偶像；以致以色列人陷於悲慘的命運裡.這並非因神不賜福他們,乃因他們在神的恩典中墮落了.
\newline
\newline
從以色列民族歷史的信仰,看我們今天的教會,我們能不加以謹慎嗎?神為我們預備極豐富的一切,是要我們在祂的恩典中長進；不要在恩典中墮落,祂要我們明白得恩典,不是沒有條件的.是的,恩典是白白賜給我們的,但是請各位注意,恩典是附有條件的；如果我們不願領受,我們就得不著.如果我們沒真正在神面前珍貴恩典,我們能否享受神的恩典呢?神選召以色列時,藉著摩西告訴他們:「如今你們若聽從我的話,遵守我的約,就要在萬民中作屬我的子民......」(出十九 5)要順從,信服,遵行,神的恩典就臨到.這些就是條件,神的恩典不是不給我們,乃因我們拒絕了.
\newline
\newline
教會在這時代,必須在神的恩典中長進,要認識我們的地位,身份,也認識我們的本份.我們一定要結果子,不然,就會被丟在火裡燒了.
\newline
\newline
以色列人在神的恩典中墮落,神審判的火已經燒在葡萄樹的枝子上,兩頭都燒了,連中間都要燒.(結十五4)根據歷史,以色列民本有十二支派,在曠野時,摩西從神領受命令,建立會幕,在四邊安營居住,每邊三個支派；會幕邊緣,東有摩西亞倫,另三邊有利未人的子孫；他們都有工作,經理會幕的事.神這樣的安排,一方面叫他們有自由,完全自主,沒有中央極權的政體.摩西也不是他們的中央領袖,唯耶和華是王.每支派努力自治,當獻祭時,他們在敬拜中合而為一.這裡我們看見一幅教會的圖畫；我們可以有不同的教會,不同的名稱,不同的宗派；如同以色列人有十二支派,大家互不干涉,各自在神面前,妥善發揮地方教會的功能.但是要合一,不是組織的合一,而是屬靈的合一,在敬拜中合一.以色列人本是這樣,到了大衛在耶路撒冷建造京城,所羅門在耶路撒冷建造聖殿時；北方支派不滿南方支派,於是省域觀念產生,以致南北分裂.這在神眼中這是件極不合理的事.宗教成為他們的政治手段,甚至北方自立在聖所拜偶像.結果,北國早於南國敗亡；一端的葡萄樹燒起來了,因神的審判臨到他們了.另一端,南國滿得神的恩福；因為神是信實的神,他曾經應許大衛家有永遠的王位；但是後來他們犯罪,所以葡萄樹枝的另一端也燒起來了,連中間都燒了；耶路撒冷在南方之北靠近北方,可說是中部,故此,耶路撒冷也要被毀.神是公義的神,萬不以有罪為無罪；神是烈火,是輕慢不得的.我們在神面前須何等謹慎!
\newline
\newline
以色列人本如荊棘,該被焚毀,但神的愛真奇妙.當摩西在何列山,看見前面有燒著的荊棘燒不掉；因為神的愛繼續在以色列人身上,神不願意消滅他們,這是神在歷史中的計劃.以色列人失敗,神仍要復興他們；神的福音仍要傳給外邦人.歷代教會失敗,神仍要保守一些忠心的人.
\newline
\newline
我們多麼需要神的憐憫!我們如同葡萄樹枝,若不結果子,有何價值呢?當柴燒也不經燒的.神是公義的神,祂要焚燒一切,連我們也要完全被焚毀的；但是神捨不得,祂屢次想將我們拯救出來；祂赦免憐憫我們,把我們拯救出來,如同從火中抽出來的柴一樣.在阿摩司書和撒迦利亞書都有:「從火中抽出來的柴.」的這一句話(摩四11,亞三12)我們蒙了神的憐憫,但我們不能止於僅僅得救,我們要進一步蒙受神的恩典.
\newline
\newline
新約有段經文,主耶穌親口說葡萄園的主人,到葡萄園要收葡萄的果子.(可十二1-12)當然這背景是論到以色列人,拒絕了神的兒子,亦即拒絕了葡萄園主的兒子,我們知道園主要在園裡找果子.神今日亦向教會要收果子,我們有沒有果子奉上呢?約翰福音第十五章是我們很熟識的經文,耶穌基督清楚說,我們是葡萄樹的枝子；若枝子不結果子就被扔在火裡燒掉因它沒有存在的價值.
\newline
\newline
神是烈火,神的公義是火.在舊約曾經論到我們在神面前,當明白神的性格.有時說耶和華的熱心必成就這事,有時說神的忿怒,有時說神是忌邪的神.這三事同屬於一,都說到烈火.祂對以色列人的心是火熱的,對全世界人的心是火熱的；祂的愛,祂的公義,祂的忿怒,都是熱切的,視乎我們在祂面前如何反應吧.讓我們真正知道神是烈火,我們須在祂面前得著祂熱切的愛,而非在祂烈怒之下,這是我們需要認識的.
\newline
\newline
耶穌在葡萄園的比喻中說:「不結果子的枝子要燒掉,結果子的要修理乾淨,使枝子結果子更多.」祂要修理,潔淨,把修理的剪刀放在結果子的枝上；讓枝子的生命力集中,果子更加甘甜豐富.
\newline
\newline
各位!教會在這時代,必須結果,否則就被扔掉.
\newline
\newline
在這個時代有不少的教會,他們作神的工作,奉主名聚會,但卻被神棄絕.「主對以弗所教會說,你若不悔改,我就把你的燈台從原處挪去.」(啟二5)因他已完全失去教會的地位和功能.我們多麼需要神在我們的教會作工.今天華人教會四分五裂,本應在神前合一；但卻如同以色列南北分裂,都是人的作為,人的見解,偏見,污穢,罪惡.我們多麼需要神的憐憫!
\newline
\newline
有些人標榜自己的屬靈,以為宗派不屬靈,四分五裂；其實宗派是神在歷史中的作為,反對宗派的,明顯更有宗派的色彩.中國教會受了太大的損失,就是因為對教會沒有清楚的認識；如果我們還在分裂的話,就必像上述葡萄樹的枝子,一端燒著,另一端連中間都燒起來.那麼,見證在哪裡?教會在這時代,不要有那偏窄似是而非的信仰,最要者,必須在神面前結果子.
\newline
\newline
教會作修理潔淨的工作,表面看來,好像是反面的,是消極的,其實是積極的.該拔的拔,該掘的掘,該拆毀的拆毀；然後進行栽植.神所作的是澈底的工作,澈底的拆毀,澈底的拔去,然後栽植.(耶一10)澈底的工作不單是表面的改善,凡要結果子的,就要修理乾淨,使枝子結果子更多.
\newline
\newline
本書第十七章的歷史背景,當年的猶大王朝,非常脆弱,他們想用政治手法,維持一時的安全.當時有兩種世界的強權,有人說,一是巴比倫,一是埃及.到底依靠哪個呢?那時猶大王朝處於進退維谷的境況,先知說依靠任何方面都是不對的.很多時候教會也有這樣的危機,信徒們也有這樣的危險；我們常感覺到無所適從.事實上,甚麼都不可靠,「惟依靠神」這是先知的忠告.
\newline
\newline
在以賽亞時代,猶大國情況危殆,因敘利亞要聯盟北國進攻亞述；猶大若與敘利亞攜手,可能擊破亞述的強權；若依靠亞述,敘利亞必來攻打他.(賽七)以賽亞卻說,甚麼都不要依靠,惟有依靠神；但他們不肯聽先知的話,以致遭受極大的損失.
\newline
\newline
請看!東風(指巴比倫)一吹地就枯乾,葡萄樹不能維持下去了.這裡我們看到教會應有的路向,教會是屬乎主的,我們要做得正確.教會在這時代要有一定的方向,我們並非留戀過去,亦非要回復使徒時代的樣式.到底教會應回到耶路撒冷的路向,或安提阿的路向呢?都不是；因神為我們開了一條又新又活的路,所以我們別留戀過去；時代背景不同,當時可用而今天則不宜.教會要更新,並非不注重過去的傳統；但是神已為我們開了一條又新又活的路.教會需要復興；若是死氣沉沉,就是你我的問題,我們都有責任.我們都要起來作神的工,神要為我們常常修理乾淨.
\newline
\newline
你我在生活中,有沒有得過神修理潔淨的經驗呢?我們若肯將自己擺在神面前,祂就修理潔淨.我們必須承認自己的軟弱脆弱無用.耶穌說我們若離了祂,就不能作甚麼.另一方面,我們若靠著神的恩典,靠著那加給我們力量的,凡事都能作(腓四13).我們就從軟弱無用的葡萄樹枝,成為堅韌有價值能結果子的枝子.不過,如果我們自大驕傲,神就不能賜福給我們並使用我們.我們的一切,都是從神領受的,無可自誇之處.另一方面,神也不賜福給那些自卑的人.基督徒常自嘆不及他人而自卑,這是不健康的靈性,我們應當謙卑,在神面前自信.當我們信心堅強完全依靠神之時,我們深信工作必有價值,能結果子.只要忠心,信服仰望神,我們所作的,必超過我們所求所想的.保羅在(林後四7)說我們有寶貝放在瓦器裡.其實貴重的物件應放在貴重器皿裡.寶貝放在瓦器裡,是個恩典的原刖；寶貝不因瓦器而貶值,反而瓦器因藏有寶貝而增值百倍.但是,榮耀的福音,奇妙的救恩,偉大的救主；經我們的見證常打了折扣,貶了價值.我們萬勿如此?
\newline
\newline
我們無所自大,然而也不可自卑；因出於神的有莫大的能力.「能力」應繙為潛力.「聖靈降臨在你們身上,你們就必得著能力.」(徒一8)「十字架的道理,在那滅亡的人為愚拙,在我們得救的人卻為神的大能.」(林前一18)保羅說:「我不以福音為恥,因為福音本是神的大能.」(羅一16)神的能力在我們身上,潛在的能力,偉大的能力,莫大的能力,莫大的潛力,全都是我們不能想像的.正像葡萄樹枝,看來沒有用處；但可以結果子；好像很脆弱,卻是很堅韌；枝上果子纍纍,彎彎負荷也不折斷.我們怎可自卑輕看自己呢,我們當靠著神的大能大力,剛強起來.「莫大」在(弗一19)說:「祂向我們這信的人所顯的能力是何等浩大」,「浩大」就是「莫大」,「基督的愛是何等長闊高深,這愛是過於人所能測度的.」(弗三18-19)就是這個意思.(林前十二31)保羅說:「我現今把最妙的道指示你們.」浩大的能力!更妙更偉大的能力!教會在這時代,需要這樣的能力.有時我們的信心太小,屬靈的容量太窄；神給我們很大的能力,我們沒有善用.教會在這時代,實在有莫大的能力；以往忽略了,現在必須注意,纔能看見神更大更奇妙的工作.
\newline
\newline
「約瑟是多結果子的樹枝,是泉旁多結果子的枝子,他的枝條探出牆外.」(創四九22)雅各晚年的靈命真是登峰造極,到達爐火純青的地步,他臨終時扶著杖頭敬拜神.多麼恭敬屬靈的風度!他說約瑟是多結果子的樹枝......枝條探出牆外.這是一幅教會增長的美麗圖畫.教會增長不僅要結果子,且要讓枝條探出牆外.神的工作須在幅度方面增長.
\newline
\newline
本書十七章末段說,以色列不但是葡萄樹,更是香柏樹栽在高處；經得起風雪考驗,神叫矮樹升高,高樹降卑.神使教會在地上,真正有尊貴榮耀的地位.神看重以色列,神的心意揀撰他們,要他們把祂的恩典帶出去；他們失敗了,但神沒有失敗.
\newline
\newline
教會在這時代,應當明白神的心意,神期望我們如多結果子的葡萄樹.有生命水的供應,我們纔能多結果子,且能發展探出牆外.神不但叫我們作葡萄樹,更要叫我們作香港樹,在屬靈高處,經得起時代的考驗.教會在這時代,是神抱著最大希望的時候；祂等待著我們,要我們好好在祂面前仰望倚靠祂；領受祂的恩典,接受祂的修理,潔淨,叫我們結果子更多.
\newline
\newline
但願主這樣地引導我們!
\newline
\newline


\newpage

\section{第54屆~港九培靈會~唐佑之~第五講}
\label{sec:549}
\textbf{生命的定律}
\newline
\newline
連結: \href{https://www.hkbibleconference.org/session-message/view/549}{\texttt{https://www.hkbibleconference.org/session-message/view/549}}
\newline
\newline
\hyperref[sec:548]{< < < PREV SERMON < < <}
~
\hyperref[sec:toc]{[返目錄]}
~
\hyperref[sec:550]{> > > NEXT SERMON > > >}
\newline
\newline
以西結書 18:1-18:4
\newline
\begin{longtable}{cl}
\hline
\hline
章節 & 經文 (和合本修訂版)\\
\hline
18:1 & \begin{tabularx}{0.7\textwidth}{X} 耶和華的話又臨到我，說： \end{tabularx} \\ \\ \relax
18:2 & \begin{tabularx}{0.7\textwidth}{X} 「你們在以色列地何以有這俗語，『父親吃了酸葡萄，兒子牙齒就酸倒』呢？ \end{tabularx} \\ \\ \relax
18:3 & \begin{tabularx}{0.7\textwidth}{X} 主耶和華說：我指著我的永生起誓，你們在以色列必不再引用這俗語。 \end{tabularx} \\ \\ \relax
18:4 & \begin{tabularx}{0.7\textwidth}{X} 看哪，所有的生命都是屬我的；父親的生命怎樣屬我，兒子的生命也照樣屬我；然而犯罪的，他必定死。 \end{tabularx} \\ \\
[1ex]
\hline
\hline
\end{longtable}
今天晚上,我們再次靠主的恩典,來領受祂的恩惠,靠著祂的恩典,來領受祂的恩典,實在是恩上加恩.
\newline
\newline
神的話語真是奇妙!祂的話一出來就有能力,一解開就有亮光；祂的話是我們生活的力量,盼望和喜樂.我們要專心仰望祂,讓祂的話真正進到我們裡面,溶化在我們生命中；叫我們的生命豐富,人生充實,在主前更有榮耀的盼望.
\newline
\newline
本章是極重要的以西結信息,也可說是以西結對以色列民族,信仰,思想,非常重要的貢獻.他提及公義的神,祂的公義是報應的公義；人若犯罪,祂就刑罰；人若回轉,祂就使人有豐盛的人生.這是我們每人在神前應負的責任,這責任是值得我們重視和正視的.本章說明極重要的生命定律.
\newline
\newline
亞當夏娃犯罪之後,神給他們知道罪惡的後果(創二17)從新約我們清楚知道,罪的工價乃是死亡(羅六23).在整本聖經中,已澈底說明神公義的原則；神是公義的神,萬不以有罪為無罪(鴻一2).神雖不喜歡刑罰人；但因祂的公義必須有審判.本章論以西結曾屢次說,神不願意人受刑罰,祂不喜歡人死亡,祂願望我們回頭而存活.可以享受祂的恩惠,在祂面前得著赦罪的恩典,得享天上各樣屬靈的福份.
\newline
\newline
在整部聖經中,讓我們清楚明白神公義的原則；如果說是生命的定律,就無人可以破壞.
\newline
\newline
摩西曾在摩押平原特別忠告以色列人須有所抉擇,若聽從神的話,順服在神面前,澈底悔改就必蒙福；相反地,若悖逆祂,不肯信祂,不肯遵行祂的話,不順從祂命令,就受咒詛.申命記清楚反覆地,說明這個主題,一邊是祝福,一邊是咒詛；以色列百姓在曠野,他們遠望基利心山和以巴路山；摩西吩咐基利心山宣告祝福的話,也在以巴路山宣告咒詛的話.
\newline
\newline
弟兄姊妹!我們人生中常有一連串的選擇與決定,如果選擇決定錯誤,將遭受很大,很多,且很深的損失；如果選擇決定得對,我們將蒙福無盡.舊約的以色列人如何蒙神賜福,同樣也遭受神的咒詛,這都不是偶然的.我一直強調,恩典是有條件的；我們如不符合恩典的條件,就不可能蒙神的恩典.並非神不給,乃因我們不肯好好地領受.
\newline
\newline
在聖經中,我們常發現罪惡的後果極其可怕；因一人犯罪,罪的惡果也要別人承擔.十誡中,神藉摩西告訴以色列人,必須尊重神也必須愛神；神必追討他的罪直到三,四代.在約書亞記,我們看見亞干犯罪,不但他一人被治死,連他家人也被治死.撒下廿一章,掃羅與基遍人所立的約破壞了,基遍人要掃羅的後裔七人,要把他們處死,這七人並沒有犯罪；但是他們卻要擔當掃羅的罪.王上廿一章載亞哈王犯罪,他不但作了很多壞事,且奪取拿伯的葡萄園.先知以利亞責備他,定他的罪,且宣佈他的子孫必承擔他流人血的罪.這些歷史指出,如果上一代犯罪,後代也要承擔.是否神不公平?神是公義的,不單對個人,也與家庭社會民族有關,犯罪的必受連累.有句俗語說:「父親吃了酸葡萄,兒子的牙齒被酸倒.」意即父親犯罪,兒子要承擔.以西結說不是神不公平,最要緊者是著重個人的責任；不是說兒子該擔當父親的罪,而是因父親犯罪,必影響兒女.他整個論調看到這個真理.個人當在神前直接負起責任,無法推諉,否認,掩飾.我們個人,不論直接或間接受別人連累,受家庭罪惡影響,最重要個人須在神前負責任.
\newline
\newline
以西結之所以傳講這信息,是因當時以色列人面臨敗亡,多人被擄外邦；耶路撒冷即將陷落,人心徨徨,前途渺茫,不知所措；他們歸咎祖宗犯罪以致於此.以西結說他們錯了,固然因祖宗犯罪,但每個人應承擔罪的責任；也許已往沒犯罪,現在也沒犯罪；但是將來呢?意即個人所作所為隨時要在神前,直接負其個人責任.神召我們,叫我們有個人向神負的責任,那是無法推諉的.該隱把亞伯殺了.神問該隱:「你的兄弟在哪裡?」神問亞當:「你在哪裡?」「那裡」非指何處,乃指當時的情況,並與神的關係如何?我們在神面前都有屬靈的責任,我們向神負責；神向我們施恩,我們不拒絕；我們領受之後,要珍貴神的恩典,要在恩典中長進,勿在恩典中墮落.該隱回答神說:「我豈是看守我兄弟的麼?」我們應看守兄弟,這是人對人的責任,是道德的.至於屬靈的責任是對神的；或者可以說是宗教另一方面的道德.神創造人,叫人有回應；祂愛我們,我們也接受祂的愛,我們也愛祂.祂施恩典給我們,我們領受了,並要在祂的恩典中長進.祂命令我們,我們應當聽從；祂賜給我們話語,我們要充分得著且實行在生活中.我們須有回應,並且要負責任.向神負的責任,是屬靈的責任.向人負的責任,是道德的責任.
\newline
\newline
有人說,舊約單講整體,集體的思想.但是,到了耶利米和以西結書,卻特別注重個人與神的關係.(耶卅一31)論及神與以色列家和猶大家另立新約,神的律法要寫在每個人的心靈裡,不再寫在石板上；所以個人與個人之間不需互相教導,乃是個人直接與神交往.神赦免人的罪,不是集體的而是個別的.耶利米書特別著重個人向神的責任,以西結書特別注重道德的責任.
\newline
\newline
本章詳述道德方面的責任.個人生活,家庭生活要聖潔；社會生活要互相關懷,對民族要負責任.以色列人要向他的民族國家負責任.
\newline
\newline
剛纔提及神如何施恩給人,我們也看到神的公義,叫人遭受很多苦難和管教；所以以西結清楚地說,父親的罪也要兒子承擔的.每一個人應向神負責,也要向人負責；向神負責是表明自己沒有犯罪,也要擔當；但若向人負責,罪惡的後果對別人就不可能沒有影響.個人所犯的罪,不可能沒有影響到家人或所接觸的人.如果我們了解這是生命的定律,我們就不難明白神為何叫一代代的受刑罰.亞干犯罪,兒子被殺；因神不能不對付這個家庭,恐這家庭留下後遺症,而禍患無窮.所以神忍痛除掉犯罪的家庭.讓我複地說,個人犯罪,別的人雖不擔當,但一定會受到影響.換言之,當我們犯罪時單獨向著神,這還較為簡單；因神是施恩的神,祂要赦免；神雖赦免,但很多罪惡的後果,會影響到我們的家人,教會,社會,民族,我們怎能擔當?我們多麼需要神的憐憫!雖然神會幫助我們,但我們必須要謹慎!
\newline
\newline
例如,大衛犯罪,他所犯的罪不堪設想!他是個合神心意的人,竟軟弱墮落到此地步!誰能說我站立得住呢?自以站立得住的,務要謹慎,免致跌倒!大衛跌倒後在神面前切實澈底悔改.神是施恩的神,祂的怒氣不過是轉眼之間；但祂的恩典,乃是一生之久.神實在赦免了大衛的罪,但罪惡的結果仍然存在.好像一顆釘子釘在木頭上,釘子拔出來後,痕跡仍留存.大衛雖蒙神赦免,但家庭中仍遭患難,有背叛,亂倫,刀劍,死亡的事發生.為甚麼神使他的家受那麼多苦難呢?難道神不饒恕他嗎?神實在饒恕了他,但罪惡的果他必須吃.我們在神面前必須謹慎!我們多麼需要神的光照,使我們可以看見自己的污穢,罪惡,虧欠,不潔；讓主的寶血在我們身上生效,纔不致繼續墮落,不致因罪影響連累了別人.教會也是如此,何等需要彼此認罪,讓主的寶血隨時潔淨我們；因為我們擔當不起那麼大的責任.
\newline
\newline
現在讓我們細思個人責任.個人須在神前負責任,任何罪都不可推諉,任何軟弱都沒有解釋的理由.世人很輕易怪責別人,當失敗時就推諉說是家庭的錯社會的錯!當成功時,一切歸功於自己的纔幹本事努力這些都是罪人的心態,但我們屬神的人,須在神前負完全責任,不可怪責別人,要完全仰望神；有時環境不夠理想,家庭背景不夠圓滿,所受的教育不夠完善.我們所處的是病態社會,但當主的恩典拯救了我,勝過了所有的一切；因祂的生命在我裡面,祂的寶血隨時潔淨我,我成了全新的人心靈更新,不再受環境,家庭,社會,學校,各方面的影響；因我們在基督裡就是新造的人,舊事已過,都變成新的了.既是如此,我們的觀念,態度就應該要更新.神要我們負完全的責任,我們每天須在神前省察,要親近祂.為善必須下工夫,付上時間的代價,多多親近神.信心不是一種觀念,而是一種關係──乃是與神的關係,這重要的事.我們無論作甚麼,必須在神面前作；因神不是單看我們的外表,也看我們的裡面.祂知道我們的內心,知道我們的動機.神是輕慢不得的,我們種的是甚麼,收的也是甚麼.神公義的報應,是我們生命的定律.我們每個人當直接向神負責,無法向神作任何推諉；不要逃避神,要面對神,要預備迎見神.因為我們終必要站在基督審判台前,特別我們信主的人,所以我們務須謹慎!我們是屬於主的,我們不在短暫的時間裡,乃活在永生裡面.有人以為死後上天堂纔得永生,其實不然；我們何時蒙恩得救,我們的永生便開始了；所以無論作甚麼,價值是存到永遠的.我們雖在地上,也學習過天上的生活；雖在短暫裡生活,但也影響直到永遠.弟兄姊妹!我們當正視重視這現實的生活,個人須向神負自己,負家庭的責任.我們要孝敬父母,可惜今日社會,孝道不再重視；但「教父母」,正列於十誡中之第五誡.有人說,應把十誡分為對神和對人兩部分.我曾說過,第五條誡是對神不是對人的.對父母不是對人；而是對神；但並非我們把父母當作神,不過在神的計劃中,父母是神在地上的代表.當然,如果我們的父母尚未信主,怎能代表神呢?但作兒女的仍然要尊敬他們.我們若在地上代表神教養兒女,這是多重大的責任；是在神面前向祂負上教養兒女的責任.如果我們失敗了,他們就會擔當我們的失敗,如果我們犯了罪,就會叫他們來擔當我們的罪惡；所以我們的失敗和罪惡,都會影響他們.有些做兒女的婚姻失敗了,其實是因父母的婚姻不美滿.我們有時因自己軟弱,無法帶領丈夫,妻子,兒女信主,還錯怪家人太硬心.在某次聚會中,有一位姊妹說:「我實在無法帶領家人信主,因為我的家庭亂七八糟.我來到教會心纔平安,把一切都忘記了,也不再想到家,我歡喜到主面前.」表面看,她很屬靈；其實,她是逃避責任；對家人不負責,逃避責任.主耶穌曾說:你獻祭之前,要先與你弟兄和好.華人教會對這方面,很不理想,因為我們信了主,對家庭卻沒有完全負起責任.小家庭不行,大家庭(教會)怎能興旺呢?教會(大家庭)不興旺,小家庭怎能有幸福呢?
\newline
\newline
很多基督徒很著重個人主義,以為個人蒙恩得救就夠了；至於有沒有參與教會聚會,有沒有在教會事奉,都無關重要.神設立教會,是要叫蒙恩得救的人結合起來,成為基督的身體,合力建立基督的身體.所以個人必須向神負責,我們不可忽略個人對團體的責任,和對教會的責任.一個亞干能使整個以色列會眾失敗,被期教會,亦因一對夫婦的奉獻不誠實,而使整個教會蒙受虧損.請記得!我們不只是個人,乃是整體,由許多個人結合起來的整體而成為教會,成為基督的身體.教會是基督的道成肉身,很具體很真實地將神表達出來.教會是很具體,很實際正確地叫人看到的.我們怎能失敗呢?當我們失敗時,福音就受到擱阻,神的真理受阻礙,我們怎能擔當起呢?
\newline
\newline
神是公義的神.人種的是甚麼,收的也是甚麼；因我們都要站在基督的審判台前.求主憐憫!求主幫助我們知道,神不喜悅我們受苦,滅亡,受管教,犯罪又影響別人；神所喜悅的是我們回頭,認罪,親近祂,在個人生活中榮耀祂.在家庭方面要負完全的責任,在教會中要忠心.以色列人要向他的民族負責任,華人教會也要向國家負責任.香港教會,許多人都沒有民族意識,我認為這並非神的旨意.在海外我常對華人弟兄姊妹說:第一我們成為神的兒女,我們蒙恩得救一定是神的旨意,因神拯救了我們.第二,你我都是華人,一定是神的旨意.我們應有民族意識.神叫我們做華人,神使全世界眾多華人；幾千年來,神保守我們這個民族；相信神的旨意要我們負起責任,正如以色列民族負起他們的責任一樣.我們要愛同胞,如同保羅愛同胞一樣,他們是骨肉之親,保羅為了福音的緣故情願被咒詛,甚至與基督分離也願意,他付上了這個代價.各位有沒有這樣的心志?有沒有愛同胞的心呢?
\newline
\newline
香港的環境,人與人之間,似乎漠不相關；但我們是華人是屬神的人,在神面前,我們應當向華人負責任,把福音帶給他們.要保持神要我們保持的文化傳統,別以為文化是世俗東西與屬靈無關.是的,文化部份是錯誤的；但部份乃是正確的.今日基督徒受了西方文化的影響,西方文化有很多世俗,阻礙我們對神的信心.
\newline
\newline
整個世界在混濁潮流中.難道我們隨波逐流嗎?不,因我們是神的兒子,是中華兒女；我們應當作中流砥柱.在神面前,我們負個人的責任,負家庭,教會,社會,民族的責任.這是神給我們之生命的定律.
\newline
\newline


\newpage

\section{第54屆~港九培靈會~唐佑之~第六講}
\label{sec:550}
\textbf{守望的吶喊}
\newline
\newline
連結: \href{https://www.hkbibleconference.org/session-message/view/550}{\texttt{https://www.hkbibleconference.org/session-message/view/550}}
\newline
\newline
\hyperref[sec:549]{< < < PREV SERMON < < <}
~
\hyperref[sec:toc]{[返目錄]}
~
\hyperref[sec:551]{> > > NEXT SERMON > > >}
\newline
\newline
以西結書 33:1-33:9
\newline
\begin{longtable}{cl}
\hline
\hline
章節 & 經文 (和合本修訂版)\\
\hline
33:1 & \begin{tabularx}{0.7\textwidth}{X} 耶和華的話臨到我，說： \end{tabularx} \\ \\ \relax
33:2 & \begin{tabularx}{0.7\textwidth}{X} 「人子啊，你要吩咐本國的百姓，對他們說：我使刀劍臨到哪一國，哪一國的百姓從他們中間選立一人，作為守望者。 \end{tabularx} \\ \\ \relax
33:3 & \begin{tabularx}{0.7\textwidth}{X} 守望者見刀劍臨到那地，若吹角警戒百姓， \end{tabularx} \\ \\ \relax
33:4 & \begin{tabularx}{0.7\textwidth}{X} 有人聽見角聲卻不受警戒，刀劍來除滅了他，這人的血必歸到自己頭上。 \end{tabularx} \\ \\ \relax
33:5 & \begin{tabularx}{0.7\textwidth}{X} 他聽見角聲，不受警戒，他的血必歸到自己身上；他若受警戒，就救了自己的命。 \end{tabularx} \\ \\ \relax
33:6 & \begin{tabularx}{0.7\textwidth}{X} 倘若守望者見刀劍臨到，卻不吹角，以致百姓未受警戒，刀劍來殺了他們中間的一個人，這人雖然因自己的罪孽而死，我卻要從守望者的手裡討他的血債。 \end{tabularx} \\ \\ \relax
33:7 & \begin{tabularx}{0.7\textwidth}{X} 「人子啊，我照樣立你作以色列家的守望者；你要聽我口中的話，替我警戒他們。 \end{tabularx} \\ \\ \relax
33:8 & \begin{tabularx}{0.7\textwidth}{X} 我對惡人說：『惡人哪，你必要死！』你若不開口警戒惡人，使他離開所行的道，這惡人必因自己的罪孽而死，我卻要從你手裡討他的血債。 \end{tabularx} \\ \\ \relax
33:9 & \begin{tabularx}{0.7\textwidth}{X} 但是你，你若警戒惡人，叫他離棄所行的道，他仍不轉離，他必因自己的罪孽而死，你卻救了自己的命。」 \end{tabularx} \\ \\
[1ex]
\hline
\hline
\end{longtable}
耶和華的話臨到先知以西結,稱他為「人子」意即脆弱的人.神對他說:「我要用你作以色列家守望的人.」在第三章,神要他吃書卷,要他受命把祂的信息,傳出來作以色列守望的人.這裡重複第三章的內容.
\newline
\newline
神不但對以西結這樣說,他對歷代的工人和教會這樣說.教會在這時代,要聽見神這樣的命令.我們不過是人,是脆弱軟弱的,不能作偉大的工作.
\newline
\newline
主耶穌基督是人子,當祂作人子時,他卻不是脆弱軟弱的；他成為人,但他是最完全的人.祂已為我們留下榜樣作了表則,教我們隨祂足蹤行.祂曾作以色列家的守望者,作全世界人類的守望者.你我也當作守望者,正如主耶穌和以西結一樣.
\newline
\newline
在舊約聖經所見,守望者的形象有兩種.一種是專責在戰爭將爆發,危急時刻高守遠望,儆醒等候.一發覺敵人,便大聲疾呼,向百姓宣佈,警誡他們危險來到；叫人不但逃避危險,還要應付戰爭.不但防守,還要攻打；要及時回頭,應付危急時期.神所要的就是這樣的守望者,每個時代都需要守望者.教會在這時代,更需要守望工作；因為面臨最危險的時期,所以我們必須在神面前等候.另一種守望者的形象非在戰時,乃在平時.聖地耶路撒冷的錫安山上,每晚祭司要作守望的工作；在漫長寂靜的黑夜中,按班次站在錫安山上守望,等候天亮.在黎明之前最黑暗的時刻,就是人們昏昏入睡的時刻；守夜者更加要儆醒,手持號角,挺胸昂首聚精神會伐步.面向東方,等待曙光出現,大聲吹角,叫醒耶路撒冷的居民到聖殿去獻早祭.這是很重要的敬拜,這也是教會應作的守望工作.我們今天在屬靈的高山上,要向人們宣告黎明的來臨.作黑夜中的守望者,每晚必須如此.本章所指的是非常的時期,而祭司的守望卻是在平常的日子.
\newline
\newline
教會在這時代,無論平常時期,或非常時期,均須作守望者.
\newline
\newline
先知以賽亞聽見從巴勒斯坦南方發出聲音:「守望的呀,夜裡如何?」(賽廿一11-12)聲調非常焦急,沉悶,懼怕,困惑,迷茫地問:「黑夜何時纔能過去?」黑夜象徵苦難.當人們在苦難中,就會發問苦難何時過去呢?
\newline
\newline
整個歷史的記載,使我們看到個漫長的黑夜；甚至我們想到在伊甸東園的黃昏,當耶和華在園中行走,天起涼風之時,人已犯了罪,逃避了神的面.從那黃昏開始,人的歷史就進到黑夜裡；人就在罪惡痛苦無知的黑夜裡,不見天明,沒有早晨,一直在黑夜中.
\newline
\newline
當以色列人出埃及之前,他們希望到曠野敬拜耶和華；但埃及王不許.大能的神降災給法老王,叫他的心回轉；但是一次又一次法老的心剛硬.到了第九次黑暗之災,(埃及人拜太陽神)埃及人之神失敗了.我們的耶和華仍是全能的神,除耶和華以外沒有別的神.那一次的黑暗災殃,使全埃及陷於黑暗中；惟有以色列人的住處,纔有亮光.因為耶和華要向全世界宣告,黑夜的消息.
\newline
\newline
主耶穌來到世界,每逢聖誕節,有美麗詩意的描寫.但是,當主耶穌降生那時,世界充滿著無邊的黑暗,人類在極大痛苦之中；以色列在異族統治的鐵蹄下,政治上沒有自由,經濟困難,社會不安寧.從耶路撒冷到耶利哥,沿途有盜賊出沒,既危險又可怕.當人類在極其失望,在苦難之黑夜中,那希望的星辰出現,神的兒子來到世界,光來到世界.約翰福音告訴我們；光來到世界,但世人不接受光；因他們行惡,拒絕光,他們情願在黑暗中.他們把耶穌釘在十字架上,當時從午正到申初,遍地黑暗；神再次宣佈黑暗的消息.雖然耶穌基督是神的兒子；但當祂為你我,為全世界人類擔當罪惡之時；父神也掩面不顧祂.那時耶穌說:「我的神阿!我的神阿!為甚麼離棄我?」主耶穌死了,埋葬了,復活了；那天早晨,就是人類在歷史中,最光明的早晨.黎明把生命的意義,重新解釋了.當耶穌升天之後,教會建立起來；信靠祂的人,就被稱為光明之子,白晝之子.教會在神的光中,跟從祂的人,就不再在黑暗中走,乃要得著生命之光.可憐今日教會的光景仍在黑夜裡.
\newline
\newline
啟示錄記載,使徒約翰被放逐到拔摩海島,在那裡有一主日；他被聖靈感動,看見天開了,人子從天而降.他說:「看哪!祂駕雲降臨,眾目要看見祂；連刺祂的人也要看見祂,地上的萬族都要因祂哀哭.」耶穌基督一定再來,要結束這世界的秩序.當人子降臨之時,祂的面貌如日頭放光；這位榮耀的人子,行走金燈台中間.燈台是夜間用的,要照亮黑暗中的人.耶穌基督設立晚餐時是在晚上,祂要教會守晚餐；教會不吃早餐,午餐,只吃晚餐,乃因要常提醒我們是在晚上.我們必須成為黑夜的守望者,我們要等待天亮,要儆醒謹慎帶著盼望等待,主必再來,決不遲延.榮耀的黎明就要來到,教會在這時代必須作守夜工作,等待主再來.彼得說:「萬物的結局近了.」千多年已經過去,現在不是更近了嗎?彼得教我們要謹慎持守,在神前等候.
\newline
\newline
我們必須謹慎,否則容易受世俗虛謊所影響.不大愛主的人易於隨從世俗.約翰說:「我們若與世俗為友,就與神為敵了.人若愛世界,愛父的心就不在他裡面.」(約壹二15)多麼危險啊!我們的環境充滿世俗成份,在濁的世俗潮流中；我們切勿隨波逐流,而應當作中流砥柱.為了主的道,我們萬勿隨從世俗,我們當謹慎持守,作忠心的守望者.較為熱心的弟兄姊妹,很容易受虛謊影響.今天有很多假先知,包括那些異端邪說,傍門左道的,也包括那些信仰純正的.信仰純正怎會虛謊呢?照本章聖經說,若非出於神直接啟示的話,照著自己的意思來說,就是虛謊.許多人歡喜聽末世預言,但我們應有末世的意識,來聽那許多的預測；我很慚愧地告訴各位,同時也警告各位,不要隨便聽那些末世的預言.因為很多人所講的是錯誤的解釋.耶穌說:主來的日子自己都不知道,惟有父知道.以往在許多年代,直到如今還沒有人知道主耶穌何時再來；我們相信主必再來,一切預言必定應驗；但人若照己意講解,那就是假先知假教訓.雖然我們所聽的,有在原則上是對的；但我們必須有屬靈的辨別力,不隨便聽信接受.要好好的在神面前追求,在教會聽神的工人為我們分解的真理,這樣纔可避免聽信虛謊.
\newline
\newline
我們在黑夜中危險中作守望者,不但要謹慎持守,也要清潔,專一愛神.我們要聖潔,否則不能作聖工,不能見我們的神；不聖潔就沒有見證,不會明白聖經的話語,就沒有屬靈的辨別力.
\newline
\newline
在黑夜中更須儆醒,因為夜間精神疲乏,昏昏欲睡.彼得說:「要儆醒禱告.」耶穌在馬太廿四章,特別論及祂再來的情形.雖不詳述細節,但叫我們務要儆醒.廿五章十個童女的比喻,童女之中五個是聰明的,五個是愚拙的；祂並沒有說五個是良善的,五個是罪惡的.所謂聰明,乃因他們的燈預備了油,而愚拙的因為沒有預備油.表面上她們都是在迎接新郎,但實際上卻有很大的差別.耶穌說:「新郎遲延未到,她們都睡著了.」聰明的童女怎能睡著呢?今天教會許多所謂聰明的童女正在睡覺,值此危急之時怎能睡覺呢?現在應是睡醒的時候.耶穌說:「你們要儆醒!」彼得說:「要儆醒禱告.」與神有親密的交往.我們的禱告,不是只想自己的需要,也非僅為教會的需要；我們當將禱告的範圍擴大,為萬人代求.我們當為所居住的城市禱告,為世界和平禱告；正如亞伯拉罕為所多瑪,蛾摩拉二城在神前祈求一樣.若城中有五十個義人,求神不要毀滅；甚至減剩十個義人,神也答應不毀滅那城.可惜全城連十個義人都沒有.神差天使搭救羅得和他的家屬,他的女婿還以為是戲言不信；羅得就帶了妻子和兩個女兒逃出.他的妻子回頭一看,就變成一根鹽柱；剩下羅得和兩女兒,而兩個女兒居然犯起亂倫的罪來.這豈不是歷史的悲劇嗎?我們多麼需要為所居的城市,為世界的平安禱告!我們不能像先知約拿,風浪大作是為了他一人的緣故,因他違背了神,不作傳福音的工作.照樣,因你我不順從神,神的兒女工人不順從神,以致世界紛亂.當遍地烽火之時,我們怎能像約拿在船底睡覺呢?非尼基的水手禱告他的假神,雖然沒有功效；但他有禱告,可惜神的僕人約拿卻不禱告神,船主叫醒約拿,說:「你這沉睡的人快起來禱告你的神吧!」這世界越來越危險,原因和理由很多,但最主要的是因為我們不禱告.我們多麼需要神的憐憫和幫助,叫我們儆醒禱告免得入了迷惑.世界安危在我們身上,我們是守望者如以西結一樣.我們要吶喊,要警誡；因為危險來到,還有人不知道不領會!先知能見人所不見,明白人所不明的；教會是先知的團體,有神的話語.我們聽明白,看清楚,有先見之明有神的話語,怎能不出聲呢?我們當吶喊,當大聲疾呼,要警誡!在此非常時期守夜,必須儆醒禱告,謹慎持守,吶喊,警誡.危險緊急嚴重的時期,不容等候,遲延,猶豫；我們必須吶喊.在這社會,有沒有聽見福音的吶喊聲呢?在這世界裡,教會有沒有發出警誡的聲音呢?若有,而人不肯悔改回轉,守望者就沒有責任；否則守望者就要負起重大而可怕的責任.
\newline
\newline
教會在這時代必須作守望者,要吶喊警誡；因為處境危殆,刻不容緩.
\newline
\newline
我們要歌唱讚美.「神使人在夜間歌唱」(伯卅五10)夜間,苦難,危險,困惑中還能歌唱嗎?惟神使人在夜間歌唱.今天我們聽見社會,世界那些嘈雜爭鬧,淫亂的聲音,給人帶來痛苦煩惱,很多人遭難,正在嘆息哭泣呻吟；聽不見聖潔的歌聲；難道我們無動於衷嗎?難道你毫不介意,毫無屬靈的感覺嗎?教會在這時代,應有敏銳屬靈的感覺,不可充耳不聞.我們要歌唱讚美,要從教會裡向外看；大家庭要歌唱,小家庭也要歌唱；我們四所聽見的都是電視,打牌的聲音；我們有責任把歌聲傳出來,從教會裡唱到教會外,從晚間唱到早晨,甚至從地上唱到天上.將來我們在天上還要歌唱,所以我們作守望工作,就是在地上學習天上的生活.保羅和西拉在監獄中歌唱,眾囚犯側耳而聽,因此幫助腓立比管獄的歸向神,成為腓立比教會最基本的會友.保羅對腓立比教會說:「你們要喜樂,我再說,你們要喜樂.」因為他實在有喜樂的經驗.
\newline
\newline
我們在守夜遇危險時,我們一面工作一面爭戰.要建立神的工作,建立神的教會,建立靈宮.爭戰是必然的,因地上有惡者屢加侵犯；我們須有勇氣,我們不但防守,還要進攻.工作和爭戰是分不開的,有時我們在教會中有工作但是沒有爭戰；以致仇敵侵襲,我們便失敗了.
\newline
\newline
教會一旦發生問題,對一般不熱心者似乎無關重要,但對那些最熱心的弟兄姊妹,卻發生問題；因為他們只有工作沒有爭戰,結果被魔鬼打敗.本應合力對付魔鬼,可是為了工作卻造成對立,彼上成為仇敵,這是多大的悲劇!我們應學習尼希米建造城牆時當眾宣佈,各人應一手拿工具,一手持兵器,邊工作邊爭戰.
\newline
\newline
教會在這時代,是爭戰的時代,是爭戰的教會,我們要重視萬世的戰爭,在每時代,教會不能安逸安穩,要隨時備戰,人人成為基督的精兵.教會在這時代,是爭戰時代的教會.
\newline
\newline
以西結是神的守望者,很多神的工人是神的守望者；很多神的兒女,也是神的守望者.
\newline
\newline
神要在我們中間尋找守望者,誰是神的守望者?神說:「要立你作守望者.」你的回應如何?
\newline
\newline
但願我們如以賽亞說:「我在這裡,請差遣我!」
\newline
\newline


\newpage

\section{第54屆~港九培靈會~唐佑之~第七講}
\label{sec:551}
\textbf{牧者的哀歌}
\newline
\newline
連結: \href{https://www.hkbibleconference.org/session-message/view/551}{\texttt{https://www.hkbibleconference.org/session-message/view/551}}
\newline
\newline
\hyperref[sec:550]{< < < PREV SERMON < < <}
~
\hyperref[sec:toc]{[返目錄]}
~
\hyperref[sec:552]{> > > NEXT SERMON > > >}
\newline
\newline
以西結書 34:1-34:6
\newline
\begin{longtable}{cl}
\hline
\hline
章節 & 經文 (和合本修訂版)\\
\hline
34:1 & \begin{tabularx}{0.7\textwidth}{X} 耶和華的話臨到我，說： \end{tabularx} \\ \\ \relax
34:2 & \begin{tabularx}{0.7\textwidth}{X} 「人子啊，你要向以色列的牧人說預言，對他們說，主耶和華如此說：禍哉！以色列的牧人只知牧養自己。牧人豈不當牧養群羊嗎？ \end{tabularx} \\ \\ \relax
34:3 & \begin{tabularx}{0.7\textwidth}{X} 你們吃肥油、穿羊毛、宰殺肥羊，卻不牧養群羊。 \end{tabularx} \\ \\ \relax
34:4 & \begin{tabularx}{0.7\textwidth}{X} 瘦弱的，你們不調養；有病的，你們不醫治；受傷的，你們未包紮；被逐的，你們不去領回；失喪的，你們不尋找；卻用暴力嚴嚴地轄制牠們。 \end{tabularx} \\ \\ \relax
34:5 & \begin{tabularx}{0.7\textwidth}{X} 牠們因無牧人就分散；既分散，就成為一切野獸的食物。 \end{tabularx} \\ \\ \relax
34:6 & \begin{tabularx}{0.7\textwidth}{X} 我的羊流落眾山之間和各高岡上，分散在全地，無人去尋，無人去找。 \end{tabularx} \\ \\
[1ex]
\hline
\hline
\end{longtable}
感謝主!耶和華是我們的牧者.
\newline
\newline
以西結接受神的命令,傳信息的內容是要責備那個牧人,因以色列的牧人失敗了,以致羊群四散,牧人沒盡其責.
\newline
\newline
人雖失敗,但神卻不失敗；耶和華要親自為牧人,耶和華的旨意要祂的兒子成為好牧人；且希望屬祂的人,跟隨好牧人的腳蹤行.
\newline
\newline
本章是一首牧人的哀歌.耶和華是牧人,為以色列的悲哀；為許多屬祂的人悲哀,為各時代的教會悲哀.求神憐憫!讓教會在這時代,完全轉向我們的牧者耶和華；讓祂親自來牧養我們；而且讓教會在這時代,能看到我們的需要.
\newline
\newline
舊約時代的牧人,不像今天教會的牧者,不是專責牧養工作.以色列人有君王牧養他們,所以這裡的牧人,是指以色列的君王；我們若略知歷史背景,就有助我們明白.當摩西領以色列人出埃及到曠野,一般觀念都認為摩西是他們的領袖；不錯,摩西在屬靈方面領導他們,但在行政上,卻非他們的領袖,也非他們的王.以色列人的王是耶和華,其他的人是耶和華的僕人；都是服事祂的,也是服事眾人的.
\newline
\newline
在耶和華為王的整個神權觀念裡,神權政治,是以色列民族的信仰.若在這方面偏差,就整個都偏差；若這方面對了,就完全正確了.所以當以色列民進入迦南美地之後,他們從半遊牧生活而進入到農業社會,他們是無法適應的,於是神就興起士師來照管他們.士師也非君王,不過是耶和華的僕人作神的工；耶和華纔是他們真正的王.可惜以色列民不尊耶和華為王!
\newline
\newline
在士師記有句重複的話:「當時以色列中沒有王,各人任意而行.」(士十七6,廿一25)那些士師也不能作耶和華的工人,故當時非常混亂.以色列人心中切望有王,到了撒母耳的時候,百姓求撒母耳為他們立王；撒母耳心中很難過,覺得他們的態度不對.我們讀這段聖經時,很容易會發生誤會,以為以色列人不可有君王；其實不然,以色列人需要君王；但這君王必須為耶和華的僕人,因為最主要的是耶和華作王.並非他們求立王不對,乃是他們的態度不對,他們不尊耶和華為王.
\newline
\newline
我們想到教會需要僕人的領袖,我們需要牧者,需要領袖；但最重要的務必尊重我們的牧長耶和華,尊祂為王.若在這方面失敗,問題就發生,正像當時以色列人的失敗一樣.當他們有了王之後,雖有先知指導那些君王,但有的君王很差.
\newline
\newline
神是信實的神,讓大衛的王位得堅立.(記於撒下七16)神雖與大衛立約,可惜他死後,以色列國南北分裂了,以致失去了合一的精神.北國完全偏離了神,而南國亦軟弱無能.以色列國王朝失敗,因他們不尊重耶和華為王.
\newline
\newline
我們可以把以色列民族,失敗的歷史經驗,應用於歷代的教會,特別是應用於今日的教會.這是我們讀舊約主要之目的.從以色列民族歷史信仰的經驗中,可以看到教會的分裂,而應用在我們各人的身上.耶和華是王,耶和華作王之時,纔有真正的平安,有真正的秩序,然後纔有和平的生活.
\newline
\newline
在本章中,先知得著神的命令,要攻擊責備那些牧人,那些君王領袖們；他們自私,沒有關懷別人,沒有真正尊耶和華為王的態度.
\newline
\newline
請看今的教會,我們有否尊耶和華為我們的王,有沒有尊基督為教會的元首,有沒有尊重神的全能呢?
\newline
\newline
以色列的牧者失敗了,他們自私,不尊重耶和華為王；只注意自己,不關懷別人.全民族敗亡的原因很多,最主要的是因為工人軟弱.
\newline
\newline
在座有很多神的工人,我們當在神前坦白,謙卑承認我們沒有作好的牧人,致令教會大受損失!瘦弱的沒有養壯,有病的沒有醫治,受傷的沒有纏裹,隸屬的沒有領回,失喪的沒有尋找；作為教會的教牧同工們!求主憐憫!我們多麼需要在神面前認罪悔改!因我們的失敗,軟弱,自私；讓好多弟兄姊妹也失敗軟弱自私；教會不興旺,以致不能復興,不強壯,也不剛強.靠著主的大能大力,我們工人須切實在神面前認罪；弟兄姊妹也要在神前認罪.有時牧人失敗,乃因信徒沒有讓牧者好好地事奉神；有時信徒軟弱,因牧人沒有好好地帶領.神為今日教會哀哭,我們有沒有聽到牧者的哀歌在教會裡?神的心焦急,因祂的工人失敗,羊群分散.受傷的不得醫治,軟弱的不得幫助,是誰之責?我們在神面前當加以省察.
\newline
\newline
在這時代環境中,牧人受了重大的壓力,飽受困難,這些都不成理由.我們若真正作神要我們作的牧養工作,必須在神面前忠心.很多弟兄姊妹,在教會中感到失望；因牧者不負責任,有時確實如此；但有時信徒對牧者的認識不夠,觀念不正確,要求不合理；不尊重牧者,沒有讓牧者有機會發展工作.這是值得我們深切思想的,常聽見有人投訴說,他教會牧者不行.這樣到底對不對?我們當省察,到底是哪方面不行呢?你有沒有為他禱告,有沒有補充他作不到之處.到底你是仰望他的牧養,還是仰望主的牧養呢?這裡聖經說:王要親自牧養我們.有時我們以為牧者無助於我們,我們就失望；我們沒有仰望耶和華作我們的牧者.有時我們沒有好好地為我們的牧人禱告,沒有尊重他,以致他失敗了.
\newline
\newline
有一次我在韓國,看見那裡教會的光景,信徒都非常尊重傳道人,那些傳道人也不一定顯得完全；但信徒們因尊重耶和華,因主的緣故,他們尊重神的僕人.我發現,信徒若尊重牧人,牧人不能不自重.這是華人教會少有的,所以牧人失敗,很可能是因信徒們不尊重他.我們有否尊重神呢?若有,我們必須尊重神的工人；並非尊重他的個人,乃是尊重他的職份.當他被尊重,他就沒法馬虎隨意自私,他曉得必須盡力而為.
\newline
\newline
掃羅仇視大衛,大衛有幾次機會可殺害他,但是卻沒這樣做；並非掃羅人格高尚使大衛尊重他,大衛尊重他乃因他是耶和華的受膏者.大衛不是尊重人而是尊重神,所以也尊重神的僕人,尊重耶和華的受膏者.
\newline
\newline
讓我再說,我們首先要尊重耶和華；因祂是我們的牧者,我們的牧長；因重祂,我們也尊重教會的牧者.若教會牧者不如理想,我們要以禱告來幫助他,而且我們當作他所忽略的；尋找失喪的,我們也有責任,失散的我們也當領回.我們應作醫治和纏裹的工作,靠著主的恩典,讓瘦弱的得以強壯.關懷別人,不僅是牧者之責,也是我們每個人所當作的；若得不到教會牧者之助,耶和華是我們的牧者,祂必親自帶領,我們必須仰望祂.
\newline
\newline
有人說,牧師喜歡強調權柄,我覺得這是不應該的；因為大家都是神的工人,每個工人應當謙卑在神前服事；若神的工人教會的牧者強調自己的權柄,恐有偏差之處.另一方面,雖然他不該強調權柄；但我們當尊重他三方面的權柄:
\newline
\newline
(一)蒙召的權柄,他是蒙神選召的；若非神選召的,當然就沒這權柄；雖在教會工作,遲早要被挪開.
\newline
\newline
(二)因受過特殊訓練,所以他有權柄.教會中雖有長執,亦有許多教會領袖們；但對於屬靈方面卻不甚了解.至於蒙神選召者,他們曾經受過神學訓練,長期造就,屬靈的事理應明白更清楚.
\newline
\newline
(三)事奉的經驗,生命的經歷,或者你說:「我們的傳道人不屬靈,經歷不行.」那就他們應當自己向神負責,因為一個真正事奉神的人,應有生命的經歷,事奉的經驗.
\newline
\newline
一個神的工人,固然不應強調權威,但是我們卻當尊重他的權威.這是教會的次序,有了次序纔有和諧.有時教會信徒的權力太大,甚至作王,要管理傳道人；他們忘記了耶和華作王,牧師和其他的,不過是僕人.信徒必須尊耶和華為大,尊重神也必定尊重神的工人；神的工人既被尊重,若他自暴自棄又自私,神怎能容許他們呢?如果他們不悔改,終必被神棄絕,所以不必信徒來對付,因我們的神有絕對的權柄.最要緊的,我們須彼此認罪.傳道人有很多的軟弱,不潔,虧欠；不但當向神認罪,也當向信徒認罪.信徒也當向傳道人認罪,這樣纔能看見聖靈的工作,看見復興的異象.
\newline
\newline
有一次,某教會請我去培靈會講道.會前牧師告訴我,那是他在該堂最後的一週,因他已提出辭職.弟兄姊妹為了對付牧師,而安排數日聚會；特別有研討會,提出許多問題,讓講員作出答覆；然後根據作答的話,來對付牧師,當面指責他.我聽了頗為難,到時我好像個審判官；但是我又不能不照聖經的話講.我惟懇切禱告仰望神,使我有智慧誠實不遷就任何人.那天的聚會人數眾多,真是來勢洶洶；可是,神作了極奇妙的工作；那班人完全沒有責備牧師.翌日,我在主日講道完畢,牧師起來講話,當眾承認自己的軟弱虧欠,流淚認罪.那時弟兄姊妹們並不因此沾沾自喜；他們也都起來認罪.那場合實在太感動人!當他們彼此認罪後,教會就因此而特別蒙福.
\newline
\newline
教會須要彼此認罪,我們都有很多的軟弱；工人也有很多的軟弱.有時因工作太忙,缺少顧及弟兄姊妹的需要；弟兄姊妹也忘了顧及工人的需要.有時工人只注意弟兄姊妹的需要,顧不到自己的需要；工人有很多需要,工人的家庭有很多需要；但弟兄姊妹沒有關懷,我們覺得有很多的虧欠.求主憐憫!求主的聖靈感動我們的心,讓我們深深感到必須認罪；不但向神認罪,而且要彼此認罪,惟有這樣,教會纔能得復興,神的工作得發展.
\newline
\newline
當我們讀到此章牧人哀歌時,好像看見神含淚為我們嘆息說:「我要親自牧養我的羊群.」牧者失敗了,教會工人失敗了,弟兄姊妹失敗了；但神沒有失敗,神是我們的牧者,要親自看顧我們.神說:「我必在美好的草場牧養他們,他們的圈必在以色列山肥美的草場吃草.」(14)正如詩篇廿三篇「躺臥在青草地上,領我在可安歇的水邊.」羊若飢餓不吃飽是不躺臥的.
\newline
\newline
感謝主!耶和華是我們的牧者,祂叫我們飽足,叫我們沒有驚駭,叫我們有平安,叫我們在溪水旁.我們一定要仰望神,尊重祂是我們的王是我們的牧者.這樣,不但我們的需要得到滿足,也會顧及別人的需要.若我們單仰望教會的牧者,我們將會失望不得滿足.我們當彼此勉勵,一同在神面前受祂牧養.
\newline
\newline
第16節說:「失喪的我必尋找,被逐的我必領回,受傷的我必纏裹,有病的,我必醫治；只是肥的壯的,我必除滅.」上述最後的一句話,我們實在不明白.在手抄本,特別是舊約繙的希臘文,乃說肥壯的,還要繼續看顧他們.當我們生命健壯時有兩種可能,一是繼續需要神保守我們；另一可能是當我們生命稍為堅強,就驕傲自大.18節說:「用蹄踐踏草地......用蹄攪渾清水.」強壯力量很大.這是今天教會的問題.軟弱的須剛強長進.有的人長進之後更加謙卑,更願意服事；但是有的人稍為長進,就以領袖自居,自以為了不起.表面上好像幫助人,其實就如肥壯的羊,把草地踏踐,把溪水攪成混濁.我曾提過,教會有問題,不是出於那些軟弱,不熱心的信徒；因根本他們不參與,問題乃出於熱心事奉者.在教會有職份有參與者,義工們長執們；按理他們生命豈非長進嗎?可是肥壯的,反而踐踏草地,攪濁清水.求主憐憫保守我們!讓我們真正知道謹慎,不然我們易於失敗.
\newline
\newline
有的人願意事奉,但卻要照他的意思作.我們別選擇來工作,當選擇作神的工人,作神的僕人.我個人最反對領袖這名詞,唯有主是我們的領袖.雖然牧師可帶領我們,但牧師也不是領袖；長執以及所有負責的弟兄姊妹,都不是領袖.我們都是工人,都應有工人的心態.工人沒有自己的選擇,沒有自己的意思；只能照主的意思,工人是要服事的,不是領導的.那麼,牧師長執及負責的弟兄姊妹,可否領導?我想可以,不過途徑方面當注意,我們是藉著服事來領導；正如主耶穌基督一樣,祂並沒有對門徒說你們要服事；乃是祂親身作服事的榜樣,為門徒洗腳.並說:你們現在不明白,將來一定會曉得.我們若照樣作,一定蒙福的.
\newline
\newline
各位!我們別選擇服事.很多人的服事有所選擇；隨己意服事,要人聽他的話.這樣算是服事嗎?他不是作僕人,而是作了主人.
\newline
\newline
教會的人太多,僕人太少；神無法工作,無法施恩,我們沒有讓神與我們同工.我們當尊耶和華為我們的王,讓基督作教會的頭；我們只作工人,工人要順服.我們當彼此服事.
\newline
\newline
教會在這時代,是服事的團體,我們每個人必須服事；務必尊耶和華為大,為聖,祂是我們的王,是我們唯一的牧者.我們若真正有此態度,教會事工纔能發展,教會興旺,得復興.願神這樣恩待我們吧!
\newline
\newline
教會在這時代,我們特別想到本章聖經末段說:「我必使他們與我山的四圍,成為福源；我也必叫時雨落下,必有福如甘霖而降.」(26節)
\newline
\newline
我們希望神向我們施恩典,我們期望恩雨大降,期望神賜福給教會.那麼,我們今日當作甚麼呢?
\newline
\newline


\newpage

\section{第54屆~港九培靈會~唐佑之~第八講}
\label{sec:552}
\textbf{復興的景象}
\newline
\newline
連結: \href{https://www.hkbibleconference.org/session-message/view/552}{\texttt{https://www.hkbibleconference.org/session-message/view/552}}
\newline
\newline
\hyperref[sec:551]{< < < PREV SERMON < < <}
~
\hyperref[sec:toc]{[返目錄]}
~
\hyperref[sec:553]{> > > NEXT SERMON > > >}
\newline
\newline
以西結書 37:1-37:10
\newline
\begin{longtable}{cl}
\hline
\hline
章節 & 經文 (和合本修訂版)\\
\hline
37:1 & \begin{tabularx}{0.7\textwidth}{X} 耶和華的手按在我身上。耶和華藉著他的靈帶我出去，把我放在平原中，平原遍滿骸骨。 \end{tabularx} \\ \\ \relax
37:2 & \begin{tabularx}{0.7\textwidth}{X} 他使我從骸骨的四圍經過，看哪，平原上面的骸骨甚多，看哪，極其枯乾。 \end{tabularx} \\ \\ \relax
37:3 & \begin{tabularx}{0.7\textwidth}{X} 他對我說：「人子啊，這些骸骨能活過來嗎？」我說：「主耶和華啊，你是知道的。」 \end{tabularx} \\ \\ \relax
37:4 & \begin{tabularx}{0.7\textwidth}{X} 他又對我說：「你要向這些骸骨說預言，對它們說：枯乾的骸骨啊，要聽耶和華的話。 \end{tabularx} \\ \\ \relax
37:5 & \begin{tabularx}{0.7\textwidth}{X} 主耶和華對這些骸骨如此說：『看哪，我必使氣息進入你們裡面，你們就要活過來。 \end{tabularx} \\ \\ \relax
37:6 & \begin{tabularx}{0.7\textwidth}{X} 我要給你們加上筋，長出肉，又給你們包上皮，使氣息進入你們裡面，你們就要活過來；你們就知道我是耶和華。』」 \end{tabularx} \\ \\ \relax
37:7 & \begin{tabularx}{0.7\textwidth}{X} 於是，我遵命說預言。正說預言的時候，有響聲，看哪，有地震；骨與骨彼此接連。 \end{tabularx} \\ \\ \relax
37:8 & \begin{tabularx}{0.7\textwidth}{X} 我觀看，看哪，骸骨上面有筋，長了肉，又包上皮，只是裡面還沒有氣息。 \end{tabularx} \\ \\ \relax
37:9 & \begin{tabularx}{0.7\textwidth}{X} 耶和華對我說：「人子啊，你要說預言，向風說預言。你要說，耶和華如此說：氣息啊，要從四方而來，吹在這些被殺的人身上，使他們活過來。」 \end{tabularx} \\ \\ \relax
37:10 & \begin{tabularx}{0.7\textwidth}{X} 於是我遵命說預言，氣息就進入骸骨，骸骨就活過來，並且用腳站起來，成為極大的軍隊。 \end{tabularx} \\ \\
[1ex]
\hline
\hline
\end{longtable}
這段經文是復興的預言,神告訴祂的工人先知以西結,當時神的靈,神的手,在先知身上；我們可以想像,他有何等感受.相信今天教會照樣有此經驗,只要我們肯仰望神.
\newline
\newline
本章著重以色列民族復興的預言.中興在望,他們帶著樂觀心態,等待神的福份來到；他們等候之時,先知對他們說了很重要的話,這是先知親身所感受到的.
\newline
\newline
耶和華的靈,把以西結帶到平原之中.那裡大概是古戰場,曾經屢遭侵犯,與仇敵激戰多次,許多陣亡者伏屍戰場,不得埋葬.當人們憑弔古戰場時,枯骨遍野可見.以西結徘徊其間,緬懷歷史的回憶,悲傷心情感慨萬千.因神的審判臨到,使以色列民族遭受深重歷史的苦難；他們經過酷劫後,情景極其淒涼荒蕪；枯乾的骸骨,說明了久歷史的遺跡.難道至此地步的民族,還能復興,還有中興的希望嗎?以人的眼光看,是絕無希望的；但在神豈有難成的事?神有奇妙的工作,計劃和作為.當以西結在那裡觀察時,神要他經歷復興的景象.教會在這時代需要復興的景象,我們須在神前仰望祂；這並非人的工作,計劃,作為；乃是神的工作,計劃,作為.教會復興就是神的工作,計劃,作為.正如神使以色列民族復興,我們應以興奮快樂的心情來接受這個異象.
\newline
\newline
「祂對我說,人子啊!」在本書中,神屢次稱以西結為人子.「人子」意即脆弱,軟弱,衰弱,無用,必死,必朽的人.這個名詞耶穌基督自己也用過.在但以理書中,人子已經不再是脆弱必朽壞的人了；甚而至有彌賽亞的含意,特別在新舊約交替的時代,有些書籍都用人子,作為彌賽亞的稱謂.所以耶穌來到世界,他常自稱是人子；一方面表明他不過是人,另一方面祂也是彌賽亞.
\newline
\newline
我鼓勵弟兄姊妹,不但要多讀新約,也要多讀舊約；多多明白新約的含義,而舊約是耶穌基督的聖經.
\newline
\newline
當以西結被稱為人子時,他意識到自己的軟弱,無知,無能,無用.各位!我們也應有此態度；從一方面看我們實在沒有用處,沒有力量,無纔能,無知識；但是神要用我們,除非你肯將自己放在神手中,像以西結一樣.
\newline
\newline
「人子阿!這些骸骨能復活麼?」照希伯來的文學,常用修詞的問號,這種問號不必答覆；如果答覆也是反面的否定的.當時以西結看見荒涼的平原,遍滿骸骨的山谷.神要祂的兒女面對現實,不要逃避現實；要明白面對的情況,明白教會的現況；不能掩飾,不能隨便；等到神的恩典來到,凡事可迎刃而解.
\newline
\newline
我們若以世俗眼光來看,世界不斷進步,物質文明日益發達,表面上極為繁榮,燈紅酒綠,花花世界,處處迷人；但若用屬靈眼光看,就完全不同.近數年來西方文壇,有所謂荒野詩人,他們有文學的修養,深入的觀察；他們以為世界物質文明,給人類帶來精神疲乏.當物質文明越發達,人心越虛空,整個世界如荒蕪的原野,沒有內容,這是文學家的見解.至於我們屬靈的基督徒如何看法?保羅在以弗所二章給我們一幅清楚的圖畫,他把世界描成荒漠墳場；他說所有的人都死在罪惡過犯之中,沒有生命的氣息,都是死人,沒有希望.我們今天真是活在死人堆中.到底有何感觸,感想呢?求主憐憫我們!教會在這時代,須有先知的見識,觀點,見解.
\newline
\newline
當以西結聽見「這些骸骨能復活麼?」這話時,他無法回答,他說:「主耶和華阿,你是知道的.」我們眼所能見的,是物質方面；屬靈的光景,我們不知道.神若施恩,世界,社會,家庭,個人纔有希望.先知以西結有屬靈的感受,他體會神所問的必有特別意思,將帶給他一線的希望.所以他說:「主阿你是知道的.」他願主的旨意成就,他既有信心,有希望,神就告訴他.
\newline
\newline
如果我們沒有信心,只看現實,我們就沒有希望.我們傳福音,作見證,有甚麼用呢?但是如果我們有信心,在神面前真正有屬靈的感受；神就向我們發命令,叫我們作大事,祂且要為我們成就大事.
\newline
\newline
神對以西結說:「你向這些骸骨發預言,說:枯乾的骸骨呀!要聽耶和華的話.」枯乾的骸骨怎能聽耶和華的話呢?但神能作成不可能的事,在神沒有難成的事.我們無法想像神的能力,我們的信心實在太小；對我們所信的神估計太小,我們屬靈的思想了解太幼稚；偉大的神,因你我的幼稚,把祂貶了值.求主憐憫,擴大我們屬靈的容量；讓我們看見,聽見,想像,希望,相信；然後神就為我們成就大事.神說:「我必給你們加上筋,使你們長肉；又將皮遮蔽你們,使氣息進入你們裡面,你就要活了.」神的話語一出來就有能力,一進入就有功效,真正帶著神的能力.先知有此經驗,有此信心.
\newline
\newline
教會在這時代,須要有以西結般的信心；尊重神的話語,聽從神的命令,順服神的旨意.讓神話語的能力能夠出來,能夠表達,能夠顯現；若不相信,能力便受限制.神的話多麼有能力,不但安定在天,亦不徒然返回,祂的話無不帶著能力的.我們若有信心,必看見神的大能.
\newline
\newline
以西結具信心聽從神的話,他大聲呼喊,奇妙的事就發生了；枯骨聯終,生筋長肉.今天神的教會也須這樣,教會復興不是單高興快樂；不單靠外表的興旺,聚會的熱鬧；最要者乃骨頭互相聯絡,肢體互相聯絡.
\newline
\newline
今天人與人之間,沒有聯絡起來,教會弟兄姊妹,沒有真正聯絡起來.怎樣聯絡呢?
\newline
\newline
首先個人在神前仰望,追求；然後肢體之間纔會聯絡起來.請看!教會的團契,弟兄姊妹雖有交通；但不是靈交,而是社交.若然,只是成為社會團體而已；神的生命受了限制.社交是以自己為中心,而靈交是以基督為中心.以自己為中心,產生嫉妒和驕傲；正如剛纔所唱詩歌「如我比你強,我就驕傲；如我不如你,我就嫉妒.」人常充滿自私,常作比較,互相競爭.可能教會與教會之間,家庭夫婦之間,彼此競爭；沒有真正的聯絡,沒有真正的交往.靈交不但顧自己,同時也關懷別人；看見別人強,別人蒙神恩典,感謝讚美主!以別人蒙神恩典為樂事.如果從這個標準來看,教會少有真正的靈交.當神的話來到,骨頭與骨頭,肢體與肢體聯絡起來,這是何等奇妙的事!我們需要仰望神話語的能力.我們越有神的話語,越願主作我們的中心；越關心別人過於自己,對人更加友愛,更加謙卑.不驕傲,沒嫉妒,沒競爭；大家同心合意在主前討祂喜歡.
\newline
\newline
各位!我們多麼需要聯絡起來,如果骨頭與骨頭不聯絡,就不能成為身體.教會是基督的身體,但是今日的教會不像基督的身體,尤其是分開分裂的,不是完整的身體.求主光照我們,叫我們在這方面仰望主的憐憫.
\newline
\newline
肢體聯絡有筋就有力量,這樣不夠,再進一步還需要長肉.肉是保護骨頭的,健壯的人皮肉豐滿.主的生命在教會,祂保護我們,使我們得起考驗磨煉,在神前蒙受祂的恩典.生命越加強,這是話語的能力；當神的話語來到教會,當我們尊重神話語之時；不但筋與筋連絡,而且得著保護.我們需要主話語的能力,保護叫我們生命堅強.
\newline
\newline
再進一步,我們要美化,需要有皮遮蓋在肉上面,使我們皮有光澤,更加光潤.以弗所第五章特別提及基督為教會捨己；教會須貞潔,沒有縐紋等類的病.皮有縐炆,不夠健康,呈衰老現象,有的人未老先衰.求主的話語美化我們的皮膚,讓我們成為聖潔無瑕疵的見證.
\newline
\newline
以弗所第四章論及神在教會,預備了很多不同的恩賜；恩賜之中特別提到使徒,先知,傳福音,牧師和教師的恩賜.教會需要各種不同恩賜的工人,給我們話語方面的供應.使徒所作的工是建立神的教會,是用神的話語,來建立先知的工作.有先知的重點,將話語來供應；使徒的工作是廣的,先知的工作是深的.廣的工作需要,深的工作也需要.傳福音是把人從罪惡中帶回來,成為神的兒女,加入教會；在教會中事奉主,在地上學習天上的生活.牧師和教師列在一起；照我個人的了解,兩種恩賜都需要,尤其是今日華人教會所最需要的恩賜.有時弟兄姊妹不大知道恩賜,華人教會雖然成立多年,但恩賜還弄不清.當我們看見一個人很有講道的恩賜時,我們認識他是誰?是信徒,是先知,是傳福音的,或是甚麼都會的.其實不然,從歷史的教訓,我曾經提及過中國有幾個奮興家,神實在藉他們恩待中國的教會；但他們卻不會建立教會,也不會牧養教會,雖然他們對教會有很大的貢獻,可是在工作上卻做壞了.求主憐憫!使我們弄清楚甚麼是恩賜,是甚麼的恩賜.有時我們覺得某人不大有恩賜,我們以為講道就是恩賜之一,其實完全不是.
\newline
\newline
今天華人教會最大需要,就是牧師和教師；他們長期的經常的供應,弟兄姊妹所需；這樣,教會纔能增長.教會不能單靠特殊的講員,弟兄姊妹不能單靠每年一次的培靈會；何況培靈會的講員,也不能真正供應弟兄姊妹的需要.真正供應許要的,真正了解我們所需要的,應是牧師和教師.各位別誤會培靈會不重要,我覺得培靈會有很大價值；培靈會目的希望每個參加者,如主席剛纔的禱告:「是要愛主愛教會.」不單聽講員講道而已,是要自己真正在主前有追求；當然我們需要補充,需要先知的講道,需要信徒們建立教會,需要傳福音的；但是牧師和教師作更澈底的工作.為甚麼要有這些恩賜的人?「為要成全信徒.」(弗四12)本來我們不明「成全」的意思,近代考古家的發現,對這些用詞更可清楚了解.「成全」是古醫學骨科的名詞,當人跌交骨脫節,就要找骨科醫生「糾正」,「成全」即「糾正」之意.那些有神恩賜,專職蒙神特別呼召的工人,有本份,有責任,有使命,作糾正工作；而弟兄姊妹要謙卑,來領受神糾正的恩典.
\newline
\newline
今日華人教會,有時不易長進；因為有的觀念態度思想作法錯了.求神的話語糾正我們,叫我們所作所為,榮耀神的名.
\newline
\newline
「成全」也為普通漁人應用的名詞,晴天他們在太陽下晒網並修補破洞.所以成全信徒即補充之意.我們的缺憾不足之處,讓神的話語來補充.
\newline
\newline
「成全」更進一步的意義,是整容外科醫生的用詞.人的容貌有缺憾,不美觀,有縐紋,可找整容師糾正美化.當然我們不是注重身體的美,我們要注意屬靈的美,教會的美,教會見證的美；聖潔美麗,將自己完全獻給主；主要我們這樣獻給祂,多麼奇妙!祂就讓我們明白祂的話語有此功效.
\newline
\newline
「成全」也是軍事名詞.當人從軍入伍,必要配備,穿上軍裝.神的話語多麼奇妙!教會在這時代,須在真道上建立自己,配備之後各盡其職,建立基督的身體.
\newline
\newline
以西結卅七章的異象,我們不但看見以色列民的復興,更看到教會的復興.我們多有福,得以明白神話語的奧妙,神的話一解開就發出亮光,叫我們得以通達；一出來就有能力,讓基督的身體建立起來.
\newline
\newline
神對以西結說:「人子阿!你要向風發預言.」舊約風和靈同義.向風發預言,並非命令神的靈而是祈求神的靈；仰望聖靈的能力,讓生命的能力進到人裡面；否則人不能活動.當先知祈求神的靈,聖靈的能力就出現彰顯,發出功效；被殺者都活了,氣息進入骸骨,骸骨活了站起來,成為極大的軍隊.真是一幅偉大復興的景象!這個景象,也要來到你我的教會中；只要我們肯仰望神,尊重神的話,聽從神的命令,順從神的旨意；聖靈將成就偉大的工作.我們毋須猶豫,毋須等候；因聖靈已預備好了,要作奇妙工作.
\newline
\newline
我們在教會,雖然有時也注重神的話語；華人教很注重話語,但不大注重聖靈的能力.話語沒有聖靈就沒有能力；所以我們聽道講道很多,但我們缺少聖靈.
\newline
\newline
數十年來,西方甚至遠東,有許多靈恩運動；當人走到理性的極端,就來另一極端,此種影響現已帶到全世界.請大家不要反對靈恩,因為靈恩乃注重聖靈的恩典.這是好的,但請注意!我並不贊聖靈恩；因為靈恩只注重聖靈而不注重話語,且非常忽略話語.有話語,沒有聖靈就沒有能力；有聖靈而沒話語,就無法讓聖靈工作.我們需要聖靈也需要話語,兩相平衡,生命就長進,教會就復興；否則就會發生很多問題.今天注重靈恩的心,對神的知識實太幼稚,只注重情緒方面.我們的信心有情緒,有理性,我們應有意志,對主的態度信心應是全能的反應；不走理性的極端,也不走情緒的極端；靈恩不是華人教會的骨骼,我們需要聖靈,也需要話語,兩者並重,然後神做奇妙的工作.
\newline
\newline
本章,我們看見死人活了,站起來成為極大的軍隊.不是極大的會眾聽道,是極大的軍隊要出去作戰.教會在這時代是爭戰的,不是安穩的,不是只安逸在主裡,享受祂的恩典.希伯來書的信息有兩方面,一方面在聖殿的幔子裡,在主寶座前享受屬靈的恩典.這是我們每人都需要的.另一方面我們也要出去營外,去受主所受的凌辱.我們也要作基督的精兵,成為極大的軍隊,向罪惡作戰,向世界作戰.讓我們靠主得勝,且得勝有餘.
\newline
\newline
「我的僕人大衛,必作他們的王,眾民必歸一個牧人......」(24)以色列民合而為一,不再為猶大國和以色列國,不再為北方南方,完全聯合起來了.人作牧者不得完全,神是完全的,耶和華親自作我們的牧人；我們同歸於一個牧人,然後纔能真正發揮功能.昨天我們講牧者的哀歌,今天講牧者的讚歌.我們要讚美我們的牧者,祂施恩牧養我們；祂走在前頭為的是我們的安全,祂願為我們捨命,祂要我們跟從祂的腳蹤行.
\newline
\newline
主說要立平安的約.永約而且是恩約.耶和華說:「我以永遠的愛愛你們.」(耶卅一3)恩愛的約,永遠的約,平安的約；叫我們想到祂的恩惠,但是請各位注意!以色列的復興,是神要恢復祂的約,神的約是聖約.「約」是雙方面的,耶和華是主動；但以色列要回應.約是兩者之間的,神願意向我施恩典,叫教會復興；祂給我們永遠的約,恩典的約,平安的約；叫我們完全歸於一個牧人名下,在祂帥領之下,成為極大的軍隊爭戰.我們有何回應?我們要向祂立約,如今是立約的時候了.祂願與我們立約,我們可以拒絕；但是我們應該接受；如果我們接受,祂要向我們施恩典,然後我們把祂所施的恩典帶出去.神要我們成為祂的軍隊,求主這樣帶領我們.
\newline
\newline
讓我們現在就與神立約吧!
\newline
\newline


\newpage

\section{第54屆~港九培靈會~唐佑之~第九講}
\label{sec:553}
\textbf{增長的影響}
\newline
\newline
連結: \href{https://www.hkbibleconference.org/session-message/view/553}{\texttt{https://www.hkbibleconference.org/session-message/view/553}}
\newline
\newline
\hyperref[sec:552]{< < < PREV SERMON < < <}
~
\hyperref[sec:toc]{[返目錄]}
~
\hyperref[sec:554]{> > > NEXT SERMON > > >}
\newline
\newline
以西結書 47:1-47:12
\newline
\begin{longtable}{cl}
\hline
\hline
章節 & 經文 (和合本修訂版)\\
\hline
47:1 & \begin{tabularx}{0.7\textwidth}{X} 他帶我回到殿門，看哪，有水從殿的門檻下面往東流出，因為這殿是朝東的。水從殿的側面，就是右邊，從祭壇的南邊往下流。 \end{tabularx} \\ \\ \relax
47:2 & \begin{tabularx}{0.7\textwidth}{X} 他帶我出北門，又領我從外邊轉到朝東的外門，看哪，水從右邊流出。 \end{tabularx} \\ \\ \relax
47:3 & \begin{tabularx}{0.7\textwidth}{X} 他手拿繩子往東出去，量了一千肘，使我涉水而過，水到腳踝。 \end{tabularx} \\ \\ \relax
47:4 & \begin{tabularx}{0.7\textwidth}{X} 他又量了一千，使我涉水而過，水就到膝；再量了一千，使我過去，水就到腰； \end{tabularx} \\ \\ \relax
47:5 & \begin{tabularx}{0.7\textwidth}{X} 又量了一千，水已成河，無法過去；因為水勢高漲成河，只能游泳，無法走過。 \end{tabularx} \\ \\ \relax
47:6 & \begin{tabularx}{0.7\textwidth}{X} 他對我說：「人子啊，你看見了嗎？」他帶我回到河邊。 \end{tabularx} \\ \\ \relax
47:7 & \begin{tabularx}{0.7\textwidth}{X} 我回到河邊時，看哪，河這邊與那邊的岸上有極多的樹木。 \end{tabularx} \\ \\ \relax
47:8 & \begin{tabularx}{0.7\textwidth}{X} 他對我說：「這水往東方流，下到亞拉巴，直到海。所流出來的水，一入海就使水變淡。 \end{tabularx} \\ \\ \relax
47:9 & \begin{tabularx}{0.7\textwidth}{X} 這兩條河所到之處，凡滋生的動物都必存活；這水流到那裡，使那裡的水變淡，因此裡面有極多的魚。這河水所到之處，百物都必存活。 \end{tabularx} \\ \\ \relax
47:10 & \begin{tabularx}{0.7\textwidth}{X} 必有漁夫站在河邊，從隱‧基底直到隱‧以革蓮，全都成了曬網的場所。那裡的魚各從其類，好像大海的魚甚多。 \end{tabularx} \\ \\ \relax
47:11 & \begin{tabularx}{0.7\textwidth}{X} 但是沼澤與池塘的水無法變淡，只能作產鹽之用。 \end{tabularx} \\ \\ \relax
47:12 & \begin{tabularx}{0.7\textwidth}{X} 河這邊與那邊的岸上必生長各類樹木，可作食物；葉子不枯乾，果子不斷絕。每月必結新果子，因為這水是從聖所流出來的。樹上的果子必作食物，葉子可以治病。」 \end{tabularx} \\ \\
[1ex]
\hline
\hline
\end{longtable}
!感謝主我們從以西結書,看到復興的景象.這是神的心意,因以色列民是神所揀選的雖然他們失敗,神卻願意他們復興更新.
\newline
\newline
教會在這時代,也要將神的心意實現出來當教會復興之時,必有增長的現象;增長時的影響,一定明顯的遍及廣大地方神的心意要復興以色列民,叫他們重新回到神面前,得著祂的恩典.他們復興時,大大的增長,而且增長的力量,影響遍及遙遠的地方.我曾提到雅各因約瑟得福,葡萄樹枝長大又長,探出牆外;其影響力無法形容.
\newline
\newline
教會在這時代必須復興,必須更新永遠陶醉在神的恩典中,是不可能的;必須讓神的能力出去,遍及各處,讓神的恩典如大海汪洋「認識耶和華榮耀的知識,要充滿遍地,好像水充滿洋海一般」哈巴谷書和以賽亞書,都有這樣的話(哈二14,賽十一9)神的恩典那麼遼闊,那樣的深,那麼豪放;但是要看我們的屬靈容量如何.神無盡的恩典,要使我們的容量滿溢,影響遍及各地.
\newline
\newline
基督徒是信靠神屬於神的人,我們的盼望在於神;神的盼望在於我們,因祂要使用我們作祂的工,完成祂目的,實現祂計劃我們靠著神的恩典,讓祂自己的能力帶領我們,一直長進,一直增長;那增長的影響,實無法描寫,無法述說偉大的神,能力偉大,恩典偉大,祂的兒女也偉大,只要我們依靠祂,就能作偉大的工我們有偉大的人生,偉大的使命,偉大的工作;神借我們完成祂偉大的計劃.
\newline
\newline
神一直帶領以西結,從四十至四八章的主題,是聖殿,敬拜;是以色列人復興的中心這裡所指聖殿不是所羅門的聖殿,因為他所建的聖殿已經毀壞,已和耶路撒冷一同敗亡,一同毀滅這裡也不是說所羅巴伯的聖殿;因為這殿乃被擄歸回,才由所羅巴伯領導以色列人建造的這聖殿乃是理想的聖殿,神要叫屬祂的人,能瞻望一所理想的聖殿──就是主耶穌基督的身體.
\newline
\newline
教會在這時代,要成為理想的聖殿,把神的榮耀表現;因這能力,能吸引多人在神前俯伏敬拜.
\newline
\newline
以西結被帶到殿門口,在祂的異象中,看見那理想的聖殿,門檻下有豐富的水流出來.巴勒斯坦地方,水是極珍貴的.當時先知所見的水流下充滿,實非物質環境所能解釋.先知看見水由殿的右邊,在祭壇的南邊往下流.這讓我們有更深切的思想.不論新約舊約,水常象徵生命.有祭壇就有生命,一方面祭壇上的生命要完全獻上,用火燒掉歸於神；另一方面,祭壇也象徵生命的火祭.舊約代不像我們現在看的清楚,因為主耶穌來到世界,我們有教會,更能明白這裡所說的是甚麼.在十字架祭壇上,主為我們完全獻上,主耶穌基督不但是大祭司且是祭品,毫不保留把自己,放在祭壇上,作為燔祭完全獻上；作為平安祭,為我們設立平安和平,叫我們能與神和好.又獻上素祭,甘心樂意為我們捨身流血；獻上贖罪祭,叫信靠祂的,罪得赦免.獻上贖愆祭,叫我們的罪完全得到赦免.這恩典是白白的,不須付任何代價；但是在祂──神的兒子耶穌,卻要付何等大的代價.我們怎能輕忽呢?我們每逢親近主,來到神施恩寶座前,十字架就活現在我們面前.讓我們仰望十字架的主,祂為我們汗流乾,血也流盡；在十字架上說:「我渴了!」當祂乾渴時,在祂腳下,在祭壇底下就流出一條生命的活水；越流越長,越深,越遠,越充實,越多；於是我們就明白這個意思,生命的河流是從祭壇底下出來的.在十字架的祭壇下,已流出一條生命的活水,而且越流越深.以西結所明白的,是民族復興的景象；我們所明白的是教會復興的景象.我們看見的是從十字架開始,從十字架那裡出發；生命的水流越流越深,越長,越遠.這就是今天晚上我們要思想的.
\newline
\newline
先知看見天使拿著準繩,量一千肘,水就漲起來；漲到踝子骨；(以色列人注重象徵的意義)踝子骨著重腳的力量.行傳第三章記,彼得約翰到聖殿去,有個瘸腿的在那裡求賙濟,「彼得說金銀我都沒有,我奉拿撒勒人耶穌的名,叫你起來行走.」瘸子立即起來行走,且跳起來,因他的踝子骨健壯了.這裡給我們明白寶貝的教訓,基督徒在屬靈方面生命長進的過程中,不是突然的,而是步步的,逐漸的.願神的恩惠如此帶領我們.
\newline
\newline
昨天晚上,我們提及靈恩運動的價值,乃注重聖靈的工作；所偏差的是當人追求聖靈時,要一下得著聖靈很大的力量.並非聖靈不能這樣幫助我們,但正常的長進,不是突然而是逐漸的.
\newline
\newline
我們需要生命逐步的長進.第一步是腳力；如果有神的恩惠,有聖靈的能力在我們身上；我們的踝子骨就健壯了,就能往前走,且力上加力.保羅說:「我們既然靠著聖靈得生,就當靠著聖靈行事.(加五25)行事原是行走的意思,可譯為「靠著聖靈行走」或「在聖靈裡面行走.」我們的生命長進是由此開始,當聖靈的能力在我們身上,就叫我們跟隨主的腳蹤,與主同行；耶穌基督是這樣的呼召祂的門徒跟從祂.舊約時代挪亞和亞伯拉罕那些人都是與神同行,跟隨祂的腳蹤行.
\newline
\newline
基督徒的生活不是靜態的,我們必須往前走,不能靜止,要一直長進走在生命道路上.如果我用稍為具體的建議,即指基督徒的生活,一定要有聖靈的能力來領導,叫我們生活,有生命的表現.
\newline
\newline
培靈會大多是青年人,成年人,壯年人,中年人在哪裡呢?可能他們太忙,或因有許多事情攔阻.從表面看,大會人數很多；但其實正如剛才楊牧師說的:「培靈會人數雖多,不過僅佔基督徒人數百分之一.」不錯!神要向我們施恩,但有多少人是真正對培靈會會感到寶貴呢?我敢說凡來參加培靈會的,神不能不施恩典賜福給我們,我們都是有福的；因為我們把握機會,來領受神所賜的恩典.神在我們每人身上,必有美好的計劃；只要我們肯完完全全地擺在祂的手中,我們就必能蒙大福.願神使用來赴會者,成為香港教會復興的核心.當教會復興之時,增長的影響不只在香港,必定要遍滿全地.
\newline
\newline
各位!你們是特別蒙愛的,所以要重視自己的身份和地位.感謝主!也有很多青年人赴會,求神感動我們的心,在青年時候未感染世俗罪惡污穢之前,保存純正愛主的心；恐怕年事漸長,受了社會世俗的影響則難以單純了.我們應注意在生活中與主同行,要保持基督徒的生活方式.第一要聖潔,貞潔,不染世俗.第二要純樸,不要隨從世俗,不要像一般人注重吃喝,輕浮的虛榮.我們應腳踏實地,作主要我們作的工,在主的恩典中有屬靈的風度.
\newline
\newline
保羅在以弗所書(弗五17),教我們不要作糊塗人,要明白主的旨意如何.「糊塗」原意湖水表面的浮萍和污物.如果我們很表面化,很幼稚,很隨從潮流；我們就沒真正的深度,就不能讓神的恩惠,在我們身上表現.若把湖面污物清除,然後陽光反射在湖水上,就顯出美麗的畫面.香港的青年很能幹,但缺少屬靈的氣質.事實上,整個環境在世俗,沒有神的信仰情形下.我們需要學習,須有屬靈的深度,不從世俗,保持基督徒純樸的生活方式；並非落後,乃為真理作見證；這世界需要基督徒福音的見證.但願香港教會的青年人,有深度的長進,與主同行,讓生命長進.活水不只流到踝子骨,還要繼續增長.
\newline
\newline
「再量千肘,水就到膝.」(四七4)水到膝蓋必有敏銳感受.當我們在神前仰望,長進再長進,漸見水流力量更大.膝蓋是以色列人,特別看重的象徵；因也們注重俯伏,屈膝敬拜在神面前.
\newline
\newline
感謝主!華人教會也著重跪下禱告.並非說禱告須拘於形式；但我相信,當人跪下時,有更多更深的虔誠.我們實在需要追求屬靈的恩典,需要禱告的操練.讓我們每天膝蓋下工夫；因為聖靈的恩典,已經到了膝蓋；若不經操練,不知如何禱告.為何有時禱告沒有功效呢?因為沒有安靜等候神,沒有安靜在神前仰望.怡性需要工夫,鍜煉身體,必須操練,何況屬靈的事!但我們實在太忙碌,環境也太嘈雜,心情難以安靜.惟求神特別幫助,讓我們勤於操練.有時我們覺得禱告如同打空氣,這正像我們沒有開電掣,電流怎能通呢?有人說我們禱告要轉上並轉進,使我們與神有直接的交託；再而完全轉向神；學習,領受,體會,漸漸感到可以完全交托.曾經有問過一個出名的傳道人:「甚麼時候我們應該禱告?」答:「當你最不願意禱告時,就是你應該禱告之時.」
\newline
\newline
當一個青年基督徒成了家,應有基督化的家庭,應有家庭禮拜,在家庭中下膝蓋工夫.華人教會,甚至長執及教會領袖,都沒有家庭禮拜.有個教會長執的兒女告訴我,父母在教會中禱告很好,在家裡卻沒有禱告.青年人應極力爭取建立家庭禮拜,這是教會常忽略的.在家庭中,夫婦應同心跪在神面前禱告,教導兒女從小學習禱告；兒女的前途,是由父母的膝蓋,為他們鋪路的.我們不要效法,那班只在教會熱心敬拜的基督徒.真正的信仰不單在禮拜堂,必須在個人生活中,也在家庭中；我們的信仰才能真正建立起來,才能教導兒女,有正確屬靈的價值感.許多基督徒不知何謂價值感,賺錢方式,生活方式無異世人；信仰沒有真正應用在生活中.許多人年青時在教會很熱心,一入社會做事,結婚以後,就逐漸冷淡；到了中年,也許他們肯回到教會；可是深受世俗渲染,已經沒有屬靈的價值感.若這樣的人,在教會當長執,教會還能復興嗎?求神憐憫!我不敢責備任何人,但很多時候所看見的事,實在令神的心難過!
\newline
\newline
教會在這時代,特別在青年人中間,(我並非注重青年人忽略中年人)我們應真正地從頭做起,重新建立信仰,使信仰深入實際生活中.
\newline
\newline
有人曾經與我討論,「基督徒能否與不信主的人結婚」的問題.根本這問題不必討論.有的人說,當然牧師是說信與不信的不能同負一軛.一個真正的基督徒,應有清晰的價值感,這並非婚姻問題,而是整個的信仰問題.如果你和個不信主的人結婚,價值感不一樣,生活方式以及許多觀念都不一樣；對兒女教育也不一樣,怎能幸福呢?若以為只不過是禮拜天的事,那就未免太幼稚,而且非常嚴重了.
\newline
\newline
教會要復興,不是單靠興奮熱鬧一場就夠了,每個人當切實在神面前,在生活中完全仰望神；要下膝蓋工夫,努力往前走.水已到踝子骨,我們當與主同行,與主同想.我們的思想,意念,理想完全屬靈,那才是最徹底的.教會在這時代,華人教會在這時代；要重新檢討以往教會的失敗,成年人壯年人快點回頭!青年人從頭開始,與主同行,與主同想.
\newline
\newline
「再量千肘,水到了腰.」以色列人思想中的「腰」有兩個意思；一是工作事奉,一是爭戰.當耶穌基督為門徒洗腳時,祂就束起腰來.這樣行動方便可以服事,爭戰時束腰,才有力量爭戰.我們須束上腰,不可馬虎隨便,要認真打起精神,振作起來.先束緊腰,然後好好工作,好好爭戰；屬靈的力量,給我們有工作的能力,爭戰的能力.我們多麼需要這樣的能力!我們整天俗事纏身,沒有時間,沒有能力；我們需要聖靈的力量撐腰.
\newline
\newline
請看!水繼續流,遮沒頭頂,無法再量；全身浸在水中,我們感到完全與主同在.水到踝子骨,與主同行,水到膝蓋,與主同想；水到腰間,與主同工,靠主得勝.當我們整個被水淹浸時,我們完全淹沒有在神的恩愛中,完全消失在神的靈裡.這是何等的感受!完全沒有自己,主在我們裡面,我們在主裡面；完全投入,在主浩大的愛中.
\newline
\newline
「祂對我說:人子阿!你看見了甚麼?祂就帶我回到河邊......見岸上有極大的樹木......水往東方流去,必流入鹽海,使水變甜.」(6-8)又苦又鹹的水變甜了.各位!如果我們真正被神的靈帶領,被神的靈充滿之時,無論何往,苦必變甜.「這河水所到之處,凡滋生的動物,都必生活.(9)死的變活了,神的能力作奇妙的工作；但是我們必須長進.
\newline
\newline
生命的水不斷地流,經幾十個世紀.教會在這時代,神帶領我們看見,生命的水流得更踴躍.當水流經我們,淹沒我們,叫我們無論往何處；苦變甜,死變活.
\newline
\newline
當摩西領以色列人經過紅海,他們讚美神,他們曉得神為他們施行大拯救；但他們行經曠野三天,沿途沒水喝,他們就立即發怨言.後來找到水卻是苦水,摩西仰望神,神指示他一棵樹,他把樹丟在水裡,水就變甜(出十五25).許多解經家都認為這是象徵新約的十字架.我相信何時有十字架,苦即變甜,死即變活.今天有許多人在痛苦中,因你我的影響,叫他們苦變成甜；有許多人死在罪惡過犯中,因你我的影響,叫他們活過來.
\newline
\newline
鹽海是世上水流最低之處,離水平線1292呎,那裡的水最深最壞,在先知異象中,最壞的變成最好；最苦的變成最甜；無生命變成有生命.神有此能力,只要我們肯仰望祂,必有效果；無論往哪裡,苦變甜,死變活,而且有正常生命的供應,每月結新果子,結果纍纍.果子當作食物,叫人飽足；葉子能醫病,叫人健康,充滿生命.
\newline
\newline
我們是否有這樣的力量?教會有沒有這樣的力量?有否這樣的影響?惟有仰望神,奇妙的事就發生了.
\newline
\newline
我們多麼希望有這樣的影響,復興是必有的；生長是必有的,影響是必有的,改變也是必有的.只要我們肯完全把自己擺在神面前,生命的力量就出來了.
\newline
\newline


\newpage

\section{第54屆~港九培靈會~唐佑之~第十講}
\label{sec:554}
\textbf{歷史的進程}
\newline
\newline
連結: \href{https://www.hkbibleconference.org/session-message/view/554}{\texttt{https://www.hkbibleconference.org/session-message/view/554}}
\newline
\newline
\hyperref[sec:553]{< < < PREV SERMON < < <}
~
\hyperref[sec:toc]{[返目錄]}
~
\hyperref[sec:536]{> > > NEXT SERMON > > >}
\newline
\newline
以西結書 1:1-1:3
\newline
\begin{longtable}{cl}
\hline
\hline
章節 & 經文 (和合本修訂版)\\
\hline
1:1 & \begin{tabularx}{0.7\textwidth}{X} 在三十年四月初五，我在迦巴魯河邊被擄的人當中，那時天開了，我看見神的異象。 \end{tabularx} \\ \\ \relax
1:2 & \begin{tabularx}{0.7\textwidth}{X} 正是約雅斤王被擄的第五年四月初五， \end{tabularx} \\ \\ \relax
1:3 & \begin{tabularx}{0.7\textwidth}{X} 在迦勒底人之地的迦巴魯河邊，耶和華的話特地臨到布西的兒子以西結祭司，耶和華的手按在他身上。 \end{tabularx} \\ \\
[1ex]
\hline
\hline
\end{longtable}
這段經文是歷史的前言.神的工作不但在歷史的環境中；我們不但想到教會過去的,將來的；也想到教會在這時代,處於歷史的苦難中.正像神藉先知以西結傳祂信息之時,以色列也在歷史的苦難中.
\newline
\newline
以色列人敗亡,很多人被擄,在失望苦難之中,神向先知施恩；讓他作神的發言人,給那失望的人希望的信息.當我們覺得佈滿雲霧黑暗之時,神的僕人看見天開了.
\newline
\newline
教會在這遍地烽火,充滿罪惡,黑暗密佈的時代；讓華人教會在這時代,抬頭仰望我們的神.我們將看見天開,看見神的異象,聽見神的話；特別臨到神的工人,神的教會及弟兄姊妹,因我們有這樣屬靈的感受.耶和華的靈,耶和華的手,降在以西結身上,降在每個弟兄姊妹身上.教會在這時代,要有屬靈的感受,感受神的靈,感受神的手,感受神的能力.
\newline
\newline
從本章4-14,我們看見神是宇宙的主,祂是創造之主,也是救贖之主.在自然界中,神用風,火,雲為祂的僕役.在聖經中,以色列人常藉著這些自然界的現象,來明白象徵的意義.以西結不但是先知,也是祭司.祭司是非常注重象徵的.在這些象徵中,我們看見神的信實.
\newline
\newline
風是聖靈的影響,雲是榮耀的景象,火是公義的能力.神是烈火,當祂公義發出時,審判隨之而出.風好像聖靈有影響的力量.(耶穌也曾經這樣說)所以在自然界,火和風都有其作用,雲,是象徵榮耀的景象.當主耶穌升天時駕雲上去；天使告訴門徒,主怎樣的駕雲升天,將來也要駕雲降臨.在啟示錄曾說:「看哪!祂駕雲降臨.」火象徵神公義的能力.風象徵聖靈的影響.風是從北方而來,在以色列人的觀念,北方是很遙遠的地方,是侵略者的出處.但是神統治萬有,神的能力浩大,足以影響很多人,很多國家,很多時代.
\newline
\newline
教會在這時代,我們看見神公義的能力,看見聖靈影響,也能夠看見神的榮耀.
\newline
\newline
以西結在被擄之地迦巴魯河邊,在外邦地,看見神的榮耀；神的榮耀不僅在耶路撒冷,也在各處,因神是統管萬有的主.我們所信靠的神,比我們所能想像的大得太多了；於是我們從自然界中,想到神恩典的領域.神不但統管萬有,而且祂的恩典也遍及各地.
\newline
\newline
我們從自然界,可以想到神的恩典,也可以聯想神在歷史中的作為,因神是歷史的主.
\newline
\newline
法國大文豪說:「神的手在歷史的巨輪上.」歷史往前進展中的一切,都有神的手在其中.神在歷史當中,在歷史之上；藉著歷史,啟示祂的作為.
\newline
\newline
「時機」「危機」或「轉機」的原意是審判,都是神審判的作為,亦可謂神的救贖.信靠祂的就不在神審判之下；違背祂的,就在祂忿怒之下.
\newline
\newline
從自然界象徵的意義,可看見實際的事.
\newline
\newline
顯出的四活物,形狀有人的臉,獅子的臉,牛的臉,鷹的臉.在以色列人的思想中,人表明智慧,獅子表明勇敢,牛表明穩健,鷹表明敏捷.我覺得這是很好的啟示,活物是服役之靈.事奉時須具備條件,有人的智慧,牛的忍耐穩健,獅子的兇猛勇敢,而且有鷹的敏捷,這可使我們了解神在歷史中的作為.教會在這時代,應當具備這四項條件,我們應有屬神的智慧,看見別人所看不見的,明白別人所不明白的.應有獅子的勇敢,向這世界宣戰,與罪惡宣戰；從事福音爭戰.我們當穩健敏捷而勿冒失,有始有終.教會在這時代,須有此精神；須有中年人,老年人,教會長者那般的穩健；青年人那般的勇敢敏捷,須有屬靈的智慧；然後同心合意往前走.
\newline
\newline
「兩個翅膀遮體」表明敬虔謙卑.當以賽亞蒙召時,兩個翅膀遮臉,表明敬虔；兩個翅膀遮腳,表明謙卑.這裡說的「遮體」就是表明這兩個意思.但沒有說「遮臉」是因為要看聖靈往哪裡去,他們就往哪裡去；這是要清楚看見聖靈的引導.「遮體」,表明活物對神的敬虔,對人的謙卑.兩個翅膀相接,表明彼此合作,合一,真正配合,同向前走.
\newline
\newline
12節清楚告訴我們,要完全正確跟隨靈的引導,纔不偏差.多麼寶貝深切的真理!教會在這時代,要走正確的路,要讓神領導我們；靈往哪裡去我們也往哪裡去.一直向前行,不轉身,沒有自己的主張,見解,選擇,打算；定睛仰望神,祂往哪裡,我們也往哪裡.與神同行,只許向前不許退後,只許成功不許失敗.
\newline
\newline
「輪上有眼睛,輪不停轉動往前走.......」輪象徵神的引導和安排.神無所不在,也無所不能,在祂沒有難事.輪上有許多眼睛,輪不斷轉動,何處都去,分秒不停.神在自然界中,在歷史中,在人心中,在教會中行動著,永不停止.
\newline
\newline
教會在這時代,必須有行動,非口講而已.當具體行動,照著聖靈的引導；或靜,或動,必須一致.因為只有一個生命,一個能力,一個本能,一個行動,一個方向,一個目標.教會在這時代,要完全仰望聖靈的帶領.以弗所書說,身體只有一個,聖靈只有一位.(弗四4)我們當一致,同心合意往前走.火的時代須往前走,不可停頓.
\newline
\newline
(22-28節)整個異象的中心是寶座,神在寶座上統治管理作王.整個歷史進展過程中,讓我們看見教會在這時代,需要看見神的寶座；在我們心中,家庭中,教會中,甚至在整個世界,神高高在寶座上.當我們尊祂為王,在寶座前俯伏時；我們就能看見神的榮耀,亦看見我們應該走的路向.
\newline
\newline
活物翅膀的響聲,像大水的聲音.今天教會見證的聲音,有沒有如大水聲音的宏大,如波濤的洶湧呢?教會有沒有作到屬靈的爭戰呢?當神發聲時,誰能站立得住呢?在人們恐懼之中,正是應把福音帶給他們的時候,讓他們肯屈服在神威榮之下.
\newline
\newline
活物行動,有時也需要站住,把翅膀垂下.應動之時活動,應等候之時等候.
\newline
\newline
有時我們不願意等候,不明白神的帶領,不肯垂下翅膀,不斷行動,掙扎,活動,非常之危險!我曾說過,香港青年很有幹勁,香港的教會有很多活動；世界各地的教會也是如此.有的活動和節目是對的,但有些卻是錯的.需要等候時,就應該等候；需要活動時,就要活動.
\newline
\newline
摩西帶領以色列人過紅海時,神為他們把海水分開,他們直往前走,不能停住；但有時遇到危險,理應快跑；而神卻吩咐摩西叫他們只管站住,等候安排.
\newline
\newline
神在自然界中,歷史中行動,歷史不斷進展.教會在歷史中,應該要走在前面領導,還是落在後面妥協呢?這是值得我們每個人深切思想的.
\newline
\newline
教會的路是全職工人的路,蒙神呼召特別作神工的人,靈往哪裡就必須緊緊跟隨.作工人的如看不見路,又怎能使教會往前走呢?教會軟弱原因很多,雖不能完全歸咎於傳道人,但傳道人的責任,是無可推諉的.教會興旺是因為工人忠心,教會軟弱也就是因為工人軟弱.工人應該清楚知道神的帶領,有智慧,勇敢,穩健,敏捷的行動.教會在這時代必須往前走,工人們要真正在神面前行動起來.我們也當為神的工人禱告,幫助關懷他們,讓他們領導我們往前走.
\newline
\newline
教會的路是上行之路,在詩篇中有上行之詩,上行則上高山.走錫安大道的人須力上加力,有時筋疲力竭,他們唱讚美詩,彼此鼓勵,一直往前走.教會在這時代和每個時代,同樣地要走向前,走向上的路.我們一定要靠著主的力量往前走,等候耶和華的必重新得力.有時教會往前走有困難,有問題,多有攔阻；但是教會要長進,並非單講教會增長的原則和方法,最要緊的,我們必須看清楚,往前向上走；困難艱苦是必然有的,但我們一定要走!
\newline
\newline
教會的路是正直的路,耶和華為自己的名,引導我走義路.(詩廿三)義路乃正直的路,不遍不倚,穩健正確平衡地走在其中.
\newline
\newline
今天華人教會,甚不平衡,重此輕彼,偏斜偏差,十分危險!求主幫助我們如何走得正當.從聖經中看,我們教會需要先知型的信仰經驗,注重神的話語；也需要祭司型的信仰經驗.教會須走祭司的路,就是好好敬拜獻祭,醫治疾病,教導律法,作差傳工作；要注重實際,勿重形式.可惜華人教會不大注意,很可能,我們敬拜,仍是隨隨便便的敬拜!我們更需要智慧型的信仰經驗,正如舊約的智慧者一樣；他們注重道德,但並非忽略了屬靈的事,因為信心沒有行為是死的.我們一直求神給這個那個,但在行為上卻沒有表現出信心來；在生活上怎能有聖潔的見證呢?智慧著重理性,但並非取代信心；而是信心裡面要有理性.信心非理性的,也非反理性的,而是超理性的.有的弟兄姊妹,以為信心毋須帶著理性,以致於迷信,那就錯了.求主憐憫!讓我們有均勻平衡的生活,讓教會所走的路,是正直正當均衡的.
\newline
\newline
教會所走的路是更新的路,不再留戀舊事；但舊的好處是應該保守的,而有些地方是應該改進的,教會應該有創新的精神.聖經說:「日光之下沒有新事,但在基督裡我們卻是新造的人,所有的都變成新的了.神給我們的恩典,每早晨都是新的；我們的外體雖然損壞,但內心卻一天新過一天.」這給我們中年人有重要的警告,我們有時太保守,任何事都不願改進；以致墨守成規,固步自封,在教會中和青年人合不來,神要我們有更新的精神,我很誠懇地勸青年人要保守一點；也很誠懇勸中年和老年人,要開明一點.太保守,以為是保守純正的信仰,結果攔阻了教會的長進；久而久之,教會就更加落後了.主耶穌基督的十字架,為我們開了一條又新又活的路；我們毋須再走舊的,死的路了.
\newline
\newline
神在每時代都有祂的工人,有些被神大大使用；他們的精神,和在神面前的忠心,是我們應該效法的.不過他們的見解,可能只適合於那時代,而不是適合於這時代的.華人教會也有很多在上代被神重用的工人,不過我們當往前走.而不是還停留在三四十年代的光景.火的時代推進我們,只許向前不許退後；神雖讓祂的工人會過去,但祂的工作還要往前,所以我們必須在神面前求長進.
\newline
\newline
教會的路是得勝的路,要爭戰纔有得勝,要得勝必須爭戰；教會必須爭戰,也必須得勝.
\newline
\newline
教會的路是合一的路,不是組織的合一,乃是屬靈的合一.神帶領我們往前走,我們不能有偏見.有人以為宗派不好,但宗派是神在歷史中的作為.有人以為獨立的教會,比宗派更屬靈；但宗派教會,卻要經年代的考驗.相反地,獨立教會卻有許多不完全,求主給我們看清楚,知道合一的路是重要的.我們不要反對宗派,因宗派是神在歷史中的作為.許多華人教會,很多基督徒,以為屬靈就不要宗派,這是極大的錯誤.
\newline
\newline
求神憐憫!使我們合而為一,在基督裡我們是一體的；因身體只有一個,聖靈只有一位；不容我們有分開紛爭的現象.我們當在屬靈方面,見證方面,敬拜方面,合而為一；不要看自己比別人強,要彼此禱告,彼此合作；還有一點點時候,那要再來的主就來,必不遲延；若有人不愛主,是可咒可詛的.
\newline
\newline
各位弟兄姊妹!我們讀過這章經文,看見神的榮耀,看見神的行動,也看見神的使者在行動著.所以教會在這時代必須往前走.
\newline
\newline
我們只有一個禱告說:「主啊!靈往哪裡去,我們也跟隨往哪裡去.」靈站住,我們也站住,當站住時,我們也像活物一樣；垂下翅膀,安靜等候；當行動時,把翅膀連接起來,直往前行,並不轉睛.我們也要像(詩二九32)說:「你開廣我心的時候,我就往你命令的道上直奔.」
\newline
\newline
求主開廣我們的心,好叫我們往祂命令的道上直奔.
\newline
\newline
起來!我們走吧!
\newline
\newline


\newpage

\section{第54屆~港九培靈會~張慕皚~第一講}
\label{sec:536}
\textbf{教會–基督的身體}
\newline
\newline
連結: \href{https://www.hkbibleconference.org/session-message/view/536}{\texttt{https://www.hkbibleconference.org/session-message/view/536}}
\newline
\newline
\hyperref[sec:554]{< < < PREV SERMON < < <}
~
\hyperref[sec:toc]{[返目錄]}
~
\hyperref[sec:537]{> > > NEXT SERMON > > >}
\newline
\newline
教會是神建立和擴展天國的器皿,神要透過教會來改變這個世代;因此我們相信,教會必須站在時代的尖端,扮演先知的角色但可惜許多時候,教會卻不能完成神所交付的責任.
\newline
\newline
在西方的農村社會中,有一種獵犬,常在田間的小路上奔馳;而在獵犬後面,則是一輛馬車很多人起初觀察,以為獵犬是在領導牠後面的馬車其實並不是,當獵犬走到十字路口時,牠便蹲在路邊的樹叢間;等馬車走過後,牠又再以極快的速度奔在馬車前面因為獵犬並不知道前面該走的方向所以,獵犬只是在那裡自我陶醉,牠以為自己在領導馬車,其實牠是欺騙了自己同樣,教會應該在時代的潮流中扮演領導的角色,這是神的心意;但我們往往卻只像那隻獵犬而已.
\newline
\newline
不久以前,曾經有人訪問過美國一些社會知名人士,問他們在美國社會中,有什麼機構團體是最具影響力的.訪問結果,他們第一個列舉的是電視台,而教會則排列第十八位.這叫我們看見,在基督教如此興旺的美國社會中,教會在社會人士的眼光中,竟然也是這樣的落伍.我不敢相信,如果香港有類似這樣的調查,在香港人眼中,教會將落在第幾位.然而我們在神面前,仍不應對教會失望,因為神在教會身上有牠美好的旨意.如果我們能夠回到教會基本的原則上,我相信神必能使用我們今日的教會.
\newline
\newline
當我們看到香港的教會時,就發現一種情形,就是香港社會所發生的問題,往往社會人士和各種不同的機構,都為社會講了很多說話,最後才見教會也略略的也提到幾句譬如說我們最近非常關心一九九七的問題,外間人士談論得非常激烈,但教會卻很少提到又譬如有很多個人道德問題,社會道德問題,諸如墮胎問題,同性戀問題;但教會絕少發言去領導這道德的潮流.香港有嚴重的青少年問題,政府為此而訓練許多社工人才,在港九各地有許多青少年中心,還有很多外展工作,都是主動去尋找青少年的.但今天教會對香港的青少年問題發表過什麼?做過什麼呢?
\newline
\newline
教會必須要成為時代的教會,領導潮流,第一樣我們要做的,就是回到聖經關於教會的基本原則上去.今天教會有一個危機,我們有很多先進處理教會問題的方法.我們用現代的工商,行政管理學,來應用在教會的行政上,管理和組織上;我們用許多公關的方法,應用在教會的公關工作上這些都是好的但假如我們只追求怎樣把這些方法應用在教會工作上,而忽略了聖經關於教會許多基本的原則；教會便不能站在時代的尖端,領導時代的潮流,被神使用得著因此我們要返回聖經的基本原則上去在今次講道會中,我會選八個值得我們留意的教會表像,看看我們必須回到的聖經原則上去.
\newline
\newline
第一個我們要思想的表像,是基督的身體教會被稱為基督的身體聖經提到有兩個基督的身體,一個是在十字架上,為我們捨身流血的基督的身體；然後神為著時代的需要,設立另一個基督的身體,就是教會.要將耶穌基督在十字架上完成的救恩,實施在這世代中.
\newline
\newline
我們要從哥林多前書十二章內,思想教會建立的三個基本原則,就是教會的成員,教會的組織,和教會的行政.
\newline
\newline
(一)教會的成員(林前十二1-3)
\newline
\newline
「弟兄們,論到屬靈的恩賜,我不願意你們不明白.他們作外邦人的時候,隨時被牽引受迷惑,去服事那啞吧偶像,這是你們知道的.所以我告訴你們,被神靈感動的,沒有說耶穌是可咒詛的;若不是被聖靈感動的,也沒有能說耶穌是主的」建立教會最重要的,就是建造的材料教會是由什麼成員,什麼肢體,什麼種類的會友組成的呢?這三節聖經告訴我們兩方面關於教會建立的材料.第二節說教會是建立在一班重生得救的人身上.哥林多教會的信徒,以前是在教會以外,是敬拜偶像的;但現在已經在撒旦掌權的黑暗世界中,被神遷移進入他愛子光明的國度中,以前他們是服事那啞吧偶像的;但現在已蒙恩得救教會是.建立在一班蒙恩得救的人身上,這是最基本的.但我們要留意的是第三節,特別是「主」字.耶穌不但是我們的救贖者,祂更是們生命的主,二者是合一的今天教會所需要的力量和材料,是一班過委身生活,完全屬於主的信徒耶穌說:!凡稱呼我主阿主阿的人,不能都進天國,惟能遵行我天父旨意的人,才能進去.我們必須在實際生活中,尊主為大,順服祂一切的旨意.教會若要重新得力,就必須返回奉獻生活的原則去.
\newline
\newline
講到委身,有兩方面是我們必須注意的聖經所講的委身,是全人的委身,全人的奉獻;包括我們的身體,思想,行為,和我們的一切,都獻給主,在耶穌的主權下,過我們的生活聖經告訴我們,這委身的生活,也是全時間的委身生活,我們不單在星期天敬拜神,過部份時間的委身生活;聖經所要求的,乃是全人的委身,和全時間的委身.其實我們的職業和工作也是事奉,我們應該把我們的職業事業化,這樣才符合全時間的委身.
\newline
\newline
個別信徒的質素,決定了教會整體的質素,這是個明顯的道理比如一千人聚會的教會,只有二百人過全人全時間的委身生活;這會就只有百分之二十是活動的.那就是說,基督的身體只有百分之二十是活動的這樣的教會,並不能為神做很多的工作所以,改良建造的材料是非常重要的有三方面的工作我們可以做.:
\newline
\newline
第一:是傳完整的福音,告訴人耶穌是他的救主和生命的主.
\newline
\newline
第二:是推動門徒的訓練,讓人作主的門徒.
\newline
\newline
第三:是不斷尋求奉獻生活的更新.
\newline
\newline
跟著我們看教會組織的基礎.(林前十二11)「這一切都是這位聖靈所運行,隨己意分給各人的.」基督身體內有不同的肢體,代表教會的弟兄姊妹有不同的恩賜,而這些恩賜都是聖靈隨己意分給各人的.所以真正組織教會的,是聖靈自己.教會是一個有組織的屬靈有機體,而組織教會的是聖靈,是他把恩賜分給各人的.但今天教會卻有很多人為的組織.我們關心人為的組織,過於關心聖靈怎樣組織教會.如果我們要回到聖經基本的組織上,我們便須發掘弟兄姊妹的恩賜,任何榮神益人的本能,都是聖靈所給的聖靈已經將恩賜賜給弟兄姊妹;我們要把他發掘出來,讓各人在教會中都能有事奉的崗位只有總動員的教會,各人都參與事奉的教會,才能產生活力,才能站在時代的尖端,被神使用.
\newline
\newline
當我們知道弟兄姊妹的恩賜時,便要把他們配搭起來.聖經講到教會的組織時,著重一很簡單的道理,就是教會的組織必須簡單而切合實際.我們必須剪去教會組織中不必要的枝葉;.今天教會發生困難,很多時候是因為教會組織過於複雜除去枝葉之後,教會便會很快前進另外,聖經給我們一個原則,就是教會組織必須按需要而增減使徒行傳第二章,選立初期教會七個執事時,他們是按需要設立和差派的.神什麼時候為摩西預備助手呢?是在他有需要的時候,不需要的時候便要放棄,不必為了保持固有傳統而保留.初期教會有凡物公用的組織,但稍後即沒有這方式,因為已經不合時宜,沒有需要了.
\newline
\newline
第三要思想的,是教會行政的基礎,十四至廿七節.所謂教會行政,就是怎樣推動教會行動的組織使能完成耶穌所交給我們的大使命.
\newline
\newline
我們留意兩個有關連的原則十四節:「身子原不是一個肢體,乃是許多肢體.」二十節:「.但如今肢體多的,身子卻是一個」把這兩節經文合起來,我們便可以得到一個很基本的原則,就是分歧中合一的原則.教會愈大,部門愈大,這是需要的.主日學部,傳道部等等,這是分歧的原則;但我們要在分歧中尋求合一.我們有不同恩賜的弟兄姊妹,但卻要尋求彼此配搭,使成一個有用的身體.如果我們在教會行政上忽略了這基本的屬靈原則,便很難推動教會的工作.
\newline
\newline
講到分歧中的合一,最大的敵人,是個人英雄主義.一有個人英雄主義的產生,便不能在分歧中尋求合一,而教會的事工便無法推展.哥林多教會有一種現象,就是弟兄姊妹注重個人的恩賜,高舉個人的恩賜,因此便產生問題,成為教會行政的致命傷.我們需要在牠裡面學習彼此相愛相顧,因為愛是合一和連結的能力.
\newline
\newline
另一個有關連的教會行政的原則,二十五節:「.免得身上分門別類,總要肢體彼此相顧」分歧中的合一所帶來的,是彼此相顧而彼此相顧最大的敵人,便是我們忽略教會中弱小的肢體,只關心照顧我們認為是重要的肢體.但保羅教導我們要特別關注那些沒有體面的肢體,給他們有體面.教會行政若要做得好,便要照這原則行.
\newline
\newline
讓我們愛教會,投身在教會的事奉中,以表示我們愛基督,樂意的事奉他.
\newline
\newline


\newpage

\section{第54屆~港九培靈會~張慕皚~第二講}
\label{sec:537}
\textbf{教會–基督的新婦}
\newline
\newline
連結: \href{https://www.hkbibleconference.org/session-message/view/537}{\texttt{https://www.hkbibleconference.org/session-message/view/537}}
\newline
\newline
\hyperref[sec:536]{< < < PREV SERMON < < <}
~
\hyperref[sec:toc]{[返目錄]}
~
\hyperref[sec:538]{> > > NEXT SERMON > > >}
\newline
\newline
昨天我們用基督的身體這表像,來講到教會的重要性.特別思想到教會建立的三方面:教會的成員,教會的組織,和教會的行政.今天我們用聖經另一個表像來講到教會的另一種特質,就是新婦.教會被形容為基督的新婦,這是講到教會尊貴的地位.因上,無論我們對現在的教會有怎樣的感受,批評或是不滿,但神卻看重教會；將最尊貴的地位給教會,所以沒有人能看輕她.如果我們愛神,也必須愛神的教會.
\newline
\newline
(約十三34-35)「我賜給你們一條新命令,乃是叫你們彼此相愛.我怎樣愛你們,你們也要怎樣相愛.你們若有彼此相愛的心,眾人因此就認出你們是我的門徒了.」(林後十一2)「我為你們起的憤恨,原是神那樣的憤恨,因為我曾把你們許配一個丈夫,要把你們如同貞潔的童女,獻給基督.」(弗五 22-24)「你們作妻子的,當順自己的丈夫,如同順服主.因為丈夫是妻子的頭,如同基督是教會的頭,祂又是教會全體的救主.教會怎樣順服基督,妻子也要怎樣凡事順服丈夫.」這幾處經文讓我們看見,教會如果是基督的新婦,他應該有一個清楚明顯的記號,就是對基督有委身的愛.新婦的表像,是代表委身的愛.
\newline
\newline
今天我們怎樣察覺一間教會是蒙神恩典,被神便用,有神同在的呢?是看他會友眾多,人才鼎鼎,組織完善嗎?不一定,乃是看他是否有對基督那種委身的愛.保羅說:「妻子怎樣順服丈夫,教會也當怎樣順服基督.家庭是以愛來連繫,丈夫要愛妻子,如同愛自己的身子.妻子順服丈夫,也是出於愛,是愛的順服,而不是勉強或假意的服從.同樣,子女在主內聽從父母,也是愛的服從,若沒有愛,這樣的服從是痛苦的,也是毫無價值的.在這個缺乏愛的世代中,教會應該回到聖經的基礎去,對基督有委身的愛.一個充滿嫉妒紛爭的教會,四分五裂的教會,都不能為主作見證.時代尖端的教會,必須留意在主面前追求愛.
\newline
\newline
我相信我們都感受到,香港是一個缺乏愛的地方.每一天的新聞,都充滿謀,搶劫,仇恨,戰爭,這些都是缺乏愛的表現.正如耶穌描寫末世的情形,說人要彼此陷害,彼此恨惡,不法的事增多,很多人的愛心,漸漸冷淡了.今天的社會,正好充滿耶穌預言的這一切事.我們必須依照聖經的教導,學習在教會中彼此相愛,并將從聖靈那裡得來的愛伸延出去；否則教會不能在這世代中起任何作用.
\newline
\newline
教會是基督的新婦,必須有委身的愛.我們來看看愛和委身這兩個有關連的真理.聖經講到愛的時候,說愛是由意志出發,以達到全人的.許多時候我們把聖經這種愛和男女情感的愛混淆.不少青年男女,在有愛情的時候便結婚,但後來日久失情,沒有感情時便離婚,因為他們的愛是以情感為出發點而已.但神所給我們的愛,并不是這樣.這並不是說,我們的愛沒有情感,而是說我們的愛以意志為出發點,而達到個人的.所以聖經所講的愛,是立志將好處帶給人.雖然別人對不起你,但你仍然愛他,仍然立志將好處帶給他.這是神的命令,我們無可選擇.主耶穌教訓我們要愛仇敵,為仇敵禱告,更何況是我們的弟兄姊妹呢!愛是有行動的,立志將好處帶給人.
\newline
\newline
聖經告訴我們,愛是一種克服罪的能力.罪是由於自我中心,這種能力使人與人分隔.當我們犯罪的時候,我們與神分隔,甚至與自己也有分隔.我們想做的事卻無力做,不願去做的卻做了,我們內心有矛盾鬥爭.罪也叫我們人與人之間有分隔,人與環境也有分隔.但神今天要用愛來克服罪惡的問題,神差祂愛子到世上,在十字架上彰顯祂至高的愛,把我們與神連結起來.當我們信主的時候,這十架的愛便澆灌在我們心內,叫我們能把這愛伸延出來,使我們能夠與人和好.
\newline
\newline
另外,愛也是教會合一的能力.教會是否有問題呢?一個很容易的診斷,就是看看教會有沒有愛.有愛便有合一；有罪便有分裂,這是必然的現象.這愛是不可或缺的,它是教會應有的本質.哥林多前書十三章特別講到愛的真理,因為哥林多教會是個屬肉體較弱的教會,他們雖然有超自然的恩賜；但在他們中間有嫉妒紛爭,所以保羅說他們中間所要追求的,最重要的還是愛.
\newline
\newline
保羅在林前十三章一開始的時候,特別講到教會有三方面,如果在這三方面缺乏愛,教會便會被神棄置.首先他說我們若能說萬人的方言,并天使的話語,這話是說到教會的弟兄姊妹雖有很高深的屬靈的追求,甚至能講萬人的方言,天使的話語；但如果沒有愛,那便成為鳴的鑼,響的鈸.那就是說,如果我們只追求個人屬靈的長進,而沒有愛的產生,我們的生命便無法產生屬靈的音響；所發出的只是嘈吵的聲音,別人聽了會刺耳.
\newline
\newline
保羅又說,我若有先知講道之能,也明白各樣的奧秘,各樣的知識,而且有全備的信,叫我能夠移山,卻沒有愛,我就算不得甚麼.這裡所講的,是教會的事奉.我們有講道的事奉,明白奧秘,有教導的事奉；又有全備的信,能夠移山.教會最困難的事,我們都能憑信心解決,但如果沒有愛,後果是非常嚴重的,就是這一切都算不得甚麼,原來的意思就是等於零.今天我們可能有很多事奉,為主勞碌奔波,像以弗所教會那樣,但缺乏愛,主便說要將這燈台從原處挪開.
\newline
\newline
保羅繼續說:我若將所有的賙濟窮人,又捨己身叫人焚燒.這是社會關懷,到達一個巔峰的狀態,但若是沒有愛,就毫無益處.愛是教會不可或缺的,甚麼時候教會缺乏愛,一切事奉都等於零.所以我們常省察自己,到底我們的事奉是否出於愛,而這愛又必須是委身的愛.
\newline
\newline
跟著我們思想委身的真理.講到委身,聖經說是接受救恩的信的一部份.那就是說,我們接受救恩時,我們的信是包括委身在內的.我們不能把耶穌是主和救贖主這兩個事實分開,因為委身是我們得救的信的一部份.得救的信是全人的信,首先我們在理性上所接受,接受聖經所講有關救恩的真理.但得救的信也包括情感,像彼得前書一章所說:「卻因信他,有說不出來,滿有榮光的大喜樂.」這是情感的回應.但得救最重要的,還是委身.我們接待耶穌在我們心中為主為王.
\newline
\newline
另外,聖經告訴我們,委身是生命成長的先決條件.如果我們屬靈的生命要長進,就必須委身.今天很多基督徒沒有長進,原因是他們一開始沒有將自己奉獻給神,沒有過一個委身奉獻的生活.在亞伯拉罕的生命中,我們可以看見這一個過程.亞伯拉罕是信心之父,他因為信神的緣故而愛神；他為了愛神而委身,他的委身是有進步的,一直進到徹底的委身.
\newline
\newline
在亞伯拉罕的生命中,有三個重要的祭壇,成為他生命長進的三個記號,也是他生命澈底委身的三個里程碑.頭一個是創世記十二章的示劍祭壇.那裡記載神呼召亞伯拉罕,要他離開家人,離鄉別井,到神所要他去的地方.今天我們也要有這樣的委身,跟隨神任祂帶領.第二個是十三章的希伯崙祭壇,在那裡他願意為神放下財富.他為愛神的緣故,不願和他侄兒羅得爭吵,在選擇土地的時候,又甘願把上好的地給羅得,自己揀選瘦的.後來神再次堅定與他所立的約.第三個更重要的祭壇,是在摩利亞地的祭壇,記載在創世記第廿二章.在那裡,神要他把獨生兒子獻上.這是神給他的考驗,其實神早已有預備,但是亞伯拉罕卻達到了他生命的最高峰,他對神有澈底委身的愛.今天,除非我們肯對基督委身,我們的生命很難成長.委身是生命成長的先決條件.
\newline
\newline
聖經告訴我們,委身就是生命的能力.教會今天所需要的,是委身出來的能力.因為委身有分隔的作用,能夠將最好的和次要的分開.我們若委身給主,便不再有機會,也不願再追求別的東西；因為我們不能又事奉神,又事奉瑪門.我們只能先求神的國和義,而神也會將我們所需用的給我們.委身是毫無條件,毫無保留,也是專一的.
\newline
\newline
今天,很多教會已我去她的本質,沒有委身的愛,所以不能見證基督.神正在尋找肯為祂委身的教會；只有我們肯委身給基督,有委身的愛時,神纔能使用.求神復興祂教會的愛心,肯把自己委身給基督.
\newline
\newline


\newpage

\section{第54屆~港九培靈會~張慕皚~第三講}
\label{sec:538}
\textbf{教會–教會的牧養(一)}
\newline
\newline
連結: \href{https://www.hkbibleconference.org/session-message/view/538}{\texttt{https://www.hkbibleconference.org/session-message/view/538}}
\newline
\newline
\hyperref[sec:537]{< < < PREV SERMON < < <}
~
\hyperref[sec:toc]{[返目錄]}
~
\hyperref[sec:540]{> > > NEXT SERMON > > >}
\newline
\newline
這兩天我們要從約翰福音十章一到十六節,這段經文中,思想教會的牧養.今天我們先看一到十節:「我實實在在的告訴你們,人進羊圈,不從門進去,倒從別處爬進去,那人就是賊,就是強盜.從門進去的,纔是羊的牧人.看門的,就給他開門；羊也聽他的聲音,他按著名叫自己的羊,把羊領出來.既放出自己的羊來,就在前頭走,羊也跟著他,因為認得他的聲音.羊不跟著主人,因為不認得他的聲音,必要逃跑.耶穌將這比喻告訴他們,但他們不明白所說的是甚麼意思.所以耶穌又對他說:我實實在在的告訴你們,我就是羊的門.凡在我以先來的,都是賊,是強盜,羊卻不聽他們.我就是門,凡從我進來的,必然得救,并且出入得草吃.盜賊來,無非要偷竊,殺害,毀壞.我來了,是要叫羊得生命,并且得的更豐盛.」
\newline
\newline
我相信大家都接受一個事實,就是現今我們所處的世代,是一個極其複雜,千變萬化的世代.因為科技的發達,帶來時代急速的變化.在這樣的情況下,很多人感到難以應付現代生活的壓力,包括從環境和從內心的掙扎出來的壓力；所以很多人精神崩潰.為了應付這種需要,政府已經有輔導和社工服務的設施.而站在時代尖端的教會,為了信徒的需要,也必須留意教會牧養的工作.教會的弟兄姊妹需要輔導,需要特別的個人關顧嗎?我相信很需要.老實說,基督徒更容易患上精神病,因為基督徒所受的壓力更大.聖經教導我們要愛神；但有人教導我們要愛世界,有很多相反的拉力在基督徒的生命中,叫我們難以應付.因此,教會需要有好的牧者,像主一樣懂得怎樣牧養羊群.
\newline
\newline
今天很多基督徒的問題,就是不肯承認問題.我們以為做基督徒是沒有問題,沒有生活壓力,沒有重擔的.固然,在聖經理想的角度來看,我們應該把一切問題交給神,過無憂無慮的生活.但事實上我們信主後,卻繼續犯罪軟弱；所以我們必須從神的家中得著牧養.基督徒必須承認教會有問題存在,我們的生命也有問題,對於問題我們要承認,要正視,要拿出來討論；否則我們在碰見問題時,寧可到外面找尋解決,那便可惜了.
\newline
\newline
現在我們從(約十1-10)來思想一個牧者應該具備的條件.一個很必然的事實,就是教會如果有好的牧者,教會的工作必然有好的成績.所以,一個好的牧者頭一樣必須具備的,就是他必須是神所差派的,留意第二節「從門進去的」.他是神所認可的,他有資格可以帶領牧養神的羊群.從門進去的,是神所真正選召的,他不必從別處爬進去,因為他不是賊.今天傳道人必須蒙召,必須蒙差遣,正如羅馬書說:「若沒有奉差遣,怎能傳道呢?」
\newline
\newline
根據聖經的教導,我們知道,一個人蒙神的呼召,首先他必須從神那裡看見他有為神作工的需要,他也會看見時代的需要.神讓他看見這時代需要他,許多失喪的靈魂需要他.當看見需要之時神會讓他有顆願意的心.我們必須有顆願意事奉神的心,等待神的呼召,而不需神用各樣的方法來迫我們.當以賽亞在異象中聽見神在討論,因為那個時代有很大的需要,神能夠差遣誰呢?有誰願意為神而去呢?以賽亞聽了並不是逃避,他乃是說:「主阿,我在這裡,請差遣我.」以賽亞願意去.當然一個人願意去,又有願意被神使用的心,神雖不一定就照他的意願去使用,可能神有另外的崗位要他站立,但先有願意的心是必須的.但願今天有更多華人的弟兄姊妹,願意等待神的隨時差派.不但我們在蒙召的時候,需要清楚知道神的差派,我們在進入工場的時候,也必須同樣清楚神的引領和差派.這樣當我們遇到困難時,便不會退縮,而會尋求神的幫助,求神負責到底.
\newline
\newline
另外,「從門進去的」,是表示傳道人願意遵照神的旨意,來做神的工作.他是神的門生,接受神的教導；他要按神的方法來作神的工.正如耶穌來到世上時他也說:「因為我從天上降下來,不是按自己的意思行,乃是要按那差我來者的意思行.」耶穌到世上,他是從門進去的,他願意完全按照神的旨意行,神的僕人必須受神話語和旨意的約束,不能為所欲為,他必須從門進去.聖經說基督徒要進的門是窄門,要走的路是窄路,而傳道人要經過的門和路,更是窄小.
\newline
\newline
賊和強盜,與從門進去的,是一個很好的對比.賊和強盜是代表我們用自己的聰明和方法去做神的工作.今天,我們在開始事奉神的時候,因為經驗不足,工作生疏,所以我們多禱告,多靠神,但過了一段日子,甚麼已經熟悉的時候,我們便很容易走上賊和強盜的方法去,憑自己的意思和聰明去做神的工作,愈來愈少禱告,愈少根據聖經的原則行事,這是一個很可能的危機.一個好的牧者,一個好的屬靈的領袖,不能夠用自己的方法,一定要完全倚靠神.
\newline
\newline
好牧人第二方面的條件,第三節「看門的,就給他開門.」很多解釋聖經的人,對看門的都有不同的解釋.我個人以為看門的可以是代表聖靈,聖靈是看門的.祂也負責開門,讓人能進到教會來.看門的主要要是給牧人開門,他主要是負責看羊,讓羊進門,也防止強盜進入.今日一個好的牧者,也必須有聖靈的同工,聖靈為他們開門.聖靈要開傳道的門,事奉的門,教會的門.耶穌在馬太福音十六章應許彼得,要將天國的鑰匙給他,這應許在使徒行傳完全實現了.五旬節聖靈降臨,為他們開天國和教會的門,叫猶太人能進入神的國度,加入神的教會.彼得站起來講道,眾人便扎心悔改,歸信基督.是否彼得有能力叫人悔改呢?不,乃是聖靈的工作.彼得只是聖靈所使用的器皿,引人進入神的教會.今天傳道人必須有聖靈的同工,工作纔有效果.教會若沒有聖靈的同工,就如同身體沒有靈魂一樣,是死的.
\newline
\newline
末世的時代,是教會的時代,也是聖靈的時代.教會一開始便有聖靈的工作,我們在使徒行傳可以清楚的看見.初期教會任何一方面的工作,都有聖靈親自的工作,聖靈親自建立教會,叫教會生長,組織教會,差派宣教士,教會有問題時,聖靈親自解決.這是初期教會的持點,完全在聖靈的管制下.今天的教會同樣需要聖靈的工作,我們不要因為太多的爭辯而害怕談聖靈,害怕追求聖靈充滿.聖經說我們必須有聖靈的充滿,有聖靈的膏抹,才能講神的信息.神的工人若要在教會中有美好的工作,必須有聖靈的同工,得著聖靈的充滿,有聖靈的膏抹.因為只有聖靈纔能為我們開傳道的門,開福音的門.
\newline
\newline
好牧人應有的第三個條件,就是要有領袖的恩賜.羊聽從他的聲音,羊跟著他走,這是領導的纔能.一個有領袖恩賜的傳道人應該有兩方面的表現.第一,他是以父母的心腸來牧養羊群.教會中常常有兩種極端的工人樣式.一種牧人是以工人的姿態出現,所謂雇工式的的事奉,這不是好牧人.另一種的極端,就是小教皇,滿有權柄,人人都要聽他,甚至歪曲事實,也要人接受.其實在兩個極端中也有真理,工人是一個服事的人,是透過事奉神來服事人,但一個好的傳道人也需要有權柄,但非小教皇的權柄.有人以為這兩個極端不能溶合,但在保羅身上,我們卻看見見兩個極端的調和.他是以父母的姿態出現,牧養群羊,既有愛心的服事,也有大的權柄.
\newline
\newline
保羅在帖撒羅尼迦前書二章那裡說:「我在你們中間存心溫柔,如同母親乳養自己的孩子.」跟著他又說:「我勸勉你們,安慰你們,囑咐你們,好像父母待自己的兒子一樣.」一個好牧人,應該將愛心的服事和權柄結合,以父母的姿態出現,他應該有愛的服事,來爭取弟兄姐妺的信服,而不是濫施權柄.當羊群得到愛心的照顧時,自然便很順服,肯跟著走,因為他是好牧人.
\newline
\newline
好牧人除了必須以父母的姿態出現之外,他還需要以神的話語來餵養羊群.那就是,傳道人必須好好的按照屬靈的話語講道,教導神的話語,以聖經輔導信徒.教會的權柄是在於神,而神又透過祂的話語來發權柄.當傳道人能好好的講解聖經,用聖經訓勉弟兄姊妹,責備信徒,這傳道人是有權柄的.其實權柄不在於他,乃在於神的話,但他都是神話語的器皿.一個不按神話語教導羊群的牧者,他的權柄是有限的.所以傳道人的權柄和講道是分不開的.
\newline
\newline
好牧人的第四個條件,是能結生命的果子,七到十節.耶穌在這裡形容他自己是生命之門,是得生命的必經之路,不但如此,他還是豐盛生命之門.今天教會所需要的傳道人,是能結生命的果子.第一,傳道人應該以榜樣來結生命的果子.傳道人有美好的生命,用生命來結果子.傳道人必須有謙卑,有好的人格,有好的榜樣,有生活的見證,這些都是生命的流露,其次纔是事奉的恩賜.當然,美好的靈性和恩賜都是必須的.保羅勸勉提摩太說:「你要謹慎自己,和自己的教訓.」這是很需要的.生命是先於教訓,靈性是最重要的,牧者必須追求有好的生命,用生命去結生命的果子.
\newline
\newline
另外,傳道人也需要用神的話去結生命的果子.因為只有神的道,是生命之道,栽種之道.要結好的果子,就必須好好傳講生命之道,純正的道理.
\newline
\newline
總括來說,今天我們看見,一個好的牧者,必須是神所差派的,有聖靈的同工,有領袖的恩賜,能夠結生命的果子.求神在這時代的教會中,興起更多這樣的傳道人.
\newline
\newline


\newpage

\section{第54屆~港九培靈會~張慕皚~第四講}
\label{sec:540}
\textbf{教會的牧養(二)}
\newline
\newline
連結: \href{https://www.hkbibleconference.org/session-message/view/540}{\texttt{https://www.hkbibleconference.org/session-message/view/540}}
\newline
\newline
\hyperref[sec:538]{< < < PREV SERMON < < <}
~
\hyperref[sec:toc]{[返目錄]}
~
\hyperref[sec:541]{> > > NEXT SERMON > > >}
\newline
\newline
今天我們繼續思想教會牧養的問題.請看(約十11-16)「我是好牧人,好牧人為羊捨命.若是雇工,不是牧人,羊也不是他自己的,他看見狼來,就撇下羊逃走,狼抓住羊,趕散羊群.雇工逃走,因他是雇工,並不顧念羊.我是好牧人,我認識我的羊,我的羊也認識我,正如父認識我,我也認識父一樣,並且我為羊捨命.我另外有羊,不是這圈裡的,我必須領他們來,他們也要聽我的聲音,並且要合成一群,歸一個牧人了.」
\newline
\newline
有人說(約十1-16)是一幅圖畫,講一個牧羊人一天的工作.早上他進入羊圈,把羊領出來,帶到草地上,中午在草地上餵養羊群.到了晚上,他把羊群帶領回來羊圈.而在黑夜時分,牧羊人所要著重的,是怎樣看守羊群,他要作防守保護的工作.同樣,教會不單需要有好的牧人,積極牧養神的羊群,同時教會更要留意怎樣防守.如果教會只有積極的牧養,而忽略防守時,就好像有許多漏洞的魚網一樣,雖然網到魚,但很快便溜走掉.因此,教會的防守工作是必須的,今天我們便來思想這個問題,看看有甚麼破壞教會的因素,然後建立防守教會的策略.
\newline
\newline
首先我們看破壞教會的因素.第一是時代黑暗.聖經說末世的時代正是黑暗的時代,黑暗是代表罪惡橫行,需要耶穌基督的真光.這可以說是環境的因素.今天的世代,是沉睡在黑暗中,善惡不分,聖俗不分.正如以賽亞書五章二十節所說:「禍哉,那些稱惡為善,稱善為惡；以暗為光,以光為暗；以苦為甜,以甜為苦的人.」這也是今日時代的一個寫照.
\newline
\newline
今天,有教會的神學家說:我們可以為愛的緣故說謊；即使是殺人,放火,也是對的,只要所做的是出於愛心.難道我們可以用愛心的為詞廢除神的律法嗎?今天,很多社會人士也說,我們要好好對待一些犯罪的人,因為他們是病人.他們犯罪犯法,是因為家庭環境,社會的因素而成,我們應該同情他們,醫治他們,而不能刑罰他們,因為他們是不好的社會制度下的犧牲品.這種善惡不分,聖俗不分風氣很容易流進教會裡；以致我們不知道聖經的道德水準是怎樣.故此,我們要在教會多講罪的問題,因為罪是可怕的事；罪不單在教會以外,更在教會之內.如果我們不追求脫離一切的罪惡,追求過一個聖潔的生活,我相信這會成為教會一個很重要的漏洞.
\newline
\newline
另外,今日人類也沉醉在物質主義和享樂主義之中.特別是香港的信徒,更容易淪落在這個危機中還不自覺.電視和大眾傳播,是今日一個很重要的影響；其實電視就和金錢一樣,是中性的東西,全視乎人怎樣利用.我們要以批判性的眼光來看今日的大眾傳媒,而不要盲目的接受.特別是電視廣告的影響；很容易把人帶進物質主義的道路去.根據統計,香港人每年平均接觸五十萬個廣告,包括電視,報章雜誌,或戶外地方.他們鼓吹物質主義,便消費者成為商品的奴隸,刺激我們的購買慾.另外,電視廣告灌輸給人一種錯誤的觀念和形像.如果我們沒有批判的眼光,便很容易接受.還有,有些電視廣告,其中含有欺騙的成份,更有些是虛構的；我們要小心,要常具屬靈批判的能力.
\newline
\newline
當這些風氣吹進教會的時候,信徒便忘記了主的話說:「你們要先求祂的國和祂的義,這些東西都要加給你們了.」很多人是先追求物質生活的享受,有剩餘的時間纔返教會聚會,試看神的家是何其荒涼,神的國怎能擴展呢?舊約神差派他的子民返國建造聖殿,但後來他們卻忘了建殿的工作,因為他們要建迼自己天花板的房屋.神便興起哈該先知向他們說話,說你們撒的種多；但收的卻少,你們吃卻不得飽,喝卻不得足；穿衣服卻不得暖,得工錢的,卻將錢放在破漏的囊中,神說你們要趕快上山取木料,建造耶和華的殿.跟著神又說:「你們盼望多得,所得的卻少；你們收到家中,我就吹去,這是為甚麼呢?因神的殿荒涼,你們各人只顧自己的房屋.」今天在物質主義的壓力下,我們作基督徒的,需要在神話語的警惕下,檢討自己.今天我為神的家,為牧養羊群,付上甚麼代價呢?
\newline
\newline
還有,今天的生活也沉睡於自信之中.人類覺得不需要神,科學是萬能的,所謂人定勝天.人用自己的方法,能解決環境的問題和自己的問題.人對自己充滿信心,沉睡在自信之中；以致今天我們很難向人傳福音.在這個自信的世代中,很多人做教會工作也太過自信；我們自信自己有恩賜,有組織的能力,有管理教會的方法；我們信自己所學到的一套,因此不能好好的信靠神.
\newline
\newline
破壞教會的另一個因素是有狼要來(十二節).耶穌在(太七15)也說:「你們要防備假先知,他們到你們這裡來,外面披著羊皮,裡面卻是殘暴的狼.」我們作教會工作的,要防備狼的破壞.保羅在(徒二十29)也提醒以弗所的信徒說:「我知道我去之後,必有兇暴的豺狼,進入你們中間,不愛惜羊群.」狼的偽裝是故意的破壞.甚麼是教會的狼群呢?任何藉宗教的名義去進行破壞神教會人或制度,都是狼.
\newline
\newline
在教會裡第一類的狼群,是異端邪教,進入神的教會進行破壞.今天的異端邪教非常盛行,在美國有三千多個以上,而且很多已經在世界各處進軍,香港也開始感受到狼群的侵襲.另一種教會中的狼群,是法利賽人形式的傳統主義者,耶穌是特別責備得嚴厲的.耶穌所反對的,不是好的傳統,而是將傳統的東西高舉；以致高過神的真理,高過基督,這是傳統主義.我們將某些制度神聖化,以致高過神的話,這便是受神所責備的.還有一些狼群,就是那些叫教會分裂的領袖.神的心意是要教會合一,任何叫教會分裂的,都是狼群.
\newline
\newline
第三種破壞神教會的因素,是雇工式的事奉(十二,十三節).這種破壞,很多時是無意的.雇工式的事奉,就是缺乏愛的事奉；當我們在事奉中只為名,為利,而缺乏愛時,這便是雇工式的.作為雇工,有時也有愛,但當愛羊群和自己的利益有抵觸時,他便寧願要自己的利益,而不要羊群了.雇工式的事奉是不能長久的,一有困難便走,擇更好更大利益的工場.雇工式的事奉是只有責任感而無權利感的,只是勉強的應付,而不是存感恩的心牧養羊群.其實這只是雇工式的工作,而不是事奉；因為真正的事奉,是付出而非得著；是建立別人而非自己,是肯付代價而非計較代價,是焚燒自己而非建立自己的.求主在教會剷除這些雇工式的事奉,又喚起更多像主基督的好牧人,就是肯為羊捨命的好牧人.
\newline
\newline
到底我們該怎樣防止呢?第一,要有完全投入的事奉(十一節).我們不顧危險,艱難,甚至不顧性命.這種完全投入的事奉,是在今天艱難的世代中,唯一防守神教會最有效的方法,否則我們無法抵受時代洪流的衝擊.怎樣纔能有效的防守呢?第二,我們要認識我們的羊(十四,十五節),如同父認識我,我也認識父一樣.這是三位一體神緊密的聯繫,互相深入的認識,這是團契生活,生命交流所產生出來的認識.今天的世代,是人際關係大崩潰的時代.我們到教會去只有團聚,而沒有真正的團契,沒有深入的認識,而傳道人對他的羊群也缺乏深入的認識.我們要認識羊群,知道他們的需要；纔能彼此分擔,共同分享,生命纔可以互相的交流和建立.否則今天的教會,就成為漏洞百出的教會了.
\newline
\newline
教會中有兩種不同形式的事奉,一種是以事工計劃為中心的事奉.在一年之初便計劃教會一年的工作,然後分給弟兄姊妹去做,配合我們的事工計劃.但另一種則像耶穌那樣的事奉.耶穌到各城各鄉傳道,醫病趕鬼,他到人群中間去認識人,知道人的需要,然後徹底的幫助人.這并不是說事工計劃不重要,只是我們要以人為中心,先了解他們；然後纔計劃一套事工計劃來符合他們的需要,這是耶穌基督好牧人的事奉.
\newline
\newline
第三個防守的策略,是要懂得招聚羊群(十六節).招聚羊群的頭一個方法是傳福音,這是對付敵人最厲害的武器.願意傳福音的教會,必然是個問題少的教會.今天教會需要積極的傳福音,纔能夠防守.另一個招聚羊群的意思是叫教會合一,這是奠基於生命的合一,和基本真理上的合一,而不是組織上的合一.除非我們有同一的信仰和生命,否則我們無法談合一.我們要尋求教會內部的合一,教會與教會間的合一,進而宗派的合一.
\newline
\newline
關於宗派問題,最少有四種態度.一種是宗派主義的態度,是不對的；那是高舉自己的宗派,過於高舉基督.一種是反宗派的宗派主義,他們反對宗派,但在反宗派中卻自成一派；而且他們所成的宗派色彩,比別的宗派更濃厚.另一種是獨立教會,只要我們不排斥別人,這原是好的,可是這樣卻很難做到教會合一.而且如果獨立教會繁殖得快,人數一多時,又會自成一派,這也不能解決宗派問題.今天我們唯一應走的路線,是健康的宗派態度.那就是互相尊重彼此之間宗派的傳統,體制；在真理上雖然不同的看法,但要尋求在分歧中合一.可以有一同開佈道會,一同為神作工的機會.否則我們便分散神的羊群,把教會弄得四分五裂,這罪是難以承擔的.我們要在神真理的亮光下,極力追求合一.
\newline
\newline
求主興起我們,肯為防守教會的緣故,有完全投入的事奉,有願意認識羊群的事奉,和懂得招聚羊群的事奉.
\newline
\newline


\newpage

\section{第54屆~港九培靈會~張慕皚~第五講}
\label{sec:541}
\textbf{教會的見證}
\newline
\newline
連結: \href{https://www.hkbibleconference.org/session-message/view/541}{\texttt{https://www.hkbibleconference.org/session-message/view/541}}
\newline
\newline
\hyperref[sec:540]{< < < PREV SERMON < < <}
~
\hyperref[sec:toc]{[返目錄]}
~
\hyperref[sec:542]{> > > NEXT SERMON > > >}
\newline
\newline
教會牧養的目的,是使弟兄姊妹能夠起來為主作見證.牧養是教會內部的工作,今天我們要思想教會對外的見證.我們要從馬太福音第五章鹽和光的表像中,思想教會在這世代的見證.(太五13-16)「你們是世上的鹽,鹽若失了味,怎能叫他再鹹呢?以後無用,不過丟在外面,被人踐踏了.你們是世上的光,城造在山上,是不能隱藏的.人點燈,不放在斗底下,是放在燈台上,就照亮一家的人.你們的光也當這樣照在人前,叫他們看見你們的好行為,便將榮耀歸給你們在天上的父.」
\newline
\newline
這幾節經文,是耶穌登山寶訓的一部份.開始的時候,耶穌先講到八福,八福就是基督徒應有的本質,講完品格之後,耶穌便題到基督徒應有的見證和事奉.這是一個很重要的關係.我們信主之後,必須有一定屬靈的品格,有了八福中的基督徒應有的本質後,便能流露出由生命而來的生活和事奉的見證.這是我們第一方面所看見的.
\newline
\newline
鹽和光在聖經中代表很多方面的真理,今天我們特別要思想的,是鹽的內灌作用.那就是說,鹽必須放進湯內,纔能產生作用.而去則有相反作用,光是不能放在斗底下的,必須放在高處,纔能照耀出去,所以光有超然的作用.鹽和光兩方面的見證,正好代表了教會兩方面的見證,我們必須同時有超然和內灌的見證,這也反映了我們的神是一位超然而內灌的神.神不單高高坐在天上,察看世人,神更是內灌的神,他領導人類的歷史,文化；甚至差派他獨生子到世上,成為罪人生活的一部份,這是神的超然和內灌.人是按照神的形像造的,因此在我們的本質中,也應該包含超然和內灌.正如耶穌說我們要內入世,但不屬世.
\newline
\newline
我們從鹽和光也思想到教會,在世上有兩方面的見證.鹽可以說是基督徒社會關懷的見證；而光則代表基督徒在世上福音的見證.教會如果要站在時代的尖端,領導世界的潮流,文化,就必須同時具備這兩種見證.
\newline
\newline
首先我們看教會社會關懷這方面的見證.在鹽的表像中,第一我們看見鹽是一種有強力的質素,鹽的濃度很高,它有強力集中的質素；因此很少的鹽,就能產生很大的影響力.基督徒社會的見證,不在乎人數有多少,而在乎我們的質素是怎樣.如果有強度的基督徒質素,雖然人數很少,卻能在社會中產生極大的影響力.基督徒的質素和記號是甚麼呢?就是那種委身的愛.
\newline
\newline
耶穌在路加福音十四章,也提到跟隨他的人應有的質素,雖然當時有很多人和他同行,但耶穌卻轉過身來,向他們講述跟隨他所必須的條件.首先是愛神勝過一切,勝過自己的父母,妻子,兒女,弟兄,姐妹,和自己的性命.耶穌的意思不是要人從他和家人中任選其一；我們要盡心愛神,也要愛家人.這真理在大多數的情況下,都是沒有衝突的.然而在某些極端的情形下要我們選擇時,我們便要選擇神.另一個跟從主應有的本質,是對人的,那就是要背起十字架來跟從主.背十字架就是捨己為人,我們要忘記自己,處處關心別人,為別人著想.第四個本質是對自己,是要撇下所有的.這樣纔配得跟從主.
\newline
\newline
在講完這一番話之後,耶穌也提到鹽,他說:鹽本是好的,但鹽若失了味,就沒有用.或丟在田裡,或堆在糞裡,都是不合適的,只好丟在外面.基督徒一定要有強力的本質,如果沒有,就像鹽失了味一樣.今天我們所用的鹽是經過化學提煉的,所以不會失味.但在耶穌的時代,鹽是從山邊或湖邊開發出來的；裡面含有許多其他的礦物質,樣子和鹽是一樣.當人買回家,放得久了,便發現這些鹽已經失了味,因為已經失去味,只剩下雜質.今天我們是否也只是鹽的雜質,而沒有味道呢?對於沒有用的鹽,耶穌說要把它丟在外面,被人踐踏,因為已經毫無用處.
\newline
\newline
在鹽的表像中,第二個教訓就是鹽必須投入.耶穌在這方面給我們很好的榜樣,祂不是只坐在天上,祂乃是投入在這黑暗的世界中,為人類死.基督徒若要有社會的影響力,他必須投入社會中.但可惜今天很多人信主愈久,他未信主的朋友便愈少；那就是說,我們愈來愈與世界分離.基督徒有責任投入今天的社會和文化中,領導今天的文化和社會.我們要在自己的職業中事奉神,更要用自己的職業事奉神.我們的職業,是教會和社會一個很重要的橋樑,我們要透過自己的職業,影響社會.我們要把自己的職業事奉化,在職業中尋求事奉主；不像是服事人,乃像是服事主.然而,更要緊的,是我們應該在所做的行業中,產生影響力；特別在一些有影響力的行業中,如教育界,大眾傳播界中發揮影響.
\newline
\newline
第三個教訓是鹽有防腐的作用,這是預防性的社會關懷,社會參與.那就是說,不必等有人在不良的社會環境中受侵害之後,我們纔幫助或解決他的問題.我們的責任,是要防止問題的發生,使人人能夠在社會中安居樂業,這是防腐的作用.教會在社會中第一樣能做的防腐工作,就是努力傳福音.很多人在研究時,往往將傳福音和社會關懷分開,認為不能混淆,但其實傳福音是一種很好的社會關懷,也是第一樣重要的社會關懷.因為所有的社會問題皆由罪起,而只有耶穌的寶血纔能解決人的罪.所以努力傳福音是必須的.
\newline
\newline
另一種的防腐工作,是將聖經所講基督化的家庭推廣出去.家庭制度是神所設立的,也遠在教會之先,它是社會最基本的單位.如果有好的家庭,自然便有好的社會,也有好的教會,好的國家.所以在教會中推行基督化的家庭是必須的,而且要把這種家庭原則向社會推廣.
\newline
\newline
還有一樣能產生防腐作用的,就是要積極辦學,這是預防性的社會關懷.一個好的基督徒老師,對學生的影響是大的；他的重要性,甚至過於教會傳道人,這是值得我們積極去做的.
\newline
\newline
第四方面,是有調味作用.調味就是叫人的生活有味道,有光彩.當人在不良的社會制度下受損害時,會很痛苦.人犯罪軟弱,也帶來人生很大的痛苦.基督徒要做鹽,在世人中間發出光彩,叫他們的生活有味道.耶穌到世上來,他與罪人為友,與罪人吃飯.教會的醫療服務,輔導服務,青年中心,都是為了幫助有需要的人,讓人重新感到味道和光彩.
\newline
\newline
跟著我們思想光彩的見證,就是教會福音的見證.社會關懷只是針對人肉身的需要,但人更重要的是他有靈魂的需要,光的見證是叫人得著基督福音的真光,從罪惡中出來,進神愛子光明的國度.這世代是黑暗的,需要福音的真光,這是我們的責任.光不能放在斗底下,乃要像城造在山上,我們也要將人完全的帶到主面前.
\newline
\newline
光的表像在福音見證上給我們甚麼教訓呢?首先,福音的見證,是真理和聖潔的見證.光在聖經中常表示真理,聖潔.我們不單要傳基督福音的真理,便應有聖潔生活的見證.我們傳福音應該是包括口傳和生活見證,二者是並重的.沒有生活的見證,別人不會相信我們口傳的.
\newline
\newline
第二,福音的見證是多姿多彩的.因為光是由六種顏色所構成,天上的彩虹,就是光的一種作用,是多姿多彩的.今天福音的見證也應該包括各方面的形式,那并不是說要將福音的形式世俗化,以迎合各人的口味.乃是說我們除了口傳外,還須配合別的方式,如幻燈,電影,戲劇,小組音樂,各種不同形式的方法,這些能吸引人來聽福音.還有,文字佈道是非要常重要和有影響力的工具.
\newline
\newline
第三,福音的見證在愈黑暗的時候是愈需要的.在黑暗的世代中,人心更渴慕要得著真光.而且,在愈黑暗的時代,傳福音就愈容易,因為人在罪惡的壓迫和捆綁下,感到需要出路,當有福音傳給他們時,他們便很容易接受.今日的世代,是人最易接受福音的世代,也是屬靈收割最偉大的時代,因為人最容易接受福音.
\newline
\newline
第四,福音的見證,是倒空自己的見證.光是將自己散發出去的,好像蠟燭一樣,是將自己焚燒,犧牲自己,以致叫人得著好處.如果我們要作光,要傳福音,也必須犧牲自己.
\newline
\newline
第五,福音的見證,是世界唯一合適的見證.我們應怎樣面對敵對的世界呢:耶穌說我們要向世界作見證,而且真理的聖靈為主作見證.這是在敵對世代中唯一合適的策略,我們既不能被它同化,也不能逃避,也不能有同樣的敵對,惟有我們向它作福音的見證,有好的生活,忍受逼迫,這是光的見證.
\newline
\newline
第六,福音的見證必須的世界性的見證.光的見證是無孔不入的,能到達很遠的地方,不斷向外散發.今天我們的眼光太過淺窄,只發散在教會內,在教會以外的四週,而不能到達遠的地方.但耶穌說他是世界的光,祂愛世人,祂吸引萬人.耶穌的福音觀念,是廣闊世界性的福音觀念.同樣,我們要遵主的吩咐,使萬族萬民作主的門徒,這責任是世界性的.我們必須要有廣闊的福音眼光.
\newline
\newline
耶穌吩咐我們要在世上作鹽作光,這兩種見證必須結合起來,成為一個的事奉.我們有福音的見證,也有社會的見證,叫神的名在這時代得著最大的榮耀.時代尖端的教會,必須有光有鹽的見證.
\newline
\newline


\newpage

\section{第54屆~港九培靈會~張慕皚~第六講}
\label{sec:542}
\textbf{天國的策略}
\newline
\newline
連結: \href{https://www.hkbibleconference.org/session-message/view/542}{\texttt{https://www.hkbibleconference.org/session-message/view/542}}
\newline
\newline
\hyperref[sec:541]{< < < PREV SERMON < < <}
~
\hyperref[sec:toc]{[返目錄]}
~
\hyperref[sec:543]{> > > NEXT SERMON > > >}
\newline
\newline
(可一14-15)「約翰下監以後,耶穌來到加利利宣傳神的福音說,日期滿了,神的國近了,你們當悔改,信福音.」(21-24)「到了迦百農,耶穌就在安息日進了會堂教訓人.眾人很希奇他的教訓,因為祂教訓他們,正像有權柄的人,不像文士.在會堂裡有一個人,被污鬼附著,他喊叫說,拿撒勒人耶穌,我們與你有甚麼相干,你來滅我們麼?我知道你是誰,乃是神的聖者.耶穌責備他們說,不要作聲,從這個人身上出來吧.污鬼叫那人抽了一陣瘋,大聲喊叫,就出來了.眾人都驚奇,以致彼此對問說,這是甚麼事,是個新道理呵.祂用權柄吩咐污鬼,連污鬼也聽了他.耶穌的名聲,就傳遍了加利利的四方.他們一出會堂,同著雅各約翰,進了西門和安得烈的家西門的岳母,正害熱病躺著,就有人告訴耶穌.耶穌進前拉著她的手,扶她起來,熱就退了,她就服事他們.天晚日落的時候,有人帶著一切害病的,和被鬼附的,來到耶穌跟前.全城人都聚集在門前.耶穌治好了許多害各樣病的人,又趕出許多鬼,不許鬼說話,因為鬼認識他.」
\newline
\newline
聖經一開始就告訴我們,神是創造天地萬物的神,但到創世記第三章時,我們看見撒旦試探亞當夏娃,叫他們犯罪離開神,因此罪便入了世界.其實罪不是從這裡開始進入天地之間,罪乃是從撒旦開始,然後進到人類中間.因為聖經說全世都臥在那惡者手下,神暫時讓撒旦成為天空的掌權者,成為黑暗世代的掌權者.但是神不讓人類永遠落在撒旦手中,在黑暗裡過活,因此神為人預備救恩.
\newline
\newline
耶穌到世上來,傳的是天國的福音.甚麼是天國呢?那並不是一塊地土.天國乃是神的統治,凡願意讓神來統治的人,神的國就在那裡,所以神的國可以在人心內,也可以在一處地方.所以天國的意思,就是神的掌權.神所最關心的,是祂國度的建立和擴展.願你的國降臨,意思就是願神的旨意行在地上,如同行在天上,這是希伯來人寫詩時的平衡體.耶穌來是為建立教會,神在創世之前已經揀選我們,叫我們成為聖潔,合乎主用.所以在創世之前,神已經為我們預備教會,耶穌來也是要建立教會.但教會只是天國的器皿,所以今天教會不應只有教會的策略,而應有國度的策略.
\newline
\newline
教會的策略,只是顧及教會內部的需要,崇拜,教導,團契等聚會,即使是向外發展,也是有限的.這只是教會的策略.但國度的策略卻是廣闊的策略,天國的眼光是廣闊的眼光.不單顧念自己教會的需要,更是顧念到整個世界的需要.我們要把神的國度擴展到天地之間.今天,神不但興起教會,成為國度主要的器皿；神更興起很多機構,成為教會的膀臂,這膀臂要伸延出去,將神的國度實現在黑暗的世代中.在香港,就有超過三十個屬靈的機構,福音的機構.我們需要有國度的眼光,心懷,策略,將神的國擴展.
\newline
\newline
我們要從耶穌的身上,看看祂怎樣推行這國度的策略.一開始的時候,他便宣傳神的福音,叫人悔改,後來又教訓人,宣講天國的道理.因此我們看見,國度的策略,第一,是宣告天國的憲法.
\newline
\newline
第二方面的策略,(23-28節)講到耶穌怎樣趕鬼,對付人類最大的敵人,這可以說是天國策略的外交和國防.一個國家的建立,除了要有好的憲法之外,更需要好的外交和國防.
\newline
\newline
第三個策略,(29-34節)講到耶穌醫好西門的岳母,并許多患病的人,這可以說是天國策略的內政和福利.任何一個國家都必須有這三方面的策略,這國家纔得穩固.
\newline
\newline
這三方面的策略,可以說是四福音耶穌言行錄的總結.耶穌一生所作,就是這三方面的工作,講道,趕鬼,醫病.
\newline
\newline
我們詳細來看看這幾個天國的策略.首先我們思想第一個策略,就是天國憲法的宣告.很多時候,我們以耶穌所講的登山寶訓為天國的憲法,是教會應有的行為準則.國家必須有憲法,天國的憲法可以包含兩方面:就是傳福音和教導,二者是分不開的.如果只有傳福音而沒有栽培,那還未履行到耶穌的大使命,因為主的使命是使萬民作祂的門徒,而門徒一詞,已經將傳福音和教導這兩個意思結合起來.我們要引人歸向耶穌,而且要使他們一生作主的門徒.傳福音加上栽培,是一個使命而已.
\newline
\newline
耶穌郅世上來,是要為神設立教會.我們看祂在約翰福音十七章分離的禱告中對神說:我在地上已經榮耀你,你所託付我的事,我已經成全了.當時耶穌還未上十字架,但耶穌說已經完成了神所託付他的事,這是指甚麼呢?耶穌在三年多的工作中,最大建樹是甚麼呢,就是訓練們徒.因此耶穌所指已經完成的事,就是指訓練門徒說的,因為跟著耶穌便為門徒禱告.而後來五旬節聖靈降臨之後,教會因著門徒的力量和忠心而得以繼續和擴展下去.因此,門徒訓練,實在是教會擴展天國的一個重要策略.
\newline
\newline
我們留意耶穌訓練門徒的方法.首先是耶穌所選的對象.在十二個門徒中,有漁夫,有像奮銳黨那些激進的份子,和不同種類的人.他們在社會中的地位并不高,而且是些年輕伙子,他們沒有受過甚麼宗教教育和訓練,也不是知名人士.但耶穌選上了他們；原因是他們是一班可以塑造的人,他們有熱心,肯順服,從來沒有反對過主,所以主選上他們.今天神也要在教會中興起門徒訓練的工作,神要揀選的,正是當日像門徒那樣的人,肯過順服的生活,有熱心,能被神塑造的.
\newline
\newline
第二,我們看耶穌訓練門徒的方法.很多時候,耶穌向群眾講道,給他們很好的教訓,因此當日跟隨耶穌的人並不少.在群眾中,有七十個是被主特別訓練,兩個兩個差派出去工作的.但耶穌近身的門徒卻只有十二個,是他經過整夜禱告揀選出來的,而在十二個門徒中,又有三個是特別與主親近的.這是耶穌訓門徒的模式和方法.在訓練門徒的過程中,我們無法給予每人相同的時間和訓練,但我們可以像主一樣,用生命悉心的栽培一班人,是肯熱心事奉,又能順服的；給他們特別的操練處,這纔是門徒訓練的方法.
\newline
\newline
第三,我們看見耶穌訓練門徒的內容,是包括深度和闊度的.栽培是成聖的功夫,叫人成聖,過成聖的生活,這種操練不但有深度,還有寬度.耶穌在馬太第五章最後一節說,所以你們要完全,像你們的天父完全一樣.這不單是我們追求成聖的最高標準,更是成聖的內容,我們不但需要個人追求深度,更要顧及寬度,怎樣對家庭,和在教會中的關係,社會道德等等,這些都是我們成聖功夫的一部份.如果教會要站在時代的尖端,便必須同時留意我們屬靈生活的深度和寬度.我們不單要關心自己教會的進展,更要關心一切屬靈的機構,神學院,福音廣播,文字工作,宣教團體等等,因為他們是教會的膀臂,能把天國憲法的道理伸展出去.
\newline
\newline
跟著我們思想天國的外交和國防.耶穌到世上來一個很主要的目的,就是要擊傷撒旦,正如創世記三章早就說:女人的後裔要傷蛇的頭,蛇要傷牠的腳跟.今天,撒旦已經被耶穌擊傷,然後牠的尾巴仍在掙扎活動；撒旦今天在世上仍有很厲害的破壞.因此教會今天要留意撒旦在世上的工作.
\newline
\newline
撒旦有兩方面的策略,首先牠叫人以自己來代替神的地位,其實這是撒旦由始至終的詭計.撒旦的墮落,是由於不肯榮耀神,要以自己取代神.人類的墮落,原因也是一樣.甚至教會的軟弱,也是這樣.如果教會倚靠神,必然注重禱告,因為每一個禱告都是表示對神的倚靠.多禱告,是表示多倚靠神；少禱告,便表示多注重自己,少靠神.今天撒旦也用盡各樣方式去擱阻人禱告,這是我們要小心防避的.
\newline
\newline
撒旦第二個策略,是叫人懷疑神的話.在引誘亞當夏娃,和試探耶穌的過程中,撒旦都是用這武器,曲解聖經,誤用神的話,和叫人疑惑神的話.但教會今天所需要的,正是神的話.神的話是教會的權柄,惟有我們在神話語的權柄下過活時；纔能抵擋撒旦對教會的攻擊.如果我們能堵塞著破口,不讓撒旦進攻,教會纔能被神使用.
\newline
\newline
最後我們思想天國的內政和福利.耶穌不單傳道,叫人的心靈得著復興,罪得赦免,耶穌也關心人肉體的需要；祂為人醫病,這是全人的關顧.今天教會的策略,也要關懷全人各方面的需要.我們不單要關心弟兄姊妹的教會生活,更要關心他們在職業上的生活,這纔是全人的關懷.除了關心弟兄姊妹在職業上的生活和見證外,還有弟兄姊妹休閒的時間,也是我們需要關心的.教會應該提供一些健康有益的活動,讓休閒的青少年人,老年人參與,這是天國的內政和福利.
\newline
\newline
教會應該有國度的策略,我們應該和其他教會,其他宗派聯合,更要和其他機構聯合,同心合意的將天國的策略推廣出去.
\newline
\newline


\newpage

\section{第54屆~港九培靈會~張慕皚~第七講}
\label{sec:543}
\textbf{教會–聖靈的居所}
\newline
\newline
連結: \href{https://www.hkbibleconference.org/session-message/view/543}{\texttt{https://www.hkbibleconference.org/session-message/view/543}}
\newline
\newline
\hyperref[sec:542]{< < < PREV SERMON < < <}
~
\hyperref[sec:toc]{[返目錄]}
~
\hyperref[sec:544]{> > > NEXT SERMON > > >}
\newline
\newline
(弗二20-22)「并且被建造在使徒和先知的根基上,有基督耶穌自己為房角石；各房靠祂聯絡得合式,漸漸成為主的聖殿.你們也靠祂同被建造成為神藉著聖靈居住的所在.」
\newline
\newline
今天我們要從聖殿這表像中,思想教會是個不斷成長的教會；是個在建造中的教會,是漸漸而成為神的聖殿.教會是在建造中,不能有固定的形式.教會若要站在時代的尖端,應該要不斷的被建造,直至能成為聖靈的居所.一個神所使用的教會,是在成長中的教會.
\newline
\newline
要達到不斷的更新,不斷成長的目的,有三件事我們必須留意的.第一是聖經的原則.任何教會要按時代的改變而改進,必須根據聖經的原則；任何與聖經的原則相違反的,一定不能改.第二是明白時代的需要.到底時代需要怎樣的教會呢?第三是聖靈的風怎樣吹.聖靈到底要在這時代怎樣帶領他的教會呢?我們要根據聖經的原則迎合時代的需要.我們要留意時代急劇的改變,那并不是說可以胡亂改變以迎合時代；而是要在時代中滿足現在人的需要,帶領他們歸向主.就如醫生雖然有良藥,但也必須了解病情,否則怎能醫好病人呢?在建造中的教會,絕對不能自滿,我們應該謙卑學習,根據聖靈的帶領,在時代中改進,以致能站在時代的尖端,被神使用.
\newline
\newline
今天的經文,讓我們看見幾方面的教訓.第一,二十節「被建造」,這是個很重要的真理.教會不是在那裡建立自己,傳道人,長執,屬靈領袖也不是在建立教會,時代尖端的教會,乃是個「被建造」的教會.誰在建造呢?是主耶穌自己.教會是被耶穌基督所建造中的.我們不要像老底嘉教會.有人形容老底嘉教會是末世教會,那是一個憑自己能力去建造的教會.
\newline
\newline
主耶穌在(太十六18)說:「我要將我的教會建造在這磐石上,陰間的權柄,不能勝過他.」這句話有很深刻的屬靈意思.耶穌說我要把我的教會建造在這磐石上,在這裡我們要留意幾方面的意思.一個教會如果是基督在那裡建造,將會是個蒙愛的教會,是主所心愛的教會,是主的新婦.今天我們的教會是否蒙主所愛呢?第二,主自己所建造的教會,是根基穩固的教會.耶穌要把教會建造在磐石上,那就是在一班肯承認耶穌是神的兒子,肯接受耶穌的主權的基督徒身上.第三,主所建造的,是有能力的教會.耶穌說陰間的權柄,不能勝過他.今天的教會,是否有傳福音的能力,有栽培的能力,有社會見證的能力,有奉獻的能力,有禱告的能力,有弟兄姊妹委身的能力呢?如果是主所建立的,必然是有能力的教會.最後,主建立的教會,是有主同在的.主所建立的,主必在其中,就像使徒行傳初期的教會,主在他們當中,特別透過聖靈的工作,在各方面帶領他的教會,今天我們有否看見主在我們當中,建立他強而有力的教會呢?
\newline
\newline
何以知道會有能力呢?有兩個很明顯的表現,一個是禱告會.禱告是能力的象徵,是得能力的源頭,教會注重禱告,表明他是個有能力的教會.另一個表現就是我們不必為教會的前途擔心,因為有主的同在和能力的覆庇,教會必然大步向前,這是我們可以肯定的.
\newline
\newline
第二方面的教訓,二十節「並且被建造在使徒和先知的根基上」.一個長進中的教會,是以使徒和先知為根基的教會.使徒和先知,是領受神的啟示,將神的啟示錄下,成為今天我們手上所有的聖經.具體來說,就是時代尖端的教會,是建立在神話語的權柄上.聖經是教會信仰和生活的根基,以及最高的權威.一切的傳統,創新,都必須符合聖經的原則,不能高抬傳統過於聖經的原則,也不能因為時代轉變,而以為聖經不合時宜,把它擱置.時代的教會必須建立在聖經的權威上.
\newline
\newline
不過,單是相信聖經的權威,并不足以建立時代的教會,我們還需要有正確的解經.很多人說是相信聖經的權威,但卻不按正意分解真理的道,只是用聖經來配合自己的論調,把宗教的權威從神身上轉移到自己身上,如此他只是利用聖經而已.因此,我們接受聖經的權威,同時也必須配合正確的釋經法.另外我們也要分辨甚麼是基督的真理,甚麼是非基要的真理.保羅在羅馬書十四章內提出很多非基要的真理,譬如說吃肉或吃素,各人的看法并不一致,我們并不要為這些爭論.但是對於基要的真理,則絕對不能妥協.譬如聖經是絕對的權威,耶穌是神的兒子,基督有完全的神聖,基督為童貞女所生等等；這些都是基要的真理,是不能妥協的.
\newline
\newline
在解經的時候,除了要有正確的解經法之外,還需要聖靈的光照.聖靈不單默示聖經,聖靈更開人心竅,叫人明白神的話.今天我們要在讀經中有所領受,必須有聖靈的照和工作.
\newline
\newline
另外,聖經的權威還需要配合全面性的教導.很多人只著重某一部份的真理,譬如說他著重四福音,便只講四福音的道理,而對聖經及社會關懷的道理置之不顧,這樣便無形中刪減了聖經的話語.所以我們雖然接受聖經的權威,同時我們也要教導全面性的真理.
\newline
\newline
還有,聖經的權威也必須有行道來配合.因為雅各書說撒旦魔鬼也相信神,但牠的信卻是戰驚的.撒旦熟悉聖經的真理,并相信它的權威,但卻沒有行為.今天如果我們只信聖經而不行道的話,聖經的權威便大受削弱.現代派和理性派用理性和科學來推翻聖經的權威,新正統派的思想,又以宗教經驗來推翻聖經的權威,對於自己有感動的經文,便承認它是神的話,否則便是人的話.甚至在福音派和保守派中間,也有人只相信聖經是局部無謬誤的,他們說聖經有科學上的謬誤,不過這倒無關痛癢,因為聖經的目的是關乎救恩,而救恩的真理是無謬誤的.對於這些,我們不能贊同.聖經是完全無誤的,我們必須像耶穌那樣,接受聖經全部的權威.
\newline
\newline
第三方面的原則和教訓,二十節「有基督耶穌自己為房角石,各房靠他聯絡得合式」.第三個長進中教會的原則,是以耶穌基督為房角石.房角石是一個標準,是建屋前所立的一塊石.以後就成為建屋的標準,和房屋能連結起來的標準.那就是說,耶穌基督是教會增長的標準,也是教會合一的標準.今天教會要增長,必須以耶穌為房角石.我們要效法耶穌那樣傳福音,到各處尋找人,而不是在教會內等人來.耶穌要教會增長之前,是先訓練門徒,門徒有口傳和生命的見證,便很容易吸引人到主前來.在合一方面,耶穌也是房角石.有愛便能合一,因為愛有連結的能力；相反地,有罪便會分裂,因為罪有破壞的力量.教會內部要先合一,然後教會與教會之間,宗派與宗派之間,也要找尋合一.
\newline
\newline
房角石也是解決教會問題的標準.哥林多教會有很多問題,保羅寫信去為他們解決；在開頭十二節之內,他一共提了十二次耶穌的名字.在保羅的思想中,是耶穌基督,惟有祂纔能解決教會的問題.今天當教會有問題時我們要以耶穌的標準去解決.
\newline
\newline
第四方面的教訓,二十一節「各房靠他聯絡得合式,漸漸成為主的聖殿.你們也靠他同被建造,成為神藉著聖靈居住的所在.」教會要建立的目標,是聖靈的居所.我們要為神建造他能夠居住的所在,這裡有很重要的屬靈意思.很多時候,教會或機構的領導人,他們努力的目的,是為建立自己的王國.我們往往所建立的,只是自己的居所；但聖經要求的,是我們要建立神的居所.教會是神的,不是我們的.
\newline
\newline
教會是聖靈的居所,這讓我們在事奉中有高貴感,我們所建造的,是聖靈的殿,何等高貴.很多人看輕教會的事奉,以為這只是一些瑣碎的雜務,其實事奉是極高貴的,因為教會是聖靈的居所.
\newline
\newline
另一方面,事奉應該有聖潔感.神的殿是聖潔的,以利的兩個兒子在神殿裡犯罪,神很嚴厲的刑罰他們.當我們建造神的殿,在教會事奉神的時候,一定要有聖潔感.神所要一班事奉祂的人,是聖潔的.
\newline
\newline
還有,事奉應該有嚴重感.我們不能馬馬虎虎的去事奉,因為建造的是神的殿.保羅說:若有人毀壞神的殿,神必要毀壞那人,因為神的殿是聖的.我們必須帶有嚴重感來事奉,這纔是對的.很多人事奉神過於隨便,我們為了自己的職業讀了很多書,還一面做一面進修；同樣,我們肯否為了教會的事奉而好好進修,接受各方面的裝備和操練.
\newline
\newline
最後,神的殿是神榮光充滿的地方.那就是說,我們在神的殿中事奉的時候,必須尋求神的榮耀.不是為我個人的好處,只要神在教會中得著榮耀,只要神的榮光是充滿這殿,便心滿意足了.聖經說是漸漸成為主的殿,意思就是我們不能滿足於教會的現狀,乃要按照聖靈的帶領,不斷改進,除去罪惡和雜質；追求更新,以致能站在時代的尖端,被主使用.
\newline
\newline
教會在建造之中,雖然今天有許多未盡善的地方,但千萬不要灰心；因為教會是在漸進之中,將來必在神的帶領下,成為榮美的殿.
\newline
\newline


\newpage

\section{第54屆~港九培靈會~張慕皚~第八講}
\label{sec:544}
\textbf{教會的使命}
\newline
\newline
連結: \href{https://www.hkbibleconference.org/session-message/view/544}{\texttt{https://www.hkbibleconference.org/session-message/view/544}}
\newline
\newline
\hyperref[sec:543]{< < < PREV SERMON < < <}
~
\hyperref[sec:toc]{[返目錄]}
~
\hyperref[sec:528]{> > > NEXT SERMON > > >}
\newline
\newline


\newpage

\section{第54屆~港九培靈會~馬有藻~第一講}
\label{sec:528}
\textbf{馬太福音導論}
\newline
\newline
連結: \href{https://www.hkbibleconference.org/session-message/view/528}{\texttt{https://www.hkbibleconference.org/session-message/view/528}}
\newline
\newline
\hyperref[sec:544]{< < < PREV SERMON < < <}
~
\hyperref[sec:toc]{[返目錄]}
~
\hyperref[sec:529]{> > > NEXT SERMON > > >}
\newline
\newline
兄弟非常感謝主!可以再回到香港來,我在港唸完中學,到美國唸大學的時候纔信主,又在唸研究院的時候纔奉獻做傳道,繼續讀神學,學習牧養教會.我大部份時間,都在德州達拉斯市工作.我在牧會時,已購置很多培靈講道集閱讀.心中很嚮往,渴望有一天能參加港九研經培靈大會.真想不到第一次來參加,便要負責做講員.我對香港的教會認識不深,這次是帶著一個很戰兢的心情來講道,請多為我代禱!
\newline
\newline
今天我們所要研究的,是馬太福音的重要,歷史的背景,目的和主題.我要從四方面來發揮.
\newline
\newline
馬太福音的重要
\newline
\newline
馬太福音不單是新約的第一卷,在新約中佔一重要地位.舊約瑪拉基書預言彌賽亞的先鋒來臨,新約第一卷書就記載他來臨的應驗.所以馬太可說是舊約的續集,是舊約和新約之間的橋樑,也是新約其它書卷的引言.如果我們明白神在世上全盤的計劃,我們就必須讀馬太福音.因為除了馬太書之外,新約其它書卷主要是講及教會的真理,信徒在教會內外生活和事奉的教訓；但是關於教會的開始,存在,目的,將來,都是先從馬太的信息下手,纔會比較清楚.
\newline
\newline
美國著名解經家J.D.McGee說:「新約書卷在次序編排,是經很嚴謹,敬虔和有學問的人慎思明辨編排纔完成.他們將馬太放在廿七卷之首位,可見是非常看重馬太,而勝過其它的書卷.」連反對基督教哲學很出名的哲學家雷那也曾說:「馬太福音,是在基督教中最重要的書卷,沒有這本書,基督教就不堪一擊了.」由此可見本書的重要了.
\newline
\newline
兄弟不是馬太專家,但我對馬太有很大興趣,神賜我三個孩子,大的叫馬太,二女的名字是馬太的女性名詞,第三個兒子是用菲律賓翻譯的馬太,所以我三個兒女皆叫馬太,連我的妻子也叫馬太,可見我是很愛馬太福音,因為這卷書是相當重要之故.
\newline
\newline
馬太福音的歷史背景
\newline
\newline
我認為在這裡先看其歷史背景,對我們很有幫助,因為我們已明白它的重要性,再看其歷史背景,和看其目的,便較為順序合理.馬太的名義是「神的禮物」,馬太又名利未,利未是「聯合」之意.可能馬太是迦百農人,他是猶太人所憎惡的稅吏；歸主後,纔成為主有力的見證人,常常向他的同僚稅吏傳福音.
\newline
\newline
當耶路撒冷教會建立後,使徒便成為教會最重要信仰的權威.凡是有關主生平的問題,使徒可給予官方式的答案.在這時候,他們感覺沒有需要記載主的言行錄,當有問題的時候,便去到使徒集合處尋找答案.但後來教會受逼迫,使徒們便開始注意如果教會被逼關門,或使徒們囚在監裡,教會便無人供給官方教訓來取代他們的權威.特別是當耶路撒冷領袖之一的司提反殉道,又加上(徒十二章)記載希律安提阿帕一世在主後四十四年殘害教會；因這些事實,他們感覺需要將主耶穌的生平記錄下來,特別是要趁使徒還存在的時候.馬太是四卷福音書中最配合歷史的一卷書.我本人相信馬太是當時教會第一本官方記錄,因為當時教會主要成員是猶太人.(徒十五章)論初期的教會開始歡迎外人,於是路加特別寫書給外邦人看,而約翰則是寫書給猶太人和外邦人看.可能當時耶路撒冷教會需要一本專給猶太人看的,有關主言行的官方記錄書.馬太本身是耶路撒冷教會的首腦之一,也是十二使徒之一；他跟隨主多年,在使徒中是學問最淵博的；如果教會要寫一本有關主的書籍,相信這職事非馬太莫屬.而且馬太的主題也是教會所關注的,因為福音先要傳給猶太人,然後傳給外邦人,大概這就是馬太的歷史背景了.
\newline
\newline
馬太福音的目的
\newline
\newline
每一本書都有它寫書的特別目的,馬太當然不例外.福音書要藉引證和辨明,來達到它的目的；所以先佈下大局和全書的架構,然後旁徵博引,使讀者容易明白本書的論題.這種按照心中藍圖來完成的方法,是新約很多書卷的特色.前芝加哥神學院院長E.D.Rechen說:「第一世紀文學特色與二十世紀的並不相似；通常二十紀在書中序言中,或在結論時,已清楚指出寫書的目的.但馬太是按照當時第一世紀的習慣,將寫書的特色,放在書的內容裡面,所以讀者要留心尋求,就能獲益.」從馬太的文學修養,句語,用字,主題,主旨和舊約的引用,層次都是很鮮明的.全書引用舊約共一二九次,可見馬太的主要對象是為猶太人；他知道猶太人所關注的是甚麼問題,這是馬太寫書所注重的內容.
\newline
\newline
猶太人處在主降臨的時候,特別關心到天國的降臨；這是神應許給大衛建立的永遠國度.但天國在何時建立和建立的條件是甚麼?這些答案都可以在本書中找到.馬太福音的重要性,全在乎主題的發揮.但若要明白這書的主題,就必須要逐步來追溯,這是閱讀馬太最大的困難.在謹慎探討下,我們可以看見這書的目的.他將這書的目的撒播在不同章節中.我自己看馬太的時候,看見至少有五個目的；而每個目的都是逐漸發揮下去,再挑選與主題有關的神蹟,來證明其主題.
\newline
\newline
(一)證明耶穌的來臨為了應驗舊約所應許的彌賽亞.神將彌賽亞賜給以色列,透過以色列使全地得福.
\newline
\newline
(二)指出耶穌的降臨,為要在地上建立神在舊約所應許給大衛的國度.(撒下七1-17),神透過先知拿單向大衛說話,這段經文指出:「大衛之約」,「家」,「國」,「位」,這段思想再出現在(路一26-33).
\newline
\newline
(三)解釋耶穌在地上建立國度的經過,及如何遭拒絕.
\newline
\newline
(四)主耶穌不能在當時建立這國度,要待他第二次再臨纔可以建立.(太廿四,廿五章)「橄欖山論題」,說明主第二次再來的目的.
\newline
\newline
(五)在主未建立國度之前,在這兩次降臨的過渡時期；說明主給予門徒的使命是甚麼?神的計劃是甚麼.
\newline
\newline
在這兩次降臨時間,神就用教會取代以色列的使命.以色列因為棄絕彌賽亞,所以救恩從他們手上轉移至外邦人去.我們今日教會就承接這偉大使命,要將福音傳遍地極.直到時候滿足,主就會再來；這是指外邦人的數目添滿了,而以色列亦全家可得救.
\newline
\newline
馬太福音的大綱
\newline
\newline
按照我們所提的,可將馬太福音分作六個分題:
\newline
\newline
(一)(太一至四11)以色列的王來了(包括耶穌的家譜,降生,早年軼事,開始工作.)
\newline
\newline
(二)(太四12至十章)以色列的國度來了(主的工作,教訓.)主用說話解釋他來臨的目的,用神蹟支持他的工作及信息.
\newline
\newline
(三)(太十一至十六章)以色列的王被棄絕了.作者逐步追尋為何彌賽亞,被當時的人及國家領袖棄絕.先提及先鋒被棄絕,再提及主被棄絕；主便宣佈再臨之間的計劃,稱為天國的奧秘.
\newline
\newline
(四)(太十七至廿五章)以色列的國度被奪去了；主怎樣訓練門徒,勇敢地進入耶路撒冷,直至受苦前的晚上.
\newline
\newline
(五)(太廿六至廿七章)以色列王被奪去了,主耶穌怎樣受苦釘在十字架上.
\newline
\newline
(六)(太廿八章)以色列王回天國了,主耶穌升天回父那裡去.
\newline
\newline
我們的主雖在天上,但祂的使命在我們當中,祂是萬王之王,也是我們生命之王,所以我們在等候主再來的時候,要在行為和生活顯明主在我們心中.
\newline
\newline


\newpage

\section{第54屆~港九培靈會~馬有藻~第二講}
\label{sec:529}
\textbf{以色列的王來了}
\newline
\newline
連結: \href{https://www.hkbibleconference.org/session-message/view/529}{\texttt{https://www.hkbibleconference.org/session-message/view/529}}
\newline
\newline
\hyperref[sec:528]{< < < PREV SERMON < < <}
~
\hyperref[sec:toc]{[返目錄]}
~
\hyperref[sec:530]{> > > NEXT SERMON > > >}
\newline
\newline


\newpage

\section{第54屆~港九培靈會~馬有藻~第三講}
\label{sec:530}
\textbf{以色列的王來了(續講一)}
\newline
\newline
連結: \href{https://www.hkbibleconference.org/session-message/view/530}{\texttt{https://www.hkbibleconference.org/session-message/view/530}}
\newline
\newline
\hyperref[sec:529]{< < < PREV SERMON < < <}
~
\hyperref[sec:toc]{[返目錄]}
~
\hyperref[sec:531]{> > > NEXT SERMON > > >}
\newline
\newline
昨天我們講到以色列的王來了.作者幾方面介紹以色列的王怎樣來到,因時間關係,今天續講(三1-四11).從第二章看見猶太人拒絕主,但主卻被外邦人所接受；作者很快跨越耶穌頭三十年時間,急於將主獻身,獻國的事介紹出來.這一段可分三部份來看.
\newline
\newline
基督的準備
\newline
\newline
基督的先鋒(三1-12)
\newline
\newline
這裡看見基督的先鋒,是用舊約的服裝,身份和權柄,來預備人的心.他一出來便宣稱:「天國近了,你們應當悔改.」天國究竟指甚麼?許多人有不同意見,首先我們看天國的意義；天國就是以天上的權柄來管理的國度,這天國可以是屬靈的或屬地的,在乎講的和聽的,彼此的了解是甚麼.既然宣講天國的是施洗約翰,他是舊約最後的先知,聽眾也是活在舊約時代的人,所以要從舊約來尋找天國的意義.神透過先知拿單將天國的應許賜給大衛王,(撒下七14-16),這應許慢慢成為舊約聖徒的盼望.透過先知,僕人等,堅定神向以色列國應許國度的來臨.所羅門之約,大,小先知,(詩八十九篇)等國度的應許,(賽二11,六十五章,彌五章,番三章,珥二章),神在以色列人歷史中,透過先知們不斷向他們堅立這約.所以「大衛的國度」成為他們熱切盼望實現的一個盼望；也是他們恆久禱告的中心.雖然國家滅亡,但在歸國時期,他們也沒有忘記這大衛的國度應許啊.在兩約之間時期,他們的心志更熾烈.兩約的文學活潑地將這天國的盼望表達出來,所以這天國從施洗約翰口中說出,他們便聯想到神應許大衛的國度.當約翰認識耶穌後,他更加努力向當時的人宣講,因為天國的王已經來到了.
\newline
\newline
天國雖受地理環境所界限,如同地上的國度；但不是失去屬靈的要求,進入這國度必須按照神對人屬靈的要求,纔可進去.人必須要悔改,這是最基本屬靈要求；人必須要重生纔可進入神的國.所以施洗約翰指出人必須悔改.在舊約中多處預告,沒有悔改的,不單不能進入天國,反有審判.施洗約翰強調的悔改水禮,靠著這水禮；看出當時的人有沒有決志.這個水禮稱為赦罪的水禮.當他們接受施洗約翰的水禮,便接納他的信息；等候天國來臨,進入國度.(徒十九3-4)幫助我們明白這水禮的意義.
\newline
\newline
基督的水禮(13-17節)
\newline
\newline
主正式開始工作前,來到約但河,承認先鋒的信息,接受他的水禮.基督的水禮引起各種不同的解釋,首先我們要明白水禮的意義.水禮基本上是「認同」的行動,一個人接受水禮,便是認同施水禮傳信息的人.施洗約翰拒絕耶穌受洗,因為他以為彌賽亞是不需要悔改的；但耶穌堅持要接受水禮.從(三15)祂自己的解釋,祂認同施洗約翰的信息:「你們應該悔改」.他要盡上諸般的義,他要有義人的名銜；纔出來工作,所以要接受水禮.否則不能被人接納.「暫且許我」,指出有時間因素,目前需要如此行,可見耶穌一生遵守律法要求.在接受水禮的過程中,神的靈便降臨主耶穌身上,應驗(賽十一章)的預言,「天上有聲音說,這是我的愛子,我所喜悅的,」應驗(詩二篇和賽四十二章)的話,這幾段都是彌賽亞預言,印證耶穌是彌賽亞；所以後來門徒約翰記載耶穌接受水禮「顯明給以色列人」(約一31).故此主的水禮有三方面的意義:
\newline
\newline
一,認同施洗約翰的信息.
\newline
\newline
二,與本國的人認同,成為神國中的一份子.
\newline
\newline
三,祂是神差遣而來的彌賽亞.
\newline
\newline
基督的試探(四1-11)
\newline
\newline
從第一至第三章,耶穌在各方面成為彌賽亞俱備十足,祂成為以色列的王是毫無疑問的.但還有一重要試驗或試探在後頭.除三一真神外,最大的力量由魔鬼而來.如果耶穌能經得起魔鬼的試探,祂便可以成為彌賽亞,(四1)指耶穌接受試探,是經過聖靈的催引,可見得到神的特許和計劃,要考驗祂的愛子能否成為以色列的彌賽亞.耶穌在三方面接受試探:石頭變成食物
\newline
\newline
魔鬼對祂說:「你若是神的兒子」,「若」在文法上,可作「既」,既是神的兒子,魔鬼承認祂有彌賽亞的權柄；但指出祂可用自己的權柄,來滿足自己肉體的需要.
\newline
\newline
從殿頂咷下
\newline
\newline
按照猶太人傳統,彌賽亞會出現在聖殿最高的地方,這是他出現的記號.魔鬼叫耶穌跳下,然後神會保護祂.
\newline
\newline
崇拜魔鬼得全地
\newline
\newline
神早已答允將全地賜予彌賽亞,但神給祂愛子全地的權柄,是要經過十字架道路.現在撒旦要引誘主,使祂不需要上十字架也可得冠冕,所以第三個試探對主特別有吸引力.
\newline
\newline
在每一次的試探和要求,我們的主都不為撒旦所誘；祂對神的旨意完全絕對順服,可見主是人人可信服的彌賽亞.這三個試探與主成為彌賽亞建立國度有關.
\newline
\newline
以色列的國度來了(四12至十章)
\newline
\newline
作者介紹了彌賽亞的可靠性,便將祂的工作握要敘述.當時有一習慣,每一個王工作之先；有一先鋒為他開路預備人心,宣告王的來臨.但這王會待先鋒結束工作之後,纔正式工作.所以(四12)是要緊轉捩點:「耶穌聽見約翰下了監.......」,作者借這句話指出兩件事:第一,主耶穌等先鋒結束工作後,纔正式工作；這是按照當時的習慣.第二,當時國家領袖對彌賽亞先鋒所傳的信息,沒有好感,便將施洗約翰下在監裡.可見他們對天國的信息耳朵發沉,心蒙脂油.於是主便來到迦百農居住,避免與國家領袖正面衝突.迦百農是主宣道的大本營,主開始祂的建國大計,主所傳的信息與施洗約翰一樣,(太四17)作者在(四23-24)用精簡的筆法將主的工作作一結論,作者介紹主三方面的工作:
\newline
\newline
教訓
\newline
\newline
在會堂裡教訓,教訓天國的福音和律法的說話,這信息與舊約天國來臨有關.
\newline
\newline
傳講天國的福音
\newline
\newline
教訓是小組性,傳講是大眾性.
\newline
\newline
醫治百姓
\newline
\newline
百姓是以色列人,當時許多不同病症的人都到耶穌那裡,要祂醫好他們,舊約多處指出,當彌賽亞出現,便會行神蹟.
\newline
\newline
這三方面指出主耶穌就是舊約所指的彌賽亞.祂將人帶到神的國度裡,藉著神蹟奇事,證明祂是真正的彌賽亞,(賽卅五5-6)稱為彌賽亞的預言.醫病趕鬼,都是彌賽亞的權柄.所以當時的人看見耶穌所行的,信的人便多了.不單加利利各地,連外邦人聚居的敘利亞,主的名聲被傳開了.
\newline
\newline
王的宣言(五至七章)
\newline
\newline
每一個王建立國度時均有他的憲章,要求,計劃.主向準天國的人宣佈,也是如此.這些條件若是接受的,天國便可實現；拒絕的,天國便不能實現.當時耶穌走遍加利利,敘利亞和低加波利,很多人跟隨祂.祂便登上一個山教訓他們.奧古斯丁稱這段為「登山寶訓」,歷世歷代信徒喜愛的經文,尤其是八福.我們分析五至七章可作兩大段:
\newline
\newline
天國子民的特徵(五1-16)
\newline
\newline
若要進入天國,需要擁有天國的特性.作者從八方面指出,每一方面都是從「有福的」開始.主教訓信徒,但聽眾中有真有假,這八福與天國的建立,及與舊約有關國度預言均吻合.「得安慰」指彌賽亞的安慰；「承受地土」指承受天國；「得飽足」指在國裡得義的飽足；「蒙憐恤」指在國度裡蒙憐恤；「得見神」指在國度裡看見神；「稱為神的兒子」指在國度裡的名稱；「承受天國」指享受國度.仔細查考,便看見八福與國度的預言有密切關係.這都是猶太人日夕所盼望的,這些福包括屬靈的,屬物質的,只有在這國度裡方可實現.
\newline
\newline
除了八福外,(13-16節)指出天國子民的責任.作「鹽」和「光」,鹽的功用固然是保持食物不腐壞,但根據(路十四34-35)耶穌的一個比喻,鹽的另一個功用是叫土地肥沃,使農產豐收.作天國的子民,一方面保持對神的忠心,另一方面使世上在屬靈上得到大豐收.光的作用剛剛相反,鹽是隱藏性；光是顯明性,兩方面皆要強調.天國的子民,一方面要作隱藏的工作,另一方面要作明顯的工作.
\newline
\newline


\newpage

\section{第54屆~港九培靈會~馬有藻~第四講}
\label{sec:531}
\textbf{以色列的國度來了(續講二)}
\newline
\newline
連結: \href{https://www.hkbibleconference.org/session-message/view/531}{\texttt{https://www.hkbibleconference.org/session-message/view/531}}
\newline
\newline
\hyperref[sec:530]{< < < PREV SERMON < < <}
~
\hyperref[sec:toc]{[返目錄]}
~
\hyperref[sec:532]{> > > NEXT SERMON > > >}
\newline
\newline
馬太福音 5:17-7:29
\newline
\begin{longtable}{cl}
\hline
\hline
章節 & 經文 (和合本修訂版)\\
\hline
5:17 & \begin{tabularx}{0.7\textwidth}{X} 「不要以為我來是要廢掉律法和先知。我來不是要廢掉，而是要成全。 \end{tabularx} \\ \\ \relax
5:18 & \begin{tabularx}{0.7\textwidth}{X} 我實在告訴你們，就是到天地都廢去，律法的一點一畫也不能廢去，直到一切都實現。 \end{tabularx} \\ \\ \relax
5:19 & \begin{tabularx}{0.7\textwidth}{X} 所以，無論誰廢掉這誡命中最小的一條，又教導人也這樣做，他在天國裡要稱為最小的。但無論誰遵行並如此教導人的，他在天國裡要稱為大。 \end{tabularx} \\ \\ \relax
5:20 & \begin{tabularx}{0.7\textwidth}{X} 我告訴你們，你們的義若不勝過文士和法利賽人的義，絕不能進天國。」 \end{tabularx} \\ \\ \relax
5:21 & \begin{tabularx}{0.7\textwidth}{X} 「你們聽過有對古人說：『不可殺人』；『凡殺人的，必須受審判。』 \end{tabularx} \\ \\ \relax
5:22 & \begin{tabularx}{0.7\textwidth}{X} 但是我告訴你們：凡向弟兄動怒的，必須受審判；凡罵弟兄是廢物的，必須受議會的審判；凡罵弟兄是白痴的，必須遭受地獄的火。 \end{tabularx} \\ \\ \relax
5:23 & \begin{tabularx}{0.7\textwidth}{X} 所以，你在祭壇上獻祭物的時候，若想起有弟兄對你懷恨， \end{tabularx} \\ \\ \relax
5:24 & \begin{tabularx}{0.7\textwidth}{X} 就要把祭物留在壇前，先去跟弟兄和好，然後來獻祭物。 \end{tabularx} \\ \\ \relax
5:25 & \begin{tabularx}{0.7\textwidth}{X} 你同告你的冤家還在路上，就要趕快與他講和，免得他把你送交給法官，法官交給警衛，你就下在監裡了。 \end{tabularx} \\ \\ \relax
5:26 & \begin{tabularx}{0.7\textwidth}{X} 我實在告訴你，就是有一個大文錢還沒有還清，你也絕不能從那裡出來。」 \end{tabularx} \\ \\ \relax
5:27 & \begin{tabularx}{0.7\textwidth}{X} 「你們聽過有話說：『不可姦淫。』 \end{tabularx} \\ \\ \relax
5:28 & \begin{tabularx}{0.7\textwidth}{X} 但是我告訴你們：凡看見婦女就動淫念的，這人心裡已經與她犯姦淫了。 \end{tabularx} \\ \\ \relax
5:29 & \begin{tabularx}{0.7\textwidth}{X} 若是你的右眼使你跌倒，就把它挖出來，丟掉。寧可失去身體中的一部分，也不讓整個身體被扔進地獄。 \end{tabularx} \\ \\ \relax
5:30 & \begin{tabularx}{0.7\textwidth}{X} 若是你的右手使你跌倒，就把它砍下來，丟掉。寧可失去身體中的一部分，也不讓整個身體下地獄。」 \end{tabularx} \\ \\ \relax
5:31 & \begin{tabularx}{0.7\textwidth}{X} 「又有話說：『無論誰休妻，都要給她休書。』 \end{tabularx} \\ \\ \relax
5:32 & \begin{tabularx}{0.7\textwidth}{X} 但是我告訴你們：凡休妻的，除非是因不貞的緣故，否則就是使她犯姦淫了；人若娶被休的婦人，也是犯姦淫了。」 \end{tabularx} \\ \\ \relax
5:33 & \begin{tabularx}{0.7\textwidth}{X} 「你們又聽過有對古人說：『不可背誓，所起的誓總要向主謹守。』 \end{tabularx} \\ \\ \relax
5:34 & \begin{tabularx}{0.7\textwidth}{X} 但是我告訴你們：甚麼誓都不可起。不可指著天起誓，因為天是神的寶座。 \end{tabularx} \\ \\ \relax
5:35 & \begin{tabularx}{0.7\textwidth}{X} 不可指著地起誓，因為地是他的腳凳；也不可指著耶路撒冷起誓，因為耶路撒冷是大君王的京城。 \end{tabularx} \\ \\ \relax
5:36 & \begin{tabularx}{0.7\textwidth}{X} 又不可指著你的頭起誓，因為你不能使一根頭髮變黑變白。 \end{tabularx} \\ \\ \relax
5:37 & \begin{tabularx}{0.7\textwidth}{X} 你們的話，是，就說是；不是，就說不是。若再多說，就是出於那惡者。」 \end{tabularx} \\ \\ \relax
5:38 & \begin{tabularx}{0.7\textwidth}{X} 「你們聽過有話說：『以眼還眼，以牙還牙。』 \end{tabularx} \\ \\ \relax
5:39 & \begin{tabularx}{0.7\textwidth}{X} 但是我告訴你們：不要與惡人作對。有人打你的右臉，連另一邊也轉過去由他打。 \end{tabularx} \\ \\ \relax
5:40 & \begin{tabularx}{0.7\textwidth}{X} 有人想要告你，要拿你的裡衣，連外衣也由他拿去。 \end{tabularx} \\ \\ \relax
5:41 & \begin{tabularx}{0.7\textwidth}{X} 有人強迫你走一里路，你就跟他走二里。 \end{tabularx} \\ \\ \relax
5:42 & \begin{tabularx}{0.7\textwidth}{X} 有求你的，就給他；有向你借貸的，不可推辭。」 \end{tabularx} \\ \\ \relax
5:43 & \begin{tabularx}{0.7\textwidth}{X} 「你們聽過有話說：『要愛你的鄰舍，恨你的仇敵。』 \end{tabularx} \\ \\ \relax
5:44 & \begin{tabularx}{0.7\textwidth}{X} 但是我告訴你們：要愛你們的仇敵，為那迫害你們的禱告。 \end{tabularx} \\ \\ \relax
5:45 & \begin{tabularx}{0.7\textwidth}{X} 這樣，你們就可以作天父的兒女了。因為他叫太陽照好人，也照壞人；降雨給義人，也給不義的人。 \end{tabularx} \\ \\ \relax
5:46 & \begin{tabularx}{0.7\textwidth}{X} 你們若只愛那愛你們的人，有甚麼賞賜呢？就是稅吏不也是這樣做嗎？ \end{tabularx} \\ \\ \relax
5:47 & \begin{tabularx}{0.7\textwidth}{X} 你們若只請你弟兄的安，有甚麼比別人強呢？就是外邦人不也是這樣做嗎？ \end{tabularx} \\ \\ \relax
5:48 & \begin{tabularx}{0.7\textwidth}{X} 所以，你們要完全，如同你們的天父是完全的。」 \end{tabularx} \\ \\ \relax
6:1 & \begin{tabularx}{0.7\textwidth}{X} 「你們要謹慎，不可故意在人面前表現虔誠，叫他們看見，若是這樣，就不能得你們天父的賞賜了。 \end{tabularx} \\ \\ \relax
6:2 & \begin{tabularx}{0.7\textwidth}{X} 「所以，你施捨的時候，不可叫人在你前面吹號，像那假冒為善的人在會堂裡和街道上所做的，故意要得人的稱讚。我實在告訴你們，他們已經得了他們的賞賜。 \end{tabularx} \\ \\ \relax
6:3 & \begin{tabularx}{0.7\textwidth}{X} 你施捨的時候，不要讓左手知道右手所做的， \end{tabularx} \\ \\ \relax
6:4 & \begin{tabularx}{0.7\textwidth}{X} 好使你隱祕地施捨；你父在隱祕中察看，必然賞賜你。」 \end{tabularx} \\ \\ \relax
6:5 & \begin{tabularx}{0.7\textwidth}{X} 「你們禱告的時候，不可像那假冒為善的人，愛站在會堂裡和十字路口禱告，故意讓人看見。我實在告訴你們，他們已經得了他們的賞賜。 \end{tabularx} \\ \\ \relax
6:6 & \begin{tabularx}{0.7\textwidth}{X} 你禱告的時候，要進入內室，關上門，向那在隱祕中的父禱告；你父在隱祕中察看，必將賞賜你。 \end{tabularx} \\ \\ \relax
6:7 & \begin{tabularx}{0.7\textwidth}{X} 你們禱告，不可像外邦人那樣重複一些空話，他們以為話多了必蒙垂聽。 \end{tabularx} \\ \\ \relax
6:8 & \begin{tabularx}{0.7\textwidth}{X} 你們不可效法他們。因為在你們祈求以前，你們所需要的，你們的父早已知道了。」 \end{tabularx} \\ \\ \relax
6:9 & \begin{tabularx}{0.7\textwidth}{X} 「所以，你們要這樣禱告：『我們在天上的父：願人都尊你的名為聖。 \end{tabularx} \\ \\ \relax
6:10 & \begin{tabularx}{0.7\textwidth}{X} 願你的國降臨；願你的旨意行在地上，如同行在天上。 \end{tabularx} \\ \\ \relax
6:11 & \begin{tabularx}{0.7\textwidth}{X} 我們日用的飲食，今日賜給我們。 \end{tabularx} \\ \\ \relax
6:12 & \begin{tabularx}{0.7\textwidth}{X} 免我們的債，如同我們免了人的債。 \end{tabularx} \\ \\ \relax
6:13 & \begin{tabularx}{0.7\textwidth}{X} 不叫我們陷入試探；救我們脫離那惡者。因為國度、權柄、榮耀，全是你的，直到永遠。阿們！』 \end{tabularx} \\ \\ \relax
6:14 & \begin{tabularx}{0.7\textwidth}{X} 你們若饒恕人的過犯，你們的天父也必饒恕你們； \end{tabularx} \\ \\ \relax
6:15 & \begin{tabularx}{0.7\textwidth}{X} 你們若不饒恕人，你們的天父也必不饒恕你們的過犯。」 \end{tabularx} \\ \\ \relax
6:16 & \begin{tabularx}{0.7\textwidth}{X} 「你們禁食的時候，不可像那假冒為善的人，臉上帶著愁容；因為他們蓬頭垢面，故意讓人看出他們在禁食。我實在告訴你們，他們已經得了他們的賞賜。 \end{tabularx} \\ \\ \relax
6:17 & \begin{tabularx}{0.7\textwidth}{X} 你禁食的時候，要梳頭洗臉， \end{tabularx} \\ \\ \relax
6:18 & \begin{tabularx}{0.7\textwidth}{X} 不要讓人看出你在禁食，只讓你隱祕中的父看見；你父在隱祕中察看，必然賞賜你。」 \end{tabularx} \\ \\ \relax
6:19 & \begin{tabularx}{0.7\textwidth}{X} 「不要為自己在地上積蓄財寶；地上有蟲子咬，能鏽壞，也有賊挖洞來偷。 \end{tabularx} \\ \\ \relax
6:20 & \begin{tabularx}{0.7\textwidth}{X} 要在天上積蓄財寶；天上沒有蟲子咬，不會鏽壞，也沒有賊挖洞來偷。 \end{tabularx} \\ \\ \relax
6:21 & \begin{tabularx}{0.7\textwidth}{X} 因為你的財寶在哪裡，你的心也在哪裡。」 \end{tabularx} \\ \\ \relax
6:22 & \begin{tabularx}{0.7\textwidth}{X} 「眼睛是身體的燈。你的眼睛若明亮，全身就光明； \end{tabularx} \\ \\ \relax
6:23 & \begin{tabularx}{0.7\textwidth}{X} 你的眼睛若昏花，全身就黑暗。你裡面的光若黑暗了，那黑暗是何等大呢！」 \end{tabularx} \\ \\ \relax
6:24 & \begin{tabularx}{0.7\textwidth}{X} 「一個人不能服侍兩個主；他不是恨這個愛那個，就是重這個輕那個。你們不能又服侍神，又服侍瑪門。」 \end{tabularx} \\ \\ \relax
6:25 & \begin{tabularx}{0.7\textwidth}{X} 「所以，我告訴你們，不要為你們的生命憂慮吃甚麼喝甚麼，或為你們的身體憂慮穿甚麼。生命不勝於飲食嗎？身體不勝於衣裳嗎？ \end{tabularx} \\ \\ \relax
6:26 & \begin{tabularx}{0.7\textwidth}{X} 你們看一看那天上的飛鳥，也不種也不收，也不在倉裡存糧，你們的天父尚且養活牠們。你們不比飛鳥貴重得多嗎？ \end{tabularx} \\ \\ \relax
6:27 & \begin{tabularx}{0.7\textwidth}{X} 你們哪一個能藉著憂慮使壽數多加一刻呢？ \end{tabularx} \\ \\ \relax
6:28 & \begin{tabularx}{0.7\textwidth}{X} 何必為衣裳憂慮呢？你們想一想野地裡的百合花是怎麼長起來的：它也不勞動也不紡線。 \end{tabularx} \\ \\ \relax
6:29 & \begin{tabularx}{0.7\textwidth}{X} 然而我告訴你們，就是所羅門極榮華的時候，他所穿戴的還不如這些花的一朵呢！ \end{tabularx} \\ \\ \relax
6:30 & \begin{tabularx}{0.7\textwidth}{X} 你們這小信的人哪！野地裡的草今天還在，明天就丟在爐裡，神還給它這樣的妝飾，何況你們呢？ \end{tabularx} \\ \\ \relax
6:31 & \begin{tabularx}{0.7\textwidth}{X} 所以，不要憂慮，說：『我們吃甚麼？喝甚麼？穿甚麼？』 \end{tabularx} \\ \\ \relax
6:32 & \begin{tabularx}{0.7\textwidth}{X} 這都是外邦人所求的。你們需要這一切東西，你們的天父都知道。 \end{tabularx} \\ \\ \relax
6:33 & \begin{tabularx}{0.7\textwidth}{X} 你們要先求神的國和他的義，這些東西都要加給你們了。 \end{tabularx} \\ \\ \relax
6:34 & \begin{tabularx}{0.7\textwidth}{X} 所以，不要為明天憂慮，因為明天自有明天的憂慮；一天的難處一天當就夠了。」 \end{tabularx} \\ \\ \relax
7:1 & \begin{tabularx}{0.7\textwidth}{X} 「你們不要評斷別人，免得你們被審判。 \end{tabularx} \\ \\ \relax
7:2 & \begin{tabularx}{0.7\textwidth}{X} 因為你們怎樣評斷別人，也必怎樣被審判；你們用甚麼量器量給人，也必用甚麼量器量給你們。 \end{tabularx} \\ \\ \relax
7:3 & \begin{tabularx}{0.7\textwidth}{X} 為甚麼看見你弟兄眼中有刺，卻不想自己眼中有梁木呢？ \end{tabularx} \\ \\ \relax
7:4 & \begin{tabularx}{0.7\textwidth}{X} 你自己眼中有梁木，怎能對你弟兄說『讓我去掉你眼中的刺』呢？ \end{tabularx} \\ \\ \relax
7:5 & \begin{tabularx}{0.7\textwidth}{X} 你這假冒為善的人！先去掉自己眼中的梁木，然後才能看得清楚，好去掉你弟兄眼中的刺。 \end{tabularx} \\ \\ \relax
7:6 & \begin{tabularx}{0.7\textwidth}{X} 不要把聖物給狗，也不要把你們的珍珠丟在豬面前，恐怕牠們踐踏了珍珠，轉過來咬你們。」 \end{tabularx} \\ \\ \relax
7:7 & \begin{tabularx}{0.7\textwidth}{X} 「你們祈求，就給你們；尋找，就找到；叩門，就給你們開門。 \end{tabularx} \\ \\ \relax
7:8 & \begin{tabularx}{0.7\textwidth}{X} 因為凡祈求的，就得著；尋找的，就找到；叩門的，就給他開門。 \end{tabularx} \\ \\ \relax
7:9 & \begin{tabularx}{0.7\textwidth}{X} 你們中間誰有兒子求餅，反給他石頭呢？ \end{tabularx} \\ \\ \relax
7:10 & \begin{tabularx}{0.7\textwidth}{X} 求魚，反給他蛇呢？ \end{tabularx} \\ \\ \relax
7:11 & \begin{tabularx}{0.7\textwidth}{X} 你們雖然不好，尚且知道拿好東西給兒女，何況你們在天上的父，他豈不更要把好東西賜給求他的人嗎？ \end{tabularx} \\ \\ \relax
7:12 & \begin{tabularx}{0.7\textwidth}{X} 所以，無論何事，你們想要人怎樣待你們，你們也要怎樣待人，因為這就是律法和先知的道理。」 \end{tabularx} \\ \\ \relax
7:13 & \begin{tabularx}{0.7\textwidth}{X} 「你們要進窄門。因為通往滅亡的門是寬的，路是大的，進去的人也多； \end{tabularx} \\ \\ \relax
7:14 & \begin{tabularx}{0.7\textwidth}{X} 通往生命的門是窄的，路是小的，找到的人也少。」 \end{tabularx} \\ \\ \relax
7:15 & \begin{tabularx}{0.7\textwidth}{X} 「你們要防備假先知。他們到你們這裡來，外面披著羊皮，裡面卻是殘暴的狼。 \end{tabularx} \\ \\ \relax
7:16 & \begin{tabularx}{0.7\textwidth}{X} 豈能在荊棘上摘葡萄呢？豈能在蒺藜裡摘無花果呢？憑著他們的果子，就可以認出他們來。 \end{tabularx} \\ \\ \relax
7:17 & \begin{tabularx}{0.7\textwidth}{X} 這樣，凡好樹都結好果子，而壞樹結壞果子。 \end{tabularx} \\ \\ \relax
7:18 & \begin{tabularx}{0.7\textwidth}{X} 好樹不能結壞果子，壞樹也不能結好果子。 \end{tabularx} \\ \\ \relax
7:19 & \begin{tabularx}{0.7\textwidth}{X} 凡不結好果子的樹就砍下來，丟在火裡。 \end{tabularx} \\ \\ \relax
7:20 & \begin{tabularx}{0.7\textwidth}{X} 所以，憑著他們的果子就可以認出他們來。」 \end{tabularx} \\ \\ \relax
7:21 & \begin{tabularx}{0.7\textwidth}{X} 「不是每一個稱呼我『主啊，主啊』的人都能進天國；惟有遵行我天父旨意的人才能進去。 \end{tabularx} \\ \\ \relax
7:22 & \begin{tabularx}{0.7\textwidth}{X} 在那日必有許多人對我說：『主啊，主啊，我們不是奉你的名傳道，奉你的名趕鬼，奉你的名行許多異能嗎？』 \end{tabularx} \\ \\ \relax
7:23 & \begin{tabularx}{0.7\textwidth}{X} 我要向他們宣告：『我從來不認識你們，你們這些作惡的人，給我走開！』」 \end{tabularx} \\ \\ \relax
7:24 & \begin{tabularx}{0.7\textwidth}{X} 「所以，凡聽了我這些話又去做的，好比一個聰明人把房子蓋在磐石上。 \end{tabularx} \\ \\ \relax
7:25 & \begin{tabularx}{0.7\textwidth}{X} 風吹，雨打，水沖，撞擊那房子，房子總不倒塌，因為根基立在磐石上。 \end{tabularx} \\ \\ \relax
7:26 & \begin{tabularx}{0.7\textwidth}{X} 凡聽了我這些話而不去做的，好比一個無知的人把房子蓋在沙土上。 \end{tabularx} \\ \\ \relax
7:27 & \begin{tabularx}{0.7\textwidth}{X} 風吹，雨打，水沖，撞擊那房子，房子就倒塌了，並且倒塌得很厲害。」 \end{tabularx} \\ \\ \relax
7:28 & \begin{tabularx}{0.7\textwidth}{X} 耶穌講完了這些話，眾人對他的教導都感到驚奇， \end{tabularx} \\ \\ \relax
7:29 & \begin{tabularx}{0.7\textwidth}{X} 因為他教導他們正像有權柄的人，不像他們的文士。 \end{tabularx} \\ \\
[1ex]
\hline
\hline
\end{longtable}
天國子民的義行(五17至七章)
\newline
\newline
1.(17-48節)主耶穌解釋舊約義行正確之意,若要成為天國子民；有些義行是天國子民必須遵守的.主耶穌來,乃是成全舊約所提及義的要求.(17-19)主耶穌來世的目的,乃是說明整個舊約律法和義的要求.主駁斥當時的猶太人,誤解律法和先知的義；結果便廢掉律法.耶穌說:「我來不是要廢掉,乃是要成全.」主在20節嚴厲地斥責文士和法利賽人的義,斷不能進天國.他們把舊約編成一套外表敬拜的系統,以為可以得救；主便列舉六個事例指出他們對律法的曲解:
\newline
\newline
21-26節禁殺之例.
\newline
\newline
27-30節姦淫之例.
\newline
\newline
31-32節休妻之例.
\newline
\newline
33-37節起誓之例.
\newline
\newline
38-42節報仇之例.
\newline
\newline
43-48節待敵之例.
\newline
\newline
猶太人對這些例有錯誤的解釋.例如:妻子把飯燒焦,便可以休妻.可見他們曲解律法,把律法的精義,變成人為和膚淺的規條.耶穌所指舊約律法是關乎人,而不是外表方面；耶穌解釋重點放在內心,動機和存心.那是全人方面,祂將舊約對義的要求作一總結束.(五48)是上文的結束,這句話引自利十九章,是整個舊約的精髓.主指出凡願意進入天國的人,均要接受這邀請,可見在地上有很高屬靈的要求.
\newline
\newline
2.(六1-七6)消極性的教導,注意「不」,「不要」字的出現,消極性方面共有七點:
\newline
\newline
六1-4施捨時不要虛假.
\newline
\newline
六5-13禱告時不要虛假.
\newline
\newline
六14-15饒恕人時不要虛假.
\newline
\newline
六16-18禁食時不要虛假.
\newline
\newline
六19-24積財時不要愚昧.
\newline
\newline
六25-34生活有需要時不要憂慮.
\newline
\newline
七1-6與人相處時不要論斷.
\newline
\newline
3.我們特別要看「主禱文」(六9-13),細心分析主禱文有六個祈求,頭三個與國度建立有關,後三個與傳講神國來臨有關.
\newline
\newline
「願人都尊你的名為聖」,在舊約神的名顯為聖是在國度裡；神的名顯現是在國度建立的時候；所以主教導人的禱告,願這國儘快來到.
\newline
\newline
「願你的國降臨」,這國乃是先鋒和主耶穌所傳的國,是神應許給大衛的國度.
\newline
\newline
「願你的旨意行在地上,如同行在天上」,主用兩個時式,將來式和現在式.換句話說神的旨意已經進行著,但尚未完全在地上實行；到國度建立時便可以行在地上,好像現在行在天上一樣.
\newline
\newline
「求賜每日的飲食」,求神供應物質.
\newline
\newline
「免我們的債」,債原文是「罪」,指屬靈的債；天國子民要學習求神赦免他們的污穢,預備主國度來臨,先要罪得赦免.
\newline
\newline
「不叫我們遇見試探,救我們脫離兇惡」,國度建立之前,會遇見許多災難.「試探」就是災難這字.舊約多處指出,國度建立前,必有災難出現,主教導他們在災難前,求主幫助他們脫離惡者和災難.
\newline
\newline
4.(七7-23)積極性教導,注意「要」字的出現,積性教導有五方面:
\newline
\newline
七7-11祈禱要真誠.
\newline
\newline
七12待人要敬愛.
\newline
\newline
七13-14要有進窄門心志.
\newline
\newline
七15-20要防備假先知.
\newline
\newline
七21-23要遵行天父旨意.
\newline
\newline
5.(七24-29)結束登山寶訓時,主耶穌給當時的人一個邀請,給他們機會成為天國的子民,「凡聽見......」是一個邀請.29節他們聽見主的教訓,覺得主很有權柄,因自古以來沒有人的教訓好像主耶穌一樣.
\newline
\newline
王的權能(八至十章)
\newline
\newline
作者介紹主耶穌的權能.行各類的神蹟,由23節指出主工作有教訓,醫病,趕鬼,教訓在五至七章已發揮了.八至九章是指醫病的部份,可見主藉著言語去顯出祂的權能；又藉著工作去證實祂的權能.八至九章記載九個神蹟,是主在地上作王最有力的憑據.注意三件要緊的事:第一,這些神蹟是在不同年日發生的,作者經過詳細挑選,證明主耶穌有建立國度的權能.第二,主的神蹟與摩西的神蹟,皆有神權柄的明證.不過摩西的神蹟是審判性,主的神蹟卻是祝福性.第三,(太十二28)是明白所有神蹟的重要經文,靠著神的靈趕鬼治病,要證明以色列的國度來了.
\newline
\newline
馬太編排他的神蹟時,常用每三個為一組.故此這九個神蹟可分作三組,在兩組之中,插入作門徒的故事來隔開.
\newline
\newline
(一)第一組:醫病神蹟
\newline
\newline
1.(八1-4)醫治長大麻瘋,這神蹟有兩大目的:第一顯出神有權能.第二從主吩咐那人去見祭司和察驗身體；可知主來地上,不是要棄掉律法,乃是要順服律法.
\newline
\newline
2.(八5-13)醫治百夫長的僕人,馬太選用這神蹟的次序,他先引述彌賽亞有醫治猶太人的權柄,現在又引述主醫治外邦人.百夫長的信心是很大的,(八10)這句話反映當時猶太領袖們對主非常硬心的；所以這神蹟指出第一,猶太人的硬心,第二,彌賽亞是關乎全地的.
\newline
\newline
3.(八14-17)醫病趕鬼,作者將彌賽亞的醫病神蹟連接到舊約的預言.指出疾病和罪有互相牽連關係,彌賽亞先赦免人的罪,纔醫治人的病.
\newline
\newline
第一個插段(八18-22)
\newline
\newline
兩個準門徒不能進入天國的原因.(18-20)因為物質享受的緣故,這裡指出天國子民要有為主吃苦的心志.「人子」是馬太出現的第一次,但九章指彌賽亞的稱號.「死人埋葬死人」指屬靈的死人埋葬屬肉體的人,第二個人以為作天國的子民還有很多機會,主責備他失去逼切感.第一個人失去犧牲感,因此兩人都不能成為天國的子民.
\newline
\newline
(二)第二組:彰顯權能的神蹟(八23-九28)
\newline
\newline
1.(八23-27)自然界,主平靜風和海.
\newline
\newline
2.(八28-34)主把加大拉的鬼趕走,鬼看見主來了,承認主是彌賽亞.(八29),因為彌賽亞的出現,國度來臨,便是他們的末日.由此可見,主有權柄勝過魔鬼的權柄.加大拉的鬼承認彌賽亞,但加大拉的居民卻拒絕了主.
\newline
\newline
3.(九1-8)主有赦罪的權能.「但叫你們知道人子有赦罪的權柄.(九6),猶太人從不敢講這種話,除非他是從神而來的；可見主是彌賽亞,要應驗(耶卅一章)赦罪之約.
\newline
\newline
第二個插段(九9-17)
\newline
\newline
指出真門徒的事,先提及馬太蒙召(九9-13),並指出國家領袖慢慢反對祂.(14節)指出禁食的問題,主將自己比喻作新郎.新郎在舊約通常指彌賽亞,主也在這裡第一次暗示要離開地上.這比喻指出舊時代要被新時代取代.
\newline
\newline
(三)第三組:復原性的神蹟(九18-34)
\newline
\newline
復原性的神蹟,強調彌賽亞有復原的能力,彌賽亞時代是復原萬物的時代(太十九28),復興時代就是彌賽亞作王的時代.
\newline
\newline
1.(九18-26)使死人復活.
\newline
\newline
2.(九27-31)瞎子承認彌賽亞為王.
\newline
\newline
3.啞吧被醫好,(34節)國家領袖把神蹟的事歸功於撒旦,但他們不能否認啞吧是被醫好了,可歎他們卻不能接受主.
\newline
\newline
(四)主差遣門徒(九35-十一1)
\newline
\newline
主耶穌不單藉著工作宣告天國來了,也差派門徒出去.(十一1)指出他自己也出去.主的憐憫(九36)是主差派的基礎,當時國家領袖已不配作牧人了；牧人是國家領袖的代名詞,於是主便差派門徒.
\newline
\newline
門徒的權柄(十1)受託的權柄與主內蘊的能力是一樣的.
\newline
\newline
門徒的數目(十1)十二是以色列的數字,因門徒的對象是猶太人.
\newline
\newline
門徒工作範圍(十5-6),在以色列地,因為天國的福音先傳給猶太人.
\newline
\newline
門徒的信息(十7-8)信息與彌賽亞及先鋒一樣,他們也有行神蹟的能力.
\newline
\newline
門徒的方式(十9-15)傳天國的福音,不要為生活憂慮,接受他們的人便蒙福,拒絕的便受審判.
\newline
\newline
門徒的危險(十16-25)預告他們,傳天國福音的,將來會遭受逼迫苦楚.讓他們先作準備,有受苦的心志,纔能有效地為主工作.
\newline
\newline
門徒的安慰(十26-42)主讓他們知道,若不是主許可,任何苦難不會臨到他們身上,當他們受逼迫時,記得有主的保守和供應,叫他們忠心為主工作,並且有主的賞賜給他們.
\newline
\newline
吩咐完了門徒,主獨自到不同地方去工作.(十一1)可見主不單吩咐別人,祂自己也親力親為去工作.
\newline
\newline


\newpage

\section{第54屆~港九培靈會~馬有藻~第五講}
\label{sec:532}
\textbf{以色列的王被棄絕了}
\newline
\newline
連結: \href{https://www.hkbibleconference.org/session-message/view/532}{\texttt{https://www.hkbibleconference.org/session-message/view/532}}
\newline
\newline
\hyperref[sec:531]{< < < PREV SERMON < < <}
~
\hyperref[sec:toc]{[返目錄]}
~
\hyperref[sec:533]{> > > NEXT SERMON > > >}
\newline
\newline
馬太由五章至第十章佈下耶穌就是彌賽亞的局.從耶穌的工作和教訓,及祂差派門徒的事上；讓讀者看見,以色列王來了,祂的國度也來了.從第十一章開始,作者指出彌賽亞所建立的國度不能成立原因.十一至十二章指出當時的人用許多方法來將主所行的神蹟不歸榮於神,而歸榮於撒旦.
\newline
\newline
以色列的王被棄絕了(十一至十六章)
\newline
\newline
基督被棄絕(十一2-30)
\newline
\newline
基督的先鋒被棄絕(2-15)
\newline
\newline
作者記載施洗約翰下在監裡有兩個目的:第一,證明當時國家領袖不歡迎主.既然不歡迎主,當然也不歡迎主的先鋒,施洗約翰被囚,是主被棄絕之前奏.第二,作者以主回答約翰門徒的問答,指出祂正是眾人期待的彌賽亞.
\newline
\newline
施洗約翰覺得主一直用憐憫待以色列人,而不是用審判待以色列人.施洗約翰初出來傳道之時,以為彌賽亞是帶著刑罰來到.(太三10-12)但他見主的作風不同,便覺得有疑惑,便打發門徒去見耶穌,主耶穌用祂自己所行的神蹟來辯明自己的身份(十一4-5),用工作代替祂的回答.當他們走後,主暗中稱讚施洗約翰,有五大要點:
\newline
\newline
1.(7-10節)約翰的偉大因他是彌賽亞的先鋒,應驗(瑪三1)的預言.
\newline
\newline
2.(11節)約翰雖然偉大,亦不及進入天國的人.因為約翰是預備天國給人進入去,從這節經文,可見天國尚未實現.
\newline
\newline
3.(12節)天國是經歷艱難纔得建立,「努力進入」原文是被動詞,有不好的含意,我認為該這樣譯:「天國是經歷難阻的,難阻的人將之奪去」.努力進入可譯作經歷難阻,努力的人就得著了.可譯作難阻的人將之奪去.天國遭遇許多難阻,尤其是法利賽人,文士等難阻,難阻的人便將天國奪去.
\newline
\newline
4.(13節)天國的總題是先知和約翰的預言,但直至約翰為止；因約翰是舊約最後的先知.
\newline
\newline
5.(14節)如果他們接受約翰的信息,天國便實現了,施洗約翰也變成以利亞了.主稱讚約翰的事奉,主給祂的先鋒五點評論,讓我們了解施洗約翰的事奉和生平.
\newline
\newline
基督棄絕當時世代.(16-24)
\newline
\newline
1.(16-19節)麻木的世代.
\newline
\newline
主用小孩子在街上玩耍,比喻當時世代.有人傳信息,但沒有人聽他們的信息.
\newline
\newline
2.(20-24節)以色列國必受責打
\newline
\newline
主宣佈這些人一定要受責罰,耶穌四處行神蹟,他們看見主的工作,但仍不悔改,將來會受主的責罰.
\newline
\newline
3.(25-30節)主給個人蒙恩得救邀請.
\newline
\newline
(25節)「聰明通達人」指國家領袖,「嬰孩」指悔改的人.既然國家領袖拒絕主,主只有向個人發出邀請,來擔負他們的重擔.
\newline
\newline
基督被棄絕過程的選例.(十二章)
\newline
\newline
作者在上文指出主棄絕祂的百姓,便舉出五個事例,證明百姓如何棄絕主和主棄絕他們.
\newline
\newline
(一)(1-8節)安息日掐麥穗.
\newline
\newline
(二)(9-21節)在安息日治病.
\newline
\newline
在這兩件事情主成功地運用舊約,指出他們屬靈的遲鈍；只注意外表禮儀,卻失去律法的精意.
\newline
\newline
(三)(22-37節)醫治被鬼附,瞎,啞的人.
\newline
\newline
反對的人將主的能力,歸給鬼王別西卜.可見他們不顧一切,出下策來反對主.主從三方面指出他們不合理的地方:
\newline
\newline
1.(25節)一城一家自相分裂必站不住,若將主的能力歸給鬼王,那麼鬼國必站立不住.
\newline
\newline
2.(26節)鬼不會趕鬼.
\newline
\newline
3.(27節)法利賽人承認他們的門徒可以趕鬼,但不能說出是靠誰的能力.
\newline
\newline
主在三方面宣佈他們將來的審判:
\newline
\newline
1.(28節)趕鬼的神蹟,是神國的明證.
\newline
\newline
2.(29節)預告在神國裡,撒旦會被捆綁.
\newline
\newline
3.(30-33節)否認主靠著聖靈行神蹟的罪；是褻瀆聖靈的罪,不可以赦免的.
\newline
\newline
(四)(38-45節)求主行神蹟
\newline
\newline
他們要求主再行神蹟,可見他們不接受主以前所行的神蹟；主便決定不再行神蹟給他們看,(除了復活的神蹟以外).這句話表示主棄絕他們,因他們先棄絕了主.
\newline
\newline
(五)(46-50節)家人反對主
\newline
\newline
主受當時百姓反對,除了祂母親外,祂的家人也反對祂；在主復活之前,祂的家人不相信祂.主指出進入天國的條件,不在乎肉身的關係,乃在乎人與神發生關係.
\newline
\newline
基督被棄絕後的宣告(十三1-52)
\newline
\newline
因為基督被棄絕,於是主做事方式也改變了.首先祂說及國度的建立,用比喻來說明.主透過比喻,釋放七方面天國的奧秘比喻.比喻有兩個目的:啟示真理和隱藏真理.(11-13節),指出「奧秘」這二字在原文是先前隱藏的,現在顯露了.有關天國的真理,主以前沒有講,現在纔顯明了.有人對這些奧秘有不同看法,有人說這些天國的奧秘,是單為猶太人說的；也有人說這些天國的比喻,是更正門徒的觀念；又有人說這些天國的奧秘,是指屬靈的天國.其實根據上下文,天國乃是施洗約翰和主所講的,這些奧秘,只不過是啟示天國另一面的真理；由於當時的人拒絕主,所以天國屬地的部份未能實現.而這裡強調的,是天國屬靈的部份,也包括教會時代.這七個比喻可分為二組,第一組是主在海邊對人說的,第二組是主在屋內說的.
\newline
\newline
撒種比喻
\newline
\newline
這比喻指出一個事實,多數的人對天國信息是遲鈍,眼瞎的,只有小部份人接受.
\newline
\newline
稗子比喻
\newline
\newline
指出在屬靈天國裡,真假摻雜,待天國在地上實現的時候,假的就會分別出來.
\newline
\newline
芥菜種比喻
\newline
\newline
指出天國由小開始,逐漸擴大至很大的領域.
\newline
\newline
麵酵比喻
\newline
\newline
指出天國裡面隱藏發展力,「全團」指天國建立的時候.
\newline
\newline
藏寶比喻
\newline
\newline
指出天國從以色列奪去後,以色列在神面前的地位,這「寶」是指以色列.
\newline
\newline
尋珠比喻
\newline
\newline
真珠指教會,指出教會在基督兩臨之間,這個過渡時期的地位.
\newline
\newline
撒網比喻
\newline
\newline
指出天國建立前所接受的審判.
\newline
\newline
最後的比喻(邀請接受天國的比喻)
\newline
\newline
指出因天國的價值非常大,要接受邀請,將新舊東西拿出來,來換這天國.
\newline
\newline
基督被棄後的工作.(十三53-十六)
\newline
\newline
天國的奧秘,是因當時國家領袖拒絕主,這個啟示刻劃了主工作方針的改變.從此主大部份時間,離開國家領袖和教導門徒,裝備他們；好叫祂返天家的時候,門徒有獨立能力去傳天國的福音.這裡可以分成三個回合:
\newline
\newline
第一回合:被拒與被受.(十三54-十四34)
\newline
\newline
(十三54-58)耶穌在自己家鄉被拒絕.
\newline
\newline
(十四1-12)主受希律的反對,先鋒的遭遇是主遭遇的前奏.主等約翰下監,纔公開廣大工作.約翰死後,主約束祂的工作,逐漸將祂的國收回.
\newline
\newline
耶穌被接受.作者提出三個猶太人接受主的事例,指出主雖然被拒絕,仍施恩憐憫:
\newline
\newline
1.(十四13-21)主喂飽五千人.
\newline
\newline
主在前面曾說過不再行神蹟,為甚麼在這裡再行神蹟?先前的神蹟,是證明身份；現在是要教導門徒,認識祂就是生命之糧,回應主禱文.
\newline
\newline
2.(十四22-33)水面上行走.
\newline
\newline
這神蹟教導門徒的信心,定睛耶穌身上,彼得不望主便沉下去了.
\newline
\newline
3.(十四34-36)在革尼撒勒醫治許多人.教導門徒的信心不要疑惑.
\newline
\newline


\newpage

\section{第54屆~港九培靈會~馬有藻~第六講}
\label{sec:533}
\textbf{以色列的王被棄絕了(續講一)}
\newline
\newline
連結: \href{https://www.hkbibleconference.org/session-message/view/533}{\texttt{https://www.hkbibleconference.org/session-message/view/533}}
\newline
\newline
\hyperref[sec:532]{< < < PREV SERMON < < <}
~
\hyperref[sec:toc]{[返目錄]}
~
\hyperref[sec:534]{> > > NEXT SERMON > > >}
\newline
\newline
馬太福音 15:1-15:39
\newline
\begin{longtable}{cl}
\hline
\hline
章節 & 經文 (和合本修訂版)\\
\hline
15:1 & \begin{tabularx}{0.7\textwidth}{X} 那時，有法利賽人和文士從耶路撒冷來見耶穌，說： \end{tabularx} \\ \\ \relax
15:2 & \begin{tabularx}{0.7\textwidth}{X} 「你的門徒為甚麼違反古人的傳統？因為他們吃飯的時候不洗手。」 \end{tabularx} \\ \\ \relax
15:3 & \begin{tabularx}{0.7\textwidth}{X} 耶穌回答他們：「你們為甚麼因你們的傳統而違反神的誡命呢？ \end{tabularx} \\ \\ \relax
15:4 & \begin{tabularx}{0.7\textwidth}{X} 神說：『當孝敬父母』；又說：『咒罵父母的，必須處死。』 \end{tabularx} \\ \\ \relax
15:5 & \begin{tabularx}{0.7\textwidth}{X} 你們倒說：『無論誰對父母說：我所當供奉你的已經作了奉獻， \end{tabularx} \\ \\ \relax
15:6 & \begin{tabularx}{0.7\textwidth}{X} 就可以不孝敬他的父親。』這就是你們藉著傳統，廢了神的話。 \end{tabularx} \\ \\ \relax
15:7 & \begin{tabularx}{0.7\textwidth}{X} 假冒為善的人哪！以賽亞指著你們所預言的說得好： \end{tabularx} \\ \\ \relax
15:8 & \begin{tabularx}{0.7\textwidth}{X} 『這百姓用嘴唇尊敬我，他們的心卻遠離我。 \end{tabularx} \\ \\ \relax
15:9 & \begin{tabularx}{0.7\textwidth}{X} 他們把人的規條當作教義教導人；他們拜我也是枉然。』」 \end{tabularx} \\ \\ \relax
15:10 & \begin{tabularx}{0.7\textwidth}{X} 耶穌叫了眾人來，對他們說：「你們要聽，也要明白。 \end{tabularx} \\ \\ \relax
15:11 & \begin{tabularx}{0.7\textwidth}{X} 從口裡進去的不玷污人，從口裡出來的才玷污人。」 \end{tabularx} \\ \\ \relax
15:12 & \begin{tabularx}{0.7\textwidth}{X} 當時，門徒進前來對他說：「法利賽人聽見這話很反感，你知道嗎？」 \end{tabularx} \\ \\ \relax
15:13 & \begin{tabularx}{0.7\textwidth}{X} 耶穌回答：「一切植物，若不是我天父栽植的，都要連根拔出來。 \end{tabularx} \\ \\ \relax
15:14 & \begin{tabularx}{0.7\textwidth}{X} 由他們吧！他們是瞎子作瞎子的嚮導；若是瞎子領瞎子，兩個人都要掉在坑裡。」 \end{tabularx} \\ \\ \relax
15:15 & \begin{tabularx}{0.7\textwidth}{X} 彼得回應他說：「請將這比喻講解給我們聽。」 \end{tabularx} \\ \\ \relax
15:16 & \begin{tabularx}{0.7\textwidth}{X} 耶穌說：「連你們也還不明白嗎？ \end{tabularx} \\ \\ \relax
15:17 & \begin{tabularx}{0.7\textwidth}{X} 難道你們不了解，凡進到口裡的，是經過肚子，又排入廁所嗎？ \end{tabularx} \\ \\ \relax
15:18 & \begin{tabularx}{0.7\textwidth}{X} 然而口裡出來的是出於心裡，這才玷污人。 \end{tabularx} \\ \\ \relax
15:19 & \begin{tabularx}{0.7\textwidth}{X} 因為出於心裡的有種種惡念，如兇殺、姦淫、淫亂、偷盜、偽證、毀謗。 \end{tabularx} \\ \\ \relax
15:20 & \begin{tabularx}{0.7\textwidth}{X} 這些才玷污人。至於不洗手吃飯，那並不玷污人。」 \end{tabularx} \\ \\ \relax
15:21 & \begin{tabularx}{0.7\textwidth}{X} 耶穌離開那裡，退到推羅、西頓境內。 \end{tabularx} \\ \\ \relax
15:22 & \begin{tabularx}{0.7\textwidth}{X} 有一個迦南婦人從那地方出來，喊著說：「主啊，大衛之子，可憐我！我女兒被鬼纏得很苦。」 \end{tabularx} \\ \\ \relax
15:23 & \begin{tabularx}{0.7\textwidth}{X} 耶穌卻一言不答。門徒進前來，求他說：「這婦人在我們後頭喊叫，請打發她走吧。」 \end{tabularx} \\ \\ \relax
15:24 & \begin{tabularx}{0.7\textwidth}{X} 耶穌回答：「我奉差遣只到以色列家迷失的羊那裡去。」 \end{tabularx} \\ \\ \relax
15:25 & \begin{tabularx}{0.7\textwidth}{X} 那婦人來拜他，說：「主啊，幫幫我！」 \end{tabularx} \\ \\ \relax
15:26 & \begin{tabularx}{0.7\textwidth}{X} 他回答：「拿孩子的餅丟給小狗吃是不妥的。」 \end{tabularx} \\ \\ \relax
15:27 & \begin{tabularx}{0.7\textwidth}{X} 婦人說：「主啊，不錯，可是小狗也吃牠主人桌上掉下來的碎屑。」 \end{tabularx} \\ \\ \relax
15:28 & \begin{tabularx}{0.7\textwidth}{X} 於是耶穌回答她說：「婦人，你的信心很大！照你所要的成全你吧。」從那時起，她的女兒就好了。 \end{tabularx} \\ \\ \relax
15:29 & \begin{tabularx}{0.7\textwidth}{X} 耶穌離開那地方，來到靠近加利利的海邊，就上山坐下。 \end{tabularx} \\ \\ \relax
15:30 & \begin{tabularx}{0.7\textwidth}{X} 有一大群人到他那裡，帶著瘸子、盲人、肢殘的、聾啞的，和好些別的病人，都放在他腳前，他就治好了他們。 \end{tabularx} \\ \\ \relax
15:31 & \begin{tabularx}{0.7\textwidth}{X} 於是眾人都驚訝，因為看見聾啞的說話，肢殘的痊癒，瘸子行走，盲人看見，他們就歸榮耀給以色列的神。 \end{tabularx} \\ \\ \relax
15:32 & \begin{tabularx}{0.7\textwidth}{X} 耶穌叫門徒來，說：「我憐憫這群人，因為他們同我在這裡已經三天，沒有吃的東西了。我不願意叫他們餓著回去，恐怕他們在路上餓昏了。」 \end{tabularx} \\ \\ \relax
15:33 & \begin{tabularx}{0.7\textwidth}{X} 門徒說：「我們在這野地，哪裡有這麼多的餅讓這許多人吃飽呢？」 \end{tabularx} \\ \\ \relax
15:34 & \begin{tabularx}{0.7\textwidth}{X} 耶穌對他們說：「你們有多少餅？」他們說：「有七個，還有幾條小魚。」 \end{tabularx} \\ \\ \relax
15:35 & \begin{tabularx}{0.7\textwidth}{X} 他就吩咐眾人坐在地上， \end{tabularx} \\ \\ \relax
15:36 & \begin{tabularx}{0.7\textwidth}{X} 拿著這七個餅和幾條魚，祝謝了，擘開，遞給門徒；門徒又遞給眾人。 \end{tabularx} \\ \\ \relax
15:37 & \begin{tabularx}{0.7\textwidth}{X} 他們都吃，並且吃飽了，收拾剩下的碎屑，裝滿了七個筐子。 \end{tabularx} \\ \\ \relax
15:38 & \begin{tabularx}{0.7\textwidth}{X} 吃的人中，男的有四千，還不算婦女和孩子。 \end{tabularx} \\ \\ \relax
15:39 & \begin{tabularx}{0.7\textwidth}{X} 耶穌叫眾人散去，就上船，來到馬加丹境內。 \end{tabularx} \\ \\
[1ex]
\hline
\hline
\end{longtable}
第二回合:被拒與被受(十五1-39)
\newline
\newline
被拒絕方面,馬太選用主與當時法利賽人辯論遺傳的事.十五1-20這段有三要點:
\newline
\newline
1.(1節)他們從耶路撒冷上來,可見他們是國家宗教的代表,來與耶穌辯論；表示這國不斷棄絕主作他們的彌賽亞.
\newline
\newline
2.(7-9節)耶穌用(賽廿九13)經文來駁斥他們曲解律法的精意.
\newline
\newline
3.(13節)暗指文士和法利賽人不是天父所栽種之物,終會連根拔出受審判.
\newline
\newline
馬太又記載主施憐憫的事蹟,十四章的對象是猶太人,但這裡的是外邦人.
\newline
\newline
1.(21-28節)耶穌與迦南婦人講道.整段事情,有兩點要留意:第一,指出以色列在神心中的地位(24節).第二,主沒有厭棄外邦人,仍施憐憫.
\newline
\newline
2.(29-31節)在加利利海邊治病,從上下文看見來接受治病的大部份是外邦人.主教導門徒一件事實,因猶太人不接受祂；主漸漸歡迎外邦人,進入祂救贖的計劃內.
\newline
\newline
3.(32-39節)餵飽四千人的事蹟,與餵飽五千人之神蹟有許多分別；最大分別是對象不同.餵飽五千人的對象是猶太人,而這裡對象是外邦人.
\newline
\newline
第三回合:被拒與被受,(十六1-28)
\newline
\newline
在被拒方面,法利賽人和撒都該人試探耶穌.(1-12節)
\newline
\newline
法利賽人和撒都該人本是世仇,但為了反對主,竟然成為好朋友.從他們的要求,可見他們否認主以往的神蹟.第一節可譯作「請他顯個從天上來的神蹟給他們看」.可見他們否認主靠聖靈,及天上的力量來行神蹟.主便駁斥他們,可觀看天文,但卻不能分辨祂所行的神蹟,便放棄他們.第4節「離開」是放棄之意.主放棄了他們之後,便警告門徒要防備這些人.
\newline
\newline
(13-28節)主被接納.
\newline
\newline
主帶領門徒到該撒利亞腓立比,考驗門徒對祂的認識；並加強天國的信息,明明告知他們,祂上十字架的真理.
\newline
\newline
1.(13-17節)有關彌賽亞的認識.
\newline
\newline
該撒利亞腓立比是外邦人之地,主問門徒對祂的認識,他們有很多的答案,但只有彼得的答案正確(16節).於是主稱讚彼得是有福的人,得神的啟示也合神心意.
\newline
\newline
2.(18-20節)有關彌賽亞國的認識.
\newline
\newline
「我要把我的教會」是關乎教會的啟示,主宣告天國的另一面孔,乃是教會的出現.磐石有三種解釋:第一,磐石指彼得；第二,磐石指彼得對主身份的認識；第三,指基督.我個人接受第三個解釋,因為在舊約文學,神是稱為磐石,而彼得後來也說教會是建立在基督磐石上(彼前二4)
\newline
\newline
教會這個字第一次在馬太出現,但這觀念對當時的人不算陌生.這字在七十士譯本常常出現,基本意思是團體,聚會,但這裡特別指一群對主有信心的門徒.從「我的」可見信徒是特別的團體.「建造」是將來式動詞,可見教會的應許尚未實現.「我要將天國的」這是使命的交託,鑰匙代表權柄和管家職份.(Stewardship),彼得有管家的責任,是將天國的信息傳揚；保證他在地上可以捆綁和釋放.這兩個詞都是法庭用語,可見彼得有合法權,宣告人的罪赦免或留下.主耶穌宣佈門徒成為新的團體,這團體使命是傳講天國福音的信息,主耶穌吩咐他們不要馬上把這消息傳開(20節),免得基督受苦的計劃被腰斬了.
\newline
\newline
3.(21-28節)啟示受苦與復活.
\newline
\newline
「從此」這是主生平很大的轉捩點,因為主認定門徒對祂的身份,有了準確的認識,所以主明示要上耶路撒冷去,為國家和全地的人受死.彼得愛主的心很強,但沒有體貼神的心意.主在三方面提醒他:第一,(24節)告訴他作主門徒必須有為主付代價的心志；第二,(27節)主必須先受苦纔有榮耀；第三(28節)告訴他們在生命結束之前,必有人看見主降臨在祂的國裡.這節聖經引起很多不同的解釋:第一,指耶穌的復活和升天；第二,指五旬節聖靈的降臨；第三,指主後七十年耶路撒冷被毀；第四,指天國建立時候的情形.我個人贊成第四種解釋,下文(十七1-7)登山變像,及(彼後一16-18)也是提及這件事.
\newline
\newline
第四回合:以色列的國度被奪去了(十七至廿五章)
\newline
\newline
因著國家領袖棄絕祂,所以主四處帶領門徒,教導,裝備,訓練他們；好叫他們將來可以繼承天國的責任.
\newline
\newline
基督被奪去前的訓練工作.(十七至二十章)
\newline
\newline
作者選出十二件事,分四組來教導門徒,天國子民的生活.
\newline
\newline
(一)(十七1-13)論天國的面貌
\newline
\newline
這件事是回應上文之事,耶穌帶門徒上山,讓他們偷窺彌賽亞國度的榮耀情形,這件事指出:
\newline
\newline
將來有國度的出現,這國度有舊約和新約聖徒；摩西,以利亞是舊約的代表；三個門徒是新約信徒的代表.
\newline
\newline
天上有聲音,證實主確是彌賽亞,(詩二篇,賽四十二章,申十八章也指證這事.)
\newline
\newline
摩西,以利亞談論之事,從(路九31)可以看見是有關主去世之事,基督的死是達到寶座的途徑.
\newline
\newline
(二)(14-21節)論信心的原則.
\newline
\newline
信心可以趕鬼,趁機指出當時世代是不信悖謬的；連門徒的信心也太小,主指出信心可以移山倒海.
\newline
\newline
(三)(22-27節)論納丁說的原則.
\newline
\newline
證實主確是世上的王.
\newline
\newline
(四)(十八1-16)論天國的大小.
\newline
\newline
主指出,凡是謙卑事奉的,便是為大.
\newline
\newline
(五)(7-2節)論絆倒人的事.
\newline
\newline
(六)(21-35節)論饒恕人的事.
\newline
\newline
(七)(十九1-12)論休妻.
\newline
\newline
(八)(13-15節)論天國的樣式.
\newline
\newline
(九)(16-26節),論進入天國的攔阻.
\newline
\newline


\newpage

\section{第54屆~港九培靈會~馬有藻~第七講}
\label{sec:534}
\textbf{以色列的國度被奪去了}
\newline
\newline
連結: \href{https://www.hkbibleconference.org/session-message/view/534}{\texttt{https://www.hkbibleconference.org/session-message/view/534}}
\newline
\newline
\hyperref[sec:533]{< < < PREV SERMON < < <}
~
\hyperref[sec:toc]{[返目錄]}
~
\hyperref[sec:535]{> > > NEXT SERMON > > >}
\newline
\newline
馬太福音 19:27-20:16
\newline
\begin{longtable}{cl}
\hline
\hline
章節 & 經文 (和合本修訂版)\\
\hline
19:27 & \begin{tabularx}{0.7\textwidth}{X} 於是彼得回應，對他說：「看哪，我們已經撇下一切跟從你了，我們會得到甚麼呢？」 \end{tabularx} \\ \\ \relax
19:28 & \begin{tabularx}{0.7\textwidth}{X} 耶穌對他們說：「我實在告訴你們，你們這些跟從我的人，到了萬物更新、人子坐在他榮耀寶座上的時候，你們也要坐在十二個寶座上，審判以色列十二個支派。 \end{tabularx} \\ \\ \relax
19:29 & \begin{tabularx}{0.7\textwidth}{X} 凡為我的名撇下房屋，或是兄弟、姊妹、父親、母親、兒女、田地的，將得著百倍，並且承受永生。 \end{tabularx} \\ \\ \relax
19:30 & \begin{tabularx}{0.7\textwidth}{X} 然而，有許多在前的，將要在後；在後的，將要在前。」 \end{tabularx} \\ \\ \relax
20:1 & \begin{tabularx}{0.7\textwidth}{X} 「因為天國好比一家的主人清早去雇人進他的葡萄園做工。 \end{tabularx} \\ \\ \relax
20:2 & \begin{tabularx}{0.7\textwidth}{X} 他和工人講定一天一個銀幣，就打發他們進葡萄園去。 \end{tabularx} \\ \\ \relax
20:3 & \begin{tabularx}{0.7\textwidth}{X} 約在上午九點鐘出去，看見市場上還有閒站的人， \end{tabularx} \\ \\ \relax
20:4 & \begin{tabularx}{0.7\textwidth}{X} 就對那些人說：『你們也進葡萄園去，我會給你們合理的工錢。』 \end{tabularx} \\ \\ \relax
20:5 & \begin{tabularx}{0.7\textwidth}{X} 他們也進去了。約在正午和下午三點鐘又出去，他也是這麼做。 \end{tabularx} \\ \\ \relax
20:6 & \begin{tabularx}{0.7\textwidth}{X} 約在下午五點鐘出去，他看見還有人站在那裡，就問他們：『你們為甚麼整天在這裡閒站呢？』 \end{tabularx} \\ \\ \relax
20:7 & \begin{tabularx}{0.7\textwidth}{X} 他們說：『因為沒有人雇我們。』他說：『你們也進葡萄園去。』 \end{tabularx} \\ \\ \relax
20:8 & \begin{tabularx}{0.7\textwidth}{X} 到了晚上，園主對工頭說：『叫工人都來，給他們工錢，從後來的起，到先來的為止。』 \end{tabularx} \\ \\ \relax
20:9 & \begin{tabularx}{0.7\textwidth}{X} 約在下午五點鐘雇的人來了，各人領了一個銀幣。 \end{tabularx} \\ \\ \relax
20:10 & \begin{tabularx}{0.7\textwidth}{X} 那些最先雇的來了，以為可以多領，誰知也是各領一個銀幣。 \end{tabularx} \\ \\ \relax
20:11 & \begin{tabularx}{0.7\textwidth}{X} 他們領了工錢，就埋怨那家的主人說： \end{tabularx} \\ \\ \relax
20:12 & \begin{tabularx}{0.7\textwidth}{X} 『我們整天勞苦受熱，那些後來的只做了一小時，你竟待他們和我們一樣嗎？』 \end{tabularx} \\ \\ \relax
20:13 & \begin{tabularx}{0.7\textwidth}{X} 主人回答其中的一人說：『朋友，我沒虧待你，你與我講定的不是一個銀幣嗎？ \end{tabularx} \\ \\ \relax
20:14 & \begin{tabularx}{0.7\textwidth}{X} 拿你的錢走吧！我樂意給那後來的和給你的一樣， \end{tabularx} \\ \\ \relax
20:15 & \begin{tabularx}{0.7\textwidth}{X} 難道我的東西不可隨我的意思用嗎？因為我作好人，你就眼紅了嗎？』 \end{tabularx} \\ \\ \relax
20:16 & \begin{tabularx}{0.7\textwidth}{X} 這樣，那在後的，將要在前；在前的，將要在後了。」 \end{tabularx} \\ \\
[1ex]
\hline
\hline
\end{longtable}
(十)(十九27至二十16),論進入天國的次序
\newline
\newline
當時有少年人問主如何可得永生?在這處得永生與進入天國,同是指一件事.主要他變賣所有的財產,但他不肯,便憂憂愁愁的走了.門徒便問誰可進入天國.主回答說:那些為主犧牲的,便可以(29節).(30節)「在前的」指以色列人；「在後的」指外邦人.因以色列人棄絕主,外邦人便後來居上了；為了介紹這件事,主在二十章用比喻來解釋.
\newline
\newline
(十一)(二十17-28)論服事的真義.
\newline
\newline
(十二)(二十29-34)醫治瞎子.
\newline
\newline
國度被奪前的事蹟.(廿一1-22)
\newline
\newline
我們要時常將(廿一43)存在心裡,馬太記載三件事,指出以色列的國度,如何從他們手上奪去,賜給能結果子的百姓.
\newline
\newline
耶穌騎驢駒進入耶路撒冷
\newline
\newline
主已訓練好門徒,便進入耶路撒冷,作最後一次將自己獻給以色列,為他們的彌賽亞.主經過一番籌劃,目的為逐步應驗彌賽亞來臨應做的事,作者在四方面指出,耶穌就是以色列的彌賽亞.
\newline
\newline
靠著舊約應驗,(4節)根據(賽六十二章)和(亞九章)的記載,彌賽亞是這樣進入國家的京都的.
\newline
\newline
當時的人將衣服鋪在地上,這行動是當時歡迎王的習慣.
\newline
\newline
用力搖棕樹枝,也是當時歡迎王的傳統.
\newline
\newline
百姓歡呼「和撒那」,這是出自(詩一一八篇),這字可意譯作「天祐我王」,「大衛的子孫」及「奉主名來的」,都是彌賽亞的別名.
\newline
\newline
耶穌進入城後,全城的人都驚動,眾人都說祂不過是一位先知,可見他們沒有接受主作彌賽亞,只有少部份的人歡迎祂.
\newline
\newline
潔淨聖殿
\newline
\newline
這事顯出基督是彌賽亞,又顯出國家領袖對祂的定論,這件事有三要點:
\newline
\newline
神的殿竟淪為賊窩,可當時國家領袖對屬靈的事眼瞎.
\newline
\newline
(13節),主引用(耶七章)顯出祂有聖殿之主人的權柄,祂趕出買賣的人行動,部份應驗(瑪三1-3)的預言.
\newline
\newline
(15-16節),有小孩子稱頌主,主引用(詩八16)指出彌賽亞國度裡,連小孩也要稱頌祂,作者也指出主棄絕他們(17節).
\newline
\newline
咒詛無花果樹
\newline
\newline
通常有幾類植物,是可以用來代表以色列的；而無花果樹是最通常的一種.這件事與主潔淨聖殿,有相連的關係.兩件事指出主棄絕以色列,這段則強調主向以色列的宣佈.無花果樹只有葉沒有果子,象徵以色列只有虛偽的外表；對主沒有信心,不接受主為彌賽亞,便受審判.
\newline
\newline
國度被奪去的明證(廿一23至廿三39)
\newline
\newline
舌戰群儒(廿一23至廿二46)
\newline
\newline
主耶穌與國家領袖辯論,先與祭司長和長老辯論；前面是象徵的行動,現在是辯論的談話.他們問耶穌的權柄從何處來?(27節)指他們不承認主的權柄,可見他們不配成為國家領袖.作者用三個比喻,來解釋這件事:
\newline
\newline
1.(28-32節)兩個兒子的比喻
\newline
\newline
這比喻指以色列當中,最受人鄙視的稅吏和娼妓,可以進入天國,反而國家領袖不肯悔改.
\newline
\newline
2.(33-46節)兇惡園戶的比喻
\newline
\newline
這國家不要彌賽亞,故此國度從他們手上被奪去了.家主是神,葡萄園是以色列,園戶殺國家領袖,僕人是歷世歷代神的僕人.園戶殺死了神所派來的僕人,最後連神的兒子也殺死,所以國度不可以交給他們；而轉交教會,教會取代以色列的地位.直至外邦人數添滿了,神纔轉向以色列,叫他們全家得救.
\newline
\newline
3.(廿二1-14)娶親的筵席
\newline
\newline
這比喻與兇惡園戶比喻很相似,而這個比喻特別指出以色列將受審判.這比喻說明兩個事實:第一,以色列受審判因他們不接受耶穌為他們的彌賽亞；第二,(8節)所召的人就是以色列人,因他們不肯接受,便開始教會時代.
\newline
\newline
4.(廿二15-22)法利賽人和希律黨
\newline
\newline
他們本是世仇,但卻同心反對主.關於納稅問題,因主不能在地上建立其國度；故此教導門徒要順服地上的國度,凡屬該撒的歸給該撒.
\newline
\newline
5.(廿二23-33)與撒都該人辯論
\newline
\newline
撒都該人既不熟悉聖經,又不曉得神的大能,(29節)指出他們在屬靈事上,不配作國家領袖.
\newline
\newline
6.(廿二34-46)與法利人辯論
\newline
\newline
法利賽人派了律法師去辯論,法利賽人將神的律法編成六百一十三條,其中二百四十八條要做的；三百六十五條不可做的.他們認為人體有二百四十八個器官,而一年之內有三百六十五日,六百一十三剛好是舊約希伯來十誡字母的總數,所以編成這套律法.這些律法師精通律例,與耶穌辯論.耶穌只用幾句說話便將整個律法中心指出來,便可見主相當熟悉律法.然後主反問他們一個問題:大衛怎可以說主是彌賽亞?這件事指出,第一,彌賽亞在大衛時代已存在；現在彌賽亞道成肉身,但他們並不相信.第二,大衛也仰望彌賽亞國度來臨,但他們卻不期望.
\newline
\newline
宣告禍哉(廿三1-36)
\newline
\newline
1.(1-12節)主教導門徒防備法利賽人.
\newline
\newline
2.(13-36節)七個禍哉
\newline
\newline
七是完整的數字,可見主對他們宣佈完整的宣判,從文學角度看,這七個禍哉是用很優美的文字寫成.頭三個禍哉指法利賽人因教義而生的禍哉；五,六,七指法利賽人因行為而生的禍哉；第四個禍哉是將兩組禍哉連貫起來.主非常嚴厲責備他們,稱他們「假冒為善」有七次之多,又稱他們作「地獄之子」(15節),「瞎眼的人」(17節),「無知的人」,「粉飾的墳墓」(27節),「毒蛇的種類」(33節),從這些字彙中,可見以色列會被責打.
\newline
\newline
為城哀哭(廿三37-39)
\newline
\newline
這件事指出五點:
\newline
\newline
1.彌賽亞只能作出愛莫能助的呼叫,「耶路撒冷」這詞重覆用法,是希伯來文強調愛的呼喚的用法,可見主極愛他們.
\newline
\newline
2.耶路撒冷故意棄絕主,雖然主多次呼召它.
\newline
\newline
3.(38節)「家」指聖殿,稱以色列為「你們」,而不指「我們」,可見主放棄他們,但不是永久的丟棄,有時間性的,(39節)表示他們將來有悔改接受彌賽亞的時候.
\newline
\newline
整段哀哭,可見主與神在歷史中的工作認同一切,(37節)看見神與彌賽亞的地位是平等的.
\newline
\newline


\newpage

\section{第54屆~港九培靈會~馬有藻~第八講}
\label{sec:535}
\textbf{以色列的國度被奪去了(續講)}
\newline
\newline
連結: \href{https://www.hkbibleconference.org/session-message/view/535}{\texttt{https://www.hkbibleconference.org/session-message/view/535}}
\newline
\newline
\hyperref[sec:534]{< < < PREV SERMON < < <}
~
\hyperref[sec:toc]{[返目錄]}
~
\hyperref[sec:501]{> > > NEXT SERMON > > >}
\newline
\newline
在最後一堂聚會中,我有一點臨別依依之感!多謝大會讓我可以在港事奉,也多謝弟兄姊妹為我的身體代禱,因時間急促,未能按照預告之講題完成,請各位原諒.
\newline
\newline
(A)國度被奪後的以色列
\newline
\newline
從上文作者清楚指出,國度怎樣從以色列手上奪去,賜給那能結果子的百姓.現在他清楚指出以色列未來的遭遇,國度已經被奪去了.這兩章是主與門徒的談話,稱為「橄欖山上的論談」,這論談的主題特別提及主奪去他們的國度,和主再來當中的遭遇.這論談有預言,比喻和勸勉.這段經文分三段來看:
\newline
\newline
(一)(廿四1-31)主再來前之豫兆
\newline
\newline
(廿三39)主清楚告訴他們不得再見祂的面,直至他們歡迎祂的時候,乃是主再來的時候,主纔能在地上建立祂的國度(太廿四廿五章)其實是詳細解釋(廿三39)那句話,他們等候主的時候,以色列會有甚麼遭遇.
\newline
\newline
1.(1-14節)主再來之普遍情形
\newline
\newline
主再來前有許多事出現,第一,有假基督出現；第二,有打仗和打仗的風聲；第三,有饑荒；第四,有地震；第五,信徒被萬民恨惡；第六,有假先知出現；第七,福音要傳遍天下.這七個預兆是主再來前普遍情形,當這些事發生愈多,主再來便愈近了.
\newline
\newline
2.(15-26節)主再來時之特殊情形
\newline
\newline
這裡是特別主再來前幾年發生之事,稱為災難時期發生之事.這段和(但九章)所提及的災難是相似的,當日子愈難過時,主再來便逼近了.(15節)指一個特別時候,災難時期有行毀壞可憎踐踏聖地,這人究竟是誰,學者有不多的解釋.有人說是提多將軍在主後七十年毀壞聖殿的情形；有人說是奮銳黨和羅馬兵對抗之情形.但一切的解釋都不合(但九章)和這裡的情形,我們不相信在歷史曾應驗這件事,我們仍待將來有應驗的時期.這「毀壞可憎的」是敵基督,將來敵基督會大大反對以色列；在聖城行很多毀壞可憎的事,這是逼近主再來的情形.
\newline
\newline
3.(27-31節)主再來時之情形
\newline
\newline
主降臨是榮光充滿的日子,而且很快；當時日月無光,風雲變色；因為榮耀的王回來了.如(詩廿四篇)所說.(28節)「屍首在那裡,鷹也必在那裡」,這裡通常有幾個解釋:
\newline
\newline
第一,指有因必有果,反對彌賽亞的必有審判的結果；
\newline
\newline
第二,有解釋鷹在天上看見死屍時,便很快飛下去,這裡形容快速之意,這種解釋可以回應(27節)指主再來好像閃電那麼快.
\newline
\newline
無論怎樣解釋,都配合人子再來的情形.(30節)「人子的兆頭」當然有很多解釋,我自己以為是指人子的榮光,主誕生時天上有榮光出現,主再來時天上也有榮光出現；人子的兆頭,讓全地的人知道祂再回來了.
\newline
\newline
(二)(廿四32至廿五30)主再來之勸勉
\newline
\newline
主耶穌預告祂再來之情形,祂稍歇自己的口氣.首先用五個比喻,來應用自己所講的神學,來勸勉門徒儆醒等候祂的再來:
\newline
\newline
1.(32-33節)無花果樹的比喻
\newline
\newline
表明主再來的日子近了.
\newline
\newline
2.(34-44節)挪亞的比喻
\newline
\newline
以前挪亞時代是怎樣,人子也是怎樣,(39節)指出不知不覺.第一個比喻著重人子近了,第二個比喻著重人不知不覺間,主便來了.
\newline
\newline
3.(45-57節)僕人的比喻
\newline
\newline
這比喻指出等候主再來時,必須有忠心；不要輕忽主再來的真理,好像那惡僕忽視主人的再來.
\newline
\newline
4.(廿五1-13)十童女的比喻
\newline
\newline
這比喻指出等候主再來,需要有適當的預備,五個聰明童女有適當準備；但五個愚拙童女沒有適當的預備.(13節)主最後作結論:「所以你們要儆醒,因為那日子,那時辰,你們不知道.」
\newline
\newline
5.(14-30節)才幹的比喻
\newline
\newline
這比喻指出等候主再來,要有忠心的事奉.那領一千銀子的僕人,因為不相信主會再來；所以不肯事奉,埋沒他的才幹,後來被主責備.
\newline
\newline
(三)(廿五31-46)主再來後的情形
\newline
\newline
這是連接(四31)的禱文,兩者可連貫起來.主腳踏在地上,會再來坐在寶座上審判萬民,建立彌賽亞國度.綿羊比作信靠天國福音的人,他們的信心有行為的支持.(如40節)人怎樣對待我的弟兄如同對待主.山羊比作拒絕天國福音的人.信靠的人在彌賽亞國度有份,拒絕的人永遠要受痛苦.
\newline
\newline
第五回合:以色列的王被奪去了(廿六至廿七章)
\newline
\newline
最後的晚餐(廿六1-35)
\newline
\newline
(一)(廿六1-19)晚餐前的事
\newline
\newline
作者比照主耶穌被棄和被愛,主被猶大出賣,但被一個女人用愛心來接納.十二個使徒中的猶大,見到跟隨主的目的成為泡影,他盼望在天國中有權勢和地位.在(6-13節)記載主被一女人用香膏膏祂,這是愛的圖畫.門徒學了甚麼是浪費,奢侈,美事,愛的見證.
\newline
\newline
(二)(20-29節)晚餐時的情形
\newline
\newline
主在晚餐時宣佈兩件事
\newline
\newline
1.有人出賣祂,主從來沒有提及有人要出賣祂,可見眾門徒覺得很震驚.這出賣的人後來覺得沒趣,便離開去作出賣主的勾當；主趁他離開後,便設立晚餐.
\newline
\newline
2.祂啟示自己是新約的主人,立約是指耶穌要應驗(耶卅一章)之新約,這裡特別提及主的死,要把救恩帶給他們.
\newline
\newline
(三)(30-35節)在路上宣佈一件事實
\newline
\newline
主提及彼得將會否認主,所有的門徒都會跌倒(31節),不得與主站立在一起.
\newline
\newline
(B)客西馬尼園的事蹟(廿六36-46)
\newline
\newline
主在客西馬尼園有很大的掙扎,但結果祂得勝,然後祂被出賣；主表現出羔羊的態度,不反抗,受苦的時候也不開口.
\newline
\newline
(C)主耶穌被審(廿六47-廿七26)
\newline
\newline
祂先在大祭司該亞法,又在巡撫彼拉多面前受審；主在受審期間,多次稱自己為彌賽亞.
\newline
\newline
(D)主耶穌被釘(廿七27-56)
\newline
\newline
他們先戲弄主,用荊棘作冠冕,穿上紫袍,用葦子作權杖.叫祂背負自己的十字架,證實自己是死囚；又是政治犯,十字架是死刑刑具.從十字架上的牌子,兵丁戲弄的用具,十字架旁的強盜,馬太一而再指出耶穌有彌賽亞的質素.主從受審到死亡有三十三個小時沒有歇息過；特別最後十四小時,擔負心靈和肉體的痛苦.主卻念念不忘旁邊悔改的強盜,和十架下的母親；可見主的愛是莫測的,祂的所行都充滿了饒恕.這時候,自然界也產生不平凡的景象；本來是烈日當空,但卻變成烏天黑地,它們好像為光明之主提出抗議,另外還有兩個神蹟出現:
\newline
\newline
(一)聖殿的幔子裂開,有何人可以撕開這幔子,除非是神蹟；這件事表明舊約時代的結束,人可以來到神的面前.
\newline
\newline
(二)當時耶路撒冷的墳墓裂開,有許多死人復活進城向多人見證,表明主的死是得勝的.
\newline
\newline
(E)主耶穌被埋葬(廿七57-66)
\newline
\newline
愛主的人把主放在墳墓,但仍在主墳墓前徘徊不去.但那些反對主的人恐怕主復活,想盡方法,要禁止門徒把主的身體偷去.
\newline
\newline
第六回合:以色列的王回天國了(廿八章)
\newline
\newline
這一章可算是跋語,復活的主與門徒進入教會的時代,這章可分三點來看:
\newline
\newline
主的復活(廿八1-10)
\newline
\newline
主勝過死亡的毒勾,陰間的權柄不能勝過祂.根據彼得在(徒二章)看,他用主的復活來證明主是王.
\newline
\newline
反復活謊言的捏造(11-15節)
\newline
\newline
(11節)見證墳墓空的,不單是信徒,也是兵丁；但他們後來受賄賂蒙蔽,把事實歪曲.
\newline
\newline
王的使命(16-20節)
\newline
\newline
主復活後向門徒顯現的有十或十一次,現在這是第八次與門徒聚集在加利利.主將大使命交託門徒,這大使命在五卷福音裡.(四福音及使徒行傳)這五卷書都用不同方式來提及,可見傳福音不是個人的專利,而是每一個屬主的人的責任.
\newline
\newline
馬太全書提及主耶穌是猶太人之王,祂也是萬王之王.雖然現在祂不能在地上成為萬王之王,但今天祂是我們心中和生命的王.若果主不能在我們身上作王,(羅五章)指出死或罪便會在我們心中作王.願一切歸屬主的人,如(羅六12-13)所講:「不要容罪在你們必死的身上作王......」,願所有信徒的一生是王的傳記,正如馬太福音是王的傳記一樣.
\newline
\newline


\newpage

\section{第55屆~港九培靈會~劉承業~第一講}
\label{sec:501}
\textbf{導論與引言}
\newline
\newline
連結: \href{https://www.hkbibleconference.org/session-message/view/501}{\texttt{https://www.hkbibleconference.org/session-message/view/501}}
\newline
\newline
\hyperref[sec:535]{< < < PREV SERMON < < <}
~
\hyperref[sec:toc]{[返目錄]}
~
\hyperref[sec:502]{> > > NEXT SERMON > > >}
\newline
\newline
傳道書 1:1-1:11
\newline
\begin{longtable}{cl}
\hline
\hline
章節 & 經文 (和合本修訂版)\\
\hline
1:1 & \begin{tabularx}{0.7\textwidth}{X} 在耶路撒冷作王、大衛的兒子、傳道者的言語。 \end{tabularx} \\ \\ \relax
1:2 & \begin{tabularx}{0.7\textwidth}{X} 傳道者說：虛空的虛空，虛空的虛空，全是虛空。 \end{tabularx} \\ \\ \relax
1:3 & \begin{tabularx}{0.7\textwidth}{X} 人一切的勞碌，就是他在日光之下的勞碌，有甚麼益處呢？ \end{tabularx} \\ \\ \relax
1:4 & \begin{tabularx}{0.7\textwidth}{X} 一代過去，一代又來，地卻永遠長存。 \end{tabularx} \\ \\ \relax
1:5 & \begin{tabularx}{0.7\textwidth}{X} 太陽上升，太陽落下，急歸所出之地。 \end{tabularx} \\ \\ \relax
1:6 & \begin{tabularx}{0.7\textwidth}{X} 風往南颳，又向北轉，不停旋轉，繞回原路。 \end{tabularx} \\ \\ \relax
1:7 & \begin{tabularx}{0.7\textwidth}{X} 江河都往海裡流，海卻不滿；江河從何處流，仍歸回原處。 \end{tabularx} \\ \\ \relax
1:8 & \begin{tabularx}{0.7\textwidth}{X} 萬事令人厭倦，人不能說盡。眼看，看不飽；耳聽，聽不足。 \end{tabularx} \\ \\ \relax
1:9 & \begin{tabularx}{0.7\textwidth}{X} 已有的事，後必再有；已行的事，後必再行。日光之下並無新事。 \end{tabularx} \\ \\ \relax
1:10 & \begin{tabularx}{0.7\textwidth}{X} 有一件事人指著說：「看，這是新的！」它在我們以前的世代早已有了。 \end{tabularx} \\ \\ \relax
1:11 & \begin{tabularx}{0.7\textwidth}{X} 已過的事，無人記念；將來的事，後來的人也不記念。 \end{tabularx} \\ \\
[1ex]
\hline
\hline
\end{longtable}
我這次選擇傳道書與大家研討,主要的是有三個原因；第一,這卷書的信息很容易被人誤解,因為若只從消極方面來看人生,在日光之下所過的生活是毫無意義的；其實,這卷書是教導我們,如何在日光之下過一個有意義的人生.第二,我發現近來有不少青年人尋求生命的意義,願意奉獻自己作神的工作；因為明白事奉神,在人生中,是最有意義的.從這卷書可以幫助我們找出何謂最有意義的人生.聞說有一個德國人,後來被告反而被捕入獄,但他仍然堅持納粹黨的思想是最好的；雖然受盡痛苦,也處之泰然.及後期滿出獄返家,他才發現原來告密者乃是他的兒子.以致令他感覺到做人沒有意義,因為連骨肉也毫無親情,故此自殺而死.起初這人雖然長久坐監也不怕,因他感覺人生有意義這就證明一個有意義的人生,對人的影響是何等的重要.為甚麼有許多人要自殺,可能只為了某些小事情,而感覺人生沒有意義.有人自殺是因失戀,也有人是為了他的理想和信仰而死,可見人生的意義,有高級和低級之分.如今我們要從傳道書中找出人生最高尚的意義何在,好叫我們做人更加有力量,更加有勁.第三,近年興起一種查經的方法──是福音性的查經；有不少教會使用這種方法,使多人得救.但可惜查經時經常苦無材料,我個人就認為傳道書是福音性查經的最好材料.雖然內文沒有提及耶穌基督,但它卻提及到我們與神的關係.這卷書的信息,可以預備人的心來接受福音.如果弟兄姊妹能夠明白這卷書,這樣也就可以與未信主的親友同事,來考查這卷書,預備他們的心了.
\newline
\newline
導言:
\newline
\newline
一,主要信息,作者指出如何過著有意義的人生,他指出有幾個人生真理體會:
\newline
\newline
a.日光之下一切都是虛空的
\newline
\newline
b.敬畏神的人,在日光之下能夠過快樂而有意義的生活；特別是可以享受神所賜予的一切；(三12-13,五18-19,八15,九7)以及有工作和事奉的機會(九10)
\newline
\newline
簡單而言,有意義的人生,是敬畏神的人生.(十二13)我覺得這信息特別適合現代的人.正如在香港,很多人覺得人生沒有樂趣,要面對各種困難和壓力,前途又不明朗,人生找不到出路.傳道書的信息,相信會幫助我們,讓我們看見一個敬畏神的人,在一個動盪,無意義,無出路的世代中,可以過著快樂和有意義的人生.另外是現代人以為人生最好是享清福,不用工作是最舒服的；但按傳道書所講的,這觀念並不正確!這並不是神的心意,因為神的心意是要每一個人都做工,能作工的,才是神所賜的福份.
\newline
\newline
二,了解傳道書的秘訣
\newline
\newline
a.認識傳道書是希伯來的智慧文學
\newline
\newline
希伯來智慧文學,最少有三個特點:第一,神是智慧的源頭,特別是箴言書常提及的,神統管一切.真正的智慧是從神而來的,認識神的人就是一個智慧的人.第二,智慧文學注重現實的問題.約伯記提及苦難,傳道書也一樣,講求現實的生活.第三,任何人都要知道怎樣去做人.本書對任何膚色的人也適合,因為它的內容普及化.希伯來人不注重理論,而注重現實生活；要在現實生活中探討和體驗,所以有學者稱智慧文學為「神學性的人類學」.本書內容提及人的生活,但這個人是不能脫離神的；表面上是以人為主,其實是以神為主.人不可脫離神而存在,必須有神,才有人生的真正意義.研究方法,是由地到天,由人到神；從現世到永世,從現實到超現實.故此研究智慧文學,是由地上帶領你到天上,由屬世帶你到屬靈.
\newline
\newline
b.傳道書常提及作者對人生的觀察
\newline
\newline
但觀察的結果不全是真理,例如(九11-12),以為作者贊成機會主義；其實在第三章他提及萬事有神的安排,他只不過在這裡指出人生百態之一；甚至有些觀察是代表作者一時的感受,而不能代表神的心意和啟示.
\newline
\newline
c.我們要認識傳道書作者的人生觀.
\newline
\newline
有人以為作者是悲觀者,因為他提及虛空的虛空等論調,也有人以為作者是樂天派,因為他說吃喝快樂.其實作者是一個現實主義者,他看清楚人生是怎樣的,而他也知道如何超脫虛空和無意義的人生,藉著神的啟示和教導,可使我們過有意義和豐盛的人生.
\newline
\newline
我們現在看(一1-11),這段經文是傳道書的引言.第一節指出作者是誰,他在耶路撒冷作王,是大衛的兒子,因此我個人相信他是所羅門.本書作者所提及的經歷,例如大有智慧,大有財富；財富無人可比,可見這人是所羅門.為何他不用自己的姓名,而說傳道者(3節),在箴言中他有出示自己的名字,這是因為箴言是教導自己的兒女,不必隱名；但現今他寫傳道書,對象是一般大眾,可能他用傳道者的名稱,是更加有份量.因為傳道者是代表神而說話,所羅門看重傳道者過於王的身份.今天我們是否看重自己基督徒的身份,抑或自己社會地位的身份呢?如果我們平時的身份.會影響到基督徒的身份,我們就必須避免.所羅門是看重自己在神面前的身份,過於作工的身份.有一位醫學博士進入神學院受造就,但每次簽名後仍加上醫學博士的銜頭；可見他雖為神學士,仍念念不忘他的過往身份.後來他因對神學院不滿,便退學了；因他仍看重世界的身份,過於在神面前的身份.究竟我們看重世上的身份,還是在基督的身份呢?
\newline
\newline
第二節,提及虛空,這名詞在全卷書出現了三十一次；從用字的次數,可見意思是相當重要的.虛空原意是一口氣或霧水,所以雅各書四章十二節的雲霧與虛空是同一個字；這名詞是代表短暫,空洞,徒勞無功,無果效,無實質.虛空的虛空是指人生非常的空虛,為何人生會空虛呢?在引言中他最少提及四個理由.第一,勞碌而無益處(3節),他在日光之下勞碌.有人解經說在日光之上便不勞碌,這種說法是不當的；因為全卷書從未提及過日光之上.作者要指出我們仍然可以在日光之下,過勞碌而有益的人生.第二,第四至六節,人生很短暫,故此人生很虛空,「一代過去,一代又來」人生變幻無窮.在歷史舞台上,看見一個朝代過去,另一個朝代又來.第五節用自然界的變幻來比喻人生的短暫,十二章7節與這裡成為迴響.因為人生會死亡,做成虛空的原因.第三,七至八節,這裡指沒有滿足的情況,人生沒有滿足,并非他沒有財富,地位,俗語說:「知足窮亦樂」.人生的虛空,不是窮苦大眾專利,連大富之人也一樣.人生虛空是人性的問題,以後傳道者提及人莫如吃喝快樂,指出人在簡單的事情,如吃飯是可以滿足,只要他有神在他裡面.第四,九至十一節,這裡指出任何事情都舊調重彈,無過去將來之分,使人厭悶之感.沒有新奇之事,難怪實存主義者提出人生苦悶；但傳道者多次提及「身後」這名詞.這名詞指將來,不一定指死之後,在傳道者的觀念中,不是沒有過去和將來之分；在十一和十二章,他提及少壯和年老之分,他指出後來一切所作的,神必審問,神會記念的.因此傳道者的觀念并非如第一章這樣的消極,他教導我們怎樣超脫虛空的人生,如何過一個有意義的人生.
\newline
\newline


\newpage

\section{第55屆~港九培靈會~劉承業~第二講　體驗人生的虛空}
\label{sec:502}
\textbf{第二講　體驗人生的虛空}
\newline
\newline
連結: \href{https://www.hkbibleconference.org/session-message/view/502}{\texttt{https://www.hkbibleconference.org/session-message/view/502}}
\newline
\newline
\hyperref[sec:501]{< < < PREV SERMON < < <}
~
\hyperref[sec:toc]{[返目錄]}
~
\hyperref[sec:503]{> > > NEXT SERMON > > >}
\newline
\newline
傳道書 1:12-2:26
\newline
\begin{longtable}{cl}
\hline
\hline
章節 & 經文 (和合本修訂版)\\
\hline
1:12 & \begin{tabularx}{0.7\textwidth}{X} 我傳道者在耶路撒冷作過以色列的王。 \end{tabularx} \\ \\ \relax
1:13 & \begin{tabularx}{0.7\textwidth}{X} 我用智慧專心探尋、考察天下所發生的一切事：神給世人何等沉重的擔子，使他們在其中勞苦！ \end{tabularx} \\ \\ \relax
1:14 & \begin{tabularx}{0.7\textwidth}{X} 我見日光之下所發生的一切事，看哪，全是虛空，全是捕風。 \end{tabularx} \\ \\ \relax
1:15 & \begin{tabularx}{0.7\textwidth}{X} 彎曲的，不能變直；缺乏的，不計其數。 \end{tabularx} \\ \\ \relax
1:16 & \begin{tabularx}{0.7\textwidth}{X} 我心裡說：「看哪，我大有智慧，勝過在我以前所有統治耶路撒冷的人；我的心也多經歷智慧和知識的事。」 \end{tabularx} \\ \\ \relax
1:17 & \begin{tabularx}{0.7\textwidth}{X} 我專心想要明白智慧，想要明白狂妄與愚昧，方知這也是捕風。 \end{tabularx} \\ \\ \relax
1:18 & \begin{tabularx}{0.7\textwidth}{X} 因為多有智慧，就多有愁煩；增加知識，就增加憂傷。 \end{tabularx} \\ \\ \relax
2:1 & \begin{tabularx}{0.7\textwidth}{X} 我心裡說：「來吧，讓我用喜樂試試你，使你享福！」看哪，這也是虛空。 \end{tabularx} \\ \\ \relax
2:2 & \begin{tabularx}{0.7\textwidth}{X} 論嬉笑，我說：「這是狂妄。」論享樂，「這有甚麼用呢？」 \end{tabularx} \\ \\ \relax
2:3 & \begin{tabularx}{0.7\textwidth}{X} 我心以智慧引導我，我心裡探究，如何用酒使身體舒暢，如何抓住愚昧，直等我看明世人在天下短暫一生中，當行何事為美。 \end{tabularx} \\ \\ \relax
2:4 & \begin{tabularx}{0.7\textwidth}{X} 我大興土木，為自己建造房屋，栽葡萄園， \end{tabularx} \\ \\ \relax
2:5 & \begin{tabularx}{0.7\textwidth}{X} 修造庭園和公園，在其中栽種各樣果樹， \end{tabularx} \\ \\ \relax
2:6 & \begin{tabularx}{0.7\textwidth}{X} 挖造水池，用以灌溉林中的幼樹。 \end{tabularx} \\ \\ \relax
2:7 & \begin{tabularx}{0.7\textwidth}{X} 我買了僕婢，也有生在家中的僕婢；又有許多牛群羊群，勝過我以前所有在耶路撒冷的人。 \end{tabularx} \\ \\ \relax
2:8 & \begin{tabularx}{0.7\textwidth}{X} 我為自己積蓄金銀，搜集各君王、各省份的財寶；又為自己得男女歌手和世人所喜愛的物，以及一個又一個的妃嬪。 \end{tabularx} \\ \\ \relax
2:9 & \begin{tabularx}{0.7\textwidth}{X} 這樣，我就日漸昌盛，勝過我以前所有在耶路撒冷的人。我的智慧仍然存留。 \end{tabularx} \\ \\ \relax
2:10 & \begin{tabularx}{0.7\textwidth}{X} 凡我眼所求的，我沒有克制它；我心所樂的，我沒有不享受。因我的心要為一切的勞碌快樂，這是我從一切勞碌中所得的報償。 \end{tabularx} \\ \\ \relax
2:11 & \begin{tabularx}{0.7\textwidth}{X} 後來，我回顧我手所經營的一切和我勞碌所做的工。看哪，全是虛空，全是捕風；在日光之下毫無益處。 \end{tabularx} \\ \\ \relax
2:12 & \begin{tabularx}{0.7\textwidth}{X} 我轉而回顧智慧、狂妄和愚昧。在王以後來的人又如何呢？不過做先前所做的就是了。 \end{tabularx} \\ \\ \relax
2:13 & \begin{tabularx}{0.7\textwidth}{X} 於是我看出智慧勝過愚昧，如同光明勝過黑暗。 \end{tabularx} \\ \\ \relax
2:14 & \begin{tabularx}{0.7\textwidth}{X} 智慧人的眼目光明，愚昧人卻在黑暗裡行。但我知道他們都有相同的遭遇。 \end{tabularx} \\ \\ \relax
2:15 & \begin{tabularx}{0.7\textwidth}{X} 我心裡就說：「愚昧人所遇見的，我也一樣遇見，那麼我何必更有智慧呢？」我心裡說：「這也是虛空。」 \end{tabularx} \\ \\ \relax
2:16 & \begin{tabularx}{0.7\textwidth}{X} 智慧人和愚昧人一樣，不會長久被人記念，因為日後都被遺忘。可嘆！智慧人和愚昧人都一樣會死亡。 \end{tabularx} \\ \\ \relax
2:17 & \begin{tabularx}{0.7\textwidth}{X} 於是我恨惡生命，因為在日光之下所發生的事我都以為煩惱，全是虛空，全是捕風。 \end{tabularx} \\ \\ \relax
2:18 & \begin{tabularx}{0.7\textwidth}{X} 我恨惡一切的勞碌，就是我在日光之下所勞碌的，因為我所得的必須留給我以後的人。 \end{tabularx} \\ \\ \relax
2:19 & \begin{tabularx}{0.7\textwidth}{X} 那人是智慧是愚昧，誰能知道呢？他竟要掌管我在日光之下用智慧勞碌所得的。這也是虛空。 \end{tabularx} \\ \\ \relax
2:20 & \begin{tabularx}{0.7\textwidth}{X} 我轉想我在日光之下所勞碌的一切工作，心就絕望。 \end{tabularx} \\ \\ \relax
2:21 & \begin{tabularx}{0.7\textwidth}{X} 因為有人用智慧、知識、靈巧勞碌工作，所得來的卻要遺留給未曾勞碌的人作產業。這也是虛空，大大不幸。 \end{tabularx} \\ \\ \relax
2:22 & \begin{tabularx}{0.7\textwidth}{X} 人一切的勞碌操心，就是他在日光之下所勞碌的，又得著了甚麼呢？ \end{tabularx} \\ \\ \relax
2:23 & \begin{tabularx}{0.7\textwidth}{X} 他日日憂慮，他的勞苦成為愁煩，連夜間心也不得休息。這也是虛空。 \end{tabularx} \\ \\ \relax
2:24 & \begin{tabularx}{0.7\textwidth}{X} 難道一個人有吃有喝，且在勞碌中享福，不是福氣嗎？我看這也是出於神的手。 \end{tabularx} \\ \\ \relax
2:25 & \begin{tabularx}{0.7\textwidth}{X} 論到吃用、享福，誰能勝過我呢？ \end{tabularx} \\ \\ \relax
2:26 & \begin{tabularx}{0.7\textwidth}{X} 神喜愛誰，就給誰智慧、知識和喜樂；惟有罪人，神使他勞苦，將他所儲藏、所堆積的歸給神所喜愛的人。這也是虛空，也是捕風。 \end{tabularx} \\ \\
[1ex]
\hline
\hline
\end{longtable}
作者在這段經文,指出他對人生有消極和積極兩種的體會:
\newline
\newline
(一)消極的體會
\newline
\newline
他首先指出人的智慧和知識都是虛空的.(一12-18,二12-17)所羅門蒙神所賜知識和智慧,他是大有智慧的人.外國有諺語說:「好似所羅門有智慧.」雖然所羅門王很有智慧,但他仍說這是虛空的.首先我們要看,何謂智慧,有人解釋智慧是成功生活的藝術；有人說是正確生活的知識,有人說是生活的經驗.可見智慧與生活相關,智慧特別強調,人與人以及人與神的關係.智慧與知識有所不同,有知識的人未必會做人,未必與神與人有美好關係.所羅門具有智慧與知識,這是很難得的；但他仍說這是虛空,到底是因為甚麼呢?
\newline
\newline
多有知識,多有愁煩(一12-14,17-18)我們常說有智慧的人,能先知先覺.但這種人也是「先天下之憂而憂,後天下之樂而樂.」有些弱智的人不會憂愁,也許神在這方面彌補他們的缺憾,反而有思想的人卻會憂愁.我們雖然知道事情的演變,但卻愛莫能助.(一15)惟有與面對的事一起憂愁,處於無奈情況之下.)
\newline
\newline
1.智慧人與愚昧人的遭遇都是虛空(二14-16)大家都會死亡,智慧的人不一定長壽,智慧與長壽無關,那麼智慧的人,比愚昧的人又強得多少呢?我以前遇見過一個可憐的媽媽,他有一兒子,在美國麻省理工學院得到博士學位,翌日便要舉行畢業典禮,可是卻在前一天的晚上因駕車失事而身亡.可見有學問的人也一樣會死亡.前兩晚在培靈會舉行之工作人員聯禱會中,聚餐時有位前輩講到一段新聞,他說:有一個在內地很出名的眼科醫生,來到香港後,卻要在酒樓做售點心員；這個人身份的急劇轉變,大家聽了也為之心酸.故此有智慧的人,遭遇不一定會比別人好.美籍猶太裔物理學家愛因斯坦,他發明相對論,臨終時卻說:「我感覺自己是一個被捆綁的人,看了真實的一眼,它飛跑了；但願我能脫離渺小知識的束縛,然後才能認識廣闊的宇宙.」他是學術方面的巨人,但仍然感覺自己的知識渺小；因此所羅門王這說法,我們也應該認同.智慧和知識都是好的,但須要服在神掌管之下,否則也都是虛空.所羅門王提醒我們具有的智商學問,如果不能帶領我們來到神的面前,可能成為我們的攔阻；因為「智者千慮,必有一失.」所以我們不能誇耀自己的知識.
\newline
\newline
2.享樂和財富,都是虛空(二1-11).這是所羅門王親身的經驗.(王上十23)「所羅門王的財寶與智慧,勝過天下的列王.」(王上十27)「王在耶路撒冷,使銀子多如石頭」(王上十一3)「所羅門有妃七百,還有嬪三百.」所羅門王享盡人間一切的福樂,我們沒有一人能及他,但他仍說:「後來我察看我手所經營的一切事,和我勞碌所成的功,誰知都是虛空,都是捕風,在日光之下毫無益處.」(傳二11),即使他有許多食物,但胃只能容納有限的食物；他的一雙眼所看有限,作為一個過來人的他,說的話是具有份量的.(二章3節9節)他指出自己仍是有智慧和清醒的.有解經家指出他故意讓自己受試驗,來看這些東西可否給他滿足.但我個人認為他不是作實驗,而是隨自己有許多享受,感覺到一切不外如是；便將自己的體驗,與讀者分享.幾年前,互愛團契曾作過一項調查,從資料顯示,指出香港青少年深受享樂主義所影響.他們響往名牌運動鞋,名廠服裝──為要得到這些享受,以致有些少女甚至不惜出賣肉體,賺取金錢,以求享受.如今這種享樂主義的風氣,便吹到教會青少年人的圈子裡了.所以數年前在洛桑的世界福音會議中,有見及此,就提出了有實行簡樸生活的必要.今天香港的教會實在有此需要,應該好好地教導信徒實行簡樸的生活.不久之前在美國有個億萬的富翁,竟然因營養不足致死,這真是天下的大笑話.其實他是害怕別人在他的食物中下毒,所以甚麼都不敢食,因餓而致死.雖然他有許多財富,但卻不能享用.俗語說:「富貴如浮雲.」豈不正對嗎?許多有錢的人,因為戰亂或政治因素,要淪落異鄉,因此財富也是虛空的.
\newline
\newline
3.勞碌也屬虛空(二18-23)作者不是怕累,而是有兩個理由感到虛空.第一,他不能長久享受勞碌的成果.(二18)財富若果能留給孝子賢孫是好的,但若留給敗家的二世祖或是敵人,便不值得了,這是所羅門王的感受.他不知道自己的兒子繼位後會怎樣.而歷史告訴我們,他的兒子羅波安作王後,苦待百姓.所羅門實在擔心後來者能否守業,所以感覺勞碌都是虛空.第二,勞碌使人憂慮愁煩.(二22-23)使人睡不安,有些人往上拚命爬,有錢後,便擔心遭人綁架；或怕廉政公署調查,所以心裡很害怕.我有一親戚,積蓄了一筆錢,就買了一層樓；他太太說他買樓後,平添了不少白髮,因為他怕租給人住人欠租,又怕火災,風災和賊劫.過勞會傷害人的身體,有人開農場養雞,每日早上在農場走一轉取雞蛋,便是午飯時候；午飯後再走一轉,又到晚飯時候.他的感覺只為雞蛋而活,生活毫無意義.
\newline
\newline
(二)積極的體會(二24-26)
\newline
\newline
作者提及虛空的突破,這裡指出有兩種人,即神所喜悅的人和罪人.我們看見有些人積蓄許多財寶,又有智慧,但他們是虛空的；然而又有些人跟他們一樣,卻可以吃用享福,為何他們有這麼大的分別呢?原因是神喜悅他們與否.罪人雖然有錢有地位,但仍是虛空捕風.其實未信主者都知道,真正福氣是由神而來,許多人在門前張貼揮春「天官賜福」.我們基督徒也該明白,福氣是由神而來.這裡提出一種討神喜悅的人生,得神喜悅的條件是甚麼呢?
\newline
\newline
1.信靠神(來十一章6節):「人非有信,就不能得神的喜悅」,今天我們只要信靠神,便能成為神所喜悅的人.神會按照我們一切的需要賜給我們,如果我們有信,便可以將罪人的身份變成為天之驕子(約一12).
\newline
\newline
2.順服神我們凡事順服神,便可得到神賜我們的福樂.天主教神學研究一項倫理問題,就是甚麼時候才是得罪神,若是小罪便要放在煉獄,需要煉清,若是大罪便不得了.今天討神喜悅,不是指我們合格與否這個標準.討神喜悅的人,是盼望得到滿分的；討神的心歡喜,凡事上順服祂.這種人才能吃用享福,連罪人所得的,也轉到他們身上(二26),但願神幫助我們作祂所喜悅的人,便能突破一般人在世間上的虛空了.
\newline
\newline


\newpage

\section{第55屆~港九培靈會~劉承業~第三講　觀察人生的虛空}
\label{sec:503}
\textbf{第三講　觀察人生的虛空}
\newline
\newline
連結: \href{https://www.hkbibleconference.org/session-message/view/503}{\texttt{https://www.hkbibleconference.org/session-message/view/503}}
\newline
\newline
\hyperref[sec:502]{< < < PREV SERMON < < <}
~
\hyperref[sec:toc]{[返目錄]}
~
\hyperref[sec:504]{> > > NEXT SERMON > > >}
\newline
\newline
傳道書 3:1-5:7
\newline
\begin{longtable}{cl}
\hline
\hline
章節 & 經文 (和合本修訂版)\\
\hline
3:1 & \begin{tabularx}{0.7\textwidth}{X} 凡事都有定期，天下每一事務都有定時。 \end{tabularx} \\ \\ \relax
3:2 & \begin{tabularx}{0.7\textwidth}{X} 生有時，死有時；栽種有時，拔出有時； \end{tabularx} \\ \\ \relax
3:3 & \begin{tabularx}{0.7\textwidth}{X} 殺戮有時，醫治有時；拆毀有時，建造有時； \end{tabularx} \\ \\ \relax
3:4 & \begin{tabularx}{0.7\textwidth}{X} 哭有時，笑有時；哀慟有時，跳舞有時； \end{tabularx} \\ \\ \relax
3:5 & \begin{tabularx}{0.7\textwidth}{X} 丟石頭有時，撿石頭有時；懷抱有時，不抱有時； \end{tabularx} \\ \\ \relax
3:6 & \begin{tabularx}{0.7\textwidth}{X} 尋找有時，失落有時；保存有時，拋棄有時； \end{tabularx} \\ \\ \relax
3:7 & \begin{tabularx}{0.7\textwidth}{X} 撕裂有時，縫補有時；沉默有時，說話有時； \end{tabularx} \\ \\ \relax
3:8 & \begin{tabularx}{0.7\textwidth}{X} 喜愛有時，恨惡有時；戰爭有時，和平有時。 \end{tabularx} \\ \\ \relax
3:9 & \begin{tabularx}{0.7\textwidth}{X} 這樣，做事的人在他所勞碌的事上得到甚麼益處呢？ \end{tabularx} \\ \\ \relax
3:10 & \begin{tabularx}{0.7\textwidth}{X} 我觀看神給世人的擔子，使他們在其中勞苦： \end{tabularx} \\ \\ \relax
3:11 & \begin{tabularx}{0.7\textwidth}{X} 神造萬物，各按其時成為美好，又將永恆安放在世人心裡；然而神從始至終的作為，人不能測透。 \end{tabularx} \\ \\ \relax
3:12 & \begin{tabularx}{0.7\textwidth}{X} 我知道，人除了終身喜樂納福，沒有一件幸福的事。 \end{tabularx} \\ \\ \relax
3:13 & \begin{tabularx}{0.7\textwidth}{X} 並且人人吃喝，在他的一切勞碌中享福，這也是神的賞賜。 \end{tabularx} \\ \\ \relax
3:14 & \begin{tabularx}{0.7\textwidth}{X} 我知道神所做的都必存到永遠；無所增添，無所減少。神這樣做，是要人在他面前存敬畏的心。 \end{tabularx} \\ \\ \relax
3:15 & \begin{tabularx}{0.7\textwidth}{X} 現今的事以前就有了，將來的事也早已有了，並且神使已過的事重新再來。 \end{tabularx} \\ \\ \relax
3:16 & \begin{tabularx}{0.7\textwidth}{X} 我又見日光之下，應有公平之處有奸惡，應有公義之處也有奸惡。 \end{tabularx} \\ \\ \relax
3:17 & \begin{tabularx}{0.7\textwidth}{X} 我心裡說：「神必審判義人和惡人，因為在那裡，各樣事務，一切工作，都有定時。」 \end{tabularx} \\ \\ \relax
3:18 & \begin{tabularx}{0.7\textwidth}{X} 我心裡說：「為世人的緣故，神考驗他們，讓他們看見自己不過像走獸一樣。」 \end{tabularx} \\ \\ \relax
3:19 & \begin{tabularx}{0.7\textwidth}{X} 因為世人遭遇的，走獸也遭遇，所遭遇的都一樣：這個怎樣死，那個也怎樣死，他們都有一樣的氣息。人不能強於走獸，全是虛空； \end{tabularx} \\ \\ \relax
3:20 & \begin{tabularx}{0.7\textwidth}{X} 都歸一處，都是出於塵土，也都歸於塵土。 \end{tabularx} \\ \\ \relax
3:21 & \begin{tabularx}{0.7\textwidth}{X} 誰知道人的氣息是往上升，走獸的氣息是下入地呢？ \end{tabularx} \\ \\ \relax
3:22 & \begin{tabularx}{0.7\textwidth}{X} 總而言之，人能夠在他經營的事上喜樂，是最好不過了，因為這是他應得的報償。他身後的事誰能領他回來看呢？ \end{tabularx} \\ \\ \relax
4:1 & \begin{tabularx}{0.7\textwidth}{X} 我轉而觀看日光之下所發生的一切欺壓之事。看哪，受欺壓的流淚，無人安慰；欺壓他們的有權勢，也無人安慰。 \end{tabularx} \\ \\ \relax
4:2 & \begin{tabularx}{0.7\textwidth}{X} 因此，我讚歎那已死的死人，勝過那還活著的活人。 \end{tabularx} \\ \\ \relax
4:3 & \begin{tabularx}{0.7\textwidth}{X} 但那尚未出生，就是未曾見過日光之下所發生之惡事的，比這兩種人更幸福。 \end{tabularx} \\ \\ \relax
4:4 & \begin{tabularx}{0.7\textwidth}{X} 我見人因彼此嫉妒而有一切的勞碌和各樣工作的成就，這也是虛空，也是捕風。 \end{tabularx} \\ \\ \relax
4:5 & \begin{tabularx}{0.7\textwidth}{X} 愚昧人抱著雙臂，自食其肉。 \end{tabularx} \\ \\ \relax
4:6 & \begin{tabularx}{0.7\textwidth}{X} 一掌滿滿而得享安靜，勝過兩掌滿滿而勞碌捕風。 \end{tabularx} \\ \\ \relax
4:7 & \begin{tabularx}{0.7\textwidth}{X} 我轉而觀看日光之下有一件虛空的事： \end{tabularx} \\ \\ \relax
4:8 & \begin{tabularx}{0.7\textwidth}{X} 有人孤單無雙，無子無兄弟，竟勞碌不息，眼目也不以財富為滿足。他說：「我勞碌，自己卻不享福，到底是為了誰呢？」這也是虛空，是極沉重的擔子。 \end{tabularx} \\ \\ \relax
4:9 & \begin{tabularx}{0.7\textwidth}{X} 兩個人總比一個人好，他們勞碌同得美好的報償。 \end{tabularx} \\ \\ \relax
4:10 & \begin{tabularx}{0.7\textwidth}{X} 若是跌倒，這人可以扶起他的同伴；倘若孤身跌倒，沒有別人扶起他來，這人就有禍了。 \end{tabularx} \\ \\ \relax
4:11 & \begin{tabularx}{0.7\textwidth}{X} 再者，二人同睡就都暖和，一人獨睡怎能暖和呢？ \end{tabularx} \\ \\ \relax
4:12 & \begin{tabularx}{0.7\textwidth}{X} 若遇敵攻擊，孤身難擋，二人就能抵擋他；三股合成的繩子不易折斷。 \end{tabularx} \\ \\ \relax
4:13 & \begin{tabularx}{0.7\textwidth}{X} 貧窮而有智慧的年輕人，勝過年老不再納諫的愚昧王， \end{tabularx} \\ \\ \relax
4:14 & \begin{tabularx}{0.7\textwidth}{X} 那人從監牢裡出來作王，在國中原是出身貧寒。 \end{tabularx} \\ \\ \relax
4:15 & \begin{tabularx}{0.7\textwidth}{X} 我見日光之下所有行走的活人，都跟隨那年輕人，就是接續作王的那位。 \end{tabularx} \\ \\ \relax
4:16 & \begin{tabularx}{0.7\textwidth}{X} 他的百姓，就是他所治理的眾人，多得無數；但後來的人還是不喜歡他。這也是虛空，也是捕風。 \end{tabularx} \\ \\ \relax
5:1 & \begin{tabularx}{0.7\textwidth}{X} 你到神的殿要謹慎你的腳步；近前聽，勝過愚昧人獻祭，他們不知道自己在作惡。 \end{tabularx} \\ \\ \relax
5:2 & \begin{tabularx}{0.7\textwidth}{X} 在神面前你不可冒失開口，也不可心急發言；因為神在天上，你在地上，所以你的話語要少。 \end{tabularx} \\ \\ \relax
5:3 & \begin{tabularx}{0.7\textwidth}{X} 事務多，令人做夢；話語多，顯出愚昧。 \end{tabularx} \\ \\ \relax
5:4 & \begin{tabularx}{0.7\textwidth}{X} 你向神許願，還願不可遲延，因他不喜歡愚昧人，你許的願應當償還。 \end{tabularx} \\ \\ \relax
5:5 & \begin{tabularx}{0.7\textwidth}{X} 你許願不還，不如不許。 \end{tabularx} \\ \\ \relax
5:6 & \begin{tabularx}{0.7\textwidth}{X} 不可放任你的口使肉體犯罪，也不可在使者面前說是錯許了。為何使神因你的聲音發怒，敗壞你手所做的呢？ \end{tabularx} \\ \\ \relax
5:7 & \begin{tabularx}{0.7\textwidth}{X} 多夢多言，其中多有虛空，你只要敬畏神。 \end{tabularx} \\ \\
[1ex]
\hline
\hline
\end{longtable}
上面我們看見作者所羅門王對人生的體驗和感受,並勸勉讀者如何過討神喜悅的生活.今天我們看看他對人生的觀察是怎樣的,從(三章1節至五章1節)
\newline
\newline
第一方面的觀察,傳道者指出人生既然事事有定時,我們何必勞碌呢?甚麼事也在神的手中,我們不需要勞苦,使自己得著許多.耶利米的禱告說:「耶和華阿,我曉得人的道路,不由自己,行路的人,也不能定自己的腳步.」(耶十23)但以理書二章21節也說:「他改變時候,日期,廢王,立王；將智慧賜與智慧人,將知識賜與聰明人.」有一本著名小說《苦帕》,提及法國一貧苦少女,後來成為法國著名宮廷貴婦；不幸漸漸命運改變,墮落成為一個窮困,不受注意的女人.故此你雖然爬得很高,始終也歸回原位.萬物皆有定期,所以毋須用自己的方法來經營勞碌,這是傳道書作者對人生虛空的觀察.
\newline
\newline
第二方面的觀察,(三章10節至15節),這段經文總意,指出人不能超越神,雖然唯物主義者高唱人定勝天；但並不表示人勝過神!他們把天比作神,其實天是指自然界.傳道書作者在這裡並不把自然界看作神.他說:「神造萬物」,神將永生放在人的心裡；所以人有不滅的靈魂,既然神能將永恆賜給人,祂自己就是永恆.整個宇宙都是神的創造,人是被造物,神是創造主,人怎能勝過神?我們常說:「謀事在人,成事在天.」人可以在地上盡其努力,但若神不批准,也是徒然.
\newline
\newline
在這段經文中,作者有下列幾個觀點的觀察,值得我們留意.
\newline
\newline
1.「我見神叫世人勞苦,使他們在其中受經練.」(三10)神叫人勞苦,目的叫人在其中受經練.既已經練,就會在人生的勞苦中而感到虛空；然後可以教導別人知道自己是有限的,人不能想得甚麼,便得甚麼.盼望與事實不是一樣的,所以神讓人勞苦,是要教訓人,使人在當中得著益處.若我們不是勞苦擔重擔,便不需要到神面前得安息.古語云:「勞則忘,逸則淫.」大衛犯罪時正是太陽平西方起床.
\newline
\newline
2.「又將永生安置在世人心裡,然而神從始至終的作為,人不能參透.」(三11),人雖然為萬物之靈,仍然是有限的,不能完全明白神的作為；人雖然是聰明,但在許多事情上仍是無知.例如為何好人會受苦,惡人會興盛；有時候我們實在不會明白.
\newline
\newline
3.「現今的事早先就有了,將來的事早已也有了,並且神使已過的事重新再來.」(三15),這裡是回應,(一9-11)但在這裡加添了一個新的意思；神使已過的事重新再來,日光之下無新事.可見神怎樣作,我們人不能改變,也不可以敵對.根據以上的觀察,傳道書作者指出人不能勝過神.
\newline
\newline
第三方面的觀察,(三章16至22),這段聖經指出人與野獸是沒有分別的,作者指出主持公義的人沒有公義,(三16-18)好像野獸一樣.弱肉強食,適者生存；誰好勝詭詐,便可生存.所以主持公義的,卻沒有公道；簡直是衣冠禽獸,在歷史上可以看見有許多的冤獄.今天報載,日本有一男子坐了三十四年監後,才發現他是清白的,這是何等冤枉.舍弟在法庭做事,他說:「打官司是屬於有錢人家的,窮人多是敗訴的.」在世界上最講公義的地方,可能是最無公義的地方,難怪作者說人生是空虛的.作者也在(三19-22)指出人與野獸一樣,也會死亡.他觀察看見,人與野獸差不多,人有獸性,也會死亡.
\newline
\newline
另外,在第四章作者觀察,人許多不愉快的經驗,是沒有公平的.
\newline
\newline
1.(四1-3)上文看見在審判的地方沒有公義,現在指出在人與人之間是沒有公平的.有人生經驗的就會體會到,人與人之間講求利益,卻沒有公平；人受欺壓,得不到安慰；奇怪的是那樣欺壓人的,自己也得不到安慰.故此作者指出,不如不做人更好,連傳道者也有這樣的感受.
\newline
\newline
2.(四4-6)人另外在處處都有嫉妒,當別人看見你比他好的時候,便眼紅起來.嫉妒會令人殺人,該隱殺亞伯,因為他嫉妒神悅納兄弟之祭品.掃羅王追殺大衛,嫉妒他的英勇又為百姓所愛戴.我讀小學五年級的時候,曾考取第一名,讀六年級時來了一個插班生；當我取書本閱讀時,他總是騷擾我；其實他比我聰明得多,但卻嫉妒我.可見小小年紀,也會嫉妒人.
\newline
\newline
第四方面的觀察,(四7-12),有人用這段聖經來支持結婚的好處；其實這段經文不是指結婚,乃是指出人有孤單的時日,孤單是很難忍受的.俗語說:「獨食難肥」從這話亦可以意味到孤單的滋味.一次在報章報導美國有一個富翁,家財有三十億美元之多；但這人悲嘆自己孤單寂寞,到五十五歲時,仍然獨身.有許多女人可能因他有錢,才會喜歡他,故此他不信任女人.他有許多別墅和勞斯萊斯名貴汽車,但沒有人與他共享,他向記者說:「雖然我很富有,但仍是形單影隻」,他非常孤寂憂傷.
\newline
\newline
(四13-16),這裡指出名望是短暫的,當你離開世界後,別人便會忘記你.考古學家及歷史學家,可以讓我們看見世上有許多歷史的偉人,但仍然有許多有名望的人因著他們的過去,竟然被人遺忘了.
\newline
\newline
第五方面的觀察,虛有其表的宗教生活.(五1-7)當一些宗教人士有宗教外表而無實質的時候；不單令人厭煩,連神也厭煩.有些時候,那些心口不一的人,會成為別人信主的攔阻.有一間教會的一位執事,擔任許多事務,崇拜時擔任招待；講道時卻坐在最後一排,與別人談話.講員請會眾打開聖經,他因老花眼,把聖經倒轉來看；我坐在他旁邊,不勝其擾,實在令人厭煩.所以五章1節提醒我們要小心自己的行為,先知撒母耳對違命的掃羅說:「聽命勝於獻祭,順從勝過羔羊的脂油.」虛偽的宗教生活,在這裡指出特別在言語方面顯露出來,言語最容易顯出自己的愚昧.箴言用許多篇幅,教導我們怎樣表達言語；雅各書的作者也指出,假如我們在言語上沒有過失,我們便是完全人.多言多語,容易得失,令人厭煩.
\newline
\newline
我們也要小心對神的言語.比如許願,作者指出許願不還願,不如不許,因為會令神發怒,祂會敗壞你手所作的.故此,我們若在神面前許願,要儘快償還.我的神學院有一位同學,他撞車昏迷,醒來時向神禱告許願奉獻給神.痊癒後,便進入神學院讀書；可惜他神學畢業後,卻不作傳道.我為這同學很擔心.一個許願而不還願的人,一生必不亨通.
\newline
\newline
我們看見傳道者對人生多方面的觀察,得出結論,看見人生不甚美好可愛.神藉使徒約翰教導我們,不要愛世界和世界上的事.今天盼望我們思想,究竟世界對我們有甚麼的吸引?讓傳道書的內容對我們說話.他智慧的觀察,因這世界真正不可靠,也不可以長存,這樣多的麻煩和痛苦,我們何必費盡力去爭取這個世界.耶穌說:「人若賺得全世界,卻喪失自己的生命,有甚麼益處呢?人還能拿甚麼換生命呢?」生命是最寶貴的.其實傳道者對人生不是完全悲觀,在(三12-14,22),讓我們看見一線曙光,特別(五章7節)「你只要敬畏神」,(三章14節),「我知道神一切所作的,都必永存,無所增添,無所減少；神這樣行,是要人在他面前存敬畏的心.」所以作者為我們提出一個對策,既然我們活在這個虛空的世界裡,難道過虛空的人生,只要我們敬畏神.敬畏神就是很尊重神,每件事讓神居首位,以神的意思和旨意為依歸.歸榮耀給神!在神面前謙卑,讓神教導我們如何做人；如何走人生的道路,這樣就是敬畏神.很多時候,我們會發疑問,自己知道神是好的,但怎能產生一個敬畏神的心?傳道書作者教我們一個方法.(三章11節)「神造萬物,各按其時成為美好,又將永生安置在世人心裡.」這裡指出神是創造宇宙的主宰,創造人類的神,祂統管我們的生命.當我們有這種思想時,自然可以產生敬畏神的心.當我們面對神的偉大,便會看見自己的微小.人不能參透,這裡不是指神是神秘的,無論我們如何尋求?人也是不知的,其實聖經是神的啟示,神願意我們認識祂.作者乃是指出神是偉大無比,儘管我們說認識神,也只是多少而已,已足夠我們去敬拜祂.當我們認識祂的偉大,也會在祂面前屈膝敬拜.(三章14節):「我知道神一切所作的,都必永存.」神是永存的,人卻很短暫,神的作為是完全的,人的作為卻是有瑕疵的,假如我們認識神的安排最好,祂是完全的,我們便會在祂面前屈膝敬拜.幾年前我在宣道會一份刊物中,看見南越一信徒,他的家園被北越炸毀,但他卻站在瓦爍前,低頭稱頌神.我相信他認識神是位不會做錯事的神,神既是完全的,偉大的；我們便會產生敬拜和景仰的心.今天在神的教會中,我們不單衹相信神,還要過敬畏神的生活,認識神的偉大；以致我們在生活中,過神要我們過的生活.
\newline
\newline
傳道者繼續提醒我們,除了敬畏神外,我們可以在今生享福喜樂.我們所享受的,是高級的福樂,與普通人的不同,所以基督徒應該是最快樂的人.(三章12節):「我知道世人,莫強如終身喜樂行善.」原來是行善最樂,幫助人是最快樂.另外,(三章22節):「故此,我見人莫強如在他經營的事上喜樂.」這就是「樂業」,基督徒在自己崗位上樂業,因為這是神的帶領,走在神的旨意中,一定很快樂；見證很自然會流露出來,使人看見基督的美麗.最影響馬丁路得是一個德國修士,叫徒勒爾,他說:「假如我不是修士,而是做鞋匠,我會做出最好和最平的鞋.」今日我們不論做任何工作,一定要做得最好,顯出我們高尚的見證.最後,(三章13節):「並且人人吃喝,在他一切勞碌中享福,這也是神的恩賜.」我們可以求神,使我們在吃喝中得到快樂.每次我們食東西時,都覺得很享受,這是何等快樂,這是勞碌中享福的一個實例,我們可以每次在吃喝時,有很好的胃口,很快樂的吃喝；每一天,我們便充滿喜樂.
\newline
\newline


\newpage

\section{第55屆~港九培靈會~劉承業~第四講　福祿壽仍屬虛空}
\label{sec:504}
\textbf{第四講　福祿壽仍屬虛空}
\newline
\newline
連結: \href{https://www.hkbibleconference.org/session-message/view/504}{\texttt{https://www.hkbibleconference.org/session-message/view/504}}
\newline
\newline
\hyperref[sec:503]{< < < PREV SERMON < < <}
~
\hyperref[sec:toc]{[返目錄]}
~
\hyperref[sec:505]{> > > NEXT SERMON > > >}
\newline
\newline
傳道書 5:8-6:12
\newline
\begin{longtable}{cl}
\hline
\hline
章節 & 經文 (和合本修訂版)\\
\hline
5:8 & \begin{tabularx}{0.7\textwidth}{X} 你若在一個地區看見窮人受欺壓，公義公平被掠奪，不要因此驚奇；有一位高過居高位的在鑒察，在他們之上還有更高的。 \end{tabularx} \\ \\ \relax
5:9 & \begin{tabularx}{0.7\textwidth}{X} 況且地的益處歸眾人，就是君王也受田地的供應。 \end{tabularx} \\ \\ \relax
5:10 & \begin{tabularx}{0.7\textwidth}{X} 喜愛銀子的，不因得銀子滿足；喜愛財富的，也不因得利益知足。這也是虛空。 \end{tabularx} \\ \\ \relax
5:11 & \begin{tabularx}{0.7\textwidth}{X} 貨物增添，吃的人也增添，物主得甚麼益處呢？不過眼看而已！ \end{tabularx} \\ \\ \relax
5:12 & \begin{tabularx}{0.7\textwidth}{X} 勞碌的人不拘吃多吃少，睡得香甜；富人的豐足卻不容他睡覺。 \end{tabularx} \\ \\ \relax
5:13 & \begin{tabularx}{0.7\textwidth}{X} 我見日光之下有一件令人憂傷的禍患，就是財主積存財富，反害自己。 \end{tabularx} \\ \\ \relax
5:14 & \begin{tabularx}{0.7\textwidth}{X} 他因遭遇不幸，財產盡失；他生了兒子，手裡卻一無所有。 \end{tabularx} \\ \\ \relax
5:15 & \begin{tabularx}{0.7\textwidth}{X} 他怎樣從母胎赤身而來，也必照樣赤身而去；他所勞碌得來的，手中分毫不能帶去。 \end{tabularx} \\ \\ \relax
5:16 & \begin{tabularx}{0.7\textwidth}{X} 這是一件令人憂傷的禍患。他來的時候怎樣，去的時候也必怎樣。他為風勞碌有甚麼益處呢？ \end{tabularx} \\ \\ \relax
5:17 & \begin{tabularx}{0.7\textwidth}{X} 並且他終身在黑暗中吃喝，多有煩惱、病痛和怒氣。 \end{tabularx} \\ \\ \relax
5:18 & \begin{tabularx}{0.7\textwidth}{X} 看哪，我所見為善為美的，就是人在神賜他一生的日子吃喝，享受日光之下勞碌得來的好處，因為這是他應得的報償。 \end{tabularx} \\ \\ \relax
5:19 & \begin{tabularx}{0.7\textwidth}{X} 而且，一個人蒙神賞賜財富與資產，又使他能享用，能獲取自己當有的報償，在他的勞碌中喜樂，這是神的賞賜。 \end{tabularx} \\ \\ \relax
5:20 & \begin{tabularx}{0.7\textwidth}{X} 他不多思念自己一生的日子，因為神使他的心充滿喜樂。 \end{tabularx} \\ \\ \relax
6:1 & \begin{tabularx}{0.7\textwidth}{X} 我見日光之下有一件禍患重壓在人身上， \end{tabularx} \\ \\ \relax
6:2 & \begin{tabularx}{0.7\textwidth}{X} 就是人蒙神賜他財富、資產和尊榮，以致他心裡所願的一樣都不缺，只是神使他不能享用，反被外人享用。這是虛空，也是禍患。 \end{tabularx} \\ \\ \relax
6:3 & \begin{tabularx}{0.7\textwidth}{X} 人若生一百個兒子，活許多歲數；他即使壽命很長，心裡卻不因福樂而滿足，又不得埋葬；我說，那流掉的胎比他倒好。 \end{tabularx} \\ \\ \relax
6:4 & \begin{tabularx}{0.7\textwidth}{X} 因為這胎虛虛而來，暗暗而去，名字被黑暗遮蔽， \end{tabularx} \\ \\ \relax
6:5 & \begin{tabularx}{0.7\textwidth}{X} 而且沒有見過天日，甚麼都不知道，這胎比那人倒享安息。 \end{tabularx} \\ \\ \relax
6:6 & \begin{tabularx}{0.7\textwidth}{X} 那人雖然活千年，再活千年，卻不能享福；眾人豈不都歸同一個地方去嗎？ \end{tabularx} \\ \\ \relax
6:7 & \begin{tabularx}{0.7\textwidth}{X} 人的勞碌都為口腹，心裡卻不知足。 \end{tabularx} \\ \\ \relax
6:8 & \begin{tabularx}{0.7\textwidth}{X} 智慧人比愚昧人有甚麼益處呢？困苦人在眾人面前知道如何行，有甚麼益處呢？ \end{tabularx} \\ \\ \relax
6:9 & \begin{tabularx}{0.7\textwidth}{X} 眼睛所看的比心裡妄想的倒好。這也是虛空，也是捕風。 \end{tabularx} \\ \\ \relax
6:10 & \begin{tabularx}{0.7\textwidth}{X} 先前所有的，早已起了名，人早知道人是如何的，不能與比自己強壯的相爭。 \end{tabularx} \\ \\ \relax
6:11 & \begin{tabularx}{0.7\textwidth}{X} 話語多，虛空也增多，這對人有甚麼益處呢？ \end{tabularx} \\ \\ \relax
6:12 & \begin{tabularx}{0.7\textwidth}{X} 人一生虛度的日子，如影兒經過，誰知道甚麼才是對他有益呢？誰能告訴他身後在日光之下會發生甚麼事呢？ \end{tabularx} \\ \\
[1ex]
\hline
\hline
\end{longtable}
我這次與大家研究傳道書,目的是盼望與大家學習如何查經,儘量深入研究每一節聖經；難解的經節,也去研究.許多時候,我們讀聖經不求甚解,遇見與我們有不合意的地方,便放棄了研究.這種讀經態度會妨礙我們靈命的成長,我們應該儘量明白神的說話；讓神的話語教導我們,改變我們.
\newline
\newline
中國人的觀念以為能夠有「福祿壽」全是最好的.但傳道書作者指出,福祿壽全未必是福氣,可能成為禍患.他指出真正的福氣乃是從神而來的,他不是提及天上屬靈的福氣；真正的福氣仍是在地上和屬物質的現今福氣,傳道書寶貴信息就是在此.我們不需要將來到天上,才可以享受這些福氣,即使現今在這邪惡的世代,也可以靠神恩典,享受現今的福氣.他特別強調屬物質方面的福氣,我們不要以為凡是屬物質的都是敗壞的,惡的；好像初期教會洛斯底派的人,認為物質是惡的.因此有人實行禁慾主義,甚至成為放縱主義.我們不能接受這種觀念,因為物質也是出於神的恩典,是從神而來的.物質並不是敗壞的,神會透過物質來賜福給我們,可見我們的神是很實際的.這段經文可以分開兩大段來看,首先指出福祿壽全可能成為禍患,跟著他指出真正的福氣是從神而來的.
\newline
\newline
(一)我們先看成為禍患的例證
\newline
\newline
作者從人生的觀察,指出福祿壽可以成為禍患.福的方面(五10)「貪愛銀子的,不因得銀子知足；貪愛豐富的,也不因得利益知足,這也是虛空.」我們常認為有錢纔是有福的,但作者指出人是不知足的,常見愈有的人就愈想有,永不知足；百萬富翁更想作千萬富翁,千萬富又想作億萬富翁.這些財富無法使人滿足,因為人是與生俱來,便不滿足的.幾年前,當香港股市攀至高峰時,人的慾望亦隨之升高想多得點錢；股市下跌時,許多人就失意便想自殺,這些都是有錢人,可見財富不能使人滿足,反而帶來憂愁禍患.如果我們無錢買股票,便毋須理會股市的漲瀉了.
\newline
\newline
1.(五11)「貨物增添,吃的人也增添,物主得甚麼益處呢?不過眼看而已.」這是指出人的外表雖然好看,比如有些人有車,住的是高尚洋房,養的是大狼狗,很可能是外表而已；未必可以帶來滿足的.收入高的人,支出也自然多；有時候住公屋的人,卻比其他的人更加有錢,這是不足為怪的事.
\newline
\newline
(五15)「他怎樣從母胎赤身而來,也必照樣赤身而去；他所勞碌得來的,手中分毫也不能帶去.」可見你所有的,只不過佔短暫的時候；赤身而來,也是赤身而去.聽說亞力山大帝臨終前吩咐人,在他的棺木開兩個洞口,伸出他的手,以示自己雖富有家財,卻不能帶走.(五13)「我見日光之下,有一宗大禍患,就是財主積存資財,反害自己.」有錢的人,可能不理會別人,他目中無神；可憐卻成為匪徒勒索的對象,錢財成為他的禍患.有些人是守財奴,想搾取別人的金錢,成了別人的禍患.(五14)也指出,即使你有兒女,你手中也是一無所有；可能兒女是敗家的,會耗盡你的家財.(五17)「並且他終身在黑暗中吃喝,多有煩惱,又有病患嘔氣.」可能因為他得的是不義之財,以致要在黑暗中吃喝,這是很悲慘的事.
\newline
\newline
(六1-2),這裡指出自己勞碌得來的成果,不能去享受,反而讓別人來使用,這不是真正的福氣.(六3)指出雖有百子千孫,但很可能也不得埋葬,這是何等可憐.傳道者指出上面許多人以為的福氣,很可能成為禍患.
\newline
\newline
2.現在我們來看祿的方面,祿指官位方面.(五8):「你若在一省之中見窮人受欺壓,並奪去公義公平的事,不要因此詫異；因有一位高過居高位的鑑察,在他們以上還有更高的.」若果你做一個官員卻貪污欺壓,會令人厭煩.但另一方面,諺云:「一山還有一山高.」有人的地位會高過你,所以我們要做一個愛黎民的父母官.
\newline
\newline
3.至於壽方面,(六6)「那人雖然活千年,再活千年,卻不享福；眾人豈不都歸一個地方去麼?」一個人怎會活到千歲?即使你有兩千歲,你仍會死亡,仍不可享福.有壽命可能是禍患.我的岳丈是一百零二歲,他常說留住了老命,卻看了更多傷心的事；當子孫不守道時,反而令他更加傷心!當看見兒孫離世,白頭人送黑頭人,這也是很悲慘的事.
\newline
\newline
當作者在上面指出許多的虛空後,在(六8-12),他提出一些小結.作者發出兩方面的感嘆,有甚麼長處呢?有甚麼益處呢?
\newline
\newline
第一方面這是在乎價值方面,不知道甚麼東西才有價值.作者指出一般人的價值觀以為福祿壽最好；但他認為這不是最好的.我們信主後,應該有另一套價值觀,與世界的人是不同的.假若我們的價值觀是與世人一樣的話；我們便難作個基督徒,會產生很多衝突矛盾.今日世界混亂,因為沒有正確的價值觀.有人認為試婚同居是可以的,有新道德觀的人,認為只要出於愛便行；殺人也是可以的,因為認為是劫富濟貧是可以接納的.幾年前,約紐停電,許多人趁機搶劫,當新聞界訪問那些被捕者時；這些人認為有錢人有,他們沒有,故搶劫是很公平的.有些基督徒在教會內有一套倫理標準,而在教會外卻又另有一套的標準.人生的空虛是因為不知道真正的價值觀.今天聖經成為必讀的書,聖經的原則,成為我們行事為人的標準.我們的神是永恆的,祂的標準也永遠一樣.
\newline
\newline
另外第二方面的感嘆(6-12「誰能告訴他身後,在日光之下有甚麼事呢?」今日不知明日事,今天是屬於你的,你不知將來的遭遇如何,因此便帶來徬徨.香港人不知道一九九七後怎樣,心裡便會徬徨不安,這種徬徨便成為人生焦慮的原因；有些人求籤占卦,是要知道將來的事.
\newline
\newline
(二)作者也提及積極方面,真正的福氣是從神而來的
\newline
\newline
(五18-20),在這段經文中,多次提及神,是神賜他吃喝,清楚告訴我們能夠生活,是神的恩典.這是很大的享受,許多恩典,我們經常領受,平凡了便不感恩.有些人有錢財,卻不可以享用,但有些人神賜他資財,可以吃喝享用.這便有很大分別,能夠享用的人是神的恩典；能夠在勞碌中享福,這是福氣.(五20)他常有愉快心情來享受.神與他一起快樂.有些人以為酒色財氣,一本萬利,坐勞斯萊斯或飛機,這才是福氣.傳道書作者指出我們在簡單的生活中,也可以享受福樂,並不需要餐餐珍饈百味,粗茶淡飯也可以享受.(五18)「我所見為善為美的,就是人在神賜他一生的日子吃喝.」他並沒有指出吃喝甚麼,可見福樂並不在乎物質的豐富或多寡,而在乎你使用時的心態.如果我們用一種感謝的心來飲水時,較你用一種憎恨的心情來吃魚翅,更加有福.(箴十七1),「設筵滿屋,大家相爭,不如有塊乾餅,大家相安.」這些屬物質的福氣,都是神的恩賜,我們不需要有許多的名譽或享受；只要我們滿足,雖然窮困,也可以享福.這是福樂的新詮解,正如箴言說:「耶和華所賜的福,使人滿足,並不加上憂慮.」
\newline
\newline


\newpage

\section{第55屆~港九培靈會~劉承業~第五講　智慧人的體驗}
\label{sec:505}
\textbf{第五講　智慧人的體驗}
\newline
\newline
連結: \href{https://www.hkbibleconference.org/session-message/view/505}{\texttt{https://www.hkbibleconference.org/session-message/view/505}}
\newline
\newline
\hyperref[sec:504]{< < < PREV SERMON < < <}
~
\hyperref[sec:toc]{[返目錄]}
~
\hyperref[sec:506]{> > > NEXT SERMON > > >}
\newline
\newline
傳道書 7:1-7:29
\newline
\begin{longtable}{cl}
\hline
\hline
章節 & 經文 (和合本修訂版)\\
\hline
7:1 & \begin{tabularx}{0.7\textwidth}{X} 名譽強如美好的膏油，人死去的日子勝過他出生的日子。 \end{tabularx} \\ \\ \relax
7:2 & \begin{tabularx}{0.7\textwidth}{X} 往喪家去，強如往宴樂的家，因為死是眾人的結局，活人必將這事放在心上。 \end{tabularx} \\ \\ \relax
7:3 & \begin{tabularx}{0.7\textwidth}{X} 憂愁強如喜笑，因為面帶愁容，終必使心喜樂。 \end{tabularx} \\ \\ \relax
7:4 & \begin{tabularx}{0.7\textwidth}{X} 智慧人的心在遭喪之家；愚昧人的心在快樂之家。 \end{tabularx} \\ \\ \relax
7:5 & \begin{tabularx}{0.7\textwidth}{X} 聽智慧人的責備，強如聽愚昧人歌唱； \end{tabularx} \\ \\ \relax
7:6 & \begin{tabularx}{0.7\textwidth}{X} 因為愚昧人的笑聲，好像鍋子下面燒荊棘的爆聲，這也是虛空。 \end{tabularx} \\ \\ \relax
7:7 & \begin{tabularx}{0.7\textwidth}{X} 勒索使智慧人變為愚妄，賄賂能敗壞人的心。 \end{tabularx} \\ \\ \relax
7:8 & \begin{tabularx}{0.7\textwidth}{X} 事情的終局強如它的起頭；存心忍耐的，勝過居心驕傲的。 \end{tabularx} \\ \\ \relax
7:9 & \begin{tabularx}{0.7\textwidth}{X} 你的心不要急躁惱怒，因為惱怒存在愚昧人的懷中。 \end{tabularx} \\ \\ \relax
7:10 & \begin{tabularx}{0.7\textwidth}{X} 不要說：為甚麼先前的日子強過現今的日子呢？你這樣問不是出於智慧。 \end{tabularx} \\ \\ \relax
7:11 & \begin{tabularx}{0.7\textwidth}{X} 智慧加上產業是美好的，對見天日的人都有益處。 \end{tabularx} \\ \\ \relax
7:12 & \begin{tabularx}{0.7\textwidth}{X} 因為智慧庇護人，好像金錢庇護人一樣；智慧能保全智慧者的生命，這就是知識的益處。 \end{tabularx} \\ \\ \relax
7:13 & \begin{tabularx}{0.7\textwidth}{X} 你要觀看神的作為，誰能使他所彎曲的變直呢？ \end{tabularx} \\ \\ \relax
7:14 & \begin{tabularx}{0.7\textwidth}{X} 順利時要喜樂；患難時當思考。神使這兩樣都發生，因此，人不知將會發生甚麼事。 \end{tabularx} \\ \\ \relax
7:15 & \begin{tabularx}{0.7\textwidth}{X} 在虛度的日子裡，我見過各樣的事情，義人在他的義中滅亡，惡人在他的惡中倒享長壽。 \end{tabularx} \\ \\ \relax
7:16 & \begin{tabularx}{0.7\textwidth}{X} 不要行義過分，也不要過於自逞智慧，何必自取敗亡呢？ \end{tabularx} \\ \\ \relax
7:17 & \begin{tabularx}{0.7\textwidth}{X} 不要行惡過分，也不要為人愚昧，何必未到期而死呢？ \end{tabularx} \\ \\ \relax
7:18 & \begin{tabularx}{0.7\textwidth}{X} 你持守這個，那個也不要鬆手才好。敬畏神的人，這一切都能兼得。 \end{tabularx} \\ \\ \relax
7:19 & \begin{tabularx}{0.7\textwidth}{X} 智慧使擁有智慧的人比城中十個官長更有能力。 \end{tabularx} \\ \\ \relax
7:20 & \begin{tabularx}{0.7\textwidth}{X} 其實世上沒有行善而不犯罪的義人。 \end{tabularx} \\ \\ \relax
7:21 & \begin{tabularx}{0.7\textwidth}{X} 人所說的話，你不要都放在心上，免得聽見你的僕人詛咒你。 \end{tabularx} \\ \\ \relax
7:22 & \begin{tabularx}{0.7\textwidth}{X} 因為你心裡知道，自己也曾屢次詛咒別人。 \end{tabularx} \\ \\ \relax
7:23 & \begin{tabularx}{0.7\textwidth}{X} 我曾用智慧試驗這一切事，我說：「要得智慧。」智慧卻離我遠。 \end{tabularx} \\ \\ \relax
7:24 & \begin{tabularx}{0.7\textwidth}{X} 萬事之理遙不可及，太深奧，誰能測透呢？ \end{tabularx} \\ \\ \relax
7:25 & \begin{tabularx}{0.7\textwidth}{X} 我轉念，一心要知道，要考察，要尋求智慧和萬事的來由，要知道邪惡為愚昧，愚昧為狂妄。 \end{tabularx} \\ \\ \relax
7:26 & \begin{tabularx}{0.7\textwidth}{X} 我發現有一種婦人比死還苦毒：她本身是陷阱，她的心是羅網，手是鎖鏈。凡蒙神喜愛的人必能躲開她；有罪的人卻被她纏住了。 \end{tabularx} \\ \\ \relax
7:27 & \begin{tabularx}{0.7\textwidth}{X} 傳道者說：「你看，我考察一件又一件，為要尋求萬事的來由，這是我所尋得的： \end{tabularx} \\ \\ \relax
7:28 & \begin{tabularx}{0.7\textwidth}{X} 我繼續尋找，卻未找到；一千當中，我找到一個男的，但在這一切當中，卻找不到一個女的。 \end{tabularx} \\ \\ \relax
7:29 & \begin{tabularx}{0.7\textwidth}{X} 你看，我所找到的只有一件，就是神造的人是正直的，但他們卻尋出許多詭計。」 \end{tabularx} \\ \\
[1ex]
\hline
\hline
\end{longtable}
過去幾日的講道,看見傳道書較為消極的方面,其實神要我們認識看透這個世界,使我們不去愛世界.另一方面在每一段裡面,給我們一線曙光；讓我們知道怎樣倚靠神,過有意義的人生.由第七章起,內容會較為積極.箴言讓我們知道認識神是智慧的開端,傳道書也是教導我們,智慧的人是一個敬畏神的人.智慧人因著他的智慧,曉得怎樣與神和人相交；在神面前怎樣過生活,這裡看見智慧人有四種體驗:
\newline
\newline
第一種體驗(七1-4),智慧的人有不同的價值觀.在(五18-20)中也看過,許多人以為福祿壽是好的,但可能成為禍患.作者漸漸帶領我們進入這內容中,在(七1-2)中他指出,死比活著更好；這與普通人的價值觀不同,智慧者認為死比活更好,他說:「人死的日子,勝過人生的日子；往遭喪的家去,強如往宴樂的家去.因為死是眾人的結局,活人也必將這事放在心上.」但他在第一節卻說:「名譽強如美好的膏油.」這句話有點不對稱,如果我們懂得希伯來詩歌體裁,便較易明白.這裡是比較的對偶,所以原意是美名勝香膏；死日勝生時,用名譽勝過膏油來襯托死比生強得多.在上面看見作者指出,因著死,人生是空虛的；故傳道者指出,假若你能認識死比生更好的話,你便不會怕死.死不能成為你的威脅,而感覺人生是空虛的.你要認識死只不過是一種轉變而已,保羅在(腓一21)說:「因我活著就是基督,我死了就有益處.」所以保羅視死如歸.一個人之死與生,是困苦年日之經歷,難怪嬰孩用哭泣來到這個世界.(七3-4)更指出憂愁勝過快樂,如果我們未曾憂愁過,便不會同情別人了.我認識一位姊妹,她的丈夫未信主,他們有一個兒子,非常活潑可愛,令他們非常歡欣.到這個孩子三歲時,因患病死亡了；那父親由此就想及死亡的問題,結果他也信了主,成為熱心的基督徒.雖然失去了兒子,但那父親卻得著永遠的生命.多年前劉師母患病,以致我常常埋怨!後來當我自我檢討時才發現到,我太看重事業,多過我的師母；在這場疾病中,我再次更新與內子的關係,使我學習到很多以前從未學習過的功課.如果我們吃過苦,我們才會學習到許多寶貴的功課,這豈不是憂愁勝過快樂的明證嗎?保羅說:「我們依著神的意思憂愁,就生出沒有後悔的懊悔來,以致得救.」(林後七10)主耶穌也說:「哀慟的人有福了,因為他們必得安慰.」我們看見智慧人,他的價值觀與普通人不同；即使遇見死亡,他的反應也很特別.
\newline
\newline
第二種體驗,智慧的優勝(七5-6),智慧人的責備,勝過愚昧人的稱許；忠言逆耳；很多時候是不容易接受的.傳道者更指出忠言是不好聽的,但愚昧人的聲音是討厭的；敢於責備你的,才是愛你的.箴言說當面的責備,強如背地的愛情.(箴廿七5)智慧人不怕接受別人的責備,他也盡責勸勉別人.(七7)指出智慧人不怕金錢的引誘,並且勝過賄賂的試探.多年前,有一位美國宣道會會友,到建道神學院講見證.提及自己是做遊艇生意,有一個官員因他接到製造二百五十艘登陸艇的訂單,要威脅他給佣金,否則會招致許多麻煩.這位弟兄堅持不付佣金,結果那位官員便諸多留難；使建造登陸艇工程受了影響,最後他因此事而破產.但他仍然分期去償還欠款,他仍然靠神去做,後來新畿內亞一群宣教士邀請他去主領培靈會,這位弟兄當年在宣道會中被選為會友代表,可見敬畏神的人不怕金錢的引誘和賄賂,仍能站穩的.(七8)「事情的終局,強如事情的起頭.」可見智慧人的結局是好的,我們的目標是要爭取最後的勝利!賽跑是我們的例子,當我們有耐心的時候,便能得到最後的勝利.很多人都聽過拔苗助長的故事,有一次我請教曾作傳道幾十年的岳父,他送給我四個字:「忍辱耐勞」,能夠忍受羞辱和勞苦,相信我們的事奉會蒙神賜福.
\newline
\newline
(七10)「不要說先前的日子,強過如今的日子,是甚麼緣故呢?你這樣問,不是出於智慧.」一個智慧人是很實際的人,不是常說及以前,乃要每日有新的領受,智慧人是實事求是的.(七11-12),智慧比產業更好,因為產業會失去,但智慧卻永遠相隨.
\newline
\newline
第三方面的體驗(七13-15),在(六12)傳道者指出人生很無意思,前途未卜,人心徬徨；但智慧人卻知道自己的前途在神的手中,因為他不會為明天自誇和驕傲,(七14)指出人生有亨通的日子,也有遭難的日子；我既不知明天的事,便不能自誇.我們既知道自己的前途在神手中,就應將自己的學業,婚姻,家庭等事放在主的手,交託神,我們便不會自誇,便不須要憂慮了.神的作為是完全永存的,這樣,我們做人便減少很多的憂慮.
\newline
\newline
最後一個體驗(七16-29),這裡指出世上無完全人.這段經文是難解釋的,有解經家指出我們應該守中庸之道,不要行義或行惡過份,但若以世上無完全人為主題,更加合理.甚麼叫做行義過份,這詞有幾方面的含意:第一,不要顯得比別人好,另外含意是不要有屬靈的驕傲；不要自逞智慧,有人解釋這句話的意思,是不要當自己是先知.作者說不要行惡過分,他承認智慧人也會犯罪,但不會為人愚昧.(七20)「時常行善而不犯罪的義人,世上實在沒有.」各位不要誤會以為自己可以行惡,只要不過份便好了.這話用的語氣,是消極方面的逃避,而不是積極方面的去行惡.他的原意是要我們謙虛一點,不要自誇比別人強；也不可認為自己是完全的,有些時候也可能偶被過犯所勝,我們只不過是一個「蒙恩的罪人而己.」我們與普通人一樣,幸得蒙神的憐憫與赦罪之恩.當我們有這種心態的時候,便不會在人面前自誇,智慧人覺得自己是無有的.最後傳道者提醒我們,(七27)「但眾女子中,沒有找到一個」,可能所羅門有許多後宮佳麗,他從自己體驗中,指出智慧人是謙卑的人；承認自己是不完全,由這裡可以看見基督教一個重要的真理；指出人是敗壞的,沒有一個義人,所以我們需要基督的拯救,神的幫助.這段經文指出不單在道德上,在知識上也沒有一個完全的人,(七23-25,28),智慧人認識自己不單在道德上未完全,也在知識上有限；在做人方面,仍有許多要學習的地方；與人合作方面,需要時常竭力追求有改進.最後傳道者在(七28)說:「看哪,一千男子中,我找到一個正直人」,原文是「看哪,一千男子中,我找到一個.」在整本聖經中,這樣的語句只有在約伯記(卅三23)中出現「一千天使中,若有一個作傳話的,與神同在,指示人所當行的事.」究竟千中的一位是誰呢?傳道者在他追求智慧,追求與神與人的關係上,他得到一個發現.他不是高抬男子,低貶女人；這位千中的一位,從聖經啟示中便知道他是耶穌基督.只有耶穌基督是完全的,耶穌是完全的神；也是完全的人,祂是我們所仰慕的.
\newline
\newline
在這一章聖經中,傳道者讓我們看見,一個智慧人,他的價值觀與普通人不同的；他看見智慧在我們人生中可以得勝的,可以面對現實和各方面的困難.他認識我們的前途在神的手中,所以他不為明日自誇或憂慮.最後他認識自己在道德和知識方面不是完全的人,他需要在對神,對人方面繼續追求,求神也讓我們有智慧者的看法和眼光.
\newline
\newline


\newpage

\section{第55屆~港九培靈會~劉承業~第六講　智慧人接納神的主權}
\label{sec:506}
\textbf{第六講　智慧人接納神的主權}
\newline
\newline
連結: \href{https://www.hkbibleconference.org/session-message/view/506}{\texttt{https://www.hkbibleconference.org/session-message/view/506}}
\newline
\newline
\hyperref[sec:505]{< < < PREV SERMON < < <}
~
\hyperref[sec:toc]{[返目錄]}
~
\hyperref[sec:507]{> > > NEXT SERMON > > >}
\newline
\newline
傳道書 8:1-9:12
\newline
\begin{longtable}{cl}
\hline
\hline
章節 & 經文 (和合本修訂版)\\
\hline
8:1 & \begin{tabularx}{0.7\textwidth}{X} 誰如智慧人呢？誰知道事情的解釋呢？人的智慧使他的臉發光，改變他臉上暴戾之氣。 \end{tabularx} \\ \\ \relax
8:2 & \begin{tabularx}{0.7\textwidth}{X} 我勸你因神誓言的緣故，當遵守王的命令。 \end{tabularx} \\ \\ \relax
8:3 & \begin{tabularx}{0.7\textwidth}{X} 不要急躁離開王的面前，不要固執行惡，因為他凡事都隨自己心意而行。 \end{tabularx} \\ \\ \relax
8:4 & \begin{tabularx}{0.7\textwidth}{X} 王的話本有權力，誰能對他說：「你在做甚麼？」 \end{tabularx} \\ \\ \relax
8:5 & \begin{tabularx}{0.7\textwidth}{X} 凡遵守命令的，必不經歷禍患；智慧人的心知道適當的時機和必經的過程。 \end{tabularx} \\ \\ \relax
8:6 & \begin{tabularx}{0.7\textwidth}{X} 各樣事務都有時機和過程，但人有苦難重壓在身。 \end{tabularx} \\ \\ \relax
8:7 & \begin{tabularx}{0.7\textwidth}{X} 他不知道將來的事，其實將來如何，誰能告訴他呢？ \end{tabularx} \\ \\ \relax
8:8 & \begin{tabularx}{0.7\textwidth}{X} 沒有人能掌握生命，將生命留住；也沒有人有權力掌管死期。這場爭戰無人能免；邪惡也不能救那行邪惡的人。 \end{tabularx} \\ \\ \relax
8:9 & \begin{tabularx}{0.7\textwidth}{X} 這一切我都見過，我專心考察日光之下所發生的一切事，有時這人管轄那人，令他受害。 \end{tabularx} \\ \\ \relax
8:10 & \begin{tabularx}{0.7\textwidth}{X} 我見惡人被埋葬；從前他們進出聖地，他們在城中的作為被人忘記。這也是虛空。 \end{tabularx} \\ \\ \relax
8:11 & \begin{tabularx}{0.7\textwidth}{X} 判罪之後不立刻執行，所以世人滿懷作惡的心思。 \end{tabularx} \\ \\ \relax
8:12 & \begin{tabularx}{0.7\textwidth}{X} 罪人雖然作惡百次，倒享長壽；然而我也知道，福樂必臨到敬畏神的人，就是在他面前心存敬畏的人。 \end{tabularx} \\ \\ \relax
8:13 & \begin{tabularx}{0.7\textwidth}{X} 惡人卻不得福樂，他的日子好像影兒不得長久，因為他不敬畏神。 \end{tabularx} \\ \\ \relax
8:14 & \begin{tabularx}{0.7\textwidth}{X} 世上有一件虛空的事，就是義人所遭遇的，反而照惡人所做的；惡人所遭遇的，反而照義人所做的。我說，這也是虛空。 \end{tabularx} \\ \\ \relax
8:15 & \begin{tabularx}{0.7\textwidth}{X} 我就稱讚快樂，原來人在日光之下，最大的福氣莫過於吃喝快樂；他在日光之下，神賜他一生的日子，要從勞碌中享受所得。 \end{tabularx} \\ \\ \relax
8:16 & \begin{tabularx}{0.7\textwidth}{X} 我專心想要明白智慧，要觀看世上所發生的事。有人晝夜不得闔眼睡覺。 \end{tabularx} \\ \\ \relax
8:17 & \begin{tabularx}{0.7\textwidth}{X} 我觀看神一切的作為，知道人不能探求日光之下所發生的事；任憑他費多少力探索，都找不出來，智慧人雖說他明白，仍不能找出來。 \end{tabularx} \\ \\ \relax
9:1 & \begin{tabularx}{0.7\textwidth}{X} 我將這一切事放在心上，詳細研究這些，就知道義人和智慧人，並他們的作為都在神手中；或是愛，或是恨，都在他們面前，但人不能知道。 \end{tabularx} \\ \\ \relax
9:2 & \begin{tabularx}{0.7\textwidth}{X} 凡臨到眾人的際遇都一樣：義人和惡人，好人，潔淨的人和不潔淨的人，獻祭的和不獻祭的，都一樣。好人如何，罪人也如何；起誓的如何，怕起誓的也如何。 \end{tabularx} \\ \\ \relax
9:3 & \begin{tabularx}{0.7\textwidth}{X} 在日光之下發生的一切事中有一件禍患，就是眾人的際遇都一樣，並且世人的心充滿了惡；活著的時候心裡狂妄，後來就歸死人那裡去了。 \end{tabularx} \\ \\ \relax
9:4 & \begin{tabularx}{0.7\textwidth}{X} 與一切活人相連的，那人還有指望，因為活著的狗勝過死了的獅子。 \end{tabularx} \\ \\ \relax
9:5 & \begin{tabularx}{0.7\textwidth}{X} 活著的人知道必死；死了的人毫無所知，也不再得賞賜，因為他們的名已被遺忘。 \end{tabularx} \\ \\ \relax
9:6 & \begin{tabularx}{0.7\textwidth}{X} 他們的愛，他們的恨，他們的嫉妒，早就消滅了。在日光之下所發生的一切事，他們永不再有份了。 \end{tabularx} \\ \\ \relax
9:7 & \begin{tabularx}{0.7\textwidth}{X} 你只管歡歡喜喜吃你的飯，心中快樂喝你的酒，因為神已經悅納你的作為。 \end{tabularx} \\ \\ \relax
9:8 & \begin{tabularx}{0.7\textwidth}{X} 你的衣服要時時潔白，你頭上也不要缺少膏油。 \end{tabularx} \\ \\ \relax
9:9 & \begin{tabularx}{0.7\textwidth}{X} 在你一生虛空的日子，就是神賜你在日光之下虛空的日子，當與你所愛的妻快活度日，因為那是你一生中在日光之下勞碌所得的報償。 \end{tabularx} \\ \\ \relax
9:10 & \begin{tabularx}{0.7\textwidth}{X} 凡你手所當做的事，要盡力去做；因為在你所必須去的陰間沒有工作，沒有謀算，沒有知識，也沒有智慧。 \end{tabularx} \\ \\ \relax
9:11 & \begin{tabularx}{0.7\textwidth}{X} 我轉而回顧日光之下，快跑的未必能贏，強壯的未必戰勝，智慧的未必得糧食，聰明的未必得財富，有學問的未必得人喜悅，全在乎各人遇上的時候和機會。 \end{tabularx} \\ \\ \relax
9:12 & \begin{tabularx}{0.7\textwidth}{X} 人不知道自己的定期。魚被險惡的網圈住，鳥被羅網捉住，禍患的時刻忽然臨到，世人陷在其中也是如此。 \end{tabularx} \\ \\
[1ex]
\hline
\hline
\end{longtable}
在第七章我們看見智慧人在生活上各方面無往而不利,作者在(八1)說:「誰如智慧人呢?誰知道事情的解釋呢?人的智慧使他的臉發光,並使他臉上的暴氣改變.」作者在這裡好像作一個總結,而又作一個開始,有智慧的人雖神采飛揚；然而智慧人並不是萬物通曉,所以他說誰知道事情的解釋呢?有時臨到我們身上遭遇,是不可以理解的.在(八2-5)所羅門提出王權的重要,我個人認為所羅門不是強調政治的王權；而是藉這王權引伸出神的權柄.既然王的權柄我們要尊重,更何況是神的權柄呢?(八3)指出君王有無上權柄,操縱國家和子民；王的說話就是法律,因此神自己更加有權柄,因為神是萬王之王,萬主之主.所羅門認識自己能夠作王,都是神所賜予.就如路加福音第七章,有迦百農的百夫長；他從自己所擁有的權力,來聯想到神的兒子耶穌基督有更大的權柄.他相信耶穌只要講一句話,他的僕人就會好了.今天我們若認識神的恩典,便會加強對神的信心.(八5)指出今日我們若果想免經歷禍患,就應該要遵守神的命令.亞當夏娃因為不遵守神的話,所以罪便進入了世界；為甚麼我們做基督徒的沒有見證?正因我們沒有遵守神的吩咐和命令.另外,傳道者在(八6-7)指出人不能控制自己的生命或際遇；他何時離開這個世界,他不能控制,秦始皇想求長生不老藥,結果他失敗了!世界上只有以利亞和以諾未死而被提,然而他們也不能控制何時被提.在(九11-12)也指出今日不知明日事,雅各書也提醒我們；明日會遭遇甚麼事,我們不知道,只有說主若願意,方可作這事或作那事.
\newline
\newline
(一)傳道者為我們提出一個原因,為何我們不可知道呢?
\newline
\newline
許多時候世界沒有公理.(八9-7)這一切我都見過,也專心查考日光之下所作的一切事,有時這人管轄那人,令人受害.」在(八10-14),他看見在世間許多時候,好人受苦,惡人興旺；惡人長壽,好人卻英年早逝,令人惋惜!長輩愛主,後輩卻不愛主；或後輩很追求,但我們父母的行為卻失榜樣.我們很難過,求神多年,也無改變.有時候種瓜不得瓜,種豆不得豆；無論如何,他提醒我們,我們的生命及際遇,是不由我們控制的,在這一段經文中,他讓我們看見,是神控制人的壽命和一切際遇.(八15)指出我們的年日在神的手中,我們應該謙卑下來,好好把握使用自己的生命.既然我們不能控制自己的際遇,便應該將自己交託給神；我們應該要敬拜神,敬畏祂,聽祂的話；因為神掌管我們的氣息和一切.談及主權的問題,香港人知得很明白,香港前途的談判,最難的地方在主權問題,這問題解決了；其它都會迎刃而解,主權問題最影響香港人的前途.在你生命中,你有沒有弄清楚你主權的問題,究竟神掌管你生命的主權,還是你自己?若果你生命的主權未清楚時,你永遠也得不到安息.我們跌倒犯罪,許多時候是因為未將生命和生活的主權交給主耶穌基督,若果你弄清楚這個問題,你其它的的問題便很容易解決.許多基督徒,因為已將生命的主權交付了耶穌基督；所以其它小問題,便有能力去解決了.我們有沒有接納神的主權呢?
\newline
\newline
(二)智慧者指出惡人的興盛是暫時的
\newline
\newline
(八1-12)斷定罪名是神判斷惡人,罪人,但神並沒有叫他立刻倒斃,潦倒；因此緣故,世人便滿心作惡.在這裡引起了我們的疑問,為何神不馬上刑罰?例如人說謊時,便他立刻腹痛,假若是這樣的話,人便不需要有自由意志了.神賦予人有自由意志,以致他愛神是心甘情願的,因他有選擇權；因此神便不立刻施刑,但最後必有審判和算賬.(八12-13),惡人得不著福樂,也不得長久的日子；惡人的興盛是屬於暫時的,例外的.我們毋須為這例外,而影響我們的忠心和愛心.
\newline
\newline
(三)智慧者知道敬畏神的人,終久必得福樂
\newline
\newline
(八12-15)(九7-9)給敬畏祂的智慧人的福樂是現實的,也是物質的.作者提出吃喝可以快樂,是因為神悅納你的作為,你與神有正常的關係.因此你便會開心吃喝；假如你心中有罪,食東西時便會沒有胃口,在這種和諧的關係中,吃喝才是真正的享受；假如我們與神有好的關係時,便能自然地吃喝快樂.另外(九8)神讓我們能夠有清潔的衣服.我們敬畏神的人,也需要注意自己的裝扮,當我們到教會聚會時,應該注意自己的服飾；不要隨隨便便,這是敬畏神的表現.(九9)指出神給我們的福樂,是包括有美滿的婚姻生活,雖然有家庭,責任便多了；但這是神在我們勞碌的事情上所給我們的分,神讓我們有家庭,可以彼此分擔困難.所以當我們的婚婚不美滿時,要思想是否我們沒敬畏神,未按神的心意來選擇配偶.若不是神所配搭的,遲早必有問題；又或者夫婦之間在愛情根基上的不穩,在基督徒夫婦間有時就會有這種情況發生.彼此不相關,沒有愛情的基礎,所以便不能真真正的快活度日.今天我們作丈夫的,可能事業心太重,未能負起作丈夫的責任；結了婚的姊妹也要以家庭為重,神的心意.是要我們與丈夫或妻子快活度日；但許多時候,我們的心意霸佔神的主權,自己做喜歡的事,便不能享受所賜的福樂.今天破裂的婚姻生活,在基督徒的家庭也可以看到,我們應該要正視這個問題.
\newline
\newline
智慧人如何充實人生的虛空.
\newline
\newline
1.智慧人要接納神的主權.特別在壽命和際遇上,若能接納神的主權我們便不會怨天尤人,靜默下來低頭敬拜了!俗語說:「塞翁失馬,焉知非福.」美國有一愛主的青年人叫卜理納,他是著名神學家愛德華滋的女婿.有一本書《拓荒英雄》提及他的生平,這青年蒙神引導,到美國印第安納工作；他拚命為神工作,二十九歲便離世.我們不明白為何神叫他這樣早便死去,但神有祂的時間；在卜理納所寫的日記中,寫得非常真摯動人,有許多人因閱讀他的日記而受感動奉獻給主.芬妮歌羅士比,是位失明的作曲家,她作了許多聖詩；雖然她自少失明,但神仍使用她；許多在靈裡失明的人,眼睛重開.我們最喜歡唱的一首聖詩,例如「在恩座前求」.這首歌作者的未婚妻,在寫詩前一天,是遇溺喪生的.作者在傷心難過時,寫下這首詩歌,幫助無數人.所以我們要接納神的主權,有時候雖然難過傷心,但要相信,這是神在管理一切.
\newline
\newline
(四)智慧人看清楚惡人的興盛是暫時的,他明白義人和惡人的結局(八12-13,17)
\newline
\newline
他承認自己只知道一些事情,但還有些事情,他是不知道的.也有許多事情,是我們不能明白的；例如苦難的問題,如果我們肯接納自己是不能全知道的時候,心境便會較平安了.
\newline
\newline
最後,(九10)我們要盡力去做事情,基督徒要做一個敬業樂業的人,無論大小事上總要盡力.在工作和事奉崗位上總要榮耀神.作榮神益人的人,既榮耀神,又可使自己得益.正如一個老闆喜歡一個盡責的僱員,會給他加薪升職,得到從人而來的好處.當我們專心在工作上時,便不會受其它問題困擾了.
\newline
\newline
傳道者指出我們要接納神的主權,也承認有許多事情,是不容易解釋的,直到我們見主面的時候,便會知清了.最後我們應在所做的工作上,盡心竭力!
\newline
\newline


\newpage

\section{第55屆~港九培靈會~劉承業~第七講　智慧面面觀}
\label{sec:507}
\textbf{第七講　智慧面面觀}
\newline
\newline
連結: \href{https://www.hkbibleconference.org/session-message/view/507}{\texttt{https://www.hkbibleconference.org/session-message/view/507}}
\newline
\newline
\hyperref[sec:506]{< < < PREV SERMON < < <}
~
\hyperref[sec:toc]{[返目錄]}
~
\hyperref[sec:508]{> > > NEXT SERMON > > >}
\newline
\newline
傳道書 9:3-10:20
\newline
\begin{longtable}{cl}
\hline
\hline
章節 & 經文 (和合本修訂版)\\
\hline
9:3 & \begin{tabularx}{0.7\textwidth}{X} 在日光之下發生的一切事中有一件禍患，就是眾人的際遇都一樣，並且世人的心充滿了惡；活著的時候心裡狂妄，後來就歸死人那裡去了。 \end{tabularx} \\ \\ \relax
9:4 & \begin{tabularx}{0.7\textwidth}{X} 與一切活人相連的，那人還有指望，因為活著的狗勝過死了的獅子。 \end{tabularx} \\ \\ \relax
9:5 & \begin{tabularx}{0.7\textwidth}{X} 活著的人知道必死；死了的人毫無所知，也不再得賞賜，因為他們的名已被遺忘。 \end{tabularx} \\ \\ \relax
9:6 & \begin{tabularx}{0.7\textwidth}{X} 他們的愛，他們的恨，他們的嫉妒，早就消滅了。在日光之下所發生的一切事，他們永不再有份了。 \end{tabularx} \\ \\ \relax
9:7 & \begin{tabularx}{0.7\textwidth}{X} 你只管歡歡喜喜吃你的飯，心中快樂喝你的酒，因為神已經悅納你的作為。 \end{tabularx} \\ \\ \relax
9:8 & \begin{tabularx}{0.7\textwidth}{X} 你的衣服要時時潔白，你頭上也不要缺少膏油。 \end{tabularx} \\ \\ \relax
9:9 & \begin{tabularx}{0.7\textwidth}{X} 在你一生虛空的日子，就是神賜你在日光之下虛空的日子，當與你所愛的妻快活度日，因為那是你一生中在日光之下勞碌所得的報償。 \end{tabularx} \\ \\ \relax
9:10 & \begin{tabularx}{0.7\textwidth}{X} 凡你手所當做的事，要盡力去做；因為在你所必須去的陰間沒有工作，沒有謀算，沒有知識，也沒有智慧。 \end{tabularx} \\ \\ \relax
9:11 & \begin{tabularx}{0.7\textwidth}{X} 我轉而回顧日光之下，快跑的未必能贏，強壯的未必戰勝，智慧的未必得糧食，聰明的未必得財富，有學問的未必得人喜悅，全在乎各人遇上的時候和機會。 \end{tabularx} \\ \\ \relax
9:12 & \begin{tabularx}{0.7\textwidth}{X} 人不知道自己的定期。魚被險惡的網圈住，鳥被羅網捉住，禍患的時刻忽然臨到，世人陷在其中也是如此。 \end{tabularx} \\ \\ \relax
9:13 & \begin{tabularx}{0.7\textwidth}{X} 我見日光之下有一樣智慧，在我看來是偉大的， \end{tabularx} \\ \\ \relax
9:14 & \begin{tabularx}{0.7\textwidth}{X} 就是有一人口稀少的小城，遇大君王前來攻擊，修築營壘，將城圍困。 \end{tabularx} \\ \\ \relax
9:15 & \begin{tabularx}{0.7\textwidth}{X} 城中有一個貧窮的智慧人，他用智慧救了那城，卻沒有人記念那窮人。 \end{tabularx} \\ \\ \relax
9:16 & \begin{tabularx}{0.7\textwidth}{X} 我就說，智慧勝過勇力；然而那貧窮人的智慧被人藐視，他的話也無人聽從。 \end{tabularx} \\ \\ \relax
9:17 & \begin{tabularx}{0.7\textwidth}{X} 寧可聽智慧人安靜的話語，不聽掌權者在愚昧人中的喊聲。 \end{tabularx} \\ \\ \relax
9:18 & \begin{tabularx}{0.7\textwidth}{X} 智慧勝過打仗的兵器；但一個罪人能敗壞許多善事。 \end{tabularx} \\ \\ \relax
10:1 & \begin{tabularx}{0.7\textwidth}{X} 死蒼蠅使做香的膏油散發臭氣；同樣，一點愚昧也能壓倒智慧和尊榮。 \end{tabularx} \\ \\ \relax
10:2 & \begin{tabularx}{0.7\textwidth}{X} 智慧人的心居右；愚昧人的心居左。 \end{tabularx} \\ \\ \relax
10:3 & \begin{tabularx}{0.7\textwidth}{X} 愚昧人的行徑顯出無知，對眾人說，他是愚昧人。 \end{tabularx} \\ \\ \relax
10:4 & \begin{tabularx}{0.7\textwidth}{X} 掌權者的怒氣若向你發作，不要離開你的本位，因為鎮定能平息大過。 \end{tabularx} \\ \\ \relax
10:5 & \begin{tabularx}{0.7\textwidth}{X} 我見日光之下有一件禍患，似乎出於統治者的錯誤， \end{tabularx} \\ \\ \relax
10:6 & \begin{tabularx}{0.7\textwidth}{X} 就是愚昧人立在高位；有錢人卻坐在低位。 \end{tabularx} \\ \\ \relax
10:7 & \begin{tabularx}{0.7\textwidth}{X} 我見僕人騎馬，王子像僕人在地上步行。 \end{tabularx} \\ \\ \relax
10:8 & \begin{tabularx}{0.7\textwidth}{X} 挖陷坑的，自己必陷在其中；拆城牆的，自己必被蛇咬。 \end{tabularx} \\ \\ \relax
10:9 & \begin{tabularx}{0.7\textwidth}{X} 開鑿石頭的，會受損傷；劈開木頭的，必遭危險。 \end{tabularx} \\ \\ \relax
10:10 & \begin{tabularx}{0.7\textwidth}{X} 鐵器鈍了，若不將刃磨快，就必多費力氣；但智慧的益處在於使人成功。 \end{tabularx} \\ \\ \relax
10:11 & \begin{tabularx}{0.7\textwidth}{X} 尚未行法術，蛇若咬人，行法術的人就得不到甚麼好處了。 \end{tabularx} \\ \\ \relax
10:12 & \begin{tabularx}{0.7\textwidth}{X} 智慧人的口說出恩言；愚昧人的嘴吞滅自己， \end{tabularx} \\ \\ \relax
10:13 & \begin{tabularx}{0.7\textwidth}{X} 他口中的話語起頭是愚昧，終局是邪惡的狂妄。 \end{tabularx} \\ \\ \relax
10:14 & \begin{tabularx}{0.7\textwidth}{X} 愚昧人多有話語。人不知將來會發生甚麼事，他身後的事誰能告訴他呢？ \end{tabularx} \\ \\ \relax
10:15 & \begin{tabularx}{0.7\textwidth}{X} 愚昧人的勞碌使自己困乏，連進城的路他也不知道。 \end{tabularx} \\ \\ \relax
10:16 & \begin{tabularx}{0.7\textwidth}{X} 邦國啊，你的君王若年少，你的群臣早晨宴樂，你就有禍了！ \end{tabularx} \\ \\ \relax
10:17 & \begin{tabularx}{0.7\textwidth}{X} 邦國啊，你的君王若是貴族之子，你的群臣按時吃喝，是為強身，不為酒醉，你就有福了！ \end{tabularx} \\ \\ \relax
10:18 & \begin{tabularx}{0.7\textwidth}{X} 因人懶惰，房頂塌下；因人手懶，房屋滴漏。 \end{tabularx} \\ \\ \relax
10:19 & \begin{tabularx}{0.7\textwidth}{X} 擺設宴席是為歡樂。酒能使人快活，錢能叫萬事應心。 \end{tabularx} \\ \\ \relax
10:20 & \begin{tabularx}{0.7\textwidth}{X} 不可詛咒君王，連起意也不可，在臥室裡也不可詛咒富人；因為空中的飛鳥必傳揚這聲音，有翅膀的必述說這事。 \end{tabularx} \\ \\
[1ex]
\hline
\hline
\end{longtable}
這段經文是十分難解的,它好像箴言一樣,喜愛使用比較；比較智慧和愚昧,組織較鬆散,故此不容易看懂.作者用比較的方法,來襯托出智慧的表現.在前幾講我曾經提過,智慧使人過成功的生活,曉得在神面前如何做人.那些未信主的人,神在他生命中是沒有地位的；但信主的人知道神鑑察我們的生活,智慧是我們過成功生活的藝術.智慧可能與知識相通,但有些時候,智慧又與知識不同；知識使我們明理,但智慧卻使我們與神與人有良好的關係.因此有知識的人不一定有智慧,智慧人也不一定有很高學位,但智慧總比知識更強.這段經文可分兩大點；第一,智慧使人過成功的生活(九13至十11),第二,怎樣表達智慧(十12-20).
\newline
\newline
(一)智慧使人過成功的生活
\newline
\newline
(九13至十11)以前我們看過智慧使我們與神,與人有良好的關係,但在這裡卻有更廣闊的意思.智慧引導我們其它有關的事,特別是政治的事.因為人的生活不單與人有關係,還有與國家,政治的關係.今日人以為有權柄便可以統管一切,所謂「槍桿子下出政權.」(九13-18),雖然智慧被人藐視,但智慧卻可以為我們做許多事.這裡提及有一城被圍困,智慧就可以拯救那城,使城裡的人得到安全.智慧勝過打仗的兵器和武力,有勇無謀是無用的.士師記提及基甸的兒子亞比米勒殺死基甸的七十個兒子後,想放火燒提備斯,卻被一女人用磨石打死；這個聰明又勇敢的婦人,就救了那城許多人的性命(士九53).另一個例子,希西家作王,亞述大軍十八萬五千人圍困大城；結果他求告,神的使者便殺死亞述大軍,解救了耶路撒冷城.在第二次世界大戰的時候,英倫三島險被德軍攻陷,但當全國人起來禱告,神令天氣急劇轉變,挽救了英國.中國有一句成語說:「原璧歸趙.」,提及戰國時候,趙國得到寶石和氏璧,舉世罕有；秦王知道後,要用十五城邑換取,趙國便派藺相如到秦國交涉.到後發現秦王食言,想併取和氏璧,藺相如便用急智得回璧玉.我們應該為執掌國家政治的領袖禱告,為地方政府高級官員禱告；使他們有敬畏神的心,有智慧來治理國事,使我們可以敬虔度日.(提前二2)假如我們在政府中擔任高位,就更當求神讓我們有智慧去參與政治的事.傳道者提醒我們,智慧使我們在政治上是有益處的,有關香港前途談判,我們應切切為雙方政府的官員及港督禱告；求神使他們明白香港人的心願,使香港能夠得到合理的地位,因為禱告能夠改變許多事情.
\newline
\newline
作者在(十1-3)舉出智慧與愚昧的分別後,他在(十4-7)提出智慧能夠幫助我們,去應付那些不好的掌權者或上司.我們可能有不好的上司或掌權者,發怨言不是智慧人應有的表現；神給我們有智慧去應付他們.掌權者是人,他們也有錯誤的地方,是否我們便離職?(十4)說:我們不要離開自己的本位.傳道者教導我們用柔和去對待上司或掌權者,可以得回他們的心.智慧的人是一個溫柔的人,回答柔和使怒消退,結出和平的果子.新約的智慧書是雅各書,在(雅13-18)說:智慧人是溫和的,使人和平的,智慧的人不單可以挽回上司的心,也使他的上司知道自己的錯誤；在那裡建立和平,使工作做得更好.有些時候,做下屬不易明白做上司的難處；在教會或基督教機構,總有領袖來帶領,若果我們不順服帶領,便有問題產生.
\newline
\newline
跟著,傳道者給我們看見,智慧使我們個人有正確的引導.(十8-9),這段經文意思是「害人反害己.」第九節看見挪移石頭,以前農地是用石頭作界限,所以這裡所指是害人之意,與上文意思一樣.在聖經中以斯帖記記載害人終害己的好例子,亞瑪力人哈曼要所有猶大人跪拜他；唯獨末底改不跪拜他,他懷恨在心,便造了一木十架吊死末底改,結果自己反被掛在木架上.其實,這是普通的常識和常理,傳道者用這提醒讀者,不要有害人的心.主耶穌也說,你們用甚麼量器量給人,也用甚麼量器量給你們.(十10)指出智慧使人事半功倍,效果大得多,你懂得將刀磨利,做事便快得多.智慧使我們與人有好關係,做事便會有效率得多.(十11)這裡指出智慧是懂得把握合宜時機,(箴廿七14)說:「清晨起來,大聲給朋友祝福的,就算是咒詛他.」,清晨大聲,許多人尚在睡眠,會令人很討厭,他說話實在不合時機.有一位姊妹家人未信主,鄰居是基督徒,但彼此沒有來往.這位姊妹的祖母去世時,鄰居便藉此來慰問；以後他們便成為親密鄰居,帶領這個家庭全家信主.這個基督徒家庭能夠把握機會,帶領人信主.
\newline
\newline
(二)如何表達我們的智慧
\newline
\newline
(十12-20),這裡主要是論及言語,怎樣在言語表達自己的智慧昵?1.少講說話,(十14)「愚昧人多有言語.」箴言書也說多言多語,難免有過；禁止嘴唇.便有智慧,」(箴十19)我在報章副刊上看見一段文章,有記者問愛因斯坦:「你成功的方程式是甚麼?」他回答說:「我成功的方程式是X=A+B+C」,X是成功,A是工作,B是休息,C是少說話.當你有工作,又有休息,而且少講話；你便走在成功的路上,這方程式值得我們回味.(十12)指出智慧人的口說出恩言,我們要說有造就人的話,要講有智慧的話,這會令人得幫助.約在二十年前,我在美國一間神學院讀書,去參加一個講座,旁邊的人問我從何處來,我便回答在某某神學院.另外的人就說:「啊,這神學院是最靠近地獄的.」原來他們是兩位新派神學院學生,這間神學院不信有天堂地獄的.在這種情況下,我真不知如何回答,但感謝神賜我智慧回答:「怎麼你們也相信有地獄嗎?」這問題令他們不知所措,可見智慧的說話可解決難處.(十20)我們要在言語上有智慧,應該做一個不咒詛別人的人,也不說人閒話.傳道者說,即使你在臥房裡講閒話,別人也會知道；故此我們對別人的批評,講閒話要謹慎小心,俗語說:「隔牆有耳」,既然如此,我們就要格外小心.2.我們如何表達自己的智慧.(十16-17),作者在這裡用兩個國家作比較,沒有人在早晨飲宴,因為宴樂通常是設在晚上的.換句話,一個智慧人要過一個有節制的生活；生活有節奏,做事有原因.英文有一個字醉酒(alcoholic),另外形容一個人有工作狂.作(Workaholic),可能我們不會做一個醉酒的人,但會造一個工作的奴隸,以致影響我們的家庭生活.有些人用工作來逃避自己的家庭,故此埋頭苦幹.摩西的律法要求土地也要有安息,對土地是好的,神讓我們六日工作,第七日休息；對我們健康有幫助,對家庭有幫助.工作的原動力,是來自神的帶領和聖靈的感動,而不是來自工作的需要.今日各階層有許多的工作機會,我們根本無可能做完所有的工作.
\newline
\newline
最後我們看(十18-19),作者在這裡用一相反論調來教導,那些好飲食的,不是智慧人,乃是愚昧人.智慧人是在勤儉上顯出來,在節儉上,使用金錢上,我們要有智慧.許多時候我們使用金錢是為面子緣故,我們應該從節儉上學習.不要為名牌貨而追求,我們若隨意揮霍,怎能見證神的榮美.我們要從享受中得到自由.近年來快餐店幫助了我們家庭不少,參加聚會後,全家在外面吃飯；到茶樓又貴,又要等位；但吃快餐只費數元,又不用花時間.智慧人在生活上和屬靈上都要有智慧,求主幫助我們,在各方面皆顯出自己是一個有智慧的人.
\newline
\newline


\newpage

\section{第55屆~港九培靈會~劉承業~第八講　當未雨綢繆}
\label{sec:508}
\textbf{第八講　當未雨綢繆}
\newline
\newline
連結: \href{https://www.hkbibleconference.org/session-message/view/508}{\texttt{https://www.hkbibleconference.org/session-message/view/508}}
\newline
\newline
\hyperref[sec:507]{< < < PREV SERMON < < <}
~
\hyperref[sec:toc]{[返目錄]}
~
\hyperref[sec:509]{> > > NEXT SERMON > > >}
\newline
\newline
傳道書 11:1-12:7
\newline
\begin{longtable}{cl}
\hline
\hline
章節 & 經文 (和合本修訂版)\\
\hline
11:1 & \begin{tabularx}{0.7\textwidth}{X} 當將你的糧食撒在水面上，因為日子久了，你必能得著它。 \end{tabularx} \\ \\ \relax
11:2 & \begin{tabularx}{0.7\textwidth}{X} 將你所擁有的分給七人，或八人，因為你不知道會有甚麼災禍臨到地上。 \end{tabularx} \\ \\ \relax
11:3 & \begin{tabularx}{0.7\textwidth}{X} 雲若滿了雨，就必傾倒在地上。樹向南倒，或向北倒，樹倒在何處，就留在何處。 \end{tabularx} \\ \\ \relax
11:4 & \begin{tabularx}{0.7\textwidth}{X} 看風的，必不撒種；望雲的，必不收割。 \end{tabularx} \\ \\ \relax
11:5 & \begin{tabularx}{0.7\textwidth}{X} 你不知道氣息如何進入孕婦的骨頭裡；照樣，造萬物之神的作為，你也無從得知。 \end{tabularx} \\ \\ \relax
11:6 & \begin{tabularx}{0.7\textwidth}{X} 早晨要撒種，晚上也不要歇手，因為你不知道哪一樣發旺；前者或後者，或兩者都一樣好。 \end{tabularx} \\ \\ \relax
11:7 & \begin{tabularx}{0.7\textwidth}{X} 光是甜美的，眼見日光是多麼好啊！ \end{tabularx} \\ \\ \relax
11:8 & \begin{tabularx}{0.7\textwidth}{X} 人活多少年，就當快樂多少年，然而也當想到黑暗的日子；因為這樣的日子必多，所要來臨的全是虛空。 \end{tabularx} \\ \\ \relax
11:9 & \begin{tabularx}{0.7\textwidth}{X} 年輕人哪，你在年少時當快樂；在年輕時使你的心歡暢，做你心所願做的，看你眼所愛看的；卻要知道，為這一切，神必審問你。 \end{tabularx} \\ \\ \relax
11:10 & \begin{tabularx}{0.7\textwidth}{X} 所以，當從心中除掉愁煩，從肉體除去痛苦；因為年少和年輕之時，全是虛空。 \end{tabularx} \\ \\ \relax
12:1 & \begin{tabularx}{0.7\textwidth}{X} 你趁著年輕、衰老的日子尚未來到，就是你所說，我毫無喜悅的那些歲月來臨之前，當記念造你的主。 \end{tabularx} \\ \\ \relax
12:2 & \begin{tabularx}{0.7\textwidth}{X} 不要等到太陽、光明、月亮、星宿變為黑暗，雨後雲又返回； \end{tabularx} \\ \\ \relax
12:3 & \begin{tabularx}{0.7\textwidth}{X} 看守房屋的發顫，強壯的屈身，推磨的婦女因人少而停工，從窗戶往外看的眼光變為昏暗； \end{tabularx} \\ \\ \relax
12:4 & \begin{tabularx}{0.7\textwidth}{X} 街門關閉，推磨的聲音微小，鳥一叫，就驚醒，唱歌女子的聲音也都微弱； \end{tabularx} \\ \\ \relax
12:5 & \begin{tabularx}{0.7\textwidth}{X} 人怕高處，路上有驚慌；杏樹開花，蚱蜢成為重擔，慾望不再挑起；因為人歸他永遠的家，弔喪的在街上往來。 \end{tabularx} \\ \\ \relax
12:6 & \begin{tabularx}{0.7\textwidth}{X} 不要等到銀鏈折斷，金罐破裂，瓶子在泉旁損壞，水輪在井口斷裂， \end{tabularx} \\ \\ \relax
12:7 & \begin{tabularx}{0.7\textwidth}{X} 塵土仍歸於地，像原來一樣，氣息仍歸於賜氣息的神。 \end{tabularx} \\ \\
[1ex]
\hline
\hline
\end{longtable}
這段經文可分三個段落,第一,傳道者勸勉讀者要行善,這是充實虛空的辦法.第二,勸勉年青讀者要快樂的過生活.第三,勸勉讀者趁年青的時候記念神.
\newline
\newline
(一)勸勉讀者要行善
\newline
\newline
(十一1-6),這段經文有人解作你要把握時機去投資,你若能勤力工作,必有好的收穫.你的投資是要多方面的.另外的說法,是認為這段提及行善,我比較贊同這種說法.因為你要將自己的糧食拿出去,撒在水面.回教徒受猶太教影響很深,可蘭經記有許多舊約的先知,它也承認耶穌是一位先知；不過模罕默德是最大的先知.他們傳說有一位皇帝的養子沉船,但因有人每日撒食物在水面上,他拾起來,得到充飢.所以回教徒習慣把食物撒在水上,因為日久必得賞賜.第二節描寫七人,八人,這是希伯來文的特色,可參考彌迦書五章五節,意思是指很多人,你要將糧食分給多人.中國人有句說話「助人為快樂之本.」你能常常幫助人,便會忘記自己的痛苦,你的人生便不感覺虛空；別人得著快樂時,你也會感覺快樂.初期教會有行善之習慣,好像百夫長哥尼流廣行善事；彼得復活了那廣行善事的多加；提前第五章記載登記的寡婦；學者認為他們是年紀老邁的寡婦,他們有行善的心方被接納.提多書三章1節,提醒眾人預備行各樣的善事.今天世界充滿痛苦患難,假若我們基督徒能夠與他們分享,神的名必定大得榮耀；受幫助的人也大得益處,他們也可以認識到這是愛的具體行動.有人稱這行動是「自動的共產」.初期教會彼此幫助缺乏的,詩篇三十七3「你當倚靠耶和華而行善,住在地上,以他的信實為糧.」香港有一教會關心貧苦人,與一間政府醫院聯絡,去幫助那些經濟有困難的人.他們撥款到醫院福利部去救濟那些人,而只須醫院送回受助者的姓名地址,以便日後教會探訪之用,這是非常好的見證.我們對所有人都一視同人去幫助他們,求神教導我們應該怎樣去做,似乎今日教會的慈惠工作做得不夠.行善可以見證主,也充實我們的人生.
\newline
\newline
(十3-6),提醒我們應該勤力工作,傳道者提出一切分內的工作,總要盡責盡力；我們不可懶惰,懶惰的人不能過一個充實的人生.第3節較難解釋,許多解經家指這是自然律,雲層厚,但有雨,人不能控制自然律；也不能明白所有的事.因此我們要將勤補拙,傳道者提醒我們,努力工作可充實空虛的人生.
\newline
\newline
(二)勸勉青年人要快樂的過生活
\newline
\newline
(十一7-10),聖經中的少年人,是指未夠三十歲的人,少年人即青年人.我們年青人,也包括其他人,應該快樂的生活.有人看見這段經文,便說傳道者是個樂天派；但神的心意是要我們在有生之年,快樂的過生活.第7節,他用眼目可以看見日光,來提醒者讀者應該喜樂,很少人為自己的眼睛可以看見,而感謝神!但若你眼部做過手術,你能重見天日,便會很快樂.我讀神學時,有一位內地會宣教士作證,他被囚在監獄的黑房三年.第一年過去,第二年,他感覺很孤單痛苦；到第三年時便開始埋怨神,懷疑神；在他最軟弱的時候,神為他開出路,使他得自由.今天我們有機會能夠看見日光,應該感謝神!享受神給我們一切的恩典,享受神在基督耶穌裡給我們的一切.假若我們信主後,並不感到享受.可能你還有許多地方未清楚,還有地方未對付；其實這種喜樂是最高尚的,我們的神就是我們最高的享受.當你與你所愛的人在一起時,乃是最高的享受；難怪傳道者說,人活多年,就當快樂多年.
\newline
\newline
為何傳道者勸勉年青人在年青時要快樂呢?
\newline
\newline
1.年青人有充沛精力智力,記憶力也好,學習能力快得多；這是神給我們美好的恩賜,你可以做自己喜歡做的事.昨天聖雅各福群會,在報章發表他們對香港青年人的調查報告,指出香港青年人「老成而無大志」,不會積極做人,盼望香港青年基督徒在人生打好基礎,以致被神使用,我們應該為年青的日子快樂.
\newline
\newline
2.作者在第10節說:「因為一生的開端,和幼年之時都是虛空的.」在我們人生初期是很痛苦的,不容易渡過,升學,戀愛,就業等.這會使人生很空虛,難怪在求偶時期的青年人都很消瘦；到中年時期,安定了便肥胖起來.我們基督徒青年人要想辦法來彌補人生的虛空,(十一9)「在幼年的日子,使你的心歡暢,行你心所願行的,看你眼所愛看的.」我們是否放縱情慾呢?傳道者繼續說下去:「為這一切事,神必審問你,所以你當從心中除掉愁煩,從肉體克去邪惡.」,基督徒不是放縱主義者,基督徒的享受和快樂,是要在神面前快樂享受；我們不能撇開神的審問,我們的享受是有責任,不是不理會別人死活.我們的自由,不能羞辱神的名.保羅在(林前十23)說:我們所行的,要有益處,要能造就人,使神的名能得榮耀,這樣的話,我們樣樣都可以做.所以傳道者提及的快樂,是有責任,而不是放縱的.我們可以盡量去享受,但我們能否在基督台前不受責備呢?
\newline
\newline
(三)應趁年青的時候記念神
\newline
\newline
(十二1-7),這裡稱神為「造你的主」,提醒讀者神與人的關係,我們應該尊敬祂,人是被造的物而已,應該將神的榮耀歸給祂.上面看過幼年的日子虛空,這裡更提及年老的日子也是虛空的.(十二7)「塵土仍歸於塵土,靈仍歸於賜靈的神.」人不是死了就一了百了,他仍要面對神；在地上無惡不作,似乎很興盛；但他死後,要面對神的審問,每個人要回到神那裡去.
\newline
\newline
作者在(2-6)用很美麗的喻意來代表人生年老的日子,在(3-4)節作者用一間朽壞的大屋來代表年老的境況.年老時手腳發顫,腰也不能伸直.推磨是形容牙齒,牙齒漸漸稀少；眼睛老花,看不清楚.街門關閉代表耳聾,推磨的響聲微小代表氣魄弱,說話也不能大聲.4節提及歌唱的女子也沒有氣力唱歌.回應上文,雀鳥一叫人就醒來,形容老人家早起,青年人卻熟睡.第5節形容老人家怕高,走路不穩.杏樹是冬天開花的樹,花先是粉紅色,後轉為白色；象徵老人家滿頭白髮,以及老人死後,形容弔喪的情況.第6節銀鍊折斷,形容老人家像油燈跌下來,因為人生是寶貴的,所以稱作銀鍊.金罐應譯作金碗,用作載油,表示載油的銀鍊折斷；打破那個金碗,這樣完結人的一生.我們在年老時,想做的事不能做,不能像年青時享受快樂的人生；吃東西又沒有牙,旅行時,又不可以爬山；有很多限制,想事奉神,卻有心無力.因此作者提醒年青人,趁著尚未年老的時候,要記念造物主.為甚麼要記念造物主,我以為有兩個原因:
\newline
\newline
第一,你會把握你年青的日子,不浪費神給你的機會,便不會太遲,趁機享受神給你的恩典.「少壯不努力,老大徒傷悲.」年輕時浪費我們的時光和機會,到年老時回顧,便感傷痛.岳飛的「滿江紅」中有話說:「莫等閒,白了少年頭,空悲切!」故此我們要趁有機會,來把握時機.我們要將黃金的時間奉獻給神,因為年老時便不中用,想獻上也無價值了.
\newline
\newline
第二,在年青有太多的抉擇和需要,要神的憐憫,否則我們便易犯錯,過失敗的人生.年老人已完成了他的責任,差不多一切已成了定局；年青人在人生中尚有許多可能性,他可能成為偉人,他尚有機會.因此更需要神的造就和引導,在我們的交友,學業,婚姻等事上,我們要將神放在我們的抉擇中,不要自己作主.(詩十六8)「我將耶和華常擺在我面前,因他在我右邊,我便不至搖動.」這是成功的秘訣,我們要尊神為大,為主,任何決定要將神放在我們的決定中；這樣,我們的人生必蒙神賜福.求主幫助我們,將傳道者的勸告放在我們的生命中.
\newline
\newline


\newpage

\section{第55屆~港九培靈會~劉承業~第九講　總結}
\label{sec:509}
\textbf{敬畏神的人生}
\newline
\newline
連結: \href{https://www.hkbibleconference.org/session-message/view/509}{\texttt{https://www.hkbibleconference.org/session-message/view/509}}
\newline
\newline
\hyperref[sec:508]{< < < PREV SERMON < < <}
~
\hyperref[sec:toc]{[返目錄]}
~
\hyperref[sec:518]{> > > NEXT SERMON > > >}
\newline
\newline
傳道書 12:8-12:14
\newline
\begin{longtable}{cl}
\hline
\hline
章節 & 經文 (和合本修訂版)\\
\hline
12:8 & \begin{tabularx}{0.7\textwidth}{X} 傳道者說：「虛空的虛空，全是虛空。」 \end{tabularx} \\ \\ \relax
12:9 & \begin{tabularx}{0.7\textwidth}{X} 再者，傳道者因有智慧，將知識教導眾人；他思量，考察，並列舉出許多箴言。 \end{tabularx} \\ \\ \relax
12:10 & \begin{tabularx}{0.7\textwidth}{X} 傳道者專心尋求可喜悅的言語，是憑正直寫的誠實話。 \end{tabularx} \\ \\ \relax
12:11 & \begin{tabularx}{0.7\textwidth}{X} 智慧人的話語如同刺棒；這些嘉言好像釘穩的釘子，都是一個牧者所賜的。 \end{tabularx} \\ \\ \relax
12:12 & \begin{tabularx}{0.7\textwidth}{X} 我兒，還有一點，你當受勸戒：著書多，沒有窮盡；讀書多，身體疲倦。 \end{tabularx} \\ \\ \relax
12:13 & \begin{tabularx}{0.7\textwidth}{X} 這些事都已聽見了，結論就是：敬畏神，謹守他的誡命，這是人當盡的本分。 \end{tabularx} \\ \\ \relax
12:14 & \begin{tabularx}{0.7\textwidth}{X} 因為人所做的事，連一切隱藏的事，無論是善是惡，神都必審問。 \end{tabularx} \\ \\
[1ex]
\hline
\hline
\end{longtable}
傳道者在總結本書時,提及他的教訓有兩方面.
\newline
\newline
(一)虛空的虛空
\newline
\newline
(十一1-8)先提到八節的話,後又重提一章2節的說話,他寫到(十二1-7)的時候,想到人老珠黃,何等悔恨!所以他便總括一方面的教訓,是虛空的虛空.另外他再提及人生積極的一方面,第9節提及他將智慧教訓眾人；他將人應該盡的本分,就是敬畏神和謹守祂的誡命.現在我們先看,傳道者憑甚麼原因來這樣去教訓人呢?他不是憑自己的臆測,第一個原因:(十二9)說:「再者,傳道者因有智慧,仍將知識教訓眾人,又默想,又查考,又陳說許多的箴言.」他承認自己有智慧,跟著他指出自己智慧來源；不是憑直覺,又不是天生的.(十二11)說:「智慧人的言語,好像刺棍；會中之師的言語,又像釘穩的釘子,都是一個牧者所賜的.」他承認他的智慧來自牧者,在舊約多處聖經提及神是牧者,(詩廿三1,詩八○1,結廿四,賽四○,)皆提及神是好牧人.到新約時,耶穌承認自己是好牧人；因此從新舊約亮光中,牧者是神自己.傳道者指出他的智慧全是由神所賜,他只不過負責將神的話傳講,帶給會眾.這也是希伯來文的特點,承認一切智慧全是由神來的.第二個原因,在(十二8)下半節,他指出自己智慧是來自默想,(詩一2)「惟喜愛耶和華的律法,晝夜思想,這人便為有福.」一個有福的人是時常默想神話的人,詩篇也有多處地方強調這點.第三個智慧的來源:(十二9)他又考查,他用神給他的聰明智慧來研究；例如文士以斯拉,他定意考究神的律法,所以以色列人歸回後,得到以斯拉的教導.原來傳道書的作者,他自己也做這種工作；承認自己的智慧來自神,來自默想,也來自考究.現代的人沒有神的話來引導,連信徒也沒有考查,研究,默想神的話語；因此在屬靈方面不夠營養,以致生命枯萎,沒有成長.許多時候連傳道人,也缺少在神話語多研究,多默想；以致按著正意分解真理的道,難怪他的羊群不能長大.
\newline
\newline
(十二12)說:「我兒,還有一層,你當受勸戒,著書多,沒有窮盡,讀書多,身體疲倦.」傳道者提醒我們不要單單讀普通的書,因為學海無涯,著書多,沒有窮盡.每一日都有許多的新書,讀書太多,也會疲倦.因此我們不要讀普通的書,他教導我們,要讀書中之書；我特別勸勉年青人,你用許多時間讀其他的書；但你用幾多時間來讀聖經.默想是得力的生活藝術,我們若在神面前等候,會使我們得力.我青年時在政府部門工作,但每個月一次抽空到太平山頂默想,對我的工作有很大幫助.現在我喜歡散步,除了能鬆弛神經外；還可以默想當日神賜給我的話.有一段時間,我每日在咭上寫下一節聖經,背誦它,默想它,使我很得幫助.由於傳道者默想神的話語,便使他得著智慧來教導眾人.
\newline
\newline
在(十二10)說:「傳道者專心尋求可喜悅的言語,是憑正直寫的誠實話.」他指出自己智慧的來源,來自正直和誠實的話,以致可以對讀者說,他所講的是好的,是他由衷之言,是他誠實的話；可能有人不大喜歡聽他的話,但這些話在後來是有功效的.他說智慧人的言語好像刺棍,在舊約時的刺棍是看牛用的.在棍上加上刺,當牛不馴服時,便用刺來管教牛.當智慧語言發出的時候,可能聽的人會覺得不好受,但最後可以教導他走在正道.另一方面,聰明的人會覺得這是對他的幫助.傳道者指出自己的說話是誠實話,並不避諱,保羅在(弗四15)教導我們「惟用愛心說誠實話.」,他講誠實話,表示他有愛心.(十二11)說:釘穩的釘子,有兩種解釋.第一種說法,這釘子是指趕牛棍上的釘子,表示耐用.第二種說法這釘子是穩妥之意.釘子不會移動,可以掛東西,表示你若接受智慧言語,你的人生好像錨一樣穩固,不會飄浮無定.所以智慧者的言語,成為我們的倚靠和幫助；因為它引導,提醒,勸勉我們,使我們知道如何走人生的道路.傳道者提出他的言語,無論好聽與否,都是有權威性的；因為來自神那裡,這些說話,是他多年在神面前研究和默想的結果.他之所以教訓人,因這是他的由衷之言,是他專心尋求,是可喜悅的言語.不是隨隨便便,風花雪月的話；所以我們每一個讀傳道書的人,都須要打開我們的心來閱讀.
\newline
\newline
(二)敬畏神和謹守祂的誡命作者極希望人能夠敬畏神
\newline
\newline
(十二13)說:「這些事都已聽見了,總意就是敬畏神,謹守祂的誡命,這是人所當盡的本分.」我解釋過敬畏神的意思,現在再複習一下敬畏神的意思；再補充以往未提及的,怎樣才是敬畏神呢?
\newline
\newline
1.無論做甚麼都盡力去做.以前我讀過一間大學,有一次在介紹學生工作辦事處的壁報板上,我看見有一幅漫畫.繪上一個青年人正在打掃,他背後棄置一些破碎的碗碟,在他面前站著位年老的主管.他說:「他們說我們負責.」,這漫畫點出負責,在不好的事情上負責.我們不但要向神負責,還要向神交賬；在一切的事情上,也包括好的,壞的；我們所做公開或隱蔽的事,有一日也要在祂面前公開的,向祂交賬和負責任.
\newline
\newline
2.敬畏神必須謹守祂的誡命.在(十二13)中說:「敬畏神是謹守祂的誡命,」謹守是小心認真地持守之意.傳道者要我們認真遵行神的話語,無論困難或容易,願意與否,也要遵行神的話.甚至可能要咬著牙關,也要遵守,這才是謹守的態度.
\newline
\newline
雅各書第一章告訴我們聽道也要行道,那裡指出為何聽道而不行道,有兩個原因:
\newline
\newline
1.單聽道而不行道,就是自己欺哄自己.
\newline
\newline
2.我們會很容易忘記所聽的,像一個人對著鏡子,很快便會忘記他的樣貌.我們住在香港,有許多機會聽道,又可參加機構辦的訓練班；但若果我們聽道而不行道的話,很快便會忘記所聽的.故此傳道者提醒我們,要謹守遵行神的話語.
\newline
\newline
敬畏神的另一個含意,在第八章中我們曾提及過,是要接納神的權柄,承認神在我們生命中完全的權柄.當我們這樣做的時候,即使我們遭遇不如意,令我們傷心難過,我們接納神是全能,全知,全愛的；我們仍可靠神的恩典,心安理得,而不怨天尤人.接納神的主權,是承認神的旨意,因神的道路是高過你的道路；雖然你暫時不明白,但會仍然相信神沒有做錯事情.我們也提及過,敬畏神的人,承認真正的福氣是從神來的；可能是將來屬靈的,也可能是屬物質現世的.每一日我們所享受的粗茶淡飯,都是神所賜的.一個敬畏神的人,不單勤儉,而且也要過有節制的生活.傳道者在整卷書強調敬畏神的人生,只有敬畏神才是有福的,能超越人生的虛空,使人生更充實.
\newline
\newline
傳道者在以往的經節中,提及如何實際的生活,他只提及今天實際生活中,來敬畏神.我們新約的信徒,該多向舊約信徒學習,我們要認識神是滿有威嚴和權能的；以致我們在祂面前,不能不屈膝敬拜.詩篇第八篇指出人算甚麼,你竟眷顧他?人雖渺小,神仍然憐憫我們,揀選我們,所以我們不能在神面前,隨隨便便的生活,按著自己意思來做人.每一個人不能脫離神而生活,這樣的人才是有福,才能過豐滿的人生.今日許多基督徒只是停留在耶穌是救主的地步,一信耶穌就可以上天堂,得永生；無疑這是正確的信仰,但我們卻不要因此而滿足.傳道書作者指出,我們不單要相信神,還要敬畏神.相信神,是承認神,但並不表示我們順服神；我們要敬畏神,敬愛祂,將祂該得的榮耀歸給祂.因祂是滿有尊榮,偉大的；我們作任何事,只要存著一個敬畏神的心,這才是有福的人生.這樣的人生,才是有意義的人生.
\newline
\newline


\newpage

\section{第55屆~港九培靈會~史保羅~第一講　教會何時需要復興}
\label{sec:518}
\textbf{第一講　教會何時需要復興}
\newline
\newline
連結: \href{https://www.hkbibleconference.org/session-message/view/518}{\texttt{https://www.hkbibleconference.org/session-message/view/518}}
\newline
\newline
\hyperref[sec:509]{< < < PREV SERMON < < <}
~
\hyperref[sec:toc]{[返目錄]}
~
\hyperref[sec:519]{> > > NEXT SERMON > > >}
\newline
\newline
啟示錄 2:1-3:22
\newline
\begin{longtable}{cl}
\hline
\hline
章節 & 經文 (和合本修訂版)\\
\hline
2:1 & \begin{tabularx}{0.7\textwidth}{X} 「你要寫信給以弗所教會的使者，說：『那右手拿著七顆星，在七個金燈臺中間行走的這樣說： \end{tabularx} \\ \\ \relax
2:2 & \begin{tabularx}{0.7\textwidth}{X} 我知道你的行為、勞碌、忍耐，也知道你不容忍惡人。你也曾察驗那自稱為使徒卻不是使徒的，看出他們是假的。 \end{tabularx} \\ \\ \relax
2:3 & \begin{tabularx}{0.7\textwidth}{X} 你能忍耐，曾為我的名勞苦而不困倦。 \end{tabularx} \\ \\ \relax
2:4 & \begin{tabularx}{0.7\textwidth}{X} 然而，有一件事我要責備你，就是你把起初的愛心拋棄了。 \end{tabularx} \\ \\ \relax
2:5 & \begin{tabularx}{0.7\textwidth}{X} 所以你要回想你是從哪裡墜落的，並且要悔改，做起初所做的工作。你若不悔改，我要到你那裡去，把你的燈臺從原處挪去。 \end{tabularx} \\ \\ \relax
2:6 & \begin{tabularx}{0.7\textwidth}{X} 然而你還有一件可取的事，就是你恨惡尼哥拉派的行為，這種行為也是我所恨惡的。 \end{tabularx} \\ \\ \relax
2:7 & \begin{tabularx}{0.7\textwidth}{X} 凡有耳朵的都應當聽聖靈向眾教會所說的話。得勝的，我必將神樂園中生命樹的果子賜給他吃。』」 \end{tabularx} \\ \\ \relax
2:8 & \begin{tabularx}{0.7\textwidth}{X} 「你要寫信給士每拿教會的使者，說：『那首先的、末後的，死過又活了的這樣說： \end{tabularx} \\ \\ \relax
2:9 & \begin{tabularx}{0.7\textwidth}{X} 我知道你的患難和貧窮—其實你卻是富足的，也知道那自稱是猶太人的所說毀謗的話，其實他們不是猶太人，而是撒但會堂的人。 \end{tabularx} \\ \\ \relax
2:10 & \begin{tabularx}{0.7\textwidth}{X} 你將要受的苦，你不用怕。看哪！魔鬼要把你們中間幾個人下在監裡，使你們受考驗，你們要遭受苦難十日。你務要至死忠心，我就賜給你那生命的冠冕。 \end{tabularx} \\ \\ \relax
2:11 & \begin{tabularx}{0.7\textwidth}{X} 凡有耳朵的都應當聽聖靈向眾教會所說的話。得勝的必不受第二次死的害。』」 \end{tabularx} \\ \\ \relax
2:12 & \begin{tabularx}{0.7\textwidth}{X} 「你要寫信給別迦摩教會的使者，說：『那有兩刃利劍的這樣說： \end{tabularx} \\ \\ \relax
2:13 & \begin{tabularx}{0.7\textwidth}{X} 我知道你的居所，就是有撒但座位之處；當我忠心的見證人安提帕在你們中間，在撒但所住的地方被殺之時，你還堅守我的名，沒有否認對我的信仰。 \end{tabularx} \\ \\ \relax
2:14 & \begin{tabularx}{0.7\textwidth}{X} 然而，有幾件事我要責備你，就是在你那裡有人服從了巴蘭的教訓；這巴蘭曾教唆巴勒將絆腳石放在以色列人面前，使他們吃祭過偶像之物，並且犯淫亂。 \end{tabularx} \\ \\ \relax
2:15 & \begin{tabularx}{0.7\textwidth}{X} 同樣，你那裡也有人服從了尼哥拉派的教訓。 \end{tabularx} \\ \\ \relax
2:16 & \begin{tabularx}{0.7\textwidth}{X} 所以，你當悔改；若不悔改，我很快就到你那裡來，用我口中的劍攻擊他們。 \end{tabularx} \\ \\ \relax
2:17 & \begin{tabularx}{0.7\textwidth}{X} 凡有耳朵的都應當聽聖靈向眾教會所說的話。得勝的，我必將那隱藏的嗎哪賜給他，並賜他一塊白石，石上寫著新的名字，除了那領受的以外，沒有人認識。』」 \end{tabularx} \\ \\ \relax
2:18 & \begin{tabularx}{0.7\textwidth}{X} 「你要寫信給推雅推喇教會的使者，說：『神的兒子，那位眼睛如火焰、雙腳像發亮的銅的這樣說： \end{tabularx} \\ \\ \relax
2:19 & \begin{tabularx}{0.7\textwidth}{X} 我知道你的行為：愛心、信心、勤勞、忍耐；又知道你末後所行的善事比起初所行的更多。 \end{tabularx} \\ \\ \relax
2:20 & \begin{tabularx}{0.7\textwidth}{X} 然而，有一件事我要責備你，就是你容忍那自稱是先知的婦人耶洗別教唆我的僕人，引誘他們犯淫亂，吃祭過偶像之物。 \end{tabularx} \\ \\ \relax
2:21 & \begin{tabularx}{0.7\textwidth}{X} 我曾給她悔改的機會，她卻不肯悔改她的淫行。 \end{tabularx} \\ \\ \relax
2:22 & \begin{tabularx}{0.7\textwidth}{X} 看吧，我要使她病倒在床上。那些與她犯姦淫的人若不悔改他們的行為，我也要使他們同受大患難。 \end{tabularx} \\ \\ \relax
2:23 & \begin{tabularx}{0.7\textwidth}{X} 我又要殺死她的兒女，眾教會就知道，我是那察看人肺腑心腸的，我要照你們的行為報應各人。 \end{tabularx} \\ \\ \relax
2:24 & \begin{tabularx}{0.7\textwidth}{X} 至於你們其餘的推雅推喇人，就是一切不隨從這教訓，不明白他們所謂撒但深奧之理的人，我告訴你們，我不會再把別的擔子放在你們身上。 \end{tabularx} \\ \\ \relax
2:25 & \begin{tabularx}{0.7\textwidth}{X} 你們只要持守那已經有的，直到我來。 \end{tabularx} \\ \\ \relax
2:26 & \begin{tabularx}{0.7\textwidth}{X} 那得勝又遵守我命令到底的，我要賜給他權柄制伏列國； \end{tabularx} \\ \\ \relax
2:27 & \begin{tabularx}{0.7\textwidth}{X} 他必用鐵杖管轄他們，如同打碎陶器， \end{tabularx} \\ \\ \relax
2:28 & \begin{tabularx}{0.7\textwidth}{X} 像我也從我父領受了權柄一樣。我又要把晨星賜給他。 \end{tabularx} \\ \\ \relax
2:29 & \begin{tabularx}{0.7\textwidth}{X} 凡有耳朵的都應當聽聖靈向眾教會所說的話。』」 \end{tabularx} \\ \\ \relax
3:1 & \begin{tabularx}{0.7\textwidth}{X} 「你要寫信給撒狄教會的使者，說：『那有神的七靈和七顆星的這樣說：我知道你的行為，就是名義上你是活的，實際上你是死的。 \end{tabularx} \\ \\ \relax
3:2 & \begin{tabularx}{0.7\textwidth}{X} 你要警醒，堅固那些剩下、快要死的，因為我發現你的行為，在我神面前沒有一樣是完全的。 \end{tabularx} \\ \\ \relax
3:3 & \begin{tabularx}{0.7\textwidth}{X} 所以，要記得你所領受和聽見的；要遵守，並要悔改。你若不警醒，我必如賊一樣來到；我幾時來到你那裡，你絕不會知道。 \end{tabularx} \\ \\ \relax
3:4 & \begin{tabularx}{0.7\textwidth}{X} 然而，在撒狄你還有幾位是未曾污穢自己衣服的，他們會穿白衣與我同行，因為他們是配穿的。 \end{tabularx} \\ \\ \relax
3:5 & \begin{tabularx}{0.7\textwidth}{X} 得勝的必這樣穿白衣，我也不從生命冊上塗去他的名；我要在我父面前，和我父的眾使者面前，宣認他的名。 \end{tabularx} \\ \\ \relax
3:6 & \begin{tabularx}{0.7\textwidth}{X} 凡有耳朵的都應當聽聖靈向眾教會所說的話。』」 \end{tabularx} \\ \\ \relax
3:7 & \begin{tabularx}{0.7\textwidth}{X} 「你要寫信給非拉鐵非教會的使者，說：『那神聖、真實的，拿著大衛的鑰匙，開了就沒有人能關，關了就沒有人能開的這樣說： \end{tabularx} \\ \\ \relax
3:8 & \begin{tabularx}{0.7\textwidth}{X} 我知道你的行為。看哪，我在你面前給你一個敞開的門，是沒有人能關的。我知道你有一點力量，也遵守我的道，沒有否認我的名。 \end{tabularx} \\ \\ \relax
3:9 & \begin{tabularx}{0.7\textwidth}{X} 那屬撒但會堂的，自稱是猶太人，其實不是猶太人，而是說謊話的，我要使他們來到你腳前下拜，使他們知道我已經愛你了。 \end{tabularx} \\ \\ \relax
3:10 & \begin{tabularx}{0.7\textwidth}{X} 因為你遵守了我堅忍的道，我也必在普天下人受試煉的時候保守你免受試煉。 \end{tabularx} \\ \\ \relax
3:11 & \begin{tabularx}{0.7\textwidth}{X} 我必快來，你要持守你所有的，免得人奪去你的冠冕。 \end{tabularx} \\ \\ \relax
3:12 & \begin{tabularx}{0.7\textwidth}{X} 得勝的，我要使他在我神的殿中作柱子，他必不再從那裡出去。我又要把我神的名和我神城的名—從天上我神那裡降下來的新耶路撒冷，和我的新名，都寫在他上面。 \end{tabularx} \\ \\ \relax
3:13 & \begin{tabularx}{0.7\textwidth}{X} 凡有耳朵的都應當聽聖靈向眾教會所說的話。』」 \end{tabularx} \\ \\ \relax
3:14 & \begin{tabularx}{0.7\textwidth}{X} 「你要寫信給老底嘉教會的使者，說：『那位阿們、誠信真實的見證者、神創造的根源這樣說： \end{tabularx} \\ \\ \relax
3:15 & \begin{tabularx}{0.7\textwidth}{X} 我知道你的行為，你也不冷也不熱；我巴不得你或冷或熱。 \end{tabularx} \\ \\ \relax
3:16 & \begin{tabularx}{0.7\textwidth}{X} 既然你如溫水，也不冷也不熱，我要從我口中把你吐出去。 \end{tabularx} \\ \\ \relax
3:17 & \begin{tabularx}{0.7\textwidth}{X} 你說：我是富足的，已經發了財，一樣都不缺，卻不知道你是困苦、可憐、貧窮、瞎眼、赤身的。 \end{tabularx} \\ \\ \relax
3:18 & \begin{tabularx}{0.7\textwidth}{X} 我勸你向我買從火中鍛鍊出來的金子，使你富足；又買白衣穿上，使你赤身的羞恥不露出來；又買眼藥抹你的眼睛，使你能看見。 \end{tabularx} \\ \\ \relax
3:19 & \begin{tabularx}{0.7\textwidth}{X} 凡我所疼愛的，我就責備管教。所以，你要發熱心，也要悔改。 \end{tabularx} \\ \\ \relax
3:20 & \begin{tabularx}{0.7\textwidth}{X} 看哪，我站在門外叩門，若有聽見我聲音而開門的，我要進到他那裡去，我與他，他與我一起吃飯。 \end{tabularx} \\ \\ \relax
3:21 & \begin{tabularx}{0.7\textwidth}{X} 得勝的，我要賜他在我寶座上與我同坐，就如我得了勝，在我父的寶座上與他同坐一般。 \end{tabularx} \\ \\ \relax
3:22 & \begin{tabularx}{0.7\textwidth}{X} 凡有耳朵的都應當聽聖靈向眾教會所說的話。』」 \end{tabularx} \\ \\
[1ex]
\hline
\hline
\end{longtable}
我能夠第二次在這培靈研經大會主講,實在感到非常高興!相信這是蓋世的培靈研經大聚會,因為有這麼多聽眾聚集在一起聽神的話；我為這聚會已經等了很久,盼望這個日子早些來到.
\newline
\newline
首先我代表加拿大的教會,特別是多倫多民眾教會,向你們問安!有許多中國人住在多倫多,相信其中一定有些人,是你們的親戚,我在那邊有時也有機會對中國人傳信息.當我來香港之前,曾經在本教會的聚會中宣佈,將於八月一日開始赴港九培靈研經大會；我問他們是否願意為這聚會禱告,全會眾約兩千人一致表示願意.我相信你們也都一定為這聚會禱告的.當神的子民開始禱告時,神就開始工作,當神開始工作,事情就開始改變.
\newline
\newline
我首先要講的是復興,何時需要復興?怎能夠復興?復興時發生甚麼?
\newline
\newline
復興,有時是關於整個教會的復興,整個教會實在需要復興.復興是普世教會的大事,可以說,教會復興並沒有影響個人.以多倫多本教會而論,全教會人數約有二千人,我們也極需要復興.
\newline
\newline
這幾個晚上,我們將講及個人復興的事,教會的每個人復興,教會才能復興.但是,復興不可能一下子臨到二千人,復興是在個人生活上發生的,是關於你和我的復興,我們不是說教會何時需要復興,而是說個人何時需要復興.我們不是說許多人何時需要復興,而是說自己何時需要復興.之後,我們要講及關於聖靈的事,有四個信息關於聖靈與個人的工作,聖靈是甚麼?聖靈在我們身上作甚麼?你我與聖靈的關係如何.末了要題到兩方面:第一,對整個世界傳道,最後要講基督化的家庭.
\newline
\newline
神要賜福給這十天的聚會,請各位代禱.
\newline
\newline
啟示錄(二章1-4)論以弗所教會,(12-15節)論別迦摩教會；第二章至第三章,主吩咐要寫信給七個教會.對教會題示將來所要發生的事.這些教會可代表這時代的某些教會,這裡論及歷史上的以弗所教會,和背道的老底嘉教會,藉以警戒歷史時代的教會.
\newline
\newline
今天晚上我所題的,是真實的教會.我們知道,以弗所地方有個教會,別迦摩地方也有個教會；在這些教會中有些人,好像你和我一樣.主說:他們中間有好的,也有不好的,因為他們發生了屬靈上的問題.
\newline
\newline
以弗所教會初期非常愛主,但後來他們離棄了起初的愛心,他們當時需要主重新感動,因為在屬靈上他們發生了問題.他們需要復興,因他們離棄了起初愛基督的心.
\newline
\newline
關於愛主的人有三個特點.
\newline
\newline
第一個特點:對所愛的人,開始的時候非常熱心,一對年青情人,起初傾談滔滔不絕,正如剛找到耶穌作他個人救主的人,最喜歡講耶穌基督,到處傳講耶穌基督.年青的基督徒最熱心,他們作見證；但是經過一段時間,就開始不如初信時的熱心了.有時我們說,我們成長很多了,在基督徒生活中成長了；但我相信這不是主耶穌所要講的,主所要講的是「你們離棄起初的愛心.」我們喪失了愛基督的熱心!
\newline
\newline
第二個特點:他們與祂有團契的渴望,我們希望在祂面前傾心吐意.年青的基督徒愛主,所以要他們到教會並非難事,他們知道若有兩三個人奉主的名聚會,主就在其中.他們樂意參加教會任何聚會,無論聚會,禱告,他們都深感主的同在.他們愛耶穌基督,所以喜歡與主說話,他們常常禱告,也喜讀聖經.他們知道讀經時主耶穌摸著了他,所以肯用功讀神的話；他們希望與主常有團契,但是年事漸長,便失去了團契的渴望.我們固然盼望在神的家中,卻常忽略了聚會,而且更逐漸減少了禱告的時間；到了年紀大些,就連讀經都少了.如果你問他:「為甚麼這樣呢?」他必說:「我作基督徒很久了,信仰上也成長了.」不過耶穌說:我們需要新的感動,需要新的復興.
\newline
\newline
第三個特點:我們盼望能夠事奉神.如果我們愛一個人,總盼望能為所愛的人做一點事.保羅在大馬色路上悔改時,他個人看見耶穌基督,他說的第一句話是「主啊!你要我作甚麼?」這是個初信主者所能接受的問題.當我們初愛主時,盼望能為主作一點事,在教會中殷勤服事；但逐漸地,我們卻失去了事奉的心志.有我們以為成長了,成熟了,已經用了很多時間事奉主了,可以放鬆一點.主說:「你是否已喪失了起初愛我的心呢?」請問在我們的一生中,是否已喪失了起初的愛心?我們必須省察,復興,我們需要神的感動；在基督徒生命中,最重要的乃是對基督的愛心.約翰福音末章記載,耶穌考問彼得的話是甚麼呢?是否耶穌要問彼得「你的教義對嗎?」「你相信這是正確的事嗎?」「你可否在這正確教義簽個名嗎?」我們所信的教義是正確的,我們所信的事是非常重要的；我們不但曉得教義的重要,我們更應當明白教義.實際上,教義並非最重要的事.或者我們要問:「彼得是否過著基督徒正確的生活?他是否去正確的地方?他所作的事對嗎?彼得的行為怎樣呢?」我們的行為是很重要的,作為神的兒女,當去正確的地方,該作正確的事,行為也是極重要的,但更重要的是耶穌問彼得:「你愛我比這些更深麼?」彼得說:「主啊!是的,你知道我愛你.」耶穌第二次又問他說:「你愛我麼?」答:「主啊!是的,你知道我愛你.」耶穌第二次又問他說:「你愛我麼?」答:「主啊!是的,你知道我愛你.」耶穌再三問彼得:「你愛我麼?」各位朋友!現在我們要問同樣的一個問題,在這廣大的基督徒聚會中,我們唱著信心的詩歌,在聚會中享受團契的喜樂；我們實在感謝神,使我們能參加這樣的聚會.但是我們是否確實愛主呢?你愛祂是否保持以往的愛呢?或許我們已經離棄了起初的愛心,若然；耶穌對你,對我說:「你們需要復興起初的愛心.」我們需要神新的感動!
\newline
\newline
主對別迦摩教會所講的另一信息,別迦摩教會有兩件事情做錯了,教會中有些人服從了巴蘭的教訓.巴蘭是個舊約時代的先知,那時候,以色列民進入應許之地,他們曾經過曠野,擊敗所有的仇敵,唯一留下的是摩押人.摩押王名叫巴勒,他看見以色列人所經過之地,覺得以軍的能力是其他軍隊所沒有的；他知道以色列軍兵不久將與本國交戰,他們奇異的軍力使他非常恐懼!所以就召見先知巴蘭,要他咒咀以色列軍兵,使他們喪失那種超越的能力.如果我們讀民數記,就知道巴蘭曾想奉命去做,三次想咒咀以色列人,但是神把他的咒語扭轉,反成祝福,他回報國王說:「神賜福給以色列民,我不能咒詛神所賜福的子民.」巴蘭向巴勒解釋這事的教訓,以色列人之所以有超越的能力,是因為他們從邪教中分別出來──遵守神給他們的誡命.和外邦人接觸時,不貪他們的東西,不參與他們的社交生活,絕對不與他們通婚.神說:「為我的緣故,要和他們分別出來,我就賜福你們,賜能力給你們.」巴蘭卻教導巴勒將絆腳石放在以色列人面前,叫他們吃祭偶像之物,行姦淫的事.因此,他們就開始與摩押人參雜,以致失去能力.這是主耶穌論到別迦摩教會的事.祂提醒基督徒應從異教中分別出來.這班基督徒知道巴蘭的故事,也知道以色列人為何失去能力；可是同樣的事卻發生在別迦摩教會；這班人生活在極其邪惡的城市,到處皆是邪教的人和廟宇；只要從其中分別出來,就可得到神所賜的能力,否則,也就失去能力.耶穌說,這就是教會中所發生的問題,是基督徒生活中所存在的問題.當我讀到這故事的時候,我省察自己的生活；我知道神要屬於祂的人從世界分別出來,基督徒的生活,與世人的生活應是不同的.當然生活中需要有些朋友,這些朋友也就是我們作見證的對象；但必須使他們看見我們與世俗的人有所不同.彼得說:「只求在神面前有無虧的良心.」(彼前三16)亦即是說:在神前有正確的良心,在人前有正確的行為,應真正為神而活.也許有人要問:「你的信心如何?為何有這樣的信心呢?未知有沒有人問過關於你的基督徒生活?有時你說,從來沒機會讓你在人前作見證；這是因你的生活與別人無異,所以沒有人和我們談論到基督的信仰.
\newline
\newline
最後,主對別迦摩教會說,有些人服從尼哥拉一黨人的教訓.他們要從信徒中分別出來,他們自命不凡,以為自己屬靈方面比他人強.他們忘記了所有的信徒在神前都是祭司,每一基督徒都是神的聖殿,坐在禮堂中的每位信徒也都是神的聖殿.在神眼中我們都是祭司,若有人認為自己屬靈上比別人強,那就是犯了別迦摩教會所發生的同樣錯謬,他們在屬靈上產生了驕傲.曾有人告訴我,他在播音方面有屬靈的經驗；我為他屬靈的經驗感謝神,我對他說,神賜福你使你有這樣的經驗.他說:和他有同樣屬靈的經驗的人也不及他.有時,我們有了屬靈的經驗,就會輕看其他沒有屬靈經驗的人.本來有好的屬靈經驗,但結果反而造成屬靈的驕傲.有時,我們從神得到屬靈的恩賜就驕傲；也許有了禱告恩賜,也許有教導叫人明白神話語的恩賜,也許有講道的恩賜,也許有唱詩的恩賜,也許有醫治的恩賜；我們運用恩賜覺得光榮驕傲,覺得比別人強,因為能作別人所不能作的事.有了屬靈的恩賜反而造成不良的結果,就是屬靈的驕傲.有時我們以老基督徒為驕傲,看輕那些初信的基督徒；有了長期屬靈的經驗,就容易發展到屬靈的驕傲.應當小心!耶穌說:「尼哥拉一黨的教訓是我所憎惡的.」如果我們自覺比他人屬靈,那是主所憎惡的.我認為值得以屬靈經驗自豪者,惟有使徒保羅,他為主多作工,多受苦；因為信仰,多次被囚；在(腓三4),他說:「其實我也可以靠肉體,若是別人想他可以靠肉體,我更可以靠著了.」在下文他列出一些可以驕傲的事情；到了這章末了,他說:「這不是說,我已經得著了,已經完全了；我乃是竭力追求,或者可以得著基督耶穌所以得著我的.」(腓三12),他曾為主受苦,他在基督徒路程中已經走了很久,將到行程的終點；可是屬靈方面,他沒有驕傲.
\newline
\newline
當我們失去起初的愛心,當我們沒有和屬世的人分別出來時,當我們有屬靈驕傲時；我們需要復興.
\newline
\newline


\newpage

\section{第55屆~港九培靈會~史保羅~第二講　教會如何可以復興}
\label{sec:519}
\textbf{第二講　教會如何可以復興}
\newline
\newline
連結: \href{https://www.hkbibleconference.org/session-message/view/519}{\texttt{https://www.hkbibleconference.org/session-message/view/519}}
\newline
\newline
\hyperref[sec:518]{< < < PREV SERMON < < <}
~
\hyperref[sec:toc]{[返目錄]}
~
\hyperref[sec:520]{> > > NEXT SERMON > > >}
\newline
\newline
歷代志下 7:12-7:14
\newline
\begin{longtable}{cl}
\hline
\hline
章節 & 經文 (和合本修訂版)\\
\hline
7:12 & \begin{tabularx}{0.7\textwidth}{X} 夜間耶和華向所羅門顯現，對他說：「我已聽了你的禱告，也選擇這地方歸我作獻祭的殿宇。 \end{tabularx} \\ \\ \relax
7:13 & \begin{tabularx}{0.7\textwidth}{X} 我若使天閉塞不下雨，或使蝗蟲吃這地的出產，或降瘟疫在我子民中， \end{tabularx} \\ \\ \relax
7:14 & \begin{tabularx}{0.7\textwidth}{X} 這稱為我名下的子民，若是謙卑自己，禱告，尋求我的面，轉離他們的惡行，我必從天上垂聽，赦免他們的罪，醫治他們的地。 \end{tabularx} \\ \\
[1ex]
\hline
\hline
\end{longtable}
歷代志下第七章12-14節這段經文,論及神對祂名下的子民說,怎能從神得到幫助.神說:如果你發生了問題,需要幫助；這裡可以找到辦法；如果你要我聽你的禱告,要求罪得赦免,要你的國家蒙福,必須要作四件事:一,自卑,二,禱告,三,尋求神的面,四,轉離惡行.
\newline
\newline
(一)自卑,這是神在很久以前,對屬祂的子民所說的話
\newline
\newline
但是神的話不是只應用在一個時代,我認為我是屬於神的子民,這就是神對我說的話；如果你認為你也屬於神,這也是神對你所說的話.我相信這是神在今晚,要對你和我所說的話,神望著我們,祂知道我們每個人的名字.神知道我名叫保羅,我聽見祂所說的話.神說:如果你想得到我的話,這裡你得到了；如果你需要我的能力,這就是辦法；如果你想我聽你的禱告,這就是你尋求的方法,你必須自卑,禱告,尋求神的面,轉離惡行.但願這段經文應用於我們每個人身上,這是關於你我個人的事,關乎每個人的事；可能這裡有幾千的聽眾,神對這大班的人說話是很容易的；不過可惜我們聽了卻沒有遵行.
\newline
\newline
兩年前我曾在這裡講過,當我讀中學時,每個學生都要學習演講；我幾乎忘記了老師教給我的話.有一天,老師對我們說:「演講時面對許多聽眾,實在是最困難的事,如果面向天花板說話,就容易得多.」有的人演講時眼睛望著窗外,有時牧師講道面對著聖經.如果想和人交通說話時眼睛必須看那人才對.老師曾教我們學習一件愚蠢的事,叫我們每個擔任演講的人站起來,各自挑選三位熟識的同學；然後用手指指著,眼睛望著,對每一位說:「我現對你說話,你要留心聽.」我永不忘記當時的情境,我望著天花板對其中一人說話,望著窗門對第二個人說話,望著自己的腳對第三人說話；結果並無一人留心聽我說的話.老師要我們練習多次.經過那次以後,已經四十年之久,我不斷在講台上；我知道我是傳講神的話,我禱告,然後張開我的眼睛,我聽唱詩,覺得非常享受.在我心中經常有一同樣的禱告,求神親自對每個聽眾說話；使他們滿覺得單獨與神同在,忘記他周圍的一切,惟獨神對他說話.我們今晚讀了這段經文,盼望各位都能採取這樣的途徑,讓我們忘記未曾來赴會的基督徒,忘記坐在我們前後左右的基督徒,忘記有人在台上講道,惟獨感到單獨與神同在.
\newline
\newline
神說:「你要復興,必須自卑.」我們要得救,必須在神前自卑,否則不能成為基督徒,因從未有一人在作基督徒之前不是罪人的.我確實知道我是個罪人的時候,我就接受耶穌基督作救主；在神面前俯伏自卑,承認自己的罪,才能成為基督徒.但是,我們常忘記得救是本乎神的恩典,也忘記我們活著也是神的恩典.藉著神的恩典,今晚你我才能夠在這裡,若非神的恩典,我就不能傳信息,你也不能走進禮拜堂,你也聽不到我所講的.基督徒的生活完全靠賴於神的恩典,要記得神說屬於祂的子民應當自卑；若不依靠神,我們工作就沒有能力,曾被神使用的人應學習這個功課.
\newline
\newline
摩西生在埃及,他成為埃及重要的人物,他與王室有關,且有很大的影響力.他很富有,他受過埃及高深教育,大有學問；他自己也知道當時的他,是何等重要的人物.有一天,他在外看見自己的同胞作奴隸,心中感到同胞的痛苦需要拯救；他認為自己有足夠的影響力,可以肩負重任,因他是王室中之一份子,到處受人敬重.所以他曾兩次想利用自己的能力和權力去幹,可是都失敗.神使他離開埃及,安放他在曠野四十年之久.在那荒涼之地,錢財權勢都盡失了；終日與羊群為伍,連羊群也非屬自己所有,是他岳父叫他看管的.在這四十年中,摩西真是一無所有,後來他在何烈山,看見燒著的荊棘.神對他說:「我要打發你去見法老,使你可以將我的百姓以色列人,從埃及領出來.」在這樣情形之下,是否他要說:「神啊!你終於找到最適當的人選,我就是你所要的那個人了；我有多方面的經驗,有很大的影響力,我可以去領百姓出來.」可是摩西回答:「我是甚麼人,竟能去見法老,將以色列人從埃及領出來呢?」摩西意思是說,他不可能做到,因影響力不夠,經驗不足,口才不好,還是另找別人吧!神對摩西說:「不管你是誰,不管你有多少能力,你只要對他們說:「我是神所差來的」他們就會聽你的話,他們知道有神的能力作你的後盾.」摩西自卑,所以神使用祂；這就是摩西在曠野學習了四十年的功課.
\newline
\newline
我是加拿大一間教會的牧師,我們每主日上午十一點聚會,每個主日我要面對兩千人傳神的道.數週前的週六晚我睡不著,我就思想主日的工作；我想到或有些人發生問題,或有些人心裡有負擔,有些人心裡傷慟,有些人婚姻上出了問題,有的人覺得心中有罪；我知道他們希望我能輔導他們.我對神說:「神啊!我沒有辦法幫助他們,也沒有辦法釋放他們的重擔；我不知道他們心裡的問題,我也無能赦免他們的罪,我未能除掉他們的罪惡感；我在主日崇拜時,也實在沒法適應他們的需要.」神對我說:「保羅!我知道你不能作,我知道你沒足夠的能力去作,我知道你沒有足夠的能力肩負他們的重擔；更無法赦免人的罪,除去任何人的罪惡感.但是,當你在講台認識自己不能作的時候,就有機會讓我幫你去作了,因聖靈將透過你去作.」在講台上可能我很軟弱,但在我裡面的聖靈卻是剛強的；或者我找不到問題的答案,但聖靈在我裡面作出了一切問題的答案.當我在講台上,我無法適應聽眾的需要時,但聖靈的能力足以適應他們的需要.當你在主日學,在詩班,在學校,在工作崗位上,你是否需要神所賜福呢?是否需要增加工作效率呢?神說:「你在我面前自卑,認識自己不能作,你才能認識神凡事都能；若然,祂必應允你的禱告.」
\newline
\newline
(二)禱告,神說:「我的子民若自卑,禱告......」首先是自卑,然後禱告
\newline
\newline
主耶穌曾講過一個故事:有兩個人到聖殿禱告,一個是法利賽人──是宗教中的人；一個是稅吏.法利賽人在聖殿中找的是顯眼地方,藉以引人注意,且高聲禱告說:「我不像別人,勒索,不義,姦淫,也不像這稅吏.......我是個非常重要的人,神啊!你一定要聽我的禱告.」耶穌說:「凡自高的,必降為卑.」至於那稅吏,卻找個隱藏的地方,不想引人注意；他俯伏在神面前,捶胸自咎,求神饒恕,因他自認是個罪人.耶穌說:「自卑的必升為高.」神不聽驕傲人的禱告,神聽自卑人的禱告；惟有自己謙卑降低,禱告才能上升到神的面前.
\newline
\newline
美國曾有位大奮興家名叫芬尼,他所講的道使許多聽眾熱烈接受；神使大復興臨到全美國,會眾渴慕救恩,急不及待；當講道還未結束,神奇妙的能力浩大,甚至使講道停止.和芬尼一起有個禱告的同伴,不斷為他禱告；芬尼講道時,同伴在密室內一直為他禱告,求神的能力臨到會場,求神帶領一切.神的話語宣揚出來,會眾中充滿了一種不可抗拒的能力.
\newline
\newline
在舊約中也有同樣的例子:約書亞和亞瑪力人爭戰,亞瑪力擁有龐大軍隊和兵器；約書亞和摩西商量對策,約書亞帶領一些人上陣.摩西和亞倫,戶珥上山頂,摩西拿著神的杖,這杖代表神的能力,神與摩西同在.在山頂高舉神的杖與神接觸；他們何時舉起執杖的手,以色列人就得勝,何時垂下,亞瑪力人就得勝.當摩西的手下沉時,亞倫和戶珥齊心合力支持摩西高舉神的杖；果然約書亞打了勝仗,以色列之得勝,並非因他們有多大的軍隊,有多好的將領；乃是打仗時有三個人不斷與神接觸.
\newline
\newline
你若需要神的能力,你必須禱告,這是復興的條件.你如果希望在生活的戰役中得勝,在基督徒生活成功,那你必須禱告.
\newline
\newline
在多倫多的本教會,一切設備很好；有空氣調節,有佷好的詩班獻詩,有很多好的主日學同工；教會中有二百四十位長老,有好的視聽器材,我為這一切感謝神!因為這使教會進行工作很大的方便.此外,我們還有很多肯禱告的弟兄姊妹；這班禱告的戰士,經常為教會一切事工禱告.如果這些東西,給我只有一個選擇的話,我寧可捨棄其他的一切；但我卻需要有禱告的戰士.神的工作,必須有禱告的支持.如果我們需要神的垂聽,需要神賜福,需要一切,就必須要禱告!」
\newline
\newline
(三)尋求神的面,聖經記載,神向尋求祂的人微笑,對不服從的人轉臉
\newline
\newline
多倫多本教會,有嬰兒的洗禮,為父牧者願意把兒女奉獻給神,讓孩童學習聖經,學習耶穌基督；希望孩童們長大之後接受耶穌基督作救主.數週前,我們把頭生九個月大的孫子奉獻給神.他是生在世界上最美麗,最聰明的孩童,像他的祖父母一樣；當我們奉獻自己的孫子時,我照樣地說:「奉父子聖靈的名奉獻給神,願神賜福你,保守你,用臉光照你,賜恩給你,神向你仰面,賜你平安.」這段經文是舊約摩西所說的話,他對亞倫說,如何為人祝福,讓神光照,在生活中讓神的面光照耀；若是罪孽,神必轉面不顧；因為罪孽使我們與神隔絕.遵行神的旨意,祂的臉必時刻向著我們.
\newline
\newline
(四)轉離惡行,這話是神對所羅門王說的
\newline
\newline
他獻殿的時候；神對他說:「這稱為我名下的子民,若是自卑,禱告,尋求我的面,轉離惡行；我必從天上垂聽,赦免他們的罪,醫治他們的地.」
\newline
\newline
(代下七14)神說這話是對著好人說的,並非對惡人說的；神看我們的內心,鑑察我們的心思意念.約翰對基督徒說:「我們若說自己無罪便是自欺,真理不在我們心裡了.」
\newline
\newline
神今天晚上也要對我們說:「親愛的基督徒!你若需要我的能力,必須轉離你的惡行.」你必須求神饒恕你的罪,讓神時刻抓住你.我們每天除了三餐謝飯之外,至少必須有兩次禱告；睡前須求神饒恕當天所犯的過錯,起床後必須求神賜給能力作正確的事,這是復興必需的條件.
\newline
\newline
如果我們能夠完成上述四個條件,今天晚上我們就得到復興了；當我們離開會場,我們就得到神的力,是我們從來未有的能力.
\newline
\newline
我們要自問:「我在生活中是否需要復興?」若然,必須在神面前自卑,要禱告,要尋求神的面,要轉離惡行.
\newline
\newline


\newpage

\section{第55屆~港九培靈會~史保羅~第三講　教會復興帶來甚麼}
\label{sec:520}
\textbf{第三講　教會復興帶來甚麼}
\newline
\newline
連結: \href{https://www.hkbibleconference.org/session-message/view/520}{\texttt{https://www.hkbibleconference.org/session-message/view/520}}
\newline
\newline
\hyperref[sec:519]{< < < PREV SERMON < < <}
~
\hyperref[sec:toc]{[返目錄]}
~
\hyperref[sec:521]{> > > NEXT SERMON > > >}
\newline
\newline
我剛才參觀過在大會中供應各種屬靈書籍的地方,這些書籍所寫關於基督徒屬靈的生活,是講台上所不能詳盡的；屬靈的書籍常常可以摸到人的心,這是講員所做不到的.書攤中有兩本書是我的著作,一是岌岌可危的當代教會.頭三晚所講的,就是這本書的首段；這本書也講及聖靈,並說方言的事；聖經所載關於聖靈充滿的經節,本書都一一提到.還有最近才譯成中文出版「約翰筆下的耶穌」這本書並非解釋全卷約翰福音,乃是提到約翰在每章中如何看耶穌.第一章約翰看耶穌是個創世者.第二章論耶穌是一個人,祂參加婚筵.第三章論耶穌和尼哥底母談話.全卷約翰福音共廿一章,每章論到耶穌.
\newline
\newline
今天晚上要講當教會復興的時候,會有甚麼事情發生.復興發生的事有種種不同,美國威爾斯那時的大復興和英國衛斯理時的大復興有所不同,但是無論如何,復興必然發生的,就是神的靈感動我們有合一的靈；如果沒有這種形出現,我們仍需復興.保羅說:「因祂使我們和睦,將兩下合而為一,拆毀了中間隔斷的牆.」(弗二14)未復興之前,各種不同的牆把我們隔開；當神的靈動工,這些牆就拆斷了.未復興之前,基督徒彼此之間有著感情的牆,有嫉妒,看別人比自己強就惱怒；當復興來到,就不再有這種情形發生了.請大家自問,是否你看見別人所有的,是自己所沒有,就生出嫉妒；看見別人做得到的事,自己做不到,就懷恨在心；見別人有所成就便憎恨呢?當復興未來到之前,教會中充滿了仇恨；神的靈臨到,這些嫉妒的城牆便拆毀了.如果你們哪一位心裡仍存嫉妒的話,那就必須要復興.神的靈不控制我們的生命,但如果心裡有嫉妒,生命就出了問題.還未復興之時,基督徒會憎恨別人,但耶穌教訓我們說,不可恨人,更要愛恨你的人.在教會中有些肢體之間,會各存成見,不喜歡交談,自以為是個好基督徒；其實這樣的基督徒,本身已有問題.可能他是個詩班成員,可能也是主日學教師；或更有可能他是個長老,牧師,在教會中任要職；但是,若本身有了問題,聖靈就不能充滿.我們實在需要從神而來新的感動,我們需要復興.
\newline
\newline
復興未出現之前,人被社會的地位隔斷,社會有貧富之懸殊,教育程度之高低；教會中也有不同之處；若某基督徒因為處身優裕而鄙視他人,那麼,他必須復興.復興給神的兒女帶來合一,破壞感情的牆就被拆毀了.
\newline
\newline
復興之後,宗派隔斷的牆也拆毀了.我個人的意見,宗派在世界上是絕對需要的,我不相信所有教會都是一樣的.各教會都有不同的基督徒,有的教會喜歡許多感情活動,有的教會的基督徒沒有甚麼活動；有些基督徒聚會時拍掌,但有些基督徒則不然.我們實在需要各種不同的教會.我要問:到底我們的宗派重要,還是我們的神重要?我的宗派或者是我的團契是否比神更重要?當神的靈感動我們的時候,我們雖然仍有不同的宗派；但是我們卻會有個同一的團契.在教會圈子裡可以找到許多同樣的例子.基甸會把聖經放在世界各酒店裡,該會有各種不同宗派的信徒,他們不談宗派,但最要緊的是將神的話傳到世界各處.基督教學生福音團契,是個各種不同宗派的團契；長老會,浸信會,宣道會,聖公會......等等的信徒都有.他們在一起,並非推動自己的宗派；而唯一目的,乃是和其他的基督徒有合一的團契.兩週前我往阿姆斯特丹,參加葛培理的佈道大會；出席者包括全世界各宗派的人,約有四千個傳道人在那裡,來自世界一百卅個國家；在會中卻未曾有人提及宗派的事,儼然屬乎一個宗派；我們彼此團契,只希望把神的訊息帶給別人；所以,當復興來到時,宗派的牆就要拆毀了.
\newline
\newline
約翰衛斯理建立了循道會,他曾夢見被帶到天門前,他和守門的天使談話,他問天使:「裡面有沒有天主教徒呢?」天使答:「沒有.」又問:「有長老會信徒嗎?」也答:「沒有.」再問:「有多少循道會的人在裡面呢?」答仍是:「沒有.」衛斯理覺得奇怪,不明所答.說:「甚麼都沒有,那麼誰在天上?」天使答:「天上所有的是一班愛神的人.」
\newline
\newline
當復興時,一切宗派的牆都拆毀了,教義的牆也都拆毀了.有的人說:教義是不重要的,但是我告訴你們,教義是絕對重要的,每個基督徒都應當認識自己所信的是甚麼,都應當知道為何要信.從聖經神的話,每個基督徒應知道所信的是誰.教義極為重要,但是很多時候,教義把人分開.一九八二年有部新基督徒百科全書發行,內容有關統計基督徒的數字,提及超過兩萬個不同的宗教團體；有些基督徒為此擔憂,為何一部聖經竟然產生諸多宗派呢?可是我們應該認識,在不同的宗派層面,基本上所信的是一樣,相信聖經中的基本真理；彼此有所不同,乃是有其他的理由.有時我們對於教會的真理,或浸禮方面有著不同的意見；有的宗派主張浸禮,有的洗禮,有的滴禮,有的教會沒有施洗.不同之處,並不影響到永遠穩固的問題；如果我已經得救了,是否再會失喪呢?又是否永遠得救呢?解釋預言有很多不同的看法.在這四個晚上,我們提到聖靈也有其他不同之處；有時我們提及聖靈同一樣的事,但是我們用不同的字眼來講聖靈；有時為這方面彼此辯論.當復興來到,一切不同之處,都成為次要；大家都注重於對付罪的問題,接受救贖.接受神的全能,接受天堂,接受聖經神的基本真理.
\newline
\newline
我小的時候,在加拿大有一次和家兄到農莊去,那時農場都用馬工作,我們看見農夫趕著載重的馬；到了山腳停下來,給馬兒休息片刻.當起程時,那隻黑馬先起步,隔數秒鐘那隻棕色馬才起步；因為起步不一致,雖努力向前仍然不動.於是農夫讓馬再休息,再次起步,棕馬在先,黑馬稍後,兩馬雖然盡力負軛,卻依舊不動.只好再休息,然後農夫命令兩馬向同一方向出動,才成功的載重上山.基督徒也常這樣,忙忙碌碌；但無濟於事,盡力傳福音,卻無人得救.很多時為了個人的興趣,盡力推動自己的小問題,以致主的工作被停留下來.當復興的時候,所有基督徒齊心合力以基督為元首,於最短時間內把福音帶給別人；對全世界的人傳說:「神愛世人,甚至將祂的獨生子賜給他們,叫一切信祂的,不至滅亡,反得永生.」除了合一之外,我們更應有熱愛人靈魂的心.
\newline
\newline
數年前我和內子同到麻省,慕迪夫婦的墳地,我們站在這位大佈道家他們的墓前；我思想,甚麼原因能夠在他身上,讓神如此重用他.慕迪傳福音的工作,對英國,美國有重大的影響,藉著他傳道,千千萬萬人得拯救.有人問慕迪:「甚麼原因能夠使你成功呢?為甚麼你所傳的道,使人對基督有反應呢?你所有的,而別人所沒有的是甚麼呢?」慕迪說:「我學習領人歸主時,我當先愛他；如果我終日批評他,責備他,我就得不到他；若真誠愛他,他就接受基督耶穌.」愛心的禱告是:「神啊!願你使我對失喪的人,有真誠的愛,願你將這愛轉成對別人的負擔.」
\newline
\newline
未知各位何時,是你最後一次對人的關懷呢?在你家庭中,或者你的父母,兄弟姊妹,未曾得救,你曾否為此整夜失眠?你曾否為家人得救的問題而夜半流淚?有時我們說要領人歸主,但卻沒這負擔,也沒有這種關懷；我們當求神使我們的負擔成為熱愛去得人,求主催促我有愛人靈魂的心去拯救別人.怎能有愛人靈魂的心呢?向人傳福音最重要的是甚麼呢?
\newline
\newline
摩西在申命記卅章19節說:「我今日呼天喚地向你作見證,我將生死禍福陳明在你面前:所以你要揀選生命,使你和你的後裔都得存活.」我們向人傳福音,為的不是多得些人參加教會,我們為主作見證,不是為叫人接受宗教；我們講道卻是要講生死的問題,禍福的問題,天堂地獄的問題；那麼,我們就有個熱愛人靈魂的心.當復興來到時,第一,信徒合一.第二,熱愛人靈魂的心.第三,有效的見證.
\newline
\newline
我曾提過美國的傳道芬尼,傳道時發生的大復興.他原是個律師,有一天,他到樹林裡,跪下在神前；他第一次接受耶穌基督,同時聖靈降臨他身上充滿他.這樣的事,對我們來說也常會發生過；當我們接受耶穌基督,然後聖靈充滿我們；但芬尼卻是得到拯救,同時也得到聖靈充滿.當他回到他的辦公室,他向周圍的人作見證,在市鎮中到處傳講耶穌；凡聽他作證的人,在當日就接受耶穌基督,也有和他們接觸之後的人,在次日也接受了耶穌.芬尼被聖靈充滿,他的見證立刻見效.
\newline
\newline
約翰福音第十二章,記載耶穌基督使拉撒路從死裡復活後的事,(9-11)是很重要的經文.約翰說,那些宗教領袖計劃要殺害耶穌,這時拉撒路剛從死裡復活；他們以為不但要殺害耶穌,連拉撒路也必須除滅!因為好些猶太人為拉撒路的緣故,回去信了耶穌.未知我們中間有多少這樣的人,又有多少人像拉撒路一樣,真正從死裡活過來；我們已經從神手中得了拯救,祂以神蹟使我們從黑暗轉向光明,因信耶穌基督,使我們從死亡進入永生.有沒有人看見我們而相信耶穌呢?我們是否有基督徒的生活,好像拉撒路一樣呢?當我們復興時,我們就有能力為主作見證,看見神大能的顯現.
\newline
\newline
昨天晚上我們從歷代志下,神提到我們復興的四個條件.如果我們具備這四個條件:第一,神要垂聽我們的禱告.第二,神要饒恕我們的罪.第三,神要醫治我們的病.神說:如果你們實行這些條件,我就在你們中間彰顯大能.
\newline
\newline
不久之前,我曾有機會到中國大陸探訪基督徒的家庭,他們述說神的大能的彰顯；講到一切神蹟的發生,也講到神醫治的能力.我聽了感到如同使徒時代的教會一樣.
\newline
\newline
在多倫多本教會,神大能的推動約有一年多,當神的靈運行,祂的大能彰顯；當我邀請別人接受主耶穌基督時,他們無需別人的推動,也無需我從中催促；很多時候有人打電話到教會,只要牧師和他們傾談,他們就歸向主了；有時從街道上來的人,牧師也有機會帶領他們信主.
\newline
\newline
希望各位有機會參加,如去年五月舉行的世界性聚會；我們希望從那聚會得到一百萬加幣,來支持百多位在外佈道的加拿大人.願神感動本教會的信徒支持這一切的工作.去年所得到的和以往所得的差不多,約為一百十三萬加幣.有個特別的情形,就是無需我催迫,請求,用壓力叫人奉獻；在民眾教會裡有神的靈感動,因此,我們常看見神大能的彰顯.
\newline
\newline


\newpage

\section{第55屆~港九培靈會~史保羅~第四講　聖靈與信徒}
\label{sec:521}
\textbf{第四講　聖靈與信徒}
\newline
\newline
連結: \href{https://www.hkbibleconference.org/session-message/view/521}{\texttt{https://www.hkbibleconference.org/session-message/view/521}}
\newline
\newline
\hyperref[sec:520]{< < < PREV SERMON < < <}
~
\hyperref[sec:toc]{[返目錄]}
~
\hyperref[sec:522]{> > > NEXT SERMON > > >}
\newline
\newline
由今天晚上開始,一連四個晚上,我們將講到聖靈的工作.今晚要講關於聖靈基本的事,大概你們曾經聽過,但我感到神的話應當多次的講.
\newline
\newline
聖靈住在每一個神的兒女裡面.主耶穌在五旬節聖靈降臨之前,講及聖靈降臨時要發生的事.(約十四15-17)耶穌說,聖靈來時要住在我們裡面.保羅說:「如果神的靈住在你們心裡,你們就不屬肉體,乃屬聖靈了；人若沒有基督的靈,就不是屬基督的.」(羅八9)在新約,有些經文重複多次提到聖靈；有的是耶穌基督說的,有的是使徒們說的.使徒多次提及我們是聖靈的殿,神的靈住在每個聽祂話語的人裡面.聖靈不住在世界任何建築物中,全世界沒有一間建築物,是為了聖靈的緣故而存在的；聖靈不住在禮拜堂,不住在廟宇裡,不住在會堂中；聖靈住在神的兒女的心裡面.這裡是一座很好的禮拜堂,我們聚集在一起；但當最後一個信徒離開禮拜堂,聖靈也跟著離去；聖靈不住在我的禮拜堂,聖靈住在我個人裡面,我是聖靈的殿.
\newline
\newline
保羅說:「豈不知你們是神的殿,神的靈住在你們裡面嗎?」(林前三16)有時我們聽人說受聖靈的洗.當然這方面是很重要的,當我們信了主,聖靈就為我們施洗,使我們進入基督的身體裡,進入神的家中；當我們成為基督徒的那日,我們便受洗歸入聖靈.「我們不拘是猶太人,是希利尼人,是為奴的,是自主的；都從一位聖靈受洗,成了一個身體,飲於一位聖靈.」(林前十二13)基督徒不須求神給他聖靈,如果你是個基督徒,你已經有聖靈了；聖經說,如果我沒有聖靈,我便不是基督徒.
\newline
\newline
聖靈不控制每個基督徒.當保羅探望撒瑪利亞信徒時,他們已聽見福音,已接受耶穌作救主；也接受了聖靈,但他們卻仍未被聖靈充滿.(徒八4-16)保羅對以弗所信徒說:「不要醉酒,酒能使人放蕩,乃要被聖靈充滿.」(弗五18)就是說,基督徒已得到聖靈,然而未被聖靈充滿.當我們接受耶穌作救主,我們就得著聖靈了.聖靈能控制我們生活的任何方面,所以每個基督徒都有聖靈居住在裡面；但是有時我們生活的某部份,沒有受到聖靈的控制.好比一所房子,居住的人信了主,聖靈就進入這屋子.但屋中有些房間,沒有讓聖靈進入,可能聖靈只到樓下尚未上樓；我們當讓聖靈進入我們一切的地方,這就是聖經所講的「被聖靈充滿.」在基督徒生活上,「被聖靈充滿」有時發生在得救的一剎那間,比如哥尼流和他的家庭,當彼得向他們傳福音時,他們就接受了；同日他們就被聖靈充滿.(徒十)但是有時「被聖靈充滿」是在人得救以後的某段時間.有的基督徒得救很久聖靈才進入裡面.使徒保羅悔改之後,經過一段時間才被聖靈充滿.行傳九章說到保羅悔改的事.保羅在大馬色的路上,悔改接受耶穌；但一直等到他進入大馬色,見亞拿尼亞之後才被聖靈充滿.當得救的同時,可能你被聖靈充滿,這是最奇妙的事；但是如果你得救後很久,從未被聖靈充滿；這也不要緊,因為以後你仍有可能被聖靈充滿.各個不同的教會,對這一點都有不同的表達；有的教會認為被聖靈充滿,是第二次的祝福,而得救時是第一次的祝福.有的教會認為聖靈充滿,是得勝的生命,在一般培靈會都有這樣說法,這似乎是培靈會的基本教訓.要得到得勝的生命,就必須被聖靈充滿；或者我們開始有得勝的生命時,我們就被聖靈充滿.有的教會認為被聖靈充滿是受聖靈的洗,第一次受聖靈的洗,是我們得救的時候.我曾說過,當我們受洗進入神的家時,而對教會來說,更有第二次受聖靈的洗.還有的教會說是聖靈充滿的生活；在培靈會上有時也這樣講,這是宣道會所用的名詞.我們實在不應該對此有所爭論,重點並非如何稱呼的問提到,最要緊的是問自己曾否被聖靈充滿；只要曾被聖靈充滿,怎樣稱呼無關重要,我們當知此事並非發生一次就有永遠的效果.有時我問基督徒:「你們有否受過聖靈充滿?」有的回答:「我被聖靈充滿,已廿五年前的事了.」有的說:「是在五年前.」不錯,可能當時他被聖靈充滿；事實上,廿五年之前被聖靈充滿,於今天晚上是不足夠的.基督徒的生命,每日需要被聖靈充滿.曾經有一次我學會了這件事,當我們第一次被聖靈充滿時,大多數的人都會記得的.我每早晨未作任何事之前,先禱告求主今日用聖靈充滿我.為了今日一切的問提到,我們需要被聖靈充滿.如果有能力作見證,今日需要被聖靈充滿.我們被聖靈充滿,不是單單像拿了個杯倒入水,我們需要繼續下去；不然,這杯水經久就變成污濁不清了.聖靈充滿,如寬闊的河流不斷流下；聖靈的河流每日潔淨我們,我們從中得到神的能力.
\newline
\newline
聖經有兩處,提及被聖靈充滿,如同醉酒一樣.彼得在五旬節之日曾說:「這些人不是喝醉酒,乃是被聖靈充滿.」保羅也說:「不要醉酒......乃要被聖靈充滿.」(弗五18)被聖靈充滿所發生的情形和喝醉酒差不多.醉酒的人,非常勇敢.多年前我在南非的約翰尼斯堡講道,有一晚,當我回酒店的路上,經過酒吧門口；看見兩人打架,其中一人體形巨大,胸膛寬闊,肌肉強壯；另一人生成細小枯廋,但當時打架卻似乎要打敗那強壯的人,哪來的勇氣呢?因他喝醉了酒.我們被聖靈充滿時,神給我們有屬神能力的勇氣.彼得和約翰醫好了在聖殿門口瘸腿的病人,當時宗教的領袖反對這件事,就控告他們；在權勢壓迫之下,他們仍然毫不畏懼,聖經說:「他們就都被聖靈充滿,放膽講論神的道.」(徒四31)舊約有位士師俄陀聶,「耶和華的靈降在他身上,他就作了以色列的士師,出去爭戰.」舊約還有一強壯的人名叫參孫,「耶和華的靈大大感動參孫......卻將獅子撕裂......」(士十四6)當人被聖靈充滿時,就好像醉酒一樣,而且對人和氣,一般喝醉酒的人,常和不相識的人談話.(徒五42)記載早期使徒的工作,在五旬節的事發生不久之後,「他們每日在殿裡,在家裡,不住的教訓人,傳耶穌是基督.」當他們被聖靈充滿,他們放膽出去向任何人傳耶穌基督.
\newline
\newline
我們沒有為耶穌作見證,乃因我們沒有被聖靈充滿；如果我們要向人講論耶穌,就必須如同醉酒一樣.通常醉酒的人是勇敢的,和睦的與人談話,並且堅持所要作的事；如果你趕走個醉酒的人,他會再回來的,無論趕走他多少次都會重返的.醉酒的人繼續作他所要作的事,永不停止；我們被聖靈充滿時,也是如此.不單有勇氣作神的工,和氣與人講耶穌基督,堅持,繼續地講耶穌基督.
\newline
\newline
保羅對以弗所長老說:「所以你們應當儆醒,記念我三年之久,晝夜不住的流淚,勸戒你們各人.」保羅說他到城市三年之久,日以繼夜地傳福音,堅持要傳福音.一個被聖靈充滿者,在屬靈方面好像喝醉酒的人,繼續堅持作神要他作的事.喝醉酒的人是勇敢的,和藹地與人談話,非常堅持的；但有時他也自知,不是全然成功,偶然也會失敗的.
\newline
\newline
在行傳廿一章,保羅當時成了羅馬帝國的囚犯,他一直被監禁,最後被害.他在耶路撒冷監獄,有機會向刑罰他的人講道,他作見證如何蒙神拯救；他講及在大馬色路上的經歷並耶穌基督的大能.但是在那次的機會,沒有人接受福音,保羅所傳的福音,聽者全無反應.在五旬節的奇遇中,當使徒傳道,有三千人接受基督；但在耶路撒冷仍有更多的人沒有接受,惟那三千人有反應.保羅也曾向阿基伯王作見證,他盡力傳講福音；但阿基伯王卻沒有接受耶穌做救主.
\newline
\newline
我曾向各位提議,作神的工需要勇氣；可是事先必須聖靈充滿才有勇氣.我們須知,向人傳講耶穌基督,未必一定有反應；這是令人憂慮的事!或者你的母親還未得救,你要向她作見證,要為她禱告,要在她面前活出基督徒的生活；但仍然不可能代她接受耶穌基督,這部份必須由她親自做的.沒有任何人可以代替她得救的,我們為耶穌作見證時,必須認識清楚這一點.無論我們如何盡力作見證,但並非人人都肯接受.
\newline
\newline
當人被聖靈充滿時,會歡喜唱歌,這不在乎有好的歌喉,曾否接受過音樂訓練,或是所唱的使人厭煩,全都不理會,總是要唱.聖經說一個被聖靈充滿的人,好像醉酒者一樣.如果我們被聖靈充滿,就像個歡喜唱歌的人一樣.「當用詩章,頌詞,靈歌,彼此對說,口唱心和的讚美主.」(弗五19)我非常欣賞這次大會的歌唱,當我們的弟兄帶領唱詩時,我十分享受.本主日我在楊濬哲牧師那禮拜堂講道,整個主日崇拜,講到信仰方面和詩歌的事.一個教會若是不歡喜唱歌,我們必須小心!當我們被聖靈充滿時,不期而然心中發出和聲共鳴.
\newline
\newline
我重複今天晚上所講的:
\newline
\newline
第一,聖靈在每個基督徒裡面,毋須追求；如果沒有聖靈的就不是基督徒.
\newline
\newline
第二,聖靈不控制每個基督徒,只要我們讓聖靈進入生活的每部份.聖靈擁有我們,我們就被聖靈充滿了.當我們被聖靈充滿時,我們具有醉酒者的特性,為神的工作勇敢去做,和氣與人講論耶穌基督；且堅持不斷地傳耶穌,甚至遇到挫折失敗,仍然堅持繼續.當我們被聖靈充滿時,神會把詩歌放在我們內心引起共鳴.
\newline
\newline


\newpage

\section{第55屆~港九培靈會~史保羅~第五講　為何未被聖靈充滿}
\label{sec:522}
\textbf{第五講　為何未被聖靈充滿}
\newline
\newline
連結: \href{https://www.hkbibleconference.org/session-message/view/522}{\texttt{https://www.hkbibleconference.org/session-message/view/522}}
\newline
\newline
\hyperref[sec:521]{< < < PREV SERMON < < <}
~
\hyperref[sec:toc]{[返目錄]}
~
\hyperref[sec:523]{> > > NEXT SERMON > > >}
\newline
\newline
今天晚上再次提及聖靈的工作.有的人未被聖靈充滿.
\newline
\newline
第一,乃因他們不曉得勝靈的事.有時我們不知道聖靈的事,是因為牧師很少對我們講；我們講了很多神的信息,講了很多耶穌基督的信息；但很多時候我們沒有講到聖靈.有些人不知道聖靈,以中性的他代表聖靈,認為不屬人性；他們知道神有位格,基督也有位格.很多人以為聖靈是種能力.很多人不認識聖靈的工作,因為聖靈沒有強迫我們讓祂控制我們每部份的工作,耶穌基督並沒有強迫我們成為一個得救的人；祂使我們有自由,選擇的權利,祂從不強迫讓祂進入你的生命中.祂說:「我站在門外扣門.」如果要祂進來,我們必須開門；聖靈也是如此,每個人得救時都有聖靈；但不是每個人都讓聖靈控制他整個的生命.保羅說:「要被聖靈充滿.」他並非說:「必須被聖靈充滿.」被聖靈充滿與否,祂給我們自由選擇.
\newline
\newline
第二,有的人不認識聖靈的工作,因他們不明白,他們以為是無關痛癢的事；有人悔改只盼望進入天堂,他們接受耶穌作救主,只因為怕下地獄.他們只求永生能得穩固,其他並不重要.是否在我們的教會中有這樣的人,他們很準時參加聚會；但在教會中從不接受任何事奉的工作,也沒有擔起任何職責.通常教會都需要招集一些義工,參加詩班,或教主日學等.照我看來教會裡只有百分之十的人,是肯做義工的,而其他百分之九十的人,就不願意被這些職務所束縛,也不願意擔當教會任何義務工作.因為既有職責,就必須要盡責去做.不錯,他們是蒙神恩典所救贖的,他們是基督徒,他們也向著天堂的路上走,是教會中好的弟兄姊妹,他們也奉獻金錢；但是他們不肯用多點時間事奉神,這樣的人對聖靈並非不知道,不過他們以為那是不關緊要的事；只要得救,卻不讓聖靈去控制他們的一生,他們是以自我來控制自己的一生.
\newline
\newline
第三,信徒未被聖靈充滿,是因生命中有罪.以弗所書第五章,提及基督徒生命中的罪之後,保羅說:「要被聖靈充滿.」信徒生命有病,必須被聖靈充滿,才能醫治.換言之,若信徒生命中有罪,就不能被聖靈充滿；聖靈不容罪孽留在信徒身上.要對付罪孽,就必須承認自己的罪；必須悔改認罪,神才除掉我們的罪.全部新約聖經都有記載,關於認罪,絕對不能將重擔委給他人.多年前我的教會有位姊妹寫信給我說:「親愛的牧師,我要向你認罪；因為我一向不斷批評你所講的道,現在我感到非常的後悔.」有的人以為這姊妹所作是很對的,但我卻不以為然；她寫了這封信給我,她自己的重擔隨即放下,但這重擔卻轉到我的背上了.我們向人認罪時要留心,別讓自己的認罪,成為別人的重擔.那位姊妹繼續參加教會的聚會,見面時和她握手；每次和她握手時,我沒有想到她認罪的事,卻想到她批評我的事；我想,未知她是否已停止對我批評?很多時我們認罪,卻把自己的重擔放在別人身上；其實,這位姊妹可以自己進入她的禱告室,跪在神的面前為著批評牧師的事認罪,從此不再批評,並為牧師代禱,這樣她豈不是很開心嗎?因她已把重擔卸給神了；沒有把重擔放我身上,我也開心.
\newline
\newline
有時,對別人認罪並非件好事情,那麼,我們對於被傷害的人應該如何?今天我在房中用了個鐘頭讀大衛和拔示巴的事,大衛犯了重罪.第一,他得罪了拔示巴,第二,得罪了拔示巴的丈夫烏利亞,他對拔示巴犯了淫亂的罪,並殺害她的丈夫烏利亞；因大衛是以色列的王,所以對全以色列民都犯了罪.先知拿單來見大衛,用手指著大衛說:「你就是那人.」大衛向神認罪,聖靈沒有記載他向任何人認罪,沒有記載他要求拔示巴饒恕,也沒記載他曾經要求烏利亞的親人饒恕,也沒有向以色列民公開認罪；聖經多處提及大衛為此向神認罪,聖經也清楚告訴我們,神赦免了他的罪.大衛犯罪向神認罪,神饒恕他的罪之後,他仍作以色列的領袖；神賜福給他比前更多.他和拔示巴結婚,所生的第一個兒子死了,第二個兒子名叫所羅門.拔示巴是大衛犯罪時的對象,先知拿單是大衛犯罪後,頂撞大衛的人；當大衛年老,他們二人同去見大衛,要大衛向以色列民宣佈所羅門是他王位的繼承人.大衛另一個兒子亞多尼雅也想繼承王位；拔示巴和拿單,都希望大衛阻止這事發生.結果是所羅門繼承了王位.大衛作以色列王時,曾經發生了許多問題；神使大衛戰勝一切的仇敵,建立王國,並使他與拔示巴所生的兒子所羅門承繼王位.請大家把這事和教會的事作個比較,當基督徒犯罪認罪之後,神就饒恕他的罪；但是教會中卻有人不願意饒恕他,甚至有人永遠記他所犯的罪,以致使他再無機會繼續事奉神.這是萬分不合理的,因為聖經說:「我們若認自己的罪,神是信實的,是公義的,必要赦免我們的罪,洗淨我們一切的不義.」(約壹一9),祂不再記念這罪,祂要把這罪投下深海中.
\newline
\newline
今天晚上,你若認你的罪,神就饒恕你的罪；如果你直接向神認罪,神必應允你,赦免你的罪.神既已饒恕你的罪,至於其他的人是無關緊要的；神既應允,祂是信實的,公義的,必然就應允.
\newline
\newline
我們如何向神認罪呢?我想大衛向神認罪,是我們最好的榜樣.詩篇第五十一篇是大衛與拔示巴犯罪後所寫的詩；希望各位把這篇詩和自己身上的罪作個對比,你就曉得在神面前怎樣懊悔認罪了.大衛在整個認罪過程中,單獨與神說話,他說:「我犯罪惟獨得罪了你.」我們犯罪都是得罪了神,並非得罪了別人；我偷了別人的東西,傷害了別人,甚至殺害了別人；我犯了律法的事,所以在我們生活中犯了罪,一定要向神認罪.
\newline
\newline
犯罪,是我們未被聖靈充滿的主要原因；有的人未被聖靈充滿,因在生活中有罪,身上的罪,未曾向神承認.如有這樣的情形,你無須對任何人說；但是如果你要得赦免,必須對神說.我們無法得到聖靈充滿,除非在神面前悔改認罪.
\newline
\newline
在座中各位!今天晚上,神有沒有對你說話呢?你心裡有沒有問題?你未被聖靈充滿,因還有未向神承認的罪.
\newline
\newline
但願你說:「我渴慕聖靈充滿,我要活出聖靈充滿的生活；我需要神控制我全人,我要做大衛所做的,在神面前將我的罪認出來.」
\newline
\newline


\newpage

\section{第55屆~港九培靈會~史保羅~第六講　聖靈在信徒生命中的工作}
\label{sec:523}
\textbf{第六講　聖靈在信徒生命中的工作}
\newline
\newline
連結: \href{https://www.hkbibleconference.org/session-message/view/523}{\texttt{https://www.hkbibleconference.org/session-message/view/523}}
\newline
\newline
\hyperref[sec:522]{< < < PREV SERMON < < <}
~
\hyperref[sec:toc]{[返目錄]}
~
\hyperref[sec:524]{> > > NEXT SERMON > > >}
\newline
\newline
今晚主席題及關於我所講的信息,特別關於明晚所要講的,是最後一次講及關於聖靈的信息──如何引致聖靈擔憂.這信息可應用於從未被聖靈充滿的弟兄姊妹,也可應用於曾經被聖靈充滿的弟兄姊妹.雖然我曾被聖靈充滿,但仍可能會令聖靈擔憂.下週一將會傳講福音的信息,我在此僅一次向未得救的朋友們,傳講福音的信息,請多帶未信者來參加；我是特別針對曾聽過福音,而仍未接受主者來講的.我會把這題目當個問題來講,夜半之後有何事發生,是關於半夜呼聲的信息.週二將講普世差傳的信息.我看見這班廣大的會眾非常特別,具有浩大的潛力,向普世作差傳的工作.在香港的中國基督徒,可在未來的廿年中把福音傳遍整個世界.該晚我也要講基督徒對冷淡信徒應負責任的信息,並講及我對非洲飢餓群眾的一點經驗.成千上萬的人整天挨餓直到睡覺的時候,翌日早晨很多人從此不醒來.當第二次世界大戰時,日本廣島在原子彈之下,隔天就會死去很多的人.週三晚是培靈會最後一晚講道的時間,將講到基督化家庭的四道牆.
\newline
\newline
這幾晚我曾到樓下書攤參觀書籍,傳記方面的書有助於基督徒的生活.神所使用那些偉人的傳記,由今天開始多了一部家父的傳記,中文書名叫《熱血滔滔》,在數小時前剛由海外運到港.我相信收集屬靈方面的傳記,可為各位帶來屬靈生命的祝福.
\newline
\newline
聖靈在信徒生命中的工作:
\newline
\newline
A,認識自己
\newline
\newline
(一)聖靈要信徒確實知道信徒是神家的一份子
\newline
\newline
有關經文:(羅八16,加四6,約壹三24,約壹四13)這幾處經文,論及聖靈在我們心中作任何工作,聖靈保證我們是屬神的一份子,聖靈保證我們有權利呼叫「阿爸父」,聖靈保證住在我們裡面,在我們心中讓我們確信是屬乎神的.這個確信不是根據我們心思所知道的,不是根據科學方面的引證,乃是聖靈給我們直覺的感覺.有時機師可能迷失方向,或因機件工作不對,或因天氣影響視力不清楚以致迷失.一個有經驗的機師在此情形之下,內心立即產生直覺轉移正確的方向,並非儀器也非眼睛指示,乃是感情上的直覺告訴他正確的途徑.有經驗的商人談生意,也有同樣的感覺,直覺會告訴他是否好的交易；並非科學方面的引證,也非眼睛所能觀察得到的,而是憑著經驗產生了直覺.當一個女孩子謀職,事先必搜集很多有關僱主的資料,她自知本身的資歷,但直覺使她認定是否適合.一般做母親的都有同樣的感覺,就是當自己的嬰孩發生問題時,她會知道的,她裡面有種感覺知道自己嬰孩的需要.照樣,在聖靈工作時,當聖靈充滿你,充滿我的時候,也讓我們知道我們乃屬神家中的一份子.祂沒證明我們是屬於神的家,沒有畫圖指明我們是屬於神的家；但祂使我們心裡深知自己是屬於神家的一份子.
\newline
\newline
(二)聖靈如何使我們確實知道我們是屬乎神家呢?
\newline
\newline
第一,否定我們的肉體.(羅八13)有的基督徒曾被神的聖靈充滿,肉體裡面給他提議,在他一生之中當服從.當肉體有此要求,通常他都會照做；但是當他被聖靈充滿時,他也和往日一樣,會有肉體的慾望；但他卻另有一個奇異的感覺,是不受肉體的操縱的,不再服從肉體的慾望；於是自知確實屬乎神的家了.聖靈使我們知道我們乃神家的一份子,不再是罪的奴僕了.
\newline
\newline
基督徒有時難免犯罪,我不問各位以往或現在有沒有罪在你心中；但是有的罪是我們經常犯的,我們也曾經想不再犯了；甚至求神讓我們不再作此犯罪的事,然而我們卻繼續不斷地犯這罪.際此培靈會,可能我們初次領悟聖靈充滿的感覺；或者像昨天晚上幾百位基督徒所感覺的一樣,承認自己的罪,讓神的聖靈充滿.也許經過了這三個晚上的聚會,你覺得沒有犯過甚麼罪；但確實感到有些奇妙的事,在身上發生,神曾經為我作了件事,聖靈確實充滿了我,因祂讓我戰勝一切的罪行.
\newline
\newline
(三)聖靈怎樣讓我確實屬乎神呢!
\newline
\newline
祂使我們有能力遵守神的誡命.「遵守神命令的,就住在神裡面......」(約壹三24)我們早就知道神的誡命是甚麼,我們好久之前,已知道神的十條誡命,也知道耶穌在山上的教訓,神在祂話語中清楚告訴我們.如果你相信聖經的話,你就知道神要你作甚麼；但是我們要遵守神的誡命常有困難.明知不可惱怒,不可嫉妒,不可憎惡別人,但卻一一違犯；明明知道神吩咐不可作的事,但卻不能遵守神的誡命.當我們被聖靈充滿時,生活有奇妙的事發生.我過基督徒生活多年,卻不能遵守神的誡命；但當被聖靈充滿之後五,六天,就有些奇妙的事要發生.在我的生命中第一次遵行神的誡命,第一次作神要我作的事,感到聖靈確實改變我的一生.神確實為我作了件事,我不但走在通道上,並和其他人站在一起,我禱告,神垂聽；我求神潔淨我的罪,祂潔淨了我.我求神的聖靈充滿我,祂就讓聖靈充滿我；當聖靈住在我們裡面,我們就開始遵行神的誡命.
\newline
\newline
(四)聖靈讓我們確實知道,我們是屬乎神,就幫助我們彼此相愛.
\newline
\newline
「我們若彼此相愛,神就住在我們裡面.」(約壹四12)有一個人求聖靈充滿,神聽了他的禱告.大概經過了兩週,奇異的事發生了.這是他有生以來,第一次對別人顯出愛心；有些人是他一向不歡喜的,但是現在他卻愛他們了,因為神改變了他的生命.不過,這並非唯一的證明,他是屬於神的家,如果我們要知道是否確實屬於神,又被聖靈充滿,這裡有張清單；有四個問題要問:1.我們是否已治死了肉體?2.對罪惡是否已得勝?3.是否遵行神的誡命?4.是否愛神的子民?
\newline
\newline
B,教導信徒
\newline
\newline
聖靈在我們裡面的第二件工作是「教導」我們.耶穌在聖靈未降臨之前曾應許聖靈要降臨,祂說:「但保惠師,就是父因我的名所要差來的聖靈；祂要將一切的事,指教你們,並且要叫你們想起我對你們所說的一切話.」(約十四26)聖靈對使徒說了兩方面的事:
\newline
\newline
(一)告訴使徒關於耶穌基督整個的啟示
\newline
\newline
耶穌曾多方教導門徒,但有的還未發生的事,耶穌不能親自教導.耶穌對門徒說了這話,是在祂受死後復活之前,耶穌告訴門徒將有這些事要發生.直到聖靈降臨,聖靈才教導門徒這重要的事,聖靈提醒門徒,耶穌在地上對他們所說的一切話.正如四本福音的作者,並非依靠自己的記憶詳細記錄所發生的事；如果這樣記述一定會有所錯誤,因為沒有人能詳記兩週前一切的話.神差聖靈提醒門徒,想起耶穌所說的一切話,聖靈直接指示門徒怎樣寫,所以聖經毫無錯誤.也許有人要說聖經有錯誤,有的人說聖經有科學上的錯誤,有的人說聖經中有不吻合之處.但是,聖經原文卻不可能有錯誤,因為是聖靈親自負責全部聖經的寫作；聖靈是神,神絕無錯誤的事.耶穌基督的啟示,是特別給門徒的,耶穌說,聖靈要教導這班門徒.這不是「應許」,若是「應許」聖靈會教導我們的.耶穌說,聖靈特別要教導這班門徒想起祂所說的一切話.耶穌的預言是對要寫福音書的門徒說的,而聖靈要教導這班門徒想起這一切的話,惟有聖靈能使門徒作要他們作的事,祂並不教導我們多寫幾本聖經；因祂已教導門徒寫了聖經,但聖靈要教導我們明白聖經.世界任何人都可讀這本聖經,無論是其他宗教的人,無神論者,各種各樣的人都可以讀聖經；但惟獨有聖靈的人才能明白聖經,聖靈使我們明白聖經的真理.
\newline
\newline
多年前我在加拿大的溫哥華住了一年之久.溫哥華是加拿大最美的城市之一,一邊是太平洋,另一邊是洛磯山脈,實在是極美的城市.但常常下雨,有霧,我在那裡一年,只有兩天好的天氣.當我廿多歲時,第一次搭火車赴溫哥華,參加英國哥倫比亞大學神學院,當時正下雨落霧；有人用車送我到達,並告訴我入住的房號.我一放下行李便開窗外望,可是甚麼都看不見,大霧矇矓,連街道都籠罩其中.到了次日清晨,我再走到窗前,大霧已消散,太陽顯出光線,那是我有生以來所見最美麗的風景；眼前浮現一片平靜的河水,船隻駱驛不絕.另一面是山脈,下面大樹林立,樹之高大乃世界罕見；有的周圍直徑達十尺之大,樹身聳入雲霄,樹梢上端有大堆積雪,其後面襯著蔚藍的天空,真是奇妙的風景!濃霧遮蔽時,任何東西都看不見.許多人好像罩著霧讀聖經,因而發生問題.我曾向他們解釋,但他們總不肯接受,惟聖靈進入,霧才能消散；然後才能初次看見美麗的山脈,初次看見碧海青天,初次看見神描寫整個的世界.
\newline
\newline
讓我再說,聖靈在信徒生命上給我們確實知道我們是屬於神的家.祂幫助我們遵守神的命令,幫助我們治死肉體的缺陷；讓我們彼此相愛,並在信徒生命中教導使徒們寫新約聖經；教導他們避免犯任何錯誤,教導我們明白聖經的真理.沒有聖靈,聖經充滿著霧,看不到聖經中美妙的山脈,海洋和聖經中神極奇妙的圖畫.
\newline
\newline
當聖靈充滿我們時,祂就會教導我們.
\newline
\newline


\newpage

\section{第55屆~港九培靈會~史保羅~第七講　如何引致聖靈擔憂}
\label{sec:524}
\textbf{第七講　如何引致聖靈擔憂}
\newline
\newline
連結: \href{https://www.hkbibleconference.org/session-message/view/524}{\texttt{https://www.hkbibleconference.org/session-message/view/524}}
\newline
\newline
\hyperref[sec:523]{< < < PREV SERMON < < <}
~
\hyperref[sec:toc]{[返目錄]}
~
\hyperref[sec:525]{> > > NEXT SERMON > > >}
\newline
\newline
今天晚上要講關於聖靈最後的信息,在整個信息中我將不斷提及(弗四15-32)這段經文.「不要叫神的聖靈擔憂,你們原是受了祂的印記,等候得贖的日子來到.」(31)這節經文和(弗五18)的經文是相對的.保羅說:「......要被聖靈充滿.」如果叫聖靈擔憂,同時要被聖靈充滿,是絕對不可能的事.聖經中有多處提及如何引致聖靈擔憂.
\newline
\newline
(一)謊言引致聖靈擔憂「所以你們要棄絕謊言,各人與鄰舍說實話；因為我們是互為肢體.」(25)
\newline
\newline
謊言有兩方面,一是清楚的謊言.我將到63歲了；如果有人問我幾歲,我說45歲,這就是清楚的說謊,人家都知道的.還有其他方法的說謊,當我們聽了一些未經證實的謠言；不明是非,就當作事實告訴別人,這就是說謊.聖經說,謊言令聖靈擔憂.
\newline
\newline
(二)生氣惱怒引致聖靈擔憂「生氣卻不要犯罪.不可含怒到日落.......」(26)
\newline
\newline
好的發怒是不犯罪的.大衛說:「求你將你的惱恨,倒在他們身上,叫你的烈怒,追上他們.」(詩六九24)聖經告訴我們,神有時是惱怒的.當摩西上山領十誡,回來時看見以色列人拜金牛犢,摩西「......便發烈怒,把兩塊版扔在山下摔碎了.」(出卅二19)神並不為此指摘摩西.在新約,耶穌也曾發烈怒,當耶穌進會堂醫治一個手枯乾的人；但法利賽人認為在安息日治病是不對的,就議論紛紛.「耶穌怒目周圍看他們,憂愁他們的心剛硬,就對那人說,伸出手來,他把手一伸,手就復了原.」(可三5)這就是所謂「生氣卻不要犯罪.」神,摩西,耶穌,發怒卻不犯罪.但聖經說有些人發怒是犯罪的.我們如何保證不發怒呢?我給大家幾個建議:我們要保證,惱怒不是針對任何人,不是自衛的惱怒.有時受人冷落,即刻還以拳頭,這就是出於自衛.我們若因人傷害了神的工作,而發怒是正確的；摩西發怒乃因以色列人犯了淫亂的罪,傷害了神.大衛求神發怒,乃因那班人妨礙神要叫大衛作的事；神叫以色列人建立國度於應許之地,當那班人防止大衛的工作；大衛就禱告求神的惱恨,倒在他們身上.耶穌基督惱怒,乃因有些人阻礙神旨意的工作；實際上,他們計劃要殺害耶穌,故意借題發揮.第二個保證生氣卻不犯罪的方法,就是解決了那件惱怒的事然後睡覺.保羅講及基督徒發怒的事就對他們說:「睡覺之前要解決惱怒的事.」當日落之後,人多數在家裡,常常對著家人發怒,外人不得而知；在外面,特別到禮拜堂,都是衣冠整齊,面露笑容,一表斯文,簡直像個聖人.保羅是個極聰明的人,他知道基督徒在禮拜堂唱詩時當然不會發怒,聽道時膝上放著聖經當然也不會發怒.所以保羅說,日落之後睡覺之前,要先解決惱怒的事.一個患了重病的人,臨終之前一定要把一生心中大事完全解決,欠人債的希望能夠還清債務；平時與人吵鬧不睦的,希望與人和好；遺下兒女的,希望能繼續供應他們；都希望臨死之前,辦妥一切事務.保羅的意思是說,我們上床睡覺之前,要作好一切的事,當作明天早晨長眠不醒了；雖然明晨我們還會醒來起床,但是保羅叫我們不要依靠那個,好讓有日時光我們仍然可以走路.
\newline
\newline
保羅說:「不可給魔鬼留地步.」(27)以致毀壞我們的見證.我想,踢足球是全世界最流行的運動,比賽時如果對方是世界聞名的球隊,必須全神貫注,特別守望；使他們沒有得球的機會,使他們無法四圍活動.照樣,對付魔鬼也是如此,魔鬼是我們對方的隊長,是全世界最有名的運動員；我們必須留心看守,時刻提防,不給機會牠任意活動,以致擊敗自己的隊伍.如果我們惱怒,就給魔鬼留地步,給牠更有機會毀滅我們.
\newline
\newline
(三)偷竊引致聖靈擔憂
\newline
\newline
偷竊最少有三種,打劫銀行是人所共知的盜賊,但打劫可能與我們無關；但是也有其他的方面,會引致我們犯上偷竊的；有時買東西找回的錢數過多,雖知道了,仍隨手放入口袋裡；這和搶劫銀行的人,一樣犯了偷竊的罪.如果我們搭巴士沒給錢,就是偷竊了巴士公司的錢,或許我們從一個國家的邊境,進入另一個家；理應經過海關檢查完稅,結果我們逃了稅,這也是犯了偷竊的罪；當加拿大人到香港,都買套西裝,因為香港的裁縫世界馳名；回到加拿大,就將新的穿在身上,把舊的放在行李箱；以此聰明方法避免海關打稅,海關人員問「有沒東西報稅?」一口清白的回答:「沒有.」這樣做就犯了說謊和偷竊兩罪,偷竊了加拿大政府的錢.有的人從僱主方面偷竊,如果上班遲到或早退,或午飯時過了鐘點,都是偷竊了僱主的時間.當我們找到一份工作,和僱主簽了合約,早晨九點開始工作,中午一小時吃中飯,下午四點半下班；如果沒有履行合約,就是破壞合約,就是偷竊僱主的時間.有時僱主讓僱員加班而沒多給工錢,就是偷竊僱員的時間.僱主和僱員各持有一份合約,沒照合約履行,則偷竊時間或金錢.或者有時從製造物件中犯了偷竊,本來每小時可以做卅件,但實際上只做了廿件.也許有人要說:「史密司先生,遲到,早退,過時是偶然有的事,何必這麼認真呢,在商業社會裡一般都是如此的.」如果我們是基督徒,就不能夠跟著一般的作風.
\newline
\newline
(四)污穢的言語引致聖靈擔憂「污穢的言語一句不可出口......」(29)
\newline
\newline
污穢的言語有數種,有時講粗口,甚至講出神的名,有時講污穢的故事；還有個方法講污穢的言語,聖經用「嚷鬧」二字(31)意思即和人爭論,強詞奪理,大聲喧嚷,以壓倒對方使其無法爭辯.聖經說這也是基督徒不該作的.還有一種污穢的言語「毀謗」,講人家的壞話,就是毀謗；若所言者非屬事實,那就更加嚴重,一言既出駟馬難追；如果偷竊,尚可有機會償還,但出口的言語,卻永遠收不回來；當然求神饒恕,神還是饒恕；不過,我們講話應當極其小心才對.
\newline
\newline
(五)惡劣的脾氣引致聖靈擔憂「一切苦毒,惱恨,忿怒,嚷鬧,毀謗,並一切的惡毒,都當從你們中間除掉.」(31)
\newline
\newline
對付惡劣脾氣有何良葯呢?我們若想打人一拳,或踢人一腳,應如何防止呢?最後的一節經文告訴我們:「並要以恩慈相待,存憐憫的心,彼此饒恕,正如神在基督裡饒恕了你們一樣.」如果我們想防止發怒,聖經說:「要以恩慈相待.」有了愛就不可能發怒.　　我們若不是被聖靈充滿,就是引致聖靈擔憂.當我們引致聖靈擔憂時,我們便喪失能力,我們的見證就不如往日之有效；不像以前被聖靈充滿了,就不再如從前知道自己乃屬乎神的家,得到聖靈的教導也就少了.
\newline
\newline


\newpage

\section{第55屆~港九培靈會~史保羅~第八講　夜半呼聲}
\label{sec:525}
\textbf{第八講　夜半呼聲}
\newline
\newline
連結: \href{https://www.hkbibleconference.org/session-message/view/525}{\texttt{https://www.hkbibleconference.org/session-message/view/525}}
\newline
\newline
\hyperref[sec:524]{< < < PREV SERMON < < <}
~
\hyperref[sec:toc]{[返目錄]}
~
\hyperref[sec:526]{> > > NEXT SERMON > > >}
\newline
\newline
現在已經到了本屆培靈會的末段時間了.最後的兩個晚上,我要將重要的信息帶給大家.明晚要講對全世界傳福音.多倫多的民眾教會支持了很多的宣教士,超過五百位宣教士分佈世界各地工作,每年所需費用達百萬以上的加拿大幣；至於本地教會所需經費,僅約75萬元.所以每年向世界差傳事工的費用,超過本地教會的費用.
\newline
\newline
如果有人只想參加一晚的聚會,那就千萬別錯過明晚的聚會,我已重複講了幾次有關於昨晚的信息.當我看見這許多青年赴會,我心裡非常感動；在這裡的青年,今晚可看見神極大的能力；相信將來的廿年中,中國有足夠的宣教士傳福音,使世界各地的人歸向主.多年來都是西方宣教士來中國傳福音,而現在是中國信徒向世界各處,未聞福音的地帶傳福音的時候了.明晚的信息,乃是關於我們對未曾聽到福音地方的人的責任.
\newline
\newline
多謝陳牧師介紹樓下書攤的書籍.亞洲歸主協會,發行了我寫的兩本書；有本家父的傳記是中華傳道會最新發行的書.這個機構初期名為華人自傳會,是中國人在中國開始的；現在工作遍佈全球,由中國宣教士和各國本地的基督徒一起工作.家父傳記,中文譯名為《熱血滔滔》書中有許多趣味的照片,其中有張我小時漂亮的像,頭髮厚而鬈曲,母親替我中間分界,兩邊披著長髮,走起路來隨風飄揚,非常美麗.有一日,母親買了很多東西提在手上,我還幫她提著手袋；在路邊玩的小孩子,特地抬起頭來問我媽媽:「這小孩子是男的還是女的?」當時使我很難為情.書中記載家父一生的事,他名為柯斯華史密司,生長在加拿大的鄉村,該地僅一間禮拜堂；他孩童時代極瘦,但他活到94歲.我相信你讀這傳記,一定帶來祝福的.
\newline
\newline
馬太福音廿五章1-13節這段經文提及半夜的事.聖經記載「夜半」共十四次,其中有一處的「夜半」是用普通的意思而言的,其他講及神的救恩或神的審判.
\newline
\newline
第一次提及夜半是猶太人過逾越節的時候,神差天使去殺埃及人的長子,使埃及人嚐到神審判的大能；同晚天使要越過以色列人的住宅,以色列人的孩童嚐到神救恩的大能；在那段經文中,有救恩也有審判,二者同在夜半發生.另一次夜半在使徒行傳記載,使徒保羅帶領獄卒歸主,當時保羅和西拉在監獄中,他們決定要舉行一次聚會.那次的聚會,我們不知道他們作些甚麼；大概是唱詩,禱告,作見證；那時獄卒正睡覺,忽然發生地震,所有的囚犯鎖鍊都鬆開.腓立比的獄卒就進到裡面,在保羅面前跪下俯伏,問道:「我應當作甚麼才能夠得救呢?」那個晚上,獄卒全家都得救.那次的地震是在夜半發生的,夜半時腓立比獄卒,嚐到神的救恩.
\newline
\newline
使徒行傳廿章記載,保羅正在一個家庭聚會,許多人聽他講道；有個青年人名猶推古,因找不到地方坐,就坐在窗台上,由窗台下望聽保羅講道.一直聽到夜半時候,青年人睡著了,由窗台跌下,仆在地上死了；於是保羅停止聚會,出去看看青年人的身體,奉主的名使他復活過來.這事發生在夜半的時候,神的大能把生命帶給一個人.
\newline
\newline
在字典中,夜半的意思是一日的終結,也是另一日的開始；是一段廿四小時的結束,也是一段廿四小時的開始；但是在我們一生中有許多段的夜半,第一段是嬰孩時代的結束,入學時代的開始.
\newline
\newline
我的第二個女兒,可說是個非常完美的孩童,她好像一般肥胖的嬰孩,手腳積肉節節,肥圓的肚子.「珍」是我們的嬰孩,我們愛她,照顧她,保護她,我們夫婦倆永不忘記,當她五歲時,我們站在前門處,她站在我們中間.她有生以來第一次上學,她對爸爸說再見,並且吻她媽媽說聲再見；然後走向學校,當我們關上門,頗有夜半之感.我們知道她再不是個嬰孩了；嬰孩時期已結束,上學時期已開始.
\newline
\newline
第二個夜半是上學的末了,畢業典禮對年青人是重要大事,無論中學,大學的畢業典禮；走上講台前領取證書學位,那時有種思想透過心中.當你第一次入學,你想到一切要學習的科目,想到你所要得的分數,想到你在學校所要做的；到了畢業時,你才知道那樣你讀過,那樣未讀過,做到或未做到,得到或未得到；夜半的鐘已經打在你心上,你能夠領受正式教育的成果,繼而開始你的工作了.
\newline
\newline
年老的男女也有一種夜半的感覺,那是發生在退休的時候.當工作的公司開會議決,給你退休時,不能不想到一生希望要做的事；年青時一切的期望及所渴望的,到了退休之日,有的希望已經達到；但是有的還未做到,夢想有的已經實現,有的永遠不能成就；工作年代的機會已經過去了,如今所走的是生命末了的階段.極其重要之夜半鐘聲已經響了,表示工作時間的終結,退休年代的開始.
\newline
\newline
今天晚上不想費時來講每段時間的終結和開始,不想詳說嬰孩時期的末了,和上學時間的開始,也不想多講工作年代的最後,和退休年代的開始.我們要講關於神「夜半」的事,就是馬太福音提及神的夜半時間；夜半乃是我們成為基督徒機會的終結.關於我們得永生與否,表示我們能否成為神家一份子的機會的終結,表示我們在神家之外之永遠,是我們受耶穌基督作救主的機會的終結；這夜半是我們結束,永遠沒有接受基督的時間了.在你我的生命中,夜半的每段時間終會發生的.
\newline
\newline
今晚我要問各位幾個問題:有的人等到夜半之後,有何事發生在他們身上呢?到了夜半的鐘聲響了,如果我們還在神家之外,如果我們還在神救恩之外,如果我們仍然在神團契生活之外,甚麼事會發生呢?
\newline
\newline
當我們年幼時,每逢外出,爸爸規定時間要我們必須回家；到了年紀漸長,可以較遲,不過,必須夜半之前回家,這是我家簡單的規則.夜半之前回家,大門沒鎖也無人過問；如果夜半之後,大門深鎖,只好叩門；弄得家人都被吵醒.一進門就得向各人解釋,遲回家的原因；不知為何,我是常常過了限定的時間數分鐘；門已上鎖了,連後門也鎖上了；我就想辦法要由窗門爬進去,可是窗門都關緊了,我只好回到前面敲門.我想爸爸一定故意使我久候門口,讓我知道時已過了夜半；我由窗內望,看見爸爸走來開門；我一進門就急忙解釋為何過了夜半才回來.事實上並無確定的限期,無論我多麼遲回家,我總可以進去的.不過,我告訴各位,這門和神的門是不同的,到了半夜,神的門就關了；那門永遠不再開的,無論你讀這段經文多少次,那門永不再開的；當神把機會的門關了,就永遠關閉.夜半的鐘聲響了,我們還不是基督徒,不屬神家的一份子,不在往天堂的路上,那就永無機會了.
\newline
\newline
在我們的生命中,有甚麼事會在夜半發生呢?我們怎知何時是夜半呢?有三件事使我們夜半的鐘聲在我們的身心響起:
\newline
\newline
第一是死亡,當我們停止呼吸,就是夜半了,肉體的生命一旦離開,就是夜半了.我們死了,接受耶穌基督作救主的機會就過去了；我們去天堂的機會就過去了,成為神家一份子的機會就過去了!
\newline
\newline
第二,夜半的鐘聲再響,主即再來.聖經清楚告訴我們,有一日耶穌基督再降臨；那時是所有的人的夜半時間,與神和好的機會就完結了.或許有人說,我還年青,活在世上還有很久的時間,離開死亡尚遠；或有人以為耶穌基督不會很快再來的.耶穌基督再臨,成為我們最後的夜半鐘聲.
\newline
\newline
第三,當我們初次聽了神的福音,聖靈感動我們,在我們裡面有種感覺,需要接受耶穌作救主；當時容易接受,因為在那情形之下很難拒絕；但為了某種理由,終而不接受.到了第二次聽道,仍然樂於接受；但基於某種原因,仍想拒絕；不過就不像第一次那麼難以拒絕了.如此情緒,在生命中一直持續下去,一次次聽了福音而拒絕接受,那就越難接受福音；因為習慣了向神拒絕,到了最後,很難使你接受歸向主.你的心愈來愈剛硬,習慣拒絕成了自然；你雖參加聚會,很歡喜唱詩,聽講道非常欣賞,但卻無動於衷,無關緊要；如果這樣,夜半時候已經來臨了,你雖繼續生活很久,終於拒絕接受耶穌作救主.或者你說:「我是這樣的,我的心太剛硬了,我拒絕耶穌太多了；甚至今晚我都不能接受耶穌,或者我生命的夜半已經到了.」然而,如果你現在回轉懇求耶穌拯救你,當機立斷,切勿延緩!你還可以得救.聖經有話說:「凡接待祂的,就是信祂名的人,祂就賜他們權柄作神的兒女.」(約一12)如果你決志求神拯救你,那麼,你的夜半時間還未來到；因為聖經說:「凡求告主名的,就必得救.」(羅十13)
\newline
\newline
「看哪!現在正是悅納的時候,現在正是拯救的日子.」(林後六2)「不要為明日自誇,因為一日要生何事,你尚且不知道.」「其實明天如何,你們還不知道,你們的生命是甚麼呢?你們原是一片雲霧,出現少時就不見了.」(雅四14)
\newline
\newline


\newpage

\section{第55屆~港九培靈會~史保羅~第九講　對冷淡信徒應負的責任}
\label{sec:526}
\textbf{第九講　對冷淡信徒應負的責任}
\newline
\newline
連結: \href{https://www.hkbibleconference.org/session-message/view/526}{\texttt{https://www.hkbibleconference.org/session-message/view/526}}
\newline
\newline
\hyperref[sec:525]{< < < PREV SERMON < < <}
~
\hyperref[sec:toc]{[返目錄]}
~
\hyperref[sec:527]{> > > NEXT SERMON > > >}
\newline
\newline
以西結書 3:17-3:19
\newline
\begin{longtable}{cl}
\hline
\hline
章節 & 經文 (和合本修訂版)\\
\hline
3:17 & \begin{tabularx}{0.7\textwidth}{X} 「人子啊，我立你作以色列家的守望者，所以你要聽我口中的話，替我警戒他們。 \end{tabularx} \\ \\ \relax
3:18 & \begin{tabularx}{0.7\textwidth}{X} 我何時指著惡人說：『他必要死』；你若不警戒他，也不勸告他，使他離開惡行，拯救他的性命，這惡人必死在罪孽之中；我卻要從你手裡討他的血債。 \end{tabularx} \\ \\ \relax
3:19 & \begin{tabularx}{0.7\textwidth}{X} 倘若你警戒惡人，他仍不轉離罪惡，也不離開惡行，他必死在罪孽之中，你卻救了自己的命。 \end{tabularx} \\ \\
[1ex]
\hline
\hline
\end{longtable}
請看(結三17-19)留意第18節末了「喪命之罪」原文作血,請大家注意!
\newline
\newline
我小時候最喜歡爸爸的書房.媽媽告訴我,當我學會爬行時,常常爬進書房.我所以歡喜爸爸的書房原因有二:第一,房中牆壁掛著許多美麗的油畫.畫面只有海和羊兩種,有的是羊和牧羊的人.我很喜歡坐在書房欣賞那些畫.
\newline
\newline
我父親常到世界落後的地方,有時離家半年之久,因那時遠行都是乘輪船,所以往返需時.父親回來時,常常帶了一些外地的東西,放置在書房裡,有次他帶了一把土人的矛回來,有次帶了一支用口吹的毒鏢回來,有次帶了一件大蛇皮回來,有次到印度帶回來一件老虎皮,把它放在地氈上,老虎的頭仰在地上.我小時候常進書房坐在老虎頭上,觀看房中所有的東西；但是有幅畫我不明其意,畫仍裝有框架,畫的背景全是黑色,黑色的後面是白色線條,我只知道那不是海,也不是羊和牧人.我覺得毫無意思,後來到了我上學讀書,開始學習；我仍常進書房,欣賞那些東西；原來我看不明白那幅畫,白的線條是些字母,拼成了字,造成句子,加上問號,我每天看見那些問句.到了我長大了,在教會也常見同樣的問話；很多時候也聽到父親講,有關那些問話的信息.那些問句是:為甚麼有的人能夠聽到福音只一次,有的人能夠聽到福音兩次以上?
\newline
\newline
耶穌餵飽五千人的時候,(聖經有多卷題及此事)聖經說五千個男人.我想其他一定還有女人和小孩,或者約有萬人吧.耶穌知道這班人生活厭倦單調,他們住在小的國家,與整個世界的局勢無關；當時巴勒斯坦沒甚麼特別的事情發生,住在鄉村的人很少有趣的事.耶穌說,這班人不是來聽我講道,乃是來看我行神蹟；他們對我的教訓沒有興趣,他們的生活一向單調,若看見神蹟必定非常興奮.第二,耶穌說:「他們困苦流離,好像群羊迷失沒有牧人一樣.」他們沒有領袖領導,他們的人生漫無目的.第三,耶穌知道他們肚子餓,因離家來到曠野已經整天了.耶穌對門徒說:「我想我們應該餵飽他們.」門徒說:「主人啊!這是很愚蠢的事,要花多少錢買食物才能餵飽他們呢?附近又沒有商店；最好還是叫他們各自回家吧!」耶穌說:「我要給他們食物.」耶穌問門徒有沒有任何食物,門徒說:「有份小孩子的午餐在我手中.」
\newline
\newline
我實在很想知道,門徒如何從小孩子拿到那份午餐,聖經沒有記載.有時我想,門徒是否靜悄悄從小孩子背後拿過來,或者用說服的方法,或者小孩子自動交出來的.這份午餐包括五個餅和兩條小魚；對這麼多人實在無濟於事.一份孩子的午餐要餵飽五千個大人哪裡可能呢?你有沒曾經對神或對自己這樣講過?或許大家都聽過全世界對福音方面的需要,有多少的人從未聽過福音；有的國家我們非常難以進入,你有否看見自己手中可奉獻給神的是甚麼?或者你小小的祈禱生活；或者你想,晝夜禱告有何用呢?世界上有許多人等待拯救,有多少財政方面的奉獻,可以作福音之用,少許的金錢不足以供普世傳福音的工作.
\newline
\newline
我曾講過,多倫多本教會每年供給百萬以上的加幣,作全球宣教工作.或者大家心目中認為這是筆大數目；但是照我的想法,全世界有許許多多的人等待福音,而願意出去工作的工人很少,且所需要的費用很大；百萬加幣正如一份小孩子的午餐,要分給五千人吃飽,實在微不足道!一點點的數目,對於浩大的需要何用之有!當我們民眾教會收宣教奉獻時,我教導主日學教師,怎樣灌輸宣教工作給小孩子們,告訴他們所收的奉獻如何用在工作上；讓每個孩子自己決定可以奉獻多少.有的每月承諾獻一毛錢,有的兩毛半錢,有的每月一元；同時,以門徒同樣的問話說:「主啊!從一個孩童那裡拿了一元二毛錢有何用呢?從一個孩童拿了三元有何用呢?一年收到十二元,能夠傳福音給幾個人呢?」回答的話是,「用這些錢可以傳福音給任何人.」
\newline
\newline
大家有沒有自問,自己所有的究竟對個人多有用呢?如果今晚我問各位,請你做個宣教士,離開家庭到落後的國家去傳福音；或者你需要學習新的語言,或者生活沒有現在過的舒服；或者你所在地方的文化完全不同,這就是宣教士的生活了.一個青年宣教士說:「我願意離開我的家庭,離開大好的前程,到神要我去的任何地方.」如果今晚我問各位:「誰願做個宣教士?」多少人要回答說:「全世界人口太多,未聽過福音的人實在太多,我這小小的生命對於偌大的需要有何用呢?」這樣的回答正像門徒對主的回答一樣,他們在主前拿起小孩子的午餐說:「這些在廣大會眾面前有何用呢?」你的奉獻,你的生命,不能夠把福音帶給世界,禱告也不夠用；我們的生命,禱告,奉獻都像那小孩的午餐,確實無濟於事.除非我們把一切所有的放在主耶穌手中；我們禱告是為耶穌,奉獻是為耶穌,當宣教士也是為耶穌；那麼,耶穌將加上祝福,使小小的午餐達到整個世界.
\newline
\newline
看!你的潛力多大啊!這裡在你的周圍,還有幾千人,請各位忘掉其他的人,我要使你認識你的生命,禱告,奉獻,能夠接觸到整個世界.
\newline
\newline
門徒怎樣把餅和魚分給會眾,請大家注意我所要說的故事對不對.門徒把餅和魚拿到左邊的第一排,按次序傳遞各人領取,在第一排來回繼續分派.
\newline
\newline
耶穌基督的教會,兩千年來把福音從有文化的地方傳出去,也以百分之九十的傳福音使者,對百分之五的人工作.本教會的工作是向著第一排的人工作,牧師,主任與我,都同樣在教會工作；每主日向同一批會眾講道,盡教會的能力適應會眾的需要.我們派發單張,探訪,也在電視中工作,派巴士出去帶未信主的人到教會來.本教會的工作至今廿五年了,我也作了廿五年了,這是神呼召我作的事；但是我也有機會到世界各地,對後面的人工作；我曾去過未聞福音的地方,因那裡的人沒有機會接受耶穌基督.
\newline
\newline
兩年前我到非洲撒哈拉沙漠的南面,那裡的部落非常缺乏食物和水,每四至五年便有一次大旱災,因此很多居民逐漸餓死,最先死的是牛,因為缺乏水和食料；每廿四小時有一萬七千人餓死.有日我經過一處,看見一個年青婦人在沙灘上側臥,用手托著頭,因無食物面臨死亡,她瘦削如柴,甚至可在十碼距離看見她胸前有幾條肋骨；她的嬰孩坐在她面前飲泣,我想他一定是無力大哭,看來快要餓死了,他薄皮包骨,非常的瘦；無論由前面或後背都可看清楚她的心臟跳動.我俯身看看婦人和嬰孩,這時他停止飲泣,似乎要喝母親的奶,可是無奶可喝,又開始飲泣；過了片刻,他又嘗試俯下喝奶,仍然沒奶喝；一次又一次地嘗試都喝不到奶.這個年青的媽媽本身將飢餓至死,情境實在悽慘可怕；但更加悲慘的乃是她親眼看見自己的嬰孩逐漸餓斃.那天我們試要幫助這對母嬰,但是已無可救藥了.
\newline
\newline
我離開了那地方搭機返加,抵家已是週六下午,在機上約卅六小時.第二天早晨,我們打開衣櫃,選擇要穿上到禮拜堂的衣服.我想到我們生活在這世界上,很多人僅有一件衣服,在加爾各答,很多人睡在街上；寒冷地帶許多人沒足夠的衣著而挨凍.禮拜日早晨,我們選擇了衣服穿上,就下樓預備吃早餐,我們吃喝可以隨意選擇.我想到非洲每日一萬七千人餓死,我們實在吃不下任何東西,只喝了半杯黑咖啡,然後就上禮拜堂.整個上午我坐在講台上,望著親愛的弟兄姊妹們,我相信他們所過的生活和牧師一樣.我不能不想到我們,為主耶穌基督的工作,甚至一毛不拔,從未作過任何犧牲的事,從未為主受過苦.有時我思想到一個現象,看見我和弟兄姊妹在天家,圍著神的寶座而坐,有整隊的男女排列在遠處,其中有黑,黃,白,紅,各種的人,他們向著懸崖走去,到了邊緣跌下去,那裡永遠黑暗,永遠與神隔絕；其中有我在非洲看見的那年青的媽媽,她已餓死了,再沒機會聽到耶穌基督的福音；我還看見,當他們臨跌下去時轉身望著我,並聽見她說:「神啊!沒有人關心我的靈魂.」於是他們永遠消失了.我心裡被譴責,我知錯了!在異象中我問神說:「神啊!我豈是看守我兄弟的麼?」神回答說:「你兄弟的血,有聲音從地裡面出來向我哀告,你兄弟的血從非洲,從日本,南美,中國,各地出來向我哀告.」我自己卻在天上；但我手中有男女的血,我沒盡力把福音帶給他們,我要為未嘗聽過福音的人負責.
\newline
\newline
既在天上,手裡有血是甚麼意思呢?我不知道為甚麼,我也不想找出為甚麼；因為在很多年前我感到有責任把福音帶給整個世界.很多年前我知道教會有責任,把福音帶給整個世界,很多年前我把自己獻給神作此工作；在神面前我說,我作為教會的牧師,我要將餘下的時間奉獻,把福音信息帶給全世界所有的人.
\newline
\newline


\newpage

\section{第55屆~港九培靈會~史保羅~第十講　如何建立基督化的家庭}
\label{sec:527}
\textbf{第十講　如何建立基督化的家庭}
\newline
\newline
連結: \href{https://www.hkbibleconference.org/session-message/view/527}{\texttt{https://www.hkbibleconference.org/session-message/view/527}}
\newline
\newline
\hyperref[sec:526]{< < < PREV SERMON < < <}
~
\hyperref[sec:toc]{[返目錄]}
~
\hyperref[sec:510]{> > > NEXT SERMON > > >}
\newline
\newline
箴言 22:15-22:15
\newline
\begin{longtable}{cl}
\hline
\hline
章節 & 經文 (和合本修訂版)\\
\hline
22:15 & \begin{tabularx}{0.7\textwidth}{X} 愚昧迷住孩童的心，用管教的杖可以遠遠趕除。 \end{tabularx} \\ \\
[1ex]
\hline
\hline
\end{longtable}
一般的房子四面都有圍牆,使我聯想到基督徒的家庭也有四幅牆,當然我所指不是物質的牆,而是應當建立在基督徒家庭的牆.
\newline
\newline
(一)建立紀律的牆
\newline
\newline
世界最聰明的人所羅門,他是以色列第三位王.他向神求智慧,他寫箴言書,在(箴廿二15)說:「愚蒙迷住孩童的心,用管教的杖可以遠遠趕除.」他說,孩童應當學習如何服從父母,如果不服從,要用紀律管教.年長的人都接受這樣的說法.基督徒的家庭對孩童更要有紀律,但不可忘記,紀律不只應用於孩童,也應用於成人,紀律應用於為父牧者,如同應用於孩童一樣.
\newline
\newline
新約聖經把士兵和基督徒作比較,保羅說:「要穿戴神所賜的全副軍裝.......」聖詩中也有多首論及基督的精兵和全副的軍裝.
\newline
\newline
任何地方都以軍事紀律為重,第二次世界大戰時,我參加加拿大陸軍,管理我們的沙展非常嚴厲,從不和氣對人說話.向我們說話以厲聲呼喊,下令出發向目標攻擊；軍令如山,我們立刻執行任務,士兵對長官唯命是從.
\newline
\newline
基督徒在神的軍隊中,神一切命令都在聖經中.一個基督徒無論年紀多大,都在神命令之下,是神軍隊中之一分子,當服從神的命令.一個基督徒每天應讀聖經,有時不愛讀經,因為聽從了自己心裡的話.服從紀律才能夠每日讀經.無論作甚麼,不要為自己的愛好去做,乃要聽神的命令去做.我們在神的軍隊裡,要服從神的紀律,聽神的話.
\newline
\newline
(二)建立好榜樣的牆
\newline
\newline
心理學家告訴我們,凡學習都是照著別人的榜樣.很多時,我們學習的心得不是從書本中來的,而是模仿別人得來的.模範就是我們學習的模型.人生最初五年,每日所見最多的是自己的父母,所以就模仿父母；學習父母的語言,模仿父母走路；如果父母以爬行代步,嬰孩也照樣模仿；有的嬰孩腳骨有點彎曲,頭部又重,學習行走較為困難,常會跌跤的,經過時常練習也就行動自如了.嬰孩常見媽媽餵奶時給以微笑,所以也就模仿微笑；可是常見孩童微笑時好像縐著眉頭；因為還未學習得好,漸漸地也就學得微笑自然了；如果媽媽餵奶時,都是縐眉頭的話,那麼嬰孩就學不到微笑了.五年時間孩童的學習,為父母都實在負著重大的責任.我的兒子孩童時代,他學我剃鬍,其實他那時哪兒有鬍鬚可剃呢?兩個女兒在五歲之前,常學習媽媽的模樣做事.要煮飯,希望學會了媽媽所做的；也希望像媽媽一樣高大,一樣漂亮.總之,兒女整天注意父母的舉動,隨時模仿；這樣,到底給你快樂或憂愁呢?你能否對你年幼的兒子說:「來吧!你跟隨我,我不是世上最強壯者,不是世上最聰明者,不是世界最富有者；但我是照著神的紀律去做；神要我做的,我做；如果你跟隨我,你所走的是正確之路」.你能否對年幼的女兒說:「我不是世上最漂亮,最好,最聰明的女人；但你須知我是走在神的道路上,我根據聖經的話生活,如你跟隨我走,這路是正確的.」或許你要對兒女說:「不要跟著我走,因我的生活不正確,不是你應該模仿的人,你還是找個別人模仿吧!如你跟隨我的你就走錯路了.」身為父牧者,對兒女應立下好榜樣；廿年後你的兒女如何,有賴於你們現在的榜樣.未知你家中有沒有好榜樣的牆呢?
\newline
\newline
(三)建立家庭崇拜的牆
\newline
\newline
如果能到教會讀聖經,是件極好的事,能參加主日崇拜,聽講解聖經是很好的；但是這些都不能代替在家裡的讀經和禱告；應當每日規定時間帶領兒女,全家一同禱告.在座中很多青年人,尚未建立自己的家庭；我提議你建立家庭時必須這樣做,從你結婚之日,開始在家裡立祭壇,夫婦一齊禱告；等有了兒女,全家一齊禱告,讓兒女們將來知道你們曾經為他們禱告.
\newline
\newline
大家都熟悉聖經中浪子的故事.耶穌說浪子遠離家庭,花費了一切所有的,落至貧窮,飢餓的地步.那時才想到父親,他知道父親經常為他禱告,這是他決定回家的主要原因.
\newline
\newline
未知各位有沒有建立家庭崇拜的生活?你的兒女有沒有聽過你禱告?當你的兒女遠離家庭,他們能否記得父母曾為他們禱告,好像浪子的父親為他兒子禱告一樣.
\newline
\newline
(四)建立救恩的牆
\newline
\newline
我們要確實知道紀律的牆,要清楚知道榜樣的牆,要知道家庭崇拜的牆,已經建立起來；但是如果家中還有人未得救,救恩的牆就不能建立.可能在座中有的人是家中唯一的基督徒,其他都未接受耶穌,你當盡力把全家帶到神的面前.要向他們作見證,每天為他們禱告；除非全家都接受耶穌,否則你的家,永遠不能成為基督化的家庭.或者你全家除了你之外都是基督徒,惟獨你至今未接受耶穌作你的救主,你未成為神家的人,你未曾求神饒恕你的罪,未曾求神救你脫離罪惡,因著你的阻礙,以致此家不能成為基督化的家庭.神的應許,是要賜福給基督徒的家庭,只要你接受耶穌,成為基督徒,神就必賜福給你的全家.
\newline
\newline
請各位細心思想自己的家庭；在你家裡,紀律的牆有否建立?好榜樣的牆有否建立?家庭崇拜生活的牆有否建立起來?救恩的牆是否建立起來呢?
\newline
\newline
相信各位都希望有個基督化的家庭,但是必須建立這四幅牆.
\newline
\newline


\newpage

\section{第55屆~港九培靈會~張子華~第一講　神對教會的呼召}
\label{sec:510}
\textbf{第一講　神對教會的呼召}
\newline
\newline
連結: \href{https://www.hkbibleconference.org/session-message/view/510}{\texttt{https://www.hkbibleconference.org/session-message/view/510}}
\newline
\newline
\hyperref[sec:527]{< < < PREV SERMON < < <}
~
\hyperref[sec:toc]{[返目錄]}
~
\hyperref[sec:511]{> > > NEXT SERMON > > >}
\newline
\newline
哥林多前書 14:40-14:40
\newline
\begin{longtable}{cl}
\hline
\hline
章節 & 經文 (和合本修訂版)\\
\hline
14:40 & \begin{tabularx}{0.7\textwidth}{X} 凡事都要規規矩矩地按著次序行。 \end{tabularx} \\ \\
[1ex]
\hline
\hline
\end{longtable}
在今次講道會中,我所揀選的經卷是哥林多前書.而今天我們要思想的經節是十四章四十節.「凡事都要規規矩矩的,按著次序行.」意思就是說,無諭大事,小事,每一件,我們都要按著神給我們的規矩和次序行出來.
\newline
\newline
哥林多前書,是保羅在以弗所住了三年之後,臨離開以弗所之時；大概是主後五十五年所寫的.而哥林多教會,是保羅在第二次宣教旅行之時所建立的.保羅除了在以弗所逗留較長時間之外,在哥林多城也住了一段時間.哥林多是一個歷史古城,也是古代一個商業中心；更是一個出名道德倫理敗壞的城市.如果把哥林多用作動詞,我「哥林多」你,意思就是我要將你差去黑暗之處.哥林多是個罪惡的城市,裡面有廟宇,有妓院,道德非常低落.但神使用保羅在這黑暗的地方,建立祂的教會,讓重生得救的基督徒；在一個黑暗罪惡包圍的環境下,可以過一個成聖的生活.
\newline
\newline
可是,在這樣的環境下,一班重生得救的基督徒,因為有過去在罪中被蒙蔽的背景,很容易便將以前的生活習慣和思想,帶進教會的生活中；以致在教會中產生種種失去神的規矩,次序和美麗的表現.保羅寫這封信主要的目的,是針對當時種種不合神規矩和次序的事,從而給這些問題提出解決和答案.這本書十分適合現今世代的需要,特別是適合香港教會的需要.香港是一個商業城市,也是一個道德淪亡的城市；我們往往可以在香港的教會中,找出一些和哥林多教會相同的問題.因此,研讀哥林多前書,有助於我們了解教會的問題,也得到解決的方法.
\newline
\newline
剛才所讀的一節經文,凡事都要規規矩矩的按著次序行,是神藉保羅給哥林多教會的一個呼召.同樣,今天神也呼召我們,要脫離那些沒有規矩的事；那些不能將神的美麗表彰出來的事；讓我們的教會,建立在一些屬主的次序和神的美麗上.今天我們先透過這節經文,對全卷書有個綜合的認識,然後從明天開始,我們才解決一些今日教會所面臨的問題.
\newline
\newline
第一,這是一個全面性的呼召.我們從三方面來思想這個呼召.保羅說:「凡事都要」,是指每一件事,不光是大事或小事,乃是每一件事,包括個人生活和教會生活的每一部份.在信徒個人生活上,在教會中,事無大小,我們都應該放在神的管理之下.保羅在書中所提到的各種問題,雖然有些可能涉及時代性的觀念:例如蒙頭等,但聖經中所提的每一件事,我們在原則上都可以應用在現今的世代上.
\newline
\newline
保羅在書中所提的問題,包括了我們在教會中的生活,在社會中的生活,在家庭中的生活,和在信仰中的生活.而每一方面的生活,在神面前都是有規矩,有次序的.譬如說,論到教會的紛爭,我們都有由意見不同而產生出來的爭論.保羅說:凡爭論都是屬肉體,沒有次序,在多方面能損害教會生活的,明天我們會思想這問題的解決.而保羅又指出,教會之所以出現紛爭,主要是因為人不尊重神的智慧,後天我們便會思想甚麼是神的智慧.教會,家庭,或個人,如果尊重神的智慧,在生活上會產生怎樣的行為.第五章講到教會發生淫亂的事,到底教會應該怎去處理這些罪惡,和懲罰行這些事的人呢?而教會經常發生的糾紛,保羅說不應該帶到法院去解決,那麼教會要怎樣處理呢?另外我們會講到基督徒的獨身和結婚問題.至於第八章所講可否吃祭偶像之物,可以提供給我們今日在生活上選擇的原則.還有聖餐問題,金錢的奉獻等,都是我們要思想的.從這一卷書中,我們認識到我們生活的各方面,沒有一樣是不需要次序的.保羅所給我們的呼召,是一個全面性的呼召.
\newline
\newline
第二,這個呼召是行動性的.聖經說:凡事都要規規矩矩的按著次序行.這個「行」字,英文是用「done」,就是已經按著次序做成.當保羅向哥林多教會發出這呼召之後,他不是要他們拿這個呼召去討論,去發表意見,他乃是要他們立即採取一個完成的行動.這是個行動性的態度,今天我們領受這個呼召的時候,也要有所行動.
\newline
\newline
當我們實行的時候,第一要有知識,第二要勤力.實行這呼召需要有知識,這是神的心意,神不希望我們在解決教會問題時,陷入無知的光景中.我們不能用無知去原諒自己,而不活出神的規矩來.當然,開始的時候我們都是嬰孩,一無所知,但這嬰兒的光景在神的憐憫,聖靈的工作,我們的追求下,很快便應脫離.我們不再是無知,乃是清楚的知道神在我們生活各方面所定的規矩.
\newline
\newline
基督徒應該在法院解決糾紛嗎?當然不可以,保羅在講到這問題時,在第六章用了很多次「豈不知」,就是怎可以不知.第十章一節更說:「弟兄們,我不願意你們不曉得,」就是我不願意看見你們無知.十二章一節又說:「弟兄們,論到屬靈的恩賜,我不願意你們不明白.」神不容許我們在無知中生活；但如果我們不追求,刻意無知,神便會放棄我們.十四章三十八節說:「若有一知道的,就由你不知道吧!」落到這樣的光景,是何等可怕的事.
\newline
\newline
另一方面,這行動性的呼召是需要我們努力去作的.面對著教會不規矩的問題,我們必須付代價的去解決,去求神改變.每一個問題的解決都不是簡單的,當教會有兩方面的爭執和意見時,雙方都必須努力付代價才能解決問題.當教會已經有人娶了他的繼母時,他必須付很大的代價才能從這罪惡中出來.在基督徒的生活中,在教會的建立中,是需要很多血汗.屬靈的成長,基督榮美的彰顯,教會有次序的表現,是沒有捷徑的.懶惰,不肯付代價,只能帶來失敗.所以保羅說我呼召你,這呼召是行動性的.我們要有足夠的知識去面對問題,更要行動去解決問題.
\newline
\newline
第三,這是個印象性的呼召.當我們以行動去解決了問題之後,會留下一個美麗的印象,一個美好的見證給社會.我們留意「規規矩矩」和「次序」這兩個詞,這兩個詞是代表美麗和責任.舊約聖經說我們要用美麗去敬拜神,這代表主人對奴隸的一個裝飾.在基督徒的生活中,我們也應有基督的裝飾,有祂的美麗和聖潔.次序是代表責任.基督徒在他的生活中,在每一件事上都應有他的先後次序.一般解經家認為這詞有軍事或軍用的含意.這代表一個將領,在攻打敵人的時候,先做了三方面的功夫；就是思想,準備和操練.然後才把軍隊差出去.保羅用這個詞,就是要我們每一間神的教會,每一個基督徒的生活,當我們面對所要做的事時,先要從聖經中經過詳細的思考,在禱告中經過仔細的計劃,在行動上經過良好的操練；這樣才能使每一件事情在屬靈的次序下井井有條的進行.
\newline
\newline
舉例來說,禱告.有多少人的禱告生活是有計劃的呢?我們怎樣禱告呢?會不會來來去去都是一樣的禱詞,沒有新的內容呢?毛病在哪裡呢?其實主要的原因是因為禱告沒有次序.一個有次序的禱告生活,是經過你的思想,計劃,把所要禱告的事項都列寫出來.這樣,當你在神面前禱告的時候,不但不會沒有禱告的事項,內容,而且你更會發覺不夠時間禱告.
\newline
\newline
很多教的祈禱會也是沒有計劃的,叫參加的人心痛.事先沒有計劃,不但沒有準備要祈禱的事項,就連聚會要唱的詩,也毫無準備.到時大家翻翻詩歌,隨便點唱幾首；大家胡亂的提出一些禱告事項,這樣便過了一週一次的祈禱會；是多麼浪費,多麼可惜.教會的禱告會沒有次序,崇拜會沒有次序,又沒有預備,這包括心靈上的預備,安排上的預備,這樣怎能討神的喜悅,怎能滿足主的心意呢?因此,我們必須在教會的事工上,在自己的生活上,將基督的美麗活出來,也要有次序的去行,這樣才能在世界上留下一個滿有印象的見證.讓我們謹記聖經的話.「凡事都要規規矩矩的按著次序行.」
\newline
\newline


\newpage

\section{第55屆~港九培靈會~張子華~第二講　教會的紛爭}
\label{sec:511}
\textbf{第二講　教會的紛爭}
\newline
\newline
連結: \href{https://www.hkbibleconference.org/session-message/view/511}{\texttt{https://www.hkbibleconference.org/session-message/view/511}}
\newline
\newline
\hyperref[sec:510]{< < < PREV SERMON < < <}
~
\hyperref[sec:toc]{[返目錄]}
~
\hyperref[sec:512]{> > > NEXT SERMON > > >}
\newline
\newline
哥林多前書 1:10-1:17
\newline
\begin{longtable}{cl}
\hline
\hline
章節 & 經文 (和合本修訂版)\\
\hline
1:10 & \begin{tabularx}{0.7\textwidth}{X} 弟兄們，我藉我們主耶穌基督的名勸你們說話要一致。你們中間不可分裂，只要一心一意彼此團結。 \end{tabularx} \\ \\ \relax
1:11 & \begin{tabularx}{0.7\textwidth}{X} 我的弟兄們，革來氏家裡的人曾對我提起你們，說你們中間有紛爭。 \end{tabularx} \\ \\ \relax
1:12 & \begin{tabularx}{0.7\textwidth}{X} 我的意思是，你們各人說：「我是屬保羅的」；「我是屬亞波羅的」；「我是屬磯法的」；「我是屬基督的。」 \end{tabularx} \\ \\ \relax
1:13 & \begin{tabularx}{0.7\textwidth}{X} 基督是分裂的嗎？保羅為你們釘了十字架嗎？你們是奉保羅的名受了洗嗎？ \end{tabularx} \\ \\ \relax
1:14 & \begin{tabularx}{0.7\textwidth}{X} 我感謝神，除了基利司布和該猶以外，我沒有給你們中的任何一個人施洗， \end{tabularx} \\ \\ \relax
1:15 & \begin{tabularx}{0.7\textwidth}{X} 免得有人說你們是奉我的名受洗的。 \end{tabularx} \\ \\ \relax
1:16 & \begin{tabularx}{0.7\textwidth}{X} 我曾為司提法那家施過洗；此外我已記不清有沒有給別人施過洗。 \end{tabularx} \\ \\ \relax
1:17 & \begin{tabularx}{0.7\textwidth}{X} 因為基督差遣我不是為施洗，而是為傳福音；並不是用智慧的言論，免得基督的十字架落了空。 \end{tabularx} \\ \\
[1ex]
\hline
\hline
\end{longtable}
今天我們開始研究哥林多教會第一個沒有次序的問題,是大大損害教會美麗的見證的,就是教會的紛爭.哥林多教會是一個富足的教會,是一個口才知識都全備的教會；在他們的恩賜,信心和盼望上,沒有一樣是不及別人的.可是一個在人看是滿有才幹和富足的教會中,卻有許多事情失去了教會的次序,因此也失去了神的榮美.第一樣保羅要責備的,就是他們中間有紛爭.
\newline
\newline
我們看一章十至十七節:「弟兄們,我藉我們主耶穌基督的名,勸你們都說一樣的話；你們中間也不可分黨,只要一心一意彼此相合.因為革來氏家裡的人,曾對我題起弟兄們來,說你們中間有紛爭.我的意思就是你們各人說,我是屬保羅的,我是屬亞波羅的,我是屬磯法的,我是屬基督的.基督是分開的麼?你們是奉保羅的名受了洗麼?我感謝神,除了基利司布並該猶以外,我沒有給你們一個人施洗.免得有人說,你們是奉我的名受洗.我也給司提反家施過洗,此外給別人施洗沒有,我卻記不清.基督差遣我,原不是為施洗,乃是為傳福音,並不用智慧的言語,免得基督的十字架落了空.」
\newline
\newline
在這裡,保羅講到了紛爭的咒詛,就是教會的紛爭所帶來的損害.以後一章十八節起到四章五節止,是保羅提出引致教會紛爭的幾個理由,而這些原因都是和神的智慧相違背的.而四章六到二十一節,就是解決紛爭的方法.今天我們只看紛爭的咒詛和解決.我們看四章六到十三節:「弟兄們,我為你們的緣故,拿這些事轉比自己和亞波羅,叫你們效法我們不可過於聖經所記,免得你們自高自大,貴重這個,輕看那個,使你與人不同的是誰呢?你有甚麼不是領受的呢?若是領受的,為何自誇,彷彿不是領受的呢?你們已經飽足了,已經豐富了,不用我們,自己就作王了.我願意你們果真作王,叫我們也得與你們一同作王.我想神把我們使徒明明列在末後,好像定死罪的囚犯,因為我們成了一台戲,給世人和天使觀看.我們為基督的緣故算是愚拙的,你們在基督裡倒是聰明的.我們軟弱,你們倒強壯；你們有榮耀,我們倒被藐視.直到如今,我們還是又飢,又渴,又赤身露體,又挨打,又沒有一定的住處.並且勞苦,親手作工.被人咒罵,我們就祝福；被人逼迫,我們就忍受；被人毀謗,我們就善勸.直到如今,人還把我們看作世界上的污穢,萬物中的渣滓.」
\newline
\newline
教會在社會一個強而有力的見證,就是我們將彼此相愛,在基督裡合而為一.其實,每一個重生得救的基督徒,都應該是合而為一的,因為我們享受同一的生命,我們有同一的聖靈作交通,也學習同一本的聖經.因此我們讀聖經,特別是詩篇第一百三十三篇,很容易看見一幅弟兄和睦同居的美麗圖畫.不過,話雖是如此,但今天在神的家中,因為人肉體的表現,卻時常產生紛爭；而且大部份的紛爭,都不是為了真理.
\newline
\newline
保羅告訴我們,教會的紛爭會帶來三個很大的咒詛和損害.第一,是任何的紛爭都會帶來教會生活的癱瘓.我們都知道,教有一個很美麗的團契生活,神將我們放在同一個教會中,叫我們能在生命中有深交,這是教會獨有的特色和生活.這種生活是建立在兩個根基上,第一是在神話語上的學習,藉著崇拜,主日學,查經班,交通聚會等,我們在神的真道上學習,漸漸同歸於一.第二,是我們在聖靈的帶領下,同心合意的禱告和事奉,這就是合一,在靈裡合一.然而,任何的糾紛都會損害這兩個重要的根基.保羅為信徒紛爭所帶來的損害而嘆息,而他也苦勸他們要一心一意的彼此相合.彼此相合,原來是補網的意思.我們要教會中作補網的工作,這樣才不致形成弟兄姊妹間的分難.
\newline
\newline
第二,紛爭所引致的咒詛,就是使教會充滿謠言.當謠言傳開時,便會由意見不合而引致敵對性的行動,同時也讓教會的見證受到虧損.華來氏並不是哥林多教會的信徒,可見他們在教會以外也聽到這謠言,而轉達給保羅.那麼紛爭所引起的謠言會叫人對教會的見證產生壞的印象.
\newline
\newline
第三,教會的紛爭,能叫教會領袖的功能受到損害,叫一班屬靈的領袖無法帶領教會向前行.因為他們把自己奉獻給人,而不是奉獻給神.像哥林多的教會,他們擁護人,稱自己是保羅派,亞波羅派,磯法派,基督派.起初他們可能只是崇尚某一個人和他的教導,但慢慢這些派卻演變成為敵對性的行動,倘若我們也把自己奉獻給人,擁護某傳道人,而輕看某傳道人,這樣,教會領袖的功能便無法發揮,陷於癱瘓了.
\newline
\newline
到底保羅對於教會的紛爭有甚麼解決的方法呢?保羅所強調的,是基督徒本身的生命.如果我們有四章六到十三節所講生命的經歷時,我們肯定不會進入任何的糾紛中.我們甚至會讓步到一個程度,即使人看我們是渣滓,我們也不會計較.
\newline
\newline
在這一段經文中,首先我們要注意的,是兩件完全對立的事,那就是以自己為中心,和以基督為中心.人若不以基督為中心,便是以自己為中心；人若以自己為中心,就不能以基督為中心.而一切的爭論和糾紛,都是因為人以自己為中心,不以基督為中心.這是生命主要的經歷.我們留意保羅在這裡提到「自高自大」這一個詞,在整卷書信中,保羅用這個詞共有七次之多.自高自大的意思,就是以自己為中心,一個人在生命中驕傲.
\newline
\newline
驕傲是非常可怕的.首先,驕傲不合乎聖經的原則.第六節說:「弟兄們,我為你們的緣故,拿這些轉比自己和亞波羅,叫他們效法我們不可過於聖經所記,免得你們自高自大,貴重這個,輕看那個.」原來聖經很清楚的告訴我們,神阻擋驕傲的人,賜恩給謙卑的人.神要粉碎人的驕傲,我們害怕對付自己的驕傲.人在糾紛中,必然是以自我為中心,有驕傲的表現.他不會讓步,也不肯讓步,乃是抵抗到底.
\newline
\newline
另外,驕傲是一個屬乎肉體,不屬靈的表現.使你與人不同的是誰呢?你有甚麼不是領受的呢?若是領受,為何自誇,彷彿不是領受的呢?一個在糾紛中驕傲的人,永遠看不見神造每一個人都是不同的.其實,我們生命中每一樣都是領受的,是從神而來的,我們必須看清這個,我們要在基本的態度上,在言語,行為,表現上,要充滿恩惠,因為我們是蒙恩的人.驕傲和恩典是相反的,一個人若明白甚麼是恩典,便絕對不會驕傲,因為他生命中的一切都是從神領受的.
\newline
\newline
還有,驕傲能叫人與人之間隔離.第八節:「你們已經飽足了,已經豐富了,不用我們,自己就作王了.我願意你們果真作王,叫我們也得與你們一同作王.」當人進入一個以自我為中心的驕傲的生活時,必然會產生一種態度,就是他不再需要任何人,就如哥林多的信徒一樣,他們連保羅也不再需要.一個驕傲的人,不只排斥別人；而且常覺得自己比別人強,因此驕傲便把人與人隔絕了.
\newline
\newline
在講完以自我為中心之後,保羅從九到十三節,便以自己的榜樣告訴信徒,一個基督徒怎樣以基督為中心；或者說一個以基督為中心的基督徒,他的生命和表現應該怎樣.保羅從第九節開始便引用一個比喻,說我們是一台戲.古時在大劇場有很多節目上演,而最後的節目,是將死囚帶出來和獅子搏鬥,其實就是餵獅子.保羅用這比喻講出,基督徒如果接受了十字架的生命,如果有十架對付的經歷,便會知道怎樣處理糾紛的事；他會讓步,寧可被人冤枉.保羅這樣的人生和態度,正是我們要追求的.十字架的經歷,帶給人精神上的痛苦,和肉身上的痛苦,但是卻使人能真正體驗到主的心腸,這正是解決教會紛爭的唯一方法.
\newline
\newline


\newpage

\section{第55屆~港九培靈會~張子華~第三講　教會與法院}
\label{sec:512}
\textbf{第三講　教會與法院}
\newline
\newline
連結: \href{https://www.hkbibleconference.org/session-message/view/512}{\texttt{https://www.hkbibleconference.org/session-message/view/512}}
\newline
\newline
\hyperref[sec:511]{< < < PREV SERMON < < <}
~
\hyperref[sec:toc]{[返目錄]}
~
\hyperref[sec:539]{> > > NEXT SERMON > > >}
\newline
\newline
哥林多前書 6:1-6:11
\newline
\begin{longtable}{cl}
\hline
\hline
章節 & 經文 (和合本修訂版)\\
\hline
6:1 & \begin{tabularx}{0.7\textwidth}{X} 你們中間有彼此爭吵的事，怎敢告到不義的人面前，而不告到聖徒面前呢？ \end{tabularx} \\ \\ \relax
6:2 & \begin{tabularx}{0.7\textwidth}{X} 你們豈不知聖徒要審判世界嗎？若世界要受你們的審判，難道你們不配審判這最小的事嗎？ \end{tabularx} \\ \\ \relax
6:3 & \begin{tabularx}{0.7\textwidth}{X} 你們豈不知我們要審判天使嗎？何況今生的事呢！ \end{tabularx} \\ \\ \relax
6:4 & \begin{tabularx}{0.7\textwidth}{X} 既是這樣，你們若有今生當審判的事，會讓教會所輕看的人來審判嗎？ \end{tabularx} \\ \\ \relax
6:5 & \begin{tabularx}{0.7\textwidth}{X} 我說這話是要使你們慚愧。難道你們中間沒有一個有智慧的人能審斷弟兄中的事嗎？ \end{tabularx} \\ \\ \relax
6:6 & \begin{tabularx}{0.7\textwidth}{X} 你們竟然有弟兄去告弟兄，而且告到不信主的人面前。 \end{tabularx} \\ \\ \relax
6:7 & \begin{tabularx}{0.7\textwidth}{X} 你們彼此告狀，這已經是你們的大錯了。為甚麼不情願受冤屈呢？為甚麼不情願吃虧呢？ \end{tabularx} \\ \\ \relax
6:8 & \begin{tabularx}{0.7\textwidth}{X} 你們反倒去冤枉人，虧負人，況且所冤枉所虧負的就是弟兄。 \end{tabularx} \\ \\ \relax
6:9 & \begin{tabularx}{0.7\textwidth}{X} 你們豈不知不義的人不能承受神的國嗎？不要自欺！無論是淫亂的、拜偶像的、姦淫的、作娼妓的，親男色的、 \end{tabularx} \\ \\ \relax
6:10 & \begin{tabularx}{0.7\textwidth}{X} 偷竊的、貪婪的、醉酒的、辱罵的、勒索的，都不能承受神的國。 \end{tabularx} \\ \\ \relax
6:11 & \begin{tabularx}{0.7\textwidth}{X} 從前你們中間也有人是這樣；但現在你們奉主耶穌基督的名，並藉著我們神的靈，已經洗淨，已經成聖，已經稱義了。 \end{tabularx} \\ \\
[1ex]
\hline
\hline
\end{longtable}
請看(林前六1-11):「你們中間有彼此相爭的事,怎敢在不義的人面前求審,不在聖徒面前求審呢?豈不知聖徒要審判世界麼?若世界為你們所審,難道你們不配審判這最小的事麼?豈不知我們要審判天使麼?何況今生的事呢?既是這樣,你們若有今生的事當審判,是派教會所輕看的人審判麼?我說這話,是要叫你們羞恥.難道你們中間沒有一個智慧人,能審判弟兄們的事麼?你們竟是弟兄與弟兄告狀,而且告在不信主的人面前.你們彼此告狀,這已經是你們的大錯了.為甚麼不情願受欺呢!為甚麼不情願吃虧呢?你們倒是欺壓人,虧負人,況且所欺壓,所虧負的,就是弟兄.
\newline
\newline
你們豈不知,不義的人不能承受神的國麼?不要自欺,無論是淫亂的,拜偶像的,姦淫的,作孌童的,親男色的,偷竊的,貪婪的,醉酒的,辱罵的,勒索的,都不能承受神的國.你們中間也有人從前是這樣,但如今你們奉主耶穌基督的名,並藉著我們神的靈,已經洗淨,成聖稱義了.」
\newline
\newline
哥林多前書的總題,是建立一個有次序有美麗的教會.保羅一開始便針對一個教會中最失去次序,最不能表彰神榮美的問題,就是他們中間有紛爭,結黨.昨天我們看過紛爭所帶給教會的咒詛,損害；其中一個最嚴重的,是叫教會生活陷於完全的癱瘓.另外也叫教會產生不必要的謠言,使教會領袖的功能消失,這是教會紛爭所引致的咒詛.面對這問題,保羅提出一個最基本生命質素的問題,來作為教會紛爭的醫治和解決.他要我們除去自高自大,以人,以己為中心,並經歷十字架道路所產生在精神上和肉體上的痛苦,寧願被人冤屈,甚至看為人間的渣滓,也都要接受.這是停止紛爭的惟一方法.
\newline
\newline
今天我們要看的紛爭,是規模更大的,就是教會中的糾紛,嚴重到一個地步；要帶到法庭去,彼此告狀.香港有不少教會,也曾做過這些最失去恩典和見證的事.在這裡,我們必須建立起一個原則,凡是在教會內的事,信徒一切的生活,都是屬於信心的生活；而信心的生活,是不能被一個沒有信心的法官去審斷；因為他沒有這個資格,這是基督徒必須持守的重要原則.教會中一切的生活,信徒之間的關係,只有一個人有資格去審判,就是主耶穌.也只有一本律法書,可以作為我們審判,定奪是非的根據的,就是聖經.我並不是否定法庭在社會中的地位,事實上法庭在社會中有重要的地位.當保羅被不信主的人冤枉時,他也曾上法庭為自己申訴.不過保羅在這裡所講的,乃是指弟兄姊妹之間,教會內部的事,若有糾紛,絕不能帶到法院.而保羅在這章聖經中,也建議了一個實際的方法去解決教會內一切的糾紛.
\newline
\newline
保羅有四個理由講出教會絕對不能在法院打官司.
\newline
\newline
(一)因為教會打官司是違背教會將來在天上所執行榮耀的職份
\newline
\newline
將來我們要做的,是審判世界,當主再來作王時,我們也要和他一同作王,並要審判世界；既然將來我們要審判世人,今天我們怎能要求世人審判我們.將來我們要審判的人,今生怎會有資格去審判我們這些屬靈的人呢?除了審判世界之外,將來我們還要審判天使.豈不知我們要審判天使麼?何況今生的事呢?既然將來我們有這樣重要的職份,可以審判一切屬靈的創造物；那麼今生有甚麼事我們不可以自己解決,反倒要求一個不信的法官去為我們判斷呢?
\newline
\newline
(二)教會打官司是錯的,因為它違背了教會對於糾紛所處理的方針
\newline
\newline
對於解決教會的糾紛,保羅說:「難道你們中間沒有一個智慧人,能審斷弟兄們的事麼?」由智慧人去處理教會的糾紛,但誰是智慧人呢?保羅認為教會應該有由信徒所組織的法庭,而且有智慧人做法官,去判決事情.教會不是一個民主的團體,乃是一個神治的團體.神是絕對的,他說怎樣就怎樣,他所說的就是真理.既然是神治的教會,神便揀選一些合祂用的僕人器皿,使教會有牧師,傳道,長老和執事等.神賜這些屬靈的領袖有權柄,讓他們可以帶領弟兄姊妹走神的路.信徒應該對神所選召的領袖們有屬靈的順服,順服他們的帶領,這是信徒的責任.
\newline
\newline
(三)教會不應該打官司,因為違背了教會應有的愛情
\newline
\newline
第七,八節:「你們彼此告狀,這已經是你們的大錯了,為甚麼不情願受欺呢?為甚麼不情願吃虧呢?你們倒是欺壓人,虧負人,況且所欺壓,所虧負的,就是弟兄.」哥林多是一個知識,口才,恩賜都全備的教會,而且也是富足的,但可惜卻缺少了信徒之間應有的愛.今天,很多教會也人才濟濟；但真正的愛在哪裡呢?當我們處在紛爭中的時候,有否問自己,十字架在哪裡?當我們在紛爭的事上堅持不下時,有否問自己愛在哪裡?到底我願不願吃虧?我願不願讓步?我願不願像主在被人誣告時,在十字架上還說:「父阿,赦免他們,因為他們所作的,他們不曉得.」教會的愛是一個犧牲,受苦,受冤屈,吃虧,被誣告的.這是關乎生命的問題.因此,打官司是違背了教會的愛情.
\newline
\newline
(四)教會打官司是錯的,因為他違背了教會的聖潔
\newline
\newline
保羅為甚麼要寫下九到十一節的經文呢?這幾節經文好像和上文不相連,但其實並不然；保羅的重點是在第十一節；他要指出,你們中間有人從前是這樣,是像第九,第十節所講的那樣敗壞.但保羅說:如今你們奉主耶穌基督的名,並藉著我們神的靈,已經洗淨,成聖,稱義了.既然已經從罪惡的環境中出來,成為聖潔,怎能再藉打官司進入這樣的人中間?他們怎會有這聖潔的標準呢?如果我們已經從罪惡污穢中出來,進入了聖潔；卻又回到不聖潔的世人那裡,求他們審判,這就是違背了教會的聖潔；枉費了主耶穌把我們從罪惡買贖出來的功勞.
\newline
\newline
如果你的教會是視產業如命,視財如命；為了產業的緣故,可以打到頭崩額裂.為了產業,為了錢財,甚至把案件帶到法官面前求審判.這樣的行為是甚麼呢?當你在污穢中,未信主之前,你對產業和錢財有一種看法；但成聖之後,你對產業和錢財又有另一種看法.這種新的價值觀,使你看錢財和產業是身外物,是神所託付給你的.但當你打官司的時候,你是把這高尚的價值觀念,帶到金錢是萬能的價值觀裡去求審判；這樣做是違背了教會的聖潔.
\newline
\newline
我很覺得,有一段日子,香港教會對社會沒有影響,而社會人士對教會也沒有特別的好感；主要的原因,是因為有幾間大教會曾經將官司打到法庭去,報紙以頭條新聞登出來.一班未信主的社會人士,便覺得教會不過是一個普通的團體；沒有分別,因為教會也打官司.這樣做是最失去恩典,最失去見證；足以影響耶穌基督的名,將耶穌的名損害.因此保羅在這章聖經內告訴我們,千萬不可以打官司.
\newline
\newline
打官司是錯的,因為他違背了我們將來要作的事；將來我們要審判世界,審判天使.打官司是錯,因為它違背了教會應該有屬靈的價值；以信心去解決教會一切的糾紛,由智慧人去處理這些案件.打官司是錯的,因為它違背了神在我們生命中所產生的愛情,這愛情是經歷過十字架而出來的；是願意吃虧,受冤屈,付最大的代價,甚至是捨己的愛情.打官司是錯的,因為它違背了教會聖潔的標準.
\newline
\newline
盼望聖靈在我們心中工作,如果有人想打官司的,應立即停止；讓我們在基督的愛情中生活,以後不再有任何糾紛是需要帶到不義的法院去.
\newline
\newline


\newpage

\section{第55屆~港九培靈會~張子華~第四講　神的智慧}
\label{sec:539}
\textbf{第四講　神的智慧}
\newline
\newline
連結: \href{https://www.hkbibleconference.org/session-message/view/539}{\texttt{https://www.hkbibleconference.org/session-message/view/539}}
\newline
\newline
\hyperref[sec:512]{< < < PREV SERMON < < <}
~
\hyperref[sec:toc]{[返目錄]}
~
\hyperref[sec:513]{> > > NEXT SERMON > > >}
\newline
\newline
哥林多前書 1:18-4:5
\newline
\begin{longtable}{cl}
\hline
\hline
章節 & 經文 (和合本修訂版)\\
\hline
1:18 & \begin{tabularx}{0.7\textwidth}{X} 因為十字架的道理，在那滅亡的人是愚拙，在我們得救的人卻是神的大能。 \end{tabularx} \\ \\ \relax
1:19 & \begin{tabularx}{0.7\textwidth}{X} 就如經上所記：「我要摧毀智慧人的智慧，廢棄聰明人的聰明。」 \end{tabularx} \\ \\ \relax
1:20 & \begin{tabularx}{0.7\textwidth}{X} 智慧人在哪裡？文士在哪裡？這世上的辯士在哪裡？神豈不是已使這世上的智慧變成愚拙了嗎？ \end{tabularx} \\ \\ \relax
1:21 & \begin{tabularx}{0.7\textwidth}{X} 既然世人憑自己的智慧不認識神，神就本著自己的智慧樂意藉著人所傳愚拙的話拯救那些信的人。 \end{tabularx} \\ \\ \relax
1:22 & \begin{tabularx}{0.7\textwidth}{X} 猶太人要的是神蹟，希臘人求的是智慧， \end{tabularx} \\ \\ \relax
1:23 & \begin{tabularx}{0.7\textwidth}{X} 我們卻是傳被釘十字架的基督，這對猶太人是絆腳石，對外邦人是愚拙； \end{tabularx} \\ \\ \relax
1:24 & \begin{tabularx}{0.7\textwidth}{X} 但對那蒙召的，無論是猶太人、希臘人，基督總是神的大能，神的智慧。 \end{tabularx} \\ \\ \relax
1:25 & \begin{tabularx}{0.7\textwidth}{X} 因為，神的愚拙總比人智慧；神的軟弱總比人強壯。 \end{tabularx} \\ \\ \relax
1:26 & \begin{tabularx}{0.7\textwidth}{X} 弟兄們哪，想一想你們的蒙召，按著人的觀點，有智慧的不多，有能力的不多，有尊貴地位的也不多。 \end{tabularx} \\ \\ \relax
1:27 & \begin{tabularx}{0.7\textwidth}{X} 但是，神揀選了世上愚拙的，為了使有智慧的羞愧；又揀選了世上軟弱的，為了使強壯的羞愧。 \end{tabularx} \\ \\ \relax
1:28 & \begin{tabularx}{0.7\textwidth}{X} 神也揀選了世上卑賤的，被人厭惡的，以及那一無所有的，為要廢掉那樣樣都有的， \end{tabularx} \\ \\ \relax
1:29 & \begin{tabularx}{0.7\textwidth}{X} 使凡血肉之軀的，在神面前，一個也不能自誇。 \end{tabularx} \\ \\ \relax
1:30 & \begin{tabularx}{0.7\textwidth}{X} 但你們得以在基督耶穌裡是本乎神，他使基督成為我們的智慧，成為公義、聖潔、救贖。 \end{tabularx} \\ \\ \relax
1:31 & \begin{tabularx}{0.7\textwidth}{X} 如經上所記：「要誇耀的，該誇耀主。」 \end{tabularx} \\ \\ \relax
2:1 & \begin{tabularx}{0.7\textwidth}{X} 弟兄們，從前我到你們那裡去，並沒有用高言大智對你們宣講神的奧祕。 \end{tabularx} \\ \\ \relax
2:2 & \begin{tabularx}{0.7\textwidth}{X} 因為我曾定了主意，在你們中間不知道別的，只知道耶穌基督並他釘十字架。 \end{tabularx} \\ \\ \relax
2:3 & \begin{tabularx}{0.7\textwidth}{X} 我在你們那裡時，又軟弱，又懼怕，又戰戰兢兢。 \end{tabularx} \\ \\ \relax
2:4 & \begin{tabularx}{0.7\textwidth}{X} 我說的話、講的道不是用委婉智慧的言語，而是以聖靈的大能來證明， \end{tabularx} \\ \\ \relax
2:5 & \begin{tabularx}{0.7\textwidth}{X} 為要使你們的信不靠著人的智慧，而是靠著神的大能。 \end{tabularx} \\ \\ \relax
2:6 & \begin{tabularx}{0.7\textwidth}{X} 然而，在成熟的人中，我們也講智慧，但不是今世的智慧，也不是今世有權有位、將要滅亡的人的智慧。 \end{tabularx} \\ \\ \relax
2:7 & \begin{tabularx}{0.7\textwidth}{X} 我們講的是從前隱藏的、神奧祕的智慧，就是神在萬世以前預定使我們得榮耀的智慧； \end{tabularx} \\ \\ \relax
2:8 & \begin{tabularx}{0.7\textwidth}{X} 這智慧，今世有權有位的人沒有一個知道，若知道，他們就不會把榮耀的主釘在十字架上了。 \end{tabularx} \\ \\ \relax
2:9 & \begin{tabularx}{0.7\textwidth}{X} 如經上所記：「神為愛他的人所預備的是眼睛未曾看見，耳朵未曾聽見，人心也未曾想到的。」 \end{tabularx} \\ \\ \relax
2:10 & \begin{tabularx}{0.7\textwidth}{X} 只有神藉著聖靈把這事向我們顯明了；因為聖靈參透萬事，就是神深奧的事也參透了。 \end{tabularx} \\ \\ \relax
2:11 & \begin{tabularx}{0.7\textwidth}{X} 除了在人裡頭的靈，誰知道人的事？照樣，除了神的靈，也沒有人知道神的事。 \end{tabularx} \\ \\ \relax
2:12 & \begin{tabularx}{0.7\textwidth}{X} 我們所領受的並不是世上的靈，而是從神來的靈，為使我們知道神把恩賜賞給我們的事。 \end{tabularx} \\ \\ \relax
2:13 & \begin{tabularx}{0.7\textwidth}{X} 我們也講說這些事，不是用人的智慧所教的言語，而是用聖靈所教的言語，用屬靈的話解釋屬靈的事。 \end{tabularx} \\ \\ \relax
2:14 & \begin{tabularx}{0.7\textwidth}{X} 然而，屬血氣的人不接受神的靈的事，他反倒以這為愚拙，並且他不能了解，因為這些事惟有屬靈的人才能領悟。 \end{tabularx} \\ \\ \relax
2:15 & \begin{tabularx}{0.7\textwidth}{X} 屬靈的人能看透萬事，卻沒有一人能看透他。 \end{tabularx} \\ \\ \relax
2:16 & \begin{tabularx}{0.7\textwidth}{X} 「誰曾知道主的心？誰會教導他？」至於我們，我們有基督的心。 \end{tabularx} \\ \\ \relax
3:1 & \begin{tabularx}{0.7\textwidth}{X} 弟兄們，我從前對你們說話，還不能把你們當作屬靈的，只能把你們當作屬肉體的，你們在基督裡僅是嬰孩。 \end{tabularx} \\ \\ \relax
3:2 & \begin{tabularx}{0.7\textwidth}{X} 我用奶餵你們，沒有用飯餵你們，因為那時你們不能吃。就是如今還是不能， \end{tabularx} \\ \\ \relax
3:3 & \begin{tabularx}{0.7\textwidth}{X} 因為你們仍是屬肉體的。你們中間有嫉妒、紛爭，這豈不是屬乎肉體，照著世人的樣子生活嗎？ \end{tabularx} \\ \\ \relax
3:4 & \begin{tabularx}{0.7\textwidth}{X} 有人說：「我是屬保羅的」；有人說：「我是屬亞波羅的」；這樣你們豈不是和世人一樣嗎？ \end{tabularx} \\ \\ \relax
3:5 & \begin{tabularx}{0.7\textwidth}{X} 亞波羅算甚麼？保羅算甚麼？我們都是神的執事，藉著我們，你們信了；這不過是照著主給各人的恩賜去做罷了。 \end{tabularx} \\ \\ \relax
3:6 & \begin{tabularx}{0.7\textwidth}{X} 我栽種了，亞波羅澆灌了，惟有神使它生長。 \end{tabularx} \\ \\ \relax
3:7 & \begin{tabularx}{0.7\textwidth}{X} 可見，栽種的算不了甚麼，澆灌的也算不了甚麼；惟有神能使它生長。 \end{tabularx} \\ \\ \relax
3:8 & \begin{tabularx}{0.7\textwidth}{X} 栽種的和澆灌的都是一樣，但將來各人要照自己的勞苦得到自己的報酬。 \end{tabularx} \\ \\ \relax
3:9 & \begin{tabularx}{0.7\textwidth}{X} 因為我們是神的同工，而你們是神的田地、神的房屋。 \end{tabularx} \\ \\ \relax
3:10 & \begin{tabularx}{0.7\textwidth}{X} 我照神所給我的恩典，好像一個聰明的工頭，立好了根基，別人在上面建造；只是各人要謹慎怎樣在上面建造。 \end{tabularx} \\ \\ \relax
3:11 & \begin{tabularx}{0.7\textwidth}{X} 因為，那已經立好的根基就是耶穌基督，此外沒有人能立別的根基。 \end{tabularx} \\ \\ \relax
3:12 & \begin{tabularx}{0.7\textwidth}{X} 若有人用金銀、寶石，草木、禾秸，在這根基上建造， \end{tabularx} \\ \\ \relax
3:13 & \begin{tabularx}{0.7\textwidth}{X} 各人的工程必將顯露，因為那日子要將它顯明，有火把它暴露出來，這火要試煉各人的工程怎樣。 \end{tabularx} \\ \\ \relax
3:14 & \begin{tabularx}{0.7\textwidth}{X} 人在那根基上所建造的工程若能保得住，他將要得賞賜。 \end{tabularx} \\ \\ \relax
3:15 & \begin{tabularx}{0.7\textwidth}{X} 人的工程若被燒了，他將損失，雖然他自己將得救，卻要像從火裡經過一樣。 \end{tabularx} \\ \\ \relax
3:16 & \begin{tabularx}{0.7\textwidth}{X} 難道不知你們是神的殿，神的靈住在你們裡面嗎？ \end{tabularx} \\ \\ \relax
3:17 & \begin{tabularx}{0.7\textwidth}{X} 若有人毀壞神的殿，神一定要毀滅那人；因為神的殿是神聖的，這殿就是你們。 \end{tabularx} \\ \\ \relax
3:18 & \begin{tabularx}{0.7\textwidth}{X} 誰都不可自欺。你們中間若有人自以為在今世有智慧，倒不如變為愚拙，好成為有智慧的。 \end{tabularx} \\ \\ \relax
3:19 & \begin{tabularx}{0.7\textwidth}{X} 因為這世界的智慧在神看來是愚拙的。如經上記著：「主使有智慧的人中了自己的詭計；」 \end{tabularx} \\ \\ \relax
3:20 & \begin{tabularx}{0.7\textwidth}{X} 又說：「主知道智慧人的意念，因為它們是虛妄的。」 \end{tabularx} \\ \\ \relax
3:21 & \begin{tabularx}{0.7\textwidth}{X} 所以，無論誰都不可誇耀人；因為萬有都是你們的， \end{tabularx} \\ \\ \relax
3:22 & \begin{tabularx}{0.7\textwidth}{X} 或保羅，或亞波羅，或磯法，或世界，或生，或死，或現今的事，或將來的事，全是你們的， \end{tabularx} \\ \\ \relax
3:23 & \begin{tabularx}{0.7\textwidth}{X} 而你們是屬基督的，基督是屬神的。 \end{tabularx} \\ \\ \relax
4:1 & \begin{tabularx}{0.7\textwidth}{X} 人應該把我們看為基督的執事，為神的奧祕的管家。 \end{tabularx} \\ \\ \relax
4:2 & \begin{tabularx}{0.7\textwidth}{X} 所求於管家的，是要他忠心。 \end{tabularx} \\ \\ \relax
4:3 & \begin{tabularx}{0.7\textwidth}{X} 我被你們評斷，或被別人評斷，我都以為是極小的事；連我自己也不評斷自己。 \end{tabularx} \\ \\ \relax
4:4 & \begin{tabularx}{0.7\textwidth}{X} 雖然我不覺得自己有錯，卻也不能因此判為無罪；審斷我的是主。 \end{tabularx} \\ \\ \relax
4:5 & \begin{tabularx}{0.7\textwidth}{X} 所以，時候未到，在主來以前甚麼都不要評斷，他要照出暗中的隱情，揭發人的動機。那時，各人要從神那裡得著稱讚。 \end{tabularx} \\ \\
[1ex]
\hline
\hline
\end{longtable}
保羅在哥林多前書中所提到的第一個教會問題,就是在教會中有紛爭.保羅用了四章的篇幅去討論這個嚴重的問題,前兩天我們已經看過教會紛爭所帶來的咒詛,與保羅所提供解決教會紛爭的方法；今天我們要講的,是引致教會紛爭的原因.從一章十八節到四章五節這一大段經文中,保羅指出教會紛爭的一個最主要原因,就是信徒高舉人的智慧,而忽略了神的智慧.如果信徒能尊重神的智慧,教會中便可以減少紛爭,甚至沒有紛爭.到底神的智慧可以成就些甚麼?而神的智慧又是甚麼呢?
\newline
\newline
首先我們來看甚麼是神的智慧.(一17-25)說:「基督差遣我,原不是為施洗,乃是為傳福音；並不用智慧的言語,免得基督的十字架落了空.因為十字架的道理,在那滅亡的人為愚拙,在我們得救的人,卻為神的大能.就如經上所記,我要滅絕智慧人的智慧,廢棄聰明人的聰明.智慧人在哪裡,文士在哪裡,這世上的辯士在哪裡.神豈不是叫這世上的智慧變成愚拙麼?世人憑自己的智慧,既不認識神,神就樂意用人所當作愚拙的道理,拯救那些信的人,這就是神的智慧了.猶太人是要神蹟,希利尼人是求智慧；我們卻是傳釘十字架的基督,在猶太人為絆腳石,在外邦人為愚拙.但在那蒙召的,無論是猶太人,希利尼人,基督總為神的能力,神的智慧.因神的愚拙總比人智慧,神的軟弱總比人強壯.」
\newline
\newline
神的智慧是甚麼呢?就是十字架的道理.因為十字架的道理,在那滅亡的人為愚拙.愚拙就是不智慧的,但在我們得救的人,卻為神的大能.神用這十字架的道理來滅絕智慧人的智慧,廢棄聰明人的聰明.神用十字架的道理去揀選那些不承認人的智慧,只承認神的智慧的人.當人認為十字架的道理是愚拙時,是不可信時,神便用人認為是愚拙的十字架的道理,去拯救一切相信的人.因此我們可以很清楚的看見,神智慧的中心,就是十字架的道理.耶穌基督到世上來,為我們釘死在十架上,這就是神的智慧.
\newline
\newline
跟著我們要看神的智慧所為我們成就的事.第一,神的智慧能為我們成就唯一的救贖.(一30)說:「但你們得在基督耶穌裡,是本乎神,神又使他成為我們的智慧,公義,聖潔,救贖.如經上所記,誇口的當指著主誇口.」既然十字架為我們成就了唯一的救贖,那麼我們在十字架下生活的人,便沒有理由可以誇口.我們已經看過,做成教會紛爭的原因,是因為人以自己為中心,而不以神為中心.可是,一個在十字架的救贖下生活的人,便不會存這樣的態度,因為他如果要誇口的話,也是指著主誇口.
\newline
\newline
神要藉著保羅特別提到智慧,是因為哥林多教會的弟兄姊妹很注重希臘人的哲學,崇尚從人所產生出來的智慧；所以保羅不單提到智慧,更把這智慧和公義,聖潔,救贖連在一起.保羅把智慧和公義連在一起,正好答覆了歷代以來人心中的問題:我怎樣才可以與神和好?以人為中心的哲學,可能會教導人不要作虧心事,不要作傷天害理的事,要多行善等,藉此尋求公義,取悅於神.但我們進入救恩的人,卻知道人若要在神面前成為公義,只有一個方法,就是要神稱我們為義.神若稱我們為義,我們更得著公義,不然,我們便不能有公義.當我們接受十字架時,神便將他的公義加在我們身上,叫我們在他面前成為公義的人.
\newline
\newline
神的智慧也和聖潔連在一起,保羅這樣說,是要答覆歷代以來人心中所存的另一個問題,就是人怎樣才能叫自己的罪得著赦免.以人為智慧的宗教哲學,企圖用各種方法去洗脫人的罪,但這一切方法在神面前都是無效的.我們從聖經知道,人得著洗罪的唯一方法,乃是藉基督耶穌的寶血.當神赦免我們時,我們便得著赦免,這是唯一的方法.罪惡的赦免,在基督徒成聖的過程中,是非常重要的.人只有在基督的十字架下,靠著基督的寶血,才可以每天過聖潔的生活.
\newline
\newline
還有,神的智慧也和救贖連在一起.所謂救贖,就是榮耀的釋放,這包括了完全擺脫罪惡的永恆觀念.保羅藉此去解答歷代以來以人為主的宗教哲學所發出的一個問題:到底有沒有天堂?到底有沒有地獄?人死後靈魂的歸宿到底何在?以人為中心的宗教和哲學,提供給我們許多有關永生的答覆.當時一班希臘哲學家,不信身體復活,他們認為人死後,靈魂便會脫離身體,所以即使復活,也與身體無關.不過,在十字架下生活的信徒,經歷過神智慧的拯救後,便深信有一天當我們離開這世界時,要親身經歷這永恆的實在性,而這永恆的生活,是再無罪惡污染的.到時我們就會真正明白甚麼是天堂,甚麼是地獄了.當耶穌再來時,我們可以得著一個完全的釋放,這就是救贖.因此,神的智慧第一樣所能成就的,就是讓我們得著永遠的生命.
\newline
\newline
第二,神的智慧將真正的和平及合一帶給教會,換句話說,神的智慧能成就教會的合一.(26-29):「弟兄們哪!可見你們蒙召的,按著肉體有智慧的不多,有能力的不多,有尊貴的也不多.神卻揀選了世上愚拙的,叫有智慧的羞愧,又揀選了世上軟弱的,叫那強壯的羞愧.神也揀選了世上卑賤的,被人厭惡的,以及那無有的,為要癈掉那有的.使一切有血氣的,在神面前一個也不能自誇.」首先,這段經文講到一個呼召.神用十字架道理的簡單性去拯救人；保羅叫哥林多信徒看看他們中間的成員,他們有智慧的不多,有能力的不多,有尊貴的也不多.即使是有,但他們和別的信徒都是從同一個門進入教會的,所以他們在教會中的地位都是一樣.保羅說這事實能促進教會的合一,當我們明白到聖經的真理,知道在基督裡信徒都是平等時；我們便成為一群不分貧富,尊貴,學識的蒙恩的肢體,這是福音所帶給我們的合一.
\newline
\newline
跟著,保羅說神用福音的尊貴性去拯救人.神揀選愚拙的,叫有智慧的羞愧；神揀選軟弱的,叫強壯的羞愧.世人看是尊貴的,有智慧的,卻不及一個願意接受福音的人,因為他揀選了屬神的智慧,這智慧是更勝一籌.雖然我們藉著十字架得了尊貴和智慧,但我們原來仍是一無可誇的,所以我們只可指著主誇口.
\newline
\newline
最後,神的智慧能使我們成為一個屬靈的人,以致我們能用屬靈的眼光去參透萬事,參透奧祕的事,參透神的心意,以致在教會中能達成合一；不會用人的智慧和人產生出來的意見,彼此敵對.(二6-16):「然而在完全的人中,我們也講智慧,但不是這世上的智慧,也不是這世上有權,有位將要敗亡之人的智慧.我們講的,乃是從前隱藏,神奧祕的智慧,就是神在萬世以前,預定使我們得榮耀的.這智慧,世上有權,有位的人,沒有一個知道的, 他們若知道,就不把榮耀的主釘在十字架上了.如經上所記,神為愛他的人所預備的,是眼睛未曾看見,耳朵未曾聽見,人心也未曾想到的.只有神藉著聖靈向我們顯明了,因為聖靈參透萬事；就是神深奧的事也參透了.除了在人裡頭的靈,誰知道人的事,像這樣,除了神的靈,也沒有人知道神的事.我們所領受的,並不是世上的靈,乃是從神來的靈,叫我們能知道神開恩賜給我們的事.並且我們講說這些事,不是用人智慧所指教的言語,乃是用聖靈所指教的言語,將屬靈的話,解釋屬靈的事.然而屬血氣的人不領會神聖靈的事,反倒以為愚拙；並且不能知道,因為這些事惟有屬靈的人才能看透.屬靈的人能看透萬事,卻沒有一人能看透了他,誰曾知道主的心去教導他呢?但我們是有基督的心了.」
\newline
\newline
基督徒有一個很特別,很寶貴的眼光,就是屬靈的眼光,是神的智慧在一個基督徒的生命中所成就一件偉大的事；這成就能讓我們以一個屬靈的眼光去參透,接納,明白神的事.我們接受了十字架,就是神的智慧之後,我們便處在一個很超越的地位,是一個屬靈的地位.有神的靈在我們裡面運行,能以明白屬靈的奧秘；就如我們裡面的靈,能明白我們內心所想的一樣.這屬靈的眼光,包括了屬靈的光照,啟示,探討和解釋.
\newline
\newline
教會為甚麼有紛爭,各持己見,互相爭持呢?是因為我們沒有成長到一個地步,讓神的靈在我們裡面作主.所以,教會若要消除紛爭,就必須讓聖靈來作主,使我們具備屬靈的眼光,能參透神的事,凡事照神的心意行,這樣,教會以神的心為心,就不會有紛爭了.
\newline
\newline


\newpage

\section{第55屆~港九培靈會~張子華~第五講　教會在倫理上的破產與管教}
\label{sec:513}
\textbf{第五講　教會在倫理上的破產與管教}
\newline
\newline
連結: \href{https://www.hkbibleconference.org/session-message/view/513}{\texttt{https://www.hkbibleconference.org/session-message/view/513}}
\newline
\newline
\hyperref[sec:539]{< < < PREV SERMON < < <}
~
\hyperref[sec:toc]{[返目錄]}
~
\hyperref[sec:514]{> > > NEXT SERMON > > >}
\newline
\newline
哥林多前書 5:1-5:13
\newline
\begin{longtable}{cl}
\hline
\hline
章節 & 經文 (和合本修訂版)\\
\hline
5:1 & \begin{tabularx}{0.7\textwidth}{X} 我確實聽說在你們中間有淫亂的事；這種淫亂連外邦人中也沒有，就是有人和他的繼母同居。 \end{tabularx} \\ \\ \relax
5:2 & \begin{tabularx}{0.7\textwidth}{X} 你們還自高自大！你們不是該覺得痛心，把做這事的人從你們中間趕出去嗎？ \end{tabularx} \\ \\ \relax
5:3 & \begin{tabularx}{0.7\textwidth}{X} 我人雖然不在你們那裡，心卻在你們那裡，好像親自與你們同在。我奉我們主耶穌的名，已經判斷了做這事的人。 \end{tabularx} \\ \\ \relax
5:4 & \begin{tabularx}{0.7\textwidth}{X} 你們聚會的時候，我的心和你們同在。你們藉著我們主耶穌的權能， \end{tabularx} \\ \\ \relax
5:5 & \begin{tabularx}{0.7\textwidth}{X} 要把這樣的人交給撒但，使他的肉體敗壞，好讓他的靈魂在主的日子可以得救。 \end{tabularx} \\ \\ \relax
5:6 & \begin{tabularx}{0.7\textwidth}{X} 你們這樣自誇是不好的。你們不知道一點麵酵能使全團發起來嗎？ \end{tabularx} \\ \\ \relax
5:7 & \begin{tabularx}{0.7\textwidth}{X} 既然你們是無酵的麵，要把舊酵除淨，好使你們成為新團；因為我們逾越節的羔羊—基督已經被殺獻為祭牲了。 \end{tabularx} \\ \\ \relax
5:8 & \begin{tabularx}{0.7\textwidth}{X} 所以，我們來守這節，不可用舊酵，就是不可用惡毒、邪惡的酵，只用純潔真實的無酵餅。 \end{tabularx} \\ \\ \relax
5:9 & \begin{tabularx}{0.7\textwidth}{X} 我先前寫信告訴過你們，不可與淫亂的人交往。 \end{tabularx} \\ \\ \relax
5:10 & \begin{tabularx}{0.7\textwidth}{X} 此話不是泛指這世上所有行淫亂的，或貪婪的，勒索的，或拜偶像的；若是這樣，你們非離開這世界不可。 \end{tabularx} \\ \\ \relax
5:11 & \begin{tabularx}{0.7\textwidth}{X} 但現在，我寫信告訴你們，若有稱為弟兄的人卻仍犯淫亂，或貪婪，或拜偶像，或辱罵，或醉酒，或勒索，這樣的人不可跟他交往，就是跟他吃飯都不可以。 \end{tabularx} \\ \\ \relax
5:12 & \begin{tabularx}{0.7\textwidth}{X} 因為審判教外的人與我何干？教內的人豈不是你們要審判嗎？ \end{tabularx} \\ \\ \relax
5:13 & \begin{tabularx}{0.7\textwidth}{X} 至於外人有神審判他們。如經上說：「要從你們中間把那邪惡的人趕出去。」 \end{tabularx} \\ \\
[1ex]
\hline
\hline
\end{longtable}
今天我們要從哥林多前書第五章來思想一個很重要的問題,就是教會在倫理道德上的破產和管教的方法.這章聖經或許很難被今天的教會接納.我想有兩個主要的原因:
\newline
\newline
(一)是一班弟兄姊妹或教會領袖,對於神的家和教會的真理,很可能有偏差
\newline
\newline
教會是屬靈的家庭,神賦與一班屬靈的領袖有從神而來的權柄；這權柄一方面是教導信徒在真理上有長進,而另一方面是維持教會應有的屬靈次序,及神在教會中應有的地位和美麗.因此當有弟兄姊妹走錯路時,作為屬靈的領袖,便應該用愛心將他們挽回.倘若他們固執不聽從,違反了神在教會中所定聖潔的標準；以致他們的生活,影響了教會的生活,保羅說我們便要以愛心和勇氣來管教這些在罪中的肢體.
\newline
\newline
(二)教會不接納這段聖經的原因,是今天很多教會,不注重神聖潔的絕對標準
\newline
\newline
我們在處理犯罪弟兄的事情上,往往以愛心或同情或忍耐為理由,等待他們自動悔改,這是錯誤的.因為神將我們從罪惡過犯中拯救出來以後,神有一個絕對聖潔的標準給他的兒女,這聖潔生活的標準,其實是為了造福我們的人生；讓我們循著聖潔標準的軌道,可以享受人生最豐盛的意義.因此,教會必須在神所賜這聖潔的標準下,維持應該有的聖潔生活.假若有人違反了這聖潔的標準,我們便不應容許這事的存在,乃要用愛心去停止這違反的事.信徒若貶低了聖潔的標準,便會造成教會不敢執行屬靈管教這重要的責任.今天我們看見有人在教會公開犯罪,但卻看不見教會用神所給的權柄,去管教這犯罪的肢體.我們很少看見這樣的事,這不是說教會沒有犯罪的人,只是教會缺乏了在這方面的認識.
\newline
\newline
保羅在這裡所講的犯罪,並不是指日常生活的軟弱,或是偶然犯罪而立即悔改的那些罪.保羅所說的,乃是有弟兄姊妹繼續不斷的在同樣的罪中生活.當我們查考這章聖經時,我們可以看見三方面的真理.
\newline
\newline
(A)我們看見教會在道德倫理上破產的嚴重性
\newline
\newline
在哥林多教會中,竟然有人娶了他的繼母,這顯然是違背了神的說話,神絕對不容許他子民中有這些事存在,這是利未記十八節所明文禁止的.而經文所提到的繼母,並不是教會的肢體.因為保羅在書信中所提的管教,只是針對娶繼母的人,而沒有管教繼母,可見這繼母並不在教會所管教的範圍內.保羅為了這亂倫的事非常傷心,因為連外邦人也不會這樣做.
\newline
\newline
這件事所帶來的嚴重性,有兩方面是值得我們注意的.
\newline
\newline
1.這是一件已經犯了,且是繼續犯下去的罪.經文說「收了」,是表示經過一個很長的程序；越過試探和意圖,而實際做了出來的一件事.保羅寫這封信時,這弟兄不單是已經娶了他的繼母,還繼續和他的繼母一起居住.一個信主的人,在道德倫理上竟然可以落到這樣敗壞的地步,這給我們一個很重要的警惕!我們雖然是一個重生得救的信徒,但我們在重生的生命中,依然可以有一個很大的可能性,陷在一個連外邦人都不會犯的罪惡中.因此,我們在聖潔的追求中,不應有絲毫的放鬆,讓撒旦留地步,或是製迼機會讓自己受誘惑.我們絕對不能相信自己不犯罪,我們必須時刻依賴聖靈,過得勝的生活.
\newline
\newline
2.保羅指出這件事的嚴重,就是當這件事發生之後,教會的人竟然不以為然.第二節說:「你們還是自高自大,並不哀痛,把行這事的人,從你們中間趕出去.」教會大家都知道這件事,可是大家都若無其事,處之泰然；認為它並不是那麼嚴重,教會並不需要採取甚麼行動,保羅說他們連哀痛也沒有.對罪的妥協,是一般教會工作不能長進一個致命的打擊.為甚麼教會沒有人對這事表哀痛呢?很可能他們將聖潔的標準降低,認為人皆會犯罪,我們應該同情和包容他們,這樣他們才不會離開教會.又或者一班領導教會的人,他們在自己的生活中也有很多敗壞；所以覺得沒有屬靈的資格,去管教這犯罪的弟兄.因此整個教會便對這事視若無睹,並不哀痛.
\newline
\newline
(B)就是倫理道德破產,教會應該有管教的責任和標準.
\newline
\newline
(五3-5):「我身子雖不在你們那裡,心卻在你們那裡,好像我親自與你們同在,已經判斷了行這事的人.就是你們聚會的時候,我的心也同在,奉我們主耶穌的名,並且我們主耶穌的權能；要把這樣的人交給撒旦,壞他的肉體,使他的靈魂,在主耶穌的日子,可以得救.」
\newline
\newline
首先我們要明白管教的意義,和管教最終的目的.保羅說要把這樣的人交給撒旦,敗壞他的肉體,這就是說,要斷絕這犯罪的肢體在教會中一切的生活；讓他不再屬於耶穌基督的團體內,而最終的目的是要使他的靈魂在主再來的時候可以得救,意思就是要挽回這弟兄.因為一個重生得救的人,當他接受管教,被逐出教會,他必然會自我反省,藉著聖靈的感動和提醒,至終他必然會真心悔改,重返教會.所以管教的目的,並不是把犯罪的肢體越逐越遠；而是藉著一個嚴厲的管教方法,讓他能悔改,重新復興.
\newline
\newline
至於管教的方法,是以耶穌基督為中心,用主的權柄,奉主的名去處理.我們不是用自己的權柄和標準；乃是以主為中心,我們以悲痛的心去管教弟兄,最終的目的是希望他們回轉.有關管教的程序,我們可以參閱(太十八15-17):「倘若你的弟兄得罪你,你就去趁著只有他和你在一處的時候,指出他的錯來.他若聽你,你便得了你的弟兄.他若不聽,你就另外帶一兩個人同去,要憑兩三人的口作見證,句句都可定準.若是不聽他們,就告訴教會,若是不聽教會,就看他像外邦人和稅吏一樣.」
\newline
\newline
教會應該負起管教的責任,教會若不管教,任由弟兄姊妹在罪中生活；其結果是可怕的,不但犯罪的肢體執迷不悟,而且更嚴重的；是叫全教會,本來是乾淨的麵都發起來.意思就是說:如果教會不堅守聖潔的標準,容讓信徒在罪中生活；那麼教會聖潔的標準便失去,教會的生活便受到很大的影響.
\newline
\newline
(C)我們從這章聖經所看見的,就是倫理道德破產而受管教之後,屬靈得釋放的秘訣
\newline
\newline
(七7-13):「你們既是無酵的麵,應當把舊酵除淨,好使你們成為新團,因為我們這逾越節的羔羊基督,已經被殺獻祭了.所以我們守這節,不可用舊酵,也不可用惡毒邪惡的酵,只用誠實真正的無酵餅.我先前寫信給你們說,不可與淫亂的人相交,這話不是指這世上一概行淫亂的,或貪婪的,勒索的,或拜偶像的.若是這樣,你們除非離開世界方可；但如今我寫信給你們說,若有稱為弟兄,是行淫亂的,或貪婪的,或拜偶像的,或辱罵的,或醉酒的,或勒索的；這樣的人不可與他相交,就是與他吃飯也不可.因為審判教外的人與我何干,教內的人豈不是你們審判的麼?至於外人有神審判他們,你們應當把那惡人從你們中間趕出去.」
\newline
\newline
在屬靈的釋放中,首先我們得到的,是一個聖潔的教會.因為教會肯執行屬靈的管教,肯堅守神聖潔的標準,這樣便成為神眼中聖潔的教會.另一方面,這管教在屬靈的釋放中,是對犯罪者一個真正愛心的輔助.愛心是有責任的；有責任的愛心,才是真正的愛心,而愛心愈大,責任便愈大.
\newline
\newline
最後,或許你會問,甚麼罪才需要接受教會的管教.關於這個,我很難一一細說.不過,我們可以根據一個原則,就是該件罪對某肢體來說,是他長期不肯悔改；影響他自己的生活,又影響到教會整個生活的時候；教會可以視乎這罪惡的嚴重性,作出管教上的責任.神是聖潔的神.神為人預備救贖,目的也是要把人從罪惡中拯救出來；但願我們能體會主的心腸,堅守神聖潔的原則；除去一切的污穢,使教會能在神面前保持聖潔.
\newline
\newline


\newpage

\section{第55屆~港九培靈會~張子華~第六講　基督徒的獨身和婚姻}
\label{sec:514}
\textbf{第六講　基督徒的獨身和婚姻}
\newline
\newline
連結: \href{https://www.hkbibleconference.org/session-message/view/514}{\texttt{https://www.hkbibleconference.org/session-message/view/514}}
\newline
\newline
\hyperref[sec:513]{< < < PREV SERMON < < <}
~
\hyperref[sec:toc]{[返目錄]}
~
\hyperref[sec:515]{> > > NEXT SERMON > > >}
\newline
\newline
哥林多前書 7:1-7:40
\newline
\begin{longtable}{cl}
\hline
\hline
章節 & 經文 (和合本修訂版)\\
\hline
7:1 & \begin{tabularx}{0.7\textwidth}{X} 關於你們信上所提的事，男人不親近女人倒好。 \end{tabularx} \\ \\ \relax
7:2 & \begin{tabularx}{0.7\textwidth}{X} 但為了避免淫亂的事，男人當各有自己的妻子，女人也當各有自己的丈夫。 \end{tabularx} \\ \\ \relax
7:3 & \begin{tabularx}{0.7\textwidth}{X} 丈夫對妻子要盡本分；妻子對丈夫也要如此。 \end{tabularx} \\ \\ \relax
7:4 & \begin{tabularx}{0.7\textwidth}{X} 妻子對自己的身體沒有主張的權柄，權柄在丈夫；丈夫對自己的身體也沒有主張的權柄，權柄在妻子。 \end{tabularx} \\ \\ \relax
7:5 & \begin{tabularx}{0.7\textwidth}{X} 夫妻不可忽略對方的需求，除非為了要專心禱告，在兩相情願下暫時分房；以後仍要同房，免得撒但趁著你們情不自禁而引誘你們。 \end{tabularx} \\ \\ \relax
7:6 & \begin{tabularx}{0.7\textwidth}{X} 我說這話是出於容忍，不是命令。 \end{tabularx} \\ \\ \relax
7:7 & \begin{tabularx}{0.7\textwidth}{X} 我願眾人像我一樣；但是各人都有來自神的恩賜，一個是這樣，一個是那樣。 \end{tabularx} \\ \\ \relax
7:8 & \begin{tabularx}{0.7\textwidth}{X} 我對沒有嫁娶的和寡婦說，他們若能維持獨身像我一樣就好。 \end{tabularx} \\ \\ \relax
7:9 & \begin{tabularx}{0.7\textwidth}{X} 但他們若不能自制，就應該嫁娶，與其慾火攻心，倒不如結婚為妙。 \end{tabularx} \\ \\ \relax
7:10 & \begin{tabularx}{0.7\textwidth}{X} 至於那已經嫁娶的，我吩咐他們—其實不是我，而是主吩咐的：妻子不可離開丈夫， \end{tabularx} \\ \\ \relax
7:11 & \begin{tabularx}{0.7\textwidth}{X} 若是離開了，不可再嫁，不然要跟丈夫復和；丈夫也不可離棄妻子。 \end{tabularx} \\ \\ \relax
7:12 & \begin{tabularx}{0.7\textwidth}{X} 我對其餘的人說—是我，不是主說—倘若某弟兄有不信的妻子，妻子也情願和他一起生活，他就不可離棄妻子。 \end{tabularx} \\ \\ \relax
7:13 & \begin{tabularx}{0.7\textwidth}{X} 妻子有不信的丈夫，丈夫也情願和她一起生活，她就不可離棄丈夫。 \end{tabularx} \\ \\ \relax
7:14 & \begin{tabularx}{0.7\textwidth}{X} 因為不信的丈夫會因著妻子成了聖潔；不信的妻子也會因著丈夫成了聖潔。不然，你們的兒女就不潔淨了，但現在他們是聖潔的。 \end{tabularx} \\ \\ \relax
7:15 & \begin{tabularx}{0.7\textwidth}{X} 倘若那不信的人要離開，就由他離開吧！無論是弟兄是姊妹，遇著這樣的事都不必拘束。神召你們原是要你們和睦。 \end{tabularx} \\ \\ \relax
7:16 & \begin{tabularx}{0.7\textwidth}{X} 你這作妻子的怎麼知道不能救你的丈夫呢？你這作丈夫的怎麼知道不能救你的妻子呢？ \end{tabularx} \\ \\ \relax
7:17 & \begin{tabularx}{0.7\textwidth}{X} 無論如何，要照主所分給各人的恩賜和神所召各人的情況生活。我在各教會裡都是這樣規定的。 \end{tabularx} \\ \\ \relax
7:18 & \begin{tabularx}{0.7\textwidth}{X} 有人受割禮後才蒙召，他就不必除去割禮的記號。有人未受割禮前蒙召，他就不必受割禮。 \end{tabularx} \\ \\ \relax
7:19 & \begin{tabularx}{0.7\textwidth}{X} 受割禮算不了甚麼，不受割禮也算不了甚麼，只要謹守神的誡命就是了。 \end{tabularx} \\ \\ \relax
7:20 & \begin{tabularx}{0.7\textwidth}{X} 各人蒙召的時候是甚麼身份，要守住這身份。 \end{tabularx} \\ \\ \relax
7:21 & \begin{tabularx}{0.7\textwidth}{X} 你是作奴隸時蒙召的嗎？不要介意；若能獲得自由，就爭取自由更好。 \end{tabularx} \\ \\ \relax
7:22 & \begin{tabularx}{0.7\textwidth}{X} 因為，蒙主呼召的奴僕是主所釋放的人；蒙主呼召的自由之人是基督的奴僕。 \end{tabularx} \\ \\ \relax
7:23 & \begin{tabularx}{0.7\textwidth}{X} 你們是重價買來的；不要作人的奴僕。 \end{tabularx} \\ \\ \relax
7:24 & \begin{tabularx}{0.7\textwidth}{X} 弟兄們，你們各人蒙召的時候是甚麼身份，要在神面前守住這身份。 \end{tabularx} \\ \\ \relax
7:25 & \begin{tabularx}{0.7\textwidth}{X} 關於未婚女子，我沒有主的命令，但我既蒙主憐憫、作為一個可信靠的人，把自己的意見告訴你們。 \end{tabularx} \\ \\ \relax
7:26 & \begin{tabularx}{0.7\textwidth}{X} 因現今的艱難，據我看來，人不如安於現狀。 \end{tabularx} \\ \\ \relax
7:27 & \begin{tabularx}{0.7\textwidth}{X} 你已經有了妻子，就不要求擺脫；你還沒有妻子，就不要想娶妻。 \end{tabularx} \\ \\ \relax
7:28 & \begin{tabularx}{0.7\textwidth}{X} 你若娶妻，並不是犯罪；未婚女子若出嫁，也不是犯罪。然而，這等人會遭受肉身上的苦難，我寧願你們免受這苦難。 \end{tabularx} \\ \\ \relax
7:29 & \begin{tabularx}{0.7\textwidth}{X} 弟兄們，我是說：時候不多了。從此以後，那有妻子的，要像沒有一樣； \end{tabularx} \\ \\ \relax
7:30 & \begin{tabularx}{0.7\textwidth}{X} 哀哭的，不像在哀哭；快樂的，不像在快樂；購買的，像一無所得； \end{tabularx} \\ \\ \relax
7:31 & \begin{tabularx}{0.7\textwidth}{X} 享受這世界的，不像在享受這世界；因為這世界的局面將要過去了。 \end{tabularx} \\ \\ \relax
7:32 & \begin{tabularx}{0.7\textwidth}{X} 我願你們一無掛慮。沒有結婚的是為主的事掛慮，想怎樣令主喜悅； \end{tabularx} \\ \\ \relax
7:33 & \begin{tabularx}{0.7\textwidth}{X} 結了婚的是為世上的事掛慮，想怎樣讓妻子喜悅， \end{tabularx} \\ \\ \relax
7:34 & \begin{tabularx}{0.7\textwidth}{X} 於是，他就分心了。沒有結婚的和未婚的女子是為主的事掛慮，為要身體和心靈都聖潔；已經出嫁的是為世上的事掛慮，想怎樣讓丈夫喜悅。 \end{tabularx} \\ \\ \relax
7:35 & \begin{tabularx}{0.7\textwidth}{X} 我說這話是為你們的益處，不是要限制你們，而是要你們做合宜的事，得以不分心地對主忠誠。 \end{tabularx} \\ \\ \relax
7:36 & \begin{tabularx}{0.7\textwidth}{X} 若有人認為自己待他的女兒不合宜，女兒也過了適婚年齡，他可以隨意處理，不算有罪，讓兩人結婚就是了。 \end{tabularx} \\ \\ \relax
7:37 & \begin{tabularx}{0.7\textwidth}{X} 倘若有人心裡堅定，沒有不得已的事，並且由得自己作主，心裡又決定了不讓女兒結婚，這樣做也好。 \end{tabularx} \\ \\ \relax
7:38 & \begin{tabularx}{0.7\textwidth}{X} 這樣看來，讓自己的女兒結婚固然是好，不讓她結婚更好。 \end{tabularx} \\ \\ \relax
7:39 & \begin{tabularx}{0.7\textwidth}{X} 丈夫活著的時候，妻子是受約束的；丈夫若長眠了，妻子就自由了，可以隨意再嫁，只是要嫁給主裡面的人。 \end{tabularx} \\ \\ \relax
7:40 & \begin{tabularx}{0.7\textwidth}{X} 然而，按我的意見，她若能守節就更有福氣。我想我自己也有神的靈的感動。 \end{tabularx} \\ \\
[1ex]
\hline
\hline
\end{longtable}
今天我們要查考哥林多前書第七章,這是一章很寶貴的聖經.起首的十一節是論到夫妻生活的模式,之後保羅便講到男女間所有的一些問題,譬如說寡居,獨身等.首先我們要思想的,是獨身的問題,就是二十五節之後的一段,論守童身的真理.
\newline
\newline
由二十五節一開始,保羅便說,論到獨身的人,我沒有主的命令.有人以為保羅自己既然是個獨身的人,所以便照自己的感受講獨身的道理,而這些道理不是主吩咐他講的.不過,這節經文的意思並不是這樣,保羅說他沒有主的命令；意思就是聖經並沒有記載過這樣的真理.舊約的聖經,有很多是提到婚姻的,但卻沒有觸及獨身；保羅是頭一個講到這問題的人,所以他並沒有可引用的話.不過他說把自己的意見告訴你們,其實他是在聖靈的感動和引領下講的；所以他所講的,其實就是神的心意.
\newline
\newline
在二十五到四十節這一段經文中,保羅要一班未婚男女對獨身有兩種看法.
\newline
\newline
1.凡是未婚的男女,絕不能視婚姻作人生最大的目標.保羅說:你沒有妻子纏著呢?就不要求妻子.保羅所指的求,不是禱告祈求,而是追求；意思就是說,為自己的終身大事禱告是好的,是合神心意的；但切不要以追求婚姻為人生最大的目標.如果結了婚,就不要再求自由,乃要在婚姻的約中生活.但若是未結婚的,就不要以結婚為人生最大的目標,因為人生還有許多別的事比結婚重要.
\newline
\newline
基督徒必須確定甚麼是我們生命的首要,這是非常重要的.保羅在這段經文中,很強調一件事實,就是我們生活在一個主快再來的日子中；因此我們不能讓任何事,來削減我們在主再來前所應該作的.假若我們以結婚為人生的大問題,很多時候,結婚便會成為我們成就神旨意的一個攔阻.所以,如果我們以遵行神的旨意為人生的目標,這是最好不過的.
\newline
\newline
而且,未結婚的人,一定要先明白結婚的真實性,對婚姻有實在的看法.一個結了婚的人,便必須在婚姻的約中生活,而這種生活會帶來很多的難處；要付上很大的代價,甚至要在某些事情上有所放棄；而且更要受一些肉身上的苦難,這些苦難是未結婚的人沒有的.婚姻的成功是要付很大代價,未婚的人必須先有所認識.為此,婚姻並不是人生所要追求的最大目標.
\newline
\newline
2.就是對婚姻要有一顆謹慎和等候神的心.(36-38節)說:「若有人以為自己待他的女兒不合宜,女兒也過了年歲,事又當行；他就可隨意辦理,不算有罪,叫二人成親就是了.倘若人心裡堅定,沒有不得已事,並且由得自己作主；心裡又決定了留下女兒不出嫁,如此行也好.這樣看來,叫自己的女兒出嫁是好,不叫她出嫁更是好.」當然這段經文是以母親為主角,但保羅勸為母的不可催女兒出嫁；所以,在獨身的時候,無論是父母或兒女,都應當以謹慎的態度等候神.
\newline
\newline
現在我們來思想基督徒的婚姻,請特別注意第二節,保羅說:「但要免淫亂事,男子當各有自己的妻子,女子也當各有自己丈夫.」我們要特別留意,「自己」一詞,當某弟兄看見某姊妹的時候,內心能說這是我的.我們看亞當怎樣娶夏娃,夏娃不是亞當找來的；乃是神看見他的需要而賜給他的.神使亞當沉睡,取下他一條肋骨來做夏娃；當亞當醒來時,看見了夏娃,便說這是我骨中的骨,肉中的肉.保羅說各人當有自己的配偶,他的意思就是這個.婚姻是神為人所預備,是由神來配合的.不過,保羅在這裡並沒有提到一些實際的方法去尋求和引證神的帶領.
\newline
\newline
但我願意提供給大家幾個簡單的引證:
\newline
\newline
(一)屬靈的引證.聖經說二人必須同負一軛,這不單是指信與不信的問題:乃是指在屬靈的生命上,能同負一軛,同走天路.
\newline
\newline
(二)實際的引證.譬如說教育程度,家庭背景,對金錢的價值觀,個性,興趣等,都是婚前彼此應該明白的.
\newline
\newline
(三)屬靈長者的引證.每一對將要結婚的青年男女,都應該找一位屬靈長者來輔導；並給予客觀的分析.屬靈長者的引證是非常重要的.
\newline
\newline
(四)你必須覺得對方有吸引力,是你心目中最美麗的.特別是女子,她必須有戀慕對方的心,要佩服他,樂意順從他.經過這個引證之後,男女二人才心悅誠服的在神面前結為夫婦；並深信對方是神預備的,且樂意毫無保留的將自己交給對方.
\newline
\newline
結婚之後,夫婦要注重三個重要的關係,這是一到十一節所提到的.
\newline
\newline
(一)是純潔或貞潔的關係.第二節:「但要免淫亂的事,男子當各有自己的妻子,女子也當各有自己的丈夫.」這代表夫婦關係的單純,貞潔.除了和自己的丈夫或妻子有性關係之外,絕不能再和別的異性有性關係；他是我唯一的配偶,這關係是純潔的.在純潔的關係中,含有無比的愛情；因為對方是自己的,所以便會珍惜,保護,不容他受損害.
\newline
\newline
(二)是伴侶的關係.三到五節:「丈夫當用合宜之分待妻子,妻子待丈夫也要如此.妻子沒有權柄主張自己的身子,乃在丈夫；丈夫也沒有權柄主張自己的身子,乃在妻子.夫妻不可彼此虧負,除非兩相情願,暫時分房；為專心禱告方可；以後仍要同房,免得撒旦趁著你們情不自禁,引誘你們.」人生最親密的伴侶,就是丈夫或妻子,因此要學習彼此體恤.對待原來是並沒有施恩的意思,乃是責任和義務的付出.夫婦二人本來是單獨的個體,各有自己的權利,嗜好,可是二人共同生活之後；他們的權利,嗜好,因為願意在責任和義務上付出代價,以致他們不在施恩的態度,而過一個和諧的生活.夫婦毫無衝突的相處,是高度的藝術,也是偉大愛情最重要的探討.
\newline
\newline
除了彼此體恤之外,夫婦間還需要有高度的性生活享受；絕不能使用對方去滿足自己的需要；而是在夫妻關係中,追求彼此都有滿足的享受.在主的祝福,在聖靈的引證,在寶血的潔淨下；夫婦的性關係,是滿足了神在夫婦的責任上其中一個主要的目的.在心靈成為一體下,我們便能在肉體的關係上,有滿足而有意義的享受.
\newline
\newline
保羅在第五節裡提到三個原則,指出夫婦在甚麼情況下可以停止性生活.第一是兩相情願.第二是暫時性,過渡性的.第三是為了人生一些重要的事,或是屬於主的事禱告.在這三個原則之下,夫婦可以分房,但這絕不是離婚的前奏.因此,分居是為了尋求神的旨意,而不是為了離開對方.
\newline
\newline
(三)是永久的關係.第十節:「至於那已經嫁娶的,我吩咐他們,其實不是我吩咐,乃是主吩咐說,妻子不可離開丈夫.」基督徒的婚姻,是永久性的.這永久的婚姻是基於兩個基礎,第一是基於神的律法.神是這麼吩咐,是這樣說的.保羅在這裡引用到基督所講的說話,是引用亞當和夏娃的關係,神所配合的,人不能分開.第二個基礎,是用到基督所講的說話,是引用亞當和夏娃的關係,神所配合的,人不能分開.第二個基礎,是基於神的愛.十一節提到和好,和好是聖經中一個重要的字,是基於神的愛情.當我們還作罪人的時候,神便叫祂的兒子為我們而死；藉著基督的救贖,我們得與神和好,這是個愛情.這愛藉著聖靈居住在我們裡面,因此婚姻是永久性的.每一對基督徒夫婦,他們的生命中有神的愛,故此沒有難處是神的愛所不能解決的；沒有裂痕是神的愛所不能彌補的.當我們堅決不要和好時,我們是貶低了基督徒婚姻的原則,而且是得罪了無辜的孩子.如果我們明白神的愛是永不改變,而且永不失敗；又把神的愛應用在婚姻上,那麼婚姻中一切的破裂,難處都會得著彌補.成功的婚姻,是雙方都必須付上無上的代價,是努力建立的結果.
\newline
\newline


\newpage

\section{第55屆~港九培靈會~張子華~第七講　基督徒在恩典中所享受的自由}
\label{sec:515}
\textbf{第七講　基督徒在恩典中所享受的自由}
\newline
\newline
連結: \href{https://www.hkbibleconference.org/session-message/view/515}{\texttt{https://www.hkbibleconference.org/session-message/view/515}}
\newline
\newline
\hyperref[sec:514]{< < < PREV SERMON < < <}
~
\hyperref[sec:toc]{[返目錄]}
~
\hyperref[sec:516]{> > > NEXT SERMON > > >}
\newline
\newline
哥林多前書 10:1-10:13
\newline
\begin{longtable}{cl}
\hline
\hline
章節 & 經文 (和合本修訂版)\\
\hline
10:1 & \begin{tabularx}{0.7\textwidth}{X} 弟兄們，我不願意你們不知道，我們的祖宗從前都在雲下，都從海中經過， \end{tabularx} \\ \\ \relax
10:2 & \begin{tabularx}{0.7\textwidth}{X} 都在雲裡、海裡受洗歸了摩西， \end{tabularx} \\ \\ \relax
10:3 & \begin{tabularx}{0.7\textwidth}{X} 並且都吃了一樣的靈糧， \end{tabularx} \\ \\ \relax
10:4 & \begin{tabularx}{0.7\textwidth}{X} 也都喝了一樣的靈水，所喝的是出於跟隨著他們的靈磐石；那磐石就是基督。 \end{tabularx} \\ \\ \relax
10:5 & \begin{tabularx}{0.7\textwidth}{X} 但他們中間多半是神不喜歡的人，所以倒斃在曠野裡了。 \end{tabularx} \\ \\ \relax
10:6 & \begin{tabularx}{0.7\textwidth}{X} 這些事都是我們的鑒戒，使我們不要貪戀惡事，像他們貪戀過的一樣。 \end{tabularx} \\ \\ \relax
10:7 & \begin{tabularx}{0.7\textwidth}{X} 也不要拜偶像，像他們中有些人曾經拜過。如經上所記：「百姓坐下吃喝，起來玩樂。」 \end{tabularx} \\ \\ \relax
10:8 & \begin{tabularx}{0.7\textwidth}{X} 我們也不可犯姦淫，像他們中有些人曾經犯過，一天就倒斃了二萬三千人。 \end{tabularx} \\ \\ \relax
10:9 & \begin{tabularx}{0.7\textwidth}{X} 也不可試探主，像他們中有些人曾試探主就被蛇咬死。 \end{tabularx} \\ \\ \relax
10:10 & \begin{tabularx}{0.7\textwidth}{X} 你們也不可發怨言，像他們中有些人曾經發過，就被毀滅者所滅。 \end{tabularx} \\ \\ \relax
10:11 & \begin{tabularx}{0.7\textwidth}{X} 這些事發生在他們身上，要作為鑒戒，而且寫下來正是要警戒我們這末世的人。 \end{tabularx} \\ \\ \relax
10:12 & \begin{tabularx}{0.7\textwidth}{X} 所以，自以為站得穩的人必須謹慎，免得跌倒。 \end{tabularx} \\ \\ \relax
10:13 & \begin{tabularx}{0.7\textwidth}{X} 你們所受的考驗無非是人所承受得了的。神是信實的，他不會讓你們遭受無法承受的考驗，在受考驗的時候，總會給你們開一條出路，讓你們能忍受得了。 \end{tabularx} \\ \\
[1ex]
\hline
\hline
\end{longtable}
今天我們要思想的經文,是(林前十1-13):「弟兄們,我不願意你們不曉得,我們的祖宗從前都在雲下,都從海中經過,都在雲裡,海裡受洗歸了摩西；並且都吃了一樣的靈食,也都喝了一樣的靈水；所喝的是出於隨著他們的磐石,那磐石就是基督.但他們中間,多半是神不喜歡的人,所以在曠野倒斃.這些事都是我們的鑑戒,叫我們不要貪戀惡事,像他們那樣貪戀的；也不要拜偶像,像他們有人拜的；如經上所記,百姓坐下吃喝,起來玩耍.我們也不要行姦淫,像他們有人行的,一天就倒斃了二萬三千人.也不要試探主,像他們有人試探的,就被蛇所滅.你們不要發怨言,像他們有發怨言的,就被滅命的所滅.他們遭遇這些事,都要作為鑑戒,並且寫在經上,正是警戒我們這末世的人.所以自己以為站得穩的,須要謹慎,免得跌倒.你們所遇見的試探,無非是人所能受的.神是信實的,必不叫你們受試探過於所能受的；在受試探的時候,總要給你們開一條出路,叫你們能忍受得住.」
\newline
\newline
保羅在這段經文中,講到基督徒在恩典中的自由生活,而這種自由的生活,很明顯是水禮的生活.保羅在這段經文的開始,給我們對水禮一個很榮耀的觀念.有人以為水禮算不得甚麼,不過是一種記號,可有可無；但在聖經的教訓中,講到水禮雖然只是一個外表的記號,但這記號卻能代表我們從一種生活進入另一種生活中.所以水禮是代表一種生活方式,而這種方式是我們與主基督聯合,是一種榮耀,滿有享受的生活方式.我們的自由,是因為耶穌基督把我們從罪中釋放出來,讓我們在基督裡得著自由.藉著洗禮,我們能進入一種新的生活方式去.
\newline
\newline
我們先來看這種生活所帶給我們四個的生活享受.
\newline
\newline
(一)是一個蒙帶領的生活
\newline
\newline
第一節說都在雲下,雲柱是代表神的使者,神的同在,神的帶領.當以色列人離開埃及,進入曠野之後,神使用雲柱和火柱去帶領他們度四十年的曠野生活.基督徒的生活,每一天都是在神的帶領眷顧下度過的,不會迷失方向,不會惆悵,因為神話語的教導,肢體的生活,牧者的教導,教會種種的生活,都讓我們時刻活在神的帶領中.
\newline
\newline
(二)天天過著蒙拯救的生活
\newline
\newline
「都從海中經過」,這是指以色列人過紅海的事.當以色列人在紅海的邊緣,前無去路,後有追兵的時候,神行奇事使紅海的水分開,讓他們下海走乾地；到埃及人追入紅海時,神又使水復合,把他們全都淹死.因此從海中經過是一個蒙拯救的生活.基督徒所過的,是釋放的生活.埃及在我們的生命中,再不能有任何權勢；而且,我們所得的救恩,不單是過去式的,並且是現在式和將來式的.那就是說,基督過去為我們所成就的救贖,使我們得著拯救；而這救恩的功效,是使我們每一天都得著釋放；並且在將來主再來的時候,我們可以得著永遠完全的救贖,這是基督徒人生高度的享受.
\newline
\newline
(三)是一個整體性的生活
\newline
\newline
注意第一節的「都」字,這都字在頭幾節經文中已經出現了五次.基督徒所享受的生活,是整體性的生活.我們有一群弟兄姊妹,能夠一起生活,一起交通,彼此關心,共同追求神的話,分享心中的喜樂,分擔生活的難處,這是整體性的基督徒生活.
\newline
\newline
(四)是一個蒙供應的生活
\newline
\newline
以色列人在曠野中,得神的供應,有水,有肉,有嗎哪；他們每天的生活,都是在倚靠神中蒙神在物質上豐富的供應.同樣,今天我們也在每天的生活中,體驗到神在各方面豐富的供應.
\newline
\newline
然而,在美好的生活享受中,我們來注意第五節:「但他們中間,多半是神不喜歡的人,所以在曠野倒斃.」這是非常可惜的事!一班日夜享受神恩典的人,竟然多半是神不喜歡；這對我們來說,是何等嚴重的警告.今天我們活在神恩典中,但我們是神所喜歡的人嗎?
\newline
\newline
神不喜歡這班人,有五個重要的原因.
\newline
\newline
1.因為他們放縱私慾,容許私慾的試探勝過他們.第六節說他們貪戀惡事,貪戀是一個私慾.當以色列人在曠野生活不久之後,他們便厭煩曠野生活的單調,枯燥,清苦,他們回想埃及時候的生活,很想回去.這班人忘了他們在主恩典中所享受的,是遠超過埃及的物質所能帶給他們的.同樣,今天有不少基督徒,起初信主的時候很興奮,很喜樂,但慢慢便感到厭倦,甚至有回到世界去的念頭.聖經說,這樣的人,是神所不喜歡的.
\newline
\newline
2.因為他們拜偶像.這明顯是指到當摩西到西乃山取十誡,遲遲未下山時百姓所作的事,他們為自己鑄造了金牛犢,向他下拜說,這就是領我們出埃及的神.基督徒在生活中若缺少了對神的享受,很容易我們便在膚淺的認識中,構成一個我們自己所喜歡的信仰標準；我們認為這就是我所信的耶穌,但其實耶穌並不是這樣的.我們所信的,是一個要求付代價,背十字架的信仰；然而這樣的付出卻帶我們進入一個豐盛的生命中.如果我們對此沒有清楚的認識,我們便很容易把神所給的標準降低,認為這就是我們所信的.我們用自己的標準生活,而忽略神的要求和標準.當我們這樣認為時,這就是我們的金牛犢,是神所不喜歡的.
\newline
\newline
3.是婚姻的試探.第八節說:「我們也不要行姦淫,像他們有人行的,一天就倒斃了二萬三千人.」這節經文所指的,是民數記第二十五章所記載的.神要求以色列的男子娶以色列的女子,但是他們卻愛上摩押的女子,和他們行姦淫的事,惹神發怒,結果神在一日之內,擊殺了他們二萬三千人.今天,神要我們在婚姻上謹慎,絕不能在婚姻上隨便而跌倒,神要我們娶一個能和我們共負一軛的對像；雖然外表可能不漂亮,不夠新潮,但是卻是愈看愈美的人.其實,真正的美麗是因為愛神,有裡面的裝飾,有從基督裡所產生出來的溫柔.婚姻的跌倒,導致神擊殺他們二萬三千人,我們不可不小心,絕不能和世人看齊,不要落在婚姻的試探和迷惑中.
\newline
\newline
4.他們落在不信主的試探中.第九節:「也不要試探主,像他們有人試探的,就被蛇所滅.」這是指以色列人在曠野生活的時候,多次及多方面對神的不信.當摩西派十二個探子去窺視迦南時；十個探子回來報惡信,認為不可上去攻取那地,他們不信神的應許；結果從埃及出來的人,除了迦勒和約書亞之外,其餘的都不得進入迦南.他們雖然對神有很多經歷,但仍多次向神表示不信,這令神不喜歡.我們有沒有這樣的不信神呢?神明說他不喜歡這樣的人.神要我們堅信祂到底,堅信他的應許,堅信祂所說過的話,堅信祂是創始成終的神.但願我們在以後人生的路程上,不會再對神投不信任的票.
\newline
\newline
5.是埋怨的試探.第十節:「你們也不要發怨言,像他們有發怨言的,就被滅命的所滅.」保羅在腓立比書告訴我們,凡事都不要發怨言,起爭論.怨言和爭論是一樣的.怨言是內心的光景,而爭論是將這心態表現出來.保羅為甚麼不要信徒發怨言呢?或者說神為甚麼不要我們發怨言呢?因為神要我們在祂裡面過著一個滿足的生活,這樣我們在內心便不會發怨言,在外表也不會起爭論.有人常發怨言,諸多不滿,事事批評,挑剔,這是神不喜歡的,這也表示我們從未在主裡享受過滿足的生活.舊約時代神擊殺那些常發怨言的以色列人,這是我們的鑑戒.怨言能阻止我們得著神的恩典,除非我們停止了一切內心的怨言,外表的爭論；才能進入以基督為滿足的生活裡,得著無上的享受.
\newline
\newline


\newpage

\section{第55屆~港九培靈會~張子華~第八講　基督的自由和社交生活}
\label{sec:516}
\textbf{第八講　基督的自由和社交生活}
\newline
\newline
連結: \href{https://www.hkbibleconference.org/session-message/view/516}{\texttt{https://www.hkbibleconference.org/session-message/view/516}}
\newline
\newline
\hyperref[sec:515]{< < < PREV SERMON < < <}
~
\hyperref[sec:toc]{[返目錄]}
~
\hyperref[sec:517]{> > > NEXT SERMON > > >}
\newline
\newline
哥林多前書 8:1-8:13
\newline
\begin{longtable}{cl}
\hline
\hline
章節 & 經文 (和合本修訂版)\\
\hline
8:1 & \begin{tabularx}{0.7\textwidth}{X} 關於祭過偶像的食物，我們曉得「我們都有知識」，但知識使人自高自大，惟有愛心能造就人。 \end{tabularx} \\ \\ \relax
8:2 & \begin{tabularx}{0.7\textwidth}{X} 若有人自以為知道甚麼，他其實仍不知道他所應當知道的。 \end{tabularx} \\ \\ \relax
8:3 & \begin{tabularx}{0.7\textwidth}{X} 若有人愛神，他就是神所認識的人了。 \end{tabularx} \\ \\ \relax
8:4 & \begin{tabularx}{0.7\textwidth}{X} 關於吃祭過偶像的食物，我們知道「偶像在世上算不得甚麼」；也知道「神只有一位，沒有別的」。 \end{tabularx} \\ \\ \relax
8:5 & \begin{tabularx}{0.7\textwidth}{X} 雖然在天上或地上有許多所謂的神明，就如他們中間有許多的神明，許多的主， \end{tabularx} \\ \\ \relax
8:6 & \begin{tabularx}{0.7\textwidth}{X} 但是我們只有一位神，就是父，萬物都出於他，我們也歸於他；並只有一位主，就是耶穌基督，萬物都是藉著他而有，我們也是藉著他而有。 \end{tabularx} \\ \\ \relax
8:7 & \begin{tabularx}{0.7\textwidth}{X} 可是，不是人人都有這知識。有人到現在因拜慣了偶像，仍以為所吃的是祭過偶像的食物；既然他們的良心軟弱，也就污穢了。 \end{tabularx} \\ \\ \relax
8:8 & \begin{tabularx}{0.7\textwidth}{X} 其實，食物不能使我們更接近神，因為我們不吃也無損，吃也無益。 \end{tabularx} \\ \\ \relax
8:9 & \begin{tabularx}{0.7\textwidth}{X} 可是，你們要謹慎，免得你們這自由竟成了軟弱人的絆腳石。 \end{tabularx} \\ \\ \relax
8:10 & \begin{tabularx}{0.7\textwidth}{X} 若有人見你這有知識的在偶像的廟裡坐席，而這人的良心是軟弱的，他豈不放膽去吃那祭過偶像的食物嗎？ \end{tabularx} \\ \\ \relax
8:11 & \begin{tabularx}{0.7\textwidth}{X} 因此，基督為他死的那軟弱弟兄，也就因你的知識沉淪了。 \end{tabularx} \\ \\ \relax
8:12 & \begin{tabularx}{0.7\textwidth}{X} 你們這樣得罪弟兄，傷了他們軟弱的良心，就是得罪基督。 \end{tabularx} \\ \\ \relax
8:13 & \begin{tabularx}{0.7\textwidth}{X} 所以，食物若使我的弟兄跌倒，我就永遠不吃肉，免得使我的弟兄跌倒了。 \end{tabularx} \\ \\
[1ex]
\hline
\hline
\end{longtable}
今天我們要思想的經文,是哥林多前書第八章全章,和第十章十四到三十三節.這兩段經文,是講到吃祭偶像之物,和基督徒的自由.
\newline
\newline
基督徒在主裡面是一個自由的人,因為基督釋放了我們.本來我們是罪的奴僕,但現今卻因著基督的釋放而得了自由.因此,我們可以根據一個原則來過活,如果我們所作的事,不會被聖靈控告良心,那就是說,和聖經的真理沒有衝突,那麼任何的事我都可以做.今天的經文告訴我們,當時哥林多的教會有一個社交的問題,到底基督徒可否吃外邦人拜過偶像的肉.因為按當時的習慣,每一次當哥林多人在廟堂裡拜完神之後,便會在廟內大擺筵席,彼此吃喝.例如喜事,嫁娶,就會到廟堂拜神,之後在廟內擺設婚宴.基督徒在未信主之時也行這樣的事,這一個社交的問題.但他們在信主之後,便發生這樣的問題,到底祭過偶像的物,他們可吃不可吃?
\newline
\newline
雖然這段經文有它的時代背景,今天我們也許不會碰到同一樣的問題,不過我們可以把它的原則應用在我們的社交生活上.譬如說,今天我們參加婚宴,可否飲酒?可否打麻雀?不少基督徒為了這樣的問題而發生爭辯,有人認為可以,有人則認為不可.保羅在這裡便給我們幾個重要原則,提醒我們在社交生活中所要注意的事,以致我們能在敗壞的社會中,在受污染的社交生活中,成為一枝獨秀,為主作美好的見證.
\newline
\newline
第一個原則:基督徒的社交生活,是受耶穌基督真理的限制.保羅說:「論到祭偶像之物,我們曉得我們都有知識.」又說:「論到吃祭偶像之物,我們知道偶像在世上算不得甚麼,也知道神只有一位,再沒有別的神.雖有稱為神的,或在天,或在地,就如那許多的神,許多的主.然而我們只有一位神,就是父,萬物都本於他,我們也歸於他,並有一位主,就是耶穌基督,萬物都是藉著祂有的,我們也是藉著祂有的.」我們的社交生活,應該受聖經真理的限制,絕不能超越真理的範圍.聖經是我們生活行事為人最高的權威和標準,當我們面對聖經的真理時,這真理應該從三方面應用在我們的生命中.
\newline
\newline
第一方面,聖經給我們一些誡命,這些誡命使我們知道一些是非黑白的真理.聖經誡命有消極性的,也有積極性的,我們必須全部接受,並且遵行,不能加己意.譬如說,耶穌在主禱文中教導我們求父叫我們不遇見試探,但如果我在生活上卻刻意把自己放進試探中,這就是違背神的旨意.又譬如說,聖經告訴我們,基督徒的身子是聖靈的殿,這身子就是指我們的身體和健康,是應該愛惜保養,用來榮耀神的.但如果有人染上某些習慣,是足以損害他的精神和健康的,他就是違背了神的旨意.這是除了誡命之外,聖靈引導我們應用的活聖經.
\newline
\newline
第二方面,神的話語也會針對我們個人的需要.譬如說,我們讀經知道神要我們往普天下去傳福音,但如果化身為學生,他暫時不能這樣做,於是神可能會針對他的情形,特別感動他節省金錢,把金錢奉獻在傳福音的事工上,藉此參與這個使命.
\newline
\newline
第二個原則:如果我們認為凡與聖經沒有抵觸的都可以行,那麼我們便會陷入一個危險中,保羅說:「知識叫人自高自大,惟有愛心能造就人.」那就是說,我們在社交生活中,除了有知識的界限之外,還有愛心的界限.我們行事為人,不能只憑知識,更要憑愛心,就是愛神和愛人的心.基督徒要有愛心的責任,如果我們只憑知識去判斷甚麼事可行不可行；而忽略當我們這樣行,還要對弟兄姊妹和對神負責任時；保羅說這樣的知識,只是帶我們進入自高自大的地步.因此我們還要留意這愛心的原則.
\newline
\newline
保羅說如果我們不是憑著愛心的原則去行事為人時,我們基督徒的自由便受到很大的污染,因為我們要考慮幾個嚴重的結果.第一個結果,是污穢了弟兄.八章七節:「但人不都有這等知識,有人到如今因拜慣了偶像,就以為所吃的是祭偶像之物,他們的良心既然軟弱,也污穢了.」留意「污穢」這詞.當有人看見你吃拜偶像之物時,保羅說這個人的良心便受到污穢.意思就是他們原來與神有密切暢通的關係；但因為你的緣故,他們與神之間的關係便受到攔阻.如果你的動作能引致這樣的結果,那是多可怕的.你攔阻和污染了一個軟弱的弟兄和神的關係.
\newline
\newline
第二個結果,是絆倒了弟兄「只是你們要謹慎,恐怕你們這自由,竟成了那軟弱人的絆腳石.」「絆腳石」原來的意思,就是當有敵人前來時,你拿石頭擲他,阻止他的前進.很多基督徒在社交生活中,如果不憑愛心的原則去作,他的一舉一動,便會成為別個弟兄姊妹靈性生活向前行的攔阻.
\newline
\newline
第三個結果,是傷了弟兄.(10-13)節:「若有人見你這有知識的,在偶像的廟裡坐席；這人的良心,若是軟弱,豈不放膽去吃那祭偶像之物麼?因此,基督為他死的那軟弱弟兄,也就因你的知識沉淪了,你們這得罪弟兄們,傷了他們軟弱的良心,就是得罪基督.」「傷了」這一個詞,全本新約聖經中只在這裡用過,是保羅唯一在這裡用過的.這詞原來的意思,是指一個孔武有力的勇士,欺負一個完全無抵抗能力的小孩.在屬靈的意義來說,如果你認為自己在聖經上有充分的知識,在神裡面是有操練的人,你所負的責任便愈大.因為神不願意看見這等在靈裡自認是剛強的人,因為行為上的不慎,而傷害了一個軟弱弟兄的心.如果這樣行的話,就是得罪了基督.因此我們每一個動作,都必須出自愛心.
\newline
\newline
第三個原則,十章十四到二十二節,我們在社交生活中,要常記念自己與主的關係.在任何的場合,任何生活圈子和動作中,我們要切記我們是和主基督有關係的人.當我們想到主基督時,我們想到與主有生命的關係.(16-17)節:「我們所祝福的杯,豈不是同領基督的血麼?我們所擘開的餅,豈不是同領基督的身體麼!我們雖多,仍是一個餅,一個身體,因為我們都是分受這一個餅.」保羅為甚麼在這裡提到一個類似聖餐的教訓呢?保羅所強調的,就是我們與主的關係.我們是分享主生命的人；我們必須堅守這生命的關係,到底我所行的,耶穌會行嗎?
\newline
\newline
另外我們與主有愛情的關係.二十二節:「我們可惹主的憤恨麼?」保羅要我們反省,到底我所行所作的會否傷主的心?會否引起主的嫉妒?會否令他不喜悅?會否損害我與主之間的愛情?還有我們與主的關係,是在光裡的關係.二十一節:「你們不能喝主的杯,又喝鬼旳杯,不能吃主的筵席,又吃鬼的筵席.」我們不能和神相交,又和鬼相交.保羅顯然覺得吃祭偶像之物有問題,雖然神只有一位,偶像算不得甚麼,可是魔鬼的國度卻是我們不能忽視的.我們與主是在光明中生活,因此凡屬於撒旦的東西,我們都要遠離.
\newline
\newline
最後一個很重要的原則,就是我們要對主負責任.三十二節是一個結論:「所以你們或吃或喝,無論作甚麼,都要為榮耀神而行.」在任何的事情上,我們都榮耀神.(腓四20)「願榮耀歸給我們的父神,直到永永遠遠,阿們.」這節經文告訴我們,基督徒在榮耀神的時候,要注意三方面.第一,在原文聖經中,在榮耀之前有一個很重要的冠詞；中文的意思就是願唯一的榮耀,這榮耀是只可以給神而不可以給別的.這是我們對神的衡量和體會所產生的一個結果,只將榮耀歸給神.第二,這是個不能停止,要一直持續到永永遠遠的榮耀.任何停止神榮耀的事,都是不對的.第三,當我將永不止息的榮耀歸給神時,這是一件實在的事.阿們就是實在,是誠心所願,願事情就這樣行的意思.這是我們對神的責任.
\newline
\newline


\newpage

\section{第55屆~港九培靈會~張子華~第九講　敬虔的崇拜生活}
\label{sec:517}
\textbf{第九講　敬虔的崇拜生活}
\newline
\newline
連結: \href{https://www.hkbibleconference.org/session-message/view/517}{\texttt{https://www.hkbibleconference.org/session-message/view/517}}
\newline
\newline
\hyperref[sec:516]{< < < PREV SERMON < < <}
~
\hyperref[sec:toc]{[返目錄]}
~
\hyperref[sec:368]{> > > NEXT SERMON > > >}
\newline
\newline
哥林多前書 11:1-11:34
\newline
\begin{longtable}{cl}
\hline
\hline
章節 & 經文 (和合本修訂版)\\
\hline
11:1 & \begin{tabularx}{0.7\textwidth}{X} 你們該效法我，像我效法基督一樣。 \end{tabularx} \\ \\ \relax
11:2 & \begin{tabularx}{0.7\textwidth}{X} 我稱讚你們，因為你們凡事記得我，又堅守我所傳授給你們的。 \end{tabularx} \\ \\ \relax
11:3 & \begin{tabularx}{0.7\textwidth}{X} 但是我要你們知道：基督是男人的頭；男人是女人的頭；神是基督的頭。 \end{tabularx} \\ \\ \relax
11:4 & \begin{tabularx}{0.7\textwidth}{X} 凡男人禱告或講道，若蒙著頭，就是羞辱自己的頭。 \end{tabularx} \\ \\ \relax
11:5 & \begin{tabularx}{0.7\textwidth}{X} 凡女人禱告或講道，若不蒙著頭，就是羞辱自己的頭，因為這就如同剃了頭髮一樣。 \end{tabularx} \\ \\ \relax
11:6 & \begin{tabularx}{0.7\textwidth}{X} 女人若不蒙著頭，就該剪了頭髮；女人若以剪髮剃髮為羞愧，就該蒙著頭。 \end{tabularx} \\ \\ \relax
11:7 & \begin{tabularx}{0.7\textwidth}{X} 男人本不該蒙著頭，因為他是神的形像和榮耀；但女人是男人的榮耀。 \end{tabularx} \\ \\ \relax
11:8 & \begin{tabularx}{0.7\textwidth}{X} 起初，男人不是由女人而出，女人卻是由男人而出。 \end{tabularx} \\ \\ \relax
11:9 & \begin{tabularx}{0.7\textwidth}{X} 而且男人不是為女人造的，女人卻是為男人造的。 \end{tabularx} \\ \\ \relax
11:10 & \begin{tabularx}{0.7\textwidth}{X} 因此，女人為天使的緣故應當在頭上有服權柄的記號。 \end{tabularx} \\ \\ \relax
11:11 & \begin{tabularx}{0.7\textwidth}{X} 然而，照主的安排，女人不可沒有男人，男人也不可沒有女人。 \end{tabularx} \\ \\ \relax
11:12 & \begin{tabularx}{0.7\textwidth}{X} 因為女人原是由男人而出，男人是藉著女人而生；但萬有都是出於神。 \end{tabularx} \\ \\ \relax
11:13 & \begin{tabularx}{0.7\textwidth}{X} 你們自己要判斷，女人禱告神，不蒙著頭合宜嗎？ \end{tabularx} \\ \\ \relax
11:14 & \begin{tabularx}{0.7\textwidth}{X} 你們的本性不也教導你們，男人若留長頭髮是他的羞辱嗎？ \end{tabularx} \\ \\ \relax
11:15 & \begin{tabularx}{0.7\textwidth}{X} 但女人留長頭髮是她的榮耀，因為這頭髮是給她蓋頭的。 \end{tabularx} \\ \\ \relax
11:16 & \begin{tabularx}{0.7\textwidth}{X} 若有人想要辯駁，我們卻沒有這樣的規矩，神的眾教會也沒有。 \end{tabularx} \\ \\ \relax
11:17 & \begin{tabularx}{0.7\textwidth}{X} 我現在吩咐你們這話不是在稱讚你們，因為你們聚會是有損無益的。 \end{tabularx} \\ \\ \relax
11:18 & \begin{tabularx}{0.7\textwidth}{X} 首先，我聽說你們教會聚會的時候有分裂的事，我也有些相信這話。 \end{tabularx} \\ \\ \relax
11:19 & \begin{tabularx}{0.7\textwidth}{X} 在你們中間必然有分門結黨的事，好使那些經得起考驗的人顯明出來。 \end{tabularx} \\ \\ \relax
11:20 & \begin{tabularx}{0.7\textwidth}{X} 你們聚會的時候，不是在吃主的晚餐， \end{tabularx} \\ \\ \relax
11:21 & \begin{tabularx}{0.7\textwidth}{X} 因為吃的時候，各人先吃自己的飯，甚至有人飢餓，有人酒醉。 \end{tabularx} \\ \\ \relax
11:22 & \begin{tabularx}{0.7\textwidth}{X} 難道你們沒有家可以吃喝嗎？還是你們藐視神的教會，使那沒有的羞愧呢？我該對你們說甚麼呢？我要稱讚你們嗎？在這事上我絕不稱讚你們！ \end{tabularx} \\ \\ \relax
11:23 & \begin{tabularx}{0.7\textwidth}{X} 我當日傳給你們的是從主所領受的。主耶穌被出賣的那一夜，拿起餅來， \end{tabularx} \\ \\ \relax
11:24 & \begin{tabularx}{0.7\textwidth}{X} 祝謝了，就擘開，說：「這是我的身體，為你們捨的；你們要如此行，為的是記念我。」 \end{tabularx} \\ \\ \relax
11:25 & \begin{tabularx}{0.7\textwidth}{X} 飯後，他也照樣拿起杯來，說：「這杯是用我的血所立的新約；你們每逢喝的時候，要如此行，來記念我。」 \end{tabularx} \\ \\ \relax
11:26 & \begin{tabularx}{0.7\textwidth}{X} 你們每逢吃這餅，喝這杯，是宣告主的死，直到他來。 \end{tabularx} \\ \\ \relax
11:27 & \begin{tabularx}{0.7\textwidth}{X} 所以，任何不按規矩吃了主的餅，喝了主的杯，就是干犯主的身體和主的血了。 \end{tabularx} \\ \\ \relax
11:28 & \begin{tabularx}{0.7\textwidth}{X} 人應該省察自己，然後吃這餅，喝這杯。 \end{tabularx} \\ \\ \relax
11:29 & \begin{tabularx}{0.7\textwidth}{X} 因為人吃喝，若不分辨是主的身體，他的吃喝就是定自己的罪了。 \end{tabularx} \\ \\ \relax
11:30 & \begin{tabularx}{0.7\textwidth}{X} 因此，在你們中間有好些軟弱的與患病的，長眠了的也不少。 \end{tabularx} \\ \\ \relax
11:31 & \begin{tabularx}{0.7\textwidth}{X} 我們若是先省察自己，就不至於受審判。 \end{tabularx} \\ \\ \relax
11:32 & \begin{tabularx}{0.7\textwidth}{X} 我們受審判的時候，就是被主管教，這樣就免得和世人一同被定罪。 \end{tabularx} \\ \\ \relax
11:33 & \begin{tabularx}{0.7\textwidth}{X} 所以，我的弟兄們，你們聚會吃晚餐的時候，要彼此等待。 \end{tabularx} \\ \\ \relax
11:34 & \begin{tabularx}{0.7\textwidth}{X} 若有人餓了，要在家裡先吃，免得你們聚會，反被定罪。其餘的事等我來的時候再安排。 \end{tabularx} \\ \\
[1ex]
\hline
\hline
\end{longtable}
崇拜是基督徒在世上一個很重要的活動,其重要的原因有幾方面.首先,它是個唯一的活動.耶穌基督對門徒有一個直接的要求,祂說:「父正要這樣的人拜他.」耶穌在約翰福音向我們表達了一個心意,就是父要我們在世上生活時,要過一個敬虔崇拜他的生活.這唯一的活動,直到將來天上,也要繼續.為此,我們在世上的時候,要不斷學習過崇拜的生活,以致將來我們在天上的時候不會陌生.
\newline
\newline
哥林多前書十一章,是我們今天要思想的經文.這章聖經可以分成兩大部份,是我們在敬虔生活中所要注意的.頭一部份是二到十六節,講到在崇拜生活中,對神必須有尊敬的態度.十七到三十四節,講到在崇拜生活中,要以主為中心,有思念和敬拜主的成份.不過,在這兩段經文中,保羅指出當我們在進入敬虔崇拜的生活以先,必須打穩一個很重要的人生根基,然後才談得上尊敬的崇拜和記念性的崇拜.
\newline
\newline
這重要的根基,綜合來說,就是信徒每天和神的關係.個人必須先有敬虔的生活,才能在整體性的崇拜中有敬虔的表現.這敬虔的態度有兩方面的表現；第一是對權柄的順服.(十一3):「我願意你們知道,基督是各人的頭,男人是女人的頭,神是基督的頭.」保羅為甚麼在論到崇拜生活時要提到這些呢?這裡提到有三個頭,基督,男人和神.神是基督的頭,在救贖的歷史來說,是差遣基督的；基督要順服神,去完成這救贖的工作.耶穌曾經說:「我的食物,就是遵行那差我來者的旨意,作成他的工.」基督的一生,對神的權柄有完全的順服.而基督是男人的頭,第七節說:「男人本不該蒙著頭,因為他是神的形像和榮耀.」男人若要在家庭和教會中作帶領的工作,就必須自己先對基督有絕對的順服.
\newline
\newline
還有一個頭,就是男人是女人的頭.其實男人和女人,在神的創造中是平等的；但這裡所指的,是神所設立的兩個機構:教會與家庭.在帶領教會和帶領家庭中,神的旨意是要男人負更大的責任,付更多的代價；走一條更艱難的路.保羅講這三個頭的時候,他的意思是要指出；無論男女,如果要過敬虔的崇拜生活,就必須有一個順服權柄的生活.因為順服神的權柄,我們也順服神在教會和家庭中所命定的權柄.我們的不滿,反叛,其就是對神反叛的一個反照.如果我們順服神,也必順服神所命定的權柄.一個不順服權柄的人,是不能進到敬虔的崇拜生活中,因為我們所敬拜的對像,是一位萬有的主宰.
\newline
\newline
除了順服權柄之外,還要追求過聖潔的生活,這也是一個很重要的根基.保羅特別用聖餐去警告哥林多的教會,餐是一個崇拜的生活,是充滿記念性的.但保羅卻警告他們,因為他們在生活上的不潔,所以招致神的審判,三十節說:「因此,在你們中間有好些軟弱的,與患病的,死的也不少.」保羅說如果他們在生活上不潔,又帶著這些不潔守主餐時,這便會帶來主的審判,懲治,甚至有軟弱,患病和死亡.
\newline
\newline
哥林多教會怎樣在守聖餐時,表現出他們平時的生活狀況呢?當保講到他們不尊重聖餐時,他講得很嚴重,(17-19)說:「我現今吩咐你們的話,不是稱讚你們,因為你們聚會,不是受益,乃是招損.第一,我聽說你們聚會的時候,彼此分門別類,我也稍微的信這話.在你們中間不免有分門結黨的事,好叫那些有經驗的人,顯明出來.」
\newline
\newline
原來哥林多教會在聖潔上,出現了兩個問題,頭一個是思想上有錯誤.這並不是說他們對真理有異端的法,只是說由於他們的思想有錯誤,因而引致分門別類和結黨的事.這分門別類的問題,正好違背了聖餐一個重要的意義,就是我們同領一個餅和一個杯的.另一個問題是有關私生活方面的.起先他們依照主設立聖餐之前所做的；他們在守聖餐之前有愛筵,大家一同先吃飯,然後才領受聖餐.這來是好的,只可惜有人慢慢染了社會的風氣,利用這段愛筵的時候,大吃大喝,失卻了記念的意義.(20-22)節:「你們聚會的時候,算不得吃主的晚餐.因為吃的時候,各人先吃自己的飯,甚至這個飢餓,那個酒醉.你們要吃喝,難道沒有家麼?還是藐視神的教會,叫那沒有的羞愧呢?我向你們可怎麼說呢?可因此稱讚你們麼?我不稱讚.」信徒把這種生活態度帶進崇拜的生活中,是多麼可恥的一件事.雖然今天的教會鮮有聖餐之前的聚餐,但保羅所指出的原則就是:我們在守聖餐之前,必須在身心靈方面有聖潔的準備,先省察自己,免得吃喝自己的罪.這決定性的因素,成為教會和個人敬拜神不可代替的基礎.因此,個人與主的關係是極重要的,能直接幫助我們過一個敬虔的崇拜生活.
\newline
\newline
當基礎打好以後,我們來到教會便沒有問題了.我們的崇拜,第一有很大的尊敬性.既然個人的生活已經有尊敬神的表現,那麼來到教會便有合宜的表現.保羅主張女人在聚會中蒙頭的用意也是在此；是表示她從心裡順服權柄.因為哥林多教會,遇見這樣的問題；當時的社會一個道德淪落的社會,有千多間妓院,因此教會面臨一個問題.其中有些女人,不願順服權柄,他們穿大紅大綠的衣服進教會,非常驕傲.保羅於是便勸勉教會的姊妹要和世俗的女子有別,便建議他們蒙頭,以外表的記號來表示內心的順服.今天,姊妹在外表上蒙頭不蒙頭,並不緊要,要緊的是有合宜的態度,有一顆敬虔的心來參加崇拜.
\newline
\newline
有幾樣事我們在崇拜中應該注意的.保羅在這裡建議用外面的表現來表達內心的敬虔,這是沒有錯的.我們要有合宜的服飾去參加合宜的活動.沒有人會穿西裝去打網球,或是穿汗衣去參加婚禮.同樣,我們在敬拜這位萬王之王,萬主之主時,我們不應有合宜的服飾去表現我們內心對主的敬拜嗎?舊約神在西乃山向以色列人顯現時,神要他們先洗衣服,預備一件乾淨的衣服,才去朝見永生神.
\newline
\newline
另外一個合宜的表現,是我們在崇拜之前心靈上的準備.我們在主日崇拜之前,先要準備好一個敬拜的心靈,才去參加崇拜.其中守時是一個很緊要的因素,先到會場,安靜自己,好好默禱,才參加崇拜.內心的敬虔,加上外表合宜的態度,才是神所喜悅的.
\newline
\newline
除了合宜的表現之外,個人還需要有適當的參與.保羅不願參加崇拜的人,成為一個旁觀者,當一個人講道時,大家都講道；當一個人禱告時,全體都一同禱告；這表示每個人的參與和投入.當牧師講道時,要留心靜聽,用心靈和精神去吸收；有人領禱時,要和他同心,一致說阿們.我們全心全意的集中精神,投入崇拜中；不東張西望,不魂遊象外,乃是全人的參與,而非旁觀者.倘若我們能全心全意的參與崇拜的每一項,就不會在崇拜中感到陌生和枯燥.
\newline
\newline
還有,一個敬虔的崇拜,應該有很深的記念性的成份.保羅特別以聖餐作為例子,說明在崇拜的生活中,我們要記念主；包括主的受死,復活；祂的一生,祂的為人,以及祂的將來.我們要以主為思念的中心,記念他.因此,崇拜和別的聚會,如祈禱會,查經班,主日學,團契等是不同的；別的聚會固然也應以耶穌為中心,但是有部份也是生命中情感的輸出.然而在崇拜中,卻只有思念主,不是只為了得著造就,更重要的將自己所有的帶給主.神當然也會藉著我們的唱詩,讚美,對祂的敬重賜福我們.因此,敬虔的崇拜是應該以主為中心的.
\newline
\newline
但願我們今天的教會,能注重敬虔的崇拜生活.不單注重教會整體的崇拜,更重要的,是個人敬虔生活的操練；要建立一個好的基礎,有順服神的心,更要學習在外表上有合宜的態度,來配合內心的敬虔.
\newline
\newline


\newpage

\section{第56屆~港九培靈會~林克己~第一講　耶和華的呼召(一節)}
\label{sec:368}
\textbf{第一講　耶和華的呼召(一節)}
\newline
\newline
連結: \href{https://www.hkbibleconference.org/session-message/view/368}{\texttt{https://www.hkbibleconference.org/session-message/view/368}}
\newline
\newline
\hyperref[sec:517]{< < < PREV SERMON < < <}
~
\hyperref[sec:toc]{[返目錄]}
~
\hyperref[sec:370]{> > > NEXT SERMON > > >}
\newline
\newline
俄巴底亞書 1:1-1:1
\newline
\begin{longtable}{cl}
\hline
\hline
章節 & 經文 (和合本修訂版)\\
\hline
1:1 & \begin{tabularx}{0.7\textwidth}{X} 俄巴底亞所見的異象。論以東說：我從耶和華那裡聽見信息，並有使者被差往列國去，說：起來吧，一同起來與以東爭戰！ \end{tabularx} \\ \\
[1ex]
\hline
\hline
\end{longtable}
各位同工同道,八十年代是人類很有挑戰的時代,科學發達,知識爆炸,我們面對的問題跟前一代的截然不同.求主用祂話語引導我們,我經過禱告,默想後,要講從俄巴底亞書看八十年代基督徒的挑戰.這位古代先知所面對的挑戰,困難以及所碰到的機會,所傳講的信息,都能帶給我們很多的幫助.可惜讀俄巴底亞書的基督徒卻不多.兄弟曾在一間教會講俄巴底亞書,但散會後,有一位姊妹說:「我在這間教會聚會了廿五年,這是頭一次,聽見講員用俄巴底亞書來教導勉勵我們.」願主解開祂的話語,使我們的靈性生活得著滋潤和亮光.
\newline
\newline
今天我們集中思想耶和華的呼召.正當以色列人在靈性退步,政治腐敗的時候；主就呼召,興起祂的僕人俄巴底亞先知,預言要發生的事.在先知所處的時代中,講出當講的話.首先我們看看這書的歷史背景,尤其是俄巴底亞傳信息的資格,環境和對象.一節:「俄巴底亞得了耶和華的默示.」本節介紹先知時,非常簡單.俄巴底亞是耶和華的僕人和器皿,但他具有甚麼資格和受過甚麼訓練呢?在舊約中,最少有十二位名叫俄巴底亞的人,他們的工作經驗,社會身份,有些聖經講得很詳細清楚,例如王上十八章說及的俄巴底亞.所發現的事蹟及對他的認識很有限,不能肯定他的家庭和教育背景.除了一節這句話承認他被耶和華所揀選作先知外；我們對他的背景差不多一無所知.揀選先知的耶和華,祂有絕對的智慧和主權；祂揀選這個人,使他終生作祂忠心的僕人.
\newline
\newline
(一)先知的資格
\newline
\newline
俄巴底亞有甚麼資格作先知,聖經中沒有提及他受過任何訓練.雖然在舊約(王下第六章)提及以利亞及以利沙,設立過先知門徒學校,門徒在一起追求認識主；但俄巴底亞先知卻沒有這種訓練.這給我們八十年代的基督徒很大的鼓勵,雖然只有教會的訓練班及聖經夜校,此外沒有其它的訓練；但只要我們用心事奉主,像俄巴底亞一樣,也可以被神大大使用.兄弟在一九六Ｏ年第一次來遠東事奉,同行的還有六位弟兄姊妹.我在七個人中是最高的,信心卻是最小的；我們當中有一位姊妹,身材最矮,但信心最大,故此我們就稱她做「摩西」.她沒有大學畢業及其他學位,但她的信心常作我們的榜樣.這位姊妹每逢作見證時,常引用(林後八12):「因為人若有願作的心,必蒙悅納,乃是照他所有的,並不是照他所無的.」她常作見證自己蒙召事奉主的時候,感覺不配及不適合,教育,恩賜,經驗方面都不能作全時間事奉,但主用這節聖經提醒我,你只要預備一顆願作的心,就可以了,我向你要求的,永不會超過你所能擺上的,我吩附你去作工,必定給你有足夠的恩典和力量.我不知聽過這位女同工這樣作見證多少次,在我作宣教事奉廿多年中,一想起這見證,便使我大得安慰和鼓勵.
\newline
\newline
基督徒常常有掙扎,因為在我們裡面有感動去做一件事,教會工作需要我們奉獻,否則別人便沒有機會信主.但另一方面,我們也感到自己的軟弱和卑微,但願從這位女宣教士「摩西」的見證中,看見神呼召我們；必定會給我們智慧和力量去做,只要我們存著一顆謙卑的心,承認自己的軟弱和限制；用信心抬起頭來仰望祂,倚靠祂,祂便能使用我們.
\newline
\newline
現在我們看看一件有趣的事,俄巴底亞的名字很有意思.聖經上人的名字與他們的工作及品格很有關係.例如保羅是微小的意思,他雖然沒有體面和健康,甚至被人藐視,在人面前微小.但因為他大有信心,在神面前便偉大.約書亞名字的意思是救主,很適合他的工作和呼召；他率領以色列人到戰場攻擊仇敵,使他們脫離敵人的手.司提反是冠冕的意思,這很適合這位初期教會第一位殉道者；他雖然為自己的見證而死,但得著生命的冠冕.俄巴底亞的名字,很有挑戰性的意義,乃是「敬拜,事奉耶和華」.有敬拜和事奉兩方面,今天基督徒,注重敬拜,但忽略事奉的人很多.雖然每主日出席崇拜,但在教會工作和事奉上,少有貢獻和參與!他們以為作基督徒是坐在一輛往天堂的巴士,牧師是司機,所有工作推到牧師身上,自己不用幫忙和參與；只享受作基督徒的福氣,不盡基督徒的責任.
\newline
\newline
相反的,有一些基督徒注重事奉,而忽略敬拜.至日常工作和教會事奉都很忙,沒有時間也沒有心,在主面前讀經等候,不自覺地生活.也不傳揚福音,一切都是靠自己,失去事奉主的力量.求主保守我們脫離這網羅,願主使我們認識到敬拜和事奉都一樣的重要,好像俄巴底亞先知一樣.他的品格和屬靈生活,都與他的名字是相稱的,難怪耶和華呼召和使用他.
\newline
\newline
(二)他的環境
\newline
\newline
俄巴底亞書寫作的日期很難肯定,大多數的解經家認為寫作日期,大概在主前第六世紀.因為猶太國的京城耶路撒冷,是在主前五八六年,被巴比倫王尼布革尼撒攻陷的.看這書的內容,提及主前五八六年的事.(一11)他責備以東人不來幫助自己的同胞,因為以東人和以色列人都是以撒的後裔.(一11)「當外人擄掠雅各的財物,外人進入他的城門,為耶路撒冷拈鬮的日子,你竟站在一旁,像與他們同夥.」,在(12-14),好像是描述攻破耶路撒冷的事,這是主的選民憂傷絕望的時代；但先知俄巴底亞仍然非常樂觀,認為主的選民經過苦難和痛苦之後；必定從苦難中平安出來,又在耶和華賜福之下蒙恩.國家政治和屬靈方面,都有一個光明的前途,在17節先知用信心說:「雅各家必得原有的產業.」21節:「國度就歸耶和華了.」這位先知,在患難和痛苦中所保持的,是樂觀和得勝的生活.我們八十年代的基督徒要養成這個習慣,無論如何受苦,所需要的,就是先知的信心和得勝的態度.保羅對這件事的回應說:「所以我們不喪膽,外體雖然毀壞,內心卻一天新似一天；我們這至暫至輕的苦楚,要為我們成就極重無比永遠的榮耀.」(林後四16-17)有一天內子在他的宣教中,遇見一個信心得勝的家庭；他們為福音緣故,被人逼迫,住在公共廁所內幾個月.他們的兒子也被強逼跪在碎玻璃上幾小時,一直到他的膝蓋流血痛苦；但是他們不畏環境的壓力和威脅,不埋怨主,也不否認主.願主幫助我們,在試探,引誘中,在前途不明之中,效法他們的榜樣；在信心和樂觀中,作俄巴底亞的屬靈後裔.
\newline
\newline
(三)他傳信息的對象
\newline
\newline
1節:「俄巴底亞得了耶和華的默示,論以東說,我從耶和華那裡得到信息.」對象是以東人,以色列人的祖宗以撒有雙生子,大的以掃,小的雅各.以色列人是雅各的後裔,以東人也是雅各的子孫；這兩種人雖然來自同一祖宗,但後代之間往往彼此為仇.這有幾個原因.但今天我們集中思想一個,其他留待以後再談.從查考創世記時,一個很有趣,意義和挑戰性的原因.當我們查考以掃和雅各的生平時,比教會遇見一個事實和統計；看見雅各有廿二次和耶和華交通,但從沒有一次記載以掃與耶和華的來往.他不敬拜,禱告神,好像一點也不承認神的存在和價值.可說以掃的生活是完全屬世界和屬血氣的.這不是說雅各沒有缺點,他在神改變他之前；兩次欺騙他的哥哥,但是問題卻發生在宗教生活的方式上,以掃因為不以「神」為重要,因此給後裔的後遺症是一個極大的影響.有句話說:「父親怎樣,兒子也怎樣.」這件事就成為以東人的絆腳石,按照一般聖經學者的研究,發現在主前六世紀,當耶和華興起先知俄巴底亞講預言的時候,以東人已成為無宗教信仰者.他們沒有神,沒有偶像,也沒有明顯的敬拜偶像.這很像我們遇見的一些人,因為沒有宗教信仰,自稱為自由思想者.兄弟未到香港以前,在馬來西亞事奉,有一位中國朋友說,在馬來西亞有兩個政黨,一個叫勞工黨,一個叫聯盟黨.他說:「我沒有宗教思想,也不支持甚麼黨,我是屬於「發財黨」；專心以發財為人生最大目標.」今日在我們周圍的人,自由思想者很多,他們不相信有神的存在.求主讓我們好像先知俄巴底亞,在自己的工作崗位上謹守；相信俄巴底亞書,會與我們同感.他這個事奉敬拜神的人,蒙召向以東人作見證,他們領受或不領受,仍要工作下去.(賽四二4):「他不灰心,也不喪膽,直到他在地上設立公理,海島都等候他的訓誨.」雖然俄巴底亞傳信息的對象,對宗教信仰沒有興趣,也毫無反應；但他忠心在工作及見證上,真是我們永垂不朽的榜樣.
\newline
\newline
最後,讓這個「父親怎樣,兒子也怎樣.」的信息中,教訓我們!我們的行為或舉止,不是別人的好榜樣,便是絆腳石.若不向人傳揚福音,吸引他們對福音有興趣；就是得罪他們,使我們輕視我們寶貴的福音.以掃屬肉體的生活,不重視神,在後裔中留下壞的影響!如果我們的靈性生活不冷不熱,便會攔阻別人來相信耶穌.有一個故事說,一個住在山上的家庭,當冬天嚴寒下雪的時候,做父親的要越過被雪遮蓋的高山去探候親友.父親因為兒子年紀小,不帶兒子同行,就獨自動身出去.但走不遠的時候,竟然聽見後面有聲音走近,原來他的兒子也跟著來.兒子說:「爸爸請你走慢一點,我會盡量跟隨你的腳蹤走,無論你往哪裡去,我也可以放膽去.」各位弟兄姊妹,我們的行為聖潔光明嗎?如果在我們走的屬靈路上,有剛信主的基督徒,看我們的生活效法我們的榜樣,他們能否這樣對我們說?「看你們的生活,你們無論往哪裡去,我知道可以放膽跟著你們去.」願主幫助我們應該以以掃為鑑戒,要效法忠心僕人俄巴底亞.在我們的時代環境中,作敬拜,事奉耶和華的僕人.
\newline
\newline


\newpage

\section{第56屆~港九培靈會~林克己~第二講　耶和華的主權}
\label{sec:370}
\textbf{第二講　耶和華的主權}
\newline
\newline
連結: \href{https://www.hkbibleconference.org/session-message/view/370}{\texttt{https://www.hkbibleconference.org/session-message/view/370}}
\newline
\newline
\hyperref[sec:368]{< < < PREV SERMON < < <}
~
\hyperref[sec:toc]{[返目錄]}
~
\hyperref[sec:372]{> > > NEXT SERMON > > >}
\newline
\newline
俄巴底亞書 1:1-1:2
\newline
\begin{longtable}{cl}
\hline
\hline
章節 & 經文 (和合本修訂版)\\
\hline
1:1 & \begin{tabularx}{0.7\textwidth}{X} 俄巴底亞所見的異象。論以東說：我從耶和華那裡聽見信息，並有使者被差往列國去，說：起來吧，一同起來與以東爭戰！ \end{tabularx} \\ \\ \relax
1:2 & \begin{tabularx}{0.7\textwidth}{X} 看哪，我要使你在列國中為最小，被人大大藐視。 \end{tabularx} \\ \\
[1ex]
\hline
\hline
\end{longtable}
「俄巴底亞得了耶和華的默示,論以東說,我從耶和華那裡聽見信息,並有使者被差往列國說:「起來罷!一同來與以東爭戰,我使你以東在列國中為最小的,被人大大藐視.」
\newline
\newline
我們從這段經文來看耶和華的主權,這是我們的時代所需要的真理,可從五方面來看:恩典,揀選,差派,歷史,報應.若將這五件放在一起,我們便發現我們的神偉大無比,超過萬有.祂的慈愛和權柄永不改變,祂能為我們各方面的前途負責,使我們坦然無懼的跟隨祂.
\newline
\newline
(一)恩典方面
\newline
\newline
我們基督徒常常提及主的恩典,唱讚美詩的時候,歌頌主恩典.聖經中也多次提及主的恩典,但究竟「恩典」兩個字是甚麼意思?與很多信徒談起,才發現他們對恩典的定義不大清楚.簡單來說,恩典是你我皆不配享受的,是耶和華的大愛.因我們是罪人,在言語,思想,行為上得罪聖潔的神；如果這位神按著公義來定我們的罪,而沒有為我們預備救法,我們誰也不能上天堂.所以聖經上提及所有神的大愛,都是我們不配享受的,這就是恩典.俄巴底亞書第一,二節所介紹的神,是在人類歷史運行的.祂想辦法把自己和自己的事,都擺在我們眼前,使我們有機會認識祂,這也是恩典.我們的神不將自己隱藏起來,祂也不默然不語,祂願意向自己創造的人類顯現,祂這樣賜恩給我們.
\newline
\newline
「俄巴底亞得了耶和華的默示」,讓我們注意「得了」這兩個字.若耶和華不給我們,我們就不得到,但在聖經裡我們可以得到一切.這裡提及的默示,是祂為自己兒女所定的計劃,安慰和引導.默示使我們發現祂為我們生命所預備的,那個完全而不能落空的計劃在這裡默示很重要,它帶著耶和華的權柄和生命,是很有能力的兩個字,表明神的話語句句都可靠.祂的先知俄巴底亞,講及猶大及以東的歷史,所看見的異象,所得到的啟示,不能不發生,這是他絕對不能抵抗的主權.對祂的兒女來說,這也是恩典.祂怎樣向他們說話,為他們各方面安排前途,在試煉和仇敵前保守他們,為他們的生命預備一個完全的計劃,這是俄巴底亞書很實際的教訓.如果我們已重生得救,是神的兒女,我們便可以將這好信息應用在自己的身上.我有一位作牧師的好朋友,他讀大學時唸地理學；他已經是基督徒,故常為自己的前途祈禱,主為他生命定下旨意.經過禱告和考慮後,他盼望畢業後能做一個教師；因為他喜歡青年人,盼望除了教學外,可以使學生認識主耶穌.但經過考試後,他的成績不大理想,不適合作教師；但他不灰心,用禱告將自己的前途,再一次交給神,求祂帶領.結果蒙主引導,全時間奉獻,讀神學的時候,每一科目都沒問題,這不是說神學比地理學容易.而乃是說明如果你讀神學,甚麼科目和課程,都不能攔阻神的計劃來實行祂的旨意?祂的恩典每一天,每一步都夠你用.
\newline
\newline
(二)揀選方面
\newline
\newline
為何神要向俄巴底亞顯現,為何在萬人中揀選了他,作祂的先知,器皿,僕人和代表.從一方面來看,俄巴底亞有適合作先知的資格,他是一個事奉,敬拜耶和華的人；但從另一方面來看,只能說是耶和華的揀選,主怎樣揀選,這是主的計劃和智慧.我承認耶和華在揀選上有絕對的主權,但我也不十分明白.我雖然很喜歡在中國人中服侍主,我也有很多中國朋友,但是神為甚麼要揀選我這樣身高的人,來做這件工作.很多時候,與替我翻譯的同工一起上講台,會眾就笑起來,難怪使徒保羅說:「我成了一台戲,給天使和世人觀看.」有一天當我上天堂的時候,我會問主,為甚麼你這麼有幽默感,揀選像我這樣高的人,在我所愛的中國人中服侍主.兄弟還有這一個問題,廿多年前,蒙召作宣教士的時候,心裡感到自己非常的不適合,一直到現在,還有這種感覺.一點也不明白為甚麼神要揀選我,想到自己廿多年來學習的國語,還是不很像樣子和不準確,怎樣不流利.當想到這位蘇格蘭宣教士馬禮遜博士,只需十二年時間,便將全本聖經翻譯成中文,我就很灰心,只有承認他是「特殊例子」.基督徒每當向裡面看,就會看見自己的缺點這麼多；須要用信心抬起頭來,用信心接受耶和華在揀選上的主權；來安慰自己,繼續為神工作.記得我自己蒙召全時間服侍主的時候,我有很多懷疑和推辭.我對主說:「主啊!你知道我各樣的缺點,沒有口才和恩賜,沒有經驗和資格,不夠聖潔,信心也不夠大,接受主的時候,沒有像保羅看見大光.主啊!我雖然不應該這樣說,也不應該這樣想,及不該消滅你的感動,但是你實在揀選錯了,請你另外揀選別人來服侍你吧!」各位兄弟姊妹,我真相信俄巴底亞先知,蒙召的時候,也有同樣的感覺,而且一生也都會有.但教會若有工作的需要,無論本地或海外,倘若我們有這方面的感動,又確實的知道,卻不去做,又沒有人去做；而這個考慮是經過時間,思想和禱告,仍然存在的.可見這個感動是來自神的靈,我們就該將這個感動告訴教會；倘若教會亦同意和承認,我們就可以放膽去做.雖然感到自己不適合也不要緊,要緊的是耶和華在揀選上的主權.因為除了耶穌外,從古代一直到現代,在這個世界,沒有一個完全的神的僕人.神既然揀選我們,祂就會負責我們各方面的需要,無論是物質或屬靈的,祂會為我們作出適合的準備.
\newline
\newline
(三)差派方面
\newline
\newline
「並有使者被差往列國去」,這位使者是誰,聖經中沒有提及,有人說他是一個人,有人說是一位天使；也有人說這句話大概是一個比方.論到以東的鄰國,怎樣圍攻以東,好像有使者招聚他們這樣做.無論如何這使者的差派,是耶和華的安排,為要在人類歷史上實行祂的旨意.這位僕人怎樣被差派,離開他的本地和崗位,到外地去傳講耶和華的信息；同樣主也在我們的時代中派僕人出去,帶著祂的消息,到祂為他們揀選的地方去；做祂為他們預備的工作,一直到他們完成了祂的大使命.我相信俄巴底亞書第一節的使者差派,可以應用在今日教會的差傳事工上.這樣可幫助我們對本地和海外工作,都發生責任感和使命感.不單要捐錢,禱告或出去；最重要是接納耶和華在差派上的主權.祂向你顯現,對你說話,叫你到何地去,請不要與祂辯論,不要與祂講價,千萬不要推辭,作耶和華差派你的一位使者,這個身份和本份,是一件尊貴榮耀的事.祂呼召你與祂同工,這不單是祂的主權,也是你的福氣；你願意放下你的一切,隨著祂的引導和感動,向祂指示你的地方,祂就能大大的使用你.
\newline
\newline
有一次兄弟返英國述職的時候,有信徒向我提及他的牧師是十分愛人靈魂的牧者.這位牧師無論在甚麼時候,一有主的感動和引導,馬上就出去探訪傳福音.有一天晚上九點鐘,這位牧師心裡有些感動,要馬上外出去探訪一對認識的夫婦,丈夫尚未信主,太太是他教會的會友,到達他們的家庭時；那位太太便接待他,因為她的丈夫尚未回家.他是一個水手,有時候要出海,於是牧師便回家了,因為明天還要去探訪.當晚十點鐘,那一個感動又回來了,他於是第二次到那對夫婦的家,但是丈夫還未回來!牧師回家後,就跪下祈禱,免得聽錯主的感動.到十一點鐘時,第三次的感動又要他到那對夫婦的家,這時婦人的丈夫剛好回來了,牧師只講了幾句話,關於聖經和祈禱的,那位水手便回答說:「我預備好了,現在接受耶穌做我的救主.」這位牧師因為承認耶和華在差派事工的主權；主吩咐他,他就去了,所以他能被神使用.
\newline
\newline
(四)歷史方面
\newline
\newline
很明顯在第一,二節看見,耶和華是人類歷史的主.無論是猶大國,以東或列國,耶和華都有主權；從祂天上的寶座,來管理他們的興旺和衰微.祂按自己美好的旨意,為凡相信祂的人,今生或來生預備光明的前途；讓我們在這個前途不明朗的時代中,從俄巴底亞的預言中,學習功課,領受教訓.在這個世界上,最偉大,最絕對的主權,都是在耶和華的手裡.讓我們把握這件事實,在這個宇宙當中最高的,最寶貴的,最大的權柄,都是在天上.
\newline
\newline
另外按照聖經的教訓,在人類歷史最後一次戰爭中；良善要得勝邪惡,我們的救主耶穌基督,要作永遠的得勝者.我們要在俄巴底亞書的教訓中,對耶和華在歷史上的主權,應該產生堅定不移的信心.耶穌說,魔鬼和福音的仇敵,所做的一切,都不能破壞祂所建立的教會.基督徒雖然經過苦難和逼迫,受試探和試煉；但是魔鬼不能奪去我們的永生和盼望,知道這點就能站立得穩,至死忠心.有了這個信仰,耶和華在歷史上有絕對的主權；美國的林肯作總統時,當時地方上以及國際間有很多的問題,令他很驚慌罣慮；那時他說:我便默想耶和華在歷史上的主權.心中就得到無限的安慰!這就是俄巴底亞先知所高舉的,從舊約到新約,從古代到近代；祂為祂的僕人所預備的都是最好的.雖然有時候我們會受到困擾,但這位耶和華會看顧我們,使我們的基督徒品格和修養上有長進.這位林肯總統,在困難的時候,最愛讀(詩三十四4):「我曾尋求耶和華,祂就應允我,救我脫離一切的恐懼.」相信你我常為將來罣慮,但讀這詩篇可大得安慰；重新得力,將大衛的經歷應用在我們身上,將這應許藏在心裡.
\newline
\newline
(五)報應方面
\newline
\newline
中國有一句話「善惡到頭終有報」,只差早或遲,「我使你以東在列國中為最小的,被人大大藐視」(二節),雖然以東與猶大乃同一祖宗所生,都是以撒的後裔.但多年來,與猶大國的仇敵同盟,逼迫猶大人.雖然如此,猶大人是神的選民,乃耶和華所揀選和所愛的,先知俄巴底亞在這裡,注重耶和華在報應上的絕對主權.雖然主好像不動工,不施行審判邪惡,都有應得的報應.亞伯拉罕為所多瑪,蛾摩拉二城,向耶和華禱告說:「將惡人與義人一樣看待,這斷不是你所行的；審判全地的主,豈不行公義嗎?」亞伯拉罕這樣禱告實在不錯,俄巴底亞書第二節,便將耶和華的審判和應許,應用在以東人身上.他們怎樣看猶大人為小的,自己也要被看為最小.這裡有一種原則和真理,就是在各時代中神的兒女,雖然暫時被仇敵壓迫玩弄,但在末後,這位聖潔公義的耶和華,必定為他們伸冤!為他們預備最好的.聖經說:「在你面前,有滿足的喜樂,在你右手中有永遠的福樂.」(詩十六11)聖經也說:「耶和華是日頭,是盾牌,要賜下恩惠和榮耀,祂未曾留下一樣的好處,不給那行動正直的人.」(詩八四11)
\newline
\newline
有一位古代教會的主教,因為為主作見證在羅馬帝王面前受審判.皇帝對主教說:「你的耶穌不過是木匠,你的木匠,現在在哪裡?現在作些甚麼呢?」主教馬上回答說:「皇上我的木匠現在在天上,忙著作木工；所作的是一副棺木,為你和你的帝國預備的.」這位主教所說的,不是出於仇恨,乃是他十分清楚,最大的主權是屬於耶和華.無論人怎樣說,神在恩典上,揀選上,差派上,歷史上和報應上；在天上和地上,永遠與愛祂的人同在.
\newline
\newline


\newpage

\section{第56屆~港九培靈會~林克己~第三講　耶和華的偉大}
\label{sec:372}
\textbf{第三講　耶和華的偉大}
\newline
\newline
連結: \href{https://www.hkbibleconference.org/session-message/view/372}{\texttt{https://www.hkbibleconference.org/session-message/view/372}}
\newline
\newline
\hyperref[sec:370]{< < < PREV SERMON < < <}
~
\hyperref[sec:toc]{[返目錄]}
~
\hyperref[sec:373]{> > > NEXT SERMON > > >}
\newline
\newline
俄巴底亞書 1:3-1:4
\newline
\begin{longtable}{cl}
\hline
\hline
章節 & 經文 (和合本修訂版)\\
\hline
1:3 & \begin{tabularx}{0.7\textwidth}{X} 你狂傲的心欺騙了你，你住在巖穴，居所在高處，心裡說：「誰能把我拉下來到地上呢？」 \end{tabularx} \\ \\ \relax
1:4 & \begin{tabularx}{0.7\textwidth}{X} 你雖如鷹高飛，在星宿之間搭窩，我必從那裡拉你下來。這是耶和華說的。 \end{tabularx} \\ \\
[1ex]
\hline
\hline
\end{longtable}
「住在山穴中,居所在高處的阿,你因狂傲自欺,心裡說,誰能將我拉下地去呢?你雖如大鷹高飛,在星宿之間搭窩,我必從那裡拉下你來,這是耶和華說的.」
\newline
\newline
有一位英國牧師寫了一本書,名為「你的神太小.」目的是提醒基督徒,我們的神耶和華實在偉大無比.有時候我們的信心很小,忘記神是永遠偉大的；遠遠超過我們所能想像的,也遠遠超越一切的問題,祂能供應我們各方面的需要.現在我們要從三個角度來看耶和華的偉大；這些基本的真理,我們不要以為簡單不重要,其實它正是我們這時代所最需要的.
\newline
\newline
(一)祂的無所不知
\newline
\newline
在這方面,我們要注意神怎樣徹底的認識這些以東人.祂既然認識他們,也必照樣認識我們；這件事對我們既有挑戰,也有提醒,又有安慰.在第三節,看見神很清楚的認識人.
\newline
\newline
第一,祂知道我們的家在那裡:「住在山穴中,居所在高處的阿!」(三節),當我們初次認識一個人,也常會問:「你住在哪裡?」,耶和華神就與我們不同,人不告訴我們,我們就不知道他住在哪裡,但神卻是無所不知的.感謝神!耶和華的名字,能帶給我們無限的安慰,可以幫助我們為自己和別人禱告.兄弟在馬來西亞傳福音的時候,有機會學習這個功課.那裡風景和環境很美麗,在我婚前的九年,住的房間不太舒適和妥當.第一次租的房間是在鄉下,在中國人的家裡；他們不會英語,所以我有好機會學習中國語；同時對中國人的文化和風俗,有一點的研究.他們很歡迎我,我也記得,當我第一次踏進房間時,我雖然不是體重的人；但那些地板,全部都震動了.有一次一位中國朋友,怕我孤單寂寞,便來陪伴我.當他踏入房間,地板便震動,他馬上改變說:「這些地板這樣震動,我不敢陪你了.」我的中國房東,對我很好,在我未曾搬來之前,按照我喜愛的顏色,在牆壁上油漆,但牆壁上有許多一口口的鐵釘,我必須把它們拔出,才可下手去油漆.第二天睡醒後便禱告讀經,才下樓洗臉；後來又發現有豬,雞和狗在我房門口睡覺.我一不小心,便仆倒在牠們身上,這些家禽的啼叫聲把全屋的人吵醒了!這村子沒有自來水,要從井裡取水.白日又沒有電,因發電機要晚上才開動；所以中午天氣最熱的時候,沒有辦法使用電風扇和冷氣機.這村子既沒有這些電用,我委實住得不大舒服,在不大安靜的環境中,真需要主的恩典,因為我實在住不慣.但感謝主,神知道這一切.神知道以東人住在那裡,也知道我們住在那裡；祂把我們放在甚麼地方,是要祂作見證.在所居住的地方,無論是舒服,安靜或方便與否；他會使我們適應,也供給我們各方面的需要.在那個村子,主為我預備了一個好的中文老師,就是房東七歲的小女兒；她每天放學後回家,進入我的房間,教我背誦那天的功課.她聽我背誦的時候,對我很嚴格,無論是口勢,發音,音調,有甚麼錯誤,她馬上就糾正.我跟她背誦的功課很多,直至今日也不忘記.耶和華是無所不知的,基督徒的環境和需要祂都知道,祂能幫助我們去解決困難,適應現實生活.耶和華是無所不知的,這話實在甜蜜!我們無論有甚麼困難和重擔,天父都知道,祂能扶持我們,按著我們的需要,加添我們的力量去承擔.
\newline
\newline
神對以東人說:「你因狂傲自欺」.驕傲和自欺,不單是以東人的問題,也是基督徒的問題,在我認識的基督徒中,誰敢說自己是夠謙卑的?雖然儘管追求謙卑,然而未達目標,卻又驕傲起來了!豈不是一個自欺的人嗎?各位弟兄姊妹,我相信基督徒的謙卑,自己很難看得出來；只有耶穌和別人,才能看出來,說出來.如果我們要追求謙卑,就必須在服侍神和服事人的態度上養成習慣,好像寫這書的俄巴底亞一樣.從第一至廿一節,一直在講預言,除了第一節提到自己得到默示外,俄巴底亞一次也沒有提及自己.他的思想集中在耶和華和以東人,他就沒有時間提到自己和高舉自己了.從先知的品格來看,信徒事奉的精神愈大,謙卑就愈大；這是耶和華的偉大,是以東人的需要.所以我們要效法這位先知,承認自己的異象,恩賜和聖潔,都是從神而來的.耶和華不給,我們就不得到；所給,也是出於祂的恩典.自己的缺點,不容易看見,需要別人的指教,才能徹底的對付；求耶穌的寶血洗淨,求聖靈的大能克服；這樣在成為聖潔上,才有所長進.
\newline
\newline
我們注意「心裡說」(三節)這三個字,耶和華所知道的,不單是人的所作所為,也是他們的思念.我們的心懷意念,所說,所做,所想,天父全部知道.以東人的作為和心中境況,耶和華都清楚；生活中,表現中,心中怎樣,耶和華都清楚.祂知怎樣對待他們,如果我們真正信主,心裡的思想和外在的表現,都一樣得到祂的喜悅.耶和華的無所不知,能帶給我們極大的安慰；因為我們心裡的一切感覺,都是祂所看見的.我們的盼望和懼怕,我們的喜樂和憂愁,都是愛我們的天父所知道的,在我們尚未認識祂以前,祂已開始為我們安排,這是基督徒很基本的福氣.
\newline
\newline
(二)祂的無所不在
\newline
\newline
這些以東人,以為能逃避神,但是因為耶和華的偉大,便逃不了.「你雖如大鷹高飛,在星宿之間搭窩.」(四節),昨天看過「惡有惡報,善有善報.」這句話,而中文有另一句話「善惡到頭終有報,高飛遠走也難逃!」,但是這些以東人不相信耶和華的無所不在,他們以為山頂上的居所和山寨,名字叫「彼特拉」(PETRA),任何人不容易爬上去攻擊他們,除非是神和他們的仇敵,想出特別的辦法,除非敵軍如大鷹高飛突然臨到,否則他們就可以平安地永遠住在那裡,不怕任何仇敵了.但是他們忘記了一件事,即使他們跑到宇宙任何一個角落；這位無所不在的神,仍然可以追趕他們.耶和華要施行審判的時候,能飛得比太空船還高,這位無所不知,無所不在的耶和華,祂所統治的範圍是包括全世界的,所有的高山,海島,都在祂主宰之中.相信這對我們基督徒很有幫助,無論在本地或海外,在甚麼環境和需要.第四節上半節的耶和華,能夠加添我們的力量,解決我們的問題,引導我們的路途,使用我們的見證.
\newline
\newline
有一位無神論者,甚麼宗教信仰也不贊成,不承認神的存在,當他遇見一位信主的小孩子,就對他說:「如果你能告訴我神在哪裡,我就給你一粒糖.」小孩子回答:「我的神偉大無比,祂創造天地,祂也充滿這個宇宙；如果你能告訴我,甚麼地方沒有神的,我就給你兩粒糖.」這小孩子的信仰,又堅定又準確.基督徒相信神是無所不在的,也可以向海外的朋友作見證,寄信的時候,附上福音單張；我們相信神怎樣在這裡向我們說話,也可以在那裡向他們說話.因為本地,海外,白晝,黑夜,對神來說都是一樣,都在祂能感動人的範圍內.在英國有一位愛主的姊妹,她丈夫年青時,在非洲工作和居住,離英國很遠.當他讀到浪子故事的比喻時,便受到聖靈的感動信主；因為他的妻子寄給他福音單張,亦多為他祈禱.讓我從自己的經驗上來舉例,有一天當我從鄉村一間教會講完道後,遇見一位年青的農夫；我與他談話的時候,他提到自己雖是基督徒,但不會祈禱.我問他:「你的禱告有困難,你清楚耶穌是你的救主嗎?」他說自己不清楚,那天晚上我不能多談,因為若趕不上尾班巴士,便要在當地住宿,我請他留下姓名地址,回家給他寫信.第二天我便寫信給他,向他介紹福音,也寫出他能接受耶穌的一篇禱文；信寫好後,便多為他禱告.結果下一次到他的教會講道,他在教會門外等候著我,滿臉光輝地說:「讀完你那封信,我便使用你寫給我的禱文,來接受耶穌作我個人的救主.現在不單禱告沒有問題,我也很喜歡作見證,下次你到其他教會講道的時候,我也想與你同行,見證主怎樣救贖了我.」
\newline
\newline
那位古代基督徒奧古斯丁,年青的時候,離家往羅馬教書.他有愛主的母親,恐怕他離開主,一定會不信主；奧古斯丁年青時,行為也很腐敗,連他爸爸也不是基督徒.當他離家後,沒有母親和好友同在,便更加走錯路.但奧古斯丁的母親認識耶和華的無所不在,多為他禱告.想不到他的兒子,到達羅馬後,與信耶穌的人做朋友,看見他們的生活,便相信主,後來成為一個偉大的基督徒.在家裡,在羅馬,也是耶和華感動人的範圍.弟兄姊妹!當我們為未信主的親友祈禱時,應該要心裡注意耶和華的無所不在,同時也相信；祂有數不完的方法和器皿,領我們所愛的人信主.
\newline
\newline
(三)祂的無所不能
\newline
\newline
第四節提到以東人在山頂,所建立的彼特拉；他們在星宿之間所搭的窩,不容易攻破的山寨.但四節下半節:「我必從那裡拉下你來」,這不是人誇口說的話,人若誇口不一定能做得到.既是無所不能的神說,無論是責備,安慰,預言,應許,就句句都可靠,句句不落空.在應許方面,有人計算聖經的應許,有三萬多次,適合你我各方面的需要.當我們靈性不冷不熱,沒有長進的時候；問題不在神方面,因為祂有這麼多應許等候我們去支取.
\newline
\newline
在佈道的工作上,人的心好像以東人,所住的山寨像彼特拉一樣；很剛硬,一點也不易為耶穌得著.但是我們的神有能力和方法,使用我們的禱告和見證,攻破堅固的營壘,使罪人向耶穌歸降順服祂,一生作耶穌的僕人.有一位反對福音的無神論者,他的態度很像俄巴底亞時代的以東人,住在自己驕傲的高山.有一天他參加一次佈道會,目的不是聽福音；乃是找機會與講員辯論.散會時向講員說:「現在我向你挑戰,你今天所講的,我認為沒有理由,故讓我們公開討論,讓這些人看清楚,信耶穌的人有用與否?」講員便接納他的挑戰,便宣佈自己的信仰,先讀:「按著定命,人人都有一死,死後且有審判」(來九27),那人便笑起來,我不相信聖經的話,你的學問有限,因為你沒有為你的信仰提出證據,講師便說:「人的話很有限,所以我便提出神的話.」無神論者不滿意他的答案,便走了.他不知不覺將那節聖經背熟了.他離開後,講員便與弟兄姊妹為這位無神論者失喪的靈魂祈禱.結果當無神論者步行回鄉下的時候,月亮發出亮光；路的兩旁都有水池,水池裡有很多青蛙,青蛙很喜歡月亮,因此聲音很大.當他回家時,沒有忘記自己所背誦的經文,青蛙的聲音好像說成「審判」兩個字,無神論者走一步,青蛙便說:「審判,審判」.這人一到家,便跪下接受耶穌做個人的救主,這是一個真實的故事.願一切榮耀歸給這無所不知,無所不在,無所不能的神.
\newline
\newline


\newpage

\section{第56屆~港九培靈會~林克己~第四講　耶和華的提醒}
\label{sec:373}
\textbf{第四講　耶和華的提醒}
\newline
\newline
連結: \href{https://www.hkbibleconference.org/session-message/view/373}{\texttt{https://www.hkbibleconference.org/session-message/view/373}}
\newline
\newline
\hyperref[sec:372]{< < < PREV SERMON < < <}
~
\hyperref[sec:toc]{[返目錄]}
~
\hyperref[sec:375]{> > > NEXT SERMON > > >}
\newline
\newline
俄巴底亞書 1:5-1:6
\newline
\begin{longtable}{cl}
\hline
\hline
章節 & 經文 (和合本修訂版)\\
\hline
1:5 & \begin{tabularx}{0.7\textwidth}{X} 盜賊若來到你那裡，小偷夜間來到，豈不是只偷他們所需要的嗎？摘葡萄的若來到你那裡，豈不留下幾串嗎？你竟全然滅絕！ \end{tabularx} \\ \\ \relax
1:6 & \begin{tabularx}{0.7\textwidth}{X} 以掃遭到搜查，他隱藏的寶物竟被尋出！ \end{tabularx} \\ \\
[1ex]
\hline
\hline
\end{longtable}
神怎樣用自己的聖潔和公義,來警告這些以東人,不讓他們永遠逼迫祂的選民,祂怎樣在以東人歷史上,要伸出大能的手,用以東人的仇敵,在這些人的生活中,來施行審判.
\newline
\newline
「盜賊若來在你那裡,或強盜夜間而來,你何竟被剪除,豈不偷竊直到夠了呢?摘葡萄的若來到你那裡,豈不剩下些葡萄呢?以掃的隱密處,何竟被搜尋,他隱藏的寶物,何竟被查出呢?」(五六節)
\newline
\newline
這裡提醒以東人,耶和華的審判,好像夜間的強盜和盜賊一樣,我們每逢讀到「強盜」的時候,便怕起來!一九六零年我剛剛作宣教士,在新加坡學國語,為了回答別人詢問我的職業,我就學了「講道的」這三個字.可是有時候,我的發音不準；結果有一天,有位姊妹問我做甚麼工作,我說了作「強盜的」.所以作宣教士需要能適應和有幽默感,雖然講得不準,但仍要繼續下去.有句話:「活到老,學到老.」在我們的日常生活,事奉和屬靈方面,基督徒上天堂前,一生都是主的學生；求主幫助我們,天天學習主的功課.主在我們想不到的時候,便好像夜間的盜賊臨到我們中間.我們所以不會感到慚愧,因為我們已有聖潔的生活和事奉；便能夠快樂地,平安地向祂交賬.聽祂對我們說:「你這又忠心又良善的僕人,你在不多的事上有忠心,我要把許多的事派你管理,你可以進來享受你主人的快樂.」
\newline
\newline
「所以你們要儆醒,因為不知道你們的主是哪一天來到.家主若知道幾更天有賊來,就必儆醒,不容人挖透房屋,這是你們所知道的,所以你們也要預備,因為你們想不到的時候人子就來了.」(太廿四42-44)另有一段聖經:「你要寫信給撒狄教會的使者說,那有神的七靈和七星的說:我知道你的行為,按名你是活的,其實是死的.你要儆醒!堅固那剩下將要衰微的,因我見你的行為,在我神面前,沒有一樣是完全的.所以要回想你是怎樣的領受,怎樣聽見的,又要遵守,並要悔改.若不儆醒,我必臨到你那裡,如同賊一樣,我幾時臨到,你也決不能知道.然而在撒狄你還有幾名是未曾污穢自己衣服的,他們要穿白衣與我同行,因為他們是配得過的.」(啟三1-4)考古學家對於撒狄的教會,有足夠的證據,因他們領受了主的話,肯徹底悔改火熱起來!歷史告訴我們,這個教會復興後,沒有人能敵擋!考古學家在一九一○至一九一四年,發現撒狄城的神廟,在牆壁上刻了一個十字架,表明不能抵擋福音的神廟,反而成為基督徒的聚會所.」
\newline
\newline
今天我們把這三段經文放在一起看(俄5-6節,太廿四42-44,啟三1-4)來研究三件事,就是:賊的來臨,賊的方法,賊的目的.在這個比方裡,要注意耶和華的提醒,不單提醒以前的以東人,也是提醒今日的基督徒.
\newline
\newline
(一)賊的來臨
\newline
\newline
俄巴底亞指出盜賊若來,而馬太福音說賊來是人不能知道的,(啟三章)說,賊幾時來到,你們不能知道.難怪俄巴底亞提醒以東人,他們需要神的提醒；我們八十年代的基督徒,也需要神的提醒.因為施行審判的耶穌基督,幾時臨到我們中間,真好像夜間來的賊一樣；在我們想不到,更加是我們絕不能知道.(俄5-6),不單是提醒,也是預言,無所不知的耶和華,雖然清楚地說自己為甚麼來和怎樣來；但這些驕傲的以東人,不怕也不承認這件事.他們不相信「善有善報,惡有惡報；不是不報,日子未到.」他們好像幾百年,一位住在德國的貴族婦人,她像俄巴底亞時代的以東人一樣,心裡驕傲,對宗教信仰沒有興趣.有一天她說:「我是自由思想家,不相信神的存在,更不相信將來的審判.但如果有一位神存在的話,我現在向祂挑戰,我現在向祂誇口；我死的時候,要埋葬在一個特別的墳墓裡,這個墳墓要用鐵和水泥做的,是特別堅固.我相信神永遠也無法打開這個墳墓,叫我站在祂面前,受祂甚麼的審判.」幾十年後,貴婦死去了,按照她的計劃,被埋葬在這個特別堅固的墳墓裡.神好像不接受她的挑戰,很多年過去,墳墓還很平安,想不到有一天,附近有一棵樹,丟下一些種子,落在墳墓細小的裂縫中,在那裡扎根,那些根便把貴婦的墳墓張開,再過幾十年,從這個墳裡長出來的,有一棵很大很美的樹；它提醒所有經過的人,神所預言的應許必定成就.要打開貴婦人的墳墓,神不需要原子彈的爆炸力,一粒小小的種子便夠了.
\newline
\newline
向以東人的審判,在歷史上主前第三世紀的人,有亞拉伯拿巴提人,來攻擊以東人,好像賊一樣的爬上去.以東人著名的彼特拉,雖建在很高的山上,有三千七百呎高,容易看到爬上來的仇敵；但是俄巴底亞所說的不錯,賊能夠夜間而來,彼特拉就被拿巴提人攻取；這些壓迫猶大的以東人,就在歷史上得到應得的報應.這件事讓耶和華的提醒應用在基督徒的生活中.我們的主耶穌,好像夜間來的賊；在我們想不到的時候,突然回來；按照聖經預言,第二次來到這世界,作罪人和聖徒的審判主.聖經提醒我們,若我們的行為不聖潔,我們所走的路不得祂的喜悅,雖然祂還要接我們上天堂；但當我們迎接祂的時候,就感到慚愧.難怪保羅寫信給腓立比教會說:「只要你們行事為人與基督的福音相稱.」(腓一27)這樣才得主的喜悅,見證也被祂使用.
\newline
\newline
有一間學校的校長,為了要學生保持書桌清潔與整齊,就舉辦了一個整潔比賽.但不公佈檢查清潔的日期,也不讓任何人知道；校長便作突擊的檢查清潔的書桌,整潔的學生便可以得到他的獎品.這比賽消息報告後,有一個學生對他的同學說:「我真盼望能夠得到這份獎品,我必須收拾好我的東西；每星期一次把書桌弄得很清潔整齊,這樣來作好準備.」他的同學回答說:「校長是在我們想不到,不能知道的時候,進來檢查我們的書桌,你每星期只有一次清潔,夠不夠呢?如果校長在你五六天內沒有收拾的時間進來檢查,你就難以得到獎品了.」那學生回答說:「你說的實在不錯,每星期一次是不夠的,那麼我每天收拾一次就好了.」那同學又說:「一天一次夠嗎?如果你早上收拾,以後將筆記簿堆上來,下午校長走進來,那你怎麼辦?」那同學思想考慮同學的話,便笑起來說:「我看只有一個有效的方法,就是一天到晚保持整齊清潔；免得我的校長,在我想不到的時候走進來.如果我沒有儆醒,他一進來便知道,我是一個沒有追求清潔整齊的人.」各位弟兄姊妹!我們等候主的再來,豈不也像那學生嗎?何等需要追求聖潔的生活,討主的喜悅,等候祂再來呀!
\newline
\newline
我們為耶穌所做的工,到底怎麼樣?做得有目的嗎?做的認真嗎?做得忠心嗎?如果我們在工作上不可靠,在行為上不聖潔,教會的工作馬馬虎虎；主突然的來臨,豈不是很可怕嗎?我們的主,在我們不儆醒的時候,忽然臨到施行審判.我們的見證和工作,都要經過祂考驗的火,還能存得住嗎?
\newline
\newline
(二)賊的方法
\newline
\newline
「盜賊若來在你那裡,或強盜夜間而來,何竟被剪除,豈不偷竊直到夠了嗎?摘葡萄的若來到你那裡,豈不剩下些葡萄嗎?」(五節),我們特別注意「豈不偷竊直到夠了嗎?」和「豈不剩下些葡萄嗎」,這裡告訴我們,賊光顧我們的家,通常不會取去我們全部的東西；賊知道他的時間和力量有限,如果他夜間進來,當然不會等到黎明才走.所以賊注意的,是那些又寶貴又容易攜帶的東西,卻剩了普通的東西.以東人注意屬世界的東西,輕視宗教的事,不渴慕屬靈真理,這方面怎樣提醒我們呢?保羅寫信給哥羅西的信徒說:「所以你們若真與基督一同復活；就當求在上面的事,那裡有基督坐在神的右邊.你們要思念上面的事,不要思念地上的事.」(西三1-2)世界上的事是暫時的,只有天上的事才是永遠的,我們所忙碌的,只有為主所做的事,才能存到永永遠遠.難怪主耶穌在登山寶訓中,教導我們一些挑戰性原則:「為自己積攢財寶在地上,地上有蟲子咬,能鏽壞,也有賊挖窟窿來偷.只有積攢財寶在天上,天上沒有蟲子咬,不能鏽壞,也沒有賊挖窟窿來偷.」(太六19-21)賊在我們想不到的時候,挖窟窿來偷,雖然剩下葡萄,賊把不方便帶走的東西剩下.但如果我們所寶貴的,全部或部份是在地上的,我們便很可憐和傷心,像那些以東人一樣.主耶穌解釋:「因為你的財寶在那裡,你的心也在那裡.」(太六22)主耶穌的話很有挑戰性,因為只要知道自己思想的是甚麼,就知道自己最寶貴的是甚麼.我們的思想,好像指南針,清楚指出我們生活的目標和方向,我們的生活和思想,最寶貴的是甚麼.我們所寶貴的,是地上的一件事或一個人?還是那位在天上,為我們受死復活升天,愛我們的耶穌呢?
\newline
\newline
有一位小女孩,很可愛,上主日學的時候,聽見這個賊和財寶的比喻；聽見後,就將耶穌的結論背熟,做她那天的金句.小妹妹回家後,也背誦給她媽媽聽.那天晚上,小妹妹睡覺的時候,她發夢,夢見自己在街上漫步,很奇怪!無論她往哪裡,都看見人的心掛在外面.她經過一間單車店,看見一個個的心掛起來,她想這些人很喜歡去旅遊.小妹妹向前走,有一間店裡面有很多好吃的東西,掛起來的心很多,小妹妹就說:「這是最喜歡吃東西的人.」她多走前幾步,第三間店有許多小孩子的玩具,小妹妹停步看看那些玩具,那些在玩具店掛起來的心是最多的.小妹妹說:「喜歡做工,讀書的,都不多,但喜歡玩樂的,十分多.」她就繼續往前行,路的兩旁再沒有店舖了,她就舉頭看；奇怪上面也有人的心掛起來,但是不很多,小妹妹就記在主日學所唸的金句:「因為你的財寶在那裡,你的心也在那裡.」她對自己說:「我的心尚未掛起來,但主耶穌這樣愛我,我知道要怎樣去愛祂,以祂為我的至寶；我的心要永遠掛在天上,多多思念我的救主,一生熱心跟隨祂.」小妹妹立定主意後,她就睡醒了.這件事提醒我們,在我們生活中；天天把心掛在天上,這樣永遠不怕賊的來臨和賊的偷竊了.
\newline
\newline
(三)賊的目的
\newline
\newline
「以掃的隱密處,何竟被搜尋,他隱藏的寶物,何竟被查出」(六節)先知俄巴底亞在這裡所描寫的賊是誰呢?相信是亞拉伯和拿巴提的軍隊.在主前五百年至三百年,常常與以東人打仗；這些仇敵作賊,來偷以東人隱藏的寶物,因為他們有力量和時間.在二百年間常常來,最後以東人抵擋不住,他們的山寨彼特拉被拿巴提人攻破!以東人真可憐,俄巴底亞說:「你何竟被剪除?從高處拉下來」,在歷史上他們便不得再復興了,這要證實了先知從耶和華所得的默示,見證聖經的來源和可靠性.因為若不是靠耶和華的啟示,先知可以開口這樣講嗎?亞拉伯和拿巴提人與普通的賊是不同,他們不單偷得愈多愈好,正如先知說:「豈不偷竊直到夠了嗎?」(五節),第六節說:他們把以東人甚麼東西也搜尋,隱藏的東西寶物也被搜去了,耶和華藉著以東的敵人所行的審判,這樣可怕和嚴厲!弟兄姊妹!這是耶和華的提醒,藉俄巴底亞書第五,六節,提醒祂在聖經和歷史的審判,何等的聖潔,何等的公道.在這個恩典的時代中,是福音和差傳的時代,當審判官還未來到,主坐在祂白色的大寶座之前；我們信祂之人,千萬不要錯過傳福音的好機會；領更多的罪人歸向耶穌,免得我們所愛的救主,在我們還沒準備好的時候,好像夜間來的賊一樣.主第二次到世上來的時候,是在我們想不到,絕不能知道的時候；忽然回來與我們算賬,我們便會感到慚愧,因我們沒有預備聖潔和熱心的行為；不能帶著無虧的良心,來見主的面.讓我們今天用謙卑的心,來領受耶和華的提醒,在我們的生活中活出來.
\newline
\newline


\newpage

\section{第56屆~港九培靈會~林克己~第五講　耶和華的警告}
\label{sec:375}
\textbf{第五講　耶和華的警告}
\newline
\newline
連結: \href{https://www.hkbibleconference.org/session-message/view/375}{\texttt{https://www.hkbibleconference.org/session-message/view/375}}
\newline
\newline
\hyperref[sec:373]{< < < PREV SERMON < < <}
~
\hyperref[sec:toc]{[返目錄]}
~
\hyperref[sec:376]{> > > NEXT SERMON > > >}
\newline
\newline
俄巴底亞書 1:7-1:7
\newline
\begin{longtable}{cl}
\hline
\hline
章節 & 經文 (和合本修訂版)\\
\hline
1:7 & \begin{tabularx}{0.7\textwidth}{X} 與你結盟的都驅趕你，直到邊界，與你和好的欺騙你，勝過你，吃你飯的人設下圈套陷害你---他卻毫無聰明。 \end{tabularx} \\ \\
[1ex]
\hline
\hline
\end{longtable}
「與你結盟的,都送你上路,直到交界.與你和好的,欺騙你,且勝過你；與你一同吃飯的,設下網羅陷害你,在你心裡毫無聰明.」
\newline
\newline
今天的分題是耶和華的警告,我們分三方面來研究；第一,有一種假的倚靠.第二,有一種假的朋友.第三,有一種假的智慧.這三樣假的事,一同稱為耶和華的警告,雖是向以東人說的,但是對你我都有教訓.我們也要靠著主,儘可能躲避開.
\newline
\newline
(一)一種假的依靠
\newline
\newline
在東方有巴比倫和亞拉伯人,在南方有埃及人,主前六世紀,這三種人的勢力很大!好像今天的美國,蘇聯和中國一樣.但是這些心裡驕傲的以東人,甚麼也不怕；他們心裡有倚靠,他們在山上建立的彼特拉,仇敵很難攻上來.「與你結盟的」(七節),就是住在迦南地的摩押人和亞捫人.在地理方面,以東地有山,有谷,高高低低,是外來人不易攻打的地帶.在歷史方面,這些與以東結盟的摩押人和亞捫人,很久以前,已認識那些地方.假如有仇敵在東方,南方進攻的話,以東人,摩押人和亞捫人就會聯合對付.他們會想辦法,或者爬上高山,或者伏在平原,或出來打仗,或隱藏洞穴；在自己最適當的時間和地方,來與敵人交戰,戰勝敵人,從人的角度來說,以東人這個倚靠,是很不錯的；從地理和歷史看,這樣為自己打算,是相當聰明的.這樣倚靠人,不需要神的幫助.以東人這種態度,叫我想起古代的尼尼微人；他們雖然有宗教信仰,跟以東人不同,但他們所拜的神很古怪.他們用一塊四十噸的大石頭,雕刻偶像立在尼尼微的城門,形像很特別.這神像有牛的身體,鳥的翅膀,人的頭；他們這個像高舉人的偉大,力量和聰明.牛的氣力很大,但人用頭腦想出的計劃,比牛更有效力；飛鳥飛得又快又高,但人用自己的腦筋解決的問題,比飛鳥更快；但他們所拜古怪的偶像,不過是假的倚靠而已.他們注重人,高舉人,以人作神來敬拜,這是尼尼微人的錯誤.按照聖經歷史家的發現,主前第六世紀的以東人,甚麼敬拜的對象也沒有；無論遇見任何困難,他們以屬世和屬世界的辦法來解決.從以東人直到今天,有許多人上了魔鬼的當；生活中只有假的倚靠,做任何事,說話,看法,辦法,都是出於人,而不倚靠神.很多信心弱,沒有信心禱告的基督徒也包括在內,不知不覺倚靠人,而停止仰望神.
\newline
\newline
讓我們看幾節聖經來對付自己的弱點,同時領受耶和華的警告,如果在我們的日常生活,家庭生活,屬靈生活和事奉生活中,有難解決的問題；讓我們注意愛我們的耶和華在聖經中,為我們預備的原則和方法,支取從上面來的力量和幫助.以賽亞先知用的兩個比方,特別適合我們的需要,「你或向左,或向右,你必聽見後邊有聲音說,這是正路,要行在其間.」(賽三十21)這節聖經談及耶和華的引導,有應許也有原則,在我們不知不覺或故意走迷路的時候；向左向右,離開主為我們安排的道路,不遵行祂為我們生命定下的旨意.我們必定有一個思想,感動或光景；使我們後面聽見聲音對我們說,這是正路要行在其間.兄弟未做父親,還沒有孩子之前,不大明白這節經文.我時常問,天父帶領我們,為甚麼要在後面對我們說話?為甚麼不在前面領路?若天父在前面引導,我們就能放心,在後面引領,好像沒有理由.我的兩個女兒在很小的時候,她們很喜歡在鄉下地方玩耍,隨便走走,當她們很開心的時候,就跑到我前面；有危險的地方,我便大聲呼喊:「你們要停步!不可再往前走,待我趕上你們；因你們不知道,前面有危險!」她們幼年時,我也不記得有多少次,是這樣來教導她們；要求她們停步,先觀察才往前走.弟兄姊妹!相信耶和華的警告,也是這樣,祂為愛我們的緣故,不想我們有一種假的倚靠.祂就作我們的父親,按著我們的需要管教我們.每逢我們要決定前途的時候,總是缺少智慧；在我們的思想和計劃中,跑得太快雖已用禱告把問題交托給神,但仍會有傾左或傾右的危險.我們這位無所不知的神,知道怎樣從後面指教自己的孩子.祂可能用一件事,一個人或直接對我們講說話；使我們認識祂的聲音來,這樣從後面叫我們停步；等候再思想一下,多禱告一次,就認識這是正路,要行在其間.祂用這種原則和應許,免得我們陷落在一種假的倚靠中,好像俄巴底亞書的以東人一樣.也免得我們的思想,計劃,好像尼尼微人所拜的牛神；因為有翅膀飛得很快,沒有認定主的時間和旨意；便不能和天父一步步的同行,不能避免一切的危險了.
\newline
\newline
先知再用另一個比方:「你豈不曾知道嗎?你豈不曾聽見嗎?永在的神耶和華創造地極的主；並不疲乏,也不困倦,他的智慧無法測度.疲乏的,他賜能力,軟弱的,他加力量；就是少年人也要疲乏困倦,強壯的也必全然跌倒.但那等候耶和華的,必重得力；他們必如鷹展翅上騰,他們奔跑,卻不困倦,行走卻不疲乏.」(賽四十29-31)先知以賽亞在31節說:「他們必如鷹展翅上騰」,好像(俄4節)「你雖如大鷹高飛」,但卻有些不同.今日的觀點,這裡有一個比較和分別；有以東人的驕傲和基督徒的能力,擺在我們面前.同時都用老鷹高飛為比方,使我們能分別假的倚靠,和真的倚靠.請注意:「展翅上騰」四個字,一隻大鷹飛到高處,所倚靠的是甚麼力量,是自己的能力嗎?還是耶和華創造世界後,所預備的風和空氣?我們只要觀看,一隻大鷹在空中高飛,便知道祂有耶和華所賜的聰明,有時候祂不震動翅膀,仍可一直高飛.祂已學會展開翅膀後,不再倚靠自己的任何力量,只要倚靠風和空氣；風氣在那裡冷,在那裡熱,都可以達到牠一直向上昇的目的.聖經裡的風是象徵聖靈,代表聖靈的充滿和能力.我們不要效法那些屬血氣的以東人,願主幫助我們,拋棄假的倚靠；展開信心的翅膀,自己被聖靈充滿.脫離屬世而過屬靈的生活,飛到與耶和華同行的高原.這樣便「重新得力」,這四個字在希伯來文,基本意思是「將人的軟弱換來神的能力」,相信這是我們取大的需要,讓我們今天注意神的警告,在假的倚靠方面,免得我們好像以東人,心裡驕傲,倚靠自己的力量,一樣上了魔鬼的當,得不著靈性生活的秘訣,就是展開信心的翅膀,被聖靈充滿之後,就一直高飛,脫離屬血氣屬世界的捆綁,有耶穌能稱讚,能使用的生活和見證.
\newline
\newline
(二)一種假的朋友
\newline
\newline
「與你結盟的......在你心裡毫無聰明.」(七節)表面和外貌看,以東人所倚靠的摩押人和亞捫人,都很不錯,對他們很客氣,注意「送你上路直到交界」這幾個字,以東人的大使和代表去拜訪的時候,摩押人和亞捫人都很歡迎招待及尊重他們,以東人便不儆醒,與這些人結盟.中國人有句話說:「畫虎畫皮難畫骨,知人口面不知心」,若這些以東人聽過這些話便會小心一點.亞捫人和摩押人,遇到以東人後,便點頭,談天說笑,一起吃飯,好像很好,很靠得住.但他們心裡很狡猾,雖然送他們上路,直到邊界,好像捨不得他們!耶和華警告這些可憐的以東人說:「與你和好的,欺騙你,且勝過你；與你一同吃飯的,設下網羅陷害你.」我們注意「與你和好」幾個字,我們基督徒與古代以東人不同；但在交友事上,需要人的指教和神的幫助.求天父幫助我們多與信主的朋友在一起,雖然他們的年齡,性格,背景,教育等,也許與我們不同；但一樣信靠和見證我們所愛的主.在屬靈的事上,大家都是一家人,基督是我們的家長,在這個屬靈的家庭裡,我們都有作弟兄姊妹的身份和責任,要學習彼此接納和了解.彼此尊重,代禱和相愛,這些事我們很熟悉；但我們要在日常生活和教會生活中實行出來.將信徒間的關係表現出來,今天固然重要,將來更加重要.為甚麼呢?簡單來說,今天我們享受的宗教信仰自由,可以聚會,有全時間的牧師和傳道人,公開自由地領會；這種情況,將來不一定可以享受得到.將來可能受逼迫,沒有公開聚會的自由,也不再有全時間的傳道人,可以自由的牧養教會.既然有這種可能性,我們應該怎樣作準備呢?第一,要養成天天讀經,禱告的習慣,學習用聖經和禱告來培養自己的靈性；將來雖沒有牧師來牧養和探訪,仍能在真道上站立得穩,日日長進.第二,我們要將聖經的應許,及有鼓勵性的經文；背熟藏在心裡,如果有人奪去我們的聖經,甚至被放在監牢裡,手裡沒有聖經,仍然能默想神的話:用我們背熟的聖經,作為我們天天的力量和亮光.第三,教會要注重聖徒的交通和溝通,如果一間教會有幾百位或幾十位會友,除了主日崇拜之外,全體會友也可以學習彼此的扶助；就是在基督徒家中分小組,用小組禱告查經,彼此鼓勵和造就.同時一起學習怎樣向外邦人見證,為未信主者禱告.這樣在我們住的地方,本地或海外,在香港或英國,將來即或沒有宗教的自由,相信教會仍有光明的前途.雖然有逼迫,還能站立得穩.看教會歷史,逼迫不單不會消滅教會,反而激勵,堅固,熬煉教會,信耶穌的人也增多.但我們在宗教自由時代中,要作聰明人,多與基督徒作朋友,作忍受苦難的準備.「兩個人總比一個人好,因為二人勞碌同得美好的果效.若是跌倒,這人可以扶起他的同伴,若是孤身跌倒,沒有別人扶起他來,這人就有禍了,再者,二人同睡,就都暖和,一人獨睡,怎能暖和呢?有人攻勝孤身一人,若有二人便能敵擋他,三股合成的繩子,不容易折斷.」(傳四9-12)這裡有同工,同行,同睡,同戰,均注重信徒之間彼此認識和扶持的重要性.無論在大聚會或小組,要學習團體生活,一齊服事主,一同跟隨主,在這個基督作家長的團體裡,真要彼此關顧弟兄姊妹.
\newline
\newline
另一方面,我們怎樣與非基督徒做朋友呢?聖經為我們的益處,也減少彼此信仰的磨擦,幫助我們設立一個榮耀神的基督化家庭,就禁止我們與未信者結婚.但與非基督徒做朋友,應怎樣解決呢?我們很易受他們影響,而漸漸倒退冷淡,我們也少有機會向他們作見證.最好的方法,要保持平衡,基督徒朋友多,非基督徒的朋友也多；這樣可以保持熱心,也可以常有機會作見證了.十幾年前,兄弟在馬來西亞一次佈道會中,領一位年青教員信主,這位弟兄信主後便非常熱心.參加每樣的聚會,一星期的每一晚上,都與基督徒在一起.我對他說:「你常與基督徒在一起,你未信主的朋友多不多?」他回答說:「感謝主,有很多.」我再問:「你最近向未信主的朋友作見證怎樣?」他很誠實的回答說:「感謝主,有很多.」我再問:「你最近向未信主的朋友作見證怎樣?」他很誠實的回答說:「這恐怕是我的問題,可能常與基督徒朋友在一起,我們未信主的朋友就愈來愈少了.」大家經過討論和禱告後,他就對我說:「我看少參加幾次聚會,多作幾次見證,這樣比較好,我可以追求長進,同時被主使用.」兩個星期後,這位青年教員,很高興對我說:「感謝主,雖然少了幾次聚會,多了幾次見證；但我今天領了一位同事信主,我看見有基督徒朋友和未信朋友的重要.」求主幫助我們,不效法以東人,他們不會交朋友,也不求耶和華的引導,結果被假的朋友拉去了,他們雖然知道,但不領受耶和華的警告.
\newline
\newline
(三)一種假的智慧
\newline
\newline
主前六世紀的以東人,有一個很大的弱點,非常狂傲,強調自己的本事和智慧,但不是耶和華關於他們所發表的意見,無所不知的耶和華,鑑察他們的心思後,對以東說:「在你心裡毫無聰明.」(七節)智慧的問題,我們今天不很留意,原因很簡單.第一,我們的時間有限.第二,耶和華的智慧,明天思想這問題,我們比較人的智慧和神的智慧,作深入研究和討論.
\newline
\newline
著名改革家馬丁路得說:「在屬靈的事上,一個十幾歲的女孩子,沒有受過良好教育,不過在家裡掃地板,洗衣服；但這女孩子只要重生得救,結果她知道關於神的事情,比一位不信主的大學教授還多.因為聖靈將屬天的智慧賜給她.」俄巴底亞書七節提到耶和華的警告,願主將這警告銘刻在我們心裡,使我們不忘記,也成為我們屬靈生活的一部份,願一切榮耀歸給主.
\newline
\newline


\newpage

\section{第56屆~港九培靈會~林克己~第六講　耶和華的智慧}
\label{sec:376}
\textbf{第六講　耶和華的智慧}
\newline
\newline
連結: \href{https://www.hkbibleconference.org/session-message/view/376}{\texttt{https://www.hkbibleconference.org/session-message/view/376}}
\newline
\newline
\hyperref[sec:375]{< < < PREV SERMON < < <}
~
\hyperref[sec:toc]{[返目錄]}
~
\hyperref[sec:378]{> > > NEXT SERMON > > >}
\newline
\newline
俄巴底亞書 1:8-1:9
\newline
\begin{longtable}{cl}
\hline
\hline
章節 & 經文 (和合本修訂版)\\
\hline
1:8 & \begin{tabularx}{0.7\textwidth}{X} 到那日，我豈不從以東除滅智慧人？從以掃山除滅聰明人？這是耶和華說的。 \end{tabularx} \\ \\ \relax
1:9 & \begin{tabularx}{0.7\textwidth}{X} 提幔哪，你的勇士必驚惶，以致以掃山的人都被殺戮剪除。 \end{tabularx} \\ \\
[1ex]
\hline
\hline
\end{longtable}


\newpage

\section{第56屆~港九培靈會~林克己~第七講　耶和華的責備}
\label{sec:378}
\textbf{第七講　耶和華的責備}
\newline
\newline
連結: \href{https://www.hkbibleconference.org/session-message/view/378}{\texttt{https://www.hkbibleconference.org/session-message/view/378}}
\newline
\newline
\hyperref[sec:376]{< < < PREV SERMON < < <}
~
\hyperref[sec:toc]{[返目錄]}
~
\hyperref[sec:379]{> > > NEXT SERMON > > >}
\newline
\newline
俄巴底亞書 1:10-1:14
\newline
\begin{longtable}{cl}
\hline
\hline
章節 & 經文 (和合本修訂版)\\
\hline
1:10 & \begin{tabularx}{0.7\textwidth}{X} 因你向兄弟雅各施暴，你必蒙羞，永被剪除。 \end{tabularx} \\ \\ \relax
1:11 & \begin{tabularx}{0.7\textwidth}{X} 當陌生人擄掠雅各的財物，當外邦人進入他的城門，為耶路撒冷抽籤分取財物的日子，你竟站在一旁，像與他們同夥。 \end{tabularx} \\ \\ \relax
1:12 & \begin{tabularx}{0.7\textwidth}{X} 你兄弟遭難的日子，你不該瞪著眼看；猶大人被滅的日子，你不該幸災樂禍；他們遭難的日子，你不該說狂傲的話。 \end{tabularx} \\ \\ \relax
1:13 & \begin{tabularx}{0.7\textwidth}{X} 我子民遭災的日子，你不該進他們的城門；他們遭災的日子，你不該瞪著眼看他們受苦；他們遭災的日子，你不該伸手搶他們的財物。 \end{tabularx} \\ \\ \relax
1:14 & \begin{tabularx}{0.7\textwidth}{X} 他們遭難的日子，你不該站在岔路口剪除他們逃脫的人，你不該交出他們的倖存者。 \end{tabularx} \\ \\
[1ex]
\hline
\hline
\end{longtable}
「因你向兄弟雅各行強暴,羞愧必遮蓋你,你也必永遠斷絕!當外人擄掠雅各的財物,外邦人進入他的城門,為耶路撒冷拈鬮的日子；你竟站在一旁,像與他們同夥,你兄弟遭難的日子,你不當因此歡樂；他們遭難的日子,你不當說狂傲的話；我民遭災的日子,你不當進他們的城門；他們遭災的日子,你不當瞪眼看著他們受苦.他們遭難的日子,你不當伸手搶他們的財物；你不當站在岔路口,剪除他們中間逃脫的；他們遭難的日子,你不當將他們剩下的人交付仇敵.」
\newline
\newline
這段經文是讓我們來研究耶和華的責備,可以從三個角度來看；第一我們應當彼此饒恕,第二我們應當彼此幫助,第三我們應當彼此保護.在分享這三點之前,我們先看一個重要的原則,當我們查考聖經的時候,我們要常常問自己；在所讀的經文裡,有沒有一些字多次出現呢?在(俄10-14節),尤其是(12-14節),有三個字「你不當」常常出現,在12節出現三次,13節出現三次,14節出現二次,一共八次.從這裡,我們可見聖經的主題是耶和華的責備,藉著「你不當」幾個字把握經文的中心思想,乃是彼此饒恕,幫助,保護.這裡信息的對象是屬血氣,屬世界的以東人；他們雖然不承認神,但可見神仍然是偉大的,有權柄能力向他們講說話.在這裡八次責備他們,使他們確實知祂是一位怎樣的神,祂使被創造的人知道要怎樣彼此來對待.
\newline
\newline
從俄巴底亞的預言,我們可以滿心感謝主；也可以在工作上得到鼓勵,因為人的心雖然驕傲,剛硬,好像這些以東人一樣.但神願意用人的說話和文字,將自己的旨意和計劃顯明出來.使這些心裡剛硬的人,在祂面前心存敬畏；使那些對福音毫無興趣的人,聽到他的聲音,而不抵擋祂的感動.蘇格蘭佈道家麥格勒年青時候,實在犯罪,放蕩,他對運動很有興趣,身體亦健康,而且成為蘇格蘭足球隊的球員.但他交上壞朋友,晚上常與他們到酒吧飲酒,時常醉酒,也常用碎玻璃樽與別人打架,又因幾次犯法,坐在囚牢多年,到年老的時候才釋放,他感覺自己浪費了生命,也因酗酒浪費他的錢財,衣服破爛不堪,他一生對宗教沒有興趣.但有一個主日早上,他在街上散步的時候,忽然聽見附近的教會基督徒唱詩的聲音,他就走進禮拜堂,坐在最後面.他雖然看見基督徒的衣服整齊清潔,感到不好意思；但牧師的講道令他大受感動.散會後他就留步與牧師談話,那天早上他很奇妙地重生得救了!他那些朋友都以為他做基督徒不會超過一星期；但神在人心裡的工作,是存到永遠的.五年後,他不但仍作基督徒,也領了一千人歸主,因為神賜給他的佈道恩賜.在蘇格蘭和英國各地,他被邀請主領佈道會,用聖經的真道,用自己的見證傳揚福音.「你不當」幾個字是向未信者說的,鼓勵我們大大傳揚福音,雖然工作有困難,人的心也剛硬,但我們所信靠的耶和華,在他們的生活中,仍能大大的得勝.
\newline
\newline
(一)我們應當彼此饒恕
\newline
\newline
「因你向兄弟雅各行強暴,羞愧必遮蓋你,你也必永遠斷絕.」(十節)這節經文的歷史背景,對我們的品格修養,很有挑戰.(創三十三1-4)有這節經文的背景,以東人的祖宗以掃,一個屬血氣的人有一個優點,就是心裡寬大,饒恕他弟弟雅各.雅各兩次欺騙他的哥哥以掃,以掃卻沒有趁機向他報仇,沒有用四百人來殺害雅各家.可惜在後來歷史上,以掃的子孫以東人,不效法他們祖宗以掃的榜樣；反而心懷仇恨,與猶大的敵人作朋友.幫助他們攻擊雅各的後裔,「向兄弟雅各行強暴」,這件事雖然是幾千年前發生的,但今天會否在我們教會中出現呢?在基督徒中間有沒有因面子和一些小事,多日不說話,也不尊重別人的意見,不願彼此謙讓和饒恕?因為兩個人的不和,便破壞全教會的相親相愛,使主的身體紛爭結黨,好像哥林多教會一樣.你有沒有在教會或團契中,由於缺少了饒恕人的愛心和精神,就失去了原有的團結,如一堆沙這麼可憐?弟兄姊妹!在你心裡有沒有仇恨,在你教會中有沒有你不肯饒恕的弟兄姊妹,這件事不是攔阻你靈命的長進嗎?不是影響你教會所作的見證嗎?願主常常幫助我們檢查自己的內心,主已經饒恕我們,我們也應饒恕別人,靠著主與有嫌隙的弟兄姊妹和好；向某一位肢體謙卑認罪,以後真有愛弟兄姊妹的心.
\newline
\newline
有一天有兩個農夫談話,一個說:「你用馬耕田,實在太辛苦了,你要用力拉韁繩,很快你便疲倦了.讓我教你一個好方法,你該訓練你的馬,好好聽從你的話；命令牠順從你的指揮,向左向右,你的工作就容易的多了.」想不到另一個農夫回答他的朋友說:「你的意思很好,但你知道我實在不願意和這匹馬說話.因在七年以前,牠踢了我一腳,這件事我永遠不忘記,也不原諒牠!因為這匹馬不應該踢我.」弟兄姊妹!那個農夫不是很可憐嗎?因為不能忘記畜牲七年以前向他所踢了一腳的事.那匹馬只做錯一次,他就為這件事害了自己七年,到底受損害的不是那匹馬,而是自己.如果基督徒像那位農夫,心存錯誤的態度；為七年,六年或兩年前發生的事,而不饒恕人；其實不是害了別人,而是害了自己.因為基督徒不肯饒恕別人,就是輕視主耶穌在十字架上所流的寶血；主耶穌受苦,也為那些釘祂的人代禱.耶和華的責備,命我們應當彼此饒恕.
\newline
\newline
(二)我們應當彼此幫助
\newline
\newline
「當外人擄掠雅各的財物......你竟站在一旁」(十一節),多數解經家以為這段經文是論到東方的巴比倫人.在主前五八六年攻破耶路撒冷,以東人雖與猶大人同祖宗所生,是猶大的親戚；注意「兄弟」,猶大人被敵人攻打時,以東人不但不來幫助,也與敵人朋比為奸.我們要將這件事,應用在教會肢體的生活中,弟兄姊妹應彼此的幫助.有人有任何需要,應盡可能供給他們.在主的工作上,若有需要,便盡本份去幫忙,免得別人的重擔太多.
\newline
\newline
我們特別留意「站在一旁」,毫無疑問,這是以東人的最大錯誤,最基本的罪惡.我們應誠實面對這幾個字,因今天在教會裡,站在一旁的基督徒很多,每一個重生得救的人,在教會工作上應該有責任感,逼切感和使命感,但我們很容易站在一旁,請問你是一個參與工作的人,還是站在一旁的人.從前在我牧養一個小教會,我們聚會時有一個問題,唱詩的時候都沒有司琴,我便報告若有弟兄姊妹懂得彈琴,願意幫忙我們,請通知我.雖然報告了兩三次,但無人出來.另外這教會講道和領會時要穿上聖袍,為了解決司琴問題,我唯有穿著聖袍拉手風琴,很不方便和不好看.過了一個時期,我就發現聚會的人中,有三位很會彈琴的弟兄姊妹,其中一位是音樂教師.他們不是沒有恩賜和貢獻,但是不肯司琴,我愈認識他們,便發現他們最大的問題是委身,他們不一定是每主日都來,天氣好,就到郊外或海灘,便沒有時間來聚會；下雨及寒冷時,不便去游泳,才來敬拜主.相信我們有時會像以東人站在一旁的態度,我們需要俄10-14 耶和華的責備.
\newline
\newline
以東人由於不願意饒恕人,也不願意幫助人；雖然別人有需要,自己也有力量和方法去幫忙,但是不動手.這是基督徒面對的一個挑戰,以為對方不可愛,我們就難以效法主去愛他.主怎樣犧牲時間與人交往.做主的工作需要愛心實踐,但我們卻缺少愛心和參與.也許我們會說,我沒有力量沒有愛心的恩賜,只能站在一旁觀望,這樣就會消滅聖靈的感動了!不錯,站在一旁實在是方便和舒適,但這只會讓那些需要救恩的人走向滅亡的道路了!同時也加重了別人的重擔.難怪神八次責備古代的以東人,也責備今天的基督徒；使我們脫離冷淡的心,要在祂工作上大發熱心!無論在本地或往海外都要去傳揚福音,只要主給我們對象和異象,我們便甘心去事奉.
\newline
\newline
有一位十八歲的姊妹,雖然信主不久,但熱心服事主有拯救靈魂的異象,縱然自己沒有傳福音的經驗,但因愛人的靈魂,不願站在一旁.她對牧師說:「教會的工作這麼多,罪人的需要這樣大,巴不得我把自己獻上在服事主的祭壇奉獻.」她的牧師馬上回答說:「你願不願意將你的聲音交託主,因為我聽過你的聲音,你有唱詩的恩賜,可以在講道前,為主獨唱,被主大大的使用.」那位姊妹很高興立刻答應了.當她第一次在教會獨唱時,散會後有一個人對她說:「你今晚的歌聲救了我的性命,最近我很灰心,幾乎要跳海自殺；但聽你的歌聲後,我知道神在生命中有計劃,我不再為將來憂慮,現在我的心大有盼望了.」弟兄姊妹這是一件真實的事,是在英國發生的.所以我們站在一旁,不單對自己,也對別人沒有益處；反而中了魔鬼的詭計,破壞了教會的工作還使教會不能增長.你心裡有感動要服事主,無論是在本地或海外,只管靠著主去,祂必會賜給你力量.你蒙召是作兒童工作,青年工作,婦女工作嗎?請不要站在一旁.如果我們把工作,推到少數人身上；聖靈分給我們各人的恩賜,又怎樣被主使用呢?
\newline
\newline
(三)我們應該彼此保護
\newline
\newline
以東人最大的罪,相信是在十四節將猶大人交付仇敵,第二大的罪,就是當同胞猶大人受害時,因此而歡樂.我們信主的人,雖然不一定像以東人驕傲,強暴；但是我們彼此對待的態度,有時候會出錯誤.當我們要各自保護自己時,在我們當中,就會生磨擦,批評和嫉妒,背後論斷別人.即使大家聚會見面的時候,表面上和好,但背後就破壞別人的名譽,議論他們的生活方式.比如不會教養孩子啦,教會長執和牧師不會講道啦；這種的態度,是不能榮耀主的.先知以賽亞預言主耶穌說:「壓傷的蘆葦,祂不折斷；將殘的燈火,祂不吹滅!」(賽四十二3)主認識我們一切的缺點,軟弱和過錯；但祂仍然安慰,幫助我們,祂更體恤我們的軟弱.但願聖靈在我們心中運行,將耶和華的責備,應用在我們生活中；在需要的地方,感動我們悔改；使我們彼此保護,不設惡計陷害弟兄姊妹.將批評變成代禱,立志要為批評我們的人,用愛心來禱告,這樣彼此饒恕,彼此幫助和彼此保護；在這三方面都領受耶和華的責備,將(俄10-14)應用在我們生活中,效法主耶穌在彼此對待的事上.
\newline
\newline


\newpage

\section{第56屆~港九培靈會~林克己~第八講　耶和華的審判}
\label{sec:379}
\textbf{第八講　耶和華的審判}
\newline
\newline
連結: \href{https://www.hkbibleconference.org/session-message/view/379}{\texttt{https://www.hkbibleconference.org/session-message/view/379}}
\newline
\newline
\hyperref[sec:378]{< < < PREV SERMON < < <}
~
\hyperref[sec:toc]{[返目錄]}
~
\hyperref[sec:381]{> > > NEXT SERMON > > >}
\newline
\newline
俄巴底亞書 1:15-1:16
\newline
\begin{longtable}{cl}
\hline
\hline
章節 & 經文 (和合本修訂版)\\
\hline
1:15 & \begin{tabularx}{0.7\textwidth}{X} 耶和華的日子臨近萬國；你所做的，人也必向你照樣做，你的報應必歸到自己頭上。 \end{tabularx} \\ \\ \relax
1:16 & \begin{tabularx}{0.7\textwidth}{X} 你們在我聖山怎樣喝了苦杯，萬國必照樣不停地喝，且喝且吞，他們就必歸於無有。 \end{tabularx} \\ \\
[1ex]
\hline
\hline
\end{longtable}
「耶和華降罰的日子臨近萬國,你怎樣行,他也必照樣向你行；你的報應必歸到你頭上.你們猶大人在我聖山上怎樣喝了苦杯；萬國也必照樣常常的喝；且喝且咽,他們就歸無有.」(15-16節)
\newline
\newline
感謝主!我們的神是信實可靠的,祂說甚麼就作甚麼.祂也是聖潔公義,審判聖徒和罪人,按照這段經文的永不改變的原則,我們分三段來研究；第一審判的範圍,第二審判的基礎,第三審判的來臨.
\newline
\newline
(一)審判的範圍
\newline
\newline
「耶和華降罰的日子臨到萬國」(15節),「萬國也必照樣常常的喝」(16節),兩次出現「萬國」,可見耶和華的偉大.舊約時的外邦人,還有一些信仰不堅定的以色列人；都有一種錯誤的觀念,以為人所拜的神一定在世界上有特別的範圍和地方,這種錯誤觀念連中國人也包括在內.(王上二十23)「亞蘭王的臣僕對亞蘭王說:以色列人的神是山神,所以他們勝過我們,但在平原與他們打仗,我們必定得勝.」這顯然有地方神的觀念,不明白耶和華無所不在,亞蘭人即敘利亞人,以為以色列人的神,好像自己所拜的神一樣；不過在自己統管狹小的範圍,就是高山,山谷和平原不是祂的範圍.在平原與以色列人打仗,山神的耶和華便不能保護他們,亞蘭人便必然打勝仗.但是我們的神有辦法來教訓這些亞蘭人,改正他們的錯誤思想,叫他們不得不承認耶和華無比的偉大,耶和華超乎萬有的榮耀.(王上二十28)「有神人來見以色列王說,耶和華如此說,亞蘭人既說我耶和華是山神,不是平原的神；所以我必將這一大群人,都交在你手中,你們就知道我是耶和華.」耶和華施行審判,高山平原都是祂的範圍.
\newline
\newline
(創廿八10-17)耶和華教導以色列人的祖宗雅各,使他永遠沒有耶和華是地方神的觀念.雅各逃避他哥哥以掃,他逃走的時候,離開家鄉迦南地,逃到外地過夜.雅各在迦南地知道耶和華是他全家所拜的神,但是在迦南地以外呢?耶和華不是小小的木頭,一是小小的地方神；乃是全地的主,世界全個角落,都是祂的範圍.在萬國中有權柄,能力,引導祂兒女的路程,保佑他們到底.
\newline
\newline
以色列人最大的祖先亞伯拉罕,怎樣認識耶和華是萬有的主,和祂審判的範圍呢?(創十八25)「審判全地的主,豈不行公義麼?」全地萬國,關於耶和華審判的範圍,亞伯拉罕和俄巴底亞有同一信仰,告訴我們世界每一角落,歷史每一階段,都由耶和華作主,作王.按照祂美好的旨意,安排,管理一切,這是基督徒的歷史觀.我們讀聖經就深深相信,無論哪一個國家,時代,社會上所有不公道的事,人際間的壓迫和偏待的事,末後都要經過耶和華聖潔的考驗；受到祂公義的審判,得到我們應得的報應.將來在神的面前,無論是公開或隱藏的事,白晝黑夜都是一樣.這樣我們可以得到鼓勵,無論我們遇見的事是何等不公平；都是這位聖潔,公義的耶和華所知道的.早晚必伸出祂大能審判的手,除掉在祂面前不對的事；在全地萬國中,將公義永恆地顯明出來,使人無可推諉.
\newline
\newline
(二)審判的基礎
\newline
\newline
「你怎樣行,他也必照樣向你行.」(15節),神這個原則,是那些行公義的人必定會贊成的,不能不說應該是這樣.聽說有一個家庭,作丈夫的對妻子及兒子缺少關心和責任,因酗酒緣故,天天到酒吧浪費金錢.有一天晚上,他睡覺時發夢,夢見有三隻貓；第一隻很肥,第二隻很瘦,第三隻瞎眼.很奇怪他第二天向家人提及自己所做的夢,想不到那位小兒子說:「我真盼望為你解夢,你所夢見的三隻貓都有特別的意思.肥貓代表酒巴老板,你天天喝他的酒,他便賺很多的錢.瘦貓是媽媽,因你浪費了很多錢,媽媽便沒有好吃,好穿的.爸爸,瞎眼的貓是你,你看不見自己的行為對兒女有害處,怎樣叫別人吃虧.」我們只要看以東人的歷史,就可以看見同一件事,他們怎樣對待別人,自己也要怎樣被對待.主前五八六年,猶大的京城耶路撒冷,被巴比倫攻破時,以東人不但不幫助自己的同胞,還站在一旁,與敵人合作.過了三百年,以東人很可憐,他們被亞拉伯人制服,卻沒有人來幫助,只能逃往猶大.俄巴底亞書這個公平的原則,耶和華的審判臨到他們的歷史.你怎樣行,別人也對你怎樣行；國家或個人,舊約或新約,都有應用.
\newline
\newline
(啟二十11-12)「我又看見一個白色的大寶座,與坐在上面的,從他面前天地都逃避,再無可見之處了.我又看見死了的人,無論大小,都站在寶座前,案卷展開了,並且另有一卷展開,就是生命冊；死了的人都憑著這些案卷所記載的,照他們所行的受審判.」俄巴底亞書與啟示錄皆注重人所行的,是審判的基礎,看出他對人,對神所有的態度怎樣.福音派基督徒強調行為是得救的表現,而不是得救的條件.但我們好像不信主者,將來照自己的行為受耶和華公義的審判；在主耶穌寶座前,要為自己的工作交賬.(俄15節)是我們每一個人應該注意的事.非基督徒方面,他們所行的,耶穌就用他們的行為,來定罪受刑罰.基督徒方面,照自己所行的受審判；但不會被定罪而滅亡,靠著耶穌的寶血解決罪的問題.以後永不再提,主耶穌鑑察並審判我們的行為,就是注意我們兩件事:我們屬靈的生活聖潔嗎?我們屬靈的工作認真嗎?如果我們在這兩件重要的事上,得祂的喜悅,就在救恩之外加上賞賜.
\newline
\newline
(林前三11-15),11節論到得救和工作的根基,就是主耶穌基督,我們的得救和工作,除祂以外,我們別無根基.這個根基在我生命中立好了,(12-15節)我們在得救的根基上只有兩樣的工程和見證.第一有屬靈愛主,愛人的目的和動機,價值和性質,都好像金銀寶石一樣.經過火,一點也不成問題,雖然有試驗,但仍永遠存得住.第二屬肉體的,高舉自己過於高舉主；討人的喜歡過於討神的喜悅；價值和性質好像草木禾楷一樣.如果是屬靈愛主的,在主看來就好像金銀寶石一樣寶貴,必定得到賞賜.在得救根基上另外加上賞賜,在得救事上另加上榮耀和喜樂.但是第二種的工程,因為動機和目的,價值和性質,都有問題.(15節)指出這樣的基督徒,雖然是得救,但因行為不聖潔,工作不忠心,見證不熱心,賞賜方面就有虧損.雖然可以上天堂享受永生,但在得救的福氣上,沒有另外加上,沒有被主稱讚和喜悅,而得到賞賜.我們要在心裡向主禱告,求主保守我們生活聖潔,工作認真,見證熱心,至死忠心!這樣在救恩外,還可得著稱讚和賞賜.
\newline
\newline
我很不喜歡批評別人和別的宗教,但很可惜,許多羅馬天主教神父和教徒,他們誤解了林前三章的意思.他們以為這裡是論及得救的問題,這裡的火是指除了耶穌的寶血,尚未潔淨罪惡.他們承認三個地方,就是有天堂,地獄和煉獄.雖然聖經從沒有一次提及煉獄,但他們強調,如果你不夠好上天堂；不夠壞下地獄,你就去煉獄.大多數不屬靈的基督徒死後要在煉獄一段時期,煉獄的火,將耶穌的寶血尚未除掉罪惡煉淨.活著的基督徒可以為去世的基督徒,獻上禱告,捐錢給教會；這樣會縮短那些去世的基督徒在煉獄的時間.改革時期,有天主教教宗的代表,來到馬丁路德的牧區,那代表叫帖實勒,他注重捐錢和買贖罪票.帖實勒講道說:「我願意你們損錢給教會,支持教宗重修羅馬大教堂的計劃,可以購買教宗祝福過的贖罪票.如此不但自己的罪得赦免,還有當你們將錢投入奉獻箱中,錢所發的響聲時；你所記念已去世的親人,馬上從煉獄釋放出去,一直到天堂去.」難怪馬丁路德忍受不住,幾日後將自己署名與教庭挑戰的九十五條款,釘在自己教會報告板上,以後作一個熱心的改革者.
\newline
\newline
(約壹二28)「小子們哪!你們要住在主裡面,這樣祂若顯現,我們就可以坦然無懼；當他來的時候,在祂面前也不至於慚愧.」這節經文注重,基督徒要住在主裡面,天天與祂親近；藉著聖經祈禱,保持與祂親密交通,活出聖潔的生命來.想到主再來時,我們可以坦然無懼!當主顯現作審判官的時候,一點也不感到慚愧.無論我們是傳道人,宣教士或平信徒,主要的是我們的忠心.在為我們的主工作上,是否熱心?在生活中和工作上,都得到主的喜悅!便可以坦然無懼,毫無畏懼也經過祂的審判.
\newline
\newline
幾十年前,英國女皇維多利亞,自己一個人在郊外散步的時候,忽然大雨,她便走到附近一間屋,問借一把雨傘.有一個婦人開門,但不認得維多利亞女皇,便拿一把破壞陳舊,差不多全爛的雨傘借給女皇.第二天女皇派侍衛穿著特別的服裝,到那婦人的家去,把雨傘還給她,侍衛對她說:「維多利亞女皇多謝你,特別派我來把雨傘還給你.」那位婦人非常震驚,臉紅地羞愧回答說:「如果我知道,昨天到我家裡來的是女皇,我一定會把最好的雨傘借給她,可惜已經錯過了這個最好的機會.」弟兄姊妹,我們的神比維多利亞女皇更加偉大!我們現在有機會把生命奉獻給祂,將最好的奉獻在祂的祭壇上；或是全時間或部份時間,在教會中事奉,我們就要忠誠熱心工作.當審判的時候,祂就照我們所行一生的事奉,稱讚我們是個又忠心又善良的僕人,在救恩以外得到賞賜.
\newline
\newline
(三)審判的來源
\newline
\newline
(俄16節)這一個預言,顯示當猶大的耶路撒冷城被攻破的時候,他們逼於要喝苦杯,這件事耶和華神必記念.過了幾百年後,神的審判便臨到壓逼猶大人的仇敵.「萬國也必照樣常常的喝」,以東人,巴比倫人,非利士人,亞摩利人,這些東方人圍繞猶大人.今天他們已歸於無有,他們已經過去,名字不再被提起.幾百年前,德國菲勒德利第一世,他對宗教信仰沒有興趣.有一天他對皇宮侍候的牧師說:「我很不喜歡教會,因為教會常要我們信聖經,聽聖經；今日我要你講一個字,來證明聖經是可信的.」牧師馬上回答說:「好,我所揀選的是『猶太人』(JEWS)」牧師所說得很對,因為那些以東人,巴比倫人等,都歸於無有,但神對猶太人在聖經所說的一切話,句句也不落空,句句也可靠,一直到我們的時代.第二次世界大戰後,中東又再出現了以色列國,聖經預言猶太人的分散,也預言他們聚合.在屬靈和恩典方面,教會是第二個以色列.今日基督徒也是神的選民,我們雖然活在一個不明朗的時代,處於犯罪作惡的社會中,但俄巴底亞所事奉敬拜的耶和華,祂認識和記念我們.祂是大有憐憫的天父,我們信主後,與祂發生父子的關係.一個永不斷絕的屬靈關係,將來的事可以交託給祂讓祂安排負責,一點也不怕將來的審判.聖經說:「如今那些在基督耶穌裡的,就不再定罪了.」
\newline
\newline


\newpage

\section{第56屆~港九培靈會~林克己~第九講　耶和華的國度}
\label{sec:381}
\textbf{第九講　耶和華的國度}
\newline
\newline
連結: \href{https://www.hkbibleconference.org/session-message/view/381}{\texttt{https://www.hkbibleconference.org/session-message/view/381}}
\newline
\newline
\hyperref[sec:379]{< < < PREV SERMON < < <}
~
\hyperref[sec:toc]{[返目錄]}
~
\hyperref[sec:1387]{> > > NEXT SERMON > > >}
\newline
\newline
俄巴底亞書 1:17-1:21
\newline
\begin{longtable}{cl}
\hline
\hline
章節 & 經文 (和合本修訂版)\\
\hline
1:17 & \begin{tabularx}{0.7\textwidth}{X} 但在錫安山必有逃脫的人，那山必成為聖；雅各家必得原有的產業。 \end{tabularx} \\ \\ \relax
1:18 & \begin{tabularx}{0.7\textwidth}{X} 雅各家必成為大火，約瑟家成為火焰；以掃家必如碎秸，遭燃燒，被吞滅，以掃家必無倖存者。這是耶和華說的。 \end{tabularx} \\ \\ \relax
1:19 & \begin{tabularx}{0.7\textwidth}{X} 他們必得尼革夫和以掃山，得謝非拉，非利士人之地，他們必得以法蓮地和撒瑪利亞地，得便雅憫和基列； \end{tabularx} \\ \\ \relax
1:20 & \begin{tabularx}{0.7\textwidth}{X} 被擄的以色列大軍必得迦南人的地，直到撒勒法，在西法拉被擄的耶路撒冷人必得尼革夫的城鎮。 \end{tabularx} \\ \\ \relax
1:21 & \begin{tabularx}{0.7\textwidth}{X} 必有一些解救者上到錫安山，審判以掃山，國度就歸耶和華了。 \end{tabularx} \\ \\
[1ex]
\hline
\hline
\end{longtable}
「在錫安山必有逃脫的人,那山也必成聖.雅各家必得原有的產業,雅各家必成為大火；約瑟家必為火燄,以掃家必如碎楷.火必將他燒著吞滅,以掃家必無餘剩的,這是耶和華說的.南地的人,必得以掃山,高原的人,必得非利士地；也得以法蓮地,和撒瑪利亞地,便雅憫人必得基列.在迦南人中被擄的耶路撒冷,必得南地的城邑；必有拯救者上到錫安山,審判以掃山,國度就歸耶和華了.」
\newline
\newline
這段經文可分三段來研究,第一是預言的成就,第二是寶貴的恩典,第三是無比的榮耀.這樣看出天父奇妙的偉大.
\newline
\newline
(一)預言的成就
\newline
\newline
(17節)「在錫安山必有逃脫的人」,這裡的「必」字,是重將來必定會發生的事,「從17-21節」一共有十四次出現.這位先知講預言,不是靠自己發表自己的意見和盼望；他所得的全部是耶和華的默示,是天父真理的旨意和計劃.
\newline
\newline
(彼後一20-21)「第一要緊的,該知道經上所有的預言,沒有可隨私意解說的；因為預言從來沒有出於人意的,乃是人被聖靈感動說出神的話來.」基督徒應該知道舊約所有的預言,不可以輕視和隨意解釋,原因又重要又簡單,這些古代的人所謂的；不是人的意思,而是聖靈神的話.我們從這論點看(俄17-21),我們的神是無所不知的,祂向先知所說的每一句話,必定在將來歷史上成就.
\newline
\newline
我們從這段聖經看預言的可靠性和準確性,來堅固自己的信心.我們揀選一些預言,來證明耶和華無所不知,無所不在,「雅各家必得原有的產業」,這句話在選民的歷史上,最少有兩次的應驗.雖然北國以色列在主前七二一年被亞述擄掠到東方,以後便不再回來.俄巴底亞預言幾十年後,南國猶大就平安地從巴比倫歸回,重修聖城耶路撒冷,得到他們的產業,這段經文第二次應驗.主後七十年耶路撒冷被羅馬人毀滅,神的選民就被分散,迄今已一千多年,在中原沒有國家和產業.但到了第二次世界大戰後,一九四八年卻正式成立以色列國；應驗俄巴底亞和舊約許多先知的預言,使猶太人得回自己的產業.從歷史上來看,我們的神將自己的應許賜給祂的選民,祂必定會記念,人也無法推翻.從一九四八年至今天,亞拉伯人有好幾次一同來攻擊以色列；要除掉這新成立的國家,但卻沒有成功.從戰役上來說,這是一件神蹟；因為猶太人與亞拉伯人打仗,猶太人的武器和人數都比不上亞拉伯人.在屬靈方面,我們是神的選民,我們雖然活在一個正在改變,前途不明朗的世代中,好像主前六世紀俄巴底亞先知的時代,但俄巴底亞那位信實可靠的神,也是我們的神.祂怎樣保守猶太人,勝過脫離他們的敵人；明顯地與他們同在,同戰,祂也一樣會看顧我們.使我們勝過撒旦擺在我們前面的試探和苦難,一直到世界的末了,享受祂的同在和保守；願主將這真理,明確地刻在我們心版上.我們的環境常常改變,每天的遭遇不同,但耶和華永不改變,我們的境況無論怎樣,祂的應許永不落空.
\newline
\newline
(18節)論到以東的結論,加上「這是耶和華說的」,既是祂的命定,就不能不發生.這裡說神的選民要成為大火,來燒盡吞滅如同碎楷的以東人,直到無餘剩的地步.在歷史和聖經裡,都有辦法研究這段經文的應驗.(可三8)「還有許多人聽見祂所作的大事,就從猶太耶路撒冷,以土買,約旦河外,並推羅西頓的四方,來到祂那裡.」在這裡我們注意「以土買」這個地方,這是俄巴底亞時的以東,耶穌降生前,以東改名為以土買,表明以東國已不存在.他們從歷史書上塗抹了,今日的報紙和地圖,承認以色列的存在,但是沒有提以東人或以東國,所以俄18節的預言.只是三百年後,在兩約時代第二世紀的戰爭,有這樣的成就:使我們承認相信天父既是人權歷史的主,最大的權柄和主權,都永遠屬祂的寶座.預言的成就,不但堅固我們的心,也能安慰我們的心懷,使我們有平安的心來倚靠耶和華,作為我們的避難所.每逢我們念主禱文的時候,常提及耶和華的國度「願你的國度降臨,願你的旨意行在地上,如同行在天上......因為國度,權柄,榮耀,全是你的,直到永遠.」
\newline
\newline
(二)寶貴的恩典
\newline
\newline
我們留心查考(俄17-21)的時候,我們發現好像有兩種矛盾互相衝突的預言.將18節和21節比較,就發現預言以東傾倒消滅,但21節還提到有拯救者.如果以東全然被大火燒盡,怎會有拯救者在以掃山審判以東人,國度就歸耶和華呢?這兩件事擺在一起,是寶貴的恩典,看見耶和華超過過犯的恩典.在兩約之間,即瑪拉基書到馬太福音中間,四百年的時期；主前第二世紀,主為猶太人興起兩位領袖.第一位猶大馬克比,他是很有才幹的軍長；在主前一六八到一六四年,四次率領猶大人勝過他們的仇敵.第二位領袖是位大祭司,名叫約翰許爾堪；他在主前一三五至一零四年,作猶大人的大祭司.在這位領袖的時期,戰場上已死了很多以東人；剩下很少的以東人,明顯地以東國沒有前途和盼望,將要歸於無有.這兩位猶大的領袖就對以東人說:「我們願意饒恕你們!也給你們機會選擇,你們可以接受割禮；表示在肉體上,接納神為我們祖宗亞伯拉罕,這個表示屬於神和歸向神的記號,歸入神的選民中.以後同我們一同享受,歸向耶和華的一切屬靈好處；因我們原是同胞,猶大人和以東人都是以撒的子孫.如果你們心裡仍懷仇恨,錯過這個機會,你們死就好了.」很奇妙!一千多年的仇恨,結果被饒恕和恩典所吞滅.所以俄18節的「吞滅」有兩方面的意思,雖然國家被吞滅,歸於無有,地名改為以土買；但是在恩典方面,馬克比和許爾堪成為21節所預言的拯救者和審判官.為以東人定下了饒恕的挑戰,在耶和華的引導下,使他們得到寶貴的恩典.只要他們願意在神的面前悔改謙卑,就可以得到神和人的赦免；不再過屬血氣和屬世界的生活,得到主在靈性上的復興和滋潤,我們可以有三方面的領受.
\newline
\newline
1.在屬靈方面,基督徒原是以東人,只有寶貴的恩典才能改變我們,使我們歸向耶穌,成為神的選民.使徒保羅注重這件事,高舉我們在耶穌裡得到的寶貴恩典；(弗二11-18),猶大人與以東人因著恩典,打成一片,這也可以應用在我們的生活中.感謝主!我們原是外邦人,好像以東人一樣,與基督無關,也沒有指望；但是主寶貴的恩典,在我們身上發生效力.雖然好像以東人,還能稱為猶大人；與神和好,在基督裡得稱為義,在教會裡成為第二個以色列享受福份.
\newline
\newline
2.在這段經文中,基督怎樣傳講和平的福音,我們作為歸向祂的人,也有責任.按照祂為我們定下不同的旨意,在教會中全時間或部份時間服事主.但願在我們的生活中,有願為祂所犧牲的,表明祂不單是我們的救主,也是我們整個生命的主.
\newline
\newline
3.我們想到兩位領袖,為我們留下饒恕敵人,寬容別人的榜樣；叫我們想到主怎樣用寶貴的恩典,白白赦免我們的罪惡.我們應該從心裡饒恕那些得罪我們的人.在14節時我們看過,現在第二次叫我們去面對,耶穌怎樣饒恕我們,我們也要饒恕別人.
\newline
\newline
有一位婦人往見醫生,診治她臉上的皮膚病；她用了醫生的藥抹臉上,但仍得不到痊癒.有一天她與醫生談話的時候,談及應該怎樣才得醫治?這婦人提到有一個她所憎惡的鄰居,她每每提到這鄰居,臉皮便發紅了；疙便顯得更多更大.醫生便對她說:「我相信你這皮膚病,不是身體的問題,而是心理的問題.你去與鄰居和好,解決你們的糾紛,看以後你的臉孔怎樣.」真奇妙!當兩人和好後,婦人第二次去見醫生,臉皮完全好了.從這個真實故事,我們可以學習一個功課,讓我們被主寶貴的恩典保持我們心懷意念,想到主赦免我們的恩典；主的寶血流出來,就有力量和感動去饒恕別人.這樣在我們的見證上,使醜陋變成佳美,好像那婦人的皮膚病得了醫治一樣.
\newline
\newline
(三)無比的榮耀
\newline
\newline
俄21節「國度就歸耶和華了」,這不單是俄巴底亞先知的結論,也是他教訓的高峰.以東人雖然不認識祂,但在將來的歷史上,特別是世界的末了,耶和華無比的榮耀,必充滿天地.
\newline
\newline
這不但是俄巴底亞的思想和觀念,我常常將這一節與(亞十四7-9)一齊看,「那日,必是耶和華所知道的,不是白晝,也不是黑夜,到了晚上才有光明.那日必有活水從耶路撒冷出來,一半往東海流,一半往西海流,冬夏都是如此.耶和華必作全地的王,那日耶和華必為獨一無二的,祂的名也是獨一無二的.」俄巴底亞與撒迦利亞所預言,描寫的,是教會從主後七十年耶路撒冷被羅馬人毀壞,一直到主耶穌第二次到世上來.撒迦利亞用「那日」兩個字,教會時代,像很長的日子,但感謝主,這日子是耶和華所知道；我們八十年代的基督徒,所遭遇的事,也包括在內.有一首詩歌「神未曾應許,天色常藍」,先知在撒迦利亞預言教會的日子,不完全是光明或黑暗；雖然教會有時候遇見逼迫和反對,天空好像有黑雲滿佈.但有時候,太陽也發光.教會的工作能順利的進行,傳揚福音的時候,歸向主的人也多.當主得勝再來時,我們八十年代的基督徒,可以用信心抬起頭來；無論環境是何等困難危險,但神為自己的兒女,所預定的計劃不會落空.擺在我們前頭的日子,一定是光明,教會無論有甚麼遭遇,從主後七十年一直到主再來,仍然有一個光明的前途.當我們所愛的主耶穌第二次再來的時候,那天有光明的晚上,祂要作萬王之王；在這個世界上,將耶和華無比的榮耀,彰顯出來.(亞十四9)「耶和華必作全地的王,那日耶和華必為獨一無二的,祂的名也是獨一無二的.」換句話說,魔鬼,邪靈,偶像,都不能敵擋耶和華無比的榮耀；到了那光明的晚上,萬膝要跪拜,萬口要承認；我們的主耶穌是獨一無二的,祂的名字也是獨一無二的.
\newline
\newline
我們看(亞十四8),這是主為今天教會所定的計劃,為你我生命中所定的旨意,我們屬靈生活應有的目標.「那日必有活水從耶路撒冷出來,一半往東海流,一半往西海流,冬夏都是如此」.主耶穌在耶路撒冷講道的時候,他在約翰福音將撒迦利亞書的預言,應用在教會的時代.主說:「信我的人,就如經上所說,從他腹中,要流出活水的江河來.」歷史學家湯恩比曾講過很有意思的話:「如果十九世紀是屬於英國人的時代,第二十世紀是屬於美國人的時代,相信廿一世紀一定是屬於中國人的時代.」看近兩百年歷史,湯恩比的話不是沒有理由的,看十九世紀,英國在維多利亞女皇時代,因工業革命,英國的影響和勢力很大.二十世紀,因為美國科學發達,現在有原子和太空的時代,但是廿一世紀會怎樣?湯恩比當然不是神的僕人,好像撒迦利亞和俄巴底亞一樣,但是他用自己的學問,發表廿一世紀的意見,這樣看歷史的潮流準確的話；如果主在二十世紀不回來的話,相信現在是華人教會作準備的一個時期.在廿一世紀,全世界都仰望中國的時候,見證大被主用,福音的活水按照(亞十四8)的預言,和(約七章)的說話,要從華人教會裡往東,往西,流出來,冬夏永流不息.藉著中國基督徒奉獻主的生活中,給世界上的人帶來救恩和生命的喜樂,這裡有團體和個人方面的挑戰.想到耶和華的國度,預言的成就,寶貴的恩典,無比的榮耀.這三方面,只要我們願意,將自己毫無保留和條件委身給耶穌,成為活水的江河.在我們被聖靈充滿的生命中流出來,東西南北,春夏秋冬,永流不息.繼續不停地為主作見證,就在(俄21)節中有份,國度就歸耶和華了.願主幫助我們,好像祂忠心的僕人俄巴底亞一樣,在我們這個在屬靈需要逼切的時代中,作敬拜,事奉耶和華的人,在我們各方面的生活中,願榮耀歸給耶和華的聖名.
\newline
\newline


\newpage

\section{第56屆~港九培靈會~楊濬哲~序}
\label{sec:1387}
\textbf{序}
\newline
\newline
連結: \href{https://www.hkbibleconference.org/session-message/view/1387}{\texttt{https://www.hkbibleconference.org/session-message/view/1387}}
\newline
\newline
\hyperref[sec:381]{< < < PREV SERMON < < <}
~
\hyperref[sec:toc]{[返目錄]}
~
\hyperref[sec:402]{> > > NEXT SERMON > > >}
\newline
\newline
八四年夏香港情況非常的動盪,因為前途還未明朗,以至人心惶惶,有不少人移民離港.是否別處地方比香港更安全呢?沒有人可以下結論給答案.
\newline
\newline
神的話語是我們腳前的燈,路上的光,我們需要神的話語來引導我們呀!同時神的話語也是我們的力量和隨時的幫助,當我們軟弱的時候,它能扶持我們,堅固我們,剛強我們.當我們心灰意冷,失意絕望的時候,它是我們唯一的安慰!苦難來無非要鍛煉我們的信心,使我們更堅心倚靠神!有生命的信徒,經過試煉更能站穩,堅守主道.
\newline
\newline
港九培靈研經會承擔著一個很沉重的任務,就是要藉著神的話語,來栽培,造就信徒.感謝神給我們看見勞苦的果效.根據八四年晚堂奮興會,決志者交回的資料顯示,收到決志表共四九一份,決志者的堂會,來自一九三間,填寫悔改復興,為主而活的佔四三九,奉獻全時間事奉者佔一二九人,決志信主者三三人,其中資料不詳者五二人.這使我們萬分鼓舞!
\newline
\newline
三位講員所傳達的信息,正合我們的需要.林克己牧師用了九堂時間,解釋俄巴底亞書,證明神是掌管明天的神,我們是神的兒女,我們更需要堅心倚靠祂.麥希真牧師以立在磐石上為主題,卻教導我們在日常生活中,怎樣去過一個健康基督徒的生活.高文牧師再度由美來港領會,他從啟示錄發掘了天上新歌的寶訓,好叫在苦難中的信徒得到無限的安慰!神是至終的勝利者,魔鬼必失敗被扔下無底坑去.那為主殉道者的靈魂,羔羊必為他們伸冤.凡被請赴羔羊婚筵者有福了!
\newline
\newline
期望大家能寶貴本書的信息,從而得到復興.
\newline
\newline
委辦會的元老卓恩高先生,已於本年六月下旬安息主懷,享年99歲.他在培靈會的勞苦與貢獻,願主記念!我們甘願跟隨他的腳蹤,忠心事主.
\newline
\newline
張有光牧師在美做了換心瓣手術,請記念為他健康祈禱!
\newline
\newline
本講道集蒙黃民弟兄,梁壽華弟兄,雷麗端姊妹三位記錄,林立基弟兄設計封面,朱建磯牧師編輯,林小娟姊妹,陳綺華姊妹協助校對,乃得完成,統此致謝!
\newline
\newline
大會將屆六十週年紀念,委辦刻正籌備慶典,敬請主內道長,不忘代禱!
\newline
\newline
願榮耀,頌讚,權能,歸與三位一體的神!阿們.
\newline
\newline


\newpage

\section{第56屆~港九培靈會~高文~第一講　神的尊榮}
\label{sec:402}
\textbf{第一講　神的尊榮}
\newline
\newline
連結: \href{https://www.hkbibleconference.org/session-message/view/402}{\texttt{https://www.hkbibleconference.org/session-message/view/402}}
\newline
\newline
\hyperref[sec:1387]{< < < PREV SERMON < < <}
~
\hyperref[sec:toc]{[返目錄]}
~
\hyperref[sec:404]{> > > NEXT SERMON > > >}
\newline
\newline
啟示錄 4:1-4:8
\newline
\begin{longtable}{cl}
\hline
\hline
章節 & 經文 (和合本修訂版)\\
\hline
4:1 & \begin{tabularx}{0.7\textwidth}{X} 這些事以後，我觀看，看見天上有一道門開著。我頭一次聽見的那好像吹號的聲音對我說：「你上這裡來，我要把此後必須發生的事指示你。」 \end{tabularx} \\ \\ \relax
4:2 & \begin{tabularx}{0.7\textwidth}{X} 我立刻被聖靈感動，見有一個寶座安置在天上，有一位坐在寶座上。 \end{tabularx} \\ \\ \relax
4:3 & \begin{tabularx}{0.7\textwidth}{X} 那坐著的，看來好像碧玉和紅寶石；又有彩虹圍著寶座，光彩好像綠寶石。 \end{tabularx} \\ \\ \relax
4:4 & \begin{tabularx}{0.7\textwidth}{X} 寶座的周圍又有二十四個座位，上面坐著二十四位長老，身穿白衣，頭上戴著金冠冕。 \end{tabularx} \\ \\ \relax
4:5 & \begin{tabularx}{0.7\textwidth}{X} 有閃電、聲音、雷轟從寶座中發出。在寶座前點著七支火炬，就是神的七靈。 \end{tabularx} \\ \\ \relax
4:6 & \begin{tabularx}{0.7\textwidth}{X} 寶座前有一個如同水晶的玻璃海。寶座的周圍，四邊有四個活物，遍體前後都長滿了眼睛。 \end{tabularx} \\ \\ \relax
4:7 & \begin{tabularx}{0.7\textwidth}{X} 第一個活物像獅子，第二個像牛犢，第三個的臉像人臉，第四個像飛鷹。 \end{tabularx} \\ \\ \relax
4:8 & \begin{tabularx}{0.7\textwidth}{X} 四個活物各有六個翅膀，遍體內外都長滿了眼睛。他們晝夜不住地說：「聖哉！聖哉！聖哉！主—全能的神；昔在、今在、以後永在！」 \end{tabularx} \\ \\
[1ex]
\hline
\hline
\end{longtable}
我在美國聽說香港居民,對前途多存有恐懼感,很多人忐忑不安；因此,我希望今年培靈會的信息,能給弟兄姊妹得到鼓勵,從神的話語得到幫助.我自己也不知道前面將會發生甚麼事情；但知道將要來的那一位正是耶穌基督,祂是掌管天地的主宰.到了有一日,普天下的人,萬膝都要向祂跪拜；故此,我們要存著信心面對前面的日子.今日教會滿受撒旦攻擊,我們所爭戰的不是血肉的對手,乃是在高處掌管天空屬靈氣的惡魔,這種邪惡勢力要盡它所能的,摧殘一切高舉耶穌基督的工作；若繼續演變下去,則未來之日,恐怕比現在更不安定,黑暗的勢力攻擊教會,將越加猛烈.從人的角度來看,為公義而爭戰,有一日公義將遭遇悲慘的命運；但是屬神的子民不是從世界的角度來看世事的發展.我們不是按著人的資源能力,來看前面的歲月,我們的眼不是望著怒海奔騰的世界；我們乃不斷仰望作工的神；無論在世上遭遇甚麼,我們必須知道神的國度必定要來臨,耶穌基督將在一切之上掌管萬有.如果想知道這個國度如何,最好的方法就是閱讀啟示錄了.
\newline
\newline
耶穌基督的啟示在啟示錄第一章一節就給我們看見全卷書的主題,要將統管萬有的基督的啟示寫出來.全書論及有關耶穌基督的啟示,或祂所給的啟示,「耶穌基督」正是啟示錄的中心；描述耶穌基督怎樣戰勝一切敵對的勢力,而建立祂的國度.
\newline
\newline
啟示錄寫出世上邪惡的勢力,如何盡其所能作摧殘的工作；而本書的中心是天上坐寶座的那位才是最終的得勝者和統治者.基督正在體驗發生在我們面前的事,一直帶領歷史的進展.在這卷書中你會常見「寶座」「國度」「權柄」等字眼,甚至有的在這卷書用了一百廿次.
\newline
\newline
「念這書上豫言的,和那些聽見又遵守其中所記載的,都是有福的.」(一3)這是神的有福應許,到了本書最末章第7節,同樣的應許再次出現.
\newline
\newline
本書滿有神豐富的寶藏,是我微小的頭腦無法解釋透徹的；有一天,這世界一切都要過去,但我相信,神所默示的每件事和話語；照著這本小小的書,神的寶貝教訓要永遠常存.我們應存尊敬的心求神聖靈的光照,神要把重要的真理教導我們.
\newline
\newline
主統管萬有:道成肉身的主在我們中間,我們就能接受祂,認識祂.約翰看見了主耶穌榮耀的異象,立即仆倒在主腳前；主叫他不用懼怕,說:「我是首先的,我是末後的......,我曾死過,現在又活了,直活到永永遠遠.......」(一17-18)這位就是駕雲降臨的主.聖經說,每個眼睛都要看見全能永遠的神.(一5-6)說祂是偉大的主,是愛我們的主,祂藉著十字架上所流的寶血,買贖我們；又使我們成為祭司,與父有交通.在這樣的異象中,約翰只能說:「願一切的尊貴,榮耀,都歸給祂,直到永永遠遠.」
\newline
\newline
讚美的詩歌:(一5-6)的讚美詩歌,在啟示錄中重新的出現；刻劃出啟示錄的主題,在父神寶座前的詩歌,又重複這個主題.聖經說,詩歌是他們向神所發出讚美的言詞.事實上,這段經文不是唱的,乃是向神說的；唱和講兩樣俱備.在啟示錄中有十四段這樣的詩歌,從第四章開始一直到十九章.這詩歌要給我們地上的歷史,屬天的描述；在地上人的制度中,天上寶座前,才有這種的詩歌.當這詩歌稱頌神時,使我們記起來,要敬拜稱頌我們的神.在整個詩歌中主要的是敬拜,帶著活祭在神面前敬拜,神是這詩歌讚美的對象.現場群眾的歌聲稱頌祂的屬靈,祂的大作為.當我們見到祂的榮光,我們受造之物才有最大的喜樂.讚美的歌聲就是天上的言語,我們現在要學習這種言語；有一日我們到了神的寶座前,才懂得這種言語.歷史學家說,這種詩歌是從未在初期教會,或是在任何崇拜時唱過的.詩歌對於民風,民族十分重要,所歌唱的真是能夠反映終日所思想的.在寶座周圍所唱的詩歌正有此作用,表明見到榮耀之主的人,心中的感受和讚美.我們若學習和他們一齊唱詩的話,我們就會感到新耶路撒冷的價值觀念了.
\newline
\newline
預言的文學家　全卷啟示錄用非常濃厚的景象和比喻解釋屬靈的事；不是用有系統的邏輯敘述,而是用幻想形式傳遞信息.此方法更易使人的感觀吸收真理,如此寫作方法在古舊時代極為通行.當人受逼迫時,傳達信息文章的寫作最適合於實際；猶太的基督信徒用此預言文學作當時被逼迫傳達信息的使用,他們早有舊約聖經作背景來明白這種預言文學.在第一世紀所有寫作之中,以啟示錄充滿著舊約聖經的影像,最少有超過60\%以上的篇幅引述舊約聖經.所以若不明白舊約聖經的背景,就無法明白啟示錄.以經解經,啟示錄達解經方法的高峰,把聖經高度的意思,搜集帶到神的寶座上.
\newline
\newline
第一章論統管萬有的基督.接著兩章論及主向七間教會發出書信.他們當時所面對的社會與今日我們一樣.到了第四章,約翰看見天上有門開了,(四2)至此,啟示錄轉向以後世界事情演變的問題上,開始了從天上望出去的角度.如今我們若從天上的角度看四圍的演變,是非常合宜的.第3節描述坐在寶座上者之公正,以及光芒萬丈,將神的尊貴,榮耀,刻劃出來；另一方面說明神的公義和救贖的計劃.　舊約時代大祭司的護心鏡乃用碧玉和紅寶石鑲在最低層處.聖經用月球代表以色列十二個支派.又有綠寶石好像天虹圍著寶座.先知以西結見神榮耀異象時,有綠色的天虹.挪亞時代洪水之後虹就出現,這表示是神立約的記號.在審判之中,我們從神再次領受活潑的盼望；天虹表示神恩典的確據,神的信實何等廣大,必守祂的應許,絕不失信.
\newline
\newline
神的能力「有閃電,聲音,雷轟,從寶座中發出.」(5)表示全能者的能力.神頒佈十誡時,在西乃山也有此情景出現.舊約以來,每有雷轟則代表全能者神的特性.「又有七盞火燈在寶座前點著,這七燈就是神的七靈.」(5)啟示錄由始至末顯示「七」的異象的出現,神的七間教會,七個金燈台,七天使有七靈,七印,七號筒,七星在基督右手中,羔羊有七角,有七眼,有七個羔羊,七個雷轟,龍的頭有七個角,七個冠冕,七座山,七個君王.「七」字在啟示錄由始至末都有.神的創造到了第七天便安息了.有七年的安息年,每第七個安息年有個禧年.逾越節距離降臨節有七個體拜,逾越節要獻上七隻羊羔.「七」象徵完整,「七靈」神的聖靈是完整的聖靈,聖靈發揮效率的完整.七盞燈在寶座前出現,「火」代表聖靈,水,潔淨與光照.聖經說,我們的神是烈火.在寶座前的光有敬畏的景象,襯托著,光亮如水晶河；玻璃海如同水晶將至高的神與其他一切被造之物隔開.
\newline
\newline
寶座中和寶座周圍有四活物,前後遍體都滿了眼睛,(6-8)望向整個前景,任何事物在七眼中無所遁形.四活物各有六個翅膀,都迅速執行神的計劃.每個活物形狀不同,第一個像獅子,第二個像牛犢,第三個臉面像人,第四個像飛鷹.(7)當以色列民在曠野時,會幕前面有四個活物；現在提醒他們想起曠野的會幕,他們聯想到與神子民的關係.四活物是最野生,最受苦,最勞力的,是最早受造之物.四活物如同指南針的四方向；說明神在整個宇宙中的範圍.聖經他處提及四活物之同時也提及撒拉法和基路伯.當以賽亞和以西結提及神的寶座,也提及寶座周圍有撒拉法,有基路伯.他們活在最接近神的領域,服侍神的靈正表明公義的服役.撒拉法在伊甸園外守衛,不准外來力量侵入.以西結提及基路伯正在向反叛的以色列人施行刑罰.先知以賽亞說:「讓撒拉法拿炭來放在嘴唇以便潔淨.」四活物在神面前敬拜,顯出對神的純潔和愛慕.到了第8節,四活物晝夜不住地稱頌神；在稱頌中他們見到寶座上全能者的性格.在至高寶座者四圍的生活形式就是這樣.他們唱:「聖哉,聖哉,聖哉,主神；是昔在,今在,以後永在的全能者.」三重的聖哉反映出三位一體真神同榮同尊.當以賽亞看見異象之時,他也呼叫:「聖哉!聖哉!聖哉!萬軍之耶和華.」
\newline
\newline
聖潔的神:聖潔的字根的意思是隔開的,分別出來的；神是聖哉,是自我滿足的,不需別的活物幫祂.祂是從任何勢力任何法律之外分別出來而永存的,從不受任何污染,聖經說我們的神有最純潔的眼目,我們就應當感到祂的純潔.我們的神是十全十美的完整.　有的人認為,何等樣的人就是拜何等樣的神；好像羅馬人拜偶像朱彼得(JUPITER)是野蠻假神,當然拜者也很野蠻.以弗所人拜淫亂女神,拜者也是如此,他們過不道德的生活.因你的生活不能高過你的神.
\newline
\newline
手潔心清的人才能來到聖潔的神面前:主來到世上,曾為我們釘死在十字架上,開了條生路,祂使滅亡的命運扭轉過來,為我們贖罪.親愛的兄姊!我們若有罪就必須及時承認；到神面前誠心誠意悔改,對罪惡絕不妥協；因為落在永生的神手中是很可怕的!我們若向神承認祂的聖潔,祂將帶我們到祂掌管的城園.全能的神成就祂的旨意,神展奧秘的膀臂,憑著大能創造天地萬物；在神沒有難事,整個宇宙的資源都在祂支配之下.
\newline
\newline
神是全能永存的神:聖經說,萬國在神手中如同一滴水,變得一無所有；祂看萬國如同天秤上的一點塵埃.世界的審判失去了生存的價值,只要神一句話,就能把世界上所有生命除去.神是十全十美的,祂有祂自己存在的原因,有祂自己的理由,有祂自己的活動；無論如何,神永遠存在.在世界上被壓迫和寄居的人們,聽到這個信息,何等大的安慰!為義受逼迫者的心得到鼓勵.神確是長久忍耐,當祂有所行動時,祂不受任何限制；因此,我們當禱告說:「願神的旨意行在地上,如同行在天上.」
\newline
\newline
法國皇帝路易十四舉行喪禮時,講員望著棺木說:「這樣華麗屬於君王的棺木裡,卻隱藏著己冷卻的屍體.」牧師向顯貴們致此哀詞,他說:「我的朋友啊!唯有神才是真正偉大!」這點真理是我們每位當明白的；世上還有人不明白「人的渺小」「神的偉大」.
\newline
\newline
天上的群眾不住地說,聖潔的神也是永遠的神；萬膝當向祂跪拜,萬口要承認,並將榮耀歸與祂.神是昔在,今在,永在的全能者.這首詩歌反映出神昔日向摩西顯現時說:「我是自有永有的.」神超越時間之上,在神無所謂繼續,開始,結局；沒有過去,沒有將來；神的名字「我是」.一切都是「我是」,時間是神創造的,人的工作如同一件衣服,終歸有變,或摺疊或拋棄；但神永遠一樣,在祂永無改變也無轉動的影兒；神是永恆的,從永遠直到永遠.
\newline
\newline
曾有一少年人在密西西比河邊,向經過的船不斷揮手,喊說:「我在這裡!」見者都笑他愚蠢；但船竟駛向河邊,船長從甲板下來,手挽著他上船,氣笛響,船開動了.岸邊的人非常驚奇；孩子說:「我肯定船必停下來的,因船長是我的爸爸.」
\newline
\newline
今天晚上,我們中間有誰心靈疲乏憂愁嗎?統管萬有的主宰,坐在天上聖潔的統治者全能的神,祂是永恆,榮耀的王,是我們的父親.祂十分愛我們,甚至將祂的獨生子賜給我們,祂邀請我們藉著耶穌基督到祂面前來,祂應許說:「凡到我這裡來的,我總不丟棄他.」
\newline
\newline


\newpage

\section{第56屆~港九培靈會~高文~第二講　金的冠冕}
\label{sec:404}
\textbf{第二講　金的冠冕}
\newline
\newline
連結: \href{https://www.hkbibleconference.org/session-message/view/404}{\texttt{https://www.hkbibleconference.org/session-message/view/404}}
\newline
\newline
\hyperref[sec:402]{< < < PREV SERMON < < <}
~
\hyperref[sec:toc]{[返目錄]}
~
\hyperref[sec:406]{> > > NEXT SERMON > > >}
\newline
\newline
啟示錄 4:9-4:10
\newline
\begin{longtable}{cl}
\hline
\hline
章節 & 經文 (和合本修訂版)\\
\hline
4:9 & \begin{tabularx}{0.7\textwidth}{X} 每逢四活物將榮耀、尊貴、感謝歸給那坐在寶座上、活到永永遠遠者的時候， \end{tabularx} \\ \\ \relax
4:10 & \begin{tabularx}{0.7\textwidth}{X} 二十四位長老就俯伏敬拜坐在寶座上活到永永遠遠的那一位，又把他們的冠冕放在寶座前，說： \end{tabularx} \\ \\
[1ex]
\hline
\hline
\end{longtable}
與主同在就是天堂　有一次我問一個人說:「你死了到天堂去嗎?」立即回答:「怎會不去呢,現在我就是住在天堂了.」基督徒現住天堂,這是千真萬確的；聖經說:聖徒與主同住在天上了.憑著信心,我們已進入主的範圍內；聖經清楚說我們是天上的國民,並且等待著主耶穌我們的救主降臨.
\newline
\newline
當思念天上的事　有一個音樂家,當他第一次演奏之後；朋友向他恭賀說:「你的演奏使我很得鼓勵!但我覺得奇怪!你的演奏好像瀰漫著上層.」他說:「因為我的老師坐在上面,我看見老師的微笑,知道他很開心.」在更深的意義上,我們更當仰望我們天上的主；我們看見主的笑容.我們在基督裡得著福份,故此,我們該明白為何聖經教導我們,應當思念上面的事.祂給我們的呼召,有何等的指望!使我們與眾聖徒一同得基業,一個承受的永生是何等的豐盛!耶穌基督的啟示寫成啟示錄,讀者正需要仰望主.
\newline
\newline
使徒約翰就用寫信的方式,寫了七封書信,給亞細亞的七個教會.第二至三章,把這七個教會的情況描繪出來；七教會當時屬靈的景況,與今日的許多教會有相似的地方.
\newline
\newline
當聽聖靈的話　約翰寫給七個教會的信,最後都勸勉說:「聖靈向眾教會所說的話,凡有耳的,就應當聽.」這是素常的,現在的呼籲；當時七教會所當聽聖靈的話,也是今時代基督徒所當聽的.
\newline
\newline
基督徒遭受窘迫　第一批讀啟示錄的基督徒,當時正面臨世上空前的困難,受羅馬帝國可怕的統治；因為基督徒不肯屈服於拜偶像的文化下,致遭受大逼迫；為了信仰,有時有許多人被殺死.有的受了不可言喻的迫害,有的被放逐於運動場上,任野獸撕裂示眾；有的用獸皮包裹最後被野獸咬死；有的身上被淋油掛上如釘十架,然後在皇帝的花園放火焚燒；在諸般的逼迫中,約翰和彼得也被窘迫至死.啟示錄約主後85-97年間面世；那時執政的王名稻米斯,此王帶來恐怖的統治.至少有四萬多的基督徒被逼迫致死；當時提摩太被以弗所的監督和一班忿怒的群眾殺死於亂石之下；約翰被放逐到拔摩海島上,與家人分離,長期在那裡作苦工.這位八十多高齡的約翰,就在形影孤單受苦的情況下寫了啟示錄.在患難之中,神仍然向他發聲,這是何等的安慰!神仍然掌管歷史.
\newline
\newline
敵對的勢力在教會中興波作浪,比在教會外的逼迫更甚；基督徒被強大壓力迫使他們放棄為基督作見證；他們向周圍一切感觀享受的文化妥協；有不少基督徒在此引誘之下受了魔鬼的捆鎖,面臨此種失敗的基督徒,啟示錄書信向他們發聲,提醒他們已失去起初的愛心,呼籲他們要悔改警惕.
\newline
\newline
啟示錄書信教我們對魔鬼狡猾的詭計當提防,叫我們確知耶穌基督才是生命的主宰.
\newline
\newline
參與天上之樂歌寶座四圍所唱的讚美詩歌一再重複耶穌是主,把天上的氣氛帶入我們的生命；盼望也帶入你安靜禱告之中,與天上詩班一同歌唱生命的詩歌,叫我們的心靈得鼓勵.但願我們把將逝的俗世纏累置之度外,不再聽那些終歸消失的世界聲音；全心全意注目於永恆的事物.「看哪!榮耀的主!」雖然我們目前所見漆黑一片；但烏雲不久即將消散,我們要見主榮面；那時我們要在榮耀裡,在主前歌頌祂.撒拉弗在神面前所唱的讚美詩,將來我們要在神的聖城歌頌.
\newline
\newline
敬拜和頌讚(四9-10)四活物將榮耀,尊貴,感謝,歸給那坐在寶座上,活到永永遠遠者.」(9節)廿四位長老參加這次的敬拜.第10節向我們介紹寶座周圍坐著廿四位身穿白衣的長老,他們頭上戴金冠冕,他們代表眾聖徒在神的寶座前.以賽亞書提及在神的寶座,有廿四位長老與神國統治,可能就是指這班人.廿四位長老與天上的群眾,與地上的祭司,晝夜不停在聖殿中獻祭.歷代志所載有廿四位利未人在地上擊鼓,彈琴,唱歌敬拜神.
\newline
\newline
「廿四」這數字很重要,代表舊約十二個以色列支派,和新約耶穌十二個門徒；把舊約和新約合而為一.廿四位長老身穿白衣,頭戴金冠冕,是屬於帝王人物,是得勝者；配得在神前與神的榮耀一同統治.耶穌曾說要與門徒一同坐在神的寶座上.這班長老俯伏在坐寶座者的面前,他們敬拜那活到永永遠遠的.「俯伏」「敬拜」加深了他們對神的敬畏.
\newline
\newline
冠冕　第四章9-10這段聖經還有另一細節,廿四位長老十分體會神的崇高性,於是他們把頭上的金冠冕放在寶座前.神曾經把冠冕,賞賜給服侍祂的人；這冠冕在文字上非指帝王式的冠冕,這裡的冠冕乃指賞賜；人的成就得到獎賞的冠冕.在較前啟示錄中提及,為神忠心至死者可得生命的冠冕.雅各曾用冠冕作比,說明基督徒受試煉之後可得生命的冠冕.彼得書信也說,基督徒受試可得榮耀的冠冕.保羅也用冠冕的詞句向忠心於主的人發出.這段聖經所說的是金的冠冕.金是貴重的金屬.當以金論神,論天上；天是神的城,天堂的街道是黃金的.這裡所說的金冠冕,是聖徒從神所領受屬天的生命；一切從罪惡污染的,今已洗淨了；在神的聖潔亮光中,火焰中,已經煉淨了.惟聖潔之物才能存留到那日,惟有純淨的真金才能留存冠冕上.他們想到自己原是不配戴金的冠冕,完全是神的恩賜；金冠冕惟神是配!他們把金冠冕放在主的寶座前歸給祂；然後由衷地向神發出讚美.
\newline
\newline
唯主是配「我們的主,我們的神,你是配得榮耀尊貴權柄的；因為你創造了萬物,並且萬物是因你的旨意被創造而有的.」(11節)這首詩歌再次強調主神是絕對配得的.同時也否定了羅馬帝王的自高自大,自封為王的傲慢.
\newline
\newline
原來羅馬帝王稻米斯在旅行時,受人跪拜稱頌說:「你是配得的!」敬拜皇帝的異端邪說宗教,向稻米斯敬拜說:「你是我們的神.」
\newline
\newline
榮耀尊貴權柄,唯天上坐寶座者才配得；因此,教會應當宣佈,與周圍一切異端邪說拜偶像的人完全無關；唯創造天地萬物之主,才是配得的.按照祂的旨意萬物因此被立,世界萬物無論天上的星辰創造之先,我們的神已經存在；除非在神掌管全地,在祂的旨意中；否則無一可以被造.　可能有人要問；為何世界還有很多苦難,很多邪惡呢?如果這位創造者神創造一切；受造的人類對四周的苦難,邪惡負有何責任呢?因為宇宙中有所謂因連論,有神,也有邪惡；邪惡的抉擇在神的身上.這首讚美詩歌深入淺出地回答這個問題.事實上,在神創造世界之先一無所有；惟有神.一切東西能夠存在乃因神的創造,並容許其存在,包括魔鬼與其同黨.神的全能,排除受造之物中的任何邪惡污穢.
\newline
\newline
神創造萬物的本意神創造原是聖潔的,但受造之物在自己的意旨裡卻反叛神,侵害了其他受造之物.　再從神的角度來看,這首詩歌發出的信息,神的工作已經開始,神的創造正反映了神的豐盛,神全美的完整.「諸天述說神的榮耀,穹蒼傳揚祂的手段.」(詩十九1)聖經說:「神看祂創造之物都是好的.」(傳三11)創造之物不是神的一部份,神本身是完整完美的；神創造的眾東西有神的記號,創造萬物是神的傑作；但是現在萬物盡皆與神分開.人敬拜不是敬拜自然界,乃是敬拜神.誰才是配得呢?豈非那位建築師那位創造者嗎?我們必須明白,人要得滿足唯一的途徑就是認識神；為著本身的榮耀我們愛祂,在向神靈修敬拜當中,我們承認神的榮耀尊貴；我們以受造之物的身份,來享受所得到的喜樂；同時我們要反映出神的榮美,神的心就因此喜樂.陽光乃從神而來,光照在地上又反映上去.神的榮光照在受造之物身上,受造之物,吸收又向發出光輝的神反映上去.在教會基本信條中,人的任務是要榮耀神；在永恆裡繼續不斷地向著神.每當神的創造為榮耀祂的本意受了曲解時,神的心就失去了喜樂.當人體貼自己所歡喜的事物時,罪惡就出現了.每當受造之物,向創造者執權柄時；神創造的原意又被曲解了.故此,不能朽壞之神的榮耀,即發生問題；在此情形之下,人就以為神是說謊者.保羅在(羅一23)說明了這個真理.當代的人敬拜羅馬帝王,淫慾泛濫；正是這段聖經所說的「今日的世代」.我們看見自我的放縱,自我的高舉；人為要爭取自立,用了許多新鮮的方法；受造之物「唯我是要」的原則,在過去和現在仍然發生.有許多淫慾橫流的電影,把兩性的事染渲曲解；使人的人生觀,注重物質的享受,人對成功的看法正反映人的價值觀念墮落.世界所謂成功人士,豈非他們擁有豐富的物質,和偉大的成就嗎?
\newline
\newline
甘心樂意的敬奉廿四位長老把頭上的金冠冕呈獻在寶座前,作為你我的好榜樣.他們毫無保留地呈獻在主面前,因此,他們的生命充滿喜樂和讚美；他們感謝神豐盛的恩典,將榮耀歸給祂.屬於神的應歸還神；從神白白領受的也應甘心樂意歸還給神,我們所奉獻的都是從神而來；除非奉還給神實無他法反映那是曾屬於你的.奉獻禮物給神,益發表示神是應當受讚美的.惟有存著愛心來呈獻給基督,為主而工作才能存到永遠.
\newline
\newline
舉世矚目的奧運會正在舉行,看到那些運動家施展技能,我也心花怒放.當得勝者接受獎牌時,各位的心豈不激動；他們所得金牌,銀牌,銅牌,一時光芒萬丈；但也會黯然失色,冠軍者亦瞬間被人遺忘了.使徒保羅說運動員努力所得的冠冕,有一天都要過去；他用賽跑得獎賞,比喻基督徒人生的賽跑,從神而得的冠冕永不消逝.保羅寫信給提摩太,(提後四7-8)離世前回顧自己的一生,他說:「但那美好的仗我已經打過了,當跑的路我已經跑盡了,所信的道我已經守住了.......」這話何等大的安慰!到了人生最後階段,回顧自己已經獻上最好的給主了；然後保羅舉目望天說:「從此以後,有公義的冠冕為我存留,就是按著公義審判的主,到了那日要賜給我的；不但賜給我,也賜給凡愛慕祂顯現的人.」我們的主要將冠冕賜給所有在祂面前忠心事奉的人；實在反映出神恩典的銀杯!當神的榮光充滿我們生命之時,你我最後將受感動,把所得冠冕,甘心樂意奉在主腳前；我們今日在世上,就應該培養這種精神.
\newline
\newline
讓我們學習天上寶座周圍的群眾,向神敬奉,同來服侍,敬拜我們的主.
\newline
\newline
救世軍的創辦人卜威廉將軍有他成功的秘訣.他是個很敬虔的領袖,他流淚說出成功的秘訣:「神得著我所有的,我看到倫敦市那些貧苦的人,我就明白神要得著我卜威廉所有的.今日如果救世軍有任何成就,乃是因神得著卜威廉向祂完全的敬拜；盡我所能發出意志的力量,我都奉獻給神了,我把所有的影響力都獻給神.」
\newline
\newline
今晚,你有否將你所有的都交給神?寶座上那位全能者的啟示發出來了,我們來獻上；讓祂得著我們所有的,這是我們唯一正確向祂的反映.無論神呼召你作甚麼,我們的神是配得所有的敬奉；我們願意獻上一切的力量,我們生命的一切要歸給祂.今晚你所得的恩典是甚麼,有沒有歸給神使用?
\newline
\newline
把你所有的奉奉獻在主的腳前,我們的心充滿渴慕神之時,我們和天上寶座前的群眾,一同歡呼讚美歌唱!
\newline
\newline
「我們的主我們的神!唯有你是配得的；為了你的榮耀,為了你的權柄,你創造萬有創造一切.」
\newline
\newline


\newpage

\section{第56屆~港九培靈會~高文~第三講　寶血買贖}
\label{sec:406}
\textbf{第三講　寶血買贖}
\newline
\newline
連結: \href{https://www.hkbibleconference.org/session-message/view/406}{\texttt{https://www.hkbibleconference.org/session-message/view/406}}
\newline
\newline
\hyperref[sec:404]{< < < PREV SERMON < < <}
~
\hyperref[sec:toc]{[返目錄]}
~
\hyperref[sec:407]{> > > NEXT SERMON > > >}
\newline
\newline
啟示錄 5:1-5:10
\newline
\begin{longtable}{cl}
\hline
\hline
章節 & 經文 (和合本修訂版)\\
\hline
5:1 & \begin{tabularx}{0.7\textwidth}{X} 我看見坐在寶座那位的右手中有書卷，正反面都寫著字，用七個印密封著。 \end{tabularx} \\ \\ \relax
5:2 & \begin{tabularx}{0.7\textwidth}{X} 我又看見一位大力的天使大聲宣告說：「有誰配展開那書卷，揭開那七個印呢？」 \end{tabularx} \\ \\ \relax
5:3 & \begin{tabularx}{0.7\textwidth}{X} 在天上、地上、地底下，沒有人能展開、能閱覽那書卷。 \end{tabularx} \\ \\ \relax
5:4 & \begin{tabularx}{0.7\textwidth}{X} 因為沒有人配展開、閱覽那書卷，我就大哭。 \end{tabularx} \\ \\ \relax
5:5 & \begin{tabularx}{0.7\textwidth}{X} 長老中有一位對我說：「不要哭。看哪，猶大支派中的獅子，大衛的根，他已得勝，能展開那書卷，揭開那七個印。」 \end{tabularx} \\ \\ \relax
5:6 & \begin{tabularx}{0.7\textwidth}{X} 我又看見寶座和四個活物，以及長老之中有羔羊站著，像是被殺的，有七個角七隻眼睛，就是神的七靈，奉差遣往普天下去的。 \end{tabularx} \\ \\ \relax
5:7 & \begin{tabularx}{0.7\textwidth}{X} 這羔羊前來，從坐在寶座上那位的右手中拿了書卷。 \end{tabularx} \\ \\ \relax
5:8 & \begin{tabularx}{0.7\textwidth}{X} 他一拿了書卷，四活物和二十四位長老就俯伏在羔羊面前，各拿著琴和盛滿了香的金爐；這香就是眾聖徒的祈禱。 \end{tabularx} \\ \\ \relax
5:9 & \begin{tabularx}{0.7\textwidth}{X} 他們唱新歌，說：「你配拿書卷，配揭開它的七印；因為你曾被殺，用自己的血從各支派、各語言、各民族、各邦國中買了人來，使他們歸於神， \end{tabularx} \\ \\ \relax
5:10 & \begin{tabularx}{0.7\textwidth}{X} 又使他們成為國民和祭司，歸於我們的神；他們將在地上執掌王權。」 \end{tabularx} \\ \\
[1ex]
\hline
\hline
\end{longtable}
七印封嚴的書卷　第五章繼續提及神寶座前的景象.當時約翰看見坐寶座的右手中有書卷,書的裡外都寫滿字,表明這本書的豐盛；為要確保書中秘密不向外宣揚,又用七印封好.這卷書記載神所命定宇宙的終局.使徒約翰聽見天使大聲呼叫要有人來揭開這卷書；但卻沒有人能打開這書.於是約翰開始哭泣.
\newline
\newline
惟有羔羊配開　有位長老安慰約翰,說:「看哪!猶大支派的獅子,大衛的根,祂已得勝,能以展開那書卷,揭開那七印.」(5)偉大的彌賽亞,在三千多年前雅各早有預言；為他兒子祝福說:「猶大啊!你弟兄們必讚美你,你手必掐住仇敵的頸項,你父親的兒子們必向你下拜.猶大是個小獅子；我兒啊!你抓了食便上去；你屈下身去,臥如公獅......直等細羅來到,萬民都必歸順.」(創四九8-9)耶西的根和寶座以後成為標準按這標準衡量.約翰所見不是獅子乃是被殺的羔羊,未曾獻過祭；因為獻祭的痕跡很清楚.那曾被殺的羔羊,現仍活著站立在那裡,說明復活的能力,我們的主,祂準備有所行動.頭上七隻角,表明全美的能力,在十字架上受死,如今成為榮耀.這位有能力者,從寶座的右手裡拿了書卷；這時候,撒拉弗和廿四位長老在羔羊面前俯伏.長老拿著琴讚美神,他們並拿著盛滿香的金爐,這香就是眾聖徒的祈禱；地上眾聖民的祈禱如乳香直升到天上,充滿了敬畏神的香氣.天上的群眾唱新歌向神稱頌,向羔羊敬拜.歌曰:「你是配拿書卷,配揭開七印；因為你曾被殺,用自己的血從各族,各方,各民,各國中買了人來.叫他們歸於神,又叫他們成為國民；作祭司,歸於神,在地上執掌王權.」(9-10)這首歌是新歌；因為基督所成就的,完全與眾不同,祂的事奉和舊約成就的工夫完全不同；因此新約帶著恩典的工作,成為新歌向祂敬拜.
\newline
\newline
寶血的功效　基督徒的生命常是更新的,充滿新的平安,新的喜樂；向著新耶路撒冷,有新的盼望.基督成就了救贖的工作,買贖了人來歸祂.「救贖」就是「買來」,買贖的代價是耶穌基督所流出的寶血；若非神的兒子在十字架上流出寶血作為贖價,這個世界就別無救法.因此十字架成了天上敬拜的基礎,基督的寶血成了虔誠敬拜的中心.
\newline
\newline
聖經共有四百六十次提及「血」,除了直接用四百六十次「血」之外；還有用「立約」「獻祭」等和血有關的名詞；聖經無一頁不提及血,和有關血的觀念.一條紅色的血繩,把全部聖經穿連在一起；使我們認識,神救贖的恩典,血和聖經十分有連帶關係.聖經說耶穌再來之時,穿著紅色袍如染過血一樣,祂的名叫做神的道.
\newline
\newline
有一次,某宣教士告訴我們,他在非洲肯雅一教會醫院所發生的事；有一個少年人到樹林割草,地上一把小刀刺傷了他的腳；他急快長途步行,趕向一教會醫院求醫.過了數小時,少年的母親也來醫院；醫生覺得非常奇怪她能找到來院的路徑,她告訴醫生,她是跟著血路來的.原來少年所經之地,留下長條血跡.
\newline
\newline
從這豐富的意義,你我找到耶穌基督乃跟從祂的腳蹤；祂每步伐都沾染著血.血的觀念,內容極其豐盛,實非我言語所能解釋透徹的.
\newline
\newline
現在讓我們從這首天上樂歌,來思想三點意義.
\newline
\newline
(1)生命
\newline
\newline
基督的血裡面有生命.不可或缺的血液,在所有高等生物體內；血把氧氣和營養料供應我們身體之需,血抵抗疾病把人體內廢物排除.血在普通人體,每分鐘經過兩次循環；從心臟細胞血的功用得到清淨.每成年人有六至七夸脫的血,每立方米厘的血,約大頭針針頭之大小,內含約六百萬個細胞,每個細胞可活120天；新陳代謝作用,每秒在人體產生二百萬個細胞.實在令人驚訝!換言之,有血才有生命.聖經用「生命之血」此名詞意思,是從血得到重生.心臟是血的中心,聖經論人的個性常提及我們的心,論整個人就用「全心」.從人生產生高貴的思想,神就除去舊的心,把新的心賜給我們；因為生命從神而來,所以對生命的血必須慎重.因此神吩咐祂的子民,不能用血作食物,要把殺牲畜的血,流盡才能吃其肉；(當然這點有衛生因素包括在內).唯一重要的原則,就是向初期猶大人教導；在利未人的律法表明出來.牲畜肉身的生命在血裡面,祭牲放在祭壇上,為我們預備了靈魂得救的方法.如果你未能到聖殿獻贖罪祭,你要取代祭司的職位在田野間獻祭,每次殺一牲畜即成救贖罪的祭.按律法規定,在田間殺了牲畜,流了血之後,要將牠埋在地裡；用泥土蓋上,如埋葬死人一樣.神早就教導祂的子民要尊重血；因血是看得見的象徵,代表神所賜的生命.聖經用「流無辜人的血」,代替「殺人」的字眼；聖經說血裡有生命,並發出聲音哀求報復.(創四10)
\newline
\newline
血最高深的意義,是指那位為世人的罪,傾倒祂的生命的主.故此,人不能輕看寶血；因為到了神的時間滿足,血就流出來.這血成為以馬內利的血,就是我們的主,被高舉釘在十字架上所流的血泉,在十字架上所獻出的寶血,成為我們得生命的寶鏡.主耶穌說:「凡吃我的肉,喝我的血的人,就有永生.」在末日我們從死裡復活,要飲主的血,就是祂在十字架上使人重新得生命的血；吸收在我們的心靈,接受神的恩典.耶穌基督的寶血,所成就的大事,於聖靈感動之下,基督成就的功勞發生在我們的心靈裡.「血」字根的來源很有趣,世界初期是象形文字；生命與基督聯合,才是血真正的意思.我們存著信心與神兒子的生命聯合,與神的性情完全有份.基督絕非把人生的一部哲學給我們,也非定行為的標準給我們,乃是將寶血賜給我們.基督徒的人生並非一套的宣召,一套的信仰；主愛我們,為我們傾倒生命,叫我們與祂的生命有份.血庫儲藏著許多不同類型的血,在一定溫度之處保全,到了需要時,才取出合適血型的血來使用.但是十字架上的血庫,比人間的血庫高貴,奇妙無比.基督血庫存有神的兒子不能朽壞,又無限制的血來供應.那賜生命的能力在十字架上有效,直到今日仍然有效；能配合任何血型的需要,向任何需要發揮功效；只要你存信心接受祂的血,進入你的心靈裡面,祂就白白賜給你.
\newline
\newline
(2)獻祭
\newline
\newline
當神賜生命之時,還有一點重要的真理,用獻祭這字來表明；在血裡生命傾倒出來,流血生命就失去而死了.在舊約獻祭禮儀上,把這種情況戲劇化,幻影幻真地描述.神命定獻祭,使以色列民明白敬拜的真義,好叫獻祭能代表人對神的最高敬拜.但神絕對不接受任何違背神所吩咐的,祂吩咐必須把毫無瑕疵的羊羔獻祭；強調要將最好,最合宜的獻給神.祭牲必須屬於獻祭人的,不可取別人的羊獻祭.因此,後來有人誤用了,把毫無瑕疵的祭牲,帶到會幕前,請祭司來獻祭；把手放在無辜的羊頭上,說出獻祭的原因,他們獻祭是出於甘心樂意,這時有個很明顯身份的轉移,羊羔代替了本身將遭遇的命運.在公共獻祭的場合,祭司代表整個民眾獻祭,祭司殺了羊羔獻上羊羔的血.如果獻祭是個別的,獻祭者就要自己殺羊羔；當血流出來,祭司拿器皿盛血,把血淋在祭壇上下.用繩把祭牲綁在祭壇角上,必須使祭牲的血流淨盡,把人向神的信服信心如活劇般地描繪出來,表示自己摒棄一切,面向著神；誠心的獻祭表明自己的內心離開罪惡,歸向神.這是舊約時代須履行的,並非義務或是例行的公事,而是表明內心的坦誠.
\newline
\newline
可能宗教最大的錯誤,是人容許儀式來代表心裡的意思.外表的儀式幫助人心靈,把握著屬靈真理,不過是途徑而非目的；若不明此點,那種儀式可說成為拜偶像；初期人敬拜神,竟然拜了偶像,就是這樣出的悲劇.人要過敬畏神的屬靈生活,常遭遇這樣的危險.
\newline
\newline
聖經一再提醒我們不要褻瀆神.主耶穌在聖殿推翻兌換銀錢的桌子,趕逐販賣牛,羊,鴿子的商人；因他們只顧圖謀利益；理應是屬靈的,卻被利用成了搾取金錢,中飽私囊的機會.本是聖潔的,他們卻褻瀆了神,神對此絕對不妥協.
\newline
\newline
今天看見教會許多的情況,我們應當接受同樣的警告!
\newline
\newline
在教會中,我們為了要表明屬靈的真理,用了很多儀式使人領受真理；比如受浸或受洗,乃表明屬靈的真實情況；我們守聖餐,惟有用禮儀來表明真理；無法用言語來表達透徹.但是神警告我們,若我們不配守聖餐而守聖餐,就是吃喝定自己的罪了.外在的儀式,是表達內心情況的途徑；但神最終乃是要求我們獻上「愛」,神要我們盡心,盡性,盡意,盡力愛祂.
\newline
\newline
獻祭最高深的意思,是「愛的獻上」.表示心裡真正與祭壇上的血認同,表明以聖潔的渴慕獻上；表示有了罪,內心願意除去罪,過愛慕聖潔的生活.神著眼於你心裡,渴慕的心意.
\newline
\newline
祭壇上的血有寶貴的意思,表示神接納你的供獻.神創造我們,巴不得祂受造之物能與祂有交通的團契.祂的聖潔,公義和愛於兩者之間得以協調,聯合；祭壇上的血,是神救人的方法.舊約所載的獻祭很多,難以言盡；但所有獻祭惟有獻祭者帶著真誠的愛心,神才悅納；一切獻祭所以有效,是因為有個最終,最新,完整的祭牲.新約之時,過逾越節的以色列人,家家戶戶獻上羊羔,耶穌被帶出耶路撒冷城門外.被釘在十字架上掛了三小時,祂的血如同泉源湧流出來；為了我們的罪,主被人掌摑,被人擊打,主吞聲忍辱,在奄奄一息之時,祂用沈重的聲音說「成了」.殿裡的幔子上下分裂兩半,這幔子原是把至聖所裡外隔開的；大祭司每年只一次預先經過潔淨才可進入.一切舊約時代所代表的禮儀,在十字架上已得到完全應驗；一切禮儀上的法典,目的已達.舊約的預表耶穌基督,至此全都經歷了；叫你的生命和十字架上的基督聯合,得以進入會幕裡面.藉著祂的寶血,使信徒能夠站立在神的祭壇前成為祭司；藉著血來到至聖所,在滿有榮光全能的神面前.我們得贖不是藉著金銀能夠朽壞的東西,不是靠祖先的遺傳；乃是藉著耶穌基督的寶血使我們得贖.我們的主是毫無瑕疵的羔羊,天上天使不住地向祂稱頌祂的寶血.
\newline
\newline
(3)立約「血」使我們想到生命,從生命再想到死亡；主守約付上生命的代價表示信實,又守約犧牲生命使人與神和好.
\newline
\newline
原始人類社會,有種禮儀表明立約,彼此用手放在一起,然後割破手腕流出血來；這種血的盟誓表示兩者之間牢不可破的友誼.據文學家研究這就是握手的起源.(賽四九16):「神說我將你銘刻在我手掌上.」正有這個意思.幾千年前,人在公眾場合上用戲劇方式表明立約,把手在眾人面前割破舉起,讓血滴在地上；使人相信所說的是真實,甚至死也要滴下去.這是今日許多地方的人宣誓時舉起右手的起源.以賽亞說,神用祂的右手發誓,用祂右邊的手臂保證.這是以賽亞的思想.無論如何,血代表約.舊約聖經,「約」字用了三百次.最美麗的約就是救贖的約.自從人在伊甸園背逆神之後,神說女人的後裔要傷蛇的頭,蛇要傷祂的腳跟.當時亞當並沒有作任何事,可換取神的應許；乃是神憑著恩典發出應許,神向人類始祖表明祂的約是何等真實；殺了一隻動物為我們始祖作毛衣；亞當和夏娃穿了這毛衣,表明神用血的約遮蓋他們.當挪亞離開了方舟出來獻祭時,神對他說話,立約不再用水審判地面；從此人類要生養眾多.天上出現的虹代表神的信實,守約,血同樣表示神的信約.神答應籍著亞伯拉罕要興起祂所選召的民族,神的應許是用血來證實,神使亞伯拉罕的後裔,成為偉大的國家；因著這個國家,這個民族,全世界萬族要得著祝福.神的子民在埃及為奴時,那次神的拯救也是帶著血的；每家要殺無瑕疵的羊羔,把祭牲的血塗在門上守逾越節的筵席；門上有血表示他們對神和對神救贖的約有信心.當摩西接受神的誡命頒發之後,就有犧牲獻祭的血流露出來,臨到書卷,臨到人的身上；摩西呼叫說:「這是神向以色列民立約的血.」全以色列民向神呼叫說:「凡神所命令的,我們都要遵行.」神的約用血封住,血所代表的一切,和代表的一切意義,在十字架上達成高峰.主耶穌和門徒吃最後晚餐時,耶穌把杯遞給門徒說:「這是我們立新約的血,為你們流出來的.」聖餐之後,他們唱了詩篇118篇歌頌羊羔；耶穌帶領門徒在聖餐桌旁唱第一節,然後門徒一同唱副歌；我們的主唱到最後一節時,提及「匠人所棄的石頭」,門徒應:「已成了房角的頭塊石頭.」「這是耶和華所定的日子.」門徒應「我們在其中要高興歡喜.」願我們藉著在世界上所獻的祭,接著門徒高聲唱詩「神啊!你的祭壇,你的祭壇,你是我的神.」又唱:「你是我的神!你的慈愛直到永遠.」守了最後晚餐,經過唱應之後；耶穌就走出去,準備被人捉拿.當我們主耶穌寶血從十架上流下之時,就是神藉著耶穌基督表示自己的愛和成全祂旨意之時；我們的神舉起祂的手,讓受造之物都看清楚,宣告神要成就的救贖,今已成就了；神守住祂的約,血留在天上向我們作見證.聖經告訴我們,基督藉著血進入至聖所,代表我們向著神,祂的血為地上完成的使命作見證,永遠在神的寶座,描述不盡神的愛心；愛心的禮物表明神永不改變,祂的話語絕對可靠.
\newline
\newline
故此,我們應該明白,為何天上在慶祝讚美神救贖的血.
\newline
\newline
多年前,當美國南北戰爭時,局勢極不安定,很多不法之徒聚集搶掠.在肯薩斯州,有一班人被稱為侵略者,他們焚燒城市.地方上組織自衛隊,執法捉拿惡徒；有一次捉了一些搶掠黨,準備處以死刑；忽有一個旁觀者指犯人中之一個,對執行任務的領隊說:「那人是我的朋友,他家裡有妻兒怪可憐的,請你釋放他,他不會是搶掠的黨人.」領隊說:「他和他們在一起,證據充足.」為朋友求情者繼續哀求,自願代刑；一再哀求釋放他的朋友；他們商量了很久,哀求者就過去代替他的朋友被槍決死了.免刑者領了代死者的屍體回到家鄉隆重安葬,為他立了墓碑寫著:我最好的朋友取代了我,為我受死.
\newline
\newline
親愛的朋友!在更深的意義,主耶穌取代了你,我本來被定罪的位置；當我們還作罪人的時候,基督為我們而死；祂把一切罪債都付清了,用大愛呼召你到祂的面前.神對你最為守約,最為友愛.現在讓我問你,你對神曾經答應的有否履行?有許多人曾經與神立約；但卻撕毀了所立的約；外表的,儀式的,你或者履行了；可是內裡的,真正的,你卻沒有履行.
\newline
\newline
請聽我的警告!千萬不要用外表儀式對神；如果你已經毀了內心的約,你必須重新與神立約.或者你已離開了神,或者你仍在四處遊蕩；今天晚上我邀請你重新與神立約.
\newline
\newline


\newpage

\section{第56屆~港九培靈會~高文~第四講　讚美的巔峰}
\label{sec:407}
\textbf{第四講　讚美的巔峰}
\newline
\newline
連結: \href{https://www.hkbibleconference.org/session-message/view/407}{\texttt{https://www.hkbibleconference.org/session-message/view/407}}
\newline
\newline
\hyperref[sec:406]{< < < PREV SERMON < < <}
~
\hyperref[sec:toc]{[返目錄]}
~
\hyperref[sec:409]{> > > NEXT SERMON > > >}
\newline
\newline
啟示錄 5:11-5:14
\newline
\begin{longtable}{cl}
\hline
\hline
章節 & 經文 (和合本修訂版)\\
\hline
5:11 & \begin{tabularx}{0.7\textwidth}{X} 我又觀看，我聽見寶座和活物及長老的周圍有許多天使的聲音；他們的數目有千千萬萬， \end{tabularx} \\ \\ \relax
5:12 & \begin{tabularx}{0.7\textwidth}{X} 大聲說：「被殺的羔羊配得權能、豐富、智慧、力量、尊貴、榮耀、頌讚。」 \end{tabularx} \\ \\ \relax
5:13 & \begin{tabularx}{0.7\textwidth}{X} 我又聽見在天上、地上、地底下、滄海裡和天地間一切所有被造之物，都說：「願頌讚、尊貴、榮耀、權勢，都歸給坐在寶座上的那位和羔羊，直到永永遠遠！」 \end{tabularx} \\ \\ \relax
5:14 & \begin{tabularx}{0.7\textwidth}{X} 四活物就說：「阿們！」眾長老也俯伏敬拜。 \end{tabularx} \\ \\
[1ex]
\hline
\hline
\end{longtable}
一息尚存仍要歌頌中世時代,有的修道院照著以色列舊約的傳統；向神獻詩永不能停,所以修士唱詩之聲此起彼落,無論晝夜都可聽到他們滿有喜樂的讚美歌聲.有一次,挪威一班侵略者來攻陷這間修道院；他們毫無僯憫地把修士殺盡,連正唱著詩的最後一個也不能倖免.有一位修士躲匿在隱密處,但因為他為了接續已停的歌聲,給敵人發現以致被殺.
\newline
\newline
當天上的撒拉弗和那班天使,頌讚的歌聲停止時；約翰看見並聽見神寶座四周又有歌聲,是眾天使加入唱天上的樂歌.
\newline
\newline
天使乃服役之靈　全卷聖經多次提及天使作為神的僕役,在敬畏神的人四周安營,保守屬神的人,在神指定他的人生道路上得到保護；有時當苦難臨到,天使作為神所差派的使者來拯救屬神的人.聖經說:「天使豈不是服役的靈,為承受救恩的人效力嗎?」「耶穌在地上工作時多次受天使服侍,天使曾來宣告耶穌降世的喜訊；帶領我們的主到埃及,也曾在曠野服侍主.主在客西馬尼園禱告,在極其痛苦之時,天使來服侍祂.當主從死裡復活時,天使將封住墓門的大石輥開,宣佈耶穌基督復活的好消息；又向門徒傳確實的消息說「主已經升天了.」聖經說:將來我們的主榮耀再臨之時,成千上萬的天使要與主一同降臨.天使是宣告審判,也是執行審判者；天使為神的選民加上印記,幫助聖徒禱告.更將新耶路撒冷向聖徒揭示,又保護聖城耶路撒冷的門戶.
\newline
\newline
第五章11節說,天使形成一個愛的圈,在撒拉弗和長老之間,他們的數目有千千萬萬,多得不可勝數!眾天使一起歌唱,聲音非常和諧合一,他們高聲唱:「曾被殺的羔羊,是配得的.......」在受造之物被救贖的子民面前,天使們宣告:曾被殺釘十字架的羔羊是配得的.天使之中有聖潔的天使,他們一直看見神救贖的計劃；他們沒有墮落過,不需要耶穌的救贖；但他們看見整個的救贖計劃時,他們分享救贖計劃成就的喜樂.他們知道耶穌救贖計劃成就之後,成千上萬數之不盡的聖徒,因此從死亡進入生命.這位救贖者,從人群中選召人；成為君尊的祭司,成為君王.所以救贖者的功勞,使被救贖的人類得到崇高的地位；甚至高於天使,可是天使不會嫉妒的.因天上根本沒有嫉妒的事,這群人得到福份,其他已得福份的人也一同喜樂.當神的座位顯明時,是帶著福份和喜樂的,所以他們一同喜樂.
\newline
\newline
天使向羔羊獻上七重的讚美　「七」表示他們的讚美是完整的,十全十美的.讚美「羔羊是配得權柄,豐富,智慧,能力,尊貴,榮耀,頌讚的.」(12)前四個讚美是向基督的屬性而發出的；後三個讚美是論人對基督應有的態度.「權柄」說明羔羊是一切力量的根源,為要祂的作為；祂一說話,事就即成,祂隨時施展祂的能力.當祂旨意發出,就無一物能抗拒祂的旨意.「豐富」:神擁有一切屬天和屬地的豐盛,整個宇宙的寶藏都屬於祂；基督有測不透的豐富.「智慧」:我們的神,我們的主,對一切事物都透徹了解；祂所有智慧明白每件事的原理,應用,發揮最大智慧的能力.祂了解要成就祂的計劃的每事項,這是智慧最高的發揮,祂捨身為人類作贖罪價.「能力」:祂的洞察力非常奇妙,祂的能力征服任何敵人.「尊貴」意思是越過一切高超的品質.我們的主隨時隨地,掌管萬有；為了祂的偉大,我們唯一的反應是「羔羊配得一切的尊貴.」「榮耀」:是屬天而來,屬神絕對的光榮,乃羔羊才配得,就是神四圍所環繞的那種光榮；惟有在至高處的榮光,見到神極高的榮耀.「頌讚」:在七重的讚美裡,頌讚如同達到金字塔的最高峰得到了福份；說明了神兒子的信實,祂成就了偉大的工作.頌讚祂,我們的喜樂達到滿溢的程度,使我們確信一切本屬於祂.在一切讚美祂偉大作為之中,與我們個人有關的,就是基督的救贖；所有的子民,都心存感恩來向祂頌讚.我們本是一無所有,如今在祂愛裡,祂樣樣的福份,都賜給我們了.
\newline
\newline
啟示錄第七章12節和第五章天使所唱的詩歌是一樣的,是一般得蒙救贖的人所唱的詩歌；其中有六個字眼相同,只是次序稍有分別.「直到永永遠遠」說明我們神的永恆性；不像世上的年代消逝,神的年日永無窮盡,祂的國度直到永遠.天使用「阿們」作為結束,同意他們稱頌讚美是真實正確的.在猶太人會堂中,會眾聽了真理都應聲「阿們.」今日世界上許多教會都有此習慣.「阿們」不但表明同意所聽的話,同時對自己也有一種約束力.有時用「阿們」作為說話的開始.「我實實在在的告訴你們」則「我阿們阿們的告訴你們.」耶穌基督鄭重說出這個真理,更加顯出祂的尊榮,權柄.新約時常於讚美詞句之後,以「阿們」結束,使整個讚美禱告更加有力.這字最偉大的使用,就是我們的主用以說明祂自己,也提到以賽亞書所用的「阿們的神」「真實的神.」(林後一20)說「神一切的應許在耶穌基督裡都是真的,都是阿們的.」第五章那位天使讚美的對象是羔羊,到了第七章,他們讚美的對象是神本身,一方面祂是創造的神,一方面祂是天使的神.天使們看見神救贖的計劃時,他們驚訝,看見成千上萬得蒙救贖的人與他們一同圍繞神寶座的四周,他們真有不可言喻的喜樂.聖經告訴我們,當世上有一個罪人悔改,天上的眾天使都要為他歡喜.因此,當整個神所救贖的教會和眾聖徒,站立在寶座四周,天使們就喜樂讚美.
\newline
\newline
家母曾告訴我,當她目送兒子永別人世時,她聽見天使的歌聲,我感到這話很可信.耶穌曾說,每個孩童都有天使作代表,仰望天父的榮面.當靈魂離開身體之後,把信徒的靈魂護送到神面前,乃天使任務之一.耶穌說的比喻,當拉撒路離世時,天使把他送到亞伯拉罕的懷抱.　美國費城有一位循道會的牧師告訴我,他母親離世的情景,當她臨終之前,因病痛不能睡在床上,經常坐在輪椅上以便照顧；在彌留期間,她對兒子說,好像看見有人進到房間來,牧師立即轉身看看；這時他母親忽然從輪椅站起來,舉起雙手,面向天使大聲說讚美神的話,遂即躺回輪椅.她得到釋放的靈魂被送到主那裡.
\newline
\newline
聖經說,當末世我們的主顯現之時,天使的詩歌又要出現；天使是神所使用的收割者來收莊稼.當號筒吹響,我們的神要從地球的四方召聚祂的選民,要將公義和邪惡分開.基督要在祂的眾天使面前作審判的工作；這班天使就是不住向神獻上讚美的群體.你我都要想到神的榮耀,加入天使歌聲的行列；在神出現的榮光裡面,我們和天使一同沐浴在神榮耀之中.
\newline
\newline
向神敬拜敬拜乃受造之物,對神極度的尊榮之回敬,敬拜預先要說明對神的信服,敬拜最深的意義不是為自己需要向神祈求,也非單因神賜福而獻上感恩；乃是整個人的心靈被神的榮耀充滿,向神五體投地俯伏敬拜.讓一切敬拜帶我們到俯伏的地步,或者是你正在講道或談道時,或者是正領受聖餐,或唱詩,或禱告,或安靜默想之時；目的是向坐在寶座上,配得統管一切的神獻上敬拜.各位!你的心靈應繼續不停如此敬拜.盡所能向神的敬拜深入我們每人思想之中,直到整個人生給神極度的榮耀奪去了.當我們真誠讚美之時,永不忘記敬拜的對象就是羔羊.當神顯出祂的恩典時,祂恩典的真道向我們顯明；藉著羔羊,我們被帶近寶座前；我們舉目仰望十字架,我們再次確切神十分愛我們.求父除去我們的眾罪.
\newline
\newline
有個旅行社專門帶人到處參觀各種新奇事物.有一天在城市公共建築物尖頂出現了特別的景象,走近看見建築物三分之二高處,有幅大石像；據說在建築期中,有個工人從上面跌下,同事們擬尋找屍體,想不到那人活活地坐在群羊上,羊被壓死了,卻救了那工人的命.工人為紀念羊救命之恩,要求在上面刻上羊像.旅遊人士聽了很受感動,想到那被救的石匠心裡的感受,正如那班天使在天上,圍繞在羔羊四圍的感受；那班天使看見基督升在寶座上,見到祂的傷痕,看見祂流的寶血,成就了整個世界的救贖；天使就大聲稱頌.
\newline
\newline
一齊頌讚神與羔羊在美國某州,有個士兵被敵方俘擄,他心想此次必嘗盡苦頭；他忽唱讚美詩,「讚美真神,萬福之源,天下萬物都該讚美.」唱到這裡,有幾十個人加入同唱「天上眾軍都當頌揚.」最後全軍齊唱「讚美聖父,聖子,聖靈,三一真神.」讚美具有感應性；當然寶座前合唱的詩班都有感應性；一群對神的敬拜頌讚引起了其他合群對神的讚美,這種模式表示共鳴,互相激勵的讚美.第五章12節讚美未完,13節又唱「在天上,地上,地底下,滄海裡,和天地間一切所有被造之物,都說,但願頌讚,尊貴,榮耀,權勢,都歸給坐寶座的和羔羊,直到永永遠遠.」這時頌讚達到整個宇宙之間,無論有理性的,無理性的,都一齊頌讚.有些學者猜測,宇宙一切甚至太陽,月亮,星宿,都加入這個讚美行列.(羅八23)說,「凡天地所有受造之物,都嘆息等候得贖的日子來臨.」唯屬神的兒女,才能享受這種的自由釋放.
\newline
\newline
有的人說,因寫作方式要表示四圍的人都來頌讚,所以提及「地底下」.有的解經家認為地底下是指在神的約之外沒有得救的人,他們死了,現在都來頌讚.這兩種解釋,聖經都支持它們；假如用後者的解釋,問題就來了.這班人未曾聽過福音,在地底下哪裡能頌讚神呢?或者是聽過福音拒絕接受的人,可能我們的神曾用自然的啟示,雖然有個限制,但都給了他們.我們可以明白,原來在今生之後,洞察明白這位真神的偉大；或者他們及時想起以前所聽過的福音知識；這樣,或者神給他們特別的機會.問題的關鍵在於那些未信耶穌者,到了這榮耀的場面是否復生向神屈膝敬拜.似乎這是最適當的解釋.神救贖的功勞極其偉大,藉十字架勝過魔鬼一切的邪惡；使曾經殺害祂刺過祂身體的人回轉歸依這位萬王之王.他們回想從前拒絕神,迫害神的兒子；以致死在地底下,但現在後悔都來加入頌讚.這樣的解釋合乎神學論點的需求.
\newline
\newline
啟示錄將神和羔羊相題並論,一次又一次地描述.當世界施行審判時,向神與羔羊獻上第一批果子,同時成為天上的光體,從神和羔羊那裡流出生命河水；耶穌基督正是看不見之神的真像,祂來到世上成就了一切人類得救贖的方法.祂是神的兒子,是三位一體的第二位,祂也是羔羊；我們常以祂救贖的工作為中心.除非藉著耶穌基督流出的寶血,否則我們永遠不能到神面前；你要認識天父,必須來到十架下向主耶穌敬拜.
\newline
\newline
永恆的讚美天使用四頌讚美,在神和羔羊寶座前稱頌:「但願頌讚,尊貴,榮耀,權勢,都歸給坐寶座的羔羊,直到永永遠遠.」(13)四頌讚美的字眼,乃讚美神的屬性；從這首詩歌我們再次看見神的永恆性,向神的讚美永不止息.向著羔羊應有不住的讚美,才是最合宜的.每逢天使唱出七重的讚美,或獻上四重的讚美時:活物都說阿們.到了第14節,在大讚美中廿四位長老俯伏敬拜神,獻上永永遠遠的頌讚.
\newline
\newline
各位!假使你能知道永恆的讚美,能享受在神永恆膀臂之中；我們的心是何等的鼓勵!這是最安全的；我們今日被救贖,不是在乎世上的人或任何財富和方法,乃在乎那位發誓要拯救我們,愛我們的那位真神.也是與主基督釘十字架的手,和永不改變的神之個性連結在一起的.
\newline
\newline
當芬蘭和蘇俄戰爭時,有一個芬蘭軍官遭遇的事蹟使他的人生絕對的改變.在一場戰爭之後,這軍官俘擄了一批敵軍,擬於次日早晨先處決七人,那七人當晚被關在冰冷的地窖.芬蘭軍官發現七死囚中有個名高士力的安詳靜坐,而其餘六人都吵鬧不堪,整夜用拳頭擊牆；高士力忽然唱聖詩:「安隱在耶穌手中,安隱在祂慈愛的懷裡,我的靈魂得到保護,我的靈魂享受祂的全美.」他一遍遍重複地唱這首詩；當他暫停時,背後有個死囚看著他說:「傻瓜!你唱詩給人聽見了.」他帶著淚眼望著他的同伴說:「這首詩歌是我母親三週之前,曾唱給我聽的,她唱關於耶穌的詩,她向耶穌禱告.」他停了片刻,鼓起勇氣望著其他的同伴說:「如果你有信仰,你隱藏你的信仰,這才是真正的懦夫.」我母親所信的神,現在是我所信的神；今晚我睡不著,我見母親的面,唱她所唱的詩.我祈求耶穌使我預備妥一切,明天站在祂的面前.我反覆思想這首詩歌,我不能私自收藏這首詩歌.」所有的同伴都說:「這是正確的,如果在這環境之下,我可以得到恩典,我從前褻瀆神,任何好的事我都踐踏了.」一個有所反應的死囚,就跪下說:「請你為我祈禱.」於是兩個蘇軍跪下彼此代禱,他們很長的禱告必定達到天上的；在這世界好像有門敞開,他們每人都感到有位肉眼看不見的臨到中間.他們都向神俯伏,他們禱告哭泣之間,永生神的聖靈充滿了所有的死囚身上,他們的談話全是屬靈的主題；恐怕皇帝,皇后都不能明白這個真理.而神給這班嬰孩,明白祂的真理；他們禱告有時唱詩,唱那首剛才所唱的安慰詩歌.他們想起從前所聽過的詩歌也唱起來.神的靈感動他們,也感動監房外看守的士兵,都情不自禁地加入團契；他們向天上讚美直到天明；當旭日初升,這班死囚被送到刑場,處決之前；他們要求不要用布蒙頭,並且要求再唱「安隱在耶穌手中」之詩；在子彈未發出之前,他們舉起雙手,眼目望天,盡其所能歌唱:「安隱在耶穌手中,在祂慈愛的懷裡得到安隱,在那裡你的愛臨到我們,我的靈魂得到全美地安息.」
\newline
\newline
我們聽了這光榮喜樂的見證,我們也要唱這首詩歌.讓我們的眼睛望著天上的寶座,和天上那班有福的群眾同聲「阿們」!
\newline
\newline


\newpage

\section{第56屆~港九培靈會~高文~第五講　苦難為何延續}
\label{sec:409}
\textbf{第五講　苦難為何延續}
\newline
\newline
連結: \href{https://www.hkbibleconference.org/session-message/view/409}{\texttt{https://www.hkbibleconference.org/session-message/view/409}}
\newline
\newline
\hyperref[sec:407]{< < < PREV SERMON < < <}
~
\hyperref[sec:toc]{[返目錄]}
~
\hyperref[sec:411]{> > > NEXT SERMON > > >}
\newline
\newline
啟示錄 6:9-6:11
\newline
\begin{longtable}{cl}
\hline
\hline
章節 & 經文 (和合本修訂版)\\
\hline
6:9 & \begin{tabularx}{0.7\textwidth}{X} 羔羊揭開第五個印的時候，我看見在祭壇底下有曾為神的道，並為作見證而被殺的人的靈魂， \end{tabularx} \\ \\ \relax
6:10 & \begin{tabularx}{0.7\textwidth}{X} 大聲喊著說：「神聖真實的主宰啊，你不審判住在地上的人，為我們所流的血伸冤，要到幾時呢？」 \end{tabularx} \\ \\ \relax
6:11 & \begin{tabularx}{0.7\textwidth}{X} 於是有白袍賜給他們各人；又有話吩咐他們還要歇息片刻，等到與他們同作僕人的，和他們的弟兄，像他們一樣被殺的人的數目湊足的時候。 \end{tabularx} \\ \\
[1ex]
\hline
\hline
\end{longtable}
我們繼續來思想天上的詩歌,這種偉大的詩歌是從啟示錄第四章開始到第十九章.我們每晚有時只講一首詩歌,有時講了幾首詩歌.第四和五章提及很多天上寶座四圍眾天軍眾天使讚美的歌聲；開始時是撒拉弗和四活物的歌頌,繼之是得蒙救贖的廿四位長老；再之,千千萬萬的天使加入唱詩；最後是宇宙間所有被造之物,天上的,地上的,地底下的,凡有生氣的活物都同來歌頌；向著活到永永遠遠,萬王之王,萬主之主的歌頌,而達到讚美的巔峰.
\newline
\newline
審判的刑罰　神掌管一切的權柄,七次審判開始,羔羊揭開七印中的第一印.(六1)每當印揭開時,地上有許多被毀壞的事發生.無論世上發生任何事情,都在神旨意範圍之內.後四印揭開時,有四個不同的人騎著四匹不同的馬.(六1-8)這幅圖畫來自撒迦利亞書,四個騎馬的人來到以色列,要把他們從苦難中拯救出來.在啟示錄,約翰所說騎馬的人,面對著世界的邪惡,第一匹馬是白的,騎馬者頭戴冠冕,手拿弓,代軍事的力量；乍看以為代表基督,細看原來是冒扮的假基督,帶著武器到處橫行要征服世界.第二匹是紅色馬,表示流血,騎馬者手持大刀準備殺人.第三人所騎的是黑色的馬,手持天平,代表食物稀少,艱苦的日子即將來到.第四匹馬是灰色,騎馬者名叫「死」,陰府也隨著他要襲擊被殺的人.這四個騎馬者帶著殺人的任務；被殺者竟達地上人口四分之一；他們要連結第五第六個慘局發生.因此,提醒人想起以西結的預言；又叫我們思想耶穌基督的預言,祂來臨之前,將要發生之事；就是世界將遭受更多的屠殺.當揭開第五印時,這幅圖畫又回到天上.
\newline
\newline
祭壇底下有殉道者的靈魂　「祭壇底上,有為神的道,並為作見證,被殺之人的靈魂；」(9)這裡使我們想到世俗人如何敵對教會；有許多基督徒為主堅守真理以致被殺.當我們向世界的末了進軍時,信徒被迫害為主殉道,將越加普遍.「祭壇」是神獻祭之處,舊約時代神吩咐以色列民建造會幕,在聖殿會幕裡面有祭壇；是依照天上的祭壇造的,所以地上有,天上也有.因此,當日獻祭時,在祭壇處人與神相會,立下恩典之約.殉道者生命之血,在神前傾倒出來,他的血發出聲音,控告殺害他們的.(11)血顯明他們對主忠心至死.在信徒死了未復活之先,這班殉道者的生命仍然存在；他們的靈魂在寶座前發出聲音,他們的靈魂仍然知道地上正發生的事.他們都很關心在世上,仍要繼續完成神的計劃；他們很關心教會,很關心地上受逼迫的信徒；他們在天上大聲呼叫說:「聖潔真實的主啊!你不審判住在地上的人,給我們伸流血的冤,要等到幾時呢?」(10)聖經中常有此問語:「主啊!惡人得勢,繁盛要到幾時呢?」神是掌管萬有的神,可以隨時停止你的子民,遭受冷血的殺害；但為何神容讓他們呢?這是個常遇見的問題.信徒歷世以來,特別在受苦時,心裡就湧出這個問題；當義人受苦之時也常發出這個問題；恐在座各位也曾經一次或多次地發過這個問題吧.
\newline
\newline
聖潔真實的主　在天上神的寶座前,一方面有讚美的聲音,一方面有呼冤的聲音.稱神為主,因為祂是創造天地萬物之主,祂有完全的主權,有無限能力.神無論作甚麼,或賞賜或收取,都是美善的.祂的聖潔就是除去一切違反愛心的事,祂的真實就是除去犯錯誤的可能.這班呼喊者的思想中,完全無疑問,最後神的旨意要勝利；雖然世界戰爭延續,但是信徒心裡確信神的旨意是十全十美的,神永遠沒錯誤.神在道德上的完整,叫我們思想這班殉道者發出的聲音；他們並非懷著報仇憎恨的心理,他們乃是誠心要求神的公義臨到.
\newline
\newline
忍耐等候　這首詩歌中所缺少的不是讚美,所需要的乃是忍耐.他們終而得到回應:「有白衣賜給他們各人,又有話對他們說,還要忍耐等候片時；等著一同作僕人的,和他們的弟兄,也像他們被殺,滿足了數目.」(11)這裡提醒我們,神救贖的計劃如今仍然進行；這種被殺的情況,要繼續直到羔羊受苦救贖的計劃,傳到地極而後止.在這些日子中,我們的主說:「要忍耐!」祂不願一人滅亡,祂願萬人悔改得救；當明白今天我們傳揚福音之時,正是許多人將在傳福音的途徑中成為殉道者之時.未完成神計劃之前,還有很多人將被殺.當神完整的教會來到時,我們就不足介意了；但從今至那日,總有人將要殉道和受苦.我們一方面要期待最終的勝利,一方面要告訴殉道的使者說:「還要安息片時.」我們當專心交託掌管全權的主,無論黑暗時代持續多久,終有一日於破曉時光見主.「白衣」表示得勝,純潔,他們穿白衣得著將來勝利的確據.殉道者的呼喊,盼望神的拯救,也反映教會的歷史.教會有沒有人發出這樣關心的問題呢?難道我們想儘快完成神的工作嗎?對神的計劃模糊不清,尚未明白神普世救贖的寬度,忘記了千萬人未聞耶穌的福音；有時我們想要把人搶救過來使他們成為聖潔,不但費時而且勞心勞力；但神面對祂永恆的目標,絕不妥協,總要完成.
\newline
\newline
為神的使命盡忠總要付上代價我們必須絕對順服神的命令.耶穌基督為我們付了生命的代價.我們遵行神的話語所付的代價,包括我們的性命；主吩咐我們:跟從祂的必須背起十字架來跟從；無論祂命定我們作甚麼就去遵行.在新約聖經,「為主作見證」和「殉道」的字眼源出於一.初期門徒為主作見證,就明白要為主殉道；當耶穌呼召他們出去為祂作見證之時,這班門徒即心裡明白,是包括為主盡忠至死.
\newline
\newline
記得數年前在南美厄瓜多爾地方,有五位宣教士為主殉道；記者到該地訪問殉道者的妻子:「為何你們的神容許這樣的事發生呢?你的丈夫出去豈不是作愛心的工作嗎?為了愛那班野人反而被他們殺害,這一點請你為我解釋.」其中一位寡婦,回應了令記者驚奇的答案「神拯救我的丈夫,免得以後他有不順服神的危險.」
\newline
\newline
基督徒不順服主,實在是可怕的事,在最深的意義說；可怕的並非死亡,因基督已經把死亡的仇敵戰敗,墳墓失去了毒勾.基督徒的人生有個聖潔的懼怕,就是怕不順服主；如果我們把自己的一生置於本身手中,這樣的道路乃違反主為我們所預備的道路.
\newline
\newline
美國有一所謂宗教刊物,批評那五位殉道的宣教士,編輯於社評中,慨嘆這班宣教士年青有為；如此殉道極為可惜,非常浪費生命.各位!這世界的智慧竟以此批評為結束,完全失去這件事的意義.　你是否以世界的物質,繁榮,安逸,享福,來衡量你的人生?或以神寶座上所有的價值觀念,來衡度你的人生價值呢?五位殉道者中有位名也利亞的,當他讀大學時寫的日記:「一個能放棄他不能控制的東西,而能獲得不能失去的永恆的人,不是傻瓜.」他未入森林之前,和他的妻子有段很長的談話,他提及或有一去不復還的可能；他們夫妻倆都下決心,無論神的旨意如何都必遵行.或者有人要說:「這樣的人生太危險,這樣的要求太高.」如果按照世界的價值觀念來說,要求太高也是不對的；但這是新約教會所能建成的材料,每個為耶穌作見證的人,必須有此特點.
\newline
\newline
神藉苦難完成祂的旨意我絕對不是說神給我們辛苦難受,也不是說神以我們受苦為樂；乃是說神無限的智慧裡面,人的苦難若彰顯,正是神配得榮耀之時；人縱所欲為,走了差路遇到危險,在神的眼中,卻要藉此完成祂的旨意,而得著榮耀.如果成就祂救贖的旨意,祂經歷苦難,請問!你我這些跟從祂的,今天豈能在安樂床上,在享福之地過我們的人生呢?我們應清楚不抗拒在神旨意中,為我們的好處,而許可臨到我們身上的苦難；我們渴望受苦難的日子過去,神等候這班受苦者最終所得的福份.若不是這些苦難,我們就不會明白神個性的深處,領略祂的恩典足夠我們用.在神的計劃中要得到真正的喜樂,要達到神的目標,須經歷長遠的路,須經過曠野遭受仇敵種種的攻擊；但我們的主,無時無刻不在我們前面帶領.在大多數情況之下,神容許苦難臨你身上,為要使你成為更豐富的人.所謂豐富,不是世上物質的繁榮,乃是神的教會所能供給我們的.當教會遭受逼迫,患難,敵對之時,就是教會最有能力最豐盛的時候.剛才我們提及那五位宣教士,當他們將福音傳入奧加族時,他們努力傳福音；但這努力並非最後的一次,因為殉道的宣教士,他們的寡婦留下在那裡學習當地的文化,向該地居民傳福音.有幾個土人信了主,成了福音的使者,向他們部落的人傳福音,使奧加族人都信了主,都曉得讚美神!他們竟然差派宣教士到河流下游地帶傳福音,結果,每個奧加族人都認識主,都能夠讚美真神.當柬浦寨地方因遭困境,難民紛紛湧入泰國時,四,五萬人一營的難民都住在用竹搭成的房子裡；也用竹築了一間禮拜堂,可容五百人在內敬拜.他們每逢主日,除了裡面坐滿五百人外,而堂外的地方上也坐了四,五千人；五年中,有兩萬五千柬浦寨人信了主.廿五年來在柬浦寨的工作,包括所有教會合共信主的人數,還不及逃入泰國柬浦寨難民之半數.在韓國教會增長的情況,是眾所週知的.而中國方面,據統計也有五千萬的基督徒.
\newline
\newline
我想香港的人,現在所顧慮的是一九九七年的問題；我們不知道將有何事發生；但是一旦有患難困苦,我們應該確信神有祂的美意；可能祂要將我們從安逸舒適逍遙自在的物質中釋放；使我們看見將來的情況,更加屈膝在主前敬拜.我們實在要感謝神,當金子在火中鍛煉時就是復興之時；須確信神容許患難臨到,都是出於祂美善的旨意.
\newline
\newline
安隱在耶穌手中　無論遭遇任何患難,我們應當安心；因為若不是經神許可,沒有一事能臨到信徒身上；所遭遇的試煉,不會過於我們所能擔當的.很多時候,世界一有危險,我們便失去信心,慌張不知所措；令人懷疑基督徒是否真心信靠他的神.我並非說世界的工作,神要我們收的莊稼我們可以不理了.當我們生命有罪未曾向主承認,我們仍知世界有患難,甚至有更多的患難；因為愛我們,聖潔真實的主,要我們曉得萬事都互相效力,叫愛神的人得益處.保羅寫信給羅馬的基督徒,後來他們親眼看見保羅為主殉道；他們表明與神之間沒有攔阻,是神的旨意臨到.所以保羅說:「在神的旨意:──萬事都互相效力,叫愛神的人得益處,就是按祂旨意被召的人.」我們若有此信心,心中才能有全美的安息.
\newline
\newline
數年前我在奧拉馬州,參加葛培理佈道大會.魯斯牧師是佈道團專門負責栽培的同工；會前週末他到紐約參加青年夏令會,和一些青年在湖中划艇；其中一位姊妹跌下水,他立即潛水拯救那位姊妹,勉力支持將她交給別人扶上岸；大家都注意被救的姊妹,卻忽略了救人的魯斯；等到要找他時已杳無蹤跡.過了兩天才找到他的浮屍.魯斯太太托她丈夫的好友打電話來報噩耗,附帶詩篇115篇第三節的經文；因魯斯牧師平時講道,慣用此節經文作為結束.當我找到此節經文,我的眼淚不住流下；我似乎不會瞭解這解經文!我腦中呈現著死者的妻子和母親,仍站在湖邊.一剎那間她們感受人生的美夢破碎了!她忽然成為無依無靠的寡婦；但是她破碎的心乃曉得仰望神.向神認罪,神透過這個慘劇仍保持祂的榮耀.這是個肯定安息在神旨意中的基督徒之經歷.「然而我們的神在天上,都隨自己的旨意行事.」(詩115-3)
\newline
\newline
在你的人生或將遭遇相似的情景時,可能在座中有人曾有此經歷；也曾從心的深處發問「神啊!為何有苦難?」巴不得你能握住神永恆的應許,得著祂永遠的恩典；好叫無論任何環境,神仍然與你在一邊.無一物能使你與神的愛隔開,不是性命,不是死亡,甚至不是天使,更不是慾望的權柄；不論高處的,低處的；無一能使你與基督的愛分離,神的愛乃在耶穌基督裡面的.所以弟兄姊妹,只要你站在神旨意中,從神領受信心；這世界一切,連情慾都要過去；但神說:「遵行我旨意,要永遠長存.」基督徒的安穩正在於此,我們信心,安心,在坐寶座永遠掌權的神手中.今天晚上,你是否安息在神旨意中?可能你心中向神呼叫:「神啊!痛苦還有多久呢?」愛你的救主說:「我永遠掌管一切,我永不離開你,也不丟棄你.」在祂的旨意中,你永遠知道神在你心中是何甜美.
\newline
\newline
今天晚上,你是否享受在神旨意中；如果你不在神旨意中,你來到祭壇上,將自己獻上；來回應這位愛你的救主,為你捨命；如果你是個負心,負義,縱己慾為榮的基督徒,難怪你有諸多痛苦.假如你完全屬於祂,你遭遇苦難,也仍然能在痛苦中得著安息.
\newline
\newline


\newpage

\section{第56屆~港九培靈會~高文~第六講　神的拯救}
\label{sec:411}
\textbf{第六講　神的拯救}
\newline
\newline
連結: \href{https://www.hkbibleconference.org/session-message/view/411}{\texttt{https://www.hkbibleconference.org/session-message/view/411}}
\newline
\newline
\hyperref[sec:409]{< < < PREV SERMON < < <}
~
\hyperref[sec:toc]{[返目錄]}
~
\hyperref[sec:412]{> > > NEXT SERMON > > >}
\newline
\newline
啟示錄 7:9-7:10
\newline
\begin{longtable}{cl}
\hline
\hline
章節 & 經文 (和合本修訂版)\\
\hline
7:9 & \begin{tabularx}{0.7\textwidth}{X} 此後，我觀看，看見有許多人，沒有人能計算，是從各邦國、各支派、各民族、各語言來的，站在寶座和羔羊面前，身穿白衣，手拿棕樹枝， \end{tabularx} \\ \\ \relax
7:10 & \begin{tabularx}{0.7\textwidth}{X} 大聲喊著說：「願救恩歸於坐在寶座上我們的神，也歸於羔羊！」 \end{tabularx} \\ \\
[1ex]
\hline
\hline
\end{longtable}
屬神者得蒙保守　啟示錄記載,世上有接踵而至的災難；但信徒能以歌唱救恩之歌的態度面對.第六章論七印開啟,開時揭開印時,有四匹馬出現；接著世界便充滿殺害摧殘的事.到了第五印揭開時,祭壇底下有殉道者的呼叫聲.第六章末了就將七印的後幾個印開啟的情形告訴我們.當第六印揭開時,好像宇宙之間的勢力都臨到人類身上；有大地震,太陽竟然變黑如黑布,月亮變紅如血；天上的星辰忽然掉在地上,好像無花果被大風搖動；整個天空捲起如書卷一樣,掉了下來.每座山都從原有位置挪移,種種現象令人十分懼怕!人們爭相逃避到山崖,他們發現有位坐在天上寶座的神,和羊羔向他們發怒；所以人類要盡快找藏身之處.這章聖經末了說:「神發怒的日子來臨了,誰能站立的住呢?」這表達人類命運的書卷最後一印開啟之前,有種現象叫人知道；神仍然看顧屬祂者的教會,在將來的審判中,真正信神者才能得到保護.第七章,天使在天上,屬於神的子民得著保護；神把他們封住保護.使我們想到古代分封習俗,君王用玉璽,戴在手上的戒子來封住,這文件所指示的就有功效；如果把玉璽置於物件上,表示該物屬於君王.在啟示錄有印蓋上的,即表示那教會基督徒是屬於神的；信徒身上的印,蓋著天上父神的名字.
\newline
\newline
願救恩歸與神和羊羔當作者約翰題到神的子民身上蓋上神的印之後,他又帶我們看天上另一景象.第9節記載,所有聖徒圍繞天上神的寶座；這班聖徒經歷很多患難,所受創傷現仍留在他們身上；他們好像長老一樣,穿著潔淨的白衣.手持棕樹枝,表示他們勝利之後,來朝見得勝尊貴的基督；祂在十字架上受死,因而取得最後勝利.天上撒拉弗四圍千萬聖徒不勝其數,熱烈的情況一望無際.神曾經應許始祖亞伯拉罕,信的人多如天上星,地上之沙；從每個民族,每個國家,各種方言合併一起；凡是神的僕人,無論古今,老少,大小都圍在神的寶座羊羔面前聚集.這些群眾大聲呼喊:「願救恩歸與坐在寶座上我們的神,也歸與羔羊.」(10)讓我把這節聖經用現代的話詳細解釋:「你是那位坐在天上滿有權柄,坐在寶座上的神,代替我們死的羔羊,正是我們得蒙救贖的原因和途徑；若非神或神的兒子呼召我們,我們絕不可能在諸多患難中得勝.若非神有恩典,我們絕不可能在神面前向神敬拜；這完全是神的作為.願將一切榮耀歸給祂.」各位!這樣的解釋能否明白呢?
\newline
\newline
有一次,一個神學生向我談論,他為何要讀神學當傳道人.這人外表粗魯,不像個傳道人；他說有一次他的太太要他一起赴奮興會,根本他毫無興趣,只為討好太太而參加.他覺得聚會很悶,當主席介紹講員時,他就開始打盹.講員用吶喊宣佈題目「救恩」,幾乎把他嚇得跌在地上；因此他就振作精神聆聽；神的靈感動他,在聚會結束之前,他信了耶穌.
\newline
\newline
今天晚上的信息,是整部聖經福音的中心.當你聽到福音,將你親愛的朋友帶到神的寶座前；當講道時,天上天使把聖經傳說:「願救恩歸於坐在寶座上的神和羔羊.」在講道未結束之前,盼望大家能夠加入這首救恩的詩歌行列.
\newline
\newline
有四件事情是要請各位必須盡力明白的:
\newline
\newline
(一)沒有救恩就必滅亡
\newline
\newline
我們已經失喪了,若沒有救恩便繼續迷失.神創造人,叫我們永遠敬拜神.意思是神給我們足夠的聰明去明白,有足夠的意志力去實行；神的心意要我們按祂聖潔的榮美,去朝見敬拜祂.但人類的始祖偏行己路,隨己欲為；從那時開始,人類一直追隨始祖的背叛,聖經說:「沒有一個義人.」(羅三10)聖經清楚告訴我們,每個生在世上的人都有罪；無論個別的集體的,都是為自己而生活,這就是罪的根源.罪惡發生乃因人違反神的命令,把歡樂物質繁榮,當作自己的神；我們理應敬拜造物之主,但人卻敬拜受造之物.有的拜人手彫刻的偶像,把世上的物質成功人物當神敬拜.有的人把認識神的永恆真理和學問比較,寧可揀選學問作自己的神,把至高神尊貴的名字當作事做.神命定有安息日,人類卻違反安息日；在家庭中違反神的誡命,沒有孝敬父母.充滿苦毒惱恨,心裡犯了殺人之罪.說到兩性之間的問題,心懷慾念,心裡犯了姦淫之罪.欺騙,撒謊,偷竊；或者你說:「我並沒做這些犯罪的事.」但是神說:「犯了誡命中之一,與罪有份,就是罪人.」不但沾染了罪,在人生的生活方式,放縱的情況是許多人生命的刻劃.神給我們很多的恩賜,可是我們濫用恩賜,離開神為我們所定的的旨意!我們使用時間的方法,反映自己尋歡作樂,比尋求敬拜神更重要.其中有樣反映罪,就是當神呼召人,人卻完全不聽；在你我心靈的深處有此傾向,有此生活.有的人不愛聽這樣的講道,這正反映你的自尊心驕傲受了侵害.讓我們面對問題,面對罪,叫我們看見,我們藐視神藉基督向我們的啟示.我們認為耶穌的教訓很好,也參加宗教的教會儀式；我自己慚愧承認從前我就是這樣假情假義的基督徒.以為我愛耶穌,我也敬拜神；但事實上我的心靈,卻沒真正認識祂.
\newline
\newline
親愛的朋友!假如你尚未得救,我了解你心裡虛空,心靈虛不自勝,沒有得到滿足!在你生命中已經失去意義,也失去了價值!儘管你如何掩飾,你的心靈深處自知已經迷失了.聖經說:如此人生的終局就是死亡.除非你扭轉局面,有了新的改變；否則你人生的終局,便是下地獄死亡.並非說你舉起拳頭,向神反叛才是失喪者；聖經說:「我們若忽略這麼大的救恩,怎能逃罪呢?耶穌說人怎樣失喪呢?就是跟著世人走上那條寬闊的大路；盲從世上大多數人所選擇的生活方式,其終局是死亡.故此,我們需要得拯救,整個世界就是失喪的世界；如果沒有救恩,人都要滅亡.
\newline
\newline
(二)耶穌基督為人類完成救贖之功
\newline
\newline
耶穌為我們受死,把救恩帶給世人.當天上的群眾以救恩的詩歌稱頌羔羊時,其中所蘊藏的真理,就表明出來；第14節進一步說明救贖的工作.請留意這段經文都是完成式的字句,說明基督為你我完成了拯救的工作.神為人類預備了救贖,要救我們脫離一切的罪.
\newline
\newline
記得有位浸信會的李牧師,敘述他初次旅遊聖地的經過.當他到了耶穌受死的地方,心裡非常激動,他走在全旅行團最前面爬上山頂；在他後面的人發現李牧師氣喘喘地,便衝上前問說:「李牧師,你以前到過這裡沒有?」他靜默片刻低聲說:「是的,千多年前我都在這裡.」親愛的朋友!一千多年前我們都在耶路撒冷.耶穌基督在那裡被釘十字架,為你和我受死.我們犯罪,在神律法審判之下,審判的終局是死亡.但神不願離棄我們,祂十分愛世人；甚至將祂的獨生子賜下,讓耶穌來擔當我們的憂傷痛苦；最後祂將你我的罪都背負上十字架.以義者來代替不義者,將你我帶到神的面前.在肉身上祂被處死,但神的聖靈使祂復活,又升到天上,坐在父的右邊統領一切.今天晚上如有任何人到神寶座面前求告祂,主耶穌必垂聽禱告.
\newline
\newline
看哪!神的羔羊洗淨了世人的罪,這是神為你我所採取的拯救行動；神以恩典來開始整個救贖的工作,也以祂的恩典來繼續.人不能用任何宗教儀式來換取神的歡心,也不能以自己的行為討神的喜悅.救恩乃是神永恆的愛,使人白白地得救.基督為我們的罪受刑罰,我們唯一能做的就是離開已往假冒的遊戲；回轉歸向神,接受主為你我完成的救恩.
\newline
\newline
我的朋友告訴我說:有個名丹尼沙的人,到朋友家裡作客,當他到達時即問:「我未到之前,你為何要找我呢?」主人說:「丹尼先生,我有一件事,還很值得告訴你的；耶穌基督為一切人類的罪受死.丹尼是個有文學修養的詩人,他聽了即說:「這消息是舊的,新的消息才算是好消息.」各位!這個消息永遠不會變舊的.天上的群眾,都來稱頌基督的救恩；永恆的神,用了無限大愛來救贖人類,天上群眾充滿歌頌之聲.
\newline
\newline
(三)神賜下救恩,把我們帶進祂兒子耶穌基督的生命裡面
\newline
\newline
救恩使人與神有個別的關係,不是我們為自己作了甚麼；這首歌說明是神親自帶給人的.耶穌用「重生」這名詞來解釋,藉著聖靈,我們真正名副其實的成了新造的人,舊事已過,一切都變成新的了.並非神把人性摧殘,事實正是相反；神挪開我們原有的本事,按照祂救贖之目的加以使用改變；因此基督把你和我按照祂的形像創造.比喻世上的領袖答應給你物質上的享受.說:「你跟隨我,我就給你新衣服.」可是神的救恩,所帶給人的福份,遠超於此；祂叫我們跟隨基督成了新人,不是片面的得了新衣而已.不是說我們從此不需要進步,我們信主之後；好像個屬靈的嬰孩開始在神裡面有更多的長進.但開始之時,我們應知和神的關係是非常確實的；神的聖靈要與你的心靈彼此作見證,叫我們知道我們是屬於神的.其中有一確據,你將得到的是從罪惡中釋放的感受.第14節以那些人的衣服,用血洗乾淨來表示,絕非暗示一個人信了耶穌,就可免去世上的誘惑；我們生活一天就必面對罪惡猛烈的引誘,有時你真會跌倒.我們若能立即承認自己的軟弱,求主赦免之後,即與主同行,繼續跟從祂；明白祂的話遵照而行.「我們與神的兒子相交」,耶穌基督的寶血就繼續洗淨我們的罪.
\newline
\newline
第15節末了,美麗地描寫神用帳幕護庇屬於神的人.這裡使我們想起舊約時代,以色列民在曠野,神命令他們建造帳幕,表示神與祂的子民同在.救恩引導我們進入蒙福的生命.
\newline
\newline
第16節提及屬於神者不再飢渴.在屬靈的意義上,心靈的飢渴到那時都得到滿足,飢渴慕義的人如今得到滿足了.耶穌說:我就是生命的糧,到我這裡來的,必定不餓,信我的,永遠不渴.(約六35)這種喜樂的高峰,寶座旁的羔羊作我們的牧者,祂要帶領我們達到生命水的泉源.主耶穌作牧人帶領群羊,主說:「我是好牧人,好牧人為羊捨命.」(約十10)好牧人帶領屬祂的群羊,到生命水的泉源；這樣的人生,難以言喻其中的福份,安慰,平安和喜樂.聖徒僅能報以感恩的眼淚；並且把我們以往一切的傷痕,難處,從我們腦中永遠拿開.真是充滿喜樂,無限榮耀的人生!我無法說盡這生命的美好.話至此,相信每個心裡感到有罪的人,或將撫心自問「我如何才能得到救恩?」
\newline
\newline
(四)白白接受神的救恩
\newline
\newline
神所賜的救恩,必須有人當禮物去接受.我們當承認自己是罪人,沒有耶穌基督,必然滅亡；神挽回頹勢,把祂的獨生子賜下為我們的罪受死,作為人的救主.神樂意把豐盛,滿有福份的生命賜給人；除非我們願意接受救恩,否則神的心愛莫能助.
\newline
\newline
耶穌在世上時,祂講道說:「神的國近了,你們當悔改信福音.」(可一15)「悔改」意思是從你犯罪的人生,轉向不再犯罪的人生；懷著破碎痛悔的心靈,從今以後遇著聖潔的人生.如果曾得罪人,應該及時與人修好關係.
\newline
\newline
有個人在鐵路局工作,未信主之前,偷了器材；信主之後,奉還所偷的.當人回轉,心裡慾望就完全改變!可能以前所為,是不知不覺的行動；然而今後,你唯一的心願是跟從主.但是,假如你心裡有所感動,知錯即改過而跟蹤主；絕非因悔改救了你,而是因肯悔改接受了神救恩的禮物.我們完全相信,配得的應歸還基督；才可以得到心靈的拯救.要拯救你,是救恩的基礎,並非是我們好的心意,也不是把信心置於教會或善良者的身上.信心必須在耶穌基督身上,因看耶穌基督的緣故；我們的神對罪人,有憐憫的心腸.當我兩個女兒小時候,我一回家便喜氣洋洋,爭先恐後來見我；我便左右手各抱一個,都安詳地躺在我的手臂上.她們毫不顧慮到會跌下來,她們知道在父親的手上是安全的.我們來到耶穌面前也是如此,祂永不放鬆祂的手,也不拒絕我們.耶穌十分愛你,你可安息在祂穩妥的膀臂上.你要向耶穌擺上信心作出決定,講員不能為你決定父母也不能為你決定；最多可以幫助你多明白真理,教會可以栽培你,但也不能為你作出決定.天使在你的四圍,但天使不能拯救你；神願你得救,神呼喚你,希望你現在信祂而得救；但神仍然讓你自己決定.
\newline
\newline
多數人都希望生活方式不變,甚至有人說:「假如你向我提出新的生活方式；我若取其所長而仍保持舊的生活方式,那就最理想.」
\newline
\newline
耶穌寶血在那裡傾倒,你必須決定你的忠心何在.神在宇宙劃出一條血路,在古舊的十字架上,耶穌基督的寶血直流到今日；在宇宙中發揮劃分成為兩部份,使我們每個耳朵能聽到福音.神為著祂所創造又是所愛的人類,將最好的禮物都給我們了.你的眼睛看見十字架時,是耶穌基督的寶血為你傾流.在神無盡的愛中,你必須作出決定,你不能再猶豫,你不能假裝若無其事,不能依照你古舊的生活渡過一生.耶穌基督的血為你傾流,你應作決定是否用愛回應祂.親愛的朋友!這是世界各人必須面對的決定,與靈魂永遠的歸宿和靈魂得救,有直接的關係.
\newline
\newline


\newpage

\section{第56屆~港九培靈會~高文~第七講　審判的確實}
\label{sec:412}
\textbf{第七講　審判的確實}
\newline
\newline
連結: \href{https://www.hkbibleconference.org/session-message/view/412}{\texttt{https://www.hkbibleconference.org/session-message/view/412}}
\newline
\newline
\hyperref[sec:411]{< < < PREV SERMON < < <}
~
\hyperref[sec:toc]{[返目錄]}
~
\hyperref[sec:413]{> > > NEXT SERMON > > >}
\newline
\newline
啟示錄 11:15-11:18
\newline
\begin{longtable}{cl}
\hline
\hline
章節 & 經文 (和合本修訂版)\\
\hline
11:15 & \begin{tabularx}{0.7\textwidth}{X} 第七位天使吹號，天上就有大聲音說：「世上的國已成了我們的主和他所立的基督的國了。他要作王直到永永遠遠！」 \end{tabularx} \\ \\ \relax
11:16 & \begin{tabularx}{0.7\textwidth}{X} 在神面前，坐在自己座位上的二十四位長老都俯伏在地上敬拜神， \end{tabularx} \\ \\ \relax
11:17 & \begin{tabularx}{0.7\textwidth}{X} 說：「今在昔在的主—全能的神啊，我們感謝你！因你執掌大權作王了。 \end{tabularx} \\ \\ \relax
11:18 & \begin{tabularx}{0.7\textwidth}{X} 外邦發怒，你的憤怒臨到了。審判死人的時候也到了；你的僕人眾先知、眾聖徒及敬畏你名的人，連大帶小得賞賜的時候到了；你毀滅那些毀滅大地者的時候也到了。」 \end{tabularx} \\ \\
[1ex]
\hline
\hline
\end{longtable}
第八章開頭,當第七印揭開時,即發生令人驚訝的事；於是天上寂靜約有二刻,那時情形十分恐怖,以致歌聲停止.天上靜寂,預備了人來迎接末次天上的景象.使人想到的日子來到之時的可怕.寂靜之間,有七位天使站在神面前,各持號筒；他們吹號之前,另有一個天使拿著金香爐,把眾聖徒向神不住的禱告都放於金香爐中；帶到神最後要施審判的台前.天使把金香爐盛滿火,把金香爐倒在地上；說明聖徒的祈禱在某種情形之下與世界的終局有關.接著有雷轟,大聲,閃電,地震,表示神在地上將有巨大的行動.前四位天使吹號時就發出恐怖的事,一切災難使宇宙之間各物,甚至日,月,星辰,都被摧殘三分之一.這不過是苦難的開始,其後還有更惡劣的情況.第九章記載,當第五,第六號筒吹響時,有更恐怖的事發生,從撒旦而來的痛苦,紛紛臨到世上居民.到了第十章,另有一個大力的天使,從天降下；手拿小書卷,這小書卷記載聖徒在末後世代將承擔的痛苦.這位天使用鄭重的聲音,宣告最後勝過一切邪惡的日子快要來臨!神的奧秘即將揭示,神的教會要進入榮耀的境界.這使我們的心感到何等甜蜜!但在神最後得勝之前,撒旦將盡牠所能帶來許多苦痛,真令人覺得恐怖!
\newline
\newline
第十一章提及將來要出現兩位見證人,就是舊約時代的摩西和以利亞.他們在末世時要忠心為神作見證,直到被敵對的勢力(魔鬼)殺害.這兩位先知死了之後,屍首沒有埋葬；但忽然之間,他們的屍體被接到天上.於是世上的逼迫就此停止.
\newline
\newline
在最後號筒吹響之前,再次見到神的寶座；天上大聲宣告,神直到永遠的統治.他們以極度歡欣的聲音,呼喊「世上的國,成了我主和主基督的國,祂要作王,直到永永遠遠.」(十一15)這次的爭戰,不再是世界上政治勢力的爭戰；乃是世上的國與國之間,要向誰效忠的爭戰.各位記得嗎?撒旦引誘耶穌向牠下拜；要將地上一切權柄都賜給耶穌.魔鬼巴不得全世界的人向牠敬拜.世界似乎一直受撒旦的勢力掌管,但千真萬確,我們的神才是真正的掌權者；到了我們的主日子來到之時,祂得到完全的勝利!勝過世界上從前魔鬼充滿詭計壟斷人那種國度.主基督宣告「天上,地下,所有權柄都已經歸給我了.」我們讀16節,神得勝的真理再次活現；24位長老再向坐寶座者俯伏敬拜.天上再次唱起聖歌,代表眾子民得救贖的長老,情不自禁又俯伏敬拜神.他們唱感恩的詩歌:「昔在,今在的主神全能者阿!我們感謝你；因你執掌大權作王了.」(17)敵對的勢力正要慶祝時,他們一而再,再而三地稱頌神偉大的聖名,這是最合宜的描寫.他們稱呼表明神是昔在,今在,永遠不改變的.永遠都是現代化的,悠久以來好像是神奧秘的統治,今被揭示了；歷代以來似乎是隱藏的,如今顯明了；祂執掌大權作王,邪惡勢力從此永遠下台了.
\newline
\newline
廿四位長老也加入寶座四圍群眾,唱出詩篇第二篇「列邦發怒,你的忿怒也臨到了；審判死人的時候也到了.」這裡所說的審判,繼續在啟示錄揭示.到了第廿章,就是審判的高峰,所有曾在世界上生存過的；這時要受審判,已死的人身體和靈魂重新聯合；從墳墓,山上,水裡,戰場,早已被遺忘的,如今重新活現.離世的人要回來,連大帶小一同站在神的寶座前,書卷和生命冊展開；按著他們生前所作所為的記錄作審判的根據.對於不信的人,和信了主又背道的人；這時是受審定罪的時候,凡名字沒記載在羔羊生命冊上的；要被扔到永遠的火湖裡,這是第二次的死.我們如果繼續以下幾章,就明白為何列邦發怒.那些忘記了神的人要逃避審判,他們巴不得山嶺倒在自己身上來逃避坐寶座者的忿怒；但是神結算的日子,人無法逃避；今生藐視神的人到那時要受神公義的審判.聖經繼續告訴我們,這樣的人靈魂要被扔到火湖裡去.那裡是為撒旦和牠的使者所預備的.到那時,一切的憐憫都停止了!所有失喪者被棄於黑暗中,孤單受苦直到永永遠遠.這種審判是神所頒佈的,並非受造之物把審判置於其他人身上.聖靈帶領人所作的禱告,要發展直到審判之日來到；貫徹公義乃是由坐寶座者執行.審判非但臨到不敬畏的人身上,第18節叫我們注意審判臨到聖徒身上.信主的人接受審判,決定他們得救與否.這班聖徒受審判,神將聖徒心裡隱藏的動機揭示；得到賞賜的人,同時也得到神的稱讚.聖經說在神燒著的火焰中,決定我們是否得賞賜.「火」顯明每個聖徒工作的性質,沒有摻入雜質,是為討神喜悅的事奉；到了審判之日都要顯明；如同煉淨的金子,要得到賞賜.相反地,為了榮耀自己的事奉,到那日都要被燒盡.聖經說這樣的人雖然得救,但卻如同從火中經過一樣.救恩並非靠自己的功勞取得,得救乃因相信耶穌基督,是祂完成的救贖工作.但是,無論如何,到了人白白接受了祂的救贖之後；就應該開步走為神作工且應當存著愛祂的心而作工；叫人看見我們的好行為,便將榮耀歸於天上的父.到了那日,非但天父得著該得的榮耀稱讚,連你也得了稱讚.第四章提及廿四位長老戴金冠冕,就是得賞賜的意思；在更深意義中,是聖城新耶路撒冷由天而降,是信了主服事祂的聖徒才有份.這就是啟示錄最後幾章所提及聖徒的賞賜.惟有一心為榮耀神的聖徒,到那日將得到這樣的賞賜；今生無論遭遇何等的苦楚,若想到將來要得無比的榮耀,便不足介意了.(羅八18)每個人的工作到那日都在神前陳明,這更加強我們在世時為主忠心的志向,我們必須為肉身所作一切的事到那日向神交賬.主說「多給誰就向誰多取.」
\newline
\newline
第19節提及廿四位長老獻上的詩歌,他們唱完詩歌,天上的殿就開了；神的約櫃在天上聖殿中出現,再次向聖徒證明這位應許我們的神,是信實的神.
\newline
\newline
或者你覺得奇怪,為何那班長老要歌頌那「審判的日子?」因為到了那日,偏行己路的人將永遠與聖徒分開；神不喜看見惡人遭報.長老們所以歡欣,並非幸災樂禍有人受刑罰；乃是為了神最終的勝利而歡欣!這個審判顯明我們的神無窮的權柄,證實了真理；證實了神兒子耶穌基督的榮耀,不停的歌聲顯明天上無限的歡樂.
\newline
\newline
茲題出三點真理:
\newline
\newline
(一)神確有恩典的約
\newline
\newline
恩典的另一面是神的忿怒,祝福與咒詛在同一個寶座流露.在今日道德淪落的社會,人對神的公義有所懷疑,人常忘記神的恩典與忿怒,是從一個寶座流露出來的真理.這真理教導我們關於神的個性,我們若褻瀆神,總有一日將遭遇報應,刑罰.公義要求天地之主,持守祂公義的責備；否則,神的道德,神的能力,神的永恆,豈非倒淨了嗎?為了對罪惡必須有刑罰這個法則,神寧可將祂獨生的兒子,被掛在十字架上,來擔當我們眾人的罪孽；因為神有正直剛強的公義,祂的公義更加顯明祂的大愛.神十分愛我們,不忍看見我們繼續行走悖逆的道路,以致錯失祂創造人的原意.故此,神把我們反叛的始祖,逐出伊甸園又將死亡帶給人類.這是神的憐憫,免得人類敗壞的情況繼續惡化；祂叫人類尋找祂的國度,這國度是罪惡死亡無法勝過的.只要我們曉悟,明白神的審判都是為我們的好處而發出.啟示錄第三章,神對老底嘉教會,發出審判的呼聲.老底嘉教會非常富有,財富不斷增加；他們自以為一無所缺,表面上好像很敬虔；實際上,聖經說他們裡面是困苦,可憐,貧窮,瞎眼,赤身的.他們只有宗教的儀式,而無屬靈的能力.他們的生命並無加強敬拜神的傾向,他們如同溫水,不冷又不熱.基督徒這種隨隨便便向神的方式,是神所厭惡的；神要從口中把他們吐出來.在聖經中,令神厭惡以至於從神口中嘔吐出來,這是最嚴厲的字句.當神在恐怖的審判發出的同時,也公佈祂的恩典和慈愛；祂叫老底嘉教會買眼葯擦眼睛,才能看見自己可憐的境況.神說:「我所疼愛的,我就責備管教他；所以你要發熱心,也要悔改.」(三19)接著,聖經出現最有恩慈的邀請:「看哪!我站在門外叩門；若有人聽見我聲音就開門的,我要進到他那裡去；我與他,他與我一同坐席.」(20)接著,神又發出令人心服的賜福應許:「得勝的,我要賜他在我寶座上與我同坐.......」(21)
\newline
\newline
各位!今晚你聽見主在門外叩門嗎?主十分愛你,正在你心門之外叩門；祂希望你有更好屬靈的生命,所以才處罰你.
\newline
\newline
(二)神把真理向我們揭示
\newline
\newline
每個人面對這真理,要負上責任向祂交賬；你越多機會領受神的道,你在主前就有越重的責任.那些從未聞耶穌福音道理的人,他們所有的,盡是社會供給他們的道理和傳統.但他們卻無法達到標準,故此他們只有到淪亡之地,而接受神的審判.猶太人已得到舊約聖經的教導,舊約聖經不但把律法和其他條例教訓以色列人,並向他們指明日後要來的基督；但他們卻抗拒不接受,以致要接受神最嚴厲的審判.有的人聽了福音拒絕接受,而且用腳踐踏神兒子的血.不信者要受審判,我們已信者也同樣要接受神的審判；我們已聽過神的道,理應有能力在神性格上繼續長進.我們學習祂,經歷祂復活的能力,在我們身上發揮出來,我們要認識祂；聖靈就越發感動提醒我們的軟弱,敗壞,過錯.
\newline
\newline
有一位機師駕駛飛機,經過海洋時忽聞機艙有特別聲音；好像機上一重要電線被咬,一旦咬斷,飛機將失事墮海.他細聽之下是老鼠所為,立即拉開機掣,使飛機升高至老鼠無法生存的高度.各位!你屬靈的高度如何?神愛你的屬靈生命,神希望你靈命更高；但未達到神的心意,神提醒你!你心靈的罪雖很微小；好像老鼠侵襲你的靈魂,阻礙你更高的認識神.
\newline
\newline
(三)審判必從神的家開始
\newline
\newline
如果你全心愛神,當神感動你知罪,你就定自己的罪；可能我們仍會跌倒,有軟弱；但我們誠心純潔向著神.基督徒的良心,應該是不得罪神,不得罪人；惟有離開一切污穢渣滓,我們生命的器皿才能被神充滿.
\newline
\newline
49年,嚇百米海島的人懇求神的復興,數週之久呼求神賜復興.有一晚,青年人在聚會中正呼求之時,讀到詩篇「誰能登耶和華的山?誰能站在聖潔的神面前?」聖經的話答:「就是手潔心清,心裡向神純潔的人.」他流淚四望會眾,說:「我們費了多時在神前求復興；但我們的心仍是污穢不清,豈非白費時光嗎?」他承認自己有罪,其他的人也有罪；所以他們一一把自己心裡的痛苦和需要說出,經過整夜直到黎明之時,這班人才真正得到復興.他們的心得到潔淨,神使用他們成為器皿,藉復興把神的福份帶給海島上成千上萬的人.
\newline
\newline
70年,我們在美國阿斯百里大學神學院也經驗如此的大復興.有一天早晨,大學舉行早會,一個高年級同學要求給機會講見證,他承認三年來過著假冒的生活；表面上他有屬靈生命,說的都是屬靈的話；但卻非心裡真實的情況.他面向同學們求寬恕,他說昨晚已在神面前與神修好關係.當他見證完畢,一個又一個起來見證；真是看見神的靈,在他們中間作工,會場中站滿準備作見證的同學；結束早會的鐘聲響了,同學們仍不離開禮堂,繼續接二連三作見證,直到午飯時間大家仍不散會.結果,學校停課,連原定的籃球比賽也取消,全校一切活動都停頓；每個人唯一要務就是與神和好,整夜不停又繼續到次日.四圍的人聞風而至,八晝夜不停地繼續,任何事物都不能攔阻此次大復興之潮.弟兄姊妹們都向神傾心吐意,有位教授也起立承認,六年來靈性倒退冷淡,他請同學來到祭壇前為老師禱告.有位牧師列出許多同學的名字說出與他們有個人緊張的關係,他的妻子也起立說自己作為師母,一向心中不快樂；甚至痛恨城市及鄉村的居民,以往對人漠不關心,她請求大家寬恕.
\newline
\newline
各位!挺身而出在神面前坦白,有時真是令人驚奇!經過復興的人,才能明白何時能得復興；就是不再認他人的罪,而是承認自己的罪.你今晚有沒有罪向神承認呢?我並非問你那些不自知的罪,而是在你生命裡有沒有你自己清楚知道的罪?神一再提醒你「這是你的破口」,神心愛你；神的聖靈希望你符合神與主耶穌聖潔的模式.當神的聖靈向你說話時,我請求你實行神的旨意在你身上,作個聖潔的器皿,合乎主用!讓神的聖靈透過你流露出去,在千萬等候生命答案的人中,你能夠成為他們的祝福.巴不得神所賜的福臨到我們的教會,但是,要從你自己開始呢!
\newline
\newline


\newpage

\section{第56屆~港九培靈會~高文~第八講　得勝者的凱歌}
\label{sec:413}
\textbf{第八講　得勝者的凱歌}
\newline
\newline
連結: \href{https://www.hkbibleconference.org/session-message/view/413}{\texttt{https://www.hkbibleconference.org/session-message/view/413}}
\newline
\newline
\hyperref[sec:412]{< < < PREV SERMON < < <}
~
\hyperref[sec:toc]{[返目錄]}
~
\hyperref[sec:415]{> > > NEXT SERMON > > >}
\newline
\newline
啟示錄 12:10-12:12
\newline
\begin{longtable}{cl}
\hline
\hline
章節 & 經文 (和合本修訂版)\\
\hline
12:10 & \begin{tabularx}{0.7\textwidth}{X} 我聽見在天上有大聲音說：「我神的救恩、能力、國度，和他所立的基督的權柄現在都來到了。因為那個在我們神面前、晝夜控告我們弟兄的，已經被摔下去了。 \end{tabularx} \\ \\ \relax
12:11 & \begin{tabularx}{0.7\textwidth}{X} 弟兄勝過那條龍是因羔羊的血，和因自己所見證的道。雖然至於死，他們也不惜自己的性命。 \end{tabularx} \\ \\ \relax
12:12 & \begin{tabularx}{0.7\textwidth}{X} 所以，諸天和住在其中的，你們都快樂吧！只是地和海有禍了！因為魔鬼知道自己的時候不多，就氣憤憤地下到你們那裡去了。」 \end{tabularx} \\ \\
[1ex]
\hline
\hline
\end{longtable}
美國初期的歷史,基督徒的夏令會,有一首詩歌是慶祝打敗了撒旦的凱旋歌.長達十八節,副歌說:「呼喊,呼喊,我們呼嘁!逐步得勝,我們呼喊,魔鬼就要哭,哈利路亞!」初期的基督徒可能較為豪放坦誠.
\newline
\newline
啟示錄將最後的得勝,向我們描述之前,(10-12)這段聖經告訴我們聖徒如何受患難,在患難中繼續得勝.第一世紀在該撒之下受逼迫的使徒,只不過反映日後繼續遭受逼迫,但至終得勝.
\newline
\newline
天上現出大異象　第十二章開頭記載有個懷孕的婦人即將生產孩子,她象徵神的選民以色列的母親,她辛苦呼喊.第二個象徵,在正要生產的婦人面前站著一條龍,代表魔鬼.魔鬼正向神表示忿怒和反抗,在婦人面前等待要吞吃她將生產的孩子.耶穌降生時,豈不是有惡勢力要殺害祂嗎?表面上似乎是人的工作,事實上乃魔鬼操縱的陰謀.
\newline
\newline
魔鬼是迷惑普天下的(9)那婦人生下的男孩子,就是彌賽亞,神命定祂統治宇宙.我們的主受死,復活,升天坐在天上父的右邊；主的得勝就顯明了.那代表蒙救贖的彌賽亞之道,要逃到曠野去,在那裡受安慰.當神的教會遭受逼迫,在那裡得到保護；世上的居民常為逃避他們的政府,但神的慈愛卻看顧屬祂的人時刻與他們同在.聖經說天上發生戰爭,魔鬼意圖在天上,神的寶座上,佔據統治的地位；這場爭戰,撒旦和牠的使者都被扔到地下去了.這個事實反映耶穌說的話:「我看見撒旦和牠的使者都墜落.」無論此事發生於何時,都足見全能者能力的彰顯；從此之後,撒旦再無計可施了.這更反映出神的教會的得勝.魔鬼既無法侵害那初生的男孩,便千方百計向神的子民進攻.魔鬼攻擊的對象是那班守神誡命,為耶穌作見證的人.(十二17)今日仍有許多敵對的勢力,向教會攻擊；如果我們不留心魔鬼的詭計,就將成為信徒的致命傷.聖經明說:基督徒的爭戰,不是屬血氣的爭戰；我們非但在肉身上爭戰,我們當向屬靈界的掌權者邪惡的勢力爭戰.因為魔鬼控制一切從地獄而來的勢力；向著神的教會和神的工作摧殘；自從人類歷史開始以來,而且還要繼續下去,我們向著龍──魔鬼爭戰.
\newline
\newline
靠主得勝　天上的聲音提醒我們,神的教會仍然是得勝的教會.撒旦是個打敗仗的敵人,在天上明亮的光中；得勝的呼聲在天上,也在地上你我生活之中.天上有大聲音說:「我神的救恩,能力,國度,並祂基督的權柄,現在都來到了；因為那在我們神面前晝夜控告我們弟兄的,已經被摔下去了.」(十二10)神戰勝了撒旦,顯明了基督徒過著得勝的生活.神已經把拯救賜給我們了,基督國度的統治已被建立.藉著耶穌的死而復活,我們蒙神拯救,蒙神膏立的人,能夠在榮耀中過生活.所以今天你不至於在失敗之中,在魔鬼恐嚇之下過生活,我們不再懼怕魔鬼的控告和恐嚇；因為那在基督耶穌裡的,就不再定罪了.這並非聖徒得到所謂消極的自由,乃是「弟兄勝過牠,是因羔羊的血,和自己所見證的道；他們雖至於死,也不愛惜性命.」(11)
\newline
\newline
第十二章11節用三言兩語說明了我們基督徒得勝的原因:
\newline
\newline
(一)羔羊的血
\newline
\newline
在天上快樂的詩歌中,一再以羔羊的血為題.血是得勝信息最主要的因素,藉著耶穌基督甘心流出的寶血；撒旦罪惡一切枷鎖都失效,藉著耶穌的寶血,我們免為魔鬼的俘擄了.
\newline
\newline
有一故事講財主想進天堂,他站在天堂門口；天使問他:「進天堂的暗號是甚麼?答了問題才可進入.」財主說:「我曾為教會多多捐款,依我的道德標準,眾所公認是毫無疑問的,當然我能夠佔個天上的位置.」天使說:「天堂的暗號不是這個.」財主只好失望地走了.接著又來個一表人才的,天使照例問他.這人說:「我服事神的教會,我是牧師,我奉神之命作了很多公義的事；我應得天上給我的榮耀.」天使說:「進天堂的暗號不是這個.」他只好走開.繼之,來了一位老婆婆,她彎彎的背顯示她多年的勞苦；她雙眼閃爍,面有榮光.天使一問她暗語,她立即舉手唱詩:「耶穌的寶血,哈利路亞,把我潔淨的血.」珍珠門立刻打開了,這個親愛的靈魂進入這美麗的城時；天上的詩班,都加入她剛才所唱「血」的詩歌.這古老的故事可能太過淺陋,但其中教訓我們應當牢記:欲進天堂,惟靠耶穌基督的寶血.
\newline
\newline
(二)自己所見證的道
\newline
\newline
他們見證的道,說出神的信息,如何向迷失的世界傳達的途徑,他們首先的責任要宣講福音.救贖的真理包含將真理傳遞出去.你毋須作個受過訓練的神學家或有恩賜的講員；你只要告訴人,自己所經歷的,所看見的和所聽見的.這並非信條,而是基督徒與主同活,自然活出的生命.據我所知,有個商人多年來用他的商店,作為耶穌福音的地方,四圍的人都到過他的雜貨店,聽他講耶穌的救恩.有時看見老年姊妹派發單張,找機會傳講耶穌；就算講幾句話也是好的.有一次,一間大學的女生,搭車時一路看聖經,旁坐的女搭客說,從來沒有人這樣在巴士上看聖經的；於是她就利用機會與旁坐者講見證.各位!這裡有個功課大家必須學習；她說:「我們的生活方式與世俗人不同,或者他們要覺得希奇為何有不同之處呢?......」她繼續講下去,旁坐者聽完了,說:「我從未聽過有人講的像你這麼動聽的,我還有個朋友坐在第二行；你所講的如果她聽見是很合適的,我可以和他換個位置,讓他坐在你旁邊聽.」這位大學生多了個機會為耶穌講見證.還未講完,前面有位男士很感興趣地說:「剛才你說甚麼,最好能講給我聽.」正講時,又有一人很好奇地說:「請你講大聲一點,讓大家可聽得到.」巴士到了總站,司機轉過頭說:「你再講多一些,我也聽.」女生心裡暗自「哈利路亞.」各位!當你心裡有負擔時,這個失喪的世界,真是願意聽到慈愛救主的福音.教會的得勝,和基督徒所做的事有特別關係.彼得對耶穌說:「主啊!你是永生神的兒子.」耶穌回頭向彼得說:「你是彼得,我要把我的教會,建造在這磐石上；陰間的權勢不能勝過祂.」(太十六18)耶穌基督的勝利是何等偉大,是有計劃的勝利.當耶穌的跟從者高舉祂,祂就呼召更多人跟從祂時,撒旦要敗退.當你決心要傳揚時,傳一個又一個,把主的福音向世界所有的人傳揚；只要基督徒用乘數表傳遞福音的方法,沒有一樣事物能攔阻我們的主耶穌；使魔鬼地獄之門被震動,魔鬼就要失敗.
\newline
\newline
(三)他們雖至於死,也不愛惜自己性命
\newline
\newline
他們甘心為基督的寶血作見證,對主忠心的程度極高,不顧一切,甚至死也不足惜；天天順服主,背起十字架；到了真正委身的時候,他就出去歌唱.
\newline
\newline
美國有一個教會,門上刻了些字「當稱謝進入祂的門,當讚美進入祂的院；當感謝祂,稱頌祂的名.」(詩一百4)在教會引述這段聖經最為合適.各位!當你進入聖殿,你的心境是否帶著唱詩歡樂呢?因為我們是祂草場的羊；所以當稱謝進入祂的門,當讚美進入祂的院.牧人把羊從草場帶到耶路撒冷進入聖殿,惟一原因進到聖殿被殺獻於神的祭壇上.這是指我們說的,我們是神草場上的羊.讓我們明白,我們得著耶穌基督的福音；我們用口唇見證主,同時我們委身給主,甚至於死都願意!我們是耶穌用祂的寶血買贖回來的,不再屬於自己了.我們將來有盼望,乃因耶穌愛我們,買贖我們.故此,聖經一再提醒信徒,靠自己是死的；向世界死了,向神卻應是活著.使徒保羅說:「我們活著,是為主而活,若死了,是為主而死；所以我們或活或死總是主的人.」(羅十四8)各位!如果我們這樣獻上自己,魔鬼便無計可施了.無論任何環境,只如此委身,就能從世界罪惡的捆鎖中得釋放.
\newline
\newline
數年前,非洲烏干達有次叛亂,時至今日這個國家仍常有叛亂.那次的叛軍到了福音站強迫牧師和校長退出來,其實他們不願離開；但在無可選擇之餘地,於是被迫上了一輛貨車開走.不多時,到了河畔,車就停在橋邊.叛軍在那裡打盹.牧師在危急關頭,仍記得問那位校長信主的程度如何?校長說:「我真正相信耶穌了.」牧師在自己的日記簿上寫給他的太太說明時間,並說:「我並不是這樣死了.」叛軍決定採取行動,命令牧師跳出貨車上橋,他只好聽從,當他行盡人生最後一程之時,他忽高聲唱:「有一地比日中更光彩,雖遙遠我因信望得見,我天父在那地常等待,有預備極安樂的房間.到日期,我樂意,與眾聖同聚會在美地；到日期,我樂意,與眾聖同聚會在美地.」雖然他在橋上被叛軍槍殺了,他的屍體被踢下河中；但楊能牧師永遠唱著這首詩歌.那位校長心想自己是第二個被害者；可是他覺得奇怪,叛軍仍在打盹；所以他就駕車逃生,他說牧師面臨死亡邊緣；還不停的唱詩；連叛軍也覺希奇.
\newline
\newline
未知我們中間哪位願意為主獻上肉身的性命,在我們心靈裡,應隨時作好準備,把身體當作活祭獻上；在耶穌基督先為我們獻上生命的亮光中,我們獻上是理所當然的.惟有如此向主獻身,我們可以加天上群眾行列同唱天上的詩歌.將來如何,現在還未顯明,但當主顯明的時候,我們要像我們的主；那時我們要親眼看見主的本像；那日萬膝要跪拜在祂的面前,萬口承認耶穌是主,歸榮耀給天父.
\newline
\newline
天上的群眾在寶座四圍高聲歡呼歌唱:「耶穌基督的救恩,權柄,能力,國度都來臨了,控告我們弟兄的惡者被摔下了.」從今以後,毋須懼怕!毋須受苦；不再受邪惡的世界恐嚇,不再被困於飄盪不定痛苦之中；因為耶穌基督已經得到最後的勝利.天上所慶祝的,就是因為這個實際的情況.你若不是活在耶穌基督的得勝裡面,你就放棄了你的權利!
\newline
\newline
親愛的弟兄姊妹!你已往曾失敗過,受世界壓力的侵害,你曾失去耶穌給你的喜樂.有的人盡所能妥協,一隻腳站在天上,一隻腳站在世界；所以你的心充滿苦惱失敗.難道你不相信神給你有美好的預備嗎?難道你不加入得勝教會的勝利歡呼嗎?聖經告訴我們,弟兄勝過祂,是因得勝羔羊的血,靠著我們所見證的道；因為我們甚至於死,也不愛惜自己的性命.你是否有此經驗?若是未有,今晚請你進入這個經驗的實際裡面,請你當機立斷,不再猶豫.你既曾經承認主耶穌的名,你今天晚上向我們的主完全投降；讓祂作你生命的主,你向祂承認曾經偏行己路,三心兩意,不冷不熱；現在決意再來跟從祂,今晚完成你的心願.進一步把自己生命的主權交給主,向主投降.
\newline
\newline


\newpage

\section{第56屆~港九培靈會~高文~第九講　萬民的聚集}
\label{sec:415}
\textbf{第九講　萬民的聚集}
\newline
\newline
連結: \href{https://www.hkbibleconference.org/session-message/view/415}{\texttt{https://www.hkbibleconference.org/session-message/view/415}}
\newline
\newline
\hyperref[sec:413]{< < < PREV SERMON < < <}
~
\hyperref[sec:toc]{[返目錄]}
~
\hyperref[sec:416]{> > > NEXT SERMON > > >}
\newline
\newline
啟示錄 15:2-15:4
\newline
\begin{longtable}{cl}
\hline
\hline
章節 & 經文 (和合本修訂版)\\
\hline
15:2 & \begin{tabularx}{0.7\textwidth}{X} 我看見彷彿有攙雜火的玻璃海；又看見那些勝了那獸和獸像，以及牠名字的數字的人，都站在玻璃海上，拿著神的豎琴。 \end{tabularx} \\ \\ \relax
15:3 & \begin{tabularx}{0.7\textwidth}{X} 他們唱神僕人摩西的歌和羔羊的歌，說：「主—全能的神啊，你的作為又偉大又奇妙！萬國之王啊，你的道路又公義又真實！ \end{tabularx} \\ \\ \relax
15:4 & \begin{tabularx}{0.7\textwidth}{X} 主啊，誰敢不敬畏你，不把榮耀歸於你的名？因為只有你是神聖的。萬民都要來，在你面前敬拜，因你公義的作為已經彰顯了。」 \end{tabularx} \\ \\
[1ex]
\hline
\hline
\end{longtable}
我們來到聖殿,能夠聽見:「大哉,聖哉,耶穌的名.」我們的心就得著鼓勵,讀啟示錄天上的詩歌時的感受正是如此.天上的羔羊得著頌讚,造成了地上事物演變的背景.
\newline
\newline
魔鬼千方百計侵害屬神的人　魔鬼明知至終要打敗仗,所以在末時代盡牠所能引誘和逼害信徒.第十三章記載,有個獸從海中上來.獸代表邪惡的勢力,不斷褻瀆神並向信徒侵害；接著,另一個獸起來,這獸代表邪惡的國家,牠的樣子,牠的角如同羔羊；是假冒的先知,牠來之目的要世上人拜第一個獸.牠很成功的,使世界上許多人,頭上戴著666獸的記號；凡不聽從的,就要被逼害以至殉道.
\newline
\newline
這件事使我們覺得極其恐怖,所以啟示錄第十四章,再次把安慰鼓勵的信息帶給我們.
\newline
\newline
唱天上的新歌　寶座的羔羊接受歌頌讚美,那班得蒙救贖的人,再次唱新歌.第十四章卻沒有記載新歌的詞句.天使逐一顯現作見證,他們的聲音,乃要呼召人敬拜創造統管萬有的真神；預言世上不敬虔的國家要衰落,(這類國家以巴比倫為代表)並且說:凡頭上有獸印記的人要遭受神的忿怒；凡對主耶穌忠心的人,為主殉道在主裡面死的人有福了.在這段經文的末了,使徒約翰看見神的兒子戴著冠冕,在世界做收割的工作；又描述一幅圖畫,有天使從殿中出來,手持鐮刀收割,又把鐮刀扔在地上；地上的莊稼就被收割了,於是收取了地上的葡萄,丟在神忿怒的大酒醡中.
\newline
\newline
第十五章開始,天上又有異象,掌管七災的七個天使一一出現.
\newline
\newline
唱摩西和羔羊的歌　天上又有讚美的歌聲,有許多人齊集在玻璃海上,他們拿著神的琴,大聲歌頌神；這些人是向神忠心,不肯向世界上的強權背叛神者低頭下拜的.啟示錄曾題及有一班保持衣裳潔白,不向獸跪拜者,就是這裡所指向神歌頌的人；聖經形容他們是初熟的果子,是純潔的,他們口不出謊言,頭腦從不思想情慾.他們稱頌神救贖的偉大作為,他們唱的是舊約神的僕人和新約羔羊之歌合成一片,律法與福音共冶一爐.摩西之歌,歌頌神把以色列民從埃及為奴之家拯救出來.當時法老的追兵到了紅海,結果全軍覆沒,屍體遍滿紅海的水中和岸上.在此敗落情形之下,出埃及記載,以色列民向神獻上摩西讚美之歌:「耶和華是我的力量,又成為我的詩歌,也成了我的拯救.你是我的神,我要稱頌你.神啊!誰能公義聖潔像你呢?我們存著敬畏的心來讚美你,你是獨行奇事者；因你的慈愛和信實,你把你的救贖的子民帶領出來,你以能力帶領他們進入你的居所.我們的耶和華從此直到永遠要統管一切.」以色列人蒙神帶領離開埃及.摩西之歌十分普遍在每個安息日,在會堂,在晚間敬拜時都唱這首詩歌；每個在神的約有份的以色列小孩子,從小就熟悉這首詩歌.它預言將來彌賽亞的拯救,神昔日如何把以色列人從埃及為奴之家帶領出來；如今神要帶領祂的信徒,到永遠安息之所.在新約帶領祂的子民離開罪惡,進入神所應許的天堂美地者,就是羔羊這位主.為記念蒙恩救贖的歷史,天上群眾高聲歌頌讚美說:「全能者啊!你的作為大哉奇哉!萬世之王啊!你的逆途義哉,誠哉!(3)他們所歌頌的是人類永遠的主宰,君王,統治萬有的真神.「主啊!誰敢不敬畏你,不將榮耀與你的名呢?因獨有你是聖的；萬民都要來在你面前敬拜,因你公義的作為,已經顯出來了.」(4)在詩詞中說明惟主是聖潔的.這種敬拜進入彌賽亞時代.「萬民都要來在你面前敬拜.」回應詩人在詩篇的話:「萬民都要來稱頌你,歸榮耀你的尊名.」
\newline
\newline
這首詩歌末段唱出另個敬拜神的原因.
\newline
\newline
神是施行公義者　神是公義的施行者.這班歌唱者唱「因你公義的作為已經顯出來了.」神絕對公義掌權,乃眾目共睹的,聖經說,我們的神在列國面前顯明祂的公義.」地極的人都看見我們的神的拯救.(賽六二2)
\newline
\newline
關於全能者的作為,茲題出三點真理:
\newline
\newline
(一)命定耶穌基督管理,統治世上各國人民,祂是宇宙的王
\newline
\newline
神對亞伯拉罕說:世界萬國,萬族要因他蒙福.先知也曾預言天上神設立的國度永遠不會摧殘,祂的政權延續永不衰敗.這個統治的世代,當人子駕雲降臨時,這國度就要開始,權柄,榮耀,國度都歸給祂.先知預言主的國度,包括所有的國,所有的民族；這些蒙福的人,要在永世裡事奉他們的王.耶穌基督生平的教訓中,實現了舊約的應許,在福音書中超過八十次,耶穌自稱為人子的.舊約先知預言將來臨到的國度,耶穌用人子來說明這應許如今要實現.耶穌自己記得祂第一次來是被殺的羔羊,但第二次再來乃作榮耀的王；祂到處傳道說:「神的國即將來臨了.」主在世時,預知神的兒子要降臨,滿有權柄；所有神的仇敵,都要被踐踏在祂腳下,這種意識時常在祂腦海中活現.
\newline
\newline
願為耶穌基督而活的人,也必須時刻期望神的國度；無論現在我們在世上遭遇何事,我們知道,公義的主要作王.祂在信徒心中作王,靠著祂的恩典,有一日我們也將與主一同榮耀中作王.我們從神的話語有此確據,所以我們毋須為前途擔憂.
\newline
\newline
有一個護士在軍人醫院中問一個正流著血的,士兵說:「你現在向死亡作好準備了嗎?」這位黑皮膚大眼睛的士兵回答說:「預備好了,因為死是為著國度.」他邊說邊用手撫胸,不多時,這位年青的士兵死了,他的手還按著自己的心.當我們以信心接受主之時,我們已實際進入主的國度中了.
\newline
\newline
(二)進了神的國度必須用信心
\newline
\newline
在這世代,我們應當告訴人進入神國度的方法,我們必作好準備；當君王再臨之前,最重要的是宣告福音的信息.耶穌說:「這福音要向萬民作見證,傳遍普天下,然後末日來到.」(太廿四14)人要認識救主的真理,且成為進入神國度的勇士.我們得著了這個信息,有責任把這好消息傳給其他的人.時間是重要的因素.有一位偉人名威可的,他說:「我們到了永恆,有足夠的時間慶祝；但是現在離太陽西下僅有數小時,時日不多,快取人吧!」耶穌為每一個人受死,每個人必須知道福音的訊息；主的寶血為你,為我而流,這血發出呼聲,呼召每個人都來歸順.有時我們傳揚福音表現懶惰的態度,我們當自問是否真正知道福音的必要性.
\newline
\newline
外國有一個惡徒名查理士比,當他將被處決時,獄中有位牧師想安慰他.牧師說完了使人得安慰的福音之後,遭受剌激的惡徒轉向牧師說:「你這樣說,你是否真信呢?你是否照你所說的而信呢?」他惱恨咬牙切齒地說:「如果你這樣說,這樣信,如果我像你信這樣真實；甚至地上有玻璃,我都要爬上去,盡我所能向其他的人宣告福音的好消息.」如果你真心相信耶穌基督為人而死的福音,豈可安坐；白白地看見同胞失喪滅亡?耶穌清楚說,若非靠著祂沒有人能到天上的寶座.世界的人必須聽福音的信息,否則毫無盼望.
\newline
\newline
美國有海外巡邏隊,是支確保安全的隊伍,他們有責任去安慰遭遇失望航海的人.有一次暴風雨之夜,船隻遇難；巡邏隊接到求助的信號,準備出海救援.其中一巡邏隊的學生,因見波浪洶湧,他大聲疾呼:「船長啊!這樣惡劣的情形,我們必定一去不復返.」船長面對這青年說:「你需要明白我們可能不會回來,但我們必須出去；因為這是海外巡邏隊創立之目的.」
\newline
\newline
當耶穌被掛在十字架上時,許多人譏笑他:「他救了別人,卻不能救自己.」(太廿七42)耶穌基督降世之目的,不是來救自己,乃是來救贖世人；祂不是來接受人的服侍,乃是來服侍人；祂傾出生命作為多人的贖價,祂說明要來尋找拯救失喪的人.這福音不是我的福音,乃是屬於祂的福音；因為基督為王,萬膝都要跪拜.在這樣的敬拜裡面,我們應當出去宣揚我們的王的信息.
\newline
\newline
現全球人口達四十五億之多,但與教會連絡的人不過一億多人而已；其中還包括了許多掛名的基督徒.按現人口增長率,到了二○○○年時,世界人口達到65億,照比率到本世紀末,基督徒人數最多不過二億多人,至十幾年後,那班未聞福音者至多相等於現在的人口.非常遺憾!西方國家教會工作太過落後,不能保持人口的增長率.在歐洲更是令人恐怖,教會一直在衰退中；教會增長之處在歐美以外的第三世界,佔三分之二的地區.除非我們竭力加速傳福音,否則將看見世界未信主者人數逾發龐大.
\newline
\newline
耶穌基督向信祂的人發出命令「往普天下去,使人作基督的門徒.」耶穌所用的字眼特別有意思,耶穌基督的門徒要效法祂,要跟從祂；祂給我們的大使命,不是只叫我們出去使人信祂；信才能進入神的國；大使命的真正意義是要去,使萬民作耶穌的門徒,這才是基督策略的精華,要得著整個世界.惟有學習耶穌基督；才能出去使人作祂的門徒.只要你對耶穌基督有認識,使人作門徒正是我們的主在世上時的工作；大使命不過說明耶穌整個人生從事的工作.這一點恐怕被多人誤解了；使人作門徒不需要從神領受特別的呼召,不需要特別的恩賜；耶穌的生活模式正是這樣不斷使人作祂的門徒；我們跟從主,我們要效法祂.我們已經認識了主的人,就當告訴別人來認識跟從我們的主,我們工作的對象是從神容許給我們認識的少數人開始；當你四圍的人從你身上學會了作基督的門徒,他們也出去叫人作基督的門徒.藉著這種繁增又繁增的方式,全世界的人都能認識我們的主.神確是呼召特別的人,使用他們的恩賜專心在教會之中.今日確有個迫切的需要,叫弟兄姊妹全時間在教會從事特別的工作,全時間作牧師作傳道；作教會的教師,作護士,醫生,工程師,宣教士等.今天晚上,神正在呼召你,把自己獻作全時間服事神的人；但是無論你是個平信徒或被教會按立的全時間工作者,兩者都必須作耶穌基督的見證人,把福音帶給世上的人.我們清楚有此目標,無論今生作甚麼,總要記得使人作耶穌的門徒.
\newline
\newline
各位!現時正是收割的時代,無論充當任何角色,讓我們一同去作福音收割的工作吧!
\newline
\newline
(三)無論主放我們於任何崗位上,我們確信神要使用我們去完成祂的大使命
\newline
\newline
當你面對審判的時候,不可使主嫌我們不忠.主在今日仍要差遣祂的僕人出去,使人作祂的門徒.假如我們不加入神的工作的隊伍,我們就是向著歷史,作反方向的事.聖經明載,歷史的終結,那日萬膝要跪拜,萬口要承認耶穌是主.在啟示錄「耶穌是主」這個主題,一而再,再而三地被慶祝,各族,各方,各民,各種方言的人都聚集；人多不勝數,他們高呼救恩稱頌神.無論天上,地上所有受造之物,都響應敬拜.從寶座的角度看下來,耶穌基督的大使命那時已經實現了,屆時天上已經開始慶祝.當我們看見世界發生的事,我們的眼不禁仰望天上；故此,我們心有詩歌,不但現在唱,將來到永恆裡還要不斷地唱；我們確實知道,我們站在勝利的一邊.
\newline
\newline
早期有位司各牧師講他在印度的經歷,有一次他忽然被一班軍士包圍,手持矛槍瞄準他的心臟；他軟弱無能,用土語唱出聖詩「大哉,聖哉耶穌的聖名,眾軍一同歌唱,準備接受土人的矛槍；他唱了一節又一節,到了唱第三節,他睜開眼睛,發現矛槍一概落在地上.土人流著眼淚,因為他們聽見詩歌中有個超越萬人之上的名字,土人要求他說明這個尊名；然後他同土人一同回到村莊,多年在那裡勞動建立教會.各位!以上事蹟使我想起我們的主要建立國度；我們傳揚福音的人,未必個個都從矛槍中得著神的拯救；無論何事發生,我們確信耶穌基督的名,直到地極都要得勝被稱揚!確有一日要接受冠冕,成為萬王之王.這時,天上的群眾唱詩讚美我們的主,說:「這位主就是萬王之王.」世上萬族都在祂面前敬拜；另一面的聖徒也要在祂腳前俯伏敬拜.
\newline
\newline
讓我們加入永恆的詩歌,呈獻冠冕,一同稱頌萬王之王.在你的人生中,有沒有這個異象來管理你呢?我們心裡現在所享受的喜樂,我們要繼續領受喜樂.
\newline
\newline


\newpage

\section{第56屆~港九培靈會~高文~第十講　羔羊的婚筵}
\label{sec:416}
\textbf{第十講　羔羊的婚筵}
\newline
\newline
連結: \href{https://www.hkbibleconference.org/session-message/view/416}{\texttt{https://www.hkbibleconference.org/session-message/view/416}}
\newline
\newline
\hyperref[sec:415]{< < < PREV SERMON < < <}
~
\hyperref[sec:toc]{[返目錄]}
~
\hyperref[sec:382]{> > > NEXT SERMON > > >}
\newline
\newline
啟示錄 19:1-19:9
\newline
\begin{longtable}{cl}
\hline
\hline
章節 & 經文 (和合本修訂版)\\
\hline
19:1 & \begin{tabularx}{0.7\textwidth}{X} 此後，我聽見好像有一大群人在天上大聲說：「哈利路亞！救恩、榮耀、權能都屬於我們的神。 \end{tabularx} \\ \\ \relax
19:2 & \begin{tabularx}{0.7\textwidth}{X} 他的判斷又真實又公義；因他判斷了那大淫婦，她用淫行敗壞了世界。神為他的僕人伸冤，向淫婦討流僕人血的罪。」 \end{tabularx} \\ \\ \relax
19:3 & \begin{tabularx}{0.7\textwidth}{X} 他們又一次說：「哈利路亞！燒淫婦的煙往上冒，直到永永遠遠。」 \end{tabularx} \\ \\ \relax
19:4 & \begin{tabularx}{0.7\textwidth}{X} 那二十四位長老和四活物就俯伏敬拜坐在寶座上的神，說：「阿們。哈利路亞！」 \end{tabularx} \\ \\ \relax
19:5 & \begin{tabularx}{0.7\textwidth}{X} 接著，有聲音從寶座出來說：「神的眾僕人哪，凡敬畏他的，無論大小，都要讚美我們的神！」 \end{tabularx} \\ \\ \relax
19:6 & \begin{tabularx}{0.7\textwidth}{X} 我聽見好像一大群人的聲音，像眾水的聲音，像大雷的聲音，說：「哈利路亞！因為主---我們的神、全能者，作王了。 \end{tabularx} \\ \\ \relax
19:7 & \begin{tabularx}{0.7\textwidth}{X} 我們要歡喜快樂，將榮耀歸給他；因為羔羊的婚期到了，他的新娘也自己預備好了， \end{tabularx} \\ \\ \relax
19:8 & \begin{tabularx}{0.7\textwidth}{X} 她蒙恩得穿明亮潔白的細麻衣：這細麻衣就是聖徒們的義行。」 \end{tabularx} \\ \\ \relax
19:9 & \begin{tabularx}{0.7\textwidth}{X} 天使對我說：「你要寫下來：凡被請赴羔羊婚宴的人有福了！」他又對我說：「這些都是神真實的話。」 \end{tabularx} \\ \\
[1ex]
\hline
\hline
\end{longtable}
在這十個晚上,我們思想天上寶座的異象.可惜因時間不足,未能把啟示錄中的異象作詳細講解；只能向大家作扼要的介紹.
\newline
\newline
在美國,凡有關啟示錄中的樂歌的知識書,我都盡量閱覽過；我也曾為這些樂歌寫本知識書.亞洲歸主協會正考慮85年把這本小冊譯成中文；假如一旦出版的話,各位可以有些詳細資料閱讀.
\newline
\newline
世界的系統和所有邪惡最後的終結這十晚,我們沒有時間看邪惡對世界的摧殘,我們集中注意力,在於天上的寶座和圍繞寶座的樂歌；但在一次寶座景象和另一次寶座景象之間,記載了世上遭受的災難和摧殘；到了末期,摧殘變本加厲.第十六至十八章所載恐怖令人震驚的情況,簡直非筆法能描述完全；我們看見世界在神審判之下的可怕.到了第十九章,聖經提醒我們,一切災難都已成過去；印已揭曉,天上號筒吹畢,七碗都已倒在地上；一直在等候死亡的景象已實現了,所有的恐怖戰兢都停止了；惡者殘暴的攻擊都過去了.
\newline
\newline
神伸討大淫婦流人血的罪　「此後,我聽見好像群眾在天上大聲說:哈利路亞,救恩,榮耀,權能,都屬乎我們的神；祂的判斷是真實公義的；因祂判斷了那用淫行敗壞世界的大淫婦,.......」(1-2)天上群眾得到榮耀,勝利凱旋的教會,因而齊聲歌唱；情不自禁地「又說,哈利路亞.燒淫婦的煙往上冒,直到永永遠遠」(3)意思是說,這個世界那些墮落的,被撒旦邪煙薰淘的惡勢力已經崩潰無遺了,整個宇宙已經脫離了罪惡及其勢力.燒盡了的煙,表明神對我們永不改變的信實.
\newline
\newline
合一的敬拜和讚美「那廿四位長老與四活物,就俯伏敬拜坐寶座的神,說,阿們,哈利路亞.有聲音從寶座出來說,神的眾僕人哪,凡敬畏祂的,無論大小,都要讚美我們的神.」(4-5)當約翰一直望著天上寶座時,他聽見好像群眾的聲音,眾水的聲音,大雷的聲音.天上的眾天軍,群眾一同稱頌神；長老,撒拉法,天使,和多不勝數的群眾,在靈裡十分合一；如眾水奔騰澎湃,衝擊海岸；大雷霹靂如萬如萬軍吶喊.這是最終結凱旋的歡呼耶穌為王,坐在寶座上的是我的主；不再是該撒了.耶穌手裡掌握著公義的權杖.「我們要歡喜快樂,將榮耀歸給祂；因羔羊婚娶的時候到了,新婦也自己預備好了.」(7)當約翰聽見天上應對的聲音,他充滿敬畏希奇,就俯伏於地.接著約翰聽見天使的聲音吩咐說:「你要寫上,凡被請赴羔羊婚筵的有福了,又對我說,這是神真實的話.」(9)神差遣天使吩咐約翰,為了地上仍生在世的人,要把所見天上寶座的異象,寫下來傳給他們知道.天上坐寶座者的聲音提醒我們,重新檢討生活的價值觀念,把優先的次序重新編排.神最關心的第一件事,就是呼召我們來敬拜祂.讓我們歡喜快樂,把頌讚榮耀歸給祂.我們的主在世時曾說:「當你受逼迫時,你要歡喜快樂；因為在天上的賞賜是大的.」天上寶座四圍的群眾,是否記得主耶穌在世時所說的話,我們不得而知；但我們當明白歸榮耀給神,再次給我們知道所有屬神的人都來頌讚祂.無論我們的身份,境況如何,能夠讚美神是我們最大的喜樂；這是神創造我們的原因.我們屬於祂,我們被創造來敬拜祂,一直到永恆,我們當俯伏在祂腳前向祂敬拜.惟有關心神的榮耀向神敬拜,才造成今日我們在世上工作的原因；惟有受造之物如此向神敬拜,才真正實現被造之目的.
\newline
\newline
一七四一年有個詩人寫了長的讚美詩,取材於聖經中讚美的話.他把稿件寄給作曲家韓德爾並附言:「這些字是神賜給我的.」當時韓德爾遭受英國貴族的欺壓；財富,生意都走下坡,甚至負債.他非常灰心喪志,疲倦,孤獨住倫敦某處,不與他人來往.當他打開郵包,全神貫注其中聖經的字句,他的眼睛被吸引,「祂被藐視,被人厭棄,祂渴望有人同情；但卻失望.沒有人真正成為祂的安慰.」韓德爾想起這些字句所指是誰；他繼續讀下去「祂信靠了神,神並沒有將祂的靈魂撇在陰間,祂將安息也賜給你.」第二批字句又在韓德爾心中引起共鳴:「奇妙,策士,全能的神,榮耀之君；我知道我救贖主活著,我歡樂!哈利路亞!」他越發感到天上樂歌之時,他執筆作曲配合那些詩詞.這位作曲家關閉自守廿四天,僕人送食物給他,他都不理會；有時拿起麵包要吃,麵包跌下地上；他靈感一來,又繼續下筆,他的僕人看見他滴下眼淚在音樂紙上；有時忽然跳躍到鋼琴邊,有時舉起手向天大聲連呼「哈利路亞!」韓德爾說:「我看見天開了,看見天上那位榮耀的主.」他被舉起到另一個世界,他從未有心靈被高舉起來的經歷.他不知道他的心靈得到極度喜樂,就在那幾天,他整個心靈充滿神的奇妙,神的榮耀,神的愛.神的榮耀完全吸引他.這位傑出的作曲家,就在那裡寫了一首舉世聞名的彌賽亞神曲.當這首詩歌演奏,唱到萬王之王,萬主之主時,英國皇帝大受感動；在演奏之中,突然站立起來；其他的聽眾也都隨著站立,直到演奏完畢才坐下.時至今日都成為慣例,彌賽亞神曲,無論演奏演唱,都把人帶到向上敬拜的境地.當我們稱頌,讚美,慶祝神的偉大時,我們的心靈分享神的榮耀和偉大.
\newline
\newline
期待榮耀之主再臨「我們要歡喜快樂,將榮耀歸給祂；因為羔羊婚娶的時候到了,新婦也自己預備好了.」(7)聖經常用「婚筵」「婚娶」代表神與祂子民的關係.到了世界的終局,榮耀的主再來執掌王權；祂吩咐聖徒參加婚筵,祂是新郎,要舉行婚娶之禮.耶穌所說的比喻中,王為兒子婚娶；請大家來赴王子的婚筵.我們必須明白當時婚娶的習例.那時,訂婚和結婚是兩件事情,其間相隔兩年之久.訂婚之後,就以夫妻稱呼；如果解除婚約,就等於離婚之嚴重性；訂了婚彼此各住父家,直到結婚之日.因此,聖經稱約瑟是馬利亞的丈夫；其實那時他們不應只訂婚而己.到了結婚之日,新郎的朋友們聚集在新郎家裡,然後浩浩蕩蕩一齊到新娘家中；新娘所有姊妹們就陪著新娘,和他們一齊到新郎的家.這樣就算完成了婚姻的儀式,大家享用豐富的婚筵,盡情歌唱快樂,不可言喻.
\newline
\newline
等到天上的新郎來臨之日,婚禮才算真正進行；那時新郎基督要把祂的新婦帶到天上,與天父同享婚筵,無限歡樂!今日的世代,教會與信徒必須忠貞為那永恆之日而活,我們等待主耶穌與天上眾天軍,一同顯現的日子；祂要接我們到天上舉行婚宴,永遠與祂在一起.這樣,你就能明白第9節所說:「凡被請赴羔羊之婚筵的有福了.」
\newline
\newline
喝新杯的日子　這段聖經把信徒比喻為新婦,被請赴結婚筵席.耶穌說:「凡在神的國一齊進食的有福了.」正是這幅圖畫.耶穌曾說:「要從東,南,西,北,各方有人來一同在天上坐席.」耶穌和門徒同進最後晚餐時說:「這杯你們要喝盡,從今以後,我不再喝這葡萄汁；直到在我父的國裡,同你們喝新的那日子.」因此,我們每逢守聖餐同時,我們期望羔羊的婚筵與主一同坐席；到那日,關於耶穌基督彌賽亞的應許都實現了；到那日；偉大的使命都完成了.到那日,神的子民從各國各方各民都被招聚而來.到那日,我們的主再舉起杯來,和門徒在羔羊婚宴中一同坐席.因此,我們今日腳步輕鬆,快樂,我們的眼睛光明滿有盼望.
\newline
\newline
教會是個正等待結婚日子的新娘,聽見天上的聲音說:「與我們一同歡樂吧!一同學習赴羔羊的婚宴吧!」這聲音有個很實際的最後形容,新娘宣佈自己經預備好了；有光明,聖潔,白色的細麻衣給她穿上,準備赴婚宴,既然新郎來了,我們準備迎接.讓我們確知,我們所穿的衣裳已蒙羔羊寶血洗滌乾淨,成為潔白的衣裳；讓我們的生活無愧配上神聖潔的管制.聖經說:「惟有清心的人才得以見神；若不聖潔,沒有人能見主.」聖靈提醒我們衣裳沾污之時,我們應及時求神潔淨,潔白地朝見神.
\newline
\newline
感謝主的恩典實在更多,當有人真誠向主承認自己的罪時,耶穌的寶血即洗淨我們一切的罪,除去我們一切的不義.因此,耶穌愛我們,獻出祂的性命,好叫將來祂的教會成為榮耀的教會；好叫我們無瑕無疵,聖潔不沾污,成為無縐紋的身體.以弗所書告訴我們,基督怎樣愛教會,同樣,丈夫應當怎樣愛妻子,妻子必須是聖潔毫無瑕疵的；基督怎樣愛我們,為我們捨命；我們今天回應的愛也必須如此深厚.
\newline
\newline
記得卅三年前我結婚時,新娘和我共享快樂之日.當我們準備結婚之前,新娘趕縫嫁衣；她的同學相助縫製.當時西方習俗,新郎婚前不得見新娘的嫁衣.我只好遵守風俗,其實心裡很想看看；有人洩漏消息給我關於新娘嫁衣的質料和款式,我巴不得先睹為快.卅三年前之六月三日終於來到,在一間很小的禮拜堂,我站在祭壇前；第一次看見我的新娘穿著美麗的嫁衣.她的父親把她送進來,最漂亮的婚紗乃前所未見,長而潔白,閃亮,毫無瑕疵；她面容顯露愛的情感,真正存著愛心的出現.她也望著我,表現我們的愛何等完整,我知道她已把一生完全付託給我了.行婚禮的誓約多麼寶貴,公開當眾宣佈我們的愛.卅三年來我越加了解誓詞的深意,有時我們共度水深火熱的日子,二人共同承擔,感到那種日子也是美的；雖然有時也有愁苦掙扎；但這些更增加我們相愛的容量.盡我所能全心,全意,全力愛她；所以婚姻十分美滿,時至今日,我們仍然繼續持守.
\newline
\newline
各位!假如我們迎娶新娘,看見婚紗濁污不堪滿有縐紋,我們將何等失望；對方不看重婚禮,因為新娘沒有把最好的帶來；從面容上發現好像對婚禮有所懷疑.那麼,我的心將何等痛苦!如果我站在祭壇前所等待的新娘,表現非全心全意愛我,我心裡的感受將會如何呢?我常以此比喻我們的主耶穌,祂希望我們眾人用最高的愛愛祂；就算我對祂的話語,祂的旨意,不甚了解；但我們應盡所能愛祂；祂給我們多少的愛,我們當捨己地把全心全意的愛,傾倒在祂身上；為我們完成將來的盼望.這就是聖經所說的「完全的約束.」我的意思並非說你現在就用完整十全十美的愛愛祂,乃是說,按你心裡所知道的有多少盡所能愛祂.耶穌告訴我們,愛就是所有先知和律法的總歸.當我們如此愛祂時,我們帶著興奮的心情等待新郎耶穌的出現；好叫號筒吹響天上出現歡呼聲之時,我已預備整齊.看哪!新郎來了,你是否聽見從天上發出的聲音呢?榮耀的神從寶座發出宣告；群眾歌聲繼續響應,所有受造之物,都向地球發出呼聲:「歡喜快樂歸榮耀給祂；羔羊的婚筵現在來臨了,祂的新婦都預備好了.」光明潔白的細麻衣,作新婦的婚紗；這細麻衣就是聖徒所行的義.「你要寫上凡被請赴羔羊之婚筵的有福了.」凡有耳的,就應當聽神的靈向眾教會所發的聲音.你有否聽見天上來的信息呢?世界淪落呼喊的聲音,你聽的太多；世界千方百計盡其能事,要把我們拉攏到地獄的深處.在你掙扎之中,你四圍的人都憎恨你.我們今日需要天上寶座和四圍的價值觀念來衡量這個世界.你是否作過敬拜神,在神面前俯伏的約裡面呢?你有否想到當主顯現時,我們見祂,我們應當時刻存著興奮的心情等待主之再臨.
\newline
\newline
各位!如果在你心靈的深處,在你基督徒的良心裡,尚未預備好來迎接主的話；讓我勸勉你,今天晚上你要預備好.如果今晚主回來,你尚未預備好朝見主；讓我懇請你,要把你的心回轉,預備齊全來迎見這位萬王之王.恐怕有許多基督徒的生活方式,將與世界共存亡；其實在你心靈裡面,你常有空虛,不安全,不滿足的感覺；當我們覺得心靈四面受敵時,正是最重要與神修好關係之時.
\newline
\newline


\newpage

\section{第56屆~港九培靈會~麥希真~第一講　信神,信人,自信?}
\label{sec:382}
\textbf{第一講　信神,信人,自信?}
\newline
\newline
連結: \href{https://www.hkbibleconference.org/session-message/view/382}{\texttt{https://www.hkbibleconference.org/session-message/view/382}}
\newline
\newline
\hyperref[sec:416]{< < < PREV SERMON < < <}
~
\hyperref[sec:toc]{[返目錄]}
~
\hyperref[sec:384]{> > > NEXT SERMON > > >}
\newline
\newline
馬太福音七章二十四至二十六節主耶穌對門徒說:「所以凡聽見我這話就去行的,好比一個聰明人,把房子蓋在磐石上.雨淋,水沖,風吹,撞著那房子,房子總不倒塌；因根基立在磐石上.凡聽見我這話不去行的,好比一個無知的人,把房子蓋在沙土上,雨淋,水沖,風吹,撞著那房子,房子就倒塌了,並且倒塌得很大.」
\newline
\newline
林前十三章說,如今常存的有信,有望,有愛,這三樣,其中最大的是愛.我們要留意到愛是最大的,但若沒有信心就沒有盼望,若沒有盼望就不會去愛.愛心是最大的,但信心是最根本的.如果你沒有信心,一切都不能實行出來.有穩固的信心,才有資格講其他.
\newline
\newline
主耶穌曾說:「你們心裡不要憂愁,你們信神,也當信我.」(約十四1)的條件,就是信神,單信神也不夠,主耶穌說:「你們也當信我.」有些基督徒以為,信神就夠了.當然,人看見宇宙奇妙美麗而信有一位主宰,又或在聖經中知道這位天父,而歸依祂,成為祂的兒女這的確是最好；但還不夠,且要信耶穌基督.(雅二19)更清楚說明這真理:「你信神只有一位,你信的不錯；鬼魔也信卻是戰驚.」一個基督徒若只說他信一位神,而不認識耶穌基督的話；他的信仰是糊里糊塗的.人信神,也接受耶穌基督做救主,他面對甚麼困難,也會有平安,不再憂愁.
\newline
\newline
單是信神是不夠的,若只加上信耶穌為你受死只算信了一半,你也要信耶穌基督從死裡復活.你若對耶穌復活不清楚,你的信心是不堅固的.我父母是基督徒,小時我也過聖誕節和復活節,但到我長大之後,我要有把握,很清楚的,知道耶穌確是從死裡復活了.當我們面對艱難甚至生命威脅時,是否仍接受耶穌復活這個事實?在這個聚會裡,我們要重新檢討自己的信心.
\newline
\newline
一九四九年我蒙召作傳道,於是就到伯特利神學院讀書.那時我曾檢討自己的信仰,特別是耶穌復活方面.祂的復活若只是一個神話,或是一個象徵性的教訓,我就不願意走上這條奉獻的道路,一生傳講這個似是而非的故事.那時我讀了一些支持和推翻復活的書籍,又在神面前禱告,我終於清楚知道耶穌確實是復活了,在這個情況下,我甘心把自己奉獻給祂.
\newline
\newline
聖經是神的話,其中最主要的根基,就是主耶穌復活一事；若這事我們不清楚,到了面對考驗的時候,我們的信心立刻就倒塌.正如主耶穌說,你聽我的話不去行,好比一個無知的人,把房子建立在沙土上,雨淋,水沖,風吹,撞著那房子,房子就倒塌,並且倒塌得很大.
\newline
\newline
約翰福音二章二十二節記載說:「所以到祂從死裡復活以後,門徒就想起祂說過這話,便信了聖經和祂所說的.」門徒聽見耶穌說,這個聖殿雖是四十六年建成的(約二19-20),但若拆毀它,祂可以在三日內重建.門徒對耶穌這番話,似懂非懂,也不問甚麼,聽了就算；但到耶穌死了之後,第三日從死裡復活時,才明白主耶穌是用聖殿表明自己的身體.門徒很清楚見到耶穌釘死在十字架的過程,羅馬兵又用槍扎祂的肋旁；證實耶穌的確死了,後來由亞利馬太約瑟,求彼拉多把祂的屍體取下,安葬在墳墓裡.門徒實在親眼看見自己所信的夫子,「神的兒子」耶穌基督是死了.但感謝神,死了的耶穌在第三日復活了,首先向抹大拉的馬利亞顯現,然後向幾個婦女顯現,跟著向彼得顯現；可能就在那個禮拜日也向雅各顯現,又在以馬忤斯路上向兩個門徒顯現,到晚上在關上門的屋裡向那十個門徒顯現.相隔一個禮拜日又向十一個門徒顯現,叫多馬伸出指頭摸祂手,伸出手來探祂的肋旁；多馬於是跪下對主說,我的主,我的神,你從死裡復活.後來在提比哩亞海邊,在加利利山上,在耶路撒冷對面橄欖山上等處向他們顯現；他們確實知道主耶穌從死裡復活了,就相信聖經所講的預言和耶穌說過的話.
\newline
\newline
今日我願意與大家重新檢討我們的信仰.門徒知道耶穌的確是復活了,就接受耶穌的話和整本聖經的話是神的啟示.耶穌復活是我們信仰的根基.聖經不是一本參考書,也不只是一本道德的書,而是神的話.好像當時的門徒那樣,到主耶穌從死裡復活之後,想起祂說過的話,便信了聖經和祂所說的.求天父幫助我們,第一,去信神,然後更進一步接受耶穌做我們的救主.(約一12)說:「凡接待祂的,就是信祂名的人,祂就賜他們權柄,作神的兒女.」但這還不夠,還要認識:耶穌第一為我們的罪死了,埋葬了；第二,祂從死裡復活.有了這些認識後,我們就相信聖經和耶穌所說的都是神的話,是我們所信,所倚靠的.這樣的信心才是真信心.這種信心不但使我們在平安時作基督徒,在面對逼迫和困難甚或生命財產損失時,也可以仍然站立得住；否則我們就好像建立在沙土上的房子那樣,完全倒塌.
\newline
\newline
跟著我要談到信人.耶穌在約翰福音十四章十二至十四節有這樣的教導:「我實實在在的告訴你們,我們作的事,信我的人也要作,並且要作比這更大的事,因為我往父那裡去.你們奉我的名無論求甚麼,就必成就,叫父因兒子得榮耀.你們若奉我的名求甚麼,我必成就.」主在這裡有一個奇妙的應許；信徒不但能作祂所作的,並且所作的是更大的.許多基督徒表面上很謙卑,但其實是他們不信主耶穌的應許.求主叫我們弟兄姊妹彼此信任,彼此同心,我們就能夠為主做大事；正如葛培理佈道會,以往作了許多大事,許多人因此信主.現代信徒有各種不同而有效的大眾傳播媒介；可以很短的時間將福音傳遍普天下,而基督所能接觸的範圍,只是巴勒斯坦這一小塊地方.
\newline
\newline
在教會裡,我們只是針對彼此的軟弱和缺點,結果沒有同心合意的禱告,彼此不信任.其實我們雖然有許多軟弱和缺點,但若願意同心禱告,神就會在我們當中成就大事；教會就能結果三十倍,六十倍,一百倍.求主幫助教會的信徒,牧師,執事等彼此信任,我們自己雖然不值得信任,但卻信主在我們當中能作大事.
\newline
\newline
我們再看另一段經文,這段經文是講「自信」,請看約翰福音十四章廿七至廿九節:「我留下平安給你們,我將我的平安賜給你們；我所賜的,不像世人所賜的.你們心裡不要憂愁,也不要膽怯.你們聽見我對你們說了,我去,還要到你們這裡來.你們若愛我,因我到父那裡去,就必喜樂,因為父是比我大的.現在事情還沒有成就,我預先告訴你們,叫你們到事情成就的時候,就可以信.」
\newline
\newline
「自信」有四方面作用:第一,有自信就能除去困難,因信主有平安喜樂,就能勝過這些困難.第二,是能克服困難,因主應許有平安和喜樂,靠主就滿有自信,勇敢的繼續去做,雖然困難仍存在,但靠著主,困難就慢慢的解決了.第三,是能接受困難,困難若無法除去或克服,但仍然可以與困難一起.第四,是通過困難,困難有時不但仍存在,而且越來越大,簡直把自己壓倒了；但因有主的平安和喜樂,就能通過這困難.我們個人的生活以至社會,家庭,教會的生活都會遇到以上四種情形；靠著主都能除去,克服,接受甚至通過種種困難.
\newline
\newline
基督徒因信以致於信.開始的時候只有少許信心,一次一次的經歷事情的成就；信心一次一次的增長,越來越大.有了信心,就有盼望,絕望的時候仍有盼望；有了盼望,我們就能有關心和力量去愛.沒有埋怨和憤恨,絕不怨天尤人,而是充滿愛心去愛.
\newline
\newline


\newpage

\section{第56屆~港九培靈會~麥希真~第二講　一個平凡基督徒的禱告生活}
\label{sec:384}
\textbf{第二講　一個平凡基督徒的禱告生活}
\newline
\newline
連結: \href{https://www.hkbibleconference.org/session-message/view/384}{\texttt{https://www.hkbibleconference.org/session-message/view/384}}
\newline
\newline
\hyperref[sec:382]{< < < PREV SERMON < < <}
~
\hyperref[sec:toc]{[返目錄]}
~
\hyperref[sec:386]{> > > NEXT SERMON > > >}
\newline
\newline
今天我們要思想基督徒的禱告生活,我們分眾人的禱告和個人的禱告這兩部份來思想.
\newline
\newline
我們先看眾人的禱告.(徒四31)「禱告完了,聚會的地方震動,他們就都被聖靈充滿,放膽講論神的道.」請注意這裡所說的震動,並不是形容眾人禱告這聲音很大,以致門窗受震動.這裡所說的,是禱告的果效；當眾人禱告完了,聚會的地方才受震動.基督徒個人的禱告固然重要,但整體性或小組的禱告也是不可缺少的.
\newline
\newline
我提出八種禱告方式,是適合團體禱告用的,這八種方式應該時常轉換,按著教會和團體的需要不停的轉換.正如我們睡覺一樣,睡姿是常常轉換的,這樣才舒服.飲食也是一樣,食物的選擇,烹調的方法也要經常轉換,這樣才不致吃厭.屬靈的事物也是一樣,領人信主有很多方式,如果只固執用一種,便很難領人信主.有的人要用四個屬靈定律,有的人要用福音橋,有人要用幸福人生的錄音帶,方式是要經常換的.禱告也一樣,需要用不同的方式,要按情況按需要時常更換.
\newline
\newline
這八種不同的禱告方式:
\newline
\newline
(一)可以用祈禱文.聖公會,信義會,循道會等都常用祈禱文.事實上,祈禱文很有用,特別是對初信的人.我們可以在團體禱告的時候,除了寫下禱告的事項,還把禱詞寫下來,分給弟兄姊妹,這對初信或初次參加祈禱會的人,十分有用.總之,教會的公禱文,或者是每次禱告會所印發的禱文,都是很有用的.
\newline
\newline
(二)輪流領禱.這是我們最傳統,最習慣和最常用的.先提出幾樣禱告事項,然後分別請幾位弟兄姊妹領禱.這種方式是最規矩,最嚴肅的,但亦很容易令弟兄姊妹散漫,心不在焉.這個我們要了解.所以輪流領禱仍然可以採用,但亦要經常加插別的方式.
\newline
\newline
(三)同聲祈禱.很多教會都愛用這個方式,不過有一點要注意的,就是切勿太大聲,否則很容易會混亂.但同聲祈禱是好的,因為大家可以同心,同聲,而且很彼此有激勵,只要用時間注意他的長處和弱處便可以了.
\newline
\newline
(四)談話式的祈禱.很多大學生都喜歡用這種方式,他們圍在一起禱告,一人講幾句,到最後作結束的那個才說奉主名求,阿們,來結束整個禱告.這個方式的禱告非常好,我們可以因著某種場合的需要,而採用這個方式.
\newline
\newline
(五)集五,三人一小組的禱告,這是教會近來流行採用的.不過有兩點我們要注意:第一是這樣的分小組不需另找地方,只要三個人彼此坐近些,大家聽到對方的禱告便成,不必另找分組的地方.第二,禱告事項一定要印出來,這樣會比較詳細方便,而且分小組禱告,代禱的事項便不應只是三五項,而最少應該有十多二十項,才有意思.分組之後,每人為三五件事禱告,整個祈禱會合起來,便會來得實際.很多時候,信徒不參加祈禱會,是覺得空洞,沒有意義.因為主席報告和解釋,已花去大部份時間,而剩下來的禱告時間,內容又不切實際；所以慢慢便沒有人來.但如果是有實質的禱告,手頭上有一系列的禱告事項,那麼禱告的意義便完全不同了.
\newline
\newline
(六)二人一小組.特色又是不必另找地方,由主席把禱告事項一樣一樣的讀出來,可以由兩人中的一人領禱,或者兩人都輪流禱告,但每樣禱告事項最好不超過一分鐘,這樣才會有禱告的實質.在禱告中,不必和天父講甚麼客套話,乃是一針見血的把要禱告的事講出來便成.
\newline
\newline
(七)連鎖的禱告.可以在教會或家庭舉行,把時間分配好,一個續一個的,每人或每單位佔十五或二十分鐘,務使教會二十四小時都有人禱告,這是非常有意思的.
\newline
\newline
(八)通宵的禱告.這通常是為教會某些特別的需要而舉行的,不過要小心一點,就是不能單由青年的弟兄姊妹參加,免得被人誤會.要有教會的長者一同參與,方可舉行.通宵禱告必須事前安排好程序,可以有唱詩,見證,讀經,分享等加在一起.而禱告的時候,又可以採用不同的方式,或整體性的,或分小組的,務求各人都有迫切的心在神面前禱告.
\newline
\newline
個人的禱告固然重要,但整體的禱告也一樣重要.怎樣才可以一同禱告,彼此火熱呢?今天我所提的,就是用不同的方式去禱告,以致禱告更有實質；可以經歷到禱告完了,大家的心靈震動,大家為主事奉的心志震動,甚至整個教會都震動.
\newline
\newline
跟著我們看個人的禱告.首先是定時的禱告,耶穌基督,彼得和保羅都有定時的禱告.我們也可以在起床,吃飯和睡覺之時都禱告,但我所指的是除了這些以外.定時禱告應該包括讚美,感謝和代求.讚美是數算神的美德,神的聖潔,公義,慈愛,永恆,不變,主權等,都是我們應該數算和讚美的.也許我們覺得很難學,但不要緊,我們可以多讀詩篇,把有關神美德的經文劃下,細細思想,這樣慢慢我們便會懂得讚美神了.
\newline
\newline
感謝是數算神的恩典,感謝神所賜給我的.而代求有兩方面可以實踐的.
\newline
\newline
第一方面是為未信主的人禱告.找一張卡,寫下七個你要向他傳福音人的姓名,然後順著次序,每天為一個禱告；又按機會向他傳福音,帶他聚會.這樣我們的工作才有目標,我們所努力的才會有收成.另外,可以用同一張卡的背後,劃上八個格,前面七個格寫下星期一到星期天的代禱事項,第八格則寫下每天都必須代禱的急切事項；這樣依照卡所編排的,每天定時禱告,有讚美,有感謝,也有代求.這樣的禱告,必然充實,而不會日久生厭,覺得味同嚼蠟而因此放棄.
\newline
\newline
第二方面,是隨時的禱告.(弗六18)「靠著聖靈,隨時多方禱告祈求.」帖前五17「不住的禱告.」神既然是無所不在的神,我們自然可以隨時隨地的向祂禱告.其實,每一件事都值得我們向主禱告.當你在街上走,看到美麗的景色時,你可以內心向神發出感謝.當你遇到困難或有疾病時,你可以立即求主醫治和幫助.即使是生活的小節,你也可以化成內心的禱告.有一個在神學院作廚師的,寫了一本書,講到他怎樣學習禱告.當他掃地時,他求主潔淨他的心,好像他掃去地上的垃圾塵埃一樣.當他擺碗筷時,他求主讓他成為能夠使人得著飽足的食物.他不是神學院的教授,他只是一個廚師；但是他學會了這個偉大的功課.今天,我們真需要學習隨時禱告與主交談.當我學會隨時隨地禱告的功課之後,我屬靈的生命得著了很大的改變.
\newline
\newline
最後是禁食禱告.(太四1-2)是基督的榜樣.(六16-17)耶穌教導門徒要禁食禱告,但不要學法利賽人那樣故意叫人看見.(太十七20-21)禁食禱告可以加添我們的信心,使我們能作大事.禁食祈禱的兩個原因就是,當我們遇到大困難,無法勝過撒旦時,便要禁食祈禱.另外,當我們有大的決定,面臨大的機會和工作時,也要禁食祈禱.(徒十三2-3)講到有小組和全教會祈禱.(徒十四23)要為教會的領袖禁食祈禱,不單為牧師傳道,更要為有困難有問題的長執祈禱；而且是禁食祈禱,這樣教會的問題便可以得著解決.
\newline
\newline


\newpage

\section{第56屆~港九培靈會~麥希真~第三講　如何讀經,而且實行出來}
\label{sec:386}
\textbf{第三講　如何讀經,而且實行出來}
\newline
\newline
連結: \href{https://www.hkbibleconference.org/session-message/view/386}{\texttt{https://www.hkbibleconference.org/session-message/view/386}}
\newline
\newline
\hyperref[sec:384]{< < < PREV SERMON < < <}
~
\hyperref[sec:toc]{[返目錄]}
~
\hyperref[sec:388]{> > > NEXT SERMON > > >}
\newline
\newline
(約五39-40)耶穌說:「你們查考聖經,因為你們以為內中有永生,給我作見證的就是這經.然而你們不肯到我這裡來得生命.」留意這兩節聖經所講的三件事,第一,是關乎永恆的事,永遠的價值,就是永生.第二,聖經是為耶穌作見證的.從創世記到啟示錄,整本聖經,都是為耶穌作見證.第三,最要緊的,是讀經之後,要來到耶穌那裡得永生.因此,我們讀經,可以知道永生,可以認識耶穌；但最重要的,是我們來到耶穌那裡得永生.
\newline
\newline
(一)我們看見神是超自然的,須要啟示祂自己.這裡有七點我要講的.
\newline
\newline
1.為甚麼我們知道聖經是神的啟示呢?因為它只有一個作者,就是神自己.雖然聖經由四十多人執筆寫成,而且頭一個和最後一個相差有一千五百年,最早的一卷大概是在非洲埃及寫成的.第二部份則在亞洲巴勒斯坦寫成,至於新約的大部份則在歐洲,希臘,意大利一帶寫成.雖然四十多位作者,在一千五百年中,分別於歐,亞,非三大陸寫成聖經；彼此不相識,但整本聖經卻是一整體,因為只有一個作者,就是神.
\newline
\newline
2.聖經就是神超自然的啟示,是因為它從頭至尾六十六卷只有一個信息.聖經由頭到尾有一條紅線穿過,這紅線就是羔羊的血.雖然有四十多位作者,又分佈於三大洲,但所寫的只有一個信息,就是羔羊的血.其中只有一個理由,就是聖經是神超自然的啟示.
\newline
\newline
3.聖經是神超自然的啟示,因為它有預言的證明.一九四八年,我的主日學老師從報紙上剪下一則新聞給我們,標題是看哪!聖經的預言應驗了.這預言就是猶太人復國.公元七十年,猶太耶路撒冷被羅馬軍隊攻破,殺了一百五十萬猶太人,猶太人從公元七十年開始便分散在全世界.但聖經卻預言有一天,神要在天的四方,把他們招聚回來,到耶路撒冷建立他們的國家.直到一九四八年五月十四日下午四時；在一種很微,四強都承認的情況下；猶太人便在巴勒斯坦恢復他們的國家,聖經的預言應驗了.
\newline
\newline
4.逼迫的證明.聖經受刀劍,理性,政治的逼迫.然而各樣可怕的逼迫之後,聖經仍然存在,而且變得更加輝煌明亮.
\newline
\newline
5.時空的證明.聖經超越時空,在農業時代是真理,到今天工業時代也是真理.落後的時代是真理,到今天太空電腦的時代仍是真理.不但超越空間.在東方,在西方,在熱帶,在北極,聖經都是真理,是放諸四海皆準的真理.
\newline
\newline
6.個人的證明.每一個基督徒,都經歷到聖經的真確,都能證明聖經是神的說話.
\newline
\newline
7.是基督的證明.耶穌親口說,聖經內中有永生,聖經是為我作見證,你們可以按著聖經的應許到我這裡來得生命.
\newline
\newline
(二)跟著我們要思想的,就是聖經是無誤的,是超自然被神保守沒有錯誤的.
\newline
\newline
1.基督的態度,基督認為聖經是無誤的.
\newline
\newline
2.聖經的態度,舊約,新約彼此認為聖經是無誤的.
\newline
\newline
3.受死的應驗.聖經對耶穌的受死講得很清楚,(賽五十二13-15),(五十三章),講到耶穌是救主,是世人所驚奇,君王所尊崇的.但同時又講到這位救主面貌憔悴,形容枯槁,為甚麼會這樣矛盾呢?八百年後,耶穌來到世上,以賽亞所講一切有關於祂的預言,全都應驗了.不單是受死的預言,連復活的預言也應驗了.這一切都表明聖經的無誤,最後是歷史的證明.歷史愈長久,科學愈發達,各種學問愈廣傳,聖經的真實性便愈加顯明.(賽二十1)提到亞述王撒珥根二千年來人都認為這記載是錯誤的,因為亞述根本沒有撒珥根這個王.但現在的考古學,卻在亞述首都尼尼微的遺址內；找到一些古代的泥版,上面竟有亞述王撒珥根的名字.他是屬於內部的王,普通的歷史書是不記載的.
\newline
\newline
(三)聖經是不矛盾的,不過有時我們的確是讀到某些看上去似乎是有矛盾的經文,那怎麼辦呢?我舉出一個例子,是有關耶穌復活的記載的.(可十六5)「他們進了墳墓,看見一個少年坐在右邊,穿著白袍,就甚驚恐.」注意這裡所說的只是一個少年人,而且是坐在右邊的.(路二十四4)這裡似乎有不同的記載,經文說:「正在猜疑之間,忽然有兩個人站在旁邊,衣服放光.」到底是一個還是兩個!是坐著還是站著呢?看上去似乎是有矛盾,其實並沒有.當時實在有兩個天使,初時是坐著的,但後來看見婦女們來,便站起；其中的一個天使向他們說話,可見記載是沒有錯誤的.
\newline
\newline
我們再看另一處有相似記載的經文,是講到加大拉,或稱為格拉森被鬼附的人.(太八28)耶穌既渡到那邊去,來到加大拉人的地方；就有兩個被鬼附的人；從墳墓裡出來迎著他,極其兇猛,甚至沒有人能從那條路上經過.」注意這裡所說的是加大拉,是有兩個被鬼附的人.再看(可五2)「他們來到海那邊,格拉森人的地方.耶穌一下船,就有一個被污鬼附著的人,從墳墓裡出來迎著祂.」有兩點是不同的,第一是地名的不同,一處記的是加大拉,另一處是格拉森；這個問題很容易解決,因為加大拉和格拉森都是同一處的地方.可是到底有幾個被鬼附的人呢!很清楚的是有兩個,但其中的一個特別兇猛,所以馬可特別記載他.可見聖經並無矛盾,並無錯誤或不合理.
\newline
\newline
還有另一處也是看上去有問題的.(路十八35)這裡記載耶穌在耶利哥醫好了一個瞎子.「耶穌將近耶利哥的時候.有一個瞎子坐在路旁討飯.」注意這裡是說耶穌將近耶利哥,就是入城的時候,有一個瞎子.我們再看(太二十30)這裡的記載便似乎有矛盾了.聖經說:「他們出耶利哥的時候,有極多的人跟隨他.有兩個瞎子坐在路旁.」到底是入城的時候遇見瞎子,還是出城的時候呢?到底是一個瞎子還是兩個瞎子呢?聖經不是似乎不對嗎?情形就像剛才的經文一樣.馬太是耶穌的門徒之一.當時馬太也在耶利哥城,所以馬太清楚的看見有兩個瞎子.但路加不是十二個門徒之一,他是後來才信主的,可能他後來才到巴勒斯坦,慢慢調查,搜集資料而寫成路加福音的；所以他的記載是後來才寫的.而且後來他只問到其中的一個瞎子,所以只記下一個.故此並無衝突.
\newline
\newline
但到底是出耶利哥還是入耶利哥之時遇見的?這個問題,二千年來並沒有答案,好像真的是有衝突.但感謝主.現代的考古學家把耶利哥城的舊址挖掘出來,發覺原來耶利哥城有新舊二城,耶穌剛好是出了舊城,要進入新城的時候,剛好是在兩城的中間.可能你會覺得這樣的解釋很牽強,但其實並不.我們看(可十46)「到了耶利哥,耶穌同門徒出耶利哥的時候,有一個討飯的瞎子,是底買的兒子巴底買,坐在路旁.」原來聖靈早已有奇妙的安排和默示,使聖經能無誤的記載下來.這節經文和今日考古學家所發現的,天衣無縫.如果不是有考古學的發現,這節經文讀起來根本是毫無意義的,因為讀的人不曉得耶利哥原來有新舊二城.原來並不是馬可記載錯誤,把出城和入城記錯,而是耶穌出了舊城,靠近新城的時候；兩城之間是通衢大道,人來人往,車水馬龍.乞丐就在那裡討飯,而馬可又特別記載了一個後來跟從主的巴底買的事蹟.
\newline
\newline
求主幫助我們,確信聖經的話,不管人怎樣批評,但我們始終堅信.最後我要講的一個例子,是(路七3)這裡講一個有信心的百夫長.「百夫長風聞耶穌的事,就託猶太人的幾個長老,去求耶穌來救他的僕人.」然後第六節「耶穌就和他們回去,離那家不遠,百夫長託幾個朋友去見耶穌,對他說:主啊!不要勞動,因你到我舍下,我不敢當.」這裡講到百夫長先請幾個長老去見耶穌,後來又請幾個朋友去.(太八5-7)耶穌進了迦百農,有一個百夫長進前來,求他說,主啊!我的僕人害癱瘓病,躺在家裡,甚是疼苦.」這裡是講到百夫長親自去請耶穌.兩處似乎是有矛盾,到底是親身去,還是派人去呢?但聖經是無誤的.有一個很清楚的解釋,就是百夫長的確是請長老和朋友去請耶穌；但意思的本身卻是百夫長本身的邀請和誠意,二者的記載根本沒有矛盾.
\newline
\newline
求主幫助我們,不但認識聖經是神超然的啟示；更確信聖經的無誤性,好叫我們站穩在這磐石上.
\newline
\newline


\newpage

\section{第56屆~港九培靈會~麥希真~第四講　今天如何能消除緊張}
\label{sec:388}
\textbf{第四講　今天如何能消除緊張}
\newline
\newline
連結: \href{https://www.hkbibleconference.org/session-message/view/388}{\texttt{https://www.hkbibleconference.org/session-message/view/388}}
\newline
\newline
\hyperref[sec:386]{< < < PREV SERMON < < <}
~
\hyperref[sec:toc]{[返目錄]}
~
\hyperref[sec:389]{> > > NEXT SERMON > > >}
\newline
\newline


\newpage

\section{第56屆~港九培靈會~麥希真~第五講　明天如何能保持堅忍}
\label{sec:389}
\textbf{第五講　明天如何能保持堅忍}
\newline
\newline
連結: \href{https://www.hkbibleconference.org/session-message/view/389}{\texttt{https://www.hkbibleconference.org/session-message/view/389}}
\newline
\newline
\hyperref[sec:388]{< < < PREV SERMON < < <}
~
\hyperref[sec:toc]{[返目錄]}
~
\hyperref[sec:391]{> > > NEXT SERMON > > >}
\newline
\newline
教會是一定受逼迫的,教會若沒有逼迫,這便是教會的問題.我們看聖經中的教會,至少遇見四種逼迫.第一種是同工及會友間的逼迫,同工彼此逼迫,會友彼此逼迫,香港的教會也有.第二種是異端邪說的逼迫.今天外國教會和香港教會都有,這兩種都是教會內部的逼迫.但聖經中的教會有兩種逼迫,就是受著猶太社會的逼迫.他們的教會有很多是猶太人；但不是猶太人的,也受猶太文化宗教的影響；他們為自己本國本民族的文化的逼迫.香港的信徒其實也有這種逼迫.年青人信了主,但家庭卻要他拜祖先,拜偶像,這些都屬於文化的逼迫.最後還有政治的逼迫.他們受猶太人的公會的逼迫,受猶太人在殖民地的政府逼迫,又受羅馬政府的逼迫,被鞭打,用石頭打死,又常常受監禁,沒收財產,甚至被刀槍所殺,被棍打死,還有被釘在十字架上的.他們不單自己受苦,家人也受牽連,被賣作奴隸.當時,教會受盡逼迫,但神卻啟示祂的使徒去安慰他們.有保羅,彼得,約翰,雅各,猶大等,都寫信去安慰他們；給他們各樣安慰的話,而其中特別有一個信息,是各個使徒都分別提到的.
\newline
\newline
這個特別的信息就是,他們雖然受盡逼迫苦楚,但不用怕,因為基督快要再來.當基督再來時,祂要作王；我們所受的一切苦就要變為榮耀與賞賜.聖經啟示祂的僕人,寫信給各個受盡最殘酷,最可怕的逼迫的教會.給他們各樣的安慰,而最大的,最終的安慰,就是耶穌基督快要再來.
\newline
\newline
雖然從理性的角度來分析,我們很難相信基督的再來會給信徒有多大的安慰；因為盡管他們日夕盼望,但基督仍遲遲未來.不過,按我個人的領受,無論是聖經的根據,或者是實際的經驗；基督的再來常常是受逼迫的信徒最大,最有效和最能幫助他們的安慰,叫他們能堅忍到底.因為一切的受苦和逼迫,將來都要變為榮耀和賞賜.耶穌再來不是平安穩妥的信徒的一些點綴品,以為耶穌的再來只是錦上添花.基督再來乃是聖靈藉使徒給與逼迫之中的信徒的安慰和鼓舞；叫他們不用懼怕,因為基督要再來,他們必須堅忍到底.
\newline
\newline
今天我們要思想的,就是基督再來的真理.我們先看十段有關的經文.(路二十一25-28),這是耶穌教訓時的說話.「日月星辰要顯出異兆,地上的邦國也有困苦,因海中波浪的響聲,就慌慌不定.天勢都要震動；人想起那將要臨到世界的事,就都嚇得魂不附體.他們要看見人子有能力,有大榮耀,駕雲降臨.一有這些事,你們就當挺身昂首,因為你們得贖的日子近了.」(約十四1-3),這是耶穌臨別時的話.「你們心裡不要憂愁,你們信神,也當信我.在我父的家裡,有許多住處,若是沒有,我就早已告訴你們了.我去原是為你們預備地方去,我若去為你們預備了地方,就必再來接你們到我那裡去.我在那裡,叫你們也在那裡.」這裡清楚講到耶穌現今去是預備地方,將來祂要再來接信徒到祂那裡.(太二十六62-64),這是審判時的話.耶穌在面臨生死時,說了這番重要的話.「大祭司就站起來,對耶穌說,你甚麼都不回答麼?這些人告你的是甚麼呢?耶穌卻不言語.大祭司對祂說,我指著永生神,叫你起誓告訴我們,你是神的兒子基督不是.耶穌對他說:你說的是.然而我告訴你們,後來你們要看見人子,坐在那權能者的右邊,駕著天上的雲降臨.」
\newline
\newline
(徒一10-11)是天使的話.「當祂往上去,他們定睛望天的時候,忽然有兩個人,身穿白衣,站在旁邊說:加利利人哪!你們為甚麼站著望天呢?這離開你們被升天的耶穌,你們見祂怎樣往天上去,祂還要怎樣來.」天使清楚的告訴門徒,基督要再來.(腓三20-21),保羅得勝靈的啟示,也講到主再來的事.「我們卻是天上的國民,並且等候救主,就是主耶穌基督,從天降臨.祂要按著那能叫萬有歸服自己的大能；將我們這卑賤的身體改變形狀,和祂自己榮耀的身體相似.(啟一7),約翰也講到同樣的話.「看哪!祂駕雲降臨,眾目要看見祂,連刺祂的人也要看見祂,地上的萬族都要因祂哀哭.這話是真實的,阿們.」(彼後三3-13),我們選讀其中的幾節:「第一要緊的,該知道在末世必有好譏誚的人,隨從自己的私慾出來譏誚說,主要降臨的應許在哪裡呢?因為從列祖睡了以來,萬物起初創造的時候仍是一樣.親愛的弟兄阿!有一件事你們不可忘記,就是主看一日如千年,千年如一日.但主的日子要像賊來到一樣；那日天必有大響聲廢去,有形質的都要被烈火銷化,地和其上的物都要燒盡了.」
\newline
\newline
雅各也寫了同樣的話(雅五7-8),「弟兄們!你們要忍耐,直到主來.看哪!農夫忍耐等候地裡寶貴的出產,直到得了秋雨春雨.你們也當忍耐,堅固你們的心,因為主來的日子近了.」除非所有使徒的話都是謊言,否則在聖靈的感動下,他們都不約而同的寫下了主再來的信息.主的兄弟猶大也有同樣的話,(14-15),「亞當的七世孫以諾,曾預言這些人說,看哪!主帶著祂的千萬聖者降臨,要在眾人身上行審判；證實那一切不敬虔的人,所妄行一切不敬虔的事,又證實不敬虔之罪人所說頂撞祂的剛愎話.」最後我們看(來九28)「像這樣,基督既然一次被獻,擔當了多人的罪,將來要向那等候祂的人第二次顯現；並與罪無關,乃是為拯救他們.」
\newline
\newline
這十段經文,清楚的論到基督再來.耶穌在教訓時,臨別時和受審判時都說到祂將要再來.另外天使,保羅,約翰,彼得,雅各,猶大,以及希伯來書的作者均論到基督再來的真理.如果今天仍有人懷疑主的再來,那真是可惜.我們要抓緊基督再來的真理,好好的預備.
\newline
\newline
不過,有人卻對基督的再來有錯誤的解釋.第一,基督再來只是當時的極端思想.這只是當時的人的信仰,是為了安慰當時的人而說的一些話；是為當代的人而說的,並不是真理.第二,基督再來就是聖靈降臨,他們把二者混為一談,說聖靈降臨之時有火,有風,教會有大復興,這就是基督再來與我們同在.但其實所有書信都是聖靈降臨之後寫的,我們的解釋不能有所偏差.第三,有人以為基督再來就是耶路撒冷的被毀,甚至是耶路撒冷被毀之後才記的,並不是預言.第四,他們解釋為基督再來是接受基督,便等於基督再來了.但聖經並不是這樣的意思,他們解釋說基督徒死後回到基督那裡,這便是基督再來.這對聖經是斷章取義,是完全錯誤的.以上五點都是人不信聖經的啟示而加以歪曲的解釋.
\newline
\newline
到底基督再來的意義是甚麼?(徒一10)說基督再來是眾目都看見的.基督升天時,是有一朵雲彩把祂接去的,當時各門徒都看得清楚；將來,主再來時也要這樣,眾目都要看見祂.(太二十四30,二十六64)(帖前四16-17,啟一7-8)都講到基督駕雲降臨,眾目要看見祂.
\newline
\newline
另外,(帖前四13-18)這裡講到基督再來有兩件的事要發生.第一,已死的信徒要復活.第二,活著的信徒要被提.我們留意聖經凡提到主再來的經文,都是當時的教會受到逼迫的.帖撒羅尼迦的教會當時大受逼迫,甚至受殺害；有的信徒因此殉道,以致還活的信徒內心憂傷,保羅在聖靈的指示下,寫出主再來的安慰信息,叫他們不用怕,因為將來無論已死或還活著的信徒,都要被提.
\newline
\newline
還有,(林前十五50-58)講到基督再來的時候,信徒無論已死的或活著的,都得著榮耀的身體,與主復活的身體相似.這是聖經清楚的教訓.
\newline
\newline
跟著,(提後四6-8)講到基督再來是要賞賜祂的門徒.當時保羅受逼迫,處境很艱難；甚至連禦寒的衣服也沒有,書本也沒有.在最困難的時候,他所得的安慰便是從基督的再來而來.當主顯現的時候,祂要把公義的冠冕賜給凡愛慕祂顯現的人；這美好的盼望,對在受苦至深的保羅來說,該是何等大的安慰啊!
\newline
\newline
(太二十五31-46)耶穌來是要審判萬民.祂要來,坐在寶座上,審判世上的萬民,這是千真萬確的事.還有(啟二十一1-8)基督再來的時候,帶給我們有新天新地.感謝神!在新天新地中,一切都要更新,而且有義居在其中；我們可以在祂裡面得到賞賜和完全的安慰,弟兄姊妹!惟有主的再來能使我們堅忍,叫我們能堅持到底.
\newline
\newline
最後,(啟二十二20-21)「證明這事的說,是了!我必快來,阿們.主耶穌啊!我願你來.願主耶穌的恩惠,常與眾聖徒同在,阿們.」
\newline
\newline


\newpage

\section{第56屆~港九培靈會~麥希真~第六講　最有效的細胞小組}
\label{sec:391}
\textbf{-- 家}
\newline
\newline
連結: \href{https://www.hkbibleconference.org/session-message/view/391}{\texttt{https://www.hkbibleconference.org/session-message/view/391}}
\newline
\newline
\hyperref[sec:389]{< < < PREV SERMON < < <}
~
\hyperref[sec:toc]{[返目錄]}
~
\hyperref[sec:392]{> > > NEXT SERMON > > >}
\newline
\newline
家是非常奇妙的,你和我都是由家庭而產生.家給我們最大的快樂,但也給我們最大的痛苦!家給我們最大的安全感,但家也給我們最大的不安感.無論你喜歡與否,你總不能脫離家庭的關係.你可以留在家中,卻視家為酒店,旅舍；你可以脫離家庭,遠走高飛,但卻不能脫離家的感情.
\newline
\newline
今天我要按著聖經的教訓,來尋找一些關於家庭的原則,和給大家一些建議.我先提出兩個聖經的原則.
\newline
\newline
(一)是全家歸主(約一10-12)
\newline
\newline
「聽見約翰的話,跟從耶穌的兩個人,一個是西門彼得的兄弟安得烈.他先找著自己的哥哥西門.對他說:我們遇見彌賽亞了(彌賽亞繙出來,就是基督).於是領他去見耶穌,耶穌看著他說,你是約翰的兒子西門；你要稱為磯法(磯法繙出來,就是彼得).」這裡記載安得烈帶領他的兄弟信主,給我們一個很好的榜樣.今天我們同樣也可以,不單可以帶弟妹信主,也可以帶兄姊信主.
\newline
\newline
(徒十24)「又次日,他們進入該撒利亞,哥尼流已經請了他的親屬,密友,等候他們.」哥尼流帶領他的親屬密友來聽福音.四十八節「於是彼得說,這些人既受了聖靈,與我們一樣,誰能禁止用水給他們施洗呢?就吩咐奉耶穌的名,給他們施洗.他們又請彼得住了幾天.」求主給我們信心,讓我們可以帶自己的親屬,密友信主.
\newline
\newline
(徒十六14-15)講到妻子先信主,然後帶領她的丈夫,兒女信主.「有一個賣紫色布疋的婦人,名叫呂底亞,是推雅推喇的人,素來敬拜神,他聽見了；主就開導他的心,叫他留心聽保羅所講的話.他和他一家既領了洗,便求我們說:你們若以為我是信主的,請到我家裡來住,於是強留我們.」呂底亞先信主,然後帶她的丈夫和兒女全家歸主.這裡有一個很重要的條件,就是我們要有好的生活見證,生命要有改變,才能領家人歸主.
\newline
\newline
另外,同一章的二十節到三十四節有另一幅圖畫,就是丈夫先信主,然後帶領妻子和全家信主的.這就是我們所熟悉的禁卒.他因保羅的見證和引導而信主的,然後他又帶他的一家信主.這同樣需要有生命的改變和生活的見證.
\newline
\newline
第二個原則是全家事奉.(可一16-20「耶穌順著加利利的海邊走,看見西門,和西門的兄弟安得烈,在海邊撒網,他們本是打魚的.耶穌對他們說:來跟從我!我要叫你們得人如得魚一樣.他們就立刻捨了網,跟從了祂.耶穌稍向前走,又見西庇太的兒子雅各,和雅各的兄弟約翰,在船上補網.耶穌隨即招呼他們,他們就把父親西庇太和雇工留在船上,跟從耶穌去了.」這裡有兩對兄弟,他們同心事奉主,這是最快樂的事.所謂事奉,不一定要做牧師,傳道,也不一定要做執事,長老,但兄弟同心事奉主是最美好的.
\newline
\newline
(路十38-42)有另一幅姊妹同心事主的圖畫,就是馬大和馬利亞一同伺候主的事.另外,也有妻子帶著丈夫事主,或是丈夫帶者妻子事主的例子的.(徒二十一8-9)丈夫帶著妻子事主的.「第二天,我們離開這裡,來到該撒利亞,就進了傳福音的腓利家裡,和他同住.他是那個執事裡的一個.他有四個女兒,都是處女,是說預言的.」腓利很會講道,也熱心佈道；他也是執事,他有四個女兒,還未結婚,但都很會講道.腓利帶著他的一家同心事主,真是十分美的.不過,不一定每個家庭都是這個模式的,(徒十八26)就是妻子帶頭,而丈夫和兒女跟著的.「他在會堂裡放膽講道,百基拉,亞居拉聽見,就接他來,將神的道給他講解更加詳細.」當時的社會,是以男性為中心的,但提到這對夫婦的時候,聖靈准許把妻子的名百基拉放在亞居拉之先.可能是妻子的性格較堅強,所以特別蒙神使用.求主給我們堅定的信心,讓我們能影響全家歸主,並全家事奉.
\newline
\newline
跟著我給大家四個提議,讓大家參考和實行.第一,我建議年青和新婚的夫婦,進行家庭崇拜.每日舉行家庭崇拜,或者說家庭的祭壇.全家聚在一起,有短短五分鐘的崇拜.其中有四個要訣是我們要把握的；第一,讀經,但不要解經,免得嚕囌.第二,祈求,也要代禱.為自己的家祈求,又要為有需要的人代禱.第三,敬拜,但不要罵人.明指和暗示都不好,只要敬拜.第四,要準時,三或五分鐘,不要太長.讀一小段聖經,然後為一兩件事代禱.每天用些少時間舉行家庭崇拜,是領全家信主的最好辦法.
\newline
\newline
(二)常常禱告
\newline
\newline
全家人聚在一起禱告,在很多情形之下是可以實行的.我提出下面十二種情況,給你參考.第一是飯前.第二是睡覺前.教導兒女在睡覺前作簡單的禱告,認罪,求主保守一夜平安.第三是起床後,簡單的將一天交給主.第四是出門前,求主保守路上平安.第五是發噩夢時.第六是受委屈時.第七是不平時.第八是有困難時.第九是犯罪時,不但是求主赦免,更求主幫助改過.第十是責打時.第十一是疾病時,十二是死亡時,這是最有效的宗教教育,幫助兒女懂得常常禱告,也是全家歸主和事主的有效途徑.
\newline
\newline
(三)是特別對年青的弟兄姊妹講的,就是教養兒女的三步曲
\newline
\newline
孩子最好是不用打；有三件事項必須注意；第一,打之前要簡單說明理由.第二,打不是報復,打是有目的的,幫助孩子記得.第三,打完之後要有兩個保證.就是保證他一定能改好的,還要保證以後一定愛他.記得,不是靠著情緒,要用理性去管理自己,保證永遠去愛他.不單是用說話,更要用行動去再三保證.我們要靠主的恩典去實踐這幾樣.
\newline
\newline
其實,神對我們正是這樣.聖經管教我們之前,必先向我們解釋.天父的管教不是報復,而是幫助我們記得.而聖經又再三向我們保證,我們是可以改的,神又向我們保證祂以永遠的愛愛我們.耶穌基督既然愛世間屬自己的人,就愛他們到底.
\newline
\newline
最後,是特別對兒女講的.今天大部份的書,只講父母要愛兒女,卻很少講兒女孝順父母.但兒女要孝順父母,這樣就能全家歸主,全家事奉.特別我要提醒一些是熱心愛主的,不管他是青年或成年.
\newline
\newline
(箴二十20)「咒罵父母的,他的燈必滅,變為漆黑的黑暗.」要記得,咒罵父母的,前途自然黑暗.(箴二十三22)「你要聽從生你的父親,你母親老了,也不可藐視她.」我們要尊敬父母,孝順父母,千萬不可欺負藐視他們.(箴三十17)「戲笑父親,藐視而不聽從母親的,他的眼睛,必為谷中的烏鴉啄出來,為鷹雛所吃.」這節經文的意思,就是說凡藐視父母的,就是有眼無珠.求主幫助我們,不管我們是少年,青年或成年；只要父母仍然在世,都必須孝順他們.
\newline
\newline
耶穌在世的等候,他對他的父母很有孝心(路二40-52).到耶穌出來做事的時候,祂仍聽從祂母親,在迦拿的婚筵上聽從祂母親的話(約二1-12)這就是孝順.最後(弗六1-3)「你們作兒女的,要在主裡聽從父母,這是理所當然的.要孝敬父母,使你得福,在世長壽,這是第一條帶應許的誡命.」
\newline
\newline
若是我們能實踐家庭崇拜,常常禱告,教導兒女和孝順父母,我們就能帶領全家歸主,並全家事主.
\newline
\newline


\newpage

\section{第56屆~港九培靈會~麥希真~第七講　金錢讓我們正視它}
\label{sec:392}
\textbf{第七講　金錢讓我們正視它}
\newline
\newline
連結: \href{https://www.hkbibleconference.org/session-message/view/392}{\texttt{https://www.hkbibleconference.org/session-message/view/392}}
\newline
\newline
\hyperref[sec:391]{< < < PREV SERMON < < <}
~
\hyperref[sec:toc]{[返目錄]}
~
\hyperref[sec:396]{> > > NEXT SERMON > > >}
\newline
\newline
有些隨便的基督徒,雖然信了主,但仍以金錢為他們的神；他們拜金錢,拜瑪門,這當然是不好的.但另一些基督徒,卻絕口不提錢,認為講到金錢便是不屬靈,不聖潔,不清白.可是聖經中講到金子約有一百五十段經文.銀子約有二百段.錢,錢財,銀錢等約有一百段.財,財富,財利,財寶,財物,貨財等約有二百段.可見聖經對金錢並不是避諱不談,反而是常常提到,要人去正視它.
\newline
\newline
我們要看三段經文,從其中認識聖經對金錢的看法.頭一段是(路十二13-21),這段經文是講到以錢財為中心的人生.「眾人中有一個人對耶穌說,夫子,請你吩咐我的兄長和我分開家業.耶穌說,你這個人,誰立我作你們斷事的官,給你們分家業呢?於是對眾人說,你們要謹慎自守,免去一切的貪心；因為人的生命,不在乎家道的豐富.就用比喻對他們說,有一個財主,田產豐富,自己心裡思想說,我的出產沒有地方收藏,怎麼辦呢?又說:我要這麼辦,要把我的倉房拆了,另蓋更大的,在那裡好收藏我一切的糧食和財物.然後要對我的靈魂說,靈魂哪!你有許多財物積存,可作多年的費用,只管安安逸逸的吃喝快樂吧!神卻對他說:無知的人哪；今夜必要你的靈魂,你所預備的,要歸誰呢?凡為自己積財,在神面前卻不富足的,也是這樣.」
\newline
\newline
財主以金錢為人生的中心,他把錢看得太重要,這很明顯是錯誤的.財主田產豐盛,這並不是錯；他要另建新的倉房也不是錯,這是好的.但他卻錯在對靈魂所說的一番話,如果他這番話是對身體說的,那就沒有問題,可惜財主分不出身體,物質的需要,和靈魂,靈性的需要.賺錢不是錯,買樓,做生意也不是錯；這是身體物質的需要,但不是屬靈的需要.所以耶穌說人如果屬靈的生命不富足,即使家道豐富,也沒有用,也是貧窮的.如果家道豐富,屬靈的生命又豐富,這就是最好的,最有福的.但如果家道不豐富,但靈性豐富,這也是幸福的.最可憐的卻是物質豐富而靈性貧窮,這是最悲哀的.我們不要做一個無知的財主,乃要做一個聰明的財主,有了財富之後；便更加追求靈性的豐富.即使沒有錢,也要追求靈性的豐富,千萬不要以金錢為人生的中心.
\newline
\newline
第二段經文,是(路十五11-16),這是以自己為中心,不認識金錢為重要的人,這也是錯的.「耶穌又說,一個人有兩個兒子.小兒子對父親說,父親!請你把我應得的家業分給我,他父親就把產業分給他們.過了不多幾日,小兒子就把他一切所有的,都收拾起來,往遠方去了.在那裡任意放蕩,浪費貲財.既耗盡了一切所有的,又遇著那地方大遭飢荒,就窮苦起來.於是去投靠那地方的一個人,那人打發他到田裡去放豬,他恨不得拿豬所吃的豆莢充飢,也沒有人給他.」
\newline
\newline
有三個原因不得不認識金錢的重要性.第一是他用的是父母的錢.第二是他有很好的機會,他的一生平步青雲.第三是他沒有人生經驗.這樣不重視金錢的態度也是不對的.因為主耶穌在(太二十五14-30)的比喻中,給我們看到領五千和領二千的,是怎樣努力的去另賺五千和二千.但被主責備的那個,是把一千埋在地裡,這是不對的.我們要有智慧的去處理金錢.
\newline
\newline
第三段經文是(路十六1-13),這是耶穌所講不義的管家的比喻.有一個管家,因為浪費主人的財物,所以被主人辭退.他在未辭退之先,便去找那些欠他主人債的,把他們所欠下的債減寫；欠一百的,改寫為欠五十,或是八十.後來主人知道了,便誇獎這不義的管家作事聰明；因為今世之子,在世事之上,較比光明之子,更加聰明.到底這個比喻該怎麼解釋呢?
\newline
\newline
通常耶穌所講的比喻,都只包含一個教訓的.但是這個比喻卻有兩個教訓,是我們要學這不義的管家,另一是我們不要學.首先,這位今世之子,懂得把握現在,預備將來；但光明之子卻不懂得.今世之子雖然用的方法不對,但他卻有遠大的眼光；這是我們要學習的,要用現有的機會去預備將來永恆的基業.以永恆為中心,這種態度才是正確的.
\newline
\newline
衛斯理約翰講出基督徒對金錢應有的態度:第一是要盡力賺錢；第二是要盡力積蓄；第三是要盡力用錢,用在永遠的目標上.而且是要趁有錢的時候便用,要把握機會用在永恆的基業上.
\newline
\newline
比喻的下一半是教訓我們不要學這不義的管家.第十節說:「人在最小的事情上忠心,在大事上也忠心；在最小的事上不義,在大事上也不義.」這管家在小事上不忠,在大事上也不忠；在小事上不義,在大事上也不義.他的前途是黑暗的,這個我們不要學.還有另一樣,第十一節:「倘若你們在不義的錢財上不忠心,誰還把那真實的錢財託付你們呢?」人若在錢財上忠心,神便會把屬靈的事交託他；但若在普通的事上不忠,神便不能把屬天的事交託我們.我們若要為主作事,就必須在普通的事上,在錢財的事上先忠心.
\newline
\newline
最後我們來看耶穌在這個比喻中所指示我們兩個賺錢的方法.第一個是十二節:「倘若你們在別人的東西上不忠心,誰還把你們自己的東西給你們呢?」這是一個賺大錢的方法,就是在別人的錢財上忠心,這樣神就把你應得的給你.人在社會上最要緊的是信用,特別是錢財上的信用；有信用之後,你便能賺到大錢.
\newline
\newline
第二個賺錢的秘訣是十三節:「一個僕人不能事奉兩個主,不是惡這個愛那個,便是重這個輕那個；你們不能又事奉主,又事奉瑪門.」要真正得到錢財,不是要去拜錢財.嗜財如命的人,是神不喜歡的,他們必定不能賺大錢.因為做生意必然有起有落；如果嗜財如命,必定不能做,做大生意的人是拿得起,放得下的.
\newline
\newline
托爾斯泰有一個很出名的故事,就是人需要多少地?這故事講到有一個人,他起先是為別人耕地,耕別人的田,而耕田所得的又要交租.於是他心想,我要買自己的田,他便克勤克儉的積下了少許錢,買了一小塊地,內心非常滿足.但不久他又覺得這塊地太小了,於是他又拿錢去多買一塊,但不久又不滿足.如是者,他在自己鄉間已買了好幾塊地,但他總覺得那塊地都太小,沒有一塊大的.
\newline
\newline
後來有人告訴他,如果他要實現夢想,就必須到鄰村去,那裡才有大的地買.於是他便賣了鄉間所有的田,把所得錢帶在身,僱了一個僕人,一同到鄰村去.找到那村的酋長之後,酋長便解釋他這裡的買地方法,他要由清晨太陽起的時候出發,把錢放在酋長腳下,然後出去行；在日落之前走回酋長腳前,伸手摸酋長的腳趾,然後他腳所行過的地便都歸他所有.於是第二天,那人便起程了,起先他走的很快,因為他想多得土地.但是走了半天,他開始累了,而且他醒覺到自己已經走的太遠,於是便趕緊回程.誰知他的身體愈來愈倦.他早上出發的地點是在山上,而現在他卻在山下；眼看太陽就要下山,而他的體力又已經耗盡.他心想,他一生的積蓄就這樣沒有了,心中真是懊悔.但僕人在山上呼喊他說,快爬上來阿,山上的太陽還沒下山哩!那人鼓起最後的氣力,手腳井用的爬上山,最後終於摸到酋長的腳趾.酋長十分高興,便把他所行過的地賣給他；然後拿著他的錢走了.可惜那人由於氣力已盡,頭一跌便死了.他的僕人於是用鋤頭為他掘了一個坑,把他埋葬,這地只佔六乘三尺,真是可憐.
\newline
\newline
弟兄姊妹!今天我們該怎樣正視金錢呢?記著,金錢不是我們人生的目標,我們應該用永恆的角度去看金錢,並且正確的使用它.
\newline
\newline


\newpage

\section{第56屆~港九培靈會~麥希真~第八講　我們承擔了一個沉重的任務}
\label{sec:396}
\textbf{第八講　我們承擔了一個沉重的任務}
\newline
\newline
連結: \href{https://www.hkbibleconference.org/session-message/view/396}{\texttt{https://www.hkbibleconference.org/session-message/view/396}}
\newline
\newline
\hyperref[sec:392]{< < < PREV SERMON < < <}
~
\hyperref[sec:toc]{[返目錄]}
~
\hyperref[sec:398]{> > > NEXT SERMON > > >}
\newline
\newline


\newpage

\section{第56屆~港九培靈會~麥希真~第九講　成年人,還是青年人?}
\label{sec:398}
\textbf{第九講　成年人,還是青年人?}
\newline
\newline
連結: \href{https://www.hkbibleconference.org/session-message/view/398}{\texttt{https://www.hkbibleconference.org/session-message/view/398}}
\newline
\newline
\hyperref[sec:396]{< < < PREV SERMON < < <}
~
\hyperref[sec:toc]{[返目錄]}
~
\hyperref[sec:491]{> > > NEXT SERMON > > >}
\newline
\newline
很多時候我們會有一個錯覺,就是以為神只呼召青年人去做傳道,一生做傳福音的工作,而不呼召成年人.其實這種觀念是錯的.神不單呼召青年人,神也呼召成年人；一些已經成長,甚至是有家室的人.今天我們要從聖經中看神怎樣呼召青年人和成年人.
\newline
\newline
首先我們看神呼召青少年人的例子.(可十四51-52)「有一個少年人,赤身披著一塊麻布,跟隨耶穌,眾人就捉拿他.他卻丟了麻布,赤身逃走了.」在耶穌被賣的那一夜,我們知道彼得跟隨耶穌,後來還三次不認主.約翰也跟從主到十字架下.但除了彼得和約翰之外,原來還有一個少年人跟從耶穌的；他披著一塊麻布,當眾人捉拿他的時候,他便丟了麻布,赤身逃走.這少年人是誰呢?聖經沒有告訴我們他的名字,但讀聖經的人,多數都認為這就是寫馬可福音的馬可.
\newline
\newline
我們或者會覺得奇怪,為甚麼在耶穌被賣的那夜裡,有個少年人會赤著身,只披一塊麻布呢?因為當晚的天氣十分寒冷,彼得要烤火取暖.但少年人當時是在野外,為甚麼會穿得這樣單薄呢?原來聖經所用的赤身,並不是說他沒有穿衣,而是說他穿了在家穿的短衣褲；然後披上一塊麻氈.另一方面,耶穌於最後的晚餐的地點,是在馬可的家中,這是聖經學者研究前文後所得的結果,幾乎是一致確定的.吃完晚餐之後,耶穌和門徒談話,禱告,然後帶同十一個門徒,離開馬可的母親馬利亞的家,往客西馬尼園去.而馬可又跟著一同去,但他只穿了家裡所穿的短衣褲,再披上一塊麻氈便去了.後來在客西馬尼園,耶穌迫切的禱告,門徒都睡著了,馬可也不例外.到兵丁來捉拿耶穌的時候,門徒驚醒,便都四散,只有彼得和約翰仍在那裡跟著,而馬可也在其中.只是兵丁卻要捉拿他,他便逃走,連麻氈也不要,便穿著短衣褲跑回家中.
\newline
\newline
教會中有很多年青人,也像馬可一樣,很願意跟從主,但可惜在遇見逼迫時,便逃走離開主,但馬可後來看見主的恩典和神蹟,便有一個復興.(徒十二11-12)「彼得醒悟過來,說,我現在真知道主差遣他的使者,救我脫離希律的手,和猶太百姓一切所盼望的.想了一想,就往那稱呼馬可的約翰的母親馬利亞家裡去,在那裡有好些人聚集禱告.」當時在馬可的家有一個特別的禱告會,因為雅各在獄中已經被殺,彼得又被囚禁,可能結局和雅各一樣,所以信徒都聚集迫切禱告.結果彼得在天使的帶領下,平安的出來.馬可自小便有屬靈的經歷,他不單是聽父母的話,他乃是親身的經歷主,所以便把自己奉獻,一生事奉主.(徒十二25)「巴拿巴和掃羅,辦完了他們供給的事,就從耶路撒冷回來,帶著稱呼馬可的約翰同去.」馬可便開始幫忙巴拿巴和掃羅,經過一段時間的訓練後,他便跟著巴拿巴和掃羅到海外去作宣教士了.
\newline
\newline
(徒十三4-5)「巴拿巴和掃羅既被聖靈差遣,就下到西流基,從那裡坐船往居比路去.到了撒拉米,就在猶太人各會堂裡傳講神的道,也有約翰,就是馬可的約翰,作他們的幫手.」馬可很熱心作海外宣教士.但可惜到十三節,保羅和他的同人,從帕弗開船,來到旁非利亞的別加,馬可就離開他們到耶路撒冷去.」這是馬可第二次的逃跑.後來主的愛激勵他,他又回轉過來了.
\newline
\newline
不過,馬可的退後引起教會有極大的事.(徒十五35-39)「過了些日子,保羅對巴拿巴說,我們可以回到從前宣傳主道的各城,看望弟兄們景況如何.巴拿巴有意,要帶著稱呼馬可的約翰同去.但保羅因為馬可從前在旁非利亞離開他們,不和他們同去作工,就以為不可帶他去.於是二人起了爭論,甚至彼此分開,巴拿巴帶著馬可,坐船往居比路去.」因著馬可一個年青人,引致巴拿巴和保羅,起了爭論.雖然這樣,神仍使用馬可.
\newline
\newline
後來,保羅發現馬可大有進步,便開始接納他.只要我們的生命有改變,一定被接納.(西四10)「與我一起坐監的西里達古問你們安.巴拿巴的表弟馬可也問你們安.(說到這馬可,你們已經受了吩咐,他若到你們那裡,你們就接待他)」以前保羅曾責備馬可,不肯與馬可同工,但現在他的態度已完全改變了.不但如此,(提後四9)保羅有更好的稱讚,他對提摩太說:「你來的時候要把馬可帶來,因為他在傳道的事上於我有益處.」不單保羅,彼得也有稱讚馬可的話,(彼前五13)「在巴比倫與你們同蒙揀選的教會問你們安.我兒子馬可也問你們安.」馬可實在成為一個神所重用的人.後來聖靈更感動他寫成馬可福音,這是四本福音書中最早寫成的,別的福音書也參考他的著作.
\newline
\newline
我們千萬不可對自己灰心；主要繼續使用我們,成為祂合用的器皿.
\newline
\newline
神不單呼召年青人.(徒廿一8-9)「第二天,我們離開那裡,來到該撒利亞,進了傳福音的腓利家裡,和他同住.他是那七個執事裡的一個.他有四個女兒,都是處女,是說預言的.」腓利有家室,有四個女兒,都是未婚的.她們雖然年輕,卻很會講道.腓利因為熱心傳福音,帶人信主,而且很會講道,特別是佈道.所以他便放下自己的事業,以後專心傳道,還影響祂的四個女兒也熱心傳道,在各種場合為主證道.
\newline
\newline
腓利是怎樣出身,怎樣奉獻作傳道的呢?(徒六2-7)告訴我們,腓利本來是一個信徒,後來教會請他作執事,幫助打理飯食的事.他十分謙卑,各樣的事務都做,後來他的恩賜便顯露出來了.腓利至少有兩樣的恩賜,就是講道和個人佈道.
\newline
\newline
我們先看他講道的恩賜.(徒八4-5)「那些分散的人,往各處去傳道.腓利下撒瑪利亞城去,宣講基督.」腓利很有口才,公開講道,他解釋道理十分詳細明白,很多人因此信主.如果你有講道的恩賜,或許神要呼召你.
\newline
\newline
腓利還有個人傳道的恩賜.(徒八26-39)講到腓利怎樣得勝靈引導,去向一個埃提阿伯的太監傳福音,結果帶領了這位埃提阿伯的財政部長信主.一個人有個人佈道的恩賜,可能就印證了他是被主呼召的.只會站在台上高聲講論,而不會作個人佈道的,很難被主使用.個人佈道是非常重要的一環,真正愛主,又愛人靈魂的人,必然有個人佈道的熱誠和效果.今天我們又有向人傳福音,有帶領人歸信主嗎?如果有的話,我們可以在主面前等候,求主呼召,使用.
\newline
\newline
年紀大並不是問題,夫婦同心事主是非常美好的,如果兒女們已長大,可以徵詢他們的意見,請他們也一同同心的事主.這樣,你就可以勇敢的踏上十字架的道路,一生專心的事奉主,放下你的事業,進入神學院,接受造就,然後出來被主使用,有人五十多歲也奉獻自己,夫婦同心的進神學院,後來被主使用.
\newline
\newline
腓利有講道的恩賜,又有個人佈道的恩賜.當聖靈吩咐他到迦薩的路上傳福音的時候,他便放下一切,跟從主去了.聖經稱腓利為「傳福音的腓利」,現代的稱呼就是佈道家.
\newline
\newline


\newpage

\section{第57屆~港九培靈會~滕近輝~第一講　詩篇中七幅復興的圖畫}
\label{sec:491}
\textbf{第一講　詩篇中七幅復興的圖畫}
\newline
\newline
連結: \href{https://www.hkbibleconference.org/session-message/view/491}{\texttt{https://www.hkbibleconference.org/session-message/view/491}}
\newline
\newline
\hyperref[sec:398]{< < < PREV SERMON < < <}
~
\hyperref[sec:toc]{[返目錄]}
~
\hyperref[sec:492]{> > > NEXT SERMON > > >}
\newline
\newline
感謝主!再次給我機會在港九培靈研經會中和各位弟兄姊妹分享主的話語.這次我們要一同領受詩篇中的信息.卅四年前我首次被邀在這大會講道,所講的題目就是詩篇；此次我又是講詩篇；但內容和那次完全不同.
\newline
\newline
願我們以預備好的心,愛主和愛主話語的心,一同來領受主的信息.
\newline
\newline
詩篇是很特別的一卷書,如果聖經沒有詩篇,我覺得是個很大的缺憾；因為在詩篇中給我們看到許多屬神者的心聲,是他們生活經驗的見證.使我們讀時常有同感,有共鳴!似乎我們要說的話,詩人已經為我們表達了.詩人的經歷和我們的經歷,多有吻合之處,所以讀了感到入心；因此,全世界的基督徒都喜讀詩篇.我們要領受詩人們屬靈的經歷,藉此得著勉勵,提醒,幫助及指引.
\newline
\newline
詩篇中第一幅復興的圖畫
\newline
\newline
時間考驗之下的復興(一O三5)詩人以老鷹羽毛豐滿的圖畫,描寫屬靈的復興；因鷹的特點就是年老時,全身羽毛換新,所以老鷹羽毛豐盛.
\newline
\newline
心理學有一個定律,人對於所熟悉之物,會逐漸失去興趣或減少興趣,愈熟悉興趣愈低.生理學上也有種現象,人工作愈久愈覺疲倦.當我們信主事奉主的年日長久,我們的熱心和屬靈的興趣就慢慢減低,負擔也減少了.於此情況下,我們需要復興.詩人以此提醒我們,願我們信主長期以後,屬靈的羽毛又豐滿起來.
\newline
\newline
約翰衛斯理臨終之言:「最好的是神與我們同在.」他到了八八高齡心中仍充滿神的同在.
\newline
\newline
北美咸美頓大學創辦人,他也是首任校長,他非常愛主,該大學造就了許多愛主獻身的青年.這位校長年老病了,當醫生坦言只有半小時的壽命,他就吩咐親人們都圍在病床邊,說:「請你們合力把我扶起跪在床上,我要禱告.」他用有生之年最後半小時,為世界歸主禱告,這是他最關心的事,這種關心一直在他心靈中到了最後的一刻.這是何等寶貝!
\newline
\newline
我們信主事奉主有多久?有沒有在時間考驗下慢慢消沉?若有,今晚主向你發聲.願主感動你幫助你,叫你復興起來.在我們奉獻的路上,每個人有奉獻的開始；那起點就變成一條線,也就是奉獻的路,一直走在這路上；然後成為全面的奉獻,不但是件事,是種工作的奉獻,乃是整全一生的奉獻.在這條奉獻的線上又有許多點,每個點是復興的點,重新,重新,重新奉獻；每次的奉獻是一次的復興,如此一直到底.就如李文斯敦臨終前的日記所寫:「耶穌基督我的王,我的主,我的神啊!我再次將自己奉獻給你.」他的奉獻是繼續的奉獻,是條線貫徹始終.你的奉獻如何?
\newline
\newline
第二幅復興的圖畫
\newline
\newline
南地的河水復流(一二六4)南地則下游地帶,北方是上游地帶是山地,山上的水形成河流南下.詩人說:「耶和華使被擄的人歸回,好像南地的河水復流.」復流之前,曾經停流,如今重新流動.停流期間情況淒涼,河床幾乎乾涸,兩畔土地樹木枯萎；當河水復流,青翠而欣欣向榮的景況重新出現,真是一幅復興的圖畫!河水停流通常有兩原因,第一因上流缺少雨水供應,乾旱情況致下游停流.我們的靈性也是如此；若我失去上面的供應,屬靈天雨不下,靈性源頭就枯乾,與神關係就枯乾,心靈少了滋潤,生命的河就停流了.所以我們必須保持上面溝通的關係.好像一輛電車開行,力量來自上面的電線,必須緊密連接；一旦斷了電線,電車立即停頓；電流供應實在是不可或缺的.照樣,我們與神的關係,是絕對少不得的；靈修是我們基本的訓練,永遠不可缺少.請看世運的選手們.當比賽前後不斷跑步,體操訓練一直不停,經常保持基本的訓練.學彈琴的人,不論技術多高,都需要經常鍛練手指,不可停止.靈修是我們的基本操練,不可停止；否則生命河流必然停頓.第二原因,河的本流被支流吸取,支流增多以致本流乾旱.
\newline
\newline
詩人說,以色列被擄到外邦,受敵人控制；現在神釋放他們重返本國,過自由生活,享受神兒女的榮耀.當他們被擄時,河流就停頓了.
\newline
\newline
當基督徒被撒旦擄掠,利用,綑綁時,生命的江河停流；若要復流,必須從敵人手中歸回,把肢體收回,不再被敵人,被罪惡利用,不再作被擄的基督徒；使撒旦在我們身上無所作為.收回之後,完全奉獻給神,作義的器皿被神使用.這樣,生命之河復流了.有時我們在生命裡,在生活中,產生了對生命力消耗的事,如同支流把我們生命本流基本的作用吸取；致水流減少甚至乾旱停流,所以我們要把這些都收回來.人生的焦點乃是先求神的國和祂的義.發揮生命力,被神使用；生命如活水的江河,能夠祝福別人.
\newline
\newline
第三幅復興的圖畫
\newline
\newline
裂口得醫治的復興(六○1-2)這段經文是詩人的禱告.他求神復興自己也復興以色列人.在字面上看,好像裂口是神造成的；因為經文說:「你使我們破敗,你向我們發怒......使地震動而且崩裂.」事實上,裂口產生乃因以色列人違背神,神的管教臨到他們.真正的原因,裂口是人造成的.
\newline
\newline
詩人大衛醒悟,代表以色列民向神禱告向神呼求:「神啊!求你醫治我們的裂口.」因主三次問他是否愛祂,但我相信還另外有理由.主知道彼得不但需要主一次地問他提醒他,乃是多次的需要,弟兄姊妹們!你我也是一樣,我們需要主多次地問,我們一次,再次地向著你所作所為有甚麼是神不喜悅的呢?你與神的關係如何?你是否正在違背神嗎?如有這樣的裂口,今晚是你復興得醫治之時；如果聖靈光照你,你和神之間有任何隔膜,攔阻,求主幫助把這些除去.如果你和弟兄姊妹間有任何隔膜,也求主助你決志與他們和好.很多時候,裂口產生乃因單方面的錯,對方就堅持必須錯方認錯並道歉；可是很希奇,在聖靈工作的實例中,常是對的一方主動要求和好；聖靈的感動使他倚靠主,與對方和好.在和好的事上,裂口的醫治上,不可講「理」,只可講「愛」,講理永遠不會和好；因雙方都認為自己有理,我們應超越「講理」進入「講愛」的境界.那時我們就放下一切的「理」和「權利」而謙卑追求和好；惟有神的愛澆灌在我們心中,裂口才能得醫治,基督徒榮耀的和睦才會產生.
\newline
\newline
第四幅復興的圖畫
\newline
\newline
「愛慕神」的復興(四二1)這是幅非常美麗的圖畫!快跑的鹿,體內水份大量消耗,因乾渴體力衰竭致行動緩慢；但到了溪水旁盡量喝水,精神振作,力量恢復!再接再厲,往前跑跳,得到復興.詩人說:「神啊!我的心切慕你,如鹿切慕溪水.」我們需要愛神的復興!
\newline
\newline
主耶穌復活之後,曾三次問彼得:「你愛我嗎?」我相信因彼得曾經三次否認主；所以為這裂口影響很大,很嚴重,詩人懇切向神祈求要得復興.
\newline
\newline
我們與神之間,與弟兄姊妹之間,有時會產生裂口,需要醫治.你和神之間有裂口嗎?主說:「主啊!是的,你知道我愛你.」每次我們這樣回答,就是愛主的復興.有時我們日常工作忙碌,忘記了主；以事奉工作代替主自己.我們應當以最高的愛獻給主,應時常挑旺裡面對主熱切的愛.在新約有三次愛的問題,就是主三次問彼得.舊約有三次神愛的呼聲,在耶利米書第三章,「背道的以色列阿回來罷.」(12)以色列民與神背道而馳,神呼召這樣的孩子回來,這是父親的呼聲.第二次的呼聲,是作丈夫的神發出的.(14)第三次的呼聲是耶和華我們的神發的.(22)我們的父,我們屬靈的丈夫,我們神耶和華.祂向我們一而再,再而三地發出愛的呼聲說:「遠離我的孩子們阿!回來罷!」我們聽到了慈愛的呼聲,聽到創造我們的偉大之神的呼聲,聽見教會新郎基督的呼聲；都是從愛裡發出來的.感謝神!有不少弟兄姊妹在冷淡退後時,再次想到神的愛,又回來,復興起來了,回到神面前,如鹿到溪水邊暢飲溪水.主等候我們重返祂的面前,盡量享受祂永恆的慈愛.
\newline
\newline
法國歷史中的路易九世皇帝,是個特別的王,他外號叫「聖徒之王」,因當時法國宮庭充滿淫亂；但這位王「情所獨鍾」的只是他的王后.結婚時他造個戒子戴在手上,刻了「神」「法國」「馬嘉烈」(王后的名)三字,他敬拜神,愛國家,愛自己的王后,他一生保持對王后真純的愛.這是法國宮庭僅有的例子,戒指常提醒他保持崇高完整真純的愛.在我們心中也應刻上「愛」字,我們愛神,愛主,永遠愛祂,把戒指戴在我們屬靈的手指上.
\newline
\newline
第五幅復興的圖畫(四十2-3)
\newline
\newline
淤泥是快沙,若失足其中,霎時使人深陷.詩人說:「當我落在淤泥的境況.神慈愛的手把我拉上來,救出來,然後使我的腳立在磐石上.我口唱新歌.」詩人所唱的是見證之歌,讚美之歌和感恩之歌.詩人大衛的一生中,實在曾經沉溺於情慾的淤泥中.一個靈性美好蒙神大慈愛的人,竟然墮落至此地步,實令人難以置信,但事實如此.感謝神!當他墮落往下沉之時,神伸展慈愛的手抓住他,把他拉上來；且使他的腳穩立磐石上.這時他醒悟悔改,回復與神的關係.大衛所唱見證讚美感恩之歌,許多人聽了都懼怕；敬畏的懼怕,他們從大衛的經歷和悔改中發現了神；產生敬畏神的心,於是回轉倚靠耶和華.罪人悔改,非但自己蒙恩,悔改的見證也幫助別人歸向神.
\newline
\newline
今日的世代,基督徒要特別小心!啟示錄十六章13節說:「一個污穢的靈好像青蛙,從龍口,獸口,並假先知的口中出來.」青蛙是污穢的癩蛤蟆,從撒旦出來,污穢的靈到普天下去.末世時候,污穢的靈充滿天下,充滿社會,叫人落在污穢的情慾裡面.時代的預言我們所見實際如此.還有(啟十七4)載:一個淫婦騎在獸背上,身穿華麗衣服,手舉金杯,那金杯中是淫亂的污穢.表示以淫亂為榮耀,這就是末世的特色.每時代都有淫亂的事；但今天卻比任何時代有所不同,今天的人將淫亂的事誇張,更引以為榮,再不是偷摸地去作.這樣的世代,使基督徒生活的環境,多麼容易受到影響,特別是青年人,不知不覺就沉溺其中,失去屬靈的判斷力,漸漸下沉,.「慾」上左是「谷」,深谷,右邊是「欠」；不足,情慾如深谷,永不滿足.下面是「心」,是清楚地指出,我們很容易沉溺淤泥中.感謝主!大衛起來了!我們也要像大衛悔改,徹底悔改像詩篇五十一篇那樣認罪.在神前那憂傷痛悔的心,神必不輕看.
\newline
\newline
第六幅復興的圖畫
\newline
\newline
事實上的復興(一○四30)神發出祂的靈,地面就更新.這是聖靈的工作.「地面更新」有兩個意思；當初神創造天地,萬物從空虛混沌顯出美麗的大自然,每年春天來到大地回春萬象更新,一切欣欣向榮.冬天過了,許多動物經過冬眠出來,恢復牠們的生命活力,整個田野美麗起來.春天和冬天的情況完全不同,正像復興的境況,我們的心靈也是這樣.有的人在事奉上進入冬眠狀態,雖然得救；但得救的生命沒表現,沒結果子,沒工作；當復興來到,進入春天,生命力活潑表現.何等美麗!
\newline
\newline
聖經記載馬可的故事,最初,保羅和馬可一同出去,到了居比路島；因前路危險,馬可就離開保羅.因此令保羅對他不滿,以至日後不再用他.但是後來馬可復興,保羅在(西四10-11)和(提後四11)就有稱讚馬可的話,說他需要馬可,馬可是他的好同工,和他一齊在患難中同工,如果你到羅馬,有個聖馬可堂,堂門口上雕刻有翅膀的獅子,象徵馬可.復興前的馬可絕不像獅子,但復興後的馬可,如獅插翼,真是美好復興的表現!
\newline
\newline
第七幅復興的圖畫(八四6)
\newline
\newline
當人經過流淚谷時,心情必然痛苦,消沉,否定一切；成為悲觀者,發怨言者,心靈破碎.但是在此過程中,神的恩典臨到,產生變化,流淚谷變成泉源之地；所流的眼淚成了秋雨,帶來秋雨之福蓋滿全谷.秋雨是與收成有關的雨,雖離收成時節遙遠,一時看不出來.很多人正當流淚時,人們看不出他將有秋雨之福；事實上秋雨,乃是與收成絕對有關的雨.神讓眼淚變成秋雨.人在痛苦中漸漸消沉,需要復興!有的人卻被痛苦打倒,不斷抱怨.
\newline
\newline
蔡蘇娟姊妹在不久之前去世,這位暗室之后享年九十四,大半生在疾病之中.當疾病突然臨到她時,在主扶持下她的眼淚真是變成秋雨,這秋雨就使她幫助了許多人；她有豐富屬靈的大收成,她的著作譯成數十種文字,在許多國家幫助了無數的人.感謝主!這畫是真實的!
\newline
\newline
今天你是否在流淚谷?你需要復興!但願聖靈將神的愛再次澆灌你心,請你注目在各各他的十字架,那是個最清楚的證明.神的愛是永恆的愛,十字架的愛一進入你心之時,一切傷口都得到醫治.
\newline
\newline


\newpage

\section{第57屆~港九培靈會~滕近輝~第二講　詩人奇妙的經歷}
\label{sec:492}
\textbf{第二講　詩人奇妙的經歷}
\newline
\newline
連結: \href{https://www.hkbibleconference.org/session-message/view/492}{\texttt{https://www.hkbibleconference.org/session-message/view/492}}
\newline
\newline
\hyperref[sec:491]{< < < PREV SERMON < < <}
~
\hyperref[sec:toc]{[返目錄]}
~
\hyperref[sec:493]{> > > NEXT SERMON > > >}
\newline
\newline
詩人的經歷,常常也是我們的經歷,所以我們要向他們學習.求保惠師聖靈在我們心中教導我們.神的作為很奇妙,常超出人意料之外.先知以賽亞說,神的道路高過我們的道路,神的意念高過我們的意念.我們常用個人的思想推測神,結果就出了差錯；我們屬靈功課重要的部份,就是學習用神的眼光來看一切的事.歷世歷代以來有許多屬靈的人學了這個功課,在他們的經驗中留下榜樣；讓我們讀聖經從中一同學習.「神用手中的巧妙引導他們.」(七八72)神的手是巧妙的,超過我們所求所想來作工,來引導我們.所以我們走靈程,要定睛主引導的手；不是注目在自己身上,環境上；也非別人身上,乃是注目在主引導巧妙的手上.我們要預備隨時發生奇妙的事,期待神奇妙的作為.詩人有何奇妙經歷,茲按時間許可讀幾處經文.
\newline
\newline
第一個奇妙的經歷,是詩人大衛的經歷(廿三5)神在敵人面前,為他擺設筵席.敵人面前和筵席,正是相反的情況；敵人面前是危險,打仗,沒平安之地；筵席是喜樂,豐富,享受之地.在神奇妙作為中兩種相反的情況連起,這是超出人意料之外.如有人常和你作對,常批評挑剔你的錯；你有福了,若你態度正確,會產生一顆珍珠；好像蚌體中的一粒沙,使牠很不舒服；就產生分泌物把沙子包圍,結果成了一顆珍珠.從前英國的格來斯頓首相,他有個特別作風；如果有人對他不好,他就特別待那人好；因此有人利用這個特點,想引他注意就特為對他不好,藉此得了好的機會.感謝神!祂的生命在我們裡面,讓我們能體會這道理,學習這功課,使我們在敵對的環境中蒙恩.記得我十二歲時,初次聽宋尚節博士講道；他手中提個裝著皮球的小籃,他拿了一個擲落地上,球跳起來.接著又拿一個更用力地拋,球跳得更高；連續一個比一個更加用力,球跳的更高.教會就是這樣,越受逼迫,越向上升,靈性越蒙恩,教會越興旺；這是事實.兩千年的教會歷史證明這事實.七年前我閱華盛頓郵報,有篇整版的報導,是位加拿大記者寫的,他到蘇聯逗留六個月,實際參察教會情況,根據他實際觀察,蘇聯基督徒人數約在三千萬至六千萬之間,約佔全國人口之1/6或1/3.我一讀之下,不敢相信；但這是一位知名的記者說的,他不可能隨便報導.真是如此嗎?神能行奇妙的事.不久前,有位主僕告訴我:福建省其中一縣,有十萬基督徒,另一城市有三間開放的教會,每主日有六至八千人聚會.去年那三間教會,約有二千人受洗,此外還有許多家庭教會,和那三間開放的教會有合作.聖靈很明顯在工作.
\newline
\newline
詩人第二個奇妙的經歷(三五13)詩人說敵人對他很不好,但是當敵人患病時,他身穿麻衣禁食刻苦己身地為他禱告,以愛心對待他.這是稀奇的事,是美麗奇妙經驗的表現；他愛仇敵,因神的愛在他心裡.神就是愛,當人真實認識神時,與神之間有生命連繫時；神的愛充滿在他裡面,他才能以愛報惡.感謝主!祂在十架上為釘死祂的人禱告；主的榜樣又在司提反的生命中重新表現.司提反被猶太人用石頭打的時候；他跪下禱告說:「主啊!不要將這罪歸到他們身上.」這正是我們的主在十架上所表現的,主充滿在司提反心中,所以他也有同樣的流露.舊約記載,約瑟抱著那些出賣他的哥哥們的頸項,流出饒恕的眼淚.約瑟是個認識神的人,聖經說他心裡有神的靈.神的靈將神的愛澆灌在他心裡,所以他以神的愛,愛那對不起他的哥哥們.在我們心中,有敵對的人嗎?或有你所不能饒恕的人嗎?若有,今晚禱告求主用祂的愛充滿你心,使你以主的愛徹底饒恕他；你自己做不到,但當主愛充滿你時,你就能做到.
\newline
\newline
詩人第三個奇妙的經歷(六六12)水火是試煉之地,所謂水深火熱.當我們有這經歷時,神巧妙的手就把我們帶到豐富之地,叫我們的靈性豐富,這是超出人意料的事.
\newline
\newline
兩年多前,我到美國以利諾州一間教會講道,春田是林肯總統的故鄉,他的墳墓也在那裡.林肯紀念碑是在華盛頓,從中我得知林肯的遭遇,他一生至少有八次重大的打擊.茲舉二例:他開始工作時與人合謀生意,受那人欺騙,以致要用十八年的時間還債.整個人生有幾個十八年呢?第二件事,他很愛一女子,當預備結婚之前,那女子死了,他非常痛苦.在這些試煉中,他用信心面對一切,結果,他學了美好的功課,得了造就,更加堅強；作好來日美國偉大總統的準備.他學了真正屬靈的見識,他帶領美國進行黑奴解放戰爭；他定了一天,使全國為解放黑奴戰爭禁食禱告仰望神.曾有人對他說:「你應求神站在我們這邊.」他答:「你怎能這樣說,我們只可求神助我們站在祂的方面.」他對真理,對神的認識很清楚.台灣的宇宙光不久將出版林肯總統傳記,由葛培理博士寫序文.在林肯總統一生中,我們看見一個有信心的人,對困難的反應.在我自己的經驗中,大概廿年前,我體弱多病,那時我忽然病倒床上三個月,表面看來是件很不幸的事；但是,感謝神!自從哪次病愈之後,身體健康,很少害病,身心大衛進步.今天回想,我仍不明白為何那次的病使我有健康絕大的轉機.神帶領我們經過水火,也帶領我們進入豐富之地.在印尼椰城的宣道堂,有一次,禮拜堂火燒,弟兄姊妹很難過,有的懷了疑惑的心,以為神連自己的禮拜堂都不能保守；信心就搖動.但是牧師很有信心,他帶領一班有信心的弟兄姊妹禱告,仰望,倚靠神；不久之後,他們建立了比前大一倍的新堂.今天聚會人數增加了三倍,每主日有四次聚會；同時還有兩處分堂,且另有差傳工作.他們經過火,而進入豐富之地；因他們認識神倚靠神,所以得到神奇妙作為成全在他們身上.
\newline
\newline
詩人第四個奇妙經歷(一一四8)在神巧妙手中,堅固的磐石變成泉源,有水流出,這是神的作為.我們曾見過心腸堅硬如石的人改變信了主.去年有個聚會在九龍塘宣道會舉行,會中有互愛團契和警察團契的詩班聯合獻唱.散會後有人稱讚是日獻唱的,是貓和老鼠的詩班；原來警察是要捉那些吸毒者；但兩個團契竟合唱讚美主!吸毒者聽了福音悔改,成了在基督裡新造的人,有新的認識,而這班信主的警察們,他們常常聚會,互相勉勵,很有愛主的心；他們一齊唱詩讚美主的拯救.這是主奇妙的工作.數年前有六間循道會聯合在九龍舉行培靈會請我講道,那晚我到達時,和六間教會同工們見面──握手；他們告訴我,其中有三位已往是吸毒者,後來信主悔改而且很愛主,作了傳道人.這是奇妙的見證.耶穌基督能改變人心,把人從罪惡痛苦中拯救出來,帶領人走上光明跟隨基督的路；並且也帶領別人走上這條路.這是何等美好!在加拿大溫哥華有次聚會之後,有個青年弟兄對我說:「滕牧師!你今晚講的我心裡有共鳴,我可以和你談談嗎?」於是我們往外面找個安靜處談話,他說當他讀港大時,他是最積極反對基督教的；後來認識了主,完全改變成為信主的人,且是個相當熱心的基督徒.有不少這樣的人,認識主之前極力強硬地反對,如同堅石一樣；但認識主後完全轉變,敬畏這位偉大慈愛的主.所以我們傳福音時,不要怕!自己的話雖沒用；但真誠的見證加上聖靈的工作就有用,就有效果.「這福音本是神的大能,要拯救一切相信的人.」我們當放膽傳福音,因主是可信的.兩千年前主曾說:「我將我的教會建立在磐石上,陰間的勢力不能勝過她.」兩千年來,這話得充份的證明,我們知道所信的是誰,我們當放膽剛強地為主作見證.主能使剛強的人改變.
\newline
\newline
詩人第五個奇妙的經歷(一一三9)神奇妙的作為使不能生育的婦人成為多子的樂母.詩人提至此,說:「你們要讚美耶和華.」撒母耳的母親不能生育,她向神懇切禱告；神垂聽了,將撒母耳賜給她,這兒子而且成為以色列傑出的領袖.在肉身方面,神能行這神蹟,在靈性方面,神更能行此神蹟.在香港有對夫婦沒有孩子,禱告了廿年也沒有孩子；有一天,婦人對她丈夫說,今天主給我個新的觀念,新的感動,叫我感到我們應辦個孤兒院,把我倆的愛擺在那些孤兒身上.她丈夫說正有同感,於是他們辦了間孤兒院,把精神力量心血,都擺在那些沒有父母可憐的孩子身上.這對夫婦有了廿幾個孩子,孩子日漸長大,有的出院；不斷有新的孩子入院,他們竟然有一百多個孩子.我相信誰都沒有見過有百多個孩子的家庭.感謝天父!他們禱告廿年沒有落空,神所賜超出他們所求所想的,神的道路高過我們的道路.所以我們當凡事交托,凡事仰望.
\newline
\newline
說話至此,我想到(賽五四1-3)這段經文和不能生育的婦女,得到孩子有關,這是歡聲的歌唱,是信心的歌唱.先知說:「你這不懷孕......未曾經過產難的要歌唱.」信心的歌唱；神要賜下很多的兒女,而且這些兒女要擴展神的地方.第二節是差傳運動的經句,二百年前神藉這節經文,在英國感動了威廉加爾；他開始用信心仰望神,推動差傳工作,把福音傳到遠處.他常傳講:「要擴展你帳幕之地,張大你居所的幔子,不要限止,要放長你的繩子,堅固你的橛子,因為你要向左向右開展；」(賽五四2-3)為了神而擴展.當以色列人聽見先知以賽亞這個信息,他們有兩種反應,一是嘲笑以賽亞的話；另外的是知道這是從神而來的信息,他們恍悟以往不擴展,如今要按神的呼聲來響應,要用信心來歌唱,以信心的步伐踏上這條路.今天也有這兩種反應.求神給我們更多的教會,能唱這差傳信心的歌.
\newline
\newline
最近有位弟兄獻身,他已畢業神學,再受訓兩個月,不久將出發到菲律賓土人地帶傳福音,這位弟兄多年患肺病,主醫治了他；他奉獻要到那偏僻的地方傳福音.那裡是荒山野嶺,出門要爬山,要學土語,無水電供應；但主成為他的力量,他知道前路艱難,但他毅然,決然憑信心走上去.相信主必大大用他,雖家人反對；但他把自己交給主,主能成全對他的呼召.感謝主!有人肯出去.希望繼續不斷有更多弟兄姊妹憑信心唱這歌.
\newline
\newline
詩人第六個奇妙經歷(九二14)年老的人仍有表現,滿有生命的汁漿,仍然發青,仍然結果子；年雖老生命力仍然旺盛.(12-13)是他們年老仍充滿生命的秘訣,他們如同棕樹和高山上筆立的香柏樹.他們栽於耶和華的殿中,發旺在神的院裡,與神有密切的關係,接近神；一直保持自己在神前,所以繼續滿有生命力.亞伯拉罕約在一百一十七歲時獻以撒,他在高齡時達到愛神的高峰.但以理九十多歲時在波斯王手下,他打開窗子面向耶路撒冷,每日三次禱告,素常一樣.請看這兩位老人家,他們都保持信心,愛神的心,儆醒的心,直到年老.青年的同道熱心,我非常快樂感謝主；但不知你們愛主熱心能持續多久?你們未曾受過考驗,三年後,結婚後,工作之後還熱心還愛主嗎?還會事奉主嗎?還會傳福音嗎?這是未知數.當然神能保守,但是我們的心志很重要.迦勒說:「神從前打發我出去時我四十歲,現在我八十五歲,四十五年如一日,那時我力量如何,現在我力量仍是如何.」今天教會需要這樣的人,站得住,能持守熱心愛心.「熱」字的寫法很有意思,上「埶」保持,下四個點代表火.水壺如果一直保持在火上,壺中水一直是熱的,因靠近火源.當人靠近主,主能保守,到年老之時仍然結果子,仍然滿了汁漿常發青,發旺在神的院子裡,因栽於神的殿中.
\newline
\newline
(林後八22),保羅作見證說:這位弟兄過去熱心,現在更加熱心；在許多事上屢次試驗他,不僅過去熱心而且更加熱心.為甚麼?保羅說:「現在他因深信你們就更加熱心了.」在他的熱心裡面有信──深信,對弟兄姊妹深信；因他對神有信心,信得過神也信得過屬神者,他相信為神而作不落空,為神的兒女作不落空,到了時候要收成,結出美好的果子,神的旨意要成全,神的時候到了,一切都要成全.感謝主!因為他有信心,所以他有熱心；真正的熱心離不開信心,沒有信心的人會失去熱心；我們開始,信心會開始,我們靈性的高峰是信心.我們靈程的經過還是靠信心；以信心為開始,為過程,為結束,由信一直有信.感謝神給我的信心,我們接受信心,在神恩典中,我們領受又領受；有了信,就必能看見神的榮耀.
\newline
\newline
主啊!我的信不足,求你幫助!
\newline
\newline
我們在跟從主上,在事奉上,在愛主的心志上,需要堅定的信心來支持!我們信主也知道所信的是誰,和主之間有生命的關係,有了一切的連繫,堅定的信心把我們和主連起來,保持我們熱心,愛心一直到底.
\newline
\newline


\newpage

\section{第57屆~港九培靈會~滕近輝~第三講　詩篇中關於事奉的信息}
\label{sec:493}
\textbf{第三講　詩篇中關於事奉的信息}
\newline
\newline
連結: \href{https://www.hkbibleconference.org/session-message/view/493}{\texttt{https://www.hkbibleconference.org/session-message/view/493}}
\newline
\newline
\hyperref[sec:492]{< < < PREV SERMON < < <}
~
\hyperref[sec:toc]{[返目錄]}
~
\hyperref[sec:494]{> > > NEXT SERMON > > >}
\newline
\newline
詩篇 90:6-90:17
\newline
\begin{longtable}{cl}
\hline
\hline
章節 & 經文 (和合本修訂版)\\
\hline
90:6 & \begin{tabularx}{0.7\textwidth}{X} 早晨發芽生長，晚上割下枯乾。 \end{tabularx} \\ \\ \relax
90:7 & \begin{tabularx}{0.7\textwidth}{X} 我們因你的怒氣而消滅，因你的憤怒而驚惶。 \end{tabularx} \\ \\ \relax
90:8 & \begin{tabularx}{0.7\textwidth}{X} 你將我們的罪孽擺在你面前，將我們的隱惡擺在你面光之中。 \end{tabularx} \\ \\ \relax
90:9 & \begin{tabularx}{0.7\textwidth}{X} 我們經過的日子，都在你震怒之下，我們度盡的年歲，好像一聲嘆息。 \end{tabularx} \\ \\ \relax
90:10 & \begin{tabularx}{0.7\textwidth}{X} 我們一生的年日是七十歲，若是強壯可到八十歲；但其中所矜誇的不過是勞苦愁煩，轉眼即逝，我們便如飛而去。 \end{tabularx} \\ \\ \relax
90:11 & \begin{tabularx}{0.7\textwidth}{X} 誰曉得你怒氣的權勢？誰因著敬畏你而曉得你的憤怒呢？ \end{tabularx} \\ \\ \relax
90:12 & \begin{tabularx}{0.7\textwidth}{X} 求你指教我們怎樣數算自己的日子，好叫我們得著智慧的心。 \end{tabularx} \\ \\ \relax
90:13 & \begin{tabularx}{0.7\textwidth}{X} 耶和華啊，我們要等到幾時呢？求你轉回，憐憫你的僕人們。 \end{tabularx} \\ \\ \relax
90:14 & \begin{tabularx}{0.7\textwidth}{X} 求你使我們早早飽得你的慈愛，好叫我們一生一世歡呼喜樂。 \end{tabularx} \\ \\ \relax
90:15 & \begin{tabularx}{0.7\textwidth}{X} 求你照著你使我們受苦的日子，和我們遭難的年歲，使我們喜樂。 \end{tabularx} \\ \\ \relax
90:16 & \begin{tabularx}{0.7\textwidth}{X} 願你的作為向你僕人們顯現，願你的榮耀向他們子孫顯明。 \end{tabularx} \\ \\ \relax
90:17 & \begin{tabularx}{0.7\textwidth}{X} 願主－我們神的恩寵歸於我們身上。願你堅立我們手所做的工，我們手所做的工，願你堅立。 \end{tabularx} \\ \\
[1ex]
\hline
\hline
\end{longtable}
剛才詩歌所唱的「事主真上算」上算兩字,我不大喜歡；因為字眼不大好,但是意思是好的.正如保羅說:「以耶穌基督為至寶.」真正的意思是以物質的獲得,和屬靈的獲得作比較；然後得個結論──事奉主有永恆的價值,是更好的選擇.也是要祂的門徒計算代價；算清楚哪樣更化算,把真正的價值觀弄清楚.跟從基督的人,和一般人有不同的價值觀,每個基督徒應該看清楚,想清楚,計算清楚；才有決心跟從基督,事奉基督,才有永恆價值的事奉.
\newline
\newline
現在讓我們來看與事奉有關的幾處經文,好像一幅幅事奉的圖畫擺在我們面前；我們要看,要想,要作選擇,要作決定.
\newline
\newline
第一幅事奉的圖畫(九十6-17)這兩節經文是個對比,從中我們有所看見,有所選擇.「早晨發芽生長,晚上割下枯乾.」這是人生短促的形容詞,人生似野地花草,早上開花,晚上枯乾；不過一日之久,何等短促!當詩人想到人生短促時,並不因此消極；他醒悟到手作的工而禱告說,神啊!「......願你堅立我們手所作的工；」詩人重複說「我們手所作的工,願你堅立.」(17)重複的話語,表達他內心感受之深.這是對比所產生的覺悟.
\newline
\newline
詩篇九十篇的第一段,用了六個形容詞表示人生的短促,且一一更加短促.「昨日」一天24小時.「早上晚上」12小時.「如睡一覺」一般人睡覺平均八小時.「夜間的一更」「更」是二小時.「好像一聲嘆息」我想最長的嘆氣至多三秒鐘.「轉眼成空」眨一眨眼,我想大概1/10秒鐘吧.每個形容詞代表的時間更短.第10節末了是總結的形容詞「我們便如飛而去.」詩人對人生短促很了解,然而他轉眼看神的事工,又看自己雙手；他要將自己的手和神的事連起來,他願伸出手來為神工作.或許他以往沒有為神工作,但現在他要伸出他的手更努力為神工作；他求神堅定他的手,也求神堅立他的工作.「堅立」有兩個意思,一是求神堅定他工作的手,使他的工作不停止繼續事奉.二是讓他的工作堅立到永遠,有永恆的價值.兩個意思合併就是詩人的禱告和心願:「神啊!我明白了,我要把手奉獻給你.求你成全,求你堅固我,使我曉得如何選擇永恆有價值的事工.不是表面的工作,乃是看清楚認真的工作.」
\newline
\newline
昨天我收到主的僕人章力生教授從美國來信,提及他每晚一點鐘睡覺五點半起床；他說睡眠比以前增加了一小時,因已達84高齡.太太勉強他最遲十二點鐘應該就寢.但他想全心寫作,彌補以往失去的時間.這是位忠心主僕的心聲,他已經寫了八十幾本書,是他心血的研究.他為主而活,年事已高仍孜孜不倦努力事奉；使我非常感動!保羅說:「你們要愛惜光陰.」「愛惜」原文是「贖回」,我們要贖回我們的光陰,被主贖回的人,不但生命被贖回,連光陰也要被贖回.每個被救贖的人,光陰也應得救贖；我們把一切都獻給主,寶貴光陰要為主用,我們要事奉主.
\newline
\newline
第一幅圖畫,我們看見右邊是花草,早晨生長,晚上枯乾.左邊是醒悟的手舉起奉獻給神,求神堅定他工作的手,繼續事奉；也求神使他工作有永恆的的價值.
\newline
\newline
第二幅事奉的圖畫(一一六16)請特別注意末句「你已經解開我的綁索.」詩人把自己當作將獻的祭牲,被綁在祭壇處的柱子,祭司來了把繩索解開,牽到祭壇旁宰殺置於祭壇上,獻祭物給神.「解開綁索」意即一切阻攔都除去了,真正預備好了；不再被綁在柱上,乃是被宰獻於祭壇上.這幅圖畫很寶貴,其中有一關鍵,綁索是否已經解開?若祭牲仍綁在柱上,雖離祭壇很近卻走不上祭壇,獻不成祭,不能實行奉獻.很多基督徒,口說奉獻,心也想奉獻,在禱告中也表示奉獻；他渴望奉獻,懷著真誠愛主的心奉獻.可是事實上他還被綁在柱子上；不能掙脫攔阻,不能克服攔阻,實際走不上奉獻的路.今晚在座如有這樣的弟兄姊妹,但願主愛把你的綁索斬斷,叫你實際擺上祭壇上.
\newline
\newline
「耶和華啊!我真是你的僕人」,詩人重複地說:「我是你的僕人.」他更進一步地說:「我是你婢女的兒子,」意即僕人的僕人.一次次重複表達他內心的感受,他感到自己不配作神的僕人,這是真正奉獻者的態度；如果人把自己的奉獻當作是立功,覺得神需要他；那就是自己還不了解自己,也不認識神.我們何等渺小,微不足道,算不得甚麼；偉大榮耀的主,那裡需要我們呢?我們豈配事奉主?但是主恩無條件臨到我們身上,呼召揀選我們這些不配的人,使我們能夠蒙恩站在恩典的地位,有權利獻上自己.我們這樣一想,真是越想越感恩,越想越感到自己不配,就從心裡湧出詩人在16節所說的話.請注意!詩人非說是僕人的兒子,乃是說婢女的兒子.在當時婢女地位低於僕人；所以意思是說:「我是你僕人的僕人的僕人.」他心裡的謙卑多層面地表達.今晚如果你奉獻,你感覺配嗎?你是否帶著立功的心態來奉獻呢?我們來到主前,正如本篇詩開頭所講的感恩；從感恩的心,產生了服事的願望.弟兄姊妹!我們為何要奉獻,要事奉?唯一原因,就是對主愛的回應,激勵,感動,溶化；叫我們很自然很快樂很真誠的奉獻.
\newline
\newline
第三幅事奉的圖畫(四十6-8)耳朵開通和遵行神的旨意連在一起,二者關係極為重要.有的聖經學者認為「開通耳朵」,在原文有特別含意是連於(出廿一5-6)兩節的話.這幅圖畫,形容有一僕人本來可得自由；但因他愛主人,寧願放棄自由,仍舊留在主人家中永遠服事主人.當他作此決定時,主人和他站在門口,用錐子穿通他的耳朵作記號,自動獻給主人作永遠的僕人.所以詩人說:「神啊!我樂意照你的的旨意行,永遠事奉你.」(8)這幅圖畫多美麗!我們事奉主是勉強或是甘心樂意呢?是自己的選擇,自己意志的決定嗎?我們是否以自己的自由放下了自由,完全為了主；這是我們大膽地用自己的自由,這是個自由的不自由.我們把自己擺在主手中,永遠服事主.「穿耳」和「愛主人」兩件事連起,然後耳朵便開通了.這僕人身上有了清楚奉獻的記號.保羅說:「人不可以攪擾我,因我身上有了基督的記號.」這記號就是保羅奉獻記號.想至此,真是希望畫家把這幅畫描繪出來；我們也希望這幅畫的主角成為自己.我們用穿耳,來回應主穿手的愛；主的手被釘釘穿,在各各他山上為我們受死!主復活的身體遺留手中的釘痕,肋膀的槍傷.復活的主對多馬說:「伸出你手探入我手中,和肋膀的釘痕.」證明復活的主身上永遠留下這兩個記號.
\newline
\newline
第四幅事奉的圖畫(六九9)詩人因關心神的殿,心裡焦急如同火燒.他愛神,所以關心神的殿,他的關心很懇切.一個愛神的人一定愛神的殿,從未有愛主的人不愛主的教會.詩人因愛神,關心神的家,神的殿,至於心裡火燒的地步.一顆燃燒的心正如約翰加爾文的印,在蓋在他的書上,信上；印彫刻著高舉的手中一顆正在燃燒的心的圖案,他經常以此印提醒自己,他願意將火熱燃燒的心獻給神.有一位傳道人置了個這樣的印,印在自己畫上,他說這幅圖畫對他幫助很大,常常提醒他.
\newline
\newline
我們每個人可以讀這節經文,說:「主啊!求你幫助我關心你的教會,在你家中事奉；當教會情況欠佳時,不是批評,不是拆檯；乃是關心流淚,心裡焦急如同火燒.永不作旁觀者,乃是作你家中的兒女,作你的僕人；你的家就是我的家,我要愛你和你的家.我願在你家裡事奉.求你把關心賜給我.」我們在教會,有時看見軟弱,看見不該有的情況,看見人的失敗；我們當在主前禱告,流淚；並非批評的時候,不是自顯為義的時候,乃是我們關心努力更加努力之時.如果神的家沒有軟弱,不需要我們,我們願意在主家中盡上力量.
\newline
\newline
第五幅事奉的圖畫(一二三1-2)「僕人的眼睛望主人的手,使女的眼睛望主母的手」有三個意思,第一,預備好了,主人的手一動,就是指揮,發號施令；僕人,使女立刻遵命,準備工作,準備行動.第二,遵照主人的手如何指揮,注目該走的路；不但要作,也要按照主人,主母的心意去作.我們事奉不但要熱心要努力,還要知道主人的心意,要了解我們的主,所定下的屬靈原則.熱心可能成為糊塗的熱心,所謂越幫越忙；越熱心越麻煩,問題越大；熱心小,毛病小,攔阻也小.這是很重要的事,必須知道主人的心意,預備好,按著主人心意而作.教會中熱心工作的人不一定缺少,但是有熱心又按著真理而事奉的人是缺少；這是主,是教會所需要的人,更是我們事奉的模式.神造了兩種特別的天使,是基路伯和撒拉拂,他們教導我們如何事奉神.基路伯的原文是知識之意,撒拉拂原文是火；有火又有知識,二者合併就是神創造事奉祂的天使的名字.是表示神對他們事奉的旨意,叫他們的事奉火熱,且有真理的知識.二者配合最為需要,我們應從這兩天使學習功課,這樣的事奉才能滿足神的心.第三,仰望神直到祂憐憫我們,仰望也有依靠之意,主人的手不但是指揮的手,也是加力量,幫助,施恩,成全的手.如果這樣的手指揮我們工作,這手也幫助我們有工作的力量.感謝主!祂是好主人,祂不但吩咐我們工作,同時也賜給我們工作的力量.所以我們不必怕順服,不必怕責任的托付；當接受主之托付時,當我們順服時,當我們說:「主啊是的」時,盡可放心,主的恩典,恩賜,同在,能力夠我們用.這三個意思是完全的,如果只有前二者而缺後者,未免令人擔心作不成主工；如果具備三者,順服的僕人,主必使用,必加力,必是成功的僕人；是神所喜悅的僕人.
\newline
\newline
第六幅事奉的圖畫(一二六6)流淚出去撒種,歡歡喜喜帶禾捆回來.我們為了事奉付上代價,有時流淚；但有一天收成,歡歡喜喜帶著禾捆回來,藏在永生的倉庫裡.永遠的收成,何等寶貴!我們常歡喜豐收,但不願付代價,然而這裡說二者互相連繫,在神的計劃中不能分開.我們要收成就必須付代價,付了代價必然有收成.二者為一,不過要記得主的另句話「那人撒種,這人收割.」有時也有此情況,別人撒種你收割；或你流淚撒種,或別人歡呼收割,你肯嗎?這樣公道嗎?當別人歡呼收割,你的感受如何?你生氣嗎?抱怨嗎?批評主不公平嗎?或是你與別人一同快樂呢?如果事奉工作是為己,你流淚撒種別人收割,就抱怨,批評,惱怒,這表示是為己工作.如果是為主工作,為了基督的國度,為了天父的喜樂,為了主的榮耀,為了主的旨意；流淚撒種以後,無論收割者是自己或別人,都能夠喜樂,都能讚美,都能滿足.這是個實驗,是證據,到底是為己或為主工作.很多時候,這人撒種那人收割.我們要準備,清楚知道我們的眼淚是秋雨；秋雨是撒種之前的雨,不是收割之前的雨.換言之,秋雨和收割距離的時候尚長.人流淚撒種,如同秋雨,看不見將來的收穫；如無信心,沒真正堅定的奉獻,經不起流淚的考驗.反言之,如果知道神在大自然中賜下秋雨,春雨,必然有收成.同樣,在屬靈的境界中,有流淚也必然有收成.聖經中提及各種的眼淚,為主流的淚,為人,為自己悔改,為受苦,為感動而流的淚都不落空,都如同秋雨；神的時候一到,都要收成.
\newline
\newline
寫彌賽亞神曲的韓德爾,其中有段「祂被藐視」他為這段經文寫譜時,心深受感動；他想到基督所付的代價,所受的痛苦,他流淚無法繼續寫作,伏案痛哭.正在此時,他的朋友來到,見此情況卻步而退.事後問他才明白原由.這位聖樂家把心靈完全投入工作,他感受基督的愛,為人受痛苦在各各他山上；他不禁流淚痛哭.他說他常聽見空中的音樂,所寫的乃是他所聽見的.這聖樂家是神所賜的,他的心靈和作品打成一片.他有個特別的禱告,求神賜他在受難節之日離開世界.我不知道他為何這樣禱告,但神終歸聽了他的禱告.他的作品彌賽亞神曲是復活節日首次獻唱的,他的一生都和主連在一起.
\newline
\newline
我們為主流淚,為主付上代價永不落空,我們流淚撒種所付的一切,有天必要收成,這收成是為主的收成.
\newline
\newline
第七幅事奉的圖畫(卅一12)原是好的器皿,破碎了,無用了.第10節把破碎和無用的原因提出.「我的力量因我的罪孽衰敗.」因罪的器皿破碎不合神用,因罪破壞了事奉和奉獻.多麼可惜!有的人原屬合用的器皿,乏善自保聖潔,落在罪中以致被破壞；但是感謝主,祂有恩典能醫治破碎的器皿,也能重造器皿.好像先知耶利米到yiu匠那裡看見重造的器皿.我們的器皿如受損傷,來到主前,讓祂醫治,讓祂復原,讓祂重造；叫我們再次將自己獻上.主愛我們到底,祂接受我們的奉獻到底,主呼召我們到祂面前,主恩待我們.
\newline
\newline
十六世紀最偉大的藝術家米蓋基羅,他的大作是大衛的彫像.那塊大理石是許多彫刻家所棄置,不用已四十年之久；偶然被米蓋基羅發現,用以彫刻大衛的像.他精心計劃,把所有缺點除去,成功了這件世界稀有的藝術珍品.
\newline
\newline
我們的主也可這樣恩等我們.
\newline
\newline
但願在座眾弟兄姊妹,都能成為主手中合用的器皿.
\newline
\newline


\newpage

\section{第57屆~港九培靈會~滕近輝~第四講　詩篇中神言的反合性功效}
\label{sec:494}
\textbf{第四講　詩篇中神言的反合性功效}
\newline
\newline
連結: \href{https://www.hkbibleconference.org/session-message/view/494}{\texttt{https://www.hkbibleconference.org/session-message/view/494}}
\newline
\newline
\hyperref[sec:493]{< < < PREV SERMON < < <}
~
\hyperref[sec:toc]{[返目錄]}
~
\hyperref[sec:495]{> > > NEXT SERMON > > >}
\newline
\newline
詩篇 119:1-119:176
\newline
\begin{longtable}{cl}
\hline
\hline
章節 & 經文 (和合本修訂版)\\
\hline
119:1 & \begin{tabularx}{0.7\textwidth}{X} 行為正直、遵行耶和華律法的，這人有福了！ \end{tabularx} \\ \\ \relax
119:2 & \begin{tabularx}{0.7\textwidth}{X} 遵守他的法度、一心尋求他的，這人有福了！ \end{tabularx} \\ \\ \relax
119:3 & \begin{tabularx}{0.7\textwidth}{X} 他們不做不義的事，但遵行他的道。 \end{tabularx} \\ \\ \relax
119:4 & \begin{tabularx}{0.7\textwidth}{X} 耶和華啊，你曾將你的訓詞吩咐我們，為要我們切實遵守。 \end{tabularx} \\ \\ \relax
119:5 & \begin{tabularx}{0.7\textwidth}{X} 但願我行事堅定，得以遵守你的律例。 \end{tabularx} \\ \\ \relax
119:6 & \begin{tabularx}{0.7\textwidth}{X} 我看重你的一切命令，就不致羞愧。 \end{tabularx} \\ \\ \relax
119:7 & \begin{tabularx}{0.7\textwidth}{X} 我學習你公義的典章，要以正直的心稱謝你。 \end{tabularx} \\ \\ \relax
119:8 & \begin{tabularx}{0.7\textwidth}{X} 我必遵守你的律例，求你不要把我全然棄絕！ \end{tabularx} \\ \\ \relax
119:9 & \begin{tabularx}{0.7\textwidth}{X} 青年要如何保持純潔呢？是要遵行你的話！ \end{tabularx} \\ \\ \relax
119:10 & \begin{tabularx}{0.7\textwidth}{X} 我曾一心尋求你，求你不要使我偏離你的命令。 \end{tabularx} \\ \\ \relax
119:11 & \begin{tabularx}{0.7\textwidth}{X} 我將你的話藏在心裡，免得我得罪你。 \end{tabularx} \\ \\ \relax
119:12 & \begin{tabularx}{0.7\textwidth}{X} 耶和華啊，你是應當稱頌的！求你將你的律例教導我！ \end{tabularx} \\ \\ \relax
119:13 & \begin{tabularx}{0.7\textwidth}{X} 我用嘴唇傳揚你口中一切的典章。 \end{tabularx} \\ \\ \relax
119:14 & \begin{tabularx}{0.7\textwidth}{X} 我喜愛你的法度，如同喜愛一切的財物。 \end{tabularx} \\ \\ \relax
119:15 & \begin{tabularx}{0.7\textwidth}{X} 我要默想你的訓詞，看重你的道路。 \end{tabularx} \\ \\ \relax
119:16 & \begin{tabularx}{0.7\textwidth}{X} 我要以你的律例為樂，我不忘記你的話。 \end{tabularx} \\ \\ \relax
119:17 & \begin{tabularx}{0.7\textwidth}{X} 求你用厚恩待你的僕人，使我存活，我就遵守你的話。 \end{tabularx} \\ \\ \relax
119:18 & \begin{tabularx}{0.7\textwidth}{X} 求你開我的眼睛，使我看出你律法中的奇妙。 \end{tabularx} \\ \\ \relax
119:19 & \begin{tabularx}{0.7\textwidth}{X} 我在地上是寄居的人，求你不要向我隱藏你的命令！ \end{tabularx} \\ \\ \relax
119:20 & \begin{tabularx}{0.7\textwidth}{X} 我時常切慕你的典章，耗盡心力。 \end{tabularx} \\ \\ \relax
119:21 & \begin{tabularx}{0.7\textwidth}{X} 受詛咒、偏離你命令的驕傲人，你已經責備他們。 \end{tabularx} \\ \\ \relax
119:22 & \begin{tabularx}{0.7\textwidth}{X} 求你除掉我所受的羞辱和藐視，因我遵守你的法度。 \end{tabularx} \\ \\ \relax
119:23 & \begin{tabularx}{0.7\textwidth}{X} 雖有掌權者坐著妄論我，你僕人卻思想你的律例。 \end{tabularx} \\ \\ \relax
119:24 & \begin{tabularx}{0.7\textwidth}{X} 你的法度也是我的喜樂，我的導師。 \end{tabularx} \\ \\ \relax
119:25 & \begin{tabularx}{0.7\textwidth}{X} 我的性命幾乎歸於塵土，求你照你的話將我救活！ \end{tabularx} \\ \\ \relax
119:26 & \begin{tabularx}{0.7\textwidth}{X} 我述說我所做的，你應允了我；求你將你的律例教導我！ \end{tabularx} \\ \\ \relax
119:27 & \begin{tabularx}{0.7\textwidth}{X} 求你使我明白你的訓詞，我要默想你的奇事。 \end{tabularx} \\ \\ \relax
119:28 & \begin{tabularx}{0.7\textwidth}{X} 我因愁苦身心耗盡，求你照你的話使我堅立！ \end{tabularx} \\ \\ \relax
119:29 & \begin{tabularx}{0.7\textwidth}{X} 求你使我離開奸詐的道路，開恩將你的律法賜給我！ \end{tabularx} \\ \\ \relax
119:30 & \begin{tabularx}{0.7\textwidth}{X} 我選擇了忠信的道路，將你的典章擺在我面前。 \end{tabularx} \\ \\ \relax
119:31 & \begin{tabularx}{0.7\textwidth}{X} 我持守你的法度；耶和華啊，求你不要叫我羞愧！ \end{tabularx} \\ \\ \relax
119:32 & \begin{tabularx}{0.7\textwidth}{X} 你使我心胸開闊的時候，我就往你命令的道路直奔。 \end{tabularx} \\ \\ \relax
119:33 & \begin{tabularx}{0.7\textwidth}{X} 耶和華啊，求你將你的律例指教我，我必遵守到底！ \end{tabularx} \\ \\ \relax
119:34 & \begin{tabularx}{0.7\textwidth}{X} 求你賜我悟性，我就遵守你的律法，且要一心遵守。 \end{tabularx} \\ \\ \relax
119:35 & \begin{tabularx}{0.7\textwidth}{X} 求你叫我遵行你的命令，因為這是我所喜愛的。 \end{tabularx} \\ \\ \relax
119:36 & \begin{tabularx}{0.7\textwidth}{X} 求你使我的心趨向你的法度，不趨向不義之財。 \end{tabularx} \\ \\ \relax
119:37 & \begin{tabularx}{0.7\textwidth}{X} 求你叫我轉眼不看虛假，使我活在你的道路中。 \end{tabularx} \\ \\ \relax
119:38 & \begin{tabularx}{0.7\textwidth}{X} 求你向敬畏你的僕人堅守你的話！ \end{tabularx} \\ \\ \relax
119:39 & \begin{tabularx}{0.7\textwidth}{X} 求你使我所懼怕的羞辱遠離我，因你的典章本為美。 \end{tabularx} \\ \\ \relax
119:40 & \begin{tabularx}{0.7\textwidth}{X} 看哪，我切慕你的訓詞，求你因你的公義賜我生命！ \end{tabularx} \\ \\ \relax
119:41 & \begin{tabularx}{0.7\textwidth}{X} 耶和華啊，求你使你的慈愛臨到我，照你的話使你的救恩臨到我， \end{tabularx} \\ \\ \relax
119:42 & \begin{tabularx}{0.7\textwidth}{X} 我就有話回答那羞辱我的，因我倚靠你的話。 \end{tabularx} \\ \\ \relax
119:43 & \begin{tabularx}{0.7\textwidth}{X} 求你叫真理的話總不離開我的口，因我仰望你的典章。 \end{tabularx} \\ \\ \relax
119:44 & \begin{tabularx}{0.7\textwidth}{X} 我要常守你的律法，直到永永遠遠。 \end{tabularx} \\ \\ \relax
119:45 & \begin{tabularx}{0.7\textwidth}{X} 我要自由而行，因我尋求了你的訓詞。 \end{tabularx} \\ \\ \relax
119:46 & \begin{tabularx}{0.7\textwidth}{X} 我要在列王面前宣講你的法度，也不致羞愧。 \end{tabularx} \\ \\ \relax
119:47 & \begin{tabularx}{0.7\textwidth}{X} 我以你的命令為樂，這命令是我所喜愛的。 \end{tabularx} \\ \\ \relax
119:48 & \begin{tabularx}{0.7\textwidth}{X} 我向我所愛的，就是你的命令高舉雙手，我也要默想你的律例。 \end{tabularx} \\ \\ \relax
119:49 & \begin{tabularx}{0.7\textwidth}{X} 求你記念你向僕人所說的話，這話使我有盼望。 \end{tabularx} \\ \\ \relax
119:50 & \begin{tabularx}{0.7\textwidth}{X} 你的話將我救活了；這是我在患難中的安慰。 \end{tabularx} \\ \\ \relax
119:51 & \begin{tabularx}{0.7\textwidth}{X} 驕傲的人極度地侮慢我，我卻未曾偏離你的律法。 \end{tabularx} \\ \\ \relax
119:52 & \begin{tabularx}{0.7\textwidth}{X} 耶和華啊，我記念你從古以來的典章，就得了安慰。 \end{tabularx} \\ \\ \relax
119:53 & \begin{tabularx}{0.7\textwidth}{X} 我因惡人離棄你的律法，怒火中燒。 \end{tabularx} \\ \\ \relax
119:54 & \begin{tabularx}{0.7\textwidth}{X} 我在世寄居，以你的律例為詩歌。 \end{tabularx} \\ \\ \relax
119:55 & \begin{tabularx}{0.7\textwidth}{X} 耶和華啊，我夜間記念你的名，我也要遵守你的律法。 \end{tabularx} \\ \\ \relax
119:56 & \begin{tabularx}{0.7\textwidth}{X} 這臨到我，是因我謹守你的訓詞。 \end{tabularx} \\ \\ \relax
119:57 & \begin{tabularx}{0.7\textwidth}{X} 耶和華是我的福分；我曾說，我要遵守你的話。 \end{tabularx} \\ \\ \relax
119:58 & \begin{tabularx}{0.7\textwidth}{X} 我一心懇求你的面，求你照你的話憐憫我！ \end{tabularx} \\ \\ \relax
119:59 & \begin{tabularx}{0.7\textwidth}{X} 我思想自己所行的道路，我的腳步就轉向你的法度。 \end{tabularx} \\ \\ \relax
119:60 & \begin{tabularx}{0.7\textwidth}{X} 我速速遵守你的命令，並不遲延。 \end{tabularx} \\ \\ \relax
119:61 & \begin{tabularx}{0.7\textwidth}{X} 惡人的繩索纏繞我，我卻沒有忘記你的律法。 \end{tabularx} \\ \\ \relax
119:62 & \begin{tabularx}{0.7\textwidth}{X} 我因你公義的典章，夜半起來稱謝你。 \end{tabularx} \\ \\ \relax
119:63 & \begin{tabularx}{0.7\textwidth}{X} 凡敬畏你、守你訓詞的人，我都與他作伴。 \end{tabularx} \\ \\ \relax
119:64 & \begin{tabularx}{0.7\textwidth}{X} 耶和華啊，遍地滿了你的慈愛；求你將你的律例教導我！ \end{tabularx} \\ \\ \relax
119:65 & \begin{tabularx}{0.7\textwidth}{X} 耶和華啊，你照你的話，善待你的僕人。 \end{tabularx} \\ \\ \relax
119:66 & \begin{tabularx}{0.7\textwidth}{X} 求你教我明辨和知識，因我信靠你的命令。 \end{tabularx} \\ \\ \relax
119:67 & \begin{tabularx}{0.7\textwidth}{X} 我未受苦以先曾經迷失，現在卻遵守你的話。 \end{tabularx} \\ \\ \relax
119:68 & \begin{tabularx}{0.7\textwidth}{X} 你本為善，所行的也善；求你將你的律例教導我！ \end{tabularx} \\ \\ \relax
119:69 & \begin{tabularx}{0.7\textwidth}{X} 驕傲的人編造謊言攻擊我，我卻要一心遵守你的訓詞。 \end{tabularx} \\ \\ \relax
119:70 & \begin{tabularx}{0.7\textwidth}{X} 他們的心蒙昧如蒙油脂，我卻喜愛你的律法。 \end{tabularx} \\ \\ \relax
119:71 & \begin{tabularx}{0.7\textwidth}{X} 我受苦是與我有益，為要使我學習你的律例。 \end{tabularx} \\ \\ \relax
119:72 & \begin{tabularx}{0.7\textwidth}{X} 你口中的律法與我有益，勝於千萬金銀。 \end{tabularx} \\ \\ \relax
119:73 & \begin{tabularx}{0.7\textwidth}{X} 你的手造了我，塑造我；求你賜我悟性學習你的命令！ \end{tabularx} \\ \\ \relax
119:74 & \begin{tabularx}{0.7\textwidth}{X} 敬畏你的人看見我就歡喜，因我仰望你的話。 \end{tabularx} \\ \\ \relax
119:75 & \begin{tabularx}{0.7\textwidth}{X} 耶和華啊，我知道你的典章是公義的；你使我受苦是以信實待我。 \end{tabularx} \\ \\ \relax
119:76 & \begin{tabularx}{0.7\textwidth}{X} 求你照著你向僕人所說的話，以慈愛安慰我。 \end{tabularx} \\ \\ \relax
119:77 & \begin{tabularx}{0.7\textwidth}{X} 求你的憐憫臨到我，使我存活，因你的律法是我的喜樂。 \end{tabularx} \\ \\ \relax
119:78 & \begin{tabularx}{0.7\textwidth}{X} 願驕傲的人蒙羞，因為他們無理傾覆我；但我要默想你的訓詞。 \end{tabularx} \\ \\ \relax
119:79 & \begin{tabularx}{0.7\textwidth}{X} 願敬畏你的人和知道你法度的人都歸向我。 \end{tabularx} \\ \\ \relax
119:80 & \begin{tabularx}{0.7\textwidth}{X} 願我的心在你的律例上完全，使我不致蒙羞。 \end{tabularx} \\ \\ \relax
119:81 & \begin{tabularx}{0.7\textwidth}{X} 我渴想你的救恩身心耗盡，我仰望你的話。 \end{tabularx} \\ \\ \relax
119:82 & \begin{tabularx}{0.7\textwidth}{X} 我因渴望你的話眼睛失明，說：「你何時安慰我呢？」 \end{tabularx} \\ \\ \relax
119:83 & \begin{tabularx}{0.7\textwidth}{X} 我雖像煙薰的皮囊，卻不忘記你的律例。 \end{tabularx} \\ \\ \relax
119:84 & \begin{tabularx}{0.7\textwidth}{X} 你僕人的年日有多少呢？你幾時向迫害我的人施行審判呢？ \end{tabularx} \\ \\ \relax
119:85 & \begin{tabularx}{0.7\textwidth}{X} 不順從你律法的驕傲人為我掘了坑。 \end{tabularx} \\ \\ \relax
119:86 & \begin{tabularx}{0.7\textwidth}{X} 你的命令盡都信實；他們無理迫害我，求你幫助我！ \end{tabularx} \\ \\ \relax
119:87 & \begin{tabularx}{0.7\textwidth}{X} 他們幾乎把我從世上除滅；但我沒有離棄你的訓詞。 \end{tabularx} \\ \\ \relax
119:88 & \begin{tabularx}{0.7\textwidth}{X} 求你照你的慈愛將我救活，我就遵守你口中的法度。 \end{tabularx} \\ \\ \relax
119:89 & \begin{tabularx}{0.7\textwidth}{X} 耶和華啊，你的話安定在天，直到永遠。 \end{tabularx} \\ \\ \relax
119:90 & \begin{tabularx}{0.7\textwidth}{X} 你的信實存到萬代；你堅立了地，地就長存。 \end{tabularx} \\ \\ \relax
119:91 & \begin{tabularx}{0.7\textwidth}{X} 天地照你的典章存到今日；萬物都是你的僕役。 \end{tabularx} \\ \\ \relax
119:92 & \begin{tabularx}{0.7\textwidth}{X} 我若不以你的律法為樂，早就在苦難中滅絕了！ \end{tabularx} \\ \\ \relax
119:93 & \begin{tabularx}{0.7\textwidth}{X} 我永不忘記你的訓詞，因你用這訓詞將我救活。 \end{tabularx} \\ \\ \relax
119:94 & \begin{tabularx}{0.7\textwidth}{X} 我是屬你的，求你救我，因我尋求了你的訓詞。 \end{tabularx} \\ \\ \relax
119:95 & \begin{tabularx}{0.7\textwidth}{X} 惡人等著要滅絕我，我卻要揣摩你的法度。 \end{tabularx} \\ \\ \relax
119:96 & \begin{tabularx}{0.7\textwidth}{X} 我看萬事盡都有限，惟有你的命令極其寬廣。 \end{tabularx} \\ \\ \relax
119:97 & \begin{tabularx}{0.7\textwidth}{X} 我何等愛慕你的律法，終日不住地思想。 \end{tabularx} \\ \\ \relax
119:98 & \begin{tabularx}{0.7\textwidth}{X} 你的命令常存在我心裡，使我比仇敵有智慧。 \end{tabularx} \\ \\ \relax
119:99 & \begin{tabularx}{0.7\textwidth}{X} 我比我的教師更通達，因我思想你的法度。 \end{tabularx} \\ \\ \relax
119:100 & \begin{tabularx}{0.7\textwidth}{X} 我比年老的更明白，因我謹守你的訓詞。 \end{tabularx} \\ \\ \relax
119:101 & \begin{tabularx}{0.7\textwidth}{X} 我阻止我的腳走一切邪路，為要遵守你的話。 \end{tabularx} \\ \\ \relax
119:102 & \begin{tabularx}{0.7\textwidth}{X} 我沒有偏離你的典章，因為你教導了我。 \end{tabularx} \\ \\ \relax
119:103 & \begin{tabularx}{0.7\textwidth}{X} 你的言語在我上膛何等甘美，在我口中比蜜更甜！ \end{tabularx} \\ \\ \relax
119:104 & \begin{tabularx}{0.7\textwidth}{X} 我藉著你的訓詞得以明白，因此，我恨惡一切虛假的行徑。 \end{tabularx} \\ \\ \relax
119:105 & \begin{tabularx}{0.7\textwidth}{X} 你的話是我腳前的燈，是我路上的光。 \end{tabularx} \\ \\ \relax
119:106 & \begin{tabularx}{0.7\textwidth}{X} 你公義的典章，我曾起誓遵守，我必按著誓言而行。 \end{tabularx} \\ \\ \relax
119:107 & \begin{tabularx}{0.7\textwidth}{X} 我極其痛苦；耶和華啊，求你照你的話將我救活！ \end{tabularx} \\ \\ \relax
119:108 & \begin{tabularx}{0.7\textwidth}{X} 耶和華啊，求你悅納我口中的讚美為甘心祭，又將你的典章教導我！ \end{tabularx} \\ \\ \relax
119:109 & \begin{tabularx}{0.7\textwidth}{X} 我的性命常在我手掌中，我卻不忘記你的律法。 \end{tabularx} \\ \\ \relax
119:110 & \begin{tabularx}{0.7\textwidth}{X} 惡人為我設下羅網，我卻沒有偏離你的訓詞。 \end{tabularx} \\ \\ \relax
119:111 & \begin{tabularx}{0.7\textwidth}{X} 我以你的法度為永遠的產業，因這是我心中所喜愛的。 \end{tabularx} \\ \\ \relax
119:112 & \begin{tabularx}{0.7\textwidth}{X} 我的心傾向你的律例，謹守到底，直到永遠。 \end{tabularx} \\ \\ \relax
119:113 & \begin{tabularx}{0.7\textwidth}{X} 心懷二意的人為我所恨；但你的律法為我所愛。 \end{tabularx} \\ \\ \relax
119:114 & \begin{tabularx}{0.7\textwidth}{X} 你是我藏身之處，是我的盾牌；我仰望你的話。 \end{tabularx} \\ \\ \relax
119:115 & \begin{tabularx}{0.7\textwidth}{X} 作惡的人哪，你們離開我吧！我要遵守我神的命令。 \end{tabularx} \\ \\ \relax
119:116 & \begin{tabularx}{0.7\textwidth}{X} 求你照你的話扶持我，使我存活，不要叫我因失望而蒙羞。 \end{tabularx} \\ \\ \relax
119:117 & \begin{tabularx}{0.7\textwidth}{X} 求你扶持我，使我得救，時常看重你的律例。 \end{tabularx} \\ \\ \relax
119:118 & \begin{tabularx}{0.7\textwidth}{X} 凡偏離你律例的人，你都輕看他們，因為他們的詭詐必歸虛空。 \end{tabularx} \\ \\ \relax
119:119 & \begin{tabularx}{0.7\textwidth}{X} 你除掉地上所有的惡人，好像除掉渣滓；因此我喜愛你的法度。 \end{tabularx} \\ \\ \relax
119:120 & \begin{tabularx}{0.7\textwidth}{X} 我因懼怕你，肉體戰慄；我害怕你的典章。 \end{tabularx} \\ \\ \relax
119:121 & \begin{tabularx}{0.7\textwidth}{X} 我行公平和公義，求你不要撇下我，交給欺壓我的人！ \end{tabularx} \\ \\ \relax
119:122 & \begin{tabularx}{0.7\textwidth}{X} 求你保證你的僕人得福，不容驕傲的人欺壓我！ \end{tabularx} \\ \\ \relax
119:123 & \begin{tabularx}{0.7\textwidth}{X} 我因盼望你的救恩和你公義的言語眼睛失明。 \end{tabularx} \\ \\ \relax
119:124 & \begin{tabularx}{0.7\textwidth}{X} 求你照你的慈愛待僕人，將你的律例教導我。 \end{tabularx} \\ \\ \relax
119:125 & \begin{tabularx}{0.7\textwidth}{X} 我是你的僕人，求你賜我悟性，得以認識你的法度。 \end{tabularx} \\ \\ \relax
119:126 & \begin{tabularx}{0.7\textwidth}{X} 這是耶和華採取行動的時候，因人廢棄了你的律法。 \end{tabularx} \\ \\ \relax
119:127 & \begin{tabularx}{0.7\textwidth}{X} 所以，我喜愛你的命令勝於金子，更勝於純金。 \end{tabularx} \\ \\ \relax
119:128 & \begin{tabularx}{0.7\textwidth}{X} 你的一切訓詞，在萬事上我都以為正直；我恨惡一切虛假的行徑。 \end{tabularx} \\ \\ \relax
119:129 & \begin{tabularx}{0.7\textwidth}{X} 你的法度奇妙，所以我一心謹守。 \end{tabularx} \\ \\ \relax
119:130 & \begin{tabularx}{0.7\textwidth}{X} 你的話一開啟就發出亮光，使愚蒙人通達。 \end{tabularx} \\ \\ \relax
119:131 & \begin{tabularx}{0.7\textwidth}{X} 我大大張口，呼吸急促，因我切慕你的命令。 \end{tabularx} \\ \\ \relax
119:132 & \begin{tabularx}{0.7\textwidth}{X} 求你轉向我，憐憫我，就像你待那些喜愛你名的人。 \end{tabularx} \\ \\ \relax
119:133 & \begin{tabularx}{0.7\textwidth}{X} 求你用你的言語使我腳步穩健，不容罪孽轄制我。 \end{tabularx} \\ \\ \relax
119:134 & \begin{tabularx}{0.7\textwidth}{X} 求你救我脫離人的欺壓，我要遵守你的訓詞。 \end{tabularx} \\ \\ \relax
119:135 & \begin{tabularx}{0.7\textwidth}{X} 求你使你的臉向僕人發光，又將你的律例教導我。 \end{tabularx} \\ \\ \relax
119:136 & \begin{tabularx}{0.7\textwidth}{X} 我的眼睛流淚成河，因為他們不守你的律法。 \end{tabularx} \\ \\ \relax
119:137 & \begin{tabularx}{0.7\textwidth}{X} 耶和華啊，你是公義的；你的典章正直！ \end{tabularx} \\ \\ \relax
119:138 & \begin{tabularx}{0.7\textwidth}{X} 你所頒佈的法度是公義的，極其可靠。 \end{tabularx} \\ \\ \relax
119:139 & \begin{tabularx}{0.7\textwidth}{X} 我的狂熱把我燒滅，因我敵人忘記你的話。 \end{tabularx} \\ \\ \relax
119:140 & \begin{tabularx}{0.7\textwidth}{X} 你的言語極其精煉，令你僕人喜愛。 \end{tabularx} \\ \\ \relax
119:141 & \begin{tabularx}{0.7\textwidth}{X} 我渺小，被人藐視，卻不忘記你的訓詞。 \end{tabularx} \\ \\ \relax
119:142 & \begin{tabularx}{0.7\textwidth}{X} 你的公義永遠公義，你的律法是確實的。 \end{tabularx} \\ \\ \relax
119:143 & \begin{tabularx}{0.7\textwidth}{X} 我遭遇患難愁苦，你的命令是我的喜樂。 \end{tabularx} \\ \\ \relax
119:144 & \begin{tabularx}{0.7\textwidth}{X} 你的法度永遠公義；求你賜我悟性，使我存活。 \end{tabularx} \\ \\ \relax
119:145 & \begin{tabularx}{0.7\textwidth}{X} 耶和華啊，我一心呼求你，求你應允我！我必謹守你的律例。 \end{tabularx} \\ \\ \relax
119:146 & \begin{tabularx}{0.7\textwidth}{X} 我向你呼求，求你救我！我要遵守你的法度。 \end{tabularx} \\ \\ \relax
119:147 & \begin{tabularx}{0.7\textwidth}{X} 天尚未亮我呼喊求救，我仰望你的話。 \end{tabularx} \\ \\ \relax
119:148 & \begin{tabularx}{0.7\textwidth}{X} 我終夜雙眼睜開，為要思想你的言語。 \end{tabularx} \\ \\ \relax
119:149 & \begin{tabularx}{0.7\textwidth}{X} 求你按你的慈愛聽我的聲音，耶和華啊，求你照你的典章將我救活！ \end{tabularx} \\ \\ \relax
119:150 & \begin{tabularx}{0.7\textwidth}{X} 追逐奸惡的人迫近了，他們遠離你的律法。 \end{tabularx} \\ \\ \relax
119:151 & \begin{tabularx}{0.7\textwidth}{X} 耶和華啊，你就在我身邊，你一切的命令是確實的！ \end{tabularx} \\ \\ \relax
119:152 & \begin{tabularx}{0.7\textwidth}{X} 我從你的法度早已知道，這法度是你永遠立定的。 \end{tabularx} \\ \\ \relax
119:153 & \begin{tabularx}{0.7\textwidth}{X} 求你看顧我的苦難，搭救我，因我不忘記你的律法。 \end{tabularx} \\ \\ \relax
119:154 & \begin{tabularx}{0.7\textwidth}{X} 求你為我的冤屈辯護，救贖我，照你的言語將我救活。 \end{tabularx} \\ \\ \relax
119:155 & \begin{tabularx}{0.7\textwidth}{X} 救恩遠離惡人，因為他們不尋求你的律例。 \end{tabularx} \\ \\ \relax
119:156 & \begin{tabularx}{0.7\textwidth}{X} 耶和華啊，你的憐憫本為大；求你照你的典章將我救活。 \end{tabularx} \\ \\ \relax
119:157 & \begin{tabularx}{0.7\textwidth}{X} 迫害我的、抵擋我的甚多，我卻沒有偏離你的法度。 \end{tabularx} \\ \\ \relax
119:158 & \begin{tabularx}{0.7\textwidth}{X} 我看見奸惡的人就憎惡，因為他們不遵守你的言語。 \end{tabularx} \\ \\ \relax
119:159 & \begin{tabularx}{0.7\textwidth}{X} 你看我何等喜愛你的訓詞！耶和華啊，求你按你的慈愛將我救活！ \end{tabularx} \\ \\ \relax
119:160 & \begin{tabularx}{0.7\textwidth}{X} 你話語的精髓是真實的，你一切公義的典章永遠長存。 \end{tabularx} \\ \\ \relax
119:161 & \begin{tabularx}{0.7\textwidth}{X} 掌權者無故迫害我，然而我的心畏懼你的話。 \end{tabularx} \\ \\ \relax
119:162 & \begin{tabularx}{0.7\textwidth}{X} 我喜愛你的言語，好像人得到許多戰利品。 \end{tabularx} \\ \\ \relax
119:163 & \begin{tabularx}{0.7\textwidth}{X} 我恨惡，憎惡虛假；惟喜愛你的律法。 \end{tabularx} \\ \\ \relax
119:164 & \begin{tabularx}{0.7\textwidth}{X} 我因你公義的典章一天七次讚美你。 \end{tabularx} \\ \\ \relax
119:165 & \begin{tabularx}{0.7\textwidth}{X} 喜愛你律法的人大有平安，任何事都不能使他們跌倒。 \end{tabularx} \\ \\ \relax
119:166 & \begin{tabularx}{0.7\textwidth}{X} 耶和華啊，我仰望你的救恩，遵行你的命令。 \end{tabularx} \\ \\ \relax
119:167 & \begin{tabularx}{0.7\textwidth}{X} 我心謹守你的法度，這法度我極其喜愛。 \end{tabularx} \\ \\ \relax
119:168 & \begin{tabularx}{0.7\textwidth}{X} 我遵守你的訓詞和法度，因我所行的道路都在你的面前。 \end{tabularx} \\ \\ \relax
119:169 & \begin{tabularx}{0.7\textwidth}{X} 耶和華啊，願我的呼求達到你面前，求你照你的話賜我悟性。 \end{tabularx} \\ \\ \relax
119:170 & \begin{tabularx}{0.7\textwidth}{X} 願我的懇求達到你面前，求你照你的言語搭救我。 \end{tabularx} \\ \\ \relax
119:171 & \begin{tabularx}{0.7\textwidth}{X} 願我的嘴唇發出讚美，因為你將律例教導我。 \end{tabularx} \\ \\ \relax
119:172 & \begin{tabularx}{0.7\textwidth}{X} 願我的舌頭歌唱你的言語，因你一切的命令盡都公義。 \end{tabularx} \\ \\ \relax
119:173 & \begin{tabularx}{0.7\textwidth}{X} 求你用你的手幫助我，因我選擇你的訓詞。 \end{tabularx} \\ \\ \relax
119:174 & \begin{tabularx}{0.7\textwidth}{X} 耶和華啊，我切慕你的救恩！你的律法是我的喜樂。 \end{tabularx} \\ \\ \relax
119:175 & \begin{tabularx}{0.7\textwidth}{X} 願我的性命存活，得以讚美你！願你的典章幫助我！ \end{tabularx} \\ \\ \relax
119:176 & \begin{tabularx}{0.7\textwidth}{X} 我走迷了路如同失喪的羊，求你尋找你的僕人，因我不忘記你的命令。 \end{tabularx} \\ \\
[1ex]
\hline
\hline
\end{longtable}
詩篇119篇正在全本聖經的中央,本篇每節都提及神的話；可謂之神的話的詩篇,本篇也是全本聖經最長的一篇,共有176節.(88-89)兩節是全篇的正中,也就是聖經中心的中心.這裡提到兩件事:「神藉著祂的愛將我救活,使我得了生命.」這是生命之道.「我就遵守你口中的法度.」這就是神的話.生命和道兩樣都是全本聖經的中心,神的話是生命之道,使我們得生命,得生命之後要遵行神的話；若無生命就不能遵行神的話,反之,單有生命而不遵行神的話,也是不夠.神的話是堅定,永恆,永不改變,永不失效的.「神的話安定在天.」這是對我們的保證,所以我們當歡歡喜喜去遵行.
\newline
\newline
今晚我們特別要注意這角度,就是神的話的反合性功效,表面似是相反的兩方面；但是合起來就構成了完美的整體.現在讓我們從本篇詩來領受這個信息.
\newline
\newline
「愛神的話」(119)「怕神的話」(120)神的話在我們裡面引起相反作用,使我們「愛」也使我們「怕」我們愛神也愛神的話,但我們得到神的話以後,又使我們產生敬畏的心.神是慈愛的,所以我們愛神；神是公義的,所以我們要敬畏神.神的話使我們產生相反的兩種作用；但是兩相配合,就構成我們和神之間完整的關係.如果人只愛神,了解神的愛；恐會有問提,以為神是溺愛的神,好像溺愛孩子的父母；可能利用神的愛以致出了毛病.基督徒只愛神是不夠的,只知道神的愛也不夠；還要知道神的公義,因此產生敬畏的心.這兩方面在第二篇連在一起非常清楚:「當存畏懼事奉耶和華,又當存戰兢而快樂.」戰兢則敬畏,快樂則在神的愛中喜樂.兩者美好地相連.在現今時代缺少敬畏的觀念,基督徒應加以學習敬畏神.
\newline
\newline
箴言書專講智慧,開頭和結束都是講敬畏耶和華.箴言書一章頭六節是緒言,第七節才是正文的首句「敬畏耶和華.」到了最後卅一章末兩節提及「敬畏耶和華的婦女.」敬畏耶和華是智慧的開端和智慧的成全.
\newline
\newline
第二方面「你的言語......何等甘美......比蜜更甜.」(103)「......你使我受苦,是以誠實待我.」神公義的判語在詩人身上工作,使他受苦.神的話是審判,是管教的話,神按照祂話的真理教導詩人,使他學習功課.神的話一方面是苦的,另方面是甜的.感謝神!祂的話又甜又苦,正是我們完全的需要.我們讀神的話,開始時只讀神的恩典,慈愛,救恩,安慰,引導,保護,恩賜.神永不改變的愛,神的憐憫；這些都覺得甜美,讀了心中極為歡喜.但是到了我們更深入讀聖經時,我們體會到神的心意,看見祂的公義,祂的管教,他甘涼性的真理.這樣就教我們感到神的話非屬簡單,繼續讀下去,神的話似乎破碎了老我；好像兩刃的利劍刺入內心,將我們的意念,心中的隱密處完全剖開,顯明出來.神的話如同手術刀刺入患處把毒瘤割掉；當時我們覺得神的話苦澀味道,但感謝神!良葯苦口,忠言逆耳.所以,神的話若與我們心中意念相反之時,我們必須放下自己,讓神的話破碎我們自己,把老我拆毀,然後重新建造我們的靈性.詩人說:「神啊!你使我受苦,是以誠實待我.」(75)76節提及神的愛.神管教我們,用祂屬靈的手術刀之後,我們得到健康喜樂,神的恩典成全在我們身上.
\newline
\newline
以西結書第三章和啟示錄第十章,都提及神的話先甜後苦,神的話好像一小書卷,先知吃時,在口中是甜的,到了肚子是苦的.到底神的話是先甜後苦還是先苦後甜呢?真正的答案兩樣都不是；乃是甜──苦──甜,不是在甜和苦之中作選擇,而是甜──苦；最後帶我們到完全的甘美裡面,苦味的作用乃成全甘美.如果神的話先甜後苦,基督徒將永遠地苦死了.我們順服神的話,就帶來心靈深處最大的平安喜樂.
\newline
\newline
第三方面(105節)「腳前的燈」和「路上的光」是兩回事,腳前的燈是我們步步所需的光；路上的光使我們看見整條的路,可看得很遠.神的話給我們整體的認識,給我們曉得人生的意義,曉得我們何來何往,曉得神全部救恩的計劃,知道神全部的旨意；知道基督徒的總方向是甚麼,知道我們最後的歸宿是甚麼.路上的光給我們看見整個的道路,看見最後的目標在哪裡；好叫我們步步地走上去.單單知道整個神的計劃不夠,還要神的話引導我們今天所要作的事.現在所當說的話,現在應有的態度,應有的動作；今天我當如何?遠和近的光相配合,就是完全的光.
\newline
\newline
聖經中有句話很難明白,說亞伯拉罕蒙神呼召,出去時還不知道往哪裡去.難道他不知道起步的方向嗎?神呼召亞伯拉罕到祂所指示之地,有了總的方向,總的意義；但他不知道這條路怎樣走,也不知道要去之地是怎樣的.既有總方向,他就按照總方向舉步然後神步步引導.我們也是一樣,我們需要整體的亮光,對聖經基本教義清楚認識；但同時還需要神的話時刻引導前面的一步.這是非常重要的.「愚昧人眼望地極.」(箴十七24)眼望地極豈非眼光遠大嗎?為何說是愚昧人呢?意思是說他只想到將來,對將來有夢想滿懷希望.但眼前他的手不作工,腳不走路,只眼望地極作白日夢；所以他的希望永不完成,只是幻想,夢想.理想和幻想,夢想,有所分別；理想必帶來行動,今日行動努力,現在一個行動努力以求達成理想,這才是真正的理想.一個理想的人應是個行動,追求,努力的人；而夢想的人卻不然,他整天睡在床上夢想.愚昧人就是這樣.感謝主!祂的話不但給我們整體的認識,也督促,引導我們腳踏實地,步步往前走,這才是基督徒的人生.所以我們每天當領受主的話.
\newline
\newline
有一次我到馬尼拉,有位弟兄駕車陪我到禮拜堂,這青年人親口對我說:「新約聖經我已看過一遍,所以我覺得無需禮拜了.」他以為看過一遍新約聖經,就很足夠不必禮拜.基督徒讀聖經並非僅知內容,每天仍需神的話,教導及提醒,每天有現實的作用.讀聖經並非因我們不曉得勝經內容,要常讀聖經,熟悉聖經；但每天仍需要神的話來提醒,才是讀聖經的真正態度,每天重新得力.
\newline
\newline
第四方面(12-15)神的話是律例,即條例.神的訓詞,神的話就是道路,道路就是原則；律例是細則.聖經把原則和細則都給我們,兩樣都是我們需要的.僅有細則,永遠不夠；無論聖經有多少頁,如果加上十倍,也不能夠把條例提得完全.但原則是活潑的,能適應各樣的處境.細則是固定,死板的,單有固定,死板的細則,不能適應人生活潑的需要；所以細則應加上原則.有的人只想要原則不要細則,那也不行,因原則的應用很籠統,易於錯誤應用,也容易錯解原意；必須兩者具備.飛機的駕駛儀器盤上有極詳盡的指標,每方向的指標有很多度數.單有總方向指標不夠,還要加上詳細的度數,然後才能正確飛行.人生處境太複雜,所以需要原則和細則兩相配合.「......耶和華的法度確定,能使愚人有智慧.」(十九7)神的話確定,明確不改變；愚昧人一看照做.若是深奧的原則,愚昧人就難以了解.智慧人多用原則,人都以自己為智慧人,所以容易產生錯誤.現在流行的「處境論理學」就是只要原則不要細則.其中有一條「只要出於愛,甚麼都可以做」的基本原則,這原則太危險啦!很多時候,人欺騙自己；因人自私心的根太深,蒙蔽欺騙自己而不自知.特別是與性有關的事,自己不能分別,把私欲當作愛,這是人性問題產生的難處.聖經非常全備,把原則和細則都給我們.
\newline
\newline
第五方面113「恨」和「愛」.神的愛使我們恨惡與神的話違反的事.神的話使我們愛與祂律法符合的事.神的話帶來恨,也帶來愛.我們恨的時候,非出於脾氣的恨,乃真理的恨.恨產生拒絕,抵抗,否定.對罪惡的恨可堅定我們對真理的愛.「你們愛耶和華的,當恨惡罪惡.」(九七10)愛神者一定恨神所恨；這恨是屬靈的恨,恨惡罪惡,遠離罪惡.當我們有這樣的恨,我們向著罪惡剛強；因神的話使我們剛強.神的話也使我們柔和,因為有愛,把柔和剛強併合一起.如果要我們自己決定何時為柔和何時為剛強,很容易錯誤；但當我們照著神的話行事,就有自然的調和.當我們效法基督,衪柔和,我們也柔和；祂剛強,我們也剛強.我們按照基督的實例去作,又根據那些實例,類推其他的例子.這樣,我們就有自然的剛強和溫柔的調和.兩者真能調和嗎?有個具體的例子,水是柔和的；但洪水卻來勢甚兇,造成災害.這種調和帶來基督徒的氣質,內方外圓,內心方正,外面態度柔和.可惜很多人常是外方內圓,很難相處,沒有原則沒有品格.
\newline
\newline
第六方面101神的話帶來禁止,給我們規範,不偏不依；給我們真理的限制,不致放縱自己,任意妄為.詩人考究神的話,結果能自由而行.(45)神的話使他得自由.詩人第一次說神的話限制他,第二次說神的話使他得自由.這是限制和自由的配合,是我們實際生活中絕對需要的.
\newline
\newline
真正的自由是有軌道的,神創造大自然有自然律,一切井井有條.天體的運行,每個星球有其軌道,按照軌道進行,一切都是美麗的；若把自然律破壞或廢除,整個大自然立刻混亂.神在大自然的創造,星球能毫無攔阻的自由運行,因有軌道.照樣,在我們靈性範圍,如要得到真正的自由,必須接受主的話,「真理必叫你得自由.」我們在真理限制之下運用真理的自由.例如駕車司機必須遵守交通規則,交通才能暢順.初習字者,老師必督促由楷書開始,一筆一劃從楷書入門,先打好基礎而後自由發揮.自由裡面有分寸,有法度,才是真正的活潑的自由.其中有基本的道理,所以聖經告訴我們,絕對不可放縱情慾.聖經「欲」字有時寫「慾」,其實「欲」已經包括七情六欲,「慾」字只限於情慾方面,含義不同.聖經中特別是保羅書信論及私欲,下面加了「心」字,其實在應該加「心」字,這樣會擴大範圍；「欲」字已經把人的自然欲望,各種各類欲望都包括在內.我們一切欲望都須受真理的控制.比如初學寫字,筆劃多的寫在一個方格裡實在不容易；但經久練習,就能書寫自如.基督徒在複雜的人生中要學習的功課,就是把我們自由的範圍,自由的軌道任意差使,不自由作奴僕.
\newline
\newline
第七方面54神的話是我們的喜樂.本篇詩多提及詩歌.以神的話為樂是最高的境界,以神為樂也以神的真理為樂,至此境界無以復加了.請看另方面,136節提及詩人流淚；因為許多人違反神的真理,他很關心,心中有負擔以致流淚.當詩人看見有人不理神的話,就焦急如同火燒；(139)因他關心神的話使他產生負擔,神的話是他的喜樂,也是他的負擔.如果人真正了解神的話,透過神的話摸到神的心,他就會有神的負擔.有人說:「在聖靈裡禱告.」這是甚麼意思呢?聖靈親自用說不出來的嘆息為我們禱告,這是聖靈關心,顧念,了解我們,這是屬靈的負擔.如果人的禱告和聖靈同感,他就是在聖靈裡禱告.保羅說:「聖靈照著神的意思禱告.」我們在聖靈裡禱告乃是體會神的心.如果我們的禱告有了負擔,那就是懇切的禱告,就是言之有物的禱告.這樣的禱告對自己是感動的,對神也是感動的.
\newline
\newline
第八方面(98-36)詩人說,神的話使我們聰明有智慧,比我們的老師,仇敵,更有智慧.98節提到仇敵,36節提到老師.神的話作我們的老師.36節說神的話使我們不貪非義之財.正當手法得來之財是個成功的商人；基督徒可以作個成功的商人,可擁有發達的企業公司；正當所得的錢財可以善為運用作有意義的事,那是很好的；但非義之財不可貪,這樣的觀念是一般世人所嘲笑的,以為這是愚蠢.
\newline
\newline
感謝神!神的話使我們有智慧也使我們愚蠢.感謝神!使我們用永恆的眼光來看一切的利益,使我們看透了事物的真正價值.當我們真有智慧之時,表面似乎是愚蠢的,正如中國有說話「大智若愚.」
\newline
\newline
我們看了這些反合的功用,神的話是我們所需要的,適合於我們人生每個場合.
\newline
\newline
還有其他經文因時間不許可不再詳述.在此,我們得了一個總結論；本篇詩提到神的話在我心中,我將神的話藏在心裡,免得我得罪神,免得我違反祂的真理,免得我成為真理之敵.
\newline
\newline


\newpage

\section{第57屆~港九培靈會~滕近輝~第五講　詩入與神的深交}
\label{sec:495}
\textbf{第五講　詩入與神的深交}
\newline
\newline
連結: \href{https://www.hkbibleconference.org/session-message/view/495}{\texttt{https://www.hkbibleconference.org/session-message/view/495}}
\newline
\newline
\hyperref[sec:494]{< < < PREV SERMON < < <}
~
\hyperref[sec:toc]{[返目錄]}
~
\hyperref[sec:496]{> > > NEXT SERMON > > >}
\newline
\newline
詩篇 22:1-22:1
\newline
\begin{longtable}{cl}
\hline
\hline
章節 & 經文 (和合本修訂版)\\
\hline
22:1 & \begin{tabularx}{0.7\textwidth}{X} 我的神，我的神，為甚麼離棄我？為甚麼遠離不救我，不聽我的呻吟？ \end{tabularx} \\ \\
[1ex]
\hline
\hline
\end{longtable}
深交就是團契,合一,愛神的心；這些都是詩人們的表現.詩篇中最顯著的特點,就是詩人們與神的深交.今晚我們要按時間許可讀一些經文,藉知詩人如何愛神,如何在愛中與神合一；詩人們與神的交通團契,何等地深刻.
\newline
\newline
詩人們愛親近神(二二1)大衛說:「人對我說,我們往耶和華的殿去,我就歡喜.」這表示大衛愛親近神,他愛親近神有兩個因素:
\newline
\newline
第一,詩人大衛愛敬拜讚美神,他說「凡在我裡面的,都要稱頌神的聖名.(一○三1)意思就是他整個人,整個生命,一切的一切都要稱頌讚美耶和華.這是詩人湧出自己的心聲.他的敬拜和讚美就是對神的事奉.
\newline
\newline
事奉有三種意義,
\newline
\newline
1.本位的事奉,就是被造者站在被造的地位敬拜創造者.
\newline
\newline
2.生活的事奉,在生活中榮耀神,彰顯神.
\newline
\newline
3.工作的事奉,完成神分派給他的工作.
\newline
\newline
以上三種事奉合在一起,才是整全的事奉.許多人都以為工作才是事奉；工作不過只是事奉的部分,我們要切記；敬拜讚美是事奉,生活見證是事奉；三個因素必須互相配合,才是完整的事奉.
\newline
\newline
第二,詩人們愛禱告,他們因愛神所以愛禱告神.大衛是全國最忙的人,他是一國之君要處理國事當然忙；但從詩篇中我們看見他是個多禱告的人.由此可知,禱告不單是時間的問題,乃是心的問題.如果你歡喜一件事,你必定爭取時間去作；如果你愛一個人你必找時間見面.讓我再說,禱告不是時間的問題,乃是我們與主的關係問題,乃是我們愛不愛神的問題.
\newline
\newline
詩人與神同感「你們愛耶和華的,都當恨惡罪惡......」(九七10)詩人愛神所愛,恨神所恨.這是新的契合.詩人願所有愛神者,都當恨惡罪惡；因為神恨罪惡.保羅說:「愛是不喜歡不義,只喜歡真理.」如果我們愛神,就當恨神所恨的不義,愛神所喜歡的真理.詩人與神之間有完全的心靈上契合.(一三九23-24)大衛有對此方面的感受,他很關心自己裡面到底是否神所歡喜的,所以他有這樣的禱告:神啊!你恨惡罪惡,求你不讓任何罪惡進入我心；神啊!我怕自己不清楚,有疏忽,求你鑑察我,時常光照我；看在我裡面有甚麼惡行沒有?(三八18)也是大衛說的話,他不但認罪,同時也為罪憂傷痛悔；他非單憑口頭認罪,乃是從憂傷痛悔的心認罪.
\newline
\newline
認罪有四步驟,第一,知罪；第二,為罪憂傷；第三,認罪得蒙赦免；第四,離罪.這四步缺一不可.合併才是真正的悔改,完全對付了罪,完全回轉向著神而行,也是與神同行.這是非常重要的.
\newline
\newline
我在美國,看見有種特別窗子可以自動開關,下雨時自動關上,雨停時自動敞開.因為水有導電作用.我們的心向著罪惡也是如此,如果有生命的電流通；神的生命在我們裡面,當罪惡來到,便有敏銳的感受,將心門關上,拒絕罪惡.一個有神豐富生命的人,自然會發出與詩人同樣的呼聲說:「我因愛神所以恨神所恨的罪惡.」(五一篇),我們看見大衛悔改的認罪,他墮落,跌倒,犯罪；但後來他悔改,他非因怕神而悔改,從11節我們看見他悔改的線索,他說「求你不要使我離開你你的面.」可見他心中愛神,怕的是離開神的面.當他犯罪時,罪使他心靈麻木,所以跌倒,做了神所恨惡的事.後來,他愛神的心復興,心靈醒悟,發現他破碎的心靈,讓神的心破碎；所以他憂傷,痛悔,認罪.這篇詩幫助我們認識,一個真正悔改認罪的人,內心感受何等深切.使大衛回頭的,是他對神的愛.
\newline
\newline
詩人與神深交的第三方面表現(八四1)詩人說:「耶和華啊!你的居所何等可愛.」廿三篇末節說,「我且要住在耶和華的殿中,直到永遠.」詩人不想離開神的殿,因他愛神的殿,詩人關心神的殿,因愛神所以愛神的殿；當神的殿有問題,他心裡焦急,痛苦如同火燒.大衛為了建殿.(一三二1-5)告訴我們,他廢寢忘食,無心作事,直到成全建造聖殿的計劃為止.他把神的殿放在首位.(代上廿九3-5)載,他為了神殿,把自己的積蓄都獻上,獻三千他連得精金,七千他連得銀子.這些都不是取國庫,而是他自己的積蓄.他說:「因我心中愛慕我神的殿,就在預備建造聖殿的材料之外,又將我自己積蓄的金銀獻上.」他因愛神的殿而奉獻.
\newline
\newline
不久之前,新界有間教會,因人數增加座位不足,教會籌備置新堂需一筆款；有一弟兄畢業港大之後找工作,他就將第一個月的薪金奉獻.當弟兄姊妹知道此事,都大受感動和鼓勵.首次的薪金,意義重大,有的人預先計劃到時如何動用；這位弟兄卻全數帶到神的家中,為神的殿,為基督的教會而用.由此證明他愛神的心.這也就是詩人的表現.
\newline
\newline
多年前我在港九研經培靈會,聽貝牧師傳信息.他是加拿大人到中國內地會的宣教士,他說很多年前有個青年,在加拿大參加差傳大會受了感動,覺得應為差傳工作奉獻,他回家後禱告,看神要他奉獻多少錢,他禱告了半小時,決定奉獻五元加幣；在當時可算為數不少.到了次日,他和女友一同出去,輕易地花費卅元卻沒經過禱告,沒加以考慮.貝牧師說,這是愛的問題.有的人為聖工奉獻引以為榮,巴不得奉獻更多；但是有人卻非常勉強,認為可免則免.兩種不同心態分別在於愛的問題上.
\newline
\newline
詩人與神深交第四方面的表現(十六3)詩人因愛神就愛神的子民.大衛說:「世上的聖民,他們又美又善,是我最喜悅的.」聖民就是以色列人.難道以色列人又美又善嗎?難道大衛不知道以色列民中有很多壞人嗎?難道不知以色列民常悖逆神嗎?唯一原因,他用神的心來想到屬神的人,用神的眼光來看神的百姓.
\newline
\newline
神用兩種眼光看屬祂的人.
\newline
\newline
第一,神在基督耶穌裡看我們,透過基督的寶血看我們；看我們聖潔毫無瑕疵,稱為聖徒；耶穌基督的寶血洗淨我們一切的罪,雖然我們罪紅如丹顏,必白如羊毛；透過耶穌基督寶血洗淨,我們站在成聖的地位,所以神愛我們.我們想到舊約的雅各,他愛妻子拉結,甚至為她工作了十四年.拉結因生產便雅憫而死.聖經說雅各愛便雅憫最深.我相信重要的原因就是當他看見便雅憫時,就想到愛妻拉結因這孩子而死；在他心中,便雅憫和拉結分不開；是他愛妻的愛子,因此對便雅憫是雙份的愛.神對我們也是這樣.耶穌基督為我們而死,當神看祂愛子基督之時,又在基督裡面看我們；祂的感覺,祂的看法就完全不同.
\newline
\newline
第二,神用未來的眼光看我們.將有一天,每個重生得救的人,要和父神永遠在一起,在永遠榮耀的天家；那時我們身體得贖,敗壞的肉體解下,穿上復活的身體.榮耀屬天的身體和罪惡完全隔開,我們真是沒有瑕疵,完全像神,進入神兒女榮耀的自由中.神用那時的眼光於現在來看我們.媽媽看孩子就是如此,不管孩子在別人眼中如何惡劣,如何沒價值,沒希望；但母親卻以希望的眼,用愛的眼來看自己的孩子；雖然抱著流淚,卻仍然不失望不絕望,繼續費盡心血在孩子身上.為了孩子好,甚至犧牲自己都願意.我們的主,我們的天父就是如此看我們.
\newline
\newline
我們從神的眼光來看,在教會中看見基督徒,弟兄姊妹；我發覺許多的虧欠,許多的軟弱,許多的問題,許多不喜悅的事,我們有何感受?我們是否愛他們呢,抑或我們也在他們軟弱之中,和整個的軟弱上的虧欠上都有份?如果我們用旁觀者的眼光來看,所得的結論,和用神的眼光,用神的心來看,所得的結論完全不同.大衛愛神,知道以色列民是神的愛所揀選的；所以也用神的心,用神的眼光,來看他的同胞.因此他真誠地說:「聖民又美又善,是我所喜悅的.」我們愛弟兄姊妹是應該的,弟兄姊妹愛我們,原諒我們,也是應該的.讓我們用這樣的眼光來看教會,看弟兄姊妹吧!
\newline
\newline
詩人與神深交第五方面的表現(卅九9)詩人們願意和神的旨意打成一片,願意接受神的旨意,來成全在他身上.對照前面第3節,原來詩人心裡非常難受,他的心發熱,激動；因有人對不起他,他受了冤枉,他用舌頭激烈地說話.在這一切之中他忽然安靜下來,他之所以接受這遭遇得著平安；因他想到所遭遇的是出於神,所以就默然不語.本來不能忍受的,現在真正,平穩,安靜地接受；並非勉強接受,乃是出於了解的接受.勉強接受壓力太大,容易使人發生胃病和精神崩潰.真正的接受完全的順服,心裡平穩安息在神旨意中.這是詩人的功課,詩人給我們學習的榜樣.惟愛神真正認識神者,才能和神的旨意打成一片；真實接受神的安排.
\newline
\newline
我們的主,當祂在世的幾個城市工作中,被人拒絕,那時主被聖靈充滿,祂有歡樂的禱告:「父啊!天地的主,我要感謝你.......」最後祂說:「因為你的美意本是如此.」祂在逆境中,在被人拒絕時,祂的心中有安息.所以主說:「凡擔重擔的人,可以到這裡來,我就使你們得享安息.你們要學我的樣式.......」有了主的心在我們裡面,主的樣式是溫柔謙卑,主的擔子是輕省的,主的軛是容易的；否則,我們負擔,負軛就負不起.感謝主!祂在我們右邊.主說:「你們當負我的軛.」主並非邀請我們來負軛,而是邀請我們來負祂的軛.「負軛」和「負主的軛」完全不同.我們負主的軛時,主在我們身邊,與我們一同負軛,主負重的一端,且以安息,柔和的心腸給我們.惟有真正交托,清楚認識主,愛主的基督徒；才能夠默默不語,把最大的壓力卸下；在冤枉待遇,痛苦遭遇之中能夠說:「父啊!是的；因為你的美意本是如此.」我們的主在十字架上所說,所表現,都是內裡安息自由的表現.
\newline
\newline
感謝主!最榮耀的勝利,十字架上的勝利,這位為我們釘十字架的主,將祂的心給我們.
\newline
\newline
與神深交第六方面的表現,詩篇第一一九篇,有18次用「愛」字,「愛神的話」.因愛神而愛神的話,與神的話打成一片.
\newline
\newline
與神深交的第七方面表現,詩人因為愛神,就追求榮耀神.詩人立了志向要榮耀神的名直到永遠.(八六12)這是詩人對神的愛的表現,為了愛神就要榮耀神.第八十六篇也是大衛寫的詩.最少有兩個證據,證明大衛的心確實如此.
\newline
\newline
第一,大衛至少有兩次機會可以殺死掃羅；但是大衛沒這樣做.雖然大衛知道,如果殺了掃羅,自己可以作王；可是大衛說:「我是誰呢?怎可伸手傷害耶和華的受膏者?」掃羅是神所膏立所揀選的；當大衛面對抉擇,或殺掃羅自己作王,或榮耀神,不殺掃羅；他終而作了清楚確定的抉擇,寧可放棄王位也要榮耀神；他看榮耀神重於王位.這話說就容易,聽起來也容易；但實行卻不容易了.大衛實在將神放在首位.這是真正愛神的表現.
\newline
\newline
第二,大衛面對巨人歌利亞的事,人皆懼怕；但是大衛勇立敵人面前.(撒上十七46)給我們看見一條線索,大衛邊跑邊呼喊:「今日耶和華必將你交在我手中......使普天下的人都知道以色列中有神.」這是榮耀神的心,要見證神的能力,見證神的真實.所以他就勇敢上陣.榮耀神的心使他無所懼怕.
\newline
\newline
詩人與神深交第八表現(四三4)「以神為樂」是真正愛神者的感受.詩人說:「到我最喜樂的神那裡.......」以這位詩人來說,神是他最喜樂的,神是他最大的喜樂.正如保羅說:「以神為樂.」這是靈性至高的境界,這境界是在愛裡達到的；惟有你愛的人才能以那人為樂,惟有你愛道才能在道中自樂；愛主才能在禱告中快樂,在事奉主的事上快樂,在親近主,禮拜時快樂；快樂從中流露.快樂無法假裝,也非聽從命令；快樂是由衷自發的.被主愛澆灌,溶化者才能以神為樂.
\newline
\newline
最後,詩人所愛的是神自己,不是神所賜的任何福份.(十六5)大衛題到「杯中的分.」坦言之,「杯」就是酒杯.大衛意思是指一般人喝酒才快樂,才滿足；就是說:「我的滿足乃是耶和華,耶和華是我杯中的酒.」大衛說:「我醒了的時候......」(十七15)不是指每晨睡醒的時候,乃是指在永恆,永世裡醒了的時候,在天家醒了的時候.「......得見祂的形象,就心滿意足了.」大衛心滿意足,非為新耶路撒冷的黃金街,璧玉城,非因天家的榮耀,富貴；乃是因為見神的面就心滿意足.如果天堂沒有神,愛神者到了那裡不會滿足.
\newline
\newline
第七十三篇的作者,和大衛同樣有「愛神自己」的感受.「除你以外,在天上我有誰呢?除你以外,在地上我也沒有所愛慕的.」(七三25)神是詩人所愛慕的最高對象.
\newline
\newline
詩人說:「我的心平穩安靜.......」滿有喜樂,平安,滿足,這是形容他在神前的好感受.他的心「......好像斷過奶的孩子在他母親的懷中；」(一三一2)斷過奶和沒斷過奶的孩子,在母親懷中有分別的；沒斷過奶的孩子在母親懷抱中,為是吃奶,吃飽了就滿足.但斷過奶的孩子,在母親懷抱中得到滿足,完全享受母親的愛；不是奶給他滿足,乃是母親的愛給他滿足.大衛說他的心在他裡面就是如此.
\newline
\newline
我們看見詩人們對神的愛,這愛的作用產生完全的合一,深交團契.何等寶貴!是我們學習的對象,是我們的功課,是我們追求的路線.
\newline
\newline


\newpage

\section{第57屆~港九培靈會~滕近輝~第六講　詩篇中關於禱告的比喻}
\label{sec:496}
\textbf{第六講　詩篇中關於禱告的比喻}
\newline
\newline
連結: \href{https://www.hkbibleconference.org/session-message/view/496}{\texttt{https://www.hkbibleconference.org/session-message/view/496}}
\newline
\newline
\hyperref[sec:495]{< < < PREV SERMON < < <}
~
\hyperref[sec:toc]{[返目錄]}
~
\hyperref[sec:497]{> > > NEXT SERMON > > >}
\newline
\newline
詩篇最重要的特點是禱告,我們從詩篇中學習禱告.如果禱告是詩篇中最重要的信息,那麼,我們就要特別注意來領受這個信息.感謝天父把聖經當中的一卷,中心的一卷,來教導我們禱告的功課,我們都是主的學生；每個主的學生都要學禱告的功課.今晚我們不但要聽,而且要學,並且要學會.我們用學生的心態,作主門徒的心態來學習.求主使我們每個人,從詩篇裡面有心得.
\newline
\newline
在詩篇中用許多比喻,把禱告的功課很生動地提出,我們要藉這些比喻,來看禱告的道理.
\newline
\newline
第一個禱告的比喻(十八28)詩人說:「你必點著我的燈.」神藉我們的禱告,點著我們的燈,所以禱告就如我們心裡的燈光,若不禱告,就是裡面的燈熄了.因此,我們想到主耶穌的話說:「人的眼睛就是身體的燈,如果眼睛明亮,全身光明；如果眼睛昏花,全身就黑暗.」接著主耶穌說:「你裡面的燈若黑暗了,那黑暗是何等大.」所以我們裡面的燈光,比肉身的燈光更重要；我們若失去禱告,心靈就進入黑暗之中.(五六13)告訴我們,在神面前有生命的光.當我們到神面前時,就是進入生命的光裡,行走生命的道路.我們藉著禱告,可到神的面前,也是進到生命的光中,走生命的道路.不禱告的人很難走生命的道路.
\newline
\newline
神在大自然創造,創造陽光也創造生物.一棵棵的植物都是向著光,向日葵最為顯著.有位植物學家作了實驗,當一棵植物長出來,他用個木盒蓋上,木盒一邊開過兩寸的小洞,讓陽光透入；結果這棵植物向著陽光這邊穿出小洞生長.我們的禱告就是向著神,我們一禱告,眼睛,心靈開始朝向神的方向.禱告是向神的回應和響往；我們禱告之時,整個心靈立刻向著神的面；因為神的面,就是屬靈的陽光.如果我們沒有禱告,心眼的方向就會改變,向著自己,向著世界.讓我再說,禱告可得定向；不禱告者,是個沒有定向的基督徒.
\newline
\newline
第二個禱告的比喻(四二1)我們現在用「需要感」的角度來看這節經文.鹿快跑,大量消耗體力水份,極需補充,所以當鹿到溪水邊,就暢飲溪水；因牠有此需要.詩人說:「神啊!我需要你,渴慕你；所以來到你面前向你禱告,向你領受,向你支取.」禱告就是從需要感來向神支取.禱告需要很重要的動力.主耶穌有個比喻說有個人夜間睡覺,忽聞扣門聲,並有哀求的聲音說:「請你開門,我需要三個餅；因有朋友從遠方來,又飢又渴,而我家空無一物可以擺上,請你起來給我三個餅.」那人說:「夜深了,我和孩子們都睡了,請你走吧!」主耶穌說:「討餅的人情詞迫切,不斷直求.」(路十一5-13)這是主對我們的啟發.我們禱告乃是因為我們面對許多的需要,和自己的需要,自己無能為力供應所需,自己缺乏,一無所有；所以帶著需要感,到主面前求.沒需要感的人,不會禱告的；迫切禱告者,都是被需要感重壓的人.惟主能供應許要.我們想到創世記中,雅各和神的使者摔跤,雅各抓著神的使者說:「如果不給我祝福,我不容你走.」他急需祝福,他要屬靈的福份.這是強烈需要感的表現.神的使者為他祝福,他得到所愛慕的,得到最好的祝福；給他新名字「以色列王子」.在神在人面前都有能力.
\newline
\newline
我們需要屬靈的福分,我們禱告,就是抓住神:抓著神的愛,抓著神的應許,抓著神的信實,抓著神的旨意.我們對說:「求你按照你的話賜下福氣；否則我不鬆手.我們禱告的手伸上去；神賜福的手工作的手就伸下來.這是必然的.「神啊!我的心切慕你,如鹿切慕溪水.」
\newline
\newline
第三個禱告的比喻(八一10)神向以色列人說:「你要大大張口,我就給你充滿.」神所拯救的人,就是屬祂的人.神向屬祂的人說:「你是屬於我的,是我的百姓；我是你的神.」這幅圖畫是巢窩裡有小鷹,用嘴對著母鷹的嘴吸取食物,嘴張的越大,吸取的食物越多.神要我們儘量張大口,祂按著我們的容量供給我們.口張得大是出於信心；如果小鷹不相信母鷹帶了食物回來,我想牠必不肯開口.我們也是一樣,神要我們以信心開口禱告,信心多大,禱告的口也張開多大.神要我們大大張口.主耶穌說:「照你的信心給你成全了.」這是主一直應用的屬靈原則,永不改變的原則；哪裡有信心,哪裡就有成全.神永遠藉著信心工作,祂永遠對我們的信心有響應；我們的信心感動神的心,推動神的手；那裡有信心,那裡就有神能力的彰顯.
\newline
\newline
喬治慕勒說:「我要全世界看一看,我要全世界的人都知道；一個完全信靠神的人,神能為他作甚麼；一個完全信靠神的人,神能藉著他作甚麼.」這是他的雄心大志.他向全世界的人證明,神向信靠祂的人,是何等施恩顯出祂的能力.他的信心沒落空,他成了信心的例證.到了他91歲時,有人作個統計,經他手的錢有一千一百萬英鎊左右.這數目在當時極其可觀,最少相應於今日之一億一千萬英鎊,合港幣十億元之多.他建立孤兒院,經費從天賜下；神手伸出供應他的需要；因為他單單仰望神,他有信心,神對他的信心有響應.
\newline
\newline
第四個禱告的比喻(一○六23)開頭小字提及破口.摩西是個禱告的人,他站在破口禱告；他的禱告補了那破口.第一個破口在城牆上,第二個在河堤上.這兩個破口位置險要.城牆的破口,敵人隨時可以進入城中.禱告可以把守破口.河有破口,會引致缺堤,造成河水泛濫,關係重大,必須及時修堵破口.聖經用破口形容禱告的作用,實在非常貼切.如不禱告,等於城牆河堤有破口；敵人進入,洪水成災.神要人站在破口禱告.這是當時摩西最重要的工作.
\newline
\newline
感謝主!我們人人有禱告的權利,可在禱告上起重要的作用.如果我們記得彼得說的:「總要儆醒禱告,因為你們的仇敵魔鬼,如同吼叫的獅子,遍地遊行,尋找可吞吃的人.」(彼前五8)撒旦遍地遊行,尋找破口,無孔不入.彼得提醒我們,撒旦圍著教會打轉,時機一到,侵入教會；也用著基督徒尋找破口,企圖侵入生活中作敗壞的工作.我們的禱告,是把破口堵住,這是神對我們極其重要的啟示.我們每個人都有缺點,這缺點就是屬靈生命的破口.我們應當謹守破口,讓缺點變成優點；弱點變成長處,軟弱變成強壯.當我們身上有傷口出現,白血球立即增加,要消滅由傷口侵入的細菌.這是神的作為,對我們的靈性也是一樣.我們務必小心謹慎,儆醒為傷口禱告,求神施恩,神施恩之後,我們的缺點便成為優點.好像患了肺病的人,更能注意自己的健康,小心照料身體,注意營養和休息.所以人們都說患過肺病者反而長命.我們如果知道自己的缺點,謙卑在神面前求神憐憫,幫助,恩待,就可以勝過缺點.神施恩之後,自己特別追求,注意弱點；因此,弱點變成優點.要把守破口,為破口禱告,為屬靈的弱點禱告.
\newline
\newline
第五個禱告的比喻(一二一1)「詩人用向山舉目」形容禱告是很好的描寫,是美麗的圖畫,也是個很好的比喻.禱告就是向山舉目.在這比喻中,山代表穩固,「祂必不叫你的腳搖動.」(3)「依靠耶和華的人,好像錫安山,永不搖動.」(三五1)當我們禱告,向山舉目,依靠神之時,我們就得堅固,站立得住.屬靈站立得住的人就是禱告的人；禱告的人,就是堅定力量的人.如果我們把銅鐵置於露天,雨淋,風吹,日晒,漸漸生鏽腐蝕.屬靈的人好像香柏樹在山上,棕樹在神的院子,在神的殿中；到了年老之時,仍然結果子滿了汁漿而發青.(九二12-14)生命力健旺,年老仍充滿生命的顏色；因為他們栽在神的殿中,在神前生活長大,在最堅固的地方.銅鐵會改變,惟在神殿中的棕樹永不改變.
\newline
\newline
基督徒若不禱告,易於冷淡,退後,跌倒,失敗,灰心,喪膽.所以我們要向山舉目,我們的力量,幫助,從造天地的耶和華而來,祂使我們堅固不搖動.
\newline
\newline
「向山舉目」第二個意思,山是高處,當我們向山舉目,就登上靈程的高處；禱告者是個向上去的人,不禱告就向下去.流水是向下的；水變成蒸氣時就向上去.我們的本性是向下去；但當我們禱告時,裡面有了向上的作用,向上的能力,使我們靈程登高.
\newline
\newline
「向山舉目」第三個意思,高山屬天.(結四O2)提及聖殿在高山頂上.其意義之一是靠近神,是屬天的,超越世界的.我們禱告向山舉目,成了屬天氣質的人,我們的性格,內在的質素改變,使我們成為屬天的人.
\newline
\newline
感謝主!禱告是我們能力的來源.
\newline
\newline
我見過一幅畫對我有很大的幫助,畫的是一個農夫在田間耕種,忽然停頓了下來,跪下禱告；卻有一位天使來用他的工具操作.他停下工作去禱告,他的工作並不受虧損；神差遣天使替他工作,且工作得更好.神動工,一切都成了.我們勞苦工作,好像門徒整夜辛苦搖櫓卻仍在原處；但主一上船,船立刻到達目的地.主動工,事事都順利了；我們不能,在主凡事都能.我們跪下禱告之時,就是主起來工作之時.
\newline
\newline
有一個主的僕人,在他的日記首頁寫著:「禱告是最重要的工作.」這是多麼寶貴,多麼真實!弟兄姊妹們!這是你的估計嗎?你的價值觀嗎?真是你的衡量嗎?你能衷心地說:「禱告是最重要的工作嗎?」禱告好像無所事事,但卻是最重要的工作；因為我們禱告時,神工作了；我們最渴慕,最需要的,神就工作了.讓我再說,神一作工甚麼都成了.
\newline
\newline
感謝主!我們向山舉目,是我們能力的來源；造天地的耶和華,成為我們的幫助.祂是創造山的,是無所不能的,無限,無量的耶和華.藉著禱告,祂將能力賜給我們.
\newline
\newline
第六個禱告的比喻(一二三1-2)這段經文是禱告的美麗比喻.前晚我們看見這個比喻用在事奉上.今晚我們要從禱告的角度來看.「僕人的眼睛怎樣望主人的手,使女的眼睛怎樣望主母的手.準備及時順服,「舉目」就是順服之意.手一指揮立刻遵命而行.
\newline
\newline
我們在禱告中等候,禱告就是順服的態度,對主說:「我預備好了,為要遵行你的旨意；願你指示,願你的手指揮,好叫我立刻遵行.」禱告之時就是神指揮之時,不禱告的人很難看見主的手指揮；不禱告的人眼花繚亂,看不見神指揮的手；只看見環境,看見世界,看見自己.禱告時靈眼看見主有能力的手,指揮的手,一指揮,就順服,順服之時就支持了這有能力的手的能力.感謝主!主的手不但指揮,同時也帶著能力並且工作；我們順服主的指揮,主能力的手加在我們身上與我們同工.我再說,人若禱告求主的心意,看主的指揮而遵行；主能力的手與他同工,工作就作成了.安提阿教會的主僕們禁食禱告,聖經說,他們禁食,禱告就是事奉.當時聖靈對他們說話.所以禱告就是得神指示的機會,我們要抓住禱告的意義.禱告時,身心靈一片敞開,說:「神啊!指示我吧!使我明白你的心意,看見你指揮的手,我願意順服.」神一定引導這樣的人,這樣的人不會走錯路；有時似乎神引他到了困境,但經過水火終而進入豐富之地.
\newline
\newline
第七個禱告的比喻(一三○6)禱告如同守夜者等候天亮.第一,這守夜者在深夜,心裡有責任感,肩負重任必須守望,注意敵人的動靜；一發覺有危險,立刻燃起烽火,吹起號筒,喚醒全城的人.守望者沉重的責任,和全城百姓的安危關係重大.第二,守望者有期待感,他等候天亮；因他知道黎明白畫將要來臨.今天世界處於深夜中,我們務當儆醒,作個守望者；當牢記主的聲音說:「總要儆醒禱告,免得入了迷惑.」主知道我們的軟弱,了解我們的環境和情況；主的意思明顯指出,若不儆醒禱告容易受迷惑.主在客西馬尼園禱告,也正是深夜時候.今天我們處於這黑暗時代,惟有儆醒禱告者才能勝過迷惑.我們當等候那榮耀的黎明,主必再來,這是祂的應許,也是我們榮耀的盼望!到了神的時候,那要來的就來,並不耽延；我們儆醒等候,作聰明的童女；切勿等主來之時,我們睡著了.
\newline
\newline
最後(一四一2)這節經文說,我們的禱告好像香,好像獻祭,都是神所喜悅的.「香」真是名副其實,發出馨香升到神的面前.從前在聖殿中有三種香味,都是最滿足神的心,是最寶貴屬靈的香氣
\newline
\newline
1.膏抹聖殿所有物件的的聖膏油,表示分別為聖,為神使用,發出香氣滿足神的心.如果教會中有許多分別為聖的基督徒,教會就滿了香氣,滿足神的心.
\newline
\newline
2.素祭的香氣,素祭預表基督完美的品格.聖經記載,素祭是細麵加入乳香,乳香的特別就是燃著了即發出香味；用細麵加入乳香做餅燃烤就發出香味,基督完美的品格何等馨香,使神得到滿足.
\newline
\newline
3.燒香的香味,代表我們的禱告升上去,神聞香味心得滿足,祂就施恩.禱告是得神施恩的最重要途徑.在我們的人生有沒有這三種香氣發出?你在教會是屬於禱告的基督徒之中嗎?在你的教會中你在禱告上起作用嗎?你在教會中呈上禱告的貢獻嗎?教會的禱告會有你在內嗎?
\newline
\newline


\newpage

\section{第57屆~港九培靈會~滕近輝~第七講　神的四種言語}
\label{sec:497}
\textbf{第七講　神的四種言語}
\newline
\newline
連結: \href{https://www.hkbibleconference.org/session-message/view/497}{\texttt{https://www.hkbibleconference.org/session-message/view/497}}
\newline
\newline
\hyperref[sec:496]{< < < PREV SERMON < < <}
~
\hyperref[sec:toc]{[返目錄]}
~
\hyperref[sec:498]{> > > NEXT SERMON > > >}
\newline
\newline
詩篇 19:1-19:14
\newline
\begin{longtable}{cl}
\hline
\hline
章節 & 經文 (和合本修訂版)\\
\hline
19:1 & \begin{tabularx}{0.7\textwidth}{X} 諸天述說神的榮耀，穹蒼傳揚他手的作為。 \end{tabularx} \\ \\ \relax
19:2 & \begin{tabularx}{0.7\textwidth}{X} 這日到那日發出言語，這夜到那夜傳出知識。 \end{tabularx} \\ \\ \relax
19:3 & \begin{tabularx}{0.7\textwidth}{X} 無言無語，也無聲音可聽。 \end{tabularx} \\ \\ \relax
19:4 & \begin{tabularx}{0.7\textwidth}{X} 它們的聲浪傳遍天下，它們的言語傳到地極。神在其中為太陽安設帳幕， \end{tabularx} \\ \\ \relax
19:5 & \begin{tabularx}{0.7\textwidth}{X} 太陽如同新郎步出洞房，又如勇士歡然奔路。 \end{tabularx} \\ \\ \relax
19:6 & \begin{tabularx}{0.7\textwidth}{X} 它從天這邊出來，繞到天那邊，沒有一物可隱藏得不到它的熱氣。 \end{tabularx} \\ \\ \relax
19:7 & \begin{tabularx}{0.7\textwidth}{X} 耶和華的律法全備，使人甦醒；耶和華的法度確定，使愚蒙人有智慧。 \end{tabularx} \\ \\ \relax
19:8 & \begin{tabularx}{0.7\textwidth}{X} 耶和華的訓詞正直，使人心快活；耶和華的命令清潔，使人眼目明亮。 \end{tabularx} \\ \\ \relax
19:9 & \begin{tabularx}{0.7\textwidth}{X} 耶和華的典章真實，全然公義，敬畏耶和華是純潔的，存到永遠， \end{tabularx} \\ \\ \relax
19:10 & \begin{tabularx}{0.7\textwidth}{X} 比金子可羨慕，比極多的純金可羨慕；比蜜甘甜，比蜂房下滴的蜜甘甜。 \end{tabularx} \\ \\ \relax
19:11 & \begin{tabularx}{0.7\textwidth}{X} 因此你的僕人受警戒，遵守這些有極大的賞賜。 \end{tabularx} \\ \\ \relax
19:12 & \begin{tabularx}{0.7\textwidth}{X} 誰能察覺自己的錯失呢？求你赦免我隱藏的過犯。 \end{tabularx} \\ \\ \relax
19:13 & \begin{tabularx}{0.7\textwidth}{X} 求你攔阻僕人不犯任意妄為的罪，不容這罪轄制我，我就完全，免犯大罪。 \end{tabularx} \\ \\ \relax
19:14 & \begin{tabularx}{0.7\textwidth}{X} 耶和華－我的磐石，我的救贖主啊，願我口中的言語，心裡的意念在你面前蒙悅納。 \end{tabularx} \\ \\
[1ex]
\hline
\hline
\end{longtable}
今晚我們要從第十九篇看神的四種言語
\newline
\newline
神的第一種言語──大自然,神藉大自然向我們講話,向我們發聲.「這日到那日發出言語,......」(2)整個大自然發出言語.「無言無語,也無聲音可聽.」(3)意即大自然的言語是無言之言,無聲之聲；非用耳朵來聽,乃用眼睛來聽.我們觀看大自然,看見何等偉大奇妙的創造,天地萬物都在發聲,我們眼見時就聽見了；我們也用心來聽,聽見整個大自然如同第一節所說「諸天述說神的榮耀,穹蒼傳揚祂的手段.」述說,傳揚都是言語的意思,整個大自然都在見證神.
\newline
\newline
現在讓我們詳細分析大自然如何見證神,如何傳說神的作為.
\newline
\newline
第一方面,大自然述說神的偉大「諸天述說神的榮耀......」(1)諸天代表整個宇宙.「他的量帶通遍天下......」(4)「他」指言語,言語範圍通遍天下；天下不單是世界的天下,乃是第一節首兩字「諸天」的天下,是指整個宇宙.
\newline
\newline
宇宙多麼偉大!科學家告訴我們,宇宙兩端需時十萬萬光年的行程；而光的速度每秒186000英里.這宇宙見證神的偉大.
\newline
\newline
第二方面,宇宙見證,述說神的榮耀,神的榮耀意義豐富,其中部分就是美麗.大自然多麼美麗!春夏秋冬四季各有不同的美麗；並且每天的美麗都在變化,有黎明的美麗,日中的美麗,黃昏的美麗,晚上的美麗；有美麗的陽光,月光,星光.大自然是最偉大藝術家的老師.所有藝術家都是向大自然學習,從大自然吸取靈感.宇宙在述說神的美,神是偉大的神,也是一位美的神.
\newline
\newline
第三方面,宇宙述說,見證神的能力,「手段」就是能力的意思.中譯這兩字用的不太好,所指屬於不好之意,用在這裡使人覺得不太妥當.其真正意義乃是能力,大自然見證神的能力.「神將大地懸在虛空.」(伯廿六7)這話令人希奇,地球多重難以想像；神以能力把地球懸在虛空,這表達神無限的能力.
\newline
\newline
第四方面,宇宙傳說見證神的信息,「這日到那日發出言語」(2節)日復一日,日夜運轉不停,恆久不變,整個天體恆常運行；星球各按恆常不變的規律運行.(耶卅三20)神用日夜的運行來見證祂的信實可靠.神說如果太陽不再升起,月亮不再升起,祂的應許就會改變.神不改變好像日夜繼續,恆常不變運行一樣,晚上我們放心睡覺,我們知道明天太陽會出來.神說:「你們對於我的信實也可這樣放心.」神偉大的創造中顯出祂的信實可靠.地球圍著太陽運行,每天有160萬英里的行程,環繞太陽一圈需時365天5時48分48秒,永遠分秒不差；甚至1-10秒也不差.我的錶算是不錯的；但兩週須較對一次.錶的機件小,容易正確運行；但地球繞太陽每天轉動160萬英里,每年365天,竟然毫無差錯.你能了解嗎?這樣精確程度真是令人不可思議.神的信實也是如此.我們相信大自然所見證神的話語.
\newline
\newline
第五方面,大自然見證神的智慧「這夜到那夜傳出知識.」(2末句)知識包括智慧在內.大自然中充滿了智慧,使我們常常驚奇.例如熱脹冷縮是物理定理；但有一例外,水結冰時膨脹,所以比重減少能浮於水面.何以有此例外?乃為保全海,河,水塘中的生物；否則水沉底下,生物都要凍死.神用無限的智慧保守祂的創造.再舉一例,蠶吐絲成繭,力大者都沒法把繭撕開；但軟弱的蛾卻能從繭中出來.照人看來,惟死心踏地封在裡面；但蛾口裡有滴強烈酸液,到了適當時吐出酸液,使繭腐蝕成個小洞,正適合蛾出來.如果洞太大,蛾出繭容易,不需掙扎,就癱瘓不能飛；若洞太小,就掙扎不出來,死在繭中.但牠們所溶的洞,不大不小,正合需要.這是偶然的嗎?不是的,這是神的智慧.大自然充滿神無限的智慧.
\newline
\newline
第六方面,大自然傳講見證神的設計「神在其間為太陽安設帳幕.」(4)從「安設」我們想到設計.在宇宙中,我們清楚看見「匠心」就是創造的智慧,設計妥善,真叫我們驚奇.例如黃蜂的巢是蜂蠟做成的,其體中產生蠟質；需要蠟之時把翅膀伸開,以每秒五千次的速度抖動,是肉眼看不出的.翅膀震動,身體發熱,毛孔開,體中的蠟溶化由毛孔流出,把全身包住,然後乾了；黃蜂用腳把體上的蠟刮下,所以每隻黃蜂造出一堆堆的蠟,用以建造自己的蜂巢.蜂巢是六角形的,據科學研究有了驚人的發現；原來六角形狀,是所有形體中最節省材料的.黃蜂可以有這樣的智慧,懂得高深的數學計算,當然是造物主給予的設計.
\newline
\newline
第七方面,大自然述說,見證神的愛「太陽如同新郎出洞房.」(5)這裡用新郎,新娘來形容,暗示愛的意義.在大自然中有愛的彰顯,神將祂愛的本性藉著創造顯明出來；如同母愛多麼偉大；感動我們.有次我由北角乘輪渡往九龍城,航程約廿分鐘.我發現前兩排座位上有一婦人緊抱嬰孩不斷親吻,無論鼻子,眼睛,嘴,耳朵,頭髮,無處不吻.照我計算,廿分鐘之中至少親吻了50多次.我頗為欣賞,心中暗笑且感動.這是母親的愛,其實那嬰孩並無漂亮可言；但在母親心目中,自己的孩子最可愛.愛是天性,天性從創造的主宰而來.祂既將愛賜給我們,難道祂自己不了解愛嗎?難道神自己沒有愛嗎?
\newline
\newline
第八方面,大自然見證神所賜的喜樂「又如勇士歡然奔路.」(5末句)「歡然」表達大自然的喜樂,大自然滿了音樂；我們聽見風聲,海浪聲,禽鳥的歌聲,動物的叫聲,各種聲音配搭成為奇妙偉大能以形容的合唱.我讀中學時,各學校有歌詠團,當時最流行的合唱曲是趙元任先生作的「海韻」.我清楚記得歌中最低音調是波濤翻騰的聲音.整個大自然在歌唱.人類音樂的天才,更是神極大的恩賜；如果人間沒有音樂,相信世界缺少很多的樂趣.音樂家實在是神給人很大的恩賜.有位知名的作曲家海頓,他作的聖曲都是活潑,充滿喜樂的；有人問他「海頓先生為甚麼你寫聖樂有這麼活潑的曲調呢?他答:「當我想到上帝時,心中自然而然充滿喜樂；所以我就這樣寫,這是我直覺的表達.」這是偉大作曲家的見證.人類的音樂天才,天性,本性,實在非常奇妙!
\newline
\newline
第九方面,「奔」運行,整個宇宙運行乃是神所推動的.古代希臘藝術家阿爾斯多德說「上帝是第一位推動者.」宇宙一切都在運行.原子中的電子一直在運行,用光的速度每秒18600里運行.這是我們永遠無法了解的.原子有內太空,天上有外太空；整個宇宙在動,神是動的推動者.
\newline
\newline
第十方面,(5節末字)「路」太陽奔路有其軌道,運行有其律例,這是萬有引力,即地心吸力.科學家認為,無一科學家知道萬有引力是甚麼；但知道萬有引力最有功效.科學家說「電」我們每天用；可是沒人知道「電」是甚麼.
\newline
\newline
第十一方面,大自然見證神周全的安排「沒有一物被隱藏不得他的熱氣.」(6)無論大的小的,天空的,地下的,生物,植物；神的安排非常周全.
\newline
\newline
第十二方面,宇宙見證神的生命「熱氣」,所有生物都需要熱氣來維持生命；熱氣是生命的作用.神創造了生命,神也預備了生命所需要的.生物是多麼奇妙的現象,當我們看見這許多情況,在這六節經文,我們深感到聽見了話語.這日到那日發出言語,諸天述說神的榮耀,穹蒼傳揚祂的手段.
\newline
\newline
我們從第一種言語可以認識神,但是對神的認識從大自然來的不夠,所以神給我們另外的話語.
\newline
\newline
神的第二種話語──聖經,聖經就是神的啟示,在聖經中神將祂的啟示,更直接清楚啟示我們,使我們更完全的認識祂.第7節至12節這段經文提及神的話有六個不同的名稱,各有其特別的意義和重點.現在讓我們很快來看:
\newline
\newline
第一個名稱,「耶和華的律法全備.」神的話被稱為律法,原文最基本的意義是「標準」.神的話是我們行事為人的標準,是我們的規範.「規」是圓形,「矩」是方形,「規矩」就是圓,方形合起的完全標準.神的話是我們思想,言語,行為的準則,是我們人生的範圍,是我們不能超越的.
\newline
\newline
第二個名稱「耶和華的法度確定.」「法度」希伯來原文是見證之意,英譯則見證神的本性,這字的原文是把神話裡的本性彰顯出來.神的話是祂本性的見證,我們讀神的話可以認識神的本性.
\newline
\newline
第三個名稱,「耶和華的訓詞正直.」「訓詞」原文主要意義是知道人的職責,每個人都有責任；作人應盡的責任是甚麼,作神的兒女責任是甚麼,神的話清楚地告訴我們.我們要擔起責任,不負責任者是不成熟的人；就算卅歲還是未算成人,六十歲也未算成人.小孩子漸漸長大,有責任感之時就是成熟了.
\newline
\newline
第四個名稱,「命令」命令必須遵行.這裡說,耶和華的法度確定,一定要遵行,就得到智慧.
\newline
\newline
第五個名稱,「耶和華的道理潔淨.」「道理」原文是「懼怕」,也就是「敬畏」.敬畏之道,我們敬畏耶和華,從神的話學習,敬畏祂,順服祂.
\newline
\newline
第六個名稱,「典章」原文是「判決」.這是神所確定的,是神的判定；我們知道以後,照神的判定行事為人,神判定哪是對哪是錯；我們照神正確的判定而行,就是行義；要避免神所判定為錯的「罪」.這是聖經的話.
\newline
\newline
聖經的話語有何功效?這裡提及五種功效:
\newline
\newline
1.「耶和華的律法全備,能甦醒人心.」:(7)叫我們的心靈醒悟,回復應有的功能.
\newline
\newline
2.「耶和華的法度,能使愚人有智慧.」:有了神的話就有智慧,就看透撒旦的詭計,不上當,不致跌入陷阱.
\newline
\newline
3.「耶和華的訓詞正直,能快活人的心.」:神的話使我們快樂,屬神者以神為樂.感謝神!真理是可樂的,是可喜悅的,叫我們的心靈快活；道不是死的,不使我們痛苦的,不是呆板的,不是沉悶乏味的；道乃是活潑,喜樂的,使我們有屬靈的奮興.
\newline
\newline
4.「耶和華的命令清潔,能明亮人的眼目.」:神的道像光,能明亮我們的道路,眼目明亮清楚看見我們的道路.
\newline
\newline
5.「況且你的僕人因此受警戒.」:神的話警戒我們,亮起紅燈,叫我們停止,免得落到危險地步.
\newline
\newline
神的第三種言語(13節)「求你攔阻僕人,不犯任意妄為的罪.」神攔阻我們不落在罪中.神用我們的良心來攔阻我們,給我們是非之心,給我們的道德感,良心執行感受.此處經文沒有提到良心；但從神在我們心裡攔阻的作用,可以自然地推論到良心.感謝天父!把良心賜給人類.
\newline
\newline
何謂良心?保羅說:「神將律法刻在人心中.」良心是從神來的.「良心」原文是「同知」兩字,神和人一同知道,神和人有同感,神的律法刻在人心裡,神的律成了人的律,我們與神同感,同時也是人與人同感；因我們同有良心,同有是非之心,同有道德感,分辨對或不對.保羅說神的律刻在我們心裡.表示良心的作用,是不改變的；但是良心受了後天因素影響,受了破壞.聖經告訴我們,人的良心受了兩種破壞:第一種破壞「人的良心,如同被熱鐵烙慣了一般.」(提前四2)違背良心一次多次就麻木無知,感覺遲鈍.到了一個地步就喪盡良心,不再受良心的責備.第二種損害,(多一15)提及良心被污穢.意即內含雜質,不純粹,作用殘缺不全,作用時有時無,有時對有時不對,受了社會風氣後天因素之影響；產生了錯誤作用的標準,錯誤的判斷；因此,良心靠不住.有人把良心的作用以圖畫表達,在茫茫大霧中有把張開的弓,弓上有箭,有兩手拉弓；但看不見人面,只見兩手拉弓準備射箭.這畫表達良心,茫茫大霧表示神秘,看不出是從哪裡來的作用,是誰拉弓,進把弓放在那裡,誰把箭上弦；雖看不見但卻有件事實,就是良心有作用.良心的箭射出人會受傷痛苦不安,就算是兇惡的人,在寂寞夜半安靜下來也會哭.在美國財政部有個特別箱子,寫著良心基金的字樣；有的人做了損害公家的事,貪了不義之財；當良心發現,就匿名寄給政府.凡是這種款項,都投入良心基金.至今該款累積三百萬美元.神的律在人心裡,良心就發現；神藉良心向人說話,叫我們曉得神是道德性的神.印度人形容良心是三角形的,平時良心沒動作,三角形安定,尖角向上；當作了不該作的事,三角形的良心開始旋轉,三個尖角不斷刺痛人心.印度人又說,一般人的良心日久都成了圓形；因為三個尖角都磨光了,圓形旋轉就沒有感覺了.這是印度人通俗的說法.良心是事實,是神的聲音.
\newline
\newline
神的第四種聲音,也是神用以攔阻人犯罪的,聖靈的聲音.神藉聖靈講話.保羅說:「有我良心被聖靈感動,給我作見證.」(羅九1)意即保羅信不過自己的良心.人的良心靠不住；因已殘缺不全,已經烙慣了,已被污穢了；功能受虧損了,所以需要聖靈的光照.聖靈的感動,良心才能漸漸恢復應有的功效.好像日規一樣,如果沒有太陽,日規毫無作用；架構還在,但沒有陽光就沒影兒,沒影兒就不能指示時間.聖靈感動光照,好像陽光一樣,聖靈之光照在我們良心上,起了作用,恢復功能；日規的影兒出現,我們便知道情況如何了.
\newline
\newline
感謝天父!給了我們聖靈的話語,大自然的話語,聖經的話語,良心的話語.這種話語四種的聲音配合要完成甚麼?請看14節,請大家一同讀「耶和華我的磐石,我的救贖主啊!願我口中的言語,心裡的意念,在你面前蒙悅納.」神的四種話語最終目的,要讓屬神的人口中的言語,心裡的意念；加上13節所說的行為,「言語,意念,行為」包括整個的人生,能得在神面前蒙悅納,能符合神的旨意,符合神的標準.
\newline
\newline
感謝天父!一個屬神者能夠裡外都脫離罪,脫離一切的罪.這裡提及六種罪,因時間關係,請各位自己作家課.
\newline
\newline
神的四種話語的功效,使一切的罪都能脫離,都能勝過.感謝天父把救恩給我們,把赦罪之恩給我們,也把祂完整的言語給我們.在我們手中,心中,我們可照神的聲音行事為人.
\newline
\newline


\newpage

\section{第57屆~港九培靈會~滕近輝~第八講　詩人的喜樂}
\label{sec:498}
\textbf{第八講　詩人的喜樂}
\newline
\newline
連結: \href{https://www.hkbibleconference.org/session-message/view/498}{\texttt{https://www.hkbibleconference.org/session-message/view/498}}
\newline
\newline
\hyperref[sec:497]{< < < PREV SERMON < < <}
~
\hyperref[sec:toc]{[返目錄]}
~
\hyperref[sec:500]{> > > NEXT SERMON > > >}
\newline
\newline
詩篇 106:5-106:19
\newline
\begin{longtable}{cl}
\hline
\hline
章節 & 經文 (和合本修訂版)\\
\hline
106:5 & \begin{tabularx}{0.7\textwidth}{X} 好使我經歷你選民的福分，享受你國民的喜樂，與你的產業一同誇耀。 \end{tabularx} \\ \\ \relax
106:6 & \begin{tabularx}{0.7\textwidth}{X} 我們與我們的祖宗一同犯罪，偏邪行惡。 \end{tabularx} \\ \\ \relax
106:7 & \begin{tabularx}{0.7\textwidth}{X} 我們的祖宗在埃及不明白你的奇事，不記念你豐盛的慈愛，反倒在紅海行了悖逆。 \end{tabularx} \\ \\ \relax
106:8 & \begin{tabularx}{0.7\textwidth}{X} 然而，他因自己的名拯救他們，為要彰顯他的大能。 \end{tabularx} \\ \\ \relax
106:9 & \begin{tabularx}{0.7\textwidth}{X} 他斥責紅海，海就乾了，帶領他們走過深海，如走曠野。 \end{tabularx} \\ \\ \relax
106:10 & \begin{tabularx}{0.7\textwidth}{X} 他拯救他們脫離恨他們之人的手，從仇敵手中救贖他們。 \end{tabularx} \\ \\ \relax
106:11 & \begin{tabularx}{0.7\textwidth}{X} 水淹沒他們的敵人，沒有一個存留。 \end{tabularx} \\ \\ \relax
106:12 & \begin{tabularx}{0.7\textwidth}{X} 那時，他們才信他的話，歌唱讚美他。 \end{tabularx} \\ \\ \relax
106:13 & \begin{tabularx}{0.7\textwidth}{X} 很快地，他們就忘了他的作為，不仰望他的指引， \end{tabularx} \\ \\ \relax
106:14 & \begin{tabularx}{0.7\textwidth}{X} 反倒在曠野起了貪婪之心，在荒地試探神。 \end{tabularx} \\ \\ \relax
106:15 & \begin{tabularx}{0.7\textwidth}{X} 他將他們所求的賜給他們，卻使他們心靈軟弱。 \end{tabularx} \\ \\ \relax
106:16 & \begin{tabularx}{0.7\textwidth}{X} 他們在營中嫉妒摩西和耶和華的聖者亞倫。 \end{tabularx} \\ \\ \relax
106:17 & \begin{tabularx}{0.7\textwidth}{X} 地就裂開，吞下大坍，掩蓋亞比蘭一夥的人。 \end{tabularx} \\ \\ \relax
106:18 & \begin{tabularx}{0.7\textwidth}{X} 有火在他們黨中點燃，有火焰燒燬了惡人。 \end{tabularx} \\ \\ \relax
106:19 & \begin{tabularx}{0.7\textwidth}{X} 他們在何烈山造了牛犢，叩拜鑄成的像， \end{tabularx} \\ \\
[1ex]
\hline
\hline
\end{longtable}
詩篇原是唱的,但我們卻是誦讀詩篇；昔日以色列人唱詩篇,在歌唱中滿有喜樂,因為歌唱和喜樂是分不開的.詩篇共150篇,每篇都是唱的,當時以色列人是在聖殿中,不斷地歌唱.
\newline
\newline
根據聖樂出版界統計,全世界印行的聖詩有50萬首.以色列人唱的150篇,而我們唱的50萬首詩.沒有任何其他題目的詩歌,和聖經題目的詩歌,數量可相比.
\newline
\newline
喜樂在詩篇中共用了150次,可說詩篇中充滿了喜樂.詩人的喜樂不是感情主義,不是為喜樂而喜樂,乃是從信心裡發出來的喜樂；不是駝鳥式的喜樂,駝鳥把頭埋在沙中幻想.詩人的喜樂,是生命裡面的流露.他們對神有信心,所以對前途有把握；他們不看暫時的患難,他們超越了環境,看見黑雲上面的太陽；在患難中仍有喜樂.
\newline
\newline
詩人患難中的喜樂(五七7)開頭介紹,大衛逃避掃羅,他有生命之危,他作這詩,第7節說:「神啊!我心堅定......我要歌頌.」第五六篇前頭介紹,大衛給非利士人拿住落在敵人手中時寫這詩,他三次重複地說:「我要讚美神的話.」第六十三篇開頭介紹,大衛在猶大曠野作了這詩；當時他更加痛苦,因他親兒押沙龍要殺他.但那時他說:「我在床上記念你,在夜更時候思念你,我的心就像飽足了骨髓肥油,我也要以歡樂的嘴唇讚美你......我就在你翅膀的蔭下歡呼.」(六三5-7)大衛仍然歡呼,仍然有信心的喜樂,這種喜樂是世界的不能給,也不能奪去的.正像主耶穌所說的平安一樣.
\newline
\newline
馬丁路德說:「神給人類最大的兩樣恩典,第一是聖經,第二是音樂.」詩篇就是這兩樣的綜合體,就是神的話和喜樂歌唱.我們看見神的話,神的真理,神心意的流露；我們就以此為樂,在喜樂中歌唱.馬丁路德寫的「堅固堡壘」一詩,全世界的基督徒都唱過.那首詩就是詩篇四十六篇,他這篇詩得到靈感,寫成一首直唱到今日的詩,已經過了四百年仍然覺得寶貴.
\newline
\newline
當美國南北戰爭停止,北軍勝利,釋放了黑奴；那時全國在波斯頓舉行感恩讚美大會.有萬人的詩班,千人的樂隊；全國最佳演奏家,歌唱家大匯集,歡喜快樂,歌唱,感恩.這是勝利的喜樂!我們聯想(啟七9-12)的異象,有千萬不勝其數的人都穿白衣,被羔羊的血洗乾淨,他們環繞寶座歌唱；詞句的部份就是彌賽亞神曲「頌讚被殺羔羊」.這詞句乃取材於聖經.感謝主!那是勝利的歌唱,喜樂的歡呼,羔羊基督完成的救恩.無數得救者環繞羔羊寶座.我們又想到啟示錄十九章第一段,羔羊勝利之時,有哈利路亞的歡呼,歌唱聲連續四次.我曾在南韓漢城,聽過萬人組成詩班的大歌唱,那是一九七四年福音大爆炸的時候；我永不忘記那時的感受,如同將來羔羊寶座前的前奏曲.
\newline
\newline
現在讓我們來看詩篇中有甚麼喜樂,詩人為何喜樂?
\newline
\newline
(一)詩人為救恩而喜樂(一一八24)「所定的日子」就是救恩的日子,我們在耶和華的救恩中高興歡喜.這救恩怎樣完成的呢?26節的預言「奉耶和華名來的,是應當稱讚的.......」因為救主耶穌基督被差遣到世上來,祂是被差來的救主,祂完成了救恩.22節告訴我們,基督完成救恩最重要的過程.耶穌基督被人丟棄,成為被棄的石頭；但是神所揀選的,神用作救恩的房角的頭塊石頭.基督建立了祂的教會藉以施行救恩,基督是教會的房角石,也是頭塊石頭.建造房屋時,其他都跟著房角石看齊,房角石是標準.基督又是頭塊石頭,這名稱的要義就是基督居首位,祂是教會的元首,所有得救者都高舉基督為王.聖經他處,房角石和頭塊石頭分開來講,是兩獨立名稱；但這裡合併而言「成了房角的頭塊石頭.」雖是一個名稱卻是兩個意義.這節經文把救恩的奇妙提出.從前以色列人中有傳說,所羅門王建造聖殿時,有個工程要物色合作聖殿房角石的石頭,遍找不獲；有人發現了一塊石頭,但他認為不合用,棄置一邊.可是始終找不到合適的,於是回到被棄的石頭那裡,才發覺是最合用的.先被丟去,後被採用,所以才有這樣的話.耶穌基督被人棄絕,被人釘死；但祂卻是神所揀選的救主.祂成全了救恩,替我們死了,犧牲生命在各各他的十字架上.
\newline
\newline
美國賓州有個墳墓,直到如今仍有很多人參觀.墓碑刻著「林肯的代死者」字樣,這是林肯總統立的.當美國南北戰爭時,北軍為解放黑奴而作戰,很多兵士犧牲.林肯總統要全國人民牢記,這些死了的人,是為活著的人付了生命的代價,他說自己也是其中之一,所以立了這墳墓刻上自己的名字.
\newline
\newline
聖經告訴我們,主基督為我們死了,祂被人棄絕,釘死.祂是神所揀選的救主,為我們犧牲生命,因祂的死我們得生.我們在這救恩中歡喜快樂.這救恩的結論:「這是耶和華所作的,在我們眼中看為希奇.」(23)當我們要靜下來思想神的救恩時,我們真是感到超出人意料之外的希奇,神的道路高過人的道路,神的意念高過人的意念.我們在這救恩中高興歡喜.每個真正知道自己得救的人,都一定喜歡.今晚,我們在此歡喜快樂,再次為奇妙的救恩感謝主,從心靈深處向代死的主發出感謝讚美的歌唱.
\newline
\newline
(二)詩人為復興而喜樂,就是得贖以後的人曾經跌倒,後來醒悟,認罪悔改,得救赦免,恢復與神美好的關係,在這時歡樂歌唱.這是認罪得赦免的快樂.詩人大衛提及歡喜快樂,因當時他認罪悔改,得到神的赦免；他未認罪之前,神的手在他身上沉重,壓傷了他的骨頭.神定罪審判的壓力在他身上,他經過憂傷痛悔以後,重新得到歡喜快樂.(五一8)大衛不認罪時,終日唉哼,心裡充滿痛苦,沒平安,沒喜樂,但認罪後感受完全改變,他得回快樂的歌唱.大衛是個有歌唱恩賜的人,但未認罪之時卻沉默下來,心裡失去平安,與神的交通破壞；認罪之後,得回喜樂平安,他的聲音又釋放出來.開始歌唱,讚美,見證,充滿了喜樂.
\newline
\newline
罪使我們與神的交通被破壞,人得救後,神的手保守他得永生.所以主說:「誰也不能從我手中將他們奪去.」誰也不能將他們從父的手中奪去.一個重生得救的人,永遠屬於主；但他若再犯罪,雖不至滅亡；卻與神之間的關係,交通,團契都被破壞,就失去生命的喜樂和平安.這時必須認罪,惟有認罪悔改,才能恢復與神的關係.
\newline
\newline
數日前我提過,真正的悔改有四個因素:知罪,為罪憂傷,認罪,離罪.這四步在詩篇五十一篇已經提過.「因為我知道我過犯；我的罪常在我面前.」(五一3)這是知罪.「神啊,憂傷悔痛的心,你必不輕看.」(五一17)這是為罪憂傷.「我向你犯罪,惟獨得罪了你,在你眼前行了這罪,......」(五一4)認罪就得到赦罪.「......塗抹我一切的罪；......」(五一14)離開罪.這四步都在這篇詩中,大衛對付了罪,脫離了罪,勝過了罪,完全的改變.從此以後,我相信大衛更加愛神,而且他對神的愛有新的質素,就是謙卑.大衛經過跌倒以後,回復愛神,謙卑在神面前.他覺得他要愛神,他悔改,都是神的憐憫,也是神的恩典,所以他充滿新的喜樂.大衛說:「從禍坑裡,從瘀泥中,把我拉上來,使我的腳站在磐石上,使我的腳步穩當.」(四十2)這時候他從痛苦中學了功課,然後他的腳步穩當；好像一隻羊經過創傷以後,更加小心.從前常離開牧人,偏行己路,遭遇危險後；被牧人救回醫治了創傷,從此他聽從牧人的話.大衛說:「我就把你的道,指教有過犯的人,罪人必歸順你.」(五一13)他跌倒的經驗成了他的信息,能幫助跌倒者重返神的面前.
\newline
\newline
司布真牧師說:「如果你發現你所愛者被人用刀刺死,當你看見兇器時有何感覺?是否很珍貴很喜歡它.收藏如同心愛之物,若這樣做,則表示你不愛被刺死者.是誰把我們的主釘死在十字架?是你,我的罪；我們的罪好像長鐵釘釘死主.如果覺得鐵釘可愛；如果看我們的罪,還貪婪不捨,那就證明我們不愛主,不愛為我們死在十字架上的主.我們愛罪過於愛主,那是世間最悲慘的事.一個蒙恩得十字架拯救的人,若愛罪過於愛主.我們愛罪過於愛主,那是世間最悲慘的事.一個蒙恩得十字架拯救的人,若愛罪過於愛主,是令主最傷痛的事.」司布真這些話,對我幫助很大,相信也幫助了許多基督徒.讓我們從這角度來看罪,當我們面對罪的試探時,讓我們想到基督的十字架.
\newline
\newline
(三)詩人因收成而喜樂(一二六6)這節經文我們前已讀過,如果是寶貴的經節,我們應該讀了又讀；可從不同的角度和重點來看.「那帶種流淚出去的,必要歡歡喜喜的帶禾捆回來.」這是收成的快樂.撒種時有困難,甚至流淚,付代價,受痛苦；但到了收割之時,那一切的代價都成過去；如今收成豐富,且儲藏在永生的倉庫裡.永生的倉庫是主耶穌用的名詞,就是永恆的倉庫,永不失去價值,永遠得著.我們曾付上代價,收成時眼淚已乾,眼淚化成歡呼；我們充滿了喜樂,帶著禾捆回家.
\newline
\newline
弟兄姊妹們!我們為主付代價,受委屈；背十字架,受痛苦,都不落空,都要收成；這是主的應許.我們在天上之時,就了解神的恩典超過我們所求所想的；我們所撒的種永不落空,種的是甚麼,收的也是甚麼.神是信實的,是創始成終的神,我們靠祂所作的,到了祂的時候,就必收成.
\newline
\newline
馬禮遜牧師,是首位來中國傳福音的基督教教士.他小時候是街上流浪的孩子,時逢英國主日學興起不久,在他故鄉有個弟兄,受感動作熱心作主日學工作,有一天,他對一位姊妹說:「可否將你家開作主日學學校,並請你擔任老師,讓街上的孩子來主日學?」這姊妹原本不想這樣做；因她覺得太難.但是受到弟兄熱心感動,終於答應了.那弟兄是開製衣廠的,凡來參加的小孩每送一套衣服以作鼓勵.那位姊妹便出去找尋街上的小孩子,馬禮遜是其中之一；但過了兩週,羅拔失蹤了,那位姊妹再去找他,他又來了.但是再過兩週,又是不見了.姊妹說:「算了吧,不必再去找他,他不會來了.」可是弟兄對她說:「姊妹啊!我心裡不知為何有個感覺,每當我看見羅拔,我感到神要用他；可否請你再去找他,告訴他,我多送套衣服給他.」姊妹只好又去找他回來,從此他再沒失蹤.後來他成了首位來華的宣教士.兩位同工為了聖工付上代價,受了為難,花了時間,費心費力,不惜金錢；可是到了神的時候,一切都不落空,都要收成.流淚撒種的必歡呼收割.
\newline
\newline
(四)詩人得釋放的喜樂(五三6)以色列人被擄到巴比倫釋放回來,從敵人手中獲得自由,真是他們歡喜快樂的一件大事.人得救後,如果不慎又落在撒旦手中,受罪的捆綁,是很悲慘的事.當日巴比倫王尼布甲尼撒到耶路撒冷,把聖殿器皿掠奪回國,放在自己偶像廟宇裡.原來是為事奉神的器皿,竟然為他事奉偶像,多麼悲慘!當屬神者被撒旦擄去作奴僕,作牠的工具,作不義的器皿,何其悲慘!當有此情形發生時,我們可以說,神的心破碎了!若一旦被釋放回歸於神,重新奉獻；挺身昂首,把自己交在神的手中；被神使用,作義的工具,何等喜樂!我們為這樣的人感謝神,這樣的人更能體會神的恩典.
\newline
\newline
北角有一間教會有個弟兄,本來吸毒,經長期掙扎無法脫離；後來信主,主釋放他脫離毒的捆綁,現在很健康,很熱心.我每次被邀請到該教會講道時,都見到他；他是我的同鄉,特別有鄉親情感.因看見主內的弟兄,經歷了神的釋放,熱心愛主追求,腳步站穩而感謝主!今年上半年,晨曦福音戒毒會,舉行十七週年感恩會,有兩位弟兄見證,一是吸毒十七年後來信主得釋放；另一個吸毒廿年,信了主得釋放.二人在台上作見證,非常寶貝；我相信他們心裡充滿喜樂,他們作見證時面容充滿喜樂.當我聽見證時,好像隱約看見他們從前吸毒的情況,面貌全然不同.這兩人多年進出監獄無數次,神拯救釋放他們,今日他們口唱感恩之歌.這是被釋放的快樂,是福音的大能.
\newline
\newline
(五)詩人樂人之樂(一○六5)樂國民的樂,這兩個樂字很有意義；第二個樂字使別人快樂,第一個樂字自己快樂.自己的快樂是從別人的快樂來的,這是非常寶貝的快樂.我們聯想新約,保羅所說的話:「故此,我們得了安慰,並且在安慰之中,因你們眾人使提多心裡暢快歡喜,我們就更加歡喜了.」(林後七13)保羅因提多歡喜,提多歡喜,保羅更歡喜,這是保羅的心.他的人生目標,把快樂帶給別人；他決志永不把憂愁痛苦帶給別人,不把自己的快樂,建築在別人的痛苦上.保羅說:「我自己定了主意,再到你們那裡去,必須大家沒有憂愁.」(林後二1)「你們眾人都以我的快樂為自己的快樂」這裡我們看見了交互的作用,看見保羅使別人快樂；別人也以保羅的快樂為快樂.請看保羅的心:「我先前心裡難過痛苦,多多流淚,寫信給你們；不是叫你們憂愁,乃是叫你們知道我格外的疼愛你們.」(林後二4)保羅為哥林多人流淚,很多人不了解他,怪責他；但後來明白他所作所然,都是為了關心疼愛他們.保羅一貫的宗旨,就是要把救恩的喜樂,神之愛的喜樂帶給別人,為了使別人喜樂,寧可自己流淚；但終有一日,他的眼淚變成歡呼的喜樂.
\newline
\newline
有首短詩使我印象深刻,題目是「祝福你點路燈的人」.古時英國,每至黃昏,有專人負責點著路燈；點著之後,燈光明亮,黑暗消失.另有一首短詩描寫一支蠟燭,燃燒時其蠟溶化滴下,如同眼淚流下,但燭光依然照耀.詩曰:「你所流的淚已經化為別人路上的光.」感謝神!這樣的人是祝福別人,也是神所賜福的人.正像神對亞伯拉罕說:「我必賜福給你,......你也要叫別人得福.」(創十二2)這是神對亞伯拉罕的後裔,及他後裔的旨意；因為神對他說:「地上萬物,都必因你的後裔得福.......」亞伯拉罕的後裔,當然是基督；但「後裔」還有其他的意義,是指以色列人,是亞伯拉罕肉身的後裔,不過也包括基督徒在內,他屬靈的後裔.「以信為本的人,就是亞伯拉罕的子孫.」(加三7)神對你我的旨意就是,神要賜福你,也要你去祝福別人.這樣,你心中就充滿喜樂.
\newline
\newline
(六)詩人以神為樂的喜樂(四三4)詩人說:神是他最喜樂的神,神就是他的喜樂.「又要以耶和華為樂,祂就將你心裡所求的賜給你.」(三七4)另一翻譯:「又要以耶和華為樂,祂就將你心裡所樂的賜給你.」這樣的翻譯是有所根據的,「求」是一種願望；心願,就是所渴望所以為樂的；當他以神為樂時,神就把他心裡所樂的賜給他.我覺得這種翻譯很有道理.聖經中給我們看見,凡尊重神的,神必尊重他.這裡我們又看見,凡以神為樂的,神也以他為樂,也賜給他心裡所樂的事.我們靈程最高峰,就是到了以神為樂的境界.快樂是自然的流露,當他以神為樂時；慈愛的父神必然把他所樂的,所寶貴,所羨慕的賜給他.神有一原則,人所渴望的,祂才賜給.主耶穌說:「沒有人把珍珠給狗,恐怕狗誤以為是石頭攻擊,回過頭來咬你.」我們的神不把祂寶貴的恩典給不渴望的人；因他不要,不接受,不重視,不珍貴.當神發現人渴望祂的話,渴望祂的真理,祂的同在,能力；渴望祂的面,渴望祂自己；神就將最好的給他,神將自己給他.
\newline
\newline


\newpage

\section{第57屆~港九培靈會~滕近輝~第九講　熱愛聖殿聖城的詩篇}
\label{sec:500}
\textbf{第九講　熱愛聖殿聖城的詩篇}
\newline
\newline
連結: \href{https://www.hkbibleconference.org/session-message/view/500}{\texttt{https://www.hkbibleconference.org/session-message/view/500}}
\newline
\newline
\hyperref[sec:498]{< < < PREV SERMON < < <}
~
\hyperref[sec:toc]{[返目錄]}
~
\hyperref[sec:499]{> > > NEXT SERMON > > >}
\newline
\newline
詩篇 84:1-84:2
\newline
\begin{longtable}{cl}
\hline
\hline
章節 & 經文 (和合本修訂版)\\
\hline
84:1 & \begin{tabularx}{0.7\textwidth}{X} 萬軍之耶和華啊，你的居所何等可愛！ \end{tabularx} \\ \\ \relax
84:2 & \begin{tabularx}{0.7\textwidth}{X} 我羨慕渴想耶和華的院宇，我的內心，我的肉體向永生神歡呼。 \end{tabularx} \\ \\
[1ex]
\hline
\hline
\end{longtable}
詩篇其中六篇有共同的特點,就是詩人表達他們對於聖殿和聖城的熱愛；同時也提及聖殿和聖城中他們最欣賞的事物.這六篇詩篇至少給我們兩方面的啟發.
\newline
\newline
第一方面,讓我們學習詩人,如何熱愛聖殿聖城,讓我們在今天能夠熱愛教會.
\newline
\newline
第二方面,我們也可在教會,把他們所欣賞的物件建造起來.
\newline
\newline
詩篇八十四篇1-2節,詩人說耶和華的居所(聖殿)何等可愛.詩人羨慕渴想耶和華的院宇,他的心腸,肉體都在聖殿中向永生神呼籲.(小字歡呼)詩人的愛慕,歡呼,渴想,都集中在神.這樣的愛真是感動我們.願我們今天也一同追求熱愛主的教會!
\newline
\newline
聖經告訴我們,教會是神所愛的,也是一般的弟兄姊妹所愛的.保羅說:「神用自己的血買贖教會.」(徒廿28)請注意!這裡不是說基督的血買贖教會,乃是說神用自己的血買贖教會.相信保羅心目中基督和神合而為一；同時,他也是加強重點於神如何愛教會的事上.基督愛教會,為教會捨己.」(弗五25)教會是主耶穌所愛的,而且愛中有代價,為教會捨己.我們也願意效法基督,愛教會為教會捨己.真愛裡面有捨己,愛教會非單憑口說,應有捨己的行動來證明.(羅八26-28)提到聖靈愛教會,用說不出來的嘆息為教會禱告.教會所有的人都在聖靈心上；好像舊約的大祭司,胸前的胸牌上有十二塊寶石,刻了十二支派的名.胸牌在心上,乃預表主耶穌,聖靈是基督的靈；基督的心也在聖靈的心中,聖靈體會基督的心；用說不出來的嘆息,為教會的弟兄姊妹禱告.「作長老的寫信給教會,說,教會就是誠心所愛的；不但我愛,是一切知道真理之人所愛的.」(約貳1)作長老的是教會的同工,是傳道人,傳道人愛教會,也是所有屬神的人.所愛的所有同道愛教會,如果教會中充滿了愛主的人,教會有福了,這教會必然興旺,會發展,會長進,會結果子,會榮耀神的名.
\newline
\newline
當詩人提到愛聖殿時,他首先最欣賞聖殿中甚麼呢?第3節提及詩人最注意,最欣賞聖殿裡的祭壇.祭壇有兩方面的意思,就主耶穌方面來說,祭壇當時是預表主耶穌獻上自己,成全救恩.就我們方面說,祭壇是我們獻上自己的地方.主為我們獻上自己,成全救恩；我們又獻上自己,為主而活,作個向主奉獻的人.主的奉獻,引起我們的奉獻；我們的奉獻,是回應主的奉獻.這兩方面合成祭壇整個的意義.詩人把祭壇形容成為一幅很生動且富詩意,富有靈意的畫.麻雀為自己找著房屋,燕子為自己找著菢雛之窩；詩人心中所注意的不是這兩種飛鳥,他心裡真正想到的是人.第4節作為解釋「如此住在你殿中的,便為有福.......」「如此」表明這不過是比喻而已,用麻雀,燕子來形容人.無論何人都在祭壇處蒙恩,從基督完成救恩方面來說,不論何等卑微都可蒙恩,麻雀微不足道,無人注意.主耶穌說:「兩個麻雀不是賣一分銀子嗎?表示很便宜,無價值,無人重視.」(太十29)
\newline
\newline
有許多人是卑微,沒人注意,沒人重視的,他們更沒地位；但是在基督的救恩裡,在神的愛中,有空處為了他們,祭壇歡迎他們,主耶穌歡迎他們.主在十架上被釘死時,祂的雙手伸開,好像在那裡表示歡迎；所有的罪人,卑微的人,都可來到十架前接受救恩.祂說過:「凡到我這裡來的,我總不丟棄他.」又說:「我若被舉起來,就要吸引萬人來歸向我.」(約六39,十二32)
\newline
\newline
在我們中間,是否還有些朋友們,未曾信主呢?主歡迎你到祂那里,在祂的救恩中有空處等候你.
\newline
\newline
在「奉獻我們自己」的意義來講,當我們奉獻以前,好像麻雀飛來飛去,沒有房屋,沒有安身之處,飄流地飛；但若來到祭壇處獻上自己,我們立刻有安居之所,有安身立命之地,便停止屬靈的飄流.基督徒還未奉獻之前,靈性是飄流的,是流浪者,如同以色列人從前在曠野飄流了四十年,多麼可惜!一個奉獻的基督徒來到祭壇處,情況立即改變；找到了安身立命之地,找到屬靈的立場,有了立足點:有了定位點,站穩了,剛強了,有定向了,人生有意義了,不再在曠野,而且進入迦南美地了,這是奉獻的作用.菢雛之窩,就是生產的地方,一個人奉獻自己,就開始屬靈結果子的生活；不奉獻的基督徒,不容易結果子,不容易有屬靈的生產.一旦走上祭壇,他的態度,他的人生,開始改變而結出屬靈美好的果子來.這是祭壇的作用,多麼寶貴!
\newline
\newline
保羅在羅馬書十二章及事奉之前,第一節先講奉獻,然後講事奉的道理.意即奉獻乃事奉的基礎.如果人沒有奉獻,不可能好好的事奉,不可能有果效的事奉,不能夠真正的事奉.必須先走上祭壇,然後才能事奉得好.(林後八2)及馬其頓的弟兄姊妹,有非常好的光景,他們表現了金錢上奉獻的厚恩.保羅為他們作見證說:「在極窮之間,還格外顯出他們樂捐的厚恩.」「他們先將自己獻給主,」(林後八5)他們先將心,將生命,把整個自己獻上；因此才有金錢,時間,恩賜,心血,以及各方面的奉獻.這是馬其頓弟兄姊妹給我們的榜樣.人若不先獻心,還沒有生命的奉獻；要他獻時間,金錢,恩賜,心血,他不甘心去作,或勉強去作；但是奉獻生命之後,整個情況改變.我們必須先建立祭壇,在我們的教會裡；如有人要開荒工作,要建立新的教會,要作織帳棚的工作；就必須先建立祭壇,帶領信主的弟兄姊妹,到奉獻的地步,讓他們了解奉獻的真理.有的人想,奉獻的真理標準太高,太難,所以不敢講奉獻的道理；那是很大的錯誤.如果教會沒有奉獻的信息,弟兄姊妹就不會奉獻生命；他們仍然是飄飛的麻雀,燕子,不會過奉獻的生活,也不會事奉主.這個屬靈先後的次序很重要,在舊約給了我們重要的功課,(拉三3)提及以色列人被擄歸回本國之後,重建聖殿,其次序是先建造祭壇,他們認為這是最重要的事.這值得我們注意,教會的建造中,我們屬靈的建造中,第一步應建立祭壇,才能夠建立好的教會；多有弟兄姊妹獻身的教會,才是多有力量的教會；獻身就是力量,是教會真正實力的所在.舊約時代,祭壇四角各有像牛角的角,聖經中牛角代表能力.基督在祭壇上流血,成全救恩；救恩永不失效!這福音本是神的大能,要救一切相信的人.永遠的能力,救人的能力,於今仍是一樣.同時,在我們方面來講,我們的奉獻也產生能力.一個奉獻的人,是個有能力的人,能打美好的仗,為主作美好的工.
\newline
\newline
請注意!保羅不僅講獻心,他也說:「你們要將身體獻上,當作活祭.」身體代表四肢百體,代表我們具體生活,我們要把具體的生活獻上.有的人只獻心,沒獻身,結果成了感情主義者；如同祭壇上有火,而無祭物,或有祭物而無火,這是祭壇上兩個不該有的情形.有人口說奉獻,禱告說奉獻,這是感情的奉獻,好像只有火；但卻沒有具體奉獻的行動,沒有祭物那是無用的.有祭物沒有火,那是冰冷的奉獻,不會滿足主的心.保羅說:「要心裡火熱,常常服事主.」心裡火熱是火,常常服事主是行動,合併一起,有火也有祭物.
\newline
\newline
日前曾提到彌賽亞神曲作者韓德爾,他作此曲時,一連廿三天閉關在房間,為他送飯者常發現他沒有進食.他整個生命,心靈,思想,感情,一致投入這種曲中,他說,當他作哈利路亞合唱時,他好像到了天門口,親耳聽見天上的歌聲,他照著所聽見的寫下；所以他說,這是神賜給他的,不是他自己寫的.我也曾提及當他寫到主受苦時,他滴下眼淚在稿紙上,他伏案上痛哭,無法繼續寫作.這是他整個人投入,完全的奉獻.所謂奉獻,就是毫無保留,整個生命投入主裡面.奉獻者必需具此精神.神要尋找完全奉獻的人,神就用能力的手抓住他；好像昔日神的手,臨到以西結一樣,一個完全奉獻者聖靈要充滿他,他便為神所用,合神所用.
\newline
\newline
穆勒先生被神大大使用時,有的人嫉妒他,批評他,說:「難道穆勒有神的專利權嗎?」穆勒聽了,反而說句非常寶貴的話:「不是穆勒有神的專利權,是神對穆勒有專利權.」意即整個穆勒的生命都在神的手裡.
\newline
\newline
美國教會歷史,有一位神大大使用的愛德華約拿單,他是耶魯大學的創辦人,他引起了新英格蘭州的大復興.這位主的僕人,在他奉獻時說:「神啊!我放棄自己一切的權利,神啊!你得到我一切的權利.」意即他再也沒有自己的要求,放下了一切,獻上了一切；一切權利都交給神.但願我們對神說:「神啊!我是個完全沒有自己權利的人,我把所有權利,都交在你手中；在你手中,握有我一切的權利.」這樣的人,主要充滿他,主的能力,聖靈的能力,要充滿他.神所用的是這樣的人.
\newline
\newline
請看第二處經文「神的城阿,有榮耀的事乃指著你說的.」(八七3)意思是說耶路撒冷聖城,充滿了神榮耀的作為,詩人很愛慕聖城,歡喜講述聖城的事蹟,在這聖城裡「有耶和華所立的根基在聖山上.」(1)神所立永恆的根基永不搖動,神所揀選,所恩待的錫安,永不動搖.這是詩人對聖城另一欣賞,他想到「神愛錫安的門,勝於愛雅各一切的住處.」(2)神愛錫安,立定了錫安一切的根基,神在錫安城作過許多榮耀的事,所以詩人很響往很愛,很渴慕聖城,他最注意的是「錫安的門.」一般意義上,門代表這城；因為必須經過門進城,所以神愛錫安的門,也就是神愛錫安城.但事實上,門還有更進一步的意義,是門內所發生的事,發生了神恩的作為.「......這一個,那一個都生在其中......」(5)「......個個生在那裡......」(4)這些代表萬民,第六節說明:「當耶和華記錄萬民的時候,祂要點出這一個生在那裡.」所以「這個那個」在這裡是特別的詞句,表示「數目多」和「個別」這兩面都寶貴,大數目固然寶貴；但神也以個別為寶貴.在錫安城耶路撒冷門內,世界許多不同國家的人,都生在那裡,這門是救恩的門.耶路撒冷在此清楚地,是教會預先的影兒；教會中有萬國的人,千千萬萬各地,各方,各民,各族都集中,在神救恩的城裡,所有的人包括在內.
\newline
\newline
我們在教會要建造救恩門,要把救恩傳揚出去.每教會必須是傳福音的教會,不斷得人,讓其他國家民族,基督國度之外的人；藉著福音進到國度裡面,構成了建立了基督偉大的國度.救恩門,福音門,要大大打開,這門要建造寬大又開放.可惜世上有許多教會的門關上,沒有福音,沒有人在那裡得救；一片荒涼,沒有福音的信息,沒有福音的愛；因此也沒有神的同在.主耶穌曾說:「你們往普天下傳福音的時候,我就常與你們同在,直到世界的末了.」主同在的應許,是和大使命的命令連在一起；當我們順從大使命去傳福音之時,主同在的應許就應驗了.傳福音的教會得到主的同在,所以興旺,蒙恩,發展.不傳福音的教會,得不到主的同在.現處境之中,我們的回應是,更加努力,更大積極,更加同心合意傳福音；全港九眾教會眾弟兄姊妹都要起來,動手建造福音門；因為今日的需要,是許多愛靈魂的人.有人說:「愛靈魂」這詞句不大合時代,只愛靈魂不愛身體,不愛人,難道得救後就不理他了嗎?不錯,其中有這樣的分別,我們關心人的全人,身體靈魂都關心,這是絕對的.我們應接受這個批評,但是,我們仍然有個最高的重點,是聖經,是主耶穌,是保羅給我們的重點「救人第一」.使人得救,與神和好,使人出死入生得永遠的生命；使人作神的兒女是第一,餘者是其次.所以我們仍是坦然地說「我們愛靈魂」,我們要救靈魂,這是最基本,最重要的第一工作.
\newline
\newline
救世軍的創辦人霍威廉,在救世軍中是個大將；他在英國開始了救世軍的工作,曾經推動福音工作四十年.福音工作和愛心的服務,打成一片.當時許多英國人都不了解他,他受盡逼迫,遭很大的壓力,遇許多的困難；但他愛靈魂的心不受影響,到處開佈道會,得到很好的果效.到了他工作四十年之後,大家對他才有認識,已往的壓力都消失了.1904年英皇愛德華七世特請他到白金漢宮,向他表示對他工作的欣賞,臨走時英皇請他在紀念冊留言,他寫「有人的雄心是藝術,有人的雄心是名譽,有人的雄心是金錢；我的雄心是靈魂.」他用Ambition或謂野心,他的野心是得靈魂.感謝天父!這樣的人,為了救人願意付代價.
\newline
\newline
為甚麼有的教會,有的弟兄姊妹努力傳福音?其不可缺的原因,就是歌羅西教會努力傳福音的原因,保羅在歌羅西書一章6節稱讚他們真知道神的恩惠.他們接受福音,真知道神的愛,救恩的愛,十字架的愛；他們真知道,真經過,真體會了；因此就發出福音的工作來.如果基督徒自己不知道得救,不曉得自己得救多麼寶貴,不曉得基督救恩十字架的意義,自己沒經歷沒體會過得救；他不可能傳福音,不會真正熱心傳福音,不會有負擔傳福音.真正知道神救恩的受恩者,才會努力傳福音,把自己所得到的和別人分享.
\newline
\newline


\newpage

\section{第57屆~港九培靈會~滕近輝~第十講　熱愛聖殿聖城的詩篇}
\label{sec:499}
\textbf{第十講　熱愛聖殿聖城的詩篇}
\newline
\newline
連結: \href{https://www.hkbibleconference.org/session-message/view/499}{\texttt{https://www.hkbibleconference.org/session-message/view/499}}
\newline
\newline
\hyperref[sec:500]{< < < PREV SERMON < < <}
~
\hyperref[sec:toc]{[返目錄]}
~
\hyperref[sec:482]{> > > NEXT SERMON > > >}
\newline
\newline
詩篇 122:1-122:9
\newline
\begin{longtable}{cl}
\hline
\hline
章節 & 經文 (和合本修訂版)\\
\hline
122:1 & \begin{tabularx}{0.7\textwidth}{X} 我喜樂，因人對我說：「我們到耶和華的殿去。」 \end{tabularx} \\ \\ \relax
122:2 & \begin{tabularx}{0.7\textwidth}{X} 耶路撒冷啊，我們的腳站在你門內。 \end{tabularx} \\ \\ \relax
122:3 & \begin{tabularx}{0.7\textwidth}{X} 耶路撒冷被建造，如同連結整齊的一座城。 \end{tabularx} \\ \\ \relax
122:4 & \begin{tabularx}{0.7\textwidth}{X} 眾支派就是耶和華的支派，上那裡去，按以色列的法度頌揚耶和華的名。 \end{tabularx} \\ \\ \relax
122:5 & \begin{tabularx}{0.7\textwidth}{X} 他們在那裡設立審判的寶座，就是大衛家的寶座。 \end{tabularx} \\ \\ \relax
122:6 & \begin{tabularx}{0.7\textwidth}{X} 你們要為耶路撒冷求平安：「願愛你的人興旺！ \end{tabularx} \\ \\ \relax
122:7 & \begin{tabularx}{0.7\textwidth}{X} 願你城中有平安！願你宮內得平靜！」 \end{tabularx} \\ \\ \relax
122:8 & \begin{tabularx}{0.7\textwidth}{X} 為我弟兄和同伴的緣故，我要說：「願你平安！」 \end{tabularx} \\ \\ \relax
122:9 & \begin{tabularx}{0.7\textwidth}{X} 為耶和華－我們神殿的緣故，我要為你求福！ \end{tabularx} \\ \\
[1ex]
\hline
\hline
\end{longtable}
一連十天的晚上,各位弟兄姊妹對主話語的渴慕,聽道的精神,兄弟得到很大的鼓勵.每晚我心裡都有壓力,我怕的是所講的不能真正幫助弟兄姊妹的靈性；當有這壓力,我就轉個方向來想；工作的是主,施恩的是主；這樣的我的心就安息下來.我們繼續仰望主,今天最後的一晚向我們講話.
\newline
\newline
詩篇中有六篇,具有共同的特點,就是詩人表達他們對於聖殿和聖城的熱愛,同時他們也了解對聖殿,聖城中,最愛慕的一些事物.這些詩篇對我們兩方面的啟發,第一方面是感動我們也愛神的教會.第二方面也感動我們去追求,在我們教會中建立詩人所愛慕的那些事物.
\newline
\newline
今晚我們的重點在於一二二和一三二篇,還有兩篇我們只能簡略一題.
\newline
\newline
詩篇第一二二篇是非常寶貴的,詩人在此表達他對聖殿,聖城的熱愛.第一節詩人說:「人對我說,我們往耶和華的殿去,我就歡喜.」詩人歡喜到神的面前,他愛神的殿；因他愛神,愛神就愛神的殿,歡歡喜喜敬拜,親近神.「耶路撒冷阿,我們的腳,站在你的門內.」(2)這裡是詩人歸屬感的表達,他是屬於神的,是神的子民；神的家就是他的家,神的殿就是也的殿.他把自己投入其中,絕對不作旁觀者；他以神的殿,神的城為榮,將自己與神的殿,神的城認同.
\newline
\newline
這樣的歸屬感是今天教會所需要的.但願每教會的同道,也有如此的歸屬感.自己和教會認同,和眾弟兄姊妹認同在一起.他是唯一美國總統在任內受洗的.我相信他以基督的教會為榮耀,他投入其中,這是個寶貴的見證.另一位美國總統卡特當總統之後,仍然在教會教主日學,他引以為榮.
\newline
\newline
弟兄姊妹!我們是屬神的,屬教會的,教會是我們的家,我們要投入其中,各人有所奉獻；將自己所領受的一切獻上,同心合意滿足主的心.請注意第二節「耶路撒冷阿,我們的腳站在你的門內.」詩人非單獨把自己放在那裡,而是和其他屬神的人放在一起,有團契的感覺.
\newline
\newline
週前,我參加教會夏令營,主題是「委身」.分題有三個,第一向神委身.第二向教會委身.第三向弟兄姊妹委身.我覺得非常好.一個向神委身者,同時一定向教會委身；向教會委身者,同時也一定向弟兄姊妹委身.我們彼此委身,我們都是屬主的人,互相有歸屬感；我們屬於神,屬於神的教會,也屬於弟兄姊妹；彼此建立,彼此關心,彼此鼓勵,彼此幫助.
\newline
\newline
在非洲傳福音的歷史,有兩個最出名的人物,一是李文斯頓,一是馬巴得；這兩位為主作很好的工作.馬巴得年長,李文斯頓年輕,當他想起事奉主之時,他住在愛丁堡；第一次講道時,他站上講台,說了幾句話之後,他的思想一片空白,講章都忘記了.他說:「弟兄姊妹,很抱謙!我講不下去了,我忘記了我的講章.」於是離開講台.這是鼎鼎大名李文斯頓事奉主的第一個經歷；在那聚會中,馬巴得也在場,見此情況,就走到李文斯頓身邊,握著他的手,對他說:「弟兄,這次的失敗,不是最後的失敗.」會後再和也談話:「一個事奉主的人不一定要講道,可以當個醫生事奉主.」後來李文斯頓攻讀醫學,畢業後也行醫也講道.兩全其美.這位主的僕人,安慰鼓勵的一個青年人；這是互相歸屬的作用,屬神者彼此連繫,彼此幫助.
\newline
\newline
第6節提及,愛耶路撒冷的人必然興旺.興旺有兩個意義,第一,蒙恩典.愛耶路撒冷的人蒙神恩待,在各方面亨通,蒙福,靈性會更興旺,整個人更興旺.第二,興旺是人數增加之意.在教會歷史,每時代,每地區,每教會,神都興起了愛祂的人,這些人靈性蒙恩,靈性有根基,愛主的心堅固,事奉的熱心繼續不斷；作了別人的榜樣.神又藉著這些人興起了更多愛主的人.「教會啊!愛你的人必然興旺.」這是歷史的事實,歷史的見證.
\newline
\newline
這些愛耶路撒冷的人「他們求」,「求」用了四次,「你們要為耶路撒冷求平安.」「願你城中平安」,「我要說,願平安在你中間.」為耶路撒冷求平安,這是很重要的.
\newline
\newline
平安有兩個意思,「願你城中平安」在敵人攻擊,侵犯之下得平安；換言之,得神保守看顧,勝過仇敵.第二個意思「平安在你中間」這平安是合而為一的,是和睦的,是彼此相愛的平安.這平安很寶貴!所有的裂口醫好了,所有的紛爭都除去了,同心合意愛神,榮耀神,見證神,事奉神；從這平安產生合一的力量,讓教會興旺.詩人很注重平安,三次連說平安.
\newline
\newline
弟兄姊妹們!教會最蒙恩典之一,就是得到平安,保持平安；如有平安,教會就能夠發展,能進步,有工作的效果,有人信主得救.何時失去平安,撒旦立刻工作；教會聖工受攔阻,受打擊,受挫折.所以詩人渴望平安,平安,平安.求主在我們眾教會賜下平安!弟兄姊妹們,同工們齊心合意把一切磨擦,一切紛爭都除去.
\newline
\newline
愛主者還有一表現,就是第四個求,「因耶和華我們聖殿的緣故,我要為你求福.」求福這個求字的原文,和上面求平安的求不同,前者是祈求,禱告的求,後者是尋求的求.兩種求都是需要的,在禱告中向神求平安,願合而為一的心充滿於教會中.大家竭力保守聖靈所賜合而為一的心；還要尋求教會的福,這福是好處,是益處；所有教會的弟兄姊妹,同心合一地求教會的益處,讓教會更蒙恩.這兩種求合起成為完全的求.
\newline
\newline
詩人如此愛慕聖殿,聖城,他們最欣賞的事物,連絡整齊的房屋「耶路撒冷被建造,如同連絡整齊的一座城.」(3)詩人居高下望,耶路撒冷城很美麗；因為房屋都是有計劃的,互相連絡整齊,很美觀.神也是欣賞教會的弟兄姊妹互相配搭,互相連繫的團契生活；教會井井有條,團契互相合作,按照聖靈所賜的恩賜,一同事奉.每人伸出手來,用神所給的恩賜來工作.神在天上看見教會良好的配搭,神的心快樂欣賞.我們應建立團契生活,讓教會如同耶路撒冷,是座連絡整體的城,我們要發掘,要運用,要增進我們的恩賜,讓恩賜作用更加發揮.別讓恩賜生鏽,別把恩賜的銀子埋沒地下；該要忠心運用,神給予我們的恩賜.保羅提及恩賜時說:「飲於一位聖靈.」他把聖靈當作泉源,時時飲取.聖靈是我們恩賜之源,我們得到聖靈的恩賜以後,要常到源頭飲取,人喝了水,精神振作.所以英文喝點水,吃點餅乾,謂之refreshment意即使精神振作.我們的恩賜也需要到恩賜的源頭聖靈那裡去重新飲取,得到振作,得到復興；不然,恩賜的作用就慢慢地消退,功效逐漸地減少.
\newline
\newline
在互相團契的生活中,第一節告訴我們美好的作用.這裡提及「我」,「人」,「耶和華」,「耶和華的殿」四方面.不但有我,還有其他屬神者,其他屬神者對我說:「我們往耶和華的殿去.」我一聽之下就歡喜.其他屬神者是愛神的人,所以歡喜到耶和華的殿去；他們也希望我去,提醒我；我聽了,愛神的心就挑旺起來,於是歡然同往.人來鼓勵我,我響應了第兄姊妹的鼓勵,一齊愛神,愛神的殿,歡喜親近神,敬拜神,事奉神,榮耀神,多麼美好!這四方面連在一起,就是團契的作用；整個的功效就是「歡喜」.屬靈的喜樂是從團契的交通而產生的.
\newline
\newline
詩人欣賞大衛家的寶座,大衛家的寶座乃預表基督的寶座.耶路撒冷聖城中有寶座.教會也應有基督的寶座,我們要一同建立基督的寶座,要擁戴基督作王,要接受基督的主權,要順服祂,高舉祂,榮耀祂；祂坐在祂國度的寶座上.保羅說:「因為神的國度不在乎這言語,乃在乎權能.」神國的要素是權能,當國民順服之時,權能就顯出來；當我們這些在基督裡面的人順服基督之時,基督的權柄就顯出來,好處就顯出來.基督的國度在於我們的順服,在於我們接受祂的主權,在於我們常對主說:「是的,是的,是的.」
\newline
\newline
彼得後書的主題論及「長進」特別四次提及主耶穌完全的名字「我們的主!救主耶穌基督.」這是彼得後書最明顯,最重要的特點,主耶穌,「主」的身份和「救主」的身份連在一起；祂不但是我們的救主,也是我們的主；我們和主是雙重的關係,第一步我們接受耶穌作救主,當我們靈性長進時,接受我們的主耶穌作我們的主人；這時,祂就是我們的主,救主耶穌基督,我們的靈性此時開始長進,從信徒進為門徒,這是我們靈程的經歷.我們要在教會中建立基督的寶座,學習作門徒,學習順服,學習否定自己,確定基督.
\newline
\newline
詩人欣賞耶路撒冷「宮內」(7句末)「願你宮內興旺.」宮內是領袖的住處,宮內興旺則領袖們蒙恩.如果要教會蒙恩,第一步要教會領袖蒙恩；如果他們的靈性蒙恩,整個教會就蒙恩.領袖是走在前頭的人,領袖有三個因素:第一,看見神的旨意.第二:遵行神的旨意.第三,推動鼓勵別人遵行神的旨意.領袖有所看見,有所推動；領袖蒙恩,宮殿興旺,全耶路撒冷蒙恩,興旺；全教會蒙恩,興旺.
\newline
\newline
求神在眾教會中興起領袖!
\newline
\newline
詩人對神的殿愛慕(一三二1-5)為了建立神的殿,寧願廢寢忘食,這是詩人的態度；他讓神的殿居首位,視為最重要；自己的一切都是次要.這種心態何等難得,結果他成功了,為神預備了建殿的材料；後來神藉著所羅門建成聖殿.
\newline
\newline
對神有強烈渴望的人,神一定用他,藉他顯出能力,顯出作為,行奇妙的事,甚行神蹟.我們的願望,渴望,非常重要.(詩一○三篇5節)那幅「如鷹返老還童」復興的圖畫,詩人指出,神「使你所願的得以知足,以致你如鷹返老還童.」「願」則願望,渴慕,追求；當我們有強烈的渴慕時,神就動工,使我們得到滿足,使我們如鷹返老還童,得到復興.每當復興蒙恩之先,一定有渴慕；否則沒有個人的復興；神在行神蹟彰顯大作為之前,我們一定先有渴慕；否則這一切都沒有.照著你所信的所渴慕的給你成全,這是神工作的二大原則.在整部聖經中,整個教會歷史中,神一直按這二大原則顯出作為.
\newline
\newline
詩人想到聖殿時,最注重的是神的約櫃；以色列人失去約櫃,他認為這是最可惜,可憐的事.後來找到了約櫃,就把神的約櫃帶進聖殿,帶進他們中間,他們得回神的約櫃.
\newline
\newline
神的約櫃基本意義有三,第一,約櫃代表神的同在.舊約歷史給我們看見,何時以色列人失去約櫃,他們就失去了神的同在,就落在敵人手中,他們的靈性敗落,國運敗落,受敵人壓制,落在敵人的權勢下.神的同在是最重要最寶貴的.教會有神的同在,甚麼都有了,失去神的同在,一切都失去.我們讀聖經,每有「神與他們同在」或「主與他們同在」一切問題解決了,教會就大大興旺,聖工就作成了.
\newline
\newline
第二個意義,約櫃代表神的引導.舊約時代以色列人在曠野,約櫃走在前面,有藍色的氈蓋頂；若居高下望,藍色的點在以色列人前面,所有的人都跟隨藍點前行,約櫃停任,以色列人也停住.(出廿五15)特別提到約櫃上面有環,環內有杠不可抽出,要一直留在環內.其明顯的意思,則隨時要抬起約櫃,準備隨時行動；聽從神隨時的命令和引導.
\newline
\newline
感謝神!神是活的神.在整個教會歷史,神一直引導,神的靈不斷工作.約櫃起行,我們都當跟著走,約櫃等於神指揮的手；神的手一動,教會就當奔走前程.三百萬以色列人齊跟隨約櫃走,浩浩蕩蕩,神在前頭,以色列人安全,成功前進.跟定主的腳步,祂是元帥,祂是先鋒；祂的路就是我們應走的路.
\newline
\newline
第三個意義,約櫃代表神與祂的百姓立約,這約就是耶和華作他們的神,他們作耶和華的子民.這約是永遠的約.
\newline
\newline
我們想到新約,當主被賣那夜「主拿起餅來擘開說,這是我的身體為你們捨的,你們每逢如此行,為的是記念我.」又拿起杯來說:「這是立新約的血,流出來使罪得赦.」主破碎的身體,流出寶血為我們立了約,主和我們,我們和主立約；主的寶血就是立約的印記.這約多麼寶貴!主為我們死,我們為主活.正是保羅所說:「一人既替眾人死,眾人就都死了,並且祂替眾人死,是叫那些活著的人,不再為自己活,乃為替他們死而復活的主活.」(林後五14-15)在這些話之前有個「想」字「原來基督的愛激勵我們；因我們想......」原文「想」字有特別意思,是經過思考以後作了選擇的決定.主一人替我們死了,叫那些活著的人(包括我在內)不再為自己活,乃為替我們死而復活的主活.這是保羅思考後邏輯的決定.英文聖經用judge表達原文的意思.各位!保羅說:「我們想.你曾經想過嗎?」保羅作此選擇和決定；如果你沒有這個選擇和決定,下面所講的和你毫無關係；因保羅有這決定,所以下面的話成為他人生的原則,這是立約.主對我們,我們與主立了約.你對主立了約嗎?你守這約嗎?這約在你有意義嗎?你是否在這約裡面?你作了這選擇的決定嗎?
\newline
\newline
這個約櫃,我們當在教會建立起來!讓主耶穌基督的愛繼續激勵我們,我們與主立了約,得主同在,得主引導.有了這三樣意義,我們甚麼都有了,這樣的教會,是前進的,是神能力彰顯的,是為基督而活的教會,如此教會真是有福了.
\newline
\newline
詩人還有所愛慕,詩人特別提及聖殿裡有明燈.聖殿的明燈一方面預表基督是世界的光.另方面預表基督徒是世界的光.明燈是我們好行為的燈,我們反照基督的品格和榮耀,在世上為主發光.教會要建立明燈,建立我們的好行為,生活的見證,基督徒在社會中,要發福音的光,福音的光一定要發!
\newline
\newline
在世運會有火炬,賽跑,接力,這都是按照二千多年前古希臘運動會的項目.那時有種比賽是三部份混合,首先所有參賽者環繞大火炬各人手執小火炬點火.第二部份是賽跑,跑時要快並保持火炬不熄滅.第三部份,一個跑完要把火炬遞給下一個跑者繼續接力.三部份合起,就是比賽的節目,時至今日成為世運會的傳統.
\newline
\newline
基督徒要執福音火炬繼續快跑,接棒並交棒,代代相繼傳福音,使福音傳遍地極；我們的火把不能熄滅,使福音的光和好行為的光同時並發.主耶穌說:「你們是世上的光,世上的鹽.」鹽是隱藏的,代表每個基督徒的好行為在社會中,是個人性質,在生活中發出作用,外表不見其作用,是實在的作用.基督徒的好品格和光,是給人看的；主說:「城造在山上不能隱藏,你們的好行為也當照在人前,叫人看見,歸榮耀於天父.」(太五16)我們有隱藏的好品格,有日常生活的好品格,也有顯露出來的好品格；這是我們愛心的行為,是基督徒團體的好行為；在整個教會團體性的服務,在社會中見證；在社會中得愛心的表現,這是要人看的,叫人知道歸榮耀給神.隱藏的,明顯的合在一起,就是我們要建立的明燈.
\newline
\newline
我們要建立祭壇,建立福音門,建立聯絡整齊的城,建立基督的寶座,建立領袖,建立約櫃,建立明燈.
\newline
\newline
詩一三七篇是詩人感情最強烈的表達,為了錫安,為了聖殿.
\newline
\newline
還有第四十八篇,以上兩篇大家可以作家課；不但參加十天的聚會,還帶了家課回去.
\newline
\newline


\newpage

\section{第57屆~港九培靈會~艾榮~第一講　如今就不被定罪了}
\label{sec:482}
\textbf{第一講　如今就不被定罪了}
\newline
\newline
連結: \href{https://www.hkbibleconference.org/session-message/view/482}{\texttt{https://www.hkbibleconference.org/session-message/view/482}}
\newline
\newline
\hyperref[sec:499]{< < < PREV SERMON < < <}
~
\hyperref[sec:toc]{[返目錄]}
~
\hyperref[sec:483]{> > > NEXT SERMON > > >}
\newline
\newline
羅馬書 8:1-8:4
\newline
\begin{longtable}{cl}
\hline
\hline
章節 & 經文 (和合本修訂版)\\
\hline
8:1 & \begin{tabularx}{0.7\textwidth}{X} 如今，那些在基督耶穌裡的人就不被定罪了。 \end{tabularx} \\ \\ \relax
8:2 & \begin{tabularx}{0.7\textwidth}{X} 因為賜生命的聖靈的律，在基督耶穌裡從罪和死的律中把你釋放出來。 \end{tabularx} \\ \\ \relax
8:3 & \begin{tabularx}{0.7\textwidth}{X} 律法既因肉體軟弱而無能為力，神就差遣自己的兒子成為罪身的樣子，為了對付罪，在肉體中定了罪， \end{tabularx} \\ \\ \relax
8:4 & \begin{tabularx}{0.7\textwidth}{X} 為要使律法要求的義，實現在我們這不隨從肉體、只隨從聖靈去行的人身上。 \end{tabularx} \\ \\
[1ex]
\hline
\hline
\end{longtable}
羅馬書第八章是一座屬靈的高山,講出神在過去,現在和將來所成就的工作,我稱這一章的題目,為貫徹始終的救恩.
\newline
\newline
今天我們要看的,是第一到第四節:「如今那些在基督耶穌裡的,就不定罪了.因為賜生命聖靈的律,在基督耶穌裡釋放了我,使我脫離罪和死的律了.律法既因肉體軟弱,有所不能行的；神就差遣自己的兒子,成為罪身的形狀,作了贖罪祭,在肉體中定了罪案.使律法的義,成就在我們這不隨從肉體,只隨從聖靈的人身上.」
\newline
\newline
這四節經文講到四個要點:第一是基督徒的地位,一到二節.第二是基督徒所得的供應,第三節.第三是神對在基督裡的人的目的,第四節上.第四是神目的的實行,第四節下.
\newline
\newline
我們先看第一,第二節所講的地位,這地位是那些在基督耶穌裡的人才得著的.一切在基督耶穌裡的人,都是相信接受耶穌為救主的人,把他們的生命交在主大能的手中.
\newline
\newline
今天,很多基督徒對信心有不同的看法和爭論.有人以為信心相信某些事情將要發生,那麼事情就真的會發生,這就叫做信心.所以人只要誠心的相信自己已經得救,他便可以得救.但聖經的真理卻不是這樣.真正的信心,是把自己完全交在神的手中,深信神會成就祂所要作的,而我們就因此得救.我們必須來到基督面前,把自己交在祂的手中.雖然我們不能完全明白神所作的,但是我們信靠祂,因為我們需要祂的幫助.所以我們把自己交在祂的手中,活在祂的管理和主權之下,這就是把人帶到基督裡的意思.
\newline
\newline
羅馬書八章一節講到那些在基督耶穌裡的人,如今就不被定罪了.這一節經文的開始,有「所以」這個意思作為開始,可以說是上面七章聖經的一個總結.「所以」如今那些在基督耶穌裡的,就不定罪了.在羅馬書第八章之前,是講到定罪的.從一章十八節到三章二十節,一共六十四節經文,講到一切活在地球之上的人,都在定罪之下；直到他歸向基督之時為止.
\newline
\newline
一章十八節講到所有的世人,都應該從大自然中認識神.人在他的良知中,應該知道善惡對錯,可是人卻拒絕依照他的良知去行.即使他們有聖經的亮光和指示,也不照著去行,所以在定罪之下.因此三章十九節說:所有人在律法之前,都定為有罪.「因為世人都犯了罪,虧缺了神的榮耀.」全能的審判者,宣告人類都在定罪之下.這是一個非常嚴重的問題.被神定罪是一件可怕的事,一章十八節說,神的忿怒,從天上顯明在一切不虔不義的人身上.
\newline
\newline
可是,從三章二十一節開始,便講到那些真心相信耶穌,把自己交主大能手中的人,便從定罪中得著釋放.耶穌能從人的身上取去罪債,祂的公義臨到人身上,叫人在神眼中得稱為義.耶穌為我們還清罪債,負起一切我們該負的責任,又把我們推薦到神面前,使神接納我們好像一個從未犯錯的人.當人在接受耶穌作救主的一剎那,耶穌為我們代罪的事便立即成就在這相信的人身上,使被定罪的人立即變成清白無罪,還可以在神面前稱義的人.這是羅馬書三到五章的真理.
\newline
\newline
第五章說,我們既得以稱義,就可以與神相和,可以站在神面前得蒙悅納.我們進入神的恩典中,接受神所給我們一切的恩惠.我們又在神榮耀的喜樂中有所享受,聖靈將神的愛澆灌在我們心中.
\newline
\newline
到羅馬書八章的時候,就把前面所講的一切恩典作個總結說:如今那些在基督耶穌裡的,就不定罪了.這包括了我們的罪得蒙赦免,在神面前得稱為義,又領受神一切的恩惠.
\newline
\newline
神不再定我們的罪,這是一個何等大的恩典,神完全赦免我們.不再記念我們的罪,並且為我們預備一個永恆快樂的天家.試想,人間的罪犯,如果一旦獲赦,已經是高興非常,更何況我們在聖潔公義的神面前,被宣判為無罪哩!
\newline
\newline
至於羅馬書八章二節,也是描述同一真理,不過是以另一些字句來表達而已.這裡講到賜生命聖靈的律,就是神的律法,把人帶到定罪,滅亡和刑罰的法則.哥林多前書十五章五十六節,講到死亡的毒勾是罪,而罪的權勢就是律法.罪之所以成為罪,是因為有律法.神在十誡和律例中表明了祂的律法；這些律法就顯出了我們的罪,而罪所帶來的就是死.羅馬書第八章二節所講的律,就是定我們有罪的神的律法,除非我們轉向神,否則我們仍在律法之下.
\newline
\newline
神的律法就如將我們放在顯微鏡之下,細心的察驗,好決定我們是否罪人,是否該受神撇棄.但羅馬書八章二節又講到另一種律,就是賜生命聖靈的律.這是在基督耶穌裡一種新的法律.我們從舊的律出來,進到神新的律中,得著生命,得著聖靈.新的律把我們從舊的律中拯救出來,這是何等的福氣,一種極寶貴的經歷.如今那些在基督耶穌裡的,就不被定罪了,因為賜生命聖靈的律,在基督耶穌裡釋放了我,使我脫離罪和死的律了.這是信徒所站新的地位,不再定罪,在基督裡得著自由.
\newline
\newline
第三節講到我們怎樣從神那裡得著供給,以致得享自由.一個在受定罪,在死亡和罪的律之中的人,怎能進到基督所賜的自由裡呢?羅馬書八章三節第一開始先講到反面,我們在耶穌裡不能得到甚麼,然後才從正面去講.第三節說:律法因肉體軟弱,有所不行的.有人以為只要他盡力去守律法,在神面前就可算為對的.但這裡清楚的說明,我們因為肉體的軟弱,是不能守全律法的.我們無法憑自己的力量去守全律法,博得神的喜悅.
\newline
\newline
不過,有另外一個方法,可以使我們在神面前稱義.神差遣祂的兒子,成為罪身的形狀,作了贖罪祭,在肉體中定了罪案.耶穌為我們受死,擔當我們一切的罪,和應受的咒詛,刑罰,使我們在神面前可以稱義,不再被定罪了.這是神所預備的方法,不是我們靠自己的力量可以得到的,但是神的供給,使我們在基督裡可以享有一個不被定罪的地位.這是基督徒經歷的開始.
\newline
\newline
到第四節,便講到前面三節所講的,都是連繫著一個目的.在神救贖的恩典中,我們有一個新的地位,而這地位和神的目的是有關係的.我們從定罪中得著釋放,在耶穌裡所賜的自由,使律法的義,成就在我們身上.有人以為我們藉著耶穌得稱為義；在稱義之後所過的生活,便與神無關.但事實上,神非常關心我們稱義之後的生活,他要我們活出神的義來.如果我們要過義的生活,就不能再靠肉體,乃要靠聖靈而活了.
\newline
\newline
明天我們要看不靠肉體生活,專靠聖靈生活的真理.我們要指出,是神的目的,叫我們成為公義的子民.
\newline
\newline
讓我們溫習我今天所講的.羅馬書八章講到救恩從開始到末了的情形；起點就是現在進到耶穌基督內,把我們的生命交在耶穌基督的手中；就如病人把自己完全交在醫生的手中一樣.當我們相信的那一剎那,靠著神的恩典,我們便從定罪中得著釋放,從罪惡的律和原則中得著釋放.
\newline
\newline
請問你有經歷過這樣的時刻嗎?你確知自己已經進入基督耶穌裡嗎?如果你已經進到耶穌基督內,那麼今天的經文便正好用在你身上,你就不再被定罪了.你從罪惡和死亡的律中得著釋放,神要領你往前走,過公義的生活.
\newline
\newline
為著神救贖我們的真理,我們要感謝祂,讓我們因著得到神的救恩而喜樂,讓我們決心遵行真理,從道理的開端,進入到公義的生活去.
\newline
\newline


\newpage

\section{第57屆~港九培靈會~艾榮~第二講　公義的實踐}
\label{sec:483}
\textbf{第二講　公義的實踐}
\newline
\newline
連結: \href{https://www.hkbibleconference.org/session-message/view/483}{\texttt{https://www.hkbibleconference.org/session-message/view/483}}
\newline
\newline
\hyperref[sec:482]{< < < PREV SERMON < < <}
~
\hyperref[sec:toc]{[返目錄]}
~
\hyperref[sec:484]{> > > NEXT SERMON > > >}
\newline
\newline
羅馬書 8:1-8:4
\newline
\begin{longtable}{cl}
\hline
\hline
章節 & 經文 (和合本修訂版)\\
\hline
8:1 & \begin{tabularx}{0.7\textwidth}{X} 如今，那些在基督耶穌裡的人就不被定罪了。 \end{tabularx} \\ \\ \relax
8:2 & \begin{tabularx}{0.7\textwidth}{X} 因為賜生命的聖靈的律，在基督耶穌裡從罪和死的律中把你釋放出來。 \end{tabularx} \\ \\ \relax
8:3 & \begin{tabularx}{0.7\textwidth}{X} 律法既因肉體軟弱而無能為力，神就差遣自己的兒子成為罪身的樣子，為了對付罪，在肉體中定了罪， \end{tabularx} \\ \\ \relax
8:4 & \begin{tabularx}{0.7\textwidth}{X} 為要使律法要求的義，實現在我們這不隨從肉體、只隨從聖靈去行的人身上。 \end{tabularx} \\ \\
[1ex]
\hline
\hline
\end{longtable}


\newpage

\section{第57屆~港九培靈會~艾榮~第三講　公義與我們的頭腦}
\label{sec:484}
\textbf{第三講　公義與我們的頭腦}
\newline
\newline
連結: \href{https://www.hkbibleconference.org/session-message/view/484}{\texttt{https://www.hkbibleconference.org/session-message/view/484}}
\newline
\newline
\hyperref[sec:483]{< < < PREV SERMON < < <}
~
\hyperref[sec:toc]{[返目錄]}
~
\hyperref[sec:485]{> > > NEXT SERMON > > >}
\newline
\newline
羅馬書 8:5-8:8
\newline
\begin{longtable}{cl}
\hline
\hline
章節 & 經文 (和合本修訂版)\\
\hline
8:5 & \begin{tabularx}{0.7\textwidth}{X} 因為，隨從肉體的人體貼肉體的事；隨從聖靈的人體貼聖靈的事。 \end{tabularx} \\ \\ \relax
8:6 & \begin{tabularx}{0.7\textwidth}{X} 體貼肉體就是死；體貼聖靈就是生命和平安。 \end{tabularx} \\ \\ \relax
8:7 & \begin{tabularx}{0.7\textwidth}{X} 因為體貼肉體就是與神為敵，對神的律法不順服，事實上也無法順服。 \end{tabularx} \\ \\ \relax
8:8 & \begin{tabularx}{0.7\textwidth}{X} 屬肉體的人無法使神喜悅。 \end{tabularx} \\ \\
[1ex]
\hline
\hline
\end{longtable}
(羅八5-8)「因為隨從肉體的人,體貼肉體的事,隨從聖靈的人,體貼聖靈的事.體貼肉體的就是死亡,體貼聖靈的乃是生命平安,原來體貼肉體的,就是與神為仇,因為不服神的律法,也是不能服.而且屬肉體的人,不能得神的喜歡.」
\newline
\newline
這段經文講到基督徒如何脫離肉體,過著順從聖靈的生活.在這裡有一個非常基本的觀念,是講到我們怎樣使用自己的頭腦,就是我們思想,感受,作決定的部份.我們怎樣使用頭腦,跟我們每天所過義的生活有絕對的關係.我們思想的焦點,關係到我們是順從肉體,抑或順從聖靈.
\newline
\newline
第五節講到有兩種可能,一是順從肉體,原來的意思是把思想都放在肉體上.一是順從聖靈,一切的思想都以屬靈為焦點.聖經所講的,不是我們偶然所想的事,乃是用心去思想鑽研的事.有人喜歡打高爾夫球,便整天想著怎樣行,又經常去打；他所看的書,講的話題,都離不開高爾夫球.我也喜歡打高爾夫球,偶然也會想到高爾夫球,但絕對不是最重要的事.經文所講的,是我們的腦不斷想到的事.
\newline
\newline
其中一個可能,是用我們的腦,專門去想肉體的事.聖經所講順從肉體,就是順從自己的私慾,是驕傲的,隨從自己的感官的.使徒約翰形容我們所處的世界,是充滿肉體的情慾,眼目的情慾,並今生的驕傲.簡單來說,是包括肉體的,感官的和驕傲的.隨從肉體就是讓思想隨從肉體所喜好的,感官覺得舒服的,驕傲的事.在我們四周有許多引誘,叫我們裡面的私慾跑出來.譬如說,一個推銷汽車零件的廣告,竟然要用一個穿得半裸的女郎來拍攝,作它的宣傳手法.時下的小說,又多描寫非道德的,敗壞的愛情.那些從事廣告事業的人,寫小說的作家；知道人裡面有順從肉體的傾向,便投其所好來招生意.當人講到兩性的吸引,身體的享受,生活的奢侈時,人裡面的情慾,便自然的發動.只要我們容許自己的頭腦,去思想這些隨從肉體的事；我們在那一剎那,便讓自己隨從了肉體.
\newline
\newline
幾年前,我和我的姐夫計劃和家人去度假,我們決定去美國北部的明尼蘇達州；那州的北部有許多湖,我們決定到那裡度假和釣魚.當假期愈來愈近的時候,我腦海中便不停的想著度假；想著釣到了大魚時的得意景象,當地的湖光山色,山明水秀的美景.雖然我仍上班,仍繼續看聖經；但腦袋卻不停的想釣魚,祈禱的時候心不在焉,想著怎樣去釣魚.總之,我整個腦袋就被這些東西佔據著.其實,放假或者說享受假期並不是錯；但這些東西卻吸引了我身體的興趣,叫我的思想隨從了肉體.
\newline
\newline
(羅八5)也講到我們所思想的,是一些以自我為中心的事.我們裡面有擁有東西的慾望,即使是嬰孩；他也抓住搖床中的玩具.當人愈大,他抓住東西的慾望便愈大.每一件我們看見的事,都想據為己有.這是我們裡面的本性.
\newline
\newline
我現在是從事編輯教會雜誌的工作,在教會的刊物中,會報導世界各地宣教士的來函.有一次我們收到一篇有趣的文章,有一對夫婦在印尼一個原始的部落作宣教的工作.幾年前,這些人的生活環境,仍處於石器時代,現在已經有了少許的進步,有些日用的物品.最近政府到這對宣教同工工作的地方,徵召當地的一些士人來修築一條馬路.政府託宣教士給當地的工資,就是肥皂,火柴,鍋,鑊等東西.可是當地的工人對宣教士給他們的工資,大衛氣憤和不滿；他們看見宣教士有照相機,錄音機,所以也們要求這些作為工錢.
\newline
\newline
人的本性不住的希望得到東西,許多人只是為了賺錢而努力.如果我們腦海中所想的,也是這些,那麼我們就是隨從肉體而生活.試想,如果有部電腦來記錄我們一星期所思想的東西,我們就會知道自己終日所想的是甚麼.很多時候,我們想的都是物質,都是想要得些甚麼,買些甚麼而已.這是體貼肉體的生活.另外,也有些人常常想著一些以自我為中心,驕傲的事.到底別人怎樣看我們,我的儀表,態度,言談是否出眾,能否得到別人的讚賞,愛戴等等.我怎樣才可以更有吸引力?更得人稱讚?這些思想,都是屬肉體的.這是第一種可能,我們的頭腦,可以思想那些滿足我們肉體的事物,以肉體的事作為思想的焦點,這就是順從肉體.
\newline
\newline
這樣思想的結局是甚麼呢?第六節說,體貼肉體的就是死.這不是身體的死.我思想到明尼蘇達州去度假,這思想絕不會導致身體的死亡.常常想著賺錢的,也不會導致身體的死亡.這裡所說的是屬靈的死亡.第七節就解釋了第六節的上半部說,原來體貼肉體的,就是與神為仇,那些專門想到如何隨從肉體,增加今生的驕傲,這些事是神不會去想的.這樣的思想,使我們離開了與神的相交,而且與神的思想背道而馳.
\newline
\newline
我生長在一個窮苦的家庭,父親只是一間工廠的小職員.為了生計,他週末便要兼職.有一次,父親找到一份髹漆的工作,在週末的時候便叫我一起去幫助.原來那天我已經約好了朋友去游泳的,但是我又必須幫忙父親,所以我人雖然在髹漆,但思想卻去了游水；結果做得很勉強,而且又覺得特別熱,在忙亂中倒得一地都是漆油.終於父親很不高興,便打發我走.父親想的是髹漆；但我想的卻是游水,兩種思想是相反.不合的.同樣,當我們想肉體的事的時候,我們就沒有想屬神的事,我們與神之間便有了對峙.
\newline
\newline
真正的生命是與神有相交的.(約十七3)說,認識神的兒子,就是永生.只有我們與神相交的時候,我們才可以經驗真正的生命.我們為自己屬靈的生命所作最好的事,就是把它帶到神面前,放在神的手之下.如果我們的頭腦充滿肉體的事,我們便與神的思想遠離,是與神為仇了.第七節說屬肉體的頭腦是不能服神的律法的.如果我們將自己的思想放在屬肉體的事上,我們最後便會進到罪惡去.時常想吃的人,終歸會變得大食,犯食.時常想著兩性之間的事的人,終歸會犯上這樣的罪.想賺錢的人,最後可能連不誠實的方法也拿出來用,目的是為了賺錢.罪將人的生命摧殘,降價.所以,思想屬肉體的事,會將我們帶進死亡去.
\newline
\newline
但我們可以有另一個可能.作為基督徒,我們可以每日體驗屬靈的事,將思想集中在屬靈的事上.所謂體貼聖靈,就是按照我們靈魂的好處的事情去思想；意思就是想及我們與神之間的關係,也表示我們想到罪惡的可怕!主的榮美,福音的真理和禱告有關的事,世人的失喪,基督徒的弟兄姊妹等等.當然,我剛才所提的這幾樣,並不代表屬靈思想的全部,但是屬靈的思想是自然包括了這些項目.這不是說,你不能再作一個好秘書,好學生,這和屬靈的思想是沒有衝突的.
\newline
\newline
我用一個例子來說明.巴里是美式足球隊一個非常傑出的球員,在一次全國性賽事中,他所得的分數,打破了以往的紀錄.在完場的時候,球證便把他們所打的那個美式足球頒贈給他.當時,群眾起立向他鼓掌,歡呼!而他在接受別人向他道賀的時候,他說,這真是一個偉大的日子.別人以為巴里是指他在球壇上的成就偉大.但巴里卻說,他所以為偉大的,是今天早上有十七個青少年,因為參加他們的祈禱會,聽了他們的見證而決志信主.雖然巴里身在球場,但是他所思想的,卻是神的事.我們可以作一個好學生,好教員,好主婦,但腦海裡卻思想神的事情.
\newline
\newline
當我們思想屬靈的事的時候,結果就是第六節所說,體貼聖靈的乃是生命平安,和死亡是剛好相反的.聖經說,人的思想如何,為人也如何.你的思想怎樣,你的行為,人生就反映出來.可惜我們太多想到自己肉體的愛好,別人對自己的看法,而少想到真正屬乎聖靈的事.我們要學習用正確的方法來使用自己的頭腦.
\newline
\newline
到底我們應該怎樣使用自己的頭腦呢?我們過去一天,一星期所想的,所關心的是甚麼呢?除非我們用正確的方法去使用頭腦,否則我們沒法去遵行神的律法；求主幫助我們,把思想放在屬靈的事上.
\newline
\newline


\newpage

\section{第57屆~港九培靈會~艾榮~第四講　煩惱的身體}
\label{sec:485}
\textbf{第四講　煩惱的身體}
\newline
\newline
連結: \href{https://www.hkbibleconference.org/session-message/view/485}{\texttt{https://www.hkbibleconference.org/session-message/view/485}}
\newline
\newline
\hyperref[sec:484]{< < < PREV SERMON < < <}
~
\hyperref[sec:toc]{[返目錄]}
~
\hyperref[sec:486]{> > > NEXT SERMON > > >}
\newline
\newline
羅馬書 8:9-8:13
\newline
\begin{longtable}{cl}
\hline
\hline
章節 & 經文 (和合本修訂版)\\
\hline
8:9 & \begin{tabularx}{0.7\textwidth}{X} 如果神的靈住在你們裡面，你們就不屬肉體，而是屬聖靈了。人若沒有基督的靈，就不是屬基督的。 \end{tabularx} \\ \\ \relax
8:10 & \begin{tabularx}{0.7\textwidth}{X} 基督若在你們裡面，身體就因罪而死，靈卻因義而活。 \end{tabularx} \\ \\ \relax
8:11 & \begin{tabularx}{0.7\textwidth}{X} 然而，使耶穌從死人中復活的神的靈若住在你們裡面，那使基督從死人中復活的，也必藉著住在你們裡面的聖靈使你們必死的身體又活過來。 \end{tabularx} \\ \\ \relax
8:12 & \begin{tabularx}{0.7\textwidth}{X} 弟兄們，這樣看來，我們不是欠肉體的債去順從肉體而活。 \end{tabularx} \\ \\ \relax
8:13 & \begin{tabularx}{0.7\textwidth}{X} 你們若順從肉體活著，必定會死；若靠著聖靈把身體的惡行處死，就必存活。 \end{tabularx} \\ \\
[1ex]
\hline
\hline
\end{longtable}
(羅八9-13)「如果神的靈住在你們心裡,你們就不屬肉體,乃屬聖靈了.人若沒有基督的靈,就不是屬基督的.基督若在你們心裡,身體就因罪而死,心靈卻因義而活.然而叫耶穌從死裡復活者的靈,若住在你們心裡,那叫基督耶穌從死裡復活的,也必藉著住在你們心裡的聖靈,使你們必死的身體又活過來.弟兄們,這樣看來,我們並不是欠肉體的債,去順從肉體活著.你們若順從肉體活著必要死,若靠著聖靈治死身體的惡行,必要活著.」
\newline
\newline
這段經文的中心字眼是「身體」和「肉體」不同,頭九節講到「肉體」的有十一次,意思是指人的本性.但十到十三節卻用「身體」,就是我們的身體.其中有三句說話,是和使用我們的身體有關的.
\newline
\newline
第十節說「我們的身體死了」,第十一節說「那叫耶穌從死裡復活的靈,使我們必死的身體又活過來.」第十三節說「我們必須治死身體的惡行」.十節所講的,是一個人生的問題,十一節是神的應許,十三節是我們使用身體的程序.今天我們要思想這三方面的真理.
\newline
\newline
首先,第十節講出一個和我們身體有關的問題,我們的身體死了.我們知道,羅馬書是寫給基督徒的,這裡所說的,是基督徒與神的正確關係,從今以後,我們所過的,是公義的生活,但我們仍在身體的軀殼內,而這身體是會死的.這是甚麼意思呢?這死的身體,原來是一個比喻.
\newline
\newline
這裡所說的死,可以包含兩方面的意思,就身體上和道德上的死.先講身體的死,我們自生下來,便朝向死亡走.我們的細胞不住的消盡,器官不住的用盡,最後整個人的氣力消失.在人生的頭二十五年中,人體增新的速度要比退化的速度快.所以頭二十五年中,我們一天一天的成長,一日比一日強壯.雖然其中也有生病和軟弱的時刻,但基本上是增長的.可是二十五歲以後,退化的程序卻比更新的程序快.比方說,人比較容易疲倦.視力開始衰退,氣力也較弱,身體也會縮小.我們正向死亡前進.這種身體的衰老會影響我們過公義的生活.比方說,我們病的時候,很難集中敬拜神,又容易發怒,難以喜樂,更不能服事神.這是身體上朽壞的過程,能攔阻我們過公義的生活.
\newline
\newline
但還有另一種死,就是道德上的死.(羅七22-23)「因為按著我裡面的意思,我是喜歡神的律,但我覺得肢體中另有個律,和我心中的律交戰,把我擄去,叫我附從肢體中犯罪的律.」保羅講到在他的肉身中,有一種罪的律,使他的身體犯罪.比方說,飲食是正常的,不是犯罪,但過量的飲食,超過身體的需求而仍暴飲暴食,那就是罪了.在婚姻中的性慾和性行為是對的；但若人不滿足在這婚約中的性生活,要在婚姻以外進行的話,那就是罪了.這是罪的律在他裡面發動,叫他有罪的行為的.
\newline
\newline
雖然我們很多時候已經可以靠神的力量去約束自己的身體,可是在我們裡面,仍然有罪的律存在,叫我們蠢蠢欲動.特別是過往我們的嗜好,會成為一種陰影,成為一種罪欲,阻礙我們過公義的生活.雖然神的義,應該每天顯在我們身上,但我們裡頭的罪律,卻成為一種攔阻.
\newline
\newline
但感謝神!第十一節有奇妙的應許,就是那叫耶穌從死裡復活者的靈,會將生命帶給我們.以前我們的身體充滿死亡,但現在那叫耶穌從死裡復活的靈,把生命帶給我們.這不是指將來耶穌再來之時的身體,乃是指現在我們擁有的身體,是現在成就的事.因此,從神的聖靈內,我們在身體和道德上都可以更新.
\newline
\newline
聖靈住在信徒心中,幫助我們在身體上有力量,可以勝過疲倦,並且可以服事神.很多人都有這樣的經歷.我們可以求聖靈加力,使我們在身體上能每日更新,勝過身體的軟弱,更有力的服事祂.雖然我們不明白神為甚麼有時會把疾病加給我們,叫我們在肉身受苦,但我們卻深信,神實在有能力醫治我們身體的疾病.使我們更能過公義的生活.
\newline
\newline
除了身體外,神還有能力幫助我們脫離道德上的死亡,以往我們受罪的束縛,在道德上死亡,但神的靈把生命帶給我們.很多醉酒的人,在信主的一剎那,神的能力臨到他們身上,使他們永遠脫離醉酒的習慣.吸毒的人也一樣,在他們信耶穌的那一刻,神就拯救他們脫離毒品了.這是聖靈所賜的新生命,在我們道德的死亡中,祂賜下新生.雖然奇蹟性的事,並不一定發生在每一個人身上,有人為要擺脫自己的壞習慣,一生掙扎努力.但無論如何,聖靈把生命,和生命的能力,賜給每一個信祂的人,使他能脫離罪和死的律.
\newline
\newline
這位使死人復活的靈,不單賜人有新的生命,更叫人內心盼望屬神的事.聖靈會讓我們傾慕神,因而擺脫罪的糾纏.因此,只要我們是屬神的人,聖靈便樂意在我們必死的身體和道德上,賜下新生命.一位基督徒的作家說,我們永遠不可以說,我要犯罪了,我不能過公義的生活了.因為聖靈將生命放在我們裡面,我們必然可以說,我能夠做到,我能夠得勝.這是第十一節的應許,聖靈把生命放在那必死的身體內.
\newline
\newline
但我們的責任是甚麼呢?第十三節說,我們要每天把身體的罪行治死.這不是說,我們要處罰自己的身體.我們要處罰的,不是自己的身體,乃是身體的惡行,我們不是靠自己的力量去治死,乃要靠聖靈.這情形就好像一些新式汽車的自動轉tie ,當你把tie 轉動小小的時候,裡面的機器便隨你所轉動的加大轉動.神在我們身上的行為,就好像這自動的轉tie ,我們有責任從罪惡轉過來,向著公義.當我們開始轉的時候,神的靈便會幫助我們完全轉向公義.所以,我們的責任是離開罪惡,向著公義.
\newline
\newline
首先,我們必須立定心志,要離開罪惡,轉向公義.神不會強迫我們行公義,我們需要立定心志,依照聖經的教導,立志過公義的生活.立志很重要,就像減肥的人一樣,如果不是下定決心,就不能成功.我們要有堅決的決心,要與罪惡絕緣,過公義的生活.約伯曾這樣說,我與我的眼睛立約,不再望任何引致我犯罪的事.我們也需要這樣,與自己的眼睛立約.今天,在我們四周圍有許多引致我們犯罪的東西,就如書籍,報刊,電視,電影等,我們要遠離這一切,要與自己的眼睛立約,不再接觸這些試探,和一切會引致我們犯罪的東西.我們要在進入試探和罪惡之前止步,不踏進罪的圈套中.聖靈不單能使我們自動轉向,他也給我們有自動的剎車掣.當我們告訴聖靈,我們不再往罪惡的道路走時,聖靈會幫助我們.
\newline
\newline
我們對自己的身體要有自律；保羅說,我是克攻己身,叫身服我.我們不能暴飲暴食,卻求神使我們身體健康.我們也不能通宵工作,卻期望明天精神奕奕.我們不能自欺,我們要有自律,對自己要愛惜,給身體有適當的飲食,運動時.我們對身體要加以屬靈的約束.
\newline
\newline
最後,我們要讓屬神的事物來充滿我們的心.如果我們裡面充滿了屬靈的事,那麼罪惡便沒有發動的餘地.
\newline
\newline
神的目的,要我們每天過著公義的生活,可是身體卻成為我們的煩惱.身體有軟弱,敗壞,損害的時候,也有許多邪情私慾,叫我們犯罪.但聖靈卻進到我們裡面,住在我們身體之內,使我們從死裡復活過來.當我們需要醫治時,神便賜下醫治.以前我們裡面有的,是犯罪的傾向,將新的生命,注入我們裡面.然而我們必須立定心志,遠離罪惡,和一切會引誘我們犯罪的地方,拒絕進入罪惡.神的靈便會把我們引進真正的公義裡.這是一些非常實際的道理,但願我們都能過一個無可指摘,討神喜悅的生活.
\newline
\newline


\newpage

\section{第57屆~港九培靈會~艾榮~第五講　神的兒女}
\label{sec:486}
\textbf{第五講　神的兒女}
\newline
\newline
連結: \href{https://www.hkbibleconference.org/session-message/view/486}{\texttt{https://www.hkbibleconference.org/session-message/view/486}}
\newline
\newline
\hyperref[sec:485]{< < < PREV SERMON < < <}
~
\hyperref[sec:toc]{[返目錄]}
~
\hyperref[sec:487]{> > > NEXT SERMON > > >}
\newline
\newline
羅馬書 8:14-8:18
\newline
\begin{longtable}{cl}
\hline
\hline
章節 & 經文 (和合本修訂版)\\
\hline
8:14 & \begin{tabularx}{0.7\textwidth}{X} 因為凡被神的靈引導的都是神的兒子。 \end{tabularx} \\ \\ \relax
8:15 & \begin{tabularx}{0.7\textwidth}{X} 你們所領受的不是奴僕的靈，仍舊害怕；所領受的是兒子名分的靈，因此我們呼叫：「阿爸，父！」 \end{tabularx} \\ \\ \relax
8:16 & \begin{tabularx}{0.7\textwidth}{X} 聖靈自己與我們的靈一同見證我們是神的兒女。 \end{tabularx} \\ \\ \relax
8:17 & \begin{tabularx}{0.7\textwidth}{X} 若是兒女，就是後嗣，是神的後嗣，和基督同作後嗣。如果我們和他一同受苦，是要我們和他一同得榮耀。 \end{tabularx} \\ \\ \relax
8:18 & \begin{tabularx}{0.7\textwidth}{X} 我認為，現在的苦楚，若比起將來要顯示給我們的榮耀，是不足介意的。 \end{tabularx} \\ \\
[1ex]
\hline
\hline
\end{longtable}
(羅八14-18)「因為凡被神的靈引導的,都是神的兒子.你們所受的,不是奴僕的心,仍舊害怕,所受的乃是兒子的心,因此我們呼叫阿爸!父!聖靈與我們的心同證我們是神的兒女.既是兒女,便是後嗣,如果我們和他一同受苦,也必和他一同得榮耀.我想現在的苦楚,若比起將來要顯於我們的榮耀,就不足介意了.」
\newline
\newline
這幾節經文提醒一些我們必須確信的事,因為這段經文是對重生得救的基督徒說的.第一,宣告神兒女的地位.第二,神應許我們得以成為神兒女的確據.第三,描寫出神兒女的地位.
\newline
\newline
首先是宣告神兒女的地位.十六節宣告我們是神的兒女,只要我們相信耶穌基督,為祂而活,我們就是神的兒女.然而我們必須把握住作神兒女這句說話的意思,因為這是一句十分重要的話.第一節說我們在基督裡不被定罪,到這裡說我們是神的兒女,這是更重要的,也是更美好的關係.第十五節是一個強烈的對比,講到我們信主之後,所受的就不是奴僕的心,仍舊害怕的來服事,乃是兒女的心.神是宇宙的創造者,祂統管一切,命令一切,如果按著我們敗壞了的性情,我們在神面前是可咒詛的,我們這些背叛神的人,對於世界許多的動亂是應該負責的,所以我們在神面前只是叛徒,是被咒詛的.
\newline
\newline
所以,當我們相信耶穌,接受了祂為救主之後,神大可以把我們當作奴僕,一個稱義的奴僕.神發命令,我們便要聽從,如果我們做錯,便要接受神忿怒的審判.救恩大可以帶來一個結果,就是人成為神的奴隸.其實,神如果接納我們做奴隸,神一樣可以維持祂的公義,因為這也是公平的事.想起我們的軟弱,渺少,想起神的偉大,想起我們的罪惡和神的聖潔,我們的無能和神的全能；所以如果我們能作神的奴隸,已經是夠好的了.
\newline
\newline
可是,神告訴我們,祂接納我們不是像一個奴隸,在祂面前存害怕的心；神乃是接納我們作祂的兒女,可以歡喜的稱祂為阿爸!父!「阿爸」是亞蘭文,是當日兒女對父親一種親密的稱呼.當日的官方語言是希臘文,但兒女回到家中,會以親切的語氣,用亞蘭話叫父親做「阿爸」.就像今天我們一般所講的是父親；可是回到家中,看見父親的時候,我們會親切的叫聲「阿爹」.我們所得的,不是奴僕的地位,戰兢害怕,對主人敬畏服從.如果我們做奴僕,已經是神的恩典,但是神竟接納我們為兒女,我們可以呼叫神做阿爸!父!這是何等有福的事,我們可以帶著無懼的心到祂面前,可以面對面的與神說話,把心事告訴祂,求祂為我們分擔.第一節說我們不被定罪,已經是有福的事,但十六節進一步的說,我們可以在神面前稱為兒女.
\newline
\newline
除了宣告我們是神的兒女之外,十六節已講到我們得以成為神的兒女的確據.這裡說,聖靈與我們的心同證我們是神的兒女.有的時候,我們並不常常感到自己是神的兒女,但事實如此.我們盡力為主而活,但有時卻不自覺的懷疑自己是否一個基督徒,這樣的懷疑是可能的.
\newline
\newline
我生長在一個基督化的家庭中,曾經有幾次佈道會,我都舉手決志,並到台前跪下認罪.卻不曉得那一次的接受,是我真正信耶穌的開始.但有人卻會清楚的記得接受主的那一刻.所以,對於那些沒有明確日期信主的人來說,他的疑惑可能是更大的.
\newline
\newline
當我們查考聖經,看見神的偉大,聖潔,又看見自己滿身罪污的時候,我們不禁會問:我怎能成為基督徒?我怎能作神的兒女呢?有人認為基督徒是不會懷疑自己與神的關係,其實不是.有時我們的確是會懷疑的.(約壹三20)是寫給基督徒的,經文說:「我們的心若責備我們,神比我們的心大,一切事沒有不知道的.」有的時候,我們的心會仍舊覺得自己被定罪,在這樣的情況下,我們就要堅信神的話,而不是信自己裡面的感覺.但疑惑是可能有的,而且它會奪去我們心裡的喜悅,平安,也阻礙我們事奉神.
\newline
\newline
另一樣更緊要的影響就是如果我們沒有得救的確據,我們靈命的長進便會受到阻滯,我們就無法過公義的生活,因為公義生活的基礎是在於我們是一個重生得救的人.因此,確信我們在神面前所站的地位是非常重要的.(羅八16)說,聖靈盡祂一切所能來證明我們是神的兒女,祂希望我們能確知自己的地位,從今不再懷疑混亂,以致阻礙自己靈命的長進.
\newline
\newline
(羅八16)沒有告訴我們聖靈怎樣使我們確知自己是神的兒女,因為聖靈所用的方法,不是千篇一律,一成不變的.有時祂會用內心的聲音,有時祂會藉著外界的事物.
\newline
\newline
首先,聖靈會用聖經向我們鮮明我們與神的關係.(林前二12)告訴我們,聖靈的賜下,叫我們能知道神施恩給我們的事.而這經文的上下文,都是講到聖經對真理的重要性的.那就是說,聖靈會藉著聖經向我們解明我們在神面前所站的地位,是神兒女的地位.聖經的真理能叫我們放心,不再疑惑,因為神的話是我們的保證.
\newline
\newline
另外,聖靈也會用過往的經歷來提醒我們.保羅講到他的得救的時候,他也常常提起過往的經歷,他是怎樣蒙神的憐憫,怎樣出死入生.我們也應當在回憶中想起神的信實,神的拯救,好叫我們今天能對神有更堅定的信心.
\newline
\newline
還有,聖靈所用的第三個方法,就是藉著我們的品行和生命的表現來幫助我們.雖然我們不應該滿足於現況,但是我們的心畢竟是渴望一些屬神的事,我們的內心,有神所賜的愛,憐憫和恩慈.這些美德,不是我們自己有的,乃是神所賜的,因為即使我們有高的標準,我們也沒有力量活出來,但神加力量給我們,使我們能有今天的表現.我們可以在自己身上看見神的改變,神的作為.
\newline
\newline
我們要讓聖靈在我們內心工作,不要消滅祂的感動.如果我們時常傾聽屬世的聲音,罪惡的引誘,又讓世俗的事侵入我們的內心,我們就很難聽見聖靈微小的聲音,因而也失去自己成為神兒女的把握了.
\newline
\newline
最後我們看看作為神兒女的結果,十七和十八節.這裡說:既是兒女,便是後嗣.這些話本來是寫給我們在羅馬帝國生活的人的,根據羅馬的律法,凡作兒女的,便自然成為繼承人,是後嗣.我們一成為神的兒女,便是後嗣,可以承繼神一切的產業.這樣的承受,不單是指今生,十八節說這是將要顯於我們的榮耀.保羅在這裡提及將來,是因為他想到基督徒在現世所受的苦難.
\newline
\newline
基督徒在世上是有苦難的,撒旦會多方的攻擊我們,陷害我們.很多人會問,為甚麼我仍要作基督徒呢?如果我不是基督徒,可能我不會受這樣多的苦.但保羅在這裡告訴我們,作為神的兒女,不單是今生的,更是將來的,將來我們要完全承受神的一切.如果我們要把將來要得的榮耀,和今天所受的苦楚比較,今天的苦難真是微不足道的.
\newline
\newline
十七節說我們是與基督同作後嗣,那就是說,基督承受多少,我們就承受多少.雖然耶穌在世的時候,飽受苦難,受人嘲笑戲弄,然而祂卻得到榮耀的身體,可以進入天堂,在父神的右邊坐下,得著權柄榮耀,可以統管萬有.而我們和基督同作後嗣,意思就是我們要承受神所承受的.我們將來要從死裡復活,得著榮耀的身體,進入天堂,在父神面前得著我們的座位,並與基督一同作王,一同審判天使,並統治萬有.事實上,我們很難用人間的詞句來描述得盡將來要承受的,因為實在是太榮耀了.
\newline
\newline
今天,凡信的人都成為神的兒女,我們有聖靈在內心,與我們的心同證我們是神的兒女.當我們面對種種困難,苦惱時,讓我們記起,我們是神的兒女,將來我們要得著何等的榮耀.弟兄姊妹,讓我們抓緊神的應許,清楚自己在神面前所站的地位,和我們將要承受的產業.
\newline
\newline


\newpage

\section{第57屆~港九培靈會~艾榮~第六講　誠懇的期望}
\label{sec:487}
\textbf{第六講　誠懇的期望}
\newline
\newline
連結: \href{https://www.hkbibleconference.org/session-message/view/487}{\texttt{https://www.hkbibleconference.org/session-message/view/487}}
\newline
\newline
\hyperref[sec:486]{< < < PREV SERMON < < <}
~
\hyperref[sec:toc]{[返目錄]}
~
\hyperref[sec:488]{> > > NEXT SERMON > > >}
\newline
\newline
羅馬書 8:19-8:25
\newline
\begin{longtable}{cl}
\hline
\hline
章節 & 經文 (和合本修訂版)\\
\hline
8:19 & \begin{tabularx}{0.7\textwidth}{X} 受造之物切望等候神的眾子顯出來。 \end{tabularx} \\ \\ \relax
8:20 & \begin{tabularx}{0.7\textwidth}{X} 因為受造之物屈服在虛空之下，不是自己願意，而是因那使它屈服的叫他如此。 \end{tabularx} \\ \\ \relax
8:21 & \begin{tabularx}{0.7\textwidth}{X} 但受造之物仍然指望從敗壞的轄制下得釋放，得享神兒女榮耀的自由。 \end{tabularx} \\ \\ \relax
8:22 & \begin{tabularx}{0.7\textwidth}{X} 我們知道，一切受造之物一同呻吟，一同忍受陣痛，直到如今。 \end{tabularx} \\ \\ \relax
8:23 & \begin{tabularx}{0.7\textwidth}{X} 不但如此，就是我們這有聖靈作初熟果子的，也是自己內心呻吟，等候得著兒子的名分，就是我們的身體得救贖。 \end{tabularx} \\ \\ \relax
8:24 & \begin{tabularx}{0.7\textwidth}{X} 我們得救是在於盼望；可是看得見的盼望就不是盼望。誰還去盼望他所看得見的呢？ \end{tabularx} \\ \\ \relax
8:25 & \begin{tabularx}{0.7\textwidth}{X} 但我們若盼望那看不見的，我們就耐心等候。 \end{tabularx} \\ \\
[1ex]
\hline
\hline
\end{longtable}
(羅八19-25)「受造之物,切望等候神的眾子顯出來.因為受造之物服在虛空之下,不是自己願意,乃是因那叫他如此的.但受造之物仍然指望脫離敗壞的轄制,得享神兒女自由的榮耀.我們知道一切受造之物,一同歎息勞苦,直到如今.不但如此,就是我們這有聖靈初結果子的,也是自己心裡歎息,等候得著兒子的名分,乃是我們的身體得贖.我們得救是在乎盼望,只是所見的盼望不是盼望,誰還盼望他所見的呢?但我們若盼望那所不見的,就必忍耐等候.」
\newline
\newline
我們一開始給羅馬書八章一個主題.貫徹始終的救恩.在這章聖經裡面,保羅先告訴我們怎樣可以在基督裡得釋放,脫離罪惡.跟著保羅把神的心意擺在我們面前,這心意就是要我們每天過公義的生活.然後,昨天我們又看見我們在神面前的地位,和所得的福氣,這是救恩的結果,不單是現在,更是伸展到將來.想到神將來要為我們成就的事,我們就屏息以待,歡欣不已.
\newline
\newline
今天思想的這段經文,是講到將來要發生的兩件事,一是關乎受造之物,一是關乎我們的身體.十九到廿一節講到一切受造之物將來要發生的,特別是指我們四周的世界,動物,山川,湖泊,星辰,日月等要發生的事；而二十三節則講到我們的身體將來要發生的.
\newline
\newline
我們先看十九到廿一節,這是講到我們四周的世界將要發生的事.這裡說,受造之物是服在虛空之下,因此它是在敗壞的的轄制之中歎息勞苦.歎息在原文的意思,是特別指產婦在生產之時所發出的呻吟疼痛.另外經文又說受造之物切望等候神眾子的顯現,當它這樣切望的時候,它就能脫離敗壞的轄制.
\newline
\newline
首先,受造之物,這整個世界,都服在虛空之下.當神創造人的時候,神把一切的權柄都賜給人.(創一28)「神就賜福給他們,又對他們說,要生養眾多,遍滿地面,治理這地,也要管理海裡的魚,空中的鳥,和地上各樣行動的活動.」亞當有管理全地的權柄,甚至能控制整個宇宙.但人因為犯罪的緣故,便失掉了智慧,和管理全地的權柄.神要限制人,免得人濫用權柄,影響宇宙.
\newline
\newline
到新約的福音書,有第二個亞當,就是耶穌基督,他大有掌權和命令宇宙的權柄,他曾經發出命令去平息風浪.我相信頭一個亞當在犯罪之前也擁有這樣的權柄,只可惜他因為犯罪而失去了.結果所有受造之物都落在漫無目的和虛空之中,保羅說現在受造之物,都落在敗壞的轄制之中.我們四周的受造之物,現今不過是敗壞,朽壞和染污了的,人已經沒有控制受造之物的能力了.二十節所講受造之物服在虛空之下的「虛空」,原來的意思就是迷惘,漫無目的.
\newline
\newline
本來,神所造的各樣物件,正等待有好的用途,一同將頌讚歸給神.但人因為犯罪,便失去使用宇宙來讚美神的能力.因此,受造之物便迷惘,空虛,不能發揮它的作用.雖然隨著科技的發展,人已經在控制宇宙上略有進展,譬如人可以登陸月球,可以控制天氣,但是比起神預先給人的能力,卻是遜色得多了.神不能給予人太多能力,因為人在道德上缺乏智慧和能力去使用宇宙的能源.
\newline
\newline
比方說,人類自懂得把原子分裂,發出原子能之後,人便幾乎立即製做原子彈來自相毀滅,人用征服自然界的能力來彼此殘殺.本來神賜人有統管宇宙的權柄,但人因為犯罪,失去了這樣權柄,以致宇宙變得漫無目的,許多受造的東西變得荒置無用.所以二十二節說受造之物歎息勞苦,自然界的災害,人為的災害,一切都使受造之物歎息不已.自然界失去控制,在歎息,因為人類失去了控制和治理他的能力.
\newline
\newline
第十九節說,受造之物,切望等候神的眾子顯出來.這是受造之物在歎息勞苦之餘的一個盼望.將來,一切屬神的兒女,一切信靠耶穌的人,要在宇宙間顯露出來.帖前四章說神的天使長要出來,神的號筒要吹響,在基督裡睡了的人要先復活,那存留在世的信徒,也要被提到空中與主相遇.在那裡,有一個盛大的婚宴,我們一切屬神的人要顯露出來,作為神的眾子.到那時,我們要完全的得贖,而且恢復資格,可以重新掌管受造之物.人本來以前所得的權柄,到那時要完全恢復.
\newline
\newline
在神的創造之物中,將來要有一次洗淨的過程.彼後三章說將有一天,天地都要變化,神要降下火,一切有形,有質的,都要被消化.到審判的日子來到,神要除去世上一切的罪,和一切與罪有關的東西.這審判的日子和羅馬書八章所說的並無出入,乃是互相吻合的.宇宙不會完全消失,乃是被神洗淨,洗淨之後仍要留存.我們作為神兒女的,到那時要領受權柄,一切的毀壞要成為過去,一切風浪要平息,一切野獸都變為馴良,沙漠要成為綠洲,天體的星球都成為一幅美麗的圖案獻給神.行星環繞時所發的聲音,要成為稱頌神的音樂.在那時,受造的一切不再歎息,乃成為讚美神的美麗詩歌.
\newline
\newline
現在所有受造之物,都屏息等候將來的榮耀,他們數算日子,要等待神盛大的婚宴,他們渴望回復神創造時的十全十美.不過,這只是受造之物的期望,還有一個更重要的,就是我們身體得贖的真理.
\newline
\newline
二十三節說我們也歎息,我們的身體仍在轄制之下,還未得贖.我們的身體也像宇宙一樣,有敗壞,有衰殘.人體有九個主要系統,人的骨骼有二百零六條,人的眼睛有五百多塊肌肉,所有這些都是會衰殘,會朽壞的.我們的眼睛失明,耳會失聰,手腳會殘廢,連口和牙齒都會受感染損壞.至於皮膚,也很容易燙傷,燒傷,或患上皮膚病.人的內臟,心,肝,脾,肺,腎等,都會衰退,甚至患病.人體內有六夸脫的血,血的濃度,血的含糖量,過高或過低都形成疾病.人可能患上的病,真是數不勝數,頭痛,發燒,嘔,肚瀉,喉痛,肺炎,感冒,關節炎等等,總之是品種繁多,不能勝數.我們的身體,實在是在許多的毀壞之下.
\newline
\newline
當人犯罪之後,人不但喪失控制大自然的權力,連人自己的身體,也招引種種敗壞的情況.我們的身體,正步向衰老和死亡.但是聖靈卻把生命賜給我們,使我們必死的身體活過來,這是十一節早就說過的.聖靈可以使我們的身體健壯,也可以醫治我們的.到二十三節說我們是聖靈初結的果子,現今聖靈所造的,只是初步的拯救,到最後我們的整個身體都要得贖.聖靈今只是讓我們初嘗將來身體得贖的喜悅.
\newline
\newline
不過,即使聖靈今天幫助我們,但我們的身體仍會衰老,死亡,所以我們仍會歎息.但二十三節給我們一個寶貴的應許,就是有一天我們的身體要得贖.(林前十五53)說我這必壞的身體,將來要變成不朽壞的,衰殘的將要變成不衰殘的.
\newline
\newline
將來,大自然和我們的身體,都是從毀壞的勢力下得贖,得蒙神的拯救,我們要得回一個真正的身體,是不會衰殘,朽壞的.而且在這新的身體內,我們仍會有個別獨特的性格.所以當耶穌再來之時,我們的身體要改變,要離開朽壞,不再衰殘.我們不再受疾病,痛楚的威脅,我們會有主一樣榮耀的身體,可以超脫地心吸力的原則,我們所唱的,好像天使的歌聲,我們的心思能完全明白真理,我們會得著神的真光作為我們的衣服.今生在病床中的弟兄姊妹,將來將不會再在病床中.今生的疲勞困倦,將來都不復存在.今生我們是步向墳墓和死亡,但將來是永遠的生命,是十全十美的.
\newline
\newline
因此,今天這幾節經文,讓我們看見受造之物和我們的身體,都有一個榮耀的盼望,縱然我們今天已在敗壞之中,但將來我們要完全的得贖.二十四和二十五節就是講到那盼望,那是在救恩之中很重要的部份.基督徒是滿有盼望的,擺在我們面前的還有許多福氣,雖然我們今天還未得著,但到將來我們得到時,那喜樂是更大的.
\newline
\newline


\newpage

\section{第57屆~港九培靈會~艾榮~第七講　聖靈的時代}
\label{sec:488}
\textbf{第七講　聖靈的時代}
\newline
\newline
連結: \href{https://www.hkbibleconference.org/session-message/view/488}{\texttt{https://www.hkbibleconference.org/session-message/view/488}}
\newline
\newline
\hyperref[sec:487]{< < < PREV SERMON < < <}
~
\hyperref[sec:toc]{[返目錄]}
~
\hyperref[sec:489]{> > > NEXT SERMON > > >}
\newline
\newline
羅馬書 8:26-8:27
\newline
\begin{longtable}{cl}
\hline
\hline
章節 & 經文 (和合本修訂版)\\
\hline
8:26 & \begin{tabularx}{0.7\textwidth}{X} 同樣，我們的軟弱有聖靈幫助。我們本不知道當怎樣禱告，但是聖靈親自用無可言喻的嘆息替我們祈求。 \end{tabularx} \\ \\ \relax
8:27 & \begin{tabularx}{0.7\textwidth}{X} 那鑒察人心的知道聖靈所體貼的，因為聖靈照著神的旨意替聖徒祈求。 \end{tabularx} \\ \\
[1ex]
\hline
\hline
\end{longtable}
(羅八26-27)「況且我們的軟弱有聖靈的幫助,我們本不曉得當怎樣禱告,只是聖靈親自用說不出來的歎息,替我們禱告.鑑察人心的,曉得勝靈的意思,因為聖靈照著神的旨意,替聖徒祈求.」
\newline
\newline
這兩節經文講到聖靈.聖靈是三位一體神的第三位,祂與聖父,聖子一同有尊貴,榮耀,權柄,能力.在聖靈身上,結集了一切神性,尊貴和崇高.然而人卻是渺小,有限的,二十六節便講到人的軟弱.「軟弱」在原文是單數,有解經家說,人只是一大堆軟弱而已.但聖靈卻幫助我們,「幫助」原來的意思就是有兩個人,一個挑著極重的擔子；但另一個就用肩膀去頂著第一個人的重擔,這叫做幫助.
\newline
\newline
在我們軟弱的時候,聖靈便來到我們旁邊,幫助我們肩負重擔,解決問題.羅馬書八章共有三次提到歎息,二十二節說受造之物一同歎息,二十三節說我們心裡歎息,現在二十六節說聖靈在歎息.受造之物為著自己而歎息,等待救贖.我們歎息,因為我們的身體太有限,我們在等待將來身體得贖.受造之物和我們都是自己歎息,但聖靈的歎息卻是為我們,祂把我們的重擔拿去,放在祂的身上.
\newline
\newline
我們這次的總題是「貫徹始終的救恩」.救恩的起點,是耶穌基督為我們成就了救贖,叫我們與神和好,因此是耶穌的工作,使我們在救恩上有了起點.然後在救恩的旅途中,聖靈一直與我們同行,向著救恩的結果來前進.
\newline
\newline
創世記二十四章記載一件很有趣的歷史事實,亞伯拉罕想為他們的兒子以撒娶妻,於是差遣他的僕人以利以謝回他的本地去,為他物色一位媳婦,最後找到了利百加.這段事實,反映出一個美麗的真理,就是耶穌基督在吸引我們.我們已經預備好要作耶穌的新婦,我們向著那羔羊的婚宴進發.以利以謝陪著利百加,從她的本國一直到亞伯拉罕的地方,一路保護她,服事她.同樣,聖靈要與我們一路同行,直向著那終局,就是羔羊的婚宴前進.在我們走在人生的路上之時,聖靈的服務是周詳的,隨時隨地的.
\newline
\newline
首先,聖靈來住在信徒心中,感動和帶領我們走在公義的生活中.聖靈給我們能力勝過肉體的傾向,使我們能逐日過著公義的生活.另外,聖靈又在我們裡面,使我們必死的身體又活過來.他注入屬靈的能力,又賜下身體的醫治.還有,聖靈又給我們確據,使我們知道自己是神的兒女,可以放心在屬靈裡長進.今天,我們看見聖靈在禱告上所作的工作.
\newline
\newline
(羅八26)講到聖靈在禱告上所給我們的幫助,是何等重要.聖經說我們本不曉得怎樣按我們的本份去禱告,我們充滿無知,不曉得如何代禱,但是聖靈卻為我們代求.意思就是說,聖靈在我們心裡作工,提醒我們該為甚麼事代禱,並且教導我們怎樣禱告,這主要是側重一些代禱的事情.聖靈主要把一種渴慕禱告的心賜給我們,祂藉著聖靈來鼓勵我們要禱告.所以,如果我們更多願意禱告,我們更要被聖靈充滿.
\newline
\newline
作神工作的人有一種危機,就是過於忙碌,以致沒有時間禱告.但是如果我們的生命,是真的被神引導的話,我們內心必然有一種催逼,要更多的禱告.聖靈不單叫我們有禱告的負擔和催逼,祂更會把需要代禱的事項告訴我們.祂可能藉著一些特別的感動來告訴我們,要為甚麼人,甚麼事禱告,但也可能藉著我們讀經的時候指示我們,或者藉別人分享他們的見證或難處的時候感動我們.因此,聖靈不但感動,祂也教導,指示我們該怎樣禱告.
\newline
\newline
另外,二十六節告訴我們,聖靈感動的結果,可能不是一些我們聽得見的聲音,可能是一些說不出來的歎息.有人解釋說,這種禱告的負擔,能使屬神的人俯伏下來,謙卑禱告,這些他們的禱告就像說不出來的歎息一樣.另一個作者說,有時當我們俯伏禱告的時候,我們的需要和重擔十分沉重,以致我們講不出說話,只能歎息.
\newline
\newline
所以我們的禱告,可以有不同層次的負擔.有時顯得簡單,我們只為自己的需要祈禱.但有時是等候在神面前,忍耐,安靜的等候.有時需要恆切的禱告,把我們的需要一再呈在神面前.然而有時我們因為心靈過於沉重,以致發不出話語,只在心靈裡默默的歎息,向神呼求.聖靈會在這種沉切的呼求中帶領我們,幫助我們.
\newline
\newline
二十七節便講到這種禱告的價值.當我們的禱告有聖靈幫助時,這禱告一定蒙神應允.因為聖靈了解到神的心思意念,並確實的了解神的旨意,所以當聖靈感動我們禱告時,這禱告一定合符神的旨意.
\newline
\newline
透過聖靈感動和幫助的禱告十分重要,能使事情改變,所禱告的能合神的心意.其實這些禱告先由神的心發出,又由聖靈將神的心意傳給我們,再由我們告訴給神,神的旨意便得以成就.今天的兩節經文正告訴我們,神所作的工作就是這樣.當我們軟弱不曉得怎樣禱告時,聖靈便來幫助我們,教導我們照神的旨意禱告.有時藉著說不出來的歎息禱告,但神卻明白我們的心意,因為是聖靈按著神的旨意加給我們的.
\newline
\newline
因此,我鼓勵你成為一個禱告人,特別是年長的兄弟姊妹,更可以從事代禱的事奉,這是非常重要的事奉,甚至比講道或別的事奉更重要,他們的賞賜是大的.我們可以從事這樣的事奉,讓聖靈指引我們,把禱告的負擔和需要放在我們心裡,有聖靈的幫助,我們可以從事這有意義的事奉.
\newline
\newline
現在我們看這兩節的經文的另一點真理,雖然我從未在參考書上看過這樣的解釋.我發現聖靈不但感動我們,叫我們代求,聖靈本身也替我們禱告.神的聖靈用說不出來的歎息的原因,是因為聖靈所用的言語,根本不是我們人間言語,所以是說不出來的.
\newline
\newline
聖經說耶穌現今是在父神的右邊,為我們代求.神的兒子長遠活著,替信徒祈求.但這裡暗示聖靈在世上也為我們禱告.二十六節說我們本不曉得該怎樣禱告,我們連怎樣為自己禱告也不知道,我們就好像一個被告,站在法官面前,不曉得怎樣為自己辯護,又好像一個小孩,不知道該吃甚麼食物.有時我們病了,不能禱告.或是太忙,沒有空禱告.但聖靈就來,把我們的重擔拿去,放在他自己的身上,然後為我們禱告,是用說不出來的歎息替我們禱告.正因為是聖靈所發的,所以神一定垂聽.
\newline
\newline
所以,聖靈除了一方面賜我們力量去為別人代求之外,聖靈還親自為我們祈求.聖靈站在我們身邊,看見我們的需要,祂要為我們禱告.特別當我們失敗,遊蕩,不討神喜歡的時候,聖靈更迫切的為我們禱告.聖靈的關心,祂的恩慈,就好像母親對兒子一樣.
\newline
\newline
雖然許多時候,我們並不自覺聖靈的工作,但聖靈就是默默的,忠心的為我們祈求.祂是不斷的,恆常的,每日的為我們禱告,幫助我們.今天,我們的生命都被罪惡所侵害過,受傷過,以往的敗壞,有很多的壞影響,依然在我們的生命中.我們充滿了軟弱,但聖靈卻願意獻身幫助我們,服事我們.這裡說神的聖靈要在我們身邊,協助我們渡過一切軟弱,一切的需要.特別在禱告的事奉上,聖靈要幫助我們.即使我們自己不曉得怎樣禱告時,聖靈也站在我們旁邊,為我們禱告.
\newline
\newline
我相信這兩節經文所給我們的教導正是這樣.既然這是真理,我們便要抓緊它,並且為此而感謝主.讓我們從聖靈的幫助中得著支持,得著力量,不單自己享受聖靈的代求,同時也接受聖靈的感動和催逼,去為別人代求.
\newline
\newline


\newpage

\section{第57屆~港九培靈會~艾榮~第八講　神為人的計劃}
\label{sec:489}
\textbf{第八講　神為人的計劃}
\newline
\newline
連結: \href{https://www.hkbibleconference.org/session-message/view/489}{\texttt{https://www.hkbibleconference.org/session-message/view/489}}
\newline
\newline
\hyperref[sec:488]{< < < PREV SERMON < < <}
~
\hyperref[sec:toc]{[返目錄]}
~
\hyperref[sec:490]{> > > NEXT SERMON > > >}
\newline
\newline
羅馬書 8:28-8:30
\newline
\begin{longtable}{cl}
\hline
\hline
章節 & 經文 (和合本修訂版)\\
\hline
8:28 & \begin{tabularx}{0.7\textwidth}{X} 我們知道，萬事都互相效力，叫愛神的人得益處，就是按他旨意被召的人。 \end{tabularx} \\ \\ \relax
8:29 & \begin{tabularx}{0.7\textwidth}{X} 因為他所預知的人，他也預定他們效法他兒子的榜樣，使他兒子在許多弟兄中作長子。 \end{tabularx} \\ \\ \relax
8:30 & \begin{tabularx}{0.7\textwidth}{X} 他所預定的人，他又召他們來；所召來的人，他又稱他們為義；所稱為義的人，他又叫他們得榮耀。 \end{tabularx} \\ \\
[1ex]
\hline
\hline
\end{longtable}
(羅八28-30),「我們曉得萬事都互相效力,叫愛神的人得益處,就是按他旨意被召的人,因為他預先所知道的人,就預先定下效法他兒子的模樣,使他兒子在許多弟兄中作長子.預先所定下的人,又召他們來,所召來的人,又稱他們為義；所稱為義的人,又叫他們得榮耀.」
\newline
\newline
這幾節經文,是把(羅八28)之前的真理來一次溫習,之後的(31-39)是讚美神的詩歌.今天我們只看(28-31)的溫習和總結.從這幾節經文中,我們來思想四個要點.我們要看二十八節所提到的人物,和所記載的應許.二十九節是神所賜的目的,到六十節是神的工作和計劃.
\newline
\newline
首先我們看廿八節所提及的各個人物.這節經文是講及一些愛神的人,就是那些選擇認識神,愛神,為神而活,心裡委身給神,討神喜悅的人.可惜今天我們對愛的觀念,已經受電視影響.他們以為情慾的衝動就是愛,情感上的一些感受,男女的喜悅就是愛.希臘文也有對這種愛的描述,但它所用的字跟愛的字是不同的.神的愛的意思,就是立定心志,為了對方的好處而犧牲自己.
\newline
\newline
有人對神的愛有很強烈的反應,他們在感情上表現得十分激烈,十分衝動,這是無可厚非的.可是,聖經所要求的並不是這樣,聖經要求我們立定心志,去愛神,去委身給神,為神而活,這才是最重要的.有的基督徒因為自覺對神缺少在感情上強烈的回應,因而自責,這是不必的.因為重要的是我們在心內決志,要決定愛神,服事神,討神的喜悅,為神而活.這節經文所講愛神的人,就是指這樣的人,而不是在感情有強烈反應的人.
\newline
\newline
對於決定愛神,服事神,為神而活的人,二十八節帶來一個應許.很多基督徒都喜愛這節金句,但有的卻誤解了其中的意思.到底這節經文真正的經文是甚麼呢?一般人對這節經文的了解,就為只要人愛神,那麼周圍所發生的事情,無論是甚麼,終歸叫愛神的人得著益處的.即使是疾病,手術,失業,失戀等,都會帶來好處.這是我們對二十八節一般的了解.
\newline
\newline
可是這節經文並不是講我們四周所發生的事,這裡所講的「萬事」,是指神所做的各樣事情,這和廿九,三十節所提到神所成就的事是有關係的.廿九節和三十節講到神的預定,選召,稱義,和叫人得榮耀.所以廿八節所說的萬事,其實是指廿九和三十節所列出的事情,就是神藉著耶穌為我們所成就各樣的事.
\newline
\newline
還有,二十八節所講的「得益處」,是指我們的拯救,救恩.所以,廿八節的解釋就是,神藉耶穌所為我們成全的各樣事情,會帶來救恩,使我們從救恩中得著好處.
\newline
\newline
比方說,當你投考大學的時候,校方會給你一本大學概覽,裡面列出教授的名單,所開科目,校內一切的設備,如圖書館,實驗室等概覽內說,如果你進了這所大學,那麼校內一切的設備,會一同運作,叫你得著最好的教育.同樣,廿九和卅十節列舉了神藉耶穌所為我們成全的萬事.而廿八節說,只要你愛神,廿九和三十節所講各樣的事情,神都必為你成就,好叫你得著祂的救恩.當然,神會使用我們周圍所發生一些不如意的事,叫我們得著好處,但這裡所說的萬事並不是指這些,只是指神的工作和計劃而已.
\newline
\newline
那麼,神為那些愛祂的人定了甚麼工作和計劃呢?三十節說神預先知道的人,神便預先定下,選召他們成為神的兒女.其實,神的心意是希望人人都得著拯救.但在世界未造之先,神就早已知道那些人會揀選服事祂.既然神預先知道,神便預早揀選他們成為得救的人.神預先決定那些選擇愛祂的人作為祂的兒女.神早已知道我們的一切,我們的出生,我們的成長,我們所遇到的一切事,以致我們的性格,喜好,背景,選擇等.神知道我們會聽到福音,又願意接受福音,神更知道誰會選擇愛祂,跟隨祂.對於這樣的人,神便預先選召他們作祂的兒女.
\newline
\newline
當人一陷在罪中時,神便立即開始祂拯救的計劃.祂知道有些人會選擇服事,所以神立即設立一個計劃,要選召他們,拯救他們.神選召以色列人,藉著他們把啟示傳給我們.祂又差遣祂的兒子到世上來,藉著祂兒子所成就的救贖把我們拯救出來.跟著神又賜下聖靈,藉聖靈把真理帶給教會,又在人中間指引人,叫人知罪悔改,叫我們得救.基督徒愈認識真理,便愈知道自己在救恩的事上,一無所有.一切都是耶穌所作的,我們毫無作為.
\newline
\newline
救恩完全是神作成的,我們只是聽見,相信,醒悟過來,接受而已.所以聖經說,是神召我們.神沒有讓我們主動的去尋找祂,相反的,神主動來找我們,召我們,忍耐的等我們接受祂.然後向所召來的人,神又稱他們為義.只要人對神有信心的回應,神便藉耶穌將那些悔改的人,都拯救出來,洗淨他們一切的罪,使我們在祂面前如同清白,未曾犯過罪的一樣.這是何等奇異的恩典.
\newline
\newline
我們得稱義的程度,就如同神的兒子耶穌一樣.神把祂兒子的衣裳賜給我們,使我們在祂面前好像祂兒子一樣的清白,潔白.神把祂兒子的公義賜給我們,叫我們穿上,而祂拿去我們污穢的衣服,放在十字架上.神定下愛祂的人要得蒙祂的拯救,祂一直作工,要選召我們歸向祂.那些蒙召的人,神又使他們稱義,變為潔白.而祂稱為義的人,祂又使他們得著榮耀.這不單指我們將來要得的榮耀,就是現在,我們已經可以得著榮耀.在我們的心內,有神所給的榮耀和喜樂,我們有與神相交的榮耀,我們有過聖潔生活的榮耀,這種榮耀還要繼續增加,擴展,到最後我們要得著新的,不朽壞的身體,可以進到天堂去,與主面對面相交,這是何等大的榮耀,而這正是神工作的計劃.
\newline
\newline
神早知道有些人會愛祂,所以祂預先定下他們得救,祂工作,好叫人能得救.而神選召信徒之後,神又稱他們為義,叫他們得榮耀.保羅在第三十節把神所為我們作的逐樣的數算出來,叫我們看見神偉大的計劃.神的命定,神的選召,神的稱義,神所賜的榮耀,這一切都是神為愛祂的人預備的.
\newline
\newline
過去幾天我們已經看過神為我們所作的奇事,今天的經文只是一個溫習和總結.但當我們數算神奇妙的工作和計劃時,我們的心又再興奮不已.
\newline
\newline
最後,我們看看神成就這一切的目的.二十九節說,神預先所定的這一切計劃是有目的的,為要使我們效法祂兒子的模樣,使祂的兒子在許多弟兄中作長子.那就是說,我們不單穿上公義的外袍,我們更有耶穌的形象,個性,道德,有祂的模樣.我們既然得稱為義,就當靠著聖靈行事為人,有神的公義和聖潔在我們的生活中.我們要效法耶穌的態度,祂對天父是順服的,祂又與父有密切的關係,祂常花很多時間與神交往,又常尋求得天父的喜悅.我們在與別人的關係,也應該像耶穌那樣無私,滿有恩慈,服事人,幫助人.還有,我們也應該在自己的生活上像耶穌,有耶穌那樣有聖潔,純潔.這是神的目的.
\newline
\newline
另外,這節經文除了講到我們現今的生活之外,還講到我們將來永遠在神面前的生活.「模樣」原文就是神的形象,(創一86)就用過這個詞,神說要按祂自己的形像樣式造人.那就是說,神要求人活在神的水平上.神造一切動物,魚,鳥,但神沒有要求這些生物活在祂水平上,神只要人活在祂的水平上,與祂有交通,過著聖潔,平安,喜樂的生活.神要人與祂一同統管萬物,分享祂所擁有的一切,這就是神的心意.
\newline
\newline
然而人因為犯罪墜落,便失去神所賜我們的一切.現今神藉著耶穌的救贖,我們可以重新得回所失去的一切；不但與神恢復關係,更得以稱為義,分享神的榮耀.這一切的工作,都為要使我們能夠效法神兒子的模樣,活在神的水平上.而在神的形像中,基督耶穌又成為我們的長兄,在許多的兒子中,耶穌是排行最先的.因此,整個救贖的目的,是神將人從犯罪的悲劇中帶出來,進到神崇高的水平的生活去.
\newline
\newline


\newpage

\section{第57屆~港九培靈會~艾榮~第九講　得勝有餘}
\label{sec:490}
\textbf{第九講　得勝有餘}
\newline
\newline
連結: \href{https://www.hkbibleconference.org/session-message/view/490}{\texttt{https://www.hkbibleconference.org/session-message/view/490}}
\newline
\newline
\hyperref[sec:489]{< < < PREV SERMON < < <}
~
\hyperref[sec:toc]{[返目錄]}
~
\hyperref[sec:472]{> > > NEXT SERMON > > >}
\newline
\newline
羅馬書 8:31-8:39
\newline
\begin{longtable}{cl}
\hline
\hline
章節 & 經文 (和合本修訂版)\\
\hline
8:31 & \begin{tabularx}{0.7\textwidth}{X} 既是這樣，我們對這些事還要怎麼說呢？神若幫助我們，誰能抵擋我們呢？ \end{tabularx} \\ \\ \relax
8:32 & \begin{tabularx}{0.7\textwidth}{X} 神既不顧惜自己的兒子，為我們眾人捨了他，豈不也把萬物和他一同白白地賜給我們嗎？ \end{tabularx} \\ \\ \relax
8:33 & \begin{tabularx}{0.7\textwidth}{X} 誰能控告神所揀選的人呢？有神稱他們為義了。 \end{tabularx} \\ \\ \relax
8:34 & \begin{tabularx}{0.7\textwidth}{X} 誰能定他們的罪呢？有基督耶穌 已經死了，而且復活了，現今在神的右邊，也替我們祈求。 \end{tabularx} \\ \\ \relax
8:35 & \begin{tabularx}{0.7\textwidth}{X} 誰能使我們與基督的愛隔絕呢？難道是患難嗎？是困苦嗎？是迫害嗎？是飢餓嗎？是赤身露體嗎？是危險嗎？是刀劍嗎？ \end{tabularx} \\ \\ \relax
8:36 & \begin{tabularx}{0.7\textwidth}{X} 如經上所記：「我們為你的緣故終日被殺；人看我們如將宰的羊。」 \end{tabularx} \\ \\ \relax
8:37 & \begin{tabularx}{0.7\textwidth}{X} 然而，靠著愛我們的主，在這一切的事上，我們已經得勝有餘了。 \end{tabularx} \\ \\ \relax
8:38 & \begin{tabularx}{0.7\textwidth}{X} 因為我深信，無論是死，是活，是天使，是掌權的，是有權能的 ，是現在的事，是將來的事， \end{tabularx} \\ \\ \relax
8:39 & \begin{tabularx}{0.7\textwidth}{X} 是高處的，是深處的，是別的受造之物，都不能使我們與神的愛隔絕，這愛是在我們的主基督耶穌裡的。 \end{tabularx} \\ \\
[1ex]
\hline
\hline
\end{longtable}
(羅八31-39)「既是這樣,還有甚麼說的呢?神若幫助我們,誰能敵擋我們呢?神既不愛惜自己的兒子為我們眾人捨了,豈不也把萬物和他一同白白的賜給我們麼?誰能控告神所揀選的人呢?有神稱他們為義了.誰能定他們的罪呢?有基督耶穌已經死了,而且從死裡復活,現今在神的右邊,也替我們祈求.誰能使我們與基督的愛隔絕呢?難道是患難麼?是困苦麼?是逼迫麼?是飢餓?是赤身露體麼?是危險麼?是刀劍麼?如經上所記,我們為你的緣故,終日被殺,人看我們如將宰的羊.然而靠著愛我們的主,在這一切的事上,已經得勝有餘了.因為我深信無論是死,是生,是天使,是掌權的,是有能力的,是現在的事,是將來的事,是高處的,是低處的,是別的受造之物,都不能叫我們與神的愛隔絕,這愛是在我們的主基督耶穌裡的.」
\newline
\newline
我們已經說過,羅馬書八章二十八到三十節,是總結全章聖經的經節.所以到第三十一節,作者說,既是這樣,還有甚麼說的呢?跟著他就發出一首偉大的讚美神的詩歌.在這幾節結束的經文中,作者向神流露出偉大的讚美.裡面的詩句,十分甜美,把要講的信息表達得十分透徹.
\newline
\newline
這段經文可以說有些軍事的作用,我們唱起來,就好像唱一首戰歌,要向前走一樣.卅一節論到戰爭,卅七節講到征服.今天我們就循這些思路去思想.
\newline
\newline
首先我們看看那些敵對我們的人.這首詩的作者並不是一個天真的人,他明白現實的難處和困境.他所認識的世界,是充滿矛盾,戰爭,敵人,對抗的.卅三節說我們作為神所揀選的,是真正面對一種勢力,這種惡的勢力就是我們的敵人.三十五節便說出我們在世上可能遇到的,有患難,困苦,逼迫,飢餓,赤身露體,危險,刀劍等.卅六節更進一步的引用舊約說「我們為你的緣故,終日被殺,人看我們如將宰的羊.」
\newline
\newline
我們在世上是有真實的敵人的,有很多試探更攻擊我們,撒旦竭盡所能要攻擊我們的信仰.這世界是背道而馳的,那些在罪中打滾的人,要用盡方法把我們拉進罪惡中.我們作基督徒的,就像一個站在彌敦道的人,所面對的人群,都是與我們背道而馳的.
\newline
\newline
有一天,一位航空工程師和朋友告訴我在高空飛行的飛機所要面對的挑戰.他說因為高空的氣溫低,而飛機的金屬在飛行時又十分熱,這種冷熱相遇的情形,使得金屬慢慢的被消耗盡,這是一個很大的問題.
\newline
\newline
同樣,我們在屬靈的生命中,愈向高處行,所受的危險和衝擊便愈大.不但這樣,我們事奉神的時候,就好像戰鬥機一樣,四圍有許多敵機想把我們擊落,撒旦用很多高射砲,從地面攻擊我們.我們四周都有敵人,作者在寫這些經文時,也寫出我們所面臨的真實敵人.
\newline
\newline
然後作者便開始往內看.卅四節講到定罪的可能性.當我們愈覺得自己虧缺神的榮耀時,便愈奇異,為甚麼神要拯救我們.我們愈覺得神的救恩是何等豐盛,希奇為甚麼神用這樣豐盛的恩典對我們.我們在神面前,因著自己的軟弱而覺得尷尬,我們的失敗控告我們,令我們內疚.這種不配和內疚便成了我們的敵人,扼殺我們.作者在向內望的時候,看到了存在於自己裡面的敵人.
\newline
\newline
因此,作者所面對的敵人,是內外夾攻的.在面對前面漫長的歲月時,作者又怎樣呢?既然周圍都充滿仇敵,危機,那麼前面的日子怎麼過呢?卅八節講到死,生,天使,掌權的,有能力的,現在的事,將來的事,高處的,低處的,別的受造之物等等.我們有想到將來的年日嗎?在今生所受的攻擊和扼制之中,我們與神的關係到底會怎樣呢?我們怎能知道在前面漫長無邊的歲月中會遭遇些甚麼事呢?
\newline
\newline
這裡說,誰能敵擋我們呢?是屬靈敵對的勢力嗎?還是我們想像不到的邪惡呢?經文所講的,是敵人對我們的一些敵對狀況,但同時,又對這些可能發生的情形提供答案.在我們四周,無論在內或外,都有成千上萬的敵軍攻擊我們,我們怎能應付呢?這裡講到神永恆的救恩所給我們的幫助.
\newline
\newline
首先作者講到我們在神面前不再被定罪,有了這地位之後,我們便可以實踐公義的生活.而且神應許我們將來要得榮耀的身體,可以永遠與祂同在.這段經文是一首凱旋的詩歌,表示絕對的信心,和得勝的確據.
\newline
\newline
在四面受敵之際,我們怎能唱得勝的凱歌呢?秘訣是在於我們有盟軍的幫助,與我們同走人生的每一步,這盟軍就是聖靈,而在經文結束之時,更提及聖父和聖子.天父是幫助我們的,只要我們在耶穌基督裡.無論四周的敵人是多麼兇猛,勢眾；無論他們的控告是甚麼,天父都會幫助我們,因為祂愛我們.三十五節說,誰能使我們與基督的愛隔絕呢?那就是說,有誰能停止祂愛我們呢?甚麼都不能.
\newline
\newline
神的心充滿了愛,神的愛在祂兒子的身上是豐盛的,這豐盛的愛湧溢出來,臨到那些在基督裡面的人.因為祂愛我們,所以祂便幫助我們,作我們的同盟.祂以永遠的愛愛我們,這是難以明白的.祂選擇愛我們,為我們犧牲,為我們而活,祂對我們的愛永不改變.
\newline
\newline
我們怎樣知道神愛我們呢?卅二節說,神既不愛惜自己的兒子,為我們眾人捨了,豈不是證明祂愛我們嗎?耶穌是父神最愛的,但神竟然把祂賜給我們,這豈不表示神是愛我們嗎?神將宇宙最珍貴的,祂的獨生愛子賜給我們.這說明了祂深愛我們,所以祂幫助我們.當我們看見神藉耶穌所為我們成就的事時,我們就知道神確實愛我們.
\newline
\newline
在眾多仇敵環繞我們之際,我們有一個偉大的同盟,就是父神.除了父神與我們同盟之外,耶穌基督也是幫助我們的.卅四節說祂為我們死了,又為我們復活了,現在正活著,為我們代求.祂以屬祂的人為祂關心和代求的對象,祂也是我們的同盟,要幫助我們.雖然在我們前進的路上,有眾多的仇敵,但父神和耶穌都與我們在一起,作我們的同盟,祂要和我們一同作戰,因為祂愛我們,而沒有在任何事物可以攔阻祂的愛.
\newline
\newline
以利沙曾經對他的僕人說,(王下六16)說,與我們同在的,比與他們同在的更多.這正好是我們的寫照.我們的同盟比我們的敵人更偉大.結果怎樣呢?卅七節說,我們是得勝有餘的人.在一切的困苦,逼迫,飢餓,赤身露體,危險之中,我們可以得勝.不管在前面有甚麼困難,敵人,但靠著那幫助我們的,我們是能得勝有餘.
\newline
\newline
雖然有仇敵攻擊之時,我們可能會憂傷,會失敗,但我們卻不會就此被消滅,或者永遠失敗.保羅在(林後四8-9)說:「我們四面受敵,卻不被困住；心裡作難,卻不至失望；遭逼迫,卻不被丟棄；打倒了,卻不至死亡.」有時我們會覺得自己的地位卑微,或是看見可怕的將來而不知所措,但是因為我們的盟軍比敵人強,所以我們不致被消滅.我們是得勝有餘的人.
\newline
\newline
當戰爭完畢時,神要邀請我們進祂的國中,到那時候便不再有戰爭,只有安息,我們可以享受一切的戰利品,享受耶穌征服敵人之後所給我們的福樂.不再有戰爭,這是何等美好的結果.雖然我們四面受敵,但因為我們有盟軍,所以勝利終於是屬於我們的.
\newline
\newline
弟兄姊妹,如果你已經是相信耶穌,是在基督裡的人,讓我們一同分享這歡樂,感謝神的大愛,並一同經歷到神為我們戰爭之後的安息.我們可以對保羅所唱的凱歌產生共鳴,因為我們也在這得勝的行列中.我們可以不應再怕前面的困苦,逼迫,飢餓,危險,刀劍,只要我們緊緊的依靠我們的盟軍,我們就必能享受最終的勝利.
\newline
\newline


\newpage

\section{第57屆~港九培靈會~鮑會園~第一講　摩利沙人彌迦──時代的先知}
\label{sec:472}
\textbf{第一講　摩利沙人彌迦──時代的先知}
\newline
\newline
連結: \href{https://www.hkbibleconference.org/session-message/view/472}{\texttt{https://www.hkbibleconference.org/session-message/view/472}}
\newline
\newline
\hyperref[sec:490]{< < < PREV SERMON < < <}
~
\hyperref[sec:toc]{[返目錄]}
~
\hyperref[sec:473]{> > > NEXT SERMON > > >}
\newline
\newline
感謝主讓我有機會在培靈會裡,與弟兄姊妹交通,當我耐心等候信息的時候,神已經向我的心講話,我與大家分享的信息,是舊約中一位不大著名的先知.但這信息卻特別寶貴.他處於一個特別的時代,神藉著彌迦傳達信息.我覺得今天的時代與先知彌迦的時代很接近,神會藉彌迦的信息,給我們今天有實際的幫助.
\newline
\newline
彌迦住在耶路撒冷一小地方,聖經繙作「摩利沙」.摩利沙離耶路撒冷南五六十華里,靠近迦特的境界,(彌一14)稱這地方為「摩利沙迦特」.摩利沙處於猶太人與迦特的地方,故站於一處獨特的地方,兩處均較重視此地.它受重視的原因,並非因它有甚麼了不起的地方,乃因這地對兩國都有重要的價值.這地有很重要的軍事價值,故猶大王常常南下巡視,看看這地方有沒有受到威脅；這地雖小,卻常常接觸到大人物.彌迦出來工作的時候,沒有人重視他的家鄉；因為在這時候,南國北國處於特別的環境,北國物質非常豐富.當時由耶羅波安二世在位執政,他將國家發展到非常興盛；故此北國的人覺得甚麼都可以做,完全靠自己的力量,沒有一點依靠神的心.而南國的一位好的君王烏西雅剛剛逝世,烏西雅王亦把南國帶領得非常興盛.而且他非常敬虔,帶領猶大人親近耶和華；但烏西雅死後,他的兒子約坦接續為王,約坦雖然未曾離開正路,但他卻是一位非常軟弱的王.約坦死後,亞哈斯繼位,情形就有了轉變；亞哈斯有本事及野心,但卻沒有敬畏神的心,使全國走上背叛的路.亞哈斯死後,希西家接續為王,他是一位非常敬虔的王,一開始便敬畏神,但由於百姓在亞哈斯時代離開神太遠,故不能把百姓帶到神面前,罪惡仍十分普遍,希西家很追求神,但臣僕卻非常敗壞.故此在這時候,南北國非常敗壞,而此時,東方又有一個強大的國家興起.亞述演變成一個強大的國家,不久之後亞述會毀滅北國以色列；就在這時候,神在北國興起何西阿先知,而南國神興起以賽亞及彌迦兩位先知,分別向國人說話.
\newline
\newline
按時間而言,彌迦出道較以賽亞稍遲,但可能兩人有段時間曾經是同工.當年南北國非常混亂,興起了很多先知；在先知中,真正敬畏神的可能只有少數.而在百姓中自稱是先知的卻有許多人,但實際上他們不過是替自己講說話.在這時期中百姓沒有敬畏神的心,而北面又有一個強盛大國來威脅.百姓中許多假先知只知為自己說話.彌迦正於這時出來工作,可以想像得到他的工作實在不容易；因為君王領袖不看重他,百姓也不接納他,好像曠野孤獨的聲音,但神卻與他同在,神將信息給他.在歷史上很少提及彌迦這個人,除本書中,只在(耶十六16-19)提及他,責備百姓及領袖；但猶大王希西家卻沒有傷害他,而尊重他為先知.雖然百姓想把先知忘記,但神卻不容許百姓忘記先知,故此先知的信息仍然有效地講說話.
\newline
\newline
我們看看彌迦作先知的特性,他非常忠心事奉神,雖然四周有人反對,他卻忠心完成神的心意；他有幾方面的特性,值得我們留意的.
\newline
\newline
(一)先知有使命感
\newline
\newline
他知道在混亂的時代,神有工作要他完成,(一18):「先知說,我必大聲哀號,赤腳露體而行,又要呼號如野狗,哀鳴如駝鳥.」神讓他知道百姓有實際的需要,他必須要完成神的工作.雖然工作有困難,但仍然要做；因為神的旨意,無論是在彌迦的時代,或任何時代都是一樣.神給我們工作,要我們去完成,不是容易的事.神對約翰說,你要把書吃到肚裡,吃的時候很甜蜜,但到肚裡時卻很苦；從神來的信息,許多時候,會令你感到很痛苦；傳的時候不容易,但神已將託付給了你,你必須要忠心去完城.神若果今天呼召我們去工作,按照我們理解,是一件很痛苦的事；但因為是從神而來的呼召和使命,我們就要去完成.
\newline
\newline
(二)我們看見先知有真正的正義感
\newline
\newline
神是公義的,如果我們從神領受信息,這信息應該是正義的,但在傳達時就不容易.在(彌三5)彌迦嚴厲責備假先知,是因為假先知使民走錯了路,故此彌迦要責備他們.審判的標準,乃是按照先知使百姓走錯路或走正路而定；照樣今天我們所傳的信息是否須要看看環境,但環境不能決定我們的信息.我們須要看周圍的人,但信息也不能因周圍的人而決定.我們須要看自己工作的能力,和生活的環境,但這些仍不能決定我們所講的信息；因為我們必須以是否使百姓走上正路來作決定.如果信息會使人走錯了路,他便是錯了.這時候要反對和責備不容易,因為已走了錯路,我們可能也如此.
\newline
\newline
(三)他與百姓有真正的認同
\newline
\newline
他站在百姓的地位；百姓所遭遇的,正是他所遭遇的.百姓遇見的痛苦和欺壓,這些事雖然不會臨到彌迦的身上,但是他仍然站在百姓的地位,去責備那些人.(一 8),赤腳露體形容被擄的情形,他們已失去主權,沒有甚麼權利可言.大概他是指著亞述擄去北國以色列民的情形而言.但是彌迦的遭遇不同,他是站在百姓痛苦的地位上,與百姓有完全的認同.百姓所犯的罪和痛苦,就是他的罪和痛苦.在聖經裡神所使用的領袖,都有這種心志.當百姓向神認罪禱告.他說:「我和我的百姓犯了罪.」其實,摩西沒有犯罪,只不過他站在百姓的地位上進言．尼希米帶領以色列人重建耶路撒冷的城牆,他看見百姓犯罪,心裡非常難過；便撕裂衣服,把爐灰撒在自己的頭上.在神的面前認罪悔改說,「我和我的百姓犯了大罪.」尼希米願一起受罰,如果要受痛苦,他一起受痛苦.在新約這意思更加清楚,主耶穌完全站在百姓的地位上；若不是如此,祂毋須上十字架,因為祂完全站在我們的地位.使徒保羅在歌羅西書第一章,說自己為代表教會受苦難,補滿基督的缺欠.沒有一位神的僕人,會覺得自己與百姓是毫無關係的；百姓的痛苦和苦難,就是他自己的痛苦和苦難.百姓的罪過,自己要承擔,在百姓悔改的時候,亦帶領百姓一同悔改.今天我們需要有這種認同感,大家都是基督的身體,一個肢體受苦,全體都受苦,在任何的情況下,我們不能看輕與基督徒的這種關係.
\newline
\newline
(四)他有勇氣在這特殊的環境下為神工作
\newline
\newline
在這本書上,他責備兩種人,責備這兩種人都不容易,(三5)指他責備的第一種人是假先知.彌迦是先知,這些領人走錯路的也是先知,現在他看見的同工走錯路,他要站在神的話語上來責備這些假先知.責備人不是容易的事,而責備自己的同輩同工,更加不容易；而且一個人常常責備他的同輩,誰會接納他呢?他會被孤立起來,這是一件很難受的事.所以要責備自己的同輩,需要有勇氣站在神話語上；無論他是甚麼人,只要他錯了,我們便要責備他.在今天我們很重視被人接納,如果被人孤立,這是很難受的；故此我們常做事與人一樣,因為這容易被人接納.無論別人做甚麼事,雖然我知道不大好,我也跟他一樣做,否則便不接納.別人怎樣打扮,我們也怎樣打扮,為的是想別人接納我.
\newline
\newline
彌迦亦有勇氣責備百姓的領袖,(三1)「雅各的首領,以色列家的官長阿,你們不當知道公平麼?(三9)「雅各家的首領,以色列家的官長阿,當聽我的話,你們厭惡公平,在一切事上屈枉正直.」,今天的時代與當時不同,今天我們處於民主時代；在報紙上常看見有人罵政府官員,甚麼罵英女皇或港督,都不是大不了的事.但在以前的中國,你敢罵皇帝,性命也保不了；在那個時代誰敢講皇帝或官長的壞話,一句壞話便有生命的危險.但彌迦有從神來的啟示和話語,他不但有勇氣責備假先知,亦有勇氣指出領袖的錯誤.彌迦先知態度不是自負,而有真正的謙卑,按詞句而言,彌迦可能是摩利沙的長老,常常因公事往耶路撒冷.他是有地位的人,但整卷書沒有提及自己的權利和地位.他只說:「耶和華如此說.」他把自己完全隱藏在神話語下,無論是顯著或卑微的工作,不是表現自己,而是神的託付.他工作目標不是表現自己,不是知道我彌迦是一個怎樣了不起的人,而我只是完成神的託付,故此他是一個謙卑的人.
\newline
\newline
現在我們簡單地看整卷書的重點,整卷書可以用「禍哉」來表達.因為百姓犯罪,所以神宣召刑罰和審判,當時百姓貪圖物質的好處,只管貪圖金錢.(二2)「他們貪圖田地就佔據,貪圖房屋便奪取.」(三5)「他們牙齒有所嚼的,他們就呼喊說平安了.」他們生活的目標就是今天,今天的人也看重物質的好處和自己的生活,前途怎樣我不管,我看重今天的生活享受.神責備百姓的第二個錯誤,就是是非不分,沒有正義的標準,(三9)「你們厭惡公平,在一切事上屈枉正直.」沒有公平正直,要的是權柄地位,我不理會合乎神的標準與否,只要我今天得到好處便好了.另外他們還有一種假冒為善的態度,看來好像很敬虔,但其實沒有神的敬虔,(三11)「首領為賄賂行審判,祭司為僱價施訓誨,先知為銀錢行占卜；他們卻依賴耶和華,說耶和華不是在我們中間麼?災禍必不臨到我們.」好像耶和華成為我們的工具,要耶和華的賜福,這是一種假敬虔.因為他們的罪過,神宣佈禍哉的信息,但彌迦的信息與其他聖經的信息,永遠不是消極的,如果神的百姓肯順服悔改,神就給他們寶貴的應許.第五章預言耶穌基督的降生,在第四章提及神施恩百姓的情況.這本書提醒百姓的錯誤,如果他們知道自己的過錯而悔改,神一定會施恩憐憫,我們需要預備自己的生活,來領受神的福氣,先知的信息,對以色列百姓非常實際,照樣對我們也很實際.
\newline
\newline
本講道集蒙黃民弟兄,梁壽華弟兄,雷麗端姊妹記錄,林雲坤弟兄封面設計,朱建磯牧師編輯,林小娟姊妹協助校對,得以如期出版,謹此致謝!
\newline
\newline
除講道集外,現場之錄音,錄影,均由福音傳播中心負責總經銷.
\newline
\newline
每年大會叨蒙九龍城浸信會,借用堂址為會場,宣道中學宣中堂借用禮堂為副會場.兩堂執事,助道部,同工,義工,熱誠招待,參予工作,均統此致謝!
\newline
\newline
在五十七屆奮興會中,得復興獻身事主者,有二百五十五人,來自一百一十八間教會.願將果子榮耀頌讚歸與父神!阿們!
\newline
\newline


\newpage

\section{第57屆~港九培靈會~鮑會園~第二講　信息的權威}
\label{sec:473}
\textbf{第二講　信息的權威}
\newline
\newline
連結: \href{https://www.hkbibleconference.org/session-message/view/473}{\texttt{https://www.hkbibleconference.org/session-message/view/473}}
\newline
\newline
\hyperref[sec:472]{< < < PREV SERMON < < <}
~
\hyperref[sec:toc]{[返目錄]}
~
\hyperref[sec:475]{> > > NEXT SERMON > > >}
\newline
\newline
彌迦書 1:1-1:7
\newline
\begin{longtable}{cl}
\hline
\hline
章節 & 經文 (和合本修訂版)\\
\hline
1:1 & \begin{tabularx}{0.7\textwidth}{X} 當猶大王約坦、亞哈斯、希西家在位的時候，耶和華的話臨到摩利沙人彌迦，他見到有關撒瑪利亞和耶路撒冷的異象。 \end{tabularx} \\ \\ \relax
1:2 & \begin{tabularx}{0.7\textwidth}{X} 萬民哪，你們都要聽！地和其上所有的，要留心聽！主耶和華要從他的聖殿指證你們的不是。 \end{tabularx} \\ \\ \relax
1:3 & \begin{tabularx}{0.7\textwidth}{X} 看哪，耶和華從他的居所出來，降臨步行地之高處。 \end{tabularx} \\ \\ \relax
1:4 & \begin{tabularx}{0.7\textwidth}{X} 眾山在他底下熔化，諸谷崩裂，如蠟熔在火中，如水沖下山坡。 \end{tabularx} \\ \\ \relax
1:5 & \begin{tabularx}{0.7\textwidth}{X} 這都是因雅各的罪過，因以色列家的罪惡。雅各的罪過在哪裡呢？豈不是在撒瑪利亞嗎？猶大的丘壇在哪裡呢？豈不是在耶路撒冷嗎？ \end{tabularx} \\ \\ \relax
1:6 & \begin{tabularx}{0.7\textwidth}{X} 因此，我必使撒瑪利亞變為田野的廢墟，用以栽植葡萄；我必把它的石頭倒在山谷，掀開它的地基。 \end{tabularx} \\ \\ \relax
1:7 & \begin{tabularx}{0.7\textwidth}{X} 城裡一切雕刻的偶像必被打碎，行淫的賞金全被火燒，我要毀滅它的一切偶像；因為從妓女的賞金積聚而來的，它們仍歸為妓女的賞金。 \end{tabularx} \\ \\
[1ex]
\hline
\hline
\end{longtable}
全本彌迦書有七章,而可分成三大段,有人把它分成四段好像特別的信息,每一段開始用同一樣說話.中文作「你們都要聽」,原文是「聽阿」,意思說神有重要的信息給你,你要留心去聽.(一2,三1,六1),看見有這三段信息,每段內容信息都有相同,但每段的信息都有中心.第一篇信息指出神的審判,第二篇神審判的原因和改正的方法.第三篇講到美好的應許.按照文字的結構和內容,好像彌迦一次過的說話；當然事實上可能是先知在三段時間的三篇信息.但寫作時經過修飾,便看來是一篇信息；但三段信息連繫得很好,我們便將它看成一篇信息分成三段來看.百姓要受審判的時代,早一段時間寫成的何西阿書,看見責備很厲害,先知彌迦也得了這種信息.(一2)看見信息的對象是萬民,但整卷書是神對以色列或猶大所施行的審判,神對百姓的審判卻呼召萬民要聽.這點很有意思,神在舊約先揀選以色列人,但神的計劃最終不是單在以色列人身上；以色列人是神的器皿,創十二章神揀選亞伯拉罕,萬民都要因他得福.林前十二章提及選民在曠野的歷史,保羅說過去的教訓成為我們鑑戒.今天有些聖經學者,單單將舊約看成以色列人的歷史,若舊約單是以色列的歷史,這些說話便沒有資格放在聖經裡.故此神對彌迦的教訓,也是要教導後來做神百姓的人.在每一段歷史的過程中,對今天教會,都有屬靈實際的教訓；這裡雖是對以色列人說話,但宣召萬民去聽.
\newline
\newline
在思想神責備的信息之前,我們應該先思想彌迦信息的根基在哪裡,(一1)「當猶大王約坦,亞哈斯,希西家在位的時候,摩利沙人彌迦得耶和華的默示.」原文是作「神的話臨到了彌迦.」這裡神的話是很具體的字,不是普通的話,乃是很具體的字.下文亦很具體指出百姓的錯誤,神將具體的話告訴彌迦,叫他傳遞給百姓.只有這時候,他才出來向百姓講說話.彌迦是百姓的領袖,常與耶路撒冷的官員來往,他在耶路撒冷可以看見全國的情形.人必須站在那地位,才有整體的觀念.在教會中我們如果是做主日學或青年部的工作,我們所關心的是主日學或青年部,其他的事便不理會；但如果你是教會的堂主任或執事會主席,你就必須關心整個教會的事情.如果你是城鎮的官員,你所關心的是城鎮的事,但中央官員就須要關心全國的事.彌迦好像站在中央政府中,去看整個民族,他知道十年左右,以色列國會被亞述滅亡了；他也好像何西阿先知,知道南國有同樣的遭遇,若不悔改,便會滅亡!他站在百姓領袖的地位,很關心這些事情.從文字來看,彌迦工作已有多年,他已是成年人,已經觀察國家很久；他知道百姓應作的事是甚麼,知道百姓應走的方向.按照政治的局勢,百姓的演變和自己的觀察,他是應該說話的；但他沒有講過話,直到神的說話臨到他.因為神的話臨到他,他便站起來.雖然百姓不喜歡聽,且會引起麻煩,但他不顧一切,因他所講的話唯一的根基,是出於神.這一點值得今天的基督徒效法,我們做事和說話的根基是甚麼,對前途和決定的根基是甚麼?若按照人的觀念,必有許多說話和意見,因今天的人十分注重理論和方法.若從心理學看人的行為,又從社會潮流看人的趨勢,這些都會影響基督徒的看法.當然我們不能脫離社會,但真正決定我們行動和方向的,不是從社會的潮流,不是從心理分析得來的.也不是從政治環境歸納得來的；基督徒生活的標準,乃是根據神的話語.
\newline
\newline
彌迦先知信息的權威是甚麼呢?他依靠整個的力量是甚麼?從這幾節聖經,可以看見他與當時的人有很大的分別.這是先知站在神的地位上所看見的,在人看來他不應當說那種說話.在(一3-4)指出神的出現,有許多處也看見同樣的話,(士五4)看見在巴拉和底波拉帶領百姓讚美耶和華的話.「耶和華阿!你從西珥出來,由以東地行走,那時地震天漏,雲也落雨,山見耶和華的面就震動,西乃山見耶和華以色列神的面,也是如此.」(賽六十四1-2)「願你裂天而降,願山在你面前震動,好火燒乾柴,又像火燒開,使你敵人知道你的名,使列國在你面前發顫.」在敵人壓迫百姓的時候,看來沒有盼望的時候,突然耶和華裂天而降,這不是肉眼可以看見的,而在我們經驗以外,耶和華臨到我們.在人不肯承認耶和華存在的時候,耶和華會突然之間臨到我們,祂的作為就顯明出來.祂的出現,會改變人的處境,今天好像很多人看不見耶和華的存在,但神會顯現祂的作為.另一方面,當百姓不承認耶和華的存在,敵人看不見耶和華的記號的時候,在他們沒有把神劃入計劃內的時候,神的作為就顯明出來.今天許多時候,可能環境和社會的壓力,或我們心理的壓力,我們在計劃決定時,根本沒有把耶和華計算在內.(雅四 13-15)「嗐,你們有話說,今天明天我們要往某城裡去,在那裡住一年,作買賣得利.其實明天如何,你們還不知道,你們的生命是甚麼呢?你們原來是一片雲霧,出現少時就不見了；你們只當說,主若願意,我們就可以活著；也可以作這事,或作那事.」我們在肉眼看不見神的時代,在人生中沒有把神計算在其內,這是多麼愚拙.當西拿基立圍困耶路撒冷的時候,他譏笑希西家王,看情形猶大百姓沒有盼望了,西拿基立不把神計算在內；他只計算自己的能力和軍隊,豈料當晚神的使者就把亞述大軍十八萬五千人殺死.如果沒有耶和華的出現,百姓便沒有盼望了.多少時候,基督徒何嘗不是這麼愚拙,因為沒有把耶和華計算在內.彌迦在此要提醒百姓,耶和華要裂天而降,祂的作為是我們不能不計算在內的.
\newline
\newline
另一方面,我們看見耶和華一出現,祂便掌管整個局勢.當暗利和亞哈作北國以色列王的時候,他們選擇了在山上建造皇宮,建造房屋；於是一座很美的城便建立起來,這就是撒瑪利亞城,是以色列國的首都.耶路撒冷同是建造在山上,形勢是逐漸斜上去的；四圍都有路直達,故此打仗時不易防守.但撒瑪利亞城是垂直的,是一座很堅固的城,敵人不容易攻破.當神一出現時,人無論怎樣也沒有辦法對祂,且在神面前無法站立得住.(一6)「所以我必使撒瑪利亞變為田野的亂堆,又作為種葡萄之處,也必將他的石頭倒在谷中,露出根基來.」以前是荒野,現在也要變為荒野,人想不到的,神便成就.神要在祂的作為中顯出祂的權柄來；祂要掌管整個局勢,我們不能不承認神的權能.今天有許多人不承認神的權柄,可能是看不見神作為的記號,但神要工作時,便可以看見祂的權柄.有人說,我們等了很久仍未看見神的作為,正如很多主忠心的僕人受逼迫,為甚麼也看不見神為他們伸冤呢?大佈道家慕迪講道時,提及神的權威,叫人應當接受耶穌.有一個自命聰明的人,在慕迪講道後來見他,表示不接受耶穌,這人說這幾年基督徒的收成都很差,有許多人竟然把田地賣給我；我的莊稼總比你們信耶穌的人好.慕迪看看他便說:「神審判的時候,也許不在秋天收莊稼的時候.」神有祂自己工作的時間,祂會彰顯權能,也會掌管一切.
\newline
\newline
我們面對這信息,應有甚麼反應呢?(一2)「萬民阿,你們都要聽.」雖然我們不知道神在甚麼時候審判,但我們知道祂一定審判(一7)「他一切雕刻的偶像必被打碎,他所得的財物必被火燒,所有的偶像我必毀滅,因為是從妓女僱價所聚來的,後必歸為妓女的僱價.」神會向他們討價,因為他們對神的呼召不忠心,照樣今天神雖然尚未審判我們,我們也要思想是否對神的託付忠心?如果我們不忠心於神的託付,有一天神會與我們計算賬目.神吩咐教會傳福音給普天下人聽,我們盡了本份未呢?神要我們得著從上面來的能力,然後要在耶路撒冷,撒瑪利亞,直到地極作祂的見證.我們對主的託付,付上多少的忠心呢?如果我們對神的託付不忠心,我們也不能逃避祂的審判.彌迦從神領受信息,他不顧一切將神的話傳達；但願今天神向我們講話,讓我們有一個不顧一切傳講神話語的心.
\newline
\newline


\newpage

\section{第57屆~港九培靈會~鮑會園~第三講　擊打和醫治}
\label{sec:475}
\textbf{第三講　擊打和醫治}
\newline
\newline
連結: \href{https://www.hkbibleconference.org/session-message/view/475}{\texttt{https://www.hkbibleconference.org/session-message/view/475}}
\newline
\newline
\hyperref[sec:473]{< < < PREV SERMON < < <}
~
\hyperref[sec:toc]{[返目錄]}
~
\hyperref[sec:476]{> > > NEXT SERMON > > >}
\newline
\newline
彌迦書 1:8-1:16
\newline
\begin{longtable}{cl}
\hline
\hline
章節 & 經文 (和合本修訂版)\\
\hline
1:8 & \begin{tabularx}{0.7\textwidth}{X} 為此我要大聲哀號，赤身赤腳行走；我要呼號如野狗，哀鳴如鴕鳥。 \end{tabularx} \\ \\ \relax
1:9 & \begin{tabularx}{0.7\textwidth}{X} 因為撒瑪利亞的創傷無法醫治，蔓延到猶大，到了我百姓的城門，直達耶路撒冷。 \end{tabularx} \\ \\ \relax
1:10 & \begin{tabularx}{0.7\textwidth}{X} 不要在迦特宣揚這事，千萬不要哭泣；卻要在伯‧亞弗拉翻滾於灰塵中。 \end{tabularx} \\ \\ \relax
1:11 & \begin{tabularx}{0.7\textwidth}{X} 沙斐的居民哪，要赤身羞愧地經過，撒南的居民不敢出門，伯‧以薛哀哭，不再支持你們。 \end{tabularx} \\ \\ \relax
1:12 & \begin{tabularx}{0.7\textwidth}{X} 瑪律的居民心甚憂急，切望得著福氣，因為災禍已從耶和華那裡臨到耶路撒冷的城門。 \end{tabularx} \\ \\ \relax
1:13 & \begin{tabularx}{0.7\textwidth}{X} 拉吉的居民哪，要用快馬套車；錫安的罪由你而起，以色列的罪過在你那裡顯出。 \end{tabularx} \\ \\ \relax
1:14 & \begin{tabularx}{0.7\textwidth}{X} 因此，你要將送別禮送到摩利設‧迦特；亞革悉的眾家族必用詭詐待以色列諸王。 \end{tabularx} \\ \\ \relax
1:15 & \begin{tabularx}{0.7\textwidth}{X} 瑪利沙的居民哪，我必使搶奪者來到你這裡；以色列的貴族必來到亞杜蘭。 \end{tabularx} \\ \\ \relax
1:16 & \begin{tabularx}{0.7\textwidth}{X} 猶大啊，為了你所喜愛的兒女，你要剪髮，剃光頭，要使你的頭光禿，如同禿鷹，因為他們被擄去離開你了。 \end{tabularx} \\ \\
[1ex]
\hline
\hline
\end{longtable}
在舊約中小先知書比較困難閱讀和解釋的,除兩三本外,其餘是用詩歌的體裁寫成的.詩歌是藝術的作品,藝術不可以用科學來分析.詩歌主要是傳遞信息,而不是逐個細節分析.但不是說詩歌的材料不準確,裡面有些材料很精確美麗；特別是哈該,他用的詞句很好,叫你聽過永不會忘記.(一10)「伯亞弗拉輥於灰塵之中.」伯亞弗拉就是灰塵的意思,在一個灰塵的城鎮,輥於灰塵之中.在猶太人中灰塵是代表悲哀,這種詞句讀後不易忘記,但不易解釋.整本彌迦書詩句很美,唱熟後不易忘記,用詞精簡.可惜當時百姓只顧生活,沒有心存記神的話語.如果彌迦向他們長篇大論,他們會左耳入右耳出；但彌迦用容易上口的詩歌,他們不需特別專心,也會記著.今天我們傳福音時,以為有些說話很重要,但聽的人可能覺得不重要.我們要把重要說話精簡地指出,特別在重要的場合；使聽的人覺得很重要,否則就錯失機會了,最近有人談論福音傳播時,指出每句話都要有內容,無用的話可不必多說.
\newline
\newline
剛才的經文指出幾個重要的信息,百姓應當為他們的罪受刑罰,然後神在刑罰會顯出恩典.首先我們思想百姓所犯的是甚麼罪,頭一個罪在(一13)「拉吉要用快馬套車.」拉吉離耶路撒冷幾十哩,在非利士人的邊界,是幾百年來猶大國重要的邊防重鎮.在軍事上非常重要,撒母耳記及列王記均提及這裡有許多馬匹軍隊；為了保守邊疆,猶大王將一切力量放在這裡,靠這地方防守耶路撒冷.故此將最好軍隊馬匹和一切的依靠,也放在這地方；可見整個百姓的信心放在軍隊上,神的僕人和百姓,不再依靠耶和華了.不再靠神的作為,而是靠自己的力量和方法；倘若神的百姓這時再把信心放在神那裡,神也不會再幫助他們了.這是多麼可怕的事,在撒母耳記上的末章,記載神要試驗百姓,激動大衛,去數點百姓,結果神刑罰他們.從經文看:「我們中間可以拿兵器的人有多少的話語中,可見大衛當時數點有多少百姓可以打仗,將盼望和依賴放在自己的軍隊上,故此神要刑罰百姓.彌迦時代也一樣,百姓把自己的信心也放在軍隊上,我有能力去對抗敵人,軍隊能保衛國家.今天教會也落在這種試探中,在計劃工作時,看自己有多少力量和方法,自己有多少經驗去工作,而忽略仰望依靠神的心.
\newline
\newline
百姓所犯的第二種罪,(一13)「錫安民的罪,由你而起.」這句話很難繙譯.天主教聖經繙作「錫安民犯罪的開端.」呂振中譯本作「錫安民犯罪的起點.」這兩個意思跟這裡的不同,應作「錫安民的罪,在這裡只是開端起點.」先知提醒百姓要思想犯罪的開始,以色列人在曠野親自體驗耶和華的真實,他們真知道耶和華的救恩.當百姓進入迦南後,看見迦南人所拜的神；上一代經驗過神真實的,已經過去了,下一代覺得迦南人的神也不錯.甚至覺得敬拜迦南人的神更好,可以打勝仗.以色列百姓漸漸忘記,他們在曠野所經驗的耶和華,效法迦南人,走上拜偶像的道路.先知提醒百姓要想念犯罪的開端,不要忘記你們最初經驗的耶和華.在今天我們也許經驗過神的真實,未信主時我們的罪擔很重,但聖靈開啟我們的心,我們在神面前認罪,罪完全赦免了.經驗了從罪釋放的喜樂自由,與主關係非常親密,但日子過了一段時間；罪得赦免的經驗已經淡忘了,蒙恩的事已有多年了.有看見四圍的人生活,他們沒有經驗神的真實,他們的生活也很好,做生意發達,讀上很好的學校；過去的經驗雖然很真實,卻被環境所勝過了.單單看見別人外表的興旺,忘記神的真實；效法四周的人生活的方式,這是多麼可怕的試探.在將來信心要受到特別考驗的時候,看見那些沒有經歷神的人,他們的生活也很好,工作和生活,比我們還好；也許我們會忘記耶和華的真實,就會受到四圍的環境所影響.以色列誠然忘記了經歷神的真實,而去敬拜四周圍人的神.
\newline
\newline
現在我們看他們犯罪的結果,因著他們犯罪,便受到神的擊打.他們以為神的刑罰和自己所犯的罪是沒有關係的,當然我們生活中的困難,不完全是得罪神.在這裡先知告訴百姓,因著他們的罪,便受兩方面的刑罰,頭一個刑罰是被擄.(一16)「猶大阿!要為你所喜愛的兒女剪除你的頭髮,使頭光禿,要大大的光禿,如同禿鷹,因為他們都被擄去離開你,耶路撒冷快要滅亡了!在未滅亡前,神警告他們的兒女會被擄去,因這是神的刑罰.亞述王西拿基立寫了一篇記錄,提及自己第一次征服耶路撒冷的事:「因為猶大王希西家不肯接受我的軛,我圍困了他的四十六座城並四周的村鎮,我用兵馬刀劍,將他完全征服；我從他的各城中,趕出二十萬零一百五十人.將他們的男女老少,騾駝,全部俘擄,將希西家囚禁在耶路撒冷,如同鳥在籠中一樣.」連他們的王希西家也被擄去,如鳥在籠中一樣,我們不能把這看為平常；因為在未發生前,神已預言,這是神的刑罰.我們不能說這些事今天不可能發生,神的百姓不對神順服,神的審判便會臨到,在今天這種歷史事實是可以重演的.
\newline
\newline
百姓所受的另一種刑罰,(一15)「以色列的尊貴人必到亞杜蘭」,「尊貴」小字作「榮耀」,以色列的榮耀必到亞杜蘭.當掃羅逼迫大衛,想用盡方法殺死大衛,大衛只可逃跑.有一次他差不多被掃羅的軍隊追上,他無法逃跑,只好走進一山洞,藏在那裡,這洞的名字叫做亞杜蘭.先知也是指這件事,大衛逃跑時要躲藏起來,否則會落在掃羅手中喪命；現在以色列的榮耀要到亞杜蘭,照樣要隱藏起來.榮耀在舊約中是很重要,榮耀基本上是神向人顯現的光榮；有神的同在,一同在神的榮耀彰顯.所羅門為神建造聖殿,建成後神與人同在,神的榮耀便充滿這殿,結果百姓不能進入殿去；以色列人的榮耀,是神自己的彰顯.神要住在百姓中,百姓可以享受與神同在的福氣；但現在以色列的榮耀卻要隱藏起來,像大衛在亞杜蘭洞裡躲藏起來一樣.神沒有忘記祂的選民,但神卻向選民掩面,把榮耀隱藏起來.神雖然愛祂的百姓,但百姓卻經驗不到神同在的事實,不能享受與神同在的喜樂；雖然與神的關係並沒有改變,不過與神的交通卻斷絕了,這是多麼可怕的刑罰!我們是屬主的人,當我們接納主之後,每一個屬主的人,都有聖靈在我們的生命中.人若沒有聖靈在生命中,便不是屬於基督的；我們有神的生命是永不改變的事實,不會因你今天做得不好便離開你.我們可以有屬靈生命的事實,但卻不能享受這種屬靈生命的喜樂,失去與主交通的甜蜜；雖然我們仍是基督徒,但作基督徒不再成為我們的喜樂.我們是得救的人,但得救的經驗不在我們的生命上發生作用了,享受不到喜樂和享受了.我們可能仍作基督徒的事,但不再覺得是享受了；基督徒不作的事,我們可能也不作,但心裡會有一種錯失的感覺,心裡很想去做.我們可能也讀聖經,與其他信徒一樣,但讀的時候,享受不到與主交通的甜蜜.如果我們覺得作基督徒是一種重擔,享受不到與神交通的甜蜜,這是神的極大刑罰,以色列因為犯罪,會被外邦人攻破,失掉與主交通的甜蜜.
\newline
\newline
但是神雖然刑罰祂的百姓,祂仍然憐憫百姓,祂會擊打,也會醫治.何西阿先知也說過同樣的話,「祂撕裂我們,也必醫治；祂打傷我們,也必纏裹.」(何六1)神永不會忘記祂的百姓,在有需要的時候,神會刑罰百姓；但又會醫治百姓；有時候刑罰是神的恩典,痛苦不容易忍受,但痛苦卻是福氣.我們發熱,肚痛是不易受,但這是恩典,因為這種記號,使我們知道自己有病；如果有病在我們裡面沒有徵兆,可能發覺此病時,便沒有希望了.如果我們在生活中能經驗因罪而來的痛苦,這是神的恩典,使我們可以甦醒過來；那麼這裡提及到百姓的刑罰,不論是被擄或神將他的榮耀隱藏起來,這些刑罰是直接從神而來的.以色列人所受的痛苦,不是單單因為外邦很強大,照樣我們與神的交通也不是單靠環境來限制,神會直接向我們施刑罰.雖然如此,神仍向我們施恩典,因神所做的一切事,目的使我們領受祂的恩典.當大衛數點百姓犯罪後,神差遣先知向大衛說,因你犯罪你要受刑罰.我給你三樣災選擇,你要落在外邦人手裡被追趕三個月呢,是有七年的饑荒臨到,還是國中有三日的瘟疫臨到；結果大衛願意落在神的手裡,因為祂有豐盛的憐憫!我們犯罪要受神的刑罰,但我們落在神的手裡,按照神的先知彰顯神的恩典.(一8)「赤腳露體而行」,彌迦在耶路撒冷工作,但他責備以色列國；以色列國與猶大國多少有些仇恨,現在以色列國受刑罰,豈不是南國的人要快樂嗎?但先知說自己也要受痛苦,並非因為有距離,也沒有因與他們沒有交通,把他們忘記.他仍然覺得大家同屬一個肢體,為他們心裡有極大的痛苦.百姓的刑罰和痛苦,就是先知的刑罰和痛苦；這種肢體的交通,巴不得激勵我們的心.每一個主內肢體,如果落在痛苦中斷不能因為大家有距離或沒有交通,便沒有同情的感覺.我們不可忘記那些受痛苦和受逼迫的人,因為他們也是我們的肢體.
\newline
\newline


\newpage

\section{第57屆~港九培靈會~鮑會園~第四講　罪中的大罪}
\label{sec:476}
\textbf{第四講　罪中的大罪}
\newline
\newline
連結: \href{https://www.hkbibleconference.org/session-message/view/476}{\texttt{https://www.hkbibleconference.org/session-message/view/476}}
\newline
\newline
\hyperref[sec:475]{< < < PREV SERMON < < <}
~
\hyperref[sec:toc]{[返目錄]}
~
\hyperref[sec:477]{> > > NEXT SERMON > > >}
\newline
\newline
彌迦書 2:1-2:13
\newline
\begin{longtable}{cl}
\hline
\hline
章節 & 經文 (和合本修訂版)\\
\hline
2:1 & \begin{tabularx}{0.7\textwidth}{X} 禍哉，那些在床上圖謀罪孽、籌劃惡事的人！天一亮，他們因手中有能力就去行惡。 \end{tabularx} \\ \\ \relax
2:2 & \begin{tabularx}{0.7\textwidth}{X} 他們看上田地就佔據，貪圖房屋便奪取；他們欺壓戶主和他的家庭，霸佔人和他的產業。 \end{tabularx} \\ \\ \relax
2:3 & \begin{tabularx}{0.7\textwidth}{X} 所以耶和華如此說：看哪，我籌劃災禍降與這家族；這災禍在你們頸項上無法解脫，你們也不能昂首而行，因為這是災禍的時刻。 \end{tabularx} \\ \\ \relax
2:4 & \begin{tabularx}{0.7\textwidth}{X} 到那日，必有人為你們唱詩歌，用悲哀的哀歌哀號，說：「我們全然敗落，我百姓的產業易主了！耶和華竟然使它離開我，我們的田地為悖逆的人所瓜分了！」 \end{tabularx} \\ \\ \relax
2:5 & \begin{tabularx}{0.7\textwidth}{X} 因此，你必無人能在耶和華的會中抽籤拉繩。 \end{tabularx} \\ \\ \relax
2:6 & \begin{tabularx}{0.7\textwidth}{X} 他們傳講說：「不可傳講；人都不可傳講這些事，羞辱不會臨到我們。」 \end{tabularx} \\ \\ \relax
2:7 & \begin{tabularx}{0.7\textwidth}{X} 雅各家啊，可這麼說嗎？耶和華沒有耐心嗎？這些事是他所行的嗎？我的言語豈不是與行動正直的人有益嗎？ \end{tabularx} \\ \\ \relax
2:8 & \begin{tabularx}{0.7\textwidth}{X} 然而，近來我的百姓興起如仇敵。你們剝去那些安然行路、不願打仗之人身上的外衣， \end{tabularx} \\ \\ \relax
2:9 & \begin{tabularx}{0.7\textwidth}{X} 把我百姓中的婦人從安樂家中趕出，又將我的榮耀從她們孩子身上永遠奪去。 \end{tabularx} \\ \\ \relax
2:10 & \begin{tabularx}{0.7\textwidth}{X} 起來，走吧！這裡並非安歇之處；因為不潔淨帶來毀壞，且是大大的毀壞。 \end{tabularx} \\ \\ \relax
2:11 & \begin{tabularx}{0.7\textwidth}{X} 若有人心存虛假，用謊言說：「我向你們傳講可得清酒和烈酒」，那人就必作這百姓的傳講者。 \end{tabularx} \\ \\ \relax
2:12 & \begin{tabularx}{0.7\textwidth}{X} 雅各家啊，我定要聚集你們，定要召集以色列的餘民，將他們安置在一處，如波斯拉的羊，又如草場上的羊群，人數眾多，大大喧嘩。 \end{tabularx} \\ \\ \relax
2:13 & \begin{tabularx}{0.7\textwidth}{X} 開路的在他們前面上去，直闖過城門，從城門出去；他們的王在前面行，耶和華在他們的前頭。 \end{tabularx} \\ \\
[1ex]
\hline
\hline
\end{longtable}
在剛才讀過的彌迦書第二章,先知指出人不看罪為大罪,只看為是當然的事；但先知卻非常著重,所以在開始時,就說禍哉!「禍哉」在原文與中文的聲音差不多,原文多「哀哉」,「悲哀」的意思.有很嚴重的災禍臨到百姓,先知心裡為百姓悲哀.這是先知的使命,傳出審判的信息；但是人卻不接受這種令人失望的信息.聽者的心裡不好受,叫人失去平安不是好事；叫人心裡受威脅,百姓的領袖也不喜歡審判及悲哀的信息.所以(一6)百姓領袖說:「你們不可說豫言.」在混亂時代,我們應設法安定人心.在第二次世界大戰時,許多報紙雜誌指出,最要緊的是安定人的心.凡能使人心裡受困擾的話要少講；叫人心裡得安慰,士氣振奮的信息可以多傳.這是百姓所要聽的,政治家也要聽的.今天的時代也一樣,我們心裡只喜歡聽叫人滿足的信息和得鼓勵的信息；而使人困擾動搖的話最好不多講.但神的信息要振動人的心,否則不會悔改轉回；今天我們喜歡的是人的信息,抑或神的信息呢?
\newline
\newline
另一方面,神呼召人叫人過分別為聖的生活.(二10)「你們起來去罷,這不是你們安息之所,因為污穢使人毀滅,而且大大毀滅.」我們要順服神的呼召,脫離污穢的事,但百姓不肯接受；卻順從虛假,滿口謊話,喝清酒濃酒說的叫人得安慰的話.令人麻醉心裡舒服,對真實失去敏感.這些假先知的信息,只叫人在威脅下得假平安,叫人在真實神的話語前妥協和不認真.今天人所喜歡接受的,是那些令人得安慰的信息；神卻叫人離開可憎惡的罪惡,叫百姓離開所居住的地方.神令百姓被擄到外邦,他們被擄是因為污穢了居住的地方,污穢使人毀滅!如果百姓能完成神的託付,百姓離開自己的地方,雖然是痛苦的事,但可避免再被罪惡所困擾,不再向罪惡妥協.教會在今天要講分別為聖的信息實在是不容易的事；因為人已漸漸被環境影響吸引,效法四周環境而生活.為了追求教會的合一,連公開不接受聖經真理的教會也接納過來.如此,亦只不過僅有合一形象,實際上卻違背真理和神的心意.個人生活也一樣,許多人知道自己所行的,不一定有真正屬靈的意義和價值；但一般人都是這樣做,你若不隨和,人覺得你很怪.以致我們的行為就與別人一樣,使人可以接納我們.神叫我們脫離罪惡的地方,無論生活思想上,要在這時代過分別為聖的生活,實在是一種的挑戰.
\newline
\newline
百姓的罪過是甚麼呢?(二1-2)「那些在床上圖謀罪孽造作奸惡的,天一發亮,因手有能力,就行出來了.他們貪圖田地就佔據,貪圖房屋便奪取；他們欺壓人,霸佔房屋和產業.」(二8-9)「我的民興起如仇敵,從那些安然經過不願打仗之人身上剝去外衣.你們將人民中的婦人,從安樂家中趕出；又將我的榮耀,從他們的小孩子盡行奪去.」他們貪圖產業房屋,簡單來說,那是貪心的罪,神為甚麼責備貪心的罪這麼厲害,因為聖經責備貪心的罪永是這麼厲害的,貪心是最易使我們走上錯誤的罪,所以十誡中用詞很精簡；但「不可貪心」卻提到有五次之多,可見貪心是多麼大的罪,保羅說貪心與拜偶像一樣,基督徒不再拜偶像,拜偶像等於否認主耶穌.當時百姓領袖有優越的地位,看見鄰居的田地便很垂涎；先知形容他們晚上在床上思想,怎樣得到那田地.在晚上沒有人看見你所作的,而思想沒有人知道,自己暗暗計劃,甚至輾轉反側睡不著；想到辦法後,天一亮後便去行動.把最沒有抵抗力的人趕出去,甚至用自己的權力把他趕出去.亞哈王想得到拿伯的葡萄園,但拿伯說這是祖宗的產業不能給你,亞哈心裡不舒服,躺在床上臉向牆,也不吃飯.我們有沒有看過小孩子得不到想要的東西時面上的表情呢?亞哈王就是這樣,結果耶洗別想辦法,將拿伯用石頭打死,奪取這葡萄園.不惜用甚麼方法,只求達到目的,這是貪心的表現.貪心的表現有甚麼意義呢?為甚麼神和先知責備貪心這麼厲害呢?我認為貪心有兩方面的意義；頭一方面,代表你思想的目標是甚麼.這些人在床上圖謀罪孽造作奸惡,思想怎樣多得房屋田地；結果被這種思想佔據,思想集中的地方會影響我們的行動.許多時候我們先用思想,想得太多,然後用行為表現出來.今天我們可能不是思想多得一些田地房屋,但我們卻思想多得一些物質的好處；把我們的思想放在物質享受上.香港是一個物質豐富的社會,生活享受比很多地方更豐富；但許多時候我們想得更多物質上的享受,越過越想多得.在商業上犯罪的都不是窮人,而是有錢人；他們已經有很多,但想得更多.希望自己的生活多點保障和依靠,將自己的思想集中在這些事情上.在教育兒女事情上也一樣,我們以為多給兒女物質享受更好；上星期我與一位姊妹談話,她家裡經濟很好,但家中有十多歲的孩子卻不聽話.這姊妹說:「他在家中想要甚麼我也供應他,電視遊戲機啦,電腦啦,我也買給他,他還想甚麼呢?」她以為物質可以滿足人的心,而實際上這孩子卻在家中反叛得很.許多時候我們也有這種思想,以為物質的好處,能滿足兒女的心,其實卻不然.
\newline
\newline
第二方面,貪心可代表人將整個盼望放在物質上,人好像得著物質,整個人的盼望就滿足了,這是多麼愚拙的事,(二3)「所以耶和華如此說,我籌劃災禍降與這族,這禍在你們的頸項上不能解脫,你也不能昂首而行,因為這時勢是惡的.」他們籌劃怎樣可以多得田地產業,但神也在籌劃怎樣刑罰他們的罪惡；而真正的籌劃和盼望,不在乎你有多少物質豐富,真正的盼望乃神在你身上怎樣計劃.今天我們若將盼望,保障放在物質上,但卻沒有神同在保障,神的災禍會臨到我們身上,無法逃脫.神仍然掌管我們的前途和計劃,如果我們的盼望放在物質上；我們的前途一點保障也沒有,真正的保障乃在乎依靠神.
\newline
\newline
百姓犯罪有兩方面的結果,第一方面百姓犯罪使眼睛蒙蔽,不能分辨是非真假,但神應許會引領和教導我們.(二7)「我耶和華的言語,豈不是與行動正直的人有益麼?」如果要主在我們生命中發生作用,我們就要從神的話中得到亮光；首先我們必須是向行動正直的人,正直可繙作正義,公義.不義的人不會分辨是非,如果我們不義便無法從神話語得到亮光.不著重正義,我們的眼睛會逐漸蒙蔽,看不出對或錯,百姓就是如此,貪圖房屋田地,以致眼睛蒙蔽.屬靈的眼光是需要慢慢鍛鍊出來的；我們若順服神的亮光,神會讓我們再進入更深的亮光中.我們若遵行神的心意,會對屬靈的亮光越來越清楚,分辨越光明.若果我們不遵行亮光的吩咐,那麼我們所得的亮光就會越來越少.在屬靈生活上,沒有一件事是獨立的；每一件事情都是代表我們屬靈生活的過程的一部份.我們遵行神的吩咐,便在屬靈上長進一步；下次看見有相似的亮光時,便更清楚明白,便容易遵行.如果今天我們得到亮光,卻沒有去遵行,下次就更加看不見亮光的意義來,今次我們沒有遵行這亮光,下次更加難去遵行了,我們需要不斷的訓練和遵行.另一方面提及正義,我們不能把舊約和新約分開,舊約新約合在一起才看得清楚.亞伯拉罕被神稱為義,因祂有相信的心,新約時保羅講得更加清楚,義人必因信得生.神做正義向人顯明,我們必須憑著信心來接受；在生活上有所得益,我們必須要有正義的生活.而正義的生活與信心是不能分離的；必須憑著信心和行動去接受.在屬靈事情上,有許多我們可用腦袋去分析；但在更深的事情上,腦袋卻領受不了.我們需要憑信心接受,憑著神的話語告訴我們；若我們單用環境或思想分析,或用科學方法分析；不錯可以分析得佷好,但卻不能看見裡面屬靈的意義.今天我們是太過依賴用理解分析,而忽略了憑信心接受聖經的教訓.比如主耶穌被釘十字架的信息,我們不可能用理智,歷史或科學方法來分析的,有些事情可以用這些方法來分析,好像你可以用科學方法,證實耶穌在二千多年前住過耶路撒冷,你也可用歷史方法,證明耶穌被羅馬人,猶太人釘在十字架上.你也可用科學方法,證明耶穌從死裡復活,但你不可用科學或歷史方法,證明耶穌為我釘死復活了.如果我們用信心接受,才可以看見耶穌的復活與受死,與我們的生命有關係.神言語是向正義的人顯明的,只有那些用信心接受神話語的人,可以得著.今天的世代,人非常重視知識和方法,著重四周環境和權利,我們所要追求的,乃是順服神的話語,耶和華的話只向那些行動正直公義的人,向那些有信心接受的人顯明,願神使我們有願意接受耶和華話語的心吧.
\newline
\newline
在(二12-13)神給我們另外的啟示,神提醒百姓,因為他們將思想,信心,集中在自己和物質上,以為有了滿足,就可解決一切問題了.但神說:「雅各家阿,我必要聚集你們,必要招聚以色列剩下的人,安置在一處,如波斯拉的羊,又如草場上的羊群.因為人數眾多,就必大大喧嘩,開路的在他們前面上去,他們直闖過城門；從城門出去,他們的王在前面行,耶和華引導他們.」你們眼目不要單顧自己,不要單想自己的滿足,因為神還有許多百姓,分散在其他地方.只有屬靈的人完全歸在一處,神的旨意才可成就,神的心意才得滿足.我們是否只關心自己教會和地方的人呢?是否單顧基督徒的工作呢?神所關心的,乃是要招聚所有的人,不是單是我們所看得見的.
\newline
\newline


\newpage

\section{第57屆~港九培靈會~鮑會園~第五講　假先知的路}
\label{sec:477}
\textbf{第五講　假先知的路}
\newline
\newline
連結: \href{https://www.hkbibleconference.org/session-message/view/477}{\texttt{https://www.hkbibleconference.org/session-message/view/477}}
\newline
\newline
\hyperref[sec:476]{< < < PREV SERMON < < <}
~
\hyperref[sec:toc]{[返目錄]}
~
\hyperref[sec:478]{> > > NEXT SERMON > > >}
\newline
\newline
彌迦書 3:1-3:12
\newline
\begin{longtable}{cl}
\hline
\hline
章節 & 經文 (和合本修訂版)\\
\hline
3:1 & \begin{tabularx}{0.7\textwidth}{X} 於是我說：雅各的領袖，以色列家的官長啊，你們要聽！你們豈不知道公平嗎？ \end{tabularx} \\ \\ \relax
3:2 & \begin{tabularx}{0.7\textwidth}{X} 你們惡善好惡，剝我百姓身上的皮，從他們的骨頭上剔肉， \end{tabularx} \\ \\ \relax
3:3 & \begin{tabularx}{0.7\textwidth}{X} 你們吃我百姓的肉，剝他們的皮，打斷他們的骨頭，如切塊下鍋，如釜中的肉。 \end{tabularx} \\ \\ \relax
3:4 & \begin{tabularx}{0.7\textwidth}{X} 到了遭災的時候，這些人要哀求耶和華，他卻不應允他們。那時，因他們所行的惡，他必轉臉離開他們。 \end{tabularx} \\ \\ \relax
3:5 & \begin{tabularx}{0.7\textwidth}{X} 論到使我百姓走入歧途的先知，他們牙齒有所嚼，就呼喊說：「平安！」誰不給他們吃，就揚言攻擊他，耶和華如此說： \end{tabularx} \\ \\ \relax
3:6 & \begin{tabularx}{0.7\textwidth}{X} 你們因此必遭遇黑夜，看不到異象；遭遇幽暗，無法占卜。太陽必向先知沉落，白晝轉為黑暗。 \end{tabularx} \\ \\ \relax
3:7 & \begin{tabularx}{0.7\textwidth}{X} 先見必抱愧，占卜的必蒙羞，他們全都摀著鬍鬚，因為神不應允他們。 \end{tabularx} \\ \\ \relax
3:8 & \begin{tabularx}{0.7\textwidth}{X} 至於我，我藉耶和華的靈，滿有能力、公平和勇氣，可以向雅各述說他的過犯，向以色列指出他的罪惡。 \end{tabularx} \\ \\ \relax
3:9 & \begin{tabularx}{0.7\textwidth}{X} 當聽這話，雅各家的領袖，以色列家的官長啊！你們厭棄公平，在一切事上屈枉正直； \end{tabularx} \\ \\ \relax
3:10 & \begin{tabularx}{0.7\textwidth}{X} 以血建立錫安，以罪孽建造耶路撒冷。 \end{tabularx} \\ \\ \relax
3:11 & \begin{tabularx}{0.7\textwidth}{X} 城裡的領袖為賄賂行審判，祭司為酬勞施訓誨，先知為銀錢行占卜；他們卻倚賴耶和華，說：「耶和華不是在我們中間嗎？災禍必不臨到我們。」 \end{tabularx} \\ \\ \relax
3:12 & \begin{tabularx}{0.7\textwidth}{X} 因此，為你們的緣故，錫安要被耕種像一塊田地，耶路撒冷要變為廢墟，這殿的山必如叢林的高處。 \end{tabularx} \\ \\
[1ex]
\hline
\hline
\end{longtable}
先知彌迦指責那些自稱先知而不是先知的人,這裡提及雅各的首領以色列家的官長；這些人站在領導上,應該帶領百姓走在神的道路,但卻領百姓走錯了路.舊約的先知不是等於今天的傳道人,雖然傳道人的工作與先知一樣,傳講神的信息；先知在特殊情況下,有特別的領受.為了那情況,神給他信息,他便應傳講這信息.照樣今天神會為屬祂的人有特殊使命,無論我們在甚麼崗位,也要忠心去完成.因為在舊約時候,並不是每個先知都是全職作先知的,他有自己的職位,但神有託負的時候便去完成神的工作.所以這可應用在每一個屬主的人身上,而不是單指傳道人.在評判先知工作前,彌迦先指出評判的標準.(三1)「你們不當知道公平嗎?」公正應作正直的標準,比如我們平時說純正的信仰,純正可以繙作這個字.神對先知說,一個人要完成神的工作,他首先要知道正確的目標和方向.這卷書有幾次提及這件事,神的話和啟示,就是標準；先知彌迦領受神的信息,照樣今天我們領受神的信息,我們應照著神標準去傳講和工作.沒有別的事可以決定我們怎樣行,單單是神的話.以色列官長不能因站在領導上,便有權決定講甚麼；不能因百姓需要的角度講甚麼,也不是百姓的好處來決定講甚麼.照樣今天沒有別的決定,我們所走的方向；不是藉著國家民族的環境和需要,不是藉著社會的風氣,乃是單單信靠神的話語,這才是純正和公平的標準.評判的第二方面標準是「惡善好惡」,你所行出來的事要符合你的標準.有些人有很好聽的目標,但行出來卻不一樣；所講的理想很高,但作出來與理想卻不符合.神提醒百姓,最終的評判標準乃是自己的行為.主耶穌說要以人的行為定善惡,按照他所結的果子可以看見；不論這人的理想是如何大或高,也不能改變神的標準.
\newline
\newline
甚麼是假先知的表現和記號呢?彌迦講出幾件事件:
\newline
\newline
1. 假先知所事奉的是他自己.他們為了自己的好處而工作,無論事奉的是甚麼,他們事奉的目標乃是自己.(三11)「首領為賄賂行審判,祭司為雇價施訓誨,先知為銀錢行占卜,他們卻依賴耶和華.」他們說依賴耶和華,其實只顧自己的利益.(二5)「論到使我民走差路的先知,他們牙齒有所嚼的,他們就呼喊說平安了；凡不供給他們吃的,他們就豫備攻擊他.」你給我有好處,我便傳講平安的信息；為了賄賂行審判,為了雇價施訓誨.一個人為自己好處工作,他評判的標準是這工價能帶給他多少好處；對我有好處,便越多努力工作,對我沒有好處便不去工作.我們的工作不一定是為得金錢,而是為別的好處,多去工作為多得別人支持和歡迎,或多得物質的好處,如果我們以多得物質的好處工作,可以想像工作的方向和取捨時,會受這些事情影響.如果我工作目標是叫人歡迎,便多講叫人喜歡的信息；在很多聚會我們看見有這種情形,傳講信息用自己的利害和生活決定.當然我們傳講信息不是不應得到本身的好處,使徒保羅說:「工人得工價是應當的.」主耶穌也有同樣的教訓,工作的有應得的報酬.但如果我們用工作報酬多少,來決定工作的方向時,我們便不能忠心傳遞神的信息.假先知是按自己報酬和好處,來決定工作的方向.
\newline
\newline
2. 使百姓走差了路.(三5)「論到使我民走差路的先知.」走差路很難解釋,原文好像醉酒一樣,一個醉酒的人會失去評判的標準,看事情不清楚,走路不正直了,失掉辨別是非的能力.一個人有目的的時候,會想法子叫人受他的影響和控制.他會控制人的方向,叫這人失掉分辨是非的能力,有很多異端興起來.好像有一韓國人文鮮明,在美國創立統一教,他自己英文也說得不好,但卻使許多美國人信服他,吸引許多人歸信他.這些人入教後,便完全失去了理智,以為過往一切都錯了,只有文鮮明的教導才對.所以要離開一切,不再承認家庭的關係,不再稱父母為父母；甚至有人離開丈夫和妻子,整個心理被這個宗教操縱,失去理智.還有一教派叫做「神的兒女」,雖然人數不太多,卻引起很大的害處；在香港這教派很活躍,雖然這幾年間多次揭發他們的錯誤,但奇怪仍然有多人跟隨.他們說用愛心作標準,要遵守耶穌的教訓要彼此相愛,要用全心甚至身體愛對方,將整個身體給對方.因此這教派男女關係非常混亂,已經有許多年青女孩子吸引進去.人失去理智,就無法分辨是非了；失掉道德標準,把羞恥當作榮耀.若要把這些人從他們中間帶領出來就要用許多工夫了.這些教派好像使用洗腦方式一樣,改變你的思想,叫你整個人被他麻醉.他們的方法非常厲害,有許多人向你傳遞錯誤的道理,結果你心裡自然會懷疑自己評判的標準,漸漸接納他們的錯誤的道理.那些道理是屬靈的,或是屬道德和政治的,那影響就大了,漸漸這些人麻醉自己.神所責備的先知,就是用錯誤道理領人走差路和麻醉別人.無論有多少人講哪一方面的真理,如果他所講的不合乎聖經的真理,那便錯了.聖經從沒有教我們用這方法領人信主,帶人信主是清清楚楚的,聖經叫我們要分辨真理；我們是照神的形像做的,神的形像是一個本能,就是可以分別對錯.這是人和其他生物的不同,人能夠分辨黑白；故此我們要運用理智分別哪一個是對,哪一個是錯.先知常說你們這樣作合理嗎?你們豈不知甚麼是公平嗎?假先知不單傳假教訓,叫人走差路,他連自己也迷惑,這是相當可憐的.假先知不一定時常欺騙人,乃是他自己心智昏迷,不能分辨是非就是了.很多時候我們也這樣麻醉自己,明明知道自己錯了也去作,做完之後便為自己找理由作藉口,多次是這樣講,有一次有位弟兄給我看一隻湯匙,這是他吃飯時從一飯店偷取的,明知這是不對；但他有理由,因當天侍應生的服務態度很差.因為是侍應所給的服務很差,所以這弟兄要報復,便帶走那湯匙.他覺得明正言順地向人炫耀；況且他不止一次是這樣,他覺得這是應當的,其實他自己已到不能分辨是非的地步了.(三11)「首領為賄賂行審判,祭司為雇價施訓誨,先知為銀錢行占卜,他們卻依賴耶和華說,耶和華不是在我們中間麼,災禍必不臨到我們.」那些假先知竟然說自己是先知,還傳講摩西的律法；其實他們錯了,因為耶和華根本不在他們中間,但他們仍說耶和華仍在我們中間.照樣,我們奉主的名建造一所大教堂和神學院,奉主的名叫許多人出去傳道,工作良好；其實所作的不合乎真理,無論我們用得多好,但卻是麻醉自己,走上錯誤的道路.外表雖然多麼興旺和有表現,但耶和華不在我們中間,這些人可憐到連自己也不知道道路和錯誤.
\newline
\newline
假先知第三個記號,就是單單傳講百姓喜歡聽的信息,經文中幾次提及這事情；講百姓聽喜歡的話,使自己被接納.但所講的不一定合乎真理,我們又是否陷入這圈套呢?以為被人接納歡迎就好呢?韓國一教會有三十萬會友,這教會多麼好呀!否則那會有這麼多會友,講道時有幾十萬人說阿利路亞,這講道一定好.但其實若所講的不合乎真理,照樣也有成千上萬的人歡迎他.異端的教訓,照樣也有成千上萬的人歡迎.有人甚至不說是神的話照樣也有人歡迎；有很多政治集會不講及神的話語,照樣亦有成千上萬的人歡迎.先知為了要吸引群眾,說快快樂樂的話語,這正是假先知的記號.
\newline
\newline
現在我們看一個真正的先知應有的記號,有幾件事值得我們留意的.(三8)「至於我,我藉耶和華的靈,滿有力量公平才能,可以向雅各說明他的過犯；向以色列指出他的罪惡.」他所傳的信息,雅各家以色列家聽不聽是另一件事；這裡的「滿有」是很重要的字,原來是很重要的動詞.彌迦這裡說:「我被公平,才能所充滿.」我現在是代表耶和華傳講信息,我首先需要神的能力充滿我；如果要神的力量充滿,就證明我工作的能力不是自己的.力量應繙作權柄.(約一12)「凡接待祂的,就是信祂名的人,祂就賜他們權柄作神的兒女.」(在這節聖經是繙作權柄兩個字.因為神賜下權柄給我們,所以我知道所說的話是從神來的,我知道我有神的同在,因為我有從神而來的權柄；有從神來的權柄不一定常常受歡迎和被接納的,但我知道所講的是對的,即使人反對不接納,也不要緊.
\newline
\newline
另一方面,他不單只充滿神的權柄,而且神讓他充滿公平.公平就是正義,正義的標準乃是神的話,神把權柄給了我；因為我有公平正義的標準,我才能傳講神的話語.不單有權柄公平,還有才能.「才能」繙得不大恰當,原文是作「勇氣」,我有從神來的權柄和標準,我還有從神來的勇氣,真正傳講神的話要有勇氣.彌迦並不是很出名的人,他竟敢公開責備以色列家的首領和長官；一個寂寂無名的人,責備領袖,需要有勇氣.
\newline
\newline
一個真正作神先知的人,他必須單單傳講神要他傳的話,因為他有這記號,他可以向雅各家說明他的過犯,向以色列家指出它的罪惡.神要他指出百姓的罪來,他便得去傳講這信息,他不能因傳這信息令百姓失望難受而不傳.他管不著這些,只單單傳講神要他講的話.如果是安慰的信息就傳講安慰的信息,如果是責備的信息,就傳責備的信息.許多時候我們要傳講福音,太容易加上自己的思想；要把福音應用在社會上,我不是說福音不可以應用在社會上；但這應用離開了神的話語,便是錯了.
\newline
\newline
先知的第三個記號是為欺壓的人說話,這裡受欺壓的人不是想像中的人,乃是真正受官長欺壓的人.那些長官從百姓身上剝皮削肉,他們被欺壓到不能生活.百姓的領袖以人血建城(三10),把百姓拉來建城；若工作不滿意,便鞭打百姓流血,百姓領袖貪圖田地房屋,但沒有人替他們講出謀公平的話.在美國,英國,他們認為工人階級會被剝削,於是有許多人自告奮勇說公平的話.其實這些工人遠超他們所得的好處,工人的福利相當大,甚至叫資本家走上破產的道路.去年英國煤礦工人大罷工,使許多國家受大損失!其實他們要求的條件不合理,政府要關掉不能出產煤的礦洞,但工會說不可以,因為關閉礦洞會使幾位工人失業；明明不合理,在這情況替工人說話是不合理的.今天有許多人不理會事實,只看階層而說話.但先知不是如此,他只為真正受欺壓的人講說話,站在真理立場上為受逼迫的人講話.是的很少人敢為真正受欺壓的人說話,我們實在需要有先知的勇氣.假先知失掉神真正的喜悅,(三4)「這些人必哀求耶和華,祂卻不應允他們.」神離開他們,不再聽他們,雖裝作敬虔的外表,也沒有用處,因為神不再在他們中間.最後(三12)「所以因你們的緣故,錫安必被耕種像一塊田,耶路撒冷必變為亂堆,這殿的山必像叢林的高處.」神要施行審判的臨到百姓,假先知和真先知的道路擺在我們面前,在這混亂的時代,求神讓我們認定所走的方向吧!
\newline
\newline


\newpage

\section{第57屆~港九培靈會~鮑會園~第六講　榮耀的盼望}
\label{sec:478}
\textbf{第六講　榮耀的盼望}
\newline
\newline
連結: \href{https://www.hkbibleconference.org/session-message/view/478}{\texttt{https://www.hkbibleconference.org/session-message/view/478}}
\newline
\newline
\hyperref[sec:477]{< < < PREV SERMON < < <}
~
\hyperref[sec:toc]{[返目錄]}
~
\hyperref[sec:479]{> > > NEXT SERMON > > >}
\newline
\newline
彌迦書 4:1-5:4
\newline
\begin{longtable}{cl}
\hline
\hline
章節 & 經文 (和合本修訂版)\\
\hline
4:1 & \begin{tabularx}{0.7\textwidth}{X} 末後的日子，耶和華殿的山必堅立，超乎諸山，高舉過於萬嶺；萬民都要流歸這山。 \end{tabularx} \\ \\ \relax
4:2 & \begin{tabularx}{0.7\textwidth}{X} 必有許多民族前往，說：「來吧，我們登耶和華的山，到雅各神的殿。他必將他的道指教我們，我們也要行他的路。」因為教誨必出於錫安，耶和華的言語必出於耶路撒冷。 \end{tabularx} \\ \\ \relax
4:3 & \begin{tabularx}{0.7\textwidth}{X} 他必在許多民族中施行審判，為遠方強盛的國斷定是非。他們要將刀打成犁頭，把槍打成鐮刀。這國不舉刀攻擊那國，他們也不再學習戰事。 \end{tabularx} \\ \\ \relax
4:4 & \begin{tabularx}{0.7\textwidth}{X} 人人都要坐在自己的葡萄樹和無花果樹下，無人使他們驚嚇；這是萬軍之耶和華親口說的。 \end{tabularx} \\ \\ \relax
4:5 & \begin{tabularx}{0.7\textwidth}{X} 萬民都奉自己神明的名行事，我們卻要奉耶和華－我們神的名而行，直到永永遠遠。 \end{tabularx} \\ \\ \relax
4:6 & \begin{tabularx}{0.7\textwidth}{X} 耶和華說：在那日，我必聚集瘸腿的，召集被趕逐的，以及我所懲治的人。 \end{tabularx} \\ \\ \relax
4:7 & \begin{tabularx}{0.7\textwidth}{X} 我要使瘸腿的成為餘民，使被趕到遠方的成為強盛之國。耶和華要在錫安山作王治理他們，從今直到永遠。 \end{tabularx} \\ \\ \relax
4:8 & \begin{tabularx}{0.7\textwidth}{X} 你，以得臺，錫安的山岡啊，先前的權柄必歸給你，耶路撒冷的國權必將歸還。 \end{tabularx} \\ \\ \relax
4:9 & \begin{tabularx}{0.7\textwidth}{X} 現在，你為何大聲呼喊呢？你中間沒有君王，你的謀士滅絕，以致疼痛抓住你，如臨產的婦人嗎？ \end{tabularx} \\ \\ \relax
4:10 & \begin{tabularx}{0.7\textwidth}{X} 錫安哪，你要疼痛生產，彷彿臨產的婦人；因為你必從城裡出來，住在田野；你要到巴比倫去，在那裡，你要蒙解救，在那裡，耶和華必救贖你脫離仇敵的手掌。 \end{tabularx} \\ \\ \relax
4:11 & \begin{tabularx}{0.7\textwidth}{X} 現在，許多國家聚集攻擊你，說：「讓錫安被玷污！讓我們親眼看到！」 \end{tabularx} \\ \\ \relax
4:12 & \begin{tabularx}{0.7\textwidth}{X} 他們卻不知道耶和華的意念，也不明白他的籌算，他聚集他們，像把禾捆聚到禾場。 \end{tabularx} \\ \\ \relax
4:13 & \begin{tabularx}{0.7\textwidth}{X} 錫安哪，起來踹穀吧！我必使你的角成為鐵，使你的蹄成為銅。你必打碎許多民族，將他們的財寶獻給耶和華，將他們的財富獻給全地的主。 \end{tabularx} \\ \\ \relax
5:1 & \begin{tabularx}{0.7\textwidth}{X} 成群的民哪，現在要聚集成隊；仇敵前來圍攻我們，要用杖擊打以色列領袖的臉頰。 \end{tabularx} \\ \\ \relax
5:2 & \begin{tabularx}{0.7\textwidth}{X} 伯利恆的以法他啊，你在猶大諸城中雖小，將來必有一位從你那裡出來，在以色列中為我作掌權者；他的根源自亙古，從太初就有。 \end{tabularx} \\ \\ \relax
5:3 & \begin{tabularx}{0.7\textwidth}{X} 因此，耶和華要將以色列人交給敵人，直到臨產的婦人生下孩子；那時，他其餘的弟兄必回到以色列人那裡。 \end{tabularx} \\ \\ \relax
5:4 & \begin{tabularx}{0.7\textwidth}{X} 他必倚靠耶和華的大能，倚靠耶和華－他神之名的威嚴，站立並牧養，使他們安然居住；因為現在他必尊大，直到地極。 \end{tabularx} \\ \\
[1ex]
\hline
\hline
\end{longtable}


\newpage

\section{第57屆~港九培靈會~鮑會園~第七講　王的被拒絕}
\label{sec:479}
\textbf{第七講　王的被拒絕}
\newline
\newline
連結: \href{https://www.hkbibleconference.org/session-message/view/479}{\texttt{https://www.hkbibleconference.org/session-message/view/479}}
\newline
\newline
\hyperref[sec:478]{< < < PREV SERMON < < <}
~
\hyperref[sec:toc]{[返目錄]}
~
\hyperref[sec:480]{> > > NEXT SERMON > > >}
\newline
\newline
彌迦書 5:5-6:5
\newline
\begin{longtable}{cl}
\hline
\hline
章節 & 經文 (和合本修訂版)\\
\hline
5:5 & \begin{tabularx}{0.7\textwidth}{X} 這位就是和平。當亞述侵入我們領土，踐踏我們宮殿時，我們就立七個牧者，八個領袖攻擊它。 \end{tabularx} \\ \\ \relax
5:6 & \begin{tabularx}{0.7\textwidth}{X} 他們要用刀劍毀壞亞述地和寧錄地的關口。當亞述侵入我們領土，踐踏我們邊境時，他必拯救我們。 \end{tabularx} \\ \\ \relax
5:7 & \begin{tabularx}{0.7\textwidth}{X} 雅各的餘民必在許多民族中，如從耶和華降下的露水，又如甘霖降在草上；他們不倚靠人，也不仰賴世人。 \end{tabularx} \\ \\ \relax
5:8 & \begin{tabularx}{0.7\textwidth}{X} 雅各的餘民必在列國中，在許多民族中，如林間百獸中的獅子，又如少壯獅子在羊群中；他若經過就必踐踏撕裂，無人搭救。 \end{tabularx} \\ \\ \relax
5:9 & \begin{tabularx}{0.7\textwidth}{X} 願你的手舉起，高過敵人！願你的仇敵都被剪除！ \end{tabularx} \\ \\ \relax
5:10 & \begin{tabularx}{0.7\textwidth}{X} 耶和華說：到那日，我必從你中間剪除馬匹，毀壞戰車； \end{tabularx} \\ \\ \relax
5:11 & \begin{tabularx}{0.7\textwidth}{X} 除滅你國中的城鎮，拆毀你一切的堡壘； \end{tabularx} \\ \\ \relax
5:12 & \begin{tabularx}{0.7\textwidth}{X} 除掉你手中的邪術，你那裡也不再有占卜的人。 \end{tabularx} \\ \\ \relax
5:13 & \begin{tabularx}{0.7\textwidth}{X} 我必從你中間除滅雕刻的偶像和柱像，你就不再跪拜自己手所造的； \end{tabularx} \\ \\ \relax
5:14 & \begin{tabularx}{0.7\textwidth}{X} 我必從你中間拔除亞舍拉，毀滅你的城鎮； \end{tabularx} \\ \\ \relax
5:15 & \begin{tabularx}{0.7\textwidth}{X} 我必在怒氣和憤怒中報應那不聽從我的列國。 \end{tabularx} \\ \\ \relax
6:1 & \begin{tabularx}{0.7\textwidth}{X} 當聽耶和華說的話：起來，向山嶺爭辯，使岡陵聽見你的聲音。 \end{tabularx} \\ \\ \relax
6:2 & \begin{tabularx}{0.7\textwidth}{X} 山嶺啊，要聽耶和華的指控！大地永久的根基啊，要聽！因耶和華控告他的百姓，與以色列爭辯。 \end{tabularx} \\ \\ \relax
6:3 & \begin{tabularx}{0.7\textwidth}{X} 「我的百姓啊，我向你做了甚麼呢？我在甚麼事上使你厭煩？你回答我吧！ \end{tabularx} \\ \\ \relax
6:4 & \begin{tabularx}{0.7\textwidth}{X} 我曾將你從埃及地領出來，從為奴之家救贖你，我差遣摩西、亞倫和米利暗在你前面帶領。 \end{tabularx} \\ \\ \relax
6:5 & \begin{tabularx}{0.7\textwidth}{X} 我的百姓啊，當記念從前摩押王巴勒如何籌算，比珥的兒子巴蘭如何回應他，當記念從什亭到吉甲所發生的事，好使你們明白耶和華公義的作為。」 \end{tabularx} \\ \\
[1ex]
\hline
\hline
\end{longtable}
在第五章開始神應許百姓,賜下一位拯救主；這件事是未來的,但在神眼中便好像現在的事.在祂眼中沒有過去和未來,因為一切皆在面前顯明了.彌賽亞按照神的應許降下來了,但百姓卻沒有接納祂.下面先知告訴我們百姓不接納祂的原因,然後告訴我們不接受的結果.第六章告訴我們,百姓的態度令神作出的反應.神說:我要與你們爭辯,這句話原是被控訴時,為自己辯護的.百姓在這裡控訴耶和華,說耶和華對我們不夠愛心,沒有按照應許的,保守百姓!你沒有向我們施恩.當百姓這樣控訴耶和華的時候,神便向百姓爭辯,神說,我在甚麼事上使你厭煩,我做錯甚麼事你要控訴我呢?人只會經驗到神豐富的恩典,又得到神的恩典是應當的；特別以色列人以為自己是神的選民,神必須要應許賜福給我們.神既是信實的,如果應許了卻不賜福,祂就不是信實的神了.不管我們好不好或有否得罪神,不管我們有否遵行神的話,我們犯甚麼錯不要緊,但神必須要守著祂信實的應許.因此他們明明拒絕了彌賽亞,但覺得神仍須要向他們施恩.這不是聖經的教訓,從摩西時代開始,申命記強調神在恩典中揀選他們,使他們在萬民中作神的見證.但他們必須要遵行神的話,他們這樣遵行神的話；就有豐富的恩典賜給他們,他們若違背神的話,神就懲罰他們.百姓要進入迦南地時,神吩咐百姓在進入迦南地後,要揀選兩批百姓；一批百姓站在基利心山,宣佈神賜福的信息,他們遵守我的話,我就會這樣賜福我的百姓；另一批百姓站在以巴路山宣佈咒詛的話,如果不遵行我的命令,我就會使咒詛到你們.這是神從起初講的,百姓應當知道神的旨意,神沒有說無論你們怎樣行,我都向你施恩；只有順服我的話,我才向你施恩.如果你們不順服,我的懲罰就臨到你們,故此神這樣懲罰祂的百姓,並沒有違背祂的應許；因百姓拒絕神預備的彌賽亞,所以百姓要受懲罰.百姓以為神向他們這樣,就要施恩；神說你把你的理由講出來,我要與百姓與以色列爭辯.神向百姓施了這麼多恩典,百姓卻轉過來拒絕神預備的彌賽亞,多麼愚拙的事,神向百姓說:「我曾將你從埃及地領出來,從作奴僕之家救贖你,我也差遣摩西亞倫和米利暗在你前面行.你們當追念摩押王巴勒所設的謀,和比珥的兒子巴蘭回答他的話；並你們從什亭到吉甲所遇見的事,好使你們知道耶和華的作為.」(六4-5)你們要看看我為你們所施的一切恩典.現在你們竟忘記了,又去走錯誤的道路；已經有這麼多證據擺在你們前面,你們仍然拒絕我的恩典.我向你們施行的審判,是否應當的呢?我們時常犯這個毛病,神今天向我們施恩,但我們卻悖逆神而受苦,便忘記神所有的恩典.如果我們在受苦時承認這是神的懲罰罪惡,祂不會縱容我們,我們也許聽過一種錯誤的教訓「普救論」,意思是普世的人都得救,因為主耶穌是慈愛的主,祂不會向人施行刑罰,在神計劃中根本沒有地獄；人離開世界後,不會到痛苦的地方.因此神會向人施行恩典,無論人犯罪或不信,因人不是靠自己的行為得救,將來神一定會救他.(六5)指出神所作的要叫人知道祂公義的作為；神不單慈愛的也是公義的.人犯罪不能不受刑罰,神所作的一切,叫我們知道祂的公義.祂有公義的標準,祂的審判可能不在今天使行；也許不在明天使行,但祂一定會按公義的標準審判.
\newline
\newline
(一)拒絕神預備的彌賽亞(五5-6)
\newline
\newline
「當亞述人進入我們的地境踐踏宮殿的時候,我們就立起七個牧者,八個首領攻擊他,他們必用刀劍毀壞亞述地和寧錄的關口.亞述人進入我們的地境踐踏的時候,他必拯救我們.」百姓以為自己有很大的能力,這裡反映北國耶羅波安二世的態度,覺得自己軍事勢力很強,可以敵擋一切的敵人,靠自己力量抵擋亞述巴比倫人,他們完全忘記神的警告,沒有一點依賴耶和華的心.(五8)「雅各餘剩的人必在多國民中,如林間百獸中的獅子；又如少壯獅子在羊群中,他若經過,就必踐踏撕裂,無人搭救.」他們以為自己很有力量,像獅子在樹木中一樣,無人能抵抗,以色列百姓以為憑自己力量,可以勝過一切.今天人以為有力量去完成目標,我憑著雙手可以創造前途,不單是不信的人如此,甚至那些自以為相信而不是信神的人也如此,以為憑自己力量去創造人間樂園；以為憑著彼此相愛,關懷社會的努力便可以成全.我們不是說基督徒不應關懷社會,但如單靠自己力量就是拒絕神.
\newline
\newline
(二)否定神在暗中的作為(五7)
\newline
\newline
雅各餘剩的人必在多國的民中,如從耶和華那裡降下的露水；又如甘霖降在草上,不仗賴人力,也不等候世人文功.」露水是很特別的東西,是古時候不能明白的力量,大雨降下你可看見,但露水看不見,只有在平靜的環境,露水才降在草上,雖然人看不見露水出現的情況,但人可清楚看見露水的果效.有些地方可能春天,夏天沒有雨水,在巴勒斯坦近海的地方就是不同,因為他們的葡萄園收成很好,完全靠露水.神在人身上施恩典的情形,可能用人眼睛看不到神工作的情況；但你用敬虔的眼睛觀察,便可看見神工作的果效.主耶穌說天國好像一點麵酵,放在麵裡......你看不見麵酵怎樣工作,但是你會看見麵發起來了.神的工作也是如此,有時叫人不能馬上看見工作的果效；但神恩典的工作,就好像天國一樣.不是很轟轟烈烈的工作,那些轟轟烈烈的工作一定很有果效；神的工作在人心裡,好像天國在人心裡一樣.在教會歷史上我們多次看見這樣情形,教會在壓力底下,工作可能有轟轟烈烈的時候,但神在人心裡發生作用；人看不見工作的時候,教會人數加多,信主的人加多,屬靈的影響力加強了,神在那裡工作.古時的人不單不明白露水的工作,而且他們把這些看作神直接工作的表現.有一本猶太人解經書,他們說露水是神特別在人間工作的恩典,他們不明白露水怎樣來,以為這是神直接的作為.外表很像沒有甚麼偉大工作,但神卻在裡面暗暗工作；不過我們要有屬靈的眼光才可看見.以色列的百姓看見露水的工作,但卻未有承認這是神的作為,他們看見百姓中復興的工作,看見百姓在逼迫中神怎樣搭救他們；不過他們看見神的恩典,卻否認神的恩典.今天有很多人看見神的作為,卻當作是人的工作和果效.
\newline
\newline
(三)依靠外邦的偶像假神百姓所犯的罪越發加深,依靠偶像不接受彌賽亞的工作(五12)
\newline
\newline
「我必從你中間除滅雕刻的偶象和柱像,你就不再跪拜自己手所造的.」這些邪術,占卜和偶像卻是你們手所造的,當時以色列百姓到了亞述國的首都,看見大馬色的偶像,也帶回去敬拜；他們說是這些神拯救我,因此要敬拜外邦神.亞哈拜巴力,耶羅波安二世敬拜大馬色的神,以為偶像能幫助他們,表示他們厭棄耶和華,覺得甚麼神和偶像都比耶和華好.今天人都有這心理,他們可以用任何理由證明聖經不可靠；不論有甚麼理由,人立刻接受這理論,而不接受聖經的理論.其實聖經的教訓,多次證明聖經是神的話是可靠的.人以為可以用其他理由來代替聖經的權威,人很容易接受自己造出來的理論和偶像,而不肯真正謙卑在神的面前.
\newline
\newline
百姓拒絕神的救恩和恩典,結果怎樣?(五15)「我也必在怒氣和忿怒中向那不聽從的列國施報,.」神講出祂的作為,這是祂作為中奇妙事情；這裡看見神的計劃和人經驗之間的關係,看見神在人中間作為的意義.有幾件事值得我們留意的.
\newline
\newline
1. 就是先知說人間的混亂情況和痛苦事情,不能與神的工作脫節.今天人不願接受這件事實,不願接受這是神的作為,人若拒絕這位和平的君,地上不會有真正的和平,人若拒絕這位掌管天地萬物的主,不可能有真正的次序,地上一定會混亂.人若果拒絕一位掌權者在天上,地上不可能有尊重權柄的現象.舊約時,百姓因拒絕神的權柄,便落在混亂的情況中.在士師記看見百姓遇見多次的混亂,弱小的國家竟然控制神的百姓,百姓為甚麼這樣軟弱?士師記說因為以色列中沒有王,百姓不肯尊耶和華為王任意而行,所以便有這些事情臨到.新約時也講及世人拒絕耶穌基督,地上會有混亂.在歷史上多次看見這種情形,當世人公開拒絕神的時候,和拒絕神權威的時候,這世界一定會混亂.主後第七,八世紀,歐洲最混亂時,整個社會包括教會,就是拒絕神的權威的時候；如果人拒絕這位和平之君,地上不可能有太平.
\newline
\newline
2.地上所以變得這麼混亂,因為人不肯聽從耶和華的話.神向那不聽祂的話的人要施行審判,世界的混亂完全控制在神手裡；神給人呼召,應該悔改,遵神而行.人如果不肯聽從的時候,神會允許混亂臨到這個世界,世人聽從與否,也決定神會否施恩給世界.
\newline
\newline
3. 神工作的原則.神給人的呼召,不是按照人的領受而定.對以色列百姓是一種呼召,對不聽從的列國,又是另一種的呼召.對敬畏神的人是一種呼召,對不認識神的人又是另一種呼召.神沒有叫那些不認識真理的人去遵行特殊的真理,神沒有叫那些不認識聖經真理的人,守某些具體的聖經教訓.不信的人不遵守主日不是問題,不信主的人不奉獻十分之一也不是問題.但神仍給他呼召,要怎樣過公義的生活,怎樣尊敬創造天地的主,如果他不遵行,社會刑罰他.但對屬神的百姓,神給他們另外的呼召,神呼召他們離棄偶像,占卜等.照樣神今天給我們與給百姓的呼召一樣,我們應按神啟示的真理過生活.神沒有叫我們對不明白的事負責任,但是我們知道明白的真理,我們就要負責任.神給我們的呼召是按我們所明白的,故此我們在神面前領受責任越多,我們的責任就越大.今天在這個混亂的時代,我們若明白神的真理,知道甚麼是敬虔正義的生活；但卻沒有按真理行,我們便會按我們所明白的負責任.主耶穌給我們的信息也一樣,多給誰就向誰多取.我們聽到神話語之後,神把祂的責任放在我們身上,我們有責任按照我們所聽到的去行.聽道不是單為滿足好奇心或某方面的需要,或單單在知識上得到滿足；我們便多得神的責任.在今天混亂的時代,我們屬神的人知道明白聖經的真理,有責任將從神所領受的,在生活中活出來；叫人知道神在人間的希望和要求,求神幫助我們盡上自己的本份.
\newline
\newline


\newpage

\section{第57屆~港九培靈會~鮑會園~第八講　清楚的準則}
\label{sec:480}
\textbf{第八講　清楚的準則}
\newline
\newline
連結: \href{https://www.hkbibleconference.org/session-message/view/480}{\texttt{https://www.hkbibleconference.org/session-message/view/480}}
\newline
\newline
\hyperref[sec:479]{< < < PREV SERMON < < <}
~
\hyperref[sec:toc]{[返目錄]}
~
\hyperref[sec:481]{> > > NEXT SERMON > > >}
\newline
\newline
彌迦書 6:6-6:16
\newline
\begin{longtable}{cl}
\hline
\hline
章節 & 經文 (和合本修訂版)\\
\hline
6:6 & \begin{tabularx}{0.7\textwidth}{X} 「我朝見耶和華，在至高神面前跪拜，當獻上甚麼呢？難道獻一歲的牛犢為燔祭來朝見他嗎？ \end{tabularx} \\ \\ \relax
6:7 & \begin{tabularx}{0.7\textwidth}{X} 耶和華豈喜悅千千的公羊，或是萬萬的油河嗎？我豈可為自己的過犯獻我的長子，為自己的罪惡獻我所親生的嗎？」 \end{tabularx} \\ \\ \relax
6:8 & \begin{tabularx}{0.7\textwidth}{X} 世人哪，耶和華已指示你何為善。他向你所要的是甚麼呢？只要你行公義，好憐憫，存謙卑的心與你的神同行。 \end{tabularx} \\ \\ \relax
6:9 & \begin{tabularx}{0.7\textwidth}{X} 耶和華向這城呼叫---看重你的名是真智慧---你們當聽懲罰和派定懲罰的人。 \end{tabularx} \\ \\ \relax
6:10 & \begin{tabularx}{0.7\textwidth}{X} 惡人家中不是仍有不義之財和惹人生氣的變小了的伊法嗎？ \end{tabularx} \\ \\ \relax
6:11 & \begin{tabularx}{0.7\textwidth}{X} 我若用不公道的天平和袋中詭詐的法碼，豈可算為清白呢？ \end{tabularx} \\ \\ \relax
6:12 & \begin{tabularx}{0.7\textwidth}{X} 城裡的有錢人遍行殘暴，其中的居民說謊話，口中的舌頭盡是詭詐。 \end{tabularx} \\ \\ \relax
6:13 & \begin{tabularx}{0.7\textwidth}{X} 因此，我也擊打你，使你受傷，因你的罪惡使你受驚駭。 \end{tabularx} \\ \\ \relax
6:14 & \begin{tabularx}{0.7\textwidth}{X} 你要吃，卻吃不飽，你的肚子仍是空空。你必被挪去，不得逃脫；如有逃脫的，我必交給刀劍。 \end{tabularx} \\ \\ \relax
6:15 & \begin{tabularx}{0.7\textwidth}{X} 你撒種，卻不得收割；踹橄欖，卻不得油抹身；有新酒，卻不得酒喝。 \end{tabularx} \\ \\ \relax
6:16 & \begin{tabularx}{0.7\textwidth}{X} 因為你遵守暗利的規條，行亞哈家一切所行的，順從他們的計謀；因此，我必使你荒涼，使你的居民遭人嗤笑，你們也必擔當我百姓的羞辱。 \end{tabularx} \\ \\
[1ex]
\hline
\hline
\end{longtable}
六章八節是這卷書的中心思想,許多人都很熟悉.在思想這段聖經前,我們先思想一兩件事.首先這段聖經是向百姓說的,因此這段聖經強調人與人的關係,而不是講及救恩的問題.一個已得救恩的人,應當怎樣在神和在人面前生活；一個已得救的人,他生活應有得救的表現.這人生活若與救恩不符合,他的救恩可能不受影響；但卻在人面前失掉見證效用.彌迦提醒屬神的人,應當怎樣在地上過生活:(六8)「世人哪!耶和華已指示你何為善,他向你所要的是甚麼呢?只要你行公義,好憐憫,存謙卑的心與你的神同行.」這裡提及三件事,就是行公義,好憐憫和存謙的心與神同在.其中第二件事是論及人與人之間的關係,只有第二件事是論及與神之間的關係.人和人之間若沒有正常的關係,就表明他與神之間的關係一定有問題.約翰告訴我們,如果你不能愛看得見的人,你怎能愛看不見的神.當然世上四周的人,都是按神的形像造的,都是我們應當愛的對象,特別是屬主的人.如果我們與人之間沒有正常關係,我們與神之間的關係一定受虧損.在(六8)說世人哪!原來是說人哪!這與先知所用的字不同,他在前面用了百姓很多次；但這裡不用百姓了,百姓是指整個屬神的百姓,是多數的名字.但先知卻不用多數的字,而用單數的字,指著每一個聽先知話的人,先知所給我們生活的標準,是對我們每一個人說的,不是對整體百姓說的.我們不理會別人的生活怎樣,我自己有個人的責任；神的話向我們個別講的,神的命令是給我個人.許多時候我們將這責任推卸周圍的環境,在一個熱心愛主的基督徒群體中,我們容易過聖潔的生活；但在四周冷淡不追求的環境下,只要我與他們差不多就可以了.先知提醒我們要對自己的生活,行為負責任,我應怎樣在神面前盡責忠心.上面提及幾件事,都是基督徒生活最基本的標準,沒有太特別和新鮮,神早已把這些話指示出來.只是我們對已知道的真理未能存著順服的態度來遵守,還要求神想要新鮮,超奇的真理,要神用超然的方法,在人或教會中行事這是錯誤的.聖經的真理已完完全全擺在我們面前,所以神不需要用超自然的作為,祂也不必給我們超奇的啟示.路加十六章提及拉撒路與財主的比喻,拉撒路在地上受苦,財主享福,但拉撒路在地上過敬畏神的生活,財主卻過著不敬虔的生活,他們離世後到兩個不同的地方.拉撒路在亞伯拉罕的懷裡享福,財主在陰間被刑罰受苦.財主向亞伯拉罕央求差遣拉撒路用指頭蘸給他一點水,因為他在火焰中極其痛苦；但亞伯拉罕說不可能,因為你我之間有深淵阻隔.財主便對亞伯拉罕說,請你差派他到世上警告世上的弟兄,叫他們悔改.亞伯拉罕說,不必了,因為他們已有摩西的律法擺在眼前.財主說若世人看見那些從永世回來的人,他們一定會容易悔改,亞伯拉罕說如果聽了摩西的話而不肯悔改,那麼就算聽見永世回來的人的話,他們也不肯聽從.這個比喻看見神已經把祂的話擺在我們面前,如果我們不肯順服已聽見的真理,今天我們要在神面前負應當的責任.
\newline
\newline
現在看神藉彌迦傳的幾個命令標準,這三件事我們簡單地思想.
\newline
\newline
(一)行公義
\newline
\newline
這處公義與別處的構造不一樣,例如在別處提及我們與神的關係,因信就在神面前稱義.這裡公義是我們在人世間過生活的標準,是公正,正直的生活；而這種生活標準已經在聖經裡啟示了.當馬丁路德把這節聖經繙成自己言語時,他譯成你應當遵守神的話,這個意思很直接.在今天混亂的世代中,我們聽見很多意見和聲音,我們很易把世人的生活,成為我們生活的標準.彌迦說:世人不論怎樣說和怎樣作,如果他所作的不合乎神的標準和神的話語,我們便有責任不去遵行.我們應當行正直的話語,按神的標準生活.一九七七年美國總統卡特就職宣誓的時候,卡特總統把手放在這節聖經上,在宣誓前他說:「我靠著神的恩典,在未來四年中行事正直,能真正行出神要我們行的標準.」在他離職時,一本雜誌報導他與內閣的談話,他表示很失望.他說:「有些事情我們知道應當怎樣做,但在這種環境和社會中卻行不通,是非常可惜的事.」在一個不尊敬神的社會裡面,沒有可能照神的話行得通.一百年前英國大主教湯普威廉,他見當時的首相時,首相表示應按登山寶訓來治理國家,但湯普威廉說不可能,怎麼一個主教說用聖經治理國家不可能呢?湯普威廉說,登山寶訓是神給屬神的人生活的標準.今天我們國家的人不全都屬神的,你不能用神的話來治理一個不屬神的社會,這說話是很普遍.但最可惜的今天許多屬神的人也不能用神的話來生活,這是多麼大的悲劇和可惜的事!今天在教會和基督徒圈子裡,許多時候也不能按神的標準來過生活,有許多不正義的行動,我們看見有某教會與某教會打官司,教會裡有欺騙的行為.神責備百姓說:「在你們中間有可惡的小升斗麼,你們仍然用不公道的天平和囊中詭詐的法碼.」在社會詭詐欺騙非常普遍,但可惜在基督徒圈子中也常常發現這樣.最近我與一位在市政局圖書館工作的弟兄談話,他工作了兩年多,而每週平均失書三十多本；甚至在基督教學校也有失書,為甚麼如此?有一位在超級市場工作的人表示,在他們營業額的百分之一或千分之三是被人所偷竊成了很普遍的事情,故此他們訂價錢時也把這損失列入預算之內.有很多方面有時不是很明顯的,但卻有許多方面卻顯明是不公平的.現時地產市道較好,他們不是按成本和利潤定價錢,而無理將價錢提高,除非你不買.專利的大機構,是按照自己的方法,可以定出很高的利潤,這都是不公道的天平和法碼；香港是這樣,很多社會也是這樣.人如果繼續這樣生活,(六15)「你必撒種卻不得收割,踹橄欖,卻不得油,踹葡萄卻不得酒喝.」因為人行事不合乎公義的標準,所以就沒有工作的果效和收成.我們當然不能用物質,金錢衡量,我們沒有按神的標準和正直的標準來生活；我們可以很努力去工作,但得不著果效!我們用盡很多人力,物力,財力,結果得不著屬靈的果效,這是多麼可悲的事情!今天我們在教會工作,用了很多力量和忙碌,但卻不見果效；有時我們看不見果效便再加點力量,但仍沒有屬靈的果效,結果我們越來越忙碌,便被忙碌所佔據.香港的基督徒人數不算少,也許佔人口的百分之三,但按整體計算可能佔百分之五,就算真正基督徒只有百分之三,有十五六萬基督徒,應該在社會上有很大的影響力又有多少呢?先知提醒我們要省察自己的行為,叫我們在這社會有應該的作用,按神所啟示正確的標準來生活.
\newline
\newline
(二)好憐憫
\newline
\newline
這句話很難繙譯,呂振中把憐憫繙成堅貞,舊約有許多地方用這字,是神與百姓立約的慈愛,是特別的慈愛,在另外的地方顯明出特別的用法,就是先知的用法.當兩個探子進入耶利哥,被妓女喇合接待,喇合把他們隱藏起來和釋放,喇合對那探子說話:「我今日恩待你們,神使你們佔領這地的時候,你們也要恩待我.」恩待就是這個字.列王記上記載以色列人攻打亞瑪力人時,遇見堅尼人；是在以色列人出埃及曾恩待以色列的人,所以百姓攻打亞瑪力人,他們特別保守堅尼人不受攻擊.神的僕人對堅尼人說:「我們遇見困難的時候,你們恩待我們」,這恩待就是這個字.我們將所有用法放在一起,便看見這個字有特別的意義,不是憐憫,可憐的意思,而是恩待,妓女喇合不是可憐探子,而是看見神與百姓同在,所以她願意把自己與百姓連在一起,站在同樣的地位.有同情認同的心,你,我之間有親密的關係,我們今天要注重憐憫,不是看見別人可憐,因為覺得自己至少比他好一點；而是我與他站在同樣的地點,因此我要完全認同他.在百姓快要被亞述滅亡的時候,神吩咐祂的百姓:你們都要被敵人壓迫,你們的責任是彼此認同；需要憐憫,心靈彼此同情,在靈裡真正彼此關心.今天許多人不承認基督特別的地方,基督徒在今日的社會裡,應有真正的彼此憐憫的心,大家是主裡的人,應有真正彼此的關心,同情.我們不是不關心社會上貧窮痛苦的人,我們應該關心他們,但我們最主要的關心對象,乃是主自己的百姓和主自己的工作,認同關心屬主的人.
\newline
\newline
(三)存謙卑的心與你的神同行
\newline
\newline
謙卑是重要的字,在構造上有特別的意義,真正的謙卑乃認識自己對神完全順服的心,不是盲目的謙卑和小看自己,不是不認識自己的權利和地位.今天我們忽略了這件事,我們要在謙卑中仍有勇氣,彼得和約翰被猶太官長捉去的時候,猶太官長說你們以後不准奉這人的名講道,彼得和約翰說:「順從人不順從神這合理嗎?」我不是不尊重在位掌權者,但我是屬神的人,我站在神話語的權威上,順服人不順服神是不合理的事,我還會奉耶穌的名傳道.我們常要有真正的謙卑,但真正謙卑的對象是神自己.我們順服神,在人面前不計較得失；不理會別人怎樣看我的成就和收穫,我不忘記我在神面前的地位和收穫.站在神百姓的立場我們要對神順服,而在人面前要謙卑,但卻不能順從人不順從神.謙卑最重要的意思,就是與神同行,與神同行很難繙出來,原文是與誰作事情,與神同行是神有權作主；我與神同行,並不表示我與神是同等地位.祂作的我跟著去作,祂走我也跟著祂走.這句話顯出我們有真正的謙卑和順服,我只不過跟著神去作,一切決定發出於祂.阿摩司書三章二節「二人若不同心,豈能同行呢?」與神同行要與神同心.與神同心的唯一方法,就是讓我順服神的心意,而不是要神遷就我；祂是主我是僕,我只有完全順服在祂的旨意,我喜歡就順服；祂的旨意我不喜歡也順服；要有這樣與神同行的心,這才是真正的謙卑.今天對世界最好的是甚麼呢?就是行公義,好憐憫,存謙卑的心與你的神同行.
\newline
\newline


\newpage

\section{第57屆~港九培靈會~鮑會園~第九講　奮興的道路}
\label{sec:481}
\textbf{第九講　奮興的道路}
\newline
\newline
連結: \href{https://www.hkbibleconference.org/session-message/view/481}{\texttt{https://www.hkbibleconference.org/session-message/view/481}}
\newline
\newline
\hyperref[sec:480]{< < < PREV SERMON < < <}
~
\hyperref[sec:toc]{[返目錄]}
~
\hyperref[sec:463]{> > > NEXT SERMON > > >}
\newline
\newline
彌迦書 7:1-7:20
\newline
\begin{longtable}{cl}
\hline
\hline
章節 & 經文 (和合本修訂版)\\
\hline
7:1 & \begin{tabularx}{0.7\textwidth}{X} 我有禍了！我好像夏日收割後的果子，又如收成之後剩餘的葡萄，沒有一掛可吃的，也沒有我心所渴想初熟的無花果。 \end{tabularx} \\ \\ \relax
7:2 & \begin{tabularx}{0.7\textwidth}{X} 地上的虔誠人滅盡了，人世間已無正直的人；他們都埋伏，為要流人的血，用羅網獵取自己的弟兄。 \end{tabularx} \\ \\ \relax
7:3 & \begin{tabularx}{0.7\textwidth}{X} 他們雙手善於作惡，君王和審判官都索取賄賂；位高的人吐出心中的慾望，彼此勾結。 \end{tabularx} \\ \\ \relax
7:4 & \begin{tabularx}{0.7\textwidth}{X} 他們當中最好的，不過像蒺藜；最正直的，不過如荊棘籬笆。你守候的日子，懲罰已經來到，他們必擾亂不安。 \end{tabularx} \\ \\ \relax
7:5 & \begin{tabularx}{0.7\textwidth}{X} 不可倚賴鄰舍，不可信靠密友；甚至對躺在你懷中的妻子也要守住你的口。 \end{tabularx} \\ \\ \relax
7:6 & \begin{tabularx}{0.7\textwidth}{X} 因為兒子藐視父親，女兒抵擋母親，媳婦抗拒婆婆，人的仇敵就是自己家裡的人。 \end{tabularx} \\ \\ \relax
7:7 & \begin{tabularx}{0.7\textwidth}{X} 至於我，我要仰望耶和華，等候那救我的神；我的神必應允我。 \end{tabularx} \\ \\ \relax
7:8 & \begin{tabularx}{0.7\textwidth}{X} 我的仇敵啊，不要向我誇耀。我雖跌倒，仍要起來；雖坐在黑暗裡，耶和華卻作我的光。 \end{tabularx} \\ \\ \relax
7:9 & \begin{tabularx}{0.7\textwidth}{X} 我要承受耶和華的惱怒，直到他為我辯護，為我伸冤，因我得罪了他；他要領我進入光明，我必得見他的公義。 \end{tabularx} \\ \\ \relax
7:10 & \begin{tabularx}{0.7\textwidth}{X} 那時我的仇敵看見這事就羞愧，他曾對我說：「耶和華－你的神在哪裡？」我必親眼見他遭報，現在，他必被踐踏，如同街上的泥土。 \end{tabularx} \\ \\ \relax
7:11 & \begin{tabularx}{0.7\textwidth}{X} 你的城牆重修的日子到了！到那日，邊界必擴展。 \end{tabularx} \\ \\ \relax
7:12 & \begin{tabularx}{0.7\textwidth}{X} 到那日，人必從亞述，從埃及的城鎮，從埃及到大河，從這海到那海，從這山到那山，都歸到你這裡。 \end{tabularx} \\ \\ \relax
7:13 & \begin{tabularx}{0.7\textwidth}{X} 然而，因居民的緣故，為了他們行事的結果，這地必然荒涼。 \end{tabularx} \\ \\ \relax
7:14 & \begin{tabularx}{0.7\textwidth}{X} 求你在迦密的樹林中，以你的杖牧放你獨居的民，你產業中的羊群；願他們像古時一樣，牧放在巴珊和基列。 \end{tabularx} \\ \\ \relax
7:15 & \begin{tabularx}{0.7\textwidth}{X} 我要顯奇事給他們看，好像出埃及地的時候一樣。 \end{tabularx} \\ \\ \relax
7:16 & \begin{tabularx}{0.7\textwidth}{X} 列國看見，雖大有勢力仍覺慚愧；他們必用手摀口，掩耳不聽。 \end{tabularx} \\ \\ \relax
7:17 & \begin{tabularx}{0.7\textwidth}{X} 他們要舔土如蛇，又如地上爬行的動物，戰戰兢兢離開他們的營寨；他們必畏懼耶和華---我們的神，也必因你而害怕。 \end{tabularx} \\ \\ \relax
7:18 & \begin{tabularx}{0.7\textwidth}{X} 有哪一個神明像你，赦免罪孽，饒恕他產業中餘民的罪過？他不永遠懷怒，喜愛施恩。 \end{tabularx} \\ \\ \relax
7:19 & \begin{tabularx}{0.7\textwidth}{X} 他必轉回憐憫我們，把我們的罪孽踏在腳下。你必將他們一切的罪投於深海。 \end{tabularx} \\ \\ \relax
7:20 & \begin{tabularx}{0.7\textwidth}{X} 你必按古時向我們列祖起誓的話，以信實待雅各，向亞伯拉罕施慈愛。 \end{tabularx} \\ \\
[1ex]
\hline
\hline
\end{longtable}
感謝主讓我過去幾天與大家分享,在今天時代,有機會思想神的話,是神的恩典.這章聖經指出神在百姓中完全的計劃,可以說是百姓整個屬靈生活的旅程.從他們最軟弱的時候開始,直至神在他們身上完全成就.
\newline
\newline
(一)百姓現在的情況
\newline
\newline
他們落在現今境況的原因,怎樣能改變他們的境況,最終使神的計劃完全成就.我們不但關心以色列百姓靈程的經過,也關心對我們的經歷把它描寫出來.現我們先看百姓的經歷,然後應用在我們身上.這裡提及百姓被擄,被征服時的痛苦,然後提到他們被擄後靈裡的改變.(七1)「哀哉!我好像夏天的果子已被收盡,又像摘了葡萄所剩下,沒有一掛可吃的.」好像一個飢餓到了應收莊稼的田地,想得一點飽足,但卻使他非常失望,以色列百姓到了神應許的地上,應得著流奶與蜜豐富之地；但結果他們被擄到外邦之地,一點也不得飽足.世人待他們的情況非常可怕,世人埋伏殺他們,沒有一點正直和關心,以色列百姓落在可憐的境況下.這段經文分節方面,第六節有一圓圈,好像成為一個段落.但從文法構造來看,這分段不大恰當；應把這圓圈放第七節的末字,第一至第七節都是說及百姓在失望中的感覺,到八節百姓轉過來向仇敵說話.(七4)「他們最好的,不過是蒺藜,最正直的,不過是荊棘籬笆；你守望者說降罰的日子已經來到,他們必擾亂不安.」百姓的信心完全失落,不再對人依靠和信賴了.(七5-6)「不要依賴鄰舍,不要信靠密友,要守住你的口.不要向你懷中的妻題說,因為兒子藐視父親,女兒抗拒母親,媳婦抗拒婆婆,人的仇敵就是自己家裡的人.」新約中主耶穌引用這說話在不同的情況裡,描述以色列百姓在欺壓中太過痛苦,所以彼此也成為仇恨；連自己最親密的也不可以信賴了.沒有人可以成為他們的幫助,人在這情況下完全孤單,是多麼痛苦可憐的生活.但到底以色列百姓還是神的選民,沒有世人可以信賴幫助的時候,他們沒有完全失望.(七7)「至於我,我要仰望耶和華,要等候那救我的神,我的神必應允我.」在人看來沒有希望的時候,以色列百姓心靈有一轉變,知道他們不可再靠自己或世人了.只有依靠耶和華.仰望和等候在原文是同義詞,事實形容兩種不同的盼望等候；一種仰望說在不知道有盼望的時候,仍然存著仰望的心,這盼望不是在可以看得見的環境來的,如同以利亞在迦密山上與假先知拜巴力比較誰是真神的時候,他的僕人看不見有下雨的徵兆,以利亞卻叫他僕人再去,一直到第七次,才看見有一片小雲從海裡上來.這正顯出以利亞那種看不見而仰望的心；在環境看來沒有神施恩的時候,仍在仰望.
\newline
\newline
(二)有把握得到所期望的神的幫助來了
\newline
\newline
我也不需再去看,很清楚堅固的信心.百姓現在有這經驗,今天我們也要應驗在自己屬靈生活上.我們有否像以色列百姓一樣,好像飢渴的人去找生命的滿足；心裡渴想卻得不到滿足,我們知道只有從神那裡才得到滿足.我們有否感覺心靈裡枯乾可憐,如同詩人在詩篇(六十三1)所說:「乾旱疲乏無水之地,我渴想你,我的心切慕你!」但心靈裡卻得不到飽足,如果一個人要靈裡復興；他先要知道自己靈裡的需要,知道自己心靈枯乾；與神距離很遠,這是我們復興的第一步.如果我們沒有覺自己靈裡乾渴飢餓,可能我們已經麻木了!也許我們有別的東西將這些枯乾佔據了,以為多作工多一點知識便可以滿足了,以為在人面前多一點表現名望便滿足了.但那些東西雖很重要,卻不能代替我們靈裡空虛飢渴的感覺.當我們知道有這需要後,就要思想有什麼能使這渴想得滿足；只有回到主面前,才可以得到滿足.因為人飢餓非因無餅,乾渴非因無水,乃因不聽耶和華的話.
\newline
\newline
以色列百姓知道自己的罪,引致靈裡空虛,罪過成為他們痛苦的原因,(七8-10)「我的仇敵阿!不要向我誇耀,我雖跌倒,卻要起來.我雖坐在黑暗裡,耶和華卻作我的光.我要忍受耶和華的惱怒臨到我,因我得罪了祂,所以我有痛苦.」只有知道自己得罪神,才會在神面前解決心靈枯乾的原因.如果我們不知道得罪神,還不會知道我們心靈枯乾的原因.在整段經文內,只有這種提到百姓認罪的事實,我得罪了耶和華,便足夠了,許多時我們犯了罪,那種罪咎感好像永遠也除不掉.我們已悔改了,也改正了自己的錯誤,甚至自己責備自己；但聖經從來沒有說,基督徒犯罪需要悔改和刑罰自己,我們不可以犯罪後責備自己.馬丁路德在未曾明白因信稱義前,他的罪擔壓到自己透不過氣來,差不多不能生活.在罪擔下非常痛苦,他脫去衣服,用鞭打自己；使身體流血,在寒冷天氣,不穿衣服睡在地板上.他用膝蓋爬上樓梯去神甫前認罪,爬到膝蓋也出血,用盡各種方法折磨自己,但心裡仍沒有平安,沒有罪得赦免的保障.聖經沒有告訴我們要為悔改而折磨自己,基督徒犯罪只可能作兩件事,1.約翰壹書一章九節:「我們若認自己的罪,神是信實的,是公義的,必要赦免我們的罪,洗淨我們一切的不義.」基督徒犯罪後就應在神面前認罪,當然我們得罪了人,也當在人面前認罪.2.我們應該靠著神的恩典離開過往所犯的罪,若得罪了人應想法子補償.聖經卻沒有告訴我們悔改時折磨自己,但很多時候神赦免了我們；我們卻不赦免自己,明明知道神赦免了我們,仍然沒有赦罪的平安!這是非常錯誤的思想,我們應用感恩的心來接納神的赦免,百姓就是這樣.
\newline
\newline
以色列百姓得復興,因為他們在神面前認罪.我們應該知道使我們痛苦的原因是罪,一切的罪基本上都是向神犯的；我們應該向神認罪,我們心靈枯乾空虛,因為與神之間交通有困難!我們要將與神之間的攔阻除掉,不能用別的東西來代替.很多時候我們心靈有罪的題,許多時候我們用別的東西來代替真正的原因；如果我們追求別的,也許可在那方向得到滿足,但是我們在神面前仍得不著心靈的平安.大衛明白這情形,他說在乾旱疲乏無水之地,心裡切想渴慕神,就表達這意義出來,大衛所以心裡疲乏乾渴,乃是因他離開神；與神之間有攔阻,真正叫我們心靈痛苦是罪.除掉罪才能叫我們得到平安,因你是永遠無法在別的地方得到滿足.在登山寶訓裡主耶穌說:「飢渴慕義的人有福了,因為他們必得飽足.」故當你飢渴慕義想念神的時候,你便得到滿足.如果你飢渴想得到食物享受和飽足,你就不單得不到心靈的滿足.也得不到你得要的東西.我們要認清我們飢渴的真正原因是什麼?今天世人和許多人過著心靈飢渴的生活,想得的雖然也得到了,但心裡卻得不到滿足.幾年前神學院有一位姊妹找我,她哥哥尚未信主,她看見他近來心情不大好；自己不能勸導哥哥,便問我是否願意跟她哥哥談談?我樂意答應了.過了一段時間,她哥哥到辦公室找我,她哥哥談及自己的經驗；我覺得這也可以代表世上大多數人的感覺.她哥哥說:「昨天晚上我悶到不得了,很孤單,於是去看電影,看過電影後,心裡更加苦悶,我想與人打架,甚至被人打一頓也好.」他以為自己苦悶孤單,因為缺少娛樂；但誤樂不能滿足他的心,他心靈苦悶孤單是因為他遠離造他的主.只有認清楚這個追求方向,我得罪了耶和華,我要追求回到耶和華的面前.當以色列百姓這樣承認罪過時,神立刻給他們一個寶貴的應許.這應許在以色列歷史真應驗了,(七11-13)「以色列阿!日子必到,你的牆垣必重修,到那日你的境界必開展,當那日,人必從亞述,從埃及的城邑,從埃及到大河,從這海到那海,從這山到那山,都歸到你這裡.然而這地因居民的緣故,又因他們行事的結果,必然荒涼.」神應許給百姓,果然應驗了,耶路撒冷要重新建造起來,城垣要重修,人要歸到耶和華的面前.在我們屬靈生活上,同樣可以應許,在俄巴底亞書也看見外邦人嘲笑百姓,也譭謗神.你們既屬耶和華,但耶和華存在的記號在哪裡呢?你們仰望耶和華,耶和華有什麼能力幫助你們呢?百姓的失敗,成了外邦人譏笑的原因,今天我們失敗時,在外邦人壓迫下,照樣心靈枯乾,生活沒有目標和盼望,在人面前失去見證,你說信耶穌可有喜樂,但你一點喜樂也沒有；耶穌可以給你力量,但你在試探中不能站住.我們失掉見證,在未信主的人面前羞辱神的名；但神在這裡有應許,你要把牆垣重新修造,讓外邦知道以色列真有神.照樣今天我們在神面前認罪,神要把我們的牆垣再重新建造,使我們在外邦人面前成為有見證的人,人不會因我們緣故譭謗我們的神.神也叫百姓從亞述,埃及歸回到耶路撒冷.歸回與神有交通,這更加寶貴!不單在外邦人面前恢復見證,而且心靈裡享受神滿足的喜樂,享受與神交通的甜蜜.這寶貴的應許給神的百姓,照樣今天也給了我們.
\newline
\newline
(三)百姓在靈程所作的
\newline
\newline
第一步先知道自己的痛苦,第二步在神面前承認自己的罪,第三步要重建牆垣,回到神面前.現在我們看第四步.從(七14-20)記載豐富的內容,簡單說是百姓向神發出讚美的聲音.「用你的杖牧放你獨居的民,就是你產業的羊群,求你容他們在巴珊和基列得食物,像古時一樣.巴珊和基列,是以色列最美好的牧場,水草非常豐富,是百姓最享受的地方,我們像羊群在巴珊和基列享福,躺在那裡讚美神的豐富.「神阿!有何神像你,赦免罪孽,饒恕你產業之餘民的罪過,不永遠懷怒,喜愛施恩,必再憐憫我們,將我們的罪孽踏在腳下；又將我們的一切投於深海,你必按古時起誓應許我們列祖的話,向雅各發誠實,向亞伯拉罕施慈愛.」(18-20)有什麼神好像禰那麼有慈愛,只要我們悔改,禰就恢復與我們的交通,只要我們認罪,你就使我們靈裡得滿足和施恩.神向百姓所施的恩何等豐富；只要我們回到神那裡,便可看見神的豐富,便要讚美神.當我們這樣讚美神的時候,便看見我們心靈的改變,以色列讚美神,對外邦仇敵的態度改變了.(六 10)看見百姓很輕視這班外邦人,要踐踏他們如同泥土一樣,但在(六17)情況便不同了.「他們必舔土如蛇,又如土中腹行的物,戰戰兢兢的出他們的營寨；他們必戰懼投降耶和華,也必因我們的神而懼怕.」這節經文很難翻譯,但這是指出外邦人在神面前得復興的經驗,全世界的人要來敬拜耶和華.外邦人要舔土如蛇,意思是他們要在神面前謙卑.如同詩篇說,外邦人要在耶和華面前屈膝,要用嘴親彌賽亞的腳,他們謙卑在主面前,敬拜真神.在先知彌迦的嘴裡,在以色列百姓的經驗裡,不再是一群可厭惡,應當受刑罰的人了.這些人是可憐的,需要在神面前謙卑,尊耶和華為主,需要接受主耶穌的救恩；如果我們在神面前靈裡得復興,我們不再關心報仇,而是關心到所有未信主的靈魂.有一本書叫做「為主受苦」,有許多人讀過,講到有人為主緣故關在監牢受很大的痛苦；面爛出血,還有人打他的臉.但那位弟兄說,我沒有感到臉上痛苦,我感到這打我的人多麼可憐；他心靈多麼黑暗需要主的救恩.這是百姓的經歷,他們受很多痛苦.後來在靈裡復興過來,讚美神的時候,為那欺壓他們的人禱告,為所有人祈求,但願他們都謙卑在神面前來敬拜我們的主.在今天混亂的時代裡,我們需要的是真正靈裡復興的經驗；我們的生命才有真正的見證,才可以帶領那些在黑暗裡的人,求神保守施恩我們,讓祂的話提醒我們罷!
\newline
\newline


\newpage

\section{第58屆~港九培靈會~吳勇~第一講}
\label{sec:463}
\textbf{教會的破口}
\newline
\newline
連結: \href{https://www.hkbibleconference.org/session-message/view/463}{\texttt{https://www.hkbibleconference.org/session-message/view/463}}
\newline
\newline
\hyperref[sec:481]{< < < PREV SERMON < < <}
~
\hyperref[sec:toc]{[返目錄]}
~
\hyperref[sec:464]{> > > NEXT SERMON > > >}
\newline
\newline
我曾在港九培靈研經會講過幾次道,有次在講台上休克.現在年紀較大,體力比前更不如前了；請弟兄姊妹為我們三位講員代禱.
\newline
\newline
此次我要講的題目,和日程表上印出的不同；但總題仍是教會與聖靈,前五天我要講「教會」,給弟兄姊妹看看今日教會的光景.後五天我要講聖靈,我恐懼戰兢,特別是講聖靈的信息；因為「聖靈」,不能夠一直在講台上講道,必須講些經歷,否則單講一些道理沒有甚麼好處.談到經歷,有些人或不贊成；因為我的經歷,不一定是你的經歷,可能產生很多誤解；但是出於神的,神必負責任.
\newline
\newline
提防破口,阿摩司書第九章,論及要堵塞破口,以免帳棚倒塌．
\newline
\newline
船有破口,海水浸入,船和海水二者分不了.城也不可破口,否則敵人便有機可乘.照樣,有了破口,外面的人流入,而裡面的人也要流出.
\newline
\newline
也許你我的破口不一,每教會的破口不一樣.我今晚要提出四點較為普遍的.
\newline
\newline
(一)基本信仰動搖
\newline
\newline
這是嚴重的破口.基本信仰不是亮光；水禮有滴的,有浸的,這是亮光問題.耶穌基督再來,有災前派,有災後派,這都是亮光問題.聖經中「國度」,馬太福音講「天國」,路加和馬可講「神的國」；有人就把天國和神的國用比喻來分開,他們講得救和得勝；得救沒有包括得勝,而得勝卻包括得救.有的人說,神的國包括天國,而天國並不包括神的國.把天國的神的國分開了.其實神的國就是天國,天國就是神國,這不過是文史背景問題.讀聖經有很多講究,要講究上下文,文法,文史背景.「天國」是文史背景,因馬太是對猶太人講的,猶太人不敢提及神；如一家庭的孩子不敢講父親的名字,認為這樣是不恭敬的.至於馬可和路加是對外邦人講,如果講天國；外邦人聽不懂,所以就用神的國.實際上,天國和神的國,不過是文史背景的問題；但許多人常為此爭辯.
\newline
\newline
基本信仰並非亮光,亮光不可容；也非制度.有人主張要制度,有聖品的制度,分高低,上下,不至紛擾,教會才有次序；儀式隆重,才能叫人肅然起敬.有人不主張聖品制度,除主耶穌外,我們大家都是弟兄；雖有前後,但一概都是弟兄；只要我們在聖靈裡面,氣氛自然隆重.有的人主張從前有,現在也應該有,這是傳統.但是有人為傳統會限制聖靈的作為,叫聖靈在我們中間沒有自由；所以我們不能死守陳規,要讓聖靈自由作工.
\newline
\newline
基本信仰並非亮光,非制度,非傳統；而是依從聖經；因神是獨一的.如果人有叛神思想,有的人心很寬,不否認別人的神是神,以為這樣才有風度,才不引人反感.這樣就牽涉基本的信仰.神子耶穌是童貞女馬利亞所生,祂為贖罪而死,祂的死像一粒麥子落在地裡的死,結出許多子粒來.所以耶穌基督的死,是從一個肉身進到許多肉身的死,是從有限進到無限的死,若死是為博愛而死,祂不過是個歷史人物,不過是人精神的寄托.這樣就牽涉基本的信仰,基本信仰動搖了；這本聖經也牽涉到基本信仰.聖經是神所默示的,沒有矛盾,沒有錯誤的.約翰壹書二章說「小子們哪......若有人犯罪......」第三章說「從神生的就不犯罪.」豈非顯著的矛盾,顯著的衝突嗎?聖經絕不矛盾也不衝突.希臘文時式分得清楚,有過去的,有現在的,也有將來的；現在式這字有繼續不繼續的意思,「從神生的就不能犯罪」的時式或是現在式的,從神生的人不再繼續不斷犯罪,不會重複地犯罪；雖然犯罪,是偶然犯罪,並非習慣性重複不斷犯罪.聖經絕無錯誤,乃是人對聖經沒有讀清楚.如果人認為聖經有錯誤,這是牽涉基本信仰.我們看見今天的教會,基本信仰動搖了,對耶穌的生起了懷疑,認為聖經有錯；因為聖經是許多人寫的,經過許多時代；聖經中有歷史,有科學,若說沒有錯誤,簡直不可思議,實無法想像的!這牽涉基本信仰.
\newline
\newline
今天的基督徒,不但要信道,也要衛道.可是,為亮光,為制度,為傳統爭辯並非衛道.
\newline
\newline
(二)信徒的靈性低落
\newline
\newline
這是教會第二個破口.我們到海灘去,如看見又髒又亂,便知潮水退了.照樣,我們看見教會又髒又亂,便知信徒的靈性低落了.
\newline
\newline
這是一種信仰的骯髒,今日世代,邪術流行,法國地方,四百多人中僅有一個醫生；但百多人中就有一個相命的,看相算命極為風行.世人不信神,籌備開店營商要擇吉日,家人去世要看風水墓地,這不叫做信仰的骯髒；但如果基督徒這樣做,就是信仰的骯髒.世人參加喪事,在遺像前鞠躬下拜,因他們沒有信神,所以不叫做信仰的骯髒.但是基督徒就不是了,因聖經記載,任何形像都要不拜；我們若明知故犯,就是沾上了信仰的骯髒.另一種是生活的骯髒.生活包羅萬象,一下無法說完,讓我們提出兩樣來思想:「謊言」,聖經記載,「是」就說「是」「不是」就說「不是」.我們說話若顛倒是非,就犯上生活的骯髒.「虛假」是非常普遍的；耶穌在世上時很憎惡法利賽人；因他們洗淨杯盤的外面,不洗淨杯盤的裡面；表裡不一,外表有謙卑的態度,內裡卻沒有謙卑的心.表面有愛心的動作,裡面卻沒愛心的念.這叫做「虛偽」,是生活上的骯髒,是信徒靈性的低落.我們還看見教會中有很多的亂,意見不同的亂,亮光的亂,權位的亂,諸般的亂,亦即老我的亂.世界上有政治問題,種族問題,宗教問題.按聖經告訴我們,這些都是私慾問題.人有老我,以自己的意見為對,自己的亮光為好；凡事要超越別人頭上,不肯低屈別人腳下.故此,教會中就產生混亂.
\newline
\newline
靈性低落,乃因信徒不追求,屬靈不長進.活在肉體中,替自己辯護,安慰自己的良心；正如嗜酒者寫酒經,強調飲酒有益可促進血液循環.嗜打麻雀者創立麻雀經,說藉此訓練聰明智慧.活在肉體的人,引用保羅所說的話「你向甚麼人就作甚麼人,好叫我得著甚麼人.」向喝酒者作喝酒的人,向打麻雀的作打麻雀的人,向戲迷作戲迷.其實,這樣不過是為肉體辯解；用聖經的話,安慰,賄賂自己的良心.讓自己犯罪,使良心不至不安；於是為非作歹,膽子包天.這是信徒靈性低落,是信徒本身的緣故.另一種是傳道人的緣故.先知以利亞,他專講方言,不講預言；因此亞哈王把他下在監獄,但這勇敢的先知,敢言人之所欲言,無所避諱.以斯拉時代,他被派出去,有七個謀士陪同,王委任他代表王的權柄.屬靈的意思,以斯拉是作出口的人,七個謀士代表寶座的七靈,講話有聖靈的能力.以斯拉何以有此福分,聖經說他是個有大德的文士；英譯他是完全人,完全意即生活無可指摘.
\newline
\newline
今日的傳道人,常不敢指責罪,不敢指責淫亂的罪,因教會的長老犯了淫亂的罪；不敢指責紛爭的罪,因教會的執事正在紛爭；不敢指責貪愛世界的罪,因信徒們都貪愛世界,一旦指責,恐怕全部跑光了.所以我們看見,靈性低落,除了信徒本身其一原因,傳道人也是個原因.
\newline
\newline
(三)迫聖靈離開教會
\newline
\newline
「...... 要按著心靈的新樣,不按著儀文的舊樣.」(羅七6)原文沒有「新」字,只講到靈樣.如果你在儀文裡面是舊的,你在靈裡面才是新的.瑪拉基有句話:「這些事何等煩瑣,並嗤之以鼻......」這些事就是屬靈的事.以賽亞也說「嚴肅會令我厭惡.」嚴肅會是屬靈的事,為何厭惡,為何嗤以鼻呢?
\newline
\newline
希西家王年代,曾經有一個大復興,在復興之中,他作了很多事,其中有一件事:「他把摩西時代留下的銅蛇打碎,成了銅塊.」(王下十八章)他只注重事物而不注重意義；銅蛇有蛇的樣子卻沒有蛇的作為；耶穌有人的樣子卻沒有人的罪,所以「能擔當人的罪」乃裡面的意義.可是猶太人對屬靈的事,只顧事物不顧意義,故此,他把銅蛇打成銅塊.
\newline
\newline
我們領聖餐,有餅有杯,我們只管事物,把餅擘開,而不管意義.這樣的光景就是儀文.教會屬靈的事,如果只重外面事物,不重裡面的意義；漸漸就會變舊,舊了就會令人厭惡.我們禮拜時,起立,唱詩,禱告,坐下,一切聽主席指示,禮拜完畢,好像空空而歸；因為那些都是儀文,每週如此,年年如此,舊了,氣氛不新鮮不活潑；教會有兩種人,其中一種人忍受不了這樣的氣氛；另一種人因在職教會,責無旁貸,只好咬著牙根上任.教會日久,樣樣陳舊,十字架的道理年年如一,種義的道理,從前和現在都是一樣；聽者講者都覺得舊了,所以就用種種方法變化花樣；找新的故事,活潑的新聞,藉此變換.甚至聚會也想辦法變花樣；有室內聚會,郊外聚會；很多教會效法商人諸多變更花樣,但傳道人無論如何都變不出甚麼,若死在儀文裡面,雖費盡全力去變也屬徒然；如果在聖靈裡,不必變換,雖舊常新,有靈的新樣.
\newline
\newline
「靈」實在很難解釋.有位教授朗誦詩篇第廿三篇,很有節奏,表情十足,聲調高低分明,讓聽者覺得悅耳.接著有位鄉下老婦也朗頌詩篇廿三篇,她不懂節奏,沒有表情,把全篇照讀；但全會場氣氛緊張,會眾非常嚴肅,甚至有人流淚.照人看來,教授的朗誦比老太太好,但教授是在經驗裡,老太太卻在靈裡.在聖靈裡是新穎的.在教會中,一首詩是舊的,在聖靈裡唱這詩是新的；一篇道理是舊的,在聖靈裡講這道是新的.離開了靈,任何東西都變成舊的了.今天教會迫聖靈離開；聖靈離開之後,教會中所有唱詩聚會講道,都變成舊的.
\newline
\newline
聖靈實在不願意離開.聖經中有許多例證.當約書亞在艾城打敗仗,便問神,戰爭為甚麼會敗呢?神說:「你們中間有人取了當滅之物.」為何神不把亞干指出來呢?因神給亞干機會,神叫約書亞給他們抽簽,先取了猶大宗族,繼而謝拉,迦米,最後制出亞干.為何不直截了當告訴約書亞這事是亞干幹的呢?因為神給人留下機會,神實在不願意離開人.
\newline
\newline
聖靈為何離開教會?我們當然想到這是罪的問題,罪迫聖靈離開.但是兄弟不以為是罪,固然罪迫聖靈離開；可是罪很容易發覺,容易對付.迫聖靈離開的是個不容易看見的；而罪迫聖靈離開乃公公然然的,好像一個家庭,父親老了,兒女厭棄討厭他,卻不敢公然迫走他,就用種種方法不聽從他,不與他接近,待他如陌生人,使老父傷心離開.我們也是這樣,不給聖靈有說話的機會,有工作的自由；使聖靈不得不離開.詩篇有句話「以別的代替.」今天教會,代替之物太多,以工作效力代替生命的見證,各教會注重工作的才能,工作的安排,要不注重生命的追求.生命的追求是如何捨己,如何背負十字架,如何講究死在我身上發動；叫生在別人身上發動,講究攻克己身,叫身服我.但是今天之教會不講究這些,只講工作不講生命,用肉體的娛樂來代替屬靈的喜樂.這樣何來屬靈的喜樂?人若順服神,不叫聖靈擔憂,便有屬靈的喜樂,彼此相愛,愛人者和被愛者都有屬靈的喜樂.
\newline
\newline
(四)對福音沒有負擔
\newline
\newline
耶穌說:「你們要聽,這是命令.」保羅也說過這樣的話,講到傳道的問題,無論得時不得時都要傳福音.意即傳福音是無理由拒絕的,時間許可或不許可都要傳,健康許不許可都要傳；因這是命令,無可推辭.有個青年人在軍隊當兵,他預備結婚,向上司請假不獲批准；就打電話請他父親寫信來,持信再去請假,仍然不准.和軍官理論,結果被監禁三天；因軍隊中只講命令不講理由.
\newline
\newline
福音是有大能的.
\newline
\newline
許多人對福音的看法,是鄉下落後地方才需要；因為這些地方沒有進步,沒有文明.但是保羅自從看見異象,發覺此種看法完全錯誤.保羅看見一馬其頓人叫他過來幫助,他就到馬其頓去；該地屬於希臘,希臘是世界著名的堡壘,有許多大哲學家,如柏拉圖,亞理斯多德等,希臘對世界文化有大的影響,所以傲視全球；既如此文明,文化既如此發達,神還叫保羅去傳福音；因為他們不得滿足,他們拜的許多假神,叫做未識之神,有萬神廟；他們不認識哪是真神,拜的全是假神,反而對真神忽略了.希臘雖有好的文化和進步的文明,實際上他們需要福音；所以神叫保羅把福音帶給他們.文化不能救人,不能救個人家庭社會；神拯救人乃是福音的大能,福音能改變一切,把邪惡改變為良善,相恨的改變為相愛的,仇敵改變為朋友,這是福音的作為.我曾提過有個大流氓被判十四年監,因他常逃獄,每逃一次便加重罪並延長刑期.有個大學畢業女生很歡喜這囚犯,一共給他寫了三百多封信,到了最後一封信,終而打動了這流氓的心,於是他成為基督徒.行為完全改變,且能幫助其他作惡的人；所以獄方就給他假釋,僅監禁了五年.他想到神的恩典,內心非常感激,為要報答神的大恩,他報名要進修神學；但是神學院收的是大學畢業生,他僅是高中畢業；因此讀不成神學,那時某工廠要請位企管主任；他參攷了此類書籍然後投考.結果在百多名求職者中被錄取,於是他當了企管主任；其實他的本意想作主工,他的女友來告訴我這些情形,我聽後頗感興趣,請他來見我；交談中他述說見證,我發現他實在信了主,神在他身上作了工作.
\newline
\newline
我就對他說,神學院不要你,你能否來作我的跟班,我收你為徒.他馬上答應,立刻辭去企管主任之職,問我當作甚麼?我說,我買了間大屋,預備給你們住,可以在那裡自修功課.他說「好極了,我年青力壯,我去打掃房子.」我駕車和他同去,過橋時差點碰倒一輛電單車,我急忙停車下來；對那人說「非常對不起,差點碰倒你,令你受驚.」我認為這樣已禮貌週到,就上車走了；不料那人開電單車直追禮拜堂門口,把我擋住去路,下車與我理論,我的跟班大抱不平,就揮拳打那人兩下,打得臉上花.在假釋期他又犯了錯誤；我勸他趕快逃避,一切由我擔當.那人忽爬起；我告訴他,揮拳的人我已經叫他走了.他就一把抓住我；禮拜堂裡正走出許多人,以為吳勇弟兄今天好厲害,把人打成這樣.我說「沒時間作解釋,急速送傷者到醫院.」結果傷者口縫了卅幾針,還把脫臼的下巴復元.醫治完畢,我說「先生,請你回家吧!」他說「那有這麼簡單,把我打成這樣,叫我回家.」他拉著我上警署,警察說「問題很簡單,把那傷人者找出來,你就沒事.」於是我寫了具結,然後回家；事後每日上警署報到.到了第六日,我很不耐煩；忽然那流氓出現,他說「是找我嗎?我是男子漢,敢作敢當,怎能把我的事嫁禍於你,請你帶我到警署,該入獄14年就14年吧!」我想,這也有道理,是做人的態度,作事要負責任.我說「我們一同禱告後就去罷.」我問他,「如果那人看見要打你怎麼辦?」他說「由他打,我絕不還手.」我們到了警署；神行了極大的奇事；我吩咐他在門口等我,我獨自進去,恰巧傷者也在裡面；他起立作揖,對我說「這幾天我看見一件事,你實在是個忠厚老實的人；本來我想敲你一筆,但我思想,敲老實忠厚人的錢,良心過不去；所以我今天要和你磋商,補償一點藥費,我就撤消控訴.」我非常高興地說「如果你肯這樣,真是感激不盡.」他經計算,要醫藥費8600元(台幣)我給他一萬元,事情就結束了.我帶打傷人的弟兄向傷者賠罪說「先生,對不起,我太鹵莽,把你打成這樣；現在我願受你責打.」那人說「不必,算我倒霉就是了.」事後流氓對我說「長老,我令你失望,我實朽木不可雕,我無希望了.」但我對他說「我相信你,相信你到底；因神的工作無不成功的.」他現在傳道,在極度迷信的地區牧養教會,僅兩年多時間,已帶領180多人信主.
\newline
\newline
福音是神的大能,福音是命令.可惜我們沒負擔,頂撞神的命令,不相信福音是有大能的,這也是教會的破口.
\newline
\newline
有了以上四個破口,帳棚怎不倒塌?教會怎不荒涼?所以阿摩司寫至此,說「我們應當堵住破口,帳棚才不致倒塌.」
\newline
\newline
我們當省察,求主光照!
\newline
\newline


\newpage

\section{第58屆~港九培靈會~吳勇~第二講}
\label{sec:464}
\textbf{沒落的教會}
\newline
\newline
連結: \href{https://www.hkbibleconference.org/session-message/view/464}{\texttt{https://www.hkbibleconference.org/session-message/view/464}}
\newline
\newline
\hyperref[sec:463]{< < < PREV SERMON < < <}
~
\hyperref[sec:toc]{[返目錄]}
~
\hyperref[sec:465]{> > > NEXT SERMON > > >}
\newline
\newline
瑪拉基書 3:2-3:2
\newline
\begin{longtable}{cl}
\hline
\hline
章節 & 經文 (和合本修訂版)\\
\hline
3:2 & \begin{tabularx}{0.7\textwidth}{X} 他來的日子，誰能當得起呢？他顯現的時候，誰能立得住呢？因為他如煉金匠的火，如漂洗者的鹼。 \end{tabularx} \\ \\
[1ex]
\hline
\hline
\end{longtable}
瑪拉基論到漂布的鹼,煉金的火.(瑪三2)馬太說:「......蠓蟲你們就濾出來......駱駝你們倒吞下去.」(太廿三24)提摩太前書三章16節論到敬虔的奧秘,被天使看見,被世人信服,被聖靈稱義,被接在榮耀裡.被聖靈稱義意即在靈性上稱義.啟示錄十九章講到穿細麻布的衣服,就是聖徒所行的義(8).廿二章論神的話不可刪,也不可加(19).三章論老底嘉教會(14-20).二至三章論及七個教會,其中示每拿教會和非拉鐵非教會是好的；其他五個教會是不好的.老底嘉是最後一個教會.請看這教會的沒落.
\newline
\newline
第一點,神對沒落教會的意見,其中有三個境界:
\newline
\newline
(1)老底嘉的名字.
\newline
\newline
(2)老底嘉不冷不熱.
\newline
\newline
(3)老底嘉教會看似富足,卻是貧窮.
\newline
\newline
老底嘉教會自欺,其實無有；看自己富有,自以為是.所以神對沒落的教會有三個警戒.
\newline
\newline
第二點,神對沒落教會的勸勉有三點講到「買」.第一買「火煉的金子」.第二買「白衣」.第三買「眼藥」,買眼藥才能看見.
\newline
\newline
神對沒落教會的警戒,老底嘉的字意,「老」意即人,「底嘉」則風俗或意見,老底嘉即人的意見,風俗.風俗含有遺傳,他們把遺傳當作道理教訓人,如洗手,洗腳都屬遺傳.風俗另一意思是裡面的喜好,照心中所喜好而行；他們把信仰變成一套外表的儀式.耶穌在世上時,和撒瑪利亞婦人談話,那女人說:「我們拜神在這個山上；你們說拜神在耶路撒冷.」耶穌說:「不是在這個山上,也不是在耶路撒冷.」山和耶路撒冷都是外面的東西:「拜神要以心靈和誠實.」耶穌來到世上,那時代的信仰已經沒落了,已經變成一套外面的儀式；所以耶穌要把他們從外面導入裡面.人喜好的神蹟,喜動不喜靜；動則注重工作；靜則注重生命；而人喜好工作,不喜好生命.耶穌在伯大尼,祂說馬大是好動,好工作的；馬利亞是好靜,好生命的.主對馬大沒有稱讚,只對馬利亞有稱讚.今天教會也是以工作代替生命.這是人喜好之一.
\newline
\newline
老底嘉教會既是講到人,他們決定事情,並非向神尋求祂的旨意,乃是大家一起尋求人的主意,非以聖經作根據；乃以人數多少作根據,因此教會中充滿了人的意見；教會成為人治而非神治.只有人在其中,神卻不在那裡,氣氛毫不嚴肅；只是發表個人的意見,沒有求告神的旨意.這是沒落教會的光景,主把這種光景公然指出.
\newline
\newline
沒落的教會不冷不熱,按理教會是基督的身體,也是基督的器皿.「教會是祂的身體,是那充滿萬有者所充滿的」(弗一23)既是這樣,充滿者如何,被充滿者也如何.耶穌在世上作甚麼,使徒教會也應作甚麼；耶穌在世上時叫瘸子起來行走,使徒的教會也叫瘸子起來行走；耶穌在世上叫死人復活,使徒也照樣叫死人復活；耶穌趕鬼,使徒保羅也趕出污鬼.所以那充滿者作甚麼,被充滿的器皿也照樣作甚麼.這是我們從使徒教會看出來的.但是老底嘉教會並非如此；因其教會中不是充滿神的萬有,乃是充滿人的萬有,人的意見才能,充滿人的肉體和血氣.耶穌禱告「願你榮耀你的兒子.」天上應說「我已經榮耀了我的名,我還要再榮耀.」這裡有兩個榮耀,第一個是耶穌來榮耀；第二個是等眾教會來榮耀.沒落的教會所充滿的不是神的東西,而是人的東西,神,耶穌怎得榮耀呢?
\newline
\newline
不冷不熱正如溫水一樣,就是「一半一半的意思」今天基督的教會,我們所見都是一半一半,道理如是.神的兒子是耶穌,名字的意思是要把人從罪中救出來；又叫以馬內利,意即神與人同在.應當傳神的兒子,我們只傳耶穌,沒有傳以馬內利,不過只傳一半.耶穌曾說過「初熟的果子」.既講初熟,理應講繼熟；初熟味道是甜的,形狀是圓的,顏色是紅的；而繼熟的樣子卻難看了.基督為我們的罪死了,又照聖經所說第三天復活了.(林前十五)福音,一個字是「死」,一個字是「活」；可是今天傳的福音,只傳死沒傳活,只有一半.還有,屬靈的事,聖餐是屬靈的事情,浸禮也是屬靈的事情；我們領聖餐把餅擘開把杯舉起,有餅有杯,這是物；吃餅飲杯,這是事,屬靈的事情,一部分是事和物；另一部分是意義.餅擘開的意義,就是耶穌的身體擘開；杯裡葡萄汁的意義,就是耶穌流下的寶血；可是人領聖餐,顧及事物而沒顧及意義,這都是一半.浸禮入水中是事,這是事物；事物之外另有意義,浸入水中與主同埋葬,從水中出來與主同復活.許多人行浸禮,只顧及事物不顧及意義,一半一半.我們聚會,用耳朵聽神的道,也要用手腳來行神的道；若聽神道不行神道,也只有一半.老底嘉教會,就是一半一半.神怎麼接受呢?
\newline
\newline
所羅門時代,有兩婦人,其中一個把孩子壓死了,彼此爭奪活的孩子；所羅門勸她們別爭論,叫人拿刀來,要把活孩子分開各得其半；活孩子的生母不忍孩子被宰開；自願把完整的活孩子給對方.一半一半的,人都不能接受,神怎能接受呢?
\newline
\newline
老底嘉教會,道理一半一半；屬靈的事也是一半一半.理應榮耀神,但擺上的卻僅一半.這是老底嘉教會被責備的第二件事.
\newline
\newline
沒落的教會自滿富足,實際貧窮.曾有人到羅馬參觀聖彼得堂,該堂富麗堂皇,金銀滿目；領人參觀的神父自誇地說:「從前彼得沒有的；現在我們有了.」參觀者中有人回答神父說:「但是從前彼得有的；你們現在卻沒有.」從前彼得進入美門時,有個生來瘸腿的求彼得賙濟,彼得說:「金銀我都沒有,只把我所有的給你；我奉拿撒勒人耶穌基督的名,叫你起來行走.」
\newline
\newline
老底嘉教會是個自欺的教會,我們應該注意的是生命,應該有的是能力；但老底嘉教會所供應的道理,只不過增加頭腦的知識卻沒有生命.我們應當講究豐盛和長進的生命；豐盛生命就是「祂必興旺,我必衰微.」耶穌在我們裡面興旺；這才是豐盛的生命,我自己在我裡面衰微,才是生命的長進.我們應該有的是能力.「......蠓蟲你們就濾出來,駱駝你們倒吞下去.」(太廿三24)蠓蟲是很小的,駱駝是巨型的.罪對付屬靈的事.
\newline
\newline
各位弟兄姊妹!我們的信心是一輩子的,悔改也是一輩子的.我們常為小罪悔改而大罪不悔改.小罪是顯而易見的,是一般人普遍有的,不至引以為恥；大罪是隱未現不為人知的.若洩漏出來,會被人譏笑；所以不敢認罪,沒有能力承認大罪.「祂來的日子,誰能當得起呢?祂顯現的時候,誰能立得住呢?因為祂如煉金之人的火,如漂布之人的鹼.」(瑪三2)這裡給我們清楚看見,布是外面的東西,表面的行為；金是含有礦物質及其他金屬,必須用火提煉.煉金的火代表人的個性；人的行為容易對付,然而個性是難以對付的.教會需要能力,才能對付個性上的軟弱.
\newline
\newline
我有個學生,個性特別,和任何人都合不來；他的心胸狹窄,脾氣暴躁；弄得學校苦無寧日.此類學生,就算成績多麼優良,畢業之後在教會怎能牧養羊群?他承認自己毛病卻無力改過.為了造就他,我就派他到工廠去,和不信主者在一起；那廠長是位弟兄,我告訴廠長；因這學生不合群,我想藉此磨煉他；過了兩週,廠長告訴我；此人來了之後全廠大亂了.於是我對這學生訓話；你到處與人不和,怎能為神工作呢?有一天,他參加一個聚會,禁食禱告；神的能力降在他身上,他完全改變了.
\newline
\newline
弟兄姊妹!教會有的應該是生命,不應該是知識；教會有的應是能力,不應單有儀式.老底嘉教會,卻非如此,而是困苦,貧窮,可憐,瞎眼,赤身,何等痛苦!教會中充滿意見的紛爭,聖經亮光的紛爭,權位的紛爭.真是痛苦可憐!好像畢士大池旁,許多瘸腿,瞎眼,血氣枯乾的.老底嘉教會正是如此!禮拜堂富麗堂皇,主卻不在其間；聚會多,卻有名無實,不能供應生命.示每拿教會也貧窮,但她是物質上的貧窮；老底嘉教會並非物質貧窮,而是靈性貧窮；信徒只喝奶不能吃乾糧；恩賜上貧窮,缺乏各樣的人才；瞎眼,對神的道,神的啟示都看不見.老底嘉教會赤身,聖經說我們應當穿上全副軍裝,用真理作帶子,公義作護心鏡,平安的福音當鞋穿,信心作盾牌,救恩作頭盔,神道作寶劍.老底嘉教會卻一無所有.
\newline
\newline
神對沒落的教會勸勉說:要悔改!要買火煉的金子,買白衣穿,買眼藥.
\newline
\newline
「買火煉的金子」就是信心；我們的信心必須經過試驗.試驗有兩種:一是貧窮苦難的試驗,一是富貴安逸的試驗.今日極權的世界就有這些試驗,試驗結果有真有假,有信有不信；試驗越多裡面越火熱越堅固；教會受此類試驗後,非真則假,非信則不信.在自由境界的試驗,這就是安逸富貴的試驗,和另一種試驗的境界完全不同.我們在香港,受的是世界的試驗,結果模糊地作信徒；不愛主,不追求,不戀慕主.基督徒從世界中間出來,按理應有所分別；但結果仍和世人一樣,基督徒理應是光,是鹽,是有功用的；可是經試驗後,卻變成無功用.世界有事,有物,我們有工作的事,有讀書的事.至於物,我們衣食住行都是物.神造亞當時,許可他有事,叫他修理看守,治理管理,這就是事.園中果子都可吃；果子是物；事是神許可的,物也是神許可的.到了有一天,事物累住我們的心；這樣的事物就謂之世界.比如我們看小說不是世界,可是一旦看到入迷,費了多半的時間,損害了身體精神,再也無暇追求神的話了；這樣,小說就叫做世界.人需要運動,因不運動膽固醇增高,身體肥胖；但如果運動太多,連週日也拼命運動,連禮拜時間也去運動；這樣就叫做世界.
\newline
\newline
各位親愛的弟兄姊妹!罪和世界哪樣可怕?罪叫人污穢,世界也叫人污穢；但是世界也會叫人很文雅的；既如此,為何世界比罪更可怕呢?因為世界能迷惑人,霸佔人.我們是自由世界的基督徒,和極權世界的基督徒不同；可是他們越試驗越堅固；我們卻發現越試驗越軟弱.聖經告訴我們,「買火煉的金子」,就是要追求經得起試驗的信心.
\newline
\newline
「買白衣穿」白衣不是基督的義；如果解釋為「基督的義」是錯誤的；因為基督的義不必買,是白白賜給的.啟示錄十九章說的白衣是信徒所行的義.「大哉,敬虔的奧秘......」(提前三16)「敬虔」,就是神用你的肉身將祂發表出來,這就是信徒所行的義.
\newline
\newline
有位太太的丈夫忽然死了,他是當廠長的,工作時被機器壓傷,以為小小的傷略為消毒就包紮起來；遂即搭飛機往日本接洽生意,次日發燒,至第三天仍不退燒.到醫院檢查,發現患了白血病,急忙打電報通知他的太太,到達時病人已奄奄一息而死；於是遺體火化,骨灰帶回台灣.這軟弱的婦人最大的兒子僅四歲,小的兩歲；她痛不欲生.住她隔壁一位本堂姊妹把這情形告訴我內子,內子就去陪她,她說太陽下山時最害怕；因為每天那時,她丈夫下班,孩子放學回家；如今丈夫不再回來了,那個時刻是最難挨過的.因此我內子每天去陪她,都在太陽下山之前到達.經過數星期,她說:「吳太太,謝謝你!時間能把感情沖淡,經過了這些日子,我已漸習慣了；這個時間不但是我忙的時候,要為兒女燒飯,也是你忙為兒女燒飯；我已可以面對現實,明天你可以不必來陪我了.」
\newline
\newline
「敬虔」就是神要藉著祂的兒女,將祂的愛活出來.「......就是神在肉身顯現,被聖靈稱義,被天使看見,被傳於外邦,被世人信服,被接在榮耀裡.」世人信服基督徒,並非你懂得禮拜,懂聖經知識；乃是你能作他所不能作的,說他所不能說的.如果你能夠靈性稱義,行為上稱義,你就被接在榮耀裡.你有基督的義嗎?當耶穌再來之日,你被提至天空；但如果單有基督的義,不一定能接在榮耀裡；當你被提到天空,主要審判你,這審判與得救無關；但與得賞有關.如果你有基督徒所行的義,在聖靈裡稱義,你就可以進而與祂一同坐席.當耶穌從天空回來,祂榮耀的再來,祂帶了千萬聖者同來；那時你若在其中,就被接在榮耀裡了.所以我們不但要有基督的義,還要在聖靈裡稱義,在行為上稱義；基督的義只能叫你得救,行為上稱義才可以叫你得賞.
\newline
\newline
「買眼藥」叫你能夠看見,看得見是很重要的.有一天,耶穌從約但河走過,約翰說:「看哪!神的羔羊.」約翰和耶穌是表兄弟,他從前沒看見；現在看見了,聖靈叫他看見.西緬也是如此；當馬利亞抱耶穌到聖殿,他接過耶穌來稱頌神為大:「主啊!你可以釋放僕人安然去世,今天我親眼看見你的救恩.」那時耶穌才降生幾天；一事不作,一言不發；西緬能知道是神的救恩,因為聖靈給他看見.
\newline
\newline
傳道人應當看見,講篇道不是說內容,組織,比喻,口才如何；應有兩個條件,一是聖經的根據.啟示錄末章所載,聖經是神的話不可加也不可減.加的意思,不能在神的話之外找話；刪的意思即不能在神的話裡否定神的話.所以傳道人講道應以聖經為根據.還有一點很要緊的,要有聖靈的解釋,就是聖靈給他有所看見；看見的問題要對,傳道人負上的責任是很重要的,尤其對基督徒更是重要.世界使我們昏昏沉沉,軟軟弱弱,所以我們需要看見；因有看見,才能站得起來.
\newline
\newline
親愛的弟兄姊妹!約翰勸勉我們要買三樣東西,買是要付代價的,我們得救是本於恩,白白得來無需代價；有的人因此誤解,既可白白得來,便無需付代價.但是這裡明說「買」,買是要付代價的.不錯,得救不需付代價,但得賞卻需付代價.我們應付何代價?有的人清早起來,把最好的時間親近神,讀神的話,這就是代價.今日世界,忙忙碌碌,雖有悠閒卻把時間犧牲了；追求事奉主就是付代價.這時代物質發展快速,地上的水喝了再喝；但有人以為這是身外物,都要過去,少賺些沒關係,不想把全身精力投入；因到頭來不過兩手空空,所以不被世界物質迷住,一切視為淡薄.這就是代價.
\newline
\newline
求主施恩!苦難的試驗,這一端雖然可怕；但是富貴的試驗那一端更加可怕.所以要追求經得起試驗的信心；不然將會被消除,被推出去.
\newline
\newline
在這時代,我們應當付代價,萬勿落到老底嘉教會地步.
\newline
\newline
要付代價!起來追求!
\newline
\newline


\newpage

\section{第58屆~港九培靈會~吳勇~第三講}
\label{sec:465}
\textbf{隨從世俗的教會}
\newline
\newline
連結: \href{https://www.hkbibleconference.org/session-message/view/465}{\texttt{https://www.hkbibleconference.org/session-message/view/465}}
\newline
\newline
\hyperref[sec:464]{< < < PREV SERMON < < <}
~
\hyperref[sec:toc]{[返目錄]}
~
\hyperref[sec:767]{> > > NEXT SERMON > > >}
\newline
\newline
今日社會秩序,越來越敗壞,犯罪率直線上升；人心愈加黑暗,人情愈加淡薄,人性愈加殘酷；這種光景需要復興.
\newline
\newline
從教會看,越來越儀式化.耶穌曾經責備說:「......你們將薄荷,茴香,芹菜,獻上十分之一；那律法上更重要的事,就是公義,憐憫,信實,反倒不行了.......」(太廿三)薄荷,芹菜,茴香等屬外面的東西,誠實公義是裡面的東西.
\newline
\newline
現在一般人的信仰,注重外面不注重裡面；社會光景如此,教會的光景也如此；所以我們感到需要復興.
\newline
\newline
舊約以斯拉講重建聖殿,尼希米講重建聖城；這兩卷書都是講重建.重建則復興,他們如何帶領復興呢?以斯拉記載:「神為要應驗藉耶利米口所說的話,就激動古列王的心.」神藉耶利米有所應許；但是應許必須有人禱告,才能應驗.當時的王下令不可拜自己的神；但以理不理王的命令,他付上代價；一天三次面向耶路撒冷禱告.尼希米也是一樣,有人報告,耶路撒冷城門被焚燒,城牆被拆毀；他就坐下哭泣,禁食禱告耶和華.所以復興是由於禱告.許多神的兒女都知道復興須要禱告,就聚集禱告；可是很希奇,始終不得復興；既然聖經歷史記載,禱告才得復興,為何我們禱告不得復興呢?先知哈巴谷向神禱告「......求你在這些年間復興你的作為.」(哈三2)復興就是神顯明出來；如要神顯明,必須把神人之間的路打通,然後神才能來到人的中間；要打通此路,必須認罪.
\newline
\newline
很多人常講禱告帶來復興；但是不知應如何禱告.從以斯拉來看,他作認罪的禱告:「我們的罪孽滅頂,罪惡滔天.」尼希米也是作認罪的禱告:「我和我父家都犯了罪.」所以我們要復興,就必須讓神顯明；如果希望神在我們中間顯明,我們就應以認罪為己鋪路.
\newline
\newline
茲題出五種罪
\newline
\newline
(一)靈修不認真的罪,我想這是很普通的事.靈修就是讀聖經禱告,既不認真,當然是讀經禱告不認真.我們讀聖經,看見整個信仰,完全是套生命的故事；神按自己的形像樣式造人,把人放在伊甸園,園中有棵生命樹；人若要有神的形象樣式.生命貫穿全部聖經,生命需要維持.植物動物有它的生命需要維持,我們信了耶穌,有屬靈的生命需要維持.約翰一章說:「太初有道,道就是神.」神把自己化成道,藉著道來維持我們屬靈的生命.所以基督徒如果不認真或不讀聖經,屬靈生命怎能維持?沒有泥土,樹木的生命怎能誰持?沒有水,魚的生命怎能維持?照樣,如果沒有神,我們屬靈生命怎能維持呢?讀聖經並非單從聖經吸收些知識,乃是從聖經中認識神；當我們讀了耶穌叫大麻瘋者得潔淨；我們便知道神是有大能的神.耶穌譴責法利賽人,我們便知道神是憎惡假冒為善的神.耶穌叫寡婦的獨生子死而復活,就看見神是位慈愛憐憫的神.
\newline
\newline
我們的生命不但用食物來維持,也用呼吸來維持；這就是我們的生命不但用讀聖經來維持,也用禱告來維持；所以許多人都用呼吸比喻禱告.人信了耶穌,必須禱告,吃飯時,有需要時,有困難時,自然而然知道要禱告；好像嬰孩出生,不需教導自會呼吸；如果人呼吸需倚靠任何物件；如倚靠鐵肺才能呼吸,這是病態,非自然現象.照樣,人有了耶穌的生命以後,不需要倚靠任何東西禱告,憑感覺禱告；歡喜快樂時就用喜樂的話告訴主；痛苦時就向主申訴苦情；這樣的禱告才是真實的禱告,不需預備,只憑自己的感覺帶到主前.
\newline
\newline
(二)認不敬畏不敬虔的罪(何七8)有句話:「以法蓮是沒有翻過的餅.」餅有兩面.(申十七20)也有句話:「我們不能偏左,也不能偏右.」並非說真理有兩面,而是說真理有這一面,有那一面.我們說耶穌是神子,也是人子.神的屬性不但有公義的一面,也有慈愛的一面；如果我們只顧慈愛的一面,而不顧公義的一面,那就是沒翻過的餅.神有可愛的一面,也有可畏的一面.家庭中的小孩子,如果不怕父親,那簡直是朽木不可雕了；若無所懼怕,甚至連神都不怕；這人簡直無可救藥了.敬畏神就是怕神,基督徒怕神乃因愛神；因不忍傷神的心.如孟子小時好玩不好學,孟母非常傷心,斷杼教子；從此孟子不敢再不好學了；因怕傷他母親的心.基督徒若敬畏神,就不敢假冒為善表裡不一了；在人前顯若正人君子,背人後卻是卑鄙小人.很多基督徒不撒謊,是非分明；因怕傷神的心；但是有很多基督徒則不然；他不怕傷神的心,所以不敬畏神.我已說過,敬虔就是神在肉身顯現,被聖靈稱義,被天使看見,被傳於外邦,被世人信服,被接在榮耀裡.(提前三16)這裡有五次「被」.所以我們看見,人被神在肉身顯現是很要緊的；基督徒被世人信服,決非因懂得宗教的儀文,也非單有聖經的知識；乃因基督徒能說,能行他們所不能的.
\newline
\newline
有次我回家時,看見一個十七,八歲的年青人在禮拜堂門口發抖；我問他「你冷嗎?餓嗎?」他一一搖頭；他手持小刀,喃喃自語「還差一點」.我覺奇怪追問,他說「剛才我以小刀想威嚇一位太太,想搶她的手提包,忽然膽怯不前；想到搶劫在台灣要被鎗斃,我差一點就被判死刑被鎗決.」我說「恭喜你；因為你沒把命賣掉；難道人命這麼不值錢嗎?你年紀輕為何作不法的事呢?」他用手拍頭說:「我真該死!我家住花蓮,我父親給我九千六百元叫我繳學費；因誤交損友,把這筆花光,因此我不得入學,無法向父親交賬.」我說「為了九千六百元險把命送掉；如今你打算怎樣?不如你下午留在這裡,等我找朋友幫忙解決.」於是我備錢和他同到學校繳交學費；總務主任問我和這青年有何關係?我說我實在說不出來；他再三要求我表明我們的關係,以便日後把這學生的成績報告給我.我不得不說「我和這青年不相識.」主任急忙招聚些年青人,對他們說明此事,讚嘆在今天人情比秋雲更薄的世代,還有這樣的事情.弟兄姊妹!我所以提及此事,是為了說明,世人所信服的,不是宗教的知識儀文；乃是你可作他所不能作的,這就是所謂「敬虔」；耶穌是慈愛憐憫的,你應把耶穌的慈愛憐憫活出來.但有許多基督徒,活著是為自己,一切都為己.亞伯拉罕所以被揀選；萬民因他得福.今日你被揀選,乃是要叫人因你得福；若你沒有讓人因你得福,你應要認罪.
\newline
\newline
(三)認愛心有虧欠的罪,人都需要愛；人難免有疾病,有挫折,有傷亡,有失敗,錯誤,糊塗,一旦落在這種光景中,極需要人的關懷,安慰,鼓勵.從愛心給人的關懷或安慰,會叫人垂下的頭抬起來,消沉的意志再次振作；叫人從死亡關頭活起來.各位!如果你在世上愛心沒有虧欠；有一日你在天上,許多人要圍住你要說:「我在地上因你愛心的關懷使我抬起頭來；使我消沉的意志振作,使我從死亡的關頭復生.」當你聽了這些話,你的心將何等快樂!但今天在教會中,我們看見許多人得不到關懷,得不到安慰；因此抬不起頭來,意志不能振作,死亡關頭無法復生；這是我們愛心有虧欠.我們活著,不是單為自己,乃為你周圍的人,為你所看見所認識的人.我們從年初至歲晚,找不到一個從你身上得益的人；這是因你愛心有虧欠.世人自私自利,我們何嘗不是如此!
\newline
\newline
(四)認懶惰的罪,人對屬世之事,不懶惰,且不敢懶惰；無論學業,事業,定要奮發殷勤,才有所建樹；在現今時代,要成功就要努力,要好成績就要奮鬥；如果希望生活改善,讓以後的日子過的比現在舒服,就一定要力爭上游,才能達到目標.世上多為學問,生活,精神體力透支,工作10小時,12小時,14/16小時；本應八小時睡覺；為了明天比今天更好,只睡四小時；甚至有的人連四小時的睡眠都沒有.我在路上看見許多人精神疲乏,眼睛發紅；人對世事不敢懶惰；但對屬靈的事是懶惰的；屬世的學業,由小學,中學,大學；屬靈方面也有此階段,從摩西五經,有外院,聖所,至聖所；摩西上山,有山腳,山腰,山上；到了新約,我們也看見有小子,少年,父老；屬靈有程度,每進一步,步步進深.聖經常用「山」字,舊約摩西看見西乃山；新約耶穌看見黑門山,新舊約都注重山；因山代表屬靈的境界,不流汗上不去,不費力達不到目的.所以我們要達到屬靈的境界,須要付上代價.有許多基督徒,屬世的代價肯付；但屬靈的代價卻不肯付.教會中,我們看見許多屬靈懶惰之輩,他們懶惰靈修,聚會,事奉.箴言書載「懶惰人說,外頭有獅子,我在街上,就必被殺.」懶人用很多理由推托自己的良心；如果你問個弟兄上週怎麼沒參加探訪?他說:「天氣冷又下雨.」另一個弟兄沒有赴禱告會；因遠方來了親戚,久不見面,無法脫身.又有個弟兄佈道會沒有來,這麼重要關乎擴展神的國度,理應大家投入才對.他說:「因頭腦昏沉,無法參加.」屬世的事,我們並不懶惰,尚且極力爭取；屬靈的事我們更應竭力爭取.對屬世的事我們肯流汗,對屬靈的事,我們反而不肯流汗.這是今日基督徒的光景.
\newline
\newline
(五)認對職責厭煩的罪,不論任職長老,執事,主日學教員或文字聖工；接受聖職,開始時感到很寶貴很新鮮,立心盡力而為,恐怕失責,辜負主之所托；但是後來遇到困難,或受同道批評,或因自己軟弱,漸漸對職責厭煩.
\newline
\newline
曾有位醫生,不論日夜都診病,賺了很多錢；但終因勞累成疾,他躺在床上思想；人的壽命最多七,八十歲,花費有限,賺了這麼多錢作甚麼?是要我作錢奴嗎?如此為人值得嗎?從此他改變作風；週六,週日和晚上不看病,晚上參加教會查經班,禱告會,週六,週日參如教會事奉；這位醫生火熱起來,他感到不能空手見神.經過一年,大家看見他的好榜樣,就選他為長老；開始時他覺職責新鮮；不敢辜負神的委託,盡力而為.後來,有兩個執事意見不和,常常吵架；他身為長老,就出面調解,勸他們和睦,才不致破壞基督的聲譽；兩執事不接受勸誡,他就用神的話責備他們；兩位執事火氣大作,反而聯合對付長老至可憐地步；長老越想越氣,想到在教會作義工,犧牲晚診時間,少賺很多錢；為的是到教會服事,一無所得,還要受氣；因此他對職責感厭煩.
\newline
\newline
教會有許多人對職責厭煩,若然,就應認罪.保羅勸導他們說:「我們在基督裡的勞苦不是徒然的.」(林前十五58)因「......可以在耶穌基督顯現的時候,得著稱讚,榮耀,尊貴.」(彼前一7)「當耶穌再來時,取去一個,撇下一個.」(太廿四)那時我們要被提；到了廿五章論及十個童女,五個聰明的,五個愚拙的；聰明的被請入坐席,愚拙的被拒於門外.聖經很有次序,廿四章論到被提,廿五章論聰明和愚拙的童女.何以有聰明愚拙之分呢?因我們被提後,要站在基督寶座前受審判；這審判與得救無關；但與得賞賜有關,有的人忠心又良善,到那日要得稱讚,要與基督同坐席；沒得稱讚者要被拒門外；滿了七年,地上的國要成為基督的國,耶穌從天降臨,帶同千萬聖者齊來,就是那班得稱讚者；那日耶穌榮耀降臨,被稱讚者,要與基督同得榮耀尊貴,與基督一同治理地上的國度,與主一同掌權.多麼尊貴!
\newline
\newline
各位弟兄姊妹!你的勞苦不至於徒然,所以對職責厭煩者應該認罪.
\newline
\newline
我們從聖經看見一件事,復興就是神臨到我們中間；荒涼就是神離開我們,神多麼不願意長久離開人；「因為祂的怒氣不過是轉眼之間；祂的恩典乃是一生之久......」(詩卅5)主來到地上,約翰為祂鋪路；如欲主回到我們中間,我們就應當為主鋪路.認罪,就是為主鋪路；約翰為主鋪路,主就降生了；如果我們為主鋪路,主即回到我們中間.我們應禱告以求復興,禱告並非求此求彼,禱告應有認罪的內容.人若回轉如聖經所載的浪子回家,他父親給他袍子穿,戒指戴,鞋子穿,成了個新人.真正認罪者將變成新人.
\newline
\newline
真正認罪者要重新奉獻.
\newline
\newline
難道從前的奉獻不真實,不算數嗎?我從前奉獻實出於心被恩感,是沒有條件沒有保留的奉獻,把自己當作祭物擺上祭壇；從前所奉獻的祭物,是亞伯的祭物,不是該隱的祭物.舊約給我們看見,該隱的祭物神不要；因他的祭物無血,神不接受無血的祭物.亞伯的祭物神是要的；因他的祭物有血,神接受有血的祭物.血是死的意思,死了凡事都沒意見；有沒有壽衣,有沒有棺材,都無意見.奉獻作傳道,是城市或鄉村沒意見,錢多或少也沒意見,因為死了.我從前奉獻,實在是亞伯祭物的奉獻,是多麼真誠；想不到久而久之,遇了工作困難,生活貧困,世界誘惑；心志漸漸轉變,落到今日地步.所以我們要重新奉獻,恢復從前的心志,作個真正悔改的人.
\newline
\newline
真正悔改者要重新調整與主的關係,就是靈修的生活──禱告和讀經；從前「奮力禱告」(賽六四)所謂奮力,不是與神奮力而是和自己或與魔鬼奮力；因為在社會謀事,工作忙碌,到了晚禱時,腰酸背痛,草率禱告了事.勝過疲乏的身體叫做奮力.有的人禱告不專心,頭腦充滿雜念.耶穌說:「禱告要進入內室.」思念混亂怎能進入內室呢?禱告應該抓住,每作一事,不告訴神絕不罷休；但現在被疲乏勝過,禱告不蒙神答應；若落到此光景,我們必須追究原因.從前我讀經必須細讀,務求得到亮光,因有光才能打屬靈的仗；但如今我讀聖經,走馬看花.落到這光景,我應該要調整.此外,還要重新調整我們的事奉.從前羨慕聖工,事奉有喜樂,有力量,在事奉中覺得輕鬆愉快,盡量安排時間出來事奉,不許可事業,家庭影響我的事奉；但今非昔比!有了多餘時間也不事奉,從前我犧牲時間,金錢來事奉,現在不但家庭影響事奉,甚至身體,情緒影響了我的事奉.為何至此光景?因為從前羨慕聖工,善工；現在卻厭惡.若落到這樣光景,我們的事奉應當調整.最後,要調整對主的地位.從聖經看,主是第一位,舊約創世記「起初神」神是第一位.再看舊約挪亞,他從方舟出來時,第一件事不是耕地,也非蓋屋,乃是築壇敬拜神,築壇是第一位.猶太人初熟之物,羊羔不是拿回家燒,乃是拿到聖殿獻給神.新約也是一樣,七日的第一日禮拜敬拜神；因神是第一位.馬利亞也是這樣,天使先告訴她要懷孕生子,她立刻想到婚姻問題；如果懷孕生子,婚姻便破壞了；另一個是律法問題；如果未婚懷孕,律法要她的命；但馬利亞對主說:「我是你的使女,我情願照著你的話成就.」因主是第一位,婚姻,性命不是第一位.各位!從前你尊主為首,慢慢地,你的事業,兒子,享受,變成第一位了.
\newline
\newline
各位弟兄姊妹!我們希望復興,必須禱告,要認罪禱告,認靈修不認真的罪,認不敬虔不敬畏罪；認愛心有虧欠的罪,對屬靈的事懶惰之罪,對職責厭煩的罪.真正認罪的人變成新人,調整與主的關係,調整對神的奉獻,調整對主的地位.我們希望主在我們中間,就必須為主鋪一條路,這條路就是認罪.
\newline
\newline


\newpage

\section{第58屆~港九培靈會~吳勇~第四講}
\label{sec:767}
\textbf{教會的大使命}
\newline
\newline
連結: \href{https://www.hkbibleconference.org/session-message/view/767}{\texttt{https://www.hkbibleconference.org/session-message/view/767}}
\newline
\newline
\hyperref[sec:465]{< < < PREV SERMON < < <}
~
\hyperref[sec:toc]{[返目錄]}
~
\hyperref[sec:467]{> > > NEXT SERMON > > >}
\newline
\newline


\newpage

\section{第58屆~港九培靈會~吳勇~第五講}
\label{sec:467}
\textbf{教會的人才}
\newline
\newline
連結: \href{https://www.hkbibleconference.org/session-message/view/467}{\texttt{https://www.hkbibleconference.org/session-message/view/467}}
\newline
\newline
\hyperref[sec:767]{< < < PREV SERMON < < <}
~
\hyperref[sec:toc]{[返目錄]}
~
\hyperref[sec:468]{> > > NEXT SERMON > > >}
\newline
\newline
未講人才之前,讓我們先來看看世代；因為知道世代如何,方知人才如何.
\newline
\newline
我用五點思想世代:－
\newline
\newline
(一)人生痛苦的世代,我想大家都有同感.從娛樂界音樂方面來看,音樂可以表達情緒；現代在台上唱歌的,又喊又叫又奔又跳.這是唱歌者,發洩內心痛苦的表現；他們有時簡直發洩到癲狂；若唱得好,千萬人都齊聲響應；因為他把人心中的痛苦發洩無遺.我們常聽到一種心臟衰竭病,心臟如同一部機器,用得太多,性能衰退.心臟為人體輸送血液,使血液循環,此外,還負擔人和事；做人多不順利,常使人苦惱麻煩,致心臟負荷沉重；今世代人的心詭詐,不易相處,又難以應付；短短幾十年人生,重擔不時起伏.除個人心中痛苦,還有家庭的痛苦；報載離婚率越來越高,原因是很多男人講究交際,多在色情場所活動.女的講究自立,舊禮教的三從四德,於今不存在了,所以常常見家庭不睦；雖生活一起,但是都冷冰冰的；除怨嘆的話,彼此無言可說；你摔我擲,非爭則吵；因此家庭氣氛,時而沉悶透不過氣,時而緊張至雞犬不寧.此外尚有兒女問題,因失教育,學業與為人,無人管教；家庭毫無溫暖,非常乏味；為夫者視回家如畏途,為妻者覺得在家中受罪；家不成家,實在痛苦!
\newline
\newline
世界經過一次,兩次的戰爭,遺留慘痛的教訓.繼續有內戰,如南北韓之戰,南北越之戰；有外戰,如兩伊戰爭,以色列和阿拉伯戰爭；現代又加上核子戰爭的威脅,星球大戰的威脅.不久之前,科學家舉行會議,聯合國作出報告；核子戰爭一旦爆發,可能世界上3/4的人被殺；文明完全摧毀,世界退至石器時代.他們又發表一篇偉論:原子塵將遮蓋太陽,氣候極大變化；地球可能回到冰河時代,人和生物都無法生存.至於戰爭的原因,有人說是經濟問題,因貧窮沒飯吃,不得不找出路,或偷或搶.國家經濟貧窮,民不聊生,就向外侵略.這是理論.實際上,我們看見中日戰爭；中國人很窮,時遭旱災,水災；日本人自明治維新後,國家進步,漸趨富足；但是中日戰爭的教訓,並非貧窮的中國打富有的日本；乃是富有的日本攻打貧窮的中國.所以戰爭並不是經濟的問題.有人說是教育的問題.兩隻狗爭骨頭,咬得皮破血流；兩虎爭隻死牛,拼到死去活來；這些都是獸的野性.進化論者講及人的野性,還未進化到頂端,尚須教育,以達文明的地步.事實並非如此,世界許多殖民地,不是野蠻者佔據文明人,而是文明者佔據野蠻人之地.戰爭原因既非經濟,亦非教育問題；而是私慾問題.政治上看見一個名詞,你為自己的利益；我也為自己的利益,大家都為自己的利益.故此,人類就無路可走,戰爭一定無可避免.這是戰爭已迫到邊緣的世代.
\newline
\newline
(三)信仰妥協的世代,我們從聖經可以找根據；但以理書載,但以理的朋友不肯妥協；尼布甲尼撒造個大像,命令行開光禮,違命者丟入火yiu 中.但以理他們是猶大人,從聖經得到教導；除神外任何偶像都不可拜；他們寧被丟進火yiu ,也不向大像下拜,這是不妥協的例子.但以理不妥協,波斯王下令,一個月內不可向自己的神下拜,違者被扔進獅子坑；但以理不遵王令,仍然一日三次,打開窗門面向耶路撒冷祈禱,與素常一樣.這也是不妥協的例子.
\newline
\newline
王明道先生寫了一本「我們是為了信仰」的書.內容是,你們所講的耶穌是從婦人所生,耶穌的死是博愛的死,耶穌的復活是精神的復活.我們所講的耶穌是從童貞女所生,耶穌的死是為贖罪而死,耶穌復活是身體的復活.我們怎能和你們站在一起呢?他因不肯妥協,被判無期徒刑；平反後才被釋放.這也是不妥協之例.
\newline
\newline
我們看見今日信仰漸趨妥協.人有層皮,分皮內皮外.城有一溝,以溝分城外城內.我們的信仰並非皮也非溝；乃是藉聖潔分別信與不信；但現在我們有何妥協呢?論喪禮,有的人主張公祭,奠祭,然後舉行安息禮拜.婚嫁的事,如不用酒,怎謂之酒席,無酒不成敬意,遭人閒言閒語；所以一邊擺酒,一邊不擺酒.這也是妥協.
\newline
\newline
(四)教會冷淡的世代「教會是基督的身體,是那充滿萬有者所充滿的.」(弗一23)萬有者怎樣,被充滿者也應怎樣；萬有者是豐富的,奇妙的,榮耀的,光明的.既是這樣,被充滿的器皿也應該是奇妙,光明,榮耀的.但是今日教會我們所見的不是這樣,乃是可憐,貧窮,黑暗,羞恥.為何被充滿者和萬有者適得其反呢?因為今之基督教會,充滿者不是這位萬有者,乃是充滿人的金錢,學問,才幹.故此,今日教會怎能不羞恥,不黑暗呢?按照聖經告訴我們,教會應該是「神的名超乎萬人之上,神的靈貫乎萬人之中,神的兒子住在眾人之內.」可是今天的教會是人作主,所以神的名就不超乎萬人之上,神的靈就不貫乎眾人之中；神的兒子就不住在眾人之內.這是教會冷淡的事實.我們不必談過去,只談現在,現在教會講更新.羅馬書講心意更新,為神傳話,言詞不可輕浮,要鄭重.聚會應更新,主日崇拜有數百,數千人；禱告聚會只有幾十人,幾個人,氣氛沉悶至極.今日教會所講的更新,是用方法來更新；可是聖經所講的更新,惟有聖靈才能更新；然而教會不講聖靈來更新,乃講人的本事,人的方法來更新.
\newline
\newline
(五)主要來的世代,這是個大題目,我沒法講很多,只以尼布甲尼撒的像略講.
\newline
\newline
耶穌的門徒問耶穌再來時有何預兆?祂說:外邦日期滿便是了.外邦日期不但指猶大亡國,這是外邦日期滿了的開始.猶大亡在巴比倫手裡.巴比倫王尼布甲尼撒,夢見個又高又大的像,金頭,銀膀臂,銅腹,鐵腿,半泥半鐵的腳,這像代表整個外邦時代,這時代愈來愈壞；從金,銀,銅,鐵直到泥.
\newline
\newline
如今世代,也是每況愈下,我們看見:
\newline
\newline
居住環境差,空氣污染,海洋污染,生態環境破壞,愈加嚴重.
\newline
\newline
道德敗壞,犯罪技術高明,手段愈辣,人心愈狠；不但殺青年人,連嬰孩也殺.至於男女關係更不得了!羅馬書講順性的用處；但卻變成逆性的用處.有男有女,女為妻,男為丈夫,這叫做順性；但今日卻不如此,男和男行可羞恥的事；其實同性戀並非現代的東西,不過早期同性戀尚有顧忌,不敢公開；但現代同性戀不以為恥,且要求合法化,爭取社會地位；已到了不知廉恥地步,人的道德壞到極處.
\newline
\newline
政治敗壞,有土地,思想,宗教,種族的衝突.
\newline
\newline
如果讀路加十九章,主人回來時,有的僕人一錠銀子賺了十錠,有的賺了五錠.十個僕人每人一錠.馬太廿五章的記載,僕人所得的銀子有多有少.路加十九章記載,每個僕人所得的銀子都是一樣.這兩處記載的分別,馬太所載是恩賜；因為恩賜有大有小.路加所載是恩典；因為恩典大家都是一樣.馬太記載主人要算賬,路加記載主人也算賬.所以當耶穌再來時,無論恩典,恩賜都要算賬；雖然恩典白白得來,但是我們怎能徒受呢?
\newline
\newline
今天的世代需要甚麼人才呢?
\newline
\newline
(一)怕神不怕人的人才,這樣的人才能忠心傳神的話；對人不討好不奉承,才能站在神的一邊；不顧自己的得失生死,才能說神要他說的,只要是神吩咐,無一言不敢發.好像耶利米作先知時,猶大人處境危險,又被巴比倫攻打；神吩咐他傳話,不是傳奮發人心的話,而是傳打擊人心的話,他吩咐猶大人要投降,把頭置於巴比倫刀下,這樣就能存活.他所傳的話,人非但不接受且不要聽.耶利米自知傳這話的後果如何；但是他仍舊傳,因為是神所吩咐傳的話.結果,一點也不錯,猶大人打他,把他下監；到了城將破時,猶大王差人領他出監,問他「神有甚麼話,快說吧!」他把從前那些話重說；因他怕神不怕人,忠心傳神的話.
\newline
\newline
如果怕人不怕神的,不會持守真道,不敢說神所吩咐的話.有的傳道人,在教會不敢講聖潔,因教會的長老不聖潔；傳道人不敢責備罪,因教會的長執犯了罪.傳道人若怕人不怕神,就無法忠於神的話語.
\newline
\newline
(二)用舌頭舔水的人才,基甸和米甸人打仗,有三萬二千人跟著他；神說:「人太多,害怕的人可以回去.」於是有二萬二千人回去,神說:「一萬人還是太多,讓他們到河邊喝水,接受試驗.」結果,用舌頭舔水的只有三百人.神說:「用這三百人去攻打米甸人.」舌頭舔水有很多解釋,其中有一解釋；水在日常生活中很重要,代表人生的需要；有種人對人生的需要永不滿足,慾無止境.有的人適可而止,就算不夠也無所謂；人對物質的需要,若沒這種態度,就貪愛世界,魔鬼用世界來迷惑你；若你沒有這種心志,魔鬼就用貧困來打擊你.人對物質若沒正確看法,在神的工作上,你沒有用場,你是會變節的.士師時代,有人問米迦家的祭司「作一家的祭司好呢?還是作一族的祭司好?」結果,他就離開米迦；他想,作一家的祭司不如作一族的祭司好；因為一族大於一家.現在美國的職業,不講情面,不講工作年歷；講的是錢之多少.若我們對物質沒有好的看法,也如米迦家裡的祭司；這種人不講求神的旨意,只講求自己的需要.至於吃苦的問題,信徒有種觀念,以為主的僕人命定要吃苦,所以很多教會不想改善神僕的生活,作為神僕,穿破的衣服,吃青菜豆腐,才是好僕人的憑據；並且如今是勢利的時代,用錢的多少衡量人的價值；所以傳道人常被鄙視,因傳道人所得的錢都是少的.還有一種不爭氣的傳道人,不以神僕自尊,反因進款少而自卑,這種人在神的工場上沒有用處.
\newline
\newline
(三)能說能行的僕人,作主的工人,有屬靈的權柄,權柄是神所賜；耶穌派門徒出去,就給他們權柄；所以鬼服人服.(本想講鬼服的見證但沒有時間)現在講人服的見證:
\newline
\newline
有位高貴的太太,當我清晨五點多鐘正在靈修時,大哭大喊來到:(她屬別的教會)我對她說:「有話慢慢說,何必哭得這麼傷心.」她說:「我們在上海是江長川證婚的,如今我們要離婚,沒法找江長川為我們辦離婚；所以我想請你代表.」我心想,傳道人向來為人證婚,沒有為人辦離婚的.我就問她為何要離婚,她說:「我丈夫年紀比我大,已經六十多歲了.」我說:「既已從年青容忍到如今,為何年紀大了才要離婚呢?」她說:「我再不能容忍下去了,自從嫁給他第九天,他就開始作壞事,至今卅八年我一直忍著.」我內人在旁說:「我剛讀約翰福音第五章,有個患癱瘓卅八年的人,耶穌叫他起來行走.」我想,今天神指示我們,神叫卅八年的癱子走路；也都能叫你卅八年淫亂的丈夫離開罪惡.」經她一講,我膽壯了；逐即打電話約她丈夫吃飯,他表示希奇；一般都是傳道人被請吃飯,沒有傳道人請人吃飯的.我說:「誰請都無所謂,不過我今天較忙,所以就在對面飯館,準時十二點吃飯,最要緊須準時.」我就對他的太太說:「已經約好你的丈夫了,讓我們看看神可行大事；你先回家去吧!」我如約到達,他早已等著,一看見我,連忙趨前握手；他的手冰冷,談話時下巴一直發抖.我看他是個做大生意的富翁,而我不過是貧窮的小傳道人；他見我卻像見了皇帝一樣.他說:「我曉得我的老伴來見你,我現在神的僕人面前,也在神面前,我不敢少說多說,我已活了這大把年紀,來日無多；見神的日子也快到了；所以我定要在見神之前；辦清楚此事；那女人我可以給張支票打發她走.」我說:「你辦得到嗎?」他說:「今天我必辦妥,我辦妥之後,你可否陪同我回家?」我們一同吃飯,禱告,而後各自回家.到了下午四時,那男人打電話告訴我事情辦好了,立刻開車接我,希望我陪他回家.當我們進他家門,我說:「江太,我給你帶個卅八年的癱子回來；因他已離淫亂.」
\newline
\newline
如果你有權柄,人會服你,鬼也會服你.權柄是很重要的.可是神的權柄賜給能說也能行的人.神給亞當權柄修理看守,給禽獸命名；因亞當服神的權柄,才能掌神的權柄.服神的權柄,就是聽神的話並且遵行,這樣才能執掌神的權柄.
\newline
\newline
大衛是個有權柄的人,因他是個能說能行的人.何以見得?太衛禱告說:「神啊!你曾使我寬廣.」(詩四1)他的心能容小人,示每咒罵他,將軍要砍侍者的頭；但大衛說:「他咒罵我是神吩咐的.」示每不過是小百姓,大衛能容小的,當然也能容大的.掃羅屢次要將他置於死地；大衛有幾次機會報復,大衛卻沒報復；因他也能容大的.神的權柄,不隨便賜給人,只賜給能說也能行的人.
\newline
\newline
(四)被聖靈充滿的人才,使徒被聖靈充滿,講道叫人扎心,雖然他們是無知的小民,但大有膽量；他們被聖靈充滿,所帶領的教會,人人凡物公用,充滿愛心,信的人恆心遵行主的道,且有神蹟奇事隨著他們；教會不斷向前增長,短短數年,從耶路撒冷直到猶太全地,並且人才輩出；勸慰子巴拿巴,傳福音的腓利,外邦使徒保羅,都是那時一一產生.聖靈充滿的人,就是這樣的人；聖靈充滿所帶領的,乃是這樣的教會.
\newline
\newline
今之教會,神的道傳出,人無動於衷；雖有學問,但是沒有膽量；人人只顧自己不顧他人,沒有愛沒有溫暖；信者只聽道而不行道,且沒有信心,沒有神蹟也沒有奇事；十年八載的教會,毫無增長,只有萎縮,人才難產.
\newline
\newline
(五)有信心的人才,信就能,不信就不能；神的能在信者身上.
\newline
\newline
有次我有感動要到阿根庭,我到領事館辦手續；看了我的護照,說:「中國人不能去.」我說:「既有領事館設立在中國地方,為何中國人不能去?」回答說:「我沒時間和你討論.」我說:「今天我很空閒,你先辦你的事,我就在此等你.」他不理我隨即走了；等了半個多小時,他從上醉酒下來,搖幌到我跟前問「你是誰?」我說；「你這麼健忘,我剛才要辦理手續到阿根庭的.」他叫我過去,用手指地球,說:「阿根庭就在這裡,台灣在那裡,要繞地球一大圈,真難得!快把護照拿出來.」我出示護照,他翻開看說「原來是傳道人,我不好意思收你的錢,可以免費簽證.」他找印鑑,因醉眼目昏花,遍找不獲.我說:「我替你找好嗎?」他答應,我進去；他吩咐我找到了印後,自己可蓋印,我就照辦；後來他寫上幾個字,因手發抖字跡歪斜；寫完了叫我蓋上另外個鋼印,說:「全部辦好了,祝你有好的旅程.」他還客氣地送我到門口,行了180度鞠躬.我拿了簽證到了阿根庭,移民局研究半天都不明白；其中有人認出寫的字是「外交官」,立即吩咐海關免查放行,且派人為我提行李；於是我輕鬆搖擺入阿根庭國境.我頓然當了天上的外交官,通行無阻.那次我在阿根庭作了很好的工作.但各位要記得!「能」不是你的事,「信」是你的事.好像當時以賽亞,神說:「我能差遣誰呢?誰肯為我去?」能是神的事,信是我們的事,肯是我們的事.神說:「誰肯為我去?」巴不得有人願把心獻給主；但要知道,不是你能而是祂能；只要你信,祂的能就臨到你的身上；因為神的能是賞賜給信的人,賞賜給肯的人.
\newline
\newline
最後我再講個見證,我信主時,實在清清楚楚,且我信主就服事主.我們住青年會地方,每日要從樓上搬凳子下樓,大家搬得滿身大汗而後禮拜.我覺得這樣不太好,應舒舒服服禮拜才是；所以我對主說我願意奉獻每週六搬好凳子,清潔並排列.有次我去搬凳子,青年會總幹事叫我聽電話,是石油公司總務處長打來的.他說:「找到你就好了,真是急死我了；我們公司總經理弟婦祝壽,托我請位傳道人講道,我卻忘記了；下午二時開始舉行,現在十一點多了,我沒法找到人,所以找你幫忙.」我大驚失色,叫我置身於一班大學教授的場合中講道.我信主剛數月；舊約聖經沒讀過,連新約尚未讀完,我講甚麼呢?他說:「你素來熱心助人,找到你沒解決不了的事.」那時我廿多歲,給他戴上高帽,我就昏三倒四了.他叫我記下地點時間；不由分說,就掛上電話.我不知如何是好,連忙回家告訴內子,她說,你為何答應?「我說我沒有答應,」她說:「趕快進去禱告」.我說:「預備講道都來不及,還要禱告,實在沒有時間.」我連忙翻聖經,從馬太,馬可,路加,約翰,翻得頭紛亂.內子看見我這光景,對我說:「講道要先禱告.」既然聖經我翻不到,我就禱告,約有一小時的禱告.神給我一節聖經；我寫了幾個大綱,隨即出發.教授家的人問我找誰,有何貴幹?我不敢說是來講道的.金教授出來了,我們彼此不相識；我問他:「是否請石油公司的總務主任找個傳道人來講道?他從我頭上看到腳底然後說:「好啦,進來吧!」我進去,滿堂卅幾位禿頭的大教授,我越看越害怕,索性閉眼禱告,時間到了,請我傳信息.這是首次經驗；沒有一個教授不感動的.於是我得了一個經驗:「能」不是我的事情,「信,肯」才是我的事情.這個世代需要這樣的人,只問你肯不肯?
\newline
\newline


\newpage

\section{第58屆~港九培靈會~吳勇~第六講}
\label{sec:468}
\textbf{聖靈的果子}
\newline
\newline
連結: \href{https://www.hkbibleconference.org/session-message/view/468}{\texttt{https://www.hkbibleconference.org/session-message/view/468}}
\newline
\newline
\hyperref[sec:467]{< < < PREV SERMON < < <}
~
\hyperref[sec:toc]{[返目錄]}
~
\hyperref[sec:469]{> > > NEXT SERMON > > >}
\newline
\newline
從今晚開始,一連五次,我們要講及聖靈.
\newline
\newline
有一次,美國第一間大學的校長到台灣講道,那大學是注重物理,校長當然是位科學頭腦的人；他開口便大聲疾呼「今天是聖靈的時代.」現在台灣有一很特別的現象；尤其是青年人來邀請講道,便指定題目要講聖靈.我想香港也不例外；因大家都想要突破,我們已經試過,終於沒法突破；方法不能突破,講道太好也沒法突破；所以大家都想,可能現在是聖靈工作的時代,許多青年都渴慕聖靈.
\newline
\newline
除今晚之外,我講道時若你心裡有問題,我會給你答覆；所以我把許多問題都找出來,盼望可以給弟兄姊妹解答.
\newline
\newline
聖靈的問題,從聖經看,約翰和路加二人很突出；他們所講的不一樣,約翰注重在你裡面的(in you)路加講的是在你的上面,(on you)約翰所講的是用氣來作表號,那天門徒聚集在一起,耶穌向他們吹一口氣.路加所講是用風來作表號；所以行傳二章說「好像一陣大風吹過.」約翰注重生命問題；路加注重工作問題.我覺得這二者不能分開,要以生命為基礎；否則聖靈的問題常會走偏,我們應該平衡；聖經說不可偏左也不可偏右；可是我們常容易注重看得見的事,不注重看不見的事.
\newline
\newline
今晚主要經文是加拉太五章,這裡共題及九個果子,按理應該是複數；聖經用字很當心,好像說神造人,造男和造女,人字是單數的；因兩人其實是一個人,女人是從男人身上出來的；好像說亞伯拉罕的後裔如天上的星,地上的沙,但後裔是單數非複數；因後裔乃指耶穌基督.這裡所講聖靈九個果子也是單數；神是三位一體,神是父,是子,是靈；神有多方面美麗.例如聖經講耶穌基督,不是用一本福音來講；因一本福音無法表彰完全的耶穌基督,四本福音才能表達；馬太說祂是王,馬可說祂是僕人,路加說祂是完全人,約翰說祂是神.啟示錄第二章生命樹一樣,活物乃指有生命之物,活物的臉面有人,獅子,牛,鷹四種.一種臉面不能表彰生命的美；四種臉面才能把生命的美表彰出來.按理,人重生得救,應當表彰基督的榮美；可是彼得後書說:「我們既有這麼大的應許,就得與神的性情有分.」「就得」意即要追求才能將基督的榮美表彰.所以「果子」這字是單數而非複數.還有,第一個果子是愛,其他的果子也都和愛有關,所以果子用單數.林前十三章說:「愛是恆久忍耐,又有恩慈；愛是不嫉妒；愛是不自誇,不張狂......」愛和美是多方面的.喜樂是愛的情緒,家庭如有愛,情緒一定很輕鬆,溫暖,喜樂；因喜樂和愛有關.和平是愛的心境.「和眾聖徒一同明白基督的愛,是何等長闊高深.」(弗三18)我以中英文聖經對照,覺得都和原文的注解不同,原文這句話沒文法；按理保羅是文學家,為何寫的沒有文法?中英文繙譯都把原文沒文法修正為有文法；其實,忠心譯者應按照原文繙譯才對；但他們自作聰明,以為原文有漏掉的字,就加上字變作有文法.事實上,愛並非文法所能表達清楚；保羅想到基督的愛,就大聲喊叫長,闊,高,深.和平是愛的心境,廣闊無邊!耶穌被釘十字架上,按理釘十字架是很痛苦,很快就會昏過去；若昏過去便無所觀看了,所以就用苦膽和酒作麻醉劑支持被釘者；釘在耶穌兩旁的強盜都接受苦膽和酒,惟耶穌不接受；強盜所以接受乃因支持不住,但耶穌支持得住,不需用麻醉劑；因為忍耐就是愛的能力.
\newline
\newline
聖靈的果子共有九個:
\newline
\newline
「仁愛」主耶穌說:「有了我的命令又遵守的,這人就是愛我的.......」(約十四21)有愛才能夠遵守命令.聖經全都是神的命令,歸納起來,就是道成肉身.道是話,話是說明,發表；神要藉人的肉身將祂說明,發表；這就是道成肉身的意思.神非但希望人的靈魂得救,且希望人的身體給祂用,用我們的身體來說明祂,發表祂.
\newline
\newline
很多人說聖經是神的話；我不敢說聖經不是神的話,但聖經分為兩部分；一部分叫做神的話,一部分叫做人的事.聖經所記載神的話不多,記載人的事十分多；例如神對亞伯拉罕說:「把以撒帶上摩利亞山獻給我.」這是神的話,次日早晨,亞伯拉罕帶以撒上山,就綁起他,拿刀要殺他.這些記載是人的事.還有,雅各在曠野,神向他顯現,說要把他枕頭所睡之地,賜給他和他的後裔為業.這是神的話,翌日,雅各用石頭作柱子,澆油其上,叫這地方做伯特利.雅各所作的叫做人的事.整本聖經都是如此；聽見了神的話,就把神的話變作人的事.
\newline
\newline
今日教會,神的話,人的事,完全無關.聖經把神的話記載,把人的事也記載；神的話是神的話,人的事也是神的話.
\newline
\newline
耶穌說:「我心裡柔和謙卑,你們當負我的軛,學我的樣式.」(太十一29)柔和就是吞聲忍氣,逆來順受；謙卑並非柔聲細語,儀表雅觀；謙卑既非聲音也非動作,乃是沒有自己.比方在聚會中高談闊論,但人家不欣賞,不贊成；那時如你感到很丟臉,很不好意思,你就是有自己.面子被人損傷,名譽受人批評；氣憤憤在心,這樣就是有自己.謙卑就是發表了意見,人不接受,心裡不覺氣憤；面子被人損傷也不存心報復；就是沒有自己.所以耶穌教我們要負祂的軛,學祂的樣式.意即要把祂的話變作你的生活,祂的話變作你的事.這就是神的吩咐,聖靈所結的第一個果子,就是遵行神的話.
\newline
\newline
「喜樂」聖經記載耶穌曾在拉撒路墓前哭；耶穌看見人在聖殿買賣牛羊鴿子,就向他們發怒；但從聖經中,幾乎沒有看見耶穌笑.好像耶穌只會哭,會怒,不會笑；其實,聖經也幾次記載耶穌喜樂.真正的喜樂,非以聲音發表,也非動作表現.人最喜樂,莫如金榜題名時,洞房花燭夜；但是喜在心頭,樂在心裡,除非以不法手段搶奪新娘回來的惡漢；就洋洋得意,翻來覆去地哈哈狂笑.真正靈裡的喜樂是在深處,是人所難以了解的,是聖靈結出來的.
\newline
\newline
「和平」有兩個解釋,一是裡面的平安,一是外面的調和.裡面的平安,就是內裡沒風浪.
\newline
\newline
蘇東坡有首詩:「稽首天中天,豪光照大千,八風吹不動,端坐紫金蓮.」現在我們只把「八風吹不動」作解釋.八風即成敗,毀譽,貧富,興衰；無論如何心不受影響.蘇東坡寫完這詩,叫書童送到過江給朋友,讓朋友欣賞這首詩,表示他是個八風吹不動的人.朋友一看,在詩首寫了一紅字,叫書童送回.蘇東坡打開一看,非常生氣,就過江與朋友算賬；書童開門說:「主人今天不見客」,他氣憤直衝進去,看見書房外寫著「八風吹不動,一屁打過江.」原來他並非八風吹不動,需是吹到他裡外上下.這個世界,有時遭患難,遇意外,患疾病；有時遭人批評,被人誤會.你的裡面如何?
\newline
\newline
聖靈所結的果子,裡面吹不動,加以外面是調和的.屬世的事上也是講調和的；音樂,藝術都應調和.屬靈也應調和,人和神調和；神人之間的路沒有阻塞,所求的才可達神那裡；神所給的可到人這裡.如果我們與神調和,多麼有福!還有與人調和,我禱告；別人的心和我共鳴,我唱詩；你響應,這樣的教會多麼美!
\newline
\newline
「忍耐」我們實在應當忍耐；因為世上有許多想不到的事.
\newline
\newline
有次在新加坡,寇世遠先生和我租了間大戲院講道.因我是新加坡人,當地人對我有感情；就刊登我的名在報上,卻沒刊登寇世遠.有倆老姊妹見報來聽我講道,她們彼此應對:「吳勇皮膚比從前白了,因新加坡氣候炎熱晒黑了,在台灣住不那麼熱,所以白一點.現在聲音不如從前宏亮.從前年紀輕,中氣足,現在年紀大不如前了.剛才他作見證說他母親得救,然後帶領他得救；十年前他說過自己得救,才帶領他母親得救.前後十年見證不一,簡直是個走江湖的傳道人.」倆姊妹因此很生氣,散會時她們向我瞪了一眼,氣憤憤走了.我覺奇怪,到底我得罪了她們甚麼?後來那姊妹回來問我:「你不是吳勇長老嗎?剛才在講台上講道,這麼快走下來了!」我說「剛才我沒有講道,是寇世遠講道.」「哦!原來他是寇世遠.」說著,連忙叫她的姐姐.如果她不回來問個清楚,真是要把我誤會一輩子,而且一傳十,十傳百,都要看我是走江湖的了.
\newline
\newline
我出來傳道時,曾對主說:「以賽亞聽見有聲音說:我可以差遣誰?誰肯為我去呢?」以賽亞有此超奇的經歷；如果你要我出來傳道,求你也給我超奇的經歷,好叫我知道是你差派的；這樣,我遭困難才不會灰心,遇打擊才不會害怕.」那時我在鐵路局作事,有人來找我的同事不遇,我就招待他；請他說明來意,等同事回來我好告訴他.他說:「請你告訴他,某人來找過他就是了.」我送到門口,和他握別時,他忽然大聲喊提摩太後書二章15節.我覺奇怪,懷疑他有神經病,我問說:「你為何喊這聖經章節呢?」他恍然地說:「我真該死,我怎會懂得勝經呢?」他隨即走了.我心裡不平安,進內打開該處聖經:「你當竭力,在神面前得蒙喜悅,作無愧的工人,按著正意分解真理的道.」神藉著不信者,念了這節經文給我；那天我沒法上班,我清楚被神呼召出來傳道.但是,各位!你要忍耐,否則你蒙召清楚也會中途變節.我出來傳道,有位老太太很照顧我；她想年青人應該培養照顧；所以直到今日我才能夠體貼年青的傳道人；因我向這位老太太學習.她的兒子是汽車公司總經理,每三天必換套新西裝；老太太把她兒子的西裝送來給我；可是他身高六尺多,我只五尺七吋.穿上去衣袖蓋過我的手,衣襟到了膝蓋.我請內人盡量想辦法改短,因為西裝是英國料子,不要實在可惜；所以我穿的西裝,上部的口袋很高,下部的口袋又小又低.當晚,我穿上改造的西裝到日本,接機的弟兄姊妹立即要帶我到西服店定做兩套衣服,我說:「不用做,我有很多衣服.」問「都是這樣的嗎?」我說:「是的,全都是英國絨料.」回答說:「料子是好的,可是樣子太看不過去.」以後,老太太更關心我,連她兒子的皮鞋也送來；因為我腳特小,穿上尚可容三隻手指的空間；但是外國皮鞋,不穿太可惜,只好盡量綁緊帶子,總算可以步行；有一天,領家庭禮拜,我快跑進巴士,人倒是上車,鞋子卻沒跟我上去；管車人見狀,叫司機停車讓我拾鞋；當我穿回鞋子時,車上的人都集中注意力在我身上.有個人說:「此人的鞋非借的就是偷的.」真叫我臉紅如血,無地自容.教會又傳說:「我們穿的是本地貨,他穿的是英國貨；我們的鞋子是土產,他穿的是洋產BALLY牌的,我們的鞋幾十元,他的鞋子千多元；背十字呆受苦,言之容易,實行卻不易.」這話傳來,我忍耐不住,召開教會全體大會,準備辭職.有位八十多歲的老人,星夜由南部趕來北部找我.說:「你要忍耐生產的痛苦.」我因此清楚蒙召,我差一點就不能忍耐,就會變節改行.聖靈所結的果子,教人應該忍耐.
\newline
\newline
「恩慈」是人得罪你,毀謗你；你能夠恩待得罪虧負你的人；沒有氣憤,報復；始終忍耐.可是一旦那得罪你的人被汽車撞了!你說:「天有眼,我不算賬,天向你算賬.」這樣就不叫做忍耐；你若能帶些水果或金錢幫助得罪你者,這就叫做恩慈.聖靈不但叫你結忍耐的果子,還叫你結恩慈的果子；善待那惡待你的人.
\newline
\newline
「良善」青年人對耶穌說:「良善的夫子阿!」耶穌說:「除神之外沒有良善的.」既然人沒有良善,那麼,良善何解?從耶穌身上看良善.有一次,猶太人捉到個犯姦淫的女子,帶到耶穌面前；說:「這女子犯罪時,我們捉到她；按猶太人律法,應用石頭打死她,你以為如何?」耶穌蹲在地上寫字,起身對他們說:「你們誰是無罪的,可以拿石頭打死她.」於是那些人一一走了.耶穌對那女人說:「我也不定你的罪；從今以後不要再犯罪.」由此可以看見兩種人,猶太人按著律法行,他們的心是律法的心；耶穌的心是恩典的心.律法之心,看這個不好,那個不對.恩典之心,不敢說這個不好,那個不對；我們應時常保持恩典的心,以律法之心對自己,以恩典之心待別人.當約翰下在監牢時,他叫人去找耶穌,說:「我們等候的是你嗎?」本來他說耶穌是神的羔羊,給耶穌解鞋帶都不配.如今說:「現在等候的是你嗎?」這是約翰對耶穌的懷疑,不滿.意思是說,如果真的是你,為何我在監裡不來救我,不來看我呢?各位弟兄姊妹!人的靈性多好,遇見苦難,有時也會軟弱!約翰軟弱之時,連耶穌都被他懷疑,被他不滿.然而耶穌說:「你們到曠野看甚麼?看被風吹動的蘆葦嗎?到皇宮看甚麼?看穿細麻衣服的人嗎?」蘆葦又長又高被風東吹西倒,沒有骨氣；約翰並非沒骨氣的人,穿細麻布衣服是養尊處優的人；約翰也非這樣的人.耶穌反而稱讚對他不滿並懷疑他的人.這就叫做良善.良善是聖靈所結的果子,我們做不到；聖靈在你身上,教你能作到.
\newline
\newline
「信實」神是信實的,祂以信實為糧.「凡祂所說的話都是是的,都是阿們的；祂怎樣說,事情就怎樣成就.」(林後)信實即守信決不食言,凡所應許的事必定照辦,不論環境許可與否一定要辦；對人的承諾,不論有理無理,都不敢推辭.神是信實的,我們也應當信實.
\newline
\newline
「溫柔」溫柔並非沒原則,沒有個性,也並非沒有人格,沒有目的,但目的並非為自己的企圖；不是討自己喜悅,目的為著討神的喜悅；為了神的喜悅,遇羞辱罵都能忍受.
\newline
\newline
「節制」必須要能約束,能控制,又安靜在神面前.若非進到內室是見不到主的.言語,舌頭要約束,是就是,不是就不是；不說閒言,也不說謊言,大言.工作有節制,不浪費,記念別人的需要.節制就是既無不及,亦不會過之.
\newline
\newline
聖靈所結九個果子,都是有表現的,都是關乎生命的.
\newline
\newline
生命和恩賜的問題,路加和約翰都說聖靈的事；約翰講生命,路加講恩賜.你說哪個要緊呢?如果只能選一樣,我寧取生命而捨恩賜；因為沒有恩賜不要緊,沒有生命就很嚴重.當然二者兼有是最好的.若無生命,恩賜不能定準.當摩西出來被神所用,他手上的杖丟在地上就變成蛇；法老術士的杖,照樣也能變蛇.摩西的杖點在水裡,河水變成血；術士也能照做.若以恩賜分辨是非真假,是靠不住,是很危險的,故此我們要以生命為基礎!
\newline
\newline
方言和神蹟的問題,聖經記載有方言,有神蹟.何為純正的信仰,若聖經有記載的,就可接受；聖經沒有的就不應接受.論到方言,人便說是靈恩派；其實,靈恩派有何不好呢?
\newline
\newline
各位弟兄姊妹!從今開始,一連幾天,我要解答很多問題；我們需要聖靈工作,沒有聖靈；我們不能突破,不能更新.求主讓我們渴慕聖靈活在我們身上,好叫基督的榮美活出來；這是我們一切的基礎.求主把這話放在我心中,叫我羨慕基督的榮美.
\newline
\newline
創世記二章有棵生命樹,啟示錄末章也有棵生命樹；生命貫穿全本聖經.創世記在河兩邊有黃金,寶石,珍珠.意即要我們的生命,如同黃金的寶貴,如同寶石般的光輝,如珍珠般的尊貴；就是神在我們身上.
\newline
\newline
恩賜是要緊的,但生命比恩賜更要緊.我們務要在此世代,讓世人看見生命的樣子；好叫人佩服,叫人羨慕.
\newline
\newline


\newpage

\section{第58屆~港九培靈會~吳勇~第七講}
\label{sec:469}
\textbf{聖靈的能力}
\newline
\newline
連結: \href{https://www.hkbibleconference.org/session-message/view/469}{\texttt{https://www.hkbibleconference.org/session-message/view/469}}
\newline
\newline
\hyperref[sec:468]{< < < PREV SERMON < < <}
~
\hyperref[sec:toc]{[返目錄]}
~
\hyperref[sec:470]{> > > NEXT SERMON > > >}
\newline
\newline
耶穌曾說祂不去,聖靈就不來；祂去了聖靈就來.聖靈既來了,我們要追求才能經歷神.我們知道,我們認識耶穌,就追求耶穌；可是今天有個怪現象,很少人認識聖靈卻追求聖靈；頭腦有很多關於聖靈的知識,生活上一點都沒有經歷.
\newline
\newline
十幾年前我被邀請去講道,要我講聖靈的題目；那時我裡面空空,腦子空空,很少關於聖靈的知識；身體也空空,很少有聖靈的經驗.我知道要講這題目,知識和經歷二者缺一不可；我只好博覽群書,逐一收集點滴.很多書寫的都是參考聖靈來的；人家可以講,我也可以講；然而經歷則非如此,人家的經歷,不一定是我的經歷；雖然我也有一樣的經歷.我開始傳道時,從聖經找到根據:「有揀選,一定要有裝備.」
\newline
\newline
出埃及記第三章記載摩西被揀選,第四章記載他被裝備；且利未記第八章有「承接聖職」四字,承接聖職要舉行禮節,就是摩西要被抹油.在新約也找到憑據,保羅在大馬色蒙召,並非蒙召就可以工作,而是必須裝備方可；神差亞拿尼亞,就先要裝備他.我從聖經找出許多根據,我既蒙召,我要工作,我必須裝備.我剛蒙恩,思想簡單,我上山禱告,有時竄進防空洞禱告,有時竄入草堆禱告；有一天,神給我經歷聖靈,這經歷現暫不提.後來我傳道,發現有件很希奇的事,當我領唱短詩時,有人前來悔改.神給我看見,不是講道給人悔改；乃是神自己叫人悔改.不是你的工作,是神自己工作；所以我說:「主啊!我有過這樣經歷；可是我聽人說,昨日的經歷不能作為今日的資本；因神的恩典每日是新的,不能單靠從前的經歷,應有現在的經歷.」我一直為此事禱告.有一天,我在鄉間講道,因患感冒聲音嘶啞,毫無氣力；和兩位瑞士人一同吃飯,我實在吃不下,就回房間休息；當我經過講台,忽打冷戰,有涕無聲；我想,今晚我怎能爬上台呢?於是入房,自然而然跪下問主:「我怎麼辦?」辦字還未出口,神給我首次說方言,我又驚又喜；那時會場已開始唱短詩,我在方言中一直與主交通.我對主說:「外面等我工作,我必須停在這裡.」當我進入會場,三百多人齊站立；我受寵若驚,我算得甚麼,他們居然對我必恭必敬.我上講台,聲音恢復了,力量也來了；會畢病人要求禱告,手踫病者,病得醫治.我滿心感謝主!我繼續帶領聚會,勢如燃燒,如水沸騰；散會之前有見證會,有個知識份子,他看見數百人中這個哭,那個哭；不明到底是何玩意.他想出來講攻擊的話,他的臉色難看,用鄙視的眼光看會場；突然間,神的靈在他身上,他跳躍,且跳的很高,方言從口而出,人完全改變.這是那次會中所見.於是我想聖經的事情,神要以利亞找以利沙,把衣服披在以利沙身上；不是以利沙要的,是以利亞作的.我想,那個要講鄙視言語的人；神的靈降在他身上,不是他求的,乃是神自己給的.「......何況天父,豈不更將聖靈給求祂的人麼?」(路十一13)我看見「更」字,我很高興；因為不求祂都給,求祂豈不更給嗎?
\newline
\newline
今晚,我只講兩點,第一,我要答覆很多的問題.
\newline
\newline
在行傳中,耶穌叫門徒要等候.若耶穌還沒有去,聖靈就不會來.叫他們等候是有道理的；但耶穌已去,聖靈已來；還叫他們等候甚麼呢?仔細來看,不論自然界或靈界,每事有時間,有次序.在自然界,一粒種子在地裡,必須等候；發芽長枝葉,開花結果,都有次序,有時間.生物界也如此,母體懷胎,嬰兒的四肢,五官都看不見；過些時候,四肢有了,五官也明了；到了二百七十天就生下了.有時間也有次序,靈界也是一樣；神造天地,並非一口氣呼成,乃是說有光,就有光,這是一天.然後把水上下分開,這是第二天.都是有次序,有時間的.所以神要我們等候.
\newline
\newline
我想,有兩個問題要等候；等候真理的問題解決,等候攔阻我問題除去.聖靈的問題,有很多真理的問題.「......不多幾日,你們要受聖靈的洗.」(徒一 5)許多人對聖靈的洗有個講法:「重生就是聖靈的洗,聖靈的洗就是重生.」大概現在很多人這樣主張；因為當時耶穌說:「過不多日你們要受聖靈的洗.」「你們」乃指使徒；難道使徒沒重生嗎?聖靈的洗和重生,不能合為一談.
\newline
\newline
有聖靈為甚麼還要追求聖靈呢?這樣的說話很多.以弗所書一章說,信就可以得著所應許的聖靈.如果沒有聖靈,你就不會認罪,不會承認耶穌是主,不會認神是父.你認罪,信耶穌是主,神是父；這就是你有聖靈的憑據.有聖靈還要追求聖靈纔合邏輯呢?神並不在邏輯裡面；如果神限於邏輯之內,就不謂之神了.我們姑且不管這些理論,從耶穌基督身上來看,祂有聖靈還要追求聖靈.當天使找到耶穌的母親馬利亞,說:「你要懷孕生子.」她立刻說:「我還沒有出嫁.」天使說:「因為你懷孕是從聖靈而來的」所以祂有聖靈,祂要開始天國的事業；祂在約但河受洗聖靈降在祂身上,好像鴿子一樣.路加四章說祂被聖靈充滿,滿有聖靈的能力.祂有聖靈又充滿聖靈.使徒也是如此,當耶穌釘在十字架以後,他們非常害怕,躲在房子裡；耶穌進去,向他們吹一口氣,說:「你們受聖靈.」到了行傳一章,耶穌說:「過不多幾日,你們受聖靈的洗.」所以,耶穌有聖靈,使徒有聖靈,還要追求聖靈,這是有根據的.
\newline
\newline
有人曾說,五旬節以前,聖靈在耶穌裡面,五旬節以後,聖靈是在身體裡面；身體就是教會.所以行傳二章,使徒被聖靈充滿,使徒是猶太人.行傳十章使徒在哥尼流家裡；他們是羅馬人,是外邦人；聖靈在他們身上.教會包括猶太人也包括外邦人；既然聖靈已經進入猶太人身體裡面,也進入外邦人身體裡面；聖靈的工就成了,聖靈的事就完成了歷史.因此,今日不必再談聖靈的洗.但是「因為這應許是給你們,和你們的兒女,並一切在遠方的人.」(徒二29)「你們」是一代,「你們的兒女」又是一代.聖靈的洗,沒時間限制；既有遠方的人,當然也有近處的人；聖靈的應許,沒有空間的限制.
\newline
\newline
希伯來書說:「各樣的洗禮.」有很多人說是指滴禮和浸禮.因為以弗所書說:「一信一洗.」所以有人只能指一樣,不能指兩樣而言.滴禮是洗禮,浸禮也是洗禮；各樣洗禮乃指水洗和靈洗.我們從聖經看,基督徒不僅是器皿,而且是器具.杯子是器皿,用來盛茶；碗是器皿,可以盛飯.基督徒是器皿,神要把祂的靈盛在裡面.我們也是器具；筆是器具,腳是器具,可以使用；所以基督徒是器皿也是器具,是信徒也是門徒；基督徒聽道還要學道,有了水洗,還要靈洗,這是有聖靈根據的.所以等候真理的問題要解決,否則你的攔阻就一直存在.等候我們攔阻的問題要解決.我們有許多攔阻的問；害怕攔阻問題,欲追求聖靈,卻追求到邪靈；我看見了多次,有人追求到邪靈的事.遇到這樣的問題該怎麼辦呢?路加十一章說,兒子求餅,父親不會給他石頭；求魚,不會給他蛇.人向神求聖靈,難道神給邪靈嗎?父親都不至於如此,何況我們的天父呢!既然這樣,為甚麼有人求聖靈反而得著邪靈呢?當西門看見使徒們按手,人就被聖靈充滿；他們吩咐鬼,鬼就被趕出去.西門看了覺很希奇,認為這是得利的門徑,就給錢使徒,請使徒把這恩賜給他.使徒對他說「你的銀錢和你一同減亡吧!」由此可見,西門居心不正,他所求的不是聖靈,而是自己的利益.求聖靈反得邪靈,問題在於居心不正；他所求的是聖靈以外的東西,所以害怕是多餘的.
\newline
\newline
理性的攔阻問題.這問題很麻煩,我們所信是頭腦想得通的,不能相信頭腦想不通的.我的毒瘤沒有取出,至今已卅六年了.有位高官為人很孝順,他父親患的和我一樣的結腸癌；他聽說我的病經已痊愈,就和我連絡來訪問我,問我治病秘方；我說我的秘方就是神蹟,他始終不相信；因他的頭腦想不通.知識份子只能接受頭腦解得通的事,這時對於理性是個很大的攔阻.基督教已漸變成早期的猶太教；猶太教有祭司有文士,祭司負責聖殿禮節的工作.
\newline
\newline
現在基督教也漸落進禮節儀文裡面.靈裡的需要,儀文不能解決,於是就落到文士裡面；文士本是研究聖經的.今日基督教研究高深學問；我並非說這些不好,不過這些未能滿足靈的需要.今日基督教走情感路線,大喊大叫,大奔大跳,發洩內裡的情感；我從前很不以為然,覺得這樣很不斯文,不成體統；但後來我想通了,我並非贊成,而是我發現神乃重感情的.何以見得?當耶穌到拉撒路墓前,祂哭了；還有,耶穌在拿因地方看見寡婦哭送已死的兒子,就動了慈心.感情是心裡面的東西,如你不敢用感情,就等於不敢用心；人若不敢用心,怎能接觸神呢?昔日大衛迎接約櫃時,在約櫃前唱歌跳舞；他的妻子米甲因此輕視他,說以色列的君王在百姓面前露體.後來米甲被詛咒,她終身不能生育.這是理性的攔阻.我並非說理性不好；但別以理性攔阻聖靈的事.
\newline
\newline
「人怕人」的攔阻問題.彼得和外邦人同吃飯時,看見猶太人來了立即走開；後來保羅看見他在此事上不對；因為福音是為猶太人也是為外邦人的；彼得看見猶太人來就避開,這和真理有所抵觸；所以保羅當眾指責他.
\newline
\newline
各位!關於真理的事,不必怕人；接受耶穌是一部分,接受聖靈也是一部分；接受耶穌又接受聖靈,才是全備的救恩.現在的人都喜歡一半一半的,從前的人也是喜歡一半一半；神叫亞伯拉罕從米所波大米出來；他走一半到吾珥就停下來,然後才繼續走下去.雅各從巴旦亞蘭出來,走到示劍就停止；他也是走了一半.還有老底嘉教會好像溫水半冷半熱,也是一半一半.各位弟兄姊妹!屬靈的事,一不作,二不休；要走就當走到底,不能只走一半.所以神等候你的問題解決,也等候你攔阻的問題解決；這是在追求之前.你在追求之時,神也在等候,等候你發覺自己軟弱無能；我們無論對自己,對整體,都何等沒有能力,在得勝上何等沒有能力!當雅各從巴旦亞蘭回迦南地,代表他走上追求的路；他需要先解決的,是他和他哥哥的問題；因他們之間有仇恨.「弟兄和睦同居......好比那貴重的油澆在亞倫的頭上......」(詩一三三1-2)問題解決了,膏油才能澆灌下來.當雅各經過以東地帶,靠近他哥哥住處時,他邊走邊跪七次,七代表完全；向他哥哥俯伏下來,直到饒恕,仇恨消失為止.但是我們常發現,一而再,總不能成功就不作了.在得勝上何等沒有能力!
\newline
\newline
人生的攔阻問題.人的一生有許多遭遇,我年青時曾經被日本人監禁,曾經患癌症；如同經過死亡的邊緣谷.人生有很多挫折,很多打擊,須有能力才能得勝.
\newline
\newline
我曾經探訪一位老將軍,他要我為他禱告,他對我說:「你禱告,神必垂聽；請你為我禱告神,叫我快快死.」我問他為何要這樣禱告?他說:「如果你這樣為我禱告,我今晚睡覺,明早就不起床了；也許我沒這樣福氣,明早起床後我就腦溢血離世；若再沒這樣福氣,我明天上街必被車撞死.」我問他為何要這樣禱告?他說:「弟兄,我問你,你豈忍心叫我過這樣的日子嗎?」他一家三口,女兒因婚姻挫折,精神分裂,崩潰；他的妻子原是小兒科醫生,因被電單車撞傷,腦子不清楚.他現年76歲了,要照料神經病的女兒和神志不清的妻子.他在得勝苦難的事上,是多麼沒有能力.我們的事奉多麼沒有能力；和他人配搭事奉,發現自己個性特別,脾氣暴躁,心胸狹窄；遇事一觸,臉紅聲厲；見人態度不好,立即吃不消；聽人說句逆耳的話,立刻受不了；雖曾受過教導,理應溫柔,應忍氣吞聲,逆來順受,道理雖懂；實際卻做不到.我們生命上沒有能力.
\newline
\newline
「我們堅固的人,應該擔代不堅固人的軟弱；」(羅十五1)我在福建時,有次出門,到沒有車的地方,我就步行；但我內人走不動,只好坐轎子.中午,轎夫停下吃飯,他們和賣飯的人談話:「所托的事辦了沒有?」「辦好了,請放心!」「堅固嗎?」我很好奇,到底所托何事邊問堅固嗎?因此,我問賣飯的,他說:「托我找配偶.」我說:「為何不問漂亮賢淑卻問堅固與否.」回答說:「他要的是個挑水,劈柴,養豬,耕田甚麼都能幹的堅固人.」「堅固」原文甚麼都能之意.但是我們發現甚麼都不能,攻克己身不能.所謂攻克己身,就是被毀謗,面不改容,心不跳,手不冰冷；被人稱讚,不洋洋自得,不得意忘形；我們謙卑不能,溫柔也不能,樣樣不能.我曾說過,能不在你,神要「從嬰孩和喫奶的口中,建立了能力......」(詩八2)弟兄姊妹!嬰孩是軟弱的意思,耶穌因軟弱被釘十字架,你應該與祂一同軟弱.聖經中有兩種軟弱,一是不好的,好像「心靈願意,肉體軟弱.」還有一種軟弱的,耶穌因軟弱被釘十字架,這種軟弱的意思,是祂對神不頂撞,不抵抗；凡出於神命令,祂完全順服.神的能就在這樣的人身上.可是我們發現了件可怕的事,我們甚至連軟弱都不能.固然不錯,神的能要能在嬰孩身上,嬰孩非常軟弱,任你抱,任你摔,毫無抵抗；但是連這一點我們都無能,這是我們的光景!我們需要聖靈的能力,教我們對神軟弱.若非聖靈工作,我們連這都做不到.
\newline
\newline
我們還有整體的需要.今日教會,雖喊叫「增長」美名,實際卻是退後；去年一百人,今年僅九十人,可能明年只有八十人.還有種是停住的,十年前後,同樣一個禮拜堂,同樣這些人.真是越看越膽寒；因這班人越來越老,教會老氣橫秋,我們對於整體也無能.
\newline
\newline
整體各有需要.個人的需要對神就是對人；對人就是對神.耶穌曾說:「我餓了,你們沒有給我吃,我在監牢,你們沒有來看我.」(太廿五42-43)他們說:「我們沒有看見你餓了,也沒看見你在監牢裡.」耶穌說:「你們既不作在這弟兄中一個最小的身上,就是不作在我身上了.」所以,你對神如何,乃是決定你對人如何.你應按聖經的教導對待人,但約翰十六章說:「只等真理的聖靈來了,他要引導你們進入一切的真理.」若非聖靈的工作,雖然腦袋知識一大堆；但生活的實際毫不見得.
\newline
\newline
今天我從《華僑日報》看見一篇攻擊基督徒的文章:「早知今日,何必當初.」論到人來信耶穌之前自由自在,信了耶穌後很不自由；一次禮拜天不到教會,牧師便追問,實在令人厭煩.
\newline
\newline
弟兄姊妹!基督徒應該進入實際,真理不是一些知識,乃是從實際活出來的；要等真理的聖靈來了,才能幫助我們進入真理.現在有許多人很苦惱,知道饒恕而不饒恕人；知道謙卑但不能謙卑,知道當順服神但代價太大無力支付.個人問題需要能力,整體問題也需要能力,處處需要能力.
\newline
\newline
我覺得基督徒的苦惱,不是知和不知的問題,乃是能和不能的問題；為何不能?因沒有能力,為何無能力?因對聖靈的問題,沒有清楚認識.「我來要把火把丟在地上,」(路十二)呂掁中牧師的聖經新譯:「他盼望能夠燒起來.」甚麼都需要火,昔日在聖殿,有火才能燒香,點火,獻祭.今日在家庭也需要火,才能燒飯,燒水.工廠也需要火,才能製造產品.我們為主作出口的也需要火,沒有火,我們只有些字句；有火才有經歷.唱詩也需要火,沒有火就沒有聲音,不能感動人,所以實在需要火.
\newline
\newline
火就是聖靈「......何況天父,豈不更將聖靈給求祂的人麼?」(路十一13)「更」意思是你沒求,祂都賜給你；你求,祂更賜給你.
\newline
\newline
求神給我們看見自己個人的需要,教會整體的需要；追求聖靈得著能力.但神在等候你真理的問題弄清楚,等候你攔阻的問題挪開；因為很多人對真理尚未清楚,很多人仍是攔阻重重；特別是理性的麻煩.
\newline
\newline
求主讓我們渴慕聖靈,其非人的才能,口才可以突破這個世代,唯聖靈可以突破
\newline
\newline


\newpage

\section{第58屆~港九培靈會~吳勇~第八講}
\label{sec:470}
\textbf{充滿聖靈}
\newline
\newline
連結: \href{https://www.hkbibleconference.org/session-message/view/470}{\texttt{https://www.hkbibleconference.org/session-message/view/470}}
\newline
\newline
\hyperref[sec:469]{< < < PREV SERMON < < <}
~
\hyperref[sec:toc]{[返目錄]}
~
\hyperref[sec:471]{> > > NEXT SERMON > > >}
\newline
\newline
今天晚上,主要的經文在以弗所書第五章.未講聖靈充滿之前,我要和昨晚一樣答些問題.
\newline
\newline
聖靈的洗,聖靈充滿,這兩件事是否一樣?行傳一章說:「你們要受聖靈的洗.」第二章說:「他們被聖靈充滿.」應洗的是受聖靈的洗,應充滿的是被聖靈充滿,為何第一章講「洗」第二章講充滿?若說一樣,為何不一致說洗或一致說充滿?人對這兩問題有很大的疑問.關於聖靈的洗,在行傳一章提過之後,就用澆灌代之；所以聖靈的澆灌就是聖靈的洗.還有「聖靈充滿」一名詞,是否即聖靈澆灌?聖靈澆灌目的在乎聖靈充滿,澆灌乃為充滿；聖靈澆灌可說是手續,聖靈充滿可說是目的；手續乃為目的.比如說買張票看球賽,買票時有人問你買票作甚麼,你說要看球賽；買票為了看球賽.追求聖靈充滿,追求聖靈澆灌；澆灌就是為看充滿,充滿澆灌,可說是同一件事,也可說是兩件事；要充滿就得被澆灌,被澆灌才能充滿.所以這三個名詞,洗是第一次,第二次用澆灌,澆灌為達到充滿；手續是為目的.
\newline
\newline
要充滿必先倒空,倒罪,倒世界,還要倒自己.腓立比第二章說,耶穌倒自己,從寶座下來馬槽,祂把權柄倒出來,祂才不從十字架下來；如果祂從十字架下來,救恩就不能完成；祂把性命也倒出.腓立比三章保羅說:「我只有一件事.」難道保羅有一件事嗎?他是希伯來生的希伯來人是法利賽人,是迦瑪列的門生；這裡已有三件事.講祂是希伯來人,這是他的血統,講他是法利賽人,這是講他的宗教；講他是迦瑪列的門生,這是講他的教育.保羅把他的血統,宗教,教育都倒了；只有一件事就是耶穌.有人說,如果能夠倒自己,就不需要被聖靈充滿；因為倒不了自己,才要被聖靈充滿,這問題該怎樣解決?應先倒空而後充滿,或先充滿而後倒空呢?我們天然的思想應先倒空然後充滿,但神不受天然的限制.日頭往前走是自然現象,住後退是反自然現象；但神不受天的限制.鐵桶的水沉下去是自然現象,浮上是反自然的現象,所以神是不受自然的限制.既然如此,到底先倒空後充滿,或先充滿後倒空?我想用浪子來作比方,浪子回家,是先洗淨自己,然後父親親他；還是父親先親他而後才洗乾淨?神的心非同人的心,按常情,應洗乾淨然後親嘴；但父親等不及他洗淨換衣,就抱住作連連親嘴；如此看來,我們骯髒也不要緊,污濁也無所謂,我們就追求聖靈充滿.
\newline
\newline
各位!聖殿中有很多物件,除了洗濯盆之外,各物皆有尺寸.猶太人要點燒香之時,必須洗淨手腳,為何洗濯盆沒有尺寸；因為洗濯盆是聖潔的.關於聖潔的事,並非沒有標準；人以為聖潔的,神不一定以為聖潔；我們從聖經記載浪子的故事,當浪子回家時,對他父親說:「我得罪了天,也得罪了你.」浪子雖然渾身骯髒；但他有潔淨的心志,心志是神所重視的.曾有個宣教士死了,義女把他埋葬後清理遺物；發現一小箱裡有封她六歲時寫給義父的信,事隔卅多年.六歲時的信,字體既不好,字句也不通順,為何宣教士看作寶貴呢?因為寶貴不在字體的字句,乃在寫信的心.照樣,神對我們也是如此,神所器重的是你乾淨的心,神看為寶貴.所以,先倒空或先充滿,此事很難解清楚；但是潔淨的心乃是後有的.這是我對此問題的答覆.
\newline
\newline
有人說,被聖靈充滿是看果子,聽方言.到底果子或方言?兄弟對此問題的看法,我們追求非果子亦非方言,乃是追求聖靈；聖靈給我們甚麼,這是聖靈的權利,非我們所能干涉的；我們不是追求果子,好叫我們生活更榮美；我們也非追求方言,叫我們的感覺更享受；這些都不是我們追求的目標,我們追求的是聖靈,追求目的為榮耀主；憑據到底是果子或是方言,我覺得都不是問題；因你追求的並非這些,神給你果子,你感謝主,給你的方言,你也感謝主.
\newline
\newline
至於方言也有問題:「他們就都被聖靈充滿,按著聖靈所賜的口才,說起別國的話來.」(徒二4)如果說的不是任何國家的話,怎麼能說是方言?這話不錯,因為聖經上的方言是種代表國家的方言；如果讀林前十二,十四章,另有種方言聽不懂的；按呂振中的繙譯,那處的方言,他用「卷舌頭」代之「方言」.有種方言是聽不懂的,這在聖經是有根據的.
\newline
\newline
「聖靈的洗」或「聖靈的浸」這四個字,好像一條船沉下水,你說「船浸了」意思是海水充滿船；我們受聖靈的浸,就被聖靈充滿在我們裡面.我們就從中向祂發出讚美,可能是聲音的讚美.所以,聽得懂的方言有之,聽不懂的方言也有,這是聖經明文的根據.
\newline
\newline
還有個問題,當被聖靈充滿之時,「忽然從天上有響聲下來,好像一陣大風吹過......又有舌頭如火焰......落在他們各人頭上.」(徒二3-4)「正是先知約珥所說的.」(徒二16)意即五旬節所發生的正是和約珥所發生的一樣.可是約珥所發生的現象,有異象,有異夢；五旬節發生的卻沒有,所以五旬節的現象,和約珥的現象完全不同.路加寫到這裡,用「正是」乃是說五旬節約珥的現象不一樣,但所發生的事正是一個.
\newline
\newline
今天人用自己的經歷來衡量別人,你被聖靈充滿,你把經歷告訴我；如果經歷不同,你說這不是被聖靈充滿.還有人竟然以別人的方言和自己不同,說這不是被聖靈充滿.人是怪的,想出來的問題也是怪的.不應該用自己的經歷去衡量別人的經歷,這樣是不合乎聖經的.
\newline
\newline
「不要作糊塗人,要明白主的旨意如何.」(弗五17)主的旨意就是被聖靈充滿.人被聖靈充滿,就與神更加親密,與神更近；真理的聖靈來了,要引導我們進入真理.上句末了的真理兩字,應繙為實際,就是說,你本來有真理知識,被聖靈充滿後,聖靈就把你真理的知識帶進真理的實際.這是多麼寶貴!人被聖靈充滿,就加入了天國的行列.被聖靈充滿,有許多美的現象；撒旦怎能放心呢?四福音講「神在一個肉身顯現.」行傳講「神在許多肉身顯現.」行傳所說的乃指耶穌,四福音所說的乃指聖靈.神在一個肉身上顯現.魔鬼已經受不了；神在許多肉身顯現,魔鬼怎吃得消呢?所以魔鬼千方百計來攻擊你.許多人被撒旦攻擊就糊塗了.我們今天的糊塗有兩種人,一是盲目批評者,他們死氣沉沉.一是盲目追求者,他們狂狂亂亂.這是今時代最普遍的現象,這兩種糊塗人,我們都不能接受,也不能贊同.我曾經看見一件有趣的事,有對夫婦到了火車站,車已開走；他們問站員:「八點零一分的火車,到底幾點鐘開?」站員莫明其妙,答「八點零一分的車就是八點零一分開.」那人又說:「現在我的錶七點五十九分,我太太的錶七點五十八分,站上的鐘是八點零三分,到底我應該看哪個錶呢?」站員說:「火車開行的時間不是根據任何人的錶,乃是根據站上的鐘點.」弟兄姊妹!我們不要作糊塗人,聖靈充滿的事,並非根據眾說紛紛,乃是根據聖經所講的.但是有許多人,聽這信這,聽那信那；莫衷一是,結果糊塗了.聖經說:「不要作糊塗人......不要醉酒......乃要被聖靈充滿.」(弗五17-18)這裡用醉酒比喻聖靈充滿.當使徒在耶路撒冷講道時,那些看見聽見的人就說他們被新酒灌滿了.被聖靈充滿的現象,就像被新酒灌滿的現象；醉酒者在酒性控制下,說話,態度舉動都異常；照樣,被聖靈充滿者在聖靈管制之下,說話,態度,舉動都異常.耶穌基督是個被聖靈充滿的；祂講道時,猶太人說,祂說話不像文士.因為祂異於文士,就是異常的意思.還有,祂的工作也異常,五餅二魚能給五千人吃飽.祂所活的也是異常,當祂背十字架上各各他,有些婦女在哭；耶穌對她們說:「不要為我哭,要為你們和你們的兒女哭.」耶穌已極度痛苦疲乏之時,還顧念別人.
\newline
\newline
各位弟兄姊妹!我們是老舊的人,我們不異常；甚至我們不正常,反常,我們說話反常；把是說成不是,做事反常；惟利是圖,以自己利益為原則,不理法律不理良心.這豈非反常?我們活的反常,活在嫉妒裡面,活在詭詐裡面,豈非反常?可是被聖靈充滿者,不活在反常也不活在正常；因為若只活在正常,和一般人沒分別.你能的別人也能,你作的別人也作,並不希奇；被聖靈充滿者,乃是活出異常,這樣才能吸引人信主.
\newline
\newline
美國南北戰爭時,有兩個好本事的將軍,李將軍是個好基督徒,白將軍是個馬虎的基督徒；他非常嫉妒李將軍,千方百計予以打擊,萬般留難；李將軍卻始終默默無聲.到了戰爭結束,李將軍辦間大學,白將軍無所事事.當時美國總統認為李將軍是國防部長的合適人選,召他進白宮,對他發表心意；但李將軍說他已從事教育；若接受國防部長之職,教育工作就半途而廢.他要求總統讓他繼續努力教育工作.總統說:「你若能推薦一位即可.」於是李將軍就推薦白將軍；總統說:「豈不是......」李將軍搶著說:「這是他與我個人的事,我推薦他是為了國家的事.」總統說:「既然如此,我願意接受你的推薦.」白將軍聽了這消息,遂即奔向李將軍家裡對他說:「在我心目中,除耶穌之外,你是最偉大的.」李將軍實在活的異常!
\newline
\newline
各位弟兄姊妹!聖經以醉酒預表聖靈充滿,接下去的經文是「當用詩章頌詞,靈歌,彼此對說......」這是論到教會生活.「凡事要奉我們主耶穌基督的名.......」(弗五19-20)這是論到個人生活.到了第六章,開始論到父母和兒女,這是家庭生活,論到人和僕人,這是社會生活.接下又論到須穿戴神所賜的全副軍裝,這是戰爭生活.不錯,聖靈充滿有許多神蹟奇事,有許多異能顯出；聖靈充滿是在日常生活顯出來.所以,以弗所第五章以後,全部講生活的問題；但是今日的基督徒,從神蹟異能去看,很少從日常家庭生活社會生活去看.
\newline
\newline
以利亞膏以利沙,當以利亞找到以利沙,用外衣披在他身上,這件外衣就是代表「膏」.以利沙得著外衣,這代表聖靈的澆灌.聖靈的澆灌容易得的.列王下,也看見一件外衣,以利亞把外衣丟下來,以利沙把外衣拾起來,用外衣打水；以利亞的神同在,水就左右分開；以利沙就過了約但河,先知們看見以利沙,便說:「感動以利亞的靈,已經感動以利沙了.」以利沙的接續,看見以利沙即等於看見以利亞；聖靈充滿,目的是看見你,等於看見基督；這是聖經上一個很基本的目的.所以列王上十九章的外衣代表聖靈澆灌,列王下二章的外衣代表聖靈充滿；王上那件外衣很容易得,王下那件外衣得之不易；因為他從吉甲開始,跟到伯特利,又跟到耶利哥,一直跟到約但；經過一城又一城,一山又一山,被聖靈灌就很著能力,能走追求的路,從吉甲到約但就是追求的路；聖靈澆灌是手續,得著能力才能走追求的路,以達到了目的──聖靈充滿.
\newline
\newline
吉甲是滾去的意思,滾去肉體的敗壞.約翰福音說「真理的聖靈來了,叫人為罪自己責備自己.」聖靈臨到你身上,就給你看見自己的光景,自己的敗壞；你成功時,自高自大；失敗時,自卑自憐；你在光處,在暗處,樣子完全不同；聖靈給你看見,別人假冒為善極其可惡,令人討厭,而自己假冒為善不覺可惡可厭；聖靈給你看見,不應有的思想你不能管制,不應有的舌頭你也管不住,口說不貪愛世界,心中卻暗貪世界；你叫別人聖潔,自己毫無聖潔；聖靈叫你看見這一切,叫你把這些從身上完全滾去.追求的路就是吉甲,滾去肉體的敗壞.
\newline
\newline
伯特利,當雅各逃避他的哥哥,在曠野時,天為他開,伯特利代表天為他開啟.
\newline
\newline
弟兄姊妹!人若經過吉甲,滾去肉體的敗壞,天就為你開啟.今日有許多神的兒女,天對他是關閉的；因他沒有經過吉甲,他對屬靈的事黑暗而不光明；如要屬靈的事光明,必須經過吉甲這一層.
\newline
\newline
耶利哥進入迦南,代表屬靈的甘甜,屬靈的豐富；進入迦南意即進入屬靈的甘甜,和屬靈的豐富；但是耶利哥倒下來,叫你不能進去.要進耶利哥就是要打仗,把耶利哥城攻塌.
\newline
\newline
弟兄姊妹!我們今天是和空中惡魔打仗,以禱告做兵器作戰.
\newline
\newline
有次,我受感動,搭飛機到印度,當時我持台灣護照,到了加爾各答出示護照時；問我是誰簽證?我說,是你們的代辦英國領事館.又問為何給你簽證?我說這應該問英國領使館.結果我護照被扣留待查；從上午等到下午無人理我；天色漸晚,我問他們:「我的問題如何解決?能入境就給我入境,否則派飛機把我送回去.」(明知無此可能)那人員說:「馬上就給你答覆.」經他們磋商,有兩個條件,第一,在加爾各答只可住五天.(除頭尾兩天實際只有三天)第二,明早九時到移民局報到.條件都寫明在護照上,然後放我入境.有位同宗接我到加爾各答浸信會神學院；當時該地旱災,學校停課,全院空無一人.天黑了,他們給我住宿的不是學生宿舍,而是樹林中一個臨時倉庫,放著一張床；傳道人對此並不計較,既來之則安之.那位吳先生大概很累,對我說明早再來就走了.林中寂無人聲,只聽見蟲叫聲；我不想吃飯,關上門禱告說:「我來此不是為觀光,這裡並無光可觀,馬路參差不齊,到處凌亂.是我自己買機票來的,並非被邀請來的,批准了五天,實際上僅有兩天,這五天作甚麼呢?主啊!?要負我的責任；我本意不是來玩的,就算有著名的皇宮,我從來都不想看,我乃是被感動來的.禱告後我睡覺了,也忘記肚子餓,不覺睡到早晨七點；忽有人扣門,起初我有點膽怯,扣門聲變本加厲,繼之大叫,我不得不開門,見兩個漆黑大漢說:「你這傢伙,弄得我們整夜不眠.」我說:「我今天剛從台灣來,和你素不相識,怎說我攪擾你的睡眠呢?」他指著我的鼻子,說:「入關是我為你簽證的.」我以為他來敲竹槓,問他有何貴幹?回答說:「本來請你明天到移民局報到,因為我們有責任,現在加爾各答有八千華人,多你一個也沒人知道,所以半夜來通知你明天不必到移民局.」我十分感謝他們.他把護照上所寫的條件擦掉,還客氣地說:「半夜打擾你,叫你只逗留五天,實在過意不去；如果你想多留下些日子也可以的.」有個蘇格蘭農夫被神感動到印度傳道,不是差會派的,所以他夜宿禮拜堂門口,全身被蚊子咬傷,每天他用一點麵粉烘乾,喝自來水度日；他問我:「你來此人地生疏怎麼傳道?你給我傳單,我替你去請人.」郊外卡巴地方有許多客家人.我到那裡傳福音,有一百五十多人信主,卡巴教會就此建立.
\newline
\newline
各位!我們以禱告與惡魔爭戰,經過耶利哥,然後才到達迦南地,這場戰爭打得過,才能進入屬靈的豐富,才能嘗到屬靈的甘甜.
\newline
\newline
來到約但河──聖靈的充滿.以西結四十七章,有「多」「少」兩個字,水到踝子骨,水到膝,到腰；水多了,人少了,等到水蓋上頭,整個人都沒了,人的「己」少了.加拉太書說情慾所結的果子有廿幾樣.「己」少了就是情慾少了,水多就是靈多,靈只有一位怎會多呢?加拉太五章,聖靈所結的果子有九樣,結出忍耐,又結出喜樂,和平.......靈多就是聖靈的果子越來越多,到約但河就是「己」埋葬然後基督的豐盛才顯明,這就是聖靈充滿.
\newline
\newline
看見以利沙就是看見以利亞,以利沙是以利亞的繼續.看見你就是看見基督,你是基督的繼續.但是以利沙有個強烈的心志.弟兄姊妹!萬勿停留在王上十九章,應進入王下二章；不能停在聖靈澆灌的階段,應進入聖靈充滿的階段,澆灌不過是手續,充滿才是目的.哪有人買了球賽的票子而不進入球場呢?如果你不進入充滿,你追求澆灌作甚麼?澆灌給你有能力,叫你往前走追求的道路,所以聖靈澆灌就是聖靈充滿.以利沙有個強烈的心志,所以他得著了.今日許多人沒有這樣的心志,所以得不著；因為神的恩是賜給「要」的人.
\newline
\newline
你有否「要」的心志呢?
\newline
\newline
耶穌曾說祂不去,聖靈就不來；祂去了聖靈就來.聖靈既來了,我們要追求才能經歷神.我們知道,我們認識耶穌,就追求耶穌；可是今天有個怪現象,很少人認識聖靈卻追求聖靈；頭腦有很多關於聖靈的知識,生活上一點都沒有經歷.
\newline
\newline
十幾年前我被邀請去講道,要我講聖靈的題目；那時我裡面空空,腦子空空,很少關於聖靈的知識；身體也空空,很少有聖靈的經驗.我知道要講這題目,知識和經歷二者缺一不可；我只好博覽群書,逐一收集點滴.很多書寫的都是參考聖靈來的；人家可以講,我也可以講；然而經歷則非如此,人家的經歷,不一定是我的經歷；雖然我也有一樣的經歷.我開始傳道時,從聖經找到根據:「有揀選,一定要有裝備.」
\newline
\newline
出埃及記第三章載摩西被揀選,第四章載他被裝備；且利未記第八章有「承接聖職」四字,承接聖職要舉行禮節,就是摩西要被抹油.在新約也找到憑據,保羅在大馬色蒙召,並非蒙召就可以工作,而是必須裝備方可；神差亞拿尼亞,就先要裝備他.我從聖經找出許多根據,我既蒙召,我要工作,我必須裝備.我剛蒙恩,思想簡單,我上山禱告,有時竄進防空洞禱告,有時竄入草堆禱告；有一天,神給我經歷聖靈,這經歷現暫不提.後來我傳道,發現有件很希奇的事,當我領唱短詩時,有人前來悔改.神給我看見,不是講道給人悔改；乃是神自己叫人悔改.不是你的工作,是神自己工作；所以我說:「主啊!我有過這樣經歷；可是我聽人說,昨日的經歷不能作為今日的資本；因神的恩典每日是新的,不能單靠從前的經歷,應有現在的經歷.」我一直為此事禱告.有一天,我在鄉間講道,因患感冒聲音嘶啞,毫無氣力；和兩位瑞士人一同吃飯,我實在吃不下,就回房間休息；當我經過講台,忽打冷戰,有涕無聲；我想,今晚我怎能爬上台呢?於是入房,自然而然跪下問主:「我怎麼辦?」辦字還未出口,神給我首次說方言,我又驚又喜；那時會場已開始唱短詩,我在方言中一直與主交通.我對主說:「外面等我工作,我必須停在這裡.」當我進入會場,三百多人齊站立；我受寵若驚,我算得甚麼,他們居然對我必恭必敬.我上講台,聲音恢復了,力量也來了；會畢病人要求禱告,手踫病者,病得醫治.我滿心感謝主!我繼續帶領聚會,勢如燃燒,如水沸騰；散會之前有見證會,有個知識份子,他看見數百人中這個哭那個哭；不明到底是何玩意.他想出來講攻擊的話,他的臉色難看,用鄙視的眼光看會場；突然間,神的靈在他身上,他跳躍,且跳的很高,方言從口而出,人完全改變.這是那次會中所見.於是我想聖經的事情,神要以利亞找以利沙,把衣服披在以利沙身上；不是以利沙要的,是以利亞作的.我想,那個要講鄙視言語的人；神的靈降在他身上,不是他求的,乃是神自己給的.「......何況天父,豈不更將聖靈給求祂的人麼?」(路十一13)我看見「更」字,我很高興；因為不求祂都給,求祂豈不更給嗎?
\newline
\newline
今晚,我只講兩點,第一,我要答覆很多的問題.
\newline
\newline
在行傳中,耶穌叫門徒要等候.若耶穌還沒有去,聖靈就不會來.叫他們等候是有道理的；但耶穌已去,聖靈已來；還叫他們等候甚麼呢?仔細來看,不論自然界或靈界,每事有時間有次序.在自然界,一粒種子在地裡,必須等候；發芽長枝葉,開花結果,都有次序有時間.生物界也如此,母體懷胎,嬰兒的四肢五官都看不見；過此時候,四肢有了,五官也明了；到了二百七十天就生下了.有時間也有次序,靈界也是一樣；神造天地,並非一口氣呼成,乃是說有光就有光,這是一天.然後把水上下分開,這是第二天.都是有次序有時間的.所以神要我們等候.
\newline
\newline
我想,有兩個問題要等候；等候真理的問題解決,等候攔阻我問題除去.聖靈的問題,有很多真理的問題.「......不多幾日,你們要受聖靈的洗.」(徒一 5)許多人對聖靈的洗有個講法:「重生就是聖靈的洗,聖靈的洗就是重生.」大概現在很多人這樣主張；因為當時耶穌說:「過不多日你們要受聖靈的洗.」「你們」乃指使徒；難道使徒沒重生嗎?聖靈的洗和重生,不能合為一談.
\newline
\newline
有聖靈為甚麼還要追求聖靈呢?這樣的說話很多.以弗所書一章說,信就可以得著所應許的聖靈.如果沒有聖靈,你就不會認罪,不會承認耶穌是主,不會認神是父.你認罪,信耶穌是主,神是父；這就是你有聖靈的憑據.有聖靈還要追求聖靈纔合邏輯呢?神並不在邏輯裡面；如果神限於邏輯之內,就不謂之神了.我們姑且不管這些理論,從耶穌基督身上來看,祂有聖靈還要追求聖靈.當天使找到耶穌的母親馬利亞,說:「你要懷孕生子.」她立刻說:「我還沒有出嫁.」天使說:「因為你懷孕是從聖靈而來的」所以祂有聖靈,祂要開始天國的事業；祂在約但河受洗聖靈降在祂身上,好像鴿子一樣.路加四章說祂被聖靈充滿,滿有聖靈的能力.祂有聖靈又充滿聖靈.使徒也是如此,當耶穌釘在十字架以後,他們非常害怕,躲在房子裡；耶穌進去,向他們吹口氣,說:「你們受聖靈.」到了行傳一章,耶穌說:「過不多幾日,你們受聖靈的洗.」所以,耶穌有聖靈,使徒有聖靈,還要追求聖靈,這是有根據的.
\newline
\newline
有人曾說,五旬節以前,聖靈在耶穌裡面,五旬節以後,聖靈是在身體裡面；身體就是教會.所以行傳二章,使徒被聖靈充滿,使徒是猶太人.行傳十章使徒在哥尼流家裡；他們是羅馬人,是外邦人；聖靈在他們身上.教會包括猶太人也包括外邦人；既然聖靈已經把猶太人進在身體裡面,也把外邦人進到身體裡面；聖靈的工就成了,聖靈的事就完成了歷史.因此,今日不必再談聖靈的洗.但是「因為這應許是給你們,和你們的兒女,並一切在遠方的人.」(徒二29)「你們」是一代,「你們的兒女」又是一代.聖靈的洗,沒時間限制；既有遠方的人,當然也有近處的人；聖靈的應許,沒有空間的限制.
\newline
\newline
希伯來書說:「各樣的洗禮.」有很多人說是指滴禮和浸禮.因為以弗所書說:「一信一洗.」所以有人只能指一樣,不能指兩樣而言.滴禮是洗禮,浸禮也是洗禮；各樣洗禮乃指水洗和靈洗.我們從聖經看,基督徒不僅是器皿,而且是器具.杯子是器皿,用來盛茶；碗是器皿,可以盛飯.基督徒是器皿,神要把祂的靈盛在裡面.我們也是器具；筆是器具,腳是器具,可以使用；所以基督徒是器皿也是器具,是信徒也是門徒；基督徒聽道還要學道,有了水洗,還要靈洗,這是有聖靈根據的.所以等候真理的問題要解決,否則你的攔阻就一直存在.等候我們攔阻的問題要解決.我們有許多攔阻的問；害怕攔阻問題,欲追求聖靈,卻追求到邪靈；我看見了多次,有人追求到邪靈的事.遇到這樣的問題該怎麼辦呢?路加十一章說,兒子求餅,父親不會給他石頭；求魚,不會給他蛇.人向神求聖靈,難道神給邪靈嗎?父親都不至於如此,何況我們的天父呢!既然這樣,為甚麼有人求聖靈反而得著邪靈呢?當西門看見使徒們按手,人就被聖靈充滿；他們吩咐鬼,鬼就被趕出去.西門看了覺很希奇,認為這是得利的門徑,就給錢使徒,請使徒把這恩賜給他.使徒對他說「你的銀錢和你一同減亡吧!」由此可見,西門居心不正,他所求的不是聖靈,而是自己的利益.求聖靈反得邪靈,問題在於居心不正；他所求的是聖靈以外的東西,所以害怕是多餘的.
\newline
\newline
理性的攔阻問題.這問題很麻煩,我們所信是頭腦想得通的,不能相信頭腦想不通的.我的毒瘤沒有取出,至今已卅六年了.有位高官為人很孝順,他父親患的和我一樣的結腸癌；他聽說我的病經已痊愈,就和我連絡來訪問我,問我治病秘方；我說我的秘方就是神蹟,他始終不相信；因他的頭腦想不通.知識份子只能接受頭腦解得通的事,這時對於理性是個很大的攔阻.基督教已漸變成早期的猶太教；猶太教有祭司有文士,祭司負責聖殿禮節的工作.
\newline
\newline
現在基督教也漸落進禮節儀文裡面.靈裡的需要,儀文不能解決,於是就落到文士裡面；文士本是研究聖經的.今日基督教研究高深學問；我並非說這些不好,不過這些未能滿足靈的需要.今日基督教走情感路線,大喊大叫大奔大跳,發洩內裡的情感；我從前很不以為然,覺得這樣很不斯文,不成體統；但後來我想通了,我並非贊成,而是我發現神乃重感情的.何以見得?當耶穌到拉撒路墓前,祂哭了；還有,耶穌在拿因地方看見寡婦哭送已死的兒子,就動了慈心.感情是心裡面的東西,如你不敢用感情,就等於不敢用心；人若不敢用心,怎能接觸神呢?昔日大衛迎接約櫃時,在約櫃前唱歌跳舞；他的妻子米甲因此輕視他,說以色列的君王在百姓面前露體.後來米甲被詛咒,她終身不能生育.這是理性的攔阻.我並非說理性不好；但別以理性攔阻聖靈的事.
\newline
\newline
「人怕人」的攔阻問題.彼得和外邦人同吃飯時,看見猶太人來了立即走開；後來保羅看見他在此事上不對；因為福音是為猶太人也是為外邦人的；彼得看見猶太人來就避開,這和真理有所抵觸；所以保羅當眾指責他.
\newline
\newline
各位!關於真理的事,不必怕人；接受耶穌是一部分,接受聖靈也是一部分；接受耶穌又接受聖靈,才是全備的救恩.現在的人都喜歡一半一半的,從前的人也是喜歡一半一半；神叫亞伯拉罕從米所波大米出來；他走一半到吾珥就停下來,然後才繼續走下去.雅各從巴旦亞蘭出來,走到示劍就停止；他也是走了一半.還有老底嘉教會好像溫水半冷半熱,也是一半一半.各位弟兄姊妹!屬靈的事,一不作,二不休；要走就當走到底,不能只走一半.所以神等候你的問題解決,也等候你攔阻的問題解決；這是在追求之前.你在追求之時,神也在等候,等候你發覺自己軟弱無能；我們無論對自己對整體,都何等沒有能力,在得勝上何等沒有能力!當雅各從巴旦亞蘭回迦南地,代表他走上追求的路；他需要先解決的,是他和他哥哥的問題；因他們之間有仇恨.「弟兄和睦同居......好比那貴重的油澆在亞倫的頭上......」(詩一三三1-2)問題解決了,膏油才能澆灌下來.當雅各經過以東地帶,靠近他哥哥住處時,他邊走邊跪七次,七代表完全；向他哥哥仆伏下來,直到饒恕,仇恨消失為止.但是我們常發現,一而再,總不能成功就不作了.在得勝上何等沒有能力!
\newline
\newline
人生的攔阻問題.人的一生有許多遭遇,我年青時曾經被日本人監禁,曾經患癌症；如同經過死亡的邊緣谷.人生有很多挫折很多打擊,須有能力才能得勝.
\newline
\newline
我曾經探訪一位老將軍,他要我為他禱告,他對我說:「你禱告,神必垂聽；請你為我禱告神,叫我快快死.」我問他為何要這樣禱告?他說:「如果你這樣為我禱告,我今晚睡覺,明早就不起床了；也許我沒這樣福氣,明早起床後我就腦溢血離世；若再沒這樣福氣,我明天上街必被車撞死.」我問他為何要這樣禱告?他說:「弟兄,我問你,你豈忍心叫我過這樣的日子嗎?」他一家三口,女兒因婚姻挫折,神經分裂,精神崩潰；他的妻子原是小兒科醫生,因被電單車撞傷,腦子不清楚.他現年76歲了,要照料神經病的女兒和神志不清的妻子.他在得勝苦難的事上,是多麼沒有能力.我們的事奉多麼沒有能力；和他人配搭事奉,發現自己個性特別,脾氣暴躁,心胸狹窄；遇事一觸,臉紅聲厲；見人態度不好,立即吃不消；聽人說句逆耳的話,立刻受不了；雖曾受過教導,理應溫柔,應忍氣吞聲逆來順受,道理雖懂；實際卻做不到.我們生命上沒有能力.
\newline
\newline
「我們堅固的人,應該擔代不堅固人的軟弱；」(羅十五1)我在福建時,有次出門,到沒有車的地方,我就步行；但我內人走不動,只好坐轎子.中午,轎夫停下吃飯,他們和賣飯的人談話:「所托的事辦了沒有?」「辦好了,請放心!」「堅固嗎?」我很好奇,到底所托何事邊問堅固嗎?因此,我問賣飯的,他說:「托我找配偶.」我說:「為何不問漂亮賢淑卻問堅固與否.」回答說:「他要的是個挑水,劈柴,養豬,耕田甚麼都能幹的堅固人.」「堅固」原文甚麼都能之意.但是我們發現甚麼都不能,攻克己身不能.所謂攻克己身,就是被毀謗,面不改容,心不跳,手不冰冷；被人稱讚,不洋洋自得,不得意忘形；我們謙卑不能,溫柔也不能樣樣不能.我曾說過,能不在你,神要「從嬰孩和喫奶的口中,建立了能力......」(詩八2)弟兄姊妹!嬰孩是軟弱的意思,耶穌因軟弱被釘十字架,你應該與祂一同軟弱.聖經中有兩種軟弱,一是不好的,好像「心靈願意,肉體軟弱.」還有一種軟弱的,耶穌因軟弱被釘十字架,這種軟弱的意思,是祂對神不頂撞,不抵抗；凡出於神命令,祂完全順服.神的能就在這樣的人身上.可是我們發現了件可怕的事,我們甚至連軟弱都不能.固然不錯,神的能要能在嬰孩身上,嬰孩非常軟弱,任你抱,任你摔,毫無抵抗；但是連這一點我們都無能,這是我們的光景!我們需要聖靈的能力,教我們對神軟弱.若非聖靈工作,我們連這都做不到.
\newline
\newline
我們還有整體的需要.今日教會,雖喊叫「增長」美名,實際卻是退後；去年一百人,今年僅九十人,可能明年只有八十人.還有種是停住的,十年前後,同樣一個禮拜堂,同樣這些人.真是越看越膽寒；因這班人越來越老,教會老氣橫秋,我們對於整體也無能.
\newline
\newline
個整體各有需要.個人的需要對神就是對人；對人就是對神.耶穌曾說:「我餓了,你們沒有給我吃,我在監牢,你們沒有來看我.」(太廿五42-43)他們說:「我們沒有看見你餓了,也沒看見你在監牢裡.」耶穌說:「你們既不作在這弟兄中一個最小的身上,就是不作在我身上了.」所以,你對神如何,乃是決定你對人如何.你應按聖經的教導對待人,但約翰十六章說:「只等真理的聖靈來了,他要引導你們進入一切的真理.」若非聖靈的工作,雖然腦袋知識一大堆；但生活的實際毫不見得.
\newline
\newline
今天我從《華僑日報》看見一篇攻擊基督徒的文章:「早知今日,何必當初.」論到人來信耶穌之前自由自在,信了耶穌後很不自由；一次禮拜天不到教會,牧師便追問,實在令人厭煩.
\newline
\newline
弟兄姊妹!基督徒應該進入實際,真理不是一些知識,乃是從實際活出來的；要等真理的聖靈來了,才能幫助我們進入真理.現在有許多人很苦惱,知道饒恕而不饒恕人；知道謙卑但不能謙卑,知道當順服神但代價太大,無力支付.個人問題需要能力,整體問題也需要能力,處處需要能力.
\newline
\newline
我覺得基督徒的苦惱,不是知和不知的問題,乃是能和不能的問題；為何不能?因沒有能力,為何無能力?因對聖靈的問題,沒有清楚認識.「我來要把火把丟在地上,」(路十二)呂掁中牧師的聖經新譯:「他盼望能夠燒起來.」甚麼都需要火,昔日在聖殿,有火才能燒香,點火,獻祭.今日在家庭也需要火,才能燒飯,燒水.工廠也需要火,才能製造產品.我們為主作出口的也需要火,沒有火,我們只有些字句；有火才有經歷.唱詩也需要火,沒有火就沒有聲音,不能感動人,所以在在需要火.
\newline
\newline
火就是聖靈「......何況天父,豈不更將聖靈給求祂的人麼?」(路十一13)「更」意思是你沒求,祂都賜給你；你求,祂更賜給你.
\newline
\newline
求神給我們看見自己個人的需要,教會整體的需要；追求聖靈得著能力.但神在等候你真理的問題弄清楚,等候你攔阻的問題挪開；因為很多人對真理尚未清楚,很多人仍是攔阻重重；特別是理性的麻煩.
\newline
\newline
求主讓我們渴慕聖靈,其非人的才能口才可以突破這個世代,唯聖靈可以突破這個世代.我們應當起來渴慕聖靈的充滿!各位!是嗎?若是,請說阿們!
\newline
\newline


\newpage

\section{第58屆~港九培靈會~吳勇~第九講}
\label{sec:471}
\textbf{聖靈的工作}
\newline
\newline
連結: \href{https://www.hkbibleconference.org/session-message/view/471}{\texttt{https://www.hkbibleconference.org/session-message/view/471}}
\newline
\newline
\hyperref[sec:470]{< < < PREV SERMON < < <}
~
\hyperref[sec:toc]{[返目錄]}
~
\hyperref[sec:455]{> > > NEXT SERMON > > >}
\newline
\newline
今天晚上已經到了大會結束的前夕.
\newline
\newline
這次非常感謝主,因為我是從美國來的,我恐怕入境的事辦的不清楚,所以戰戰兢兢.參加大會感恩會次日,我便回美國,領十天學生聚會；雖然僅有幾十人,但我仍要如期舉行,這裡的聚會有數千聽眾,但數十聽眾我仍要講道.請弟兄姊妹們繼續代禱,求主托住我.
\newline
\newline
這次講道,關於教會,我講了五個題目,關於聖靈,已經講了三個題目；深信很多人心裡有問題,我心中有了準備；如果大家要作文章,神給我的話,也一定要大膽地傳.請大家為此聚會禱告.我講到聖靈的果子,我很著重以生活為基礎.我又講到對聖靈的追求.昨晚講到聖靈的充滿.今晚我們要看個異象.
\newline
\newline
摩西出來為神使用,神給他看見異象,他看見荊棘在火焰中沒有被燒毀.荊棘不是耐火的材料,著火就燒毀了.荊棘代表肉體的東西,火是屬靈的試驗；神給摩西看見肉體的東西燃燒,就算有多大的本事都不是材料.這是神叫摩西看見此異象的目的.雅各看見異象,就把那地叫伯特利,後來把石頭澆油；由此可見石頭澆油才是伯特利的材料.單單石頭不能作伯特利的材料,單單油也不能作伯特利的材料.石頭代表人,伯特利代表教會；單單人不能作教會的材料；人的裡面有神,才能作教會的材料.
\newline
\newline
聖靈的工作可分為三點來思想:
\newline
\newline
第一點:骸骨,就是死.耶穌第一次講教會,講到陰間的權勢(太十六18)死就是陰間的權勢,所以今日的教會,要遇到陰間權勢的攻擊.昔日的猶太人曾經遭遇三種仇敵,這三種仇敵都是陰間的權勢,就是埃及,巴比倫,亞述,許多猶太人都被這三個仇敵殺死.
\newline
\newline
埃及代表屬地的豐富；因有尼羅河灌溉兩岸,出產豐富.猶太人在曠野時,想到埃及的豐富,有菜,蒜.屬地的豐富要扼殺屬靈的生命；屬地的豐富盡是虛空.在猶太人的飯廳常掛兩幅圖畫,一是初生嬰孩握緊兩手掌,一是人離世時手掌鬆開,猶太人以此警戒自己.人來到世界,要賺得世界,爭奪世界；但無論賺的,爭的多少,最後皆空.雖然屬地的豐富是虛空,但世人費盡精神,時間,為屬地的豐富投資；卻沒多餘精神和時間追求屬靈的事；因此,屬地的豐富就把我們屬靈的生命扼殺了.
\newline
\newline
第二種仇敵代表巴比倫,當時他們用磚造塔,不是用石頭.(創十一)石頭代表神工,磚代表人工.今之基督教會,用人工代替神工,就把屬靈生命扼殺了.
\newline
\newline
何謂人工代替神工?教會是基督的身體,身體是肢體,應互相配搭,我們從神工講究,應該從生命講究,就是要走十字架的道路.耶穌被審判時,彼拉多查不出耶穌的罪,要把他釋放；但猶太人極力要除掉他.所以十字架是除掉的意思.但是今之基督教會,所講究不是神工乃是人工,講究公共關係,交際的藝術,用人工代替神工；現在教會做事,每個國家辦事,必先計劃用多少錢；若以神工講究,我們應向神求；若以人工講究,我們用捐冊向人要.今日教會的治理,若講究神工,我們應用膝蓋來禱告神；但講究人工及注重方法和組織.巴比倫以人工代替神,所以也把屬靈生命扼殺了.
\newline
\newline
亞述和猶太人是同宗,可說他們很接近.今天和我們最接近的,並非我們的朋友,甚至不是我們的配偶；和我們最接近的是人的肉體,無論起居飲食,都和我們在一起.亞述代表肉體的邪情私慾,這些扼殺了屬靈的生命.邪情有七情六慾,有喜怒哀樂；並非我們不可發怒,耶穌在聖殿也曾發怒.摩西在看見猶太人拜金牛時也曾發怒,把石版摔碎；但他們的怒非為私慾乃為神的榮耀發怒.我們常為面子被人損傷而怒,為利益被人侵犯而怒；這是私慾,是為肉體需要而怒.我們每人都有所需要,需要明天比今天更好,更舒服,更富裕；這是神所許可,也是良心許可的；可是人常用不正當手段,甚至偷,搶以達到滿足肉體的需要.
\newline
\newline
從前猶太人變成骸骨,一片死亡,就是被埃及,巴比倫,亞述殺了.
\newline
\newline
今天從教會光景看,看見一片死亡.有次,我請教一位醫生怎樣才宣佈死亡?他說:
\newline
\newline
(1)心臟停止跳動；
\newline
\newline
(2)呼吸停止；
\newline
\newline
(3)瞳孔放大；
\newline
\newline
(4)血壓零度.
\newline
\newline
如果從這些衡量死亡,我們從屬靈的光景來看,心臟是全身樞紐；人的生命,我們的樞紐,我想是感覺.瑪拉基書載,猶太人一直問神:「你說你愛我們,你何時愛我們；你說我們污穢你的壇,我們在哪兒污穢你的壇呢?」神愛他們,他們沒有感覺；他們污穢了神的壇也沒有感覺.這是猶太人屬靈的光景,沒有感覺,麻木不仁；從瑪拉基書到馬太福音,間隔四百年都在黑暗中.各位!如果你對屬靈事情沒有感覺,沒有反應,這是死亡現象.有天我說,呼吸代表禱告.嬰孩出世,自然曉得呼吸,不必教導.基督徒重生得救,很自然地,用膳要禱告,睡覺要禱告,有的基督徒除了以主禱文禱告,就不會禱告；這樣的人不會呼吸,等於死亡.瞳孔放大就看不見了.在以賽亞書有句話:「不是日光,不是月光,乃是耶和華的光.」我們看聖經,不是靠張三,李四的光,乃是靠聖靈的光；如果只見白紙黑字,這就是瞳孔放大,屬靈一片黑暗,近乎死亡.
\newline
\newline
血壓零度.血壓是全身的動力,動力沒有了,與肉身爭戰,與世界爭戰,與空中惡魔爭戰都沒有動力,這也近乎死亡.
\newline
\newline
當時以西結看見遍地骸骨,就是死亡光景太普遍.
\newline
\newline
今日教會,照樣有安排,計劃,聚會,活動；可是教會中卻沒有生命,沒有活潑,沒有新鮮；雖有聚會,會中聽不見主的聲音.如果有主的聲音,神的話是精義,沒有主的聲音,神的話是字句；精義叫人活,字句叫人死.我們常發現,基督徒裡面有感覺,來參加聚會,覺得多點知識,多點熱鬧,除此一無所獲；不聚會也死,越聚會越死；參加活動,只是肉體的動,沒有聖靈的動.體貼肉體者死,體貼靈者才生.所以以西結看見遍地骸骨；神也給我們看見到處都是骸骨,到處死亡.神要叫骸骨活起來；但神作每事,必要藉著人.耶穌把水變酒,要使人挑水,因為神的事,一定要透過人,這人就是以西結.弟兄姊妹!今天既有諸多死亡,神需要像以西結這樣的人.
\newline
\newline
以西結是個有靈的人,他並非有學問,才幹,口才的人；雖有這些是好的,但最要緊是有靈.從聖經查考,無論新約,舊約,凡被神使用的都是有靈的人.我們看士師記那些士師,每興起一位士師,耶和華的靈降在他身上.新約也是如此,保羅在大馬色路上,他曾問主:「我應當為你作甚麼?」主沒答覆他作這個,作那個,主說:「你作的時候,人就告訴你.」後來亞拿尼亞先把手放在他身上,叫他得著裝備,叫他有靈.
\newline
\newline
各位!我們不是作學會的工作,也不是作社會的工作,我們是作為屬靈的工作；有靈才能作靈的工作.怎樣才有靈呢?我這裡也講揀選,人若被揀選必須被裝備.出埃及記三章,摩西被揀選,第四章摩西被裝備；舊約如此,新約也如是,耶穌揀選使徒,祂升天,天國的事交給使徒；使徒必須被裝備,才能挑擔這重任.揀選是非常重要的,出來作聖工的人,揀選有兩方面,一方面是神揀選你.神揀選摩西,當摩西生下來,他的母親沒法藏起他,就放在箱裡隨河水流下；神安排法老的女兒救起摩西.神對他必有安排,安排摩西在法老皇宮中,得了埃及的學問.
\newline
\newline
兄弟也被神揀選,神對我有安排；然後給我操練,當我信主的第二個禮拜,講台上宣佈下禮拜開辦主日學.年青人哪個參加請舉手,全場舉手唯我一人,所以目標特別顯著.他們把我也看得特別寶貴,主席下來和我握手擁抱,對我說:「下主日十點禮拜,我們九時主日學,請你準時到達.」我們共有七人,都是廿多歲的青年,另有個卅幾歲的主日學教員,他講聖經給我們聽.接下去的主日,他被工作機構調派別處,我們沒有老師,大家面面相覷；後來推舉一位代替老師.那人說:「不成,我信主僅三年,不但舊約沒有讀過；新約也沒讀畢,怎能當主日學教員.」又推舉一位,他說:「你信主已經三年,我只有兩年,更加不行.」後來推舉到我來,我說:「我信主只有三個禮拜.」我想,如果大家都推辭,主日學不是要解散嗎?我是新加坡人,大概臉皮較厚,所以我就接受了.於是我站起來,把那天我所讀的馬太福音,給他們講了一遍；講完之後,竟然有人說:「講的不錯啊!」我實在不敢相信；但是個個都說:「是的」,還請我下禮拜繼續擔任.回家我告訴內人,她非常驚奇說:「你舊約沒讀過,新約連馬太福音都未讀完.」我說:「因沒有老師,只得充充場面；他們還請我下禮拜再講,不知怎麼辦呢?」她說:「從明天開始,你就和太陽競賽.」我不明白,她說:「舊約出埃及記,猶太人拾瑪哪是在太陽未出來之前,太陽一出瑪哪就化了.太陽未出來你就開始讀聖經.」自從那天起,除了生病,我必不讓太陽比我早起,至今已41年了,有揀選,有安排,神對摩西有揀選,神對他就有安排；對我有揀選,就為我安排園地,叫我從中學習；但是單神揀選不夠,還要你揀選神.
\newline
\newline
各位!希伯來書十一章載,摩西放棄皇室的快樂,寧可和神的百姓同受苦害,這就是揀選神.有許多人被神揀選；但不能出來,因被很多東西包圍,結果胎死腹中.被前途,工作,艱難,害怕所包圍,雖有被神揀選的心志；但沒有揀選神,到最後,揀選不成功.
\newline
\newline
以西結是有靈的人,是被神揀選的人,他也是揀選神的人.
\newline
\newline
如果神揀選你,你也揀選神,這樣的揀選才算成功；如果揀選成功,你就被裝備,裝備就是靈.
\newline
\newline
以西結是個有異象的人,異象是很要緊的,從新約來看,彼得有異象,他才知道福音是為猶太人,也是為外邦人.保羅看異象,他才知道往馬其頓去,才知道該走的路線,異象對傳道人很要緊,傳道人站在講台上講話,有裡面的話,有外面的話；先有裡面的話,再用外面的話來解釋裡面的話.傳道人若沒有裡面的話,要用甚麼外面的話呢?裡面的話是神用異象的啟示.我常見為主作工者,沒有裡面的話,所以也就沒外面的話.有啟示才有負擔,負擔是從啟示來的,從啟示而來的負擔不勉強,是甘心的.
\newline
\newline
聖經上的人物怎樣得異象呢?以西結在迦巴魯河邊禱告時看見異象.但以理是預言的鎖匙；如果沒有但以理書,這預言就解不開；要把預言解開,必須讀但以理書；但以理所以有這些預言,乃是從異象來的；他得異象,乃是從禱告來的.新約人物也一樣,彼得看見異象,是在硝皮匠家禱告時看見的.約翰在拔摩海島,禱告時看見異象.以西結是個有異象的人,是個會禱告的人.人若不會禱告,就沒有異象.
\newline
\newline
以西結是個有信心的人.神問以西結說:「人子啊!這些骸骨能復活麼?」他回答說:「神阿,你是最知道的.」(英譯)意即這些骸骨能否復活在乎神,一切在乎神的能力,神的作為；神的能力和作為,要在信的人身上顯出來.以西結因著信,神的能力在他身上.
\newline
\newline
一九五一年我患了結腸癌,已達蔓延階段,癌細胞散開；開了刀,醫生發現病勢已向上發展,胃部整個被毒瘤包圍；因此,醫生不敢動,只好立即縫合,希望這樣可以多活幾個月.到了第三天,教會開會決定應將此事告訴我內人,好讓她作心理準備；當醫生告訴她時,她暈過去,那時她才29歲,最大的孩子四歲多,最小的不到一歲.有位作長老的姊妹,上樓抱起我內人,對她說:「可憐呀!這麼年輕!」意即這麼年輕將成寡婦；內人聽了就哭起來.這時,神的話安慰她:「神能照著運行在我們心裡的大力,充充足足的成就一切,超過我們所求的所想的.」(弗三20)於是她哭聲停止,擦乾眼淚,說:「神啊!我感謝你,你能就夠了.」她知道神能,憂愁消失了,回復從前正常的光景.到了第十一天,當我睡醒時,正要讀聖經時,對內人說:「有人翻過我的聖經」她說:「你剛睡醒,視力不清,有風吹過,你說有人翻你的聖經.」我說:「是的,現已停了,你拿聖經過來給我看.」平時讀經有重要的句子,用了紅筆劃,其中我看見三句話,
\newline
\newline
(1)耶穌在前頭走.
\newline
\newline
(2)主要用他.
\newline
\newline
(3) 解開.
\newline
\newline
我下意識自言自語「耶穌在我前頭走,行過死蔭幽谷也不怕遭害.如果主要用我；我還有年月存活.既然解開,我不必綁在床上.」我對內人說:「回家去吧!」我抱著個漲滿水的大腹出院.結果,水消了,毒瘤化了；我出來傳道直到如今.所以,神的能,人的信；人的信,神的能.可是信的另方面就是順服,我們常說:「信而順服.」
\newline
\newline
以利亞是個順服的人,神叫他在基立溪旁,他遵命；神叫他見亞哈,亞哈要對付他；性命堪虞,但雖至於死,他也順服.
\newline
\newline
為主所用的人肯順服主,是很要緊的.我內人從前在監牢任職十八年,後來監牢要從城市搬遷到鄉村；本來上班僅需十分鐘路程.搬後,來回最少需三小時；家中小孩子多,教會會友也很多,作為傳道人的內助,忙碌情形可以想像到.曾有一年,她因忙碌過度,神經衰弱；有次出門探訪,忽然忘記家在那裡；中午沒回家吃飯,太陽將下山也沒回來；後來一名大學生,發現她臉色蒼白,在大街小巷穿竄.問她作甚麼,她說找不到自己家在何方,幸虧那學生領她回家,她神經衰弱竟至此地步!因監牢要搬遷,我就和她磋商,工作十八年了,讓別人去擔當吧!我這樣說是為了減少她的負擔,免得身體惡化.她問我理由何在?我不答覆,後來一直繼續工作到最後一週；她講完道之後,許多犯人都圍繞她,有的給她倒茶,有的給她扇子；她工作了十八年,犯人中有無期徒刑的,有販毒的,有判刑幾十年的；大家和她感情很好,有個犯人問:「吳太太,下週搬遷,你還來嗎?」她想:「如果我說來,丈夫叫我不要來,如我說不來；神感動我來,到底聽從丈夫或聽從神呢?神大,但丈夫也不小；於是她回答「再說吧!」這時,犯人都哭了.她的心被感動地說:「我來,無論如何辛苦我都來.」結果,神居首位.
\newline
\newline
以西結是個有靈,有異象,有信心,而且肯順服神的人；這樣的人才合神所用.他使死的骸骨變成活的軍隊,首先他的工作,是使骸骨連絡起來.骨頭是身體裡面的東西；教會就是身體,有身體才能彰顯基督榮耀,才能執行基督的計劃.所以撒旦從古至今,極力攻擊這個身體,要叫這個身體──教會與社會同化；使教會失去功用,自形腐化,身體變成屍體.骨頭是身體內的東西,常因裡面有問題而致腐化；有時因裡面意見不同,彼此不滿,互相誤會,批評,攻擊,種種原因致身體分散.但是我們感謝神的恩典,一個被神所用,有靈,有異象,有信心,能順服的工人,當他作工時,能把內裡的問題挖出來,好叫這個世上的身體能聚集起來.
\newline
\newline
第二件工作,骨頭合起之後,加筋長肉,「加」「長」,就是培養,筋是連絡的；「愛」連絡全德,把愛培養起來.教會把愛培養是很重要的.「......愛是不計算人的惡」(林前十三5)(人的惡這幾字是譯者加上的)不計算我待人的恩,也不計算人對我的惡,對我的批評,攻擊,一概不計算；這樣,身體才不會受傷.被神用的僕人,要把愛培養起來,教會才能建立.長肉,肉是健康的現象,也是生命的現象,要把生命培養起來.生命就是領受了神的道,道在腦子裡,在日常生活中,這樣,道才叫做生命.
\newline
\newline
今天的基督教會,要注重的不是建築物,也不在乎音樂,其實音樂很要緊,但教會並不以此吸引人.紐約有個特別的建築物,如果講台上有根針掉下,連坐在最後邊的人都可清楚聽見；那個特別的建築物不需要擴音器,但是教會吸引人也非用這個.教會若滿了恩典,自然會吸引人.怎樣才有恩典呢?「弟兄和睦同居...... 好比亞倫頭上的膏油,又好比黑門的甘露.......」(詩一三三)我們知道,膏油就是聖靈,箴言十九章說甘露是神的恩典.甘露從黑門山來,黑門是毀滅的意思；如果你有毀滅就有甘露；毀滅是人與人之間的問題.我曾提及約瑟在埃及作宰相,後來他的兄弟發現,原來宰相是他們的兄弟約瑟；當年他們把他賣了,這仇他怎不報?所以大家一齊向他求饒恕,但是約瑟說:「從前你們意欲害我；但神的旨意是好的,為要成就今日的光景.」他兄弟們的無情,兇惡,將約瑟全都毀滅了,所以甘露降下,他作了埃及宰相.還有大衛,他是當時民族的救星,國家的英雄,對猶大國有很大的貢獻；不料掃羅嫉妒,不容他,也不讓他；他千方百計要害他的命,這時大衛可說是與掃羅有極大的冤仇.雖有數次機會報復,大衛卻不報復；將軍妻為他報復,甚至大衛阻止；因他和掃羅之間的仇恨已經毀滅,所以甘露就下來,他作了以色列的君王.各位!如果教會中彼此相安,心靈相通,教會就滿了甘露,滿了神的恩典；教會自然就吸引人.被神所用的人,不但培養愛,也培養生命；所謂生命的解釋,神的道裝在腦袋,神的道就變作生活；這樣就滿了恩典.
\newline
\newline
還有一道,「皮」就蓋起來,皮是表面的一層,可以保護裡面；皮是聖潔的,帶領人聖潔,帶領人與世俗有分別.所以以西結被神所用,挖出裡面的問題,培養裡面的愛和生命,帶領他們過聖潔生活.
\newline
\newline
還有一點,「氣息」,有氣息才會活起來.教會裡面,不在於有無聚會,有無活動,乃在於有氣息,才是活的.氣息就是靈,有靈的教會才是活的教會.使徒的教會是有靈的教會；因為有靈就有引導,所以把腓利引導到曠野,福音傳到外邦.有靈就有啟示,才知道福音為猶大人,也為外邦人.有靈就有大能的明證；行傳這卷書,不是單用道講給人聽；也用大能的明證顯給人看.我們教會有個小孩子,一天發燒看醫生,據說是感冒,服藥後燒仍不退；在家裡邊走邊跌,他母親說:「為何不好好走?」孩子說:「我眼睛看不清楚.」母親急忙帶他看小兒科醫生,發現有隻眼睛瞎了；就提議送台大醫院診治.一經檢查,腦子有瘤,壓住視線神經；醫生說再觀察兩天；如瘤繼續發展,就當儘速開刀.繼之又發現孩子不能小便；因瘤壓住神經又起了變化,醫生決定為他開刀；問醫生開刀的成功率多少,答一半一半.母親急不及待,打電話給我內人,問她應否開刀?我內人很難答覆,對她說:「你求問主吧,因我們的主是活的；祂會給你答覆.」她禁食禱告,主給她清楚的啟示:「耶和華是醫治人的.」她立即到醫院辦理出院手續；醫生對她大罵一頓說:「開刀還有五十分的希望,不開刀一分希望都沒有；我未曾見過這樣的母親,連五十分的希望也不爭取.」她知道和醫生無法講清楚,終於帶孩子出院.當晚,孩子的眼睛看得見了,次日完全恢復正常,到醫院ｘ光檢查,瘤子不見了.這是大能的明證.
\newline
\newline
如果你熟讀行傳,有靈就有對付罪惡的能力；這是今日基督教會最難的.有次,教會見證會,有的人出口成章,有的人兩句話都說不清；有個司機太太起身作見證,講了廿分鐘,大家都不知道她講甚麼；其中有位弟兄,忍無可忍說:「張太太,你在講甚麼?浪費大家的時間,坐下吧!」張太不好意思,臉紅了,繼之臉又變青,突然嘩啦大哭,躺在地上滾.三百多人的聚會,她在地上邊滾邊哭邊喊叫「主啊!我被人侮辱.」當時我是主席,我說:「弟兄姊妹!大家都跪下禱告.」我擬以會眾禱告聲壓過張太太哭鬧的聲音；不料哭聲變本加厲,但是卻無法超越會眾的禱告聲；後來我聽見霹霹啪啪之聲,難道她打人嗎?我張眼一看,原來她打的是自己,並說:「我該死,我作了魔鬼的差役,使我又哭又鬧,攪擾你的教會,主啊!我怎變成這種人呢?」之後,那叫她坐下的弟兄,揚聲禱告「主啊!這把火是我點的,若要懲罰,請懲罰我.」他遂即到張太面前求她饒恕；張太說:「弟兄!今天若非你的話,我還不自知生命的程度如何；因著你的話把我生命的程度完全顯露,看見自己多麼不行,多麼敗壞.」一個最壞的聚會,變成了最好的聚會.有靈才能對付罪,有靈,甚麼都有.
\newline
\newline
弟兄姊妹!教會需要聖靈的工作,沒有聖靈就沒有引導,沒有啟示；沒有對付罪的能力,沒有大能的明證.所以一個神僕被神用來作此工作；因為有靈之故,當以斯拉和他們查經,一打開聖經,猶太人都站起來.「站起來」意即整裝待發,立即遵行.
\newline
\newline
今日之基督教會,好像昔日之畢士大池旁躺著的瞎子,瘸子.但是以西結所帶領的教會,成為一支耶和華的軍隊；能與埃及打那屬靈豐富的仗,能與巴比倫打那以神工代替人工的仗；與亞述打那肉體敗壞的仗.
\newline
\newline
各位弟兄姊妹!要站起來,成為一支耶和華的軍隊；但是,一定要聖靈的工作才能站起來；要聖靈的工作,要好像以西結這樣屬靈的人；神把靈裝備給他,不但神揀選他,還要他揀選神；如果只有神揀選你,而你不揀選神,這個揀選不會成功的.
\newline
\newline
求主今天晚上,教我們渴慕聖靈的工作；不但神揀選我們,我們也揀選神；但是選神的人需要付上代價.
\newline
\newline


\newpage

\section{第58屆~港九培靈會~張慕皚~第一講}
\label{sec:455}
\textbf{基督教「時間觀」}
\newline
\newline
連結: \href{https://www.hkbibleconference.org/session-message/view/455}{\texttt{https://www.hkbibleconference.org/session-message/view/455}}
\newline
\newline
\hyperref[sec:471]{< < < PREV SERMON < < <}
~
\hyperref[sec:toc]{[返目錄]}
~
\hyperref[sec:456]{> > > NEXT SERMON > > >}
\newline
\newline
提摩太後書 3:1-3:17
\newline
\begin{longtable}{cl}
\hline
\hline
章節 & 經文 (和合本修訂版)\\
\hline
3:1 & \begin{tabularx}{0.7\textwidth}{X} 你該知道，末世必有艱難的日子來到。 \end{tabularx} \\ \\ \relax
3:2 & \begin{tabularx}{0.7\textwidth}{X} 那時人會專愛自己，貪愛錢財，自誇，狂傲，毀謗，違背父母，忘恩負義，心不聖潔， \end{tabularx} \\ \\ \relax
3:3 & \begin{tabularx}{0.7\textwidth}{X} 沒有親情，抗拒和解，好說讒言，不能節制，性情兇暴，不愛良善， \end{tabularx} \\ \\ \relax
3:4 & \begin{tabularx}{0.7\textwidth}{X} 賣主賣友，任意妄為，自高自大，愛好宴樂，不愛神， \end{tabularx} \\ \\ \relax
3:5 & \begin{tabularx}{0.7\textwidth}{X} 有敬虔的外貌，卻背棄了敬虔的實質，這等人你要避開。 \end{tabularx} \\ \\ \relax
3:6 & \begin{tabularx}{0.7\textwidth}{X} 他們當中有人潛入別人家裡，操縱無知的婦女；這些婦女被罪惡壓制，被各樣的私慾引誘， \end{tabularx} \\ \\ \relax
3:7 & \begin{tabularx}{0.7\textwidth}{X} 雖然常常學習，終久無法達到明白真理的地步。 \end{tabularx} \\ \\ \relax
3:8 & \begin{tabularx}{0.7\textwidth}{X} 從前雅尼和佯庇怎樣反對摩西，這等人也怎樣抵擋真理；他們的心地敗壞，信仰經不起考驗。 \end{tabularx} \\ \\ \relax
3:9 & \begin{tabularx}{0.7\textwidth}{X} 然而，他們沒有進步，因為他們的愚昧必在眾人面前顯露出來，像那兩人一樣。 \end{tabularx} \\ \\ \relax
3:10 & \begin{tabularx}{0.7\textwidth}{X} 但你已經追隨了我的教導、行為、志向、信心、寬容、愛心、忍耐， \end{tabularx} \\ \\ \relax
3:11 & \begin{tabularx}{0.7\textwidth}{X} 以及我在安提阿、以哥念、路司得所遭遇的迫害和苦難。我忍受了何等的迫害！但從這一切苦難中，主都把我救了出來。 \end{tabularx} \\ \\ \relax
3:12 & \begin{tabularx}{0.7\textwidth}{X} 其實，凡立志在基督耶穌裡敬虔度日的，也都將受迫害。 \end{tabularx} \\ \\ \relax
3:13 & \begin{tabularx}{0.7\textwidth}{X} 只是作惡的和騙人的將變本加厲，迷惑人也被人迷惑。 \end{tabularx} \\ \\ \relax
3:14 & \begin{tabularx}{0.7\textwidth}{X} 至於你，你要持守所學習的和所確信的，因為你知道是跟誰學的， \end{tabularx} \\ \\ \relax
3:15 & \begin{tabularx}{0.7\textwidth}{X} 並且知道你從小明白聖經，這聖經能使你因在基督耶穌裡的信有得救的智慧。 \end{tabularx} \\ \\ \relax
3:16 & \begin{tabularx}{0.7\textwidth}{X} 聖經都是神所默示的，於教訓、督責、使人歸正、教導人學義都是有益的， \end{tabularx} \\ \\ \relax
3:17 & \begin{tabularx}{0.7\textwidth}{X} 叫屬神的人得以完全，預備行各樣的善事。 \end{tabularx} \\ \\
[1ex]
\hline
\hline
\end{longtable}
感謝主,帶領培靈會進入第五十八屆事奉,深信神在今年培靈會向我們施恩,許多弟兄姊妹付上代價代禱!這表明不是我們能做些甚麼,而是藉著禱告,聖靈在此工作,使我們能領受神所預備的.
\newline
\newline
我以「基督徒的人生觀」為主題.因為我深信,在今日末世時代中,聖經告訴我們,世界潮流向教會信徒沖激,許多基督徒看人生,是沒有分別的,(提後三章)；但保羅提醒我們說,在末世的時候,必有危險的日子來到.我們不知道下面的經文,是指教會外面,抑或內部的情況:「因為那時人要專顧自己,貪愛錢財,自誇狂傲,謗讟,違背父母,忘恩負義,心不聖潔,無親情,不解怨,好說讒言,不能自約,性情兇暴,不愛良善,賣主賣友,任意忘為,自高自大,愛宴樂不愛神,有敬虔的外貌,卻背了敬虔的實意.」許多人談到這裡的時候,認為是末世時代；這是教會中所看到的情形,因為許多屬神的人,在末世時候,他們的生活方式和非基督徒很難分別.
\newline
\newline
我們要再一次思想,基督徒的人生觀根據聖經的教訓,是一個怎樣的人生觀呢?生活方式與不信的人有甚麼分別呢?如無分別我們就不能成為榮耀的教會了.神的教會是神從世界分別出來的人,我們的生活和人生觀必須和不信的人有所分別；所以今次我就選了人生觀的八方面,與大家在八個早上一同查考聖經.今天早上,我和大家思想的是──「基督徒的時間觀」.我們應如何看時間呢?
\newline
\newline
時間是一樣很難捉摸的事物.奧古斯丁說:「時間是甚麼?若沒有人問我,我知道；如果有人要我解釋,我就不知道.」另一個學者說:「時間是最難下定義的,亦是充滿矛盾的東西.過去的時間,已是不存在,將來的時間,亦未曾來到；而現在的時間呢!當我們要解釋何謂時間,就已經成為過去,像天上的閃電一樣,一出現,就立刻消失.」我們很有同感,時間是很難捉摸的.
\newline
\newline
在四福音中,我們可以看到耶穌基督對時間很重視,也很敏感.耶穌的時間觀應該成為我們基督徒的時間觀.(約翰福音九章四節)耶穌說:「趁著白日,我們必須作那差我來者的工；黑夜將到,就沒有人能作工了.」上文第八章記載耶穌和敵人有很激烈的爭辯,以致敵人要追殺祂.在逃避敵人追殺的時候,耶穌停下來,在路邊醫好一個盲眼的人；後來又可以見到耶穌醫好這人心靈的瞎眼,我們稱此為全人的福音.不過我們要留意的是九章四節,耶穌教導我們對時間的看法:「趁著白日,我們必須作那差我來者的工,黑夜將到,就沒有人能作工了.」
\newline
\newline
(一)時間就是生命
\newline
\newline
耶穌在此告訴我們,時間就是生命.「白日」是代表人還有肉體生命,還有人生的時間,不要等到黑夜來臨的時候.在此「黑夜」代表死亡.趁著還有生命的時候,作那差我來者的工.神起初做人的時候是沒有死亡的,肉體生命沒有時間的限制,人墮落之後,罪和時間聯合起來成為人類厲害的敵人.人一出生,就面向墳墓,年紀愈大,墳墓離我們愈近.罪和時間使我們精神,力量,和美貌消耗.時間就是生命.有一位哲學家說:「生命是時間的累積,浪費時間就是浪費生命,是慢性的自殺.」
\newline
\newline
雖然時間就是生命,但聖經也說,生命不單有量,而且也有質的.耶穌在這裡所講的是量方面.時間是生命的量,而工作就是我們生命的質.聖經講的永生,不但是指量方面,永永遠遠活在下去；其實聖經說到信和不信的人都永遠活下去；而不同的地方是──信的人在天堂永遠享受神賜我們的福樂；而不信的人卻在地獄裡受永遠的痛苦.聖經裡所注重的永生是生命的質素,當然也講到量方面,即時間方面.耶穌講論祂帶給我們豐盛生命就是這生命的質素.主活在世上只不過是三十三年,在這段時間中,祂為我們完成了救恩,祂著重生命的質素,但同時祂沒有忽略生命的量.祂知道自己在世的時日不多,所以就把握每一段時間來事奉神.
\newline
\newline
基督徒的時間觀和非基督徒的有甚麼分別呢?如果二者都同意時間就是生命,基督徒不同的地方就是,我們基督徒的生命是屬主的；因此我們的時間也是主的,我們不是雇工,雇工大部份的時間是自己的,我們是主的奴僕,我們整個生命是屬於主的,時間全由主人來支配.我們這樣奉獻是因為主買贖我們,我們應該全然屬主,但神不強迫我們,所以保羅就這樣說:「我以神的慈悲勸你們,將身體獻上......」這是愛的回應.我們若不將生命獻給主,主有時會任憑我們,這是一種很嚴重的刑罰.有些基督徒說是將生命獻給主,但大部份的時間屬於自己的,這是很常見的矛盾.我們常把時間分成三部份,一部份是在教會的事奉,是屬主的；一部份是在工廠或寫字樓,時間是屬僱主的,第三部份是空閒的時間,是屬自己的.但一個屬神的人卻不能把時間這樣劃分,我們生命屬主,所有的時間也是屬主.(雅四13-17)說:「你們有話說,今天,明天我們要往某城裡去,在那裡住一年,作買賣得利；其實明天如何,你們還不知道.你們的生命是甚麼呢?你們原來是一片雲霧,出現少時就不見了.你們只當說,主若願意,我們就可以活著,也可以作這事,或作那事.現今你們竟以張狂誇口,凡這樣誇口都是惡的.」神就這樣指責信徒企圖將時間掌握自己手中.我們應該將生命和時間都交在主手中.
\newline
\newline
(二)時間就是恩賜
\newline
\newline
神給耶穌三十三年時間,耶穌知道不能再求多一點.所以祂說,祂趁著還有時間就要去做神所交付的工作.從這方面看,可見時間是神的恩賜.神賜給各人的時間都不同；有些人長些,有些人短一些.人生的計劃,是神為我們安排,預備的.(弗二10)說:「我們原是祂的工作,在基督耶穌裡造成的,為要我們行善,就是神所預備叫我們行的.」「工作」在這裡原文的意思可翻作神的精心傑作.像保羅說:我們好像一團泥,在神這位窯匠手中,塑造成一個器皿,合祂使用.神不斷在我們的生命的時間中塑造我們；問題是我們有沒有容許神來塑造.(詩三十一15)說:「我終身的事在你手中.」有些英文譯本譯作:「我一生的時間在你手中.」從耶穌身上就可以見到這真理,祂知道自己的時間,所以祂有時也說祂的時候還沒有到,就不做某事；如在迦拿的筵席中,就是一例.耶穌每一件事都有祂的時候,祂必須按照神的時間,安排而活.
\newline
\newline
基督徒的人生觀,就是時間和生命都是屬於神的,所以要尋求神的心意,看神要我們如何使用祂所賜給我們的時間；在時間中,我們要作甚麼,其實神已經為我們預備好.那麼我們怎樣知道神在我們生命中的心意呢?神的僕人陶恕給我們一些指示.他說:「如果聖經所要我們去做的事情,我們就要去做；聖經禁止的,就不要去做；如果聖經沒有明文規定的；只要我們存心凡事榮耀神,叫人得著益處,我們就去做,這都是在神的旨意中.」
\newline
\newline
(三)時間就是珍寶
\newline
\newline
耶穌說,「趁著白日」.時間是很寶貴的,所以要珍惜.所羅門說:「你趁著年幼,衰敗的日子尚未來到......當記念造你的主.」摩西說:「求你指教我們怎樣數算自己的日子.」保羅說:「要愛惜光陰.」時間是非常寶貴的產業和資源,這種資源和其他的資源不同；這種資源是不能開源節流的,因為時間一過去就沒有了.自然界許多東西都可以控制,但時間的消失是不能控制的.
\newline
\newline
我們不要走到兩個極端,一個就是不好好的把握時間；到頭來,我們就成為時間的工具.時間若是珍寶的話,我們就要好好的投資；這就首先要了解我們一生要做些甚麼.如果神要你一生做家庭主婦,你就投資做家庭主婦；如果神要你做工程師,你就要一生投資做工程師.耶穌很清楚自己來世間要作甚麼,祂說:「趁著白日,我們必須作那差我來者的工.」祂很清楚神差祂到世上來做甚麼.所以祂能將時間好好的投資,耶穌說:「我要按那差我來者的意思行.」祂了解自己一生的目標是甚麼.各位要知道神要我們一生做甚麼,特別年青人,你們或者不知道,今晚回家就要好好的等候在神的前面；問神─祂要我一生投資在甚麼事情上呢?
\newline
\newline
(四)時間就是機會
\newline
\newline
耶穌說,我要趁著時間所帶來的機會,好好的做神要我做的事情.有人在某地方,需要我的醫治,我若不醫治,機會就會失去.
\newline
\newline
弟兄姊妹,特別是年青人,現在有時間和我們父母相處的時候；就要好好孝敬我們的父母,時間一過去,機會就會失去.我們又要把握每一個事奉神的機會,好好的事奉神,不要等那樣,等這樣.我們不知神給我們多少時間,所以機會就要把握,我們如何把握時間給我們的機會呢?
\newline
\newline
第一,我們要編排事情的輕重緩急.神要我們做甚麼,我們就要專心的去做.要將人生時間的次序安排得好,否則我們就不能好好的運用我們的時間.耶穌在(太六33)教導我們說:「先求神的國和神的義,這一切東西都要加給你們了.」這是我們生命優先的次序,神的國和義是優先的.我們做甚麼行業都要將神的國和義放在先,其他的一切可以不要,但神的國和義最寶貴.
\newline
\newline
第二,要過簡樸的生活.人的生活若過份多姿多彩的話,就浪費很多時間.我們有甚麼東西是不必要的,就要除去．許多東西都會浪費我們許多時間．都市生活太富裕,使我們許多時間浪費在不必要的事物上.
\newline
\newline
第三,有效的運用時間.我們必須保持身體健康,身體軟弱的時候,就會花費許多時間,也做不出甚麼來.適當的運動和休息,不是浪費時間的,反而叫人更有效率.
\newline
\newline
第四,要知道自己的恩賜,去做我們能力可以做的事.許多人不自量力,去做不是自己恩賜的事,這樣就要花特別的時間,而且做得不好,壓力又大.
\newline
\newline
第五,要懂得安排一日,或一個月的時間,要計劃如何使用.要做的事情要預先計劃,時間運用就會更充份.
\newline
\newline
最後一點是,我們要學習婉拒別人的要求.自己做得到的,當然要幫助人,但自己能力達不到時,就要有愛心的婉拒.
\newline
\newline
時間就是機會,神給你甚麼機會,就不要等待,要立刻去做.
\newline
\newline
基督徒要明白,我們的時間在神的手中,也要明白神的旨意；然後把握神所給我們每一個機會,每一分,每一秒.
\newline
\newline


\newpage

\section{第58屆~港九培靈會~張慕皚~第二講}
\label{sec:456}
\textbf{基督教「工作觀」}
\newline
\newline
連結: \href{https://www.hkbibleconference.org/session-message/view/456}{\texttt{https://www.hkbibleconference.org/session-message/view/456}}
\newline
\newline
\hyperref[sec:455]{< < < PREV SERMON < < <}
~
\hyperref[sec:toc]{[返目錄]}
~
\hyperref[sec:457]{> > > NEXT SERMON > > >}
\newline
\newline
創世記 1:26-1:28
\newline
\begin{longtable}{cl}
\hline
\hline
章節 & 經文 (和合本修訂版)\\
\hline
1:26 & \begin{tabularx}{0.7\textwidth}{X} 神說：「我們要照著我們的形像，按著我們的樣式造人，使他們管理海裡的魚、天空的鳥、地上的牲畜和全地，以及地上爬的一切爬行動物。」 \end{tabularx} \\ \\ \relax
1:27 & \begin{tabularx}{0.7\textwidth}{X} 神就照著他的形像創造人，照著神的形像創造他們；他創造了他們，有男有女。 \end{tabularx} \\ \\ \relax
1:28 & \begin{tabularx}{0.7\textwidth}{X} 神賜福給他們，神對他們說：「要生養眾多，遍滿這地，治理它；要管理海裡的魚、天空的鳥和地上各樣活動的生物。」 \end{tabularx} \\ \\
[1ex]
\hline
\hline
\end{longtable}
我們應當怎樣按照神的話完成神交付我們的職業呢?(創一26)記載:
\newline
\newline
「神說:我們要照著我們的形像,按著我們的樣式造人；使他們管理海裡的魚,空中的鳥,地上的牲畜,和全地,並地上所爬的一切昆蟲.」28節又記載說:「神就賜福給他們,又對他們說,要生養眾多,遍滿地面,治理這地；也要管理海裡的魚,空中的鳥,和地上各樣行動的活物.」
\newline
\newline
(傳三13)說:「並且人人吃喝,在他一切勞碌中享福,這也是神的恩賜.」有些譯本是這樣譯的:「在他一切工作裡面,他得到滿足,這就是神的恩賜.」
\newline
\newline
(加六10)說:「所以有了機會,就當向眾人行善,向信徒一家的人更當這樣.」
\newline
\newline
最後請看(西三22-25):「你們作僕人的,要凡事聽從你們肉身的主人；不要只在眼前事奉,像是討人喜歡,總要存心誠實敬畏主.無論作甚麼,都要從心裡作,像是給主作的,不是給人作的；因為知道從主那裡必得著基業為賞賜,你們所事奉的乃是主基督.那行不義的,必受不義的報應,主並不偏待人.」
\newline
\newline
有人說,人生好像一本書.神給我們人生的時間,可說是生命的紙張；而我們的工作,就是人生之書的字句.從一本書的文字,可以給我們看到這本書的意思.那麼我們人生的意義從何而找到呢?人生的意義乃是在我們的工作那裡找到的.
\newline
\newline
據統計,今天的人平均壽命可到七十歲,其中有二十年是用來睡覺的.另外二十五年是工作的時間,所餘下的三十五年是用來吃東西,穿衣,聽電話......等等.甚至工作的時間佔我們整個人生的二十年.求主教導我們珍惜這二十年的時間.
\newline
\newline
在這個都市化,工業化的城市裡,許多人不能在工作中找到意義.我們在工作中感到煩躁,沒有味道.這成為今日人類包括許多基督徒在內,對工作的一個很嚴重的問題.在此情形中,我們只有兩種做法:一是改變工作,另外就是去改變我們自己.本人認為,後者是我們需要去做的.因此,當我們感到工作沒有意義的時候,應當怎樣做呢?就是要改變人生的態度,轉換工作不是辦法.如果對人生工作的態度不正確的話,換多少份工對你仍沒有甚麼幫助.我們現在來看著聖經,如何教導我們看職業吧!
\newline
\newline
一．工作是人性的表達
\newline
\newline
在創世記的第一章可以見到,神造人的時候,神把我們造成是必須工作的,這是我們本性的一部份.如果不工作的話,人的本性便不能得以發揮.神造人把人造成必須工作是反映神的本性,因為創世記那裡說,人是按照神那樣,是一位工作者.聖經一開始就講到神創造天地的工作,神不但造天設地,創世記二章八節說,神在東方的伊甸立了一個園子.中文「立了園子」一句可以譯作「種植了一個園子」.神在第一章處,神說有甚麼,就有甚麼；但神種植一個園子的時候,就要像園丁那樣,細心地種植；把園子打理得美麗,好叫人的始祖在其中得以過舒適的生活.
\newline
\newline
神創造工作完成後,在第七日歇了工作,所歇的是創造的工作；神護理那創造物的工作,是繼續下去的,所以神是不停止工作的.人卻六日工作,一日休息,但神不是這樣的.約翰福音五章記載有人指責耶穌在安息日治病,但耶穌說:「我父作事直到如今,我也作事.」(17節)聖經又記載,世界各地各國,大事小事,都在神掌握之中；祂是管理萬國的神,更為我們預備救贖的工作.
\newline
\newline
神造人,好像他自己一樣,是必須工作的.古代希臘把人分做休閒階級和奴隸；前者是貴族,不需要工作的,奴僕就要工作.我們中國人多少也有這種心理,希望不須要工作,享清福.聖經的觀念並不是這樣,主耶穌來到世界作一位工作者,是一個木匠.亞當,夏娃未犯罪前是在伊甸園子作修理看守的工作,就是在園子裡工作.這工作相當繁重,但當時他們未犯罪,工作就會滿有喜樂和滿足感,後來犯了罪,終身勞苦纔得吃.我們要知道,這是人墮落後帶來的惡果,工作本身原不是給人帶來痛苦的.
\newline
\newline
另一方面,我們要避免一個錯誤的觀念,以為工作是沒有意思的,只不過是為了賺錢過生活.許多基督徒認為日常的工作沒有意思,返教會工作才是事奉,才有意思.我們要了解,神造人的時候,是把我們造成必須工作的.
\newline
\newline
二．工作是恩賜的發揮
\newline
\newline
在工作的時候,就是發揮恩賜.廣義的說,恩賜就是(雅一17)所說的:「各種美善的恩賜和各樣全備的賞賜,都是從上頭來的.」我們重生得救後,在教會中事奉,神給我們的才幹,是恩賜；而那與生俱來,神給我們的才幹,是恩賜；而那些與生俱來,神給我們的才能,能在社會中任職的,也是神所給的恩賜；因為各樣美善的恩賜,都是由上頭而來的.以上所講的,我們可以將之分開為事奉的恩賜和職業的恩賜.神是要我們按照祂所給的工作恩賜發揮我們自己.我們要問自己,究竟神給了我們甚麼工作的恩賜?神是按著工作的才幹來用我們的,若要知道自己適合做甚麼工作,進入甚麼行業；你要先知道,你的工作恩賜是甚麼；神會根據這些帶領你,祂不會叫人做那些他才幹以外的工作.這是最重要的,因為人生的意義是在工作中找到的.
\newline
\newline
我們選擇職業的時候,不要以該職位是否能給我們帶來社會地位,帶來更多的金錢等作為條件；而是要看看神給我們的才幹是否適合擔當那個職位.作父母的,也要用這個原則來幫助自己的兒女,發展他們前途；不要勉強,以免誤了他們的一生.懷才不遇,是最可惜的；而能發揮所長,就是最寶貴的.若按著自己的才幹,恩賜來工作的話,就會帶來很大的滿足感.
\newline
\newline
同時,依照這個原則,工作是無分貴賤的,只要按著自己的恩賜來發揮就行.所以無論我們是做甚麼職業,只要靠著神給我們的能力,我們就會找到滿足感.約瑟在埃及做奴隸,他就發揮那奴隸的恩賜；神就使他步步高升,實現自我.
\newline
\newline
三．工作是社會服務
\newline
\newline
(加六10)說:「我們要向眾人行善」意思我們的工作,是一種社會服務和參與.
\newline
\newline
神造人,把我們造成一個群體的人.人必須在群體中成長.我們需要周圍的人,而我們的工作就是要服事周圍的人.基督徒不能單在教會中愛神,基督徒同時也要愛周圍的人.加拉太書那裡說,主內弟兄姊妹是我們優先關心的對象,但也要關心我們周圍的人.
\newline
\newline
聖經中有一些例子,有些人以為在聖殿敬拜神,就可以討神的喜悅.在以賽亞書第一章可以見到,神嚴厲斥責當時的以色列人；因為他們只懂在聖殿中敬拜神,獻祭,守嚴肅會等；而不知關心周圍的人的需要.不在社會中行公義,不為社會上的孤兒,寡婦伸冤辯屈；他們本沒有服務的精神.神所要的是愛神,愛人的結合.人若說愛神,但卻恨弟兄；使徒約翰說,這種愛是虛偽的.我們不能只回教會敬拜而不關心社會.
\newline
\newline
參政,抗議,設立福利中心等等固然是基督徒應有的關懷社會行動；不過更重要的是,基督徒要透過本身的職業來建設我們的社區.基督徒如果在本身的職業做得好的話,那就是最有效的社會關懷.許多人覺得工作枯燥無味,就是因為沒有這種真理的體認.工作就是服務人群.
\newline
\newline
在今日的工業社會中,注重分工,許多時看不到工作的意義,如何服務人群更看不到,如果人人認識自己的工作是對國家社會有很大的貢獻；人的工作態度就會改變.舉一個例,做生意的人不應單顧賺錢,還應有社會服務的責任；不然,容我說一句,他們的動機和出黃色刊物的人差不多,因為大家的動機都是賺錢.所以做生意的人要留意,自己公司的工作或工廠的出品,對社會的安定繁榮是否有幫助?是否給別人帶來許多就業的機會.在政府裡,在報界中,基督徒都要求神給我們恩典做得好；在自己權力範圍中影響本身的行業,這是最好的社會關懷和參與.
\newline
\newline
四．工作是國度事奉
\newline
\newline
歌羅西書第三章告訴我們說,工作是一種國度的事奉.基督徒不單要服務社會,也要建立神的國度.我們的工作是一種事奉,若能這樣存心去想,工作纔具有價值呢!
\newline
\newline
歌羅西書第三章那裡,是保羅對當時做奴隸的人說話.作奴隸是很慘的,但保羅對他們說:「要存敬畏的心去做,好像為主而做,而不是為人而做一樣.因為你們所事奉的仍是你們的主.」做奴隸的工作也是事奉神.所以我們今天應該深深的感受到,任何一種正當的職業,我們都要存事奉神的心去做.
\newline
\newline
(弗四1)說:「既然蒙召,行事為人,就當與蒙召的恩相稱.」意思是說,神呼召我們去過一個全時間事奉神的生活,整個生活都是事奉神,包括職業和工作.馬丁路德告訴我們,工作無聖俗之分；不是在教會做的才是聖工,而在教會以外做的,就是世俗之工.馬丁路德說,我們不應有這樣的看法.工作是聖還是俗是在於是否有神的呼召.他又說,如果一個修女,沒有神的呼召,她的修道和服務工作是最世俗不過的；同時一個女工,有神的呼召,在別人家庭中服事人,這女工的工作就是最神聖的.馬丁路德所講的是很對的.保羅說做奴隸的工作,也是服事神的工作；何況我們不是做奴隸,而是做文員,做工人等,就更應該是事奉神.做家庭主婦,其實也是一個偉大的事奉；養育兒女,可以對社會國家有貢獻；例如聖經中的哈拿,他養育撒母耳,為撒母耳祈禱,她屬靈的生命影響撒母耳；後來撒母耳在政治上,社會上,宗教上改變了當時的世代.歷史上許多改變當時世代的偉人,如奧古斯丁,約翰衛斯理等；都有一位影響他們的母親,這些母親忠心她們的工作,認定這些工作是事奉神.今日我們無論做甚麼職業,都要認定我們乃是在那裡事奉主,所做的都是主的工作；這樣的話,我們真正的主人,正如聖經所講的,乃是主耶穌基督.無論你做縫紉工人,或在餐室當侍應生或是其他工作；你所做的是神透過你的上司給你作的.所以,所謂無意義的工作,就成為有意義了；因為那是服事神的工作,值得全心,悉力以赴.既然是神的工作,我們就會盡心而不擔心；因為成功與失敗,都是神的事；即如僱員不會擔心生意如何,因為生意是老板的.
\newline
\newline
我們將來是要面對審判,審判的根據是甚麼呢?根據是我們所做的是否為主而做.工作為主而做的意思是,第一,體會到所做的,是主自己的工作；第二,在工作中要按照聖經的原則榮神益人；第三,必須在本身行業中見證主；第四,工作要出於愛神和愛人.如果我們能這樣工作的話,我們的工作是滿有意義的.
\newline
\newline


\newpage

\section{第58屆~港九培靈會~張慕皚~第三講}
\label{sec:457}
\textbf{基督教「休閒觀」}
\newline
\newline
連結: \href{https://www.hkbibleconference.org/session-message/view/457}{\texttt{https://www.hkbibleconference.org/session-message/view/457}}
\newline
\newline
\hyperref[sec:456]{< < < PREV SERMON < < <}
~
\hyperref[sec:toc]{[返目錄]}
~
\hyperref[sec:458]{> > > NEXT SERMON > > >}
\newline
\newline
申命記 5:12-5:15
\newline
\begin{longtable}{cl}
\hline
\hline
章節 & 經文 (和合本修訂版)\\
\hline
5:12 & \begin{tabularx}{0.7\textwidth}{X} 「『當守安息日為聖日，正如耶和華－你神所吩咐的。 \end{tabularx} \\ \\ \relax
5:13 & \begin{tabularx}{0.7\textwidth}{X} 六日要勞碌做你一切的工， \end{tabularx} \\ \\ \relax
5:14 & \begin{tabularx}{0.7\textwidth}{X} 但第七日是向耶和華－你的神當守的安息日。這一日，你和你的兒女、僕婢、牛、驢、牲畜，以及你城裡寄居的客旅，都不可做任何的工，使你的僕婢可以和你一樣休息。 \end{tabularx} \\ \\ \relax
5:15 & \begin{tabularx}{0.7\textwidth}{X} 你要記念你在埃及地作過奴僕，耶和華－你的神用大能的手和伸出來的膀臂領你從那裡出來。因此，耶和華－你的神吩咐你守安息日。 \end{tabularx} \\ \\
[1ex]
\hline
\hline
\end{longtable}
「當照耶和華你神所吩咐的,守安息日為聖日.六日要勞碌作你一切的工,但第七日,是向耶和華你神當守的安息日.這一日你和你的兒女,僕婢,牛,驢,牲畜,並在你城裡寄居的客旅；無論何工都不可作,使你的僕婢可以和你一樣安息.你也要記念你在埃及地作過奴僕.耶和華你神用大能的手,和伸出來的膀臂,將你從那裡領出來；因此,耶和華你的神吩咐你守安息日.」(申五12-15).請再看(出廿8-11):「當記念安息日,守為聖日.六日要勞碌作你一切的工.但第七日是向耶和華你神當守的安息日.這一日你和你的兒女,僕婢,牲畜,並你城裡寄居的客旅,無論何工都不可作.因為六日之內,耶和華造天,地,海,和其中的萬物,第七日便安息.所以耶和華賜福與安息日,定為聖日.」
\newline
\newline
申命記和出埃及記,兩本書都記載「要守安息日」,但其中有兩個不同卻有關連的原因.申命記是把守安息日的原因和神帶領以色列人出埃及的事連在一起說的；但出埃及記卻把守安息日的原因和神的創造連在一起,這幫助我們去思想.請看另一段經文:「使徒聚集到耶穌那裡,將一切所作的事,所傳的道,全告訴祂.祂就說,你們來同我暗暗的到曠野地方歇一歇.這是因為來往的人多,他們連喫飯也沒有工夫.他們就坐船,暗暗的往曠野地方去.」(可六30-32)
\newline
\newline
在二十世紀的時代,我們不能不承認我們有很多休閒的時間,但我們往往很容易忽略了.若我們仔細去數算,我們發覺在我們手上實在很多休閒的時間,然而,我們又聽到不少弟兄姊妹埋怨時間不夠,沒有時間事奉神......等等,這都顯示我們沒有智慧運用時間.另一方面,我們又似乎在聖經裡找不著如何運用休閒的真理.不少人認為聖經記載的時代,是以工作為中心的時代,沒有提及如何運用休閒的時間.其實聖經記載了很多關於運用休閒時間的原則；而且現今的時代,慢慢變成為一個以休閒為中心的社會.不過,我們又發覺在休閒的時代裡,犯罪率也愈來愈多,這和休閒時間增加有很密切關係.
\newline
\newline
「休閒時間」是甚麼意思?有人認為非工作時間就是休閒；但研究休閒時間的人,認為非工作時間不一定是休閒；有意義而又帶來一個豐盛人生的非工作時間,才是休閒的時間.有些年青人在非工作時間,幹一些不法勾當,這不算是休閒.也有人寓娛樂於工作中,工作與娛樂分不開,就如某些政治家一般；甚至科學家也廢寢忘餐地不停工作,他們感到很快樂.在這樣的時代,究竟聖經如何教導我們運用休閒呢?
\newline
\newline
一．屬靈時間
\newline
\newline
基督徒要認定神賜予的休閒時間,乃是屬靈的時間.在出埃及記和申命記中提及安息日,這是神所稱為「聖」日.在新約時代,我們守的主日,就是主的日子.屬靈時間的意思,就是在神主權之下,我們過日子的時間.因此,我們必須除去一個重要的誤解,就是休閒的時間是不屬靈.存有這樣誤解的人,容易產生工作狂的傾向,這是一種病態.這種人若沒有工作做,便覺得痛苦.所以我們要好好擁有休閒時間而禱告；求神幫助我們過一個有意義的休閒,好好享受這個屬靈的時間.
\newline
\newline
神呼召我們過一個全時間的基督徒生活,我們每一分,秒都是屬於神的；所以,我們的休閒時間也在耶穌基督的主權之下.不然,我們很容易在休閒的時間裡,幹不法的勾當.大衛的例子成為我們的鑑戒,他有不少休閒的時間,所以他很遲起床,結果他受到引誘犯罪.參孫是另一個例子,他是一個好色之徒,本來士師是有很多工作做的,但他卻有不少時間接觸異性,這都製造受試探的機會.我們當要求神幫助我們好好運用休閒；否則,我們容易在這些空閒時間裡犯罪,離棄神.還有,我們必須在主再來時,向神交賬.
\newline
\newline
二．心靈安息的時間
\newline
\newline
「六日要勞碌作你的工作.但第七日是向耶和華你神當守的安息日.」(出廿9-10)
\newline
\newline
「安息」,是指心靈裡的安息.猶太人把他們安息日和假期的娛樂連結在一起.他們過節期是連群結隊,一齊吃喝,歡喜快樂.所以,猶太人把敬拜,娛樂和休息是連在一起的.不過,我們今日不再守安息日,而守主日.因為聖經說:「所以不拘在飲食上,或節期,月朔,安息日,都不可讓人論斷你們.這些原是後事的影兒.那形體卻是基督.」(西二16-17).我們因為耶穌基督的救恩,而進入那真正的安息裡.
\newline
\newline
(出廿11):「因為六日之內,耶和華造天,地,海,和其中的萬物,第七日便安息.所以耶和華賜福與安息日,定為聖日.」神的休息並非靜止的休息.耶穌也曾說:「我父作事直到如今,我也作事.」(約七17),所以耶穌在安息日也醫病.神六日工作,但在第七日祂仍主持天地萬物運作的工作,祂乃進入豐盛生活的工作裡.而申命記更把出埃及的事,連在一起說,因為這預表了耶穌基督來到世上拯救罪人,以致我們能夠進入祂的安息裡.這種神創造的安息,和出埃及入迦南的安息,都預表了我們在神裡面心靈的安息,那就是和諧,喜樂,和平安的生活,藉著我們的敬拜,我們得著心靈和諧的安息.
\newline
\newline
世上不少人都追求心靈上的安息,但他們得不著.因為他們追求的方向錯誤,以為安靜地度假便可安寧,其實這都不行的.甚至基督徒也不例外,不少人把世俗的觀念帶進教會聚會裡；把教會氣氛變成世俗娛樂化,結果失去神的安息.惟有一個以神為中心的聚會,我們心靈才能得著安息.
\newline
\newline
三．恢復身體精神的時間
\newline
\newline
「我願你凡事興盛,身體健壯,正如你的靈魂興盛一樣.」(約翰三書 2).我們不單要追求心靈的安息,也要尋求我們身體,精神得到安息.我們的身體乃是聖靈的殿,我們必須愛惜.
\newline
\newline
如何恢復身體,精神的力量呢?如果一個人的工作是較用腦的,他需要多運動；而勞力的人則需要休息,而不是去運動.因此,我們要好好運用休閒時間,恢復身體,精神的力量.
\newline
\newline
四．生活有意義的時間
\newline
\newline
申命記把出埃及的事蹟連在安息日一起說,是預表了耶穌基督來到世上拯救罪人,以致我們能夠進入祂的安息裡；這表示我們都可因信基督,而進入那豐盛的生活裡.「我來了,為要叫人得生命並且得的更豐盛.」(約十10).我們所過的生活,必須有意義的去追求.
\newline
\newline
如何能過一個有意義的生活呢?第一,我們必須有娛樂.在迦拿婚宴裡,基督與朋友同歡樂,又飲又食(約二1-11).神是厚賜百物的神,祂願意我們享受祂所賜的東西,過豐盛享受的人生.基督徒的娛樂,是有一定原則的；首先,我們要考慮這種娛樂是否能叫我們重新得力.其次,我們乃是享受神的恩典,而不是放縱情慾.此外,我們的娛樂會否損害我們身心靈的健康呢?還有,我們的娛樂會否令人跌倒?對人有沒有幫助呢?並且,我們的娛樂是否會控制我們呢?又是否會僭奪了我們家庭生活呢?
\newline
\newline
五．事奉的時間
\newline
\newline
「事奉」這字與「敬拜」是同一個字,所以我們的事奉和敬拜是一樣的.我們要用非辦公時間,向同事傳福音.耶穌基督在畢士大的池子旁邊,醫治有病的人,他在休閒的時間裡,不忘工作.(約五1-9)
\newline
\newline
教會要幫助弟兄姊妹過一個休閒的生活,教導他們如何參與事奉,如何娛樂.這才是一個全人的教導.這也是教會需要去教導弟兄姊妹的事工,以致弟兄姊妹在休閒時間裡,更能榮耀神,事奉神.
\newline
\newline


\newpage

\section{第58屆~港九培靈會~張慕皚~第四講}
\label{sec:458}
\textbf{基督教「財富觀」}
\newline
\newline
連結: \href{https://www.hkbibleconference.org/session-message/view/458}{\texttt{https://www.hkbibleconference.org/session-message/view/458}}
\newline
\newline
\hyperref[sec:457]{< < < PREV SERMON < < <}
~
\hyperref[sec:toc]{[返目錄]}
~
\hyperref[sec:459]{> > > NEXT SERMON > > >}
\newline
\newline
提摩太前書 6:6-6:10
\newline
\begin{longtable}{cl}
\hline
\hline
章節 & 經文 (和合本修訂版)\\
\hline
6:6 & \begin{tabularx}{0.7\textwidth}{X} 其實，敬虔加上知足就是大利。 \end{tabularx} \\ \\ \relax
6:7 & \begin{tabularx}{0.7\textwidth}{X} 因為我們沒有帶甚麼到世上來，也不能帶甚麼去； \end{tabularx} \\ \\ \relax
6:8 & \begin{tabularx}{0.7\textwidth}{X} 只要有衣有食，我們就該知足。 \end{tabularx} \\ \\ \relax
6:9 & \begin{tabularx}{0.7\textwidth}{X} 但那些想要發財的人就陷在誘惑、羅網和許多無知有害的慾望中，使人沉淪，以致敗壞和滅亡。 \end{tabularx} \\ \\ \relax
6:10 & \begin{tabularx}{0.7\textwidth}{X} 貪財是萬惡之根。有人因貪戀錢財而背離信仰，用許多愁苦把自己刺透了。 \end{tabularx} \\ \\
[1ex]
\hline
\hline
\end{longtable}
(提前六6-10):「然而敬虔加上知足的心,便是大利了.因為我們沒有帶甚麼到世上來,也不能帶甚麼去.只要有衣有食,就當知足.但那些想要發財的人,就陷在迷惑,落在網羅,和許多無知有害的私慾裡,叫人沉在敗壞和滅亡中.貪財是萬惡之根,有人貪戀錢財,就被引誘離了真道,用許多愁苦把自己刺透了.」(17-19):「你要囑咐那些今世富足的人,不要自高；也不要倚靠無定的錢財,只要倚靠那厚賜百物給我們享受的神.又要囑咐他們行善,在好事上富足,甘心施捨,樂意供給人；為自己積成美好的根基,預備將來,叫他們持定那真正的生命.」
\newline
\newline
聖經在講到錢財的時候,往往提醒我們要十分小心.因為金錢一方面的確很有用,可以成就很多事；但如果用得不當,也可能毀掉我們的生命.那麼,基督徒應該怎樣看錢財呢?
\newline
\newline
一．錢財是神的賞賜
\newline
\newline
17節神是那厚賜百物給我們享用的,這當然也包括了錢財.從創業的角度來看,一切資源都是屬於神.(詩廿四1):「地和其中所充滿的,世界和住在其間的,都屬耶和華.」一切都是神的,神是萬有之首,是由神來的.神把祂所擁有的賜給人,並且交託人管理.因此,財富是神所賜的,而神也容許人有私產.十誡中的不可偷盜,就包含這個意思,因為這些私產是屬於別人的,我們就不能霸佔.
\newline
\newline
當神容許世上有貧,富的分別,利未記講到獻祭的時候,特別說明有錢的可以獻牛羊,但貧窮的則可以獻班鳩,雛鴿.在獻祭的制度中,神也容許有貧富的分別.聖經也講到神使人貧窮,也使人富足；他使人卑微,也使人高貴.人是神的管家,但只是暫時的,今生的；將來我們要向神交賬,所以我們要小心,要作神百般恩賜的好管家.我們所有的一切,生命,時間,財富,兒女,一切都是神所賜的；我們是神所給的這一切的管家.
\newline
\newline
錢財是神的賞賜,是神賞給人的福.亞伯拉罕,約伯,所羅門等,他們都得到神所賜的財富.(箴十22):「耶和華所賜的福,使人富足,並不加上憂慮.」我們從神所得的富足和世人所得的不同；世人多有錢財,就多有憂慮；但我們從神所得的富足,是不加上憂慮的.當我們殷勤作工的時候,神便會將財富賜給我們,而我們也可以把它積蓄起來.
\newline
\newline
不過,在神眾多的賞賜中,錢財並不是最好的；其他很多神所賞賜的,比錢財更好.我們千萬不可在價值觀上產生錯誤的觀念；以為某人有錢,必然是神把最好的賜給他.但是某人沒有錢,可能神把更好的賞賜給他.在聖經的價值觀上,錢財並不是最好的.
\newline
\newline
箴言說,智慧比金子更寶貴.(箴十六16)說:「得智慧勝似得金子,選聰明強如選銀子.」公義在價值觀上也比錢財寶貴.(箴十六8)說:「多有財利,行事不義；不如少有財利,行事公義.」另一方面,愛情也比金子更寶貴,雅歌書就有這樣的觀念.更重要的,神的話語比極多的精金更寶貴,比蜂房下滴的蜜更甘甜.馬丁路德說:「財富是神給人最渺小的禮物,它遠比不上神的話語,也比不上身體的美貌,健康,或是頭腦的恩賜；如智慧,知識,才能等.然而人卻日夜為錢勞碌,所以神便把錢財給那些愚蠢的人,而且只給他們錢財.」或者馬丁路德的話較為偏激,但他的價值觀和聖經的很符合.
\newline
\newline
神最大的恩賜,就是把耶穌基督給我們,祂是活水,是生命的糧.我們得到祂,就得到一切.耶穌是遠超乎神所給我們的一切恩賜之上.
\newline
\newline
二．財富可能帶來嚴重危機
\newline
\newline
留意聖經說到貪財所帶來的禍害,而不是金錢本身.金錢本身是中性的,但是如果我們對金錢有不正確的態度；誤用金錢,這樣金錢便會帶來聖經所說的一切禍害.當我們貪財,以錢財為我們的偶像,為我們人生追求的目標；若我們把一生的生命,時間都投資下去的話,這種貪心便帶來無窮的禍害.
\newline
\newline
很多時候,金錢會成為我們跟從主的攔阻,少年的官就是一個很好的例子；他為了不願放下所擁有的極多財富去跟從主；結果便憂憂愁愁的離開了耶穌.所以,錢財會叫我們遠離神.無知的財主所想到的,也只是他自己；在他所講的說話,他沒有提及神或別人；他所講的,只是自己.
\newline
\newline
另外,一個貪財的人,也會使自己與弟兄分離,就是叫人與人分離.財主和拉撒路在這方面給我們很大的教訓.財主天天奢華宴樂；但他門外的乞丐拉撒路,卻只吃他所吃剩的東西.他是自私的,從來沒有想到他門口那個有需要的窮人,這是他的罪.結果他死後到了陰間受苦.雖然財主沒有作甚麼壞事；但是他那種自私,在神看來已經是罪.
\newline
\newline
還有,錢財會使人陷於迷惑和愁苦的網羅中,將我們刺透.所以,如果我們不善用金錢,我們不單享受不到金錢的好處,還會讓金錢毀了我們的一生.世界上有許多人,就是受愛財所累.
\newline
\newline
三．貧富決定於人的態度
\newline
\newline
貧富不是在於我們擁有多少,乃是在於我們對金錢的態度.聖經說,貪財是萬惡之根.第一樣我們要除去的,就是貪心.我們知道,進取心是好的,但是要侵佔別人,這種自私的心,就是貪心,他會用不正當的手段去得到財富,這樣的貪心就是罪.一個貪心的人,永遠不會滿足.
\newline
\newline
但是第六節卻給我們一個積極的教導:「然而敬虔加上知足的心,便是大利了.」知足的人,會看見自己所擁有的,但不知足的人,只看見別人所擁有的.所以知足的人,永遠不會貧窮,而不知足的人,永遠不會富足.
\newline
\newline
聖經說,敬虔加上知足,就是大利.敬虔就是愛主,有主就有一切,我們一切的好處都不在神以外,這是敬虔的人才會有的態度.神既不愛惜他的兒子,為我們眾人捨了,豈不也把萬物和他一同白白的賜給我們麼?所以我們若是愛主,就會擁有一切.
\newline
\newline
箴言三十章記載亞古珥有一個祈求,他求神使他不貧窮,也不富足,只賜給他需用的飲食.因為他若是太貧窮,就會偷竊；但若是太富足,就會不認神；所以最好是不窮也不富,有衣有食就夠了.這也可以成為我們的禱告,我們應該學習像主那樣,過一個簡樸的生活,有日用的所需就好了.培根形容金錢就好像行軍的背囊那樣,太重會走不動,太輕則不能維生.求神幫助我們對金錢有正確的態度.
\newline
\newline
四．財富必須善用
\newline
\newline
我們要運用財富的時候,必須認識財富是最好的僕人；但也是最可怕的主人.我們要成為一個佔有金錢的人,而不是成為一個被金錢佔有的人.很多人的失敗,是在於他們被金錢支配,以金錢為人生的目標；而不是以金錢為工具去使用.我們要懂得享用神給我們的金錢.卻不要自私地去揮霍,像那無知的財主那樣.我們要學習保羅隨遇而安的態度,懂得怎樣處豐富,也懂得怎樣處卑賤.神給我們富足,我們可以享受,也分給有需要的人；神讓我們貧窮,我們也不發怨言.耶穌在地上生活的時候,也是這樣.耶穌在迦拿的婚筵中,和別人一起享用歡樂,祂也應邀去某些人家中吃飯,作客.可是耶穌自己卻連枕首的地方都沒有.
\newline
\newline
我們要用錢財幫助有需要的人,(提前六18-19):「又要囑咐他們行善,在好事上富足；甘心施捨,樂意供給人；為自己積成美好的根基,預備將來,叫他們持定真正的生命.」我們要在神面前作個富足的人,把我們的一切投資在永恆的事上.耶穌在不義的管家的比喻上,講到這個不義的管家,雖然不義,尚且懂得為永恆打算,更何況我們呢?
\newline
\newline
講到金錢的奉獻,我們只以舊約的什一奉獻作為原則,這其實是最起碼的；但我們在新約的恩典時代,更要學習完全的奉獻；把自己的一切,連同生命都奉獻給主.神不是看我們獻多少,祂乃是看我們對主的心,和所盡的力.耶穌在稱讚窮寡婦的時候,所給我們的準則就是這樣；到底我們為自己留下多少,又為神獻上多少呢?求主幫助我們懂得運用今世所得的財富,去為神盡力,為自己的永恆打算.
\newline
\newline


\newpage

\section{第58屆~港九培靈會~張慕皚~第五講}
\label{sec:459}
\textbf{基督教「成功觀」}
\newline
\newline
連結: \href{https://www.hkbibleconference.org/session-message/view/459}{\texttt{https://www.hkbibleconference.org/session-message/view/459}}
\newline
\newline
\hyperref[sec:458]{< < < PREV SERMON < < <}
~
\hyperref[sec:toc]{[返目錄]}
~
\hyperref[sec:460]{> > > NEXT SERMON > > >}
\newline
\newline
歷代志下 1:1-1:17
\newline
\begin{longtable}{cl}
\hline
\hline
章節 & 經文 (和合本修訂版)\\
\hline
1:1 & \begin{tabularx}{0.7\textwidth}{X} 大衛的兒子所羅門鞏固他的國度；耶和華－他的神與他同在，使他極其尊大。 \end{tabularx} \\ \\ \relax
1:2 & \begin{tabularx}{0.7\textwidth}{X} 所羅門吩咐全以色列，就是千夫長、百夫長、審判官、全以色列的眾領袖和族長前來。 \end{tabularx} \\ \\ \relax
1:3 & \begin{tabularx}{0.7\textwidth}{X} 所羅門率領全會眾往基遍的丘壇去，因那裡有神的會幕，就是耶和華的僕人摩西在曠野所造的。 \end{tabularx} \\ \\ \relax
1:4 & \begin{tabularx}{0.7\textwidth}{X} 只是神的約櫃，大衛已經從基列‧耶琳接到他所預備的地方，因他曾在耶路撒冷為約櫃支搭了帳幕， \end{tabularx} \\ \\ \relax
1:5 & \begin{tabularx}{0.7\textwidth}{X} 把戶珥的孫子，烏利的兒子比撒列所造的銅壇擺在基遍耶和華的會幕前。所羅門和會眾求告耶和華。 \end{tabularx} \\ \\ \relax
1:6 & \begin{tabularx}{0.7\textwidth}{X} 所羅門上到耶和華面前會幕的銅壇那裡，在壇上獻一千祭牲為燔祭。 \end{tabularx} \\ \\ \relax
1:7 & \begin{tabularx}{0.7\textwidth}{X} 當夜，神向所羅門顯現，對他說：「你願我賜你甚麼，你可以求。」 \end{tabularx} \\ \\ \relax
1:8 & \begin{tabularx}{0.7\textwidth}{X} 所羅門對神說：「你曾向我父親大衛大施慈愛，使我接續他作王。 \end{tabularx} \\ \\ \relax
1:9 & \begin{tabularx}{0.7\textwidth}{X} 耶和華神啊，現在求你實現向我父親大衛所應許的話；因你立我作這百姓的王，他們如同地上的塵沙那樣多。 \end{tabularx} \\ \\ \relax
1:10 & \begin{tabularx}{0.7\textwidth}{X} 現在，求你賜我智慧聰明，好在這百姓面前出入；不然，誰能判斷你這麼多的百姓呢？」 \end{tabularx} \\ \\ \relax
1:11 & \begin{tabularx}{0.7\textwidth}{X} 神對所羅門說：「你有這心意，不求資財、豐富、尊榮，也不求滅絕恨你之人的性命，又不求長壽；我既立你作我百姓的王，你只求智慧聰明，好審判我的百姓， \end{tabularx} \\ \\ \relax
1:12 & \begin{tabularx}{0.7\textwidth}{X} 我必賜你智慧聰明，也必賜你資財、豐富、尊榮，在你以前的列王未曾有過，在你以後也不會再有。」 \end{tabularx} \\ \\ \relax
1:13 & \begin{tabularx}{0.7\textwidth}{X} 於是，所羅門從基遍丘壇會幕前回到耶路撒冷，治理以色列。 \end{tabularx} \\ \\ \relax
1:14 & \begin{tabularx}{0.7\textwidth}{X} 所羅門聚集戰車騎兵；他有一千四百輛戰車，一萬二千名騎兵，安置在屯車城，在耶路撒冷的王那裡。 \end{tabularx} \\ \\ \relax
1:15 & \begin{tabularx}{0.7\textwidth}{X} 王在耶路撒冷使金銀多如石頭，香柏木多如謝非拉的桑樹。 \end{tabularx} \\ \\ \relax
1:16 & \begin{tabularx}{0.7\textwidth}{X} 所羅門的馬是從埃及和科威運來的，是王的商人按著定價從科威買來的。 \end{tabularx} \\ \\ \relax
1:17 & \begin{tabularx}{0.7\textwidth}{X} 他們從埃及進口戰車，每輛六百舍客勒銀子，馬每匹一百五十舍客勒；赫人眾王和亞蘭諸王的戰車和馬，也是經由他們的手出口的。 \end{tabularx} \\ \\
[1ex]
\hline
\hline
\end{longtable}
今天我們要從歷代志下第一章裡面,思想基督徒對成功應該有的觀念.追求成功是人人都有的心理,但可惜自人類犯罪之後；人類追求成就便往往成了不正當的野心.基督徒所求的成就,如果是為榮神益人的話,那是好的；但如果是出於自私,要征服別人；奴役別人的話；那是不該有的野心.
\newline
\newline
另外,人對成功的標準,也是各有不同.比方有人打劫得手之後,同行可能視他為成功；但社會人士卻視他為敗類.基督徒和非基督徒成功的觀念,也是不同的.今天不少基督徒之所以失敗,是因為沒有從神那裡而來成功的觀念.我們以世俗對成功的觀念來衡量自己的成就,結果這就成為我們最大的失敗.
\newline
\newline
如果說聖經中的成功人物；很多人都會想到所羅門；到底所羅門是怎樣成功呢!雖然所羅門的晚年,因為沉迷女色,以致一敗塗地.但無論如何,他確有過輝煌的成就,是神所賜給他的.今天我們要有所羅門成功之處,來作為我們的借鏡.大衛有很多兒子,但神特別賜福給所羅門；使他國位堅固,且甚為尊大.為甚麼呢!
\newline
\newline
一．他是謙卑的人
\newline
\newline
成功的人是謙卑的.我們可以把所羅門和大衛的其他兩個兒子作比較.押沙龍在他父在位的時候,已經驕傲,要奪王位,結果神把他丟棄,最後他是凄涼的死在樹上.另一是亞多尼雅,他想趁父親年老時奪取王位,但結果他不能繼承父位.因為一方面神是已經定了把王位傳給所羅門,但另一方面也由於他們太驕傲.
\newline
\newline
但是,所羅門從沒有爭奪過王位；只有先知拿單迫切的去求王,要早日膏立這兒子為王.所羅門沒有爭奪王位,但是神把王位賜給他.在箴言裡面,所羅門多次提到驕傲的人,是神所恨惡的,只有謙卑的人才是蒙福的人.比方說(箴十一2)「驕傲來,羞恥也來；謙遜人卻有智慧.」(箴十六18)「驕傲在敗壞以先,狂心在跌倒之前.」(箴廿九23)「人的高傲,必使他卑下；心裡謙遜的,必得尊榮.」
\newline
\newline
所羅門知道甚麼是謙卑,也知道謙卑的人的結果.耶穌也清楚的告訴我們,凡自高的,必降為卑；自卑的必升為高.一個成功的人,在神眼中必須是一個謙卑的人.慕安德烈說:「謙卑就是心如止水,無任何的奢望.」別人無論如何對待他；都不會不滿或沾沾自喜.
\newline
\newline
一個肯投靠神,在神面前謙卑的人；在神看來已經是成功的.奧古斯丁說:「驕傲叫天使變為魔鬼,而謙卑叫魔鬼成為天使.」成熟的基督徒必然是謙卑的,就像成熟的麥子那樣；愈是成熟,頭便愈向下垂.謙卑的人是忘記自己,高舉基督的.他在自我的評估中,會感到自己的渺小,這是一個成功的人.箴言講到神所憎惡的有七種人,頭一種就是眼目高傲的,這是神所最憎惡的.所以驕傲的人,在神面前就已經是失敗的,不能使用的.
\newline
\newline
摩西四十歲之時,心高氣傲；以為自己有能力去拯救以色列人.但是神沒有用他,卻把他帶到曠野,在四十年的牧羊生涯中；讓摩西學會謙卑之後,神才呼召他,使用他.大委身為一國之君,但是他卻稱自己為一隻小羊；耶和華才是引導和牧養他一生的牧者.
\newline
\newline
一個謙卑的人,他會深深認識到神的恩典.所羅門深知王位是神給他的,他是在神的恩典中生活的.今天,我們不單要知道這得救是本乎恩,我們一生的生活也活在神的恩典中；祂把我們本來毫無價值的一生,變為有價值.保羅說:「我今日成了何等樣的人,是蒙神的恩才成的.」我們要在神的恩典中尋求成就.神的恩典加上我們的努力,才能成就.
\newline
\newline
二．他是敬虔的人
\newline
\newline
歷代志下一2-6,講到所羅門帶領千夫長,百夫長,審判官,首領,族長等人,到邱壇那裡獻一千畜牲為燔祭.當他得王位之後,他做的頭一件事,就是帶領國內一班重要的官長,一同向神敬拜.在神看來,一個成功的人,就是一個敬虔的人.
\newline
\newline
所羅門所獻的燔祭,原來是把祭物燒盡,讓馨香之氣升上去的意思.這代表基督徒將自己奉獻給神,這是最美好的敬拜.保羅說,要將身體獻上,當作活祭,就是這個意思.這是理所當然,討神喜悅的事奉；而這事奉基本上就是一種敬拜.一個真正敬畏神的人,是必然把自己獻上,為主而活的.保羅講到活祭的時候,他說我們是向罪和自己死,但向神卻是活著的.
\newline
\newline
大衛就是這樣把自己的兒子獻給神,(代上廿九19)說:「又求你賜我兒子所羅門誠實的心,遵守你的命令,法度,律例.」大衛為所羅門所祈求的,是要他的兒子跟隨主,活在神的旨意中.所羅門在他一生的大部份日子中,得蒙神大大賜福；是因為他樂意行在神的旨意中,照神的旨意過活.這樣的人在神眼中是成功的,而且這成功是必然的,肯定的.
\newline
\newline
聖經中有很多這樣的例子.撒母耳的母親哈拿──她謙卑等候在神面前,得神大大賜福；後來她又願意把自己頭生的長子獻回給神.但以理的一個朋友為了活在神的旨意中,不聽王命去拜金像；雖然被扔在火湖中,但是火湖卻是最安全的地方.所以遵行神旨意的人,必然是亨通,順利,和最安全的.但以理自己也是一樣,他活在神的旨意中；雖然被扔在獅子坑中,但仍然得到安全.反觀約拿,他因為不服從神的旨意,以致情況每況愈下,十分可憐!
\newline
\newline
可是,當我們切實遵行神旨意的時候,我們會發覺環境對我們不利；反而那些世人,卻事事亨通.不過,詩篇三十七篇卻在這方面給我們很好的教導.詩人說:「不要為作惡的,心懷不平；也不要向那行不義,生出嫉妒.因為他們如草快被割下,又如青菜快要枯乾.」第五節:「當將你的事交託耶和華；並倚靠祂,祂就必成全.」第七節:「你當默然倚靠耶和華,耐性等候祂；不要因那道路通達的,和那惡謀成就的,心懷不平.」
\newline
\newline
三．他是先求神的國和義的人
\newline
\newline
(代下一7-11)講到神允許所羅門求甚麼就給他,所羅門於是求神給他智慧,好讓他能治理國民；使人民得益處,使神得著榮耀.神十分喜悅他的禱告,連他沒有求的,神都給他.一個成功的人,不是為求自己的利益,乃是先求神的國和義.耶穌也是這樣向我們發出這樣的挑戰,叫我們先求神的國和義,他把其他一切都給我們.這才是神眼中所認為的成功者.
\newline
\newline
有很多人,在眼中看是成功,但在神看來則失敗.無知的財主就是一個好例子,人看他很富有,很成功；但因為他只知為自己,沒有理會神,所以在神看來是失敗的.這不是說,我們不須努力上進,不須為學業,事業盡力；這乃是說我們要以神的國和義,作為我們人生最高和唯一的目標.以前我們以學位,以錢財為目標,但現在我們所努力的,所爭取的,是為了神的國度,這兩者是不同的.
\newline
\newline
先求神的國和義,也包括了建立我們的品格.神的義是神的聖潔,我們要追求的聖潔.無知的財主建造的只是倉庫,但他沒有建立品格,結果是失敗的.亞伯拉罕和羅得的不同,就是在此.羅得只顧自己,沒有建立品格,他為自己選一塊肥美的地,雖然產業是得到了,但是神並不賜給他.
\newline
\newline
四．他是忠心善用恩賜的人
\newline
\newline
14-17節講到所羅門擁有許多神所賜的恩典,但他沒有把所得的收起來；他乃是好好的運用,是忠心善用恩賜的人.雖然各人所賜不同,但是只要我們忠心的善用,神對我們的稱讚是一樣的.神看重的是忠心,而不是人的表現.將來有些寂寂無聞的人,會得到主大大的獎賞；是因為他們忠心的善用神給他的恩賜.站穩崗位,不貪慕虛榮,高位,乃是忠心的事奉,這是最重要的.
\newline
\newline
忠心是甚麼呢?忠心這個詞,原來是指羅馬時代負責在艦隊中搖船的奴僕.他們站在船艙,分成兩排的,聽著領隊的口號而操作；當口號快,他們便要快,口號慢,他們就要慢.在神面前忠心的人,就是我們仰望耶穌基督,他說快就快,慢就慢,停就停,我們是照主所指揮的去做,這就是忠心了.
\newline
\newline
五．他是勝過逆境考驗的人
\newline
\newline
所羅門是從逆境中出來的,特別是在大衛年老,未選立繼承人的那段期間,所羅門的處境是危險的,他可能會被爭奪權位的人所殺.傳道書讓我們了解到所羅門所經歷的一些患難和痛苦,人生是虛空的,無常的.但是他經得起逆境的考驗.
\newline
\newline
基督徒的人生,絕不是一帆風順的.會有痛苦,挫折,患難；但成功的基督徒,可以勝過逆境,而不是被逆境所勝.人生必有幽谷,但是成功的基督徒是在逆境中靠主的恩典經過的.很多時候,神把痛苦患難加給我們,是對我們有益的,是為叫我們更深的認識和經歷神.約伯是一個很好的例子,他在經過磨煉之後,說:「從前我風聞有你,現在我親眼看見你.」這是他對神的認識,因著痛苦而加深了.
\newline
\newline
有時神會救我們脫離患難,就像尼尼微城的人的經歷,以色列人經過紅海的經歷.但有時神並不這樣做,但以理就要經歷苦難,他的三個朋友也經歷苦難,神要他們在經歷苦難的時候,才把他們救出來.有時神或者像對保羅那樣,給我們一根刺,藉此而讓我們經歷神更大的能力和福氣,但我們要小心,不要用自己的方法,不合神旨意的方法,去脫離患難,這只會令我們更苦.我們要等候神,要靠神的力量去度過逆境.
\newline
\newline
在神眼中,一個成功的人,是謙卑的人,是敬虔的人,是先求神的國和義的人；是忠心善用恩賜的人,能勝過逆境的人,這和世人對成功的看法,是多麼的不同；我們所求的,應該是主裡的成功.
\newline
\newline


\newpage

\section{第58屆~港九培靈會~張慕皚~第六講}
\label{sec:460}
\textbf{基督教「家庭觀」}
\newline
\newline
連結: \href{https://www.hkbibleconference.org/session-message/view/460}{\texttt{https://www.hkbibleconference.org/session-message/view/460}}
\newline
\newline
\hyperref[sec:459]{< < < PREV SERMON < < <}
~
\hyperref[sec:toc]{[返目錄]}
~
\hyperref[sec:461]{> > > NEXT SERMON > > >}
\newline
\newline
以弗所書 5:22-5:22
\newline
\begin{longtable}{cl}
\hline
\hline
章節 & 經文 (和合本修訂版)\\
\hline
5:22 & \begin{tabularx}{0.7\textwidth}{X} 作妻子的，你們要順服自己的丈夫，如同順服主。 \end{tabularx} \\ \\
[1ex]
\hline
\hline
\end{longtable}
(弗五22,六4):「你們作妻子的,當順服自己的丈夫,如同順服主.因為丈夫是妻子的頭,如同基督是教會的頭,他又是教會全體的救主.教會怎樣順服基督,妻子也要怎樣凡事順服丈夫.你們作丈夫的,要愛你們的妻子,正如基督愛教會,為教會捨己.要用水藉著道,把教會洗淨,成為聖潔.可以獻給自己,作個榮耀的教會；毫無玷污皺紋等類的病,乃是聖潔沒有瑕疵的.丈夫也當照樣愛妻子,如同愛自己的身子；愛妻子,便是愛自己了.從來沒有人恨惡自己的身子,總是保養顧惜；正像基督待教會一樣,同我們是他身上的肢體.為這個緣故,人要離開父母,與妻子連合,二人成為一體.這是極大的奧秘.但我是指著基督和教會說的.然而你們各人都當愛妻子,如同愛自己一樣,妻子也當敬重他的丈夫.你們作兒女的,要在主裡聽從父母,這是理所當然的.要孝敬父母,使你得福,在世長壽；這是第一條帶應許的誡命.你們作父親的,不要惹兒女的氣,只要照著主的教訓和警戒,養育他們.」
\newline
\newline
二十世紀有一個很大的改變,就是家庭制度的改變.有專家提出有五方面顯著的改變.一是離婚率的增加.二是丈夫和父親的影響力下降.三是婚外情的普遍和增加.四是職業母親增加.五是家庭成員人數下降.另外又有專家預測傳統從聖經出來的家庭制度,已經到了將沒落的時代.對於這樣的報導,我們作為基督徒的,有甚麼反應呢?我們不能同意一些消極的看法.婚姻其實是神為祂的榮耀和人的幸福而設立的,而家庭制度是社會最基本的單位.家庭制度如果好和興旺；則社會,教會,以至國家都會好和興旺；反過來說也是一樣.所以基督徒應該很留意自己的家庭生活,和在世代的見證.除了個人見證之外,還需要基督徒家庭在這黑暗世代中見證.家庭生活的質素,可以決定整個世代文明的質素.
\newline
\newline
今天我們要從聖經中看四個基督化家庭的特點:
\newline
\newline
一．主裡面的群體:
\newline
\newline
剛才所讀的經文中,有十一次提到耶穌基督；可見保羅十分著重家庭與耶穌基督的關係.一個基督化家庭,必須是一個在主裡面的群體；意思就是以耶穌在聖經中所給我們的原則,作為我們組織家庭的依歸.
\newline
\newline
首先,基督徒必須選擇基督徒為配偶,這個原則十分重要,信與不信的,不能同負一軛.(林前七39)對丈夫已離世的寡婦的勸勉是,如果你們再嫁人的話；就必須要嫁給在主裡面的人,以致成為一個在主裡面的群體.
\newline
\newline
倘若有些家庭,只是父親信主,兒女未信；或者只是兒女信主,父母未信.對於這樣的家庭,保羅在(林前七章)也有這樣的教導；就是家庭中若有不信主的妻子,或兒女,可以因為丈夫信主而成為聖潔.意思就是如果家庭中有信主的成員,這些信主的成員便應該有一種影響力；用主的原則去影響這個家庭,以致這個家庭也能以主的原則去過生活.這是信主的成員,在他們家中應有的目標.
\newline
\newline
今天,雖然教會裡有一個頗嚴重的問題,就是男少女多；但是我們不能就此便放鬆原則.我們只能多禱告,求主將更多得救的男子加入教會中.不過,若是教會中一旦有弟兄姊妹,與未信的人結了婚；在教會和弟兄姊妹的勸阻下,仍一意孤行的話；雖然這樣的做法是錯,違反了聖經的教導；但是教會仍然要用愛心去接納他們,幫助他們；而不應該使本來是信主的失落.這是對教會說的話,而不是向未結婚的弟兄姊妹說的.事實上,信和不信的結婚,後果是不堪設想的.
\newline
\newline
另一方面,一個主裡面的家庭,是接受耶穌的主權,以主居首位的.當我們結婚的時候,應該把家庭奉獻給主,認主在家庭中居首位.當然丈夫是一家之主,在丈夫之上,耶穌才是家庭中真正的一家之主.
\newline
\newline
還有一樣要注意的,就是主裡面的家庭,除了藉婚姻生兒育女之外；更應該在主裡有團契,有交通.因為家庭中應該有很好的屬靈交往.這是在主裡面的家庭的特點.我們不要讓電視影響和取代我們和家庭成員的交通.在主裡面的家庭,必須有家庭的祭壇；一同讀經,禱告,有事奉的目標.要讓外人到這家庭的時候,能感覺到有基督的馨香之氣.當然這不是在於牆上有沒有掛基督是我家之主的牌；主要的是在於我們的生活,有沒有聖經,詩歌,屬靈的雜誌在家中,這一點十分重要.
\newline
\newline
消極方面,一個主裡面的家庭,應該尋求脫離世俗的影響,不要學習世俗人的生活方式；反而要用基督化家庭的生活方式去影響他們.世俗化的家庭,只講生活的享受,我們不能這樣.我們要有正確的價值觀,要以兒女的教育為我們所關心的目標.須知道在有墮落的青少年之前,早已有墮落的父母.當我們忽略照管他們的時候,他們便墮落了.我們要著重兒女道德和靈性的教育,而不是只迫他們向上爬就算.
\newline
\newline
二．有次序的群體
\newline
\newline
妻子要順服丈夫,丈夫是妻子的頭,兒女要聽從父母,聖經十分強調群體中的次序,家庭是這樣,教會也是這樣.在教會中,基督是各人的頭,男人是女人的頭,神是基督的頭.當我們相信耶穌,加入身體之後,無論是男是女,妻子,丈夫,我們的屬靈地位,以及和主的關係,都是平等的.在本質上,從神的創造中,男女是平等的,沒分彼此的.可是在工作的關係上,行政的體系中,我們便必須有上下之分,這不是文化問題,而是三位一體神的本質和給我們的榜樣.
\newline
\newline
三位一體的神,無論是聖父,聖子,或聖靈,都有一樣的豐盛,榮美,是完全一樣的.但當耶穌,就是聖子,到世上來執行神的使命的時候,父是比祂大.但在本質上,子與父是原為一的,而且父作事直到如今,子也作事.為此,猶太人覺得十分氣憤,要拿石頭打死祂.由此可見本質和工作上的分別.我們所講的家庭中的次序,也只是在工作關係上的次序.在家庭的關係上,妻子應該順服丈夫；這個「順服」的詞,原來是指軍隊的跟隨,委身,效忠,一致的行動.教會怎樣順服基督,妻子也要怎樣順服丈夫；是出於自願的,甘心的,應該的,而且是凡事的.當然凡事並不包括一些違反真理的事.不過,妻子的順服,是一種愛的回應,因為丈夫愛妻子,正如基督愛教會.當妻子感受到丈夫的愛,關心和照顧的時候,她的順服就是一種愛的順服.
\newline
\newline
主的兒女,要在主裡聽從父母和孝敬父母.這裡的聽從和妻子順服丈夫是有分別的.聽從包含有勉強的成份,聖經沒有叫妻子聽從丈夫,只是順服丈夫.但作兒女的,是指在父母照顧下未成年的兒女,需要凡事聽從父母.但這並不是說,成年之後的兒女不要聽從父母；他仍需要一生孝敬和聽從父母.但是父母在兒女成長之後,便應該尊重兒女,讓兒女也有獨立的思想,他只是提醒和勸導他們.只是保羅在這裡的要求,是特別給那些在父母教養之下,仍未成年的兒女.除了聽從之外,還要孝敬父母；看重父母,看他們是有份量的人.兒女應該發自內心,對父母有敬重,尊重.聖經裡看重家庭的次序,要把次序和紀律放在家庭的群體中,讓權柄,紀律和次序都能清楚的表明出來.所以兒女孝敬,聽從父母,是十分重要的.舊約說咒罵父母的.要把他治死；要當眾行刑,這是十分嚴重的罪行,到今日神的準則仍是一樣.
\newline
\newline
三．教育的群體
\newline
\newline
教育兒女是父母的責任,而且是最重要的,保羅說父母教養子女的時候,不要惹兒女的氣.這在當日的社會中,是一個很重要的教導.因為在保羅當日的社會中,父親有絕對的權柄.如果兒女一生下來,抱去給父親,而作父親的轉身就走的話；這就是表示他不接納這個孩子,人便會把他丟在街上,作垃圾處理.而孩子的一生,都在他父親的權下,即使兒子作了官,但父親若要他死,也都可以；所以,在這樣的文化背景下父母很容易惹兒女的氣,惡待他們的子女.但聖經教導我們應該有好的家庭關係,父親要尊重自己的兒女,不要惹他們的氣.這是孩子也應該有的權利:他們應該有生存的權利,,有得到教養的權利,不被虐待的權利,公平待遇的權利,和被無條件接納的權利.
\newline
\newline
父母教育子女,不應只像學校那樣,只給他們謀生的技能；更重要的,父母應該給他們有品格的操練,屬靈和道德的操練,使他們懂得怎樣做人,這是更重要的.道德的教育和靈性的訓練,讓他們得著主的生命,是父母的責任.
\newline
\newline
四．愛的群體
\newline
\newline
經文雖然講到丈夫愛妻子,但前面我們已經講過,妻子順服丈夫也是出於愛,不然就沒有價值.而兒女聽從和孝敬父母也必須出於愛.家庭是愛的群體,夫婦之間,父母和子女之間,都有愛的存在.而且正維繫彼此的感情的,是從神而來的愛；是無條件的,是以意志為出發的,而不純粹是感情的愛.我們要立志將好處帶給對方.夫妻間要有愛,而父母給子女的,應該是愛的培育,甚至動物間也是如此.
\newline
\newline
耶穌在約翰福音十四章講到愛就是同在.我們不要單顧供給兒女物質的需要,而忽略他們所需要的；其實是我們和他們在一起.另外,愛是需要表達的,否則對方就不知道.但我們所給孩子的愛；不是溺愛,乃是嚴愛,是包含有管教的,是公義的.這是基督化家庭的特點.
\newline
\newline


\newpage

\section{第58屆~港九培靈會~張慕皚~第七講}
\label{sec:461}
\textbf{基督教「政府觀」}
\newline
\newline
連結: \href{https://www.hkbibleconference.org/session-message/view/461}{\texttt{https://www.hkbibleconference.org/session-message/view/461}}
\newline
\newline
\hyperref[sec:460]{< < < PREV SERMON < < <}
~
\hyperref[sec:toc]{[返目錄]}
~
\hyperref[sec:462]{> > > NEXT SERMON > > >}
\newline
\newline
羅馬書 13:1-13:7
\newline
\begin{longtable}{cl}
\hline
\hline
章節 & 經文 (和合本修訂版)\\
\hline
13:1 & \begin{tabularx}{0.7\textwidth}{X} 在上有權柄的，人人要順服，因為沒有權柄不是來自神的。掌權的都是神所立的。 \end{tabularx} \\ \\ \relax
13:2 & \begin{tabularx}{0.7\textwidth}{X} 所以，抗拒掌權的就是抗拒神所立的；抗拒的人必自招審判。 \end{tabularx} \\ \\ \relax
13:3 & \begin{tabularx}{0.7\textwidth}{X} 作官的原不是要使行善的懼怕，而是要使作惡的懼怕。你願意不懼怕掌權的嗎？只要行善，你就可得他的稱讚； \end{tabularx} \\ \\ \relax
13:4 & \begin{tabularx}{0.7\textwidth}{X} 因為他是神的用人，是與你有益的。你若作惡，就該懼怕，因為他不是徒然佩劍；他是神的用人，為神的憤怒，報應作惡的。 \end{tabularx} \\ \\ \relax
13:5 & \begin{tabularx}{0.7\textwidth}{X} 所以，你們必須順服，不但是因神的憤怒，也是因著良心。 \end{tabularx} \\ \\ \relax
13:6 & \begin{tabularx}{0.7\textwidth}{X} 你們納糧也為這個緣故，因他們是神的僕役，專管這事。 \end{tabularx} \\ \\ \relax
13:7 & \begin{tabularx}{0.7\textwidth}{X} 凡人所當得的，就給他。當得糧的，給他納糧；當得稅的，給他上稅；當懼怕的，懼怕他；當恭敬的，恭敬他。 \end{tabularx} \\ \\
[1ex]
\hline
\hline
\end{longtable}
請大家翻開羅馬書第十三章.我們將會從這章經文去瞭解基督徒對政府的看法.我們先來看第1-7節:「在上有權柄的,人人當順服他；因為沒有權柄不是出於神的.凡掌權的都是神所命的,所以抗拒掌權的；就是抗拒神的命.抗拒的必自取刑罰,作官的原不是叫行善的懼怕,乃是叫作惡的懼怕.你願意不懼怕掌權的麼,你只要行善,就可得他的稱讚；因為他是神的用人,是與你有益的.你若作惡,卻當懼怕,因為他不是空空的佩劍.他是神的用人,是伸冤的,刑罰那作惡的.所以你們必須順服,不但是因為刑罰,也是因為良心.你們納糧,也為這個緣故.因他們是神的差役.常常特管這事.凡人所當得的,就給他.當得糧的,給他納糧；當得稅的,給他上稅.當懼怕的懼怕他；當恭敬的,恭敬他.」
\newline
\newline
羅馬書主要分開兩部份.第一部份講述信仰,特別是神的救恩,即第一章至第十一章.由第十二章至第十六章是這卷書的第二部份講述基督徒的生活.因此,羅馬書包括了基督徒的信仰和生活的討論.第十二章開始說到,我們先要將我們的生命奉獻給神,將身體獻上,當作活祭.然後,再講到我們在教會裡應有的生活,人與人之間的關係.跟著,第十三章清楚說明基督徒與政府應有的關係.有人誤會基督徒不應談及政府和政治的問題,其實聖經裡提及不少關於政府,政治的問題.保羅在羅馬書提到的不是地方政府的問題,而是一般性,神的原則；關於基督徒公民責任,我們如何與一般性的政府建立關係.因此,我們不應把地方政治,或國家政治帶到教會裡.但我們相信聖經裡,舊約,新約,摩西的教訓,耶穌的教訓,保羅的教訓及普通書信裡；都提及很多很多關於政府,和基督徒公民的責任的教導.我們都應該去領受,去實踐.特別是羅馬書第十三章,提醒我們如何去做一個好公民；如何與政府建立關係.以下是我們有幾方面的真理,需要接受的.
\newline
\newline
一．基督徒生活在兩個國度裡
\newline
\newline
保羅說我們是神國度裡的子民,但又同時身處在世上的國度裡,所以基督徒有雙重的身份.英國有一位神學家使徒德,他講及以弗所書時,要我們留意保羅是寫信給以弗所的聖徒；就是在基督耶穌裡有忠心的人,在這裡提及基督徒有兩個住址,一是以弗所,另一是在基督耶穌裡.保羅在(林前二1)說:「寫信給在哥林多神的教會,就是在基督耶穌裡成聖,蒙召作聖徒的.」這裡也說明基督徒有雙重的身份.其實,主耶穌也有提及這方面的真理,曾有法利賽人問耶穌應否納稅給該撒.耶穌回答:「神的物應歸神,該撒的應歸給該撒.」(參太廿二15-22)我們是地上的公民,所以應該納稅.若我們留意「原文」,我們會看出問耶穌的人是用「給」錢該撒是否應該?但耶穌則用「還給」該撒,這表示這錢本屬該撒的.保羅在羅馬書第十三章七節中提及「上稅」一詞,也就是耶穌所用的字「還給」政府.從地上的國度來看,我們應該盡上本份.從神的國度看,我們亦要盡上對神的本份.耶穌用銀錢的意思,是說明銀錢上有該撒的形像,而人是按著神的形像造的,所以人是屬於神的.人最終的效忠應該屬於我們的神,沒有任何一個政府能夠取替神的地位.但耶穌說到我們有雙重身份時,便要有雙重的責任.因此,基督徒應該逃避兩種極端的態度.第一種是避世主義,就如中古時代,早期教會的隱士,過著修道院的生活.雖然現在沒有這種生活,但當基督徒只有神國度的生活,而欠缺社會方面的生活,則可以說是一種避世主義.此外,我們亦應避免另一個極端,即世俗化運動.世俗化運動把教會世俗化,即要求教會只有社會的生活而沒有神國度的生活,這種極端亦是不對的.
\newline
\newline
當我們說及基督徒身處兩個國度時,我們必須明白兩個國度建立的基礎是不相同的.神的國度是建立在「愛」的基礎上,而世上的國度主要是建立在「權力」上.意大利有一位出名的政治家曾說:「一切政府主要的基礎是在於好的法律和好的武器,這是一個好政府的基礎.」但耶穌卻不是這樣說,雖然猶太人想擁護耶穌為王,甚至祂的門徒也有如此盼望；所以耶穌臨上十字架前,雅各,約翰也來找耶穌.希望耶穌得國時,他們兩人可以坐在耶穌左右；但耶穌立即更正他們的思想.耶穌說外邦人是有君王治理,有大臣操權管理,耶穌說明在外邦的國度裡,是講「權柄,管束」.但在門徒中,誰要為大的,必須作你們的用人,為首的,必作眾人的僕人.正如耶穌的榜樣:「人子來,不是要受人服侍,乃是要服侍人,並且作多人的贖價.」所以神的國度出現於教會裡,是建立在「愛」的原則上.此外,兩種國度的原則是不可掉用；例如,耶穌說有人打你的右臉,連左臉也轉過來由他打.有人要拿你的裡衣,連外衣也由他拿去,這是國度的原則.但應用在國家政治上就行不通,若蘇聯向美國要求取去德州,美國就說連加州也給你,那就不行了.相反來說,教會長者,長老若在教會裡講權力,權柄,這就不合乎神國度.教會的原則,是建立在「愛」的服侍上.
\newline
\newline
保羅認為僕人的服侍和權柄的結合,是在父母這形像上.一個屬靈長者在教會裡服侍.就如父母般的愛心去服侍,所以教會就把權柄交給他們.如果一個傳道人服侍教會,以權力壓人,就必定失敗.因為這是地上的原則,而非神國度的原則.這是基督徒對政府的第一個觀念.
\newline
\newline
二．一切政權來自神
\newline
\newline
「在上有權柄,人人當順服他；因為沒有權柄不是出於神的,凡掌權的都是神所命的.」(羅十三1)這裡重複表明一切政權是出於神.「在上掌權的」是指一般政府,即能履行神所指定政府應有的功能和任務.聖經沒有特別指明,哪一種政府是神所不許可的；只有神權政治的摩西時代,是由神直接統治的.從教會歷史和聖經原則中看,我們發現民主政制的理想,有很多是從聖經裡出來的.不過,民主政制也有限制,並非完全,沒有瘕疵.諸如多數服從少數的例子,就不一定是對的.聖經中,十二探子偵探迦南地後,只有兩人贊成進軍迦南地；其他的人就不敢前進了,結果是哪方對呢?是少數對.因此,我們要小心民主政制方面的危機和限制.
\newline
\newline
一切政權是出於神的,原因何在?因為神是創造宇宙萬物的神,所以地上一切東西都是屬於神.神不單創造,祂還帶領,照料祂所創造的一切.主耶穌在安息日治病,祂是對的.正如祂說:「我父作事到如今,我也作事.」神在安息日停了創造工作,但祂仍照管世界.主耶穌說祂是神,所以在安息日裡,祂也醫治.「國權是耶和華的,祂是管理萬國的.」(詩廿二28)所以,神是有權收回政權的.例如神廢掉尼布革尼撒王(但四章),在以斯帖記中,神的手掌管歷史等等,這都表明王的心在耶和華手中,好像隴溝的水,隨意流轉(箴廿一1)
\newline
\newline
三．政府是神管理國家的器皿
\newline
\newline
「因為他是神的用人.」「你們納糧,也為這個緣故；因他們是神的差役,常常特管這事.」(羅十三4,6).一切政府,官員都是神的差役僕人.舊約說及古列王,也是神所膏立的.
\newline
\newline
政府是神的僕人,為要維持法紀.正如羅馬書十三章三至四節所說:「作官的原不是叫行善的懼怕,乃是叫作惡的懼怕.你願意不懼怕掌權的嗎?你只要行善,就可得稱讚；因為他是神的用人,是與你有益的.你若作惡,卻當懼怕,因為他不是空空的佩劍,他是神的用人.」這裡說到政府有賞善罰惡的功能；但可惜我們發現現有的政府,不能符合聖經的標準.有人以為犯法的人是因為生理的原因,或是心理作用；所以不少時候,我們當殺人犯是病人,我們要醫治他.不過,聖經卻不是這樣看,聖經看人犯法就是一個罪人,在法律面前是犯人.政府有責任令社會安定,繁榮和守法紀.據最近調查,一般年青人最關心的是治安的問題,這是政府的責任.
\newline
\newline
另一方面,政府也是保護孤苦無助的人,即福利上的責任,這也是我們納稅的用途.政府需要向神,向人負責任.還有,一個正常的政權是需要將政權和道德結合的,有政權,就要有道德.若政權和道德脫節,這就產生很大的破壞力,也是天地間最可怕的事,如希特拉擁有政權,卻缺乏道德,因而引起很可怕的戰爭後果.
\newline
\newline
四．基督徒應盡公民的責任
\newline
\newline
政府盡了它的責任,身為公民的也必盡上責任.作公民的,先要學習順服.「在上有權柄的,人人當順服他」(羅十三1),我們不可抗拒掌權的,這包括對政府的規律和要求,「順服」是含有嚴格,積極的意思.此外,我們亦要留意:「人與職位」的分開.有時某一個官長的行為很不合聖經,但在其職位上,我們仍須尊重他.對於他頒佈的律例,我們仍須尊重.主耶穌在(太廿三2-3)說:「文士和法利賽人,坐在摩西的位上；凡他們可吩咐你們的,你們都要謹守,遵行.但不要效法他們的行為,因為他們能說不能行.」耶穌要求我們把「人和職位」分開.「凡人所當得的,就給他；當得糧的,給他納糧；當得稅的,給他上稅；當懼怕的,懼怕他；當恭敬的,恭敬他.」(羅十三7)
\newline
\newline
我們需要欣賞政府對我們一切的好處,令我們在生活上可以安居,特別是香港政府,它在道路建設上,醫療制度......上,都給我們不少好處.聖經教導我們,應當尊重的,我們要尊重她.在上掌權者是服侍我們,是神的僕人,我們要尊敬他們.
\newline
\newline
另外,公民的責任是納稅.耶穌也曾有榜樣,叫我們納稅；透過彼得釣魚,得了稅銀去納稅.(太十七24-27)
\newline
\newline
還有,參政也是公民的責任.參加選舉,被選,都是對政府的支持.柏拉圖認為:「智慧人若不參與政治的懲罰,就是在壞人的領導下過生活.」有人認為參政是黑暗,不好的,我們要潔身自愛.其實,這樣正是「參政」的理由.我們是世上的光,需要在世人前發光,正如昔日的但以理一樣.不過,我們必須對另一重要的真理,就是教會和個人面對政治是有分別的.我們不應用教會名義參政,因為教會非政治團體.耶穌用「該撒的物,應給該撒；神的物,歸於神」的意思,多少是包含政教要分離,要獨立的教導.教會乃一個有組織的團體,必須與有制度的政府分開,但基督徒個人方面,卻必須積極參與社會行動.不過,香港教會卻可以服務社會；透過辦學,醫院,青少年中心,服務社群.至於涉及政治方面,教會則應該避開.神給教會的責任是傳福音,把神的大使命充份實踐.基督徒因為生活在兩個國度裡,他一方面,以個人名義參與政府的運作.而教會必須有信心,相信神是建立萬國的神,福音在任何情況下,都必須傳開.因為主耶穌曾說:「祂要把教會建造在這磐石上,陰間的權柄,不能勝過她.」(太十六18)從教會歷史中,我們看到福音不因惡劣環境而遜色.神的旨意是；「祂安排我們在那一地方,不論環境如何,我們要在那地方忠心事奉神盡上我們作公民的責任.」神的旨意是最重要的.
\newline
\newline


\newpage

\section{第58屆~港九培靈會~張慕皚~第八講}
\label{sec:462}
\textbf{基督教「豐盛人生觀」}
\newline
\newline
連結: \href{https://www.hkbibleconference.org/session-message/view/462}{\texttt{https://www.hkbibleconference.org/session-message/view/462}}
\newline
\newline
\hyperref[sec:461]{< < < PREV SERMON < < <}
~
\hyperref[sec:toc]{[返目錄]}
~
\hyperref[sec:447]{> > > NEXT SERMON > > >}
\newline
\newline
約翰福音 10:10-10:10
\newline
\begin{longtable}{cl}
\hline
\hline
章節 & 經文 (和合本修訂版)\\
\hline
10:10 & \begin{tabularx}{0.7\textwidth}{X} 盜賊來，無非要偷竊、殺害、毀壞；我來了，是要羊得生命，並且得的更豐盛。 \end{tabularx} \\ \\
[1ex]
\hline
\hline
\end{longtable}
基督徒豐盛的人生乃主耶穌的應許.「我來了,是要叫人得生命,並且得的更豐盛.」(約十10).基督徒豐盛的人生要在主耶穌裡才可得著.
\newline
\newline
基督徒的豐盛人生觀與非基督徒的豐盛人生觀有何不同的地方?非基督徒豐盛的人生觀,乃是追求名利,他們認為名譽,地位,就是表示他們的豐盛.而基督徒的豐盛人生觀又怎樣呢?基督徒也會追求名利,地位；但卻在名利,地位之上,因為基督徒有神.若果沒有神,名利富貴都變成空虛的了.
\newline
\newline
當我們看詩篇九十篇時,我們發覺第1-11節,摩西描述了人生的短暫和空虛；然後第12-17節中,包含了六個禱告；也代表基督徒追求奔走的方向,以致我們生命得到充實,有一個豐盛的生命.
\newline
\newline
一．追求智慧
\newline
\newline
「求你指教我們怎樣數算自己的日子,好叫我們得著智慧的心.」(詩九十12)「智慧」乃是從神而來的,不是世俗的.智慧與知識是有關連的,亦有分別的；知識多是頭腦方面,而智慧是頭腦知識適當的運用.摩西教導我們要有從神而來的智慧去生活.摩西在第1-11節,先指出人的空虛和短暫,以致我們要有智慧去過豐盛的人生.許多人沒有這種智慧,以為他們能夠永遠地活下去,從來不想到人會有結束的一天.
\newline
\newline
人生的短暫如摩西所說:「如水沖去」,「如睡一覺」「如生長的草,早晨發芽生長,晚上割下枯乾.」(詩九十5-6),又「好像一聲嘆息」(詩九十9).人生就是如此短暫,無奈.還有人生是「勞苦愁煩,轉眼成空」(詩九十10).人生為何如此短暫,空虛呢?原來因為人活在罪的壓迫之下,如詩篇九十11說:「誰曉得你怒氣的權勢,誰按著你該受的敬畏曉得你的忿怒呢?」,人生活在怒氣下,人生就是空虛.人沒有神,就生出焦慮,失望,悲觀來.本世紀的實存主義所表達的人生,就如摩西這裡所說,人生就是空虛,苦悶,無奈,短暫的,人需要認識這種智慧.但摩西更進一步要求認識神的智慧.
\newline
\newline
認識神,認識耶和華是智慧的開端.可惜今日不少人看見人生的變幻空虛,仍不歸回神的懷抱；他們尋求星座,通勝,算命,交鬼來指導他們；甚至有些基督徒也學效,這是非常危險的事.我們應該把我們的生命,交在神的手中,就如大衛一樣,他願意跟從神.
\newline
\newline
另一方面,摩西認為智慧也有愛惜光陰的意思.數算日子,就是愛惜光陰.「要愛惜光陰」(弗五16),「愛惜」有買贖的意味.光陰不愛惜,就會浪費,又會過去.因此我們要好好把握時間,用在神的事工上.
\newline
\newline
二．追求靈交
\newline
\newline
「耶和華阿,我們要等到幾時呢?求你轉回,為你的僕人後悔.」(詩九十13).摩西求神不要再向我們發怒,求神除去祂的怒氣；我們願意在神面前認罪,以致我們可以與神有靈交.
\newline
\newline
罪帶來空虛,分離.這是多可怕的事啊!感謝神,藉耶穌在十字架的代罪,我們可以重歸神的懷抱,神與我們同在.舊約中,大祭司每年一次代表我們入聖所,與神有密切的靈交,這是代表性的.但新的時代,有耶穌成為我們代罪的羔羊,以致我們可以直接與神有密切的相交.舊約時代,我們要進入聖所,但今日,我們的身體是聖靈的殿,神住在我們心中,與我們同在.現在雖然有物質的阻隔,但我們仍可與神有相交,這是非常寶貴的.我們要過實現主與我們同在的生活.
\newline
\newline
我們要學習每時每刻與神親近,這是豐盛人生必須追求的方向.
\newline
\newline
三．追求慈愛
\newline
\newline
「求你使我們早早飽得你的慈愛,好叫我們一生一世歡呼喜樂.」(詩九十14)求神使我們得到祂的慈愛.罪是有分隔的能力,而愛卻有連繫的能力.愛叫我們建立「樹與枝子」的關係,神在我們裡面,我們在神裡面.愛是維持生命的能力,所以大衛說:「愛比生命更好」(詩六十三3).人生有量而無質是非常痛苦的,而愛是生命的質,所以,人生需要愛來滋潤.將來在地獄的人,他們也有生命；但他們沒有質,沒有愛,所以們是痛苦的.故此,我們需要追求神的慈愛.
\newline
\newline
如何能追求得到神的慈愛呢?耶穌教導我們:「有了我的命令又遵守的,愛我的必蒙我父愛他,我也要愛他,並且要向他顯現.」(約十四21)遵守神命令的人,必蒙神的慈愛和尊重.大衛對神的順服依靠,以致帶來他的宣認:「我一生一世必有恩惠慈愛隨著我,我且要住在耶和華的殿中,直到永遠.」(詩廿三6)神向那遵守的人施行慈愛,神的慈愛令我們充滿快樂.
\newline
\newline
四．追求喜樂
\newline
\newline
「求你照著你使我們受苦的日子,和我們遭難的年歲,叫我們喜樂.」(詩九十15).這裡的「喜樂」與第十四節的「歡呼喜樂」是不同的.第十四節的「歡呼喜樂」乃是感性的歡呼；而這裡的「喜樂」非情緒性的,即使在受苦的境況裡,仍有喜樂.
\newline
\newline
不少時候,我們喜歡在受苦的時候,尋求原因；但神卻不一定讓我們知道受苦的原因.神容許我們受苦,總有祂的意思,不是我們如此渺小卑微的人所能瞭解的,不過,我們可以學習仰望神.
\newline
\newline
「神阿,我的心切慕你,如鹿切慕溪水.」(詩四十二1).大衛在痛苦中仰望神,這提醒我們在受苦時,不要單尋求原因,乃要切慕神,等候神.「神阿,你是我的神,我要切切的尋求你.在乾旱疲乏無水之地,我渴想你,我的心切慕你.」(詩六十三1).在痛苦中,我們仰賴神,這可幫助我們忘記痛苦.
\newline
\newline
保羅的榜樣值得我們參考.他在羅馬坐監時,寫信給腓立比教會,充份表現他積極的人生態度,凡事從好處著想.這樣纔容易幫助我們脫離痛苦的環境,在情感上,也較為舒暢.「喜樂」是含有意志上的意思,即保羅在腓立比書第四章四節所說:「你們要靠主常常喜樂,我再說,你們要喜樂.」我們要凡事從好處想,便有積極樂觀的態度,也容易經過痛苦的環境.
\newline
\newline
五．追求能力
\newline
\newline
「願你的作為向你僕人顯現,願你的榮耀向他們子孫顯明.」(詩九十16).我們要追求神的能力.在這裡要留意:摩西盼望神的作為(即神的能力)向你僕人(眾數)顯現.
\newline
\newline
不少時候,我們祈求神賜我能力,使我的事奉有效；賜福我的教會更加興旺；賜福神學院的工作上更加有效的事奉你,但我們看到摩西的禱告,並不是那樣自私的,他願意神的能力向眾僕人彰顯,只要神的能力向這世代彰顯,他便能得到滿足.這也是我們需要去學習的禱告態度.
\newline
\newline
我們如何能得到神的能力?透過奉獻的生活,服事教會,神便顯出祂的能力.我們需要看見一班徹底奉獻自己,在生活見證上顯出祂的能力的人.在二十世紀有一位神的僕人叨雷,他帶來教會長期性的復興,不論他去到那裡,那裡便有復興.他得能力的秘訣就是:(一)有班人願意對付自己的罪,過聖潔的生活.(二)這班人用心禱告,看到神的能力從天而降.(三)這班人願意隨時讓主使用,成為神的僕人.只要任何地方教會遵照叨雷的秘訣,這裡便有長期性的復興.所以我們要去追求神的能力,讓神的能力彰顯.
\newline
\newline
六．追求成就
\newline
\newline
「願主我們神的榮美,歸於我們身上.願你堅立我們手所作的工.我們手所作的工,願你堅立.」(詩九十17).摩西兩次提及「願神堅立我們的工作」,這裡說人生應追求的成就.沒有人願意有一個失敗的人生,但基督徒的成功在哪裡呢?摩西提及「願神堅立我們手所作的工」,即求神幫助我們的工作有長久的價值.正如主耶穌說:「把房屋建在磐石上,而非沙土上.」這就是一個豐盛的人生所必須追求的.
\newline
\newline
豐盛的人生是滿足,有成就的.所以我們要學習問:「我們現在所做的工作是否有永恆的價值?」這個態度是必須有的.我們做事要有忠心,不要與人比較,神會堅立祂的工作.
\newline
\newline
求主幫助我們好好追求豐盛人生的六方面.(一)追求智慧,(二)追求靈交,(三)追求慈愛,(四)追求喜樂,(五)追求能力,(六)追求成就.
\newline
\newline


\newpage

\section{第58屆~港九培靈會~焦源濂~第一講　領袖的鑑戒}
\label{sec:447}
\textbf{第一講　領袖的鑑戒}
\newline
\newline
連結: \href{https://www.hkbibleconference.org/session-message/view/447}{\texttt{https://www.hkbibleconference.org/session-message/view/447}}
\newline
\newline
\hyperref[sec:462]{< < < PREV SERMON < < <}
~
\hyperref[sec:toc]{[返目錄]}
~
\hyperref[sec:448]{> > > NEXT SERMON > > >}
\newline
\newline
路得記 1:1-1:5
\newline
\begin{longtable}{cl}
\hline
\hline
章節 & 經文 (和合本修訂版)\\
\hline
1:1 & \begin{tabularx}{0.7\textwidth}{X} 士師統治的時候，國中有饑荒。在猶大的伯利恆，有一個人帶著妻子和兩個兒子往摩押地去寄居。 \end{tabularx} \\ \\ \relax
1:2 & \begin{tabularx}{0.7\textwidth}{X} 這人名叫以利米勒，他的妻子名叫拿娥米；他兩個兒子，一個名叫瑪倫，一個名叫基連，都是猶大伯利恆的以法他人。他們到了摩押地，就住在那裡。 \end{tabularx} \\ \\ \relax
1:3 & \begin{tabularx}{0.7\textwidth}{X} 後來拿娥米的丈夫以利米勒死了，剩下她和兩個兒子。 \end{tabularx} \\ \\ \relax
1:4 & \begin{tabularx}{0.7\textwidth}{X} 兩個兒子娶了摩押女子，一個名叫俄珥巴，第二個名叫路得，在那裡住了約有十年。 \end{tabularx} \\ \\ \relax
1:5 & \begin{tabularx}{0.7\textwidth}{X} 瑪倫和基連二人也死了，剩下拿娥米，沒有丈夫，也沒有兒子。 \end{tabularx} \\ \\
[1ex]
\hline
\hline
\end{longtable}
路得記是發生在士師時代,這時代的事跡,絕大部份已記在士師記內.不過在士師時代這樣混亂敗壞的局面中,有路得記這個故事,實在是這時代中一片突出的乾淨樂土.
\newline
\newline
路得記是描寫一個外邦女子如何蒙神揀選,表明在舊約時代中,神並沒有棄絕外邦人.另一方面,路得記中有一個愛情故事,而事實上,整卷路得記,都是講到愛的事.不但是波阿斯與路得之間夫妻的愛；也有婆媳之愛,主僕之愛......等；換言之,這卷書實在充滿了愛.哪裡有愛,哪裡就有神.(約翰一書)說:「住在愛中的,就是住在神裡面,神也住在他裡面.」因此,在這樣的愛裡面,帶來了世人的希望──路得記的末了,是大衛王的出現.
\newline
\newline
以色列人經過了四百多年的士師時代,期間充滿了敗壞,黑暗,絕望,然而在大衛王出現後,士師時代便一去不返.在這處,我們應該看到士師記給今天的世代有很清楚的教訓.大衛王是預表耶穌基督,任何時代唯一的希望,只有耶穌基督；基督是藉著一小群的人被帶到世人中間,路得記是根治士師時代危機的一卷書.士師記是寫敗壞的根源,路得記寫的是醫治的根源.
\newline
\newline
士師記中,敗壞的根源就是,「以色列中沒有王,各人任意而行.」士師記中的以色列人,不是一些沒有信仰的人,乃是一些不尊重信仰的人；他們並沒有把信仰實踐在生活中.他們雖然信神,卻仍拜著很多的假神,不將耶和華看為他們的神.在這種自由而行,任意妄為的時代；帶給世界敗壞和絕望,四百多年也找不到希望.神雖然聽他們的禱告而興起士師,但是,以色列人並沒有從根本上看到自己的弊病；只是在受外邦壓迫到最痛苦時才呼求神,神憐憫拯救他們,痛苦解除後,他們又遠離神.信仰鬆懈,生活再度敗壞,能力亦跟著消失!然後,再愛外邦欺壓,再呼求神；士師再次被興起.......這種因素循環不息；如是者,有七次之多的循環記在士師記裡面,沒有出路可言.
\newline
\newline
一個人若沒有從根本發現其信仰問題的所在,便找不到脫離苦難的真正出路.但在路得記中,神又給我們另一方向.
\newline
\newline
故此,可見神在士師時代工作的兩條路,首先,神用十三個士師,套用今天的話來形容,他們都是時代的英雄豪傑,例如底波拉,參孫等是特出的人物.可惜他們只是短暫地出現,不能被神長久地使用；在這種情況下,神同時揀選了不受人注意的拿俄米,路得和波阿斯；他們可被神使用,將以色列人從根本的痛苦和軟弱中救出來!兩者所不同的,在於前者的表現,是在其才能和恩賜上.其實並沒有很特殊的靈性,例如參孫的靈性並不好,生活見證更是差,以致神不能完全地使用他.後者如路得記中的人物,他們的生命和生活都能與信仰緊緊結合；尤其是當眾人背道的時候,他們仍能謹守信仰於生活中.住在愛裡面的,也就住在神裡面；神也住在他們裡面,這樣的人才是世界的希望.路得記裡的一個小家庭,和波阿斯的田園,都是一小群人所組成；這個家庭與教會,是神要給世界希望的所在.請不要忘記,家庭在任何時代,都是神工作的基地；有形的教會可被摧毀,但是有屬靈的家庭,是撒旦所不能摧毀的.教會和家庭一樣,若滿有神同在,陰間的權柄不能勝過她.如今我們要一起探討路得記中所要追求的方向.但願每一個神的僕人,不要單單在個人的恩賜上追求,更要注意的是個人與神之間關係的密切,這是信仰與生活的配合.
\newline
\newline
今天要來看一個人(1-5節)──以利米勒.大家都承認他是一個好人,又是個好丈夫.至於何者才是好丈夫,當然各人的要求不同.按照聖經的標準,好丈夫應是妻子的頭:「頭」的意思是在家庭中帶頭負責任.丈夫還有第二個資格,就是要愛妻子,以利米勒是不是這樣的一個人呢?
\newline
\newline
以利米勒這個名字誠然是美,意思是「神是我的王」.這並非一個偶然的名字,若與士師時代相提並論便看出了；當時士師時代沒有王,人們也不要神作王,各人任意而行.現在有一個人願意在這時代中,讓神作他的王,這個人的確是有信仰,有理想的人；他是那個時代的人,但又不屬於那個時代.再看他的為人,參考第一節:「在猶大伯利恆有一個人帶著妻子和兩個兒子往摩押地去寄居.」這裡的「帶」字,表明他是一個帶頭的人,一個負責任的丈夫,也是一個願意負擔子的人；他關心自己的妻兒,帶他們離開伯利恆,在苦難中為妻兒找一條出路.「愛」是他的特性,又能做頭,又有愛心,這不是合乎聖經標準的丈夫嗎?結果,以利米勒把妻兒帶到絕境,自己及兩個兒子更死在摩押地.這裡可帶出了問題,有些人不愛家人,不能帶頭,使自己的家人受苦；有些人有愛心,又做帶頭的,但結果還是使自己的家人受苦,原因何在?是否聖經的真理不能實行?神的話是永遠靠得住的?
\newline
\newline
在此要了解一件事,作為一個領袖,應怎樣負起責任?是靠自己抑或是在神裡面?這是一個值得三思的問題,愛自己的家人,卻不一定使家人幸福,這便要進一步問這個「愛」的意思.
\newline
\newline
作牧師最感頭痛的是夫妻吵架,常在半夜聽到投訴的電話,開車到他家中,勸導一番；後來夫妻相好,便覺得回教會是件麻煩的事,頗浪費時間；過了不久又來找牧師了,因為他們二人又打架了,及至和好後,牧師便成了討厭人物.我終於悟了一件事,這樣的家庭,永遠不會相愛,因為夫妻相愛時,便把神丟在一邊,遇到困難也就不能解決；故此,我們應該明白只有在神裡面才有希望.
\newline
\newline
我稍為分享一點自己很得教訓的事.小兒與一位傳道人林弟兄的女兒結婚了；事情發生於七八年,我被林弟兄從遠東邀請往紐約的夏令營講道,小兒當時在紐約讀書,本來打算到他那裡住宿.甫下飛機,林弟兄便接我們一家人到他家中,還請我們留下住宿.夏令營後,我回到遠東.幾年後,我還往美國定居,小兒告知我們他已有女朋友；細問之下,原來是林弟兄的女兒,兩人是上次我們在林弟兄家裡作客時認識的,我才知道那個夏令營還有一個後果存在.兒子可能以為我們一定很高興；可是,我們在當時沒有甚麼反應.照理由說他們的往來應該是門當戶對的一回事；同是傳道人的子女,又是中國人,還冀盼些甚麼?但是,我們總認為父親愛主,女兒不一定愛主,所以一直沒有表態.
\newline
\newline
又過了幾年,兒子終於在一次我們四人同坐車子時,向我們夫婦發出挑戰,並且點名逐個提問,要我們立即表示意見,這時不得不回答；我答他們如果清楚是神的心意,我便不反對,兒子很高興.再問內子時,內子回答:「你們交往以來的這幾年,若比以前更愛主,便贊成；沒有進步呢,我則有所保留；若是比前退步,我便反對.」我的答案雖然很屬靈,內子的卻更實際.相信這一席話給他倆有一定的影響；後來他們結婚了,的確能在愛主的事上同心合意.一對新婚夫妻,當然很相愛,只是,要注意這個「愛」是出於天然之愛,還是在神裡面的愛.以利米勒雖愛妻兒,由於這「愛」是出於天然的愛,便成為他全家人受苦的根源.與神的關係密切,是一切幸福因素之源.
\newline
\newline
我們且看看以利米勒失敗的原因.他本是很有理想的基督徒,問題在環境的轉變；假如伯利恆沒有饑荒,以利米勒是人人看好的神的兒子.
\newline
\newline
現實也一樣,真正有好品質的信徒,不但在為人方面要好,更重要的是屬靈的水準要高；這表現在應付困難危機的時候,比他人更能承受打擊,更沉著,更鎮定地應付危機.以利米勒為何在環境轉變時會失敗呢?相信這失敗是平時所種下的根.縱然他的名字是「以神為王」,但是在平日他並沒有真正「以神為王」的操練.在沒有面臨抉擇的情況下,他並不曉得自己已經深受世人的為人作風所影響,以致在逆境出現時,他裡面的軟弱便出現.故此,我們在平安時應學習(羅十二2-3)保羅說:「所以弟兄們,我以神的慈悲勸你們,將身體獻上,當作活祭；是聖潔的,是神所喜悅的.你們如此事奉,乃是理所當然的.不要效法這個世界,只要心意更新而變化；叫你們察驗何為神的善良,純全,可喜悅的旨意.」這是均衡生活的操練.平安時沒有操練,困難出現時也就不能勝過!
\newline
\newline
幾年前我曾在這裡講過約伯記,我總感到約伯是一個好的丈夫及父親；也是一個屬靈的領袖.他在危機中,能屹立不倒應付危機,全賴他平常的操練.現在不妨再讀這幾節經文.(伯一45)
\newline
\newline
「他的兒子,按著日子,各在自己家裡設擺筵宴,就打發人去,請了他們的三個姐妹來,與他們一同喫喝.筵宴的日子過了,約伯打發人去叫他們自潔.他清早起來,按著他們眾人的數目獻燔祭.因為他說,恐怕我兒子犯了罪,心中棄掉神,約伯常常這樣行.」
\newline
\newline
經文中提到,約伯的十個子女能彼此相愛,這是一個很美的見證；宴樂時經常提醒子女,有明顯的教導；另外,又暗中逐個提名為他們禱告,怕子女暗中犯罪.這些都是經常做的事.
\newline
\newline
在信仰與生活上能保持合一的信徒,常常有討神喜悅的心；突然的困難臨到時,能夠站立得住.所以,一個屬靈的領袖,不但要有突出的才能,美好的愛心,更重要的是有美好的靈性.
\newline
\newline
我們當中或有做父親的和做領袖的,你的靈性怎樣?不要以為愛家人,盡責任便可,切記要有美好的靈性,一切的美德在基督裡面.否則,就是有最大的聰明和愛心,也不一定能使人幸福.有很多信徒,在危機時更愈發緊緊地跟從主.
\newline
\newline
(詩六十三)是大衛逃難時在曠野所寫.他說:我在曠野乾旱疲乏無水之地渴想神.他說神啊!我的心緊緊地跟隨你.從前我跟隨你,在危機時我更要緊緊地跟隨你.這是一個真懂得脫離患難,美好靈性的人.我們若忘記這一點,那就必定會走向另一條路.最後請看(賽五十10-11)
\newline
\newline
「你們中間誰是敬畏耶和華,聽從他僕人之話的,這人行在暗中,沒有亮光.當倚靠耶和華的名,仗賴自己的神,凡你們點火,用火把圍繞自己的；可以行在你們的火焰裡,並你們所點的火把中.這是我手所定的,你們必躺在悲慘之中.」
\newline
\newline


\newpage

\section{第58屆~港九培靈會~焦源濂~第二講　好人的災難}
\label{sec:448}
\textbf{第二講　好人的災難}
\newline
\newline
連結: \href{https://www.hkbibleconference.org/session-message/view/448}{\texttt{https://www.hkbibleconference.org/session-message/view/448}}
\newline
\newline
\hyperref[sec:447]{< < < PREV SERMON < < <}
~
\hyperref[sec:toc]{[返目錄]}
~
\hyperref[sec:449]{> > > NEXT SERMON > > >}
\newline
\newline
路得記 1:1-1:15
\newline
\begin{longtable}{cl}
\hline
\hline
章節 & 經文 (和合本修訂版)\\
\hline
1:1 & \begin{tabularx}{0.7\textwidth}{X} 士師統治的時候，國中有饑荒。在猶大的伯利恆，有一個人帶著妻子和兩個兒子往摩押地去寄居。 \end{tabularx} \\ \\ \relax
1:2 & \begin{tabularx}{0.7\textwidth}{X} 這人名叫以利米勒，他的妻子名叫拿娥米；他兩個兒子，一個名叫瑪倫，一個名叫基連，都是猶大伯利恆的以法他人。他們到了摩押地，就住在那裡。 \end{tabularx} \\ \\ \relax
1:3 & \begin{tabularx}{0.7\textwidth}{X} 後來拿娥米的丈夫以利米勒死了，剩下她和兩個兒子。 \end{tabularx} \\ \\ \relax
1:4 & \begin{tabularx}{0.7\textwidth}{X} 兩個兒子娶了摩押女子，一個名叫俄珥巴，第二個名叫路得，在那裡住了約有十年。 \end{tabularx} \\ \\ \relax
1:5 & \begin{tabularx}{0.7\textwidth}{X} 瑪倫和基連二人也死了，剩下拿娥米，沒有丈夫，也沒有兒子。 \end{tabularx} \\ \\ \relax
1:6 & \begin{tabularx}{0.7\textwidth}{X} 拿娥米與兩個媳婦起身，要從摩押地回去，因為她在摩押地聽見耶和華眷顧自己的百姓，賜糧食給他們。 \end{tabularx} \\ \\ \relax
1:7 & \begin{tabularx}{0.7\textwidth}{X} 她和兩個媳婦就起行，離開所住的地方，上路回猶大地去。 \end{tabularx} \\ \\ \relax
1:8 & \begin{tabularx}{0.7\textwidth}{X} 拿娥米對兩個媳婦說：「你們各自回娘家去吧！願耶和華恩待你們，像你們待已故的人和我一樣。 \end{tabularx} \\ \\ \relax
1:9 & \begin{tabularx}{0.7\textwidth}{X} 願耶和華使你們各自在新的丈夫家中得歸宿！」於是拿娥米與她們親吻，她們就放聲大哭， \end{tabularx} \\ \\ \relax
1:10 & \begin{tabularx}{0.7\textwidth}{X} 對她說：「不，我們要與你一同回你的百姓那裡去。」 \end{tabularx} \\ \\ \relax
1:11 & \begin{tabularx}{0.7\textwidth}{X} 拿娥米說：「我的女兒啊，回去吧！為何要跟我去呢？我還能生兒子作你們的丈夫嗎？ \end{tabularx} \\ \\ \relax
1:12 & \begin{tabularx}{0.7\textwidth}{X} 我的女兒啊，回去吧！我年紀老了，不能再有丈夫。就算我還有希望，今夜有丈夫，而且也生了兒子， \end{tabularx} \\ \\ \relax
1:13 & \begin{tabularx}{0.7\textwidth}{X} 你們豈能等著他們長大呢？你們能守住自己不嫁人嗎？我的女兒啊，不要這樣。我比你們更苦，因為耶和華伸手擊打我。」 \end{tabularx} \\ \\ \relax
1:14 & \begin{tabularx}{0.7\textwidth}{X} 兩個媳婦又放聲大哭，俄珥巴與婆婆吻別，但是路得卻緊跟著拿娥米。 \end{tabularx} \\ \\ \relax
1:15 & \begin{tabularx}{0.7\textwidth}{X} 拿娥米說：「看哪，你嫂嫂已經回她的百姓和她的神明那裡去了，你也跟你嫂嫂回去吧！」 \end{tabularx} \\ \\
[1ex]
\hline
\hline
\end{longtable}
上次我們談到,士師時代雖然充滿敗壞,但神仍看顧祂的兒女,在那時代興起救恩,神用了兩種人,先是那些英雄豪傑,滿有才能的人；再用一些平凡的和命運不幸的人.士師在工作上事奉,路得則用生活來事奉；工作的事奉有一時的果效,惟有生命與生活的事奉才有長遠的功效.因此,士師的影響是暫時的,但路得等人物的影響是永久的.因為在他們的家中出了大衛王,當大衛出現時,士師時代便一去不返了.大衛是個預表耶穌的人,當他將生命來事奉的時候；便將耶穌帶到這個世界,惟有耶穌能使世界有希望.同時,這裡告訴我們兩種事奉的方式；士師的事奉,是一種個人主義的事奉；而路得記裡面的事奉,是少數人同心合意的事奉.那裡有一個小家庭,不是一個人,而是幾個人；在波阿斯的田莊,他們是以色列人的希望.可見每一個人的家庭,都是神所使用的所在；每一個教會都是神要賜福給這世界的根據地.這些家庭和教會有甚麼特點呢?路得記告訴我們它唯一的特點就是愛.住在愛裡面,就是住在神的裡面；神也住在他們裡面,這就是神賜福世界的辦法.因此,在家庭教會之中,那些作領袖的人便有極大的責任；他們的本質和才能,不但要有美好的生活和愛心；並且這愛心和生活,不是出於天然,而是在基督裡面.一個屬靈的領袖,不單是一個好人,並且要有屬靈的美德；有美好的靈性,能夠承擔起危機出現時的壓力；在回顧茫茫的時候,認清解救的方向.
\newline
\newline
我們今天要講第二個人──拿俄米.拿俄米是一個好人,一個好妻子.她的名字叫做甜,是個甜太太；她是一個好婆婆,連她的媳婦都稱讚她.雖然她又寡,又老,又醜,而且很窮；這樣的老人家,竟能得到兩個媳婦的疼愛.
\newline
\newline
我在美國牧養教會,裡面很多都是在美國結婚的年青夫婦.有時候,弟兄姊妹常常找我,要求我為他們禱告；原因是家姑將要到訪,禱告是希望家姑早點離去.這個問題拿俄米就不會遇到.她雖然一無所有,只會累人,但她的兩個媳婦也捨不得離棄她,這就是以證明她是一個名副其實的好人.因此,讓我們來效法拿俄米.
\newline
\newline
現在要學習怎樣做一個甜人,不管你們從前甜不甜,但是盼望你們經過這次聚會,身上多少添了些拿俄米的味道,有甜味出現.
\newline
\newline
要作甜人,首先要學習口甜.讀了今天的經文,相信大家都會被拿俄米的話所感動,她的話多麼甘甜.人是言語的動物,藉著言語,人與人能互相了解,溝通,藉著言語,把人與人之間的關係帶到一個更深的層次.愛一個人,起碼會用言語來表達愛意,因此我們中國人說「談戀愛」.我甚少看見談戀愛的人是不談話的,也很少看見相愛的夫妻是不說話的；婚後多年的夫妻仍然有許多話要說,這說明這對夫婦的關係很好.然而,說話的內容同樣重要.一8:「你們各人回娘家去罷,願耶和華恩待你們,像你們恩待已死的人和我一樣.」這話真美,不單為她們祝福,並且為她們所作的好事念念不忘.換上別的家姑,恐怕會認為兒子結婚不到幾年就死掉,這個媳婦定是剋夫相了.拿俄米記念媳婦對兒子的好處,不把莫須有的罪名歸咎媳婦身上,實在難得.要做好的基督徒,首先要學習說甜言.夫婦在新婚時說過甜蜜話後,現在有沒有再講呢,試想想上一次對配偶說甜話是甚麼時候呢?對最親近的人也不能甜；那麼竟對哪些人才會說甜話呢?希望藉著今次聚會,向最親近,或在我們周遭的人,我們的晚輩說甜蜜的話,不要常常指責他們.年青人若要被神所用,首先要學習講甜話.保羅勸提摩太凡事要作年青人的榜樣；最重要就是在言語上有好的表現.中國人在這方面特別缺乏；在外國,那些售貨員十分有禮貌,不單笑面迎人,就是你煩了她半天,沒有購物而離去；她仍要說謝謝,令顧客感到舒服.相反地去到中國人的商店,好像受刑罰一般；相當可怕,真是敬而遠之.不少人在教會裡好久也不說一句話,更談不上說甜言；弟兄姊妹在教會內,若能以言語互相勉勵,自然就能吸引未信主的朋友進來.
\newline
\newline
拿俄米第二個特點,就是心甜.光是口甜心不甜,長久以後就會令人反感!聽那些阿諛奉承,油嘴滑舌的人說話；真恨不得給他兩記耳光.那些話聽來叫人倒胃口.惟有口甜心又甜的,才能教人感到舒暢.如何心甜呢?就要注意在受苦時候的表現；苦難能把人的心試驗出來.人在受苦時,一向都是怨天尤人,嫉妒別人享福,恨上帝使我受苦；責怪周遭的人不了解,不同情我,不幫助我......一天到晚在痛苦裡面徘徊!心變得非常狹隘,思想有偏差,人變得非常自私,心便不夠甜.拿俄米受了大苦,她並未變得孤獨；沒有只為自己打算,也想到別人.我相信她要回伯利恆的時候,心裡其實盼望兩個媳婦與她一起；她們已走了一段路,只是當她想到兩個年青人,還有她們的將來,她就不願把她倆帶到痛苦裡去.拿俄米實在是個心甜的人,平時待人好,受苦的時候也待人好；這纔真是好,富足時能幫助人,貧窮時能幫助人,那就更好.因此,我們中間若有弟兄姊妹受苦,切記不要灰心,要緊記神在這時要知道你的心如何,同時也讓你知道自己；如果你有一個甜心,能感動人,亦能感動神.
\newline
\newline
拿俄米不但口甜,心甜,更重要的是人甜.心甜的人亦不一定人甜.相反,有的人心甜但卻不被人了解,甚至不為人接受.我相信每個做兒女的都承認母親的心是甜的,但卻往往受不了,擔當不起母親善意的嚕囌.有些妻子愛丈夫如同愛兒子,甚麼都要管；吃的,穿的,都要管,怎樣用錢也要管；丈夫卻往往不能忍受這樣的妻子.人甜就是人見人愛,不論到那裡,都自然能吸引別人.就如有蜜的花,定有蜜蜂在附近徘徊不散；如果花沒有蜜,蜜蜂便不會飛前去.拿俄米就像滿載蜂蜜的花一樣,兩個媳婦也趕不走,能做到這樣,就真是為人的完美；為何一個人心甜而兒女和配偶卻未能領受呢?乃是因為她為人有不少缺點,不夠完美,只有完美的東西才能吸引人.中國人有一句話:「窮在路邊無人問,富在深山有遠親.」窮困時途人看見會躲得老遠,富足之時,親戚也多起來.就如很多人自大陸出來,再回家鄉探親；親朋戚友忽然間多起來,因為你身上帶有港幣.拿俄米身無分文,兩個媳婦還要跟從她；照理說在這種時候,誰放她在眼內呢?在美國很多人鬧離婚,常為了爭取子女的撫養權而對簿公堂；卻從沒有聽聞為人子女的為了領養父母而爭個你死我活,只有你推我讓.拿俄米只是家姑而非母親,尚且得兩個媳婦的愛戴；可見她為人的完美.
\newline
\newline
一個人在言語行為上的見證非常重要.我們在主內,可看清楚神為何要使用拿俄米.別以為這樣受大苦的一個女子沒有多大的用處；中間有弟兄姊妹一無所有嗎?你年紀老邁嗎?不要以為一生沒有希望；在人看來可能是這樣,但我們的神,知道誰是真正愛祂的；祂的眼目遍察全地,要顯大能幫助那些向他的心存誠實的人.
\newline
\newline
拿俄米這樣好的人,為何神讓她受苦呢?這個極大的奧秘；也是個時常困惑我們的事,必須要看看內裡的原因.時間的關係,留待下次再講.
\newline
\newline


\newpage

\section{第58屆~港九培靈會~焦源濂~第三講}
\label{sec:449}
\textbf{神愛中的人}
\newline
\newline
連結: \href{https://www.hkbibleconference.org/session-message/view/449}{\texttt{https://www.hkbibleconference.org/session-message/view/449}}
\newline
\newline
\hyperref[sec:448]{< < < PREV SERMON < < <}
~
\hyperref[sec:toc]{[返目錄]}
~
\hyperref[sec:450]{> > > NEXT SERMON > > >}
\newline
\newline
路得記 1:5-1:18
\newline
\begin{longtable}{cl}
\hline
\hline
章節 & 經文 (和合本修訂版)\\
\hline
1:5 & \begin{tabularx}{0.7\textwidth}{X} 瑪倫和基連二人也死了，剩下拿娥米，沒有丈夫，也沒有兒子。 \end{tabularx} \\ \\ \relax
1:6 & \begin{tabularx}{0.7\textwidth}{X} 拿娥米與兩個媳婦起身，要從摩押地回去，因為她在摩押地聽見耶和華眷顧自己的百姓，賜糧食給他們。 \end{tabularx} \\ \\ \relax
1:7 & \begin{tabularx}{0.7\textwidth}{X} 她和兩個媳婦就起行，離開所住的地方，上路回猶大地去。 \end{tabularx} \\ \\ \relax
1:8 & \begin{tabularx}{0.7\textwidth}{X} 拿娥米對兩個媳婦說：「你們各自回娘家去吧！願耶和華恩待你們，像你們待已故的人和我一樣。 \end{tabularx} \\ \\ \relax
1:9 & \begin{tabularx}{0.7\textwidth}{X} 願耶和華使你們各自在新的丈夫家中得歸宿！」於是拿娥米與她們親吻，她們就放聲大哭， \end{tabularx} \\ \\ \relax
1:10 & \begin{tabularx}{0.7\textwidth}{X} 對她說：「不，我們要與你一同回你的百姓那裡去。」 \end{tabularx} \\ \\ \relax
1:11 & \begin{tabularx}{0.7\textwidth}{X} 拿娥米說：「我的女兒啊，回去吧！為何要跟我去呢？我還能生兒子作你們的丈夫嗎？ \end{tabularx} \\ \\ \relax
1:12 & \begin{tabularx}{0.7\textwidth}{X} 我的女兒啊，回去吧！我年紀老了，不能再有丈夫。就算我還有希望，今夜有丈夫，而且也生了兒子， \end{tabularx} \\ \\ \relax
1:13 & \begin{tabularx}{0.7\textwidth}{X} 你們豈能等著他們長大呢？你們能守住自己不嫁人嗎？我的女兒啊，不要這樣。我比你們更苦，因為耶和華伸手擊打我。」 \end{tabularx} \\ \\ \relax
1:14 & \begin{tabularx}{0.7\textwidth}{X} 兩個媳婦又放聲大哭，俄珥巴與婆婆吻別，但是路得卻緊跟著拿娥米。 \end{tabularx} \\ \\ \relax
1:15 & \begin{tabularx}{0.7\textwidth}{X} 拿娥米說：「看哪，你嫂嫂已經回她的百姓和她的神明那裡去了，你也跟你嫂嫂回去吧！」 \end{tabularx} \\ \\ \relax
1:16 & \begin{tabularx}{0.7\textwidth}{X} 路得說：「不要勸我離開你，轉去不跟隨你。你往哪裡去，我也往哪裡去；你在哪裡住，我也在哪裡住；你的百姓就是我的百姓；你的神就是我的神。 \end{tabularx} \\ \\ \relax
1:17 & \begin{tabularx}{0.7\textwidth}{X} 你死在哪裡，我也死在哪裡，葬在哪裡。只有死能使你我分離；不然，願耶和華重重懲罰我！」 \end{tabularx} \\ \\ \relax
1:18 & \begin{tabularx}{0.7\textwidth}{X} 拿娥米見路得決意要跟自己去，就不再對她說甚麼了。 \end{tabularx} \\ \\
[1ex]
\hline
\hline
\end{longtable}
上一講我們談到拿俄米,她的名字叫做甜；甜言,甜心,甜人.她不但是個好家姑,好妻子,更是個基督徒應效法的好模範.
\newline
\newline
今天我們繼續上次還未談完的問題.既然拿俄米是這樣的一個好人,為甚麼神不看顧她,使她受大苦.在幾年間,丈夫,兩個兒子都死去；她年紀又老邁,受到如此重大的打擊,餘下的光陰還有甚麼希望呢?當日滿滿的去,有丈夫,有兒子,有很多的金銀財寶,還有別人所沒有的機會,可以脫離饑荒；加上有家庭的溫暖,實在令人羨慕.只是,不到十年光景,便甚麼都失去,空空的回來.錢財去了,丈夫,兒子失去了,自己也老了,還有甚麼希望呢?相信拿俄米的心中對神充滿疑問:神啊你為何這樣待我,我並非是一個沒有好行為的基督徒,我的丈夫也是,他愛妻子愛兒子；就是我們走錯了一點,但你為何如此嚴重的打擊我.
\newline
\newline
對於這個問題,我想許多弟兄姊妹都有同感.為何神要苦待這樣好的基督徒?基督徒和非信徒都會這樣問,上帝有眼睛的嗎?為何讓好人受苦?拿俄米這樣的好人信上帝有甚麼好結果呢?這上帝可信的嗎?因此,有不少基督徒因受苦就跌倒.不單這樣,就是那些未受苦的基督徒,看見人家受苦,也疑問叢生,還未受苦就先跌倒了,不信主的人則更不用說.
\newline
\newline
神究竟看不看顧好人?實在是一個很深奧的問題.幾年前我講過約伯記,那裡可提供一些答案；不過相信好人受苦這個問題,雖然有了答案,但仍然是一個問題.我們頭腦上的知識很清楚,但遇上實際事情的時候便糊里糊塗.這便需要神的話不斷的提醒.
\newline
\newline
對於好人受苦這問題,世上的人有他們的說法,不信神的人會說:這有甚麼稀奇呢?自古紅顏多薄命.從前拿俄米太幸福了.現在命運倒轉變了,因此她要受苦了.我們又怎樣解釋呢?對於苦難的問題,只有神的話才能給予我們正確的答案.拿俄米的受苦,至少有兩個原因；第一個原因,就是為別人受苦,人是群居的動物,生活在一起,必定會產生互相的影響；關係愈接近的人,影響便更厲害,這是無法避免的.其次,人如果犯了罪,甚至有了錯誤；錯誤亦有它的影響力,周遭的人一定受到牽連.罪有傳染性,譬如三藩市是一個同性戀的大本營；那裡的同性戀者非常活躍,蔓延出去,甚至見怪不怪了!他們每年還有一次大遊行.罪令人痛苦,這是罪的後遺症.
\newline
\newline
聖經告訴我們犯罪必然會有苦果.當我們活在這個世界,周遭的人一舉一動都必影響我們,甚至於影響整個世界.比方說,丈夫沉迷賭博,賭錢大半是要賠的；做他的妻子便苦了,他輸了錢,你就要捱窮,連兒女也要吃苦頭.再說,夫妻二人打架,雖然只是他倆人的事,但兒子在旁觀看,亦會非常難過!很多年青人都喜愛看電視的摔角節目,打得愈精彩便愈高興.有一天在家見父母打架了,兒子絕對不會像看電視一般去欣賞了；因為打架的人是他父母,故此就影響到他的心靈.戰爭的結果,帶來家破人亡,妻離子散,剩下多少的孤兒寡婦,為何有這樣的災難,這些人有甚麼罪呀?他們是無辜的,只是受到別人錯誤判斷的影響,受到領袖野心的影響；因此世上沒有平安,是罪惡使人失去平安.神應許我們,把祂的百姓從罪裡拯救出來,不但把人從罪惡中救出來,也必須把人從罪惡的環境裡救出來.在世上我們有苦難,但感謝神!在主內有平安.我們不能絕對避免苦難,但可以減少苦難.這便要揀選在主內的環境,追求主內的生活,為我們的城市求平安!
\newline
\newline
拿俄米受苦的第二個原因,是為她自己受苦.她三番四次的說:耶和華使我受大苦,耶和華伸手攻擊我；從這話就知道她在苦難中感受到神的管教.她內心有罪疚,到底這是甚麼我們不知道；我們只能猜想以利米勒是一個好丈夫,可能在苦難未到時,拿俄米向他建議全家遷到一個安全的地方.不少丈夫是受妻子擺佈的,本來他不會這樣做,只因著愛妻子的關係,便照辦了.在我的教會,常發覺新婚的妻子總是要丈夫添一所大房子.丈夫本來只要一間小房便足夠；但妻子卻要一幢大房子,於是丈夫為子妻子的緣故,便做牛做馬.我認為拿俄米並不是這樣的人；我想她是沒有主張的人.苦難來的時候,以利米勒為了妻子的緣故,要逃難到摩押地去.雖然那地方不大好,曾經拜過摩洛神,他知道神是會不高興的.不過只是暫居數天,災難過後便可回去,相信沒有大問題；便請准妻子的意思,若果妻子反對,他就會卻步.誰知拿俄米是一個甜太太,只知馴服地追隨丈夫的做法.當拿俄米回想當初的情景,丈夫問自己的時候,自己並沒有按照神的律法回答他,卻完全把責任放在丈夫身上；因此後來受苦自己也要負一部份的責任.好人不一定是好基督徒,除非這個好是按著聖經的原則,在神的權柄之下.可是,今天很多基督徒都喪失了這一種觀念.
\newline
\newline
一次,我們教會一對年青人結婚,我跟他們證婚以後,他們就到了夏威夷去度蜜月.一星期後,他們回來了,我便問他倆有沒有吵架；他們答說這麼愉快,不可能會發生爭執的.我卻認為他們有時候吵架才是正常的,恐怕他倆變成同一鼻孔出氣,喪失原則；這樣對他倆不好,對於別人也不好.舉個例說,若丈夫有一個不好的缺點；有一天禮拜時,牧師事前並不曉得他有這個缺點；但在講道時剛巧打中了他的缺點,他可能誤會了這篇道就是牧師對他說的,埋怨牧師不在私下直言,卻在講臺上罵自己.回家後告訴太太,做妻子的也跟著生氣起來,因為她愛自己的丈夫,便站在他那邊；這是好是壞呢?按世人來說,這是好事,但按照神的真理來說這是不好的.
\newline
\newline
我們今天看見許多人受苦,往往就是因為放棄了原則.有一位弟兄全家移民到美國,本來他是一位大學教授,是德國的博士,在台灣是一位相當有名氣的人物.但是他太太提心吊膽,怕台灣不安全,千方百計往美國去.唯一的辦法是在美國投資了,於是他連大學教授也不當了,在美國買了一間餐館,做起餐館老闆來,而他的妻子也成了女招待.那位教授本是為了妻兒著想而跑到美國去,結果他的才能全不能發揮,重擔多得不得了.後來他心臟病發作就死了,留下孤寡.我並非不贊成人家移居外地,而是強調行動的動機和能力是來自何處.若果是從自己出發便要自己負責任；若是神帶領的,神便會負責；人帶領的,便由人去負責.拿俄米就是由丈夫負了她的責任,結果是何等的不幸.因此,為何好人會受苦呢?可能是自己太好了,為遷就自己所愛的人,把神放在一旁.若果是因為這樣而受苦,就不能怪神了.
\newline
\newline
拿俄米畢竟是個愛神的人,故此,神沒有忘記她,仍把她的生命保留；在她覺得自己一無所有的時候,她就從苦境裡轉回.一位又窮,又老,又寡的女人,在舊社會裡究竟有甚麼希望呢?我們可以看見神是她的將來,神要她再次起來,而且起來後還有美好的前途等著她.現在讓我們看看,神怎樣使拿俄米從苦境中轉回.
\newline
\newline
(得一6)「他就與兩個兒婦起身要從摩押地歸回,因為他在摩押地,聽見耶和華眷顧自己的百姓,賜糧食與他們.」
\newline
\newline
這是她從苦境轉回的起點.她在摩押地住了十年,甚至在她丈夫,兒子死了後,她仍住在那裡.她不想回到猶大地去,直到有一天她聽見耶和華賜福與她的百姓.我不相信猶大地的饑荒有十年之久,為何神早已賜福給祂的百姓,而她卻全不知道呢?內裡包含了一些教訓.苦難和罪惡不但令我們肉身受苦,更可怕的是我們的心靈也受到傷害；以致灰心喪膽,心胸狹隘.因此苦難能使人看不見前途,只看見目前；看見苦難,卻看不到一位能解決苦難的神.這就是苦難的可怕!由於這樣,不少受苦的人不肯禱告,不肯聚會,不想到施恩的神,一直在苦難的深淵裡.當我們想到神,想到苦難背後,想到心靈的需要,那便有福了,這是我們需要學習的.
\newline
\newline
去年,在我教會的夏令會中,有一位弟兄負責註冊的工作,他白天要上班,但他為了註冊的工作便提早半天下班；希望早些到營地把註冊手續辦好.當他下班後,他便驅車往營地去.他邊駕車,邊禱告求神讓他開車順利.但因為大路擠塞的緣故,所以他要繞小路,開快車.他便向神禱告,希望不會被警察抄牌.開了幾步,警車已到,檢控他超速駕駛,罰了錢.及後,他心裡便埋怨,為何愈禱告愈不靈驗；不禱告還可以,一禱告便出事了.平時上班開快車也沒事,但今次為了事奉神卻出事了,真是沒有道理呀!忽然,他想起自己時常開快車,但都沒有給警察遇著,卻從來不感謝神,而今次開快車被抄牌便立即埋怨神.由此可見,人受苦時思想很容易變得不平常；如果我們心靈不正常,那怎能找出努力的方向呢?
\newline
\newline
詩篇廿三,有一句話很寶貴──「祂使我的靈魂甦醒」.相信很多人讀這篇詩的時候,都只注意第二節──「祂使我躺臥在青草地上,領我在可安歇的水邊.」搬新屋子,加薪了,真是青草地,安歇的水邊；然而,人不能管吃,管睡便過一生,人要走路的.詩篇廿三所以美,主要是在第三節以後.詩人找著了人生的路向.人與動物所不同的,是人不但有肉體,還有靈魂；肉體是暫時,靈魂卻是永恆的；所以人有義路──一條永生的路.只是我們的靈魂常在沉睡,神為愛我們的緣故,使我們的靈魂甦醒.怎樣甦醒呢?「耶和華的律法能甦醒人心.」這樣,拿俄米耳朵忽然開了,聽見「耶和華賜福祂的百姓」,於是,她的靈魂甦醒,路便出現了.本來她以為只有死路一條,現在聽見耶和華的聲音,便起來回去了.盼望今天若有受苦的弟兄姊妹,你的靈魂要甦醒過來,看神怎樣把你帶到奇妙的道路上.拿俄米聽見神的話,立即不顧一切代價,起來往前走.這也是我們所需要的,聽見要有所行動,這是不容易的.若果我是拿俄米,起來後絕不回伯利恆,實在沒有顏臉再見故友；這樣寒酸地回去,我吃不消人家的眼光.聖經形容拿俄米回到伯利恆時,全城的婦女很哄動,「這是拿俄米嗎?這還是甜人嗎?......」這些話好聽嗎?她自己也不好意思,自稱為「瑪拉」,不再是昔日的拿俄米,一輩子苦到底.但是苦也要回來,為了投靠神,聽從神的話；不怕別人的譏笑,這種心志神是悅納的.難道神只有管教,沒有施恩嗎?我絕不信神會不看顧一個愛神又有好行為的人.神留下她的生命,也就是顧念她的好處；並且給她另一次機會,賜她美好的獎賞.不但將來給她有好的結局,在她一回頭時,神便讓她結了最美的果子──路得,為神得最美好的人,這個路得是神所專用的人；是拿俄米受苦時所結的果子；故此,路得可承受一切的壓力和打擊.
\newline
\newline
我們帶人信主時,常告訴人自己信主後,生意發達,妻子聽話,兒子不生病了.以致很多人為得好處而信耶穌,生意本來還好,後來卻失敗了；兒子向來不生病,現在卻病死了；那便不再信耶穌,因為我們告訴人信耶穌後會凡事亨通順利.這裡有一個人信主後似乎沒有好結果,她親口說「耶和華降禍與我」.在她的生命中,是多麼的痛苦,但她還能結果子,藉著她從神而來的行為,信心和愛心,結出路得這個果子,成為全人類的祝福.好人的苦難,是神的恩典.
\newline
\newline


\newpage

\section{第58屆~港九培靈會~焦源濂~第四講}
\label{sec:450}
\textbf{愛的跟隨}
\newline
\newline
連結: \href{https://www.hkbibleconference.org/session-message/view/450}{\texttt{https://www.hkbibleconference.org/session-message/view/450}}
\newline
\newline
\hyperref[sec:449]{< < < PREV SERMON < < <}
~
\hyperref[sec:toc]{[返目錄]}
~
\hyperref[sec:451]{> > > NEXT SERMON > > >}
\newline
\newline
路得記 1:8-1:18
\newline
\begin{longtable}{cl}
\hline
\hline
章節 & 經文 (和合本修訂版)\\
\hline
1:8 & \begin{tabularx}{0.7\textwidth}{X} 拿娥米對兩個媳婦說：「你們各自回娘家去吧！願耶和華恩待你們，像你們待已故的人和我一樣。 \end{tabularx} \\ \\ \relax
1:9 & \begin{tabularx}{0.7\textwidth}{X} 願耶和華使你們各自在新的丈夫家中得歸宿！」於是拿娥米與她們親吻，她們就放聲大哭， \end{tabularx} \\ \\ \relax
1:10 & \begin{tabularx}{0.7\textwidth}{X} 對她說：「不，我們要與你一同回你的百姓那裡去。」 \end{tabularx} \\ \\ \relax
1:11 & \begin{tabularx}{0.7\textwidth}{X} 拿娥米說：「我的女兒啊，回去吧！為何要跟我去呢？我還能生兒子作你們的丈夫嗎？ \end{tabularx} \\ \\ \relax
1:12 & \begin{tabularx}{0.7\textwidth}{X} 我的女兒啊，回去吧！我年紀老了，不能再有丈夫。就算我還有希望，今夜有丈夫，而且也生了兒子， \end{tabularx} \\ \\ \relax
1:13 & \begin{tabularx}{0.7\textwidth}{X} 你們豈能等著他們長大呢？你們能守住自己不嫁人嗎？我的女兒啊，不要這樣。我比你們更苦，因為耶和華伸手擊打我。」 \end{tabularx} \\ \\ \relax
1:14 & \begin{tabularx}{0.7\textwidth}{X} 兩個媳婦又放聲大哭，俄珥巴與婆婆吻別，但是路得卻緊跟著拿娥米。 \end{tabularx} \\ \\ \relax
1:15 & \begin{tabularx}{0.7\textwidth}{X} 拿娥米說：「看哪，你嫂嫂已經回她的百姓和她的神明那裡去了，你也跟你嫂嫂回去吧！」 \end{tabularx} \\ \\ \relax
1:16 & \begin{tabularx}{0.7\textwidth}{X} 路得說：「不要勸我離開你，轉去不跟隨你。你往哪裡去，我也往哪裡去；你在哪裡住，我也在哪裡住；你的百姓就是我的百姓；你的神就是我的神。 \end{tabularx} \\ \\ \relax
1:17 & \begin{tabularx}{0.7\textwidth}{X} 你死在哪裡，我也死在哪裡，葬在哪裡。只有死能使你我分離；不然，願耶和華重重懲罰我！」 \end{tabularx} \\ \\ \relax
1:18 & \begin{tabularx}{0.7\textwidth}{X} 拿娥米見路得決意要跟自己去，就不再對她說甚麼了。 \end{tabularx} \\ \\
[1ex]
\hline
\hline
\end{longtable}
今天要看神所重用的一個人──路得.她是一個好人,但好人不一定是神所要用的.常聽人說,信仰是勸人為善,基督徒也常有這種觀念.然而,聖經給我們有更深的要求,這個「好」要在神眼中看為「好」.(弗二10)「我們原是祂的工作,在基督耶穌裡造成的；為叫我們行善,就是神所豫備叫我們行的.」「善」是神要我們作的事.路得是士師時代所看重的人,四百年的士師時代充滿敗壞,有誰可被神所用呢?參孫,底波拉,耶弗他等,都只能被神暫時使用；路得才是神所重用的,她到底有甚麼好處呢?
\newline
\newline
路得是軟弱的女子,又是個外邦人,更是寡婦.神所要用的是她的品質,在苦難的考驗中顯出她真實的美麗來.可以說,路得不但是好人,更是得勝的好人,在神的道路上得勝.今天所讀的經文,顯然是神看中她的起點.拿俄米本著好意有三次的勸告,也就是三次的考驗,而路得全然得勝.
\newline
\newline
第一次的考驗是記載於8-9節.拿俄米叫她倆個兒婦回娘家,是釋放她們自由；自由是一個考驗,也是神給人最好的禮物.曾經很多人為自由說了很多感動人的話,例如:「不自由,毋寧死!」另外,著名的「生命誠可貴,愛情價更高；若為自由故,兩者皆可拋.」可知道人看自由比生命,愛情更為寶貴,人人都會珍惜!不自由或失去自由的人,更渴慕得著自由,就如小孩子,坐牢的,還有一種人,是身為媳婦的人.在美國,很多作媳婦的聽見家姑要來,便非常緊張.去年加州的卜其利發生一件事,有一個中國女子在校被人強暴後而被殺.我認識這個女子的父親,他是波士頓麻省理工學院的著名教授；為何這女子有這樣好的家庭,卻要從波士頓這種文化城,走到西部的加州呢?相信是受時代的風氣所影響,稍為長大便不擇手段地要得著自由,這是沒有原則的自由.
\newline
\newline
現在拿俄米主動地給她們自由.各位姊妹想想,這是福音還是禍音?相信大家都會為有這樣開明的家姑而感謝神.然而,這處給我們看見非常重要的事──自由有時候也可以成為考驗；這是一件可怕的事,誤用了自由,走到錯誤的道路也不知道.路得和俄耳巴都勝過了自由的考驗,她們選擇不自由的路.
\newline
\newline
過往我常在查經班裡事奉,前來參加查經的人有些是自由,有些卻是不自由的,原因是很多人在周末沒有節目,全宿舍空空如也；尤其是在新年時,其他同學都有節目,自己摸摸口袋,實在不自由!故此才到一些基督徒的圈子來,一起參加查經班,這是因為不自由才來的.畢業後,也就說再見,從此再沒有來了.有些人立志奉獻給主,是因為得了癌症,不自由了,只要神醫治他,便一定做傳道.也有人非常愛主,只因為失戀；若果女朋友愛他,便神也不要了.在教會事奉也一樣,不是自己自由選擇,是迫不得已被人選中逼上梁山的；下次怎樣也不想再當選.也有些人為了家庭生活而拒絕事奉；這些人都是太愛自己的自由,懂得用自己的自由來事奉主,才是可貴的.
\newline
\newline
在我們一生中,自由時多時少；年紀成長的和有錢的人則多自由.健康的,本事大,條件好之時便多得自由；自由愈多,考驗也愈多,失敗的機會也多.這樣看來,自由雖是好,但也同時帶來考驗.神給了我們自由,然而,我們如何能勝過自由的考驗,以及分辨哪種自由應放棄,哪種自由可以取呢?這便要向路得及俄耳巴來學這種秘訣.
\newline
\newline
(得一10)說:「不然,我們必與你一同回你本國去.」
\newline
\newline
自由可以考驗人生是否有固定目標.不懂得使用自由的人,往往沒有人生的目標及方向.有一回,小兒在夜半發高燒,我立即駕車出去買藥；走了沒遠,便看見很多人在打架,我一向非常愛看打架,若是在平常時候,我看到一大群人扭作一團打架,一定會停下來看.但是,現在我有一個目標,要買藥救我的兒子,故此,我就懂得使用我的自由.有些小孩整天遊玩,有些則會先完成功課才玩耍；雖然看見門口很多小孩玩得很高興,但他有自己的目標,便懂得運用自由.交異性朋友是一種自由,但是一個人訂婚後,既有立定的目標,再拈惹其他的女子,恐怕未婚妻也不能忍受；因為目標已定,雖可交異性朋友,但自由已受到一定的限制.又例如我是傳道人,平常喜歡吃大蒜,這是我個人的自由；但是今天晚上我要去探訪,那就不能吃大蒜；免得人家難以忍受我的口氣.因此,有目標時便懂得運用自由.
\newline
\newline
何謂自由呢?主說:「真理必叫你們得自由.」自由並非為所欲為,乃是有一種力量叫人做應做的事,這便是真自由.今天神要求我們愛祂,是用自由的選擇來愛祂.耶穌說:「若有人要跟從我,就當捨己,天天背起他的十字架來跟從我.」這裡的「若」字很重要,就是說你可以跟從或不跟從.但如果要跟從呢,就要甘心樂意地將自己獻上.希望各位兄弟姊妹思想一下,你做了這麼多年基督徒,有沒有基督徒應有的志向?保羅說:「我奔跑不是像沒有定向的,我打拳不是像打空氣,我只有一個目標,就是要得神從上而召我來得的獎賞.」如果你有這樣的心志,便不會埋怨身為基督徒不能做這個那個,因為你心裡非常清楚應怎樣做.
\newline
\newline
經過了第一個考驗,拿俄米又給了第二個考驗；讓兩位兒媳思想跟從的意義.她們跟從拿俄米,既浪費青春,又增加彼此的痛苦.相信拿俄米也希望她倆能跟從她；但是想到兩位年青人跟從自己只有受苦,便把這個現實的問題提出；這是一個現實的考驗.有何好處?有甚麼好結果?熱心是好,但不要熱昏了頭......就如很多年青人一樣.現實確是一個問題,因為我們是活在一個現實的社會中,有現實的需要,常勸人信主,人家會問信耶穌有何好處；事奉有甚麼利益,牧師事奉還有薪水,但做執事不但沒有薪水,還要多出錢；這個現實,是很大的困難.
\newline
\newline
我在一間美國人的教會聚會,有一個執事是開車行的,他參加教會很多活動.有一天,在一個長執會上,他說:「牧師,我明白每個信徒都應奉獻十分之一!但奉獻是不能單以金錢計算,時間也是金錢.每個人看時間有不同的價值.例如我是開車行的,禮拜天做禮拜約二小時,可能已失去賣出五部車的機會.這其實已是一種奉獻.」可見這位執事的賬項算得很清楚；故此,做禮拜時,他心裡正記掛著那些賣不出的車子,最後想想還是算了,當為奉獻罷.
\newline
\newline
我們若要看現實的話,路得和俄珥巴的現實是更現實的問題.對於古代的女子來說,一生人之中只靠丈夫和兒子,否則,要生存實在很困難.拿俄米知道自己沒有希望,但兩個年青的媳婦前面還有一大段日子,必須考慮這現實的問題.所以拿俄米的勸告,我們不能說她不對,只能在神的裡面有更美的現實,不過當時她們還不知道而已；這對路得來說,是很厲害的考驗.俄珥巴立即便清醒了,惟有路得還是那樣的不現實.如何才能對現實的問題有一個正確的觀念,並且能勝過這考驗呢?
\newline
\newline
(得一14):「兩個兒媳婦又放聲而哭,俄珥巴與婆婆親嘴而別,只是路得捨不得拿俄米.」
\newline
\newline
「捨不得」便是愛的意思,只有這種力量勝過現實,因為「愛」是超現實的,太現實的人便不懂愛情.試想一個男子在談戀愛時,每樣事情都要與對方計較,寫信一定要對方也回同等字數的信；上餐館抑或是婚後的開支,都要求對方與自己付出同等的價值.請問誰人願意跟這樣的人來往,這等如與一座機器相愛一樣,毫無意義.愛情感動人的地方,在於它是超現實的；對方根本不漂亮,別人不覺其特出之處,他會認為她與別不同.人家說「愛是美酒」便是如此,非常香令人陶醉；人在愛情中,有很多幻想,是一幅美麗的圖畫.
\newline
\newline
有一回,我到醫院探訪一位產婦,恭喜她生了個好兒子(其實她的兒子很醜),這個姊妺開心得很,還告訴我說她的兒子很愛國；看見美國人便哭,看見中國人便笑；母親對兒子的愛是超現實的.若果父母愛子女不超現實的話,便會生下孩子後,先訂下合約；說明將來兒子應如何回報母親,然後才好撫養他的兒女；假如這樣,小孩可能長不大了.父母必須超現實,好像欠了子女的債一樣,只有自己會錯,兒子不會有錯的時候,稍為做得不足,父母往往責怪自己,這便是愛情.
\newline
\newline
路得是個超現實的人,拿俄米又何嘗不是?否則,她一定叫兩個媳婦留下,好等以後生病有人照應,生活上也有人可以倚靠.現在,她因愛兩個媳婦便超乎現實.事實上,最不現實的是我們的主耶穌基督.祂若要講現實的話,大可在上十字架前,先看看哪些人可被救贖；這樣,祂便會免上十字架,因為沒有人配得!祂的代贖受苦也只有白費愛心,祂釘十字架為了救人；下面站著的人還嘲諷祂若是神子,便請跳下來.若果耶穌是現實的話,一定跳下來再跳上去給那些人看看,那十字架便沒有意義了.感謝神!祂榮耀所發出的光輝,從亙古到永遠常存的那位,祂甚至愛我們,為我們捨己.十字架感動我的地方,是愈久愈新；因為我們愈來愈感到自己不配被愛,愈感到及發現自己的敗壞；但是我的主為我掛在十字架上──超現實的愛.祂叫我們真正地愛祂,豈是多餘的嗎?
\newline
\newline
我們再看第三個考驗,就是孤單的考驗.第二個考驗有一個人跌倒,現在剩下路得一人.這是很大的考驗.我們每個人都喜歡找朋友,沒有朋友的人是非常痛苦的!若我們一輩子都有好朋友,在最需要別人扶持時連好朋友也沒有,實在是一件極難當的事情.試想有誰人像俄珥巴對路得來說是那樣親密的呢?二人同樣年齡,種族,性格,同樣的家庭及命運,就是嫁給猶太人及後又死了丈夫......等.如果說何時是路得最需要俄珥巴的話,恐怕是到了猶大地後.然而,一直在身旁的朋友,現在忽然退後,這個打擊甚大!對於一個事奉的人來說,孤單是很大的考驗,往往當痛苦達到最深時,也是最孤單的時候.約伯在失去一切時,撒旦對他還有一個厲害的打擊；就是妻子的離開,並叫他棄掉神,自己死去罷,是多麼沉重的打擊!耶穌在十字架上最沉重的打擊也是孤單的體驗,沒有人安慰祂,門徒全都離開祂,故此祂呼喊:「我的神!你為何離棄我?」一個人在這種時候仍能得勝,相信祝福也跟著來到.所有其他的試探和試煉都不能使他動搖.路得的秘訣是這裡的一句話:「我並不孤單.」當然,最好是俄珥巴能一起走.若不然則路得要做抉擇.感謝神!路得選了拿俄米,同時也是揀選上帝的福份.在世上我們常希望有更多的朋友,但有一天朋友因為不跟從主而離開了；盼望我們能像路得那樣,揀選那位受苦的主.有祂才有一切,好像路得有了婆婆,便得著從神而來上帝的賜福.故此,實際上來說路得並不是孤單的.
\newline
\newline


\newpage

\section{第58屆~港九培靈會~焦源濂~第五講}
\label{sec:451}
\textbf{神的引導}
\newline
\newline
連結: \href{https://www.hkbibleconference.org/session-message/view/451}{\texttt{https://www.hkbibleconference.org/session-message/view/451}}
\newline
\newline
\hyperref[sec:450]{< < < PREV SERMON < < <}
~
\hyperref[sec:toc]{[返目錄]}
~
\hyperref[sec:452]{> > > NEXT SERMON > > >}
\newline
\newline
路得記 1:15-1:18
\newline
\begin{longtable}{cl}
\hline
\hline
章節 & 經文 (和合本修訂版)\\
\hline
1:15 & \begin{tabularx}{0.7\textwidth}{X} 拿娥米說：「看哪，你嫂嫂已經回她的百姓和她的神明那裡去了，你也跟你嫂嫂回去吧！」 \end{tabularx} \\ \\ \relax
1:16 & \begin{tabularx}{0.7\textwidth}{X} 路得說：「不要勸我離開你，轉去不跟隨你。你往哪裡去，我也往哪裡去；你在哪裡住，我也在哪裡住；你的百姓就是我的百姓；你的神就是我的神。 \end{tabularx} \\ \\ \relax
1:17 & \begin{tabularx}{0.7\textwidth}{X} 你死在哪裡，我也死在哪裡，葬在哪裡。只有死能使你我分離；不然，願耶和華重重懲罰我！」 \end{tabularx} \\ \\ \relax
1:18 & \begin{tabularx}{0.7\textwidth}{X} 拿娥米見路得決意要跟自己去，就不再對她說甚麼了。 \end{tabularx} \\ \\
[1ex]
\hline
\hline
\end{longtable}
前一講提到路得受到三方面的考驗──自由的考驗,現實的考驗和孤單的考驗.在神來看,這是祝福而不是考驗,但是我們知道對於路得來說都是考驗.
\newline
\newline
談到自由的考驗,自由愈多,考驗也愈多,要勝過自由的考驗,則須有確定的人生目標.就如路得說:「我定意要跟從你.」有志向才懂得使用自由.基督徒應有基督徒的目標及志向；沒有目標,便會錯用自由.另外是現實的考驗.對於女子來說,婚姻及子女是最現實的問題.不錯,我們是活在現實中,有些現實是不可少的；然而,不要以為只憑對現實的眼光便能解決問題.愛神的人雖然超現實,神仍能把我們放在更美好的現實裡面.雅各和以掃便是一個例子.以掃是個很現實的人,長子名份對他來說既看不見也摸不著,實在沒有意義.雅各則是超現實的人,他渴望得著長子名份；有一天,終於用一碗紅豆湯來交換以掃的長子名份.在以掃看來,當時的紅豆湯是最現實的,以致出賣了長子名份也不覺.今天,我們已看見誰才是真正現實的人.雅各的後裔一直是神保守和賜福的對象；以掃的後裔在何處呢?俄珥巴和路得也是一樣,當時俄珥巴很重視現實,但是今天她的名字在哪裡呢?路得當時不現實,然而今天我們發現她的名字仍在現實中.
\newline
\newline
這樣看來,現實有兩種,一種是眼中的現實,一種是真正的現實.我們若憑自己的眼光,則只能活在眼前的現實中,若活在神的現實裡面,那才是真正和永不會成為過去的現實.一個不愛主的人,永遠不知道如何面對他的現實的.
\newline
\newline
至於孤單的考驗,也許我們對其中的領會不太深刻；然而對於愛主的弟兄姊妹來說,有一天將令經歷知心朋友的離去,最同心的信徒都離去．在做基督徒的過程中,有時大家都熱心追求,誰知說不定當一個時代的轉變來到,很多人都退後離開主；連最好的朋友,信徒或甚至是牧師也要跌倒.孤單的考驗,如何能勝過呢?切記在主的裡面提醒自己,孩單雖然是一件痛苦的事,同時也是一個選擇的問題.就像路得要作的選擇,結果她選擇了拿俄米,認為有婆婆便夠了.這樣路得便輕易地勝過考驗.她向婆婆作了大膽的保證,無論到哪裡都願意跟從,這也是滿有愛心的奉獻.故此,拿俄米也就不再勸她了.從此便決定了路得的一生,也決定了以色列命運的開始；甚至決定了全世界人類的將來.因為路得是大衛的祖先,也是耶穌基督肉身的先祖.
\newline
\newline
從路得對婆婆所說的一段話中,看見她不是說要為婆婆作甚麼,乃是表達出她與婆婆的關係.這個關係也是神所重視的.人所注意的是工作.然而服事神與服事人是不同的,在工作中,僱主所重視的是員工的工作能力而非個人的品格,抑或是個人與僱主的關係；下屬在家中常常打妻子也不要緊,只要工作上有表現便足夠.服事神則不同,神所重視的是事奉祂的人與祂的關係；這方面好,工作一定也做得好,這是聖經中服事的意義.(約十二26)說:「若有人服事我,就當跟從我,我在那裡,服事我的人也要在那裡.若有人服事我,我父必尊重他.」這是主耶穌的話,服事祂的人,不是說要做祂的工,乃是要跟從祂.工作雖然重要,關係更加重要；這樣的服事,不但主滿足,連父也要尊他.弟兄姊妹!你到教會事奉,是否只注意工作或是注意追求與主的關係呢?關係若不好,工作就常常做得不好?很多做主工作服事主的人,他們的人緣非常惡劣.教會的問題往往非由信徒引起的,多數是由於那些服事神作為領袖者的問題.太熱心作主的工作卻忽略了與主的關係,這樣的事奉有何意義呢?
\newline
\newline
多年前我在菲律賓,參加了一個有總統在座的聚會.當時,大家正在客廳熱鬧非常,忽然,氣氛緊張了起來,原來是總統到了.我看見有很多輛車子停下,從車上跳下了一群非常高大神勇的人,下車後便站在兩邊；然後車上跳下另一個人,這個人瘦且黑,個子又小,這便是總統；他走起路非常神氣,那些高大的侍從則必恭必敬地跟隨總統.我心想這些人很寫意啊!總統到哪裡他們也到哪裡,總統乘車或搭乘飛機出國,他們也一樣；好像甚麼事都不必做,非常快樂.其實,表面上看來他們像不必做任何事,只進進出出地跟著總統,實際上甚麼都要做.總統吩咐他們出外買東西,他們立刻便要出去,若有人要暗殺總統,這些人便把生命也拚去保護總統.服事神便是這樣,所要追求的是與神的關係,而不是只注意實際的工作.
\newline
\newline
路得對拿俄米所說的一段話中,有幾個特點,盼望我們一起學習.
\newline
\newline
路得對拿俄米說:「你往哪裡去,我也往哪裡去.你在哪裡住宿,我也在哪裡住宿.你的國就是我的國,你的神就是我的神.」緊緊的跟隨,亦步亦趨,形影不離,隨時聽命.弟兄姊妹試思想一下,是誰在指引你的路,甚麼力量使你參加該教會?找工作時是甚麼原則支持你選擇該份工作呢?這些問題非基督徒也許很容易便解決了,但你是基督徒,若希望神喜悅,這便不是簡單的問題.今天很多基督徒都是人往哪裡,我也往哪裡去；只跟隨好朋友,或是牧師,結果人家跌倒我也跟著跌倒.故此,要把跟隨的對象弄清楚；只有主才不會錯,祂是我們應該跟隨的對象.有時候,我們有另一種態度,便是要主跟隨我們,我到哪裡去,主也要到哪裡.用這種態度事奉主,休想得到神的賜福.有些人婚姻失敗,便怨神不賜福自己的婚姻；應細想當初結婚時主批准了嗎?神未批准便結婚,後來婚姻發生問題難道要神負責嗎?一個服事主的人有沒有前途,在這一點上便可看清楚.
\newline
\newline
保羅是一個事奉主的人,我認為他在歸主時就與人不同；參(徒廿二8-9):「我回答說,主啊!你是誰?他說,我就是你所逼迫的拿撒勒人耶穌.與我同行的人,看見了那光,卻沒有聽明那位對我說話的聲音.」這是保羅的經驗.主向他顯現時,他發出兩個問題,首先問主:「你是誰.」保羅在這已稱耶穌為主,但是他希望更清楚這位用大光把他打倒在地的主,祂的名字是甚麼；主回答他以後,他立即便把耶穌當為主.跟著,保羅隨即問第二個問題:「我當作甚麼,」這兩個問題很重要,很多人連第一個問題也未弄清楚,所以信仰上糊里糊塗.有些人雖然知道主是誰,但是沒有關心過,「我當作甚麼」這個問題.保羅是一個非常有主張的人,他很聰明很有才能；難道他一信主便一無所能,不知道自己應做甚麼嗎?他乃是一得救便將自己的一生全然奉獻,做主要他做的事.他不再憑自己的聰明才智來做決定,而是注重他與主的關係；緊緊地與主連結,可見他得著事奉的秘訣.他知道一件事,那就是:主開始,主負責任；人開始,人負責任.這也是成功的秘訣.有些人事奉看見神蹟,有些人事奉只有開頭而不能成功,關鍵在於開始時錯了,只有將主權放在主的帶領下,主必為自己的名引導義路.這是一個人善用生命的秘訣,唯有這樣,才有永恆的價值.故此,保羅在臨終時唱凱歌:「當跑的路我已跑盡了.」何謂跑的路,與他開始時問主「我當作甚麼」那句話有很大的關係.他所跑的,是主要他跑的路；在一生的路上有主的同在.臨終時回顧所走過的路,有無限的甜蜜和安慰,能夠歡然見主.盼望我們能學習常常跟從主.
\newline
\newline
兄弟今日在這裡講過,心裡有很多感觸,也充滿感謝!我在一九六二年從中個大陸出來,當時實在非常灰心和喪膽.由於經歷一連串的事故,我對事奉主感到非常懼怕；雖然蒙神的恩典,奇妙地把我從中國釋放出來,但是我於前途有自己的打算,決定不再當傳道人,實在傳怕了.當我在澳門會見妻子,想到我們能夠在好像不可能的情況下相逢,心裡充滿了快樂和感恩.那天晚上,我跪在旅館的房中向神獻上感恩的禱告.我說:「神啊!我感謝你,因為你將生命再一次還給我,若不是你,我怎能夠有這樣快樂的團聚.」當我在那裡感謝時,心裡有一個感動；你這個感謝有多少份量?你感謝我,就用一句話便能表達你的真情感謝嗎?我呼召你傳道,你不肯聽我的話,現在要走自己的道路.我給你有機會再渡你的一生,你的感謝有甚麼意思?」我心裡受了責備,在那個時刻向神作出第二次的奉獻.我明白除了奉獻來感謝主之外,沒有其他途逕可以向主真正表達我的感謝.這次的奉獻,我非常清楚只有一條道路,就是做主要我做的事.
\newline
\newline
後來,我回香港.岳父家那邊全家都是做出入口生意的,剛好我在大學唸的是經濟科,未信主時很希望做生意,及後奉獻做傳道便沒有這念頭.當時,我的親戚邀請我加入他們的公司,因為他們正需要我這種人才；若果沒有在澳門時的禱告,我一定會認為這是神的安排了.現在我則不敢接受這項邀請.我的岳父是個熱心的信徒,在教會裡當長老,他叫我好好的向神禱告；為前面作打算,看看神有何帶領.並給我一些零用錢,囑咐我不必介意,他全力支持我.真是非常感謝,我打開一看,那時是一九六二年,岳父一下子給了我六佰圓；由於這筆錢是岳父給我的,退回去就不對了.但是,我心裡知道,以後不能再向他人要錢了,既然做傳道,便要定睛看神如何帶領,做傳道還要親人供養,那簡直太沒有見證了.故此,我決定搬出來住,完全獨立,與妻子一同向神求一間房子.感謝神,終於給我找到了房子；由於先前有弟兄王國顯送了一套西裝給我,房東看見我西裝筆挺,便很樂意租房子給我；當然,我入伙時她心冷了一大半,因為我連睡板也是自己用木塊亂釘起來的.
\newline
\newline
一切安頓下來,跟著便要找工作,有教會的領袖見了我,只是說為我祈禱,並沒有解決我的需要；又有教會請我做幹事,我推辭了.道路愈來愈沒有希望.後來有人請我到建道神學院代課,也算是有一點的收入；可是出入長洲的交通費也不廉宜.經濟又愈來愈拮据,無論怎樣窮困；房租一定按時交,以免被趕走,有失見證!況且錢就快沒有了.有一天,早上離家時衣袋裡只剩兩塊多錢；我扣除來回的船費,將剩下的幾角錢交給內子,作為買菜的錢；並叫她做我最喜歡吃的菜(豆腐豆芽菜),那天,正當放學後我要走時,有一個學生跑上來,說有一封信,是靈實醫院的一位護士長託她交給我的.我打開來看,裡面只有壹佰圓的現鈔.事實上,我到醫院團契不過講了一次道,那個護士長我並不認識,她卻記得我,就在這種時候託人從調景嶺送到長洲來,我歡喜得不得了.當天晚上,我買了「叉燒」回家,有一點肉做菜；不再吃「但以理」餐了,豐豐富富地享受神的恩典.
\newline
\newline
然而,工作方面還是沒有著落,就在這個時候,我收到菲律賓薛玉光院長從菲律賓來的電報,邀請我到菲律賓聖經學校任教.收到電報時,雖然在環境上有一條出路,但是我不敢貿然接受；我剛從大陸來港數個月,到菲律賓不是一件簡單的事；何況當時菲律賓正在排華反共.在等候中,我想起臨離開國內時,有一位屬靈的老前輩叫我到香港後,一定要前往菲律賓,當時我不以為意；也不知他為何會這樣說,現在他的話卻不謀而合了.使我感到其中有神的帶領,於是便答應了薛院長的邀請.怎知,答應了以後,香港的路卻開了；有大教會請我擔任牧師.我只好拒絕他們,因我已決定前往菲律賓；那些兄弟姊妹認為我一定不能去,因為最近菲律賓領事館不會簽證給華人進境的.然而,幾個月後領事館沒有查問我從那裡來港,便糊塗地簽證批准我入境.當飛機到達菲律賓時,海關的官員一看見我是華人,卻認為我的護照是偽造的.但礙於國家的形象,只能扣留我的護照,再進行調查；這樣,我們便在那裡做教書的工作.當我們在那裡作工時,心裡充滿著歡欣喜樂,因為知道我們所站的地方,我所做的事情,是神要我來,也是神要我做的.
\newline
\newline
一直到今天,我回頭看自己走過的道路,若有一些不值一提的事,都是我自己任意妄為的結果.一些我感到歡喜甜蜜的事,都是因為「祂往那裡去,我也往那裡去」.故此,弟兄姊妹,在服事主的過程中,要注意與神的關係.這不但關係工作的成敗,也關係到人生的福氣和禍氣.
\newline
\newline
嘆氣,一聲嘆息,因為我們都是在虛空裡建造.
\newline
\newline


\newpage

\section{第58屆~港九培靈會~焦源濂~第六講}
\label{sec:452}
\textbf{希望的中心}
\newline
\newline
連結: \href{https://www.hkbibleconference.org/session-message/view/452}{\texttt{https://www.hkbibleconference.org/session-message/view/452}}
\newline
\newline
\hyperref[sec:451]{< < < PREV SERMON < < <}
~
\hyperref[sec:toc]{[返目錄]}
~
\hyperref[sec:453]{> > > NEXT SERMON > > >}
\newline
\newline
路得記 1:15-1:18
\newline
\begin{longtable}{cl}
\hline
\hline
章節 & 經文 (和合本修訂版)\\
\hline
1:15 & \begin{tabularx}{0.7\textwidth}{X} 拿娥米說：「看哪，你嫂嫂已經回她的百姓和她的神明那裡去了，你也跟你嫂嫂回去吧！」 \end{tabularx} \\ \\ \relax
1:16 & \begin{tabularx}{0.7\textwidth}{X} 路得說：「不要勸我離開你，轉去不跟隨你。你往哪裡去，我也往哪裡去；你在哪裡住，我也在哪裡住；你的百姓就是我的百姓；你的神就是我的神。 \end{tabularx} \\ \\ \relax
1:17 & \begin{tabularx}{0.7\textwidth}{X} 你死在哪裡，我也死在哪裡，葬在哪裡。只有死能使你我分離；不然，願耶和華重重懲罰我！」 \end{tabularx} \\ \\ \relax
1:18 & \begin{tabularx}{0.7\textwidth}{X} 拿娥米見路得決意要跟自己去，就不再對她說甚麼了。 \end{tabularx} \\ \\
[1ex]
\hline
\hline
\end{longtable}
前一講我們說過路得服事婆婆,是先注意她與婆婆的關係.在路得的話裡有兩個特質,昨天已講過第一點──「你往哪裡去,我也往哪裡去.」這是一個緊緊的跟隨.現在要談第二點──「你的國就是我的國,你的神就是我的神.」這句話的意義重大,可看出她跟隨的徹底；因為神的國是一個人行動隱藏的因素.
\newline
\newline
俄珥巴與路得原有很多相似的地方,俄珥巴愛婆婆並不下於路得；但到了現實問題的出現,我們就看到她倆分歧的地方.不錯,我們可以說路得更愛她婆婆；但為何在這時她的愛心才見顯明呢?拿俄米給了一個很好的答案.(一15),拿俄米說:「看哪你嫂子已回她的本國和她所拜的神那裡去了.」拿俄米對現實的看法是受她的信仰來決定的,她怎樣相信就對現實有相應的看法.故此,何以見得路得更愛婆婆,與及她對現實有另一個看法呢?因為她說:「你的國就是我的國,你的神就是我的神.」因此她還未到拿俄米的家,信仰和國籍也都變了；而俄珥巴卻改變了:那就是愛摩押的神.關鍵就在這裡,同樣一個考驗,兩個人有不同的反應；好像在海邊看見的帆船,相同的風向；有的船往東走,有的往西走；這就是受到藏於水底下的舵所影響.舵能利用風使船走向不同的方向,就像隱藏的信仰產生的力量,不知不覺就會反映出來.
\newline
\newline
相信大家都能從電視收看到奧運會,當蘇聯與中國對賽時,相信大家都會捧中國的場；但當北京與香港對賽,抑或是台灣跟北京對賽時,這就可能會出現不少意見.在美國有不少美籍華人,看見美國與蘇聯對賽時自然捧美國的場,然而當美國與中國對賽時情況便不同了,他們的感情很自然地受到固有的國籍和民族所影響.今天,很多人嘴裡的信仰不一定是真信仰,他們信金錢,錢在那裡心也往那處跑；雖然嘴裡說是信耶穌,卻和拜天官賜福,財神爺的人沒有兩樣.真正有根基的信仰,才是有其行事為人的方向.保羅說我們行事為人是憑著信心而不是憑眼見,我們信的是所望之事的實底,未見之事的確據.雖然我們看不見神,但我們仍然信祂；在環境並不平安的時候,我們內心有平安.洪水泛濫的時候,耶和華坐在天上為王.(詩二九10)祂必賜平安的福與嚮往祂的人.大家都是很好的基督徒,但當潮流轉變的時候,教會內有不少人跟隨潮流；那時候你跟不跟潮流?就要看你信仰的根基是否穩固.(約六66)主回過頭來便問門徒說,當眾人都離去的時候,你們也要離去嗎?彼得站出來說:「主呀,你有永生之道,我們還歸從誰呢?」(約六68)彼得能夠在潮流退的時候,仍然站得住,是因他有實際的信仰.保羅傳福音受了大苦,他不但沒有倒下,並且認為這是恩典,因他知道自己所信的是誰.他明白進入神的國,必須經歷許多的苦難；能為神的國而受苦是榮耀.就如很多人感到為國犧牲是光榮,鬧革命是榮耀.一個愛神的人為了永遠榮耀且必要降臨的國度受一點痛苦,一點的犧牲；將一點愛心獻給那萬王之王,斷不會感到羞恥.
\newline
\newline
弟兄姊妹!你為神的國度付出多少?當神國降臨時,你是這國度的精英,抑或是這世界的俘虜?你若願意在神國度佔一席位,那便要在世上度一個不屬於這世界的生活.
\newline
\newline
再看第三方面──「你在哪裡死,我也在哪裡死.除非死能使你我分離；不然願耶和華重重的降罰給我.」堅定不移的跟隨,不是一時的熱心,而是長久的態度；不單是家姑在世時善待她；就是去世後也願意留在那裡.這是一個多麼堅定的奉獻.神感動我們奉獻時,希望我們不要有所保留,能像路得那樣堅定自己奉獻的立場.
\newline
\newline
另一方面,這也是最美麗的奉獻.路得的奉獻不但表現在性質美,更在於時間上的美；在拿俄米一無所有的時候,她徹底地追隨,把自己奉獻給婆婆.相信基督再來的時候,許多基督徒都會說:「主啊,我要跟從你,你到哪裡去我也到哪裡去.」只因為那裡有榮耀,這種時候的跟隨就一點也不美麗了.故此,路得的跟隨在時間上是最好的了.若有一天跟隨主的路上,看見耶穌不受人歡迎,看見祂被許多人離棄；如果你要有美的跟從,那就是時候了.我們的主在這世界是不受歡迎的,愛主的人少；真正願意跟從主的人更少.主在世上就像拿俄米一樣,我們愛很多東西,就是不愛耶穌；但若有人愛主；用實際的行動愛祂；主早已知悉,心裡有無限的喜樂和安慰.
\newline
\newline
請續看(一22至二7)路得真正跟婆婆回到伯利恆後,一連串奇妙的事情就發生了.試留意幾處經節,(一22)「回到伯利恆,正是動手割大麥的時候.」(二3)「恰巧到了以利米勒本族的人波阿斯的田裡.」(二4)「波阿斯正從伯利恆來到.」在這裡我們看見天時,地利,人和的因素同時出現了；因此他們不久便結為夫妻.假設在今天的婚禮,一定有人會問,關於他倆戀愛的經過,相信路得只能有一個解釋；那就是神的帶領.神不喜歡我們做糊塗人,世界的人沒有時間觀念；也不懂人生意義；我們是神所救贖的,祂喜歡引導我們.(羅馬書八14)「凡被神引導的都是神的兒子.」至於我們怎樣才能明白神的引導呢?相信不少人有這樣的問題,我就簡單地把這原則談一談.
\newline
\newline
首先,要注意的是人對神的態度.就像拿俄米和路得從前看不見神的引導；直到她們回頭,神的引導才顯明出來,因此人對神的態度是神引導的最基本條件.當我們內心不願做某件事的時候,姑勿論父母怎樣大聲叫,我們也不願意；就是做也是敷衍了事.保羅說得好,要明白神的旨意,就不要效法世界,並要把身體獻上作活祭；一生討主的喜悅,滿足祂作為自己的目標.人愛甚麼就容易被其吸引,愛賭博的人很容易被賭博吸引,愛吸毒就被毒品吸引；那麼一個人愛主,主就容易吸引他.這就是路得回頭以後,主能引導他的基本條件.另外,主怎樣在這基礎上帶領她呢?(二2)「摩押女子路得對拿俄米說,容我往田間去；我蒙誰的恩,就在誰的身後捨取麥穗.」這行動好像是路得自己想出來的,實際上是被聖經的真理所引導.請參考利未記(十九9-10)「你們的地收割莊稼,不可割盡田角,也不可拾取所遺落的......要留給窮人和寄居的.我是耶和華你們的神.」另外在申命記(廿四19)也有同樣的記載.可以說路得的行動是出於聖經明顯的教導,以致出現不明顯的事,就如申命記(廿九29)所說:「隱秘的事是屬耶和華我們神的,惟有明顯的事,是永遠屬我們和我們子孫的；好叫我們遵行這律法上的一切話.」
\newline
\newline
因此神的引導有兩種,一是明顯；二是隱藏.無論我們怎樣聰明,也無法了解掌握.但神是公平的,他說明顯的事是屬於我們的,因此我們要負責任,奧秘的事則由神負責.路得來到伯利恆,明顯的事就是可以在收割時到別人的田莊拾取麥穗.她按照聖經的真理去實行,盡了被神引導的責任；至於不明顯的事,例如與波阿斯結婚,就是神的事情,神會掌管一切.故此,我們要做的就是遵行神的話.
\newline
\newline
(詩篇廿五5)「求你以你的真理引導我,教訓我,因為你是救我的神,我終日等待你.」12「誰敬畏耶和華,耶和華必指示他當選擇的道路.」14「耶和華與敬畏祂的人親密,祂必得將自己的約指示他們.」親近並敬畏神的人就能明白神奧秘的引導,這是我們一生要學習的.藉著這樣的操練,不但能把事情辦好,更需要的是讓我們與神的關係更好.大衛是神所喜悅的人,聖經裡說他是凡事尋求神心意的人,在每件事上他都先親近神.
\newline
\newline


\newpage

\section{第58屆~港九培靈會~焦源濂~第七講}
\label{sec:453}
\textbf{安息的尋求}
\newline
\newline
連結: \href{https://www.hkbibleconference.org/session-message/view/453}{\texttt{https://www.hkbibleconference.org/session-message/view/453}}
\newline
\newline
\hyperref[sec:452]{< < < PREV SERMON < < <}
~
\hyperref[sec:toc]{[返目錄]}
~
\hyperref[sec:454]{> > > NEXT SERMON > > >}
\newline
\newline
路得記 1:22-2:7
\newline
\begin{longtable}{cl}
\hline
\hline
章節 & 經文 (和合本修訂版)\\
\hline
1:22 & \begin{tabularx}{0.7\textwidth}{X} 拿娥米從摩押地回來了，她的媳婦摩押女子路得跟她在一起。她們到了伯利恆，正是開始收割大麥的時候。 \end{tabularx} \\ \\ \relax
2:1 & \begin{tabularx}{0.7\textwidth}{X} 拿娥米有一個親戚，是她丈夫以利米勒本族的人，名叫波阿斯，是個大財主。 \end{tabularx} \\ \\ \relax
2:2 & \begin{tabularx}{0.7\textwidth}{X} 摩押女子路得對拿娥米說：「讓我到田裡去拾取麥穗，我在誰的眼中蒙恩，就跟在誰的身後。」拿娥米說：「女兒啊，你去吧。」 \end{tabularx} \\ \\ \relax
2:3 & \begin{tabularx}{0.7\textwidth}{X} 路得就去了。她來到田間，在收割的人身後拾取麥穗。她恰巧來到以利米勒本族的人波阿斯那塊田裡。 \end{tabularx} \\ \\ \relax
2:4 & \begin{tabularx}{0.7\textwidth}{X} 看哪，波阿斯正從伯利恆來，對收割的人說：「願耶和華與你們同在！」他們對他說：「願耶和華賜福給你！」 \end{tabularx} \\ \\ \relax
2:5 & \begin{tabularx}{0.7\textwidth}{X} 波阿斯對監督收割的僕人說：「那是誰家的女子？」 \end{tabularx} \\ \\ \relax
2:6 & \begin{tabularx}{0.7\textwidth}{X} 監督收割的僕人回答說：「她是摩押女子，跟隨拿娥米從摩押地回來的。 \end{tabularx} \\ \\ \relax
2:7 & \begin{tabularx}{0.7\textwidth}{X} 她說：『請你容許我拾取麥穗，在收割的人身後撿禾捆中掉落的麥穗。』她就來了，從早晨直到如今，除了在屋子裡坐一會兒，她都留在這裡。」 \end{tabularx} \\ \\
[1ex]
\hline
\hline
\end{longtable}
上一講我們談到神的引導.人有神的引導,一生就充滿甜的回憶.詩篇廿三就是在神引導下喜樂的見證；大衛看耶和華為他的牧者,他知道自己在人生的問題上是軟弱的羊,因此,他一生不敢靠自己的聰明才智.這篇詩所以感動人,乃因是每一位信徒共有的詩.神是大衛的牧者,也是我們每一個人的牧者,問題在於我們有否接受祂的引導.路得有神的引導,便覺得人生如一首甜美的詩歌.我們若願意有神的引導,就應注意對神的態度；以討神的喜悅,作為人生的目的.
\newline
\newline
另一方面,路得到伯利恆後,神的話向她指出方向；就是窮人在收割時,有拾取麥穗的權利,她就按照聖經的引導而行.今天很多人常對神明的話置諸不理,在面對極大困難時,才求神的引導；結果得不著答案.因為平常時不服從神的帶領,不管神的榮耀；有重大事情臨到就更不會順從神了.
\newline
\newline
然而,並非每一件事,聖經都有清楚的指示.這樣,我們又怎知神的引導呢?在路得身上,除了有聖經的引導外；還有就是拿俄米的指導.基督徒除了聖經之外,還有神的聖靈作為引導.神的靈住在我們心裡,感動我們明白真理；這就是我們在困難中,能找到方向的重要途徑.至於我們怎能分辨?聖靈的引導,還是自己的感動呢?這就要看聖靈的性質是甚麼.首先,聖靈來了是為榮耀主,因此要留意能令主得榮耀的感動.聖靈來是要我們為罪,為義,為審判,而責備自己；因此當人犯罪或與人的關係發生困難時,內心有感動叫我們自責；叫我們看清自己的錯誤向人求和,這就是神的引導.另一方面,聖靈來了能叫人結出美好的果子,若有感動能使我們美好的品格；把主完美的性格顯露出來,這也是聖靈的引導.前面講到原則,跟著要談的是具體的操練.
\newline
\newline
第一方面我們必須要明白聖經.像路得那樣,她知道聖經裡面有記載收割時,窮人可到別人的田裡拾麥穗；我想在上一講之前,很多人還不知道聖經中有這樣的根據.跑到伯利恆就只有餓死的份兒.路得不是一個掛名的以色列人,她口裡說:你的神就是我的神,你的國就是我的國.因此她將神的話放在心中,是一個實實在在投靠神的以色列人.反之很多受洗的基督徒,信主多年也不留意神的話；相信這是有名無實的基督徒.要尋求神的帶領,便要先下功夫,將神的話研究清楚.
\newline
\newline
此外,要有活潑的信心,亦即是有行動的信心；缺少行為的信心是死的,路得雖有神的話,但要照著來行也不容易.一方面有環境的困難；在士師時代的混亂局面,不乏搶殺,姦淫的事件,客旅也受到人的欺侮.路得是一個無依無靠的摩押女子,來到輕視外邦人的猶大地；一個人孤單地謀生活,是件不簡單的事.從這裡,就可看到路得的信心來.中國人常說貧困,一個人窮,就會被困起來；但路得窮而不困,因此她能突破環境.人在逆境時常會自憐；路得並不自卑,她既投靠耶和華,便不理會別人的看法.她看自己是神的孩子,享有神孩子的權利.
\newline
\newline
另一方面,路得很積極的為人.前面說過的還不夠,神盼望我們有美好的為人.田莊還有其他人在拾麥穗,但只有路得給監管收割的人留下美好的印象；以致當波阿斯聽到報告,心裡就受感動了.神的引導不但要有明顯的原則和美好的靈性,還必須有實際美好的生活去配合,這樣神的引導就清楚了.
\newline
\newline
或許我們會問,為何按著聖經的原則行,也沒有經驗神的奇妙,我相信不是前面原則錯誤；或是個人信心不足的話,很可能就是沒有生活的配合.這一點我們常失敗了也不知道原因.有一回我讀到但以理書時很受感動,姑勿論但以理的志向如何偉大；有一點值得注意的,是他的為人非常柔和謙卑.這樣的配合,才能彰顯神的榮耀.但以理書(一9-12)那兩個「求」字是很美麗的兩個字.一般人尤其是年青人在那種情況下怎樣表達自己的意向呢?「我不去,因為這是鬼餐,我不沾污自己,要斬首也不去.......」我相信話還未說完,頭便被砍掉了；本來這話是很合理,被殺也會感到光榮,事實上我們說他因自己的愚昧而死.故此,要學習注意自己在言語,生活,為人方面的見證；這樣,神必有奇妙的事,臨到我們身上.
\newline
\newline
相信很多人都羨慕路得的結局,但卻認為自己沒有路得那樣的機緣；諸如她的好婆婆和她為人的美麗.其實這只是懶惰人的藉口.(啟三8)神給我們一個應許──「你有一點力量,也當遵行我的道,也沒有棄絕我名；我要給你一個敞開的門,是沒有人能關的.」我們看到但以理和路得的經驗,他們在任何事情上都遵守主的道,雖然面臨困境；還是以主的名為重.記得有一個很有意思的笑話,我小時候喜歡看武俠小說,很佩服別人功夫高強,甚麼鐵布衫,金鐘罩等；但我都不羨慕,我只想練到身上有電力發出來,能使人跌倒,叫牆分開；好像很容易,只可惜這力量非出自我,而是從外來的.在屬靈的事上,這是個事實；不少信徒有神的同在便有敞開的門.甚麼時候人能以神的名為重,就必看見一個敞開的門,有主帶領著經過一個一個的難關,這是我向各位作的一個誠實的見證.
\newline
\newline
神引導路得來到伯利恆這個地方,也就是路得記的中心地點；神在士師時代,藉著這地方解決以色列國的困難.因此伯利恆的田莊是那個時代希望的中心,是亂世的樂土；多麼奇妙,豐富,美麗.伯利恆是一個小城,卻是出大事的地方；日後有一位從亙古太初就有的出來,要為神執掌主權直到永遠；這就是耶穌基督.伯利恆的田莊是今天教會應當有的榜樣,教會就是這黑暗,混亂,魔鬼掌權世界裡的一片樂土,是神的兒女享受神的愛之所在.很多人流離找不到方向,但像路得一樣投靠在耶和華翅膀下,就滿得神的賞賜.聖經首次提到教會,在馬太福音十六章「我要把我的教會建造在這磐石上,陰間的權柄也不能勝過它.」這是教會產生的影響.以弗所書講到教會,最後論到爭戰,教會雖小,卻是對抗陰間權勢的地方.故此,主從耶路撒冷開始建立教會；因那時耶路撒冷是摩鬼的巢穴,是其力量具體表現的地方.然而,第一次教會聚會就有三千人悔改.再看安提阿教會,那裡只有幾個先知和教師,信徒也沒有耶路撒冷的教會那麼多；但這小地方因聽從主,差派巴拿巴及保羅出外傳福音.這本是一間小教會發起的一件小事情；卻是改變人類歷史的事,因為福音的廣傳是從那處開始.
\newline
\newline
歷史上有許多運動,興起時都是轟轟烈烈的,但影響力很快便消失了；譬如安提阿事件之前的四五百年青人大大的軍事運動；希臘亞歷山大帝東征西討,大軍所到之處勢如破竹；一直到東方的印度洋邊.今天我們回顧一下,亞歷山大的事業已不復存在,但巴拿巴及保羅在神的引領下,影響力到今天還存在.
\newline
\newline
今天我們的人雖少,卻是神希望的所在,別人看我們只是一個宗教團體；但時代的希望是在我們教會身上,讓我們不要小看自己.在神的眼中整個世界都是腐敗的,祂要拯救這世界；工具就是這個小小的教會.祂把希望放在我們身上；時間關係,下一講再繼續.
\newline
\newline


\newpage

\section{第58屆~港九培靈會~焦源濂~第八講}
\label{sec:454}
\textbf{愛的真諦}
\newline
\newline
連結: \href{https://www.hkbibleconference.org/session-message/view/454}{\texttt{https://www.hkbibleconference.org/session-message/view/454}}
\newline
\newline
\hyperref[sec:453]{< < < PREV SERMON < < <}
~
\hyperref[sec:toc]{[返目錄]}
~
\hyperref[sec:1269]{> > > NEXT SERMON > > >}
\newline
\newline
路得記 2:16-2:21
\newline
\begin{longtable}{cl}
\hline
\hline
章節 & 經文 (和合本修訂版)\\
\hline
2:16 & \begin{tabularx}{0.7\textwidth}{X} 你們還要從捆裡抽一些出來，留給她拾取，不可責備她。」 \end{tabularx} \\ \\ \relax
2:17 & \begin{tabularx}{0.7\textwidth}{X} 這樣，路得在田間拾取麥穗，直到晚上。她把所拾取的麥穗打了約有一伊法的大麥。 \end{tabularx} \\ \\ \relax
2:18 & \begin{tabularx}{0.7\textwidth}{X} 路得把所拾取的帶進城去給婆婆看，又把她吃飽所剩的拿出來，給了婆婆。 \end{tabularx} \\ \\ \relax
2:19 & \begin{tabularx}{0.7\textwidth}{X} 婆婆問她說：「你今日在哪裡拾取麥穗？在哪裡做工呢？願那照顧你的得福。」路得告訴婆婆，她在誰那裡做工，說：「我今日在一個名叫波阿斯的人那裡做工。」 \end{tabularx} \\ \\ \relax
2:20 & \begin{tabularx}{0.7\textwidth}{X} 拿娥米對媳婦說：「願那人蒙耶和華賜福，因為他不斷地恩待活人死人。」拿娥米又對她說：「那人是我們本族的人，是一個可以贖我們產業的至親。」 \end{tabularx} \\ \\ \relax
2:21 & \begin{tabularx}{0.7\textwidth}{X} 摩押女子路得說：「他還對我說：『你要緊跟著我的僕人拾取麥穗，直到他們把我所有的莊稼收割完畢。』」 \end{tabularx} \\ \\
[1ex]
\hline
\hline
\end{longtable}
上一講提到伯利恆,尤以伯利恆中波阿斯的田莊為主.它是混亂時代的一塊樂土,也是在神永恆計劃中要影響世界的所在,是人禍福的根源.另外,我們也談到教會應有的形象,可惜,今天的世代混亂,教會更加混亂,不能為世界提供出路和希望.今天我們要就波阿斯田莊的特點,來思想教會需要怎樣建立?
\newline
\newline
伯利恆之意乃「糧食之家」──在這裡有糧食的供應.伯利恆有饑荒便產生很可怕的後果.以利米勒是一個好人,他害怕饑荒而往他處找尋出路,最後只有受苦.教會本是天上糧食供應的所在.人活著並非單靠食物,乃是靠神口裡所出的一切話.今天有多人痛苦,並不因為物質的糧食不足,阿摩司先知說得好,人飢餓非因無餅,乾渴非因無水,乃因不聽耶和華的話.有些國家的人沒有聖經,那種痛苦比身體的苦更甚,因為失去了力量的源頤.今天我們有神的話語,作為我們生命力量的根源,故此,我們應當做一個渴慕神話語的人.
\newline
\newline
另一方面,伯利恆的田莊裡有非常美妙的人際關係.教會團體與世界團體不同,她具有最奇妙的人際關係,任何世界的團體,若要發揮其效用,則必須要有簡單的人際關係；諸如醫生公會,商業公會,老人會,童軍會......不合資格入會的就會引起問題.這也是個罪惡世界所產生的現象,品流不同的人就無法相處.主把人從罪惡中拯救出來,在地上設立教會,為要證明神是真的.在教會內,所有相信基督的人都可以被容納；故此教會內的份子是最複雜的,教會若只有接納簡單的會友,那就不是教會了.試看這個田莊的見證,那些監管收割的人雖是僕人；但他們與波阿斯主僕間的關係很融洽.波阿斯一到田莊,他們就彼此祝福問安.路得頭一天到田莊,他們對她的背景已瞭如指掌；波阿斯也是這樣.在教會內,新人往往在頭一次到時會受到歡迎,但過後隨即被冷落；在此我們真要小心自己的一舉一動.
\newline
\newline
前年美國羅省的華人第一浸信會,發生了一件事.一次早堂崇拜,原打算與一間小教會聯合聚會,那小教會便派了十位弟兄前來參加崇拜,怎料因交通的關係而遲到了.當他們在考慮進去與否時,招待極力邀請他們留下來,最後,終於五個人留下參加該堂崇拜,其他的則留待午堂崇拜再來.在聚會中途,忽然有一瘋子用槍射殺牧師和主席,還要射殺下面的會眾.這時那五人中有一個是休班警員,身上有佩槍,於是他開槍射殺了兇手.
\newline
\newline
我們看到這位警員留下來就是救了不少人的生命,由此可見,作招待的做得好帶來的後果是我們無法想像的.做禮拜不單是與神的關係；更重要的就是做人的工作.聖經上說:「我們原是神的工作」.換言之,(弗二10)人就是神的工作,無論我們作多少事,但若與人關係不好,也就沒有真正為神工作,婆媳之間本是很困難的關係,然而,在這田莊裡,拿俄米與路得婆媳間更有不同的種族,相處卻很和諧.另外,波阿斯是大財主,路得則一無所有；他倆能結合,生活非常美滿；可見田莊裡的人際關係實在很美.
\newline
\newline
另一個特點,這是個按照聖經經營的田莊.它能按照神的話語開放給窮人,在田莊裡,一切最高的權威是神的話語.這是值得今天教會學習的.今天的教會有一個很大的缺欠,許多人不明白教會的真理；用人的方法來治理教會,由於人是敗壞的,因此總有矛盾出現.若能以神的話語來治理教會,則有共同追求的目標,在真道上同歸於一.
\newline
\newline
另一方面,波阿斯是個大財主,可說是大能的財主,這預表了我們主耶穌基督一個很好的形像.波阿斯是財主,但也是大能者,英雄；很多人有錢無能,有災難時便跑掉,天天提心吊膽.我們的主是完美的,祂是一個財主,並不是在於金錢豐厚,乃是祂有豐盛的生命美德.所謂「人的生命不在於家道豐富,乃在於有生命的美德」.這種財富,取之不盡,用之不竭,並可使別人也富足.故此沒有人被神看為義人,惟有義者耶穌基督才能使我們稱義.祂為我們死在十字架上,為我們流出寶血；一位全然可愛的主,代替我們的不義.基督最大的能力是祂有愛心,波阿斯也是.世上沒有其他力量,比愛更偉大和更有力量.原來有一個人比波阿斯更有資格去贖路得的田產,可是想到她是摩押女子,恐怕對自己有妨礙便放棄了.波阿斯願意犧牲個人名譽,不顧後果而贖路得的田產；不但讓路得在自己的田莊拾取麥穗,並且連自己也願意給她.弟兄姊妹!你來到教會為要得甚麼?是否為求內心平安而來?抑或是牧師講得動聽,教會裡的人又與外間的不同呢?這些都是好處但只是麥穗而已.主有更大的好處要給我們,這就是祂自己.可惜這麼大的愛卻不被人了解,我們只注意那些細微而不值得注意的東西,但卻忽視了主自己的心.
\newline
\newline
此外,這田莊是由屬靈的家庭所組成,裡面充滿了仁愛,這可從波阿斯與路得的結合而見之.聖經裡載有很多婚事,令人感動的有以撒和利百加,只可惜兩人在年老時沒有好的見證.波阿斯與路得的婚事記載很短,卻感動人；可以喻基督與教會的關係.一般美滿的婚姻,往往有兩個重要因素；首先要門當戶對,否則無從溝通配搭.另外還要有時間談戀愛和彼此了解.波阿斯和路得兩人的婚姻,看來都欠缺這些因素.波阿斯是大財主,又是富有品格的猶大人,年紀也較大,路得是摩押女子,是個一無所有的年輕寡婦；表面看來二人實不相配.另外,二人相識不久就結婚,路得回到伯利恆是收割大麥之時,結婚時則剛割完小麥,前後不過七星期；這樣閃電式的婚姻卻能美滿.這裡我們且來看波阿斯的為人；他實在是所有弟兄的榜樣.
\newline
\newline
波阿斯是標準的丈夫,他不但有錢有地位,更有仁愛和謙卑,做事細微而週到；他為路得作了周詳的安排,他吩咐僕人不許欺負路得......各方面的需要他都想到了.另一方面,波阿斯有言語和行動的配合.他不但安排照顧路得,還以甜言來感動激勵路得,他對路得說:「你來投靠耶和華以色列神的翅膀,願你滿得祂的賞賜.」這是鼓勵路得在主裡要堅持下去；這話一說,路得更受感動了.她說:「我雖然不及你的一個使女,你還是用慈愛的話安慰了我的心.」很多弟兄愛妻子,但不大講話,須知有愛的言語,才能打動對方的心.愛不但在身體方面的照顧,還需要心靈上的安慰扶持.此外,在吃飯時,波阿斯特地烘餅給路得吃；又吩咐割麥穗的人故意在路得前面掉些麥穗下來,怪不得路得拾了大綑麥穗回來,也不知道是波阿斯暗中做的.這種明處和暗處的照顧,是丈夫對妻子應有的照顧.除此以外,波阿斯是一個成熟的人.他有極豐富的情感,卻能同時控制自己.他很愛路得但不表態；直到那一個晚上發現路得睡在他腳下,才曉得路得對自己也有愛意.只是他對路得說:「你只管在這裡睡到天亮,因為還有一個人比我和你之間的關係更接近,如果他不娶你的話,我指著永生的耶和華起誓,我今天必定要為你辦成這件事情.」如此熱情的人,在最豐富的時間仍然謹守神的律法,實在難得.這樣美的人結為夫妻,就成了最美的婚姻,是愛的神蹟,也是神的工作.
\newline
\newline
我總覺得家庭與教會是不可分割關係；家庭是神工作的目的,創造之時就設立家庭,在宇宙的將來,父的家裡也是家庭.因此若家庭有主的同在,世界就有希望,撒旦也無從施計,凡是危機時代,屬靈的偉人都是由屬靈的家庭造就出來,這是時代的希望.
\newline
\newline


\newpage

\section{第59屆~港九培靈會~戴紹曾~第一講}
\label{sec:1269}
\textbf{亞伯拉罕}
\newline
\newline
連結: \href{https://www.hkbibleconference.org/session-message/view/1269}{\texttt{https://www.hkbibleconference.org/session-message/view/1269}}
\newline
\newline
\hyperref[sec:454]{< < < PREV SERMON < < <}
~
\hyperref[sec:toc]{[返目錄]}
~
\hyperref[sec:1271]{> > > NEXT SERMON > > >}
\newline
\newline
兄弟有機會再次來香港赴港九培靈研經會,深感榮幸；有機會和許牧師,金牧師一同配搭,覺得很高興；謝謝主的恩典,讓我們能夠同心合意興旺福音.我們來自不同地方,我們的樣子也不一樣；但我相信,我們有同一的心,要作耶穌基督忠心的門徒,我們的禱告是求神的國降臨；當我們看見那些破口在神的家,我們有個心巴不得能成為堵破口者.
\newline
\newline
詩篇一○六篇和以西結廿二章,二處經文都題及「破口」兩字.詩一○六23題到「摩西站在當中」,下面小字有「破口」兩字.在當時有個破口,摩西堵住那個破口.
\newline
\newline
可能有些青年和弟兄姊妹們覺得奇怪,神的國怎會有破口?我們讀聖經就慢慢明白了:神的百姓遠離祂之時,那就是破口──他們漸漸冷淡,效法世界,陷於罪惡中時,那就是破口；神的刑罰會臨到祂的國,臨到祂的百姓；也許刑罰是藉著敵人的攻擊,或許是藉著大自然的災害,或其他的方法.神要審判祂的百姓；可是,如果有人站起來為神的國代禱,帥領神的百姓歸回神的面前,使他們謙卑回到神的腳前；神的忿怒將會離開他們.作此事者就是站在破口的人.
\newline
\newline
摩西是個很明顯的例子.以色列人在西乃山,犯罪作惡,拜金牛犢,說:「這是帶領我們出埃及的神.」神要刑罰祂的百姓,神的忿怒要臨到他們；但是摩西害怕,他禱告說:「神啊!這是祂的百姓.」摩西藉著禱告站在破口.
\newline
\newline
以西結書所題的例子,更為明顯.神說祂的百姓犯罪；廿二章題到先知,祭司,國家首領,甚至眾民,「先知同謀背叛,他們如獅子抓撕掠物,那些先知吞滅人民,搶奪財寶,結果使這地多有寡婦.」(25)神的僕人竟然作惡如此!我們看見國度的破口.神說:「其中的祭司強解我的律法,褻瀆我的聖物,不分別聖的和俗的,也不使人分辨潔淨的和不潔淨的,又遮眼不顧我的安息日.」(26)這是祭司的所為,「其中的首領彷彿豺狼抓撕掠物,殺人流血,傷害人命,要得不義之財.」(27)正像先知阿摩司所講的.當時社會腐敗至此!作首領的沒有按公義帶領國家,理應為民榜樣,而他們卻為自己的享受,地位,前途,家庭而作惡多端,不顧百姓,以自我為中心,他們同流合污.「國內眾民一味的欺壓慣行搶奪,虧負困苦窮乏的,背理欺壓寄居的.」(29)連一般的百姓都敗壞了,他們的生活和那些領導階層不相上下.因此,神說:「我在他們中間找一人重修牆垣,在我面前為這國站在破口防堵,使我不滅絕這國；卻找不著一個.」(30)
\newline
\newline
神的國有許多破口,神尋找堵破口者.巴不得在座青年有此心志；靠主的幫助,為神的榮耀,為神的家,為神的國,願作堵破口者.
\newline
\newline
年前我在新加坡,因早兩年間經濟不景,以致許多人相當困難；所以基督徒舉行座談會,討論基督徒應採取的行動.講台上,有的人講經濟為何不景氣,有的說政府應採取措施解決問題.到了最後,有位陳弟兄,他是政府議員,他很愛主,在國會常勇敢為主作見證,他說:「剛才眾說紛耘,莫衷一是,我現在根據神的話告訴大家,我們國家為何經歷諸多痛苦呢?」他的解釋和以西結書廿二章大同小異,他指出基督徒的家裡,教會中,許多敗壞的情況.我聽了,感到好像舊約的先知；他呼籲大家要悔改,基督徒應該回到神面前.
\newline
\newline
希伯來書第十一章,特別題到「亞伯拉罕因著信,蒙召的時候,就遵命出去,往將來要得為業的地方去,出去的時候,還不知往哪裡去.他因著信,就在所應許之地作客；好像在異地居住帳棚,與那同蒙一個應許的以撒,雅各一樣.因為他等候那座有根基的城,就是神所經營所建造的.」(8-10)亞伯拉罕是個站破口的人,是堵破口者.
\newline
\newline
如果我們細讀創世記,我們會發現有兩個可怕悲劇的循環.神創造天地,神愛世人,愛亞當,夏娃；可是神所造的人遠離了祂,違背祂的話,結果,神的忿怒臨到他們；不僅臨到亞當和夏娃；也臨到當時全世界的人；洪水泛濫之後,神所救的.唯挪亞一家人,可惜他的子孫也犯罪作惡；甚至他們設法防備神的刑罰；因他們恐怕神要毀滅世界,他們一方面偏行己路,一方面想盡辦法防備神的審判；那時,神把他們分散世界各地.聖經記載,神的忿怒兩次臨到世界；此時,亞伯拉罕就出現在歷史的舞台上,他站在破口,他如何堵破口呢?
\newline
\newline
亞伯拉罕對認識真神有巨大貢獻,在當時有許多人,也可以說,當時所有的人不認識獨一無二的真神,他們所拜的是自己手所造的,他們不認識這位創造天地的主宰,這方面的真理世人不了解,不知有位真神,祂是真實信實的.
\newline
\newline
亞伯拉罕成為神手中的器皿,救贖世人的媒介.
\newline
\newline
神揀選亞伯拉罕,從他的子孫,耶穌基督降生在人間；從他的地位,我們看見神使用他是何等奇妙,不只傳遞真神的事實；神而且藉著他將救世主送到我們中間.
\newline
\newline
亞伯拉罕一生之中,有很多事蹟可作我們的榜樣,在許多的考驗,許多的事上,都能幫助我們,無論是他的軟弱,他作為我們的警戒；他的得勝,作為我們的勉勵,他實在是個偉大的人,是舊約聖經的偉人.如果我們細讀新約聖經,題到亞伯拉罕至少有74次:是全舊約人物無一可比擬的,可能有大衛或摩西:但亞伯拉罕確是個偉大人物,猶太人稱他為國父.耶穌題到他,有時耶穌責備猶太人:「你們說亞伯拉罕是你們的父,但你們所行的為何不像亞伯拉罕呢?若你們真是亞伯拉罕子孫,理當以亞伯拉罕的心為心,他何等信神,順服神.」司提反也講到亞伯拉罕的信心,講到他受割禮的事,講到撒拉的墳墓,講到亞伯拉罕在迦南沒有立足之地,唯有一地埋葬他的妻子.這也是偉大的見證.保羅也提及亞伯拉罕,許多神學思想都傳講亞伯拉罕的信心,特別是因信稱義的真理.希伯來書第十一章也是論及這位近行神旨意的亞伯拉罕.創世記第十二章開始記載亞伯拉罕的生平.
\newline
\newline
我覺得很有意思,亞伯拉罕的生平,一開始頭三個字「耶和華」.各位在座的青年!我們應認清一件事,這位偉大的亞伯拉罕,他的開始是基於耶和華,他一切長處在乎神:並非他自己的努力,自己的聰明:乃是真神在他身上的作為,他走第一步時,神在前面呼召他.(創十二l)神的恩典在我們身上也是如此,並非我們先作甚麼:乃是我們的神先動工.
\newline
\newline
亞伯拉罕蒙召時有幾方面,我們當特別注意,神呼召他要離開本地,本族,父家.
\newline
\newline
可能我們中間許多朋友,覺得這樣的信息很動聽:因為面對一九九七,叫我出來,正合我意.兩三年前,我在溫哥華,華人學生冬令會,有一青年對我講他心中的困擾,到底留在溫哥華,還是回香港呢?他明知神要用他:但他想在加拿大,神也可以用他.
\newline
\newline
當時亞伯拉罕離開吾珥,該處是文化相當復興之地,是商業中心:留在那裡多麼舒適,將去的是個偏僻的地方.事實上,他不曉得神召他往何處去；神呼召他離開本地,本族,父家,最主要的是神要用他,要他尊神於首位.亞伯拉罕愛國家,愛同胞,愛父家；但我相信,他也愛他的生命:當神呼召他時,神對他說:「你是否願意將我置於首位?在你自己的國家,財物,家庭之上,你是否願意將自己擺在我的手中?」
\newline
\newline
在你我身上,不一定有這樣的考驗:但我相信,原則上一定會有:因為神在我們心中要居首位.各位青年,你可記得耶穌問彼得說:「你愛我比這些更深嗎?」神呼召亞伯拉罕出來,就是這個意思.
\newline
\newline
當我十幾歲時,被困在日本集中營,有位英國宣教士,綽號蘇格蘭飛人,因他善跑.一九二四年他參加世運,恰巧參賽項目在禮拜天,而他是個虔誠基督徒,注重主日:他向教練提出因禮拜天他不能參賽:可是項目已定,無法更妀,經教練再三勸導,他卻堅持不參加:甚至當時威爾斯公爵相勸說:「難道你不愛國嗎?」答說:「我愛國,但我更愛我的神.」無論如何,都不能勸服他:後來只好讓他放棄原定在主日舉行的百咪賽跑,而參加另日舉行的四百咪賽跑.準備比賽前,有個選手遞交字條給他,內容是撒母耳記上的話:「尊重我的,我必尊重他.」賽跑結果,他不但得冠軍,且破世界記錄.到了第二年,他離英赴華,在天津傳福音廿年之久,抗戰期間,他把妻兒先遣返加拿大,打算稍後返國,可是卻被日軍捉到集中營,我在營中認識他,他對我的影響很深,他愛國家,但是他更愛神的國:所以神吩咐他要出來,他就順服.
\newline
\newline
求主幫助我們,叫我們像亞伯拉罕一樣,遵命而跟隨他的神,服在神的手下:因為神對他說:「要往我指示你的地方去.」各位青年!今天我們是否願意服在神手下?雖然前路不得而知:可是偉大的神呼召我們時,我們要順服.有時我們想,神要呼召我們出來是件艱難的事:可是彼得說:「你們為此蒙召,好叫你們承受福氣.」
\newline
\newline
蒙召者承受福氣.神對亞伯拉罕說:「出來吧!跟隨我吧!我要指示你當走的路,我要賜福給你.」神的呼召必定帶著祝福:但何時來,何種祝福,有很多不同的情形:可能是物質方面的祝福,更可能是屬靈的祝福:可能是今天的,是明天的,是將來的,或是永遠的祝幅.
\newline
\newline
神使我們成為別人蒙福的器皿.神對亞伯拉罕說:「你也要叫別人蒙福,地上的萬族都臨到我們,我們不可自私享受,主的恩典應當流露出去；如同主耶穌說:「......人若渴了.可以到我這裡來喝.信我的人......從他的腹中要流出活水的江河來.」(約七37-38)感謝主!不單是止我們的渴,不單祝福我們,也要使別人蒙福.撒瑪利亞婦人正打水之時,耶穌告訴她活水的道理；她得到以後,回到自己的村莊；本來那些人拒絕她,這次她回去；因為她認識了耶穌,自己得到滿足；神的恩典從她裡面流露出去,她把村莊的人帶到耶穌面前.我們意料不到這人能作個人佈道,她有過五個丈夫,現在和她一起生活的並非她的丈夫.她要作宣教士根本沒可能；可是她認識了耶穌,不但自己得到滿足,她的一生也全面改變；而且她開始叫別人蒙福.保羅對提摩太說:「你自己所領受的,你應當告訴別人,傳給別人.」在座弟兄姊妹!你曾否想過,神的恩典臨到你身上,不單為自己的享受；神要你的家,你四圍的人,你的學校,全香港,全世界,藉著你也蒙福.
\newline
\newline
神的呼召也帶著考驗,神呼召亞伯拉罕出來,他的生平自始至終受考驗,年青到年老一直受考驗.迦南是流奶與蜜之地；可是他一進迦南地就遇到飢荒.
\newline
\newline
當保羅得到馬其頓異象,由亞洲傳道到歐洲時,第一站是腓立比,他帶領一些人信耶穌,也趕鬼；但忽然被下監,當晚在監牢裡,他曉得順服神,必須步步經過考驗.
\newline
\newline
亞伯拉罕也是如此,等了漫長的時間才看見神的應許應驗；到了最後,甚至神要他把兒子獻上.各位青年!我們看見亞伯拉罕也有軟弱,失敗,這些可作我們的警戒.亞伯拉罕開始出吾珥到哈蘭,就停頓下來,只走了一半的路,是否他不肯前往呢?是否其他的人如羅得對他有影響呢?我們不得而知,但聖經清楚記載,神叫亞伯拉罕出來往迦南地去.為何半途停頓?他到了迦南地,遇饑荒,相當困難；他出來時信神,為甚麼在迦南地不能照樣信神呢?他還要離開迦南下埃及；這樣一來,帶給他很多的困擾.到了埃及,他不但離開了神的道路,他瞞騙別人,險些害了他的妻子,可能在那裡夏甲開始跟著他們.他在埃及發了財；羅得也發了財,接著他們就有紛爭,因為那一次的軟弱,帶給他許多痛苦；神應許賜兒子給他,卻等了好久,他急不及待,想盡辦法要幫助神完成神的應許,一而再的向神要求用另外的方法,家僕可否作我的後裔,可否承受神的應許?神說:「不,雖然他忠心；但不是我的方法.」撒拉對亞伯拉罕說:「耶和華使我不能生育,求你和我的使女同房,或者我可以因她得孩子.」多少時候,人想用自己的方法解決神的困難,用盡方法要幫助神,卻沒想到神的時候還未到.感謝神!神憐憫亞伯拉罕.
\newline
\newline
亞伯拉罕能順服神,羅得雖然自私,過著腐敗的生活,亞伯拉罕能夠饒恕他.到了最大的考驗,神吩咐他把獨生子獻上,他也甘心順服.事實上,神所要的不是他的兒子,神要的是亞伯拉罕的心；當他把兒子擺在祭壇上,神的聲音再次出現,神知道亞伯拉罕的奉獻是完全的奉獻.
\newline
\newline
各位青年!有一件我認為極寶貴的事必須要講,我們看見亞伯拉罕的許多得勝；但有一事常被人忽略,就是他如何關心他的家,他信神吩咐他出來,他是個愛家的人.創世記廿四章載,亞伯拉罕年紀老邁,他的妻子撒拉已經去世；他想到兒子婚姻的事,吩咐僕人不可在拜偶像之地為兒子找配偶,也不可帶兒子回老家；因為神呼召他出來,應許將這地賜給他和他的後裔.他吩咐僕人回故鄉,物色一個以他的心為心的人,也就是他肯出來肯順服神；他也盼望兒子和媳婦有同樣的心,他清楚認識,這個家是事奉神的家,他要兒子知道,神如何地帶領他；要媳婦像自己有絕對順服神的心.
\newline
\newline
在座為父牧者,巴不得你們有亞伯拉罕的心!多少時候我們為子女考慮婚姻問題,我們所考慮的不像亞伯拉罕所考慮的條件.亞伯拉罕實在是個偉大的父親,他所認識的神；聖經說:亞伯拉罕的神是信實,慈愛,天天與亞伯拉罕同行,是以馬內利的神.在許多考驗中,神與亞伯拉罕同在,祂是全能的神,是使無變為有的神,是公義的神,是聽禱告的神,是永恆的神.
\newline
\newline
感謝主!亞伯拉罕深覺神是創造天地的主宰,他將一切都擺在神的手中；亞伯拉罕就成為一個堵住破口的人.
\newline
\newline
弟兄姊妹!亞伯拉罕有信心,他的信心寄托在偉大的神身上.
\newline
\newline


\newpage

\section{第59屆~港九培靈會~戴紹曾~第二講}
\label{sec:1271}
\textbf{約瑟}
\newline
\newline
連結: \href{https://www.hkbibleconference.org/session-message/view/1271}{\texttt{https://www.hkbibleconference.org/session-message/view/1271}}
\newline
\newline
\hyperref[sec:1269]{< < < PREV SERMON < < <}
~
\hyperref[sec:toc]{[返目錄]}
~
\hyperref[sec:1274]{> > > NEXT SERMON > > >}
\newline
\newline
亞伯拉罕是個偉大的見證人,當全世界的人不認識獨一無二的真神時,神就興起了亞伯拉罕；不但把神學的心理擺在人間,更把活的見證帶了出來.他在日常生活中經歷神,他相信神,他是信心之父；昨天晚上,我們所注意的,不是他信心的偉大,而是他信心的對象,這位真神是偉大的.
\newline
\newline
約瑟也是為神國堵破口的人,他的生平非常重要；我們看聖經的篇幅可以斷定其重要性.我們都曉得亞伯拉罕是偉人；創世記記載亞伯拉罕和記載約瑟的篇幅是一樣的,可見寫這段歷史的人認為年青的約瑟,他的見證,他的生平之重要.
\newline
\newline
雅各的眾子中,唯約瑟得到雙重的祝福,他有兩個兒子作了以色列的支派,唯約瑟有此福份.新約聖經一連七次題到約瑟的事.
\newline
\newline
約瑟生長於哈蘭,他是拉結所生的兒子；在他年青時,他們全家搬到迦南地,所以他對哈蘭地的生活相當熟悉,迦南地的情形,他也了解.約瑟是他父親所疼愛的兒子；因父親愛他過於愛其他的兒子,所以引起其他兒子的不滿；這並非約瑟之錯,但卻叫他遭受很大的困難.他夢見神將來要使用他,天真地告訴哥哥們,引致他們懷恨在心；有時他在父親面前報告哥哥們作的壞事,他們更不高興.後來哥哥們把他賣到埃及為奴,因被誣告而下監；到了神的時候,約瑟被高舉,不但救了他的家,也救了一個國家.
\newline
\newline
約瑟的生平,我相信大家都很熟悉；他所經歷的考驗相當嚴格.詩篇一○五篇題到此事,作者說,神用祂的話來考驗約瑟.這位青年所作所為,非常忠心；正如主耶穌所說的,如果你們在小事上忠心,將來在大事上也忠心；如果在小事上開始不忠心,將來誰敢把大事交給你呢?約瑟的父親派他看守羊群,他忠心去做,叫他送信,看看哥哥們的情況如何；雖然路途危險；因為那一帶的人,恨惡雅各的家,隨時有被殺害之虞；可是當雅各吩咐他去時,這青年遵命而行.
\newline
\newline
各位青年!你是否學會一個功課,如果你不能順服父母,就很難順服真神,世上的父,我們不順服:天上的父,我們怎能順服呢?當然,若世上的父親教我們違背聖經的真理,那我們要先聽從神的話.可是若在自己的家沒有學習順服,我們在神的家怎能順服呢?
\newline
\newline
約瑟在小事上順服主,他下到埃及,在波提乏的家管理職務,小心翼翼,甚至在監牢,他的舉止,也使獄長看出他的精明,具有行政才幹,所以委以重任.神與約瑟同在,一切事情都順利,他懂得忠心做事,無論在人前人後,大事小事,都是忠心耿耿.
\newline
\newline
約瑟經歷的考驗:
\newline
\newline
第一方面的考驗:常受虐待,被人冤枉.他的哥哥們恨他,想把他置於死地；結果他們不敢動刀,藉著別人賣他到埃及為奴.約瑟到了埃及為奴.約瑟到了埃及,忠心在波提乏家裡做事:豈料波提乏的妻子誣告他,以致無辜入獄:事實上他問心無愧,卻被強指為惡徒.約瑟在監裡安慰其他的同囚者,還幫助他們解夢；他為酒政解夢說:「三天之內,法老必叫你官復原職.但你得好處的時侯,求你記念我.」酒政滿口承諾；可是出監後卻把約瑟忘記了.約瑟兩年之久,痛苦地過監牢生活.約瑟雖經百般考驗,他始終持守信仰,沒有放棄對神的信心.他深信神是慈愛全能.
\newline
\newline
各位青年!雖然你我沒有像約瑟的遭遇；但我發現許多青年基督徒,他們第一個考驗是從家庭開始,信了耶穌之後,在家裡經歷信仰的考驗；可能父母為兒女有所打算,反對他們走奉獻的道路,於是考驗就臨到了.親愛的弟兄姊妹!當我們遇到冤枉的事,被人誤會,這種考驗實在難受!當人把我們忘記,讓我們無辜經歷痛苦,你能否經得起考驗?當不如意的事臨到,我們常會感覺為何有此遭遇,難道神沒能力嗎?為何我有病?為何我有困難?為何我遭反對?神啊?你在哪裡?
\newline
\newline
第二方面的考驗:約瑟遇到甚大的引誘,是較狡猾的考驗,其實也是很自然的一種考驗,相信在座的青年,弟兄姊妹都會有這樣的考驗.約瑟生成俊美,這樣的青年較容易遇到考驗,約瑟年十七,八歲,難道他沒有性慾嗎?雖然我們說他是偉人,我們覺得他是偉大的榜樣；但他有人的性情,和我們一樣.他在波提乏家裡所遇到,實在是嚴重,厲害的考驗.波提乏的妻子是他的主母,聖經說她以目送情給約瑟；天天考驗他.人第一次得勝,若不小心,很容易鬆懈.保羅說:「自己以為站立得穩的,須要謹慎!免得跌倒.」(林前十12)主耶穌受洗,從水裡一上來,天開了,有天父的聲音說:「這是我的愛子.」聖靈也降臨在他身上.屬靈方面,這是高峰；但是接著,耶穌進到曠野,經過很大的考驗.是的,我們經過了屬靈高潮以後,容易遇到更大的考驗.約瑟很聰明,他曉得危機四伏,他也知道自己的軟弱,他不敢也不肯玩火.許多基督徒不及約瑟的謹慎；有時我在世界各地旅行講道,在一些基督徒家裡發現有黃色小說,我心裡非常難過,尤其青年人看這類的書；有的家庭牆上掛了那種不可見人的裸體畫,許多的淫穢對青年是極大的考驗.約瑟平時不隨便進主人的屋子,除非有事情要辦,他曉得自己的軟弱,不給魔鬼留地步.關於這方面,我們應當謹慎,要防魔鬼的詭計.我想,當時約瑟遇考驗,他很容易解釋,很容易自圓其說；埃及的道德觀和迦南的道德觀根本不一樣,亞伯拉罕敬畏神,以撒敬畏神,我父親雅各也敬畏神；但此地的人都不認識真神,大可以同流合污.今日許多青年的道德觀是根據多數而定；然而約瑟的信仰,他的道德觀有很清楚的根據,他可以解釋,他的哥哥猶大也犯淫亂的.如果你讀創世記卅八和卅九章,是要給我們今時代的人強烈的對比；一方面我們看見猶大的失敗,一方面看見約瑟的得勝；我們看見有人失敗,還看見也有人得勝.其實,約瑟盡管可以說是那女人採取主動,她是主人,年長又有權柄,她怎麼吩咐,僕人應唯命是從；而且在埃及,無人相識,不必掩人耳目,行動可隨便,可以自由；男孩子自然有性慾,任其發揮,身體是神創造的,聽其自然.可能這些思想都在約瑟腦海中；甚至到了一個地步,魔鬼趁機作祟,慫恿,引誘,叫人懷疑神的話.好像亞當,夏娃受引誘,當蛇開始試探時,對他們說:難道神真的說嗎?神知道你吃了這果子,眼睛會明亮,神不要你有此享受.
\newline
\newline
今天有的青年感覺,若作個真基督徒,太不自由；其實,在耶穌基督裡面,我們才知道甚麼叫做真自由.
\newline
\newline
約瑟經過多方的引誘,仍然持守純潔,可見神在人間可以保守一個青年過聖潔的生活,幫助我們過祂所喜悅的生活,神是全能的,祂呼召我們要聖潔.帖撒羅尼迦前書五24告訴我們:「那召你們的本是信實的,祂必成就這事.」感謝主!靠著自己沒有法子過聖潔的生活；但是約瑟靠著神才能持守聖潔.
\newline
\newline
第三方面的考驗:約瑟被高舉,有重大的責任交在他手裡,在家有重任,他盡忠職守；在波提乏的家,在監牢,都負重任,且忠心完成.後來他作了全埃及的首相,他仍謙卑盡責,不論管錢管事,他都很忠心,不隨便,不貪污,是非清楚,公私分明,人事管理方面也很有智慧.
\newline
\newline
約瑟被哥哥們賣到埃及,經過十三年的痛苦,從十七歲開始走下坡路,直到卅歲；這是在人的解釋；但是在神的解釋,這階段是裝備時期,神要裝備祂所重用的僕人.按人看,約瑟走的是下坡路,被賣到埃及,在波提乏家裡為奴,被主母誣告,被下監牢,被人遣忘,經過十三年的痛苦.約瑟懂得饒恕人.他的哥哥們也有十三年的痛苦；但我寧可跟約瑟走那十三年的下坡路,不願跟他的哥哥們走他們痛苦的道路；表面上看他們沒有痛苦,但他們的內心非常痛苦,他們遭遇困難,良心發動,想起已往在弟弟身上作了壞事,必是神施行報應；至於約瑟卻沒把十三年的痛苦記在心裡,他真是衷心饒恕人.神要用的人,就像約瑟一樣.
\newline
\newline
如果今天你坐在神面前,但在你心深處對人不滿,可能你覺得父母太對不起你,你不肯饒恕他們.兩年前兄弟在東馬之沙巴地方聚會,會畢有一女子和我傾談,她說她沒辦法饒恕父親；因為父親害她.或者你不能饒恕的非你父親,非你的孩子而是別人,是教會中的弟兄姊妹,哦!各位!我們不饒恕人,神怎饒恕我們呢?我們將成為被綑綁的人；因我們不肯釋放別人,自己就變成不饒恕別人的奴僕.當約瑟升高,有權柄,有地位之時,他仍然持守他的謙卑.約瑟成功的秘訣何在?他怎能得勝?怎能持守他對神的信,持守他的純潔,過敬虔的生活呢?約瑟生平之中,他知道神是真實的,他真是實踐神與他同在.
\newline
\newline
第一,約瑟先求神的喜悅,不求自己的享受,自己的喜好,他以討神喜悅為前題,這是他成功的秘訣.當約瑟遇到引誘；波提乏的妻子要與他同寢,約瑟對她說:「我怎能作這大惡得罪神呢?」(創卅九9末句)他沒想到肉體的享受,他的伯父以掃曾犯了這樣的錯；當他餓得將死之時,就將長子的名份賣了,為的是換取碗紅豆湯.但是約瑟卻不為這些肉體的享受,作惡得罪神.
\newline
\newline
第二,他看見神的大能並依靠神的能力,正如撒迦利亞書中神說:「不是倚靠勢力,不是倚靠才能,乃是倚靠我的靈,方能成事.」亞四6約瑟有此觀念,他自知無能,他站在法老王面前,法老說:「我聽見人說,你能解夢.」他答:「這不在乎我.」他知道是神的能力在他身上.約瑟看見神的能力,一切完全倚靠神的能力.
\newline
\newline
多少時候,主給我們恩賜才幹,我們不知不覺以為是自己的.保羅告訴哥林多教會,一切所有的是從神那裡來的.
\newline
\newline
今晨,我在恩泉堂講道,我想到馬禮遜,他是來華的第一個基督教宣教士.當他從英國要登輪啟程時有人問他到哪裡去?去作甚麼?答:「到中國去傳福音.」那人就嗤笑他說:「你以為你能影響那古老的國家嗎?」他答:「我不能,神能.」馬禮遜有這樣的看見.
\newline
\newline
在座青年!如果我們靠自己的本事,靠自己的能力；我們常會失敗；當我們倚靠神的大能大力,我們會看見神得勝.
\newline
\newline
第三,約瑟知道神與他同在,所以不論環境如何惡劣,心中仍有平安喜樂；因為他在主裡面.「約瑟住在他主人埃及人的家中；耶和華與他同在,他就百事順利.」(創卅九2)「凡在約瑟手下的事,司獄一概不察,因為耶和華與約瑟同在,耶和華使他所作的,盡都順利.」(23)約瑟深深體驗到以馬內利；神與他同在.
\newline
\newline
猶太人和中國人有同樣的習慣,給孩子起名都有很深的含意,不像外國人的名字都沒有意義,在字典也查不出來.約瑟的長子名叫嗎哪西,意即「神使我忘記我的困苦.」他看見小嗎哪西時,神使他多年痛苦盡消.活在神裡面的人,實在有此享受!並非沒痛苦,人非草木,難道當時他離開父家,心裡有平安嗎?難道他不思家嗎?然而神使他忘記一切的痛苦.約瑟次子以法蓮,意即「神使我昌盛.」約瑟得勝,有地位,有權柄,發財,他不以為是自己的事,他知道一切都是出於神；因他活在神裡面.
\newline
\newline
第四,約瑟知道神有祂的計劃,他年青時作夢,深知神有祂的計劃,到了神的時候,神必成全；所以約瑟安心等待神的時候.這是約瑟得勝的秘訣.他服在神計劃之下.
\newline
\newline
很多基督徒有他自己的計劃,盼望神的計劃符合自己的計劃,否則就不理會神而偏行己路；可是神的道路高過我們的道路.
\newline
\newline
約瑟曉得神有美好的計劃,他將自己完全擺在神的手中,將一切事都交托神；因他曉得神必成就.
\newline
\newline
一九四○年我們全家在山東的煙台,那時我父母傳福音的工作已經多方受阻；因日軍已佔領中國沿海一帶,進入我父母工作的地區.我父親決定把全家帶到美國；因為留在中國太危險,且妨礙工作,船票已定；但家父懂得禱告,在那段日子不住地禱告,他對主說「不是照著我的計劃,乃是照著你的計劃.」到了時候,家父有感動,覺得不應該離開中國,他把四個孩子留在內地會學校,他們到陝西,西安,五年半之久沒和孩子們見面.
\newline
\newline
弟兄姊妹!有時我們覺得門已開了；約瑟或者遇到這樣的試探,波提乏的家門開了,騎上埃及的馬回老家迦南地去,不需受此監牢之苦.約瑟卻順服在神計劃之下,他信神的應許.到了他臨終之時,約瑟對他的弟兄們說:「我要死了,但神必定看顧你們,領你們從這地上去,到祂起初所應許......之地......你們要把我的骸骨從這裡搬上去.」(創五十24-25)約瑟知道人終歸過去；但神的話,神所應許不會過去,約瑟信神的話.神的話安定在天,永不改變,不會過去,昨日,今日,直到永遠不改變.約瑟求神的喜悅；他知道神的能力,他服在神能力之下；他知道神與他同在,所以能經得起極其困難的環境；他知道神有祂的美意,他說:「神啊!不是照著我的意思,乃是照你的美意.」約瑟信神的話,所以救了他的家,也救了全埃及人.當時有很大的破口,神興起約瑟,一家將滅亡滅種之時,神興起約瑟,埃及將亡國遭遇大饑荒之時,神興起祂重用的僕人──約瑟.
\newline
\newline
求主幫助我們思想神的話,將神的話記在心中.
\newline
\newline
各位青年!約瑟可作我們的榜樣,我們的神沒有改變；只要按照這些原則過我們的人生,你也要成為堵破口的人.
\newline
\newline


\newpage

\section{第59屆~港九培靈會~戴紹曾~第三講}
\label{sec:1274}
\textbf{摩西的生平}
\newline
\newline
連結: \href{https://www.hkbibleconference.org/session-message/view/1274}{\texttt{https://www.hkbibleconference.org/session-message/view/1274}}
\newline
\newline
\hyperref[sec:1271]{< < < PREV SERMON < < <}
~
\hyperref[sec:toc]{[返目錄]}
~
\hyperref[sec:1276]{> > > NEXT SERMON > > >}
\newline
\newline
謝謝主的恩典帶領我們到了第三個晚上,昨晚我們實在看見神的恩典在我們中間；不只是許多青年聽道並受感動,也有約180位到前面來；在主面前要對付,要獻身.今晨,其中的一位告訴我,他昨晚實在得釋放；已往總覺得很捆綁,有不良習慣；但經禱告,主斷開那些鎖練.他昨晚回家,和他未信主的父親提起受洗的問題,出乎意料,他父親沒發怒也沒反對.我們當為他禱告.願神在我們中間工作.
\newline
\newline
在人類歷史上,當神的百姓遠離神,犯罪作惡時；神的忿怒定會臨到他們身上,但是神常興起一些人,他們用禱告把極不理想的環境扭轉.第一個晚上,我們看見神所造的人,雖然他們明明看見神的創造,神的傑作,應當知道有位真神；可是他們不認識神而偏行己路,此時,神興起亞伯拉罕,把神的真理帶進人間,他是位見證人,見證宇宙中有位真神；祂不只創造天地,祂更愛世人的人.亞伯拉罕作為神手中的器皿,藉著神的僕人,藉著他的後代,耶穌基督降生,作世人的救主.亞伯拉罕實在是位為神的國堵破口者.
\newline
\newline
昨晚我們看見約瑟的見證,他的父親在迦南地,是神所應許的；可是約瑟的哥哥們不認識神,他們走自己的路,神的刑罰要臨到他們.約瑟下埃及,表面看似是悲劇,但約瑟有屬靈的眼光,他完全知道自己是神手中的人,所以最後他對哥哥們說:「差我到這裡來的不是你們,乃是出於神；因為祂要救人,要救我們的家,要救埃及.」約瑟實在是位偉大的堵破口者.
\newline
\newline
摩西在舊約之中實在是偉人.聖經記載,摩西面對面與神說話,好像朋友一般.這話極其寶貴,從他身上給我們很多學習的功課.摩西是神所興起的領袖,雖然他不願意.我想在座有些青年,和摩西一樣,覺得自己不能作甚麼.可是神仍然要揀選使用.
\newline
\newline
詩篇一○六23告訴我們,當神的忿怒要臨到以色列民,神所揀選的摩西站在破口.摩西是堵破口者,率領以色列民出埃及；神藉摩西的手將律法交給以色列人.因時間關係,不能細講這方面,但十條誡命確是非常重要的真理,可作人生道德的基礎；不僅在摩西時代,歷代以來,十誡成為道德的最高標準,是全世界所認同的.摩西也教導他的同胞如何敬畏神,並且教導他們認識神的本性,帶領他的同胞一直走到迦南地邊界.
\newline
\newline
摩西生平可分為三個階段,他活在世上一百廿年.所以有人就把他從一至四十歲,為第一階段,第二和第三階段,也是四十年.也有人解,這三個四十年完全不一樣,摩西的態度不一樣,第一個四十年,摩西的態度是「我能」,第二個四十年是「我不能」,經過失敗,他不敢再抬頭；但是摩西終於進到第三個階段「神能」.
\newline
\newline
在座中可能有如摩西的情形,過去曾想事奉主卻失敗了,因靠自己；而現在你也落在「我不能」的階段.
\newline
\newline
當內地會要從中國大陸撤退時,神興起孫德生牧師,帶領內地會在東亞一帶事奉；謝謝主!他直到今日仍忠心事奉神,他有點像迦勒；迦勒到了85歲,說:「四十年前我的力量如何,現在我的力量也如何.」數日前接孫德生老弟兄的信,他仍到世界各處為主作見證；相信在座中多有從他的書,他的見證得益不少.
\newline
\newline
摩西第一個四十年的階段,他的家庭背景,表面上沒有甚麼,他的父親在埃及為奴；當時法老下令,凡以色列人所生的男孩要殺盡,法老所作的正和神的應許相反.第一,法老不要以色列人生養眾多；但神應許亞伯拉罕的後裔要像天上的星,海邊的沙,多不勝數.第二,法老怕以色列人強盛,與敵聯合攻擊而後脫離埃及.但是約瑟曾對他的同胞說:「有一天神必帶領我們離開埃及,別忘記把我的骸骨帶回老家.」神已應許將迦南地賜給亞伯拉罕和他的後裔；法老卻極力加以攔阻.摩西的父母很虔誠,這對敬畏神的夫婦,真是把當時的局勢扭轉.希伯來書十一23記載信心的偉人.其中有這對奴僕,無人注目；可是他們向神有絕對的信心,他們也愛小兒子,為他冒了自己性命危險,不怕王命,保存小摩西的性命,本是不可能的事,但神是掌管一切；歷史並非掌握在法老的手,雖然法老自以為是；實際上,歷史掌握於神,神用奇妙方法救了摩西；這對敬畏神的父母,他們用敬虔的心教導兒子.由此可見,父母在家庭的影響實在很大,他們真是如同約書亞說:「至於我和我家,我們必定事奉耶和華.」
\newline
\newline
基督徒的父母,千萬別小看自己,請看這對為奴的夫婦,他們的影響多大!今天我們在自己的家,也當扮演這樣的角色,不要把這責任推給教會,推給主日學,推給青年團契.父母當負起責任,作子女的榜樣,教導孩子學習敬畏耶和華.
\newline
\newline
摩西教育方面的背景,司提反題到摩西:「摩西學了埃及人一切的學問.」(徒七22)真是出人意料!一個奴僕的兒子竟然上了埃及最有名的大學；在神奇妙計劃中,摩西學了埃及人一切的學問,成為一個哲學家,他很有口才,行政本領的力量很大,更是個偉大的工程師；他也會寫詩,他是將軍；他在埃及,神實在恩待他.
\newline
\newline
前些日子,我和朋友談到日本大學的情形,我想,在香港,新加坡,台灣,都差不多,進大學之前,生活相當緊張；因為被錄取者很少,所以應考者都很拼命；一旦考取,十分鬆懈,校園生活放蕩隨便.我想摩西在埃及和他的同學也是如此；因為當時社會腐敗,淫亂不法的事也多,摩西在此環境之下易受影響,是同流合污呢?或是中流砥柱呢?這是第一方面的危險.保羅提醒我們「不要效法世界.」可是許多青年在大學時,入社會時,開始效法世界:世界的印,印在他們的心,印在他們的生活上.第二方面的危險,學問多,容易自高自大,因對世界,對大自然較為了解,所以傲視一切,覺得不需要神,可以獨闖天下.實在很危險!並非我們不需要高深的學問:但要求主保守我們有謙卑的心.
\newline
\newline
摩西五經,作者大概不是摩西:因為當中有句「世上最謙卑的人是摩西.」最少這句話不是摩西寫的.路易斯寫的地獄來鴻,大魔鬼給小魔鬼寫信,教導如何試探基督徒,後來大鬼發現小鬼敗北:大鬼寫信責備小鬼,並教牠進行下一步試探,叫那基督徒謙卑,使他謙卑到了個地步,他就驕傲起來.魔鬼實在很狡滑,常使人謙卑,然後叫人因自己的謙卑而驕傲.學問高者更要小心!求主保守我們有謙卑的心.
\newline
\newline
第三件事,摩西須作終身的決定,那是改變他一生方向的決定.希伯來書十一章24-26論及這事.我發現有三個K(國語音)音:肯,可,看,意即摩西不肯稱為法老女兒之子,他寧可和神的百姓同受苦害,也不願暫時享受罪中之樂.他看為基督的凌辱,比埃及的財物更寶貴.摩西怎會有此決定?何來原動力?動機是甚麼?他的態度願意遵行神的旨意,願意把神的旨意置於首位,願意服在神的手下,願意尋求神的國和神的義.同時,他看見同胞的痛苦,動了憐憫的心,他關心同胞,愛神也愛同胞:所以摩西作了個非常難得的決定；因為他的四圍都是世界,很容易被吸引,摩西勝過世界的秘訣在於信心,他拒絕法老皇宮的地位,拒絕暫時的喜樂；罪中之樂,財寶,享受,他一概不要,連埃及這個富有的國家,他都拒絕,希伯來書十一章兩句話說明其所以然,「因他想望所要得的賞賜.」「想望」並非眼前的事,是遙遠的事,這是成熟的,屬靈者的決定；不成熟者只顧眼前,摩西是想望所要得的賞賜.第二,摩西看見了那不能看見的主,這是很寶貴的,他唯一看見主,其餘一切他都不要；因為一切都要過去,只是暫時的,也不能徹底解決他內心的問題,更不能解決來生的問題.以掃和摩西完全相反,以掃沒有看見主,只想到自己將死,死了一切都完結；神的應許也沒有了,長子的名分與自己無關.但是摩西看見主,所以有重要決定.可惜摩西在第一階段失敗了；因他用自己的方法,用血氣進行計劃,雖然他知道神要用他；他耐心等候.摩西看見埃及人欺侮他的同胞,他熱血奔騰,急不及待；他雖有埃及人的一切學問,但到底他是以色列人,他終於把那可惡的埃及人打死:他動了血氣,所以不能成就神杓事.「因為人的怒氣,並不能成就神的義.」(雅一20)摩西需要很多的學習.
\newline
\newline
摩西第二個四十年神給他長時間在平靜地方學習,他對神有更深一步的認識.我想,摩西和一般常人一樣,只覺得行動要緊,一停頓下來就焦急.神把摩西帶到平靜之地,使他單獨得神的教導.
\newline
\newline
保羅得救之後有三年之久在阿拉伯曠野.
\newline
\newline
我又想到宋尚節,在他還未出來傳道之前:神也給他頗長的時間安靜.當時人都說他發瘋,把他關在精神病院,就在那時他詳讀神的話:與人隔絕,時刻與主同在.
\newline
\newline
摩西在曠野,神教導他,這是很重要的學習:神裝備祂的僕人.摩西如同約伯,可以說:「從前風問有你,現在親眼看見你.」摩西經過四十年的學習,清楚地,徹底地認識創造天地的主宰.
\newline
\newline
當我讀完大學時,很有負擔要傳福音,許多朋友勸我不必讀神學,從速出去傳福音,不用花費時間裝備.感謝主!主攔阻我,使我有安靜的時候學習.
\newline
\newline
摩西早年所接觸的都是博學者,在曠野都是普通的牧羊人:摩西向他們學習:因為將來需要帶領一班平凡的,為奴的人.摩西離開埃及皇宮後,進入工場,有很好的學習.我覺得神學院也應有這樣的學習:我們頭腦複雜,如果能向一般工人學習,可能福音更有效地傳開.摩西已往頭腦充滿理論,都是屬靈的知識:這時他深深體驗神的真實,他看見神的真實:這時神呼召摩西.
\newline
\newline
神的呼召是摩西第二階段的最高峰,四十年快過去之時,有一天摩西在曠野放羊,他看見很奇妙的事；在曠野,荊棘遍地叢生不以為奇；但摩西看見非凡的事,忽然荊棘被焚燒而產生一種榮燿,並沒有被火焚毀:摩西所見不過是普通東西,但有了火就成為榮耀的東西.他想到神的作為,平凡的東西在神手中,就成為榮耀的東西:從神而來的榮耀和力量,在神手中,荊棘和樹都沒有燒毀,沒有受害.摩西學習,神的手在人的身上帶著榮燿,不會叫人受損,摩西學習,他看見神的慈愛.
\newline
\newline
神吩咐摩西脫下鞋來,這裡有兩件事我們必須注意:第一,聖潔之地,必須脫下鞋來,神要摩西認識神的聖潔.另有一意思,脫下鞋子,不需自己行動,乃需神動工.神對摩西說:「你手中拿的是甚麼?」當時摩西滿口的解釋,推辭說:「我是甚麼人,憑甚麼回到埃及,我要講甚麼,誰叫我去呢；我去傳甚麼信息呢?我憑甚麼權柄呢?我沒有口才.神啊!你要差遣誰就差遣誰.」摩西想盡藉口.我覺得摩西很可愛,因為我們也是這樣.神說:「摩西!我要用你.」神發怒了.希望在座青年,如果遇到這樣的情形,不要一味推辭；恐怕到了時候,神的忿怒臨到.神問摩西:「你手裡拿的是甚麼?」摩西手裡的東西,不是大學的書,不是埃及之物；乃是一根牧羊的杖.神說:「把杖丟在地上,你手中之物算不得甚麼；把你自己交給我,我要使用你.」把手裡所有的交給神,讓神分別為聖；微小的東西將成為多人蒙福的東西,正如五餅二魚使五千人飽足.
\newline
\newline
各位!你手裡拿的是甚麼?你是否願意擺在主手裡?
\newline
\newline
卅五年前,當我在港大學國語,初出來工作,主日在九龍禮拜.有次主日,牧師講的一個見證,是我一輩子都忘不了,現在和大家分享；那位牧師說及國內一位弟兄,他又聾又啞,覺得自己不能作甚麼；他禱告,想到有件事他能作,所以每清早他從村中走到河邊,那裡沒有橋,來往必須涉水過河,這位老弟兄免費背著人們過河,每次完成任務,他就給人一份福音單張；因此帶領許多人信耶穌,可能比村中牧師所帶領的更多；他又聾又啞,沒有機會求學,沒機會學習其他職業；主愛在他身上,他用這簡單的方法表達基督的愛,帶領許多人信耶穌.
\newline
\newline
你手中持有的是甚麼?如果過去你推辭,自覺沒有甚麼,以為主不能用你；我認為這是自我欺騙.我們面對世上各樣事情,做買賣,就業,有時野心很大,要創自己的天下；可是一旦聽見主說:「你手裡拿的是甚麼?」你回答說:「沒有甚麼,你不會用這個的,只好留下我自己用.」求主憐憫!摩西完全把自己擺在主的手中,今日你的光景如何?是否在「我能」的時候?你要向摩西學習；如果你要成為堵破口者,必須曉得「神能」「我不能」.
\newline
\newline


\newpage

\section{第59屆~港九培靈會~戴紹曾~第四講}
\label{sec:1276}
\textbf{摩西的見證}
\newline
\newline
連結: \href{https://www.hkbibleconference.org/session-message/view/1276}{\texttt{https://www.hkbibleconference.org/session-message/view/1276}}
\newline
\newline
\hyperref[sec:1274]{< < < PREV SERMON < < <}
~
\hyperref[sec:toc]{[返目錄]}
~
\hyperref[sec:1278]{> > > NEXT SERMON > > >}
\newline
\newline
感謝讚美主的恩典,數日來天氣晴朗,讓我們不受攔阻赴會；我也很感謝主,每晚帶領許多青年聚集.我們用禱告的心聽主的話,深信主的美意將我們聚在一起；不是聽講員的話,乃是聽主的話.
\newline
\newline
舊約的聖經,神的國有破口,神興起人堵住破口；新約聖經,教會也有破口；在教會歷史中,教會也有破口；到了時候,神興起祂揀選的人堵破口.從舊約聖經,我們看見亞伯拉罕,約瑟,摩西,這樣的偉人.教會歷史,我們看見神興起馬丁路得,約翰衛斯理,宋尚節,王明道,都是神所興起的忠心僕人；在不同時代,面對不同的問題堵破口.
\newline
\newline
在以西結書,神尋找堵破口者,神說:「我在他們中間,尋找一人重修牆垣,在我面前為這國站在破口防堵,使我不滅絕這國,卻找不著一個.」(結廿二30)這句話很可怕,創造天地之主宰尋找合乎祂用的人,卻找不著一個.
\newline
\newline
是否現今時代也是這樣?神正在尋找合乎祂用的器皿要堵破口,會不會找不著一個呢?
\newline
\newline
昨晚我們講到摩西的見證,他的生平可分為三個階段:前四十年在埃及,我們也看見他的家庭和教育的背景,以及各方面的裝備,並且他一生最重要的選擇,那選擇是摩西生平的轉捩點.第二個四十年是在曠野安靜地學習神,安靜中有機會多思想神的存在,神的真理.摩西在那段時間對神的認識愈加清楚,他看見燒著卻並沒燒毀的荊棘；他看見普通之物,有了神的同在就成為榮耀之物.我們也是極普通的人,只要有神同在就不一樣了；如同保羅所說:「若有人在基督裡,他就是新造的人；舊事己過,都變成新的了.」摩西在曠野看見荊棘,荊棘不過是極普通的植物；但當神充滿荊棘,就完全不一樣了；且神的榮耀對荊棘並沒有損壞,沒有燒毀.神的手在我們身上也是如此.摩西走近荊棘,聽見神的聲音,吩咐他脫下鞋子,因所站的是聖地,摩西此時學習神的聖潔；並且他知道有行動的時候也有安靜的時候；接著,神對他說:「你手裡拿的是甚麼?」神接過摩西手裡的杖.杖算不得甚麼；但在神手中神可以使用,真是奇妙的事!神吩咐摩西應當如何進行,而摩西竭盡所能推卸責任；但他覺得一生是個見證.神的命令帶著成功的能力,並非神的命令臨到我們之後祂便袖手旁觀；當神的命令臨到,同時也帶著成功和成就這事的能力.
\newline
\newline
今天晚上我們要思想摩西的第三個四十年.
\newline
\newline
摩西一生之三分二是在裝備,只有三分之一完成神給他的使命,此中有深刻的意義,在這四十年的階段有很多的事情,今晚不可能都講.摩西從曠野回埃及,向法老要求給神的百姓離開埃及,法老硬心不肯；神降十災臨到埃及和法老,經十災之後,法老才勉強給摩西領以色列民出埃及.他們在曠野,經歷神的同在,也經歷神的回應,在西乃山有神的啟示,神將偉大的十條誡命交給摩西,摩西教導同胞謹守十誡,並教導他們如何敬畏神；他們造了會幕,其中有很多意義.以色列人出埃及之後,常發怨言,沒有感謝的心,真是忘恩負義!埋怨神,說:「我們吃的,喝的,穿的在哪裡?」今天早晨,金牧師清楚告訴我們,這些東西固然要緊；但更重要的是先求神的國.以色列人忘記了救他們的神,在西乃山下鑄了個金牛犢,說:「這是領我們出埃及的神.」想不到他們竟然如此!
\newline
\newline
在座中蒙恩的弟兄姊妹,常常忘記主在我們身上的恩典；主為我們流寶血,受盡十字架的痛苦,我們都忘記了.
\newline
\newline
摩西帶領以色列人到迦南地邊界,這是神的計劃,是神的應許,是整個歷史的去向；約瑟臨終之時也提及此事,每一代的青年聽父母講故事,他們滿懷希望有一天承受神所應許的迦南地；但不料當以色列民來到迦南地的邊境,畏縮不前；因害怕那地的人高大強壯,城牆堅固；他們打算另立一人帶領他們回埃及去.
\newline
\newline
摩西經歷七種不同的考驗,一個堵破口者是經得起考驗的.從多章經文歸納七件事值得我們注意:
\newline
\newline
第一,摩西面對妥協的考驗,神之目的和計劃要救以色列民出埃及,且應許把迦南地賜給他們；但是法老說他們可以在埃及作神的選民,可以敬拜他們的神.(今天教會中有許多人,不肯與世界分手.)摩西堅守立場,認為作神的選民,不單敬拜神,還應當離開埃及.法老答應給以色列人離開；可是仍然要留在近處.我想今天也有很多人遇到同樣的考驗,或有人對你說:「你信耶穌就信,何必這樣熱心.」教會面臨同樣的危機,有時我們只想到自己跟隨主,沒有帶領全家歸向主,事奉主；為父母的基督徒,不看重兒女靈性的問題,信或不信讓兒女自由決定；但是卻嚴厲教導兒女做人處世的事.為兒女籌劃接受優等教育,教導兒女生財之路；他們只關心兒女屬世的事,而屬靈方面漠不關心.法老只答應讓男丁離開埃及,婦孺及其他須留下.摩西說:「我們神家,神國的人都要走.」摩西不肯妥協.法老又說:「你們男女老少都可以走；不過,羊群牛群要留下.」我們中間有許多人,身心獻給主,但事業和工作留屬自己,禮拜時屬於主,上班時是自己的事.
\newline
\newline
摩西不肯留下任何東西在埃及,凡所有的都要帶走；他遇考驗,總不妥協,他不半途而廢,也不敢取巧行走捷徑.他抱著破釜沉舟的精神,將全人,全家獻上給主；一切所有全都獻給主,絲毫不留下.
\newline
\newline
在座的青年們!求主給我們有這樣的心志.「不妥協」也是王明道的見證.求主給我們堅定不移的信心.
\newline
\newline
第二,摩西面對不可能的事.一九八七年我們面對著一九九七年,似乎無路可走,前途茫茫.各位!我們看見摩西,他的異象似乎不可能,這個生在埃及奴僕之家的人怎能帶領一切為奴的走出奴僕之家呢?他們沒有軍隊,沒有錢財也沒有勢力,怎可能完成這事呢?真是夢想,狂想!各位!神呼召我們,就是呼召我們作不可能作的事.感謝主!多少時候,我們選擇的可能斷定我們一生的方向；多少時候,神要我們選的是件很難的選擇；我們平時按著自己的軟弱選擇易行的路.我們看見個很具體的例子,當摩西率領以色列人出埃及,剛走不遠,前面有紅海,後面有埃及的追兵；實在進退維谷,前無去路,後無盼望；但是神為摩西開路.保羅說:「你們所遇見的試探,無非是人所能受的......」有時我們自己覺得受不了,沒法得勝；可是保羅告訴我們,當遇到這樣的事情時,主一定為我們「開一條出路,叫我們能忍受得住.」(林前十13)當以色列人前面無路,主為他們開路,海水分開,他們走乾地過去.這本是不可能的事,然而神與他們同在.也許今晚你到禮拜堂來時,你覺得在你前面是不可能的事；主要你作的你作不到.但若全心依靠主,你不能作；主能作,祂要為你開條出路；你將看見前面紅海的水分開.摩西在曠野,帶領二百萬同胞前行,民生問題天天要解決,曠野一無所有；但神是信實的,供應他們一切的需要.親愛的弟兄姊妹!我們的神沒有改變；如果我們先求衣,食,住,神的國怎麼降臨呢?我們若先求神的國,神即刻為我們預備吃的,穿的,住的.想起這一百廿年神在內地會的恩典,不但在戴德生的時代,神供應一切事奉祂的人的需要,今天在東南亞千位同工,神繼續供應.摩西的神是我們的神,當日在曠野供應那麼多的人；神是以利亞的神,是耶和華以勒,祂必供應我們的需要.只要我們真心相信主,順服主,完成主給我們的使命,抓住祂的信實當作我們的食物；祂要賜給我們屬靈的能力,工作成功的能力.戴德生曾說:「在天國的事,神將大事交給我們,要我們完成,平常三個階段,第一階段「不可能.」第二階段「困難」.第三階段「完成」.巴不得弟兄姊妹們可以經歷這個真實；不可能的事變成困難,困難的事終於達到完成.
\newline
\newline
第三,摩西能接受別人的建議,有的時候,精明有把握且受過高等教育的領袖,不易採納別人的意見；但是出埃及記十八章記載,以色列人到了西乃山,摩西的岳父葉忒羅來了兩天,就發現事有蹊蹺,他看見女婿這個青年從埃及帶領全以色列民出來,理應值得驕傲；但他是聰明人,也很關心女婿,他指出摩西所作的事不聰明；因摩西獨當一面為眾以色列民解決所有的問題,無論大小事都需摩西處理；所以他提議應當找些聰明人互相分工合作.摩西接受這個提議,實在是件很美的事.申命記卅四7記載,摩西死時一百廿歲,眼目沒有昏花,精神沒有衰敗.因為他懂得接受別人的意見.
\newline
\newline
新約聖經的彼得,當耶路撒冷教會人數增多時,彼得和其他使徒照顧不來,以致教會發生不愉快的事,有些人被忽略,怨聲四起,教會面臨分裂；彼得了解這個困難,說:我們的責任,乃是以傳道祈禱為事,行政方面的事,應由其他執事負責；所以就選舉七個執事,以負專責.於是問題得以解決,教會日益發展,人數增多.因為有好的領袖,知道怎樣分工合作,把責任交給能勝任的人.今日華人教會需要這樣的心,教會中不應由一人抓權,一手包辦,許多教會的問題都由此而起；若像摩西謙卑自知力量有限,別人也有恩賜；當時摩西有見及此,他帶領二百萬以色列民；神給他智慧,接納別人的建議.這也是一種考驗,摩西經得起第三個考驗.
\newline
\newline
第四,禱告的重要,有時危險臨到,我們也許覺得行動須快才能解決問題；但摩西重視禱告,主與他面對面說話,如同朋友.他的禱告不限於有需要,他的禱告有感恩,有讚美,有敬拜,有認罪,有代禱,真是圓滿的禱告!
\newline
\newline
我們的禱告只知道求此求彼,我們雖是教會中人,卻無異於外邦人；我們有需要時才禱告.俗語云:「平時不燒香,臨急抱佛腳.」教會中也有類似的事.求主憐憫!
\newline
\newline
摩西真正知道禱告的重要,當亞瑪力人來攻擊出埃及的以色列民,摩西打發約書亞帶領以軍與亞瑪力軍交戰,他退到山上禱告,以色列人便得勝.第二次我們看見摩西站在破口,(出卅二)就是詩篇所講:「摩西站在他們中間.」他知道神的忿怒要臨到祂的百姓,就站在破口,用禱告堵住破口.巴不得今天華人教會更多的弟兄姊妹,用禱告堵破口.
\newline
\newline
在出埃及記卅三章中,摩西的禱告有三點:第一,神的救贖,神對摩西說:「你的百姓,」摩西不接受這句話,說:「這百姓是你用大力和大能的手,從埃及地領出來的.」
\newline
\newline
第二,神的聖經,摩西說為甚麼使埃及人議論,神領他們出去原有惡意,是要降禍與他們,到了曠野將他們從地上除滅.這樣豈不是羞辱神的聖名嗎?外邦人要說,你有能力救他們脫離埃及,卻沒有能力帶他們進迦南地；有能力作應許而難道會落空嗎?親愛的弟兄姊妹!摩西禱告神:「神啊!為你名之故......」摩西所關心的是神的名.
\newline
\newline
第三,神是信實的,摩西禱告神:「求你記念你僕人亞伯拉罕,以撒,雅各；你曾指著自己起誓,必使你們的後裔像天上的星那樣多；如果現在滅絕這些百姓,你的應許豈不歸於徒然嗎?」摩西抓住神的應許呼求,神聽了他的禱告.摩西的禱告堵住了破口,他懂得禱告的重要,挽回神的忿怒.
\newline
\newline
在座的青年們!為你的親人禱告,為你全家歸主向主禱告,摩西的神是你的神,祂如何垂聽摩西的禱告；讓我們抓住主的信實,仰望主的名,記念主的救贖,主不但為你,為我死在十字架上,為你的父母,為你的全家.巴不得今天晚上,我們用禱告堵住破口,捆綁魔鬼撒旦；看見主的憐憫.讓我們學習摩西的禱告!
\newline
\newline
第五,摩西勝過野心的考驗,這事和剛才所講的有連帶關係,此次考驗好像直接從神而來,好像神要考驗祂的僕人,神對摩西說,由得我吧,我要滅絕這百姓,然後叫你的後裔,成為大國,不是亞伯拉罕作國父,是摩西作國父；不是亞伯拉罕的子孫,是摩西的子孫.摩西以神的名為重,不看重自己的名；若神不赦免以色列民,情願被除掉生命冊上的名.這裡我們看見摩西得勝.在神的國裡太多人有野心了；所建立的不是神的國,是建立自己的地盤,自己的國.摩西說:「神啊!不是我的名,乃是你的名,願你名得榮耀!」摩西是極謙和的人,經得起野心的考驗.
\newline
\newline
第六,摩西勝過嫉妒的考驗,民數記十一章記載,有一青年向約書亞報告,有兩人被靈感動說預言.約書亞真可愛,他很擁護摩西,認為摩西最偉大；他對摩西說:「有二人被靈感動說預言,請禁止他們.」摩西對他說:「你為我的緣故,願耶和華把祂的靈降在他們身上.」感謝主!摩西看見聖靈作工作時不願禁止,且盼望他們每個人都受靈的感動；因他曉得勝靈真正作工時,無人能干預.
\newline
\newline
記得十年前兄弟在這裡,特別題到巴拿巴從耶路撒冷走到安提阿時,看見一群無名英雄,到那裡為主作見證,主賜福他們的見證,很多人歸向主；巴拿巴毫不妒忌,相反地,他為了看見神的恩典而歡喜,不是他自己的工作,是這班沒有受過神學訓練的弟兄姊妹的工作,他看見神使用他們,歡喜快樂.
\newline
\newline
摩西在曠野時,看見神的靈在他們中間,他不願意干涉,他巴不得更多的人真正受聖靈感動.如果你帶領青年工作,主日學,看見神的恩典在別人身上,不要妒忌,而當感謝；要像摩西一樣,不容約書亞干涉,巴不得眾百姓都受聖靈感動.當然摩西所指是聖靈真正的工作.我們看見以色列以往的歷史,利未支派獻祭用其他的火和祭物,神的忿怒臨到倆祭司身上.
\newline
\newline
第七,物色和訓練接捧人,這對以色列的領袖是個很大的考驗,對教會也是一種考驗,對今日華人教會更是考驗.怎樣物色接班人,我想今天中國大陸需要工人,在大陸之外的華人教會很需要忠心事主的工人.個人事奉主不難；但到了工作將告一段落時,有沒那樣的勇氣,有沒寬大的心胸,物色最好的人選?摩西為此事禱告,(民廿四)向神要求接棒人:因他曉得不可能活到永遠,總有一天要過去:雖然他盼望能進迦南地,他有一次禱告,神不許；他覺得需要接棒人,他把以色列人帶到一個地步,另有領袖要帶他們往前走:摩西禱告時,神指著約書亞給他,說:「我的手在他身上,我要將我的靈賜給他；你要教導他,勉勵他,為他作榜樣.」摩西很有智慧,常帶領約書亞,好像他的門徒一樣:這是耶穌教導門徒的方法,馬可福音一章記載,耶穌揀選十二門徒,叫他們常和祂同在,向祂學習,且打發他們出去工作.摩西就是用此方怯,以身作則:約書亞看見這位偉大的領袖禱告的生活,謙卑的態度,果斷的精神,和他的信心,約書亞悉心學習.今日教會也應學習這樣的作法.約書亞終於順利地接棒.
\newline
\newline
我們今天看見摩西所遇見的考驗,這些考驗有時會臨到我們.妥協的考驗,不可能事情的考驗,接受別人意見的考驗,禱告重要的考驗.禱告是否在我們一生佔為優先?野心的考驗,嫉妒的考驗,預備接棒人的考驗.摩西經得起這些考驗:但摩西也有他的軟弱,失敗,因一次的發怒,沒有在百姓之中榮耀神,神不讓摩西進入迦南地:摩西禱告,神對他說:「別再提此事,你不能進迦南美地去.」或者你覺得這是悲劇,摩西終身事奉神,跟著天上的異象教導他們,帶領他們到迦南邊界,摩西卻不能進去:可是他的接棒人,在摩西給他的基礎上前往,完成神的美意.這是成功的見證,雖然摩西本身沒有進迦南地:可是承受他的工作者,在摩西的工作上,把以色列人帶進去:還有,過了兩千年,摩西所預言的基督來到世上,耶穌帶領門徒上山,向他們改變形像,摩西和以利亞在耶穌的兩旁,摩西終於進了迦南地,看見神的榮躍.
\newline
\newline
親愛的弟兄姊妹!你我前面的道路將有許多考驗,堵破囗者靠主能經得起一切的考驗.求主給我們信心,但願主感動摩西的靈,也感動我們.
\newline
\newline


\newpage

\section{第59屆~港九培靈會~戴紹曾~第五講}
\label{sec:1278}
\textbf{約書亞和迦勒}
\newline
\newline
連結: \href{https://www.hkbibleconference.org/session-message/view/1278}{\texttt{https://www.hkbibleconference.org/session-message/view/1278}}
\newline
\newline
\hyperref[sec:1276]{< < < PREV SERMON < < <}
~
\hyperref[sec:toc]{[返目錄]}
~
\hyperref[sec:1279]{> > > NEXT SERMON > > >}
\newline
\newline
今天晚上,我們要看兩位將軍:約書亞和迦勒,他們是很齊心的同工,且也是堵破口者,他們時常彼此勉勵.他倆都是生長在埃及,在摩西帶領之下出埃及:只有這兩人從埃及出來,經過曠野,進入迦南地:其他的人都倒斃在曠野:這兩位堵破口者,終於承受神所應許之地.在以色列支派中,南部的猶大和中部的以法蓮,這兩支派所領受的地較多:迦勒屬於猶大支派,約書亞是以法蓮支派:神把重要的地賜給他們,非無原因:他們忠心,有屬靈的雄心,非屬世的野心.神賜福他們.
\newline
\newline
我們可以從約書亞和迦勒學習重要的屬靈爭戰原則:因為今日的基督徒也參加屬靈的戰場,我們是基督的精兵:所以保羅寫信給那時代的基督徒,要他們作剛強的人,要穿上全副軍裝,拿著信德當盾牌.信心就是盾牌:魔鬼射箭時,信心可以吞滅一切的火箭.(弗六)信是屬靈的武器,主給我們信心,主所賜的全副軍裝:我們自己原來沒有這些,是主所賜的:如果我們沒有,應該向主求,我們才能得剛強,才能打勝仗.
\newline
\newline
當約書亞和迦勒奉差做探子,(民十三,十四)根據申命記第一章,我們才曉得,打發十二探子窺探聖地,是出於百姓的主意,所以摩西就對百姓說,我們不如打發十二個人進去,每支派一個代表,看看該地如何.到底他們的動機何在,是否要看神所形容流奶與蜜之地是不是真實,那裡的人怎樣,該地值不值得爭取?所以百姓出此主意.過了四十天,那十二個人回來,抬了一些水果:有葡萄,無花果,石榴,都是迦南地的出產.十個人走出來,說,這地果然照神所形容的,真是流奶與蜜之地,這些豐富的出產足以瞪明,「然而......」.當我讀大學唸心理學,我覺得很簡單,至今都忘記了:可是我記得一件事:yes,but,not,他們先說「是」,繼說「但是」結果「不行」.雖然那麼好的地方,「可惜」,「然而」,這些字充滿了不信.我想,今日的基督徒有此情形,認為神的應許不錯,可是我多麼軟弱,我的家庭,環境,學校等等的壓力,社會惡劣:神的話不錯,可是.......十個探子報告,敵人身體高大,城牆堅固,我們渺小,又無武器,我們不過是一群從埃及出來的奴僕,也未經作戰的特別訓練,我們在他們面前如同蚱蜢,被腳一踩就完了.這個報告多麼可憐!他們最大的錯,就是眼中無神,他們所看的是人,是環境,是自已,他們沒有仰望神,沒有想到神的應許,神的能力.至於約書亞和迦勒,覺得那十個人所報告都沒有錯:但他倆充滿了信心,雖然情形惡劣；可是神若幫助,誰能敵擋呢?他倆不過是少數人,卻沒受多數人的影響；他們看透了迦南地,信心曉得神能幫助.他們必然勝過迦南人,他們說,迦南的巨人將成為我們的食物.他們有這樣大的信心,因他們認識主,他們站住以信心堵塞破口,在那不信的世代,他們的眼睛能看見他們所看不到的主:正如亞伯拉罕一樣.
\newline
\newline
在以色列人的歷史中,有個相當有名的先知──以利沙在多丹地方:有一次,他的僕人清晨起來,發現敵人昨夜進攻,圍繞村莊.敵人主要的目標是以利沙；僕人非常懼怕,回來報告,敵人已經圍繞多丹城.以利沙禱告神,求神開僕人的眼睛.
\newline
\newline
今天晚上,兄弟有這樣的禱告,求主開我們眾人的眼睛.在座的青年們,弟兄姊妹們!求主給我們看見,與我們同在的,比敵對我們的還多；反對我們的雖然很厲害,看上去我們沒有辦法；可是以信心的眼睛,我們曉得必能得勝.
\newline
\newline
神的話是我們必需的武器,在屬靈的爭戰上,神的話非常重要；所以保羅說:「拿著聖靈的寶劍,就是神的話.」
\newline
\newline
當約書亞出來要作全以色列國的領袖,第一章告訴我們,摩西已經死了,以色列人站在約但河東岸,約書亞帥領以色列民準備進迦南地；神將得勝的秘訣告訴約書亞:「這律法書不可離開你的口,總要晝夜思想......如此你的道路就可以亨通,凡事順利.」(書一8)成功的秘訣在耶和華,大衛也有此認識.他說:「惟喜愛耶和華的律法,晝夜思想,這人便為有福.」(詩一2)如果你要真正蒙福,要把神的話藏在心裡.詩篇一一九篇,大衛問個問題,然後自己回答:「少年人用甚麼潔淨他的行為呢?是要遵行你的話.」可能在座青年一直思想這個問題,今天淫亂的世代,腐敗的社會,青年怎樣才能過聖潔的生活呢?大衛的回答是要遵行神的話；大衛還說:「我將你的話藏在心裡,免得我得罪你.」感謝主!這是必勝的秘訣.
\newline
\newline
在座青年有沒有把神的話常擺在心裡呢?
\newline
\newline
主耶穌是我們的榜樣,祂受浸之後進到曠野,撒旦一而再,再而三引誘試探耶穌；每次主耶穌用神的話──聖靈的寶劍抵擋魔鬼,勝過牠,叫牠離開.
\newline
\newline
神的話帶著能力,大衛說:「你的話是我腳前的燈,是我路上的光.」神的話在屬靈爭戰時使我們得勝；使神的工作在我們手裡成功.
\newline
\newline
今晨金牧師提到中國教會的復興；我覺得復興重要因素之一是神的話.
\newline
\newline
一百八十年前有位青年馬禮遜,離開英國來華,他是基督教來華第一個教士；如果你對會史熟悉,你就曉得他不是來華的第一位宣教士.早在唐朝,景教的教士已經到了長安,他們作了翻譯的工作；可是沒把整本聖經譯成中文,這是第七世紀；到了第十二世紀,羅馬天主教法蘭西修會,也派代表到北京,那是元代之時,作翻譯工作,且設立孤兒院；可是整本聖經沒譯成中文.到了第十六世紀,耶穌會的宣教士利馬竇也來了,明朝晚年,他們作了很多工作,有很多方面我們可以跟他們學習；可惜他們沒有把整本聖經譯成中文.從第七世紀至十九世紀,雖不斷有宣教士來華；直到馬禮遜來時才有整本聖經譯成中文.我們感謝主,雖然馬禮遜在工作上有許多缺欠；但這事對華人教會有極大的貢獻.景教,法蘭西聖會,耶穌會,有他們的時期,到了時候,全部離開,工作幾乎消失了.到了一九四九年,華人教會的宣教士都要離華,許多牧師被迫做普通的事,教會面對很大的考驗,那時候約有百萬基督徒；今日有人說是四百萬.記得六年前兄弟在這裡,有晚,我作了報告；後來有位青年來找我,他很可愛很直爽地說:「戴牧師,你講錯了；浙江省現有基督徒大約是五百萬,這是最近的統計.」今天全中國信徒,有人說是五千萬,我不曉得是否有這麼多；不過,我知道現在比卅六至卅八年前多了不少倍.在文化大革命時,許多弟兄姊妹們的聖經被沒收；聽說在廈門,紅衛兵燒聖經三日三夜之久.在星洲,我有一手抄本,是西安一個老弟兄送給我的,文革期間,他親手抄錄羅馬書,他借人家的聖經,他愛神的話.在風雨中,基督徒靠神的話站穩,屬靈的戰爭需要神的話.
\newline
\newline
親愛的弟兄姊妹!我們需要神的話,一九九七年我們需要神的話；不但聽神的話,更要遵行神的話.耶穌開始講道時說:單聽神的話不夠,要聽也要行,這樣就好像個聰明的建築師,把你的靈魂的柱子建造在磐石上,雨淋,風吹,水沖,也不至倒塌；因建造在神的話語上.由此可見信心的重要,神的話之重要.
\newline
\newline
我們是否尊主為主?約書亞第五章記載,約書亞是全國的領袖,他看見耶利哥城有個人站著,他不知此人是耶利哥人或以色列人；他就走過去對那人說:「你是誰?是幫助我們還是幫助我們的敵人呢?」答:「我是萬軍之耶和華軍隊的元帥.」祂吩咐約書亞脫下腳上的鞋子,意思要約書亞認清楚,不是他作元帥,而是他跟著這位偉大的主,主怎樣吩咐,他怎樣作；他要服權柄,要尊主為主.我想,在約書亞一生之中,這真是個轉捩點；過去他跟隨摩西,如此偉大的領袖；可是他不能依靠這些,他有許多經歷；現在他要認清楚,在屬靈爭戰中,不是我們作主,乃是神作主.所以耶穌說:「不是憑口說主啊主啊!」
\newline
\newline
許多基督徒,憑著嘴講主；可是生活是為自己,為己而活；依靠自己的計劃,不是靠主,不是順服主.神說:「約書亞!要把鞋子脫下,要順服權柄.」求主幫助我們!
\newline
\newline
我們在這幾章聖經中,可以看到順服的重要,也看到不順服的可怕,我們看見亞干在耶利哥沒有順服神,也看見約書亞完全順服；神給他的命令,可說是莫名其妙的命令；約書亞是將軍,神吩咐他一連七天要圍繞耶利哥城列隊行走；他們是打仗的,卻吩咐他們繞城.
\newline
\newline
親愛的青年們!有時我們不明白神為何叫我們作這事,有時,神要我們作的事,實非我們所能理解的；但是約書亞完全順服.
\newline
\newline
有件事我覺得很有意思,當以色列人圍繞耶利哥時,他們不可出聲；平日常發怨言意見多多的以色列民,那時只好閉口不言；他們默默無聲地圍繞耶利哥城.順服神的話,雖然不了解,但順服神的話必不會錯；只要我們肯順服神,到了時候,神必給我們看見.至於亞干,或許他也不明白神為何禁止拿那些東西.大好的機會,耶利哥人被殺,留下漂亮的衣物,金銀,財寶,正合從曠野出來的人的需要；幾十年過著窮困的日子,亞干起了貪心,私自留下衣服,金銀,以為無人知道；然而神知道,神是輕慢不得的.
\newline
\newline
各位弟兄姊妹請注意!你我的行動,不但影響整個教會；你屬靈的光景,不單是你個人的事,也會影響別人,你軟弱,你為魔鬼留地步,隨從自己的情慾；不但害自己也害別人.數月前,美國雜誌刊登一些牧師的失敗,相信基督徒都非常難過；教會有名的牧師,教會的領袖,竟淪到那樣地步!各位!我不是講那些牧師,我是講我們自己,沒有一個人是為自己活的,你的影響比你所想的大得多.你沒有信心,你不順服,你講不該講的話,別人會受影響.求主憐憫!你看,亞干這人,他不肯尊主為主,他有自己的計劃；結果,全國受影響,卅六人死亡.
\newline
\newline
我們是否願意將自己完全奉獻給主,讓主改變我們的生命,給我們新的生命,給我們新的道德標準；不是照著世界的,乃是照著主的；那就是尊主為主,接受祂在我們身上的計劃；讓祂使用我們的時間,才能,一切一切.
\newline
\newline
禱告的重要與禱告的能力,在屬靈的爭戰,禱告太重要.約書亞在出埃及記十七章,已經看見禱告的重要.摩西在山上,用禱告托住平原的軍隊；我想那時約書亞清楚看見,和亞瑪力人交戰,毫無力量,毫無盼望能打勝仗；一定是摩西在山上禱告；可是,我們發現約書亞忘記禱告,他們勝過耶利哥之後,第二座城是艾城,他們沒有禱告,因為覺得這小地方,很容易打勝仗；他們如果先禱告,就會知道有問題.到了第七,八,九章講那些基遍人,事實上他們是巴勒斯坦人；但他們欺騙約書亞,說他們是從很遠處而來的；約書亞沒有禱告,結果就和他們立約,次日才發現自己上了當,他自以為有把握,其實他受騙了.到了第十章,我們發現約書亞學會禱告的重要,他和亞摩利人五個王打仗之時,他不住地禱告求問主,應不應該出發,何時出發,從何處出發?一舉一動,他都禱告,就在和敵人交戰的當天,他還禱告說:「主啊!日子不夠長,願延長時間叫太陽停留不要平西；因為要完成今日的爭戰.」約書亞明白禱告的重要.感謝主!
\newline
\newline
弟兄姊妹!主帶領我們上屬靈的戰場,我們需要以禱告勝過魔鬼；過主所喜悅的生活.
\newline
\newline
約書亞和以色列人所面對的工作,是他們面對許多未得之地.神對約書亞說:「你年紀老了,還有許多未得之地.」(書十三1)「約書亞對以色列人說......神所賜給你們的地,你們耽延不去得,要到幾時呢!」在此,我們看見迦勒,他有美好的見證.
\newline
\newline
迦勒的生平和摩西一樣,似乎可分為三個階段,年青時他作探子窺探聖地.中年時和以色列民在曠野飄流.我們沒有看見他,我相信他有很多的學習.第三階段85歲,他發現能力和45年前一樣.他有很美的見證,請各位特別注意三件事:
\newline
\newline
第一,迦勒是位力量持久的人,他不僅有個好的開始,他真是剛強到底.許多青年開始時跑得好,可是一踏入社會,就業,有家庭,漸漸就軟弱離開主；而迦勒並非虎頭蛇尾的人,他從開始到年老,力量一樣剛強.好像保羅作見證說:「那美好的仗我打過了,當跑的路我跑盡了,所信的道我守住了.」
\newline
\newline
第二,迦勒願意接受艱難的任務,摩西打發他窺探聖地,當他回來時,他們倆作信心的報告.他知道神與他們同在.到了85歲,他還是肯接受艱難的任務,他對約書亞說:摩西給我們的應許；現在我來了,我的力量不減當年,前面這座山,請給我罷!我參巧地圖,希伯崙是迦南地最高的地方,比迦密山還高,是敵人最堅固之地,要把山上敵人趕走實不容易；但迦勒不怕艱難任務.
\newline
\newline
第三,迦勒盼望全家參與工作,好像約書亞所說:「至於我和我家,我們必定事奉耶和華.」
\newline
\newline
迦勒考慮女兒的婚姻,他用奇特的方法進行.他報告,有哪位屬於同一支派的男孩,敢出來把希伯崙山上的敵人趕走,爭取那地,我就把女兒給他.大概他從中物色和自己一樣有勇氣,不怕困難,有信心的人選；他自己挺身而出,也盼望全家都出來.他如何信主,也盼望女兒女婿都信主；他怎樣順服主,盼望下一代同樣順服主,他所考慮的不但是他這一代,他也很關心下一代；所以他發出挑戰.俄陀聶奪取了那城,他是第一個士師；迦勒就把女兒押撒給他為妻.押撒野心很大,不但要上面的地,連下面的地也要.我想這是屬靈的意思,和她的父親一樣.
\newline
\newline
親愛的弟兄姊妹!今天我們看到這個見證,和我們有何關係?我們是否是個力量持久,不怕艱難的人?是否關心我們的全家,關心下一代?是否要完成神給我們的使命?這是迦勒的見證.我們還沒有把神所應許的全地取過來；未得之地,要決志爭取過來；未得之地,要決心爭取.今天基督徒面對許多未得之地,在香港有許多同胞不認識耶穌；我們有沒有想到這些未得之地,有沒有像迦勒的心志?巴不得神的國降臨,神要我們得的,我們完全得著.復活的主對我們說:「天上地下所有的權柄都賜給我,所以你們去,使萬民作我的門徒.」這是主給我們的命令.以色列人要得那塊地；我們要得全世界,把世上的人都帶到主的面前,使祂的國降臨,使祂交給我們的任務能早日完成.可惜今天千千萬萬的人不認識耶穌,否認主的教會,少有像迦勒這樣的人,不怕艱難.
\newline
\newline
感謝主!我認識一位孔醫生,他是新加坡的弟兄,他放下好的工作,離星赴非洲作福音醫療工作；好像迦勒一樣,不怕艱辛,他的薪水很少,且那地區不安全,生活簡樸；可是為了神的國,為了完成主給我們的使命,他去了.
\newline
\newline
親愛的弟兄姊妹!我們肯不肯?屬靈爭戰的四個秘訣,可幫助我們得著未得之地,我們肯不肯站起來,說:「主啊!幫助我學會打屬靈的勝仗；為要完成這個大使命.」
\newline
\newline


\newpage

\section{第59屆~港九培靈會~戴紹曾~第六講}
\label{sec:1279}
\textbf{撒母耳}
\newline
\newline
連結: \href{https://www.hkbibleconference.org/session-message/view/1279}{\texttt{https://www.hkbibleconference.org/session-message/view/1279}}
\newline
\newline
\hyperref[sec:1278]{< < < PREV SERMON < < <}
~
\hyperref[sec:toc]{[返目錄]}
~
\hyperref[sec:1280]{> > > NEXT SERMON > > >}
\newline
\newline
在以色列的歷史,神的百姓常常遇到破口；神興起亞伯拉罕,約瑟,摩西,約書亞和迦勒要堵破口；雖然他們出生的時間不一樣,當時的需要也不同；可是神都使用他們.
\newline
\newline
當我們看撒母耳生平時,也許我們會感到那時處於政治方面的危機；因在他們西邊的非利士人常常進攻威脅以色列人；但是破口不在此,真正的破口不是非利士的威脅,也非敵人的壓力.如果要知道破口是甚麼,我們必須了解以色列人的光景；我們若不注意他們內在的情形,社會,道德的墮落,他們屬靈的光景；只看他們的歷史,只看敵人的威脅；我們會忽略那破口.
\newline
\newline
當撒母耳起來時,他主要的是面對百姓屬靈方面的需要；把神的道傳給他們.撒母耳是士師和君王中間的橋樑,他是祭司；雖然他不是亞倫的後裔,他是利未的後裔.利未有四個兒子,暗蘭是亞倫的父親,第二個兒子生了一子名叫可拉,此人嫉妒摩西,起來反對；神的忿怒臨到可拉.撒母耳是可拉的後代,雖然可拉大大失敗；可是神仍然關心他的後代,也使用他的後代,撒母耳是個很美的例子,他是祭司,他的孫子希幔會唱歌；在大衛設立聖殿詩歌時,希幔負重任.我們看見,從失敗,神就興起祂的榮耀.
\newline
\newline
撒母耳是祭司也是先知,掃羅的僕人說,「這城裡有一位神人,是眾人所尊重的,凡他所說的全部都應驗.」(撒上九6)甚至新約聖經提到撒母耳說他預言耶穌的時候,聖靈降臨的事情,(徒三24)可見他是偉大的先知,當時不叫先知,叫做先見；因他先看見了.
\newline
\newline
以色列國當時的破口到底有何現象?我已講過,外來的壓力是非利士人的攻擊,許多以色列人害怕這些敵人；但這不是最重要的事情.士師記四次提到,當時「以色列國沒有王」無人領導國家,政治方面是個真空,屬靈方面更是如此!以色列中沒有王,各人任意而行,這是很可怕的情況；事實上國王太多了,而每個王都任意而行,為所欲為,不顧他人,不關心他人的需要；當時也沒有屬靈的領袖,以利老先生,在他年青時也許有其魄力,有其貢獻,可惜聖經隻字不題；當以利在聖經上出現時,年紀已經老邁,聖經描寫這位老先生,都是靜坐他的座位；大概年紀老了行動不便.年紀大是自然的現象,每個人都必走上此路；可是到此階段,許多工作不能進行,已不能扮演原來的角色了.
\newline
\newline
聖經記載以利尊重兒子過於尊重神,這話很可怕!如果我們尊重兒女或父母過於尊敬神,一定會發生問題.這個領袖怕得罪兒子,怕惹他們的氣,他沒想到影響別人；雖然以利知道兩個兒子犯罪作惡,他沒有阻止他們；沒盡為父的責任,以致他們隨己意而行.以利對兒子極其寬容；但對別人卻非常嚴格.有一天,以利在會幕坐於自己座位,他看見一婦人進來,他眼花耳聾,只見婦人嘴唇抖動,但聽不見聲音；所以懷疑這婦人喝醉酒.因為這樣的事素常發生；可見當時社會的腐敗.以利不問青紅皂白,破口就責備哈拿到底要醉酒到幾時?過了一段時間,撒母耳聽見主的聲音,次日早晨,以利和這青年講話:「昨天晚上主向你講的話是甚麼.你要告訴我；如果你不全部照實講,嚴重的後果便會臨到你.」以利對撒母耳那麼嚴格,對哈拿也那麼厲害；而對自己的兒子卻不干涉不攔阻.
\newline
\newline
我們為父母的,有沒給兒女作榜樣,教導,攔阻,鼓勵他們,盡父母的本分,在神的家,我們有沒有做到這些呢?
\newline
\newline
以利的兩個兒子,何弗尼和非尼哈很壞；第二章記載,他們不認識耶和華.我想,第二代,第三代的基督徒,都是同樣的危險.以利認識神,一輩子要順服神,可是他兩個兒子,繼承父業,根本沒有生命；雖是以利的兒子,但是他們不屬於神.可能有些生長在基督化家庭的青年,父母是虔誠的基督徒,認識主,愛主,會禱告,喜愛神的話；可是自己至今尚未重生得救,父母給了好的榜樣,在教會常有機會聽道；可是自己還沒有屬靈生命.兄弟一直到十七歲處於這樣的環境,我父母是宣教士,我生在中國內地會學校讀書；所以在家,在學校,都有機會讀經禱告,可是我沒有生命；一直到我從日本集中營出來那天,在安徽參加佈道會,我看見一些青年歸向耶穌,我受了感動；他們是第一次聽福音的,他們接受耶穌；我多次聽過福音,在我心深處有聲音說:「你是該死的.」意思是,主說:「你生長在這樣的家庭,上這樣好的學校；抗戰時期許多人死亡,我保存你的性命,有機會和你父母團聚,你不懂得感謝,還不肯將心歸我.」我看見中國青年歸向主之時,我受感動,十七歲時接受耶穌作我的救主.我常講「我的同胞」因為我生在中國,重生在中國,是在中國青年歸向主之時,我才認識耶穌.朋友!你生長在基督化的家庭；如果你沒有生命,你真是要接受耶穌作你個人的救主；你父母的信仰不是你的,他們的信仰可以感動你的心；可是不能代替你去相信.我們真是需要謙卑在主面前,接受耶穌作我們個人的救主.
\newline
\newline
以利兩個兒子,這兩個不認識耶和華的祭司,藐視神的祭物；神責備以利和他兩個兒子,聖經用「服己」來形容,他們利用宗教信仰發財,他們只想到自己的享受.眾弟兄姊妹們把祭物帶來會幕；因為他們有罪,深深感覺神的忿怒會臨到,盼望這些祭司作為中保,接受他們的祭物獻在祭壇上,幫助他們的罪得赦免,與神和好.真是可憐的光景!這兩個祭司到聖所,不是為了神和人之間的中保,乃是為自己的肚腹,為了肥己.他們貪污,在聖潔的神聖所裡與來敬拜的婦女犯姦淫.
\newline
\newline
今天在教會事奉神的人,身為牧師,長老,執事,我們扮演的到底是甚麼角色?是幫助人親近神嗎?還是只顧自己的前途,自己的享受?
\newline
\newline
各位!我們看見當時國家的破口,宗教墮落,信仰和生活完全脫節.今日基督徒多有同樣的情形,禮拜天,使徒聖經背得滾瓜爛熟；到了禮拜一上街,上班,上學,忘的一乾二淨,信仰和生活脫節了.
\newline
\newline
何弗尼和非尼哈不聽父親的話,不孝敬父母,犯了第一條帶應許的誡命.在當時,宗教的形式最重要；可是信仰的實際完全沒有了,形式漂亮,實際沒有了.我們在信仰,在宗教,有很多的符號,十字架代表甚麼?舊約的約櫃代表甚麼?這些都是符號；但這些符號已經失去內容,只成了一種形式,沒有真實,沒有意義.
\newline
\newline
當非利士人進攻以色列,以色列人害怕；有人提議請何弗尼非尼哈把約櫃搬出來戰場上,這樣,非利士人定不能勝.約櫃是神與人,人與神立約的符號,約櫃很漂亮,外面都是金的；但價值不在於金子,不在於漂亮；價值乃在於神與人,人與神立約；當人偏行己路,約櫃還有何用呢?這約我們已不承認,已經不要了；各人照自己的意思過人生,已經沒有約了.可是敵人也很迷信,他們害怕,因為他們懂得以色列人的歷史,甚至比他們更清楚；他們聽見以色列人喊叫,約櫃出來了.照敵人的判斷,帶以色列人出埃及的神與他們同在,祂怎樣救他們脫離埃及,經過曠野,進入迦南地；這位神現與他們同在,我們沒有辦法了；只好大家勇敢吧!可是那天敵人卻未遭遇到困難,很容易把以色列人打敗,兩個祭司被打死了,約櫃被擄去了.事實上,約櫃已失去意義.在舊約聖經,我們看見掃羅也有這樣的觀念,他和亞瑪力人打仗以後,沒有順服神；把上好的牛羊帶回來,當撒母耳來見掃羅時,掃羅對他講了很屬靈的話:「你吩咐我作的事我都作完了......我們把最好的牛羊留下要獻給神.」撒母耳對掃羅說,你錯了,神要的不是你的祭物,神要的是你的順服,順服勝過祭物.彌迦書六8告訴我們:「世人哪!耶和華已指示你何為善；祂向你所要的是甚麼呢!只要你行公義,好憐憫,存謙卑的心,與你的神同行.」
\newline
\newline
在教會,我們多少時候有同樣的迷信,有的人受過浸；可是沒有重生,沒新的生命；水不能代表新的生命；因為沒有接受耶穌作個人的救主.有的女孩子頸項掛個十字架,十字架是順服耶穌的意思,追隨為我們死在十字架上的耶穌.如果我們天天犯罪,明知故犯,掛個十字架不過是裝飾品,毫無意義.有的人,把聖經擱置,沒有讀經,也沒有遵行聖經的話；聖經不可能把神的恩典帶給我們；如果我們像以斯拉一樣,專心查考神的話,遵行神的話,教導神的話,這樣才有意思.
\newline
\newline
在以利那時,神的話稀少；甚至可以說神的榮耀離開了；以色列人大聲喊叫,可惜他們沒有能力,沒有辦法勝過敵人.請看!面對這些破口,神如何興起撒母耳並使用他.有解經家告訴我們,在整本舊約聖經,以色列人達到第一個高峰就是摩西帥領以色列人離開埃及,第二個高峰就是撒母耳時代,神如何地使用撒母耳,並將當時的情形扭轉過來.
\newline
\newline
「耶和華的言語稀少,不常有默示.」(撒上三1)「耶和華又在示羅顯現,因為耶和華又將自己的話默示撒母耳......」(撒上三21)在此,我們又看見一位堵破口者.
\newline
\newline
撒母耳的父母有虔誠敬畏神的態度.父母影響撒母耳太重要,真有決定性.我們曉得孟母三遷,對兒子有很大的影響.哈拿對撒母耳的影響當然很大.西諺說:「搖搖籃之手,統治天下.」這句話值得各位父母思想；這些人忠心敬拜神,雖然當時四圍是腐敗的環境,雖然以利和哈拿,到了節期,都忠心到會所去；可是哈拿帶著沉重的心情；因為好像神沒有聽她的禱告,當她很痛苦之時,以利還責備她,真是火上加油!可是哈拿仍謙卑對以利說話,她傾心吐意,把心中隱憂向神申訴.在第二章,我們看見哈拿真是個會禱告的母親.第一章記載她心中有負擔,生活有困難,她把這些問題都帶到神面前.正如詩篇所說:「當將你的事交托耶和華,並倚靠祂,祂就必成全.」(卅七5)哈拿明白了,把心裡的負搪帶到神面前,她的禱告充滿了讚美,她九次說神是主,她知道神是歷史的主,是國家的主,是她家中的主,她尊神為主:祂是救主,是聖潔的,是獨一無二的,可靠的,全知全能的,供應一切所需要的；審判世人的主.神聽了哈拿的禱告,照祂所求的把撒母耳賜給她.我想,撒母耳出世以後,哈拿或有許多的想法；這是我的獨生子:示羅地腐敗,以利兩個兒子太壞,我沒法實行所許的願,相信神會了解,不如把撒母耳留在家裡,自己好好照顧他.哈拿並非這樣的人,她怎樣許願必照看行:她的信心,相信神聽了她的禱告,把兒子賜給她,她更相信神必保守她的兒子,她把兒子帶到以利的家,她相信神是信實的,神必保全一切.
\newline
\newline
撒母耳個人的生活是很美麗的見證,無論他四圍環境如何惡劣,聖殿中有淫亂的事:可是撒母耳持守他的聖潔,過神所喜悅的生活,忠心事奉神.他也尊敬以利老先生,雖然以利兩個兒子不聽話,可是撒母耳卻很尊敬他.有一天,撒母耳聽見神的聲音,他再三地跑到老先生那裡,以為老先生叫他:他的言行舉止都有很好的表現,雖然老人家有很多的困難,到底他是神的祭司,撒母耳尊敬他,伺候他.
\newline
\newline
撒母耳也順服神,神所說的,他都作了,他留心聽神的話.在示羅地方,他聽見主的聲音,在伯利恆地方,當他要選個王代替掃羅,他的心親近主,敬聽主的話；他得到神的啟示,聽了就順服.主與撒母耳同在.
\newline
\newline
撒母耳一切公開的事奉　當神向他顯現,神對他說話,然後他忠心地把神的話傳達給以色列人(三21)他不折不扣,不保留地把神的話都釋放出來:撒母耳召集他的同胞.向他們傳道,要他們悔改,說:「你們若一心歸順耶和華,」(七3)一心就是全心,歸向並順服,要歸回神的家,要順服神,過去約櫃已經失去意義了；可是現在要重新與神立約,要一心歸順耶和華:要把外邦的神一概除掉.親愛的弟兄姊妹!這裡有個佈道大會,有個奮興大會,撒母耳呼召百姓要歸向神,要離開罪惡,離開偶像,除去那些代替神在他們心裡的事,要專心一意事奉神.
\newline
\newline
宋尚節在中國教會歷史中,他在中國內地,和海外地方,大大蒙神使用.「宋尚節傳」的英文版本,使徒德先生寫序言,他提及神使用宋尚節的四個重要原因:第一,宋尚節是完全奉獻的人.有人說:「若沒有大的捨去,就沒有大的使用.」我相信宋尚節的見證也是這樣:他如何放棄美國的生活,放棄他在大學的寫作,放棄教書的機會:甚至他的父母對他很失望,可能他的同學也是如此,他把文憑丟到大海:並非他看不起學問,而是他真正破釜沉舟地要跟隨主:因他曉得這些東西是試探,他願意不保留地追隨基督.神使用這樣奉獻的人.
\newline
\newline
第二個原因,宋尚節懂得能力的來源──神的話.他用很多時間禱告,他曉得十字架的重要:在自己的身上,生活,信息:十字架是能力的來源,還有聖靈都是能力的來源.
\newline
\newline
第三個原因,使徒德告訴我們,宋尚節十分真實,毫不虛假,他不怕得罪人:但他又常常得罪人,他問宣教士,問牧師得救了沒有.記得家父告訴我,他到河南開封,有位上官醫生還未信耶穌；我父母特請他來聽宋尚節講道,他根據聖經講「罪」,那醫生愈聽愈氣,會畢,他對我父親說:「我不再來啦,這人講道由始至終都在講我的事,他為何如此得罪我?家父對他說:「宋尚節」不認識你,也沒人把你的事告訴他.」宋尚節不怕得罪人,只怕得罪神:好像約瑟一樣,他要的是神的喜悅.我想,撒母耳也是如此.
\newline
\newline
第四,宋尚節一切的工作都是透過教會.神使用宋尚節.
\newline
\newline
各位!今天神在香港,尋找合乎祂用的人,要堵破囗.我們的社會和撒母耳時代的社會不一樣,那時是農業社會,情形較簡單:今之社會非常複雜,是科技社會；但是很多方面也有相同之處,社會的腐敗,神百姓的需要,今日亦然.
\newline
\newline
撒母耳的父母都很敬畏神,懂得禱告,許過的願一定要償還.撒母耳無論是他個人的生活,外面的事奉,大大蒙神使用:因為他要討神喜悅,他得著神的話,然後在自己生活中活出來:別人看見他,都喜歡他:因他十分真實.路加形容耶穌,說「祂蒙眾人的喜悅.」撒母耳也是這樣,神和人都喜悅他:因為他的信仰是真實的.
\newline
\newline
弟兄姊妹!當我們面對將來,面對許多的破囗,面對今日的社會,讓我們看清楚,第一步是要撫心自問,問題是否在自己裡面?孔子說:「修身而後家齊,家齊而後國治,國治而後天下平.」天下平是從修身開始.如果我們想,社會那麼多問題,九七那麼多的問號,卻沒想到自己的需要,可能我們將犯上何弗尼,非尼哈的錯誤,那是何等可怕!可是如果我們願意在神光照之下,對付罪,跟從主,忠心把主的話傳開:神會用我們堵住破囗.
\newline
\newline


\newpage

\section{第59屆~港九培靈會~戴紹曾~第七講}
\label{sec:1280}
\textbf{大衛}
\newline
\newline
連結: \href{https://www.hkbibleconference.org/session-message/view/1280}{\texttt{https://www.hkbibleconference.org/session-message/view/1280}}
\newline
\newline
\hyperref[sec:1279]{< < < PREV SERMON < < <}
~
\hyperref[sec:toc]{[返目錄]}
~
\hyperref[sec:1281]{> > > NEXT SERMON > > >}
\newline
\newline
舊約聖經充滿了見證；男人和女人的見證,敬畏神的見證.這些人如何堵破口你
\newline
\newline
當全世界的人對真神失去認識之時,神興起亞伯拉罕,把獨一無二的真神介紹給全世界；這不只是神學方面很大的貢獻,更是敬畏神的人,怎樣體驗神的之真實的見證.亞伯拉罕日常生活中體驗神的同在,神是信實的；亞伯拉罕是信心之父,他認識信實的神,抓住神的信實作為一生的基礎.
\newline
\newline
巴不得我們也以神的信實作為我們生活和事奉的基礎.
\newline
\newline
我們也看見約瑟,摩西,約書亞,迦勒,撒母耳,這些偉人都見證神的能力；並且神使用他們,在相當困難情形之下；別人沒有辦法之時,神興起這些人---堵住破口的人.
\newline
\newline
除了耶穌之外,聖經論生平的篇幅,是以大衛為最多；亞伯拉罕只有十四章,約瑟大約也有十四章；可是敘述大衛的生平,有62章之多.如果今晚要細講大衛的生平,恐怕講至明晨都講不完；所以我們只能走馬看花,看些重要的教訓.聖經真是可愛且寶貴,因為十分真實；描寫一個人,很清楚地把其長處和短處都記載；直截了當把人的軟弱都述說出來,不但報喜同時也要報憂；記載大衛的生平也是如此.
\newline
\newline
聖經說大衛是合乎神心意的人,同時也提及他大大的失敗.(撒下十一-廿四)到底哪個失敗最大,我們不必問；兩個都可怕,一是淫亂的罪,一是驕傲的罪；但是神對驕傲的罪比淫亂的罪刑罰更重；並非說,對淫亂方面可以隨便一點,乃是提醒我們,心裡的驕傲是神所恨惡的.保羅說:「我真苦啊!我發現在我裡面有個律,立志行善由得我,只是行出來由不得我,我一生很高的標準要如何跑,如何作,可是作不出來；我很苦啊!好像在我身上有死的身體.」(羅七)我想,大衛生平時,也看見有此困難；好像有兩個大衛,一是偉大的大衛,一是失敗的大衛.
\newline
\newline
神的憐憫在大委身上很清楚地彰顯出來.大衛是個相當有才幹的人；但是才幹卻帶來困難和危險.
\newline
\newline
如果你是很聰明的青年,應該感謝神,也求主保守你；因為那些才幹,能力,一不小心,就會令我們作更壞的事.
\newline
\newline
以下略題大衛的失敗.
\newline
\newline
聖經並沒說大衛沒有罪,很清楚地把他的責任,他得罪神的事都提出；並且三千年來,全世界的人都可知道大衛的失敗,這是很大的刑罰.但是,感謝主!三千年來,我們也看見神的拯救.聖經告訴我們,大衛失敗終知回頭.「苦海無邊,回頭是岸」,大衛願意回頭.好像彼得回頭,向主認罪.猶大就沒有這樣的行動了.
\newline
\newline
詩篇五一篇,是大衛犯罪之後的禱告.相信許多弟兄姊妹從這章經文大得幫助.大衛說:「我向你犯罪......我是在罪孽裡生的......神啊......求你為我造清潔的心,使我裡面重新有正直的靈.不要丟棄我,使我離開你的面；不要從我收回你的聖靈.」(4,5,10,11)他看見神從掃羅身上收回聖靈,他知道這是何等可怕的事,他說:求你使我仍得救恩之樂；因罪中之樂是暫時的,會過去的,不會滿足我心,我要的是救恩之樂；我要的也不是感覺,我要的是主自己.我犯了罪,主啊,不要離開我.大衛的禱告,認罪,實在很寶貴!神聽了大衛的禱告,赦免他的罪.
\newline
\newline
各位青年!請你不要忽略罪所帶來的後果,有時我們說神是慈愛的,的確這是真實的.神聽我們的禱告；我們認罪,神赦免我們；可是,多少時候也要吃罪所帶來一切的後果.
\newline
\newline
大衛犯罪的後果,女兒被污辱,兒子背叛,他最疼愛的兒子造反,全國也亂了,從撒母耳記下第十一章以後,全國也亂了,由於一人的影響.你別以為自己犯的是小錯,不足輕重；但聖經明明警告我們,罪有它的後果.神是輕慢不得的,人種的甚麼,收的也是甚麼.
\newline
\newline
求主用大衛的見證來提醒,警告我們吧.
\newline
\newline
平常我們講及大衛的罪,我們很少提及拔示巴的罪；其實她也有責任,傍晚時候在平頂洗澡,極不合宜,真是給魔鬼留地步!她是故意或無意；我們不願置評,聖經並不題,至少我們知道她真是給魔鬼留了地步.在座的姊妹!我們應當處處謹慎,無論衣著,動作,都不要給魔鬼留地步.並非說大衛沒錯,大衛大錯特錯；但拔示巴也有她的責任.
\newline
\newline
我們千萬不要玩火,相信大家不會忘記約瑟的見證；他所以能保守純潔,因他不願意玩火；正在有責任之時,更要慎防試探.這是我們曾經題過的見證.
\newline
\newline
大衛堵破口,大衛和歌利亞的故事是我們所熟悉的,今晚不必詳述；可是此見證充滿著寶貴屬靈的教訓.我們從兩方面來看大衛,光看他如何經過考驗,然後看關於他自己的看見和動機,他能力的來源.
\newline
\newline
當掃羅犯罪作惡,(撒上十六)神的靈就離開他；當撒母耳膏大衛,神的靈就降臨在大委身上,大衛是個有聖靈的人,堵破口的人；若沒有聖靈,實在不能堵破口,可是請注意!有聖靈的人,被聖靈充滿的人,還要經過很多考驗；並非你一旦被聖靈充滿,你就長大成熟；還需要聖靈步步帶領前往.
\newline
\newline
大衛經過長時間的考驗.他既是個被聖靈充滿的人,為何神長時間考驗他呢!我認為這是他被裝備的時候；好像約瑟經過很長的時間才能完成神給他的使命.大衛需要對神有更準確的認識；所以神從各方面來考驗他.
\newline
\newline
第一,大衛從他的家人經過考驗,大衛在自己家中有他的工作,也有他業餘的時間.一個人能否運用業餘的時間,從中可斷定他能否成功.孫德生老弟兄曾說:「許多人因誤用悠閒時間,以致失去了神的最好.」我想,大衛能彈琴,他一定用很多時間學習；所以詩篇中充滿了大衛的歌,他邊彈邊唱.大衛要迎接歌利亞的挑戰,他用當時農業社會農夫的武器,機弦是極其簡單的武器,可是他能瞄之極準.我想他一定用了很多工夫鍛鍊.
\newline
\newline
基督徒的時間是神給的,我們當作個忠心的管家；不只錢財,恩賜,才幹；尤其是時間,如果天天虛度光陰,浪費時間；怎能使用你堵破口呢?
\newline
\newline
大衛有音樂的天份,有一天,他被選進皇宮；那時因掃羅違背神,走自己的路；神的靈離開了掃羅.另外一個靈捆綁他,需要有人彈琴；他們就把大衛找來彈琴.大衛工餘回家,幫理家務,他的父親和哥哥都是放羊的,過著農民生活；年青的大衛雖然在皇宮作音樂家,但是他工餘放羊如常.
\newline
\newline
在座間許多青年,有機會受良好教育；是你父母所沒有的機會,你在家裡是否高高在上,有沒輕看你的父母的態度?有沒存著要別人伺候的態度?如果我們要學習裝備以堵破口,我們應當謙卑伺候別人.我想到許多青年的父母還未信耶穌,如果父母發現已信耶穌的兒女輕視他們,不願意謙卑；那怎麼能夠吸引父母信耶穌呢?是否影響父母仍拜偶像,走他們迷信的路?
\newline
\newline
大衛謙卑,當他父親派他到前線去這是很危險的任務；非利士人攻打以色列人,在途中很可能遇敵；但大衛遵從父命.
\newline
\newline
一個堵破口的人,如果不學習聽話,在家中,在教會都不聽話,我們怎能聽神的話呢?
\newline
\newline
大衛到前線去,他三個哥哥考驗他,罵他湊熱鬧,好奇心,不負責任.實在冤枉!
\newline
\newline
多少時候,我們信耶穌,首要的考驗是來自自己的家庭；可能是父母.
\newline
\newline
大衛把許多美味的食物送到哥哥們那裡,他們非但不感謝,還忘恩負義地責備大衛.事實上大衛是奉父親之命；他們罵他不負責任；其實他在出發之前,早已安排把羊群交給別人看管,事先已盡了責任.有時,人對基督徒也是如此；如果人家講的對,我們應該謙卑接受.當大衛的哥哥以利押誣告他時,他閉口不言,轉過身與別人講話.這是多美的見證!他覺得毋須為自己解釋表白,免得浪費時間；前面受敵,有更重要的事要做.
\newline
\newline
數年前,兄弟在台灣中華福音神學院時,有個姊妹遇到很大的考驗；每年招生時,許多想獻身的青年常遇家庭的考驗.開學前夕,有位姊妹來對我說:「明早家父要來領我回家,不准我讀神學.」她告訴我,當她諗完大學,就在自己家的公司作事已有兩年；可是一直感到神揀選她,需要獻身；她父親雖是基督徒,卻加以攔阻.記得那晚我們跪在神面前禱告,這姊妹哭了,說:「主啊!我怎麼辦呢?你明明叫我要孝敬父母；但是我心靈深處知道你揀選我,叫我一生事奉你.」禱告完畢,她回宿舍；次晨,她父親來了.我想,我是外國人,恐怕使他產生另一種感覺,所以請我們的輔導主任陳牧師代我出面,他用了半小時多和這先生談話,越談聲音越大,後來說一定要見院長；原來他是天津人,具有爽直脾氣,他說事情很簡單,若上午十二點之前女兒跟我回家,那就沒事；否則要上法院脫離父女關係.我大吃一驚,他是基督徒,他女兒為了順服主；他居然出此狂言.兄弟不僅一次遇此情形,我覺得很可怕,基督徒的父母竟然攔阻兒女奉獻；我費了多少唇舌都沒有用,最後我對他說:「先生,我應向你道歉,昨晚我才知道你們家裡有此困難；如果早知道,必登門共商解決辦法；現在請你把女兒暫留在這裡,下午三時我護送她回家,我們一起商量,一齊禱告,到時怎樣決定就怎樣作.」感謝主,他答應了.我看見神在行神蹟,當天下午,全校為此事禱告,我和陳牧師到他們家,他的太太話語連珠達半小時,使我們無法講話.女孩子的男朋友,因知道她的父母反對她奉獻；用神的話給她寫信:「愛父母過於愛我的,不配作我的門徒.」他們發現了此信,加罪於他同謀造反；在那一年,這男孩子也進神學院；他們恐怕此對男女不經許可私行結婚.我向他們解釋,我們有校規,在學年中不許結婚.於是他們才鬆了口氣.那天,神真是使用陳牧師,他向這對基督徒父母說:「今年我的女兒也奉獻,她母親在緬甸去世,她住在宿舍,睡在上床,你的女兒可睡在下床；我如何關心我的女兒,同樣關心你的女兒.」陳牧師這句智慧的話,真是把這件困難的事扭轉.他們答應了.不到兩年,我主持婚禮,我充滿喜樂；台前的青年就是那兩位,父母都贊成.今天這倆青年在台北信義會忠心事奉主,作學生工作,大大蒙神使用.回憶當時,是何等的困難!
\newline
\newline
大衛也遇到社會和國王的考驗.當時的婦女唱歌詞引起掃羅嫉妒.她們唱:「掃羅殺死千千,大衛殺死萬萬.」掃羅聽了大不高興.大衛到了前線,又另有考驗,他看見同胞都懼怕,無一人有勇氣,他們都戰戰兢兢面對敵人.掃羅沒有出來領導他的百姓,當他開始作王時,基列雅比人遇亞捫人攻擊,那時掃羅卻站著發抖.堵破口者大衛說:「你們不要害怕,神與我們同在；我去堵這個破口.」掃羅聽見大衛的話,就召見他；大衛站在掃羅面前,掃羅對他說:「你太年青.」這實在不錯；在教會歷史,神常常用年青者作年長者不敢作的事.感謝主!
\newline
\newline
戴德生來華時,他還不到廿歲.我想到蘇恩佩姊姊；她在新加坡,在台灣,在香港；神如何使用她,她在年青之時,短短的時間她願意擺上.
\newline
\newline
掃羅對大衛說:「你太年青,你可以再等；但神要用你.」保羅對提摩太說:「不要讓人小看你年輕,你要作眾人的榜樣.」掃羅又說:「你沒有經驗.」「年青」是對的；但神也要用年青的人；「你沒有經驗.」掃羅這話講錯了.大衛說:「掃羅王啊!並不然,我有很多經驗,我雖然年青,雖然不在戰場,因我不是軍人；但我在日常生活中經歷到神的真實,我是牧羊的,當獅子,猛獸來抓我的小羊；神幫助我勝過,把牠們打死了.」大衛的詩:「耶和華是我的牧者,我必不至缺乏.」他缺乏,神會供應.「我雖然經過死蔭的幽谷,也不怕遭害；因為你與我同在.」大衛知道神與他同在,神會供應他的需要.你怕出來奉獻嗎?你擔心生活怎麼辦?大衛告訴我們,我必不至缺乏；因為創造天地之主供應我的需要,祂是耶和華以勒,祂既能供應以利亞的需要；以利亞的神也是我的神,祂永不改變.大衛真是有經驗,所以他告訴掃羅:「我並非沒有經驗；神在曠野地方,在我進行普通工作時,祂顯出大能；我知道祂是真實的,因我有經驗,所以我有信心,有更大的信心,我可以去.」可憐得很!掃羅常常矛盾.口說:「好啊,願主與你同在.」把自己的盔甲給大衛穿上.但是,大衛說了智慧的話；大衛的經驗是最寶貴的經驗,他沒有靠血肉的經驗,沒有靠人的膀臂,方法.他說:「請你收回吧!」隨即走了出去.
\newline
\newline
大衛面對歌利亞的考驗,大衛看見歌利亞個子高大,還有帶兵器者在他身旁.歌利亞大聲罵陣以色列的軍隊:「你們是掃羅的僕人......(他看不起大衛)......你當我是狗嗎?持著牧羊的杖想把我趕走嗎?你來吧!我把你的肉給飛鳥吃,今天我要打死你.」很大的聲音幾乎地都震動；但是大衛屹立不動面對敵人,毫不懼怕；因主與他同在.
\newline
\newline
一九八○年兄弟有機會回中國大陸,我看見許多我父母的朋友,他們告訴我,在過去的年日,特別在文革期間,最寶貴的經文是詩篇廿三篇第4節大衛說的話.
\newline
\newline
大衛是位堵破口者,他的看見真是屬靈人的看見,他曉得神與他同在一定使他得勝；他的信心看見了那看不見的主.希伯來書十一章告訴我們,大衛和那些領袖得勝是因他們有信心.歌利亞說他們是掃羅的僕人；但是大衛的看見和領受並非如此,大衛說:「你為何怒罵耶和華的軍隊?」這話非常寶貴!「神的國,神的百姓」當時神的百姓實在不像神的百姓,不像軍隊；如果你讀十七11,十三6-7形容這支軍隊戰戰兢兢地上戰場；很多已成逃兵,躲藏在洞穴；餘下的人都非常害怕.為何大衛說是耶和華的軍隊呢?因為大衛所見並非那些發抖的弟兄,他感到神在他們中間,他有屬靈的看見.在神的國,我們應有如此的看見；有時,我們看所屬教會人數很少,覺得領導教會的人沒有才幹,或沒有恩賜；兄姊間沒有了不起的人物,這群人不能作甚麼.你看!耶穌所揀選的十二門徒,有的是打魚的,怎能把福音傳到地極呢?你看見的是人,你有否看見主在他們中間?主說:「父怎樣差遣我,現在我照樣差遣你們去使萬民作我的門徒,我必與你同在,感謝主,天上地下所有的權柄都賜給我了；我與你們同在.」撒迦利亞說:「不是倚靠勢力,不是倚靠才能,乃是倚靠神的靈,方能成事.」感謝主!大衛有屬靈的看見.
\newline
\newline
今天晚上,我真是願意禱告,求主開我們的心眼,叫我們看見,我們這群人是耶和華的軍隊；神若幫助我們,誰能敵擋我們呢?讚美主!我們這群人是屬耶和華的；不是掃羅的僕人,不屬於任何人；神在我們中間,人看不起我們,看我們無能,我們四十晝夜不敢出來,我們沒有力量；可是我們仍然是耶和華的軍隊.感謝主!神要用華人教會,大大使用華人教會,使普天下的人認識神.
\newline
\newline
大衛要作兩個見證,他的動機真寶貝!
\newline
\newline
第一個見證是給全世界的人接受這個功課並學習.這位堵破口的大衛,要全世界的人知道以色列中間有位真神,要高舉真神；有亞伯拉罕的心志；大衛也是如此,不是高抬自己,自己的家,自己的國；他所關心的是神的國,他願意高舉真神.請問:你活著為的是甚麼,目的何在?是為自己而活嗎?是否要別人認識你所認識的神呢?大衛在放羊時,體驗到主的真實,知道主的寶貴；他藉著那個戰場,要全世界人都認識神.大衛是堵破口者.
\newline
\newline
第二個動機,大衛要同胞都知道得勝的秘訣,他自己年青,卻想教導同胞知道得勝非靠武器,乃靠神的大能；大衛懂得國家的歷史,他認識真神,他曉得神與他同在.
\newline
\newline
求主給我們這樣的看見,也給我們這樣的動機.
\newline
\newline
謝謝主!大衛兩千年之後,他的子孫降生在人間,祂名叫耶穌.這位有時軟弱有時剛強的大衛,藉著他的後裔,藉著他的血統；神將祂的兒子降在我們中間.雖然大衛曾經失敗；但他悔改；神要使用他,當時用他,千年之後還繼續用他.我們看見,耶穌的得勝是最大的得勝,因著祂得勝,我們都能得救,我們都能靠主得勝,蒙主的使用.
\newline
\newline
各位弟兄姊妹!我們看見這位堵破口者,能經得起考驗,不懼怕.
\newline
\newline
求主給我們兩個動機:第一,高抬真神,高舉耶穌.第二,讓眾人知道得勝的秘訣不在乎自己,乃在乎主.
\newline
\newline


\newpage

\section{第59屆~港九培靈會~戴紹曾~第八講}
\label{sec:1281}
\textbf{以斯帖}
\newline
\newline
連結: \href{https://www.hkbibleconference.org/session-message/view/1281}{\texttt{https://www.hkbibleconference.org/session-message/view/1281}}
\newline
\newline
\hyperref[sec:1280]{< < < PREV SERMON < < <}
~
\hyperref[sec:toc]{[返目錄]}
~
\hyperref[sec:1282]{> > > NEXT SERMON > > >}
\newline
\newline
在教會歷史中,主大大使用姊妹.真是意想不到；今天全世界的宣教士,大概有三分之二是女的.有人說,因為弟兄不肯奉獻；所以主多多地揀選姊妹,這實在是弟兄們的羞恥!
\newline
\newline
這幾天晚上,我們在思想堵破口者；有男的有女的堵破口者,舊約時代有,新約時代也有,一直到今日,教會歷史都有.
\newline
\newline
「...... 焉知你得了王后位分,不是為現今的機會麼?」(斯四14)這節經文大家都很熟悉.英譯稍為不同,「你得到王國是為了這個時候.」到底是甚麼時候呢?
\newline
\newline
第一,神掌管的時候,這是我們必須明白的,當我們讀以斯帖記,特別是第一章；我們覺得亞哈隨魯王掌管的地方很多,他自以為他是掌管一切的,表面看來也是這樣；他管轄的領土從印度直到非洲,達一百廿七個省,(美國只有五十一省)他請客不是一天一頓的筵席,而是七天的筵席,吃喝都是最豐富的食物；他住的也是豪華的皇宮,每有客到,他親自領客人參觀,並顯示他所有的財富.然而我們在第一章發現亞哈隨魯王,連妻子都不能管,王后不肯聽從他的話；還有,內心惡毒的哈曼,他也不能管；他常翻覆床上不能入睡,他的力量實在有限；至於哈曼,自以為了不起；雖然他在王一人之下,但權威在萬人之上；要誰向他跪拜就跪拜,要誰去就揮之即去,哈曼的確管理很多事情；但是人的力量終於有限.這個時候是真神所掌管的時候,何以見得?
\newline
\newline
其一,以斯帖是個被擄到異國的猶大女子,她有雙重的悲劇,她不但被擄,且是個父母俱亡的孤女.這樣的女子竟然成為王后,實在是神蹟,是神的安排.
\newline
\newline
其二,末底改正當那時有守門工作,同時,他有兩個同事,這兩人恨亞哈隨魯王,想謀害國王；可是陰謀被末底改發覺,就向王報告.這一切都是出於神；如果其他的人報告,情形就完全不同.所以我們看見,神預備各人的情形是很奇妙的.哈曼進行要攻擊末底改,他完全失敗了；事實上,哈曼不是攻擊末底改,乃是攻擊神.神是歷史的主宰.
\newline
\newline
保羅曾說:「時候滿足,神差遣兒子到世上來.」或者我們要問為何等這麼久呢?神把一切事情都安排好了,時候好了,耶穌到世上來.
\newline
\newline
各位!一九六七年中國大陸文革爆發時,我們的神坐在寶座；一九七七年文革過去,我們的神仍然坐在寶座上；一九八七年今天晚上,我們的主坐在寶座上,十年之後的一九九七年,我們的神仍然坐在祂的寶座上.我一直有個禱告,巴不得今天晚上,主給我們以賽亞的異象.以賽亞有一天帶著沉重的心情進入聖殿；因為烏西雅王死了,國家的寶座上沒有王了.以賽亞在聖殿,神讓他看見異象；看見天上的寶座,神坐在寶座上.他看見聖潔,公義,慈愛,全能的神；神與他同在.求主開我們的心眼,看見我們的主坐在寶座上.第二,敵人攻擊的時候,末底改是個好人,他忠心盡責辦理一切的事.他負責保護國王,發現有人要害國王；這位敬畏神的人,立即盡忠職守；他向王報告,救了王的性命.
\newline
\newline
基督徒,敬畏神的人是愛國者,愛護執掌權柄的人.
\newline
\newline
一九八○年兄弟第一次回到離開卅四年之久的中國大陸,到我的出生地.見到我父母許多朋友,和他們有很好的交通.我聽了個見證:在河南一帶有很多人民公社的人信了耶穌,有百分之六十至百分之七十之多；他們說,甚至連那些幹部都不反對；因他們發現基督徒說話誠實可靠.不撒謊,不偷竊,且認真作工,完成國家所定的生產目標.這個好見證,我聽了大受感動.你工作的見證有沒這樣好?如果我們對國家對公司不忠心,偷懶,撒謊,不誠實；未免羞辱我們的神,虧缺神的榮耀.
\newline
\newline
末底改還有另方面的見證(斯三2)亞哈隨魯王下令,所有的人當向哈曼跪拜；但末底改不肯跪拜,因他曉得哈曼不過是人,凡尊敬在上一切保護國王的事他都願意做；但是和自己信仰發生衝突時他不能做.末底改面臨信仰考驗.
\newline
\newline
早期教會,弟兄姊妹遇很大的考驗,他們敬畏神是沒有困難的；一旦開口要為耶穌基督作見證,特別是要見證基督從死裡復活；那些宗教領袖非常反對,不准他們講耶穌的名字,也不准講耶穌復活.彼得和約翰說:「順從神,不順從人,是應該的.」(徒五:29)聖經的教訓,在合理情況之下我們應當順從執政掌權的,服在權柄之下.
\newline
\newline
許多基督徒,該順服的卻不順服；不該順服的反而順服.應該繳稅,沒有按照應繳稅項付稅,這是國民當盡的本分；但是很多基督徒想盡辦法瞞稅.在上的人不准奉耶穌的名傳福音,我們就閉口不傳,那是應該的嗎?
\newline
\newline
末底改愛他的王,保護王的性命；可是當王的命令與自己信仰衝突時,他站穩.
\newline
\newline
我曾閱讀香港基督教週報,有人說香港基督徒不可向同胞傳福音.我不明白,我們的同胞,我們骨肉之親為何不可向他們傳福音?聖經教我們要關心家人.我們應該關心同胞,應該順服神的命令.復活的耶穌對我們說:「你們要去,使萬民作我的門徒.」求主幫助我們!我們在香港許多工作尚未作完,我不否認,在香港還有千千萬萬的同胞未認識耶穌,我們應當盡力進行,向他們傳福音；工廠中的同胞,以及未有機會接受好的教育的人；無論在香港在大陸的同胞,我們都應該向他們傳福音；這是復活之主的命令.求主幫助我們,效法末底改和以斯帖的榜樣；看見敵人的攻擊,末底改站立起來.他的信仰成為哈曼攻擊的重要因素.各位!我們這群基督徒的信仰和其他的人不一樣,信了耶穌,我們成為新造的人；當然,我們不能抱著高高在上的態度.今晨許牧師講道,特別題及約翰十七章耶穌的禱告,主沒有把我們從世界帶出來,把我們留在世人中間,祂保守我們說:「他們不屬世界.」感謝主!這是神在我們身上的作為.
\newline
\newline
末底改和別人不一樣,乃因他有信仰,他敬畏神；怕得罪神,他一心一意要討神的喜悅:所以他的人生觀與眾不同.我們如何?求主給我們以此為榜樣.
\newline
\newline
哈曼向末底改進攻,你以為是攻擊末底改嗎?錯了,他根本是攻擊神.主坐在寶座上,攻擊我們則攻擊神.你應有這樣的眼光.大衛就有這樣的看見,歌利亞攻擊的,不是這群發抖的以色列人；乃是攻擊萬軍之耶和華.
\newline
\newline
哈曼也是攻擊神；如果他的陰謀成功,神的應許會落空.神應許亞伯拉罕的後裔如同天上星星之多；如果哈曼把以色列人都消滅,神的應許在哪裡?猶大人全都滅亡,誰歸回聖地呢?神老早就藉先知預言七十年要歸回.神在舊約聖經中應許彌賽亞降臨；若猶太人滅絕,彌賽亞怎樣來呢?大衛的後裔怎來呢?由此可見,哈曼攻擊的是神,不是人.感謝主,應該給我們一個新的看見；他們不是反對我們,乃是反對全能的神；他們怎能成功呢?主耶穌在世上,也曾遇撒旦攻擊.撒旦極其狡猾,聖經描寫牠的作風,有時如同吼叫的獅子,有時好像光明的天使；在不知不覺之中來引誘我們,向我們挑戰擾亂我們；所有跟隨耶穌的人常會遇到撒旦的攻擊.在教會歷史,我們發現確有此事實.當主耶穌升天以後,一直到大約第四世紀,直到羅馬皇帝君士坦丁接受耶穌之時,基督徒處的是惡劣的環境,大受逼迫.今晨有位姊妹給我看個符號是一條魚的形狀,她問我這是否早期基督徒用來代表信仰的符號?是的,那時因他們受逼迫,要讓其他基督徒知道自己是有信仰的人,希臘文這字有很深的意義,
\newline
\newline
第一個字母代表耶穌,
\newline
\newline
第二個字母代表基督,
\newline
\newline
第三個字母代表神,
\newline
\newline
第四個字母代表兒子,
\newline
\newline
第五個字母代表救主,
\newline
\newline
所以用一條魚代表他們的信仰.當時基督徒受極大的逼迫,就好像希伯來書十一章末幾節經文所形容的.到了第四世紀完全不同了,作基督徒是件時髦的事；因為皇帝已經信主,且是位虔誠的基督徒；所以一下子人人都作了基督徒,沒有逼迫了；但是那時,一個真正的基督徒所受的考驗是很難的；撒旦不再像吼叫的獅子,而是像光明的天使來引誘人,使人放棄信仰,外表仍是基督徒；可是內心已不愛主了,教會要用很多時間討論基本信仰.
\newline
\newline
聖經告訴我們,當撒旦進攻時,我們要抵擋,聖經應許我們,撒旦必離開我們而逃跑.求主給我們這樣的勇氣,靠主的大能大力,當仇敵進攻,我們就要抵擋；用神的話,靠聖靈的能力得勝牠.
\newline
\newline
第三,敬畏神的人採取行動的時候,末底改就是這樣的人,他知道當時的情形,他也認識他的神,他對以斯帖說:「此時你若閉口不言,猶大人必從別處得解脫,蒙拯救......」末底改有信心,他曉得,如果人失敗,神必興起別人來拯救祂的百姓.
\newline
\newline
當大衛開始作王時有很多同胞歸向他,「以薩迦支派,有二百族長,都通達時務,知道以色列人所當行的.」(代上十二32)求主給我們今時代的基督徒,有這樣的聰明,給我們通達時務,了解這個時候；更知道我們應該怎樣採取行動.請看!以斯帖面對這樣的考驗,她怎樣來堵破口呢?我想,她有很多不同的反應,很多的可能；衪許她在宮中過的舒服安全,這些問題不會臨到她,她可以安逸過皇宮的生活,她不理這些；或者當她聽到消息,她心中充滿懼怕,不知所措,她沒有行動；或者她想,不得已時可以一走了之,皇宮中有的是馬和馬車,可是她不走這路.或者她覺得沒有前途,這些人遲早都要滅亡；因哈曼勢力強大,他得勢之日,沒有盼望；吃喝快樂吧!趁機享受吧!不理明天,不理一九九七,反正勢必如此,有機會就過享受的生活吧!但是以斯帖謙卑順服末底改的命令.
\newline
\newline
各位青年!可能屬靈的長輩關心你,我們應該接受他們的輔導；如果自高自大,神不會使用我們,我們絕對不會堵住破口.
\newline
\newline
以斯帖是個懂得禱告能力的人,她送信給末底改,要他禁食禱告,並答應遵命去見亞哈隨魯王.由此,我們看見禱告的重要,她真是明白禱告的功效.聖經說:「義人祈禱所發的力量是大有功效的.」(雅五16)新約時代基督徒禱告,彼得從監牢出來.敬畏神的以斯帖禱告,以色列人得釋放.這是禱告的能力.
\newline
\newline
巴不得我們在未來的數週,用禱告托住包樂佈道大會；魔鬼要攻擊我們的同胞,不要我們的親友認識耶穌,束縛他們在魔鬼的監獄中；可是復活的主要釋放他們.我們當用禱告托住神的僕人,托住一切籌備的工作.
\newline
\newline
以斯帖為了同胞把自己擺上,她抱著犧牲的精神採取行動,不顧自身的安危去進行工作.這是多麼寶貴的榜樣!今天有很多基督徒,只顧自己的性命,享受,前途安全了；但是以斯帖卻不然,她罔顧自己,違例入宮見國王.感謝主給以斯帖有這樣的心,給她的同胞帶來更豐盛的生命.
\newline
\newline
哈曼是個以自我為中心的人,一切為己,並且驕傲,一切計劃只為自己的前途；可是,結果不但他自己死,且連累了他的全家走進死亡的道路.這裡有個強烈的對比,以斯帖不顧自己的性命,為主之故,為同胞之故,願意自己擺上；她得救了,得生命了；同胞得了更豐盛的生命.
\newline
\newline
「現今的機會」與你與我都有關係.神把你安排在如今的崗位上,你要學習為主發光.在學校的,可以影響同學；教書的,可以影響同事認識耶穌,讓他們看見主愛在你身上,而吸引到主跟前.如果你是公務員,是否你在一切工作上有基督徒的見證,秉公行事,在光明中行呢?做生意的弟兄姊妹!你是否在光明中行?還是和外人同流合污?他們有兩套的賬目,我們是否也如此?開工廠,對待工人,是否按公義?有否想如何影響他們信耶穌?有否找出時間讓傳道人,牧師,基督徒,向他們作見證?近日我聽到個很好的見證；內地會在港有從韓國來的宣教士,和一些韓國人在一起,後來有一對夫婦信了耶穌,他們大概是開工廠的,他們發現這位青年宣教士在香港,不是專向韓國人傳福音,而是向中國人傳福音；所以就請他到工廠傳福音.後來開工廠的太太對宣教士的太太提議在工廠開個查經班,專向女士們傳福音；由廠方發給加班工資,讓許多工人有機會聽福音.你是醫生,你是護士,主的愛有沒有從你身上流露出去?你今天的地位,這機會是主給你的.
\newline
\newline
主曾講「好撒瑪利亞人的故事」:有一人從耶路撒冷下耶利哥,途中遭人打得半死,躺在地上；有個祭司經過,有個利末人經過,看見這位有需要的人；他們視而不見,一走了之.我想他們一定有諸多顧慮,若幫助他,對自己有性命之危；而且很忙,無暇顧及,並可能會惹了一些麻煩.祭司和利未人,理應是敬畏神的,卻失去憐憫的心；反而一個撒瑪利亞人看見了,盡量的幫助這個人.
\newline
\newline
各位!現今的機會是主給我們的,祂與我們同在,賜給我們生命的能力；祂給我們恩賜,才幹,祂要使用我們.我們是否願意將一切擺在主的手中?
\newline
\newline
求主用以斯帖的見證感動我們的心,讓我們效法她的榜樣,當年她真是成為堵破口者；今日你是否願意?你能否說:「主啊!你如何用以斯帖為她的同胞,求你今天用我；你如何用我,我不知道；但我願意被你使用.」
\newline
\newline


\newpage

\section{第59屆~港九培靈會~戴紹曾~第九講}
\label{sec:1282}
\textbf{但以理}
\newline
\newline
連結: \href{https://www.hkbibleconference.org/session-message/view/1282}{\texttt{https://www.hkbibleconference.org/session-message/view/1282}}
\newline
\newline
\hyperref[sec:1281]{< < < PREV SERMON < < <}
~
\hyperref[sec:toc]{[返目錄]}
~
\hyperref[sec:1283]{> > > NEXT SERMON > > >}
\newline
\newline
在這次會期間,我們看見很多信仰上的偉人,他們都經歷考驗,有的是很嚴厲的考驗；可是他們終於靠神能堵住破口.
\newline
\newline
亞伯拉罕能幫助全世界的人,知道有位獨一無二的真神.我們看見約瑟真是討神的喜歡.摩西知道神是他的朋友；神的同在使很普通的變成榮耀的；他把手中小小的東西擺上神的手裡,神就大大的使用.約書亞認識神,知道祂是萬軍耶和華；所以他就服在神權柄之下.神興起撒母耳；本來神的話稀少,很少有默示；可是這位願意傾聽神聲音的人,他領受了神的話,然後就傳給百姓.他堵住破口.我們看見大衛,靠著神的名戰勝歌利亞.我們也應學習如何倚靠主得勝的名過得勝的生活.以斯帖願意為同胞冒性命的危險；神用她榮耀自己的名.
\newline
\newline
現在我用數分鐘時間作個見證,我想講一個愛情故事.我的家族本不認識神,兩百年前在衛斯理時代,戴德生的曾祖父才開始認識神,他本是最反對遊行佈道的人；每當衛斯理和同工到達,他很反對.有時他口袋裡帶著蕃茄和臭蛋,進行破壞露天佈道聚會；可是,神的東西飛得更快,神的話打動了他的心.約書亞記廿四15是我們家族八代最喜愛的經文.有一天,當戴雅各準備破壞聚會時,神的話打動他的心.約書亞說:「至於我和我家,我們必定事奉耶和華.」他不願意思想這句話；然而當神的靈把一節經文放在我們心裡時,我們實在無法逃脫；直到他回家都不會忘記.過了段時間,當戴雅各結婚之日,清早他離家到偏僻處安靜下來,那時他未信主；但他想在這重要的日子應該思考,(在座男子請別等結婚之日才思考,這樣太遲了.)這節經文又在他的腦海中,無論如何擺脫不了；於是在聖靈感動之下,他接受耶穌,他跪在主前認罪悔改；之後,他發現婚禮時間已到,急忙回家換衣服即往禮拜堂直奔；朋友見他來到,才敲響鐘聲開始婚禮；其實原定時間已過.婚禮結束,他在茶會中作見證；大家以為他起立講些笑話,料想不到那天他真是作了見證,宣佈遲到的原因,他在田間如何受聖靈感動,悔改,信耶穌.當時眾人之中最覺驚奇的大概是他的新娘；她一定想,難道嫁了個衛斯理的傳道人嗎?但是亦無可奈何；因婚禮已完成,後悔不及.戴雅各帶了新娘回家,就開始禱告；「主啊!我已經對你說,至於我和我家,我們必定事奉耶和華.現在只有我個人,還不算我的家,求主感動新娘的心.」可是他愈禱告,新娘的心愈剛硬；有一天他回到家中,看見新娘坐著悶悶不樂；心想她真是需要耶穌；但經許久禱告,她還未接受；於是他採取實際行動,他把新娘抱到樓上,傾訴心聲；他先跪在床邊,用雙手壓在她身上,讓她不至避開,他開始禱告:「主啊!我已經禱告這麼久,至於我和我家,我們必定事奉耶和華.」他邊禱告邊哭,同時手加壓力在新娘身上；想不到新娘也哭了,她心有所感,哭了一會兒也就開始禱告.那天她也接受了耶穌!我並非介紹新的個人佈道法；如果你們要採取這方法,我不敢負責任.可是有件事我們應注意,兩百年前的一個決定,影響到今天.可能這幾天早堂,晚堂的聚會,你作了重要的決定；如果還有兩百年的歷史,(未知主是否在這兩百年之間回來,也很有可能.)弟兄姊妹作了今天的決定,也會有那樣長遠的影響；不只你一個人,連你的家及其他的人也在內.
\newline
\newline
戴德生來華已是百卅年前的事,他是單身來的,那時他廿歲,有時他非常寂寞,他盼望有個同伴和他一起事奉.他在英國有個女朋友,彼此也通信,女孩子音樂天份很高,他認為她適合在中國傳福音；可是這位小姐卻不以為然.今晨我在生命堂曾提及一事,有一時期,戴德生從上海到汕頭,和威廉芬配搭作內地福音工作.這位蘇格蘭人是名佈道家,天路歷程的中文最早版本是他翻譯的；他倆一齊配搭進入內地,戴德生病倒了,同伴就送他回汕頭轉乘小船回上海；病的相當嚴重,當時這航程需要十天,病勢漸漸恢復,到了上海,不料住所失火,盡都付之一炬.當時還沒有內地會,支持他的委員沒有按時送生活費來,他帶著病且身邊一無所有；忽然收到女朋友從英國寄來的信,非常高興,心想必有佳音,信寫:「德生,你太愛中國,你為何不回來?如果你回英國牧養教會,一樣地事奉主豈不很好嗎,若你回來,我父親一定答應我們的婚事；若你堅持在中國傳福音,我們就此分手.」各位青年!在中國傳福音,和在英國傳福音不一樣.如果神的旨意要你在香港事奉主,你無論如何解釋,你到別處事奉主,是不一樣的.神吩咐約拿往尼尼微城去,他逃避往他施去,這絕不一樣的.但是,各位青年!我們的情感怎麼辦?如何決定呢?我看到戴德生的日記,那時他在上海,當他接英國女友來信後次日早晨,他上街道去,並且買船票；他出去分發口袋裡所有福音單張,勸人信耶穌；他曉得這是神為他安排的工場,他不離開,他寧願和英國的女友分手；因為彼此不能同心,他堅決留在中國；這是主給他的使命.奇妙得很!主是信實的,想不到他在寧波找到配偶,有個孤女名馬利亞,她父母原在馬來西亞一帶向華僑傳福音,也是來自英國,父母已死在工作地區,女兒來到寧波,和另一位宣教士配搭工作.神奇妙地帶給戴德生一個愛中國人的另一半；雖然一起配搭的時日不長；因為馬利亞的壽命很短,可是他倆同心合意,把內地會建立起來,內地會只有五年的歷史；戴德生所愛的馬利亞就回天家了.有的弟兄姊妹看戴德生傳的電影後問我:「你是哪裡來的?」因為影片中好像這家庭沒有孩子；可是我祖父還在,他是馬利亞的長子.
\newline
\newline
今天晚上,我們要講但以理的見證；盼望弟兄姊妹們不是覺得聽了這些故事很有趣,這樣未免沒有意義,我們要看見神如何用你和你的家；看見神的作為才有意義.
\newline
\newline
但以理是神所使用的青年,他的見證,大家都很熟悉,他如何在獅子洞裡；他三個朋友在火中.
\newline
\newline
但以理是個力量持久的人,如同迦勒一樣.他曾經在五個王統治下作顧問,從尼布甲尼撒直到古列之時,他繼續作國家顧問.巴比倫亡國,瑪代興起；有的王發瘋了,有的王逼迫敬畏神的人；在皇宮中也有陰謀,情勢瞬息萬變；可是但以理仍能成為中流砥柱,站穩立場,為神作見證.正如約翰一書二章所論:「這世界,和其上的情慾,都要過去,惟獨遵行神旨意的,是永遠常存.」(約壹二17)但以理就是這樣的人,無論國度興亡,但以理是永遠常存的人.所以王對但以理說:「你素常事奉的神.」連別人都看出但以理是個堅固站穩的人.當但以理在獅子坑中,王在上面對他說:「你素常事奉的神有沒有救你?」但以理說:「願王萬歲,神與我同在.」但以理素常事奉的神保守他.
\newline
\newline
一個堵破口者,必經許多考驗,我們看了這些偉人,他們都經歷嚴格考驗；神叫亞伯拉罕將獨生子獻在祭壇上.神揀選大衛,很多年掃羅追趕他,想殺害他；大衛經歷這樣的考驗.昨晚我們提到以斯帖的見證,她身歷性命之危,這是重大的考驗.
\newline
\newline
讓我們來看但以理的考驗:
\newline
\newline
面對國際局勢的考驗　當時的以色列國被擄,相繼猶大國也被擄,耶路撒冷城被尼布甲尼撒攻破掃蕩,但以理眼見聖城的遭災；巴比倫軍隊進入聖城污穢聖殿,把聖殿的器皿帶走,收入他神的廟裡.(但一2)年青的但以理也被擄；國破家亡,成了無家可歸的人.但以理曾否思想:到底我所敬拜的神是否全能?若然,耶路撒冷為何有此遭遇,聖殿為何有此遭遇,聖殿為何遭敵人污辱,敬畏神者為何也被擄；但以理是否懷疑神的全能?
\newline
\newline
當我們遭遇世界瓦解,四圍搖動,我們有何感想?是否懷疑神的全能?如果神是全能的,也不是慈愛的；才讓這些事臨到我.各位!求主給我們屬靈的眼光.屬靈的人遭遇考驗,實非神不愛我們,神要裝備合乎祂用的人；你若不經嚴格考驗,怎能有豐富的,有果效的事奉.
\newline
\newline
但以理信心的考驗,當我們讀但以理書,我感到很有意思；有的事情他能夠接受,有的事情他站穩拒絕；他身處異國,被擄到巴比倫,環境完全不同；該作的和不該作的是甚麼?憑何作決定?以何為準繩呢?第一和第二章告訴我們,他順著當時的情形而作,巴比倫的教育他可以接受；可能那是種不信神的教育,有很多唯物論思想；但尚可接受,並非一切盡吞下去；國王命令的訓練和教育他接受了,巴比倫的語言,文學他都研究；王給他一個巴比倫的名字,他也接受；並且他在尼布甲尼撒王皇宮作事.他雖是以色列人,他愛自己的國家；可是那時因他服在神審判之下.神藉先知告訴以色列人,因他們犯罪作惡,神的審判要臨到他們.但以理聽先知的話,知道必須70年才能歸回聖地；所以那段期間他服在神權柄之下.但是,另有一些事情,但以理的立場非常清楚,凡牽涉他的信仰,牽涉他與主之間的關係,他的禱告,見證,諸事他有立場.「但以理卻立志,不用王的膳,和王所飲的酒,玷污自己.」(一8)因他知道那些東西曾經祭過廟裡的偶像,所以他概不接受.到了第三章,國王造了個金偶像,命令各人跪拜.但以理的三個朋友受了但以理的影響,他們都立志不玷污自己,不跪拜；他們拒絕,不願意低頭.有一次,大利烏王下令不准向神禱告,只能向王祈求.可是但以理仍在自己家裡「一日三次,雙膝跪在他神的面前,禱告感謝,與素常一樣.」(六10)只要不牽涉信仰,不影響和主之間的關係一定要保守,這是最寶貴的.
\newline
\newline
讀了王明道弟兄的「五十年來」一書,我特別注意抗戰時期王明道弟兄在北京,當日軍進到北京時,他們動員組織基督教團,和日本一樣有個教團；日方要所有牧師和教會的人都參加,他們到王明道處,表明來意；王明道對他們說:「我不參加!和沒有信仰,不相信聖經,不信耶穌是童貞女馬利亞所生的；不信耶穌復活,不信耶穌寶血能洗淨我們的罪的人聯合,我不作；同時帶著政治色彩,我不參加.」過了一週,他們又來了；但王老弟兄已有準備,很客氣地請他們進入禮拜堂的前面；他們嚇了一跳,因為有副棺材放在講台上.他對他們說:「我的棺材已經買了；你們要鎗決我,請便吧!無論如何我不參加教團.」
\newline
\newline
信仰受影響時要站穩,我們與主之間的關係遭受威脅時,我們要珍惜與主之間的關係；今天我們太忽略這一點.我們當如約瑟,他對波提乏的妻子說:「我怎可作這事得罪我的神.」他所珍惜的是與神之間的關係.
\newline
\newline
但以理當然不喜歡得罪人；可是他更怕得罪神；他成功秘訣何在?「如此,這但以理當大利烏王在位的時候,和波斯王古列王在位的時候,大大亨通.」(六28)但以理成功的秘訣:－
\newline
\newline
第一,但以理很清楚在事情開始時擺出立場,不保留,不打折扣,不含糊,清清楚楚把敬畏神的立場擺出來.
\newline
\newline
有許多青年從學校進入社會時,因沒有把基督徒的立場擺出來,結果,當試探引誘臨到,已經失去了立場.步入新場合,寧可把自己的立場擺出,我是基督徒,若要我行詭詐,我不幹；不聖潔的事,我不作.這是基督徒的立場.
\newline
\newline
第二,但以理不怕採取行動,他立刻接受挑戰.
\newline
\newline
第三,但以理找些志同道合的人,這非常重要；當然主與他同在,這是他成功最重要的秘訣；他曉得屬靈的交通,屬靈的團契；人有時遇到考驗,站立不穩；可是兩個弟兄同站在一處,互相幫助,站立得穩,滿有力量.但以理和三個朋友的見證:「但以理回到他的居所,將這事告訴他的同伴哈拿尼雅,米沙利,亞撒利雅,」(二17)他和屬靈的朋友談論.提摩太後書二21-22告訴我們:「人若自潔,脫離卑賤的事,就必作貴重的器皿,成為聖潔,合乎主用,預備行各樣的善事.你要逃避少年的私慾,同那清心禱告主的人追求公義,信德,仁愛,和平.」同那清心禱告主的人一同追求,行神喜悅的路.但以理懂成功的秘訣,明白團契的重要.
\newline
\newline
我想到衛斯理在牛津大學讀書時,他曉得許多教授不認識神,他就組織神聖會,同學們和他一起讀經禱告.
\newline
\newline
第四,但以理也知道禱告的重要,「我便禁食,披麻蒙灰,定意向主神祈禱懇求.」(九3)但以理在禱告中追求神的啟示.天使長加百列對但以理說,本來廿一天之前你開始禱告時我要來；可是波斯的魔君攔阻去路,等到米迦勒來時,去路才通.
\newline
\newline
禱告是我們屬靈的武器,因時間關係,我們不能詳述但以理的禱告.請各位特別注意,當你有需要時,將你的需要；無論是物質的,知識的,力量的,保守的各方面的需要都擺在主前.但以理向神禱告,當神答應他賜給他時,他懂得感謝.
\newline
\newline
但以理的禱告是感恩的禱告(二20-23)他在日常生活中禱告.兩年前,我聽印度著名解經家威廉士講道,這位弟兄實在蒙主使用,他在內地會總會帶領靈修時間,講及使徒行傳教會的禱告有三種情形:第一種禱告是例行的禱告.但以理有此禱告,他在日常生活中有固定時間在主前禱告,第二種,在壓力之下的禱告.第三種,優先恕罪的禱告.但以理更會為他的同胞禱告.整個第九章是但以理的禱告,他曉得神已應許過了七十年他的同胞可以從被擄之地歸回耶路撒冷聖地,他曉得歸回之日不能帶罪回去,他就為同胞禱告.但以理的禱告,好像蘇格蘭洛克斯的禱告一樣:「求主將蘇格蘭賜給我；否則就任憑我死吧!」他盼望他的同胞早日歸向耶穌.但以理也有此禱告.
\newline
\newline
希望各位回家之後讀但以理第九章,再三思想他的禱告生活.
\newline
\newline
謝謝主的恩典,我覺得我們在外面的弟兄姊妹,可以有很多的學習；有時我們太驕傲,想早日去告訴國內同胞如何傳福音.事實上,我們應當在他們腳前學習,如何在苦難中為主而活.當兄弟首次回到河南省的開封,已是離了老家卅四年,我碰到一個同伴,聽見他的見證,大受感動.有位老弟兄,原是內地會醫院的總務,他說,有位姊妹,是我父母從前的學生,在文革期間,她安慰了很多人,她按家逐戶探訪幫助人.有個幹部的妻子發瘋,服中藥,西藥都無效,十八年一直被關鎖在小房間內；如果放她出來,就到處亂嚷,衣服不整,披頭散髮；過的真是非人生活!有日,瘋女人忽然哭了,說:「耶穌能救我嗎?」幹部莫名其妙,不知耶穌是誰,有人告訴他,耶穌就是基督徒所相信的那位.幹部隨即把妻子帶到那位姊妹的家,說明一切情形,並對姊妹說:「你是信耶穌的,你要救她.」姊妹對幹部說:「把你的妻子留在我這裡,你先回去；讓我來負責.」她立刻為病人洗澡更衣,給她吃東西,愛她,為她禁食禱告.不到一個月,幹部的妻子清醒過來,完全復原；不但精神,身體復原,並且接受耶穌基督作她個人的救主,成為新造的人；她回家後,首先帶領丈夫信主.感謝主的恩典,真是禱告的能力!
\newline
\newline
親愛的弟兄姊妹!但以理是懂得禱告的人,這位站在破口者,相信神,知道神是永恆的神,祂的國度永遠常存；但以理親眼看見巴比倫過去,瑪代過去；一個個的王過去,可是他知道神的國永遠不會過去,他是永遠國度的國民,他認識國王.這位堵破口者,帶給神的百姓很大的復興,當他們歸回聖地；因有這中流砥柱,很多人認識真神,他能夠站穩,事奉他的神.
\newline
\newline
你要不要?願意不願意?把你自己獻給主.
\newline
\newline


\newpage

\section{第59屆~港九培靈會~戴紹曾~第十講}
\label{sec:1283}
\textbf{尼希米}
\newline
\newline
連結: \href{https://www.hkbibleconference.org/session-message/view/1283}{\texttt{https://www.hkbibleconference.org/session-message/view/1283}}
\newline
\newline
\hyperref[sec:1282]{< < < PREV SERMON < < <}
~
\hyperref[sec:toc]{[返目錄]}
~
\hyperref[sec:1388]{> > > NEXT SERMON > > >}
\newline
\newline
我應該謝謝大家,在這十天當中,我相信講者比聽者得更多的福氣,這是神的恩典.也謝謝陳牧師和我一起配搭.相信許多弟兄姊妹用禱告來托住這個講台,這是禱告的能力.
\newline
\newline
培靈大會實在有神的帶領.本屆的題目都是針對著神的國.這個永恆的國,不改變的國,榮耀的國度；誰才能進去呢?就是門徒.
\newline
\newline
我們也看見神的國常有破口；因為人遠離神.這幾天我們一直在看舊約聖經幾位堵破口者.在一年前,我所思想的題目不是「堵破口者」我本想講主禱文,每天晚上講一句；但是禱告中感到不通,最後才決定這十天所講的題目.我覺得勝靈的帶領實在很奇妙!
\newline
\newline
未講尼希米見證之前,我先講些關於我家的見證:家父生於蘇格蘭,他是孿生子,身體非常軟弱,醫生告訴我的祖母,這兩個孩子活不到六小時,過了六小時,醫生說活不到六天；過了一個星期,醫生說活不到六個月；可是當六個月時,我的祖父母已經把他們帶到中國.家父活到84歲才回天家,他也是在內地會學校讀書,畢業後離開山東煙台到上海.他是個藥劑師的學徒；他發現老板的女兒很可愛,就開始和她談戀愛；但老板很不高興,因覺得戴某沒有甚麼出色,所以禁止女兒和他來往；可是他倆仍很相愛,常密約外出相會；正在那時,神帶領一位宣教士由日本來滬舉行佈道大會；家父赴會,神的話打動他的心.羅馬書一章,保羅提及許多的罪說:「世人都犯了罪.」這章長論罪的經文,其中有一罪「不孝敬父母.」聖靈在家父心中說話:「你是孝敬父母的兒子,很聽從父母的話；可是現在你常帶那女孩子出去；明知她被父親禁止和你在一起,雖然你們沒行差踏錯；但原則上你叫她違背父母的話,叫她不孝敬父母.」家父受聖靈責備,當晚作了決定,從此不再見這女子,除非她的父親許可.同時,家父真正得救了,已往他雖生長在宣教士的家；但他不認識主；佈道會那晚,耶穌基督成為他的救主；他的一生改變了.實在很奇妙!有個時期他在河南開封內地醫院工作,後來到美國進修大學.神奇妙的帶領,他認識他祖父戴德生的神也是他的神.他離開上海時,他的旅費不足,只夠由上海搭船到西雅圖；如要到達讀書地方,費用沒有；當他到西雅圖數日,都不向任何人啟齒,他只對主說:「你知道我的需要,求你為我預備.」感謝主!當他打算去買火車票時,有人把信封交在他手中；神的供應極其奇妙!信封中的錢正好購買火車票.亞伯拉罕的神是我們的神,以利亞的神,今日也是我們的神.家父進了大學,在那裡認識了家母,他們唸完大學便結婚,之後,來中國傳福音,在河南,陝西,廿年之久傳講主的福音.一九四六年我父家要離華赴美,我們四個孩子已經出了集中營；和父母分離了五年半才團聚,就在那段期間我認識了主耶穌,在皖北地方,我接受耶穌作我個人救主.到美國,本擬在紐約教會學校讀書；可是當我們到了美國,從西海岸要往東走之時,忽接消息說紐約州沒有房子住,學校卻是沒有問題；我們分開五年半剛相聚,父母不想把我們安置在校舍；所以臨時改變計劃,我們沒有去紐約,停留在麥芝根,那裡有座房子,也有教會學校；我父母深感這是神為我們預備的.我在麥芝根讀大學二年級時,大學一年級有位小姐,她生長在麥芝根,她小的時候,有一主日上午牧師講道時,給青年有機會獻身,這女孩到前面,當她禱告,心裡聽見主問她說:「你真是全心一意把自己奉獻給我嗎?若然,你是否願意作個宣教士?」她心裡回應主說:「主啊,若你許可,我願意作個宣教士.」這件事就暫時擱置下來.當她唸高中時,有個宣教士到她們的禮拜堂,是剛從中國來的,報告在華傳福音的情形,再次給弟兄姊妹和青年獻身機會,那晚,這女孩子深感神要她在華人中傳福音,她獻身.想不到她開始在台灣傳福音,她的姓和那位宣教士一樣,因為她嫁給那宣教士的兒子,她就是我的母親.當內子萊茵讀大學時,我們彼此不認識；主步步帶領我們,當時完全不知道她有心在華人中傳福音；事實上,我一直拒絕在華人中傳福音；因為我心想一定許多人認為我是盲從,曾祖父,祖父,父親,都在華人中傳福音；所以我心中產生一種強烈的抵抗；我想讀法律.但是,神真奇妙,當我在光明中行時,帶領我認識萊茵,然後我才知道神在她身上的帶領；那時神也讓我清楚知道應該在華人中傳福音,不是因著祖先的關係,乃是主自己此時要我在華人中傳福音.謝謝主的恩典,神所配合的事,實在奇妙!我們看見亞伯拉罕,神為他的兒子以撒所預備的；我想在我一生事奉之中,神有如此奇妙的預備.今天當我這見證要結束時,我願意補充幾句話:常有人對我說,你們戴家對華人教會的貢獻如何如何.......我覺得不應該這樣講；如果我們是順服神的人,一切是神的恩典；並且我覺得應該反過來講,以我戴紹曾而言,我從華人教會所得的祝福,遠超我給華人的恩典,這話是由我心深處講的.我的兒子繼宗,十一年前來香港,時在第一屆華福會.他生長於台灣,讀的是中文學校；雖然他長大在宣教士之家,卻不認識主.有一天,林治平弟兄來找我,他是宇宙光主編,他說藝術團準備在華福會中演「庚子年」一劇,問我可否幫忙參加演出,劇中需要一位宣教士,請我扮演.我不反對；但我在中華福音神學院實在很忙,我提議請我的兒子代替,他個子比我高比我壯；林弟兄連忙說好.林弟兄邀請他,他立刻答應.數星期他們一直在排練；聖靈開始在他心裡作工,一九七六年正式在此講台上,和其他演員一同演出,其中對白:「義和團的人要來,需作個很重要的決定,到底要不要逃避,要不要離開中國?......如果在這裡的中國基督徒死,我和他們一齊死,我不走.」當時他講這話,並非口頭講而已,他心裡對主說:「主啊!我願意一生奉獻給你.」所以我說,我們家從中國教會所領受的恩典實在很豐富.那年繼宗回到學校時,有人問我:「戴繼宗為甚麼變了,和過去不同了.」我說:「是的,他重生了,變成一個新造的人.」謝謝主的恩典,他今年五月底從神學院畢業,六月底在美國東北波士頓華人教會,和李壽全牧師配搭,學習事奉主,過兩年,他盼望在台灣和香港華人之中傳福音.神的恩典很奇妙!約書亞說:「至於我和我家,我們必定事奉耶和華.」
\newline
\newline
求主給在座每位青年這樣的心志.
\newline
\newline
今天晚上,我們不可能講完尼希米的見證；因為只有幾十分鐘的時間.
\newline
\newline
尼希米的確是位堵住破口的人,請各位特別注意!工作的開始是從心裡的負擔開始.
\newline
\newline
第一,尼希米是位關心的人,這位原來在非常安全的地方,工作單位相當重要,地位很高的尼希米；他和但以理一樣,在巴比倫一帶工作.四圍大多數人非猶大人,他所服事的王也非猶大人；可是有一天他看見自己的同胞,便問:「耶路撒冷的光景如何?」由此可見,他的心絕非忘本的人；雖然他處於安全之地,他想知道耶路撒冷的光景；神百姓的情形.當他聽到報告,心裡非常難過,他願接受報告；他雖然明知所報告的必是耶路撒冷城極壞的光景.
\newline
\newline
我不知道今天一般的基督徒,是否更清楚更深刻地了解世界的需要,許多未聽過福音者,走滅亡道路者的光景,他們在何處,我們是否了解?有時我們從不過問,似乎與自己無關.
\newline
\newline
第二,尼希米以禱告為工作的開始,當他聽到有關於耶路撒冷的光景,他並非一時之感動而已；他四個月之久禱告且禁食,因所聽的報告大大感動他的心.尼希米的禱告,如何地崇拜,如何地認罪,為同胞代禱,向主祈求幫助,求主在他採取行動時,把當說的話放在他口裡,引領他的腳步,走主要他走的路.
\newline
\newline
尼希米的禱告作我們的榜樣,盼望大家讀尼希米記,有關四個月之久的長篇禱告.
\newline
\newline
神的工作,堵住破口的工作需要禱告；沒有禱告的基礎,工作不會成功；得不到神的啟示和帶領,工作由始至終是靠禱告.
\newline
\newline
第三,工作開始時尼希米採取行動,當他聽到報告,可能有很多的思想,很多不同的反應；或者他想,禱告了四個月,已盡了本份,貢獻不少了,可以卸任了；也許他想,這事與我無關,我在皇宮是王的主將,擔任重要的工作.因為當時的王害怕人以毒藥謀害,所以擔任酒政是保護王性命的人,有的當酒政者甚至是宰相.尼希米向王請假回耶路撒冷,他毫不顧慮,王是否准許,是否因此大發雷霆,以後回來是否仍保留工作地位?王發現尼希米心中有事,問他要求甚麼?尼希米立即回答:「王啊!求你差遣打發我回去.」尼希米要親自回祖國,他不以金錢派別人代替；雖然重建城牆需要鉅款.我發現有不少青年以錢換生命；神所要的是你的生命,你的一生；可是我們對主說:「別人去,我可以奉獻金錢來支持.」神要的是尼希米自己；他說:「主,我在這裡.」他對王說:「求你打發我去.」於是工作就開始了.
\newline
\newline
我們看見堵破口工作進行之時,有兩件事情大家注意.有很多的禱告,我已說過尼希米的禱告；但這裡有兩件事,一是尼希米的計劃和組織.二是敵人的反對.在教會,有時一些基督徒錯誤的觀念,以為有計劃有組織就不屬靈；這是錯的.如果我們的計劃和組織代替聖靈的工作,這工作必不成功；若在聖靈引導之下冷靜思想,訂計劃,工作有組織,有條有理；因我們的神是有計劃,有組織,有條理的神,我們看見,神創造天地的次序,建造會幕的計劃,聖殿工作的進行,以色列支派的排列；就知道神是有計劃的神.為何我們覺得,根據臨時的感動比較屬靈,而冷靜下來,祈求聖靈帶領我們的思想,為何不屬靈呢?
\newline
\newline
尼希米是個有計劃的人,在他四個月禱告之時,他一面禱告,一面思想.何以見得?因為王問他要離開多久?需要甚麼?他一口回答,可見他曾經思想過此些問題,早有準備.一個堵破口者,必須有好計劃；並且工作必須有好的組織,計劃和組織都由禱告中聖靈的感動而產生.他向王要詔書,如果到了耶路撒冷才想要詔書便來不及了,必被敵人攔阻他的工作.王給他詔書,他所要走的路就通行無阻了,建造城牆所需材料也都備齊了.第二章講到尼希米半夜出去巡視,雖然很危險；但他覺得必須親身了解耶路撒冷的需要.一個堵破口者,為了工作需要,應親自有認識,然後才能向同胞講話；次日他招聚以色列的領袖,向他們說話,激勵他們的心,動員他的同胞；若不事先觀察,實在作不到.他把黑暗的光景告訴同胞,他不願敷衍地說不成問題；因他曉得共處於極度惡劣環境之下,所以就開門見山地對他們說:耶路撒冷的光景如何如何,然後他開始作見證.在第三章,兩次有句同樣的話:「我神施恩的手怎樣幫助我.」(三8-18)他提及神的手在他身上,他為何要離開書珊的王宮,他告訴同胞,這異象是出於神的.
\newline
\newline
各位青年!如果神帶領你一生事奉祂；相信許多弟兄姊妹已有這樣的帶領.當困難來到,能夠藉著神揀選你之時,見證神如何將異象給你,你在一切逼迫困難之下可以站穩；因為你知道是遵命而來的,是出於神,是神的手帶領.
\newline
\newline
尼希米身居高職安閒的位置,居然回耶路撒冷,絕非出於自己.以斯拉也是用同樣的這句話:「神的恩手在我身上幫助了我.」這句話決定了三件事情,第一是神的供應,當尼希米和以斯拉從遠方到耶路撒冷之時；神用奇妙的方法把金銀供應他們.神的供應,神的恩手!
\newline
\newline
兄弟站在這裡,真是願意見證神是信實的.內地會一百廿多年的歷史,每個月神供應祂僕人生活的需要,今天全球內地會同工有千人之多,是不向別人開口,這是神知道的,神感動弟兄姊妹的心,實在很奇妙,我們數月來新加坡總部重建辦事處,兩年前當這異象來時,我們覺得工作擴大了,同工增多,特別是亞洲的同工；需要擴建辦事處；但當時一個錢也沒有,工程進行時,每月主的供應都超過我們所當付出的,到了工程的尾聲,距完工前三個月,我們又看見主的供應.
\newline
\newline
我想到新加坡有個華人女孩名蓮達,她有很美的見證.她讀神學期間,有天從家裡出來,鄰居的馬來人太太對蓮達說:「最近我觀察你的生活和一般人不同；我們的女兒可以和你作朋友嗎?」蓮達非常高興；因為她正盼望帶領她信耶穌,如今終於實現了,馬來女孩歸入了主的名下.後來,有一次她到蓮達家裡,看見蓮達的母親獨坐,似乎非常難過,就問她是否因不久蓮達將到外地當宣教士,無法按月供應生活而擔憂?說:「是的.」馬來女孩對她說:「蓮達走了以後,以往她供給你多少,讓我來負責.」這是神的供應.神是信實的,神要你作個堵破口者,難道你看不見祂是信實的嗎?祂施恩的手在你身上時,祂供應你一切的需要,在人才方面也供應,在路途中也保守你.以斯拉和尼希米都有此經驗,神的恩典在他們身上.
\newline
\newline
尼希米進行工作之時,敵人來攻擊,內憂外患.這幾天,我們看見神所揀選的人,當他們堵破口時,都處於風雨中,在相當考驗之下.尼希米也是如此,人嗤笑他,威脅他,甚至採取行動和他打架；可是尼希米以禱告勝過他們.可能你讀第三章覺得索然無味,諸多名字,諸多的門,諸多專用名詞,毫無意思!我倒覺得很有意思；因看見尼希米的計劃,看見神兒女的同心.
\newline
\newline
教會中,常常個人主義太強,基督徒領袖各自為政；好像許牧師所題的,有時宗派主義太濃厚,超過一切.我們沒有想到神的國,我們所見的是自己的範圍,沒看見我們所建造的是神的國,我們所堵的破口是神國的工作.孫中山先生說中國像一盤散沙.多少時候教會豈不如此,教會中信徒之間不和睦,領袖之間互相輕看,教會團體彼此不和,彼此競爭.
\newline
\newline
尼希米帶領同胞,同心合意重建聖城,把破口都堵住了.這是神的工作.
\newline
\newline
尼希米帶領同胞在屬靈的奮興中,城牆完成,尼希米帶領同胞讀神的話,以斯拉講解聖經.到了第九章,全以色列民伏在神的面前,認罪悔改；神大大作工,外面工作已經完成,內在的事正要修理.到了第十三章,尼希米看見社會不對的地方,看見神的百姓沒有分別為聖,聖殿中有不潔之物,許多人沒有奉獻十分之一,導致神的工作不能順利進行,許多神的百姓忘記安息日,甚至信與不信者通婚.尼希米看見有聖殿,有城牆；可是他知道更重要的,是神的兒女的心,要完全向著主,完全歸主.過去,神的忿怒臨到他們身上,如今歸回,重建聖城,聖殿,尼希米帶領他們徹底悔改.我們看見神賜福,看見工作的發展,看見福音傳開.尼希米是個堵破口者.
\newline
\newline
今天你看見教會的光景,看見許多的破口,你也知道自己屬靈的光景,你願不願作個堵破口者?我不曉得我們還有多久,有沒有十年,有沒有更長時間?
\newline
\newline
在新加坡有個弟兄給我作個小見證.兩年前薛醫生到中國西安參觀.那天是禮拜六,他到秦始皇陵之後,他問嚮導,西安有沒有禮拜堂,他想去禮拜.那人說該地似乎沒有基督教會.於是薛醫生心有所感到那邊開個禮拜堂,每主日有聚會；當晚他上床睡覺,心裡不安,他想到西安是大都市,有很多同胞卻沒有教會；他在睡夢中看見一個穿短褲的中國男孩子走過來,就問他是誰叫甚麼名字?答說:「我是耶穌.」薛弟兄一驚,耶穌問他:「你要我何時回來?」他馬上說:「不要現在回來,我們還未預備好,還有,你所交給我們的工作還未作完.」耶穌說:「好吧!我再給你一代,我必與你同在.」我不知道薛弟兄的夢是否帶著預言的意思.主對我們說:「還有一代.」
\newline
\newline
今天在華人之中和華人之外,還有千千萬萬人未認識耶穌；在我們的地位,我們的工作,我們是否不關心,我們是否願意像尼希米說:「主啊!請差遣我.」
\newline
\newline


\newpage

\section{第59屆~港九培靈會~楊濬哲~序}
\label{sec:1388}
\textbf{序}
\newline
\newline
連結: \href{https://www.hkbibleconference.org/session-message/view/1388}{\texttt{https://www.hkbibleconference.org/session-message/view/1388}}
\newline
\newline
\hyperref[sec:1283]{< < < PREV SERMON < < <}
~
\hyperref[sec:toc]{[返目錄]}
~
\hyperref[sec:1236]{> > > NEXT SERMON > > >}
\newline
\newline
港九培靈研經會,創辦於一九二八年,迄一九八八年,適逢是六十週年了.六十年,已超過了半個世紀,在這萬變的世代中,永不改變的惟有是主耶穌基督.
\newline
\newline
港九培靈會是個老人了,飽歷滄桑,挫折,戰爭,窮乏,搬遷.時至今日,可以說仍然是老當益壯,老而不凋,老練成熟!惟一可誇耀的,只有主的恩典；我們可引證(約一:16)的話:「從祂豐滿的恩典裡,我們都領受了,而且恩上加恩.」前人得復興,留下來的佳果,一直延伸到今天.在這幾十年來,神多方,多次的,藉著祂的僕人和使女,傳遞信息,發出呼聲；或為悔改,或為復興,或為奉獻.使數以千計的基督徒同蒙恩福,奠下了信仰穩固的基礎.教會亦從此日見興旺,增長.主恩實在數算不盡!
\newline
\newline
本講道集係第五十九屆大會講道記錄,三位講員所傳達的,都是時代的信息；經輯錄成集,以保留文獻.深盼讀者能珍惜分享.
\newline
\newline
今日教會倡導門徒訓練風氣甚盛,蒙許道良牧師,從約翰福音中,以神學觀點,介紹出門徒之模式,並闡述門徒的含義.金新宇牧師移民加拿大已數年,在教會忠心事主.他的信息是傳說神國的榮耀,給我們國度有更清楚的認識,激勵信徒必須先求神的國,和愛慕神的國.戴紹曾牧師列舉出聖經重要的人物,如亞伯拉罕,約瑟,摩西,約書亞,迦勒,撒母耳,大衛等,如何為神國防堵破口.讀者若細閱本書,重溫信息,靈命必能獲得造就.
\newline
\newline
本講道集蒙游偉業弟兄,林小娟姊妹,雷麗端姊妹記錄,朱建磯牧師編輯；游偉業,林小娟二位兄姊協助校對.謹此致謝!
\newline
\newline
六十屆大會為表示向神感恩起見,每晚均有特備音樂節目歌頌.首末兩晚蒙城浸及欣樂聯合詩班合唱.黃永熙博士,張美萍女士分別指揮.其餘獨唱或合唱者,有李嘉琦先生,陳嘉璐女醫生,費明儀女士,招梁碧冕女士,何金秀莉女士,趙麗儀女士,楊淑貞小姐,陳之霞女士,張蓮女士等.願他們的歌聲,成為嘴唇的果子,佳美的祭品,蒙神悅納!
\newline
\newline
大會蒙香港聖經公會,特為六十週年記念,印製一萬本聖經,以五折廉價發售；藉以鼓勵信徒愛慕聖經,研讀聖經及篤信聖經.委辦對聖經公會此舉,表示深切謝意.
\newline
\newline
本會多年來擬向當局申請註冊有限公司,在楊世傑律師協助下,本年已順利完成；深感主恩浩翰!本會並已獲稅務局批准,超過百元之收據,可作減低稅額之用.
\newline
\newline
六十屆大會籌備倉卒,蒙九龍城浸信會借出場地,同工,義工,助道會,各兄姊協助,招待,維持秩序.福音傳播中心同工,協助寄發宣傳品,及現場錄音,錄影.基督教週報,時代論壇報,刊豋新聞稿及廣告,廣事宣傳,以收效益.統此誌謝!
\newline
\newline


\newpage

\section{第59屆~港九培靈會~許道良~第一講}
\label{sec:1236}
\textbf{似是綑綁}
\newline
\newline
連結: \href{https://www.hkbibleconference.org/session-message/view/1236}{\texttt{https://www.hkbibleconference.org/session-message/view/1236}}
\newline
\newline
\hyperref[sec:1388]{< < < PREV SERMON < < <}
~
\hyperref[sec:toc]{[返目錄]}
~
\hyperref[sec:1237]{> > > NEXT SERMON > > >}
\newline
\newline
約翰福音成書的目的,是要帶領讀者認識道成肉身的基督,相信接受祂而得永生.這種相信並非迷信,乃是基於証據,經過了解而作的決定.同時,非只口中的認信,而是不斷更新持守的決定.約翰看「跟從主作門徒」並非兩種不同的呼召,而是程度上的分別而已.
\newline
\newline
作門徒是要重活基督的生命,這是一種蓄意的選擇──過著符合真理的生活方式,有模倣基督的心態.希望透過這八篇信息,編織一幅生命的圖畫,使我們對門徒的觀念有更清楚的領會.
\newline
\newline
符類福音常將門徒提昇至與眾不同的地位；這在約翰福音則有所不同,門徒只是眾多信主者中被篩剩的,是一群願意跟從主到底的人.約翰重視人的選擇,是以較少寫主呼召門徒的情況.此外,約翰另一個特點,就是描寫神的榮耀.四福音中,惟有約翰沒有記載登山變像的事蹟；但卻在頁首提到「見過主的榮光」.很多解經家認為這是指登山變像,我卻認為約翰是別有用心,他描寫神榮耀的彰顯.但並非如舊約那種光輝燦爛,使人望而卻步的；神的榮耀乃在曠野,病痛,貧乏及苦難中彰顯的,最終更在各各他山上彰顯.換言之,神的榮耀是透過一個完全順服的生命之一舉一動中,得以彰顯.當神的旨意能毫無攔阻地行使出來的時候,便是神榮耀的彰顯.這似是一個矛盾──在人的卑微,羞辱中見神的榮耀；這卻是道成肉身要表達的信息.耶穌要門徒繼續這種工作,以卑微順服,處理自我,使人得見神的榮耀.
\newline
\newline
約翰福音是受著猶太文化的影響,猶太文學的表達有如樂章中,主題可以在不同地方出現,但能清楚帶出信息；甚至不同的主題互相交錯,為使形象更清晰.所以我這次不會注重經節的講解,而會將相同的觀念合併,以觀察清楚的圖像.
\newline
\newline
今天我們要思想一個題目──似是綑綁
\newline
\newline
未信的人常認為基督徒是最沒有自由的人,既不能飲酒跳舞,也不能賭馬,滿身鎖鏈.故此一直不嚮往基督徒的生活.門徒的意思,是享有最少自由幅度的人.有時候,我們舉頭望向外面的世界,也會埋怨自身太多束縛.
\newline
\newline
家裡有一頭狗,牠不愛受束縛,開著大門時,牠自由進出,非常快樂；可是,門關起來便不得了,牠在屋內徘徊著,很不自在.事實上,開門或關門牠也在同一間屋子,只是,門關起來牠便有心理威脅.然而,沒有原則的自由,最終仍帶來束縛.就如船在海中航行,若失了舵,不但沒有自由,反而受另一種壓力所約束.風浪潮水成為其新的鎖鏈,只有當舵能正當運行,它便能超脫束縛.
\newline
\newline
跟從主的人,是失去自由,抑或是真正享受自由的人呢?耶穌希望我們的「信」進入更高的領域,能真正結出永恆的果子.「相信」不過是一個人生方向的開始,嚴格來說算不得甚度,除非跟著有生命的委身.故此,約翰提到「信的人多,跟從的人少」.一個自稱為信徒的人,實在需要委身來印証.
\newline
\newline
經文中,耶穌輿猶太人談及自由的問題.然而,他們卻以為這信息應對奴僕說,不明白原來自己正在綑縛中.耶穌暗示他們是失去自由的人,雖然相信上帝,卻與祂有距離.可見,第一種束縛人的,就是無知,無知的人難以被說服.有曰:「秀才遇著兵,有理說不清.」秀才雖比大兵有知識,卻因為大兵的無知而不能說服他.人不明白自己的無知,實在是可悲的現象.肴臘智者蘇格拉底,認為自己所以勝過別人,是因為他比別人更明白自己的無知.
\newline
\newline
最近,播道神學院來了一批短期宣教士,他們對香港的環境很陌生.有一天,一位姊妹丟了錢包,由於裡面有我的姓名及電話號碼:警局便來電詢問我,原來有人拾了錢包送到警局.於是,我便問那位姊妹是否掉了錢包:她感到很可笑,認為自己沒有可能掉了錢包的；後來她到警局,才嚇然發現那是自己的錢包.還好,裡面的証件及現金都沒有失掉.可見,很多時侯我們自己無知迷失,還不自知:很奇怪,前幾天這位姊妹又發生同樣的事,這次她便明白自己又丟了錢包.
\newline
\newline
猶太人認為他們是亞伯拉罕的後裔,從來未曾作別人的奴僕.事實上,倘若他們深思之後,便不會這樣說.當我們回顧以色列人的歷史時便發現他們曾在埃及,巴比倫,及亞述等國為奴:即使在當時,也是受羅馬的管治.以色列的國家特徵是「亡國奴」.只是他們卻不肯面對現實,恃看自己是亞伯拉罕的後裔,神應許之民:不承認受政治的壓迫,永久為奴的光景:只盼望有一天彌賽亞來到,他們可得到永遠的釋放.他們故意抗拒事實的真相,陶醉於自己的理想:人往往是這樣,是故意喜歡掩飾自己的無知的.
\newline
\newline
我認識一位美國籍的姊妹,她在美國教人煮中國菜.後來到香港旅遊,一直不敢上中國菜館:害怕所吃的中國菜與自己所教的不同,故此,她來港有好幾天,都沒有吃過中國菜.無知使人害怕面對真相,恐怕它對自己產生威脅,然而,真理始終向愚昧無知挑戰.「無知」是自我囚禁,自我封鎖:就如一個瞎眼的人,受自己的「瞎」所封鎖而矢去自由.
\newline
\newline
另一方面,耶穌又指出另一事實,(八34)「所有犯罪的,就是罪的奴僕.」這些猶太人不但無知,且在罪的困鎖中.這「罪」不單指一件行為,抑或是一句話:而是指一種「罪」的生活方式,一種作風,「罪」是有連瑣性的,就如說謊,要不斷地繼續說下去,才能遮掩第一句謊話.於是,「罪」的行為便像一個網般裍鎖看人.罪又像一種病毒,蠶食人的身體,使人失去應有的自由.故此耶穌形容犯罪的人,就如奴僕一樣,受罪的管制,失去自由.
\newline
\newline
很多人犯罪,還以為自己享有自由.可是,不久這種所謂自由使成為綑綁.人失去了不犯罪的自由,必須順從肉禮的私慾來行事.以為暢所欲為,實是身不由己,慣常犯罪的人,不覺罪的束縛,除非真正看見自由的真相:才能明白自己的景況.基督要釋放人,使人享受自由的人生,這就要通過真理的途徑,(八36)「你們必曉得真理,真理必叫你們得以自由.」(八36)「所以天父的兒子若叫你們自由,你們就真自由了.」這段經文有兩個層面的意思,一方面提到啟示的內容,另方面論及基督的釋放.以下要從這兩方面思想真理的角色.
\newline
\newline
耶穌指示猶太人得自由的出路,但猶太人怕受束縛而不接納.只是,耶穌清楚說明只有在真理的束縛下,才真正得自由.首先,真理成為生活上的指引,(八31)提到「我的道」,這與真理是相同的意思:意即耶穌所做示的話,並祂整個生活的彰顯:這是真理的流露,是從上而來的客觀資料.「遵守我的道」的「遵守」在有些文字上,譯作「住在這裡,」意即留在這裡.我們的生活,應在真理的約束及保護下,才有真正的自由.耶穌要人「曉得真理,」受其約束,以得自由,「曉得」意即了解,這比理性的探討更深入,就是以生命去體驗基督所說的話；有了知識,又要加上生命的印証.當真理溶化於生命,便得自由.詩篇中詩人也認為當他在律法中行走時,便得到自由.駕駛汽車的人雖不喜歡交通燈,但沒有它則難以維持秩序.
\newline
\newline
另一方面,耶穌還給了另一種的釋放.祂說:「天父的兒子若叫你們自由,你們就真自由了.」外在的引導,需要配合內在生命的質素.「天父的兒子」是神所差派的彌賽亞,要完成救贖的那三位一體的真神.基督本有神的權柄,叫人得自由.祂成為贖罪祭,藉祂的死將罪惡的鎖鏈割斷:使人從罪惡死亡中得拯救.故此,我們若要享受真正的自由,便先要有被神兒子改變的生命:接納祂所完成的救贖.真正的信徒,是一個被罪釋放,受真理引導的人.
\newline
\newline
耶穌要人享受自由的人生,這可從消極與積極兩方面來看.在消極方面,我們從無知的境況中被釋放:既接受真理,便脫離蒙蔽的階段.西方歷史中,有啟蒙時期,說明人從無知,迷信中被理性所啟迪.可惜,這啟蒙時期變了另一種的束縛,就是將理性推至極端.世界上有很多宗教,也是在尋求脫離無知的束縛.例如印度教中,有很多虔誠的教徒,希望透過內省,了解真理的面孔和形像,而從中得到解脫.又如佛祖釋迦牟尼,他利用很長的時期來悟道,希望了解真理,洞察人生的真相:最後,得出以「苦」為真理的結論.故此,他設「聖」「聖諦」及「八正道」來解決苦的問題,其實這些不過是一種主親心理的解釋,並沒有真正的解脫.耶穌基督是永恆的道,親自成為肉身,親自活出真理的模式,將真理向人做示.當人接受這真理,便脫離無知的困苦.從此,我們知道那些看得見的,只是短暫的事物,看不見的才是永恆.當人明白這點時,便超越了人的無知,不會受這些短暫事物所束縛.
\newline
\newline
積極方面,我們看見在基督裡的自由,是從罪中釋放的自由；就如加拉太書所說,在基督裡面的人,不受罪所壓制.然而,我們卻常認為自由是消極的自由,是一種「Freedom from」,而非「Freedom to」,聖經則提到兩方面的意義,跟從主的人,並非要避免做某些事而已,而是一群有自由的人,重新達到神當初做人的目的及理想,反映創造主的完美,這才是真正的自由.
\newline
\newline
此外,跟從主的人,不但享有自由,更要將內在自由的潛質活出來；有行出真善美,活出基督樣式的自由.我每次經過機場,常感到驚奇,一架巨大無比的七四七噴射機,竟能離開地面進入天空,真是奇蹟,這種時刻,地心吸力雖仍在,但有更大的力量──氣體動力去抗拒地心吸力.於是,七四七直插入雲霄,在天空中自由飛翔；這是在真理引導下的自由.飛機師的位置上有很多鏡,這是定律的象徵,當機師留意定律時,這飛機便有真正的自由.最近有蘇聯的飛機墮毀,調查結果發現是機師不看儀器,只靠自己的飛行經驗,就在那時,飛機便失去其應有的自由.
\newline
\newline
主的門徒是被真理釋放的人,生活上有真理的引導,生命上有基督的釋放.不要說你是主的門徒,卻未領略真正的自由.假如這一講大家未能領會；那以下的則更難領會.主要我們不但相信,而且要得到自由,祂希望徹底地釋放我們,使我們的生命真能活出神原來創造人的目的.跟隨主的腳步,被真理引導,似是一個綑綁,其實是真正的自由.
\newline
\newline


\newpage

\section{第59屆~港九培靈會~許道良~第二講}
\label{sec:1237}
\textbf{落在地裡的一粒麥子}
\newline
\newline
連結: \href{https://www.hkbibleconference.org/session-message/view/1237}{\texttt{https://www.hkbibleconference.org/session-message/view/1237}}
\newline
\newline
\hyperref[sec:1236]{< < < PREV SERMON < < <}
~
\hyperref[sec:toc]{[返目錄]}
~
\hyperref[sec:1238]{> > > NEXT SERMON > > >}
\newline
\newline
相信各位都聽過一個童話,說及有一個人問一面鏡子,希望知道世上誰是最可愛的人.雖然是個童話,卻反映出人對自己的看法；環顧四周,誰也不及鏡子裡的自己可愛.
\newline
\newline
有一回,我將車子泊在神學院門口,一位女學生經過,把車的側鏡扭轉到自己的方向,以整理頭髮,她亦知道我坐在車裡觀察她,她只陶醉於欣賞鏡中的自己.然後,她沒有把車鏡扭好就施施然離去.於是,我下車請她把車鏡扭回原位.人時常感到自己是最可愛,只顧個人的利益.自我是最重要的,今天市面上充斥各種商品,都是為了保養顧惜自我；例如各種化裝品,減肥的用品等,都是為了滿足自我的需要.各種廣告也是向我們的弱點挑戰,以提高人的享受為依歸,因此那些廣告都很成功.另外,很多銀行也盡量提供貸款,讓人有更理想的自我生活.然而,這些都不過是生命的化裝品而已,並非真正的生命.假如將生命投資在這些事物上,最終會喪失生命.
\newline
\newline
很多基督徒自稱是主的門徒,卻沒有真正對付自我.十年前的那粒種子,今天仍是那個狀況.「一粒麥子若不落在地裡死了仍舊是一粒.」這是個永恆的真理.經文的上文提到有些希利尼人來見耶穌；他們或許想見神蹟或聽哲理.然而,耶穌只說一句話:「人子得榮耀的時候到了.」跟著便請他們看一粒死在地裡的麥子,以這為榮耀；在一個絕對順服神的生命裡,看見神的榮耀.這不過是以自然界的現象,來解釋深奧的屬靈真理.對於農夫來說,這現象平常得很.
\newline
\newline
我在美國唸書的時候,曾到朋友的農莊裡過感恩節.他帶我參觀各種農業工具,又帶我到榖倉；裡面堆滿種子,是最優秀的種子,留著早春時播種用的.不曉得箇中道理的門外漢,一定認為這是浪費；為何不把這些肥壯的麥子去皮磨成麵粉呢?這比拋在泥土上好些!然而,奇蹟就在那裡出現,經過太陽的熱能,雨水的滋潤,小小的麥子便開始萌芽.至仲秋出現一片金黃色的麥田,原來是一粒麥子變成了千千的麥子,成為一望無際的麥田.事實上,農夫每次撒種都是一個冒險,因為他們不能絕對保證那些種子會發芽生長.將種子磨成粉當然有用途,但這用途非常有限；最多可供應一家幾天的食用.而將麥子撒在地裡,卻有無限量的用途.當我陪同這位農夫坐在收割機上時,看見他臉上那滿足的笑容,這就是最好的報酬.
\newline
\newline
然而,這段經文又如何解釋呢?耶穌說:「愛惜自己生命的,就失喪生命；在這世上恨惡自己生命的,就保守生命到永生.」這就是剛才那段經文的屬靈解釋.查考原文時,會發現前三個與最後一個「生命」,有不同的意義.(路十二15)「人的生命,不在乎家道豐富.」這裡的「生命」與前面最後一個「生命」是同字.此外,(路十二22)「不要為生命憂慮喫甚麼,為身體憂慮穿甚麼.」這個「生命」是另一種生命.前者的「生命」是會過去,因為它與物質世界,形象世界有密切關係,這也是基督要我們放棄的一種生命；我們卻往往背道而馳,只執著形象世界及物質世界的生命.經文中,「愛惜」和「恨惡」是兩個不同的執著點.「愛惜」是關注,將整個注意力集中於其上；「恨惡」則是忽略及輕視之.這種執著是要不留情面的,因為人無法兼顧兩種生命；欲得永生屬靈的生命,就必須要放棄那短暫物質的生命.要討神喜悅就不能討自己的喜愛.基督讓我們認識這種抉擇,並看見其結果.只是,我們有自由選擇,卻沒有自由抗拒這選擇所帶來的結果.耶穌勸我們作一個明智的選擇,若看重今生,永恆的生命就必定受到虧損.
\newline
\newline
有一回,一位佈道家到越南講道.有一位執事自豪地對該位佈道家說:「我做了基督徒二十年.」重複了數遍,這執事在教會的名聲其實不大好；他希望藉著二十年這數字令對方留意自己.只是,旁邊的人聽他這話都感到羞恥；做了二十年的信徒竟仍有舊生命的彰顯.以為「二十年」是個榮耀,認識他的人卻看為羞恥.一粒麥子二十年了仍舊是一粒,從未死去.
\newline
\newline
基督的生命就是個榜樣,若從現世的價值觀來看,祂是個失敗的人,完全缺乏現世人所謂成功的跡象.連枕首的地方也沒有,死時身上只有一套衣褲,還要借別人的墳墓埋葬.這種輕看今生的得失,正反映了「麥子落在地上死了」的心態.祂一生就是過著捨己的生涯,十字架上被釘死不過是捨已最巔峰的表達.從人的角度看祂是一事無成,但十字架上祂說了一句很重要的話:「成了」.這粒麥子落在地裡死了,就結出許多的子粒來；在復活節的清晨我們看見那爆炸性的能力；今天在座各位也能證明這粒麥子的後果,一粒絕對順服的麥子更改變了整個人類的命運.
\newline
\newline
另一方面(十二27)耶穌說:「我現在心裡憂愁.」可見祂心裡有很大的痛苦及掙扎,但很快又提醒自己正是為這個時候而來的,然而最後的考驗實在叫人不易接受；將自己的意願放下,對每一個人來說,都不是一件容易的事情.只是,耶穌順服說:「父啊,願你榮耀你的名.」這粒麥子要死了,為要彰顯父神的榮耀.當時,有聲音說:「我已經榮耀了我的名,還要再榮耀.」從基督順服的生涯裡,神已經得了榮耀,但是祂還要再得到榮耀.以後的榮耀比先前的榮耀更加燦爛；耶穌就憑著這句話而得到力量,走上各各他的道路.然而,祂不但自己走這條路,同時還要求門徒跟從祂走.(十二26)「若有人服事我,就當跟從我.」這是個呼召,要求我們把自己的生命擺上,放下自我.這種捨己的精神,捨己的決定往往是與自己本性相違背的.
\newline
\newline
事實上,「罪」基本上是過份看重自己的需要,只把注意力放在個人身上,而忽視了神的要求.就如(賽五十三6)說:「各人偏行己路.」「罪」意思是──「我」比神更重要.做基督的門徒最大的攔阻,不是家人的反對,或朋友的嘲諷；而是自己,不能放棄自我,與主一起受死,埋葬以致復活,是一種至終持守的經歷.過於顧惜自己的人,不配稱為主的門徒,跟隨主是要徹底的放棄自己.另一方面,這捨己本身並非最終目的,而是透過捨己以成全更高的理想.
\newline
\newline
相傳錫克教的第十祖有一次希望重振其教,卻發現關鍵就在於缺乏一班委身的人；只靠宣講學說及遊行宣傳,不能重振宗教的聲威.因此,他希望從跟隨他的人中；呼召五個對他絕對效忠並委身的人.第一個出來了,教主就帶他入了帳棚,教主出來時手上握著一把淌著血的劍,接著有第二位委身者出來,教主也帶他入帳棚,出來時劍仍是淌著血的,跟著第三位......如是者,直到第五位.這五個人對教主是絕對委身的,最後那五個人出來了,卻沒有受傷,那些血原來是宰殺山羊的血；這位教主就從五個人的委身裡重新建立錫克教；最後成為很有影響力的宗教.雖然他們不認識真神,但那種委身的心情實在值得佩服.我們認識並相信一位真神,然而我們委身的心志到甚麼程度呢?今日香港的教會需要更多委身的基督徒,但數目不需要多,假若今天在座有五個人真正對主順服的話,相信那爆炸的能力,可以改變整個社會.
\newline
\newline
我們常說要跟從主,但因為過於看重自我而攔阻了這個工作；神的榮耀便不能像在耶穌身上一樣彰顯出來.我們一方面高呼服事主,另一方面卻高舉自己的要求.另外,我們的事奉帶有很多條件；要服事主,則先要看時間方便與否,要奉獻則要看有沒有餘下的金錢.傳道人也不例外,很多人委身做傳道,卻很看重自己的意願,致使很多教會感嘆傳道人一代不如一代.當然有些教會故意將額外的重擔加在傳道人身上,也不是正確的做法.一個全心服從主的人不應計較任何條件；當事奉主帶著條件的時候,無形衝散及否定了事奉主的心志；變成了商場上的交易,這不是真正的跟從.要作主的門徒,則必須像麥子落在地裡死了,否則很難結出子粒來.很多傳道人事奉了很多年,卻沒有多少改變,那粒麥子未死透；以致教會很多事工都無法推行.基督所關心的是神的榮耀,我們也應讓自己死去以服事我們的主,並須完成神在我們身上的托負.自然界自然地彰顯神的榮耀；唯有強烈意志的人遮蓋了神的榮耀.
\newline
\newline
教會歷史上有一位被神大大使用的人──佐治懷特腓,神透過他帶來大復興.他講道時那種挑戰令人難以抗拒.當時一位政要去聽道時故意不帶備金錢,因為怕聽道時受感動而把所有的錢奉獻.只是,聽道後他終於要向旁邊的人借錢奉獻.為何這麼多人自稱是傳道人,神沒有大大的使用他們呢?原來,我們從佐治的生命中,看見一粒麥子在地裡面死了.有一回,他在某地講了幾堂道,到晚上已非常疲倦,有一家庭招待了他去住宿.夜深了,有人從老遠來要聽他講道,主人家便代他推辭；擾攘中佐治醒了,明白了究竟後他就給他們講了兩個小時.佐治疲憊極了,但他從來不計較.天亮時,主人發現佐治已離開了世界；神已接他回到天堂,神的榮耀在這粒麥子身上完全彰顯出來!還有,我們所熟悉的宋尚節博士,從人的角度看是多麼浪費,好的學位,好的機會,為何不好好保存呢?神卻讓這粒麥子落在地裡死了,使它有爆炸性的能力.
\newline
\newline
今天我們不是要尋找千千萬萬的麥子,而是尋找一粒絕對順服的麥子.我們是否願意超越自己的有限呢?讓眾人在我們身上看見神的榮耀.耶穌當時的話,今日依然有效:「一粒麥子不落在地裡死了仍舊是一粒；若是死了,就結出許多子粒來.」香港正等候這一粒麥子,中國大陸正在等候這粒麥子；全世界也在等候這粒麥子,但願神在這時代裡,興起祂所重用的人.
\newline
\newline


\newpage

\section{第59屆~港九培靈會~許道良~第三講}
\label{sec:1238}
\textbf{主人如此 何況僕人}
\newline
\newline
連結: \href{https://www.hkbibleconference.org/session-message/view/1238}{\texttt{https://www.hkbibleconference.org/session-message/view/1238}}
\newline
\newline
\hyperref[sec:1237]{< < < PREV SERMON < < <}
~
\hyperref[sec:toc]{[返目錄]}
~
\hyperref[sec:1239]{> > > NEXT SERMON > > >}
\newline
\newline
(路廿二)記載了耶穌為門徒設立了最後的筵席,飯廳裡一切都預備了,我們可以想像門徒陸續入席；由於沒有劃定位置,他們就爭論誰坐在主的旁邊.沒有注意入口處擺放了一個盤,一壺水,一條手巾；這些都是僕人為主人洗腳的用具.飯前,耶穌除掉外衣,捲起衣袖,把水倒在盤裡；取了毛巾,開始為門徒洗腳.我們可以想像當時門徒愕然的景象,每個人都猜測耶穌的用意,可能開始感到羞愧,但又沒有勇氣詢問.
\newline
\newline
今天我們有很多聚首的機會,例如執事會,團契,小組交通,查經班.......我們的心態會否仍是重演當日門徒的境況呢?每次出席聚會的時候,是否想到自己要得到甚麼,別人留意我的需要嗎?尊重我的意見嗎?若感到自己在教會未被尊重,就想辦法離開,找另一間教會.不錯,每個人都有很多需要,有物質上,精神上或靈性上的需要.以致教會內有太多的需要,都沒有人去滿足.要跟從主,就要學習祂的榜樣,耶穌要門徒學習彼此洗腳,做一些非常卑微的工作,但卻能滿足其他兄弟姊妹的需要.
\newline
\newline
首先,我們看看基督那虛己的預表.(約十三1)提到耶穌愛屬祂的人愛到底.祂知道自己將要上十字架,有時間的限制,只有爭取機會以表達愛的觀念.祂主動去愛別人,沒有計較那些人配不配受自己服事.我們的主,萬王之王,紆尊降貴,因為愛的緣故,願意把這一切放下.這萬王之王應該是受人尊崇,受人愛戴崇拜的；但祂卻作了這卑微的服侍.基督離開榮美的天家,來到卑賤的世上,立定志向過一生服侍的生涯；不是尋找服事的人,而是尋找服事的對象.四福音中常看見耶穌接觸那些受社會輕視的人,如乞丐,瞎子,長大麻瘋及患長年疾病的人,甚至死屍.以前的印度人分為四級,以祭司最高級,祭司們絕不能接觸社會上的賤民,以防受到污染,影響自己的救恩.我們的主絕不是這種祭司,祂來到世上,從不以階級的角度去看待別人.中國歷史上也有皇帝微服出巡的事蹟,但也不會接觸社會上低層的人.基督來到世界,根本沒有人察覺祂是神；甚至祂的門徒也需要一段時間才明白過來.耶穌並不是故意掩飾自己的身份,只是祂放下自己的身份,以致別人根本看不出祂就是萬王之王；這也是祂在別離時,要教導門徒的功課,祂拿起毛巾來洗門徒那骯髒的腳.這些人中,有包括那在眾人面前否認祂的彼得,並出賣祂的猶大.當時,沒有一人配得接受耶穌那卑微的服事,基督一生就是過著奴僕的生涯,雖然將來有一天祂要作王；但在世的時候是做一個奴僕,祂沒有自己的主權,只有被人利用,甚至吃飯的時間都沒有.這種奴僕的心態,在耶穌日常生活中表露無遺；彼得知道自己不配被耶穌洗腳,便拒絕接受；我相信他一半是基於客氣的,因此他沒有搶了耶穌手上的毛巾,替耶穌洗腳.他只是拒絕耶穌為他洗腳,可見,彼得並未真正了解這行動背後的含意.
\newline
\newline
此外,我們可從另一個層面了解洗腳的意義.在當時的文化來說,只有奴僕向尊貴的人洗腳；基督便透過這謙卑的服事,帶出非常重要的屬靈意義.耶穌說:「我所作的,你如今不知道,後來必明白.」這暗示行動背後還有另一個意義.又說:「我若不洗你,你就與我無分了.」表明卑微服事的背後,還有潔淨的象徵.一個罪人若未被潔淨,根本不能親近一位聖潔的神.在神的國度裡也沒有份兒.彼得聽了耶穌的話,便立即有了反應；要求主也洗他的頭和手,希望與基督完全有分.只是,耶穌進一步解釋說:「凡洗過澡的人,只要把腳一洗,全身就乾淨了.」這或許是聖經上所講的洗禮,我們既是基督內被潔淨,洗禮不是重複的行動,一次洗了就足夠.然而耶穌又暗示已受洗的人,也需要做一些潔淨的工作.祂說:「你們是乾淨的,然而不都是乾淨的.」這是雙關語,暗示中間有一位猶大,是未被潔淨；同時也暗示,我們的生命並非完全潔淨,需要作一些潔淨的功夫.(來九11-14)清楚說我們不是靠水來洗禮以得潔淨,這不過是外表的象徵,印證我們生命上已有的改變.然而,我們得以潔淨是靠基督的血.因此,耶穌的卑微,順服,是要從「血」的角度去看,門徒是看見水；耶穌都要他們明白,真正的卑微,犧牲不是水,而是用血來表達.基督愛門徒愛到底,連自己的寶血也願意傾出來,生命也獻上了；因為若不流血,罪就不得赦免.潔淨的背後,是基督生命的付上.十字架的工作,不但一次潔淨我們,這赦罪的泉源還繼續洗淨我們信主後的過錯.神只為人類預備了一個潔淨的泉源──十字架.歷代基督徒常犯了一個毛病,雖然靠著十字架的寶血得到赦免,卻希望靠著律法繼續潔淨的工作.神並沒有為人預備兩個潔淨的泉源,基督在十字架上所成就的救贖,不但能赦免我們信主前的一切罪過；同時,也能夠赦免我們以後的過錯.這就是耶穌要門徒明白的.
\newline
\newline
彼得要學習一個功課,就是承認自己的需要,願意讓主洗腳,這也是我們應學習的,我們信主後同樣會犯錯,但卻不應將自己的軟弱,錯失加以掩飾,或找藉口解釋；應像彼得那樣,讓主洗腳.最近有弟兄決定到北美讀神學,臨走時問我,要做一個傳道人,需要注意甚麼；我告訴他,應時常注意自己.因為我們常看見別人的軟弱,過失,卻忽略了自己的生命；身為傳道人,若忽略了自己的生命,事奉就會失去能力.若能在神面前保持潔淨的生命,神在任何艱難的環境,也會使用我們.耶穌洗門徒的腳,是要他們保持潔淨,因為有重要的使命託付他們.
\newline
\newline
另一方面,洗腳只是個象徵的儀式,不及內裡的心態來得重要.當日猶大雖也被耶穌洗腳,卻仍是污穢的.今天,許多人接受洗禮,但並非每個人都真正被潔淨.表面上的罪,雖能爭取別人的饒恕,但不等於內心有真正的懺悔.我們認罪的動機其實相當複雜,可能為了逃避責任,或希望別人接納；甚至是透過認罪來掩飾自己其他的罪行,這些都不是真正的認罪.在整個預表裡,耶穌要門徒學習彼此洗腳.洗腳象徵謙卑的服事,祂謙卑地服事他們；使他們知道怎樣用謙卑的心態彼此服事.只是,服事與謙卑往往不一定能連在一起,很多人服事主,卻少有謙卑的服事.
\newline
\newline
以前有位執事,在教會裡相當活躍,卻不能與其他執事和洽相處；開會時常命令別人做事,偶不如意他便大聲喝罵.一回,家父在聚會時將奉獻的程序稍作調動,這其實沒有多大影響；然而執事在台前大聲指責家父做錯事；沒有徵求他同意便更改聚會的程序.這位執事有服事,卻是存著大老闆的心態去服事.真正的服事,是不見其人,但見其工作的果效,看見需要得到滿足,污穢得到潔淨.可惜,很多人不過是透過服事來抬高自己；使人看見自己的重要性,這是在服事上常出現的試探.耶穌要人彼此洗腳,但不要人沒有存謙卑的心,因為這種服事只能完成事工,而不能建立生命.唯有存著謙卑服事的心態,才能接觸別人生命的需要,在這接觸中建立別人,激勵人一同走天路.
\newline
\newline
可惜在教會團契內,我們看見許多弟兄姊妹被神饒恕了,卻一直未被其他弟兄姊妹接納；這裡我們應當學習的功課.主既然饒恕他,洗淨他,我們那有權去記恨呢?多少團契生活為著人沒有彼此饒恕而受到虧損,基督不單用水洗我們的腳,更用血來洗我們的生命；我們豈敢不饒恕其他弟兄姊妹呢?我是誰呢?我不過是在十字架影下被潔淨的罪人,我們怎樣對待犯錯的弟兄姊妹呢?真正的門徒,手裡不是握著法官的小鎚；而是握著一條毛巾,隨時為人洗腳.耶穌說:「我是你們的主,你們的夫子,尚且洗你們的腳,你們也當彼此洗腳；我給你們作了榜樣,叫你們照著我向你們所作的去作.你們既然知道這事,若是去行就有福了.」
\newline
\newline


\newpage

\section{第59屆~港九培靈會~許道良~第四講}
\label{sec:1239}
\textbf{沒有間隔的羊圈}
\newline
\newline
連結: \href{https://www.hkbibleconference.org/session-message/view/1239}{\texttt{https://www.hkbibleconference.org/session-message/view/1239}}
\newline
\newline
\hyperref[sec:1238]{< < < PREV SERMON < < <}
~
\hyperref[sec:toc]{[返目錄]}
~
\hyperref[sec:1240]{> > > NEXT SERMON > > >}
\newline
\newline
一直以來,猶太人都以為自己是神的選民,只有他們才被神照顧餵養,以為神計劃裡只有猶太人的羊圈.但感謝神,這並非神計劃裡的最後一章.(約十)耶穌提及祂另有羊不是在羊圈裡面,這對猶太人來說是很大的威脅,多個世紀他們以為羊圈內只有猶太人,沒有外邦人.耶穌甚至說另外的羊,要與現在的羊群合一,同歸於一個牧人.第一世紀的猶太人基督徒經過很大的考驗,才稍為領悟這功課.然而,這問題到今天仍然成為我們的挑戰.
\newline
\newline
雖然主所牧養的羊圈是沒有間隔的,但人將其分隔成很多小羊圈,彼此間缺乏溝通,團契,甚至互相攻擊.環顧今天的教會,可謂宗派林立,有些沒有宗派背景的,則批評有宗派的人,事實上,他們主要的分別在於沒有歷史的根,無宗派的教會也有自己獨特的儀式,組織及事奉方式.我並非反對宗派的設立,因這純屬必然,是文化地域的不同而形成的.然而,我反對宗派主義,他們以自己的宗派為最純正,最是合乎神心意的教會.當我們高舉自己時,已違背了神的心意.我關心的不是外表組織的問題,而是彼此間合一的精神,使世人看見一個教會,一個福音,一個救贖計劃.約十七章是大祭司的禱告,我們的主儼如進入至聖所向神懇求,為門徒及因門徒的道理而信之人禱告,表示主很關注合一的問題,暗示合一在實踐上的困難.同時也看出其重要性.
\newline
\newline
首先,我們看合一的意義.耶穌給我們建立合一的模式,就正如聖父,聖子的合一那樣.合一並不是指劃一,基督道成肉身是有形有體,父神的靈性是無形無體的,故此父子不是劃一,而是在某一層面,範圍的合一.門徒也一樣,每個都有自己的特性,並不劃一.今天一些異端都強調劃一,只容許一個意見,一種解釋,以領導人為依歸,這是一種統治的手段,不是神的心意.主不要求我們劃一──有相同名稱,相同組織,相同行政方式.......
\newline
\newline
合一也不是聯合,聯合通常是指組織上的結合,為了工作上的方便及效率,當這目的達到後,組織便解散.在聯合行動中,各人仍是以自己的利益為重:同時也缺乏生命的接觸.可說是基於自己的利益,而加入聯合組織,聯合國便是一個例子.可惜教會歷史裡的合一運動,都不能超越聯合的性質.
\newline
\newline
父與子的合一,是生命上的聯繫,父在子裡面,子在父裡面,主希望門徒有這種合一,在毫無攔阻的交流中,兩方面的心意達成一致,父的旨意等於子的願望,人看見子就如看見父.父已將榮耀賜予子,意即基督完全彰顯神的榮耀,與神合一,看見基督則如見神的榮耀.因此,基督要將榮耀賜予門徒,好使他們彰顯祂的屬性.當重點放在神的榮耀上,合一便顯得簡單,否則,當我們標榜自己時,便無法談合一.耶穌要他們關注神的榮燿,才能實踐合一.保羅在以弗所書形容合一,是以身體各肢體來解釋,這例子很貼切,身體各肢有獨特的功能,但卻形成一個身體,這是主希望我們達到的合一.
\newline
\newline
至於合一應建立在何種基礎上呢?耶穌所說的合一,不可建立於人為的基礎上.十七20耶穌說,「我不但為這二人祈求,也為那些因他們的話信我的人祈求,使他們都合而為一」.可見,合一應建立在屬靈的基礎上.人因聽了純正的福音而有生命的改變,福音只有一個,就是門徒所宣講的福音,這因信稱義的福音;完全是基督的恩典,並非人勞苦得來的.耶穌說合一要建基於這基礎上.今天很多異端邪教都鼓勵人信耶穌,這不過是掩飾而已.摩門教也談耶穌,但他們所宣講的,不是門徒所講的福音,不少真理基礎膚淺的便被吸引而離真道.各位應警惕,當一些人所傳播的信息,不是使徒當天宣講的,則不可能與其談合一,經文中的合一,是一群歸入基督的人,與基督有生命的關係,這聯繫成為門徒之間合一的基礎.換言之,失去與基督的關係,也就不能談合一,因為大家的生命質素不一樣,人生方向及最終的命運都不同.然則,與神的關係存在的話,合一的實踐,甚至不需我們努力也能經歷得到.
\newline
\newline
最近,一位姊妹從美國到中國旅遊,回程路經香港,她聽過我及內子的名字但素末謀面,旅途中她很焦急盼望見我們.相見後我們一見如故.因為與主有親密的關係,便嚮往見主內的弟兄姊妹.這種合一的實踐毫不費力,很自然便表達出來.然而我們應小心,一方面不可在不同的基礎上隨便論合一,另一方面,不應以人為的方法擴大基督所立的基礎.今天很多合一的攔阻是人為的問題,在相信基督為救主外,還加上些條件,例如接納獨特的儀式,傳統,教會條例等人為的規範,致使我們在主內不能經歷合一的寶貴.
\newline
\newline
有些教會認為合一暫時無法實踐,便以此為藉口而推搪合一的責任,推說合一要到天國後才能實現.實際上,天國已進入人間,屬天國的子民今天就可實踐合一.以下兩節經文,主很關心一件事情；21「叫世人可以信你差了我來.」23「叫世人知道你差了我來.」由此可見合一這事實是可以叫人知道的.也就是說,合一的屬靈基礎是無形的,但合一可以透過某種形式表達出來,使世上相信神差了耶穌來.講到合一時,常只強調內涵而忽略形式,抑或兩者相反,至讓我們看到兩方面的平衡,合一有生命的連結,是真正屬靈的內涵,同時也有外表的形式.若有合一的心志,卻從來沒有行動,這心志很值得懷疑.正如青年人談戀愛,他們心志絕對不會隱藏,愛慕對方,則千方百計找適當的形式表達其心志.
\newline
\newline
基督徒最少可用三個層面表達合一.首先,是團契的表達,團契意即信徒生命彼此交通,交流.香港只是彈丸之地,交通方便,因此,應多製造生命分享的機會,不可被動地期待團契的出現,而應主動爭取團契的機會,例如探訪附近的教會,使生命能接觸生命.這溝通使我們產生彼此建立,彼此提醒的作用.這就是主的心意.
\newline
\newline
另一層面是合一的見證,在今天道德破產的時代,信徒應對罪惡有一明確的立場.社會上充滿不公義之事,屬主的人應有合一的見證斥責罪惡.一盞燈雖有獨特的功能,但在社會上其功能非常有限,若將這些燈集合起來,讓整個社會看見我們有明礭的合一立場,那一定會產生巨大的影響.教會內不需要隱藏在小羊圈的燈,因為燈是要照耀黑暗地方的.對於不道德的行為,我們應有合一的見證,提倡公義的原則,使別人知道神是慈愛,又是公義的神.
\newline
\newline
另外是事奉的層面.我們在神面前領受一個共同的使命,不要時常想霸佔地盤,把自己的羊蓋上印章,我們應共同承擔使命.香港有很多事工已開始有合一的表現,七六年葛培理佈道會時,合一的能力遍及整個香港社會.年底有包樂佈道會,希望大家爭取合一事奉的美好機會.
\newline
\newline
合一的實踐是個成長的過程,不是一朝一夕的事奉.以弗所書四章也提到這真理.主沒有希望合一使所有事情變為一致,這是我們的目標,也是主的方向；因此,在過程中有不協調的情況是理所當然的,不要在發現困難時便放棄.記得小女學走路時很困難,要靠我們慢慢鼓勵,悉心教導才有進展,假若當初我們因她身體的不協調而放棄.恐怕今天她還不懂步行.合一也像這成長過程,當我們踏出第一步時,也就成為第二步的力量,直到討神喜悅.
\newline
\newline
為何耶穌如此看重合一呢?21「使他們都合而為一,......叫世人可以信你差了我來.」在舊約時代,先知被接納與否取決於他是否神所差遣的.先知若不是神所差遣的,則沒有人會接納他的話；可見差遣的重要性.以賽亞書五十三章講到基督,說祂是神所膏立公義的僕,會成為人類的救贖者.耶穌的身份,使命,與合一的見證有著密切的關係.約十二44提到,我們信主耶穌,也就是信差祂來的父.約十三20又說:人接待門徒即接待基督,接待父所差來的子,也就是接待差子的父.因此,合一的問題與基督救贖的計劃有密切關係.十七3「認識你獨一的真神,且認識你所差來的耶穌基督,這就是永生.」合一的問題與整個人類的救贖有直接的關係.基督要人看見祂為人預備的救贖,祂是為此被差遣到世上,我們也為此被差到世界去；因此,整個救贖計劃與合一見證有密切關係,當我們四分五裂時,便不能講合一的信息,世人也不會接受.
\newline
\newline
二千年來,基督的門徒建立了無數的教會,醫院,孤兒院......差派了多少宣教士,按立了多少牧師,餵飽了多少飢餓的人,安慰了數不清受傷的心靈,然而,合一的心志,合一的表現仍然是最脆弱的一環,信徒之間仍有間隔,這會否成為今天許多人未聽見福音的原因呢?世人會否因為我們的分裂及彼此漠視,而留在羊圈之外呢?求主幫助我們拆除這些間隔,因我們原屬一個牧人,就是那為我們捨命的基督,我們的主,我們的牧人.
\newline
\newline


\newpage

\section{第59屆~港九培靈會~許道良~第五講}
\label{sec:1240}
\textbf{葡萄樹上的枝子}
\newline
\newline
連結: \href{https://www.hkbibleconference.org/session-message/view/1240}{\texttt{https://www.hkbibleconference.org/session-message/view/1240}}
\newline
\newline
\hyperref[sec:1239]{< < < PREV SERMON < < <}
~
\hyperref[sec:toc]{[返目錄]}
~
\hyperref[sec:1243]{> > > NEXT SERMON > > >}
\newline
\newline
基督徒的生命絕不是消極枯燥,缺乏色彩的,基督徒應有豐盛的生命.神賜與我們的生命,不但是沒有疾病或瑕疪,乃是積極結出果子的生命.我曉得一些基督徒,他們雖沒有犯錯,但他們的生命對人並沒有幫助,可算是一些沒有果實的枝子.
\newline
\newline
要得著豐盛的生命,能結出果實,是有條件的,並不能單靠自己而達到.都市化的社會,常令人以為自己不需要他人的支持；試想,我們在大廈裡住了十數年,卻不認識自己的鄰舍.於是,我們便產生了錯覺,以為單憑己力便能解決自己的需要.其實,人是需要相當大的支持系統.生理上我們需要空氣,水份等養料來維持；可惜在超級市場內,我們被精美的包裝蒙蔽了,看不到農夫的勞苦,工人的辛勞；只看見金錢帶來的滿足.我們的生命得以維持,必須靠他人的支持,當我們和那些支持系統脫節後,生命便出現危機.屬靈的生命也不例外,有人以為信主後可靠自己過基督徒的生活,這不是門徒的心態；因為當我們與屬靈的支持系統脫離關係,生命就會枯萎.耶穌說:「你們若多結果子,就是我的門徒了.」耶穌心目中的門徒,是要多結果子,具有豐盛生命的人.舊約詩篇第一篇,指出義人是那些按時候結果子的人.單有外表並不足夠,神要我們的生命具有實質.詩九十二14說:義人到年老的時候仍要多結果子.這是真門徒的生命.葉子只可被人觀賞；但果子卻可滿足人的飢餓.耶穌希望我們的生命能滿足別人的需要.
\newline
\newline
現在讓我們思想生命的源頭,(十五1)「我就是真葡萄樹.」耶穌說祂就是主幹,是供應所有支幹的命脈.生命就在耶穌裡面,若離開這生命的源頭,生命便會枯萎.歌羅西書也說一切豐盛皆於基督內.耶穌是我們得生命的源頭,就是在我們得救後,也要靠主維持我們的生命.若以為得救後可靠自己的力量行事,這樣便遠離了豐盛生命的源頭.
\newline
\newline
要得著豐盛的生命,有何秘訣呢?(十五1-2)「我父是栽培的人,凡屬我不結果子的枝子,祂就剪去；凡結果子的,祂就修理乾淨,使枝子結果子更多.」(十五4)「你們要常在我裡面,我也常在你們裡面,枝子若不常在葡萄樹上,自己就不能結果子；你們若不常在我裡面,也是這樣.」這裡清楚說明兩個秘訣,第一是神的作為,然後是人的責任.
\newline
\newline
神的作為是我們無法做的工作,正如有些人的生命結滿果子,有些則只有葉子；關鍵在於有神的工作,這在(十五2-6)可看清楚.「凡屬我不結果子的枝子,祂就剪去.」(十五6)「人若不常在我裡面,就像枝子丟在外面枯乾；人拾起來,扔在火裡燒了.」神有兩方面的工作,一是剪除,另外是修理.前者第二節的經文,其中一個解釋是；有些人自稱為基督徒,事實並不是；故此剪去的是那些原不屬於基督的人.我卻不同意這解釋,因為聖經清楚提及「凡屬我的......」,是指那些真正屬基督,卻不結果子的人.另一個解釋,認為剪去是指工作上神審判的意思,並不影響個人得著救恩.然而,第六節說到那被拋棄到外面的枝子,被火焚燒,這正是題到那最後嚴厲的審判；是對付生命的問題.與工作無關.我們可從幾方面了解經文；前面題到那些「屬主的人」或「在主裡面的人」,這些人曾經有一段時期與主建立了生命的關係,後來卻放棄了,以致生命枯萎；最終被神所剪除.雖然,(十28)有應許「我又賜給他們永生,他們永不滅亡,誰也不能從我手裡把他們奪去.」這是千真萬確的應許,但是,我們要留意,信主時無人強迫我們；神看重我們的自由,是以神也不會故意攔阻我們離開主.不錯,沒有人可將我們從基督手中搶過來；但若是我們故意離開主,神便任憑我們.如此則沒有救恩為我們預備.希伯來書亦提到有人曾嚐過救恩的滋味,後來故意離棄主,如此這枝子就會枯乾,最後被剪除.神的恩典是嚴厲的,忽冷忽熱的教會,祂要從口中吐出.
\newline
\newline
另一方面,神另外一個作為會令我們十分安慰.「凡結果子的,祂就修理乾淨,使枝子結果子更多.」一方面神有剪除的工作,另外卻有修理的工作.為何基督徒的生命要修理呢?因為生命不是塑膠花,基督徒的生命有其內在的潛力,是會成長的生命.在成長過程中或許有不合乎神心意的地方,那便要剪除.神修剪是因為祂看重我們,知道我們的生命可以結出果子.不少人進神學院,以為四年的課程是個甜蜜的夏令會,或者是安靜的退修營,一切掙扎已成為過去,殊不知道這就是神開始對付他們的時候.神既悅納了他們那要結果子的心志,跟著剪刀就揮動了；我們看到在激流痛苦中,一個個生命成長起來；成為比前更加燦爛,美麗的生命,這是我們在神學院感到最為滿足的事.神的剪刀也會修剪我們的生命,要我們學習順服,結更多的果子.神所修剪的不一定是罪惡的事情,而是一些攔阻我們生命結出果子的枝葉.很多年青人會問,談戀愛有何不對?談戀愛並不是錯事,但在人生某個階段；戀愛會影響神在我們生命裡要做的工作；看報紙,看電視,下棋......這些都不一定錯；但當這些事情,娛樂攔阻我們的生命結出果子,神就會修剪；要我們放棄,最終目的就是要我們結出果子.若願意接受神的修剪,生命的果子要比前多三十倍,六十倍甚至一百倍之多.很多人顧惜自己,不願被神修剪,神也就順他們的心意,但他們便成為沒有果實的生命.
\newline
\newline
此外,神有神的工作,我們也有自己應當承擔的職責.(十五4)「你們要常在我裡面,我也常在你們裡面；枝子若不常在葡萄樹上,自己就不能結果子,你們若不常在我裡面,也是這樣.」耶穌說祂會常在我們裡面,這是很確實的應許；問題是我們會否在主的裡面.「常在主裡面」就是恆久持守與主的關係.有些人信主後並不願意與神保持緊密的關係,而與祂保持一段距離；生活上並非全沒基督徒的樣式,但不是「常在主裡面.」當我們與主的關係疏遠時,不但不結果子,同時會枯萎；我們時常會遭遇很多試探,引誘,使我們漸漸與主疏遠或甚至離開主.例如進入大學,有新的知識,便開始懷疑信仰；出來做事沒有時間,便與教會疏離,面臨結婚時,為著對方而放棄了與神緊密的關係.弟兄姊妹,不要被引誘,陷入試探中而跌倒,繼而離開主；因為離開主,我們不能結果子.
\newline
\newline
至於在實際情況下,如何把「常在主裡面」這句話實踐呢?我認為有兩方面:
\newline
\newline
首先,要有與主相交的時刻,一個安靜的靈修.安靜的時刻能提醒我們在繁忙的生活中,安靜地等候神的話語,就像馬利亞坐在主的跟前等候祂一樣.我們時常作要求的禱告,很忙碌地說話,沒有讓主有機會插進我們忙碌的生命裡.我們需要安靜以聽到主的聲音,與主相交；讓主的生命影響我們的生命,這就是安靜的靈修.
\newline
\newline
其次,要與主同行,這包括我們日常一切活動,無論一舉一動等言談舉止,都在主的範圍裡面.要時刻學習在主面前過日子,這種生活並不是憑著感受,而是一個信心的生活,在主裡面學習信心的功課.主要求我們常在祂裡面,與主有緊密的關係；羊圈外有各種引誘,我們只能在某個範圍內生活,支取所需的營養.聖經裡,神十分欣賞能與祂同行的人,只可惜為數的不多.神要我們如何與祂同行呢?照彌迦先知所說:「神要我們行公義,好憐憫,存謙卑的心,與你的神同行.」(彌六8)神要我們在生活上表彰祂的榮美,更重要的是存謙卑的心與神同行,「常在主裡面」是生活上任何方面,包括職業,學校,商界,市場等問題,都要在神裡面.當我們的生命與主有這廣闊層面的聯繫時,結果子將是極其自然的效果.
\newline
\newline
我們的生命結果子,並不是為了自己去享受,而為使別人得飽足.聖經上有多處經文提到生命的果子.如(弗五9,加五22)都提到生命的果子；然而,生命的果子與事奉的果子有連帶關係.神不但期望我們的生命酷似基督；同時,也希望透過我們的生命餵飽其他有需要的人；當人接觸我們時,也感受到我們的生命力.
\newline
\newline
我唸神學時曾參加一間教會,這教會不大,但每逢主日都擠滿人,我發現關鍵在這教會的牧師,他每次講道時都能餵養人的需要；他並非是個華麗的講者,但在他話語裡流靈出生命的果實.因此教會雖細小,但卻擠滿了人.很多飢渴的人來到他面前得到飽足.這位牧師的秘訣,是定期有一天完全不見客人,整天在神面前安靜,與神靈交；由於他與神有緊密的關係,便能源源不絕地供應他人.這就是神要我們過的那種生活.不是為自己而活,而是結出果實以滿足人的需要.
\newline
\newline
另一方面,我們多結果子,神便得到榮耀,這就是基督徒人生的最高目標.我們並不是為自己活,而是為神而活.(太七)提到「憑著果子就可以認出他們來」,這些果實,是要印證神在我們生命上的一些作為,使人從我們的生命見神.要彰顯神的榮美,還有目標比這更高嗎?基督一生的目的是要榮耀神；只要能達到榮耀神的目的,祂任何代價也願意付上,基督曾說:「我來了是要叫羊得生命,並且得的更豐盛.」(約十10)我們雖然屬主,脫離了罪惡；神更要我們進入祂那豐盛,這就是耶穌基督的門徒.我們的生命,就盛載著神恩典的器皿,使多人藉此得嚐主恩的滋味.
\newline
\newline
今天我們要省察自己屬靈的境況,可能已有多年忽略個人的生命；以為與主保持不冷不熱的關係已經非常足夠.間中到教會,間中禱告,間中有靈修,這都不是神所喜悅的.屬靈的生命不是一種遊戲或娛樂,可以隨時放下；而是一種生死的問題.當被引誘要離開主時,請記著有一天我們要面對神可怕的審判.應在神面前好好地省察.另外,若果發覺自己受很嚴重的考驗,經歷艱難的試煉,不要灰心放棄,可能神要在你的生命裡做一些修剪的工作；使生命更為蓬勃,結出更多的果實.因此,不要抗拒神的工作,要甘於放棄自己,順服神的心意.神從不浪費我們的眼淚,不會白白讓我們受苦；因祂是慈悲憐憫,祂是充滿智慧又是信實的主.我們願意成為哪種生命呢?蓬勃豐盛的生命,抑或成為枯枝呢?
\newline
\newline


\newpage

\section{第59屆~港九培靈會~許道良~第六講}
\label{sec:1243}
\textbf{沒有選擇的命令}
\newline
\newline
連結: \href{https://www.hkbibleconference.org/session-message/view/1243}{\texttt{https://www.hkbibleconference.org/session-message/view/1243}}
\newline
\newline
\hyperref[sec:1240]{< < < PREV SERMON < < <}
~
\hyperref[sec:toc]{[返目錄]}
~
\hyperref[sec:1245]{> > > NEXT SERMON > > >}
\newline
\newline
今天要講到主給門徒的命令.命令是個無可選擇的吩咐,要作主的門徒,則非接受不可；以下要講的是愛的命令.
\newline
\newline
在這個時代,談戀愛的自由比以前大得多；不少年青人很年輕就談「戀愛」.由於父母家庭不加約束,因此,到適婚時期已累積了不少戀愛經驗；照理說他們婚姻成功率理應高過上代,事實卻相反.許多人以為自己懂得愛,卻根本不明白愛是甚麼一回事；對於現今的社會,「愛的教育」比「「性教育」來得更迫切.這不但要有理論,還需要示範.我相信屬基督的團體,應該最有資格示範「愛的教育」,這也是基督要門徒作的.耶穌臨別時,沒有叫門徒建立紀念碑,也沒有為門徒設計一個嚴謹的組織,沒有教導門徒怎樣賣廣告宣傳自己......卻吩咐他們要彼此相愛.這就成為門徒獨特的標誌,我們若有這標誌,眾人必看出我們是基督的門徒了.我建議對預備受洗的弟兄姊妹,不但要考問信德,也應考問愛心.信心是個人得救的問題,愛心則是關乎別人的需要.基督的門徒不能只顧自己得救的問題,也應付出自己的生命,使別人得祝福.故此,我們接受新會友時要更加嚴謹,以免我們的標誌被沖淡.
\newline
\newline
耶穌說:「我賜給你們一條新命令,乃是叫你們彼此相愛.」青年人常說墮入愛河,意即身不由己以致滑入愛河中,這是羅曼蒂克情感的愛；這種愛不可被命令,來時不受命令,消失時也不能用命令去挽回.然而,耶穌不是談情感的愛,而是意志的愛；是自己做決定去愛別人.兩種截然不同的愛,前者基於別人的吸引；後者則基於自己的決定,是可以被命令的.因為重點在自己身上.很多人談戀愛先看自己的需要,挑選愛的對象時,要選一個能滿足我理想的才愛；所以很快便轉移愛的對象.因為有崇高的理想,理想的丈夫或是理想的太太；相處日久便發現對方不合自己的標準.甚至婚後仍有這想法,實際上這種愛是自私的掩飾.耶穌所講的愛是以對方為中心,這種愛不是一種享受,權利,乃是付出,是責任的愛.我們抱玩可愛的嬰孩,是一種享受；與廿四小時照顧這嬰孩有很大分別,這才是意志的愛.
\newline
\newline
此外,耶穌要我們明白這不但是命令,也是一個共同的責任,祂不是把這重擔交托給我們中間的一個人,乃是叫我們彼此相愛；這崇高的愛要像網狀互相聯繫才能實現.可惜,在教會裡常只有幾個人付出愛,弄至筋疲力竭,負荷過重,甚至精神崩潰.愛的擔子不能由一小撮人承擔,耶穌要我們一起承擔這愛的責任.另一方面,耶穌說這是一個新命令.舊約律法的總綱也是愛,然而只要求做到愛人如己,這已是非常崇高的愛.耶穌說:「我要你們彼此相愛,像我愛你們一樣.」基督愛我們以致犧牲自己,這愛超越了舊約的愛.因此,是一個新命令,一個更高的要求.耶穌說世上沒有比這更偉大的愛,當我們還作罪人的時候,祂就為我們死；因此基督的愛,就成為我們相愛的標準.
\newline
\newline
首先,主的愛是犧牲的愛,是完全付出的愛.其次,主的愛表示了生命的交流.耶穌視門徒為自己的朋友,不是主僕的關係:乃是朋友的關係,朋友之間有親密的交流.因此,門徒可以毫無攔阻地向神祈求:雖然他們有軟弱,時常跌倒:只是耶穌視他們為朋友,這是個完全的接衲.可見,耶穌的愛是個主動的行動,祂說:「不是你們揀選我,而是我揀選了你們.」祂主動揀選了他們成為愛的對象,甚至派遣他們去結果子.
\newline
\newline
除了基督的愛作為標準外,另一個客觀的標準,就是主的話語.(約十五10)「你們若遵守了我的命令,就常在我的愛裡；正如我遵守了我父的命令,常在祂的愛裡.」「我的命令」就是祂所吩付的一切話:基督所啟示的就是客觀的真理,也是絕對的真理.「愛」要在真理的引導下實踐出來.神有劃定的範圍,當我們超出了這範圍,就沒法經歷祂的愛:因為神不會為愛的緣故而違背祂聖潔的本性,在愛裡不能容認罪惡的存在,這也是神對我們的要求.故此,我們應遵守主的話來愛弟兄姊妹,這是純潔的愛.現今的處境倫理學高舉另一種愛；認為愛是沒有原則,沒有真理,只要有愛,則任何情況下都可付出愛.然而,若缺乏客觀真理的引導,愛就變成主觀的決定:受環境及當前處境所控制,這不是真正的愛.我們應在真理裡面找到愛的方向.愛和順服似是兩回事,郤彼此關連,一方面要順服主的命令,另一方面要用愛去表達這個順服.我們遵守主的命令表示對祂的愛,然後有足夠的能力愛弟兄姊妹.以弗所書中,保羅勸勉丈夫愛自己的妻子,也吩咐妻子順服自己的丈夫.也許姊妹們會感到不公平,實際上愛和順服,二者箇中的精義是一樣的.保羅要求丈夫愛妻子到達忘我的程度,同時,也要求妻子順服丈夫到沒有自我的境界.因此愛和順服的精義都是完全以對方為出發點的意思.
\newline
\newline
另外,基督不但給我們命令,也讓我們知道怎樣去支取愛的力量.祂吩咐我們要常在祂的愛裡面,才能得到愛的力量.「常在主的愛裡面」,也就是要留在主愛的範圍內:不可隨便轉移愛的對象,我們要學習在主的愛內,無論順境或逆境,與主保持親密的關係:初信主時體驗主的愛,信主後直到今天都在主內享受祂的愛.以弗所教會出了問題,就是因為離棄了起初的愛:雖然他們有很多工作,但在神的眼中卻沒有多大價值.我們彼此相愛,是因為主的愛洋溢而自然流露:這種愛不是我們本性所有的,故此要常在主的愛裡,才能夠彼此相愛.另一方面,我們彼此相愛,因為主先愛我們,我們相愛是回應主愛的行動:既是這樣,我們就不應計較對方的反應.若果我們清楚這觀念時,就可以去愛一個不可愛的人或甚至敵人；因為這是對主愛的回應,是主的愛充滿我們的自然流露.
\newline
\newline
許多人喜歡講「愛」,卻難付諸於行動:因為愛使我們將生命敝開,失了防備,就容易受傷害.故此,我們在這方面非常吝嗇.耶穌要我們真正地實踐愛:雖然我們可能受到傷害,被利用,但仍要去愛.實踐愛時,要注意畿點:
\newline
\newline
(一)要存著犧牲的心態　愛不是接受,而是付出.愛也不是平等交易,否則我們也勝不遇墨家的倫理標準.墨家講「投之以桃,報之以李.」這是平等交易的愛.犧牲自我的心態,就是說我一切的擁有物也可以犧牲,甚至生命也可以擺上.
\newline
\newline
(二)耶穌要門徒彼此相愛　這愛就是要他們有生命上的交流.愛永不是孤島,愛應能接觸其他的生命,使其他人也經歷這愛.基督徒的圈子裡有團契,團契的真義是生命的分享,交流.今天我們談到溝通時,往往只重技巧,只注意表達自己.其實,更重要的是願意去聆聽別人的傾訴:與快樂的人同樂,與哀傷的人同哭,這才是愛.
\newline
\newline
幾年前,哥哥給我長途電話,原來他的女友在海外宣教不幸身亡:接到電話後我不知所措,內子提醒我前往探望.看見他時,我沒有特別的話跟他說,只是在旁感受他的哀傷.另一回,在我事奉的教會中,一對熱心的老年夫婦,得悉兒子心臟病發而死:這兒子已四十歲,還沒有結婚,對父母非常孝順.這事對老夫婦的打擊可想而知.當時,我以牧者身份前去探望,我無法說出一句話:只有坐在床邊與他們一起流淚,足有兩小時.過後感到自己一無所成,然而,數星期後,這對夫婦感謝我為他們所做的一切:就是在他們哀傷時,感到有愛的同在.在這生命的交流接觸中,哀傷的心靈得到治療.這時代,我們太講求技巧,效率,弟兄姊妹間只有事工的關係,缺乏了生命的分享,不能一起分擔憂愁,分享喜樂.耶穌要我們彼此相愛,最少要做到生命的交流.我們應珍惜或製造相聚的機會,使我們的生命能彼此植根,互相交流.
\newline
\newline
(三)愛的實踐包括主動的表示 :我們應主動去愛他人,不是等待別人的愛,才付出自己的愛；當我們主動去愛時,也能激勵別人將愛擺上.
\newline
\newline
今天,我們跟從主的人有何標誌,使人看出我們是主的門徒呢?耶穌吩咐我們要有愛的標誌,是毫無保留,彼此相愛的圖像.猶太拉比有個動人的故事,話說一對兄弟以務農為生,弟弟有妻兒,哥哥則未結婚.收成時,哥哥想起弟弟有家室之累,自已則單身一人,小量穀種已夠.因此每晚 一袋穀物,拋在弟弟的穀倉內.另一方面,弟弟看到自己榖倉儲滿糧食；想到自己年老時有兒女照顧,哥哥則不同,他應該要有更多的儲蓄.因此,弟弟每晚都背一袋榖物到哥哥的倉庫裡.一晚復一晚,兩兄弟都發現,自己倉的儲糧都沒有減少.一天晚上,他們在禾場上相遇,彼此都背一袋榖倉物,才恍然大悟；兄弟倆人相擁而哭,為著彼此的愛心而受感動.神就在那地方建立自己的殿.雖然這只是個相傳的故事,但卻表示出弟兄姊妹彼此相愛,在神眼中是極美麗的圖像.
\newline
\newline
耶穌說,我給你們一條新命令,就是要你們彼此相愛；我怎樣愛你們,你們也應怎樣彼此相愛.
\newline
\newline


\newpage

\section{第59屆~港九培靈會~許道良~第七講}
\label{sec:1245}
\textbf{肩負使命的客旅}
\newline
\newline
連結: \href{https://www.hkbibleconference.org/session-message/view/1245}{\texttt{https://www.hkbibleconference.org/session-message/view/1245}}
\newline
\newline
\hyperref[sec:1243]{< < < PREV SERMON < < <}
~
\hyperref[sec:toc]{[返目錄]}
~
\hyperref[sec:1249]{> > > NEXT SERMON > > >}
\newline
\newline
基督徒與非基督徒有何分別呢?別人是過路人,是客旅,是寄居的,我們也是.神既拯救了我們,又為何要把我們留於苦難的世界呢?很多信徒在水深火熱時向神呼求:「主啊!你何時再來?」我相信神沒有錯,神不會故意讓我們經歷苦難；苦難本身沒有意義,除非為了更高的目標去承擔苦難.因為神仍深深愛著被罪惡綑綁的罪人,希望他們早日脫離災難.表面上信徒與非信徒沒有分別,我們吃喝嫁娶,生兒育女,發展事業,就業.......唯一不同的是在星期日做禮拜.別人看見我們的分別嗎?可惜,信徒在社會中是一群隱藏的人；自稱是光,是鹽,卻不能在社會上起作用.我們生存的目標及意義,與世人有何分別呢?最高的理想不過是充滿我們的禮拜堂.然而,主的心意是要我們去充滿整個世界.主耶穌說:「父怎樣差了我,我也照樣差遣你們到世上去.」因此,我們在世上不是旅遊,視察,觀光,而是要完成神的托負──一個見證主的托負.
\newline
\newline
誰是肩負使命,見證主的人呢?14節耶穌說這些門徒不屬於這世界,現今的世界被罪惡滲透,仍在惡者掌握中,是抗拒神旨意的勢力.這世界有自己的價值,方向及作風,信徒不屬於這世界,因此在世界上不受人歡迎.正如主來到世上,遭遇到很大的抗拒,因為光來到世界,成為黑暗的挑戰.故此,世界憎恨我們,不要以為詫異,世界的價值觀建立在短暫的物質,屬主之人的價值觀是與永恆掛勾；有聖潔標準,不與罪人同流合污.正如彼得說我們是寄居的.帶著使命進入世界,應提醒自己不迎合世界的潮流.世界若因我們的公義,聖潔及愛心而恨我們,則不需為自己辯護；但若因我們對苦難不關心而恨惡我們,則應自行省察.若為義而受逼迫,便是有福的人.
\newline
\newline
我們不屬於這世界,是否就屬於自己呢?(十七6-9)說我們是屬神的人,神將我們賜與基督,因此不能再以自己的喜好而定人生的目標.我們與世人的分別,在於我們是被拯救,被赦免的人.因此每個屬主的都應當是一個見證人.只是今天出現一個可悲的現象,許多人將見證職業化,只出錢請人來為我們作見證.故此只有一小撮人是專業見證人；這樣,我們在世界的影響也就極之微弱.作見證是每個蒙恩人的責任,這呼籲是主向每個屬祂的人發出的.一方面,這是個報恩的行動.另一方面,這是唯一能影響整個世界的方法.若將見證縮窄在一小撮人身上；那便會限制我們活動的範圍.在極權的國家,作見證很受政權約束；不是任何時間或每個人都可作見證.然而,今天我們有自由超越地域,文化為主作見證,如果每個基督徒醫生,律師,校長,主婦,商家......都將見證帶入自己的圈子,則在整個世界會產生鉅大的影響力.
\newline
\newline
經文中也讓我們明白見證的先決條件.15節「我不求你叫他們離開世界,只求你保守他們脫離那惡者.」我們知道自己正處於敵人的陣地,有自身的軟弱及限制.因此在大祭司的禱文中,耶穌求神保守這些見證人脫離惡者的勢力,免受其騷擾而影響他們的使命.以弗所書形容基督徒的生活就是一埸爭戰,與空中靈界的勢力爭戰.若單憑我們有限的智慧及力量,則難以取勝.若果你對信仰認真的話,便會感受敵人的攻擊,直至無法抵擋,憑己力去應付這些引誘,試探,便會跌倒,灰心,對生命,對神起了懷疑.倘若沒有神的保守,相信大部份人都不會在這裡了.耶穌向神祈求,保守我們脫離那惡者,耶穌很認識那敵人,祂親身經過試探.要勝過試探並不簡單；若沒有神的保守,我們會跌倒；多次的跌倒,使我們慚愧!會放棄為主作見證,對自己,對神失去信心.不錯,撒旦的攻擊很凌厲,但要記得耶穌的禱告,我們有神的保守；祂是我們永遠的保障,堅固的堡壘,所有權柄都在祂手裡,我們可放膽到世界為祂作見證.
\newline
\newline
除了有神的保守,還要有真理的裝備.主說:「我要將你的道賜給他們,求你用真理使他們成聖.」真理的作用不在增加知識,乃是塑造一個合神心意的新生命；使生命有其抵抗力.真理在我們生命裡不斷更新,使我們達到神的標準.今天神學上常探討真理本色化(contextualization)的問題,這只是第一步；最重要是將真理融化於自己生命裡,這便是(actualization).耶穌以真理叫人成聖,也就是說真理經已融化在我們生命裡,每個信徒都應積極追求這真理,以應用於生命上.在香港,聽道的恩賜很明顯,然而,滿腦子的真理,是否能令我們成聖呢?若不能的話,我們未有資格成為見證人.耶穌希望透過真理,建立一個穩固的生命.這生命本身就是有力的見證,也是個信息.林後三章說基督徒的生命是基督的荐信,由聖靈而寫的信,當別人看到我們的生命,就如見到基督的形像一樣.
\newline
\newline
另一方面,經文中也提到我們見證的場所.耶穌說:「我不求你叫他們離開世界,你怎樣差我到世上,我也照樣差他們到世上.」整個世界是我們見證的場所.一向以來,我們以為見證只是圍內人的事,所受的裝備也只為教會的事工,但這導致一個誤會,以為教會之內的是事奉；教會以外是私人生活範圍；以為教會之內的事奉,便代替了教會外的見證.神的心意是要我們進入世界作祂的見證人,教會的事奉只是個起點.聽聞美國一個油田,有大量機器日夜操作,將油泵上地面；鄰居問油田主人,為何只生產油而不出賣呢?油田的主人則答說,泵上來的油也不夠機器的消耗,那有剩餘的油出售呢!教會的情形是否一樣呢?神的心意是要我們有剩餘的油給別人.教會的團契生活是使人重新得力,生命更新,更有力量進入世界.雖然世界的科枝進步,令人以為不再需要神；然而,地鐵站看到的,是一張張無表情的臉孔；背後都有痛苦的經歷,虛空的心境.多少人在黑暗裡等待光,更多人渴望得到真正的愛；主就是要差我們進入這世界中.因此,我們居住的地方,我們的職業,學校都是我們的燈臺；神要我們在那裡燃起見證的火炬.
\newline
\newline
感謝神!今天香港信徒都有優良的職業水準,有好醫生,好校長,好教師,但更應達到好的基督徒醫生,教師,校長,在職業上流出基督的馨香.今天許多不同行業都有團契,希望透過團契生活彼此鼓勵,提醒.不同行業都有其獨特的困難.各人把心願擺在一起,共同商討可行的見證途徑.台灣的「的士團契」便很有新意,他們利用播放錄音帶來與乘客談道,這不是每個行業都能做到.每個崗位都有自己的挑戰,神要我們在自己獨特的崗位上,找一個適當的方法見證神的榮美.主耶穌要我們進入人群中為祂作見證,但我們卻只拉人進來,沒有膽量出去；教會的佈道會中,往往只有百分之十是未信的人.當我們清楚自己在世界的目的,便應校正人生的目標,在我們活動的範圍見證神.
\newline
\newline
接著我們看見證的方式.18節「你怎樣差我到世上,我也照樣差他們到世上.」耶穌是以道成肉身的形式出現,過人的生活；以表明神榮耀.因此,在耶穌被差派的模式中,最少有兩個層面:
\newline
\newline
首先是具體的行動,香港人有一句俗語:「講耶穌」,意即只有口講而沒有行動,因此世人以此話來取笑我們.我們不但講,還應活出基督；具體行動可維護公義,抵擋罪惡的蔓延.就如耶穌潔淨聖殿,不單講理論,也用行動表示.祂又用行動為人驅魔,彰顯神的榮耀；當信仰真的影響我們的生命時,我們整個生命的行動,就都是明確的見證了.為真理,公義,應有具體的行動啊!
\newline
\newline
另外,就是愛的關懷.耶穌在世上表現出一個愛的生命,心中充滿憐憫；這都是愛的自然流露.基督徒謙卑的服事及愛心的表達,不是希望他人回教會,這應該由對方自己決定.我們不可存著動機做這事.因為未信者眼睛是雪亮的；愛若另有目的,他們是不會接納的.他們甚至嘲笑我們利用手段,動機不良.然而有愛的關懷,是因基督的愛先充滿我們而有的自然流露.別人也能感受到我們的關懷及憐憫.
\newline
\newline
除此之外,我們也需要一個清楚的宣講,二者互相配搭.一方面有行動的表示,另一方面有清楚的解釋.當神帶領以色列人出埃及時,在行動背後也有解釋；使人看出是神的作為.舊約時常提醒人要記念神的作為,神的律法；兩方面的配搭,成為清楚的信息.我們宣講時,不應以福音為恥,跟人談甚麼事情都容易,就是談福音最困難!但若不清楚地宣講,別人便不明白我們生命的意義及生命的改變,很多基督徒以不明顯地表達信仰而自豪,福音預工是要做,但應帶入撒種的工作.鬆土的工作只是暫時性,最終要做收成的工作.耶穌宣講天國的信息時用了各種智慧,使他人容易明白.我們講道理時常以為愈深奧便愈有涵養或學問,使簡單的道變得複雜；耶穌講道擅於就地取材,天上飛鳥,地上花朵,都能把深奧道理表達出來.
\newline
\newline
我們在世上肩負重要的使命,要將神的榮美表達,將基督全備的救恩宣揚.教會不過是個訓練所,退修營,為裝備我們滲透整個社會.雖然我們見證的性質,工作崗位可能不同；但我們都是主的見證人.
\newline
\newline
一個人在海峽中救了一些沉船的水手,他發現該海峽是危險地帶,便築了一座燈塔.後來,有些人受這人的勇敢及他的志願所感動；加入了志願拯溺的工作,很快便組織了拯溺會.接著還開設拯溺班,拯溺研討會,但卻慢慢忽略了拯溺工作.後期更將會所易名為拯溺俱樂部；以後有人問起為何叫拯溺俱樂部,會員便解釋這是歷史包袱；不知為何以前用拯溺這詞.但現在已沒有拯溺工作,會員可坐在會所裡,看見船隻向下沉,也感覺與自己沒有關係.船在沉,但拯溺俱樂部的會員仍在享受自己的團契.這不就是今天的教會嗎?開始時為拯溺,今天卻變了俱樂部.不錯,我們有甜蜜的交通及團契,但這是否神的心意呢?弟兄姊妹,你們的燈台還有見證的嗎?你們的鹽是否已失了味?
\newline
\newline


\newpage

\section{第59屆~港九培靈會~許道良~第八講}
\label{sec:1249}
\textbf{愛的十字架}
\newline
\newline
連結: \href{https://www.hkbibleconference.org/session-message/view/1249}{\texttt{https://www.hkbibleconference.org/session-message/view/1249}}
\newline
\newline
\hyperref[sec:1245]{< < < PREV SERMON < < <}
~
\hyperref[sec:toc]{[返目錄]}
~
\hyperref[sec:1251]{> > > NEXT SERMON > > >}
\newline
\newline
馬太福音有記載耶穌在加利利海邊,呼召彼得的情況,當時彼得不加考慮便放下漁具跟從主.我想他沒料到自己要走一條崎嶇的路,其中有勝利的歡笑,也有失敗的痛哭.彼得曾在該撒利亞腓立比,宣佈基督教最重要的教義:耶穌就是基督,是永生神的兒子.他曾在山上看見主的變像,這是超越屬世的限制,進入榮耀的經歷.他曾與主經歷豐盛,受人抬舉；也曾經歷貧乏,受人逼迫.我相信彼得不會忘記他在客西馬尼園的軟弱,還有最痛心的是,在主最需要他支持時,竟然三次不認主.後來,基督受死,埋葬且復活了；而彼得及其他門徒仍驚魂未定,非常迷惘.雖然他們看見並摸過復活的主,卻仍未得著能力；直至聖靈降臨,他們的生命才真正地改變.彼得就是一個如此失敗的人.然而,他卻是我們跟從主的人一個真實的寫照.
\newline
\newline
經文中記載復活的主,如何與門徒恢復那破裂的關係；耶穌主動為他們預備早餐,這是與最後晚餐互相呼應的筵席.晚餐後門徒失敗離棄主,現在主則預備另一個筵席；以彌補他們的過失.反之,門徒卻為此而產生陌生感,非常尷尬；因被他們離棄的主,竟再次邀請他們進餐.但主看重的,不是過去的失敗,乃是放在前面的使命及托付.因此,祂要清楚門徒們有否一個基本條件──愛主的心志,成為他們事奉的動機.事實上,撫心自問,我們事奉神的動機非常複雜；在愛主之外,還有其他動機.以為是為主受苦,其實只是一種不正常的殉道心理.我年輕時,決心要被主使用,想像自己淒慘地在深山野嶺孤單獨行；這是我事奉的最高理想.卻沒有想到在深山野嶺,那有事奉的對象.只以為事奉神愈孤單,愈飢餓愈好.有受苦的心志並非是主動找尋痛苦；乃是說:我願意為主的緣故,接受別人對我的處理.另一方面,我們事奉神,可能是發覺社會上,無法找到更好的職業；做傳道人有聽眾,得人尊重.然而,沒有屬靈的質素,便走上事奉的路,終會被掏汰.
\newline
\newline
因此,耶穌要查問彼得的心志.那時候,彼得對基督的愛已有深入的了解；他看見主毫無保留地為他死,在主的身上經歷了饒恕,醫治及接納.在這特殊的時刻,主的呼召臨到彼得.雖然他那時正與門徒們重操故業,但耶穌卻要彼得作一個抉擇.事實上,跟從主永遠是一個選擇,甚至奉獻了若干年,仍常在比較的情況下作出選擇.假若跟從主,是在無可奈何及被迫下而決定；例如升學機會,就業機會十分渺茫才跟從主.這種跟從實在值得懷疑.「你愛我比這些更深麼?」這是個需要時常更新的選擇.耶穌要彼得明白這決定的含意,因此祂三次查問彼得的心志；是要他明白這個問題的重要性,使他記著這是事奉最基本的動機,讓他留下一個不可磨滅的印象.
\newline
\newline
這段經文,在原文有兩個不同的「愛」字,因此很多人認為有不同的意思.他們認為耶穌先用較高標準的「愛」來問彼得,但因為彼得做不到,故才改用另一個較低層次的愛提問彼得.我則不同意此種解釋,因為在約翰福音內,兩個愛字(Agape)常交替使用,都表達同一意思,並沒有顯著的分別.可參考(約十四23；十七23；十六27)這些經文提及神對世人的愛,其中兩個「愛」字都有使用.故此,我認為這段經文的重點不在於兩個「愛」字的分別,而在於這問題的重複性,耶穌要彼得清楚自己事奉的動機,否則不能持守未來艱苦的事奉.
\newline
\newline
對於耶穌的提問,彼得說:「主啊你知道我愛你.」雖然他曾否認主,令人失望,但他知道主清楚他有痛悔的心；故此願意將過往生命中的碎片交給主,再次決意跟從主.這決定並非保證以後不再失敗,只是神看重我們有真誠的心志.這也是我們與非信徒的分別；基督徒會軟弱做錯事,不同的是；信主的人有一個目標,心志.雖會跌倒,但有痛悔的心；能再次起來跟從主,走屬天的路.這心志成為我們人生的目標,並校正我們日常一切的行為.此外,愛不是一句說話,而是一個生命,需要從行動上表達出來.耶穌不但要求彼得有愛的心志,還要他有愛的承擔；因此主對彼得說:「你餵養我的羊.」耶穌立刻將工作交托彼得,將祂心愛的羊群托予彼得照顧.耶穌深知這工作的艱難,因為每隻羊都有強烈的自我意志,然而都是主所愛的對象.祂盼望這群羊能安睡於羊圈,不受豺狼猛獸的威脅.因此祂不會隨便將這工作交給人.因為它不是一種職業,而需要有愛的承擔,因著愛主的緣故,也愛屬主的人.在這方面,摩西是一個好例子,雖然百姓多次埋怨他,又想用石頭砸死他,但摩西對他們的關懷在禱告中表露無遺；甚至願意為百姓代受神的懲罰.因此神將牧養的職責交給摩西,這就是愛的承擔；倘若我們中間有牧者,對羊沒有憐憫的心賜及愛的關懷,就應當在神面前好好省察.否則,不但無法發揮牧者的功能,同時也會失去當初事奉的衝勁.
\newline
\newline
主將愛的承擔交給彼得後,便講出愛的代價.這代價是事奉神的人不愛聽的.我們常常關心報酬的問題,但耶穌呼召人工作時,沒有說明會回報甚麼；卻指出對方要付上甚麼代價.我相信這就是成功與失敗牧者的關鍵,耶穌讓彼得知道他要付出很大的代價.「年少時候,要束上帶子隨意往來.」若跟從主,這年輕的自由便受到約束,因為年老時要伸出雙手受到綑縛,失去自由.然而,這卻是榮耀神必經之路.越戰時有一位隨軍牧師,隨軍到前線照顧他們的靈性；也因此救出許多垂死的士兵.他在寫信給家人時說:「我能在人生的戰場上作主的僕人,實在是榮幸.」我們身為牧者的知道羊的需要嗎?我們願付上自由的代價,去認識我們的羊嗎?十幾年前我到香港時,發現教會中大多是青年人,心裡實在興奮；但是,今天在教會又發現大部份是青年人,十多年前的青年人,今天在哪裡呢?身為牧者的,知道他們是從哪個破口流失了呢?
\newline
\newline
相傳在尼羅王時代,基督徒受到逼害而逃離羅馬城.當時彼得也在逃難潮中,主向他顯現說:「彼得,彼得,羅馬城內還有一個空的十字架.」彼得醒悟過來,慚愧地回到羅馬城,最後被捉拿,要受十字架之刑；但他感到不配像耶穌那樣釘十字架；便要求倒釘十字架.這雖是個傳說,但作者寫約翰福音時已知彼得的命運.耶穌說:「你年老的時候要伸出手來,別人要把你束上,帶你到不願意去的地方.」彼得放下一切跟從主就得到這些,最後以死來榮耀神,正如一粒落在地裡死了的麥子.保羅的經歷也一樣,林後六章可見他的掙扎,保羅有很多需要,不但被人忽略,還要去照顧很多有需要的羊群；他醫治很多創傷的心靈,但自己心靈的創傷卻無人問津；這就是愛的代價.基督徒所講的愛,不是一朵玫瑰花,而是一個十字架.
\newline
\newline
有一回,一位年老的美國宣教士,從非洲回到自己的祖國,同船有剛狩獵回來的羅斯福總統；岸上擠湧著大批人群及儀仗隊,熱鬧非常,在歡迎總統回來.當總統離去後,宣教士提著行李,獨自站在甲板上觀看,便開始流淚,他說:「父神啊,我遵守你的呼召,到非洲偏僻之地宣教.我在那裡埋葬了我的妻子,兒女,現在獨自回到自己的家裡,卻沒有一人來歡迎我.」那時,他彷彿聽見有聲音說:「我兒,你還沒有回到家裡.」傳教士將自己一生擺上,得到甚麼呢?一無所有；然而神卻透過這種方式來榮耀自己的名,這就是愛的代價了.今天,許多服事神的僕人付上自己的一切,他們的代價可能在今世沒有人明白,要到來世榮耀的國度才知曉.二千年來許多牧者忘記了自我,十字架的道路真的染滿血跡.他們攀越山嶺,走遍大城小鎮；尋找迷失的羊,為我們留下了佳美的腳蹤.然而在這行列裡也有不少人失落,離棄了主,這些缺口直至今天仍需要人填補.
\newline
\newline
公元四世紀,東羅馬國王下令所有人民,向當地的神靈獻祭.當時,黑海南面城裡有一支軍隊；有四十位基督徒士兵,一致抗拒王命.於是他們被捕,並判死刑；要在寒冷的天氣,赤身走向湖心,直至凍死.但是,這四十個士兵說:「你們可擁有我們的服裝,我們的身體,但我們仍決心追隨基督.」一個寒冷的晚上,死刑就執行；他們赤身步向湖中,一面行,一面唱:「基督的四十勇士,你既將生命賜予我們,我們永不離開你.」很有勇氣地步向湖中,聲音逐漸微弱.深夜時,有一人爬出來,走向溫暖的浴室,口裡慚愧地說:「基督的三十九勇士.」但因溫差太大就倒斃了.這時,有一人腳脫下制服,走向湖中心,邊走邊說:「基督的四十勇士.」第二天,在湖中撈起四十具屍體,其中有一個是獄長的屍體,他在監守那四十個士兵時,自己的生命也得到改變,這個缺口已由獄長補上了.
\newline
\newline
今天香港還有很多缺口,很多羊群無人照顧.一九九七的問題,不應是我們的恐懼,反之,應視它為一個挑戰.現在香港有八百間教會,我們可否在神面前立定志願,使香港在一九九七年時,有一千九百九十九間教會呢?誰說這一代的傳道人不如上一代呢?就讓我們開始去更正這種說法；我們奉獻給神是個真誠的奉獻.主在廿世紀的時代,仍然呼召尋找彼得,一個合乎祂心意的彼得.
\newline
\newline


\newpage

\section{第59屆~港九培靈會~金新宇~第一講}
\label{sec:1251}
\textbf{神榮耀的統治}
\newline
\newline
連結: \href{https://www.hkbibleconference.org/session-message/view/1251}{\texttt{https://www.hkbibleconference.org/session-message/view/1251}}
\newline
\newline
\hyperref[sec:1249]{< < < PREV SERMON < < <}
~
\hyperref[sec:toc]{[返目錄]}
~
\hyperref[sec:1254]{> > > NEXT SERMON > > >}
\newline
\newline
我們知道神的國度的主要意義,是神掌權,神要作王；而神掌權的領域──在舊約是祂的子民；在新約是在信徒的心中.神的國度是永存的,是從永遠到永遠；而神為憐憫世人,就藉著祂獨生子耶穌基督的救贖工作,讓神的國進入人間.今天我們看見神的國,完全是耶穌救贖工作的果效.
\newline
\newline
從詩篇一四五篇的內容,給我們清楚了解神的統治.我們看見神是一位榮耀的君王,祂的國度是榮耀的國度.這篇詩是大衛王在詩篇中最後寫的一篇,也是他唯一的讚美詩.大衛如何讚美神,每一位信徒也必須讚美神,並且要將如何讚美成為我們信仰生命中的首位.我們仔細研讀這篇詩時,就知道是一篇敬拜用的詩.猶太人十分看重這篇字母詩,在一年的崇拜生活中,最少有三次唸這篇詩,每次都在早晚唸誦敬拜讚美神.初期教會的猶太信徒,在復活節後第七個主日的聖靈降臨節──(即五旬節),必定用這篇詩歌頌神；所以這首詩實在對敬畏神的人非常重要.我們不妨從幾個角度思想這篇詩的重要性:
\newline
\newline
(一)神是偉大的神,本篇第3節說:「耶和華本為大,該受大讚美,其大無法測度.」我們的神是王,是大神；而從1-6節更看見:祂信實可靠,是我們的保障；祂是神,祂是王,是宇宙的統治者,祂掌管著明天.
\newline
\newline
(二)神是滿有恩惠,慈愛的神,一個大有能力的君王不一定愛民如子；而我們這位偉大的神,卻很愛我們.第7-10,14-20節,大衛說明神何等愛我們；所以大衛在詩篇開始時,就以:「我的神,我的王,我尊崇你,我要永永遠遠稱頌你的名」向神讚美；而在結束時,則以:「我的口要說出讚美耶和華的話,惟願凡有血氣的,都永永遠遠稱頌祂的聖名.」來稱頌神.神是滿有恩惠,慈愛,不像世上的君王,雖有能力卻沒有慈愛,又或是有慈愛卻沒有能力.神的愛是透過祂愛子的死,表明出來,讓渺小的人類得到拯救.
\newline
\newline
(三)神榮耀的國度,神榮耀的統治,在這首詩表露無遺!大衛是君王,但他看見他的神,是萬王之王,萬主之主；祂是君王,卻服在神的王權之下.所以我們應看見我們的神是偉大的君王,而祂是永恆國度的大君王!大衛亦看見自己不過是人!他在(詩篇一四四4)說:「人好像一口氣,他的年日,如同影兒快快過去.」他看見自己的光景,比較起神永遠的國度,內心不禁湧出無限的讚美!
\newline
\newline
不單是猶太人時常樂於背誦,他在(詩篇一四五4)強調:「這代要對那代頌讚你的作為,也要傳揚你的大能.」大衛自己讚美神,也呼喚歷世歷代的子民都要如此讚美；神的教會也要教導信徒讚美神,代代相傳要讚美神,這是我們對神應有的敬拜.
\newline
\newline
我們的神是偉大無比的.(詩九十五3):「因耶和華為大神,為大王,超乎萬神之上.」寫詩的人亦有同樣的看法.地上每一個角落,也在祂統治之下；每一個時代,神都是在作王.
\newline
\newline
祂的作為,祂的大能,祂的威嚴和尊榮,祂的大德.讓我們再深入思想這幾方面:
\newline
\newline
(一)神創造的偉大,宇宙的偉大,在詩篇第八篇被大衛刻劃出來:「我觀看你指頭所造的天,並你所陳設的月亮星宿,便說──人算甚麼,你竟顧念祂!」他發現神的創造,不是一點點的.當時還未有望遠鏡,只憑眼看；今天用長程望遠鏡看時,這個宇宙實在比大衛所看的更大,而且我們亦未能完全看見所有.今天我們單從太陽系中,就發現有一百億銀河!從地上看每個銀河時,有多少星呢?平均有一千億,這樣我們實在看見神的偉大!
\newline
\newline
(二)神愛的偉大,從創世記到啟示錄,我們都看見神的作為,把祂的子民從埃及藉羔羊拯救出來；看哪!我們逾越節的羔羊基督,已經被殺獻祭.我們看見祂偉大的愛,祂偉大的作為；而聖經清楚記載祂必再來；世上有最後一次的大戰,主耶穌騎在白馬上,擊殺列國,勝而又勝!從啟示錄十九章我們看見神最後的勝利,一切敵擋神的權勢,罪惡,都完全消滅,公義要在我們中間彰顯；我們要看見新天新地,有義在其中.祂是我們的神,是大神,祂的能力是超過我們所能想像的.
\newline
\newline
我年紀已很大,經歷神的偉大比青年人會多一點,又蒙神垂聽禱告,我經歷神是又真又活的.「在神豈有難成的事嗎!」這是亞伯拉罕的經歷；保羅在羅馬書四17說:「亞伯拉罕所信的,是那叫死人復活,使無變為有的神.」世上有能力的君王可能是暴君,歷史上有很多這樣的例子；但我們的神卻是滿有慈愛的神.
\newline
\newline
(詩一四五7):「他們記念你的大恩,就要傳出來,並要歌唱你的公義.」這真是寶貴,神將祂的大恩傾倒出來,讓人以讚美稱頌作回應.神的大恩與公義,都與神救贖的作為有關.保羅說:「神既不愛惜祂的兒子,為我們眾人捨了,豈不把萬物和祂一同白白的賜給我們麼?」(羅八32)這是保羅從心底說出讚美的話.聖經在舊約中的「公義」,有「信實」的含意,表示神的應許信實可靠；祂與亞伯拉罕立約,與大衛立約,現今在基督裡與我們立約.
\newline
\newline
(詩一四五17)「耶和華在祂一切所行的,無不公義,在祂一切所作的,都有慈愛.」第8節:「耶和華有恩惠,有憐憫,不輕易發怒,大有慈愛.」這是出埃及時神的表現；在人不配得,而神卻主動愛我們.我們沒有誰能配得神的恩典,但祂卻臨到我們,給我們憐憫──憐憫是神給人最深切的關懷.祂不輕易或永遠的發怒,祂乃是大有慈愛.
\newline
\newline
「神愛世人」,這樣的愛是無可測度和比較的.在舊約先知約拿的經歷中,有一條蟲子咬死了一棵蓖麻,將他在蓖麻下遮蔭的享受取去,他就發怒,神卻對他說:「尼尼微這大城,其中不能分辨左手右手的有十二萬多人,並有許多牲畜.我豈能不愛惜呢?」(拿四11)神愛的範圍廣闊,實在偉大.(詩一四五9):「耶和華善待萬民,祂的慈悲,覆庇祂一切所造的.」祂不只愛世人,更愛一切祂所造成的.現在世界敗壞的光景──空氣染污,海水染污,噪音的染污,酸雨的出現,這些事都不是神所喜悅的事；祂起初讓我們住在空氣清新的大自然中,所以環境的保護,是神所喜悅的；而因著人的犯罪,萬物都服在虛空之下,受敗壞的轄制,但是神仍然愛護祂所創造的,到主耶穌再來作王的時候,萬物都要復興(羅八20-25).大衛在(一四五9)說:「耶和華阿,你一切所造的,都要稱謝你.」而(詩篇六五13):「草場以羊群為衣,谷中也滿了五榖,這一切都歡呼歌唱.」這些沒有生命的東西,都來讚美神；我們這位偉大的神,看顧著祂整個造成的宇宙.
\newline
\newline
(三)我們要讚美神.大衛在(一四五14-17)表示他是不止息地讚美,要永遠讚美──神的公義,神的信實,神的扶持,神的供給,常與人親近.我們應該接受神的王權,順服祂,敬畏祂,愛慕祂,求告祂和感謝讚美祂.
\newline
\newline
在末了大衛又讚美神的國度.地上的國度是有領域的；但神的國卻不須要在地上,而全世界都是神的.雖然最後很多人背叛神!神仍在各處掌權,神容忍他們,繼續給機會他們悔改,神的心意本是要萬人得救,明白真道；這是你與我的責任,要把福音傳遍天下.今天神國度的統治在各基督徒的心中,所以人看不見.今天神的國度是屬靈的,隱藏的；但有一天當主再來的時候,大家都要看見神的國度榮耀的臨到.大衛在(代上二九10-12)對以色列會眾說稱頌神的話:「耶和華阿,尊大,能力,榮耀,強勝,威嚴,都是你的.......我們的神阿,現在我們稱謝你,讚美你榮耀的名.」
\newline
\newline
那時大衛是王者之尊,他的兒子所羅門接續他為王.他看見神的恩典,內心發出無限的讚美,這是一首偉大的詩歌.當我聯想起主禱文的話:「國度,權柄,榮耀,全是你的,直到永永遠遠.」國度!永遠的國是神的!神實在偉大:「大能,恩惠,慈愛,憐憫和公義.」實在顯出祂的榮耀.所以大衛在(一四五11-13)說:「傳說你國的榮耀,談論你的大能......存到萬代.」我們要說出神國度的榮耀,更願祂早日降臨.今天我們信主的人,經歷神的國在我們心中,接受神的統治；洪水泛濫的時候,耶和華坐著為王.地上不論發生任何事情,世界充滿各種痛苦和問題；我們信主的人要在夜間歌唱,唱信心的歌,看見神的權能.
\newline
\newline
世人的罪叫人失去榮耀和尊貴,祇有基督再來,神國降臨,一切的問題方能解決；我們的神不會失敗,祂是大神.我們要盼望等候,教會在地上要禱告神國的降臨,要傳福音,讓神的國臨到許多沒有主的人.
\newline
\newline
大衛說:「我的神我的王阿,我要尊崇你,我要永永遠遠稱頌你的名.」感謝主,讓我們成為讚美的人；讓我們看見所信的是誰!我們要天天稱頌你,也要永永遠遠讚美你的名!
\newline
\newline


\newpage

\section{第59屆~港九培靈會~金新宇~第二講}
\label{sec:1254}
\textbf{耶穌基督是王}
\newline
\newline
連結: \href{https://www.hkbibleconference.org/session-message/view/1254}{\texttt{https://www.hkbibleconference.org/session-message/view/1254}}
\newline
\newline
\hyperref[sec:1251]{< < < PREV SERMON < < <}
~
\hyperref[sec:toc]{[返目錄]}
~
\hyperref[sec:1256]{> > > NEXT SERMON > > >}
\newline
\newline
「你壽數滿足,與你列祖同睡的時候,我必使你的後裔接續你的位,我也必堅定他的國.他必為我的名建造殿宇,我必堅定他的國位,直到永遠.」這個應許是最後在耶穌基督身上成就,祂要成為永遠的王.
\newline
\newline
(腓二5-11)講及主耶穌第一次來的身份和使命；保羅藉此勉勵腓立比教會的信徒.「要以基督耶穌的心為心,以祂的態度為態度,以祂的心思為心思.祂虛己成為奴僕的形像,成為人的樣式,更自己卑微,存心順服,死在十字架上.」這是主耶穌第一次來的光景；成為卑微的拿撒勒人耶穌,預備迎接死亡:「所以神將祂升為至高,又賜給祂那超乎萬名之上的名,叫一切在天上的,地上的和地底下的,因耶穌的名無不屈膝,無不口稱耶穌基督為主,使榮耀歸與父神.」這是主耶穌第二次榮耀再來的預告.
\newline
\newline
馬太福音二章講及耶穌被生下來時,東方博士來到耶路撒冷,問及猶太人的王在哪裡?他們看見祂的星特意來拜祂.不錯,祂來是要做王.猶太人等候了數百年,盼望他們的王──彌賽亞來到,拯救他們從仇敵的手中出來,他們不斷地等待.主耶穌來就是彌賽亞來到,祂要作猶太人的王,也是要作我們的王.主耶穌被賣的那一夜,祂在羅馬巡撫彼拉多面前受審.彼拉多問耶穌:「你是猶太人的王嗎?」主回答說:「你說我是王.我為此而生,也是為此來到世間,特為給真理作見證.」(約十八37)主耶穌來是作王,把神的國度帶到人間,並統治我們.祂的國度是屬靈的,是隱藏的,是人所看不見,更是非物質的；正是馬太福音十三章論及天國奧秘的精意.
\newline
\newline
感謝主祂還要再來,那個時候是有能力駕雲而來,在地上建立祂的國度,神榮耀的國度要臨到我們.
\newline
\newline
路加福音十七章裡法利賽人問耶穌:「神的國甚麼時候來到?」主耶穌回答說:「神的國就在你們中間(心裡).」雖然世界的人看不見神的國,但在我們信徒中,祂是我們的王.我們要相信祂,順服祂,愛祂,跟從祂,我們更要把自己奉獻給祂(羅十二1),將身體獻上,成為活祭,去討神喜悅.保羅在哥林多前書又說:「豈不知你們的身子就是聖靈的殿麼?這聖靈是從神而來,住在你們裡頭的,並且你們不是自己的人.因為你們是重價買來的,所以要在你們的身子上榮耀神.」(林前六19-20)我們必須遵行神的旨意,將自己獻給神,讓神在我們生命中作王.耶穌為我們的罪死了,埋葬了又復活了；使徒彼得在五旬節的時候為主作見證:「故此,以色列全家當確實的知道,你們釘在十字架上的這位耶穌,神已經立祂為主,為基督了.(徒二36)祂是主,祂是王,祂是基督,是神差來的彌賽亞.」約翰在啟示錄中亦有記述,「並那誠實作見證的,從死裡首先復活,為世上君王,元首的耶穌基督.(啟一5)」祂實在是我們的元首,是萬王之王,萬主之主.
\newline
\newline
我們從三方面來看耶穌基督是王:
\newline
\newline
(一)從舊約預言看
\newline
\newline
舊約清楚指出祂是要來的彌賽亞,祂要坐在大衛的王位上,從(撒下七12-13)神給大衛的應許,我們就知道主耶穌來是成就這應許.大衛的後裔要永遠作王,並且要一脈相承,而主耶穌正是大衛的子孫.
\newline
\newline
在新約福音書中,時常講及耶穌是大衛的子孫,是表明祂是彌賽亞.馬太第九章講及兩個瞎子跟著耶穌,喊叫說:「大衛的子孫,可憐我們罷!」耶穌就表明彌賽亞的能力,把他們醫治了.詩篇一一○1:「耶和華對我主說,你坐在我的右邊,等我使你仇敵作你的腳凳.」這節聖經在新約至少被引用廿七次,有直接或間接用.我主就是彌賽亞,被所有人接納為王,王的仇敵遲早被征服.不論在詩篇或先知書,都講述基督為王.在路加福音第四章,耶穌來到拿撒勒,在安息日進會堂唸聖經.有人把先知彌賽亞的書交給祂,祂就打開,找到一處寫著說:「主的靈在我的身上,因為祂用膏膏我,叫我傳福音給貧窮的人.差遣我報告被擄的得釋放,瞎眼的得看見,叫那受壓制的得自由,報告神悅納人的禧年.」主耶穌讀完後把書關起來,坐下來對會堂裡的人說:「今天這經應驗在你們耳中了.」我們知道這段經文是從以賽亞書六十一章來的,是指著彌賽亞而說的,他是主前七百多年預言這件事,主穌穌證明關乎彌賽亞的預言應驗在他們中間,特別報告神悅納人的禧年；以賽亞書六十一2:「報告耶和華的恩年,和我們神報仇的日子.」主第一次來不是為報仇而來,不是為審判；而是把恩典帶到人間,把國度帶來.當祂第二次來時,要榮耀臨到,要審判.所以主耶穌來到人間,應驗了彌賽亞的預言.
\newline
\newline
(二)從耶穌的家譜看
\newline
\newline
馬太福音第一章論及耶穌的家譜,也是要證明耶穌是彌賽亞.猶太人注重家譜,將家譜放在公會中；猶太歷史家約瑟夫的自傳中,就把家譜寫上,用以證明他是神的選民,清楚祖先的歷史.同樣地,中國人也喜歡家譜,目的在誇耀祖先的光榮,或是表明身家清白,甚至有些人把家譜改弄,以增加聲望.
\newline
\newline
我們看這裡耶穌的家譜,就發現是王的家譜,祂是神派來的彌賽亞；因為彌賽亞必定是亞伯拉罕和大衛的子孫,這追溯至(創十二1)與亞伯拉罕所立的約和撒下七章與大衛立的約.所以家譜一開始就證明:「亞伯拉罕的後裔,大衛的子孫,耶穌基督的家譜.」馬太講完家譜後,在太一7說:「這樣,從亞伯拉罕到大衛,共有十四代；從遷至巴比倫的時候到基督,又有十四代.」這樣把家譜分成三個十四代；馬太這樣分法也許是表明大衛的字母是DWD拼成,D=4,W=6,D=4,加起來是十四.三個十四代表明主耶穌與大衛有密切的關係,每一個代表明不同的立約和應許,即亞伯拉罕,大衛的約和新約,這三個應許更完全應驗在耶穌基督身上.主就是彌賽亞.
\newline
\newline
家譜的順序是一個生一個,可是到16節說:「雅各生約瑟,就是馬利亞的丈夫.那稱為基督的耶穌,是從馬利亞生的.」清楚表明耶穌不是約瑟生的,而在法律上卻被認可；主耶穌是童貞女馬利亞所生,所以祂沒有原罪,是完全人,而能擔當我們世人的罪.祂成為彌賽亞拯救我們；因為神使無罪的,替我們成為罪,好叫我們在祂裡面成為神的義.
\newline
\newline
(三)從耶穌在地上所行的神蹟奇事看
\newline
\newline
耶穌在四本福音書的工作見證中,證明祂是神,是王.(太四23):「耶穌走遍加利利,在各會堂裡教訓人,傳天國的福音,醫治百姓各樣的病症.」這節聖經說明祂三件主要的工作──教導,傳揚,醫治；這是彌賽亞三重的工作.祂行了很多神蹟奇事,叫風浪平靜.詩篇八十九與一百零七篇說明,祇有神才能叫風浪止息.祂行的神蹟奇事是證明祂有王權.(太十二28):「我若靠著神的靈趕鬼,這就是神的國臨到你們了.」馬太福音更講及神的國度五十多次,在整本福音書更有百多次講及國度,表明耶穌是天國的王,祂有王權去醫病,趕鬼,行使各種神蹟奇事.
\newline
\newline
馬丁路得說:「主耶穌第一次來臨時,祂頭上的荊棘是冠冕,十字架是寶座；當主再來時,天上,地下,一切都屈膝,口稱基督為主.」聖經學者研究神的兒子的名字,約有二百零八個,所以神要給祂萬名之上的名字,是主.舊約的主是指耶和華；主是神,主亦是基督.我們今天盼望主再來,將來就必定實現.(約十二32):「我若從地上被舉起來,就要吸引萬人來歸我.」神要所有的人歸向祂,吸引是有強迫的意思；有一天,主再來時,一切人都要屈膝,無論願意否,都要口稱耶穌為主.我們信主的人甘心快樂稱耶穌為主,我們應有好的生活,讓主耶穌在我們生命中作王,作主,遵行神的旨意,討神喜悅,榮耀主名!
\newline
\newline


\newpage

\section{第59屆~港九培靈會~金新宇~第三講}
\label{sec:1256}
\textbf{先求祂的國 — 國度的重要性}
\newline
\newline
連結: \href{https://www.hkbibleconference.org/session-message/view/1256}{\texttt{https://www.hkbibleconference.org/session-message/view/1256}}
\newline
\newline
\hyperref[sec:1254]{< < < PREV SERMON < < <}
~
\hyperref[sec:toc]{[返目錄]}
~
\hyperref[sec:1257]{> > > NEXT SERMON > > >}
\newline
\newline
主耶穌是我們的救主,是我們生命中的王,我們必須給主在我們生命中居首位.主耶穌說:「你們要先求祂的國」.神的國是神的王權,要管理我們的生命.主耶穌在這段經文中(太六22-34)好像忽視了我們物質上的需要──衣,食,住的供應；其實祂並不是這樣的用意,而是用比較的方法突出「神的國和義」更重要.我們心中所想望的,工作的動機,內裡最深的願望,應該在神的國中.
\newline
\newline
耶穌沒有忽略我們肉體的需要.在(約六)變餅變魚給五千人吃飽的事蹟,充份表現出祂關顧人的整全需要；聖經的其他篇章,例如雅各書就教導:「若是弟兄,或是姐妹,赤身露體,又缺了日用的飲食......卻不給他們身體所需用的,這有甚麼益處呢?」(雅二15-16)表明信徒要在行動上幫助肢體解決生活上的需要.我們要將神的國和義放在生命的首位.英國的劍橋七傑,是內地會差派來華的宣教士,其中有一位是Cassel,他帶往中國的行李上都有兩個字──God First神居首位；他這樣做,真是讓神居首位.所以他在中國有美好的事奉,天父必看顧他和其他的需要.神的手掌管托住萬有,在地上祂亦掌權作王；若人不以神居首位,地上必充滿混亂,苦難.在聯合國開始組織時,決定不將神放在憲法中,世界至今亦持續在混亂鬥爭中.
\newline
\newline
「神的國」對基督徒非常重要,我們不單在情感上追求；更重要在實踐上,在神的工作中有份.有一個傳道人清晨起床時,就這樣禱告:「基督阿!早晨!我愛你.今日你有甚麼特別工作要做呢?我盼望能夠在你的工作上有份.感謝神!阿們!」他就是這樣形容神的國有不同的名稱,馬太福音用──天國；其他的福音書卻用神的國,啟示錄是用──我主和我主基督的國,這些名稱都是一樣,因為馬太福音的對象是寫給猶太人；所以避諱不用「神」直稱神的名字,故此用天國來稱呼,其中稱「天國」有三十四次,而稱「神國」只有四次.在(太四17)說:「天國近了,你們應當悔改.」而在(可一15)卻說:「日期滿了,神的國近了,你們當悔改,信福音.」從這兩段的經文比較,顯然意思是相同的.
\newline
\newline
(一)神的國是我們追求的目標
\newline
\newline
耶穌出來作傳道工作時,教導門徒的重點在──宣揚神的國(路九1-2).祂自己在馬太福音最少兩次說明祂工作的中心,是宣揚神國的福音(太四23,九35).主耶穌基督在地上,不斷宣講神國的福音,四卷福音書有一百多次講及神的國,而講及教會只有三次,這更肯定神的國在福音書的重要性!我們今天必須傳講天國的福音,不要局限福音只講天堂地獄.天國的福音是要人接受基督成為生命的管理者,生命的王,不但將來可得榮耀的盼望；而在今天,就可得平靜,安息和喜樂,享受神的同在.有一年我在澳洲教學,在復活節時很多同學和教授都放假離開學校,我一個卻孤單地在大學的校園個多星期,那時我有機會多些時間讀經,祈禱和讚美神；我深深感受與神同在,十分甜蜜,肉身雖在人間,但心靈卻置身於天堂中,我享受與神同在的快樂!
\newline
\newline
主耶穌說:「你們要先求祂的國和祂的義.這些東西都要加給你們了.」(太六33)我們環顧世界的痛苦和災難,都是因為人棄絕和厭棄神,更使萬物都受到影響,服在虛空之下,受敗壞的轄制(羅八20-21)；人的罪影響遍及萬物.有人說:「這個世界再不能讓人居住!」空氣,水,噪音,酸雨等等的污染問題,都影響我們的居住環境.最近幾年來的愛滋病浪潮,使世界陷入恐慌和危機中,現在據統計有一千萬的愛滋病帶菌者,有些人是無辜被傳染而招致染病,也許是由輸血的緣故.今年與我相識的一位老弟兄要施手術,女兒時常陪伴左右,要為父親輸血以便進行手術,因為她不敢動用血庫的血,以免父親染上愛滋病.這樣的世界,你看可憐不可憐!這個世界實在沒有甚麼盼望!
\newline
\newline
此外還有經濟的問題,和平的問題,戰爭的問題,雖然人不是不努力去解決問題,但努力也是束手無策!每年聯合國都不斷頒發諾貝爾獎金給有貢獻的學者,獎勵他們能對不同的領域,學術或問題上探討解決的方法；可是問題始終存在,人費盡心血也未有關鍵性的突破,將問題圓滿解決.人到底是有限的!唯有神掌權,甚麼問題都可解決!祇有在耶穌基督再來的時候,一切的問題方可解決!所以聖經告訴我們,必須讓神在我們生命中居首位,接受主的王權,我們的生活方能和諧穩妥；無論環境如何惡劣,神與我們同在,給我們安息,賜我們盼望.
\newline
\newline
從以色列史中,我們知道本來是應該由神統治他們,但許多以色列人卻沒有給神作王；士師時期是以色列人最黑暗的時代,是他們失敗的歷史；其中四次講及國家的淫亂,道德的墮落,受到敵人的欺壓.究竟是甚麼原因呢?聖經說:「那時以色列中沒有王,各人任意而行.」(士廿一25)(撒上十二12)表示「其實耶和華你們的神是你們的王.」但以色列人沒有以神為王,他們沒有聽從神的話,結果以色列軟弱.以色列最好的王是神自己.我們今天最理想的王是耶穌基督,我們等候祂來,萬物都要復興,把公義和平帶到人間.(撒上八4-5)以色列的長老見撒母耳說:「......現在求你為我們立一個王治理我們,像列國一樣?」你看可憐嗎?神所揀選的子民沒有讓神作王,沒有求神的國；他們最要緊的是,與其他國家看齊,他們有好的藉口:「撒母耳你年紀老了,你的兒子又不行你的道,我們要立一王.」撒母耳不喜歡；神對撒母耳說:「你只管依從,因為他們不是厭棄你,乃是厭棄我,不要我作他們的王.」(撒上八7)結果以色列人非常可憐,為王的掃羅開始時很謙卑,後來驕傲起來,在吉甲犯罪獻祭,沒有聽撒母耳叫他等候的話,違背先知的話即違背神的話,神要將他的王位賜給更合神心意的大衛.後來大衛為王,統一以色列國,雖然他亦有軟弱失敗；但國家仍富強起來,及至所羅門,國家影響力遍及近東.可惜他後來娶淫亂拜偶像之外邦女子為妃嬪,使國家陷入拜偶像的危機中,他死後國家分為南北國,並且互相殘殺.當國家行耶和華眼中看為善的事情,國家就興旺,無奈以色列和猶大的王始終軟弱犯罪；最後以色列人國破家亡,被擄到各處,過著亡國奴的生活.這是歷史的鑑戒和寫照.以色列人悲慘的遭遇是因他們沒有以耶和華為王.美國開始立國時,國民都敬畏神,一八六四年開始發行兩分錢的銀幣,上面寫著──我們信靠神,那時美國富強起來,因為他們以神為王.今天你我是神的子民,我們必須讓神在我們心中作王,讓耶穌基督成為我們生命的主宰.
\newline
\newline
(二)神的國度是我們生命的優先
\newline
\newline
我們要先求神的國和義,要全心去求,迫切盡心去求,更要以祂的國為一切的優先.我們不應過自我中心的生活,我們的老我必須要死,讓基督在其中作王,不以自己為主.我們要不斷求神的國和義,我們的生活和事奉對象是神自己.(來十一6)說:「人非有信,就不能得神的喜悅；因為到神面前的人,必須信有神,且信他賞賜那尋求他的人.」我們要天天殷勤切切的尋求神,不是在逆境或有難處需要時才去親近祂,我們要有信心並要討神喜悅,我們的對象是神自己.現今世界充滿敵擋神的權勢,基督徒面對這場屬靈的爭戰,必須要與神保持良好的關係,靠著神的大能大力,面對那屬靈氣的惡魔.軍隊在外面打仗,必須與基地保持好的聯繫,方能取得補給,不斷有力量爭戰,基督徒也是一樣,必須倚靠就近神方能得勝,孤軍作戰是不行的.我們要在教會中生活,忠心事奉,討神喜悅,過著聖潔公義的生活,不斷遵行祂的旨意.我們不能以追求世界的物質生活和情慾為滿足；人若賺得全世界,賠上自己的生命,有甚麼益處呢?人還拿甚麼換生命呢?(太十六26)很多宣教士,醫生在落後,設備不足的地區工作,他們從早到晚不斷工作,目的是要滿足神的心意；藉著醫療宣教,見證神的大愛,傳福音搶救失喪的靈魂.他們雖然勞苦,但滿有喜樂,感受神的同在,豐盛的生命在他身上顯露出來.因為他們願意先求神的國和神的義!他們不是尋求自己的滿足!他們如以賽亞先知一樣,看見神的榮耀,並且順服尊主為大,願意對神回應:「我在這裡,請差遣我!」親愛的弟兄姊妹,您願意先求神的國和義嗎?您願意以祂為生命的優先嗎?
\newline
\newline


\newpage

\section{第59屆~港九培靈會~金新宇~第四講}
\label{sec:1257}
\textbf{神的國 — 現在和將來}
\newline
\newline
連結: \href{https://www.hkbibleconference.org/session-message/view/1257}{\texttt{https://www.hkbibleconference.org/session-message/view/1257}}
\newline
\newline
\hyperref[sec:1256]{< < < PREV SERMON < < <}
~
\hyperref[sec:toc]{[返目錄]}
~
\hyperref[sec:1261]{> > > NEXT SERMON > > >}
\newline
\newline
讓我們從馬太福音第十三章來看,國度的現在和將來的光景.將來的國度,在聖經中有預言,雖然詳細的情況我們不知道,但實際是好得無比的；當主再來的時候,榮耀的國度便臨到.可是在今天,神的國度又如何呢?
\newline
\newline
馬太福音第十三章主耶穌講及天國的奧秘,又提及七個比喻,讓我們知道天國用甚麼方式在我們中間；其中在比喻的幾個形容字眼,除了第一個撒種的比喻外,都是如此的──天國好像.......這是很重要的信息,論及有關天國的真理；在舊約時代先知都不知道,而現今主耶穌正清楚的講出來；神的國度在我們心中,是在每一位信徒的心田作王.若我們詳細研究,可從符類福音中了解神的國:(太十三,可四,和路八).最特別是上述三段經文論及第一個比喻時,開始都沒有──「天國好像......」的字眼來形容!這是甚麼原因呢?因為第一個比喻說明天國如何開始在我們中間.我們未信主前神未能在我們心中掌權,我們必須首先得到神的話語,撒種是神的話臨到人的心田,發生福音的功效；我們亦完全順服在神的旨意中生活,天國就自然開始在我們心中.
\newline
\newline
(一)撒種的比喻
\newline
\newline
撒種的比喻除了說明天國在我們心中外,也可應用在其他的場合中；雖然焦點是說明天國的開始,但整本聖經的其他真埋,我們亦應如撒種一般地接受.而這比喻是表明一個大的原則,去接受神的話語.
\newline
\newline
我們必須了解其中說明的四種心田,表示對神話語反應的不同,而我們應該像好的田地,去接納祂,更要明白,遵行神的話.福音派信仰的教會都相信,當人相信接受主,神的國度已到人的心中；可是論及將來國度的觀點,就有不同的看法.而主要都承認天國已經開始在信徒的心中,是屬靈的,非物質的和隱藏的.此外,我們正在等候主再來,與眾聖徒,天使一同來到地上,這是看得見和豐滿的見證.在一九七四年洛桑會議中,德國的神學教授Geerhardus Vos說:「今天的國度可以說是恩典的國度；因為耶穌把恩典帶到我們中間；當主耶穌再來的時候,稱之為榮耀的國度.」其實這是同一個神的國度,而不是稱之為現在和將來的國度；卻從國度的本質來區分,稱為恩典和榮耀的國度.這福音要傳遍天下,然後末期才到──這福音就是恩典的福音,叫信主的人可進到神的國度.不錯,現在正是恩典的國度.當主再來,要審判世人,審判大巴比倫,大淫婦,假先知,敵基督和所有敵擋神的權勢；這樣恩典的門要關閉,然後榮耀的國度臨到.
\newline
\newline
馬太福音十三章用此比喻來解釋.本來比喻是要叫人更明白；但這裡的比喻恰巧相反,是要叫人不明白!所以門徒追問時,主便特別向他們解釋:「天國的奧秘,只叫你們知道,不叫他們知道.凡有的,還要加給他,......也要奪去.」(太十三11-12)為甚麼耶穌要這樣做呢?而祂繼續引用先知以賽亞的話:「你們聽是要聽見,卻不明白......心裡明白,回轉過來,我就醫治他們.」(太十三14-15)很多時這令人費解!原來在舊約中,我們知道以色列人的心剛硬,似乎是麻木了,所以這是神對以色列人的審判,叫他們聽是要聽見,卻不明白；看是要看見,卻不曉得.到耶穌在世工作的時候,傳天國福音,醫病趕鬼,但卻被以色列人反對,拒絕,有些人甚至指責祂是靠鬼王別西卜趕鬼.他們不肯相信,心裡剛硬,實在是到達愚頑的地步；正如以色列人昔日的光景一樣.天國的奧秘當然深奧難明,但由於主耶穌的講解和聖靈的光照,我們現在亦可明白.
\newline
\newline
撒種的比喻表示人對神話語的回應.我們的心田,是頭三個或是第四個?我們不但聽道,更要明白,遵守,這樣會有一百倍或六十倍或三十倍的收成；撒種的人是主自己,種子是神的道,地土是人的心田.
\newline
\newline
(二)稗子的比喻
\newline
\newline
主耶穌說:「天國好像人撒好種在田裡,及至人睡覺的時候,有仇敵來,將稗子撒在麥子裡,就走了.」(太十三24-25)稗子就是惡者之子,撒稗子的仇敵就是魔鬼；撒好種的是人子──主自己,田地是世界,好種就是天國之子.這就是現今的世界有天國的子民和惡者之子,同樣的存在；或者有些應用在教會內,亦有這種情形出現.魔鬼不停地在工作,破壞神國度的工作,我們必須小心警惕!我們是在屬靈的爭戰中.麥子和稗子開始生長時,是無從分別的；及至長大後,才被識別出來,初期教會亦有假先知,假師傅,假使徒和假的弟兄,在歌羅西書,彼得後書和約翰壹書,都有這些虛假見證的事實；把似是而非的道理混在信徒中,目的是要使神的教會軟弱,所以今天的教會或天國,就有這樣的難處,有魔鬼千方百計來破壞.事奉主的僕人應指出異端和錯誤的說法,把真理說明辯白,讓信徒領受清楚,講道的人負很大的責任.(雅三1)「我的弟兄們,不要多人作師傅,因為曉得我們要受更重的判斷.」我們為主傳話的人必須在主面前等候,得著話語去栽培建立教會.馬丁路德形容教會時,是用爭戰的教會,表示教會要在靈界爭戰.教會更是真理的柱石和根基.至於那些稗子,解經家形容是有毒的麥子,而魔鬼時常把毒素引進教會中；叫多人不明白真理.
\newline
\newline
(三)芥菜種的比喻
\newline
\newline
「天國好像一粒芥菜種,有人拿去種在田裡.這原是百種裡的最小的......天上的飛鳥來宿在他的枝上.」(太十三31-32)天國從很小開始,如巴勒斯坦地最微小的芥菜種一樣；長大後卻非同小可,大約有三公尺高；今天神的國亦會如此擴張.主耶穌臨離世主持聖餐時只有十二個門徒,祂升天後在馬可樓上有一百二十人；到五旬節時有三千人信主,後來又加多五千；直到今天,神的國度不斷擴張,教會在各族,各方,各民中被建立起來.雖然魔鬼設法攻擊攔阻,亦無法阻攔國度的擴張.
\newline
\newline
(四)麵酵的比喻
\newline
\newline
「天國好像麵酵,有婦人拿來,藏在三斗麵裡,直等全團都發起來.」(太十三33)這是指出神國度內裡的發展；芥菜種的發展是外面的,看得見的.神的話語,福音是神的大能,要改變相信接受的人.麵酵滲透整個麵團；福音的種子被接納,明白和遵行時,就改變生命.很多神使用的僕人使女,在未信主前犯罪失喪,信主後卻完全被改變,如保羅,奧古斯丁.達爾文是進化主義者,在一八八三年在Southsea Island時,發現那裡的人野蠻,落後,沒有文化,這些人真是沒有盼望；可是過了三十四年,他再到那地方,發覺那裡的人完全改變,其中有禮拜堂,他們生活有大的改變,並且充滿喜樂.他非常驚奇,後來他發現他離開後,有人到那裡傳福音,福音的大能改變了那些人.神的道種,在歷世歷代都發生內在改變的功效；馬丁路德帶出因信稱義的真理,把教會挽回；約翰衛斯理在英國度德敗壞的時代,藉著成聖的真理把教會復興過來,堵住破口.我們熟識的宋尚節,王明道,亦同樣使用神的道,作拯救和建立教會的偉大工作.基督徒應作精兵,將福音傳出去.
\newline
\newline
(五)藏寶於田與(六)尋珠的比喻
\newline
\newline
「天國好像寶貝藏在地裡......天國又好像買賣人尋找好珠子.」(太十三44-45)天國非常寶貴,我們要追求神的國度.我們發現天國的寶貝後,要盡力甚至變賣一切,去換取天國的成就.當時人沒有保險庫放寶藏,只有埋在地裡,當人死了,寶藏無人知道；有農夫在種田時,發現地中有寶藏,便把地買下來,這是人不會注意到的寶貝.天國就是這樣,要我們先求神的國和義,這是天國的重要性和寶貝.
\newline
\newline
(七)撒網的比喻
\newline
\newline
「天國好像網撒在海裡,聚攏各樣水族.網既滿了......丟在火爐裡,在那裡必要哀哭切齒了」(太十三47-50).這個比喻與稗子的比喻有同樣的性質,到了一天,義人與惡人要被分開；即麥子與稗子或好的魚與不好的,都會被分開.神是公義聖潔的,祂會按公義施行審判,到末世的時候,有最後的審判.這是世界最後的結局.
\newline
\newline
所以我們在神的國度中,要從七個比喻了解現在國度的光景,更要追求努力去擴展天國；
\newline
\newline
傳福音,面對撒旦魔鬼一切的攻擊,付代價,滿足神的心意,使主早日再來,開始永遠榮耀的國度!
\newline
\newline


\newpage

\section{第59屆~港九培靈會~金新宇~第五講}
\label{sec:1261}
\textbf{國度與教會}
\newline
\newline
連結: \href{https://www.hkbibleconference.org/session-message/view/1261}{\texttt{https://www.hkbibleconference.org/session-message/view/1261}}
\newline
\newline
\hyperref[sec:1257]{< < < PREV SERMON < < <}
~
\hyperref[sec:toc]{[返目錄]}
~
\hyperref[sec:1264]{> > > NEXT SERMON > > >}
\newline
\newline
過去我們從不同層面思想國度的信息,今天我們要在國度與教會的關係上探討.我們建立神的教會與擴展神的國度是否相同?
\newline
\newline
(一)國度與教會的分別
\newline
\newline
從教會歷史的立場看,是有不同的看法.最早的教父主張神的國度是屬靈的──在信徒的心中,教會在地上要表彰的國度.後來君士坦丁統一了羅馬,把基督教成為羅馬的國教；教會歷史家該撒利亞猶西比則主張,羅馬帝國是基督教的國家,是神統治的國家；所以這便是天國──教會在國度之中.可是奧古斯丁不同意這看法,因為羅馬皇是個黑暗的君王,羅馬國度更是個黑暗的國度,國度應該真正由神統治,國度就是爭戰的教會,他贊同國度相等於教會；而天主教亦接受他這種看法.在基督教中的解經家VOS根據(太十六18-19)主耶穌說:「我要把我的教會建造在這磐石上,陰間的權柄不能勝過他.我要把天國的鑰匙給你......」這裡他強調教會好像一所屋子的建造,有門並須要鑰匙,天國鑰匙是教會的鑰匙,所以他總結教會就是國度.可是在基督教中有很多人不同意這種看法──國度不等於教會,兄弟亦有這種看法.教會與國度不同,因為教會是一群蒙恩的人在一起,教會是被神統治的一群人；國度是神的權柄,神的統治；所以兩個是不同的.但我們相信教會與國度有密切的關係.主耶穌在地上從來沒有對門徒說:「你們就是國度.」我們看見教會與基督一同掌權.主耶穌說:「我要把教會建造在這磐石上,陰間的權柄不能勝過他.我要把天國的鑰匙給你.」彼得在五旬節時使用鑰匙,傳揚福音,叫三千人信主,把福音帶到人群中,教會的信徒是神統治的領域.
\newline
\newline
教會是基督的新婦,基督是天國的王,新婦與基督實在關係密切.教會是神在地上的代理人,是神國度在地上的工具；教會為要彰顯神的國和權能.教會的任務是要傳揚神國的福音,要拓展神的國度,而不是建立神的國度.「天國近了,你們要悔改,信福音!」這裡指明天國是臨到的.當人信主時,國度就臨到他的身上.我們要推廣神的國度,因為教會是神國度的基地.我們建立基督的身體──教會,使新婦預備好,所以我們必須作工,興旺福音.
\newline
\newline
(二)教會在地上所扮演的角色
\newline
\newline
我們看見主耶穌在地上,是顯出祂是天國的王；祂行了很多神蹟奇事,證明祂的王權.我們最終的目的是要得到神的國度,而我們必須禱告──願你的國降臨.教會在這二千年來不斷這樣禱告,神的國要降臨到我們中間,這禱告的盼望在於求主耶穌早日再來,神榮耀的國度來到.在啟示錄約翰的禱告,祂聽見主的聲音:「是的,我必快來!」他回答說:「阿們,我願你來.」願你的國降臨與願你快來好像不同,其實是指同一件事.主在世時傳講主要是國度的信息,四福音講論一百多次；但論及教會只有三次,因為那時教會還未被建立,而以色列人卻一直盼望彌賽亞來到.四福音主要是國度的信息.舊約也是以國度為主要的信息.
\newline
\newline
及後在書信中,似乎表面上有轉變,沒有強調國度；在以弗所書有九次講教會,一次講國度；哥林多前書有廿二次講教會,五次講國度,書信中重點在教會；使徒行傳則由國度的信息轉至教會的信息──國度八次,教會廿二次.在新約書信結束的啟示錄中,又回到國度的信息.所以貫串整本聖經,國度的信息是十分重要.在使徒行傳開始時主與門徒相處四十天,主要是談論神國的事；在結束時,保羅在(徒廿八31)放膽傳講神國的道.
\newline
\newline
書信為甚麼講教會比較多呢?(1)主耶穌在世,天國的王在地上；復活的基督有能力住在信徒的心中,主藉著信徒和教會,主在其中.教會是基督的身體,是那充滿萬有者所充滿的,基督就如在地上.弟兄姊妹,我們在地上的責任是何等的重,世上的人看不見主,卻看見我們,我們應怎樣榮耀主!基督徒的軟弱實在對國度的神的名有很大的虧損.有一個傳道人說,當耶穌的肉身不在我們中間,跟隨祂的人必定替代祂,成為祂的代表,這是使徒門和現今信徒的責任；我們被委托作王的代表.主耶穌復活後曾說:「父怎樣差遣我,我要照樣差遣你們.」所以教會在地是主的全權代表.
\newline
\newline
我們看以弗所書第一章,清楚告訴我們,(弗19-22)「並知道祂向我們這信的人所顯的能力,是何等浩大,就是照祂在基督身上,所運行的大能大力......又將萬有服在祂的腳下中使祂為教會作萬有之首.教會是祂的身體,是那充滿萬有者所充滿的.」天國的王藉著祂的身體──教會在地上.基督是頭,而頭必須有身體方能行動,去工作.教會是代替基督在地上工作.這是書信強調解釋教會的原因.教會在地上要彰顯主的權能.我們奉主的名禱告,能發揮莫大的權能.(2)可能在教會開始時常講天國,會導致人誤會.在使徒行傳十七章,保羅和西拉被指為攪亂天下的.他們被指責是違背該撒的命令,說另有一主耶穌.所以講及國度時會帶來羅馬政府的反感.這是書信很少論及國度的另一原因.書信的作者便用教會為強調的主題.
\newline
\newline
「教會是基督的身體,是充滿萬有者所充滿的.」我們的主是充滿萬有,萬有要靠著基督去滿足其需要；而神的統治,是讓信徒和教會得益處.神要充滿身體,基督要充滿萬有；整個宇宙為著教會,要被基督充滿.教會必須充滿基督的生命,神的愛,聖靈及基督的恩賜和神的同在；所以那充滿萬有者要充滿祂的教會.在原文字義上「充滿」有兩個特別的意義:
\newline
\newline
(一)繙成「豐滿」──即基督的豐滿充滿了我們.
\newline
\newline
(二)繙成「補滿」──即教會要補滿基督的豐滿.
\newline
\newline
難道基督還不夠完全嗎?那麼為何要教會去補滿呢?回應這兩個深入的意義,我們不妨深入思想.
\newline
\newline
(一)為甚麼基督要充滿祂的教會?因為教會要彰顯神的國度:基督這位天國的王在地上,彰顯神的國度；基督這位天國的王在地上,彰顯了神國的權柄和表徵；而今天信徒有基督的同在和能力,亦可見證表彰神的國度.(弗三19下)「便叫神一切所充滿的,充滿了你們.」得著一切的豐盛即得著神的愛和能力.此外,神亦給教會恩賜；以弗所書第四章論及神賜給教會不同恩賜的人；林前十二章論及聖靈將恩賜給予個別信徒.(弗三10)「為要藉著教會,使天上執政的掌權的,現在得知神百般的智慧.」總括來說,教會要使用神所賜予的豐盛,宣揚神完全的榮耀.這是信徒在地上的工作.
\newline
\newline
(二)補滿基督的豐滿有兩方面的含義:
\newline
\newline
(1)基督充滿教會,要教會有力量完成祂在地上的工作；主再來前有很多工作要完成,教會靠著聖靈完成工作.有一個神話形容當基督帶著身體回到天上時,有釘痕的手和在地上受苦的形像,天使加百列問祂:「你在地上受了很多苦,是否地上的人都知道你如何的愛他們呢?」主回答:「不是.只有很少人知道!但我留在地上的幾個門徒,他們會傳給萬民,叫世人知道我的救贖.」加百列說:「如果他們忘記或不遵行呢?你是否有其他計劃?」主說:「沒有!」所以天國的的擴展,在乎教會中的你與我；天使不能代替信徒,因為祂們沒有經歷十字架的大愛,寶血救贖的能力.所以教會在地上是基督在地上,完成祂救贖的工作.
\newline
\newline
(2)教會在地上是基督在地上.(西一24)保羅說:「現在我為你們受苦,倒覺歡樂,並且為基督的身體,就是為教會,要在我肉身上補滿患難的缺欠.」我們要不單補滿基督的豐滿,更要補滿基督患難的缺欠.然而十字架受苦受死的救贖工作已經成就,我們不能參與,若基督在地上亦會受不信者的恨惡和逼迫,我們是基督徒,世界亦會恨惡我們.教會在地上亦要繼續代替基督受苦；因為教會與基督聯合,教會是基督的身體,逼迫教會即逼迫基督.基督徒又與基督聯合,我們肯定要面對逼迫,受苦,補滿基督患難的缺欠.
\newline
\newline
教會是基督的新婦,羔羊的妻子.我們與基督的關係十分密切.(詩八5):「神叫人比天使微小一點,並賜他榮耀尊貴的冠冕.」這是大衛對神看人的認識.原文中的「天使」,聖經小字是「神」.其實稍次於神,這表示神看重人,在人的地位,身份和受造的本質上,是超過一切的.雖然人永遠不能與神同等,但實在有生命的人是奇妙尊貴,被神看重.所以特別是蒙恩得救的信徒更要彰顯神的榮耀.主復活後將一切權柄賜給我們,教會有大使命必須遵行,要把天國的子民和好的種子,撒在世界各處,到末日要顯出來,新婦亦預備好了.教會是十分重要的團體,是一個屬靈的有機體,為神的國度爭戰,面對撒旦的權勢.基地必須被建立起來,教會必須堅強起來,面對爭戰,同心合意,興旺福音.教會在耶和華軍隊的元帥帶領下,要打美好的仗.(啟十九7-8):「我們要歡喜快樂,將榮耀歸給祂.因為羔羊婚宴的時候到了,新婦也自己預備好了,就蒙恩得穿光明潔白的細麻衣,這細麻衣就是聖徒所行的義.」我們要存這盼望,直到主再來.
\newline
\newline


\newpage

\section{第59屆~港九培靈會~金新宇~第六講}
\label{sec:1264}
\textbf{怎樣進神的國}
\newline
\newline
連結: \href{https://www.hkbibleconference.org/session-message/view/1264}{\texttt{https://www.hkbibleconference.org/session-message/view/1264}}
\newline
\newline
\hyperref[sec:1261]{< < < PREV SERMON < < <}
~
\hyperref[sec:toc]{[返目錄]}
~
\hyperref[sec:1265]{> > > NEXT SERMON > > >}
\newline
\newline
我們過去是死在罪惡過犯中,沒有神,更沒有指望；今天神藉著基督十字架的救贖,拯救了我們,使我們升高,成為羔羊的妻子,基督的身子.教會在地上要擴展神的國度,基地必須建立穩固,所以教會要被建立是非常重要的.我們事奉這位萬王之王,萬主之主,是要透過教會成為基地,才能有好的事奉；我們必須忠心於自己的所屬教會,事奉,奉獻,看重建立教會的團契.雖然我們要有國度的眼光胸懷,但重點在教會,並超越教會的牆,伸展到各處；教會更應支持福音機構,服事眾教會,進入得著更多的未得之民.我們最重要的目的,是拓展神的國度,榮耀神,彰顯神,彰顯神的權能.
\newline
\newline
今天我們要思想如何進入神的國.約翰福音第三章,論及耶穌與尼哥底母說:「我實實在在的告訴你,人若不重生,就不能見神的國.」(約三3)又在5節說:「我實實在在的告訴你,人若不是從水和聖靈生的,就不能進神的國.」我們必須重生,方能看見和進入神的國.我們要進世界的國度是易事,只要買機票,花時間和金錢一行便可,無論任何名勝古蹟,都能一覽無遺.如果我們要看神的國,只有一條路──必須重生.如果我們要一嚐生命水,享受生命樹的果子,我們必須重生.重生是神的作為,不是自己任何方法和功勞；重生須要由上賜下,須要由神而生,完全是神的恩典.可是人亦有責任,對神有回應.(弗二8)「你們得救是本乎恩,也因著信；這並不是出於自己,乃是神所賜的.」我們應有以下的回應:
\newline
\newline
(一)相信與悔改
\newline
\newline
真實的信心是悔改.這不是表面的後悔或懊悔,而是要改變.改──在罪惡生活中出來.悔改是心思和態度的改變.原來我們的心向著世界,向著罪惡,背向神；現在卻轉過來,心向著神,行神所喜悅的事,追求公義,聖潔.如果祇是相信而沒有悔改,就只有頭腦思想上的相信,我們的心仍未歸向神,聖靈就不能重生我們.(西一13):「祂救了我們脫離黑暗的權勢,把我們遷到祂愛子的國裡.」這是神的作為,神的恩典.祂救我們脫離罪惡的權勢,國度；而我們必須立志離開罪惡黑暗的權勢,才能進入愛子的國度.這一方面是神的工作,而一方面人要悔改相信.「我們心裡相信,就可以稱義；口裡承認,就可以得救.」(羅十10)傳福音方法中的屬靈四定律非常好,使人明白福音；可是頭腦思想上的相信,卻不一定保證人可進入神的國度.心裡相信是要悔改,決心離開罪惡.神重用的傳道人陶恕博士寫了一本書,其中講及一個州長進到監獄去,獄中無人知曉他是州長,他找到一位年青的囚犯,想赦免他,便對他說:「若你能離開監獄,你會如何做呢?」那青年回答:「第一件要殺定我罪的法官.」那州長心裡很難過,他又如何能赦免這樣的人.
\newline
\newline
我們必須悔改,痛恨我們所犯的過錯,罪行.(林前六9-11):「你們豈不知,不義的人不能承受神的國麼?......並藉著我們神的靈,已經洗淨,成聖稱義了.」主耶穌開始傳道工作時,大聲疾呼說:「天國近了,你們應當悔改.」將來在新聖城耶路撒冷,「凡不潔淨的,並那行可憎與虛謊之事的,總不得進那城,只有名字寫在羔羊生命冊上的纔得進去.」(啟廿一27)以賽亞先知論及赦免的道理時說:「你們的罪雖像硃紅,必變成雪白；雖紅如丹顏,必白如羊毛.」(賽一18下)這是何等大的赦免!但在赦免前,他提及:「你們要洗濯,自潔；從我眼前除掉你們的惡行,要止住作惡.學習行善,尋求公平.......」(賽一16-17)這是我們的責任.
\newline
\newline
(二)進入窄門
\newline
\newline
(太七13-14):「你們要進窄門.因為引到滅亡,那門是寬的,路是大的,進去的人也多.引到永生,那門是窄的,路是小的,找著的人也少.」如何進入生命呢?進入生命是進入神的國,進入生命和神的國必須「進窄門」.今天我們可進入神的國,享受生命的豐盛,並要走小路.這條小路窄門要我們付代價,摒除一切的罪惡,放下物質一切的引誘,放下自私自義,放下不肯赦免人的罪,我們更要捨己.主耶穌在這裡的教訓是給人挑戰,讓人去揀選；這裡有兩個門,兩條路,有兩批人,有兩個不同的結局.主耶穌要人作出揀選,不是單在思想上,而是有實際的行動.我們如何作國度的公民呢?我們如何得生命呢?我們必須揀選.我們一生中有很多的選揀,決定,有些是不太重要的,如吃喝玩樂；可是有些對前途有很大的影響:如職業,配偶.但最重要的揀選是關乎耶穌基督與祂的國度,這不但影響我們前面的日子,更對我們將來的永遠有密切的關係,是永死與永生的問題.
\newline
\newline
在舊約聖經中,神多次提醒祂的子民,藉著先知吩咐他們要揀選.在曠野中,摩西要求以色列人作一個整全的選擇,是關於他們生與死的選擇:(申卅19-20)「我今日呼天喚地向你作見證,我將生死,禍福陳明在你面前,所以你要揀選生命,使你和你的後裔都得存活,且愛耶和華你的神,聽從祂的話,要靠祂,因為祂是你的生命,你的日子長久,也在乎祂.」他要以色列人揀選神,因為神是我們的生命,而神的話帶給我們生命,聽神的話就可得著生命.以色列人進入迦南地之後,約書亞也要求以色列人揀選:(書廿四14-15)「現在你們要敬畏耶和華,誠心實意的事奉他.將你們列祖在大河那邊和在埃及所事奉的神除掉,去事奉耶和華,......至於我,和我家,我們必定事奉耶和華.」這是事奉偶像或是事奉神自己的揀選.在迦密山頂,以利亞問神的子民,給他們機會揀選:(王上十八21)「你們心持兩意要到幾時呢?若耶和華是神,就當順從耶和華；若巴力是神,就當順從巴力.眾民一言不答.」我們從聖經中看見,神不斷要我們揀選.今天神要我們選擇,是否願意得著基督,順從神的話,得著生命.我們實在是沒有其他的選擇.
\newline
\newline
中國人喜歡中庸之道,走中間的路線；可是神卻要求我們抉擇.做天國的子民是不容易的,我們要走完這條窄路.除祂以外,別無拯救,主是道路,真理和生命,必須藉著祂；祂更是門,通往永生之門.當我們走永生之路時,我們順服聖靈的帶領,過分別為聖的生活,有很多影響見證的習慣,生活,如抽煙,醉酒,賭博等,我們都不參與,也許很多未信者會批評我們太窄了!你們不須要這樣做!我信主後熱心事奉,很多親友時常對我說:「有信仰是很好的,可是不要太過迷!你要順應潮流,不可太古板.」但我們不能夠如此,更不能活在罪中.我們要勝過文士和法利賽人單顧外面的儀文,沒有敬虔的特質；我們裡外都應符合聖徒的體統,過聖潔公義的生活.西印度群島有一個人,寧願信回回教,不接受基督教,因為他信的教是寬大的,可以讓人與罪一同進入裡面；但你們信基督的路,實在太窄；要將一切的罪放在背後,實在太難了.寬門大路很吸引人,受人歡喜,更合潮流；想做就去做,沒有約束,可以自由；但最後卻帶來滅亡,被審判.進入小路窄門要付很大的代價,要捨己,背十字架跟從神,完全獻給神使用!這個「小」字有受苦難和逼迫的意思.在(徒十四22)保羅說:「堅固門徒的心,勸他們恆守所信的道.又說,我們進入神的國,必須經歷許多艱難.」這裡的「艱難」與「小」字的字根相同.跟從神也許要放下我們的尊嚴.例如在(王下五1-19)講及亞蘭王的元帥乃縵的事蹟:他貴為尊大的元帥,卻患了大麻瘋,這病可謂藥石無靈；後來聽從一位以色列的小女子介紹,找先知以利沙求醫.以利沙叫他到約但河洗七次,病就可不藥而癒；可是乃縵心中氣忿,最後放下尊嚴聽先知吩咐,病果真好了.我們進神的國,也許要像乃縵一樣,謙卑下來,甚而放下自我的尊嚴!(路十三24):「耶穌對眾人說:你們要努力進窄門.」如打仗一樣,努力的爭取；又如運動員參加世運會,努力地要取得金牌一樣!所以我們要努力放下一切不合神心意的東西,專心尋求,遵行神的旨意,享受神的豐盛,榮神益人,結滿果子.
\newline
\newline


\newpage

\section{第59屆~港九培靈會~金新宇~第七講}
\label{sec:1265}
\textbf{誰是天國的子民}
\newline
\newline
連結: \href{https://www.hkbibleconference.org/session-message/view/1265}{\texttt{https://www.hkbibleconference.org/session-message/view/1265}}
\newline
\newline
\hyperref[sec:1264]{< < < PREV SERMON < < <}
~
\hyperref[sec:toc]{[返目錄]}
~
\hyperref[sec:1267]{> > > NEXT SERMON > > >}
\newline
\newline
今天我們思想,在神的國度中,誰是天國的子民?這個主題與昨天的信息有關連.作天國的子民,必須要讓神在我們的生命中作主,包括認識神是怎樣的一位神?我們應如何生活?天國的主耶穌基督的性格如何?我們應該如何效法祂的榜樣?
\newline
\newline
在主耶穌登山寶訓的講論中,論及八福的真理時,就表現出神的屬性,我們要怎樣效法.此外我們要盡天國子民的責任,發揮應有的影響力,有好的生活,彰顯生命的力量.主耶穌說:「你們的光也當這樣照在人前,叫他們看見你們的好行為,便將榮耀歸給你們在天上的父.」(太五16)
\newline
\newline
(一)「八福」的要求
\newline
\newline
(1)「八福」是神生命的表徵,亦是神子民生活的標準,是我們支取生命的泉源.神的公義,表明祂公義的管治,亦是我們過公義生活的準則.(詩九7-9)說:「耶和華坐著為王,直到永遠；祂已經為審判設擺祂的寶座,祂要按公義審判世界,按正直判斷萬民.......」公義是神的本性,祂的國也是公義的.耶穌基督亦是公義的彌賽亞,我們必須過公義的生活.約翰對神的認識,在約翰一書二29記載:「你們若知道祂是公義的,就知道凡行公義之人都是祂所生的.」我們要識別信徒的生命,從他的表現和行為中,就可對他有基本的認識.保羅在羅馬書十四17說:「因為神的國不在乎吃喝,只在乎公義,和平,並聖靈中的喜樂.同彼得在(彼後三13):「但我們照祂的應許,盼望新天新地,有義居在其中.」這些門徒的生命中,以神的義為生活的中心.我們亦應效法他們,從世界中分別出來,過公義的生活.我們的神亦有其他屬性,如憐憫,愛等；我們應該像神.
\newline
\newline
登山寶訓的八福,從(太五3-10)說出在國度裡的生活的標準,這是我們很熟悉的；我們不妨從國度的焦點來看,了解其中的意義.第一福至八福講及子民的性格.第一福虛心的人有福了,可譯成「靈裡貧窮」的人有福了,天國現在是他們的；第八福為義受逼迫的人有福了,因為天國現在是他們的.感謝神,我們全心投靠祂,天國就是我們的；我們可以現在就享受救恩和其他很多福份.可是在第二至第七個福,就有些不同,「我們必得......」,這似乎是將來的事!當我們詳細研究希臘文的文法時,發現這是形容現在至將來的表達,是不斷的持續.所以第二至七福,應該包含在第一至第八福之間,是天國應有的福份；這種將來的表達,亦表示那些福份將來會完全得到.這些應許,是由神確實的保證,我們必定得到的.
\newline
\newline
(二)「八福」的應許
\newline
\newline
(1)「虛心的人有福了,天國是他們的.」(太五3)我們希望有份於神的國度,必須有決心進入神的國；要進入神國必須虛心放下自己,不要過自我為中心及自義的生活.我們要靈裡貧窮,如乞丐需要食物一樣,期望求神的憐憫,仰望神的恩典.「照我本相,無善足陳」,是我們的光景.從前有一個藝術家,想畫一幅浪子回頭的故事,他便想找一個模特兒,剛好遇見一個乞丐,正合他作品的需要,他便邀請那乞丐明天見他,預備把那人作模特兒；可是第二天那人來到,並且穿得很整潔,可憐相亦沒有了,他幾乎不認出那人,那人怎能作模特兒呢?進入神的國必須是虛心的,靈裡貧窮的人.
\newline
\newline
(2)「哀慟的人有福了,因為他們必得安慰.」(太五4)看見自己屬靈的光景,為自己難過,更為普世失喪的人活在罪惡過犯中難過；這樣哀慟的人,必定得到安慰.在新天新地臨到時,我們便得到完全的安慰,神要擦去我們的眼淚,此時不再有死亡,痛苦和悲哀.
\newline
\newline
(3)「溫柔的人有福了,因為他們必承受地土.」(太五5)我們的主耶穌是溫柔的人,祂在世時說:「我心裡柔和謙卑,你們當負我的軛,學我的樣式；這樣,你們心裡就必得享安息.」(太十一29)祂進入耶路撒冷的時候,是騎著驢駒子,應驗先知撒迦利亞的預言:「看哪!你的王來到你這裡,祂是公義的,並且施行拯救,謙謙和和的騎著驢,就是騎著驢的駒子.」(亞九9下)祂是如此溫柔,是我們學效的榜樣.溫柔的人必得地土.在舊約(詩卅七11):「但謙卑人必承受地土,以豐盛的平安為樂.」香港人從事地產行業的,都想有地土；大財主通常有地和產業,這裡所指的不是物質上的地土,溫柔人不是活在世俗中尋求；更不會貪愛世界,但他們是從神那裡領受,感受到滿足.世上的人卻沒有真正的滿足.
\newline
\newline
(4)「飢渴慕義的人有福了,因為他們必得飽足.」(太五6)他們追求公義,如飢如渴.天國是公義的國度,我們要活出公義的生活；並要追求神的國和神的公義.主耶穌在太五20說:「我告訴你們,你們的義,若不勝於文士和法利賽人的義,斷不能進天國.」耶穌用嚴厲警告的字眼:「我告訴你們......」,表明真實的公義生活不在乎外表的守法,而要有實質,從內心敬虔流露出來,持守真理,是表裡一致的公義.耶穌更用六個例子來指出法利賽人和文士的不義:如他們那強調的「不可殺人」,「不可姦淫」,「不可背誓」,「以眼還眼,以牙還牙」,「當愛你的鄰舍,恨你的仇敵」等事實中,其中還有很多不足夠的地方,未能如天父要求的完全!今天聖靈會幫助我們,超過我們能力範圍的實踐,過更合神旨意的公義生活.「使律法的義,成就在我們這不隨從肉體,只隨從聖靈的人身上.」(羅八4)這是更有內涵的義者生活.
\newline
\newline
(5)「憐恤人的人有福了,因為他們必蒙憐恤.」(太五7)憐恤包括赦免一切得罪自己的人,而且是對一切有需要的人加以憐憫.主在地上,看見很多困苦流離的人,就憐憫他們.祂要幫助一切受欺壓的人,患難的人.神看顧孤兒寡婦,為一切受逼迫的人伸冤.福音不單傳給中等階層或高階層的人,更是要傳給低階層的窮苦人；我們目標不是傳社會福音,卻是在傳救贖的福音,拯救失喪的靈魂.
\newline
\newline
(6)「清心的人有福了,因為他們必得見神.」(太五8)我們做事要有單純專一的動機,純潔地追求認識神.專一的委身事奉神,神的國度和權柄.這種內裡的純正是非常重要的.(啟廿二3下):「在城裡有神和羔羊的寶座,祂的僕人都要事奉祂,也要見祂的面.」那時人要面對面見神.今天我們亦可經歷神的同在,神的愛和神的能力.
\newline
\newline
(7)「使人和睦的人有福了,因為他們必稱為神的兒子.」(太五9)天國的王來到地上是藉著祂在十字架上的死,祂流出的血,叫我們與神和好,亦讓猶太人與外邦人和好,亦叫人與人之間和好.猶太人拉比強調:「一個人最高的任務,是能叫人與人之間有好的關係.」我們要將和平帶到人間,這種和平是以公義為基礎,而公義的果效就是平安.這是和平福音所生發的影響.信徒用生活同證和平的福音,是神國子民的表現.
\newline
\newline
(8)「為義受逼迫的人有福了,因為天國是他們的.」(太五10)當我們高舉公義,活出義的生活,必定要受人的逼迫.彼得在書信中曾說:「你們為義受苦是有福的.」我們這樣生活,必定遭遇逼迫.(提後三12下):「凡立志在基督耶穌裡敬虔度日的,也要受逼迫.」
\newline
\newline
我們從「八福」中看到天國子民的標準,必須有如此品格,並且活出這種見證,方能影響這個世界.主耶穌繼續說:「你們是世上的鹽和光.」我們要為天國作見證.
\newline
\newline
基督徒有雙重公民的身份；我們一方面是地上的公民,另一方面更是天國的公民,有當盡的責任.我們更應順從神,如鹽生發防腐的作用；又如光是放在燈臺上,照亮四周的人.這好比基督徒的美德生活,要影響社會.在十八世紀時代英國社會道德敗壞,神興起約翰衛斯理作復興教會的工作,並且更新改進社會的道德,使英國在十九世紀成為國際強國.主耶穌的愛,赦免,美德,要透過我們彰顯出來,用生命見證所傳的福音；光是外顯的工作,鹽是內在隱藏的工作,是我們作見證時一同發揮的影響.你們的光也當這樣照在人前,叫他們看見你們的好行為,便將榮耀歸給你們在天上的父.願我們都成為天國子民,發揮福音的真光,照耀黑暗中的人.
\newline
\newline


\newpage

\section{第59屆~港九培靈會~金新宇~第八講}
\label{sec:1267}
\textbf{願你的國降臨 — 國度與差傳}
\newline
\newline
連結: \href{https://www.hkbibleconference.org/session-message/view/1267}{\texttt{https://www.hkbibleconference.org/session-message/view/1267}}
\newline
\newline
\hyperref[sec:1265]{< < < PREV SERMON < < <}
~
\hyperref[sec:toc]{[返目錄]}
~
\hyperref[sec:1193]{> > > NEXT SERMON > > >}
\newline
\newline
神的兒女二千年來,唸主禱文時,不斷求神的國降臨；可是,「願你的國降臨」到底有甚麼意思呢?雖然神的國今日已經來臨──主耶穌第一次降世,就是天國的王來到,神的王權統治已實現於信徒心中；但神的國度最豐滿,榮耀的來臨,要等待主第二次再來才實現.
\newline
\newline
我們禱告:「願你的國降臨」也就是指神的國最後榮耀,完全實現在地上,並得勝一切仇敵；消滅一切罪惡,把公義,和平和喜樂,帶給我們.所以「願你的國降臨」與使徒約翰在啟示錄(廿二20)禱告說:「阿們!主阿!我願你來!」有密切的關連；約翰是聽到主說:「是的,我必快來!」；而他迫切的說:「主耶穌阿,我願你來!」.「主再來」和「神的國降臨」是不能分開的.
\newline
\newline
「願你的國降臨」是指末後發生的事.今天祇有部份信徒讓基督作王,這個世界仍然沒有歸主；而我們所盼望的,是主再來時,祂要在所有人身上作王.這是我們應有的盼望和禱告.當主教導門徒禱告時說:「所以你們禱告,可以這樣說:『我們在天上的父......』門徒沒有問:『你的國』是指甚麼說的?」他們明白:「因為猶太人在等待彌賽亞的來臨,建立祂的王國；得勝仇敵,脫離羅馬人的統治；使以色列人再次享受比過去更榮耀的王國.」
\newline
\newline
主耶穌升天後,門徒在(徒一6)問主說:「主阿,你復興以色列國,就在這時候嗎?」主耶穌沒有否認祂要復興以色列國；但這必須是主再來時,才會實現的.所以今天我們當靠著聖靈的能力,為主作見證,傳福音,作差傳事工,擴展神的國度；使神的王權能被更多人接納.這樣,我們必須在下面努力:
\newline
\newline
(一)我們當繼續不斷禱告──願神的國早日來臨；
\newline
\newline
(二)我們當把天國的福音,傳遍天下,給萬民作見證,然後末期才到；
\newline
\newline
(三)我們應等候主再來,持守敬虔,聖潔的生活,來促使主早日來臨.
\newline
\newline
我們在這裡可詳細思想這三方面:
\newline
\newline
(一)我們應該禱告:
\newline
\newline
(太九35-38)主對門徒說:「要收的莊稼多,作工的人少,所以你們當求莊稼的主打發工人出去,收祂的莊稼.」今天這種情況,比過去任何時期更嚴重.Ruth Siemens在「世界的機會」一九八七年一月號說:「全世界百分之九十五的宣教工作,在全球百分之十七的地方進行,大部份的國家是禁止宣教士入境工作的；而那些開放的國家,當地人對福音全無接觸,必須最少有五十萬位宣教士作跨越文化的工作,方能滿足這個時代的需要.」
\newline
\newline
聯合國人口活動基金組織估計:「人類第五十億名嬰兒將在今年一九八七年七月十一日誕生,比較起一九五○年的二十億,多了一倍半.而世界人口每分鐘增長一百五十人,每天增長廿二萬；試問今天,宣教士每天增加多少?可補足廣大工作的需要!」
\newline
\newline
我們當求莊稼的主是對的.可是,我們只有禱告而沒有行動,也是沒有用的.馬丁路德當年的同工,開始時只在禱告中支持他改教工作；有一晚夢中見一片很大的禾場,單有馬丁路德在收割,他便醒悟過來,不單為他禱告,更加入改教的工作行列,一同收割.
\newline
\newline
(二)這天國的福音,要傳遍天下,對萬民作見證,然後末期才來到.(太四14)
\newline
\newline
我們傳福音的目的,不祇是把得救的人加給教會,不祇是讓神的國和王權能擴展到多人心中；更是要將福音傳遍天下,對萬民作見證,然後末期才到.我們當然盼望更多人能接受主,但也要普遍地將福音傳到世界.
\newline
\newline
羅馬書十一章保羅講到今天的福音要臨到外邦人,等到外邦人的數目添滿:「外邦人的數目添滿」,是指多少人?聖經說:「只有神知道.」主耶穌曾說:「我另外有羊,不是這圈裡的,我必須領他們來,他們也要聽我的聲音,並且要合成一群,同歸一個牧人.」(約十16)我們當把主另一個羊圈的羊帶進主的羊圈；福音實在已託付我們,我們當盡力忠心去傳揚.
\newline
\newline
但要期待主等二次來臨,祇在本地傳福音,領人歸主,還是不夠的.這福音必須傳遍天下,給萬民作見證.故此今天我們必須強調差傳工作；這是主託付教會和信徒的大使命(太廿八18-20).祂是主,祂是救世主(約三16)；因為神愛世人!那些拒絕福音的人,當主再來臨時,他們不是不明不白的受審判.好像挪亞方舟的時代,方舟門關了,洪水就來了；那在方舟外的人,由挪亞傳道時計算,神容忍等待了一百二十年,但人仍然不信,他們不能說神並未警告他們悔改.今天末後的日子,與挪亞的日子是一樣的.(路十七26-27).
\newline
\newline
傳福音或差傳工作是為了擴張神的國度,預備主榮耀再臨.(西一13)說明我們重生得救時,就已經脫離了黑暗的權勢──撒旦的權勢；而(約壹五19)又說明撒旦是那惡者,是世界的王.這是傳福音的目的:把重生得救的人帶進神的國度,這是神的作為和聖靈的工作；擴展神的國度,教會和我們必定與撒旦作屬靈的爭戰.我們可奉主的名,捆綁撒旦,叫牠離開,更可向撒旦宣告基督的得勝,粉碎牠一切作為.
\newline
\newline
但撒旦在地上並未被消滅.今天我們並未把全人類從撒旦的國度,爭取到基督的國度.我們傳福音祇是要求人作一個抉擇:聖靈給他們自由接受或拒絕；所以我們傳福音或作差傳工作,並非要這個世界基督化,這有待主再來時才能實現.今天我們做的福音工作是使神的國與撒旦的國更明顯的對立；我們從啟示錄看出,在主再來之前,罪惡更猖獗,神的審判和忿怒,更厲害臨到這個罪惡世界；一直到主再來的時候,罪惡才全被消除,撒旦和牠的使者,最後被扔在硫磺的火湖裡.
\newline
\newline
我們差傳的使命應有這榮耀的異象!可帶給我們力量和鼓舞.我們當有國度的眼光,天國子民的胸懷；今天神的國度要擴張,這是差傳的重要目標.如果不是為了神的國度,教會作差傳事工就缺少意義:在海外有人得救,對本地教會的得救人數並未增加.但擴展神的國度,是我們最終的目的.
\newline
\newline
(三)等候主再來,持守敬虔,聖潔的生活,促使主早日來臨──即神榮耀國度的彰顯.
\newline
\newline
這是我們應有的禱告:「願你的國降臨」!我們不知這榮耀的國度何時來臨?但知道祂必快來.我們怎樣等候主再來?進入祂榮耀的國度?我們怎樣迎見我們的王?(約壹二28)「小子們哪!你們原住在主裡面,這樣,祂若顯現,我們就可以坦然無懼,當祂來的時候,在祂面前也不至於慚愧.」(彼後三11-12)「這一切既然都要如此銷化,你們為人該當怎樣聖潔?怎樣敬虔?切切仰望神的日子來到.」
\newline
\newline
今天猶太人不接受基督,他們暫時被神棄絕,救恩臨到我們的外邦人；我們當有聖潔,敬虔的生活,等候主再來,也催促主早日再來,帶來神榮耀的國度!
\newline
\newline


\newpage

\section{第60屆~港九培靈會~劉承業~第一講}
\label{sec:1193}
\textbf{得勝的先決條件}
\newline
\newline
連結: \href{https://www.hkbibleconference.org/session-message/view/1193}{\texttt{https://www.hkbibleconference.org/session-message/view/1193}}
\newline
\newline
\hyperref[sec:1267]{< < < PREV SERMON < < <}
~
\hyperref[sec:toc]{[返目錄]}
~
\hyperref[sec:1194]{> > > NEXT SERMON > > >}
\newline
\newline
約書亞記 1:1-1:9
\newline
\begin{longtable}{cl}
\hline
\hline
章節 & 經文 (和合本修訂版)\\
\hline
1:1 & \begin{tabularx}{0.7\textwidth}{X} 耶和華的僕人摩西死了以後，耶和華對摩西的助手嫩的兒子約書亞說： \end{tabularx} \\ \\ \relax
1:2 & \begin{tabularx}{0.7\textwidth}{X} 「我的僕人摩西死了。現在你要起來，和眾百姓過這約旦河，往我所要賜給以色列人的地去。 \end{tabularx} \\ \\ \relax
1:3 & \begin{tabularx}{0.7\textwidth}{X} 凡你們腳掌所踏之地，我都照我所應許摩西的話賜給你們了。 \end{tabularx} \\ \\ \relax
1:4 & \begin{tabularx}{0.7\textwidth}{X} 從曠野和這黎巴嫩，直到大河，就是幼發拉底河，赫人的全地，又到大海日落的方向，都要作你們的疆土。 \end{tabularx} \\ \\ \relax
1:5 & \begin{tabularx}{0.7\textwidth}{X} 你一生的日子，必無人能在你面前站立得住。我怎樣與摩西同在，也必照樣與你同在；我必不撇下你，也不丟棄你。 \end{tabularx} \\ \\ \relax
1:6 & \begin{tabularx}{0.7\textwidth}{X} 你當剛強壯膽，因為你必使這百姓承受那地為業，就是我向他們列祖起誓要給他們的地。 \end{tabularx} \\ \\ \relax
1:7 & \begin{tabularx}{0.7\textwidth}{X} 只要剛強，大大壯膽，謹守遵行我僕人摩西所吩咐你的一切律法，不可偏離左右，使你無論往哪裡去，都可以順利。 \end{tabularx} \\ \\ \relax
1:8 & \begin{tabularx}{0.7\textwidth}{X} 這律法書不可離開你的口，總要晝夜思想，好使你謹守遵行這書上所寫的一切話。如此，你的道路就可以亨通，凡事順利。 \end{tabularx} \\ \\ \relax
1:9 & \begin{tabularx}{0.7\textwidth}{X} 我豈沒有吩咐你嗎？你當剛強壯膽，不要懼怕，也不要驚惶，因為你無論往哪裡去，耶和華你的神必與你同在。」 \end{tabularx} \\ \\
[1ex]
\hline
\hline
\end{longtable}


\newpage

\section{第60屆~港九培靈會~劉承業~第二講}
\label{sec:1194}
\textbf{信心的榜樣}
\newline
\newline
連結: \href{https://www.hkbibleconference.org/session-message/view/1194}{\texttt{https://www.hkbibleconference.org/session-message/view/1194}}
\newline
\newline
\hyperref[sec:1193]{< < < PREV SERMON < < <}
~
\hyperref[sec:toc]{[返目錄]}
~
\hyperref[sec:1195]{> > > NEXT SERMON > > >}
\newline
\newline
約書亞記 2:1-2:14
\newline
\begin{longtable}{cl}
\hline
\hline
章節 & 經文 (和合本修訂版)\\
\hline
2:1 & \begin{tabularx}{0.7\textwidth}{X} 嫩的兒子約書亞從什亭暗中派兩個人作探子，說：「你們去窺探那地和耶利哥。」於是二人去了，來到一個名叫喇合的妓女家裡，在那裡睡覺。 \end{tabularx} \\ \\ \relax
2:2 & \begin{tabularx}{0.7\textwidth}{X} 有人告訴耶利哥王說：「看哪，今夜有以色列人到這裡來窺探此地。」 \end{tabularx} \\ \\ \relax
2:3 & \begin{tabularx}{0.7\textwidth}{X} 耶利哥王派人到喇合那裡， 說：「你要交出那來到你這裡、進了你家的人，因為他們來是要窺探全地。」 \end{tabularx} \\ \\ \relax
2:4 & \begin{tabularx}{0.7\textwidth}{X} 但女人已把二人藏起來，卻說：「那兩個人確實到我這裡來過，他們從哪裡來，我卻不知道。 \end{tabularx} \\ \\ \relax
2:5 & \begin{tabularx}{0.7\textwidth}{X} 天黑、要關城門的時候，他們就出去了。他們往哪裡去我也不知道。你們趕快去追他們，就必追上。」 \end{tabularx} \\ \\ \relax
2:6 & \begin{tabularx}{0.7\textwidth}{X} 其實，這女人已經領二人上了屋頂，把他們藏在她擺列在屋頂的亞麻梗中。 \end{tabularx} \\ \\ \relax
2:7 & \begin{tabularx}{0.7\textwidth}{X} 那些人就往約旦河的路上追趕他們，直到渡口。追趕他們的人一出去，城門就關了。 \end{tabularx} \\ \\ \relax
2:8 & \begin{tabularx}{0.7\textwidth}{X} 二人還沒有睡之前，女人就上屋頂，到他們那裡， \end{tabularx} \\ \\ \relax
2:9 & \begin{tabularx}{0.7\textwidth}{X} 對他們說：「我知道耶和華已經把這地賜給你們了，並且我們也都懼怕你們。這地所有的居民在你們面前都融化了。 \end{tabularx} \\ \\ \relax
2:10 & \begin{tabularx}{0.7\textwidth}{X} 因為我們聽見你們出埃及的時候，耶和華怎樣在你們前面使紅海的水乾了，並且你們怎樣處置約旦河東的兩個亞摩利王，西宏和噩，把他們完全消滅。 \end{tabularx} \\ \\ \relax
2:11 & \begin{tabularx}{0.7\textwidth}{X} 我們一聽見就膽戰心驚，人人因你們的緣故勇氣全失。耶和華－你們的神是天上地下的神。 \end{tabularx} \\ \\ \relax
2:12 & \begin{tabularx}{0.7\textwidth}{X} 現在我既然恩待你們，求你們指著耶和華向我起誓，你們也要恩待我的父家。請你們給我一個確實的憑據， \end{tabularx} \\ \\ \relax
2:13 & \begin{tabularx}{0.7\textwidth}{X} 要救活我的父母、兄弟、姊妹，和所有屬他們的，拯救我們的性命脫離死亡。」 \end{tabularx} \\ \\ \relax
2:14 & \begin{tabularx}{0.7\textwidth}{X} 那二人對她說：「我們願意以性命來替你們死。你們若不洩漏我們這件事，當耶和華將這地賜給我們的時候，我們必以慈愛和誠信待你。」 \end{tabularx} \\ \\
[1ex]
\hline
\hline
\end{longtable}
(書二1-14)
\newline
\newline
約書亞記裡面充滿信心的經歷,非常著重信心.神不但要以色列人的十二支派憑信心分地,更要他們憑信爭戰.今天的重點在信心與生活見證,並思想兩個信心人物──喇合和迦勒的生活.
\newline
\newline
(一)喇合
\newline
\newline
從約書亞記第二章的背景中,記載她是一位生活在耶利哥的外邦女子,她是個妓女；無論她是被迫或是自甘墮落而淪為妓女,也反映出她低下的道德生活和教育水平.可是,她卻是大衛王的祖父波阿斯的母親,又是耶穌的先祖.這是由於她有極大的信心,超越了出身背景的卑微,反成為舉足輕重的人物.他更在神面前蒙大福,成為聖經記載偉大的屬靈信心偉人.
\newline
\newline
今日信心的危機就是現代人的危機,何況香港人更是要面對九七現實的轉變呢!聖經所論及的信心危機,不單是政治前途的問題,更涉及永恆的事實,並且是永恆與神的關係.我們分析喇合的信心:
\newline
\newline
(1)喇合聽而相信(書二10-11)
\newline
\newline
「因為我們聽見.......耶和華──你們的神本是上天下地的神.」她雖然未見過神,只是聽聞,就深深相信了.信心的事實就如真實看見一樣.就如(來十一1):「信就是所望之事的實底；是未見之事的確據.」對有信心的人來說,就是聽聞如同看見一樣.在希伯來書十一13:「這些人都是存著信心死的,並沒有得著所應許的,卻從遠處望見,且歡喜迎接.」神應許以色列的先祖,他們的後代如海邊的沙,天上的星那麼多!又說他們必得著迦南地為業；可是到他們死後,他們仍未能得到.這裡形容他們從遠處望見,就如得著一樣,歡喜迎接.在舊約先知書的預言中所用的動詞,大部份是用完成式的語法,表示未發生的事已經完成了；藉以表達先知對神的信心,神要成就的事是必定要成就的.神啟示要成就的事,就用完成式表示已發生了.有些人信耶穌要有證據,或者要看到神蹟才會相信,可是耶穌說:「......那沒有看見就信的有福了.」(約一29下)喇合就是這種沒有看見而相信的人.她聽聞以色列全能神的作為,就相信,這樣的人實在是有福.她的蒙福是因為對神的信心.今日我們信耶穌的人,也是因聽聞見證或道理而相信的.但我們尚未去經歷,故此必須要慢慢經歷聖經真理的真實,又不斷的經歷,建立我們的信仰.信道是從聽道來的,蒙大福的人是聽而相信的.
\newline
\newline
(2)喇合相信掌管上天下地的神(書一11)
\newline
\newline
神是掌管宇宙萬物的主宰,祂不是限制於教會內,卻是滿有威榮.祂掌管自然界的一切,又掌管人類過去,現在,未來的歷史,祂更掌管人類的命運.試問你所相信的神是否某一個時代的神或是地方的神呢?這並不是聖經所啟示的神.我以前唸神學時一位宣教師的老師,曾經講述她父母早期來華從事宣教工作的見證.當時華人信主十分困難,因要冒生命危險的代價；而那對宣教士夫婦工作的地區,又有很多土匪出沒.有一回他們收到土匪的恐嚇信,迫令他們在短期離開那裡,不准他們傳福音；否則在某時某日來放火,燒毀他們的住宅.他們只好禱告交托給神,求神保守他們的家庭和生命,他們並沒有被嚇倒.到了指定的時間,他們安然無恙；可是他們卻接到一封更嚴厲的信,又要在某時某日連禮拜堂一拼燒去.他們只好與教會的弟兄姊妹一同禱告,求神為他們解決困難,顯出祂的榮耀來!到了指定的時間,沒有任何事發生,只是收到第三封信,更揚言要連其他信徒的房屋也一同燒毀.他們都勇敢用信心面對危難,大家同心繼續禱告,求神保守.就在土匪要來的那天晚上,突然颳起狂風雷暴和大雨,結果土匪們沒有前來燒屋,事後也沒有來擾亂,也許是神大能的彰顯,將土匪們嚇走.因為我們的神是掌管上天下地的神,事情的成敗在乎我們對神的認識和信心.我們的神觀是怎樣?我們應效法喇合,堅定的信靠祂.
\newline
\newline
(3)喇合的信心是有行為表現的
\newline
\newline
按雅各書論及信心,是應該有行為的,否則這種信心是死的.就如喇合的信心,產生出行動,她接待,收留和隱藏以色列的探子,她深信那位以色列人相信的耶和華是神.我們是信耶穌的人,我們有沒有見證福音的行為呢?我們的見證有否反映主耶穌的榮美呢?現代的信徒可能有很大的掙扎,喇合雖然沒有好的背景；可是她的信心產生行為,這種行為是冒險和付上生命為代價的.假若耶利哥王知道她窩藏以色列的探子,她就必死無疑!很多時我們信主不是在安舒自由的環境,也許是在逼迫試煉中；無論甚麼環境,我們都應為著我們的信仰,有所表現和付上代價.喇合未相信神之前可能私德也是不好,但她竟然有這麼大的信心.我如今在事奉的小教會中,有一群年長愛主的姊妹,經常禱告托住教會的工作.他們特別看重週三和週五晚上的禱告會；平時他們在家庭也時常禱告,神就與這間教會同在.在這群姊妹中特別有一位,她沒有機會受教育,但非常熱愛讀經,並且特別有信心代禱；每星期都抽一天禁食禱告,逢主日早上很早就到教會預備各種工作.人未必需要很有學問或恩賜,但神可能給我們很大的信心,激勵其他的人,甚而教會的傳道人.信耶穌的人不應有自卑感,在神的眼中同樣看為寶貴,神樂意使用一切願意事奉的人.最低限度我們可以祈禱,成為禱告的勇士.
\newline
\newline
(4)喇合信心的果效
\newline
\newline
我們應該檢討我們的信心,是否成為他人的激勵和幫助,正如喇合一樣呢?她的信心使她全家得救,並且使她成為耶穌基督的先祖.請問各位家中是否有人未信呢?我們應該帶領家人信主,又要效法喇合的信心；發起全家歸主運動使教會不斷增長!喇合在屬靈上產生帶頭作用,我們應該從家裡開始,落實在我們接觸的親友中傳福音,引領他們信主.我們的宣教工場,就是我們的家庭.
\newline
\newline
喇合更影響以色列人的信心(書二23-24):「探子根據她的話對約書亞說:『耶和華果然將那全地交在我們手中；那地的一切居民在我們面前心都消化了.』」當探子回報時,說出這見證,(等別是2-4節對約書亞的說話),明顯是喇合的信心和說話；她這些積極的信息,帶給以色列人極大的鼓勵,特別是對神的僕人約書亞.當我們深入思想時,發覺約書亞記第二章喇合的事件,對攻打耶利哥城有重大的意義.神用一位外邦妓女的信心,去鼓勵祂的子民,這實在是奇妙的事.今日基督徒傳福音最大的阻力,是沒有生活的見證,因為我們的信心不能成為別人的幫助和鼓勵哩.
\newline
\newline
(二)迦勒
\newline
\newline
迦勒是十二支派中猶大支派的一位族長,當摩西差派十二位探子了解迦南情況時,迦勒與約書亞是其中的兩個,並且同時是報喜訊的.他們的信心不斷鼓勵以色列人進迦南,而當時二十歲以上的人只有約書亞和迦勒可以踏進迦南.雖然神選立約書亞為摩西的接棒人,迦勒卻謙卑地到約書亞面前,求取神應許要賜予的希伯崙地業.
\newline
\newline
迦勒的信心在(書十四6-14)有清楚的記載,請看(10-12)「自從耶和華對摩西說這話的時候,耶和華照他所應許的,使我存活這四十五年；看哪,現今我八十五歲了,我還是強壯,像摩西打發我去的那天一樣；無論是我爭戰,是出入,我的力量那時如何,現在還是如何.......」他真是老當益壯,四十年如一日.他對神的應許,充滿信心,更相信那是不會落空的.神的應許在成就前,他繼續仰望,這是他生命的原動力,滿有生存的意志.今日你靠甚麼活著?是你所擁有的,抑或是神的應許,神的話呢?神的應許和話語是永不落空的；從迦勒對神話語的信靠,實在生發無窮的力量,這是他的生存意志的基礎.以色列人四十年在曠野行走,雖然發生了很多動亂,但他沒有受到影響.穆勒從德國移民到英國,他養活了很多孤兒.他開辦孤兒院的工作,從不募捐,他只向神祈求,仰望神的供應.他實在是個有信心的人,有一節聖經,是他的座右銘,那就是詩篇五十五22:「你要把你的重擔卸給耶和華,他必撫養你；他永不叫義人動搖.」這節經文對我的幫助也很大,在不同的譯本中,「撫養」可解作「支持」；就算我們的擔子有多麼重大!神必撫養,支持我們.穆勒就靠這節聖經過他的一生.今天,我們應有神的話語才能堅定我們的信心,並且成為我們生活的動力.迦勒就是那樣堅信神的應許,數十年沒有改變,生命充滿力量.
\newline
\newline
迦勒除了信靠神的應許外,更專心跟隨神.約書亞記提及他專心跟隨有三處的經文:(8,9和14節).第一次是他自己的認信,第二次是摩西的見證,第三次是聖靈的默示──聖經作者的記載.信耶穌具體的表現是跟從耶穌.我們並不是跟從人或是潮流,人會有軟弱使人失望；但神卻永不改變.今日基督徒最大的問題是沒有專心跟隨神,我們有時三心兩意,很容易受環境的影響.有時更「花心」,不能專心跟隨神；靈性就時常起伏不定.有時也會腳踏兩船,或者成為「騎牆派」的人；這樣便使不信者輕視我們,因此就輕視基督教或我們所信的神,我現在分享一個見證；我教會所在附近有一個核桃溪Walnut Creek,大約從三藩市去需時兩小時的車程,那裡有一個有名的劃則師,經營一所則師樓,有關他的見證也曾在「中信」出版的刊物中介紹,他名字是Gerry Loving,他的見證有五個重點:
\newline
\newline
1.十二年前他學習將自己的事業交托給神
\newline
\newline
2.他學習無論做甚麼事情都要誠實
\newline
\newline
3.他凡事藉著禱告交托給神,求問神
\newline
\newline
4.他要遵行什一奉獻,這是包括他任何的收益
\newline
\newline
5.他要領人歸主,後來他帶領了公司一位職員信主；稍後那人更成為他股東之一,使他的業務發展得更蓬勃.
\newline
\newline
這個弟兄所寫的見證題目是:「退出騎牆派」.一個專心跟從神的人是不會做「騎牆派」或「一腳踏兩船」；又要愛神,而又要愛世界的.專心跟隨神的人必定將他的工作交托給神,凡事誠實,不斷禱告仰望神.在奉獻上盡責任,並且不斷帶領人認識耶穌基督.他衡量和決定事情沒有兩種準則,他祇按聖經的標準行事,不論在任何的場合,環境,他始終如一,就是站穩在神話語的原則上.
\newline
\newline
今日我們要像喇合一樣,信心,信仰都能影響他人；並且先從家裡開始,延伸到其他的人身上.
\newline
\newline


\newpage

\section{第60屆~港九培靈會~劉承業~第三講}
\label{sec:1195}
\textbf{憑信的戰爭}
\newline
\newline
連結: \href{https://www.hkbibleconference.org/session-message/view/1195}{\texttt{https://www.hkbibleconference.org/session-message/view/1195}}
\newline
\newline
\hyperref[sec:1194]{< < < PREV SERMON < < <}
~
\hyperref[sec:toc]{[返目錄]}
~
\hyperref[sec:1196]{> > > NEXT SERMON > > >}
\newline
\newline
約書亞記 3:1-3:11
\newline
\begin{longtable}{cl}
\hline
\hline
章節 & 經文 (和合本修訂版)\\
\hline
3:1 & \begin{tabularx}{0.7\textwidth}{X} 約書亞清早起來，和以色列眾人起行，離開什亭，來到約旦河，過河以前住在那裡。 \end{tabularx} \\ \\ \relax
3:2 & \begin{tabularx}{0.7\textwidth}{X} 過了三天，官長走遍營中， \end{tabularx} \\ \\ \relax
3:3 & \begin{tabularx}{0.7\textwidth}{X} 吩咐百姓說：「當你們看見利未家的祭司抬著耶和華－你們神的約櫃的時候，你們就要起行離開所住的地方，跟著約櫃走， \end{tabularx} \\ \\ \relax
3:4 & \begin{tabularx}{0.7\textwidth}{X} 使你們知道所當走的路，因為這條路是你們從來沒有走過的。只是你們要與約櫃相隔約二千肘，不可太靠近約櫃。」 \end{tabularx} \\ \\ \relax
3:5 & \begin{tabularx}{0.7\textwidth}{X} 約書亞吩咐百姓說：「你們要使自己分別為聖，因為明天耶和華必在你們中間行奇事。」 \end{tabularx} \\ \\ \relax
3:6 & \begin{tabularx}{0.7\textwidth}{X} 約書亞對祭司說：「你們抬起約櫃，在百姓的前面過去。」於是他們抬起約櫃，走在百姓前面。 \end{tabularx} \\ \\ \relax
3:7 & \begin{tabularx}{0.7\textwidth}{X} 耶和華對約書亞說：「從今日起，我必使你在以色列眾人眼前被尊為大，使他們知道我怎樣與摩西同在，也必照樣與你同在。 \end{tabularx} \\ \\ \relax
3:8 & \begin{tabularx}{0.7\textwidth}{X} 你要吩咐抬約櫃的祭司說：『你們到了約旦河的水邊，要在約旦河中站著。』」 \end{tabularx} \\ \\ \relax
3:9 & \begin{tabularx}{0.7\textwidth}{X} 約書亞對以色列人說：「你們近前，到這裡來，聽耶和華－你們神的話。」 \end{tabularx} \\ \\ \relax
3:10 & \begin{tabularx}{0.7\textwidth}{X} 約書亞說：「你們因這事會知道永生的神在你們中間，他必從你們面前趕出迦南人、赫人、希未人、比利洗人、革迦撒人、亞摩利人、耶布斯人。 \end{tabularx} \\ \\ \relax
3:11 & \begin{tabularx}{0.7\textwidth}{X} 看哪！全地之主的約櫃必在你們的前面過去，到約旦河裡。 \end{tabularx} \\ \\
[1ex]
\hline
\hline
\end{longtable}
今天與大家思想以色列人的爭戰,這是對我們在事奉上的原則上有關係的.神首先藉著外邦女子喇合的信心,去感動激勵以色列人的兩個探子；使他們回去報積極的喜信,讓以色列民及約書亞得到鼓舞.(書二24):「探子對約書亞說:『耶和華果然將那全地交在我們手中；那地的一切居民在我們面前心都消化.』喇合有信心,這兩個探子也有信心；這樣就鼓勵幫助了其他的以色列人.神就這樣預備了以色列人的信心,預備好他們爭戰.
\newline
\newline
第二步的預備,就是藉著過約但河的過渡.神使河水分開,命以色列人行過河底的乾地,讓他們認定耶和華是掌管上天下地的神.當時除了約書亞和迦勒外,其他的人大都在20歲以下,未曾有渡過紅海的經歷；並且大部份是在曠野出生的.今日神要這新一代的以色列人,看見祂大能的神蹟,使他們的信心大大加強；從而有足夠的力量去爭戰.約書亞在第四章末段,要以色列人在河中取石為記,就是要表達這個意思.(書四23-24):「因為耶和華──你們的神在你們前面使約但河的水乾了,等著你們過來,就如耶和華──你們的神從前在我們前面使紅海乾了,等著我們過來一樣,要使地上萬民都知道,耶和華的手大有能力,也要使你們永遠敬畏耶和華──你們的神」.
\newline
\newline
神是全能全知的,神在以色列人未出戰前,就預備方法,裝備他們；雖然他們曾經不知怎樣過河,但神卻給他們的神蹟,使他們勇敢的前進,信心堅定.在座當中很多是參與事奉的弟兄姊妹；或者有些是剛剛開始事奉,甚而考慮全時間事奉神的；這是很榮耀的事,但卻不是容易的.在一個未開始事奉神的人身上,神已開始動祂的善工,不斷地造就他；使他能勝任愉快,擔負責任.在事奉的經歷中往往不是一帆風順的,甚至有很多困難；但在困難中神卻要介入改變處境,使我們更深入的認識祂的作為和奇妙信實.我記得在我讀神學時,學院有佈道團往醫院探訪；我最失敗是在那裡見到一些病人,不知如何入手開始接獨他們──當病人見到我們拿著單張和聖經接近他們時；他們就用毛巾掩面,表示拒絕的態度,我們只好無奈地離開.後來神容許我妻子病倒,有很長的時間在病床上度過；神就使用這種環境使我學習接觸和服侍病人,或是軟弱的人,成為他們的幫助.今日無論你在怎樣的困難中,可能神在其中有祂的計劃,並且要裝備你,讓你更有效地事奉神.
\newline
\newline
我們來看(書三9-11):「約書亞對以色列人說:『你們近前來,聽耶和華──你們神的話』.約書亞說:『看哪,普天下主的約櫃必在你們前頭過去,到約但河裡,因此你們就知道在你們中間有永生神,並且祂必在你們面前趕出迦南人,赫人......耶布斯人』.」從這幾節經文中,我們看見過約但河與征服迦南地有非常密切的關係.他們先經歷了這個過河的神蹟奇事,讓他們認識神的大能,再使用他們作成趕逐滅絕迦南人的工作.過約但河是全體以色列人的事實,經驗,不論男女老幼,神要全以色列人都在事奉的行列中；事奉者背後要有他們的家族,家人去支持他們,事奉者不能孤軍作戰.所以我們週圍若有親屬事奉神,我們應該盡力為他們代禱?成為他們的支持和鼓勵.以色列人整體的過河經歷,成為他們整體信心的造就.我們若有家庭的,必然要全家作戰,發揮隊工的果效,神的工作就更加興旺.
\newline
\newline
我們現在來看約書亞記第六章,思想以色列人圍攻耶利哥,歸納屬靈爭戰兩個重要的原則:
\newline
\newline
(一)聽從順服神的話
\newline
\newline
在第三章我們知道約書亞已經囑咐以色列人,要他們看見約櫃,就要他們跟在後面,他們就可知道當走的路.今日我們走人生的路,需要聽神的帶領,指引,隨著聖靈的感動和聖經的原則生活；否則我們便沒有資格為神工作.我們不是為了自己的喜好,或一時的興緻去做,因為這樣不是事奉神,充其量只是參與一些宗教的活動而已.事奉神實質上是一種敬拜,不是按私意去親近神；而是用心靈誠實敬拜,這是神的原則.到底神怎樣吩咐以色列人圍攻耶利哥城呢!
\newline
\newline
第一日至第六日,一聲不響地繞城行一週,約櫃在他們當中.今日我們事奉,需要有神的同在──約櫃在他們當中就表現要與事奉者同在.我們的事奉可能是在沒有人注意的地方或是隱藏的事奉,有些卻是為人所觸目的位置,無論如何,他們在神的面前都是事奉神.今日我們的事奉不是要誇耀自己的重要,或是表揚自己的名聲,身價,無論是在教會做名牧；或在小禮拜堂工作,神看我們的標準是在我們是否忠心,並且給予我們同樣的獎賞.今日我們的事奉,就我的看法和盼望,是如約瑟一樣,無論是卑微處坐監或為奴僕,都是神的同工,並且見證神的同在；所以在任何的崗位,哪一種的教會,哪一個角落,我們所求的只是神的同在,就算是最卑微的打掃清潔工作,也是討神喜悅,榮神益人.就如我的先祖父,是在教會中作打掃的工作,因而信了主；並且帶領了整個家屬信主.當我們聽從神的吩咐,就看見神的同在.以色列人圍繞耶利哥城之辦法,按普遍的情理看來；實在有點兒愚昧,甚至是有點荒謬；就如我們發展教會單是靠祈禱和傳福音時.有些人就會抨擊這些方法是落伍的,不切合時代的步伐和需要!可是依靠神,聽從神的說話是萬世不變的,按神的心意去事奉,是超越環境,時間和場合的,也是最有效的方法.
\newline
\newline
第二方面他們要遵守──不能取當滅之物.(書六18-19):「至於你們,務要謹慎,不可取那當滅的物,恐怕你們取了那當滅的物就連累以色列人的全營,使全營受咒詛.惟有金子,銀子,和銅鐵的器皿都要歸耶和華為聖,必入耶和華的庫中.」用現在的字眼來講,就是不能取那些「違禁品」,所有敵人所用的東西──包括一切有生氣的,都是以色列人的「違禁品」,都是應該滅絕的.在古代國家對戰爭的看法,一般視為聖戰.所以未出戰時要拜祭他們的神靈,又是代表神靈和為他們的神去戰爭；並且要把敵人全部都滅絕,將擄物就獻給他們的神靈.以色列人當時的方法不是特別殘忍；從考古學看來,得知當時迦南人的風俗習慣是非常的淫穢,假若以色列人與他們一同生活,就會受他們的影響,從歷史角度看以色列人的問題,也是受迦南人的影響,以致他們的道德宗教都受到污染,因而變質.
\newline
\newline
在以色列人繞城六日之後,到第七日有所行動和他們不能取當滅之物的事情上,表現出一個屬靈的原則,是要屬神的人必須按神的吩咐行事.
\newline
\newline
(二)相信神的話和行為
\newline
\newline
書六2:「耶穌華曉諭約書亞說:『我已經把王,並大能的勇士,都交在你手中.』」我們必須相信神的話和作為!假如以色列人不相信約書亞這番說話,他們就不能成就神的工作.他們相信神的話,他們繞城而行並沒有發怨言,並且順從,沒有其他意見.到第七天中,他們就吹號角,眾百姓就大聲呼喊,城牆就向內塌陷,他們就攻入耶利哥.他們必須有這樣的信心.唯有如此,他們方能得勝.當他們信靠神的作為時,不是袖手旁觀,有行動；而是配合神的工作,彼此有參與,不論是祭司或拿兵器的以色列人.換言之即全職的傳道人或帶職的信徒,都應一起參與；但真正打仗的不是以色列的軍隊,而是神親自為以色列爭戰.當然以色列人有參戰,而神卻親自同工同在,使以色列大獲全勝.
\newline
\newline
我們從以上分析以色列人圍攻耶利哥城,可以歸納出爭戰或得勝事奉的原則:
\newline
\newline
(1)是按神的吩咐和心意而行
\newline
\newline
我有一個敬愛的師長,是劉福群牧師.他從中國內地逃難來到香港,神吩咐他兩件事:第一是建道神學院要復課,第二要編譯青年聖歌.他就遵照神的吩咐,在餘下的時光完成這兩件事工,並且有很好的成績.雖然他當時在很困難的環境中,不知從何開始,他始終如一的堅持到底,神就不斷帶領使用他.我相信很多人曾經看過「十字架與彈簧刀」這本書,書中的作者是在賓汐凡尼亞州一個郊區的小城鎮牧養教會,神感動帶領他去紐約大城市中的貧民區和吸毒者中做基層的福音工作.在人看來他沒有那些條件去承擔這種工作.第一他無大城市工作的經驗,第二他更沒有基層福音工作的經驗.可是他順服神的帶領,結果神與他同工,那裡的事工發展得很快,有很好的成效.所以我們能順服神的帶領,去承擔任何的事奉崗位時,神必定負責他的需用.(太十一29)耶穌說:「我心裡柔和謙卑,你們當負我的軛,學我的樣式；這樣,你們心裡就必得享安息.」甚麼是「我的軛呢?」「軛」是代表工作,即當一個人在神為他安排的工作崗位事奉時,那人已可享安息.當然那人學主的樣式,必須心裡柔和謙卑.這是主給我們的應許,所以各位弟兄姊妹,我們應反省；我們的事奉是否有喜樂和享安息,我們就知道我們是在背負主的軛,當我們忠心去做時,主應許我們必得享安息.
\newline
\newline
神會供給我們一切的需要,祂會給恩典,給能力；甚至給恩賜我們,讓我們可勝任神交付我們的工作.
\newline
\newline
(2)相信神為我們爭戰和作事
\newline
\newline
我記得一位剛畢業的神學生,他來找我,請教我有甚麼進工場的提示心得；我知道這位弟兄比較缺乏自信心,所以我就與他看一段經文.就是(腓四13):「我靠著那加給我力量的,凡事都能作.」他事奉一年後回母院分享早會的事奉經歷時,就分享這節聖經.雖然短短一年,他看見他的事業有神的同在,對於自信心方面也加強了很多.我內心非常高興,但這節聖經同樣對我有很大的幫助.特別我覺得這幾年在北美的事奉中,我更感受到:「唯有靠著那加給我力量的,凡事都能作.」北美教會在同一間教會已有不同的文化；有操不同語言的人,就如要說英語,國語和粵語.自己是教會唯一按立的牧師,所以甚麼東西也要做,一個人要做幾個人的工作.在不同文化,背景,方言的人；若不是靠著那加給我力量的神,我實在無能力面對這重大的挑戰.我很喜歡一節聖經,希望送給那些願意事奉主的肢體的,就是(約五17):「我父做事直到如今,我也做事.」那是耶穌醫好了卅八年的癱子後,與法利賽人辯論時所說的話.我們的主不斷在作事,祂會為你為我去作事──當我們祈求的時候,就得著；當我們尋找的時候,就給我們尋見,叩門的時候,就給我們開門.
\newline
\newline
我以前有一位同工,他去了別處,有時寄一些信來,他很強調神不斷為他作事,他事奉神是看神的作為.我對他的領受很有共鳴.今日在事奉工場中,責任是那麼重大,我們自己有那麼多的難處和軟弱,而且還要去得勝,如果不是靠著神為我們行事,我們如何能擔當呢?尤其是一些在教會事奉的同工,要任勞任怨,別人對他們的要求又特別高,很多時要做到百分之一百的完全；又好像生活在一個金魚缸中,給人觀察.正如保羅所說:「我們好像一台戲,給人觀看.」今日非基督徒也如此看待一些基督徒,看他們怎樣生活,基督徒也是如一台戲一樣.有時我們會被誤會,受別人冤枉；信徒間發生糾紛時,傳道人要做中間人,夾在其中,真是有氣無法宣洩,我們又如何呢!我時常緊記一節聖經,是(伯十九25):「我知道我的救贖主活著.」這裡的「救贖主」原意是「伸辯者」:因為約伯投訴無門,本來他三個朋友是來安慰他,可是他們不諒解他,反而加給他很大的壓力.所以他在信心中如此仰望神:「末了必站在地上.」當你在受委屈時,神不單是拯救我們的神,祂更會為我們伸辯；祂不單在末日才為我們伸辯,就是現在,祂會為我的正直而伸辯.當我們如此了解認識神,我們在事奉中就算遇上任何困難,都能靠主經過,並且得勝.在事奉的崗位上,面對一切超乎我們能力的事情,我們只要等候神,仰望神,就可以解決.在我中學時讀過一首由一位盲人John Milton所寫的詩,其中有一句講及他失明:「那些只是站著和等候的人,也是事奉的人.」特別因他是盲人,他甚麼也不能作,他只有祈禱,他也是在事奉中.
\newline
\newline
我們不要忘記摩西一句很重要的話,(出四14):「耶和華必為你們爭戰；你們只管靜默,不要作聲.」就在以色列人將要過紅海前,摩西勉勵以色列人的話.有一年我收到一本畢業刊,裡面有一位神學院的同工寫了一段話:「抱著為主的心志,現在回來告訴同學,二十年來是神為我作事.」很多時我們事奉,以為犧牲很多,其實我們沒有做甚麼!只是神透過我們這個器皿,去作成祂的工作,這就是屬靈的爭戰,就是憑信的爭戰.當我們相信神,又信靠祂的作為,並按祂的方法和心意而行,我們更凡事順利亨通.求神憐憫我們這不配的器皿,讓我們經歷神的大能力,給神使用.
\newline
\newline


\newpage

\section{第60屆~港九培靈會~劉承業~第四講}
\label{sec:1196}
\textbf{神要作元帥}
\newline
\newline
連結: \href{https://www.hkbibleconference.org/session-message/view/1196}{\texttt{https://www.hkbibleconference.org/session-message/view/1196}}
\newline
\newline
\hyperref[sec:1195]{< < < PREV SERMON < < <}
~
\hyperref[sec:toc]{[返目錄]}
~
\hyperref[sec:1200]{> > > NEXT SERMON > > >}
\newline
\newline
約書亞記 5:1-5:9
\newline
\begin{longtable}{cl}
\hline
\hline
章節 & 經文 (和合本修訂版)\\
\hline
5:1 & \begin{tabularx}{0.7\textwidth}{X} 約旦河西亞摩利人的眾王和靠海迦南人的眾王，聽見耶和華在以色列人前面使約旦河的水乾了，直到他們過了河，眾王因以色列人的緣故都膽戰心驚，勇氣全失。 \end{tabularx} \\ \\ \relax
5:2 & \begin{tabularx}{0.7\textwidth}{X} 那時，耶和華對約書亞說：「你要造火石刀，第二次為以色列人行割禮。」 \end{tabularx} \\ \\ \relax
5:3 & \begin{tabularx}{0.7\textwidth}{X} 約書亞就造了火石刀，在哈爾拉勒山為以色列人行割禮。 \end{tabularx} \\ \\ \relax
5:4 & \begin{tabularx}{0.7\textwidth}{X} 約書亞行割禮的原因是這樣：從埃及出來的眾百姓，所有能打仗的男丁，出了埃及以後，都死在曠野的路上。 \end{tabularx} \\ \\ \relax
5:5 & \begin{tabularx}{0.7\textwidth}{X} 這些從埃及出來的眾百姓都受過割禮；但是那些出埃及以後，在曠野的路上所生的眾百姓卻沒有受過割禮。 \end{tabularx} \\ \\ \relax
5:6 & \begin{tabularx}{0.7\textwidth}{X} 以色列人在曠野走了四十年，直到那從埃及出來，全國能打仗的人都消滅了，因為他們沒有聽從耶和華的話。耶和華曾向他們起誓，必不容許他們看見耶和華向他們列祖起誓要給我們的地，就是流奶與蜜之地。 \end{tabularx} \\ \\ \relax
5:7 & \begin{tabularx}{0.7\textwidth}{X} 他們的子孫，就是耶和華興起接續他們的，都沒有受過割禮；因為在路上他們沒有受割禮，約書亞就為他們行割禮。 \end{tabularx} \\ \\ \relax
5:8 & \begin{tabularx}{0.7\textwidth}{X} 全國的人都受了割禮，留在營中自己的地方，直到痊癒。 \end{tabularx} \\ \\ \relax
5:9 & \begin{tabularx}{0.7\textwidth}{X} 耶和華對約書亞說：「我今日將埃及的羞辱從你們身上除掉了。」因此，那地方名叫吉甲，直到今日。 \end{tabularx} \\ \\
[1ex]
\hline
\hline
\end{longtable}
(書五1-9,13-15)
\newline
\newline
今天要與各位思想第五章的真理,剛才讀出的經文,是關於神要以色列人行割禮的事情；因為在曠野出生的以色列人,大都未受割禮.割禮是神與亞伯拉罕立約時,要求以色列人必定要做的,而這種禮儀,是表示立約的證據,叫他們永遠不會忘記他們的先祖是曾經與神立約,在創世記十七章便詳細記載這割禮立約的背景.(書五9)耶和華對約書亞說:「我今日將埃及的羞辱從你們身上輥去了.」這次約書亞為以色列人行割禮.是與埃及的羞辱有關,若我照這經文與(西二11)所講及,這割禮是把肉體的情慾除掉.這次割禮給以色列人屬靈的意義,是要以色列人不再貪戀埃及給他們的一切.「埃及」,是暗指「世界」,即屬世的事物,人若貪戀那些,就無法參與這場屬靈的戰爭,更無法與撒旦和世界空中的掌權者爭戰.就算他們肯爭戰,也必是有能力,而且會妥協,並很容易與世俗為友.年老的使徒約翰在(約壹二15)提醒我們:「不要愛世界和世界上的事.人若愛世界,愛父的心就不在他裡面了.」在約書亞的時代,神要整個以色列的民族除去埃及的羞辱,更命令他們不再留戀那些事物,要他們一心的歸向神,可能給他們爭戰的得勝.接續割禮的事件後,(在13-15)節是約書亞本人的親身經歷的事蹟.
\newline
\newline
自從以色列人過了約但河後,約書亞的聲名大振,在(書四14):「當那日,耶和華使約書亞在以色列人眾人眼前尊大.在他平生的日子,百姓敬畏他,像從前敬畏摩西一樣.」之後,神吩咐約書亞為以色列人行割禮,當他們痊癒後就要準備攻取重要的耶利哥城.(書五13):「約書亞靠近耶利哥的時候,舉目觀看.」他可能親自視察耶利哥的形勢,好謀算進攻的戰略.「不料,有一個人手裡有拔出來的刀,對面站立.」約書亞很有膽色地走近那人,問他說:「你是幫助我們呢?是幫助我們的敵人呢?」那人回答說:「不是的,我來是要作耶和華軍隊的元帥.」那人的回答是表明他的身份,他來你要作耶和華軍隊的元帥.約書亞一聽到這句話,就俯伏在地下拜.究竟這手裡拿刀的人是誰呢?聖經學者們有不同的解釋,有些人認為他是神差派來的使者或天使長；有些則認為是神自己在肉身上顯現.第一個說法是不對的,因為當約書亞知道祂的身份,就下跪拜祂,若果是天使的話,那麼天使是不能接受人的跪拜的；而唯有神自己,方可成為人敬拜的對象.從這樣看來,第二種說法是較合理.究竟在這個時候,神為甚麼要向約書亞顯現呢?其實神要約書亞知道一件事,雖然他是以色列人最高的元帥,可是實則是神,在以色列背後,真正是以色列人的最高元帥,或者我們會誤解約書亞帶領以色列人過約但河後,便驕傲起來?但聖經沒有顯示出約書亞有這樣的弱點.也許約書亞並沒有自高自大,但神可能給以後讀聖經的人,讓他們知道,無論他們在人的眼中如何被人尊崇,有何等大的恩賜能力；他們都應服在神權柄之下,謙卑下拜.無論是約書亞或是任何的人,若要被神使用,我們一定要謙卑下來,並且承認神是我們的主,我們在他面前只不過是僕人而已.當他們有這種態度和認識時,我們方能被神更大使用.在保羅的身上,亦有同樣的經歷.(林後十二7)說:他自己身上有一根刺,是神給他的,免得他過於自高.雖然那刺是出於撒旦的差役,卻是神准許的.保羅有各種的恩賜,上過三層天,又得到神的啟示,神也很使用他,但神卻給他一根刺,使他能謙卑在神面前,繼續被神使用.神在以色列人攻打迦南地的大城之前,特別要以色列人和約書亞知道,真正以色列人的元帥,是神自己,而不是約書亞.
\newline
\newline
根據(出十二41):「以色列人是耶和華的軍隊.」這個信息是你與我需要得著的.我們必須承認神是我們最高的權威,在我們一切所行的事上,都要認定神自己；否則在任何的屬靈爭戰中,都會失敗.從這段經文中我們可以糾正一些錯誤的神觀.或者有時我們祇看見神是我們的幫助者；我們需要,有困難就尋求神；沒有需要的時候就遠離神.這樣我們不知不覺地做了主人,例如我們家裡很多事要做時,就請一個菲傭來幫助,但她始終是工人,是個幫助者；我們不請她時,就要她辭工.因此我們不應單視神為我們的幫助者.我們可能用很多方法和其他的幫助而替代和忽略了神.這也許是現代人最大的問題!大家都聽過阿拉丁神燈的神話:「當有需要時,就擦那神燈,有巨人出來,為擦他出來的主人做事,這樣那巨人是有求必應的.」我們的神雖然是全能的神,我們不應有這樣錯誤的思想,以免利用神,得罪神.耶穌基督不是那樣的神,我們必須尊神為我們的主,我們的王.正確的觀念應該如保羅所說的:「我是耶穌的僕人.」這裡的「僕人」有「奴僕」的含意,意即無自己的主權,完全由神管理,支配,受著祂所指揮.神是我們的主宰,我們的前途及一切,都在祂手中,任憑祂的帶領,指示.舊約聖經論及奴僕時,在(出廿一1-6):你在百姓面前所要立的典章是這樣:「你若買希伯來人作奴僕,......他就永遠服事主人.」真正的奴僕是永遠屬於主人的,並且服事主人；連他的妻子和兒女,都是屬於他的主人.今日如果我們真正認識我們在耶穌基督面前的身份,我們雖然是神的兒女,也是主人的奴僕.(羅馬書六章及第八章),我們願意將自由交給耶穌；好像舊約那個奴僕,十分敬愛他的主人,所以六年的工作完畢,第七年他應該可以離去.但那僕人願意終身追隨服事他的主人,主人便帶他到審判官那裡；用錐子穿他的耳朵,那奴僕就可以永遠服事主人.今日我們就如這所說的奴僕,我們有自由可以犯罪,但因神那麼愛我,並且拯救了我,我願意將生命主權交出來,永遠屬於神.也許有人看過邊雲波弟兄所寫的小冊「獻給無名的傳道者」,內中寫著:「是自己底手,甘心放下享受,是自己底腳,甘心到這道路上來奔跑.」每一個跟從耶穌的人,是自己甘心情願放下那份自由,將我們的主權交給耶穌.
\newline
\newline
(羅十二1):「所以弟兄們,我以神的慈悲勸你們,將身體獻上,當作活祭；是聖潔的,是神所喜悅的,你們如此事奉乃是理所當然的.」「祭品」是沒有自己主權的,祭牲是主人所支配的；今日我們在神面前,就好像獻上給神的祭品一樣,將自己交給耶穌,這是基於神的大愛.在羅馬書一1開始,神就表明祂偉大的救恩,神在基督耶穌裡所做的一切,所以我們對祂的愛,應該有所回應.我們可以活生生行上祭壇,願意被神收納,帶領,造就和使用我們；可是這個活祭亦可從祭壇走下來,神沒有勉強我們這樣做.這活祭還是一個意思,表示我們獻上是有一個實質的,是一個獻上的生活,用生活作為獻上的基礎,更是一種生命的見證.
\newline
\newline
我記得年青的時候,在香港社會福利處工作,雖然自己生長在一個基督教家庭,但只是傳統地接受這種宗教信仰.有一日我不知為甚麼有這種思想,我是否需要凡事依靠神?我自己不覺自問,為甚麼一時間會那麼宗教化?因為神離我實際上太遠了!自己祇有一笑置之.但過了幾個星期後,這個問題仍不斷重複徘徊在我的思想裡,我何必凡事依靠神,這是老人家的老生常談.有甚麼問題自己擔當便是.結果有一次放工的時候,在黃大仙巴士站附近,為這個問題思想了個多小時,我覺得必須要解決這個問題,不被它再煩擾我.當時神給我看見自己十分不可靠,並且被很多事物所限制,神要我靠祂是要賜福給我,後來我領會依靠神不是以往的錯誤看法,甚麼都不作!依靠神讓神帶領我們行一條正路；今日的青年人應該好自為之,否則很容易成為迷途的人.耶穌說:「我是道路,真理,生命,若不是藉著我,沒有人能到父那裡去.」(約14止)我們的一生需要有神成為我們的引導者.若我們沒有將我們的生命主權交出來時,神就不能帶領我們的前路,否則我們就會成為浪子,受很大的痛苦和損失；後來我不單只凡事依靠神,並且要追求凡事榮耀神,討神喜悅,這樣才是跟從神的人.自從我那次決定順服,凡事依靠神,神就是不斷引導指引我,我的生命直到如今,也蒙神帶領；三十年前那一天的決定,我到現在也沒有後悔,耶穌時常與我同在,我不斷經歷神的大能大力,並且對神的旨意有更多的了解和認識.(羅十二2)跟著成為活祭的真理講論後,指出:「不要效法這個世界,只要心意更新而變化,叫你們察驗何為神的善良,純全可喜悅的旨意.」很多時我們希望尋求神的旨意卻是模糊不清,是因為我們沒有將生命的主權交給祂,過活祭的生活；其二是我們效法這個世界,沒有心意更新而變化,導致我們不能明瞭神的旨意,今日的教會不能復興,是因為很多信徒沒有將生命的主權交給耶穌,很多人以為(羅十二1-2)是神對那些全職事奉者或神學生的教訓,其實那段經文是給所有信徒的.「弟兄們」是指所有信主的人.我們不論是否參與事奉,或在社會中作見證,若要發揮更大的影響力；就必須正視生命主權的前題,我們必須肯交出主權給神.今天香港面對的問題,正是主權問題；這個問題影響每一個住在香港的人!我們有否尊主為大?我們很多屬靈上的問題都在這關鍵性的主權問題上!
\newline
\newline
我們再看(書五14下-15):「約書亞就俯伏在地下拜,說:『我主有甚麼吩咐僕人.』耶和華軍隊的元帥對約書亞說:『把你腳上的鞋脫下來,因為你所站的地方是聖的.』約書亞就照著行了.約書亞在經文中表現順服的回應態度,從他回應中,可知道他對神完全的順服:
\newline
\newline
(一)敬拜的態度
\newline
\newline
他俯服在地下拜.我們無論在甚麼時候,環境中,我們為人都應存著對神的敬畏,因為我們是活在神面前.大衛說:「我將神擺在我的面前.」我們無論在哪裡,都不能逃避神,我們應樂意親近敬畏神,大衛的得勝秘訣就是有親近和敬畏神的心態.保羅亦說:「你們或吃或喝,都要為榮耀神而作.」這是一種隨時隨地的心態,要面對神,一切的作事,講話,思想有神的鑑察,審問!大衛在(詩十六2下)說:「你是我的主,我的好處不在你以外.」有一首詩就用這節聖經唱出來,使我們認識神的主權,當我們有成就時可歸榮耀給神.另外有一節聖經是施洗約翰所講的話:「他必興旺,我必衰微.」(約三20)
\newline
\newline
(二)尋求主的吩咐
\newline
\newline
「主有甚麼吩咐僕人」,今日你有沒有尋求主的吩咐?你有沒有從靈修知道主的吩咐?我們要在神面前,等候神發施號令,就不會有太多的差錯,做事亦不缺乏自信.我們有神的指引就會很安心和放膽做事,就算別人反對或輕視那些工作,我們亦能快樂地完成,並且能討神喜悅.
\newline
\newline
(三)放棄權利,移交權柄
\newline
\newline
「把你腳上的鞋脫下來,因為你所站的地方是聖的.」約書亞照著做了.為甚麼要脫鞋呢?這裡有深奧的意思.「脫鞋」意即放棄權利,移交權柄.請大家看(得四7-10):「從前,在以色列人中要定奪甚麼事?或贖回,或交易,這人就脫鞋給那人.......」路得和拿俄米是寡婦,回到猶太伯利恆,遇到波阿斯那個大財主,要求波阿斯盡近親的本份,娶路得為妻而為其親屬留後代；但有一個人比波阿斯更親的,所以那放棄娶路得之權利,方能讓波阿斯娶路得.在城門口那親屬決定放棄權利,便脫鞋表示將權利移交波阿斯,去娶路得.元帥要約書亞脫鞋,因那地是聖的──是屬神的.表示約書亞應將自己權利交給神.摩西在何烈山上見到異象,神亦吩咐他脫鞋,因他所站的地是聖的.神要求約書亞有一個行動的表示,要將權柄交給神.雖然他在以色列人的眼中是最偉大的領袖,擁有無上的權威,他卻要服在神之下,受耶和華──神的指揮.
\newline
\newline
我認識一位唸完博士學位的姊妹,她受感動要往台灣工業女中宣教,並且參加了台灣工業福音團契,為了解女工的生活和心態,她自己也當起女工來.在面試時,女工不能與上師握手,彼此的階級分得很清楚,雖然她的家人很希望她留在美國,但她知道神帶領她去那裡服事那些女工,台灣的工業福音工作有很大的需要.這位姊妹真的將她的主權交出來,任由神的差遣指引；我曾經與她一同事奉,彼此為同工,覺得她非常好.
\newline
\newline
戴德生先生往中國宣教之前,曾經有一位未婚妻,不贊成他前來中國傳道；後來他為了順服神,於是彼此分手,他便獨自來中國宣教.一個願意將主權交托給神的人,必定有行動上的表現.我們必須在一切的抉擇中,讓神作主作王,交朋,婚姻等決定,您有沒有在神的面前,等候祂說:「我主我神,有甚麼吩咐我做呢?」我們是否願意,如約書亞一樣,按神的心意行事,給神更大的使用呢?
\newline
\newline


\newpage

\section{第60屆~港九培靈會~劉承業~第五講}
\label{sec:1200}
\textbf{聖潔的重要}
\newline
\newline
連結: \href{https://www.hkbibleconference.org/session-message/view/1200}{\texttt{https://www.hkbibleconference.org/session-message/view/1200}}
\newline
\newline
\hyperref[sec:1196]{< < < PREV SERMON < < <}
~
\hyperref[sec:toc]{[返目錄]}
~
\hyperref[sec:1202]{> > > NEXT SERMON > > >}
\newline
\newline
約書亞記 7:1-7:9
\newline
\begin{longtable}{cl}
\hline
\hline
章節 & 經文 (和合本修訂版)\\
\hline
7:1 & \begin{tabularx}{0.7\textwidth}{X} 以色列人在當滅之物上犯了罪。猶大支派中，謝拉的曾孫，撒底的孫子，迦米的兒子亞干取了當滅之物，耶和華的怒氣就向以色列人發作。 \end{tabularx} \\ \\ \relax
7:2 & \begin{tabularx}{0.7\textwidth}{X} 約書亞從耶利哥派人往伯特利東邊，靠近伯‧亞文的艾城去，對他們說：「你們上去窺探那地。」那些人就上去窺探艾城。 \end{tabularx} \\ \\ \relax
7:3 & \begin{tabularx}{0.7\textwidth}{X} 他們回到約書亞那裡，對他說：「眾百姓不必都上去，只要二、三千人上去就能攻取艾城；不必勞動眾百姓都上去，因為他們人少。」 \end{tabularx} \\ \\ \relax
7:4 & \begin{tabularx}{0.7\textwidth}{X} 於是百姓中約有三千人上那裡去，但他們竟在艾城的人面前逃跑。 \end{tabularx} \\ \\ \relax
7:5 & \begin{tabularx}{0.7\textwidth}{X} 艾城的人擊殺他們約三十六人，從城門前追趕他們，直到示巴琳，在下坡的地方擊敗他們。他們都膽戰心驚，融化如水。 \end{tabularx} \\ \\ \relax
7:6 & \begin{tabularx}{0.7\textwidth}{X} 約書亞和以色列的長老就撕裂衣服，在耶和華的約櫃前臉伏於地，直到晚上。他們把灰撒在頭上。 \end{tabularx} \\ \\ \relax
7:7 & \begin{tabularx}{0.7\textwidth}{X} 約書亞說：「唉！主耶和華啊，你為甚麼領這百姓過約旦河，把我們交在亞摩利人手中，使我們滅亡呢？我們不如住在約旦河的那邊！ \end{tabularx} \\ \\ \relax
7:8 & \begin{tabularx}{0.7\textwidth}{X} 主啊，求求你，以色列人既在仇敵面前轉身逃跑，我還有甚麼可說的呢？ \end{tabularx} \\ \\ \relax
7:9 & \begin{tabularx}{0.7\textwidth}{X} 迦南人和這地所有的居民聽見了就必圍困我們，把我們的名從地上除去。那時，你為你至大的名要怎樣做呢？」 \end{tabularx} \\ \\
[1ex]
\hline
\hline
\end{longtable}
(書七1-9)
\newline
\newline
(一)犯罪的惡果
\newline
\newline
帖前四3:「神的旨意是要你們成為聖潔,遠避淫行.」(彼前一16):「你們要聖潔,因為我是聖潔的.」神雖然愛罪人,但祂卻恨惡罪惡,並且是一位忌邪施報的神.我們剛才讀約書亞第七章時,見到以色列人失敗,這失敗之前是一次輝煌的勝利,第六章與第七章記載的事奉形成一個強烈的對比.第六章所攻陷的耶利哥城,比較七章的艾城堅固,強大,可是以色列人卻大獲全勝,反而在艾城遇到慘痛的挫敗──當以色列人的探子偵察艾城的虛實,發覺艾城很少,所居住的人也很小:「眾民不必都上去,只要二,三千人上去就能攻取艾城；不必勞累眾民都去,因為那裡的人少.」以色列人敗在艾城手下,不是輕敵,是第七章一節:「以色列人在當滅的物上犯了罪......迦米的兒子亞干取了當滅的物；神的怒氣就向以色列人發作.」神吩咐以色列人不要取耶利哥任何的東西,但亞干卻犯罪,取了當滅之物,據為己有,所以神的怒氣向他們發作,以色列人就失敗了.我們再看10-12節:「耶和華吩咐約書亞說:『起來!你為何這樣俯伏在地上呢?以色列人犯了罪,違背了我所吩咐他們的約,取了當滅的物；又偷竊,又作詭詐,又把那當滅的物放在他們的家具裡.因此,以色列人在仇敵面前站不住.他們在仇敵面前轉背逃跑,是因成了被咒詛的；你們若不把當滅的物從你們中間除掉,我就不再與你們同在了.』以色列人中有人犯罪,失去神的同在,這是以色列人失敗最重要的原因.我們從以色列人的歷史中,看見這個原則更加明顯；他們敬畏神,神就與他們同在,使用他們.當他們敬拜偶像,犯罪跌倒,就失去神的同在,被外邦人所征服.他們甚至被擄到外邦,成為國破家亡的一群.整本舊約聖經都記載,假若人不順服神,敬拜神,潔淨保守自己；人就軟弱犯罪,生活一敗塗地.罪對一個人的屬靈生命,如疾病對肉體自然生命一樣；又如一個人帶病去打仗,他就毫無力量,自身也難保!一個靈裡犯罪的人,不肯認罪悔改,他就無法面對魔鬼撒旦,更無法得勝!他沒有能力服事主,更沒有能力作美好的見證!我們看見很多事奉神的人,他們開始事奉的時候神很使用他們；後來因為一些事情犯罪跌倒,結果神就離棄他們,他們也不再為主所使用,聽講有一位事奉神的工人,他犯了姦淫,他在一間基督教學校事奉,每次參加聚會就坐在最後一行之一角,他不是參與聚會敬拜神,而是在翻查閱報紙,他根本無心聽道,他與神沒有正常的關係,因為他在靈裡有罪,後來他就不知去向了.
\newline
\newline
一個人犯了罪,又不願悔改對付罪,就必定大大影響他的見證和事奉的能力.最近在美國有幾位出名的電視佈道家,本來他們都很受人歡迎,使人得到很多益處.可是他們私生活上不檢點,令他們失敗跌倒；因此使很多支持愛護他們的肢體失望,還給教外的人士有機會藉此毀謗神的名.所以我們每一位願意事奉,作主精兵的人,必須小心謹慎我們的見證生活!或許我們在事奉上滿有果效,有各樣的恩賜成就各種事工；但忽略了聖潔的生活,我們就會一方面去建造,另一方卻在拆毀!甚而拆毀比建造更大.醫生在施手術之前須作好各樣準備,進行徹底的消毒,然後進入手術室；而在手術中,也必須嚴格全面地消毒,使病人不會被任何病毒細菌所感染,因為這是一個拯救生命的工作.傳福音拯救靈魂也是關乎拯救生命的工作,這不但關係到人短暫的生命,而是關乎永恆生命的,所以事奉者必須自潔.(賽五十11下):「你們扛抬耶和華器皿的人哪.」有一個工人不小心把主人心愛的花瓶弄破,他就把裂口補好,若不細察看是看不見裂縫的；可是,這個秘密被一個工人看見,為了遮掩這件過失,他便要替人多做很多工作,以致那人不告發這事情；後來那人更諸多苛求,使他透不過氣來,他好像成了那人的奴隸!這是因為他一直沒有清理打破花瓶的事.今日若我們犯錯,得罪神之後,又不願意承認,我們就會成了罪惡的奴隸；魔鬼撒旦就在我們身上控告我們.我們的把柄也在牠手裡,我們就不斷受指責,並且慢慢失去力量；神的道不能在我們身上有出路.每一個城市都有一些黑市居民,他們被迫要接受很多惡劣的待遇,他們的工作和生活都在違法進行,時常有被告發的危機,而被解出境.為甚麼他們的遭遇那麼悽慘!因為他們是非法入境者,所以一天這個問題沒有解決,他們就要悲慘度日.如果我們有罪,被神光照,就要認罪悔改,方能享受正常基督徒生活.若我們繼續在罪中生活,不願自潔,我們必定被撒旦所箝制,在罪中被捆綁,神的能力就不能在我們身上發揮,我們亦失去爭戰的能力.
\newline
\newline
(二)一個犯罪可影響全群
\newline
\newline
我們繼續(書七24-26):「眾人把謝拉的曾孫亞干和那銀子,那件衣服,那條金子,並亞干的兒女,牛,驢,羊,帳棚以及他所有的,都帶到亞割谷去.說:『你為甚麼連累我們呢?今日耶和華必叫你受連累.』:因此那地方叫亞割就是連累的意思,直到今日.」一個人犯罪,影響整個以色列民族!
\newline
\newline
大衛一個人犯姦淫,影響他整個家族,先知拿單傳遞神的責備說話:「刀劍必永不離開你的家!」(撒下十二10下)所以他的兒子暗嫩就嫖污了同父異母的妹妹他瑪；他瑪的兄長押沙龍就報復,把暗嫩殺死；後來押沙龍叛變又被殺；結果他兩個親生的兒子就死了.他與拔示巴所生的第一個兒子又死去；最後,在他晚年時,他的另一個兒子亞多尼雅作反,結果被殺.以色列人的律法規定,若偷盜被定罪,尚且要倍償最少四倍；神是很公義的!雖然大衛是一個合神心意的王,神不會縱容他犯罪,並要按他的公義去追討一切犯罪者.
\newline
\newline
在以色列的歷史記載中,有一個基本的通則,就是王若行神眼中看為正的事,他們是好的王,全國就復興受惠,若王行得不好,甚而犯罪,就會帶來國家民族的災殃.所以一個人犯罪,就影響整個團體.為何一個人犯罪會影響全體呢?這又是否合理呢?若按常理:「一人做事,一人當」,為何王犯罪而他的子民都要受罰呢?亞干犯罪,理應單罰亞干,為何又牽連他全家呢?其實在這裡有一個很重要的真理；我們彼此在主裡並耶穌基督的教會中,我們有互為肢體的關係.我們雖然有很多不同,例如不同背景,不同經歷,不同看法,不同時代等......但我們在基督裡,卻成為一個身體；我們彼此相稱弟兄姊妹,在靈裡是同一個家庭.
\newline
\newline
我們彼此有親密的關係,生命上的關係；我們身體上任何一個部份有傷痛,就會影響全身!而一個肢體受苦,全部都受苦；一個肢體得榮耀,全體就得榮耀.我們要明白這個真理,方能更有效實踐彼此相愛,互相接納；否則我們便會相咬相吞.就算在身體上有缺點,我們亦會接納,並且保護他,愛惜他；將加倍的體面加上去,彌補他的不足.故此我們就不會分門別類,生發不同的批評!甚至互相敵視.教會是神的家,假若用一個家庭角度去看,家中一個人有問題,全家就受影響!就如一個家人患了重病,全家在經濟上,情緒上,精神上都要彼此承擔,受著一定的影響.若家人中有人吸毒犯罪,被下在監獄,全家都丟臉子.若家裡有人有好表現,得到很高的榮譽,彼此就會歡喜快樂.我們彼此是息息相關,互相影響；我們在行為生活上要謹慎小心,若自己犯罪不潔,就會影響其他主內的肢體,神就不能與我們整體同在.就如舊約中的該隱,他曾說:「我豈是看守我兄弟的嗎?」(創四9下)我們對弟兄姊妹有屬靈上責任.雖然大家都已信主,但卻要藉著彼此信心的見證和生活,互相激勵,一同長進,並且作戰.
\newline
\newline
(三)不要忽視微小的罪和罪的連鎖性
\newline
\newline
請看(書七11):「以色列人犯了罪,違背了我所吩咐他們的約,取了當滅的物；又偷竊,又行詭詐,......」所以一件的罪,引發另一件的罪,再引致第三件.......亞干的犯罪,就是一個清楚的寫照.他私自留下當滅之物,詐作若無其事,他也許裝作很屬靈；其實他內裡充滿詭詐,更是一個虛偽的人.但神知道他隱藏的罪.那些物件本來不是他的,他卻私取,這就是偷竊.例如有一些人一同去「做世界」,作不法之事,事情被發覺,就殺人滅口,為爭取贓物的利益,又把同黨殺掉,像這樣的人,罪惡就一件連一件的犯了；逾犯逾多而不能自拔!你在小事上犯罪,犯罪的膽量就會大起來,逐漸犯上更大的罪.俗語說:「少時偷針,大時偷金」,這是心態上的改變,進而成為習慣.有一個故事,講及一個中東人,在沙漠中生活,在晚上睡覺時,他的駱駝向他講話:「主人主人,我在外很冷,恐怕我很快著涼生病,就不能在明天起行!可否讓我將鼻子伸入帳幕,好使我不致冷倒?」主人便答應牠的要求.過了片刻,牠又講話:「可否讓我的頸項進入?」主人又答應了.及後牠不斷要求,要使前腳,全身和後腳都一同進入帳幕裡,最後整個帳幕就倒塌下來.這個故事給我們知道,罪惡引誘是會得寸進尺,起初是慢慢地工作,就如羅得往所多瑪挪移帳棚一樣,是漸進的,後來更住在那裡,這是他沾染的過程.我在外國進修時,有一節聖經給我很重要的提醒,那是(路廿二31):「主又說:『西門!西門!撒旦想要得著你們,好篩你們像篩麥子一樣.』」耶穌在彼得未跌倒不認主前曾提醒他,撒旦攻擊他們如篩麥子一樣.若你在農村生活就會見過農夫篩稻米,把榖篩去,留下粗壯的米粒.農夫篩稻米也是慢慢來,就如魔鬼撒旦對付我們一樣!牠不會一下子叫我們離棄神,或是很大的反叛,而是讓我們不自覺時漸漸遠離神；牠經常用懷柔手段,開始時影響我們的靈修生活,使我們不重視讀經祈禱,又輕看參與聚會,對講員諸多評價和不滿,慢慢我們就會減少參加聚會；或在聚會中心不在焉,最後甚而停止聚會,漸漸地我們的靈性就退步.
\newline
\newline
我記得在讀書的年代是寄宿的,有一年與一位弟兄同房；那位弟兄要負責清潔的工作,有一次有人送他一大罐植物油,他便放那罐油在廚房,給大家共用.在幾天前我曾經買了一罐較小的植物油,我那時就起了貪念,把自己的油收起來,單是用那位弟兄拿出來公用的油.在這件小事上,我實在有軟弱,晚上靈修時神就責備我,我心裡感到很不平安；並感到不能專心讀經和禱告,心靈中似乎很不舒服.過了一段時間,當我反省時,神就指出我的貪心,有自己的東西不用,而佔人的便宜.那晚我就祈禱認罪,並且在第二天一早起來時,就向那位同學道歉,承認自己的過錯.經過這事後,我的靈修就沒有甚麼困難.回想起來,我深深感到神憐憫我,當我們在思想,態度或行為上有不對的地方,聖靈就指出我隱而未現的罪；除去我與神相交的攔阻,及與其他肢體關係的過失.所以我們不要忽視很微小的罪,並且要正視和對付一切微小的罪.
\newline
\newline
(四)聖潔的標準
\newline
\newline
教會或信徒對聖潔有很多不同的看法,就聖經的真理去思想,我們可以從舊約利未記得很多寶貴的指引.請看(利十一44-45):「我是耶和華──你們的神；所以你們要成為聖潔,因為我是聖潔的.......我是把你們從埃及地領出來的耶和華,要作你們的神；所以你們要聖潔,因為我是聖潔的.」
\newline
\newline
1.聖潔是與神保持正常的關係──這是聖潔的基礎.聖徒的成聖,表示聖徒是屬於神的；聖物就是屬神的物件.我們既是屬神的,神在我們身上就有祂的主權,去管理,支配和使用我們；在我們的生活上,就能聖潔,這是地位上的聖潔.
\newline
\newline
2.聖潔是分別為聖──這是本質上的聖潔.請看(利廿26:「你們要歸我為聖,因為我──耶和華是聖的,並叫你們與萬民有分別,使你們作我的民.」我們要與其他的人有所分別!這是聖潔的標準:(雅四4):「......豈不知與世俗為友就是與神為敵嗎?所以凡想要與世俗為友的,就是與神為敵了.」在那裡雅各重複說明世俗與神是對立的；今天我們的生活要堅持聖經的原則；在(林後六14):「你們和不信的原不相配,不要同負一軛.」這裡是論及合作的事情,不論是一同做生意或尋找婚姻的伴侶,或是教會事奉的原則；信主與不信的合作,是聖經所否定.在教會的建立和發展上,一切不合乎聖經原則的屬世手段和方法,理應禁止使用.我過往事奉的教會,每一年的聖誕節都送一些罐頭食物給一些貧民區的居民,及一些聖誕禮物給兒童；當時有些人提議要發動整個城中,募捐更多的金錢物資去作成這件事；我不贊成這做法,因信與不信不能合作,我們應持守這原則.我們可去信各教會,用基督徒的名義去做,去作見證,我們也可按我們的力量去工作.我們在個人生活和事奉中,不能忽略別人所不注重的小罪,很多時候,小小的事情就會成為大錯.
\newline
\newline
最後請看(書八1-2):「耶和華對約書亞說:不要懼怕,也不要驚惶.......你要在城後設下伏兵.」這裡我們看見以色列人按著神的吩咐,把亞干的事件處理,除去連累他們的罪；神就再一次給予約書亞及以色列民應許,將艾城的王,人民,和地都交在他們手中.並且指示他們作戰的策略,要他們在城後設下伏兵.當我們肯像以色列人一樣,願意認罪悔改時,我們便得到神的教導和指引；我們就知道如何爭戰,並且經歷得勝.大衛雖然犯了姦淫的罪,在詩篇五十一篇表露他悔罪的心態,他經歷神不會看輕憂傷的靈,神就饒恕及再使用他.當教會能全體認罪悔改時,神就賜下大復興；或者領袖或個人願意認罪悔改,復興仍然是可能的.有一間教會的一位信徒,他在崇拜聚會後被神感動,承認自己的罪；並在那時站立向會眾承認,他如何假冒為善.奉獻時只是空手放進奉獻袋中,實在是犯罪；請大家原諒他,並且他亦在神面前悔改.想不到又有一人起來,承認自己同樣犯了這罪,再有其他的人就承認所犯的其他罪,又願意認罪悔改.自此以後,那個教會得到很大的復興.今天我們從以色列人艾城之役,得到一個得勝的原則,就是我們要聖潔,竭力保守與神的關係,主動的分別為聖,順服神的主權；願我們都成為合神心意的子民和教會,得著大的復興.
\newline
\newline


\newpage

\section{第60屆~港九培靈會~劉承業~第六講}
\label{sec:1202}
\textbf{屬靈爭戰的策略}
\newline
\newline
連結: \href{https://www.hkbibleconference.org/session-message/view/1202}{\texttt{https://www.hkbibleconference.org/session-message/view/1202}}
\newline
\newline
\hyperref[sec:1200]{< < < PREV SERMON < < <}
~
\hyperref[sec:toc]{[返目錄]}
~
\hyperref[sec:1205]{> > > NEXT SERMON > > >}
\newline
\newline
以弗所書 6:14-6:20
\newline
\begin{longtable}{cl}
\hline
\hline
章節 & 經文 (和合本修訂版)\\
\hline
6:14 & \begin{tabularx}{0.7\textwidth}{X} 所以，要站穩了，用真理當作帶子束腰，用公義當作護心鏡遮胸， \end{tabularx} \\ \\ \relax
6:15 & \begin{tabularx}{0.7\textwidth}{X} 又用和平的福音當作預備走路的鞋穿在腳上。 \end{tabularx} \\ \\ \relax
6:16 & \begin{tabularx}{0.7\textwidth}{X} 此外，要拿信德當作盾牌，用來撲滅那惡者一切燒著的箭。 \end{tabularx} \\ \\ \relax
6:17 & \begin{tabularx}{0.7\textwidth}{X} 要戴上救恩的頭盔，拿著聖靈的寶劍—就是神的道。 \end{tabularx} \\ \\ \relax
6:18 & \begin{tabularx}{0.7\textwidth}{X} 要靠著聖靈，隨時多方禱告祈求，並要為此警醒不倦，為眾聖徒祈求。 \end{tabularx} \\ \\ \relax
6:19 & \begin{tabularx}{0.7\textwidth}{X} 也要為我祈求，讓我有口才，能放膽開口講明福音的奧祕， \end{tabularx} \\ \\ \relax
6:20 & \begin{tabularx}{0.7\textwidth}{X} 我為這福音的奧祕作了帶鐵鏈的使者，讓我能照著當盡的本分放膽宣講。 \end{tabularx} \\ \\
[1ex]
\hline
\hline
\end{longtable}
乙．屬靈爭戰的策略
\newline
\newline
(一)禱告的事奉(弗六14-20)
\newline
\newline
以上的五天,我與各位思想屬靈爭戰的先決條件,亦能在這段經文看到同樣的原則；如遵行神的話,就比對14節「用真理當作帶子束腰」及17節「拿著聖靈的寶劍,就是神的道」；論及信心的使用,就比對16節「拿著信德當作籐牌,可以滅盡那惡者一切的火箭」；論及聖潔的重要,亦可比對14節下「用公義當作護心鏡遮胸」.今天我們要特別思想弗六18-20,禱告的事奉:「靠著聖靈,隨時多方禱告祈求」.
\newline
\newline
論及屬靈爭戰的策略,提及禱告的事奉在約書亞的記載中可能不太多.正如第九章指出約書亞及以色列人在與基遍人立約的事上,並沒有求問耶和華.(書九14)結果他們被基遍人欺騙.在第十章,他們大敗亞摩利人的五王,在12節約書亞就向神禱告,叫日月停止:「日頭啊,你要停在基遍；月亮啊,你要止在亞雅崙谷.於是日頭停留,月亮止住,直等國民向敵人報仇.」在這件事件中,有一個非常重要的總結在14節:「在這日以前,這日以後,耶和華聽人禱告,沒有像這日的,是因耶和華為以色列爭戰.」禱告與爭戰,是不能分割的.在使徒時代,初期教會期間,(徒六4):「但我們要以專心祈禱傳道為事.」使徒們看事奉時,重點就在祈禱和傳道上,並且是祈禱在先,更表明祈禱事奉的重要.
\newline
\newline
在路二36-37記載一位女先知亞拿,雖然少年時就守寡,但她卻一生在聖殿中事奉神,主要的工作是──禁食祈求,晝夜事奉神.她事奉神的年日,更是八十四年之久,並不斷儆醒,等候救贖主彌賽亞的來臨；最後,她就親眼看這位救贖主──耶穌,她實在是有福!建道神學院在梧州的時期是聖經學院,學生分男女院住宿；女院長是羅聖愛女士,她亦是培靈研經大會的創辦人之一.她退休後便到一所老人院安度晚年.她立志為每一位在建道就讀的神學生禱告；所以每年學期開課,她便取得就讀同學的名單,每一天為那些學生提名禱告.當同學畢業時；就必定去探望羅女士以表謝意.這位長者雖然年紀老邁,仍然願意為神工作,並且在禱告事奉上,付上極大的代價.
\newline
\newline
今天在各地事奉,被神重用的僕人,使女,都是十分看重禱告.禱告事奉最好的榜樣,是我們的主耶穌.希伯來書形容祂是大祭司,是長遠活著,為我們不斷禱告.當我們軟弱的時候,耶穌這位大祭司和聖靈,都會為我們代禱,使我們除去懼怕,大發熱心,神沒有忘記和撇棄我們,由一八五八至五九年在基督教歷史中有一次大復興,是與祈禱有直接關係的；當時教會在一個非常軟弱的光景中,有一位年青的傳道人被派往紐約工作,並且開始一個中午的禱告會,讓下班的人士參加,第一次的聚會在開始後二十分鐘才有第一位參加,三十分鐘有第二位來,當天的禱告會共有四位參加；以後這個祈禱繼續下去,並且有很多人參加.聚會地方擠滿了人群,後來整個紐約市有很多這種午間祈禱會；慢慢其他的城市,都有這樣的祈禱會.在祈禱會中很多人蒙恩得救,並分享見證,得著更新和復興.在那年內,全美國最少有一百萬人蒙恩得救,亦有一百萬人得著復興；原因是有一個人開始了祈禱會,帶來莫大的影響.這次的復興並沒有出名的傳道人帶領,只有一個人忠心地開始祈禱.在美國東北有一所威廉士大學,其中有幾位青年學生,不斷同心為海外宣教工作祈禱.有一天他們祈禱時下大雨,他們就跑到附近的草堆去避雨,並且在乾草堆中禱告.這次禱告會神的靈大大工作,他們都被感動要委身去推動,並且投身於海外宣教工作.他們就開始了宣教歷史中重要的乾草堆祈禱團運動,並且促使了海外宣教運動的推行.
\newline
\newline
佈道家慕迪有一次去英國倫敦一間教會講道,聚會後有一位姊妹回去被另一位臥病中的姊妹,詢問剛才聚會的情況.那位姊妹原來過去不斷為教會祈禱,求神帶領慕迪前來那裡講道.那天她就為慕迪晚間聚會禁食祈禱,那天晚上的聚會,聖靈大大工作；比對日間的聚會,雖然是相同的講員和同一群的會眾,會眾的反應就完全不同.這是因為有一位臥病姊妹的禁食祈禱,以致聚會的果效有重大的分別.
\newline
\newline
在美國羅省唐人街一間人數眾多的華人教會,每星期三晚都有四,五百人參加祈禱會,教會非常興旺；可惜其中一位牧師和一位執事給一名神經漢開槍射殺,經過這事後,禱告會人數不但沒有減少,而且不斷增多,接近有一千人參加.
\newline
\newline
很多弟兄姊妹非常希望他們的教會得著復興.今天早上的早禱會也為香港教會的復興祈禱.香港教會雖然做很多的工作,表面看來十分蓬勃；可是若要問教會屬靈的根是否夠深?屬靈的深度是否夠厚?我們若這樣反省的時候,我們必須承認我們的軟弱!香港有一間非常復興的教會每逢星期六早上,有一個為教會復興的禁食祈禱會,在每年除夕有一個通宵祈禱會,那一間教會起初是一群年長的姊妹,用祈禱來托住教會的建立工作.
\newline
\newline
葛培理佈道團非常蒙神使用,有一次我去參觀他們的總部,發現他們辦事井井有條,同工非常愛主,同心事奉,而他自己亦蒙神重用；最後我發覺他們更重要的是注重禱告.佈道團去一個地方工作前,就先有宣傳和組織祈禱小組的同工先到,為佈道工作預備鋪路,並且成立很多三人祈禱小組,為所要邀請的人不斷代禱.幾個月前就求神施恩,拯救那些赴會者.所以在聚會中神就大大工作,感動很多人歸向神.這種佈道的果效,實在有很多各種的預備；特別是禱告的事奉,付了很大的代價.在香港及其他地方,開始了三元福音倍進運動,它其實是一種傳福音的訓練方法,在香港就叫它「三福」訓練,這種訓練有一個特點,就是以三人一組去作傳福音；其中兩個人要作其他一人的禱告支持者.不斷用愛心,忍耐,發揮恩賜,作逐家探訪佈道的工作.香港有學園傳道會,曾經舉辦過不少大型的佈道會.有一次我有機會探訪他們的總部,其中他們有一間祈禱室,是不斷舉行二十四小時的連鎖祈禱事奉；為著各地的工場及機構的需要禱告,神就使用學園傳道會,在世界各地拓展神的國度,這是他們重禱告之事奉.
\newline
\newline
兄弟過去幾年有機會在教會學習牧養的事奉,由全時間作神學教育,轉任全時間的牧養工作.在適應中覺得事務繁多,特別在我首先工作的教會,非常複雜；面對這種環境,使我更加要依靠仰望神,特別在禱告上,有很多學習事奉的基礎,是在膝蓋祈禱上.在主日學的教導工作中,當我教授希伯來書,更研究耶穌基督是有大祭司的職份,並且長遠活著,為我們代禱.今日我這個傳道人,亦有祭司的職份,為眾信徒祈求.所以我就每天為每一個信徒或長執提名禱告,求神在他們的生命中做奇妙的工作,求神帶領教會,堵塞破口,遠離罪惡；並且每星期抽一天時間禁食禱告.又在每星期兩三次與師母為教會各種事情禱告.另外每天晚上的家庭崇拜,都要彼此代禱,又為教會的工作代禱.在教會事奉之始,我就傳講禱告的信息,鼓勵弟兄姊妹勤於禱告,並且一同前來參加教會的祈禱會.在教會中,又發動組織祈禱鍊隊,將五,六個弟兄姊妹分成一條祈禱鍊,每隊各一隊長.若其中有任何需要代禱的事情,就可以打電話給祈禱鍊隊的隊長；隊長就打電話給各隊員及其他隊長,再轉給各隊員.所通知的信息須用筆寫清楚記下來,並且要保密.每隊員每天要祈禱十五分,並以一星期為期,交托給神,並且彼此分擔重擔.若有特殊事項發生,我們就在每天晚上指定的時間,一同代禱,使每一位弟兄姊妹能夠參與事奉.若有重要的事情須代禱,就用表格給各人選擇分時段作廿四小時的連鎖祈禱.教會亦設有代禱卡,只要弟兄姊妹把需要代禱的事情寫上,經收集後,同工就會在當天為他們禱告.其中更有為病人祈禱,或為來信主的親友,甚至其他的人的需要.在崇拜前亦有半小時祈禱的時間,為聚會和參與事奉者祈禱,讓神使用事奉者,帶來美好的果效.
\newline
\newline
禱告的事奉是既簡單,又能讓每一位信徒都可參與,是最有效的事奉.不論我們的健康情況如何,或沒有甚麼大的恩賜,但每一個基督徒有同樣的權利,來到神面前,奉耶穌的名祈求,使神的應許應驗.現今香港有很多參與事奉或與事奉有關的運動,其中最缺乏是背後禱告的支持；很多人關心基本法!但有沒有一個為基本法建議的禱告會呢?想到九七前後的情況,有沒有經常性為這些適應祈禱呢?我們有信徒更新運動；可是,我們是否有信徒更新運動的禱告會呢?在近期的教會歷史中,韓國有很大的復興,韓國的信徒非常注重禱告.今日香港教會亦可向韓國的信徒和教會學習；今天香港教會若要得到復興,必須要在禱告中仰望,祈求神奇妙的作為,此外沒有更好的方法.
\newline
\newline
主耶穌曾說:「我的殿必稱為萬國禱告的殿.」(可十一17)這是教會不單是有肢體生活,傳福音的活動,更應該有禱告的生活,為萬國,萬民代求.今天教會是否看重禱告的職事呢?現今的教會對話語的職事很重視；可是,對禱告的職事,時常會忽略的.教會是領人去禱告的地方,唯有教會重視禱告,教會才能發揮應有的功能.今日我們不能做甚麼!只有仰望那位我們向祂禱告的神,事情才能有所改變.有關禱告的事奉,請容許我有些建議:
\newline
\newline
一．鼓勵信徒自己組織祈禱小組──按照他們合適的時間地點,作經常性的禱告事奉.在教會各樣聖工的推行中,必須加上禱告小組,用禱告支持各種工作的推動,並且讓全教會在禱告上參與事奉.例如在探訪工作中,可留下一組人特別為各區的探訪工作代禱,用禱告托住工作和成全事工.又例如在培靈會舉行時,有一些肢體可在聚會的進行中不斷代禱,讓神更使用各種的聚會.
\newline
\newline
二．要重整教會祈禱會的內容,質素──教會應重視禱告會,多作會前的準備,對時間的編排,內容的表達都應詳加思考,讓祈禱不是傳道人或幾位老弱的弟兄姊妹所參與,我們應積極鼓勵人參與,讓信徒明白祈禱會的意義和重要性.當會眾自己去祈禱時就經歷神的同在,和靈命的更新,讓參與禱告者不是聚會的座客,而是一個真正事奉的禱告人.在禱告中為神而爭戰,為福音的廣傳而努力,為神的榮耀而爭戰,他們是禱告勇士.或者可以專設一些為病人提名禱告的聚會,讓神成就祂的醫治工作,彰顯各種神蹟奇事.可否在各個家庭組織祭壇,並且在家庭中進行祈禱.今日雖然我們很忙碌,但當我們專注祈禱的事奉時,問題會有所改變.家庭透過家庭祭壇的建立,能促進家人的溝通,彼此的關顧和彼此相愛,更經歷神更大的恩典,認識禱告的權利和福氣.
\newline
\newline
今日我們希望神更新我們,我們就要從自己的禱告生活開始.我們必須增加禱告的時間,讓禱告生活成為每一位弟兄姊妹的重點.傳道人可能很多時間預備講章,但有多少時間為那聚會祈禱呢?今日真是希望各位在禱告事奉上有更新,禱告不是一項輔助品,而是一項重要的事奉,是一切事奉的基礎.
\newline
\newline


\newpage

\section{第60屆~港九培靈會~劉承業~第七講}
\label{sec:1205}
\textbf{攻守的並重}
\newline
\newline
連結: \href{https://www.hkbibleconference.org/session-message/view/1205}{\texttt{https://www.hkbibleconference.org/session-message/view/1205}}
\newline
\newline
\hyperref[sec:1202]{< < < PREV SERMON < < <}
~
\hyperref[sec:toc]{[返目錄]}
~
\hyperref[sec:1209]{> > > NEXT SERMON > > >}
\newline
\newline
約書亞記 11:16-11:23
\newline
\begin{longtable}{cl}
\hline
\hline
章節 & 經文 (和合本修訂版)\\
\hline
11:16 & \begin{tabularx}{0.7\textwidth}{X} 約書亞奪了那全地，就是山區、整個尼革夫、歌珊全地、低地、亞拉巴、以色列的山區和山下的低地， \end{tabularx} \\ \\ \relax
11:17 & \begin{tabularx}{0.7\textwidth}{X} 從上西珥的哈拉山，直到黑門山下面黎巴嫩平原的巴力‧迦得。他擒獲了那裡的眾王，把他們殺死。 \end{tabularx} \\ \\ \relax
11:18 & \begin{tabularx}{0.7\textwidth}{X} 約書亞和這些王作戰了很長的一段日子。 \end{tabularx} \\ \\ \relax
11:19 & \begin{tabularx}{0.7\textwidth}{X} 除了希未人基遍的居民之外，沒有一城與以色列人講和，都是以色列人作戰奪來的。 \end{tabularx} \\ \\ \relax
11:20 & \begin{tabularx}{0.7\textwidth}{X} 因為耶和華的意思是要使他們的心剛硬，來與以色列人作戰，好使他們全被殺滅，不蒙憐憫，反被除滅，正如耶和華所吩咐摩西的。 \end{tabularx} \\ \\ \relax
11:21 & \begin{tabularx}{0.7\textwidth}{X} 那時約書亞來到，剪除了住山區、希伯崙、底璧、亞拿伯、整個猶大山區和以色列山區的亞衲族人。約書亞把他們和他們的城鎮盡都毀滅。 \end{tabularx} \\ \\ \relax
11:22 & \begin{tabularx}{0.7\textwidth}{X} 以色列人的地中沒有留下一個亞衲族人，只有一些還留在迦薩、迦特和亞實突。 \end{tabularx} \\ \\ \relax
11:23 & \begin{tabularx}{0.7\textwidth}{X} 這樣，約書亞照著耶和華所吩咐摩西的一切話奪了那全地，就按著以色列支派所得的份把地分給他們為業。於是國中太平，沒有戰爭了。 \end{tabularx} \\ \\
[1ex]
\hline
\hline
\end{longtable}
經文:(書十一16-23)
\newline
\newline
約書亞帶領以色列人攻打迦南,(書十二24)記載他們一共殺了卅一個王；在第六章攻陷耶利哥,第十,十一章又攻陷南地,其後又北伐；第十二章就交代他們的爭戰暫時完畢.神雖然應許他們得著迦南美地,但他們必須要爭戰,攻取每一個地方.可是神在他們攻城前,已用各種方法,造就,操練他們.首先是藉著過約但河事件,堅定他們的信心；稍後又藉著行割禮,提醒他們與神立約的關係.他們是屬神的子民,神必定負責他們一切的需要；神又藉著摩西的律例,誡命和典章,給他們標準去生活；只要他們能遵守立約的條文,他們便得到更大的賜福和應許.
\newline
\newline
從以色列歷史中反映出──神不單看重造就工作,亦同樣重視爭戰的工作；這種爭戰的工作就如現今的傳福音工作一樣.「守」是造就和栽培,而「攻」則是爭戰和傳福音.從約書亞記,我們希望歸納出攻守並重的原則,也是教會得勝的策略!注重栽培造就和傳福音的工作.在主耶穌所賜的大使命(太廿八18-22),亦包括這個重要的策略:造就和宣教.「所以你們要去,使萬民作我的門徒.」使萬民作主門徒,是包括這兩方面的行動.在大使命的經文中,其中主要的動詞不是「去」,也不是「教訓」,而是「作主門徒」.耶穌最重要的吩咐是「使萬民作主門徒」,這是我們要向所有人傳福音的使命；傳福音和造就是連結在一起,也是教會的兩大支柱.
\newline
\newline
一間興旺的教會,是針對這兩方面平衡發展.栽培造就和訓練,是注重質的增長；傳福音和宣教,是注重量的增長.質量兩方面不斷的增長,才能滿足主的心意.有些教會比較內向,就側重造就,少注重傳福音,當然發展較穩定,可是量的增長卻緩慢,教會對外面的失喪靈魂關心不足.有一間教會,他的牧者不喜歡所牧養的教會太大,認為超過一百人的教會就會發生很多問題,他就是小教會路線；兄弟並不贊成這種看法,若因怕發生問題,而不去增長,也許是有點自私.例如:百多人的教會需多請一個同工時,若彼此配搭不好,就容易發生問題:一個人一手包辦,我行我素,似乎更易辦事.人多時會影響親切感,意見又會增加,更難於調和,這也許是信心不足.相反地,當神賜給教會得救人數天天加增,我們應當相信,神也會保守弟兄姊妹有一個親切的感情；並且同工又能合作,熱心事奉神,在神的恩典中克服問題和困難.另外可能有一些教會,只著重傳福音,卻不重視栽培造就的事工,教會似乎在量的增長有顯著的表現,在質的建立上,就不大理想.結果造成很多屬靈的嬰孩出現,教會根基不穩；並且容易走上屬世的路線,導致名存實亡.教會只是脹大,所有的是不健康的增長,又會百病叢生.我認識一間教會,因信徒靈命成長鬆散,加入教會時沒有經過靈程的評定,就被接納成為會友；初入教會不久就擔任重要的事奉角色,雖然教會發展得很快,實質卻沒有屬靈的根基.結果很多屬血氣的事情就發生,造成分裂的現象.我認識一位曾在機構事奉的同工,後來轉往一間教會事奉,他採用了機構的原則作福音工作,在他帶領下教會開始增長；當細察他工作方式時,我感到很難受.因為當他接觸一些未信者時,那人願意一同禱告,他就宣佈那人已經重生得救,並且很快就給他們施洗.這樣工作表面上看來很蓬勃,但實際上是疏忽了栽培工作；信徒屬靈的根基淺薄,一段時間後工作就漸走下坡.兩個禮拜前我曾經與這位同工聯絡交通,知道他可能會離職,這實在是一件可惜的事.
\newline
\newline
我們在神的教會事奉,不能太內向或太外向,必須有適當的均衡,並且攻守並重,注意質量平衡的增長.讓我們在這裡再深入思想這兩方面的工作:
\newline
\newline
(一)「守」的策略
\newline
\newline
栽培工作的焦點,是大使命重要的中心工作,「使萬民作門徒」.教會的事工中心,是在信徒的生命中做建立的工作,讓他們成為主的門徒；就如主耶穌在(太十六24)說:「若有人要跟從我,就當捨己,背起他的十字架來跟從我.」他們應該是捨己,願意付代價來跟從神,事奉神,這樣教會方有真正的復興動力.在整個傳福音的過程中,要讓慕道者真實知道怎樣去信主,怎樣去跟主,怎樣才是主的門徒.又向他們說明所要付的代價,才接受他們參加慕道班,洗禮班,以致他們能全面的預備,加入投身於教會,成為會友,基督的肢體.我們不能因遷就人的需要,而降低屬靈的要求和標準.現在有很多教會或傳道人標榜一種成功神學,就是信主後甚麼都好!甚麼都順利!並且神是有求必應的.這種財富和健康的福音,只是部份的真理；這是種似乎富吸引力或近似利誘的片面宣講,會做成更大的功利主義.這未必是聖靈的真正工作,我們必須真正相信神的大能,會改變一切的人.一個人信主,是聖靈親自的工作使人認識自己的內裡需要,自我的罪,及經歷神的赦免和慈愛,並且將自己的生命交托給神.在我事奉的教會中有一位年長的姊妹,她的女兒在溫哥華生病過世；她的女婿是一位回教徒,不經過這次事件,她的女婿是絕對不會信主的.當時他傷心難過,偶然在妻子的遺物中發現一本岳母的聖經,他好奇打開聖經,就翻到詩篇廿三篇,又留心細閱,特別對最尾的兩節,感受很深.為了更了解這些經文,他便去找另一位信主的親戚；那位基督徒便帶他往教會,給予他幫助.過了個多月,他就決志信主；後來更安慰傷心的岳母,那位年長的姊妹就介紹他給我認識.在這事件中,神已預備好他的心,並且再次禱告接受神,現在他在溫哥華一間教會中熱心事奉神.
\newline
\newline
主耶穌在大使命中說:「凡我所吩咐你們的,都教訓他們遵守.」這是一種教導的工作.現今的教會,一般都缺乏教導的工作,就我所接觸的教會,不論是在香港或海外,教導的事工都不夠強.求主幫助我們加強講台的信息,充份豐富查經的聚會；並且讓我們事奉神的人能多多在神的話語上下工夫,去餵養群羊.若我們只是單顧傳福音,而忽略栽培,造就和訓練工作；就好像教會的前門大開,有很多人進來,但教會的後門卻開得更大,那些從後門走了的人,以後要重返教會就更困難.而且他們會有很多誤解,藉口和評價,慢慢就失落了.我們必須加強栽培,切實相交相愛,作出美好的見證；讓初信者透過與成屬靈命的基督徒交往,就更被激發愛主.
\newline
\newline
在「守」的策略中,我們更應訓練信徒,使他們投入教會的事奉,給予機會,發揮他們的恩賜.若信徒在教會中閒散,就如一些人單是進食,沒有工作,就漸漸肥胖起來,慢慢失去工作的能力；透過訓練,讓信徒不斷操練,事奉,他們就更長進.所以無論是大小崗位的事奉,都應給予信徒機會,讓他們被動員起來.可惜有些教會,因人手不足,又不加以任何的訓練,讓初信或經驗淺薄的肢體不斷工作；即使初信時他們很熱心,可是不懂如何做；受到挫折失敗後,便會灰心倒退,甚而影響到靈命成長.有一位弟兄曾向我投訴,他的教會沒有給他訓練的機會,就要他擔任事奉；他形容自己是「爛頭卒」,心中充滿掙扎.故此,教會必須教導和訓練,使弟兄姊妹知道如何事奉,這樣才能事半功倍,去享受事奉的得勝和喜樂哩!
\newline
\newline
(二)「攻」的策略
\newline
\newline
傳福音不但有本地的佈道,更有海外的差傳工作.兄弟很虧欠因時間上的關係,經驗上的不足；雖然知道差傳工作很重要,卻只能集中在傳福音事工上與各位思想.近年香港教會興起了一股佈道的熱潮,所以教會能不斷增長,很多神學院招生滿額,無法取錄更多的投考者；這是香港教會的福氣,我們為此感謝神.但是要造就,訓練信徒,差派他們往各地傳福音,這才是人人事奉的良方.
\newline
\newline
上一講論及禱告的事奉時,可以看見禱告也是一個美好事奉的方法.不論參加者的年齡的高低,恩賜多寡,每一位都可藉禱告將事情交託給神,並且托住教會的工作.傳福音的工作正如禱告事奉一樣,可讓弟兄姊妹人人參與.最近中國信徒佈道會出版了一本季刊,名叫「傳」.其中一期的小標題是:「口傳,身傳,差傳.」我們每一個人都可以藉著日常生活,去見證神.有一個口去講神的榮美和作為,有工作的能力去謀生；在經濟上支持神的工作.我們需要傳福音,因為這是主的吩咐；主耶穌臨離開世界時,清楚囑咐門徒:「你們要去,使萬民作我的門徒.」神要使所有的人都成為門徒,今日我們要去,不論遇見怎樣的人,都有一個責任,是屬靈上,更是福音上的負擔.就如明天陳牧師講論的保羅,他是一位福音的戰士.按保羅自述:「我若不傳福音,我就有禍了!為我骨肉之親,就算與基督分離,我也願意!」他有愛人靈魂的心,他常常感嘆,欠了人福音的債.很多時我都用這種欠借感來勉勵自己；實在有很多的福音債,我們並未償還!我們真是希望見主面的時候,能盡量還清這些福音的債,那時就不會羞愧了.
\newline
\newline
傳福音是每一個基督徒的天職,因為我們要作主的見證人；信徒要在耶路撒冷,猶太全地,撒瑪利亞；直到地極,作主的見證!我們可以在家庭,工作的場所,為主作見證.宣道會的創辦人宣信博士曾經說過:「每一個基督徒在某處,都是一個宣教士.」我們傳福音,可以促進主的再來.(太廿四14):「這天國的福音要傳遍天下,對萬民作見證,然後末期才來到.」我們很盼望主早日回來,結束這世界的一切苦難紛爭；但福音仍未傳遍天下,神仍然憐憫那些未得救的失喪靈魂.傳福音與我們個人靈命的增長和教會的興旺有密切的關係；當我們成為一個福音的使者時,我們就更經歷神的同在,豐富我們對神的認識.很多教會看重傳福音的工作,就充滿生氣,十分興旺.有一個機構,它的創辦人是肯尼第,是「三福」的發起人.起初他的教會很少人,工作開始時充滿困難,後來他就發起傳福音的工作；編寫「三元福音倍進法」,並且廣泛推廣到其他教會.他不單令自己的教會不斷的增長,更讓其他教會得到傳福音的復興!他們工作的特點是一個人傳福音,背後有兩位代禱者.有一個退休的牧師,被一間只有九個人的教會邀請,到那裡工作.他就訓練人傳福音,一個一個不斷地訓練,結果幾年後,教會非常興旺；償還所付的欠債,並且能供給他們兩夫婦需用,到各處領會.他雖然出外,但教會的弟兄姊妹又很自發地事奉,並且不斷地增長.這是我們所要羨慕的成功.
\newline
\newline
這位牧者聽從主的吩咐,「使萬民作主門徒」,不但帶領人認識主,更建立他們,裝備他們；讓他們能在教會中事奉,並且生生不息,繼續的傳福音,才是一個成功的教會.
\newline
\newline
在我的教會中也設立了聖工小組,訓練教導他們分工,努力事奉；弟兄姊妹又響應參與,並且知道如何事奉,教會就充滿動力和不斷增長.任何一個教會願意傳福音,就會不斷成長增長.當我們對傳福音有負擔,又肯去做,組織更多小組,去傳福音,代禱,交通,立約,就成為使命的福音小組,增長就持續.有時我們覺得自己很軟弱,也不懂得說話技巧,好像我也是個拙口笨舌的人；雖然如此,神亦能使用.王下五章記載以色列人被亞蘭人所侵略,當時亞蘭的元帥乃縵,患了大麻瘋；但從以色列中被擄來的小女子,卻對她主母說:「巴不得我主人去見撒瑪利亞地先知,必能治好他的大麻瘋.」她只講了一句說話,就令乃縵元帥急往以色列尋求醫治；後來得到以利沙的幫助而認識神.乃縵元帥得醫治後見以利沙,見證說:「如今我知道了,除了以色列之外,普天下沒有神......(王下五16)」因她一句話,就讓神在外邦國家大得榮耀!
\newline
\newline
我認識一位曾患上精神病的弟兄,在患病期間遍尋各地名醫,病卻沒有起色；反而日益嚴重,接近精神分裂的地步.但偶然認識一位基督徒的中醫師,那位醫生就向他講見證；並說明唯有神才能醫治他,且介紹他到教會找我.我了解他的情況後,便引導他認識耶穌,明白福音；他信主後很奇妙地不藥而癒,這完全是神的工作.後來他家中各人相繼信主,並且他時常作見證,說明他過去患病的事,是有使命的；因為神用疾病作間接拯救的工作.他不單自己得益,他的家人,朋友,甚至聽他見證的人,都得到很大的幫助,願意接受主耶穌.那位基督徒醫師只講了一句語:「你要信耶穌,祂可以幫助你.」神就負責施行拯救.讓我們時常作見證,完成神的工作.
\newline
\newline
我們更要攻守並重,時常事奉,保持靈性的更新,樂傳福音,注意質量均衡的增長.
\newline
\newline


\newpage

\section{第60屆~港九培靈會~劉承業~第八講}
\label{sec:1209}
\textbf{領袖的質素}
\newline
\newline
連結: \href{https://www.hkbibleconference.org/session-message/view/1209}{\texttt{https://www.hkbibleconference.org/session-message/view/1209}}
\newline
\newline
\hyperref[sec:1205]{< < < PREV SERMON < < <}
~
\hyperref[sec:toc]{[返目錄]}
~
\hyperref[sec:1211]{> > > NEXT SERMON > > >}
\newline
\newline
神雖然應許以色列人得迦南美地為業,然而,他們要由領袖帶領,憑信爭戰,去承受所得的地方.神就興起迦勒,約書亞等領袖去帶領他們；不斷得勝,承受神賜給的產業.今日我們靠主作聖工,謙卑地屈膝祈求；凡事交托神,讓神作成祂自己的工.可是,神卻要使用人,不單是全職的傳道人；還有教會的長執領袖,和每一位信徒,共同分擔神的工作.今天與大家思想領袖的質素,是廣義上的領袖,是所有願意參與一同事奉的一群神的用人.無論是在特殊崗位事奉,或是作隱藏的工作,或是預備事奉的信徒；都可作參考,使我們更能揀選一些優秀的領袖.我們將從約書亞的身上,思想領袖的質素.
\newline
\newline
一間教會或機構,若能有優秀的領袖作領導,是神特別賜福他們的明證；假若一個組織沒有優秀的領袖,就算在下的人如何努力,拼命工作,在工作收效上,必然大打折扣.很多時因領袖的問題,以致其他的弟兄姊妹都灰心喪志,受到很大的虧損,工作上只有事倍而功半的果效.神的工作若有好的領袖帶領,工作就必定興旺,有很大發展的潛力.
\newline
\newline
當談及領袖時,首要是他們已經重生得救.請看(書一8):「這律法書不可離開你的口,總要晝夜思想,好使你謹守遵行這書上所寫的一切話.」這是神在約書亞開始繼承摩西,作以色列人的領袖時,向他囑咐的一段重要的話.屬靈領袖的質素是:
\newline
\newline
(一)要謹守遵行神的話
\newline
\newline
我們從約書亞身上,實在看見他謹守遵行神的話.每一次神吩咐他辦理的事情,都辦得非常妥當.例如在不同的戰役,神所指示的進攻方法,他都依計行事.其他事情如過約旦河,憑信分地,揀立逃城,他都尋找神的心意辦理.神非常看重祂的工人,是否按祂的心意而行?神就在(書一8)明確指出謹守遵行神的話語的重要；並且在吩咐後加上一個寶貴的應許,來鼓勵,幫助他.(書一8)下的應許是:「如此,你的道路就可以亨通,凡事順利.」在謹守遵行神話語之前,我們必須認識和信靠神的話.
\newline
\newline
約書亞記提及另一個出色的領袖是迦勒,他也是深信神的話和應許,所以他在四十五年的生活中,不斷專心跟從神,並且在八十五高齡仍滿有信心和力量,要攻取神曾經應許給他的希伯崙城.今天我們仍要信靠神的應許去生活；這樣我們才能樹立良好的榜樣,在工作上為神所使用.假若我們擁有神的話語而沒有遵行,就容易流於只用人為的方法去作屬靈的事工.教會就很快被世俗的方法和風氣所影響；並且陷入危機中,承受很大的虧損.教會沒有神的話語為標準,就很難確立正確的方向,事奉的路線亦會出現問題.教會領袖若沒有按神話語的標準作領導,就如瞎子領瞎子,結果很容易發生危機.
\newline
\newline
(書七15)說明約書亞是真實地遵照神的話語,去懲治犯罪的亞干:「亞干有當滅的物在他那裡,他和他所有的必被火焚燒；因他違背了耶和華的約,又因他在以色列中行了愚妄的事.」神吩咐約書亞將取了當滅的物的人和所有的都用火焚燒.也許他在執行這件刑罰上是非常之困難,因為在人情上有一定的難處；當我們謹守神的話,持守公義和聖潔的原則時,就容易得罪人,或被人指責「沒有人情味」!不夠愛心和寬容.在神的工作和教會事奉中,領袖們必須謹守遵行神的話,不論會招來甚麼的後果,也要持守到底.當我們要選立一些領袖時,除了看他們是否熱心；更要評定他們對聖經和真理的了解,使有生命和知識的人作領袖,方能使聖工有所發展,否則就弄巧反拙.
\newline
\newline
(二)接受神的主權
\newline
\newline
(書五14):他回答說:「不是的,我來是要作耶和華軍隊的元帥.」約書亞就俯伏在地下拜,說:『我主有甚麼話吩咐僕人.』我們要尊主為大,接受神的主權在我們身上,並且願意謙卑順服,這是約書亞第二個優點.假若領袖是自驕自傲,即使他有很大的恩賜,神也不會使用他!約書亞在未出戰時,眾人都公認他是以色列軍隊的元帥.但神在這裡向他顯現,指出以色列真正的元帥是神自己；約書亞真正接納神的權柄,在神的引導下爭戰才得勝.聖經上所記述的領袖；他們能作領袖的原因,完全是神的恩典和神的揀選.保羅就是一個明顯的例子,一個人事奉神的時候,神就將恩賜,恩典加在其身上,以致他能作成神的工作.保羅在(林後四5)說明自己的心聲:「我們原不是傳自己,乃是傳基督耶穌為主,並且自己因耶穌作你們的僕人.」兄弟曾經在基督教學校中工作,有一個時期要教授聖經；當中在教授馬可福音第十章時候,有一個特別的經歷,希望與各位分享.那一天我在解釋(可十42-44):「耶穌叫他們來,對他們說:『你們知道,外邦人有尊為君主的,治理他們,有大臣操權管束他們；只是在你們中間,不是這樣.你們中間,誰願為大,就必作你們的用人；在你們中間,誰願為首,就必作眾人的僕人.』」基督教中的耶穌,正在勸勉門徒,不要爭論誰為大；他們當中誰願為大的,誰人願為首領,就要作眾人的僕人.這裡耶穌說明,屬靈的領袖是眾人的用人和僕人.在我投入於講解時,同學中的班長突然回答說:「阿Sir,我知道了!其實耶穌在那裡所說的是:「那裡沒有『阿頭』,只有祂是『阿頭』.」我心裡想想他這番話的時候,也覺得他說得很對.我自己為首時,就好像工人,甚麼時候都沒有人為首,只有耶穌基督為首.就算我們在教會中,仍然有人做主席,會長,團長等職份.但那些人如聖經所說,應存著謙卑的態度；甚至成為『阿尾』,有這種不同的心態才是.結果就沒有真正的首領,而是神自己做首領,神亦興起使用有這種事奉心態的人,成就祂的工作.我們要接納神的主權在我們身上,我們就將選擇權向神開放；聽憑祂的指引,察驗祂的旨意.因此保羅用活祭來形容每一個願意跟從神的人,我們是祭品,沒有自己的主權,是任由獻祭者大祭司的安排.活祭是以獻身的生活為基礎,以為主而活的心態,來見證榮耀神.
\newline
\newline
在神的教會中,我們需要接受神的主權,聽從神的指揮,吩咐,這樣的人才能真正謙卑,彼此順服.今日教會中很多領袖都在爭先為首,以自己的意見為權威,令其他的人很難配搭工作；失去團隊事奉的動力,只是在賣弄獨腳戲,但很快就獨力難支.
\newline
\newline
(三)是個禱告戰士
\newline
\newline
(書十12):「當耶和華將亞摩利人交付以色列人的日子,約書亞就禱告耶和華.......」我們必須要在密室中禱告,如約書亞的榜樣,時常禱告.他禱告生活的特點是在得勝的時候禱告,不是在失敗中才回到神那裡.我們時常以為在灰心,失意,受挫和失敗,最後才去信靠神,向神禱告.但在平安穩妥的日子,我們就忘記神,很少向神祈求!我們在得勝時更要禱告,以感謝為祭獻給神,並且將榮耀頌讚歸給神,我們方能得到更大,更全備的勝利!
\newline
\newline
我認識一位能幹的弟兄,他是一位知名的心理學家,可是他不相信禱告；後來他自己反省,為何很多華人教會都不接納他,並向我訴苦.我覺得一個不相信禱告的基督徒,往往是靠己力去事奉,在神的家中工作,會帶來很多的危機.若然我們要事奉神,必須要在膝頭上下工夫,方能與神同工.
\newline
\newline
(四)家庭的見證和支持
\newline
\newline
(書廿四15):「若是你們以事奉耶和華為不好,今日就可以選擇所要事奉的:是你們列祖在大河那邊所事奉的神呢?是你們所住這地的亞摩利人的神呢；至於我和我家,我們必定事奉耶和華.」在經句的下半節:「至於我和我家,我們必定事奉耶和華.」是大家非常熟識,甚而在結婚時亦會採用,作為主要的金句.這句聖經不單代表約書亞的決心,更是他的見證.我們每一個事奉神的人,都應從家裡開始作見證.英文的諺語說:「我們的愛心,應從家裡開始!」
\newline
\newline
提摩太前書第三章論及執事,監督和長老,都說明屬靈領袖的要求,必須好好地管理他們的家和兒女；他們應該在家中活出美好的見證.幫助影響他們的家人,並且進而影響幫助其他的人.也許有些領袖在外面被人稱讚,卻在家庭中生活被人指責,他們必然不能討神喜悅.
\newline
\newline
我們事奉神的人,在家庭中應有美好的見證,發揮影響力；特別在靈性生活上,要成為家人的好榜樣,並且得到家人的支持；同心合意地事奉,便能發揮大的果效.這樣沒有後顧之憂,事奉就能更加得力.兄弟所事奉的教會,若有人被推薦為執事,或擔任特別的事奉,假若他們是已婚的,必是夫婦兩人同被邀請,約見交通；就算他們只是單方面參與事奉,我們亦會取得他配偶的意願.知道他們能作出支持,才給予事奉的職事；這是因為我們深信,事奉必須有家庭的支持,尤其是配偶的支持；這是非常重要的；若全家能一同事奉,那樣的見證是最美好的,也能發揮更大的力量.
\newline
\newline
在今天的教會中,很多弟兄姊妹都希望看見一些事奉的榜樣,特別在基督化家庭榜樣上的美好見證.我們在揀選事奉神的人時,必須認真了解他們在家庭中的生活,影響力和見證,因為這是重要的領袖質素.有些可能在外面非常熱心事奉,但在家裡可能疏於照顧家人,令他們有很大的虧損.兄弟十分感謝神,雖然自己有很多短處,但我的家庭卻很支持我的事奉.就算我們在外面遇到任何的困難挫折；我們都可以帶回自己的家中,為各種事情禱告.一起用禱告仰望神,並且能重新得力,透過全家的事奉,更能體會神更大的恩典.
\newline
\newline
(五)領袖的產生
\newline
\newline
在約書亞記中,除了以上四種條件質素外,我們更知道神會興起領袖.因為屬靈的領袖,不是生出來的,而是訓練出來的.我們環顧領袖的缺乏,有時深嘆無處可尋；可是我們若注重訓練,將來必定有領袖出來,承接工作的責任.也許中國人較易忌才,未能好好栽培發展後輩；更不能實行和盼望青出於藍而勝於藍的原則；我們必須裁培鼓勵後輩,又看重領袖訓練的工作.
\newline
\newline
按兄弟的了解,領袖訓練可分為三個階段:
\newline
\newline
(1)真理的領受
\newline
\newline
若沒有真理作基礎,但有熱心和恩賜,或屬世的才華,就很難成為優秀的屬靈領袖.(提後二2)保羅勉勵提摩太:「你在許多見證人面前聽見我所教訓的,也要交托那忠心的能教導別人的人.」領袖必須有穩固的真理基礎,願意追求認識神的話,方能參與各方面的事奉.在兄弟牧養教會的經歷中,很注意觀察參加查經班或主日學的學員,對研讀神話語的反應,是否真正有心追求；然後經過一段時間,能被那些長進的人接納,讓他們接受一些事奉.我認識一位同工,他每到一間教會事奉,都用禱告的心,希望揀選栽培幾位的弟兄姊妹作領袖,彼此配搭事奉.
\newline
\newline
(2)有事奉的操練
\newline
\newline
事奉神不單在知識上,更應在實際參與和操練上,給予機會,不斷操練；很多教會的領袖只是會議的參與者,可能單用口講,未能有實際的事奉,只是紙上談兵.今天我們要訓練屬靈領袖,必須給予其機會操練和學習,鼓勵他們成長,成為信徒的榜樣.
\newline
\newline
在過去兄弟發展的教會,時常設立一些聖工小組去推動協助工作；並在其中作教導和訓練,使他們漸漸懂得如何去事奉,以後就可成為教會的領袖.
\newline
\newline
有些教會有一些叫提摩太計劃,就是要年青的人參與事奉.造就他們成為將來的接班人,帶領教會的工作,不斷更新成長.
\newline
\newline
(3)門徒訓練
\newline
\newline
現今很多教會都注意幾方面工作,造就建立成熟的生命,成為主的門徒.這樣方能成為教會優秀的領袖,在事奉神的能力上就不斷加強.耶穌基督就是這樣花了三年多的時間訓練十二個門徒,成就向普世傳福的計劃.
\newline
\newline
(4)注意兒童和青少年工作
\newline
\newline
如果我們要訓練領袖,就必須從根本做起,並要從基督教教育入手,這是最重要的一環.這些工作是現今教會很多時候會忽略的.「十年樹木,百年樹人」,很多有份量的領袖必須是花時間慢慢栽培而成,欲速則不達.如果領袖沒有根基,當他站在領導的位置上,就會把弱點暴露出來,成為很大的問題；當我們平衡各種事工之前,應該先建立根基.
\newline
\newline
求主幫助我們的教會,不單看眼前,更應為未來打算,作好訓練的工作,迎接未來更大的發展和挑戰.
\newline
\newline


\newpage

\section{第60屆~港九培靈會~劉承業~第九講}
\label{sec:1211}
\textbf{同心與配搭}
\newline
\newline
連結: \href{https://www.hkbibleconference.org/session-message/view/1211}{\texttt{https://www.hkbibleconference.org/session-message/view/1211}}
\newline
\newline
\hyperref[sec:1209]{< < < PREV SERMON < < <}
~
\hyperref[sec:toc]{[返目錄]}
~
\hyperref[sec:1216]{> > > NEXT SERMON > > >}
\newline
\newline
先父在生的時候,他的行為也不大好,家中的子女對他很畏懼.當他中午回來午睡時,子女們紛紛走避；直至第二次世界大戰爆發時,他就完全的改變.成為一個慈愛,顧家的父親；他親自買菜燒飯給我們吃,他改變後的影子,時常停留在我的腦海中.他在晚上又講故事給我們聽,至今他所講的故事,我仍然記憶猶新.為何他會判若兩人,有那樣重大的改變?正因為面對已爆發戰事,他覺得自己已時日無多了!
\newline
\newline
今日我們住在香港這個非常安定繁榮,表面看來是平安穩妥的地方；但若從屬靈的角度去看時,實在這是一個爭戰的時代.(彼前五8-9):「務要謹守,儆醒.因為你們的仇敵魔鬼,如同吼叫的獅子；遍地遊行,尋找可吞吃的人.你們要用堅固的信心抵擋他.......」今日一些蒙福的教會,都是同心合意,興旺福音的,就如腓立比教會就是一個很好的寫照；初期教會時常同心合意,在殿裡或家中擘餅,存著歡喜誠實的心用飯；然後,主就預備將得救的人數加給他們.今日我們都盼望我們個人的靈性,我們的團契和教會得到復興和更新,其中最重要的關鍵是今天的信息!同心與配搭.
\newline
\newline
約書亞時代,人民是在靈性較高的層次上；因為他們能彼此同心配搭,證明同心就是力量的真理.撒旦最懼怕的是教會的肢體能夠同心,因為牠知道任何的逼迫,都不能令教會妥協低頭；唯有分裂,讓他們失去同心的追求,牠便能有機可乘了.我們從(書一16)看看以色列人如何同心:「他們回答約書亞說:『你所吩咐我們行的,我們都必行；你所差遣我們去的,我們都必去.』」同心的表現是在支持,接納,並在順服領袖們的引導上流露出來!約書亞之所以能被神這麼大的使用,因為以色列人全體都支持他,(書一18)更說明,若以色列人不遵從約書亞命令,結果如何呢?「無論甚麼人違背你的命令,不聽從你所吩他的一切話,就必治死他.你只要剛強壯膽!」當教會的信徒能夠支持他們的傳道人和領袖們時,他們能夠剛強壯膽.有好的牧人而無信徒支持,他就像無兵司令,不能發展作用.教會單有好的執事也不行,因為她不能靠幾個人去建立,必須全教會一同去面對這場屬靈的爭戰.
\newline
\newline
我們身為支持者,不能因領袖有恩賜,才能,或領袖與我們有特殊密切的關係,而大力支持他,與他合作.因為人是有缺點,難免犯錯,就如約書亞在第九章與基遍人立約的事上,就造成很大的錯誤.當時以色列人並無埋怨約書亞,他們都很支持接立約書亞這位領袖；在分地的事情上,也看見以色列人很順服他的決定,在利益財產上的處理分配,是很容易出現磨擦不和,但在本書十八章中,約書亞吩咐分地的原則,並採用神決斷的方法──拈鬮,去處理分配,所以他們都接納這等方法,沒有表示不滿.他們接立約書亞,因為他們知道神揀選使用他,他們又順服神管治的權柄.
\newline
\newline
迦勒的事奉也是一樣,他更謙卑地到約書亞面前求取希伯崙的地業；雖然他的背景和閱歷可與約書亞看齊,平起平坐；但他卻非常尊重約書亞的決策,定奪,因為他接納神在約書亞身上的權柄,而沒有因不滿而策動猶大支派叛變.當神興起領袖去帶領我們時,我們都應支持,接納,以致順服他,使神的工作因我們的同心配搭而興旺起來.
\newline
\newline
對領袖的支持,可從他們的欣賞,愛戴中表現出來；這樣的表達,可帶來很大的鼓勵.領袖的工作往往是吃力不討好的,很容易就被人誤解,批評；教會中很多的信徒就會對領袖的工作產生逃避的心態,不肯承接領導的職事.在教會中事奉的領袖,同工,也許要任勞任怨,不怕人的批評和誹議；並要厚起面皮,面對各種的問題.也不怕失去面子,要為真理而站穩持守到底.
\newline
\newline
我認識一位一年一任的執事,他在任之時似乎受了很多委屈；所以他有時就嘆氣,要在卸任後對批評他的人還以顏色.他這種態度雖然不正確,須要被糾正；但他在教會中為領袖所受的苦楚,也是很多事奉者的心聲.這種現象很普遍,令很多願意承擔的領袖裹足不前.我們應該設法去支持欣賞我們的領袖,使他們有信心勇敢為神作事.我們更應為領袖們代禱,使他們更有力量事奉.
\newline
\newline
我們在(書廿二)可看到以色列人如何同心.那裡記載流便,迦得,嗎哪西半支派的人打完仗,預備過約但河返回他們所得的地業生活.那時他們發覺若過河後,恐怕其他支派或將來他們的子孫,可能忘記他們曾參與爭戰的歷史；便在河邊築一座壇,以作記念,以便給將來的子孫謹記敬畏耶和華.可是其他九個半支派卻誤解他們的做法,以為他們所造的壇是拜其他偶像,放棄敬拜事奉神.他們就同心合意的為神大發熱心,要興師問罪,阻止他們的做法.我們若要同心,就要先與神同心,追求聖潔,同心抗拒異端,並不妥協；這是為真理,為神的心意,為神的性情而力爭,因為這是同心的基礎.我們這樣同心,是穩固的,不會受環境或別人的疏擺而影響；求神幫助我們有這樣穩固的基礎,去建立同心的行動.
\newline
\newline
同心的另一個表現在(書廿二14-15).兩個半支派的戰士打仗完畢,要返回家園；他們過去曾經冒死地爭戰,離開他們的家園,又與其他九個半支派同心作戰.他們不是單顧自己,而是為整體以色列人的好處,作出了犧牲.我們要真正同心,是要放棄自私心的；並且看別人比自己強,這樣我們才能擴闊我們的心胸.以愛心出發,顧及別人的需要,成為神更大的祝福.在教會中我們應強調彼此相愛,顧及到有需要的人；可能教會慈惠部,是教會中較弱的一環.因未能發展對他人有所幫助,因為我們少付出和施與.實在我們周圍很多人,在屬靈和肉身上都有很大的需要,我們亦要關顧別人的整全需要.
\newline
\newline
我們同心,並不是說甚麼都是一樣.在以弗所書第四章論及的同心,是說明在不同的崗位；職份或恩賜中,我們應要同心配搭,有同一個目標和見證,去建立基督的身體.我們雖然在學問,年齡,背景,經歷,文化等事上有很多不同；但我們的目標仍是相同的,是要建立基督的身體；彼此接納,諒解.在不同職份和功能中,我們必須注重彼此的配搭和協調.我們從(弗四11-12)就看到第一個配搭,是傳道人與弟兄姊妹的配搭；今日教會發生很多問題就在這關係上,並且彼此不熟知對的職責,生發很多磨擦.
\newline
\newline
今天論及配搭,是從(弗411-12):「他所賜的,有使徒,有先知,有傳福音的,有牧師和教師,為要成全信徒,各盡其職,建立基督的身體.」這裡論及神賜給教會各種恩賜或職務.為了「成全聖徒,各盡其職」.這裡的「成全」兩個字,有裝備的意思,是裝備,造就信徒；預備這些信徒作神使用的器皿,去作成神的聖工.教會的傳道和牧師,就為要裝備成全聖徒,作主的工作,建立基督的身體.這些被裝備成全好的信徒,就去建立基督的身體──教會.教會中的傳道人與信徒,必須建立起團隊的精神,彼此配搭事奉；在原文翻譯上,成全信徒後之逗號可略去.使成全聖徒各盡其職連結一起,更加強全體配搭的重要；這樣教會的工作方能興旺,並且發揮更大的動力.美國史丹福大學旁邊有一間半島聖經教會,這間教會就是從五個家庭開始發展；牧養該會的牧師,在那裡工作了三十多年.該會出版一本書名為──「以愛相繫」.內容特別論及教會每一個弟兄姊妹,如何配搭事奉,以愛相繫.其中有弟兄姊妹奉獻,起來事奉,並且被教會其他家庭支持,得著迅速的發展.
\newline
\newline
我們如何去配搭呢?從(弗四13-16)說出三方面的合一.第一是在真理上的合一.第13節:「直等到我們眾人在真道上同歸於一,認識神的兒子,得以長大成人,滿有基督長成的身量.」我們一定要在真理上同受一樣造就和教訓；特別是在研究解釋聖經上,不能稍有偏差,或是標奇立異.否則就會成為異端,如摩門教和復臨安息日會相似,不能自拔.我們要在真理上合一,領受純正的教訓.第二是在靈交上的合一:第15節要連於元首耶穌基督,與神有一個良好的關係.我們若尊主為大,接受神的主權,我們就不論背景上的差距,仍然能夠交通,互相尊重.在靈裡不能交通,可能是彼此不能坦誠,開放自己；沒有用愛心說誠實話.我們不應掩飾自己,並且要保持良好的開放態度,接納與任何按真理生活的肢體相交.另一方面我們要凡事長進,若一人長進,一人放鬆生活,不願追求；彼此就會有差距,在靈裡相交就出現問題.能在靈裡相交,我們就能同心合一.第三是在事工上的合一；第16節百節各按各職,各人在不同崗位和職份上努力成長,事奉；照著各體的功用彼此相助,便叫身體漸漸增長,在愛中建立自己.神的教會需要增長,被建立起來；就必須要人人起來,並且各人一同參與事奉.我們在神的家中,要使肢體百節各按各職；我的靈性有沒有長進,事奉有沒有更投入,得力?我們應時作反省,保持一種戰爭的心態,作耶穌基督的精兵.我們要為神的家做得更好,並且要為神作大事,好像現代宣教之父威廉克里的盼望:「為神祈求大事,為神成就大事.」
\newline
\newline
今天我們面對陰險可怖的魔鬼,更要小心提防；同心協力的建立神的家和教會,使神得著最高的榮耀.
\newline
\newline


\newpage

\section{第60屆~港九培靈會~唐佑之~第一講}
\label{sec:1216}
\textbf{安慰}
\newline
\newline
連結: \href{https://www.hkbibleconference.org/session-message/view/1216}{\texttt{https://www.hkbibleconference.org/session-message/view/1216}}
\newline
\newline
\hyperref[sec:1211]{< < < PREV SERMON < < <}
~
\hyperref[sec:toc]{[返目錄]}
~
\hyperref[sec:1220]{> > > NEXT SERMON > > >}
\newline
\newline
以賽亞書 40:1-40:10
\newline
\begin{longtable}{cl}
\hline
\hline
章節 & 經文 (和合本修訂版)\\
\hline
40:1 & \begin{tabularx}{0.7\textwidth}{X} 你們的神說：「要安慰，安慰我的百姓。 \end{tabularx} \\ \\ \relax
40:2 & \begin{tabularx}{0.7\textwidth}{X} 要對耶路撒冷說安慰的話，向它宣告，它的戰爭已結束，它的罪孽已赦免；它為自己一切的罪，已從耶和華手中加倍受罰。」 \end{tabularx} \\ \\ \relax
40:3 & \begin{tabularx}{0.7\textwidth}{X} 有聲音呼喊著：「要在曠野為耶和華預備道路，在沙漠為我們的神修直大道。 \end{tabularx} \\ \\ \relax
40:4 & \begin{tabularx}{0.7\textwidth}{X} 一切山窪都要填滿，大小山岡都要削平；陡峭的要變為平坦，崎嶇的必成為平原。 \end{tabularx} \\ \\ \relax
40:5 & \begin{tabularx}{0.7\textwidth}{X} 耶和華的榮耀必然顯現，凡有血肉之軀的都一同看見，因為這是耶和華親口說的。」 \end{tabularx} \\ \\ \relax
40:6 & \begin{tabularx}{0.7\textwidth}{X} 有聲音說：「你喊叫吧！」我說：「我喊叫甚麼呢？」凡有血肉之軀的盡都如草，他的一切榮美像野地的花。 \end{tabularx} \\ \\ \relax
40:7 & \begin{tabularx}{0.7\textwidth}{X} 耶和華吹一口氣，草就枯乾，花也凋謝。百姓誠然是草； \end{tabularx} \\ \\ \relax
40:8 & \begin{tabularx}{0.7\textwidth}{X} 草必枯乾，花必凋謝，惟有我們神的話永遠立定。 \end{tabularx} \\ \\ \relax
40:9 & \begin{tabularx}{0.7\textwidth}{X} 報好信息的錫安哪，要登高山；報好信息的耶路撒冷啊，要極力揚聲。揚聲不要懼怕，對猶大的城鎮說：「看哪，你們的神！」 \end{tabularx} \\ \\ \relax
40:10 & \begin{tabularx}{0.7\textwidth}{X} 看哪，主耶和華必以大能臨到，他的膀臂必為他掌權；看哪，他的賞賜在他那裡，他的報應在他面前。 \end{tabularx} \\ \\
[1ex]
\hline
\hline
\end{longtable}
在港九培靈研經會60週年的盛會中,我們大家都存著感謝和快樂的心.剛才詩班所唱詩歌,我們深切地回應,因為裡面有敬拜,歌頌,感謝；更有祈求,我們祈求的是:仰望聖靈的能力在我們每人身上發生奇妙的能力.我們竭力呼喊；求主復興我們!願主恩待!讓我們在主前真正蒙福,且蒙福更深,相信主必賜福如同已往的60年一樣.在這60年當中,我們民族經歷了最大而危險的憂患時代；神恩待我們,以祂的真理向我們施恩.主耶穌說:「父啊!我真是願意脫離這個時代.」主所處的時代也是個憂患時代；但主說:「我來乃是為這時代而生.」感謝主!我們處身時代憂患的環境中,神不但賜福給我們,還要我們把這福份帶出去,成就福音的使命.
\newline
\newline
本書四十:1-11可說是頌讚大合唱的序曲,從中我們聽到四種聲音,1-2節恩惠的聲音,3-5節真理的聲音,6-8節應許的聲音,9-11節見證的聲音.主若幫助我們,我們必能聽見這些聲音；聖靈向我們所說的話,凡有耳的就應當聽.
\newline
\newline
基督徒生活,最重要的藝術是靈性的藝術,我們須用受教育者的耳朵來聽；神說:「你們要安慰我的百姓,要對耶路撒冷說安慰的話,又向他宣告,他爭戰的日子已滿了,他的罪孽赦免了,他為自己的一切罪,從耶和華手中加倍受罰.」安慰是福音的內容,是福音使者的經驗,是福音使命應有的態度.我們在這憂患時代最需要的是神的安慰,我們負著福音的使命務將安慰帶給人.安慰這字在舊約用的很奇妙；原意是喘氣,嘆息,包含憂傷和說不出的心意和情緒.創世記第六章,耶和華看見人犯罪,心裡難過,後悔造人.
\newline
\newline
我們的主,我們的神是受苦的神,當教會受苦時；神為我們,為教會,為世人焦急；我們能否體會祂迫切的心賜?「你們的神說,你們要安慰安慰我的百姓.」強烈加重語氣,表明神對我們有無限的心意.「你們」是多數式,因為神覺得祂所愛的人實在太親密；在希伯來原文是「愛的親熱.」偉大的神為我們憂傷,所以安慰我們；祂知道世上有很多憂患,祂賜給我們安慰的福音.「你們的神」「我們的百姓」何等密切的關係!聖經中每逢題及「我的百姓」都帶著熱烈之感；但是神的百姓實在太令神失望.何西阿書載:「羅阿米不是我們的百姓.」(何9)神十分失望,神揀選以色列民,並非因為他們是優秀的民族,非因他們有優秀的文化；神乃是揀選世上最軟弱之民,藉以顯明祂的能力.照樣,神揀選我們作祂的兒女,非因我們比別人好；乃是要在這些軟弱的人身上顯示出祂的大能力.我們多麼需要神的憐憫!可惜神的百姓失敗了,沒有把福音帶出去,沒有福音的見證；雖是神的兒女；但神要審判,這是歷史的見證,神不僅叫我們蒙受祂的恩惠,神也要我們把福音帶出去；若無福音的見證,審判就要臨到我們身上；因為審判要從神的家起首,教會要受審判,因神是公義的神；但是神不捨得,所以為此憂傷；祂對耶路撒冷說安慰的話,耶路撒冷是大軍的京城,是神的居所,是聖潔的所在；或說耶路撒冷是我們的教會,神針對教會說安慰的話,這裡的安慰有所不同；對錫安的安慰帶著傷感；而此處安慰的意思是「摸到我們的心」言者並非空洞的理論,非理想；乃是祂的話能夠摸到我們每個人的心；因為祂願意向我們說話,聖靈所說的話,凡有耳的就應當聽.
\newline
\newline
關於「安慰的話」有很多不同的翻譯,有譯為溫和的話,恩慈的話,撫慰的話,誠心順服的話,鼓勵的話.神要向我們說話,我們聽見了沒有?求主幫助!我們要聽,像個受教者.聽,不是隨便聽,要細心,專心,誠心,全心,全信地聽；聽而接受,聽而信；信道是從聽道來的；聽而信從祂,倚靠祂.今日教會一般基督徒中,以為信心就是信靠.如果我們信靠神,是好的；但有時我們連倚靠神的心都沒有.信心不單是信靠,且是信從.我們多麼需要神憐憫!「要對耶路撒冷說安慰的話,又向祂宣告」意即不單說而且成就,這是希伯來原文的含意.「草必枯乾,花必凋殘,惟有我們神的話,必永遠立定.」(8)神的話多麼寶貝!
\newline
\newline
「宣告說,他爭戰的日子已滿了.」爭戰意即服役,苦役.人生如服苦役!約伯記題及人生的虛空.然而日子滿了,神的恩典臨到我們了,教會的復興就在眼前；但是,我們到底有此認識嗎?有否為此而努力?「他的罪孽赦免了」神應許我們救贖之恩,饒恕了我們的罪孽.「他為自己的一切罪,從耶和華手中加倍受罰.」難道神這麼殘忍嗎?所謂加倍就是「雙重」一方面是我們內裡的感受,另方面是我們外面的見證；當我們的罪得赦免,裡面有此經驗而表達出來,此乃福音的使命.願神恩待,向我們施恩!
\newline
\newline
以色列人雖是神的子民,但是審判從神的家起首；他們被擄到外邦,但這日子已經過去；人失敗,神沒有失敗；人軟弱以致沒見證的能力；但深信主能復興我們,赦免我們,除掉我們的痛苦；當我們得此經驗,然後把此經驗帶出去,這就是福音的內容,也是福音的使命.
\newline
\newline
以色列人第一次曾經從埃及為奴之地出來.第二次出埃及,不一定是從埃及國出來,而是重新從罪之奴隸中出來,從撒旦權勢下出來.今日我們教會需要第二次出埃及的經驗,第一次出埃及乃從罪惡之中出來；但曾幾何時我們冷淡軟弱!我們的禱告是「我們需要復興,需要第二次出埃及的經驗.」求主幫助我們認識在此憂患時代,必須將福音使命帶出去；還有一點點的時候,那要求的就來,並不遲延.保羅說:「若有人不愛主,這人是可咒可詛的.」主必要來,這是我們的心願!我們曾經失敗,軟弱,沒有討神的喜悅；但從今開始,我們務要重新振作.求主復興我們!求主安慰的靈安慰我們,差遣我們；叫我們將安慰的福音帶出去,讓許多人蒙受主的恩典.
\newline
\newline
主幫助我們每個人,每一家庭,每一教會.這是我們的禱告.
\newline
\newline


\newpage

\section{第60屆~港九培靈會~唐佑之~第二講}
\label{sec:1220}
\textbf{重新得力}
\newline
\newline
連結: \href{https://www.hkbibleconference.org/session-message/view/1220}{\texttt{https://www.hkbibleconference.org/session-message/view/1220}}
\newline
\newline
\hyperref[sec:1216]{< < < PREV SERMON < < <}
~
\hyperref[sec:toc]{[返目錄]}
~
\hyperref[sec:1222]{> > > NEXT SERMON > > >}
\newline
\newline
以賽亞書 41:1-41:29
\newline
\begin{longtable}{cl}
\hline
\hline
章節 & 經文 (和合本修訂版)\\
\hline
41:1 & \begin{tabularx}{0.7\textwidth}{X} 眾海島啊，在我面前靜默；萬民要重新得力，讓他們近前來陳述，我們可以彼此辯論。 \end{tabularx} \\ \\ \relax
41:2 & \begin{tabularx}{0.7\textwidth}{X} 誰從東方興起一人，憑公義召他來到腳前？誰將列國交給他，使他管轄列王，把他們如灰塵交與他的刀，如風吹的碎秸交與他的弓？ \end{tabularx} \\ \\ \relax
41:3 & \begin{tabularx}{0.7\textwidth}{X} 他追趕君王，安然走過，快速地腳不落地。 \end{tabularx} \\ \\ \relax
41:4 & \begin{tabularx}{0.7\textwidth}{X} 誰做成這事，從起初宣召歷代呢？就是我－耶和華！我是首先的，也與末後的同在。 \end{tabularx} \\ \\ \relax
41:5 & \begin{tabularx}{0.7\textwidth}{X} 眾海島看見就都害怕，地極也都戰兢，他們近前來； \end{tabularx} \\ \\ \relax
41:6 & \begin{tabularx}{0.7\textwidth}{X} 各人互相幫助，對弟兄說：「壯膽吧！」 \end{tabularx} \\ \\ \relax
41:7 & \begin{tabularx}{0.7\textwidth}{X} 木匠鼓勵銀匠，用鎚子打光的鼓勵打砧的，對銲工說：「銲得好！」又用釘子釘穩，免得它倒下。 \end{tabularx} \\ \\ \relax
41:8 & \begin{tabularx}{0.7\textwidth}{X} 惟你以色列，我的僕人，雅各，我所揀選的，我朋友亞伯拉罕的後裔， \end{tabularx} \\ \\ \relax
41:9 & \begin{tabularx}{0.7\textwidth}{X} 你是我從地極領來，從地角召來的，我對你說：「你是我的僕人；我揀選你，並不棄絕你。」 \end{tabularx} \\ \\ \relax
41:10 & \begin{tabularx}{0.7\textwidth}{X} 你不要害怕，因為我與你同在；不要驚惶，因為我是你的神。我必堅固你，幫助你，用我公義的右手扶持你。 \end{tabularx} \\ \\ \relax
41:11 & \begin{tabularx}{0.7\textwidth}{X} 看哪，凡向你發怒的都抱愧蒙羞，與你相爭的必如無有，並要滅亡。 \end{tabularx} \\ \\ \relax
41:12 & \begin{tabularx}{0.7\textwidth}{X} 與你爭鬥的，你要尋找他們，卻遍尋不著；與你爭戰的必如無有，成為虛無。 \end{tabularx} \\ \\ \relax
41:13 & \begin{tabularx}{0.7\textwidth}{X} 因為我耶和華－你的神必攙扶你的右手，對你說：「不要害怕！我必幫助你。」 \end{tabularx} \\ \\ \relax
41:14 & \begin{tabularx}{0.7\textwidth}{X} 蟲子雅各，以色列人哪，不要害怕！我必幫助你；救贖你的是以色列的聖者。這是耶和華說的。 \end{tabularx} \\ \\ \relax
41:15 & \begin{tabularx}{0.7\textwidth}{X} 看哪，我使你成為全新的打穀機，齒輪銳利；你要把山嶺打得粉碎，使岡陵如同糠秕。 \end{tabularx} \\ \\ \relax
41:16 & \begin{tabularx}{0.7\textwidth}{X} 你要簸揚它們，風要將它們吹去；旋風要颳散它們。你卻要以耶和華為喜樂，因以色列的聖者誇耀。 \end{tabularx} \\ \\ \relax
41:17 & \begin{tabularx}{0.7\textwidth}{X} 困苦貧窮人尋找水，卻尋不著；他們因口渴，舌頭乾燥。我－耶和華必應允他們，我---以色列的神必不離棄他們。 \end{tabularx} \\ \\ \relax
41:18 & \begin{tabularx}{0.7\textwidth}{X} 我要在光禿的高地開江河，在谷中開泉源；我要使沙漠變為水池，使乾地變為湧泉。 \end{tabularx} \\ \\ \relax
41:19 & \begin{tabularx}{0.7\textwidth}{X} 我要在曠野栽植香柏樹、皂莢樹、番石榴樹，和野橄欖樹。在沙漠一同栽上松樹、杉樹，和黃楊樹， \end{tabularx} \\ \\ \relax
41:20 & \begin{tabularx}{0.7\textwidth}{X} 好叫人看見，知道，思想，明白；這是耶和華親手做的，是以色列的聖者所造的。 \end{tabularx} \\ \\ \relax
41:21 & \begin{tabularx}{0.7\textwidth}{X} 耶和華說：「你們要呈上你們的案件。」雅各的君王說：「你們要提出你們的理由。」 \end{tabularx} \\ \\ \relax
41:22 & \begin{tabularx}{0.7\textwidth}{X} 讓它們近前來，告訴我們將來要發生甚麼事！你們要說明先前發生的事，好讓我們思索；或者告訴我們將來的事，使我們得知事情的結局。 \end{tabularx} \\ \\ \relax
41:23 & \begin{tabularx}{0.7\textwidth}{X} 你們要指明未來的事，使我們知道你們是神明！你們或降福，或降禍，好使我們驚奇，一同觀看。 \end{tabularx} \\ \\ \relax
41:24 & \begin{tabularx}{0.7\textwidth}{X} 看哪，你們屬乎虛無，你們的作為也屬虛空；那選擇你們的是可憎惡的。 \end{tabularx} \\ \\ \relax
41:25 & \begin{tabularx}{0.7\textwidth}{X} 我從北方興起一人，他從日出之地而來，是求告我名的；他必踩踏掌權者，如踩踏泥土，又如陶匠踹泥一般。 \end{tabularx} \\ \\ \relax
41:26 & \begin{tabularx}{0.7\textwidth}{X} 有誰從起初宣佈這事，使我們知道呢？有誰從先前指明，使我們說「他是對的」呢？沒有人宣佈，沒有人指明，也沒有人聽見你們的話。 \end{tabularx} \\ \\ \relax
41:27 & \begin{tabularx}{0.7\textwidth}{X} 我首先對錫安說，看哪，他們在此！我要將一位報好信息的賜給耶路撒冷。 \end{tabularx} \\ \\ \relax
41:28 & \begin{tabularx}{0.7\textwidth}{X} 然而我觀看，並無一人；我詢問的時候，他們中間也沒有謀士可回答。 \end{tabularx} \\ \\ \relax
41:29 & \begin{tabularx}{0.7\textwidth}{X} 看哪，他們盡是麻煩，所做的工都屬虛無；所鑄的偶像是風，是虛空。 \end{tabularx} \\ \\
[1ex]
\hline
\hline
\end{longtable}
在時代的憂患中,我們須有福音的使命.先知屬靈的意識；耶和華是歷史之主,每個歷史時代中；神一定要施恩,一定要成就祂救贖的工作.當我們思想自然界,我們知道神是創造的主宰,當我們思想歷史的領域,我們知道神是救贖的主；在歷史中,神的工作是救贖,祂在教會為我們預備救贖的恩典,要我們傳揚救贖的恩典,這就是福音的使命.先知的意識能夠思想憂患的感受,這種感受不單先知僅有,非他們思想出來的,也非他們的敏感而有的,乃是神的心意,祂有無限的憂患,祂為我們受苦；十字架不過是表達神的苦,所以先知是憂患的意識；受苦的主為我們的苦難而受苦；先知有此感受時便表達出來.
\newline
\newline
本書第四十至五十五章,先知不但把憂患的意識表達,用直接的情緒表達非常透徹,美麗,帶著淚,嘆息,也帶著歡笑的聲音.當我們讀先知書時,我們也和先知一樣的悲嘆,體會,流淚,嘆息.
\newline
\newline
昨晚講四種聲音,因時間關係沒有詳細解釋,現在重複一下.
\newline
\newline
恩惠的聲音(四十1-2)赦免的恩惠,神是赦免的神.「主耶和華啊,你若究察罪孽,誰能站立得住呢?但在你有赦免之恩.......」(詩一三○:3-4)恩惠的聲音特別強調赦免的恩典,有赦罪的恩典也有正式的安慰,使我們唱出得救的樂歌.
\newline
\newline
真理的聲音(3-5)「道成肉身,住在我們中間,充充滿滿的有恩典,有真理.......」(約14)先知第二個訊息著重預言方面的真理.每逢我們題及預言,就論到將來要發生的事,將來要發生的並非怪異的事；而是神按著祂奇妙的恩典,把祂的真理祂的計劃帶出來.
\newline
\newline
應許的聲音(6-8)神的應許和真理永遠立定.
\newline
\newline
見證的聲音(9-11)亦即福音宣道的聲音.
\newline
\newline
昨晚我們已經思想了首兩節.「有人聲喊著說:『在曠野預備耶和華的路.』」或作「曠野有人喊著說『當預備耶和華的路.』」我認為後者意思較正確.曠野是貧乏無水荒亂無倫之地；但是竟然有人聲,有安慰,有真理的聲音.這實在值得我們反覆思想.
\newline
\newline
神為我們預備的,不是曠野而是樂園；但人所要的不是樂園而是曠野.耶穌來到世上,要變曠野為樂園,這是復興的信息.我們須在神前,認識樂園,恢復樂園.現在世界科技雖發達,但這世界就是曠野；十幾年來西方文壇,荒原詩人,描述這世界是荒漠的田野.充滿許多苦難,這是今日教會所能體會的,弟兄姊妹在曠野生活卻沒有憂患的感覺,在世俗環境中隨波逐流；我們沒有在世界狂潮裡作中流砥柱,沒有傳福音的悟性.求神幫助我們!讓我們體會曠野應有的聲音,聽見曠野的聲音是甚麼聲音.
\newline
\newline
路加第九章載,當耶穌在山上變像時,有摩西和以利亞出現；耶穌說關於祂去世的事.「去世」在路加記載就是出埃及.當人子主耶穌去世,將罪人拯救出來；我們必須有出埃及的經驗；今日教會需要有新的出埃及經驗；我們要從罪惡,軟弱,冷淡,災難中出來；這是教會所需要的,是真正復興的意義.當以色列人從被擄之地出來時,需要人領導他們,他們過紅海,經曠野,後來又經過約但河,這是他們首次出埃及的經驗；他們到了乾旱曠野地方,神引導他們,除去一切的障礙,把高低的變為平坦,崎嶇的變為平原,這就是其中所描寫的意義.今日教會需要這樣的工作.當主耶穌剛到世上,曠野的人聖約翰的聲音,他大聲疾呼「天國近了.」他是福音使者,為主耶穌預備道路.當主耶穌第二次再來時,教會你,我都要作先鋒,成為福音使者,如同約翰在曠野發出人聲,而且我們必須把一切障礙物除去,迎接耶穌再次到世上時永遠作王.這是我們每時每刻應該注意的事.我們有沒為耶穌第二次再來作應作的事?主禱文中「願你的國降臨,願你的旨意行在地上,如同行在天上.」這是多麼重要的話!當我們講到主再來的事,就想到很多事將要出現；預言要應驗雖然不錯,但大部份都不對；因為沒好好研究聖經,我們常講主再來的預兆；可是最重要的反而不談,主耶穌說天國福音要傳遍普天下,然後末期才來到.(太24)有時我們講及基督再來的真理,帶著好奇的態度來分析真理,再來看正確的實據；主耶穌再來最重要者在於福音要傳遍天下,這是你我應負的責任,應該作的事,為主預備道路.主快再來了,教會沒有好好地傳福音,沒有好好的為主預備道路.求主憐憫!叫我們積極,正面,樂觀的等候主再來；我們常注意消極方面的時局等等......而不積極去傳福音；時代的憂患越來越嚴重,教會傳福音的使命也愈加重大,沙漠地要修平我們神的道,一切山窪都要填滿,大小山岡都要削平,高高低低的要改為平坦,神的福音才可以作此工作.
\newline
\newline
聖經說,神要使驕傲的人降卑,卑微的人升高.這就是福音的內容,是救恩的功效.神提拔我,從低微的地步救拔我,這是原則,是實際的現象.歷史方面來看,邦國興衰無定；但這一切都要過去,人的成功失敗也都要過去.神叫我們預備道路,但有時崎嶇難走,有時遇高山的困難.傳福音,教會的事,我們的見證；要作中流砥柱並非易事,唯求主幫助!主在地上設立教會,為要我們享受祂的恩典,教我們在祂恩典中操練,經過困難,危險,付上代價,有足夠的經驗,然後我們知道主的力量非常奇妙；因為主的能力在軟弱的人身上顯得完全.自己的努力,靠主的恩典,還必須付代價,必須努力與神同工,神在歷史上之目的必然達到,救贖的恩典是神在歷史上的目的.
\newline
\newline
「耶和華的榮耀必然復現.」「榮耀」希臘人譯為救恩,意思更為顯明,事實上,神在救恩上顯明祂的榮耀.神的榮耀就是神的公義和聖潔,當救恩實現,神的公義就表現,聖潔也表現出來.「榮耀」更重要的意思是「同在」,有神的同在；摩西在山上,神指示他會幕的樣式；會幕是神與人相會之處；當摩西獻祭敬拜時,神的榮耀充滿會幕.何處有敬拜就有神的同在,何處有神同在就有神的榮耀.所羅門建造聖殿.奉獻和獻祭,神的榮光就充滿聖殿；後來以色列人失敗,聖殿被毀,以色列人被擄到外邦,當他們回國,重建聖殿.哈該書說:「這殿後來的榮耀比先前的榮耀更大.」以後還有更大的,當耶穌再次來到世上「道成了肉身,住在我們中間......」(約一:14)「住」原意是搭帳棚.耶穌基督就是會幕,藉著祂,我們與神相會.耶穌說到聖殿毀壞,三天內要再建立,(約二)猶太人以為不可能,後來才明白這是指耶穌的身體而言；因為主基督是神的榮耀,要表現耶和華救恩的能力.耶穌死後埋葬,身體仍在地上,有神的同在,榮耀,救恩,如以弗所書說教會就是基督的身體,是充滿萬有者所充滿的.充滿了神的榮耀；但還有更大的榮耀,林前五章提到每個信徒都是聖殿.我們都是神的殿,都是聖靈的殿,是屬靈的,聖潔的,我們有神的榮耀.「耶和華的榮耀必然顯現.」在今日之教會顯現,在你我身上的顯現；當我們把福音帶出去時,把救恩實現出來；耶和華的榮耀必然顯現,「凡有血氣的,必一同看見,因為這是耶和華親口說的.」這指全世界的人都有機會明白福音,得著救恩；若他們拒絕,是他們的責任；但若我們不傳福音就有禍了；因我們沒讓他們明白.每個世人都應該看見,因為這是耶和華親口說的,神不但「說」也必「成就」；不過你我必須付上福音的代價.
\newline
\newline
神應許的聲音(6)「有人聲說,你喊叫罷......凡有血氣的,盡都如草,他的美容,都像野地的花.」美容應譯為信實.人的信實如花草,容易凋殘.世界不可靠,人靠不住,不能依賴人的信實,惟神是信實的.保羅講過,我們縱然失信,但神是信實的.「草必枯乾,花必凋殘,因為耶和華的氣吹在其上.」(7)耶和華的氣吹在人的鼻孔裡,那人就成為有靈的活人.(創二)但是第7節所說的「氣」是神的怒氣,指神的審判,忿怒；當熱風吹起,花便凋殘零落；當審判之日,所有世上的人都要滅亡；因世界和其上的情慾都要過去,惟獨遵行神旨意的,才能永遠常存.人和世上一切事物都不可靠,唯神話可靠,神的話永遠立定,這是神的應許,必然實現.我們多得神的話,就更有保障,更有確據；我們得鼓勵,得安慰；神的話是我們生活最主要的中心.在時代的憂患中,神的話是我們傳福音的力量,福音的內容是神的話；當我們把神的話帶出來時,必有功效；因為神的話安定在天,永不改變,神的話沒一句不帶著能力的,可惜我們對神的話認識太少,讀的,講的,遵行的太少.這是今天我們的失敗,也是教會的軟弱.我們都需要神的話!
\newline
\newline
傳福音報好信息:「報好信息給錫安的阿,你要登高山；報好信息給耶路撒冷的阿,你要極力揚聲......」(9)現代中文聖經譯本較為正確,原意是「錫安阿,要報好信息,耶路撒冷阿,要報好信息.」如何報好的信息?要登高山,我們須有山頂上的經驗,與神有親密交通的經驗；上得高,看得遠,看見福音的需要,明白福音的使命；我們須在神前等候,就能聽見神的聲音,明白神的命令,知道神給我們的囑咐.登高山意即自動的,不勉強,不盲從；這是神的心意.讓我們靠著主的恩典,明白福音的需要.傳福音的人!敬拜神的人!不要懼怕!在愛裡沒有懼怕；要歡樂,要極力揚聲,傳揚神救恩奇妙；我們不能不感謝,歌頌,我們必須有此興奮態度.
\newline
\newline
「看哪,你們的神」(9末句)神來到了,這是很重要的宣告.報好消息能保持教會的青春.教會是基督的新婦,基督用水把教會洗乾淨,把我們的縐紋等類都除掉.(弗五)新婦為何有縐紋呢?因今天教會充滿了縐紋等類的病,處於衰老之中.要報好信息,傳福音,教會就有活力,才能保持教會的青春,表明教會是基督的新婦.
\newline
\newline
「主耶和華必像大能者臨到,他的膀臂必為他掌權......」(10)神的膀臂是大能者的膀臂,另方面是最溫和的膀臂.「他必像牧人牧養自己的羊群,用膀臂聚集羊羔抱在懷中,慢慢引導那乳養小羊的.」(11)我們也許不明白,怎麼神的膀臂大能又溫和嗎?我們的神真是奇妙?我們都是祂的小羊,抱在祂懷裡且慢慢引導,祂知道我們走不快,知道我們軟弱,祂何等大的忍耐,這是神奇妙的大恩典!我們當在神前感謝讚美祂,因祂實在有無限的恩惠,祂也是大能者,且是位慈愛的牧人.
\newline
\newline
「疲乏的,他賜能力,軟弱的,他加力量；就是少年人也要疲乏困倦,強壯的也必全然跌倒；但那等候耶和華的,必重新得力,他們如鷹展翅上騰,他們奔跑卻不困倦,行走卻不疲乏.」(29-31)我們是小羊,祂用膀臂抱著我們；但我們也是鷹鳥,要展翅上騰.「等候耶和華的,必重新得力.」而後我們就能飛翔.有些弟兄姊妹呆板地等候；但有的覺得等候太消極.名佈道家慕勒先生有次搭船,半夜輪船著火；有的人不知所措,有的基督徒禱告；有的基督徒認為等禱告完了,火也燒透了；當前之務是急須滅火.有的人就問慕勒,應先禱告或應先救火呢?他說:「邊救火邊禱告.」今日基督徒卻是邊救火邊等候.我們須等候,禱告,仰望；但我們一定要行動,要學習,要長進.申命記卅二章摩西提及神訓練以色列民,如母鷹訓練小鷹.母鷹喂養小鷹,若長此以往,小鷹永不長進飛翔.直到母鷹毀壞巢窩,用大力的膀臂攜帶小鷹高空飛翔,有時放下膀臂,讓小鷹拼命掙扎,然後加以救援；如此一次次地訓練小鷹.你我都應在神恩典中受訓；祂愛我們,但祂有時叫我們不得安逸,遭遇憂患,我們才能學習飛翔,我們勤於學習,真正在神前等候,我們就重新得力,如鷹展翅上騰,到那時,我們奔跑卻不困倦,行走卻不疲乏.如何「飛」「跑」「走」?這似乎次序顛倒；其實這裡有很重要的意義:飛速度快；跑需較長的時間；走更需要力量.我們屬靈的長進,不需要太快而需要穩健；在福音道路上行走要步步穩健,靠主恩典,向前又向上.
\newline
\newline


\newpage

\section{第60屆~港九培靈會~唐佑之~第三講}
\label{sec:1222}
\textbf{見證}
\newline
\newline
連結: \href{https://www.hkbibleconference.org/session-message/view/1222}{\texttt{https://www.hkbibleconference.org/session-message/view/1222}}
\newline
\newline
\hyperref[sec:1220]{< < < PREV SERMON < < <}
~
\hyperref[sec:toc]{[返目錄]}
~
\hyperref[sec:1223]{> > > NEXT SERMON > > >}
\newline
\newline
以賽亞書 43:1-43:28
\newline
\begin{longtable}{cl}
\hline
\hline
章節 & 經文 (和合本修訂版)\\
\hline
43:1 & \begin{tabularx}{0.7\textwidth}{X} 雅各啊，創造你的耶和華，以色列啊，造成你的那位，現在如此說：「你不要害怕，因為我救贖了你；我曾提你的名召你，你是屬我的。 \end{tabularx} \\ \\ \relax
43:2 & \begin{tabularx}{0.7\textwidth}{X} 你從水中經過，我必與你同在，你渡過江河，水必不漫過你；你在火中行走，也不被燒傷，火焰必不燒著你身。 \end{tabularx} \\ \\ \relax
43:3 & \begin{tabularx}{0.7\textwidth}{X} 因為我是耶和華－你的神，是以色列的聖者－你的救主；我使埃及作你的贖價，使古實和西巴代替你。 \end{tabularx} \\ \\ \relax
43:4 & \begin{tabularx}{0.7\textwidth}{X} 因我看你為寶貝為尊貴；又因我愛你，所以使人代替你，使萬民替換你的生命。 \end{tabularx} \\ \\ \relax
43:5 & \begin{tabularx}{0.7\textwidth}{X} 你不要害怕，因我與你同在；我必領你的後裔從東方來，又從西方召集你。 \end{tabularx} \\ \\ \relax
43:6 & \begin{tabularx}{0.7\textwidth}{X} 我要對北方說，交出來！對南方說，不可扣留！要將我的兒子從遠方帶來，將我的女兒從地極領回， \end{tabularx} \\ \\ \relax
43:7 & \begin{tabularx}{0.7\textwidth}{X} 就是凡稱為我名下的人，是我為自己的榮耀創造的，是我所塑造，所做成的。」 \end{tabularx} \\ \\ \relax
43:8 & \begin{tabularx}{0.7\textwidth}{X} 你要將有眼卻瞎、有耳卻聾的民都帶出來！ \end{tabularx} \\ \\ \relax
43:9 & \begin{tabularx}{0.7\textwidth}{X} 任憑萬國聚集，任憑萬民會合。他們當中誰能說明，並將先前的事指示我們呢？讓他們帶來見證，顯明他們有理，看是否聽見的人會說：「果然是真的。」 \end{tabularx} \\ \\ \relax
43:10 & \begin{tabularx}{0.7\textwidth}{X} 你們是我的見證，是我所揀選的僕人，為了要使你們知道，且信服我，又明白我就是耶和華。在我以前沒有任何被造的真神，在我以後也必沒有。這是耶和華說的。 \end{tabularx} \\ \\ \relax
43:11 & \begin{tabularx}{0.7\textwidth}{X} 我，惟有我是耶和華；除我以外沒有救主。 \end{tabularx} \\ \\ \relax
43:12 & \begin{tabularx}{0.7\textwidth}{X} 我曾指示，我曾拯救，我曾說明，並沒有外族的神明在你們中間。你們是我的見證，我是神。這是耶和華說的。 \end{tabularx} \\ \\ \relax
43:13 & \begin{tabularx}{0.7\textwidth}{X} 自有日子以來，我就是神，誰也不能救人脫離我的手。我要行事，誰能逆轉呢？ \end{tabularx} \\ \\ \relax
43:14 & \begin{tabularx}{0.7\textwidth}{X} 耶和華---你們的救贖主、以色列的聖者如此說：「因你們的緣故，我已派遣人到巴比倫去；要使迦勒底人都如難民，坐自己素來宴樂的船下來。 \end{tabularx} \\ \\ \relax
43:15 & \begin{tabularx}{0.7\textwidth}{X} 我是耶和華－你們的聖者，是創造以色列的，是你們的君王。」 \end{tabularx} \\ \\ \relax
43:16 & \begin{tabularx}{0.7\textwidth}{X} 那在滄海中開道，在大水中開路， \end{tabularx} \\ \\ \relax
43:17 & \begin{tabularx}{0.7\textwidth}{X} 使戰車、馬匹、軍兵、勇士一同出來，使他們仆倒，不再起來，使他們滅沒，好像熄滅之燈火的耶和華如此說： \end{tabularx} \\ \\ \relax
43:18 & \begin{tabularx}{0.7\textwidth}{X} 「你們不要追念從前的事，也不要思想古時的事。 \end{tabularx} \\ \\ \relax
43:19 & \begin{tabularx}{0.7\textwidth}{X} 看哪，我要行一件新事，如今就要顯明，你們豈不知道嗎？我必在曠野開道路，在沙漠開江河。 \end{tabularx} \\ \\ \relax
43:20 & \begin{tabularx}{0.7\textwidth}{X} 野地的走獸要尊敬我，野狗和鴕鳥也必尊敬我。因我使曠野有水，使沙漠有河，好賜給我的百姓、我的選民喝。 \end{tabularx} \\ \\ \relax
43:21 & \begin{tabularx}{0.7\textwidth}{X} 這百姓是我為自己造的，為要述說我的美德。」 \end{tabularx} \\ \\ \relax
43:22 & \begin{tabularx}{0.7\textwidth}{X} 「雅各啊，你並沒有求告我；以色列啊，你倒厭煩我。 \end{tabularx} \\ \\ \relax
43:23 & \begin{tabularx}{0.7\textwidth}{X} 你並沒有將你的羊帶來獻給我做燔祭，也沒有用牲祭尊敬我；我未曾因素祭使你操勞，也沒有因乳香使你厭煩。 \end{tabularx} \\ \\ \relax
43:24 & \begin{tabularx}{0.7\textwidth}{X} 你沒有用銀子為我買香菖蒲，也沒有用祭物的油脂使我飽足；倒使我因你的罪惡操勞，使我因你的罪孽厭煩。 \end{tabularx} \\ \\ \relax
43:25 & \begin{tabularx}{0.7\textwidth}{X} 我，惟有我為自己的緣故塗去你的過犯，我也不再記得你的罪惡。 \end{tabularx} \\ \\ \relax
43:26 & \begin{tabularx}{0.7\textwidth}{X} 你儘管提醒我，讓我們來辯論；儘管陳述，自顯為義。 \end{tabularx} \\ \\ \relax
43:27 & \begin{tabularx}{0.7\textwidth}{X} 你的始祖犯罪，你的師傅違背我； \end{tabularx} \\ \\ \relax
43:28 & \begin{tabularx}{0.7\textwidth}{X} 因此，我要凌辱聖所的領袖，使雅各遭毀滅，使以色列受辱罵。」 \end{tabularx} \\ \\
[1ex]
\hline
\hline
\end{longtable}
我們相信神的救恩奇妙,聖靈的心意奇妙.「眾海島阿,當在我面前靜默,眾民當重新得力,都要近前來才可以說話......」(1)海島是遙遠的地方；神的話大有能力,廣泛地深入人心.「我耶和華是首先的,也與末後的同在.」(四十一:4)全本聖經讓我們知道,神是為我們信心創始成終的主,祂是首先的,也是末後的,是永永遠遠的；我們當在祂面前俯服敬拜；因祂是生命的主.
\newline
\newline
本章有許多重要的話；但我們在本章只能看到神的能力.「我是......我必......我要......」神有祂的心意要成就在我們身上.「我是」是很重要的字眼；有人說約翰福音中有七次「我是」.耶穌基督特別表明祂的身份也表明祂的能力,祂曾說祂是好牧人,是羊的門,是葡萄樹等等.「我是耶和華」在舊約中反覆述說這句話,表明神的恩典和祂恩典的啟示,也表明祂啟示的恩典；每逢我們聽見這話,就知道祂有說不盡的恩典.許多人以為舊約聖經是講律法；這是嚴重的錯誤；其實舊約非常注重恩典.當神叫摩西曉諭以色列民如何守十誡,神並沒有吩咐他當講或不當講的是甚麼.很多人只注重十誡中「不可」「當」卻沒注意前面重要的話──「我是耶和華,曾經把你們從埃及為奴之地領出來.」(出廿)神先說:「我是」然後說當作甚麼.既是如此的大恩典.我們行事為人就應與蒙召之恩相稱.舊約有律法；但律法之前必有恩典；「我是耶和華」真是表明神的作為和恩惠!
\newline
\newline
「你不要害怕,因為我與你同在；不要驚惶,因為「我是」你的神；我必堅固你,我必幫助你,我必用公義的右手扶持你.」(10)「因為我耶和華你的神,必攙扶你的右手,對你說,不要害怕,我必幫助你.」(13)然後神說了些尖銳的對比:「你這蟲雅各......不要害怕；......」(14)「看哪,我已使你成為有快齒打糧的新器具......」(15)先把他們看如軟弱的小虫,後又看他們如同快齒的新器具.
\newline
\newline
我們一方面承認我們的軟弱,微小,膚淺,幼稚；但另方面要承認我們的偉大,有用,能作大事；並非我們能作甚麼,乃是神藉著我們作.耶穌說離了祂,甚麼都不能作.但保羅說,靠著那加給我們力量的,凡事都能作；保羅又告訴我們,我們有寶貝在瓦器裡面.這句話似有矛盾,寶貝是有價值的,怎能放在沒價值的容器裡呢?這就是神奇妙的恩典!我們多麼無用,幼稚；但我們多麼有用,多麼偉大；因主是偉大的；我們所得的救恩是偉大的,主的能力是偉大的；所以我們有了寶貝在瓦器內,就絕然不同了；不過,這裡有兩個可能:第一,瓦器太卑微,寶貝放在瓦器裡就貶值了；神的恩典偉大,到了我們身上就把神的恩典和能力限制了,因我們的卑微把主的恩典限制了,結果我們沒把福音帶出去；我們帶出去的,人家不接受；因我們已把主的生命貶了價值.但另一可能,因瓦器裡有寶貝,使瓦器聲價百倍.保羅說:「這莫大的能力不是出於我們,乃是出於神.」感謝主!我們有這恩典,有這能力；心中充滿快樂.
\newline
\newline
「我要在淨光的高處開江河,在谷中開泉源,我要使沙漠變為水池,使乾地變為湧泉.」(18)「......我要將一位報好信息的賜給耶路撒冷.(27)在本章我們看見神的能力真是偉大,祂是歷史的主,歷史每個細節都表達神救贖的能力.
\newline
\newline
「誰從東方興起一人......」(2)「我從北方興起一人......」(25)東方是波斯,當時興起波斯王古列,准許以色列人歸回祖國；並非偶然的,也不是外邦君王作的事；乃是出於神的旨意.神是掌管歷史的主,祂的旨意必然成就；甚至外邦也在神管治之下.神准許波斯國,從東方戰勝至北方的巴比倫.歷史的每個時間,都是經過神之許可,定有祂的旨意和目的,就是耶和華的名；天國的福音,要傳遍普天下.末世混亂的事,是反面的看法；其實,基督徒應視正面.神准許罪惡紛亂猖狂於世,為要達到救贖的目的；使福音傳遍普天下而後主耶穌再來.如果我們了解此點,我們看見時代的憂患,非但沒消極悲觀的感覺,反而樂觀並且積極；我們知道主快再來,我們有更重大的使命.福音的功效更大,能力更偉大；為此更應興奮快樂傳福音,所以我們應以正面態度勝過反面的態度.
\newline
\newline
在四十:9也有一處「看哪,你們的神.」四十二章開頭:「看哪,我的僕人......」這裡把我們的注意力集中在神的僕人；祂是報好信息的,是我們的神,我們的主.祂降世為人,是貧窮,卑微,受苦的僕人；祂為我們作成奇妙的事.主來到世上,非以役人,乃役於人；且捨命作多人的贖價.我們思想這位僕人,祂怎樣,我們也應該怎樣；在這憂患時代中,應如何承擔福音的使命.
\newline
\newline
誰是神的僕人?「惟你以色列我的僕,雅各我所揀選的......」(四十一:8)原來的僕人是以色列民,是整體的民族,或說民族的整體.神揀選以色列民,非因他們是優秀民族,有優秀文化；唯一的答案,是神的恩典.神所以揀選我們,唯有一解釋,也就是神的恩典.所以當我們來到主前,想到我們是何等樣的人竟蒙了神的恩典,我們就當感謝讚美.
\newline
\newline
神的揀選有其目的,為叫被揀選者服事祂.揀選是為事奉,為服事作神的僕人.照我們所了解,僕人是整個的以色列民；但不是每個以色列人都仰望神.從整體來說,以色列人失敗了,他們沒有好好的傳福音,所以神的審判臨到他們.因為審判要從神的家起首,神公平不偏待人；審判後,敬畏神的餘民；神留下保守看顧,幫助剛強,復興他們,為民族復興的核心.這班敬畏神的餘民僅屬少數.神讓他們從被擄之地歸回祖國,得以虔誠敬拜神；直等到主基督降世之時.主耶穌是敬虔的猶太人,是神獨一的兒子,救恩也從祂而來.耶穌說:「一粒麥子不落在地裡死了仍舊是一粒,若是死了就結出好多子粒來.」這意思教我們明白,僕人就是主耶穌,再進一步,從集體到個人；從個人到集體,你我都在此範圍內.因祂死了結出許多子粒來,把許多弟兄姊妹帶進榮耀裡.
\newline
\newline
主耶穌是神的僕人,是受苦的僕人,是彌賽亞,是基督,然後叫教會每個人都成為彌賽亞的團契.換言之,教會是僕人的團體,我們都是僕人.基督是唯一受苦的僕人；我們當效法祂作僕人,跟隨祂腳步,學習祂榜樣.若我們善作僕人,教會就興旺,就明白時代的憂患；把福音帶出去,因為僕人是福音的使者.求主幫助我們!
\newline
\newline
「看哪,我的僕人,我所扶持,所揀選,心裡所喜悅的,我已將我的靈賜給他......」扶持意即抓住.主耶穌知道自己是神的僕人,緊緊地抓住神.我們也要抓住祂,祂就扶持我們.我們也當服事祂,凡事不憑自己的意思；所有計劃,喜好,理想,要完全照主的心意服事祂.這樣,祂就扶持我們.
\newline
\newline
主耶穌是神所喜悅的,天上有聲音說:「這是我的愛子,我所喜悅的.」耶穌得父神的喜悅,唯一原則就是完全順服倚靠神.主從水裡上來時聖靈像鴿子降在祂身上；這是幅最完美的圖畫!鴿子與羔羊!主來到世上,約翰說:「看哪,神的羔羊.」羔羊是順服的,鴿子是聖靈的意思；鴿子降在羔羊身上,聖靈降在順服者身上.
\newline
\newline
我們要得勝靈的能力,要追求,禱告.有人以為要大聲禱告,聖靈才能降臨.是的,我相信聖靈願意照我們追求的心意,向我們施恩.但最要緊並非大聲禱告,也非熱鬧唱和；而是我們完全順服的態度.我們既蒙恩,就成為祂的僕人,如同主是祂的僕人一樣必須服事主人,僕人是主人所揀選的,必須完全順服；而後聖靈才臨到我們身上.但是不能因一時的充滿而滿足；有了聖靈的能力,就當把公理傳給外邦.公理的意思是「公義」和「審判」.神公義的審判要臨到世界,人被聖靈光照感動時；我們傳福音給他,他就為罪,為義,為審判,自己責備自己.這就是傳福音的內容.當耶穌來到世上,祂要把公理傳給外邦,工作之先,祂得著聖靈的能力；若沒有能力,工作就沒有果效.我們實在需要神的憐憫!
\newline
\newline
神的工作是安靜的工作「他不喧嚷,不揚聲,也不使街上聽見祂的聲音.」(2)聖工是內裡深入的工作.我們需要這樣的恩典!當神創造之時,不囂張,不喧嚷；當神拯救之時,也沒喧嚷之聲.很多奇妙的工作都是神在安靜中成就的.我們不是單看外表,看熱鬧,滿足人的好奇,而是讓聖靈作深入奇妙的工作.祂用的是極溫和,柔和,忍耐的方法,因為主知道自己所作的是甚麼.神在工作；我們若信心堅定,我們就柔和,不囂張,不誇耀；奇妙的事就發生了.
\newline
\newline
壓傷的蘆葦祂不折斷,將殘的燈火祂不吹滅.蘆葦非常脆弱,這描寫世人的情況；如四十章所說:「凡有血氣的,盡都如草.」人是脆弱的,容易折斷容易受傷；但是神有恩惠,這位僕人幫助壓傷的蘆葦沒折斷.我們傳福音也當如此,遍地人群如同壓傷的蘆葦；我們不能折斷.神差遣你和我；所有蒙恩者,應用親熱的手包裹破碎的心靈.燈火初點時火力薄弱滿了黑煙,容易熄滅；但是主的心意不讓它熄滅；這是僕人的工作.我們也當效法主的榜樣去幫助人,不讓生命的火花熄滅.有的人年青時很多理想,很大努力；但這是現實的,殘酷的,迅速失去了一切希望；好像將殘的燈火.今日世人因沒有得到福音的好處,沒有得到神的恩典,好像將殘的燈火.我們必須幫助他們；惟主的福音才能使他們希望的火花重新燃起；因福音是神的大能,要救一切相信的人.福音的能力實在太奇妙!我們當步主後塵,效法主在世上時所作；祂醫治病者,赦免罪人,安慰傷心者,扶持軟弱者；祂用此方法傳福音.很多人需要福音,我們能否安靜,溫柔地幫助他們?聖靈要在他們身上作奇妙工作；讓我們有敏銳的眼睛,時常看見人的需要,把福音傳給他們.壓傷的蘆葦,將殘的燈火；不單是對世人,對需要福音者說的；也是對福音使者說的.你,我是福音的使者,既已得了神的恩典；但多少時候我們軟弱,如將殘的燈火,壓傷的蘆葦.神用親熱的手幫助,更新,復興,醫治,加力量給我門,叫我們重新振作.軟弱變為剛強,為要叫我們真正成為福音的使者.人生在世,常經不起罪惡的引誘,受不了苦難的威脅和打擊而冷淡,我們多麼軟弱!求主幫助,叫我們重新掙扎,缺油燈火不摧殘.
\newline
\newline
我們需要聖靈的油,撒迦利亞書提到燈台兩邊有兩棵橄欖樹.橄欖有油,可以隨時豐足供應燈的需要.你我都需聖靈的供應,信心的火才能發旺,火能燒淨污穢.我們需要復興的火,聖靈的能力,屬靈的恩惠；我們在神前要徹底悔改,除去一切骯髒,讓火光明亮.有時我們不明白,為何神叫基督徒遭遇困難,經歷許多苦難甚至軟弱；但是,我們當知道,何時軟弱,何時就會剛強.因為神的能力,在軟弱者身上顯得完全.如果沒有疾病,怎見得神的醫治；如果沒有失敗,怎知道神的能力；沒有缺乏,怎知神的豐富呢?沒有憂傷,怎知神的安慰?我們在疾病,患難,病苦中才能體會神的恩典.不但夠我們用,且超過我們所求所想的.當我們軟弱時,到主面前,主能復興我們,叫我們重新得力,鼓起福音的勇氣作屬靈的爭戰,把福音帶出去.
\newline
\newline
神真正的僕人「他不灰心,也不喪膽......」(4)(原意不吹滅不折斷)人失敗,絕望；但神沒失敗,沒失望；神鼓勵,幫助,復興我們.叫我們不灰心,還要努力工作；神以無限的愛愛我們,祂的恩典一直感動我們,引導我們完全歸向祂.弟兄姊妹!要振作!不等待,不遲延；因為神有說不盡的恩賜和無限的心意.
\newline
\newline
第8節教我們明白,僕人不但有真正的能力,得到工作的方法,知道目的和使命.這僕人是眾民的中保(原文作約)神與以色列民立約,有新約,有舊約.約表明神的恩典一定實現,但我們當信從祂,把福音帶給外邦人.因為我們與神立約,有神恩惠的約；因著你我之故,神拯救他們.請勿忽略今世尊貴的福份!如果我們失敗了；他們就得不到拯救.我們倚靠神,我們是「眾民的中保,作外邦的光.」神給我們的路向很清楚,我們必須向前,不許退後.破釜沉舟前進,不但向前且要向上,如鷹展翅上騰.邁向萬邦,基督徒是教會福音工作的路向.華人教會差傳工作向外邦.「看哪,先前的事已經成就......」(9)神已經作成了救贖工作.「現在我將新事說明.」當拿但業見耶穌時,耶穌說:「你坐在無花果樹底下,我已經看見了你.」這話的含意說明了拿但業心裡的理想.舊約彌迦書說彌賽亞來的時候,地上沒有戰爭,只有和平,人都要坐在無花果樹和葡萄樹底下.換言之,耶穌知道拿但業的理想使他覺得驚奇,耶穌還對他說:「你要看見更大的事.」耶穌曾告訴門徒祂所作的,信祂的人也要作,並且要作更大的事.(約十四:12)這就是新事.神要我們作新事,看見更大的事,作更大的事；因神有更大的計劃,更大的能力,更大的福音奇妙工作.但我們必須把自己完全交給神,好像僕人一樣,在神前毫無保留,不猶豫,徹底順服,甘心樂意,把全人獻上祭壇給神.
\newline
\newline


\newpage

\section{第60屆~港九培靈會~唐佑之~第四講}
\label{sec:1223}
\textbf{裝備}
\newline
\newline
連結: \href{https://www.hkbibleconference.org/session-message/view/1223}{\texttt{https://www.hkbibleconference.org/session-message/view/1223}}
\newline
\newline
\hyperref[sec:1222]{< < < PREV SERMON < < <}
~
\hyperref[sec:toc]{[返目錄]}
~
\hyperref[sec:1224]{> > > NEXT SERMON > > >}
\newline
\newline
以賽亞書 44:1-49:26
\newline
\begin{longtable}{cl}
\hline
\hline
章節 & 經文 (和合本修訂版)\\
\hline
44:1 & \begin{tabularx}{0.7\textwidth}{X} 「我的僕人雅各，我所揀選的以色列啊，現在你當聽。 \end{tabularx} \\ \\ \relax
44:2 & \begin{tabularx}{0.7\textwidth}{X} 那位造你，使你在母腹中成形，並要幫助你的耶和華如此說：我的僕人雅各，我所揀選的耶書崙哪，不要害怕！ \end{tabularx} \\ \\ \relax
44:3 & \begin{tabularx}{0.7\textwidth}{X} 因為我要把水澆灌乾渴的地方，使水湧流在乾旱之地。我要將我的靈澆灌你的後裔，使我的福臨到你的子孫。 \end{tabularx} \\ \\ \relax
44:4 & \begin{tabularx}{0.7\textwidth}{X} 他們要在草叢中生長，如溪水旁的柳樹。 \end{tabularx} \\ \\ \relax
44:5 & \begin{tabularx}{0.7\textwidth}{X} 這個要說：『我是屬耶和華的』，那個要以雅各的名自稱，又有一個在手上寫著：『歸耶和華』，並自稱為以色列。」 \end{tabularx} \\ \\ \relax
44:6 & \begin{tabularx}{0.7\textwidth}{X} 耶和華－以色列的君王，以色列的救贖主－萬軍之耶和華如此說：「我是首先的，也是末後的；除我以外再沒有神。 \end{tabularx} \\ \\ \relax
44:7 & \begin{tabularx}{0.7\textwidth}{X} 自從古時我設立了人，誰能像我宣告，指明，又為自己陳說呢？讓他指明未來的事和必成的事吧！ \end{tabularx} \\ \\ \relax
44:8 & \begin{tabularx}{0.7\textwidth}{X} 你們不要恐懼，也不要害怕。我豈不是從上古就告訴並指示你們了嗎？你們是我的見證人！除我以外，豈有神呢？誠然沒有磐石，就我所知，一個也沒有！」 \end{tabularx} \\ \\ \relax
44:9 & \begin{tabularx}{0.7\textwidth}{X} 製造偶像的人盡都虛空，他們所喜悅的全無益處；偶像的見證人毫無所見，毫無所知，以致他們羞愧。 \end{tabularx} \\ \\ \relax
44:10 & \begin{tabularx}{0.7\textwidth}{X} 誰製造神像，鑄造偶像？這些都是無益的。 \end{tabularx} \\ \\ \relax
44:11 & \begin{tabularx}{0.7\textwidth}{X} 看哪，他的同夥都必羞愧。工匠不過是人，任他們聚集，任他們站立吧！他們都必懼怕，一同羞愧。 \end{tabularx} \\ \\ \relax
44:12 & \begin{tabularx}{0.7\textwidth}{X} 鐵匠用工具在火炭上工作，用鎚打出形狀，用他有力的膀臂來錘。他因飢餓而無力氣；因未喝水而疲倦。 \end{tabularx} \\ \\ \relax
44:13 & \begin{tabularx}{0.7\textwidth}{X} 木匠拉線，用筆劃出樣子，用鉋子鉋成形狀，又用圓規劃了模樣。他仿照人的體態，做出美妙的人形，放在廟裡。 \end{tabularx} \\ \\ \relax
44:14 & \begin{tabularx}{0.7\textwidth}{X} 他砍伐香柏樹，又取杉樹和橡樹，在樹林中讓它茁壯；或栽種松樹，得雨水滋潤長大。 \end{tabularx} \\ \\ \relax
44:15 & \begin{tabularx}{0.7\textwidth}{X} 這樹，人可用以生火；他拿一些來取暖，又搧火烤餅，而且做神像供跪拜，做雕刻的偶像向它叩拜。 \end{tabularx} \\ \\ \relax
44:16 & \begin{tabularx}{0.7\textwidth}{X} 他將一半的木頭燒在火中，用它烤肉來吃；吃飽了，就自己取暖說：「啊哈，我暖和了，我看到火了！」 \end{tabularx} \\ \\ \relax
44:17 & \begin{tabularx}{0.7\textwidth}{X} 然後又用剩下的一半做了一個神明，就是雕刻的偶像，向這偶像俯伏叩拜，向它禱告說：「求你拯救我，因你是我的神明。」 \end{tabularx} \\ \\ \relax
44:18 & \begin{tabularx}{0.7\textwidth}{X} 他們既無知，又不思想；因為耶和華蒙蔽他們的眼，使他們看不見，塞住他們的心，使他們不明白。 \end{tabularx} \\ \\ \relax
44:19 & \begin{tabularx}{0.7\textwidth}{X} 沒有一個心裡醒悟，有知識，有聰明，能說：「我曾拿一部分用火燃燒，在炭火上烤餅，也烤肉來吃。這剩下的，我豈要做可憎之像嗎？我豈可向木頭叩拜呢？」 \end{tabularx} \\ \\ \relax
44:20 & \begin{tabularx}{0.7\textwidth}{X} 他以灰塵為食，心裡迷糊，以致偏邪，不能自救，也不能說：「我右手中豈不是有虛謊嗎？」 \end{tabularx} \\ \\ \relax
44:21 & \begin{tabularx}{0.7\textwidth}{X} 雅各啊，要思念這些事；以色列啊，你是我的僕人。我造了你，你是我的僕人，以色列啊，我必不忘記你。 \end{tabularx} \\ \\ \relax
44:22 & \begin{tabularx}{0.7\textwidth}{X} 我塗去你的過犯，像厚雲消散；塗去你的罪惡，如薄霧消失。你當歸向我，因我救贖了你。 \end{tabularx} \\ \\ \relax
44:23 & \begin{tabularx}{0.7\textwidth}{X} 諸天哪，應當歌唱，因為耶和華成就這事。地的深處啊，應當歡呼；眾山哪，要出聲歌唱；樹林和其中所有的樹木啊，你們都當歌唱！因為耶和華救贖了雅各，並要因以色列榮耀自己。 \end{tabularx} \\ \\ \relax
44:24 & \begin{tabularx}{0.7\textwidth}{X} 從你在母腹中就造了你，你的救贖主－耶和華如此說：「我－耶和華創造萬物，獨自鋪張諸天，親自展開大地； \end{tabularx} \\ \\ \relax
44:25 & \begin{tabularx}{0.7\textwidth}{X} 我使虛謊的預兆失效，愚弄占卜的人，使智慧人退後，使他的知識變為愚拙； \end{tabularx} \\ \\ \relax
44:26 & \begin{tabularx}{0.7\textwidth}{X} 卻使我僕人的話站得住，成就我使者的籌算。我論耶路撒冷說：『必有人居住』；論猶大的城鎮說：『必被建造，我必重建其中的廢墟。』 \end{tabularx} \\ \\ \relax
44:27 & \begin{tabularx}{0.7\textwidth}{X} 我對深淵說：『乾了吧！我要使你的江河乾涸』； \end{tabularx} \\ \\ \relax
44:28 & \begin{tabularx}{0.7\textwidth}{X} 論居魯士說：『他是我的牧人，他要成就我所喜悅的，下令建造耶路撒冷，發命令立穩聖殿的根基。』」 \end{tabularx} \\ \\ \relax
45:1 & \begin{tabularx}{0.7\textwidth}{X} 耶和華對所膏的居魯士如此說，他的右手我曾攙扶，使列國降服在他面前，列王的腰帶我曾鬆開，使城門在他面前敞開，不得關閉： \end{tabularx} \\ \\ \relax
45:2 & \begin{tabularx}{0.7\textwidth}{X} 「我要在你前面行，修平崎嶇之地。我必打破銅門，砍斷鐵閂。 \end{tabularx} \\ \\ \relax
45:3 & \begin{tabularx}{0.7\textwidth}{X} 我要將暗中的寶物和隱藏的財富賜給你，使你知道提名召你的就是我－耶和華，以色列的神。 \end{tabularx} \\ \\ \relax
45:4 & \begin{tabularx}{0.7\textwidth}{X} 因我的僕人雅各，我所揀選的以色列，我提名召你；你雖不認識我，我也加給你名號。 \end{tabularx} \\ \\ \relax
45:5 & \begin{tabularx}{0.7\textwidth}{X} 我是耶和華，再沒有別的了；除了我以外再沒有神。你雖不認識我，我必給你束腰。 \end{tabularx} \\ \\ \relax
45:6 & \begin{tabularx}{0.7\textwidth}{X} 從日出之地到日落之處使人都知道除我以外，沒有別的。我是耶和華，再沒有別的了。 \end{tabularx} \\ \\ \relax
45:7 & \begin{tabularx}{0.7\textwidth}{X} 我造光，又造暗；施平安，又降災禍；做成這一切的是我－耶和華。 \end{tabularx} \\ \\ \relax
45:8 & \begin{tabularx}{0.7\textwidth}{X} 「諸天哪，要如雨傾盆而降，雲要降下公義，地要裂開，救恩湧出，使公義也一同滋長；這都是我－耶和華造的。」 \end{tabularx} \\ \\ \relax
45:9 & \begin{tabularx}{0.7\textwidth}{X} 「那與造他的主爭論的人有禍了！他不過是地上瓦塊中的一片。泥土豈可對塑造它的說：『你做的是甚麼？你所做的物怎麼沒有把手呢？』 \end{tabularx} \\ \\ \relax
45:10 & \begin{tabularx}{0.7\textwidth}{X} 有人對父親說，『你生的是甚麼』，對母親說，『你生產的是甚麼』；這人有禍了！」 \end{tabularx} \\ \\ \relax
45:11 & \begin{tabularx}{0.7\textwidth}{X} 耶和華－以色列的聖者，就是造以色列的如此說：「難道我孩子的未來，你們能質問我，我手的工作，你們可以吩咐我嗎？ \end{tabularx} \\ \\ \relax
45:12 & \begin{tabularx}{0.7\textwidth}{X} 我造大地，又創造人在地上。我親手鋪張諸天，天上萬象也是我所任命的。 \end{tabularx} \\ \\ \relax
45:13 & \begin{tabularx}{0.7\textwidth}{X} 我憑公義興起居魯士，又要修直他一切的道路。他必建造我的城，釋放我被擄的民，不為工價，也不為獎賞。」這是萬軍之耶和華說的。 \end{tabularx} \\ \\ \relax
45:14 & \begin{tabularx}{0.7\textwidth}{X} 耶和華如此說：「埃及的出產和古實的貨物必歸你；身量高大的西巴人，他們必過來歸你，為你所有。他們必帶著鎖鏈過來跟隨你，向你下拜，祈求你說：『神真是在你中間，再沒有別的，沒有別的神。』」 \end{tabularx} \\ \\ \relax
45:15 & \begin{tabularx}{0.7\textwidth}{X} 救主－以色列的神啊，你誠然是隱藏自己的神。 \end{tabularx} \\ \\ \relax
45:16 & \begin{tabularx}{0.7\textwidth}{X} 製造偶像的都要抱愧蒙羞，他們要一同歸於慚愧。 \end{tabularx} \\ \\ \relax
45:17 & \begin{tabularx}{0.7\textwidth}{X} 惟有以色列必蒙耶和華拯救，得永遠的救恩。你們必不蒙羞，也不抱愧，直到永世無盡。 \end{tabularx} \\ \\ \relax
45:18 & \begin{tabularx}{0.7\textwidth}{X} 耶和華如此說，他創造諸天，他是神；他造了地，形成它，堅固它，並非創造它為荒涼，而是要給人居住：「我是耶和華，再沒有別的。 \end{tabularx} \\ \\ \relax
45:19 & \begin{tabularx}{0.7\textwidth}{X} 我不在隱密黑暗之地說話，也沒有對雅各的後裔說，『你們尋求我是徒然的』，我－耶和華所講的是公義，所說的是正直。」 \end{tabularx} \\ \\ \relax
45:20 & \begin{tabularx}{0.7\textwidth}{X} 「你們從列國逃脫的人，要一同聚集前來。那些抬著雕刻的木偶、祈求不能救人之神明的，毫無知識。 \end{tabularx} \\ \\ \relax
45:21 & \begin{tabularx}{0.7\textwidth}{X} 你們要近前來說明，讓他們彼此商議。誰從古時指明這事？誰從上古述說它？不是我－耶和華嗎？除了我以外，再沒有神；我是公義的神，又是救主；除了我以外，再沒有別的了。 \end{tabularx} \\ \\ \relax
45:22 & \begin{tabularx}{0.7\textwidth}{X} 「地的四極都當轉向我，就必得救；因為我是神，再沒有別的。 \end{tabularx} \\ \\ \relax
45:23 & \begin{tabularx}{0.7\textwidth}{X} 我指著自己起誓，公義從我的口發出，這話並不返回：『萬膝必向我跪拜，萬口必憑我起誓。』 \end{tabularx} \\ \\ \relax
45:24 & \begin{tabularx}{0.7\textwidth}{X} 「人論我說，『公義、能力，惟獨在乎耶和華。人必歸向他，凡向他發怒的都必蒙羞。 \end{tabularx} \\ \\ \relax
45:25 & \begin{tabularx}{0.7\textwidth}{X} 以色列的後裔必因耶和華得稱為義，並要彼此誇耀。』」 \end{tabularx} \\ \\ \relax
46:1 & \begin{tabularx}{0.7\textwidth}{X} 彼勒叩拜，尼波屈身；巴比倫的偶像馱在走獸和牲畜背上。你們所抬的成了重馱，使牲畜疲乏。 \end{tabularx} \\ \\ \relax
46:2 & \begin{tabularx}{0.7\textwidth}{X} 這些神明一同屈身叩拜，不能救自己，反倒遭人擄去。 \end{tabularx} \\ \\ \relax
46:3 & \begin{tabularx}{0.7\textwidth}{X} 雅各家，以色列家所有的餘民哪，你們自從生下就蒙我抱，自出母胎便由我來背，你們都要聽從我。 \end{tabularx} \\ \\ \relax
46:4 & \begin{tabularx}{0.7\textwidth}{X} 直到你們年老，我不改變；直到你們髮白，我仍扶持。我已造你，就必背你；我必抱你，也必拯救。 \end{tabularx} \\ \\ \relax
46:5 & \begin{tabularx}{0.7\textwidth}{X} 你們將誰與我相比，與我相等，將誰與我相較，使我們相似呢？ \end{tabularx} \\ \\ \relax
46:6 & \begin{tabularx}{0.7\textwidth}{X} 他們從錢囊中倒出金子，用天平秤出銀子，雇銀匠造成神像，他們又俯伏，又叩拜。 \end{tabularx} \\ \\ \relax
46:7 & \begin{tabularx}{0.7\textwidth}{X} 他們抬起神像，扛在肩上，安置在定處，使它站立，不離本位；人呼求它，它卻不回答，也無法救人脫離災難。 \end{tabularx} \\ \\ \relax
46:8 & \begin{tabularx}{0.7\textwidth}{X} 你們當記得這事，立定心意。叛逆的人哪，要留心思想。 \end{tabularx} \\ \\ \relax
46:9 & \begin{tabularx}{0.7\textwidth}{X} 要追念上古的事，因為我是神，並無別的；我是神，沒有能與我相比的。 \end{tabularx} \\ \\ \relax
46:10 & \begin{tabularx}{0.7\textwidth}{X} 我從起初就指明末後的事，從古時便言明未成的事，說：「我的籌算必立定；凡我所喜悅的，我必成就。」 \end{tabularx} \\ \\ \relax
46:11 & \begin{tabularx}{0.7\textwidth}{X} 我召鷙鳥從東方來，召那成就我籌算的人從遠方來。我已說出，就必成就；我已謀定，也必做成。 \end{tabularx} \\ \\ \relax
46:12 & \begin{tabularx}{0.7\textwidth}{X} 你們這些心中頑固、遠離公義的人，要聽從我。 \end{tabularx} \\ \\ \relax
46:13 & \begin{tabularx}{0.7\textwidth}{X} 我使我的公義臨近，它已不遠。我的救恩必不遲延。我要為以色列－我的榮耀在錫安施行救恩。 \end{tabularx} \\ \\ \relax
47:1 & \begin{tabularx}{0.7\textwidth}{X} 少女巴比倫 哪，下來坐在塵埃；迦勒底啊，沒有寶座，要坐在地上；你不再稱為柔弱嬌嫩。 \end{tabularx} \\ \\ \relax
47:2 & \begin{tabularx}{0.7\textwidth}{X} 要用磨磨麵，揭去面紗，脫去長裙，露腿渡河。 \end{tabularx} \\ \\ \relax
47:3 & \begin{tabularx}{0.7\textwidth}{X} 你的下體必被露出；你的羞辱必被看見。我要報復，誰也不寬容。 \end{tabularx} \\ \\ \relax
47:4 & \begin{tabularx}{0.7\textwidth}{X} 我們的救贖主是以色列的聖者，他的名為萬軍之耶和華。 \end{tabularx} \\ \\ \relax
47:5 & \begin{tabularx}{0.7\textwidth}{X} 迦勒底啊，你要靜坐，進入黑暗中，因你不再稱為萬國之后。 \end{tabularx} \\ \\ \relax
47:6 & \begin{tabularx}{0.7\textwidth}{X} 我向我的百姓發怒，使我的產業受凌辱，將他們交在你手中；然而你毫不憐憫他們，連老年人你也加極重的軛。 \end{tabularx} \\ \\ \relax
47:7 & \begin{tabularx}{0.7\textwidth}{X} 你說：「我必永遠為后。」你不將這事放在心上，也不思想事情的結局。 \end{tabularx} \\ \\ \relax
47:8 & \begin{tabularx}{0.7\textwidth}{X} 你這專好宴樂、以為地位穩固的，現在當聽這話。你心中說：「惟有我，除我以外再沒有別的。我必不致寡居，也不經歷喪子之痛。」 \end{tabularx} \\ \\ \relax
47:9 & \begin{tabularx}{0.7\textwidth}{X} 哪知，喪子、寡居這兩件事一日之間忽然臨到你；你雖多行邪術、廣施魔咒，這兩件事必全然臨到你身上。 \end{tabularx} \\ \\ \relax
47:10 & \begin{tabularx}{0.7\textwidth}{X} 你倚靠自己的惡行，說：「無人看見我。」你的智慧聰明使你走偏，你心裡說：「惟有我，除我以外再沒有別的了。」 \end{tabularx} \\ \\ \relax
47:11 & \begin{tabularx}{0.7\textwidth}{X} 但災禍臨到你，你不知如何驅除；災害落在你身上，你也無法除掉，你所不知道的毀滅必忽然臨到你身上。 \end{tabularx} \\ \\ \relax
47:12 & \begin{tabularx}{0.7\textwidth}{X} 儘管使用從幼年就施行的魔符和眾多的邪術吧！或許有些幫助，或許可以致勝。 \end{tabularx} \\ \\ \relax
47:13 & \begin{tabularx}{0.7\textwidth}{X} 你籌劃太多，以致疲倦。讓那些觀天象，看星宿，在初一說預言的都起來，救你脫離所要臨到你的事！ \end{tabularx} \\ \\ \relax
47:14 & \begin{tabularx}{0.7\textwidth}{X} 看哪，他們要像碎秸被火焚燒，無法救自己脫離火焰的魔掌；沒有炭火可以取暖，你也不能坐在火旁。 \end{tabularx} \\ \\ \relax
47:15 & \begin{tabularx}{0.7\textwidth}{X} 你所操勞的事都像這樣；從你幼年以來與你交易的都各奔己路，沒有一人來救你。 \end{tabularx} \\ \\ \relax
48:1 & \begin{tabularx}{0.7\textwidth}{X} 雅各家，稱為以色列名下，從猶大的源頭而出的啊，你們指著耶和華的名起誓，提說以色列的神，卻不憑誠信，也不憑公義； \end{tabularx} \\ \\ \relax
48:2 & \begin{tabularx}{0.7\textwidth}{X} 你們自稱為聖城之民，倚靠名為萬軍之耶和華－以色列的神；現在，當聽我言： \end{tabularx} \\ \\ \relax
48:3 & \begin{tabularx}{0.7\textwidth}{X} 「先前的事，我自古已說明，已從我口而出，是我所指示的；我瞬間行事，事便成就。 \end{tabularx} \\ \\ \relax
48:4 & \begin{tabularx}{0.7\textwidth}{X} 因為我知道你是頑梗的；你的頸項是鐵的，你的額頭是銅的。 \end{tabularx} \\ \\ \relax
48:5 & \begin{tabularx}{0.7\textwidth}{X} 所以，我自古就給你說明，在事未成以先指示你，免得你說：『這些事是我的偶像所行的，是我雕刻的偶像和鑄造的神像所命定的。』 \end{tabularx} \\ \\ \relax
48:6 & \begin{tabularx}{0.7\textwidth}{X} 「你既已聽見，現在要察看這一切；你們不是要說明嗎？從今以後，我要指示你新事，就是你所不知道的隱密事。 \end{tabularx} \\ \\ \relax
48:7 & \begin{tabularx}{0.7\textwidth}{X} 這事是現今造的，並非自古就有，在今日以先，你未曾聽見；免得你說：『看哪，這事我早已知道了。』 \end{tabularx} \\ \\ \relax
48:8 & \begin{tabularx}{0.7\textwidth}{X} 誠然你未曾聽見，也未曾知道；你的耳朵從來未曾開通。我原知道你行事極其詭詐，你自從出母胎以來，就稱為悖逆的。 \end{tabularx} \\ \\ \relax
48:9 & \begin{tabularx}{0.7\textwidth}{X} 「我為我的名暫且忍怒，為了我的榮耀向你容忍，不將你剪除。 \end{tabularx} \\ \\ \relax
48:10 & \begin{tabularx}{0.7\textwidth}{X} 看哪，我熬煉你，卻不像熬煉銀子；你在苦難的火爐中，我試煉你。 \end{tabularx} \\ \\ \relax
48:11 & \begin{tabularx}{0.7\textwidth}{X} 我為自己的緣故必做這事，我豈能被褻瀆？我必不將我的榮耀歸給別神。 \end{tabularx} \\ \\ \relax
48:12 & \begin{tabularx}{0.7\textwidth}{X} 雅各－我所選召的以色列啊，當聽從我：我是耶和華，我是首先的，也是末後的。 \end{tabularx} \\ \\ \relax
48:13 & \begin{tabularx}{0.7\textwidth}{X} 我親手立了地的根基，以右手鋪張諸天；我一召喚，天地就都立定。」 \end{tabularx} \\ \\ \relax
48:14 & \begin{tabularx}{0.7\textwidth}{X} 「你們都當聚集而聽，偶像之中誰曾說明這些事？耶和華愛他，他必向巴比倫成就耶和華的旨意，耶和華的膀臂也要加在迦勒底人身上。 \end{tabularx} \\ \\ \relax
48:15 & \begin{tabularx}{0.7\textwidth}{X} 我，惟有我曾說過，我選召他，領他來，他的道路必亨通。 \end{tabularx} \\ \\ \relax
48:16 & \begin{tabularx}{0.7\textwidth}{X} 你們要接近我來聽這話，我從起初就未曾在隱密之處說話，萬事之始，我就在那裡。」現在，主耶和華差遣了我，帶著他的靈而來。 \end{tabularx} \\ \\ \relax
48:17 & \begin{tabularx}{0.7\textwidth}{X} 耶和華－你的救贖主，以色列的聖者如此說：「我是耶和華－你的神，我教導你，使你得益處，指引你當走的路。 \end{tabularx} \\ \\ \relax
48:18 & \begin{tabularx}{0.7\textwidth}{X} 甚願你聽從我的命令，你的平安就會如河水，你的公義如海浪， \end{tabularx} \\ \\ \relax
48:19 & \begin{tabularx}{0.7\textwidth}{X} 你的後裔必多如海沙，你腹中所生的必多如沙粒。他的名絕不從我面前剪除，也不滅絕。」 \end{tabularx} \\ \\ \relax
48:20 & \begin{tabularx}{0.7\textwidth}{X} 你們要從巴比倫出來，從迦勒底人中逃脫，以歡呼的聲音宣告，將這事傳揚到地極，說：耶和華救贖了他的僕人雅各！ \end{tabularx} \\ \\ \relax
48:21 & \begin{tabularx}{0.7\textwidth}{X} 他引導他們經過沙漠，他們卻未嘗乾渴；他為他們使水從磐石流出，磐石裂開，水就湧出。 \end{tabularx} \\ \\ \relax
48:22 & \begin{tabularx}{0.7\textwidth}{X} 耶和華說：「惡人必不得平安！」 \end{tabularx} \\ \\ \relax
49:1 & \begin{tabularx}{0.7\textwidth}{X} 眾海島啊，當聽從我！遠方的眾民哪，要留心聽！自出母胎，耶和華就選召我；自出母腹，他就稱呼我的名。 \end{tabularx} \\ \\ \relax
49:2 & \begin{tabularx}{0.7\textwidth}{X} 他使我的口如快刀，把我藏在他手蔭之下；又使我成為磨利的箭，把我藏在他箭袋之中； \end{tabularx} \\ \\ \relax
49:3 & \begin{tabularx}{0.7\textwidth}{X} 對我說：「你是我的僕人以色列；我必因你得榮耀。」 \end{tabularx} \\ \\ \relax
49:4 & \begin{tabularx}{0.7\textwidth}{X} 我卻說：「我勞碌是徒然，我盡力是虛無虛空。耶和華誠然以公平待我，我的賞賜在我的神那裡。」 \end{tabularx} \\ \\ \relax
49:5 & \begin{tabularx}{0.7\textwidth}{X} 現在耶和華說話，他從我出母胎，就造我作他的僕人，要使雅各歸向他，使以色列聚集在他那裡。耶和華看我為尊貴，我的神是我的力量。 \end{tabularx} \\ \\ \relax
49:6 & \begin{tabularx}{0.7\textwidth}{X} 他說：「你作我的僕人，使雅各眾支派復興，使以色列中蒙保存的人歸回；然而此事尚小，我還要使你作萬邦之光，使你施行我的救恩，直到地極。」 \end{tabularx} \\ \\ \relax
49:7 & \begin{tabularx}{0.7\textwidth}{X} 救贖主－以色列的聖者耶和華對那被人藐視、本國憎惡、統治者奴役的如此說：「君王看見就站起來，領袖也要下拜；這都是因信實的耶和華，因揀選你的以色列的聖者。」 \end{tabularx} \\ \\ \relax
49:8 & \begin{tabularx}{0.7\textwidth}{X} 耶和華如此說：「在悅納的時候，我應允了你；在拯救的日子，我幫助了你。我要保護你，要藉著你與百姓立約，為了復興遍地，使人承受荒蕪之地為業； \end{tabularx} \\ \\ \relax
49:9 & \begin{tabularx}{0.7\textwidth}{X} 對那被捆綁的人說：『出來吧！』對在黑暗裡的人說：『顯現吧！』他們在路上必得飲食，在光禿的高地必有食物。 \end{tabularx} \\ \\ \relax
49:10 & \begin{tabularx}{0.7\textwidth}{X} 他們不飢不渴，炎熱和烈日必不傷害他們；因為憐憫他們的必引導他們，領他們到水泉旁邊。 \end{tabularx} \\ \\ \relax
49:11 & \begin{tabularx}{0.7\textwidth}{X} 我必在眾山開闢路徑，大道也要填高。 \end{tabularx} \\ \\ \relax
49:12 & \begin{tabularx}{0.7\textwidth}{X} 看哪，他們從遠方來；有些從北方來，有些從西方來，有些從色弗尼地來。」 \end{tabularx} \\ \\ \relax
49:13 & \begin{tabularx}{0.7\textwidth}{X} 諸天哪，應當歡呼！大地啊，應當快樂！眾山哪，應當揚聲歌唱！因為耶和華已經安慰他的百姓，他要憐憫他的困苦之民。 \end{tabularx} \\ \\ \relax
49:14 & \begin{tabularx}{0.7\textwidth}{X} 錫安說：「耶和華離棄了我，主忘記了我。」 \end{tabularx} \\ \\ \relax
49:15 & \begin{tabularx}{0.7\textwidth}{X} 婦人焉能忘記她吃奶的嬰孩，不憐憫她所生的兒子？即或有忘記的，我卻不忘記你。 \end{tabularx} \\ \\ \relax
49:16 & \begin{tabularx}{0.7\textwidth}{X} 看哪，我將你銘刻在我掌上，你的城牆常在我眼前。 \end{tabularx} \\ \\ \relax
49:17 & \begin{tabularx}{0.7\textwidth}{X} 建立你的勝過毀壞你的，使你荒廢的必都離你而去。 \end{tabularx} \\ \\ \relax
49:18 & \begin{tabularx}{0.7\textwidth}{X} 你舉目向四圍觀看，他們都聚集來到你這裡。我指著我的永生起誓，你定要以他們為妝飾佩戴，帶著他們，像新娘一樣。這是耶和華說的。 \end{tabularx} \\ \\ \relax
49:19 & \begin{tabularx}{0.7\textwidth}{X} 至於你荒廢淒涼之處，並你被毀壞之地，如今居民必嫌太窄，吞滅你的必離你遙遠。 \end{tabularx} \\ \\ \relax
49:20 & \begin{tabularx}{0.7\textwidth}{X} 你要再聽見喪失子女後所生的兒女說：「這地方我居住太窄，請你給我地方居住。」 \end{tabularx} \\ \\ \relax
49:21 & \begin{tabularx}{0.7\textwidth}{X} 那時你心裡必說：「我既喪子不育，被擄，飄流在外，誰給我生了這些？誰將他們養大呢？看哪，我被撇下獨自一人時，他們都在哪裡呢？」 \end{tabularx} \\ \\ \relax
49:22 & \begin{tabularx}{0.7\textwidth}{X} 主耶和華如此說：「看哪，我必向列國舉手，向萬民豎立大旗；他們必將你的兒子抱在懷中帶來，將你的女兒背在肩上扛來。 \end{tabularx} \\ \\ \relax
49:23 & \begin{tabularx}{0.7\textwidth}{X} 列王必作你的養父，王后必作你的乳母。他們必以臉伏地，向你下拜，並舔你腳上的塵土。你就知道我是耶和華，等候我的必不致羞愧。」 \end{tabularx} \\ \\ \relax
49:24 & \begin{tabularx}{0.7\textwidth}{X} 勇士搶去的豈能奪回？被殘暴者擄掠的豈能得解救呢？ \end{tabularx} \\ \\ \relax
49:25 & \begin{tabularx}{0.7\textwidth}{X} 但耶和華如此說：「就是勇士所擄掠的，也可以奪回；殘暴者所搶的，也可以得解救。與你相爭的，我必與他相爭，我也要拯救你的兒女。 \end{tabularx} \\ \\ \relax
49:26 & \begin{tabularx}{0.7\textwidth}{X} 我必使那欺壓你的吃自己的肉，飲自己的血，如喝甜酒喝醉一樣。凡有血肉之軀的都必知道我－耶和華是你的救主，是你的救贖主，是雅各的大能者。」 \end{tabularx} \\ \\
[1ex]
\hline
\hline
\end{longtable}
我們在聖經中清楚看到時代憂患與福音使命的歷史背景,也看到神這樣的心意.耶和華是憂患的神,祂常為人的罪惡而痛苦；看見人的滅亡而焦急.我們的主是憂患之主,祂多受痛苦常經憂患；所以福音是付上極大的代價才成就的.
\newline
\newline
從教會歷史看,當人感到時代憂患之時,教會才能興旺,才能感到福音使命的重大.感謝主!我們生在憂患的時代!有的人生於安樂而長於憂患；有的人生於憂患而長於憂患.為何有諸多憂患?因為神要恩待我們,教我們能得福音的好處,體會人心對福音的需要；於是教會建立,當教會注意福音的使命並且實行福音說憂患本身是種裝備；沒有憂患的裝備,我們不知福音的重要；所以我們當感謝讚美神的恩待.有次我讀了篇禱告文,有一教會遭大逼迫經歷很多困難,弟兄姊妹卻熱心愛主.雖然那時環境受限制,但福音沒有受限制,神的恩典沒有受限制；而且更奇妙地叫多人蒙恩得救.事後憂患過去了,大家得以自由敬拜,也充分活動；但是教會火熱的狀況也不再有了.這篇禱告文:「求主給我們更多的憂患.」我們並非要求苦難,但神藉著歷史的環境,叫我們能看見教會時代的使命.
\newline
\newline
本書四十九章可說是第二首受苦僕人之詩.四十八章起題及三方面的裝備:
\newline
\newline
(一),思想,觀念,知識的裝備.
\newline
\newline
(二),屬靈的裝備,讓聖靈來裝備我們,
\newline
\newline
(三),行動的裝備,以行動表現,邊作邊學習；邊行動邊準備；因神要托付我們更大的事,更奇妙的工作,更重大的使命.這些都是福音的工作和使命.求主憐憫幫助!讓我們不單明白其中的意思,信息；且對主說:「主啊!求你裝備我,叫我作好準備作你的工.」
\newline
\newline
第四十八:3,6-8,神特別向以色列人說話；因他們是神所揀選之僕,要照神的心意把福音帶出去,他們需要裝備；否則他們就成了悖逆的人,以色列人從他們的歷史經驗中,知道他們一直蒙神賜福而感謝神；但漸漸就忽略了神的恩典；所以神不斷激勵他們.
\newline
\newline
勿忘神在我們身上所成就的恩典!先前的事很重要,必須記得；否則損失極大.詩一O三:3說:「不可忘記他們的一切恩惠.」保羅說:「要忘記背後,努力前面.」(腓三:13)是的,有的事我們應當忘記；但神的恩典切不可忘記.然而,我們最大的失敗是牢記應該忘記的；忘記應該牢記的.要記得神的恩典有歷史的賡續性,神是首先也是末後的；神開始,繼續,完成.保羅說:「因為萬有都是本於他,倚靠他,歸於他.......」(羅十一:36)這裡我們看到屬靈的循環.
\newline
\newline
神藉以賽亞特別提醒以色列民:「主說,早先的事我從古時說明.......」我覺得今天社會文化有很大的危機；青年人常不願意記念從前的事,以為已往的已不合時宜.我認為應當注意傳統,聖經就是傳統,我們怎能否認傳統呢?先前之事當記得；但另一方面,我們不能單留戀過去的事.換言之,我們必須更新；若專講傳統,使我們固步自封裹足不前；教會沒有進步是因墨守成規.好處應當保留,不好的應當改變應當更新；我們應注意傳統也應有突破.只有突破沒有傳統,那未免太膚淺；(一般青年人如此)只有傳統沒有更新教會就沒有進步,(這是中,老年人的通病)教會復興,一方面要尊重傳統,一方面要注意更新.
\newline
\newline
教會的進步須在觀念上,思想上,知識上裝備,這還不夠；還須在隱密處等候神的恩典.到神前親近,得著祂的話語,神要把祂的靈來感動我們(16).我們必須在屬靈上真正裝備,才有行動上的裝備.請看17-19,我們在思想上有了準備,在屬靈恩典上有了成就,在主的恩典和知識上有了長進；我們必須行動,因為神已為我們預備了道路.耶穌釘十字架,「藉著祂給我們開了一條又新又活的路.」叫我們知道當走的路.約書亞帶領以色列人過約但河時,教導他們要跟隨約櫃走.因為約櫃象徵神的同在和引導；約櫃是神的話,有神的律法.
\newline
\newline
要往前走!走又新又活的路,我們需要靠主的恩典；這條路是平安公義的路.平安並非無危險困難和阻礙.平安意即秩序與和諧,神創造天地井然有條,把光和暗,天和地,海洋和陸地,生物和非生物,植物和動物分開；動物中又把高等和低等分開.高等動物之中,把人和其他動物分開；不但做分開工作,也做聚攏的工作各從其類,建立完美的秩序,所以有和諧.神的救贖也是如此,祂做分開的工作,把我們從罪惡黑暗中分別出來,使我們成為光明白晝之子,成為聖潔的義人.神實在恩待我們!把我們分別出來成為屬祂的人.祂也作聚攏的工作各從其類；這就是教會的真理,教會就是把我們分別出來再聚攏一起,這就是秩序,有了秩序就有和諧,沒有和諧,就沒有平安.神拯救了我們,我們的內心由紛亂而至有秩序,和諧,平安,快樂.平安就是完整的意思；不是片段零碎,而是完全完美沒缺陷的.主如此恩待我們；我們應當走祂的路,遵祂命令,完全仰望祂；讓祂領導,讓我們知道當走的路；以至我們的平安好像河水一樣.有首詩歌:「平安好像江河」作曲者遭遇家庭變故,他仰望神；神賜他平安,讓他真正與神有極其密切的關係.當我們有此裝備,福音見證就有能力；很多人容易接受,教會成為福音見證之所.
\newline
\newline
我們都應該作神的僕人以色列民,彌賽亞的國民,主耶穌,教會的僕人；凡蒙恩得救的人都是神的僕人.茲將這些方面連繫思想,從以色列民族的整體,提及彌賽亞個人,再提到教會的整體；你我教會的整體.
\newline
\newline
海島是遙遠之地要宣告福音的能力,傳福音的使者,必須向普天下傳福音.「眾海島阿,當聽我言；遠方的眾民哪,留心而聽；自我出胎,耶和華就選召我,自出母腹,他就提我的名.」賽四十九:1這叫我們不能不想到耶利米書一章,當神選召耶利米,對他說:「我未將你造在腹中,我已曉得你,你未出母胎,我已分別你為聖,我已派你作列國的先知.(5)各位!你我在神的計劃和心意中,早就要你我成為祂的僕人；用祂的恩典用祂的愛揀選我們,叫我們成為福音的使者,不但自己得福音的好處,且要把福音帶出去,讓多人得著.當耶利米聽見耶和華所說的話,非常害怕,「就說,主耶和華阿,我不知怎樣說,因為我是年幼的.」年幼這字,在希伯來文中是2-40歲.這裡的年幼非指年齡,乃指經驗.耶利米自覺經驗不足；但是神特別恩待他,「於是耶和華伸手按我的口說,我已將當說的話傳給你.」神要耶利米作先知.先知不是單講將來的事,也著重當代的信息,兩方面並重.當時,神給耶利米看見兩個異象,第一,他看見一根杏樹枝.在那非常憂患時代,以色列民已到滅亡邊緣；如同嚴冬一切都枯萎了.杏樹枝宣告春天來臨.「耶和華說,你看得不錯；因為我留意保守我的話,使得成就.」神的保守,看顧,施恩,注意,使耶利米得到很大的鼓勵,滿有希望.福音是宣告基督的恩惠,讓我們得著祂的生命.我們須有耶利米的經驗.我們每個人都在神的計劃中,未出母胎,神已規定了.神呼召耶利米；神讓祂的兒子到世上來；也叫你我成為祂的僕人,作福音的使者.
\newline
\newline
耶利米看見的第二個異象:「一個燒開的鍋從北而傾,」災害來到了!我們傳福音也應有此認識；人若不接受福音,是多麼可怕危險的事!為甚麼我們要傳福音呢?因為人為禍了,快滅亡了!我們在神面前,須有屬靈的眼光,看到人的悲哀,危險；世界快滅亡了,教會若不及時傳福音,就會滅亡.
\newline
\newline
「他使我的口如快刀......」(2)口很重要；但我們常用來說閒話,理應用來傳福音才對.英國有個文學家說,全世界有兩種人,一種是滔滔不絕,專講廢話的人；一種是有話講,有信息傳,有見證作；但他們卻閉口不言.這是否指基督徒呢?我們多麼需要神開我們的口傳福音,讓每個人的口都成為先知,讓教會成為先知的團體.口如快刀意即「攻」不是「守」,要為福音爭戰.我們在教會常守不攻,不敢講話,不敢表達立場,我們的見證在哪裡?宗教改革時,馬丁路德的口號「教會是爭戰的教會.」我們須爭戰,要伸出刀來；但刀不鋒利,我們應要操練.所以神「將我藏在他手蔭之下」保護,造就,在隱密處得神所賜豐盛奇妙的恩典.這可說是希伯來書的內容；先在幔子內然後到營外,在幔子裡在隱密處與神親密交往.有高深屬靈的經驗,有聖靈的能力佔有,潔淨,感動,充滿我們；然後出幔子,到營外去為王工作.「又使我成為磨亮的箭,」射箭要有目的,不中軸就無用.新約聖經特別提到罪,罪即不中軸.今天教會任意放箭卻未達目的；最悲慘的就像何西阿書所載,以色列人如同反背的弓射向自己.教會時有此現象,弟兄姊妹彼此放箭,弄錯了方向.真是求神憐憫!放箭須注意目的和方向；我們也必須把箭磨亮,去鏽除塵；我們必須在主恩典中操練,在神前得到潔淨,清潔合乎主用；預備行各樣的善事.
\newline
\newline
我們肯否讓主來磨亮?磨亮需要一番工夫,經艱苦的過程.我們當將自己完全交給主,主要使用我們,「使我成為磨亮的箭,將我藏在他箭袋之中.」(2)我們被磨亮之後,主把我們隱藏,暫不使用.也許我們覺得沒機會服事主為主工作；其實不然,在箭袋中似乎是隱藏的；但主隨時可以使用.我們有沒這個準備.
\newline
\newline
中文舊約聖經常有「我在這裡」四字,英文只有here I am三字,希伯來文只有一字「看我」.我們能否對主說:「主啊!我在這裡,看我,我已準備好了；你已磨亮我,潔淨我；把我放在箭袋裡.我願意讓你使用.」
\newline
\newline
很多時候,我們覺得可以工作,要熱心傳福音；可是工作沒效果,說:「我勞碌是徒然,我盡力是虛無虛空,」「然而我當得的理必在耶和華那裡,我的賞賜必在我神那裡.」(4)工作不一定有表面的效果,而表面的效果不一定是實際的效果.我們作主工,不是看效果如何,而是看自己的態度如何,工作情形如何.神要的是忠心,只要我們忠心去幹,其他一切神必負責.
\newline
\newline
神看顧施恩,使我們作眾民的中保,我們須把福音帶出去,神已為我們開了大道.「看哪,這些從遠方來,這些從北方,從西方來,這些從秦國來.」(12)早期的教會學者,認為秦國是我們中國,但聖經註釋原文是希尼.我想這是正確.在一九四七年死海古卷發現,也有這樣的字眼,指的是西方.我真是盼望秦國是指中國；但我相信神要使用華人.世界各地有華人,神要華人在時代憂患中傳福音；華人好像以色列人受過很多患難；我們知道何謂憂患,這樣,我們傳福音就更有力量.
\newline
\newline
福音使者應發的聲音「諸天哪,應當歡呼；大地阿,應當快樂；眾山哪,應當發聲歌唱；因為耶和華已經安慰他的百姓,也要憐恤他困苦之民.」(13)神是安慰的神,安慰是福音的內容；但是我們怎樣來唱歌呢?歷代記載希西家燔祭一獻,就唱出耶和華的歌.沒有燔祭,沒奉獻,沒擺上就沒歌唱.
\newline
\newline
我們願意歌唱,就當毫無保留將自己擺上.
\newline
\newline


\newpage

\section{第60屆~港九培靈會~唐佑之~第五講}
\label{sec:1224}
\textbf{待命}
\newline
\newline
連結: \href{https://www.hkbibleconference.org/session-message/view/1224}{\texttt{https://www.hkbibleconference.org/session-message/view/1224}}
\newline
\newline
\hyperref[sec:1223]{< < < PREV SERMON < < <}
~
\hyperref[sec:toc]{[返目錄]}
~
\hyperref[sec:1226]{> > > NEXT SERMON > > >}
\newline
\newline
以賽亞書 50:1-50:11
\newline
\begin{longtable}{cl}
\hline
\hline
章節 & 經文 (和合本修訂版)\\
\hline
50:1 & \begin{tabularx}{0.7\textwidth}{X} 耶和華如此說：「我休了你們的母親，她的休書在哪裡呢？我將你們賣給了我哪一個債主呢？看哪，你們被賣是因你們的罪孽；你們的母親被休，是因你們的過犯。 \end{tabularx} \\ \\ \relax
50:2 & \begin{tabularx}{0.7\textwidth}{X} 我來的時候，為何沒有人呢？我呼喚的時候，為何無人回應呢？我的膀臂豈是過短、不能救贖嗎？我豈無拯救之力嗎？看哪，我一斥責，海就乾了；我使江河變為曠野，其中的魚因無水腥臭，乾渴而死。 \end{tabularx} \\ \\ \relax
50:3 & \begin{tabularx}{0.7\textwidth}{X} 我使諸天以黑暗為衣，以麻布為遮蓋。」 \end{tabularx} \\ \\ \relax
50:4 & \begin{tabularx}{0.7\textwidth}{X} 主耶和華賜我受教者的舌頭，使我知道怎樣用言語扶助疲乏的人。主每天早晨喚醒，喚醒我的耳朵，使我能聽，像受教者一樣。 \end{tabularx} \\ \\ \relax
50:5 & \begin{tabularx}{0.7\textwidth}{X} 主耶和華開啟我的耳朵，我並未違背，也未退後。 \end{tabularx} \\ \\ \relax
50:6 & \begin{tabularx}{0.7\textwidth}{X} 人打我的背，我任他打；人拔我兩頰的鬍鬚，我由他拔；人侮辱我，向我吐唾沫，我並不掩面。 \end{tabularx} \\ \\ \relax
50:7 & \begin{tabularx}{0.7\textwidth}{X} 主耶和華必幫助我，所以我不抱愧。我硬著臉面好像堅石，也知道我必不致蒙羞。 \end{tabularx} \\ \\ \relax
50:8 & \begin{tabularx}{0.7\textwidth}{X} 稱我為義的與我相近；誰與我爭論，讓我們來對質；誰與我作對，讓他近前來吧！ \end{tabularx} \\ \\ \relax
50:9 & \begin{tabularx}{0.7\textwidth}{X} 看哪，主耶和華必幫助我，誰能定我有罪呢？看哪，他們都要像衣服漸漸破舊，被蛀蟲蛀光。 \end{tabularx} \\ \\ \relax
50:10 & \begin{tabularx}{0.7\textwidth}{X} 你們當中有誰是敬畏耶和華，聽從他僕人的話語，卻行在黑暗中，沒有亮光的，當倚靠耶和華的名，仰賴自己的神。 \end{tabularx} \\ \\ \relax
50:11 & \begin{tabularx}{0.7\textwidth}{X} 看哪，你們當中所有點火、以火把圍繞自己的人，當行走在你們的火焰裡，並你們所點的火把中。這是我親手為你們定的：你們必躺臥在悲慘之中。 \end{tabularx} \\ \\
[1ex]
\hline
\hline
\end{longtable}
本書40-50章論及以色列民族的復興,也應用在教會的復興上.舊約聖經讓我們明白以色列民的歷史,明白基督救恩的奇妙；藉以應用在今天教會與教會中個人的生活上.
\newline
\newline
以色列民歷史的經驗中,有很多憂患；神藉這些憂患,使他們明白祂的恩典.叫他們不但自己享受神的恩典,同時也要把神的恩典帶給別人.
\newline
\newline
神是歷史之主,每個歷史細節都經過神所允許；不單是以色列的歷史,也是指著整個世界的歷史.神在歷史之中,在歷史之上,神統管萬有；耶和華坐著為王,甚至大水泛濫之時,歷史環境極其惡劣之時,神仍坐著為王.這是基督徒應有的認識,而且神也藉歷史啟示出來,讓人知道祂是耶和華獨一的真神.
\newline
\newline
歷史是神與人的對話,叫人知道救恩的需要,知道人在罪惡的環境中；未信者應接受救恩,已信者要擔負傳福音的使命.
\newline
\newline
歷史是神審判的事實,當歷史有危機,或時機,或轉機時,就是說明神公義的審判.我們面對一切轉機,危機,時機之時,必須擔負福音使命.
\newline
\newline
以色列人在憂患中,感到失望,煩惱,前途艱苦,他們已不知所措.人在患難中,有兩種可能的反應；一是感謝神,因為知道神藉歷史的環境叫我們更明白教會當負的使命.另一反應是發怨言,埋怨神不看顧,就像昔日的以色列人一樣.所以神藉以賽亞告訴他們說:「婦人焉能忘記她喫奶的嬰孩,不憐恤她所生的兒子；即或有忘記的,我卻不忘記你.」(四十九:15)感謝主!這不但是聖經所載的真理,也是歷史的事實.在歷代的教會,神沒有忘記；可能教會遭遇極大的困難逼迫,但是神不會忘記；神是看顧人的神.耶穌復活的當晚,門關了；但耶穌站在門徒中間,對他們說:「願你們平安.」神不會忘記祂的教會,不會忘記祂的百姓,也不會忘記我們.這裡,先知用婦女的情形來描寫,是聖經中少有的,聖經常用父親代表神,很少講到母親.講到神是丈夫,君主,牧人,都是男性表現.為何不用女性呢?因為歷史文化背景,當時以色列的鄰邦有很多假神,有很多不健全甚至淫穢的思想；所以聖經竭力避免說神是母親是婦女；並非否認婦女的尊貴；相反地,在此我們特別看見先知強調母性的光輝.到了本書66章,再用母性描寫神:「母親怎樣安慰兒子,神就照樣安慰你們......」(13)先知在必要時用女性更高貴地,把神的恩典和慈愛表達,用最好的比方來說明神不會忘記我們.神要以色列人建造會幕,因為會幕是人與神相會地方.所羅門建聖殿完成,他禱告時聽見神的聲音說:「......我的眼,我的心,也必常在那裡.」神的眼是看顧的眼,神的心是關懷的心,常在聖殿內,在敬拜的地方,在教會,在我們敬拜之中.
\newline
\newline
神一定不會忘記我們「看哪,我將你銘刻在我掌上,你的牆垣常在我眼前.」(16)以色列人的思想中,手掌非常重要,最要的事寫在手掌上.神一定記得我們,關懷,注意；甚至勝過我們自己的注意,這是多麼大的福份!我們在神手掌之中,蒙保守,看顧；一直在主恩愛中.
\newline
\newline
「牆垣」意即保護.以色列人被擄,歸回,建造聖殿；尼希米再造城牆,這樣才是完全的保障.
\newline
\newline
教會需要城牆,雖然我們傳福音對外希望開放,希望更多人進來.我們應常開教會的門,但我們需要牆；否則教會無保障,沒保護,罪惡世俗的力量進入.家庭也是如此,需要牆垣,家有安全.我們的心需要牆垣,罪惡就不能進來,求主保守我們!
\newline
\newline
我們需要神的看顧,神看顧祂的聖地,就是應許之地,是神賜給以色列人的.申命記極強調只有一位神,一個子民,一種敬拜,一個土地,就是神所賜之地.神帶領以色列民從曠野進入迦南美地,他們出了埃及,若不進入迦南地,他們所得的救恩只不過一半.基督徒常有此現象,從罪惡中出來；但是沒有進入神的真理和聖潔,沒進入神所應許之美地.我們只得了一半,並非神不給我們；我們努力過約但河,爭戰,可是沒有進入.救恩不需要代價,是白白得來的；但進一步要得主的恩典,得神的聖潔；好像以色列人進到迦南美地還要爭戰.約書亞特別說明基督徒應作屬靈爭戰；可惜以色列人失敗了；神所賜流奶與蜜之地,成了荒廢之地；但是神要復興.神對我們也是這樣,祂要我們蒙受福份,賜給我們豐富的產業.詩十六篇述說神所賜多麼美好的基業.可惜淒涼荒廢!我們多麼需要神的憐憫!教會復興是神的工作；但我們必須在神面前仰望.「至於你荒廢淒涼之處,並你被毀壞之地.現今眾城居住必顯為太窄,吞滅你的必離你遙遠.」(19)
\newline
\newline
教會需要復興,教會不復興,座位空置；教會復興,座無虛設甚至太窄.我們應當努力福音工作.今日北美教會,有許多香港移民,在加拿大東部,多倫多,溫哥華幾乎每架飛機都是香港移民；所以教會的人特別多,他們都很高興,然而我卻略懷傷感!香港教會人數是否少了呢?如果大家努力,教會復興,人數反而要增加.各位!我們有這樣的異象嗎?相信神必恩待我們,向我們施恩；神必領導我們.
\newline
\newline
神是得勝的神,「主耶和華說,我必向列國舉手,向萬民豎立大旗......列王......王后......他們必......向你下拜,並舔你腳上的塵土......」(22-23)這裡表明他們是失敗者.我們有沒勇氣和福音爭戰?神說:「與你們相爭的我必與他們相爭.我要拯救你的兒女......凡有血氣的,必都知道我耶和華是你的救主,是你的救贖主,是雅各的大能者.」(25-26)我們是得勝者,靠主得勝且得勝得有餘；因為我們的主是得勝的主.很多時候只嘆息我們的軟弱,這是不對的；我們當為主得勝的能力而誇勝,而讚美!應有積極的態度.靠主的話,軟弱變為剛強,主就使用我們.我們雖是小蟲；但要成為打糧快齒的器具.我們有時輕看神,也看見自己太微小；其實,神多麼偉大!我們也多麼偉大!我們屬於神,是神的兒女,祂的能力,在我們軟弱的人身上顯得完全.祂是我們的主,我們倚靠祂必然得勝.當我們有得勝希望之時,必須記得我們是神的僕人,我們需要教會,學習,等候；靠祂的恩典經受時代的憂患,我們必須付上福音的代價.
\newline
\newline
第五十:1-3是聖父的聲音,4-9是聖子的聲音,10-11是聖靈的聲音.這時我們特別要注意聖子的聲音.主耶穌是神的兒子,當祂來到世上,存心卑微,全心順服.腓立比二章說祂是受苦的僕人.我們能從先知以賽亞的預言看出.可能以賽亞傳信息時,都不明白到底受苦之僕是誰；然而我們有新約的記載,比舊約時代的人了解更多.主耶穌一切事跡聖經都有記載,都應驗了先知的預言,那受苦的僕人就是神的兒子.保羅在腓立比書二:5說:「你們要以基督的心為心.」基督的心是僕人的心,是兒子的心；我們若需要祂的心,必須效法祂的榜樣,跟隨祂的腳步.
\newline
\newline
一．受教的舌頭(五十:4)第三首受苦僕人的詩,受苦僕人應有好的態度:以學習的態度受教.受教者本是指門徒而言.四福音題門徒或類似的字眼共250次.近年來我們注意如何學習作主的門徒.門徒需要受教且需要操練,有受教者的舌頭.在言語上若沒過失那就是完全的人.我們每天當在主前思想舌頭有沒得罪人,得罪神?讓神來教導,使我們知道如何用言語扶助疲乏的人.神給我們舌頭,目的要我們說造就人的話；要我們見證福音的大能,述說其真理的奇妙.主耶穌在世上,祂沒有一句話不是帶著神的愛的.我們真是需要操練舌頭,去幫助人,因為很多人心靈破碎.神沒有打發天使,卻打發我們用親熱的手,包裹破碎的心靈；神沒有吩咐天使,卻叫我們去安慰傷心的人.當我們有苦處,我們得神的安慰,也要安慰別人.(林後)我們是安慰使者.神說:「你要安慰安慰我的百姓.」
\newline
\newline
各位!如果你真的成就福音的使命,也就是成就安慰的使命.今天很多人需要安慰；我們當仰望神,從祂得著安慰的話.約伯受苦時,他的三個朋友對他所講的話,對是對的；可是沒有愛心,沒有同情也沒有安慰.教會中弟兄姊妹需要安慰.當我們把福音帶出去時,我們需要安慰人.
\newline
\newline
二．聆聽的耳朵:「主每早晨題醒題醒我的耳朵.使我能聽,像受教者一樣.」(4)我們不但須有受教的舌頭,且須有聆聽的耳朵.聽道容易.但行道太難.「聽道而不行道的,就像人對著鏡子看見自己本來的面目；看見,走後,隨即忘了他的相貌如何.」(雅一:23-24)雅各這話,使我發出會心的微笑；他是個男子漢,一般男子不大照鏡子的,所以照了之後,遂即忘記自己的相貌.女的就不然,照了又照,細心地照.所以我們聽道,要好像姊妹照鏡子一樣,不是單欣賞自己的美貌；且要找出自己的缺陷.我們在主前思想,省察,知道自己的軟弱,不夠美,不夠青春；需要主的恩典臨到我們.「以色列阿,你要聽......」(申六:4)留心聽,詳細聽,專心聽,全心聽；聽受,聽信,聽從；我們需要有聆聽的耳朵；我們來到神前安靜在主前；單獨與神交往,體會祂的同在,聽見祂的聲音.靈修很重要,作神僕人的,不是照己意行事；要明白主人的意思,所以我們每天都要等候祂的吩咐.主耶穌在世上,每清晨上山聆聽神的聲音,然後祂才說話,教導,實行.我們當學習作主的真正僕人.耶穌基督被釘在十字架之前,彼得削掉大祭司長的僕人的一個耳朵；耶穌即行神跡醫治那人的耳朵,免得將來在主審判台前藉詞因耳朵被削而不能聽；聖靈所說的話,凡有耳的,都應當聽.「主耶和華開通我的耳朵,我並沒有違背,也沒有退後.」(5)我們有聆聽的耳朵,為的是順服神；不違背,不膽怯,不退後,不冷淡.
\newline
\newline
三．痛苦的背項「人打我的背,我任他打,人拔我腮頰的鬍鬚,我由他拔；辱我吐我,我並不掩面.」(6)受苦的僕人甘心受辱.申命記廿五:9說罪犯是該被人打背的,是極大的侮辱；因為這是神的旨意.耶穌曾告訴門徒(當時也有其他聽道者)「有人打你右臉,連左臉也轉過來由他打.」我覺得被打沒反抗已夠忍耐,還要轉過另一邊來由他打,我實在做不到.基督徒常覺聖經的話很對；但太過理想,唯耶穌做得到.我們聽道而不行道,思想沒有結合行為,沒有把福音見證帶出來.主耶穌所說的,祂自己實行了；且要我們也去行.當耶穌被捉拿受審時,被人打:「耶穌說,我若說的不是,你可以指證那不是；我若說的是,你為甚麼打我呢?」(約十八:23)耶穌被打,被釘死在十字架,付出最重的代價；祂非消極地忍受,乃是積極地承受；祂說:如果我不釘十字架,我可以差遣十二營的天使來營救我.」耶穌甘心樂意被釘十字架,因為看見擺在前面的喜樂,就輕看十字架的苦難,祂為的是福音,救恩的緣故.
\newline
\newline
各位!這能否成為我們的經驗?我們需要靠神的恩典,化消極為積極；化被動為主動,化勉強為樂意.為何在時代的憂患中,我們還沒有承受福音的使命,福音是需要苦難來成就,也需要苦難來傳揚.我們傳福音沒有力量；因為我們沒有憂患,或不覺有憂患,或不肯付代價.
\newline
\newline
四．堅硬的臉面「主耶和華必幫助我,所以我不抱愧；我硬著臉面好像堅石,我也知道我必不蒙羞.」(7)第9節也有「主耶和華要幫助我們.」當主耶穌最後一次到耶路撒冷,祂為城哀哭.當主看見無花果樹不結果子,祂的心多麼難過；但祂極堅定,面向耶路撒冷.
\newline
\newline
我們也當硬著臉面好像堅石,堅定,確定地往前走；當主的旨意向我們顯明,當看見福音是我們的責任和使命時；毫不猶疑,毋須等待,勇往直前,好像主耶穌一樣.
\newline
\newline
有人說第三首僕人之詩是主耶穌的客西馬尼園經驗.聖經說耶穌進入客西馬尼園；祂再往前走.主進一步地順服,奉獻,仰望,敬拜；然後祂跪下,向天父禱告:「父阿!倘若可行,求你叫這杯離開我；然而不要照我的意思,而要照你的意思.」我否認基督徒常說禱告能改變萬事.耶穌禱告求父將苦杯拿掉,耶穌就不用釘十字架；那麼我們哪能得主的救恩?我們當說禱告不能改變萬事.因為神有絕對的權能,要照祂的旨意成就；不受人的支配.禱告只能改變禱告者；越禱告越順服,越交托,越沒有己見,完全照主的旨意.我認為客西馬尼園的禱告才是真正的主禱文.馬太福音第六章,其實並非主的禱告,乃是主教導門徒的禱告；因為耶穌無罪,無需在神前認罪.只能說這是信徒禱文,家禱文；真正的主禱文是耶穌在客西馬尼園的禱告「......不要照我的意思,只要照你的意思.」耶穌真正「待命」祂等待主的命令.
\newline
\newline
各位!要成就福音的使命,必須等待主的命令,如何等待?客西馬尼園的經驗,完全倚靠,交托,順服的禱告:「願主的旨意成就!」
\newline
\newline


\newpage

\section{第60屆~港九培靈會~唐佑之~第六講}
\label{sec:1226}
\textbf{更新}
\newline
\newline
連結: \href{https://www.hkbibleconference.org/session-message/view/1226}{\texttt{https://www.hkbibleconference.org/session-message/view/1226}}
\newline
\newline
\hyperref[sec:1224]{< < < PREV SERMON < < <}
~
\hyperref[sec:toc]{[返目錄]}
~
\hyperref[sec:1228]{> > > NEXT SERMON > > >}
\newline
\newline
以賽亞書 51:1-52:15
\newline
\begin{longtable}{cl}
\hline
\hline
章節 & 經文 (和合本修訂版)\\
\hline
51:1 & \begin{tabularx}{0.7\textwidth}{X} 追求公義、 尋求耶和華的人哪，當聽從我！你們要追想自己是從哪塊磐石鑿出，從哪個巖穴挖掘而來； \end{tabularx} \\ \\ \relax
51:2 & \begin{tabularx}{0.7\textwidth}{X} 要追想你們的祖宗亞伯拉罕和生你們的撒拉；因為我選召亞伯拉罕時，他只有一個人，但我賜福給他，使他增多。 \end{tabularx} \\ \\ \relax
51:3 & \begin{tabularx}{0.7\textwidth}{X} 耶和華已經安慰錫安，安慰了錫安一切的廢墟，使曠野如伊甸，使沙漠像耶和華的園子；其中必有歡喜、快樂、感謝，和歌唱的聲音。 \end{tabularx} \\ \\ \relax
51:4 & \begin{tabularx}{0.7\textwidth}{X} 我的民哪，要留心聽我，我的國啊，要向我側耳；因為訓誨必從我而出，我必使我的公理成為萬民之光。 \end{tabularx} \\ \\ \relax
51:5 & \begin{tabularx}{0.7\textwidth}{X} 我的公義臨近，我的救恩發出。我的膀臂要審判萬民，眾海島都要等候我，倚賴我的膀臂。 \end{tabularx} \\ \\ \relax
51:6 & \begin{tabularx}{0.7\textwidth}{X} 你們要向天舉目，觀看下面的地；天必像煙雲消散，地必如衣服漸漸破舊；其上的居民也要如此死亡。惟有我的救恩永遠長存，我的公義也不廢掉。 \end{tabularx} \\ \\ \relax
51:7 & \begin{tabularx}{0.7\textwidth}{X} 知道公義、將我的訓誨存在心中的人哪，當聽從我！不要怕人的辱罵，也不要因人的毀謗驚惶。 \end{tabularx} \\ \\ \relax
51:8 & \begin{tabularx}{0.7\textwidth}{X} 因為他們必像衣服被蛀蟲蛀；像羊毛被蟲子咬。惟有我的公義永遠長存，我的救恩直到萬代。 \end{tabularx} \\ \\ \relax
51:9 & \begin{tabularx}{0.7\textwidth}{X} 耶和華的膀臂啊，興起，興起！以能力為衣穿上，像古時的年日，像上古的世代一樣興起！從前砍碎拉哈伯、刺透大魚的，不是你嗎？ \end{tabularx} \\ \\ \relax
51:10 & \begin{tabularx}{0.7\textwidth}{X} 使海與深淵的水乾涸，在海的深處開路，使救贖的民走過的，不是你嗎？ \end{tabularx} \\ \\ \relax
51:11 & \begin{tabularx}{0.7\textwidth}{X} 耶和華救贖的民必歸回，歌唱來到錫安；永恆的喜樂必歸到他們頭上。他們必得著歡喜快樂，憂傷嘆息盡都逃避。 \end{tabularx} \\ \\ \relax
51:12 & \begin{tabularx}{0.7\textwidth}{X} 我，惟有我是安慰你們的。你是誰，竟怕那必死的人，怕那生命如草的世人， \end{tabularx} \\ \\ \relax
51:13 & \begin{tabularx}{0.7\textwidth}{X} 卻忘記鋪張諸天、立定地基、造你的耶和華？你因欺壓者圖謀毀滅所發的暴怒，終日害怕，其實那欺壓者的暴怒在哪裡呢？ \end{tabularx} \\ \\ \relax
51:14 & \begin{tabularx}{0.7\textwidth}{X} 被擄的即將得釋放，不至於死而下入地府，也不致缺乏食物。 \end{tabularx} \\ \\ \relax
51:15 & \begin{tabularx}{0.7\textwidth}{X} 我是耶和華－你的神，我攪動大海，使海中的波浪澎湃，萬軍之耶和華是我的名。 \end{tabularx} \\ \\ \relax
51:16 & \begin{tabularx}{0.7\textwidth}{X} 我已將我的話放在你口中，用我的手影遮蔽你，為要安定諸天，立定地基，並對錫安說：「你是我的百姓。」 \end{tabularx} \\ \\ \relax
51:17 & \begin{tabularx}{0.7\textwidth}{X} 耶路撒冷啊，興起，興起！站起來！你從耶和華手中喝了他憤怒的杯，那使人東倒西歪的杯，直到喝盡。 \end{tabularx} \\ \\ \relax
51:18 & \begin{tabularx}{0.7\textwidth}{X} 她所生育的孩子中，沒有一個攙她的；她所撫養的孩子中，沒有一個扶她的。 \end{tabularx} \\ \\ \relax
51:19 & \begin{tabularx}{0.7\textwidth}{X} 這雙重的災難臨到你，有誰憐憫你呢？破壞和毀滅，饑荒和戰爭臨到，我如何能安慰你呢？ \end{tabularx} \\ \\ \relax
51:20 & \begin{tabularx}{0.7\textwidth}{X} 你的孩子發昏，在各街頭躺臥，如同網羅裡的羚羊，滿了耶和華的憤怒，滿了你神的斥責。 \end{tabularx} \\ \\ \relax
51:21 & \begin{tabularx}{0.7\textwidth}{X} 因此，你這困苦卻非因酒而醉的，當聽這話， \end{tabularx} \\ \\ \relax
51:22 & \begin{tabularx}{0.7\textwidth}{X} 你的主，耶和華，就是為他百姓辯護的神如此說：「看哪，我已從你手中接過那使人東倒西歪的杯，就是我憤怒的杯，你必不再喝。 \end{tabularx} \\ \\ \relax
51:23 & \begin{tabularx}{0.7\textwidth}{X} 我必將這杯遞在苦待你的人手中。他們曾對你說：『你屈身，任我們踐踏過去吧！』你就以背為地，又如街道，任人走過。 \end{tabularx} \\ \\ \relax
52:1 & \begin{tabularx}{0.7\textwidth}{X} 錫安哪，興起！興起！穿上你的能力！聖城耶路撒冷啊，穿上你華美的衣服！因為從今以後，未受割禮、不潔淨的必不再進入你中間。 \end{tabularx} \\ \\ \relax
52:2 & \begin{tabularx}{0.7\textwidth}{X} 耶路撒冷啊，抖去塵埃，起來坐在王位上！被擄的錫安哪，解開你頸上的鎖鏈！ \end{tabularx} \\ \\ \relax
52:3 & \begin{tabularx}{0.7\textwidth}{X} 耶和華如此說：「你們白白地被賣，也必不用銀子贖回。」 \end{tabularx} \\ \\ \relax
52:4 & \begin{tabularx}{0.7\textwidth}{X} 主耶和華如此說：「先前我的百姓下到埃及，在那裡寄居，末後又有亞述人欺壓他們。 \end{tabularx} \\ \\ \relax
52:5 & \begin{tabularx}{0.7\textwidth}{X} 我的百姓既是白白地被擄，如今我在這裡做甚麼呢？這是耶和華說的。轄制他們的人歡呼，我的名終日不斷受褻瀆，這是耶和華說的。 \end{tabularx} \\ \\ \relax
52:6 & \begin{tabularx}{0.7\textwidth}{X} 因此，我的百姓必認識我的名；在那日，他們必知道說這話的就是我。看哪，是我！」 \end{tabularx} \\ \\ \relax
52:7 & \begin{tabularx}{0.7\textwidth}{X} 在山上報佳音，傳平安，報好信息，傳揚救恩，那人的腳蹤何等佳美啊！他對錫安說：「你的神作王了！」 \end{tabularx} \\ \\ \relax
52:8 & \begin{tabularx}{0.7\textwidth}{X} 聽啊，你守望之人的聲音，他們揚聲一同歡唱；因為他們必親眼看見耶和華返回錫安。 \end{tabularx} \\ \\ \relax
52:9 & \begin{tabularx}{0.7\textwidth}{X} 耶路撒冷的廢墟啊，要出聲一同歡唱；因為耶和華安慰了他的百姓，救贖了耶路撒冷。 \end{tabularx} \\ \\ \relax
52:10 & \begin{tabularx}{0.7\textwidth}{X} 耶和華在萬國眼前露出聖臂，地的四極都要看見我們神的救恩。 \end{tabularx} \\ \\ \relax
52:11 & \begin{tabularx}{0.7\textwidth}{X} 離開吧！離開吧！你們要從巴比倫出來。你們扛抬耶和華器皿的人哪，不要沾染不潔淨之物，離去時務要保持潔淨。 \end{tabularx} \\ \\ \relax
52:12 & \begin{tabularx}{0.7\textwidth}{X} 你們出來必不致匆忙，也不致奔逃；因為耶和華要在你們前頭行，以色列的神必作你們的後盾。 \end{tabularx} \\ \\ \relax
52:13 & \begin{tabularx}{0.7\textwidth}{X} 看哪，我的僕人行事必有智慧，他必被高升，高舉，升到至高之處。 \end{tabularx} \\ \\ \relax
52:14 & \begin{tabularx}{0.7\textwidth}{X} 許多人因他驚奇---他的面貌比別人憔悴，他的外表比世人枯槁--- \end{tabularx} \\ \\ \relax
52:15 & \begin{tabularx}{0.7\textwidth}{X} 同樣，他也必使許多國家驚奇，君王要向他閉口。未曾傳給他們的，他們必看見；未曾聽見過的事，他們要明白。 \end{tabularx} \\ \\
[1ex]
\hline
\hline
\end{longtable}
第五十一章開頭:「你們要追想被鑿而出的磐石,被挖而出的巖穴.要追想你們的祖先亞伯拉罕」神叫亞伯拉罕從迦勒底的吾珥出來,離開那拜偶像之地到未知之地,叫他家族因他得福.這是家族的更新.神叫以色列人出埃及,離開拜偶像之地,使這民族得以更新.以色列民第二次出埃及,如五十二:11:「你們離開罷,離開罷,從巴比倫出來.」離開是更新的途徑,不離開就不能更新.神叫以色列人離開巴比倫被擄之地,偶像之地,使他們得著更新.我們也應離開不該在之地方,不該作之事,不應有的思想,我們就得著更新.亞伯拉罕離開吾珥以後,這個家族就成為敬拜的家族,使以色列人常追想亞伯拉罕和列祖們敬拜的經驗.敬拜是更新唯一必須的條件.雅各曾離家出走,在伯特利遇見神；他說那地是神的殿,天的門；他在神前敬拜,每當他軟弱時,神叫他追想伯特利.本來他走的路不很正確,但當他回來時,在示劍地方；神叫他必須離開,叫他起來到伯特利去.我們應當到神所要我們去的地方,要追想神的恩典,追想蒙恩的經驗,追想在教會中所得的福份.這是我們更新的第一步.
\newline
\newline
伊甸的更新,「耶和華已經安慰錫安,和錫安一切的荒場,使曠野像伊甸.」(五:3)神為我們所預備的是伊甸園；可是人犯罪,不要樂園而要曠野.我們在人類歷史中,看的都是曠野.西方文壇荒原詩人,描寫整個世界是一片曠野.世界花綠美麗,實際是墳埸；所有的人都「死在罪惡過犯之中,」(弗二:1)沒有生命氣息,只有死亡和絕望.神的恩典要我們更新,生命更新；把曠野變為樂園,樂園失而復得.
\newline
\newline
在真理上更新,公理是神公義審判的道理.歷史是神與人的對話,或謂神公義的作為.把審判的消息帶出來,叫人歸向祂；因為神審判的消息乃「萬民之光」(4)「因為訓誨必從我而出,」(4)律法只能使人知罪,人不能靠律法得救.我們常覺得律法是反面的意思,其實並非我們想像中那麼不好；不過當時保羅為保證猶太人對律法的錯解,但我們並沒律法的錯誤,也沒猶太人的背景,為何我們要先進到猶太人的錯誤裡面,然後再出來呢?這是因為我們沒有讀清楚全本聖經.從前宣教士把真理帶出來,西方思想受了希臘的影響,知道辨證的方法；但華人不需要這種背景,我們可以直接讀.律法是神的訓誨,是神的旨意,是神的話語.很多時候因為我們不注重律法,不注重神的訓誨；所以我們在恩典中墮落了.我們說不要律法,只要恩典；恩典是白白賜予的,雖然我們不好,神仍然恩待我們.其實律法裡面有恩典,恩典裡面有律法；神的慈愛中有嚴厲,有公義,神是輕慢不得的.如果我們能夠真正注意這方面,我們一定會注意信心和信心的行為.求主幫助!我們需要神的訓誨,需要神的話語；或說我們需要神的律法,我們在神面前應該端正,謹慎,敬虔；在主前好好地聽祂的話,遵從祂的旨意.
\newline
\newline
伊甸的更新把我們整個的生命改變,叫我們也在真理上更新,因為神審判的信息成為萬民之光.神說:「我的公義臨近,我的救恩發出,我的膀臂要審判萬民.」(5)「耶和華的膀臂阿,興起,興起,以能力為衣穿上.」(9)五十一和五十二兩章,屢次提興起,加重語氣,表明意義深長.耶和華有大能的膀臂,祂是行動的神,祂說了,事就成了.
\newline
\newline
在禱告上更新,我們禱告常有錯誤的觀念.只有禱告,沒有行動,這樣等於逃避；不想,不敢,也不願意做甚麼,一味禱告.我們禱告定要對主說:「我願意照祂命令去做.」有時我們不知事情如何做,是否應該做；弟兄姊妹來問我,我也不知如何是好；只叫他們禱告.這是逃避,教會的長者,傳道人竟然也逃避；沒有真正在神的真理上下工夫,自己不知怎樣幫助弟兄姊妹,竟以禱告掩飾.我認為這是不可饒恕的罪.我們應該禱告,讀經,尋求神的旨意.若禱告而無行動,則逃避；若行動無禱告,則悖逆；根本不明白神的旨意；這都是罪,是神所不喜悅的.我們應該在神前迫切禱告,求神給我們行動的能力,遵行神旨意的能力.傳福音工作,是神的旨意；但我們必須迫切禱告,仰望神.
\newline
\newline
神是得勝的神「從前砍碎拉哈伯.」(10)(拉哈伯是從前巴比倫埃及神話中的怪物)神勝過一切,不論地上的,地底的,字宙的萬物.我們不要常嘆息自己的軟弱,要靠主得勝而誇口；發出哈利路亞的讚美.主是得勝的王,萬王之王；祂是主,萬主之主；祂「使海的深處變為贖民經過之路.」(10)「耶和華救贖的民必歸回,歌唱來到錫安；」(11)「惟有我是安慰你們的；」(12)
\newline
\newline
聖經中提到「我」是多麼親切,和我們有密切關係.「安息」多麼親熱的字眼!福音的信息就是安慰；安慰是神的嘆息,神激烈的情緒,神關心的情懷；有說不盡的心意,要我們蒙受祂的福份,神要摸到我們的心.在新約「安慰」即像母親懷中的孩子享受慈母的撫摸,安慰,看顧.神說:「惟有我,是安慰你們的；你是誰?」這話實在太寶貝!當我們禱告,默想,做事之時；神要問「你是誰?」我覺得以賽亞照神的囑咐問人是誰很有意義.「人算甚麼?」(詩八:4)人不是東西,人是人；除了罵人才說「人是甚麼」,應當說:「人是誰?」當神問:「你是誰?」我心裡充滿感謝,流淚對主說:「今天我成為何等樣人,乃因你的恩而成,我是屬神的兒女,是神的僕人,我多麼尊榮!」我們有沒有尊重神給我們這尊榮的身份?當路得的婆婆叫她在夜間去找波阿斯,(得)並非叫她做不可告人的事；乃是要顧及波阿斯的面子,因路得是外邦人；也許波阿斯不願娶她為妻.波斯願意接納她；她回家,婆婆說:「女兒阿,怎麼樣了.」這是很好的中文意譯.直譯:「你是誰?」路得的婆婆問她「你是誰?」她回家時臉發光輝,因她前途蒙神的恩惠,找到了歸宿.我們在神前也是這樣,我們的樣子,身份,沒有改變；當神問:「你是誰?」我們當存快樂感謝的態度說:「神啊!我是屬你的,我永遠屬於你.」
\newline
\newline
神要審判,「耶路撒冷阿,興起,興起,站起來,你從耶和華手中喝了祂忿怒之杯,」(17)神賜能力給祂的僕人,給祂的百姓之前,神要審判；因審判從神的家起首；如果沒管教,沒審判；神的百姓就不配,也不能作福音的使者.如果我們真正的興起,我們需要神赦免之恩.神先要審判教會,而後審判全世界；這時就是教會把福音帶出去的時候,當天國福音傳遍天下；主等候快再來；祂尚未來,並非祂不願意來；而是耶穌寬容,希望人人都得救.
\newline
\newline
教會,弟兄姊妹,你,我,我們都尚未準備,沒有興起,沒有真正在認罪悔改,沒接受神的審判.「主耶和華阿,你若究察罪孽,誰能站得住呢?但在你有赦免之恩,」(詩一○三)神的「怒氣不過是轉眼之間,祂的恩典乃是一生之久；一宿雖然有哭泣,早晨便必歡呼.」(詩卅:3)我們有痛苦艱難,有時可能是神的審判；但我們若認罪悔改,審判要過去的；如果你真正願意在祂裡面得著更新,神必恩待.
\newline
\newline
五十二1是正面的話,他們求神的膀臂興起,幫助他們有能力.他們經過艱難,患難,管教,審判,得到赦罪的平安,如今真正興起,披上神所賜的能力,穿上華美的衣服.華美衣服原意即新娘的禮服.教會是基督的新婦；我們要清潔,華美,專一地獻給主；祂用清潔的道完全洗淨,叫我們沒有皺紋等類的疾病,叫我們保持美麗青春,不衰老；穿上華美之衣等新郎.新郎快來了!你準備好了嗎?華美衣服乃指聖靈的能力.「耶和華的靈降在基甸身上.」(士六:34意譯)直譯是「神把祂的靈穿在基甸身上.」意即神的靈就是基甸的衣服；也是我們的衣服,我們穿上神的靈,就有見證的能力,實際作福音的使命.神看我們多寶貝!我們不夠美,神要我們穿上華美的衣服；不是外表的,物質的；而是神要我們真正得著聖靈的能力,成為見證人.「聖靈降臨在你們身上,你們就必得著能力......」(徒一:8)保羅說:「十字架的道理,在那滅亡的人為愚拙,在我們得救的人卻為神的大能.」(林前一:18)「我們有這寶貝放在瓦器裡,要顯明這莫大的能力,是出於神,不是出於我們.」(林後四:7)我們需要聖靈的能力,才能夠傳福音,承擔福音使命.神為叫我們在時代的憂患中有此體驗.神說:「所以我的百姓必知道我的名,到那日他們必知道說這話的就是我；看哪,是我.」(五十二:6)神願意與我們同在,看顧,保守,潔淨,使用,充滿我們,神已準備好；我們必須準備,回應主:「我在這裡,我願隨時給你使用.」我們當像磨亮鋒利的箭在箭袋裡,隨時可用.我們必須更新,在神前等候；工作更新,有工作能力；工作異象更新,看得清楚,作得實際,果敢地幹.
\newline
\newline
在報信息和傳福音的事工上更新,第7節給我們有此看見,對福音重新認識,重視,重新宣告神作主了.神是萬王之王,有絕對主權,祂作王,全地有平安,萬民都戰抖,大水泛濫之時,耶和華坐著為王,祂是絕對的主,是歷史之王,掌管一切.「這人的腳登山何等美.」(7未句)這裡的原意較我們所了解的更豐富；神藉以賽亞給我們明白,作為神的兒女要打扮,不是靠化裝品和時尚衣服,而是以聖靈的能力作我們華美之衣；有佳美的信息作我們腳上的裝飾,這還不夠美,更需要神來裝備上更新,在主前一同歌唱.
\newline
\newline
作時代的守望者「聽阿,你守望之人的聲音,他們揚起聲來,一同歌唱；因為耶和華歸回錫安的時候,他們必親眼看見.」(8)誰是守望者?一是軍事上的守望,站在守望樓上窺探敵人出現,就吹角警告.這是先知工作,我們都有責任,作時代的守望者.時代充滿了憂患,我們發出聲音,讓人知道若不悔改就會滅亡.教會需要守望的工作!守望者還無有祭司的工作.在耶路撒冷聖殿外,祭司要按次輪值,守望等候天亮；黎明之前黑夜更深,疲倦愈加；但守望的祭司必須加倍儆醒,睜眼面向東方；等曙光一現即吹號角,叫醒全耶路撒冷的人.
\newline
\newline
各位弟兄姊妹!我們當作黑夜的守望者,漫長黑夜怎能昏昏欲睡呢?天將亮,榮耀晨曦將臨,耶穌將榮耀再來；還有一點時候要來的就來,不再遲延!故此,我們急需把福音帶出去；但自身先要對付,要準備,要裝備；在神前清清潔潔.「你們離開罷,要從其中出來......務要自潔.」(11)自潔作主工,作主見證,傳主福音.當約書亞領以色列人過約但河時,對他們說:「要自潔,因為明天耶和華在我們中間行奇事.」奇事即神蹟,是人所作不到的；但必需的條件,務要自潔；今天的自潔,保證明天的奇事.神要在我們中間行奇事,給我們更新的恩典；然而要自潔,合乎主用,預備行各樣善事.耶和華必在前頭行,以色列的神必作你們的後盾；如牧羊人在羊的前頭走,且有恩惠慈愛隨著祂的羊群.有人前面走,牧羊狗後面保護.我們安穩一直向前走,不許退後；只許得勝,不許失敗.
\newline
\newline
主與我們同在!我們必需自潔!
\newline
\newline


\newpage

\section{第60屆~港九培靈會~唐佑之~第七講}
\label{sec:1228}
\textbf{承受}
\newline
\newline
連結: \href{https://www.hkbibleconference.org/session-message/view/1228}{\texttt{https://www.hkbibleconference.org/session-message/view/1228}}
\newline
\newline
\hyperref[sec:1226]{< < < PREV SERMON < < <}
~
\hyperref[sec:toc]{[返目錄]}
~
\hyperref[sec:1229]{> > > NEXT SERMON > > >}
\newline
\newline
以賽亞書 53:1-53:12
\newline
\begin{longtable}{cl}
\hline
\hline
章節 & 經文 (和合本修訂版)\\
\hline
53:1 & \begin{tabularx}{0.7\textwidth}{X} 我們所傳的有誰信呢？耶和華的膀臂向誰顯露呢？ \end{tabularx} \\ \\ \relax
53:2 & \begin{tabularx}{0.7\textwidth}{X} 他在耶和華面前生長如嫩芽，像根出於乾地。他無佳形美容使我們注視他，也無美貌使我們仰慕他。 \end{tabularx} \\ \\ \relax
53:3 & \begin{tabularx}{0.7\textwidth}{X} 他被藐視，被人厭棄；多受痛苦，常經憂患。他被藐視，好像被人掩面不看的一樣，我們也不尊重他。 \end{tabularx} \\ \\ \relax
53:4 & \begin{tabularx}{0.7\textwidth}{X} 他誠然擔當我們的憂患，背負我們的痛苦；我們卻以為他受責罰，是被神擊打苦待。 \end{tabularx} \\ \\ \relax
53:5 & \begin{tabularx}{0.7\textwidth}{X} 他為我們的過犯受害，為我們的罪孽被壓傷。因他受的懲罰，我們得平安；因他受的鞭傷，我們得醫治。 \end{tabularx} \\ \\ \relax
53:6 & \begin{tabularx}{0.7\textwidth}{X} 我們都如羊走迷，各人偏行己路；耶和華使我們眾人的罪孽都歸在他身上。 \end{tabularx} \\ \\ \relax
53:7 & \begin{tabularx}{0.7\textwidth}{X} 他被欺壓受苦，卻不開口；他像羔羊被牽去宰殺，又像羊在剪毛的人手下無聲，他也是這樣不開口。 \end{tabularx} \\ \\ \relax
53:8 & \begin{tabularx}{0.7\textwidth}{X} 因受欺壓和審判，他被奪去，誰能想到他的世代呢？因為他從活人之地被剪除，為我百姓的罪過他被帶到死裡。 \end{tabularx} \\ \\ \relax
53:9 & \begin{tabularx}{0.7\textwidth}{X} 他雖然未行殘暴，口中也沒有詭詐，人還使他與惡人同穴，與財主同墓。 \end{tabularx} \\ \\ \relax
53:10 & \begin{tabularx}{0.7\textwidth}{X} 耶和華的旨意要壓傷他，使他受苦。當他的生命作為贖罪祭時，他必看見後裔，他的年日必然長久。耶和華所喜悅的事，必在他手中亨通。 \end{tabularx} \\ \\ \relax
53:11 & \begin{tabularx}{0.7\textwidth}{X} 因自己的勞苦，他必看見光就心滿意足。因自己的認識，我的義僕使許多人得稱為義，他要擔當他們的罪孽。 \end{tabularx} \\ \\ \relax
53:12 & \begin{tabularx}{0.7\textwidth}{X} 因此，我要使他與位大的同份，與強盛的均分擄物。因為他傾倒自己的生命，以致於死，也列在罪犯之中。他卻擔當多人的罪，為他們的過犯代求。 \end{tabularx} \\ \\
[1ex]
\hline
\hline
\end{longtable}
有人說耶利米書是聖父的書；因為從中可以聽到父神的聲音,關懷之父站在山頭上呼叫；「背道的兒女啊!回來罷!」有人說以賽亞書是聖子的書或福音書,其中充滿關於耶穌基督是彌賽亞的道理.以西結書是聖靈的書提到耶和華的靈或耶和華的手在先知身上.
\newline
\newline
我們從先知書中能看見聖父,聖子,聖靈,奇妙的工作.先知的信息是為當代,也是為每個時代；因為神的話永遠有功效:當代有當代的福音和信息,但在每時代,神都可以應用祂的話語在每個人身上.神的話沒有一句不帶著能力的.
\newline
\newline
本書由40-55章共16章,新約羅馬書也有十六章的經文,所以有人說本書這16章可說是舊約的羅馬書；其中論及基本救恩,集中於救主身上.當我們讀這16章聖經記載,我們就能明白我們的救主,更認識仰望祂,更信靠信服祂.甚至有人說全卷以賽亞書66章,好像全部聖經共66卷；舊約有39卷,本書於一至卅九章都講主基督再來的預言.從40章開始,我們如同讀福音書；在曠野有人聲喊說,要修正預備主的道.主耶穌來時,約翰為祂預備道路.本書第40章正像福音書開章.當我們讀到本書第66章,論及新天新地,就像啟示錄所載一樣.從40-66章正如新約全部27卷,中間第53章,論及救主的十字架,永遠是我們的榮耀；這不但是舊約的最重要的真理,而且新約屢次應用這些話.我們讀聖經,讀到神的真理極其偉大,如大海汪洋,更覺得我們知道的實在太少；讀了神的真理,內心感受擴大,我們當求主擴大我們屬靈的容量,更多地容納神的話語.有時我們好像小孩子在海邊,用小貝殼盛水；洋洋自得.我們讀聖經,誠然所得無幾；但若真正得著,也能充分使我們享受主的恩典.主的真理實在太奇妙!求神真理的靈光照引導我們!
\newline
\newline
第53章是最主要的經文.受苦的道理是十字架的奧秘；我們知道這是指主而言,因為有新約的憑證；好像使徒行傳第八章,埃提阿伯太監讀聖經不明白；神感動腓利教導他.
\newline
\newline
神的僕人到底是誰?首先我們想到神揀選以色列整體民族作祂的僕人.僕人的工作是把福音,把主的道理傳揚.但是,僕人不單指民族的整體,也指個人；舊約彌賽亞是新約的基督,事實上,僕人也是指著基督.進一步我們知道基督的身體就是教會,教會的整體就是神的僕人,我們都是屬主的；祂如何,我們也當如何；祂作僕人,我們也作僕人；我們應有作僕人的樣式,心態,殷勤的態度.
\newline
\newline
主耶穌是神的僕人,我們也是神的僕人,我們重新為此事奉,服事.主是受苦之僕；如果僕人須要受苦,你我也必須受苦,成為受苦之僕；如果僕人須承受憂患,你我也要承受憂患.今日是憂患時代,我們必須體會時代的憂患,然後才知道在這時代傳福音的工作.以色列民受托作此工作,他們失敗,但神並沒有失敗,甚至幫助鼓勵他們.叫他們在審判之後還有餘數；餘數中他們要成為民族復興之核心,眾民的中保,外邦之光.教會也當如此.
\newline
\newline
主耶穌是我們的榜樣,是我們的標竿；我們跟隨這位中保,完全跟隨祂；如果祂承受憂患,我們也要如此.我們在主面前,明白,得著,享受祂的恩典,我們須把祂的恩典帶出來.這就是我們讀本書最主要的意思.
\newline
\newline
有時我們只談53章,其實這首詩應從五十二:13開始.神說:「我的僕人行事必有智慧.」這裡記有「看哪,」42章一開頭就有「看哪,我的僕人,」表明集中注意力看主──神的僕人.在未提及主耶穌這位受苦的僕人,先提這位僕人要得尊貴榮耀.
\newline
\newline
僕人需要工作,行動；教會必須有傳福音的行動,有服事人；事奉神的行動.我們不過是僕人,沒有自己的主張,計劃,選擇；惟聽命主人.主,這位受苦的僕人行事有智慧,祂有屬靈的見解,是主人的見解；僕人有智慧,表明他完全照主人的意思而作.智慧另譯行事通達,意即遵從主人而作,一定成功亨通的；不是表面,乃是實際；神一定要屬祂的人作成他的事.意即祂被高舉上升為至高,這是屬靈的原則,當我們降卑時就升高,(弗四:9)主耶穌尚且這樣；我們豈不更須要神的恩典給我們降卑?我們從聖經多處看,均須要體會神奇妙的旨意阿.
\newline
\newline
舊約以色列民等待彌賽亞來臨,彌賽亞是出自大衛家,是位顯赫的君王,有權威,有公義,帶來和平.本書第七章論及童女生子,名叫以馬內利.第九章載,有一子為我們而生,政權必擔在他肩上,他是奇妙的策士,和平之君,全能的神.第十一章論及祂出生卑微,好像耶西的根；但沒有提及大衛和他的父親耶西,先知預言的這位；是因有神的靈與祂同在,所以成了非常重要偉大的君王.從40章開始,我們看見這位彌賽亞並非一位威武的君王；他是位受苦之僕,這是許多人不易明白的事.耶穌基督降世,說:「我來非以役人,乃役於人；且要捨命,作多人的贖價.」僕人是卑微不受重視的,主道成肉身所作成的工；當時一般人都不明白他以僕人姿態出現；甚至門徒都不了解；還希望耶穌做君王時,坐在祂左右邊.耶穌一直教導他們,大的要服事小的；雖然如此,他們仍不明白.當耶穌離世之前為門徒洗腳.約翰福音十三章是極重要的真理；如果細讀,可知耶穌自知離世歸父的時候到了,還有一點時候,未免有緊急感.我們屬神的僕人應有緊急感,趁有今天,快來服事主.耶穌不但有緊急感,還有尊貴感.祂是神子.從父而來,還要回到榮耀裡去；也就是在祂最尊貴之時,成為最卑微的僕人.這是一般人難以了解的.世界一般情況,地位低者要服事地位高的；但在神家中全然不同,大的要服事小的,所以主為門徒洗腳；主知道,門徒不知道,可能我們也不知道.那麼,我們就不能作主工；因不明白主的心意.門徒當時的確不知道,所以西門彼得對主說:「主阿,你洗我的腳麼?」耶穌答:「我所作的,你如今不知道,後來必明白.」
\newline
\newline
很多時候,我們在神的工作上,在教會工作中,很多不清楚,不了解的,自以為知道；其實不知道；所以犯了很多錯誤,乃因我們不知道主的心意,不會作主的僕人.有次,耶穌醫治瞎眼的,問他看見了嗎?瞎子說:「本來甚麼都看不見,現在看見了；看見人好像樹木,還會走動.」他看見了,可是看的不對,不正確,更不徹底；因為人非樹木,且樹木不會行走.但是我們不能說他是瞎眼的.今天我們好像看見了,在屬靈方面看見,可是看得太膚淺,太幼稚,似是而非,不正確,不透徹,故此我們失敗了.
\newline
\newline
當耶穌基督作神的僕人時,大家覺非常驚奇,完全出乎人意料之外.
\newline
\newline
我們生在充滿世俗的社會環境,把世俗混入教會；自以為對,其實錯了；教會大部份人對屬靈的事非常錯誤.我們多麼需要神的憐憫!有的人以為熱心,每主日聚會,有奉獻,有事奉,自覺滿意；但若看得清楚一點,看見自己的軟弱,幼稚,冷淡,和污穢,實在太差了!主問你看見了嗎?我們常以為過得去就算了.
\newline
\newline
作僕人應該真正明白主人的意思,在約翰福音13章,耶穌只會作僕人,不全作領袖.教會領袖越多,意見越多,問提也越多.教會需要僕人,僕人越多,教會越興旺.耶穌親自替門徒洗腳,是真正屬靈的操練；當耶穌替門徒洗腳,他們就明白了.耶穌說:「知道又遵行的有福了.」教會很多工作都太過用人的方法,我們需要神的恩典,我們每個人都要做僕人.「執事」希臘文意即在灰塵當中.僕人一聽到主人命令立即快跑,灰塵揚起.執事是做事的.另外的意思就是要留心聽主人的意思.耶穌離棄尊榮,謙卑作僕人,人都不明白.「許多人因祂驚奇.」(14)「這樣祂必定洗淨許多國民；君王要向神閉口......」(15)耶穌是僕人,有無限的權威.我們須學習主的榜樣；祂是最尊貴卻成為最卑微的.教會的長執最尊貴；但是一定要成為最卑微的.神的僕人是最尊貴的；但一定要最卑微.你我都是神的僕人,要謙卑!讓神向我們施恩!傳憂患的信息,在憂患要把信息帶出來.「我們所傳的,(或作所傳與我們的)有誰信呢?耶和華的膀臀向誰顯露呢?」(五十三:1)「十字架的道理在滅亡的人是愚拙的,但在得救的人是神的大能.」有誰懂得呢?十架的道理有時成為絆腳石,有的人很難接受.「我們所傳的有誰信呢?」這是先知的嘆息,也是僕人的感嘆.當耶穌向尼哥底母傳時他不明白,他學問淵博,懂得猶太的宗教思想；耶穌說:「我對你們說地上的事,你尚且不信,若說天上的事如何能信呢?」(約三:12)傳福音的工作須有恆忍的態度.在憂患時代,人需要福音但可悲人拒絕了所需要的；因他們不明白真道,沒有聖靈的光照；需要我們禱告,行動；切實把福音帶給我們.我們需要仰望聖靈的能力,在憂患中傳福音,多麼需要神賜給我們智慧!
\newline
\newline
第五十二:13:「僕人行事必有智慧.」給我們想到但以理的第12章的真理.但以理很有智慧,使多人歸向神,好像天上的星明亮；因主賜給祂的僕人.傳福音很艱難,但若有神的同在,有聖靈的感動；有傳福音的熱心,神一定叫我們做同樣的工作；因為在這憂患的時代更需要福音.
\newline
\newline
培靈研經會已60年了,在這憂患時代,60年前中國軍伐政變,非常紛亂；繼之抗戰發生,事後一直有內亂,我們時在憂患中.感謝主!特別恩待我們!英國歷史家說,廿一世紀是中國人的世紀,屆時中國人口增加,對全世界影響很大.故此,主預備了華人信徒來工作.如果細讀舊約聖經,可體會以色列人經歷許多憂患；同時也在憂患中把福音帶出去.今天神給我們的托負也是這樣.
\newline
\newline
一個僕人「他在耶和華面前生長如嫩芽,像根出於乾地.」(五十三:2)嫩芽生命脆弱,需要水份,極難在乾地成長,使這是神的恩典.人若不經患難,不會知道福音的重要,不能體會恩典的奇妙,不會知道神的真理充滿神的心意.神恩待華人,恩待華人教會.華人教會有重大責任,因為很難成長.
\newline
\newline
先知書預言基督降世,生於憂患.當時社會極其混亂,希律下令殺害全以色列的嬰孩.在羅馬政權下沒有自由,在政治,社會,經濟,宗教,不安環境下,神叫祂的兒子降世；基督靠著神的恩典,在神保護,保守之下,在患難中生長.
\newline
\newline
我們需要成長,但在憂患中成長似乎不可能；因多方環境都於我們不利,然而我們在神面前仍然得以成長.我們有否靠著神的恩典,我們有多少時間真正在神面前?有個基督徒活到75歲,據統計,他在生活中,除去睡覺,工作,吃飯,穿衣及進洗手間等和坐車旅行......之外,只剩下半年的時間；一生75年只用半年敬拜神.
\newline
\newline
主耶穌在極度困難情況中成長,靠著神的恩典,祂以神的信實為糧,以神旨意當飯；成長直到各各他的十字架,在那裡祂成為生命樹,結出許多的子粒.生命樹必須經艱苦的過程；你我也必須成長,不一定有外面的美容；真正的美貌是聖潔.主是何等奇妙的救主!
\newline
\newline
我們想到神的恩典,體會神給我們的福份.今天到處需要福音,香港地人情淡薄,令人有冷落感；仔細研究,發現許多人充滿失敗,失敗的態度；對世界一切憤慨,對事持反面態度,這是病態的社會.不單香港如此,實在到處皆是.所有人對福音需要愈大.神的兒女在哪裡?福音使者在哪裡?怎能無動於衷?怎能裹足不前?怎能遲延?還有一點點時候,快做傳福音工作!很多華人移民各處,任何角落都有華人；深信神的旨意藉此傳開福音.巴黎有十萬華人,只有三個教會,數百信徒.南美,北美有很大需要；當然我們非好高務遠.就在我們住處,神若叫我們目前留此,作傳福音工作,時候也不多!趕快工作!神對我們有很大的計劃,有無限的心意；我們當在神前真正仰望倚靠祂,完全奉獻給祂.
\newline
\newline


\newpage

\section{第60屆~港九培靈會~唐佑之~第八講}
\label{sec:1229}
\textbf{捨上}
\newline
\newline
連結: \href{https://www.hkbibleconference.org/session-message/view/1229}{\texttt{https://www.hkbibleconference.org/session-message/view/1229}}
\newline
\newline
\hyperref[sec:1228]{< < < PREV SERMON < < <}
~
\hyperref[sec:toc]{[返目錄]}
~
\hyperref[sec:1231]{> > > NEXT SERMON > > >}
\newline
\newline
以賽亞書 53:1-53:12
\newline
\begin{longtable}{cl}
\hline
\hline
章節 & 經文 (和合本修訂版)\\
\hline
53:1 & \begin{tabularx}{0.7\textwidth}{X} 我們所傳的有誰信呢？耶和華的膀臂向誰顯露呢？ \end{tabularx} \\ \\ \relax
53:2 & \begin{tabularx}{0.7\textwidth}{X} 他在耶和華面前生長如嫩芽，像根出於乾地。他無佳形美容使我們注視他，也無美貌使我們仰慕他。 \end{tabularx} \\ \\ \relax
53:3 & \begin{tabularx}{0.7\textwidth}{X} 他被藐視，被人厭棄；多受痛苦，常經憂患。他被藐視，好像被人掩面不看的一樣，我們也不尊重他。 \end{tabularx} \\ \\ \relax
53:4 & \begin{tabularx}{0.7\textwidth}{X} 他誠然擔當我們的憂患，背負我們的痛苦；我們卻以為他受責罰，是被神擊打苦待。 \end{tabularx} \\ \\ \relax
53:5 & \begin{tabularx}{0.7\textwidth}{X} 他為我們的過犯受害，為我們的罪孽被壓傷。因他受的懲罰，我們得平安；因他受的鞭傷，我們得醫治。 \end{tabularx} \\ \\ \relax
53:6 & \begin{tabularx}{0.7\textwidth}{X} 我們都如羊走迷，各人偏行己路；耶和華使我們眾人的罪孽都歸在他身上。 \end{tabularx} \\ \\ \relax
53:7 & \begin{tabularx}{0.7\textwidth}{X} 他被欺壓受苦，卻不開口；他像羔羊被牽去宰殺，又像羊在剪毛的人手下無聲，他也是這樣不開口。 \end{tabularx} \\ \\ \relax
53:8 & \begin{tabularx}{0.7\textwidth}{X} 因受欺壓和審判，他被奪去，誰能想到他的世代呢？因為他從活人之地被剪除，為我百姓的罪過他被帶到死裡。 \end{tabularx} \\ \\ \relax
53:9 & \begin{tabularx}{0.7\textwidth}{X} 他雖然未行殘暴，口中也沒有詭詐，人還使他與惡人同穴，與財主同墓。 \end{tabularx} \\ \\ \relax
53:10 & \begin{tabularx}{0.7\textwidth}{X} 耶和華的旨意要壓傷他，使他受苦。當他的生命作為贖罪祭時，他必看見後裔，他的年日必然長久。耶和華所喜悅的事，必在他手中亨通。 \end{tabularx} \\ \\ \relax
53:11 & \begin{tabularx}{0.7\textwidth}{X} 因自己的勞苦，他必看見光就心滿意足。因自己的認識，我的義僕使許多人得稱為義，他要擔當他們的罪孽。 \end{tabularx} \\ \\ \relax
53:12 & \begin{tabularx}{0.7\textwidth}{X} 因此，我要使他與位大的同份，與強盛的均分擄物。因為他傾倒自己的生命，以致於死，也列在罪犯之中。他卻擔當多人的罪，為他們的過犯代求。 \end{tabularx} \\ \\
[1ex]
\hline
\hline
\end{longtable}
以賽亞書第五十三章,可說是舊約中最偉大的一章,如果沒有這章,聖經真理就不完全,如果信仰沒有十字架,我們的信仰就完全沒有內容.沒有耶穌基督,我們的信仰的對象是虛無的.
\newline
\newline
我們來到主前,思想這偉大的篇章,我想盡所能給各位解釋；但我越想越覺得自己在神面前是誰呢?怎可能講解這偉大的真理?我的思想有限,怎能講得明白?我的言語更有限,更不能表達這篇偉大的真理,實在需要聖靈的光照!讓我們靠聖靈的引導,一步步來思想如此偉大的真理,美妙的詩歌；深刻的意思吧.我們在主前像那瞎子,耶穌對他說:「我能為你作甚麼?瞎子說:「我要能看見.」我們需要看見,看見我們的救主,看見祂的十字架,看見祂受苦的奧秘,看見憂患的意義；清楚看見我們應有的福音使命；看得清楚徹底,明白,真正有屬靈的眼光.
\newline
\newline
我們所傳的有誰信呢?這是先知的感嘆.歷代許多神僕傳福音,遭遇種種困難；很多時候,福音沒有人接受,並非神沒有能力,也非真理不清楚,而是人沒接受的力量.多麼需要神開人的眼睛,首先讓神開我們的眼睛；要人接受,自己先要接受；要人相信,自己也先要相信.但是若我們信心方面沒實際的長進,相反地有時退後；特別對熟悉的聖經內容,十架道理的深奧奇妙,自以為知道,其實不然.這樣的福音怎能叫人相信呢?自己必須有堅強的信心,十架的道理,永是我的榮耀,不以福音為恥；福音的能力才能發展,福音信息才能傳開,真正成為神使用之僕.我們知道神的話有能力,但人的回應卻大有區別.有次我在某教會講道,有人希望我向位姊妹的丈夫傳道,領他信主；我先了解此人參加禮拜經已十年,為的是討太太的歡心.他雖明白許多聖經的道理,可是他不相信；我問他是否基督徒,他說:「一半,因為太太是基督徒.」我說基督徒沒有一半的,因為神只有兒女,沒有女婿,媳婦,也沒有孫子；惟個人接受神的救恩.神的話大有能力,但是這聽道十年的人還不信主.這使我體會一件事,好像太陽同樣晒著一塊泥和一塊牛油；泥土越晒越硬,因水份蒸發了,牛油一晒即溶.作福音使者,要真正仰望神,要付上代價；別讓聽福音者像泥土越晒越硬.基督徒也可能,有此危險,當蒙恩開始,熱心追求,每篇道都非常覺得寶貝,每篇信息都是新的；漸漸就覺得平常,更有甚者是淡然無味.若是這樣,非常危險!當他打開心門,會前,當主席的老先生為我禱告,求神賜給特別的信息是前所未有聽過的.我說我講的是聖經的話,都是你們已聽過的.他所以這樣禱告是因聽道太多麻木了.如果我們沒有堅強的信心,也不能幫助別人有堅強的信心.我們多麼需要謙卑的心來領受.
\newline
\newline
「我們務要認識耶和華,竭力追求認識祂......」(何六5)我們雖然認識神,但是太膚淺,太幼稚,太表面,認識不夠；我們當進一步認識神,然後帶著信心把信息傳揚；叫別人相信,說:「我因信所以如此說話.」
\newline
\newline
差傳意即先要肯定,神在基督裡向我們肯定,我們這些罪人靠主寶血得救；我們在基督裡也要向神肯定,我們在神前越肯定,在人前也越肯定.求主幫助我們常常溫習,思想十字架的道理.另一方面,對很多來信主的人；我們要有信心,有信心的行為和見證；充滿信心的態度來到神面前,耶和華的膀臂也會向凡信靠祂的人顯露.神大能的膀臂,要普及全世界,因神愛世人,愛教會,也愛我個人.「我已經與基督同釘十字架,現在活著的,不再是我,乃是基督在我裡面活著......」(加二:20)基督愛我,為我捨己,我個人蒙受神的恩典,主的愛為我們捨上,我們當感謝讚美主.「基督愛教會,為教會捨己.」(弗五:25)主不但愛我們,也愛教會,教會是主的寶血買贖來的；祂沒有捨上,教會就沒基督徒了.祂不捨上,我們今天在哪裡呢?神愛世人,將獨生子賜下.愛是「捨上」,「給與.」神沒有留下一樣的好處不給我們.這是五十三章教導我們的.
\newline
\newline
生命寶貴,神的生命更加寶貴!神創造人時,有生命樹,有分別善惡樹；但人不要生命,要善惡,所以人墮落了.在人類歷史中,說明了人墮落的狀況.直到主耶穌降世,就像嫩芽從乾地生長,雖然緩慢卻是壯健的；雖艱難,但繼續不斷長大,到了十字架上完全捨上.讓我們看見生命樹,並且有生命的果子；救恩白白的賜給我們,毋須付上任何代價.但神的兒子付了極重的代價,以致祂無佳形美容；憂患壓傷了祂,但是實際上祂裡面充滿聖潔,全然美麗,可羨慕的,我們得著祂十架的恩典；我們認識祂,才能欣賞祂的美貌.祂榮耀的美貌,反應在我們身上；使我們有基督的榮美,才能作基督的使者.摩西在西乃山上與神親密交通四十晝夜,當他下山,臉上發光而不自知.我們當在神前追求聖潔,多多親近祂；禱告,讀聖經；體會祂的同在.這樣一定可得神所給予我們的美.而後真正發揮神恩惠的功效.
\newline
\newline
我們的主是憂患之子,是受苦之僕；人犯罪以致連累祂.「他被藐視,好像被人掩面不看的一樣.」(五十三:3)這句話有三種翻譯,第一種如上述；人藐視,掩面不看他.第二,他自己不願給人看.因為祂的痛苦無人了解,人只有侮辱藐視祂；不明白祂迫切的心賜.第三,好像被神掩面不看一樣.祂是神子,為人擔當罪孽；神不忍看祂的兒子,神是公義的,不以有罪為無罪.當耶穌在十字架上,遍地黑暗；好像神在哀哭,不忍祂的兒子受苦.
\newline
\newline
弟兄姊妹!當我們來到主前,是否瞻仰祂的榮美?世人不重視,我們當重視,當仰望祂,思念祂的愛,祂的恩典,福份；我們充滿快樂,不能不俯伏敬拜；因祂是配得稱頌,榮耀,讚美和感謝的；祂願意我們看祂,因祂知道我們有恭敬,仰望的心意,有倚靠,順服的心.
\newline
\newline
神在天上為你,我,為地上人的靈魂痛苦,祂說:「人子來,不是受人服事,乃是要服事人；並且要捨命,作多人的贖價.」主在十字架說:「我的神!我的神!為甚麼離棄我?」因為神不看他,「祂誠然擔當我們的憂患,背負我們的痛苦；」可惜人不知道!「我們卻以為祂受責罰,被神擊打苦待了.」(4)聖經告訴我們,有罪就有苦難,有刑罰；然而耶穌是無罪的,祂卻站在罪人的地位上.在約伯記,當約伯受苦時,他的朋友來安慰他,反而更增加他的痛苦；他們認為約伯受苦是因犯罪.雖然不能說約伯完全清潔；但他覺得若將罪和苦,相題並論不一定完全對的.有次門徒問耶穌:「那人瞎眼是因他自己犯罪呢?還是他的父母犯罪?」耶穌說:「不是他犯罪,也不是他的父母犯罪；」在這裡使我們明白受苦的奧秘就是十字架的奧秘.若犯罪受苦,罪有應得；種的是甚麼,收的也是甚麼.耶穌受苦,乃因世人都犯了罪,虧缺了神的榮耀；祂本可不接受十字架,因祂沒有罪,祂是神的兒子,祂所以甘心捨上,是為擔當我們的罪.罪把痛苦帶來,可是苦難不一定表明有罪.當我們明白,耶穌受苦是為愛你,愛我的緣故,我們惟有低頭敬拜,抬頭仰望為我們受苦的主.
\newline
\newline
從第4節開始,我們看到很重要的字眼,「他」擔當「我們」的憂患,背負的痛苦.主為我們的過犯受責罰,為我們的罪孽壓傷,「祂」「我們」之中何等豐富的道理!我們是該死,該滅亡的人,在世上沒有指望；祂指望；祂是神子,是萬王之王,萬主之主；竟然成為我們的救主.祂受刑罰,我們得平安；祂受鞭傷,我們得醫治.多麼奇妙的恩典!多麼偉大的救恩!人若忽略這麼大的救恩,怎能逃罪呢?
\newline
\newline
我們當有堅強的信心,把福音帶出去；世界將要滅亡,我們若不起來傳福音；社會直趨毀滅.保羅說:「不傳福音就有禍了.」
\newline
\newline
因祂受的刑罰我們得平安,平安意即完整.未信主之前,我們的生命支離破碎；得救之後,生命,思想,情感,意志完整了.我們應該表裡一致,信心行為一致,個人與家庭生活一致,小家庭,大家庭,教會一致,是完整的平安,表明教會的平安.世界有苦難；但是耶穌說:「你們可以放心,因我已勝過世界.」「放心」亦即平安.保羅說:「讓基督的平安在我們心裡作主.」祂是和平之君,是我們的平安.
\newline
\newline
因祂受鞭傷我們得醫治,醫治即從病患中出來.當以色列人從埃及出來,他們在紅海岸邊唱耶和華的歌,讚美神奇妙的帶領,後來他們經三天的路程找不到水喝,他們就發怨言；摩西禱告；神指示他把樹放在苦水中,水即變甜,神對他們說:「因為我是醫治你們的.」不能喝的苦水變甜,就是醫治；祂治好我們的病成為健康,就是醫治.耶穌曾說:「健康的人用不著醫生,有病人的人才用得著,我來是找罪人.」祂的醫治是救恩的道理,叫我們的生命都改變了；甚至在我們艱苦的環境,在憂患之中,都感到主何等甜美!我們每天都需要醫治,需要神的憐憫.
\newline
\newline
「我們都如羊走迷,各人偏行己路；耶和華使我們眾人的罪孽都歸在他身上.」(6)這失落的世界,迷失的人到處皆是；但是神的愛尋找我們回來.
\newline
\newline
「他被欺壓......卻不開口......又像羊在剪毛的人手下無聲,他也是這樣不開口；」(7-8)這裡兩次說祂不開口.祂完全順服神的旨意,緘默代表祂的順服.在約伯記也題到這道理,當時神在旋風中向約伯說話；約伯用手掩口不敢講話,表明他順服.當耶穌受審,大祭司彼拉多問祂,祂都不開口；祂非恃神子自負而表現出不害怕的態度,也非帶著冷漠猶豫的態度；乃因祂完全順服父神的旨意.我們也當如此,若是出於神的,就當默然不語,在神面前順服.
\newline
\newline
有人受苦是自作自受,罪有應得,有人受苦是受了別人的連累；然而基督受苦,是甘心樂意為我們的罪而捨上,義的代替不義,絕對承受神的命令.基督徒不肯付福音的代價,只要十字架的道路.耶穌為我死,我能得生,耶穌流寶血我才能得潔淨.十字架的道路要犧牲,要捨上；許多基督徒都不願意.耶穌說若有人要跟從祂,就當背起十字架來跟從.保羅和西拉在腓立比傳道,被官府捉拿,挨打,下監牢；為福音受苦,但他們懂得,走十架的道路必須犧牲,完全獻上,身雖受苦猶覺甘甜.他們歌頌讚美主,從心坎裡唱出美妙的詩歌；眾囚犯都側耳而聽,忽然地大震動,監門全開,囚犯鎖鍊鬆開；獄卒大驚,俯伏請示當怎樣行才可得救?說:「當信主耶穌,你和你一家都必得救.」為了引人歸主,他們甘心受苦；基督救恩即成就了；我們受苦,救恩就傳開了.
\newline
\newline
傳福音工作困難重重,必須付上福音的代價!
\newline
\newline
受苦的主為我捨上,我捨上甚麼呢?
\newline
\newline


\newpage

\section{第60屆~港九培靈會~唐佑之~第九講}
\label{sec:1231}
\textbf{展望}
\newline
\newline
連結: \href{https://www.hkbibleconference.org/session-message/view/1231}{\texttt{https://www.hkbibleconference.org/session-message/view/1231}}
\newline
\newline
\hyperref[sec:1229]{< < < PREV SERMON < < <}
~
\hyperref[sec:toc]{[返目錄]}
~
\hyperref[sec:1233]{> > > NEXT SERMON > > >}
\newline
\newline
以賽亞書 54:1-54:17
\newline
\begin{longtable}{cl}
\hline
\hline
章節 & 經文 (和合本修訂版)\\
\hline
54:1 & \begin{tabularx}{0.7\textwidth}{X} 你這不懷孕、不生育的，要歡呼；你這未曾經過產難的，要歡呼，揚聲呼喊；因為被遺棄的婦人，比有丈夫的人兒女更多；這是耶和華說的。 \end{tabularx} \\ \\ \relax
54:2 & \begin{tabularx}{0.7\textwidth}{X} 要擴張你帳幕之地，伸展你居所的幔子，不要縮回；要放長你的繩子，堅固你的橛子。 \end{tabularx} \\ \\ \relax
54:3 & \begin{tabularx}{0.7\textwidth}{X} 因為你要向左向右開展，你的後裔必得列國為業，又使荒廢的城鎮有人居住。 \end{tabularx} \\ \\ \relax
54:4 & \begin{tabularx}{0.7\textwidth}{X} 不要懼怕，因你必不致蒙羞；不要抱愧，因你必不致受辱。你必忘記年輕時的羞愧，不再記得守寡的恥辱。 \end{tabularx} \\ \\ \relax
54:5 & \begin{tabularx}{0.7\textwidth}{X} 因為造你的是你的丈夫，萬軍之耶和華是他的名；救贖你的是以色列的聖者，他必稱為全地之神。 \end{tabularx} \\ \\ \relax
54:6 & \begin{tabularx}{0.7\textwidth}{X} 耶和華召你，如同召回心中憂傷遭遺棄的婦人，就是年輕時所娶被遺棄的妻子；這是你的神說的。 \end{tabularx} \\ \\ \relax
54:7 & \begin{tabularx}{0.7\textwidth}{X} 我離棄你不過片時，卻要大施憐憫將你尋回。 \end{tabularx} \\ \\ \relax
54:8 & \begin{tabularx}{0.7\textwidth}{X} 我因漲溢的怒氣，一時向你轉臉，但我要以永遠的慈愛憐憫你；這是耶和華－你的救贖主說的。 \end{tabularx} \\ \\ \relax
54:9 & \begin{tabularx}{0.7\textwidth}{X} 這事於我有如挪亞的洪水；我怎樣起誓不再使挪亞的洪水淹沒全地，也照樣起誓不再向你發怒，且不斥責你。 \end{tabularx} \\ \\ \relax
54:10 & \begin{tabularx}{0.7\textwidth}{X} 大山可以挪開，小山可以遷移，但我的慈愛必不離開你，我平安的約也不遷移；這是憐憫你的耶和華說的。 \end{tabularx} \\ \\ \relax
54:11 & \begin{tabularx}{0.7\textwidth}{X} 你這受困苦、被暴風捲走、不得憐憫的城，看哪，我必以灰泥來做你的石頭，以藍寶石立你的根基， \end{tabularx} \\ \\ \relax
54:12 & \begin{tabularx}{0.7\textwidth}{X} 又以紅寶石造你的女牆，以晶瑩的珠玉造你的城門，以珍貴的寶石造你四圍的邊界。 \end{tabularx} \\ \\ \relax
54:13 & \begin{tabularx}{0.7\textwidth}{X} 你的兒女都要領受耶和華的教導，你的兒女必大享平安。 \end{tabularx} \\ \\ \relax
54:14 & \begin{tabularx}{0.7\textwidth}{X} 你必因公義得堅立，必遠離欺壓，毫不懼怕；你必遠離驚嚇，驚嚇必不臨近你。 \end{tabularx} \\ \\ \relax
54:15 & \begin{tabularx}{0.7\textwidth}{X} 若有人攻擊你，這非出於我；凡攻擊你的，必因你仆倒。 \end{tabularx} \\ \\ \relax
54:16 & \begin{tabularx}{0.7\textwidth}{X} 看哪，我造了那吹炭火、打造合用兵器的鐵匠；我也造了那殘害人、行毀滅的人。 \end{tabularx} \\ \\ \relax
54:17 & \begin{tabularx}{0.7\textwidth}{X} 凡為攻擊你而造的兵器必無效用；在審判時興起用口舌攻擊你的，你必駁倒他。這是耶和華僕人的產業，是他們從我所得的義；這是耶和華說的。 \end{tabularx} \\ \\
[1ex]
\hline
\hline
\end{longtable}
以賽亞書五十三章歸納起來,受苦的僕人有三方面的重點:
\newline
\newline
(一)憂患之主為我們擔當痛苦.祂背負重擔,擔當罪孽.替代我們受苦,因此我們得平安,得醫治,得一切的恩惠.當祂願意承擔苦難,憂患之時,救恩就在我們身上發生功效；藉著我們叫很多人蒙恩.今天晚上,我們來敬拜為我們在十架上捨身流血的主,溫習祂的憂患,才能明白何謂憂患；我們敬拜仰望祂,就能明白苦難有時是替代的意思,願長承受憂患.
\newline
\newline
(二)受苦的僕人,緘默的主.祂願意順服父神,所以祂不願意說甚麼；因祂知道所信的是誰,順服的對象是誰,祂完全順服,甘願替代.
\newline
\newline
人若不甘願受苦,就怨天尤人.憂患是神托付我們的.正像神曾經托付祂的兒子一樣.事實上,受苦是種受托的經驗,每當我們講受托作管家,就想到自己是時間,金錢,福音,才幹的管家,我們都是受托的,無一不是領受來的,不值得自誇.可是很少聽見苦難也是受托的；不信主者絕不會明白,他們只會憤慨,墮落,自暴自棄；但是我們因為明白苦難是甚麼,十架的道理是甚麼.也知道受苦的奧秘,和憂患的意義；所以我們因受托而受苦.若你靈性程度不夠,神不會將苦難交給你.當我們長進,願意付代價時,我們就會因受托而受苦,把福音的能力表現出來.我們不要怕受苦,因聖經說:「萬事都互相效力,叫愛神的人得益處.」你若真正追求主道,真正愛主；你當在主前感謝,讚美.憂患是愛主者生命的經歷,我們需要十字架的功勞,都認為十字架很寶貝；但是我們都不願走十架的道路,因為要犧牲,要獻上一切給神,要付上代價.
\newline
\newline
(三)主耶穌願意付上代價,為要成就神的旨意.祂成就救人之恩,祂不但承受,且完全獻上.主受審問時不開口,並非傲慢,冷漠；也非懼怕.雖然祂是神子降世為人,有人的軟弱,懼怕；但祂果敢地承受十架的苦難,完全順服神；所以祂毋須解釋,毋須推辭；惟默然在神前順服,在人前受苦,把自己完全捨上.主是代罪羔羊「被牽到宰殺之地,又像羊在剪毛的人手下無聲,祂也是這樣不開口.」(7)
\newline
\newline
我們是主的小羊,羊毛太長,羊就懶惰,骯髒,墮落深坑,迷失了.主是我們的牧人「為自己的名引導我們走義路.」(詩23)祂要教我們走正直,正當,正確的路.今天在大都市,教會群羊體胖,毛長；所以神要對付,要管教,要剪羊毛；為使我們在祂面前,蒙受更多,更豐盛的恩典.神的恩典原是訓練,管教；我們常需要有此經驗.如果只有健康而沒有疾病；只有平安而沒有不安；只有豐富而沒有缺乏；只有歡樂而沒有痛苦；我們的靈性就會下降,沒有力量傳福音,以為經常到教會就自滿自足.
\newline
\newline
神藉耶利米說:「你們這些敬拜神......說,這是耶和華的殿,是耶和華的殿,是耶和華的殿.」連講三次,以為到聖殿敬拜就毫無問題.這是錯誤的思想,只有外表沒有徹底；只有信心沒有行為；只接受福音的好處,沒有接受福音的能力；落到這樣地步,多麼需要神的憐憫!我們必須在神前謙卑順服,軟弱時,讓神訓練,使我們明白祂的恩典；有病求神醫治,教我們體會神是醫治的神；缺乏時才體會神的豐富,有痛苦才體會憂患對基督徒靈命的重要,能得長進.阿拉伯語有云:「只有太陽沒有雨水,只能造成一片沙漠而無綠洲.」只有歡樂沒有痛苦,就會導致靈性軟弱.神為了造就我們,就叫我們去體會憂患.讓我們想到主的憂患,想到主的代價；而後我們也願意承受憂患,為了別人需要救恩願付上代價.
\newline
\newline
有位韓國傳道人,他小時候在北韓,那時基督徒受逼迫；他們把信者和不信者區分,然後用機關槍掃射信主者.很多人因此就否認了耶穌.那孩子要為主作見證,站在基督徒的一邊；但是被軍兵一腳踢開；因而蒙主保守；長大後,他傳信息千篇一律說:「福音,救恩是主用寶血成就的；所以福音,救恩也要用血來傳揚.」他真是付上福音的代價!教會歷史家塔古良說:「殉道者的血是教會的種子.」這句話常被人引用.若無殉道者,沒有人為主受苦受死；今天哪裡有你和我?我們都滅亡了,哪裡有教會?耶穌說:「你們就是我的見證.」見證原意是殉道者.
\newline
\newline
當我們讀五十三章,想到主完全捨上.「耶和華卻定意將他壓傷,使他受苦,耶和華以他為贖罪祭......」(五十三:10)贖罪祭就是獻上羔羊代罪,是舊約獻祭的條例.新約希伯來書第九章開始,說明耶穌基督是贖罪祭也是贖愆祭；意思即是我們自己不知有罪；耶穌卻為我們擔當了.
\newline
\newline
耶穌基督也獻上平安祭「因祂是我們的和睦,將兩下合而為一,拆毀了中間隔斷的牆；」(弗二:14)使我們與神也與人和好,把兩下歸為一體；本來猶太人和外邦人分開,但是在主裡面完全在一起.
\newline
\newline
耶穌基督也是素祭,表明祂樂意順服；祂也是燔祭,毫不保留,完全獻上.
\newline
\newline
贖罪祭,贖愆祭,平安祭,素祭,燔祭；耶穌基督都為我們獻上了；但是祂要我們跟隨祂的腳蹤行,叫我們繼續獻祭；並非祂獻的不夠,乃是要我們學習祂的榜樣,繼續為福音使命努力.
\newline
\newline
我們要平安祭,基督為我們受死,成就了和睦「又將勸人和他和好的職分賜給我.」傳福音就是和好的工作,素祭也是我們必須獻的.「甘心樂意的捐獻是主所喜悅的.」(林後九:7)奉獻意即甘心樂意；不勉強,不作難；積極地.這是對的,但是我們當獻額外的奉獻；燔祭是必要的,而素祭不是必要的；可是一般教會把林後八章所提的奉獻錯解為主要的奉獻.主要的奉獻是燔祭,是羅馬書十二章:「將身體獻上,當作活祭.」不為自己有所保留,因為一切都屬於神,出於神,歸給神.完全獻上,是全然也是必然的,無一人可以逃避.要承受福音工作,必須捨上自己；我們完全屬於祂,必須承認祂的主權,順服祂的旨意.雖然舊約獻祭已經過去,但現在仍然有效；耶穌基督都已為我們獻上.當然我們不配獻贖罪祭和贖愆祭；我們當獻燔祭,完全獻上自己；甘心樂意獻上素祭,獻平安祭,叫人與神和好；為福音而活,作福音的見證.
\newline
\newline
主耶穌為我們釘死在十架上,最後的話:「成了.」也有人說是「我將我的靈魂交在父的手裡.」我覺得「成了」是極重要的話；祂受苦,成就了救贖的功效.
\newline
\newline
無窮生命的大能,「他必看見後裔延長年日,耶和華喜悅的事,必在他手中亨通.」(10)有的人喜歡把五十三章分成段落:1-2節是祂的感嘆,3-6指祂的憂患.7-9指祂的受死,10-11指祂復活.主復活得勝了仇敵,魔鬼,罪惡.「擄物」即戰利品.主在十架上為罪犯代求:「將命傾倒,以致於死,」(12)祂完全捨上,成就救贖工作.」祂必看見自己勞苦的功效,便心滿意足.」(11)勞苦的功效有不同的翻譯:「祂必在光中看見勞苦的功效.」因為當祂受苦時,天地黑暗；但祂眼睛明亮,看見神的旨意,看到十架的功效.祂早就說過,一粒麥子若不落在地裡死了,仍是一粒,死了就能結出許多子粒來.愛能滿足,祂看見自己勞苦的功效便心滿意足,「他因那擺在前面的喜樂,就輕看羞辱,忍受了十架的苦難.」(來十二:2)如果我們真正愛人的靈魂,關懷他；當他願意接受主,那時的快樂實無法描述；有愛才有滿足.主愛我們,只要我們有一點回應,一點信心；一點依靠；一點順服；一點仰望；一點事奉；一點聖潔；主便得滿足.但是多少時候我們令主傷心!雖然我們幼稚,若果有一點回應；主非常欣賞.如果我們愛主,對主的恩典就感到滿足；若不愛主,常與人作比較,覺得主給別人的恩典多過自己.有時我們說:「萬事都互相效力,叫愛神的人得益處.」這要看愛神的程度如何；若愛的不夠,體會主的愛也不夠；我們若不能體會主的愛,就發怨言.主愛我們,愛是要勞苦,要付代價；愛不謹是獲得,也是給與的.聖經說:「我們愛,因為神先愛我們.」我們就明白,原來基督的愛這樣向我們顯明了.十字架充分說明了愛的意義,沒有愛,不會滿足；沒勞苦就沒有愛,再進一步,若無喜樂也不會勞苦.由此,我們看到主的滿足.但是,我們能否常讓主滿足,或偶然讓主滿足?但曾幾何時,我們又讓主傷懷,失望,嘆息.但願我們都能令主滿足!當主滿足,我們更加滿足.我覺得基督徒都不滿足；因我們沒滿足主的心,沒明白主滿足的心意.
\newline
\newline
第四十五章有兩首美妙的詩歌,提及主勞苦的功效,這位勞苦的僕人,付了勞苦的代價,常經憂患,多受痛苦；因此產生了勞苦的功效.「要發聲歌唱,揚聲歡呼；因為......兒女更多.」(1)意即以色列民必定復興.若把這節經文應用在今日的環境；相信這是講及教會復興,很多人要蒙恩得救.神創造世界,生養眾多,是神的計劃和目的.有的人認為實行家庭計劃不合聖經真理.我在北美還鼓勵弟兄姊妹多生育孩子,多受宗教教育；因為需要更多的見證人.我並非說家庭計劃不需要；事實上,是因為能源缺乏,糧食問題,教育等多方的問題.但新加坡,美國因節育關係,老者老,幼者幼；中間幾代接不上來.總之最要緊者,我們能從神的賜福,想到神的教會的蒙福.舊約神所賜的福有三方面:
\newline
\newline
第一,多子多孫,若無孩子就是神的咒詛.例如撒母耳的母親也希望生孩子.
\newline
\newline
第二,神的賜福在於家道豐富.以賽亞書五十三章預言耶穌與人同埋,與財主同葬.財主與惡人是同樣意思；舊約大部份是這樣說,並非說有錢財就表明有罪惡；其實財富是神所賜的福.相反地,貧窮是因為懶惰不工作.
\newline
\newline
第三,著仇敵的城門,意即和平安居樂業.我想,這三點當應用在教會中:
\newline
\newline
第一,教會要增長,多子多孫；使多人歸向神；數量增加質量增進.
\newline
\newline
第二,家道豐富,五十四:11開始提及很多寶物.
\newline
\newline
第三,得神的產業.「你的後裔必得多國為業.」(3)「這是耶和華僕人的產業.」(17)以屬物質的產業解釋屬靈的事；教會是神的產業.我們須爭戰,得地業.我們從事福音爭戰,教會屬靈方面豐富,有很多產業.我們靠主得勝,得著和平.
\newline
\newline
神如何賜福以色列民,神也賜福我們的教會；因福音的緣故,得救人數天天增加,教會興旺,屬靈方面非常豐富.因福音的緣故,魔鬼敗退；聖靈在教會中照常顯大,因為基督是我們的主,我們的頭.
\newline
\newline
「要擴張你帳幕之地,張大你居所的幔子不要限止,要放長你的繩子,堅固你的橛子.」(2)橛子若堅固帳幕就不致倒塌；繩子放長,帳幕擴大,是教會復興的現象；增長的秘訣.必須擴大屬靈的容量；讓神的恩典更豐富賜下.這裡給我們看見差傳宣教的異象「向左向右開展!得多國為業!......」(3)願神賜福!
\newline
\newline
華人遍佈全球各地,有不少基督徒；但教會還不夠,宣教士不夠,傳道人更少.多麼需要神呼召我們!留在香港的,神叫我們在香港作見證；神可能差我們到遠方,無論近處,遠處,我們必須把福音傳揚.時代的憂患,我們的福音使命不能逃脫,必須承受,捨上,然後才能看見遠象,看見異象.天國福音要傳遍普天下,以往是西方人,現在是東方人.我常覺得神在歷史文化環境中良好地預備我們,我發現中國人學習語言迅速；語言家認為西方人舌頭太圓,日本人舌頭太方,所以舌頭轉動不靈.中國人的舌頭不圓又不方,有利於學習任何語言.西方人常有民族優越感；我們卻有民族自卑感.基督徒萬勿自卑,但要謙卑.非洲人歡迎黃種人向他們傳道而不歡迎白種人,甚至討厭白種人；因為我們的樣子和他們差不多,所以我們傳福音特別有效.現在時機已到,求主開我們的眼睛,向左向右都能看見而且開展,我們應有福音隊伍勇往向前且向上.華人生活條件不高,較易適應我們宣教的條件實在太好!我們肯去嗎?神對以賽亞說:「我可以差遣誰呢,誰肯為我去呢?」你能否對主說:「我就在這裡,看我!」
\newline
\newline
請各位注意以下經文表達之意:──「這是耶和華說的.」(6)「這是耶和華你的救贖主說的.」(8)「這是憐恤你的耶和華說的.」(10)「這是耶和華說的.」(17)
\newline
\newline
平常讀聖經,都不覺得以上的句子有何特別；無論任何文字都沒法翻譯出其意.「說」並非一般說話的說,應是微聲細語.在先知的文字著作中,常有「耶和華如此說.」表明是嚴整的話語,也是很嚴重很嚴肅的事.是大聲的宣告,可是信息來了的「這是耶和華說的」,乃是用低微輕柔的聲音說的；表明神的安慰,如同父母安慰兒女.神盼望我們受感動,對祂有回應,對祂有奉獻的心.
\newline
\newline
讓我們承受憂患,把自己完全捨上!
\newline
\newline
當我們安靜之時,就能聽見神的聲音.
\newline
\newline


\newpage

\section{第60屆~港九培靈會~唐佑之~第十講}
\label{sec:1233}
\textbf{歌頌}
\newline
\newline
連結: \href{https://www.hkbibleconference.org/session-message/view/1233}{\texttt{https://www.hkbibleconference.org/session-message/view/1233}}
\newline
\newline
\hyperref[sec:1231]{< < < PREV SERMON < < <}
~
\hyperref[sec:toc]{[返目錄]}
~
\hyperref[sec:1389]{> > > NEXT SERMON > > >}
\newline
\newline
以賽亞書 55:1-55:13
\newline
\begin{longtable}{cl}
\hline
\hline
章節 & 經文 (和合本修訂版)\\
\hline
55:1 & \begin{tabularx}{0.7\textwidth}{X} 來！你們所有乾渴的，都當來到水邊；沒有銀錢的也可以來。你們都來，買了吃；不用銀錢，不付代價，就可買酒和奶。 \end{tabularx} \\ \\ \relax
55:2 & \begin{tabularx}{0.7\textwidth}{X} 你們為何花錢買那不是食物的東西，用勞碌得來的買那無法使人飽足的呢？你們要留意聽從我的話，就能吃那美物，得享肥甘，心中喜樂。 \end{tabularx} \\ \\ \relax
55:3 & \begin{tabularx}{0.7\textwidth}{X} 當側耳而聽，來到我這裡；要聽，就必存活。我要與你們立永約，就是應許給大衛那可靠的慈愛。 \end{tabularx} \\ \\ \relax
55:4 & \begin{tabularx}{0.7\textwidth}{X} 看哪，我已立他作萬民的見證，立他作萬民的君王和發令者。 \end{tabularx} \\ \\ \relax
55:5 & \begin{tabularx}{0.7\textwidth}{X} 看哪，你要召集素不認識的國民，素不認識的國民要奔向你；這都因耶和華---你的神，因以色列的聖者已經榮耀了你。 \end{tabularx} \\ \\ \relax
55:6 & \begin{tabularx}{0.7\textwidth}{X} 當趁耶和華可尋找的時候尋找他，在他接近的時候求告他。 \end{tabularx} \\ \\ \relax
55:7 & \begin{tabularx}{0.7\textwidth}{X} 惡人當離棄自己的道路，不義的人應除掉自己的意念。歸向耶和華，耶和華就必憐憫他；當歸向我們的神，因為他必廣行赦免。 \end{tabularx} \\ \\ \relax
55:8 & \begin{tabularx}{0.7\textwidth}{X} 我的意念非同你們的意念，我的道路非同你們的道路。這是耶和華說的。 \end{tabularx} \\ \\ \relax
55:9 & \begin{tabularx}{0.7\textwidth}{X} 天怎樣高過地，照樣，我的道路高過你們的道路，我的意念高過你們的意念。 \end{tabularx} \\ \\ \relax
55:10 & \begin{tabularx}{0.7\textwidth}{X} 雨雪從天而降，並不返回，卻要滋潤土地，使地面發芽結實，使撒種的有種，使要吃的有糧。 \end{tabularx} \\ \\ \relax
55:11 & \begin{tabularx}{0.7\textwidth}{X} 我口所出的話也必如此，絕不徒然返回，卻要成就我的旨意，達成我差它的目的。 \end{tabularx} \\ \\ \relax
55:12 & \begin{tabularx}{0.7\textwidth}{X} 你們必歡歡喜喜出來，平平安安蒙引導。大山小山必在你們面前歡呼，田野的樹木也都拍掌。 \end{tabularx} \\ \\ \relax
55:13 & \begin{tabularx}{0.7\textwidth}{X} 松樹長出，代替荊棘；番石榴長出，代替蒺藜。這要為耶和華留名，作為永不磨滅的證據。 \end{tabularx} \\ \\
[1ex]
\hline
\hline
\end{longtable}
當我們充滿感謝與讚美時,我們實在要歌頌神.以賽亞書五十五,是兩首歌頌的詩:內容是信息與見證.我們有信息必要宣告；有神跡應該見證.
\newline
\newline
要發出福音的呼喊,第一首詩歌,(1-9)歌頌一個很重要的信息,從中我們能聽見福音的呼喊,亦即呼喊邀請.
\newline
\newline
教會要發出福音的呼喊:「一切乾渴的都當就近水來.」福音的好處白白可得；不但有水,也有酒和奶；水象徵生命的滿足,耶穌說人若喝這水就永遠不渴.主給我們這恩典,我們得了也要分給別人.酒象徵生命的喜樂.我們得到救恩就唱出得救的樂歌,充滿喜樂；因為主有奇妙的恩典,甚至在憂患中,我們仍然喜樂；如同主耶穌,祂看擺在面前的喜樂就輕看羞辱；忍受或承受十字架的苦難.你我也當有這樣的喜樂,舉起救恩之杯歌頌神,在神面前有滿足的喜樂.乃象徵生命的滋養,使我們生長健壯.
\newline
\newline
這章聖經很古老,但是先知的信息沒淹埋在歷史廢墟中,反而在時代的走廊發出巨響；在每時代,每角落,每個人心中發生共鳴.今天晚上,我們讀這段聖經,應當充滿感謝；我們如今成了何等樣的人,得著神的恩典何其大.我們要感謝歌頌祂,把歌頌的信息帶出去,這就是我們福音的使命.時代嚴重有很多憂患,人需要生命的滿足,喜樂,滋養；我們只須宣告,毋須付任何代價；因主已付上重大的代價.我們當感謝讚美祂,求主憐憫；叫我們不僅蒙恩,也應感謝,更要報恩.
\newline
\newline
「我必與你們立永約,就是應許大衛那可靠的恩典.」(3)節經文非常重要,若我們片面地讀,未免太可惜!神說祂用永遠的愛愛我們,為我們立永約.在聖經中,有神顯明的歷史；從縱切面來看,我們知道歷史的事實；從橫切面看,我們知道歷史的意義.神在歷史之中,也在歷史之上,藉著歷史將祂的公義表明.於是,我們想到歷史時,特別是想到歷史的意義時；「約」是中心的思想,聖經是永遠的約；是神發動的,是神的旨意,是恩惠的約；叫我們蒙受恩惠.聖經中最明顯的約,第一是當挪亞甫出方舟,首先建立祭壇敬拜神；神對他說話,與他立約,不再把洪水降在地上.(創九)神的公義有更多慈愛；陰雨過去,天虹出現,雲彩中太陽脫穎而出；神的恩惠在銀色的邊緣顯明了.多麼奇妙恩惠的約!虹是多采多姿五色繽紛的；我們在神前蒙受的恩典也是如此.「各樣美善的恩賜；和各樣全備的賞賜,都是從父那裡降下來的.」(雅一:17)光若分開就成為不同顏色的虹,如同百般奇妙的恩典.我們想到挪亞的約,就想到神的恩典.
\newline
\newline
當以色列民出埃及到了西乃曠野,神藉摩西向他們曉諭:若聽從神的話遵守神的約,就在萬民中作神的子民；否則得不著神的恩典.恩典本無條件,因為憑自己不能作甚麼；但另方面恩典是有條件的；若不信靠,不信服,不仰望神,就得不到祂的恩典.
\newline
\newline
我們又想大衛的約.神應許大衛的家和國,必在祂面前永遠堅立；國位也必堅定,國位也必有堅定,直到永遠.以色列的歷史,猶太南國的王,無一不出於大衛的家；甚至被擄之後,仍然是大衛家的.神的信實何等廣大,於是我們明白,舊約彌賽亞預言的基礎.彌賽亞是出自大衛家的,是榮耀,顯赫,公義的君王；人都要在祂面前俯伏敬拜.這是本書一至卅九章的內容.四十章開始,這位彌賽亞不是榮耀的君王,而是受苦的僕人；卑微,憂患,被藐視,被忽略；這就是十字架的奧秘.耶穌必須受死,我們才能活；祂受苦是替代人的罪,祂在神前完全順服.祂的受苦是甘願的受苦,在任何情況之下祂保持緘默.主的苦難也是為救贖而受苦,主為我們死了然後救恩才能成就.在四十至五十五章,我們了解這位彌賽亞就是我們的基督.如此為我們受苦受死,我們才明白甚麼叫做憂患；才知道應當作甚麼.我們是受托者,是受苦的管家；我們受苦使別人得益處,叫人藉我們得到神的恩典；這就是福音使命的內容,也是應該真正有生活的見證.
\newline
\newline
大衛的約之後,還有一個約.「耶和華說日子將到,我要與以色列家和猶太家另立新約.」(耶卅一:31)這可說是新約前的新約；舊約中的新約.此約有三重點:第一,直接認識神.第二,敬畏神的恩典.第三,神的話語在我們裡面.這也是約翰一書的內容,屬靈的知識,屬靈的相信,屬靈的潔淨.於是我們知道新約是甚麼,當主耶穌被賣那一夜,祂拿起杯來,說:「這是我用血立的新約.」神的計劃是有連續性的,是永遠的約.神與我們立約,我們當中守於約；別人因我們向神的約而得到「約」的恩典──「可靠的恩典.」(3)這裡的恩典是多數的,希伯來文多數的字眼有很多的意思,不單指量也指質.我想是指諸般各樣的奇妙恩典；比這個更多,表明極其奇妙,超乎我們所求所想的.是神為祂所愛的人預備的,是我們眼所未曾見,耳所未曾聽,心所未曾想過的；是可靠的恩典.「可靠」這字在舊約聖經很多,有時是憐憫,有時是恩慈,有時是誠實.「你的誠實,極其廣大.」(哀三:23)就是這裡可靠的意思；我們縱然失信,但神是信實的.神的信實多麼奇妙!
\newline
\newline
第一個福音呼喊之聲,我們要邀請人來得主的恩典,我們既已得,也當報知別人來；凡蒙恩得救的人,在時代憂患中,應作福音的見證!
\newline
\newline
第二個福音呼喊之聲「趁耶和華可尋找的時候尋找他,相近的時候求告他!」(6)在古抄本中,不是說「時候」乃是說「地點.」我想,兩者可以放在一起；時候會過去的,但教會是基督的身體,是我們的家庭；在那裡可尋找耶和華,在你我的身上可以尋找耶和華,這就是福音的見證.尋求主,求告主名的,就必得救,這是福音的應許.
\newline
\newline
第三個福音呼喊之聲「惡人當離棄自己的道路,不義的人當除掉自己的意念,歸向耶和華,耶和華就必憐恤他.當歸向神,因為神必廣行赦免.」(7)我們希望別人離開罪,自己先要離開罪；希望別人得赦免,我必要在神前求赦免的恩典.好多時候我們失敗!在歷史中有次大復興,有位名人叫芬尼,他引用聖經:「要開墾荒地......公義如雨水在你們身上.」(十:12)意即我們當靠神的恩典,但自己要徹底認罪悔改,讓聖靈光照；我們需要十架寶血來潔淨,叫我們清潔合乎主用,預備行各樣善事.芬尼題及亞干犯的罪,(作不應該做的)和省略的罪,應作而不作)他特別注意基督徒很多省略的罪；應傳的不傳,應作見證不作；應該事奉,服事都不去作.實在太多省略的罪!以致我們的心田荒廢；雖有恩雨降下,但潮濕不過只在表面,沒有滲透泥土.
\newline
\newline
「來罷,我們歸向耶和華,祂撕裂我們,也必醫治,祂打傷我們,也必纏裡.」(何六:1)神要管教,因為審判是從神的家起首；如果歸向神,祂必出現如晨光.
\newline
\newline
神沒有離開我們,乃是我們離開神.耶穌講浪子回頭的比喻(路十五)浪子離開父親背道而馳,越走越遠,墮落越深；當他醒回轉向父家,他越走越近.太陽本身是不動的；但地球轉離太陽,黑夜便來臨.主一直等待,憐憫我們,期望我們切實悔改,回轉歸向神,而後神就憐恤我們.「當歸向神,神必廣行赦免.」(7末句)
\newline
\newline
「耶和華說,我的意念,非同你們的意念,我的道路,非同你們的道路.」(9)神進一步叫我們體驗祂的恩典和旨意.「意念」不單指神的思想,也包括神的計劃和目的；在時代,社會,家庭,教會,個人的歷史中,神都有設計.我們等候祂,祂就向我們顯明祂的旨意；叫我們蒙福,蒙恩,而後把恩典帶出來.這麼簡單,重要,奇妙!我們必須在神前切實悔改,切實順服,不是憑自己的思想和判斷,乃是照著神的意念,行神的道路.有人說,人的盡頭是神的起頭.當人陷於絕境,神顯明祂的作為.我覺得這話適合於未信者,因他們不知道倚靠神,等走投無路時才想到神.基督徒應當以「神的開始就是我們的開頭」才對,別等到盡頭才讓神開頭；以致失去了許多的福份.凡事,凡物都本於祂,倚靠祂；因為祂是我們信心創始成終的主.我們要完全倚靠祂,特別在時代憂患中,我們不知何往,進退維谷；然而主必定指引我們的路,叫我們別倚靠自己的聰明.主是我們的好牧人,祂在前面走,羊跟從祂.求主憐憫,叫我們跟隨主的腳蹤行.
\newline
\newline
第二首歌頌的詩,可說從10節開始；未講之前我先補充,第8節「耶和華說」這句話應該放在第9節末了.舊約聖經大概有360多次「這是耶和華說的,」在耶利米書最多,有120多次.這句話即微聲細語,帶著無限的憐憫,安慰,同情,關懷；因祂的溫和柔聲,叫我偉大,長進,尊崇.這是大衛和詩人的經驗,是主的恩典.
\newline
\newline
「雨雪從天而降,並不返回,卻滋潤土地,使地上發芽結實!」(10)「看哪,弟兄和睦同居,是何等的善,何等的美.......好比黑門的甘露,」(詩一三三)這首詩是雨和雪最好的解釋.黑門山很高,終年積雪,當和暖的陽光普照,雪隨春風吹拂,滋潤黑門的農田.我們要開懇荒地,必須在神前完全破碎,然後神供應恩雨,滋潤我們的心,浸透了神的一切恩典,結出聖靈的果子.神藉先知所說的話,絕不徒然返回,成就主所喜悅的.(11)
\newline
\newline
「你們必歡歡喜喜而出來,平平安安蒙引導；」(12)這是以色列人第二次出埃及.第一次出埃及,從為奴之地出來；第二次從巴比倫罪惡之地出來；神引導他們,向他們行神蹟奇事.「松樹長出代替荊棘；番石榴長出,代替蒺藜.」(13)當摩西蒙召時,看見燒著的荊棘而沒有燒毀.神是烈火,荊棘理應燒毀.以色列民本應滅亡；但神奇妙的愛在十架上表明.於是我們明白,沒用的,該除掉的,該燒盡的都以松樹和番石榴代替了.神無限奇妙的愛叫我們得厚恩；從自然領域看見恩典的境界,「大山小山必在你們面前發聲歌唱；田野的樹木也都拍掌.」(12)大山可以挪開,小山可以遷移；但神的慈愛必不離開你.
\newline
\newline
我們這些蒙恩者為何緘默?我們應該歌唱,隨時隨地歌唱得救的樂歌,福音的信息,十架的真理；從教會裡唱到教會外,從黑夜唱到天亮.教會在黑夜中見燈台,我們是黑夜的守望者；天快亮了,主快再來；榮耀的黎明快將來到,雖夜間充滿憂患痛苦；神使人在夜間歌唱.人生短暫,必須及時歌唱；從地上唱到天上,不斷歌頌.在羔羊的寶座邊,我們將看見很多主的寶血所救贖的人.
\newline
\newline
福音見證的行列,主復活之後,有更多人參加福音的行列,每個時代；尤其憂患時代苦難的環境中,更多人參加,其中有你和我嗎?但願今晚對主說:「主啊!感謝你,我要歌頌你；我一定要參加福音的行列.」
\newline
\newline


\newpage

\section{第60屆~港九培靈會~楊濬哲~序}
\label{sec:1389}
\textbf{序}
\newline
\newline
連結: \href{https://www.hkbibleconference.org/session-message/view/1389}{\texttt{https://www.hkbibleconference.org/session-message/view/1389}}
\newline
\newline
\hyperref[sec:1233]{< < < PREV SERMON < < <}
~
\hyperref[sec:toc]{[返目錄]}
~
\hyperref[sec:1172]{> > > NEXT SERMON > > >}
\newline
\newline
去年是港九培靈研經會六十週年紀念,這是一個非常值得感恩的慶典.
\newline
\newline
多謝九龍城浸信會,歷年來均借出堂址為會場,使大會能於八月初旬,逐年順利舉行.感謝神!賜福給每一屆的大會,都能夠如期開會,圓滿閉會.
\newline
\newline
憶六十屆大會為表示向神感恩起見,首末兩晚均蒙城浸及欣樂聯合詩班,獻唱多首大型聖詩.並蒙黃永熙博士,張美萍女士分別擔任指揮.其餘每晚均有合唱或獨唱.計有互愛團契弟兄,陳嘉璐女醫生,費明儀女士,招梁碧冕女士,何金秀莉女士,趙麗儀女士,楊淑貞小姐,陳之霞女士,張蓮女士等.歌聲令聚會倍增喜樂,藉他們嘴唇的讚美,成為佳美的祭品,蒙神喜悅!
\newline
\newline
香港聖經公會,特為大會印製一萬本紀念版聖經,以五折優惠發售.大會表示深切謝意.
\newline
\newline
本講道集乃去年之講道記錄.研經會由陳終道牧師主講,總題:「存到永遠.」引述聖經中多位名人之事蹟勉眾.講道會由劉承業牧師主講,總題:屬靈的戰爭.(從約書亞記看一個人和教會得勝的生活,和事奉的原則)奮興會由唐佑之牧師主講.總題:時代憂患與福音使命.(以賽亞書40-55章之研讀)
\newline
\newline
本講道集蒙林小娟姊妹,游偉業弟兄,朱雷麗端姊妹記錄；朱建磯牧師編輯,林小娟姊妹校對.謹此致謝!
\newline
\newline
在大會結束之感恩會中,曾題到幾件值得感恩的事；
\newline
\newline
(1)陳終道牧師於大會期間曾感冒咳嗽,蒙主醫治如常講道.
\newline
\newline
(2)聖經公會特為大會印製之一萬本紀念版聖經,全部售罄.
\newline
\newline
(3)大會印製一萬本六十週年紀念特刊,全部送完.
\newline
\newline
(4)大會費用需款二十四萬元,全部足用.
\newline
\newline
(5)赴會統計達六萬人次.
\newline
\newline
(6)為主而活與全時間奉獻之青年,多達五百八十餘人.
\newline
\newline
(7)赴會者除市區外,遠道而來者,有來自香港仔,筲箕灣,離島,屯門,荃灣,元朗,西貢等地.足見信徒愛慕主道的心,何等迫切!
\newline
\newline
每屆大會均蒙九龍城浸信會,同公義工,助道會,各兄姊協助招待,維持秩序.福音傳播中心同工,寄發宣傳品,現場錄音錄影.基督教週報刊登新聞稿及廣告,廣事宣傳,藉收效益,蒙各堂兄姊踴躍赴會並奉獻,統此致謝!
\newline
\newline


\newpage

\section{第60屆~港九培靈會~陳終道~第一講}
\label{sec:1172}
\textbf{第一個殉道士亞伯}
\newline
\newline
連結: \href{https://www.hkbibleconference.org/session-message/view/1172}{\texttt{https://www.hkbibleconference.org/session-message/view/1172}}
\newline
\newline
\hyperref[sec:1389]{< < < PREV SERMON < < <}
~
\hyperref[sec:toc]{[返目錄]}
~
\hyperref[sec:1173]{> > > NEXT SERMON > > >}
\newline
\newline
創世記 4:1-4:8
\newline
\begin{longtable}{cl}
\hline
\hline
章節 & 經文 (和合本修訂版)\\
\hline
4:1 & \begin{tabularx}{0.7\textwidth}{X} 那人和他妻子夏娃同房，夏娃就懷孕，生了該隱 ，她說：「我靠耶和華得了一個男的。」 \end{tabularx} \\ \\ \relax
4:2 & \begin{tabularx}{0.7\textwidth}{X} 她又生了該隱的弟弟亞伯。亞伯是牧羊的；該隱是耕地的。 \end{tabularx} \\ \\ \relax
4:3 & \begin{tabularx}{0.7\textwidth}{X} 過了一些日子，該隱拿地裡的出產為供物獻給耶和華； \end{tabularx} \\ \\ \relax
4:4 & \begin{tabularx}{0.7\textwidth}{X} 亞伯也把他羊群中頭生的和羊的脂肪獻上。耶和華看中了亞伯和他的供物， \end{tabularx} \\ \\ \relax
4:5 & \begin{tabularx}{0.7\textwidth}{X} 卻看不中該隱和他的供物。該隱就非常生氣，沉下臉來。 \end{tabularx} \\ \\ \relax
4:6 & \begin{tabularx}{0.7\textwidth}{X} 耶和華對該隱說：「你為甚麼生氣呢？你為甚麼沉下臉來呢？ \end{tabularx} \\ \\ \relax
4:7 & \begin{tabularx}{0.7\textwidth}{X} 你若做得對，豈不仰起頭來嗎？你若做得不對，罪就伏在門前。它想要控制你，你卻要制伏它。」 \end{tabularx} \\ \\ \relax
4:8 & \begin{tabularx}{0.7\textwidth}{X} 該隱與他弟弟亞伯說話 。 二人正在田間時，該隱起來攻擊他弟弟亞伯，把他殺了。 \end{tabularx} \\ \\
[1ex]
\hline
\hline
\end{longtable}
聖經上有幾處經文談到亞伯,除(創四1-8)外,還有:(太廿三34-35,來十一4,十二24)都有提及.單看創世記,難以明白亞伯原來是個殉道者,但在馬太福音,耶穌清楚表示亞伯是一位殉道者；因為他為義而受逼迫,故被列為殉道者.
\newline
\newline
希伯來書十一章可說是「信心英雄表」；其中羅列了很多的信心偉人,亞伯則排於首位,他因為所獻的祭,比該隱的更美,被嫉妒殺害了.但是,亞伯所流的血至今仍舊說話.現在,就讓我們看第一位殉道士,這信心偉人亞伯所說的究竟是甚麼,以及他的一生所要闡明的道理.
\newline
\newline
首先,亞伯的一生表明了:人非有信不能得神的喜悅.亞伯獻祭是因著信,就在他所獻的祭物上,神為他作見證.神只悅納亞伯的祭物,而不悅納該隱的祭物,其實是有原因的.
\newline
\newline
亞伯因著信而獻祭,他所信的就是為人贖罪的羔羊.(創三章)提及亞當犯罪後,用無花果樹葉編衣蔽體,但是卻不能遮蓋他的羞恥.於是,神殺了牲畜,以牠的皮做衣服,遮蓋了亞當的羞恥.亞伯信那要來的羔羊,祂可以擔當人的罪,祂的義像衣服,可以遮蓋人的羞恥.利未記摩西所宣佈的律法,指明所有祭物都要用頭生羔羊的脂油.可見,亞伯所獻的是合乎神的心意,又是符合律法的原則.事實上,並不是牛羊能為人贖罪,而是牛羊所象徵的基督才能贖我們的罪.施洗約翰介紹耶穌是神的羔羊,背負世人罪孽的；猶太人的錯誤,就是把牛羊當作真正能贖罪的.但不明白,只有牛羊所象徵的耶穌基督,惟有祂才能贖我們的罪.
\newline
\newline
第二,亞伯的殉道說明了兩種救法.神悅納亞伯用羔羊獻的祭,是要表明祂將要完成的救法.舊約中,獻祭時用的祭牲要符合很多的條件,諸如頭生的,沒有殘疾的,沒有受傷的......這樣嚴格的獻祭,表明了要求的基督,祂是完全沒有瑕疵的,只有祂才有資格承擔人的罪.
\newline
\newline
有一回我赴宴,主人家嗜吃肥豬肉,他很好客,不停地將肥肉送在我面前；忽略了我是不吃肥豬肉的,他雖然好心一片,卻得了反效果.亞伯獻的祭是神所悅納的,而該隱所獻的則相反,他以地上的出產獻予神.亞伯所獻的是流血的祭,這正合乎救恩的原則；若不流血,罪就不得赦免.罪要先受刑罰.才可得赦免.神以那羔羊來替我們受刑罰,那羔羊就是基督,除祂以外別無拯救,因為在天下人間,沒有賜下別的名,我們可靠著得救!祂有資格為全人類承擔罪的刑罰,因祂能滿足神公義的要求.
\newline
\newline
該隱獻的是素祭,在舊約,素祭是要與其他有血的祭物同獻上的.故此該隱的祭物不得神悅納；他只憑自己喜歡的而獻.不少拜菩薩的人也是為著求自己的利益,黑社會壞份子在行事前也會先求關帝保佑.這樣,他們的神只不過是被利用以達到人的目的,該隱正表明了這種功利的思想.
\newline
\newline
第三,亞伯及該隱的獻祭,又表明兩種不同的信心及行為.亞伯因信被神稱為義,不是因為行為被稱義；該隱欲以自己的行為來討神喜悅.結果,亞伯的祭物得神悅納,但卻被兄長所殺害.後來,神追討該隱的罪,照著該隱所作的而審判他.靠行為稱義的,結果沒有好行為；真正因信心領受神的生命,就有傾向善的能力.基督徒裡面有神的生命,無論受多大打擊,仍能充滿喜樂.
\newline
\newline
第四,亞伯的生平說明了兩種事奉.該隱也作屬靈的事,但這些事與事奉無關；他的獻祭是為了表揚自己,不是求討神的喜悅.因此,亞伯和該隱正代表了兩種不同的人生觀.一種是討神喜悅的人生,祂為我們死是叫我們活著的人不再為自己活；而是為那替我們死而復活的主而活,這是人生的意義.世上的情慾,虛榮,聲望,地位,金錢都會轉眼成空,不能存至永恆；因為他求神的喜悅.該隱表現了另一種人生觀,他所作的是為了自己,不是事奉神.後來,更因為亞伯做得比他更好,要將他除去.同樣,很多基督徒聽道若受感動,便願意接受真理的修改,而結果子,若不接受,則更會反駁,挑人的錯失來掩護自己的過失.因此,保羅說基督的香氣能叫人活,也能叫人死；不但自己的生命受傷,更會傷害別人.這是兩種人生觀.
\newline
\newline
最後,亞伯的血更說明了神是賞善罰惡的神.亞伯照自己所信的擺上羔羊獻予神,更獻上自己的生命.聖經上明顯地指出該隱是屬那惡者(約壹三12).神是公義的,信仰不是按家族,團體計算的,是按個別而論.該隱雖是亞當的兒子,卻是屬那惡者的.
\newline
\newline
另一方面,希伯來書(十二24)中,更說明了基督的血所講的,比亞伯的血所講的更美.亞伯是標準的殉道士,為著求神的喜悅,毫無怨言地給親兄殺死,他的血含有定罪的成份.然而,基督的血卻是宣告赦免的信息.亞伯雖好,但他的好只足夠顯出別人的不好；基督的血是代贖的血,讓我們可以靠著祂流血的功勞,使罪得赦免.因此祂的血所說的,比亞伯的更美.
\newline
\newline
感謝神,今天我們有基督,我們有好的行為,目的不是要別人被定罪,而是要引領人到可得赦罪的救主面前,叫人知道他需要一位赦罪的救主──耶穌基督,唯有祂才能叫人的罪得赦免.
\newline
\newline
創世記的篇幅是寶貴的,作者記載亞伯獻祭及被殺的經過；是要說明這事的重要,表明了救恩的基本原理.你和我在此試想一下,當了基督徒至今,有否像亞伯那樣,曾為基督做過甚麼事是可以存到永恆的呢?
\newline
\newline


\newpage

\section{第60屆~港九培靈會~陳終道~第二講}
\label{sec:1173}
\textbf{劃時代的挪亞}
\newline
\newline
連結: \href{https://www.hkbibleconference.org/session-message/view/1173}{\texttt{https://www.hkbibleconference.org/session-message/view/1173}}
\newline
\newline
\hyperref[sec:1172]{< < < PREV SERMON < < <}
~
\hyperref[sec:toc]{[返目錄]}
~
\hyperref[sec:1174]{> > > NEXT SERMON > > >}
\newline
\newline
創世記 5:28-5:29
\newline
\begin{longtable}{cl}
\hline
\hline
章節 & 經文 (和合本修訂版)\\
\hline
5:28 & \begin{tabularx}{0.7\textwidth}{X} 拉麥活到一百八十二歲，生了一個兒子， \end{tabularx} \\ \\ \relax
5:29 & \begin{tabularx}{0.7\textwidth}{X} 給他起名叫挪亞，說：「在耶和華所詛咒的地上，這個兒子必使我們從工作和手中的勞苦得到安慰。」 \end{tabularx} \\ \\
[1ex]
\hline
\hline
\end{longtable}
今天我們要看挪亞的生命.提起挪亞,就不期然想到他的方舟,但聖經上最先提到挪亞,卻描述他在接受大使命前,是家庭中的孝子,為父母分擔勞苦,使父母得到安慰(創五).要知道自己能否作神使用的器皿；先要看自己在家中,是否是一個盡責任的兒子.
\newline
\newline
第一,西方有一個錯誤的觀念常對(創二章)「人要離開父母與妻子連合」有所誤解.其實,這句話真正的意義是表明家庭中,夫妻關係會取代父母與子女的關係；並不是說子女毋須照顧父母.經文指出的意思,是兒女已成家立室,以後與丈夫或妻子的關係之親密,更甚於父母.挪亞跟父母同住,負擔他們的操作勞苦,直至進入方舟也是一樣.猶太人從舊約至新約,對家族的觀念是十分強的,不但兒女孝順父母,甚至孫子孫女也要孝順祖父母.挪亞是一位孝順的兒子,在家庭生活上他是成功的.
\newline
\newline
第二,談到挪亞的方舟(彼前三20)指出那使基督從死裡復活的聖靈,也是曾藉著挪亞在那時代,向人傳福音的靈.那福音就是挪亞的方舟,表明基督和祂的救恩.方舟經過洪水,就像人受浸時,水所表明的意義一樣,象徵我們與基督一同埋葬一同復活.挪亞當時對救恩是不大清楚明白,他只是依照神的命令造了方舟；只是方舟經過洪水,就表明了在基督裡面的人能經得起神的審判,以致能得救.故此,在舊約時那些有信心的人所信的,與我們今天相信的大致相同；分別只是舊約信的在於表明救恩的原理而已.
\newline
\newline
洪水時代可能也有其他的船,只有挪亞的方舟能經得起洪水的考驗,表明這救法是神所設計的.基督來到世上承擔人的罪,是神的計劃.祂只是一次把自己作為贖罪祭,就能承受到永遠；因此,舊約的人所信的,就是那要來臨的耶穌基督.(創廿二18)神對亞伯拉罕的應許中,其中「萬國」,是顯明福音是為萬國而設立.「後裔」是指耶穌基督.(加三16),意即萬國要因耶穌基督而得福.挪亞及亞伯拉罕所信的,都是要來臨的基督,只是用不同的預表來解釋.
\newline
\newline
第三,挪亞的時代,是邪惡淫亂的時代.那時代的人又吃,又喝,又嫁,又娶,不知不覺洪水就把他們吞滅.吃喝嫁娶本是正常的事,經文的意思,是指那時代的人只注重吃喝嫁娶,飲食,男女關係,慾望成為那時代的特徵.然而,聖經提及進入方舟的有八個人,就是挪亞一家八口,挪亞有三個兒子,這一家八口說明挪亞一家的夫妻關係,在那時代分別為聖,一夫一妻的,他的三個兒子也是.事實上,挪亞只是照神所指示而建造方舟,按神的旨意生活行事；但因他們的敬虔,愛主,他的見證就定了那時代的罪(來十一7).故此,不要以為我們現在所處的時代艱難,就無法作一個得勝者,這不是理由,每個時代都有它艱難之處.
\newline
\newline
第四,挪亞的信心有他特別之處,使他能在那時代得勝.他的信心有幾個特點:
\newline
\newline
(一)他的信心是長期孤立在邪惡的世代中,他沒有同流合污.神宣佈要滅那世代時,他仍存義道,用一百廿年建造方舟.他的信心經過長期,孤單地在邪惡的世代中成全,完全不受影響.
\newline
\newline
(二)他的信心能產生實際的行為.信心使他在家庭生活上成功,他的思想,心思並沒有被那世代的潮流所改變.信心使挪亞做了空前絕後的事；雖然聖經沒有明文提到方舟在山上建造,但我們可想而知,因為洪水來臨時,水位不斷上升,然後將方舟飄浮起來.在山頂上造方舟,旁人定必引為笑柄,但他仍堅持相信神所指示的沒有錯,藉著信心完成神的託付.
\newline
\newline
(三)他的信心是忍耐的信心.建造方舟一百二十年是忍耐,在方舟上生活一年多也是忍耐.與眾多動物一同困在方舟裡生活絕不好受.洪水退卻時,挪亞也沒有立即離開方舟,仍三番四次等候神的命令才出去.他忍耐等候神的旨意,直至成功.最後挪亞出方舟,首先做的事就是獻祭給神,向神表示感恩,歸榮耀給神.聖經上提及挪亞將各類潔淨的牲畜,獻在壇上作為燔祭,耶和華聞馨香之氣,決定不因人的緣故咒詛地.馨香之氣表明挪亞對神的感恩,他完成了偉大的使命,先行獻祭,歸榮耀給神.今天,神委託我們做某件事,我們完成後會否將榮耀完全歸於神呢?此外,挪亞的獻祭,更是先遵行神的命令,完成使命然後才獻上；並不是只有獻祭而沒有遵行神的命令.
\newline
\newline
最後,挪亞經過很多考驗,長久的忍耐,相信,完成神所吩咐的,然後,他成為新世界的主人.神將更大的應許給他,跟他重新立約；在挪亞以後,神將動物也賜給人作食物.挪亞造方舟,進入新世代的真正意義,是神要證明人的生命已經敗壞,必須有救恩.人的生命已趨向於惡,就算將所有惡人滅絕,剩下好人,結果也是一樣敗壞,包括挪亞本人.挪亞雖然是完全的人,但在成功後他也醉酒,並且赤身露體.然後也就演變到巴別塔叛亂的事.
\newline
\newline
因此,神要藉著歷史來教訓我們,當罪進入世界時,罪也進入人的生命,無論環境如何變新都是徒然的.人犯罪並不是因環境的緣故,我們需要的,是生命的主給我們的新生命.挪亞方舟在神整個救贖計劃裡,說明一個事實,就是人不能自救,人的生命已敗壞,人需要屬神的生命,只有神能給我們.這是挪亞所完成的見證,他的使命帶給我們的信息,使我們明白人無法從肉體中有任何的良善,只有在耶穌基督裡,才能將祂的榮美活出來.
\newline
\newline


\newpage

\section{第60屆~港九培靈會~陳終道~第三講}
\label{sec:1174}
\textbf{隨處得勝的約瑟}
\newline
\newline
連結: \href{https://www.hkbibleconference.org/session-message/view/1174}{\texttt{https://www.hkbibleconference.org/session-message/view/1174}}
\newline
\newline
\hyperref[sec:1173]{< < < PREV SERMON < < <}
~
\hyperref[sec:toc]{[返目錄]}
~
\hyperref[sec:1175]{> > > NEXT SERMON > > >}
\newline
\newline
創世記 37:1-37:36
\newline
\begin{longtable}{cl}
\hline
\hline
章節 & 經文 (和合本修訂版)\\
\hline
37:1 & \begin{tabularx}{0.7\textwidth}{X} 雅各住在迦南地，就是他父親寄居的地。 \end{tabularx} \\ \\ \relax
37:2 & \begin{tabularx}{0.7\textwidth}{X} 這是雅各的事蹟。約瑟十七歲與他哥哥們一同牧羊。他是個少年，與他父親的妾辟拉和悉帕的兒子們常在一起。約瑟把他們的惡行報給父親。 \end{tabularx} \\ \\ \relax
37:3 & \begin{tabularx}{0.7\textwidth}{X} 以色列愛約瑟過於其他的兒子，因為約瑟是他年老生的；他給約瑟做了一件長袍。 \end{tabularx} \\ \\ \relax
37:4 & \begin{tabularx}{0.7\textwidth}{X} 哥哥們見父親愛約瑟過於他們，就恨約瑟，不與他說友善的話。 \end{tabularx} \\ \\ \relax
37:5 & \begin{tabularx}{0.7\textwidth}{X} 約瑟做了一個夢，告訴他哥哥們，他們就更加恨他。 \end{tabularx} \\ \\ \relax
37:6 & \begin{tabularx}{0.7\textwidth}{X} 約瑟對他們說：「請聽我做的這個夢： \end{tabularx} \\ \\ \relax
37:7 & \begin{tabularx}{0.7\textwidth}{X} 看哪，我們在田裡捆禾稼；看哪，我的捆起來站著；看哪，你們的捆圍著我的捆下拜。」 \end{tabularx} \\ \\ \relax
37:8 & \begin{tabularx}{0.7\textwidth}{X} 他的哥哥們對他說：「難道你真的要作我們的王嗎？難道你真的要統治我們嗎？」他們就因他的夢和他的話更加恨他。 \end{tabularx} \\ \\ \relax
37:9 & \begin{tabularx}{0.7\textwidth}{X} 後來他又做了另一個夢，告訴他哥哥們說：「看哪，我又做了一個夢；看哪，太陽、月亮和十一顆星都向我下拜。」 \end{tabularx} \\ \\ \relax
37:10 & \begin{tabularx}{0.7\textwidth}{X} 約瑟告訴他父親和哥哥們，他父親就責備他說：「你做的這是甚麼夢！難道我和你母親、你的兄弟真的要俯伏在地，來向你下拜嗎？」 \end{tabularx} \\ \\ \relax
37:11 & \begin{tabularx}{0.7\textwidth}{X} 他的哥哥們都嫉妒他，他父親卻把這事存在心裡。 \end{tabularx} \\ \\ \relax
37:12 & \begin{tabularx}{0.7\textwidth}{X} 約瑟的哥哥們到示劍去放他們父親的羊。 \end{tabularx} \\ \\ \relax
37:13 & \begin{tabularx}{0.7\textwidth}{X} 以色列對約瑟說：「你哥哥們不是在示劍放羊嗎？來，我派你到他們那裡去。」約瑟對他說：「我在這裡。」 \end{tabularx} \\ \\ \relax
37:14 & \begin{tabularx}{0.7\textwidth}{X} 以色列對他說：「你去看看你哥哥們是否平安，羊群是否平安，再回來告訴我。」於是他派約瑟出希伯崙谷，約瑟就往示劍去了。 \end{tabularx} \\ \\ \relax
37:15 & \begin{tabularx}{0.7\textwidth}{X} 有人遇見他，看哪，他在田野走迷了路。那人問他說：「你找甚麼？」 \end{tabularx} \\ \\ \relax
37:16 & \begin{tabularx}{0.7\textwidth}{X} 他說：「我找我的哥哥們，請告訴我，他們在哪裡放羊。」 \end{tabularx} \\ \\ \relax
37:17 & \begin{tabularx}{0.7\textwidth}{X} 那人說：「他們已經離開這裡走了，我聽見他們說：『我們往多坍去。』」約瑟就去追哥哥們，在多坍找到了他們。 \end{tabularx} \\ \\ \relax
37:18 & \begin{tabularx}{0.7\textwidth}{X} 他們遠遠看見他，趁他還沒有走近他們，就圖謀要殺死他。 \end{tabularx} \\ \\ \relax
37:19 & \begin{tabularx}{0.7\textwidth}{X} 他們彼此說：「看哪！那做夢的來了。 \end{tabularx} \\ \\ \relax
37:20 & \begin{tabularx}{0.7\textwidth}{X} 現在，來吧！我們把他殺了，丟在一個坑裡，就說有惡獸把他吃了。我們且看他的夢將來怎麼樣。」 \end{tabularx} \\ \\ \relax
37:21 & \begin{tabularx}{0.7\textwidth}{X} 呂便聽見了，要救約瑟脫離他們的手，說：「我們不可害他的性命。」 \end{tabularx} \\ \\ \relax
37:22 & \begin{tabularx}{0.7\textwidth}{X} 呂便又對他們說：「不可流他的血，可以把他丟在這曠野的坑裡，不可下手害他。」呂便要救他脫離他們的手，把他還給他父親。 \end{tabularx} \\ \\ \relax
37:23 & \begin{tabularx}{0.7\textwidth}{X} 約瑟到了他哥哥們那裡，他們就剝去他的外衣，就是他身上那件長袍。 \end{tabularx} \\ \\ \relax
37:24 & \begin{tabularx}{0.7\textwidth}{X} 他們抓住他，把他丟在坑裡。那坑是空的，裡頭沒有水。 \end{tabularx} \\ \\ \relax
37:25 & \begin{tabularx}{0.7\textwidth}{X} 他們坐下吃飯，舉目觀看，看哪，有一群以實瑪利人從基列來，用駱駝馱著香料、乳香、沒藥，要帶下埃及去。 \end{tabularx} \\ \\ \relax
37:26 & \begin{tabularx}{0.7\textwidth}{X} 猶大對他的兄弟們說：「我們殺我們的弟弟，遮掩他的血有甚麼好處呢？ \end{tabularx} \\ \\ \relax
37:27 & \begin{tabularx}{0.7\textwidth}{X} 來，我們把他賣給以實瑪利人，不要下手害他，因為他是我們的弟弟，我們的骨肉。」他的兄弟們就聽從了他。 \end{tabularx} \\ \\ \relax
37:28 & \begin{tabularx}{0.7\textwidth}{X} 那時，有些米甸的商人從那裡經過，就把約瑟從坑裡拉上來。他們以二十塊銀子把約瑟賣給以實瑪利人，他們就把約瑟帶到埃及去了。 \end{tabularx} \\ \\ \relax
37:29 & \begin{tabularx}{0.7\textwidth}{X} 呂便回到坑旁，看哪，約瑟不在坑裡，就撕裂自己的衣服， \end{tabularx} \\ \\ \relax
37:30 & \begin{tabularx}{0.7\textwidth}{X} 回到他兄弟們那裡，說：「孩子不在了。我往哪裡去才好呢？」 \end{tabularx} \\ \\ \relax
37:31 & \begin{tabularx}{0.7\textwidth}{X} 於是，他們宰了一隻公山羊，拿了約瑟的那件外衣染上了血， \end{tabularx} \\ \\ \relax
37:32 & \begin{tabularx}{0.7\textwidth}{X} 派人把長袍送到他們的父親那裡，說：「我們發現這個， 請認一認，是不是你兒子的外衣？」 \end{tabularx} \\ \\ \relax
37:33 & \begin{tabularx}{0.7\textwidth}{X} 他認出來，就說：「這是我兒子的外衣，惡獸把他吃了，約瑟一定被撕碎了！」 \end{tabularx} \\ \\ \relax
37:34 & \begin{tabularx}{0.7\textwidth}{X} 雅各就撕裂衣服，腰間圍上麻布，為他兒子哀傷了多日。 \end{tabularx} \\ \\ \relax
37:35 & \begin{tabularx}{0.7\textwidth}{X} 他的兒女都起來安慰他，他卻不肯受安慰，說：「我必哀傷著下陰間，到我兒子那裡。」約瑟的父親就為他哀哭。 \end{tabularx} \\ \\ \relax
37:36 & \begin{tabularx}{0.7\textwidth}{X} 米甸人把約瑟賣到埃及，給法老的官員，就是護衛長波提乏。 \end{tabularx} \\ \\
[1ex]
\hline
\hline
\end{longtable}
聖經中的約瑟是一個得勝者,他無論作何事都是得勝者,因為他的一生都有耶和華的同在,是一個真正的得勝者.
\newline
\newline
約瑟的一生,可分為四個階段:
\newline
\newline
(一)在他居家的時期,即是他17歲之前的歲月.
\newline
\newline
(二)被賣在波提乏家中作奴隸之時期.
\newline
\newline
(三)被冤枉入獄作囚犯之年月.
\newline
\newline
(四)被升高,作全埃及宰相之時.約瑟的生平可從這四個主要階段來看.
\newline
\newline
現先從第一階段了解約瑟居家時,他的生活及見證.雅各偏愛約瑟,因約瑟為拉結所生；然而,問題是因雅各的偏心,就給約瑟帶來了許多禍患,約瑟不但沒有恃寵生嬌,還與哥哥一起牧羊.雖然約瑟得雅各偏愛,但卻有十個哥哥憎恨他,這使約瑟在家庭的處境非常困難.另一方面,約瑟雖然與哥哥一起牧羊,但卻不在哥哥的惡行上有份.(三172):「約瑟將他哥哥的惡行,報給他們的父親.」約瑟膽敢將哥哥們的惡行告知父親,足以證明他沒有參與他們的惡行,否則一定給哥哥們抓著把柄.可見,約瑟深得父親寵愛,但沒有因此而被寵壞；他一方面幫助哥哥們牧羊；另一方面又將他們的惡行告知父親,這並非為離間父親與哥哥們的關係,乃為雅各家的好處.雅各一家人在當時還是寄居在人家的地方,當時的迦南地並未屬於雅各家.
\newline
\newline
我在東南亞一帶生活了很多年,接觸到很多華僑,他們中間很多非常富有,只是常常受人欺侮.因此,大多數華僑都很愛國,當聽見中國強盛起來,便立即變賣所有產業,回國居住；心裡巴不得祖國可以為他們出一口氣.
\newline
\newline
(一)雅各家在僑居之地並不安全.約瑟將哥哥們的惡行告知父親,並非為討父親的歡心；而是維護雅各家在迦南的處境.否則約瑟不會故意惹哥哥們的憎厭.約瑟這樣做,並非與哥哥們不和,約瑟實在是愛他們.卅七章記載約瑟奉父親之命前往找尋哥哥們,本來雅各只是叫他往示劍找他們.約瑟找不著,便再往多坍去,終於找著了哥哥.為何雅各擔心十個兒子在示劍的安全呢?從創卅四章可知,雅各家與示劍人結下了冤仇.故此,雅各對於示劍人一向提心吊膽；打發兒子到那裡放羊,也可能是迫不得已.雅各打發約瑟前去找哥哥,可能那時他們已過了應該回家的時分.本來,十個哥哥們到示劍已不平安,然則約瑟一人前去更是不安全了,約瑟大可去到示劍便截回,然後告訴父親找不著哥哥們便可.只是,約瑟繼續前往,直至找著哥哥們.可見,約瑟也是擔心哥哥們的安全,他愛哥哥們.故此,他將哥哥們的惡行告知父親,是因為愛的緣故；不想哥哥們再作惡事,以免惹別人的反感.同時,也因為約瑟敬畏神,不忍心看見別人做一些不敬畏神的事.我們看見約瑟在家庭過著得勝的生活,他愛家人,但不包庇罪惡,反而指責罪惡；並不出於個人的私心或怨恨.他愛哥哥們,另一方面,又不想他們行差踏錯,這便是愛.他不惜冒著生命危險,去找那些正計劃殺他的哥哥們.最後被賣在以實瑪利人的手中,淪為奴僕的身份.
\newline
\newline
故此,有說約瑟的經歷很像我們的主耶穌基督,祂來尋找一些憎恨祂的罪人,最終被人出賣,釘上十字架.
\newline
\newline
(二)約瑟在波提乏家中的時期.記載於卅九章,先參看1-4節.雖然聖經沒有記載約瑟在波提乏家中住了多少年日,但可推想大概；被賣與升做宰相期間,約有十三年,估計有三,四年左右在監裡,也就是說有八,九年的時間是在波提乏家中度過.初到波提乏家中不過是17歲的童子,是奴隸的身份.主人所以能信任他,是因為經過很多事情,都能看見耶和華與他同在；使他所辦的事情盡都順利.這個階段的處境比較在家裡時更難,在家中是父親的驕兒,現在是人家的奴僕.然而,約瑟卻能扭轉劣境,因為祂敬畏神,神與他同在,以致主人信任他,提昇他為管家.
\newline
\newline
然而,約瑟在波提乏家中遇著很大的試探；波提乏的妻子──約瑟的主母企圖引誘他.約瑟堅決地拒絕她,請參考(三十九7-18).我們不說約瑟被冤枉之事,而看約瑟所面對的試探.這個試探並非容易對付.試想,約瑟天天在家裡都要面對主母,管家往往要向主母交代；其次,約瑟孤身一個青年人在外,遠離家庭,沒有家人或朋友的約束；同時,當時又沒有第三者在場.這種情況使人很容易接受試探.然而,約瑟如何勝過試探呢?只因為約瑟怕得罪神,他真正地活在神的眼前.今天,基督徒要做得勝者,你是否活在神的眼前,抑或活在人的輿論中?假若是活在神的眼中,無論人家看見與否,都是一樣.這是約瑟勝過試探的秘訣.中國人有曰:君子慎獨,也就是說,君子在獨處時要特別謹慎.
\newline
\newline
另一方面,約瑟更面對重大的考驗,他雖然拒絕了試探,卻反被主母誣告他,因而被放在監裡.相信約瑟屢次拒絕試探時,一定深信他沒有犯罪,神也必定會救他.但結果約瑟勝試探,卻沒有好下場,神也沒有救他；使他白白蒙上這種罪名.以前是驕兒,接著做奴僕,現在是囚犯,身份不斷下降.然而,約瑟並沒有怨天尤人,相信他心裡一定有很多疑問,不明白神為何讓他受苦,只是,他仍堅信神有祂的美意.故此,他沒有怨言,雖在監獄中,仍然有耶和華的同在(創三十九23).神不救他出獄,不為他伸冤,但是,卻與他同在.因此,不要以為我們的處境不好,便不敬畏神,不愛主,一再失敗,沒有喜樂......但願約瑟的經歷成為我們的榜樣.
\newline
\newline
一位姊妹,未信主時常騙取其母,以得多些零用.信主後,知道這樣做是不當的,便不敢再作.然而,她心裡一直禱告,希望神因她不再欺騙母親,而感動母親加添零用.結果,幾個月過去,禱告都沒有實現,她便以此詢問牧師.結果該堂牧師請她檢討自己的動機:她拒絕犯罪,是因為恨惡罪惡,抑或是為了得到一些好處呢?喜愛公義又恨惡罪惡,這才是體會神的心意.
\newline
\newline
約瑟拒絕試探,並非要神拯救他,或給他一些好處,是因為他敬畏神；怕得罪神,因而拒絕主母.故此,就是在監獄中,也有神的同在.假若神真是聽他的禱告,為他伸冤,頂多也只能回復管家的身份.神暫不拯救他,讓他繼續受很多苦,這樣,才能當上埃及的宰相,不但管理波提乏的家,而是做法老的管家,管理全埃及的人.故此,神並非不聽禱告,只是神遲延成就祂的旨意,是有祂的美意.只藉著法老的一個夢,便使約瑟擢升為宰相.
\newline
\newline
(三)我們看看約瑟做了宰相後,他的靈性生活如何.多年前有則新聞,記載一位做苦力的人中了彩票,在一夜間成為富翁,先行享受一番,把剩下的錢開辦酒樓,因不善經營,在兩年間虧蝕了所有資本.那時,酒肉朋友也都離他而去；只好再回頭做苦力.但可惜兩年間的富足生活,不但養胖了他,同時使他喪失了從前的力氣；不能搬東西了.他走投無路,頓覺人生乏味,便投海自盡.可見,一個人從貧窮的境況忽然做了富翁,並未能受得住這遭遇.同樣,約瑟本是囚犯,卻升做宰相,他如何應付呢?這個試探很大,然而,約瑟做了宰相,也如常地敬畏神.能伸能屈,堪稱為大丈夫.
\newline
\newline
以下請查看數段經文:(創四十一51)記載約瑟給長子命名為嗎哪西,意即神使他忘了一切困苦和我父的全家.換言之,這名字有感謝神的恩典之意.52節記載他給次子起名為以法蓮,意即神使他在受苦的地方昌盛.同時有感恩的意思.
\newline
\newline
四十二18約瑟說:「我是敬畏神的.」
\newline
\newline
四十三23約瑟考驗哥哥們是否善待便雅憫.其家宰說:「是你們的神,和你們父親的神,賜給你們財寶在你他的口袋裡.」證明了約瑟的家宰也認識神的作為,同時知道他主人是敬畏神的人.
\newline
\newline
五十章記載雅各死後,約瑟的哥哥們恐怕約瑟會報仇,以前只是因為父親的顏臉才暫且放過他們,現在約瑟只要下一道命令,便可將他們置諸死地,然而,約瑟安慰他們說:「從前你們的意思是要害我,但神的意思原是好的,要保全許多人的性命,成就今日的光景.」
\newline
\newline
(四)約瑟的遺囑交代要將他的骸骨搬回迦南,葬在列祖的墓地.這說明了約瑟雖然在埃及為宰相,但他沒有忘記神的應許,終有一天以色列人要歸回迦南地.
\newline
\newline
約瑟17歲被賣,卅歲做宰相,110歲逝世,做宰相的年日也有八十年之久,被賣往埃及的日子這麼長,並沒有改變他的心志,也沒有動搖他的信仰.今天,很多人移居外地,只要八個星期便可將一個人改變.約瑟離開家庭時只有17歲,可知他在17歲前,已建立起他的信仰基礎,以及人生目標.故此,他能夠經受得起在埃及漫長的日子,富貴榮華,絲毫不搖動他的心志.蓋棺定論,直至死的那天,他的心仍然想望迦南──神應許之地.他不知道為何神讓他經過那麼多坎坷的道路；但他認識一件事,就是他何時敬畏神,神就與他同在.因此,約瑟被賣在埃及,成就了神計劃及旨意.約瑟把他父家接到埃及,使他們生養眾多,成為以後以色列立國的根基.在神要設立屬祂的國度之計劃中,約瑟作了這個貢獻.這是因為約瑟一生敬畏神,凡事信靠神,不知不覺間,成就了神在他身上的計劃.
\newline
\newline
今天,我們應該把自己放在神全能的手中；讓祂在我們的身上,完成祂要成就的旨意.
\newline
\newline


\newpage

\section{第60屆~港九培靈會~陳終道~第四講}
\label{sec:1175}
\textbf{神所尊重的摩西}
\newline
\newline
連結: \href{https://www.hkbibleconference.org/session-message/view/1175}{\texttt{https://www.hkbibleconference.org/session-message/view/1175}}
\newline
\newline
\hyperref[sec:1174]{< < < PREV SERMON < < <}
~
\hyperref[sec:toc]{[返目錄]}
~
\hyperref[sec:1176]{> > > NEXT SERMON > > >}
\newline
\newline
民數記 12:1-12:9
\newline
\begin{longtable}{cl}
\hline
\hline
章節 & 經文 (和合本修訂版)\\
\hline
12:1 & \begin{tabularx}{0.7\textwidth}{X} 摩西娶了古實女子為妻。米利暗和亞倫因他娶了古實女子就批評他， \end{tabularx} \\ \\ \relax
12:2 & \begin{tabularx}{0.7\textwidth}{X} 他們說：「難道耶和華只與摩西說話嗎？他不也與我們說話嗎？」耶和華聽見了。 \end{tabularx} \\ \\ \relax
12:3 & \begin{tabularx}{0.7\textwidth}{X} 摩西為人極其謙和，勝過地面上的任何人。 \end{tabularx} \\ \\ \relax
12:4 & \begin{tabularx}{0.7\textwidth}{X} 忽然，耶和華對摩西、亞倫和米利暗說：「你們三個人都出來，到會幕這裡。」他們三個人就出來了。 \end{tabularx} \\ \\ \relax
12:5 & \begin{tabularx}{0.7\textwidth}{X} 耶和華在雲柱中降臨，停在會幕門口，叫亞倫和米利暗。二人就出來， \end{tabularx} \\ \\ \relax
12:6 & \begin{tabularx}{0.7\textwidth}{X} 耶和華說：「你們要聽我的話：你們中間若有先知，我－耶和華必在異象中向他顯現，在夢中與他說話； \end{tabularx} \\ \\ \relax
12:7 & \begin{tabularx}{0.7\textwidth}{X} 但我的僕人摩西不是這樣，他在我全家是盡忠的。 \end{tabularx} \\ \\ \relax
12:8 & \begin{tabularx}{0.7\textwidth}{X} 我與他面對面說話，清清楚楚，不用謎語，他甚至看見我的形像。你們為何批評我的僕人摩西而不懼怕呢？」 \end{tabularx} \\ \\ \relax
12:9 & \begin{tabularx}{0.7\textwidth}{X} 耶和華向他們怒氣發作，就離開了。 \end{tabularx} \\ \\
[1ex]
\hline
\hline
\end{longtable}
舊約中,亞伯拉罕,摩西及大衛可說是最受以色列人敬重的人,其中以摩西尤甚.摩西對神話語的領受,為舊約先知之最.甚至在新約時,猶太人亦以摩西的律法來指控耶穌,可見摩西在猶太人心目中有重要的地位.現在我們會在舊約中,挑選摩西所獨有的榮耀來談.
\newline
\newline
第一,摩西被稱為舊約的中保(加三16-17),這是摩西特有的尊榮,是律法的經手人.在西乃山時,神宣佈與以色列人立約,要以色列人自潔準備立約.然後,摩西將神的約言向他們宣告,又將牛的血一半灑在約書上；一半灑在以色列人身上,這樣,舊約便成立了.聖經有分新舊約,新約為基督所設立,舊約則經由摩西而成立.舊約的設立,是教導以色列人如何敬拜神,叫人知罪.舊約的律法非叫人得救,故稱為訓蒙的師傅,至新約時,律法才得以完成.律法字句的背後,是要人明白救恩.摩西是律法的中保,但他不能擔保百姓會得救.律法就如血壓計,只能讓人知道血壓情況如何,但不能叫血壓變為高或低,又如ｘ光,只能使人知道有毛病,但沒有醫治功能.律法的中保不能替人承擔罪,律法叫人知弊端,但人要自行找醫治方法.因此,新約的中保──基督替人成全律法,擔保人可以得救,祂替人還債,使人有得救的把握.
\newline
\newline
第二,經文中(十二7)神印證摩西為神的全家盡忠.意即為全以色列人忠心,並不是為其支派或私人的利益,亞倫的兩個兒子,第一天上任時已被神處罰而死,雖是悲傷的事,但摩西囑咐亞倫,不能表示悲哀.只能完全站在神的一邊,其感情不能為自己而抒發,要完全好神所好.摩西在神的工作上全無偏私,完全為神的全家而忠心.今天,在教會中,人往往忽略神的全家,只看見自己的小圈子.見別人工作好便加以譭謗,此亦是亞倫,米利暗與摩西的分別,他們二人有私心.故此,神特別稱讚摩西,他為神的全家盡忠.
\newline
\newline
第三,神稱讚摩西「為人極其謙和,勝過世上的眾人」(十二3),聖經上只有一人得神如此讚賞.摩西並非沒有脾氣之人,他多次發怒,最嚴重的是與神立約後,下山即看見百姓拜金牛犢；此舉已毀了剛立的約,摩西一怒之下,便摔碎法版.此外,摩西還多次向百姓發怒.因此,以色列人多次要以石頭擲他,又多次埋怨.發生何事,便埋怨摩西.然而,摩西為人極其謙和,在眾多以色列人的煩瑣事中,在許多的怨言中,作為一個領導委實不容易.因此,神稱讚摩西的溫和,是世上沒有人能及的.
\newline
\newline
第四,摩西行了最多懲罰性的神蹟,整本聖經中,以利亞曾行大神蹟,使天三年零六個月不下雨,這是懲罰性的神蹟.摩西則行很多類似性質的神蹟.諸如:埃及中的十災,最後一災更是殺埃及人的長子.
\newline
\newline
過紅海時,法老的追兵葬身紅海.以後,對於反對者的懲罰神蹟也很多,這在先知中,沒有人能及摩西.
\newline
\newline
第五,只有摩西可見到神的背面.摩西因以色列人拜金牛犢之事,上山求神的赦免.又求神向他顯現,但神對他說:沒有人可以看見我.只是,神讓摩西可得見祂的背面.這看來似乎有矛盾之處,舊約亞伯拉罕,以撒,雅各等都會在夢中,得到神向他們顯現；先知但以理,以西結等亦曾見神的榮耀,以致全身仆倒.神的榮耀並非人所能承當,聖經明說人見神的榮耀就必死.按神本來的榮耀,是無人有資格可以見的；故此,神以不同程度的榮耀向不同的人顯現.然而,無人可見神的面.而摩西則有所不同,神特別讓他得見祂的背面.這是摩西獨有的尊榮,表明他與神關係之親密.
\newline
\newline
第六,摩西的面貌發光.摩西從山上下來面容發光,以色列人不敢就近他,以致摩西要用帕子掩面才與百姓說話.顯明摩西與神的關係特別親密.今天,我們若要在這時代站立得住,就必定要與神有親密的關係．摩西能成為眾多神僕之中,領受神最多訊息的人,是因他與神有親密關係.林後三章保羅曾引用摩西來講新約的職事,舊約的職事摩西有神的榮光,而摩西是表明基督,今天我們因基督的緣故,得以在神家裡事奉,承擔新約的職事.這職事比摩西的榮光更大；因今天事奉基督,是在祂的榮光中,成為發光的信徒.我們並非有榮耀,是因主有榮耀；只要我們抓緊目標,服事我們的主人；便自然有榮耀,能得著尊貴的身份.
\newline
\newline
第七,神與摩西面對面明說(十二8).意即神對摩西說的話句句清楚明確.整本摩西五經,每句話都是清楚明說,如「不可殺人」,「不可貪心」,全是權威的啟示.神給摩西最多啟示,作為與他同在的證據.就如神的啟示不到臨老祭司以利,而臨到童子撒母耳一樣.因以利不會尊重神,故此,耶和華的言語稀少,這是很悲慘的現象.試問今天,我們是否落到「耶和華言語稀少」的境況呢?教會不尊主為大,庇護罪惡,則神的訊息不會臨到,只會出現屬靈的饑荒.
\newline
\newline
第八,摩西是神親自指明,作為預表基督的先知.申命記十九18記載,神要在弟兄中間,興起一位先知像摩西,這「先知」是指基督.施洗約翰初出來傳道時,猶太人問他是否那「先知」?使徒出來傳道,亦引用申命記的經文,以證明基督確是那先知(徒三22).只有摩西,神指著他來預表基督.
\newline
\newline
最後,只有摩西,得神親自為他安葬.申命記載述神將摩西葬在摩押地.現在有人逝世,得政府要員,太平紳士等扶柩,已是非常光采的事.世上有何人,是萬主之主親自為他埋葬的呢?這是摩西事奉主所得到的光榮.摩西一生受盡很多冤屈,以色列人在他生前極之痛恨他,但死後卻又非常敬愛他；以他為光榮,摩西若像亞倫般任由百姓擺佈,就必不得眾人之敬愛.摩西雖然溫柔,但堅持原則及真理的方向.故此,他死後以色列人幾乎視之為偶像.
\newline
\newline
此外,摩西與以利亞有同一榮耀,便是基督登山變像時,與主一起講論祂離世之事.事奉世上之事,會與世事一起成為過去,然而,摩西的事奉能存到永恆.摩西一生最光榮的年日,是八十歲以後的日子,可謂大器晚成.重要的是他能將餘下的時日全為主用.
\newline
\newline


\newpage

\section{第60屆~港九培靈會~陳終道~第五講}
\label{sec:1176}
\textbf{看見異象的以賽亞}
\newline
\newline
連結: \href{https://www.hkbibleconference.org/session-message/view/1176}{\texttt{https://www.hkbibleconference.org/session-message/view/1176}}
\newline
\newline
\hyperref[sec:1175]{< < < PREV SERMON < < <}
~
\hyperref[sec:toc]{[返目錄]}
~
\hyperref[sec:1177]{> > > NEXT SERMON > > >}
\newline
\newline
以賽亞書 6:1-6:13
\newline
\begin{longtable}{cl}
\hline
\hline
章節 & 經文 (和合本修訂版)\\
\hline
6:1 & \begin{tabularx}{0.7\textwidth}{X} 當烏西雅王崩的那年，我看見主坐在高高的寶座上。他的衣裳下襬遮滿聖殿。 \end{tabularx} \\ \\ \relax
6:2 & \begin{tabularx}{0.7\textwidth}{X} 上有撒拉弗侍立，各有六個翅膀：兩個翅膀遮臉，兩個翅膀遮腳，兩個翅膀飛翔， \end{tabularx} \\ \\ \relax
6:3 & \begin{tabularx}{0.7\textwidth}{X} 彼此呼喊說：「聖哉！聖哉！聖哉！萬軍之耶和華；他的榮光遍滿全地！」 \end{tabularx} \\ \\ \relax
6:4 & \begin{tabularx}{0.7\textwidth}{X} 因呼喊者的聲音，門檻的根基震動，殿裡充滿了煙雲。 \end{tabularx} \\ \\ \relax
6:5 & \begin{tabularx}{0.7\textwidth}{X} 那時我說：「禍哉！我滅亡了！因為我是嘴唇不潔的人，住在嘴唇不潔的民中，又因我親眼看見大君王－萬軍之耶和華。」 \end{tabularx} \\ \\ \relax
6:6 & \begin{tabularx}{0.7\textwidth}{X} 有一撒拉弗向我飛來，手裡拿著燒紅的炭，是用火鉗從壇上取下來的， \end{tabularx} \\ \\ \relax
6:7 & \begin{tabularx}{0.7\textwidth}{X} 用炭沾我的口，說：「看哪，這炭沾了你的嘴唇，你的罪孽便除掉，你的罪惡就赦免了。」 \end{tabularx} \\ \\ \relax
6:8 & \begin{tabularx}{0.7\textwidth}{X} 我聽見主的聲音說：「我可以差遣誰呢？誰肯為我們去呢？」我說：「我在這裡，請差遣我！」 \end{tabularx} \\ \\ \relax
6:9 & \begin{tabularx}{0.7\textwidth}{X} 他說：「你去告訴這百姓說：『你們聽了又聽，卻不明白；看了又看，卻不曉得。』 \end{tabularx} \\ \\ \relax
6:10 & \begin{tabularx}{0.7\textwidth}{X} 要使這百姓心蒙油脂，耳朵發沉，眼睛昏花；恐怕他們眼睛看見，耳朵聽見，心裡明白，回轉過來，就得醫治。」 \end{tabularx} \\ \\ \relax
6:11 & \begin{tabularx}{0.7\textwidth}{X} 我就說：「主啊，這到幾時為止呢？」他說：「直到城鎮荒涼，無人居住，房屋空無一人，土地極其荒蕪； \end{tabularx} \\ \\ \relax
6:12 & \begin{tabularx}{0.7\textwidth}{X} 耶和華將人遷到遠方，國內被撇棄的土地很多。 \end{tabularx} \\ \\ \relax
6:13 & \begin{tabularx}{0.7\textwidth}{X} 國內剩下的人若還有十分之一，也必被吞滅。然而如同大樹與橡樹，雖被砍伐，殘幹卻仍存留，聖潔的苗裔是它的殘幹。」 \end{tabularx} \\ \\
[1ex]
\hline
\hline
\end{longtable}
以賽亞在烏西雅王駕崩的時候看見異象.烏西雅總算是一位好王,在位的年日亦頗長；但可惜稍有成就便心驕氣傲,擅進聖殿燒香；而長了大痲瘋,故此,不能正式執政,只好由他的長子執掌王權.當國家漸走下坡時,一個大致上稱得上好的王卻駕崩,就會使人更落入失望中.何況,接任繼位的約坦王更壞得多,亞哈斯王就更不用說了.當時,北方以色列國已接近滅亡,令到南國大感到悲觀.這時,先知以賽亞看見異象,他看見神坐在高高的寶座上.烏西雅王雖然晚節不保；但坐在寶座上的神,卻是永遠不改變的.
\newline
\newline
舊約中,除祭司以外,連王帝也不能在聖殿燒香.(代下二十六16-20)清楚記載烏西雅王擅進聖殿燒香的事；因此,他長了大痲瘋.祭司的選立,除了只指利未支派外,還須指明要由亞倫的子孫才可以選立；故此,舊約中很多人羨慕作祭司的事奉.連大衛王也不例外,他更渴慕能在神的殿中事奉.烏西雅王因為心高氣傲要燒香,結果長了大痲瘋.耶羅波安所犯的罪,除了設金牛犢代替神之外；還選立凡民作祭司,即按立亞倫的子孫以外的人做祭司,實在是罪上加罪了.其實事奉神的權利並非每個世俗的人都能得到的.因此,今天我們因基督能成為祭司的國度,有權利事奉神；就不要輕忽這個權利了.何況,這在舊約,是很多君王所羨慕,甚至要爭奪也得不到的福份.烏西雅王雖然作了不該作之事,但他到底也是羨慕這福份.可惜,今天的基督徒雖有這等權利,卻不曉得愛惜!
\newline
\newline
然後,以賽亞描寫他所看見的異象.他看見神高高的坐在寶座上,衣裳下垂遮滿聖殿.我相信以賽亞必定是在聖殿中,敬拜神的時候看見異象；然而,他所看見的卻是天上的聖殿,只有天上的聖殿才有寶座.衣裳遮滿聖殿,有兩個解釋:其一是指神的榮耀充滿聖殿,「衣裳」其實是象徵一種榮耀,充滿聖殿.我個人則更喜歡另一種解釋,我相信先知所看見的,是真正的衣裳.這不但以賽亞曾看見,使徒約翰也曾有同樣的看見；還有其他先知,看見人子身穿白衣,同樣是看見衣裳.衣裳遮滿聖殿；意即連同地上的聖殿,也在神的衣裳蔭庇下,先知在聖殿中敬拜,顯得自己極其渺小.地上的君王已成為過去,然而天上的神卻永遠存在；祂的蔭庇,眷顧,榮耀和能力,要保護屬祂的人.地上有一九九七,天上的神卻沒有所謂一九九七；因為祂從亙古到永遠,坐在寶座上；永遠不變,祂的衣裳遮滿了聖殿.
\newline
\newline
此外,「其中有撒拉弗」,有六個翅膀,分別遮臉,遮腳及飛翔；表現出不配的含意,可見就是最接近神的人,也不敢面對神的榮耀.這異象特別之處,是天使們彼此呼喊說:聖哉,聖哉,聖哉,這是他們對所事奉之神的形容詞.有人說,將來在天堂的生活一定是非常的煩悶,要整天唱詩讚美,必定缺少樂趣.但是,從天使們彼此呼喊讚美神,卻表現出有未盡讚美說話之感.相信大家也認識有些老人家,他們在人前不斷訴說兒孫的趣事,永遠不會生厭；還總感到未能講盡種種可愛的事,每次提起都像新事一般.天使之表現,就像怎樣呼喊讚美,也無法唱盡神的榮美及聖潔；故此,他們還不斷地唱著,互相呼喊讚美神.
\newline
\newline
接著,以賽亞看見聖殿門檻的根基震動,殿充滿了煙雲.這個聖殿料必是指地上的殿.於是,以賽亞說:「禍哉!我滅亡了.因為我是嘴唇不潔的人,又住在嘴唇不潔的民中.」其實,天使沒指責以賽亞；但從天使的讚美聲中,先知已感覺到,自己實在是嘴唇不潔的人.從天使的頌聲中,先知已領悟了神是何等的聖潔.於是,他將自己與天使比較；卻只能說:禍哉,我滅亡了.先知深深感受到自己的不配.試問,我們曾否有這種體會和光照,使人從心靈醒悟的體驗呢?事實上一個人若沒有聖靈的光照,就不會看見自己的罪,也不會去認罪.
\newline
\newline
何以先知特別提到嘴唇不潔的罪呢?雅各書曾說:若有人在言語上沒有過失,這人便是完全人.事實上,在我們一切生活行為上,說話是佔最多比重的行為；每個人都是說話多過做事,因為開口說話在所有行動中,是最方便的一種.故此,以賽亞強調自己是嘴唇不潔之人,又是住在嘴唇不潔之民中.意思是不但自己會講不潔的話；同時又會受人的影響而多說不潔的話.例如是非,譭謗的說話,只會愈傳愈多,愈講愈難聽.故此,以賽亞看見異象時,就說:禍哉,我滅亡了.當先知有這種醒覺時,神便潔淨他的嘴唇；天使從壇上取下紅炭,沾先知的嘴,便除掉罪惡.這「壇」必定是指獻燔祭的壇,這壇首次點著的火,是神從天上降下的.因此,燔祭壇又是預表基督的十字架；祂將自己獻上,如馨香的祭,為人贖罪.
\newline
\newline
我們憑甚麼來潔淨自己呢?是憑基督贖罪的功勞.只有將自己完全獻上,讓燔祭壇的火燃燒我們,才有資格為主傳講訊息.嘴唇的影響遍及人的全身；當我們懂得謹慎言語,為神而用時,這嘴唇才是潔淨的.以賽亞被燔祭壇的火沾他的嘴唇,他便得潔淨.因為一個人能在言語上謹慎,他所受的造就,將是整個靈性生命也包括在內.
\newline
\newline
其後,先知又聽到呼召──「我可以差遣誰呢,誰肯為我們去呢.」真正照神的話語,而發出的警告,責備或勸勉,是不受人歡迎的.以賽亞也心有同感.故此,神再次對他發出呼召,因為願意去的並不是很多人.尤其是這個世代,人從小便給寵壞了,是不受責備的世代.以賽亞的使命更是非常艱巨.因為他的聽眾雖是再三聽道,也不受感動的.(賽六9-10)新約中,保羅曾引這段經文向猶太人傳講.(徒二十八23-29)形容那些聽眾的硬心,是「恐怕」──害怕因聽道而受感動.然而,神卻不改變祂的信息,表明神一再的忍耐；仍然差遣祂的僕人來勸告硬心的百姓.這便是以賽亞所肩負的使命；他的使命艱巨,訊息無人歡迎,但仍然要忠心去傳講.
\newline
\newline
「我可以差遣誰呢?誰肯為我們去呢?」以賽亞回答:「我在這裡,請差遣我.」求主施恩,讓我們在神面前看見自己的不配,又看見神的偉大；使我們能潔淨自己的嘴唇,有膽量在這悖逆的世代中,為福音作美好的見證.
\newline
\newline


\newpage

\section{第60屆~港九培靈會~陳終道~第六講}
\label{sec:1177}
\textbf{嫉惡如仇的施洗約翰}
\newline
\newline
連結: \href{https://www.hkbibleconference.org/session-message/view/1177}{\texttt{https://www.hkbibleconference.org/session-message/view/1177}}
\newline
\newline
\hyperref[sec:1176]{< < < PREV SERMON < < <}
~
\hyperref[sec:toc]{[返目錄]}
~
\hyperref[sec:1178]{> > > NEXT SERMON > > >}
\newline
\newline
馬太福音 3:1-3:17
\newline
\begin{longtable}{cl}
\hline
\hline
章節 & 經文 (和合本修訂版)\\
\hline
3:1 & \begin{tabularx}{0.7\textwidth}{X} 在那些日子，施洗的約翰出來，在猶太的曠野宣講： \end{tabularx} \\ \\ \relax
3:2 & \begin{tabularx}{0.7\textwidth}{X} 「你們要悔改！因為天國近了。」 \end{tabularx} \\ \\ \relax
3:3 & \begin{tabularx}{0.7\textwidth}{X} 這人就是以賽亞先知所說的：「在曠野有聲音呼喊著：預備主的道，修直他的路。」 \end{tabularx} \\ \\ \relax
3:4 & \begin{tabularx}{0.7\textwidth}{X} 這約翰身穿駱駝毛的衣服，腰束皮帶，吃的是蝗蟲和野蜜。 \end{tabularx} \\ \\ \relax
3:5 & \begin{tabularx}{0.7\textwidth}{X} 那時，耶路撒冷、全猶太和全約旦河地區的人，都到約翰那裡去， \end{tabularx} \\ \\ \relax
3:6 & \begin{tabularx}{0.7\textwidth}{X} 承認他們的罪，在約旦河裡受他的洗。 \end{tabularx} \\ \\ \relax
3:7 & \begin{tabularx}{0.7\textwidth}{X} 約翰看見許多法利賽人和撒都該人也來受洗，就對他們說：「毒蛇的孽種啊，誰指示你們逃避那將要來的憤怒呢？ \end{tabularx} \\ \\ \relax
3:8 & \begin{tabularx}{0.7\textwidth}{X} 你們要結出果子來，和悔改的心相稱。 \end{tabularx} \\ \\ \relax
3:9 & \begin{tabularx}{0.7\textwidth}{X} 不要自己心裡說：『我們有亞伯拉罕為祖宗。』我告訴你們，神能從這些石頭中給亞伯拉罕興起子孫來。 \end{tabularx} \\ \\ \relax
3:10 & \begin{tabularx}{0.7\textwidth}{X} 現在斧子已經放在樹根上，凡不結好果子的樹就砍下來，丟在火裡。 \end{tabularx} \\ \\ \relax
3:11 & \begin{tabularx}{0.7\textwidth}{X} 我是用水給你們施洗，叫你們悔改；但那在我以後來的，能力比我更大，我就是給他提鞋子也不配，他要用聖靈與火給你們施洗。 \end{tabularx} \\ \\ \relax
3:12 & \begin{tabularx}{0.7\textwidth}{X} 他手裡拿著簸箕，要揚淨他的穀物，把麥子收在倉裡，把糠用不滅的火燒盡。」 \end{tabularx} \\ \\ \relax
3:13 & \begin{tabularx}{0.7\textwidth}{X} 當時，耶穌從加利利來到約旦河，到了約翰那裡，請約翰為他施洗。 \end{tabularx} \\ \\ \relax
3:14 & \begin{tabularx}{0.7\textwidth}{X} 約翰想要阻止他，說：「我應該受你的洗，你怎麼到我這裡來呢？」 \end{tabularx} \\ \\ \relax
3:15 & \begin{tabularx}{0.7\textwidth}{X} 耶穌回答他：「暫且這樣做吧，因為我們理當這樣履行全部的義。」於是約翰就依了他。 \end{tabularx} \\ \\ \relax
3:16 & \begin{tabularx}{0.7\textwidth}{X} 耶穌受了洗，隨即從水裡上來。天忽然為他開了，他看見神的靈降下，彷彿鴿子落在他身上。 \end{tabularx} \\ \\ \relax
3:17 & \begin{tabularx}{0.7\textwidth}{X} 這時，天上有聲音說：「這是我的愛子，我所喜愛的。」 \end{tabularx} \\ \\
[1ex]
\hline
\hline
\end{longtable}
基督降生之前,先知預言有使者為祂預備道路(賽四十3),今天,我們要從幾方面來看這使者──施洗約翰的生平.
\newline
\newline
第一,要看施洗約翰的出生.他的出生有其特別的地方,他是唯一有先知預言他會來臨的先知.除以賽亞外,瑪拉基也有預言.可見,施洗約翰的誕生早在神的計劃中.他跟我們一樣都是普通人,若果他在神計劃中,是要來為耶穌預備道路,則我們今天在世並非偶然,也有所擔負著的使命.
\newline
\newline
當天使向撒迦利亞報信時,形容施洗約翰將有以利亞的心志,能力,以後我們也看到約翰的確如以利亞般,勇敢指責罪惡.此外,聖經上形容他在母腹中就被聖靈充滿；路加記載馬利亞往見以利沙伯時,施洗約翰在母腹中,見到主的母親就跳動歡喜,被聖靈充滿,表明這人自母腹就被神揀選了.
\newline
\newline
第二,就是施洗約翰的洗禮.先要注意的,是他的洗禮與舊約中的潔淨禮無關.約翰的洗禮不需要到祭司或聖殿那裡,行潔淨禮及獻祭等手續.約翰的洗禮是從神的啟示來的.他的洗禮,作用在於叫人悔改,信那在他以後要來的基督.因此,他宣講:「神的國近了,你們要悔改,信福音.」使徒行傳(十九1-7)講到在以弗所的幾位門徒,未曾受聖靈,只受了悔改的禮,但這還未能得救,還須相信接受基督為救主,才能得救,因此,保羅奉基督的名為他們施洗.
\newline
\newline
約翰未到世上,不但作基督的先鋒為人預備道路,連他所施行的洗禮都是個預備.當時,耶穌也到約翰處受洗.祂是希望藉著施洗約翰所行的洗禮,成為後來信徒的標記──公開承記自己的信仰.
\newline
\newline
第三,談約翰的工作.他的工作方式新奇,他身穿毛衣,吃蝗蟲野蜜,到曠野傳道.然而,猶太全地,約但河岸的人都走到曠野聽他的道.他以先知的姿態出現,這不是為了做作,而是神為他預備,從他的出生以致出來傳道的表現,都是這樣神奇,他與基督不同,耶穌的出現是平凡的,別人看輕祂,不尊重祂這個平凡的人.約翰的傳道是為了修直主的道路,將人心裡抵擋基督的心態預備好；因此,約翰可謂是最早做福音預工的人.
\newline
\newline
另外,約翰很厲害,直接勇敢地指責罪惡.他稱前來受洗的法利賽人及撒都該人為:毒蛇的種類,並告訴他們,凡不結果子的樹,都要被砍下來燒掉(太三7-10).罵人「毒蛇的種類」是非同小可的；碰巧耶穌也是這樣責備法利賽人.約翰責備他們,是因他們前來受洗並不是為著悔改,而是裝作謙卑及悔改的模樣.意即不能以虛偽的宗教儀式,來逃避將來的忿怒；除非從心裡真正悔改,否則神不會受外表的表演所愚弄.
\newline
\newline
今天的人是不易接受忠誠勸告的.有一回,一位弟兄的父母向我哭訴他們作執事的兒子不孝順,我便勸告那弟兄；他表面上順從了,還帶同父母與我吃晚飯.但他後來離開了原來的教會,在街上也不跟我打招呼,如同陌路人；事實上,他沒有接受我的勸告.現在的教育也看重鼓勵,稱讚；小時候有父母的教導,但長大後,有能力自立,供養父母,父母就少管了.有成就後,成為別人的上司,則更加少受人勸告及責備.耶穌曾說:人都說你好的時候,你就有禍了.因為這使人愈來愈不曉得自己的本相,愈來愈驕傲.沒有人告誡,每個人都說好話,但若細想一下,這些都是謊言,便會感到虛空.
\newline
\newline
第四,約翰是最忠心高舉基督的人.約翰出來傳道時,他受群眾的擁護及尊敬.當時,不少人懷疑約翰就是基督,有人詢問他,但他堅定否認.他是個專心高舉基督的人.他說:「祂必興旺,我必衰微.」其實,細想一下,這是個極大的試探,因為當時人人皆信任敬重他；而且還未認識耶穌基督.但是,約翰一點也沒有偷竊基督的榮耀,他將榮耀完全歸於基督.他在世上工作時間雖短,但卻完全地完成了神所託付的使命──把基督介紹給世人,為祂預備道路,指證人要悔改,相信基督.因此,耶穌稱許他說:「凡婦人所生的,沒有人大過施洗約翰.」眾先知能為耶穌講預言已是無限光榮；但他們無福份見到耶穌.約翰卻作了耶穌的先鋒,親自替耶穌施洗,這是何等大的福氣.按他所蒙的福,他得到眾先知渴望的福份,按他的使命來說,他已完全了自己的使命；他完全將基督舉薦出來,讓世人認識,完成使命後生命就結束,很短的傳道時間,卻有很大的成就.
\newline
\newline
今天,我們怎樣作基督的先鋒呢?基督第一次已來了,但祂允諾要再來.並說:「我必快來」,當他再來時,祂要我們作祂的先鋒.羅馬書十二章勸奉信徒,「要將身體獻上,當作活祭」,神要用我們這肉身活著的人,祂不用天使,因為天使沒有我們的身體；神要我們這在肉身活著的人來彰顯祂的榮美.聖經說:有人問你心中盼望的緣由,就要常作準備,以溫柔敬畏的心回答各人(彼前三18).這就是福音預工,就是基督的先鋒,因為我們有盼望,基督將要再來.
\newline
\newline
現在「不法的事增多,很多人愛心冷淡」,這「不法的事」不是偏重教外,而是指教內,教外的無所謂「愛心」；不認識十字架的人所說的愛,全是條件交換.這「愛心」是指教內的,今日的現象正是基督快來的預兆；讓我們作今日的施洗約翰吧.基督第一次來臨,祂只需要一個先鋒；但祂再來臨時,祂要千千萬萬的先鋒.為他預備道路,將人僵硬的心,反叛神的心解開.
\newline
\newline


\newpage

\section{第60屆~港九培靈會~陳終道~第七講}
\label{sec:1178}
\textbf{勇敢殉道的司提反}
\newline
\newline
連結: \href{https://www.hkbibleconference.org/session-message/view/1178}{\texttt{https://www.hkbibleconference.org/session-message/view/1178}}
\newline
\newline
\hyperref[sec:1177]{< < < PREV SERMON < < <}
~
\hyperref[sec:toc]{[返目錄]}
~
\hyperref[sec:1179]{> > > NEXT SERMON > > >}
\newline
\newline
使徒行傳 6:3-6:6
\newline
\begin{longtable}{cl}
\hline
\hline
章節 & 經文 (和合本修訂版)\\
\hline
6:3 & \begin{tabularx}{0.7\textwidth}{X} 所以弟兄們，當從你們中間選出七個有好名聲、滿有聖靈和智慧，我們派他們管理這事。 \end{tabularx} \\ \\ \relax
6:4 & \begin{tabularx}{0.7\textwidth}{X} 至於我們，我們要專注於祈禱和傳道的事奉。」 \end{tabularx} \\ \\ \relax
6:5 & \begin{tabularx}{0.7\textwidth}{X} 這話使全會眾都喜悅，就揀選了司提反—他是一個滿有信心和聖靈的人；他們又揀選了腓利、伯羅哥羅、尼迦挪、提門、巴米拿，並皈依猶太教的安提阿人尼哥拉， \end{tabularx} \\ \\ \relax
6:6 & \begin{tabularx}{0.7\textwidth}{X} 叫他們站在使徒面前，使徒禱告後，就為他們按手。 \end{tabularx} \\ \\
[1ex]
\hline
\hline
\end{longtable}
司提反是初期教會選出的首批執事.經文中,我們清楚知道,歷史第一位殉道者並不是傳道人,而是執事司提反.執事除了是指教會中的職份外,還有另一種意思；保羅也時常喜歡以「執事」來自稱；意思是服事人的服役者；謙虛地,風塵僕僕地服事神的人.此外,保羅的用法也含有職份的意義；保羅雖然不是被選的執事,但他的心態是屬神的,他在神面前有同樣的負擔和託付.領受了神交給他的責任,常把自己看成是應該服事人,將福音奧秘向人宣講的執事.
\newline
\newline
司提反未獲選執事前,在教會裡他也是向神負責任的.職份有它的好處,因為是別人給了我們職份,既要為這職份而忠心；不單對神有交代；同時,也要對人有所交代.假若我們在教會裡並不是執事,沒有職份,卻仍留意神家裡的需要,按自己的能力而愛神的家,向神盡忠,既不是因為別人給了職份,不是受著職份的驅策而負責；而是為愛主的緣故,愛護神的家,主動去填補教會的漏洞.這些人雖沒有執事的職份,但在神面前已經有執事的心志,是事奉神的人,司提反獲選為執事,正表明了他在平時已是個忠心事奉神的人.
\newline
\newline
現在,我們要看看司提反為人的幾個特別之處；第一,有好的名聲.這是七位被選執事的首要條件.所謂好名聲,並非俗世的名聲,而是有好的生活見證,美好的信心,品德；具備各種因信仰而表現的美德,因而得到別人的稱讚及敬佩.不但在教會內有愛主的表現,在家庭及社會中,在工作的崗位上能盡責,才稱得上有好名聲.
\newline
\newline
第二,被聖靈充滿.當時被選執事的主要任務是管理飯食,因他們過著凡物公用的生活.希伯來書第十章追述初期的希伯來信徒,因信主的緣故家業被人搶去,他們甘心樂意忍受這些痛苦.使徒行傳第四章,就是記述基督徒受猶太教徒嚴酷的逼迫；一部份信徒的家產被人強搶；仍有家產的信徒則把餘下的賣掉,分給別人,過著大家庭的生活.那時,信徒數目最少也有一萬之眾,這樣大的團體,需要有專門管理飯食的人.於是,彼得向眾門徒說,他們撇下神的道來管理飯食是不合宜的.必須在信徒中間,選出七個執事來管理飯食.要知道,管理這麼多人的飯食殊不容易,各人口味不同,同樣一窩飯；有人嫌軟,有人嫌硬,無論怎樣也有人埋怨的.因此,管理飯食也需要有聖靈充滿,要任勞任怨,能擔當起別人的怨言.就像在教會做事,做好了別人不知道,做得不好就要受別人責備；若不是聖靈充滿,實在是無法忍受,擔當不起的事.
\newline
\newline
第三,大有信心.門徒數目日漸增多,管理飯食的人就要用心去面對他們的需要,同時,也要應付口糧缺乏的難題.事實上,若果在飯食的事情上有信心,作別的事情也能有信心倚靠神.不要以為飯食是小事情,不少人為此而跌倒,基督接受的第一個試探,就是將石頭變為食物.我們若曉得在日用的飲食上倚靠神,在別的事上也有信心.
\newline
\newline
第四,熟識聖經.第七章中,司提反在毫無準備的辯論中,由亞伯拉罕至基督幾千年的歷史,他不但非常精通,能扼要地說出,且明白舊約的精義乃是預言基督.今天的基督徒太忽視聖經,缺乏天國的文化,只有世俗的文化,忽略聖經的重要,往往片段地了解聖經的內容；所以才會產生不少的異端.基督徒不熟識明白聖經,便不會有分辨能力；無法以聖經的原則建立基督徒的人生觀,以及獨特的生活方式.當時教會能夠吸引人,因為他們能顯出基督徒文化與世俗的分別,按聖經的觀念來建立自己的人生；因此,旁人就奇怪這些人哪裡來的巨大力量及吸引力,能聚成一個如此相愛的團體.
\newline
\newline
第五,勇敢指責罪惡.(徒七51)記載司提反的話非常厲害,以致群眾聽見就恨到咬牙切齒,用石頭攻擊他.然而,司提反是按照聖靈的感動而說話,神要藉著他責備那些硬心抗拒聖靈的人.司提反雖然是負責管理飯食的執事,但他也為福音作見證,按真理說話.在教會內他是執事,管理飯食,照顧弟兄姊妹的生活需要；另一方面,他仍有著一個當然的責任,是所有基督徒應該有的,就是為神的真理作見證.
\newline
\newline
第六,司提反的面貌好像天使的面貌；意即他的面貌帶有屬神的善良,屬天的美好.況且是那些逼害他的人所說的；可見他連責備人時的面貌也像天使.所謂「相由心生」,司提反被石頭攻擊,他最後的祈禱就是求神「不要將這罪歸他們」,這話就像基督臨終時為罪人祈禱一般.可見這位新約的殉道士,面對死亡時非常從容,絕不勉強.當我精讀這段聖經時,就懷疑當一個人面對死亡時,能否像司提反這樣坦然無懼呢.後來,我才真正感受到,如果我們的心是向神坦然,則能對死毫不畏懼.以下是我的見證:
\newline
\newline
大概五年前,我要接受一項心臟手術,雖然在西方這手術是司空見慣的,但也算是大手術,正常也要留院三,四星期.但我內心感受到一種平安,是我這一生也未曾享受過的；我完全不感到懼怕,反而是別人為我擔心.手術的過程十分複雜,先要將心臟取出急凍,再由大腿抽出血管駁上心臟,這種手術具有相當的危險性.但神給我平安,使我對死亡毫不畏懼,完全不需要安慰.正因這樣,我痊癒的速度十分快,不消十天便出院,第十五天就站在台上作見證了.感謝神,我真正感受到主所賜的平安,這平安是自然的；只要預備好,向神坦然,平安就自然來臨.我們是主的人,我們預備好了嗎?若果今天神就召我們見祂,我們能否坦然見祂?我們活了這麼久,作了這許多年的基督徒.曾否作過何事是可存到永遠；以致當我們見主時,也可以坦然無懼呢?
\newline
\newline


\newpage

\section{第60屆~港九培靈會~陳終道~第八講}
\label{sec:1179}
\textbf{外邦使徒保羅}
\newline
\newline
連結: \href{https://www.hkbibleconference.org/session-message/view/1179}{\texttt{https://www.hkbibleconference.org/session-message/view/1179}}
\newline
\newline
\hyperref[sec:1178]{< < < PREV SERMON < < <}
~
\hyperref[sec:toc]{[返目錄]}
~
\hyperref[sec:1180]{> > > NEXT SERMON > > >}
\newline
\newline
使徒行傳 9:1-9:9
\newline
\begin{longtable}{cl}
\hline
\hline
章節 & 經文 (和合本修訂版)\\
\hline
9:1 & \begin{tabularx}{0.7\textwidth}{X} 掃羅不斷用威嚇兇悍的口氣向主的門徒說話。他去見大祭司， \end{tabularx} \\ \\ \relax
9:2 & \begin{tabularx}{0.7\textwidth}{X} 要求發信給大馬士革的各會堂，若是找著信奉這道的人，無論男女，都准他捆綁帶到耶路撒冷。 \end{tabularx} \\ \\ \relax
9:3 & \begin{tabularx}{0.7\textwidth}{X} 掃羅在途中，將到大馬士革的時候，忽然有一道光從天上下來，四面照射著他， \end{tabularx} \\ \\ \relax
9:4 & \begin{tabularx}{0.7\textwidth}{X} 他就仆倒在地，聽見有聲音對他說：「掃羅！掃羅！你為甚麼迫害我？」 \end{tabularx} \\ \\ \relax
9:5 & \begin{tabularx}{0.7\textwidth}{X} 他說：「主啊！你是誰？」主說：「我就是你所迫害的耶穌。 \end{tabularx} \\ \\ \relax
9:6 & \begin{tabularx}{0.7\textwidth}{X} 起來！進城去，你應該做的事，必有人告訴你。」 \end{tabularx} \\ \\ \relax
9:7 & \begin{tabularx}{0.7\textwidth}{X} 同行的人站在那裡，說不出話來，因為他們聽見聲音，卻看不見人。 \end{tabularx} \\ \\ \relax
9:8 & \begin{tabularx}{0.7\textwidth}{X} 掃羅從地上起來，睜開眼睛，竟不能看見甚麼。有人拉他的手，領他進了大馬士革。 \end{tabularx} \\ \\ \relax
9:9 & \begin{tabularx}{0.7\textwidth}{X} 他三天甚麼都看不見，也不吃也不喝。 \end{tabularx} \\ \\
[1ex]
\hline
\hline
\end{longtable}
保羅是希伯來裔的羅馬籍人.他宣稱自己是便雅憫支派的人,又是第八日受割禮的.表明自己是在一個嚴格,完全按照律法禮儀的家庭出生的.當時,保羅在著名的教法師迦瑪列門下受教,品學優異.經文講述主耶穌向保羅顯現,耶穌雖沒有直接指正保羅的罪；但在主的光照中,他才明白到自己所逼迫的,乃是主耶穌.於是,他真心認罪悔改.一個人要成為基督徒的時候,他也必定需要真正的知罪悔改,這是聖靈光照的結果.使人看見自己是罪人,需要接納基督作為救主；保羅很清楚地有了這經驗.
\newline
\newline
另外,保羅的得救是與蒙召一起的,他得救時,已聽到耶穌吩咐他到外邦人中傳福音,使他們從黑暗歸向光明.使徒行傳9-15章之間,相隔了約十四年,其中他在亞拉伯住了三年.(加一15-18)新約中有十三封書信的信息,對福音的真理,救恩的原理,及新舊約的關係,講解得很清楚；相信是這三年中,神所給保羅的啟示和教導.可見保羅得救,蒙召及受訓練都是十分特別的,他對舊約極深造詣的背景,都是神預備他將來成為福音使者的重要操練.
\newline
\newline
後來,他在本鄉大數城的地方住了十年,跟著才到安提阿做教師,做二,三年牧養工作.然後,聖靈才差派巴拿巴與他一同出外佈道；那時他才被稱為使徒.雖然他得救時已清楚蒙召,要作外邦的使徒；但他一直默默地耕耘,直至神的時候到了,才奉派作外邦的使徒.保羅是以:講教義,講律法與恩典的關係,講基督救贖的原理而著名.另外,他的文字事奉也是特出的恩賜；後來,他在工作中顯出神蹟能力,醫治病人,懲罰罪惡；這都是在他作外邦使徒後才顯明的恩賜.在使徒中,保羅為主受苦是最多的,他以百般的忍耐,作為自己是使徒身份的憑據.還有,他最擅長開荒工作,在完全沒有聽過福音的地方建立教會.真正是開荒的大使徒.以上種種的恩賜,都是使徒行傳十三章以後的事.在這之前,他謙卑安靜地在不出名的教會中默默耕耘.
\newline
\newline
有些弟兄姊妹喜歡知道自己的恩賜,然後才在教會中按自己的恩賜事奉.究竟恩賜是在事奉中知道的呢?還是先知道了才能事奉呢?不少基督徒與其說知道自己的恩賜在先,倒不如說是知道自己的興趣罷了!其實,事奉不是根據興趣,而是根據神的安排,恩賜是特殊的恩典,與事奉是分不開的.不信基督的人也有他們的特長,可能比基督徒還要突出；天賦才能未必一定就是恩賜；除非將這些才能獻給神使用,才能成為神恩賜的一部份.要在教會內事奉,才是由聖靈所賜；為事奉的需要而給予的,才是恩賜.保羅被差遣時,才顯出使徒所需的恩賜；這一切,神都按照他工作的需要賜予他.
\newline
\newline
另一方面,我們看保羅對自己同胞的愛心.一個人要有愛靈魂的心,才能傳福音；如果我們不愛自己同胞,就更不會愛遠方不認識的人.所以必須先愛附近的,然後才到遠處的；這是人感情自然的發展.保羅是外邦的使徒,但他受到的逼迫,主要是來自猶太人.然而,保羅深愛猶太同胞,甚至願意為著他們的得救,自己與基督分離也在所不計.會堂是猶太人聚會的地方,保羅幾次入猶太人的會堂,向他們傳福音；但每次都受著猶太人的逼迫.雖然這樣,保羅在外邦傳道,也是先向猶太人傳講,保羅是奉召作為外邦的使徒.他本可以不向猶太人傳福音,也免受他們帶起的逼迫；因為保羅所傳的福音,令猶太人可誇的地方,完全一筆勾消.他們以為自己是神唯一的選民,便看不起外邦人.保羅所傳講的,使他們看見所有的人都是罪人；連猶太人也包括在內,這是他們不能接受的.保羅雖有羅馬籍,但他沒有忘記自己是希伯來人,他負上為希伯來人傳福音的責任.今天,我們有否愛同胞的心呢?我們與生俱來都是中國人,我們對同胞就有當然的責任；全世界最多的是中國人,只有中國人向中國人傳福音是最有效的.
\newline
\newline
保羅為基督受苦比任何人更多.哥林多後書十一章保羅說:「我比他們多受勞苦,多下監牢,受鞭打是過重的,冒死是屢次有的,被猶太人打五次......又屢次行遠路,遭江河的的危險,曠野的危險,海中的危險,假弟兄的危險,受勞碌,受困苦,多次不得睡,又飢又渴,多次不得食,受寒冷,赤身露體,除了這外面的事,還有為眾教會掛心的事,天天壓在我身上.有誰軟弱我不軟弱呢?有誰跌到我不焦急呢?」相信我們連百分之一也比不上保羅.只是,保羅受這麼多苦,仍始終忠心不移.試想,我們捱了打還有膽再傳福音嗎?保羅向主絕對忠心；因此,到了生命的盡頭,他能放膽地說:「那美好的仗我已經打過了,當跑的路我已經跑盡了,所信的道我已經守住了.從此以後,有公義的冠冕為我存留,就是按著公義審判的主,到了那日要賜給我的,不但賜給我,也賜給凡愛慕祂顯現的人.」(提後四7-8)
\newline
\newline
我們每個人都有自己當打的仗,當跑的路,並非隨意地跑；而是跑主所要我跑的路.保羅的一生,並非活得很長久；他殉道時,大約是主後六十多年,而他又是比較年輕的時候蒙召,雖然聖經沒有明顯記載他的年歲,但預料他年紀並不太大,因為司提反殉道時,保羅只是個替人看管衣裳的少年人(徒七58).然而,他一生工作的果效卻存到今日.每當我們提到使徒保羅時,沒有法子不覺得慚愧；保羅為著愛主愛人的靈魂,願意付上代價.求主今日施恩與我們,在神的話語上,在傳福音的工作上,讓我們效法保羅.我們今天的遭遇比起保羅好得多了.願意我們能有愛主愛人靈魂的心,在行完我們當行的路時；有一天,能像保羅所說,從此以後,有公義的冠冕為我存留.
\newline
\newline


\newpage

\section{第60屆~港九培靈會~陳終道~第九講}
\label{sec:1180}
\textbf{永垂千古的平凡人}
\newline
\newline
連結: \href{https://www.hkbibleconference.org/session-message/view/1180}{\texttt{https://www.hkbibleconference.org/session-message/view/1180}}
\newline
\newline
\hyperref[sec:1179]{< < < PREV SERMON < < <}
~
\hyperref[sec:toc]{[返目錄]}
~
\hyperref[sec:1124]{> > > NEXT SERMON > > >}
\newline
\newline
這一講我們要看幾位在聖經中不太特出,但蒙神使用的人物.
\newline
\newline
(一)主的母親馬利亞
\newline
\newline
(路一38)記馬利亞回答天使時,說了一句寶貴的話:「我是主的使女,情願照你的話成就在我身上.」假如馬利亞不說這話,聖靈就不會藉著她使基督道成肉身.相信我們很少思想這句話的嚴重性；事實上,馬利亞接受使命是要準備付代價的,可能會帶來巨大的逼迫.童女懷孕,必要冒著任何人間的艱難,誤會,委屈,閒言蜚語；只是,馬利亞已預備好心靈去接受這一切.當我們說要奉獻自己給主使用,心裡怎樣想呢?要像葛培理,宋尚節那樣受主重用嗎?我們時常憧憬會受到重視,歡迎；可有想過,奉獻給主使用時,要面對無法估計的逼迫與反對；這時,我們會否仍讓主使用?馬利亞願意,聖靈就在那刻感孕她.馬利亞奉獻的結果,並沒有遭遇她所預計的困難.神已為她安排,解除了她可能要面對的困擾；天使跟著向約瑟報夢,約瑟迎娶了馬利亞.當馬利亞預備要為主受苦時,神沒有使她受到超過她所能承擔的苦.
\newline
\newline
馬利亞的奉獻是空前絕後的,這是何等光榮的奉獻,就是出於一個真純奉獻的心而完成的.她不是大使徒,不會講道,沒有政治地位,又不富裕,只是木匠的妻子而已.一位平凡的女子,她在屬靈方面的貢獻,卻沒有人能與之相比較.今天,我們雖然沒有她的福份,但我們有沒有可奉獻的呢?
\newline
\newline
(二)朝見耶穌的博士
\newline
\newline
博士朝見耶穌是我們耳熟能詳的故事.聖經上說,他們是看到星而來的；根據我個人的看法,那夥星絕不是天上普通的星,因為就是月亮的光度,也不可能指出耶穌誕生的正確位置,我們也不能藉月亮的光找浸信會.先知巴蘭曾預言:「有星出於雅各」,這「星」就是指基督,因此,我相信他們是從舊約中得到啟示而來的.博士清楚了解啟示；他們要朝拜的是嬰孩耶穌；可貴之處是他們帶著貴重的物品,從老遠到來,見到的是生於馬槽的嬰孩.但他們並沒有失望,仍將黃金,乳香,沒藥毫無保留地獻上.不但沒有得到任何好處而空手回去,還要冒著被希律王追殺的危險.今天,不少人奉獻若得不到神的賜福,立刻就埋怨了,這種奉獻有問題.博士的奉獻,是把握了那歷史性的一刻,神要他們作那件事,他們做了,歷史上再沒有別人能得到那樣的機會和地位,他們的事蹟就存到永永遠遠.
\newline
\newline
(三)奉獻五餅二魚的小孩
\newline
\newline
五餅二魚餵飽五千男丁是耶穌最出名的神蹟之一,但這神蹟的關鍵,卻在於那小孩獻出五餅二魚才達成.那時,耶穌向群眾講道,他們事前是毫無準備,沒有帶備口糧的.黃昏了,眾人都餓了,其中只有一位孩子有五餅二魚；試想,成年人尚且感到餓,何況是小孩子呢?相信那小孩也想像得到；自己拿了僅有的五餅二魚出來,分給這麼多人,自己定要捱餓.但是,他還是把五餅二魚奉獻出來.這五餅二魚在小孩子的手裡只夠他自己吃,但在耶穌的手中則不單餵飽五千人,更剩下零碎的十二籃.此外,這神蹟還有特別之處；當時,天快黑了,耶穌叫他們坐下,然後望天祝謝；擘開餅分給眾人.神蹟的重點,在於將餅分給眾人時,餅不會比原來的少,反而愈來愈多.換句話說,這神蹟是透過他們中間無私的分享而顯出來的.反之,若有人接到餅,拿了自己的份量後,沒有將餘下的遞給別人；這樣,神蹟就停在那裡,無法繼續了.我們可有接到別人的餅後便留下來,沒有繼續分與他人享用呢?
\newline
\newline
(四)伯大尼的馬利亞
\newline
\newline
我第一次唸這段經文時,很奇怪何以馬利亞用香膏澆耶穌有重大的價值.基督來到世上為人受死,以致復活,完成救恩,這是基督在世的使命.耶穌曾經多次向門徒解釋祂要受苦.第三天復活,但門徒仍不明白,也不重視祂的話.馬利亞顯然知道耶穌要受死及復活,因此,她早就盡她所能,買極貴重的真哪噠香膏獻給耶穌.在耶穌坐席時,她打破玉瓶,把香膏澆在主身上.門徒批評她浪費,耶穌卻為她申辯,稱讚她預先為主澆上香膏.馬利亞知道耶穌將要受死,而門徒卻不明白,她給予耶穌的安慰大於那香膏的價值.以後,雖有尼哥底母及抹大拉的馬利亞獻上香膏,但耶穌真正享受的,就只有這瓶真哪噠香膏.
\newline
\newline
(五)為主背十字架的古利奈人西門
\newline
\newline
在耶穌前往各各他的路上,兵丁找著路過的古利奈人西門,強迫他為耶穌背十字架.在聖經中,這是唯一提到有人被迫背上十字架.其餘的一向是人自願選擇,耶穌上十字架受苦也是自願的.但是,西門卻是被兵丁抓著,背負十字架.相信在途中必定有不少怨言.但到了各各他,看到耶穌被釘在十字架,看見天地從午正到申初忽然黑暗,看見耶穌斷氣時天崩地裂,不禁想到被釘十字架的耶穌確是神的兒子.西門愈想愈覺得自己能為耶穌背十字架,實在是幸運,光榮,他不禁想起自己算是甚麼,竟能無意地為主背十字架；再也沒有人有這等機會.西門不是出於自願,也沒有馬利亞的細心,但他卻為耶穌背負十架；成為空前絕後的偉事.
\newline
\newline
以上所說都是平凡人,但他們所做的卻能永垂千古.我們不要少看自己的平凡,若我們誠心誠意愛主,毫無私心,真誠的奉獻,存這心態事奉；我們在神的家裡都有自己獨特的貢獻.只要存著愛主的心,甘願為主擺上一切的心,就足夠了.弟兄姊妹,所有平凡的人,都可以在神的國度裡有不平凡的貢獻!問題是我們是否願意在以後的日子中對主說:我們所做的一切不是為自己,不是為人的榮耀,為教會的虛榮,為地位,權勢.......而是一心一意地做在主的身上,神就有特別使用我們的地方.
\newline
\newline


\newpage

\section{第61屆~港九培靈會~劉銳光~第一講}
\label{sec:1124}
\textbf{同心}
\newline
\newline
連結: \href{https://www.hkbibleconference.org/session-message/view/1124}{\texttt{https://www.hkbibleconference.org/session-message/view/1124}}
\newline
\newline
\hyperref[sec:1180]{< < < PREV SERMON < < <}
~
\hyperref[sec:toc]{[返目錄]}
~
\hyperref[sec:1125]{> > > NEXT SERMON > > >}
\newline
\newline
使徒行傳 1:14-1:14
\newline
\begin{longtable}{cl}
\hline
\hline
章節 & 經文 (和合本修訂版)\\
\hline
1:14 & \begin{tabularx}{0.7\textwidth}{X} 這些人和幾個婦人，包括耶穌的母親馬利亞，和耶穌的兄弟，都同心合意地恆切禱告。 \end{tabularx} \\ \\
[1ex]
\hline
\hline
\end{longtable}


\newpage

\section{第61屆~港九培靈會~劉銳光~第二講}
\label{sec:1125}
\textbf{相愛}
\newline
\newline
連結: \href{https://www.hkbibleconference.org/session-message/view/1125}{\texttt{https://www.hkbibleconference.org/session-message/view/1125}}
\newline
\newline
\hyperref[sec:1124]{< < < PREV SERMON < < <}
~
\hyperref[sec:toc]{[返目錄]}
~
\hyperref[sec:1127]{> > > NEXT SERMON > > >}
\newline
\newline
使徒行傳 2:43-2:47
\newline
\begin{longtable}{cl}
\hline
\hline
章節 & 經文 (和合本修訂版)\\
\hline
2:43 & \begin{tabularx}{0.7\textwidth}{X} 眾人都心存敬畏；使徒們又行了許多奇事神蹟。 \end{tabularx} \\ \\ \relax
2:44 & \begin{tabularx}{0.7\textwidth}{X} 信的人都聚在一處，凡物公用， \end{tabularx} \\ \\ \relax
2:45 & \begin{tabularx}{0.7\textwidth}{X} 又賣了田產和家業，照每一個人所需要的分給他們。 \end{tabularx} \\ \\ \relax
2:46 & \begin{tabularx}{0.7\textwidth}{X} 他們天天同心合意恆切地在聖殿裡敬拜，且在家中擘餅，存著歡喜坦誠的心用飯， \end{tabularx} \\ \\ \relax
2:47 & \begin{tabularx}{0.7\textwidth}{X} 讚美神，得全體百姓的喜愛。主將得救的人天天加給他們。 \end{tabularx} \\ \\
[1ex]
\hline
\hline
\end{longtable}


\newpage

\section{第61屆~港九培靈會~劉銳光~第三講}
\label{sec:1127}
\textbf{禱告}
\newline
\newline
連結: \href{https://www.hkbibleconference.org/session-message/view/1127}{\texttt{https://www.hkbibleconference.org/session-message/view/1127}}
\newline
\newline
\hyperref[sec:1125]{< < < PREV SERMON < < <}
~
\hyperref[sec:toc]{[返目錄]}
~
\hyperref[sec:1130]{> > > NEXT SERMON > > >}
\newline
\newline
使徒行傳 1:14-1:15
\newline
\begin{longtable}{cl}
\hline
\hline
章節 & 經文 (和合本修訂版)\\
\hline
1:14 & \begin{tabularx}{0.7\textwidth}{X} 這些人和幾個婦人，包括耶穌的母親馬利亞，和耶穌的兄弟，都同心合意地恆切禱告。 \end{tabularx} \\ \\ \relax
1:15 & \begin{tabularx}{0.7\textwidth}{X} 那時，有許多人聚會，約有一百二十名，彼得在弟兄中間站起來，說： \end{tabularx} \\ \\
[1ex]
\hline
\hline
\end{longtable}


\newpage

\section{第61屆~港九培靈會~劉銳光~第四講}
\label{sec:1130}
\textbf{能力}
\newline
\newline
連結: \href{https://www.hkbibleconference.org/session-message/view/1130}{\texttt{https://www.hkbibleconference.org/session-message/view/1130}}
\newline
\newline
\hyperref[sec:1127]{< < < PREV SERMON < < <}
~
\hyperref[sec:toc]{[返目錄]}
~
\hyperref[sec:1132]{> > > NEXT SERMON > > >}
\newline
\newline
使徒行傳 1:8-1:33
\newline
\begin{longtable}{cl}
\hline
\hline
章節 & 經文 (和合本修訂版)\\
\hline
1:8 & \begin{tabularx}{0.7\textwidth}{X} 但聖靈降臨在你們身上，你們就必得著能力，並要在耶路撒冷、猶太全地和撒瑪利亞，直到地極，作我的見證。」 \end{tabularx} \\ \\ \relax
1:9 & \begin{tabularx}{0.7\textwidth}{X} 說了這些話，他們正看的時候，他被接上升，有一朵雲彩從他們眼前把他接去。 \end{tabularx} \\ \\ \relax
1:10 & \begin{tabularx}{0.7\textwidth}{X} 他升上去的時候，他們定睛望天，看哪，有兩個人身穿白衣站在他們旁邊， \end{tabularx} \\ \\ \relax
1:11 & \begin{tabularx}{0.7\textwidth}{X} 說：「加利利人哪，你們為甚麼站著望天呢？這離開你們被接升天的耶穌，你們見他怎樣升上天去，他也要怎樣來臨。」 \end{tabularx} \\ \\ \relax
1:12 & \begin{tabularx}{0.7\textwidth}{X} 有一座山，名叫橄欖山，離耶路撒冷不遠，有安息日可行走的路程。那時，門徒從那裡回耶路撒冷去， \end{tabularx} \\ \\ \relax
1:13 & \begin{tabularx}{0.7\textwidth}{X} 他們一進城，就上了所住的樓房；在那裡有彼得、約翰、雅各、安得烈、腓力、多馬、巴多羅買、馬太、亞勒腓的兒子雅各、激進黨的西門，和雅各的兒子猶大。 \end{tabularx} \\ \\ \relax
1:14 & \begin{tabularx}{0.7\textwidth}{X} 這些人和幾個婦人，包括耶穌的母親馬利亞，和耶穌的兄弟，都同心合意地恆切禱告。 \end{tabularx} \\ \\ \relax
1:15 & \begin{tabularx}{0.7\textwidth}{X} 那時，有許多人聚會，約有一百二十名，彼得在弟兄中間站起來，說： \end{tabularx} \\ \\ \relax
1:16 & \begin{tabularx}{0.7\textwidth}{X} 「諸位弟兄，聖經的話必須應驗。聖經中，聖靈曾藉大衛的口預先說到那領人來拿耶穌的猶大； \end{tabularx} \\ \\ \relax
1:17 & \begin{tabularx}{0.7\textwidth}{X} 他本來算是我們中的一個，並且得了這一份使徒的職任。 \end{tabularx} \\ \\ \relax
1:18 & \begin{tabularx}{0.7\textwidth}{X} 這人用他不義的代價買了一塊田，以後身子仆倒，肚腹崩裂，腸子都流出來。 \end{tabularx} \\ \\ \relax
1:19 & \begin{tabularx}{0.7\textwidth}{X} 住在耶路撒冷的人都知道這事，所以按著他們當地的話把那塊田叫亞革大馬，就是「血田」的意思。 \end{tabularx} \\ \\ \relax
1:20 & \begin{tabularx}{0.7\textwidth}{X} 因為《詩篇》上寫著：「願他的住處變為廢墟，無人在內居住。」又說：「願別人得他的職分。」 \end{tabularx} \\ \\ \relax
1:21 & \begin{tabularx}{0.7\textwidth}{X} 所以，主耶穌在我們中間出入的整段時間， \end{tabularx} \\ \\ \relax
1:22 & \begin{tabularx}{0.7\textwidth}{X} 就是從約翰施洗起，直到主離開我們被接上升的日子為止，必須從那常與我們一起的人中，立一位與我們同作耶穌復活的見證。」 \end{tabularx} \\ \\ \relax
1:23 & \begin{tabularx}{0.7\textwidth}{X} 於是他們推舉兩個人，就是那叫巴撒巴，又稱為猶士都的約瑟，和馬提亞。 \end{tabularx} \\ \\ \relax
1:24 & \begin{tabularx}{0.7\textwidth}{X} 眾人禱告說：「主啊，你知道萬人的心，求你從這兩個人中指明你所揀選的是哪一位， \end{tabularx} \\ \\ \relax
1:25 & \begin{tabularx}{0.7\textwidth}{X} 去得這使徒的職任；這職位猶大已經丟棄，往自己的地方去了。」 \end{tabularx} \\ \\ \relax
1:26 & \begin{tabularx}{0.7\textwidth}{X} 於是眾人為他們搖籤，搖出馬提亞來；他就和十一個使徒同列。 \end{tabularx} \\ \\ \relax
2:1 & \begin{tabularx}{0.7\textwidth}{X} 五旬節那日到了，他們全都聚集在一起。 \end{tabularx} \\ \\ \relax
2:2 & \begin{tabularx}{0.7\textwidth}{X} 忽然，有響聲從天上下來，好像一陣大風吹過，充滿了他們所坐的整座屋子； \end{tabularx} \\ \\ \relax
2:3 & \begin{tabularx}{0.7\textwidth}{X} 又有舌頭如火焰向他們顯現，分開落在他們每個人身上。 \end{tabularx} \\ \\ \relax
2:4 & \begin{tabularx}{0.7\textwidth}{X} 他們都被聖靈充滿，就按著聖靈所賜的口才說起別國的話來。 \end{tabularx} \\ \\ \relax
2:5 & \begin{tabularx}{0.7\textwidth}{X} 那時，有從天下各國來的虔誠的猶太人，住在耶路撒冷。 \end{tabularx} \\ \\ \relax
2:6 & \begin{tabularx}{0.7\textwidth}{X} 這聲音一響，許多人都來聚集，各人因為聽見門徒用他們各自的鄉談說話，就甚納悶， \end{tabularx} \\ \\ \relax
2:7 & \begin{tabularx}{0.7\textwidth}{X} 都詫異驚奇說：「看哪，這些說話的不都是加利利人嗎？ \end{tabularx} \\ \\ \relax
2:8 & \begin{tabularx}{0.7\textwidth}{X} 我們每個人怎麼聽見他們說我們生來所用的鄉談呢？ \end{tabularx} \\ \\ \relax
2:9 & \begin{tabularx}{0.7\textwidth}{X} 我們帕提亞人、瑪代人、以攔人，和住在美索不達米亞、猶太、加帕多家、本都、亞細亞、 \end{tabularx} \\ \\ \relax
2:10 & \begin{tabularx}{0.7\textwidth}{X} 弗呂家、旁非利亞、埃及的人，並靠近古利奈的利比亞一帶地方的人，僑居的羅馬人， \end{tabularx} \\ \\ \relax
2:11 & \begin{tabularx}{0.7\textwidth}{X} 包括猶太人和皈依猶太教的人， 克里特人和阿拉伯人，都聽見他們用我們的鄉談講論神的大作為。」 \end{tabularx} \\ \\ \relax
2:12 & \begin{tabularx}{0.7\textwidth}{X} 眾人就都驚奇困惑，彼此說：「這是甚麼意思呢？」 \end{tabularx} \\ \\ \relax
2:13 & \begin{tabularx}{0.7\textwidth}{X} 還有人譏誚，說：「他們是灌滿了新酒吧！」 \end{tabularx} \\ \\ \relax
2:14 & \begin{tabularx}{0.7\textwidth}{X} 彼得和十一個使徒站起來，他就高聲向眾人說：「猶太人和所有住在耶路撒冷的人哪，這件事你們要知道，要側耳聽我的話。 \end{tabularx} \\ \\ \relax
2:15 & \begin{tabularx}{0.7\textwidth}{X} 這些人並不像你們所想的喝醉了，因為現在才早晨九點鐘。 \end{tabularx} \\ \\ \relax
2:16 & \begin{tabularx}{0.7\textwidth}{X} 這正是藉著先知約珥所說的： \end{tabularx} \\ \\ \relax
2:17 & \begin{tabularx}{0.7\textwidth}{X} 『神說：在末後的日子，我要將我的靈澆灌凡血肉之軀的。你們的兒女要說預言；你們的少年要見異象；你們的老人要做異夢。 \end{tabularx} \\ \\ \relax
2:18 & \begin{tabularx}{0.7\textwidth}{X} 在那些日子，我要把我的靈澆灌，甚至給我的僕人和婢女，他們要說預言。 \end{tabularx} \\ \\ \relax
2:19 & \begin{tabularx}{0.7\textwidth}{X} 在天上，我要顯出奇事，在地下，我要顯出神蹟，有血，有火，有煙霧。 \end{tabularx} \\ \\ \relax
2:20 & \begin{tabularx}{0.7\textwidth}{X} 太陽要變為黑暗，月亮要變為血，這都在主大而光榮的日子未到以前。 \end{tabularx} \\ \\ \relax
2:21 & \begin{tabularx}{0.7\textwidth}{X} 那時，凡求告主名的都必得救。』 \end{tabularx} \\ \\ \relax
2:22 & \begin{tabularx}{0.7\textwidth}{X} 「以色列人哪，你們要聽我這些話：拿撒勒人耶穌就是神以異能、奇事、神蹟向你們證明出來的人，這些事是神藉著他在你們中間施行，正如你們自己知道的。 \end{tabularx} \\ \\ \relax
2:23 & \begin{tabularx}{0.7\textwidth}{X} 他既按著神確定的旨意和預知被交與人，你們就藉著不法之人的手把他釘在十字架上，殺了。 \end{tabularx} \\ \\ \relax
2:24 & \begin{tabularx}{0.7\textwidth}{X} 神卻將死的痛苦解除，使他復活了，因為他原不能被死拘禁。 \end{tabularx} \\ \\ \relax
2:25 & \begin{tabularx}{0.7\textwidth}{X} 大衛指著他說：『我看見主常在我眼前，他在我右邊，使我不至於動搖。 \end{tabularx} \\ \\ \relax
2:26 & \begin{tabularx}{0.7\textwidth}{X} 所以我心裡歡喜，我的舌頭快樂，而且我的肉身要安居在指望中。 \end{tabularx} \\ \\ \relax
2:27 & \begin{tabularx}{0.7\textwidth}{X} 因你必不將我的靈魂撇在陰間，也不讓你的聖者見朽壞。 \end{tabularx} \\ \\ \relax
2:28 & \begin{tabularx}{0.7\textwidth}{X} 你已將生命的道路指示我，必使我在你面前充滿快樂。』 \end{tabularx} \\ \\ \relax
2:29 & \begin{tabularx}{0.7\textwidth}{X} 「諸位弟兄，先祖大衛的事，我可以坦然地對你們說：他死了，也埋葬了，而且他的墳墓直到今日還在我們這裡。 \end{tabularx} \\ \\ \relax
2:30 & \begin{tabularx}{0.7\textwidth}{X} 既然大衛是先知，他知道神曾向他起誓，要從他的後裔中立一位坐在他的寶座上。 \end{tabularx} \\ \\ \relax
2:31 & \begin{tabularx}{0.7\textwidth}{X} 他預先看見了，就講論基督的復活，說：『他不被撇在陰間；他的肉身也不見朽壞。』 \end{tabularx} \\ \\ \relax
2:32 & \begin{tabularx}{0.7\textwidth}{X} 這耶穌，神已經使他復活了，我們都是這事的見證人。 \end{tabularx} \\ \\ \relax
2:33 & \begin{tabularx}{0.7\textwidth}{X} 他既被高舉在神的右邊，又從父受了所應許的聖靈，就把你們所看見所聽見的，澆灌下來。 \end{tabularx} \\ \\ \relax
2:34 & \begin{tabularx}{0.7\textwidth}{X} 大衛並沒有升到天上，但他自己說：『主對我主說：你坐在我的右邊， \end{tabularx} \\ \\ \relax
2:35 & \begin{tabularx}{0.7\textwidth}{X} 等我使你的仇敵作你的腳凳。』 \end{tabularx} \\ \\ \relax
2:36 & \begin{tabularx}{0.7\textwidth}{X} 故此，以色列全家當確實知道，你們釘在十字架上的這位耶穌，神已經立他為主，為基督了。」 \end{tabularx} \\ \\ \relax
2:37 & \begin{tabularx}{0.7\textwidth}{X} 眾人聽見這話，覺得扎心，就對彼得和其餘的使徒說：「諸位弟兄，我們該怎樣做呢？」 \end{tabularx} \\ \\ \relax
2:38 & \begin{tabularx}{0.7\textwidth}{X} 彼得對他們說：「你們各人要悔改，奉耶穌基督的名受洗，使你們的罪得赦免，就會領受所賜的聖靈。 \end{tabularx} \\ \\ \relax
2:39 & \begin{tabularx}{0.7\textwidth}{X} 因為這應許是給你們和你們的兒女，並一切在遠方的人，就是給所有主—我們的神所召來的人。」 \end{tabularx} \\ \\ \relax
2:40 & \begin{tabularx}{0.7\textwidth}{X} 彼得還用更多別的話作見證，勸勉他們說：「你們當救自己脫離這彎曲的世代。」 \end{tabularx} \\ \\ \relax
2:41 & \begin{tabularx}{0.7\textwidth}{X} 於是領受他話的人，都受了洗；那一天，門徒約添了三千人。 \end{tabularx} \\ \\ \relax
2:42 & \begin{tabularx}{0.7\textwidth}{X} 他們都專注於使徒的教導和彼此的團契，擘餅和祈禱。 \end{tabularx} \\ \\ \relax
2:43 & \begin{tabularx}{0.7\textwidth}{X} 眾人都心存敬畏；使徒們又行了許多奇事神蹟。 \end{tabularx} \\ \\ \relax
2:44 & \begin{tabularx}{0.7\textwidth}{X} 信的人都聚在一處，凡物公用， \end{tabularx} \\ \\ \relax
2:45 & \begin{tabularx}{0.7\textwidth}{X} 又賣了田產和家業，照每一個人所需要的分給他們。 \end{tabularx} \\ \\ \relax
2:46 & \begin{tabularx}{0.7\textwidth}{X} 他們天天同心合意恆切地在聖殿裡敬拜，且在家中擘餅，存著歡喜坦誠的心用飯， \end{tabularx} \\ \\ \relax
2:47 & \begin{tabularx}{0.7\textwidth}{X} 讚美神，得全體百姓的喜愛。主將得救的人天天加給他們。 \end{tabularx} \\ \\ \relax
3:1 & \begin{tabularx}{0.7\textwidth}{X} 下午三點鐘禱告的時候，彼得和約翰上聖殿去。 \end{tabularx} \\ \\ \relax
3:2 & \begin{tabularx}{0.7\textwidth}{X} 一個從母腹裡就是瘸腿的人正被人抬來，他們天天把他放在聖殿的一個叫美門的門口，求進聖殿的人施捨。 \end{tabularx} \\ \\ \relax
3:3 & \begin{tabularx}{0.7\textwidth}{X} 他看見彼得、約翰將要進聖殿，就求他們施捨。 \end{tabularx} \\ \\ \relax
3:4 & \begin{tabularx}{0.7\textwidth}{X} 彼得和約翰定睛看他，彼得說：「看著我們！」 \end{tabularx} \\ \\ \relax
3:5 & \begin{tabularx}{0.7\textwidth}{X} 那人就注目看他們，指望從他們得著甚麼。 \end{tabularx} \\ \\ \relax
3:6 & \begin{tabularx}{0.7\textwidth}{X} 彼得卻說：「金銀我都沒有，但我把我有的給你：奉拿撒勒人耶穌基督的名起來行走！」 \end{tabularx} \\ \\ \relax
3:7 & \begin{tabularx}{0.7\textwidth}{X} 於是彼得拉著他的右手，扶他起來；他的腳和踝骨立刻健壯了， \end{tabularx} \\ \\ \relax
3:8 & \begin{tabularx}{0.7\textwidth}{X} 就跳起來，站著，又開始行走。他跟他們進了聖殿，邊走邊跳，讚美神。 \end{tabularx} \\ \\ \relax
3:9 & \begin{tabularx}{0.7\textwidth}{X} 百姓都看見他又行走，又讚美神， \end{tabularx} \\ \\ \relax
3:10 & \begin{tabularx}{0.7\textwidth}{X} 認得他是那素常坐在聖殿的美門口求人施捨的，就因他所遇到的事滿心驚訝詫異。 \end{tabularx} \\ \\ \relax
3:11 & \begin{tabularx}{0.7\textwidth}{X} 那人正在稱為所羅門的廊下，拉住彼得和約翰，大家都覺得很驚訝，一齊跑到他們那裡。 \end{tabularx} \\ \\ \relax
3:12 & \begin{tabularx}{0.7\textwidth}{X} 彼得看見，就對百姓說：「以色列人哪，為甚麼因這事而驚訝呢？為甚麼定睛看我們，以為我們憑自己的能力和虔誠使這人行走呢？ \end{tabularx} \\ \\ \relax
3:13 & \begin{tabularx}{0.7\textwidth}{X} 亞伯拉罕的神、以撒的神、雅各的神，就是我們列祖的神，已經榮耀了他的僕人耶穌，這耶穌就是你們交付官府的那位，彼拉多決定要釋放他時，你們卻在彼拉多面前棄絕了他。 \end{tabularx} \\ \\ \relax
3:14 & \begin{tabularx}{0.7\textwidth}{X} 你們棄絕了那聖潔公義者，反而要求釋放一個兇手給你們。 \end{tabularx} \\ \\ \relax
3:15 & \begin{tabularx}{0.7\textwidth}{X} 你們殺了那生命的創始者，神卻叫他從死人中復活；我們都是這事的見證人。 \end{tabularx} \\ \\ \relax
3:16 & \begin{tabularx}{0.7\textwidth}{X} 因信他的名，他的名使你們所看見所認識的這人健壯了；正是他所賜的信心使這人在你們眾人面前完全好了。 \end{tabularx} \\ \\ \relax
3:17 & \begin{tabularx}{0.7\textwidth}{X} 「如今，弟兄們，我知道你們做這事是出於無知，你們的官長也是如此。 \end{tabularx} \\ \\ \relax
3:18 & \begin{tabularx}{0.7\textwidth}{X} 但神藉著眾先知的口預先宣告過基督將要受害的事，就這樣應驗了。 \end{tabularx} \\ \\ \relax
3:19 & \begin{tabularx}{0.7\textwidth}{X} 所以，你們當悔改歸正，使你們的罪得以塗去， \end{tabularx} \\ \\ \relax
3:20 & \begin{tabularx}{0.7\textwidth}{X} 這樣，那安舒的日子就必從主面前來到；主也必差遣所預定給你們的基督耶穌來臨。 \end{tabularx} \\ \\ \relax
3:21 & \begin{tabularx}{0.7\textwidth}{X} 他必須留在天上，直到萬物復興的時候，就是神自古藉著聖先知的口所說的。 \end{tabularx} \\ \\ \relax
3:22 & \begin{tabularx}{0.7\textwidth}{X} 摩西曾說：『主—你們的神要從你們弟兄中給你們興起一位先知像我，凡他向你們所說的一切，你們都要聽從。 \end{tabularx} \\ \\ \relax
3:23 & \begin{tabularx}{0.7\textwidth}{X} 凡不聽從那先知的，必將從民中滅絕。』 \end{tabularx} \\ \\ \relax
3:24 & \begin{tabularx}{0.7\textwidth}{X} 從撒母耳以來和後繼的眾先知，凡說預言的，也都曾宣告這些日子。 \end{tabularx} \\ \\ \relax
3:25 & \begin{tabularx}{0.7\textwidth}{X} 你們是先知的子孫，也是神與你們祖宗所立之約的子孫，就是對亞伯拉罕說：『地上萬族都將因你的後裔得福。』 \end{tabularx} \\ \\ \relax
3:26 & \begin{tabularx}{0.7\textwidth}{X} 神既興起他的僕人，就先差他到你們這裡來，賜福給你們，使各人回轉，離開你們的邪惡。」 \end{tabularx} \\ \\ \relax
4:1 & \begin{tabularx}{0.7\textwidth}{X} 彼得和約翰正向百姓說話的時候，祭司們、守殿官和撒都該人來了， \end{tabularx} \\ \\ \relax
4:2 & \begin{tabularx}{0.7\textwidth}{X} 就很煩惱，因為使徒們教導百姓，傳揚在耶穌的事上證明有死人復活， \end{tabularx} \\ \\ \relax
4:3 & \begin{tabularx}{0.7\textwidth}{X} 於是下手拿住他們；因為天已經晚了，就把他們押在拘留所到第二天。 \end{tabularx} \\ \\ \relax
4:4 & \begin{tabularx}{0.7\textwidth}{X} 但聽道的人有許多信了，男人的數目約有五千。 \end{tabularx} \\ \\ \relax
4:5 & \begin{tabularx}{0.7\textwidth}{X} 第二天，官長、長老和文士在耶路撒冷聚集， \end{tabularx} \\ \\ \relax
4:6 & \begin{tabularx}{0.7\textwidth}{X} 又有亞那大祭司、該亞法、約翰、亞歷山大，和大祭司的親族都在那裡。 \end{tabularx} \\ \\ \relax
4:7 & \begin{tabularx}{0.7\textwidth}{X} 他們叫使徒站在中間，問他們：「你們憑甚麼能力，奉誰的名做這事呢？」 \end{tabularx} \\ \\ \relax
4:8 & \begin{tabularx}{0.7\textwidth}{X} 那時，彼得被聖靈充滿，對他們說：「民間的官長和長老啊， \end{tabularx} \\ \\ \relax
4:9 & \begin{tabularx}{0.7\textwidth}{X} 倘若今日我們被查問是因為在殘障的人身上所行的善事，就是這人怎麼得了痊癒， \end{tabularx} \\ \\ \relax
4:10 & \begin{tabularx}{0.7\textwidth}{X} 那麼，你們大家和以色列全民都當知道，站在你們面前的這人得痊癒，是因你們所釘在十字架、神使他從死人中復活的拿撒勒人耶穌基督的名。 \end{tabularx} \\ \\ \relax
4:11 & \begin{tabularx}{0.7\textwidth}{X} 這位耶穌是：『你們匠人所丟棄的石頭，已成了房角的頭塊石頭。』 \end{tabularx} \\ \\ \relax
4:12 & \begin{tabularx}{0.7\textwidth}{X} 除他以外，別無拯救，因為在天下人間，沒有賜下別的名，我們可以靠著得救。」 \end{tabularx} \\ \\ \relax
4:13 & \begin{tabularx}{0.7\textwidth}{X} 他們見彼得、約翰的膽量，又看出他們原是沒有學問的平民，就很驚訝，認出他們曾是跟耶穌一起的； \end{tabularx} \\ \\ \relax
4:14 & \begin{tabularx}{0.7\textwidth}{X} 又看見那治好了的人和他們一同站著，就無話可駁。 \end{tabularx} \\ \\ \relax
4:15 & \begin{tabularx}{0.7\textwidth}{X} 於是他們吩咐他們兩人從議會退出，就彼此商議， \end{tabularx} \\ \\ \relax
4:16 & \begin{tabularx}{0.7\textwidth}{X} 說：「我們當怎樣辦這兩個人呢？因為他們誠然行了一件明顯的神蹟，凡住在耶路撒冷的人都知道，我們也不能否認。 \end{tabularx} \\ \\ \relax
4:17 & \begin{tabularx}{0.7\textwidth}{X} 但為避免這事越發在民間傳揚，我們必須威嚇他們，叫他們不可再奉這名對任何人講論。」 \end{tabularx} \\ \\ \relax
4:18 & \begin{tabularx}{0.7\textwidth}{X} 於是他們叫了兩人來，禁止他們，再不可奉耶穌的名講論或教導人。 \end{tabularx} \\ \\ \relax
4:19 & \begin{tabularx}{0.7\textwidth}{X} 彼得和約翰回答他們說：「聽從你們，不聽從神，在神面前合理不合理，你們自己判斷吧！ \end{tabularx} \\ \\ \relax
4:20 & \begin{tabularx}{0.7\textwidth}{X} 我們所看見所聽見的，我們不能不說。」 \end{tabularx} \\ \\ \relax
4:21 & \begin{tabularx}{0.7\textwidth}{X} 官長為百姓的緣故，想不出任何法子懲罰他們，只好威嚇一番就把他們釋放了；這是因眾人為了所行的奇事都歸榮耀與神。 \end{tabularx} \\ \\ \relax
4:22 & \begin{tabularx}{0.7\textwidth}{X} 原來經歷這神蹟醫好的人有四十多歲了。 \end{tabularx} \\ \\ \relax
4:23 & \begin{tabularx}{0.7\textwidth}{X} 二人既被釋放，就到自己的人那裡去，把祭司長和長老所說的話都告訴他們。 \end{tabularx} \\ \\ \relax
4:24 & \begin{tabularx}{0.7\textwidth}{X} 他們聽見了，就同心合意地高聲向神說：「主宰啊！你是那創造天、地、海和其中萬物的； \end{tabularx} \\ \\ \relax
4:25 & \begin{tabularx}{0.7\textwidth}{X} 你曾藉著聖靈託你僕人—我們祖宗大衛的口說：『外邦為甚麼擾動？萬民為甚麼謀算虛妄的事？ \end{tabularx} \\ \\ \relax
4:26 & \begin{tabularx}{0.7\textwidth}{X} 地上的君王都站穩，臣宰也聚集一處，要對抗主，對抗主的受膏者。』 \end{tabularx} \\ \\ \relax
4:27 & \begin{tabularx}{0.7\textwidth}{X} 希律和本丟‧彼拉多，同外邦人和以色列民，果然在這城裡聚集，要攻打你所膏的聖僕耶穌， \end{tabularx} \\ \\ \relax
4:28 & \begin{tabularx}{0.7\textwidth}{X} 做了你手和你旨意所預定必成就的事。 \end{tabularx} \\ \\ \relax
4:29 & \begin{tabularx}{0.7\textwidth}{X} 主啊，現在求你鑒察，他們的威嚇，使你僕人放膽講你的道， \end{tabularx} \\ \\ \relax
4:30 & \begin{tabularx}{0.7\textwidth}{X} 伸出你的手來，讓醫治、神蹟、奇事藉著你聖僕耶穌的名行出來。」 \end{tabularx} \\ \\ \relax
4:31 & \begin{tabularx}{0.7\textwidth}{X} 他們禱告完了，聚會的地方震動；他們都被聖靈充滿，放膽傳講神的道。 \end{tabularx} \\ \\ \relax
4:32 & \begin{tabularx}{0.7\textwidth}{X} 許多信徒都一心一意，沒有一人說他的任何東西是自己的，都是大家公用。 \end{tabularx} \\ \\ \relax
4:33 & \begin{tabularx}{0.7\textwidth}{X} 使徒以大能見證主耶穌復活；眾人也都蒙了大恩。 \end{tabularx} \\ \\ \relax
4:34 & \begin{tabularx}{0.7\textwidth}{X} 他們當中沒有一個缺乏的，因為凡有田產房屋的都賣了，把所賣的錢拿來， \end{tabularx} \\ \\ \relax
4:35 & \begin{tabularx}{0.7\textwidth}{X} 放在使徒腳前，照每人所需要的，分給每人。 \end{tabularx} \\ \\ \relax
4:36 & \begin{tabularx}{0.7\textwidth}{X} 有一個利未人，名叫約瑟，使徒稱他為巴拿巴（巴拿巴翻出來就是安慰之子），生在塞浦路斯。 \end{tabularx} \\ \\ \relax
4:37 & \begin{tabularx}{0.7\textwidth}{X} 他有田地，也賣了，把錢拿來，放在使徒腳前。 \end{tabularx} \\ \\ \relax
5:1 & \begin{tabularx}{0.7\textwidth}{X} 有一個人，名叫亞拿尼亞，同他的妻子撒非喇，賣了田產， \end{tabularx} \\ \\ \relax
5:2 & \begin{tabularx}{0.7\textwidth}{X} 把錢私自留下一部分，他的妻子也知道，其餘的部分拿來放在使徒腳前。 \end{tabularx} \\ \\ \relax
5:3 & \begin{tabularx}{0.7\textwidth}{X} 彼得說：「亞拿尼亞！為甚麼撒但充滿了你的心，使你欺騙聖靈，把賣田地的錢私自留下一部分呢？ \end{tabularx} \\ \\ \relax
5:4 & \begin{tabularx}{0.7\textwidth}{X} 田地還沒有賣，不是你自己的嗎？既賣了，錢不是你作主嗎？你怎麼心裡會想這樣做呢？你不是欺騙人，是欺騙神！」 \end{tabularx} \\ \\ \relax
5:5 & \begin{tabularx}{0.7\textwidth}{X} 亞拿尼亞一聽見這些話，就仆倒，斷了氣；所有聽見的人都非常懼怕。 \end{tabularx} \\ \\ \relax
5:6 & \begin{tabularx}{0.7\textwidth}{X} 有些年輕人起來，把他裹好，抬出去埋葬了。 \end{tabularx} \\ \\ \relax
5:7 & \begin{tabularx}{0.7\textwidth}{X} 約過了三小時，他的妻子進來，還不知道所發生的事。 \end{tabularx} \\ \\ \relax
5:8 & \begin{tabularx}{0.7\textwidth}{X} 彼得對她說：「你告訴我，你們賣田地的錢就是這些嗎？」她說：「就是這些。」 \end{tabularx} \\ \\ \relax
5:9 & \begin{tabularx}{0.7\textwidth}{X} 彼得對她說：「你們為甚麼同謀來試探主的靈呢？你看，埋葬你丈夫之人的腳已到門口，他們也要把你抬出去。」 \end{tabularx} \\ \\ \relax
5:10 & \begin{tabularx}{0.7\textwidth}{X} 她立刻仆倒在彼得腳前，斷了氣。那些年輕人進來，見她已經死了，就把她抬出去，埋在她丈夫旁邊。 \end{tabularx} \\ \\ \relax
5:11 & \begin{tabularx}{0.7\textwidth}{X} 全教會和所有聽見這些事的人都非常懼怕。 \end{tabularx} \\ \\ \relax
5:12 & \begin{tabularx}{0.7\textwidth}{X} 主藉使徒的手在民間行了許多神蹟奇事；他們都同心合意地聚集在所羅門的廊下。 \end{tabularx} \\ \\ \relax
5:13 & \begin{tabularx}{0.7\textwidth}{X} 其餘的人沒有一個敢接近他們，百姓卻尊重他們。 \end{tabularx} \\ \\ \relax
5:14 & \begin{tabularx}{0.7\textwidth}{X} 信主的人越發增添，連男帶女都很多， \end{tabularx} \\ \\ \relax
5:15 & \begin{tabularx}{0.7\textwidth}{X} 甚至有人將病人抬到街上，放在床上或褥子上，好讓彼得走過來的時候，或者影子投在一些人身上。 \end{tabularx} \\ \\ \relax
5:16 & \begin{tabularx}{0.7\textwidth}{X} 還有許多人帶著病人和被污靈纏磨的，從耶路撒冷四圍的城鎮來，他們全都得了醫治。 \end{tabularx} \\ \\ \relax
5:17 & \begin{tabularx}{0.7\textwidth}{X} 於是，大祭司採取行動，他和他所有一起的人，就是撒都該派的人，滿心忌恨， \end{tabularx} \\ \\ \relax
5:18 & \begin{tabularx}{0.7\textwidth}{X} 就下手拿住使徒，把他們押在公共拘留所內。 \end{tabularx} \\ \\ \relax
5:19 & \begin{tabularx}{0.7\textwidth}{X} 但在夜間主的使者開了監門，領他們出來，說： \end{tabularx} \\ \\ \relax
5:20 & \begin{tabularx}{0.7\textwidth}{X} 「你們去，站在聖殿裡，把這生命的一切話講給百姓聽。」 \end{tabularx} \\ \\ \relax
5:21 & \begin{tabularx}{0.7\textwidth}{X} 使徒聽了這話，天將亮的時候就進聖殿裡去教導人。大祭司和他一起的人來了，叫齊議會的人和以色列人的眾長老，然後派人到監牢裡去把使徒提出來。 \end{tabularx} \\ \\ \relax
5:22 & \begin{tabularx}{0.7\textwidth}{X} 但差役到了，不見他們在監裡，就回來稟報， \end{tabularx} \\ \\ \relax
5:23 & \begin{tabularx}{0.7\textwidth}{X} 說：「我們看見監牢關得很緊，警衛也站在門外，但打開門來，裡面一個人都不見。」 \end{tabularx} \\ \\ \relax
5:24 & \begin{tabularx}{0.7\textwidth}{X} 守殿官和祭司長聽了這些話，心裡困惑，不知這事將來如何。 \end{tabularx} \\ \\ \relax
5:25 & \begin{tabularx}{0.7\textwidth}{X} 有一個人來稟報說：「你們押在監裡的人，現在站在聖殿裡教導百姓。」 \end{tabularx} \\ \\ \relax
5:26 & \begin{tabularx}{0.7\textwidth}{X} 於是守殿官和差役去帶使徒來，並沒有用暴力，因為怕百姓用石頭打他們。 \end{tabularx} \\ \\ \relax
5:27 & \begin{tabularx}{0.7\textwidth}{X} 他們把使徒帶來了，就叫他們站在議會前。大祭司問他們， \end{tabularx} \\ \\ \relax
5:28 & \begin{tabularx}{0.7\textwidth}{X} 說：「我們不是嚴嚴地禁止你們，不可奉這名教導人嗎？看，你們倒把你們的道理充滿了耶路撒冷，想要叫這人的血歸到我們身上！」 \end{tabularx} \\ \\ \relax
5:29 & \begin{tabularx}{0.7\textwidth}{X} 彼得和眾使徒回答：「我們必須順從神，勝於順從人。 \end{tabularx} \\ \\ \relax
5:30 & \begin{tabularx}{0.7\textwidth}{X} 你們掛在木頭上殺害的耶穌，我們祖宗的神已經使他復活了。 \end{tabularx} \\ \\ \relax
5:31 & \begin{tabularx}{0.7\textwidth}{X} 神把他高舉在自己的右邊，使他作元帥，作救主，使以色列人得以悔改，並且罪得赦免。 \end{tabularx} \\ \\ \relax
5:32 & \begin{tabularx}{0.7\textwidth}{X} 我們是這些事的見證人；神賜給順從的人的聖靈也為這些事作見證。」 \end{tabularx} \\ \\ \relax
5:33 & \begin{tabularx}{0.7\textwidth}{X} 議會的人聽了極其惱怒，想要殺他們。 \end{tabularx} \\ \\ \relax
5:34 & \begin{tabularx}{0.7\textwidth}{X} 但有一個法利賽人，名叫迦瑪列，是眾百姓所敬重的律法教師，他在議會中站起來，吩咐人把使徒暫且帶到外面去， \end{tabularx} \\ \\ \relax
5:35 & \begin{tabularx}{0.7\textwidth}{X} 然後對眾人說：「以色列人哪，對於這些人，你們應當小心怎樣處理。 \end{tabularx} \\ \\ \relax
5:36 & \begin{tabularx}{0.7\textwidth}{X} 從前杜達出現，自命不凡，附從他的人數約有四百；他被殺後，附從他的人全都散了，歸於無有。 \end{tabularx} \\ \\ \relax
5:37 & \begin{tabularx}{0.7\textwidth}{X} 此後，登記戶籍的時候，又有加利利的猶大出現，引誘百姓跟從他，他也滅亡，附從他的人也都四散了。 \end{tabularx} \\ \\ \relax
5:38 & \begin{tabularx}{0.7\textwidth}{X} 現在，我勸你們不要管這些人，任憑他們吧！他們所謀所為若是出於人，必要敗壞； \end{tabularx} \\ \\ \relax
5:39 & \begin{tabularx}{0.7\textwidth}{X} 若是出於神，你們就不能敗壞他們，恐怕你們倒是攻擊神了。」議會的人被他說服了， \end{tabularx} \\ \\ \relax
5:40 & \begin{tabularx}{0.7\textwidth}{X} 就叫使徒來，把他們打了，又吩咐他們不可奉耶穌的名講道，然後把他們釋放了。 \end{tabularx} \\ \\ \relax
5:41 & \begin{tabularx}{0.7\textwidth}{X} 他們歡歡喜喜地離開議會，因他們算配為這名受辱。 \end{tabularx} \\ \\ \relax
5:42 & \begin{tabularx}{0.7\textwidth}{X} 他們就每日在聖殿裡，在家裡，不住地教導人，傳耶穌是基督的福音。 \end{tabularx} \\ \\ \relax
6:1 & \begin{tabularx}{0.7\textwidth}{X} 那些日子，門徒增多，有說希臘話的猶太人向希伯來人發怨言，因為在日常的供給上忽略了他們的寡婦。 \end{tabularx} \\ \\ \relax
6:2 & \begin{tabularx}{0.7\textwidth}{X} 十二使徒叫眾門徒來，說：「我們撇下神的道去管理飯食，是不合宜的。 \end{tabularx} \\ \\ \relax
6:3 & \begin{tabularx}{0.7\textwidth}{X} 所以弟兄們，當從你們中間選出七個有好名聲、滿有聖靈和智慧，我們派他們管理這事。 \end{tabularx} \\ \\ \relax
6:4 & \begin{tabularx}{0.7\textwidth}{X} 至於我們，我們要專注於祈禱和傳道的事奉。」 \end{tabularx} \\ \\ \relax
6:5 & \begin{tabularx}{0.7\textwidth}{X} 這話使全會眾都喜悅，就揀選了司提反—他是一個滿有信心和聖靈的人；他們又揀選了腓利、伯羅哥羅、尼迦挪、提門、巴米拿，並皈依猶太教的安提阿人尼哥拉， \end{tabularx} \\ \\ \relax
6:6 & \begin{tabularx}{0.7\textwidth}{X} 叫他們站在使徒面前，使徒禱告後，就為他們按手。 \end{tabularx} \\ \\ \relax
6:7 & \begin{tabularx}{0.7\textwidth}{X} 神的道興旺起來；在耶路撒冷門徒數目增加得很多，也有許多祭司聽從了這信仰。 \end{tabularx} \\ \\ \relax
6:8 & \begin{tabularx}{0.7\textwidth}{X} 司提反滿有恩惠和能力，在民間行了大奇事和神蹟。 \end{tabularx} \\ \\ \relax
6:9 & \begin{tabularx}{0.7\textwidth}{X} 當時有從稱為「自由人」會堂，並古利奈、亞歷山大會堂來的人，還有些從基利家、亞細亞來的人，起來和司提反辯論。 \end{tabularx} \\ \\ \relax
6:10 & \begin{tabularx}{0.7\textwidth}{X} 司提反是以智慧和聖靈說話，眾人抵擋不住， \end{tabularx} \\ \\ \relax
6:11 & \begin{tabularx}{0.7\textwidth}{X} 就收買人來說：「我們聽見他說褻瀆摩西和神的話。」 \end{tabularx} \\ \\ \relax
6:12 & \begin{tabularx}{0.7\textwidth}{X} 他們又煽動百姓、長老和文士，就突然來捉拿他，把他帶到議會去， \end{tabularx} \\ \\ \relax
6:13 & \begin{tabularx}{0.7\textwidth}{X} 設下假見證，說：「這個人不斷地說話，侮辱神聖的地方和律法。 \end{tabularx} \\ \\ \relax
6:14 & \begin{tabularx}{0.7\textwidth}{X} 我們曾聽見他說，這拿撒勒人耶穌要毀壞這地方，也要改變摩西所交給我們的規矩。」 \end{tabularx} \\ \\ \relax
6:15 & \begin{tabularx}{0.7\textwidth}{X} 在議會裡坐著的人都定睛看他，見他的面貌好像天使的面貌。 \end{tabularx} \\ \\ \relax
7:1 & \begin{tabularx}{0.7\textwidth}{X} 大祭司說：「果真有這些事嗎？」 \end{tabularx} \\ \\ \relax
7:2 & \begin{tabularx}{0.7\textwidth}{X} 司提反說：「諸位父老弟兄請聽！從前我們的祖宗亞伯拉罕在美索不達米亞，還沒有住在哈蘭的時候，榮耀的神向他顯現， \end{tabularx} \\ \\ \relax
7:3 & \begin{tabularx}{0.7\textwidth}{X} 對他說：『你要離開本地和親族，往我所要指示你的地去。』 \end{tabularx} \\ \\ \relax
7:4 & \begin{tabularx}{0.7\textwidth}{X} 他就離開迦勒底人的地方，住在哈蘭。他父親死了以後，神使他從那裡搬到你們現在所住的地方。 \end{tabularx} \\ \\ \relax
7:5 & \begin{tabularx}{0.7\textwidth}{X} 在這裡神並沒有給他產業，連立足的地方都沒有，但應許要將這地賜給他和他的後裔為業，雖然那時他還沒有兒子。 \end{tabularx} \\ \\ \relax
7:6 & \begin{tabularx}{0.7\textwidth}{X} 神這樣說：『他的後裔必寄居外邦，那裡的人要使他們作奴隸，苦待他們四百年。』 \end{tabularx} \\ \\ \relax
7:7 & \begin{tabularx}{0.7\textwidth}{X} 神又說：『但我要懲罰使他們作奴隸的那國。以後他們要出來，在這地方事奉我。』 \end{tabularx} \\ \\ \relax
7:8 & \begin{tabularx}{0.7\textwidth}{X} 神又賜他割禮的約。於是亞伯拉罕生了以撒，在第八日給他行了割禮；後來以撒生雅各，雅各生十二位先祖。 \end{tabularx} \\ \\ \relax
7:9 & \begin{tabularx}{0.7\textwidth}{X} 「先祖嫉妒約瑟，把他賣到埃及去，神卻與他同在， \end{tabularx} \\ \\ \relax
7:10 & \begin{tabularx}{0.7\textwidth}{X} 救他脫離一切苦難，又使他在埃及王法老面前蒙恩，又有智慧。法老派他作埃及國的宰相兼管法老的全家。 \end{tabularx} \\ \\ \relax
7:11 & \begin{tabularx}{0.7\textwidth}{X} 後來全埃及和迦南遭遇饑荒和大災難，我們的祖宗絕了糧。 \end{tabularx} \\ \\ \relax
7:12 & \begin{tabularx}{0.7\textwidth}{X} 雅各聽見在埃及有糧，就打發我們的祖宗初次往那裡去。 \end{tabularx} \\ \\ \relax
7:13 & \begin{tabularx}{0.7\textwidth}{X} 第二次約瑟與兄弟們相認，法老才認識他的家族。 \end{tabularx} \\ \\ \relax
7:14 & \begin{tabularx}{0.7\textwidth}{X} 約瑟就打發人，請父親雅各和全族七十五個人都來。 \end{tabularx} \\ \\ \relax
7:15 & \begin{tabularx}{0.7\textwidth}{X} 於是雅各下了埃及，後來他和我們的祖宗都死在那裡； \end{tabularx} \\ \\ \relax
7:16 & \begin{tabularx}{0.7\textwidth}{X} 他們又被遷到示劍，葬於亞伯拉罕在示劍用銀子從哈抹子孫買來的墳墓裡。 \end{tabularx} \\ \\ \relax
7:17 & \begin{tabularx}{0.7\textwidth}{X} 「當神應許亞伯拉罕的日期將到的時候，以色列人在埃及人丁興旺， \end{tabularx} \\ \\ \relax
7:18 & \begin{tabularx}{0.7\textwidth}{X} 直到另一位不認識約瑟的王興起統治埃及。 \end{tabularx} \\ \\ \relax
7:19 & \begin{tabularx}{0.7\textwidth}{X} 他用詭計待我們的宗族，苦待我們的祖宗，強迫他們丟棄嬰孩，使嬰孩不能存活。 \end{tabularx} \\ \\ \relax
7:20 & \begin{tabularx}{0.7\textwidth}{X} 就在那時，摩西生了下來，神看為俊美，在父親家裡被撫養了三個月。 \end{tabularx} \\ \\ \relax
7:21 & \begin{tabularx}{0.7\textwidth}{X} 他被丟棄的時候，法老的女兒拾了去，當自己的兒子撫養。 \end{tabularx} \\ \\ \relax
7:22 & \begin{tabularx}{0.7\textwidth}{X} 摩西學了埃及人一切的學問，說話辦事都有才能。 \end{tabularx} \\ \\ \relax
7:23 & \begin{tabularx}{0.7\textwidth}{X} 「他到了四十歲，心中起意去看望他的弟兄以色列人。 \end{tabularx} \\ \\ \relax
7:24 & \begin{tabularx}{0.7\textwidth}{X} 他見他們中的一個人受冤屈，就庇護他，為那被壓迫的人報仇，打死了那埃及人。 \end{tabularx} \\ \\ \relax
7:25 & \begin{tabularx}{0.7\textwidth}{X} 他以為他的弟兄們必明白神是藉他的手搭救他們，他們卻不明白。 \end{tabularx} \\ \\ \relax
7:26 & \begin{tabularx}{0.7\textwidth}{X} 第二天，他遇見有人在打架，就想勸他們和好，說：『二位，你們是弟兄，為甚麼彼此欺負呢？』 \end{tabularx} \\ \\ \relax
7:27 & \begin{tabularx}{0.7\textwidth}{X} 那欺負鄰舍的人把他推開，說：『誰立你作我們的領袖和審判官呢？ \end{tabularx} \\ \\ \relax
7:28 & \begin{tabularx}{0.7\textwidth}{X} 難道你要殺我像昨天殺那埃及人一樣嗎？』 \end{tabularx} \\ \\ \relax
7:29 & \begin{tabularx}{0.7\textwidth}{X} 摩西聽見這話就逃走了，寄居於米甸地，在那裡生了兩個兒子。 \end{tabularx} \\ \\ \relax
7:30 & \begin{tabularx}{0.7\textwidth}{X} 「過了四十年，在西奈山的曠野，有一位天使在荊棘的火焰中向摩西顯現。 \end{tabularx} \\ \\ \relax
7:31 & \begin{tabularx}{0.7\textwidth}{X} 摩西見了那異象，覺得很驚訝，正往前觀看的時候，有主的聲音說： \end{tabularx} \\ \\ \relax
7:32 & \begin{tabularx}{0.7\textwidth}{X} 『我是你列祖的神，就是亞伯拉罕、以撒、雅各的神。』摩西戰戰兢兢，不敢觀看。 \end{tabularx} \\ \\ \relax
7:33 & \begin{tabularx}{0.7\textwidth}{X} 主對他說：『把你腳上的鞋脫下來，因為你所站的地方是聖地。 \end{tabularx} \\ \\ \relax
7:34 & \begin{tabularx}{0.7\textwidth}{X} 我的百姓在埃及所受的困苦，我確實看見了；他們悲嘆的聲音，我也聽見了。我下來要救他們。現在，你來，我要差你往埃及去。』 \end{tabularx} \\ \\ \relax
7:35 & \begin{tabularx}{0.7\textwidth}{X} 「這摩西就是有人曾棄絕他說『誰立你作我們的領袖和審判官』的，神卻藉那在荊棘中顯現的天使的手差派他作領袖，作解救者。 \end{tabularx} \\ \\ \relax
7:36 & \begin{tabularx}{0.7\textwidth}{X} 這人領以色列人出來，在埃及地，在紅海，在曠野的四十年間行了奇事神蹟。 \end{tabularx} \\ \\ \relax
7:37 & \begin{tabularx}{0.7\textwidth}{X} 這人是摩西，就是那曾對以色列人說『神要從你們弟兄中給你們興起一位先知像我』的。 \end{tabularx} \\ \\ \relax
7:38 & \begin{tabularx}{0.7\textwidth}{X} 這人是那曾在曠野的會眾中和西奈山上，與那對他說話的天使同在，又與我們祖宗同在的，他領受了活潑的聖言傳給我們。 \end{tabularx} \\ \\ \relax
7:39 & \begin{tabularx}{0.7\textwidth}{X} 我們的祖宗不肯聽從，反棄絕他，他們的心轉向埃及， \end{tabularx} \\ \\ \relax
7:40 & \begin{tabularx}{0.7\textwidth}{X} 對亞倫說：『你為我們造神明，在我們前面引路，因為領我們出埃及地的這個摩西，我們不知道他遭遇了甚麼事。』 \end{tabularx} \\ \\ \relax
7:41 & \begin{tabularx}{0.7\textwidth}{X} 那時，他們造了一個牛犢，又拿祭物獻給那像，為自己手所做的工作歡躍。 \end{tabularx} \\ \\ \relax
7:42 & \begin{tabularx}{0.7\textwidth}{X} 但是神轉臉不顧，任憑他們祭拜天上的日月星辰，正如先知書上所寫的：『以色列家啊，你們四十年間在曠野，何曾將犧牲和祭物獻給我？ \end{tabularx} \\ \\ \relax
7:43 & \begin{tabularx}{0.7\textwidth}{X} 你們抬著摩洛的帳幕和理番---你們神明的星，就是你們所造為要敬拜的像。因此，我要把你們遷到巴比倫外去。』 \end{tabularx} \\ \\ \relax
7:44 & \begin{tabularx}{0.7\textwidth}{X} 「我們的祖宗在曠野，有作證的會幕，是神吩咐摩西照著他所看見的樣式做的。 \end{tabularx} \\ \\ \relax
7:45 & \begin{tabularx}{0.7\textwidth}{X} 這帳幕，我們的祖宗同約書亞相繼承受了，當神在他們面前趕走外邦人的時候，他們把這帳幕搬進承受為業之地，直存到大衛的日子。 \end{tabularx} \\ \\ \relax
7:46 & \begin{tabularx}{0.7\textwidth}{X} 大衛在神面前蒙恩，祈求為雅各的家預備居所。 \end{tabularx} \\ \\ \relax
7:47 & \begin{tabularx}{0.7\textwidth}{X} 但卻是所羅門為神造成殿宇。 \end{tabularx} \\ \\ \relax
7:48 & \begin{tabularx}{0.7\textwidth}{X} 其實，至高者並不住人手所造的，就如先知所言： \end{tabularx} \\ \\ \relax
7:49 & \begin{tabularx}{0.7\textwidth}{X} 『主說：天是我的寶座，地是我的腳凳。你們要為我造怎樣的殿宇？哪裡是我安歇的地方呢？ \end{tabularx} \\ \\ \relax
7:50 & \begin{tabularx}{0.7\textwidth}{X} 這一切不都是我手所造的嗎？』 \end{tabularx} \\ \\ \relax
7:51 & \begin{tabularx}{0.7\textwidth}{X} 「你們這硬著頸項，心與耳未受割禮的人哪，時常抗拒聖靈！你們的祖宗怎樣，你們也怎樣。 \end{tabularx} \\ \\ \relax
7:52 & \begin{tabularx}{0.7\textwidth}{X} 先知中有哪一個不是受你們祖宗的迫害呢？他們把預先宣告那義者要來的人殺了。如今你們成了那義者的出賣者和兇手了。 \end{tabularx} \\ \\ \relax
7:53 & \begin{tabularx}{0.7\textwidth}{X} 你們領受了天使所傳佈的律法，竟不遵守。」 \end{tabularx} \\ \\ \relax
7:54 & \begin{tabularx}{0.7\textwidth}{X} 眾人聽見這些話，心中極其惱怒，向司提反咬牙切齒。 \end{tabularx} \\ \\ \relax
7:55 & \begin{tabularx}{0.7\textwidth}{X} 但司提反滿有聖靈，定睛望天，看見神的榮耀，又看見耶穌站在神的右邊， \end{tabularx} \\ \\ \relax
7:56 & \begin{tabularx}{0.7\textwidth}{X} 就說：「我看見天開了，人子站在神的右邊。」 \end{tabularx} \\ \\ \relax
7:57 & \begin{tabularx}{0.7\textwidth}{X} 眾人大聲喊叫，摀著耳朵，齊心衝向他， \end{tabularx} \\ \\ \relax
7:58 & \begin{tabularx}{0.7\textwidth}{X} 把他推到城外，用石頭打他。作見證的人把他們的衣裳放在一個名叫掃羅的青年腳前。 \end{tabularx} \\ \\ \relax
7:59 & \begin{tabularx}{0.7\textwidth}{X} 他們正用石頭打司提反的時候，他呼求說：「主耶穌啊，求你接納我的靈魂！」 \end{tabularx} \\ \\ \relax
7:60 & \begin{tabularx}{0.7\textwidth}{X} 然後他跪下來，大聲喊著：「主啊，不要將這罪歸於他們！」說了這話，就長眠了。 \end{tabularx} \\ \\ \relax
8:1 & \begin{tabularx}{0.7\textwidth}{X} 掃羅也贊同處死他。從那一天開始，耶路撒冷的教會遭受到大迫害，除了使徒以外，眾門徒都分散在猶太和撒瑪利亞各處。 \end{tabularx} \\ \\ \relax
8:2 & \begin{tabularx}{0.7\textwidth}{X} 有些虔誠的人把司提反埋葬了，為他大大哀哭。 \end{tabularx} \\ \\ \relax
8:3 & \begin{tabularx}{0.7\textwidth}{X} 掃羅卻殘害教會，挨家挨戶地進去，拉著男女關在監裡。 \end{tabularx} \\ \\ \relax
8:4 & \begin{tabularx}{0.7\textwidth}{X} 那些分散的人往各地去傳福音的道。 \end{tabularx} \\ \\ \relax
8:5 & \begin{tabularx}{0.7\textwidth}{X} 腓利下撒瑪利亞城去，向當地人宣講基督。 \end{tabularx} \\ \\ \relax
8:6 & \begin{tabularx}{0.7\textwidth}{X} 眾人都聚精會神，同心合意地聽腓利所說的話，一邊聽他的話，一邊看他所行的神蹟。 \end{tabularx} \\ \\ \relax
8:7 & \begin{tabularx}{0.7\textwidth}{X} 因為有許多人被污靈附著，那些污靈大聲呼叫，從他們身上出來；還有許多癱瘓的、瘸腿的都得了醫治。 \end{tabularx} \\ \\ \relax
8:8 & \begin{tabularx}{0.7\textwidth}{X} 那城裡，有極大的喜樂。 \end{tabularx} \\ \\ \relax
8:9 & \begin{tabularx}{0.7\textwidth}{X} 有一個人名叫西門，向來在那城裡行邪術，自命為大人物，使撒瑪利亞的居民驚奇。 \end{tabularx} \\ \\ \relax
8:10 & \begin{tabularx}{0.7\textwidth}{X} 所有的人，從小到大都聽從他，說：「這個人就是神的能力，那稱為大能者的。」 \end{tabularx} \\ \\ \relax
8:11 & \begin{tabularx}{0.7\textwidth}{X} 他們聽從他，因他很久以來用邪術使他們驚奇。 \end{tabularx} \\ \\ \relax
8:12 & \begin{tabularx}{0.7\textwidth}{X} 當他們信了腓利所傳神國的福音和耶穌基督的名，連男帶女都受了洗。 \end{tabularx} \\ \\ \relax
8:13 & \begin{tabularx}{0.7\textwidth}{X} 西門自己也信了；既受了洗，就常與腓利在一處，看見他所行的神蹟和大異能，就覺得很驚奇。 \end{tabularx} \\ \\ \relax
8:14 & \begin{tabularx}{0.7\textwidth}{X} 在耶路撒冷的使徒聽見撒瑪利亞人領受了神的道，就打發彼得和約翰到他們那裡去。 \end{tabularx} \\ \\ \relax
8:15 & \begin{tabularx}{0.7\textwidth}{X} 兩個人下去，就為他們禱告，要讓他們領受聖靈， \end{tabularx} \\ \\ \relax
8:16 & \begin{tabularx}{0.7\textwidth}{X} 因為聖靈還沒有降在他們任何一個人身上，他們只奉主耶穌的名受了洗。 \end{tabularx} \\ \\ \relax
8:17 & \begin{tabularx}{0.7\textwidth}{X} 於是使徒按手在他們頭上，他們就領受了聖靈。 \end{tabularx} \\ \\ \relax
8:18 & \begin{tabularx}{0.7\textwidth}{X} 西門看見使徒一按手，就有聖靈賜下，就拿錢給使徒， \end{tabularx} \\ \\ \relax
8:19 & \begin{tabularx}{0.7\textwidth}{X} 說：「請把這權柄也給我，使我手按著誰，誰就可以領受聖靈。」 \end{tabularx} \\ \\ \relax
8:20 & \begin{tabularx}{0.7\textwidth}{X} 彼得對他說：「你的銀子和你一同滅亡吧！因為你想神的恩賜是可以用錢買的。 \end{tabularx} \\ \\ \relax
8:21 & \begin{tabularx}{0.7\textwidth}{X} 你在這道上無份無關；因為你在神面前心懷不正。 \end{tabularx} \\ \\ \relax
8:22 & \begin{tabularx}{0.7\textwidth}{X} 你要為你這樣的惡而悔改，祈求主，或者你心裡的意念可得赦免。 \end{tabularx} \\ \\ \relax
8:23 & \begin{tabularx}{0.7\textwidth}{X} 我看出你正在苦膽之中，被不義捆綁著。」 \end{tabularx} \\ \\ \relax
8:24 & \begin{tabularx}{0.7\textwidth}{X} 西門回答說：「請你們為我求主，使你們所說的，沒有一樣臨到我身上。」 \end{tabularx} \\ \\ \relax
8:25 & \begin{tabularx}{0.7\textwidth}{X} 使徒既作了見證，並且宣講了主的道，就回耶路撒冷去，一路在撒瑪利亞好些村莊傳揚福音。 \end{tabularx} \\ \\ \relax
8:26 & \begin{tabularx}{0.7\textwidth}{X} 有主的一個使者對腓利說：「起來！向南走，往那從耶路撒冷下迦薩的路上去。」那路是曠野。 \end{tabularx} \\ \\ \relax
8:27 & \begin{tabularx}{0.7\textwidth}{X} 腓利就起身去了。不料，有一個埃塞俄比亞人，是個有大權的太監，在埃塞俄比亞女王甘大基的手下總管銀庫，他上耶路撒冷去禮拜。 \end{tabularx} \\ \\ \relax
8:28 & \begin{tabularx}{0.7\textwidth}{X} 回程中，他坐在車上，正念著以賽亞先知的書， \end{tabularx} \\ \\ \relax
8:29 & \begin{tabularx}{0.7\textwidth}{X} 聖靈對腓利說：「你去！靠近那車走。」 \end{tabularx} \\ \\ \relax
8:30 & \begin{tabularx}{0.7\textwidth}{X} 腓利就跑到太監那裡，聽見他正在念以賽亞先知的書，就說：「你明白你所念的嗎？」 \end{tabularx} \\ \\ \relax
8:31 & \begin{tabularx}{0.7\textwidth}{X} 他說：「沒有人指教我，怎能明白呢？」於是他請腓利上車，與他同坐。 \end{tabularx} \\ \\ \relax
8:32 & \begin{tabularx}{0.7\textwidth}{X} 他所念的那段經文是這樣：「他像羊被牽去宰殺，又像羔羊在剪毛的人手下無聲，他也是這樣不開口。 \end{tabularx} \\ \\ \relax
8:33 & \begin{tabularx}{0.7\textwidth}{X} 他卑微的時候，得不到公義的審判，誰能述說他的身世？因為他的生命從地上被奪去。」 \end{tabularx} \\ \\ \relax
8:34 & \begin{tabularx}{0.7\textwidth}{X} 太監回答腓利說：「請問，先知說這話是指誰，是指自己，還是指別人呢？」 \end{tabularx} \\ \\ \relax
8:35 & \begin{tabularx}{0.7\textwidth}{X} 腓利就開口，從這段經文開始，對他傳講耶穌的福音。 \end{tabularx} \\ \\ \relax
8:36 & \begin{tabularx}{0.7\textwidth}{X} 二人正沿路往前走，到了有水的地方，太監說：「看哪！這裡有水，有甚麼能阻止我受洗呢？」 \end{tabularx} \\ \\ \relax
8:37 & \begin{tabularx}{0.7\textwidth}{X} 腓利說：『你若是一心相信，就可以。』他回答：『我信耶穌基督是神的兒子。』 \end{tabularx} \\ \\ \relax
8:38 & \begin{tabularx}{0.7\textwidth}{X} 於是他吩咐把車停下來，腓利和太監二人一同下到水裡，腓利就給他施洗。 \end{tabularx} \\ \\ \relax
8:39 & \begin{tabularx}{0.7\textwidth}{X} 他們從水裡上來，主的靈把腓利提了去，太監再也看不見他了，就歡歡喜喜地上路。 \end{tabularx} \\ \\ \relax
8:40 & \begin{tabularx}{0.7\textwidth}{X} 後來有人在亞鎖都遇見腓利；他走遍那地方，在各城宣揚福音，一直到凱撒利亞。 \end{tabularx} \\ \\ \relax
9:1 & \begin{tabularx}{0.7\textwidth}{X} 掃羅不斷用威嚇兇悍的口氣向主的門徒說話。他去見大祭司， \end{tabularx} \\ \\ \relax
9:2 & \begin{tabularx}{0.7\textwidth}{X} 要求發信給大馬士革的各會堂，若是找著信奉這道的人，無論男女，都准他捆綁帶到耶路撒冷。 \end{tabularx} \\ \\ \relax
9:3 & \begin{tabularx}{0.7\textwidth}{X} 掃羅在途中，將到大馬士革的時候，忽然有一道光從天上下來，四面照射著他， \end{tabularx} \\ \\ \relax
9:4 & \begin{tabularx}{0.7\textwidth}{X} 他就仆倒在地，聽見有聲音對他說：「掃羅！掃羅！你為甚麼迫害我？」 \end{tabularx} \\ \\ \relax
9:5 & \begin{tabularx}{0.7\textwidth}{X} 他說：「主啊！你是誰？」主說：「我就是你所迫害的耶穌。 \end{tabularx} \\ \\ \relax
9:6 & \begin{tabularx}{0.7\textwidth}{X} 起來！進城去，你應該做的事，必有人告訴你。」 \end{tabularx} \\ \\ \relax
9:7 & \begin{tabularx}{0.7\textwidth}{X} 同行的人站在那裡，說不出話來，因為他們聽見聲音，卻看不見人。 \end{tabularx} \\ \\ \relax
9:8 & \begin{tabularx}{0.7\textwidth}{X} 掃羅從地上起來，睜開眼睛，竟不能看見甚麼。有人拉他的手，領他進了大馬士革。 \end{tabularx} \\ \\ \relax
9:9 & \begin{tabularx}{0.7\textwidth}{X} 他三天甚麼都看不見，也不吃也不喝。 \end{tabularx} \\ \\ \relax
9:10 & \begin{tabularx}{0.7\textwidth}{X} 那時，在大馬士革有一個門徒，名叫亞拿尼亞。主在異象中對他說：「亞拿尼亞！」他說：「主啊，我在這裡。」 \end{tabularx} \\ \\ \relax
9:11 & \begin{tabularx}{0.7\textwidth}{X} 主對他說：「起來！往那叫直街的路去，在猶大的家裡，去找一個大數人，名叫掃羅；他正在禱告， \end{tabularx} \\ \\ \relax
9:12 & \begin{tabularx}{0.7\textwidth}{X} 在異象中看見了一個人，名叫亞拿尼亞，進來為他按手，讓他能再看得見。」 \end{tabularx} \\ \\ \relax
9:13 & \begin{tabularx}{0.7\textwidth}{X} 亞拿尼亞回答：「主啊，我聽見許多人講到這個人，說他怎樣在耶路撒冷多多苦待你的聖徒， \end{tabularx} \\ \\ \relax
9:14 & \begin{tabularx}{0.7\textwidth}{X} 並且他在這裡有從祭司長得來的權柄，要捆綁一切求告你名的人。」 \end{tabularx} \\ \\ \relax
9:15 & \begin{tabularx}{0.7\textwidth}{X} 主對他說：「你只管去。他是我所揀選的器皿，要在外邦人、君王和以色列人面前宣揚我的名。 \end{tabularx} \\ \\ \relax
9:16 & \begin{tabularx}{0.7\textwidth}{X} 我也要指示他，為我的名必須受許多的苦難。」 \end{tabularx} \\ \\ \relax
9:17 & \begin{tabularx}{0.7\textwidth}{X} 亞拿尼亞就去了，進入那家，把手按在掃羅身上，說：「掃羅弟兄，在你來的路上向你顯現的主，就是耶穌，打發我來，叫你能再看得見，又被聖靈充滿。」 \end{tabularx} \\ \\ \relax
9:18 & \begin{tabularx}{0.7\textwidth}{X} 掃羅的眼睛上立刻好像有鱗一般的東西掉下來，他就能再看得見，於是他起來，受了洗， \end{tabularx} \\ \\ \relax
9:19 & \begin{tabularx}{0.7\textwidth}{X} 吃過飯體力就恢復了。 \end{tabularx} \\ \\ \relax
9:20 & \begin{tabularx}{0.7\textwidth}{X} 立刻在各會堂裡傳揚耶穌，說他是神的兒子。 \end{tabularx} \\ \\ \relax
9:21 & \begin{tabularx}{0.7\textwidth}{X} 凡聽見的人都很驚奇，說：「在耶路撒冷殘害求告這名的不就是這個人嗎？他不是到這裡來要捆綁他們，帶到祭司長那裡去嗎？」 \end{tabularx} \\ \\ \relax
9:22 & \begin{tabularx}{0.7\textwidth}{X} 但掃羅越發有能力，駁倒住在大馬士革的猶太人，證明耶穌是基督。 \end{tabularx} \\ \\ \relax
9:23 & \begin{tabularx}{0.7\textwidth}{X} 過了好些日子，猶太人商議要殺掃羅， \end{tabularx} \\ \\ \relax
9:24 & \begin{tabularx}{0.7\textwidth}{X} 但他們的計謀被掃羅知道了。他們晝夜在城門守候著要殺他。 \end{tabularx} \\ \\ \relax
9:25 & \begin{tabularx}{0.7\textwidth}{X} 他的門徒就在夜間用筐子把他從城牆上縋了下去。 \end{tabularx} \\ \\ \relax
9:26 & \begin{tabularx}{0.7\textwidth}{X} 掃羅到了耶路撒冷，想與門徒結交，大家卻都怕他，不信他是門徒。 \end{tabularx} \\ \\ \relax
9:27 & \begin{tabularx}{0.7\textwidth}{X} 只有巴拿巴接待他，領他去見使徒，把他在路上怎麼看見主，主怎麼向他說話，他在大馬士革怎麼奉耶穌的名放膽傳道，都述說出來。 \end{tabularx} \\ \\ \relax
9:28 & \begin{tabularx}{0.7\textwidth}{X} 於是掃羅在耶路撒冷同門徒出入來往，奉主的名放膽傳道， \end{tabularx} \\ \\ \relax
9:29 & \begin{tabularx}{0.7\textwidth}{X} 並和說希臘話的猶太人講論辯駁，他們卻想法子要殺他。 \end{tabularx} \\ \\ \relax
9:30 & \begin{tabularx}{0.7\textwidth}{X} 弟兄們知道了，就帶他下凱撒利亞，送他往大數去。 \end{tabularx} \\ \\ \relax
9:31 & \begin{tabularx}{0.7\textwidth}{X} 那時，猶太、加利利、撒瑪利亞各處的教會都得平安，建立起來，凡事敬畏主，蒙聖靈的安慰，人數逐漸增多。 \end{tabularx} \\ \\ \relax
9:32 & \begin{tabularx}{0.7\textwidth}{X} 彼得在眾信徒中到處奔波的時候，也到了住在呂大的聖徒那裡。 \end{tabularx} \\ \\ \relax
9:33 & \begin{tabularx}{0.7\textwidth}{X} 他在那裡遇見一個人，名叫以尼雅，得了癱瘓，在褥子上躺了八年。 \end{tabularx} \\ \\ \relax
9:34 & \begin{tabularx}{0.7\textwidth}{X} 彼得對他說：「以尼雅，耶穌基督醫好你了，起來！整理你的褥子吧。」他立刻就起來了。 \end{tabularx} \\ \\ \relax
9:35 & \begin{tabularx}{0.7\textwidth}{X} 凡住呂大和沙崙的人都看見了他，就歸向主。 \end{tabularx} \\ \\ \relax
9:36 & \begin{tabularx}{0.7\textwidth}{X} 在約帕有一個女門徒，名叫大比大，翻出來的意思是多加；她廣行善事，多施賙濟。 \end{tabularx} \\ \\ \relax
9:37 & \begin{tabularx}{0.7\textwidth}{X} 當時，她患病死了，有人把她清洗後，停在樓上。 \end{tabularx} \\ \\ \relax
9:38 & \begin{tabularx}{0.7\textwidth}{X} 呂大原與約帕相近；門徒聽見彼得在那裡，就派兩個人去見他，央求他說：「請快到我們那裡去，不要耽延。」 \end{tabularx} \\ \\ \relax
9:39 & \begin{tabularx}{0.7\textwidth}{X} 彼得就起身和他們同去。他到了，就有人領他上樓。眾寡婦都站在彼得旁邊哭，拿多加與她們同在時所做的內衣外衣給他看。 \end{tabularx} \\ \\ \relax
9:40 & \begin{tabularx}{0.7\textwidth}{X} 彼得叫她們都出去，然後跪下禱告，轉身對著屍體說：「大比大，起來！」她就睜開眼睛，看見彼得，就坐了起來。 \end{tabularx} \\ \\ \relax
9:41 & \begin{tabularx}{0.7\textwidth}{X} 彼得伸手扶她起來，叫那些聖徒和寡婦都進來，把多加活活地交給他們。 \end{tabularx} \\ \\ \relax
9:42 & \begin{tabularx}{0.7\textwidth}{X} 這事傳遍了約帕，就有許多人信了主。 \end{tabularx} \\ \\ \relax
9:43 & \begin{tabularx}{0.7\textwidth}{X} 此後，彼得在約帕一個皮革匠西門的家裡住了好些日子。 \end{tabularx} \\ \\ \relax
10:1 & \begin{tabularx}{0.7\textwidth}{X} 在凱撒利亞有一個人名叫哥尼流，是意大利營的百夫長。 \end{tabularx} \\ \\ \relax
10:2 & \begin{tabularx}{0.7\textwidth}{X} 他是個虔誠人，他和全家都敬畏神。他多多賙濟百姓，常常向神禱告。 \end{tabularx} \\ \\ \relax
10:3 & \begin{tabularx}{0.7\textwidth}{X} 有一天，約在下午三點鐘，他在異象中清楚看見神的一個使者進來，到他那裡，對他說：「哥尼流。」 \end{tabularx} \\ \\ \relax
10:4 & \begin{tabularx}{0.7\textwidth}{X} 哥尼流定睛看他，驚惶地說：「主啊，甚麼事？」天使對他說：「你的禱告和你的賙濟已達到神面前，蒙記念了。 \end{tabularx} \\ \\ \relax
10:5 & \begin{tabularx}{0.7\textwidth}{X} 現在你要派人往約帕去，請一位稱為彼得的西門來。 \end{tabularx} \\ \\ \relax
10:6 & \begin{tabularx}{0.7\textwidth}{X} 他住在一個皮革匠西門的家裡，房子就在海邊。」 \end{tabularx} \\ \\ \relax
10:7 & \begin{tabularx}{0.7\textwidth}{X} 向他說話的天使離開後，哥尼流叫了兩個僕人和常伺候他的一個虔誠的兵來， \end{tabularx} \\ \\ \relax
10:8 & \begin{tabularx}{0.7\textwidth}{X} 把一切的事都講給他們聽，然後就派他們往約帕去。 \end{tabularx} \\ \\ \relax
10:9 & \begin{tabularx}{0.7\textwidth}{X} 第二天，他們走路將近那城，約在正午，彼得上房頂去禱告。 \end{tabularx} \\ \\ \relax
10:10 & \begin{tabularx}{0.7\textwidth}{X} 他覺得餓了，想要吃。那家的人正預備飯的時候，彼得魂遊象外， \end{tabularx} \\ \\ \relax
10:11 & \begin{tabularx}{0.7\textwidth}{X} 看見天開了，有一塊好像大布的東西降下，四角吊著縋在地上， \end{tabularx} \\ \\ \relax
10:12 & \begin{tabularx}{0.7\textwidth}{X} 裡面有地上各樣四腳的走獸、爬蟲和天上的飛鳥。 \end{tabularx} \\ \\ \relax
10:13 & \begin{tabularx}{0.7\textwidth}{X} 又有聲音對他說：「彼得，起來！宰了吃。」 \end{tabularx} \\ \\ \relax
10:14 & \begin{tabularx}{0.7\textwidth}{X} 彼得卻說：「主啊，絕對不可！凡污俗和不潔淨的東西，我從來沒有吃過。」 \end{tabularx} \\ \\ \relax
10:15 & \begin{tabularx}{0.7\textwidth}{X} 第二次有聲音再對他說：「神所潔淨的，你不可當作污俗的。」 \end{tabularx} \\ \\ \relax
10:16 & \begin{tabularx}{0.7\textwidth}{X} 這樣一連三次，那東西隨即收回天上去了。 \end{tabularx} \\ \\ \relax
10:17 & \begin{tabularx}{0.7\textwidth}{X} 正當彼得心裡困惑，不知所看見的異象是甚麼意思時，哥尼流所差來的人已經找到了西門的家，站在門外， \end{tabularx} \\ \\ \relax
10:18 & \begin{tabularx}{0.7\textwidth}{X} 喊著問有沒有一位稱為彼得的西門住在這裡。 \end{tabularx} \\ \\ \relax
10:19 & \begin{tabularx}{0.7\textwidth}{X} 彼得還在思考那異象的時候，聖靈對他說：「有三個人來找你。 \end{tabularx} \\ \\ \relax
10:20 & \begin{tabularx}{0.7\textwidth}{X} 起來，下去，跟他們同去，不要疑惑，因為是我差他們來的。」 \end{tabularx} \\ \\ \relax
10:21 & \begin{tabularx}{0.7\textwidth}{X} 於是彼得下去見那些人，說：「我就是你們要找的人，你們是為了甚麼緣故在這裡？」 \end{tabularx} \\ \\ \relax
10:22 & \begin{tabularx}{0.7\textwidth}{X} 他們說：「百夫長哥尼流是個義人，敬畏神，為猶太全民族所稱讚。他蒙一位聖天使指示，叫他請你到他家裡去，要聽你講話。」 \end{tabularx} \\ \\ \relax
10:23 & \begin{tabularx}{0.7\textwidth}{X} 彼得就請他們進去住宿。次日，他起身和他們同去，還有約帕的幾個弟兄跟他一起去。 \end{tabularx} \\ \\ \relax
10:24 & \begin{tabularx}{0.7\textwidth}{X} 又次日，他進入凱撒利亞，哥尼流已經請了他的親朋好友在等候他們。 \end{tabularx} \\ \\ \relax
10:25 & \begin{tabularx}{0.7\textwidth}{X} 彼得一進去，哥尼流就迎接他，俯伏在他腳前拜他。 \end{tabularx} \\ \\ \relax
10:26 & \begin{tabularx}{0.7\textwidth}{X} 但是彼得拉他起來，說：「你起來，我自己也不過是人。」 \end{tabularx} \\ \\ \relax
10:27 & \begin{tabularx}{0.7\textwidth}{X} 彼得和他一邊說話一邊進去，見有好些人聚集， \end{tabularx} \\ \\ \relax
10:28 & \begin{tabularx}{0.7\textwidth}{X} 就對他們說：「你們知道，猶太人和別國的人結交來往本是不合規矩的，但神已經指示我，無論甚麼人都不可看作污俗或不潔淨的。 \end{tabularx} \\ \\ \relax
10:29 & \begin{tabularx}{0.7\textwidth}{X} 所以，我一被邀請，沒有推辭就來了。現在請問，你們為甚麼叫我來呢？」 \end{tabularx} \\ \\ \relax
10:30 & \begin{tabularx}{0.7\textwidth}{X} 哥尼流說：「四天前，這個時候，我在家中守著下午三點鐘的禱告，忽然有一個人穿著明亮的衣裳站在我面前， \end{tabularx} \\ \\ \relax
10:31 & \begin{tabularx}{0.7\textwidth}{X} 說：『哥尼流，你的禱告已蒙垂聽，你的賙濟在神面前已蒙記念了。 \end{tabularx} \\ \\ \relax
10:32 & \begin{tabularx}{0.7\textwidth}{X} 你要派人往約帕去，請那稱為彼得的西門來，他住在海邊一個皮革匠西門的家裡。』 \end{tabularx} \\ \\ \relax
10:33 & \begin{tabularx}{0.7\textwidth}{X} 所以我立刻派人去請你。你來了真好。現在我們都在神面前，要聽主吩咐你的一切話。」 \end{tabularx} \\ \\ \relax
10:34 & \begin{tabularx}{0.7\textwidth}{X} 彼得開口說：「我真的看出神是不偏待人的。 \end{tabularx} \\ \\ \relax
10:35 & \begin{tabularx}{0.7\textwidth}{X} 不但如此，在各國中那敬畏他而行義的人都為他所悅納。 \end{tabularx} \\ \\ \relax
10:36 & \begin{tabularx}{0.7\textwidth}{X} 神藉著耶穌基督—他是萬有的主—傳和平的福音，把這道傳給以色列人。 \end{tabularx} \\ \\ \relax
10:37 & \begin{tabularx}{0.7\textwidth}{X} 這話在約翰傳揚洗禮以後，從加利利起，傳遍了猶太。 \end{tabularx} \\ \\ \relax
10:38 & \begin{tabularx}{0.7\textwidth}{X} 神怎樣以聖靈和能力膏了拿撒勒人耶穌，這都是你們知道的。他到處奔波，行善事，醫好凡被魔鬼壓制的人，因為神與他同在。 \end{tabularx} \\ \\ \relax
10:39 & \begin{tabularx}{0.7\textwidth}{X} 他在猶太人之地和耶路撒冷所行的一切事，有我們作見證人。他們竟把他掛在木頭上殺了。 \end{tabularx} \\ \\ \relax
10:40 & \begin{tabularx}{0.7\textwidth}{X} 第三天，神使他復活，使他顯現出來； \end{tabularx} \\ \\ \relax
10:41 & \begin{tabularx}{0.7\textwidth}{X} 不是顯現給所有的人看，而是顯現給神預先所揀選為他作見證的人看，就是我們這些在他從死人中復活以後和他同吃同喝的人。 \end{tabularx} \\ \\ \relax
10:42 & \begin{tabularx}{0.7\textwidth}{X} 他吩咐我們傳道給眾人，證明他是神所立定，要作審判活人、死人的審判者。 \end{tabularx} \\ \\ \relax
10:43 & \begin{tabularx}{0.7\textwidth}{X} 眾先知也為這人作見證：凡信他的人，必藉著他的名得蒙赦罪。」 \end{tabularx} \\ \\ \relax
10:44 & \begin{tabularx}{0.7\textwidth}{X} 彼得還在說這些話的時候，聖靈降在一切聽道的人身上。 \end{tabularx} \\ \\ \relax
10:45 & \begin{tabularx}{0.7\textwidth}{X} 那些奉割禮的信徒和彼得同來，見聖靈的恩賜也澆在外邦人身上，就都驚奇； \end{tabularx} \\ \\ \relax
10:46 & \begin{tabularx}{0.7\textwidth}{X} 因聽見他們說方言，稱讚神為大。於是彼得回答： \end{tabularx} \\ \\ \relax
10:47 & \begin{tabularx}{0.7\textwidth}{X} 「這些人既受了聖靈，跟我們一樣，誰能阻止用水給他們施洗呢？」 \end{tabularx} \\ \\ \relax
10:48 & \begin{tabularx}{0.7\textwidth}{X} 他就吩咐奉耶穌基督的名給他們施洗。於是他們請彼得住了幾天。 \end{tabularx} \\ \\ \relax
11:1 & \begin{tabularx}{0.7\textwidth}{X} 使徒和在猶太的眾弟兄聽到外邦人也領受了神的道。 \end{tabularx} \\ \\ \relax
11:2 & \begin{tabularx}{0.7\textwidth}{X} 等到彼得上了耶路撒冷，那些奉割禮的信徒和他爭辯， \end{tabularx} \\ \\ \relax
11:3 & \begin{tabularx}{0.7\textwidth}{X} 說：「你竟進入未受割禮之人當中，和他們一同吃飯！」 \end{tabularx} \\ \\ \relax
11:4 & \begin{tabularx}{0.7\textwidth}{X} 彼得就開始把這事逐一向他們解釋，說： \end{tabularx} \\ \\ \relax
11:5 & \begin{tabularx}{0.7\textwidth}{X} 「我在約帕城裡禱告的時候，魂遊象外，看見異象，有一塊好像大布的東西降下，四角吊著從天縋下，直來到我跟前。 \end{tabularx} \\ \\ \relax
11:6 & \begin{tabularx}{0.7\textwidth}{X} 我定睛觀看，見內中有地上四腳的牲畜、野獸、爬蟲和天上的飛鳥。 \end{tabularx} \\ \\ \relax
11:7 & \begin{tabularx}{0.7\textwidth}{X} 我還聽見有聲音對我說：『彼得，起來！宰了吃。』 \end{tabularx} \\ \\ \relax
11:8 & \begin{tabularx}{0.7\textwidth}{X} 我說：『主啊，絕對不可！凡污俗或不潔淨的東西從來沒有進過我的口。』 \end{tabularx} \\ \\ \relax
11:9 & \begin{tabularx}{0.7\textwidth}{X} 第二次，有聲音從天上回答：『神所潔淨的，你不可當作污俗的。』 \end{tabularx} \\ \\ \relax
11:10 & \begin{tabularx}{0.7\textwidth}{X} 這樣一連三次，然後一切就都收回天上去了。 \end{tabularx} \\ \\ \relax
11:11 & \begin{tabularx}{0.7\textwidth}{X} 正當那時，有三個從凱撒利亞差來見我的人，站在我們所住的屋子門前。 \end{tabularx} \\ \\ \relax
11:12 & \begin{tabularx}{0.7\textwidth}{X} 聖靈吩咐我和他們同去，不要疑惑，還有這六位弟兄也跟我一起去，我們進了那人的家。 \end{tabularx} \\ \\ \relax
11:13 & \begin{tabularx}{0.7\textwidth}{X} 那人就告訴我們，他如何看見一位天使站在他家裡，說：『你派人往約帕去，請那稱為彼得的西門來， \end{tabularx} \\ \\ \relax
11:14 & \begin{tabularx}{0.7\textwidth}{X} 他有話要告訴你，因這些話你和你的全家都可以得救。』 \end{tabularx} \\ \\ \relax
11:15 & \begin{tabularx}{0.7\textwidth}{X} 我一開始講話，聖靈就降在他們身上，正像當初降在我們身上一樣。 \end{tabularx} \\ \\ \relax
11:16 & \begin{tabularx}{0.7\textwidth}{X} 我就想起主的話如何說：『約翰用水施洗，但你們要在聖靈裡受洗。』 \end{tabularx} \\ \\ \relax
11:17 & \begin{tabularx}{0.7\textwidth}{X} 既然神給他們恩賜，像在我們信主耶穌基督的時候給了我們一樣，我是誰，能攔阻神嗎？」 \end{tabularx} \\ \\ \relax
11:18 & \begin{tabularx}{0.7\textwidth}{X} 眾人聽見這些話，就不說話了，只歸榮耀給神，說：「這樣看來，神也賜恩給外邦人，使他們悔改得生命了。」 \end{tabularx} \\ \\ \relax
11:19 & \begin{tabularx}{0.7\textwidth}{X} 那些因司提反的事遭患難而四處分散的門徒，直走到腓尼基、塞浦路斯和安提阿。他們不向別人講道，只向猶太人講。 \end{tabularx} \\ \\ \relax
11:20 & \begin{tabularx}{0.7\textwidth}{X} 但內中有塞浦路斯和古利奈人，他們到了安提阿也向希臘人傳講主耶穌的福音。 \end{tabularx} \\ \\ \relax
11:21 & \begin{tabularx}{0.7\textwidth}{X} 主的手與他們同在，信而歸主的人數很多。 \end{tabularx} \\ \\ \relax
11:22 & \begin{tabularx}{0.7\textwidth}{X} 這風聲傳到耶路撒冷教會的人耳中，他們就打發巴拿巴到安提阿去。 \end{tabularx} \\ \\ \relax
11:23 & \begin{tabularx}{0.7\textwidth}{X} 他到了那裡，看見神所賜的恩就歡喜，勸勉眾人要立定心志，恆久靠主。 \end{tabularx} \\ \\ \relax
11:24 & \begin{tabularx}{0.7\textwidth}{X} 這巴拿巴原是個好人，滿有聖靈和信心，於是有許多人歸服了主。 \end{tabularx} \\ \\ \relax
11:25 & \begin{tabularx}{0.7\textwidth}{X} 他又往大數去找掃羅， \end{tabularx} \\ \\ \relax
11:26 & \begin{tabularx}{0.7\textwidth}{X} 找著了，就帶他到安提阿去。他們足有一年和教會一同聚集，教導了許多人。門徒稱為「基督徒」是從安提阿開始的。 \end{tabularx} \\ \\ \relax
11:27 & \begin{tabularx}{0.7\textwidth}{X} 當那些日子，有幾位先知從耶路撒冷下到安提阿。 \end{tabularx} \\ \\ \relax
11:28 & \begin{tabularx}{0.7\textwidth}{X} 內中有一位，名叫亞迦布，站起來，藉著聖靈指示普天下將有大饑荒；這事在克勞第年間果然實現了。 \end{tabularx} \\ \\ \relax
11:29 & \begin{tabularx}{0.7\textwidth}{X} 於是門徒決定，照各人的力量捐錢，送去供給住在猶太的弟兄。 \end{tabularx} \\ \\ \relax
11:30 & \begin{tabularx}{0.7\textwidth}{X} 他們就這樣做了，託巴拿巴和掃羅的手送到眾長老那裡。 \end{tabularx} \\ \\ \relax
12:1 & \begin{tabularx}{0.7\textwidth}{X} 約在那時候，希律王下手苦待教會中的一些人， \end{tabularx} \\ \\ \relax
12:2 & \begin{tabularx}{0.7\textwidth}{X} 用刀殺了約翰的哥哥雅各。 \end{tabularx} \\ \\ \relax
12:3 & \begin{tabularx}{0.7\textwidth}{X} 他見猶太人喜歡這事，也去拿住彼得。那時候正是除酵節期間。 \end{tabularx} \\ \\ \relax
12:4 & \begin{tabularx}{0.7\textwidth}{X} 希律捉了彼得，押在監裡，交給四班士兵看守，每班四個人，企圖要在逾越節後把他提出來，當著百姓辦他。 \end{tabularx} \\ \\ \relax
12:5 & \begin{tabularx}{0.7\textwidth}{X} 於是彼得被囚在監裡，教會卻為他切切禱告神。 \end{tabularx} \\ \\ \relax
12:6 & \begin{tabularx}{0.7\textwidth}{X} 希律將要提他出來的前一夜，彼得被兩條鐵鏈鎖著，睡在兩個士兵當中；門前還有警衛看守。 \end{tabularx} \\ \\ \relax
12:7 & \begin{tabularx}{0.7\textwidth}{X} 忽然，有主的一個使者顯現，牢房裡有光照耀；天使拍彼得的肋旁，叫醒了他，說：「快起來！」鐵鏈就從他手上脫落下來。 \end{tabularx} \\ \\ \relax
12:8 & \begin{tabularx}{0.7\textwidth}{X} 天使對他說：「束上腰帶，穿上鞋子。」他就照著做了。天使又對他說：「披上外衣，跟我來。」 \end{tabularx} \\ \\ \relax
12:9 & \begin{tabularx}{0.7\textwidth}{X} 彼得就出來跟著他走，不知道天使所做是真的，以為見了異象。 \end{tabularx} \\ \\ \relax
12:10 & \begin{tabularx}{0.7\textwidth}{X} 他們經過了第一層和第二層監牢，就來到往城內的鐵門，那門就自動給他們開了。他們出來，走過一條街，忽然天使離開他去了。 \end{tabularx} \\ \\ \relax
12:11 & \begin{tabularx}{0.7\textwidth}{X} 彼得清醒過來，說：「現在我真知道主差遣他的使者，救我脫離希律的手，和猶太人所期待的一切。」 \end{tabularx} \\ \\ \relax
12:12 & \begin{tabularx}{0.7\textwidth}{X} 他明白了，就到那稱為馬可的約翰的母親馬利亞家去，在那裡已有好些人聚集禱告。 \end{tabularx} \\ \\ \relax
12:13 & \begin{tabularx}{0.7\textwidth}{X} 彼得敲外門時，有一個使女，名叫羅大，出來應門， \end{tabularx} \\ \\ \relax
12:14 & \begin{tabularx}{0.7\textwidth}{X} 認出是彼得的聲音，歡喜得顧不了開門，就跑進去報信，說彼得站在門外。 \end{tabularx} \\ \\ \relax
12:15 & \begin{tabularx}{0.7\textwidth}{X} 他們對她說：「你瘋了！」使女堅持真有其事。他們說：「那是他的天使。」 \end{tabularx} \\ \\ \relax
12:16 & \begin{tabularx}{0.7\textwidth}{X} 彼得不停地敲門；他們開了門，一見是他，就很驚奇。 \end{tabularx} \\ \\ \relax
12:17 & \begin{tabularx}{0.7\textwidth}{X} 彼得做個手勢，要他們不作聲，就告訴他們主怎樣領他出監；又說：「你們要把這些事告訴雅各和眾弟兄。」然後，他離開往別處去了。 \end{tabularx} \\ \\ \relax
12:18 & \begin{tabularx}{0.7\textwidth}{X} 到了天亮，士兵中起了不少騷動，不知道彼得到哪裡去了。 \end{tabularx} \\ \\ \relax
12:19 & \begin{tabularx}{0.7\textwidth}{X} 希律找他，找不著，就審問警衛，下令帶走他們處死。後來希律離開猶太，下凱撒利亞去，住在那裡。 \end{tabularx} \\ \\ \relax
12:20 & \begin{tabularx}{0.7\textwidth}{X} 希律向推羅和西頓的人發怒。他們那一帶地方是從王的土地供應糧食的，因此就託了王的內侍大臣伯拉斯都的情，一心來求和。 \end{tabularx} \\ \\ \relax
12:21 & \begin{tabularx}{0.7\textwidth}{X} 希律在所定的日子，穿上朝服，坐在位上，對他們演講。 \end{tabularx} \\ \\ \relax
12:22 & \begin{tabularx}{0.7\textwidth}{X} 民眾一直喊著：「這是神明的聲音，不是人的聲音。」 \end{tabularx} \\ \\ \relax
12:23 & \begin{tabularx}{0.7\textwidth}{X} 希律不歸榮耀給神，所以主的使者立刻擊打他，他被蟲咬，就斷了氣。 \end{tabularx} \\ \\ \relax
12:24 & \begin{tabularx}{0.7\textwidth}{X} 神的道日見興旺，越發廣傳。 \end{tabularx} \\ \\ \relax
12:25 & \begin{tabularx}{0.7\textwidth}{X} 巴拿巴和掃羅完成了供給的事，就回到耶路撒冷，帶著稱為馬可的約翰同去。 \end{tabularx} \\ \\ \relax
13:1 & \begin{tabularx}{0.7\textwidth}{X} 在安提阿的教會中，有幾位先知和教師，就是巴拿巴和稱為尼結的西面、古利奈人路求，與希律分封王一起長大的馬念，和掃羅。 \end{tabularx} \\ \\ \relax
13:2 & \begin{tabularx}{0.7\textwidth}{X} 他們在事奉主和禁食的時候，聖靈說：「要為我分派巴拿巴和掃羅去做我召他們做的工作。」 \end{tabularx} \\ \\ \relax
13:3 & \begin{tabularx}{0.7\textwidth}{X} 於是他們禁食禱告後，給巴拿巴和掃羅按手，然後派遣他們走了。 \end{tabularx} \\ \\ \relax
13:4 & \begin{tabularx}{0.7\textwidth}{X} 他們既蒙聖靈差遣，就下到西流基，從那裡坐船往塞浦路斯去， \end{tabularx} \\ \\ \relax
13:5 & \begin{tabularx}{0.7\textwidth}{X} 到了撒拉米，就在猶太人各會堂裡宣講神的道，也有約翰作他們的幫手。 \end{tabularx} \\ \\ \relax
13:6 & \begin{tabularx}{0.7\textwidth}{X} 他們走遍全島，直到帕弗，在那裡遇見一個術士—猶太人的假先知，名叫巴耶穌。 \end{tabularx} \\ \\ \relax
13:7 & \begin{tabularx}{0.7\textwidth}{X} 這人常和士求‧保羅省長在一起。士求‧保羅是個通達人，他請巴拿巴和掃羅來，要聽神的道。 \end{tabularx} \\ \\ \relax
13:8 & \begin{tabularx}{0.7\textwidth}{X} 只是術士以呂馬(他的名字翻出來就是行法術的意思)敵對使徒，設法使省長遠離這信仰。 \end{tabularx} \\ \\ \relax
13:9 & \begin{tabularx}{0.7\textwidth}{X} 掃羅，又名保羅，被聖靈充滿，定睛看他， \end{tabularx} \\ \\ \relax
13:10 & \begin{tabularx}{0.7\textwidth}{X} 說：「你這充滿各樣詭詐奸惡，魔鬼的兒子，一切正義的仇敵，你還不停止扭曲主的正道嗎？ \end{tabularx} \\ \\ \relax
13:11 & \begin{tabularx}{0.7\textwidth}{X} 現在你看，主的手臨到你身上，你會瞎眼，暫時看不見日光。」立刻迷濛和黑暗籠罩著他，他到處摸索，求人拉著手領他。 \end{tabularx} \\ \\ \relax
13:12 & \begin{tabularx}{0.7\textwidth}{X} 省長看見所發生的事就信了，因對主的教導感到驚奇。 \end{tabularx} \\ \\ \relax
13:13 & \begin{tabularx}{0.7\textwidth}{X} 保羅和他的同伴從帕弗開船，來到旁非利亞的別加，約翰卻離開他們，回耶路撒冷去了。 \end{tabularx} \\ \\ \relax
13:14 & \begin{tabularx}{0.7\textwidth}{X} 他們從別加往前行，來到彼西底的安提阿。在安息日，他們進了會堂就坐下。 \end{tabularx} \\ \\ \relax
13:15 & \begin{tabularx}{0.7\textwidth}{X} 在讀完了律法和先知的書，會堂主管們叫人過去，對他們說：「二位弟兄，你們若有甚麼勸勉眾人的話，請說。」 \end{tabularx} \\ \\ \relax
13:16 & \begin{tabularx}{0.7\textwidth}{X} 保羅就站起來，做個手勢，說：「諸位以色列人和一切敬畏神的人，請聽。 \end{tabularx} \\ \\ \relax
13:17 & \begin{tabularx}{0.7\textwidth}{X} 這以色列民的神揀選了我們的祖宗，當百姓寄居埃及的時候抬舉他們，用大能的手領他們從那地出來。 \end{tabularx} \\ \\ \relax
13:18 & \begin{tabularx}{0.7\textwidth}{X} 他在曠野容忍他們，約有四十年。 \end{tabularx} \\ \\ \relax
13:19 & \begin{tabularx}{0.7\textwidth}{X} 他消滅了迦南地七族的人後，把那地分給他們為業， \end{tabularx} \\ \\ \relax
13:20 & \begin{tabularx}{0.7\textwidth}{X} 約有四百五十年。此後，他給他們設立士師，直到撒母耳先知的時候。 \end{tabularx} \\ \\ \relax
13:21 & \begin{tabularx}{0.7\textwidth}{X} 從那時起，他們要求立一個王，神就將便雅憫支派中基士的兒子掃羅給他們作王，共四十年。 \end{tabularx} \\ \\ \relax
13:22 & \begin{tabularx}{0.7\textwidth}{X} 他廢了掃羅之後，就興起大衛作他們的王，又為他作見證說：『我尋得耶西的兒子大衛，他是合我心意的人，他要遵行我一切的旨意。』 \end{tabularx} \\ \\ \relax
13:23 & \begin{tabularx}{0.7\textwidth}{X} 從這人的後裔中，神已經照著所應許的為以色列人興起一位救主，就是耶穌。 \end{tabularx} \\ \\ \relax
13:24 & \begin{tabularx}{0.7\textwidth}{X} 在他沒有出來以前，約翰已向以色列全民宣講悔改的洗禮。 \end{tabularx} \\ \\ \relax
13:25 & \begin{tabularx}{0.7\textwidth}{X} 約翰快走完他的人生路程時，說：『你們以為我是誰？我不是；但是有一位在我以後來的，我就是解他腳上的鞋帶也不配。』 \end{tabularx} \\ \\ \relax
13:26 & \begin{tabularx}{0.7\textwidth}{X} 「諸位弟兄—亞伯拉罕的子孫和你們中間敬畏神的人哪，這救世的道是傳給我們的。 \end{tabularx} \\ \\ \relax
13:27 & \begin{tabularx}{0.7\textwidth}{X} 耶路撒冷的居民和他們的官長，因為不認識這基督，也不明白每安息日所讀的先知的書，把他定了死罪，正應驗了先知的預言。 \end{tabularx} \\ \\ \relax
13:28 & \begin{tabularx}{0.7\textwidth}{X} 雖然他們查不出他有該死的罪狀，還是要求彼拉多把他殺了。 \end{tabularx} \\ \\ \relax
13:29 & \begin{tabularx}{0.7\textwidth}{X} 他們既實現了經上指著他所記的一切話，就從木頭上把他取下來，放在墳墓裡。 \end{tabularx} \\ \\ \relax
13:30 & \begin{tabularx}{0.7\textwidth}{X} 神卻使他從死人中復活。 \end{tabularx} \\ \\ \relax
13:31 & \begin{tabularx}{0.7\textwidth}{X} 有許多日子，他向那些從加利利同他上耶路撒冷的人顯現，這些人如今在民間成為他的見證人。 \end{tabularx} \\ \\ \relax
13:32 & \begin{tabularx}{0.7\textwidth}{X} 我們報好信息給你們，就是那應許祖宗的話， \end{tabularx} \\ \\ \relax
13:33 & \begin{tabularx}{0.7\textwidth}{X} 神已經向我們這些作他們兒女的應驗，使耶穌復活了。正如《詩篇》第二篇上記著：『你是我的兒子，我今日生了你。』 \end{tabularx} \\ \\ \relax
13:34 & \begin{tabularx}{0.7\textwidth}{X} 論到神使他從死人中復活，不再歸於朽壞，他曾這樣說：『我必將所應許大衛那聖潔、可靠的恩典賜給你們。』 \end{tabularx} \\ \\ \relax
13:35 & \begin{tabularx}{0.7\textwidth}{X} 所以他也在另一篇說：『你必不讓你的聖者見朽壞。』 \end{tabularx} \\ \\ \relax
13:36 & \begin{tabularx}{0.7\textwidth}{X} 大衛在世的時候，遵行了神的旨意就長眠了，歸到他祖宗那裡，已見朽壞； \end{tabularx} \\ \\ \relax
13:37 & \begin{tabularx}{0.7\textwidth}{X} 惟獨神使他復活的那一位，他並未見朽壞。 \end{tabularx} \\ \\ \relax
13:38 & \begin{tabularx}{0.7\textwidth}{X} 所以弟兄們，你們當知道：赦罪的道是由這人傳給你們的， \end{tabularx} \\ \\ \relax
13:39 & \begin{tabularx}{0.7\textwidth}{X} 你們靠摩西的律法在不得稱義的一切事上，每一個信靠這位耶穌的都得稱義了。 \end{tabularx} \\ \\ \relax
13:40 & \begin{tabularx}{0.7\textwidth}{X} 所以，你們要小心，免得先知書上所說的臨到你們： \end{tabularx} \\ \\ \relax
13:41 & \begin{tabularx}{0.7\textwidth}{X} 『要觀看，你們這些藐視的人，要驚訝，要滅亡，因為在你們的日子，我行一件事，雖有人告訴你們，你們總是不信。』」 \end{tabularx} \\ \\ \relax
13:42 & \begin{tabularx}{0.7\textwidth}{X} 他們走出會堂的時候，眾人請他們在下一個安息日再講這些話給他們聽。 \end{tabularx} \\ \\ \relax
13:43 & \begin{tabularx}{0.7\textwidth}{X} 散會以後，有許多猶太人和敬虔的皈依猶太教的人跟從了保羅和巴拿巴。二人對他們講話，勸他們務要恆久倚靠神的恩典。 \end{tabularx} \\ \\ \relax
13:44 & \begin{tabularx}{0.7\textwidth}{X} 到下一個安息日，全城的人幾乎都聚集起來，要聽主的道。 \end{tabularx} \\ \\ \relax
13:45 & \begin{tabularx}{0.7\textwidth}{X} 但猶太人看見這麼多的人，就滿心嫉妒，辯駁保羅所說的話，並且毀謗他。 \end{tabularx} \\ \\ \relax
13:46 & \begin{tabularx}{0.7\textwidth}{X} 於是保羅和巴拿巴放膽說：「神的道本應先傳給你們；只因你們棄絕這道，斷定自己不配得永生，我們就轉向外邦人。 \end{tabularx} \\ \\ \relax
13:47 & \begin{tabularx}{0.7\textwidth}{X} 因為主曾這樣吩咐我們：『我已經立你作萬邦之光，使你施行我的救恩，直到地極。』」 \end{tabularx} \\ \\ \relax
13:48 & \begin{tabularx}{0.7\textwidth}{X} 外邦人聽見這話很歡喜，讚美主的道，凡被指定得永生的人都信了。 \end{tabularx} \\ \\ \relax
13:49 & \begin{tabularx}{0.7\textwidth}{X} 於是主的道傳遍了那一帶地方。 \end{tabularx} \\ \\ \relax
13:50 & \begin{tabularx}{0.7\textwidth}{X} 但猶太人挑唆虔敬尊貴的婦女和城內有名望的人，迫害保羅和巴拿巴，把他們趕出境外。 \end{tabularx} \\ \\ \relax
13:51 & \begin{tabularx}{0.7\textwidth}{X} 二人對著眾人跺掉腳上的塵土，然後往以哥念去了。 \end{tabularx} \\ \\ \relax
13:52 & \begin{tabularx}{0.7\textwidth}{X} 門徒滿心喜樂，又被聖靈充滿。 \end{tabularx} \\ \\ \relax
14:1 & \begin{tabularx}{0.7\textwidth}{X} 同樣的事也發生在以哥念。保羅和巴拿巴進了猶太人的會堂，在那裡講道，所以有很多猶太人和希臘人都信了。 \end{tabularx} \\ \\ \relax
14:2 & \begin{tabularx}{0.7\textwidth}{X} 但那不順從的猶太人煽動外邦人，使他們心裡仇恨弟兄。 \end{tabularx} \\ \\ \relax
14:3 & \begin{tabularx}{0.7\textwidth}{X} 二人在那裡住了好些日子，倚靠主放膽講道，主藉他們的手施行神蹟奇事，證明他恩惠的道。 \end{tabularx} \\ \\ \relax
14:4 & \begin{tabularx}{0.7\textwidth}{X} 城裡的眾人卻分裂了：有依附猶太人的，有依附使徒的。 \end{tabularx} \\ \\ \relax
14:5 & \begin{tabularx}{0.7\textwidth}{X} 那時，外邦人、猶太人和他們的官長，一齊擁上來，要凌辱使徒，用石頭打他們。 \end{tabularx} \\ \\ \relax
14:6 & \begin{tabularx}{0.7\textwidth}{X} 使徒知道了，就逃到呂高尼的路司得和特庇兩個城，以及周圍地方去， \end{tabularx} \\ \\ \relax
14:7 & \begin{tabularx}{0.7\textwidth}{X} 在那裡繼續傳福音。 \end{tabularx} \\ \\ \relax
14:8 & \begin{tabularx}{0.7\textwidth}{X} 路司得城裡有一個兩腳無力的人，他從母腹裡就是瘸腿的，老是坐著，從來沒有走過。 \end{tabularx} \\ \\ \relax
14:9 & \begin{tabularx}{0.7\textwidth}{X} 他聽保羅講道；保羅定睛看他，見他有信心，可得痊癒， \end{tabularx} \\ \\ \relax
14:10 & \begin{tabularx}{0.7\textwidth}{X} 就大聲說：「起來！兩腳站直。」那人就跳起來，開始行走。 \end{tabularx} \\ \\ \relax
14:11 & \begin{tabularx}{0.7\textwidth}{X} 眾人看見保羅所做的事，就用呂高尼話大聲說：「有神明藉著人形降臨在我們中間了。」 \end{tabularx} \\ \\ \relax
14:12 & \begin{tabularx}{0.7\textwidth}{X} 於是他們稱巴拿巴為宙斯，稱保羅為希耳米，因為他總是帶頭說話。 \end{tabularx} \\ \\ \relax
14:13 & \begin{tabularx}{0.7\textwidth}{X} 城外有宙斯廟的祭司牽著牛，拿著花環，來到門前，要同眾人一起獻祭。 \end{tabularx} \\ \\ \relax
14:14 & \begin{tabularx}{0.7\textwidth}{X} 巴拿巴和保羅二位使徒聽見，就撕開衣裳，跳進眾人中間，喊著： \end{tabularx} \\ \\ \relax
14:15 & \begin{tabularx}{0.7\textwidth}{X} 「諸位，為甚麼做這些事呢？我們也是人，性情和你們一樣。我們傳福音給你們，是要你們離棄這些虛妄的事，歸向那創造天、地、海和其中萬物的永生的神。 \end{tabularx} \\ \\ \relax
14:16 & \begin{tabularx}{0.7\textwidth}{X} 他在從前的世代，任憑萬國各行其道； \end{tabularx} \\ \\ \relax
14:17 & \begin{tabularx}{0.7\textwidth}{X} 然而他未嘗不為自己留下證據來，就如常行善事，從天降雨，賞賜豐年，使你們飲食飽足，滿心喜樂。」 \end{tabularx} \\ \\ \relax
14:18 & \begin{tabularx}{0.7\textwidth}{X} 二人說了這些話，總算攔住眾人不獻祭給他們。 \end{tabularx} \\ \\ \relax
14:19 & \begin{tabularx}{0.7\textwidth}{X} 但有些猶太人，從安提阿和以哥念來，挑唆眾人，並且用石頭打保羅，以為他死了，就把他拖到城外。 \end{tabularx} \\ \\ \relax
14:20 & \begin{tabularx}{0.7\textwidth}{X} 當門徒圍著他的時候，他站了起來，走進城去。第二天，保羅同巴拿巴往特庇去。 \end{tabularx} \\ \\ \relax
14:21 & \begin{tabularx}{0.7\textwidth}{X} 保羅和巴拿巴對那城裡的人傳了福音，使好些人成為門徒後，又回路司得、以哥念、安提阿去， \end{tabularx} \\ \\ \relax
14:22 & \begin{tabularx}{0.7\textwidth}{X} 堅固門徒的心，勸他們持守他們的信仰，說：「我們進入神的國，必須經歷許多艱難。」 \end{tabularx} \\ \\ \relax
14:23 & \begin{tabularx}{0.7\textwidth}{X} 二人在各教會中選立了長老，禁食禱告後，把他們交託給他們所信的主。 \end{tabularx} \\ \\ \relax
14:24 & \begin{tabularx}{0.7\textwidth}{X} 二人經過彼西底來到旁非利亞， \end{tabularx} \\ \\ \relax
14:25 & \begin{tabularx}{0.7\textwidth}{X} 在別加講了道，就下亞大利去， \end{tabularx} \\ \\ \relax
14:26 & \begin{tabularx}{0.7\textwidth}{X} 從那裡坐船回安提阿去。當初，眾人就在這地方，把他們交託在神的恩典中，要完成現在所做的工。 \end{tabularx} \\ \\ \relax
14:27 & \begin{tabularx}{0.7\textwidth}{X} 他們一到那裡，就聚集了會眾，述說神藉他們所行的一切事，並且神怎樣為外邦人開了信道的門。 \end{tabularx} \\ \\ \relax
14:28 & \begin{tabularx}{0.7\textwidth}{X} 二人在那裡同門徒住了一段日子。 \end{tabularx} \\ \\ \relax
15:1 & \begin{tabularx}{0.7\textwidth}{X} 有幾個人從猶太下來，教導弟兄們說：「你們若不按照摩西的規矩受割禮，不能得救。」 \end{tabularx} \\ \\ \relax
15:2 & \begin{tabularx}{0.7\textwidth}{X} 保羅和巴拿巴跟他們發生了激烈的爭執和辯論；大家就決定指派保羅、巴拿巴和本會的幾個人，為所辯論的事上耶路撒冷去見使徒和長老。 \end{tabularx} \\ \\ \relax
15:3 & \begin{tabularx}{0.7\textwidth}{X} 於是教會為他們送行。他們經過腓尼基、撒瑪利亞，沿途敘說外邦人歸主的事，使眾弟兄都非常歡喜。 \end{tabularx} \\ \\ \relax
15:4 & \begin{tabularx}{0.7\textwidth}{X} 他們到了耶路撒冷，教會、使徒和長老都接待他們，他們就述說神同他們所做的一切事。 \end{tabularx} \\ \\ \relax
15:5 & \begin{tabularx}{0.7\textwidth}{X} 惟有幾個法利賽派的信徒起來，說：「必須給外邦人行割禮，吩咐他們遵守摩西的律法。」 \end{tabularx} \\ \\ \relax
15:6 & \begin{tabularx}{0.7\textwidth}{X} 使徒和長老聚集商議這事。 \end{tabularx} \\ \\ \relax
15:7 & \begin{tabularx}{0.7\textwidth}{X} 辯論了許久後，彼得站起來，對他們說：「諸位弟兄，你們知道神早已在你們中間揀選了我，讓外邦人從我口中得聽福音之道，而且相信。 \end{tabularx} \\ \\ \relax
15:8 & \begin{tabularx}{0.7\textwidth}{X} 知道人心的神也為他們作了見證，賜聖靈給他們，正如給我們一樣； \end{tabularx} \\ \\ \relax
15:9 & \begin{tabularx}{0.7\textwidth}{X} 又藉著信潔淨了他們的心，他們和我們之間並沒有甚麼分別。 \end{tabularx} \\ \\ \relax
15:10 & \begin{tabularx}{0.7\textwidth}{X} 現在你們為甚麼試探神，要把我們祖宗和我們所不能負的軛放在門徒的頸項上呢？ \end{tabularx} \\ \\ \relax
15:11 & \begin{tabularx}{0.7\textwidth}{X} 相反地，我們相信，我們得救是因主耶穌的恩典，和他們一樣。」 \end{tabularx} \\ \\ \relax
15:12 & \begin{tabularx}{0.7\textwidth}{X} 眾人都默默無聲，聽巴拿巴和保羅述說神藉著他們在外邦人中所行的神蹟和奇事。 \end{tabularx} \\ \\ \relax
15:13 & \begin{tabularx}{0.7\textwidth}{X} 他們講完了，雅各回答說：「諸位弟兄，請聽我說。 \end{tabularx} \\ \\ \relax
15:14 & \begin{tabularx}{0.7\textwidth}{X} 剛才西門述說神當初怎樣眷顧外邦人，從他們中間選取人民歸於自己的名下； \end{tabularx} \\ \\ \relax
15:15 & \begin{tabularx}{0.7\textwidth}{X} 眾先知的話也與這意思相符合。 \end{tabularx} \\ \\ \relax
15:16 & \begin{tabularx}{0.7\textwidth}{X} 正如經上所寫的：『此後，我要回來，重新修造大衛倒塌了的帳幕，從廢墟中重新修造，把它建立起來， \end{tabularx} \\ \\ \relax
15:17 & \begin{tabularx}{0.7\textwidth}{X} 使剩餘的人，就是凡稱我名的外邦人，都尋求主。 \end{tabularx} \\ \\ \relax
15:18 & \begin{tabularx}{0.7\textwidth}{X} 這話是自古以來顯明這些事的主說的。』 \end{tabularx} \\ \\ \relax
15:19 & \begin{tabularx}{0.7\textwidth}{X} 所以，我的意見是不可難為那歸向神的外邦人； \end{tabularx} \\ \\ \relax
15:20 & \begin{tabularx}{0.7\textwidth}{X} 但是要寫信吩咐他們禁戒偶像所玷污的東西、血和勒死的牲畜，禁戒淫亂。 \end{tabularx} \\ \\ \relax
15:21 & \begin{tabularx}{0.7\textwidth}{X} 因為歷代以來，摩西的書在各城都有人宣講，每逢安息日，也在會堂裡誦讀。」 \end{tabularx} \\ \\ \relax
15:22 & \begin{tabularx}{0.7\textwidth}{X} 那時，使徒、長老和全教會認為應從他們中間揀選人，差他們和保羅、巴拿巴一同到安提阿去，所揀選的就是稱為巴撒巴的猶大和西拉。這二人在弟兄中是領袖。 \end{tabularx} \\ \\ \relax
15:23 & \begin{tabularx}{0.7\textwidth}{X} 他們帶去的信說：「使徒和作長老的弟兄們向安提阿、敘利亞、基利家外邦眾弟兄問安。 \end{tabularx} \\ \\ \relax
15:24 & \begin{tabularx}{0.7\textwidth}{X} 我們聽說，有幾個人從我們這裡出去，用一些話騷擾你們，使你們的心困惑，其實我們並沒有吩咐他們。 \end{tabularx} \\ \\ \relax
15:25 & \begin{tabularx}{0.7\textwidth}{X} 我們認為，既然我們同心定意，就揀選幾個人，派他們同我們所親愛的巴拿巴和保羅到你們那裡去。 \end{tabularx} \\ \\ \relax
15:26 & \begin{tabularx}{0.7\textwidth}{X} 這二人曾為我主耶穌基督的名不顧自己的性命。 \end{tabularx} \\ \\ \relax
15:27 & \begin{tabularx}{0.7\textwidth}{X} 所以我們派猶大和西拉去，他們也會親口述說這些事。 \end{tabularx} \\ \\ \relax
15:28 & \begin{tabularx}{0.7\textwidth}{X} 因為聖靈和我們決定除了這幾件重要的事，不將別的重擔放在你們身上， \end{tabularx} \\ \\ \relax
15:29 & \begin{tabularx}{0.7\textwidth}{X} 就是禁戒偶像所玷污的東西、血和勒死的牲畜，禁戒淫亂。這幾件你們若能自己禁戒就好了。祝你們安康！」 \end{tabularx} \\ \\ \relax
15:30 & \begin{tabularx}{0.7\textwidth}{X} 他們既奉了差遣就下安提阿去，聚集會眾，把書信交給他們。 \end{tabularx} \\ \\ \relax
15:31 & \begin{tabularx}{0.7\textwidth}{X} 眾人念了，因為信上鼓勵的話而感到欣慰。 \end{tabularx} \\ \\ \relax
15:32 & \begin{tabularx}{0.7\textwidth}{X} 猶大和西拉自己也是先知，就用許多話勸勉弟兄，堅固他們。 \end{tabularx} \\ \\ \relax
15:33 & \begin{tabularx}{0.7\textwidth}{X} 二人住了些日子，弟兄們打發他們平平安安地回到差遣他們的人那裡去。 \end{tabularx} \\ \\ \relax
15:34 & \begin{tabularx}{0.7\textwidth}{X} 惟有西拉決定仍住在那裡。 \end{tabularx} \\ \\ \relax
15:35 & \begin{tabularx}{0.7\textwidth}{X} 但保羅和巴拿巴仍留在安提阿，和許多別的人一同教導，並傳揚主的道。 \end{tabularx} \\ \\ \relax
15:36 & \begin{tabularx}{0.7\textwidth}{X} 過了些日子，保羅對巴拿巴說：「讓我們回到從前宣揚主道的各城，看看弟兄們的情況如何。」 \end{tabularx} \\ \\ \relax
15:37 & \begin{tabularx}{0.7\textwidth}{X} 巴拿巴有意要帶稱為馬可的約翰同去； \end{tabularx} \\ \\ \relax
15:38 & \begin{tabularx}{0.7\textwidth}{X} 但保羅認為不宜帶他去，因為馬可從前在旁非利亞離開他們，不和他們一起工作。 \end{tabularx} \\ \\ \relax
15:39 & \begin{tabularx}{0.7\textwidth}{X} 於是二人起了爭執，甚至彼此分手。巴拿巴帶著馬可，坐船往塞浦路斯去； \end{tabularx} \\ \\ \relax
15:40 & \begin{tabularx}{0.7\textwidth}{X} 保羅則揀選了西拉，也出發了，蒙弟兄們把他交於主的恩典中。 \end{tabularx} \\ \\ \relax
15:41 & \begin{tabularx}{0.7\textwidth}{X} 他就走遍了敘利亞、基利家，堅固眾教會。 \end{tabularx} \\ \\ \relax
16:1 & \begin{tabularx}{0.7\textwidth}{X} 後來，保羅來到特庇，又到路司得。在那裡有一個門徒，名叫提摩太，是信主的猶太婦人的兒子，他父親卻是希臘人。 \end{tabularx} \\ \\ \relax
16:2 & \begin{tabularx}{0.7\textwidth}{X} 路司得和以哥念的弟兄都稱讚他。 \end{tabularx} \\ \\ \relax
16:3 & \begin{tabularx}{0.7\textwidth}{X} 保羅要帶他同去，只因那些地方的猶太人都知道他父親是希臘人，就給他行了割禮。 \end{tabularx} \\ \\ \relax
16:4 & \begin{tabularx}{0.7\textwidth}{X} 他們經過各城，把耶路撒冷使徒和長老所決定的規條交給門徒遵守。 \end{tabularx} \\ \\ \relax
16:5 & \begin{tabularx}{0.7\textwidth}{X} 於是眾教會信心越發堅固，人數天天增加。 \end{tabularx} \\ \\ \relax
16:6 & \begin{tabularx}{0.7\textwidth}{X} 因為聖靈禁止他們在亞細亞講道，他們就經過弗呂家、加拉太一帶地方。 \end{tabularx} \\ \\ \relax
16:7 & \begin{tabularx}{0.7\textwidth}{X} 到了每西亞的邊界，他們想要往庇推尼去，耶穌的靈卻不許。 \end{tabularx} \\ \\ \relax
16:8 & \begin{tabularx}{0.7\textwidth}{X} 他們就越過每西亞，下特羅亞去。 \end{tabularx} \\ \\ \relax
16:9 & \begin{tabularx}{0.7\textwidth}{X} 夜間，有異象向保羅顯現。有一個馬其頓人站著求他說：「請你過來，到馬其頓來幫助我們！」 \end{tabularx} \\ \\ \relax
16:10 & \begin{tabularx}{0.7\textwidth}{X} 保羅既看見這異象，我們就立即設法往馬其頓去，認為神呼召我們傳福音給那裡的人。 \end{tabularx} \\ \\ \relax
16:11 & \begin{tabularx}{0.7\textwidth}{X} 我們從特羅亞開船，直行駛到撒摩特喇，第二天到了尼亞坡里； \end{tabularx} \\ \\ \relax
16:12 & \begin{tabularx}{0.7\textwidth}{X} 從那裡來到腓立比，就是馬其頓這一帶的一個重要城市，也是羅馬的駐防城。我們在這城裡住了幾天。 \end{tabularx} \\ \\ \relax
16:13 & \begin{tabularx}{0.7\textwidth}{X} 在安息日，我們出城門，到了河邊，知道那裡有一個禱告的地方，我們就坐下來對那些聚會的婦女講道。 \end{tabularx} \\ \\ \relax
16:14 & \begin{tabularx}{0.7\textwidth}{X} 有一個賣紫色布的婦人，名叫呂底亞，是推雅推喇城的人，素來敬拜神。她在聽著，主就開導她的心，使她留心聽保羅所講的話。 \end{tabularx} \\ \\ \relax
16:15 & \begin{tabularx}{0.7\textwidth}{X} 她和她一家都領了洗，就求我們說：「你們若以為我是真心信主的，請到我家裡來住。」於是她堅決請我們留下。 \end{tabularx} \\ \\ \relax
16:16 & \begin{tabularx}{0.7\textwidth}{X} 後來，我們往那禱告的地方去時，有一個被占卜的靈附身的使女迎面走來，她使用法術使她的主人們發了大財。 \end{tabularx} \\ \\ \relax
16:17 & \begin{tabularx}{0.7\textwidth}{X} 她跟隨保羅和我們，喊著說：「這些人是至高神的僕人，對你們傳講救人的道路。」 \end{tabularx} \\ \\ \relax
16:18 & \begin{tabularx}{0.7\textwidth}{X} 她一連好幾天這樣喊叫，保羅就心中厭煩，轉身對那靈說：「我奉耶穌基督的名吩咐你從她身上出來！」那靈立刻出來了。 \end{tabularx} \\ \\ \relax
16:19 & \begin{tabularx}{0.7\textwidth}{X} 使女的主人們見發財的指望沒有了，就揪住保羅和西拉，拉他們到市上去見官； \end{tabularx} \\ \\ \relax
16:20 & \begin{tabularx}{0.7\textwidth}{X} 又帶他們到行政官長們面前，說：「這些騷擾我們城的，他們是猶太人， \end{tabularx} \\ \\ \relax
16:21 & \begin{tabularx}{0.7\textwidth}{X} 竟傳佈我們羅馬人所不可接受、不可遵守的規矩。」 \end{tabularx} \\ \\ \relax
16:22 & \begin{tabularx}{0.7\textwidth}{X} 群眾就一齊起來攻擊他們。官長們吩咐撕開他們的衣裳，用棍子打； \end{tabularx} \\ \\ \relax
16:23 & \begin{tabularx}{0.7\textwidth}{X} 打了許多棍，就把他們下在監裡，囑咐獄警嚴緊看守。 \end{tabularx} \\ \\ \relax
16:24 & \begin{tabularx}{0.7\textwidth}{X} 獄警領了這樣的命令，就把他們下在內監，兩腳拴在木架上。 \end{tabularx} \\ \\ \relax
16:25 & \begin{tabularx}{0.7\textwidth}{X} 約在半夜，保羅和西拉正在禱告，唱詩讚美神，眾囚犯也側耳聽著的時候， \end{tabularx} \\ \\ \relax
16:26 & \begin{tabularx}{0.7\textwidth}{X} 忽然，地大震動，甚至監牢的地基都搖動了，監門立刻全開，眾囚犯的鎖鏈也都解開了。 \end{tabularx} \\ \\ \relax
16:27 & \begin{tabularx}{0.7\textwidth}{X} 獄警一醒，看見監門全開，以為囚犯已經逃走，就拔刀要自殺。 \end{tabularx} \\ \\ \relax
16:28 & \begin{tabularx}{0.7\textwidth}{X} 保羅大聲呼叫：「不要傷害自己！我們都在這裡。」 \end{tabularx} \\ \\ \relax
16:29 & \begin{tabularx}{0.7\textwidth}{X} 獄警叫人拿燈來，就衝進去，戰戰兢兢地俯伏在保羅和西拉面前。 \end{tabularx} \\ \\ \relax
16:30 & \begin{tabularx}{0.7\textwidth}{X} 然後獄警領他們出來，說：「二位先生，我必須做甚麼才可以得救？」 \end{tabularx} \\ \\ \relax
16:31 & \begin{tabularx}{0.7\textwidth}{X} 他們說：「當信主耶穌，你和你一家都必得救。」 \end{tabularx} \\ \\ \relax
16:32 & \begin{tabularx}{0.7\textwidth}{X} 他們就把主的道講給他和他全家的人聽。 \end{tabularx} \\ \\ \relax
16:33 & \begin{tabularx}{0.7\textwidth}{X} 當夜，就在那時候，獄警把他們帶去，洗他們的傷；他和他所有的家人立刻都受了洗。 \end{tabularx} \\ \\ \relax
16:34 & \begin{tabularx}{0.7\textwidth}{X} 於是獄警領他們上自己的家裡去，給他們擺上飯。他和全家的人，因為信了神，都滿心喜樂。 \end{tabularx} \\ \\ \relax
16:35 & \begin{tabularx}{0.7\textwidth}{X} 到了天亮，官長們打發差役來，說：「釋放那兩個人吧。」 \end{tabularx} \\ \\ \relax
16:36 & \begin{tabularx}{0.7\textwidth}{X} 獄警就把這些話告訴保羅：「官長們打發人來，要釋放你們，現在可以出監，平平安安去吧。」 \end{tabularx} \\ \\ \relax
16:37 & \begin{tabularx}{0.7\textwidth}{X} 保羅卻說：「我們是羅馬人，並沒有定罪，他們竟在公眾面前打了我們，又把我們下在監裡；現在要私下趕我們出去嗎？這不行！叫他們自己來領我們出去吧！」 \end{tabularx} \\ \\ \relax
16:38 & \begin{tabularx}{0.7\textwidth}{X} 差役把這些話回稟官長們；官長們聽見他們是羅馬人，就害怕了， \end{tabularx} \\ \\ \relax
16:39 & \begin{tabularx}{0.7\textwidth}{X} 於是來勸他們，領他們出來，請他們離開那城。 \end{tabularx} \\ \\ \relax
16:40 & \begin{tabularx}{0.7\textwidth}{X} 二人出了監牢，往呂底亞家裡去，見了弟兄們，勸慰他們一番，就離開了。 \end{tabularx} \\ \\ \relax
17:1 & \begin{tabularx}{0.7\textwidth}{X} 保羅和西拉經過暗妃坡里、亞波羅尼亞，來到帖撒羅尼迦，在那裡有猶太人的會堂。 \end{tabularx} \\ \\ \relax
17:2 & \begin{tabularx}{0.7\textwidth}{X} 保羅照他素常的規矩進去，一連三個安息日，根據聖經與他們辯論， \end{tabularx} \\ \\ \relax
17:3 & \begin{tabularx}{0.7\textwidth}{X} 講解和說明基督必須受害，從死人中復活；又說：「我所傳給你們的這位耶穌就是基督。」 \end{tabularx} \\ \\ \relax
17:4 & \begin{tabularx}{0.7\textwidth}{X} 他們中間有些人聽了勸，就跟從保羅和西拉，還有許多虔敬的希臘人，尊貴的婦女也不少。 \end{tabularx} \\ \\ \relax
17:5 & \begin{tabularx}{0.7\textwidth}{X} 但不信的猶太人心裡嫉妒，聚集了些市井流氓，搭夥成群，煽動全城的人闖進耶孫的家，要把保羅和西拉帶到民眾那裡。 \end{tabularx} \\ \\ \relax
17:6 & \begin{tabularx}{0.7\textwidth}{X} 那些人找不著他們，就把耶孫和幾個弟兄拉到地方官那裡，喊叫著：「這些攪亂天下的人也到這裡來了， \end{tabularx} \\ \\ \relax
17:7 & \begin{tabularx}{0.7\textwidth}{X} 耶孫竟收留他們。這些人都違背凱撒的命令，說另有一個王耶穌。」 \end{tabularx} \\ \\ \relax
17:8 & \begin{tabularx}{0.7\textwidth}{X} 眾人和地方官聽見這些話，就惶恐了， \end{tabularx} \\ \\ \relax
17:9 & \begin{tabularx}{0.7\textwidth}{X} 於是收了耶孫和其餘的人的保證金後，釋放了他們。 \end{tabularx} \\ \\ \relax
17:10 & \begin{tabularx}{0.7\textwidth}{X} 當夜，弟兄們立刻送保羅和西拉往庇哩亞去；二人到了，就進入猶太人的會堂。 \end{tabularx} \\ \\ \relax
17:11 & \begin{tabularx}{0.7\textwidth}{X} 這地方的猶太人比帖撒羅尼迦的人開明，熱心領受這道，天天查考聖經，要知道這道是否真實。 \end{tabularx} \\ \\ \relax
17:12 & \begin{tabularx}{0.7\textwidth}{X} 所以，他們中間有許多信了，又有希臘的尊貴婦人，男人也不少。 \end{tabularx} \\ \\ \relax
17:13 & \begin{tabularx}{0.7\textwidth}{X} 但帖撒羅尼迦的猶太人知道保羅又在庇哩亞傳神的道，就往那裡去，煽動挑撥群眾。 \end{tabularx} \\ \\ \relax
17:14 & \begin{tabularx}{0.7\textwidth}{X} 於是，弟兄們立刻送保羅到海邊去，西拉和提摩太卻仍留在庇哩亞。 \end{tabularx} \\ \\ \relax
17:15 & \begin{tabularx}{0.7\textwidth}{X} 護送保羅的人帶他到了雅典，他們領了保羅的命令，叫西拉和提摩太趕快到他那裡來，然後回去了。 \end{tabularx} \\ \\ \relax
17:16 & \begin{tabularx}{0.7\textwidth}{X} 保羅在雅典等候他們的時候，看見滿城都是偶像，就心裡非常難過。 \end{tabularx} \\ \\ \relax
17:17 & \begin{tabularx}{0.7\textwidth}{X} 於是他在會堂裡與猶太人和虔敬的人，以及每日在市場上所遇見的人辯論。 \end{tabularx} \\ \\ \relax
17:18 & \begin{tabularx}{0.7\textwidth}{X} 還有伊壁鳩魯和斯多亞兩派的哲學家也與他爭辯。有的說：「這胡言亂語的要說甚麼？」有的說：「他似乎是宣傳外邦鬼神的。」這是因保羅傳講耶穌與復活的福音。 \end{tabularx} \\ \\ \relax
17:19 & \begin{tabularx}{0.7\textwidth}{X} 他們就把他帶到亞略巴古，說：「你所講的這新學說，我們也可以知道嗎？ \end{tabularx} \\ \\ \relax
17:20 & \begin{tabularx}{0.7\textwidth}{X} 因為你有些奇怪的事傳到我們耳中，我們想知道這些事是甚麼意思。」 \end{tabularx} \\ \\ \relax
17:21 & \begin{tabularx}{0.7\textwidth}{X} 原來所有的雅典人和居住在那裡的外國人都無暇管別的事，只是談談或聽聽新聞。 \end{tabularx} \\ \\ \relax
17:22 & \begin{tabularx}{0.7\textwidth}{X} 保羅站在亞略巴古當中，說：「諸位雅典人！我看你們凡事很敬畏鬼神。 \end{tabularx} \\ \\ \relax
17:23 & \begin{tabularx}{0.7\textwidth}{X} 我到處走走的時候，仔細觀察你們所敬拜的，發現一座壇，上面寫著『獻給未識之神明』。你們所不認識而敬拜的，我現在向你們宣告： \end{tabularx} \\ \\ \relax
17:24 & \begin{tabularx}{0.7\textwidth}{X} 他是創造宇宙和其中萬物的神；他既是天地的主，就不住在人手所造的殿宇裡， \end{tabularx} \\ \\ \relax
17:25 & \begin{tabularx}{0.7\textwidth}{X} 也不用人手去服侍，好像缺少甚麼似的；自己倒將生命、氣息、萬物賜給萬人。 \end{tabularx} \\ \\ \relax
17:26 & \begin{tabularx}{0.7\textwidth}{X} 他從一人造出萬族，居住在全地面上，並且預先定準他們的年限和所住的疆界， \end{tabularx} \\ \\ \relax
17:27 & \begin{tabularx}{0.7\textwidth}{X} 為要使他們尋求神，或者可以揣摩而找到他，其實他離我們各人不遠。 \end{tabularx} \\ \\ \relax
17:28 & \begin{tabularx}{0.7\textwidth}{X} 我們生活、行動、存在都在於他。就如你們的詩人也有人說：『我們也是他所生的。』 \end{tabularx} \\ \\ \relax
17:29 & \begin{tabularx}{0.7\textwidth}{X} 既然我們是神所生的，就不應該以為神的神性像人用手藝和心思所雕刻的金、銀、石像一般。 \end{tabularx} \\ \\ \relax
17:30 & \begin{tabularx}{0.7\textwidth}{X} 世人蒙昧無知的時候，神並不追究，如今卻吩咐各處的人都要悔改。 \end{tabularx} \\ \\ \relax
17:31 & \begin{tabularx}{0.7\textwidth}{X} 因為他已經定了日子，要藉著他所設立的人按公義審判天下，並且使他從死人中復活，給萬人作可信的憑據。」 \end{tabularx} \\ \\ \relax
17:32 & \begin{tabularx}{0.7\textwidth}{X} 眾人聽見死人復活的話，就有人譏誚他；又有人說：「我們會再聽你講這事。」 \end{tabularx} \\ \\ \relax
17:33 & \begin{tabularx}{0.7\textwidth}{X} 於是保羅從他們當中出去了。 \end{tabularx} \\ \\ \relax
17:34 & \begin{tabularx}{0.7\textwidth}{X} 但有幾個人依附他，信了主，其中有亞略巴古的議員丟尼修，和一個名叫大馬哩的婦人，還有幾個與他們一起的人。 \end{tabularx} \\ \\ \relax
18:1 & \begin{tabularx}{0.7\textwidth}{X} 這些事以後，保羅離開雅典，來到哥林多。 \end{tabularx} \\ \\ \relax
18:2 & \begin{tabularx}{0.7\textwidth}{X} 他遇見一個生在本都的猶太人，名叫亞居拉。不久前，他帶著妻子百基拉從意大利來，因為克勞第命令所有的猶太人都離開羅馬。保羅去投靠他們。 \end{tabularx} \\ \\ \relax
18:3 & \begin{tabularx}{0.7\textwidth}{X} 他們本是製造帳棚為業。保羅因與他們同業，就和他們同住，一同做工。 \end{tabularx} \\ \\ \relax
18:4 & \begin{tabularx}{0.7\textwidth}{X} 每逢安息日，保羅在會堂裡辯論，勸導猶太人和希臘人。 \end{tabularx} \\ \\ \relax
18:5 & \begin{tabularx}{0.7\textwidth}{X} 西拉和提摩太從馬其頓來的時候，保羅正專心傳道，向猶太人證明耶穌是基督。 \end{tabularx} \\ \\ \relax
18:6 & \begin{tabularx}{0.7\textwidth}{X} 當他們抗拒他、毀謗他的時候，他就抖掉衣裳的灰塵，對他們說：「你們的罪歸到你們自己的頭上，與我無干。從今以後，我要往外邦人那裡去。」 \end{tabularx} \\ \\ \relax
18:7 & \begin{tabularx}{0.7\textwidth}{X} 於是他離開那裡，到了一個人的家裡，他名叫提多‧猶士都，是敬拜神的人，他的家靠近會堂。 \end{tabularx} \\ \\ \relax
18:8 & \begin{tabularx}{0.7\textwidth}{X} 會堂的主管基利司布和全家都信了主，還有許多哥林多人聽了就信，而且受了洗。 \end{tabularx} \\ \\ \relax
18:9 & \begin{tabularx}{0.7\textwidth}{X} 夜間，主在異象中對保羅說：「不要怕，只管講，不要沉默， \end{tabularx} \\ \\ \relax
18:10 & \begin{tabularx}{0.7\textwidth}{X} 有我與你同在，沒有人會下手害你，因為在這城裡有許多屬我的人。」 \end{tabularx} \\ \\ \relax
18:11 & \begin{tabularx}{0.7\textwidth}{X} 保羅在那裡住了一年六個月，將神的道教導他們。 \end{tabularx} \\ \\ \relax
18:12 & \begin{tabularx}{0.7\textwidth}{X} 到迦流作亞該亞省長的時候，猶太人齊心起來攻擊保羅，拉他到法庭， \end{tabularx} \\ \\ \relax
18:13 & \begin{tabularx}{0.7\textwidth}{X} 說：「這個人教唆人不按著律法敬拜神。」 \end{tabularx} \\ \\ \relax
18:14 & \begin{tabularx}{0.7\textwidth}{X} 保羅剛要開口，迦流對猶太人說：「你們這些猶太人哪！如果是為冤枉或奸惡的事，我理當耐性聽你們。 \end{tabularx} \\ \\ \relax
18:15 & \begin{tabularx}{0.7\textwidth}{X} 既然你們所爭論的是關乎用字、名目和你們的律法，你們自己去辦吧！這樣的事我不願意審問。」 \end{tabularx} \\ \\ \relax
18:16 & \begin{tabularx}{0.7\textwidth}{X} 於是，他把他們逐出法庭。 \end{tabularx} \\ \\ \relax
18:17 & \begin{tabularx}{0.7\textwidth}{X} 眾人就揪住會堂的主管所提尼，在法庭前打他。這些事迦流都不管。 \end{tabularx} \\ \\ \relax
18:18 & \begin{tabularx}{0.7\textwidth}{X} 保羅又住了好些日子，就辭別了弟兄，坐船到敘利亞去。百基拉、亞居拉和他同去。他因為許過願，就在堅革哩剃了頭髮。 \end{tabularx} \\ \\ \relax
18:19 & \begin{tabularx}{0.7\textwidth}{X} 到了以弗所，保羅就把他們留在那裡，自己進了會堂，和猶太人辯論。 \end{tabularx} \\ \\ \relax
18:20 & \begin{tabularx}{0.7\textwidth}{X} 眾人請他多住些日子，他沒有答應， \end{tabularx} \\ \\ \relax
18:21 & \begin{tabularx}{0.7\textwidth}{X} 就辭別他們，說：「神若許可，我還要回到你們這裡來。」於是他上船離開以弗所。 \end{tabularx} \\ \\ \relax
18:22 & \begin{tabularx}{0.7\textwidth}{X} 他在凱撒利亞下了船，上耶路撒冷去問候教會，隨後下安提阿去。 \end{tabularx} \\ \\ \relax
18:23 & \begin{tabularx}{0.7\textwidth}{X} 他在那裡住了些日子，又離開了那裡，逐一經過加拉太和弗呂家各地方，堅固眾門徒。 \end{tabularx} \\ \\ \relax
18:24 & \begin{tabularx}{0.7\textwidth}{X} 有一個生在亞歷山大的猶太人，名叫亞波羅，來到以弗所，他很有口才，很會講解聖經。 \end{tabularx} \\ \\ \relax
18:25 & \begin{tabularx}{0.7\textwidth}{X} 這人已經在主的道路上受了訓練，心裡火熱，精確地講論和教導耶穌的事；可是他只知道約翰的洗禮。 \end{tabularx} \\ \\ \relax
18:26 & \begin{tabularx}{0.7\textwidth}{X} 他開始在會堂裡放膽講道；百基拉、亞居拉聽見，就接他來，將神的道路給他更精確地講解。 \end{tabularx} \\ \\ \relax
18:27 & \begin{tabularx}{0.7\textwidth}{X} 他想要往亞該亞去，弟兄們就勉勵他，並寫信請門徒們接待他，他到了那裡，多多幫助那些蒙恩信主的人， \end{tabularx} \\ \\ \relax
18:28 & \begin{tabularx}{0.7\textwidth}{X} 因為他在公眾面前極力駁倒猶太人，引聖經證明耶穌是基督。 \end{tabularx} \\ \\ \relax
19:1 & \begin{tabularx}{0.7\textwidth}{X} 亞波羅在哥林多的時候，保羅經過了內陸地區，來到以弗所，在那裡他遇見幾個門徒， \end{tabularx} \\ \\ \relax
19:2 & \begin{tabularx}{0.7\textwidth}{X} 問他們：「你們信的時候領受了聖靈沒有？」他們說：「沒有，我們連甚麼是聖靈都沒有聽過。」 \end{tabularx} \\ \\ \relax
19:3 & \begin{tabularx}{0.7\textwidth}{X} 保羅說：「這樣，你們受的是甚麼洗呢？」他們說：「是受了約翰的洗。」 \end{tabularx} \\ \\ \relax
19:4 & \begin{tabularx}{0.7\textwidth}{X} 保羅說：「約翰所施的是悔改的洗禮，他告訴百姓當信那在他以後要來的那位，就是耶穌。」 \end{tabularx} \\ \\ \relax
19:5 & \begin{tabularx}{0.7\textwidth}{X} 他們聽見這話以後，就奉主耶穌的名受洗。 \end{tabularx} \\ \\ \relax
19:6 & \begin{tabularx}{0.7\textwidth}{X} 保羅給他們按手，聖靈就降在他們身上，他們開始說方言 和說預言。 \end{tabularx} \\ \\ \relax
19:7 & \begin{tabularx}{0.7\textwidth}{X} 他們約有十二個人。 \end{tabularx} \\ \\ \relax
19:8 & \begin{tabularx}{0.7\textwidth}{X} 保羅進會堂，一連三個月放膽講道，辯論神國的事，勸導眾人。 \end{tabularx} \\ \\ \relax
19:9 & \begin{tabularx}{0.7\textwidth}{X} 後來，有些人心裡剛硬不信，在眾人面前毀謗這道；保羅就離開他們，也叫門徒與他們分開，就在推喇奴的講堂天天辯論。 \end{tabularx} \\ \\ \relax
19:10 & \begin{tabularx}{0.7\textwidth}{X} 這樣有兩年之久，使一切住在亞細亞的，無論是猶太人是希臘人，都聽見主的道。 \end{tabularx} \\ \\ \relax
19:11 & \begin{tabularx}{0.7\textwidth}{X} 神藉保羅的手行了些奇異的神蹟， \end{tabularx} \\ \\ \relax
19:12 & \begin{tabularx}{0.7\textwidth}{X} 甚至有人從保羅身上拿走手巾或圍裙放在病人身上，病就消除了，邪靈也出去了。 \end{tabularx} \\ \\ \relax
19:13 & \begin{tabularx}{0.7\textwidth}{X} 那時，有幾個巡迴各處念咒趕鬼的猶太人，擅自利用主耶穌的名，向那些被邪靈所附的人說：「我奉保羅所傳的耶穌命令你們出來！」 \end{tabularx} \\ \\ \relax
19:14 & \begin{tabularx}{0.7\textwidth}{X} 做這事的是猶太祭司長士基瓦的七個兒子。 \end{tabularx} \\ \\ \relax
19:15 & \begin{tabularx}{0.7\textwidth}{X} 但邪靈回答他們：「耶穌我知道，保羅我也認識，你們卻是誰呢？」 \end{tabularx} \\ \\ \relax
19:16 & \begin{tabularx}{0.7\textwidth}{X} 被邪靈所附的人就撲到他們身上，制伏他們，勝過他們，使他們赤著身子，受了傷，從那房子裡逃出去了。 \end{tabularx} \\ \\ \relax
19:17 & \begin{tabularx}{0.7\textwidth}{X} 凡住在以弗所的，無論是猶太人是希臘人，都知道這件事，也都懼怕；主耶穌的名從此就更被尊為大了。 \end{tabularx} \\ \\ \relax
19:18 & \begin{tabularx}{0.7\textwidth}{X} 許多已經信的人來承認並公開自己所行的事。 \end{tabularx} \\ \\ \relax
19:19 & \begin{tabularx}{0.7\textwidth}{X} 又有許多平素行邪術的人把他們的書都拿來，堆積在眾人面前焚燒。他們計算書價，得知共值五萬塊銀錢。 \end{tabularx} \\ \\ \relax
19:20 & \begin{tabularx}{0.7\textwidth}{X} 這樣，主的道大大興旺，而且普遍傳開了。 \end{tabularx} \\ \\ \relax
19:21 & \begin{tabularx}{0.7\textwidth}{X} 這些事過後，保羅心裡決定要經過馬其頓、亞該亞，就往耶路撒冷去。他說：「我到了那裡以後，也必須到羅馬去看看。」 \end{tabularx} \\ \\ \relax
19:22 & \begin{tabularx}{0.7\textwidth}{X} 於是他差遣兩個助手提摩太和以拉都往馬其頓去，自己暫時留在亞細亞。 \end{tabularx} \\ \\ \relax
19:23 & \begin{tabularx}{0.7\textwidth}{X} 那時，因這道路而起的騷動不小。 \end{tabularx} \\ \\ \relax
19:24 & \begin{tabularx}{0.7\textwidth}{X} 有一個銀匠，名叫底米丟，是製造亞底米神銀龕的，他使從事這手藝的人生意發達。 \end{tabularx} \\ \\ \relax
19:25 & \begin{tabularx}{0.7\textwidth}{X} 他聚集他們和同行的工人，說：「諸位，你們知道我們是倚靠這生意發財的。 \end{tabularx} \\ \\ \relax
19:26 & \begin{tabularx}{0.7\textwidth}{X} 你們看到，也聽見這保羅不但在以弗所，也幾乎在亞細亞全地，引誘迷惑了許多人，說：『人手所做的不是神明。』 \end{tabularx} \\ \\ \relax
19:27 & \begin{tabularx}{0.7\textwidth}{X} 這樣，不僅我們這行業陷入被藐視的危險，就是大女神亞底米的廟也要被人輕看，連亞細亞全地和普天下所敬拜的女神的威望也受損害了。」 \end{tabularx} \\ \\ \relax
19:28 & \begin{tabularx}{0.7\textwidth}{X} 眾人聽見，就怒氣沖沖，喊著說：「大哉，以弗所人的亞底米！」 \end{tabularx} \\ \\ \relax
19:29 & \begin{tabularx}{0.7\textwidth}{X} 於是滿城都騷動起來。眾人抓住與保羅同行的馬其頓人該猶和亞里達古，齊心衝進劇場。 \end{tabularx} \\ \\ \relax
19:30 & \begin{tabularx}{0.7\textwidth}{X} 保羅想要進到民眾那裡，門徒卻不許他去。 \end{tabularx} \\ \\ \relax
19:31 & \begin{tabularx}{0.7\textwidth}{X} 連亞細亞的幾位官員，是保羅的朋友，也打發人來勸他不要冒險到劇場裡去。 \end{tabularx} \\ \\ \relax
19:32 & \begin{tabularx}{0.7\textwidth}{X} 聚集的人亂成一團，有的喊這個，有的喊那個，大半不知道為了甚麼聚集。 \end{tabularx} \\ \\ \relax
19:33 & \begin{tabularx}{0.7\textwidth}{X} 猶太人把亞歷山大推出去，人群中有人慫恿他，他就做手勢，要向民眾申訴。 \end{tabularx} \\ \\ \relax
19:34 & \begin{tabularx}{0.7\textwidth}{X} 但他們一認出他是猶太人，大家就異口同聲喊著：「大哉，以弗所人的亞底米！」約喊了兩小時。 \end{tabularx} \\ \\ \relax
19:35 & \begin{tabularx}{0.7\textwidth}{X} 城裡的書記官安撫了群眾後，說：「以弗所人哪，誰不知道以弗所人的城是看守大亞底米的廟和從宙斯那裡落下來的像的守護者呢？ \end{tabularx} \\ \\ \relax
19:36 & \begin{tabularx}{0.7\textwidth}{X} 既然這些事是駁不倒的，你們就要安靜下來，不可妄動。 \end{tabularx} \\ \\ \relax
19:37 & \begin{tabularx}{0.7\textwidth}{X} 你們把這些人帶來，他們並沒有偷竊廟中之物，也沒有褻瀆我們的女神。 \end{tabularx} \\ \\ \relax
19:38 & \begin{tabularx}{0.7\textwidth}{X} 如果底米丟和他同行的手藝人有控告的事，自有公堂，也有省長，他們可以彼此控告。 \end{tabularx} \\ \\ \relax
19:39 & \begin{tabularx}{0.7\textwidth}{X} 你們若有別的事請求，可以在合法的集會裡解決。 \end{tabularx} \\ \\ \relax
19:40 & \begin{tabularx}{0.7\textwidth}{X} 今日的擾亂本是無緣無故的，有被控告的危險。這次的騷動，我們也說不出理由來。」 \end{tabularx} \\ \\ \relax
19:41 & \begin{tabularx}{0.7\textwidth}{X} 他說完這些話，就叫眾人散會。 \end{tabularx} \\ \\ \relax
20:1 & \begin{tabularx}{0.7\textwidth}{X} 騷亂平定以後，保羅請門徒來，勸勉了他們，就辭別他們，往馬其頓去。 \end{tabularx} \\ \\ \relax
20:2 & \begin{tabularx}{0.7\textwidth}{X} 他走遍那一帶地方，用許多話勸勉門徒，然後來到希臘， \end{tabularx} \\ \\ \relax
20:3 & \begin{tabularx}{0.7\textwidth}{X} 在那裡住了三個月。他快要坐船往敘利亞去的時候，猶太人設計害他，他就決定從馬其頓回去。 \end{tabularx} \\ \\ \relax
20:4 & \begin{tabularx}{0.7\textwidth}{X} 同他到亞細亞去的，有庇哩亞人畢羅斯的兒子所巴特，帖撒羅尼迦人亞里達古和西公都，還有特庇人該猶和提摩太，又有亞細亞人推基古和特羅非摩。 \end{tabularx} \\ \\ \relax
20:5 & \begin{tabularx}{0.7\textwidth}{X} 這些人先走，在特羅亞等候我們。 \end{tabularx} \\ \\ \relax
20:6 & \begin{tabularx}{0.7\textwidth}{X} 過了除酵節的日子，我們從腓立比開船，五天以後到了特羅亞，和他們相會，在那裡住了七天。 \end{tabularx} \\ \\ \relax
20:7 & \begin{tabularx}{0.7\textwidth}{X} 七日的第一日，我們聚會擘餅的時候，保羅因次日要起行，就為他們講道，直講到半夜。 \end{tabularx} \\ \\ \relax
20:8 & \begin{tabularx}{0.7\textwidth}{X} 我們聚會的那座樓上有好些燈火。 \end{tabularx} \\ \\ \relax
20:9 & \begin{tabularx}{0.7\textwidth}{X} 有一個少年，名叫猶推古，坐在窗口上，沉沉入睡。保羅講了多時，少年睡熟了，從三層樓上掉下去，扶起來時已經死了。 \end{tabularx} \\ \\ \relax
20:10 & \begin{tabularx}{0.7\textwidth}{X} 保羅下去，伏在他身上，抱著他，說：「你們不要慌亂，他還有氣呢！」 \end{tabularx} \\ \\ \relax
20:11 & \begin{tabularx}{0.7\textwidth}{X} 保羅又上樓去，擘餅，吃了，再講了許久，直到天亮才離開。 \end{tabularx} \\ \\ \relax
20:12 & \begin{tabularx}{0.7\textwidth}{X} 他們把那活過來的孩子帶走，大家得到很大的安慰。 \end{tabularx} \\ \\ \relax
20:13 & \begin{tabularx}{0.7\textwidth}{X} 我們先上船，起航往亞朔去，想要在那裡接保羅；因為他是這樣安排的，他自己本來打算要走陸路。 \end{tabularx} \\ \\ \relax
20:14 & \begin{tabularx}{0.7\textwidth}{X} 他既在亞朔與我們相會，我們就接他上船，來到米推利尼。 \end{tabularx} \\ \\ \relax
20:15 & \begin{tabularx}{0.7\textwidth}{X} 我們從那裡開船，第二天到了基阿的對岸；再下一天，在撒摩靠岸，又過了一天，到了米利都。 \end{tabularx} \\ \\ \relax
20:16 & \begin{tabularx}{0.7\textwidth}{X} 因為保羅早已決定要越過以弗所，免得在亞細亞耽延，他急忙前行，假如可能的話，在五旬節前能趕到耶路撒冷。 \end{tabularx} \\ \\ \relax
20:17 & \begin{tabularx}{0.7\textwidth}{X} 保羅從米利都打發人往以弗所去，請教會的長老來。 \end{tabularx} \\ \\ \relax
20:18 & \begin{tabularx}{0.7\textwidth}{X} 他們來了，保羅對他們說：「你們自己知道，自從我到亞細亞的第一天，我怎樣跟你們相處， \end{tabularx} \\ \\ \relax
20:19 & \begin{tabularx}{0.7\textwidth}{X} 怎樣凡事謙卑，以眼淚服侍主，又因猶太人的謀害經歷試煉。 \end{tabularx} \\ \\ \relax
20:20 & \begin{tabularx}{0.7\textwidth}{X} 你們也知道，凡對你們有益的，我沒有一樣隱瞞不說的，或在公眾面前，或在每一個人的家裡，我都教導你們， \end{tabularx} \\ \\ \relax
20:21 & \begin{tabularx}{0.7\textwidth}{X} 不論猶太人和希臘人，我都已證明他們當在神面前悔改，信靠我們的主耶穌。 \end{tabularx} \\ \\ \relax
20:22 & \begin{tabularx}{0.7\textwidth}{X} 現在我被聖靈催迫要往耶路撒冷去，雖然不知道在那裡會遭遇甚麼事， \end{tabularx} \\ \\ \relax
20:23 & \begin{tabularx}{0.7\textwidth}{X} 但知道聖靈在各城裡向我指證，說有捆鎖與患難等著我。 \end{tabularx} \\ \\ \relax
20:24 & \begin{tabularx}{0.7\textwidth}{X} 我卻不以性命為念，只要走完我的路程，完成我從主耶穌所領受的職分，為神恩典的福音作見證。 \end{tabularx} \\ \\ \relax
20:25 & \begin{tabularx}{0.7\textwidth}{X} 「我素常在你們中間到處傳講神的國；現在我知道，你們眾人以後不會再見到我的面了。 \end{tabularx} \\ \\ \relax
20:26 & \begin{tabularx}{0.7\textwidth}{X} 所以我今日向你們作證，你們中間無論何人死亡，罪不在我。 \end{tabularx} \\ \\ \relax
20:27 & \begin{tabularx}{0.7\textwidth}{X} 因為神一切的旨意，我並沒有退縮不傳給你們的。 \end{tabularx} \\ \\ \relax
20:28 & \begin{tabularx}{0.7\textwidth}{X} 聖靈立你們作全群的監督，你們就當為自己謹慎，也為全群謹慎，牧養神的教會，就是他用自己血所買來的。 \end{tabularx} \\ \\ \relax
20:29 & \begin{tabularx}{0.7\textwidth}{X} 我知道，在我離開以後必有兇暴的豺狼進入你們中間，不顧惜羊群。 \end{tabularx} \\ \\ \relax
20:30 & \begin{tabularx}{0.7\textwidth}{X} 就是你們中間也必有人起來，說悖謬的話，要引誘門徒跟從他們。 \end{tabularx} \\ \\ \relax
20:31 & \begin{tabularx}{0.7\textwidth}{X} 所以你們要警醒，記念我三年之久，晝夜不斷地流淚勸戒你們各人。 \end{tabularx} \\ \\ \relax
20:32 & \begin{tabularx}{0.7\textwidth}{X} 現在我把你們交託給神和他恩惠的道；這道能建立你們，使你們和一切成聖的人同得基業。 \end{tabularx} \\ \\ \relax
20:33 & \begin{tabularx}{0.7\textwidth}{X} 我未曾貪圖一個人的金、銀或衣服。 \end{tabularx} \\ \\ \relax
20:34 & \begin{tabularx}{0.7\textwidth}{X} 你們自己知道，我靠兩隻手工作來供給我和同工的需用。 \end{tabularx} \\ \\ \relax
20:35 & \begin{tabularx}{0.7\textwidth}{X} 我凡事給你們作榜樣，叫你們知道應當這樣勞苦，扶助軟弱的人，又當記念主耶穌的話，說：『施比受更為有福。』」 \end{tabularx} \\ \\ \relax
20:36 & \begin{tabularx}{0.7\textwidth}{X} 保羅說完了這些話，就和大家跪下來禱告。 \end{tabularx} \\ \\ \relax
20:37 & \begin{tabularx}{0.7\textwidth}{X} 眾人痛哭，抱著保羅的頸項跟他親吻。 \end{tabularx} \\ \\ \relax
20:38 & \begin{tabularx}{0.7\textwidth}{X} 叫他們最傷心的，就是他說「以後不會再見到我的面」那句話。於是他們送他上船去了。 \end{tabularx} \\ \\ \relax
21:1 & \begin{tabularx}{0.7\textwidth}{X} 我們離別了眾人，就開船直航到哥士，第二天到了羅底，又從那裡到帕大喇。 \end{tabularx} \\ \\ \relax
21:2 & \begin{tabularx}{0.7\textwidth}{X} 我們遇見一隻船要往腓尼基去，就上船起航。 \end{tabularx} \\ \\ \relax
21:3 & \begin{tabularx}{0.7\textwidth}{X} 我們望見塞浦路斯，就從南邊行過，往敘利亞去，在推羅上岸，因為船要在那裡卸貨。 \end{tabularx} \\ \\ \relax
21:4 & \begin{tabularx}{0.7\textwidth}{X} 我們在那裡找到了一些門徒，就住了七天。他們藉著聖靈的感動，告訴保羅不要上耶路撒冷去。 \end{tabularx} \\ \\ \relax
21:5 & \begin{tabularx}{0.7\textwidth}{X} 幾天之後，我們又出發前行。他們眾人同妻子兒女都送我們到城外，我們都跪在灘上禱告，彼此辭別。 \end{tabularx} \\ \\ \relax
21:6 & \begin{tabularx}{0.7\textwidth}{X} 我們上了船，他們就回家去了。 \end{tabularx} \\ \\ \relax
21:7 & \begin{tabularx}{0.7\textwidth}{X} 我們從推羅行完航程，來到了多利買，問候那裡的弟兄，和他們同住了一天。 \end{tabularx} \\ \\ \relax
21:8 & \begin{tabularx}{0.7\textwidth}{X} 第二天，我們離開那裡，來到凱撒利亞，就進了傳福音的腓利家裡，和他同住；他是那七個執事裡的一個。 \end{tabularx} \\ \\ \relax
21:9 & \begin{tabularx}{0.7\textwidth}{X} 他有四個女兒，都是未出嫁的，都會說預言。 \end{tabularx} \\ \\ \relax
21:10 & \begin{tabularx}{0.7\textwidth}{X} 我們在那裡多住了好幾天，有一個先知，名叫亞迦布，從猶太下來。 \end{tabularx} \\ \\ \relax
21:11 & \begin{tabularx}{0.7\textwidth}{X} 他到了我們這裡，就拿保羅的腰帶，捆上自己的手腳，說：「聖靈這樣說：『猶太人在耶路撒冷要如此捆綁這腰帶的主人，把他交在外邦人手裡。』」 \end{tabularx} \\ \\ \relax
21:12 & \begin{tabularx}{0.7\textwidth}{X} 我們聽見這些話，就跟當地的人苦勸保羅不要上耶路撒冷去。 \end{tabularx} \\ \\ \relax
21:13 & \begin{tabularx}{0.7\textwidth}{X} 於是保羅回答：「你們為甚麼這樣痛哭，使我心碎呢？我為主耶穌的名，不但被人捆綁，就是死在耶路撒冷也是願意的。」 \end{tabularx} \\ \\ \relax
21:14 & \begin{tabularx}{0.7\textwidth}{X} 既然保羅不聽勸，我們就住了口，只說：「願主的旨意成就。」 \end{tabularx} \\ \\ \relax
21:15 & \begin{tabularx}{0.7\textwidth}{X} 過了這幾天，我們收拾行李上耶路撒冷去。 \end{tabularx} \\ \\ \relax
21:16 & \begin{tabularx}{0.7\textwidth}{X} 有凱撒利亞的幾個門徒和我們同去，帶我們到一個早期的門徒塞浦路斯人拿孫的家裡，請我們與他同住。 \end{tabularx} \\ \\ \relax
21:17 & \begin{tabularx}{0.7\textwidth}{X} 我們到了耶路撒冷，弟兄們歡歡喜喜地接待我們。 \end{tabularx} \\ \\ \relax
21:18 & \begin{tabularx}{0.7\textwidth}{X} 第二天，保羅同我們去見雅各；所有的長老也都在場。 \end{tabularx} \\ \\ \relax
21:19 & \begin{tabularx}{0.7\textwidth}{X} 保羅向他們問安，然後將神用他在外邦人中所做的事奉，一一述說了。 \end{tabularx} \\ \\ \relax
21:20 & \begin{tabularx}{0.7\textwidth}{X} 他們聽見了，就歸榮耀給神，對保羅說：「弟兄，你看猶太人中有數以萬計的信徒，而他們都是熱心於律法的人。 \end{tabularx} \\ \\ \relax
21:21 & \begin{tabularx}{0.7\textwidth}{X} 他們曾聽見人說，你教導所有在外邦的猶太人離棄摩西，對他們說，不要給孩子行割禮，也不要遵守規矩。 \end{tabularx} \\ \\ \relax
21:22 & \begin{tabularx}{0.7\textwidth}{X} 眾人必聽見你來了，這可怎麼辦呢？ \end{tabularx} \\ \\ \relax
21:23 & \begin{tabularx}{0.7\textwidth}{X} 你就照著我們的話做吧！我們這裡有四個人，都有願在身。 \end{tabularx} \\ \\ \relax
21:24 & \begin{tabularx}{0.7\textwidth}{X} 你帶他們去，與他們一同行潔淨的禮，替他們繳納規費，讓他們得以剃頭。這樣，眾人就會知道，先前所聽見關於你的事都是假的；而且也知道，你自己為人循規蹈矩，遵行律法。 \end{tabularx} \\ \\ \relax
21:25 & \begin{tabularx}{0.7\textwidth}{X} 至於信主的外邦人，我們已經根據我們的決議寫信，叫他們要禁戒偶像所玷污的東西、血和勒死的牲畜，禁戒淫亂。」 \end{tabularx} \\ \\ \relax
21:26 & \begin{tabularx}{0.7\textwidth}{X} 於是保羅帶著那四個人，第二天與他們一同行了潔淨禮，進了聖殿，報告潔淨期滿的日子，等候祭司為他們各人獻上祭物。 \end{tabularx} \\ \\ \relax
21:27 & \begin{tabularx}{0.7\textwidth}{X} 那七日將完，從亞細亞來的猶太人看見保羅在聖殿裡，就煽動所有的群眾，下手拿住他， \end{tabularx} \\ \\ \relax
21:28 & \begin{tabularx}{0.7\textwidth}{X} 喊著：「以色列人哪，來幫忙！這就是在各處教導眾人糟蹋我們百姓、律法和這地方的人。不但如此，他還帶了希臘人進聖殿，污穢了這聖地。」 \end{tabularx} \\ \\ \relax
21:29 & \begin{tabularx}{0.7\textwidth}{X} 這話是因他們曾看見以弗所人特羅非摩跟保羅一起在城裡，以為保羅帶他進了聖殿。 \end{tabularx} \\ \\ \relax
21:30 & \begin{tabularx}{0.7\textwidth}{X} 於是全城都騷動，百姓一齊跑來，拿住保羅，拉他出聖殿，殿門立刻都關了。 \end{tabularx} \\ \\ \relax
21:31 & \begin{tabularx}{0.7\textwidth}{X} 他們正想要殺他，有人報信給營裡的千夫長，說耶路撒冷全城都亂了。 \end{tabularx} \\ \\ \relax
21:32 & \begin{tabularx}{0.7\textwidth}{X} 千夫長立刻帶著士兵和幾個百夫長，跑下去到他們那裡。他們見了千夫長和士兵，就停下來不打保羅。 \end{tabularx} \\ \\ \relax
21:33 & \begin{tabularx}{0.7\textwidth}{X} 於是千夫長上前拿住他，吩咐用兩條鐵鏈捆鎖，又問他是甚麼人，做了甚麼事。 \end{tabularx} \\ \\
[1ex]
\hline
\hline
\end{longtable}


\newpage

\section{第61屆~港九培靈會~劉銳光~第五講}
\label{sec:1132}
\textbf{見證}
\newline
\newline
連結: \href{https://www.hkbibleconference.org/session-message/view/1132}{\texttt{https://www.hkbibleconference.org/session-message/view/1132}}
\newline
\newline
\hyperref[sec:1130]{< < < PREV SERMON < < <}
~
\hyperref[sec:toc]{[返目錄]}
~
\hyperref[sec:1133]{> > > NEXT SERMON > > >}
\newline
\newline
使徒行傳 1:21-1:22
\newline
\begin{longtable}{cl}
\hline
\hline
章節 & 經文 (和合本修訂版)\\
\hline
1:21 & \begin{tabularx}{0.7\textwidth}{X} 所以，主耶穌在我們中間出入的整段時間， \end{tabularx} \\ \\ \relax
1:22 & \begin{tabularx}{0.7\textwidth}{X} 就是從約翰施洗起，直到主離開我們被接上升的日子為止，必須從那常與我們一起的人中，立一位與我們同作耶穌復活的見證。」 \end{tabularx} \\ \\
[1ex]
\hline
\hline
\end{longtable}


\newpage

\section{第61屆~港九培靈會~劉銳光~第六講}
\label{sec:1133}
\textbf{喜樂}
\newline
\newline
連結: \href{https://www.hkbibleconference.org/session-message/view/1133}{\texttt{https://www.hkbibleconference.org/session-message/view/1133}}
\newline
\newline
\hyperref[sec:1132]{< < < PREV SERMON < < <}
~
\hyperref[sec:toc]{[返目錄]}
~
\hyperref[sec:1135]{> > > NEXT SERMON > > >}
\newline
\newline
使徒行傳 2:46-2:47
\newline
\begin{longtable}{cl}
\hline
\hline
章節 & 經文 (和合本修訂版)\\
\hline
2:46 & \begin{tabularx}{0.7\textwidth}{X} 他們天天同心合意恆切地在聖殿裡敬拜，且在家中擘餅，存著歡喜坦誠的心用飯， \end{tabularx} \\ \\ \relax
2:47 & \begin{tabularx}{0.7\textwidth}{X} 讚美神，得全體百姓的喜愛。主將得救的人天天加給他們。 \end{tabularx} \\ \\
[1ex]
\hline
\hline
\end{longtable}


\newpage

\section{第61屆~港九培靈會~劉銳光~第七講}
\label{sec:1135}
\textbf{信仰純正}
\newline
\newline
連結: \href{https://www.hkbibleconference.org/session-message/view/1135}{\texttt{https://www.hkbibleconference.org/session-message/view/1135}}
\newline
\newline
\hyperref[sec:1133]{< < < PREV SERMON < < <}
~
\hyperref[sec:toc]{[返目錄]}
~
\hyperref[sec:1138]{> > > NEXT SERMON > > >}
\newline
\newline
使徒行傳 9:20-9:22
\newline
\begin{longtable}{cl}
\hline
\hline
章節 & 經文 (和合本修訂版)\\
\hline
9:20 & \begin{tabularx}{0.7\textwidth}{X} 立刻在各會堂裡傳揚耶穌，說他是神的兒子。 \end{tabularx} \\ \\ \relax
9:21 & \begin{tabularx}{0.7\textwidth}{X} 凡聽見的人都很驚奇，說：「在耶路撒冷殘害求告這名的不就是這個人嗎？他不是到這裡來要捆綁他們，帶到祭司長那裡去嗎？」 \end{tabularx} \\ \\ \relax
9:22 & \begin{tabularx}{0.7\textwidth}{X} 但掃羅越發有能力，駁倒住在大馬士革的猶太人，證明耶穌是基督。 \end{tabularx} \\ \\
[1ex]
\hline
\hline
\end{longtable}


\newpage

\section{第61屆~港九培靈會~劉銳光~第八講}
\label{sec:1138}
\textbf{傳福音}
\newline
\newline
連結: \href{https://www.hkbibleconference.org/session-message/view/1138}{\texttt{https://www.hkbibleconference.org/session-message/view/1138}}
\newline
\newline
\hyperref[sec:1135]{< < < PREV SERMON < < <}
~
\hyperref[sec:toc]{[返目錄]}
~
\hyperref[sec:1139]{> > > NEXT SERMON > > >}
\newline
\newline


\newpage

\section{第61屆~港九培靈會~劉銳光~第九講}
\label{sec:1139}
\textbf{聖靈}
\newline
\newline
連結: \href{https://www.hkbibleconference.org/session-message/view/1139}{\texttt{https://www.hkbibleconference.org/session-message/view/1139}}
\newline
\newline
\hyperref[sec:1138]{< < < PREV SERMON < < <}
~
\hyperref[sec:toc]{[返目錄]}
~
\hyperref[sec:1115]{> > > NEXT SERMON > > >}
\newline
\newline
使徒行傳 1:2-1:2
\newline
\begin{longtable}{cl}
\hline
\hline
章節 & 經文 (和合本修訂版)\\
\hline
1:2 & \begin{tabularx}{0.7\textwidth}{X} 直到他藉著聖靈吩咐所揀選的使徒後，被接上升的日子為止。 \end{tabularx} \\ \\
[1ex]
\hline
\hline
\end{longtable}


\newpage

\section{第61屆~港九培靈會~張子華~第一講}
\label{sec:1115}
\textbf{導論}
\newline
\newline
連結: \href{https://www.hkbibleconference.org/session-message/view/1115}{\texttt{https://www.hkbibleconference.org/session-message/view/1115}}
\newline
\newline
\hyperref[sec:1139]{< < < PREV SERMON < < <}
~
\hyperref[sec:toc]{[返目錄]}
~
\hyperref[sec:1116]{> > > NEXT SERMON > > >}
\newline
\newline


\newpage

\section{第61屆~港九培靈會~張子華~第二講}
\label{sec:1116}
\textbf{亞伯拉罕的信心 (一)}
\newline
\newline
連結: \href{https://www.hkbibleconference.org/session-message/view/1116}{\texttt{https://www.hkbibleconference.org/session-message/view/1116}}
\newline
\newline
\hyperref[sec:1115]{< < < PREV SERMON < < <}
~
\hyperref[sec:toc]{[返目錄]}
~
\hyperref[sec:1117]{> > > NEXT SERMON > > >}
\newline
\newline


\newpage

\section{第61屆~港九培靈會~張子華~第三講}
\label{sec:1117}
\textbf{亞伯拉罕的信心 (二)}
\newline
\newline
連結: \href{https://www.hkbibleconference.org/session-message/view/1117}{\texttt{https://www.hkbibleconference.org/session-message/view/1117}}
\newline
\newline
\hyperref[sec:1116]{< < < PREV SERMON < < <}
~
\hyperref[sec:toc]{[返目錄]}
~
\hyperref[sec:1118]{> > > NEXT SERMON > > >}
\newline
\newline


\newpage

\section{第61屆~港九培靈會~張子華~第四講}
\label{sec:1118}
\textbf{亞伯拉罕的信心 (三)}
\newline
\newline
連結: \href{https://www.hkbibleconference.org/session-message/view/1118}{\texttt{https://www.hkbibleconference.org/session-message/view/1118}}
\newline
\newline
\hyperref[sec:1117]{< < < PREV SERMON < < <}
~
\hyperref[sec:toc]{[返目錄]}
~
\hyperref[sec:1119]{> > > NEXT SERMON > > >}
\newline
\newline


\newpage

\section{第61屆~港九培靈會~張子華~第五講}
\label{sec:1119}
\textbf{亞伯拉罕的信心 (四)}
\newline
\newline
連結: \href{https://www.hkbibleconference.org/session-message/view/1119}{\texttt{https://www.hkbibleconference.org/session-message/view/1119}}
\newline
\newline
\hyperref[sec:1118]{< < < PREV SERMON < < <}
~
\hyperref[sec:toc]{[返目錄]}
~
\hyperref[sec:1120]{> > > NEXT SERMON > > >}
\newline
\newline


\newpage

\section{第61屆~港九培靈會~張子華~第六講}
\label{sec:1120}
\textbf{以撒的順服}
\newline
\newline
連結: \href{https://www.hkbibleconference.org/session-message/view/1120}{\texttt{https://www.hkbibleconference.org/session-message/view/1120}}
\newline
\newline
\hyperref[sec:1119]{< < < PREV SERMON < < <}
~
\hyperref[sec:toc]{[返目錄]}
~
\hyperref[sec:1121]{> > > NEXT SERMON > > >}
\newline
\newline


\newpage

\section{第61屆~港九培靈會~張子華~第七講}
\label{sec:1121}
\textbf{雅各的真我作為}
\newline
\newline
連結: \href{https://www.hkbibleconference.org/session-message/view/1121}{\texttt{https://www.hkbibleconference.org/session-message/view/1121}}
\newline
\newline
\hyperref[sec:1120]{< < < PREV SERMON < < <}
~
\hyperref[sec:toc]{[返目錄]}
~
\hyperref[sec:1122]{> > > NEXT SERMON > > >}
\newline
\newline
創世記 32:22-32:30
\newline
\begin{longtable}{cl}
\hline
\hline
章節 & 經文 (和合本修訂版)\\
\hline
32:22 & \begin{tabularx}{0.7\textwidth}{X} 他夜間起來，帶著兩個妻子，兩個婢女和十一個孩子，過了雅博渡口。 \end{tabularx} \\ \\ \relax
32:23 & \begin{tabularx}{0.7\textwidth}{X} 他帶著他們，送他們過河，他所有的一切也都過去， \end{tabularx} \\ \\ \relax
32:24 & \begin{tabularx}{0.7\textwidth}{X} 只剩下雅各一人。有一個人來和他摔跤，直到黎明。 \end{tabularx} \\ \\ \relax
32:25 & \begin{tabularx}{0.7\textwidth}{X} 那人見自己勝不過他，就摸了他的大腿窩一下。雅各的大腿窩就在和那人摔跤的時候扭了。 \end{tabularx} \\ \\ \relax
32:26 & \begin{tabularx}{0.7\textwidth}{X} 那人說：「天快亮了，讓我走吧！」雅各說：「你不給我祝福，我就不讓你走。」 \end{tabularx} \\ \\ \relax
32:27 & \begin{tabularx}{0.7\textwidth}{X} 那人說：「你叫甚麼名字？」他說：「雅各。」 \end{tabularx} \\ \\ \relax
32:28 & \begin{tabularx}{0.7\textwidth}{X} 那人說：「你的名字不要再叫雅各，要叫以色列，因為你與神和人較力，都得勝了。」 \end{tabularx} \\ \\ \relax
32:29 & \begin{tabularx}{0.7\textwidth}{X} 雅各問他說：「請告訴我你的名字。」那人說：「何必問我的名字呢？」於是他在那裡為雅各祝福。 \end{tabularx} \\ \\ \relax
32:30 & \begin{tabularx}{0.7\textwidth}{X} 雅各就給那地方起名叫毗努伊勒 ，說：「我面對面見了神，我的性命仍得保全。」 \end{tabularx} \\ \\
[1ex]
\hline
\hline
\end{longtable}


\newpage

\section{第61屆~港九培靈會~張子華~第八講}
\label{sec:1122}
\textbf{約瑟的升高}
\newline
\newline
連結: \href{https://www.hkbibleconference.org/session-message/view/1122}{\texttt{https://www.hkbibleconference.org/session-message/view/1122}}
\newline
\newline
\hyperref[sec:1121]{< < < PREV SERMON < < <}
~
\hyperref[sec:toc]{[返目錄]}
~
\hyperref[sec:1123]{> > > NEXT SERMON > > >}
\newline
\newline


\newpage

\section{第61屆~港九培靈會~張子華~第九講}
\label{sec:1123}
\textbf{約瑟的完美生命}
\newline
\newline
連結: \href{https://www.hkbibleconference.org/session-message/view/1123}{\texttt{https://www.hkbibleconference.org/session-message/view/1123}}
\newline
\newline
\hyperref[sec:1122]{< < < PREV SERMON < < <}
~
\hyperref[sec:toc]{[返目錄]}
~
\hyperref[sec:1141]{> > > NEXT SERMON > > >}
\newline
\newline


\newpage

\section{第61屆~港九培靈會~韓力生~第一講}
\label{sec:1141}
\textbf{請看新約教會 (哥林多前書概覽)}
\newline
\newline
連結: \href{https://www.hkbibleconference.org/session-message/view/1141}{\texttt{https://www.hkbibleconference.org/session-message/view/1141}}
\newline
\newline
\hyperref[sec:1123]{< < < PREV SERMON < < <}
~
\hyperref[sec:toc]{[返目錄]}
~
\hyperref[sec:1145]{> > > NEXT SERMON > > >}
\newline
\newline
哥林多前書 1:1-1:7
\newline
\begin{longtable}{cl}
\hline
\hline
章節 & 經文 (和合本修訂版)\\
\hline
1:1 & \begin{tabularx}{0.7\textwidth}{X} 奉神旨意，蒙召作基督耶穌使徒的保羅，同弟兄所提尼， \end{tabularx} \\ \\ \relax
1:2 & \begin{tabularx}{0.7\textwidth}{X} 寫信給在哥林多神的教會—就是在基督耶穌裡成聖、蒙召作聖徒的—以及所有在各處求告我主耶穌基督之名的人。基督是他們的主，也是我們的主。 \end{tabularx} \\ \\ \relax
1:3 & \begin{tabularx}{0.7\textwidth}{X} 願恩惠、平安從我們的父神並主耶穌基督歸給你們！ \end{tabularx} \\ \\ \relax
1:4 & \begin{tabularx}{0.7\textwidth}{X} 我常為你們感謝我的神，因神在基督耶穌裡所賜給你們的恩惠。 \end{tabularx} \\ \\ \relax
1:5 & \begin{tabularx}{0.7\textwidth}{X} 因為你們在他裡面凡事富足，具有各種口才、各樣知識， \end{tabularx} \\ \\ \relax
1:6 & \begin{tabularx}{0.7\textwidth}{X} 正如我為基督作的見證在你們心裡得以堅固， \end{tabularx} \\ \\ \relax
1:7 & \begin{tabularx}{0.7\textwidth}{X} 以致你們在恩賜上一無欠缺，切切等候我們主耶穌基督的顯現。 \end{tabularx} \\ \\
[1ex]
\hline
\hline
\end{longtable}


\newpage

\section{第61屆~港九培靈會~韓力生~第二講}
\label{sec:1145}
\textbf{人如何尋找神 (上)}
\newline
\newline
連結: \href{https://www.hkbibleconference.org/session-message/view/1145}{\texttt{https://www.hkbibleconference.org/session-message/view/1145}}
\newline
\newline
\hyperref[sec:1141]{< < < PREV SERMON < < <}
~
\hyperref[sec:toc]{[返目錄]}
~
\hyperref[sec:1146]{> > > NEXT SERMON > > >}
\newline
\newline
哥林多前書 1:10-1:31
\newline
\begin{longtable}{cl}
\hline
\hline
章節 & 經文 (和合本修訂版)\\
\hline
1:10 & \begin{tabularx}{0.7\textwidth}{X} 弟兄們，我藉我們主耶穌基督的名勸你們說話要一致。你們中間不可分裂，只要一心一意彼此團結。 \end{tabularx} \\ \\ \relax
1:11 & \begin{tabularx}{0.7\textwidth}{X} 我的弟兄們，革來氏家裡的人曾對我提起你們，說你們中間有紛爭。 \end{tabularx} \\ \\ \relax
1:12 & \begin{tabularx}{0.7\textwidth}{X} 我的意思是，你們各人說：「我是屬保羅的」；「我是屬亞波羅的」；「我是屬磯法的」；「我是屬基督的。」 \end{tabularx} \\ \\ \relax
1:13 & \begin{tabularx}{0.7\textwidth}{X} 基督是分裂的嗎？保羅為你們釘了十字架嗎？你們是奉保羅的名受了洗嗎？ \end{tabularx} \\ \\ \relax
1:14 & \begin{tabularx}{0.7\textwidth}{X} 我感謝神，除了基利司布和該猶以外，我沒有給你們中的任何一個人施洗， \end{tabularx} \\ \\ \relax
1:15 & \begin{tabularx}{0.7\textwidth}{X} 免得有人說你們是奉我的名受洗的。 \end{tabularx} \\ \\ \relax
1:16 & \begin{tabularx}{0.7\textwidth}{X} 我曾為司提法那家施過洗；此外我已記不清有沒有給別人施過洗。 \end{tabularx} \\ \\ \relax
1:17 & \begin{tabularx}{0.7\textwidth}{X} 因為基督差遣我不是為施洗，而是為傳福音；並不是用智慧的言論，免得基督的十字架落了空。 \end{tabularx} \\ \\ \relax
1:18 & \begin{tabularx}{0.7\textwidth}{X} 因為十字架的道理，在那滅亡的人是愚拙，在我們得救的人卻是神的大能。 \end{tabularx} \\ \\ \relax
1:19 & \begin{tabularx}{0.7\textwidth}{X} 就如經上所記：「我要摧毀智慧人的智慧，廢棄聰明人的聰明。」 \end{tabularx} \\ \\ \relax
1:20 & \begin{tabularx}{0.7\textwidth}{X} 智慧人在哪裡？文士在哪裡？這世上的辯士在哪裡？神豈不是已使這世上的智慧變成愚拙了嗎？ \end{tabularx} \\ \\ \relax
1:21 & \begin{tabularx}{0.7\textwidth}{X} 既然世人憑自己的智慧不認識神，神就本著自己的智慧樂意藉著人所傳愚拙的話拯救那些信的人。 \end{tabularx} \\ \\ \relax
1:22 & \begin{tabularx}{0.7\textwidth}{X} 猶太人要的是神蹟，希臘人求的是智慧， \end{tabularx} \\ \\ \relax
1:23 & \begin{tabularx}{0.7\textwidth}{X} 我們卻是傳被釘十字架的基督，這對猶太人是絆腳石，對外邦人是愚拙； \end{tabularx} \\ \\ \relax
1:24 & \begin{tabularx}{0.7\textwidth}{X} 但對那蒙召的，無論是猶太人、希臘人，基督總是神的大能，神的智慧。 \end{tabularx} \\ \\ \relax
1:25 & \begin{tabularx}{0.7\textwidth}{X} 因為，神的愚拙總比人智慧；神的軟弱總比人強壯。 \end{tabularx} \\ \\ \relax
1:26 & \begin{tabularx}{0.7\textwidth}{X} 弟兄們哪，想一想你們的蒙召，按著人的觀點，有智慧的不多，有能力的不多，有尊貴地位的也不多。 \end{tabularx} \\ \\ \relax
1:27 & \begin{tabularx}{0.7\textwidth}{X} 但是，神揀選了世上愚拙的，為了使有智慧的羞愧；又揀選了世上軟弱的，為了使強壯的羞愧。 \end{tabularx} \\ \\ \relax
1:28 & \begin{tabularx}{0.7\textwidth}{X} 神也揀選了世上卑賤的，被人厭惡的，以及那一無所有的，為要廢掉那樣樣都有的， \end{tabularx} \\ \\ \relax
1:29 & \begin{tabularx}{0.7\textwidth}{X} 使凡血肉之軀的，在神面前，一個也不能自誇。 \end{tabularx} \\ \\ \relax
1:30 & \begin{tabularx}{0.7\textwidth}{X} 但你們得以在基督耶穌裡是本乎神，他使基督成為我們的智慧，成為公義、聖潔、救贖。 \end{tabularx} \\ \\ \relax
1:31 & \begin{tabularx}{0.7\textwidth}{X} 如經上所記：「要誇耀的，該誇耀主。」 \end{tabularx} \\ \\
[1ex]
\hline
\hline
\end{longtable}


\newpage

\section{第61屆~港九培靈會~韓力生~第三講}
\label{sec:1146}
\textbf{人如何尋找神 (下)}
\newline
\newline
連結: \href{https://www.hkbibleconference.org/session-message/view/1146}{\texttt{https://www.hkbibleconference.org/session-message/view/1146}}
\newline
\newline
\hyperref[sec:1145]{< < < PREV SERMON < < <}
~
\hyperref[sec:toc]{[返目錄]}
~
\hyperref[sec:1147]{> > > NEXT SERMON > > >}
\newline
\newline
哥林多前書 2:1-2:16
\newline
\begin{longtable}{cl}
\hline
\hline
章節 & 經文 (和合本修訂版)\\
\hline
2:1 & \begin{tabularx}{0.7\textwidth}{X} 弟兄們，從前我到你們那裡去，並沒有用高言大智對你們宣講神的奧祕。 \end{tabularx} \\ \\ \relax
2:2 & \begin{tabularx}{0.7\textwidth}{X} 因為我曾定了主意，在你們中間不知道別的，只知道耶穌基督並他釘十字架。 \end{tabularx} \\ \\ \relax
2:3 & \begin{tabularx}{0.7\textwidth}{X} 我在你們那裡時，又軟弱，又懼怕，又戰戰兢兢。 \end{tabularx} \\ \\ \relax
2:4 & \begin{tabularx}{0.7\textwidth}{X} 我說的話、講的道不是用委婉智慧的言語，而是以聖靈的大能來證明， \end{tabularx} \\ \\ \relax
2:5 & \begin{tabularx}{0.7\textwidth}{X} 為要使你們的信不靠著人的智慧，而是靠著神的大能。 \end{tabularx} \\ \\ \relax
2:6 & \begin{tabularx}{0.7\textwidth}{X} 然而，在成熟的人中，我們也講智慧，但不是今世的智慧，也不是今世有權有位、將要滅亡的人的智慧。 \end{tabularx} \\ \\ \relax
2:7 & \begin{tabularx}{0.7\textwidth}{X} 我們講的是從前隱藏的、神奧祕的智慧，就是神在萬世以前預定使我們得榮耀的智慧； \end{tabularx} \\ \\ \relax
2:8 & \begin{tabularx}{0.7\textwidth}{X} 這智慧，今世有權有位的人沒有一個知道，若知道，他們就不會把榮耀的主釘在十字架上了。 \end{tabularx} \\ \\ \relax
2:9 & \begin{tabularx}{0.7\textwidth}{X} 如經上所記：「神為愛他的人所預備的是眼睛未曾看見，耳朵未曾聽見，人心也未曾想到的。」 \end{tabularx} \\ \\ \relax
2:10 & \begin{tabularx}{0.7\textwidth}{X} 只有神藉著聖靈把這事向我們顯明了；因為聖靈參透萬事，就是神深奧的事也參透了。 \end{tabularx} \\ \\ \relax
2:11 & \begin{tabularx}{0.7\textwidth}{X} 除了在人裡頭的靈，誰知道人的事？照樣，除了神的靈，也沒有人知道神的事。 \end{tabularx} \\ \\ \relax
2:12 & \begin{tabularx}{0.7\textwidth}{X} 我們所領受的並不是世上的靈，而是從神來的靈，為使我們知道神把恩賜賞給我們的事。 \end{tabularx} \\ \\ \relax
2:13 & \begin{tabularx}{0.7\textwidth}{X} 我們也講說這些事，不是用人的智慧所教的言語，而是用聖靈所教的言語，用屬靈的話解釋屬靈的事。 \end{tabularx} \\ \\ \relax
2:14 & \begin{tabularx}{0.7\textwidth}{X} 然而，屬血氣的人不接受神的靈的事，他反倒以這為愚拙，並且他不能了解，因為這些事惟有屬靈的人才能領悟。 \end{tabularx} \\ \\ \relax
2:15 & \begin{tabularx}{0.7\textwidth}{X} 屬靈的人能看透萬事，卻沒有一人能看透他。 \end{tabularx} \\ \\ \relax
2:16 & \begin{tabularx}{0.7\textwidth}{X} 「誰曾知道主的心？誰會教導他？」至於我們，我們有基督的心。 \end{tabularx} \\ \\
[1ex]
\hline
\hline
\end{longtable}


\newpage

\section{第61屆~港九培靈會~韓力生~第四講}
\label{sec:1147}
\textbf{建造的代價}
\newline
\newline
連結: \href{https://www.hkbibleconference.org/session-message/view/1147}{\texttt{https://www.hkbibleconference.org/session-message/view/1147}}
\newline
\newline
\hyperref[sec:1146]{< < < PREV SERMON < < <}
~
\hyperref[sec:toc]{[返目錄]}
~
\hyperref[sec:1149]{> > > NEXT SERMON > > >}
\newline
\newline
哥林多前書 3:1-3:23
\newline
\begin{longtable}{cl}
\hline
\hline
章節 & 經文 (和合本修訂版)\\
\hline
3:1 & \begin{tabularx}{0.7\textwidth}{X} 弟兄們，我從前對你們說話，還不能把你們當作屬靈的，只能把你們當作屬肉體的，你們在基督裡僅是嬰孩。 \end{tabularx} \\ \\ \relax
3:2 & \begin{tabularx}{0.7\textwidth}{X} 我用奶餵你們，沒有用飯餵你們，因為那時你們不能吃。就是如今還是不能， \end{tabularx} \\ \\ \relax
3:3 & \begin{tabularx}{0.7\textwidth}{X} 因為你們仍是屬肉體的。你們中間有嫉妒、紛爭，這豈不是屬乎肉體，照著世人的樣子生活嗎？ \end{tabularx} \\ \\ \relax
3:4 & \begin{tabularx}{0.7\textwidth}{X} 有人說：「我是屬保羅的」；有人說：「我是屬亞波羅的」；這樣你們豈不是和世人一樣嗎？ \end{tabularx} \\ \\ \relax
3:5 & \begin{tabularx}{0.7\textwidth}{X} 亞波羅算甚麼？保羅算甚麼？我們都是神的執事，藉著我們，你們信了；這不過是照著主給各人的恩賜去做罷了。 \end{tabularx} \\ \\ \relax
3:6 & \begin{tabularx}{0.7\textwidth}{X} 我栽種了，亞波羅澆灌了，惟有神使它生長。 \end{tabularx} \\ \\ \relax
3:7 & \begin{tabularx}{0.7\textwidth}{X} 可見，栽種的算不了甚麼，澆灌的也算不了甚麼；惟有神能使它生長。 \end{tabularx} \\ \\ \relax
3:8 & \begin{tabularx}{0.7\textwidth}{X} 栽種的和澆灌的都是一樣，但將來各人要照自己的勞苦得到自己的報酬。 \end{tabularx} \\ \\ \relax
3:9 & \begin{tabularx}{0.7\textwidth}{X} 因為我們是神的同工，而你們是神的田地、神的房屋。 \end{tabularx} \\ \\ \relax
3:10 & \begin{tabularx}{0.7\textwidth}{X} 我照神所給我的恩典，好像一個聰明的工頭，立好了根基，別人在上面建造；只是各人要謹慎怎樣在上面建造。 \end{tabularx} \\ \\ \relax
3:11 & \begin{tabularx}{0.7\textwidth}{X} 因為，那已經立好的根基就是耶穌基督，此外沒有人能立別的根基。 \end{tabularx} \\ \\ \relax
3:12 & \begin{tabularx}{0.7\textwidth}{X} 若有人用金銀、寶石，草木、禾秸，在這根基上建造， \end{tabularx} \\ \\ \relax
3:13 & \begin{tabularx}{0.7\textwidth}{X} 各人的工程必將顯露，因為那日子要將它顯明，有火把它暴露出來，這火要試煉各人的工程怎樣。 \end{tabularx} \\ \\ \relax
3:14 & \begin{tabularx}{0.7\textwidth}{X} 人在那根基上所建造的工程若能保得住，他將要得賞賜。 \end{tabularx} \\ \\ \relax
3:15 & \begin{tabularx}{0.7\textwidth}{X} 人的工程若被燒了，他將損失，雖然他自己將得救，卻要像從火裡經過一樣。 \end{tabularx} \\ \\ \relax
3:16 & \begin{tabularx}{0.7\textwidth}{X} 難道不知你們是神的殿，神的靈住在你們裡面嗎？ \end{tabularx} \\ \\ \relax
3:17 & \begin{tabularx}{0.7\textwidth}{X} 若有人毀壞神的殿，神一定要毀滅那人；因為神的殿是神聖的，這殿就是你們。 \end{tabularx} \\ \\ \relax
3:18 & \begin{tabularx}{0.7\textwidth}{X} 誰都不可自欺。你們中間若有人自以為在今世有智慧，倒不如變為愚拙，好成為有智慧的。 \end{tabularx} \\ \\ \relax
3:19 & \begin{tabularx}{0.7\textwidth}{X} 因為這世界的智慧在神看來是愚拙的。如經上記著：「主使有智慧的人中了自己的詭計；」 \end{tabularx} \\ \\ \relax
3:20 & \begin{tabularx}{0.7\textwidth}{X} 又說：「主知道智慧人的意念，因為它們是虛妄的。」 \end{tabularx} \\ \\ \relax
3:21 & \begin{tabularx}{0.7\textwidth}{X} 所以，無論誰都不可誇耀人；因為萬有都是你們的， \end{tabularx} \\ \\ \relax
3:22 & \begin{tabularx}{0.7\textwidth}{X} 或保羅，或亞波羅，或磯法，或世界，或生，或死，或現今的事，或將來的事，全是你們的， \end{tabularx} \\ \\ \relax
3:23 & \begin{tabularx}{0.7\textwidth}{X} 而你們是屬基督的，基督是屬神的。 \end{tabularx} \\ \\
[1ex]
\hline
\hline
\end{longtable}


\newpage

\section{第61屆~港九培靈會~韓力生~第五講}
\label{sec:1149}
\textbf{明確的判斷 — 動機}
\newline
\newline
連結: \href{https://www.hkbibleconference.org/session-message/view/1149}{\texttt{https://www.hkbibleconference.org/session-message/view/1149}}
\newline
\newline
\hyperref[sec:1147]{< < < PREV SERMON < < <}
~
\hyperref[sec:toc]{[返目錄]}
~
\hyperref[sec:1151]{> > > NEXT SERMON > > >}
\newline
\newline
哥林多前書 4:1-4:21
\newline
\begin{longtable}{cl}
\hline
\hline
章節 & 經文 (和合本修訂版)\\
\hline
4:1 & \begin{tabularx}{0.7\textwidth}{X} 人應該把我們看為基督的執事，為神的奧祕的管家。 \end{tabularx} \\ \\ \relax
4:2 & \begin{tabularx}{0.7\textwidth}{X} 所求於管家的，是要他忠心。 \end{tabularx} \\ \\ \relax
4:3 & \begin{tabularx}{0.7\textwidth}{X} 我被你們評斷，或被別人評斷，我都以為是極小的事；連我自己也不評斷自己。 \end{tabularx} \\ \\ \relax
4:4 & \begin{tabularx}{0.7\textwidth}{X} 雖然我不覺得自己有錯，卻也不能因此判為無罪；審斷我的是主。 \end{tabularx} \\ \\ \relax
4:5 & \begin{tabularx}{0.7\textwidth}{X} 所以，時候未到，在主來以前甚麼都不要評斷，他要照出暗中的隱情，揭發人的動機。那時，各人要從神那裡得著稱讚。 \end{tabularx} \\ \\ \relax
4:6 & \begin{tabularx}{0.7\textwidth}{X} 弟兄們，為你們的緣故，我拿這些事應用到我自己和亞波羅身上，讓你們從我們學到「不可過於聖經所記」這話的意思，免得你們自高自大，看重這個，看輕那個。 \end{tabularx} \\ \\ \relax
4:7 & \begin{tabularx}{0.7\textwidth}{X} 使你與人不同的是誰呢？你所有的有哪一個不是領受的呢？若是領受的，為何自誇，彷彿不是領受的呢？ \end{tabularx} \\ \\ \relax
4:8 & \begin{tabularx}{0.7\textwidth}{X} 你們已經飽足了，已經富足了，用不著我們，自己就作王了。我願意你們果真作王，讓我們也可以與你們一同作王！ \end{tabularx} \\ \\ \relax
4:9 & \begin{tabularx}{0.7\textwidth}{X} 我想，神把我們作使徒的明顯地列在末後，好像定死罪的囚犯，因為我們成了一臺戲，給世界、天使和眾人觀看。 \end{tabularx} \\ \\ \relax
4:10 & \begin{tabularx}{0.7\textwidth}{X} 我們為基督的緣故成為愚拙的；你們在基督裡倒是聰明的。我們軟弱，你們倒強壯；你們有榮耀，我們倒被藐視。 \end{tabularx} \\ \\ \relax
4:11 & \begin{tabularx}{0.7\textwidth}{X} 直到如今，我們還是又飢又渴，又赤身露體，又挨打，又到處漂泊， \end{tabularx} \\ \\ \relax
4:12 & \begin{tabularx}{0.7\textwidth}{X} 並且勞碌，親手做工；被人咒罵，我們就祝福；被人迫害，我們就忍受； \end{tabularx} \\ \\ \relax
4:13 & \begin{tabularx}{0.7\textwidth}{X} 被人毀謗，我們就勸導。直到如今，人還把我們看作世上的污穢，萬物中的渣滓。 \end{tabularx} \\ \\ \relax
4:14 & \begin{tabularx}{0.7\textwidth}{X} 我寫這些話，不是要使你們羞愧，而是要警戒你們，好像我所愛的兒女一樣。 \end{tabularx} \\ \\ \relax
4:15 & \begin{tabularx}{0.7\textwidth}{X} 雖然你們在基督裡有無數的導師，卻沒有許多父親，因我是在基督耶穌裡用福音生了你們。 \end{tabularx} \\ \\ \relax
4:16 & \begin{tabularx}{0.7\textwidth}{X} 所以，我求你們要效法我。 \end{tabularx} \\ \\ \relax
4:17 & \begin{tabularx}{0.7\textwidth}{X} 因此，我已差提摩太到你們那裡去。他在主裡面是我親愛和忠心的兒子；他要提醒你們，我在基督耶穌裡怎樣行事為人，在各處各教會中怎樣教導人。 \end{tabularx} \\ \\ \relax
4:18 & \begin{tabularx}{0.7\textwidth}{X} 有些人以為我不到你們那裡去而自高自大。 \end{tabularx} \\ \\ \relax
4:19 & \begin{tabularx}{0.7\textwidth}{X} 但是，主若准許，我會很快到你們那裡去；我所要知道的，不是那些自高自大者的言語，而是他們的權能。 \end{tabularx} \\ \\ \relax
4:20 & \begin{tabularx}{0.7\textwidth}{X} 因為神的國不在乎言語，而在乎權能。 \end{tabularx} \\ \\ \relax
4:21 & \begin{tabularx}{0.7\textwidth}{X} 你們願意怎麼樣呢？要我帶著棍子到你們那裡去呢，還是帶著慈愛溫柔的心呢？ \end{tabularx} \\ \\
[1ex]
\hline
\hline
\end{longtable}


\newpage

\section{第61屆~港九培靈會~韓力生~第六講}
\label{sec:1151}
\textbf{明確的判斷 — 行為}
\newline
\newline
連結: \href{https://www.hkbibleconference.org/session-message/view/1151}{\texttt{https://www.hkbibleconference.org/session-message/view/1151}}
\newline
\newline
\hyperref[sec:1149]{< < < PREV SERMON < < <}
~
\hyperref[sec:toc]{[返目錄]}
~
\hyperref[sec:1153]{> > > NEXT SERMON > > >}
\newline
\newline
哥林多前書 5:1-5:13
\newline
\begin{longtable}{cl}
\hline
\hline
章節 & 經文 (和合本修訂版)\\
\hline
5:1 & \begin{tabularx}{0.7\textwidth}{X} 我確實聽說在你們中間有淫亂的事；這種淫亂連外邦人中也沒有，就是有人和他的繼母同居。 \end{tabularx} \\ \\ \relax
5:2 & \begin{tabularx}{0.7\textwidth}{X} 你們還自高自大！你們不是該覺得痛心，把做這事的人從你們中間趕出去嗎？ \end{tabularx} \\ \\ \relax
5:3 & \begin{tabularx}{0.7\textwidth}{X} 我人雖然不在你們那裡，心卻在你們那裡，好像親自與你們同在。我奉我們主耶穌的名，已經判斷了做這事的人。 \end{tabularx} \\ \\ \relax
5:4 & \begin{tabularx}{0.7\textwidth}{X} 你們聚會的時候，我的心和你們同在。你們藉著我們主耶穌的權能， \end{tabularx} \\ \\ \relax
5:5 & \begin{tabularx}{0.7\textwidth}{X} 要把這樣的人交給撒但，使他的肉體敗壞，好讓他的靈魂在主的日子可以得救。 \end{tabularx} \\ \\ \relax
5:6 & \begin{tabularx}{0.7\textwidth}{X} 你們這樣自誇是不好的。你們不知道一點麵酵能使全團發起來嗎？ \end{tabularx} \\ \\ \relax
5:7 & \begin{tabularx}{0.7\textwidth}{X} 既然你們是無酵的麵，要把舊酵除淨，好使你們成為新團；因為我們逾越節的羔羊—基督已經被殺獻為祭牲了。 \end{tabularx} \\ \\ \relax
5:8 & \begin{tabularx}{0.7\textwidth}{X} 所以，我們來守這節，不可用舊酵，就是不可用惡毒、邪惡的酵，只用純潔真實的無酵餅。 \end{tabularx} \\ \\ \relax
5:9 & \begin{tabularx}{0.7\textwidth}{X} 我先前寫信告訴過你們，不可與淫亂的人交往。 \end{tabularx} \\ \\ \relax
5:10 & \begin{tabularx}{0.7\textwidth}{X} 此話不是泛指這世上所有行淫亂的，或貪婪的，勒索的，或拜偶像的；若是這樣，你們非離開這世界不可。 \end{tabularx} \\ \\ \relax
5:11 & \begin{tabularx}{0.7\textwidth}{X} 但現在，我寫信告訴你們，若有稱為弟兄的人卻仍犯淫亂，或貪婪，或拜偶像，或辱罵，或醉酒，或勒索，這樣的人不可跟他交往，就是跟他吃飯都不可以。 \end{tabularx} \\ \\ \relax
5:12 & \begin{tabularx}{0.7\textwidth}{X} 因為審判教外的人與我何干？教內的人豈不是你們要審判嗎？ \end{tabularx} \\ \\ \relax
5:13 & \begin{tabularx}{0.7\textwidth}{X} 至於外人有神審判他們。如經上說：「要從你們中間把那邪惡的人趕出去。」 \end{tabularx} \\ \\
[1ex]
\hline
\hline
\end{longtable}


\newpage

\section{第61屆~港九培靈會~韓力生~第七講}
\label{sec:1153}
\textbf{明確的判斷 — 爭執}
\newline
\newline
連結: \href{https://www.hkbibleconference.org/session-message/view/1153}{\texttt{https://www.hkbibleconference.org/session-message/view/1153}}
\newline
\newline
\hyperref[sec:1151]{< < < PREV SERMON < < <}
~
\hyperref[sec:toc]{[返目錄]}
~
\hyperref[sec:1155]{> > > NEXT SERMON > > >}
\newline
\newline
哥林多前書 6:1-6:20
\newline
\begin{longtable}{cl}
\hline
\hline
章節 & 經文 (和合本修訂版)\\
\hline
6:1 & \begin{tabularx}{0.7\textwidth}{X} 你們中間有彼此爭吵的事，怎敢告到不義的人面前，而不告到聖徒面前呢？ \end{tabularx} \\ \\ \relax
6:2 & \begin{tabularx}{0.7\textwidth}{X} 你們豈不知聖徒要審判世界嗎？若世界要受你們的審判，難道你們不配審判這最小的事嗎？ \end{tabularx} \\ \\ \relax
6:3 & \begin{tabularx}{0.7\textwidth}{X} 你們豈不知我們要審判天使嗎？何況今生的事呢！ \end{tabularx} \\ \\ \relax
6:4 & \begin{tabularx}{0.7\textwidth}{X} 既是這樣，你們若有今生當審判的事，會讓教會所輕看的人來審判嗎？ \end{tabularx} \\ \\ \relax
6:5 & \begin{tabularx}{0.7\textwidth}{X} 我說這話是要使你們慚愧。難道你們中間沒有一個有智慧的人能審斷弟兄中的事嗎？ \end{tabularx} \\ \\ \relax
6:6 & \begin{tabularx}{0.7\textwidth}{X} 你們竟然有弟兄去告弟兄，而且告到不信主的人面前。 \end{tabularx} \\ \\ \relax
6:7 & \begin{tabularx}{0.7\textwidth}{X} 你們彼此告狀，這已經是你們的大錯了。為甚麼不情願受冤屈呢？為甚麼不情願吃虧呢？ \end{tabularx} \\ \\ \relax
6:8 & \begin{tabularx}{0.7\textwidth}{X} 你們反倒去冤枉人，虧負人，況且所冤枉所虧負的就是弟兄。 \end{tabularx} \\ \\ \relax
6:9 & \begin{tabularx}{0.7\textwidth}{X} 你們豈不知不義的人不能承受神的國嗎？不要自欺！無論是淫亂的、拜偶像的、姦淫的、作娼妓的，親男色的、 \end{tabularx} \\ \\ \relax
6:10 & \begin{tabularx}{0.7\textwidth}{X} 偷竊的、貪婪的、醉酒的、辱罵的、勒索的，都不能承受神的國。 \end{tabularx} \\ \\ \relax
6:11 & \begin{tabularx}{0.7\textwidth}{X} 從前你們中間也有人是這樣；但現在你們奉主耶穌基督的名，並藉著我們神的靈，已經洗淨，已經成聖，已經稱義了。 \end{tabularx} \\ \\ \relax
6:12 & \begin{tabularx}{0.7\textwidth}{X} 「凡事我都可行」，但不是凡事都有益處。「凡事我都可行」，但無論哪一件，我都不受它的轄制。 \end{tabularx} \\ \\ \relax
6:13 & \begin{tabularx}{0.7\textwidth}{X} 「食物是為肚腹，肚腹是為食物」；但神要使這兩樣都毀壞。身體不是為淫亂，而是為主；主也是為身體。 \end{tabularx} \\ \\ \relax
6:14 & \begin{tabularx}{0.7\textwidth}{X} 神已經使主復活，也要用他自己的能力使我們復活。 \end{tabularx} \\ \\ \relax
6:15 & \begin{tabularx}{0.7\textwidth}{X} 你們豈不知道你們的身體是基督的肢體嗎？我可以把基督的肢體作為娼妓的肢體嗎？絕對不可！ \end{tabularx} \\ \\ \relax
6:16 & \begin{tabularx}{0.7\textwidth}{X} 你們豈不知道與娼妓苟合的，就是與她成為一體嗎？因為主說：「二人要成為一體。」 \end{tabularx} \\ \\ \relax
6:17 & \begin{tabularx}{0.7\textwidth}{X} 但與主聯合的，就是與主成為一靈。 \end{tabularx} \\ \\ \relax
6:18 & \begin{tabularx}{0.7\textwidth}{X} 你們要遠避淫行。人所犯的，無論甚麼罪，都在身體以外；惟有行淫的，是得罪自己的身體。 \end{tabularx} \\ \\ \relax
6:19 & \begin{tabularx}{0.7\textwidth}{X} 你們豈不知道你們的身體是聖靈的殿嗎？這聖靈是從神而來，住在你們裡面的。而且你們不是屬自己的人， \end{tabularx} \\ \\ \relax
6:20 & \begin{tabularx}{0.7\textwidth}{X} 因為你們是重價買來的。所以，要在你們的身體上榮耀神。 \end{tabularx} \\ \\
[1ex]
\hline
\hline
\end{longtable}


\newpage

\section{第61屆~港九培靈會~韓力生~第八講}
\label{sec:1155}
\textbf{獨身或結婚 — 神的呼召}
\newline
\newline
連結: \href{https://www.hkbibleconference.org/session-message/view/1155}{\texttt{https://www.hkbibleconference.org/session-message/view/1155}}
\newline
\newline
\hyperref[sec:1153]{< < < PREV SERMON < < <}
~
\hyperref[sec:toc]{[返目錄]}
~
\hyperref[sec:1159]{> > > NEXT SERMON > > >}
\newline
\newline
哥林多前書 7:1-7:40
\newline
\begin{longtable}{cl}
\hline
\hline
章節 & 經文 (和合本修訂版)\\
\hline
7:1 & \begin{tabularx}{0.7\textwidth}{X} 關於你們信上所提的事，男人不親近女人倒好。 \end{tabularx} \\ \\ \relax
7:2 & \begin{tabularx}{0.7\textwidth}{X} 但為了避免淫亂的事，男人當各有自己的妻子，女人也當各有自己的丈夫。 \end{tabularx} \\ \\ \relax
7:3 & \begin{tabularx}{0.7\textwidth}{X} 丈夫對妻子要盡本分；妻子對丈夫也要如此。 \end{tabularx} \\ \\ \relax
7:4 & \begin{tabularx}{0.7\textwidth}{X} 妻子對自己的身體沒有主張的權柄，權柄在丈夫；丈夫對自己的身體也沒有主張的權柄，權柄在妻子。 \end{tabularx} \\ \\ \relax
7:5 & \begin{tabularx}{0.7\textwidth}{X} 夫妻不可忽略對方的需求，除非為了要專心禱告，在兩相情願下暫時分房；以後仍要同房，免得撒但趁著你們情不自禁而引誘你們。 \end{tabularx} \\ \\ \relax
7:6 & \begin{tabularx}{0.7\textwidth}{X} 我說這話是出於容忍，不是命令。 \end{tabularx} \\ \\ \relax
7:7 & \begin{tabularx}{0.7\textwidth}{X} 我願眾人像我一樣；但是各人都有來自神的恩賜，一個是這樣，一個是那樣。 \end{tabularx} \\ \\ \relax
7:8 & \begin{tabularx}{0.7\textwidth}{X} 我對沒有嫁娶的和寡婦說，他們若能維持獨身像我一樣就好。 \end{tabularx} \\ \\ \relax
7:9 & \begin{tabularx}{0.7\textwidth}{X} 但他們若不能自制，就應該嫁娶，與其慾火攻心，倒不如結婚為妙。 \end{tabularx} \\ \\ \relax
7:10 & \begin{tabularx}{0.7\textwidth}{X} 至於那已經嫁娶的，我吩咐他們—其實不是我，而是主吩咐的：妻子不可離開丈夫， \end{tabularx} \\ \\ \relax
7:11 & \begin{tabularx}{0.7\textwidth}{X} 若是離開了，不可再嫁，不然要跟丈夫復和；丈夫也不可離棄妻子。 \end{tabularx} \\ \\ \relax
7:12 & \begin{tabularx}{0.7\textwidth}{X} 我對其餘的人說—是我，不是主說—倘若某弟兄有不信的妻子，妻子也情願和他一起生活，他就不可離棄妻子。 \end{tabularx} \\ \\ \relax
7:13 & \begin{tabularx}{0.7\textwidth}{X} 妻子有不信的丈夫，丈夫也情願和她一起生活，她就不可離棄丈夫。 \end{tabularx} \\ \\ \relax
7:14 & \begin{tabularx}{0.7\textwidth}{X} 因為不信的丈夫會因著妻子成了聖潔；不信的妻子也會因著丈夫成了聖潔。不然，你們的兒女就不潔淨了，但現在他們是聖潔的。 \end{tabularx} \\ \\ \relax
7:15 & \begin{tabularx}{0.7\textwidth}{X} 倘若那不信的人要離開，就由他離開吧！無論是弟兄是姊妹，遇著這樣的事都不必拘束。神召你們原是要你們和睦。 \end{tabularx} \\ \\ \relax
7:16 & \begin{tabularx}{0.7\textwidth}{X} 你這作妻子的怎麼知道不能救你的丈夫呢？你這作丈夫的怎麼知道不能救你的妻子呢？ \end{tabularx} \\ \\ \relax
7:17 & \begin{tabularx}{0.7\textwidth}{X} 無論如何，要照主所分給各人的恩賜和神所召各人的情況生活。我在各教會裡都是這樣規定的。 \end{tabularx} \\ \\ \relax
7:18 & \begin{tabularx}{0.7\textwidth}{X} 有人受割禮後才蒙召，他就不必除去割禮的記號。有人未受割禮前蒙召，他就不必受割禮。 \end{tabularx} \\ \\ \relax
7:19 & \begin{tabularx}{0.7\textwidth}{X} 受割禮算不了甚麼，不受割禮也算不了甚麼，只要謹守神的誡命就是了。 \end{tabularx} \\ \\ \relax
7:20 & \begin{tabularx}{0.7\textwidth}{X} 各人蒙召的時候是甚麼身份，要守住這身份。 \end{tabularx} \\ \\ \relax
7:21 & \begin{tabularx}{0.7\textwidth}{X} 你是作奴隸時蒙召的嗎？不要介意；若能獲得自由，就爭取自由更好。 \end{tabularx} \\ \\ \relax
7:22 & \begin{tabularx}{0.7\textwidth}{X} 因為，蒙主呼召的奴僕是主所釋放的人；蒙主呼召的自由之人是基督的奴僕。 \end{tabularx} \\ \\ \relax
7:23 & \begin{tabularx}{0.7\textwidth}{X} 你們是重價買來的；不要作人的奴僕。 \end{tabularx} \\ \\ \relax
7:24 & \begin{tabularx}{0.7\textwidth}{X} 弟兄們，你們各人蒙召的時候是甚麼身份，要在神面前守住這身份。 \end{tabularx} \\ \\ \relax
7:25 & \begin{tabularx}{0.7\textwidth}{X} 關於未婚女子，我沒有主的命令，但我既蒙主憐憫、作為一個可信靠的人，把自己的意見告訴你們。 \end{tabularx} \\ \\ \relax
7:26 & \begin{tabularx}{0.7\textwidth}{X} 因現今的艱難，據我看來，人不如安於現狀。 \end{tabularx} \\ \\ \relax
7:27 & \begin{tabularx}{0.7\textwidth}{X} 你已經有了妻子，就不要求擺脫；你還沒有妻子，就不要想娶妻。 \end{tabularx} \\ \\ \relax
7:28 & \begin{tabularx}{0.7\textwidth}{X} 你若娶妻，並不是犯罪；未婚女子若出嫁，也不是犯罪。然而，這等人會遭受肉身上的苦難，我寧願你們免受這苦難。 \end{tabularx} \\ \\ \relax
7:29 & \begin{tabularx}{0.7\textwidth}{X} 弟兄們，我是說：時候不多了。從此以後，那有妻子的，要像沒有一樣； \end{tabularx} \\ \\ \relax
7:30 & \begin{tabularx}{0.7\textwidth}{X} 哀哭的，不像在哀哭；快樂的，不像在快樂；購買的，像一無所得； \end{tabularx} \\ \\ \relax
7:31 & \begin{tabularx}{0.7\textwidth}{X} 享受這世界的，不像在享受這世界；因為這世界的局面將要過去了。 \end{tabularx} \\ \\ \relax
7:32 & \begin{tabularx}{0.7\textwidth}{X} 我願你們一無掛慮。沒有結婚的是為主的事掛慮，想怎樣令主喜悅； \end{tabularx} \\ \\ \relax
7:33 & \begin{tabularx}{0.7\textwidth}{X} 結了婚的是為世上的事掛慮，想怎樣讓妻子喜悅， \end{tabularx} \\ \\ \relax
7:34 & \begin{tabularx}{0.7\textwidth}{X} 於是，他就分心了。沒有結婚的和未婚的女子是為主的事掛慮，為要身體和心靈都聖潔；已經出嫁的是為世上的事掛慮，想怎樣讓丈夫喜悅。 \end{tabularx} \\ \\ \relax
7:35 & \begin{tabularx}{0.7\textwidth}{X} 我說這話是為你們的益處，不是要限制你們，而是要你們做合宜的事，得以不分心地對主忠誠。 \end{tabularx} \\ \\ \relax
7:36 & \begin{tabularx}{0.7\textwidth}{X} 若有人認為自己待他的女兒不合宜，女兒也過了適婚年齡，他可以隨意處理，不算有罪，讓兩人結婚就是了。 \end{tabularx} \\ \\ \relax
7:37 & \begin{tabularx}{0.7\textwidth}{X} 倘若有人心裡堅定，沒有不得已的事，並且由得自己作主，心裡又決定了不讓女兒結婚，這樣做也好。 \end{tabularx} \\ \\ \relax
7:38 & \begin{tabularx}{0.7\textwidth}{X} 這樣看來，讓自己的女兒結婚固然是好，不讓她結婚更好。 \end{tabularx} \\ \\ \relax
7:39 & \begin{tabularx}{0.7\textwidth}{X} 丈夫活著的時候，妻子是受約束的；丈夫若長眠了，妻子就自由了，可以隨意再嫁，只是要嫁給主裡面的人。 \end{tabularx} \\ \\ \relax
7:40 & \begin{tabularx}{0.7\textwidth}{X} 然而，按我的意見，她若能守節就更有福氣。我想我自己也有神的靈的感動。 \end{tabularx} \\ \\
[1ex]
\hline
\hline
\end{longtable}


\newpage

\section{第61屆~港九培靈會~韓力生~第九講}
\label{sec:1159}
\textbf{恩賜的分配及規則}
\newline
\newline
連結: \href{https://www.hkbibleconference.org/session-message/view/1159}{\texttt{https://www.hkbibleconference.org/session-message/view/1159}}
\newline
\newline
\hyperref[sec:1155]{< < < PREV SERMON < < <}
~
\hyperref[sec:toc]{[返目錄]}
~
\hyperref[sec:1162]{> > > NEXT SERMON > > >}
\newline
\newline
哥林多前書 12:1-12:31
\newline
\begin{longtable}{cl}
\hline
\hline
章節 & 經文 (和合本修訂版)\\
\hline
12:1 & \begin{tabularx}{0.7\textwidth}{X} 弟兄們，關於屬靈的恩賜，我不願意你們不明白。 \end{tabularx} \\ \\ \relax
12:2 & \begin{tabularx}{0.7\textwidth}{X} 你們知道，你們作外邦人的時候，隨事被引誘，受了迷惑去拜不會出聲的偶像。 \end{tabularx} \\ \\ \relax
12:3 & \begin{tabularx}{0.7\textwidth}{X} 所以，我要你們知道，被神的靈感動的，沒有人會說「耶穌該受詛咒」；若不是被聖靈感動的，也沒有人能說「耶穌是主」。 \end{tabularx} \\ \\ \relax
12:4 & \begin{tabularx}{0.7\textwidth}{X} 恩賜有許多種，卻是同一位聖靈所賜。 \end{tabularx} \\ \\ \relax
12:5 & \begin{tabularx}{0.7\textwidth}{X} 事奉有許多種，卻是事奉同一位主。 \end{tabularx} \\ \\ \relax
12:6 & \begin{tabularx}{0.7\textwidth}{X} 工作有許多種，卻是同一位神在萬人中運行萬事。 \end{tabularx} \\ \\ \relax
12:7 & \begin{tabularx}{0.7\textwidth}{X} 聖靈彰顯在各人身上，是要使人得益處。 \end{tabularx} \\ \\ \relax
12:8 & \begin{tabularx}{0.7\textwidth}{X} 有人藉著聖靈領受智慧的言語；有人也靠著同一位聖靈領受知識的言語； \end{tabularx} \\ \\ \relax
12:9 & \begin{tabularx}{0.7\textwidth}{X} 又有人由同一位聖靈領受信心；還有人由同一位聖靈領受醫病的恩賜； \end{tabularx} \\ \\ \relax
12:10 & \begin{tabularx}{0.7\textwidth}{X} 又有人能行異能，又有人能作先知，又有人能辨別諸靈，又有人能說方言，又有人能翻方言。 \end{tabularx} \\ \\ \relax
12:11 & \begin{tabularx}{0.7\textwidth}{X} 這一切都是由惟一的、同一位聖靈所運行，隨著自己的旨意分給各人的。 \end{tabularx} \\ \\ \relax
12:12 & \begin{tabularx}{0.7\textwidth}{X} 就如身體是一個，卻有許多肢體，身體的肢體雖多，仍是一個身體；基督也是這樣。 \end{tabularx} \\ \\ \relax
12:13 & \begin{tabularx}{0.7\textwidth}{X} 我們無論是猶太人是希臘人，是為奴的是自主的，都從一位聖靈受洗成了一個身體，並且共享這位聖靈。 \end{tabularx} \\ \\ \relax
12:14 & \begin{tabularx}{0.7\textwidth}{X} 身體原不只是一個肢體，而是許多肢體。 \end{tabularx} \\ \\ \relax
12:15 & \begin{tabularx}{0.7\textwidth}{X} 假如腳說：「我不是手，所以不屬於身體」，它不能因此就不屬於身體。 \end{tabularx} \\ \\ \relax
12:16 & \begin{tabularx}{0.7\textwidth}{X} 假如耳朵說：「我不是眼睛，所以不屬於身體」，它也不能因此就不屬於身體。 \end{tabularx} \\ \\ \relax
12:17 & \begin{tabularx}{0.7\textwidth}{X} 假如全身是眼睛，聽覺在哪裡呢？假如全身是耳朵，嗅覺在哪裡呢？ \end{tabularx} \\ \\ \relax
12:18 & \begin{tabularx}{0.7\textwidth}{X} 但現在神隨自己的意思把肢體一一安置在身體上了。 \end{tabularx} \\ \\ \relax
12:19 & \begin{tabularx}{0.7\textwidth}{X} 假如全都是一個肢體，身體在哪裡呢？ \end{tabularx} \\ \\ \relax
12:20 & \begin{tabularx}{0.7\textwidth}{X} 但現在肢體雖多，身體還是一個。 \end{tabularx} \\ \\ \relax
12:21 & \begin{tabularx}{0.7\textwidth}{X} 眼睛不能對手說：「我用不著你。」頭也不能對腳說：「我用不著你。」 \end{tabularx} \\ \\ \relax
12:22 & \begin{tabularx}{0.7\textwidth}{X} 不但如此，身上的肢體，人以為軟弱的，更是不可缺少的； \end{tabularx} \\ \\ \relax
12:23 & \begin{tabularx}{0.7\textwidth}{X} 身上的肢體，我們認為不體面的，越發給它加上體面；我們不雅觀的，越發裝飾得雅觀。 \end{tabularx} \\ \\ \relax
12:24 & \begin{tabularx}{0.7\textwidth}{X} 我們雅觀的肢體自然用不著裝飾；但神配搭這身子，把加倍的體面給那有缺欠的肢體， \end{tabularx} \\ \\ \relax
12:25 & \begin{tabularx}{0.7\textwidth}{X} 免得身體不協調，總要肢體彼此照顧。 \end{tabularx} \\ \\ \relax
12:26 & \begin{tabularx}{0.7\textwidth}{X} 假如一個肢體受苦，所有的肢體就一同受苦；假如一個肢體得光榮，所有的肢體就一同快樂。 \end{tabularx} \\ \\ \relax
12:27 & \begin{tabularx}{0.7\textwidth}{X} 你們是基督的身體，並且各自都是肢體。 \end{tabularx} \\ \\ \relax
12:28 & \begin{tabularx}{0.7\textwidth}{X} 神在教會所設立的：第一是使徒；第二是先知；第三是教師；其次是行異能的；再次是醫病的恩賜，幫助人的，治理事的，說方言的。 \end{tabularx} \\ \\ \relax
12:29 & \begin{tabularx}{0.7\textwidth}{X} 難道個個都是使徒嗎？難道個個都是先知嗎？難道個個都是教師嗎？難道個個都是行異能的嗎？ \end{tabularx} \\ \\ \relax
12:30 & \begin{tabularx}{0.7\textwidth}{X} 難道個個都是有醫病的恩賜嗎？難道個個都是說方言的嗎？難道個個都是翻方言的嗎？ \end{tabularx} \\ \\ \relax
12:31 & \begin{tabularx}{0.7\textwidth}{X} 你們要追求那更大的恩賜。我現今把最妙的道指示你們。 \end{tabularx} \\ \\
[1ex]
\hline
\hline
\end{longtable}


\newpage

\section{第61屆~港九培靈會~韓力生~第十講}
\label{sec:1162}
\textbf{神聖的潤滑劑}
\newline
\newline
連結: \href{https://www.hkbibleconference.org/session-message/view/1162}{\texttt{https://www.hkbibleconference.org/session-message/view/1162}}
\newline
\newline
\hyperref[sec:1159]{< < < PREV SERMON < < <}
~
\hyperref[sec:toc]{[返目錄]}
~
\hyperref[sec:1085]{> > > NEXT SERMON > > >}
\newline
\newline
哥林多前書 13:1-13:13
\newline
\begin{longtable}{cl}
\hline
\hline
章節 & 經文 (和合本修訂版)\\
\hline
13:1 & \begin{tabularx}{0.7\textwidth}{X} 我若能說人間的方言，甚至天使的語言，卻沒有愛，我就成為鳴的鑼、響的鈸一般。 \end{tabularx} \\ \\ \relax
13:2 & \begin{tabularx}{0.7\textwidth}{X} 我若有先知講道的能力，也明白各樣的奧祕，各樣的知識，而且有齊備的信心，使我能夠移山，卻沒有愛，我就算不了甚麼。 \end{tabularx} \\ \\ \relax
13:3 & \begin{tabularx}{0.7\textwidth}{X} 我若將所有的財產救濟窮人，又犧牲自己的身體讓人誇讚，卻沒有愛，仍然對我無益。 \end{tabularx} \\ \\ \relax
13:4 & \begin{tabularx}{0.7\textwidth}{X} 愛是恆久忍耐；又有恩慈；愛是不嫉妒；愛是不自誇，不張狂， \end{tabularx} \\ \\ \relax
13:5 & \begin{tabularx}{0.7\textwidth}{X} 不做害羞的事，不求自己的益處，不輕易發怒，不計算人的惡， \end{tabularx} \\ \\ \relax
13:6 & \begin{tabularx}{0.7\textwidth}{X} 不喜歡不義，只喜歡真理； \end{tabularx} \\ \\ \relax
13:7 & \begin{tabularx}{0.7\textwidth}{X} 凡事包容，凡事相信，凡事盼望，凡事忍耐。 \end{tabularx} \\ \\ \relax
13:8 & \begin{tabularx}{0.7\textwidth}{X} 愛是永不止息。先知講道之能終必歸於無有；說方言之能終必停止；知識也終必歸於無有。 \end{tabularx} \\ \\ \relax
13:9 & \begin{tabularx}{0.7\textwidth}{X} 我們現在所知道的有限，先知所講的也有限， \end{tabularx} \\ \\ \relax
13:10 & \begin{tabularx}{0.7\textwidth}{X} 等那完全的來到，這有限的必消逝。 \end{tabularx} \\ \\ \relax
13:11 & \begin{tabularx}{0.7\textwidth}{X} 我作孩子的時候，說話像孩子，心思像孩子，意念像孩子；既長大成人，就把孩子的事丟棄了。 \end{tabularx} \\ \\ \relax
13:12 & \begin{tabularx}{0.7\textwidth}{X} 我們現在是對著鏡子觀看，模糊不清；到那時，就要面對面了。我如今所認識的有限，到那時就全認識，如同主認識我一樣。 \end{tabularx} \\ \\ \relax
13:13 & \begin{tabularx}{0.7\textwidth}{X} 如今常存的有信，有望，有愛這三樣，其中最大的是愛。 \end{tabularx} \\ \\
[1ex]
\hline
\hline
\end{longtable}


\newpage

\section{第62屆~港九培靈會~劉少康~第一講}
\label{sec:1085}
\textbf{重建的時機 — 尼希米的呼籲}
\newline
\newline
連結: \href{https://www.hkbibleconference.org/session-message/view/1085}{\texttt{https://www.hkbibleconference.org/session-message/view/1085}}
\newline
\newline
\hyperref[sec:1162]{< < < PREV SERMON < < <}
~
\hyperref[sec:toc]{[返目錄]}
~
\hyperref[sec:1086]{> > > NEXT SERMON > > >}
\newline
\newline


\newpage

\section{第62屆~港九培靈會~劉少康~第二講}
\label{sec:1086}
\textbf{見證的重建 — 尼希米的領導}
\newline
\newline
連結: \href{https://www.hkbibleconference.org/session-message/view/1086}{\texttt{https://www.hkbibleconference.org/session-message/view/1086}}
\newline
\newline
\hyperref[sec:1085]{< < < PREV SERMON < < <}
~
\hyperref[sec:toc]{[返目錄]}
~
\hyperref[sec:1087]{> > > NEXT SERMON > > >}
\newline
\newline


\newpage

\section{第62屆~港九培靈會~劉少康~第三講}
\label{sec:1087}
\textbf{敬拜的重建 — 以斯拉的行動}
\newline
\newline
連結: \href{https://www.hkbibleconference.org/session-message/view/1087}{\texttt{https://www.hkbibleconference.org/session-message/view/1087}}
\newline
\newline
\hyperref[sec:1086]{< < < PREV SERMON < < <}
~
\hyperref[sec:toc]{[返目錄]}
~
\hyperref[sec:1088]{> > > NEXT SERMON > > >}
\newline
\newline


\newpage

\section{第62屆~港九培靈會~劉少康~第四講}
\label{sec:1088}
\textbf{心靈的重建 — 以斯拉的心志}
\newline
\newline
連結: \href{https://www.hkbibleconference.org/session-message/view/1088}{\texttt{https://www.hkbibleconference.org/session-message/view/1088}}
\newline
\newline
\hyperref[sec:1087]{< < < PREV SERMON < < <}
~
\hyperref[sec:toc]{[返目錄]}
~
\hyperref[sec:1089]{> > > NEXT SERMON > > >}
\newline
\newline


\newpage

\section{第62屆~港九培靈會~劉少康~第五講}
\label{sec:1089}
\textbf{生活的重建 — 以斯拉的責備}
\newline
\newline
連結: \href{https://www.hkbibleconference.org/session-message/view/1089}{\texttt{https://www.hkbibleconference.org/session-message/view/1089}}
\newline
\newline
\hyperref[sec:1088]{< < < PREV SERMON < < <}
~
\hyperref[sec:toc]{[返目錄]}
~
\hyperref[sec:1090]{> > > NEXT SERMON > > >}
\newline
\newline


\newpage

\section{第62屆~港九培靈會~劉少康~第六講}
\label{sec:1090}
\textbf{委身的重建 — 以斯帖的榜樣}
\newline
\newline
連結: \href{https://www.hkbibleconference.org/session-message/view/1090}{\texttt{https://www.hkbibleconference.org/session-message/view/1090}}
\newline
\newline
\hyperref[sec:1089]{< < < PREV SERMON < < <}
~
\hyperref[sec:toc]{[返目錄]}
~
\hyperref[sec:1091]{> > > NEXT SERMON > > >}
\newline
\newline


\newpage

\section{第62屆~港九培靈會~劉少康~第七講}
\label{sec:1091}
\textbf{優次的重建 — 哈該的信息}
\newline
\newline
連結: \href{https://www.hkbibleconference.org/session-message/view/1091}{\texttt{https://www.hkbibleconference.org/session-message/view/1091}}
\newline
\newline
\hyperref[sec:1090]{< < < PREV SERMON < < <}
~
\hyperref[sec:toc]{[返目錄]}
~
\hyperref[sec:1092]{> > > NEXT SERMON > > >}
\newline
\newline


\newpage

\section{第62屆~港九培靈會~劉少康~第八講}
\label{sec:1092}
\textbf{信心的重建 — 撒迦利亞的挑戰}
\newline
\newline
連結: \href{https://www.hkbibleconference.org/session-message/view/1092}{\texttt{https://www.hkbibleconference.org/session-message/view/1092}}
\newline
\newline
\hyperref[sec:1091]{< < < PREV SERMON < < <}
~
\hyperref[sec:toc]{[返目錄]}
~
\hyperref[sec:1093]{> > > NEXT SERMON > > >}
\newline
\newline


\newpage

\section{第62屆~港九培靈會~劉少康~第九講}
\label{sec:1093}
\textbf{愛心的重建 — 瑪拉基的勸勉}
\newline
\newline
連結: \href{https://www.hkbibleconference.org/session-message/view/1093}{\texttt{https://www.hkbibleconference.org/session-message/view/1093}}
\newline
\newline
\hyperref[sec:1092]{< < < PREV SERMON < < <}
~
\hyperref[sec:toc]{[返目錄]}
~
\hyperref[sec:1094]{> > > NEXT SERMON > > >}
\newline
\newline


\newpage

\section{第62屆~港九培靈會~杜偉力~第一講}
\label{sec:1094}
\textbf{往下扎根}
\newline
\newline
連結: \href{https://www.hkbibleconference.org/session-message/view/1094}{\texttt{https://www.hkbibleconference.org/session-message/view/1094}}
\newline
\newline
\hyperref[sec:1093]{< < < PREV SERMON < < <}
~
\hyperref[sec:toc]{[返目錄]}
~
\hyperref[sec:1096]{> > > NEXT SERMON > > >}
\newline
\newline
詩篇 51:1-51:19
\newline
\begin{longtable}{cl}
\hline
\hline
章節 & 經文 (和合本修訂版)\\
\hline
51:1 & \begin{tabularx}{0.7\textwidth}{X} 神啊，求你按你的慈愛恩待我！按你豐盛的憐憫塗去我的過犯！ \end{tabularx} \\ \\ \relax
51:2 & \begin{tabularx}{0.7\textwidth}{X} 求你將我的罪孽洗滌淨盡，潔除我的罪！ \end{tabularx} \\ \\ \relax
51:3 & \begin{tabularx}{0.7\textwidth}{X} 因為我知道我的過犯；我的罪常在我面前。 \end{tabularx} \\ \\ \relax
51:4 & \begin{tabularx}{0.7\textwidth}{X} 我向你犯罪，惟獨得罪了你，在你眼前行了這惡，以致你責備的時候顯為公義，判斷的時候顯為清白。 \end{tabularx} \\ \\ \relax
51:5 & \begin{tabularx}{0.7\textwidth}{X} 看哪，我是在罪孽裡生的，在我母親懷胎的時候就有了罪。 \end{tabularx} \\ \\ \relax
51:6 & \begin{tabularx}{0.7\textwidth}{X} 你所喜愛的是內心的誠實；求你在我隱密處使我得智慧。 \end{tabularx} \\ \\ \relax
51:7 & \begin{tabularx}{0.7\textwidth}{X} 求你用牛膝草潔淨我，我就乾淨；求你洗滌我，我就比雪更白。 \end{tabularx} \\ \\ \relax
51:8 & \begin{tabularx}{0.7\textwidth}{X} 求你使我得聽歡喜快樂的聲音，使你所壓傷的骨頭可以踴躍。 \end{tabularx} \\ \\ \relax
51:9 & \begin{tabularx}{0.7\textwidth}{X} 求你轉臉不看我的罪，塗去我一切的罪孽。 \end{tabularx} \\ \\ \relax
51:10 & \begin{tabularx}{0.7\textwidth}{X} 神啊，求你為我造清潔的心，使我裡面重新有正直的靈。 \end{tabularx} \\ \\ \relax
51:11 & \begin{tabularx}{0.7\textwidth}{X} 不要丟棄我，使我離開你的面；不要從我收回你的聖靈。 \end{tabularx} \\ \\ \relax
51:12 & \begin{tabularx}{0.7\textwidth}{X} 求你使我重得救恩之樂，以樂意的靈來扶持我， \end{tabularx} \\ \\ \relax
51:13 & \begin{tabularx}{0.7\textwidth}{X} 我就把你的道指教有過犯的人，罪人必歸順你。 \end{tabularx} \\ \\ \relax
51:14 & \begin{tabularx}{0.7\textwidth}{X} 神啊，你是拯救我的神；求你救我脫離流人血的罪！我的舌頭就高唱你的公義。 \end{tabularx} \\ \\ \relax
51:15 & \begin{tabularx}{0.7\textwidth}{X} 主啊，求你使我嘴唇張開，我的口就傳揚讚美你的話！ \end{tabularx} \\ \\ \relax
51:16 & \begin{tabularx}{0.7\textwidth}{X} 你本不喜愛祭物，若喜愛，我就獻上；燔祭你也不喜悅。 \end{tabularx} \\ \\ \relax
51:17 & \begin{tabularx}{0.7\textwidth}{X} 神所要的祭就是憂傷的靈；神啊，憂傷痛悔的心，你必不輕看。 \end{tabularx} \\ \\ \relax
51:18 & \begin{tabularx}{0.7\textwidth}{X} 求你隨你的美意善待錫安，建造耶路撒冷的城牆。 \end{tabularx} \\ \\ \relax
51:19 & \begin{tabularx}{0.7\textwidth}{X} 那時，你必喜愛公義的祭和燔祭，全牲的燔祭；那時，人必將公牛獻在你壇上。 \end{tabularx} \\ \\
[1ex]
\hline
\hline
\end{longtable}


\newpage

\section{第62屆~港九培靈會~杜偉力~第二講}
\label{sec:1096}
\textbf{往下扎根}
\newline
\newline
連結: \href{https://www.hkbibleconference.org/session-message/view/1096}{\texttt{https://www.hkbibleconference.org/session-message/view/1096}}
\newline
\newline
\hyperref[sec:1094]{< < < PREV SERMON < < <}
~
\hyperref[sec:toc]{[返目錄]}
~
\hyperref[sec:1098]{> > > NEXT SERMON > > >}
\newline
\newline


\newpage

\section{第62屆~港九培靈會~杜偉力~第三講}
\label{sec:1098}
\textbf{往下扎根}
\newline
\newline
連結: \href{https://www.hkbibleconference.org/session-message/view/1098}{\texttt{https://www.hkbibleconference.org/session-message/view/1098}}
\newline
\newline
\hyperref[sec:1096]{< < < PREV SERMON < < <}
~
\hyperref[sec:toc]{[返目錄]}
~
\hyperref[sec:1102]{> > > NEXT SERMON > > >}
\newline
\newline


\newpage

\section{第62屆~港九培靈會~杜偉力~第四講}
\label{sec:1102}
\textbf{往下扎根}
\newline
\newline
連結: \href{https://www.hkbibleconference.org/session-message/view/1102}{\texttt{https://www.hkbibleconference.org/session-message/view/1102}}
\newline
\newline
\hyperref[sec:1098]{< < < PREV SERMON < < <}
~
\hyperref[sec:toc]{[返目錄]}
~
\hyperref[sec:1105]{> > > NEXT SERMON > > >}
\newline
\newline


\newpage

\section{第62屆~港九培靈會~杜偉力~第五講}
\label{sec:1105}
\textbf{向上結果}
\newline
\newline
連結: \href{https://www.hkbibleconference.org/session-message/view/1105}{\texttt{https://www.hkbibleconference.org/session-message/view/1105}}
\newline
\newline
\hyperref[sec:1102]{< < < PREV SERMON < < <}
~
\hyperref[sec:toc]{[返目錄]}
~
\hyperref[sec:1107]{> > > NEXT SERMON > > >}
\newline
\newline


\newpage

\section{第62屆~港九培靈會~杜偉力~第六講}
\label{sec:1107}
\textbf{向上結果}
\newline
\newline
連結: \href{https://www.hkbibleconference.org/session-message/view/1107}{\texttt{https://www.hkbibleconference.org/session-message/view/1107}}
\newline
\newline
\hyperref[sec:1105]{< < < PREV SERMON < < <}
~
\hyperref[sec:toc]{[返目錄]}
~
\hyperref[sec:1109]{> > > NEXT SERMON > > >}
\newline
\newline


\newpage

\section{第62屆~港九培靈會~杜偉力~第七講}
\label{sec:1109}
\textbf{向上結果}
\newline
\newline
連結: \href{https://www.hkbibleconference.org/session-message/view/1109}{\texttt{https://www.hkbibleconference.org/session-message/view/1109}}
\newline
\newline
\hyperref[sec:1107]{< < < PREV SERMON < < <}
~
\hyperref[sec:toc]{[返目錄]}
~
\hyperref[sec:1112]{> > > NEXT SERMON > > >}
\newline
\newline


\newpage

\section{第62屆~港九培靈會~杜偉力~第八講}
\label{sec:1112}
\textbf{向上結果}
\newline
\newline
連結: \href{https://www.hkbibleconference.org/session-message/view/1112}{\texttt{https://www.hkbibleconference.org/session-message/view/1112}}
\newline
\newline
\hyperref[sec:1109]{< < < PREV SERMON < < <}
~
\hyperref[sec:toc]{[返目錄]}
~
\hyperref[sec:1113]{> > > NEXT SERMON > > >}
\newline
\newline


\newpage

\section{第62屆~港九培靈會~杜偉力~第九講}
\label{sec:1113}
\textbf{向上結果}
\newline
\newline
連結: \href{https://www.hkbibleconference.org/session-message/view/1113}{\texttt{https://www.hkbibleconference.org/session-message/view/1113}}
\newline
\newline
\hyperref[sec:1112]{< < < PREV SERMON < < <}
~
\hyperref[sec:toc]{[返目錄]}
~
\hyperref[sec:1114]{> > > NEXT SERMON > > >}
\newline
\newline


\newpage

\section{第62屆~港九培靈會~杜偉力~第十講}
\label{sec:1114}
\textbf{向上結果}
\newline
\newline
連結: \href{https://www.hkbibleconference.org/session-message/view/1114}{\texttt{https://www.hkbibleconference.org/session-message/view/1114}}
\newline
\newline
\hyperref[sec:1113]{< < < PREV SERMON < < <}
~
\hyperref[sec:toc]{[返目錄]}
~
\hyperref[sec:1284]{> > > NEXT SERMON > > >}
\newline
\newline
詩篇 51:1-51:19
\newline
\begin{longtable}{cl}
\hline
\hline
章節 & 經文 (和合本修訂版)\\
\hline
51:1 & \begin{tabularx}{0.7\textwidth}{X} 神啊，求你按你的慈愛恩待我！按你豐盛的憐憫塗去我的過犯！ \end{tabularx} \\ \\ \relax
51:2 & \begin{tabularx}{0.7\textwidth}{X} 求你將我的罪孽洗滌淨盡，潔除我的罪！ \end{tabularx} \\ \\ \relax
51:3 & \begin{tabularx}{0.7\textwidth}{X} 因為我知道我的過犯；我的罪常在我面前。 \end{tabularx} \\ \\ \relax
51:4 & \begin{tabularx}{0.7\textwidth}{X} 我向你犯罪，惟獨得罪了你，在你眼前行了這惡，以致你責備的時候顯為公義，判斷的時候顯為清白。 \end{tabularx} \\ \\ \relax
51:5 & \begin{tabularx}{0.7\textwidth}{X} 看哪，我是在罪孽裡生的，在我母親懷胎的時候就有了罪。 \end{tabularx} \\ \\ \relax
51:6 & \begin{tabularx}{0.7\textwidth}{X} 你所喜愛的是內心的誠實；求你在我隱密處使我得智慧。 \end{tabularx} \\ \\ \relax
51:7 & \begin{tabularx}{0.7\textwidth}{X} 求你用牛膝草潔淨我，我就乾淨；求你洗滌我，我就比雪更白。 \end{tabularx} \\ \\ \relax
51:8 & \begin{tabularx}{0.7\textwidth}{X} 求你使我得聽歡喜快樂的聲音，使你所壓傷的骨頭可以踴躍。 \end{tabularx} \\ \\ \relax
51:9 & \begin{tabularx}{0.7\textwidth}{X} 求你轉臉不看我的罪，塗去我一切的罪孽。 \end{tabularx} \\ \\ \relax
51:10 & \begin{tabularx}{0.7\textwidth}{X} 神啊，求你為我造清潔的心，使我裡面重新有正直的靈。 \end{tabularx} \\ \\ \relax
51:11 & \begin{tabularx}{0.7\textwidth}{X} 不要丟棄我，使我離開你的面；不要從我收回你的聖靈。 \end{tabularx} \\ \\ \relax
51:12 & \begin{tabularx}{0.7\textwidth}{X} 求你使我重得救恩之樂，以樂意的靈來扶持我， \end{tabularx} \\ \\ \relax
51:13 & \begin{tabularx}{0.7\textwidth}{X} 我就把你的道指教有過犯的人，罪人必歸順你。 \end{tabularx} \\ \\ \relax
51:14 & \begin{tabularx}{0.7\textwidth}{X} 神啊，你是拯救我的神；求你救我脫離流人血的罪！我的舌頭就高唱你的公義。 \end{tabularx} \\ \\ \relax
51:15 & \begin{tabularx}{0.7\textwidth}{X} 主啊，求你使我嘴唇張開，我的口就傳揚讚美你的話！ \end{tabularx} \\ \\ \relax
51:16 & \begin{tabularx}{0.7\textwidth}{X} 你本不喜愛祭物，若喜愛，我就獻上；燔祭你也不喜悅。 \end{tabularx} \\ \\ \relax
51:17 & \begin{tabularx}{0.7\textwidth}{X} 神所要的祭就是憂傷的靈；神啊，憂傷痛悔的心，你必不輕看。 \end{tabularx} \\ \\ \relax
51:18 & \begin{tabularx}{0.7\textwidth}{X} 求你隨你的美意善待錫安，建造耶路撒冷的城牆。 \end{tabularx} \\ \\ \relax
51:19 & \begin{tabularx}{0.7\textwidth}{X} 那時，你必喜愛公義的祭和燔祭，全牲的燔祭；那時，人必將公牛獻在你壇上。 \end{tabularx} \\ \\
[1ex]
\hline
\hline
\end{longtable}


\newpage

\section{第62屆~港九培靈會~楊濬哲~序}
\label{sec:1284}
\textbf{序}
\newline
\newline
連結: \href{https://www.hkbibleconference.org/session-message/view/1284}{\texttt{https://www.hkbibleconference.org/session-message/view/1284}}
\newline
\newline
\hyperref[sec:1114]{< < < PREV SERMON < < <}
~
\hyperref[sec:toc]{[返目錄]}
~
\hyperref[sec:1076]{> > > NEXT SERMON > > >}
\newline
\newline
90年第62屆港九培靈研經大會,能如期舉行,圓滿閉會,實在是父神莫大的厚恩.在該日子裡伊拉克進軍科威特,引致中東局勢緊張；雖然各國譴責伊拉克,但侯賽因一意孤行,促成海灣之戰.石油股市,均受波動!
\newline
\newline
在一個動盪不安的世代裡,惟有神的話語是我們足以倚靠的.港九培靈會仍然擔任著一個重要的角色,就是忠心傳揚神的話.
\newline
\newline
本屆研經會請鮑會園牧師主講.總題:生命與生活──腓立比書的研究.他強調神所看的不是教會的大小,不是表面的成就多或寡；而是真理的根基是否穩固,生命的實踐是否真實.只有生命的根基建立好了,無論將來的環境怎樣改變,或者面臨前所未有的挑戰,我們就仍站立得住.劉少康牧師主講講道會.總題:來吧!我們重新建造.內容以選民被擄至巴比倫七十年期滿,神應許他們回國重建聖殿,聖城；恢復對神的信心與愛心勉眾.該怎樣面對這個時代?對自己屬靈生命是否滿意?是否滿意教會發展?我們必須重建我們的靈命,重建我們的信心和愛神的心；好叫我們聽見呼聲說:「基督徒呀!站出來.」杜偉力牧師原籍印度.憶已故委辦元老張吉盛先生,生前曾言；盼望港九培靈會,能夠有一屆請到一位印度籍牧師來領會就好了；這多年來的禱告,神果然成就了他的心願.杜牧師主領晚堂奮興會.總題:往下扎根,向上結果.他強調這個信息對香港教會的情況很適用的.香港基督徒面對九七問題,作出各樣的計劃準備；印度教會也為大家祈禱.在這情形之下,神希望祂的子民在任何景況中,在陽光或在風雨之下；在痛苦,逼迫,苦難或繁榮之時,必須結出果子來.杜牧師在信息中,把印度這個古國的文化,宗教,教會的背景,都帶了出來；使我們對神愛世人,願意萬人得救的旨意,更為深刻銘骨!
\newline
\newline
本屆大會赴會人數,略遜於往年；但決志人數,反而增多,合共728人.(決志信主者五人,悔改更新為主而活者565人,全時間事奉者138人,資料不詳者20人).赴會者來自港九市區,以及離島,筲箕灣,香港仔,鴨利洲,田灣,華富,元朗,屯門,上水,大埔,西貢,將軍澳等地.
\newline
\newline
杜偉力牧師於感恩會中致詞謂:十日之聚會,已使我投入成為你們的一份子,有機會與你們一起同工事奉,感到十分興奮.每晚見到會眾飢渴赴會,帶聖經寫筆記,令我難以忘懷.大會能維持62年,誠屬不易；時代變遷,使許多屬靈聚會漸滲入娛樂節目,變成世俗化!培靈會能忠心傳主話,主必能保守你們,劉少康牧師致詞謂；我帶著戰兢之心參予事奉,因堅浸重建而有感傳重建之信息；先有領受,然後分享.大會期間,相信有不少代禱的支持者,至感欣慰!鮑會園牧師謂:我們三人講題均不一,但能有美好的配搭,十分奇妙!我對大眾傳播有負擔,欣聞培靈會之聲帶能藉廣播傳遞至遠方,對許多沒有教牧同工之教會,必大有幫助.
\newline
\newline
委辦會元老張約翰先生及司庫梁冠倫先生,畢生忠心事主；為培靈事工,尤為盡力；雖被主接去,仍留下典範,使吾人知所效法!相信主裡的勞苦,不是徒然的!感謝神!又興起城浸執事會主席,陳世英先生加入為委辦,共勷聖工.
\newline
\newline
本講道集,乃62屆大會之講道記錄.蒙王大衛弟兄,記錄研經會.梁家麗姊妹,梁莊儀姊妹,歐瑞蘭姊妹,林蘊詩姊妹,周婉玲姊妹,姚麗妮姊妹,李秋華姊妹,林小玲姊妹,鄭坤英姊妹等記錄講道會.雷麗端姊妹記錄奮興會.朱建磯牧師整理編校,乃得如期出版,謹此致謝!
\newline
\newline
每屆大會,均蒙九龍城浸信會,借用為會場；執事,同工,義工及各兄姊協助招待,維持秩序.福音傳播中心,錄音,錄影；基督教週報刊登新聞及廣告；神托之家,耀安工場,寄發宣傳品；彼此配搭廣事宣傳,藉收效益.蒙各教會兄姊踴躍赴會及奉獻.統此致謝!
\newline
\newline


\newpage

\section{第62屆~港九培靈會~鮑會園~第一講}
\label{sec:1076}
\textbf{現實的智慧──分辨是非的禱告}
\newline
\newline
連結: \href{https://www.hkbibleconference.org/session-message/view/1076}{\texttt{https://www.hkbibleconference.org/session-message/view/1076}}
\newline
\newline
\hyperref[sec:1284]{< < < PREV SERMON < < <}
~
\hyperref[sec:toc]{[返目錄]}
~
\hyperref[sec:1077]{> > > NEXT SERMON > > >}
\newline
\newline
腓立比書 1:9-1:19
\newline
\begin{longtable}{cl}
\hline
\hline
章節 & 經文 (和合本修訂版)\\
\hline
1:9 & \begin{tabularx}{0.7\textwidth}{X} 我所禱告的就是：要你們的愛心，在知識和各樣見識上，不斷增長， \end{tabularx} \\ \\ \relax
1:10 & \begin{tabularx}{0.7\textwidth}{X} 使你們能分辨是非，在基督的日子作真誠無可指責的人， \end{tabularx} \\ \\ \relax
1:11 & \begin{tabularx}{0.7\textwidth}{X} 更靠著耶穌基督結滿仁義的果子，歸榮耀稱讚給神。 \end{tabularx} \\ \\ \relax
1:12 & \begin{tabularx}{0.7\textwidth}{X} 弟兄們，我要你們知道，我所遭遇的事反而使福音更興旺， \end{tabularx} \\ \\ \relax
1:13 & \begin{tabularx}{0.7\textwidth}{X} 以致御營全軍和其餘的人都知道我是為基督的緣故受捆鎖的； \end{tabularx} \\ \\ \relax
1:14 & \begin{tabularx}{0.7\textwidth}{X} 而且那在主裡的弟兄，多半都因我受的捆鎖而篤信不疑，越發放膽無所懼怕地傳道。 \end{tabularx} \\ \\ \relax
1:15 & \begin{tabularx}{0.7\textwidth}{X} 有些人傳基督是出於嫉妒紛爭；有些人是出於好意。 \end{tabularx} \\ \\ \relax
1:16 & \begin{tabularx}{0.7\textwidth}{X} 後者是出於愛心，知道我奉差遣是為福音辯護的。 \end{tabularx} \\ \\ \relax
1:17 & \begin{tabularx}{0.7\textwidth}{X} 前者傳基督是出於自私，動機不純，企圖要加增我捆鎖的苦楚。 \end{tabularx} \\ \\ \relax
1:18 & \begin{tabularx}{0.7\textwidth}{X} 這又何妨呢？或是假意或是真心，無論如何，只要基督被傳開了，為此我就歡喜。我還要歡喜， \end{tabularx} \\ \\ \relax
1:19 & \begin{tabularx}{0.7\textwidth}{X} 因為我知道，這事藉著你們的祈禱和耶穌基督的靈的幫助，終必使我得到釋放。 \end{tabularx} \\ \\
[1ex]
\hline
\hline
\end{longtable}
過去,神給我很多機會在培靈會與大家一同交通,但愈多機會思想神的話,我的心情就愈加嚴肅.我從來沒有像這次培靈會那麼緊張,因為在這時代,具有一個非常嚴肅的責任.
\newline
\newline
我一直都很喜歡腓立比這卷書信,過去也曾有許多機會在不同的場合思想本書的某段經文,但從未將整卷書信連在一起來思想.因此,請大家特為我們禱告,願腓立比書以亮光幫助我們,也盼望神藉祂的話賜福予我們.
\newline
\newline
簡介
\newline
\newline
腓立比教會是比較特別的教會,在保羅經常工作的教會中,腓立比教會算是比較小的；雖然它所處的地理位置很重要,如同歐洲,亞洲的大門一樣.但這城市卻不很大,猶太人數目也不多；且很少受猶太教影響,人也比較單純.所以,當保羅一來傳道,那些人很快就接受了福音,教會也很快就建立起來了,而且在真理上有很好的根基.雖然保羅在腓立比的工作時間較在以弗所,哥林多都短,但信徒在真理上的領受非常深入,知道怎樣把真理應用於生活上,難怪保羅常提到馬其頓教會是眾教會的榜樣.這個事實給予我們很大的激勵,因為神所看的不是教會有很大,或表面的成就多好；而是真理的根基是否穩固,生命的實踐是否真實.
\newline
\newline
腓立比書也是很特別的書信,它不像羅馬書那樣作系統教導,也沒有加拉太書中的錯誤信仰以及歌羅西書中的異端思想；更沒有理論性的教訓,卻是保羅自己生命中的親身體驗.除了哥林多後書外,此書是保羅提到個人經驗最多的一卷書,因為保羅已將真理融會在他整個生命裡,並以此來建立弟兄姊妹的生命.其實這是給我們有更深入的思想:生命的質素,才是被神使用的基本原因,只有生命的根基建立好了,無論將來的環境怎樣改變,或者面臨前所未有的挑戰,我們就仍站得住.
\newline
\newline
渴慕的禱告
\newline
\newline
從(腓一8)看出保羅禱告時,是切切的想念眾人,「想念」可譯作「渴慕」,「渴慕」的對象可能是指腓立比人,也可能是指禱告的內容；意思是他深切願望所祈求的真能成就在他們中間.這裡提醒我們；禱告並不在乎好聽,乃在於是否有渴慕的心將事情帶到神面前.若不是這樣,禱告就失去了目的,祈求當然也沒有果效.
\newline
\newline
愛心長進的禱告
\newline
\newline
「我所禱告的,就是要你們的愛心,在知識和各樣見識上,多而又多」.(一9)這是禱告蒙允,生命蒙福的先決條件.「多而又多」是指繼續不斷的增長, 衡量屬靈生命成長的標準不是知識,也不是見識；乃是愛心是否有長進.若愛心沒有長進,那麼其他方面的成就也不算得甚麼,一切的果效也沒有用處(林前十三 2-3).若以工作果效來說,舊約先知約拿是最成功的,因為他只在尼尼微城走了一遍,宣告了一個信息,合城上下就認罪悔改；但約拿卻走到城外觀看,盼望神毀滅該城；結果直到今天,約拿還是成為受責的先知.沒有愛心的工作算不得甚麼；神要的不是工作表現,乃是愛心的長進.這裡有兩方面值得思想的；
\newline
\newline
(一)愛主的心:「若有人不愛主,這人可詛可咒」.(林前十六22)人若沒有愛主的心,就不配作神的工作；若果追求只為了責任,工作或學習,只為了表現；那麼就不算是事奉了.若不是出於愛神的心,一切屬靈活動都必失卻價值.主耶穌在加利利海邊問彼得:「你愛我比這些更深麼」?彼得回答說:「主啊,你知道我愛你」就因這心志,主差彼得餵養小羊.若彼得以自己的決心,經驗或恩賜作考慮,而不是基於對主的愛；那麼他就沒有資格來餵養主的羊了.
\newline
\newline
(二)愛人的心:保羅勸勉我們不但要有愛主的心,也要有愛人的心.不論在教會裡站在任何崗位,無論是牧者或普通的弟兄姊妹,我們若不愛看得見的肢體；那麼愛主的心也是不真實的.
\newline
\newline
見識長進的禱告
\newline
\newline
「見識」意指屬靈的眼光,能看見事物的真正意義；這不單是明白真理,而且還知道真理與我們生命的關係.
\newline
\newline
「見識」的另一個解釋,記載在(代上十二32):「通達時務,知道以色列人所當行的」就是見識了,「通達時務」是指明白在這環境中,作為神的百姓應當如何行.我們在追求屬靈的事上也要有這樣的見識,不單要洞察局勢,更重要的是知道在這情況下該怎樣行；才合乎神的心意,對屬靈生命的見證才有幫助.
\newline
\newline
在未來的環境裡我們可能要面對時勢,社會,人情,甚至教會的改變；若沒有屬靈的見識,我們就不知道在這劇變的情況下應當怎樣行.
\newline
\newline
分辨是非的禱告
\newline
\newline
保羅不單為眾人的愛心在知識和見識上長進而禱告,更渴望他們能分辨是非(一10).「分辨是非」在正常的觀念裡,可能只是好壞,善惡的分別；但經文中聖經的小字卻把原來文字的意思表達出來了──「分辨」是指喜愛美好的事,不僅是知道,還有願意遵行的心,對基督徒來說,分辨事物的對錯並不困難,而腓立比的信徒也沒有這個需要；反而他們在真理和生活上都有很好的表現.保羅用不著在這方面為他們操心.在此,保羅乃是勉勵他們要追求喜愛美好的事情,把更好與次好的事情區別出來；按整體書信的教訓,我們可把第一節的「分別是非」歸納作下列幾方面:
\newline
\newline
(一)分辨真正美好的事:許多事情表面上看起來很好,但對具屬靈見識的人來說,就能很容易從真正美好的事情中分別出來；有屬靈見識的人不會單看表面上的吸引和美好.例如:有人以為某些禱告的動作看來很屬靈,於是就跟別人學起來了；其實,這些表面的功夫是非常容易學會的,但卻與屬靈生命完全沒有關係. 所以不要把假的屬靈記號,看作生命追求的目標.
\newline
\newline
(二)分辨真正重要的事......在生活中有許多美好的事情,不一定具真正的屬靈價值.現代人生活緊張,節奏急速,傳道人也是整天忙著開會,忙著應付緊急的事情；忽略了生命中最重要,最具價值的事情；這包括對自己靈命的儆醒,對弟兄姊妹屬靈方面的守望責任等.如果我們不懂得分辨甚麼才是真正美好, 重要的事,那麼就是沒有屬靈見識了.
\newline
\newline
(三)分辨事情的限度:有些事情確具美好,也有價值,但不一定要把全部的時間放在自己喜愛的「重要」事上,而忽略了自己在屬靈生命上應負的責任.年青人往往毫無節制地做自己喜歡做的事.有些學生平常忽略認真讀書的責任,到臨考時就慌忙請年老宣教士為她禱告；而自己也一邊喝汽水,一邊禱告.欲以不停的禱告來逃避讀書的責任,這是極其沒有見識的.
\newline
\newline
(四)分辨工作的方法:要用最好的方法來完成神交托的工作,不要為了一時的利益或方便而妥協,更不應用不同的標準來作自己的事和教會的事.在做一切的事情之前,都要仔細地想清楚,這是否最好方法.
\newline
\newline
(五)分辨生活的試探:生活中某些表現在世人看來只是小缺點,但若不小心防備,可能鑄成大錯而破壞整個生命.雅歌書第二章描寫書拉密女為王捉拿葡萄園裡的小狐狸,因為小狐狸雖不傷人；且看來可愛,但卻把園中剛開花的莊稼毀壞了；同樣,在我們的生活中也要防備小試探,一點一點懶惰,自私,驕傲,貪婪......都可能叫我們失去見證.
\newline
\newline
總括來說,我們要在屬靈生命上儆醒追求,不但愛心上有長進,屬靈的見識上也要有長進,因此能分辨是非,作成熟的基督徒.
\newline
\newline


\newpage

\section{第62屆~港九培靈會~鮑會園~第二講}
\label{sec:1077}
\textbf{最好的選擇──活著就是基督}
\newline
\newline
連結: \href{https://www.hkbibleconference.org/session-message/view/1077}{\texttt{https://www.hkbibleconference.org/session-message/view/1077}}
\newline
\newline
\hyperref[sec:1076]{< < < PREV SERMON < < <}
~
\hyperref[sec:toc]{[返目錄]}
~
\hyperref[sec:1078]{> > > NEXT SERMON > > >}
\newline
\newline
腓立比書 1:20-1:30
\newline
\begin{longtable}{cl}
\hline
\hline
章節 & 經文 (和合本修訂版)\\
\hline
1:20 & \begin{tabularx}{0.7\textwidth}{X} 這就是我所切慕、所盼望的：沒有一事能使我羞愧；反倒凡事坦然無懼，無論是生是死，總要讓基督在我身上照常顯大。 \end{tabularx} \\ \\ \relax
1:21 & \begin{tabularx}{0.7\textwidth}{X} 因為我活著就是基督，死了就有益處。 \end{tabularx} \\ \\ \relax
1:22 & \begin{tabularx}{0.7\textwidth}{X} 但是，我在肉身活著，若能有工作的成果，我就不知道該挑選甚麼。 \end{tabularx} \\ \\ \relax
1:23 & \begin{tabularx}{0.7\textwidth}{X} 我處在兩難之間：我情願離世與基督同在，因為這是好得無比的； \end{tabularx} \\ \\ \relax
1:24 & \begin{tabularx}{0.7\textwidth}{X} 然而，我為你們肉身活著更加要緊。 \end{tabularx} \\ \\ \relax
1:25 & \begin{tabularx}{0.7\textwidth}{X} 既然我這樣深信，就知道仍要留在世間，且與你們眾人一起存留，使你們在所信的道上又長進又喜樂， \end{tabularx} \\ \\ \relax
1:26 & \begin{tabularx}{0.7\textwidth}{X} 為了我再到你們那裡時，你們在基督耶穌裡的誇耀越發加增。 \end{tabularx} \\ \\ \relax
1:27 & \begin{tabularx}{0.7\textwidth}{X} 最重要的是：你們行事為人要與基督的福音相稱，這樣，無論我來見你們，或不在你們那裡，都可以聽到你們的景況，知道你們同有一個心志，站立得穩，為福音的信仰齊心努力， \end{tabularx} \\ \\ \relax
1:28 & \begin{tabularx}{0.7\textwidth}{X} 絲毫不怕敵人的威脅；以此證明他們會沉淪，你們會得救，這是出於神。 \end{tabularx} \\ \\ \relax
1:29 & \begin{tabularx}{0.7\textwidth}{X} 因為你們蒙恩，不但得以信服基督，而且要為他受苦。 \end{tabularx} \\ \\ \relax
1:30 & \begin{tabularx}{0.7\textwidth}{X} 你們的爭戰，就與你們曾在我身上見過、現在所聽到的是一樣的。 \end{tabularx} \\ \\
[1ex]
\hline
\hline
\end{longtable}
從這段經文中可以看出,當時保羅正面對極大的考驗:身處監獄而不知前途如何.但他藉著腓立比信徒的禱告和耶穌基督之靈的幫助,終必叫他得救.保羅這裡所說的「得救」不單是指從獄中得釋放,更是渴望神在他生命中完成救贖的目的.「得救」的意義相當豐富,不僅是脫離罪惡進入永生；而且還包含生命成為聖潔,甚至將來與主同得榮耀的身體得贖,都包括在救恩的意義裡(羅馬書八章).但在這段經文中,保羅說的「得救」是指在生命中完成神救贖之目的,也是解釋救恩最基本的意義.
\newline
\newline
腓一20「照著我的切慕所盼望的,沒有一事叫我羞愧」.保羅並沒有因著所遭遇的事情,而感到羞愧；彼得書信也教導信徒不要為生活中的困難,考驗,譏笑,辱罵等而覺羞恥.該羞恥的只有我們的罪,不是因犯罪受刑罰而羞恥；乃是因犯罪得罪神而羞恥.亞當犯罪後感到羞恥,只是因著赤身露體及罪惡被揭露,所以也就躲藏起來,這種面對罪惡的態度是不正確的.
\newline
\newline
20節的下半節「無論是生,是死,總叫基督在我身上照常顯大」這裡所表達的意思很美,深信無論生活中遭遇任何事情,都是神在我身上的計劃；不是我故意生或死,或藉任何方法來高抬自己；相反,乃是存順服的心,讓基督藉著我的遭遇而顯為大.高舉的是基督不是我,主動的是神不是我.保羅就是為了這個目的來走生命的路程,也只有從這角度看21節的經文,才能顯出其意義來.
\newline
\newline
腓一21「我活著就是基督,我死了就有益處」,是整段經文的中心思想.
\newline
\newline
(一)我活著就是基督
\newline
\newline
「我活著就是基督」不是保羅自誇稱自己為基督,而是聲稱他活著,就是讓神在他身上成就自己的旨意.這裡,我們要從三方面來思想,基督徒在世上存在的目的是甚麼?
\newline
\newline
(一)活出基督的生命:羅馬書第八章指出神按著祂的旨意揀選我們；以弗所書則說我們生活的目的要長成基督的身量；還有一些經文叫我們的生命要活像基督,像神的兒子一樣.活像主耶穌,就是基督徒生命最基本的追求目的.
\newline
\newline
基督徒的生命不應以成就些甚麼作為目標,而要以追求生命的聖潔作為基本的使命.聖經常用「果子」來形容我們的生命,果子不是做來的,乃是結出來的,甚麼樣的生命,就結出甚麼果子來.因此,不要把生命追求的重點放錯,若果我們有了基督長成的身量,有了基督成形的生命；那麼,我們就自然能結出生命的果實來,而在生活中也自然有美好的表現.
\newline
\newline
(二)與恩相稱的生活:「只要你們行事為人與基督的福音相稱」(一27)是指我們行事為人表明基督的福音.「基督的福音」就是把主的生活充滿了恩典,「道成肉身住在我們中間,充充滿滿的有恩典有真理」(約一14)雖然我們沒有資格,能力來領受神的恩典；但神白白地將恩典給予我們,為的就是要叫我們顯出恩典和生命與生活.
\newline
\newline
恩典的生活必須與恩典的態度相配合,甚麼才算是有恩典呢?「恩典,這個字的原意是很特別的,英文字是指「有風度的生活」.例如:約翰福音記載有人把犯姦淫的婦人拉到耶穌面前(八1-11),但主耶穌的處理卻顯出滿有恩典的態度,他低著頭在地上畫字,沒有譏笑或鄙視,也沒有自以為義,這是風度.另外,當馬大埋怨馬利亞,沒有幫她一同服侍時；主耶穌沒有責備馬大或馬利亞.只對馬大說:「馬利亞已經選擇那上好的福份,是不能奪去的」.沒有責備人,這也是風度.還有,當保羅和巴拿巴帶著年青的馬可一起出外佈道,而馬可也許缺乏受害的經驗就中途離去；這事叫保羅很失望.結果在第二次旅程開始時,保羅極力反對巴拿巴還要帶馬可一起出外；而巴拿巴卻沒有因此指責保羅,也沒有為馬可辯護；於是就獨自帶同馬可踏上第二次佈道旅程.結果,保羅到了晚年時,托人叫馬可回來；因馬可在傳福音的事上是與他有益處.可見,巴拿巴在重要的關頭把馬可挽回來,沒有輕視,也沒有撇棄,這就是有風度的生活.今日在基督教的圈子裡,是多麼需要有這種生命的風度.既然領受了神豐富的恩典,就在日常生活中,要活出與恩典相稱的見證來.
\newline
\newline
(三)齊心努力傳福音:要過一個積極傳福音的生活,特別在今天這個時代,既然活著就是基督；就要齊心努力地為福音作見證.未來香港要面對改變和挑戰,也許傳道人工作受到限制,很多傳道的工作就要落在信徒的身上；基督徒要負起更大的責任.因此,今日平信徒訓練是非常重要的事,但教會若把太多的注意力放在門徒訓練,其目的只是為了植堂,訓練或栽培等事工；而忽略了傳福音的責任,那麼我們的生命目標可能放錯了,可能忽略了生命中最重要的使命.
\newline
\newline
「我活著就是基督」一方面要追求聖潔的生命,活出基督長成的身量,過與恩典相稱的生活；另一方面也要積極地為福音作見證,這就是救恩最基本的意義.
\newline
\newline
(二)我死了就有益處
\newline
\newline
這不是指消極逃避的觀念,反而具有兩方面積極的意義.
\newline
\newline
(一)與主有親密的交通:基督徒所渴望的就是永遠與主同在,與主有親密的交通,但只有我們面對面與主相見,才能達到這種親密的關係.保羅為著與主有親密的交通而渴望與主永遠同在,是積極,正確的；但我們若為其他的目的而希望與主同在,那麼我們的動機可能錯了.
\newline
\newline
(二)聖潔不犯罪的生活:保羅在世上追求聖潔的生活已幾十年了,但他仍然達不到完全不犯罪的地步.反而說「我所願意的善,我反不作；我所不願意的惡,我倒去作」.(羅七19).其實,我們今生的考驗就是為了預備將來與主面對面聖潔的生活；雖然今天我們的生活會軟弱,失敗；但將來得著完全的救贖而與主永遠同在,我們的生活就永遠不再犯罪了.假若沒有分辨善惡樹的考驗,亞當也沒有犯罪；不過他的生命即使是上天堂也經不起考驗.今日我們雖然犯了罪,但經過基督的救贖,神給予我們新造的生命,這就是將來永遠不再犯罪的生命了,所以,保羅說:「我死了就有益處」是指離開塵世進入永不犯罪的世界；這是非常美好的盼望.
\newline
\newline
既然如此,保羅為甚麼又說:「我正在兩難之間」(23節)呢?原因是在當時遭禁的情況下,保羅也做世上的工作,所以離開世界與主同在是他個人最大的喜好和盼望的；但為了腓立比信徒的好處,他在肉身活著更是要緊的(24節).保羅為了教會,為了神給他的托付,為了成就神的旨意,他甘心樂意將自己所喜愛,所願意的放下,而情願留在世上,這是不容易的選擇,也是蒙福的選擇.
\newline
\newline
今天,在我們的生活中,也同樣要面對許多的抉擇,我們是否願意為了弟兄姊妹的需要,為了教會的緣故,而放棄自己的喜好的.保羅的生命蒙神賜福的原因,神是因為他甘心樂意為主付上代價,作出犧牲.我們也願意這樣效法他嗎?
\newline
\newline


\newpage

\section{第62屆~港九培靈會~鮑會園~第三講}
\label{sec:1078}
\textbf{努力的方向──同心合意的事奉}
\newline
\newline
連結: \href{https://www.hkbibleconference.org/session-message/view/1078}{\texttt{https://www.hkbibleconference.org/session-message/view/1078}}
\newline
\newline
\hyperref[sec:1077]{< < < PREV SERMON < < <}
~
\hyperref[sec:toc]{[返目錄]}
~
\hyperref[sec:1079]{> > > NEXT SERMON > > >}
\newline
\newline
腓立比書 2:1-2:4
\newline
\begin{longtable}{cl}
\hline
\hline
章節 & 經文 (和合本修訂版)\\
\hline
2:1 & \begin{tabularx}{0.7\textwidth}{X} 所以，在基督裡若有任何勸勉，若有任何愛心的安慰，若有任何聖靈的團契，若有任何慈悲憐憫， \end{tabularx} \\ \\ \relax
2:2 & \begin{tabularx}{0.7\textwidth}{X} 你們就要意志相同，愛心相同，有一致的心思，一致的想法，使我的喜樂得以滿足。 \end{tabularx} \\ \\ \relax
2:3 & \begin{tabularx}{0.7\textwidth}{X} 凡事不可自私自利，不可貪圖虛榮；只要心存謙卑，各人看別人比自己強。 \end{tabularx} \\ \\ \relax
2:4 & \begin{tabularx}{0.7\textwidth}{X} 各人不要單顧自己的事，也要顧別人的事。 \end{tabularx} \\ \\
[1ex]
\hline
\hline
\end{longtable}
保羅寫腓立比書的目的,就是勉勵腓立比的弟兄姊妹,要在生活中表明生命的道.生命若失去見證,一切的工作也失去果效；因此,在第二章中保羅提醒信徒要同心合意完成基督的使命.
\newline
\newline
「腓二2」在中文聖經的繙譯上,好像是表達三個不同的思想；但實際上保羅要表達的只有一個中心要旨:同心合意興旺福音；也是整卷腓立比書的中心思想.
\newline
\newline
腓立比教會熱心於靈命的追求,其生命亦有長進,難道他們中間還有不同心的事情發生嗎?「腓四2」描述友阿爹和循都基兩位熱心愛主的女執事,她們與保羅一齊為福音勞苦.可能在生活中亦有很好的見證,既然如此她們又怎會不能同心合意地事奉呢?在今日的教會裡,我們也常常看見弟兄姊妹不同心的現象；雖然大家都熱心事主,追求成長；在大事上沒有失見證,但在小事上卻時常鬧意見.有時可能因著不同的看法,不同的性格,不同的背景或不同的做事方法而導致不同心.求主幫助我們不要讓小事攔阻同心合意的事奉；不要太介意別人的性格是否率直,保守或做事是否快,慢；要緊的是努力追求,同心合意地傳揚福音.這裡保羅提出之兩方面是值得我們一齊努力的方向.
\newline
\newline
(一)愛心相同
\newline
\newline
「腓二1-2」「所以在基督裡若有甚麼勸勉,愛心有甚麼安慰,聖靈有甚麼交通,心中有甚麼慈悲憐憫,你們就要意念相同,愛心相同,有一樣的心思,有 一樣的意念,使我的喜樂可以滿足.」這裡的「意念相同」和「有一樣的意念」都是副詞,用作形容追求「愛心相同,有一樣的心思」的態度；換句話說,保羅要信徒努力追求的不是意念相同,而是愛心相同,目標相同；有一樣的愛心,有一樣的追求方向.
\newline
\newline
「愛心相同」是指甚麼呢?是提醒我們要同樣地愛主,愛教會或愛使徒嗎?但經文中卻沒有勉勵我們如何愛主,只在第5節勸勉我們當效法主耶穌謙卑的榜樣.因此,這裡的「愛心相同」很可能是指弟兄姊妹之間要有相同的愛心；若我要別人怎麼愛我,那我就要怎樣愛別人；否則的話,我根本沒有要求別人愛我的權利.作傳道的,不能因傳道的身份而要求得到多點愛；平信徒也不能說因不懂得如何表達愛心,或者對教會沒有歸屬感而不去愛.我們要在教會得到愛,那麼我們就要先付出相同的愛；若果我們不拿出自己的愛心,卻單要求別人來愛自己,可以說我們根本沒有這個權利.整個教會要同心合意的事奉,首先要建立的就是「愛心相同」,學習在教會中主動表達愛心.
\newline
\newline
(二)一樣的心思
\newline
\newline
這裡的「心思」是指心中所掛念的,思想中看為重要的事情.也即是我們生命中所追求的方向及目標.教會中弟兄姊妹的心思若不一樣,彼此追求的目標若不一致；那麼就無法完成同心合意的事奉.例如說在教會裡有人看重團契,有人看重詩班,有人看重主日學；每個單位各顧各的,彼此沒有連繫；甚至互相指責對方,這樣的事奉又怎能稱得上「同心合意」呢?
\newline
\newline
有一樣的心思,即追求同一的方向,這是本段經文的中心思想.但基督徒應該追求的方向是甚麼呢?「腓二15-16」講得很清楚:「在這彎曲悖謬的世代,作神無暇疵的兒女.你們顯在這世代中,好像明光照耀；將生命的道表明出來.」,這就是我們整體追求的方向和目標,我們所努力的也只是為了這件事.既然教會有了同一的方向,同一的挑戰；那麼即使彼此看法不同,意見不同,但思想還是相同的.撒旦攻擊教會的最好方法,就是使教會失去同一的目標和方向,因此我們要時刻保持「一樣的心思」.
\newline
\newline
早期教會的修道士,為了更好地親近主而放棄許多的權利或享受；甚至連個人的婚姻也給放下了.既然他們有了清楚的方向和堅定的心志,那麼在一起事奉必定是同心合意的；但事實上卻剛好相反:他們常常為著眼前的小事而發生爭執,甚至彼此對抗；眼前的小事使他們失去共同追求的方向.取而代之的,是地位,權利,名譽的爭奪.在今天這個世代,我們同樣也需要有清晰的追求目標,否則的話,也必定會失去合一的見證.我們是為了甚麼來參加教會事奉呢?是單為了交通, 分享等目的嗎?我們的前面若沒有一樣的挑戰,那麼我們的生活就一定失去方向；只有共同面對同一的挑戰,追求同心合意的事奉,我們才能衝破困難；在這世代中如同明光照耀,將生命的道表明出來.
\newline
\newline
存心謙卑
\newline
\newline
「腓二3」「凡事不可結黨,不可貪圖虛浮的榮耀.只要存心謙卑,各人看別人比自己強.」教會的事奉若是為了貪圖虛浮的榮耀,同心合意是絕對沒有可能的.在一間教會裡有兩位女執事,開始時只是彼此間有點小誤會,後來一年多時間大家互不理睬,也沒有說話.結果雙方均漸有悔意,但為了虛浮的榮耀；雙方互不相讓,大家都不肯謙卑下來,不肯主動地先向對方打開話題.「貪圖虛浮榮耀」的事奉,是絕對不能同心合意的事奉,更不是存心謙卑的事奉.
\newline
\newline
有時候,謙卑是所有基督徒品格中最難學的功課,但我們若是為了追求同一的目標,而且又真正體會到主耶穌赦罪之恩典；深深感受到自己生命的不配,竭力追求作長大成熟的人.這樣,我們才能學會謙卑的功課；因為只有愈長進的人才愈謙卑,使徒保羅就是一個很好的例子:在哥林多前書,保羅自稱他是「使徒中最小的,不配稱為使徒」「林前十五9」,這是謙卑的表現；後來他屬靈生命再長進了,他不僅自覺是眾聖徒中最小的；更承認「在罪人中是個罪魁」.(提前一15) 由此可見,愈經歷神的恩典,生命會愈長進；我們就愈覺到自己的不配.既然主這樣愛我,我又有甚麼可以誇口的呢?別人看不起我,不尊重我,又有甚麼關係呢? 我是個卑微的人,別人尊重我與否,並不影響我在神面前的地位；更不影響我自己應該努力追求的方向.
\newline
\newline
保羅在本段經文中,提出三方面是基督徒應該追求的,只有這樣,我們才能同心合意地事奉!
\newline
\newline


\newpage

\section{第62屆~港九培靈會~鮑會園~第四講}
\label{sec:1079}
\textbf{完全的榜樣──以基督的心為心}
\newline
\newline
連結: \href{https://www.hkbibleconference.org/session-message/view/1079}{\texttt{https://www.hkbibleconference.org/session-message/view/1079}}
\newline
\newline
\hyperref[sec:1078]{< < < PREV SERMON < < <}
~
\hyperref[sec:toc]{[返目錄]}
~
\hyperref[sec:1080]{> > > NEXT SERMON > > >}
\newline
\newline
腓立比書 2:5-2:11
\newline
\begin{longtable}{cl}
\hline
\hline
章節 & 經文 (和合本修訂版)\\
\hline
2:5 & \begin{tabularx}{0.7\textwidth}{X} 你們當以基督耶穌的心為心： \end{tabularx} \\ \\ \relax
2:6 & \begin{tabularx}{0.7\textwidth}{X} 他本有神的形像，卻不堅持自己與神同等； \end{tabularx} \\ \\ \relax
2:7 & \begin{tabularx}{0.7\textwidth}{X} 反倒虛己，取了奴僕的形像，成為人的樣式；既有人的樣子， \end{tabularx} \\ \\ \relax
2:8 & \begin{tabularx}{0.7\textwidth}{X} 就謙卑自己，存心順服，以至於死，且死在十字架上。 \end{tabularx} \\ \\ \relax
2:9 & \begin{tabularx}{0.7\textwidth}{X} 所以神把他升為至高，又賜給他超乎萬名之上的名， \end{tabularx} \\ \\ \relax
2:10 & \begin{tabularx}{0.7\textwidth}{X} 使一切在天上的、地上的和地底下的，因耶穌的名，眾膝都要跪下， \end{tabularx} \\ \\ \relax
2:11 & \begin{tabularx}{0.7\textwidth}{X} 眾口都要宣認：耶穌基督是主，歸榮耀給父神。 \end{tabularx} \\ \\
[1ex]
\hline
\hline
\end{longtable}
保羅在腓立比書第二章1-4節,勸勉弟兄姊妹要同心合意地事奉主,而同心合意的重要條件是存心謙卑,在這裡他引用主耶穌自我謙卑的例子來解釋這個真理.聖經中許多重要真理都是以舉例的形式來表達.例如:羅馬書第五章講到因信稱義時,亦論及亞當「原罪」的問題,同樣,這裡提到主耶穌謙卑榜樣時,也解釋了道成肉身(約1-14)的真理.雖然如此,本章5-11節的中心思想,主要是表達以主耶穌的謙卑作為效法的榜樣.但主謙卑的目的卻是為了把人類從罪惡中拯救出來,完成救贖人類的使命.同樣,今天我們存心謙卑的目的也應該是要藉著我們的生命表現把人帶到神面前來,而不是要得別人的稱讚.
\newline
\newline
因此,(腓二6-7)就是解釋主耶穌如何謙卑自己來完成救贖的工作,其中包含非常重要的神學思想.基本上,主道成肉身顯明了兩方面的真理:一是放下或犧牲了一些東西,二是加上或取了一些東西.
\newline
\newline
(一)反例虛己
\newline
\newline
「祂本有神的形像,不以自己與神同等為強奪的」(6節)這裡的「本有」在原文應作「本是」,即「祂本是神的形像」.「本有」的東西是身外的,不是本身的一部份,即使失去了也影響不到不了其本性,但主耶穌卻「本是」神的形像.
\newline
\newline
「神的形像」在這裡不是指一般的照片,塑像等一些與我們相似的東西,因為「形像」這個字在其他經文中,是常見的.
\newline
\newline
但全部都用作形容外表的形像,全本聖經只有這裡的用法最特別,這裡的「形像」是指本性,即神的本性,而不是神的樣子.在神學上,我們或稱為神的屬性,例如:神是個靈,神是無所不知,無所不能,無所不在......我們就是用這些特性來形容神.因為這是神不可缺少的本性或形像,不是人的眼睛可以察看或衡量的.「祂本有神的形像」的原意是指主耶穌具有神的本性──祂就是神,本性若損失就不再是神了.所以,主耶穌本是神的形像,也永遠是神的形像,因祂是完全的神,具有完全的形像,而非部份,祂更沒有放下自己的任何本性或神性.
\newline
\newline
「反倒虛己」(7節)的意思就是「倒空自己」,主倒空了些甚麼呢?有人說,主道成肉身後,祂把無所不知,無所不能,無所不在等神性倒空了,於是他在世上受到限制.主耶穌若是真的倒空了部份神性,那麼祂再也不是神,怎能稱為我們的救主呢?祂有資格完成救贖,就是因為祂是完全的神,又是完全的人,這是絕對不可妥協的真理.「虛己」這個字在新約中只有保羅採用過五次,其用法(羅四14；林前一17；林後九3,15；腓二7)全是指某些東西失掉了功效,而不是指倒空某些具體的東西,換句話說,主耶穌不是倒空自己,乃是謙卑自己,祂怎麼謙卑自己呢?祂「不以自己與神同等為強奪的」(6節以下),而「與神同等」 在文法上是副詞,用作形容祂存在的方式.主耶穌本是神的形像,具有神的榮耀,權柄和地位,並受天使的敬拜；但道成肉身後,祂就不在榮耀中表現自己,祂改變的是存在的方式,而非祂的本性.這就好像中國古時的皇帝微服出巡一樣；雖然祂脫下朝服,穿上民裝,但卻絲毫沒有改變他原有的身份.同樣,主耶穌暫時放棄天上的存在方式,甘心樂意地放下自己應當享受的權柄,不在榮耀中表現自己；這就是不「強奪」的意思,也是道成肉身的目的.
\newline
\newline
深願今天我們也能效法主耶穌的榜樣,為了福音的緣故,把應當享受,擁有的權利放棄.其實,很多時候我們喜歡抓緊自己的權利,今天社會上一些只顧自己利益的請願,罷工等風氣也漸漸影響著教會及福音機構；希望我們真能以基督的心為心,以祂作我們的榜樣,甘心放下自己應得的權利.馬太福音第四章記載主耶穌經過四十晝夜的禁食；撒旦前來試探祂,要祂用自己的方法去滿足人的基本需要,因為這是人應當有的享受和權利.但主的回答:「人活著不是單靠食物」,意思是指既然神的旨意叫祂那時不吃；祂就甘心放棄自己的權利,為要完成道成肉身拯救人類的目的.
\newline
\newline
今日,我們願意為主犧牲些甚麼呢?現代社會物質豐裕,我們不需要作任何犧牲,就能完成一些表面的責任.例如不用為了奉獻而缺衣忍餓；有人以為通宵祈禱或放棄精彩電視節目來聽道就是為主犧牲,其實我們還不明白甚麼叫做犧牲:有位兄姊為了到落後的地方傳福音而甘願放棄外國的護照,過著與當地人完全相同的生活；結果在短短兩年時間內成立了三個小家庭教會.這種為福音的緣故而付上代價,才算是真正的犧牲.但這種放棄物質的享受卻遠不及主的犧牲,因祂不僅放棄榮耀的地位,而且還過者不被尊重的生活.你為主放棄了些甚麼呢?
\newline
\newline
(二)取了奴僕的形像
\newline
\newline
主耶穌不但「反倒虛己」放棄了神榮耀的存在方式,而且還「取了奴僕的形像,成為人的樣式」(7節)祂為了完成救贖工作,在神性上又加上人性,加上人性一切的限制,把不必要的擔子加在自己身上,榮耀的神成為軟弱的人；這就好比一個年青力壯的大漢病倒了,吃喝樣樣要別人照料,其心情可想像是多麼的難受.主耶穌為了拯救人類而道成肉身,降生在卑微的馬槽；事事要人服侍,樣樣受限制.然而祂卻甘心忍受這些痛苦.我們甘心為主背負十字架嗎?「十字架」不單是指一些考驗或苦難,便是指甘為福音而受苦,甚至犧牲自己.初期教會教父坡利甲,就是為了堅持信仰而不向該撒像下拜；結果被羅馬政府綁在樹上活活燒死.他說: 「八十多年來,我的主從來沒有背棄我,我現在又怎能為保存生命而否認主呢!」宣教歷史上還有許許多多為主,為福音的緣故而甘願受苦和犧牲的見證.我們是否願意效法主完全謙卑的榜樣,以基督的心為心?雖然我們不能像主那樣放棄神的榮耀而道成肉身；但我們應有主的心志──為了完成神的旨意,為了成為神合用的器皿；甘願犧牲自己應有的權利,負起艱苦的擔子,完成神所交托的使命!
\newline
\newline


\newpage

\section{第62屆~港九培靈會~鮑會園~第五講}
\label{sec:1080}
\textbf{工作的目的──作成得救的工夫}
\newline
\newline
連結: \href{https://www.hkbibleconference.org/session-message/view/1080}{\texttt{https://www.hkbibleconference.org/session-message/view/1080}}
\newline
\newline
\hyperref[sec:1079]{< < < PREV SERMON < < <}
~
\hyperref[sec:toc]{[返目錄]}
~
\hyperref[sec:1081]{> > > NEXT SERMON > > >}
\newline
\newline
腓立比書 2:12-2:18
\newline
\begin{longtable}{cl}
\hline
\hline
章節 & 經文 (和合本修訂版)\\
\hline
2:12 & \begin{tabularx}{0.7\textwidth}{X} 我親愛的，這樣看來，你們向來是順服的，不但我在你們那裡，就是我現在不在你們那裡的時候更是順服的，就當恐懼戰兢完成你們自己得救的事； \end{tabularx} \\ \\ \relax
2:13 & \begin{tabularx}{0.7\textwidth}{X} 因為是神在你們心裡運行，使你們又立志又實行，為要成就他的美意。 \end{tabularx} \\ \\ \relax
2:14 & \begin{tabularx}{0.7\textwidth}{X} 你們無論做甚麼事，都不要發怨言起爭論， \end{tabularx} \\ \\ \relax
2:15 & \begin{tabularx}{0.7\textwidth}{X} 好使你們無可指責，誠實無偽，在這彎曲悖謬的世代作神無瑕疵的兒女。你們在這世代中要像明光照耀， \end{tabularx} \\ \\ \relax
2:16 & \begin{tabularx}{0.7\textwidth}{X} 將生命的道顯明出來，使我在基督的日子得以誇耀我沒有白跑，也沒有徒勞。 \end{tabularx} \\ \\ \relax
2:17 & \begin{tabularx}{0.7\textwidth}{X} 我以你們的信心為供獻的祭物，我若被澆獻在其上也是喜樂，並且與你們眾人一同喜樂。 \end{tabularx} \\ \\ \relax
2:18 & \begin{tabularx}{0.7\textwidth}{X} 你們也要照樣喜樂，並且與我一同喜樂。 \end{tabularx} \\ \\
[1ex]
\hline
\hline
\end{longtable}
腓立比書從第二章開始就解釋如何在生活中追求完成神的旨意,在這彎曲悖謬的世代將生命的道表明出來.以生命來傳揚神的道,這不單是保羅對腓立比信徒的訓勉；也是我們今日的使命.
\newline
\newline
本段經文看似很特別,與保羅其他的教訓並不相符:「當恐懼戰兢,作成你們得救的工夫」(二12下)得救不是本乎恩又因著信嗎?人怎能作成得救的工夫呢?保羅已在加拉太書再三提醒信徒「既靠聖靈入門,如今是靠肉身成全麼」(加三3)意思是指救恩完全是神的工作,人在「得救的工夫」沒有任何的功勞,一點也幫不了神的忙；人不能藉著任何方式來賺取救恩,甚麼時候人若要以自己的善事,功勞自誇,或者要靠律法稱義,那麼他就是與基督隔絕,從恩典中墜落了(加五4).既然如此,這裡「得救的工夫」又指的是甚麼呢?
\newline
\newline
(一)見證主
\newline
\newline
首先,我們先思想一下「救恩」的真正意義.基本上來說:「救恩」包含與神,與人兩方面的關係,而且這兩種關係是相對的.在與神的關係上來說:「救恩」是指罪惡得赦免,將來不受刑罰,得著永生的盼望等等.這方面的功勞完全是藉著耶穌作成的；人在神面前只是蒙恩的罪人而己,沒有任何的工夫可做,也沒有任何的功勞可誇；但在與人的關係上來說,「救恩」則是指重生得救的表現；或者說一個得救的人,在人面前活出得救的見證來.在生活中追求聖潔,吸引人到主面前來,這就是人應當努力完成的工夫或使命.保羅在提多書中說:「祂為我們捨了自己,要贖我們脫離罪惡」(二14)這是指第一方面說的,然後他又補充:「潔淨我們,特作自己的子民,熱心為善」就是指第二方面說的.既使我們蒙了救恩,就當有蒙恩的表現,許多時候我們明白第一方面與神的關係,但卻忽略第二方面對人應付的責任.保羅就是勉勵信徒要在這方面積極追求「恐懼戰兢作成得救的工夫」.
\newline
\newline
有的時候,我們過份注重將來的救恩,來生的福氣,而忽略現在的救恩,今生的福氣；結果造成雖然得救,但屬靈生命卻很幼嫩,生活異常貧乏的現像.這也是羅馬書第八章所說的,我們要追求得著「兒子的名份」.中國人不大著重名份,認為「有名無實」沒有用；但聖經中的名份卻是代表著屬靈的福氣,名份若失掉了,福氣也沒有了.舊約以掃把長子名份賣給雅各後,雖然他事實上仍是長子,但卻失去長子應有的福份.今天我們藉著耶穌的救恩得著「兒子的名份」,是否享受神兒子的福氣和喜樂呢?同樣,很多基督徒雖然已經得著救恩的事實；但卻沒有享受到救恩的福氣.整天為家庭,工作掛慮,擔憂,心中沒有平安,臉上也沒有喜樂.......保羅提醒我們要在這方面樍極地追求,不但有得救的生命,更要有得救的表現,「恐懼戰兢,作成得救的工夫」就是指得著救恩的福氣,追求聖潔的生活.
\newline
\newline
為甚麼我們的生命若不「恐懼戰兢」地努力追求,隨時會失去聖潔的生活及見證呢?主要的原因有兩方面:一是外在的敵人,二是內在的敵人.
\newline
\newline
(1)外在的撒旦:撒旦對基督徒的攻擊是厲害的,牠隨時找機會來試探,控告我們；牠知道我們的強弱之處,要在我們的軟弱上叫我們跌倒.有的人在某方面可能特別剛強(例如:不吸煙),但不等於在其他方面也是同樣剛強:有的人很溫柔,但缺乏愛心；有的人有愛心,但缺乏謙卑；有的人謙卑,但缺乏容忍的心.不要以為生命的某一方面得勝了,就忽略了其他方面的儆醒,要隨時小心,免陷入撒旦的試探裡.
\newline
\newline
(2)內在的老我:保羅做了幾十年的基督徒,還這樣慨嘆道:「我所願意的善,我反不作.我所不願意的惡,我倒去作.若我去作所不願意作的,就不是我作的,乃是住在我裡頭的罪作的」.(羅七19-20)可見老我是多麼可怕!許多時候,我們的事奉工作表面上確是做得很好；但內心卻沾沾自喜,驕傲不已；甚至以此來顯出別人的不是,高抬自己.人心比萬物都詭詐,壞到極處；我們的工作表現是很好,但動機卻是錯誤的.要隨時防備老我的攻擊,不要給它任何表現自己的機會；「恐懼戰兢作成得救的工夫」.
\newline
\newline
(二)順服主
\newline
\newline
「這樣看來,我親愛的弟兄,你們既是常順服的；不但我在你們那里,就是我如今不在你們那里,更是順服的.」(二12上)這段經文的原意不是稱讚,乃是命令；而在中文繙譯裡有兩個「順服」,但在原文中只有一個,因此這樣的繙譯可能更貼切一點:「你們應當順服,我在那裡是要這樣；我不在時也要這樣」.順服不是做給別人看的,其目的是建造自己的屬靈生命,既是為了建立自己的生命；那麼我們就不應只做表面的工夫,只有眼前的事奉.在世俗社會,許多人為了得到上司的欣賞,就努力做好表面的工夫.因此,作工的需要督工,沒有監督就不作工；學生需要舍監,沒有監察就不守規矩.......這些都是世人做事的態度, 基督徒卻不應該這樣；我們當「恐懼戰兢作成得救的工夫」,無論在任何環境下,都要過著完全順服主的生活；叫我們的生命因此得以建立.
\newline
\newline
(三)渴慕主
\newline
\newline
「因為你們立志行事,都是神在你們心裡運行,為要完成祂的美意.」(一13)這裡的「立志行事」就是指追求救恩,建造屬靈生命的工夫.「立志行事都是神在你們心裡運行」表明是神主動來感動我們去追求祂,不是單靠我們的決心,或者自己的力量,為的是要完成祂的美意.在得救的事上,我們不能靠自己的功勞；同樣,在屬靈追求的事上,我們亦不能靠自己的努力；乃是主感動我們的心.雅歌書第一章很生動地描寫所羅門王與書拉密女的關係；這裡的君王可以代表耶穌基督,而女子則代表追求的基督徒.當書拉密女面對面看見所羅門王的美好,內裡就產生愛慕,追求的心,但當她願意與王親近時,卻又發現環境,家庭成了她的攔阻.於是他對王說:「願你吸引我,我就快跑跟隨你.」屬靈生命的追求也是一樣,是神藉聖靈主動吸引,感動我們,叫我們立志行事都成為神的美意.使我們的生命能完成神交托的使命,因為聖靈的能力在我們心裡成了追求的動力,生活的力量.
\newline
\newline
既然聖靈是我們的力量,甚至遠勝過我們的力量；那麼我們的生活還得努力追求嗎?聖靈作事有一定的原則和方法:雖然祂能幫助我們過聖潔的生活,但祂不勉強我們,不違背祂自己的旨意.聖靈若任意發揮其力量,人是擔當不了的,正如使徒行傳第五章的亞拿尼亞和撒非喇因欺哄聖靈而遭殺一樣.聖靈所用的方法是感動我們的心,當我們願意把主權交出多少,順服多少,聖靈的力量也就發揮多少.舊約創世記中,羅得在所多瑪,蛾摩拉城接待兩位天使；而城裡的壞人就來攻破房門,就在羅得跟他們理論之際,天使把他拉到後面,結果問題就解決了.請留意當羅得願意順服天使的命令而退到後面,天使就發揮其能力；否則單靠羅得一己之力來抵擋,就一定失敗.
\newline
\newline
我們屬靈生命追求聖潔也是一樣,甚麼時候我們願意完全順服聖靈的掌管,甚麼時候就能經歷聖靈的力量.聖靈不僅重生我們,更要賜予屬靈的福氣,幫助我們經過敵人的攻擊和老我的試探；在我們的生命作成得救的工夫,成就神的美意.
\newline
\newline


\newpage

\section{第62屆~港九培靈會~鮑會園~第六講}
\label{sec:1081}
\textbf{最高的願望 — 認識基督和祂的大能}
\newline
\newline
連結: \href{https://www.hkbibleconference.org/session-message/view/1081}{\texttt{https://www.hkbibleconference.org/session-message/view/1081}}
\newline
\newline
\hyperref[sec:1080]{< < < PREV SERMON < < <}
~
\hyperref[sec:toc]{[返目錄]}
~
\hyperref[sec:1082]{> > > NEXT SERMON > > >}
\newline
\newline
腓立比書 3:1-3:11
\newline
\begin{longtable}{cl}
\hline
\hline
章節 & 經文 (和合本修訂版)\\
\hline
3:1 & \begin{tabularx}{0.7\textwidth}{X} 末了，我的弟兄們，你們要靠主喜樂。我把這些話再寫給你們，對我並不困難，對你們卻是妥當的。 \end{tabularx} \\ \\ \relax
3:2 & \begin{tabularx}{0.7\textwidth}{X} 應當防備犬類，防備作惡的，防備妄自行割的。 \end{tabularx} \\ \\ \relax
3:3 & \begin{tabularx}{0.7\textwidth}{X} 因為真受割禮的，就是我們這藉著神的靈敬拜、以基督耶穌為誇耀、不依靠肉體的。 \end{tabularx} \\ \\ \relax
3:4 & \begin{tabularx}{0.7\textwidth}{X} 其實，我也可以靠肉體；若是別人以為他可以依靠肉體，我更可以。 \end{tabularx} \\ \\ \relax
3:5 & \begin{tabularx}{0.7\textwidth}{X} 我出生後第八天受割禮；我是以色列族、便雅憫支派的人，是希伯來人所生的希伯來人。就律法說，我是法利賽人； \end{tabularx} \\ \\ \relax
3:6 & \begin{tabularx}{0.7\textwidth}{X} 就熱心說，我是迫害教會的；就律法上的義說，我是無可指責的。 \end{tabularx} \\ \\ \relax
3:7 & \begin{tabularx}{0.7\textwidth}{X} 只是我先前以為對我是有益的，我現在因基督的緣故而當作是有損的。 \end{tabularx} \\ \\ \relax
3:8 & \begin{tabularx}{0.7\textwidth}{X} 不但如此，我已把萬事當作是有損的，因我以認識我主基督耶穌為至寶。我為他已經丟棄萬事，看作糞土，為要贏得基督， \end{tabularx} \\ \\ \relax
3:9 & \begin{tabularx}{0.7\textwidth}{X} 並且得以在他裡面，不是有自己因律法而得的義，而是有信基督的義，就是基於信，從神而來的義， \end{tabularx} \\ \\ \relax
3:10 & \begin{tabularx}{0.7\textwidth}{X} 使我認識基督，知道他復活的大能，並且知道和他一同受苦，效法他的死， \end{tabularx} \\ \\ \relax
3:11 & \begin{tabularx}{0.7\textwidth}{X} 或許我也得以從死人中復活。 \end{tabularx} \\ \\
[1ex]
\hline
\hline
\end{longtable}
在(腓三1-11)的經文中,保羅繼續解釋他生命追求的方向:好像明光照耀,將生命的道表現出來.這裡保羅用兩個很強烈的對照,來解釋自己得救前後所追求的方向.在得救前,他是憑自己的功勞,靠行律法而稱義；得救後,他把以前所依靠,甚至誇口的,都當作糞土丟棄了.放下自己,對我們來說是多麼困難；但只有這樣,我們才能得著基督.我們不能一邊恃著自己的功勞,一邊找緊主的恩典；否則生命的追求,一定不會有真正的果效.
\newline
\newline
(一)認識基督
\newline
\newline
保羅在第10節再次表明生命追求的目的:「使我認識基督,曉得祂復活的大能；並且曉得和祂一同受苦,效法祂的死.」保羅為主工作了三十多年,在未寫腓立比書之前,保羅已在亞洲的西部及歐洲的部份地區建立了許多教會；而且他也完成了哥林多前後書,羅馬書,加拉太書,甚至以弗所書的寫作；保羅寫了這麼多書信,怎能還說要認識基督?難道他是沒有真正認識基督嗎?如果連保羅都說不認識基督,恐怕我們也不敢自稱認識基督了.
\newline
\newline
其實,「認識」這個字在聖經中有三種不同的用法,每種用法均有不同的意義.或者說藉著不同的「認識」方法可以得到不同的知識.第一種是稱為間接的認識:我們對歷史的知識多是屬於這一類.雖然我們沒有經過歷史事件,但從歷史資料中仍可得到這種知識.第二種是藉著思想而來的知識:在生活中有許多這樣的例子,很多事情想通了就明白,想不通無論怎樣解釋都不明白.甚麼叫「第四度空間」呢?我們明白單方,立方的樣子,但四次方,五次方......的樣子又是怎樣的呢?這全都是不能想像；也無法經驗到的知識.但假若我們肯動腦筋,仍然可以得到這類知識,聖經中許多神學問題也是屬於這一類.第三類是憑著經驗而來的知識:這種知識只有經驗過才能明白.假如有人從未吃過橙,即使他讀過形容橙味的文章；或者想像其味道是怎樣的,但若不是親自吃過；他仍然還是不知道橙的真正味道.保羅所說的「認識基督」,就是這種的知識和經驗.
\newline
\newline
保羅過去對主已有第一種的知識:他與主的門徒有很好的交通,甚至有人說他親眼見過主,但不肯相信.保羅也有第二種的知識:他思想過許多神學的問題,也知道很多屬靈的奧秘；明白道成肉身,三位一體,救贖,教會等真理.甚至藉著神的啟示,他也能著書立說.但是,他單有間接,或者想通的知識還是不夠,他還要在生命中親身經驗所認識的主耶穌.經歷祂的大能和大愛,經歷祂的聖潔,無所不在......,「使我認識基督」,就是保羅生命追求的方向.
\newline
\newline
今天,我們也像保羅一樣追求這種的知識,不單有聽來的,學來的或其他頭腦上的知識和理論；更重要的是在生命中經驗所明白的真理,經歷主耶穌.只有這樣,我們的生命才能在考驗中站立得住；才能面對各種的挑戰.約翰福音第九章記載主耶穌醫好生來瞎眼的人；而法利賽人就要控告祂在安息日治病,但這曾是瞎眼的人卻確信主醫治的大能,認定自己經驗中的知識,而且知識是永遠不會失掉的.因為他在生命中真正經驗了主耶穌,我們要像瞎子一樣追求這樣的知識.單有頭腦的知識是不夠的,甚至因為有些人知道了一些聖經知識,內心反而變得更加剛硬,更加拒絕,知識並不能叫人真正歸回主,只有生命中更深入地經驗主,更真實地追求「認識基督」.
\newline
\newline
(二)曉得祂復活的大能
\newline
\newline
保羅不但要追求「認識基督」,更要「曉得祂復活的大能」.這裡指的不是主耶穌從死裡復活的大能；因為這種能力只有在主的身上才能彰顯.我們永遠也沒有資格經歷這種大能；這裡的「大能」乃是指我們在生活中靠主勝過罪惡捆綁的能力.基督若沒有復活,我們就仍在罪裡(林前十五17)；但基督復活了,就能幫助我們脫離罪惡的勢力.
\newline
\newline
基督徒在生活中經常要面對罪惡的挑戰,可惜有的時候,我們在罪惡中失敗了.基督徒信主以後,即使他有脫離罪惡的經驗,享受到罪擔得釋放的喜樂；但罪惡還是常常纏繞著我們,找機會來試探我們,要叫我們在試探中跌倒失敗.
\newline
\newline
約翰福音第十一章記載主耶穌用超然的能力使拉撒路從死裡復活,當拉撒路自墓中走出來的時候,雖然他有生命,但卻不能生活；因為「手腳裹著布,臉上包著頭巾.」所以耶穌吩咐人說:「解開,叫他走!」(十一44)這是非常生動,強而有力的比喻；當我們信主的一剎那,我們確是有了新的生命,但也許我們仍不能過完全得救的生活；因為還有許多罪惡像裹屍布一樣纏繞著我們.許多人明明知道發脾氣是不對的,但卻被自己的脾氣勝過而大發雷霆,事後當然也感到難過；並立志悔改,結果呢?正如保羅所說的:「立志為善由得我,只是行出來由不得我」.若是靠自己的能力,我們必然發現自己不斷重蹈覆轍,生活在罪惡的捆綁之下.因為我們沒有能力掙脫罪惡的權勢,只有主耶穌復活的大能,砍斷一切纏繞我們的罪惡；拯救我們脫離罪惡的力量,依靠聖靈的能力,就是過得勝生活的秘訣.甚麼時候我們不靠自己的能力而依靠主復活的大能,我們就能得著勝過罪惡的力量,追求聖潔的生活.
\newline
\newline
(三)和祂一同受苦,效法祂的死
\newline
\newline
為甚麼基督徒追求長進,卻反要受苦呢?基本的原因就是當我們追求更深地認識主時,撒旦一定會更加厲害地攻擊我們；目的就是要叫我們在軟弱,試探中失敗.因此,我們愈追求長進,撒旦的攻擊也愈厲害,我們內心的掙扎也愈大；遇著這種情況,我們不要灰心,難過；因為內心有掙扎就是表示我們生命中有力量要抵擋撒旦的攻擊.相反,我們若在試探中失敗而心裡毫無掙扎的話；那麼可能要更加小心,也許我們的生命對撒旦來說毫無威脅.牠根本不用來試探,攻擊我們,而我們也感覺不到撒旦的試探；又或者我們的生命中已失去抵擋撒旦的力量.既然沒有力量,自然也不用掙扎.因此,生命愈長進,生活中可能面對愈大,愈多的掙扎； 所以我們要效法主的受苦心志,經驗祂復活的大能,真正追求更認識主.
\newline
\newline
另一方面,保羅不僅要與主同受苦,而且還要效法祂的死,這是甚麼意思呢?(約翰福音十二32)記載:「我若從地上被舉起來,就要吸引萬人來歸我.」原來主耶穌受死的目的,就是要把人帶到神面前來；門徒也要效法主這種受死的心志,藉著我們對主的順服,把人帶到主的面前來.
\newline
\newline
在香港這個時代,我們更多需要藉著順服主的生活,讓人更加認識我們的主.若我們不願意順服主,或者不肯為順服主而付上代價,這是香港教會的悲劇!從 七十年代到八十年代初期,報考本港神學院的青年人很多,幾乎每一間神學院均有額滿的現象；最近二,三年的報考人數卻顯著下降了.是因為政治的局勢冷卻了我們的事主熱忱?還是前途,環境的不安動搖了我們跟隨主的心志?保羅生命追求的目的,成為我們今天每一個人的挑戰,特別是蒙主呼召的青年人,我們是否願意完全順服神的旨意,在生命中更認識基督；經驗祂復活的大能,勝過罪惡和試探；並效法主受苦,受死的心志,為福音的緣故付上代價,目的就是要把人帶到神面前來.
\newline
\newline


\newpage

\section{第62屆~港九培靈會~鮑會園~第七講}
\label{sec:1082}
\textbf{追求的目標 — 凡事作完全人}
\newline
\newline
連結: \href{https://www.hkbibleconference.org/session-message/view/1082}{\texttt{https://www.hkbibleconference.org/session-message/view/1082}}
\newline
\newline
\hyperref[sec:1081]{< < < PREV SERMON < < <}
~
\hyperref[sec:toc]{[返目錄]}
~
\hyperref[sec:1083]{> > > NEXT SERMON > > >}
\newline
\newline
腓立比書 3:12-3:16
\newline
\begin{longtable}{cl}
\hline
\hline
章節 & 經文 (和合本修訂版)\\
\hline
3:12 & \begin{tabularx}{0.7\textwidth}{X} 這不是說我已經得著了，已經完全了；而是竭力追求，或許可以得著基督耶穌所要我得著的。 \end{tabularx} \\ \\ \relax
3:13 & \begin{tabularx}{0.7\textwidth}{X} 弟兄們，我不是以為自己已經得著了；我只有一件事，就是忘記背後，努力面前的， \end{tabularx} \\ \\ \relax
3:14 & \begin{tabularx}{0.7\textwidth}{X} 向著標竿直跑，要得神在基督耶穌裡從上面召我來得的獎賞。 \end{tabularx} \\ \\ \relax
3:15 & \begin{tabularx}{0.7\textwidth}{X} 所以，我們中間凡是成熟的人，總要存這樣的心；若在甚麼事上存別樣的心，神也會把這些事指示你們。 \end{tabularx} \\ \\ \relax
3:16 & \begin{tabularx}{0.7\textwidth}{X} 然而，我們達到甚麼地步，就當照這個地步行。 \end{tabularx} \\ \\
[1ex]
\hline
\hline
\end{longtable}
前面提到保羅所追求的方向,於是放棄一切來認識基督,曉得祂的大能,效法祂的死；目的就是在靈裡更深的經驗主耶穌,使屬靈生命不斷的長進,成熟.保羅不僅認定生命追求的方向；也清楚追求的目標,這亦是本段經文的討論的主題.
\newline
\newline
(一)竭力追求
\newline
\newline
生命需要不斷地長進,不長進的生命一定有問題；肉體的生命如是,屬靈生命亦是一樣.「這不是說,我已經得著了,已經完全了」(三12)難道保羅生命追求上幾十年仍沒有「得著」,還不夠「完全」嗎?事實確是如此,特別在羅馬書第七章,保羅坦白地承認他心裡知道甚麼是好的,也願意遵行神的旨意,只可惜行不出來.這裡給予我們一個非常重要的真理,就是基督徒雖然要在屬靈生命上不斷長進；但在今生卻永遠達不到完全的地步.因此,若有人認為自己已經完全了, 可能他還不明白「完全」的真正意義,或者他把「完全」的標準降低了,誤以為做到某些生活上的表現就達到完全的地步,這都是真理的誤解而造成的.
\newline
\newline
保羅很清楚生命追求的目標,也明白衡量生命的標準,所以他知道自己還沒有完全.(腓三12)可以譯作:「得著基督耶穌所要我得的」.換句話說,是基督耶穌為我的生命定下目標,要我努力追求達到這標準；我們要用主的標準來衡量自己,而不是自己定一個容易達到的目標,並以達到這標準自以為滿足.許多基督徒的生命追求就是是這樣,為自己定一個最起碼的標準；不偷,不搶,沒有在人前羞辱主名,星期日去主日崇拜,又有十一奉獻......以為達到這標準已差不多了.在屬靈生命上不要更深的追求.但保羅提醒我們在生命的追求上,絕對不能馬馬虎虎地就自覺滿足；我們要達到基督耶穌要我們達到的標準,只有這樣的生命才能更加完全.要達到主的標準,首先我們的生命就應該對主完全的順服,生活要完全的聖潔.若果我們有一個這樣的標準,那麼前面就永遠有清晰的挑戰和追求的目標.
\newline
\newline
(二)完全人
\newline
\newline
既然我們每一個人都要竭力地追求長進,甚至連保羅也自認還沒有達到完全的地步；那麼經文第15節是指著誰說的呢?「所以我們中間凡是完全人,總要存這樣的心」(三15上)不同的解經家對這節經文均有特別的解釋,在瞭解「完全人」是指著誰說的之前,我們先要明白「完全人」的準確意思.
\newline
\newline
「完全人」這個字在新約中使用過很多次,其字根與第12節的「完全」是一樣的,但兩者的句子構造卻不一樣,因此就顯出「完全人」是具有特別的意義: (來五14)「惟獨長大成人的,才能吃乾糧」中的「長大成人」,(林前十四)「在惡事上要作嬰孩,在心志上總要作大人」中的「大人」(弗四13)「在真道上同歸於一,認識神的兒子,得以長大成人,滿有基督長成的身量.」中的「長大成人」,都和腓立比書的「完全人」是一樣的意思.總括來說,這個字在許多經文中都是指屬靈生命的成長,因此腓立比書的「完全人」若譯作「長大成人」,意思就清楚了；而(三15)則可解作:「我們中間凡是長大成熟的人,總要有這樣的追求心志」.簡言之,保羅強調的是要在屬靈生命上追求作長大成熟的人,而不是指生活中達到百分之百完全的人；特別在哥林多前書第三章,保羅勉勵信徒不要在屬靈上作嬰孩,應該追求長大成熟.
\newline
\newline
根據希伯來書第五,六章的形容,再加上對一般嬰孩的認識,基本上,我們可以明白屬靈嬰孩的表現是怎樣的.首先,嬰孩只顧自己的需要,他餓了就哭,但別人挨餓他從不關心；此外,嬰孩也不知如何照顧自己的需要,即使冰箱裡有牛奶,他也不懂得拿來喝.屬靈嬰孩的表現也是一樣,凡事只為自己著想,不知道如何在靈裡追求；也不曉得如何在神的話語得餵養.還有,嬰孩喜歡表現自己,若沒有表演的機會就不高興.同樣,今天許多教會中的屬靈嬰孩,事奉作工是為要得人的稱讚；若沒有機會表現自己或得不到稱讚,心裡就不高興.這些會是屬靈嬰孩的表現,保羅鼓勵我們要作長大成熟的人.
\newline
\newline
(三)長大成熟的人
\newline
\newline
怎樣才可以成為長大成熟的人呢?首先最重要的是,要把長大成熟的根基建立在屬靈的真理上,在生活中不斷地追求和遵行神的話.
\newline
\newline
(1)追求的心:「若在甚麼事上,存別樣的心,神也必以此指示你們.」(三15下)只有存追求主的心,即使有時方向錯了,主的恩典仍會幫助,保守我們,糾正我們走回正路.但我們若在任何事上存別樣的心,神就必定指示我們,改正我們的錯誤.「指示」是個很重要的字,在新約裡通常譯作「啟示」即神把真理指示給我們.有了真理,我們的生命才能長大成熟；追求的方向才是正確.因此,我們追求的根基必須單單建立在真理上,不能靠人的恩賜或經驗,更不能憑好聽的理論或情感；否則的話,我們的生命必定會出現許多的問題.
\newline
\newline
今日的世代充滿了許多人類的思想和理論,要分辨哪些是真理確實不容易,如同約翰福音第十二章記載主耶穌聽見神向祂說話,但站在祂旁邊的門徒卻以為是雷聲(約十二29)因為他們不能分辨是神的聲音,還是打雷的聲音.另外,舊約也記載大衛王蒙神帶領與敵人作戰,由於敵人十分詭詐,所以神吩咐大衛:「等到天晚時,你聽見樹梢上有腳聲,我就指示你該怎樣行.」山谷裡的樹梢常被風吹動而發出聲音,如何聽見「樹梢上有腳步的聲音」呢?是的,別人分辨不出來的,但大衛卻能做得到.為甚麼主耶穌能聽見神的聲音,而其他門徒及兵丁卻聽不見呢?主要的原因是因為他們的屬靈耳朵在平時訓練有素,他們聽慣了神的聲音,並且常存順服的心.今天我們也要訓練屬靈的耳朵能聽見神的聲音,無論是讀書,聽道,走路或者工作時,都要追求聽神的聲音.平時常訓練的耳朵,加上順服的心志,就 能幫助我們分辨神的聲音,這樣我們也能按著真理來追求長大成熟.
\newline
\newline
(2)遵行的心:「然而我們到了甚麼地步,就當照著甚麼地步行.」(三16)在神面前領受真理到了甚麼地步,就要按著甚麼地步行,這是屬靈生命長進的重要原則.我們若按著所明白的真理去行,神就啟示更深的真理；我們若再按這真理去行,那麼我們的屬靈生命就一步一步向前邁進了.但假若我們領受了真理, 卻又不願意遵行真理,或生活的表現也跟不上,那麼這領受就對我們毫無幫助；這樣的話,神也不再給我們有新的指示了.因此,屬靈的知識與生活的追求是互相配合的,真理的知識有多少,生活的改變也要有多少；接著真理的啟發,一步一步的遵從而行,這樣我們的屬靈生命就自然漸漸長進,作長大成熟的「完全人」.
\newline
\newline


\newpage

\section{第62屆~港九培靈會~鮑會園~第八講}
\label{sec:1083}
\textbf{生活的盼望 — 作天上的國民}
\newline
\newline
連結: \href{https://www.hkbibleconference.org/session-message/view/1083}{\texttt{https://www.hkbibleconference.org/session-message/view/1083}}
\newline
\newline
\hyperref[sec:1082]{< < < PREV SERMON < < <}
~
\hyperref[sec:toc]{[返目錄]}
~
\hyperref[sec:1084]{> > > NEXT SERMON > > >}
\newline
\newline
腓立比書 3:17-4:3
\newline
\begin{longtable}{cl}
\hline
\hline
章節 & 經文 (和合本修訂版)\\
\hline
3:17 & \begin{tabularx}{0.7\textwidth}{X} 弟兄們，你們要一同效法我，也當留意看那些效法我們榜樣的人。 \end{tabularx} \\ \\ \relax
3:18 & \begin{tabularx}{0.7\textwidth}{X} 因為，我屢次告訴你們，現在又流淚告訴你們：許多人行事是基督十字架的仇敵。 \end{tabularx} \\ \\ \relax
3:19 & \begin{tabularx}{0.7\textwidth}{X} 他們的結局就是滅亡。他們的神明是自己的肚腹；他們以自己的羞辱為光榮，專以地上的事為念。 \end{tabularx} \\ \\ \relax
3:20 & \begin{tabularx}{0.7\textwidth}{X} 我們卻是天上的國民，並且等候救主，就是主耶穌基督從天上降臨。 \end{tabularx} \\ \\ \relax
3:21 & \begin{tabularx}{0.7\textwidth}{X} 他要按著那能使萬有歸服自己的大能，把我們這卑賤的身體改變形狀，和他自己榮耀的身體相似。 \end{tabularx} \\ \\ \relax
4:1 & \begin{tabularx}{0.7\textwidth}{X} 我所親愛、所想念的弟兄們，你們就是我的喜樂，我的冠冕。我親愛的，你們應當靠主站立得穩。 \end{tabularx} \\ \\ \relax
4:2 & \begin{tabularx}{0.7\textwidth}{X} 我勸友阿蝶和循都基要在主裡同心。 \end{tabularx} \\ \\ \relax
4:3 & \begin{tabularx}{0.7\textwidth}{X} 我也求你這真實同負一軛的，要幫助這兩個女人，因為她們在福音上曾與我、革利免和我其餘的同工一同勞苦，他們的名字都在生命冊上。 \end{tabularx} \\ \\
[1ex]
\hline
\hline
\end{longtable}
保羅在前面的經文中提醒我們,屬靈的生命要有追求的方向和目標,但單有靈裡的追求還不夠,這裡他鼓勵我們還要在實際的生活中把屬靈生命表現出來.
\newline
\newline
保羅在這段經文中,用了兩幅強烈對比的圖畫,來描繪兩種截然不同的基督徒.第一幅圖畫是第三章的(17-19節),保羅所形容的是一些明白真理但生活見證卻不好的人,對於這些人,保羅指出他們的結局就是沉淪,而第二幅圖畫則是第20節所指的天上的國民.這些人即使生活上或有軟弱,但他們生命的追求方向正確了,其生活仍然在人前可成為有力的見證.在腓四章中保羅提到幾個人的名字,他們事實上是天上的國民,雖然在生活,事奉中有軟弱；但他們的見證仍然被神所悅納.因此,我們下面將詳細比較這兩幅圖畫.
\newline
\newline
(一)十字架的仇敵
\newline
\newline
「因為有許多人行事,是基督十字架的仇敵.我屢次告訴你們,現在又流淚地告訴你們:他們的結局就是沉淪；他們的神就是自己的肚腹,他們以自己的羞辱為榮耀,專以地上的事為念.」(三18-19)雖然保羅這裡沒有清楚指明「十字架的仇敵」是誰；但大概是指著具體的人說的.而這些人不會是一般不認識基督的人,很可能是指一些明白真理,又領受了某些真理的基督徒,他們不會公開地拒絕十字架的道理；但其生活表現卻一點也不像基督徒.經文第19節就是論及他們幾方面的表現,但在解釋上卻有點困難,因為這裡保羅好像是指著三件事情來說:
\newline
\newline
(一)他們的神就是自己的肚腹；
\newline
\newline
(二)他們以自己的羞辱為榮耀；
\newline
\newline
(三)專以地上的事為念.
\newline
\newline
但這三件事情從句子的構造上來看,並不是平衡的；整段的主要意思是在第三句話上,這句話包括了這些基督徒的特色和生活表現,專以地上的事為念；成了基督十字架仇敵的一個記號.他們並非拒絕十字架的真理,也不像哥林多教會那樣犯姦淫的罪,也沒有像哥羅西教會那樣追求世上的學問,更沒有敬拜偶像；然而, 他們整個人的心思所關注的全是地上的事.生命所追求的也是地上的事,「專以地上的事為念」.
\newline
\newline
這些基督徒雖然口裡聲稱認識主,但其生活表現卻與所承認的主相違反.他們可以每主日守禮拜,能講屬靈的話語；也會做屬靈的事情,但這一切背後的動機,卻是為了爭取地上的表現,為要得人前的榮耀.(雅四4)告訴我們；凡想與世俗為友的,就是與神為敵的；這也是保羅在這裡的教訓.因此,無論我們作甚麼,若我們心裡的追求的目標是地上的事,做事的原則也是以地上的標準作抉擇；那麼,我們可能隨時把主耶穌撇掉,隨時會把信仰放棄.在啟示錄第三章中,主耶穌責備不冷不熱的老底嘉教會,以致主說:「必從我口中把你吐出去」(三16).「不冷不熱」的人並非不知道甚麼是好的,甚麼是壞,而是經過考慮之後,認為屬靈的事雖非不重要,但也不是最重要,於是只在方便時追求,不方便就不追求了.
\newline
\newline
我的家鄉是在北平,義和團的時候,人家稱信耶穌的外國人為「大毛子」,信耶穌的中國人則為「二毛子」.後來「大毛子」離開了,義和團的人就開始專對付信耶穌的「二毛子」.當時的基督徒十分驚慌,害怕遭受逼迫.於是,有人為了避難,就在家中擺放一個牌位,牌子的一面是十字架,另一面卻是祖宗的靈位.義和團來了,就放祖宗靈位；走了,就改掛十字,看情況隨時更改牌位的方向.
\newline
\newline
「專以地上的事為念」的人為了自己的方便,隨時都可以作出不承認主的事情；難怪保羅說這些人的結局是沉淪.因為既然他們隨時可以否認主,那麼主也隨時可以否認他們.
\newline
\newline
(二)天國的子民
\newline
\newline
從20節開始,保羅論及天國子民三方面的特徵:在生活中的盼望,生命對主完全的順服,以及作世人的榜樣,見證主要國度.
\newline
\newline
(1)生活的盼望:「我們卻是天上的國民,並且等候救主,就是主耶穌基督,從上降臨.」(三20)基督徒是天國的子民,因此我們的盼望是天國的降臨,並等候主的再來.主耶穌說:「在世上你們有苦難,但你們可以放心,我已經勝了世界:」(約十六33)將來有一天主再來的時候,一切的困難和不公平的事情都要除掉,公義就會顯明；而我們就與主同在榮耀裡,一齊分享祂的榮耀,這是我們的喜樂；亦是每一個基督徒在苦難的盼望.
\newline
\newline
(2)生命的順服:「祂要按著那能叫萬有歸服自己以大能,將我們這卑賤的身體改變形狀,和祂自己榮耀的身體相似.」(三21)基督徒有盼望的原因, 是因主要從天上降臨,接我們到天上與祂永遠同在,作天國的公民.既然如此,天國的子民就非常有意義了,因為主耶穌是我們的主,祂是萬主之主,萬王之王.在這國度裡,祂是絕對的公義和慈愛；同樣也要求天國的子民對祂絕對的順服.
\newline
\newline
今天我們可能忽略了主在我們身上應有的權威,許多時候我們對主的態度可能太隨便了.我們雖然口裡稱主耶穌是我們的主,但在現今講求人權的社會,有時我們也竟然在主面前講求自己的權利來.既然作為天國的子民,那麼我們就應該在萬王之王面前沒有任何的權利可講,在生命中追求對主完全的順服；而不是想著自己的喜好或者環境的變遷,乃是主在我身上的旨意,並得著祂所要我得的；若沒有這樣對主完全順服的心,我們就不配進入神的國.
\newline
\newline
(3)世人的榜樣:「我所親愛,所想念的弟兄們,你們就是我的喜樂,我的冠冕.我親愛的弟兄,你們應當靠主站立得穩.」(四1)既然基督徒是天國的公民,在世上是耶穌基督的代表,那麼就該對世界有應負的責任.世人看不見天國是甚麼樣子,只能從基督徒的身上看到天上的國度,因此天國的公民就需要把這個國度見證出來,在生命中活出天國子民的樣式,建立天上的國度.
\newline
\newline
在早期的殖民地政府,宗主國享有特別的權利,聽聞在南洋一帶的社會,英國公民從不坐三等火車；認為只有坐頭等才不致失卻國家的面子.雖然公民以身份是屬世的,但世人還是這麼尊重,謹慎；作為天國的子民,我們的生活就必須與我們的身份相稱,否則就容易失去見證.很多時候,我們常忘卻了自己的身份,稍遇到不如意的事情就大發脾氣,難怪別人就因此指責我們:基督徒也不外如是吧!許多事情雖然本身並不一定是錯誤的,但若與我們的身份不相符的話,我們寧願不作.保羅在(腓三17)說:「你們要效法我,也當留意看那些照我們榜樣行的人」.今日縱使世界的道德標準日漸低落,但世人對基督徒的要求是比較高的,因此我們不能輕看自己的身份；也不應忘卻自己的職責,在生命中活出應有的見證.在人前作出應有的榜樣；因為我們是天國的子民.
\newline
\newline


\newpage

\section{第62屆~港九培靈會~鮑會園~第九講}
\label{sec:1084}
\textbf{堅守的動力 — 靠主站立得穩}
\newline
\newline
連結: \href{https://www.hkbibleconference.org/session-message/view/1084}{\texttt{https://www.hkbibleconference.org/session-message/view/1084}}
\newline
\newline
\hyperref[sec:1083]{< < < PREV SERMON < < <}
~
\hyperref[sec:toc]{[返目錄]}
~
\hyperref[sec:1035]{> > > NEXT SERMON > > >}
\newline
\newline
腓立比書 4:4-4:13
\newline
\begin{longtable}{cl}
\hline
\hline
章節 & 經文 (和合本修訂版)\\
\hline
4:4 & \begin{tabularx}{0.7\textwidth}{X} 你們要靠主常常喜樂。我再說，你們要喜樂。 \end{tabularx} \\ \\ \relax
4:5 & \begin{tabularx}{0.7\textwidth}{X} 要讓眾人知道你們謙讓的心。主已經近了。 \end{tabularx} \\ \\ \relax
4:6 & \begin{tabularx}{0.7\textwidth}{X} 應當一無掛慮，只要凡事藉著禱告、祈求和感謝，將你們所要的告訴神。 \end{tabularx} \\ \\ \relax
4:7 & \begin{tabularx}{0.7\textwidth}{X} 神所賜那超越人所能了解的平安，必在基督耶穌裡，保守你們的心懷意念。 \end{tabularx} \\ \\ \relax
4:8 & \begin{tabularx}{0.7\textwidth}{X} 末了，弟兄們，凡是真實的、凡是可敬的、凡是公義的、凡是清潔的、凡是可愛的、凡是有美名的，若有甚麼德行，若有甚麼稱讚，你們都要留意。 \end{tabularx} \\ \\ \relax
4:9 & \begin{tabularx}{0.7\textwidth}{X} 你們從我所學習的，所領受的，所聽見的，所看見的事，你們都要繼續去做，賜平安的神就必與你們同在。 \end{tabularx} \\ \\ \relax
4:10 & \begin{tabularx}{0.7\textwidth}{X} 我靠主大大喜樂，因為你們關懷我的心如今又表現了出來；其實你們一直都關懷我，只是沒有機會罷了。 \end{tabularx} \\ \\ \relax
4:11 & \begin{tabularx}{0.7\textwidth}{X} 我並不是因缺乏而說這話，因為我已經學會無論在甚麼景況都可以知足。 \end{tabularx} \\ \\ \relax
4:12 & \begin{tabularx}{0.7\textwidth}{X} 我知道怎樣處卑賤，也知道怎樣處豐富；或飽足或飢餓，或有餘或缺乏，任何事情，任何景況，我都得了祕訣。 \end{tabularx} \\ \\ \relax
4:13 & \begin{tabularx}{0.7\textwidth}{X} 我靠著那加給我力量的，凡事都能做。 \end{tabularx} \\ \\
[1ex]
\hline
\hline
\end{longtable}
腓立比教會在屬靈上是很蒙福的,但當時他們受到敵人的驚嚇(一28)；而保羅在寫本書信的時候,他也正身處監牢,生死未卜.當生命面臨極大考驗的時刻,在種種的壓力之下,腓立比書卻充滿喜樂的信息.
\newline
\newline
(一)靠主常常喜樂
\newline
\newline
(腓四4):「你們要靠主常常喜樂.我再說,你們要喜樂.」這裡的「常常」是指「繼續不斷」的意思,保羅勉勵腓立比信徒無論在任何的考驗和壓力中,都當靠主常常喜樂；這種喜樂不是指外表或環境上的；乃是指親身經驗到從主而來的真正喜樂.根據羅馬書的教導,基督徒不僅要在患難中有喜樂,並且更要為著患難而常常喜樂(五3),這是多麼特殊的教訓!其實,基督徒的喜樂不是建立在外在的環境上,乃是明白神要藉著苦難在我們身上成就祂的美意.聖經常常給我們這樣的見證,就是把我們外在環境所遭遇的事情,與內在的屬靈生命的感受分開:「你使我心裡快樂,勝過那豐收五榖新酒的人」(詩四7)按人的標準,物質的豐富就能叫人歡天喜地,但詩人卻說神賜予的喜樂比一切都寶貴.因為從肉身,物質或環境而來的喜樂都是暫時的,若果將來的一切都改變了,那麼我們心中的喜樂也必隨之消失.因此,保羅提醒我們不但要靠主喜樂,更要靠著考驗和患難而喜樂,深信神的計劃是最美好的.
\newline
\newline
(二)謙讓的心
\newline
\newline
腓四5「當叫眾人知道你們謙讓的心,主已經近了.」新約中「謙讓」這個字共用過五次,可以另繙譯作「和平」,「溫柔」,意思就是指在主面前謙卑,溫柔而得著和平的態度.但這種態度是對哪些人來說的呢?這裡保羅指的可能是第三章的「十架的仇敵」,或者是走上異端的人,由於這些人的真理知識有限,生命方面可能追求錯了,因而為教會帶來許多困難,甚至逼迫教會.保羅勸勉我們對這些人要有謙讓的心,不要因他們的態度來對待他們,反而更要用溫柔,謙卑的心來體諒他們的軟弱,明白他們在真理,思想及知識上的無知之處,包容他們或者因著某些環境中的困難而軟弱跌倒.
\newline
\newline
保羅勉勵腓立比信徒要靠主常常喜樂；又當存謙讓的心,原因是「主已經近了」.在原文的解釋中,這句話只有「主近」的意思:一方面可能是指時間上的主再來；另一方面也可能是指空間上主的同在.根據經文的文法,後者的解釋比較合理,就是說主隨時在我們的身邊,我們只管坦然無懼地來到主的施恩寶座前,為要得憐恤,蒙恩惠作隨時的幫助(來四16).所以,即使我們的生命受到逼迫或考驗,當自己的方法解決不了,肉體忍受不了的時候,聖靈成了隨時的幫助；其實,在我們承受重擔的時候,聖靈一直在旁扶持,是在我們遇著危險或不能擔當時.因此,我們不但要常在主裡常常喜樂,也要存謙讓的心在主裡面對考驗,讓主成為隨時的幫助,我們就能在追求成熟的生活中站立得住.
\newline
\newline
(三)應當一無掛慮
\newline
\newline
(腓四6)「應當一無掛慮,只要凡事藉著禱告,祈求和感謝,將你們所要的告訴神.」中文的「一無掛慮」繙譯得很恰當,這裡不是指甚麼都不關心,甚麼都不管,或者一切教會行政,個人事業及家庭生活等都不理,只理屬靈的事情；這裡所指的是把一切的事情都帶到主的面前,不要為任何的事情而掛慮.試想當我們樣樣都理會,事事都掛慮時,我們的心又怎能安心地倚靠神呢?換句話說,基督徒應該關心生活上的事情,也應當盡個人的本份,但不能讓這些事分散我們倚靠神的心,馬太福音說:「你們要先求祂的國和祂的義,這些東西都要加給你們了.所以不要為明天憂慮,因為明天自有明天的憂慮,一天的難處一天當就夠了.」(太六 33-34)天空的飛鳥不是甚麼都不關心,但不因為食物而擔心掛慮,只要盡本份就一定找到.同樣,我們生活中所遭遇的一切事情,都有神獨特的計劃和安排, 我們只要把所計劃,思想,關心的──凡事告訴神.凡事包括團契,主日學的工作,個人的靈性,日常的生活,學業,生意等,都藉著禱告交托給神.
\newline
\newline
許多基督徒以為神只理會屬靈的事,而家庭的經濟,個人的生活在神看來並不重要,這是錯誤的觀念.實際上,神關心的不僅是我們的靈魂得救,也關心我們的身體得贖；因為身體是神的殿,靈魂要在肉體中才能事奉神；不單是在今生,來生也是一樣.身體與靈魂分離是不正常的現象,所以聖經說我們要「等候」身體 「得贖」,意思是將來靈魂和身體都要同蒙拯救,復活的身體再與靈魂復合後才有資格與神永遠同在,藉復活的身體來事奉神；因此,身體和靈魂是同等重要的,若果我們在禱告中將身體和靈魂都交托給主,我們就不必為任何事掛慮了.
\newline
\newline
(四)出人意外的平安
\newline
\newline
(腓四7)「神所賜出人意外的平安,必在基督耶穌裡保守你們的心懷意念.」有人解釋「出人意外的平安」是指超乎人理智可以明白的平安,即是人同耶穌基督的救贖而享受與神和好的平安(賽五十三5)；但這裡所指的不單是人在神面前的平安；在主裡完全得安息的平安,這平安是人無法理解的.主耶穌說:「我留下平安給你們,我將我的平安賜給你們.我所賜的不像世人所賜的.你們心裡不要憂愁,也不要膽怯」(約十四27),神所賜的平安不是可以用人的標準衡量的,主耶穌在十字架的苦難面前仍有平安,因祂確知是走在神的旨意中.同樣,腓立比教會若行在神的旨意中,即使面臨考驗和壓力,他們仍可靠著主所賜的平安而一無掛慮.
\newline
\newline
(五)你們都要思念
\newline
\newline
(腓四8)「弟兄們,我還有未盡的話,凡是真實的,可敬的,公義的,清潔的,可愛的,有美名的,若有甚麼德行,若有甚麼稱讚,這些事你們都要思念.」保羅勉勵我們要思想美好,屬靈,得神喜悅的事情,而不是思念環境的驚嚇,肉體的痛苦或不公平的遭遇.一個人常思念的,就成為內心所關注的,他的生命也必因此受影響.我們所思念的若盡是屬肉體的事,那麼,我們就一定是屬世界的人；(腓三19)亦說明「專以地上的事為念」者,不能體貼神的心意.心中的思念影響我們整個的屬靈方向,你用甚麼充滿你的思想呢?
\newline
\newline
(六)凡事都能作
\newline
\newline
(腓四13)「我靠著那加給我力量的,凡事都能作.」在實際生活經驗中,我們發現雖然靠主的力量,但有的事仍然辦不到；到底這裡的「凡事都能作」是指著甚麼說的?其實,這裡的「凡事」就是指前面經文所講到的靠主站立得穩(1節),靠主常常喜樂(4節),以及靠主隨時知足(12節)；我們不是靠著自己,乃是靠著主的力量在生活上存心謙讓,一無掛慮；因為主隨時幫助我們,我們在祂的旨意中,能夠做一切祂要我們所作的,又能靠著祂在追求的事上站立得穩.
\newline
\newline
為甚麼我們要靠主才站立得穩呢?因為世上有兩個陣勢相爭,而基督徒不單是站在主的一邊與撒旦爭戰,更是在基督裡已經戰勝撒旦；因此,我們必須常常把守勝利的成果,不要離開得勝的地位,求主保守我們,無論將來要面對任何環境,我們要在主裡站立得穩.
\newline
\newline


\newpage

\section{第63屆~港九培靈會~劉東崑~第一講}
\label{sec:1035}
\textbf{神為何創造宇宙萬物}
\newline
\newline
連結: \href{https://www.hkbibleconference.org/session-message/view/1035}{\texttt{https://www.hkbibleconference.org/session-message/view/1035}}
\newline
\newline
\hyperref[sec:1084]{< < < PREV SERMON < < <}
~
\hyperref[sec:toc]{[返目錄]}
~
\hyperref[sec:1036]{> > > NEXT SERMON > > >}
\newline
\newline
以弗所書 1:9-1:9
\newline
\begin{longtable}{cl}
\hline
\hline
章節 & 經文 (和合本修訂版)\\
\hline
1:9 & \begin{tabularx}{0.7\textwidth}{X} 照自己在基督裡所立定的美意，使我們知道他旨意的奧祕， \end{tabularx} \\ \\
[1ex]
\hline
\hline
\end{longtable}
「都是照他自己所豫定的美意,叫我們知道他旨意的奧秘.」(弗一9)
\newline
\newline
神願意我們知道祂旨意的奧秘,也就是祂在萬世以前所豫定的美意.愈有智慧的人,做事愈有計劃.神是智慧的總根源,做事更有大計劃；對於宇宙萬有有祂豫定的美意,過去是隱藏的,如今都讓我們明白過來.(參弗三9)歷世歷代這奧秘都是隱藏的,一直到聖靈啟示祂的聖使徒和先知,人才慢慢明白神為何及如何創造宇宙萬有；神對於祂的創造有何目的,這些都是屬於神旨意的奧秘.所以我們要慢慢把它分解明白出來,好使神的計劃成就在我們身上.
\newline
\newline
剛才我們所讀的聖經,哥林多前書第二章一至七節,前五節講到我們先要得著救恩,保羅以前到哥林多的時候,沒有講及神的奧秘,因為那時,當地的人還沒有得救.但是保羅再到哥林多的時候,當地的人已經得救了,被聖靈重生了,就可以明白聖經中的奧秘了.神奧秘的事需要聖靈幫助我們才能參透(林前二 9-10).假如我們未曾重生得救；靈性未曾與神接觸,就無法領受神的奧秘.我們查考神的奧秘,需要迫切的禱告,把靈打開；聖靈賜我們悟性,智慧,使我們能夠看見神預定的美意.
\newline
\newline
我們基督徒的經歷可分為兩部份,首先是得救恩──我們得救是本乎恩,也因著信；只要信就可得救了.但是得救只不過是一個開始,神更有祂的美意,就是在萬世以前預定我們得榮耀；要得榮耀就不是僅僅相信便行,乃要竭力的追求,還要操練,學習,脫去舊人,穿上新人.當我們把得救恩的事情都弄清楚以後,就要追求得榮耀,得榮耀乃是神豫定的美意(林前二6-7).因為要得榮耀,所以我們要明白神旨意的奧秘.我們這幾天的查經,就是要說明神豫定的美意,亦即是說明神旨意的奧秘.
\newline
\newline
(一)認識神
\newline
\newline
中國人的「神」字有甚麼意義?有說聰明正直之謂「神」,這說法不是頂好.孟子曰:「大而化之之謂聖,聖而不可知之謂神.」所以神在聖經中有另一個名稱就是至聖者.神比大而化之的還要大,神是偉大得令人難以明白；廟裡頭的造像都很渺小,不足以稱為神.我們所信的神大得不得了,其大無法測度(詩一四五3).神為大,我們不能全知(伯卅六26),究竟神有多大呢?且看主耶穌的親自介紹(約十29).主耶穌告訴我們:「神比萬有都大」, 萬有就是指天上地上一切所有的.我們對這宇宙萬有不能明白有多大,我只能解釋一二,作為對各位的一點提醒而已.科學家告訴我們:地球屬於太陽系,太陽系有九大行星,圍繞太陽運轉,太陽乃一恆星,屬銀河系,銀河系裡頭像太陽般發光的星有三千億,小星星不計其數,像地球般不發光的行星無法勝數,尚未計在其中. 銀河系佔的空間直徑有十萬光年,光每秒鐘跑十八萬六千二百六十二哩(即三十萬公里),光每年跑六億萬哩,十萬光年所跑的哩數則是一個天文數字,人不可以想像；然而,這個銀河系,只不過是天上無數星雲裡的一個而已.到底天上有多少星雲呢?現在人還沒法測度,只是用盡了一切的觀察力,和先進的望遠鏡,才能看到 七十億光年.在七十億光年裡,大概有六十億銀河系,它們各自組在一起成為一個重力系統；這是天文學家不斷觀測出來的.我們現存的這個銀河系乃是跟其他十六 星雲一起運轉的.有的時候是五個銀河系一起在轉,並且他們是向外擴散的,有些時候卻又是五百個銀河系一起在轉.所以說,在七十億光年範圍之內,有六十多億銀河系可以估計.根據愛因斯坦的算法,這個物質宇宙估計有七百億光年之大,人類只能觀測到它的十分之一.物質宇宙不過是神創造諸宇宙裡的一個(來一2)； 上下左右謂之「宇」,古往今來謂之「宙」,「宇宙」換言之即是空間和時間.神創造諸宇宙,不是僅僅創造物質宇宙；還有非物質宇宙,將來我們要脫離這個物質宇宙；往那非物質的宇宙去,這是神旨意的奧秘.
\newline
\newline
神比萬有都大,僅僅一個物質宇宙我們已無法觀測,神創造諸宇宙,其大無法測度,不能全知,這是我們對神的一個認識.親愛的弟兄姊妹,在你我的思想觀念裡,在你對神的想像中,必須有一概念:(一)神是大而化之,聖而不可知.(二)神是靈(約四24),靈跟物質不同,無方積,無重量,無形像；動能制靜,靜能制動；不能化驗,不能分析,因它不是物質.物質剛好相反,一定佔一個地方,有體積,有重量,有形像；動者恆動,靜者恆靜,能夠化驗,能夠分析.神是個靈,雖是不能見卻又是至活的；凡是看不見的東西,又能接受的,就叫做信.神實在太大,讓我們看也看不見；像我們這些眼光如豆,視覺世界有限的人的時候曾喊著說:「以利,以利,拉馬撒巴各大尼!」以利的意思是「能力的總匯」,另聖經有多處翻譯稱呼神為大能者,全能者都是同一的意思.神所有的能力,人都看不見,所以神要把祂的能力從自隱的變成彰顯的.宇宙的創造,主要為彰顯神能(詩一五1),這就是神創造的目的.
\newline
\newline
神是愛的總源(詩六十二12,約壹四16),一切的愛都在祂裡頭.神是靈,祂滿有能力,卻又是自隱的；假如祂不發表,不顯明,祂的愛等於零,祂的能力也看不見(詩卅六9).神是生命的總源頭,滿有能力,慈愛；神要把祂一切隱藏的豐盛都表明出來,這就是神創造的目的.神內裡不單只有能力,更有豐富的智和知識(羅十一33).人類為何要尋求知識呢?這就是因為人內裡有一個靈,所以人異於禽獸,有一個求知的慾望(箴二10).神是知識的源頭,祂要把知識發揮出來,讓人的靈去求知,慢慢的多知道神,所以到了末後,認識耶和華的知識必充滿全地,現在的科學家們愈來愈認識神；因為他們的知識漸漸進步,藉著研究宇宙萬物愈多的瞭解神.
\newline
\newline
太初有道,道就是發表,神一把祂的能力,智慧和豐盛發表,宇宙萬有就創造出來了.神的發表有兩大類；其一表明神外在的榮耀,諸天述說神的榮耀,穹蒼傳揚祂的手段(詩十九1)整個宇宙是一個力場,今早你從家裡跑到這個會場；如果吃了早點,你會感到滿有能量；你坐在這兒聽道,整個腦細胞都在動,你聽得愈仔細,神經細胞就活動得愈快,消耗你的能量就愈厲害,因為裡頭的力量用光了；不能久坐要趕快起來去補充食物.整個宇宙都在能力的運動之下,因為神藉創造宇宙萬有去顯明祂的能力,而這外在的顯明尚嫌未夠,祂的內在還有許多美德(詩一○六2).
\newline
\newline
本來神造人,是要在人身上表明祂的美德；可惜其後被撒旦破壞,沒有一人達到這個目的,都落在罪裡.後來那道成肉身的來了,祂把神一切美德都表明出來,祂就是耶穌基督.在太初,祂是出來的道,一千九百多年前,祂道成肉身,來表明神一切的美德(約一1-2,14,18).神是隱藏的,祂藉著道創造萬有,神外在的榮耀都彰顯出來了,神內在還有許多細緻的美德,只有主耶穌道成肉身才能把神完全的表明出來.所以神創造的目的在於把祂自己表明,正因神的內容豐富,祂要藉著萬有來表明祂外在的榮耀,又要藉著我們這些人來表明祂內在的榮耀.
\newline
\newline


\newpage

\section{第63屆~港九培靈會~劉東崑~第二講}
\label{sec:1036}
\textbf{神如何創造宇宙萬物}
\newline
\newline
連結: \href{https://www.hkbibleconference.org/session-message/view/1036}{\texttt{https://www.hkbibleconference.org/session-message/view/1036}}
\newline
\newline
\hyperref[sec:1035]{< < < PREV SERMON < < <}
~
\hyperref[sec:toc]{[返目錄]}
~
\hyperref[sec:1037]{> > > NEXT SERMON > > >}
\newline
\newline
以弗所書 3:2-3:5
\newline
\begin{longtable}{cl}
\hline
\hline
章節 & 經文 (和合本修訂版)\\
\hline
3:2 & \begin{tabularx}{0.7\textwidth}{X} 想必你們曾聽見神賜恩給我，把關切你們的職分託付我， \end{tabularx} \\ \\ \relax
3:3 & \begin{tabularx}{0.7\textwidth}{X} 用啟示讓我知道福音的奧祕，正如我以前略略寫過的。 \end{tabularx} \\ \\ \relax
3:4 & \begin{tabularx}{0.7\textwidth}{X} 你們讀了，就會知道我深深了解基督的奧祕； \end{tabularx} \\ \\ \relax
3:5 & \begin{tabularx}{0.7\textwidth}{X} 這奧祕在以前的世代沒有讓人知道，像如今藉著聖靈向他的聖使徒和先知啟示一樣， \end{tabularx} \\ \\
[1ex]
\hline
\hline
\end{longtable}


\newpage

\section{第63屆~港九培靈會~劉東崑~第三講}
\label{sec:1037}
\textbf{神對宇宙萬有作如何安排}
\newline
\newline
連結: \href{https://www.hkbibleconference.org/session-message/view/1037}{\texttt{https://www.hkbibleconference.org/session-message/view/1037}}
\newline
\newline
\hyperref[sec:1036]{< < < PREV SERMON < < <}
~
\hyperref[sec:toc]{[返目錄]}
~
\hyperref[sec:1038]{> > > NEXT SERMON > > >}
\newline
\newline
以弗所書 3:9-3:11
\newline
\begin{longtable}{cl}
\hline
\hline
章節 & 經文 (和合本修訂版)\\
\hline
3:9 & \begin{tabularx}{0.7\textwidth}{X} 又使眾人都明白甚麼是歷代以來隱藏在創造萬物之神裡的奧祕， \end{tabularx} \\ \\ \relax
3:10 & \begin{tabularx}{0.7\textwidth}{X} 為要在現今藉著教會使天上執政的、掌權的知道神百般的智慧。 \end{tabularx} \\ \\ \relax
3:11 & \begin{tabularx}{0.7\textwidth}{X} 這是照著神在我們主基督耶穌裡所完成的永恆的計劃。 \end{tabularx} \\ \\
[1ex]
\hline
\hline
\end{longtable}


\newpage

\section{第63屆~港九培靈會~劉東崑~第四講}
\label{sec:1038}
\textbf{神為何創造人類}
\newline
\newline
連結: \href{https://www.hkbibleconference.org/session-message/view/1038}{\texttt{https://www.hkbibleconference.org/session-message/view/1038}}
\newline
\newline
\hyperref[sec:1037]{< < < PREV SERMON < < <}
~
\hyperref[sec:toc]{[返目錄]}
~
\hyperref[sec:1039]{> > > NEXT SERMON > > >}
\newline
\newline
以賽亞書 43:7-43:7
\newline
\begin{longtable}{cl}
\hline
\hline
章節 & 經文 (和合本修訂版)\\
\hline
43:7 & \begin{tabularx}{0.7\textwidth}{X} 就是凡稱為我名下的人，是我為自己的榮耀創造的，是我所塑造，所做成的。」 \end{tabularx} \\ \\
[1ex]
\hline
\hline
\end{longtable}


\newpage

\section{第63屆~港九培靈會~劉東崑~第五講}
\label{sec:1039}
\textbf{神如何創造人類}
\newline
\newline
連結: \href{https://www.hkbibleconference.org/session-message/view/1039}{\texttt{https://www.hkbibleconference.org/session-message/view/1039}}
\newline
\newline
\hyperref[sec:1038]{< < < PREV SERMON < < <}
~
\hyperref[sec:toc]{[返目錄]}
~
\hyperref[sec:1040]{> > > NEXT SERMON > > >}
\newline
\newline


\newpage

\section{第63屆~港九培靈會~劉東崑~第六講}
\label{sec:1040}
\textbf{神在人身上有預定的美意}
\newline
\newline
連結: \href{https://www.hkbibleconference.org/session-message/view/1040}{\texttt{https://www.hkbibleconference.org/session-message/view/1040}}
\newline
\newline
\hyperref[sec:1039]{< < < PREV SERMON < < <}
~
\hyperref[sec:toc]{[返目錄]}
~
\hyperref[sec:1041]{> > > NEXT SERMON > > >}
\newline
\newline
以弗所書 1:9-1:11
\newline
\begin{longtable}{cl}
\hline
\hline
章節 & 經文 (和合本修訂版)\\
\hline
1:9 & \begin{tabularx}{0.7\textwidth}{X} 照自己在基督裡所立定的美意，使我們知道他旨意的奧祕， \end{tabularx} \\ \\ \relax
1:10 & \begin{tabularx}{0.7\textwidth}{X} 要照著所安排的，在時機成熟的時候，使天上、地上、一切所有的，都在基督裡面同歸於一。 \end{tabularx} \\ \\ \relax
1:11 & \begin{tabularx}{0.7\textwidth}{X} 我們也在他裡面得了基業；這原是那位隨己意行萬事的神照著自己的旨意所預定的， \end{tabularx} \\ \\
[1ex]
\hline
\hline
\end{longtable}


\newpage

\section{第63屆~港九培靈會~劉東崑~第七講}
\label{sec:1041}
\textbf{神如何成就祂的美意}
\newline
\newline
連結: \href{https://www.hkbibleconference.org/session-message/view/1041}{\texttt{https://www.hkbibleconference.org/session-message/view/1041}}
\newline
\newline
\hyperref[sec:1040]{< < < PREV SERMON < < <}
~
\hyperref[sec:toc]{[返目錄]}
~
\hyperref[sec:1042]{> > > NEXT SERMON > > >}
\newline
\newline
馬太福音 22:14-22:14
\newline
\begin{longtable}{cl}
\hline
\hline
章節 & 經文 (和合本修訂版)\\
\hline
22:14 & \begin{tabularx}{0.7\textwidth}{X} 因為被召的人多，選上的人少。」 \end{tabularx} \\ \\
[1ex]
\hline
\hline
\end{longtable}


\newpage

\section{第63屆~港九培靈會~劉東崑~第八講}
\label{sec:1042}
\textbf{神安排時間,祂如何藉暫時成就永遠}
\newline
\newline
連結: \href{https://www.hkbibleconference.org/session-message/view/1042}{\texttt{https://www.hkbibleconference.org/session-message/view/1042}}
\newline
\newline
\hyperref[sec:1041]{< < < PREV SERMON < < <}
~
\hyperref[sec:toc]{[返目錄]}
~
\hyperref[sec:1043]{> > > NEXT SERMON > > >}
\newline
\newline
傳道書 3:1-3:2
\newline
\begin{longtable}{cl}
\hline
\hline
章節 & 經文 (和合本修訂版)\\
\hline
3:1 & \begin{tabularx}{0.7\textwidth}{X} 凡事都有定期，天下每一事務都有定時。 \end{tabularx} \\ \\ \relax
3:2 & \begin{tabularx}{0.7\textwidth}{X} 生有時，死有時；栽種有時，拔出有時； \end{tabularx} \\ \\
[1ex]
\hline
\hline
\end{longtable}


\newpage

\section{第63屆~港九培靈會~劉東崑~第九講}
\label{sec:1043}
\textbf{人如何經歷時間進入永遠}
\newline
\newline
連結: \href{https://www.hkbibleconference.org/session-message/view/1043}{\texttt{https://www.hkbibleconference.org/session-message/view/1043}}
\newline
\newline
\hyperref[sec:1042]{< < < PREV SERMON < < <}
~
\hyperref[sec:toc]{[返目錄]}
~
\hyperref[sec:1392]{> > > NEXT SERMON > > >}
\newline
\newline
以弗所書 3:2-3:11
\newline
\begin{longtable}{cl}
\hline
\hline
章節 & 經文 (和合本修訂版)\\
\hline
3:2 & \begin{tabularx}{0.7\textwidth}{X} 想必你們曾聽見神賜恩給我，把關切你們的職分託付我， \end{tabularx} \\ \\ \relax
3:3 & \begin{tabularx}{0.7\textwidth}{X} 用啟示讓我知道福音的奧祕，正如我以前略略寫過的。 \end{tabularx} \\ \\ \relax
3:4 & \begin{tabularx}{0.7\textwidth}{X} 你們讀了，就會知道我深深了解基督的奧祕； \end{tabularx} \\ \\ \relax
3:5 & \begin{tabularx}{0.7\textwidth}{X} 這奧祕在以前的世代沒有讓人知道，像如今藉著聖靈向他的聖使徒和先知啟示一樣， \end{tabularx} \\ \\ \relax
3:6 & \begin{tabularx}{0.7\textwidth}{X} 就是外邦人在基督耶穌裡，藉著福音，得以同為後嗣，同為一體，同為蒙應許的人。 \end{tabularx} \\ \\ \relax
3:7 & \begin{tabularx}{0.7\textwidth}{X} 我作了這福音的僕役，是照著神的恩賜，是照他運行的大能賜給我的。 \end{tabularx} \\ \\ \relax
3:8 & \begin{tabularx}{0.7\textwidth}{X} 雖然我比眾聖徒中最小的還小，他還賜我這恩典，讓我把基督那測不透的豐富傳給外邦人， \end{tabularx} \\ \\ \relax
3:9 & \begin{tabularx}{0.7\textwidth}{X} 又使眾人都明白甚麼是歷代以來隱藏在創造萬物之神裡的奧祕， \end{tabularx} \\ \\ \relax
3:10 & \begin{tabularx}{0.7\textwidth}{X} 為要在現今藉著教會使天上執政的、掌權的知道神百般的智慧。 \end{tabularx} \\ \\ \relax
3:11 & \begin{tabularx}{0.7\textwidth}{X} 這是照著神在我們主基督耶穌裡所完成的永恆的計劃。 \end{tabularx} \\ \\
[1ex]
\hline
\hline
\end{longtable}


\newpage

\section{第63屆~港九培靈會~楊濬哲~序}
\label{sec:1392}
\textbf{序}
\newline
\newline
連結: \href{https://www.hkbibleconference.org/session-message/view/1392}{\texttt{https://www.hkbibleconference.org/session-message/view/1392}}
\newline
\newline
\hyperref[sec:1043]{< < < PREV SERMON < < <}
~
\hyperref[sec:toc]{[返目錄]}
~
\hyperref[sec:1066]{> > > NEXT SERMON > > >}
\newline
\newline
九一年第六十三屆港九培靈研經大會,能如期舉行,圓滿閉會；深感父神莫大厚恩,在動盪王不安的世代裡,使我們仍有祂的話語,足以倚靠.
\newline
\newline
本屆研經會,請劉東崑牧師主講.總題:「神旨意的奧秘」.劉牧師以主題之前段講述救恩,後段講述如何追求得著榮耀.對於神在宇宙萬有中的旨意,歷世歷代都是隱藏的；直到聖靈啟示祂的聖徒和先知,人才得以漸漸明白神旨意的奧秘.
\newline
\newline
午堂講道會,請黃子嘉牧師主講.總題:「討主喜悅的生活.」黃牧師謂,基督徒在世上每天所過的生活,最重要的是討主喜悅；有美好的見證,才能帶領親友外人信主.因此,要從聖經中不同之角度,探討神所喜悅的生活方式；使我們能勝過世界物質之引誘,活出美好的見證來.
\newline
\newline
晚堂奮興會,請麥希真牧師主講.總題:「渴慕基督.」麥牧師指出,現今世人都在團團亂轉,失去方向.基督徒生活在這急速變化的世界中,必須有一個軸心,基督就是我們的軸心；我們環繞著這個軸心,才不致出軌.現代的教會雖然有良好的神學思想,和工作理論；但最好仍是多講論耶穌基督.如果我們認定祂是我們的軸心,那麼,無論世界怎樣變,我們都不會失落.
\newline
\newline
在晚堂奮興會中,講員呼召時,決志信主者有十七人,悔改更新為主而活者二百二十六人,全時間事奉者三十二人,合共為二百七十六人.赴會者來自港九,新界,離島各地,足見他們愛慕主道的心,是何等的迫切!
\newline
\newline
黃子嘉牧師於感恩會中致詞謂:我患了感冒,雖然連日勞累,反而痊愈了.我希奇港九培靈研經會,能夠辦理得如此成功,十天聚會並不簡單,實在要感謝主的恩典；我必繼續為你們代禱!麥希真牧師致詞謂,返港似回到家中,看見會眾能專心赴會,尤以青年人佔大部份,深信教會是有前途的.麥師母亦云:我得格外恩典,能隨麥牧師來赴會,回憶四十至五十年代,我與麥牧師,年年例必參加聚會；麥牧師也是培靈會的果子,在聽道時蒙召奉獻的.劉牧師因趕赴台北講道,未能出席感恩會.
\newline
\newline
本講道集,是六十三屆大會的講道記錄.蒙鄭坤英姊妹記錄研經會.王大衛弟兄記錄講道會.雷麗端姊妹記錄奮興會.朱建磯牧師整理編校,方能如期出版,謹此致謝!
\newline
\newline
歷屆大會,均蒙九龍城浸信會,借用為會場；執事,同工,義工及各兄姊協助招待,維持秩序.福音傳播中心錄音錄影基督教週報刊登新聞稿及廣告；時代論 壇,星島晚報,刊登廣告；神托之家耀安工場,寄發宣傳品；彼此配搭,廣事宣傳,藉增效益.更蒙各教會兄姊踴躍赴會及奉獻,統致謝忱!
\newline
\newline


\newpage

\section{第63屆~港九培靈會~麥希真~第一講}
\label{sec:1066}
\textbf{基督的原初}
\newline
\newline
連結: \href{https://www.hkbibleconference.org/session-message/view/1066}{\texttt{https://www.hkbibleconference.org/session-message/view/1066}}
\newline
\newline
\hyperref[sec:1392]{< < < PREV SERMON < < <}
~
\hyperref[sec:toc]{[返目錄]}
~
\hyperref[sec:1067]{> > > NEXT SERMON > > >}
\newline
\newline


\newpage

\section{第63屆~港九培靈會~麥希真~第二講}
\label{sec:1067}
\textbf{基督的門徒}
\newline
\newline
連結: \href{https://www.hkbibleconference.org/session-message/view/1067}{\texttt{https://www.hkbibleconference.org/session-message/view/1067}}
\newline
\newline
\hyperref[sec:1066]{< < < PREV SERMON < < <}
~
\hyperref[sec:toc]{[返目錄]}
~
\hyperref[sec:1068]{> > > NEXT SERMON > > >}
\newline
\newline


\newpage

\section{第63屆~港九培靈會~麥希真~第三講}
\label{sec:1068}
\textbf{基督的天國}
\newline
\newline
連結: \href{https://www.hkbibleconference.org/session-message/view/1068}{\texttt{https://www.hkbibleconference.org/session-message/view/1068}}
\newline
\newline
\hyperref[sec:1067]{< < < PREV SERMON < < <}
~
\hyperref[sec:toc]{[返目錄]}
~
\hyperref[sec:1069]{> > > NEXT SERMON > > >}
\newline
\newline


\newpage

\section{第63屆~港九培靈會~麥希真~第四講}
\label{sec:1069}
\textbf{基督的教會}
\newline
\newline
連結: \href{https://www.hkbibleconference.org/session-message/view/1069}{\texttt{https://www.hkbibleconference.org/session-message/view/1069}}
\newline
\newline
\hyperref[sec:1068]{< < < PREV SERMON < < <}
~
\hyperref[sec:toc]{[返目錄]}
~
\hyperref[sec:1070]{> > > NEXT SERMON > > >}
\newline
\newline


\newpage

\section{第63屆~港九培靈會~麥希真~第五講}
\label{sec:1070}
\textbf{基督的教導}
\newline
\newline
連結: \href{https://www.hkbibleconference.org/session-message/view/1070}{\texttt{https://www.hkbibleconference.org/session-message/view/1070}}
\newline
\newline
\hyperref[sec:1069]{< < < PREV SERMON < < <}
~
\hyperref[sec:toc]{[返目錄]}
~
\hyperref[sec:1071]{> > > NEXT SERMON > > >}
\newline
\newline


\newpage

\section{第63屆~港九培靈會~麥希真~第六講}
\label{sec:1071}
\textbf{基督的愛顧}
\newline
\newline
連結: \href{https://www.hkbibleconference.org/session-message/view/1071}{\texttt{https://www.hkbibleconference.org/session-message/view/1071}}
\newline
\newline
\hyperref[sec:1070]{< < < PREV SERMON < < <}
~
\hyperref[sec:toc]{[返目錄]}
~
\hyperref[sec:1072]{> > > NEXT SERMON > > >}
\newline
\newline


\newpage

\section{第63屆~港九培靈會~麥希真~第七講}
\label{sec:1072}
\textbf{基督的審判}
\newline
\newline
連結: \href{https://www.hkbibleconference.org/session-message/view/1072}{\texttt{https://www.hkbibleconference.org/session-message/view/1072}}
\newline
\newline
\hyperref[sec:1071]{< < < PREV SERMON < < <}
~
\hyperref[sec:toc]{[返目錄]}
~
\hyperref[sec:1073]{> > > NEXT SERMON > > >}
\newline
\newline


\newpage

\section{第63屆~港九培靈會~麥希真~第八講}
\label{sec:1073}
\textbf{基督的再來}
\newline
\newline
連結: \href{https://www.hkbibleconference.org/session-message/view/1073}{\texttt{https://www.hkbibleconference.org/session-message/view/1073}}
\newline
\newline
\hyperref[sec:1072]{< < < PREV SERMON < < <}
~
\hyperref[sec:toc]{[返目錄]}
~
\hyperref[sec:1074]{> > > NEXT SERMON > > >}
\newline
\newline


\newpage

\section{第63屆~港九培靈會~麥希真~第九講}
\label{sec:1074}
\textbf{基督的受死}
\newline
\newline
連結: \href{https://www.hkbibleconference.org/session-message/view/1074}{\texttt{https://www.hkbibleconference.org/session-message/view/1074}}
\newline
\newline
\hyperref[sec:1073]{< < < PREV SERMON < < <}
~
\hyperref[sec:toc]{[返目錄]}
~
\hyperref[sec:1075]{> > > NEXT SERMON > > >}
\newline
\newline


\newpage

\section{第63屆~港九培靈會~麥希真~第十講}
\label{sec:1075}
\textbf{基督的復活}
\newline
\newline
連結: \href{https://www.hkbibleconference.org/session-message/view/1075}{\texttt{https://www.hkbibleconference.org/session-message/view/1075}}
\newline
\newline
\hyperref[sec:1074]{< < < PREV SERMON < < <}
~
\hyperref[sec:toc]{[返目錄]}
~
\hyperref[sec:1044]{> > > NEXT SERMON > > >}
\newline
\newline


\newpage

\section{第63屆~港九培靈會~黃子嘉~第一講}
\label{sec:1044}
\textbf{盼望的生活}
\newline
\newline
連結: \href{https://www.hkbibleconference.org/session-message/view/1044}{\texttt{https://www.hkbibleconference.org/session-message/view/1044}}
\newline
\newline
\hyperref[sec:1075]{< < < PREV SERMON < < <}
~
\hyperref[sec:toc]{[返目錄]}
~
\hyperref[sec:1045]{> > > NEXT SERMON > > >}
\newline
\newline
哥林多後書 5:9-5:10
\newline
\begin{longtable}{cl}
\hline
\hline
章節 & 經文 (和合本修訂版)\\
\hline
5:9 & \begin{tabularx}{0.7\textwidth}{X} 所以，無論是住在身內或住在身外，我們都立了志向要得主的喜悅。 \end{tabularx} \\ \\ \relax
5:10 & \begin{tabularx}{0.7\textwidth}{X} 因為我們眾人必須站在基督審判臺前受審，為使各人按著本身所行的，或善或惡受報。 \end{tabularx} \\ \\
[1ex]
\hline
\hline
\end{longtable}


\newpage

\section{第63屆~港九培靈會~黃子嘉~第二講}
\label{sec:1045}
\textbf{信心的生活}
\newline
\newline
連結: \href{https://www.hkbibleconference.org/session-message/view/1045}{\texttt{https://www.hkbibleconference.org/session-message/view/1045}}
\newline
\newline
\hyperref[sec:1044]{< < < PREV SERMON < < <}
~
\hyperref[sec:toc]{[返目錄]}
~
\hyperref[sec:1046]{> > > NEXT SERMON > > >}
\newline
\newline
希伯來書 11:39-11:39
\newline
\begin{longtable}{cl}
\hline
\hline
章節 & 經文 (和合本修訂版)\\
\hline
11:39 & \begin{tabularx}{0.7\textwidth}{X} 這些人都是因信獲得了讚許，卻仍未得著所應許的， \end{tabularx} \\ \\
[1ex]
\hline
\hline
\end{longtable}


\newpage

\section{第63屆~港九培靈會~黃子嘉~第三講}
\label{sec:1046}
\textbf{愛心的生活}
\newline
\newline
連結: \href{https://www.hkbibleconference.org/session-message/view/1046}{\texttt{https://www.hkbibleconference.org/session-message/view/1046}}
\newline
\newline
\hyperref[sec:1045]{< < < PREV SERMON < < <}
~
\hyperref[sec:toc]{[返目錄]}
~
\hyperref[sec:1047]{> > > NEXT SERMON > > >}
\newline
\newline
約翰一書 3:16-3:23
\newline
\begin{longtable}{cl}
\hline
\hline
章節 & 經文 (和合本修訂版)\\
\hline
3:16 & \begin{tabularx}{0.7\textwidth}{X} 基督為我們捨命，我們從此就知道何為愛；我們也當為弟兄捨命。 \end{tabularx} \\ \\ \relax
3:17 & \begin{tabularx}{0.7\textwidth}{X} 凡有世上財物的，看見弟兄缺乏，卻關閉了惻隱的心，神的愛怎能住在他裡面呢？ \end{tabularx} \\ \\ \relax
3:18 & \begin{tabularx}{0.7\textwidth}{X} 孩子們哪，我們相愛，不要只在言語或舌頭上，總要以行為和真誠表現出來。 \end{tabularx} \\ \\ \relax
3:19 & \begin{tabularx}{0.7\textwidth}{X} 從這一點，我們會知道，我們是出於真理的，並且我們在神面前可以安心， \end{tabularx} \\ \\ \relax
3:20 & \begin{tabularx}{0.7\textwidth}{X} 即使我們的心責備自己，神比我們的心大，他知道一切。 \end{tabularx} \\ \\ \relax
3:21 & \begin{tabularx}{0.7\textwidth}{X} 親愛的，我們的心若不責備我們，在神面前就可以坦然無懼了。 \end{tabularx} \\ \\ \relax
3:22 & \begin{tabularx}{0.7\textwidth}{X} 我們一切所求的，就從他得著，因為我們遵守他的命令，行他所喜悅的事。 \end{tabularx} \\ \\ \relax
3:23 & \begin{tabularx}{0.7\textwidth}{X} 神的命令就是：我們要信他兒子耶穌基督的名，並且照他所賜給我們的命令彼此相愛。 \end{tabularx} \\ \\
[1ex]
\hline
\hline
\end{longtable}


\newpage

\section{第63屆~港九培靈會~黃子嘉~第四講}
\label{sec:1047}
\textbf{成聖的生活}
\newline
\newline
連結: \href{https://www.hkbibleconference.org/session-message/view/1047}{\texttt{https://www.hkbibleconference.org/session-message/view/1047}}
\newline
\newline
\hyperref[sec:1046]{< < < PREV SERMON < < <}
~
\hyperref[sec:toc]{[返目錄]}
~
\hyperref[sec:1048]{> > > NEXT SERMON > > >}
\newline
\newline
羅馬書 1:7-1:7
\newline
\begin{longtable}{cl}
\hline
\hline
章節 & 經文 (和合本修訂版)\\
\hline
1:7 & \begin{tabularx}{0.7\textwidth}{X} 我寫信給你們在羅馬、為神所愛、蒙召作聖徒的眾人。願恩惠、平安從我們的父神和主耶穌基督歸給你們！ \end{tabularx} \\ \\
[1ex]
\hline
\hline
\end{longtable}


\newpage

\section{第63屆~港九培靈會~黃子嘉~第五講}
\label{sec:1048}
\textbf{得勝的生活}
\newline
\newline
連結: \href{https://www.hkbibleconference.org/session-message/view/1048}{\texttt{https://www.hkbibleconference.org/session-message/view/1048}}
\newline
\newline
\hyperref[sec:1047]{< < < PREV SERMON < < <}
~
\hyperref[sec:toc]{[返目錄]}
~
\hyperref[sec:1049]{> > > NEXT SERMON > > >}
\newline
\newline
約翰福音 16:33-16:39
\newline
\begin{longtable}{cl}
\hline
\hline
章節 & 經文 (和合本修訂版)\\
\hline
16:33 & \begin{tabularx}{0.7\textwidth}{X} 我對你們說了這些事，是要使你們在我裡面有平安。在世上你們有苦難，但你們要有勇氣，我已經勝過世界。」 \end{tabularx} \\ \\ \relax
17:1 & \begin{tabularx}{0.7\textwidth}{X} 耶穌說了這些話，就舉目望天，說：「父啊，時候到了，願你榮耀你的兒子，使兒子也榮耀你； \end{tabularx} \\ \\ \relax
17:2 & \begin{tabularx}{0.7\textwidth}{X} 因為你曾賜給他權柄掌管凡血肉之軀的，使他把永生賜給你所賜給他的人。 \end{tabularx} \\ \\ \relax
17:3 & \begin{tabularx}{0.7\textwidth}{X} 認識你—獨一的真神，並且認識你所差來的耶穌基督，這就是永生。 \end{tabularx} \\ \\ \relax
17:4 & \begin{tabularx}{0.7\textwidth}{X} 我在地上已經榮耀你，你交給我做的工作，我已完成了。 \end{tabularx} \\ \\ \relax
17:5 & \begin{tabularx}{0.7\textwidth}{X} 父啊，現在求你使我在你面前得榮耀，就是在未有世界以前，我同你享有的榮耀。 \end{tabularx} \\ \\ \relax
17:6 & \begin{tabularx}{0.7\textwidth}{X} 「你從世上賜給我的人，我已把你的名顯明給他們。他們本是你的，你把他們賜給我，他們也遵守了你的道。 \end{tabularx} \\ \\ \relax
17:7 & \begin{tabularx}{0.7\textwidth}{X} 現在他們知道，你所賜給我的一切都是從你那裡來的； \end{tabularx} \\ \\ \relax
17:8 & \begin{tabularx}{0.7\textwidth}{X} 因為你所賜給我的話，我已經賜給他們，他們也領受了，又確實知道，我是從你出來的，並且信你差了我來。 \end{tabularx} \\ \\ \relax
17:9 & \begin{tabularx}{0.7\textwidth}{X} 我為他們祈求，不為世人祈求，卻為你所賜給我的人祈求，因他們本是你的。 \end{tabularx} \\ \\ \relax
17:10 & \begin{tabularx}{0.7\textwidth}{X} 凡是我的都是你的，你的也是我的，並且我因他們得了榮耀。 \end{tabularx} \\ \\ \relax
17:11 & \begin{tabularx}{0.7\textwidth}{X} 我到你那裡去；我不再留在世上，他們卻在世上。聖父啊，求你因你的名，就是你所賜給我的名，保守他們，使他們像我們一樣合而為一。 \end{tabularx} \\ \\ \relax
17:12 & \begin{tabularx}{0.7\textwidth}{X} 我與他們同在的時候，我奉你的名，就是你所賜給我的名，保守了他們，我也護衛了他們；其中除了那滅亡之子，沒有一個滅亡的，好使經上的話得以應驗。 \end{tabularx} \\ \\ \relax
17:13 & \begin{tabularx}{0.7\textwidth}{X} 現在我到你那裡去，我在世上說這些話，是要他們心裡充滿了我的喜樂。 \end{tabularx} \\ \\ \relax
17:14 & \begin{tabularx}{0.7\textwidth}{X} 我已把你的道賜給他們；世界恨他們，因為他們不屬世界，正如我不屬世界一樣。 \end{tabularx} \\ \\ \relax
17:15 & \begin{tabularx}{0.7\textwidth}{X} 我不求你把他們從世上接走，只求你保全他們，使他們脫離那惡者。 \end{tabularx} \\ \\ \relax
17:16 & \begin{tabularx}{0.7\textwidth}{X} 他們不屬世界，正如我不屬世界一樣。 \end{tabularx} \\ \\ \relax
17:17 & \begin{tabularx}{0.7\textwidth}{X} 求你用真理使他們成聖；你的道就是真理。 \end{tabularx} \\ \\ \relax
17:18 & \begin{tabularx}{0.7\textwidth}{X} 你怎樣差我到世上，我也照樣差他們到世上。 \end{tabularx} \\ \\ \relax
17:19 & \begin{tabularx}{0.7\textwidth}{X} 我為他們的緣故使自己分別為聖，為要使他們也因真理成聖。 \end{tabularx} \\ \\ \relax
17:20 & \begin{tabularx}{0.7\textwidth}{X} 「我不但為這些人祈求，也為那些藉著他們的話信我的人祈求， \end{tabularx} \\ \\ \relax
17:21 & \begin{tabularx}{0.7\textwidth}{X} 使他們都合而為一。正如父你在我裡面，我在你裡面，使他們也在我們裡面，好讓世人信是你差我來的。 \end{tabularx} \\ \\ \relax
17:22 & \begin{tabularx}{0.7\textwidth}{X} 你所賜給我的榮耀，我已賜給他們，使他們合而為一，像我們合而為一。 \end{tabularx} \\ \\ \relax
17:23 & \begin{tabularx}{0.7\textwidth}{X} 我在他們裡面，你在我裡面，使他們完完全全合而為一，讓世人知道是你差我來的，也知道你愛他們，如同愛我一樣。 \end{tabularx} \\ \\ \relax
17:24 & \begin{tabularx}{0.7\textwidth}{X} 父啊，我在哪裡，願你所賜給我的人也同我在哪裡，使他們看見你所賜給我的榮耀，因為創世以前，你已經愛我了。 \end{tabularx} \\ \\ \relax
17:25 & \begin{tabularx}{0.7\textwidth}{X} 公義的父啊，世人未曾認識你，我卻認識你，這些人也知道是你差我來的。 \end{tabularx} \\ \\ \relax
17:26 & \begin{tabularx}{0.7\textwidth}{X} 我已讓他們認識你的名，還要讓他們認識，好讓你愛我的愛在他們裡面，我也在他們裡面。」 \end{tabularx} \\ \\ \relax
18:1 & \begin{tabularx}{0.7\textwidth}{X} 耶穌說了這些話，就同門徒出去，過了汲淪溪。在那裡有一個園子，他和門徒進去了。 \end{tabularx} \\ \\ \relax
18:2 & \begin{tabularx}{0.7\textwidth}{X} 出賣耶穌的猶大也知道那地方，因為耶穌和門徒屢次在那裡聚集。 \end{tabularx} \\ \\ \relax
18:3 & \begin{tabularx}{0.7\textwidth}{X} 猶大領了一隊兵，以及祭司長和法利賽人的聖殿警衛，拿著燈籠、火把和兵器來到園裡。 \end{tabularx} \\ \\ \relax
18:4 & \begin{tabularx}{0.7\textwidth}{X} 耶穌知道將要臨到自己的一切事，就出來對他們說：「你們找誰？」 \end{tabularx} \\ \\ \relax
18:5 & \begin{tabularx}{0.7\textwidth}{X} 他們回答他：「拿撒勒人耶穌。」耶穌對他們說：「我就是。」出賣他的猶大也同他們站在一起。 \end{tabularx} \\ \\ \relax
18:6 & \begin{tabularx}{0.7\textwidth}{X} 耶穌一對他們說「我就是」，他們就退後，倒在地上。 \end{tabularx} \\ \\ \relax
18:7 & \begin{tabularx}{0.7\textwidth}{X} 他又問他們：「你們找誰？」他們說：「拿撒勒人耶穌。」 \end{tabularx} \\ \\ \relax
18:8 & \begin{tabularx}{0.7\textwidth}{X} 耶穌回答：「我已經告訴你們，我就是。你們若找的是我，就讓這些人走吧。」 \end{tabularx} \\ \\ \relax
18:9 & \begin{tabularx}{0.7\textwidth}{X} 這要應驗耶穌說過的話：「你所賜給我的人，我一個也不失落。」 \end{tabularx} \\ \\ \relax
18:10 & \begin{tabularx}{0.7\textwidth}{X} 西門‧彼得帶著一把刀，就拔出來，把大祭司的僕人砍了一刀，削掉了他的右耳，那僕人名叫馬勒古。 \end{tabularx} \\ \\ \relax
18:11 & \begin{tabularx}{0.7\textwidth}{X} 於是耶穌對彼得說：「收刀入鞘吧！我父給我的杯，我豈可不喝呢？」 \end{tabularx} \\ \\ \relax
18:12 & \begin{tabularx}{0.7\textwidth}{X} 那隊兵、千夫長和猶太人的警衛拿住耶穌，把他捆綁了， \end{tabularx} \\ \\ \relax
18:13 & \begin{tabularx}{0.7\textwidth}{X} 先帶到亞那面前，因為他是那年的大祭司該亞法的岳父。 \end{tabularx} \\ \\ \relax
18:14 & \begin{tabularx}{0.7\textwidth}{X} 這該亞法就是從前向猶太人忠告說「一個人替百姓死是有利的」那個人。 \end{tabularx} \\ \\ \relax
18:15 & \begin{tabularx}{0.7\textwidth}{X} 西門‧彼得跟著耶穌，另一個門徒也跟著；那門徒是大祭司所認識的，他就同耶穌進了大祭司的院子。 \end{tabularx} \\ \\ \relax
18:16 & \begin{tabularx}{0.7\textwidth}{X} 彼得卻站在門外。大祭司所認識的那個門徒出來，對看門的使女說了一聲，就領彼得進去。 \end{tabularx} \\ \\ \relax
18:17 & \begin{tabularx}{0.7\textwidth}{X} 那看門的使女對彼得說：「你不也是這人的門徒嗎？」他說：「我不是。」 \end{tabularx} \\ \\ \relax
18:18 & \begin{tabularx}{0.7\textwidth}{X} 僕人和警衛因為天冷生了炭火，站在那裡取暖；彼得也同他們站著取暖。 \end{tabularx} \\ \\ \relax
18:19 & \begin{tabularx}{0.7\textwidth}{X} 於是，大祭司盤問耶穌有關他的門徒和他教導的事。 \end{tabularx} \\ \\ \relax
18:20 & \begin{tabularx}{0.7\textwidth}{X} 耶穌回答他：「我一向都是公開地對世人講話，我常在會堂和聖殿裡，就是猶太人聚集的地方教導人，我私下並沒有講甚麼。 \end{tabularx} \\ \\ \relax
18:21 & \begin{tabularx}{0.7\textwidth}{X} 你為甚麼問我呢？去問那些聽過我講話的人，我所說的，他們都知道。」 \end{tabularx} \\ \\ \relax
18:22 & \begin{tabularx}{0.7\textwidth}{X} 耶穌說了這些話，旁邊站著的一個警衛打了他一耳光，說：「你這樣回答大祭司嗎？」 \end{tabularx} \\ \\ \relax
18:23 & \begin{tabularx}{0.7\textwidth}{X} 耶穌回答他：「假如我說的不對，你指證不對的地方；假如我說的對，你為甚麼打我呢？」 \end{tabularx} \\ \\ \relax
18:24 & \begin{tabularx}{0.7\textwidth}{X} 於是亞那把耶穌綁著押解到大祭司該亞法那裡。 \end{tabularx} \\ \\ \relax
18:25 & \begin{tabularx}{0.7\textwidth}{X} 西門‧彼得正站著取暖，有人對他說：「你不也是他的門徒嗎？」彼得不承認，說：「我不是。」 \end{tabularx} \\ \\ \relax
18:26 & \begin{tabularx}{0.7\textwidth}{X} 大祭司的一個僕人，是被彼得削掉耳朵那人的親屬，說：「我不是看見你同他在園子裡嗎？」 \end{tabularx} \\ \\ \relax
18:27 & \begin{tabularx}{0.7\textwidth}{X} 彼得又不承認，立刻雞就叫了。 \end{tabularx} \\ \\ \relax
18:28 & \begin{tabularx}{0.7\textwidth}{X} 他們把耶穌從該亞法那裡押解到總督府。那時是清早。他們自己卻不進總督府，恐怕染了污穢，不能吃逾越節的宴席。 \end{tabularx} \\ \\ \relax
18:29 & \begin{tabularx}{0.7\textwidth}{X} 於是彼拉多出來，到他們那裡，說：「你們告這人是為甚麼事呢？」 \end{tabularx} \\ \\ \relax
18:30 & \begin{tabularx}{0.7\textwidth}{X} 他們回答他說：「這人若不作惡，我們就不會把他交給你了。」 \end{tabularx} \\ \\ \relax
18:31 & \begin{tabularx}{0.7\textwidth}{X} 彼拉多對他們說：「你們自己帶他去，按著你們的律法問他吧！」猶太人說：「我們沒有殺人的權柄。」 \end{tabularx} \\ \\ \relax
18:32 & \begin{tabularx}{0.7\textwidth}{X} 這是要應驗耶穌所說，指自己將要怎樣死的話。 \end{tabularx} \\ \\ \relax
18:33 & \begin{tabularx}{0.7\textwidth}{X} 於是彼拉多又進了總督府，叫耶穌來，對他說：「你是猶太人的王嗎？」 \end{tabularx} \\ \\ \relax
18:34 & \begin{tabularx}{0.7\textwidth}{X} 耶穌回答：「這話是你說的，還是別人論到我時對你說的呢？」 \end{tabularx} \\ \\ \relax
18:35 & \begin{tabularx}{0.7\textwidth}{X} 彼拉多回答：「難道我是猶太人嗎？你的同胞和祭司長把你交給我。你做了甚麼事呢？」 \end{tabularx} \\ \\ \relax
18:36 & \begin{tabularx}{0.7\textwidth}{X} 耶穌回答：「我的國不屬於這世界；我的國若屬於這世界，我的部下就會為我戰鬥，使我不至於被交給猶太人。只是我的國不屬於這世界。」 \end{tabularx} \\ \\ \relax
18:37 & \begin{tabularx}{0.7\textwidth}{X} 於是彼拉多對他說：「那麼，你是王了？」耶穌回答：「是你說我是王。我為此而生，也為此來到世界，為了給真理作見證。凡屬真理的人都聽我的話。」 \end{tabularx} \\ \\ \relax
18:38 & \begin{tabularx}{0.7\textwidth}{X} 彼拉多對他說：「真理是甚麼呢？」 \end{tabularx} \\ \\ \relax
18:39 & \begin{tabularx}{0.7\textwidth}{X} 但你們有個規矩，在逾越節要我給你們釋放一個人，你們要我給你們釋放這猶太人的王嗎？」 \end{tabularx} \\ \\ \relax
18:40 & \begin{tabularx}{0.7\textwidth}{X} 他們又再喊著說：「不要這人！要巴拉巴！」這巴拉巴是個強盜。 \end{tabularx} \\ \\ \relax
19:1 & \begin{tabularx}{0.7\textwidth}{X} 於是，彼拉多命令把耶穌帶去鞭打了。 \end{tabularx} \\ \\ \relax
19:2 & \begin{tabularx}{0.7\textwidth}{X} 士兵用荊棘編了冠冕，戴在他頭上，給他穿上紫袍， \end{tabularx} \\ \\ \relax
19:3 & \begin{tabularx}{0.7\textwidth}{X} 又走到他面前，說：「萬歲，猶太人的王！」他們就打他耳光。 \end{tabularx} \\ \\ \relax
19:4 & \begin{tabularx}{0.7\textwidth}{X} 彼拉多又出來對眾人說：「看，我帶他出來見你們，讓你們知道我查不出他有甚麼罪狀。」 \end{tabularx} \\ \\ \relax
19:5 & \begin{tabularx}{0.7\textwidth}{X} 耶穌出來，戴著荊棘冠冕，穿著紫袍。彼拉多對他們說：「看哪，這個人！」 \end{tabularx} \\ \\ \relax
19:6 & \begin{tabularx}{0.7\textwidth}{X} 祭司長和聖殿警衛看見他，就喊著說：「釘十字架！釘十字架！」彼拉多對他們說：「你們自己把他帶去釘十字架吧！我查不出他有甚麼罪狀。」 \end{tabularx} \\ \\ \relax
19:7 & \begin{tabularx}{0.7\textwidth}{X} 猶太人回答他：「我們有律法，按照律法，他是該死的，因為他自以為是神的兒子。」 \end{tabularx} \\ \\ \relax
19:8 & \begin{tabularx}{0.7\textwidth}{X} 彼拉多聽見這話，越發害怕， \end{tabularx} \\ \\ \relax
19:9 & \begin{tabularx}{0.7\textwidth}{X} 又進了總督府，對耶穌說：「你是哪裡來的？」耶穌卻不回答。 \end{tabularx} \\ \\ \relax
19:10 & \begin{tabularx}{0.7\textwidth}{X} 於是彼拉多對他說：「你不對我說話嗎？難道你不知道我有權柄釋放你，也有權柄把你釘十字架嗎？」 \end{tabularx} \\ \\ \relax
19:11 & \begin{tabularx}{0.7\textwidth}{X} 耶穌回答他：「若不是從上頭賜給你的，你就毫無權柄辦我，所以，把我交給你的那人罪更重了。」 \end{tabularx} \\ \\ \relax
19:12 & \begin{tabularx}{0.7\textwidth}{X} 從此，彼拉多想要釋放耶穌，無奈猶太人喊著說：「你若釋放這個人，你就不是凱撒的忠臣。凡自立為王的就是背叛凱撒。」 \end{tabularx} \\ \\ \relax
19:13 & \begin{tabularx}{0.7\textwidth}{X} 彼拉多聽見這些話，就帶耶穌出來，到了一個地方，叫「鋪華石處」，希伯來話叫厄巴大，就在那裡坐堂。 \end{tabularx} \\ \\ \relax
19:14 & \begin{tabularx}{0.7\textwidth}{X} 那日是逾越節的預備日，約在正午。彼拉多對猶太人說：「看哪，你們的王！」 \end{tabularx} \\ \\ \relax
19:15 & \begin{tabularx}{0.7\textwidth}{X} 他們就喊著：「除掉他！除掉他！把他釘十字架！」彼拉多對他們說：「要我把你們的王釘十字架嗎？」祭司長回答：「除了凱撒，我們沒有王。」 \end{tabularx} \\ \\ \relax
19:16 & \begin{tabularx}{0.7\textwidth}{X} 於是彼拉多把耶穌交給他們去釘十字架。 \end{tabularx} \\ \\ \relax
19:17 & \begin{tabularx}{0.7\textwidth}{X} 耶穌背著自己的十字架出來，到了一個地方，名叫「髑髏地」，希伯來話叫各各他。 \end{tabularx} \\ \\ \relax
19:18 & \begin{tabularx}{0.7\textwidth}{X} 他們就在那裡把他釘在十字架上，還有兩個人和他一同被釘，一邊一個，耶穌在中間。 \end{tabularx} \\ \\ \relax
19:19 & \begin{tabularx}{0.7\textwidth}{X} 彼拉多又寫了一個牌子，釘在十字架上，寫的是：「猶太人的王，拿撒勒人耶穌。」 \end{tabularx} \\ \\ \relax
19:20 & \begin{tabularx}{0.7\textwidth}{X} 有許多猶太人念這牌子，因為耶穌被釘十字架的地方靠近城，而且牌子是用希伯來、羅馬、希臘三種文字寫的。 \end{tabularx} \\ \\ \relax
19:21 & \begin{tabularx}{0.7\textwidth}{X} 猶太人的祭司長就對彼拉多說：「不要寫『猶太人的王』，要寫『那人說：我是猶太人的王』。」 \end{tabularx} \\ \\ \relax
19:22 & \begin{tabularx}{0.7\textwidth}{X} 彼拉多回答：「我寫了就寫了。」 \end{tabularx} \\ \\ \relax
19:23 & \begin{tabularx}{0.7\textwidth}{X} 士兵把耶穌釘在十字架上以後，把他的衣服拿來分為四份，每人一份。他們又拿他的內衣，這件內衣沒有縫，是上下一片織成的。 \end{tabularx} \\ \\ \relax
19:24 & \begin{tabularx}{0.7\textwidth}{X} 他們就彼此說：「我們不要撕開，我們抽籤，看是誰的。」這要應驗經上的話說：「他們分了我的外衣，為我的內衣抽籤。」士兵果然做了這些事。 \end{tabularx} \\ \\ \relax
19:25 & \begin{tabularx}{0.7\textwidth}{X} 站在耶穌十字架旁邊的，有他的母親、姨母、革羅罷的妻子馬利亞，和抹大拉的馬利亞。 \end{tabularx} \\ \\ \relax
19:26 & \begin{tabularx}{0.7\textwidth}{X} 耶穌見母親和他所愛的那門徒站在旁邊，就對母親說：「母親，看，你的兒子！」 \end{tabularx} \\ \\ \relax
19:27 & \begin{tabularx}{0.7\textwidth}{X} 又對那門徒說：「看，你的母親！」從那刻起，那門徒就接她到自己家裡去了。 \end{tabularx} \\ \\ \relax
19:28 & \begin{tabularx}{0.7\textwidth}{X} 這事以後，耶穌知道各樣的事已經成了，為使經上的話應驗，就說：「我渴了。」 \end{tabularx} \\ \\ \relax
19:29 & \begin{tabularx}{0.7\textwidth}{X} 有一個盛滿了醋的罐子放在那裡，他們就拿海綿蘸滿了醋，綁在牛膝草上，送到他嘴邊。 \end{tabularx} \\ \\ \relax
19:30 & \begin{tabularx}{0.7\textwidth}{X} 耶穌嘗了那醋，說：「成了！」就低下頭，斷了氣。 \end{tabularx} \\ \\ \relax
19:31 & \begin{tabularx}{0.7\textwidth}{X} 因為這日是預備日，又因為那安息日是個大日子，猶太人就來求彼拉多叫人打斷他們的腿，把他們搬走，免得屍首在安息日留在十字架上。 \end{tabularx} \\ \\ \relax
19:32 & \begin{tabularx}{0.7\textwidth}{X} 於是士兵來，把第一個人的腿，和與耶穌同釘的另一個人的腿，都打斷了。 \end{tabularx} \\ \\ \relax
19:33 & \begin{tabularx}{0.7\textwidth}{X} 當他們來到耶穌那裡，見他已經死了，就沒有打斷他的腿。 \end{tabularx} \\ \\ \relax
19:34 & \begin{tabularx}{0.7\textwidth}{X} 然而有一個士兵拿槍扎他的肋旁，立刻有血和水流出來。 \end{tabularx} \\ \\ \relax
19:35 & \begin{tabularx}{0.7\textwidth}{X} 看見這事的人作了見證—他的見證是真的，他知道自己所說的是真的—好讓你們也信。 \end{tabularx} \\ \\ \relax
19:36 & \begin{tabularx}{0.7\textwidth}{X} 這些事發生，為要應驗經上的話：「他的骨頭一根也不可折斷。」 \end{tabularx} \\ \\ \relax
19:37 & \begin{tabularx}{0.7\textwidth}{X} 另有經文也說：「他們要仰望自己所扎的人。」 \end{tabularx} \\ \\ \relax
19:38 & \begin{tabularx}{0.7\textwidth}{X} 這些事以後，亞利馬太的約瑟來求彼拉多，要把耶穌的身體領去。他是耶穌的門徒，只因怕猶太人，就暗地裡作門徒。彼拉多准許了，他就把耶穌的身體領走。 \end{tabularx} \\ \\ \relax
19:39 & \begin{tabularx}{0.7\textwidth}{X} 尼哥德慕也來了，就是先前夜裡去見耶穌的那位，他帶著約一百斤的沒藥和沉香。 \end{tabularx} \\ \\ \relax
19:40 & \begin{tabularx}{0.7\textwidth}{X} 他們照猶太人喪葬的規矩，用細麻布加上香料，把耶穌的身體裹好了。 \end{tabularx} \\ \\ \relax
19:41 & \begin{tabularx}{0.7\textwidth}{X} 在耶穌釘十字架的地方有一個園子，園子裡有一座新墓穴，是從來沒有葬過人的。 \end{tabularx} \\ \\ \relax
19:42 & \begin{tabularx}{0.7\textwidth}{X} 因為那天是猶太人的預備日，而那墳墓又在附近，他們就把耶穌安放在那裡。 \end{tabularx} \\ \\ \relax
20:1 & \begin{tabularx}{0.7\textwidth}{X} 七日的第一日清早，天還黑的時候，抹大拉的馬利亞來到墳墓，看見石頭已從墳墓挪開了， \end{tabularx} \\ \\ \relax
20:2 & \begin{tabularx}{0.7\textwidth}{X} 就跑來見西門‧彼得和耶穌所愛的那個門徒，對他們說：「有人從墳墓裡把主移走了，我們不知道他們把他放在哪裡。」 \end{tabularx} \\ \\ \relax
20:3 & \begin{tabularx}{0.7\textwidth}{X} 彼得和那門徒就出來，往墳墓去。 \end{tabularx} \\ \\ \relax
20:4 & \begin{tabularx}{0.7\textwidth}{X} 兩個人同跑，那門徒比彼得跑得快，先到了墳墓， \end{tabularx} \\ \\ \relax
20:5 & \begin{tabularx}{0.7\textwidth}{X} 低頭往裡看，看見細麻布還放在那裡，只是沒有進去。 \end{tabularx} \\ \\ \relax
20:6 & \begin{tabularx}{0.7\textwidth}{X} 西門‧彼得隨後也到了，進了墳墓，看見細麻布放在那裡， \end{tabularx} \\ \\ \relax
20:7 & \begin{tabularx}{0.7\textwidth}{X} 又看見耶穌的裹頭巾沒有和細麻布放在一起，是另在一處捲著。 \end{tabularx} \\ \\ \relax
20:8 & \begin{tabularx}{0.7\textwidth}{X} 然後先到墳墓的那門徒也進去，他看見就信了。 \end{tabularx} \\ \\ \relax
20:9 & \begin{tabularx}{0.7\textwidth}{X} 他們還不明白聖經所說耶穌必須從死人中復活的意思。 \end{tabularx} \\ \\ \relax
20:10 & \begin{tabularx}{0.7\textwidth}{X} 於是兩個門徒回自己的住處去了。 \end{tabularx} \\ \\ \relax
20:11 & \begin{tabularx}{0.7\textwidth}{X} 馬利亞卻站在墳墓外面哭。她哭的時候，低頭往墳墓裡看， \end{tabularx} \\ \\ \relax
20:12 & \begin{tabularx}{0.7\textwidth}{X} 看見兩個天使穿著白衣，在安放耶穌身體的地方坐著，一個在頭，一個在腳。 \end{tabularx} \\ \\ \relax
20:13 & \begin{tabularx}{0.7\textwidth}{X} 天使對她說：「婦人，你為甚麼哭？」她對他們說：「因為有人把我主移走了，我不知道他們把他放在哪裡。」 \end{tabularx} \\ \\ \relax
20:14 & \begin{tabularx}{0.7\textwidth}{X} 說了這些話，她轉過身來，看見耶穌站在那裡，卻不知道他是耶穌。 \end{tabularx} \\ \\ \relax
20:15 & \begin{tabularx}{0.7\textwidth}{X} 耶穌問她：「婦人，你為甚麼哭？你找誰？」馬利亞以為他是看園子的，就對他說：「先生，若是你把他移了去，請告訴我，你把他放在哪裡，我去把他移回來。」 \end{tabularx} \\ \\ \relax
20:16 & \begin{tabularx}{0.7\textwidth}{X} 耶穌對她說：「馬利亞。」馬利亞轉過身來，用希伯來話對他說：「拉波尼！」（「拉波尼」就是老師的意思。） \end{tabularx} \\ \\ \relax
20:17 & \begin{tabularx}{0.7\textwidth}{X} 耶穌對她說：「不要拉住我，因為我還沒有升上去見我的父。你到我弟兄那裡去告訴他們，我要升上去見我的父，也是你們的父，見我的神，也是你們的神。」 \end{tabularx} \\ \\ \relax
20:18 & \begin{tabularx}{0.7\textwidth}{X} 抹大拉的馬利亞就向門徒報信：「我已經看見了主。」她又把主對她說的話告訴他們。 \end{tabularx} \\ \\ \relax
20:19 & \begin{tabularx}{0.7\textwidth}{X} 那日（就是七日的第一日）晚上，門徒因怕猶太人，所在的地方門都關了。耶穌來，站在當中，對他們說：「願你們平安！」 \end{tabularx} \\ \\ \relax
20:20 & \begin{tabularx}{0.7\textwidth}{X} 說了這話，他把手和肋旁給他們看。門徒一看見主就喜樂了。 \end{tabularx} \\ \\ \relax
20:21 & \begin{tabularx}{0.7\textwidth}{X} 於是耶穌又對他們說：「願你們平安！父怎樣差遣了我，我也照樣差遣你們。」 \end{tabularx} \\ \\ \relax
20:22 & \begin{tabularx}{0.7\textwidth}{X} 說了這話，他向他們吹一口氣，說：「領受聖靈吧！ \end{tabularx} \\ \\ \relax
20:23 & \begin{tabularx}{0.7\textwidth}{X} 你們赦免誰的罪，誰的罪就得赦免；你們不赦免誰的罪，誰的罪就不得赦免。」 \end{tabularx} \\ \\ \relax
20:24 & \begin{tabularx}{0.7\textwidth}{X} 那十二使徒中，有個叫低土馬的多馬，耶穌來的時候，他沒有和他們在一起。 \end{tabularx} \\ \\ \relax
20:25 & \begin{tabularx}{0.7\textwidth}{X} 其他的門徒就對他說：「我們已經看見主了。」多馬卻對他們說：「除非我看見他手上的釘痕，用我的指頭探入那釘痕，用我的手探入他的肋旁，我絕不信。」 \end{tabularx} \\ \\ \relax
20:26 & \begin{tabularx}{0.7\textwidth}{X} 過了八日，門徒又在屋裡，多馬也和他們在一起。門都關了，耶穌來，站在當中，說：「願你們平安！」 \end{tabularx} \\ \\ \relax
20:27 & \begin{tabularx}{0.7\textwidth}{X} 然後他對多馬說：「把你的指頭伸到這裡來，看看我的手；把你的手伸過來，探入我的肋旁。不要疑惑，總要信！」 \end{tabularx} \\ \\ \relax
20:28 & \begin{tabularx}{0.7\textwidth}{X} 多馬回答，對他說：「我的主！我的神！」 \end{tabularx} \\ \\ \relax
20:29 & \begin{tabularx}{0.7\textwidth}{X} 耶穌對他說：「你因為看見了我才信嗎？那沒有看見卻信的有福了。」 \end{tabularx} \\ \\ \relax
20:30 & \begin{tabularx}{0.7\textwidth}{X} 耶穌在他門徒面前另外行了許多神蹟，沒有記錄在這書上。 \end{tabularx} \\ \\ \relax
20:31 & \begin{tabularx}{0.7\textwidth}{X} 但記載這些事是要使你們信耶穌是基督，是神的兒子，並且使你們信他，好因著他的名得生命。 \end{tabularx} \\ \\ \relax
21:1 & \begin{tabularx}{0.7\textwidth}{X} 這些事以後，耶穌在提比哩亞海邊又向門徒顯現。他怎樣顯現記在下面。 \end{tabularx} \\ \\ \relax
21:2 & \begin{tabularx}{0.7\textwidth}{X} 西門‧彼得、叫低土馬的多馬、加利利的迦拿人拿但業、西庇太的兩個兒子，和另外兩個門徒，都在一起。 \end{tabularx} \\ \\ \relax
21:3 & \begin{tabularx}{0.7\textwidth}{X} 西門‧彼得對他們說：「我打魚去。」他們對他說：「我們也和你一起去。」他們就出去，上了船；那一夜並沒有打著甚麼。 \end{tabularx} \\ \\ \relax
21:4 & \begin{tabularx}{0.7\textwidth}{X} 天剛亮的時候，耶穌站在岸上，門徒卻不知道他是耶穌。 \end{tabularx} \\ \\ \relax
21:5 & \begin{tabularx}{0.7\textwidth}{X} 耶穌就對他們說：「孩子們！你們有吃的沒有？」他們回答他：「沒有。」 \end{tabularx} \\ \\ \relax
21:6 & \begin{tabularx}{0.7\textwidth}{X} 耶穌對他們說：「你們把網撒在船的右邊，就會得到。」於是他們撒下網去，竟拉不上來了，因為魚很多。 \end{tabularx} \\ \\ \relax
21:7 & \begin{tabularx}{0.7\textwidth}{X} 耶穌所愛的那門徒對彼得說：「是主！」那時西門‧彼得赤著身子，一聽見是主，就束上外衣，跳進海裡。 \end{tabularx} \\ \\ \relax
21:8 & \begin{tabularx}{0.7\textwidth}{X} 其餘的門徒因離岸不遠，約有二百肘，就坐著小船把那網魚拉過來。 \end{tabularx} \\ \\ \relax
21:9 & \begin{tabularx}{0.7\textwidth}{X} 他們上了岸，看見那裡有炭火，上面有魚和餅。 \end{tabularx} \\ \\ \relax
21:10 & \begin{tabularx}{0.7\textwidth}{X} 耶穌對他們說：「把剛才打的魚拿幾條來。」 \end{tabularx} \\ \\ \relax
21:11 & \begin{tabularx}{0.7\textwidth}{X} 西門‧彼得就上船，把網拉到岸上，網裡滿了大魚，共一百五十三條；雖然魚這樣多，網卻沒有破。 \end{tabularx} \\ \\ \relax
21:12 & \begin{tabularx}{0.7\textwidth}{X} 耶穌對他們說：「你們來吃早飯。」門徒中沒有一個敢問他：「你是誰？」因為他們知道他是主。 \end{tabularx} \\ \\ \relax
21:13 & \begin{tabularx}{0.7\textwidth}{X} 耶穌走過來，拿餅給他們，也照樣拿魚給他們。 \end{tabularx} \\ \\ \relax
21:14 & \begin{tabularx}{0.7\textwidth}{X} 耶穌從死人中復活後向門徒顯現，這是第三次。 \end{tabularx} \\ \\ \relax
21:15 & \begin{tabularx}{0.7\textwidth}{X} 他們吃完了早飯，耶穌對西門‧彼得說：「約翰的兒子西門，你愛我比這些更深嗎？」彼得對他說：「主啊，是的，你知道我愛你。」耶穌說：「你餵養我的小羊。」 \end{tabularx} \\ \\ \relax
21:16 & \begin{tabularx}{0.7\textwidth}{X} 耶穌第二次又對他說：「約翰的兒子西門，你愛我嗎？」彼得對他說：「主啊，是的，你知道我愛你。」耶穌說：「你牧養我的羊。」 \end{tabularx} \\ \\ \relax
21:17 & \begin{tabularx}{0.7\textwidth}{X} 耶穌第三次對他說：「約翰的兒子西門，你愛我嗎？」彼得因為耶穌第三次對他說「你愛我嗎」，就憂愁，對耶穌說：「主啊，你無所不知，你知道我愛你。」耶穌說：「你餵養我的羊。 \end{tabularx} \\ \\ \relax
21:18 & \begin{tabularx}{0.7\textwidth}{X} 我實實在在地告訴你，你年輕的時候，自己束上帶子，隨意往來；但年老的時候，你要伸出手來，別人要把你束上，帶你到不願意去的地方。」 \end{tabularx} \\ \\ \relax
21:19 & \begin{tabularx}{0.7\textwidth}{X} 耶穌說這話，是指彼得會怎樣死來榮耀神。說了這話，耶穌對他說：「你跟從我吧！」 \end{tabularx} \\ \\ \relax
21:20 & \begin{tabularx}{0.7\textwidth}{X} 彼得轉過身來，看見耶穌所愛的那門徒跟著，就是在晚餐時靠著耶穌胸膛說「主啊，出賣你的是誰」的那門徒。 \end{tabularx} \\ \\ \relax
21:21 & \begin{tabularx}{0.7\textwidth}{X} 彼得看見他，就問耶穌：「主啊，這個人怎樣呢？」 \end{tabularx} \\ \\ \relax
21:22 & \begin{tabularx}{0.7\textwidth}{X} 耶穌對他說：「假如我要他等到我來的時候還在，跟你有甚麼關係呢？你跟從我吧！」 \end{tabularx} \\ \\ \relax
21:23 & \begin{tabularx}{0.7\textwidth}{X} 於是這話在弟兄中間流傳，說那門徒不死。其實，耶穌不是說他不死，而是對彼得說：「假如我要他等到我來的時候還在，跟你有甚麼關係呢？」 \end{tabularx} \\ \\ \relax
21:24 & \begin{tabularx}{0.7\textwidth}{X} 這門徒就是為這些事作見證、並且記載這些事的，我們知道他的見證是真的。 \end{tabularx} \\ \\ \relax
21:25 & \begin{tabularx}{0.7\textwidth}{X} 耶穌所行的事還有許多，若是一一都寫出來，我想，就是全世界也容不下所要寫的書。 \end{tabularx} \\ \\
[1ex]
\hline
\hline
\end{longtable}


\newpage

\section{第63屆~港九培靈會~黃子嘉~第六講}
\label{sec:1049}
\textbf{奉獻的生活}
\newline
\newline
連結: \href{https://www.hkbibleconference.org/session-message/view/1049}{\texttt{https://www.hkbibleconference.org/session-message/view/1049}}
\newline
\newline
\hyperref[sec:1048]{< < < PREV SERMON < < <}
~
\hyperref[sec:toc]{[返目錄]}
~
\hyperref[sec:1063]{> > > NEXT SERMON > > >}
\newline
\newline
羅馬書 12:1-12:2
\newline
\begin{longtable}{cl}
\hline
\hline
章節 & 經文 (和合本修訂版)\\
\hline
12:1 & \begin{tabularx}{0.7\textwidth}{X} 所以，弟兄們，我以神的慈悲勸你們，將身體獻上當作活祭，是聖潔的，是神所喜悅的，你們如此事奉乃是理所當然的。 \end{tabularx} \\ \\ \relax
12:2 & \begin{tabularx}{0.7\textwidth}{X} 不要效法這個世界，只要心意更新而變化，叫你們察驗何為神的善良、純全、可喜悅的旨意。 \end{tabularx} \\ \\
[1ex]
\hline
\hline
\end{longtable}


\newpage

\section{第63屆~港九培靈會~黃子嘉~第七講}
\label{sec:1063}
\textbf{身體的生活}
\newline
\newline
連結: \href{https://www.hkbibleconference.org/session-message/view/1063}{\texttt{https://www.hkbibleconference.org/session-message/view/1063}}
\newline
\newline
\hyperref[sec:1049]{< < < PREV SERMON < < <}
~
\hyperref[sec:toc]{[返目錄]}
~
\hyperref[sec:1064]{> > > NEXT SERMON > > >}
\newline
\newline
羅馬書 12:1-12:8
\newline
\begin{longtable}{cl}
\hline
\hline
章節 & 經文 (和合本修訂版)\\
\hline
12:1 & \begin{tabularx}{0.7\textwidth}{X} 所以，弟兄們，我以神的慈悲勸你們，將身體獻上當作活祭，是聖潔的，是神所喜悅的，你們如此事奉乃是理所當然的。 \end{tabularx} \\ \\ \relax
12:2 & \begin{tabularx}{0.7\textwidth}{X} 不要效法這個世界，只要心意更新而變化，叫你們察驗何為神的善良、純全、可喜悅的旨意。 \end{tabularx} \\ \\ \relax
12:3 & \begin{tabularx}{0.7\textwidth}{X} 我憑著所賜我的恩對你們每一位說：不要把自己看得太高，要照著神所分給各人的信心來衡量，看得合乎中道。 \end{tabularx} \\ \\ \relax
12:4 & \begin{tabularx}{0.7\textwidth}{X} 正如我們一個身子上有好些肢體，肢體也不都有一樣的用處。 \end{tabularx} \\ \\ \relax
12:5 & \begin{tabularx}{0.7\textwidth}{X} 這樣，我們許多人在基督裡是一個身體，互相聯絡作肢體。 \end{tabularx} \\ \\ \relax
12:6 & \begin{tabularx}{0.7\textwidth}{X} 按著所得的恩典，我們各有不同的恩賜：或說預言，要按著信心的程度說預言； \end{tabularx} \\ \\ \relax
12:7 & \begin{tabularx}{0.7\textwidth}{X} 或服事的，要專一服事；或教導的，要專一教導； \end{tabularx} \\ \\ \relax
12:8 & \begin{tabularx}{0.7\textwidth}{X} 或勸勉的，要專一勸勉；施捨的，要誠實；治理的，要殷勤；憐憫人的，要樂意。 \end{tabularx} \\ \\
[1ex]
\hline
\hline
\end{longtable}


\newpage

\section{第63屆~港九培靈會~黃子嘉~第八講}
\label{sec:1064}
\textbf{結果子的生活}
\newline
\newline
連結: \href{https://www.hkbibleconference.org/session-message/view/1064}{\texttt{https://www.hkbibleconference.org/session-message/view/1064}}
\newline
\newline
\hyperref[sec:1063]{< < < PREV SERMON < < <}
~
\hyperref[sec:toc]{[返目錄]}
~
\hyperref[sec:1065]{> > > NEXT SERMON > > >}
\newline
\newline
歌羅西書 1:9-1:10
\newline
\begin{longtable}{cl}
\hline
\hline
章節 & 經文 (和合本修訂版)\\
\hline
1:9 & \begin{tabularx}{0.7\textwidth}{X} 因此，我們自從聽見的日子就不住地為你們禱告和祈求，願你們滿有一切屬靈的智慧和悟性，真正知道神的旨意， \end{tabularx} \\ \\ \relax
1:10 & \begin{tabularx}{0.7\textwidth}{X} 好使你們行事為人對得起主，凡事蒙他喜悅，在一切善事上結果子，對神的認識更有長進。 \end{tabularx} \\ \\
[1ex]
\hline
\hline
\end{longtable}


\newpage

\section{第63屆~港九培靈會~黃子嘉~第九講}
\label{sec:1065}
\textbf{宣教的生活}
\newline
\newline
連結: \href{https://www.hkbibleconference.org/session-message/view/1065}{\texttt{https://www.hkbibleconference.org/session-message/view/1065}}
\newline
\newline
\hyperref[sec:1064]{< < < PREV SERMON < < <}
~
\hyperref[sec:toc]{[返目錄]}
~
\hyperref[sec:1017]{> > > NEXT SERMON > > >}
\newline
\newline
馬太福音 12:18-12:40
\newline
\begin{longtable}{cl}
\hline
\hline
章節 & 經文 (和合本修訂版)\\
\hline
12:18 & \begin{tabularx}{0.7\textwidth}{X} 「看哪，我所揀選的僕人，我所親愛，心所喜悅的；我要將我的靈賜給他，他必將公理傳給外邦。 \end{tabularx} \\ \\ \relax
12:19 & \begin{tabularx}{0.7\textwidth}{X} 他不爭吵，不喧嚷，街上也沒有人聽見他的聲音。 \end{tabularx} \\ \\ \relax
12:20 & \begin{tabularx}{0.7\textwidth}{X} 壓傷的蘆葦，他不折斷，將殘的燈火，他不吹滅，直到他使公理得勝。 \end{tabularx} \\ \\ \relax
12:21 & \begin{tabularx}{0.7\textwidth}{X} 外邦人都要仰望他的名。」 \end{tabularx} \\ \\ \relax
12:22 & \begin{tabularx}{0.7\textwidth}{X} 當時，有人把一個被鬼附，又盲又啞的人帶到耶穌那裡，耶穌醫治他，那啞巴就能說話，又能看見。 \end{tabularx} \\ \\ \relax
12:23 & \begin{tabularx}{0.7\textwidth}{X} 眾人都驚奇，說：「這不是大衛之子嗎？」 \end{tabularx} \\ \\ \relax
12:24 & \begin{tabularx}{0.7\textwidth}{X} 但法利賽人聽見，就說：「這個人趕鬼，無非是靠著鬼王別西卜罷了。」 \end{tabularx} \\ \\ \relax
12:25 & \begin{tabularx}{0.7\textwidth}{X} 耶穌知道他們的心思，就對他們說：「一國自相紛爭，必定荒蕪；一城一家自相紛爭，必立不住。 \end{tabularx} \\ \\ \relax
12:26 & \begin{tabularx}{0.7\textwidth}{X} 若撒但趕出撒但，就是自相紛爭，他的國怎能立得住呢？ \end{tabularx} \\ \\ \relax
12:27 & \begin{tabularx}{0.7\textwidth}{X} 我若靠著別西卜趕鬼，你們的子弟趕鬼又靠著誰呢？這樣，他們要作你們的判官。 \end{tabularx} \\ \\ \relax
12:28 & \begin{tabularx}{0.7\textwidth}{X} 我若靠著神的靈趕鬼，那麼，神的國就已臨到你們了。 \end{tabularx} \\ \\ \relax
12:29 & \begin{tabularx}{0.7\textwidth}{X} 人怎能進壯士家裡搶奪他的東西呢？除非先綁住那壯士，否則無法搶奪他的家。 \end{tabularx} \\ \\ \relax
12:30 & \begin{tabularx}{0.7\textwidth}{X} 不跟我一起的，就是反對我；不與我一起收聚的，就是在拆散。 \end{tabularx} \\ \\ \relax
12:31 & \begin{tabularx}{0.7\textwidth}{X} 所以我告訴你們，人一切的罪和褻瀆的話都可得赦免，但是褻瀆聖靈，總不得赦免。 \end{tabularx} \\ \\ \relax
12:32 & \begin{tabularx}{0.7\textwidth}{X} 凡說話干犯人子的，還可得赦免；但是說話干犯聖靈的，今世來世總不得赦免。」 \end{tabularx} \\ \\ \relax
12:33 & \begin{tabularx}{0.7\textwidth}{X} 「你們知道樹好，果子也好；又知道樹壞，果子也壞；因為看果子就可以知道樹。 \end{tabularx} \\ \\ \relax
12:34 & \begin{tabularx}{0.7\textwidth}{X} 毒蛇的孽種啊，你們既是惡人，怎能說出好話來呢？因為心裡所充滿的，口裡就說出來。 \end{tabularx} \\ \\ \relax
12:35 & \begin{tabularx}{0.7\textwidth}{X} 善人從他所存的善發出善來；惡人從他所存的惡發出惡來。 \end{tabularx} \\ \\ \relax
12:36 & \begin{tabularx}{0.7\textwidth}{X} 我告訴你們，凡是人所說的閒話，在審判的日子，要句句供出來； \end{tabularx} \\ \\ \relax
12:37 & \begin{tabularx}{0.7\textwidth}{X} 因為要憑你的話定你為義，也要憑你的話定你有罪。」 \end{tabularx} \\ \\ \relax
12:38 & \begin{tabularx}{0.7\textwidth}{X} 當時，有幾個文士和法利賽人對耶穌說：「老師，我們想請你顯個神蹟給我們看看。」 \end{tabularx} \\ \\ \relax
12:39 & \begin{tabularx}{0.7\textwidth}{X} 耶穌回答他們：「邪惡淫亂的世代求看神蹟，除了先知約拿的神蹟以外，再沒有神蹟給他們看了。 \end{tabularx} \\ \\ \relax
12:40 & \begin{tabularx}{0.7\textwidth}{X} 約拿三日三夜在大魚肚腹中，同樣，人子也要三日三夜在地裡面。 \end{tabularx} \\ \\
[1ex]
\hline
\hline
\end{longtable}


\newpage

\section{第64屆~港九培靈會~張慕皚~第一講}
\label{sec:1017}
\textbf{虛心 — 靈命成長的起點}
\newline
\newline
連結: \href{https://www.hkbibleconference.org/session-message/view/1017}{\texttt{https://www.hkbibleconference.org/session-message/view/1017}}
\newline
\newline
\hyperref[sec:1065]{< < < PREV SERMON < < <}
~
\hyperref[sec:toc]{[返目錄]}
~
\hyperref[sec:1018]{> > > NEXT SERMON > > >}
\newline
\newline
馬太福音 5:1-5:3
\newline
\begin{longtable}{cl}
\hline
\hline
章節 & 經文 (和合本修訂版)\\
\hline
5:1 & \begin{tabularx}{0.7\textwidth}{X} 耶穌看見這一群人，就上了山，坐下後，門徒到他跟前來， \end{tabularx} \\ \\ \relax
5:2 & \begin{tabularx}{0.7\textwidth}{X} 他開口教導他們說： \end{tabularx} \\ \\ \relax
5:3 & \begin{tabularx}{0.7\textwidth}{X} 「心靈貧窮的人有福了！因為天國是他們的。 \end{tabularx} \\ \\
[1ex]
\hline
\hline
\end{longtable}


\newpage

\section{第64屆~港九培靈會~張慕皚~第二講}
\label{sec:1018}
\textbf{哀慟 — 洞悉人性的真面目}
\newline
\newline
連結: \href{https://www.hkbibleconference.org/session-message/view/1018}{\texttt{https://www.hkbibleconference.org/session-message/view/1018}}
\newline
\newline
\hyperref[sec:1017]{< < < PREV SERMON < < <}
~
\hyperref[sec:toc]{[返目錄]}
~
\hyperref[sec:1019]{> > > NEXT SERMON > > >}
\newline
\newline
馬太福音 5:4-5:4
\newline
\begin{longtable}{cl}
\hline
\hline
章節 & 經文 (和合本修訂版)\\
\hline
5:4 & \begin{tabularx}{0.7\textwidth}{X} 哀慟的人有福了！因為他們必得安慰。 \end{tabularx} \\ \\
[1ex]
\hline
\hline
\end{longtable}


\newpage

\section{第64屆~港九培靈會~張慕皚~第三講}
\label{sec:1019}
\textbf{溫柔 — 合乎中度的自我形象}
\newline
\newline
連結: \href{https://www.hkbibleconference.org/session-message/view/1019}{\texttt{https://www.hkbibleconference.org/session-message/view/1019}}
\newline
\newline
\hyperref[sec:1018]{< < < PREV SERMON < < <}
~
\hyperref[sec:toc]{[返目錄]}
~
\hyperref[sec:1020]{> > > NEXT SERMON > > >}
\newline
\newline
馬太福音 5:5-5:5
\newline
\begin{longtable}{cl}
\hline
\hline
章節 & 經文 (和合本修訂版)\\
\hline
5:5 & \begin{tabularx}{0.7\textwidth}{X} 謙和的人有福了！因為他們必承受土地。 \end{tabularx} \\ \\
[1ex]
\hline
\hline
\end{longtable}


\newpage

\section{第64屆~港九培靈會~張慕皚~第四講}
\label{sec:1020}
\textbf{飢渴慕義 — 尋求活在神的旨意中}
\newline
\newline
連結: \href{https://www.hkbibleconference.org/session-message/view/1020}{\texttt{https://www.hkbibleconference.org/session-message/view/1020}}
\newline
\newline
\hyperref[sec:1019]{< < < PREV SERMON < < <}
~
\hyperref[sec:toc]{[返目錄]}
~
\hyperref[sec:1021]{> > > NEXT SERMON > > >}
\newline
\newline
馬太福音 5:6-5:6
\newline
\begin{longtable}{cl}
\hline
\hline
章節 & 經文 (和合本修訂版)\\
\hline
5:6 & \begin{tabularx}{0.7\textwidth}{X} 飢渴慕義的人有福了！因為他們必得飽足。 \end{tabularx} \\ \\
[1ex]
\hline
\hline
\end{longtable}


\newpage

\section{第64屆~港九培靈會~張慕皚~第五講}
\label{sec:1021}
\textbf{憐恤 — 愛的實踐}
\newline
\newline
連結: \href{https://www.hkbibleconference.org/session-message/view/1021}{\texttt{https://www.hkbibleconference.org/session-message/view/1021}}
\newline
\newline
\hyperref[sec:1020]{< < < PREV SERMON < < <}
~
\hyperref[sec:toc]{[返目錄]}
~
\hyperref[sec:1022]{> > > NEXT SERMON > > >}
\newline
\newline
馬太福音 5:2-5:2
\newline
\begin{longtable}{cl}
\hline
\hline
章節 & 經文 (和合本修訂版)\\
\hline
5:2 & \begin{tabularx}{0.7\textwidth}{X} 他開口教導他們說： \end{tabularx} \\ \\
[1ex]
\hline
\hline
\end{longtable}


\newpage

\section{第64屆~港九培靈會~張慕皚~第六講}
\label{sec:1022}
\textbf{清心 — 邁向屬靈最高境界}
\newline
\newline
連結: \href{https://www.hkbibleconference.org/session-message/view/1022}{\texttt{https://www.hkbibleconference.org/session-message/view/1022}}
\newline
\newline
\hyperref[sec:1021]{< < < PREV SERMON < < <}
~
\hyperref[sec:toc]{[返目錄]}
~
\hyperref[sec:1023]{> > > NEXT SERMON > > >}
\newline
\newline
馬太福音 5:8-5:8
\newline
\begin{longtable}{cl}
\hline
\hline
章節 & 經文 (和合本修訂版)\\
\hline
5:8 & \begin{tabularx}{0.7\textwidth}{X} 清心的人有福了！因為他們必得見神。 \end{tabularx} \\ \\
[1ex]
\hline
\hline
\end{longtable}


\newpage

\section{第64屆~港九培靈會~張慕皚~第七講}
\label{sec:1023}
\textbf{使人和睦 — 跟隨基督的腳蹤行}
\newline
\newline
連結: \href{https://www.hkbibleconference.org/session-message/view/1023}{\texttt{https://www.hkbibleconference.org/session-message/view/1023}}
\newline
\newline
\hyperref[sec:1022]{< < < PREV SERMON < < <}
~
\hyperref[sec:toc]{[返目錄]}
~
\hyperref[sec:1024]{> > > NEXT SERMON > > >}
\newline
\newline
馬太福音 5:9-5:9
\newline
\begin{longtable}{cl}
\hline
\hline
章節 & 經文 (和合本修訂版)\\
\hline
5:9 & \begin{tabularx}{0.7\textwidth}{X} 締造和平的人有福了！因為他們必稱為神的兒子。 \end{tabularx} \\ \\
[1ex]
\hline
\hline
\end{longtable}


\newpage

\section{第64屆~港九培靈會~張慕皚~第八講}
\label{sec:1024}
\textbf{為義受逼迫 — 背十字架跟隨主}
\newline
\newline
連結: \href{https://www.hkbibleconference.org/session-message/view/1024}{\texttt{https://www.hkbibleconference.org/session-message/view/1024}}
\newline
\newline
\hyperref[sec:1023]{< < < PREV SERMON < < <}
~
\hyperref[sec:toc]{[返目錄]}
~
\hyperref[sec:1025]{> > > NEXT SERMON > > >}
\newline
\newline
馬太福音 5:10-5:12
\newline
\begin{longtable}{cl}
\hline
\hline
章節 & 經文 (和合本修訂版)\\
\hline
5:10 & \begin{tabularx}{0.7\textwidth}{X} 為義受迫害的人有福了！因為天國是他們的。 \end{tabularx} \\ \\ \relax
5:11 & \begin{tabularx}{0.7\textwidth}{X} 「人若因我辱罵你們，迫害你們，捏造各樣壞話毀謗你們，你們就有福了！ \end{tabularx} \\ \\ \relax
5:12 & \begin{tabularx}{0.7\textwidth}{X} 要歡喜快樂，因為你們在天上的賞賜是很多的。在你們以前的先知，人也是這樣迫害他們。」 \end{tabularx} \\ \\
[1ex]
\hline
\hline
\end{longtable}


\newpage

\section{第64屆~港九培靈會~戴紹曾~第一講}
\label{sec:1025}
\textbf{誠心盼望得見耶穌}
\newline
\newline
連結: \href{https://www.hkbibleconference.org/session-message/view/1025}{\texttt{https://www.hkbibleconference.org/session-message/view/1025}}
\newline
\newline
\hyperref[sec:1024]{< < < PREV SERMON < < <}
~
\hyperref[sec:toc]{[返目錄]}
~
\hyperref[sec:1026]{> > > NEXT SERMON > > >}
\newline
\newline
路加福音 2:25-2:38
\newline
\begin{longtable}{cl}
\hline
\hline
章節 & 經文 (和合本修訂版)\\
\hline
2:25 & \begin{tabularx}{0.7\textwidth}{X} 那時，在耶路撒冷有一個人，名叫西面；這人又公義又虔誠，素常盼望以色列的安慰者來到，又有聖靈在他身上。 \end{tabularx} \\ \\ \relax
2:26 & \begin{tabularx}{0.7\textwidth}{X} 他得了聖靈的啟示，知道自己未死以前必看見主所立的基督。 \end{tabularx} \\ \\ \relax
2:27 & \begin{tabularx}{0.7\textwidth}{X} 他受了聖靈的感動，進入聖殿，正遇見耶穌的父母抱著孩子進來，要照律法的規矩而行。 \end{tabularx} \\ \\ \relax
2:28 & \begin{tabularx}{0.7\textwidth}{X} 西面就把他抱過來，稱頌神說： \end{tabularx} \\ \\ \relax
2:29 & \begin{tabularx}{0.7\textwidth}{X} 「主啊，如今可以照你的話，容你的僕人安然去世； \end{tabularx} \\ \\ \relax
2:30 & \begin{tabularx}{0.7\textwidth}{X} 因為我的眼睛已經看見你的救恩， \end{tabularx} \\ \\ \relax
2:31 & \begin{tabularx}{0.7\textwidth}{X} 就是你在萬民面前所預備的： \end{tabularx} \\ \\ \relax
2:32 & \begin{tabularx}{0.7\textwidth}{X} 是啟示外邦人的光，是你民以色列的榮耀。」 \end{tabularx} \\ \\ \relax
2:33 & \begin{tabularx}{0.7\textwidth}{X} 孩子的父母因論耶穌的這些話就驚訝。 \end{tabularx} \\ \\ \relax
2:34 & \begin{tabularx}{0.7\textwidth}{X} 西面給他們祝福，又對孩子的母親馬利亞說：「這孩子被立，是要叫以色列中許多人跌倒，許多人興起；又要成為毀謗的對象， \end{tabularx} \\ \\ \relax
2:35 & \begin{tabularx}{0.7\textwidth}{X} 叫許多人心裡的意念顯露出來；你自己的心也要被劍刺透。」 \end{tabularx} \\ \\ \relax
2:36 & \begin{tabularx}{0.7\textwidth}{X} 又有位女先知，名叫亞拿，是亞設支派法內力的女兒，年紀已經老邁，從童女出嫁，同丈夫住了七年， \end{tabularx} \\ \\ \relax
2:37 & \begin{tabularx}{0.7\textwidth}{X} 就寡居了，現在已經八十四歲。她不離開聖殿，禁食祈求，晝夜事奉神。 \end{tabularx} \\ \\ \relax
2:38 & \begin{tabularx}{0.7\textwidth}{X} 正當那時，她進前來感謝神，對一切盼望耶路撒冷得救贖的人講論這孩子的事。 \end{tabularx} \\ \\
[1ex]
\hline
\hline
\end{longtable}


\newpage

\section{第64屆~港九培靈會~戴紹曾~第二講}
\label{sec:1026}
\textbf{博士遠來拜見耶穌}
\newline
\newline
連結: \href{https://www.hkbibleconference.org/session-message/view/1026}{\texttt{https://www.hkbibleconference.org/session-message/view/1026}}
\newline
\newline
\hyperref[sec:1025]{< < < PREV SERMON < < <}
~
\hyperref[sec:toc]{[返目錄]}
~
\hyperref[sec:1027]{> > > NEXT SERMON > > >}
\newline
\newline
馬太福音 2:1-2:12
\newline
\begin{longtable}{cl}
\hline
\hline
章節 & 經文 (和合本修訂版)\\
\hline
2:1 & \begin{tabularx}{0.7\textwidth}{X} 在希律作王的時候，耶穌生在猶太的伯利恆。有幾個博學之士從東方來到耶路撒冷，說： \end{tabularx} \\ \\ \relax
2:2 & \begin{tabularx}{0.7\textwidth}{X} 「那生下來作猶太人之王的在哪裡？我們在東方看見他的星，特來拜他。」 \end{tabularx} \\ \\ \relax
2:3 & \begin{tabularx}{0.7\textwidth}{X} 希律王聽見了，就心裡不安；耶路撒冷全城的人也都不安。 \end{tabularx} \\ \\ \relax
2:4 & \begin{tabularx}{0.7\textwidth}{X} 他就召集了祭司長和民間的文士，問他們：「基督該生在哪裡？」 \end{tabularx} \\ \\ \relax
2:5 & \begin{tabularx}{0.7\textwidth}{X} 他們說：「在猶太的伯利恆。因為有先知記著： \end{tabularx} \\ \\ \relax
2:6 & \begin{tabularx}{0.7\textwidth}{X} 『猶大地的伯利恆啊，你在猶大諸城中並不是最小的；因為將來有一位統治者要從你那裡出來，牧養我以色列民。』」 \end{tabularx} \\ \\ \relax
2:7 & \begin{tabularx}{0.7\textwidth}{X} 於是，希律暗地裡召了博學之士來，查問那星是甚麼時候出現的， \end{tabularx} \\ \\ \relax
2:8 & \begin{tabularx}{0.7\textwidth}{X} 就派他們往伯利恆去，說：「你們去仔細尋訪那小孩子，找到了就來報信，我也好去拜他。」 \end{tabularx} \\ \\ \relax
2:9 & \begin{tabularx}{0.7\textwidth}{X} 他們聽了王的話就去了。忽然，在東方所看到的那顆星在前面引領他們，一直行到小孩子所在地方的上方就停住了。 \end{tabularx} \\ \\ \relax
2:10 & \begin{tabularx}{0.7\textwidth}{X} 他們看見那星，就非常歡喜； \end{tabularx} \\ \\ \relax
2:11 & \begin{tabularx}{0.7\textwidth}{X} 進了房子，看見小孩子和他母親馬利亞，就俯伏拜那小孩子，揭開寶盒，拿出黃金、乳香、沒藥，作為禮物獻給他。 \end{tabularx} \\ \\ \relax
2:12 & \begin{tabularx}{0.7\textwidth}{X} 因為在夢中得到主的指示，不要回去見希律，他們就從別的路回自己的家鄉去了。 \end{tabularx} \\ \\
[1ex]
\hline
\hline
\end{longtable}


\newpage

\section{第64屆~港九培靈會~戴紹曾~第三講}
\label{sec:1027}
\textbf{先鋒引人候見耶穌}
\newline
\newline
連結: \href{https://www.hkbibleconference.org/session-message/view/1027}{\texttt{https://www.hkbibleconference.org/session-message/view/1027}}
\newline
\newline
\hyperref[sec:1026]{< < < PREV SERMON < < <}
~
\hyperref[sec:toc]{[返目錄]}
~
\hyperref[sec:1028]{> > > NEXT SERMON > > >}
\newline
\newline
約翰福音 1:29-1:34
\newline
\begin{longtable}{cl}
\hline
\hline
章節 & 經文 (和合本修訂版)\\
\hline
1:29 & \begin{tabularx}{0.7\textwidth}{X} 第二天，約翰看見耶穌來到他那裡，就說：「看哪，神的羔羊，除去世人的罪的！ \end{tabularx} \\ \\ \relax
1:30 & \begin{tabularx}{0.7\textwidth}{X} 這就是我曾說『那在我以後來的先於我，因為在我以前，他已經存在』的那一位。 \end{tabularx} \\ \\ \relax
1:31 & \begin{tabularx}{0.7\textwidth}{X} 我先前不認識他，如今我來用水施洗，為要使他顯明給以色列人。」 \end{tabularx} \\ \\ \relax
1:32 & \begin{tabularx}{0.7\textwidth}{X} 約翰又作見證說：「我曾看見聖靈彷彿鴿子從天降下，停留在他的身上。 \end{tabularx} \\ \\ \relax
1:33 & \begin{tabularx}{0.7\textwidth}{X} 我先前不認識他，可是那差我來用水施洗的對我說：『你看見聖靈降下來，停留在誰的身上，誰就是用聖靈施洗的。』 \end{tabularx} \\ \\ \relax
1:34 & \begin{tabularx}{0.7\textwidth}{X} 我看見了，所以作證：這一位是神的兒子。」 \end{tabularx} \\ \\
[1ex]
\hline
\hline
\end{longtable}


\newpage

\section{第64屆~港九培靈會~戴紹曾~第四講}
\label{sec:1028}
\textbf{夜裡先生訪見耶穌}
\newline
\newline
連結: \href{https://www.hkbibleconference.org/session-message/view/1028}{\texttt{https://www.hkbibleconference.org/session-message/view/1028}}
\newline
\newline
\hyperref[sec:1027]{< < < PREV SERMON < < <}
~
\hyperref[sec:toc]{[返目錄]}
~
\hyperref[sec:1029]{> > > NEXT SERMON > > >}
\newline
\newline
約翰福音 3:1-3:21
\newline
\begin{longtable}{cl}
\hline
\hline
章節 & 經文 (和合本修訂版)\\
\hline
3:1 & \begin{tabularx}{0.7\textwidth}{X} 有一個法利賽人，名叫尼哥德慕，是猶太人的官。 \end{tabularx} \\ \\ \relax
3:2 & \begin{tabularx}{0.7\textwidth}{X} 這人夜裡來見耶穌，對他說：「拉比，我們知道你是由神那裡來作老師的；因為你所行的神蹟，若沒有神同在，無人能行。」 \end{tabularx} \\ \\ \relax
3:3 & \begin{tabularx}{0.7\textwidth}{X} 耶穌回答他說：「我實實在在地告訴你，人若不重生，就不能見神的國。」 \end{tabularx} \\ \\ \relax
3:4 & \begin{tabularx}{0.7\textwidth}{X} 尼哥德慕對他說：「人已經老了，如何能重生呢？豈能再進母腹生出來嗎？」 \end{tabularx} \\ \\ \relax
3:5 & \begin{tabularx}{0.7\textwidth}{X} 耶穌回答：「我實實在在地告訴你，人若不是從水和聖靈生的，就不能進神的國。 \end{tabularx} \\ \\ \relax
3:6 & \begin{tabularx}{0.7\textwidth}{X} 從肉身生的就是肉身；從靈生的就是靈。 \end{tabularx} \\ \\ \relax
3:7 & \begin{tabularx}{0.7\textwidth}{X} 我說『你們必須重生』，你不要驚訝。 \end{tabularx} \\ \\ \relax
3:8 & \begin{tabularx}{0.7\textwidth}{X} 風隨著意思吹，你聽見風的聲音，卻不知道是從哪裡來，往哪裡去；凡從聖靈生的也是如此。」 \end{tabularx} \\ \\ \relax
3:9 & \begin{tabularx}{0.7\textwidth}{X} 尼哥德慕問他：「怎麼能有這些事呢？」 \end{tabularx} \\ \\ \relax
3:10 & \begin{tabularx}{0.7\textwidth}{X} 耶穌回答，對他說：「你是以色列人的老師，還不明白這些事嗎？ \end{tabularx} \\ \\ \relax
3:11 & \begin{tabularx}{0.7\textwidth}{X} 我實實在在地告訴你，我們所說的是我們知道的，我們所見證的是我們見過的，你們卻不領受我們的見證。 \end{tabularx} \\ \\ \relax
3:12 & \begin{tabularx}{0.7\textwidth}{X} 我對你們說地上的事，你們尚且不信，若對你們說天上的事，如何能信呢？ \end{tabularx} \\ \\ \relax
3:13 & \begin{tabularx}{0.7\textwidth}{X} 除了從天降下的人子，沒有人升過天。 \end{tabularx} \\ \\ \relax
3:14 & \begin{tabularx}{0.7\textwidth}{X} 摩西在曠野怎樣舉蛇，人子也必須照樣被舉起來， \end{tabularx} \\ \\ \relax
3:15 & \begin{tabularx}{0.7\textwidth}{X} 要使一切信他的人都得永生。 \end{tabularx} \\ \\ \relax
3:16 & \begin{tabularx}{0.7\textwidth}{X} 「神愛世人，甚至將他獨一的兒子賜給他們，叫一切信他的人不致滅亡，反得永生。 \end{tabularx} \\ \\ \relax
3:17 & \begin{tabularx}{0.7\textwidth}{X} 因為神差他的兒子到世上來，不是要定世人的罪，而是要使世人因他得救。 \end{tabularx} \\ \\ \relax
3:18 & \begin{tabularx}{0.7\textwidth}{X} 信他的人不被定罪；不信的人已經被定罪了，因為他不信神獨一兒子的名。 \end{tabularx} \\ \\ \relax
3:19 & \begin{tabularx}{0.7\textwidth}{X} 光來到世上，世人因自己的行為是惡的，不愛光，倒愛黑暗，這就定了他們的罪。 \end{tabularx} \\ \\ \relax
3:20 & \begin{tabularx}{0.7\textwidth}{X} 凡作惡的人都恨惡光，不來接近光，恐怕他的行為被暴露。 \end{tabularx} \\ \\ \relax
3:21 & \begin{tabularx}{0.7\textwidth}{X} 但實行真理的人就來接近光，為要顯明他的行為是靠神而行的。」 \end{tabularx} \\ \\
[1ex]
\hline
\hline
\end{longtable}


\newpage

\section{第64屆~港九培靈會~戴紹曾~第五講}
\label{sec:1029}
\textbf{井旁婦人遇見耶穌}
\newline
\newline
連結: \href{https://www.hkbibleconference.org/session-message/view/1029}{\texttt{https://www.hkbibleconference.org/session-message/view/1029}}
\newline
\newline
\hyperref[sec:1028]{< < < PREV SERMON < < <}
~
\hyperref[sec:toc]{[返目錄]}
~
\hyperref[sec:1030]{> > > NEXT SERMON > > >}
\newline
\newline
約翰福音 4:27-4:42
\newline
\begin{longtable}{cl}
\hline
\hline
章節 & 經文 (和合本修訂版)\\
\hline
4:27 & \begin{tabularx}{0.7\textwidth}{X} 正在這時，門徒回來了。他們對耶穌正在和一個婦人說話感到驚訝，可是沒有人說：「你要甚麼？」或說：「你為甚麼和她說話？」 \end{tabularx} \\ \\ \relax
4:28 & \begin{tabularx}{0.7\textwidth}{X} 那婦人留下水罐，往城裡去，對眾人說： \end{tabularx} \\ \\ \relax
4:29 & \begin{tabularx}{0.7\textwidth}{X} 「你們來看！有一個人把我素來所做的一切事都說了出來，難道這個人就是基督嗎？」 \end{tabularx} \\ \\ \relax
4:30 & \begin{tabularx}{0.7\textwidth}{X} 他們就出城，來到耶穌那裡。 \end{tabularx} \\ \\ \relax
4:31 & \begin{tabularx}{0.7\textwidth}{X} 就在這個時候，門徒求耶穌說：「拉比，請吃吧。」 \end{tabularx} \\ \\ \relax
4:32 & \begin{tabularx}{0.7\textwidth}{X} 耶穌對他們說：「我有食物吃，是你們不知道的。」 \end{tabularx} \\ \\ \relax
4:33 & \begin{tabularx}{0.7\textwidth}{X} 門徒就彼此說：「難道有人拿甚麼給他吃了嗎？」 \end{tabularx} \\ \\ \relax
4:34 & \begin{tabularx}{0.7\textwidth}{X} 耶穌對他們說：「我的食物就是要遵行差我來那位的旨意，完成他的工作。 \end{tabularx} \\ \\ \relax
4:35 & \begin{tabularx}{0.7\textwidth}{X} 你們不是說『到收割的時候還有四個月』嗎？我告訴你們，舉目向田觀看，莊稼熟了，可以收割了。 \end{tabularx} \\ \\ \relax
4:36 & \begin{tabularx}{0.7\textwidth}{X} 收割的人已經得工錢，為永生儲存五穀，使撒種的和收割的一同快樂。 \end{tabularx} \\ \\ \relax
4:37 & \begin{tabularx}{0.7\textwidth}{X} 『那人撒種，這人收割』，這話可見是真的。 \end{tabularx} \\ \\ \relax
4:38 & \begin{tabularx}{0.7\textwidth}{X} 我差你們去收你們所沒有辛勞的；別人辛勞，你們享受他們辛勞的成果。」 \end{tabularx} \\ \\ \relax
4:39 & \begin{tabularx}{0.7\textwidth}{X} 那城裡有好些撒瑪利亞人信了耶穌，因為那婦人作見證，說：「他把我素來所做的一切事都說了出來。」 \end{tabularx} \\ \\ \relax
4:40 & \begin{tabularx}{0.7\textwidth}{X} 於是撒瑪利亞人來見耶穌，求他在他們那裡住下，他就在那裡住了兩天。 \end{tabularx} \\ \\ \relax
4:41 & \begin{tabularx}{0.7\textwidth}{X} 因為耶穌的話，信的人就更多了。 \end{tabularx} \\ \\ \relax
4:42 & \begin{tabularx}{0.7\textwidth}{X} 他們對那婦人說：「現在我們信，不再是因為你的話，而是我們親自聽見了，知道這人真是世界的救主。」 \end{tabularx} \\ \\
[1ex]
\hline
\hline
\end{longtable}


\newpage

\section{第64屆~港九培靈會~戴紹曾~第六講}
\label{sec:1030}
\textbf{變像山上只見耶穌}
\newline
\newline
連結: \href{https://www.hkbibleconference.org/session-message/view/1030}{\texttt{https://www.hkbibleconference.org/session-message/view/1030}}
\newline
\newline
\hyperref[sec:1029]{< < < PREV SERMON < < <}
~
\hyperref[sec:toc]{[返目錄]}
~
\hyperref[sec:1031]{> > > NEXT SERMON > > >}
\newline
\newline
馬可福音 9:1-9:8
\newline
\begin{longtable}{cl}
\hline
\hline
章節 & 經文 (和合本修訂版)\\
\hline
9:1 & \begin{tabularx}{0.7\textwidth}{X} 耶穌又對他們說：「我實在告訴你們，站在這裡的，有人在沒經歷死亡以前，必定看見神的國帶著能力臨到。」 \end{tabularx} \\ \\ \relax
9:2 & \begin{tabularx}{0.7\textwidth}{X} 過了六天，耶穌帶著彼得、雅各、約翰，領他們悄悄地上了高山。他在他們面前變了形像， \end{tabularx} \\ \\ \relax
9:3 & \begin{tabularx}{0.7\textwidth}{X} 衣服放光，極其潔白，地上漂布的人沒有一個能漂得那樣白。 \end{tabularx} \\ \\ \relax
9:4 & \begin{tabularx}{0.7\textwidth}{X} 有以利亞和摩西向他們顯現，並且與耶穌說話。 \end{tabularx} \\ \\ \relax
9:5 & \begin{tabularx}{0.7\textwidth}{X} 彼得對耶穌說：「拉比，我們在這裡真好！我們來搭三座棚，一座為你，一座為摩西，一座為以利亞。」 \end{tabularx} \\ \\ \relax
9:6 & \begin{tabularx}{0.7\textwidth}{X} 彼得不知道說甚麼才好，因為他們很害怕。 \end{tabularx} \\ \\ \relax
9:7 & \begin{tabularx}{0.7\textwidth}{X} 有一朵雲彩來遮蓋他們，又有聲音從雲彩裡出來，說：「這是我的愛子，你們要聽從他！」 \end{tabularx} \\ \\ \relax
9:8 & \begin{tabularx}{0.7\textwidth}{X} 門徒連忙向周圍觀看，不再看見任何人，只見耶穌同他們在一起。 \end{tabularx} \\ \\
[1ex]
\hline
\hline
\end{longtable}


\newpage

\section{第64屆~港九培靈會~戴紹曾~第七講}
\label{sec:1031}
\textbf{討飯瞎子看見耶穌}
\newline
\newline
連結: \href{https://www.hkbibleconference.org/session-message/view/1031}{\texttt{https://www.hkbibleconference.org/session-message/view/1031}}
\newline
\newline
\hyperref[sec:1030]{< < < PREV SERMON < < <}
~
\hyperref[sec:toc]{[返目錄]}
~
\hyperref[sec:1032]{> > > NEXT SERMON > > >}
\newline
\newline
馬可福音 10:46-10:52
\newline
\begin{longtable}{cl}
\hline
\hline
章節 & 經文 (和合本修訂版)\\
\hline
10:46 & \begin{tabularx}{0.7\textwidth}{X} 他們到了耶利哥。耶穌同門徒並許多人離開耶利哥的時候，有一個討飯的盲人，是底買的兒子巴底買，坐在路旁。 \end{tabularx} \\ \\ \relax
10:47 & \begin{tabularx}{0.7\textwidth}{X} 他聽見是拿撒勒的耶穌，就喊了起來，說：「大衛之子耶穌啊，可憐我吧！」 \end{tabularx} \\ \\ \relax
10:48 & \begin{tabularx}{0.7\textwidth}{X} 有許多人責備他，不許他作聲，他卻越發喊著：「大衛之子啊，可憐我吧！」 \end{tabularx} \\ \\ \relax
10:49 & \begin{tabularx}{0.7\textwidth}{X} 耶穌就站住，說：「叫他過來。」他們就叫那盲人，對他說：「放心，起來！他在叫你啦。」 \end{tabularx} \\ \\ \relax
10:50 & \begin{tabularx}{0.7\textwidth}{X} 盲人就丟下衣服，跳起來，走到耶穌那裡。 \end{tabularx} \\ \\ \relax
10:51 & \begin{tabularx}{0.7\textwidth}{X} 耶穌回答他說：「你要我為你做甚麼？」盲人對他說：「拉波尼，我要能看見。」 \end{tabularx} \\ \\ \relax
10:52 & \begin{tabularx}{0.7\textwidth}{X} 耶穌對他說：「你去吧！你的信救了你。」盲人立刻看得見，就在路上跟隨耶穌。 \end{tabularx} \\ \\
[1ex]
\hline
\hline
\end{longtable}


\newpage

\section{第64屆~港九培靈會~戴紹曾~第八講}
\label{sec:1032}
\textbf{樹上撒該想見耶穌}
\newline
\newline
連結: \href{https://www.hkbibleconference.org/session-message/view/1032}{\texttt{https://www.hkbibleconference.org/session-message/view/1032}}
\newline
\newline
\hyperref[sec:1031]{< < < PREV SERMON < < <}
~
\hyperref[sec:toc]{[返目錄]}
~
\hyperref[sec:1033]{> > > NEXT SERMON > > >}
\newline
\newline
路加福音 19:1-19:10
\newline
\begin{longtable}{cl}
\hline
\hline
章節 & 經文 (和合本修訂版)\\
\hline
19:1 & \begin{tabularx}{0.7\textwidth}{X} 耶穌進了耶利哥，要從那裡經過。 \end{tabularx} \\ \\ \relax
19:2 & \begin{tabularx}{0.7\textwidth}{X} 有一個人名叫撒該，作稅吏長，是個財主。 \end{tabularx} \\ \\ \relax
19:3 & \begin{tabularx}{0.7\textwidth}{X} 他要看看耶穌是怎樣的人，只因人多，他的身材又矮，所以看不見。 \end{tabularx} \\ \\ \relax
19:4 & \begin{tabularx}{0.7\textwidth}{X} 於是他跑到前頭，爬上桑樹，要看耶穌，因為耶穌要從那裡經過。 \end{tabularx} \\ \\ \relax
19:5 & \begin{tabularx}{0.7\textwidth}{X} 耶穌到了那裡，抬頭一看，對他說：「撒該，快下來！今天我必須住在你家裡。」 \end{tabularx} \\ \\ \relax
19:6 & \begin{tabularx}{0.7\textwidth}{X} 他就急忙下來，歡歡喜喜地接待耶穌。 \end{tabularx} \\ \\ \relax
19:7 & \begin{tabularx}{0.7\textwidth}{X} 眾人看見，都私下議論說：「他竟然到罪人家裡去住宿。」 \end{tabularx} \\ \\ \relax
19:8 & \begin{tabularx}{0.7\textwidth}{X} 撒該站著對主說：「主啊，我把所有的一半給窮人；我若勒索了誰，就還他四倍。」 \end{tabularx} \\ \\ \relax
19:9 & \begin{tabularx}{0.7\textwidth}{X} 耶穌對他說：「今天救恩到了這家，因為他也是亞伯拉罕的子孫。 \end{tabularx} \\ \\ \relax
19:10 & \begin{tabularx}{0.7\textwidth}{X} 人子來是要尋找和拯救失喪的人。」 \end{tabularx} \\ \\
[1ex]
\hline
\hline
\end{longtable}


\newpage

\section{第64屆~港九培靈會~戴紹曾~第九講}
\label{sec:1033}
\textbf{希利尼人求見耶穌}
\newline
\newline
連結: \href{https://www.hkbibleconference.org/session-message/view/1033}{\texttt{https://www.hkbibleconference.org/session-message/view/1033}}
\newline
\newline
\hyperref[sec:1032]{< < < PREV SERMON < < <}
~
\hyperref[sec:toc]{[返目錄]}
~
\hyperref[sec:1034]{> > > NEXT SERMON > > >}
\newline
\newline
約翰福音 12:20-12:36
\newline
\begin{longtable}{cl}
\hline
\hline
章節 & 經文 (和合本修訂版)\\
\hline
12:20 & \begin{tabularx}{0.7\textwidth}{X} 那時，上來過節禮拜的人中，有幾個希臘人。 \end{tabularx} \\ \\ \relax
12:21 & \begin{tabularx}{0.7\textwidth}{X} 他們來見加利利的伯賽大人腓力，請求他說：「先生，我們想見耶穌。」 \end{tabularx} \\ \\ \relax
12:22 & \begin{tabularx}{0.7\textwidth}{X} 腓力去告訴安得烈，然後安得烈同腓力去告訴耶穌。 \end{tabularx} \\ \\ \relax
12:23 & \begin{tabularx}{0.7\textwidth}{X} 耶穌回答他們說：「人子得榮耀的時候到了。 \end{tabularx} \\ \\ \relax
12:24 & \begin{tabularx}{0.7\textwidth}{X} 我實實在在地告訴你們，一粒麥子不落在地裡死了，仍舊是一粒；若是死了，就結出許多子粒來。 \end{tabularx} \\ \\ \relax
12:25 & \begin{tabularx}{0.7\textwidth}{X} 愛惜自己性命的，就喪失性命；那恨惡自己在這世上的性命的，要保全性命到永生。 \end{tabularx} \\ \\ \relax
12:26 & \begin{tabularx}{0.7\textwidth}{X} 若有人服事我，就當跟從我；我在哪裡，服事我的人也要在哪裡；若有人服事我，我父必尊重他。」 \end{tabularx} \\ \\ \relax
12:27 & \begin{tabularx}{0.7\textwidth}{X} 「我現在心裡憂愁，我說甚麼才好呢？說『父啊，救我脫離這時候』嗎？但我正是為這時候來的。 \end{tabularx} \\ \\ \relax
12:28 & \begin{tabularx}{0.7\textwidth}{X} 父啊，願你榮耀你的名！」於是有聲音從天上來，說：「我已經榮耀了我的名，還要再榮耀。」 \end{tabularx} \\ \\ \relax
12:29 & \begin{tabularx}{0.7\textwidth}{X} 站在旁邊的眾人聽見，就說：「打雷了。」另有的說：「有天使對他說話。」 \end{tabularx} \\ \\ \relax
12:30 & \begin{tabularx}{0.7\textwidth}{X} 耶穌回答說：「這聲音不是為我，而是為你們來的。 \end{tabularx} \\ \\ \relax
12:31 & \begin{tabularx}{0.7\textwidth}{X} 現在正是這世界受審判的時候；現在這世界的統治者要被趕出去。 \end{tabularx} \\ \\ \relax
12:32 & \begin{tabularx}{0.7\textwidth}{X} 我從地上被舉起來的時候，我要吸引萬人來歸我。」 \end{tabularx} \\ \\ \relax
12:33 & \begin{tabularx}{0.7\textwidth}{X} 耶穌這話是指自己將要怎樣死說的。 \end{tabularx} \\ \\ \relax
12:34 & \begin{tabularx}{0.7\textwidth}{X} 眾人就回答他：「我們聽見律法書上說，基督是永存的；你怎麼說，人子必須被舉起來呢？這人子是誰呢？」 \end{tabularx} \\ \\ \relax
12:35 & \begin{tabularx}{0.7\textwidth}{X} 耶穌對他們說：「光在你們中間為時不多了，應該趁著有光的時候行走，免得黑暗臨到你們；那在黑暗裡行走的，不知道往何處去。 \end{tabularx} \\ \\ \relax
12:36 & \begin{tabularx}{0.7\textwidth}{X} 你們趁著有光，要信從這光，使你們成為光明之子。」 \end{tabularx} \\ \\
[1ex]
\hline
\hline
\end{longtable}


\newpage

\section{第64屆~港九培靈會~戴紹曾~第十講}
\label{sec:1034}
\textbf{復活清晨去見耶穌}
\newline
\newline
連結: \href{https://www.hkbibleconference.org/session-message/view/1034}{\texttt{https://www.hkbibleconference.org/session-message/view/1034}}
\newline
\newline
\hyperref[sec:1033]{< < < PREV SERMON < < <}
~
\hyperref[sec:toc]{[返目錄]}
~
\hyperref[sec:1391]{> > > NEXT SERMON > > >}
\newline
\newline
約翰福音 20:1-20:23
\newline
\begin{longtable}{cl}
\hline
\hline
章節 & 經文 (和合本修訂版)\\
\hline
20:1 & \begin{tabularx}{0.7\textwidth}{X} 七日的第一日清早，天還黑的時候，抹大拉的馬利亞來到墳墓，看見石頭已從墳墓挪開了， \end{tabularx} \\ \\ \relax
20:2 & \begin{tabularx}{0.7\textwidth}{X} 就跑來見西門‧彼得和耶穌所愛的那個門徒，對他們說：「有人從墳墓裡把主移走了，我們不知道他們把他放在哪裡。」 \end{tabularx} \\ \\ \relax
20:3 & \begin{tabularx}{0.7\textwidth}{X} 彼得和那門徒就出來，往墳墓去。 \end{tabularx} \\ \\ \relax
20:4 & \begin{tabularx}{0.7\textwidth}{X} 兩個人同跑，那門徒比彼得跑得快，先到了墳墓， \end{tabularx} \\ \\ \relax
20:5 & \begin{tabularx}{0.7\textwidth}{X} 低頭往裡看，看見細麻布還放在那裡，只是沒有進去。 \end{tabularx} \\ \\ \relax
20:6 & \begin{tabularx}{0.7\textwidth}{X} 西門‧彼得隨後也到了，進了墳墓，看見細麻布放在那裡， \end{tabularx} \\ \\ \relax
20:7 & \begin{tabularx}{0.7\textwidth}{X} 又看見耶穌的裹頭巾沒有和細麻布放在一起，是另在一處捲著。 \end{tabularx} \\ \\ \relax
20:8 & \begin{tabularx}{0.7\textwidth}{X} 然後先到墳墓的那門徒也進去，他看見就信了。 \end{tabularx} \\ \\ \relax
20:9 & \begin{tabularx}{0.7\textwidth}{X} 他們還不明白聖經所說耶穌必須從死人中復活的意思。 \end{tabularx} \\ \\ \relax
20:10 & \begin{tabularx}{0.7\textwidth}{X} 於是兩個門徒回自己的住處去了。 \end{tabularx} \\ \\ \relax
20:11 & \begin{tabularx}{0.7\textwidth}{X} 馬利亞卻站在墳墓外面哭。她哭的時候，低頭往墳墓裡看， \end{tabularx} \\ \\ \relax
20:12 & \begin{tabularx}{0.7\textwidth}{X} 看見兩個天使穿著白衣，在安放耶穌身體的地方坐著，一個在頭，一個在腳。 \end{tabularx} \\ \\ \relax
20:13 & \begin{tabularx}{0.7\textwidth}{X} 天使對她說：「婦人，你為甚麼哭？」她對他們說：「因為有人把我主移走了，我不知道他們把他放在哪裡。」 \end{tabularx} \\ \\ \relax
20:14 & \begin{tabularx}{0.7\textwidth}{X} 說了這些話，她轉過身來，看見耶穌站在那裡，卻不知道他是耶穌。 \end{tabularx} \\ \\ \relax
20:15 & \begin{tabularx}{0.7\textwidth}{X} 耶穌問她：「婦人，你為甚麼哭？你找誰？」馬利亞以為他是看園子的，就對他說：「先生，若是你把他移了去，請告訴我，你把他放在哪裡，我去把他移回來。」 \end{tabularx} \\ \\ \relax
20:16 & \begin{tabularx}{0.7\textwidth}{X} 耶穌對她說：「馬利亞。」馬利亞轉過身來，用希伯來話對他說：「拉波尼！」（「拉波尼」就是老師的意思。） \end{tabularx} \\ \\ \relax
20:17 & \begin{tabularx}{0.7\textwidth}{X} 耶穌對她說：「不要拉住我，因為我還沒有升上去見我的父。你到我弟兄那裡去告訴他們，我要升上去見我的父，也是你們的父，見我的神，也是你們的神。」 \end{tabularx} \\ \\ \relax
20:18 & \begin{tabularx}{0.7\textwidth}{X} 抹大拉的馬利亞就向門徒報信：「我已經看見了主。」她又把主對她說的話告訴他們。 \end{tabularx} \\ \\ \relax
20:19 & \begin{tabularx}{0.7\textwidth}{X} 那日（就是七日的第一日）晚上，門徒因怕猶太人，所在的地方門都關了。耶穌來，站在當中，對他們說：「願你們平安！」 \end{tabularx} \\ \\ \relax
20:20 & \begin{tabularx}{0.7\textwidth}{X} 說了這話，他把手和肋旁給他們看。門徒一看見主就喜樂了。 \end{tabularx} \\ \\ \relax
20:21 & \begin{tabularx}{0.7\textwidth}{X} 於是耶穌又對他們說：「願你們平安！父怎樣差遣了我，我也照樣差遣你們。」 \end{tabularx} \\ \\ \relax
20:22 & \begin{tabularx}{0.7\textwidth}{X} 說了這話，他向他們吹一口氣，說：「領受聖靈吧！ \end{tabularx} \\ \\ \relax
20:23 & \begin{tabularx}{0.7\textwidth}{X} 你們赦免誰的罪，誰的罪就得赦免；你們不赦免誰的罪，誰的罪就不得赦免。」 \end{tabularx} \\ \\
[1ex]
\hline
\hline
\end{longtable}


\newpage

\section{第64屆~港九培靈會~楊濬哲~序}
\label{sec:1391}
\textbf{序}
\newline
\newline
連結: \href{https://www.hkbibleconference.org/session-message/view/1391}{\texttt{https://www.hkbibleconference.org/session-message/view/1391}}
\newline
\newline
\hyperref[sec:1034]{< < < PREV SERMON < < <}
~
\hyperref[sec:toc]{[返目錄]}
~
\hyperref[sec:1009]{> > > NEXT SERMON > > >}
\newline
\newline
第六十四屆港九培靈研經大會,如期於九二年八月一日至十日,假九龍城浸信會舉行.感謝神的恩典,得以順利圓滿閉會.
\newline
\newline
研經會蒙播道神學院院長許道良牧師主講,總題:「朝聖客旅指南──彼得前書今詮」.許牧師指出,在現今世代中,許多基督徒經不起時代的衝擊而失去信仰,漸漸忘記屬天永恆的盼望,這是十分危險的.為此,特別提醒信徒重視自己的身份和使命,從世界分別出來,有活潑之盼望,有不朽壞的生命,有不衰殘的基業,充滿喜樂地為主而活.
\newline
\newline
講道會請建道神學院院長,張慕凱牧師主講,總題:「從八福看靈命的操練」.目的乃希望基督徒認識屬靈的操練,並以聖經及耶穌基督為中心,作為屬靈操練的路向.在這主題中又特別強調:每一個基督徒都應該具備作為天國子民屬靈生命之特質.所追求的道德標準,應以「八福」中的八種美德為準則.
\newline
\newline
奮興會由海外基督使團香港署理主任戴紹曾牧師主講,總題:「我們願意見耶穌.」戴牧師在每晚講道前,都講述戴德生家族的見證,令會眾信心大得激勵! 戴牧師覺得:現今社會政治在改變,經濟不穩定,道德墮落,人心惶惶,沒有歸屬感.因此,我們要尋找耶穌,因為祂是歷史之主宰.在十晚的分題中,從不同的人物,角度,見證,去思想,學習,從而得著寶貴之經驗.盼望耶穌真正成為基督徒的異象.
\newline
\newline
讀者若能細閱本書,自然能獲益良多.
\newline
\newline
大會期間,為應付躍踴赴會者,城浸地方全部開放；除正堂外,加設十三部閉路電視轉播.書攤方面,設有廿三個攤位,方便赴會者購買.
\newline
\newline
在晚堂奮興會中,講員呼召.決志信主者,有廿四人,悔改更新為主而活者八十九人,全時間事奉者九十六人,合共為二○九人.
\newline
\newline
本年內我們雖然痛失元老委辦鄭德音牧師,但也為他的佳美腳蹤與事奉而感謝神!在缺乏人手時,主又為我們預備了樓恩德校長,與柳鎮平牧師答應加入委辦會,使我們能同心事奉主!
\newline
\newline
本講道集,是六十四屆大會的講道記錄.蒙鄭坤英姊妹記錄研經會；吳品嫻姊妹記錄講道會；雷麗端姊妹記錄奮興會；朱建磯牧師整理編校；西門何治能先生設計封面,乃能如期出版,謹此致謝!
\newline
\newline
本屆大會蒙九龍城浸信會,借用堂址為會場；執事,同工,義工,及各弟兄姊妹協助招待,維持秩序.福音傳播中心負責錄音,基督教週報刊登新聞稿及廣告；時代論壇,明報,刊登廣告；神托之家耀安工場,寄發宣傳品；彼此配搭,以廣宣傳,藉增效益.更蒙各教會兄姊,主內同工同道,踴躍赴會及奉獻,統此致謝!
\newline
\newline


\newpage

\section{第64屆~港九培靈會~許道良~第一講}
\label{sec:1009}
\textbf{曙光遙望}
\newline
\newline
連結: \href{https://www.hkbibleconference.org/session-message/view/1009}{\texttt{https://www.hkbibleconference.org/session-message/view/1009}}
\newline
\newline
\hyperref[sec:1391]{< < < PREV SERMON < < <}
~
\hyperref[sec:toc]{[返目錄]}
~
\hyperref[sec:1010]{> > > NEXT SERMON > > >}
\newline
\newline
彼得前書並非一本容易明白的書卷,很多信徒誤以為此書教人面對苦難,為主而死,只合遭受逼迫的國家使用.其實,此書較多篇幅集中於為主而活的課題上,期望此番研經,能把各位的焦點稍稍移開,不單只思想到為主而死更考慮到怎樣為主而活.
\newline
\newline
歷代以來,在很多宗教沒自由的國家裡,不少基督徒經不起為主而死的考驗；出賣主,失卻了信仰,中國教會也有這樣的例子.當然,在中國教會的歷史裡, 不乏有殉道者的血,鮮明地見證了十字架道理的真實.然而我個人憂慮的是；在宗教自由的國家裡,會有更多信徒經不起為主而活的考驗.在未有統計數字的推斷下,我懷疑「為主而活」與「為主而死」兩種考驗,前者的傷亡率可能遠比後者的傷亡率還高.
\newline
\newline
我們看見在宗教自由的國家裡,許多人信了主,但往往又把持不定,很輕易的就在世俗的衝擊下放棄了信仰,出賣了主.彼得前書雖是彼得在第一世紀寫給小亞細亞教會的信,作者的時間距今差不多有二千年之久,但其中所教導和勸勉的,卻是教會今日要正視的問題.譬如書中說到永恆的焦點,如何影響我們今世的生活和心態；書中又提及,如何在一個情慾泛濫的社會中過聖潔生活.這卷書言及信仰落實到政治和婚姻生活秘訣,苦難的意義何在,我們怎樣才能過非表面化的肢體生活.這書又教人如何牧養神的教會,又如何在這滿佈陷阱的世代裡避免失足.我們今日來看這書,發覺它的主題何其適切教會的需要；神不單給予小亞細亞教會這一些信息,也為今日我們的教會預備了這些信息.
\newline
\newline
主耶穌升天以後,聖靈在五旬節降臨,教會遂於極端困難中建立起來,當時有成千上萬的人陸續歸主；可惜,耶穌基督的聖名對當代的世人而言,依然礙眼, 不受重視,甚至成了眼中釘.我相信神是透過逼迫使門徒離開耶路撒冷,把福音帶至撒瑪利亞,甚至直到地極.神藉著祂的僕人,就如保羅,巴拿巴,路加,馬可等等；他們在多處建立教會,主要對象是猶太人和外邦人.
\newline
\newline
彼得前書的讀者,大部份來自保羅從未踏足的小亞細亞地區.其中以猶太人居多,也有少數的外邦人,他們都是歸向了主的,只是因為政治和經濟的理由而居於猶太地以外,他們以小亞細亞為暫居之地,落葉歸根,總有一天,要回歸耶路撒冷的.然而,彼得透過這本書信,把他們的視線,從地上轉移到天上去.彼得前書還要處理一個無可避免的矛盾；就是說基督徒一方面既是屬天上的,是神國的公民,一方面卻仍得居住在地上.屬天和屬地象徵著兩個不同的取向,兩個不同的人生觀和價值觀.屬天的方面是聖潔的,榮耀的,永恆的；屬地的方面是污穢的,黑暗的,短暫的.一個代表福樂,另一個卻充滿苦難.基督徒往往好像剛剛離開各各他,卻仍未看見復活清晨的曙光.無可否認,我們今日都處於上述的這種張力之中；有些基督徒以為與世俗認同,就可免除這種張力；他們名義上稱自己為基督徒, 生活方面實際上與世人無異.對於此等基督徒而言,人生其實哪有張力?有些基督徒以為完全逃避社會的接觸,就可保持一己的聖潔.其實,最具破壞性的,不是外在的問題,乃是發自內心的邪情私慾；這正如耶穌基督在馬可福音第七章所論心裡的污穢.希望透過這幾堂的聚會,我們可以知道怎樣去化解上述的矛盾,從而明白我們身份的獨特性,再次肯定存活於世的使命.
\newline
\newline
彼得並非循一邏輯性的方式,去表達他的思想,故我們很難為彼得前書去寫大綱,但此書之主題甚為鮮明.故此,我們會從一段講解到另一段,作主題式的探討；盼望聖靈隨著這八堂的研讀,校正我們的心態和我們的人生觀；於此動盪不安的時代裡,為主打一場美好的大仗.
\newline
\newline
今天,我們先來看彼得前書第一主題:神在天上為我們預備永恆的基業.許多時候,人生就像迷魂陣一樣,充滿一大堆無答案的問題,不少人發出了問題卻得不著答案.這甚至連基督徒也不例外,他們提出問題,讓牧師和傳道人來解答,終究也是尋不著答案.世間上充滿著無法解決的難題,信徒彷彿在暴風雨當中搖搖欲墜；有被苦難的洪流淹沒,或被世俗潮流沖走的可能；或在試鍊的火焰中,信心銷毀.若要在這長期掙扎和搏鬥中堅持下去,實在要付很大的代價.究竟這樣做是否值得呢?為主的名去堅持信念,犧牲性命,或使生意受虧損,是否值得呢?我們彷彿就在漫長的黑夜裡尚未窺見曙光.落在這樣的光景中,人多半有許多的懷疑與埋怨.然而,在這人類歷史的水平線之上,曙光已經初露,人生是有盼望的,對基督徒來說,永恆的基業就在目前.只要我們認知人生有盼望,那能叫我們忍受人生幾許的艱難困苦；許多的壓力,我們也能承擔.在漆黑隧道裡踽踽前進的路人,只要知道再往前摸索一段路程,就能見到一個花香常漫的公園；他就願意奔跑這黑暗的路程.人在病痛當中,倘若有醫生應許一個月後必能痊癒,就算再痛苦,他也能承受.有次一位姊妹因為脊骨疼痛,以致坐臥不安,整天在病榻上痛苦呻吟；後經醫生診斷認為需要接受手術,手術完畢再探訪她的時候問她是否痛楚；她說從前的痛是腐化的痛,如今的痛是康復的痛,知道有一天這痛楚會完全的消除.
\newline
\newline
人生需要有盼望,一位生長於中國大陸的女士,千方百計的要到香港來,親友問她懂何手藝?她說到港以後就算照舊以製掃帚為業,也幹得有盼望；原因是她現行活著無甚意義.人生若無盼望,生活也就毫無意義,任何困難都變得難以忍受.更危險者,人在無盼望中會放縱自己；沉迷於罪中之樂.究竟基督徒與世人有何分別?剛才所讀的經文給我們一些指示:基督徒是一群對永恆有把握的人,從外表而言,他們的衣食住行,讀書就業,婚姻家庭等與世人似乎沒有差別；他們也經歷到悲歡離合,生老病死.但彼得提醒他的讀者,基督徒在世是群寄居的過路客.他們在表面上可能與世人沒有多大的分別,但卻在取向與心態上必然有別.基督徒是重生得救,神在萬人中揀選成為祂兒女的,他們藉聖靈得以歸屬於神並分別為聖.他們靠主耶穌的寶血洗淨,能夠坦然無懼與主相交,他們是重生的,一切舊的污穢都得以徹底潔淨,新生命展開了,這是永恆的,不朽壞的.所以說,在重生的人之中,永恆的種子經已萌芽.故彼得說:我們似乎雖與世人沒分別,這是就形式而言,但在生命上卻有很大的分野.我們好些時與青年人討論婚姻的問題,提到信與不信的差別,很多人以為是形式的分別；倘若不信的一方肯陪同已信的一方去崇拜,主日學,團契和查經班,那不就行了嗎?其實,這只是形式的問題,信與不信者主要是生命的分別.重生的人不再懼怕死亡的毒鉤,罪惡的綑縛,魔鬼的折磨和威嚇,屬神的人有活潑的盼望.他們不像局外人那樣守候生命的結束和審判的來臨；他們可以含笑以面對死亡.因為他們有著活生生,不被消滅的盼望；這盼望有根有基,不似夢想,有實際的經歷.
\newline
\newline
多少時候,人的盼望極其脆弱,前陣子我們看到黎巴嫩的人質,他們的家屬一旦尋得人質的蛛絲馬蹟,就已經充滿盼望；稍有蹟象顯示親屬尚未離世,就已極其興奮.人若能把握著某些東西,盼望就有根據.但我們確實發現世事少有把握,人以為抓著線索,得著財產以及名位,前途就有保障,其實,這些東西轉眼間也就煙消雲散了.彼得所指活潑的盼望,乃是不斷在心靈湧現,就算在苦難中也不會消滅,即使在逼迫中也不能消滅旳.彼得進一步指出:基督徒有一不能朽壞的基業, 此基業乃是不玷污,不衰殘的.猶太人的宿願乃是在巴勒斯坦建立家園,這是根據神對亞伯拉罕旳應許,到今日他們依舊有這心態和印象；所以不少猶太人,現今仍要從蘇聯返回巴勒斯坦.然而,巴勒斯坦仍稱不上是他們最終的家鄉；這地屢遭戰亂,血腥遍野,畢竟是個被假神玷污之地.耶路撒冷聖殿曾被異教玷污,這個猶太人的基業可說是毫無保障的.就算到了今天,以色列國仍是在備戰的狀態；這國家並非神所應許不朽壞的基業.早些時候,香港人經歷了一場銀行風波,那與美國近年銀行倒閉事件有相似之處；許多人把自己一生積蓄存放在銀行裡,希望年老時有錢可用.可是,當銀行一旦出事,他們的把握,他們的基業,都只能變作殘酷的回憶.
\newline
\newline
耶穌基督在受難之前,門徒都很徬徨失落,但耶穌卻安慰他們說:「我去原是為你們預備地方去(那是一處不能朽壞,不會衰殘和玷污的基業).」基督徒知道神為我們預備一永恆的基業,這一信念是否出於阿Q精神,自我催眠哩?彼得卻斷言不是,他確認基督徒有永恆的把握和重生的經歷,是建基於耶穌基督復活的客觀事實之上.這事實並非一主觀願望,因祂確實戰勝了死亡,擁有絕對的權柄,成為人類歷史上一個極大的轉捩點.正因為基督復活的緣故,我們得著重生與盼望.
\newline
\newline
所以,每當我們信心搖動的時候,應再次注目於基督復活的事實之上.初期教會宣講基督的時候,總離不開祂復活的事實.倘若我們把基督復活這事實挪開就比眾人更加可憐；因為不信主的人,可以享受今生,自由自在,我們卻似乎受著諸多的束縛.然而,基督確已復活,這是我們永恆盼望的基礎!請勿容讓魔鬼來欺騙你,把一些不純正的思想放在你裡頭.基督復活了,那朝向永恆的生命已經展開,以前的生命受罪的影響而有期限,復活的生命卻是永遠的生命.只要在基督裡,我們即已開始經歷這樣的生命.過去,神已經立下良好的基礎,今日,神照樣有特殊的作為,好叫我們知道有一天可以享受祂為我們預備的基業.
\newline
\newline
今天所讀的經文讓我們看到:神有祂的能力足以保守我們,必能得著所預備到末世要顯現的救恩.並不是憑個人堅毅的決心和神學的資格可以走上這路,乃是靠神能力的保守.倘若靠一己的能力,相信我們永遠沒有機會承受神完整的救恩.有一次,一位青年跟牧師去搭火車,由於全速前進的關係,車廂之內震盪不定.青年人開始跟牧師談話,他坦白對牧師說已信主有一段時間,但他沒法經得起生活上的考驗與壓力；在世俗與信仰的張力之間,他坦言再沒法支持下去,要考慮放棄他的信仰.牧師很有智慧的答道:你信否我可以叫這柄小刀直立於聖經上?(當時車廂仍然搖晃不定),青年人回答說:當然不,沒可能!牧師於是用手持著小刀立於聖經之上.青年人說:這個不算,因你用手持刀,這方才成.牧師說:「我從沒說過要用甚麼方法去使小刀立在聖經上,只是我用手把持著；就是在這震動不定的車廂內,小刀也能立在聖經上.」青年人於是恍然大悟:他自己徒然掙扎,無可能叫他站立得住,唯有靠著神的承託與扶持,乃能堅立.所以,很多時候基督徒的經歷就像保羅所寫:「似乎受責罰卻是不致喪命,似乎憂愁卻是常常喜樂,似乎貧窮卻是叫許多人富足的,似乎一無所有卻是樣樣都有.」的一般,生活中好像有許多矛盾一樣,有時靈性低沉至一光景甚至乎要去賣主；然而,神的手卻保守著我們.彼得讓我們一方面看見美麗的圖像,另一方面也看到今生的困惑,彼得說:「但如今,在百般的試煉中暫時憂愁.」(彼前一6)現今我們在百般試煉之中,信仰遭受衝擊；外在有逼迫,內在有情慾的試探.基督徒許多時候有這樣的感受；就是在未信主之前,生命中少有這些掙扎.如今反而多了掣肘,凡事都要思前想後,謀定而動.未信主以前與同事間的相處愉快無比,信主以後同事卻把我視作怪人看待. 未信主以前,父母對我甚為疼惜,信主以後,竟稱我為不孝子.各位弟兄姊妹,不知大家有沒有上述的感受.我們要知道:耶穌基督所帶來福音的信息必然與這時代的文化有抵觸,某一些信息的本身甚至可以成為這潮流中的絆腳石,它教我們以另一種價值觀去看萬事.主耶穌所傳講的福音與我們的罪性明顯地有衝突,我們信主以後,良知比從前敏感多了,加上聖靈的提醒,自然會感覺得有壓力；信仰在衝擊的狀況中,撒旦不單止與我們敵對,更是節節進攻,毫不放鬆.故此我們要學會如何防止魔鬼的攻勢.
\newline
\newline
彼得教導我們接納苦難的事實,基督徒多時落在熾熱的煉爐之中,切記小心勿在羅籐樹下自憐,或是把琴掛在柳樹上.彼得再進一步提示我們一個真理,就是當我們信仰受衝擊之時,真正屬神的人,因有信心的關係就仍然大有喜樂.事實上,信徒也會遭遇憂愁,但只是暫時性的,所以切勿在暫時的經歷中決定永恆的取 向；許多人好容易在暫時的壓迫和痛苦之中放棄了最寶貴的東西一那永恆的家鄉.相反地,彼得卻教導我們要從永恆的焦點看暫時的苦楚,正因如此,我們方才明白,保羅和西拉為何在監牢裡仍能歌唱,因為他們的焦點放在永恆之上.我們並不否認信徒在世有苦難,信仰也不免受到衝擊,無論在公司裡,家庭裡,甚至在街市上,隨時隨地都有可能遭受考驗.我們比較相信一些能聽,能見的東西,然而,信心卻是一種內在的認知和確信；所以不能夠憑我們的感官,觸覺去體驗神的工作. 而且,這信心有時可能攙雜了我們個人的成見,別人的理論和意見,這信心混雜了一些不潔淨的思想,故此這信心需要受考驗,使煉淨至純真的地步,好在主基督耶穌顯現的時候得到榮耀和稱讚.神要求我們對祂的信心不是幼稚的,充滿雜質的,祂藉著各種的經歷去煉淨我們,使我們對神的信成為一種毫無保留,完全體驗神救贖大恩的信心.
\newline
\newline
彼得前書第一章八至十二節指出:當我們信仰受衝擊信心遭考驗之際,反要更加起來堅守我們的信仰.「因信他就有說不出來,滿有榮光的大喜樂.」(彼前 一8)這經節裡的信字是現在式,即謂不單止過去相信救主基督生活的事實,更要有現在進行式的相信；恐怕很多基督徒以過去的信心去迎接目前的衝擊,實在很不足夠.因為過去的信只足夠讓人面對當時的考驗,我們須有現在的信來承擔今日的衝擊,正因如此,我們在苦難中仍然大有喜樂.我們所信不是些空泛的理念,乃是信神僕人和先知考察過的資料,這些舊約的記載,差不多透過一千五百年的時間,將神救恩的計劃逐步向人講解,先知一直在考察這些事實,在他們的察驗中發現基 督的受苦和復活並非偶然.主耶穌釘身十架並非一無奈的行動,亦非湊巧,乃有神的計劃在其中.神在創世之前就預備好,再經神僕小心的察驗；結果都一一應驗, 這就成為我們信心的根據.這救贖計劃非人所構思,乃是神將一些本來隱藏的奧秘向我們啟示出來.所以我們信神並非迷信,乃是信神的話語和祂的啟示,在此過程中,聖靈施展工作,打開我們的眼目,讓我們看見神真理的奧妙.
\newline
\newline
彼得提醒我們在困惑與試煉之中,不要胡思亂想,求神問卜,迷信命理,風水等；乃要重新察驗信心的基礎.今日基督徒開始在這方面有所動搖.每天只要打開報紙,就有人根據你的星相和生肖,教導你如何趨吉避凶,當基督徒對於應否閱讀這些專欄存有疑問的時候,就是信仰受考驗之時,應當立刻重拾信仰的確據.彼得與我們探討這些問題,意在教我們重整人生觀,「所以要約束你們的心,謹慎自守.......」(彼前一13)正因我們有天上永恆的基業,又有活潑的盼望和把握, 所以要約束我們的思想,認真對待自己的信仰和守著永恆的確據；心思意念都要來一次嚴謹的重整,從一個永恆的焦點來判斷萬事；包括我們的人生觀,價值觀和道德觀.我們的人生方向,價值判斷,善惡標準等方面的真理,容後幾堂詳細講論.然而彼得囑咐我們謹慎自守,是特別指著道德方面的儆醒自律而言.世人只有今生的盼望,香港市民更要面對九七的衝擊,故多半只管吃喝玩樂.所以基督徒要謹慎自守,約束自己.在苦難中,世人只管痛哭流淚,失去盼望.一般人更只顧現實, 追逐金錢名利,他們可說是為所欲為,想做就去做；卻不知道這一切連所隱藏的,神必審問.基督徒對這等事情應有何想法?彼得教我們要專心盼望耶穌基督顯現的時候所帶來給你們的恩,「專心」意謂全神貫注,不容有任何騷擾臨到,就如奧林匹克運動員等候發號施令前那種專注一樣.彼得用這樣的字眼來比喻我們盼望耶穌 基督顯現所帶來的恩.主的再來也就成了我們的集中點；不是說各人都要放棄自己的職業,乃是心態上不要以為愈忙就愈成功(今天許多的傳道人也沒法逃避這試探,連基督徒也沒有時間去關心其他的弟兄姊妹).究竟信徒眾多事務之中,有多少項與永恆扯得上關係?也許,在座弟兄姊妹當中不少人落在試煉之中,期望你們不要放棄,因為遙望曙光,黎明將至.希伯來書告訴我們:許多基督徒都是存著信心而死的,暫時還未得到所應許的,卻是從遠處看見且歡喜迎接；又承認在世不過是客旅,是寄居的,卻羨慕一個天上更美的家鄉.所以,弟兄姊妹,請把下垂的手,發酸的腿挺起來,繼續走這天路,有一日,必然經歷耶穌基督全然的救贖.
\newline
\newline


\newpage

\section{第64屆~港九培靈會~許道良~第二講}
\label{sec:1010}
\textbf{手潔心清}
\newline
\newline
連結: \href{https://www.hkbibleconference.org/session-message/view/1010}{\texttt{https://www.hkbibleconference.org/session-message/view/1010}}
\newline
\newline
\hyperref[sec:1009]{< < < PREV SERMON < < <}
~
\hyperref[sec:toc]{[返目錄]}
~
\hyperref[sec:1011]{> > > NEXT SERMON > > >}
\newline
\newline
彼得前書 1:16-1:16
\newline
\begin{longtable}{cl}
\hline
\hline
章節 & 經文 (和合本修訂版)\\
\hline
1:16 & \begin{tabularx}{0.7\textwidth}{X} 因為經上記著：「你們要成為聖，因為我是神聖的。」 \end{tabularx} \\ \\
[1ex]
\hline
\hline
\end{longtable}


\newpage

\section{第64屆~港九培靈會~許道良~第三講}
\label{sec:1011}
\textbf{擇善棄惡}
\newline
\newline
連結: \href{https://www.hkbibleconference.org/session-message/view/1011}{\texttt{https://www.hkbibleconference.org/session-message/view/1011}}
\newline
\newline
\hyperref[sec:1010]{< < < PREV SERMON < < <}
~
\hyperref[sec:toc]{[返目錄]}
~
\hyperref[sec:1012]{> > > NEXT SERMON > > >}
\newline
\newline


\newpage

\section{第64屆~港九培靈會~許道良~第四講}
\label{sec:1012}
\textbf{贏取汝心}
\newline
\newline
連結: \href{https://www.hkbibleconference.org/session-message/view/1012}{\texttt{https://www.hkbibleconference.org/session-message/view/1012}}
\newline
\newline
\hyperref[sec:1011]{< < < PREV SERMON < < <}
~
\hyperref[sec:toc]{[返目錄]}
~
\hyperref[sec:1013]{> > > NEXT SERMON > > >}
\newline
\newline


\newpage

\section{第64屆~港九培靈會~許道良~第五講}
\label{sec:1013}
\textbf{為義受苦}
\newline
\newline
連結: \href{https://www.hkbibleconference.org/session-message/view/1013}{\texttt{https://www.hkbibleconference.org/session-message/view/1013}}
\newline
\newline
\hyperref[sec:1012]{< < < PREV SERMON < < <}
~
\hyperref[sec:toc]{[返目錄]}
~
\hyperref[sec:1014]{> > > NEXT SERMON > > >}
\newline
\newline


\newpage

\section{第64屆~港九培靈會~許道良~第六講}
\label{sec:1014}
\textbf{責無旁貸}
\newline
\newline
連結: \href{https://www.hkbibleconference.org/session-message/view/1014}{\texttt{https://www.hkbibleconference.org/session-message/view/1014}}
\newline
\newline
\hyperref[sec:1013]{< < < PREV SERMON < < <}
~
\hyperref[sec:toc]{[返目錄]}
~
\hyperref[sec:1015]{> > > NEXT SERMON > > >}
\newline
\newline


\newpage

\section{第64屆~港九培靈會~許道良~第七講}
\label{sec:1015}
\textbf{善牧楷模}
\newline
\newline
連結: \href{https://www.hkbibleconference.org/session-message/view/1015}{\texttt{https://www.hkbibleconference.org/session-message/view/1015}}
\newline
\newline
\hyperref[sec:1014]{< < < PREV SERMON < < <}
~
\hyperref[sec:toc]{[返目錄]}
~
\hyperref[sec:1016]{> > > NEXT SERMON > > >}
\newline
\newline


\newpage

\section{第64屆~港九培靈會~許道良~第八講}
\label{sec:1016}
\textbf{步步為營}
\newline
\newline
連結: \href{https://www.hkbibleconference.org/session-message/view/1016}{\texttt{https://www.hkbibleconference.org/session-message/view/1016}}
\newline
\newline
\hyperref[sec:1015]{< < < PREV SERMON < < <}
~
\hyperref[sec:toc]{[返目錄]}
~
\hyperref[sec:983]{> > > NEXT SERMON > > >}
\newline
\newline


\newpage

\section{第65屆~港九培靈會~周永健~第一講}
\label{sec:983}
\textbf{神顧念我們──我下來是要救他們}
\newline
\newline
連結: \href{https://www.hkbibleconference.org/session-message/view/983}{\texttt{https://www.hkbibleconference.org/session-message/view/983}}
\newline
\newline
\hyperref[sec:1016]{< < < PREV SERMON < < <}
~
\hyperref[sec:toc]{[返目錄]}
~
\hyperref[sec:984]{> > > NEXT SERMON > > >}
\newline
\newline
出埃及記 1:1-2:25
\newline
\begin{longtable}{cl}
\hline
\hline
章節 & 經文 (和合本修訂版)\\
\hline
1:1 & \begin{tabularx}{0.7\textwidth}{X} 以色列的眾兒子各帶著家眷，和雅各一同來到埃及，他們的名字如下： \end{tabularx} \\ \\ \relax
1:2 & \begin{tabularx}{0.7\textwidth}{X} 呂便、西緬、利未、猶大、 \end{tabularx} \\ \\ \relax
1:3 & \begin{tabularx}{0.7\textwidth}{X} 以薩迦、西布倫、便雅憫、 \end{tabularx} \\ \\ \relax
1:4 & \begin{tabularx}{0.7\textwidth}{X} 但、拿弗他利、迦得、亞設。 \end{tabularx} \\ \\ \relax
1:5 & \begin{tabularx}{0.7\textwidth}{X} 凡從雅各生的，共有七十人。那時，約瑟已經在埃及。 \end{tabularx} \\ \\ \relax
1:6 & \begin{tabularx}{0.7\textwidth}{X} 約瑟和他所有的兄弟，以及那一代的人都死了。 \end{tabularx} \\ \\ \relax
1:7 & \begin{tabularx}{0.7\textwidth}{X} 然而，以色列人生養眾多，繁衍昌盛，極其強盛，遍滿了那地。 \end{tabularx} \\ \\ \relax
1:8 & \begin{tabularx}{0.7\textwidth}{X} 有一位不認識約瑟的新王興起，統治埃及。 \end{tabularx} \\ \\ \relax
1:9 & \begin{tabularx}{0.7\textwidth}{X} 他對自己的百姓說：「看哪，以色列人的百姓比我們還多，又比我們強盛。 \end{tabularx} \\ \\ \relax
1:10 & \begin{tabularx}{0.7\textwidth}{X} 來吧，讓我們機巧地待他們，恐怕他們增多起來，將來若有戰爭，他們就聯合我們的仇敵來攻擊我們，然後離開這地去了。」 \end{tabularx} \\ \\ \relax
1:11 & \begin{tabularx}{0.7\textwidth}{X} 於是埃及人派監工管轄他們，用勞役苦待他們。他們為法老建造儲貨城，就是比東和蘭塞。 \end{tabularx} \\ \\ \relax
1:12 & \begin{tabularx}{0.7\textwidth}{X} 可是越苦待他們，他們就越發增多，更加繁衍，埃及人就因以色列人愁煩。 \end{tabularx} \\ \\ \relax
1:13 & \begin{tabularx}{0.7\textwidth}{X} 埃及人嚴厲地強迫以色列人做工， \end{tabularx} \\ \\ \relax
1:14 & \begin{tabularx}{0.7\textwidth}{X} 使他們因苦工而生活痛苦；無論是和泥，是做磚，是做田間各樣的工，一切的工埃及人都嚴厲地對待他們。 \end{tabularx} \\ \\ \relax
1:15 & \begin{tabularx}{0.7\textwidth}{X} 埃及王又對希伯來的接生婆，一個名叫施弗拉，另一個名叫普阿的說： \end{tabularx} \\ \\ \relax
1:16 & \begin{tabularx}{0.7\textwidth}{X} 「你們為希伯來婦人接生，臨盆的時候要注意，若是男的，就把他殺了，若是女的，就讓她活。」 \end{tabularx} \\ \\ \relax
1:17 & \begin{tabularx}{0.7\textwidth}{X} 但是接生婆敬畏神，不照埃及王的吩咐去做，卻讓男孩活著。 \end{tabularx} \\ \\ \relax
1:18 & \begin{tabularx}{0.7\textwidth}{X} 埃及王召了接生婆來，對她們說：「你們為甚麼做這事，讓男孩活著呢？」 \end{tabularx} \\ \\ \relax
1:19 & \begin{tabularx}{0.7\textwidth}{X} 接生婆對法老說：「因為希伯來婦人與埃及婦人不同；希伯來婦人健壯，接生婆還沒有到，她們已經生產了。」 \end{tabularx} \\ \\ \relax
1:20 & \begin{tabularx}{0.7\textwidth}{X} 神恩待接生婆；以色列人增多起來，極其強盛。 \end{tabularx} \\ \\ \relax
1:21 & \begin{tabularx}{0.7\textwidth}{X} 接生婆因為敬畏神，神就叫她們成立家室。 \end{tabularx} \\ \\ \relax
1:22 & \begin{tabularx}{0.7\textwidth}{X} 法老吩咐他的眾百姓說：「把所生的每一個男孩都丟到尼羅河裡去，讓所有的女孩存活。」 \end{tabularx} \\ \\ \relax
2:1 & \begin{tabularx}{0.7\textwidth}{X} 有一個利未家的人娶了一個利未女子為妻。 \end{tabularx} \\ \\ \relax
2:2 & \begin{tabularx}{0.7\textwidth}{X} 那女人懷孕，生了一個兒子，見他俊美，就把他藏了三個月， \end{tabularx} \\ \\ \relax
2:3 & \begin{tabularx}{0.7\textwidth}{X} 後來不能再藏，就取了一個蒲草箱，抹上柏油和樹脂，將孩子放在裡面，把箱子擱在尼羅河邊的蘆葦中。 \end{tabularx} \\ \\ \relax
2:4 & \begin{tabularx}{0.7\textwidth}{X} 孩子的姊姊遠遠站著，要知道他究竟會怎樣。 \end{tabularx} \\ \\ \relax
2:5 & \begin{tabularx}{0.7\textwidth}{X} 法老的女兒來到尼羅河邊洗澡，她的女僕們在河邊行走。她看見在蘆葦中的箱子，就派一個使女把它拿來。 \end{tabularx} \\ \\ \relax
2:6 & \begin{tabularx}{0.7\textwidth}{X} 她打開箱子，看見那孩子。看哪，男孩在哭，她就可憐他，說：「這是希伯來人的一個孩子。」 \end{tabularx} \\ \\ \relax
2:7 & \begin{tabularx}{0.7\textwidth}{X} 孩子的姊姊對法老的女兒說：「我去叫一個希伯來婦人來作奶媽，替你乳養這孩子，好嗎？」 \end{tabularx} \\ \\ \relax
2:8 & \begin{tabularx}{0.7\textwidth}{X} 法老的女兒對她說：「去吧！」那女孩就去叫了孩子的母親來。 \end{tabularx} \\ \\ \relax
2:9 & \begin{tabularx}{0.7\textwidth}{X} 法老的女兒對她說：「你把這孩子抱去，替我乳養這孩子，我必給你工錢。」那婦人就把孩子接過來，乳養他。 \end{tabularx} \\ \\ \relax
2:10 & \begin{tabularx}{0.7\textwidth}{X} 孩子長大了，婦人把他帶到法老的女兒那裡，就作了她的兒子。她給孩子起名叫摩西，說：「因我把他從水裡拉出來。」 \end{tabularx} \\ \\ \relax
2:11 & \begin{tabularx}{0.7\textwidth}{X} 過了一段日子，摩西長大了，他出去到他同胞那裡，看見他們的勞役。他看見一個埃及人打他的同胞，一個希伯來人。 \end{tabularx} \\ \\ \relax
2:12 & \begin{tabularx}{0.7\textwidth}{X} 他左右觀看，見沒有人，就把埃及人打死了，藏在沙土裡。 \end{tabularx} \\ \\ \relax
2:13 & \begin{tabularx}{0.7\textwidth}{X} 第二天他出去，看哪，有兩個希伯來人在打架，他就對那兇惡的人說：「你為甚麼打你同族的人呢？」 \end{tabularx} \\ \\ \relax
2:14 & \begin{tabularx}{0.7\textwidth}{X} 那人說：「誰立你作我們的領袖和審判官呢？難道你要殺我，像殺那埃及人一樣嗎？」摩西就懼怕，說：「這事一定是讓人知道了。」 \end{tabularx} \\ \\ \relax
2:15 & \begin{tabularx}{0.7\textwidth}{X} 法老聽見這事，就設法要殺摩西。於是摩西逃走，躲避法老，到了米甸地，坐在井旁。 \end{tabularx} \\ \\ \relax
2:16 & \begin{tabularx}{0.7\textwidth}{X} 米甸的祭司有七個女兒；她們來打水，打滿了槽，要給父親的羊群喝水。 \end{tabularx} \\ \\ \relax
2:17 & \begin{tabularx}{0.7\textwidth}{X} 有一些牧羊人來，把她們趕走，摩西卻起來幫助她們，取水給她們的羊群喝。 \end{tabularx} \\ \\ \relax
2:18 & \begin{tabularx}{0.7\textwidth}{X} 她們回到父親流珥那裡；他說：「今日你們為何這麼快就回來了呢？」 \end{tabularx} \\ \\ \relax
2:19 & \begin{tabularx}{0.7\textwidth}{X} 她們說：「有一個埃及人來救我們脫離牧羊人的手，他甚至打水給我們的羊群喝。」 \end{tabularx} \\ \\ \relax
2:20 & \begin{tabularx}{0.7\textwidth}{X} 他對女兒們說：「那人在哪裡？你們為甚麼撇下他呢？去請他來吃飯吧！」 \end{tabularx} \\ \\ \relax
2:21 & \begin{tabularx}{0.7\textwidth}{X} 摩西願意和那人同住， 那人就把女兒西坡拉給摩西為妻。 \end{tabularx} \\ \\ \relax
2:22 & \begin{tabularx}{0.7\textwidth}{X} 西坡拉生了一個兒子，摩西給他起名叫革舜，因他說：「我在外地作了寄居者。」 \end{tabularx} \\ \\ \relax
2:23 & \begin{tabularx}{0.7\textwidth}{X} 過了許多年，埃及王死了。以色列人因做苦工，就嘆息哀求；他們因苦工所發出的哀聲達於神。 \end{tabularx} \\ \\ \relax
2:24 & \begin{tabularx}{0.7\textwidth}{X} 神聽見他們的哀聲，就記念他與亞伯拉罕、以撒、雅各所立的約。 \end{tabularx} \\ \\ \relax
2:25 & \begin{tabularx}{0.7\textwidth}{X} 神看顧以色列人，神是知道的。 \end{tabularx} \\ \\
[1ex]
\hline
\hline
\end{longtable}
出埃及記是一本很奇妙的書,它不單講到以色列人為何在埃及作奴隸,更論到神如何領他們離開埃及地.出埃及記不但講歷史的過去,同時裡面也有很多豐富 屬靈的信息；講到神有拯救的能力與恩典,以色列人作神子民尊貴的身份與地位.這卷書對今日香港的信徒與教會,有高度的適切性；希望藉這八天的研經,幫助大家更能深入了解這書的信息,認識及明白神的心意.引導我們在這時代活出神子民的見證.
\newline
\newline
聖經首五卷書,「創,出,利,民,申」稱為摩西五經.這五卷書其實並非五卷不同的書,而是一卷書分成五章,五個段落；每一章都有它的中心信息及主題.創世記是講神的創造與萬物的起源；出埃及記是講神的拯救；利未記是講成聖；民數記是講神的引導與管教；申命記是講神的信實與供應.這五卷書稱為摩西五經.今次的研經會,我們特別去看出埃及記,是講神的拯教,和祂的恩典與能力.以色列人作為神子民之後的身份,地位及使命,希望對我們香港的教會,信徒有同樣重要的屬靈信息.我們先看出埃及記一1-10,二23-25,在這兩章聖經中,我們會用三個題目去查考:
\newline
\newline
(一)以色列人在埃及做了奴隸
\newline
\newline
為何以色列人會去到埃及呢?如果我們要明白背景,必須先讀創世記,創世記與出埃及記是分不開的,並不是兩卷書,是上,下卷.創世記告訴我們,雅各與他的家人下到埃及,是因他們所居住的迦南發生飢荒,就移民去埃及.因雅各的兒子約瑟被兄長賣到埃及地作奴隸,後來神賜福他,使他被提升為埃及的宰相,治理埃及全地.最後因迦南地飢荒,就接了他父親雅各和他的兄弟來到埃及地居住.由於約瑟是埃及的宰相,所以他的家人備受優待,住在歌珊地,這地非常肥沃.但當我們讀到出埃及記時,整個情景都很不同了.出埃及記一1-7.「約瑟和他的弟兄,並那一代的人,都死了.」(6)過了一段很長的時間,四百三十年(出十二 40)約瑟那一代的人都死了,在這麼長的時間,會有甚麼的變化?
\newline
\newline
1.七節:「以色列人生養眾多,並且繁茂,極其強盛,滿了那地.」以色列人從雅各一家的七十五人,經過了四百多年,發展到人丁眾多,繁茂強盛,滿了埃及地.應驗了神曾對亞伯拉罕的應許:他的子孫君王從你而立,國度從你而出.」以色列人在埃及地,經過了四百多年,成了一個強大的民族.
\newline
\newline
2.八節:告訴我們有更大而重要的改變,就是有一個不認識約瑟的新王起來治理埃及.所謂一朝天子一朝臣,不承認約瑟曾經給埃及的貢獻,就改變了對以色列人的態度.以色列人不是本地人,是寄居者；所以他們受到歧視,對付與逼害.為甚麼呢?這是很特別有趣的,請大家留心看(出一9-10).這裡論兩樣事:神賜福以色列人,人多強盛,這是一個事實.但埃及人與法老擔心打仗時,以色列人不會幫他們,倒去幫助敵人,那時他們就會受虧損.但這不是事實,不過是將事實進一步加上自己的推測,看為事實而已.以色列人在埃及地是一等良民,對國家很有貢獻,而且安份守己,但埃及王沒有這樣想,反而恐懼擔心,缺乏安全感.以色列人強盛,反構成了埃及的威脅；神賜福給以色列人,竟成了法老王逼迫他們的一個原因.有時神給我們的恩典,福份,會引致別人的妒忌,甚至逼害.人 就是如此分不清事實與幻想；當受到威脅時,就會做一些事情去對付.法老將這班本來寄居的以色列人逼害,作奴隸管轄他們.這是約瑟意想不到的事,自己做埃及的宰相,一心接待父親兄弟,逃避迦南的飢荒,以為可以從此安居樂業；想不到經過多代之後,他的後裔竟變成奴隸.人生的變幻,是我們不能掌握的.所以當我們遇上變幻時,不要灰心,失望,只要相信神在其中,一定有祂的美意.以色列人受到埃及王法老的逼害,這逼害有四個方式,循序漸進地變本加厲.
\newline
\newline
(1)十一節:埃及人派督工監視,控制,管轄他們,用盡方法去打擊苦害他們.
\newline
\newline
(2)十三,十四節:嚴嚴使以色列人作工,成為奴隸.為法老建造兩座積貨城,就是比東和蘭塞,將加倍的重擔苦害他們.
\newline
\newline
(3)十二節:用毒棘的手段,命希伯來的收生婆,殺死以色列人所出生的男嬰.
\newline
\newline
(4)二十二節:不用收生婆,授權給埃及人去執法,凡以色列人初生的男嬰,要丟在河裡淹死.這簡直不是一個法治的社會,而是人治的社會了.
\newline
\newline
神的子民以色列人受盡了痛苦與逼害.弟兄姊妹!我們所接觸的事物,和所處的環境,原本有好條件,但可能都會不斷轉變.雅各約瑟怎樣也想不到,他們的子孫會淪為奴隸；若早知如此,他們就不往埃及去或早點離開埃及了.同樣今日在我們的生活,家庭,教會,社會中,會有許多新的事物,會帶來始料不及的轉變； 以致給我們遇上許多的苦惱,痛苦,逼迫,歧視,這時候我們怎樣去應付呢?
\newline
\newline
(二)以色列人向神哀求,發出呼聲
\newline
\newline
「使他們因作苦工覺得命苦.」(出一14).當人遇上不幸時都自認命苦,與別人比較是很自然的.幸好以色列人不停留在這地步.雖然作苦工覺得命苦, 他們會向神哀求(出二23).如果以色列人作苦工覺得命苦,只曉得自怨自艾,這是一條死路；但以色列人嘆息哀求神,這就變成一條出路.我們遭遇環境不幸時,雖然感到痛苦,但能發出嘆息哀求神,這才是出路.在這世上,遇上人生種種遭遇,我們會有兩種反應,而聖經給我們的信息,這兩種的反應是同樣的真實而存在的:
\newline
\newline
1.讚美神,歌頌神.為何我們要歌頌讚美神呢?因為神給我們各樣的恩典,醫治,幫助,祂的愛拯救我,我感謝祂!我要讚美,歌頌神；詩篇充滿著這樣的讚美,所以我們的歌頌和讚美,是從喜樂表現出來的聲音.
\newline
\newline
2.哀嘆哀歌.人生有痛苦有悲傷,失敗,死亡,面對著這些遭遇,人從心底裡向神發出嘆息哀求,這是痛苦裡的聲音.基督徒可否發出哀歌呢?聖經中有許多的哀歌,耶利米哀歌是其中一首,詩篇中也有不少的哀歌.詩人的遭遇向神發出哀求:「神啊!你不看顧我要到幾時呢?」這些哀歌是指向神,是信心的表現和記號,引導我們親近神.很多人以為基督徒一定要讚美歌頌神,不可有哀歌,如果發出哀聲嘆息,是表示沒有信心,不信靠神.其實基督徒讚美歌頌,也不能避開哀歌；作基督徒不一定是成功,快樂的.即使最快樂的人,有時都會遇上痛苦,憂傷；請不要誤會,聖經告訴我們要常常喜樂,不住的禱告,凡事謝恩,這是不錯的； 聖經要我們從心發出向神的倚靠,這種倚靠是用讚美歌頌；但有時是透過哀歌.感謝神給我們有這樣的自由和權利,在我們痛苦,灰心,失望當中而感到命苦；但反而因這些痛苦遭遇可到神面前嘆息,哀聲達於神,這是神給我們一個去親近祂倚靠祂的途徑.一個健康基督徒的生命,應是能夠讚美稱頌,又能哀嘆發出哀聲的.
\newline
\newline
(三)神聽我們的哀求
\newline
\newline
最奇妙的是我們不單能向神發出祈求哀嘆,更重要的是這些祈求是不會徒然的,因為神顧念屬祂的人,神聽我們的祈求,特別是我們的哀聲.神如何回應以色 列人的嘆息呢?「神聽見他們的哀聲,就記念祂與亞伯拉罕,以撒,雅各所立的約.神看顧以色列人,也知道他們的苦情.」(出二24-25)聖經中用了四個動詞,去形容神的反應:
\newline
\newline
(1)神聽見,
\newline
\newline
(2)神記念,
\newline
\newline
(3)神看顧,
\newline
\newline
(4)神知道.
\newline
\newline
感謝神,不論我們落在甚麼光景中,神聽見,記得,看見又知道,這是很大的安慰；所以到最後,神呼召摩西的時候,「耶和華說:我的百姓在埃及所受的困苦,我實在看見了；他們因受督工的轄制所發的哀聲,我也聽見了.我原知道他們的痛苦.」(出三7),神有行動,所以第三章7節是回應第二章24-25節, 神見以色列人所受的困苦,「我下來是要救他們脫離埃及人的手.」(出三8)耶穌基督下來要救我們,他死在十字架上流血捨命,救我們脫離罪惡死亡；神如何顧念以色列人,祂今日同樣顧念一切屬祂的人.
\newline
\newline
環境的改變帶給我們很多影響,你現今處身在一個甚麼的情況下呢?是否像昔日以色列人一樣?以色列人的經驗也提醒了我們,他們不單覺自己命苦,更加因所遭遇,逼使他們嘆息哀求親近神.當我們遇到不順利不如意的事,有沒有驅使我們更加去親近,祈求和倚靠神呢?我們不妨向神發出哀聲祈求,因為祂是會聽的. 最後我們要肯定神是顧念我們的,我們要將一切的重擔卸給神,因祂顧念我們.這不過是開始,跟著我們會看見神的行動:「我下來是要救他們」,祂就開始了祂偉大的拯救工作.第三章我們會讀到神呼召,興起一位偉大的領袖摩西,帶領以色列人出埃及；神今日同樣呼召我們去事奉祂,完成祂所交托給我們的使命.
\newline
\newline


\newpage

\section{第65屆~港九培靈會~周永健~第二講}
\label{sec:984}
\textbf{神呼召我們事奉祂──我要打發你們去見法老}
\newline
\newline
連結: \href{https://www.hkbibleconference.org/session-message/view/984}{\texttt{https://www.hkbibleconference.org/session-message/view/984}}
\newline
\newline
\hyperref[sec:983]{< < < PREV SERMON < < <}
~
\hyperref[sec:toc]{[返目錄]}
~
\hyperref[sec:985]{> > > NEXT SERMON > > >}
\newline
\newline
出埃及記 3:1-6:30
\newline
\begin{longtable}{cl}
\hline
\hline
章節 & 經文 (和合本修訂版)\\
\hline
3:1 & \begin{tabularx}{0.7\textwidth}{X} 摩西牧放他岳父米甸祭司葉特羅的羊群，他領羊群往曠野的那一邊去，到了神的山，就是何烈山。 \end{tabularx} \\ \\ \relax
3:2 & \begin{tabularx}{0.7\textwidth}{X} 耶和華的使者在荊棘的火焰中向他顯現。摩西觀看，看哪，荊棘在火中焚燒，卻沒有燒燬。 \end{tabularx} \\ \\ \relax
3:3 & \begin{tabularx}{0.7\textwidth}{X} 摩西說：「我要轉過去看這大異象，這荊棘為何沒有燒燬呢？」 \end{tabularx} \\ \\ \relax
3:4 & \begin{tabularx}{0.7\textwidth}{X} 耶和華見摩西轉過去看，神就從荊棘裡呼叫他說：「摩西！摩西！」他說：「我在這裡。」 \end{tabularx} \\ \\ \relax
3:5 & \begin{tabularx}{0.7\textwidth}{X} 神說：「不要靠近這裡。把你腳上的鞋脫下來，因為你所站的地方是聖地」。 \end{tabularx} \\ \\ \relax
3:6 & \begin{tabularx}{0.7\textwidth}{X} 他又說：「我是你父親的神，是亞伯拉罕的神，以撒的神，雅各的神。」摩西蒙上臉，因為怕看神。 \end{tabularx} \\ \\ \relax
3:7 & \begin{tabularx}{0.7\textwidth}{X} 耶和華說：「我確實看見了我百姓在埃及所受的困苦，我也聽見了他們因受監工苦待所發的哀聲；我確實知道他們的痛苦。 \end{tabularx} \\ \\ \relax
3:8 & \begin{tabularx}{0.7\textwidth}{X} 我下來是要救他們脫離埃及人的手，領他們從那地上來，到美好與寬闊之地，到流奶與蜜之地，就是迦南人、赫人、亞摩利人、比利洗人、希未人、耶布斯人之地。 \end{tabularx} \\ \\ \relax
3:9 & \begin{tabularx}{0.7\textwidth}{X} 現在，看哪，以色列人的哀聲達到我這裡，我也看見埃及人怎樣欺壓他們。 \end{tabularx} \\ \\ \relax
3:10 & \begin{tabularx}{0.7\textwidth}{X} 現在，你去，我要差派你到法老那裡，把我的百姓以色列人從埃及領出來。」 \end{tabularx} \\ \\ \relax
3:11 & \begin{tabularx}{0.7\textwidth}{X} 摩西對神說：「我是甚麼人，竟能去見法老，把以色列人從埃及領出來呢？」 \end{tabularx} \\ \\ \relax
3:12 & \begin{tabularx}{0.7\textwidth}{X} 神說：「我必與你同在。這就是我差派你去，給你的憑據：你把百姓從埃及領出來之後，你們必在這山上事奉神。」 \end{tabularx} \\ \\ \relax
3:13 & \begin{tabularx}{0.7\textwidth}{X} 摩西對神說：「看哪，我到以色列人那裡，對他們說：『你們祖宗的神差派我到你們這裡來。』他們若對我說：『他叫甚麼名字？』我要對他們說甚麼呢？」 \end{tabularx} \\ \\ \relax
3:14 & \begin{tabularx}{0.7\textwidth}{X} 神對摩西說：「我是自有永有的」；又說：「你要對以色列人這樣說：『那自有永有的差派我到你們這裡來。』」 \end{tabularx} \\ \\ \relax
3:15 & \begin{tabularx}{0.7\textwidth}{X} 神又對摩西說：「你要對以色列人這樣說：『耶和華－你們祖宗的神，就是亞伯拉罕的神，以撒的神，雅各的神差派我到你們這裡來。』這是我的名，直到永遠；這也是我的稱號，直到萬代。 \end{tabularx} \\ \\ \relax
3:16 & \begin{tabularx}{0.7\textwidth}{X} 你去召集以色列的長老，對他們說：『耶和華－你們祖宗的神，就是亞伯拉罕的神，以撒的神，雅各的神向我顯現，說：我實在眷顧了你們，眷顧你們在埃及的遭遇。 \end{tabularx} \\ \\ \relax
3:17 & \begin{tabularx}{0.7\textwidth}{X} 我也曾說：要把你們從埃及的困苦中領出來，往迦南人、赫人、亞摩利人、比利洗人、希未人、耶布斯人的地去，就是到流奶與蜜之地。』 \end{tabularx} \\ \\ \relax
3:18 & \begin{tabularx}{0.7\textwidth}{X} 他們必聽你的話。你和以色列的長老要到埃及王那裡，對他說：『耶和華－希伯來人的神向我們顯現，現在求你讓我們往曠野去，走三天的路程，為要向耶和華我們的神獻祭。』 \end{tabularx} \\ \\ \relax
3:19 & \begin{tabularx}{0.7\textwidth}{X} 我知道若不用大能的手，埃及王不會放你們走。 \end{tabularx} \\ \\ \relax
3:20 & \begin{tabularx}{0.7\textwidth}{X} 因此，我必伸出我的手，在埃及施行我一切的神蹟，擊打這地，然後，他才放你們走。 \end{tabularx} \\ \\ \relax
3:21 & \begin{tabularx}{0.7\textwidth}{X} 我必使埃及人看得起你們，你們離開的時候就不至於空手而去。 \end{tabularx} \\ \\ \relax
3:22 & \begin{tabularx}{0.7\textwidth}{X} 每一個婦女必向她的鄰舍，以及寄居在她家裡的女人，索取金器、銀器和衣裳，給你們的兒女穿戴。這樣你們就掠奪了埃及人。」 \end{tabularx} \\ \\ \relax
4:1 & \begin{tabularx}{0.7\textwidth}{X} 摩西回答說：「看哪！他們不會信我，也不會聽我的話，因為他們必說：『耶和華並沒有向你顯現。』」 \end{tabularx} \\ \\ \relax
4:2 & \begin{tabularx}{0.7\textwidth}{X} 耶和華對摩西說：「你手裡的是甚麼？」他說：「是杖。」 \end{tabularx} \\ \\ \relax
4:3 & \begin{tabularx}{0.7\textwidth}{X} 耶和華說：「把它丟在地上！」他一丟在地上，杖就變成一條蛇；摩西逃走避開牠。 \end{tabularx} \\ \\ \relax
4:4 & \begin{tabularx}{0.7\textwidth}{X} 耶和華對摩西說：「伸出手來，拿住牠的尾巴---摩西就伸出手，抓住牠，牠就在摩西的手掌中變為杖--- \end{tabularx} \\ \\ \relax
4:5 & \begin{tabularx}{0.7\textwidth}{X} 為了要使他們信耶和華他們祖宗的神，就是亞伯拉罕的神，以撒的神，雅各的神，曾向你顯現了。」 \end{tabularx} \\ \\ \relax
4:6 & \begin{tabularx}{0.7\textwidth}{X} 耶和華又對他說：「把手放進懷裡。」他就把手放進懷裡。當他把手抽出來，看哪，手竟然長了痲瘋，像雪一樣白。 \end{tabularx} \\ \\ \relax
4:7 & \begin{tabularx}{0.7\textwidth}{X} 耶和華說：「把手放回懷裡---他就把手放回懷裡。當他把手從懷裡再抽出來，看哪，手復原了，與全身的肉一樣--- \end{tabularx} \\ \\ \relax
4:8 & \begin{tabularx}{0.7\textwidth}{X} 倘若他們不信你，也不聽第一個神蹟的聲音，他們會信第二個神蹟的聲音。 \end{tabularx} \\ \\ \relax
4:9 & \begin{tabularx}{0.7\textwidth}{X} 倘若他們不信這兩個神蹟，不聽你的話，你就從尼羅河裡取些水，倒在乾的地上。你從尼羅河裡所取的水必在乾地上變成血。」 \end{tabularx} \\ \\ \relax
4:10 & \begin{tabularx}{0.7\textwidth}{X} 摩西對耶和華說：「主啊，求求你，我並不是一個能言善道的人，以前這樣，就是你對僕人說話以後也是這樣，因為我是拙口笨舌的。」 \end{tabularx} \\ \\ \relax
4:11 & \begin{tabularx}{0.7\textwidth}{X} 耶和華對他說：「誰造人的口呢？誰使人口啞、耳聾、目明、眼瞎呢？豈不是我－耶和華嗎？ \end{tabularx} \\ \\ \relax
4:12 & \begin{tabularx}{0.7\textwidth}{X} 現在，去吧，我必賜你口才，指教你應當說的。」 \end{tabularx} \\ \\ \relax
4:13 & \begin{tabularx}{0.7\textwidth}{X} 摩西說：「主啊，求求你，你要藉著誰的手，就差派誰去吧！」 \end{tabularx} \\ \\ \relax
4:14 & \begin{tabularx}{0.7\textwidth}{X} 耶和華的怒氣向摩西發作，說：「你不是有一個哥哥利未人亞倫嗎？我知道他是個能言善道的人。看哪，他正出來迎接你。他一見到你，心裡就歡喜。 \end{tabularx} \\ \\ \relax
4:15 & \begin{tabularx}{0.7\textwidth}{X} 你要跟他說話，把話放在他的口裡，我要賜你口才，也要賜他口才，又要教你們做當做的事。 \end{tabularx} \\ \\ \relax
4:16 & \begin{tabularx}{0.7\textwidth}{X} 他要替你向百姓說話；他要當你的口，你要當他的神。 \end{tabularx} \\ \\ \relax
4:17 & \begin{tabularx}{0.7\textwidth}{X} 你手裡要拿這杖，用它來行神蹟。」 \end{tabularx} \\ \\ \relax
4:18 & \begin{tabularx}{0.7\textwidth}{X} 於是，摩西回到他岳父葉特羅那裡，對他說：「請你讓我回埃及我同胞那裡，看他們還在不在。」葉特羅對摩西說：「平平安安地去吧！」 \end{tabularx} \\ \\ \relax
4:19 & \begin{tabularx}{0.7\textwidth}{X} 耶和華在米甸對摩西說：「你要回埃及去，因為那些尋索你命的人都死了。」 \end{tabularx} \\ \\ \relax
4:20 & \begin{tabularx}{0.7\textwidth}{X} 摩西就帶著妻子和兩個兒子，讓他們騎上驢，回埃及地去。摩西手裡拿著神的杖。 \end{tabularx} \\ \\ \relax
4:21 & \begin{tabularx}{0.7\textwidth}{X} 耶和華對摩西說：「你回到埃及去的時候，要留意將我交在你手中的一切奇事行在法老面前。但我要任憑他的心剛硬，他必不放百姓走。 \end{tabularx} \\ \\ \relax
4:22 & \begin{tabularx}{0.7\textwidth}{X} 你要對法老說：『耶和華如此說：以色列是我的兒子，我的長子。 \end{tabularx} \\ \\ \relax
4:23 & \begin{tabularx}{0.7\textwidth}{X} 我對你說過：放我的兒子走，好事奉我。你還是不肯放他走。看哪，我要殺你頭生的兒子。』」 \end{tabularx} \\ \\ \relax
4:24 & \begin{tabularx}{0.7\textwidth}{X} 在路上住宿的地方，耶和華遇見摩西，想要殺他。 \end{tabularx} \\ \\ \relax
4:25 & \begin{tabularx}{0.7\textwidth}{X} 西坡拉就拿一塊火石，割下她兒子的包皮，碰觸摩西的腳，說：「你真是我血的新郎了。」 \end{tabularx} \\ \\ \relax
4:26 & \begin{tabularx}{0.7\textwidth}{X} 這樣，耶和華才放了他。那時，西坡拉說：「你因割禮就是血的新郎了」。 \end{tabularx} \\ \\ \relax
4:27 & \begin{tabularx}{0.7\textwidth}{X} 耶和華對亞倫說：「你往曠野去迎接摩西。」他就去，在神的山遇見摩西，就親他。 \end{tabularx} \\ \\ \relax
4:28 & \begin{tabularx}{0.7\textwidth}{X} 摩西將耶和華差派他所說的話和吩咐他所行的神蹟都告訴了亞倫。 \end{tabularx} \\ \\ \relax
4:29 & \begin{tabularx}{0.7\textwidth}{X} 摩西和亞倫就去召集以色列的眾長老。 \end{tabularx} \\ \\ \relax
4:30 & \begin{tabularx}{0.7\textwidth}{X} 亞倫將耶和華對摩西所說的一切話述說了一遍，又在百姓眼前行了那些神蹟， \end{tabularx} \\ \\ \relax
4:31 & \begin{tabularx}{0.7\textwidth}{X} 百姓就信了。他們聽見耶和華眷顧以色列人，鑒察他們的困苦，就低頭敬拜。 \end{tabularx} \\ \\ \relax
5:1 & \begin{tabularx}{0.7\textwidth}{X} 後來，摩西和亞倫去對法老說：「耶和華－以色列的神這樣說：『放我的百姓走，好讓他們在曠野向我守節。』」 \end{tabularx} \\ \\ \relax
5:2 & \begin{tabularx}{0.7\textwidth}{X} 法老說：「耶和華是誰，要我聽他的話，讓以色列人去？我不認識耶和華，也不放以色列人走！」 \end{tabularx} \\ \\ \relax
5:3 & \begin{tabularx}{0.7\textwidth}{X} 他們說：「希伯來人的神已向我們顯現了。求你讓我們往曠野去，走三天的路程，向耶和華我們的神獻祭，免得他用瘟疫、刀劍攻擊我們。」 \end{tabularx} \\ \\ \relax
5:4 & \begin{tabularx}{0.7\textwidth}{X} 埃及王對他們說：「摩西、亞倫！你們為甚麼叫百姓不做工呢？去，服你們的勞役吧！」 \end{tabularx} \\ \\ \relax
5:5 & \begin{tabularx}{0.7\textwidth}{X} 他又說：「看哪，這地的以色列人如今這麼多，你們竟然叫他們歇下勞役！」 \end{tabularx} \\ \\ \relax
5:6 & \begin{tabularx}{0.7\textwidth}{X} 當天，法老吩咐監工和工頭說： \end{tabularx} \\ \\ \relax
5:7 & \begin{tabularx}{0.7\textwidth}{X} 「你們不可照以前一樣提供草給百姓做磚，要叫他們自己去撿草。 \end{tabularx} \\ \\ \relax
5:8 & \begin{tabularx}{0.7\textwidth}{X} 他們平時做磚的數目，你們仍舊向他們要，一點不可減少，因為他們是懶惰的，所以才呼求說：『讓我們去向我們的神獻祭。』 \end{tabularx} \\ \\ \relax
5:9 & \begin{tabularx}{0.7\textwidth}{X} 你們要把更重的工作加在這些人身上，使他們在其中勞碌，不去理會謊言。」 \end{tabularx} \\ \\ \relax
5:10 & \begin{tabularx}{0.7\textwidth}{X} 監工和工頭出來對百姓說：「法老這樣說：『我不給你們草， \end{tabularx} \\ \\ \relax
5:11 & \begin{tabularx}{0.7\textwidth}{X} 你們自己在哪裡能找到草，就往哪裡去找吧！但你們的工作一點也不可減少。』」 \end{tabularx} \\ \\ \relax
5:12 & \begin{tabularx}{0.7\textwidth}{X} 於是，百姓分散在埃及全地，撿碎秸當草用。 \end{tabularx} \\ \\ \relax
5:13 & \begin{tabularx}{0.7\textwidth}{X} 監工催逼他們，說：「你們每天要做完一天的工，與先前有草一樣。」 \end{tabularx} \\ \\ \relax
5:14 & \begin{tabularx}{0.7\textwidth}{X} 法老的監工擊打他們所派的以色列工頭，說：「為甚麼昨天和今天你們沒有按照以前做磚的數目，完成你們的工作呢？」 \end{tabularx} \\ \\ \relax
5:15 & \begin{tabularx}{0.7\textwidth}{X} 以色列人的工頭來哀求法老說：「為甚麼這樣待你的僕人呢？ \end{tabularx} \\ \\ \relax
5:16 & \begin{tabularx}{0.7\textwidth}{X} 監工不把草給僕人，並且對我們說：『做磚吧！』看哪，你僕人挨了打，其實是你百姓的錯。」 \end{tabularx} \\ \\ \relax
5:17 & \begin{tabularx}{0.7\textwidth}{X} 法老卻說：「懶惰，你們真是懶惰！所以你們說：『讓我們去向耶和華獻祭吧。』 \end{tabularx} \\ \\ \relax
5:18 & \begin{tabularx}{0.7\textwidth}{X} 現在，去做工吧！草是不會給你們，磚卻要如數交納。」 \end{tabularx} \\ \\ \relax
5:19 & \begin{tabularx}{0.7\textwidth}{X} 以色列人的工頭聽見「你們每天做磚的工作一點也不可減少」，就知道惹上禍了。 \end{tabularx} \\ \\ \relax
5:20 & \begin{tabularx}{0.7\textwidth}{X} 他們離開法老出來，正遇見摩西和亞倫站在那裡等候他們， \end{tabularx} \\ \\ \relax
5:21 & \begin{tabularx}{0.7\textwidth}{X} 就向他們說：「願耶和華鑒察你們，施行判斷，因為你們使我們在法老和他臣僕面前有了臭名，把刀遞在他們手中來殺我們。」 \end{tabularx} \\ \\ \relax
5:22 & \begin{tabularx}{0.7\textwidth}{X} 摩西回到耶和華那裡，說：「主啊，你為甚麼苦待這百姓呢？為甚麼差派我呢？ \end{tabularx} \\ \\ \relax
5:23 & \begin{tabularx}{0.7\textwidth}{X} 自從我到法老那裡，奉你的名說話，他就苦待這百姓，你卻一點也沒有拯救你的百姓。」 \end{tabularx} \\ \\ \relax
6:1 & \begin{tabularx}{0.7\textwidth}{X} 耶和華對摩西說：「現在你必看見我向法老所行的事，使他因我大能的手放以色列人走，因我大能的手把他們趕出他的地。」 \end{tabularx} \\ \\ \relax
6:2 & \begin{tabularx}{0.7\textwidth}{X} 神吩咐摩西，對他說：「我是耶和華。 \end{tabularx} \\ \\ \relax
6:3 & \begin{tabularx}{0.7\textwidth}{X} 我從前向亞伯拉罕、以撒、雅各顯現為全能的神；至於我的名耶和華，我未曾讓他們知道。 \end{tabularx} \\ \\ \relax
6:4 & \begin{tabularx}{0.7\textwidth}{X} 我要與他們堅立我的約，要把迦南地，他們寄居的地賜給他們。 \end{tabularx} \\ \\ \relax
6:5 & \begin{tabularx}{0.7\textwidth}{X} 我聽見以色列人被埃及人奴役的哀聲，我就記念我的約。 \end{tabularx} \\ \\ \relax
6:6 & \begin{tabularx}{0.7\textwidth}{X} 所以你要對以色列人說：『我是耶和華；我要除去埃及人加給你們的勞役，救你們脫離他們的奴役。我要用伸出來的膀臂，藉嚴厲的懲罰救贖你們。 \end{tabularx} \\ \\ \relax
6:7 & \begin{tabularx}{0.7\textwidth}{X} 我要以你們為我的百姓，我也要作你們的神。我除去埃及人加給你們的勞役，你們就知道我是耶和華你們的神。 \end{tabularx} \\ \\ \relax
6:8 & \begin{tabularx}{0.7\textwidth}{X} 我起誓應許給亞伯拉罕、以撒、雅各的地，我要領你們進去，將那地賜給你們為業。我是耶和華。』」 \end{tabularx} \\ \\ \relax
6:9 & \begin{tabularx}{0.7\textwidth}{X} 摩西把這話告訴以色列人，但是他們因心裡愁煩，又因苦工，就不肯聽摩西的話。 \end{tabularx} \\ \\ \relax
6:10 & \begin{tabularx}{0.7\textwidth}{X} 耶和華吩咐摩西說： \end{tabularx} \\ \\ \relax
6:11 & \begin{tabularx}{0.7\textwidth}{X} 「你去對埃及王法老說，讓以色列人離開他的地。」 \end{tabularx} \\ \\ \relax
6:12 & \begin{tabularx}{0.7\textwidth}{X} 摩西在耶和華面前說：「看哪，以色列人尚且不聽我，法老怎麼會聽我這不會講話的人呢？」 \end{tabularx} \\ \\ \relax
6:13 & \begin{tabularx}{0.7\textwidth}{X} 耶和華吩咐摩西和亞倫，命令他們到以色列人和埃及王法老那裡，把以色列人從埃及地領出來。 \end{tabularx} \\ \\ \relax
6:14 & \begin{tabularx}{0.7\textwidth}{X} 以色列人族長的名字如下：以色列長子呂便的兒子是哈諾、法路、希斯倫、迦米；這是呂便的家族。 \end{tabularx} \\ \\ \relax
6:15 & \begin{tabularx}{0.7\textwidth}{X} 西緬的兒子是耶母利、雅憫、阿轄、雅斤、瑣轄，和迦南女子生的兒子掃羅；這是西緬的家族。 \end{tabularx} \\ \\ \relax
6:16 & \begin{tabularx}{0.7\textwidth}{X} 以下是利未的兒子按著家譜的名字：革順、哥轄、米拉利。利未一生的歲數是一百三十七歲。 \end{tabularx} \\ \\ \relax
6:17 & \begin{tabularx}{0.7\textwidth}{X} 革順的兒子按著家族是立尼、示每。 \end{tabularx} \\ \\ \relax
6:18 & \begin{tabularx}{0.7\textwidth}{X} 哥轄的兒子是暗蘭、以斯哈、希伯倫、烏薛。哥轄一生的歲數是一百三十三歲。 \end{tabularx} \\ \\ \relax
6:19 & \begin{tabularx}{0.7\textwidth}{X} 米拉利的兒子是抹利和母示；這是利未按著家譜的家族。 \end{tabularx} \\ \\ \relax
6:20 & \begin{tabularx}{0.7\textwidth}{X} 暗蘭娶了他父親的妹妹約基別為妻，她為他生了亞倫和摩西。暗蘭一生的歲數是一百三十七歲。 \end{tabularx} \\ \\ \relax
6:21 & \begin{tabularx}{0.7\textwidth}{X} 以斯哈的兒子是可拉、尼斐、細基利。 \end{tabularx} \\ \\ \relax
6:22 & \begin{tabularx}{0.7\textwidth}{X} 烏薛的兒子是米沙利、以利撒反、西提利。 \end{tabularx} \\ \\ \relax
6:23 & \begin{tabularx}{0.7\textwidth}{X} 亞倫娶了亞米拿達的女兒，拿順的妹妹，以利沙巴為妻，她為他生了拿答、亞比戶、以利亞撒、以他瑪。 \end{tabularx} \\ \\ \relax
6:24 & \begin{tabularx}{0.7\textwidth}{X} 可拉的兒子是亞惜、以利加拿、亞比亞撒；這是可拉的家族。 \end{tabularx} \\ \\ \relax
6:25 & \begin{tabularx}{0.7\textwidth}{X} 亞倫的兒子以利亞撒娶了普鐵的一個女兒為妻，她為他生了非尼哈。這是利未人按著家族的族長。 \end{tabularx} \\ \\ \relax
6:26 & \begin{tabularx}{0.7\textwidth}{X} 這就是曾聽見耶和華說「把以色列人按著隊伍從埃及地領出來」的亞倫和摩西， \end{tabularx} \\ \\ \relax
6:27 & \begin{tabularx}{0.7\textwidth}{X} 對埃及王法老說要將以色列人從埃及領出來的，也是這摩西和亞倫。 \end{tabularx} \\ \\ \relax
6:28 & \begin{tabularx}{0.7\textwidth}{X} 當耶和華在埃及地對摩西說話的時候， \end{tabularx} \\ \\ \relax
6:29 & \begin{tabularx}{0.7\textwidth}{X} 耶和華對摩西說：「我是耶和華；我對你所說的一切話，你都要告訴埃及王法老。」 \end{tabularx} \\ \\ \relax
6:30 & \begin{tabularx}{0.7\textwidth}{X} 摩西在耶和華面前說：「看哪，我是不會講話的人，法老怎麼會聽我呢？」 \end{tabularx} \\ \\
[1ex]
\hline
\hline
\end{longtable}


\newpage

\section{第65屆~港九培靈會~周永健~第三講}
\label{sec:985}
\textbf{神用大能的手行奇事──埃及人就知道我是耶和華}
\newline
\newline
連結: \href{https://www.hkbibleconference.org/session-message/view/985}{\texttt{https://www.hkbibleconference.org/session-message/view/985}}
\newline
\newline
\hyperref[sec:984]{< < < PREV SERMON < < <}
~
\hyperref[sec:toc]{[返目錄]}
~
\hyperref[sec:986]{> > > NEXT SERMON > > >}
\newline
\newline
出埃及記 7:1-12:51
\newline
\begin{longtable}{cl}
\hline
\hline
章節 & 經文 (和合本修訂版)\\
\hline
7:1 & \begin{tabularx}{0.7\textwidth}{X} 耶和華對摩西說：「我使你在法老面前像神一樣，你的哥哥亞倫是你的代言人。 \end{tabularx} \\ \\ \relax
7:2 & \begin{tabularx}{0.7\textwidth}{X} 凡我所吩咐你的，你都要說。你的哥哥亞倫要對法老說，讓以色列人離開他的地。 \end{tabularx} \\ \\ \relax
7:3 & \begin{tabularx}{0.7\textwidth}{X} 我要使法老的心固執，我也要在埃及地多行神蹟奇事。 \end{tabularx} \\ \\ \relax
7:4 & \begin{tabularx}{0.7\textwidth}{X} 法老必不聽從你們，因此我要伸手嚴厲地懲罰埃及，把我的軍隊，就是我的百姓以色列人從埃及地領出來。 \end{tabularx} \\ \\ \relax
7:5 & \begin{tabularx}{0.7\textwidth}{X} 我伸手攻擊埃及，把以色列人從他們中間領出來的時候，埃及人就知道我是耶和華。」 \end{tabularx} \\ \\ \relax
7:6 & \begin{tabularx}{0.7\textwidth}{X} 摩西和亞倫就去做；他們照耶和華吩咐的去做了。 \end{tabularx} \\ \\ \relax
7:7 & \begin{tabularx}{0.7\textwidth}{X} 摩西和亞倫與法老說話的時候，摩西八十歲，亞倫八十三歲。 \end{tabularx} \\ \\ \relax
7:8 & \begin{tabularx}{0.7\textwidth}{X} 耶和華對摩西和亞倫說： \end{tabularx} \\ \\ \relax
7:9 & \begin{tabularx}{0.7\textwidth}{X} 「法老若吩咐你們說：『你們行一件奇事吧！』你就對亞倫說：『把杖丟在法老面前！杖會變成蛇。』」 \end{tabularx} \\ \\ \relax
7:10 & \begin{tabularx}{0.7\textwidth}{X} 摩西和亞倫到法老那裡去，照耶和華所吩咐的去做。亞倫把杖丟在法老和他臣僕面前，杖就變成蛇。 \end{tabularx} \\ \\ \relax
7:11 & \begin{tabularx}{0.7\textwidth}{X} 法老也召了智慧人和行邪術的人來，這些埃及術士也用邪術照樣做。 \end{tabularx} \\ \\ \relax
7:12 & \begin{tabularx}{0.7\textwidth}{X} 他們各人丟下自己的杖，杖就變成蛇；但亞倫的杖吞了他們的杖。 \end{tabularx} \\ \\ \relax
7:13 & \begin{tabularx}{0.7\textwidth}{X} 法老心裡剛硬，不聽摩西和亞倫，正如耶和華所說的。 \end{tabularx} \\ \\ \relax
7:14 & \begin{tabularx}{0.7\textwidth}{X} 耶和華對摩西說：「法老心硬，不肯放百姓走。 \end{tabularx} \\ \\ \relax
7:15 & \begin{tabularx}{0.7\textwidth}{X} 明天早晨你要到法老那裡去，看哪，他出來往水邊去，你要到尼羅河邊去迎見他，手裡拿著那根變過蛇的杖。 \end{tabularx} \\ \\ \relax
7:16 & \begin{tabularx}{0.7\textwidth}{X} 你要對他說：『耶和華－希伯來人的神差派我到你這裡，說：放我的百姓走，到曠野事奉我。看哪，到如今你還是不聽。 \end{tabularx} \\ \\ \relax
7:17 & \begin{tabularx}{0.7\textwidth}{X} 耶和華如此說：看哪，我要用我手裡的杖擊打尼羅河中的水，水就變成血；這樣，你就知道我是耶和華。 \end{tabularx} \\ \\ \relax
7:18 & \begin{tabularx}{0.7\textwidth}{X} 河裡的魚必死，河也要發臭，埃及人就厭惡喝這河裡的水。』」 \end{tabularx} \\ \\ \relax
7:19 & \begin{tabularx}{0.7\textwidth}{X} 耶和華對摩西說：「你要對亞倫說：『拿你的杖，伸出你的手在埃及所有的水上，在他們的江、河、池塘，所有水聚集的地方上，叫水變成血。在埃及全地，無論在木器中，石器中，都必有血。』」 \end{tabularx} \\ \\ \relax
7:20 & \begin{tabularx}{0.7\textwidth}{X} 摩西和亞倫就照耶和華所吩咐的去做。亞倫在法老和他臣僕眼前舉杖擊打尼羅河裡的水，河裡的水都變成血了。 \end{tabularx} \\ \\ \relax
7:21 & \begin{tabularx}{0.7\textwidth}{X} 河裡的魚死了，河也臭了，埃及人就不能喝這河裡的水；埃及遍地都有了血。 \end{tabularx} \\ \\ \relax
7:22 & \begin{tabularx}{0.7\textwidth}{X} 但是，埃及的術士也用邪術照樣做了；法老心裡剛硬，不聽摩西和亞倫，正如耶和華所說的。 \end{tabularx} \\ \\ \relax
7:23 & \begin{tabularx}{0.7\textwidth}{X} 法老轉身回宮去，並不把這事放在心上。 \end{tabularx} \\ \\ \relax
7:24 & \begin{tabularx}{0.7\textwidth}{X} 所有的埃及人都沿著尼羅河邊挖掘，要找水喝，因為他們不能喝河裡的水。 \end{tabularx} \\ \\ \relax
7:25 & \begin{tabularx}{0.7\textwidth}{X} 耶和華擊打尼羅河後，過了七天。 \end{tabularx} \\ \\ \relax
8:1 & \begin{tabularx}{0.7\textwidth}{X} 耶和華對摩西說：「你要到法老那裡，對他說：『耶和華如此說：放我的百姓走，好事奉我。 \end{tabularx} \\ \\ \relax
8:2 & \begin{tabularx}{0.7\textwidth}{X} 你若不肯放他們走，看哪，我必以青蛙之災擊打你的疆土。 \end{tabularx} \\ \\ \relax
8:3 & \begin{tabularx}{0.7\textwidth}{X} 尼羅河要滋生青蛙；這青蛙要上來進你的宮殿和你的臥房，上你的床榻，進你臣僕的房屋，上你百姓的身上，進你的爐灶和你的揉麵盆。 \end{tabularx} \\ \\ \relax
8:4 & \begin{tabularx}{0.7\textwidth}{X} 這些青蛙要跳上你、你百姓和你眾臣僕的身上。』」 \end{tabularx} \\ \\ \relax
8:5 & \begin{tabularx}{0.7\textwidth}{X} 耶和華對摩西說：「你要對亞倫說：『伸出你手裡的杖在江、河、池塘上，把青蛙帶上埃及地來。』」 \end{tabularx} \\ \\ \relax
8:6 & \begin{tabularx}{0.7\textwidth}{X} 亞倫伸手在埃及的眾水上，青蛙就上來，遮滿了埃及地。 \end{tabularx} \\ \\ \relax
8:7 & \begin{tabularx}{0.7\textwidth}{X} 術士也用他們的邪術照樣去做，把青蛙帶上埃及地。 \end{tabularx} \\ \\ \relax
8:8 & \begin{tabularx}{0.7\textwidth}{X} 法老召摩西和亞倫來，說：「請你們祈求耶和華使這些青蛙離開我和我的百姓，我就讓這百姓去向耶和華獻祭。」 \end{tabularx} \\ \\ \relax
8:9 & \begin{tabularx}{0.7\textwidth}{X} 摩西對法老說：「悉聽尊便，告訴我何時為你、你臣僕和你的百姓祈求，使青蛙被剪除，離開你和你的宮殿，只留在尼羅河裡。」 \end{tabularx} \\ \\ \relax
8:10 & \begin{tabularx}{0.7\textwidth}{X} 他說：「明天。」摩西說：「就照你的話吧，為要叫你知道沒有像耶和華我們神的， \end{tabularx} \\ \\ \relax
8:11 & \begin{tabularx}{0.7\textwidth}{X} 青蛙必會離開你、你宮殿、你臣僕和你的百姓，只留在尼羅河裡。」 \end{tabularx} \\ \\ \relax
8:12 & \begin{tabularx}{0.7\textwidth}{X} 於是摩西和亞倫離開法老出去。摩西為了青蛙的事呼求耶和華，因為他帶來青蛙攪擾法老。 \end{tabularx} \\ \\ \relax
8:13 & \begin{tabularx}{0.7\textwidth}{X} 耶和華就照摩西的請求去做；在屋裡、院中、田間的青蛙都死了。 \end{tabularx} \\ \\ \relax
8:14 & \begin{tabularx}{0.7\textwidth}{X} 眾人把青蛙聚攏成堆，地就發出臭氣。 \end{tabularx} \\ \\ \relax
8:15 & \begin{tabularx}{0.7\textwidth}{X} 但法老見災禍舒緩了，就硬著心，不聽從他們，正如耶和華所說的。 \end{tabularx} \\ \\ \relax
8:16 & \begin{tabularx}{0.7\textwidth}{X} 耶和華對摩西說：「你要對亞倫說：『伸出你的杖擊打地上的塵土，使塵土在埃及全地變成蚊子。』」 \end{tabularx} \\ \\ \relax
8:17 & \begin{tabularx}{0.7\textwidth}{X} 他們就照樣做了。亞倫伸出他手裡的杖，擊打地上的塵土，人和牲畜身上就有了蚊子；埃及全地的塵土都變成蚊子了。 \end{tabularx} \\ \\ \relax
8:18 & \begin{tabularx}{0.7\textwidth}{X} 術士也用邪術要照樣產生蚊子，卻做不成。於是人和牲畜的身上都有了蚊子。 \end{tabularx} \\ \\ \relax
8:19 & \begin{tabularx}{0.7\textwidth}{X} 術士對法老說：「這是神的手指。」法老心裡剛硬，不聽摩西和亞倫，正如耶和華所說的。 \end{tabularx} \\ \\ \relax
8:20 & \begin{tabularx}{0.7\textwidth}{X} 耶和華對摩西說：「你要清早起來，站在法老面前。看哪，法老來到水邊，你就對他說：『耶和華如此說：放我的百姓走，好事奉我。 \end{tabularx} \\ \\ \relax
8:21 & \begin{tabularx}{0.7\textwidth}{X} 你若不放我的百姓走，看哪，我要派成群的蒼蠅到你、你臣僕和你百姓身上，進你的宮殿；埃及人的房屋和他們所住的地都要滿了成群的蒼蠅。 \end{tabularx} \\ \\ \relax
8:22 & \begin{tabularx}{0.7\textwidth}{X} 那一日，我必把我百姓所住的歌珊地分別出來，使那裡沒有成群的蒼蠅，好叫你知道我---耶和華是在全地之中。 \end{tabularx} \\ \\ \relax
8:23 & \begin{tabularx}{0.7\textwidth}{X} 我要施行救贖，區隔我的百姓和你的百姓。明天必有這神蹟。』」 \end{tabularx} \\ \\ \relax
8:24 & \begin{tabularx}{0.7\textwidth}{X} 耶和華就這樣做了。大群的蒼蠅進入法老的宮殿和他臣僕的房屋；在埃及全地，地就因這成群的蒼蠅毀壞了。 \end{tabularx} \\ \\ \relax
8:25 & \begin{tabularx}{0.7\textwidth}{X} 法老召了摩西和亞倫來，說：「去，在此地向你們的神獻祭。」 \end{tabularx} \\ \\ \relax
8:26 & \begin{tabularx}{0.7\textwidth}{X} 摩西說：「這樣做是不妥的，因為我們要獻給耶和華－我們神的祭物是埃及人所厭惡的；看哪，我們在埃及人眼前獻他們所厭惡的，他們豈不拿石頭打死我們嗎？ \end{tabularx} \\ \\ \relax
8:27 & \begin{tabularx}{0.7\textwidth}{X} 我們要遵照耶和華－我們神所吩咐我們的，往曠野去，走三天路程，向他獻祭。」 \end{tabularx} \\ \\ \relax
8:28 & \begin{tabularx}{0.7\textwidth}{X} 法老說：「我可以放你們走，在曠野向耶和華－你們的神獻祭，只是不可走得太遠。你們要為我祈禱。」 \end{tabularx} \\ \\ \relax
8:29 & \begin{tabularx}{0.7\textwidth}{X} 摩西說：「看哪，我要從你這裡出去祈求耶和華，使成群的蒼蠅明天離開法老、法老的臣僕和法老的百姓；法老卻不可再欺騙，不讓百姓去向耶和華獻祭。」 \end{tabularx} \\ \\ \relax
8:30 & \begin{tabularx}{0.7\textwidth}{X} 於是摩西離開法老，去祈求耶和華。 \end{tabularx} \\ \\ \relax
8:31 & \begin{tabularx}{0.7\textwidth}{X} 耶和華就照摩西的請求去做，使成群的蒼蠅離開法老、他的臣僕和他的百姓，一隻也沒有留下。 \end{tabularx} \\ \\ \relax
8:32 & \begin{tabularx}{0.7\textwidth}{X} 但這一次法老又硬著心，不放百姓走。 \end{tabularx} \\ \\ \relax
9:1 & \begin{tabularx}{0.7\textwidth}{X} 耶和華對摩西說：「你要到法老那裡，對他說：『耶和華－希伯來人的神如此說：放我的百姓走，好事奉我。 \end{tabularx} \\ \\ \relax
9:2 & \begin{tabularx}{0.7\textwidth}{X} 你若不肯放他們走，仍要強留他們， \end{tabularx} \\ \\ \relax
9:3 & \begin{tabularx}{0.7\textwidth}{X} 看哪，耶和華的手必以嚴重的瘟疫加在你田間的牲畜上，就是在馬、驢、駱駝、牛群和羊群的身上。 \end{tabularx} \\ \\ \relax
9:4 & \begin{tabularx}{0.7\textwidth}{X} 耶和華卻要分別以色列的牲畜和埃及的牲畜，凡屬以色列人的，一隻都不死。』」 \end{tabularx} \\ \\ \relax
9:5 & \begin{tabularx}{0.7\textwidth}{X} 耶和華就設定時間，說：「明天耶和華必在此地行這事。」 \end{tabularx} \\ \\ \relax
9:6 & \begin{tabularx}{0.7\textwidth}{X} 第二天，耶和華行了這事。埃及的牲畜全都死了，只是以色列人的牲畜，一隻都沒有死。 \end{tabularx} \\ \\ \relax
9:7 & \begin{tabularx}{0.7\textwidth}{X} 法老派人去，看哪，以色列人的牲畜連一隻都沒有死。可是法老硬著心，不放百姓走。 \end{tabularx} \\ \\ \relax
9:8 & \begin{tabularx}{0.7\textwidth}{X} 耶和華對摩西和亞倫說：「你們從爐裡滿滿捧出爐灰，摩西要在法老眼前把它撒在空中。 \end{tabularx} \\ \\ \relax
9:9 & \begin{tabularx}{0.7\textwidth}{X} 這灰要在埃及全地變成塵土，使埃及全地的人和牲畜身上起泡生瘡。」 \end{tabularx} \\ \\ \relax
9:10 & \begin{tabularx}{0.7\textwidth}{X} 摩西和亞倫取了爐灰，站在法老面前。摩西把它撒在空中，人和牲畜的身上就起泡生瘡了。 \end{tabularx} \\ \\ \relax
9:11 & \begin{tabularx}{0.7\textwidth}{X} 因為這瘡，術士在摩西面前站立不住，術士和所有埃及人的身上都生了瘡。 \end{tabularx} \\ \\ \relax
9:12 & \begin{tabularx}{0.7\textwidth}{X} 但耶和華任憑法老的心剛硬，不聽摩西和亞倫，正如耶和華對摩西所說的。 \end{tabularx} \\ \\ \relax
9:13 & \begin{tabularx}{0.7\textwidth}{X} 耶和華對摩西說：「你要清早起來，站在法老面前，對他說：『耶和華－希伯來人的神如此說：放我的百姓走，好事奉我。 \end{tabularx} \\ \\ \relax
9:14 & \begin{tabularx}{0.7\textwidth}{X} 因為這一次我要使一切的災禍臨到你自己，你臣僕和你百姓的身上，為要叫你知道在全地沒有像我的。 \end{tabularx} \\ \\ \relax
9:15 & \begin{tabularx}{0.7\textwidth}{X} 現在，我若伸手用瘟疫攻擊你和你的百姓，你就會從地上除滅了。 \end{tabularx} \\ \\ \relax
9:16 & \begin{tabularx}{0.7\textwidth}{X} 然而，我讓你存活，是為了要使你看見我的大能，並要使我的名傳遍全地。 \end{tabularx} \\ \\ \relax
9:17 & \begin{tabularx}{0.7\textwidth}{X} 可是你仍然向我的百姓自高自大，不放他們走。 \end{tabularx} \\ \\ \relax
9:18 & \begin{tabularx}{0.7\textwidth}{X} 看哪，明天大約這時候，我必使大量的冰雹降下，這是從埃及立國直到如今沒有出現過的。 \end{tabularx} \\ \\ \relax
9:19 & \begin{tabularx}{0.7\textwidth}{X} 現在，你要派人把你的牲畜和你田間一切所有的帶去躲避；任何在田間，無論是人是牲畜沒有回到屋內的，冰雹必降在他們身上，他們就必死。』」 \end{tabularx} \\ \\ \relax
9:20 & \begin{tabularx}{0.7\textwidth}{X} 法老的臣僕中，懼怕耶和華這話的，就讓他的奴僕和牲畜逃進屋裡。 \end{tabularx} \\ \\ \relax
9:21 & \begin{tabularx}{0.7\textwidth}{X} 但那不把耶和華這話放在心上的，就把他的奴僕和牲畜留在田裡。 \end{tabularx} \\ \\ \relax
9:22 & \begin{tabularx}{0.7\textwidth}{X} 耶和華對摩西說：「你向天伸出你的手，使冰雹降在埃及全地，降在埃及地的人和牲畜身上，以及田間各樣菜蔬上。」 \end{tabularx} \\ \\ \relax
9:23 & \begin{tabularx}{0.7\textwidth}{X} 摩西向天伸杖，耶和華就打雷下雹，有火降到地上；耶和華下雹在埃及地上。 \end{tabularx} \\ \\ \relax
9:24 & \begin{tabularx}{0.7\textwidth}{X} 那時，有雹，也有火在雹中閃爍，極其嚴重；自從埃及立國以來，全地沒有像這樣的。 \end{tabularx} \\ \\ \relax
9:25 & \begin{tabularx}{0.7\textwidth}{X} 在埃及全地，冰雹擊打田間所有的人和牲畜，擊打一切的菜蔬，也打壞了田間一切的樹木。 \end{tabularx} \\ \\ \relax
9:26 & \begin{tabularx}{0.7\textwidth}{X} 惟獨以色列人所住的歌珊地沒有冰雹。 \end{tabularx} \\ \\ \relax
9:27 & \begin{tabularx}{0.7\textwidth}{X} 法老差派人去召摩西和亞倫來，對他們說：「這一次我犯罪了。耶和華是公義的；我和我的百姓是邪惡的。 \end{tabularx} \\ \\ \relax
9:28 & \begin{tabularx}{0.7\textwidth}{X} 請你們祈求耶和華，因神的雷轟和冰雹已經夠了。我要放你們走，你們不用再留下來了。」 \end{tabularx} \\ \\ \relax
9:29 & \begin{tabularx}{0.7\textwidth}{X} 摩西對他說：「我一出城就向耶和華舉起雙手；雷必止住，也不再有冰雹，叫你知道地是屬於耶和華的。 \end{tabularx} \\ \\ \relax
9:30 & \begin{tabularx}{0.7\textwidth}{X} 至於你和你的臣僕，我知道你們仍然不敬畏耶和華神。」 \end{tabularx} \\ \\ \relax
9:31 & \begin{tabularx}{0.7\textwidth}{X} 那時，亞麻和大麥被摧毀了，因為大麥已經吐穗，亞麻也開了花。 \end{tabularx} \\ \\ \relax
9:32 & \begin{tabularx}{0.7\textwidth}{X} 只是小麥和粗麥沒有被摧毀，因為它們還沒有長成。 \end{tabularx} \\ \\ \relax
9:33 & \begin{tabularx}{0.7\textwidth}{X} 摩西離開法老出了城，向耶和華舉起雙手，雷和雹就止住，雨也不再下在地上了。 \end{tabularx} \\ \\ \relax
9:34 & \begin{tabularx}{0.7\textwidth}{X} 法老見雨、雹、雷止住，又再犯罪；他和他的臣僕都硬著心。 \end{tabularx} \\ \\ \relax
9:35 & \begin{tabularx}{0.7\textwidth}{X} 法老的心剛硬，不放以色列人走，正如耶和華藉著摩西所說的。 \end{tabularx} \\ \\ \relax
10:1 & \begin{tabularx}{0.7\textwidth}{X} 耶和華對摩西說：「你要到法老那裡，因我使他硬著心，也使他臣僕硬著心，為要在他們中間顯出我的這些神蹟來， \end{tabularx} \\ \\ \relax
10:2 & \begin{tabularx}{0.7\textwidth}{X} 並要叫你將我嚴厲對付埃及的事，和在他們中間所行的神蹟，傳於兒子和孫子的耳中，好叫你們知道我是耶和華。」 \end{tabularx} \\ \\ \relax
10:3 & \begin{tabularx}{0.7\textwidth}{X} 摩西和亞倫就到法老那裡，對他說：「耶和華－希伯來人的神這樣說：『你在我面前不肯謙卑要到幾時呢？放我的百姓走，好事奉我。 \end{tabularx} \\ \\ \relax
10:4 & \begin{tabularx}{0.7\textwidth}{X} 你若不肯放我的百姓走，看哪，明天我要使蝗蟲進入你的境內， \end{tabularx} \\ \\ \relax
10:5 & \begin{tabularx}{0.7\textwidth}{X} 遮滿地面，甚至地也看不見了。牠們要吃那冰雹後所剩，就是留給你們的；並且要吃那生長在田間的一切樹木。 \end{tabularx} \\ \\ \relax
10:6 & \begin{tabularx}{0.7\textwidth}{X} 你的宮殿和你眾臣僕的房屋，以及一切埃及人的房屋，都要被蝗蟲佔滿；你祖宗和你祖宗的祖宗在世以來，直到今日都沒有見過。』」摩西就轉身離開法老出去。 \end{tabularx} \\ \\ \relax
10:7 & \begin{tabularx}{0.7\textwidth}{X} 法老的臣僕對法老說：「這傢伙成為我們的羅網要到幾時呢？讓這些人去事奉耶和華－他們的神吧！埃及快要滅亡了，你還不知道嗎？」 \end{tabularx} \\ \\ \relax
10:8 & \begin{tabularx}{0.7\textwidth}{X} 於是摩西和亞倫被召回來見法老。法老對他們說：「去，事奉耶和華－你們的神吧！但要去的是哪些人呢？」 \end{tabularx} \\ \\ \relax
10:9 & \begin{tabularx}{0.7\textwidth}{X} 摩西說：「我們要帶著年老的和年少的同去，要帶著我們的兒子和女兒，以及我們的羊群牛群一起去，因為我們要向耶和華守節。」 \end{tabularx} \\ \\ \relax
10:10 & \begin{tabularx}{0.7\textwidth}{X} 法老對他們說：「願耶和華與你們同在吧！我若讓你們帶著你們的孩子同去，看，災禍就在你們面前！ \end{tabularx} \\ \\ \relax
10:11 & \begin{tabularx}{0.7\textwidth}{X} 不可都去！你們壯年人去事奉耶和華吧，因為這是你們所求的。」於是法老把他們從自己面前趕出去。 \end{tabularx} \\ \\ \relax
10:12 & \begin{tabularx}{0.7\textwidth}{X} 耶和華對摩西說：「你向埃及地伸出你的手，使蝗蟲上到埃及地，吃地上冰雹後所剩一切的植物。」 \end{tabularx} \\ \\ \relax
10:13 & \begin{tabularx}{0.7\textwidth}{X} 摩西就向埃及地伸杖；整整一晝一夜，耶和華使東風颳在埃及地上，到了早晨，東風把蝗蟲颳了來。 \end{tabularx} \\ \\ \relax
10:14 & \begin{tabularx}{0.7\textwidth}{X} 蝗蟲上到埃及全地，落在埃及全境，非常厲害；蝗蟲這麼多，是空前絕後的。 \end{tabularx} \\ \\ \relax
10:15 & \begin{tabularx}{0.7\textwidth}{X} 蝗蟲遮滿地面，地上一片黑暗。牠們吃盡了地上一切的植物和冰雹過後所剩樹上的果子。埃及全地，無論是樹木，是田間的植物，連一點綠的也沒有留下。 \end{tabularx} \\ \\ \relax
10:16 & \begin{tabularx}{0.7\textwidth}{X} 於是法老急忙召了摩西和亞倫來，說：「我得罪了耶和華－你們的神，又得罪了你們。 \end{tabularx} \\ \\ \relax
10:17 & \begin{tabularx}{0.7\textwidth}{X} 現在求你，就這一次，饒恕我的罪，祈求耶和華－你們的神救我脫離這次的死亡。」 \end{tabularx} \\ \\ \relax
10:18 & \begin{tabularx}{0.7\textwidth}{X} 摩西就離開法老，去祈求耶和華。 \end{tabularx} \\ \\ \relax
10:19 & \begin{tabularx}{0.7\textwidth}{X} 耶和華轉變風向，使強勁的西風吹來，把蝗蟲颳起，吹入紅海；在埃及全境連一隻也沒有留下。 \end{tabularx} \\ \\ \relax
10:20 & \begin{tabularx}{0.7\textwidth}{X} 但耶和華任憑法老的心剛硬，不放以色列人走。 \end{tabularx} \\ \\ \relax
10:21 & \begin{tabularx}{0.7\textwidth}{X} 耶和華對摩西說：「你向天伸出你的手，使黑暗籠罩埃及地；這黑暗甚至可以摸得到。」 \end{tabularx} \\ \\ \relax
10:22 & \begin{tabularx}{0.7\textwidth}{X} 摩西向天伸出他的手，濃密的黑暗就籠罩了埃及全地三天之久。 \end{tabularx} \\ \\ \relax
10:23 & \begin{tabularx}{0.7\textwidth}{X} 三天內，人人彼此看不見，誰也不敢起身離開原地；但所有以色列人住的地方卻有光。 \end{tabularx} \\ \\ \relax
10:24 & \begin{tabularx}{0.7\textwidth}{X} 法老就召摩西來，說：「去，事奉耶和華吧！只是你們的羊群牛群要留下來。你們的孩子可以和你們同去。」 \end{tabularx} \\ \\ \relax
10:25 & \begin{tabularx}{0.7\textwidth}{X} 摩西說：「你必須把祭物和燔祭牲交在我們手中，讓我們可以向耶和華我們的神獻祭。 \end{tabularx} \\ \\ \relax
10:26 & \begin{tabularx}{0.7\textwidth}{X} 我們的牲畜也要與我們同去，連一蹄也不留下，因為我們要從牲畜中挑選來事奉耶和華－我們的神。未到那裡之前，我們還不知道要用甚麼來事奉耶和華。」 \end{tabularx} \\ \\ \relax
10:27 & \begin{tabularx}{0.7\textwidth}{X} 但耶和華任憑法老的心剛硬，法老不肯放他們走。 \end{tabularx} \\ \\ \relax
10:28 & \begin{tabularx}{0.7\textwidth}{X} 法老對摩西說：「離開我去吧！你要小心，不要再見我的面，因為再見我面的那日，你就必死！」 \end{tabularx} \\ \\ \relax
10:29 & \begin{tabularx}{0.7\textwidth}{X} 摩西說：「就照你說的，我也不要再見你的面了！」 \end{tabularx} \\ \\ \relax
11:1 & \begin{tabularx}{0.7\textwidth}{X} 耶和華對摩西說：「我要再降一個災禍給法老和埃及，然後他必讓你們離開這裡。他放你們走的時候，一定會趕你們全都離開這裡。 \end{tabularx} \\ \\ \relax
11:2 & \begin{tabularx}{0.7\textwidth}{X} 你要傳於百姓耳中，叫他們男的女的各向鄰舍索取金器銀器。」 \end{tabularx} \\ \\ \relax
11:3 & \begin{tabularx}{0.7\textwidth}{X} 耶和華使埃及人看得起他的百姓，並且摩西在埃及地，在法老臣僕和百姓眼中看為偉大。 \end{tabularx} \\ \\ \relax
11:4 & \begin{tabularx}{0.7\textwidth}{X} 摩西說：「耶和華如此說：『約到半夜，我必出去走遍埃及。 \end{tabularx} \\ \\ \relax
11:5 & \begin{tabularx}{0.7\textwidth}{X} 凡在埃及地，從坐寶座的法老到推磨的婢女所生的長子，以及一切頭生的牲畜，都必死。 \end{tabularx} \\ \\ \relax
11:6 & \begin{tabularx}{0.7\textwidth}{X} 埃及全地必有大大的哀號，這將是空前絕後的。 \end{tabularx} \\ \\ \relax
11:7 & \begin{tabularx}{0.7\textwidth}{X} 至於以色列人中，無論是人是牲畜，連狗也不敢向他們吠叫，使你們知道耶和華區別埃及和以色列。』 \end{tabularx} \\ \\ \relax
11:8 & \begin{tabularx}{0.7\textwidth}{X} 你所有的這些臣僕都要下到我這裡，向我下拜說：『請你和跟從你的百姓都離開吧！』然後我才離開。」於是，摩西氣憤憤地離開法老出去了。 \end{tabularx} \\ \\ \relax
11:9 & \begin{tabularx}{0.7\textwidth}{X} 耶和華對摩西說：「法老必不聽你們，為了要使我在埃及地多行奇事。」 \end{tabularx} \\ \\ \relax
11:10 & \begin{tabularx}{0.7\textwidth}{X} 摩西和亞倫在法老面前行了這一切奇事，但耶和華任憑法老的心剛硬，不讓以色列人離開他的地。 \end{tabularx} \\ \\ \relax
12:1 & \begin{tabularx}{0.7\textwidth}{X} 耶和華在埃及地對摩西和亞倫說： \end{tabularx} \\ \\ \relax
12:2 & \begin{tabularx}{0.7\textwidth}{X} 「你們要以本月為正月，為一年之首。 \end{tabularx} \\ \\ \relax
12:3 & \begin{tabularx}{0.7\textwidth}{X} 你們要吩咐以色列全會眾說：本月初十，各人要按著家庭取羔羊，一家一隻羔羊。 \end{tabularx} \\ \\ \relax
12:4 & \begin{tabularx}{0.7\textwidth}{X} 若一家的人太少，吃不了一隻羔羊，就要按照人數和隔壁的鄰舍共取一隻；你們要按每人的食量來估算羔羊。 \end{tabularx} \\ \\ \relax
12:5 & \begin{tabularx}{0.7\textwidth}{X} 你們要從綿羊或山羊中取一隻無殘疾、一歲的公羔羊， \end{tabularx} \\ \\ \relax
12:6 & \begin{tabularx}{0.7\textwidth}{X} 要把牠留到本月十四日；那日黃昏的時候，以色列全會眾要把羔羊宰了。 \end{tabularx} \\ \\ \relax
12:7 & \begin{tabularx}{0.7\textwidth}{X} 他們要取一些血，塗在他們吃羔羊的房屋兩邊的門框上和門楣上。 \end{tabularx} \\ \\ \relax
12:8 & \begin{tabularx}{0.7\textwidth}{X} 當晚要吃羔羊的肉；要用火烤了，與無酵餅和苦菜一起吃。 \end{tabularx} \\ \\ \relax
12:9 & \begin{tabularx}{0.7\textwidth}{X} 不可吃生的，或用水煮的，要把羔羊連頭帶腿和內臟用火烤了吃。 \end{tabularx} \\ \\ \relax
12:10 & \begin{tabularx}{0.7\textwidth}{X} 一點也不可留到早晨；若有留到早晨的，要用火燒了。 \end{tabularx} \\ \\ \relax
12:11 & \begin{tabularx}{0.7\textwidth}{X} 你們要這樣吃羔羊：腰間束帶，腳上穿鞋，手中拿杖，快快地吃。這是耶和華的逾越。 \end{tabularx} \\ \\ \relax
12:12 & \begin{tabularx}{0.7\textwidth}{X} 因為那夜我要走遍埃及地，把埃及地一切頭生的，無論是人是牲畜，都擊殺了；我要對埃及所有的神明施行審判。我是耶和華。 \end{tabularx} \\ \\ \relax
12:13 & \begin{tabularx}{0.7\textwidth}{X} 這血要在你們所住的房屋上作記號；我一見這血，就逾越你們。我擊打埃及地的時候，災殃必不臨到你們身上施行毀滅。」 \end{tabularx} \\ \\ \relax
12:14 & \begin{tabularx}{0.7\textwidth}{X} 「你們要記念這日，世世代代守這日為耶和華的節日，作為你們永遠的定例。 \end{tabularx} \\ \\ \relax
12:15 & \begin{tabularx}{0.7\textwidth}{X} 你們要吃無酵餅七日。第一日要把酵從你們各家中除去，因為從第一日到第七日，任何吃有酵之物的，必從以色列中剪除。 \end{tabularx} \\ \\ \relax
12:16 & \begin{tabularx}{0.7\textwidth}{X} 第一日當有聖會，第七日也當有聖會。在這兩日，任何工作都不可做，只能預備各人的食物，這是惟一可做的工作。 \end{tabularx} \\ \\ \relax
12:17 & \begin{tabularx}{0.7\textwidth}{X} 你們要守除酵節，因為我在這一日把你們的軍隊從埃及地領出來。所以，你們要世世代代守這日，立為永遠的定例。 \end{tabularx} \\ \\ \relax
12:18 & \begin{tabularx}{0.7\textwidth}{X} 從正月十四日晚上，直到二十一日晚上，你們要吃無酵餅。 \end{tabularx} \\ \\ \relax
12:19 & \begin{tabularx}{0.7\textwidth}{X} 在你們各家中，七日之內不可有酵，因為凡吃有酵之物的，無論是寄居的，是本地的，必從以色列的會中剪除。 \end{tabularx} \\ \\ \relax
12:20 & \begin{tabularx}{0.7\textwidth}{X} 任何有酵的物，你們都不可吃；在你們一切的住處要吃無酵餅。」 \end{tabularx} \\ \\ \relax
12:21 & \begin{tabularx}{0.7\textwidth}{X} 於是，摩西召了以色列的眾長老來，對他們說：「你們要為家人取羔羊，把逾越的羔羊宰了。 \end{tabularx} \\ \\ \relax
12:22 & \begin{tabularx}{0.7\textwidth}{X} 要拿一把牛膝草，蘸盆裡的血，把盆裡的血塗在門楣上和兩邊的門框上。直到早晨你們誰也不可出自己家裡的門。 \end{tabularx} \\ \\ \relax
12:23 & \begin{tabularx}{0.7\textwidth}{X} 因為耶和華要走遍埃及，施行擊殺，他看見血在門楣上和兩邊的門框上，耶和華就必逾越那門，不讓滅命者進你們的家，施行擊殺。 \end{tabularx} \\ \\ \relax
12:24 & \begin{tabularx}{0.7\textwidth}{X} 你們要守這命令，作為你們和你們子孫永遠的定例。 \end{tabularx} \\ \\ \relax
12:25 & \begin{tabularx}{0.7\textwidth}{X} 日後，你們到了耶和華所應許賜給你們的那地，就要守這禮儀。 \end{tabularx} \\ \\ \relax
12:26 & \begin{tabularx}{0.7\textwidth}{X} 你們的兒女對你們說：『這禮儀是甚麼意思呢？』 \end{tabularx} \\ \\ \relax
12:27 & \begin{tabularx}{0.7\textwidth}{X} 你們就說：『這是獻給耶和華逾越節的祭物。當耶和華擊殺埃及人的時候，他逾越了以色列人在埃及的房屋，救了我們各家。』」於是百姓低頭敬拜。 \end{tabularx} \\ \\ \relax
12:28 & \begin{tabularx}{0.7\textwidth}{X} 以色列人就去做；他們照耶和華吩咐摩西和亞倫的去做了。 \end{tabularx} \\ \\ \relax
12:29 & \begin{tabularx}{0.7\textwidth}{X} 到了半夜，耶和華把埃及地所有頭生的，就是從坐寶座的法老，到關在牢裡的人的長子，以及一切頭生的牲畜，盡都殺了。 \end{tabularx} \\ \\ \relax
12:30 & \begin{tabularx}{0.7\textwidth}{X} 法老和他眾臣僕，以及所有的埃及人，都在夜間起來了。在埃及有大大的哀號，因為沒有一家不死人的。 \end{tabularx} \\ \\ \relax
12:31 & \begin{tabularx}{0.7\textwidth}{X} 夜間，法老召了摩西和亞倫來，說：「起來！你們和以色列人，都離開我的百姓出去，照你們所說的，去事奉耶和華吧！ \end{tabularx} \\ \\ \relax
12:32 & \begin{tabularx}{0.7\textwidth}{X} 照你們所說的，連羊群牛群也帶走，也為我祝福吧！」 \end{tabularx} \\ \\ \relax
12:33 & \begin{tabularx}{0.7\textwidth}{X} 埃及人催促百姓趕快離開那地，因為埃及人說：「我們都快死了。」 \end{tabularx} \\ \\ \relax
12:34 & \begin{tabularx}{0.7\textwidth}{X} 百姓就拿著沒有發酵的生麵，把揉麵盆包在衣服中，扛在肩上。 \end{tabularx} \\ \\ \relax
12:35 & \begin{tabularx}{0.7\textwidth}{X} 以色列人照摩西的話去做，向埃及人索取金器、銀器和衣裳。 \end{tabularx} \\ \\ \relax
12:36 & \begin{tabularx}{0.7\textwidth}{X} 耶和華使埃及人看得起他的百姓，埃及人就給了他們所要的。他們就掠奪了埃及人。 \end{tabularx} \\ \\ \relax
12:37 & \begin{tabularx}{0.7\textwidth}{X} 以色列人從蘭塞起程，往疏割去。除了小孩，步行的男人約有六十萬。 \end{tabularx} \\ \\ \relax
12:38 & \begin{tabularx}{0.7\textwidth}{X} 又有許多不同族群的人，以及眾多的羊群牛群，和他們一同上去。 \end{tabularx} \\ \\ \relax
12:39 & \begin{tabularx}{0.7\textwidth}{X} 他們用埃及帶出來的生麵烤成無酵餅。這生麵是沒有發酵的；因為他們被催促離開埃及，不能耽延，就沒有為自己預備食物。 \end{tabularx} \\ \\ \relax
12:40 & \begin{tabularx}{0.7\textwidth}{X} 以色列人住在埃及共四百三十年。 \end{tabularx} \\ \\ \relax
12:41 & \begin{tabularx}{0.7\textwidth}{X} 正滿四百三十年的那一天，耶和華的全軍從埃及地出來了。 \end{tabularx} \\ \\ \relax
12:42 & \begin{tabularx}{0.7\textwidth}{X} 這是向耶和華守的夜，他領他們出埃及地；這是以色列眾人世世代代要向耶和華守的夜。 \end{tabularx} \\ \\ \relax
12:43 & \begin{tabularx}{0.7\textwidth}{X} 耶和華對摩西和亞倫說：「逾越節的條例是這樣：外邦人不可吃這羔羊。 \end{tabularx} \\ \\ \relax
12:44 & \begin{tabularx}{0.7\textwidth}{X} 但是你們用銀子買來，又受過割禮的奴僕可以吃。 \end{tabularx} \\ \\ \relax
12:45 & \begin{tabularx}{0.7\textwidth}{X} 寄居的和雇工都不可吃。 \end{tabularx} \\ \\ \relax
12:46 & \begin{tabularx}{0.7\textwidth}{X} 應當在一個屋子裡吃，不可把肉帶到屋外，骨頭一根也不可折斷。 \end{tabularx} \\ \\ \relax
12:47 & \begin{tabularx}{0.7\textwidth}{X} 以色列全會眾都要守這禮儀。 \end{tabularx} \\ \\ \relax
12:48 & \begin{tabularx}{0.7\textwidth}{X} 若有外人寄居在你那裡，要向耶和華守逾越節，他所有的男子務要先受割禮，然後才可以當他是本地人，容許他守這禮儀。但未受割禮的都不可吃這羔羊。 \end{tabularx} \\ \\ \relax
12:49 & \begin{tabularx}{0.7\textwidth}{X} 本地人和寄居在你們中間的外人當守同一個條例。」 \end{tabularx} \\ \\ \relax
12:50 & \begin{tabularx}{0.7\textwidth}{X} 以色列眾人就去做，他們照耶和華吩咐摩西和亞倫的去做了。 \end{tabularx} \\ \\ \relax
12:51 & \begin{tabularx}{0.7\textwidth}{X} 正當那日，耶和華將以色列人按著他們的隊伍從埃及地領了出來。 \end{tabularx} \\ \\
[1ex]
\hline
\hline
\end{longtable}


\newpage

\section{第65屆~港九培靈會~周永健~第四講}
\label{sec:986}
\textbf{神使我們得勝──耶和華必為你們爭戰}
\newline
\newline
連結: \href{https://www.hkbibleconference.org/session-message/view/986}{\texttt{https://www.hkbibleconference.org/session-message/view/986}}
\newline
\newline
\hyperref[sec:985]{< < < PREV SERMON < < <}
~
\hyperref[sec:toc]{[返目錄]}
~
\hyperref[sec:987]{> > > NEXT SERMON > > >}
\newline
\newline
出埃及記 14:1-15:27
\newline
\begin{longtable}{cl}
\hline
\hline
章節 & 經文 (和合本修訂版)\\
\hline
14:1 & \begin{tabularx}{0.7\textwidth}{X} 耶和華吩咐摩西說： \end{tabularx} \\ \\ \relax
14:2 & \begin{tabularx}{0.7\textwidth}{X} 「你吩咐以色列人轉回，要在比‧哈希錄前面，密奪和海的中間，巴力‧洗分的前面安營。你們要在對面，靠近海邊安營。 \end{tabularx} \\ \\ \relax
14:3 & \begin{tabularx}{0.7\textwidth}{X} 以色列人這樣做，法老必說：『他們在此地迷了路，曠野把他們困住了。』 \end{tabularx} \\ \\ \relax
14:4 & \begin{tabularx}{0.7\textwidth}{X} 我要任憑法老的心剛硬，他要追趕他們。我必在法老和他全軍身上得榮耀，埃及人就知道我是耶和華。」於是以色列人照樣做了。 \end{tabularx} \\ \\ \relax
14:5 & \begin{tabularx}{0.7\textwidth}{X} 有人報告埃及王說：「百姓逃跑了！」法老和他的臣僕對百姓改變了心意，說：「我們放以色列人走，不再服事我們，我們怎麼會做這種事呢？」 \end{tabularx} \\ \\ \relax
14:6 & \begin{tabularx}{0.7\textwidth}{X} 法老就預備戰車，帶領他的軍兵同去， \end{tabularx} \\ \\ \relax
14:7 & \begin{tabularx}{0.7\textwidth}{X} 他帶了六百輛特選的戰車和埃及所有的戰車，每輛都有軍官。 \end{tabularx} \\ \\ \relax
14:8 & \begin{tabularx}{0.7\textwidth}{X} 耶和華任憑埃及王法老的心剛硬，他就追趕以色列人；以色列人卻抬起頭來出去了。 \end{tabularx} \\ \\ \relax
14:9 & \begin{tabularx}{0.7\textwidth}{X} 埃及人追趕他們，法老一切的馬匹、戰車、戰車長，與軍兵就在海邊上，靠近比‧哈希錄，在巴力‧洗分的前面，在他們安營的地方追上了。 \end{tabularx} \\ \\ \relax
14:10 & \begin{tabularx}{0.7\textwidth}{X} 法老逼近的時候，以色列人舉目，看哪，埃及人追來了，就非常懼怕，以色列人向耶和華哀求。 \end{tabularx} \\ \\ \relax
14:11 & \begin{tabularx}{0.7\textwidth}{X} 他們對摩西說：「難道埃及沒有墳地，你要把我們帶來死在曠野嗎？你為甚麼這樣待我們，將我們從埃及領出來呢？ \end{tabularx} \\ \\ \relax
14:12 & \begin{tabularx}{0.7\textwidth}{X} 我們在埃及豈沒有對你說過，不要攪擾我們，讓我們服事埃及人嗎？因為服事埃及人總比死在曠野好。」 \end{tabularx} \\ \\ \relax
14:13 & \begin{tabularx}{0.7\textwidth}{X} 摩西對百姓說：「不要怕，要站穩，看耶和華今天向你們所要施行的拯救，因為你們今天所看見的埃及人必永遠不再看見了。 \end{tabularx} \\ \\ \relax
14:14 & \begin{tabularx}{0.7\textwidth}{X} 耶和華必為你們爭戰，你們要安靜！」 \end{tabularx} \\ \\ \relax
14:15 & \begin{tabularx}{0.7\textwidth}{X} 耶和華對摩西說：「你為甚麼向我哀求呢？你吩咐以色列人往前走。 \end{tabularx} \\ \\ \relax
14:16 & \begin{tabularx}{0.7\textwidth}{X} 你舉手向海伸杖，把水分開。以色列人要下到海中，走在乾地上。 \end{tabularx} \\ \\ \relax
14:17 & \begin{tabularx}{0.7\textwidth}{X} 看哪，我要任憑埃及人的心剛硬，他們就跟著下去。我要在法老和他的全軍、戰車、戰車長身上得榮耀。 \end{tabularx} \\ \\ \relax
14:18 & \begin{tabularx}{0.7\textwidth}{X} 我在法老和他的戰車、戰車長身上得榮耀的時候，埃及人就知道我是耶和華。」 \end{tabularx} \\ \\ \relax
14:19 & \begin{tabularx}{0.7\textwidth}{X} 在以色列營前行走的神的使者移動，走到他們後面；雲也從他們的前面移動，站在他們後面。 \end{tabularx} \\ \\ \relax
14:20 & \begin{tabularx}{0.7\textwidth}{X} 它來到埃及營和以色列營的中間：一邊有雲和黑暗，另一邊它照亮夜晚，整夜彼此不得接近。 \end{tabularx} \\ \\ \relax
14:21 & \begin{tabularx}{0.7\textwidth}{X} 摩西向海伸手，耶和華就用強勁的東風，使海水在一夜間退去，海就成了乾地；水分開了。 \end{tabularx} \\ \\ \relax
14:22 & \begin{tabularx}{0.7\textwidth}{X} 以色列人下到海中，走在乾地上，水在他們左右成了牆壁。 \end{tabularx} \\ \\ \relax
14:23 & \begin{tabularx}{0.7\textwidth}{X} 埃及人追趕他們，法老一切的馬匹、戰車和戰車長都跟著下到海中。 \end{tabularx} \\ \\ \relax
14:24 & \begin{tabularx}{0.7\textwidth}{X} 破曉時分，耶和華從雲柱、火柱中瞭望埃及的軍兵，使埃及的軍兵混亂。 \end{tabularx} \\ \\ \relax
14:25 & \begin{tabularx}{0.7\textwidth}{X} 他使他們的車輪脫落，難以前行，埃及人說：「我們從以色列人面前逃跑吧！因耶和華為他們作戰，攻擊埃及了。」 \end{tabularx} \\ \\ \relax
14:26 & \begin{tabularx}{0.7\textwidth}{X} 耶和華對摩西說：「你要向海伸手，使水回流到埃及人，他們的戰車和戰車長身上。」 \end{tabularx} \\ \\ \relax
14:27 & \begin{tabularx}{0.7\textwidth}{X} 摩西就向海伸手，到了天亮的時候，海恢復原狀。埃及人逃避水的時候，耶和華把他們推入海中。 \end{tabularx} \\ \\ \relax
14:28 & \begin{tabularx}{0.7\textwidth}{X} 海水回流，淹沒了戰車和戰車長，以及那些跟著以色列人下到海中的法老全軍，連一個也沒有剩下。 \end{tabularx} \\ \\ \relax
14:29 & \begin{tabularx}{0.7\textwidth}{X} 以色列人卻在海中走乾地，水在他們的左右成了牆壁。 \end{tabularx} \\ \\ \relax
14:30 & \begin{tabularx}{0.7\textwidth}{X} 那一日，耶和華拯救以色列脫離埃及人的手。以色列人看見埃及人死在海邊。 \end{tabularx} \\ \\ \relax
14:31 & \begin{tabularx}{0.7\textwidth}{X} 以色列人看見耶和華向埃及人所施展的大能，百姓就敬畏耶和華，並且信服耶和華和他的僕人摩西。 \end{tabularx} \\ \\ \relax
15:1 & \begin{tabularx}{0.7\textwidth}{X} 那時，摩西和以色列人向耶和華唱這歌，說：「我要向耶和華歌唱，因他大大得勝，將馬和騎馬的投在海中。 \end{tabularx} \\ \\ \relax
15:2 & \begin{tabularx}{0.7\textwidth}{X} 耶和華是我的力量，是我的詩歌，他也成了我的拯救。這是我的神，我要讚美他；我父親的神，我要尊崇他。 \end{tabularx} \\ \\ \relax
15:3 & \begin{tabularx}{0.7\textwidth}{X} 耶和華是戰士；耶和華是他的名。 \end{tabularx} \\ \\ \relax
15:4 & \begin{tabularx}{0.7\textwidth}{X} 「法老的戰車、軍兵，他已拋在海中；法老精選的軍官都沉於紅海。 \end{tabularx} \\ \\ \relax
15:5 & \begin{tabularx}{0.7\textwidth}{X} 深水淹沒他們；他們好像石頭墜到深處。 \end{tabularx} \\ \\ \relax
15:6 & \begin{tabularx}{0.7\textwidth}{X} 耶和華啊，你的右手施展能力，大顯榮耀；耶和華啊，你的右手摔碎仇敵。 \end{tabularx} \\ \\ \relax
15:7 & \begin{tabularx}{0.7\textwidth}{X} 你大發威嚴，摧毀了你的敵人；你發出烈怒，吞滅他們如同碎秸。 \end{tabularx} \\ \\ \relax
15:8 & \begin{tabularx}{0.7\textwidth}{X} 因你鼻中的氣，水就聚成堆，大水豎立如壘，海的中心深水凝結。 \end{tabularx} \\ \\ \relax
15:9 & \begin{tabularx}{0.7\textwidth}{X} 仇敵說：『我要追趕，我要追上，我要分擄物，在他們身上滿足我的心願，我要拔刀，親手毀滅他們。』 \end{tabularx} \\ \\ \relax
15:10 & \begin{tabularx}{0.7\textwidth}{X} 你用風一吹，海水就淹沒他們；他們像鉛沉在大水之中。 \end{tabularx} \\ \\ \relax
15:11 & \begin{tabularx}{0.7\textwidth}{X} 「耶和華啊，眾神明中，誰能像你？誰能像你，至聖至榮，可頌可畏，施行奇事！ \end{tabularx} \\ \\ \relax
15:12 & \begin{tabularx}{0.7\textwidth}{X} 你伸出右手，地就吞滅他們。 \end{tabularx} \\ \\ \relax
15:13 & \begin{tabularx}{0.7\textwidth}{X} 「你以慈愛引領你所救贖的百姓；你以能力引導他們到你的聖所。 \end{tabularx} \\ \\ \relax
15:14 & \begin{tabularx}{0.7\textwidth}{X} 萬民聽見就戰抖；疼痛抓住非利士的居民。 \end{tabularx} \\ \\ \relax
15:15 & \begin{tabularx}{0.7\textwidth}{X} 那時，以東的族長驚惶，摩押的英雄被戰兢抓住，迦南所有的居民都融化。 \end{tabularx} \\ \\ \relax
15:16 & \begin{tabularx}{0.7\textwidth}{X} 驚駭恐懼臨到他們；耶和華啊，因你膀臂的大能，他們如石頭寂靜不動，等候你百姓過去，等候你所贖的百姓過去。 \end{tabularx} \\ \\ \relax
15:17 & \begin{tabularx}{0.7\textwidth}{X} 你要將他們領進去，栽在你產業的山上，耶和華啊，就是你為自己所造的住處，主啊，就是你手所建立的聖所。 \end{tabularx} \\ \\ \relax
15:18 & \begin{tabularx}{0.7\textwidth}{X} 耶和華必作王，直到永永遠遠！」 \end{tabularx} \\ \\ \relax
15:19 & \begin{tabularx}{0.7\textwidth}{X} 法老的馬匹、戰車和戰車長下到海中，耶和華使海水回流到他們身上；以色列人卻走在海中的乾地上。 \end{tabularx} \\ \\ \relax
15:20 & \begin{tabularx}{0.7\textwidth}{X} 那時，米利暗女先知，亞倫的姊姊，手裡拿著鈴鼓；眾婦女也跟她出去打鼓跳舞。 \end{tabularx} \\ \\ \relax
15:21 & \begin{tabularx}{0.7\textwidth}{X} 米利暗回應他們：「你們要歌頌耶和華，因他大大得勝，將馬和騎馬的投在海中。」 \end{tabularx} \\ \\ \relax
15:22 & \begin{tabularx}{0.7\textwidth}{X} 摩西領以色列人從紅海起程，到了書珥的曠野，在曠野走了三天，找不到水。 \end{tabularx} \\ \\ \relax
15:23 & \begin{tabularx}{0.7\textwidth}{X} 到了瑪拉，他們不能喝瑪拉的水，因為水是苦的；所以那地名叫瑪拉。 \end{tabularx} \\ \\ \relax
15:24 & \begin{tabularx}{0.7\textwidth}{X} 百姓就向摩西發怨言，說：「我們喝甚麼呢？」 \end{tabularx} \\ \\ \relax
15:25 & \begin{tabularx}{0.7\textwidth}{X} 摩西呼求耶和華，耶和華指示他一棵樹。他把樹丟在水裡，水就變甜了。耶和華在那裡為他們定了律例、典章，在那裡考驗他們。 \end{tabularx} \\ \\ \relax
15:26 & \begin{tabularx}{0.7\textwidth}{X} 他說：「你若留心聽從耶和華－你神的話，行我眼中看為正的事，側耳聽我的誡令，遵守我一切的律例，我就不將所加於埃及人的疾病加在你身上，因為我是醫治你的耶和華。」 \end{tabularx} \\ \\ \relax
15:27 & \begin{tabularx}{0.7\textwidth}{X} 他們到了以琳，在那裡有十二股水泉，七十棵棕樹；他們就在那裡的水邊安營。 \end{tabularx} \\ \\
[1ex]
\hline
\hline
\end{longtable}


\newpage

\section{第65屆~港九培靈會~周永健~第五講}
\label{sec:987}
\textbf{神將我們分別為聖──如鷹將你們背在翅膀上帶來歸我}
\newline
\newline
連結: \href{https://www.hkbibleconference.org/session-message/view/987}{\texttt{https://www.hkbibleconference.org/session-message/view/987}}
\newline
\newline
\hyperref[sec:986]{< < < PREV SERMON < < <}
~
\hyperref[sec:toc]{[返目錄]}
~
\hyperref[sec:988]{> > > NEXT SERMON > > >}
\newline
\newline
出埃及記 19:1-19:25
\newline
\begin{longtable}{cl}
\hline
\hline
章節 & 經文 (和合本修訂版)\\
\hline
19:1 & \begin{tabularx}{0.7\textwidth}{X} 以色列人出埃及地以後，第三個月的初一，就在那一天他們來到了西奈的曠野。 \end{tabularx} \\ \\ \relax
19:2 & \begin{tabularx}{0.7\textwidth}{X} 他們從利非訂起程，來到西奈的曠野，在那裡的山下安營。 \end{tabularx} \\ \\ \relax
19:3 & \begin{tabularx}{0.7\textwidth}{X} 摩西到神那裡，耶和華從山上呼喚他說：「你要這樣告訴雅各家，對以色列人說： \end{tabularx} \\ \\ \relax
19:4 & \begin{tabularx}{0.7\textwidth}{X} 『我向埃及人所行的事，你們都看見了， 我如鷹將你們背在翅膀上，帶你們來歸我。 \end{tabularx} \\ \\ \relax
19:5 & \begin{tabularx}{0.7\textwidth}{X} 如今你們若真的聽從我的話，遵守我的約，就要在萬民中作屬我的子民，因為全地都是我的。 \end{tabularx} \\ \\ \relax
19:6 & \begin{tabularx}{0.7\textwidth}{X} 你們要歸我作祭司的國度，為神聖的國民。』這些話你要告訴以色列人。」 \end{tabularx} \\ \\ \relax
19:7 & \begin{tabularx}{0.7\textwidth}{X} 摩西去召了百姓中的長老來，將耶和華吩咐他的話當面告訴他們。 \end{tabularx} \\ \\ \relax
19:8 & \begin{tabularx}{0.7\textwidth}{X} 百姓都同聲回答：「凡耶和華所說的，我們一定遵行。」摩西就將百姓的話回覆耶和華。 \end{tabularx} \\ \\ \relax
19:9 & \begin{tabularx}{0.7\textwidth}{X} 耶和華對摩西說：「看哪，我要在密雲中臨到你那裡，叫百姓在我與你說話的時候可以聽見，就可以永遠相信你了。」於是，摩西將百姓的話稟告耶和華。 \end{tabularx} \\ \\ \relax
19:10 & \begin{tabularx}{0.7\textwidth}{X} 耶和華對摩西說：「你往百姓那裡去，使他們今天明天分別為聖，又叫他們洗衣服。 \end{tabularx} \\ \\ \relax
19:11 & \begin{tabularx}{0.7\textwidth}{X} 第三天要預備好，因為第三天耶和華要在眾百姓眼前降臨在西奈山。 \end{tabularx} \\ \\ \relax
19:12 & \begin{tabularx}{0.7\textwidth}{X} 你要在山的周圍給百姓劃定界限，說：『你們當謹慎，不可上山去，也不可摸山的邊界。凡摸這山的，必被處死。 \end{tabularx} \\ \\ \relax
19:13 & \begin{tabularx}{0.7\textwidth}{X} 不可用手碰他，要用石頭打死，或射死；無論是人是牲畜，都不可活。』到角聲拉長的時候，他們才可到山腳來。」 \end{tabularx} \\ \\ \relax
19:14 & \begin{tabularx}{0.7\textwidth}{X} 摩西下山到百姓那裡去，使他們分別為聖，他們就洗衣服。 \end{tabularx} \\ \\ \relax
19:15 & \begin{tabularx}{0.7\textwidth}{X} 他對百姓說：「第三天要預備好；不可親近女人。」 \end{tabularx} \\ \\ \relax
19:16 & \begin{tabularx}{0.7\textwidth}{X} 到了第三天早晨，山上有雷轟、閃電和密雲，並且角聲非常響亮，營中的百姓盡都戰抖。 \end{tabularx} \\ \\ \relax
19:17 & \begin{tabularx}{0.7\textwidth}{X} 摩西率領百姓出營迎見神，都站在山下。 \end{tabularx} \\ \\ \relax
19:18 & \begin{tabularx}{0.7\textwidth}{X} 西奈山全山冒煙，因為耶和華在火中降臨山上。山的煙霧上騰，彷彿燒窯，整座山劇烈震動。 \end{tabularx} \\ \\ \relax
19:19 & \begin{tabularx}{0.7\textwidth}{X} 角聲越來越響，摩西說話，神以聲音回答他。 \end{tabularx} \\ \\ \relax
19:20 & \begin{tabularx}{0.7\textwidth}{X} 耶和華降臨在西奈山頂上，耶和華召摩西上山頂，摩西就上去了。 \end{tabularx} \\ \\ \relax
19:21 & \begin{tabularx}{0.7\textwidth}{X} 耶和華對摩西說：「你下去警告百姓，免得他們闖過來看耶和華，就會有許多人死亡。 \end{tabularx} \\ \\ \relax
19:22 & \begin{tabularx}{0.7\textwidth}{X} 那些親近耶和華的祭司也要把自己分別為聖，免得耶和華忽然出來擊殺他們。」 \end{tabularx} \\ \\ \relax
19:23 & \begin{tabularx}{0.7\textwidth}{X} 摩西對耶和華說：「百姓不能上西奈山，因為你已經警告我們說：『要在山的周圍劃定界限，使山成聖。』」 \end{tabularx} \\ \\ \relax
19:24 & \begin{tabularx}{0.7\textwidth}{X} 耶和華對他說：「下去吧，你要和亞倫一起上來；只是祭司和百姓不可闖上來到耶和華這裡，免得耶和華忽然出來擊殺他們。」 \end{tabularx} \\ \\ \relax
19:25 & \begin{tabularx}{0.7\textwidth}{X} 於是，摩西下到百姓那裡告訴他們。 \end{tabularx} \\ \\
[1ex]
\hline
\hline
\end{longtable}


\newpage

\section{第65屆~港九培靈會~周永健~第六講}
\label{sec:988}
\textbf{神體恤我們的軟弱──耶和華是在我們中間不是}
\newline
\newline
連結: \href{https://www.hkbibleconference.org/session-message/view/988}{\texttt{https://www.hkbibleconference.org/session-message/view/988}}
\newline
\newline
\hyperref[sec:987]{< < < PREV SERMON < < <}
~
\hyperref[sec:toc]{[返目錄]}
~
\hyperref[sec:989]{> > > NEXT SERMON > > >}
\newline
\newline
出埃及記 15:1-17:16
\newline
\begin{longtable}{cl}
\hline
\hline
章節 & 經文 (和合本修訂版)\\
\hline
15:1 & \begin{tabularx}{0.7\textwidth}{X} 那時，摩西和以色列人向耶和華唱這歌，說：「我要向耶和華歌唱，因他大大得勝，將馬和騎馬的投在海中。 \end{tabularx} \\ \\ \relax
15:2 & \begin{tabularx}{0.7\textwidth}{X} 耶和華是我的力量，是我的詩歌，他也成了我的拯救。這是我的神，我要讚美他；我父親的神，我要尊崇他。 \end{tabularx} \\ \\ \relax
15:3 & \begin{tabularx}{0.7\textwidth}{X} 耶和華是戰士；耶和華是他的名。 \end{tabularx} \\ \\ \relax
15:4 & \begin{tabularx}{0.7\textwidth}{X} 「法老的戰車、軍兵，他已拋在海中；法老精選的軍官都沉於紅海。 \end{tabularx} \\ \\ \relax
15:5 & \begin{tabularx}{0.7\textwidth}{X} 深水淹沒他們；他們好像石頭墜到深處。 \end{tabularx} \\ \\ \relax
15:6 & \begin{tabularx}{0.7\textwidth}{X} 耶和華啊，你的右手施展能力，大顯榮耀；耶和華啊，你的右手摔碎仇敵。 \end{tabularx} \\ \\ \relax
15:7 & \begin{tabularx}{0.7\textwidth}{X} 你大發威嚴，摧毀了你的敵人；你發出烈怒，吞滅他們如同碎秸。 \end{tabularx} \\ \\ \relax
15:8 & \begin{tabularx}{0.7\textwidth}{X} 因你鼻中的氣，水就聚成堆，大水豎立如壘，海的中心深水凝結。 \end{tabularx} \\ \\ \relax
15:9 & \begin{tabularx}{0.7\textwidth}{X} 仇敵說：『我要追趕，我要追上，我要分擄物，在他們身上滿足我的心願，我要拔刀，親手毀滅他們。』 \end{tabularx} \\ \\ \relax
15:10 & \begin{tabularx}{0.7\textwidth}{X} 你用風一吹，海水就淹沒他們；他們像鉛沉在大水之中。 \end{tabularx} \\ \\ \relax
15:11 & \begin{tabularx}{0.7\textwidth}{X} 「耶和華啊，眾神明中，誰能像你？誰能像你，至聖至榮，可頌可畏，施行奇事！ \end{tabularx} \\ \\ \relax
15:12 & \begin{tabularx}{0.7\textwidth}{X} 你伸出右手，地就吞滅他們。 \end{tabularx} \\ \\ \relax
15:13 & \begin{tabularx}{0.7\textwidth}{X} 「你以慈愛引領你所救贖的百姓；你以能力引導他們到你的聖所。 \end{tabularx} \\ \\ \relax
15:14 & \begin{tabularx}{0.7\textwidth}{X} 萬民聽見就戰抖；疼痛抓住非利士的居民。 \end{tabularx} \\ \\ \relax
15:15 & \begin{tabularx}{0.7\textwidth}{X} 那時，以東的族長驚惶，摩押的英雄被戰兢抓住，迦南所有的居民都融化。 \end{tabularx} \\ \\ \relax
15:16 & \begin{tabularx}{0.7\textwidth}{X} 驚駭恐懼臨到他們；耶和華啊，因你膀臂的大能，他們如石頭寂靜不動，等候你百姓過去，等候你所贖的百姓過去。 \end{tabularx} \\ \\ \relax
15:17 & \begin{tabularx}{0.7\textwidth}{X} 你要將他們領進去，栽在你產業的山上，耶和華啊，就是你為自己所造的住處，主啊，就是你手所建立的聖所。 \end{tabularx} \\ \\ \relax
15:18 & \begin{tabularx}{0.7\textwidth}{X} 耶和華必作王，直到永永遠遠！」 \end{tabularx} \\ \\ \relax
15:19 & \begin{tabularx}{0.7\textwidth}{X} 法老的馬匹、戰車和戰車長下到海中，耶和華使海水回流到他們身上；以色列人卻走在海中的乾地上。 \end{tabularx} \\ \\ \relax
15:20 & \begin{tabularx}{0.7\textwidth}{X} 那時，米利暗女先知，亞倫的姊姊，手裡拿著鈴鼓；眾婦女也跟她出去打鼓跳舞。 \end{tabularx} \\ \\ \relax
15:21 & \begin{tabularx}{0.7\textwidth}{X} 米利暗回應他們：「你們要歌頌耶和華，因他大大得勝，將馬和騎馬的投在海中。」 \end{tabularx} \\ \\ \relax
15:22 & \begin{tabularx}{0.7\textwidth}{X} 摩西領以色列人從紅海起程，到了書珥的曠野，在曠野走了三天，找不到水。 \end{tabularx} \\ \\ \relax
15:23 & \begin{tabularx}{0.7\textwidth}{X} 到了瑪拉，他們不能喝瑪拉的水，因為水是苦的；所以那地名叫瑪拉。 \end{tabularx} \\ \\ \relax
15:24 & \begin{tabularx}{0.7\textwidth}{X} 百姓就向摩西發怨言，說：「我們喝甚麼呢？」 \end{tabularx} \\ \\ \relax
15:25 & \begin{tabularx}{0.7\textwidth}{X} 摩西呼求耶和華，耶和華指示他一棵樹。他把樹丟在水裡，水就變甜了。耶和華在那裡為他們定了律例、典章，在那裡考驗他們。 \end{tabularx} \\ \\ \relax
15:26 & \begin{tabularx}{0.7\textwidth}{X} 他說：「你若留心聽從耶和華－你神的話，行我眼中看為正的事，側耳聽我的誡令，遵守我一切的律例，我就不將所加於埃及人的疾病加在你身上，因為我是醫治你的耶和華。」 \end{tabularx} \\ \\ \relax
15:27 & \begin{tabularx}{0.7\textwidth}{X} 他們到了以琳，在那裡有十二股水泉，七十棵棕樹；他們就在那裡的水邊安營。 \end{tabularx} \\ \\ \relax
16:1 & \begin{tabularx}{0.7\textwidth}{X} 以色列全會眾從以琳起程，在出埃及之後第二個月十五日到了以琳和西奈中間，汛的曠野。 \end{tabularx} \\ \\ \relax
16:2 & \begin{tabularx}{0.7\textwidth}{X} 以色列全會眾在曠野向摩西和亞倫發怨言。 \end{tabularx} \\ \\ \relax
16:3 & \begin{tabularx}{0.7\textwidth}{X} 以色列人對他們說：「我們寧願在埃及地死在耶和華手中！那時我們坐在肉鍋旁，吃餅得飽。你們卻將我們領出來，到這曠野，要叫這全會眾都餓死啊！」 \end{tabularx} \\ \\ \relax
16:4 & \begin{tabularx}{0.7\textwidth}{X} 耶和華對摩西說：「看哪，我要從天降食物給你們。百姓可以出去，每天收集當天的分量。這樣，我就可以考驗他們是否遵行我的指示。 \end{tabularx} \\ \\ \relax
16:5 & \begin{tabularx}{0.7\textwidth}{X} 到第六天，他們預備食物，所收集的分量要比每天所收的多一倍。」 \end{tabularx} \\ \\ \relax
16:6 & \begin{tabularx}{0.7\textwidth}{X} 摩西和亞倫對以色列眾人說：「到了晚上，你們就知道是耶和華將你們從埃及地領出來的。 \end{tabularx} \\ \\ \relax
16:7 & \begin{tabularx}{0.7\textwidth}{X} 早晨，你們要看見耶和華的榮耀，因為耶和華聽見你們向他所發的怨言了。我們算甚麼，你們竟然向我們發怨言呢？」 \end{tabularx} \\ \\ \relax
16:8 & \begin{tabularx}{0.7\textwidth}{X} 摩西又說：「耶和華晚上必給你們肉吃，早晨必給你們食物得飽，因為耶和華已經聽見你們向他所發的怨言。我們算甚麼呢？你們的怨言不是向我們發的，而是向耶和華發的。」 \end{tabularx} \\ \\ \relax
16:9 & \begin{tabularx}{0.7\textwidth}{X} 摩西對亞倫說：「你對以色列全會眾說：『你們來到耶和華面前，因為他已經聽見你們的怨言了。』」 \end{tabularx} \\ \\ \relax
16:10 & \begin{tabularx}{0.7\textwidth}{X} 亞倫正對以色列全會眾說話的時候，他們轉向曠野，看哪，耶和華的榮光在雲中顯現。 \end{tabularx} \\ \\ \relax
16:11 & \begin{tabularx}{0.7\textwidth}{X} 耶和華吩咐摩西說： \end{tabularx} \\ \\ \relax
16:12 & \begin{tabularx}{0.7\textwidth}{X} 「我已經聽見以色列人的怨言了。你要對他們說：『到黃昏的時候，你們要吃肉，早晨也必有食物得飽。你們就知道我是耶和華－你們的神。』」 \end{tabularx} \\ \\ \relax
16:13 & \begin{tabularx}{0.7\textwidth}{X} 到了晚上，有鵪鶉上來，遮滿營地；早晨，營地周圍有一層露水。 \end{tabularx} \\ \\ \relax
16:14 & \begin{tabularx}{0.7\textwidth}{X} 那一層露水蒸發之後，看哪，曠野的表面出現了小圓物，好像地上的薄霜一樣。 \end{tabularx} \\ \\ \relax
16:15 & \begin{tabularx}{0.7\textwidth}{X} 以色列人看見了，不知道是甚麼，就彼此說：「這是甚麼？」摩西對他們說：「這是耶和華給你們吃的食物。 \end{tabularx} \\ \\ \relax
16:16 & \begin{tabularx}{0.7\textwidth}{X} 耶和華所吩咐的是這樣：『你們每個人要按自己的食量收集，各人要為帳棚裡的人收集，按照人口數每個人一俄梅珥。』」 \end{tabularx} \\ \\ \relax
16:17 & \begin{tabularx}{0.7\textwidth}{X} 以色列人就照樣去做；有的收多，有的收少。 \end{tabularx} \\ \\ \relax
16:18 & \begin{tabularx}{0.7\textwidth}{X} 用俄梅珥量一量，多收的沒有餘，少收的也沒有缺；各人都按著自己的食量收集。 \end{tabularx} \\ \\ \relax
16:19 & \begin{tabularx}{0.7\textwidth}{X} 摩西對他們說：「任何人都不可以把所收的留到早晨。」 \end{tabularx} \\ \\ \relax
16:20 & \begin{tabularx}{0.7\textwidth}{X} 然而他們不聽從摩西，當中有人把食物留到早晨，食物就生蟲發臭了。摩西就向他們發怒。 \end{tabularx} \\ \\ \relax
16:21 & \begin{tabularx}{0.7\textwidth}{X} 他們每日早晨按著各人的食量收集；太陽一發熱，食物就融化了。 \end{tabularx} \\ \\ \relax
16:22 & \begin{tabularx}{0.7\textwidth}{X} 到第六天，他們收集了雙倍的食物，每個人二俄梅珥。會眾的官長來告訴摩西， \end{tabularx} \\ \\ \relax
16:23 & \begin{tabularx}{0.7\textwidth}{X} 摩西對他們說：「耶和華吩咐：『明天是安息日，是向耶和華守的聖安息日。你們要烤的就烤，要煮的就煮，所剩下的都留到早晨。』」 \end{tabularx} \\ \\ \relax
16:24 & \begin{tabularx}{0.7\textwidth}{X} 他們就照摩西的吩咐把剩下的留到早晨，這些食物既不發臭，裡頭也沒有生蟲。 \end{tabularx} \\ \\ \relax
16:25 & \begin{tabularx}{0.7\textwidth}{X} 摩西說：「你們今天就吃這些吧！因為今天是向耶和華守的安息日，你們在野外必找不著食物了。 \end{tabularx} \\ \\ \relax
16:26 & \begin{tabularx}{0.7\textwidth}{X} 六天可以收集，第七天是安息日，這一天甚麼也沒有了。」 \end{tabularx} \\ \\ \relax
16:27 & \begin{tabularx}{0.7\textwidth}{X} 第七天，百姓中有人出去收，甚麼也找不著。 \end{tabularx} \\ \\ \relax
16:28 & \begin{tabularx}{0.7\textwidth}{X} 耶和華對摩西說：「你們不肯遵守我的誡令和教導，要到幾時呢？ \end{tabularx} \\ \\ \relax
16:29 & \begin{tabularx}{0.7\textwidth}{X} 你們看，耶和華既然將安息日賜給你們，所以第六天他就賜給你們兩天的食物，第七天各人都要留在自己的地方，不許任何人從這裡出去。」 \end{tabularx} \\ \\ \relax
16:30 & \begin{tabularx}{0.7\textwidth}{X} 於是百姓在第七天安息了。 \end{tabularx} \\ \\ \relax
16:31 & \begin{tabularx}{0.7\textwidth}{X} 以色列家給這食物取名叫嗎哪，它的樣子像芫荽子，顏色是白的，吃起來像和蜜的薄餅。 \end{tabularx} \\ \\ \relax
16:32 & \begin{tabularx}{0.7\textwidth}{X} 摩西說：「耶和華所吩咐的是這樣：『要裝滿一俄梅珥的嗎哪留給你們的後代，使他們可以看見我領你們出埃及地的時候，在曠野所給你們吃的食物。』」 \end{tabularx} \\ \\ \relax
16:33 & \begin{tabularx}{0.7\textwidth}{X} 摩西對亞倫說：「你拿一個罐子，裝滿一俄梅珥的嗎哪，存在耶和華面前，留給你們的後代。」 \end{tabularx} \\ \\ \relax
16:34 & \begin{tabularx}{0.7\textwidth}{X} 耶和華怎麼吩咐摩西，亞倫就照樣做，把嗎哪存留作見證。 \end{tabularx} \\ \\ \relax
16:35 & \begin{tabularx}{0.7\textwidth}{X} 以色列人吃嗎哪共四十年，直到進入有人居住的地方；他們吃嗎哪，直到迦南地的邊境。 \end{tabularx} \\ \\ \relax
16:36 & \begin{tabularx}{0.7\textwidth}{X} 一俄梅珥是一伊法的十分之一。 \end{tabularx} \\ \\ \relax
17:1 & \begin{tabularx}{0.7\textwidth}{X} 以色列全會眾遵照耶和華的吩咐，從汛的曠野一段一段地往前行。他們在利非訂安營，但百姓沒有水喝。 \end{tabularx} \\ \\ \relax
17:2 & \begin{tabularx}{0.7\textwidth}{X} 百姓就與摩西爭鬧，說：「給我們水喝吧！」摩西對他們說：「你們為甚麼與我爭鬧呢？你們為甚麼試探耶和華呢？」 \end{tabularx} \\ \\ \relax
17:3 & \begin{tabularx}{0.7\textwidth}{X} 百姓在那裡口渴要喝水，就向摩西發怨言，說：「你為甚麼把我們從埃及領出來，使我們和我們的兒女，以及牲畜都渴死呢？」 \end{tabularx} \\ \\ \relax
17:4 & \begin{tabularx}{0.7\textwidth}{X} 摩西就呼求耶和華說：「我要怎樣對待這百姓呢？他們差一點就要拿石頭打死我了。」 \end{tabularx} \\ \\ \relax
17:5 & \begin{tabularx}{0.7\textwidth}{X} 耶和華對摩西說：「你帶著以色列的幾個長老，走在百姓前面，手裡拿著你先前擊打尼羅河的杖，去吧！ \end{tabularx} \\ \\ \relax
17:6 & \begin{tabularx}{0.7\textwidth}{X} 看哪，我要在何烈的磐石那裡，站在你面前。你要擊打磐石，水就會從磐石流出來，給百姓喝。」摩西就在以色列的長老眼前這樣做了。 \end{tabularx} \\ \\ \relax
17:7 & \begin{tabularx}{0.7\textwidth}{X} 他給那地方起名叫瑪撒，又叫米利巴，因為以色列人在那裡爭鬧，並且試探耶和華，說：「耶和華是否在我們中間呢？」 \end{tabularx} \\ \\ \relax
17:8 & \begin{tabularx}{0.7\textwidth}{X} 那時，亞瑪力來到利非訂，和以色列爭戰。 \end{tabularx} \\ \\ \relax
17:9 & \begin{tabularx}{0.7\textwidth}{X} 摩西對約書亞說：「你為我們選出人來，出去和亞瑪力爭戰。明天我要站在山頂上，手裡拿著神的杖。」 \end{tabularx} \\ \\ \relax
17:10 & \begin{tabularx}{0.7\textwidth}{X} 於是，約書亞照著摩西對他所說的話去做，和亞瑪力爭戰。摩西、亞倫和戶珥都上了山頂。 \end{tabularx} \\ \\ \relax
17:11 & \begin{tabularx}{0.7\textwidth}{X} 摩西何時舉手，以色列就得勝；何時垂手，亞瑪力就得勝。 \end{tabularx} \\ \\ \relax
17:12 & \begin{tabularx}{0.7\textwidth}{X} 但摩西的雙手沉重，他們就搬一塊石頭來放在他下面，他就坐在上面。亞倫與戶珥扶著他的手，一個在這邊，一個在那邊，他的手就穩住，直到日落。 \end{tabularx} \\ \\ \relax
17:13 & \begin{tabularx}{0.7\textwidth}{X} 約書亞用刀打敗了亞瑪力和他的百姓。 \end{tabularx} \\ \\ \relax
17:14 & \begin{tabularx}{0.7\textwidth}{X} 耶和華對摩西說：「你要把這事記錄在書上作紀念，又念給約書亞聽：我要把亞瑪力的名字從天下全然塗去。」 \end{tabularx} \\ \\ \relax
17:15 & \begin{tabularx}{0.7\textwidth}{X} 摩西築了一座壇，起名叫「耶和華尼西」。 \end{tabularx} \\ \\ \relax
17:16 & \begin{tabularx}{0.7\textwidth}{X} 他說：「我指著耶和華的寶座發誓，耶和華必世世代代和亞瑪力爭戰。」 \end{tabularx} \\ \\
[1ex]
\hline
\hline
\end{longtable}


\newpage

\section{第65屆~港九培靈會~周永健~第七講}
\label{sec:989}
\textbf{教導我們──我從天上和你們說話了}
\newline
\newline
連結: \href{https://www.hkbibleconference.org/session-message/view/989}{\texttt{https://www.hkbibleconference.org/session-message/view/989}}
\newline
\newline
\hyperref[sec:988]{< < < PREV SERMON < < <}
~
\hyperref[sec:toc]{[返目錄]}
~
\hyperref[sec:990]{> > > NEXT SERMON > > >}
\newline
\newline
出埃及記 25:1-25:40
\newline
\begin{longtable}{cl}
\hline
\hline
章節 & 經文 (和合本修訂版)\\
\hline
25:1 & \begin{tabularx}{0.7\textwidth}{X} 耶和華吩咐摩西說： \end{tabularx} \\ \\ \relax
25:2 & \begin{tabularx}{0.7\textwidth}{X} 「你要吩咐以色列人獻禮物給我。凡甘心樂意獻給我的禮物，你們都可以收下。 \end{tabularx} \\ \\ \relax
25:3 & \begin{tabularx}{0.7\textwidth}{X} 要從他們收的禮物是：金、銀、銅， \end{tabularx} \\ \\ \relax
25:4 & \begin{tabularx}{0.7\textwidth}{X} 藍色、紫色、朱紅色紗，細麻，山羊毛， \end{tabularx} \\ \\ \relax
25:5 & \begin{tabularx}{0.7\textwidth}{X} 染紅的公羊皮、精美的皮料，金合歡木， \end{tabularx} \\ \\ \relax
25:6 & \begin{tabularx}{0.7\textwidth}{X} 點燈的油，做膏油的香料、做香的香料， \end{tabularx} \\ \\ \relax
25:7 & \begin{tabularx}{0.7\textwidth}{X} 紅瑪瑙與寶石，可以鑲嵌在以弗得和胸袋上。 \end{tabularx} \\ \\ \relax
25:8 & \begin{tabularx}{0.7\textwidth}{X} 他們要為我造聖所，使我住在他們中間。 \end{tabularx} \\ \\ \relax
25:9 & \begin{tabularx}{0.7\textwidth}{X} 你們要按照我指示你的，帳幕和其中一切器具的樣式，照樣去做。」 \end{tabularx} \\ \\ \relax
25:10 & \begin{tabularx}{0.7\textwidth}{X} 「他們要用金合歡木做一個櫃子，長二肘半，寬一肘半，高一肘半。 \end{tabularx} \\ \\ \relax
25:11 & \begin{tabularx}{0.7\textwidth}{X} 你要把它裡裡外外包上純金，四圍要鑲上金邊。 \end{tabularx} \\ \\ \relax
25:12 & \begin{tabularx}{0.7\textwidth}{X} 要鑄造四個金環，安在櫃子的四腳上；這邊兩個環，那邊兩個環。 \end{tabularx} \\ \\ \relax
25:13 & \begin{tabularx}{0.7\textwidth}{X} 要用金合歡木做兩根槓，包上金子。 \end{tabularx} \\ \\ \relax
25:14 & \begin{tabularx}{0.7\textwidth}{X} 要把槓穿過櫃旁的環，以便抬櫃。 \end{tabularx} \\ \\ \relax
25:15 & \begin{tabularx}{0.7\textwidth}{X} 這槓要留在櫃的環內，不可抽出來。 \end{tabularx} \\ \\ \relax
25:16 & \begin{tabularx}{0.7\textwidth}{X} 要把我所要賜給你的法版放在櫃裡。 \end{tabularx} \\ \\ \relax
25:17 & \begin{tabularx}{0.7\textwidth}{X} 要用純金做一個櫃蓋，長二肘半，寬一肘半。 \end{tabularx} \\ \\ \relax
25:18 & \begin{tabularx}{0.7\textwidth}{X} 要造兩個用金子錘出的基路伯，從櫃蓋的兩端錘出它們。 \end{tabularx} \\ \\ \relax
25:19 & \begin{tabularx}{0.7\textwidth}{X} 這端錘出一個基路伯，那端錘出一個基路伯；從櫃蓋的兩端錘出兩個基路伯。 \end{tabularx} \\ \\ \relax
25:20 & \begin{tabularx}{0.7\textwidth}{X} 二基路伯的翅膀要向上張開，用翅膀遮住櫃蓋，臉要彼此相對；基路伯的臉要朝向櫃蓋。 \end{tabularx} \\ \\ \relax
25:21 & \begin{tabularx}{0.7\textwidth}{X} 要把櫃蓋安在櫃的上邊，又要把我所要賜給你的法版放在櫃裡。 \end{tabularx} \\ \\ \relax
25:22 & \begin{tabularx}{0.7\textwidth}{X} 我要在那裡與你相會，並要從法版之櫃的櫃蓋上，兩個基路伯的中間，將我要吩咐以色列人的一切事告訴你。」 \end{tabularx} \\ \\ \relax
25:23 & \begin{tabularx}{0.7\textwidth}{X} 「你要用金合歡木做一張供桌，長二肘，寬一肘，高一肘半， \end{tabularx} \\ \\ \relax
25:24 & \begin{tabularx}{0.7\textwidth}{X} 把它包上純金，四圍鑲上金邊。 \end{tabularx} \\ \\ \relax
25:25 & \begin{tabularx}{0.7\textwidth}{X} 供桌的四圍各做一掌寬的邊緣，邊緣周圍要鑲上金邊。 \end{tabularx} \\ \\ \relax
25:26 & \begin{tabularx}{0.7\textwidth}{X} 要為供桌做四個金環，把環安在四個桌腳的四角上。 \end{tabularx} \\ \\ \relax
25:27 & \begin{tabularx}{0.7\textwidth}{X} 環要靠近邊緣，以便穿槓抬供桌。 \end{tabularx} \\ \\ \relax
25:28 & \begin{tabularx}{0.7\textwidth}{X} 要用金合歡木做兩根槓，包上金子，用來抬供桌。 \end{tabularx} \\ \\ \relax
25:29 & \begin{tabularx}{0.7\textwidth}{X} 要用純金做桌上的盤、碟，以及澆酒祭的壺和杯。 \end{tabularx} \\ \\ \relax
25:30 & \begin{tabularx}{0.7\textwidth}{X} 要把供餅擺在桌上，常在我面前。」 \end{tabularx} \\ \\ \relax
25:31 & \begin{tabularx}{0.7\textwidth}{X} 「你要造一座用純金錘出的燈臺。燈臺的座、幹、杯、花萼和花瓣，都要和燈臺接連一塊。 \end{tabularx} \\ \\ \relax
25:32 & \begin{tabularx}{0.7\textwidth}{X} 燈臺兩旁要伸出六根枝子：這邊三根，那邊三根。 \end{tabularx} \\ \\ \relax
25:33 & \begin{tabularx}{0.7\textwidth}{X} 這邊枝子上有三個杯，形狀像杏花，有花萼有花瓣；那邊枝子上也有三個杯，形狀像杏花，有花萼有花瓣。從燈臺伸出來的六根枝子都是如此。 \end{tabularx} \\ \\ \relax
25:34 & \begin{tabularx}{0.7\textwidth}{X} 燈臺本身要有四個杯，形狀像杏花，有花萼有花瓣。 \end{tabularx} \\ \\ \relax
25:35 & \begin{tabularx}{0.7\textwidth}{X} 燈臺的第一對枝子下面有花萼，燈臺的第二對枝子下面有花萼，燈臺的第三對枝子下面也有花萼；燈臺伸出的六根枝子都是如此。 \end{tabularx} \\ \\ \relax
25:36 & \begin{tabularx}{0.7\textwidth}{X} 花萼和枝子都要和燈臺接連一塊，全是從一塊純金錘出來的。 \end{tabularx} \\ \\ \relax
25:37 & \begin{tabularx}{0.7\textwidth}{X} 要做燈臺的七盞燈，燈要點燃，照亮前面。 \end{tabularx} \\ \\ \relax
25:38 & \begin{tabularx}{0.7\textwidth}{X} 要用純金做燈剪和燈盤。 \end{tabularx} \\ \\ \relax
25:39 & \begin{tabularx}{0.7\textwidth}{X} 做燈臺和這一切的器具要用一他連得純金。 \end{tabularx} \\ \\ \relax
25:40 & \begin{tabularx}{0.7\textwidth}{X} 要謹慎，照著在山上指示你的樣式去做。」 \end{tabularx} \\ \\
[1ex]
\hline
\hline
\end{longtable}


\newpage

\section{第65屆~港九培靈會~周永健~第八講}
\label{sec:990}
\textbf{神與我們同在──使我可以住在他們中間}
\newline
\newline
連結: \href{https://www.hkbibleconference.org/session-message/view/990}{\texttt{https://www.hkbibleconference.org/session-message/view/990}}
\newline
\newline
\hyperref[sec:989]{< < < PREV SERMON < < <}
~
\hyperref[sec:toc]{[返目錄]}
~
\hyperref[sec:991]{> > > NEXT SERMON > > >}
\newline
\newline
出埃及記 25:1-31:18
\newline
\begin{longtable}{cl}
\hline
\hline
章節 & 經文 (和合本修訂版)\\
\hline
25:1 & \begin{tabularx}{0.7\textwidth}{X} 耶和華吩咐摩西說： \end{tabularx} \\ \\ \relax
25:2 & \begin{tabularx}{0.7\textwidth}{X} 「你要吩咐以色列人獻禮物給我。凡甘心樂意獻給我的禮物，你們都可以收下。 \end{tabularx} \\ \\ \relax
25:3 & \begin{tabularx}{0.7\textwidth}{X} 要從他們收的禮物是：金、銀、銅， \end{tabularx} \\ \\ \relax
25:4 & \begin{tabularx}{0.7\textwidth}{X} 藍色、紫色、朱紅色紗，細麻，山羊毛， \end{tabularx} \\ \\ \relax
25:5 & \begin{tabularx}{0.7\textwidth}{X} 染紅的公羊皮、精美的皮料，金合歡木， \end{tabularx} \\ \\ \relax
25:6 & \begin{tabularx}{0.7\textwidth}{X} 點燈的油，做膏油的香料、做香的香料， \end{tabularx} \\ \\ \relax
25:7 & \begin{tabularx}{0.7\textwidth}{X} 紅瑪瑙與寶石，可以鑲嵌在以弗得和胸袋上。 \end{tabularx} \\ \\ \relax
25:8 & \begin{tabularx}{0.7\textwidth}{X} 他們要為我造聖所，使我住在他們中間。 \end{tabularx} \\ \\ \relax
25:9 & \begin{tabularx}{0.7\textwidth}{X} 你們要按照我指示你的，帳幕和其中一切器具的樣式，照樣去做。」 \end{tabularx} \\ \\ \relax
25:10 & \begin{tabularx}{0.7\textwidth}{X} 「他們要用金合歡木做一個櫃子，長二肘半，寬一肘半，高一肘半。 \end{tabularx} \\ \\ \relax
25:11 & \begin{tabularx}{0.7\textwidth}{X} 你要把它裡裡外外包上純金，四圍要鑲上金邊。 \end{tabularx} \\ \\ \relax
25:12 & \begin{tabularx}{0.7\textwidth}{X} 要鑄造四個金環，安在櫃子的四腳上；這邊兩個環，那邊兩個環。 \end{tabularx} \\ \\ \relax
25:13 & \begin{tabularx}{0.7\textwidth}{X} 要用金合歡木做兩根槓，包上金子。 \end{tabularx} \\ \\ \relax
25:14 & \begin{tabularx}{0.7\textwidth}{X} 要把槓穿過櫃旁的環，以便抬櫃。 \end{tabularx} \\ \\ \relax
25:15 & \begin{tabularx}{0.7\textwidth}{X} 這槓要留在櫃的環內，不可抽出來。 \end{tabularx} \\ \\ \relax
25:16 & \begin{tabularx}{0.7\textwidth}{X} 要把我所要賜給你的法版放在櫃裡。 \end{tabularx} \\ \\ \relax
25:17 & \begin{tabularx}{0.7\textwidth}{X} 要用純金做一個櫃蓋，長二肘半，寬一肘半。 \end{tabularx} \\ \\ \relax
25:18 & \begin{tabularx}{0.7\textwidth}{X} 要造兩個用金子錘出的基路伯，從櫃蓋的兩端錘出它們。 \end{tabularx} \\ \\ \relax
25:19 & \begin{tabularx}{0.7\textwidth}{X} 這端錘出一個基路伯，那端錘出一個基路伯；從櫃蓋的兩端錘出兩個基路伯。 \end{tabularx} \\ \\ \relax
25:20 & \begin{tabularx}{0.7\textwidth}{X} 二基路伯的翅膀要向上張開，用翅膀遮住櫃蓋，臉要彼此相對；基路伯的臉要朝向櫃蓋。 \end{tabularx} \\ \\ \relax
25:21 & \begin{tabularx}{0.7\textwidth}{X} 要把櫃蓋安在櫃的上邊，又要把我所要賜給你的法版放在櫃裡。 \end{tabularx} \\ \\ \relax
25:22 & \begin{tabularx}{0.7\textwidth}{X} 我要在那裡與你相會，並要從法版之櫃的櫃蓋上，兩個基路伯的中間，將我要吩咐以色列人的一切事告訴你。」 \end{tabularx} \\ \\ \relax
25:23 & \begin{tabularx}{0.7\textwidth}{X} 「你要用金合歡木做一張供桌，長二肘，寬一肘，高一肘半， \end{tabularx} \\ \\ \relax
25:24 & \begin{tabularx}{0.7\textwidth}{X} 把它包上純金，四圍鑲上金邊。 \end{tabularx} \\ \\ \relax
25:25 & \begin{tabularx}{0.7\textwidth}{X} 供桌的四圍各做一掌寬的邊緣，邊緣周圍要鑲上金邊。 \end{tabularx} \\ \\ \relax
25:26 & \begin{tabularx}{0.7\textwidth}{X} 要為供桌做四個金環，把環安在四個桌腳的四角上。 \end{tabularx} \\ \\ \relax
25:27 & \begin{tabularx}{0.7\textwidth}{X} 環要靠近邊緣，以便穿槓抬供桌。 \end{tabularx} \\ \\ \relax
25:28 & \begin{tabularx}{0.7\textwidth}{X} 要用金合歡木做兩根槓，包上金子，用來抬供桌。 \end{tabularx} \\ \\ \relax
25:29 & \begin{tabularx}{0.7\textwidth}{X} 要用純金做桌上的盤、碟，以及澆酒祭的壺和杯。 \end{tabularx} \\ \\ \relax
25:30 & \begin{tabularx}{0.7\textwidth}{X} 要把供餅擺在桌上，常在我面前。」 \end{tabularx} \\ \\ \relax
25:31 & \begin{tabularx}{0.7\textwidth}{X} 「你要造一座用純金錘出的燈臺。燈臺的座、幹、杯、花萼和花瓣，都要和燈臺接連一塊。 \end{tabularx} \\ \\ \relax
25:32 & \begin{tabularx}{0.7\textwidth}{X} 燈臺兩旁要伸出六根枝子：這邊三根，那邊三根。 \end{tabularx} \\ \\ \relax
25:33 & \begin{tabularx}{0.7\textwidth}{X} 這邊枝子上有三個杯，形狀像杏花，有花萼有花瓣；那邊枝子上也有三個杯，形狀像杏花，有花萼有花瓣。從燈臺伸出來的六根枝子都是如此。 \end{tabularx} \\ \\ \relax
25:34 & \begin{tabularx}{0.7\textwidth}{X} 燈臺本身要有四個杯，形狀像杏花，有花萼有花瓣。 \end{tabularx} \\ \\ \relax
25:35 & \begin{tabularx}{0.7\textwidth}{X} 燈臺的第一對枝子下面有花萼，燈臺的第二對枝子下面有花萼，燈臺的第三對枝子下面也有花萼；燈臺伸出的六根枝子都是如此。 \end{tabularx} \\ \\ \relax
25:36 & \begin{tabularx}{0.7\textwidth}{X} 花萼和枝子都要和燈臺接連一塊，全是從一塊純金錘出來的。 \end{tabularx} \\ \\ \relax
25:37 & \begin{tabularx}{0.7\textwidth}{X} 要做燈臺的七盞燈，燈要點燃，照亮前面。 \end{tabularx} \\ \\ \relax
25:38 & \begin{tabularx}{0.7\textwidth}{X} 要用純金做燈剪和燈盤。 \end{tabularx} \\ \\ \relax
25:39 & \begin{tabularx}{0.7\textwidth}{X} 做燈臺和這一切的器具要用一他連得純金。 \end{tabularx} \\ \\ \relax
25:40 & \begin{tabularx}{0.7\textwidth}{X} 要謹慎，照著在山上指示你的樣式去做。」 \end{tabularx} \\ \\ \relax
26:1 & \begin{tabularx}{0.7\textwidth}{X} 「你要用十幅幔子做帳幕。這些幔子要用搓的細麻和藍色、紫色、朱紅色紗織成，並且以刺繡的手藝繡上基路伯。 \end{tabularx} \\ \\ \relax
26:2 & \begin{tabularx}{0.7\textwidth}{X} 每幅幔子要長二十八肘，每幅幔子寬四肘，全部的幔子尺寸都要一樣。 \end{tabularx} \\ \\ \relax
26:3 & \begin{tabularx}{0.7\textwidth}{X} 這五幅幔子要彼此相連；那五幅幔子也彼此相連。 \end{tabularx} \\ \\ \relax
26:4 & \begin{tabularx}{0.7\textwidth}{X} 在這一組相連幔子的末幅邊上要縫藍色的鈕環；在另一組相連幔子的末幅邊上也要照樣做。 \end{tabularx} \\ \\ \relax
26:5 & \begin{tabularx}{0.7\textwidth}{X} 這幅幔子上要縫五十個鈕環，另一組相連幔子的末幅上也縫五十個鈕環，環環相對。 \end{tabularx} \\ \\ \relax
26:6 & \begin{tabularx}{0.7\textwidth}{X} 要做五十個金鉤，用鉤子使幔子彼此相連，成為一個帳幕。 \end{tabularx} \\ \\ \relax
26:7 & \begin{tabularx}{0.7\textwidth}{X} 「你要用山羊毛織十一幅幔子來作帳幕的罩棚。 \end{tabularx} \\ \\ \relax
26:8 & \begin{tabularx}{0.7\textwidth}{X} 每幅幔子要長三十肘，每幅幔子寬四肘；十一幅幔子的尺寸都要一樣。 \end{tabularx} \\ \\ \relax
26:9 & \begin{tabularx}{0.7\textwidth}{X} 要把五幅幔子連成一幅，又把六幅幔子連成一幅，這第六幅幔子要在罩棚的前面摺上去。 \end{tabularx} \\ \\ \relax
26:10 & \begin{tabularx}{0.7\textwidth}{X} 在這一組相連幔子的末幅邊上要縫五十個鈕環；在另一組相連幔子的末幅邊上也縫五十個鈕環。 \end{tabularx} \\ \\ \relax
26:11 & \begin{tabularx}{0.7\textwidth}{X} 要做五十個銅鉤，鉤在鈕環中，使罩棚相連成為一個。 \end{tabularx} \\ \\ \relax
26:12 & \begin{tabularx}{0.7\textwidth}{X} 罩棚幔子餘下垂著的，那餘下的半幅要垂在帳幕的背面。 \end{tabularx} \\ \\ \relax
26:13 & \begin{tabularx}{0.7\textwidth}{X} 罩棚的幔子兩旁所餘下的，這邊一肘，那邊一肘，要垂在帳幕的兩邊，蓋住帳幕。 \end{tabularx} \\ \\ \relax
26:14 & \begin{tabularx}{0.7\textwidth}{X} 要用染紅的公羊皮做罩棚的蓋，再用精美皮料做外層的蓋。 \end{tabularx} \\ \\ \relax
26:15 & \begin{tabularx}{0.7\textwidth}{X} 「你要用金合歡木做豎立帳幕的木板， \end{tabularx} \\ \\ \relax
26:16 & \begin{tabularx}{0.7\textwidth}{X} 木板要長十肘，每塊板寬一肘半， \end{tabularx} \\ \\ \relax
26:17 & \begin{tabularx}{0.7\textwidth}{X} 每塊板有兩個榫頭可以彼此銜接。帳幕一切的板都要這樣做。 \end{tabularx} \\ \\ \relax
26:18 & \begin{tabularx}{0.7\textwidth}{X} 你要做帳幕的木板：南面，就是面向南方的那一邊，要做二十塊板。 \end{tabularx} \\ \\ \relax
26:19 & \begin{tabularx}{0.7\textwidth}{X} 在這二十塊板底下要做四十個帶卯眼的銀座；兩個卯眼接連這塊板上的兩個榫頭，另外兩個卯眼接連那塊板上的兩個榫頭。 \end{tabularx} \\ \\ \relax
26:20 & \begin{tabularx}{0.7\textwidth}{X} 帳幕的第二邊，就是北面，也要做二十塊板， \end{tabularx} \\ \\ \relax
26:21 & \begin{tabularx}{0.7\textwidth}{X} 和四十個帶卯眼的銀座；這塊板底下有兩個卯眼，那塊板底下也有兩個卯眼。 \end{tabularx} \\ \\ \relax
26:22 & \begin{tabularx}{0.7\textwidth}{X} 帳幕的後面，就是西面，要做六塊板。 \end{tabularx} \\ \\ \relax
26:23 & \begin{tabularx}{0.7\textwidth}{X} 帳幕後面的角落要做兩塊板。 \end{tabularx} \\ \\ \relax
26:24 & \begin{tabularx}{0.7\textwidth}{X} 下端的板是成雙的，上端要連在一起，直到頂端的第一個環子；兩塊板都要這樣，做成兩個角落。 \end{tabularx} \\ \\ \relax
26:25 & \begin{tabularx}{0.7\textwidth}{X} 一共有八塊板和十六個帶卯眼的銀座；這塊板底下有兩個卯眼，那塊板底下也有兩個卯眼。 \end{tabularx} \\ \\ \relax
26:26 & \begin{tabularx}{0.7\textwidth}{X} 「你要用金合歡木做橫木：為帳幕這面的板做五根橫木， \end{tabularx} \\ \\ \relax
26:27 & \begin{tabularx}{0.7\textwidth}{X} 為帳幕那面的板做五根橫木，又為帳幕後面，就是朝西的板做五根橫木。 \end{tabularx} \\ \\ \relax
26:28 & \begin{tabularx}{0.7\textwidth}{X} 板腰間的橫木，要從一頭通到另一頭。 \end{tabularx} \\ \\ \relax
26:29 & \begin{tabularx}{0.7\textwidth}{X} 板要包上金子，又要做板上的金環來套橫木；橫木也要包上金子。 \end{tabularx} \\ \\ \relax
26:30 & \begin{tabularx}{0.7\textwidth}{X} 要照著在山上所指示你的樣式，把帳幕豎立起來。 \end{tabularx} \\ \\ \relax
26:31 & \begin{tabularx}{0.7\textwidth}{X} 「你要用藍色、紫色、朱紅色紗，和搓的細麻織幔子，以刺繡的手藝繡上基路伯。 \end{tabularx} \\ \\ \relax
26:32 & \begin{tabularx}{0.7\textwidth}{X} 要把幔子掛在四根包金的金合歡木柱子上，柱子有金鉤，並且安在四個帶卯眼的銀座上。 \end{tabularx} \\ \\ \relax
26:33 & \begin{tabularx}{0.7\textwidth}{X} 要把幔子垂掛在鉤子上，把法櫃抬進幔子內；這幔子要將聖所和至聖所隔開。 \end{tabularx} \\ \\ \relax
26:34 & \begin{tabularx}{0.7\textwidth}{X} 又要把櫃蓋安在至聖所內的法櫃上， \end{tabularx} \\ \\ \relax
26:35 & \begin{tabularx}{0.7\textwidth}{X} 把供桌安在幔子的外面，供桌在北面，燈臺在帳幕的南面，和供桌相對。 \end{tabularx} \\ \\ \relax
26:36 & \begin{tabularx}{0.7\textwidth}{X} 「你要用藍色、紫色、朱紅色紗，和搓的細麻，以刺繡的手藝為帳幕織門簾。 \end{tabularx} \\ \\ \relax
26:37 & \begin{tabularx}{0.7\textwidth}{X} 要用金合歡木為簾子做五根柱子，包上金子。柱子有金鉤，又為柱子鑄造五個帶卯眼的銅座。」 \end{tabularx} \\ \\ \relax
27:1 & \begin{tabularx}{0.7\textwidth}{X} 「你要用金合歡木做祭壇，長五肘，寬五肘，這壇是正方形的，高三肘。 \end{tabularx} \\ \\ \relax
27:2 & \begin{tabularx}{0.7\textwidth}{X} 要在壇的四角做四個翹角，與壇接連一塊；要把壇包上銅。 \end{tabularx} \\ \\ \relax
27:3 & \begin{tabularx}{0.7\textwidth}{X} 要做桶子來盛壇上的灰，又要做鏟子、盤子、肉叉和火盆；壇上一切的器具都要用銅做。 \end{tabularx} \\ \\ \relax
27:4 & \begin{tabularx}{0.7\textwidth}{X} 要為壇做一個銅網，在網的四角做四個銅環， \end{tabularx} \\ \\ \relax
27:5 & \begin{tabularx}{0.7\textwidth}{X} 把網安在壇四圍的邊的下面，使網垂到壇的半腰。 \end{tabularx} \\ \\ \relax
27:6 & \begin{tabularx}{0.7\textwidth}{X} 又要用金合歡木為壇做槓，包上銅。 \end{tabularx} \\ \\ \relax
27:7 & \begin{tabularx}{0.7\textwidth}{X} 這槓要穿過壇兩旁的環子，用來抬壇。 \end{tabularx} \\ \\ \relax
27:8 & \begin{tabularx}{0.7\textwidth}{X} 要用板做壇，壇的中心是空的，都照著在山上所指示你的樣式做。」 \end{tabularx} \\ \\ \relax
27:9 & \begin{tabularx}{0.7\textwidth}{X} 「你要做帳幕的院子。南面，就是面向南方的那一邊，要用搓的細麻做院子的帷幔，長一百肘， \end{tabularx} \\ \\ \relax
27:10 & \begin{tabularx}{0.7\textwidth}{X} 院子要有二十根柱子，二十個帶卯眼的銅座。要用銀做柱子的鉤和箍。 \end{tabularx} \\ \\ \relax
27:11 & \begin{tabularx}{0.7\textwidth}{X} 北面的長度也一樣，帷幔長一百肘，要有二十根柱子，二十個帶卯眼的銅座。要用銀做柱子的鉤和箍。 \end{tabularx} \\ \\ \relax
27:12 & \begin{tabularx}{0.7\textwidth}{X} 院子的西面有帷幔，寬五十肘，帷幔要有十根柱子，十個帶卯眼的座。 \end{tabularx} \\ \\ \relax
27:13 & \begin{tabularx}{0.7\textwidth}{X} 院子的東面，就是面向東方的那一邊，寬五十肘。 \end{tabularx} \\ \\ \relax
27:14 & \begin{tabularx}{0.7\textwidth}{X} 一邊的帷幔有十五肘，要有三根柱子，三個帶卯眼的座。 \end{tabularx} \\ \\ \relax
27:15 & \begin{tabularx}{0.7\textwidth}{X} 另一邊的帷幔也有十五肘，要有三根柱子，三個帶卯眼的座。 \end{tabularx} \\ \\ \relax
27:16 & \begin{tabularx}{0.7\textwidth}{X} 院子的門要有二十肘長的簾子，用藍色、紫色、朱紅色紗，和搓的細麻，以刺繡的手藝織成；要有四根柱子，四個帶卯眼的座。 \end{tabularx} \\ \\ \relax
27:17 & \begin{tabularx}{0.7\textwidth}{X} 院子四圍一切的柱子都要用銀子箍著，要用銀做柱子的鉤子，用銅做帶卯眼的座。 \end{tabularx} \\ \\ \relax
27:18 & \begin{tabularx}{0.7\textwidth}{X} 院子要長一百肘，寬五十肘，高五肘。要用搓的細麻做帷幔，用銅做帶卯眼的座。 \end{tabularx} \\ \\ \relax
27:19 & \begin{tabularx}{0.7\textwidth}{X} 帳幕中各樣用途的器具，以及帳幕一切的橛子和院子裡一切的橛子，都要用銅做。」 \end{tabularx} \\ \\ \relax
27:20 & \begin{tabularx}{0.7\textwidth}{X} 「你要吩咐以色列人，把搗成的純橄欖油拿來給你，用以點燈，使燈經常點著； \end{tabularx} \\ \\ \relax
27:21 & \begin{tabularx}{0.7\textwidth}{X} 在會幕中法櫃前的幔子外，亞倫和他的兒子要從晚上到早晨，在耶和華面前照管這燈。這要成為以色列人世世代代永遠的定例。」 \end{tabularx} \\ \\ \relax
28:1 & \begin{tabularx}{0.7\textwidth}{X} 「你要從以色列人中，叫你的哥哥亞倫和他的兒子拿答、亞比戶、以利亞撒、以他瑪一同親近你，作事奉我的祭司。 \end{tabularx} \\ \\ \relax
28:2 & \begin{tabularx}{0.7\textwidth}{X} 你要為你哥哥亞倫做聖衣，以示尊嚴和華美。 \end{tabularx} \\ \\ \relax
28:3 & \begin{tabularx}{0.7\textwidth}{X} 要吩咐一切心中有智慧的，就是我用智慧的靈所充滿的人，為亞倫做衣服，使他分別為聖，作事奉我的祭司。 \end{tabularx} \\ \\ \relax
28:4 & \begin{tabularx}{0.7\textwidth}{X} 所要做的是胸袋、以弗得、外袍、織成的內袍、禮冠和腰帶。他們要為你哥哥亞倫和他的兒子做聖衣，使他們作祭司事奉我。 \end{tabularx} \\ \\ \relax
28:5 & \begin{tabularx}{0.7\textwidth}{X} 要用金色、藍色、紫色、朱紅色紗，和細麻去縫製。」 \end{tabularx} \\ \\ \relax
28:6 & \begin{tabularx}{0.7\textwidth}{X} 「他們要用金色、藍色、紫色、朱紅色紗，和搓的細麻，以刺繡的手藝做以弗得。 \end{tabularx} \\ \\ \relax
28:7 & \begin{tabularx}{0.7\textwidth}{X} 以弗得當有兩條肩帶，接上兩端，使它相連。 \end{tabularx} \\ \\ \relax
28:8 & \begin{tabularx}{0.7\textwidth}{X} 以弗得的精緻帶子，要以一樣的手藝，用金色、藍色、紫色、朱紅色紗，和搓的細麻縫製，與以弗得接連在一起。 \end{tabularx} \\ \\ \relax
28:9 & \begin{tabularx}{0.7\textwidth}{X} 要取兩塊紅瑪瑙，在上面刻以色列兒子的名字： \end{tabularx} \\ \\ \relax
28:10 & \begin{tabularx}{0.7\textwidth}{X} 六個名字在一塊寶石上，六個名字在另一塊寶石上，都按照他們出生的次序。 \end{tabularx} \\ \\ \relax
28:11 & \begin{tabularx}{0.7\textwidth}{X} 要以雕刻寶石的手藝，如同刻印章，把以色列兒子的名字刻在這兩塊寶石上，並把寶石鑲在金槽裡。 \end{tabularx} \\ \\ \relax
28:12 & \begin{tabularx}{0.7\textwidth}{X} 要把這兩塊寶石安在以弗得的兩條肩帶上，為以色列人作紀念石。亞倫要在耶和華面前把他們的名字帶在兩肩上，作為紀念。 \end{tabularx} \\ \\ \relax
28:13 & \begin{tabularx}{0.7\textwidth}{X} 要用金子做兩個槽， \end{tabularx} \\ \\ \relax
28:14 & \begin{tabularx}{0.7\textwidth}{X} 再用純金打兩條鏈子，像編成的繩子一樣，把這編成的金鏈扣在槽上。」 \end{tabularx} \\ \\ \relax
28:15 & \begin{tabularx}{0.7\textwidth}{X} 「你要以刺繡的手藝做一個決斷的胸袋，和做以弗得的方法一樣，用金色、藍色、紫色、朱紅色紗，和搓的細麻縫製。 \end{tabularx} \\ \\ \relax
28:16 & \begin{tabularx}{0.7\textwidth}{X} 胸袋是正方形的，疊成兩層，長一虎口，寬一虎口。 \end{tabularx} \\ \\ \relax
28:17 & \begin{tabularx}{0.7\textwidth}{X} 要在上面鑲四行寶石：第一行是紅寶石、紅璧璽、紅玉； \end{tabularx} \\ \\ \relax
28:18 & \begin{tabularx}{0.7\textwidth}{X} 第二行是綠寶石、藍寶石、金剛石； \end{tabularx} \\ \\ \relax
28:19 & \begin{tabularx}{0.7\textwidth}{X} 第三行是紫瑪瑙、白瑪瑙、紫晶； \end{tabularx} \\ \\ \relax
28:20 & \begin{tabularx}{0.7\textwidth}{X} 第四行是水蒼玉、紅瑪瑙、碧玉。這些都要鑲在金槽中。 \end{tabularx} \\ \\ \relax
28:21 & \begin{tabularx}{0.7\textwidth}{X} 這些寶石要有以色列十二個兒子的名字，如同刻印章，每一顆有自己的名字，代表十二個支派。 \end{tabularx} \\ \\ \relax
28:22 & \begin{tabularx}{0.7\textwidth}{X} 要在胸袋上用純金打鏈子，像編成的繩子一樣。 \end{tabularx} \\ \\ \relax
28:23 & \begin{tabularx}{0.7\textwidth}{X} 要為胸袋做兩個金環，把這兩個環安在胸袋的兩端。 \end{tabularx} \\ \\ \relax
28:24 & \begin{tabularx}{0.7\textwidth}{X} 要把那兩條編成的金鏈繫在胸袋兩端的兩個環上。 \end{tabularx} \\ \\ \relax
28:25 & \begin{tabularx}{0.7\textwidth}{X} 又要把鏈子的另外兩端扣在兩個槽上，安在以弗得前面的肩帶上。 \end{tabularx} \\ \\ \relax
28:26 & \begin{tabularx}{0.7\textwidth}{X} 要做兩個金環，安在胸袋的兩端，在以弗得裡面的邊上。 \end{tabularx} \\ \\ \relax
28:27 & \begin{tabularx}{0.7\textwidth}{X} 再做兩個金環，安在以弗得前面兩條肩帶的下邊，靠近接縫處，在以弗得精緻帶子的上面。 \end{tabularx} \\ \\ \relax
28:28 & \begin{tabularx}{0.7\textwidth}{X} 要用藍色的帶子把胸袋的環與以弗得的環繫住，使胸袋綁在以弗得的精緻帶子上，不致鬆脫。 \end{tabularx} \\ \\ \relax
28:29 & \begin{tabularx}{0.7\textwidth}{X} 亞倫進聖所的時候，要把刻著以色列兒子名字的決斷胸袋帶著，放在心上，在耶和華面前常作紀念。 \end{tabularx} \\ \\ \relax
28:30 & \begin{tabularx}{0.7\textwidth}{X} 又要將烏陵和土明放在決斷胸袋裡；亞倫進到耶和華面前的時候，要放在心上。這樣，亞倫在耶和華面前要把以色列人的決斷胸袋常常帶著，放在心上。」 \end{tabularx} \\ \\ \relax
28:31 & \begin{tabularx}{0.7\textwidth}{X} 「你要做以弗得的外袍，顏色全是藍的。 \end{tabularx} \\ \\ \relax
28:32 & \begin{tabularx}{0.7\textwidth}{X} 袍上方的中間要留一個領口，領口周圍的領邊要以手藝編織而成，好像鎧甲的領口，免得破裂。 \end{tabularx} \\ \\ \relax
28:33 & \begin{tabularx}{0.7\textwidth}{X} 袍子下襬，就是下襬的周圍要用藍色、紫色、朱紅色紗做石榴，周圍的石榴中間要有金鈴鐺： \end{tabularx} \\ \\ \relax
28:34 & \begin{tabularx}{0.7\textwidth}{X} 一個金鈴鐺一個石榴，一個金鈴鐺一個石榴，在袍子下襬的周圍。 \end{tabularx} \\ \\ \relax
28:35 & \begin{tabularx}{0.7\textwidth}{X} 亞倫供職的時候要穿這袍。他進入聖所到耶和華面前，以及出來的時候，袍上的鈴聲必被聽見，使他不至於死。 \end{tabularx} \\ \\ \relax
28:36 & \begin{tabularx}{0.7\textwidth}{X} 「你要用純金做一面牌，如同刻印章，在上面刻『歸耶和華為聖』。 \end{tabularx} \\ \\ \relax
28:37 & \begin{tabularx}{0.7\textwidth}{X} 要用藍色的帶子把牌繫在禮冠上，在禮冠的正前面。 \end{tabularx} \\ \\ \relax
28:38 & \begin{tabularx}{0.7\textwidth}{X} 這牌必在亞倫的額上，亞倫要擔當干犯聖物的罪孽；這聖物是以色列人在一切聖禮物上所分別為聖的。這牌要常在他的額上，使他們可以在耶和華面前蒙悅納。 \end{tabularx} \\ \\ \relax
28:39 & \begin{tabularx}{0.7\textwidth}{X} 要用細麻編織內袍，用細麻做禮冠，又以刺繡的手藝做腰帶。 \end{tabularx} \\ \\ \relax
28:40 & \begin{tabularx}{0.7\textwidth}{X} 「你要為亞倫的兒子做內袍、腰帶、頭巾，以示尊嚴和華美。 \end{tabularx} \\ \\ \relax
28:41 & \begin{tabularx}{0.7\textwidth}{X} 要把這些給你哥哥亞倫和他的兒子穿戴，又要膏他們，授予聖職，使他們分別為聖，作事奉我的祭司。 \end{tabularx} \\ \\ \relax
28:42 & \begin{tabularx}{0.7\textwidth}{X} 要用細麻布給他們做褲子來遮掩下體，從腰間直到大腿。 \end{tabularx} \\ \\ \relax
28:43 & \begin{tabularx}{0.7\textwidth}{X} 亞倫和他兒子進入會幕，或接近祭壇，在聖所供職的時候要穿上褲子，免得擔當罪孽而死。這要成為亞倫和他後裔永遠的定例。」 \end{tabularx} \\ \\ \relax
29:1 & \begin{tabularx}{0.7\textwidth}{X} 「這是你使他們分別為聖，作事奉我的祭司時要做的事：取一頭公牛犢，兩隻無殘疾的公綿羊， \end{tabularx} \\ \\ \relax
29:2 & \begin{tabularx}{0.7\textwidth}{X} 無酵餅、用油調和的無酵餅，和抹油的無酵薄餅；這些餅都要用細麥麵做成。 \end{tabularx} \\ \\ \relax
29:3 & \begin{tabularx}{0.7\textwidth}{X} 這些餅要裝在一個籃子裡，用籃子帶來，又把公牛和兩隻公綿羊牽來。 \end{tabularx} \\ \\ \relax
29:4 & \begin{tabularx}{0.7\textwidth}{X} 要帶亞倫和他兒子到會幕的門口，用水洗他們。 \end{tabularx} \\ \\ \relax
29:5 & \begin{tabularx}{0.7\textwidth}{X} 要拿服裝，給亞倫穿上內袍和以弗得的外袍，以及以弗得，又帶上胸袋，束上以弗得精緻的帶子。 \end{tabularx} \\ \\ \relax
29:6 & \begin{tabularx}{0.7\textwidth}{X} 要把禮冠戴在他頭上，將聖冕加在禮冠上， \end{tabularx} \\ \\ \relax
29:7 & \begin{tabularx}{0.7\textwidth}{X} 把膏油倒在他頭上膏他。 \end{tabularx} \\ \\ \relax
29:8 & \begin{tabularx}{0.7\textwidth}{X} 要帶他的兒子來，給他們穿上內袍。 \end{tabularx} \\ \\ \relax
29:9 & \begin{tabularx}{0.7\textwidth}{X} 要給亞倫和他的兒子束上腰帶，裹上頭巾，他們就憑永遠的定例得祭司的職分。又要授聖職給亞倫和他的兒子。 \end{tabularx} \\ \\ \relax
29:10 & \begin{tabularx}{0.7\textwidth}{X} 「你要把公牛牽到會幕前，亞倫和他的兒子要按手在公牛的頭上。 \end{tabularx} \\ \\ \relax
29:11 & \begin{tabularx}{0.7\textwidth}{X} 你要在耶和華面前，在會幕的門口宰這公牛。 \end{tabularx} \\ \\ \relax
29:12 & \begin{tabularx}{0.7\textwidth}{X} 要取些公牛的血，用指頭抹在祭壇的四個翹角上，把其餘的血全倒在壇的底座上。 \end{tabularx} \\ \\ \relax
29:13 & \begin{tabularx}{0.7\textwidth}{X} 要把所有包著內臟的脂肪、肝上的網油、兩個腎和腎上的脂肪，都燒在壇上。 \end{tabularx} \\ \\ \relax
29:14 & \begin{tabularx}{0.7\textwidth}{X} 只是公牛的肉、皮、糞都要在營外用火焚燒；這牛是贖罪祭。 \end{tabularx} \\ \\ \relax
29:15 & \begin{tabularx}{0.7\textwidth}{X} 「你要牽一隻公綿羊來，亞倫和他兒子要按手在這羊的頭上。 \end{tabularx} \\ \\ \relax
29:16 & \begin{tabularx}{0.7\textwidth}{X} 你要宰這羊，把血灑在祭壇的周圍。 \end{tabularx} \\ \\ \relax
29:17 & \begin{tabularx}{0.7\textwidth}{X} 再把羊切成肉塊，洗淨內臟和腿，連肉塊和頭放在一處。 \end{tabularx} \\ \\ \relax
29:18 & \begin{tabularx}{0.7\textwidth}{X} 要把全羊燒在壇上。這是獻給耶和華的燔祭，是獻給耶和華馨香的火祭。」 \end{tabularx} \\ \\ \relax
29:19 & \begin{tabularx}{0.7\textwidth}{X} 「你要把第二隻公綿羊牽來，亞倫和他兒子要按手在這羊的頭上。 \end{tabularx} \\ \\ \relax
29:20 & \begin{tabularx}{0.7\textwidth}{X} 你要宰這羊，取些血抹在亞倫的右耳垂和他兒子的右耳垂上，又抹在他們右手的大拇指和右腳的大腳趾上，然後把其餘的血灑在壇的周圍。 \end{tabularx} \\ \\ \relax
29:21 & \begin{tabularx}{0.7\textwidth}{X} 你要取些膏油和壇上的血，彈在亞倫和他的衣服上，以及他兒子和他們的衣服上；亞倫和他的衣服，他兒子和他們的衣服都成為聖了。 \end{tabularx} \\ \\ \relax
29:22 & \begin{tabularx}{0.7\textwidth}{X} 「你要取這羊的脂肪，肥尾巴、包著內臟的脂肪、肝上的網油、兩個腎、腎上的脂肪和右腿，這是聖職禮所獻的公綿羊； \end{tabularx} \\ \\ \relax
29:23 & \begin{tabularx}{0.7\textwidth}{X} 再從耶和華面前那裝無酵餅的籃子中取一個餅、一個油餅和一個薄餅， \end{tabularx} \\ \\ \relax
29:24 & \begin{tabularx}{0.7\textwidth}{X} 把它們都放在亞倫的手和他兒子的手上，在耶和華面前搖一搖，作為搖祭。 \end{tabularx} \\ \\ \relax
29:25 & \begin{tabularx}{0.7\textwidth}{X} 然後，你要從他們手中接過來，放在燔祭上，一起燒在壇上，作為耶和華面前馨香之氣；這是獻給耶和華的火祭。 \end{tabularx} \\ \\ \relax
29:26 & \begin{tabularx}{0.7\textwidth}{X} 「你要取亞倫聖職禮所獻公綿羊的胸，在耶和華面前搖一搖，作為搖祭；這份就是你的。 \end{tabularx} \\ \\ \relax
29:27 & \begin{tabularx}{0.7\textwidth}{X} 那搖祭的胸和舉祭的腿，就是聖職禮獻公綿羊時所搖的、所舉的，你要使它們分別為聖，是歸給亞倫和他兒子的。 \end{tabularx} \\ \\ \relax
29:28 & \begin{tabularx}{0.7\textwidth}{X} 這是亞倫和他子孫憑永遠的定例從以色列人中所應得的；因為這是舉祭，是從以色列人的平安祭中取出，作為獻給耶和華的舉祭。 \end{tabularx} \\ \\ \relax
29:29 & \begin{tabularx}{0.7\textwidth}{X} 「亞倫的聖衣要傳給他的子孫，使他們在受膏和承接聖職的時候穿上。 \end{tabularx} \\ \\ \relax
29:30 & \begin{tabularx}{0.7\textwidth}{X} 他的子孫接續他當祭司的，每逢進入會幕在聖所供職的時候，要穿這聖衣七天。 \end{tabularx} \\ \\ \relax
29:31 & \begin{tabularx}{0.7\textwidth}{X} 「你要拿聖職禮所獻的公綿羊，在聖處煮牠的肉。 \end{tabularx} \\ \\ \relax
29:32 & \begin{tabularx}{0.7\textwidth}{X} 亞倫和他兒子要在會幕的門口吃這羊的肉和籃子裡的餅。 \end{tabularx} \\ \\ \relax
29:33 & \begin{tabularx}{0.7\textwidth}{X} 他們要吃那些用來贖罪之物，好承接聖職，使他們分別為聖。外人不可吃，因為這是聖物。 \end{tabularx} \\ \\ \relax
29:34 & \begin{tabularx}{0.7\textwidth}{X} 那聖職禮所獻的肉或餅，若有剩餘留到早晨，就要把剩下的用火燒了，不可再吃，因為這是聖物。 \end{tabularx} \\ \\ \relax
29:35 & \begin{tabularx}{0.7\textwidth}{X} 「你要這樣照我一切所吩咐的，向亞倫和他兒子行授聖職禮七天。 \end{tabularx} \\ \\ \relax
29:36 & \begin{tabularx}{0.7\textwidth}{X} 為了贖罪，每天要獻一頭公牛為贖罪祭。你要為祭壇贖罪，使壇潔淨，並要用膏抹壇，使壇成為聖。 \end{tabularx} \\ \\ \relax
29:37 & \begin{tabularx}{0.7\textwidth}{X} 要為壇贖罪七天，使壇成為聖，壇就成為至聖。凡觸摸壇的都成為聖。」 \end{tabularx} \\ \\ \relax
29:38 & \begin{tabularx}{0.7\textwidth}{X} 「這是你要獻在壇上的：每天不可間斷地獻兩隻一歲的羔羊； \end{tabularx} \\ \\ \relax
29:39 & \begin{tabularx}{0.7\textwidth}{X} 早晨獻第一隻羔羊，黃昏獻第二隻羔羊。 \end{tabularx} \\ \\ \relax
29:40 & \begin{tabularx}{0.7\textwidth}{X} 獻第一隻羔羊時，要同時獻上十分之一伊法細麵，調和四分之一欣搗成的油，再獻四分之一欣酒作澆酒祭。 \end{tabularx} \\ \\ \relax
29:41 & \begin{tabularx}{0.7\textwidth}{X} 黃昏你獻第二隻羔羊，要照早晨的素祭和同獻的澆酒祭獻上，作為獻給耶和華馨香的火祭。 \end{tabularx} \\ \\ \relax
29:42 & \begin{tabularx}{0.7\textwidth}{X} 這要在耶和華面前，在會幕的門口，作為你們世世代代經常獻的燔祭。我要在那裡與你們相會，和你說話。 \end{tabularx} \\ \\ \relax
29:43 & \begin{tabularx}{0.7\textwidth}{X} 我要在那裡與以色列人相會，會幕就要因我的榮耀成為聖。 \end{tabularx} \\ \\ \relax
29:44 & \begin{tabularx}{0.7\textwidth}{X} 我要使會幕和祭壇分別為聖，也要使亞倫和他的兒子分別為聖，作事奉我的祭司。 \end{tabularx} \\ \\ \relax
29:45 & \begin{tabularx}{0.7\textwidth}{X} 我要住在以色列人中，作他們的神。 \end{tabularx} \\ \\ \relax
29:46 & \begin{tabularx}{0.7\textwidth}{X} 他們必知道我是耶和華－他們的神，是將他們從埃及地領出來的，為要住在他們中間。我是耶和華－他們的神 。」 \end{tabularx} \\ \\ \relax
30:1 & \begin{tabularx}{0.7\textwidth}{X} 「你要用金合歡木做一座燒香的壇， \end{tabularx} \\ \\ \relax
30:2 & \begin{tabularx}{0.7\textwidth}{X} 長一肘，寬一肘，這壇是正方形的，高二肘。壇的四個翹角與壇接連一塊。 \end{tabularx} \\ \\ \relax
30:3 & \begin{tabularx}{0.7\textwidth}{X} 要把壇的上面與壇的四圍，以及壇的四個翹角包上純金；又要在壇的四圍鑲上金邊。 \end{tabularx} \\ \\ \relax
30:4 & \begin{tabularx}{0.7\textwidth}{X} 要在壇的兩個對側，金邊下面做兩個金環，用來穿槓抬壇。 \end{tabularx} \\ \\ \relax
30:5 & \begin{tabularx}{0.7\textwidth}{X} 要用金合歡木做槓，包上金子。 \end{tabularx} \\ \\ \relax
30:6 & \begin{tabularx}{0.7\textwidth}{X} 要把壇放在法櫃前的幔子外，對著法櫃上的櫃蓋，就是我與你相會的地方。 \end{tabularx} \\ \\ \relax
30:7 & \begin{tabularx}{0.7\textwidth}{X} 亞倫要在壇上燒芬芳的香；每早晨整理燈的時候，他都要燒這香。 \end{tabularx} \\ \\ \relax
30:8 & \begin{tabularx}{0.7\textwidth}{X} 黃昏點燈的時候，亞倫也要燒這香。這是你們世世代代在耶和華面前常燒的香。 \end{tabularx} \\ \\ \relax
30:9 & \begin{tabularx}{0.7\textwidth}{X} 在這壇上不可燒別樣的香，不可獻燔祭、素祭，也不可獻澆酒祭。 \end{tabularx} \\ \\ \relax
30:10 & \begin{tabularx}{0.7\textwidth}{X} 亞倫每年一次要為壇的四個翹角贖罪。他每年一次要用贖罪祭的血為壇贖罪，作為世世代代的定例。這壇在耶和華面前是至聖的。」 \end{tabularx} \\ \\ \relax
30:11 & \begin{tabularx}{0.7\textwidth}{X} 耶和華吩咐摩西說： \end{tabularx} \\ \\ \relax
30:12 & \begin{tabularx}{0.7\textwidth}{X} 「你數點以色列人，計算人頭時，被數的每一個人要把他生命的贖價獻給耶和華，免得災殃在數點中臨到他們。 \end{tabularx} \\ \\ \relax
30:13 & \begin{tabularx}{0.7\textwidth}{X} 每一個被數的人要按照聖所的舍客勒，付半舍客勒，一舍客勒是二十季拉；這半舍客勒是獻給耶和華的禮物。 \end{tabularx} \\ \\ \relax
30:14 & \begin{tabularx}{0.7\textwidth}{X} 每一個被數的人，就是二十歲以上的，要將這禮物獻給耶和華。 \end{tabularx} \\ \\ \relax
30:15 & \begin{tabularx}{0.7\textwidth}{X} 富有的不必多付，貧窮的也不可少出，各人都要獻半舍客勒給耶和華，作你們生命的贖價。 \end{tabularx} \\ \\ \relax
30:16 & \begin{tabularx}{0.7\textwidth}{X} 你要向以色列人收這贖罪的銀子，用在會幕的事工。這要在耶和華面前為以色列人作紀念，作你們生命的贖價。」 \end{tabularx} \\ \\ \relax
30:17 & \begin{tabularx}{0.7\textwidth}{X} 耶和華吩咐摩西說： \end{tabularx} \\ \\ \relax
30:18 & \begin{tabularx}{0.7\textwidth}{X} 「你要用銅做洗濯盆和盆座，用來洗濯。要將盆放在會幕和祭壇的中間，盆裡盛水。 \end{tabularx} \\ \\ \relax
30:19 & \begin{tabularx}{0.7\textwidth}{X} 亞倫和他的兒子要用這盆洗手洗腳。 \end{tabularx} \\ \\ \relax
30:20 & \begin{tabularx}{0.7\textwidth}{X} 他們進會幕，或是走近壇前供職，獻火祭給耶和華的時候，必須用水洗濯，免得死亡； \end{tabularx} \\ \\ \relax
30:21 & \begin{tabularx}{0.7\textwidth}{X} 他們要洗手洗腳，免得死亡。這是亞倫和他的後裔世世代代永遠的定例。」 \end{tabularx} \\ \\ \relax
30:22 & \begin{tabularx}{0.7\textwidth}{X} 耶和華吩咐摩西說： \end{tabularx} \\ \\ \relax
30:23 & \begin{tabularx}{0.7\textwidth}{X} 「你要取上等的香料，就是五百舍客勒流質的沒藥、二百五十香肉桂、二百五十香菖蒲， \end{tabularx} \\ \\ \relax
30:24 & \begin{tabularx}{0.7\textwidth}{X} 和五百桂皮，都按照聖所的舍客勒；再取一欣橄欖油， \end{tabularx} \\ \\ \relax
30:25 & \begin{tabularx}{0.7\textwidth}{X} 以做香的方法調和製成聖膏油，它就成為聖膏油。 \end{tabularx} \\ \\ \relax
30:26 & \begin{tabularx}{0.7\textwidth}{X} 要用這膏油抹會幕和法櫃， \end{tabularx} \\ \\ \relax
30:27 & \begin{tabularx}{0.7\textwidth}{X} 供桌和供桌的一切器具，燈臺和燈臺的器具，以及香壇、 \end{tabularx} \\ \\ \relax
30:28 & \begin{tabularx}{0.7\textwidth}{X} 燔祭壇和壇的一切器具，洗濯盆和盆座。 \end{tabularx} \\ \\ \relax
30:29 & \begin{tabularx}{0.7\textwidth}{X} 你要使這些分別為聖，成為至聖；凡觸摸它們的都成為聖。 \end{tabularx} \\ \\ \relax
30:30 & \begin{tabularx}{0.7\textwidth}{X} 要膏亞倫和他的兒子，使他們分別為聖，作事奉我的祭司。 \end{tabularx} \\ \\ \relax
30:31 & \begin{tabularx}{0.7\textwidth}{X} 你要吩咐以色列人說：『你們要世世代代以這油為我的聖膏油。 \end{tabularx} \\ \\ \relax
30:32 & \begin{tabularx}{0.7\textwidth}{X} 不可把這油倒在別人身上，也不可用配製這膏油的方法製成同樣的膏油。這膏油是聖的，你們要以它為聖。 \end{tabularx} \\ \\ \relax
30:33 & \begin{tabularx}{0.7\textwidth}{X} 凡調和與此類似的膏油，或將它膏在別人身上的，這人要從百姓中剪除。』」 \end{tabularx} \\ \\ \relax
30:34 & \begin{tabularx}{0.7\textwidth}{X} 耶和華吩咐摩西說：「你要取香料，就是拿他弗、施喜列、喜利比拿，這些香料再加純乳香，每樣都要相同的分量。 \end{tabularx} \\ \\ \relax
30:35 & \begin{tabularx}{0.7\textwidth}{X} 你要用這些加上鹽，以配製香料的方法，製成純淨又神聖的香。 \end{tabularx} \\ \\ \relax
30:36 & \begin{tabularx}{0.7\textwidth}{X} 要取一點這香，搗成細的粉，放在會幕中的法櫃前，就是我和你相會的地方。你們要以這香為至聖。 \end{tabularx} \\ \\ \relax
30:37 & \begin{tabularx}{0.7\textwidth}{X} 你們不可用這配製的方法為自己做香；要以這香為聖，歸於耶和華。 \end{tabularx} \\ \\ \relax
30:38 & \begin{tabularx}{0.7\textwidth}{X} 為要聞香味而配製同樣的香的，這人要從百姓中剪除。」 \end{tabularx} \\ \\ \relax
31:1 & \begin{tabularx}{0.7\textwidth}{X} 耶和華吩咐摩西說： \end{tabularx} \\ \\ \relax
31:2 & \begin{tabularx}{0.7\textwidth}{X} 「你看，我已經題名召猶大支派中戶珥的孫子，烏利的兒子比撒列。 \end{tabularx} \\ \\ \relax
31:3 & \begin{tabularx}{0.7\textwidth}{X} 我以神的靈充滿他，使他有智慧，有聰明，有知識，能做各樣的工， \end{tabularx} \\ \\ \relax
31:4 & \begin{tabularx}{0.7\textwidth}{X} 能設計圖案，用金、銀、銅製造各物， \end{tabularx} \\ \\ \relax
31:5 & \begin{tabularx}{0.7\textwidth}{X} 又能雕刻鑲嵌用的寶石，雕刻木頭，做各樣的工。 \end{tabularx} \\ \\ \relax
31:6 & \begin{tabularx}{0.7\textwidth}{X} 看哪，我委派但支派中亞希撒抹的兒子亞何利亞伯與他同工。凡心裡有智慧的，我更要賜給他們智慧的心，能做我所吩咐的一切， \end{tabularx} \\ \\ \relax
31:7 & \begin{tabularx}{0.7\textwidth}{X} 就是會幕、法櫃和其上的櫃蓋、會幕中一切的器具、 \end{tabularx} \\ \\ \relax
31:8 & \begin{tabularx}{0.7\textwidth}{X} 供桌和供桌的器具、純金的燈臺和燈臺的一切器具、香壇、 \end{tabularx} \\ \\ \relax
31:9 & \begin{tabularx}{0.7\textwidth}{X} 燔祭壇和壇的一切器具、洗濯盆與盆座、 \end{tabularx} \\ \\ \relax
31:10 & \begin{tabularx}{0.7\textwidth}{X} 供祭司職分用的精緻禮服，亞倫祭司的聖衣和他兒子的衣服， \end{tabularx} \\ \\ \relax
31:11 & \begin{tabularx}{0.7\textwidth}{X} 以及膏油和聖所用的芬芳的香。他們都要照我所吩咐的一切去做。」 \end{tabularx} \\ \\ \relax
31:12 & \begin{tabularx}{0.7\textwidth}{X} 耶和華對摩西說： \end{tabularx} \\ \\ \relax
31:13 & \begin{tabularx}{0.7\textwidth}{X} 「你要吩咐以色列人說：『你們務要守我的安息日，因為這是你我之間世世代代的記號，叫你們知道我是耶和華，是使你們分別為聖的。 \end{tabularx} \\ \\ \relax
31:14 & \begin{tabularx}{0.7\textwidth}{X} 你們要守安息日，以它為聖日。凡干犯這日的，必被處死；凡在這日做工的，那人必從百姓中剪除。 \end{tabularx} \\ \\ \relax
31:15 & \begin{tabularx}{0.7\textwidth}{X} 六日要做工，但第七日是向耶和華守完全安息的安息聖日。凡在安息日做工的，必被處死。』 \end{tabularx} \\ \\ \relax
31:16 & \begin{tabularx}{0.7\textwidth}{X} 以色列人要守安息日，世世代代守安息日為永遠的約。 \end{tabularx} \\ \\ \relax
31:17 & \begin{tabularx}{0.7\textwidth}{X} 這是我和以色列人之間永遠的記號，因為六日之內耶和華造天地，第七日就安息舒暢。」 \end{tabularx} \\ \\ \relax
31:18 & \begin{tabularx}{0.7\textwidth}{X} 耶和華在西奈山和摩西說完了話，就把兩塊法版交給他，是神用指頭寫的石版。 \end{tabularx} \\ \\
[1ex]
\hline
\hline
\end{longtable}


\newpage

\section{第65屆~港九培靈會~張慕皚~第一講}
\label{sec:991}
\textbf{教會 — 末世的金燈臺}
\newline
\newline
連結: \href{https://www.hkbibleconference.org/session-message/view/991}{\texttt{https://www.hkbibleconference.org/session-message/view/991}}
\newline
\newline
\hyperref[sec:990]{< < < PREV SERMON < < <}
~
\hyperref[sec:toc]{[返目錄]}
~
\hyperref[sec:992]{> > > NEXT SERMON > > >}
\newline
\newline
啟示錄 1:9-1:20
\newline
\begin{longtable}{cl}
\hline
\hline
章節 & 經文 (和合本修訂版)\\
\hline
1:9 & \begin{tabularx}{0.7\textwidth}{X} 我—約翰就是你們的弟兄，在耶穌裡和你們一同在患難、國度、忍耐裡有份的，為神的道，並為給耶穌作的見證，曾在那名叫拔摩的海島上。 \end{tabularx} \\ \\ \relax
1:10 & \begin{tabularx}{0.7\textwidth}{X} 有一主日我被聖靈感動，聽見在我後面有大聲音如吹號， \end{tabularx} \\ \\ \relax
1:11 & \begin{tabularx}{0.7\textwidth}{X} 說：「把你所看見的寫在書上，寄給以弗所、士每拿、別迦摩、推雅推喇、撒狄、非拉鐵非、老底嘉那七個教會。」 \end{tabularx} \\ \\ \relax
1:12 & \begin{tabularx}{0.7\textwidth}{X} 我轉過身來要看看是誰的聲音在跟我說話。我一轉過來，看見了七個金燈臺； \end{tabularx} \\ \\ \relax
1:13 & \begin{tabularx}{0.7\textwidth}{X} 在燈臺中間有一位好像人子的，身穿垂到腳的長袍，胸間束著金帶。 \end{tabularx} \\ \\ \relax
1:14 & \begin{tabularx}{0.7\textwidth}{X} 他的頭與髮皆白，如白羊毛，如雪；他的眼睛好像火焰， \end{tabularx} \\ \\ \relax
1:15 & \begin{tabularx}{0.7\textwidth}{X} 雙腳好像在爐中鍛鍊得發亮的銅，聲音好像眾水的聲音。 \end{tabularx} \\ \\ \relax
1:16 & \begin{tabularx}{0.7\textwidth}{X} 他右手拿著七顆星，從他口中吐出一把兩刃的利劍，面貌好像烈日放光。 \end{tabularx} \\ \\ \relax
1:17 & \begin{tabularx}{0.7\textwidth}{X} 我看見了他，就仆倒在他腳前，像死人一樣。他用右手按著我說：「不要怕。我是首先的，是末後的， \end{tabularx} \\ \\ \relax
1:18 & \begin{tabularx}{0.7\textwidth}{X} 又是永活的。我曾死過，看哪，我是活著的，直到永永遠遠；並且我拿著死亡和陰間的鑰匙。 \end{tabularx} \\ \\ \relax
1:19 & \begin{tabularx}{0.7\textwidth}{X} 所以，你要把所看見的事、現在的事和以後將發生的事，都寫下來。 \end{tabularx} \\ \\ \relax
1:20 & \begin{tabularx}{0.7\textwidth}{X} 至於你所看見、在我右手中的七顆星和那七個金燈臺的奧祕就是：七顆星是七個教會的使者，七個燈臺是七個教會。」 \end{tabularx} \\ \\
[1ex]
\hline
\hline
\end{longtable}


\newpage

\section{第65屆~港九培靈會~張慕皚~第二講}
\label{sec:992}
\textbf{以弗所教會 — 愛的光輝}
\newline
\newline
連結: \href{https://www.hkbibleconference.org/session-message/view/992}{\texttt{https://www.hkbibleconference.org/session-message/view/992}}
\newline
\newline
\hyperref[sec:991]{< < < PREV SERMON < < <}
~
\hyperref[sec:toc]{[返目錄]}
~
\hyperref[sec:993]{> > > NEXT SERMON > > >}
\newline
\newline
啟示錄 2:1-2:7
\newline
\begin{longtable}{cl}
\hline
\hline
章節 & 經文 (和合本修訂版)\\
\hline
2:1 & \begin{tabularx}{0.7\textwidth}{X} 「你要寫信給以弗所教會的使者，說：『那右手拿著七顆星，在七個金燈臺中間行走的這樣說： \end{tabularx} \\ \\ \relax
2:2 & \begin{tabularx}{0.7\textwidth}{X} 我知道你的行為、勞碌、忍耐，也知道你不容忍惡人。你也曾察驗那自稱為使徒卻不是使徒的，看出他們是假的。 \end{tabularx} \\ \\ \relax
2:3 & \begin{tabularx}{0.7\textwidth}{X} 你能忍耐，曾為我的名勞苦而不困倦。 \end{tabularx} \\ \\ \relax
2:4 & \begin{tabularx}{0.7\textwidth}{X} 然而，有一件事我要責備你，就是你把起初的愛心拋棄了。 \end{tabularx} \\ \\ \relax
2:5 & \begin{tabularx}{0.7\textwidth}{X} 所以你要回想你是從哪裡墜落的，並且要悔改，做起初所做的工作。你若不悔改，我要到你那裡去，把你的燈臺從原處挪去。 \end{tabularx} \\ \\ \relax
2:6 & \begin{tabularx}{0.7\textwidth}{X} 然而你還有一件可取的事，就是你恨惡尼哥拉派的行為，這種行為也是我所恨惡的。 \end{tabularx} \\ \\ \relax
2:7 & \begin{tabularx}{0.7\textwidth}{X} 凡有耳朵的都應當聽聖靈向眾教會所說的話。得勝的，我必將神樂園中生命樹的果子賜給他吃。』」 \end{tabularx} \\ \\
[1ex]
\hline
\hline
\end{longtable}


\newpage

\section{第65屆~港九培靈會~張慕皚~第三講}
\label{sec:993}
\textbf{士每拿教會 — 苦難的光輝}
\newline
\newline
連結: \href{https://www.hkbibleconference.org/session-message/view/993}{\texttt{https://www.hkbibleconference.org/session-message/view/993}}
\newline
\newline
\hyperref[sec:992]{< < < PREV SERMON < < <}
~
\hyperref[sec:toc]{[返目錄]}
~
\hyperref[sec:994]{> > > NEXT SERMON > > >}
\newline
\newline
啟示錄 2:8-2:11
\newline
\begin{longtable}{cl}
\hline
\hline
章節 & 經文 (和合本修訂版)\\
\hline
2:8 & \begin{tabularx}{0.7\textwidth}{X} 「你要寫信給士每拿教會的使者，說：『那首先的、末後的，死過又活了的這樣說： \end{tabularx} \\ \\ \relax
2:9 & \begin{tabularx}{0.7\textwidth}{X} 我知道你的患難和貧窮—其實你卻是富足的，也知道那自稱是猶太人的所說毀謗的話，其實他們不是猶太人，而是撒但會堂的人。 \end{tabularx} \\ \\ \relax
2:10 & \begin{tabularx}{0.7\textwidth}{X} 你將要受的苦，你不用怕。看哪！魔鬼要把你們中間幾個人下在監裡，使你們受考驗，你們要遭受苦難十日。你務要至死忠心，我就賜給你那生命的冠冕。 \end{tabularx} \\ \\ \relax
2:11 & \begin{tabularx}{0.7\textwidth}{X} 凡有耳朵的都應當聽聖靈向眾教會所說的話。得勝的必不受第二次死的害。』」 \end{tabularx} \\ \\
[1ex]
\hline
\hline
\end{longtable}


\newpage

\section{第65屆~港九培靈會~張慕皚~第四講}
\label{sec:994}
\textbf{別迦摩教會 — 聖潔生活的光輝}
\newline
\newline
連結: \href{https://www.hkbibleconference.org/session-message/view/994}{\texttt{https://www.hkbibleconference.org/session-message/view/994}}
\newline
\newline
\hyperref[sec:993]{< < < PREV SERMON < < <}
~
\hyperref[sec:toc]{[返目錄]}
~
\hyperref[sec:995]{> > > NEXT SERMON > > >}
\newline
\newline
啟示錄 2:12-2:17
\newline
\begin{longtable}{cl}
\hline
\hline
章節 & 經文 (和合本修訂版)\\
\hline
2:12 & \begin{tabularx}{0.7\textwidth}{X} 「你要寫信給別迦摩教會的使者，說：『那有兩刃利劍的這樣說： \end{tabularx} \\ \\ \relax
2:13 & \begin{tabularx}{0.7\textwidth}{X} 我知道你的居所，就是有撒但座位之處；當我忠心的見證人安提帕在你們中間，在撒但所住的地方被殺之時，你還堅守我的名，沒有否認對我的信仰。 \end{tabularx} \\ \\ \relax
2:14 & \begin{tabularx}{0.7\textwidth}{X} 然而，有幾件事我要責備你，就是在你那裡有人服從了巴蘭的教訓；這巴蘭曾教唆巴勒將絆腳石放在以色列人面前，使他們吃祭過偶像之物，並且犯淫亂。 \end{tabularx} \\ \\ \relax
2:15 & \begin{tabularx}{0.7\textwidth}{X} 同樣，你那裡也有人服從了尼哥拉派的教訓。 \end{tabularx} \\ \\ \relax
2:16 & \begin{tabularx}{0.7\textwidth}{X} 所以，你當悔改；若不悔改，我很快就到你那裡來，用我口中的劍攻擊他們。 \end{tabularx} \\ \\ \relax
2:17 & \begin{tabularx}{0.7\textwidth}{X} 凡有耳朵的都應當聽聖靈向眾教會所說的話。得勝的，我必將那隱藏的嗎哪賜給他，並賜他一塊白石，石上寫著新的名字，除了那領受的以外，沒有人認識。』」 \end{tabularx} \\ \\
[1ex]
\hline
\hline
\end{longtable}


\newpage

\section{第65屆~港九培靈會~張慕皚~第五講}
\label{sec:995}
\textbf{推雅推喇教會 — 職業的光輝}
\newline
\newline
連結: \href{https://www.hkbibleconference.org/session-message/view/995}{\texttt{https://www.hkbibleconference.org/session-message/view/995}}
\newline
\newline
\hyperref[sec:994]{< < < PREV SERMON < < <}
~
\hyperref[sec:toc]{[返目錄]}
~
\hyperref[sec:996]{> > > NEXT SERMON > > >}
\newline
\newline
啟示錄 2:18-2:29
\newline
\begin{longtable}{cl}
\hline
\hline
章節 & 經文 (和合本修訂版)\\
\hline
2:18 & \begin{tabularx}{0.7\textwidth}{X} 「你要寫信給推雅推喇教會的使者，說：『神的兒子，那位眼睛如火焰、雙腳像發亮的銅的這樣說： \end{tabularx} \\ \\ \relax
2:19 & \begin{tabularx}{0.7\textwidth}{X} 我知道你的行為：愛心、信心、勤勞、忍耐；又知道你末後所行的善事比起初所行的更多。 \end{tabularx} \\ \\ \relax
2:20 & \begin{tabularx}{0.7\textwidth}{X} 然而，有一件事我要責備你，就是你容忍那自稱是先知的婦人耶洗別教唆我的僕人，引誘他們犯淫亂，吃祭過偶像之物。 \end{tabularx} \\ \\ \relax
2:21 & \begin{tabularx}{0.7\textwidth}{X} 我曾給她悔改的機會，她卻不肯悔改她的淫行。 \end{tabularx} \\ \\ \relax
2:22 & \begin{tabularx}{0.7\textwidth}{X} 看吧，我要使她病倒在床上。那些與她犯姦淫的人若不悔改他們的行為，我也要使他們同受大患難。 \end{tabularx} \\ \\ \relax
2:23 & \begin{tabularx}{0.7\textwidth}{X} 我又要殺死她的兒女，眾教會就知道，我是那察看人肺腑心腸的，我要照你們的行為報應各人。 \end{tabularx} \\ \\ \relax
2:24 & \begin{tabularx}{0.7\textwidth}{X} 至於你們其餘的推雅推喇人，就是一切不隨從這教訓，不明白他們所謂撒但深奧之理的人，我告訴你們，我不會再把別的擔子放在你們身上。 \end{tabularx} \\ \\ \relax
2:25 & \begin{tabularx}{0.7\textwidth}{X} 你們只要持守那已經有的，直到我來。 \end{tabularx} \\ \\ \relax
2:26 & \begin{tabularx}{0.7\textwidth}{X} 那得勝又遵守我命令到底的，我要賜給他權柄制伏列國； \end{tabularx} \\ \\ \relax
2:27 & \begin{tabularx}{0.7\textwidth}{X} 他必用鐵杖管轄他們，如同打碎陶器， \end{tabularx} \\ \\ \relax
2:28 & \begin{tabularx}{0.7\textwidth}{X} 像我也從我父領受了權柄一樣。我又要把晨星賜給他。 \end{tabularx} \\ \\ \relax
2:29 & \begin{tabularx}{0.7\textwidth}{X} 凡有耳朵的都應當聽聖靈向眾教會所說的話。』」 \end{tabularx} \\ \\
[1ex]
\hline
\hline
\end{longtable}


\newpage

\section{第65屆~港九培靈會~張慕皚~第六講}
\label{sec:996}
\textbf{撒狄教會 — 生命更新的光輝}
\newline
\newline
連結: \href{https://www.hkbibleconference.org/session-message/view/996}{\texttt{https://www.hkbibleconference.org/session-message/view/996}}
\newline
\newline
\hyperref[sec:995]{< < < PREV SERMON < < <}
~
\hyperref[sec:toc]{[返目錄]}
~
\hyperref[sec:997]{> > > NEXT SERMON > > >}
\newline
\newline
啟示錄 3:1-3:6
\newline
\begin{longtable}{cl}
\hline
\hline
章節 & 經文 (和合本修訂版)\\
\hline
3:1 & \begin{tabularx}{0.7\textwidth}{X} 「你要寫信給撒狄教會的使者，說：『那有神的七靈和七顆星的這樣說：我知道你的行為，就是名義上你是活的，實際上你是死的。 \end{tabularx} \\ \\ \relax
3:2 & \begin{tabularx}{0.7\textwidth}{X} 你要警醒，堅固那些剩下、快要死的，因為我發現你的行為，在我神面前沒有一樣是完全的。 \end{tabularx} \\ \\ \relax
3:3 & \begin{tabularx}{0.7\textwidth}{X} 所以，要記得你所領受和聽見的；要遵守，並要悔改。你若不警醒，我必如賊一樣來到；我幾時來到你那裡，你絕不會知道。 \end{tabularx} \\ \\ \relax
3:4 & \begin{tabularx}{0.7\textwidth}{X} 然而，在撒狄你還有幾位是未曾污穢自己衣服的，他們會穿白衣與我同行，因為他們是配穿的。 \end{tabularx} \\ \\ \relax
3:5 & \begin{tabularx}{0.7\textwidth}{X} 得勝的必這樣穿白衣，我也不從生命冊上塗去他的名；我要在我父面前，和我父的眾使者面前，宣認他的名。 \end{tabularx} \\ \\ \relax
3:6 & \begin{tabularx}{0.7\textwidth}{X} 凡有耳朵的都應當聽聖靈向眾教會所說的話。』」 \end{tabularx} \\ \\
[1ex]
\hline
\hline
\end{longtable}


\newpage

\section{第65屆~港九培靈會~張慕皚~第七講}
\label{sec:997}
\textbf{非拉鐵非教會 — 福音的光輝 (一)}
\newline
\newline
連結: \href{https://www.hkbibleconference.org/session-message/view/997}{\texttt{https://www.hkbibleconference.org/session-message/view/997}}
\newline
\newline
\hyperref[sec:996]{< < < PREV SERMON < < <}
~
\hyperref[sec:toc]{[返目錄]}
~
\hyperref[sec:998]{> > > NEXT SERMON > > >}
\newline
\newline
啟示錄 3:7-3:13
\newline
\begin{longtable}{cl}
\hline
\hline
章節 & 經文 (和合本修訂版)\\
\hline
3:7 & \begin{tabularx}{0.7\textwidth}{X} 「你要寫信給非拉鐵非教會的使者，說：『那神聖、真實的，拿著大衛的鑰匙，開了就沒有人能關，關了就沒有人能開的這樣說： \end{tabularx} \\ \\ \relax
3:8 & \begin{tabularx}{0.7\textwidth}{X} 我知道你的行為。看哪，我在你面前給你一個敞開的門，是沒有人能關的。我知道你有一點力量，也遵守我的道，沒有否認我的名。 \end{tabularx} \\ \\ \relax
3:9 & \begin{tabularx}{0.7\textwidth}{X} 那屬撒但會堂的，自稱是猶太人，其實不是猶太人，而是說謊話的，我要使他們來到你腳前下拜，使他們知道我已經愛你了。 \end{tabularx} \\ \\ \relax
3:10 & \begin{tabularx}{0.7\textwidth}{X} 因為你遵守了我堅忍的道，我也必在普天下人受試煉的時候保守你免受試煉。 \end{tabularx} \\ \\ \relax
3:11 & \begin{tabularx}{0.7\textwidth}{X} 我必快來，你要持守你所有的，免得人奪去你的冠冕。 \end{tabularx} \\ \\ \relax
3:12 & \begin{tabularx}{0.7\textwidth}{X} 得勝的，我要使他在我神的殿中作柱子，他必不再從那裡出去。我又要把我神的名和我神城的名—從天上我神那裡降下來的新耶路撒冷，和我的新名，都寫在他上面。 \end{tabularx} \\ \\ \relax
3:13 & \begin{tabularx}{0.7\textwidth}{X} 凡有耳朵的都應當聽聖靈向眾教會所說的話。』」 \end{tabularx} \\ \\
[1ex]
\hline
\hline
\end{longtable}


\newpage

\section{第65屆~港九培靈會~張慕皚~第八講}
\label{sec:998}
\textbf{非拉鐵非教會 — 福音的光輝 (二)}
\newline
\newline
連結: \href{https://www.hkbibleconference.org/session-message/view/998}{\texttt{https://www.hkbibleconference.org/session-message/view/998}}
\newline
\newline
\hyperref[sec:997]{< < < PREV SERMON < < <}
~
\hyperref[sec:toc]{[返目錄]}
~
\hyperref[sec:999]{> > > NEXT SERMON > > >}
\newline
\newline
啟示錄 3:7-3:13
\newline
\begin{longtable}{cl}
\hline
\hline
章節 & 經文 (和合本修訂版)\\
\hline
3:7 & \begin{tabularx}{0.7\textwidth}{X} 「你要寫信給非拉鐵非教會的使者，說：『那神聖、真實的，拿著大衛的鑰匙，開了就沒有人能關，關了就沒有人能開的這樣說： \end{tabularx} \\ \\ \relax
3:8 & \begin{tabularx}{0.7\textwidth}{X} 我知道你的行為。看哪，我在你面前給你一個敞開的門，是沒有人能關的。我知道你有一點力量，也遵守我的道，沒有否認我的名。 \end{tabularx} \\ \\ \relax
3:9 & \begin{tabularx}{0.7\textwidth}{X} 那屬撒但會堂的，自稱是猶太人，其實不是猶太人，而是說謊話的，我要使他們來到你腳前下拜，使他們知道我已經愛你了。 \end{tabularx} \\ \\ \relax
3:10 & \begin{tabularx}{0.7\textwidth}{X} 因為你遵守了我堅忍的道，我也必在普天下人受試煉的時候保守你免受試煉。 \end{tabularx} \\ \\ \relax
3:11 & \begin{tabularx}{0.7\textwidth}{X} 我必快來，你要持守你所有的，免得人奪去你的冠冕。 \end{tabularx} \\ \\ \relax
3:12 & \begin{tabularx}{0.7\textwidth}{X} 得勝的，我要使他在我神的殿中作柱子，他必不再從那裡出去。我又要把我神的名和我神城的名—從天上我神那裡降下來的新耶路撒冷，和我的新名，都寫在他上面。 \end{tabularx} \\ \\ \relax
3:13 & \begin{tabularx}{0.7\textwidth}{X} 凡有耳朵的都應當聽聖靈向眾教會所說的話。』」 \end{tabularx} \\ \\
[1ex]
\hline
\hline
\end{longtable}


\newpage

\section{第65屆~港九培靈會~張慕皚~第九講}
\label{sec:999}
\textbf{老底嘉教會 — 長久興旺的光輝 (一)}
\newline
\newline
連結: \href{https://www.hkbibleconference.org/session-message/view/999}{\texttt{https://www.hkbibleconference.org/session-message/view/999}}
\newline
\newline
\hyperref[sec:998]{< < < PREV SERMON < < <}
~
\hyperref[sec:toc]{[返目錄]}
~
\hyperref[sec:1000]{> > > NEXT SERMON > > >}
\newline
\newline
啟示錄 3:14-3:22
\newline
\begin{longtable}{cl}
\hline
\hline
章節 & 經文 (和合本修訂版)\\
\hline
3:14 & \begin{tabularx}{0.7\textwidth}{X} 「你要寫信給老底嘉教會的使者，說：『那位阿們、誠信真實的見證者、神創造的根源這樣說： \end{tabularx} \\ \\ \relax
3:15 & \begin{tabularx}{0.7\textwidth}{X} 我知道你的行為，你也不冷也不熱；我巴不得你或冷或熱。 \end{tabularx} \\ \\ \relax
3:16 & \begin{tabularx}{0.7\textwidth}{X} 既然你如溫水，也不冷也不熱，我要從我口中把你吐出去。 \end{tabularx} \\ \\ \relax
3:17 & \begin{tabularx}{0.7\textwidth}{X} 你說：我是富足的，已經發了財，一樣都不缺，卻不知道你是困苦、可憐、貧窮、瞎眼、赤身的。 \end{tabularx} \\ \\ \relax
3:18 & \begin{tabularx}{0.7\textwidth}{X} 我勸你向我買從火中鍛鍊出來的金子，使你富足；又買白衣穿上，使你赤身的羞恥不露出來；又買眼藥抹你的眼睛，使你能看見。 \end{tabularx} \\ \\ \relax
3:19 & \begin{tabularx}{0.7\textwidth}{X} 凡我所疼愛的，我就責備管教。所以，你要發熱心，也要悔改。 \end{tabularx} \\ \\ \relax
3:20 & \begin{tabularx}{0.7\textwidth}{X} 看哪，我站在門外叩門，若有聽見我聲音而開門的，我要進到他那裡去，我與他，他與我一起吃飯。 \end{tabularx} \\ \\ \relax
3:21 & \begin{tabularx}{0.7\textwidth}{X} 得勝的，我要賜他在我寶座上與我同坐，就如我得了勝，在我父的寶座上與他同坐一般。 \end{tabularx} \\ \\ \relax
3:22 & \begin{tabularx}{0.7\textwidth}{X} 凡有耳朵的都應當聽聖靈向眾教會所說的話。』」 \end{tabularx} \\ \\
[1ex]
\hline
\hline
\end{longtable}


\newpage

\section{第65屆~港九培靈會~張慕皚~第十講}
\label{sec:1000}
\textbf{老底嘉教會 — 長久興旺的光輝 (二)}
\newline
\newline
連結: \href{https://www.hkbibleconference.org/session-message/view/1000}{\texttt{https://www.hkbibleconference.org/session-message/view/1000}}
\newline
\newline
\hyperref[sec:999]{< < < PREV SERMON < < <}
~
\hyperref[sec:toc]{[返目錄]}
~
\hyperref[sec:1390]{> > > NEXT SERMON > > >}
\newline
\newline
啟示錄 3:14-3:22
\newline
\begin{longtable}{cl}
\hline
\hline
章節 & 經文 (和合本修訂版)\\
\hline
3:14 & \begin{tabularx}{0.7\textwidth}{X} 「你要寫信給老底嘉教會的使者，說：『那位阿們、誠信真實的見證者、神創造的根源這樣說： \end{tabularx} \\ \\ \relax
3:15 & \begin{tabularx}{0.7\textwidth}{X} 我知道你的行為，你也不冷也不熱；我巴不得你或冷或熱。 \end{tabularx} \\ \\ \relax
3:16 & \begin{tabularx}{0.7\textwidth}{X} 既然你如溫水，也不冷也不熱，我要從我口中把你吐出去。 \end{tabularx} \\ \\ \relax
3:17 & \begin{tabularx}{0.7\textwidth}{X} 你說：我是富足的，已經發了財，一樣都不缺，卻不知道你是困苦、可憐、貧窮、瞎眼、赤身的。 \end{tabularx} \\ \\ \relax
3:18 & \begin{tabularx}{0.7\textwidth}{X} 我勸你向我買從火中鍛鍊出來的金子，使你富足；又買白衣穿上，使你赤身的羞恥不露出來；又買眼藥抹你的眼睛，使你能看見。 \end{tabularx} \\ \\ \relax
3:19 & \begin{tabularx}{0.7\textwidth}{X} 凡我所疼愛的，我就責備管教。所以，你要發熱心，也要悔改。 \end{tabularx} \\ \\ \relax
3:20 & \begin{tabularx}{0.7\textwidth}{X} 看哪，我站在門外叩門，若有聽見我聲音而開門的，我要進到他那裡去，我與他，他與我一起吃飯。 \end{tabularx} \\ \\ \relax
3:21 & \begin{tabularx}{0.7\textwidth}{X} 得勝的，我要賜他在我寶座上與我同坐，就如我得了勝，在我父的寶座上與他同坐一般。 \end{tabularx} \\ \\ \relax
3:22 & \begin{tabularx}{0.7\textwidth}{X} 凡有耳朵的都應當聽聖靈向眾教會所說的話。』」 \end{tabularx} \\ \\
[1ex]
\hline
\hline
\end{longtable}


\newpage

\section{第65屆~港九培靈會~楊濬哲~序}
\label{sec:1390}
\textbf{序}
\newline
\newline
連結: \href{https://www.hkbibleconference.org/session-message/view/1390}{\texttt{https://www.hkbibleconference.org/session-message/view/1390}}
\newline
\newline
\hyperref[sec:1000]{< < < PREV SERMON < < <}
~
\hyperref[sec:toc]{[返目錄]}
~
\hyperref[sec:1001]{> > > NEXT SERMON > > >}
\newline
\newline
九三年第六十五屆港九培靈研經大會,能如期舉行,並圓滿閉會,實是父神奇妙的浩恩.
\newline
\newline
本屆早堂研經會由中國神學研究院周永健牧師主講.總題:「作神的子民──出埃及記的信息」.周牧師從出埃及記,講解神對以色列人之揀選和拯救,賦與他們子民的身份,頒佈律法,並吩咐他們建造會幕等事蹟,進而抽絲剝繭地闡釋:神愛我們世人,為我們預備救贖計劃,更願意我們信徒按祂心意過得勝生活,引導我們活出神子民的見證,並願常與我們同在,叫我們最終得應許,蒙恩福.
\newline
\newline
午堂講道會,由香港神學院院長褚永華牧師主講.總題:「耶穌的比喻」.透過這些比喻,褚牧師具體地點出神的心意,祂對教會,個人生命和事奉的要求,當頭棒喝地喚醒我們不要閒站,乃當盡自己的責任,在這時代為真理迸發真光.
\newline
\newline
晚堂奮興會,無巧不成書,也是院長級人士,乃建道神學院院長張慕皚牧師.本來原定的講員是第三次應邀之多倫多民眾教會堂主任史保羅牧師,但於七月初傳來緊急信函,謂因大病未癒,須遵醫生忠告休養,故未克依約來港赴會,歉甚云云.頃接此信,實感彷徨不已,一面急請其他委辦會成員代禱,一面打電話到世界各地接洽講員,可惜時間實在太急,許多人都婉拒了!但感謝神,終於為我們預備了張慕皚牧師,心頭大石才得放下,連忙感激讚美父神的恩眷!張牧師之總題: 「從啟示錄七教會看末世教會的見證」.他透過基督對七教會的說話,指出信徒在這黑暗末世中所面對的引誘和挑戰,必須靠主儆醒,不斷更新,委身基督,剛強站穩,為鹽為光,力傳福音,以報主恩,並得獎賞.有人說:「一個人的生命力如何,從他在緊急中質和量的表現,就可看出.」誠然,這讚賞的評語送給張牧師並不為過,因為要在不足一個月時間內,預備十篇連續的奮興會講章殊不簡單；張牧師卻做到了,而且內容踏實,深入得巨細無遺,淺出得人人易明.從三年來晚堂聚會人數最穩定及最多的情況,就可見一斑.
\newline
\newline
十天總赴會者共近四萬人次,除了港,九,新界及離島信徒外,更有來自中國大陸,英國,美國及加拿大的肢體.而在晚堂奮興會中,講員呼召時,共有5人 決志信主,另悔改更新為主而活者有508人,全時間奉獻事主者212人,合共725人.當看見這些作工的果效時,實感主恩浩瀚,更感在主裡的勞苦,並不徒然.
\newline
\newline
在過往一年內,委辦丘育靈牧師蒙主寵召,安息主懷,丘牧師在委辦會中貢獻良多,如今病失良才,實堪惋惜.幸丘牧師自發現患有癌病,即以其充沛靈力和堅毅信心,並神的加恩和加力,與病魔抗拒多年,至死忠心,實留下可敬可佩之風範.此外,朱建磯牧師伉儷亦曾先後入院治病,其他委辦均年事漸高,實極需主內兄姊多多代禱記念.另一方面,繼前年添兩位新委辦後,今年邀得余也魯教授及曾立華牧師加入委辦會,協力策劃大會事工,誠感主恩!
\newline
\newline
在此,我要特別向委辦朱建磯牧師致最深謝意,因他數十年如一日,默默地承擔著統籌宣傳印刷事宜,尤其是整理編校每年講道集的工作,實功不可抹；惜因精神關係,遂於去年辭退編校一職,由黃崇護先生接任,並蒙吳品嫻姊妹,鄭坤英姊妹及雷麗端姊妹等繼續分別為各堂聚會筆錄講章,西門何治能弟兄設計封面,使第六十五屆之講道集如期出版,謹此致謝!
\newline
\newline
本屆大會仍蒙九龍城浸信會,借用全部場地設施,並動員執事,同工,義工及各兄姊協助招待,維持秩序,保安,場務等.又蒙福音傳播中心協助宣傳,影音轉播,錄音,基督教週報刊登新聞稿及廣告,時代論壇刊登廣告,神托會耀安工場寄發宣傳品；彼此配搭,廣事宣傳,藉增效益.更蒙各教會牧者響應鼓勵,主內各兄姊踊躍赴會及奉獻,統致謝忱!
\newline
\newline


\newpage

\section{第65屆~港九培靈會~褚永華~第一講}
\label{sec:1001}
\textbf{浪子的比喻}
\newline
\newline
連結: \href{https://www.hkbibleconference.org/session-message/view/1001}{\texttt{https://www.hkbibleconference.org/session-message/view/1001}}
\newline
\newline
\hyperref[sec:1390]{< < < PREV SERMON < < <}
~
\hyperref[sec:toc]{[返目錄]}
~
\hyperref[sec:1002]{> > > NEXT SERMON > > >}
\newline
\newline
路加福音 15:11-15:32
\newline
\begin{longtable}{cl}
\hline
\hline
章節 & 經文 (和合本修訂版)\\
\hline
15:11 & \begin{tabularx}{0.7\textwidth}{X} 耶穌又說：「一個人有兩個兒子。 \end{tabularx} \\ \\ \relax
15:12 & \begin{tabularx}{0.7\textwidth}{X} 小兒子對父親說：『父親，請你把我應得的家業分給我。』他父親就把財產分給他們。 \end{tabularx} \\ \\ \relax
15:13 & \begin{tabularx}{0.7\textwidth}{X} 過了不多幾天，小兒子把他一切所有的都收拾起來，往遠方去了。在那裡，他任意放蕩，浪費錢財。 \end{tabularx} \\ \\ \relax
15:14 & \begin{tabularx}{0.7\textwidth}{X} 他耗盡了一切所有的，又恰逢那地方有大饑荒，就窮困起來。 \end{tabularx} \\ \\ \relax
15:15 & \begin{tabularx}{0.7\textwidth}{X} 於是他去投靠當地的一個居民，那人打發他到田裡去放豬。 \end{tabularx} \\ \\ \relax
15:16 & \begin{tabularx}{0.7\textwidth}{X} 他恨不得拿豬所吃的豆莢充飢，也沒有人給他甚麼吃的。 \end{tabularx} \\ \\ \relax
15:17 & \begin{tabularx}{0.7\textwidth}{X} 他醒悟過來，就說：『我父親有多少雇工，糧食有餘，我倒在這裡餓死嗎？ \end{tabularx} \\ \\ \relax
15:18 & \begin{tabularx}{0.7\textwidth}{X} 我要起來，到我父親那裡去，對他說：父親！我得罪了天，又得罪了你， \end{tabularx} \\ \\ \relax
15:19 & \begin{tabularx}{0.7\textwidth}{X} 從今以後，我不配稱為你的兒子，把我當作一個雇工吧。』 \end{tabularx} \\ \\ \relax
15:20 & \begin{tabularx}{0.7\textwidth}{X} 於是他起來，往他父親那裡去。相離還遠，他父親看見，就動了慈心，跑去擁抱著他，連連親他。 \end{tabularx} \\ \\ \relax
15:21 & \begin{tabularx}{0.7\textwidth}{X} 兒子對他說：『父親！我得罪了天，又得罪了你，從今以後，我不配稱為你的兒子。』 \end{tabularx} \\ \\ \relax
15:22 & \begin{tabularx}{0.7\textwidth}{X} 父親卻吩咐僕人：『快把那上好的袍子拿出來給他穿，把戒指戴在他指頭上，把鞋穿在他腳上， \end{tabularx} \\ \\ \relax
15:23 & \begin{tabularx}{0.7\textwidth}{X} 把那肥牛犢牽來宰了，我們來吃喝慶祝； \end{tabularx} \\ \\ \relax
15:24 & \begin{tabularx}{0.7\textwidth}{X} 因為我這個兒子是死而復活，失而復得的。』他們就開始慶祝。 \end{tabularx} \\ \\ \relax
15:25 & \begin{tabularx}{0.7\textwidth}{X} 「那時，大兒子正在田裡。他回來，離家不遠時，聽見奏樂跳舞的聲音， \end{tabularx} \\ \\ \relax
15:26 & \begin{tabularx}{0.7\textwidth}{X} 就叫一個僮僕來，問是甚麼事。 \end{tabularx} \\ \\ \relax
15:27 & \begin{tabularx}{0.7\textwidth}{X} 僮僕對他說：『你弟弟回來了，你父親因為他無災無病地回來，把肥牛犢宰了。』 \end{tabularx} \\ \\ \relax
15:28 & \begin{tabularx}{0.7\textwidth}{X} 大兒子就生氣，不肯進去，他父親出來勸他。 \end{tabularx} \\ \\ \relax
15:29 & \begin{tabularx}{0.7\textwidth}{X} 他對父親說：『你看，我服侍你這麼多年，從來沒有違背過你的命令，而你從來沒有給我一隻小山羊，叫我和朋友們一同快樂。 \end{tabularx} \\ \\ \relax
15:30 & \begin{tabularx}{0.7\textwidth}{X} 但你這個兒子和娼妓吃光了你的財產，他一回來，你倒為他宰了肥牛犢。』 \end{tabularx} \\ \\ \relax
15:31 & \begin{tabularx}{0.7\textwidth}{X} 父親對他說：『兒啊！你常和我同在，我所有的一切都是你的； \end{tabularx} \\ \\ \relax
15:32 & \begin{tabularx}{0.7\textwidth}{X} 可是你這個弟弟是死而復活，失而復得的，所以我們理當歡喜慶祝。』」 \end{tabularx} \\ \\
[1ex]
\hline
\hline
\end{longtable}


\newpage

\section{第65屆~港九培靈會~褚永華~第二講}
\label{sec:1002}
\textbf{撒種的比喻}
\newline
\newline
連結: \href{https://www.hkbibleconference.org/session-message/view/1002}{\texttt{https://www.hkbibleconference.org/session-message/view/1002}}
\newline
\newline
\hyperref[sec:1001]{< < < PREV SERMON < < <}
~
\hyperref[sec:toc]{[返目錄]}
~
\hyperref[sec:1003]{> > > NEXT SERMON > > >}
\newline
\newline
馬可福音 4:1-4:20
\newline
\begin{longtable}{cl}
\hline
\hline
章節 & 經文 (和合本修訂版)\\
\hline
4:1 & \begin{tabularx}{0.7\textwidth}{X} 耶穌又在海邊教導人。有一大群人到他那裡聚集，他只好上船坐下。船在海裡，眾人都靠近海，站在岸上。 \end{tabularx} \\ \\ \relax
4:2 & \begin{tabularx}{0.7\textwidth}{X} 耶穌就用許多比喻教導他們。在教導的時候，他對他們說： \end{tabularx} \\ \\ \relax
4:3 & \begin{tabularx}{0.7\textwidth}{X} 「你們聽啊，有一個撒種的出去撒種。 \end{tabularx} \\ \\ \relax
4:4 & \begin{tabularx}{0.7\textwidth}{X} 他撒的時候，有的落在路旁，飛鳥來把它吃掉了。 \end{tabularx} \\ \\ \relax
4:5 & \begin{tabularx}{0.7\textwidth}{X} 有的落在土淺的石頭地上，因為土不深，很快就長出苗來， \end{tabularx} \\ \\ \relax
4:6 & \begin{tabularx}{0.7\textwidth}{X} 太陽出來一曬，因為沒有根就枯乾了。 \end{tabularx} \\ \\ \relax
4:7 & \begin{tabularx}{0.7\textwidth}{X} 有的落在荊棘裡，荊棘長起來，把它擠住了，就結不出果實。 \end{tabularx} \\ \\ \relax
4:8 & \begin{tabularx}{0.7\textwidth}{X} 又有的落在好土裡，就發芽長大，結出果實，有三十倍的，有六十倍的，有一百倍的。」 \end{tabularx} \\ \\ \relax
4:9 & \begin{tabularx}{0.7\textwidth}{X} 耶穌又說：「有耳可聽的，就應當聽！」 \end{tabularx} \\ \\ \relax
4:10 & \begin{tabularx}{0.7\textwidth}{X} 耶穌獨自一人的時候，跟隨他的人和十二使徒問他這些比喻的意思。 \end{tabularx} \\ \\ \relax
4:11 & \begin{tabularx}{0.7\textwidth}{X} 耶穌對他們說：「神國的奧祕只讓你們知道，若是對外人講，凡事就用比喻， \end{tabularx} \\ \\ \relax
4:12 & \begin{tabularx}{0.7\textwidth}{X} 要他們看了又看，卻看不清，聽了又聽，卻不明白，免得他們回轉過來，獲得赦免。」 \end{tabularx} \\ \\ \relax
4:13 & \begin{tabularx}{0.7\textwidth}{X} 耶穌又對他們說：「你們不明白這比喻嗎？這樣怎能明白一切的比喻呢？ \end{tabularx} \\ \\ \relax
4:14 & \begin{tabularx}{0.7\textwidth}{X} 撒種的人所撒的就是道。 \end{tabularx} \\ \\ \relax
4:15 & \begin{tabularx}{0.7\textwidth}{X} 那撒在路旁的種子，就是人聽了道，撒但立刻來，把撒在他們心裡的道奪了去。 \end{tabularx} \\ \\ \relax
4:16 & \begin{tabularx}{0.7\textwidth}{X} 那撒在石頭地上的，就是人聽了道，立刻歡喜領受， \end{tabularx} \\ \\ \relax
4:17 & \begin{tabularx}{0.7\textwidth}{X} 因心裡沒有根，不過是暫時的，一旦為道遭受患難或迫害，立刻就跌倒。 \end{tabularx} \\ \\ \relax
4:18 & \begin{tabularx}{0.7\textwidth}{X} 還有那撒在荊棘裡的，就是人聽了道， \end{tabularx} \\ \\ \relax
4:19 & \begin{tabularx}{0.7\textwidth}{X} 後來有世上的憂慮、錢財的迷惑，和別樣的私慾進來，把道擠住了，結不出果實。 \end{tabularx} \\ \\ \relax
4:20 & \begin{tabularx}{0.7\textwidth}{X} 那撒在好土裡的，就是人聽了道，領受了，並且結了果實，有三十倍的，有六十倍的，有一百倍的。」 \end{tabularx} \\ \\
[1ex]
\hline
\hline
\end{longtable}


\newpage

\section{第65屆~港九培靈會~褚永華~第三講}
\label{sec:1003}
\textbf{財主與拉撒路的比喻}
\newline
\newline
連結: \href{https://www.hkbibleconference.org/session-message/view/1003}{\texttt{https://www.hkbibleconference.org/session-message/view/1003}}
\newline
\newline
\hyperref[sec:1002]{< < < PREV SERMON < < <}
~
\hyperref[sec:toc]{[返目錄]}
~
\hyperref[sec:1004]{> > > NEXT SERMON > > >}
\newline
\newline
路加福音 16:19-16:31
\newline
\begin{longtable}{cl}
\hline
\hline
章節 & 經文 (和合本修訂版)\\
\hline
16:19 & \begin{tabularx}{0.7\textwidth}{X} 「有一個財主穿著紫色袍和細麻布衣服，天天奢華宴樂。 \end{tabularx} \\ \\ \relax
16:20 & \begin{tabularx}{0.7\textwidth}{X} 又有一個討飯的，名叫拉撒路，渾身長瘡，被人放在財主門口， \end{tabularx} \\ \\ \relax
16:21 & \begin{tabularx}{0.7\textwidth}{X} 想得財主桌子上掉下來的碎食充飢，甚至還有狗來舔他的瘡。 \end{tabularx} \\ \\ \relax
16:22 & \begin{tabularx}{0.7\textwidth}{X} 後來那討飯的死了，被天使帶去放在亞伯拉罕的懷裡。財主也死了，並且埋葬了。 \end{tabularx} \\ \\ \relax
16:23 & \begin{tabularx}{0.7\textwidth}{X} 他在陰間受苦，舉目遠遠地望見亞伯拉罕，又望見拉撒路在他懷裡， \end{tabularx} \\ \\ \relax
16:24 & \begin{tabularx}{0.7\textwidth}{X} 他就喊著說：『我祖亞伯拉罕哪，可憐我吧！請打發拉撒路來，用指頭尖蘸點水，涼涼我的舌頭，因為我在這火焰裡，極其痛苦。』 \end{tabularx} \\ \\ \relax
16:25 & \begin{tabularx}{0.7\textwidth}{X} 亞伯拉罕說：『孩子啊，你該回想你生前享過福，拉撒路也同樣受過苦，如今他在這裡得安慰，你卻受痛苦。 \end{tabularx} \\ \\ \relax
16:26 & \begin{tabularx}{0.7\textwidth}{X} 除此之外，在你們和我們之間，有深淵隔開，以致人要從這邊過到你們那邊是不可能的；要從那邊過到這邊也是不可能的。』 \end{tabularx} \\ \\ \relax
16:27 & \begin{tabularx}{0.7\textwidth}{X} 財主說：『我祖啊，既然這樣，求你打發拉撒路到我父家去， \end{tabularx} \\ \\ \relax
16:28 & \begin{tabularx}{0.7\textwidth}{X} 因為我還有五個兄弟，他可以警告他們，免得他們也來到這痛苦的地方。』 \end{tabularx} \\ \\ \relax
16:29 & \begin{tabularx}{0.7\textwidth}{X} 亞伯拉罕說：『他們有摩西和先知的話可以聽從。』 \end{tabularx} \\ \\ \relax
16:30 & \begin{tabularx}{0.7\textwidth}{X} 他說：『不！我祖亞伯拉罕哪，假如有一個人從死人中到他們那裡去，他們一定會悔改。』 \end{tabularx} \\ \\ \relax
16:31 & \begin{tabularx}{0.7\textwidth}{X} 亞伯拉罕對他說：『如果他們不聽從摩西和先知的話，就是有人從死人中復活，他們也不會信服的。』」 \end{tabularx} \\ \\
[1ex]
\hline
\hline
\end{longtable}


\newpage

\section{第65屆~港九培靈會~褚永華~第四講}
\label{sec:1004}
\textbf{好撒瑪利亞人的比喻}
\newline
\newline
連結: \href{https://www.hkbibleconference.org/session-message/view/1004}{\texttt{https://www.hkbibleconference.org/session-message/view/1004}}
\newline
\newline
\hyperref[sec:1003]{< < < PREV SERMON < < <}
~
\hyperref[sec:toc]{[返目錄]}
~
\hyperref[sec:1005]{> > > NEXT SERMON > > >}
\newline
\newline
路加福音 10:25-10:37
\newline
\begin{longtable}{cl}
\hline
\hline
章節 & 經文 (和合本修訂版)\\
\hline
10:25 & \begin{tabularx}{0.7\textwidth}{X} 有一個律法師起來試探耶穌，說：「老師！我該做甚麼才可以承受永生？」 \end{tabularx} \\ \\ \relax
10:26 & \begin{tabularx}{0.7\textwidth}{X} 耶穌對他說：「律法上寫的是甚麼？你是怎樣念的呢？」 \end{tabularx} \\ \\ \relax
10:27 & \begin{tabularx}{0.7\textwidth}{X} 他回答說：「你要盡心、盡性、盡力、盡意愛主—你的神，又要愛鄰如己。」 \end{tabularx} \\ \\ \relax
10:28 & \begin{tabularx}{0.7\textwidth}{X} 耶穌對他說：「你回答得正確，你這樣做就會得永生。」 \end{tabularx} \\ \\ \relax
10:29 & \begin{tabularx}{0.7\textwidth}{X} 那人要證明自己有理，就對耶穌說：「誰是我的鄰舍呢？」 \end{tabularx} \\ \\ \relax
10:30 & \begin{tabularx}{0.7\textwidth}{X} 耶穌回答：「有一個人從耶路撒冷下耶利哥去，落在強盜手中。他們剝去他的衣裳，把他打個半死，丟下他走了。 \end{tabularx} \\ \\ \relax
10:31 & \begin{tabularx}{0.7\textwidth}{X} 偶然有一個祭司從那條路下來，看見他就從另一邊過去了。 \end{tabularx} \\ \\ \relax
10:32 & \begin{tabularx}{0.7\textwidth}{X} 又有一個利未人來到那裡，看見他，也照樣從另一邊過去了。 \end{tabularx} \\ \\ \relax
10:33 & \begin{tabularx}{0.7\textwidth}{X} 可是，有一個撒瑪利亞人路過那裡，看見他就動了慈心， \end{tabularx} \\ \\ \relax
10:34 & \begin{tabularx}{0.7\textwidth}{X} 上前用油和酒倒在他的傷處，包裹好了，扶他騎上自己的牲口，帶他到旅店裡去，照應他。 \end{tabularx} \\ \\ \relax
10:35 & \begin{tabularx}{0.7\textwidth}{X} 第二天，他拿出兩個銀幣來，交給店主，說：『請你照應他，額外的費用，我回來時會還你。』 \end{tabularx} \\ \\ \relax
10:36 & \begin{tabularx}{0.7\textwidth}{X} 你想，這三個人哪一個是落在強盜手中那人的鄰舍呢？」 \end{tabularx} \\ \\ \relax
10:37 & \begin{tabularx}{0.7\textwidth}{X} 他說：「是憐憫他的。」耶穌對他說：「你去，照樣做吧！」 \end{tabularx} \\ \\
[1ex]
\hline
\hline
\end{longtable}


\newpage

\section{第65屆~港九培靈會~褚永華~第五講}
\label{sec:1005}
\textbf{十個童女的比喻}
\newline
\newline
連結: \href{https://www.hkbibleconference.org/session-message/view/1005}{\texttt{https://www.hkbibleconference.org/session-message/view/1005}}
\newline
\newline
\hyperref[sec:1004]{< < < PREV SERMON < < <}
~
\hyperref[sec:toc]{[返目錄]}
~
\hyperref[sec:1006]{> > > NEXT SERMON > > >}
\newline
\newline
馬太福音 25:1-25:13
\newline
\begin{longtable}{cl}
\hline
\hline
章節 & 經文 (和合本修訂版)\\
\hline
25:1 & \begin{tabularx}{0.7\textwidth}{X} 「那時，天國好比十個童女拿著燈出去迎接新郎。 \end{tabularx} \\ \\ \relax
25:2 & \begin{tabularx}{0.7\textwidth}{X} 其中有五個是愚拙的，五個是聰明的。 \end{tabularx} \\ \\ \relax
25:3 & \begin{tabularx}{0.7\textwidth}{X} 愚拙的拿著燈，卻沒有帶油； \end{tabularx} \\ \\ \relax
25:4 & \begin{tabularx}{0.7\textwidth}{X} 聰明的拿著燈，又盛了油在器皿裡。 \end{tabularx} \\ \\ \relax
25:5 & \begin{tabularx}{0.7\textwidth}{X} 新郎遲延的時候，她們都打盹，睡著了。 \end{tabularx} \\ \\ \relax
25:6 & \begin{tabularx}{0.7\textwidth}{X} 半夜有人喊：『看，新郎來了，你們出來迎接他。』 \end{tabularx} \\ \\ \relax
25:7 & \begin{tabularx}{0.7\textwidth}{X} 那些童女就都起來挑亮她們的燈。 \end{tabularx} \\ \\ \relax
25:8 & \begin{tabularx}{0.7\textwidth}{X} 愚拙的對聰明的說：『請分點油給我們，因為我們的燈要滅了。』 \end{tabularx} \\ \\ \relax
25:9 & \begin{tabularx}{0.7\textwidth}{X} 聰明的回答：『恐怕不夠你我用的；你們還是自己到賣油的那裡去買吧。』 \end{tabularx} \\ \\ \relax
25:10 & \begin{tabularx}{0.7\textwidth}{X} 她們去買的時候，新郎到了。那預備好了的，與他進去共赴婚宴，門就關了。 \end{tabularx} \\ \\ \relax
25:11 & \begin{tabularx}{0.7\textwidth}{X} 其餘的童女隨後也來了，說：『主啊，主啊，給我們開門！』 \end{tabularx} \\ \\ \relax
25:12 & \begin{tabularx}{0.7\textwidth}{X} 他卻回答：『我實在告訴你們，我不認識你們。』 \end{tabularx} \\ \\ \relax
25:13 & \begin{tabularx}{0.7\textwidth}{X} 所以，你們要警醒，因為那日子，那時辰，你們不知道。」 \end{tabularx} \\ \\
[1ex]
\hline
\hline
\end{longtable}


\newpage

\section{第65屆~港九培靈會~褚永華~第六講}
\label{sec:1006}
\textbf{交銀僕人的比喻}
\newline
\newline
連結: \href{https://www.hkbibleconference.org/session-message/view/1006}{\texttt{https://www.hkbibleconference.org/session-message/view/1006}}
\newline
\newline
\hyperref[sec:1005]{< < < PREV SERMON < < <}
~
\hyperref[sec:toc]{[返目錄]}
~
\hyperref[sec:1007]{> > > NEXT SERMON > > >}
\newline
\newline
馬太福音 25:14-25:30
\newline
\begin{longtable}{cl}
\hline
\hline
章節 & 經文 (和合本修訂版)\\
\hline
25:14 & \begin{tabularx}{0.7\textwidth}{X} 「天國又好比一個人要出外遠行，就叫了僕人來，把他的家業交給他們。 \end{tabularx} \\ \\ \relax
25:15 & \begin{tabularx}{0.7\textwidth}{X} 他按著各人的才幹，給他們銀子：一個給了五千，一個給了二千，一個給了一千，就出外遠行去了。 \end{tabularx} \\ \\ \relax
25:16 & \begin{tabularx}{0.7\textwidth}{X} 那領五千的立刻拿去做買賣，另外賺了五千。 \end{tabularx} \\ \\ \relax
25:17 & \begin{tabularx}{0.7\textwidth}{X} 那領二千的也照樣另賺了二千。 \end{tabularx} \\ \\ \relax
25:18 & \begin{tabularx}{0.7\textwidth}{X} 但那領一千的去掘開地，把主人的銀子埋藏了。 \end{tabularx} \\ \\ \relax
25:19 & \begin{tabularx}{0.7\textwidth}{X} 過了許久，那些僕人的主人來了，和他們算賬。 \end{tabularx} \\ \\ \relax
25:20 & \begin{tabularx}{0.7\textwidth}{X} 那領五千的又帶著另外的五千來，說：『主啊，你交給我五千。請看，我又賺了五千。』 \end{tabularx} \\ \\ \relax
25:21 & \begin{tabularx}{0.7\textwidth}{X} 主人說：『好，你這又善良又忠心的僕人，你在少許的事上忠心，我要派你管理許多的事，進來享受你主人的快樂吧！』 \end{tabularx} \\ \\ \relax
25:22 & \begin{tabularx}{0.7\textwidth}{X} 那領二千的也進前來，說：『主啊，你交給我二千。請看，我又賺了二千。』 \end{tabularx} \\ \\ \relax
25:23 & \begin{tabularx}{0.7\textwidth}{X} 主人說：『好，你這又善良又忠心的僕人，你在少許的事上忠心，我要派你管理許多的事，進來享受你主人的快樂吧！』 \end{tabularx} \\ \\ \relax
25:24 & \begin{tabularx}{0.7\textwidth}{X} 那領一千的也進前來，說：『主啊，我知道你，你是個嚴厲的人：沒有種的地方也要收割，沒有播的地方也要收穫， \end{tabularx} \\ \\ \relax
25:25 & \begin{tabularx}{0.7\textwidth}{X} 我就害怕，去把你的一千銀子埋藏在地裡。請看，你的銀子在這裡。』 \end{tabularx} \\ \\ \relax
25:26 & \begin{tabularx}{0.7\textwidth}{X} 他的主人回答他說：『你這又惡又懶的僕人，你既知道我沒有種的地方也要收割，沒有播的地方也要收穫， \end{tabularx} \\ \\ \relax
25:27 & \begin{tabularx}{0.7\textwidth}{X} 就該把我的銀子放給兌換銀錢的人，到我來的時候可以連本帶利收回。 \end{tabularx} \\ \\ \relax
25:28 & \begin{tabularx}{0.7\textwidth}{X} 把他這一千奪過來，給那有一萬的。 \end{tabularx} \\ \\ \relax
25:29 & \begin{tabularx}{0.7\textwidth}{X} 因為凡有的，還要加給他，叫他有餘；沒有的，連他所有的也要奪過來。 \end{tabularx} \\ \\ \relax
25:30 & \begin{tabularx}{0.7\textwidth}{X} 把這無用的僕人丟在外面黑暗裡，在那裡他要哀哭切齒了。』」 \end{tabularx} \\ \\
[1ex]
\hline
\hline
\end{longtable}


\newpage

\section{第65屆~港九培靈會~褚永華~第七講}
\label{sec:1007}
\textbf{無花果樹的比喻}
\newline
\newline
連結: \href{https://www.hkbibleconference.org/session-message/view/1007}{\texttt{https://www.hkbibleconference.org/session-message/view/1007}}
\newline
\newline
\hyperref[sec:1006]{< < < PREV SERMON < < <}
~
\hyperref[sec:toc]{[返目錄]}
~
\hyperref[sec:1008]{> > > NEXT SERMON > > >}
\newline
\newline
馬可福音 11:12-11:21
\newline
\begin{longtable}{cl}
\hline
\hline
章節 & 經文 (和合本修訂版)\\
\hline
11:12 & \begin{tabularx}{0.7\textwidth}{X} 第二天，他們從伯大尼出來，耶穌餓了。 \end{tabularx} \\ \\ \relax
11:13 & \begin{tabularx}{0.7\textwidth}{X} 他遠遠地看見一棵無花果樹，樹上有葉子，就過去，看是不是在樹上可以找到甚麼。他到了樹下，竟找不到甚麼，只有葉子，因為不是無花果的季節。 \end{tabularx} \\ \\ \relax
11:14 & \begin{tabularx}{0.7\textwidth}{X} 耶穌就對樹說：「從今以後，永沒有人吃你的果子。」他的門徒都聽到了。 \end{tabularx} \\ \\ \relax
11:15 & \begin{tabularx}{0.7\textwidth}{X} 他們來到耶路撒冷。耶穌一進聖殿，就趕出在聖殿裡做買賣的人，推倒兌換銀錢之人的桌子和賣鴿子之人的凳子； \end{tabularx} \\ \\ \relax
11:16 & \begin{tabularx}{0.7\textwidth}{X} 也不許人拿著器具從聖殿裡經過。 \end{tabularx} \\ \\ \relax
11:17 & \begin{tabularx}{0.7\textwidth}{X} 他教導他們說：「經上不是記著：『我的殿要稱為萬國禱告的殿嗎？你們倒使它成為賊窩了。』」 \end{tabularx} \\ \\ \relax
11:18 & \begin{tabularx}{0.7\textwidth}{X} 祭司長和文士聽見這話，就想法子要除掉耶穌，卻又怕他，因為眾人都對他的教導感到驚奇。 \end{tabularx} \\ \\ \relax
11:19 & \begin{tabularx}{0.7\textwidth}{X} 每天晚上，他們都到城外去。 \end{tabularx} \\ \\ \relax
11:20 & \begin{tabularx}{0.7\textwidth}{X} 早晨，他們從那裡經過，看見無花果樹連根都枯乾了。 \end{tabularx} \\ \\ \relax
11:21 & \begin{tabularx}{0.7\textwidth}{X} 彼得想起耶穌的話來，就對他說：「拉比，你看！你所詛咒的無花果樹已經枯乾了。」 \end{tabularx} \\ \\
[1ex]
\hline
\hline
\end{longtable}


\newpage

\section{第65屆~港九培靈會~褚永華~第八講}
\label{sec:1008}
\textbf{葡萄園的比喻}
\newline
\newline
連結: \href{https://www.hkbibleconference.org/session-message/view/1008}{\texttt{https://www.hkbibleconference.org/session-message/view/1008}}
\newline
\newline
\hyperref[sec:1007]{< < < PREV SERMON < < <}
~
\hyperref[sec:toc]{[返目錄]}
~
\hyperref[sec:973]{> > > NEXT SERMON > > >}
\newline
\newline
馬太福音 20:1-20:16
\newline
\begin{longtable}{cl}
\hline
\hline
章節 & 經文 (和合本修訂版)\\
\hline
20:1 & \begin{tabularx}{0.7\textwidth}{X} 「因為天國好比一家的主人清早去雇人進他的葡萄園做工。 \end{tabularx} \\ \\ \relax
20:2 & \begin{tabularx}{0.7\textwidth}{X} 他和工人講定一天一個銀幣，就打發他們進葡萄園去。 \end{tabularx} \\ \\ \relax
20:3 & \begin{tabularx}{0.7\textwidth}{X} 約在上午九點鐘出去，看見市場上還有閒站的人， \end{tabularx} \\ \\ \relax
20:4 & \begin{tabularx}{0.7\textwidth}{X} 就對那些人說：『你們也進葡萄園去，我會給你們合理的工錢。』 \end{tabularx} \\ \\ \relax
20:5 & \begin{tabularx}{0.7\textwidth}{X} 他們也進去了。約在正午和下午三點鐘又出去，他也是這麼做。 \end{tabularx} \\ \\ \relax
20:6 & \begin{tabularx}{0.7\textwidth}{X} 約在下午五點鐘出去，他看見還有人站在那裡，就問他們：『你們為甚麼整天在這裡閒站呢？』 \end{tabularx} \\ \\ \relax
20:7 & \begin{tabularx}{0.7\textwidth}{X} 他們說：『因為沒有人雇我們。』他說：『你們也進葡萄園去。』 \end{tabularx} \\ \\ \relax
20:8 & \begin{tabularx}{0.7\textwidth}{X} 到了晚上，園主對工頭說：『叫工人都來，給他們工錢，從後來的起，到先來的為止。』 \end{tabularx} \\ \\ \relax
20:9 & \begin{tabularx}{0.7\textwidth}{X} 約在下午五點鐘雇的人來了，各人領了一個銀幣。 \end{tabularx} \\ \\ \relax
20:10 & \begin{tabularx}{0.7\textwidth}{X} 那些最先雇的來了，以為可以多領，誰知也是各領一個銀幣。 \end{tabularx} \\ \\ \relax
20:11 & \begin{tabularx}{0.7\textwidth}{X} 他們領了工錢，就埋怨那家的主人說： \end{tabularx} \\ \\ \relax
20:12 & \begin{tabularx}{0.7\textwidth}{X} 『我們整天勞苦受熱，那些後來的只做了一小時，你竟待他們和我們一樣嗎？』 \end{tabularx} \\ \\ \relax
20:13 & \begin{tabularx}{0.7\textwidth}{X} 主人回答其中的一人說：『朋友，我沒虧待你，你與我講定的不是一個銀幣嗎？ \end{tabularx} \\ \\ \relax
20:14 & \begin{tabularx}{0.7\textwidth}{X} 拿你的錢走吧！我樂意給那後來的和給你的一樣， \end{tabularx} \\ \\ \relax
20:15 & \begin{tabularx}{0.7\textwidth}{X} 難道我的東西不可隨我的意思用嗎？因為我作好人，你就眼紅了嗎？』 \end{tabularx} \\ \\ \relax
20:16 & \begin{tabularx}{0.7\textwidth}{X} 這樣，那在後的，將要在前；在前的，將要在後了。」 \end{tabularx} \\ \\
[1ex]
\hline
\hline
\end{longtable}


\newpage

\section{第66屆~港九培靈會~沈保羅~第一講}
\label{sec:973}
\textbf{祂到底是誰}
\newline
\newline
連結: \href{https://www.hkbibleconference.org/session-message/view/973}{\texttt{https://www.hkbibleconference.org/session-message/view/973}}
\newline
\newline
\hyperref[sec:1008]{< < < PREV SERMON < < <}
~
\hyperref[sec:toc]{[返目錄]}
~
\hyperref[sec:974]{> > > NEXT SERMON > > >}
\newline
\newline


\newpage

\section{第66屆~港九培靈會~沈保羅~第二講}
\label{sec:974}
\textbf{雙重的考慮}
\newline
\newline
連結: \href{https://www.hkbibleconference.org/session-message/view/974}{\texttt{https://www.hkbibleconference.org/session-message/view/974}}
\newline
\newline
\hyperref[sec:973]{< < < PREV SERMON < < <}
~
\hyperref[sec:toc]{[返目錄]}
~
\hyperref[sec:975]{> > > NEXT SERMON > > >}
\newline
\newline


\newpage

\section{第66屆~港九培靈會~沈保羅~第三講}
\label{sec:975}
\textbf{死亡,苦難,害怕}
\newline
\newline
連結: \href{https://www.hkbibleconference.org/session-message/view/975}{\texttt{https://www.hkbibleconference.org/session-message/view/975}}
\newline
\newline
\hyperref[sec:974]{< < < PREV SERMON < < <}
~
\hyperref[sec:toc]{[返目錄]}
~
\hyperref[sec:976]{> > > NEXT SERMON > > >}
\newline
\newline


\newpage

\section{第66屆~港九培靈會~沈保羅~第四講}
\label{sec:976}
\textbf{那時 : 誤會,挫折,拒絕}
\newline
\newline
連結: \href{https://www.hkbibleconference.org/session-message/view/976}{\texttt{https://www.hkbibleconference.org/session-message/view/976}}
\newline
\newline
\hyperref[sec:975]{< < < PREV SERMON < < <}
~
\hyperref[sec:toc]{[返目錄]}
~
\hyperref[sec:977]{> > > NEXT SERMON > > >}
\newline
\newline


\newpage

\section{第66屆~港九培靈會~沈保羅~第五講}
\label{sec:977}
\textbf{這是我的愛子}
\newline
\newline
連結: \href{https://www.hkbibleconference.org/session-message/view/977}{\texttt{https://www.hkbibleconference.org/session-message/view/977}}
\newline
\newline
\hyperref[sec:976]{< < < PREV SERMON < < <}
~
\hyperref[sec:toc]{[返目錄]}
~
\hyperref[sec:978]{> > > NEXT SERMON > > >}
\newline
\newline


\newpage

\section{第66屆~港九培靈會~沈保羅~第六講}
\label{sec:978}
\textbf{祂的人生}
\newline
\newline
連結: \href{https://www.hkbibleconference.org/session-message/view/978}{\texttt{https://www.hkbibleconference.org/session-message/view/978}}
\newline
\newline
\hyperref[sec:977]{< < < PREV SERMON < < <}
~
\hyperref[sec:toc]{[返目錄]}
~
\hyperref[sec:979]{> > > NEXT SERMON > > >}
\newline
\newline


\newpage

\section{第66屆~港九培靈會~沈保羅~第七講}
\label{sec:979}
\textbf{在那裡吃}
\newline
\newline
連結: \href{https://www.hkbibleconference.org/session-message/view/979}{\texttt{https://www.hkbibleconference.org/session-message/view/979}}
\newline
\newline
\hyperref[sec:978]{< < < PREV SERMON < < <}
~
\hyperref[sec:toc]{[返目錄]}
~
\hyperref[sec:980]{> > > NEXT SERMON > > >}
\newline
\newline


\newpage

\section{第66屆~港九培靈會~沈保羅~第八講}
\label{sec:980}
\textbf{戰勝驚恐}
\newline
\newline
連結: \href{https://www.hkbibleconference.org/session-message/view/980}{\texttt{https://www.hkbibleconference.org/session-message/view/980}}
\newline
\newline
\hyperref[sec:979]{< < < PREV SERMON < < <}
~
\hyperref[sec:toc]{[返目錄]}
~
\hyperref[sec:981]{> > > NEXT SERMON > > >}
\newline
\newline


\newpage

\section{第66屆~港九培靈會~沈保羅~第九講}
\label{sec:981}
\textbf{地平線的那一邊}
\newline
\newline
連結: \href{https://www.hkbibleconference.org/session-message/view/981}{\texttt{https://www.hkbibleconference.org/session-message/view/981}}
\newline
\newline
\hyperref[sec:980]{< < < PREV SERMON < < <}
~
\hyperref[sec:toc]{[返目錄]}
~
\hyperref[sec:982]{> > > NEXT SERMON > > >}
\newline
\newline


\newpage

\section{第66屆~港九培靈會~沈保羅~第十講}
\label{sec:982}
\textbf{愛的比較}
\newline
\newline
連結: \href{https://www.hkbibleconference.org/session-message/view/982}{\texttt{https://www.hkbibleconference.org/session-message/view/982}}
\newline
\newline
\hyperref[sec:981]{< < < PREV SERMON < < <}
~
\hyperref[sec:toc]{[返目錄]}
~
\hyperref[sec:954]{> > > NEXT SERMON > > >}
\newline
\newline


\newpage

\section{第66屆~港九培靈會~蘇穎智~第一講}
\label{sec:954}
\textbf{在地如在天的敬拜}
\newline
\newline
連結: \href{https://www.hkbibleconference.org/session-message/view/954}{\texttt{https://www.hkbibleconference.org/session-message/view/954}}
\newline
\newline
\hyperref[sec:982]{< < < PREV SERMON < < <}
~
\hyperref[sec:toc]{[返目錄]}
~
\hyperref[sec:955]{> > > NEXT SERMON > > >}
\newline
\newline
啟示錄 4:1-5:14
\newline
\begin{longtable}{cl}
\hline
\hline
章節 & 經文 (和合本修訂版)\\
\hline
4:1 & \begin{tabularx}{0.7\textwidth}{X} 這些事以後，我觀看，看見天上有一道門開著。我頭一次聽見的那好像吹號的聲音對我說：「你上這裡來，我要把此後必須發生的事指示你。」 \end{tabularx} \\ \\ \relax
4:2 & \begin{tabularx}{0.7\textwidth}{X} 我立刻被聖靈感動，見有一個寶座安置在天上，有一位坐在寶座上。 \end{tabularx} \\ \\ \relax
4:3 & \begin{tabularx}{0.7\textwidth}{X} 那坐著的，看來好像碧玉和紅寶石；又有彩虹圍著寶座，光彩好像綠寶石。 \end{tabularx} \\ \\ \relax
4:4 & \begin{tabularx}{0.7\textwidth}{X} 寶座的周圍又有二十四個座位，上面坐著二十四位長老，身穿白衣，頭上戴著金冠冕。 \end{tabularx} \\ \\ \relax
4:5 & \begin{tabularx}{0.7\textwidth}{X} 有閃電、聲音、雷轟從寶座中發出。在寶座前點著七支火炬，就是神的七靈。 \end{tabularx} \\ \\ \relax
4:6 & \begin{tabularx}{0.7\textwidth}{X} 寶座前有一個如同水晶的玻璃海。寶座的周圍，四邊有四個活物，遍體前後都長滿了眼睛。 \end{tabularx} \\ \\ \relax
4:7 & \begin{tabularx}{0.7\textwidth}{X} 第一個活物像獅子，第二個像牛犢，第三個的臉像人臉，第四個像飛鷹。 \end{tabularx} \\ \\ \relax
4:8 & \begin{tabularx}{0.7\textwidth}{X} 四個活物各有六個翅膀，遍體內外都長滿了眼睛。他們晝夜不住地說：「聖哉！聖哉！聖哉！主—全能的神；昔在、今在、以後永在！」 \end{tabularx} \\ \\ \relax
4:9 & \begin{tabularx}{0.7\textwidth}{X} 每逢四活物將榮耀、尊貴、感謝歸給那坐在寶座上、活到永永遠遠者的時候， \end{tabularx} \\ \\ \relax
4:10 & \begin{tabularx}{0.7\textwidth}{X} 二十四位長老就俯伏敬拜坐在寶座上活到永永遠遠的那一位，又把他們的冠冕放在寶座前，說： \end{tabularx} \\ \\ \relax
4:11 & \begin{tabularx}{0.7\textwidth}{X} 「我們的主，我們的神，你配得榮耀、尊貴、權柄，因為你創造了萬物，萬物因你的旨意被創造而存在。」 \end{tabularx} \\ \\ \relax
5:1 & \begin{tabularx}{0.7\textwidth}{X} 我看見坐在寶座那位的右手中有書卷，正反面都寫著字，用七個印密封著。 \end{tabularx} \\ \\ \relax
5:2 & \begin{tabularx}{0.7\textwidth}{X} 我又看見一位大力的天使大聲宣告說：「有誰配展開那書卷，揭開那七個印呢？」 \end{tabularx} \\ \\ \relax
5:3 & \begin{tabularx}{0.7\textwidth}{X} 在天上、地上、地底下，沒有人能展開、能閱覽那書卷。 \end{tabularx} \\ \\ \relax
5:4 & \begin{tabularx}{0.7\textwidth}{X} 因為沒有人配展開、閱覽那書卷，我就大哭。 \end{tabularx} \\ \\ \relax
5:5 & \begin{tabularx}{0.7\textwidth}{X} 長老中有一位對我說：「不要哭。看哪，猶大支派中的獅子，大衛的根，他已得勝，能展開那書卷，揭開那七個印。」 \end{tabularx} \\ \\ \relax
5:6 & \begin{tabularx}{0.7\textwidth}{X} 我又看見寶座和四個活物，以及長老之中有羔羊站著，像是被殺的，有七個角七隻眼睛，就是神的七靈，奉差遣往普天下去的。 \end{tabularx} \\ \\ \relax
5:7 & \begin{tabularx}{0.7\textwidth}{X} 這羔羊前來，從坐在寶座上那位的右手中拿了書卷。 \end{tabularx} \\ \\ \relax
5:8 & \begin{tabularx}{0.7\textwidth}{X} 他一拿了書卷，四活物和二十四位長老就俯伏在羔羊面前，各拿著琴和盛滿了香的金爐；這香就是眾聖徒的祈禱。 \end{tabularx} \\ \\ \relax
5:9 & \begin{tabularx}{0.7\textwidth}{X} 他們唱新歌，說：「你配拿書卷，配揭開它的七印；因為你曾被殺，用自己的血從各支派、各語言、各民族、各邦國中買了人來，使他們歸於神， \end{tabularx} \\ \\ \relax
5:10 & \begin{tabularx}{0.7\textwidth}{X} 又使他們成為國民和祭司，歸於我們的神；他們將在地上執掌王權。」 \end{tabularx} \\ \\ \relax
5:11 & \begin{tabularx}{0.7\textwidth}{X} 我又觀看，我聽見寶座和活物及長老的周圍有許多天使的聲音；他們的數目有千千萬萬， \end{tabularx} \\ \\ \relax
5:12 & \begin{tabularx}{0.7\textwidth}{X} 大聲說：「被殺的羔羊配得權能、豐富、智慧、力量、尊貴、榮耀、頌讚。」 \end{tabularx} \\ \\ \relax
5:13 & \begin{tabularx}{0.7\textwidth}{X} 我又聽見在天上、地上、地底下、滄海裡和天地間一切所有被造之物，都說：「願頌讚、尊貴、榮耀、權勢，都歸給坐在寶座上的那位和羔羊，直到永永遠遠！」 \end{tabularx} \\ \\ \relax
5:14 & \begin{tabularx}{0.7\textwidth}{X} 四活物就說：「阿們！」眾長老也俯伏敬拜。 \end{tabularx} \\ \\
[1ex]
\hline
\hline
\end{longtable}


\newpage

\section{第66屆~港九培靈會~蘇穎智~第二講}
\label{sec:955}
\textbf{浩劫當前 — 六印之災}
\newline
\newline
連結: \href{https://www.hkbibleconference.org/session-message/view/955}{\texttt{https://www.hkbibleconference.org/session-message/view/955}}
\newline
\newline
\hyperref[sec:954]{< < < PREV SERMON < < <}
~
\hyperref[sec:toc]{[返目錄]}
~
\hyperref[sec:956]{> > > NEXT SERMON > > >}
\newline
\newline
啟示錄 6:1-7:17
\newline
\begin{longtable}{cl}
\hline
\hline
章節 & 經文 (和合本修訂版)\\
\hline
6:1 & \begin{tabularx}{0.7\textwidth}{X} 我看見羔羊揭開七個印中第一個印的時候，聽見四活物中的一個活物，聲音如雷，說：「你來！」 \end{tabularx} \\ \\ \relax
6:2 & \begin{tabularx}{0.7\textwidth}{X} 我就觀看，看見一匹白馬，騎在馬上的拿著弓，並有冠冕賜給他。他出來征服，勝而又勝。 \end{tabularx} \\ \\ \relax
6:3 & \begin{tabularx}{0.7\textwidth}{X} 羔羊揭開第二個印的時候，我聽見第二個活物說：「你來！」 \end{tabularx} \\ \\ \relax
6:4 & \begin{tabularx}{0.7\textwidth}{X} 就另有一匹馬出來，是紅色的；有權柄賜給了那騎馬的，要從地上奪去太平，使人彼此相殺；他又接受了一把大刀。 \end{tabularx} \\ \\ \relax
6:5 & \begin{tabularx}{0.7\textwidth}{X} 羔羊揭開第三個印的時候，我聽見第三個活物說：「你來！」我就觀看，看見一匹黑馬；騎在馬上的，手裡拿著天平。 \end{tabularx} \\ \\ \relax
6:6 & \begin{tabularx}{0.7\textwidth}{X} 我聽見在四個活物中似乎有聲音說：「一個銀幣買一升麥子，一個銀幣買三升大麥；油和酒不可糟蹋。」 \end{tabularx} \\ \\ \relax
6:7 & \begin{tabularx}{0.7\textwidth}{X} 羔羊揭開第四個印的時候，我聽見第四個活物說：「你來！」 \end{tabularx} \\ \\ \relax
6:8 & \begin{tabularx}{0.7\textwidth}{X} 我就觀看，看見一匹灰色馬；騎在馬上的，名字叫作「死」，陰間也隨著他；有權柄賜給他們，可以用刀劍、饑荒、瘟疫、野獸，殺害地上四分之一的人。 \end{tabularx} \\ \\ \relax
6:9 & \begin{tabularx}{0.7\textwidth}{X} 羔羊揭開第五個印的時候，我看見在祭壇底下有曾為神的道，並為作見證而被殺的人的靈魂， \end{tabularx} \\ \\ \relax
6:10 & \begin{tabularx}{0.7\textwidth}{X} 大聲喊著說：「神聖真實的主宰啊，你不審判住在地上的人，為我們所流的血伸冤，要到幾時呢？」 \end{tabularx} \\ \\ \relax
6:11 & \begin{tabularx}{0.7\textwidth}{X} 於是有白袍賜給他們各人；又有話吩咐他們還要歇息片刻，等到與他們同作僕人的，和他們的弟兄，像他們一樣被殺的人的數目湊足的時候。 \end{tabularx} \\ \\ \relax
6:12 & \begin{tabularx}{0.7\textwidth}{X} 羔羊揭開第六個印的時候，我看見地大震動，太陽變黑像粗麻布，整個月亮變紅像血， \end{tabularx} \\ \\ \relax
6:13 & \begin{tabularx}{0.7\textwidth}{X} 天上的星辰墜落在地上，如同無花果樹被大風搖動，落下未熟的果子一樣。 \end{tabularx} \\ \\ \relax
6:14 & \begin{tabularx}{0.7\textwidth}{X} 天就裂開，好像書卷被捲起來；山嶺海島都被移動離開原位。 \end{tabularx} \\ \\ \relax
6:15 & \begin{tabularx}{0.7\textwidth}{X} 地上的君王、臣宰、將軍、富戶、壯士，和一切為奴的、自主的，都藏在山洞和巖石穴裡， \end{tabularx} \\ \\ \relax
6:16 & \begin{tabularx}{0.7\textwidth}{X} 向山和巖石說：「倒在我們身上吧！把我們藏起來，躲避坐寶座者的臉面和羔羊的憤怒； \end{tabularx} \\ \\ \relax
6:17 & \begin{tabularx}{0.7\textwidth}{X} 因為他們遭憤怒的大日子到了，誰能站得住呢？」 \end{tabularx} \\ \\ \relax
7:1 & \begin{tabularx}{0.7\textwidth}{X} 此後，我看見四位天使站在地的四角，執掌地上四方的風，使風不吹在地上、海上和各種樹上。 \end{tabularx} \\ \\ \relax
7:2 & \begin{tabularx}{0.7\textwidth}{X} 我又看見另有一位天使從日出之地上來，拿著永生神的印。他向那得到權柄能傷害地和海的四位天使大聲喊著， \end{tabularx} \\ \\ \relax
7:3 & \begin{tabularx}{0.7\textwidth}{X} 說：「你們不可傷害地、海和樹林，等我們在我們神眾僕人的額上蓋了印。」 \end{tabularx} \\ \\ \relax
7:4 & \begin{tabularx}{0.7\textwidth}{X} 我聽見以色列人各支派中受印的數目有十四萬四千； \end{tabularx} \\ \\ \relax
7:5 & \begin{tabularx}{0.7\textwidth}{X} 猶大支派中受印的有一萬二千；呂便支派中有一萬二千；迦得支派中有一萬二千； \end{tabularx} \\ \\ \relax
7:6 & \begin{tabularx}{0.7\textwidth}{X} 亞設支派中有一萬二千；拿弗他利支派中有一萬二千；瑪拿西支派中有一萬二千； \end{tabularx} \\ \\ \relax
7:7 & \begin{tabularx}{0.7\textwidth}{X} 西緬支派中有一萬二千；利未支派中有一萬二千；以薩迦支派中有一萬二千； \end{tabularx} \\ \\ \relax
7:8 & \begin{tabularx}{0.7\textwidth}{X} 西布倫支派中有一萬二千；約瑟支派中有一萬二千；便雅憫支派中受印的有一萬二千。 \end{tabularx} \\ \\ \relax
7:9 & \begin{tabularx}{0.7\textwidth}{X} 此後，我觀看，看見有許多人，沒有人能計算，是從各邦國、各支派、各民族、各語言來的，站在寶座和羔羊面前，身穿白衣，手拿棕樹枝， \end{tabularx} \\ \\ \relax
7:10 & \begin{tabularx}{0.7\textwidth}{X} 大聲喊著說：「願救恩歸於坐在寶座上我們的神，也歸於羔羊！」 \end{tabularx} \\ \\ \relax
7:11 & \begin{tabularx}{0.7\textwidth}{X} 眾天使都站在寶座和眾長老，以及四個活物的周圍，俯伏在寶座前，敬拜神， \end{tabularx} \\ \\ \relax
7:12 & \begin{tabularx}{0.7\textwidth}{X} 說：「阿們！頌讚、榮耀、智慧、感謝、尊貴、權能、力量都歸於我們的神，直到永永遠遠。阿們！」 \end{tabularx} \\ \\ \relax
7:13 & \begin{tabularx}{0.7\textwidth}{X} 長老中有一位回應我說：「這些穿白衣的是誰？是從哪裡來的？」 \end{tabularx} \\ \\ \relax
7:14 & \begin{tabularx}{0.7\textwidth}{X} 我對他說：「我主啊，你是知道的。」他向我說：「這些人是從大患難中出來的，他們曾用羔羊的血把衣裳洗得潔白。 \end{tabularx} \\ \\ \relax
7:15 & \begin{tabularx}{0.7\textwidth}{X} 所以，他們在神寶座前，晝夜在他殿中事奉他；那坐在寶座上的要用帳幕覆庇他們。 \end{tabularx} \\ \\ \relax
7:16 & \begin{tabularx}{0.7\textwidth}{X} 他們不再飢，不再渴；太陽必不傷害他們，任何炎熱也不傷害他們， \end{tabularx} \\ \\ \relax
7:17 & \begin{tabularx}{0.7\textwidth}{X} 因為寶座中的羔羊必牧養他們，領他們到生命水的泉源；神必擦去他們一切的眼淚。」 \end{tabularx} \\ \\
[1ex]
\hline
\hline
\end{longtable}


\newpage

\section{第66屆~港九培靈會~蘇穎智~第三講}
\label{sec:956}
\textbf{禍不單行 — 六號之災}
\newline
\newline
連結: \href{https://www.hkbibleconference.org/session-message/view/956}{\texttt{https://www.hkbibleconference.org/session-message/view/956}}
\newline
\newline
\hyperref[sec:955]{< < < PREV SERMON < < <}
~
\hyperref[sec:toc]{[返目錄]}
~
\hyperref[sec:957]{> > > NEXT SERMON > > >}
\newline
\newline
啟示錄 8:1-9:21
\newline
\begin{longtable}{cl}
\hline
\hline
章節 & 經文 (和合本修訂版)\\
\hline
8:1 & \begin{tabularx}{0.7\textwidth}{X} 羔羊揭開第七個印的時候，天上寂靜約有半小時。 \end{tabularx} \\ \\ \relax
8:2 & \begin{tabularx}{0.7\textwidth}{X} 我看見那站在神面前的七位天使，有七枝號賜給他們。 \end{tabularx} \\ \\ \relax
8:3 & \begin{tabularx}{0.7\textwidth}{X} 另有一位天使拿著金香爐來，站在祭壇旁邊；有許多香賜給他，要和眾聖徒的祈禱一同獻在寶座前的金壇上。 \end{tabularx} \\ \\ \relax
8:4 & \begin{tabularx}{0.7\textwidth}{X} 那香的煙和眾聖徒的祈禱從天使的手中一同升到神面前。 \end{tabularx} \\ \\ \relax
8:5 & \begin{tabularx}{0.7\textwidth}{X} 天使拿著香爐，盛滿了壇上的火，倒在地上；就有雷轟、響聲、閃電、地震。 \end{tabularx} \\ \\ \relax
8:6 & \begin{tabularx}{0.7\textwidth}{X} 拿著七枝號筒的七位天使預備好要吹號。 \end{tabularx} \\ \\ \relax
8:7 & \begin{tabularx}{0.7\textwidth}{X} 第一位天使吹號，就有冰雹和火攙著血扔在地上；地的三分之一和樹的三分之一被燒掉了，一切的青草也被燒掉了。 \end{tabularx} \\ \\ \relax
8:8 & \begin{tabularx}{0.7\textwidth}{X} 第二位天使吹號，就有像火燒著的大山扔在海中；海的三分之一變成血， \end{tabularx} \\ \\ \relax
8:9 & \begin{tabularx}{0.7\textwidth}{X} 海中有生命的被造之物死了三分之一，船隻也毀壞了三分之一。 \end{tabularx} \\ \\ \relax
8:10 & \begin{tabularx}{0.7\textwidth}{X} 第三位天使吹號，就有燒著的大星好像火把從天上墜下來，落在江河的三分之一和眾水的泉源上。 \end{tabularx} \\ \\ \relax
8:11 & \begin{tabularx}{0.7\textwidth}{X} 這星名叫「苦艾」；眾水的三分之一變為苦艾，許多人因水變苦而死了。 \end{tabularx} \\ \\ \relax
8:12 & \begin{tabularx}{0.7\textwidth}{X} 第四位天使吹號，太陽的三分之一、月亮的三分之一、星辰的三分之一都被擊打，以致日月星的三分之一變黑了，白晝的三分之一沒有光，黑夜也是這樣。 \end{tabularx} \\ \\ \relax
8:13 & \begin{tabularx}{0.7\textwidth}{X} 我觀看，聽見一隻在空中飛的鷹大聲說：「禍哉！禍哉！禍哉！地上的居民哪，其餘的三位天使快要吹號了！」 \end{tabularx} \\ \\ \relax
9:1 & \begin{tabularx}{0.7\textwidth}{X} 第五位天使吹號，我就看見一顆星從天上墜落到地上；有無底坑的鑰匙賜給它。 \end{tabularx} \\ \\ \relax
9:2 & \begin{tabularx}{0.7\textwidth}{X} 它開了無底坑，就有煙從坑裡往上冒，好像大火爐的煙；太陽和天空都因這煙昏暗了。 \end{tabularx} \\ \\ \relax
9:3 & \begin{tabularx}{0.7\textwidth}{X} 有蝗蟲從煙中出來，飛到地上，有權柄賜給牠們，好像地上的蠍子有權柄一樣。 \end{tabularx} \\ \\ \relax
9:4 & \begin{tabularx}{0.7\textwidth}{X} 牠們奉命不可傷害地上的草、各樣綠色植物和各種樹木，惟獨可傷害額上沒有神印記的人； \end{tabularx} \\ \\ \relax
9:5 & \begin{tabularx}{0.7\textwidth}{X} 但是不許蝗蟲害死他們，只可使他們受痛苦五個月；這痛苦就像人被蠍子螫了的痛苦一樣。 \end{tabularx} \\ \\ \relax
9:6 & \begin{tabularx}{0.7\textwidth}{X} 在那些日子，人求死，卻死不了；想死，死卻避開他們。 \end{tabularx} \\ \\ \relax
9:7 & \begin{tabularx}{0.7\textwidth}{X} 蝗蟲的形狀好像預備上陣的戰馬一樣，頭上戴的好像金冠冕，臉面好像男人的臉面， \end{tabularx} \\ \\ \relax
9:8 & \begin{tabularx}{0.7\textwidth}{X} 頭髮像女人的頭髮，牙齒像獅子的牙齒； \end{tabularx} \\ \\ \relax
9:9 & \begin{tabularx}{0.7\textwidth}{X} 牠們胸前有甲，好像鐵甲；又有翅膀的響聲，好像許多車馬奔跑上陣的聲音。 \end{tabularx} \\ \\ \relax
9:10 & \begin{tabularx}{0.7\textwidth}{X} 牠們有尾巴像蠍子，長著毒刺，尾巴上的毒刺有能力傷害人五個月。 \end{tabularx} \\ \\ \relax
9:11 & \begin{tabularx}{0.7\textwidth}{X} 牠們有無底坑的使者作牠們的王，按著希伯來話名叫亞巴頓，希臘話名叫亞玻倫。 \end{tabularx} \\ \\ \relax
9:12 & \begin{tabularx}{0.7\textwidth}{X} 第一樣災禍過去了；看哪，還有兩樣災禍要來。 \end{tabularx} \\ \\ \relax
9:13 & \begin{tabularx}{0.7\textwidth}{X} 第六位天使吹號，我聽見有聲音從神面前金壇的四角發出來， \end{tabularx} \\ \\ \relax
9:14 & \begin{tabularx}{0.7\textwidth}{X} 吩咐那吹號的第六位天使，說：「把那捆綁在幼發拉底大河的四個使者釋放了。」 \end{tabularx} \\ \\ \relax
9:15 & \begin{tabularx}{0.7\textwidth}{X} 那四個使者就被釋放；他們原是預備好，在特定的年、月、日、時，要殺人類的三分之一。 \end{tabularx} \\ \\ \relax
9:16 & \begin{tabularx}{0.7\textwidth}{X} 騎兵有二億；他們的數目我聽見了。 \end{tabularx} \\ \\ \relax
9:17 & \begin{tabularx}{0.7\textwidth}{X} 我在異象中看見那些馬和騎馬的：騎馬的穿著火紅、紫瑪瑙及硫磺色的胸甲；馬的頭好像獅子的頭，有火、有煙、有硫磺從馬的口中噴出來。 \end{tabularx} \\ \\ \relax
9:18 & \begin{tabularx}{0.7\textwidth}{X} 從馬的口中所噴出來的火、煙和硫磺這三樣災害殺了人類的三分之一。 \end{tabularx} \\ \\ \relax
9:19 & \begin{tabularx}{0.7\textwidth}{X} 馬的能力在於牠們的口和尾巴；牠們的尾巴像蛇，有頭，用頭來傷害人。 \end{tabularx} \\ \\ \relax
9:20 & \begin{tabularx}{0.7\textwidth}{X} 其餘未曾被這些災難所殺的人仍舊不為自己手所做的悔改，還是去拜鬼魔和那些不能看、不能聽、不能走，用金、銀、銅、木、石所造的偶像。 \end{tabularx} \\ \\ \relax
9:21 & \begin{tabularx}{0.7\textwidth}{X} 他們也不為自己所犯的那些兇殺、邪術、淫亂、偷竊的事悔改。 \end{tabularx} \\ \\
[1ex]
\hline
\hline
\end{longtable}


\newpage

\section{第66屆~港九培靈會~蘇穎智~第四講}
\label{sec:957}
\textbf{能帶給人出路的人}
\newline
\newline
連結: \href{https://www.hkbibleconference.org/session-message/view/957}{\texttt{https://www.hkbibleconference.org/session-message/view/957}}
\newline
\newline
\hyperref[sec:956]{< < < PREV SERMON < < <}
~
\hyperref[sec:toc]{[返目錄]}
~
\hyperref[sec:959]{> > > NEXT SERMON > > >}
\newline
\newline
啟示錄 10:1-11:19
\newline
\begin{longtable}{cl}
\hline
\hline
章節 & 經文 (和合本修訂版)\\
\hline
10:1 & \begin{tabularx}{0.7\textwidth}{X} 我又看見另一位大力的天使從天降下，披著雲彩，頭上有彩虹，臉面像太陽，兩腳像火柱。 \end{tabularx} \\ \\ \relax
10:2 & \begin{tabularx}{0.7\textwidth}{X} 他手裡拿著展開的小書卷。他右腳踏海，左腳踏地， \end{tabularx} \\ \\ \relax
10:3 & \begin{tabularx}{0.7\textwidth}{X} 大聲呼喊，好像獅子吼叫。呼喊完了，就有七個雷發出聲音。 \end{tabularx} \\ \\ \relax
10:4 & \begin{tabularx}{0.7\textwidth}{X} 七個雷發聲後，我正要寫出來，就聽見從天上有聲音說：「七個雷所說的，你要封上，不可寫出來。」 \end{tabularx} \\ \\ \relax
10:5 & \begin{tabularx}{0.7\textwidth}{X} 我所看見的那踏海踏地的天使向天舉起右手， \end{tabularx} \\ \\ \relax
10:6 & \begin{tabularx}{0.7\textwidth}{X} 指著創造天和天上之物、地和地上之物、海和海中之物、直活到永永遠遠的那位起誓，說：「不再有時日了。」 \end{tabularx} \\ \\ \relax
10:7 & \begin{tabularx}{0.7\textwidth}{X} 但在第七位天使要吹號的日子，神的奧祕就要成全了，正如神向他僕人眾先知所宣告的。 \end{tabularx} \\ \\ \relax
10:8 & \begin{tabularx}{0.7\textwidth}{X} 我先前從天上所聽見的那聲音又吩咐我說：「你去，把那踏海踏地之天使手中展開的小書卷拿過來。」 \end{tabularx} \\ \\ \relax
10:9 & \begin{tabularx}{0.7\textwidth}{X} 我就走到天使那裡，對他說，請他把小書卷給我。他對我說：「你拿去，把它吃光。它會使你肚子發苦，然而在你口中會甘甜如蜜。」 \end{tabularx} \\ \\ \relax
10:10 & \begin{tabularx}{0.7\textwidth}{X} 於是我從天使手中把小書卷接過來，把它吃光了，在我口中果然甘甜如蜜，吃了以後，我肚子覺得發苦。 \end{tabularx} \\ \\ \relax
10:11 & \begin{tabularx}{0.7\textwidth}{X} 天使們對我說：「你必須指著許多民族、邦國、語言、君王再說預言。」 \end{tabularx} \\ \\ \relax
11:1 & \begin{tabularx}{0.7\textwidth}{X} 有一根蘆葦，像丈量的杖，賜給我；且有話說：「起來！將神的殿和祭壇，以及在殿中禮拜的人，都量一量。 \end{tabularx} \\ \\ \relax
11:2 & \begin{tabularx}{0.7\textwidth}{X} 只是殿外的院子不用量，因為這是要給外邦人的；他們將踐踏聖城四十二個月。 \end{tabularx} \\ \\ \relax
11:3 & \begin{tabularx}{0.7\textwidth}{X} 「我要賜權柄給我那兩個見證人，穿著粗麻衣說預言一千二百六十天。」 \end{tabularx} \\ \\ \relax
11:4 & \begin{tabularx}{0.7\textwidth}{X} 他們就是那站在世界之主面前的兩棵橄欖樹和兩個燈臺。 \end{tabularx} \\ \\ \relax
11:5 & \begin{tabularx}{0.7\textwidth}{X} 若有人想要害他們，就有火從他們口中噴出來，燒滅仇敵；凡想要害他們的都必須這樣被殺。 \end{tabularx} \\ \\ \relax
11:6 & \begin{tabularx}{0.7\textwidth}{X} 這二人有權柄關閉天空，使他們說預言的日子不下雨；又有權柄使水變為血，並且能隨時隨意用各樣的災害擊打大地。 \end{tabularx} \\ \\ \relax
11:7 & \begin{tabularx}{0.7\textwidth}{X} 他們作完見證的時候，那從無底坑裡上來的獸要跟他們交戰，並且得勝，把他們殺了。 \end{tabularx} \\ \\ \relax
11:8 & \begin{tabularx}{0.7\textwidth}{X} 他們的屍首將倒在大城的街道上；這城按著靈意叫所多瑪，又叫埃及，就是他們的主釘十字架的地方。 \end{tabularx} \\ \\ \relax
11:9 & \begin{tabularx}{0.7\textwidth}{X} 從各民族、支派、語言、邦國中有人觀看他們的屍首三天半，又不許人把屍首安放在墳墓裡。 \end{tabularx} \\ \\ \relax
11:10 & \begin{tabularx}{0.7\textwidth}{X} 住在地上的人會因他們而歡喜快樂，互相饋送禮物，因為這兩位先知曾使住在地上的人受痛苦。 \end{tabularx} \\ \\ \relax
11:11 & \begin{tabularx}{0.7\textwidth}{X} 過了這三天半，有生命的氣息從神那裡進入他們裡面，他們就站起來；看見他們的人都大大懼怕。 \end{tabularx} \\ \\ \relax
11:12 & \begin{tabularx}{0.7\textwidth}{X} 兩位先知聽見有大聲音從天上對他們說：「上這裡來。」他們就駕著雲上了天，他們的仇敵也看見了。 \end{tabularx} \\ \\ \relax
11:13 & \begin{tabularx}{0.7\textwidth}{X} 正在那時候，地大震動，城倒塌了十分之一；因地震而死的有七千人，其餘的都恐懼，歸榮耀給天上的神。 \end{tabularx} \\ \\ \relax
11:14 & \begin{tabularx}{0.7\textwidth}{X} 第二樣災禍過去了；看哪，第三樣災禍快到了。 \end{tabularx} \\ \\ \relax
11:15 & \begin{tabularx}{0.7\textwidth}{X} 第七位天使吹號，天上就有大聲音說：「世上的國已成了我們的主和他所立的基督的國了。他要作王直到永永遠遠！」 \end{tabularx} \\ \\ \relax
11:16 & \begin{tabularx}{0.7\textwidth}{X} 在神面前，坐在自己座位上的二十四位長老都俯伏在地上敬拜神， \end{tabularx} \\ \\ \relax
11:17 & \begin{tabularx}{0.7\textwidth}{X} 說：「今在昔在的主—全能的神啊，我們感謝你！因你執掌大權作王了。 \end{tabularx} \\ \\ \relax
11:18 & \begin{tabularx}{0.7\textwidth}{X} 外邦發怒，你的憤怒臨到了。審判死人的時候也到了；你的僕人眾先知、眾聖徒及敬畏你名的人，連大帶小得賞賜的時候到了；你毀滅那些毀滅大地者的時候也到了。」 \end{tabularx} \\ \\ \relax
11:19 & \begin{tabularx}{0.7\textwidth}{X} 於是，神天上的聖所開了，在他聖所中，他的約櫃出現了；隨後有閃電、響聲、雷轟、地震、大冰雹。 \end{tabularx} \\ \\
[1ex]
\hline
\hline
\end{longtable}


\newpage

\section{第66屆~港九培靈會~蘇穎智~第五講}
\label{sec:959}
\textbf{神選民的仇敵}
\newline
\newline
連結: \href{https://www.hkbibleconference.org/session-message/view/959}{\texttt{https://www.hkbibleconference.org/session-message/view/959}}
\newline
\newline
\hyperref[sec:957]{< < < PREV SERMON < < <}
~
\hyperref[sec:toc]{[返目錄]}
~
\hyperref[sec:960]{> > > NEXT SERMON > > >}
\newline
\newline
啟示錄 12:1-13:18
\newline
\begin{longtable}{cl}
\hline
\hline
章節 & 經文 (和合本修訂版)\\
\hline
12:1 & \begin{tabularx}{0.7\textwidth}{X} 天上出現了一個大兆頭：有一個婦人身披太陽，腳踏月亮，頭戴十二顆星的冠冕； \end{tabularx} \\ \\ \relax
12:2 & \begin{tabularx}{0.7\textwidth}{X} 她懷了孕，在生產的陣痛中疼痛地喊叫。 \end{tabularx} \\ \\ \relax
12:3 & \begin{tabularx}{0.7\textwidth}{X} 天上又出現了另一個兆頭：有一條大紅龍，有七個頭十個角；七個頭上戴著七個冠冕。 \end{tabularx} \\ \\ \relax
12:4 & \begin{tabularx}{0.7\textwidth}{X} 牠的尾巴拖拉著天上星辰的三分之一，把它們摔在地上。然後龍站在那將要生產的婦人面前，等她生產後要吞吃她的孩子。 \end{tabularx} \\ \\ \relax
12:5 & \begin{tabularx}{0.7\textwidth}{X} 婦人生了一個男孩子，就是將來要用鐵杖管轄萬國的；她的孩子被提到神和他寶座那裡去。 \end{tabularx} \\ \\ \relax
12:6 & \begin{tabularx}{0.7\textwidth}{X} 婦人就逃到曠野，在那裡有神給她預備的地方，使她在那裡被供養一千二百六十天。 \end{tabularx} \\ \\ \relax
12:7 & \begin{tabularx}{0.7\textwidth}{X} 天上發生了爭戰。米迦勒同他的使者與龍作戰，龍同牠的使者也起來應戰， \end{tabularx} \\ \\ \relax
12:8 & \begin{tabularx}{0.7\textwidth}{X} 牠們都打敗了，天上再也沒有牠們的地方。 \end{tabularx} \\ \\ \relax
12:9 & \begin{tabularx}{0.7\textwidth}{X} 大龍就是那古蛇，名叫魔鬼，又叫撒但，是迷惑普天下的；牠被摔在地上，牠的使者也一同被摔下去。 \end{tabularx} \\ \\ \relax
12:10 & \begin{tabularx}{0.7\textwidth}{X} 我聽見在天上有大聲音說：「我神的救恩、能力、國度，和他所立的基督的權柄現在都來到了。因為那個在我們神面前、晝夜控告我們弟兄的，已經被摔下去了。 \end{tabularx} \\ \\ \relax
12:11 & \begin{tabularx}{0.7\textwidth}{X} 弟兄勝過那條龍是因羔羊的血，和因自己所見證的道。雖然至於死，他們也不惜自己的性命。 \end{tabularx} \\ \\ \relax
12:12 & \begin{tabularx}{0.7\textwidth}{X} 所以，諸天和住在其中的，你們都快樂吧！只是地和海有禍了！因為魔鬼知道自己的時候不多，就氣憤憤地下到你們那裡去了。」 \end{tabularx} \\ \\ \relax
12:13 & \begin{tabularx}{0.7\textwidth}{X} 龍見自己被摔在地上，就迫害那生男孩子的婦人。 \end{tabularx} \\ \\ \relax
12:14 & \begin{tabularx}{0.7\textwidth}{X} 於是有大鷹的兩個翅膀賜給婦人，讓她能飛到曠野，到自己的地方，躲避那蛇。她在那裡受供養一載二載半載。 \end{tabularx} \\ \\ \relax
12:15 & \begin{tabularx}{0.7\textwidth}{X} 蛇在婦人背後，從口中噴出水來，像河一樣，要將婦人沖走。 \end{tabularx} \\ \\ \relax
12:16 & \begin{tabularx}{0.7\textwidth}{X} 地卻幫助了婦人，開口吞了從龍口噴出來的水。 \end{tabularx} \\ \\ \relax
12:17 & \begin{tabularx}{0.7\textwidth}{X} 於是龍向婦人發怒，去與她其餘的兒女作戰，就是與那些遵守神命令、為耶穌作見證的。那時龍站在海邊沙灘上。 \end{tabularx} \\ \\ \relax
13:1 & \begin{tabularx}{0.7\textwidth}{X} 我又看見一隻獸從海裡上來，有十個角七個頭；在十個角上戴著十個冠冕，七個頭上有褻瀆的名號。 \end{tabularx} \\ \\ \relax
13:2 & \begin{tabularx}{0.7\textwidth}{X} 我所看見的獸，形狀像豹，腳像熊的腳，口像獅子的口。那條龍將自己的能力、座位和大權柄都給了牠。 \end{tabularx} \\ \\ \relax
13:3 & \begin{tabularx}{0.7\textwidth}{X} 我看見獸的七個頭中，有一個似乎受了致命傷，那傷卻醫好了。全地的人都很驚訝，跟從了那隻獸。 \end{tabularx} \\ \\ \relax
13:4 & \begin{tabularx}{0.7\textwidth}{X} 他們都拜那條龍，因為牠把自己的權柄給了獸；又拜那隻獸，說：「誰能比這隻獸，誰能與牠交戰呢？」 \end{tabularx} \\ \\ \relax
13:5 & \begin{tabularx}{0.7\textwidth}{X} 龍又賜給那隻獸說誇大褻瀆話的口，又賜給牠權柄可以任意行事四十二個月。 \end{tabularx} \\ \\ \relax
13:6 & \begin{tabularx}{0.7\textwidth}{X} 那獸就開口向神說褻瀆的話，褻瀆神的名和他的帳幕，就是那些住在天上的。 \end{tabularx} \\ \\ \relax
13:7 & \begin{tabularx}{0.7\textwidth}{X} 牠又被准許與聖徒作戰，並且得勝，也賜給牠權柄，可以制伏各支派、各民族、各語言、各邦國。 \end{tabularx} \\ \\ \relax
13:8 & \begin{tabularx}{0.7\textwidth}{X} 凡住在地上、名字從創世以來沒有記在被殺羔羊的生命冊上的人都要拜牠。 \end{tabularx} \\ \\ \relax
13:9 & \begin{tabularx}{0.7\textwidth}{X} 凡有耳朵的都聽吧！ \end{tabularx} \\ \\ \relax
13:10 & \begin{tabularx}{0.7\textwidth}{X} 該被擄掠的，必被擄掠；該被刀殺的，必被刀殺。在此，聖徒要有耐心和信心。 \end{tabularx} \\ \\ \relax
13:11 & \begin{tabularx}{0.7\textwidth}{X} 我又看見另一隻獸從地裡上來。牠有兩個角如同羔羊，說話好像龍。 \end{tabularx} \\ \\ \relax
13:12 & \begin{tabularx}{0.7\textwidth}{X} 牠在第一隻獸面前施行第一隻獸所有的權柄，並且使地和住在地上的人拜那致命傷被醫好了的第一隻獸。 \end{tabularx} \\ \\ \relax
13:13 & \begin{tabularx}{0.7\textwidth}{X} 這隻獸又行大奇事，甚至在人面前使火從天降在地上。 \end{tabularx} \\ \\ \relax
13:14 & \begin{tabularx}{0.7\textwidth}{X} 牠得了權柄在第一隻獸面前能行奇事，迷惑住在地上的人，告訴他們要為那受過刀傷還活著的獸造個像。 \end{tabularx} \\ \\ \relax
13:15 & \begin{tabularx}{0.7\textwidth}{X} 又有權柄賜給牠，讓那隻獸的像有生氣，並且能說話，又使所有不拜獸像的人都被殺害。 \end{tabularx} \\ \\ \relax
13:16 & \begin{tabularx}{0.7\textwidth}{X} 牠又使眾人，無論大小、貧富，自主的、為奴的，都在右手上，或是在額上，打一個印記； \end{tabularx} \\ \\ \relax
13:17 & \begin{tabularx}{0.7\textwidth}{X} 這樣，除了那有印記，有獸的名或有獸名數字的，都不得買或賣。 \end{tabularx} \\ \\ \relax
13:18 & \begin{tabularx}{0.7\textwidth}{X} 在此，要有智慧：讓有悟性的人解開獸的數目吧，因為這是一個人的數字，那數字是六百六十六。 \end{tabularx} \\ \\
[1ex]
\hline
\hline
\end{longtable}


\newpage

\section{第66屆~港九培靈會~蘇穎智~第六講}
\label{sec:960}
\textbf{窮途末路的地球 — 七碗之災}
\newline
\newline
連結: \href{https://www.hkbibleconference.org/session-message/view/960}{\texttt{https://www.hkbibleconference.org/session-message/view/960}}
\newline
\newline
\hyperref[sec:959]{< < < PREV SERMON < < <}
~
\hyperref[sec:toc]{[返目錄]}
~
\hyperref[sec:961]{> > > NEXT SERMON > > >}
\newline
\newline
啟示錄 14:1-16:21
\newline
\begin{longtable}{cl}
\hline
\hline
章節 & 經文 (和合本修訂版)\\
\hline
14:1 & \begin{tabularx}{0.7\textwidth}{X} 我又觀看，看見羔羊站在錫安山，和他在一起的有十四萬四千人，都有他的名和他父親的名寫在額上。 \end{tabularx} \\ \\ \relax
14:2 & \begin{tabularx}{0.7\textwidth}{X} 我聽見從天上有聲音，像眾水的聲音和大雷的聲音，我所聽見的聲音好像琴師所彈的琴聲。 \end{tabularx} \\ \\ \relax
14:3 & \begin{tabularx}{0.7\textwidth}{X} 他們在寶座前，和在四活物及眾長老前唱新歌，除了從地上買來的那十四萬四千人以外，沒有人能學這歌。 \end{tabularx} \\ \\ \relax
14:4 & \begin{tabularx}{0.7\textwidth}{X} 這些人未曾沾染婦女，他們原是童身。羔羊無論往哪裡去，他們都跟隨他。他們是從人間買來的，作為初熟的果子歸給神和羔羊。 \end{tabularx} \\ \\ \relax
14:5 & \begin{tabularx}{0.7\textwidth}{X} 在他們口中找不出謊言，他們是沒有瑕疵的。 \end{tabularx} \\ \\ \relax
14:6 & \begin{tabularx}{0.7\textwidth}{X} 我又看見另一位天使在空中飛翔，有永遠的福音要傳給住在地上的人，就是各邦國、各支派、各語言、各民族。 \end{tabularx} \\ \\ \relax
14:7 & \begin{tabularx}{0.7\textwidth}{X} 他大聲說：「要敬畏神，把榮耀歸給他，因為他施行審判的時候已經到了。要敬拜那創造天、地、海和水源的主。」 \end{tabularx} \\ \\ \relax
14:8 & \begin{tabularx}{0.7\textwidth}{X} 另有第二位天使接著說：「傾覆了！那曾叫列國喝淫亂、烈怒之酒的大巴比倫傾覆了！」 \end{tabularx} \\ \\ \relax
14:9 & \begin{tabularx}{0.7\textwidth}{X} 另有第三位天使接著他們，大聲說：「若有人拜那隻獸和獸像，在額上或在手上受了印記， \end{tabularx} \\ \\ \relax
14:10 & \begin{tabularx}{0.7\textwidth}{X} 他也必喝神烈怒的酒；這酒是斟在神憤怒的杯中的純酒。他要在聖天使和羔羊面前，在火與硫磺之中受痛苦。 \end{tabularx} \\ \\ \relax
14:11 & \begin{tabularx}{0.7\textwidth}{X} 使他們受痛苦的煙往上冒，直到永永遠遠。那些拜獸和獸像，受了牠名字的印記的人，晝夜不得安寧。」 \end{tabularx} \\ \\ \relax
14:12 & \begin{tabularx}{0.7\textwidth}{X} 在此，遵守神命令和堅信耶穌真道的聖徒要有耐心。 \end{tabularx} \\ \\ \relax
14:13 & \begin{tabularx}{0.7\textwidth}{X} 我聽見從天上有聲音說：「你要寫下：從今以後，在主裡死去的人有福了。」聖靈說：「是的，他們要從自己的勞苦中得安息，因為工作的成果永隨著他們。」 \end{tabularx} \\ \\ \relax
14:14 & \begin{tabularx}{0.7\textwidth}{X} 我又觀看，看見有一片白雲，雲上坐著一位好像是人子的，頭上戴著金冠冕，手裡拿著鋒利的鐮刀。 \end{tabularx} \\ \\ \relax
14:15 & \begin{tabularx}{0.7\textwidth}{X} 另有一位天使從聖所出來，向那坐在雲上的大聲喊著：「伸出你的鐮刀來收割吧，因為收割的時候已經到了，地上的莊稼已經熟透了。」 \end{tabularx} \\ \\ \relax
14:16 & \begin{tabularx}{0.7\textwidth}{X} 於是那坐在雲上的把鐮刀向地上揮去，地上的莊稼就收割了。 \end{tabularx} \\ \\ \relax
14:17 & \begin{tabularx}{0.7\textwidth}{X} 另有一位天使從天上的聖所出來，他也拿著鋒利的鐮刀。 \end{tabularx} \\ \\ \relax
14:18 & \begin{tabularx}{0.7\textwidth}{X} 另有一位天使從祭壇出來，是有權柄管火的，向那拿著鋒利鐮刀的大聲喊著說：「伸出鋒利的鐮刀來，收取地上葡萄樹的果子，因為葡萄熟透了。」 \end{tabularx} \\ \\ \relax
14:19 & \begin{tabularx}{0.7\textwidth}{X} 那天使就把鐮刀向地上揮去，收取了地上的葡萄，扔進神憤怒的大醡酒池裡。 \end{tabularx} \\ \\ \relax
14:20 & \begin{tabularx}{0.7\textwidth}{X} 那醡酒池在城外被踹踏，有血從醡酒池裡流出來，漲到馬的嚼環那麼高，約有一千六百斯他迪那麼遠。 \end{tabularx} \\ \\ \relax
15:1 & \begin{tabularx}{0.7\textwidth}{X} 我看見在天上有另一兆頭，大而且奇，就是七位天使掌管末了的七種災難，因為神的烈怒在這七種災難中發盡了。 \end{tabularx} \\ \\ \relax
15:2 & \begin{tabularx}{0.7\textwidth}{X} 我看見彷彿有攙雜火的玻璃海；又看見那些勝了那獸和獸像，以及牠名字的數字的人，都站在玻璃海上，拿著神的豎琴。 \end{tabularx} \\ \\ \relax
15:3 & \begin{tabularx}{0.7\textwidth}{X} 他們唱神僕人摩西的歌和羔羊的歌，說：「主—全能的神啊，你的作為又偉大又奇妙！萬國之王啊，你的道路又公義又真實！ \end{tabularx} \\ \\ \relax
15:4 & \begin{tabularx}{0.7\textwidth}{X} 主啊，誰敢不敬畏你，不把榮耀歸於你的名？因為只有你是神聖的。萬民都要來，在你面前敬拜，因你公義的作為已經彰顯了。」 \end{tabularx} \\ \\ \relax
15:5 & \begin{tabularx}{0.7\textwidth}{X} 此後，我看見在天上那存放法櫃的聖所開了。 \end{tabularx} \\ \\ \relax
15:6 & \begin{tabularx}{0.7\textwidth}{X} 那掌管七種災難的七位天使從聖所出來，穿著潔白明亮的細麻衣，胸間束著金帶。 \end{tabularx} \\ \\ \relax
15:7 & \begin{tabularx}{0.7\textwidth}{X} 四個活物中，有一個把盛滿了活到永永遠遠之神烈怒的七個金碗給了那七位天使。 \end{tabularx} \\ \\ \relax
15:8 & \begin{tabularx}{0.7\textwidth}{X} 聖所中充滿了神的榮耀和權能而來的煙。沒有人能進入聖所，直等到那七位天使降完了七種災難。 \end{tabularx} \\ \\ \relax
16:1 & \begin{tabularx}{0.7\textwidth}{X} 我聽見有大聲音從聖所裡出來，向那七位天使說：「你們去，把盛著神烈怒的七碗傾倒在地上。」 \end{tabularx} \\ \\ \relax
16:2 & \begin{tabularx}{0.7\textwidth}{X} 第一位天使去，把碗傾倒在地上，就有又臭又毒的瘡生在那些有獸的印記和拜獸像的人身上。 \end{tabularx} \\ \\ \relax
16:3 & \begin{tabularx}{0.7\textwidth}{X} 第二位天使把碗傾倒在海裡，海就變成像死人的血一樣，海裡所有的活物都死了。 \end{tabularx} \\ \\ \relax
16:4 & \begin{tabularx}{0.7\textwidth}{X} 第三位天使把碗傾倒在河流和水源裡，水就變成血了。 \end{tabularx} \\ \\ \relax
16:5 & \begin{tabularx}{0.7\textwidth}{X} 我聽見掌管眾水的天使說：「昔在、今在的聖者啊，你做的判斷公義； \end{tabularx} \\ \\ \relax
16:6 & \begin{tabularx}{0.7\textwidth}{X} 因他們曾流過聖徒與先知的血，現在你給他們血喝，這是他們該受的。」 \end{tabularx} \\ \\ \relax
16:7 & \begin{tabularx}{0.7\textwidth}{X} 我又聽見祭壇中有聲音說：「是的，主—全能的神啊，你的判斷又真實又公義！」 \end{tabularx} \\ \\ \relax
16:8 & \begin{tabularx}{0.7\textwidth}{X} 第四位天使把碗傾倒在太陽上，使太陽可用火烤人。 \end{tabularx} \\ \\ \relax
16:9 & \begin{tabularx}{0.7\textwidth}{X} 人被炎熱所烤，就褻瀆那有權掌管這些災難的神的名，他們沒有悔改，也沒有把榮耀歸給神。 \end{tabularx} \\ \\ \relax
16:10 & \begin{tabularx}{0.7\textwidth}{X} 第五位天使把碗傾倒在獸的座位上，獸的國就變成黑暗。人因疼痛而咬自己的舌頭； \end{tabularx} \\ \\ \relax
16:11 & \begin{tabularx}{0.7\textwidth}{X} 又因所受的疼痛和生的瘡，就褻瀆天上的神，也沒有為他們的行為悔改。 \end{tabularx} \\ \\ \relax
16:12 & \begin{tabularx}{0.7\textwidth}{X} 第六位天使把碗傾倒在大幼發拉底河上，河水就乾了，為要給從日出之地所來的眾王預備道路。 \end{tabularx} \\ \\ \relax
16:13 & \begin{tabularx}{0.7\textwidth}{X} 我又看見三個污穢的靈，好像青蛙，從龍的口、獸的口和假先知的口中出來。 \end{tabularx} \\ \\ \relax
16:14 & \begin{tabularx}{0.7\textwidth}{X} 他們本是鬼魔的靈，施行奇事，到普天下眾王那裡去，召集他們在全能者神的大日子作戰。 \end{tabularx} \\ \\ \relax
16:15 & \begin{tabularx}{0.7\textwidth}{X} 看哪，我來像賊一樣。那警醒、穿著衣服的人有福了；他不至於赤身而行，給人看見他的羞恥。 \end{tabularx} \\ \\ \relax
16:16 & \begin{tabularx}{0.7\textwidth}{X} 於是，那三個鬼魔把眾王聚集在希伯來話叫作哈米吉多頓的地方。 \end{tabularx} \\ \\ \relax
16:17 & \begin{tabularx}{0.7\textwidth}{X} 第七位天使把碗傾倒在空中，就有大聲音從聖所的寶座上出來，說：「成了！」 \end{tabularx} \\ \\ \relax
16:18 & \begin{tabularx}{0.7\textwidth}{X} 又有閃電、響聲、雷轟、大地震，自從地上有人以來沒有這樣大、這樣厲害的地震。 \end{tabularx} \\ \\ \relax
16:19 & \begin{tabularx}{0.7\textwidth}{X} 那大城裂為三段，列國的城也都倒塌了。神記起了大巴比倫城，把那盛自己烈怒的酒杯遞給她。 \end{tabularx} \\ \\ \relax
16:20 & \begin{tabularx}{0.7\textwidth}{X} 各海島都逃避了，眾山也不見了。 \end{tabularx} \\ \\ \relax
16:21 & \begin{tabularx}{0.7\textwidth}{X} 又有大冰雹從天掉落在人身上，每一個約重一他連得，以致人因冰雹的災難而褻瀆神，因為那災難太大了。 \end{tabularx} \\ \\
[1ex]
\hline
\hline
\end{longtable}


\newpage

\section{第66屆~港九培靈會~蘇穎智~第七講}
\label{sec:961}
\textbf{必有後「福! 覆!」的大巴比倫}
\newline
\newline
連結: \href{https://www.hkbibleconference.org/session-message/view/961}{\texttt{https://www.hkbibleconference.org/session-message/view/961}}
\newline
\newline
\hyperref[sec:960]{< < < PREV SERMON < < <}
~
\hyperref[sec:toc]{[返目錄]}
~
\hyperref[sec:962]{> > > NEXT SERMON > > >}
\newline
\newline
啟示錄 17:1-19:21
\newline
\begin{longtable}{cl}
\hline
\hline
章節 & 經文 (和合本修訂版)\\
\hline
17:1 & \begin{tabularx}{0.7\textwidth}{X} 拿著七個碗的七位天使中，有一位前來對我說：「來，我要讓你看那坐在眾水之上的大淫婦所要受的懲罰； \end{tabularx} \\ \\ \relax
17:2 & \begin{tabularx}{0.7\textwidth}{X} 地上的君王都曾與她行淫，住在地上的人也喝醉了她淫亂的酒。」 \end{tabularx} \\ \\ \relax
17:3 & \begin{tabularx}{0.7\textwidth}{X} 我在聖靈感動下，被天使帶到曠野去，我看見一個女人騎在朱紅色的獸上；那隻獸有七個頭十個角，遍體有褻瀆的名號。 \end{tabularx} \\ \\ \relax
17:4 & \begin{tabularx}{0.7\textwidth}{X} 那女人穿著紫色和朱紅色的衣服，用金子、寶石、珍珠作妝飾，手拿著金杯，杯中盛滿了可憎之物和她淫亂的污穢。 \end{tabularx} \\ \\ \relax
17:5 & \begin{tabularx}{0.7\textwidth}{X} 在她額上寫著奧祕的名字，說：「大巴比倫，世上的淫婦和一切可憎之物的母。」 \end{tabularx} \\ \\ \relax
17:6 & \begin{tabularx}{0.7\textwidth}{X} 我又看見那女人喝醉了聖徒的血和為耶穌作見證的人的血。我看見她，非常詫異。 \end{tabularx} \\ \\ \relax
17:7 & \begin{tabularx}{0.7\textwidth}{X} 天使對我說：「你為甚麼詫異呢？我要把這女人和馱著她那七頭十角的獸的奧祕告訴你。 \end{tabularx} \\ \\ \relax
17:8 & \begin{tabularx}{0.7\textwidth}{X} 你曾看見的獸，以前有，現在沒有，將來要從無底坑裡上來，又歸於沉淪。凡住在地上、名字從創世以來沒有記在生命冊上的人看見那隻獸都要詫異，因為牠以前有，現在沒有，以後再有。 \end{tabularx} \\ \\ \relax
17:9 & \begin{tabularx}{0.7\textwidth}{X} 在此要有智慧的心思：那七個頭就是女人所坐的七座山；他們又是七個王， \end{tabularx} \\ \\ \relax
17:10 & \begin{tabularx}{0.7\textwidth}{X} 五個已經倒了，一個還在，一個還沒有來到；他來的時候必須只暫時停留。 \end{tabularx} \\ \\ \relax
17:11 & \begin{tabularx}{0.7\textwidth}{X} 那以前有、現在沒有的獸就是第八個，他也和那七個同列，正歸於沉淪。 \end{tabularx} \\ \\ \relax
17:12 & \begin{tabularx}{0.7\textwidth}{X} 你曾看見的那十個角就是十個王；他們還沒有得到國度，但他們要和那隻獸同得權柄作王一個時辰。 \end{tabularx} \\ \\ \relax
17:13 & \begin{tabularx}{0.7\textwidth}{X} 他們同心把自己的能力權柄交給那隻獸。 \end{tabularx} \\ \\ \relax
17:14 & \begin{tabularx}{0.7\textwidth}{X} 他們將與羔羊作戰，羔羊必勝過他們，因為羔羊是萬主之主、萬王之王，而同羔羊在一起的是蒙召、被選、忠心的人。」 \end{tabularx} \\ \\ \relax
17:15 & \begin{tabularx}{0.7\textwidth}{X} 天使又對我說：「你所看見那淫婦坐的眾水，就是許多民族、人民、邦國、語言。 \end{tabularx} \\ \\ \relax
17:16 & \begin{tabularx}{0.7\textwidth}{X} 你所看見的那十個角與獸必恨這淫婦，他們要使她孤獨赤身，又要吃她的肉，用火將她燒盡。 \end{tabularx} \\ \\ \relax
17:17 & \begin{tabularx}{0.7\textwidth}{X} 因為神使諸王同心執行他的旨意，把他們自己的國交給那隻獸，直等到神的話都應驗了。 \end{tabularx} \\ \\ \relax
17:18 & \begin{tabularx}{0.7\textwidth}{X} 你所看見的那女人就是管轄地上眾王的大城。」 \end{tabularx} \\ \\ \relax
18:1 & \begin{tabularx}{0.7\textwidth}{X} 此後，我看見另一位有大權柄的天使從天降下，地由於他的榮耀而發光。 \end{tabularx} \\ \\ \relax
18:2 & \begin{tabularx}{0.7\textwidth}{X} 他以強而有力的聲音喊著說：「傾覆了！大巴比倫傾覆了！她成了鬼魔的住處，各樣污穢之靈的巢穴，各樣污穢之鳥的窩，各樣污穢可憎之獸的出沒處。 \end{tabularx} \\ \\ \relax
18:3 & \begin{tabularx}{0.7\textwidth}{X} 因為列國都喝了她淫亂大怒的酒；地上的君王和她行淫；地上的商人因她極度奢華而發了財。」 \end{tabularx} \\ \\ \relax
18:4 & \begin{tabularx}{0.7\textwidth}{X} 我又聽見另一個聲音從天上說：「我的民哪，從那城出來吧！免得和她在罪上有份，受她所受的災殃； \end{tabularx} \\ \\ \relax
18:5 & \begin{tabularx}{0.7\textwidth}{X} 因她的罪惡滔天，神已經記得她的不義。 \end{tabularx} \\ \\ \relax
18:6 & \begin{tabularx}{0.7\textwidth}{X} 她怎樣待人，也要怎樣待她，按她所行的加倍地報應她；用她調酒的杯加倍調給她喝。 \end{tabularx} \\ \\ \relax
18:7 & \begin{tabularx}{0.7\textwidth}{X} 她怎樣榮耀自己，怎樣奢華，也要使她照樣痛苦悲哀。因她心裡說：『我坐了皇后的位，並不是寡婦，絕不至於悲哀。』 \end{tabularx} \\ \\ \relax
18:8 & \begin{tabularx}{0.7\textwidth}{X} 所以在一天之內，她的災殃要一齊來到，就是死亡、悲哀、饑荒。她將被火燒盡，因為審判她的主神大有能力。」 \end{tabularx} \\ \\ \relax
18:9 & \begin{tabularx}{0.7\textwidth}{X} 地上的君王，與她行淫、一同奢華的，看見燒她的煙，就必為她哭泣哀號； \end{tabularx} \\ \\ \relax
18:10 & \begin{tabularx}{0.7\textwidth}{X} 因怕她的痛苦，就遠遠地站著，說：「禍哉，禍哉，這大城！堅固的巴比倫城啊！一時之間，你的審判要來到了。」 \end{tabularx} \\ \\ \relax
18:11 & \begin{tabularx}{0.7\textwidth}{X} 地上的商人也都為她哭泣悲哀，因為沒有人再買他們的貨物了； \end{tabularx} \\ \\ \relax
18:12 & \begin{tabularx}{0.7\textwidth}{X} 這貨物就是金、銀、寶石、珍珠、細麻布、絲綢、紫色和朱紅色衣料、各樣香木、各樣象牙的器皿、各樣極寶貴的木頭和銅、鐵、大理石的器皿， \end{tabularx} \\ \\ \relax
18:13 & \begin{tabularx}{0.7\textwidth}{X} 和肉桂、豆蔻、香料、香膏、乳香、酒、油、細麵、麥子、牛、羊、馬、馬車，以及奴隸、人口。 \end{tabularx} \\ \\ \relax
18:14 & \begin{tabularx}{0.7\textwidth}{X} 「你所貪愛的果子離開了你；你一切的珍饈美味和華美的物件都從你那裡毀滅，絕對見不到了。」 \end{tabularx} \\ \\ \relax
18:15 & \begin{tabularx}{0.7\textwidth}{X} 販賣這些貨物、藉著她發財的商人，因怕她的痛苦，就遠遠地站著哭泣悲哀， \end{tabularx} \\ \\ \relax
18:16 & \begin{tabularx}{0.7\textwidth}{X} 說：「禍哉，禍哉，這大城！她穿著細麻、紫色、朱紅色的衣服，用金子、寶石、珍珠為妝飾。 \end{tabularx} \\ \\ \relax
18:17 & \begin{tabularx}{0.7\textwidth}{X} 一時之間，這麼多的財富就歸於無有了。」所有的船長和到處航海的，水手以及所有靠海為業的，都遠遠地站著， \end{tabularx} \\ \\ \relax
18:18 & \begin{tabularx}{0.7\textwidth}{X} 看見燒她的煙，就喊著說：「有哪一個城能跟這大城比呢？」 \end{tabularx} \\ \\ \relax
18:19 & \begin{tabularx}{0.7\textwidth}{X} 於是他們把灰塵撒在頭上，哭泣悲哀地喊著說：「禍哉，禍哉，這大城！凡有船在海中的，都因她的珍寶成了富足。她在一時之間就成為荒蕪。 \end{tabularx} \\ \\ \relax
18:20 & \begin{tabularx}{0.7\textwidth}{X} 天哪，眾聖徒、眾使徒、眾先知啊！你們都要因她歡喜，因為神已經在她身上為你們伸了冤。」 \end{tabularx} \\ \\ \relax
18:21 & \begin{tabularx}{0.7\textwidth}{X} 有一位大力的天使舉起一塊石頭，好像大磨石，扔在海裡，說：「巴比倫大城也必這樣猛力地被扔下去，絕對見不到了。 \end{tabularx} \\ \\ \relax
18:22 & \begin{tabularx}{0.7\textwidth}{X} 彈琴、歌唱、吹笛、吹號的聲音，在你中間絕對聽不見了；各行手藝的技工在你中間絕對見不到了；推磨的聲音在你中間絕對聽不見了； \end{tabularx} \\ \\ \relax
18:23 & \begin{tabularx}{0.7\textwidth}{X} 燈臺的光在你中間絕對不再照耀了；新郎和新娘的聲音在你中間絕對聽不見了。你的商人原來是地上的顯要；萬國也被你的邪術迷惑了。 \end{tabularx} \\ \\ \relax
18:24 & \begin{tabularx}{0.7\textwidth}{X} 先知、聖徒和地上一切被殺的人的血都在這城裡找到了。」 \end{tabularx} \\ \\ \relax
19:1 & \begin{tabularx}{0.7\textwidth}{X} 此後，我聽見好像有一大群人在天上大聲說：「哈利路亞！救恩、榮耀、權能都屬於我們的神。 \end{tabularx} \\ \\ \relax
19:2 & \begin{tabularx}{0.7\textwidth}{X} 他的判斷又真實又公義；因他判斷了那大淫婦，她用淫行敗壞了世界。神為他的僕人伸冤，向淫婦討流僕人血的罪。」 \end{tabularx} \\ \\ \relax
19:3 & \begin{tabularx}{0.7\textwidth}{X} 他們又一次說：「哈利路亞！燒淫婦的煙往上冒，直到永永遠遠。」 \end{tabularx} \\ \\ \relax
19:4 & \begin{tabularx}{0.7\textwidth}{X} 那二十四位長老和四活物就俯伏敬拜坐在寶座上的神，說：「阿們。哈利路亞！」 \end{tabularx} \\ \\ \relax
19:5 & \begin{tabularx}{0.7\textwidth}{X} 接著，有聲音從寶座出來說：「神的眾僕人哪，凡敬畏他的，無論大小，都要讚美我們的神！」 \end{tabularx} \\ \\ \relax
19:6 & \begin{tabularx}{0.7\textwidth}{X} 我聽見好像一大群人的聲音，像眾水的聲音，像大雷的聲音，說：「哈利路亞！因為主---我們的神、全能者，作王了。 \end{tabularx} \\ \\ \relax
19:7 & \begin{tabularx}{0.7\textwidth}{X} 我們要歡喜快樂，將榮耀歸給他；因為羔羊的婚期到了，他的新娘也自己預備好了， \end{tabularx} \\ \\ \relax
19:8 & \begin{tabularx}{0.7\textwidth}{X} 她蒙恩得穿明亮潔白的細麻衣：這細麻衣就是聖徒們的義行。」 \end{tabularx} \\ \\ \relax
19:9 & \begin{tabularx}{0.7\textwidth}{X} 天使對我說：「你要寫下來：凡被請赴羔羊婚宴的人有福了！」他又對我說：「這些都是神真實的話。」 \end{tabularx} \\ \\ \relax
19:10 & \begin{tabularx}{0.7\textwidth}{X} 我就俯伏在他腳前要拜他。他對我說：「千萬不可！我和你，以及那些為耶穌作見證的弟兄同是僕人。你要敬拜神。」因為那些為耶穌作見證的人有預言的靈。 \end{tabularx} \\ \\ \relax
19:11 & \begin{tabularx}{0.7\textwidth}{X} 後來我看見天開了。有一匹白馬，騎在馬上的稱為「誠信」、「真實」，他審判和爭戰都憑著公義。 \end{tabularx} \\ \\ \relax
19:12 & \begin{tabularx}{0.7\textwidth}{X} 他的眼睛如火焰，頭上戴著許多冠冕；他身上寫著一個名字，除了他自己沒有人知道。 \end{tabularx} \\ \\ \relax
19:13 & \begin{tabularx}{0.7\textwidth}{X} 他穿著浸過血的衣服；他的名稱為「神之道」。 \end{tabularx} \\ \\ \relax
19:14 & \begin{tabularx}{0.7\textwidth}{X} 眾天軍都騎著白馬，穿著又白又潔淨的細麻衣跟隨他。 \end{tabularx} \\ \\ \relax
19:15 & \begin{tabularx}{0.7\textwidth}{X} 有利劍從他口中出來，用來擊打列國。他要用鐵杖管轄他們，並且要踹全能神烈怒的醡酒池。 \end{tabularx} \\ \\ \relax
19:16 & \begin{tabularx}{0.7\textwidth}{X} 在他衣服和大腿上寫著「萬王之王，萬主之主」的名號。 \end{tabularx} \\ \\ \relax
19:17 & \begin{tabularx}{0.7\textwidth}{X} 我又看見一位天使站在太陽中，向天空一切的飛鳥大聲喊著說：「你們聚集來赴神的大宴席， \end{tabularx} \\ \\ \relax
19:18 & \begin{tabularx}{0.7\textwidth}{X} 為要吃君王的肉、將軍的肉、壯士的肉、馬和騎士的肉、一切自主的和為奴的，以及尊貴的和卑賤的肉。」 \end{tabularx} \\ \\ \relax
19:19 & \begin{tabularx}{0.7\textwidth}{X} 我又看見那獸和地上的君王，和他們的軍隊都聚集，要與白馬騎士和他的軍隊作戰。 \end{tabularx} \\ \\ \relax
19:20 & \begin{tabularx}{0.7\textwidth}{X} 那獸被擒拿了；那在獸面前曾行奇事、迷惑了接受獸的印記和拜獸像的人的假先知，也與獸同被擒拿。他們兩個就活生生地被扔進燒著硫磺的火湖裡， \end{tabularx} \\ \\ \relax
19:21 & \begin{tabularx}{0.7\textwidth}{X} 其餘的人被白馬騎士口中吐出來的劍殺了；所有的飛鳥都吃飽了他們的肉。 \end{tabularx} \\ \\
[1ex]
\hline
\hline
\end{longtable}


\newpage

\section{第66屆~港九培靈會~蘇穎智~第八講}
\label{sec:962}
\textbf{大結局 — 有人快樂有人愁}
\newline
\newline
連結: \href{https://www.hkbibleconference.org/session-message/view/962}{\texttt{https://www.hkbibleconference.org/session-message/view/962}}
\newline
\newline
\hyperref[sec:961]{< < < PREV SERMON < < <}
~
\hyperref[sec:toc]{[返目錄]}
~
\hyperref[sec:963]{> > > NEXT SERMON > > >}
\newline
\newline
啟示錄 19:1-20:15
\newline
\begin{longtable}{cl}
\hline
\hline
章節 & 經文 (和合本修訂版)\\
\hline
19:1 & \begin{tabularx}{0.7\textwidth}{X} 此後，我聽見好像有一大群人在天上大聲說：「哈利路亞！救恩、榮耀、權能都屬於我們的神。 \end{tabularx} \\ \\ \relax
19:2 & \begin{tabularx}{0.7\textwidth}{X} 他的判斷又真實又公義；因他判斷了那大淫婦，她用淫行敗壞了世界。神為他的僕人伸冤，向淫婦討流僕人血的罪。」 \end{tabularx} \\ \\ \relax
19:3 & \begin{tabularx}{0.7\textwidth}{X} 他們又一次說：「哈利路亞！燒淫婦的煙往上冒，直到永永遠遠。」 \end{tabularx} \\ \\ \relax
19:4 & \begin{tabularx}{0.7\textwidth}{X} 那二十四位長老和四活物就俯伏敬拜坐在寶座上的神，說：「阿們。哈利路亞！」 \end{tabularx} \\ \\ \relax
19:5 & \begin{tabularx}{0.7\textwidth}{X} 接著，有聲音從寶座出來說：「神的眾僕人哪，凡敬畏他的，無論大小，都要讚美我們的神！」 \end{tabularx} \\ \\ \relax
19:6 & \begin{tabularx}{0.7\textwidth}{X} 我聽見好像一大群人的聲音，像眾水的聲音，像大雷的聲音，說：「哈利路亞！因為主---我們的神、全能者，作王了。 \end{tabularx} \\ \\ \relax
19:7 & \begin{tabularx}{0.7\textwidth}{X} 我們要歡喜快樂，將榮耀歸給他；因為羔羊的婚期到了，他的新娘也自己預備好了， \end{tabularx} \\ \\ \relax
19:8 & \begin{tabularx}{0.7\textwidth}{X} 她蒙恩得穿明亮潔白的細麻衣：這細麻衣就是聖徒們的義行。」 \end{tabularx} \\ \\ \relax
19:9 & \begin{tabularx}{0.7\textwidth}{X} 天使對我說：「你要寫下來：凡被請赴羔羊婚宴的人有福了！」他又對我說：「這些都是神真實的話。」 \end{tabularx} \\ \\ \relax
19:10 & \begin{tabularx}{0.7\textwidth}{X} 我就俯伏在他腳前要拜他。他對我說：「千萬不可！我和你，以及那些為耶穌作見證的弟兄同是僕人。你要敬拜神。」因為那些為耶穌作見證的人有預言的靈。 \end{tabularx} \\ \\ \relax
19:11 & \begin{tabularx}{0.7\textwidth}{X} 後來我看見天開了。有一匹白馬，騎在馬上的稱為「誠信」、「真實」，他審判和爭戰都憑著公義。 \end{tabularx} \\ \\ \relax
19:12 & \begin{tabularx}{0.7\textwidth}{X} 他的眼睛如火焰，頭上戴著許多冠冕；他身上寫著一個名字，除了他自己沒有人知道。 \end{tabularx} \\ \\ \relax
19:13 & \begin{tabularx}{0.7\textwidth}{X} 他穿著浸過血的衣服；他的名稱為「神之道」。 \end{tabularx} \\ \\ \relax
19:14 & \begin{tabularx}{0.7\textwidth}{X} 眾天軍都騎著白馬，穿著又白又潔淨的細麻衣跟隨他。 \end{tabularx} \\ \\ \relax
19:15 & \begin{tabularx}{0.7\textwidth}{X} 有利劍從他口中出來，用來擊打列國。他要用鐵杖管轄他們，並且要踹全能神烈怒的醡酒池。 \end{tabularx} \\ \\ \relax
19:16 & \begin{tabularx}{0.7\textwidth}{X} 在他衣服和大腿上寫著「萬王之王，萬主之主」的名號。 \end{tabularx} \\ \\ \relax
19:17 & \begin{tabularx}{0.7\textwidth}{X} 我又看見一位天使站在太陽中，向天空一切的飛鳥大聲喊著說：「你們聚集來赴神的大宴席， \end{tabularx} \\ \\ \relax
19:18 & \begin{tabularx}{0.7\textwidth}{X} 為要吃君王的肉、將軍的肉、壯士的肉、馬和騎士的肉、一切自主的和為奴的，以及尊貴的和卑賤的肉。」 \end{tabularx} \\ \\ \relax
19:19 & \begin{tabularx}{0.7\textwidth}{X} 我又看見那獸和地上的君王，和他們的軍隊都聚集，要與白馬騎士和他的軍隊作戰。 \end{tabularx} \\ \\ \relax
19:20 & \begin{tabularx}{0.7\textwidth}{X} 那獸被擒拿了；那在獸面前曾行奇事、迷惑了接受獸的印記和拜獸像的人的假先知，也與獸同被擒拿。他們兩個就活生生地被扔進燒著硫磺的火湖裡， \end{tabularx} \\ \\ \relax
19:21 & \begin{tabularx}{0.7\textwidth}{X} 其餘的人被白馬騎士口中吐出來的劍殺了；所有的飛鳥都吃飽了他們的肉。 \end{tabularx} \\ \\ \relax
20:1 & \begin{tabularx}{0.7\textwidth}{X} 我又看見一位天使從天降下，手裡拿著無底坑的鑰匙和一條大鐵鏈。 \end{tabularx} \\ \\ \relax
20:2 & \begin{tabularx}{0.7\textwidth}{X} 他抓住那龍，那古蛇，就是魔鬼、撒但，把牠捆綁了一千年， \end{tabularx} \\ \\ \relax
20:3 & \begin{tabularx}{0.7\textwidth}{X} 扔在無底坑裡，把無底坑關閉，用印封上，使牠不再迷惑列國，等到那一千年滿了。這些事以後，牠必須暫時被釋放。 \end{tabularx} \\ \\ \relax
20:4 & \begin{tabularx}{0.7\textwidth}{X} 我又看見一些寶座，坐在上面的有審判的權柄賜給他們。我又看見那些因為給耶穌作見證，並為神之道被斬首的人的靈魂，和沒有拜過那獸與獸像、也沒有在額上和手上打過牠印記的人的靈魂。他們都復活了，與基督一同作王一千年。 \end{tabularx} \\ \\ \relax
20:5 & \begin{tabularx}{0.7\textwidth}{X} 這是頭一次的復活。其餘的死人還沒有復活，直等那一千年滿了。 \end{tabularx} \\ \\ \relax
20:6 & \begin{tabularx}{0.7\textwidth}{X} 在頭一次復活有份的有福了，聖潔了！第二次的死在他們身上沒有權柄，但他們要作神和基督的祭司，也要與基督一同作王一千年。 \end{tabularx} \\ \\ \relax
20:7 & \begin{tabularx}{0.7\textwidth}{X} 那一千年滿了，撒但會從監牢裡被釋放， \end{tabularx} \\ \\ \relax
20:8 & \begin{tabularx}{0.7\textwidth}{X} 出來要迷惑地上四方的列國，就是歌革和瑪各，使他們聚集爭戰。他們的人數多如海沙。 \end{tabularx} \\ \\ \relax
20:9 & \begin{tabularx}{0.7\textwidth}{X} 他們上來佈滿了全地，圍住聖徒的營與蒙愛的城，就有火從天降下，燒滅了他們。 \end{tabularx} \\ \\ \relax
20:10 & \begin{tabularx}{0.7\textwidth}{X} 那迷惑他們的魔鬼被扔進硫磺的火湖裡，就是那獸和假先知所在的地方，他們會晝夜受折磨，直到永永遠遠。 \end{tabularx} \\ \\ \relax
20:11 & \begin{tabularx}{0.7\textwidth}{X} 我又看見一個白色的大寶座和那坐在上面的；天和地都從他面前逃避，再也找不到它們的位置了。 \end{tabularx} \\ \\ \relax
20:12 & \begin{tabularx}{0.7\textwidth}{X} 我又看見死了的人，無論大小，都站在寶座前。案卷都展開了，並另有一卷展開，就是生命冊。死了的人都憑著這些案卷所記載的，照他們所行的受審判。 \end{tabularx} \\ \\ \relax
20:13 & \begin{tabularx}{0.7\textwidth}{X} 於是海交出其中的死人，死亡和陰間也交出其中的死人；他們都照各人所行的受審判。 \end{tabularx} \\ \\ \relax
20:14 & \begin{tabularx}{0.7\textwidth}{X} 死亡和陰間也被扔進火湖裡，這火湖就是第二次的死。 \end{tabularx} \\ \\ \relax
20:15 & \begin{tabularx}{0.7\textwidth}{X} 凡名字沒有記在生命冊上的人，就被扔進火湖裡。 \end{tabularx} \\ \\
[1ex]
\hline
\hline
\end{longtable}


\newpage

\section{第66屆~港九培靈會~蘇穎智~第九講}
\label{sec:963}
\textbf{更美的家}
\newline
\newline
連結: \href{https://www.hkbibleconference.org/session-message/view/963}{\texttt{https://www.hkbibleconference.org/session-message/view/963}}
\newline
\newline
\hyperref[sec:962]{< < < PREV SERMON < < <}
~
\hyperref[sec:toc]{[返目錄]}
~
\hyperref[sec:964]{> > > NEXT SERMON > > >}
\newline
\newline
啟示錄 21:1-22:21
\newline
\begin{longtable}{cl}
\hline
\hline
章節 & 經文 (和合本修訂版)\\
\hline
21:1 & \begin{tabularx}{0.7\textwidth}{X} 我又看見一個新天新地，因為先前的天和先前的地已經過去了，海也不再有了。 \end{tabularx} \\ \\ \relax
21:2 & \begin{tabularx}{0.7\textwidth}{X} 我又看見聖城，新耶路撒冷由神那裡，從天而降，預備好了，就如新娘打扮整齊，等候丈夫。 \end{tabularx} \\ \\ \relax
21:3 & \begin{tabularx}{0.7\textwidth}{X} 我聽見有大聲音從寶座出來，說：「看哪，神的帳幕在人間！他要和他們同住，他們要作他的子民。神要親自與他們同在。 \end{tabularx} \\ \\ \relax
21:4 & \begin{tabularx}{0.7\textwidth}{X} 神要擦去他們一切的眼淚；不再有死亡，也不再有悲哀、哭號、痛苦，因為先前的事都過去了。」 \end{tabularx} \\ \\ \relax
21:5 & \begin{tabularx}{0.7\textwidth}{X} 那位坐在寶座上的說：「看哪，我把一切都更新了！」他又說：「你要寫下來，因為這些話是可信靠的，是真實的。」 \end{tabularx} \\ \\ \relax
21:6 & \begin{tabularx}{0.7\textwidth}{X} 他又對我說：「成了！我是阿拉法，我是俄梅戛；我是開始，我是終結。我要把生命的泉水白白賜給那口渴的人喝。 \end{tabularx} \\ \\ \relax
21:7 & \begin{tabularx}{0.7\textwidth}{X} 得勝的要承受這些為業；我要作他的神，他要作我的兒子。 \end{tabularx} \\ \\ \relax
21:8 & \begin{tabularx}{0.7\textwidth}{X} 至於膽怯的、不信的、可憎的、殺人的、淫亂的、行邪術的、拜偶像的和一切說謊話的人，他們將在燒著硫磺的火湖裡有份；這是第二次的死。」 \end{tabularx} \\ \\ \relax
21:9 & \begin{tabularx}{0.7\textwidth}{X} 拿著七個金碗、盛滿末後七種災禍的七位天使中，有一位來對我說：「你來，我要給你看新娘，就是羔羊的妻子。」 \end{tabularx} \\ \\ \relax
21:10 & \begin{tabularx}{0.7\textwidth}{X} 我在聖靈感動下，天使帶我到一座高大的山，給我看由神那裡、從天而降的聖城耶路撒冷， \end{tabularx} \\ \\ \relax
21:11 & \begin{tabularx}{0.7\textwidth}{X} 這城有神的榮耀，它光輝如同極貴的寶石，好像碧玉，明如水晶。 \end{tabularx} \\ \\ \relax
21:12 & \begin{tabularx}{0.7\textwidth}{X} 它有高大的牆，有十二個門，門上有十二位天使，門上又寫著以色列人十二個支派的名字。 \end{tabularx} \\ \\ \relax
21:13 & \begin{tabularx}{0.7\textwidth}{X} 東邊有三個門，北邊有三個門，南邊有三個門，西邊有三個門。 \end{tabularx} \\ \\ \relax
21:14 & \begin{tabularx}{0.7\textwidth}{X} 城牆有十二個根基，根基上有羔羊十二使徒的名字。 \end{tabularx} \\ \\ \relax
21:15 & \begin{tabularx}{0.7\textwidth}{X} 那對我說話的天使拿著金的蘆葦當尺，要量那城、城門和城牆。 \end{tabularx} \\ \\ \relax
21:16 & \begin{tabularx}{0.7\textwidth}{X} 城是四方的，長寬一樣。天使用蘆葦量那城，共有一萬二千斯他迪，長、寬、高都是一樣。 \end{tabularx} \\ \\ \relax
21:17 & \begin{tabularx}{0.7\textwidth}{X} 他又量了城牆，按著人的尺寸，就是天使的尺寸，共有一百四十四肘。 \end{tabularx} \\ \\ \relax
21:18 & \begin{tabularx}{0.7\textwidth}{X} 牆是碧玉造的；城是純金的，如同明淨的玻璃。 \end{tabularx} \\ \\ \relax
21:19 & \begin{tabularx}{0.7\textwidth}{X} 城牆的根基是用各樣寶石修飾的：第一個根基是碧玉，第二是藍寶石，第三是綠瑪瑙，第四是綠寶石， \end{tabularx} \\ \\ \relax
21:20 & \begin{tabularx}{0.7\textwidth}{X} 第五是紅瑪瑙，第六是紅寶石，第七是黃璧璽，第八是水蒼玉，第九是紅璧璽，第十是翡翠，第十一是紫瑪瑙，第十二是紫晶。 \end{tabularx} \\ \\ \relax
21:21 & \begin{tabularx}{0.7\textwidth}{X} 十二個門是十二顆珍珠；每一個門是一顆珍珠造的。城內的街道是純金的，好像透明的玻璃。 \end{tabularx} \\ \\ \relax
21:22 & \begin{tabularx}{0.7\textwidth}{X} 我沒有看見城內有殿，因主—全能者神和羔羊就是城的殿。 \end{tabularx} \\ \\ \relax
21:23 & \begin{tabularx}{0.7\textwidth}{X} 那城內不用日月光照，因為有神的榮耀光照，又有羔羊為城的燈。 \end{tabularx} \\ \\ \relax
21:24 & \begin{tabularx}{0.7\textwidth}{X} 列國要藉著城的光行走；地上的君王要把自己的榮耀帶給那城。 \end{tabularx} \\ \\ \relax
21:25 & \begin{tabularx}{0.7\textwidth}{X} 城門白晝總不關閉，在那裡沒有黑夜。 \end{tabularx} \\ \\ \relax
21:26 & \begin{tabularx}{0.7\textwidth}{X} 人要將列國的榮耀尊貴帶給那城。 \end{tabularx} \\ \\ \relax
21:27 & \begin{tabularx}{0.7\textwidth}{X} 凡不潔淨的，和那行可憎與虛謊之事的人，都不得進那城，只有名字寫在羔羊生命冊上的才得進去。 \end{tabularx} \\ \\ \relax
22:1 & \begin{tabularx}{0.7\textwidth}{X} 天使又讓我看一道生命水的河，明亮如水晶，從神和羔羊的寶座流出來， \end{tabularx} \\ \\ \relax
22:2 & \begin{tabularx}{0.7\textwidth}{X} 經過城內街道的中央；在河的兩邊有生命樹，結十二樣 的果子，每月都結果子；樹上的葉子可作醫治萬民之用。 \end{tabularx} \\ \\ \relax
22:3 & \begin{tabularx}{0.7\textwidth}{X} 以後不再有任何詛咒。在城裡將有神和羔羊的寶座。他的僕人都要事奉他， \end{tabularx} \\ \\ \relax
22:4 & \begin{tabularx}{0.7\textwidth}{X} 也要見他的面。他的名字將寫在他們的額上。 \end{tabularx} \\ \\ \relax
22:5 & \begin{tabularx}{0.7\textwidth}{X} 不再有黑夜；他們也不需要燈光或日光，因為主神要光照他們。他們要作王，直到永永遠遠。 \end{tabularx} \\ \\ \relax
22:6 & \begin{tabularx}{0.7\textwidth}{X} 天使又對我說：「這些話是可信靠的，是真實的。主，就是賜靈感給眾先知的神，差遣他的使者，要將必須快要發生的事指示他的眾僕人。」 \end{tabularx} \\ \\ \relax
22:7 & \begin{tabularx}{0.7\textwidth}{X} 「看哪，我必快來！凡遵守這書上預言的有福了。」 \end{tabularx} \\ \\ \relax
22:8 & \begin{tabularx}{0.7\textwidth}{X} 這些事是我－約翰所聽見所看見的。當我聽見看見時，就俯伏在指示我的天使腳前要拜他。 \end{tabularx} \\ \\ \relax
22:9 & \begin{tabularx}{0.7\textwidth}{X} 他對我說：「千萬不可！我與你和你的弟兄眾先知，以及那些守這書上的話的人，同是作僕人。你要敬拜神。」 \end{tabularx} \\ \\ \relax
22:10 & \begin{tabularx}{0.7\textwidth}{X} 他又對我說：「不可封了這書上的預言，因為時候近了。 \end{tabularx} \\ \\ \relax
22:11 & \begin{tabularx}{0.7\textwidth}{X} 不義的，讓他仍舊不義；污穢的，讓他仍舊污穢；為義的，讓他仍舊為義；聖潔的，讓他仍舊聖潔。」 \end{tabularx} \\ \\ \relax
22:12 & \begin{tabularx}{0.7\textwidth}{X} 「看哪，我必快來！賞罰在我，要照每個人所行的報應他。 \end{tabularx} \\ \\ \relax
22:13 & \begin{tabularx}{0.7\textwidth}{X} 我是阿拉法，我是俄梅戛；我是首先的，我是末後的；我是開始，我是終結。」 \end{tabularx} \\ \\ \relax
22:14 & \begin{tabularx}{0.7\textwidth}{X} 那些洗淨自己衣服的有福了！他們可得權柄到生命樹那裡，也能從門進城。 \end{tabularx} \\ \\ \relax
22:15 & \begin{tabularx}{0.7\textwidth}{X} 城外有犬類、行邪術的、淫亂的、殺人的、拜偶像的，以及所有喜愛和行虛謊的人。 \end{tabularx} \\ \\ \relax
22:16 & \begin{tabularx}{0.7\textwidth}{X} 「我－耶穌差遣我的使者，為了眾教會向你們證明這些事。我是大衛的根，是他的後裔；我是明亮的晨星。」 \end{tabularx} \\ \\ \relax
22:17 & \begin{tabularx}{0.7\textwidth}{X} 聖靈和新娘都說：「來！」聽見的人也要說：「來！」口渴的人也要來，願意的人都可以白白取生命的水喝。 \end{tabularx} \\ \\ \relax
22:18 & \begin{tabularx}{0.7\textwidth}{X} 我警告一切聽見這書上預言的人：若有人在這預言上加添甚麼，神必將記在這書上的災禍加在他身上。 \end{tabularx} \\ \\ \relax
22:19 & \begin{tabularx}{0.7\textwidth}{X} 這書上的預言，若有人刪去甚麼，神必從這書上所記的生命樹和聖城刪去他的份。 \end{tabularx} \\ \\ \relax
22:20 & \begin{tabularx}{0.7\textwidth}{X} 證明這些事的說：「是的，我必快來！」阿們！主耶穌啊，我願你來！ \end{tabularx} \\ \\ \relax
22:21 & \begin{tabularx}{0.7\textwidth}{X} 願主耶穌的恩惠與眾聖徒同在。阿們！ \end{tabularx} \\ \\
[1ex]
\hline
\hline
\end{longtable}


\newpage

\section{第66屆~港九培靈會~韋理信~第一講}
\label{sec:964}
\textbf{在動盪轉變的社會中需要一個至善至美的標準}
\newline
\newline
連結: \href{https://www.hkbibleconference.org/session-message/view/964}{\texttt{https://www.hkbibleconference.org/session-message/view/964}}
\newline
\newline
\hyperref[sec:963]{< < < PREV SERMON < < <}
~
\hyperref[sec:toc]{[返目錄]}
~
\hyperref[sec:965]{> > > NEXT SERMON > > >}
\newline
\newline
彼得後書 1:3-1:3
\newline
\begin{longtable}{cl}
\hline
\hline
章節 & 經文 (和合本修訂版)\\
\hline
1:3 & \begin{tabularx}{0.7\textwidth}{X} 神的神能已把一切關乎生命和虔敬的事賜給我們，因我們認識那用自己榮耀和美德召我們的神。 \end{tabularx} \\ \\
[1ex]
\hline
\hline
\end{longtable}


\newpage

\section{第66屆~港九培靈會~韋理信~第二講}
\label{sec:965}
\textbf{至善至美的標準之基礎 — 神寶貴極大的應許}
\newline
\newline
連結: \href{https://www.hkbibleconference.org/session-message/view/965}{\texttt{https://www.hkbibleconference.org/session-message/view/965}}
\newline
\newline
\hyperref[sec:964]{< < < PREV SERMON < < <}
~
\hyperref[sec:toc]{[返目錄]}
~
\hyperref[sec:966]{> > > NEXT SERMON > > >}
\newline
\newline
彼得後書 1:4-1:4
\newline
\begin{longtable}{cl}
\hline
\hline
章節 & 經文 (和合本修訂版)\\
\hline
1:4 & \begin{tabularx}{0.7\textwidth}{X} 因此，他已把又寶貴又極大的應許賜給我們，使我們既脫離世上從情慾來的敗壞，就得分享神的本性。 \end{tabularx} \\ \\
[1ex]
\hline
\hline
\end{longtable}


\newpage

\section{第66屆~港九培靈會~韋理信~第三講}
\label{sec:966}
\textbf{在信心德行上至善至美的標準}
\newline
\newline
連結: \href{https://www.hkbibleconference.org/session-message/view/966}{\texttt{https://www.hkbibleconference.org/session-message/view/966}}
\newline
\newline
\hyperref[sec:965]{< < < PREV SERMON < < <}
~
\hyperref[sec:toc]{[返目錄]}
~
\hyperref[sec:967]{> > > NEXT SERMON > > >}
\newline
\newline
提摩太後書 1:5-1:5
\newline
\begin{longtable}{cl}
\hline
\hline
章節 & 經文 (和合本修訂版)\\
\hline
1:5 & \begin{tabularx}{0.7\textwidth}{X} 我記得你無偽的信心，這信心先存在你外祖母羅以和你母親友妮基的心裡，我深信也存在你的心裡。 \end{tabularx} \\ \\
[1ex]
\hline
\hline
\end{longtable}


\newpage

\section{第66屆~港九培靈會~韋理信~第四講}
\label{sec:967}
\textbf{在知識上至善至美的標準}
\newline
\newline
連結: \href{https://www.hkbibleconference.org/session-message/view/967}{\texttt{https://www.hkbibleconference.org/session-message/view/967}}
\newline
\newline
\hyperref[sec:966]{< < < PREV SERMON < < <}
~
\hyperref[sec:toc]{[返目錄]}
~
\hyperref[sec:968]{> > > NEXT SERMON > > >}
\newline
\newline
彼得後書 1:5-1:5
\newline
\begin{longtable}{cl}
\hline
\hline
章節 & 經文 (和合本修訂版)\\
\hline
1:5 & \begin{tabularx}{0.7\textwidth}{X} 正因這緣故，你們要分外地努力。有了信心，又要加上德行；有了德行，又要加上知識； \end{tabularx} \\ \\
[1ex]
\hline
\hline
\end{longtable}


\newpage

\section{第66屆~港九培靈會~韋理信~第五講}
\label{sec:968}
\textbf{在節制上至善至美的標準}
\newline
\newline
連結: \href{https://www.hkbibleconference.org/session-message/view/968}{\texttt{https://www.hkbibleconference.org/session-message/view/968}}
\newline
\newline
\hyperref[sec:967]{< < < PREV SERMON < < <}
~
\hyperref[sec:toc]{[返目錄]}
~
\hyperref[sec:969]{> > > NEXT SERMON > > >}
\newline
\newline
彼得後書 1:6-1:6
\newline
\begin{longtable}{cl}
\hline
\hline
章節 & 經文 (和合本修訂版)\\
\hline
1:6 & \begin{tabularx}{0.7\textwidth}{X} 有了知識，又要加上節制；有了節制，又要加上忍耐；有了忍耐，又要加上虔敬； \end{tabularx} \\ \\
[1ex]
\hline
\hline
\end{longtable}


\newpage

\section{第66屆~港九培靈會~韋理信~第六講}
\label{sec:969}
\textbf{在忍耐上至善至美的標準}
\newline
\newline
連結: \href{https://www.hkbibleconference.org/session-message/view/969}{\texttt{https://www.hkbibleconference.org/session-message/view/969}}
\newline
\newline
\hyperref[sec:968]{< < < PREV SERMON < < <}
~
\hyperref[sec:toc]{[返目錄]}
~
\hyperref[sec:970]{> > > NEXT SERMON > > >}
\newline
\newline
彼得後書 1:6-1:6
\newline
\begin{longtable}{cl}
\hline
\hline
章節 & 經文 (和合本修訂版)\\
\hline
1:6 & \begin{tabularx}{0.7\textwidth}{X} 有了知識，又要加上節制；有了節制，又要加上忍耐；有了忍耐，又要加上虔敬； \end{tabularx} \\ \\
[1ex]
\hline
\hline
\end{longtable}


\newpage

\section{第66屆~港九培靈會~韋理信~第七講}
\label{sec:970}
\textbf{在虔敬上至善至美的標準}
\newline
\newline
連結: \href{https://www.hkbibleconference.org/session-message/view/970}{\texttt{https://www.hkbibleconference.org/session-message/view/970}}
\newline
\newline
\hyperref[sec:969]{< < < PREV SERMON < < <}
~
\hyperref[sec:toc]{[返目錄]}
~
\hyperref[sec:971]{> > > NEXT SERMON > > >}
\newline
\newline
彼得後書 1:6-1:6
\newline
\begin{longtable}{cl}
\hline
\hline
章節 & 經文 (和合本修訂版)\\
\hline
1:6 & \begin{tabularx}{0.7\textwidth}{X} 有了知識，又要加上節制；有了節制，又要加上忍耐；有了忍耐，又要加上虔敬； \end{tabularx} \\ \\
[1ex]
\hline
\hline
\end{longtable}


\newpage

\section{第66屆~港九培靈會~韋理信~第八講}
\label{sec:971}
\textbf{在愛弟兄和愛眾人上至善至美的標準}
\newline
\newline
連結: \href{https://www.hkbibleconference.org/session-message/view/971}{\texttt{https://www.hkbibleconference.org/session-message/view/971}}
\newline
\newline
\hyperref[sec:970]{< < < PREV SERMON < < <}
~
\hyperref[sec:toc]{[返目錄]}
~
\hyperref[sec:972]{> > > NEXT SERMON > > >}
\newline
\newline
彼得後書 1:7-1:7
\newline
\begin{longtable}{cl}
\hline
\hline
章節 & 經文 (和合本修訂版)\\
\hline
1:7 & \begin{tabularx}{0.7\textwidth}{X} 有了虔敬，又要加上愛弟兄的心；有了愛弟兄的心，又要加上愛眾人的心。 \end{tabularx} \\ \\
[1ex]
\hline
\hline
\end{longtable}


\newpage

\section{第66屆~港九培靈會~韋理信~第九講}
\label{sec:972}
\textbf{如何達到至善至美的標準}
\newline
\newline
連結: \href{https://www.hkbibleconference.org/session-message/view/972}{\texttt{https://www.hkbibleconference.org/session-message/view/972}}
\newline
\newline
\hyperref[sec:971]{< < < PREV SERMON < < <}
~
\hyperref[sec:toc]{[返目錄]}
~
\hyperref[sec:945]{> > > NEXT SERMON > > >}
\newline
\newline
彼得後書 1:8-1:11
\newline
\begin{longtable}{cl}
\hline
\hline
章節 & 經文 (和合本修訂版)\\
\hline
1:8 & \begin{tabularx}{0.7\textwidth}{X} 你們有了這幾樣，再繼續增長，就必使你們在認識我們的主耶穌基督上，不至於懶散和不結果子了。 \end{tabularx} \\ \\ \relax
1:9 & \begin{tabularx}{0.7\textwidth}{X} 沒有這幾樣的人就是瞎眼，是短視，忘了他過去的罪已經得了潔淨。 \end{tabularx} \\ \\ \relax
1:10 & \begin{tabularx}{0.7\textwidth}{X} 所以，弟兄們，要更加努力，使你們的蒙召和被選堅定不移。你們實行這幾樣，就永不失腳。 \end{tabularx} \\ \\ \relax
1:11 & \begin{tabularx}{0.7\textwidth}{X} 這樣，必叫你們豐豐富富地得以進入我們主－救主耶穌基督永遠的國度。 \end{tabularx} \\ \\
[1ex]
\hline
\hline
\end{longtable}


\newpage

\section{第67屆~港九培靈會~滕近輝~第一講}
\label{sec:945}
\textbf{國度真理的重要性}
\newline
\newline
連結: \href{https://www.hkbibleconference.org/session-message/view/945}{\texttt{https://www.hkbibleconference.org/session-message/view/945}}
\newline
\newline
\hyperref[sec:972]{< < < PREV SERMON < < <}
~
\hyperref[sec:toc]{[返目錄]}
~
\hyperref[sec:946]{> > > NEXT SERMON > > >}
\newline
\newline


\newpage

\section{第67屆~港九培靈會~滕近輝~第二講}
\label{sec:946}
\textbf{舊約的國度}
\newline
\newline
連結: \href{https://www.hkbibleconference.org/session-message/view/946}{\texttt{https://www.hkbibleconference.org/session-message/view/946}}
\newline
\newline
\hyperref[sec:945]{< < < PREV SERMON < < <}
~
\hyperref[sec:toc]{[返目錄]}
~
\hyperref[sec:947]{> > > NEXT SERMON > > >}
\newline
\newline


\newpage

\section{第67屆~港九培靈會~滕近輝~第三講}
\label{sec:947}
\textbf{國度的奧秘}
\newline
\newline
連結: \href{https://www.hkbibleconference.org/session-message/view/947}{\texttt{https://www.hkbibleconference.org/session-message/view/947}}
\newline
\newline
\hyperref[sec:946]{< < < PREV SERMON < < <}
~
\hyperref[sec:toc]{[返目錄]}
~
\hyperref[sec:948]{> > > NEXT SERMON > > >}
\newline
\newline


\newpage

\section{第67屆~港九培靈會~滕近輝~第四講}
\label{sec:948}
\textbf{國度的原則}
\newline
\newline
連結: \href{https://www.hkbibleconference.org/session-message/view/948}{\texttt{https://www.hkbibleconference.org/session-message/view/948}}
\newline
\newline
\hyperref[sec:947]{< < < PREV SERMON < < <}
~
\hyperref[sec:toc]{[返目錄]}
~
\hyperref[sec:949]{> > > NEXT SERMON > > >}
\newline
\newline


\newpage

\section{第67屆~港九培靈會~滕近輝~第五講}
\label{sec:949}
\textbf{國度的擴展}
\newline
\newline
連結: \href{https://www.hkbibleconference.org/session-message/view/949}{\texttt{https://www.hkbibleconference.org/session-message/view/949}}
\newline
\newline
\hyperref[sec:948]{< < < PREV SERMON < < <}
~
\hyperref[sec:toc]{[返目錄]}
~
\hyperref[sec:950]{> > > NEXT SERMON > > >}
\newline
\newline


\newpage

\section{第67屆~港九培靈會~滕近輝~第六講}
\label{sec:950}
\textbf{國度的勝利}
\newline
\newline
連結: \href{https://www.hkbibleconference.org/session-message/view/950}{\texttt{https://www.hkbibleconference.org/session-message/view/950}}
\newline
\newline
\hyperref[sec:949]{< < < PREV SERMON < < <}
~
\hyperref[sec:toc]{[返目錄]}
~
\hyperref[sec:951]{> > > NEXT SERMON > > >}
\newline
\newline


\newpage

\section{第67屆~港九培靈會~滕近輝~第七講}
\label{sec:951}
\textbf{國度的價值觀}
\newline
\newline
連結: \href{https://www.hkbibleconference.org/session-message/view/951}{\texttt{https://www.hkbibleconference.org/session-message/view/951}}
\newline
\newline
\hyperref[sec:950]{< < < PREV SERMON < < <}
~
\hyperref[sec:toc]{[返目錄]}
~
\hyperref[sec:952]{> > > NEXT SERMON > > >}
\newline
\newline


\newpage

\section{第67屆~港九培靈會~滕近輝~第八講}
\label{sec:952}
\textbf{國度的本體與外延}
\newline
\newline
連結: \href{https://www.hkbibleconference.org/session-message/view/952}{\texttt{https://www.hkbibleconference.org/session-message/view/952}}
\newline
\newline
\hyperref[sec:951]{< < < PREV SERMON < < <}
~
\hyperref[sec:toc]{[返目錄]}
~
\hyperref[sec:953]{> > > NEXT SERMON > > >}
\newline
\newline


\newpage

\section{第67屆~港九培靈會~滕近輝~第九講}
\label{sec:953}
\textbf{國度與得救的異同}
\newline
\newline
連結: \href{https://www.hkbibleconference.org/session-message/view/953}{\texttt{https://www.hkbibleconference.org/session-message/view/953}}
\newline
\newline
\hyperref[sec:952]{< < < PREV SERMON < < <}
~
\hyperref[sec:toc]{[返目錄]}
~
\hyperref[sec:936]{> > > NEXT SERMON > > >}
\newline
\newline


\newpage

\section{第67屆~港九培靈會~焦源濂~第一講}
\label{sec:936}
\textbf{受苦之前的約伯 — 幸福的真諦}
\newline
\newline
連結: \href{https://www.hkbibleconference.org/session-message/view/936}{\texttt{https://www.hkbibleconference.org/session-message/view/936}}
\newline
\newline
\hyperref[sec:953]{< < < PREV SERMON < < <}
~
\hyperref[sec:toc]{[返目錄]}
~
\hyperref[sec:937]{> > > NEXT SERMON > > >}
\newline
\newline


\newpage

\section{第67屆~港九培靈會~焦源濂~第二講}
\label{sec:937}
\textbf{大禍臨頭的約伯 — 靈界的奧秘}
\newline
\newline
連結: \href{https://www.hkbibleconference.org/session-message/view/937}{\texttt{https://www.hkbibleconference.org/session-message/view/937}}
\newline
\newline
\hyperref[sec:936]{< < < PREV SERMON < < <}
~
\hyperref[sec:toc]{[返目錄]}
~
\hyperref[sec:938]{> > > NEXT SERMON > > >}
\newline
\newline


\newpage

\section{第67屆~港九培靈會~焦源濂~第三講}
\label{sec:938}
\textbf{面對災難的約伯 — 真人的真像 (一)}
\newline
\newline
連結: \href{https://www.hkbibleconference.org/session-message/view/938}{\texttt{https://www.hkbibleconference.org/session-message/view/938}}
\newline
\newline
\hyperref[sec:937]{< < < PREV SERMON < < <}
~
\hyperref[sec:toc]{[返目錄]}
~
\hyperref[sec:939]{> > > NEXT SERMON > > >}
\newline
\newline


\newpage

\section{第67屆~港九培靈會~焦源濂~第四講}
\label{sec:939}
\textbf{面對災難的約伯 — 真人的真像 (二)}
\newline
\newline
連結: \href{https://www.hkbibleconference.org/session-message/view/939}{\texttt{https://www.hkbibleconference.org/session-message/view/939}}
\newline
\newline
\hyperref[sec:938]{< < < PREV SERMON < < <}
~
\hyperref[sec:toc]{[返目錄]}
~
\hyperref[sec:940]{> > > NEXT SERMON > > >}
\newline
\newline


\newpage

\section{第67屆~港九培靈會~焦源濂~第五講}
\label{sec:940}
\textbf{面對災難的約伯 — 真人的真像 (三)}
\newline
\newline
連結: \href{https://www.hkbibleconference.org/session-message/view/940}{\texttt{https://www.hkbibleconference.org/session-message/view/940}}
\newline
\newline
\hyperref[sec:939]{< < < PREV SERMON < < <}
~
\hyperref[sec:toc]{[返目錄]}
~
\hyperref[sec:941]{> > > NEXT SERMON > > >}
\newline
\newline


\newpage

\section{第67屆~港九培靈會~焦源濂~第六講}
\label{sec:941}
\textbf{面對災難的約伯 — 真人的真像 (四)}
\newline
\newline
連結: \href{https://www.hkbibleconference.org/session-message/view/941}{\texttt{https://www.hkbibleconference.org/session-message/view/941}}
\newline
\newline
\hyperref[sec:940]{< < < PREV SERMON < < <}
~
\hyperref[sec:toc]{[返目錄]}
~
\hyperref[sec:942]{> > > NEXT SERMON > > >}
\newline
\newline


\newpage

\section{第67屆~港九培靈會~焦源濂~第七講}
\label{sec:942}
\textbf{面對災難的約伯 — 家庭的危機}
\newline
\newline
連結: \href{https://www.hkbibleconference.org/session-message/view/942}{\texttt{https://www.hkbibleconference.org/session-message/view/942}}
\newline
\newline
\hyperref[sec:941]{< < < PREV SERMON < < <}
~
\hyperref[sec:toc]{[返目錄]}
~
\hyperref[sec:943]{> > > NEXT SERMON > > >}
\newline
\newline


\newpage

\section{第67屆~港九培靈會~焦源濂~第八講}
\label{sec:943}
\textbf{知友陪伴的約伯 — 肢體生活的榜樣}
\newline
\newline
連結: \href{https://www.hkbibleconference.org/session-message/view/943}{\texttt{https://www.hkbibleconference.org/session-message/view/943}}
\newline
\newline
\hyperref[sec:942]{< < < PREV SERMON < < <}
~
\hyperref[sec:toc]{[返目錄]}
~
\hyperref[sec:944]{> > > NEXT SERMON > > >}
\newline
\newline


\newpage

\section{第67屆~港九培靈會~焦源濂~第九講}
\label{sec:944}
\textbf{知友陪伴的約伯 — 友誼的珍貴}
\newline
\newline
連結: \href{https://www.hkbibleconference.org/session-message/view/944}{\texttt{https://www.hkbibleconference.org/session-message/view/944}}
\newline
\newline
\hyperref[sec:943]{< < < PREV SERMON < < <}
~
\hyperref[sec:toc]{[返目錄]}
~
\hyperref[sec:927]{> > > NEXT SERMON > > >}
\newline
\newline


\newpage

\section{第67屆~港九培靈會~鮑會園~第一講}
\label{sec:927}
\textbf{福音的託付}
\newline
\newline
連結: \href{https://www.hkbibleconference.org/session-message/view/927}{\texttt{https://www.hkbibleconference.org/session-message/view/927}}
\newline
\newline
\hyperref[sec:944]{< < < PREV SERMON < < <}
~
\hyperref[sec:toc]{[返目錄]}
~
\hyperref[sec:928]{> > > NEXT SERMON > > >}
\newline
\newline
羅馬書 1:1-1:17
\newline
\begin{longtable}{cl}
\hline
\hline
章節 & 經文 (和合本修訂版)\\
\hline
1:1 & \begin{tabularx}{0.7\textwidth}{X} 基督耶穌的僕人保羅，蒙召為使徒，奉派傳神的福音。 \end{tabularx} \\ \\ \relax
1:2 & \begin{tabularx}{0.7\textwidth}{X} 這福音是神從前藉眾先知，在聖經上所應許的。 \end{tabularx} \\ \\ \relax
1:3 & \begin{tabularx}{0.7\textwidth}{X} 論到他兒子－我主耶穌基督，按肉體說，是從大衛後裔生的； \end{tabularx} \\ \\ \relax
1:4 & \begin{tabularx}{0.7\textwidth}{X} 按神聖的靈說，因從死人中復活，用大能顯明他是神的兒子。 \end{tabularx} \\ \\ \relax
1:5 & \begin{tabularx}{0.7\textwidth}{X} 我們從他蒙恩受了使徒的職分，為他的名在萬國中使人因信而順服， \end{tabularx} \\ \\ \relax
1:6 & \begin{tabularx}{0.7\textwidth}{X} 其中也有你們這蒙召屬耶穌基督的人。 \end{tabularx} \\ \\ \relax
1:7 & \begin{tabularx}{0.7\textwidth}{X} 我寫信給你們在羅馬、為神所愛、蒙召作聖徒的眾人。願恩惠、平安從我們的父神和主耶穌基督歸給你們！ \end{tabularx} \\ \\ \relax
1:8 & \begin{tabularx}{0.7\textwidth}{X} 首先，我靠著耶穌基督，為你們眾人感謝我的神，因你們的信德傳遍了天下。 \end{tabularx} \\ \\ \relax
1:9 & \begin{tabularx}{0.7\textwidth}{X} 我在他兒子的福音上，用心靈所事奉的神可以見證，我怎樣不住地提到你們， \end{tabularx} \\ \\ \relax
1:10 & \begin{tabularx}{0.7\textwidth}{X} 在我的禱告中常常懇求，或許照神的旨意，最終我能毫無阻礙地往你們那裡去。 \end{tabularx} \\ \\ \relax
1:11 & \begin{tabularx}{0.7\textwidth}{X} 因為我迫切地想見你們，要把一些屬靈的恩賜分給你們，使你們得以堅固， \end{tabularx} \\ \\ \relax
1:12 & \begin{tabularx}{0.7\textwidth}{X} 也可以說，我在你們中間，因你我彼此的信心而同得安慰。 \end{tabularx} \\ \\ \relax
1:13 & \begin{tabularx}{0.7\textwidth}{X} 弟兄們，我不願意你們不知道，我屢次計劃往你們那裡去，要在你們中間得些果子，如同在其餘的外邦人中一樣，只是到如今仍有攔阻。 \end{tabularx} \\ \\ \relax
1:14 & \begin{tabularx}{0.7\textwidth}{X} 無論是希臘人、未開化的人、聰明人、愚拙人，我都欠他們的債， \end{tabularx} \\ \\ \relax
1:15 & \begin{tabularx}{0.7\textwidth}{X} 所以願意盡我的力量把福音也傳給你們在羅馬的人。 \end{tabularx} \\ \\ \relax
1:16 & \begin{tabularx}{0.7\textwidth}{X} 我不以福音為恥；這福音本是神的大能，要救一切相信的，先是猶太人，後是希臘人。 \end{tabularx} \\ \\ \relax
1:17 & \begin{tabularx}{0.7\textwidth}{X} 因為神的義正在這福音上顯明出來；這義是本於信，以至於信。如經上所記：「義人必因信得生。」 \end{tabularx} \\ \\
[1ex]
\hline
\hline
\end{longtable}


\newpage

\section{第67屆~港九培靈會~鮑會園~第二講}
\label{sec:928}
\textbf{世人的需要}
\newline
\newline
連結: \href{https://www.hkbibleconference.org/session-message/view/928}{\texttt{https://www.hkbibleconference.org/session-message/view/928}}
\newline
\newline
\hyperref[sec:927]{< < < PREV SERMON < < <}
~
\hyperref[sec:toc]{[返目錄]}
~
\hyperref[sec:929]{> > > NEXT SERMON > > >}
\newline
\newline
羅馬書 1:18-2:15
\newline
\begin{longtable}{cl}
\hline
\hline
章節 & 經文 (和合本修訂版)\\
\hline
1:18 & \begin{tabularx}{0.7\textwidth}{X} 原來，神的憤怒從天上顯明在一切不虔不義的人身上，就是那些行不義壓制真理的人。 \end{tabularx} \\ \\ \relax
1:19 & \begin{tabularx}{0.7\textwidth}{X} 神的事情，人所能知道的，原顯明在人心裡，因為神已經向他們顯明。 \end{tabularx} \\ \\ \relax
1:20 & \begin{tabularx}{0.7\textwidth}{X} 自從造天地以來，神的永能和神性是明明可知的，雖然眼不能見，但藉著所造之物就可以了解看見，叫人無可推諉。 \end{tabularx} \\ \\ \relax
1:21 & \begin{tabularx}{0.7\textwidth}{X} 因為，他們雖然知道神，卻不把他當作神榮耀他，也不感謝他。他們的思想變為虛妄，無知的心昏暗了。 \end{tabularx} \\ \\ \relax
1:22 & \begin{tabularx}{0.7\textwidth}{X} 他們自以為聰明，反成了愚昧， \end{tabularx} \\ \\ \relax
1:23 & \begin{tabularx}{0.7\textwidth}{X} 將不能朽壞之神的榮耀變為偶像，仿照必朽壞的人、飛禽、走獸、爬蟲的形像。 \end{tabularx} \\ \\ \relax
1:24 & \begin{tabularx}{0.7\textwidth}{X} 所以，神任憑他們隨著心裡的情慾行污穢的事，以致彼此羞辱自己的身體。 \end{tabularx} \\ \\ \relax
1:25 & \begin{tabularx}{0.7\textwidth}{X} 他們將神的真實變為虛謊，去敬拜事奉受造之物，不敬奉那造物的主—主是可稱頌的，直到永遠。阿們！ \end{tabularx} \\ \\ \relax
1:26 & \begin{tabularx}{0.7\textwidth}{X} 因此，神任憑他們放縱可羞恥的情慾。他們的女人把自然的關係變成違反自然的； \end{tabularx} \\ \\ \relax
1:27 & \begin{tabularx}{0.7\textwidth}{X} 男人也是如此，放棄了和女人自然的關係，慾火攻心，男的和男的彼此貪戀，行可恥的事，就在自己身上受這逆性行為當得的報應。 \end{tabularx} \\ \\ \relax
1:28 & \begin{tabularx}{0.7\textwidth}{X} 他們既然故意不認識神，神就任憑他們存扭曲的心，做那些不該做的事， \end{tabularx} \\ \\ \relax
1:29 & \begin{tabularx}{0.7\textwidth}{X} 裝滿了各樣不義、邪惡、貪婪、惡毒，滿心是嫉妒、兇殺、紛爭、詭詐、毒恨，又是毀謗的、 \end{tabularx} \\ \\ \relax
1:30 & \begin{tabularx}{0.7\textwidth}{X} 說人壞話的、怨恨神的、侮辱人的、狂傲的、自誇的、製造是非的、忤逆父母的、 \end{tabularx} \\ \\ \relax
1:31 & \begin{tabularx}{0.7\textwidth}{X} 頑梗不化的、言而無信的、無情無義的、不憐憫人的。 \end{tabularx} \\ \\ \relax
1:32 & \begin{tabularx}{0.7\textwidth}{X} 他們雖知道神判定做這樣事的人是該死的，然而他們不但自己去做，還贊同別人去做。 \end{tabularx} \\ \\ \relax
2:1 & \begin{tabularx}{0.7\textwidth}{X} 所以，你這評斷人的人哪，無論你是誰，都無可推諉。你在甚麼事上評斷人，就在甚麼事上定自己的罪。因你這評斷人的，自己所做的卻和別人一樣。 \end{tabularx} \\ \\ \relax
2:2 & \begin{tabularx}{0.7\textwidth}{X} 我們知道這樣做的人，神必公平地審判他。 \end{tabularx} \\ \\ \relax
2:3 & \begin{tabularx}{0.7\textwidth}{X} 你這個人哪，你評斷做這樣事的人，自己所做的卻和別人一樣，你以為能逃脫神的審判嗎？ \end{tabularx} \\ \\ \relax
2:4 & \begin{tabularx}{0.7\textwidth}{X} 還是你藐視他豐富的恩慈、寬容、忍耐，不知道他的恩慈是領你悔改嗎？ \end{tabularx} \\ \\ \relax
2:5 & \begin{tabularx}{0.7\textwidth}{X} 你竟放任你剛硬不悔改的心，為自己累積憤怒！在憤怒的日子，神公義的審判要顯示出來。 \end{tabularx} \\ \\ \relax
2:6 & \begin{tabularx}{0.7\textwidth}{X} 他要照各人的行為報應各人。 \end{tabularx} \\ \\ \relax
2:7 & \begin{tabularx}{0.7\textwidth}{X} 凡恆心行善，尋求榮耀、尊貴和不能朽壞的，就有永生報償他們； \end{tabularx} \\ \\ \relax
2:8 & \begin{tabularx}{0.7\textwidth}{X} 但是那些自私自利、不順從真理、反順從不義的人，就有惱恨、憤怒報應他們。 \end{tabularx} \\ \\ \relax
2:9 & \begin{tabularx}{0.7\textwidth}{X} 他要把患難、困苦加給一切作惡的人，先是猶太人，後是希臘人； \end{tabularx} \\ \\ \relax
2:10 & \begin{tabularx}{0.7\textwidth}{X} 卻把榮耀、尊貴、平安加給一切行善的人，先是猶太人，後是希臘人。 \end{tabularx} \\ \\ \relax
2:11 & \begin{tabularx}{0.7\textwidth}{X} 因為神不偏待人。 \end{tabularx} \\ \\ \relax
2:12 & \begin{tabularx}{0.7\textwidth}{X} 凡在律法之外犯了罪的，將在律法之外滅亡；凡在律法之內犯了罪的，將按律法受審判。 \end{tabularx} \\ \\ \relax
2:13 & \begin{tabularx}{0.7\textwidth}{X} 原來在神面前，不是聽律法的為義，而是行律法的稱義。 \end{tabularx} \\ \\ \relax
2:14 & \begin{tabularx}{0.7\textwidth}{X} 沒有律法的外邦人若順著本性行律法上的事，他們雖然沒有律法，自己就是自己的律法。 \end{tabularx} \\ \\ \relax
2:15 & \begin{tabularx}{0.7\textwidth}{X} 他們顯明律法的功用刻在他們心裡，他們的良心一同作證—他們的內心掙扎，有時自責，有時為自己辯護。 \end{tabularx} \\ \\
[1ex]
\hline
\hline
\end{longtable}


\newpage

\section{第67屆~港九培靈會~鮑會園~第三講}
\label{sec:929}
\textbf{救恩的完成}
\newline
\newline
連結: \href{https://www.hkbibleconference.org/session-message/view/929}{\texttt{https://www.hkbibleconference.org/session-message/view/929}}
\newline
\newline
\hyperref[sec:928]{< < < PREV SERMON < < <}
~
\hyperref[sec:toc]{[返目錄]}
~
\hyperref[sec:930]{> > > NEXT SERMON > > >}
\newline
\newline
羅馬書 3:21-3:31
\newline
\begin{longtable}{cl}
\hline
\hline
章節 & 經文 (和合本修訂版)\\
\hline
3:21 & \begin{tabularx}{0.7\textwidth}{X} 但如今，神的義在律法之外已經顯明出來，有律法和先知為證： \end{tabularx} \\ \\ \relax
3:22 & \begin{tabularx}{0.7\textwidth}{X} 就是神的義，因信耶穌基督加給一切信的人。這並沒有分別， \end{tabularx} \\ \\ \relax
3:23 & \begin{tabularx}{0.7\textwidth}{X} 因為世人都犯了罪，虧缺了神的榮耀， \end{tabularx} \\ \\ \relax
3:24 & \begin{tabularx}{0.7\textwidth}{X} 如今卻蒙神的恩典，藉著在基督耶穌裡的救贖，就白白地得稱為義。 \end{tabularx} \\ \\ \relax
3:25 & \begin{tabularx}{0.7\textwidth}{X} 神設立耶穌作贖罪祭，是憑耶穌的血，藉著信，要顯明神的義；因為他用忍耐的心寬容人先前所犯的罪， \end{tabularx} \\ \\ \relax
3:26 & \begin{tabularx}{0.7\textwidth}{X} 好使今時顯明他的義，讓人知道他自己為義，也稱信耶穌的人為義。 \end{tabularx} \\ \\ \relax
3:27 & \begin{tabularx}{0.7\textwidth}{X} 既是這樣，哪裡可誇口呢？沒有可誇的。是藉甚麼法呢？功德嗎？不是！是藉信主之法。 \end{tabularx} \\ \\ \relax
3:28 & \begin{tabularx}{0.7\textwidth}{X} 所以我們認定，人稱義是因著信，不在於律法的行為。 \end{tabularx} \\ \\ \relax
3:29 & \begin{tabularx}{0.7\textwidth}{X} 難道神只是猶太人的嗎？不也是外邦人的嗎？是的，他也是外邦人的神。 \end{tabularx} \\ \\ \relax
3:30 & \begin{tabularx}{0.7\textwidth}{X} 既然神是一位，他就要本於信稱那受割禮的為義，也要藉著信稱那未受割禮的為義。 \end{tabularx} \\ \\ \relax
3:31 & \begin{tabularx}{0.7\textwidth}{X} 這樣，我們藉著信廢了律法嗎？絕對不是！更是鞏固律法。 \end{tabularx} \\ \\
[1ex]
\hline
\hline
\end{longtable}


\newpage

\section{第67屆~港九培靈會~鮑會園~第四講}
\label{sec:930}
\textbf{救恩的功效}
\newline
\newline
連結: \href{https://www.hkbibleconference.org/session-message/view/930}{\texttt{https://www.hkbibleconference.org/session-message/view/930}}
\newline
\newline
\hyperref[sec:929]{< < < PREV SERMON < < <}
~
\hyperref[sec:toc]{[返目錄]}
~
\hyperref[sec:931]{> > > NEXT SERMON > > >}
\newline
\newline
羅馬書 5:1-5:11
\newline
\begin{longtable}{cl}
\hline
\hline
章節 & 經文 (和合本修訂版)\\
\hline
5:1 & \begin{tabularx}{0.7\textwidth}{X} 所以，我們既因信稱義，就藉著我們的主耶穌基督得以與神和好。 \end{tabularx} \\ \\ \relax
5:2 & \begin{tabularx}{0.7\textwidth}{X} 我們又藉著他，因信得以進入現在所站立的這恩典中，並且歡歡喜喜盼望神的榮耀。 \end{tabularx} \\ \\ \relax
5:3 & \begin{tabularx}{0.7\textwidth}{X} 不但如此，就是在患難中也是歡歡喜喜的，因為知道患難生忍耐， \end{tabularx} \\ \\ \relax
5:4 & \begin{tabularx}{0.7\textwidth}{X} 忍耐生老練，老練生盼望， \end{tabularx} \\ \\ \relax
5:5 & \begin{tabularx}{0.7\textwidth}{X} 盼望不至於落空，因為神的愛，已藉著所賜給我們的聖靈，澆灌在我們心裡。 \end{tabularx} \\ \\ \relax
5:6 & \begin{tabularx}{0.7\textwidth}{X} 我們還軟弱的時候，基督就在特定的時刻為不敬虔之人死。 \end{tabularx} \\ \\ \relax
5:7 & \begin{tabularx}{0.7\textwidth}{X} 為義人死，是少有的；為仁人死，或者有敢做的。 \end{tabularx} \\ \\ \relax
5:8 & \begin{tabularx}{0.7\textwidth}{X} 惟有基督在我們還作罪人的時候為我們死，神的愛就在此向我們顯明了。 \end{tabularx} \\ \\ \relax
5:9 & \begin{tabularx}{0.7\textwidth}{X} 現在我們既靠著他的血稱義，就更要藉著他得救，免受神的憤怒。 \end{tabularx} \\ \\ \relax
5:10 & \begin{tabularx}{0.7\textwidth}{X} 因為我們作仇敵的時候，尚且藉著神兒子的死得以與神和好，既已和好，就更要因他的生得救了。 \end{tabularx} \\ \\ \relax
5:11 & \begin{tabularx}{0.7\textwidth}{X} 不但如此，我們既藉著我們的主耶穌基督得以與神和好，也就藉著他以神為樂。 \end{tabularx} \\ \\
[1ex]
\hline
\hline
\end{longtable}


\newpage

\section{第67屆~港九培靈會~鮑會園~第五講}
\label{sec:931}
\textbf{救恩的結果}
\newline
\newline
連結: \href{https://www.hkbibleconference.org/session-message/view/931}{\texttt{https://www.hkbibleconference.org/session-message/view/931}}
\newline
\newline
\hyperref[sec:930]{< < < PREV SERMON < < <}
~
\hyperref[sec:toc]{[返目錄]}
~
\hyperref[sec:932]{> > > NEXT SERMON > > >}
\newline
\newline
羅馬書 6:1-6:23
\newline
\begin{longtable}{cl}
\hline
\hline
章節 & 經文 (和合本修訂版)\\
\hline
6:1 & \begin{tabularx}{0.7\textwidth}{X} 這樣，我們要怎麼說呢？我們可以仍在罪中使恩典增多嗎？ \end{tabularx} \\ \\ \relax
6:2 & \begin{tabularx}{0.7\textwidth}{X} 絕對不可！我們向罪死了的人，豈可仍在罪中活著呢？ \end{tabularx} \\ \\ \relax
6:3 & \begin{tabularx}{0.7\textwidth}{X} 難道你們不知道，我們這受洗歸入基督耶穌的人，就是受洗歸入他的死嗎？ \end{tabularx} \\ \\ \relax
6:4 & \begin{tabularx}{0.7\textwidth}{X} 所以，我們藉著洗禮歸入死，和他一同埋葬，是要我們行事為人都有新生的樣子，像基督藉著父的榮耀從死人中復活一樣。 \end{tabularx} \\ \\ \relax
6:5 & \begin{tabularx}{0.7\textwidth}{X} 我們若與他合一，經歷與他一樣的死，也將經歷與他一樣的復活。 \end{tabularx} \\ \\ \relax
6:6 & \begin{tabularx}{0.7\textwidth}{X} 我們知道，我們的舊人和他同釘十字架，使罪身滅絕，叫我們不再作罪的奴隸， \end{tabularx} \\ \\ \relax
6:7 & \begin{tabularx}{0.7\textwidth}{X} 因為已死的人是脫離了罪。 \end{tabularx} \\ \\ \relax
6:8 & \begin{tabularx}{0.7\textwidth}{X} 我們若與基督同死，我們信也必與他同活， \end{tabularx} \\ \\ \relax
6:9 & \begin{tabularx}{0.7\textwidth}{X} 因為知道基督既從死人中復活，就不再死，死也不再作他的主了。 \end{tabularx} \\ \\ \relax
6:10 & \begin{tabularx}{0.7\textwidth}{X} 他死了，是對罪死，只這一次；他活，是對神活著。 \end{tabularx} \\ \\ \relax
6:11 & \begin{tabularx}{0.7\textwidth}{X} 這樣，你們也要看自己對罪是死的，在基督耶穌裡對神卻是活的。 \end{tabularx} \\ \\ \relax
6:12 & \begin{tabularx}{0.7\textwidth}{X} 所以，不要讓罪在你們必死的身上掌權，使你們順從身體的私慾。 \end{tabularx} \\ \\ \relax
6:13 & \begin{tabularx}{0.7\textwidth}{X} 也不要把你們的肢體獻給罪作不義的工具，倒要像從死人中活著的人，把自己獻給神，並把你們的肢體獻給神作義的工具。 \end{tabularx} \\ \\ \relax
6:14 & \begin{tabularx}{0.7\textwidth}{X} 罪必不能作你們的主，因你們不在律法之下，而是在恩典之下。 \end{tabularx} \\ \\ \relax
6:15 & \begin{tabularx}{0.7\textwidth}{X} 那又怎麼樣呢？我們在恩典之下，不在律法之下，就可以犯罪嗎？絕對不可！ \end{tabularx} \\ \\ \relax
6:16 & \begin{tabularx}{0.7\textwidth}{X} 難道你們不知道，你們獻自己作奴僕，順從誰就作誰的奴僕嗎？或作罪的奴隸，以至於死；或作順服的奴僕，以至於成義。 \end{tabularx} \\ \\ \relax
6:17 & \begin{tabularx}{0.7\textwidth}{X} 感謝神！因為你們從前雖然作罪的奴隸，現在卻從心裡順服了所傳給你們教導的典範。 \end{tabularx} \\ \\ \relax
6:18 & \begin{tabularx}{0.7\textwidth}{X} 你們既從罪裡得了釋放，就作了義的奴僕。 \end{tabularx} \\ \\ \relax
6:19 & \begin{tabularx}{0.7\textwidth}{X} 我因你們肉體的軟弱，就以人的觀點來說。你們從前怎樣把肢體獻給不潔不法作奴隸，以至於不法；現在也要照樣將肢體獻給義作奴僕，以至於成聖。 \end{tabularx} \\ \\ \relax
6:20 & \begin{tabularx}{0.7\textwidth}{X} 因為你們作罪的奴隸時，不被義所約束。 \end{tabularx} \\ \\ \relax
6:21 & \begin{tabularx}{0.7\textwidth}{X} 那麼，你們現在所看為羞恥的事，當時有甚麼果子呢？那些事的結局就是死。 \end{tabularx} \\ \\ \relax
6:22 & \begin{tabularx}{0.7\textwidth}{X} 但如今，你們既從罪裡得了釋放，作了神的奴僕，就結出果子，以至於成聖，那結局就是永生。 \end{tabularx} \\ \\ \relax
6:23 & \begin{tabularx}{0.7\textwidth}{X} 因為罪的工價乃是死；惟有神的恩賜，在我們的主基督耶穌裡，乃是永生。 \end{tabularx} \\ \\
[1ex]
\hline
\hline
\end{longtable}


\newpage

\section{第67屆~港九培靈會~鮑會園~第六講}
\label{sec:932}
\textbf{實際的經驗}
\newline
\newline
連結: \href{https://www.hkbibleconference.org/session-message/view/932}{\texttt{https://www.hkbibleconference.org/session-message/view/932}}
\newline
\newline
\hyperref[sec:931]{< < < PREV SERMON < < <}
~
\hyperref[sec:toc]{[返目錄]}
~
\hyperref[sec:933]{> > > NEXT SERMON > > >}
\newline
\newline
羅馬書 7:1-7:25
\newline
\begin{longtable}{cl}
\hline
\hline
章節 & 經文 (和合本修訂版)\\
\hline
7:1 & \begin{tabularx}{0.7\textwidth}{X} 弟兄們，我對你們這些明白律法的人說，你們豈不知道律法約束人是在他活著的時候嗎？ \end{tabularx} \\ \\ \relax
7:2 & \begin{tabularx}{0.7\textwidth}{X} 就如女人有了丈夫，丈夫還活著，她就被律法約束；丈夫若死了，她就從丈夫的律法中解脫了。 \end{tabularx} \\ \\ \relax
7:3 & \begin{tabularx}{0.7\textwidth}{X} 所以丈夫還活著，她若跟了別的男人，就叫淫婦；丈夫若死了，她就脫離了律法，雖然跟了別的男人，也不是淫婦。 \end{tabularx} \\ \\ \relax
7:4 & \begin{tabularx}{0.7\textwidth}{X} 我的弟兄們，這樣說來，你們藉著基督的身體對律法也是死了，使你們歸於另一位，就是歸於那從死人中復活的，為要使我們結果子給神。 \end{tabularx} \\ \\ \relax
7:5 & \begin{tabularx}{0.7\textwidth}{X} 因為我們屬肉體的時候，那因律法而生犯罪的慾望在我們肢體中發動，以致結出死亡的果子。 \end{tabularx} \\ \\ \relax
7:6 & \begin{tabularx}{0.7\textwidth}{X} 但如今，我們既然在捆綁我們的律法上死了，就從律法中解脫，使我們服侍主，要按著聖靈的新樣，不按著儀文的舊樣。 \end{tabularx} \\ \\ \relax
7:7 & \begin{tabularx}{0.7\textwidth}{X} 這樣，我們要怎麼說呢？律法是罪嗎？絕對不是！但是，若不是藉著律法，我就不知何為罪；若不是律法說「不可貪心」，我就不知何為貪心。 \end{tabularx} \\ \\ \relax
7:8 & \begin{tabularx}{0.7\textwidth}{X} 然而，罪趁著機會，藉著誡命，使各樣的貪心在我裡頭發動，因為沒有律法，罪是死的。 \end{tabularx} \\ \\ \relax
7:9 & \begin{tabularx}{0.7\textwidth}{X} 以前沒有律法的時候，我是活的；但是誡命來到，罪活起來， \end{tabularx} \\ \\ \relax
7:10 & \begin{tabularx}{0.7\textwidth}{X} 我就死了。那本該叫人活的誡命反而叫我死。 \end{tabularx} \\ \\ \relax
7:11 & \begin{tabularx}{0.7\textwidth}{X} 因為罪趁著機會，藉著誡命誘惑我，並且藉著誡命殺了我。 \end{tabularx} \\ \\ \relax
7:12 & \begin{tabularx}{0.7\textwidth}{X} 這樣看來，律法是聖的，誡命也是聖的、義的、善的。 \end{tabularx} \\ \\ \relax
7:13 & \begin{tabularx}{0.7\textwidth}{X} 那麼，那善的是叫我死嗎？絕對不是！叫我死的是罪。罪藉著那善的叫我死，為要顯出這真是罪，以致罪藉著誡命更顯出是惡極了。 \end{tabularx} \\ \\ \relax
7:14 & \begin{tabularx}{0.7\textwidth}{X} 我們原知道律法是屬靈的，我卻是屬肉體的，是已經賣給罪了。 \end{tabularx} \\ \\ \relax
7:15 & \begin{tabularx}{0.7\textwidth}{X} 因為我所做的，我自己不明白。我所願意的，我並不做；我所恨惡的，我反而去做。 \end{tabularx} \\ \\ \relax
7:16 & \begin{tabularx}{0.7\textwidth}{X} 如果我所做的是我所不願意的，我得承認律法是善的。 \end{tabularx} \\ \\ \relax
7:17 & \begin{tabularx}{0.7\textwidth}{X} 事實上，這不是我做的，而是住在我裡面的罪做的。 \end{tabularx} \\ \\ \relax
7:18 & \begin{tabularx}{0.7\textwidth}{X} 我也知道，住在我裡面的，就是我肉體之中，沒有善。因為立志為善由得我，只是行出來由不得我。 \end{tabularx} \\ \\ \relax
7:19 & \begin{tabularx}{0.7\textwidth}{X} 我所願意的善，我不去做；我所不願意的惡，我反而去做。 \end{tabularx} \\ \\ \relax
7:20 & \begin{tabularx}{0.7\textwidth}{X} 如果我去做我不願意做的，就不是我做的，而是住在我裡面的罪做的。 \end{tabularx} \\ \\ \relax
7:21 & \begin{tabularx}{0.7\textwidth}{X} 我覺得有個律，就是我願意行善的時候，就有惡纏著我。 \end{tabularx} \\ \\ \relax
7:22 & \begin{tabularx}{0.7\textwidth}{X} 因為，按著我裡面的人，我喜歡神的律， \end{tabularx} \\ \\ \relax
7:23 & \begin{tabularx}{0.7\textwidth}{X} 但我看出肢體中另有個律和我內心的律交戰，把我擄去，使我附從那肢體中罪的律。 \end{tabularx} \\ \\ \relax
7:24 & \begin{tabularx}{0.7\textwidth}{X} 我真苦啊！誰能救我脫離這必死的身體呢？ \end{tabularx} \\ \\ \relax
7:25 & \begin{tabularx}{0.7\textwidth}{X} 感謝神，靠著我們的主耶穌基督就能！這樣看來，一方面，我內心順服神的律，另一方面，肉體卻順服罪的律了。 \end{tabularx} \\ \\
[1ex]
\hline
\hline
\end{longtable}


\newpage

\section{第67屆~港九培靈會~鮑會園~第七講}
\label{sec:933}
\textbf{勝利的秘訣}
\newline
\newline
連結: \href{https://www.hkbibleconference.org/session-message/view/933}{\texttt{https://www.hkbibleconference.org/session-message/view/933}}
\newline
\newline
\hyperref[sec:932]{< < < PREV SERMON < < <}
~
\hyperref[sec:toc]{[返目錄]}
~
\hyperref[sec:934]{> > > NEXT SERMON > > >}
\newline
\newline
羅馬書 8:1-8:25
\newline
\begin{longtable}{cl}
\hline
\hline
章節 & 經文 (和合本修訂版)\\
\hline
8:1 & \begin{tabularx}{0.7\textwidth}{X} 如今，那些在基督耶穌裡的人就不被定罪了。 \end{tabularx} \\ \\ \relax
8:2 & \begin{tabularx}{0.7\textwidth}{X} 因為賜生命的聖靈的律，在基督耶穌裡從罪和死的律中把你釋放出來。 \end{tabularx} \\ \\ \relax
8:3 & \begin{tabularx}{0.7\textwidth}{X} 律法既因肉體軟弱而無能為力，神就差遣自己的兒子成為罪身的樣子，為了對付罪，在肉體中定了罪， \end{tabularx} \\ \\ \relax
8:4 & \begin{tabularx}{0.7\textwidth}{X} 為要使律法要求的義，實現在我們這不隨從肉體、只隨從聖靈去行的人身上。 \end{tabularx} \\ \\ \relax
8:5 & \begin{tabularx}{0.7\textwidth}{X} 因為，隨從肉體的人體貼肉體的事；隨從聖靈的人體貼聖靈的事。 \end{tabularx} \\ \\ \relax
8:6 & \begin{tabularx}{0.7\textwidth}{X} 體貼肉體就是死；體貼聖靈就是生命和平安。 \end{tabularx} \\ \\ \relax
8:7 & \begin{tabularx}{0.7\textwidth}{X} 因為體貼肉體就是與神為敵，對神的律法不順服，事實上也無法順服。 \end{tabularx} \\ \\ \relax
8:8 & \begin{tabularx}{0.7\textwidth}{X} 屬肉體的人無法使神喜悅。 \end{tabularx} \\ \\ \relax
8:9 & \begin{tabularx}{0.7\textwidth}{X} 如果神的靈住在你們裡面，你們就不屬肉體，而是屬聖靈了。人若沒有基督的靈，就不是屬基督的。 \end{tabularx} \\ \\ \relax
8:10 & \begin{tabularx}{0.7\textwidth}{X} 基督若在你們裡面，身體就因罪而死，靈卻因義而活。 \end{tabularx} \\ \\ \relax
8:11 & \begin{tabularx}{0.7\textwidth}{X} 然而，使耶穌從死人中復活的神的靈若住在你們裡面，那使基督從死人中復活的，也必藉著住在你們裡面的聖靈使你們必死的身體又活過來。 \end{tabularx} \\ \\ \relax
8:12 & \begin{tabularx}{0.7\textwidth}{X} 弟兄們，這樣看來，我們不是欠肉體的債去順從肉體而活。 \end{tabularx} \\ \\ \relax
8:13 & \begin{tabularx}{0.7\textwidth}{X} 你們若順從肉體活著，必定會死；若靠著聖靈把身體的惡行處死，就必存活。 \end{tabularx} \\ \\ \relax
8:14 & \begin{tabularx}{0.7\textwidth}{X} 因為凡被神的靈引導的都是神的兒子。 \end{tabularx} \\ \\ \relax
8:15 & \begin{tabularx}{0.7\textwidth}{X} 你們所領受的不是奴僕的靈，仍舊害怕；所領受的是兒子名分的靈，因此我們呼叫：「阿爸，父！」 \end{tabularx} \\ \\ \relax
8:16 & \begin{tabularx}{0.7\textwidth}{X} 聖靈自己與我們的靈一同見證我們是神的兒女。 \end{tabularx} \\ \\ \relax
8:17 & \begin{tabularx}{0.7\textwidth}{X} 若是兒女，就是後嗣，是神的後嗣，和基督同作後嗣。如果我們和他一同受苦，是要我們和他一同得榮耀。 \end{tabularx} \\ \\ \relax
8:18 & \begin{tabularx}{0.7\textwidth}{X} 我認為，現在的苦楚，若比起將來要顯示給我們的榮耀，是不足介意的。 \end{tabularx} \\ \\ \relax
8:19 & \begin{tabularx}{0.7\textwidth}{X} 受造之物切望等候神的眾子顯出來。 \end{tabularx} \\ \\ \relax
8:20 & \begin{tabularx}{0.7\textwidth}{X} 因為受造之物屈服在虛空之下，不是自己願意，而是因那使它屈服的叫他如此。 \end{tabularx} \\ \\ \relax
8:21 & \begin{tabularx}{0.7\textwidth}{X} 但受造之物仍然指望從敗壞的轄制下得釋放，得享神兒女榮耀的自由。 \end{tabularx} \\ \\ \relax
8:22 & \begin{tabularx}{0.7\textwidth}{X} 我們知道，一切受造之物一同呻吟，一同忍受陣痛，直到如今。 \end{tabularx} \\ \\ \relax
8:23 & \begin{tabularx}{0.7\textwidth}{X} 不但如此，就是我們這有聖靈作初熟果子的，也是自己內心呻吟，等候得著兒子的名分，就是我們的身體得救贖。 \end{tabularx} \\ \\ \relax
8:24 & \begin{tabularx}{0.7\textwidth}{X} 我們得救是在於盼望；可是看得見的盼望就不是盼望。誰還去盼望他所看得見的呢？ \end{tabularx} \\ \\ \relax
8:25 & \begin{tabularx}{0.7\textwidth}{X} 但我們若盼望那看不見的，我們就耐心等候。 \end{tabularx} \\ \\
[1ex]
\hline
\hline
\end{longtable}


\newpage

\section{第67屆~港九培靈會~鮑會園~第八講}
\label{sec:934}
\textbf{生活的動力}
\newline
\newline
連結: \href{https://www.hkbibleconference.org/session-message/view/934}{\texttt{https://www.hkbibleconference.org/session-message/view/934}}
\newline
\newline
\hyperref[sec:933]{< < < PREV SERMON < < <}
~
\hyperref[sec:toc]{[返目錄]}
~
\hyperref[sec:935]{> > > NEXT SERMON > > >}
\newline
\newline
羅馬書 12:1-12:21
\newline
\begin{longtable}{cl}
\hline
\hline
章節 & 經文 (和合本修訂版)\\
\hline
12:1 & \begin{tabularx}{0.7\textwidth}{X} 所以，弟兄們，我以神的慈悲勸你們，將身體獻上當作活祭，是聖潔的，是神所喜悅的，你們如此事奉乃是理所當然的。 \end{tabularx} \\ \\ \relax
12:2 & \begin{tabularx}{0.7\textwidth}{X} 不要效法這個世界，只要心意更新而變化，叫你們察驗何為神的善良、純全、可喜悅的旨意。 \end{tabularx} \\ \\ \relax
12:3 & \begin{tabularx}{0.7\textwidth}{X} 我憑著所賜我的恩對你們每一位說：不要把自己看得太高，要照著神所分給各人的信心來衡量，看得合乎中道。 \end{tabularx} \\ \\ \relax
12:4 & \begin{tabularx}{0.7\textwidth}{X} 正如我們一個身子上有好些肢體，肢體也不都有一樣的用處。 \end{tabularx} \\ \\ \relax
12:5 & \begin{tabularx}{0.7\textwidth}{X} 這樣，我們許多人在基督裡是一個身體，互相聯絡作肢體。 \end{tabularx} \\ \\ \relax
12:6 & \begin{tabularx}{0.7\textwidth}{X} 按著所得的恩典，我們各有不同的恩賜：或說預言，要按著信心的程度說預言； \end{tabularx} \\ \\ \relax
12:7 & \begin{tabularx}{0.7\textwidth}{X} 或服事的，要專一服事；或教導的，要專一教導； \end{tabularx} \\ \\ \relax
12:8 & \begin{tabularx}{0.7\textwidth}{X} 或勸勉的，要專一勸勉；施捨的，要誠實；治理的，要殷勤；憐憫人的，要樂意。 \end{tabularx} \\ \\ \relax
12:9 & \begin{tabularx}{0.7\textwidth}{X} 愛，不可虛假；惡，要厭惡；善，要親近。 \end{tabularx} \\ \\ \relax
12:10 & \begin{tabularx}{0.7\textwidth}{X} 愛弟兄，要相親相愛；恭敬人，要彼此推讓； \end{tabularx} \\ \\ \relax
12:11 & \begin{tabularx}{0.7\textwidth}{X} 殷勤，不可懶惰。要靈裡火熱；常常服侍主。 \end{tabularx} \\ \\ \relax
12:12 & \begin{tabularx}{0.7\textwidth}{X} 在盼望中要喜樂；在患難中要忍耐；禱告要恆切。 \end{tabularx} \\ \\ \relax
12:13 & \begin{tabularx}{0.7\textwidth}{X} 聖徒有缺乏，要供給；異鄉客，要殷勤款待。 \end{tabularx} \\ \\ \relax
12:14 & \begin{tabularx}{0.7\textwidth}{X} 要祝福迫害你們的，要祝福，不可詛咒。 \end{tabularx} \\ \\ \relax
12:15 & \begin{tabularx}{0.7\textwidth}{X} 要與喜樂的人同樂；要與哀哭的人同哭。 \end{tabularx} \\ \\ \relax
12:16 & \begin{tabularx}{0.7\textwidth}{X} 要彼此同心，不要心高氣傲，倒要俯就卑微的人。不要自以為聰明。 \end{tabularx} \\ \\ \relax
12:17 & \begin{tabularx}{0.7\textwidth}{X} 不要以惡報惡，眾人以為美的事要留心去做。 \end{tabularx} \\ \\ \relax
12:18 & \begin{tabularx}{0.7\textwidth}{X} 若是可行，總要盡力與眾人和睦。 \end{tabularx} \\ \\ \relax
12:19 & \begin{tabularx}{0.7\textwidth}{X} 各位親愛的，不要自己伸冤，寧可給主的憤怒留地步，因為經上記著：「主說：『伸冤在我，我必報應。』」 \end{tabularx} \\ \\ \relax
12:20 & \begin{tabularx}{0.7\textwidth}{X} 不但如此，「你的仇敵若餓了，就給他吃；若渴了，就給他喝。因為你這樣做，就是把炭火堆在他的頭上。」 \end{tabularx} \\ \\ \relax
12:21 & \begin{tabularx}{0.7\textwidth}{X} 不要被惡所勝，反要以善勝惡。 \end{tabularx} \\ \\
[1ex]
\hline
\hline
\end{longtable}


\newpage

\section{第67屆~港九培靈會~鮑會園~第九講}
\label{sec:935}
\textbf{教會的使命}
\newline
\newline
連結: \href{https://www.hkbibleconference.org/session-message/view/935}{\texttt{https://www.hkbibleconference.org/session-message/view/935}}
\newline
\newline
\hyperref[sec:934]{< < < PREV SERMON < < <}
~
\hyperref[sec:toc]{[返目錄]}
~
\hyperref[sec:911]{> > > NEXT SERMON > > >}
\newline
\newline
羅馬書 15:20-15:33
\newline
\begin{longtable}{cl}
\hline
\hline
章節 & 經文 (和合本修訂版)\\
\hline
15:20 & \begin{tabularx}{0.7\textwidth}{X} 這樣，我立了志向，不在基督的名已經傳揚過的地方傳福音，免得建造在別人的根基上； \end{tabularx} \\ \\ \relax
15:21 & \begin{tabularx}{0.7\textwidth}{X} 卻如經上所記：「未曾傳給他們的，他們必看見；未曾聽見過的事，他們要明白。」 \end{tabularx} \\ \\ \relax
15:22 & \begin{tabularx}{0.7\textwidth}{X} 因此我多次被攔阻，不能到你們那裡去。 \end{tabularx} \\ \\ \relax
15:23 & \begin{tabularx}{0.7\textwidth}{X} 但如今，在這一帶再沒有可傳的地方，而且這許多年來，我迫切想去你們那裡， \end{tabularx} \\ \\ \relax
15:24 & \begin{tabularx}{0.7\textwidth}{X} 盼望到西班牙去的時候經過，得見你們，先與你們彼此交往，心裡稍得滿足，然後蒙你們為我送行。 \end{tabularx} \\ \\ \relax
15:25 & \begin{tabularx}{0.7\textwidth}{X} 但如今我要到耶路撒冷去，供應聖徒的需要。 \end{tabularx} \\ \\ \relax
15:26 & \begin{tabularx}{0.7\textwidth}{X} 因為馬其頓和亞該亞人樂意湊出一些捐款給耶路撒冷聖徒中的窮人。 \end{tabularx} \\ \\ \relax
15:27 & \begin{tabularx}{0.7\textwidth}{X} 這固然是他們樂意的，其實也算是所欠的債；因為外邦人既然分享了他們靈性上的好處，就當把肉體上的需用供給他們。 \end{tabularx} \\ \\ \relax
15:28 & \begin{tabularx}{0.7\textwidth}{X} 等我辦完了這事，把這筆捐款交付給他們，我就要路過你們那裡，到西班牙去。 \end{tabularx} \\ \\ \relax
15:29 & \begin{tabularx}{0.7\textwidth}{X} 我也知道去你們那裡的時候，我將帶著基督豐盛的恩典去。 \end{tabularx} \\ \\ \relax
15:30 & \begin{tabularx}{0.7\textwidth}{X} 弟兄們，我藉著我們的主耶穌基督，又藉著聖靈的愛，勸你們與我一同竭力為我祈求神， \end{tabularx} \\ \\ \relax
15:31 & \begin{tabularx}{0.7\textwidth}{X} 使我脫離在猶太不順從的人，也讓我在耶路撒冷的事奉可蒙聖徒悅納， \end{tabularx} \\ \\ \relax
15:32 & \begin{tabularx}{0.7\textwidth}{X} 並使我照著神的旨意歡歡喜喜地到你們那裡，與你們同得安息。 \end{tabularx} \\ \\ \relax
15:33 & \begin{tabularx}{0.7\textwidth}{X} 願賜平安的神與你們眾人同在。阿們！ \end{tabularx} \\ \\
[1ex]
\hline
\hline
\end{longtable}


\newpage

\section{第68屆~港九培靈會~吳勇~第一講}
\label{sec:911}
\textbf{不把神當神}
\newline
\newline
連結: \href{https://www.hkbibleconference.org/session-message/view/911}{\texttt{https://www.hkbibleconference.org/session-message/view/911}}
\newline
\newline
\hyperref[sec:935]{< < < PREV SERMON < < <}
~
\hyperref[sec:toc]{[返目錄]}
~
\hyperref[sec:912]{> > > NEXT SERMON > > >}
\newline
\newline
約翰福音 17:3-17:3
\newline
\begin{longtable}{cl}
\hline
\hline
章節 & 經文 (和合本修訂版)\\
\hline
17:3 & \begin{tabularx}{0.7\textwidth}{X} 認識你—獨一的真神，並且認識你所差來的耶穌基督，這就是永生。 \end{tabularx} \\ \\
[1ex]
\hline
\hline
\end{longtable}


\newpage

\section{第68屆~港九培靈會~吳勇~第二講}
\label{sec:912}
\textbf{離開與回來}
\newline
\newline
連結: \href{https://www.hkbibleconference.org/session-message/view/912}{\texttt{https://www.hkbibleconference.org/session-message/view/912}}
\newline
\newline
\hyperref[sec:911]{< < < PREV SERMON < < <}
~
\hyperref[sec:toc]{[返目錄]}
~
\hyperref[sec:913]{> > > NEXT SERMON > > >}
\newline
\newline


\newpage

\section{第68屆~港九培靈會~吳勇~第三講}
\label{sec:913}
\textbf{三個願望}
\newline
\newline
連結: \href{https://www.hkbibleconference.org/session-message/view/913}{\texttt{https://www.hkbibleconference.org/session-message/view/913}}
\newline
\newline
\hyperref[sec:912]{< < < PREV SERMON < < <}
~
\hyperref[sec:toc]{[返目錄]}
~
\hyperref[sec:914]{> > > NEXT SERMON > > >}
\newline
\newline
出埃及記 40:34-40:34
\newline
\begin{longtable}{cl}
\hline
\hline
章節 & 經文 (和合本修訂版)\\
\hline
40:34 & \begin{tabularx}{0.7\textwidth}{X} 那時，雲彩遮蓋會幕，耶和華的榮光充滿了帳幕。 \end{tabularx} \\ \\
[1ex]
\hline
\hline
\end{longtable}


\newpage

\section{第68屆~港九培靈會~吳勇~第四講}
\label{sec:914}
\textbf{屬肉體與屬靈}
\newline
\newline
連結: \href{https://www.hkbibleconference.org/session-message/view/914}{\texttt{https://www.hkbibleconference.org/session-message/view/914}}
\newline
\newline
\hyperref[sec:913]{< < < PREV SERMON < < <}
~
\hyperref[sec:toc]{[返目錄]}
~
\hyperref[sec:915]{> > > NEXT SERMON > > >}
\newline
\newline
創世記 25:22-25:22
\newline
\begin{longtable}{cl}
\hline
\hline
章節 & 經文 (和合本修訂版)\\
\hline
25:22 & \begin{tabularx}{0.7\textwidth}{X} 胎兒們在她腹中彼此相爭，她就說：「若是如此，我為甚麼會這樣呢？」她就去求問耶和華。 \end{tabularx} \\ \\
[1ex]
\hline
\hline
\end{longtable}


\newpage

\section{第68屆~港九培靈會~吳勇~第五講}
\label{sec:915}
\textbf{生命的守護者}
\newline
\newline
連結: \href{https://www.hkbibleconference.org/session-message/view/915}{\texttt{https://www.hkbibleconference.org/session-message/view/915}}
\newline
\newline
\hyperref[sec:914]{< < < PREV SERMON < < <}
~
\hyperref[sec:toc]{[返目錄]}
~
\hyperref[sec:916]{> > > NEXT SERMON > > >}
\newline
\newline
民數記 17:8-17:8
\newline
\begin{longtable}{cl}
\hline
\hline
章節 & 經文 (和合本修訂版)\\
\hline
17:8 & \begin{tabularx}{0.7\textwidth}{X} 第二天，摩西進到法櫃的帳幕去，看哪，利未族亞倫的杖已經發芽，長了花苞，開了花，也結出熟的杏子！ \end{tabularx} \\ \\
[1ex]
\hline
\hline
\end{longtable}


\newpage

\section{第68屆~港九培靈會~吳勇~第六講}
\label{sec:916}
\textbf{火與水的異象}
\newline
\newline
連結: \href{https://www.hkbibleconference.org/session-message/view/916}{\texttt{https://www.hkbibleconference.org/session-message/view/916}}
\newline
\newline
\hyperref[sec:915]{< < < PREV SERMON < < <}
~
\hyperref[sec:toc]{[返目錄]}
~
\hyperref[sec:917]{> > > NEXT SERMON > > >}
\newline
\newline
以西結書 47:5-47:5
\newline
\begin{longtable}{cl}
\hline
\hline
章節 & 經文 (和合本修訂版)\\
\hline
47:5 & \begin{tabularx}{0.7\textwidth}{X} 又量了一千，水已成河，無法過去；因為水勢高漲成河，只能游泳，無法走過。 \end{tabularx} \\ \\
[1ex]
\hline
\hline
\end{longtable}


\newpage

\section{第68屆~港九培靈會~吳勇~第七講}
\label{sec:917}
\textbf{兩個安息}
\newline
\newline
連結: \href{https://www.hkbibleconference.org/session-message/view/917}{\texttt{https://www.hkbibleconference.org/session-message/view/917}}
\newline
\newline
\hyperref[sec:916]{< < < PREV SERMON < < <}
~
\hyperref[sec:toc]{[返目錄]}
~
\hyperref[sec:918]{> > > NEXT SERMON > > >}
\newline
\newline
腓立比書 2:6-2:8
\newline
\begin{longtable}{cl}
\hline
\hline
章節 & 經文 (和合本修訂版)\\
\hline
2:6 & \begin{tabularx}{0.7\textwidth}{X} 他本有神的形像，卻不堅持自己與神同等； \end{tabularx} \\ \\ \relax
2:7 & \begin{tabularx}{0.7\textwidth}{X} 反倒虛己，取了奴僕的形像，成為人的樣式；既有人的樣子， \end{tabularx} \\ \\ \relax
2:8 & \begin{tabularx}{0.7\textwidth}{X} 就謙卑自己，存心順服，以至於死，且死在十字架上。 \end{tabularx} \\ \\
[1ex]
\hline
\hline
\end{longtable}


\newpage

\section{第68屆~港九培靈會~吳勇~第八講}
\label{sec:918}
\textbf{約翰的補網}
\newline
\newline
連結: \href{https://www.hkbibleconference.org/session-message/view/918}{\texttt{https://www.hkbibleconference.org/session-message/view/918}}
\newline
\newline
\hyperref[sec:917]{< < < PREV SERMON < < <}
~
\hyperref[sec:toc]{[返目錄]}
~
\hyperref[sec:919]{> > > NEXT SERMON > > >}
\newline
\newline


\newpage

\section{第68屆~港九培靈會~吳勇~第九講}
\label{sec:919}
\textbf{復興的器皿}
\newline
\newline
連結: \href{https://www.hkbibleconference.org/session-message/view/919}{\texttt{https://www.hkbibleconference.org/session-message/view/919}}
\newline
\newline
\hyperref[sec:918]{< < < PREV SERMON < < <}
~
\hyperref[sec:toc]{[返目錄]}
~
\hyperref[sec:920]{> > > NEXT SERMON > > >}
\newline
\newline


\newpage

\section{第68屆~港九培靈會~吳勇~第十講}
\label{sec:920}
\textbf{人與路}
\newline
\newline
連結: \href{https://www.hkbibleconference.org/session-message/view/920}{\texttt{https://www.hkbibleconference.org/session-message/view/920}}
\newline
\newline
\hyperref[sec:919]{< < < PREV SERMON < < <}
~
\hyperref[sec:toc]{[返目錄]}
~
\hyperref[sec:902]{> > > NEXT SERMON > > >}
\newline
\newline
使徒行傳 13:22-13:22
\newline
\begin{longtable}{cl}
\hline
\hline
章節 & 經文 (和合本修訂版)\\
\hline
13:22 & \begin{tabularx}{0.7\textwidth}{X} 他廢了掃羅之後，就興起大衛作他們的王，又為他作見證說：『我尋得耶西的兒子大衛，他是合我心意的人，他要遵行我一切的旨意。』 \end{tabularx} \\ \\
[1ex]
\hline
\hline
\end{longtable}


\newpage

\section{第68屆~港九培靈會~歐達~第一講}
\label{sec:902}
\textbf{先知的險境}
\newline
\newline
連結: \href{https://www.hkbibleconference.org/session-message/view/902}{\texttt{https://www.hkbibleconference.org/session-message/view/902}}
\newline
\newline
\hyperref[sec:920]{< < < PREV SERMON < < <}
~
\hyperref[sec:toc]{[返目錄]}
~
\hyperref[sec:903]{> > > NEXT SERMON > > >}
\newline
\newline
約翰福音 4:43-4:45
\newline
\begin{longtable}{cl}
\hline
\hline
章節 & 經文 (和合本修訂版)\\
\hline
4:43 & \begin{tabularx}{0.7\textwidth}{X} 過了那兩天，耶穌離開那地方，往加利利去。 \end{tabularx} \\ \\ \relax
4:44 & \begin{tabularx}{0.7\textwidth}{X} 因為耶穌自己作過見證說：「先知在自己的家鄉是沒有人尊敬的。」 \end{tabularx} \\ \\ \relax
4:45 & \begin{tabularx}{0.7\textwidth}{X} 到了加利利，加利利人都歡迎他，因為他們也上耶路撒冷去過節，曾經看過他在節期間所做的一切事。 \end{tabularx} \\ \\
[1ex]
\hline
\hline
\end{longtable}


\newpage

\section{第68屆~港九培靈會~歐達~第二講}
\label{sec:903}
\textbf{與主默契}
\newline
\newline
連結: \href{https://www.hkbibleconference.org/session-message/view/903}{\texttt{https://www.hkbibleconference.org/session-message/view/903}}
\newline
\newline
\hyperref[sec:902]{< < < PREV SERMON < < <}
~
\hyperref[sec:toc]{[返目錄]}
~
\hyperref[sec:904]{> > > NEXT SERMON > > >}
\newline
\newline
詩篇 24:1-24:3
\newline
\begin{longtable}{cl}
\hline
\hline
章節 & 經文 (和合本修訂版)\\
\hline
24:1 & \begin{tabularx}{0.7\textwidth}{X} 地和其中所充滿的，世界和住在其中的，都屬耶和華。 \end{tabularx} \\ \\ \relax
24:2 & \begin{tabularx}{0.7\textwidth}{X} 他把地建立在海上，安定在江河之上。 \end{tabularx} \\ \\ \relax
24:3 & \begin{tabularx}{0.7\textwidth}{X} 誰能登耶和華的山？誰能站在他的聖所？ \end{tabularx} \\ \\
[1ex]
\hline
\hline
\end{longtable}


\newpage

\section{第68屆~港九培靈會~歐達~第三講}
\label{sec:904}
\textbf{正直忠誠}
\newline
\newline
連結: \href{https://www.hkbibleconference.org/session-message/view/904}{\texttt{https://www.hkbibleconference.org/session-message/view/904}}
\newline
\newline
\hyperref[sec:903]{< < < PREV SERMON < < <}
~
\hyperref[sec:toc]{[返目錄]}
~
\hyperref[sec:905]{> > > NEXT SERMON > > >}
\newline
\newline
詩篇 15:1-15:5
\newline
\begin{longtable}{cl}
\hline
\hline
章節 & 經文 (和合本修訂版)\\
\hline
15:1 & \begin{tabularx}{0.7\textwidth}{X} 耶和華啊，誰能寄居你的帳幕？誰能居住你的聖山？ \end{tabularx} \\ \\ \relax
15:2 & \begin{tabularx}{0.7\textwidth}{X} 就是行為正直、做事公義、心裡說實話的人。 \end{tabularx} \\ \\ \relax
15:3 & \begin{tabularx}{0.7\textwidth}{X} 他不以舌頭讒害人，不惡待朋友，也不隨夥毀謗鄰舍。 \end{tabularx} \\ \\ \relax
15:4 & \begin{tabularx}{0.7\textwidth}{X} 他眼中藐視匪類，卻尊重那敬畏耶和華的人。他發了誓，雖然自己吃虧也不更改。 \end{tabularx} \\ \\ \relax
15:5 & \begin{tabularx}{0.7\textwidth}{X} 他不放債取利，不受賄賂以害無辜。做這些事的人必永不動搖。 \end{tabularx} \\ \\
[1ex]
\hline
\hline
\end{longtable}


\newpage

\section{第68屆~港九培靈會~歐達~第四講}
\label{sec:905}
\textbf{融和合一}
\newline
\newline
連結: \href{https://www.hkbibleconference.org/session-message/view/905}{\texttt{https://www.hkbibleconference.org/session-message/view/905}}
\newline
\newline
\hyperref[sec:904]{< < < PREV SERMON < < <}
~
\hyperref[sec:toc]{[返目錄]}
~
\hyperref[sec:906]{> > > NEXT SERMON > > >}
\newline
\newline
約翰福音 4:19-4:24
\newline
\begin{longtable}{cl}
\hline
\hline
章節 & 經文 (和合本修訂版)\\
\hline
4:19 & \begin{tabularx}{0.7\textwidth}{X} 婦人對他說：「先生，我看你是一位先知。 \end{tabularx} \\ \\ \relax
4:20 & \begin{tabularx}{0.7\textwidth}{X} 我們的祖宗在這山上敬拜神，你們倒說，應當敬拜的地方是在耶路撒冷。」 \end{tabularx} \\ \\ \relax
4:21 & \begin{tabularx}{0.7\textwidth}{X} 耶穌對她說：「婦人，你要信我。時候將到，你們敬拜父，既不在這山上，也不在耶路撒冷。 \end{tabularx} \\ \\ \relax
4:22 & \begin{tabularx}{0.7\textwidth}{X} 你們所敬拜的，你們不知道；我們所敬拜的，我們知道，因為救恩是從猶太人出來的。 \end{tabularx} \\ \\ \relax
4:23 & \begin{tabularx}{0.7\textwidth}{X} 時候將到，現在就是了，那真正敬拜父的，要用心靈和誠實敬拜他，因為父要這樣的人敬拜他。 \end{tabularx} \\ \\ \relax
4:24 & \begin{tabularx}{0.7\textwidth}{X} 神是靈，所以敬拜他的必須用心靈和誠實敬拜他。」 \end{tabularx} \\ \\
[1ex]
\hline
\hline
\end{longtable}


\newpage

\section{第68屆~港九培靈會~歐達~第五講}
\label{sec:906}
\textbf{竭誠事主}
\newline
\newline
連結: \href{https://www.hkbibleconference.org/session-message/view/906}{\texttt{https://www.hkbibleconference.org/session-message/view/906}}
\newline
\newline
\hyperref[sec:905]{< < < PREV SERMON < < <}
~
\hyperref[sec:toc]{[返目錄]}
~
\hyperref[sec:907]{> > > NEXT SERMON > > >}
\newline
\newline
哥林多前書 12:1-12:7
\newline
\begin{longtable}{cl}
\hline
\hline
章節 & 經文 (和合本修訂版)\\
\hline
12:1 & \begin{tabularx}{0.7\textwidth}{X} 弟兄們，關於屬靈的恩賜，我不願意你們不明白。 \end{tabularx} \\ \\ \relax
12:2 & \begin{tabularx}{0.7\textwidth}{X} 你們知道，你們作外邦人的時候，隨事被引誘，受了迷惑去拜不會出聲的偶像。 \end{tabularx} \\ \\ \relax
12:3 & \begin{tabularx}{0.7\textwidth}{X} 所以，我要你們知道，被神的靈感動的，沒有人會說「耶穌該受詛咒」；若不是被聖靈感動的，也沒有人能說「耶穌是主」。 \end{tabularx} \\ \\ \relax
12:4 & \begin{tabularx}{0.7\textwidth}{X} 恩賜有許多種，卻是同一位聖靈所賜。 \end{tabularx} \\ \\ \relax
12:5 & \begin{tabularx}{0.7\textwidth}{X} 事奉有許多種，卻是事奉同一位主。 \end{tabularx} \\ \\ \relax
12:6 & \begin{tabularx}{0.7\textwidth}{X} 工作有許多種，卻是同一位神在萬人中運行萬事。 \end{tabularx} \\ \\ \relax
12:7 & \begin{tabularx}{0.7\textwidth}{X} 聖靈彰顯在各人身上，是要使人得益處。 \end{tabularx} \\ \\
[1ex]
\hline
\hline
\end{longtable}


\newpage

\section{第68屆~港九培靈會~歐達~第六講}
\label{sec:907}
\textbf{潔淨純全}
\newline
\newline
連結: \href{https://www.hkbibleconference.org/session-message/view/907}{\texttt{https://www.hkbibleconference.org/session-message/view/907}}
\newline
\newline
\hyperref[sec:906]{< < < PREV SERMON < < <}
~
\hyperref[sec:toc]{[返目錄]}
~
\hyperref[sec:908]{> > > NEXT SERMON > > >}
\newline
\newline
詩篇 119:9-119:9
\newline
\begin{longtable}{cl}
\hline
\hline
章節 & 經文 (和合本修訂版)\\
\hline
119:9 & \begin{tabularx}{0.7\textwidth}{X} 青年要如何保持純潔呢？是要遵行你的話！ \end{tabularx} \\ \\
[1ex]
\hline
\hline
\end{longtable}


\newpage

\section{第68屆~港九培靈會~歐達~第七講}
\label{sec:908}
\textbf{門徒樣式}
\newline
\newline
連結: \href{https://www.hkbibleconference.org/session-message/view/908}{\texttt{https://www.hkbibleconference.org/session-message/view/908}}
\newline
\newline
\hyperref[sec:907]{< < < PREV SERMON < < <}
~
\hyperref[sec:toc]{[返目錄]}
~
\hyperref[sec:909]{> > > NEXT SERMON > > >}
\newline
\newline
帖撒羅尼迦前書 1:6-1:10
\newline
\begin{longtable}{cl}
\hline
\hline
章節 & 經文 (和合本修訂版)\\
\hline
1:6 & \begin{tabularx}{0.7\textwidth}{X} 你們成為效法我們，更效法主的人，因聖靈所激發的喜樂，在大患難中領受了真道， \end{tabularx} \\ \\ \relax
1:7 & \begin{tabularx}{0.7\textwidth}{X} 從此你們作了馬其頓和亞該亞所有信主的人的榜樣。 \end{tabularx} \\ \\ \relax
1:8 & \begin{tabularx}{0.7\textwidth}{X} 因為主的道已經從你們那裡傳播出去，你們向神的信心不只在馬其頓和亞該亞，就是在各處也都傳開了，所以不用我們說甚麼話。 \end{tabularx} \\ \\ \relax
1:9 & \begin{tabularx}{0.7\textwidth}{X} 因為他們自己已經傳講我們是怎樣進到你們那裡，你們是怎樣離棄偶像，歸向神來服侍那又真又活的神， \end{tabularx} \\ \\ \relax
1:10 & \begin{tabularx}{0.7\textwidth}{X} 等候他兒子從天降臨，就是神使他從死人中復活的那位救我們脫離將來憤怒的耶穌。 \end{tabularx} \\ \\
[1ex]
\hline
\hline
\end{longtable}


\newpage

\section{第68屆~港九培靈會~歐達~第八講}
\label{sec:909}
\textbf{福音使命}
\newline
\newline
連結: \href{https://www.hkbibleconference.org/session-message/view/909}{\texttt{https://www.hkbibleconference.org/session-message/view/909}}
\newline
\newline
\hyperref[sec:908]{< < < PREV SERMON < < <}
~
\hyperref[sec:toc]{[返目錄]}
~
\hyperref[sec:910]{> > > NEXT SERMON > > >}
\newline
\newline
哈巴谷書 1:5-1:5
\newline
\begin{longtable}{cl}
\hline
\hline
章節 & 經文 (和合本修訂版)\\
\hline
1:5 & \begin{tabularx}{0.7\textwidth}{X} 你們要向列國觀看，注意看，要驚奇，再驚奇！因為在你們的日子，有一件事發生，儘管有人說了，你們還是不信。 \end{tabularx} \\ \\
[1ex]
\hline
\hline
\end{longtable}


\newpage

\section{第68屆~港九培靈會~歐達~第九講}
\label{sec:910}
\textbf{牧者的禱告}
\newline
\newline
連結: \href{https://www.hkbibleconference.org/session-message/view/910}{\texttt{https://www.hkbibleconference.org/session-message/view/910}}
\newline
\newline
\hyperref[sec:909]{< < < PREV SERMON < < <}
~
\hyperref[sec:toc]{[返目錄]}
~
\hyperref[sec:893]{> > > NEXT SERMON > > >}
\newline
\newline


\newpage

\section{第68屆~港九培靈會~蕭壽華~第一講}
\label{sec:893}
\textbf{跟從耶穌到底是甚麼意思?}
\newline
\newline
連結: \href{https://www.hkbibleconference.org/session-message/view/893}{\texttt{https://www.hkbibleconference.org/session-message/view/893}}
\newline
\newline
\hyperref[sec:910]{< < < PREV SERMON < < <}
~
\hyperref[sec:toc]{[返目錄]}
~
\hyperref[sec:894]{> > > NEXT SERMON > > >}
\newline
\newline


\newpage

\section{第68屆~港九培靈會~蕭壽華~第二講}
\label{sec:894}
\textbf{天國子民的品格與見證}
\newline
\newline
連結: \href{https://www.hkbibleconference.org/session-message/view/894}{\texttt{https://www.hkbibleconference.org/session-message/view/894}}
\newline
\newline
\hyperref[sec:893]{< < < PREV SERMON < < <}
~
\hyperref[sec:toc]{[返目錄]}
~
\hyperref[sec:895]{> > > NEXT SERMON > > >}
\newline
\newline


\newpage

\section{第68屆~港九培靈會~蕭壽華~第三講}
\label{sec:895}
\textbf{超越文士的義(一)}
\newline
\newline
連結: \href{https://www.hkbibleconference.org/session-message/view/895}{\texttt{https://www.hkbibleconference.org/session-message/view/895}}
\newline
\newline
\hyperref[sec:894]{< < < PREV SERMON < < <}
~
\hyperref[sec:toc]{[返目錄]}
~
\hyperref[sec:896]{> > > NEXT SERMON > > >}
\newline
\newline


\newpage

\section{第68屆~港九培靈會~蕭壽華~第四講}
\label{sec:896}
\textbf{超越文士的義 (二)}
\newline
\newline
連結: \href{https://www.hkbibleconference.org/session-message/view/896}{\texttt{https://www.hkbibleconference.org/session-message/view/896}}
\newline
\newline
\hyperref[sec:895]{< < < PREV SERMON < < <}
~
\hyperref[sec:toc]{[返目錄]}
~
\hyperref[sec:897]{> > > NEXT SERMON > > >}
\newline
\newline


\newpage

\section{第68屆~港九培靈會~蕭壽華~第五講}
\label{sec:897}
\textbf{天國子民隱藏的操練}
\newline
\newline
連結: \href{https://www.hkbibleconference.org/session-message/view/897}{\texttt{https://www.hkbibleconference.org/session-message/view/897}}
\newline
\newline
\hyperref[sec:896]{< < < PREV SERMON < < <}
~
\hyperref[sec:toc]{[返目錄]}
~
\hyperref[sec:898]{> > > NEXT SERMON > > >}
\newline
\newline


\newpage

\section{第68屆~港九培靈會~蕭壽華~第六講}
\label{sec:898}
\textbf{天國子民的人生目標}
\newline
\newline
連結: \href{https://www.hkbibleconference.org/session-message/view/898}{\texttt{https://www.hkbibleconference.org/session-message/view/898}}
\newline
\newline
\hyperref[sec:897]{< < < PREV SERMON < < <}
~
\hyperref[sec:toc]{[返目錄]}
~
\hyperref[sec:899]{> > > NEXT SERMON > > >}
\newline
\newline


\newpage

\section{第68屆~港九培靈會~蕭壽華~第七講}
\label{sec:899}
\textbf{天國子民的「眼光」}
\newline
\newline
連結: \href{https://www.hkbibleconference.org/session-message/view/899}{\texttt{https://www.hkbibleconference.org/session-message/view/899}}
\newline
\newline
\hyperref[sec:898]{< < < PREV SERMON < < <}
~
\hyperref[sec:toc]{[返目錄]}
~
\hyperref[sec:900]{> > > NEXT SERMON > > >}
\newline
\newline


\newpage

\section{第68屆~港九培靈會~蕭壽華~第八講}
\label{sec:900}
\textbf{天國子民的關係}
\newline
\newline
連結: \href{https://www.hkbibleconference.org/session-message/view/900}{\texttt{https://www.hkbibleconference.org/session-message/view/900}}
\newline
\newline
\hyperref[sec:899]{< < < PREV SERMON < < <}
~
\hyperref[sec:toc]{[返目錄]}
~
\hyperref[sec:901]{> > > NEXT SERMON > > >}
\newline
\newline


\newpage

\section{第68屆~港九培靈會~蕭壽華~第九講}
\label{sec:901}
\textbf{天國子民的委身}
\newline
\newline
連結: \href{https://www.hkbibleconference.org/session-message/view/901}{\texttt{https://www.hkbibleconference.org/session-message/view/901}}
\newline
\newline
\hyperref[sec:900]{< < < PREV SERMON < < <}
~
\hyperref[sec:toc]{[返目錄]}
~
\hyperref[sec:874]{> > > NEXT SERMON > > >}
\newline
\newline


\newpage

\section{第69屆~港九培靈會~周永健~第一講}
\label{sec:874}
\textbf{確認主權 經歷信實}
\newline
\newline
連結: \href{https://www.hkbibleconference.org/session-message/view/874}{\texttt{https://www.hkbibleconference.org/session-message/view/874}}
\newline
\newline
\hyperref[sec:901]{< < < PREV SERMON < < <}
~
\hyperref[sec:toc]{[返目錄]}
~
\hyperref[sec:875]{> > > NEXT SERMON > > >}
\newline
\newline


\newpage

\section{第69屆~港九培靈會~周永健~第二講}
\label{sec:875}
\textbf{持定身份 聖潔自守}
\newline
\newline
連結: \href{https://www.hkbibleconference.org/session-message/view/875}{\texttt{https://www.hkbibleconference.org/session-message/view/875}}
\newline
\newline
\hyperref[sec:874]{< < < PREV SERMON < < <}
~
\hyperref[sec:toc]{[返目錄]}
~
\hyperref[sec:876]{> > > NEXT SERMON > > >}
\newline
\newline


\newpage

\section{第69屆~港九培靈會~周永健~第三講}
\label{sec:876}
\textbf{恆常禱告 認罪祈求}
\newline
\newline
連結: \href{https://www.hkbibleconference.org/session-message/view/876}{\texttt{https://www.hkbibleconference.org/session-message/view/876}}
\newline
\newline
\hyperref[sec:875]{< < < PREV SERMON < < <}
~
\hyperref[sec:toc]{[返目錄]}
~
\hyperref[sec:877]{> > > NEXT SERMON > > >}
\newline
\newline


\newpage

\section{第69屆~港九培靈會~周永健~第四講}
\label{sec:877}
\textbf{不拜別神 專心靠主}
\newline
\newline
連結: \href{https://www.hkbibleconference.org/session-message/view/877}{\texttt{https://www.hkbibleconference.org/session-message/view/877}}
\newline
\newline
\hyperref[sec:876]{< < < PREV SERMON < < <}
~
\hyperref[sec:toc]{[返目錄]}
~
\hyperref[sec:878]{> > > NEXT SERMON > > >}
\newline
\newline


\newpage

\section{第69屆~港九培靈會~周永健~第五講}
\label{sec:878}
\textbf{謙卑自省 莫虧主恩}
\newline
\newline
連結: \href{https://www.hkbibleconference.org/session-message/view/878}{\texttt{https://www.hkbibleconference.org/session-message/view/878}}
\newline
\newline
\hyperref[sec:877]{< < < PREV SERMON < < <}
~
\hyperref[sec:toc]{[返目錄]}
~
\hyperref[sec:879]{> > > NEXT SERMON > > >}
\newline
\newline


\newpage

\section{第69屆~港九培靈會~周永健~第六講}
\label{sec:879}
\textbf{回歸國土 同心建造}
\newline
\newline
連結: \href{https://www.hkbibleconference.org/session-message/view/879}{\texttt{https://www.hkbibleconference.org/session-message/view/879}}
\newline
\newline
\hyperref[sec:878]{< < < PREV SERMON < < <}
~
\hyperref[sec:toc]{[返目錄]}
~
\hyperref[sec:880]{> > > NEXT SERMON > > >}
\newline
\newline


\newpage

\section{第69屆~港九培靈會~周永健~第七講}
\label{sec:880}
\textbf{合力互助 共度危難}
\newline
\newline
連結: \href{https://www.hkbibleconference.org/session-message/view/880}{\texttt{https://www.hkbibleconference.org/session-message/view/880}}
\newline
\newline
\hyperref[sec:879]{< < < PREV SERMON < < <}
~
\hyperref[sec:toc]{[返目錄]}
~
\hyperref[sec:881]{> > > NEXT SERMON > > >}
\newline
\newline


\newpage

\section{第69屆~港九培靈會~周永健~第八講}
\label{sec:881}
\textbf{聆聽主道 樂意遵行}
\newline
\newline
連結: \href{https://www.hkbibleconference.org/session-message/view/881}{\texttt{https://www.hkbibleconference.org/session-message/view/881}}
\newline
\newline
\hyperref[sec:880]{< < < PREV SERMON < < <}
~
\hyperref[sec:toc]{[返目錄]}
~
\hyperref[sec:882]{> > > NEXT SERMON > > >}
\newline
\newline


\newpage

\section{第69屆~港九培靈會~陳黔開~第一講}
\label{sec:882}
\textbf{教會 — 你既愛之,又…}
\newline
\newline
連結: \href{https://www.hkbibleconference.org/session-message/view/882}{\texttt{https://www.hkbibleconference.org/session-message/view/882}}
\newline
\newline
\hyperref[sec:881]{< < < PREV SERMON < < <}
~
\hyperref[sec:toc]{[返目錄]}
~
\hyperref[sec:883]{> > > NEXT SERMON > > >}
\newline
\newline


\newpage

\section{第69屆~港九培靈會~陳黔開~第二講}
\label{sec:883}
\textbf{敬拜 — 參與或投入}
\newline
\newline
連結: \href{https://www.hkbibleconference.org/session-message/view/883}{\texttt{https://www.hkbibleconference.org/session-message/view/883}}
\newline
\newline
\hyperref[sec:882]{< < < PREV SERMON < < <}
~
\hyperref[sec:toc]{[返目錄]}
~
\hyperref[sec:884]{> > > NEXT SERMON > > >}
\newline
\newline


\newpage

\section{第69屆~港九培靈會~陳黔開~第三講}
\label{sec:884}
\textbf{佈道 — 使萬民作門徒}
\newline
\newline
連結: \href{https://www.hkbibleconference.org/session-message/view/884}{\texttt{https://www.hkbibleconference.org/session-message/view/884}}
\newline
\newline
\hyperref[sec:883]{< < < PREV SERMON < < <}
~
\hyperref[sec:toc]{[返目錄]}
~
\hyperref[sec:885]{> > > NEXT SERMON > > >}
\newline
\newline


\newpage

\section{第69屆~港九培靈會~陳黔開~第四講}
\label{sec:885}
\textbf{教導 — 學無止境}
\newline
\newline
連結: \href{https://www.hkbibleconference.org/session-message/view/885}{\texttt{https://www.hkbibleconference.org/session-message/view/885}}
\newline
\newline
\hyperref[sec:884]{< < < PREV SERMON < < <}
~
\hyperref[sec:toc]{[返目錄]}
~
\hyperref[sec:886]{> > > NEXT SERMON > > >}
\newline
\newline


\newpage

\section{第69屆~港九培靈會~陳黔開~第五講}
\label{sec:886}
\textbf{禱告—誰曾與神款款深談呢?}
\newline
\newline
連結: \href{https://www.hkbibleconference.org/session-message/view/886}{\texttt{https://www.hkbibleconference.org/session-message/view/886}}
\newline
\newline
\hyperref[sec:885]{< < < PREV SERMON < < <}
~
\hyperref[sec:toc]{[返目錄]}
~
\hyperref[sec:887]{> > > NEXT SERMON > > >}
\newline
\newline


\newpage

\section{第69屆~港九培靈會~陳黔開~第六講}
\label{sec:887}
\textbf{恩賜 — 你尚待發掘的是甚麼?}
\newline
\newline
連結: \href{https://www.hkbibleconference.org/session-message/view/887}{\texttt{https://www.hkbibleconference.org/session-message/view/887}}
\newline
\newline
\hyperref[sec:886]{< < < PREV SERMON < < <}
~
\hyperref[sec:toc]{[返目錄]}
~
\hyperref[sec:888]{> > > NEXT SERMON > > >}
\newline
\newline


\newpage

\section{第69屆~港九培靈會~陳黔開~第七講}
\label{sec:888}
\textbf{肢體 — 你屬誰?}
\newline
\newline
連結: \href{https://www.hkbibleconference.org/session-message/view/888}{\texttt{https://www.hkbibleconference.org/session-message/view/888}}
\newline
\newline
\hyperref[sec:887]{< < < PREV SERMON < < <}
~
\hyperref[sec:toc]{[返目錄]}
~
\hyperref[sec:889]{> > > NEXT SERMON > > >}
\newline
\newline


\newpage

\section{第69屆~港九培靈會~陳黔開~第八講}
\label{sec:889}
\textbf{受死 — 有益?死地?}
\newline
\newline
連結: \href{https://www.hkbibleconference.org/session-message/view/889}{\texttt{https://www.hkbibleconference.org/session-message/view/889}}
\newline
\newline
\hyperref[sec:888]{< < < PREV SERMON < < <}
~
\hyperref[sec:toc]{[返目錄]}
~
\hyperref[sec:890]{> > > NEXT SERMON > > >}
\newline
\newline


\newpage

\section{第69屆~港九培靈會~黃子嘉~第一講}
\label{sec:890}
\textbf{福音是神彰顯榮耀與奧秘}
\newline
\newline
連結: \href{https://www.hkbibleconference.org/session-message/view/890}{\texttt{https://www.hkbibleconference.org/session-message/view/890}}
\newline
\newline
\hyperref[sec:889]{< < < PREV SERMON < < <}
~
\hyperref[sec:toc]{[返目錄]}
~
\hyperref[sec:891]{> > > NEXT SERMON > > >}
\newline
\newline
羅馬書 1:1-1:32
\newline
\begin{longtable}{cl}
\hline
\hline
章節 & 經文 (和合本修訂版)\\
\hline
1:1 & \begin{tabularx}{0.7\textwidth}{X} 基督耶穌的僕人保羅，蒙召為使徒，奉派傳神的福音。 \end{tabularx} \\ \\ \relax
1:2 & \begin{tabularx}{0.7\textwidth}{X} 這福音是神從前藉眾先知，在聖經上所應許的。 \end{tabularx} \\ \\ \relax
1:3 & \begin{tabularx}{0.7\textwidth}{X} 論到他兒子－我主耶穌基督，按肉體說，是從大衛後裔生的； \end{tabularx} \\ \\ \relax
1:4 & \begin{tabularx}{0.7\textwidth}{X} 按神聖的靈說，因從死人中復活，用大能顯明他是神的兒子。 \end{tabularx} \\ \\ \relax
1:5 & \begin{tabularx}{0.7\textwidth}{X} 我們從他蒙恩受了使徒的職分，為他的名在萬國中使人因信而順服， \end{tabularx} \\ \\ \relax
1:6 & \begin{tabularx}{0.7\textwidth}{X} 其中也有你們這蒙召屬耶穌基督的人。 \end{tabularx} \\ \\ \relax
1:7 & \begin{tabularx}{0.7\textwidth}{X} 我寫信給你們在羅馬、為神所愛、蒙召作聖徒的眾人。願恩惠、平安從我們的父神和主耶穌基督歸給你們！ \end{tabularx} \\ \\ \relax
1:8 & \begin{tabularx}{0.7\textwidth}{X} 首先，我靠著耶穌基督，為你們眾人感謝我的神，因你們的信德傳遍了天下。 \end{tabularx} \\ \\ \relax
1:9 & \begin{tabularx}{0.7\textwidth}{X} 我在他兒子的福音上，用心靈所事奉的神可以見證，我怎樣不住地提到你們， \end{tabularx} \\ \\ \relax
1:10 & \begin{tabularx}{0.7\textwidth}{X} 在我的禱告中常常懇求，或許照神的旨意，最終我能毫無阻礙地往你們那裡去。 \end{tabularx} \\ \\ \relax
1:11 & \begin{tabularx}{0.7\textwidth}{X} 因為我迫切地想見你們，要把一些屬靈的恩賜分給你們，使你們得以堅固， \end{tabularx} \\ \\ \relax
1:12 & \begin{tabularx}{0.7\textwidth}{X} 也可以說，我在你們中間，因你我彼此的信心而同得安慰。 \end{tabularx} \\ \\ \relax
1:13 & \begin{tabularx}{0.7\textwidth}{X} 弟兄們，我不願意你們不知道，我屢次計劃往你們那裡去，要在你們中間得些果子，如同在其餘的外邦人中一樣，只是到如今仍有攔阻。 \end{tabularx} \\ \\ \relax
1:14 & \begin{tabularx}{0.7\textwidth}{X} 無論是希臘人、未開化的人、聰明人、愚拙人，我都欠他們的債， \end{tabularx} \\ \\ \relax
1:15 & \begin{tabularx}{0.7\textwidth}{X} 所以願意盡我的力量把福音也傳給你們在羅馬的人。 \end{tabularx} \\ \\ \relax
1:16 & \begin{tabularx}{0.7\textwidth}{X} 我不以福音為恥；這福音本是神的大能，要救一切相信的，先是猶太人，後是希臘人。 \end{tabularx} \\ \\ \relax
1:17 & \begin{tabularx}{0.7\textwidth}{X} 因為神的義正在這福音上顯明出來；這義是本於信，以至於信。如經上所記：「義人必因信得生。」 \end{tabularx} \\ \\ \relax
1:18 & \begin{tabularx}{0.7\textwidth}{X} 原來，神的憤怒從天上顯明在一切不虔不義的人身上，就是那些行不義壓制真理的人。 \end{tabularx} \\ \\ \relax
1:19 & \begin{tabularx}{0.7\textwidth}{X} 神的事情，人所能知道的，原顯明在人心裡，因為神已經向他們顯明。 \end{tabularx} \\ \\ \relax
1:20 & \begin{tabularx}{0.7\textwidth}{X} 自從造天地以來，神的永能和神性是明明可知的，雖然眼不能見，但藉著所造之物就可以了解看見，叫人無可推諉。 \end{tabularx} \\ \\ \relax
1:21 & \begin{tabularx}{0.7\textwidth}{X} 因為，他們雖然知道神，卻不把他當作神榮耀他，也不感謝他。他們的思想變為虛妄，無知的心昏暗了。 \end{tabularx} \\ \\ \relax
1:22 & \begin{tabularx}{0.7\textwidth}{X} 他們自以為聰明，反成了愚昧， \end{tabularx} \\ \\ \relax
1:23 & \begin{tabularx}{0.7\textwidth}{X} 將不能朽壞之神的榮耀變為偶像，仿照必朽壞的人、飛禽、走獸、爬蟲的形像。 \end{tabularx} \\ \\ \relax
1:24 & \begin{tabularx}{0.7\textwidth}{X} 所以，神任憑他們隨著心裡的情慾行污穢的事，以致彼此羞辱自己的身體。 \end{tabularx} \\ \\ \relax
1:25 & \begin{tabularx}{0.7\textwidth}{X} 他們將神的真實變為虛謊，去敬拜事奉受造之物，不敬奉那造物的主—主是可稱頌的，直到永遠。阿們！ \end{tabularx} \\ \\ \relax
1:26 & \begin{tabularx}{0.7\textwidth}{X} 因此，神任憑他們放縱可羞恥的情慾。他們的女人把自然的關係變成違反自然的； \end{tabularx} \\ \\ \relax
1:27 & \begin{tabularx}{0.7\textwidth}{X} 男人也是如此，放棄了和女人自然的關係，慾火攻心，男的和男的彼此貪戀，行可恥的事，就在自己身上受這逆性行為當得的報應。 \end{tabularx} \\ \\ \relax
1:28 & \begin{tabularx}{0.7\textwidth}{X} 他們既然故意不認識神，神就任憑他們存扭曲的心，做那些不該做的事， \end{tabularx} \\ \\ \relax
1:29 & \begin{tabularx}{0.7\textwidth}{X} 裝滿了各樣不義、邪惡、貪婪、惡毒，滿心是嫉妒、兇殺、紛爭、詭詐、毒恨，又是毀謗的、 \end{tabularx} \\ \\ \relax
1:30 & \begin{tabularx}{0.7\textwidth}{X} 說人壞話的、怨恨神的、侮辱人的、狂傲的、自誇的、製造是非的、忤逆父母的、 \end{tabularx} \\ \\ \relax
1:31 & \begin{tabularx}{0.7\textwidth}{X} 頑梗不化的、言而無信的、無情無義的、不憐憫人的。 \end{tabularx} \\ \\ \relax
1:32 & \begin{tabularx}{0.7\textwidth}{X} 他們雖知道神判定做這樣事的人是該死的，然而他們不但自己去做，還贊同別人去做。 \end{tabularx} \\ \\
[1ex]
\hline
\hline
\end{longtable}


\newpage

\section{第69屆~港九培靈會~黃子嘉~第二講}
\label{sec:891}
\textbf{福音是神救人大能與大恩}
\newline
\newline
連結: \href{https://www.hkbibleconference.org/session-message/view/891}{\texttt{https://www.hkbibleconference.org/session-message/view/891}}
\newline
\newline
\hyperref[sec:890]{< < < PREV SERMON < < <}
~
\hyperref[sec:toc]{[返目錄]}
~
\hyperref[sec:892]{> > > NEXT SERMON > > >}
\newline
\newline
羅馬書 1:1-1:32
\newline
\begin{longtable}{cl}
\hline
\hline
章節 & 經文 (和合本修訂版)\\
\hline
1:1 & \begin{tabularx}{0.7\textwidth}{X} 基督耶穌的僕人保羅，蒙召為使徒，奉派傳神的福音。 \end{tabularx} \\ \\ \relax
1:2 & \begin{tabularx}{0.7\textwidth}{X} 這福音是神從前藉眾先知，在聖經上所應許的。 \end{tabularx} \\ \\ \relax
1:3 & \begin{tabularx}{0.7\textwidth}{X} 論到他兒子－我主耶穌基督，按肉體說，是從大衛後裔生的； \end{tabularx} \\ \\ \relax
1:4 & \begin{tabularx}{0.7\textwidth}{X} 按神聖的靈說，因從死人中復活，用大能顯明他是神的兒子。 \end{tabularx} \\ \\ \relax
1:5 & \begin{tabularx}{0.7\textwidth}{X} 我們從他蒙恩受了使徒的職分，為他的名在萬國中使人因信而順服， \end{tabularx} \\ \\ \relax
1:6 & \begin{tabularx}{0.7\textwidth}{X} 其中也有你們這蒙召屬耶穌基督的人。 \end{tabularx} \\ \\ \relax
1:7 & \begin{tabularx}{0.7\textwidth}{X} 我寫信給你們在羅馬、為神所愛、蒙召作聖徒的眾人。願恩惠、平安從我們的父神和主耶穌基督歸給你們！ \end{tabularx} \\ \\ \relax
1:8 & \begin{tabularx}{0.7\textwidth}{X} 首先，我靠著耶穌基督，為你們眾人感謝我的神，因你們的信德傳遍了天下。 \end{tabularx} \\ \\ \relax
1:9 & \begin{tabularx}{0.7\textwidth}{X} 我在他兒子的福音上，用心靈所事奉的神可以見證，我怎樣不住地提到你們， \end{tabularx} \\ \\ \relax
1:10 & \begin{tabularx}{0.7\textwidth}{X} 在我的禱告中常常懇求，或許照神的旨意，最終我能毫無阻礙地往你們那裡去。 \end{tabularx} \\ \\ \relax
1:11 & \begin{tabularx}{0.7\textwidth}{X} 因為我迫切地想見你們，要把一些屬靈的恩賜分給你們，使你們得以堅固， \end{tabularx} \\ \\ \relax
1:12 & \begin{tabularx}{0.7\textwidth}{X} 也可以說，我在你們中間，因你我彼此的信心而同得安慰。 \end{tabularx} \\ \\ \relax
1:13 & \begin{tabularx}{0.7\textwidth}{X} 弟兄們，我不願意你們不知道，我屢次計劃往你們那裡去，要在你們中間得些果子，如同在其餘的外邦人中一樣，只是到如今仍有攔阻。 \end{tabularx} \\ \\ \relax
1:14 & \begin{tabularx}{0.7\textwidth}{X} 無論是希臘人、未開化的人、聰明人、愚拙人，我都欠他們的債， \end{tabularx} \\ \\ \relax
1:15 & \begin{tabularx}{0.7\textwidth}{X} 所以願意盡我的力量把福音也傳給你們在羅馬的人。 \end{tabularx} \\ \\ \relax
1:16 & \begin{tabularx}{0.7\textwidth}{X} 我不以福音為恥；這福音本是神的大能，要救一切相信的，先是猶太人，後是希臘人。 \end{tabularx} \\ \\ \relax
1:17 & \begin{tabularx}{0.7\textwidth}{X} 因為神的義正在這福音上顯明出來；這義是本於信，以至於信。如經上所記：「義人必因信得生。」 \end{tabularx} \\ \\ \relax
1:18 & \begin{tabularx}{0.7\textwidth}{X} 原來，神的憤怒從天上顯明在一切不虔不義的人身上，就是那些行不義壓制真理的人。 \end{tabularx} \\ \\ \relax
1:19 & \begin{tabularx}{0.7\textwidth}{X} 神的事情，人所能知道的，原顯明在人心裡，因為神已經向他們顯明。 \end{tabularx} \\ \\ \relax
1:20 & \begin{tabularx}{0.7\textwidth}{X} 自從造天地以來，神的永能和神性是明明可知的，雖然眼不能見，但藉著所造之物就可以了解看見，叫人無可推諉。 \end{tabularx} \\ \\ \relax
1:21 & \begin{tabularx}{0.7\textwidth}{X} 因為，他們雖然知道神，卻不把他當作神榮耀他，也不感謝他。他們的思想變為虛妄，無知的心昏暗了。 \end{tabularx} \\ \\ \relax
1:22 & \begin{tabularx}{0.7\textwidth}{X} 他們自以為聰明，反成了愚昧， \end{tabularx} \\ \\ \relax
1:23 & \begin{tabularx}{0.7\textwidth}{X} 將不能朽壞之神的榮耀變為偶像，仿照必朽壞的人、飛禽、走獸、爬蟲的形像。 \end{tabularx} \\ \\ \relax
1:24 & \begin{tabularx}{0.7\textwidth}{X} 所以，神任憑他們隨著心裡的情慾行污穢的事，以致彼此羞辱自己的身體。 \end{tabularx} \\ \\ \relax
1:25 & \begin{tabularx}{0.7\textwidth}{X} 他們將神的真實變為虛謊，去敬拜事奉受造之物，不敬奉那造物的主—主是可稱頌的，直到永遠。阿們！ \end{tabularx} \\ \\ \relax
1:26 & \begin{tabularx}{0.7\textwidth}{X} 因此，神任憑他們放縱可羞恥的情慾。他們的女人把自然的關係變成違反自然的； \end{tabularx} \\ \\ \relax
1:27 & \begin{tabularx}{0.7\textwidth}{X} 男人也是如此，放棄了和女人自然的關係，慾火攻心，男的和男的彼此貪戀，行可恥的事，就在自己身上受這逆性行為當得的報應。 \end{tabularx} \\ \\ \relax
1:28 & \begin{tabularx}{0.7\textwidth}{X} 他們既然故意不認識神，神就任憑他們存扭曲的心，做那些不該做的事， \end{tabularx} \\ \\ \relax
1:29 & \begin{tabularx}{0.7\textwidth}{X} 裝滿了各樣不義、邪惡、貪婪、惡毒，滿心是嫉妒、兇殺、紛爭、詭詐、毒恨，又是毀謗的、 \end{tabularx} \\ \\ \relax
1:30 & \begin{tabularx}{0.7\textwidth}{X} 說人壞話的、怨恨神的、侮辱人的、狂傲的、自誇的、製造是非的、忤逆父母的、 \end{tabularx} \\ \\ \relax
1:31 & \begin{tabularx}{0.7\textwidth}{X} 頑梗不化的、言而無信的、無情無義的、不憐憫人的。 \end{tabularx} \\ \\ \relax
1:32 & \begin{tabularx}{0.7\textwidth}{X} 他們雖知道神判定做這樣事的人是該死的，然而他們不但自己去做，還贊同別人去做。 \end{tabularx} \\ \\
[1ex]
\hline
\hline
\end{longtable}


\newpage

\section{第69屆~港九培靈會~黃子嘉~第三講}
\label{sec:892}
\textbf{福音是神賜人因信而稱義也賜人平安喜樂}
\newline
\newline
連結: \href{https://www.hkbibleconference.org/session-message/view/892}{\texttt{https://www.hkbibleconference.org/session-message/view/892}}
\newline
\newline
\hyperref[sec:891]{< < < PREV SERMON < < <}
~
\hyperref[sec:toc]{[返目錄]}
~
\hyperref[sec:921]{> > > NEXT SERMON > > >}
\newline
\newline
羅馬書 1:1-1:32
\newline
\begin{longtable}{cl}
\hline
\hline
章節 & 經文 (和合本修訂版)\\
\hline
1:1 & \begin{tabularx}{0.7\textwidth}{X} 基督耶穌的僕人保羅，蒙召為使徒，奉派傳神的福音。 \end{tabularx} \\ \\ \relax
1:2 & \begin{tabularx}{0.7\textwidth}{X} 這福音是神從前藉眾先知，在聖經上所應許的。 \end{tabularx} \\ \\ \relax
1:3 & \begin{tabularx}{0.7\textwidth}{X} 論到他兒子－我主耶穌基督，按肉體說，是從大衛後裔生的； \end{tabularx} \\ \\ \relax
1:4 & \begin{tabularx}{0.7\textwidth}{X} 按神聖的靈說，因從死人中復活，用大能顯明他是神的兒子。 \end{tabularx} \\ \\ \relax
1:5 & \begin{tabularx}{0.7\textwidth}{X} 我們從他蒙恩受了使徒的職分，為他的名在萬國中使人因信而順服， \end{tabularx} \\ \\ \relax
1:6 & \begin{tabularx}{0.7\textwidth}{X} 其中也有你們這蒙召屬耶穌基督的人。 \end{tabularx} \\ \\ \relax
1:7 & \begin{tabularx}{0.7\textwidth}{X} 我寫信給你們在羅馬、為神所愛、蒙召作聖徒的眾人。願恩惠、平安從我們的父神和主耶穌基督歸給你們！ \end{tabularx} \\ \\ \relax
1:8 & \begin{tabularx}{0.7\textwidth}{X} 首先，我靠著耶穌基督，為你們眾人感謝我的神，因你們的信德傳遍了天下。 \end{tabularx} \\ \\ \relax
1:9 & \begin{tabularx}{0.7\textwidth}{X} 我在他兒子的福音上，用心靈所事奉的神可以見證，我怎樣不住地提到你們， \end{tabularx} \\ \\ \relax
1:10 & \begin{tabularx}{0.7\textwidth}{X} 在我的禱告中常常懇求，或許照神的旨意，最終我能毫無阻礙地往你們那裡去。 \end{tabularx} \\ \\ \relax
1:11 & \begin{tabularx}{0.7\textwidth}{X} 因為我迫切地想見你們，要把一些屬靈的恩賜分給你們，使你們得以堅固， \end{tabularx} \\ \\ \relax
1:12 & \begin{tabularx}{0.7\textwidth}{X} 也可以說，我在你們中間，因你我彼此的信心而同得安慰。 \end{tabularx} \\ \\ \relax
1:13 & \begin{tabularx}{0.7\textwidth}{X} 弟兄們，我不願意你們不知道，我屢次計劃往你們那裡去，要在你們中間得些果子，如同在其餘的外邦人中一樣，只是到如今仍有攔阻。 \end{tabularx} \\ \\ \relax
1:14 & \begin{tabularx}{0.7\textwidth}{X} 無論是希臘人、未開化的人、聰明人、愚拙人，我都欠他們的債， \end{tabularx} \\ \\ \relax
1:15 & \begin{tabularx}{0.7\textwidth}{X} 所以願意盡我的力量把福音也傳給你們在羅馬的人。 \end{tabularx} \\ \\ \relax
1:16 & \begin{tabularx}{0.7\textwidth}{X} 我不以福音為恥；這福音本是神的大能，要救一切相信的，先是猶太人，後是希臘人。 \end{tabularx} \\ \\ \relax
1:17 & \begin{tabularx}{0.7\textwidth}{X} 因為神的義正在這福音上顯明出來；這義是本於信，以至於信。如經上所記：「義人必因信得生。」 \end{tabularx} \\ \\ \relax
1:18 & \begin{tabularx}{0.7\textwidth}{X} 原來，神的憤怒從天上顯明在一切不虔不義的人身上，就是那些行不義壓制真理的人。 \end{tabularx} \\ \\ \relax
1:19 & \begin{tabularx}{0.7\textwidth}{X} 神的事情，人所能知道的，原顯明在人心裡，因為神已經向他們顯明。 \end{tabularx} \\ \\ \relax
1:20 & \begin{tabularx}{0.7\textwidth}{X} 自從造天地以來，神的永能和神性是明明可知的，雖然眼不能見，但藉著所造之物就可以了解看見，叫人無可推諉。 \end{tabularx} \\ \\ \relax
1:21 & \begin{tabularx}{0.7\textwidth}{X} 因為，他們雖然知道神，卻不把他當作神榮耀他，也不感謝他。他們的思想變為虛妄，無知的心昏暗了。 \end{tabularx} \\ \\ \relax
1:22 & \begin{tabularx}{0.7\textwidth}{X} 他們自以為聰明，反成了愚昧， \end{tabularx} \\ \\ \relax
1:23 & \begin{tabularx}{0.7\textwidth}{X} 將不能朽壞之神的榮耀變為偶像，仿照必朽壞的人、飛禽、走獸、爬蟲的形像。 \end{tabularx} \\ \\ \relax
1:24 & \begin{tabularx}{0.7\textwidth}{X} 所以，神任憑他們隨著心裡的情慾行污穢的事，以致彼此羞辱自己的身體。 \end{tabularx} \\ \\ \relax
1:25 & \begin{tabularx}{0.7\textwidth}{X} 他們將神的真實變為虛謊，去敬拜事奉受造之物，不敬奉那造物的主—主是可稱頌的，直到永遠。阿們！ \end{tabularx} \\ \\ \relax
1:26 & \begin{tabularx}{0.7\textwidth}{X} 因此，神任憑他們放縱可羞恥的情慾。他們的女人把自然的關係變成違反自然的； \end{tabularx} \\ \\ \relax
1:27 & \begin{tabularx}{0.7\textwidth}{X} 男人也是如此，放棄了和女人自然的關係，慾火攻心，男的和男的彼此貪戀，行可恥的事，就在自己身上受這逆性行為當得的報應。 \end{tabularx} \\ \\ \relax
1:28 & \begin{tabularx}{0.7\textwidth}{X} 他們既然故意不認識神，神就任憑他們存扭曲的心，做那些不該做的事， \end{tabularx} \\ \\ \relax
1:29 & \begin{tabularx}{0.7\textwidth}{X} 裝滿了各樣不義、邪惡、貪婪、惡毒，滿心是嫉妒、兇殺、紛爭、詭詐、毒恨，又是毀謗的、 \end{tabularx} \\ \\ \relax
1:30 & \begin{tabularx}{0.7\textwidth}{X} 說人壞話的、怨恨神的、侮辱人的、狂傲的、自誇的、製造是非的、忤逆父母的、 \end{tabularx} \\ \\ \relax
1:31 & \begin{tabularx}{0.7\textwidth}{X} 頑梗不化的、言而無信的、無情無義的、不憐憫人的。 \end{tabularx} \\ \\ \relax
1:32 & \begin{tabularx}{0.7\textwidth}{X} 他們雖知道神判定做這樣事的人是該死的，然而他們不但自己去做，還贊同別人去做。 \end{tabularx} \\ \\
[1ex]
\hline
\hline
\end{longtable}


\newpage

\section{第69屆~港九培靈會~黃子嘉~第四講}
\label{sec:921}
\textbf{福音是神使人稱義與成聖}
\newline
\newline
連結: \href{https://www.hkbibleconference.org/session-message/view/921}{\texttt{https://www.hkbibleconference.org/session-message/view/921}}
\newline
\newline
\hyperref[sec:892]{< < < PREV SERMON < < <}
~
\hyperref[sec:toc]{[返目錄]}
~
\hyperref[sec:922]{> > > NEXT SERMON > > >}
\newline
\newline
羅馬書 6:1-6:23
\newline
\begin{longtable}{cl}
\hline
\hline
章節 & 經文 (和合本修訂版)\\
\hline
6:1 & \begin{tabularx}{0.7\textwidth}{X} 這樣，我們要怎麼說呢？我們可以仍在罪中使恩典增多嗎？ \end{tabularx} \\ \\ \relax
6:2 & \begin{tabularx}{0.7\textwidth}{X} 絕對不可！我們向罪死了的人，豈可仍在罪中活著呢？ \end{tabularx} \\ \\ \relax
6:3 & \begin{tabularx}{0.7\textwidth}{X} 難道你們不知道，我們這受洗歸入基督耶穌的人，就是受洗歸入他的死嗎？ \end{tabularx} \\ \\ \relax
6:4 & \begin{tabularx}{0.7\textwidth}{X} 所以，我們藉著洗禮歸入死，和他一同埋葬，是要我們行事為人都有新生的樣子，像基督藉著父的榮耀從死人中復活一樣。 \end{tabularx} \\ \\ \relax
6:5 & \begin{tabularx}{0.7\textwidth}{X} 我們若與他合一，經歷與他一樣的死，也將經歷與他一樣的復活。 \end{tabularx} \\ \\ \relax
6:6 & \begin{tabularx}{0.7\textwidth}{X} 我們知道，我們的舊人和他同釘十字架，使罪身滅絕，叫我們不再作罪的奴隸， \end{tabularx} \\ \\ \relax
6:7 & \begin{tabularx}{0.7\textwidth}{X} 因為已死的人是脫離了罪。 \end{tabularx} \\ \\ \relax
6:8 & \begin{tabularx}{0.7\textwidth}{X} 我們若與基督同死，我們信也必與他同活， \end{tabularx} \\ \\ \relax
6:9 & \begin{tabularx}{0.7\textwidth}{X} 因為知道基督既從死人中復活，就不再死，死也不再作他的主了。 \end{tabularx} \\ \\ \relax
6:10 & \begin{tabularx}{0.7\textwidth}{X} 他死了，是對罪死，只這一次；他活，是對神活著。 \end{tabularx} \\ \\ \relax
6:11 & \begin{tabularx}{0.7\textwidth}{X} 這樣，你們也要看自己對罪是死的，在基督耶穌裡對神卻是活的。 \end{tabularx} \\ \\ \relax
6:12 & \begin{tabularx}{0.7\textwidth}{X} 所以，不要讓罪在你們必死的身上掌權，使你們順從身體的私慾。 \end{tabularx} \\ \\ \relax
6:13 & \begin{tabularx}{0.7\textwidth}{X} 也不要把你們的肢體獻給罪作不義的工具，倒要像從死人中活著的人，把自己獻給神，並把你們的肢體獻給神作義的工具。 \end{tabularx} \\ \\ \relax
6:14 & \begin{tabularx}{0.7\textwidth}{X} 罪必不能作你們的主，因你們不在律法之下，而是在恩典之下。 \end{tabularx} \\ \\ \relax
6:15 & \begin{tabularx}{0.7\textwidth}{X} 那又怎麼樣呢？我們在恩典之下，不在律法之下，就可以犯罪嗎？絕對不可！ \end{tabularx} \\ \\ \relax
6:16 & \begin{tabularx}{0.7\textwidth}{X} 難道你們不知道，你們獻自己作奴僕，順從誰就作誰的奴僕嗎？或作罪的奴隸，以至於死；或作順服的奴僕，以至於成義。 \end{tabularx} \\ \\ \relax
6:17 & \begin{tabularx}{0.7\textwidth}{X} 感謝神！因為你們從前雖然作罪的奴隸，現在卻從心裡順服了所傳給你們教導的典範。 \end{tabularx} \\ \\ \relax
6:18 & \begin{tabularx}{0.7\textwidth}{X} 你們既從罪裡得了釋放，就作了義的奴僕。 \end{tabularx} \\ \\ \relax
6:19 & \begin{tabularx}{0.7\textwidth}{X} 我因你們肉體的軟弱，就以人的觀點來說。你們從前怎樣把肢體獻給不潔不法作奴隸，以至於不法；現在也要照樣將肢體獻給義作奴僕，以至於成聖。 \end{tabularx} \\ \\ \relax
6:20 & \begin{tabularx}{0.7\textwidth}{X} 因為你們作罪的奴隸時，不被義所約束。 \end{tabularx} \\ \\ \relax
6:21 & \begin{tabularx}{0.7\textwidth}{X} 那麼，你們現在所看為羞恥的事，當時有甚麼果子呢？那些事的結局就是死。 \end{tabularx} \\ \\ \relax
6:22 & \begin{tabularx}{0.7\textwidth}{X} 但如今，你們既從罪裡得了釋放，作了神的奴僕，就結出果子，以至於成聖，那結局就是永生。 \end{tabularx} \\ \\ \relax
6:23 & \begin{tabularx}{0.7\textwidth}{X} 因為罪的工價乃是死；惟有神的恩賜，在我們的主基督耶穌裡，乃是永生。 \end{tabularx} \\ \\
[1ex]
\hline
\hline
\end{longtable}


\newpage

\section{第69屆~港九培靈會~黃子嘉~第五講}
\label{sec:922}
\textbf{福音是神賜人榮美與榮耀}
\newline
\newline
連結: \href{https://www.hkbibleconference.org/session-message/view/922}{\texttt{https://www.hkbibleconference.org/session-message/view/922}}
\newline
\newline
\hyperref[sec:921]{< < < PREV SERMON < < <}
~
\hyperref[sec:toc]{[返目錄]}
~
\hyperref[sec:923]{> > > NEXT SERMON > > >}
\newline
\newline
羅馬書 8:1-8:39
\newline
\begin{longtable}{cl}
\hline
\hline
章節 & 經文 (和合本修訂版)\\
\hline
8:1 & \begin{tabularx}{0.7\textwidth}{X} 如今，那些在基督耶穌裡的人就不被定罪了。 \end{tabularx} \\ \\ \relax
8:2 & \begin{tabularx}{0.7\textwidth}{X} 因為賜生命的聖靈的律，在基督耶穌裡從罪和死的律中把你釋放出來。 \end{tabularx} \\ \\ \relax
8:3 & \begin{tabularx}{0.7\textwidth}{X} 律法既因肉體軟弱而無能為力，神就差遣自己的兒子成為罪身的樣子，為了對付罪，在肉體中定了罪， \end{tabularx} \\ \\ \relax
8:4 & \begin{tabularx}{0.7\textwidth}{X} 為要使律法要求的義，實現在我們這不隨從肉體、只隨從聖靈去行的人身上。 \end{tabularx} \\ \\ \relax
8:5 & \begin{tabularx}{0.7\textwidth}{X} 因為，隨從肉體的人體貼肉體的事；隨從聖靈的人體貼聖靈的事。 \end{tabularx} \\ \\ \relax
8:6 & \begin{tabularx}{0.7\textwidth}{X} 體貼肉體就是死；體貼聖靈就是生命和平安。 \end{tabularx} \\ \\ \relax
8:7 & \begin{tabularx}{0.7\textwidth}{X} 因為體貼肉體就是與神為敵，對神的律法不順服，事實上也無法順服。 \end{tabularx} \\ \\ \relax
8:8 & \begin{tabularx}{0.7\textwidth}{X} 屬肉體的人無法使神喜悅。 \end{tabularx} \\ \\ \relax
8:9 & \begin{tabularx}{0.7\textwidth}{X} 如果神的靈住在你們裡面，你們就不屬肉體，而是屬聖靈了。人若沒有基督的靈，就不是屬基督的。 \end{tabularx} \\ \\ \relax
8:10 & \begin{tabularx}{0.7\textwidth}{X} 基督若在你們裡面，身體就因罪而死，靈卻因義而活。 \end{tabularx} \\ \\ \relax
8:11 & \begin{tabularx}{0.7\textwidth}{X} 然而，使耶穌從死人中復活的神的靈若住在你們裡面，那使基督從死人中復活的，也必藉著住在你們裡面的聖靈使你們必死的身體又活過來。 \end{tabularx} \\ \\ \relax
8:12 & \begin{tabularx}{0.7\textwidth}{X} 弟兄們，這樣看來，我們不是欠肉體的債去順從肉體而活。 \end{tabularx} \\ \\ \relax
8:13 & \begin{tabularx}{0.7\textwidth}{X} 你們若順從肉體活著，必定會死；若靠著聖靈把身體的惡行處死，就必存活。 \end{tabularx} \\ \\ \relax
8:14 & \begin{tabularx}{0.7\textwidth}{X} 因為凡被神的靈引導的都是神的兒子。 \end{tabularx} \\ \\ \relax
8:15 & \begin{tabularx}{0.7\textwidth}{X} 你們所領受的不是奴僕的靈，仍舊害怕；所領受的是兒子名分的靈，因此我們呼叫：「阿爸，父！」 \end{tabularx} \\ \\ \relax
8:16 & \begin{tabularx}{0.7\textwidth}{X} 聖靈自己與我們的靈一同見證我們是神的兒女。 \end{tabularx} \\ \\ \relax
8:17 & \begin{tabularx}{0.7\textwidth}{X} 若是兒女，就是後嗣，是神的後嗣，和基督同作後嗣。如果我們和他一同受苦，是要我們和他一同得榮耀。 \end{tabularx} \\ \\ \relax
8:18 & \begin{tabularx}{0.7\textwidth}{X} 我認為，現在的苦楚，若比起將來要顯示給我們的榮耀，是不足介意的。 \end{tabularx} \\ \\ \relax
8:19 & \begin{tabularx}{0.7\textwidth}{X} 受造之物切望等候神的眾子顯出來。 \end{tabularx} \\ \\ \relax
8:20 & \begin{tabularx}{0.7\textwidth}{X} 因為受造之物屈服在虛空之下，不是自己願意，而是因那使它屈服的叫他如此。 \end{tabularx} \\ \\ \relax
8:21 & \begin{tabularx}{0.7\textwidth}{X} 但受造之物仍然指望從敗壞的轄制下得釋放，得享神兒女榮耀的自由。 \end{tabularx} \\ \\ \relax
8:22 & \begin{tabularx}{0.7\textwidth}{X} 我們知道，一切受造之物一同呻吟，一同忍受陣痛，直到如今。 \end{tabularx} \\ \\ \relax
8:23 & \begin{tabularx}{0.7\textwidth}{X} 不但如此，就是我們這有聖靈作初熟果子的，也是自己內心呻吟，等候得著兒子的名分，就是我們的身體得救贖。 \end{tabularx} \\ \\ \relax
8:24 & \begin{tabularx}{0.7\textwidth}{X} 我們得救是在於盼望；可是看得見的盼望就不是盼望。誰還去盼望他所看得見的呢？ \end{tabularx} \\ \\ \relax
8:25 & \begin{tabularx}{0.7\textwidth}{X} 但我們若盼望那看不見的，我們就耐心等候。 \end{tabularx} \\ \\ \relax
8:26 & \begin{tabularx}{0.7\textwidth}{X} 同樣，我們的軟弱有聖靈幫助。我們本不知道當怎樣禱告，但是聖靈親自用無可言喻的嘆息替我們祈求。 \end{tabularx} \\ \\ \relax
8:27 & \begin{tabularx}{0.7\textwidth}{X} 那鑒察人心的知道聖靈所體貼的，因為聖靈照著神的旨意替聖徒祈求。 \end{tabularx} \\ \\ \relax
8:28 & \begin{tabularx}{0.7\textwidth}{X} 我們知道，萬事都互相效力，叫愛神的人得益處，就是按他旨意被召的人。 \end{tabularx} \\ \\ \relax
8:29 & \begin{tabularx}{0.7\textwidth}{X} 因為他所預知的人，他也預定他們效法他兒子的榜樣，使他兒子在許多弟兄中作長子。 \end{tabularx} \\ \\ \relax
8:30 & \begin{tabularx}{0.7\textwidth}{X} 他所預定的人，他又召他們來；所召來的人，他又稱他們為義；所稱為義的人，他又叫他們得榮耀。 \end{tabularx} \\ \\ \relax
8:31 & \begin{tabularx}{0.7\textwidth}{X} 既是這樣，我們對這些事還要怎麼說呢？神若幫助我們，誰能抵擋我們呢？ \end{tabularx} \\ \\ \relax
8:32 & \begin{tabularx}{0.7\textwidth}{X} 神既不顧惜自己的兒子，為我們眾人捨了他，豈不也把萬物和他一同白白地賜給我們嗎？ \end{tabularx} \\ \\ \relax
8:33 & \begin{tabularx}{0.7\textwidth}{X} 誰能控告神所揀選的人呢？有神稱他們為義了。 \end{tabularx} \\ \\ \relax
8:34 & \begin{tabularx}{0.7\textwidth}{X} 誰能定他們的罪呢？有基督耶穌 已經死了，而且復活了，現今在神的右邊，也替我們祈求。 \end{tabularx} \\ \\ \relax
8:35 & \begin{tabularx}{0.7\textwidth}{X} 誰能使我們與基督的愛隔絕呢？難道是患難嗎？是困苦嗎？是迫害嗎？是飢餓嗎？是赤身露體嗎？是危險嗎？是刀劍嗎？ \end{tabularx} \\ \\ \relax
8:36 & \begin{tabularx}{0.7\textwidth}{X} 如經上所記：「我們為你的緣故終日被殺；人看我們如將宰的羊。」 \end{tabularx} \\ \\ \relax
8:37 & \begin{tabularx}{0.7\textwidth}{X} 然而，靠著愛我們的主，在這一切的事上，我們已經得勝有餘了。 \end{tabularx} \\ \\ \relax
8:38 & \begin{tabularx}{0.7\textwidth}{X} 因為我深信，無論是死，是活，是天使，是掌權的，是有權能的 ，是現在的事，是將來的事， \end{tabularx} \\ \\ \relax
8:39 & \begin{tabularx}{0.7\textwidth}{X} 是高處的，是深處的，是別的受造之物，都不能使我們與神的愛隔絕，這愛是在我們的主基督耶穌裡的。 \end{tabularx} \\ \\
[1ex]
\hline
\hline
\end{longtable}


\newpage

\section{第69屆~港九培靈會~黃子嘉~第六講}
\label{sec:923}
\textbf{福音是神要人奉獻作活祭}
\newline
\newline
連結: \href{https://www.hkbibleconference.org/session-message/view/923}{\texttt{https://www.hkbibleconference.org/session-message/view/923}}
\newline
\newline
\hyperref[sec:922]{< < < PREV SERMON < < <}
~
\hyperref[sec:toc]{[返目錄]}
~
\hyperref[sec:924]{> > > NEXT SERMON > > >}
\newline
\newline
羅馬書 12:1-12:21
\newline
\begin{longtable}{cl}
\hline
\hline
章節 & 經文 (和合本修訂版)\\
\hline
12:1 & \begin{tabularx}{0.7\textwidth}{X} 所以，弟兄們，我以神的慈悲勸你們，將身體獻上當作活祭，是聖潔的，是神所喜悅的，你們如此事奉乃是理所當然的。 \end{tabularx} \\ \\ \relax
12:2 & \begin{tabularx}{0.7\textwidth}{X} 不要效法這個世界，只要心意更新而變化，叫你們察驗何為神的善良、純全、可喜悅的旨意。 \end{tabularx} \\ \\ \relax
12:3 & \begin{tabularx}{0.7\textwidth}{X} 我憑著所賜我的恩對你們每一位說：不要把自己看得太高，要照著神所分給各人的信心來衡量，看得合乎中道。 \end{tabularx} \\ \\ \relax
12:4 & \begin{tabularx}{0.7\textwidth}{X} 正如我們一個身子上有好些肢體，肢體也不都有一樣的用處。 \end{tabularx} \\ \\ \relax
12:5 & \begin{tabularx}{0.7\textwidth}{X} 這樣，我們許多人在基督裡是一個身體，互相聯絡作肢體。 \end{tabularx} \\ \\ \relax
12:6 & \begin{tabularx}{0.7\textwidth}{X} 按著所得的恩典，我們各有不同的恩賜：或說預言，要按著信心的程度說預言； \end{tabularx} \\ \\ \relax
12:7 & \begin{tabularx}{0.7\textwidth}{X} 或服事的，要專一服事；或教導的，要專一教導； \end{tabularx} \\ \\ \relax
12:8 & \begin{tabularx}{0.7\textwidth}{X} 或勸勉的，要專一勸勉；施捨的，要誠實；治理的，要殷勤；憐憫人的，要樂意。 \end{tabularx} \\ \\ \relax
12:9 & \begin{tabularx}{0.7\textwidth}{X} 愛，不可虛假；惡，要厭惡；善，要親近。 \end{tabularx} \\ \\ \relax
12:10 & \begin{tabularx}{0.7\textwidth}{X} 愛弟兄，要相親相愛；恭敬人，要彼此推讓； \end{tabularx} \\ \\ \relax
12:11 & \begin{tabularx}{0.7\textwidth}{X} 殷勤，不可懶惰。要靈裡火熱；常常服侍主。 \end{tabularx} \\ \\ \relax
12:12 & \begin{tabularx}{0.7\textwidth}{X} 在盼望中要喜樂；在患難中要忍耐；禱告要恆切。 \end{tabularx} \\ \\ \relax
12:13 & \begin{tabularx}{0.7\textwidth}{X} 聖徒有缺乏，要供給；異鄉客，要殷勤款待。 \end{tabularx} \\ \\ \relax
12:14 & \begin{tabularx}{0.7\textwidth}{X} 要祝福迫害你們的，要祝福，不可詛咒。 \end{tabularx} \\ \\ \relax
12:15 & \begin{tabularx}{0.7\textwidth}{X} 要與喜樂的人同樂；要與哀哭的人同哭。 \end{tabularx} \\ \\ \relax
12:16 & \begin{tabularx}{0.7\textwidth}{X} 要彼此同心，不要心高氣傲，倒要俯就卑微的人。不要自以為聰明。 \end{tabularx} \\ \\ \relax
12:17 & \begin{tabularx}{0.7\textwidth}{X} 不要以惡報惡，眾人以為美的事要留心去做。 \end{tabularx} \\ \\ \relax
12:18 & \begin{tabularx}{0.7\textwidth}{X} 若是可行，總要盡力與眾人和睦。 \end{tabularx} \\ \\ \relax
12:19 & \begin{tabularx}{0.7\textwidth}{X} 各位親愛的，不要自己伸冤，寧可給主的憤怒留地步，因為經上記著：「主說：『伸冤在我，我必報應。』」 \end{tabularx} \\ \\ \relax
12:20 & \begin{tabularx}{0.7\textwidth}{X} 不但如此，「你的仇敵若餓了，就給他吃；若渴了，就給他喝。因為你這樣做，就是把炭火堆在他的頭上。」 \end{tabularx} \\ \\ \relax
12:21 & \begin{tabularx}{0.7\textwidth}{X} 不要被惡所勝，反要以善勝惡。 \end{tabularx} \\ \\
[1ex]
\hline
\hline
\end{longtable}


\newpage

\section{第69屆~港九培靈會~黃子嘉~第七講}
\label{sec:924}
\textbf{福音是神要建立榮耀教會}
\newline
\newline
連結: \href{https://www.hkbibleconference.org/session-message/view/924}{\texttt{https://www.hkbibleconference.org/session-message/view/924}}
\newline
\newline
\hyperref[sec:923]{< < < PREV SERMON < < <}
~
\hyperref[sec:toc]{[返目錄]}
~
\hyperref[sec:925]{> > > NEXT SERMON > > >}
\newline
\newline
羅馬書 12:1-12:21
\newline
\begin{longtable}{cl}
\hline
\hline
章節 & 經文 (和合本修訂版)\\
\hline
12:1 & \begin{tabularx}{0.7\textwidth}{X} 所以，弟兄們，我以神的慈悲勸你們，將身體獻上當作活祭，是聖潔的，是神所喜悅的，你們如此事奉乃是理所當然的。 \end{tabularx} \\ \\ \relax
12:2 & \begin{tabularx}{0.7\textwidth}{X} 不要效法這個世界，只要心意更新而變化，叫你們察驗何為神的善良、純全、可喜悅的旨意。 \end{tabularx} \\ \\ \relax
12:3 & \begin{tabularx}{0.7\textwidth}{X} 我憑著所賜我的恩對你們每一位說：不要把自己看得太高，要照著神所分給各人的信心來衡量，看得合乎中道。 \end{tabularx} \\ \\ \relax
12:4 & \begin{tabularx}{0.7\textwidth}{X} 正如我們一個身子上有好些肢體，肢體也不都有一樣的用處。 \end{tabularx} \\ \\ \relax
12:5 & \begin{tabularx}{0.7\textwidth}{X} 這樣，我們許多人在基督裡是一個身體，互相聯絡作肢體。 \end{tabularx} \\ \\ \relax
12:6 & \begin{tabularx}{0.7\textwidth}{X} 按著所得的恩典，我們各有不同的恩賜：或說預言，要按著信心的程度說預言； \end{tabularx} \\ \\ \relax
12:7 & \begin{tabularx}{0.7\textwidth}{X} 或服事的，要專一服事；或教導的，要專一教導； \end{tabularx} \\ \\ \relax
12:8 & \begin{tabularx}{0.7\textwidth}{X} 或勸勉的，要專一勸勉；施捨的，要誠實；治理的，要殷勤；憐憫人的，要樂意。 \end{tabularx} \\ \\ \relax
12:9 & \begin{tabularx}{0.7\textwidth}{X} 愛，不可虛假；惡，要厭惡；善，要親近。 \end{tabularx} \\ \\ \relax
12:10 & \begin{tabularx}{0.7\textwidth}{X} 愛弟兄，要相親相愛；恭敬人，要彼此推讓； \end{tabularx} \\ \\ \relax
12:11 & \begin{tabularx}{0.7\textwidth}{X} 殷勤，不可懶惰。要靈裡火熱；常常服侍主。 \end{tabularx} \\ \\ \relax
12:12 & \begin{tabularx}{0.7\textwidth}{X} 在盼望中要喜樂；在患難中要忍耐；禱告要恆切。 \end{tabularx} \\ \\ \relax
12:13 & \begin{tabularx}{0.7\textwidth}{X} 聖徒有缺乏，要供給；異鄉客，要殷勤款待。 \end{tabularx} \\ \\ \relax
12:14 & \begin{tabularx}{0.7\textwidth}{X} 要祝福迫害你們的，要祝福，不可詛咒。 \end{tabularx} \\ \\ \relax
12:15 & \begin{tabularx}{0.7\textwidth}{X} 要與喜樂的人同樂；要與哀哭的人同哭。 \end{tabularx} \\ \\ \relax
12:16 & \begin{tabularx}{0.7\textwidth}{X} 要彼此同心，不要心高氣傲，倒要俯就卑微的人。不要自以為聰明。 \end{tabularx} \\ \\ \relax
12:17 & \begin{tabularx}{0.7\textwidth}{X} 不要以惡報惡，眾人以為美的事要留心去做。 \end{tabularx} \\ \\ \relax
12:18 & \begin{tabularx}{0.7\textwidth}{X} 若是可行，總要盡力與眾人和睦。 \end{tabularx} \\ \\ \relax
12:19 & \begin{tabularx}{0.7\textwidth}{X} 各位親愛的，不要自己伸冤，寧可給主的憤怒留地步，因為經上記著：「主說：『伸冤在我，我必報應。』」 \end{tabularx} \\ \\ \relax
12:20 & \begin{tabularx}{0.7\textwidth}{X} 不但如此，「你的仇敵若餓了，就給他吃；若渴了，就給他喝。因為你這樣做，就是把炭火堆在他的頭上。」 \end{tabularx} \\ \\ \relax
12:21 & \begin{tabularx}{0.7\textwidth}{X} 不要被惡所勝，反要以善勝惡。 \end{tabularx} \\ \\
[1ex]
\hline
\hline
\end{longtable}


\newpage

\section{第69屆~港九培靈會~黃子嘉~第八講}
\label{sec:925}
\textbf{福音是神要建立美善社會}
\newline
\newline
連結: \href{https://www.hkbibleconference.org/session-message/view/925}{\texttt{https://www.hkbibleconference.org/session-message/view/925}}
\newline
\newline
\hyperref[sec:924]{< < < PREV SERMON < < <}
~
\hyperref[sec:toc]{[返目錄]}
~
\hyperref[sec:926]{> > > NEXT SERMON > > >}
\newline
\newline
羅馬書 12:1-12:21
\newline
\begin{longtable}{cl}
\hline
\hline
章節 & 經文 (和合本修訂版)\\
\hline
12:1 & \begin{tabularx}{0.7\textwidth}{X} 所以，弟兄們，我以神的慈悲勸你們，將身體獻上當作活祭，是聖潔的，是神所喜悅的，你們如此事奉乃是理所當然的。 \end{tabularx} \\ \\ \relax
12:2 & \begin{tabularx}{0.7\textwidth}{X} 不要效法這個世界，只要心意更新而變化，叫你們察驗何為神的善良、純全、可喜悅的旨意。 \end{tabularx} \\ \\ \relax
12:3 & \begin{tabularx}{0.7\textwidth}{X} 我憑著所賜我的恩對你們每一位說：不要把自己看得太高，要照著神所分給各人的信心來衡量，看得合乎中道。 \end{tabularx} \\ \\ \relax
12:4 & \begin{tabularx}{0.7\textwidth}{X} 正如我們一個身子上有好些肢體，肢體也不都有一樣的用處。 \end{tabularx} \\ \\ \relax
12:5 & \begin{tabularx}{0.7\textwidth}{X} 這樣，我們許多人在基督裡是一個身體，互相聯絡作肢體。 \end{tabularx} \\ \\ \relax
12:6 & \begin{tabularx}{0.7\textwidth}{X} 按著所得的恩典，我們各有不同的恩賜：或說預言，要按著信心的程度說預言； \end{tabularx} \\ \\ \relax
12:7 & \begin{tabularx}{0.7\textwidth}{X} 或服事的，要專一服事；或教導的，要專一教導； \end{tabularx} \\ \\ \relax
12:8 & \begin{tabularx}{0.7\textwidth}{X} 或勸勉的，要專一勸勉；施捨的，要誠實；治理的，要殷勤；憐憫人的，要樂意。 \end{tabularx} \\ \\ \relax
12:9 & \begin{tabularx}{0.7\textwidth}{X} 愛，不可虛假；惡，要厭惡；善，要親近。 \end{tabularx} \\ \\ \relax
12:10 & \begin{tabularx}{0.7\textwidth}{X} 愛弟兄，要相親相愛；恭敬人，要彼此推讓； \end{tabularx} \\ \\ \relax
12:11 & \begin{tabularx}{0.7\textwidth}{X} 殷勤，不可懶惰。要靈裡火熱；常常服侍主。 \end{tabularx} \\ \\ \relax
12:12 & \begin{tabularx}{0.7\textwidth}{X} 在盼望中要喜樂；在患難中要忍耐；禱告要恆切。 \end{tabularx} \\ \\ \relax
12:13 & \begin{tabularx}{0.7\textwidth}{X} 聖徒有缺乏，要供給；異鄉客，要殷勤款待。 \end{tabularx} \\ \\ \relax
12:14 & \begin{tabularx}{0.7\textwidth}{X} 要祝福迫害你們的，要祝福，不可詛咒。 \end{tabularx} \\ \\ \relax
12:15 & \begin{tabularx}{0.7\textwidth}{X} 要與喜樂的人同樂；要與哀哭的人同哭。 \end{tabularx} \\ \\ \relax
12:16 & \begin{tabularx}{0.7\textwidth}{X} 要彼此同心，不要心高氣傲，倒要俯就卑微的人。不要自以為聰明。 \end{tabularx} \\ \\ \relax
12:17 & \begin{tabularx}{0.7\textwidth}{X} 不要以惡報惡，眾人以為美的事要留心去做。 \end{tabularx} \\ \\ \relax
12:18 & \begin{tabularx}{0.7\textwidth}{X} 若是可行，總要盡力與眾人和睦。 \end{tabularx} \\ \\ \relax
12:19 & \begin{tabularx}{0.7\textwidth}{X} 各位親愛的，不要自己伸冤，寧可給主的憤怒留地步，因為經上記著：「主說：『伸冤在我，我必報應。』」 \end{tabularx} \\ \\ \relax
12:20 & \begin{tabularx}{0.7\textwidth}{X} 不但如此，「你的仇敵若餓了，就給他吃；若渴了，就給他喝。因為你這樣做，就是把炭火堆在他的頭上。」 \end{tabularx} \\ \\ \relax
12:21 & \begin{tabularx}{0.7\textwidth}{X} 不要被惡所勝，反要以善勝惡。 \end{tabularx} \\ \\
[1ex]
\hline
\hline
\end{longtable}


\newpage

\section{第69屆~港九培靈會~黃子嘉~第九講}
\label{sec:926}
\textbf{福音是神要託付榮耀使命}
\newline
\newline
連結: \href{https://www.hkbibleconference.org/session-message/view/926}{\texttt{https://www.hkbibleconference.org/session-message/view/926}}
\newline
\newline
\hyperref[sec:925]{< < < PREV SERMON < < <}
~
\hyperref[sec:toc]{[返目錄]}
~
\hyperref[sec:726]{> > > NEXT SERMON > > >}
\newline
\newline
羅馬書 1:1-1:32
\newline
\begin{longtable}{cl}
\hline
\hline
章節 & 經文 (和合本修訂版)\\
\hline
1:1 & \begin{tabularx}{0.7\textwidth}{X} 基督耶穌的僕人保羅，蒙召為使徒，奉派傳神的福音。 \end{tabularx} \\ \\ \relax
1:2 & \begin{tabularx}{0.7\textwidth}{X} 這福音是神從前藉眾先知，在聖經上所應許的。 \end{tabularx} \\ \\ \relax
1:3 & \begin{tabularx}{0.7\textwidth}{X} 論到他兒子－我主耶穌基督，按肉體說，是從大衛後裔生的； \end{tabularx} \\ \\ \relax
1:4 & \begin{tabularx}{0.7\textwidth}{X} 按神聖的靈說，因從死人中復活，用大能顯明他是神的兒子。 \end{tabularx} \\ \\ \relax
1:5 & \begin{tabularx}{0.7\textwidth}{X} 我們從他蒙恩受了使徒的職分，為他的名在萬國中使人因信而順服， \end{tabularx} \\ \\ \relax
1:6 & \begin{tabularx}{0.7\textwidth}{X} 其中也有你們這蒙召屬耶穌基督的人。 \end{tabularx} \\ \\ \relax
1:7 & \begin{tabularx}{0.7\textwidth}{X} 我寫信給你們在羅馬、為神所愛、蒙召作聖徒的眾人。願恩惠、平安從我們的父神和主耶穌基督歸給你們！ \end{tabularx} \\ \\ \relax
1:8 & \begin{tabularx}{0.7\textwidth}{X} 首先，我靠著耶穌基督，為你們眾人感謝我的神，因你們的信德傳遍了天下。 \end{tabularx} \\ \\ \relax
1:9 & \begin{tabularx}{0.7\textwidth}{X} 我在他兒子的福音上，用心靈所事奉的神可以見證，我怎樣不住地提到你們， \end{tabularx} \\ \\ \relax
1:10 & \begin{tabularx}{0.7\textwidth}{X} 在我的禱告中常常懇求，或許照神的旨意，最終我能毫無阻礙地往你們那裡去。 \end{tabularx} \\ \\ \relax
1:11 & \begin{tabularx}{0.7\textwidth}{X} 因為我迫切地想見你們，要把一些屬靈的恩賜分給你們，使你們得以堅固， \end{tabularx} \\ \\ \relax
1:12 & \begin{tabularx}{0.7\textwidth}{X} 也可以說，我在你們中間，因你我彼此的信心而同得安慰。 \end{tabularx} \\ \\ \relax
1:13 & \begin{tabularx}{0.7\textwidth}{X} 弟兄們，我不願意你們不知道，我屢次計劃往你們那裡去，要在你們中間得些果子，如同在其餘的外邦人中一樣，只是到如今仍有攔阻。 \end{tabularx} \\ \\ \relax
1:14 & \begin{tabularx}{0.7\textwidth}{X} 無論是希臘人、未開化的人、聰明人、愚拙人，我都欠他們的債， \end{tabularx} \\ \\ \relax
1:15 & \begin{tabularx}{0.7\textwidth}{X} 所以願意盡我的力量把福音也傳給你們在羅馬的人。 \end{tabularx} \\ \\ \relax
1:16 & \begin{tabularx}{0.7\textwidth}{X} 我不以福音為恥；這福音本是神的大能，要救一切相信的，先是猶太人，後是希臘人。 \end{tabularx} \\ \\ \relax
1:17 & \begin{tabularx}{0.7\textwidth}{X} 因為神的義正在這福音上顯明出來；這義是本於信，以至於信。如經上所記：「義人必因信得生。」 \end{tabularx} \\ \\ \relax
1:18 & \begin{tabularx}{0.7\textwidth}{X} 原來，神的憤怒從天上顯明在一切不虔不義的人身上，就是那些行不義壓制真理的人。 \end{tabularx} \\ \\ \relax
1:19 & \begin{tabularx}{0.7\textwidth}{X} 神的事情，人所能知道的，原顯明在人心裡，因為神已經向他們顯明。 \end{tabularx} \\ \\ \relax
1:20 & \begin{tabularx}{0.7\textwidth}{X} 自從造天地以來，神的永能和神性是明明可知的，雖然眼不能見，但藉著所造之物就可以了解看見，叫人無可推諉。 \end{tabularx} \\ \\ \relax
1:21 & \begin{tabularx}{0.7\textwidth}{X} 因為，他們雖然知道神，卻不把他當作神榮耀他，也不感謝他。他們的思想變為虛妄，無知的心昏暗了。 \end{tabularx} \\ \\ \relax
1:22 & \begin{tabularx}{0.7\textwidth}{X} 他們自以為聰明，反成了愚昧， \end{tabularx} \\ \\ \relax
1:23 & \begin{tabularx}{0.7\textwidth}{X} 將不能朽壞之神的榮耀變為偶像，仿照必朽壞的人、飛禽、走獸、爬蟲的形像。 \end{tabularx} \\ \\ \relax
1:24 & \begin{tabularx}{0.7\textwidth}{X} 所以，神任憑他們隨著心裡的情慾行污穢的事，以致彼此羞辱自己的身體。 \end{tabularx} \\ \\ \relax
1:25 & \begin{tabularx}{0.7\textwidth}{X} 他們將神的真實變為虛謊，去敬拜事奉受造之物，不敬奉那造物的主—主是可稱頌的，直到永遠。阿們！ \end{tabularx} \\ \\ \relax
1:26 & \begin{tabularx}{0.7\textwidth}{X} 因此，神任憑他們放縱可羞恥的情慾。他們的女人把自然的關係變成違反自然的； \end{tabularx} \\ \\ \relax
1:27 & \begin{tabularx}{0.7\textwidth}{X} 男人也是如此，放棄了和女人自然的關係，慾火攻心，男的和男的彼此貪戀，行可恥的事，就在自己身上受這逆性行為當得的報應。 \end{tabularx} \\ \\ \relax
1:28 & \begin{tabularx}{0.7\textwidth}{X} 他們既然故意不認識神，神就任憑他們存扭曲的心，做那些不該做的事， \end{tabularx} \\ \\ \relax
1:29 & \begin{tabularx}{0.7\textwidth}{X} 裝滿了各樣不義、邪惡、貪婪、惡毒，滿心是嫉妒、兇殺、紛爭、詭詐、毒恨，又是毀謗的、 \end{tabularx} \\ \\ \relax
1:30 & \begin{tabularx}{0.7\textwidth}{X} 說人壞話的、怨恨神的、侮辱人的、狂傲的、自誇的、製造是非的、忤逆父母的、 \end{tabularx} \\ \\ \relax
1:31 & \begin{tabularx}{0.7\textwidth}{X} 頑梗不化的、言而無信的、無情無義的、不憐憫人的。 \end{tabularx} \\ \\ \relax
1:32 & \begin{tabularx}{0.7\textwidth}{X} 他們雖知道神判定做這樣事的人是該死的，然而他們不但自己去做，還贊同別人去做。 \end{tabularx} \\ \\
[1ex]
\hline
\hline
\end{longtable}


\newpage

\section{第70屆~港九培靈會~曾立華~第一講}
\label{sec:726}
\textbf{為追求誠信正直感的禱告}
\newline
\newline
連結: \href{https://www.hkbibleconference.org/session-message/view/726}{\texttt{https://www.hkbibleconference.org/session-message/view/726}}
\newline
\newline
\hyperref[sec:926]{< < < PREV SERMON < < <}
~
\hyperref[sec:toc]{[返目錄]}
~
\hyperref[sec:727]{> > > NEXT SERMON > > >}
\newline
\newline


\newpage

\section{第70屆~港九培靈會~曾立華~第二講}
\label{sec:727}
\textbf{在困境中勇於表達誠信正直的品格}
\newline
\newline
連結: \href{https://www.hkbibleconference.org/session-message/view/727}{\texttt{https://www.hkbibleconference.org/session-message/view/727}}
\newline
\newline
\hyperref[sec:726]{< < < PREV SERMON < < <}
~
\hyperref[sec:toc]{[返目錄]}
~
\hyperref[sec:728]{> > > NEXT SERMON > > >}
\newline
\newline
腓立比書 1:12-1:20
\newline
\begin{longtable}{cl}
\hline
\hline
章節 & 經文 (和合本修訂版)\\
\hline
1:12 & \begin{tabularx}{0.7\textwidth}{X} 弟兄們，我要你們知道，我所遭遇的事反而使福音更興旺， \end{tabularx} \\ \\ \relax
1:13 & \begin{tabularx}{0.7\textwidth}{X} 以致御營全軍和其餘的人都知道我是為基督的緣故受捆鎖的； \end{tabularx} \\ \\ \relax
1:14 & \begin{tabularx}{0.7\textwidth}{X} 而且那在主裡的弟兄，多半都因我受的捆鎖而篤信不疑，越發放膽無所懼怕地傳道。 \end{tabularx} \\ \\ \relax
1:15 & \begin{tabularx}{0.7\textwidth}{X} 有些人傳基督是出於嫉妒紛爭；有些人是出於好意。 \end{tabularx} \\ \\ \relax
1:16 & \begin{tabularx}{0.7\textwidth}{X} 後者是出於愛心，知道我奉差遣是為福音辯護的。 \end{tabularx} \\ \\ \relax
1:17 & \begin{tabularx}{0.7\textwidth}{X} 前者傳基督是出於自私，動機不純，企圖要加增我捆鎖的苦楚。 \end{tabularx} \\ \\ \relax
1:18 & \begin{tabularx}{0.7\textwidth}{X} 這又何妨呢？或是假意或是真心，無論如何，只要基督被傳開了，為此我就歡喜。我還要歡喜， \end{tabularx} \\ \\ \relax
1:19 & \begin{tabularx}{0.7\textwidth}{X} 因為我知道，這事藉著你們的祈禱和耶穌基督的靈的幫助，終必使我得到釋放。 \end{tabularx} \\ \\ \relax
1:20 & \begin{tabularx}{0.7\textwidth}{X} 這就是我所切慕、所盼望的：沒有一事能使我羞愧；反倒凡事坦然無懼，無論是生是死，總要讓基督在我身上照常顯大。 \end{tabularx} \\ \\
[1ex]
\hline
\hline
\end{longtable}


\newpage

\section{第70屆~港九培靈會~曾立華~第三講}
\label{sec:728}
\textbf{恰當的正直生活形態}
\newline
\newline
連結: \href{https://www.hkbibleconference.org/session-message/view/728}{\texttt{https://www.hkbibleconference.org/session-message/view/728}}
\newline
\newline
\hyperref[sec:727]{< < < PREV SERMON < < <}
~
\hyperref[sec:toc]{[返目錄]}
~
\hyperref[sec:729]{> > > NEXT SERMON > > >}
\newline
\newline
腓立比書 1:21-1:30
\newline
\begin{longtable}{cl}
\hline
\hline
章節 & 經文 (和合本修訂版)\\
\hline
1:21 & \begin{tabularx}{0.7\textwidth}{X} 因為我活著就是基督，死了就有益處。 \end{tabularx} \\ \\ \relax
1:22 & \begin{tabularx}{0.7\textwidth}{X} 但是，我在肉身活著，若能有工作的成果，我就不知道該挑選甚麼。 \end{tabularx} \\ \\ \relax
1:23 & \begin{tabularx}{0.7\textwidth}{X} 我處在兩難之間：我情願離世與基督同在，因為這是好得無比的； \end{tabularx} \\ \\ \relax
1:24 & \begin{tabularx}{0.7\textwidth}{X} 然而，我為你們肉身活著更加要緊。 \end{tabularx} \\ \\ \relax
1:25 & \begin{tabularx}{0.7\textwidth}{X} 既然我這樣深信，就知道仍要留在世間，且與你們眾人一起存留，使你們在所信的道上又長進又喜樂， \end{tabularx} \\ \\ \relax
1:26 & \begin{tabularx}{0.7\textwidth}{X} 為了我再到你們那裡時，你們在基督耶穌裡的誇耀越發加增。 \end{tabularx} \\ \\ \relax
1:27 & \begin{tabularx}{0.7\textwidth}{X} 最重要的是：你們行事為人要與基督的福音相稱，這樣，無論我來見你們，或不在你們那裡，都可以聽到你們的景況，知道你們同有一個心志，站立得穩，為福音的信仰齊心努力， \end{tabularx} \\ \\ \relax
1:28 & \begin{tabularx}{0.7\textwidth}{X} 絲毫不怕敵人的威脅；以此證明他們會沉淪，你們會得救，這是出於神。 \end{tabularx} \\ \\ \relax
1:29 & \begin{tabularx}{0.7\textwidth}{X} 因為你們蒙恩，不但得以信服基督，而且要為他受苦。 \end{tabularx} \\ \\ \relax
1:30 & \begin{tabularx}{0.7\textwidth}{X} 你們的爭戰，就與你們曾在我身上見過、現在所聽到的是一樣的。 \end{tabularx} \\ \\
[1ex]
\hline
\hline
\end{longtable}


\newpage

\section{第70屆~港九培靈會~曾立華~第四講}
\label{sec:729}
\textbf{同心培育基督謙卑正直的氣質}
\newline
\newline
連結: \href{https://www.hkbibleconference.org/session-message/view/729}{\texttt{https://www.hkbibleconference.org/session-message/view/729}}
\newline
\newline
\hyperref[sec:728]{< < < PREV SERMON < < <}
~
\hyperref[sec:toc]{[返目錄]}
~
\hyperref[sec:730]{> > > NEXT SERMON > > >}
\newline
\newline
腓立比書 2:1-2:11
\newline
\begin{longtable}{cl}
\hline
\hline
章節 & 經文 (和合本修訂版)\\
\hline
2:1 & \begin{tabularx}{0.7\textwidth}{X} 所以，在基督裡若有任何勸勉，若有任何愛心的安慰，若有任何聖靈的團契，若有任何慈悲憐憫， \end{tabularx} \\ \\ \relax
2:2 & \begin{tabularx}{0.7\textwidth}{X} 你們就要意志相同，愛心相同，有一致的心思，一致的想法，使我的喜樂得以滿足。 \end{tabularx} \\ \\ \relax
2:3 & \begin{tabularx}{0.7\textwidth}{X} 凡事不可自私自利，不可貪圖虛榮；只要心存謙卑，各人看別人比自己強。 \end{tabularx} \\ \\ \relax
2:4 & \begin{tabularx}{0.7\textwidth}{X} 各人不要單顧自己的事，也要顧別人的事。 \end{tabularx} \\ \\ \relax
2:5 & \begin{tabularx}{0.7\textwidth}{X} 你們當以基督耶穌的心為心： \end{tabularx} \\ \\ \relax
2:6 & \begin{tabularx}{0.7\textwidth}{X} 他本有神的形像，卻不堅持自己與神同等； \end{tabularx} \\ \\ \relax
2:7 & \begin{tabularx}{0.7\textwidth}{X} 反倒虛己，取了奴僕的形像，成為人的樣式；既有人的樣子， \end{tabularx} \\ \\ \relax
2:8 & \begin{tabularx}{0.7\textwidth}{X} 就謙卑自己，存心順服，以至於死，且死在十字架上。 \end{tabularx} \\ \\ \relax
2:9 & \begin{tabularx}{0.7\textwidth}{X} 所以神把他升為至高，又賜給他超乎萬名之上的名， \end{tabularx} \\ \\ \relax
2:10 & \begin{tabularx}{0.7\textwidth}{X} 使一切在天上的、地上的和地底下的，因耶穌的名，眾膝都要跪下， \end{tabularx} \\ \\ \relax
2:11 & \begin{tabularx}{0.7\textwidth}{X} 眾口都要宣認：耶穌基督是主，歸榮耀給父神。 \end{tabularx} \\ \\
[1ex]
\hline
\hline
\end{longtable}


\newpage

\section{第70屆~港九培靈會~曾立華~第五講}
\label{sec:730}
\textbf{誠信正直的領袖}
\newline
\newline
連結: \href{https://www.hkbibleconference.org/session-message/view/730}{\texttt{https://www.hkbibleconference.org/session-message/view/730}}
\newline
\newline
\hyperref[sec:729]{< < < PREV SERMON < < <}
~
\hyperref[sec:toc]{[返目錄]}
~
\hyperref[sec:733]{> > > NEXT SERMON > > >}
\newline
\newline
腓立比書 2:12-2:30
\newline
\begin{longtable}{cl}
\hline
\hline
章節 & 經文 (和合本修訂版)\\
\hline
2:12 & \begin{tabularx}{0.7\textwidth}{X} 我親愛的，這樣看來，你們向來是順服的，不但我在你們那裡，就是我現在不在你們那裡的時候更是順服的，就當恐懼戰兢完成你們自己得救的事； \end{tabularx} \\ \\ \relax
2:13 & \begin{tabularx}{0.7\textwidth}{X} 因為是神在你們心裡運行，使你們又立志又實行，為要成就他的美意。 \end{tabularx} \\ \\ \relax
2:14 & \begin{tabularx}{0.7\textwidth}{X} 你們無論做甚麼事，都不要發怨言起爭論， \end{tabularx} \\ \\ \relax
2:15 & \begin{tabularx}{0.7\textwidth}{X} 好使你們無可指責，誠實無偽，在這彎曲悖謬的世代作神無瑕疵的兒女。你們在這世代中要像明光照耀， \end{tabularx} \\ \\ \relax
2:16 & \begin{tabularx}{0.7\textwidth}{X} 將生命的道顯明出來，使我在基督的日子得以誇耀我沒有白跑，也沒有徒勞。 \end{tabularx} \\ \\ \relax
2:17 & \begin{tabularx}{0.7\textwidth}{X} 我以你們的信心為供獻的祭物，我若被澆獻在其上也是喜樂，並且與你們眾人一同喜樂。 \end{tabularx} \\ \\ \relax
2:18 & \begin{tabularx}{0.7\textwidth}{X} 你們也要照樣喜樂，並且與我一同喜樂。 \end{tabularx} \\ \\ \relax
2:19 & \begin{tabularx}{0.7\textwidth}{X} 我靠主耶穌希望很快能差提摩太去見你們，好讓我知道你們的事而心裡得著安慰。 \end{tabularx} \\ \\ \relax
2:20 & \begin{tabularx}{0.7\textwidth}{X} 因為我沒有別人與我同心，真正關懷你們的事。 \end{tabularx} \\ \\ \relax
2:21 & \begin{tabularx}{0.7\textwidth}{X} 其他的人都求自己的事，並不求耶穌基督的事。 \end{tabularx} \\ \\ \relax
2:22 & \begin{tabularx}{0.7\textwidth}{X} 但你們知道提摩太是經得起考驗的，他與我為了福音一同服侍，待我像兒子待父親一樣。 \end{tabularx} \\ \\ \relax
2:23 & \begin{tabularx}{0.7\textwidth}{X} 所以，我一看出我的事怎樣了結，我希望立刻差他去， \end{tabularx} \\ \\ \relax
2:24 & \begin{tabularx}{0.7\textwidth}{X} 但我靠著主自信我不久也會去。 \end{tabularx} \\ \\ \relax
2:25 & \begin{tabularx}{0.7\textwidth}{X} 然而，我想必須差以巴弗提到你們那裡去。他是我的弟兄、同工和戰友，是你們差遣來供應我需要的。 \end{tabularx} \\ \\ \relax
2:26 & \begin{tabularx}{0.7\textwidth}{X} 他很想念你們眾人，並且極其難過，因為你們聽見他病了。 \end{tabularx} \\ \\ \relax
2:27 & \begin{tabularx}{0.7\textwidth}{X} 他真的生病了，幾乎要死。然而神憐憫他，不但憐憫他，也憐憫我，免得我憂上加憂。 \end{tabularx} \\ \\ \relax
2:28 & \begin{tabularx}{0.7\textwidth}{X} 所以，我更要盡快送他回去，好讓你們再見到他而喜樂，我也可以減少憂愁。 \end{tabularx} \\ \\ \relax
2:29 & \begin{tabularx}{0.7\textwidth}{X} 故此，你們要在主裡歡歡喜喜地接待他，而且要尊重這樣的人， \end{tabularx} \\ \\ \relax
2:30 & \begin{tabularx}{0.7\textwidth}{X} 因他為做基督的工作不顧性命，幾乎至死，為要補足你們供應我不夠的地方。 \end{tabularx} \\ \\
[1ex]
\hline
\hline
\end{longtable}


\newpage

\section{第70屆~港九培靈會~曾立華~第六講}
\label{sec:733}
\textbf{獲得誠信正直質素的秘訣}
\newline
\newline
連結: \href{https://www.hkbibleconference.org/session-message/view/733}{\texttt{https://www.hkbibleconference.org/session-message/view/733}}
\newline
\newline
\hyperref[sec:730]{< < < PREV SERMON < < <}
~
\hyperref[sec:toc]{[返目錄]}
~
\hyperref[sec:735]{> > > NEXT SERMON > > >}
\newline
\newline
腓立比書 3:1-3:16
\newline
\begin{longtable}{cl}
\hline
\hline
章節 & 經文 (和合本修訂版)\\
\hline
3:1 & \begin{tabularx}{0.7\textwidth}{X} 末了，我的弟兄們，你們要靠主喜樂。我把這些話再寫給你們，對我並不困難，對你們卻是妥當的。 \end{tabularx} \\ \\ \relax
3:2 & \begin{tabularx}{0.7\textwidth}{X} 應當防備犬類，防備作惡的，防備妄自行割的。 \end{tabularx} \\ \\ \relax
3:3 & \begin{tabularx}{0.7\textwidth}{X} 因為真受割禮的，就是我們這藉著神的靈敬拜、以基督耶穌為誇耀、不依靠肉體的。 \end{tabularx} \\ \\ \relax
3:4 & \begin{tabularx}{0.7\textwidth}{X} 其實，我也可以靠肉體；若是別人以為他可以依靠肉體，我更可以。 \end{tabularx} \\ \\ \relax
3:5 & \begin{tabularx}{0.7\textwidth}{X} 我出生後第八天受割禮；我是以色列族、便雅憫支派的人，是希伯來人所生的希伯來人。就律法說，我是法利賽人； \end{tabularx} \\ \\ \relax
3:6 & \begin{tabularx}{0.7\textwidth}{X} 就熱心說，我是迫害教會的；就律法上的義說，我是無可指責的。 \end{tabularx} \\ \\ \relax
3:7 & \begin{tabularx}{0.7\textwidth}{X} 只是我先前以為對我是有益的，我現在因基督的緣故而當作是有損的。 \end{tabularx} \\ \\ \relax
3:8 & \begin{tabularx}{0.7\textwidth}{X} 不但如此，我已把萬事當作是有損的，因我以認識我主基督耶穌為至寶。我為他已經丟棄萬事，看作糞土，為要贏得基督， \end{tabularx} \\ \\ \relax
3:9 & \begin{tabularx}{0.7\textwidth}{X} 並且得以在他裡面，不是有自己因律法而得的義，而是有信基督的義，就是基於信，從神而來的義， \end{tabularx} \\ \\ \relax
3:10 & \begin{tabularx}{0.7\textwidth}{X} 使我認識基督，知道他復活的大能，並且知道和他一同受苦，效法他的死， \end{tabularx} \\ \\ \relax
3:11 & \begin{tabularx}{0.7\textwidth}{X} 或許我也得以從死人中復活。 \end{tabularx} \\ \\ \relax
3:12 & \begin{tabularx}{0.7\textwidth}{X} 這不是說我已經得著了，已經完全了；而是竭力追求，或許可以得著基督耶穌所要我得著的。 \end{tabularx} \\ \\ \relax
3:13 & \begin{tabularx}{0.7\textwidth}{X} 弟兄們，我不是以為自己已經得著了；我只有一件事，就是忘記背後，努力面前的， \end{tabularx} \\ \\ \relax
3:14 & \begin{tabularx}{0.7\textwidth}{X} 向著標竿直跑，要得神在基督耶穌裡從上面召我來得的獎賞。 \end{tabularx} \\ \\ \relax
3:15 & \begin{tabularx}{0.7\textwidth}{X} 所以，我們中間凡是成熟的人，總要存這樣的心；若在甚麼事上存別樣的心，神也會把這些事指示你們。 \end{tabularx} \\ \\ \relax
3:16 & \begin{tabularx}{0.7\textwidth}{X} 然而，我們達到甚麼地步，就當照這個地步行。 \end{tabularx} \\ \\
[1ex]
\hline
\hline
\end{longtable}


\newpage

\section{第70屆~港九培靈會~曾立華~第七講}
\label{sec:735}
\textbf{破壞誠信正直的因素}
\newline
\newline
連結: \href{https://www.hkbibleconference.org/session-message/view/735}{\texttt{https://www.hkbibleconference.org/session-message/view/735}}
\newline
\newline
\hyperref[sec:733]{< < < PREV SERMON < < <}
~
\hyperref[sec:toc]{[返目錄]}
~
\hyperref[sec:737]{> > > NEXT SERMON > > >}
\newline
\newline
腓立比書 3:17-4:1
\newline
\begin{longtable}{cl}
\hline
\hline
章節 & 經文 (和合本修訂版)\\
\hline
3:17 & \begin{tabularx}{0.7\textwidth}{X} 弟兄們，你們要一同效法我，也當留意看那些效法我們榜樣的人。 \end{tabularx} \\ \\ \relax
3:18 & \begin{tabularx}{0.7\textwidth}{X} 因為，我屢次告訴你們，現在又流淚告訴你們：許多人行事是基督十字架的仇敵。 \end{tabularx} \\ \\ \relax
3:19 & \begin{tabularx}{0.7\textwidth}{X} 他們的結局就是滅亡。他們的神明是自己的肚腹；他們以自己的羞辱為光榮，專以地上的事為念。 \end{tabularx} \\ \\ \relax
3:20 & \begin{tabularx}{0.7\textwidth}{X} 我們卻是天上的國民，並且等候救主，就是主耶穌基督從天上降臨。 \end{tabularx} \\ \\ \relax
3:21 & \begin{tabularx}{0.7\textwidth}{X} 他要按著那能使萬有歸服自己的大能，把我們這卑賤的身體改變形狀，和他自己榮耀的身體相似。 \end{tabularx} \\ \\ \relax
4:1 & \begin{tabularx}{0.7\textwidth}{X} 我所親愛、所想念的弟兄們，你們就是我的喜樂，我的冠冕。我親愛的，你們應當靠主站立得穩。 \end{tabularx} \\ \\
[1ex]
\hline
\hline
\end{longtable}


\newpage

\section{第70屆~港九培靈會~曾立華~第八講}
\label{sec:737}
\textbf{整全誠信正直的生活}
\newline
\newline
連結: \href{https://www.hkbibleconference.org/session-message/view/737}{\texttt{https://www.hkbibleconference.org/session-message/view/737}}
\newline
\newline
\hyperref[sec:735]{< < < PREV SERMON < < <}
~
\hyperref[sec:toc]{[返目錄]}
~
\hyperref[sec:856]{> > > NEXT SERMON > > >}
\newline
\newline
腓立比書 4:2-4:23
\newline
\begin{longtable}{cl}
\hline
\hline
章節 & 經文 (和合本修訂版)\\
\hline
4:2 & \begin{tabularx}{0.7\textwidth}{X} 我勸友阿蝶和循都基要在主裡同心。 \end{tabularx} \\ \\ \relax
4:3 & \begin{tabularx}{0.7\textwidth}{X} 我也求你這真實同負一軛的，要幫助這兩個女人，因為她們在福音上曾與我、革利免和我其餘的同工一同勞苦，他們的名字都在生命冊上。 \end{tabularx} \\ \\ \relax
4:4 & \begin{tabularx}{0.7\textwidth}{X} 你們要靠主常常喜樂。我再說，你們要喜樂。 \end{tabularx} \\ \\ \relax
4:5 & \begin{tabularx}{0.7\textwidth}{X} 要讓眾人知道你們謙讓的心。主已經近了。 \end{tabularx} \\ \\ \relax
4:6 & \begin{tabularx}{0.7\textwidth}{X} 應當一無掛慮，只要凡事藉著禱告、祈求和感謝，將你們所要的告訴神。 \end{tabularx} \\ \\ \relax
4:7 & \begin{tabularx}{0.7\textwidth}{X} 神所賜那超越人所能了解的平安，必在基督耶穌裡，保守你們的心懷意念。 \end{tabularx} \\ \\ \relax
4:8 & \begin{tabularx}{0.7\textwidth}{X} 末了，弟兄們，凡是真實的、凡是可敬的、凡是公義的、凡是清潔的、凡是可愛的、凡是有美名的，若有甚麼德行，若有甚麼稱讚，你們都要留意。 \end{tabularx} \\ \\ \relax
4:9 & \begin{tabularx}{0.7\textwidth}{X} 你們從我所學習的，所領受的，所聽見的，所看見的事，你們都要繼續去做，賜平安的神就必與你們同在。 \end{tabularx} \\ \\ \relax
4:10 & \begin{tabularx}{0.7\textwidth}{X} 我靠主大大喜樂，因為你們關懷我的心如今又表現了出來；其實你們一直都關懷我，只是沒有機會罷了。 \end{tabularx} \\ \\ \relax
4:11 & \begin{tabularx}{0.7\textwidth}{X} 我並不是因缺乏而說這話，因為我已經學會無論在甚麼景況都可以知足。 \end{tabularx} \\ \\ \relax
4:12 & \begin{tabularx}{0.7\textwidth}{X} 我知道怎樣處卑賤，也知道怎樣處豐富；或飽足或飢餓，或有餘或缺乏，任何事情，任何景況，我都得了祕訣。 \end{tabularx} \\ \\ \relax
4:13 & \begin{tabularx}{0.7\textwidth}{X} 我靠著那加給我力量的，凡事都能做。 \end{tabularx} \\ \\ \relax
4:14 & \begin{tabularx}{0.7\textwidth}{X} 然而，你們能和我分擔憂患是一件好事。 \end{tabularx} \\ \\ \relax
4:15 & \begin{tabularx}{0.7\textwidth}{X} 腓立比人哪，你們也知道我開始傳福音、離開馬其頓的時候，在收支的事上，除了你們以外，並沒有別的教會和我分擔。 \end{tabularx} \\ \\ \relax
4:16 & \begin{tabularx}{0.7\textwidth}{X} 就是我在帖撒羅尼迦，你們也一再差人來供給我的需用。 \end{tabularx} \\ \\ \relax
4:17 & \begin{tabularx}{0.7\textwidth}{X} 我並不求甚麼饋贈，只求你們的果子不斷增多，歸在你們的賬上。 \end{tabularx} \\ \\ \relax
4:18 & \begin{tabularx}{0.7\textwidth}{X} 但我已經如數收到，並且有餘；我已經充足，因我從以巴弗提受了你們的饋贈，當作極美的香氣，為神所接納、所喜悅的祭物。 \end{tabularx} \\ \\ \relax
4:19 & \begin{tabularx}{0.7\textwidth}{X} 我的神必照他榮耀的豐富，在基督耶穌裡，使你們一切所需用的都充足。 \end{tabularx} \\ \\ \relax
4:20 & \begin{tabularx}{0.7\textwidth}{X} 願榮耀歸給我們的父神，直到永永遠遠。阿們！ \end{tabularx} \\ \\ \relax
4:21 & \begin{tabularx}{0.7\textwidth}{X} 請問候在基督耶穌裡的各位聖徒。跟我一起的眾弟兄都問候你們。 \end{tabularx} \\ \\ \relax
4:22 & \begin{tabularx}{0.7\textwidth}{X} 眾聖徒都問候你們，特別在凱撒家裡的人問候你們。 \end{tabularx} \\ \\ \relax
4:23 & \begin{tabularx}{0.7\textwidth}{X} 願主耶穌基督的恩與你們的靈同在！ \end{tabularx} \\ \\
[1ex]
\hline
\hline
\end{longtable}


\newpage

\section{第70屆~港九培靈會~林三綱~第一講}
\label{sec:856}
\textbf{所羅門 (耶底底亞) 生平}
\newline
\newline
連結: \href{https://www.hkbibleconference.org/session-message/view/856}{\texttt{https://www.hkbibleconference.org/session-message/view/856}}
\newline
\newline
\hyperref[sec:737]{< < < PREV SERMON < < <}
~
\hyperref[sec:toc]{[返目錄]}
~
\hyperref[sec:857]{> > > NEXT SERMON > > >}
\newline
\newline


\newpage

\section{第70屆~港九培靈會~林三綱~第二講}
\label{sec:857}
\textbf{靈命經歷(1) 兩個心願(上)}
\newline
\newline
連結: \href{https://www.hkbibleconference.org/session-message/view/857}{\texttt{https://www.hkbibleconference.org/session-message/view/857}}
\newline
\newline
\hyperref[sec:856]{< < < PREV SERMON < < <}
~
\hyperref[sec:toc]{[返目錄]}
~
\hyperref[sec:858]{> > > NEXT SERMON > > >}
\newline
\newline


\newpage

\section{第70屆~港九培靈會~林三綱~第三講}
\label{sec:858}
\textbf{靈命經歷(1) 兩個心願(下)}
\newline
\newline
連結: \href{https://www.hkbibleconference.org/session-message/view/858}{\texttt{https://www.hkbibleconference.org/session-message/view/858}}
\newline
\newline
\hyperref[sec:857]{< < < PREV SERMON < < <}
~
\hyperref[sec:toc]{[返目錄]}
~
\hyperref[sec:859]{> > > NEXT SERMON > > >}
\newline
\newline


\newpage

\section{第70屆~港九培靈會~林三綱~第四講}
\label{sec:859}
\textbf{靈命經歷 (2) 躥山越嶺}
\newline
\newline
連結: \href{https://www.hkbibleconference.org/session-message/view/859}{\texttt{https://www.hkbibleconference.org/session-message/view/859}}
\newline
\newline
\hyperref[sec:858]{< < < PREV SERMON < < <}
~
\hyperref[sec:toc]{[返目錄]}
~
\hyperref[sec:860]{> > > NEXT SERMON > > >}
\newline
\newline


\newpage

\section{第70屆~港九培靈會~林三綱~第五講}
\label{sec:860}
\textbf{靈命經歷 (3) 夜間的華轎}
\newline
\newline
連結: \href{https://www.hkbibleconference.org/session-message/view/860}{\texttt{https://www.hkbibleconference.org/session-message/view/860}}
\newline
\newline
\hyperref[sec:859]{< < < PREV SERMON < < <}
~
\hyperref[sec:toc]{[返目錄]}
~
\hyperref[sec:861]{> > > NEXT SERMON > > >}
\newline
\newline


\newpage

\section{第70屆~港九培靈會~林三綱~第六講}
\label{sec:861}
\textbf{靈命經歷 (4) 患難與榮美}
\newline
\newline
連結: \href{https://www.hkbibleconference.org/session-message/view/861}{\texttt{https://www.hkbibleconference.org/session-message/view/861}}
\newline
\newline
\hyperref[sec:860]{< < < PREV SERMON < < <}
~
\hyperref[sec:toc]{[返目錄]}
~
\hyperref[sec:862]{> > > NEXT SERMON > > >}
\newline
\newline


\newpage

\section{第70屆~港九培靈會~林三綱~第七講}
\label{sec:862}
\textbf{治國建殿}
\newline
\newline
連結: \href{https://www.hkbibleconference.org/session-message/view/862}{\texttt{https://www.hkbibleconference.org/session-message/view/862}}
\newline
\newline
\hyperref[sec:861]{< < < PREV SERMON < < <}
~
\hyperref[sec:toc]{[返目錄]}
~
\hyperref[sec:863]{> > > NEXT SERMON > > >}
\newline
\newline


\newpage

\section{第70屆~港九培靈會~林三綱~第八講}
\label{sec:863}
\textbf{獻殿}
\newline
\newline
連結: \href{https://www.hkbibleconference.org/session-message/view/863}{\texttt{https://www.hkbibleconference.org/session-message/view/863}}
\newline
\newline
\hyperref[sec:862]{< < < PREV SERMON < < <}
~
\hyperref[sec:toc]{[返目錄]}
~
\hyperref[sec:864]{> > > NEXT SERMON > > >}
\newline
\newline


\newpage

\section{第70屆~港九培靈會~焦源濂~第一講}
\label{sec:864}
\textbf{一個與弟兄迥別的青年}
\newline
\newline
連結: \href{https://www.hkbibleconference.org/session-message/view/864}{\texttt{https://www.hkbibleconference.org/session-message/view/864}}
\newline
\newline
\hyperref[sec:863]{< < < PREV SERMON < < <}
~
\hyperref[sec:toc]{[返目錄]}
~
\hyperref[sec:865]{> > > NEXT SERMON > > >}
\newline
\newline
創世記 37:1-37:36
\newline
\begin{longtable}{cl}
\hline
\hline
章節 & 經文 (和合本修訂版)\\
\hline
37:1 & \begin{tabularx}{0.7\textwidth}{X} 雅各住在迦南地，就是他父親寄居的地。 \end{tabularx} \\ \\ \relax
37:2 & \begin{tabularx}{0.7\textwidth}{X} 這是雅各的事蹟。約瑟十七歲與他哥哥們一同牧羊。他是個少年，與他父親的妾辟拉和悉帕的兒子們常在一起。約瑟把他們的惡行報給父親。 \end{tabularx} \\ \\ \relax
37:3 & \begin{tabularx}{0.7\textwidth}{X} 以色列愛約瑟過於其他的兒子，因為約瑟是他年老生的；他給約瑟做了一件長袍。 \end{tabularx} \\ \\ \relax
37:4 & \begin{tabularx}{0.7\textwidth}{X} 哥哥們見父親愛約瑟過於他們，就恨約瑟，不與他說友善的話。 \end{tabularx} \\ \\ \relax
37:5 & \begin{tabularx}{0.7\textwidth}{X} 約瑟做了一個夢，告訴他哥哥們，他們就更加恨他。 \end{tabularx} \\ \\ \relax
37:6 & \begin{tabularx}{0.7\textwidth}{X} 約瑟對他們說：「請聽我做的這個夢： \end{tabularx} \\ \\ \relax
37:7 & \begin{tabularx}{0.7\textwidth}{X} 看哪，我們在田裡捆禾稼；看哪，我的捆起來站著；看哪，你們的捆圍著我的捆下拜。」 \end{tabularx} \\ \\ \relax
37:8 & \begin{tabularx}{0.7\textwidth}{X} 他的哥哥們對他說：「難道你真的要作我們的王嗎？難道你真的要統治我們嗎？」他們就因他的夢和他的話更加恨他。 \end{tabularx} \\ \\ \relax
37:9 & \begin{tabularx}{0.7\textwidth}{X} 後來他又做了另一個夢，告訴他哥哥們說：「看哪，我又做了一個夢；看哪，太陽、月亮和十一顆星都向我下拜。」 \end{tabularx} \\ \\ \relax
37:10 & \begin{tabularx}{0.7\textwidth}{X} 約瑟告訴他父親和哥哥們，他父親就責備他說：「你做的這是甚麼夢！難道我和你母親、你的兄弟真的要俯伏在地，來向你下拜嗎？」 \end{tabularx} \\ \\ \relax
37:11 & \begin{tabularx}{0.7\textwidth}{X} 他的哥哥們都嫉妒他，他父親卻把這事存在心裡。 \end{tabularx} \\ \\ \relax
37:12 & \begin{tabularx}{0.7\textwidth}{X} 約瑟的哥哥們到示劍去放他們父親的羊。 \end{tabularx} \\ \\ \relax
37:13 & \begin{tabularx}{0.7\textwidth}{X} 以色列對約瑟說：「你哥哥們不是在示劍放羊嗎？來，我派你到他們那裡去。」約瑟對他說：「我在這裡。」 \end{tabularx} \\ \\ \relax
37:14 & \begin{tabularx}{0.7\textwidth}{X} 以色列對他說：「你去看看你哥哥們是否平安，羊群是否平安，再回來告訴我。」於是他派約瑟出希伯崙谷，約瑟就往示劍去了。 \end{tabularx} \\ \\ \relax
37:15 & \begin{tabularx}{0.7\textwidth}{X} 有人遇見他，看哪，他在田野走迷了路。那人問他說：「你找甚麼？」 \end{tabularx} \\ \\ \relax
37:16 & \begin{tabularx}{0.7\textwidth}{X} 他說：「我找我的哥哥們，請告訴我，他們在哪裡放羊。」 \end{tabularx} \\ \\ \relax
37:17 & \begin{tabularx}{0.7\textwidth}{X} 那人說：「他們已經離開這裡走了，我聽見他們說：『我們往多坍去。』」約瑟就去追哥哥們，在多坍找到了他們。 \end{tabularx} \\ \\ \relax
37:18 & \begin{tabularx}{0.7\textwidth}{X} 他們遠遠看見他，趁他還沒有走近他們，就圖謀要殺死他。 \end{tabularx} \\ \\ \relax
37:19 & \begin{tabularx}{0.7\textwidth}{X} 他們彼此說：「看哪！那做夢的來了。 \end{tabularx} \\ \\ \relax
37:20 & \begin{tabularx}{0.7\textwidth}{X} 現在，來吧！我們把他殺了，丟在一個坑裡，就說有惡獸把他吃了。我們且看他的夢將來怎麼樣。」 \end{tabularx} \\ \\ \relax
37:21 & \begin{tabularx}{0.7\textwidth}{X} 呂便聽見了，要救約瑟脫離他們的手，說：「我們不可害他的性命。」 \end{tabularx} \\ \\ \relax
37:22 & \begin{tabularx}{0.7\textwidth}{X} 呂便又對他們說：「不可流他的血，可以把他丟在這曠野的坑裡，不可下手害他。」呂便要救他脫離他們的手，把他還給他父親。 \end{tabularx} \\ \\ \relax
37:23 & \begin{tabularx}{0.7\textwidth}{X} 約瑟到了他哥哥們那裡，他們就剝去他的外衣，就是他身上那件長袍。 \end{tabularx} \\ \\ \relax
37:24 & \begin{tabularx}{0.7\textwidth}{X} 他們抓住他，把他丟在坑裡。那坑是空的，裡頭沒有水。 \end{tabularx} \\ \\ \relax
37:25 & \begin{tabularx}{0.7\textwidth}{X} 他們坐下吃飯，舉目觀看，看哪，有一群以實瑪利人從基列來，用駱駝馱著香料、乳香、沒藥，要帶下埃及去。 \end{tabularx} \\ \\ \relax
37:26 & \begin{tabularx}{0.7\textwidth}{X} 猶大對他的兄弟們說：「我們殺我們的弟弟，遮掩他的血有甚麼好處呢？ \end{tabularx} \\ \\ \relax
37:27 & \begin{tabularx}{0.7\textwidth}{X} 來，我們把他賣給以實瑪利人，不要下手害他，因為他是我們的弟弟，我們的骨肉。」他的兄弟們就聽從了他。 \end{tabularx} \\ \\ \relax
37:28 & \begin{tabularx}{0.7\textwidth}{X} 那時，有些米甸的商人從那裡經過，就把約瑟從坑裡拉上來。他們以二十塊銀子把約瑟賣給以實瑪利人，他們就把約瑟帶到埃及去了。 \end{tabularx} \\ \\ \relax
37:29 & \begin{tabularx}{0.7\textwidth}{X} 呂便回到坑旁，看哪，約瑟不在坑裡，就撕裂自己的衣服， \end{tabularx} \\ \\ \relax
37:30 & \begin{tabularx}{0.7\textwidth}{X} 回到他兄弟們那裡，說：「孩子不在了。我往哪裡去才好呢？」 \end{tabularx} \\ \\ \relax
37:31 & \begin{tabularx}{0.7\textwidth}{X} 於是，他們宰了一隻公山羊，拿了約瑟的那件外衣染上了血， \end{tabularx} \\ \\ \relax
37:32 & \begin{tabularx}{0.7\textwidth}{X} 派人把長袍送到他們的父親那裡，說：「我們發現這個， 請認一認，是不是你兒子的外衣？」 \end{tabularx} \\ \\ \relax
37:33 & \begin{tabularx}{0.7\textwidth}{X} 他認出來，就說：「這是我兒子的外衣，惡獸把他吃了，約瑟一定被撕碎了！」 \end{tabularx} \\ \\ \relax
37:34 & \begin{tabularx}{0.7\textwidth}{X} 雅各就撕裂衣服，腰間圍上麻布，為他兒子哀傷了多日。 \end{tabularx} \\ \\ \relax
37:35 & \begin{tabularx}{0.7\textwidth}{X} 他的兒女都起來安慰他，他卻不肯受安慰，說：「我必哀傷著下陰間，到我兒子那裡。」約瑟的父親就為他哀哭。 \end{tabularx} \\ \\ \relax
37:36 & \begin{tabularx}{0.7\textwidth}{X} 米甸人把約瑟賣到埃及，給法老的官員，就是護衛長波提乏。 \end{tabularx} \\ \\
[1ex]
\hline
\hline
\end{longtable}


\newpage

\section{第70屆~港九培靈會~焦源濂~第二講}
\label{sec:865}
\textbf{一個「作夢」的聖徒}
\newline
\newline
連結: \href{https://www.hkbibleconference.org/session-message/view/865}{\texttt{https://www.hkbibleconference.org/session-message/view/865}}
\newline
\newline
\hyperref[sec:864]{< < < PREV SERMON < < <}
~
\hyperref[sec:toc]{[返目錄]}
~
\hyperref[sec:866]{> > > NEXT SERMON > > >}
\newline
\newline
創世記 37:1-37:36
\newline
\begin{longtable}{cl}
\hline
\hline
章節 & 經文 (和合本修訂版)\\
\hline
37:1 & \begin{tabularx}{0.7\textwidth}{X} 雅各住在迦南地，就是他父親寄居的地。 \end{tabularx} \\ \\ \relax
37:2 & \begin{tabularx}{0.7\textwidth}{X} 這是雅各的事蹟。約瑟十七歲與他哥哥們一同牧羊。他是個少年，與他父親的妾辟拉和悉帕的兒子們常在一起。約瑟把他們的惡行報給父親。 \end{tabularx} \\ \\ \relax
37:3 & \begin{tabularx}{0.7\textwidth}{X} 以色列愛約瑟過於其他的兒子，因為約瑟是他年老生的；他給約瑟做了一件長袍。 \end{tabularx} \\ \\ \relax
37:4 & \begin{tabularx}{0.7\textwidth}{X} 哥哥們見父親愛約瑟過於他們，就恨約瑟，不與他說友善的話。 \end{tabularx} \\ \\ \relax
37:5 & \begin{tabularx}{0.7\textwidth}{X} 約瑟做了一個夢，告訴他哥哥們，他們就更加恨他。 \end{tabularx} \\ \\ \relax
37:6 & \begin{tabularx}{0.7\textwidth}{X} 約瑟對他們說：「請聽我做的這個夢： \end{tabularx} \\ \\ \relax
37:7 & \begin{tabularx}{0.7\textwidth}{X} 看哪，我們在田裡捆禾稼；看哪，我的捆起來站著；看哪，你們的捆圍著我的捆下拜。」 \end{tabularx} \\ \\ \relax
37:8 & \begin{tabularx}{0.7\textwidth}{X} 他的哥哥們對他說：「難道你真的要作我們的王嗎？難道你真的要統治我們嗎？」他們就因他的夢和他的話更加恨他。 \end{tabularx} \\ \\ \relax
37:9 & \begin{tabularx}{0.7\textwidth}{X} 後來他又做了另一個夢，告訴他哥哥們說：「看哪，我又做了一個夢；看哪，太陽、月亮和十一顆星都向我下拜。」 \end{tabularx} \\ \\ \relax
37:10 & \begin{tabularx}{0.7\textwidth}{X} 約瑟告訴他父親和哥哥們，他父親就責備他說：「你做的這是甚麼夢！難道我和你母親、你的兄弟真的要俯伏在地，來向你下拜嗎？」 \end{tabularx} \\ \\ \relax
37:11 & \begin{tabularx}{0.7\textwidth}{X} 他的哥哥們都嫉妒他，他父親卻把這事存在心裡。 \end{tabularx} \\ \\ \relax
37:12 & \begin{tabularx}{0.7\textwidth}{X} 約瑟的哥哥們到示劍去放他們父親的羊。 \end{tabularx} \\ \\ \relax
37:13 & \begin{tabularx}{0.7\textwidth}{X} 以色列對約瑟說：「你哥哥們不是在示劍放羊嗎？來，我派你到他們那裡去。」約瑟對他說：「我在這裡。」 \end{tabularx} \\ \\ \relax
37:14 & \begin{tabularx}{0.7\textwidth}{X} 以色列對他說：「你去看看你哥哥們是否平安，羊群是否平安，再回來告訴我。」於是他派約瑟出希伯崙谷，約瑟就往示劍去了。 \end{tabularx} \\ \\ \relax
37:15 & \begin{tabularx}{0.7\textwidth}{X} 有人遇見他，看哪，他在田野走迷了路。那人問他說：「你找甚麼？」 \end{tabularx} \\ \\ \relax
37:16 & \begin{tabularx}{0.7\textwidth}{X} 他說：「我找我的哥哥們，請告訴我，他們在哪裡放羊。」 \end{tabularx} \\ \\ \relax
37:17 & \begin{tabularx}{0.7\textwidth}{X} 那人說：「他們已經離開這裡走了，我聽見他們說：『我們往多坍去。』」約瑟就去追哥哥們，在多坍找到了他們。 \end{tabularx} \\ \\ \relax
37:18 & \begin{tabularx}{0.7\textwidth}{X} 他們遠遠看見他，趁他還沒有走近他們，就圖謀要殺死他。 \end{tabularx} \\ \\ \relax
37:19 & \begin{tabularx}{0.7\textwidth}{X} 他們彼此說：「看哪！那做夢的來了。 \end{tabularx} \\ \\ \relax
37:20 & \begin{tabularx}{0.7\textwidth}{X} 現在，來吧！我們把他殺了，丟在一個坑裡，就說有惡獸把他吃了。我們且看他的夢將來怎麼樣。」 \end{tabularx} \\ \\ \relax
37:21 & \begin{tabularx}{0.7\textwidth}{X} 呂便聽見了，要救約瑟脫離他們的手，說：「我們不可害他的性命。」 \end{tabularx} \\ \\ \relax
37:22 & \begin{tabularx}{0.7\textwidth}{X} 呂便又對他們說：「不可流他的血，可以把他丟在這曠野的坑裡，不可下手害他。」呂便要救他脫離他們的手，把他還給他父親。 \end{tabularx} \\ \\ \relax
37:23 & \begin{tabularx}{0.7\textwidth}{X} 約瑟到了他哥哥們那裡，他們就剝去他的外衣，就是他身上那件長袍。 \end{tabularx} \\ \\ \relax
37:24 & \begin{tabularx}{0.7\textwidth}{X} 他們抓住他，把他丟在坑裡。那坑是空的，裡頭沒有水。 \end{tabularx} \\ \\ \relax
37:25 & \begin{tabularx}{0.7\textwidth}{X} 他們坐下吃飯，舉目觀看，看哪，有一群以實瑪利人從基列來，用駱駝馱著香料、乳香、沒藥，要帶下埃及去。 \end{tabularx} \\ \\ \relax
37:26 & \begin{tabularx}{0.7\textwidth}{X} 猶大對他的兄弟們說：「我們殺我們的弟弟，遮掩他的血有甚麼好處呢？ \end{tabularx} \\ \\ \relax
37:27 & \begin{tabularx}{0.7\textwidth}{X} 來，我們把他賣給以實瑪利人，不要下手害他，因為他是我們的弟弟，我們的骨肉。」他的兄弟們就聽從了他。 \end{tabularx} \\ \\ \relax
37:28 & \begin{tabularx}{0.7\textwidth}{X} 那時，有些米甸的商人從那裡經過，就把約瑟從坑裡拉上來。他們以二十塊銀子把約瑟賣給以實瑪利人，他們就把約瑟帶到埃及去了。 \end{tabularx} \\ \\ \relax
37:29 & \begin{tabularx}{0.7\textwidth}{X} 呂便回到坑旁，看哪，約瑟不在坑裡，就撕裂自己的衣服， \end{tabularx} \\ \\ \relax
37:30 & \begin{tabularx}{0.7\textwidth}{X} 回到他兄弟們那裡，說：「孩子不在了。我往哪裡去才好呢？」 \end{tabularx} \\ \\ \relax
37:31 & \begin{tabularx}{0.7\textwidth}{X} 於是，他們宰了一隻公山羊，拿了約瑟的那件外衣染上了血， \end{tabularx} \\ \\ \relax
37:32 & \begin{tabularx}{0.7\textwidth}{X} 派人把長袍送到他們的父親那裡，說：「我們發現這個， 請認一認，是不是你兒子的外衣？」 \end{tabularx} \\ \\ \relax
37:33 & \begin{tabularx}{0.7\textwidth}{X} 他認出來，就說：「這是我兒子的外衣，惡獸把他吃了，約瑟一定被撕碎了！」 \end{tabularx} \\ \\ \relax
37:34 & \begin{tabularx}{0.7\textwidth}{X} 雅各就撕裂衣服，腰間圍上麻布，為他兒子哀傷了多日。 \end{tabularx} \\ \\ \relax
37:35 & \begin{tabularx}{0.7\textwidth}{X} 他的兒女都起來安慰他，他卻不肯受安慰，說：「我必哀傷著下陰間，到我兒子那裡。」約瑟的父親就為他哀哭。 \end{tabularx} \\ \\ \relax
37:36 & \begin{tabularx}{0.7\textwidth}{X} 米甸人把約瑟賣到埃及，給法老的官員，就是護衛長波提乏。 \end{tabularx} \\ \\
[1ex]
\hline
\hline
\end{longtable}


\newpage

\section{第70屆~港九培靈會~焦源濂~第三講}
\label{sec:866}
\textbf{人間辛酸備嘗}
\newline
\newline
連結: \href{https://www.hkbibleconference.org/session-message/view/866}{\texttt{https://www.hkbibleconference.org/session-message/view/866}}
\newline
\newline
\hyperref[sec:865]{< < < PREV SERMON < < <}
~
\hyperref[sec:toc]{[返目錄]}
~
\hyperref[sec:867]{> > > NEXT SERMON > > >}
\newline
\newline
創世記 37:1-37:36
\newline
\begin{longtable}{cl}
\hline
\hline
章節 & 經文 (和合本修訂版)\\
\hline
37:1 & \begin{tabularx}{0.7\textwidth}{X} 雅各住在迦南地，就是他父親寄居的地。 \end{tabularx} \\ \\ \relax
37:2 & \begin{tabularx}{0.7\textwidth}{X} 這是雅各的事蹟。約瑟十七歲與他哥哥們一同牧羊。他是個少年，與他父親的妾辟拉和悉帕的兒子們常在一起。約瑟把他們的惡行報給父親。 \end{tabularx} \\ \\ \relax
37:3 & \begin{tabularx}{0.7\textwidth}{X} 以色列愛約瑟過於其他的兒子，因為約瑟是他年老生的；他給約瑟做了一件長袍。 \end{tabularx} \\ \\ \relax
37:4 & \begin{tabularx}{0.7\textwidth}{X} 哥哥們見父親愛約瑟過於他們，就恨約瑟，不與他說友善的話。 \end{tabularx} \\ \\ \relax
37:5 & \begin{tabularx}{0.7\textwidth}{X} 約瑟做了一個夢，告訴他哥哥們，他們就更加恨他。 \end{tabularx} \\ \\ \relax
37:6 & \begin{tabularx}{0.7\textwidth}{X} 約瑟對他們說：「請聽我做的這個夢： \end{tabularx} \\ \\ \relax
37:7 & \begin{tabularx}{0.7\textwidth}{X} 看哪，我們在田裡捆禾稼；看哪，我的捆起來站著；看哪，你們的捆圍著我的捆下拜。」 \end{tabularx} \\ \\ \relax
37:8 & \begin{tabularx}{0.7\textwidth}{X} 他的哥哥們對他說：「難道你真的要作我們的王嗎？難道你真的要統治我們嗎？」他們就因他的夢和他的話更加恨他。 \end{tabularx} \\ \\ \relax
37:9 & \begin{tabularx}{0.7\textwidth}{X} 後來他又做了另一個夢，告訴他哥哥們說：「看哪，我又做了一個夢；看哪，太陽、月亮和十一顆星都向我下拜。」 \end{tabularx} \\ \\ \relax
37:10 & \begin{tabularx}{0.7\textwidth}{X} 約瑟告訴他父親和哥哥們，他父親就責備他說：「你做的這是甚麼夢！難道我和你母親、你的兄弟真的要俯伏在地，來向你下拜嗎？」 \end{tabularx} \\ \\ \relax
37:11 & \begin{tabularx}{0.7\textwidth}{X} 他的哥哥們都嫉妒他，他父親卻把這事存在心裡。 \end{tabularx} \\ \\ \relax
37:12 & \begin{tabularx}{0.7\textwidth}{X} 約瑟的哥哥們到示劍去放他們父親的羊。 \end{tabularx} \\ \\ \relax
37:13 & \begin{tabularx}{0.7\textwidth}{X} 以色列對約瑟說：「你哥哥們不是在示劍放羊嗎？來，我派你到他們那裡去。」約瑟對他說：「我在這裡。」 \end{tabularx} \\ \\ \relax
37:14 & \begin{tabularx}{0.7\textwidth}{X} 以色列對他說：「你去看看你哥哥們是否平安，羊群是否平安，再回來告訴我。」於是他派約瑟出希伯崙谷，約瑟就往示劍去了。 \end{tabularx} \\ \\ \relax
37:15 & \begin{tabularx}{0.7\textwidth}{X} 有人遇見他，看哪，他在田野走迷了路。那人問他說：「你找甚麼？」 \end{tabularx} \\ \\ \relax
37:16 & \begin{tabularx}{0.7\textwidth}{X} 他說：「我找我的哥哥們，請告訴我，他們在哪裡放羊。」 \end{tabularx} \\ \\ \relax
37:17 & \begin{tabularx}{0.7\textwidth}{X} 那人說：「他們已經離開這裡走了，我聽見他們說：『我們往多坍去。』」約瑟就去追哥哥們，在多坍找到了他們。 \end{tabularx} \\ \\ \relax
37:18 & \begin{tabularx}{0.7\textwidth}{X} 他們遠遠看見他，趁他還沒有走近他們，就圖謀要殺死他。 \end{tabularx} \\ \\ \relax
37:19 & \begin{tabularx}{0.7\textwidth}{X} 他們彼此說：「看哪！那做夢的來了。 \end{tabularx} \\ \\ \relax
37:20 & \begin{tabularx}{0.7\textwidth}{X} 現在，來吧！我們把他殺了，丟在一個坑裡，就說有惡獸把他吃了。我們且看他的夢將來怎麼樣。」 \end{tabularx} \\ \\ \relax
37:21 & \begin{tabularx}{0.7\textwidth}{X} 呂便聽見了，要救約瑟脫離他們的手，說：「我們不可害他的性命。」 \end{tabularx} \\ \\ \relax
37:22 & \begin{tabularx}{0.7\textwidth}{X} 呂便又對他們說：「不可流他的血，可以把他丟在這曠野的坑裡，不可下手害他。」呂便要救他脫離他們的手，把他還給他父親。 \end{tabularx} \\ \\ \relax
37:23 & \begin{tabularx}{0.7\textwidth}{X} 約瑟到了他哥哥們那裡，他們就剝去他的外衣，就是他身上那件長袍。 \end{tabularx} \\ \\ \relax
37:24 & \begin{tabularx}{0.7\textwidth}{X} 他們抓住他，把他丟在坑裡。那坑是空的，裡頭沒有水。 \end{tabularx} \\ \\ \relax
37:25 & \begin{tabularx}{0.7\textwidth}{X} 他們坐下吃飯，舉目觀看，看哪，有一群以實瑪利人從基列來，用駱駝馱著香料、乳香、沒藥，要帶下埃及去。 \end{tabularx} \\ \\ \relax
37:26 & \begin{tabularx}{0.7\textwidth}{X} 猶大對他的兄弟們說：「我們殺我們的弟弟，遮掩他的血有甚麼好處呢？ \end{tabularx} \\ \\ \relax
37:27 & \begin{tabularx}{0.7\textwidth}{X} 來，我們把他賣給以實瑪利人，不要下手害他，因為他是我們的弟弟，我們的骨肉。」他的兄弟們就聽從了他。 \end{tabularx} \\ \\ \relax
37:28 & \begin{tabularx}{0.7\textwidth}{X} 那時，有些米甸的商人從那裡經過，就把約瑟從坑裡拉上來。他們以二十塊銀子把約瑟賣給以實瑪利人，他們就把約瑟帶到埃及去了。 \end{tabularx} \\ \\ \relax
37:29 & \begin{tabularx}{0.7\textwidth}{X} 呂便回到坑旁，看哪，約瑟不在坑裡，就撕裂自己的衣服， \end{tabularx} \\ \\ \relax
37:30 & \begin{tabularx}{0.7\textwidth}{X} 回到他兄弟們那裡，說：「孩子不在了。我往哪裡去才好呢？」 \end{tabularx} \\ \\ \relax
37:31 & \begin{tabularx}{0.7\textwidth}{X} 於是，他們宰了一隻公山羊，拿了約瑟的那件外衣染上了血， \end{tabularx} \\ \\ \relax
37:32 & \begin{tabularx}{0.7\textwidth}{X} 派人把長袍送到他們的父親那裡，說：「我們發現這個， 請認一認，是不是你兒子的外衣？」 \end{tabularx} \\ \\ \relax
37:33 & \begin{tabularx}{0.7\textwidth}{X} 他認出來，就說：「這是我兒子的外衣，惡獸把他吃了，約瑟一定被撕碎了！」 \end{tabularx} \\ \\ \relax
37:34 & \begin{tabularx}{0.7\textwidth}{X} 雅各就撕裂衣服，腰間圍上麻布，為他兒子哀傷了多日。 \end{tabularx} \\ \\ \relax
37:35 & \begin{tabularx}{0.7\textwidth}{X} 他的兒女都起來安慰他，他卻不肯受安慰，說：「我必哀傷著下陰間，到我兒子那裡。」約瑟的父親就為他哀哭。 \end{tabularx} \\ \\ \relax
37:36 & \begin{tabularx}{0.7\textwidth}{X} 米甸人把約瑟賣到埃及，給法老的官員，就是護衛長波提乏。 \end{tabularx} \\ \\
[1ex]
\hline
\hline
\end{longtable}


\newpage

\section{第70屆~港九培靈會~焦源濂~第四講}
\label{sec:867}
\textbf{受大的打擊,更深的扎根}
\newline
\newline
連結: \href{https://www.hkbibleconference.org/session-message/view/867}{\texttt{https://www.hkbibleconference.org/session-message/view/867}}
\newline
\newline
\hyperref[sec:866]{< < < PREV SERMON < < <}
~
\hyperref[sec:toc]{[返目錄]}
~
\hyperref[sec:868]{> > > NEXT SERMON > > >}
\newline
\newline
創世記 39:1-40:23
\newline
\begin{longtable}{cl}
\hline
\hline
章節 & 經文 (和合本修訂版)\\
\hline
39:1 & \begin{tabularx}{0.7\textwidth}{X} 約瑟被帶下埃及去。有一個埃及人波提乏，是法老的官員，是護衛長，他從那些帶約瑟下來的以實瑪利人手中把約瑟買了去。 \end{tabularx} \\ \\ \relax
39:2 & \begin{tabularx}{0.7\textwidth}{X} 約瑟在他埃及主人的家中，耶和華與他同在，他是一個通達的人。 \end{tabularx} \\ \\ \relax
39:3 & \begin{tabularx}{0.7\textwidth}{X} 他主人見耶和華與他同在，又見耶和華使他手裡所辦的事都順利， \end{tabularx} \\ \\ \relax
39:4 & \begin{tabularx}{0.7\textwidth}{X} 約瑟就在主人眼前蒙恩，伺候他主人，主人派他管理家務，把一切所有的都交在他手裡。 \end{tabularx} \\ \\ \relax
39:5 & \begin{tabularx}{0.7\textwidth}{X} 自從主人派約瑟管理家務和他一切所有的，耶和華就因約瑟的緣故賜福給那埃及人的家；凡家裡和田間一切所有的，都蒙耶和華賜福。 \end{tabularx} \\ \\ \relax
39:6 & \begin{tabularx}{0.7\textwidth}{X} 波提乏把他一切所有的都交在約瑟手中，除了自己所吃的食物，其他的事一概不知。約瑟英俊健美。 \end{tabularx} \\ \\ \relax
39:7 & \begin{tabularx}{0.7\textwidth}{X} 這些事以後，約瑟主人的妻子以目送情給約瑟，說：「你與我同寢吧！」 \end{tabularx} \\ \\ \relax
39:8 & \begin{tabularx}{0.7\textwidth}{X} 約瑟拒絕，對他主人的妻子說：「看哪，一切家務我主人一概不知，他把所有的都交在我手裡。 \end{tabularx} \\ \\ \relax
39:9 & \begin{tabularx}{0.7\textwidth}{X} 在這家裡沒有人比我更大，除你以外，他也沒有留下一樣不交給我，因為你是他的妻子。我怎能行這麼大的惡，得罪神呢？」 \end{tabularx} \\ \\ \relax
39:10 & \begin{tabularx}{0.7\textwidth}{X} 她天天這樣對約瑟說，約瑟卻不聽從她，不與她同寢，也不和她在一起。 \end{tabularx} \\ \\ \relax
39:11 & \begin{tabularx}{0.7\textwidth}{X} 有一天，約瑟進屋裡去辦事，家裡沒有一個人在那屋子裡， \end{tabularx} \\ \\ \relax
39:12 & \begin{tabularx}{0.7\textwidth}{X} 婦人就拉住他的衣服，說：「你與我同寢吧！」約瑟把衣服留在她手裡，逃出外面去了。 \end{tabularx} \\ \\ \relax
39:13 & \begin{tabularx}{0.7\textwidth}{X} 婦人看見約瑟把衣服留在她手裡逃到外面， \end{tabularx} \\ \\ \relax
39:14 & \begin{tabularx}{0.7\textwidth}{X} 就叫了家裡的人來，對他們說：「看，他帶了一個希伯來人到我們這裡戲弄我們。他到我這裡來，要與我同寢，我就大聲喊叫。 \end{tabularx} \\ \\ \relax
39:15 & \begin{tabularx}{0.7\textwidth}{X} 他聽見我放聲大喊，就把他的衣服留在我這裡，逃出外面去了。」 \end{tabularx} \\ \\ \relax
39:16 & \begin{tabularx}{0.7\textwidth}{X} 婦人把約瑟的衣服放在身邊，直到他主人回家， \end{tabularx} \\ \\ \relax
39:17 & \begin{tabularx}{0.7\textwidth}{X} 就用這樣的話對他說：「你帶到我們這裡來的那希伯來僕人進來要調戲我， \end{tabularx} \\ \\ \relax
39:18 & \begin{tabularx}{0.7\textwidth}{X} 我放聲大喊，他就把衣服留在我身邊，逃到外面。」 \end{tabularx} \\ \\ \relax
39:19 & \begin{tabularx}{0.7\textwidth}{X} 主人聽見他妻子對他說的話，說：「你的僕人就是這樣對待我」，就非常生氣。 \end{tabularx} \\ \\ \relax
39:20 & \begin{tabularx}{0.7\textwidth}{X} 約瑟的主人把他抓起來，關在監獄裡，就是王的囚犯被關的地方。於是約瑟在那裡坐牢。 \end{tabularx} \\ \\ \relax
39:21 & \begin{tabularx}{0.7\textwidth}{X} 但耶和華與約瑟同在，向他施恩，使他在監獄長的眼前蒙恩。 \end{tabularx} \\ \\ \relax
39:22 & \begin{tabularx}{0.7\textwidth}{X} 監獄長就把監獄裡所有的囚犯都交在約瑟手下；在那裡的一切事都由他處理。 \end{tabularx} \\ \\ \relax
39:23 & \begin{tabularx}{0.7\textwidth}{X} 任何交在約瑟手中的事，監獄長一概不察，因為耶和華與約瑟同在，耶和華使他所做的都順利。 \end{tabularx} \\ \\ \relax
40:1 & \begin{tabularx}{0.7\textwidth}{X} 這些事以後，埃及王的司酒長和司膳長得罪了他們的主埃及王。 \end{tabularx} \\ \\ \relax
40:2 & \begin{tabularx}{0.7\textwidth}{X} 法老就對司酒長和司膳長兩個官員發怒， \end{tabularx} \\ \\ \relax
40:3 & \begin{tabularx}{0.7\textwidth}{X} 把他們關在護衛長府內的監獄裡，就是約瑟被囚的地方。 \end{tabularx} \\ \\ \relax
40:4 & \begin{tabularx}{0.7\textwidth}{X} 護衛長把他們交給約瑟，約瑟就伺候他們。他們被關了一段日子。 \end{tabularx} \\ \\ \relax
40:5 & \begin{tabularx}{0.7\textwidth}{X} 關在監獄裡的這兩個人，就是埃及王的司酒長和司膳長，在同一個晚上各自做了一個夢，每個夢都有自己的解釋。 \end{tabularx} \\ \\ \relax
40:6 & \begin{tabularx}{0.7\textwidth}{X} 到了早晨，約瑟來到他們那裡看他們，看哪，他們很憂愁。 \end{tabularx} \\ \\ \relax
40:7 & \begin{tabularx}{0.7\textwidth}{X} 他就問一同關在他主人府內法老的官員，說：「你們今日為甚麼面帶愁容呢？」 \end{tabularx} \\ \\ \relax
40:8 & \begin{tabularx}{0.7\textwidth}{X} 他們對他說：「我們各自做了一個夢，卻沒有人能講解。」約瑟對他們說：「解夢不是出於神嗎？請你們把夢告訴我。」 \end{tabularx} \\ \\ \relax
40:9 & \begin{tabularx}{0.7\textwidth}{X} 司酒長就把夢告訴約瑟，對他說：「在我的夢中，看哪，有一棵葡萄樹在我面前， \end{tabularx} \\ \\ \relax
40:10 & \begin{tabularx}{0.7\textwidth}{X} 樹上有三根枝子。枝子發了芽，開了花，結出串串成熟的葡萄。 \end{tabularx} \\ \\ \relax
40:11 & \begin{tabularx}{0.7\textwidth}{X} 法老的杯在我手中，我就拿葡萄擠在法老的杯裡，把杯遞到他手中。」 \end{tabularx} \\ \\ \relax
40:12 & \begin{tabularx}{0.7\textwidth}{X} 約瑟對他說：「夢的解釋是這樣：三根枝子就是三天； \end{tabularx} \\ \\ \relax
40:13 & \begin{tabularx}{0.7\textwidth}{X} 三天之內，法老要讓你抬起頭來，叫你官復原職。你仍要遞杯在法老的手中，像先前作他的司酒長一樣。 \end{tabularx} \\ \\ \relax
40:14 & \begin{tabularx}{0.7\textwidth}{X} 但你得福的時候，請你記得我，向我施慈愛，在法老面前提起我，救我出這監牢。 \end{tabularx} \\ \\ \relax
40:15 & \begin{tabularx}{0.7\textwidth}{X} 我實在是從希伯來人之地被拐來的，我在這裡也沒有做過甚麼，好叫他們把我關在牢裡。」 \end{tabularx} \\ \\ \relax
40:16 & \begin{tabularx}{0.7\textwidth}{X} 司膳長見夢解得好，就對約瑟說：「在我夢中，看哪，我頭上頂著三個裝餅的籃子； \end{tabularx} \\ \\ \relax
40:17 & \begin{tabularx}{0.7\textwidth}{X} 最上面的籃子裡有為法老烤的各樣食物，有飛鳥來吃我頭上籃子裡的食物。」 \end{tabularx} \\ \\ \relax
40:18 & \begin{tabularx}{0.7\textwidth}{X} 約瑟說：「夢的解釋是這樣：三個籃子就是三天； \end{tabularx} \\ \\ \relax
40:19 & \begin{tabularx}{0.7\textwidth}{X} 三天之內，法老要讓你抬起頭來，身首異處，把你掛在木架上，必有飛鳥來吃你身上的肉。」 \end{tabularx} \\ \\ \relax
40:20 & \begin{tabularx}{0.7\textwidth}{X} 到了第三天，正是法老的生日，他為眾臣僕擺設宴席，使司酒長和司膳長從眾臣僕中抬起頭來， \end{tabularx} \\ \\ \relax
40:21 & \begin{tabularx}{0.7\textwidth}{X} 讓司酒長官復原職，仍舊遞杯在法老手中， \end{tabularx} \\ \\ \relax
40:22 & \begin{tabularx}{0.7\textwidth}{X} 卻把司膳長掛起來，正如約瑟向他們所講解的。 \end{tabularx} \\ \\ \relax
40:23 & \begin{tabularx}{0.7\textwidth}{X} 然而，司酒長不記得約瑟，竟忘了他。 \end{tabularx} \\ \\
[1ex]
\hline
\hline
\end{longtable}


\newpage

\section{第70屆~港九培靈會~焦源濂~第五講}
\label{sec:868}
\textbf{突然青雲直上}
\newline
\newline
連結: \href{https://www.hkbibleconference.org/session-message/view/868}{\texttt{https://www.hkbibleconference.org/session-message/view/868}}
\newline
\newline
\hyperref[sec:867]{< < < PREV SERMON < < <}
~
\hyperref[sec:toc]{[返目錄]}
~
\hyperref[sec:869]{> > > NEXT SERMON > > >}
\newline
\newline
創世記 41:1-41:57
\newline
\begin{longtable}{cl}
\hline
\hline
章節 & 經文 (和合本修訂版)\\
\hline
41:1 & \begin{tabularx}{0.7\textwidth}{X} 過了兩年，法老做夢，看哪，自己站在尼羅河邊， \end{tabularx} \\ \\ \relax
41:2 & \begin{tabularx}{0.7\textwidth}{X} 看哪，有七頭母牛從尼羅河裡上來，長相俊美，肌肉肥壯，在蘆葦中吃草。 \end{tabularx} \\ \\ \relax
41:3 & \begin{tabularx}{0.7\textwidth}{X} 看哪，隨後又有七頭母牛從尼羅河裡上來，長相醜陋，肌肉乾瘦，與那七頭母牛一同站在河邊。 \end{tabularx} \\ \\ \relax
41:4 & \begin{tabularx}{0.7\textwidth}{X} 這長相醜陋，肌肉乾瘦的七頭母牛吃了那長相俊美又肥壯的七頭母牛。法老就醒了。 \end{tabularx} \\ \\ \relax
41:5 & \begin{tabularx}{0.7\textwidth}{X} 他又睡著，第二次做夢，看哪，一株麥桿長了七個穗子，又肥大又佳美， \end{tabularx} \\ \\ \relax
41:6 & \begin{tabularx}{0.7\textwidth}{X} 看哪，隨後又長出七個穗子，又細弱又被東風吹焦了。 \end{tabularx} \\ \\ \relax
41:7 & \begin{tabularx}{0.7\textwidth}{X} 這細弱的穗子吞了那七個又肥大又飽滿的穗子。法老醒了，看哪，是個夢。 \end{tabularx} \\ \\ \relax
41:8 & \begin{tabularx}{0.7\textwidth}{X} 到了早晨，法老心裡不安，就派人把埃及所有的術士和智慧人都召來。法老把所做的夢告訴他們，但是沒有人能為法老解夢。 \end{tabularx} \\ \\ \relax
41:9 & \begin{tabularx}{0.7\textwidth}{X} 那時司酒長對法老說：「我今日想起我的罪來。 \end{tabularx} \\ \\ \relax
41:10 & \begin{tabularx}{0.7\textwidth}{X} 從前法老對臣僕發怒，把我和司膳長關在護衛長府內的監牢裡。 \end{tabularx} \\ \\ \relax
41:11 & \begin{tabularx}{0.7\textwidth}{X} 我們兩人在同一晚上各做一夢，每個夢都有各自的解釋。 \end{tabularx} \\ \\ \relax
41:12 & \begin{tabularx}{0.7\textwidth}{X} 同我們在一起有一個希伯來的年輕人，是護衛長的僕人。我們告訴他，他就為我們解夢，照著各人的夢講解。 \end{tabularx} \\ \\ \relax
41:13 & \begin{tabularx}{0.7\textwidth}{X} 後來事情正如他給我們講解的實現了，我官復原職，司膳長被掛起來了。」 \end{tabularx} \\ \\ \relax
41:14 & \begin{tabularx}{0.7\textwidth}{X} 於是法老派人去召約瑟，他們就急忙把他從牢裡提出來。他就剃頭刮臉，換衣服，進到法老面前。 \end{tabularx} \\ \\ \relax
41:15 & \begin{tabularx}{0.7\textwidth}{X} 法老對約瑟說：「我做了一個夢，沒有人能講解。我聽人說，你聽了夢就能講解。」 \end{tabularx} \\ \\ \relax
41:16 & \begin{tabularx}{0.7\textwidth}{X} 約瑟回答法老說：「這不在乎我。神必應允法老平安。」 \end{tabularx} \\ \\ \relax
41:17 & \begin{tabularx}{0.7\textwidth}{X} 法老對約瑟說：「在我的夢中，看哪，我站在尼羅河邊， \end{tabularx} \\ \\ \relax
41:18 & \begin{tabularx}{0.7\textwidth}{X} 看哪，有七頭母牛從尼羅河裡上來，肌肉肥壯，外形俊美，在蘆葦中吃草。 \end{tabularx} \\ \\ \relax
41:19 & \begin{tabularx}{0.7\textwidth}{X} 看哪，隨後又有七頭母牛上來，虛弱，外形很醜陋，肌肉又乾瘦，在埃及全地，我沒有見過這樣醜陋的牛。 \end{tabularx} \\ \\ \relax
41:20 & \begin{tabularx}{0.7\textwidth}{X} 這乾瘦又醜陋的母牛吃了那先前的七頭肥母牛， \end{tabularx} \\ \\ \relax
41:21 & \begin{tabularx}{0.7\textwidth}{X} 進了肚子以後卻看不出已經進了肚子，那醜陋的長相仍舊和先前一樣。我就醒了。 \end{tabularx} \\ \\ \relax
41:22 & \begin{tabularx}{0.7\textwidth}{X} 我又在夢中觀看，看哪，一株麥桿長了七個穗子，又飽滿又佳美， \end{tabularx} \\ \\ \relax
41:23 & \begin{tabularx}{0.7\textwidth}{X} 看哪，隨後又長出七個穗子，枯槁，細弱，又被東風吹焦了。 \end{tabularx} \\ \\ \relax
41:24 & \begin{tabularx}{0.7\textwidth}{X} 這些細弱的穗子吞了那七個佳美的穗子。我告訴術士，卻沒有人能為我講解。」 \end{tabularx} \\ \\ \relax
41:25 & \begin{tabularx}{0.7\textwidth}{X} 約瑟對法老說：「法老的夢是同一個。神已把要做的事指示法老了。 \end{tabularx} \\ \\ \relax
41:26 & \begin{tabularx}{0.7\textwidth}{X} 七頭好母牛是七年，七個佳美的穗子也是七年，這是同一個夢。 \end{tabularx} \\ \\ \relax
41:27 & \begin{tabularx}{0.7\textwidth}{X} 那隨後上來的七頭乾瘦又醜陋的母牛是七年；那七個空心，被東風吹焦的穗子也一樣，都是七個荒年。 \end{tabularx} \\ \\ \relax
41:28 & \begin{tabularx}{0.7\textwidth}{X} 這就是我對法老所說，神已把要做的事顯明給法老了。 \end{tabularx} \\ \\ \relax
41:29 & \begin{tabularx}{0.7\textwidth}{X} 看哪，必有七個大豐年來到埃及全地， \end{tabularx} \\ \\ \relax
41:30 & \begin{tabularx}{0.7\textwidth}{X} 隨後又有七個荒年，甚至埃及地的人都忘了先前的豐收，這地必被饑荒所滅。 \end{tabularx} \\ \\ \relax
41:31 & \begin{tabularx}{0.7\textwidth}{X} 因為那後來的饑荒非常嚴重，就不覺得這地先前有豐收。 \end{tabularx} \\ \\ \relax
41:32 & \begin{tabularx}{0.7\textwidth}{X} 至於法老兩次做夢，是因為神已經確定這事，神必速速成就。 \end{tabularx} \\ \\ \relax
41:33 & \begin{tabularx}{0.7\textwidth}{X} 現在，請法老選一個聰明又有智慧的人，委派他治理埃及地。 \end{tabularx} \\ \\ \relax
41:34 & \begin{tabularx}{0.7\textwidth}{X} 請法老這樣做，委派官員治理這地，在七個豐年的期間，徵收埃及地出產的五分之一， \end{tabularx} \\ \\ \relax
41:35 & \begin{tabularx}{0.7\textwidth}{X} 叫他們聚集未來豐年一切的糧食，積存五穀歸在法老的手下作糧食，儲藏在各城裡。 \end{tabularx} \\ \\ \relax
41:36 & \begin{tabularx}{0.7\textwidth}{X} 這糧食可以為這地作儲備，為了埃及地要來的七個荒年，免得這地被饑荒所滅。」 \end{tabularx} \\ \\ \relax
41:37 & \begin{tabularx}{0.7\textwidth}{X} 這事在法老和他眾臣僕眼中都覺得好。 \end{tabularx} \\ \\ \relax
41:38 & \begin{tabularx}{0.7\textwidth}{X} 法老對臣僕說：「像這樣的人，有神的靈在他裡面，我們豈能找得著呢？」 \end{tabularx} \\ \\ \relax
41:39 & \begin{tabularx}{0.7\textwidth}{X} 法老對約瑟說：「神既指示你這一切事，就沒有人像你這樣聰明又有智慧。 \end{tabularx} \\ \\ \relax
41:40 & \begin{tabularx}{0.7\textwidth}{X} 你可以治理我的家；我的百姓都必服從你口中的命令。惟獨在寶座上，我比你大。」 \end{tabularx} \\ \\ \relax
41:41 & \begin{tabularx}{0.7\textwidth}{X} 法老又對約瑟說：「看，我委派你治理埃及全地。」 \end{tabularx} \\ \\ \relax
41:42 & \begin{tabularx}{0.7\textwidth}{X} 法老就脫下手上帶印的戒指，戴在約瑟的手上，給他穿上細麻衣，把金鏈戴在他的頸項上， \end{tabularx} \\ \\ \relax
41:43 & \begin{tabularx}{0.7\textwidth}{X} 又給約瑟坐他的副座車，在他前面有人呼叫說：「跪下。」於是，法老委派他治理埃及全地。 \end{tabularx} \\ \\ \relax
41:44 & \begin{tabularx}{0.7\textwidth}{X} 法老對約瑟說：「我是法老，若沒有你的命令，埃及全地的人都不可擅自辦事。」 \end{tabularx} \\ \\ \relax
41:45 & \begin{tabularx}{0.7\textwidth}{X} 法老給約瑟起名叫撒發那特‧巴內亞，又將安城的祭司波提‧非拉的女兒亞西納給他為妻。約瑟就出去治理埃及地。 \end{tabularx} \\ \\ \relax
41:46 & \begin{tabularx}{0.7\textwidth}{X} 約瑟在埃及王法老面前侍立的時候年三十歲。約瑟從法老面前出去，巡行埃及全地。 \end{tabularx} \\ \\ \relax
41:47 & \begin{tabularx}{0.7\textwidth}{X} 七個豐年之內，地的出產極其豐盛， \end{tabularx} \\ \\ \relax
41:48 & \begin{tabularx}{0.7\textwidth}{X} 約瑟聚集埃及地七年一切的糧食，把糧食積存在各城裡，就是把各城周圍田地的糧食都積存在該城裡。 \end{tabularx} \\ \\ \relax
41:49 & \begin{tabularx}{0.7\textwidth}{X} 約瑟積存的五穀很多，如同海邊的沙，無法計算，數也數不清。 \end{tabularx} \\ \\ \relax
41:50 & \begin{tabularx}{0.7\textwidth}{X} 荒年未到以前，安城的祭司波提‧非拉的女兒亞西納為約瑟生了兩個兒子。 \end{tabularx} \\ \\ \relax
41:51 & \begin{tabularx}{0.7\textwidth}{X} 約瑟給長子起名叫瑪拿西，因為他說：「神使我忘了一切的困苦和我父的全家。」 \end{tabularx} \\ \\ \relax
41:52 & \begin{tabularx}{0.7\textwidth}{X} 他給次子起名叫以法蓮，因為他說：「神使我在受苦的地方興盛。」 \end{tabularx} \\ \\ \relax
41:53 & \begin{tabularx}{0.7\textwidth}{X} 埃及地的七個豐年一過， \end{tabularx} \\ \\ \relax
41:54 & \begin{tabularx}{0.7\textwidth}{X} 七個荒年就來了，正如約瑟所說的。各地都有饑荒，惟獨埃及全地有糧食。 \end{tabularx} \\ \\ \relax
41:55 & \begin{tabularx}{0.7\textwidth}{X} 等到埃及全地也有了饑荒，眾百姓就向法老哀求糧食。法老對所有的埃及人說：「你們到約瑟那裡去，凡他所說的，你們都要做。」 \end{tabularx} \\ \\ \relax
41:56 & \begin{tabularx}{0.7\textwidth}{X} 當時饑荒遍滿了全地，約瑟就開了各處的糧倉，賣糧食給埃及人。埃及地的饑荒非常嚴重。 \end{tabularx} \\ \\ \relax
41:57 & \begin{tabularx}{0.7\textwidth}{X} 各地的人都去埃及，到約瑟那裡買糧食，因為全地的饑荒非常嚴重。 \end{tabularx} \\ \\
[1ex]
\hline
\hline
\end{longtable}


\newpage

\section{第70屆~港九培靈會~焦源濂~第六講}
\label{sec:869}
\textbf{舊夢重提}
\newline
\newline
連結: \href{https://www.hkbibleconference.org/session-message/view/869}{\texttt{https://www.hkbibleconference.org/session-message/view/869}}
\newline
\newline
\hyperref[sec:868]{< < < PREV SERMON < < <}
~
\hyperref[sec:toc]{[返目錄]}
~
\hyperref[sec:870]{> > > NEXT SERMON > > >}
\newline
\newline
創世記 41:1-42:38
\newline
\begin{longtable}{cl}
\hline
\hline
章節 & 經文 (和合本修訂版)\\
\hline
41:1 & \begin{tabularx}{0.7\textwidth}{X} 過了兩年，法老做夢，看哪，自己站在尼羅河邊， \end{tabularx} \\ \\ \relax
41:2 & \begin{tabularx}{0.7\textwidth}{X} 看哪，有七頭母牛從尼羅河裡上來，長相俊美，肌肉肥壯，在蘆葦中吃草。 \end{tabularx} \\ \\ \relax
41:3 & \begin{tabularx}{0.7\textwidth}{X} 看哪，隨後又有七頭母牛從尼羅河裡上來，長相醜陋，肌肉乾瘦，與那七頭母牛一同站在河邊。 \end{tabularx} \\ \\ \relax
41:4 & \begin{tabularx}{0.7\textwidth}{X} 這長相醜陋，肌肉乾瘦的七頭母牛吃了那長相俊美又肥壯的七頭母牛。法老就醒了。 \end{tabularx} \\ \\ \relax
41:5 & \begin{tabularx}{0.7\textwidth}{X} 他又睡著，第二次做夢，看哪，一株麥桿長了七個穗子，又肥大又佳美， \end{tabularx} \\ \\ \relax
41:6 & \begin{tabularx}{0.7\textwidth}{X} 看哪，隨後又長出七個穗子，又細弱又被東風吹焦了。 \end{tabularx} \\ \\ \relax
41:7 & \begin{tabularx}{0.7\textwidth}{X} 這細弱的穗子吞了那七個又肥大又飽滿的穗子。法老醒了，看哪，是個夢。 \end{tabularx} \\ \\ \relax
41:8 & \begin{tabularx}{0.7\textwidth}{X} 到了早晨，法老心裡不安，就派人把埃及所有的術士和智慧人都召來。法老把所做的夢告訴他們，但是沒有人能為法老解夢。 \end{tabularx} \\ \\ \relax
41:9 & \begin{tabularx}{0.7\textwidth}{X} 那時司酒長對法老說：「我今日想起我的罪來。 \end{tabularx} \\ \\ \relax
41:10 & \begin{tabularx}{0.7\textwidth}{X} 從前法老對臣僕發怒，把我和司膳長關在護衛長府內的監牢裡。 \end{tabularx} \\ \\ \relax
41:11 & \begin{tabularx}{0.7\textwidth}{X} 我們兩人在同一晚上各做一夢，每個夢都有各自的解釋。 \end{tabularx} \\ \\ \relax
41:12 & \begin{tabularx}{0.7\textwidth}{X} 同我們在一起有一個希伯來的年輕人，是護衛長的僕人。我們告訴他，他就為我們解夢，照著各人的夢講解。 \end{tabularx} \\ \\ \relax
41:13 & \begin{tabularx}{0.7\textwidth}{X} 後來事情正如他給我們講解的實現了，我官復原職，司膳長被掛起來了。」 \end{tabularx} \\ \\ \relax
41:14 & \begin{tabularx}{0.7\textwidth}{X} 於是法老派人去召約瑟，他們就急忙把他從牢裡提出來。他就剃頭刮臉，換衣服，進到法老面前。 \end{tabularx} \\ \\ \relax
41:15 & \begin{tabularx}{0.7\textwidth}{X} 法老對約瑟說：「我做了一個夢，沒有人能講解。我聽人說，你聽了夢就能講解。」 \end{tabularx} \\ \\ \relax
41:16 & \begin{tabularx}{0.7\textwidth}{X} 約瑟回答法老說：「這不在乎我。神必應允法老平安。」 \end{tabularx} \\ \\ \relax
41:17 & \begin{tabularx}{0.7\textwidth}{X} 法老對約瑟說：「在我的夢中，看哪，我站在尼羅河邊， \end{tabularx} \\ \\ \relax
41:18 & \begin{tabularx}{0.7\textwidth}{X} 看哪，有七頭母牛從尼羅河裡上來，肌肉肥壯，外形俊美，在蘆葦中吃草。 \end{tabularx} \\ \\ \relax
41:19 & \begin{tabularx}{0.7\textwidth}{X} 看哪，隨後又有七頭母牛上來，虛弱，外形很醜陋，肌肉又乾瘦，在埃及全地，我沒有見過這樣醜陋的牛。 \end{tabularx} \\ \\ \relax
41:20 & \begin{tabularx}{0.7\textwidth}{X} 這乾瘦又醜陋的母牛吃了那先前的七頭肥母牛， \end{tabularx} \\ \\ \relax
41:21 & \begin{tabularx}{0.7\textwidth}{X} 進了肚子以後卻看不出已經進了肚子，那醜陋的長相仍舊和先前一樣。我就醒了。 \end{tabularx} \\ \\ \relax
41:22 & \begin{tabularx}{0.7\textwidth}{X} 我又在夢中觀看，看哪，一株麥桿長了七個穗子，又飽滿又佳美， \end{tabularx} \\ \\ \relax
41:23 & \begin{tabularx}{0.7\textwidth}{X} 看哪，隨後又長出七個穗子，枯槁，細弱，又被東風吹焦了。 \end{tabularx} \\ \\ \relax
41:24 & \begin{tabularx}{0.7\textwidth}{X} 這些細弱的穗子吞了那七個佳美的穗子。我告訴術士，卻沒有人能為我講解。」 \end{tabularx} \\ \\ \relax
41:25 & \begin{tabularx}{0.7\textwidth}{X} 約瑟對法老說：「法老的夢是同一個。神已把要做的事指示法老了。 \end{tabularx} \\ \\ \relax
41:26 & \begin{tabularx}{0.7\textwidth}{X} 七頭好母牛是七年，七個佳美的穗子也是七年，這是同一個夢。 \end{tabularx} \\ \\ \relax
41:27 & \begin{tabularx}{0.7\textwidth}{X} 那隨後上來的七頭乾瘦又醜陋的母牛是七年；那七個空心，被東風吹焦的穗子也一樣，都是七個荒年。 \end{tabularx} \\ \\ \relax
41:28 & \begin{tabularx}{0.7\textwidth}{X} 這就是我對法老所說，神已把要做的事顯明給法老了。 \end{tabularx} \\ \\ \relax
41:29 & \begin{tabularx}{0.7\textwidth}{X} 看哪，必有七個大豐年來到埃及全地， \end{tabularx} \\ \\ \relax
41:30 & \begin{tabularx}{0.7\textwidth}{X} 隨後又有七個荒年，甚至埃及地的人都忘了先前的豐收，這地必被饑荒所滅。 \end{tabularx} \\ \\ \relax
41:31 & \begin{tabularx}{0.7\textwidth}{X} 因為那後來的饑荒非常嚴重，就不覺得這地先前有豐收。 \end{tabularx} \\ \\ \relax
41:32 & \begin{tabularx}{0.7\textwidth}{X} 至於法老兩次做夢，是因為神已經確定這事，神必速速成就。 \end{tabularx} \\ \\ \relax
41:33 & \begin{tabularx}{0.7\textwidth}{X} 現在，請法老選一個聰明又有智慧的人，委派他治理埃及地。 \end{tabularx} \\ \\ \relax
41:34 & \begin{tabularx}{0.7\textwidth}{X} 請法老這樣做，委派官員治理這地，在七個豐年的期間，徵收埃及地出產的五分之一， \end{tabularx} \\ \\ \relax
41:35 & \begin{tabularx}{0.7\textwidth}{X} 叫他們聚集未來豐年一切的糧食，積存五穀歸在法老的手下作糧食，儲藏在各城裡。 \end{tabularx} \\ \\ \relax
41:36 & \begin{tabularx}{0.7\textwidth}{X} 這糧食可以為這地作儲備，為了埃及地要來的七個荒年，免得這地被饑荒所滅。」 \end{tabularx} \\ \\ \relax
41:37 & \begin{tabularx}{0.7\textwidth}{X} 這事在法老和他眾臣僕眼中都覺得好。 \end{tabularx} \\ \\ \relax
41:38 & \begin{tabularx}{0.7\textwidth}{X} 法老對臣僕說：「像這樣的人，有神的靈在他裡面，我們豈能找得著呢？」 \end{tabularx} \\ \\ \relax
41:39 & \begin{tabularx}{0.7\textwidth}{X} 法老對約瑟說：「神既指示你這一切事，就沒有人像你這樣聰明又有智慧。 \end{tabularx} \\ \\ \relax
41:40 & \begin{tabularx}{0.7\textwidth}{X} 你可以治理我的家；我的百姓都必服從你口中的命令。惟獨在寶座上，我比你大。」 \end{tabularx} \\ \\ \relax
41:41 & \begin{tabularx}{0.7\textwidth}{X} 法老又對約瑟說：「看，我委派你治理埃及全地。」 \end{tabularx} \\ \\ \relax
41:42 & \begin{tabularx}{0.7\textwidth}{X} 法老就脫下手上帶印的戒指，戴在約瑟的手上，給他穿上細麻衣，把金鏈戴在他的頸項上， \end{tabularx} \\ \\ \relax
41:43 & \begin{tabularx}{0.7\textwidth}{X} 又給約瑟坐他的副座車，在他前面有人呼叫說：「跪下。」於是，法老委派他治理埃及全地。 \end{tabularx} \\ \\ \relax
41:44 & \begin{tabularx}{0.7\textwidth}{X} 法老對約瑟說：「我是法老，若沒有你的命令，埃及全地的人都不可擅自辦事。」 \end{tabularx} \\ \\ \relax
41:45 & \begin{tabularx}{0.7\textwidth}{X} 法老給約瑟起名叫撒發那特‧巴內亞，又將安城的祭司波提‧非拉的女兒亞西納給他為妻。約瑟就出去治理埃及地。 \end{tabularx} \\ \\ \relax
41:46 & \begin{tabularx}{0.7\textwidth}{X} 約瑟在埃及王法老面前侍立的時候年三十歲。約瑟從法老面前出去，巡行埃及全地。 \end{tabularx} \\ \\ \relax
41:47 & \begin{tabularx}{0.7\textwidth}{X} 七個豐年之內，地的出產極其豐盛， \end{tabularx} \\ \\ \relax
41:48 & \begin{tabularx}{0.7\textwidth}{X} 約瑟聚集埃及地七年一切的糧食，把糧食積存在各城裡，就是把各城周圍田地的糧食都積存在該城裡。 \end{tabularx} \\ \\ \relax
41:49 & \begin{tabularx}{0.7\textwidth}{X} 約瑟積存的五穀很多，如同海邊的沙，無法計算，數也數不清。 \end{tabularx} \\ \\ \relax
41:50 & \begin{tabularx}{0.7\textwidth}{X} 荒年未到以前，安城的祭司波提‧非拉的女兒亞西納為約瑟生了兩個兒子。 \end{tabularx} \\ \\ \relax
41:51 & \begin{tabularx}{0.7\textwidth}{X} 約瑟給長子起名叫瑪拿西，因為他說：「神使我忘了一切的困苦和我父的全家。」 \end{tabularx} \\ \\ \relax
41:52 & \begin{tabularx}{0.7\textwidth}{X} 他給次子起名叫以法蓮，因為他說：「神使我在受苦的地方興盛。」 \end{tabularx} \\ \\ \relax
41:53 & \begin{tabularx}{0.7\textwidth}{X} 埃及地的七個豐年一過， \end{tabularx} \\ \\ \relax
41:54 & \begin{tabularx}{0.7\textwidth}{X} 七個荒年就來了，正如約瑟所說的。各地都有饑荒，惟獨埃及全地有糧食。 \end{tabularx} \\ \\ \relax
41:55 & \begin{tabularx}{0.7\textwidth}{X} 等到埃及全地也有了饑荒，眾百姓就向法老哀求糧食。法老對所有的埃及人說：「你們到約瑟那裡去，凡他所說的，你們都要做。」 \end{tabularx} \\ \\ \relax
41:56 & \begin{tabularx}{0.7\textwidth}{X} 當時饑荒遍滿了全地，約瑟就開了各處的糧倉，賣糧食給埃及人。埃及地的饑荒非常嚴重。 \end{tabularx} \\ \\ \relax
41:57 & \begin{tabularx}{0.7\textwidth}{X} 各地的人都去埃及，到約瑟那裡買糧食，因為全地的饑荒非常嚴重。 \end{tabularx} \\ \\ \relax
42:1 & \begin{tabularx}{0.7\textwidth}{X} 雅各見埃及有糧，就對兒子們說：「你們為甚麼彼此對看呢？」 \end{tabularx} \\ \\ \relax
42:2 & \begin{tabularx}{0.7\textwidth}{X} 他又說：「看哪，我聽見埃及有糧，你們可以下到那裡，從那裡為我們買些糧來，我們就可以存活，不至於死。」 \end{tabularx} \\ \\ \relax
42:3 & \begin{tabularx}{0.7\textwidth}{X} 於是，約瑟的十個哥哥都下去，到埃及買糧食。 \end{tabularx} \\ \\ \relax
42:4 & \begin{tabularx}{0.7\textwidth}{X} 至於約瑟的弟弟便雅憫，雅各沒有派他和哥哥們同去，因為雅各說：「恐怕他遭難。」 \end{tabularx} \\ \\ \relax
42:5 & \begin{tabularx}{0.7\textwidth}{X} 以色列的兒子們來了，在前來的人當中，為要買糧食，因為迦南地也有饑荒。 \end{tabularx} \\ \\ \relax
42:6 & \begin{tabularx}{0.7\textwidth}{X} 當時在埃及地掌權的人是約瑟，賣糧給各地眾百姓的就是他。約瑟的哥哥們來了，臉伏於地，向他下拜。 \end{tabularx} \\ \\ \relax
42:7 & \begin{tabularx}{0.7\textwidth}{X} 約瑟看見他哥哥們，就認出他們，卻對他們裝作陌生人，向他們說嚴厲的話，對他們說：「你們從哪裡來？」他們說：「我們從迦南地來買糧。」 \end{tabularx} \\ \\ \relax
42:8 & \begin{tabularx}{0.7\textwidth}{X} 約瑟認得他哥哥們，他們卻不認得他。 \end{tabularx} \\ \\ \relax
42:9 & \begin{tabularx}{0.7\textwidth}{X} 約瑟想起從前所做的那兩個夢，就對他們說：「你們是奸細，你們來是要窺探這地的虛實。」 \end{tabularx} \\ \\ \relax
42:10 & \begin{tabularx}{0.7\textwidth}{X} 他們對他說：「我主啊，不是的，僕人們是來買糧的。 \end{tabularx} \\ \\ \relax
42:11 & \begin{tabularx}{0.7\textwidth}{X} 我們都是同一個人的兒子，我們是誠實的人。僕人們並不是奸細。」 \end{tabularx} \\ \\ \relax
42:12 & \begin{tabularx}{0.7\textwidth}{X} 約瑟對他們說：「不，你們一定是窺探這地的虛實來的。」 \end{tabularx} \\ \\ \relax
42:13 & \begin{tabularx}{0.7\textwidth}{X} 他們說：「僕人們本是兄弟十二人，我們都是迦南地同一個人的兒子。看哪，最小的今日在我們父親那裡，有一個不在了。」 \end{tabularx} \\ \\ \relax
42:14 & \begin{tabularx}{0.7\textwidth}{X} 約瑟對他們說：「我剛才對你們說過了，你們是奸細！ \end{tabularx} \\ \\ \relax
42:15 & \begin{tabularx}{0.7\textwidth}{X} 我指著法老的性命起誓，若是你們最小的弟弟不到這裡來，你們就不可以離開這裡；這樣你們就可以證實自己了。 \end{tabularx} \\ \\ \relax
42:16 & \begin{tabularx}{0.7\textwidth}{X} 要派你們當中的一個人去，把你們的弟弟帶來。至於你們，都要關在這裡，好證實你們的話是不是真的。若不是，我指著法老的性命起誓，你們一定是奸細。」 \end{tabularx} \\ \\ \relax
42:17 & \begin{tabularx}{0.7\textwidth}{X} 於是約瑟把他們一起都關在監裡三天。 \end{tabularx} \\ \\ \relax
42:18 & \begin{tabularx}{0.7\textwidth}{X} 第三天，約瑟對他們說：「我是敬畏神的，你們這麼做就可以活。 \end{tabularx} \\ \\ \relax
42:19 & \begin{tabularx}{0.7\textwidth}{X} 如果你們是誠實的人，留你們兄弟中的一個關在監牢裡，你們帶糧食回去，救你們家的饑荒， \end{tabularx} \\ \\ \relax
42:20 & \begin{tabularx}{0.7\textwidth}{X} 再把你們最小的弟弟帶到我這裡來。如此，你們的話就是真的了，你們也不至於死。」他們就照樣做了。 \end{tabularx} \\ \\ \relax
42:21 & \begin{tabularx}{0.7\textwidth}{X} 他們彼此說：「我們在弟弟身上實在犯了罪。他哀求我們的時候，我們看見他的痛苦，卻不肯聽，所以這場苦難臨到我們。」 \end{tabularx} \\ \\ \relax
42:22 & \begin{tabularx}{0.7\textwidth}{X} 呂便回答他們說：「我不是對你們說過，不可傷害那孩子嗎？只是你們不肯聽，看哪，他的血在追討了。」 \end{tabularx} \\ \\ \relax
42:23 & \begin{tabularx}{0.7\textwidth}{X} 他們不知道約瑟在聽，因為在他們之間有傳譯官。 \end{tabularx} \\ \\ \relax
42:24 & \begin{tabularx}{0.7\textwidth}{X} 約瑟轉身離開他們，哭了一場，又回來對他們說話，就從他們中間抓了西緬，在他們眼前捆綁他。 \end{tabularx} \\ \\ \relax
42:25 & \begin{tabularx}{0.7\textwidth}{X} 約瑟吩咐人把他們的器皿裝滿糧食，把各人的銀子退還在各人的袋裡，又給他們路上需用的食物。人就為他們這樣做了。 \end{tabularx} \\ \\ \relax
42:26 & \begin{tabularx}{0.7\textwidth}{X} 他們把糧食馱在驢上，離開那裡去了。 \end{tabularx} \\ \\ \relax
42:27 & \begin{tabularx}{0.7\textwidth}{X} 到了住宿的地方，有一個人打開袋子，要拿飼料餵驢，就看見自己的銀子，看哪，仍在袋口上。 \end{tabularx} \\ \\ \relax
42:28 & \begin{tabularx}{0.7\textwidth}{X} 他對兄弟們說：「我的銀子退回來了，看哪，還在我袋子裡！」他們戰戰兢兢，心都快跳出來了，彼此說：「神向我們做的是甚麼呢？」 \end{tabularx} \\ \\ \relax
42:29 & \begin{tabularx}{0.7\textwidth}{X} 他們來到迦南地他們的父親雅各那裡，把所遭遇的事都告訴他，說： \end{tabularx} \\ \\ \relax
42:30 & \begin{tabularx}{0.7\textwidth}{X} 「那地的主對我們說嚴厲的話，把我們當作窺探那地的奸細。 \end{tabularx} \\ \\ \relax
42:31 & \begin{tabularx}{0.7\textwidth}{X} 我們對他說：『我們是誠實的人，並不是奸細。 \end{tabularx} \\ \\ \relax
42:32 & \begin{tabularx}{0.7\textwidth}{X} 我們本是兄弟十二人，都是同一個父親的兒子，有一個不在了，最小的今日和我們父親在迦南地。』 \end{tabularx} \\ \\ \relax
42:33 & \begin{tabularx}{0.7\textwidth}{X} 那地的主對我們說：『只有這樣我才知道你們是誠實的人：留你們兄弟中的一個在我這裡，你們帶糧食回去，救你們家的饑荒， \end{tabularx} \\ \\ \relax
42:34 & \begin{tabularx}{0.7\textwidth}{X} 再把你們最小的弟弟帶到我這裡來，我就知道你們不是奸細，是誠實的人。然後，我就把你們的兄弟交還你們，你們也可以在此地做買賣。』」 \end{tabularx} \\ \\ \relax
42:35 & \begin{tabularx}{0.7\textwidth}{X} 後來他們倒空袋子，看哪，各人的銀囊都在袋子裡。他們和父親看見銀囊就都害怕。 \end{tabularx} \\ \\ \relax
42:36 & \begin{tabularx}{0.7\textwidth}{X} 他們的父親雅各對他們說：「你們害我喪失了我的兒子：約瑟不在了，西緬也不在了，你們還要帶走便雅憫！這些事都臨到我身上了。」 \end{tabularx} \\ \\ \relax
42:37 & \begin{tabularx}{0.7\textwidth}{X} 呂便對他父親說：「我若不帶他回來給你，你可以殺我的兩個兒子。只管把他交在我手裡，我必帶他回來給你。」 \end{tabularx} \\ \\ \relax
42:38 & \begin{tabularx}{0.7\textwidth}{X} 雅各說：「我的兒子不可與你們一同下去。他哥哥死了，只剩下他。他若在你們行走的路上遭難，你們就害我白髮蒼蒼、悲悲慘慘下陰間去了。」 \end{tabularx} \\ \\
[1ex]
\hline
\hline
\end{longtable}


\newpage

\section{第70屆~港九培靈會~焦源濂~第七講}
\label{sec:870}
\textbf{卻沒有一個人能看透他}
\newline
\newline
連結: \href{https://www.hkbibleconference.org/session-message/view/870}{\texttt{https://www.hkbibleconference.org/session-message/view/870}}
\newline
\newline
\hyperref[sec:869]{< < < PREV SERMON < < <}
~
\hyperref[sec:toc]{[返目錄]}
~
\hyperref[sec:871]{> > > NEXT SERMON > > >}
\newline
\newline
創世記 42:1-43:34
\newline
\begin{longtable}{cl}
\hline
\hline
章節 & 經文 (和合本修訂版)\\
\hline
42:1 & \begin{tabularx}{0.7\textwidth}{X} 雅各見埃及有糧，就對兒子們說：「你們為甚麼彼此對看呢？」 \end{tabularx} \\ \\ \relax
42:2 & \begin{tabularx}{0.7\textwidth}{X} 他又說：「看哪，我聽見埃及有糧，你們可以下到那裡，從那裡為我們買些糧來，我們就可以存活，不至於死。」 \end{tabularx} \\ \\ \relax
42:3 & \begin{tabularx}{0.7\textwidth}{X} 於是，約瑟的十個哥哥都下去，到埃及買糧食。 \end{tabularx} \\ \\ \relax
42:4 & \begin{tabularx}{0.7\textwidth}{X} 至於約瑟的弟弟便雅憫，雅各沒有派他和哥哥們同去，因為雅各說：「恐怕他遭難。」 \end{tabularx} \\ \\ \relax
42:5 & \begin{tabularx}{0.7\textwidth}{X} 以色列的兒子們來了，在前來的人當中，為要買糧食，因為迦南地也有饑荒。 \end{tabularx} \\ \\ \relax
42:6 & \begin{tabularx}{0.7\textwidth}{X} 當時在埃及地掌權的人是約瑟，賣糧給各地眾百姓的就是他。約瑟的哥哥們來了，臉伏於地，向他下拜。 \end{tabularx} \\ \\ \relax
42:7 & \begin{tabularx}{0.7\textwidth}{X} 約瑟看見他哥哥們，就認出他們，卻對他們裝作陌生人，向他們說嚴厲的話，對他們說：「你們從哪裡來？」他們說：「我們從迦南地來買糧。」 \end{tabularx} \\ \\ \relax
42:8 & \begin{tabularx}{0.7\textwidth}{X} 約瑟認得他哥哥們，他們卻不認得他。 \end{tabularx} \\ \\ \relax
42:9 & \begin{tabularx}{0.7\textwidth}{X} 約瑟想起從前所做的那兩個夢，就對他們說：「你們是奸細，你們來是要窺探這地的虛實。」 \end{tabularx} \\ \\ \relax
42:10 & \begin{tabularx}{0.7\textwidth}{X} 他們對他說：「我主啊，不是的，僕人們是來買糧的。 \end{tabularx} \\ \\ \relax
42:11 & \begin{tabularx}{0.7\textwidth}{X} 我們都是同一個人的兒子，我們是誠實的人。僕人們並不是奸細。」 \end{tabularx} \\ \\ \relax
42:12 & \begin{tabularx}{0.7\textwidth}{X} 約瑟對他們說：「不，你們一定是窺探這地的虛實來的。」 \end{tabularx} \\ \\ \relax
42:13 & \begin{tabularx}{0.7\textwidth}{X} 他們說：「僕人們本是兄弟十二人，我們都是迦南地同一個人的兒子。看哪，最小的今日在我們父親那裡，有一個不在了。」 \end{tabularx} \\ \\ \relax
42:14 & \begin{tabularx}{0.7\textwidth}{X} 約瑟對他們說：「我剛才對你們說過了，你們是奸細！ \end{tabularx} \\ \\ \relax
42:15 & \begin{tabularx}{0.7\textwidth}{X} 我指著法老的性命起誓，若是你們最小的弟弟不到這裡來，你們就不可以離開這裡；這樣你們就可以證實自己了。 \end{tabularx} \\ \\ \relax
42:16 & \begin{tabularx}{0.7\textwidth}{X} 要派你們當中的一個人去，把你們的弟弟帶來。至於你們，都要關在這裡，好證實你們的話是不是真的。若不是，我指著法老的性命起誓，你們一定是奸細。」 \end{tabularx} \\ \\ \relax
42:17 & \begin{tabularx}{0.7\textwidth}{X} 於是約瑟把他們一起都關在監裡三天。 \end{tabularx} \\ \\ \relax
42:18 & \begin{tabularx}{0.7\textwidth}{X} 第三天，約瑟對他們說：「我是敬畏神的，你們這麼做就可以活。 \end{tabularx} \\ \\ \relax
42:19 & \begin{tabularx}{0.7\textwidth}{X} 如果你們是誠實的人，留你們兄弟中的一個關在監牢裡，你們帶糧食回去，救你們家的饑荒， \end{tabularx} \\ \\ \relax
42:20 & \begin{tabularx}{0.7\textwidth}{X} 再把你們最小的弟弟帶到我這裡來。如此，你們的話就是真的了，你們也不至於死。」他們就照樣做了。 \end{tabularx} \\ \\ \relax
42:21 & \begin{tabularx}{0.7\textwidth}{X} 他們彼此說：「我們在弟弟身上實在犯了罪。他哀求我們的時候，我們看見他的痛苦，卻不肯聽，所以這場苦難臨到我們。」 \end{tabularx} \\ \\ \relax
42:22 & \begin{tabularx}{0.7\textwidth}{X} 呂便回答他們說：「我不是對你們說過，不可傷害那孩子嗎？只是你們不肯聽，看哪，他的血在追討了。」 \end{tabularx} \\ \\ \relax
42:23 & \begin{tabularx}{0.7\textwidth}{X} 他們不知道約瑟在聽，因為在他們之間有傳譯官。 \end{tabularx} \\ \\ \relax
42:24 & \begin{tabularx}{0.7\textwidth}{X} 約瑟轉身離開他們，哭了一場，又回來對他們說話，就從他們中間抓了西緬，在他們眼前捆綁他。 \end{tabularx} \\ \\ \relax
42:25 & \begin{tabularx}{0.7\textwidth}{X} 約瑟吩咐人把他們的器皿裝滿糧食，把各人的銀子退還在各人的袋裡，又給他們路上需用的食物。人就為他們這樣做了。 \end{tabularx} \\ \\ \relax
42:26 & \begin{tabularx}{0.7\textwidth}{X} 他們把糧食馱在驢上，離開那裡去了。 \end{tabularx} \\ \\ \relax
42:27 & \begin{tabularx}{0.7\textwidth}{X} 到了住宿的地方，有一個人打開袋子，要拿飼料餵驢，就看見自己的銀子，看哪，仍在袋口上。 \end{tabularx} \\ \\ \relax
42:28 & \begin{tabularx}{0.7\textwidth}{X} 他對兄弟們說：「我的銀子退回來了，看哪，還在我袋子裡！」他們戰戰兢兢，心都快跳出來了，彼此說：「神向我們做的是甚麼呢？」 \end{tabularx} \\ \\ \relax
42:29 & \begin{tabularx}{0.7\textwidth}{X} 他們來到迦南地他們的父親雅各那裡，把所遭遇的事都告訴他，說： \end{tabularx} \\ \\ \relax
42:30 & \begin{tabularx}{0.7\textwidth}{X} 「那地的主對我們說嚴厲的話，把我們當作窺探那地的奸細。 \end{tabularx} \\ \\ \relax
42:31 & \begin{tabularx}{0.7\textwidth}{X} 我們對他說：『我們是誠實的人，並不是奸細。 \end{tabularx} \\ \\ \relax
42:32 & \begin{tabularx}{0.7\textwidth}{X} 我們本是兄弟十二人，都是同一個父親的兒子，有一個不在了，最小的今日和我們父親在迦南地。』 \end{tabularx} \\ \\ \relax
42:33 & \begin{tabularx}{0.7\textwidth}{X} 那地的主對我們說：『只有這樣我才知道你們是誠實的人：留你們兄弟中的一個在我這裡，你們帶糧食回去，救你們家的饑荒， \end{tabularx} \\ \\ \relax
42:34 & \begin{tabularx}{0.7\textwidth}{X} 再把你們最小的弟弟帶到我這裡來，我就知道你們不是奸細，是誠實的人。然後，我就把你們的兄弟交還你們，你們也可以在此地做買賣。』」 \end{tabularx} \\ \\ \relax
42:35 & \begin{tabularx}{0.7\textwidth}{X} 後來他們倒空袋子，看哪，各人的銀囊都在袋子裡。他們和父親看見銀囊就都害怕。 \end{tabularx} \\ \\ \relax
42:36 & \begin{tabularx}{0.7\textwidth}{X} 他們的父親雅各對他們說：「你們害我喪失了我的兒子：約瑟不在了，西緬也不在了，你們還要帶走便雅憫！這些事都臨到我身上了。」 \end{tabularx} \\ \\ \relax
42:37 & \begin{tabularx}{0.7\textwidth}{X} 呂便對他父親說：「我若不帶他回來給你，你可以殺我的兩個兒子。只管把他交在我手裡，我必帶他回來給你。」 \end{tabularx} \\ \\ \relax
42:38 & \begin{tabularx}{0.7\textwidth}{X} 雅各說：「我的兒子不可與你們一同下去。他哥哥死了，只剩下他。他若在你們行走的路上遭難，你們就害我白髮蒼蒼、悲悲慘慘下陰間去了。」 \end{tabularx} \\ \\ \relax
43:1 & \begin{tabularx}{0.7\textwidth}{X} 那地的饑荒非常嚴重。 \end{tabularx} \\ \\ \relax
43:2 & \begin{tabularx}{0.7\textwidth}{X} 他們從埃及帶來的糧食吃完了，父親對他們說：「你們再去給我們買些糧來。」 \end{tabularx} \\ \\ \relax
43:3 & \begin{tabularx}{0.7\textwidth}{X} 猶大對他說：「那人嚴厲地警告我們說：『你們的弟弟若不和你們同來，你們就不要來見我的面。』 \end{tabularx} \\ \\ \relax
43:4 & \begin{tabularx}{0.7\textwidth}{X} 你若派我們的弟弟跟我們同去，我們就下去給你買糧； \end{tabularx} \\ \\ \relax
43:5 & \begin{tabularx}{0.7\textwidth}{X} 你若不派他去，我們就不下去，因為那人對我們說：『你們的弟弟若不和你們同來，你們就不要來見我的面。』」 \end{tabularx} \\ \\ \relax
43:6 & \begin{tabularx}{0.7\textwidth}{X} 以色列說：「你們為甚麼這樣害我，告訴那人你們還有弟弟呢？」 \end{tabularx} \\ \\ \relax
43:7 & \begin{tabularx}{0.7\textwidth}{X} 他們說：「那人詳細問到我們和我們的家人，說：『你們的父親還在嗎？你們還有兄弟嗎？』我們就按著他的這些話告訴他，我們怎麼知道他會說：『把你們的弟弟帶下來』呢？」 \end{tabularx} \\ \\ \relax
43:8 & \begin{tabularx}{0.7\textwidth}{X} 猶大又對他父親以色列說：「請派這年輕人和我同去，我們就動身前去，好叫我們和你，以及我們的孩子都得存活，不至於死。 \end{tabularx} \\ \\ \relax
43:9 & \begin{tabularx}{0.7\textwidth}{X} 我為他擔保，你可以從我手中要人，我若不帶他回來交在你面前，我就對你永遠擔當這罪。 \end{tabularx} \\ \\ \relax
43:10 & \begin{tabularx}{0.7\textwidth}{X} 我們若沒有耽擱，現在第二趟都回來了。」 \end{tabularx} \\ \\ \relax
43:11 & \begin{tabularx}{0.7\textwidth}{X} 父親以色列對他們說：「如果必須如此，你們要這樣做：把本地土產中最好的乳香、蜂蜜、香料、沒藥、堅果、杏仁各取一點，放在器皿裡，帶下去送給那人作禮物。 \end{tabularx} \\ \\ \relax
43:12 & \begin{tabularx}{0.7\textwidth}{X} 手裡要帶雙倍的銀子，把退還在你們袋口的銀子親手帶回去；或許那是個失誤。 \end{tabularx} \\ \\ \relax
43:13 & \begin{tabularx}{0.7\textwidth}{X} 帶著你們的弟弟，動身再去見那人。 \end{tabularx} \\ \\ \relax
43:14 & \begin{tabularx}{0.7\textwidth}{X} 願全能的神使你們在那人面前蒙憐憫，放你們另一個兄弟和便雅憫回來。我若要失喪兒子，就喪了吧！」 \end{tabularx} \\ \\ \relax
43:15 & \begin{tabularx}{0.7\textwidth}{X} 於是，他們拿著那些禮物，手裡也帶雙倍的銀子，並且帶著便雅憫，動身下到埃及，站在約瑟面前。 \end{tabularx} \\ \\ \relax
43:16 & \begin{tabularx}{0.7\textwidth}{X} 約瑟見便雅憫和他們同來，就對管家說：「把這些人領到屋裡。要宰殺牲畜，預備宴席，因為中午這些人要跟我吃飯。」 \end{tabularx} \\ \\ \relax
43:17 & \begin{tabularx}{0.7\textwidth}{X} 那人就照約瑟所說的去做，領他們進約瑟的屋裡。 \end{tabularx} \\ \\ \relax
43:18 & \begin{tabularx}{0.7\textwidth}{X} 這些人因為被領到約瑟的屋裡，就害怕，說：「領我們到這裡來，必是因為當初退還在我們袋裡的銀子，要設計害我們，抓我們去當奴隸，搶奪我們的驢。」 \end{tabularx} \\ \\ \relax
43:19 & \begin{tabularx}{0.7\textwidth}{X} 他們就挨近約瑟的管家，在屋子門口和他說話， \end{tabularx} \\ \\ \relax
43:20 & \begin{tabularx}{0.7\textwidth}{X} 說：「我主啊，求求你，我們當初下來，真的是要買糧食。 \end{tabularx} \\ \\ \relax
43:21 & \begin{tabularx}{0.7\textwidth}{X} 後來到了住宿的地方，我們打開袋子，看哪，各人的銀子還在自己的袋口上，銀子的分量一點不少。現在我們親手把它帶回來， \end{tabularx} \\ \\ \relax
43:22 & \begin{tabularx}{0.7\textwidth}{X} 我們手裡又帶了另外的銀子來買糧食。我們不知道是誰把銀子放在我們袋裡的。」 \end{tabularx} \\ \\ \relax
43:23 & \begin{tabularx}{0.7\textwidth}{X} 他說：「你們平安！不要害怕，是你們的神和你們父親的神把財寶放在你們的袋裡。你們的銀子，我已經收了。」他就把西緬帶出來，交給他們。 \end{tabularx} \\ \\ \relax
43:24 & \begin{tabularx}{0.7\textwidth}{X} 那人領這些人進約瑟的屋裡，給他們水洗腳，又給他們飼料餵驢。 \end{tabularx} \\ \\ \relax
43:25 & \begin{tabularx}{0.7\textwidth}{X} 他們預備好禮物，等候約瑟中午來，因為他們聽說他們要在那裡吃飯。 \end{tabularx} \\ \\ \relax
43:26 & \begin{tabularx}{0.7\textwidth}{X} 約瑟來到家裡，他們就把手中的禮物拿進屋裡給他，俯伏在地，向他下拜。 \end{tabularx} \\ \\ \relax
43:27 & \begin{tabularx}{0.7\textwidth}{X} 約瑟問他們安，又說：「你們的父親，就是你們所說的那位老人家平安嗎？他還在嗎？」 \end{tabularx} \\ \\ \relax
43:28 & \begin{tabularx}{0.7\textwidth}{X} 他們說：「你僕人，我們的父親平安，他還在。」於是他們低頭下拜。 \end{tabularx} \\ \\ \relax
43:29 & \begin{tabularx}{0.7\textwidth}{X} 約瑟舉目看見他同母的弟弟便雅憫，就說：「你們向我所說那最小的弟弟就是這位嗎？」又說：「我兒啊，願神賜恩給你！」 \end{tabularx} \\ \\ \relax
43:30 & \begin{tabularx}{0.7\textwidth}{X} 約瑟愛弟之情激動，就急忙找個地方去哭。他進入自己的房間，哭了一場。 \end{tabularx} \\ \\ \relax
43:31 & \begin{tabularx}{0.7\textwidth}{X} 他洗了臉出來，勉強忍住，就說：「開飯吧！」 \end{tabularx} \\ \\ \relax
43:32 & \begin{tabularx}{0.7\textwidth}{X} 他們為約瑟單獨擺了一席，為那些人又擺了一席，也為和約瑟同吃飯的埃及人另擺了一席，因為埃及人不和希伯來人一同吃飯；那是埃及人所厭惡的。 \end{tabularx} \\ \\ \relax
43:33 & \begin{tabularx}{0.7\textwidth}{X} 兄弟們被安排在約瑟面前坐席，都按著長幼的次序，這些人彼此感到詫異。 \end{tabularx} \\ \\ \relax
43:34 & \begin{tabularx}{0.7\textwidth}{X} 約瑟把他面前的食物分給他們，但便雅憫所得的比別人多五倍。他們就喝酒，和約瑟一同暢飲。 \end{tabularx} \\ \\
[1ex]
\hline
\hline
\end{longtable}


\newpage

\section{第70屆~港九培靈會~焦源濂~第八講}
\label{sec:871}
\textbf{一群真弟兄}
\newline
\newline
連結: \href{https://www.hkbibleconference.org/session-message/view/871}{\texttt{https://www.hkbibleconference.org/session-message/view/871}}
\newline
\newline
\hyperref[sec:870]{< < < PREV SERMON < < <}
~
\hyperref[sec:toc]{[返目錄]}
~
\hyperref[sec:872]{> > > NEXT SERMON > > >}
\newline
\newline
創世記 44:1-45:28
\newline
\begin{longtable}{cl}
\hline
\hline
章節 & 經文 (和合本修訂版)\\
\hline
44:1 & \begin{tabularx}{0.7\textwidth}{X} 約瑟吩咐管家說：「按照他們的驢子所能馱的，把這些人的袋子裝滿糧食，再把各人的銀子放在各人的袋口上， \end{tabularx} \\ \\ \relax
44:2 & \begin{tabularx}{0.7\textwidth}{X} 我的杯，就是那個銀杯，要和買糧的銀子一同放在最年輕的那個人的袋口上。」管家就照約瑟所說的話去做了。 \end{tabularx} \\ \\ \relax
44:3 & \begin{tabularx}{0.7\textwidth}{X} 天一亮，這些人和他們的驢子就被送走了。 \end{tabularx} \\ \\ \relax
44:4 & \begin{tabularx}{0.7\textwidth}{X} 他們出城走了不遠，約瑟對管家說：「起來，去追那些人，追上了就對他們說：『你們為甚麼以惡報善呢？ \end{tabularx} \\ \\ \relax
44:5 & \begin{tabularx}{0.7\textwidth}{X} 這不是我主人用來飲酒，確實用它來占卜的嗎？你們這麼做是不對的！』」 \end{tabularx} \\ \\ \relax
44:6 & \begin{tabularx}{0.7\textwidth}{X} 管家追上他們，把這些話對他們說了。 \end{tabularx} \\ \\ \relax
44:7 & \begin{tabularx}{0.7\textwidth}{X} 他們對他說：「我主為甚麼說這樣的話呢？你僕人們絕不會做這樣的事。 \end{tabularx} \\ \\ \relax
44:8 & \begin{tabularx}{0.7\textwidth}{X} 看哪，我們從前在袋口上發現的銀子，尚且從迦南地帶來還你，我們又怎麼會從你主人家裡偷竊金銀呢？ \end{tabularx} \\ \\ \relax
44:9 & \begin{tabularx}{0.7\textwidth}{X} 你僕人中無論在誰那裡找到杯子，就叫他死，我們也要作我主的奴隸。」 \end{tabularx} \\ \\ \relax
44:10 & \begin{tabularx}{0.7\textwidth}{X} 管家說：「現在就照你們的話做吧！在誰那裡找到杯子，誰就作我的奴隸，其餘的人都沒有罪。」 \end{tabularx} \\ \\ \relax
44:11 & \begin{tabularx}{0.7\textwidth}{X} 於是他們各人急忙把袋子卸在地上，各人打開自己的袋子。 \end{tabularx} \\ \\ \relax
44:12 & \begin{tabularx}{0.7\textwidth}{X} 管家就搜查，從年長的開始到年幼的為止，那杯竟在便雅憫的袋子裡找到了。 \end{tabularx} \\ \\ \relax
44:13 & \begin{tabularx}{0.7\textwidth}{X} 他們就撕裂衣服，各人把馱子抬在驢上，回城去了。 \end{tabularx} \\ \\ \relax
44:14 & \begin{tabularx}{0.7\textwidth}{X} 猶大和他兄弟們來到約瑟的屋裡，約瑟還在那裡，他們就在他面前俯伏於地。 \end{tabularx} \\ \\ \relax
44:15 & \begin{tabularx}{0.7\textwidth}{X} 約瑟對他們說：「你們做的是甚麼事呢？你們豈不知像我這樣的人必懂得占卜嗎？」 \end{tabularx} \\ \\ \relax
44:16 & \begin{tabularx}{0.7\textwidth}{X} 猶大說：「我們對我主能說甚麼呢？還有甚麼話可說呢？我們還能為自己表白嗎？神已經查出你僕人的罪孽了。看哪，我們與那在他手中找到杯子的人都是我主的奴隸。」 \end{tabularx} \\ \\ \relax
44:17 & \begin{tabularx}{0.7\textwidth}{X} 約瑟說：「我絕不能做這樣的事！誰的手中找到杯子，誰就作我的奴隸。至於你們，可以平平安安上到你們父親那裡去。」 \end{tabularx} \\ \\ \relax
44:18 & \begin{tabularx}{0.7\textwidth}{X} 猶大挨近他，說：「我主啊，求求你，讓僕人說一句話給我主聽，不要向僕人發烈怒，因為你如同法老一樣。 \end{tabularx} \\ \\ \relax
44:19 & \begin{tabularx}{0.7\textwidth}{X} 我主曾問僕人們說：『你們有父親、兄弟沒有？』 \end{tabularx} \\ \\ \relax
44:20 & \begin{tabularx}{0.7\textwidth}{X} 我們對我主說：『我們有父親，他已經年老，還有他老年所生的一個小兒子。他哥哥死了，他的母親只剩下他一個孩子，父親也疼愛他。』 \end{tabularx} \\ \\ \relax
44:21 & \begin{tabularx}{0.7\textwidth}{X} 你對僕人說：『把他帶下到我這裡來，讓我親眼看看他。』 \end{tabularx} \\ \\ \relax
44:22 & \begin{tabularx}{0.7\textwidth}{X} 我們對我主說：『這年輕人不能離開他父親，若是離開，父親就會死。』 \end{tabularx} \\ \\ \relax
44:23 & \begin{tabularx}{0.7\textwidth}{X} 你對僕人說：『你們最小的弟弟若不和你們一同下來，你們就不要來見我的面。』 \end{tabularx} \\ \\ \relax
44:24 & \begin{tabularx}{0.7\textwidth}{X} 我們上到你僕人，我們父親那裡，就把我主的話告訴了他。 \end{tabularx} \\ \\ \relax
44:25 & \begin{tabularx}{0.7\textwidth}{X} 後來，我們的父親說：『你們再去給我買些糧來。』 \end{tabularx} \\ \\ \relax
44:26 & \begin{tabularx}{0.7\textwidth}{X} 我們說：『我們不能下去。最小的弟弟若和我們同去，我們就可以下去。因為，最小的弟弟若不和我們同去，我們必不能見那人的面。』 \end{tabularx} \\ \\ \relax
44:27 & \begin{tabularx}{0.7\textwidth}{X} 你僕人，我父親對我們說：『你們知道我的妻子給我生了兩個兒子。 \end{tabularx} \\ \\ \relax
44:28 & \begin{tabularx}{0.7\textwidth}{X} 一個離開我走了，我說他必是被野獸撕碎了，直到如今我再沒有見過他； \end{tabularx} \\ \\ \relax
44:29 & \begin{tabularx}{0.7\textwidth}{X} 現在你們又要把這個從我面前帶走。倘若他遭難，那麼你們就害我白髮蒼蒼、悲悲慘慘下陰間去了。』 \end{tabularx} \\ \\ \relax
44:30 & \begin{tabularx}{0.7\textwidth}{X} 如今我回到你僕人，我父親那裡，若沒有這年輕人和我們同去，我父親的命是與這年輕人的命相連的， \end{tabularx} \\ \\ \relax
44:31 & \begin{tabularx}{0.7\textwidth}{X} 當我們的父親看見沒有了這年輕人，他就會死。這樣，我們就害你僕人，我們的父親白髮蒼蒼、悲悲慘慘下陰間去了。 \end{tabularx} \\ \\ \relax
44:32 & \begin{tabularx}{0.7\textwidth}{X} 僕人曾向我父親為這年輕人擔保，說：『我若不帶他回來交給父親，我就在父親面前永遠擔當這罪。』 \end{tabularx} \\ \\ \relax
44:33 & \begin{tabularx}{0.7\textwidth}{X} 現在，求你把僕人留下，代替這年輕人作我主的奴隸，讓這年輕人和他哥哥們一同上去。 \end{tabularx} \\ \\ \relax
44:34 & \begin{tabularx}{0.7\textwidth}{X} 若這年輕人不和我一起，我怎能上到我父親那裡呢？恐怕我要看到災禍臨到我父親了。」 \end{tabularx} \\ \\ \relax
45:1 & \begin{tabularx}{0.7\textwidth}{X} 約瑟在所有侍立在他旁邊的人面前情不自禁，就喊叫說：「每一個人都離開我，出去吧！」約瑟和兄弟相認的時候沒有一人站在他那裡。 \end{tabularx} \\ \\ \relax
45:2 & \begin{tabularx}{0.7\textwidth}{X} 他放聲大哭，埃及人聽見了，法老家中的人也聽見了。 \end{tabularx} \\ \\ \relax
45:3 & \begin{tabularx}{0.7\textwidth}{X} 約瑟對他兄弟們說：「我就是約瑟。我的父親還在嗎？」他兄弟們不敢回答他，因為他們在他面前都很驚惶。 \end{tabularx} \\ \\ \relax
45:4 & \begin{tabularx}{0.7\textwidth}{X} 約瑟又對他兄弟們說：「靠近我一點。」他們就近前來。他說：「我是被你們賣到埃及的兄弟約瑟。 \end{tabularx} \\ \\ \relax
45:5 & \begin{tabularx}{0.7\textwidth}{X} 現在，不要因為把我賣到這裡而憂傷，對自己生氣，因為神差我在你們以先來，為要保全性命。 \end{tabularx} \\ \\ \relax
45:6 & \begin{tabularx}{0.7\textwidth}{X} 現在這地的饑荒已經二年了，還有五年不能耕種，沒有收成。 \end{tabularx} \\ \\ \relax
45:7 & \begin{tabularx}{0.7\textwidth}{X} 神差我在你們以先來，為要給你們在世上存留餘種，大施拯救，保全你們的性命。 \end{tabularx} \\ \\ \relax
45:8 & \begin{tabularx}{0.7\textwidth}{X} 這樣看來，差我到這裡來的不是你們，而是神。他又使我如同法老之父，作他全家之主，和埃及全地掌權的人。 \end{tabularx} \\ \\ \relax
45:9 & \begin{tabularx}{0.7\textwidth}{X} 你們要趕緊上到我父親那裡，對他說：『你兒子約瑟這樣說：神已立我作全埃及之主，請你下到我這裡來，不要耽擱。 \end{tabularx} \\ \\ \relax
45:10 & \begin{tabularx}{0.7\textwidth}{X} 你和你的兒子孫子，羊群牛群，以及一切所有的，都可以住在歌珊地，與我相近。 \end{tabularx} \\ \\ \relax
45:11 & \begin{tabularx}{0.7\textwidth}{X} 我要在那裡奉養你，因為還有五年的饑荒，免得你和你的家屬，以及一切所有的，都陷入窮困中。』 \end{tabularx} \\ \\ \relax
45:12 & \begin{tabularx}{0.7\textwidth}{X} 看哪，你們的眼睛和我弟弟便雅憫的眼睛都看見，是我親口對你們說話。 \end{tabularx} \\ \\ \relax
45:13 & \begin{tabularx}{0.7\textwidth}{X} 你們要把我在埃及一切的尊榮和你們所有看見的事情都告訴我父親，也要趕緊請我父親下到這裡來。」 \end{tabularx} \\ \\ \relax
45:14 & \begin{tabularx}{0.7\textwidth}{X} 於是約瑟伏在他弟弟便雅憫的頸項上哭，便雅憫也在他的頸項上哭。 \end{tabularx} \\ \\ \relax
45:15 & \begin{tabularx}{0.7\textwidth}{X} 他又親眾兄弟，伏著他們哭。過後，他的兄弟就和他說話。 \end{tabularx} \\ \\ \relax
45:16 & \begin{tabularx}{0.7\textwidth}{X} 這消息傳到法老的宮裡，說：「約瑟的兄弟們來了。」法老和他的臣僕眼中都看為好。 \end{tabularx} \\ \\ \relax
45:17 & \begin{tabularx}{0.7\textwidth}{X} 法老對約瑟說：「你要吩咐你的兄弟們說：『你們要這樣做：把馱子抬在牲口上，動身到迦南地去， \end{tabularx} \\ \\ \relax
45:18 & \begin{tabularx}{0.7\textwidth}{X} 請你們的父親和你們的家屬都到我這裡來，我要把埃及地的美物賜給你們，你們也要吃這地肥美的出產。』 \end{tabularx} \\ \\ \relax
45:19 & \begin{tabularx}{0.7\textwidth}{X} 你要吩咐他們：『要這樣做：從埃及地帶著車輛去，把你們的孩子和妻子，以及你們的父親都接來。 \end{tabularx} \\ \\ \relax
45:20 & \begin{tabularx}{0.7\textwidth}{X} 你們的眼不要顧惜你們的家具，因為埃及全地的美物都是你們的。』」 \end{tabularx} \\ \\ \relax
45:21 & \begin{tabularx}{0.7\textwidth}{X} 以色列的兒子們就照樣做了。約瑟遵照法老的吩咐，給他們車輛和路上需用的食物。 \end{tabularx} \\ \\ \relax
45:22 & \begin{tabularx}{0.7\textwidth}{X} 他又給所有哥哥每人一套衣服， 卻給便雅憫三百銀子，五套衣服。 \end{tabularx} \\ \\ \relax
45:23 & \begin{tabularx}{0.7\textwidth}{X} 他也送給父親十匹公驢，馱著埃及的美物，以及十匹母驢，馱著給他父親在路上需用的穀物、餅和糧食。 \end{tabularx} \\ \\ \relax
45:24 & \begin{tabularx}{0.7\textwidth}{X} 於是約瑟送他的兄弟們回去，對他們說：「你們不要在路上爭吵。」 \end{tabularx} \\ \\ \relax
45:25 & \begin{tabularx}{0.7\textwidth}{X} 他們從埃及上去，來到迦南地他們的父親雅各那裡， \end{tabularx} \\ \\ \relax
45:26 & \begin{tabularx}{0.7\textwidth}{X} 告訴他說：「約瑟還活著，並且作了埃及全地掌權的人。」雅各心裡冰涼，因為不信他們。 \end{tabularx} \\ \\ \relax
45:27 & \begin{tabularx}{0.7\textwidth}{X} 他們就把約瑟對他們所說一切的話都告訴了他。他看見約瑟派來接他的車輛，他們父親雅各的靈就甦醒了。 \end{tabularx} \\ \\ \relax
45:28 & \begin{tabularx}{0.7\textwidth}{X} 以色列說：「夠了！我的兒子約瑟還活著，我要趁我未死之前去見他。」 \end{tabularx} \\ \\
[1ex]
\hline
\hline
\end{longtable}


\newpage

\section{第70屆~港九培靈會~焦源濂~第九講}
\label{sec:872}
\textbf{甚麼是真的成就?}
\newline
\newline
連結: \href{https://www.hkbibleconference.org/session-message/view/872}{\texttt{https://www.hkbibleconference.org/session-message/view/872}}
\newline
\newline
\hyperref[sec:871]{< < < PREV SERMON < < <}
~
\hyperref[sec:toc]{[返目錄]}
~
\hyperref[sec:873]{> > > NEXT SERMON > > >}
\newline
\newline
創世記 47:1-47:31
\newline
\begin{longtable}{cl}
\hline
\hline
章節 & 經文 (和合本修訂版)\\
\hline
47:1 & \begin{tabularx}{0.7\textwidth}{X} 約瑟進去告訴法老說：「我的父親和我的兄弟帶著羊群牛群，以及他們一切所有的，從迦南地來了。看哪，他們正在歌珊地。」 \end{tabularx} \\ \\ \relax
47:2 & \begin{tabularx}{0.7\textwidth}{X} 約瑟從他所有兄弟中挑選五個人，引他們到法老面前。 \end{tabularx} \\ \\ \relax
47:3 & \begin{tabularx}{0.7\textwidth}{X} 法老對約瑟的兄弟說：「你們是做甚麼的？」他們對法老說：「你僕人是牧羊的，我們和我們的祖宗都是這樣。」 \end{tabularx} \\ \\ \relax
47:4 & \begin{tabularx}{0.7\textwidth}{X} 他們又對法老說：「迦南地的饑荒非常嚴重，僕人的羊群沒有牧草，所以我們來到這地寄居。現在求你准許僕人住在歌珊地。」 \end{tabularx} \\ \\ \relax
47:5 & \begin{tabularx}{0.7\textwidth}{X} 法老對約瑟說：「你的父親和你的兄弟到你這裡來了， \end{tabularx} \\ \\ \relax
47:6 & \begin{tabularx}{0.7\textwidth}{X} 埃及地都在你面前，只管讓你父親和你兄弟住在最好的地，他們可以住在歌珊地。你若知道他們中間有能幹的人，就派他們看管我的牲畜。」 \end{tabularx} \\ \\ \relax
47:7 & \begin{tabularx}{0.7\textwidth}{X} 約瑟帶他父親雅各來，站在法老面前，雅各就為法老祝福。 \end{tabularx} \\ \\ \relax
47:8 & \begin{tabularx}{0.7\textwidth}{X} 法老對雅各說：「你平生的年日是多少呢？」 \end{tabularx} \\ \\ \relax
47:9 & \begin{tabularx}{0.7\textwidth}{X} 雅各對法老說：「我在世寄居的年日是一百三十年，我一生的歲月又短又苦，比不上我祖先在世寄居的年日。」 \end{tabularx} \\ \\ \relax
47:10 & \begin{tabularx}{0.7\textwidth}{X} 雅各又為法老祝福，就從法老面前退出去了。 \end{tabularx} \\ \\ \relax
47:11 & \begin{tabularx}{0.7\textwidth}{X} 約瑟安頓他的父親和兄弟，遵照法老的命令，把埃及境內最好的地，就是蘭塞地，給他們作為產業。 \end{tabularx} \\ \\ \relax
47:12 & \begin{tabularx}{0.7\textwidth}{X} 約瑟用糧食供給他父親和兄弟們，以及他父親全家的人，照扶養親屬的人口供給。 \end{tabularx} \\ \\ \relax
47:13 & \begin{tabularx}{0.7\textwidth}{X} 饑荒非常嚴重，全地都絕了糧，埃及地和迦南地都因饑荒耗損了。 \end{tabularx} \\ \\ \relax
47:14 & \begin{tabularx}{0.7\textwidth}{X} 約瑟收集了埃及地和迦南地所有的銀子，就是眾人買糧的銀子，約瑟就把那些銀子都帶到法老的宮裡。 \end{tabularx} \\ \\ \relax
47:15 & \begin{tabularx}{0.7\textwidth}{X} 埃及地和迦南地的銀子都花光了，埃及眾人到約瑟那裡，說：「我們的銀子都用完了，求你給我們糧食吧！我們為甚麼要死在你面前呢？」 \end{tabularx} \\ \\ \relax
47:16 & \begin{tabularx}{0.7\textwidth}{X} 約瑟說：「銀子若是用完了，可以把你們的牲畜賣給我，我就以你們的牲畜換糧食給你們。」 \end{tabularx} \\ \\ \relax
47:17 & \begin{tabularx}{0.7\textwidth}{X} 於是他們把牲畜帶到約瑟那裡，約瑟就拿糧食換了他們的馬、羊、牛、驢；那一年他因換他們一切的牲畜，用糧食養活他們。 \end{tabularx} \\ \\ \relax
47:18 & \begin{tabularx}{0.7\textwidth}{X} 那一年過去，第二年他們又來到約瑟那裡，對他說：「不瞞我主，我們的銀子都花光了，牲畜也都歸於我主了。我們在我主面前，除了自己的身體和土地以外，一無所剩。 \end{tabularx} \\ \\ \relax
47:19 & \begin{tabularx}{0.7\textwidth}{X} 你為甚麼要眼看著我們人死地荒呢？求你用糧食買我們和我們的地，我們和我們的地就要為法老效力。求你給我們種子，使我們可以存活，不致死亡，土地也不致荒蕪。」 \end{tabularx} \\ \\ \relax
47:20 & \begin{tabularx}{0.7\textwidth}{X} 於是，約瑟為法老買了埃及所有的土地，埃及人因饑荒所迫，都賣了自己的田地；那些地都歸給法老了。 \end{tabularx} \\ \\ \relax
47:21 & \begin{tabularx}{0.7\textwidth}{X} 至於百姓，從埃及邊界的一端到另一端，約瑟使他們作奴隸。 \end{tabularx} \\ \\ \relax
47:22 & \begin{tabularx}{0.7\textwidth}{X} 只有祭司的土地，約瑟沒有買，因為祭司從法老領取薪俸，靠法老的薪俸過活，所以沒有賣自己的土地。 \end{tabularx} \\ \\ \relax
47:23 & \begin{tabularx}{0.7\textwidth}{X} 約瑟對百姓說：「看哪，我今日為法老買了你們和你們的土地。看哪，這些種子是給你們的，你們可以耕種土地。 \end{tabularx} \\ \\ \relax
47:24 & \begin{tabularx}{0.7\textwidth}{X} 將來收割的時候，你們要把五分之一納給法老，另外四分可以給你們作田地的種子，作你們和你們全家大小的食物。」 \end{tabularx} \\ \\ \relax
47:25 & \begin{tabularx}{0.7\textwidth}{X} 他們說：「你救了我們的性命，願我們在我主眼前蒙恩，我們情願作法老的奴隸。」 \end{tabularx} \\ \\ \relax
47:26 & \begin{tabularx}{0.7\textwidth}{X} 於是約瑟為埃及的土地立下定例，直到今日，就是收成的五分之一要歸法老。惟獨祭司的土地例外，不歸於法老。 \end{tabularx} \\ \\ \relax
47:27 & \begin{tabularx}{0.7\textwidth}{X} 以色列人住在埃及境內的歌珊地。他們在那裡得了產業，並且生養眾多。 \end{tabularx} \\ \\ \relax
47:28 & \begin{tabularx}{0.7\textwidth}{X} 雅各住在埃及地十七年。雅各一生的年日是一百四十七年。 \end{tabularx} \\ \\ \relax
47:29 & \begin{tabularx}{0.7\textwidth}{X} 以色列的死期快到了，就叫了他兒子約瑟來，對他說：「我若在你眼前蒙恩，把你的手放在我大腿底下，以慈愛和誠實向我承諾，必不將我葬在埃及。 \end{tabularx} \\ \\ \relax
47:30 & \begin{tabularx}{0.7\textwidth}{X} 我與我祖先同睡的時候，你要將我帶出埃及，把我葬在他們所葬的地方。」約瑟說：「我必遵照你的吩咐去做。」 \end{tabularx} \\ \\ \relax
47:31 & \begin{tabularx}{0.7\textwidth}{X} 雅各說：「你向我起誓吧！」約瑟就向他起了誓。於是以色列在床頭敬拜。 \end{tabularx} \\ \\
[1ex]
\hline
\hline
\end{longtable}


\newpage

\section{第70屆~港九培靈會~焦源濂~第十講}
\label{sec:873}
\textbf{神的意思原是好的}
\newline
\newline
連結: \href{https://www.hkbibleconference.org/session-message/view/873}{\texttt{https://www.hkbibleconference.org/session-message/view/873}}
\newline
\newline
\hyperref[sec:872]{< < < PREV SERMON < < <}
~
\hyperref[sec:toc]{[返目錄]}
~
\hyperref[sec:741]{> > > NEXT SERMON > > >}
\newline
\newline
創世記 50:1-50:26
\newline
\begin{longtable}{cl}
\hline
\hline
章節 & 經文 (和合本修訂版)\\
\hline
50:1 & \begin{tabularx}{0.7\textwidth}{X} 約瑟伏在他父親的臉上，在他臉上哭，又親他。 \end{tabularx} \\ \\ \relax
50:2 & \begin{tabularx}{0.7\textwidth}{X} 約瑟吩咐伺候他的醫生們用香料塗他父親，醫生就用香料塗了以色列。 \end{tabularx} \\ \\ \relax
50:3 & \begin{tabularx}{0.7\textwidth}{X} 四十天滿了，就是塗香料所規定的日子滿了。埃及人為他哀哭了七十天。 \end{tabularx} \\ \\ \relax
50:4 & \begin{tabularx}{0.7\textwidth}{X} 過了哀悼的日子，約瑟對法老家中的人說：「我若在你們眼前蒙恩，請你們對法老說： \end{tabularx} \\ \\ \relax
50:5 & \begin{tabularx}{0.7\textwidth}{X} 『我父親曾叫我起誓說：看哪，我快要死了，你要將我葬在迦南地，在我為自己所掘的墳墓裡。』現在求你准我上去葬我父親，然後我必回來。」 \end{tabularx} \\ \\ \relax
50:6 & \begin{tabularx}{0.7\textwidth}{X} 法老說：「你可以上去，照你父親叫你起的誓，將他安葬。」 \end{tabularx} \\ \\ \relax
50:7 & \begin{tabularx}{0.7\textwidth}{X} 於是約瑟上去葬他父親。與他一同上去的有法老的眾臣僕和法老家中的長老，以及埃及地所有的長老， \end{tabularx} \\ \\ \relax
50:8 & \begin{tabularx}{0.7\textwidth}{X} 還有約瑟的全家和他的兄弟們，以及他父親的家屬；只留下他們的孩子和羊群牛群在歌珊地。 \end{tabularx} \\ \\ \relax
50:9 & \begin{tabularx}{0.7\textwidth}{X} 又有車輛和駕駛兵和他一同上去，隊伍非常龐大。 \end{tabularx} \\ \\ \relax
50:10 & \begin{tabularx}{0.7\textwidth}{X} 他們到了約旦河東亞達的禾場，就在那裡大大地號咷痛哭。約瑟為他父親哀哭了七天。 \end{tabularx} \\ \\ \relax
50:11 & \begin{tabularx}{0.7\textwidth}{X} 迦南的居民看見亞達禾場上的哀哭，就說：「這是埃及人一場極大的哀哭。」因此那地方名叫亞伯‧麥西，是在約旦河東。 \end{tabularx} \\ \\ \relax
50:12 & \begin{tabularx}{0.7\textwidth}{X} 雅各的兒子們遵照父親的吩咐去辦了， \end{tabularx} \\ \\ \relax
50:13 & \begin{tabularx}{0.7\textwidth}{X} 他們把他送到迦南地，葬在幔利對面的麥比拉田間的洞裡；那田是亞伯拉罕向赫人以弗崙買來作墳地的產業。 \end{tabularx} \\ \\ \relax
50:14 & \begin{tabularx}{0.7\textwidth}{X} 約瑟葬了他父親以後，就和他的兄弟，以及所有同他上去葬他父親的人，都回埃及去了。 \end{tabularx} \\ \\ \relax
50:15 & \begin{tabularx}{0.7\textwidth}{X} 約瑟的哥哥們見父親死了，就說：「也許約瑟仍然懷恨我們，會照我們從前待他一切的惡，重重報復我們。」 \end{tabularx} \\ \\ \relax
50:16 & \begin{tabularx}{0.7\textwidth}{X} 他們就傳口信給約瑟說：「你父親未死之前曾吩咐說： \end{tabularx} \\ \\ \relax
50:17 & \begin{tabularx}{0.7\textwidth}{X} 『你們要對約瑟這樣說：從前你哥哥們惡待你，你要饒恕他們的過犯和罪惡。』現在求你饒恕你父親的神之僕人們的過犯。」他們對約瑟說了這話，約瑟就哭了。 \end{tabularx} \\ \\ \relax
50:18 & \begin{tabularx}{0.7\textwidth}{X} 他的哥哥們又來俯伏在他面前，說：「看哪，我們是你的奴隸。」 \end{tabularx} \\ \\ \relax
50:19 & \begin{tabularx}{0.7\textwidth}{X} 約瑟對他們說：「不要怕，我豈能代替神呢？ \end{tabularx} \\ \\ \relax
50:20 & \begin{tabularx}{0.7\textwidth}{X} 從前你們的意思是要害我，但神的意思原是好的，要使許多百姓得以存活，成就今日的光景。 \end{tabularx} \\ \\ \relax
50:21 & \begin{tabularx}{0.7\textwidth}{X} 現在你們不要害怕，我必養活你們和你們的孩子。」於是約瑟安慰他們，講了使他們安心的話。 \end{tabularx} \\ \\ \relax
50:22 & \begin{tabularx}{0.7\textwidth}{X} 約瑟和他父親的家屬都住在埃及。約瑟活了一百一十年。 \end{tabularx} \\ \\ \relax
50:23 & \begin{tabularx}{0.7\textwidth}{X} 約瑟看到以法蓮第三代的子孫。瑪拿西的孫子，瑪吉的兒子，出生時都放在約瑟的膝上。 \end{tabularx} \\ \\ \relax
50:24 & \begin{tabularx}{0.7\textwidth}{X} 約瑟對他的兄弟說：「我快要死了，但神必定看顧你們，領你們從這地上去，到他起誓應許給亞伯拉罕、以撒、雅各之地。」 \end{tabularx} \\ \\ \relax
50:25 & \begin{tabularx}{0.7\textwidth}{X} 約瑟叫以色列的子孫起誓：「神必定眷顧你們，你們要把我的骸骨從這裡帶上去。」 \end{tabularx} \\ \\ \relax
50:26 & \begin{tabularx}{0.7\textwidth}{X} 約瑟死了，那時他一百一十歲。人用香料塗了他，把他收殮在棺材裡，停放在埃及。 \end{tabularx} \\ \\
[1ex]
\hline
\hline
\end{longtable}


\newpage

\section{第71屆~港九培靈會~劉承業~第一講}
\label{sec:741}
\textbf{詩人為何要讚美?}
\newline
\newline
連結: \href{https://www.hkbibleconference.org/session-message/view/741}{\texttt{https://www.hkbibleconference.org/session-message/view/741}}
\newline
\newline
\hyperref[sec:873]{< < < PREV SERMON < < <}
~
\hyperref[sec:toc]{[返目錄]}
~
\hyperref[sec:742]{> > > NEXT SERMON > > >}
\newline
\newline


\newpage

\section{第71屆~港九培靈會~劉承業~第二講}
\label{sec:742}
\textbf{「讚美詩」的讚美}
\newline
\newline
連結: \href{https://www.hkbibleconference.org/session-message/view/742}{\texttt{https://www.hkbibleconference.org/session-message/view/742}}
\newline
\newline
\hyperref[sec:741]{< < < PREV SERMON < < <}
~
\hyperref[sec:toc]{[返目錄]}
~
\hyperref[sec:743]{> > > NEXT SERMON > > >}
\newline
\newline


\newpage

\section{第71屆~港九培靈會~劉承業~第三講}
\label{sec:743}
\textbf{「感恩詩」的讚美}
\newline
\newline
連結: \href{https://www.hkbibleconference.org/session-message/view/743}{\texttt{https://www.hkbibleconference.org/session-message/view/743}}
\newline
\newline
\hyperref[sec:742]{< < < PREV SERMON < < <}
~
\hyperref[sec:toc]{[返目錄]}
~
\hyperref[sec:744]{> > > NEXT SERMON > > >}
\newline
\newline


\newpage

\section{第71屆~港九培靈會~劉承業~第四講}
\label{sec:744}
\textbf{「個人哀歌」的讚美}
\newline
\newline
連結: \href{https://www.hkbibleconference.org/session-message/view/744}{\texttt{https://www.hkbibleconference.org/session-message/view/744}}
\newline
\newline
\hyperref[sec:743]{< < < PREV SERMON < < <}
~
\hyperref[sec:toc]{[返目錄]}
~
\hyperref[sec:745]{> > > NEXT SERMON > > >}
\newline
\newline


\newpage

\section{第71屆~港九培靈會~劉承業~第五講}
\label{sec:745}
\textbf{「團體哀歌」的讚美}
\newline
\newline
連結: \href{https://www.hkbibleconference.org/session-message/view/745}{\texttt{https://www.hkbibleconference.org/session-message/view/745}}
\newline
\newline
\hyperref[sec:744]{< < < PREV SERMON < < <}
~
\hyperref[sec:toc]{[返目錄]}
~
\hyperref[sec:746]{> > > NEXT SERMON > > >}
\newline
\newline


\newpage

\section{第71屆~港九培靈會~劉承業~第六講}
\label{sec:746}
\textbf{苦難中的讚美}
\newline
\newline
連結: \href{https://www.hkbibleconference.org/session-message/view/746}{\texttt{https://www.hkbibleconference.org/session-message/view/746}}
\newline
\newline
\hyperref[sec:745]{< < < PREV SERMON < < <}
~
\hyperref[sec:toc]{[返目錄]}
~
\hyperref[sec:747]{> > > NEXT SERMON > > >}
\newline
\newline


\newpage

\section{第71屆~港九培靈會~劉承業~第七講}
\label{sec:747}
\textbf{逆境中的讚美}
\newline
\newline
連結: \href{https://www.hkbibleconference.org/session-message/view/747}{\texttt{https://www.hkbibleconference.org/session-message/view/747}}
\newline
\newline
\hyperref[sec:746]{< < < PREV SERMON < < <}
~
\hyperref[sec:toc]{[返目錄]}
~
\hyperref[sec:748]{> > > NEXT SERMON > > >}
\newline
\newline


\newpage

\section{第71屆~港九培靈會~劉承業~第八講}
\label{sec:748}
\textbf{如何實踐讚美}
\newline
\newline
連結: \href{https://www.hkbibleconference.org/session-message/view/748}{\texttt{https://www.hkbibleconference.org/session-message/view/748}}
\newline
\newline
\hyperref[sec:747]{< < < PREV SERMON < < <}
~
\hyperref[sec:toc]{[返目錄]}
~
\hyperref[sec:24]{> > > NEXT SERMON > > >}
\newline
\newline


\newpage

\section{第71屆~港九培靈會~滕近輝~第一講}
\label{sec:24}
\textbf{真救主:七證}
\newline
\newline
連結: \href{https://www.hkbibleconference.org/session-message/view/24}{\texttt{https://www.hkbibleconference.org/session-message/view/24}}
\newline
\newline
\hyperref[sec:748]{< < < PREV SERMON < < <}
~
\hyperref[sec:toc]{[返目錄]}
~
\hyperref[sec:23]{> > > NEXT SERMON > > >}
\newline
\newline
約翰福音 4:39-4:42
\newline
\begin{longtable}{cl}
\hline
\hline
章節 & 經文 (和合本修訂版)\\
\hline
4:39 & \begin{tabularx}{0.7\textwidth}{X} 那城裡有好些撒瑪利亞人信了耶穌，因為那婦人作見證，說：「他把我素來所做的一切事都說了出來。」 \end{tabularx} \\ \\ \relax
4:40 & \begin{tabularx}{0.7\textwidth}{X} 於是撒瑪利亞人來見耶穌，求他在他們那裡住下，他就在那裡住了兩天。 \end{tabularx} \\ \\ \relax
4:41 & \begin{tabularx}{0.7\textwidth}{X} 因為耶穌的話，信的人就更多了。 \end{tabularx} \\ \\ \relax
4:42 & \begin{tabularx}{0.7\textwidth}{X} 他們對那婦人說：「現在我們信，不再是因為你的話，而是我們親自聽見了，知道這人真是世界的救主。」 \end{tabularx} \\ \\
[1ex]
\hline
\hline
\end{longtable}
剛才我從台上望下去,立刻留意到一位頭髮全白了的姊妹,相信她參加港九培靈研經會四十多年了!我是一九五一年開始參加的,在大會中見到過三代的年輕人.我可以說是大會的見證人之一,我要見證一件事:歲月流逝,但《聖經》的道沒有改變,培靈會的信息也沒有改變!因為人的需要也沒有改變!在這四十九年中,我親眼看見很多很多的青年人奉獻自己給主,感謝神!我們所傳的信息不需要改變,因為神的道是萬古長新的!五十年來,社會的情況改變很多,思潮改變很多,環境改變很多,但是我們仍然講同樣的信息,因為可以滿足人心靈的基本的,最重要的需要!五十年在歷史中算不得甚麼,事實上二千年來主的道沒有改變,古老的福音詩歌,古老的福音歷久常新!二千年來,人類改變了多少?社會改變了多少?文化變了多少?理論改變了多少?多少的理論興起又消退,起來又過去,起來又過去?但只有古老的《聖經》仍然是生命之道,我們今天仍然講二千年前主耶穌所講的道,仍然講千年前保羅所講的道!我們發現這道仍然是今天的人所要的!從講台上望下去,青年人所佔的比例相當高,但願主得到這一代的基督徒!興起來!信主,愛主,見證主!作主賜福的器皿!
\newline
\newline
今天我們要從《約翰福音》中領受信息,《約》中有一個字用了三十次,那就是「真」字.用了這麼多次,表示《約》很注重「真」這個信息；再就是「實實在在告訴你們」,也用了三十次；「見證」兩字用了四十二次之多.這三個字眼基本的意義都一樣:「見證是真的,是神實實在在告訴我們的真道」,這就是《約》的中心信息.「知道這真是救世主」(約四42),耶穌是真的救主,我們都可以從心裡作見證,我們信祂,因祂是真的救世主!我們對救主的信心有把握嗎?有的基督徒信了,但卻半信半疑；有的信了,但對信仰還沒有把握,又講不出確實的憑據來!今晚我們就要看看耶穌基督是真救主的可信的憑據.
\newline
\newline
《約》中常講到「你們可以信」,那就是可以真正相信的理由.(一7)中「這人」即施洗約翰,他是一個品格高尚的人,為真理犧牲了自己的生命.當時希律王娶了自己的弟婦為妻,一般的人敢怒不敢言,約翰起來當面指責希律王:「你娶你兄弟的妻子是不合理的.」(可六18)希律王是一個暴君,他知道這樣做等於送命,但仍然很坦誠地,很勇敢地把真理講出來；結果希律發怒將他斬首,他是為真理殉道的一個人.這樣的人為耶穌作見證,他的見證應該是可信的!
\newline
\newline
另一個同樣性質的見證就是耶穌基督的母親馬利亞.她也是信耶穌的,相信耶穌是無罪的救主!大家想想:知子莫如母,馬利亞對她的兒子耶穌認識得很清楚,看到祂從小長大,如果祂有問題,有罪,有過錯,那麼她不會信耶穌是無罪的救主!但她信,就表示她沒有發現耶穌有罪!如果耶穌有罪,馬利亞可以說:「你們不要信祂了,因為你們認識祂不清楚；我認識得清楚,我才不信呢!」馬利亞信耶穌,這件事實是一個很強的見證!「因耶穌的話,信的人就更多了.」
\newline
\newline
(約四41)人們怎知耶穌的話是真的呢?是藉著一個撒馬利亞婦人的改變!這婦人是一個墮落的女人,曾有過五個丈夫,當時的恐怕是第六個；這樣的一個女人竟然悔改了!看見她的改變,聽見她訴說為甚麼改變:因她聽見主講生命活水的道!她就出來作見證,那些聽見的人受感動,他們在想耶穌講的話,就了解了!原來他們聽見不信,但現看見一個信的人之人生改變,這樣奇妙地,徹底地大改變,他們受了感動!就信了耶穌!他們信耶穌的話有一個過程:原來聽見不信,但經過撒馬利亞婦人的見證之後,再想,再聽,才相信；這婦人改變的見證使他們深深地思想:耶穌基督可信嗎?祂是救主嗎?得到信仰的結論:是的,祂是救主!因為祂拯救,改變了世界上千千萬萬的罪人!這是何等奇妙的事實!
\newline
\newline
大家都知道「救世軍」的創辦人是布威廉,他是一個很注重傳福音,佈道工作的人.他的佈道會作法特別:他有一個樂隊,每個隊員必須具備一條件:原是墮落的人,信主之後人生改變.每次傳福音時樂隊演奏,演奏前,演奏中,及演奏完畢分別由一個隊員作見證,聽到的人深深感受到:這樣墮落的人能改變嗎?竟然能改變!聽到這活生生見證的人紛紛感動悔改信主,使福音的工作愈來愈蒙恩.這情形二千年來在世界各地常常不斷發生,這是心靈改變!這是人生的改變!這是因為聽見了主耶穌所講的生命之道而接受!
\newline
\newline
主耶穌說,《舊約聖經》是祂的見證(約五39),因為在《舊約》中有許多關於祂的預言,這些預言都應驗了!實在令我們感到驚奇,《舊約》的預言真是奇妙!就(賽五三9)來說吧,這裡預言主耶穌基督死後葬在甚麼地方,不可理解的是此處提到兩個地方:其一,與惡人同埋,其二,誰知死時與財主同葬.怎麼能這樣呢?一個人死後怎能葬在兩地呢?除非財主即惡人,惡人即財主；但事實上卻不是,原來這預言的應驗有一個很奇妙的過程.今日基督徒都知道,原來耶穌和兩個強盜同釘十字架,本來是三個人死後葬在一起:本來與惡人同埋；誰知出人意料的事發生了,一個財主,名叫約瑟,他原來暗暗信耶穌,看見主耶穌釘十字架的情景深受感動,他的信仰堅定了,毅然站出來,請求將耶穌的身體領走,葬在他家族的墳墓中:與財主同葬.我們看到事實的曲折發展都預言在這裡,仔細想想是不是覺得很希奇?!《舊約》的預言令我們驚奇.請再看(12下)的另一個預言:這裡講到救主將祂的生命好像水從水罐裡倒出來一樣；這個預言同樣應驗了,當祂釘十字架時身體有七處流血,主耶穌生命的血流完了,將生命倒空,以致於死.又講到主耶穌被釘死在兩個強盜之間,卻為罪犯代求；我們都記得主耶穌在十字架上為那兩強盜的禱告.請看,發生在十字架上的情形這裡講得多清楚!好像事後寫的一樣!這是神事先的預言!所以耶穌基督釘十字架是神恩典救贖的計劃,一步一步地成就應驗!不是偶然的,我們看見神的愛,神的心思在其中.
\newline
\newline
(約七31)中說到有許多的人信了耶穌,因為祂行神蹟!關於主所行的神蹟,我提出以下兩點:
\newline
\newline
1．主耶穌從來沒有藉神蹟作宣揚的意思,不是藉著神蹟吸引人注意,每次耶穌行神蹟都是出於憐憫,出於愛心,是出於愛的神蹟.
\newline
\newline
2．神蹟即神的作為,神能創造這偉大的宇宙,難道不能行神蹟嗎?事實上宇宙是一個大神蹟,人的身體亦是神蹟中的神蹟!人身體的構造太奇妙了,說它是神蹟,一點也不言過其實!整個的宇宙都是神蹟,如果神創造了宇宙,祂在歷史中行幾個神蹟算得了甚麼呢?
\newline
\newline
所以,神能行神蹟是一個很合理的結論.耶穌基督在世上用愛心行了很多神蹟,許多人看見受感動,產生了對主耶穌的信心,今日世界上許多第一流的科學家也相信神蹟.
\newline
\newline
我們再看另一處主耶穌是救主的證據:(八29-30)中說到有許多人信祂,為甚麼呢?因為主耶穌常做神喜悅的事,即合乎倫理道德的事,善事.神的本性即真理,仁義,聖潔,那裡有這樣的善事,神就喜悅.換句話說,耶穌基督在歷史裡面是一個道德最高尚的人.德國文豪歌德說:在《四福音》裡所反映出耶穌基督之高尚品格,崇高道德,是整個人類歷史中獨特的,相信將來沒有人能超過的.耶穌主張:「你們願意人怎樣待你們,你們也要怎樣待人.」(路六31)莫等別人做了你才做,你先自己做到.這是創造性的道德!主耶穌本人的道德使祂成為可信仰的救主,難怪當時許多人聽見,看見他表現就信了.
\newline
\newline
主耶穌說:「如今事情還沒有成就,我要先告訴你們,叫你們到事情成就的時候可以信我是基督.」(約十三19),祂把將來要發生的事預先說明,到時候大家看見預言成就便可以相信.前面我們看到《舊約》中的預言奇妙地應驗,在這裡我們看到主耶穌自己說預言；主耶穌說到祂被釘十字架時,就要吸引萬人來歸向祂(十二32),那時祂已面對十字架,但沒有人知道結果怎樣,但主耶穌預先暗示出來:十字架要成為許多人的救恩!誰能這樣預言呢?被釘死在十字架上,在人眼中是一件完全失敗的事,但主耶穌說祂被釘之時,就吸引世上千千萬萬的人來跟從祂!這是講誑語嗎?不!我們今天看見耶穌基督的預言一點沒錯!
\newline
\newline
「這天國的福音要傳遍天下,對萬民作見證,然後末期纔來到.」(太廿四14)福音傳遍天下就是基督再來的日子,基督救恩的福音要傳遍天下.這個加利利人耶穌,這個木匠之子耶穌,這個沒有地位,學問的耶穌,這個極其平凡的耶穌,竟然宣告:祂的道理要傳遍天下!如果主耶穌基督沒有事先預言,我們可以說世界上無論何事都可發生,主耶穌事先預言祂的道,祂的天國福音要傳遍天下!今天我們親眼看見,應驗了!
\newline
\newline
這些證據對我個人來講都是很強的證據,堅定了我的信仰.希望在座的每一位弟兄姊妹都有堅強的信心,都能從心底說:「我深知我信的是誰.」我們的信仰很確定,在任何時代,任何潮流,任何沖激之下都能站立得住:我是信耶穌的,我跟從耶穌,我見證耶穌!傳揚救恩,願使多人得救,耶穌是真救主!
\newline
\newline


\newpage

\section{第71屆~港九培靈會~滕近輝~第二講}
\label{sec:23}
\textbf{真葡萄樹}
\newline
\newline
連結: \href{https://www.hkbibleconference.org/session-message/view/23}{\texttt{https://www.hkbibleconference.org/session-message/view/23}}
\newline
\newline
\hyperref[sec:24]{< < < PREV SERMON < < <}
~
\hyperref[sec:toc]{[返目錄]}
~
\hyperref[sec:22]{> > > NEXT SERMON > > >}
\newline
\newline
一九五二年,我第一次在港九培靈研經會上講道,三十歲的我,當時是最年青的講員,而今次我卻是最年長的講員!人生真是過得很快啊!我藉此提醒在座的年青的弟兄姊妹們:趁早把自己當作活祭,獻上給神!開始為主而活,跟從,見證,事奉,愛主!你們還有很多年可以為主而活,那是何等地美好!見主面時,也是何等的快樂!五十年前我參加在英格蘭南部舉行的一個有五千人的大型英文培靈會(其中約有一半是大學生),這個培靈會繼續到現在已有一百二十二年的歷史了.二十年前我再去時人數已增加到七千人,青年人更多,在聚會中有許多人獻身給主.最近一次在美國舉行的大型的學生宣教大會有一萬九千大學生從不同的地方到那裡聚集,聖靈大大作工,有五千一百多位大學生獻身作宣教士!我們常說今天這時代青年人很難信主,就算信了也很難熱心；但事實恰好相反,去年底,美國一出名差傳刊物報到:在美國十幾間大學裡學生們出現了一個復興的現象,上千的大學生在一起懇切,整晚禱告,這是神藉著聖靈的工作,是時代的潮流所不能攔阻的!所以我們對香港的青年弟兄姊妹寄予很大的期望,神一定會更多,更深的在青年人中間作工.青年,壯年的弟兄姊妹們,主耶穌寶貴你們,要用你們,願主得著你們!求神興起新一代的基督徒來為祂大有作為!我們等待復興來到香港.
\newline
\newline
(約十五1)中為甚麼要用這「真」字呢?原因:
\newline
\newline
1．當時在聖殿的門口雕刻了葡萄樹,很多猶太人在門口出出入入,看慣了那雕刻.主耶穌說:那只是一個象徵,而我是真葡萄樹.
\newline
\newline
2．先知以賽亞說到以色列人好像一顆壞的葡萄樹,不能結好葡萄,乃是結了野葡萄(賽五4),這是以色列人讓神失望,傷心的事,是以色列的失敗!
\newline
\newline
主耶穌針對當時他們處於律法時代,沒有得到真正的屬靈的生命的那情況,說:「我是真葡萄樹,你們是枝子.」這是多麼寶貴的一個比喻:枝子連在葡萄樹上.主耶穌提到三種枝子:
\newline
\newline
1．不結果子的枝子(約十五2):這是指掛名的,沒有主的生命的,沒有重生得救的基督徒,這樣的人和主沒有生命的聯繫,故不能結果子!這樣的枝子主會剪去.教會裡也有這樣的人,他們的名字在會友的名冊上,但沒有在羔羊的生命冊上,他們與主的救恩沒有關係!這是可悲的,所以我們每一個教會的會友要自我審查:我是否一個掛名的,表面屬主而沒有結果子的基督徒呢?我們經過審查後要醒悟過來,然後真正重新得救,這才是屬靈真境的開始!我們有一位真救主,祂賜給我們真生命.我們的靈魂得救,得著永遠的生命,才能談到屬靈的真境.
\newline
\newline
2．多結果子的枝子(5):因為它們連在葡萄樹上,它們常在主裡面,主常在它們裡面.「常」和「裡面」這兩個字眼在這章經文裡各用了十二次；「裡面」是生命的聯繫,生命的交通,生命的交流.一個人內在的屬靈的關係非常密切,這人就能多結果子.主是能源,離開了主,我們不能結屬靈的果子!枝子若不連在葡萄樹上,就不能結果子!
\newline
\newline
3．結果子更多的枝子(2下):我們看到神把枝子修理乾淨,將枝子上的雜枝剪掉,因雜枝會消耗枝子的生命汁漿,生命汁漿消耗掉就不能結果子,有經驗的園丁一定要常常清理雜枝,讓樹的生命汁漿集中在結果子上.主耶穌要我們人生有一個生命的焦點,然後把一些沒有意義的事剪掉,這些事消耗我們的生命力,消耗我們的時間,令我們分神,令我們的人生沒有焦點,沒有作用!神所大用的慕迪先生,口袋裡常帶著兩個放大鏡,一個是完整的,他說:「當陽光照在這鏡上產生一個焦點,這焦點很高熱度,能引起燃燒.」另一個是裂開兩半的,他說:「這也是一放大鏡,但陽光照射下不起作用.」很多基督徒在靈性上一無所成,不能結果子,就是因為他的人生沒有一個真正的焦點!請看社會上的成功人士,哪一個沒有集中努力的焦點呢?這是一般的成就,屬靈的成就上也一樣,求主幫助我們將人生中沒有意義的事除掉.主耶穌用祂的道來潔淨枝子(3),主的道是殺虫劑,把樹枝上的害虫消滅,讓樹枝健康,結果子.主的道潔淨我們的人生,常常有些罪的小虫爬到我們身上,這些小虫會長大,這些罪惡的事情會使我們的力量衰敗,讓我們失去了結果子的能力,這是撒旦的工作,主的道要將我們洗淨,洗去一切攔阻我們的罪,讓我們有力量,有屬靈的作用!
\newline
\newline
當枝子連在樹上的時候,就產生了四種的交流:
\newline
\newline
1．能力的交流(5-6):離開了主耶穌,我們就不能結果子.枝子連在樹上,就能結果子.枝子離開了葡萄樹,就沒有能力結果子,這是必然的!結果子的力量不是在枝子裡面,而是在樹裡面.多年前,我曾在九龍一間教會的團契講道,講道前有一姊妹作見證,在我心中留下深刻印象；她在一所基督教的醫院裡面作護士,工作很忙,很多時候病人的要求不能一一滿足,工作忙,壓力大令到她緊張,很容易發脾氣,但脾氣過後很難過,她想:這些病人需要神的愛,如果他們不能在我這基督徒身上得到,他們從何處得到呢?她很後悔,但後悔之後又發脾氣,發了脾氣又後悔,她甚至對自己感到絕望!有一天在聽道時聽到傳道人講這段《聖經》,知道離開了主,自己不能作甚麼,她領會了.以後每天工作之前跪在神前懇切禱告:「主啊,我沒有愛,您就是愛!求您將您的愛充滿在我心中,讓我用您的愛來服侍這些病人!」從此,她經歷到:主的愛透過她發出來.離開了葡萄樹,枝子就會枯乾,枯乾的枝子當然不能結果子!很多年前九龍塘宣道會在一個聚會中,「互愛團契」和「香港皇家警察以諾團契」兩詩班聯合獻唱；「互愛」是由一班戒了毒的人,有的人是坐監,在監獄聽聞福音,悔改信主,並在主裡面得到更新的人組成.而「以諾」是警察,本來執法者要捉犯罪者,但現在都信了主,可以聯合在一起,他們的人生都有了新的改變,一起來讚美真神!真是奇妙的改變!一個人信主以後要繼續連在主裡面,源源地支取主的生命,保持自己結果子的力量,這是能力的交流.我們與主相連,主的能力進來運行,結出果子奉獻給主,這是多麼美麗的交流!
\newline
\newline
2．道的交流(7):主的話就是主的道.「願意」這兩個字表示內心的願望,渴慕,追求；這是主的道在人心裡產生的回應,這心中的渴慕,追求變成禱告,「祈求」是內心渴慕,願望追求的表達.神的回應降下來,神的話進來產生願望,成為禱告,禱告上升,主耶穌就給我們成就!這就是道的交流.枝子連在葡萄樹上產生這交流,這樣交流不久之後,這人就道化,他和主的道,主的話打成一片,主的道使他變化,這是何等的寶貴!
\newline
\newline
3．愛的交流(9-12):這一段中不繼地講到「裡面」,主耶穌叫我們在祂裡面,祂也常在我們裡面；枝子在樹裡面,樹的生命在枝子裡面；這「裡面」表示一個交流,不斷地交流.愛的交流,這是個甚麼樣的愛呢?從(13)中我們看到「捨命」是主的愛,祂為我們犧牲生命,在十字架上付了最重的代價!這是犧牲,捨己,無比,長闊高深的愛!我們受感動,受激勵,回應這愛!主是我們最好的朋友,為我們捨命,這是最高的感情,我們用感情愛主沒錯,但不止於此,(10,15)中還談到道義的愛,如果我們不守主的誡命,不聽主的話,這愛是不完全的.再就是執行的愛,了解主的心意；主耶穌說,我的心事都告訴了你們,你們懂得,這是知心之交.主愛我們永不改變,但在很多時候,主會不會對我們有遺憾呢?祂關心的事我們不關心,祂認為重要的事,我們不看重,我們的心好像與主有距離,我們沒有了解主的心,不知道祂的旨意,心意,仍然按照自己的心意生活,為人......我們回應,接受的愛時,以上各方面都很重要!
\newline
\newline
4．喜樂的交流(11):主的喜樂進到我們心中,成為我們的喜樂.這是豐滿的喜樂!當主耶穌在兩個城市講道被拒絕之時,被聖靈感動,歡樂禱告(路十21),祂在逆境中歡樂,這是祂不改變的喜樂!主將這喜樂給了我們,讓我們在任何環境中保持這喜樂；這喜樂能使人滿足,富足!在基督裡我們有屬天的喜樂,藉著基督充滿在我們心中!這不是一時的快樂,這是永恆的喜樂和滿足!主的喜樂進來,充滿心中,變作敬拜,感謝,讚美升上去,這交流多麼地寶貴!
\newline
\newline


\newpage

\section{第71屆~港九培靈會~滕近輝~第三講}
\label{sec:22}
\textbf{真自由}
\newline
\newline
連結: \href{https://www.hkbibleconference.org/session-message/view/22}{\texttt{https://www.hkbibleconference.org/session-message/view/22}}
\newline
\newline
\hyperref[sec:23]{< < < PREV SERMON < < <}
~
\hyperref[sec:toc]{[返目錄]}
~
\hyperref[sec:21]{> > > NEXT SERMON > > >}
\newline
\newline
約翰福音 8:31-8:36
\newline
\begin{longtable}{cl}
\hline
\hline
章節 & 經文 (和合本修訂版)\\
\hline
8:31 & \begin{tabularx}{0.7\textwidth}{X} 耶穌對信他的猶太人說：「你們若繼續遵守我的道，就真是我的門徒了。 \end{tabularx} \\ \\ \relax
8:32 & \begin{tabularx}{0.7\textwidth}{X} 你們將認識真理，真理會使你們自由。」 \end{tabularx} \\ \\ \relax
8:33 & \begin{tabularx}{0.7\textwidth}{X} 他們回答他：「我們是亞伯拉罕的後裔，從來沒有作過誰的奴隸，你怎麼說『會使你們自由』呢？」 \end{tabularx} \\ \\ \relax
8:34 & \begin{tabularx}{0.7\textwidth}{X} 耶穌回答他們：「我實實在在地告訴你們，所有犯罪的人就是罪的奴隸。 \end{tabularx} \\ \\ \relax
8:35 & \begin{tabularx}{0.7\textwidth}{X} 奴隸不能永遠住在家裡；兒子才永遠住在家裡。 \end{tabularx} \\ \\ \relax
8:36 & \begin{tabularx}{0.7\textwidth}{X} 所以，神的兒子若使你們自由，你們就真正自由了。」 \end{tabularx} \\ \\
[1ex]
\hline
\hline
\end{longtable}
事奉主近五十年來,統計了一下,連我自己都不相信:大概有四千篇講章.所以開始事奉主時遇到難處要感謝神,如果我們中間有傳道人在講道上有難處,那要謙卑地追求,學習,追求,學習,神仍然可以用你,神會繼續地供應,這近五十年來講道的恩賜,神賜給我了!這是我的見證.
\newline
\newline
(約八31-36)中提到真自由,真自由是甚麼樣的自由呢?主耶穌提到真理必叫你們得以自由,即是說真自由是不自由的自由,因為自由受真理的限制,不合真理的事不能做,在真理的軌道上,原則上來自由!離開了真理沒有真自由,好像火車沿軌道前行,離開軌道寸步難行,立刻停止一樣.開車要看交通燈號,有紅燈的限制,如果一意孤行,就會撞車,進醫院,躺在病床上包裹傷處,失去了一切的自由!主耶穌講的話非常真切,按照這個意義來講,神自己也不自由,「祂不能背乎自己」(提後二13),就是不能違背自己的本性!神的本性是:真理,仁義,聖潔,神無論做甚麼,一定按照祂的本性,所以神也不自由.基督自由嗎?如果自由的話,為甚麼上十字架?!十字架自由嗎?感謝神!十字架是愛的自由,我們的主把祂的愛完全的彰顯出來,大膽地運用了自己的自由.
\newline
\newline
真自由相對的是假自由,那是撒旦允許人的自由,撒旦舉起一面旗幟:自由.結果卻是自由的不自由!撒旦說:你們到我這裡來,想作甚麼便做甚麼,你可以放縱你的一切,不必理別人是好是壞,完全沒有拘束,這樣放縱的自由最後帶來痛苦,帶來奴役,最後成為罪的奴僕.這是撒旦揮著假自由的旗號,吸引了世上許多人來跟隨,它有為所欲為的自由.
\newline
\newline
真自由的原則是甚麼呢?
\newline
\newline
一．在真理原則下的自由:我們立定心志:不合真理的事我們不做,違反神旨意的事我們不做.主耶穌說祂就是「道路,真理,生命」(約十四6),祂的道,祂的表現就是真理的光.
\newline
\newline
二．在愛心原則下的自由:「弟兄們,你們蒙召是要得自由,只是不可將你們的自由當作放縱情慾的機會,總要用愛心互相服事.」(加五13)不是放縱情慾,乃是用愛心互相服事.愛是我們真的自由,愛把我們從仇恨中釋放出來,仇恨把許多人綑綁；希特勒因為仇恨殺害了六百萬猶太人,你看仇恨是何等可怕!有一次我在夏令會中講道,另一位講員是羅海倫宣教士,她的見證非常感動人,我讀著,感到十字架的愛真是活化出來,她在非洲剛果扎伊爾作醫療宣教士,將自己的愛心傾倒出來幫助當地人,在扎伊爾發生內戰時羅醫生遭軍人強姦,她的心破碎了!她心裡有人性自然墮落裡面出來的恨!回國後感覺到不能再回去,直到有一天,她面對基督十字架,忽然十字架的愛感動她的心,彷彿聽到主耶穌在十字架上那禱告:「父阿,赦免他們......!」那十字架的代價,那十字架愛的流露溶化了她的心!她在那裡哭,哭,心得到了醫治.她得到愛中的釋放,於是再回去,繼續幫助扎伊爾人.恨,綑綁世上多少人,如果今天晚上在座有人心中有恨沒有消解,願你進入真自由,讓基督十字架的愛消除你的恨!
\newline
\newline
三．在服事原則下的自由:在服務裡將我們的自由充分發揮.(約十三)中提到主離世歸父的時刻到了,祂就為門徒洗腳.祂是老師,是主,但祂為門徒洗腳!有一次,我看到一副世界名畫－那是耶穌替彼得洗腳,畫家所表達的耶穌的姿態是一隻腳跪在地上替彼得洗.我想:尊貴的主,榮耀的主,祂是主!竟然如此謙卑地為卑微的人洗腳!這也是一件非常感動我們的事,主耶穌大膽地使用了祂的自由,留下一個永恆的榜樣給我們.
\newline
\newline
四．為神而作的自由:「你們雖是自由的,卻不可藉著自由遮蓋惡毒,總要作神的僕人.」(彼前二16)不可用自由來遮蓋惡毒,不可將自由當作幌子來表演給人看,外表美其名曰「自由」,裡面卻是不可告人的東西!這就是撒旦所說的自由；彼得說:「你們雖是自由的,卻不可藉著自由遮蓋惡毒.」一著名作家說:「自由,自由,世界上多少壞事是藉著你的名作出來的.」彼得又說:「總要作神的僕人」即是說:為了神從心底去做,用服事神的心來做一切事.這樣才是動機純正,這樣就脫離了自我,名利的綑綁.保羅說:「你們從前雖然作罪的奴僕,現今卻從心裡順服了所傳給你們道理的模範.」(羅六17)要從心裡遵循神的旨意,從心裡順服神的道.做神的僕人,為神而作,這種心態是自由的一個很重要的原則.自我綑綁很厲害；我們很容易以表演的心態去做事,要榮耀自己.我們要藉著從心裡順服神的旨意,以奉獻者的心態來得勝!有人覺得順服神的旨意不自由；諸多限制,想擺脫,擺「我用慈繩愛索牽引他們.」(何十一4)繩,好像綁住人,但是是慈繩；索,好像鎖住人,但是是愛索,是慈愛的繩索!神的旨意限制我們,但神的旨意中有祂的愛!當時的以色列人不想要這愛的限制,「我的帳棚毀壞；我的繩索折斷.」(耶十20),神感到祂的百姓掙斷祂的鎖出去了,百姓掙斷了神的慈繩愛索,要自由!《耶》頭幾章中有三次神發出慈愛的呼聲:「回來吧,回來吧!」一個認識神的人,接受神旨意的約束,那是愛的約束,在這約束裡有真自由!這是我們的天父要我們注意的一個原則.
\newline
\newline
五．在經義原則下的自由:「主的靈在那裡,那裡就得以自由.」(林後三17)上文講到律法的字句,字句叫人死,經義叫人活.「經義」是甚麼呢?即是聖靈.聖靈原來真正的意思,是在律法的條文的字句裡面；一個人真正了解《聖經》的經義以後,知道經義是從聖靈感動而來,所以聖靈是經義,真理的靈.聖靈在一個人心中運行之時,把人帶進屬靈的自由中,脫離了字句表面的意思,脫了條文,脫離了條規,進入了生命真義之中.感謝主!這裡面有自由!在今天這個時代,很多青年人要自由,其中包括基督徒.如果在追求自由的過程中面對的都是條規,條規,他會反抗!「反」是今天青年人的特色.規條引起「反」的心態,所以我們不能用規條來限制現代人的自由,只有聖靈帶領人進入真道的經義裡(經義即聖靈原來的意思),經義與人的生命結成一片,讓他在生命中看見亮光,了解真理,這時聖靈使他得自由.所以神的靈在那裡,那裡有自由!這是聖靈的工作引發出來的內在的生命的自由,這才是真自由!一件事很多人做,一個屬主的青年人可能不做；因為它不合真理,不合神的旨意,不合神的道!「有自由不做」,這又是大膽地運用了自由.從聖靈而得到的釋放,只有這一方法能令今天的基督徒順服神的道,這是真正的路線.正如(啟九1-5)中所講,無底坑裡有煙冒上來,(「坑」是象徵撒旦所住的地方)令到日頭和天空都昏暗了,整個世界烏煙瘴氣,在煙中有帶有蠍子能力的蝗虫飛出來,但牠們不可能傷到額頭上有印記的基督徒!基督徒有神的印記,有神的生命在裡面,撒旦不能傷害到他；這是一個榮耀的自由!脫離罪惡綑綁的生命是真自由的一個要素,這生命是靠聖靈運行來加強的.那裡有神的靈,那裡有自由!讓一切的真理變成活潑的原則在我們的心中自由運行,不是死板的字句,不是條規!不是「不要」,「不要」,乃是裡面有一種生命的動力,生命的感受,生命的引導,生命的價值觀!這是在這時代的非常重要的認識.
\newline
\newline
六．在造就人的原則下的自由:「凡事都可行,但不都有益處.凡事都可行,但不都造就人.」(林前十23)基督徒不是為自己活著,要常常想到別人,要幫助別人,建立別人,作神手中賜福的器皿,造就人!在這個原則裡去運用自由.
\newline
\newline
七．在榮耀神的原則下的自由:「無論做甚麼,都要為榮耀神而行.」(林前十31)不能榮耀神的不能做,在這個原則之中來自由.
\newline
\newline
這七個原則非常基本,真自由建立在七個原則的根基之上,這是神兒女榮耀的真自由!我們蒙召是要得自由,蒙召的自由就是這樣的自由!讓我們學習這些原則,藉這些原則表現出基督徒的真自由.這時代的潮流改變不了我們,在任何潮流的沖激下,我們都要掌握這七項真自由,站穩在這基礎上,就可以向潮流誇勝,向撒旦誇勝,向罪惡誇勝,最後向死亡誇勝!在罪惡的天羅地網中,基督給我們開了一條出路.正如(約壹五19)中所講:「全世界都臥在那惡者手下.」天羅地網把世界包圍,但是基督到世界來為我們開了一條自由的出路,通向永恆的榮耀!感謝神!神的兒子叫你自由,你就真自由了!
\newline
\newline


\newpage

\section{第71屆~港九培靈會~滕近輝~第四講}
\label{sec:21}
\textbf{真以色列人}
\newline
\newline
連結: \href{https://www.hkbibleconference.org/session-message/view/21}{\texttt{https://www.hkbibleconference.org/session-message/view/21}}
\newline
\newline
\hyperref[sec:22]{< < < PREV SERMON < < <}
~
\hyperref[sec:toc]{[返目錄]}
~
\hyperref[sec:20]{> > > NEXT SERMON > > >}
\newline
\newline
約翰福音 1:43-1:46
\newline
\begin{longtable}{cl}
\hline
\hline
章節 & 經文 (和合本修訂版)\\
\hline
1:43 & \begin{tabularx}{0.7\textwidth}{X} 又過了一天，耶穌想要往加利利去。他找到腓力，就對他說：「來跟從我！」 \end{tabularx} \\ \\ \relax
1:44 & \begin{tabularx}{0.7\textwidth}{X} 這腓力是伯賽大人，是安得烈和彼得的同鄉。 \end{tabularx} \\ \\ \relax
1:45 & \begin{tabularx}{0.7\textwidth}{X} 腓力找到拿但業，對他說：「摩西在律法書上所寫的，和眾先知所記的那一位，我們遇見了，就是約瑟的兒子拿撒勒人耶穌。」 \end{tabularx} \\ \\ \relax
1:46 & \begin{tabularx}{0.7\textwidth}{X} 拿但業對他說：「拿撒勒還能出甚麼好的嗎？」腓力說：「你來看。」 \end{tabularx} \\ \\
[1ex]
\hline
\hline
\end{longtable}
一個真中國人是具有中國血統的人,一個真以色列人有雙重的血統:在他們身體裡面有先祖亞伯拉罕的血統,在他們的靈性裡有亞伯拉罕屬靈的血統.甚麼是亞伯拉罕屬靈的血統呢?即他的屬靈的特別素質所構成的.拿但業,主耶穌說他是個真以色列人,因為主耶穌在他身上發現了真以色列人的特徵,他是亞伯拉罕真正的子孫!我們來看看亞伯拉罕屬靈的特別素質；這些素質在今天真以色列人身上顯明出來.
\newline
\newline
一．亞伯拉罕是信心之父
\newline
\newline
亞伯拉罕是以信為本的人,他信神,他的信心使他走一條信仰之路；亞伯拉罕得到神的呼召,離開了他的故鄉吾珥(加勒底,後來的巴比倫),他的信心有具體的表現,信仰成了人生的道路.真正的信仰不單單是頭腦的知識,乃是我們所走的人生道路；凡信這道理的人走的是人生美好的一條道路:「我卻不以性命為念,也不看為寶貴,只要行完我的路程,成就我從主耶穌所領受的職事,證明神恩惠的福音.」(使廿24),亞伯拉罕的信仰就是這樣的信仰,他與神的關係是藉著信心而產生的,他信神,所以神向他顯現；也可以說,神向他顯現和他信神同時發生.
\newline
\newline
我們的信仰是怎樣來的呢?是神的靈在我們心裡將神向我們顯現,我們的心靈看見了,感覺到神的存在,那是聖靈的工作的最重要的一剎那,而且是藉著道(神的話語).所以神向亞伯拉罕顯現之時向他講話,他用信心回應了神的話,按照神的引導起來走一條新的道路.
\newline
\newline
二．亞伯拉罕是一個有使命的人
\newline
\newline
「我必叫你成為大國.我必賜福給你,叫你的名為大；你也要叫別人得福.」(創十二2)神賜福給亞伯拉罕沒有到此為止,神要叫他蒙福之後要使別人得福,要做一個賜福的器皿,和別人分享他所得的福分；亞伯拉罕真是擔起這個使命,他領受了神所賜的福,又與他人分享!「並且地上萬國都必因你的後裔得福.」(廿二18)這使命不單止是給亞伯拉罕一個人的,他是神賜福的器皿,他的後裔也是神賜福的器皿.器皿即渠道,神的恩典透過渠道臨到了別人!亞伯拉罕的後裔分三種:
\newline
\newline
(1)以色列人,
\newline
\newline
(2)耶穌基督,
\newline
\newline
(3)那以信為本的人(加三7).
\newline
\newline
神給亞伯拉罕的使命也是神給我們的使命,耶穌基督當然完成了這個使命；世上千千萬萬的人因祂付代價而蒙恩得救.以色列人有沒有做賜福的器皿呢?在這件事上以色列人是失敗的,是演繹的悲劇,他們沒有擔起這使命,這是很可惜的事!這成為我們的警戒.我們這些以信為本的基督徒要接受以色列人失敗的警戒,要彼此勉勵,擔起這使命.我們要和別人分享神的恩,神的愛,神的道,在這使命上依靠神再接再厲.凡是有使命感的人是亞伯拉罕的後裔,我們在這事上應該常常學習.弟兄姊妹們,神的旨意是要你,我擔起這使命,成為賜福的器皿,就是蒙福之後和人分享這福.
\newline
\newline
三．亞伯拉罕是一個分別出來歸於神的人
\newline
\newline
亞伯拉罕身上有一個記號:割禮.(創十七)提到亞伯拉罕和他家每一個男子都受了割禮,這割禮表示他們和其他人不同:他們是屬神的,要和世上的人有所分別,割禮是分別為聖的記號.今天所有屬神的人應該有這屬靈的記號.保羅說:「因為外面作猶太人的,不是真猶太人；外面肉身的割禮,也不是真割禮.惟有裡面作的,才是真猶太人；真割禮也是心裡的,在乎靈,不在乎儀文.這人的稱讚不是從人來的,乃是從神來的.」(羅二28-29)甚麼是真割禮呢?是指屬靈的割禮,是分別為聖之意.「情慾的事都是顯而易見的,就如姦淫,污穢,邪蕩,拜偶像,邪術,仇恨,爭競,忌恨,惱怒,結黨,紛爭,異端,嫉妒,醉酒,荒宴等類.我從前告訴你們,現在又告訴你們,行這樣事的人必不能承受神的國.」(加五19-21)「情慾」即是社會中的各種罪惡,要將這些割去,這就是屬靈的割禮.「耶和華說:『看哪,日子將到,我要刑罰一切受過割禮,心卻未受割禮的.』」(耶九25)神要刑罰身體受過割禮而心未受過割禮的人,以色列人每個男丁一生下來就受割禮,但神說身體受過割禮而心未受過割禮的,仍然在神的審判之下.我們的耳朵,口,手,腳,眼睛,思想......整個人都要受割禮,不做,不看,不聽,不想神不喜悅的事,正如保羅所講:「將身體獻上,當作活祭,是聖潔的,是神所喜悅的.」(羅十二1)保羅不是說將心獻上,乃是說將身體獻上；如果有人說:我把心獻給神了；手是我的,想做甚麼就做甚麼；眼睛是我的,想看甚麼就看甚麼；口是我的,想講甚麼就講甚麼；腳是我的,想走甚麼路就走甚麼路；結果奉獻是空的,沒有具體的內容.我們要整個人受具體的割禮,屬靈的割禮,分別為聖的割禮.這是亞伯拉罕屬靈的特質,也是我們屬靈的特質；所以我們真是亞伯拉罕的子孫,也是真以色列人!
\newline
\newline
四．亞伯拉罕與神立約
\newline
\newline
「約」是雙方面的.神對亞伯拉罕說:「你當在我面前作完全人.」(創十七1)神應許使亞伯拉罕後裔極多極多,成為多國的父,並為他改了名,將他原來的「亞伯蘭」改為「亞伯拉罕」.這是雙方面的約,神履行了一方面的責任,亞伯拉罕也履行了另一方的責任.今天我們都是亞伯拉罕屬靈的子孫,全世界每一個民族中都有許多屬主的人,這是神作成的；亞伯拉罕呢?他也靠著神追求,努力作一個完全人.這約後來有兩個階段:
\newline
\newline
(1)神藉著律法與以色列人立約:在以色列人方面的責任,是遵循神的道做完全人.在神方面的責任,是祝福以色列人,使他們成為大國,成為神恩典的出口.但是,律法的約失敗了!(2)立新約:「耶和華說:『日子將到,我要與以色列家和猶大家另立新約,不像我拉著他們祖宗的手,領他們出埃及地的時候,與他們所立的約.我雖作他們的丈夫,他們卻背了我的約.......我要將我的律法放在他們裡面,寫在他們心上.我要作他們的神,他們要作我的子民.』」(耶卅一31-33)這新約,神寫在屬神的人的裡面,他們的心上,也就是內在的律法!成了屬靈的渠道,內心的原則,不是條文,不是儀式.舊約律法漸漸變成了外面的條文,儀式,現在新約是寫在心裡,成為心裡的原則,成為內心的指導,在內心遵循神的旨意.不再是外面的字句,而是心中的神旨了.所以凡是將神的話接收在心版上的,就是真以色列人了!「無論作甚麼,都要從心裡作,像是給主作的,不是給人作的.」(西三23)「你們從前雖然作罪的奴僕,現今卻從心裡順服了所傳給你們道理的模範.」(羅六17)我們要從心裡順服,心裡事奉,要口唱心和地讚美神,要用心靈和誠實來敬拜神,這是內在的生命,內在的生命之律!讓我們得到釋放,進入屬靈的實際中,成為亞伯拉罕真正的子孫.
\newline
\newline
五．亞伯拉罕獻「活祭」
\newline
\newline
亞伯拉罕將他的獨生愛子以撒,當作祭物獻給神,他作這件事真是不容易.神要試驗他,在這過程中,可想而知他是多麼痛苦,他心中一定有問號:「神啊,以撒是您應許給我的,為甚麼又要收回去?」他不明白,但他信得過神,就照著做；當獻上祭壇時,神的使者來了,叫他不要傷害這童子,又說:「地上的萬國都必因你的後裔得福.」(創廿二18)這「後裔」當時是指以撒而講的,神的意思是說:「我把你獻給我的歸還你,並且加上祝福!」這是神常常做的事,祂要求我們奉獻,奉獻生命,奉獻金錢,奉獻時間,奉獻心血來事奉祂,傳福音,分享神的恩典；當我們這樣獻上的時候,神悅納,然後賜下更豐富的恩典!亞伯拉罕獻上以撒,神給回他,他並未失去,而且有更有意義的事在後面,神藉這件事來預表神把祂的獨生愛子賜給人類.今天亞伯拉罕真正的子孫,真以色列人,也要像亞伯拉罕一樣獻上「活祭」!亞伯拉罕獻了愛子為「活祭」,今天我們要把自己當作「活祭」獻給神(羅十二1-2),因為神的慈悲,神的愛激勵了我們,感動了我們,為了回應神的愛,就把自己獻上,當作「活祭」!亞伯拉罕獻上兒子也是為了回應神的愛,所以真以色列人的一個標誌就是愛神,以致將自己獻給神!
\newline
\newline
以色列人在《舊約》時代竟然把耶和華的祭壇拆掉,來建立巴力的祭壇,先知以利亞面對這處境,起來大聲疾呼,將以色列人聚集在迦密山頂,一個人面對四百五十拜巴力的先知,另外還有耶洗別的先知共八百五十個.在這處境下,以利亞剛強,一心建立被廢去的耶和華祭壇,重新將祭獻上,在禱告中得到神悅納,從天降下火來燃燒祭物,以利亞大聲呼籲人們歸向真神,以色列人復興起來了!結果巴力的先知被殺.巴力這宗教與現今的潮流有密切的關係,在考古的著作中,巴力的偶像是一個騎在馬上的裸體的女人,她一手拿蛇,一手拿一枝白色的水仙花,代表巴力教四要點:
\newline
\newline
(1)騎馬代表巴力是戰神,
\newline
\newline
(2)裸女雙乳代表豐收,是農神,
\newline
\newline
(3)一隻手中的蛇代表它宗教的部份,
\newline
\newline
(4)另一隻手中的水仙花,象徵自由性愛.今天這世界是戰爭的世界,物質主義的世界,要豐收再豐收,今天的世界的潮流是撒旦大行其道,講自由,講性愛的時代.巴力教的這四特徵影響著當時的以色列人.
\newline
\newline
感謝神!祂藉先知以利亞喚醒了他們,在以利亞的信心和心志的號召下,以色列人復興了.今天我們也需要以利亞的心志能力的復興,我們拆去巴力的祭壇,重建對神的祭壇!在這個祭壇獻上愛的祭物,為回應神的愛而獻上自己為神而活.
\newline
\newline


\newpage

\section{第71屆~港九培靈會~滕近輝~第五講}
\label{sec:20}
\textbf{真光}
\newline
\newline
連結: \href{https://www.hkbibleconference.org/session-message/view/20}{\texttt{https://www.hkbibleconference.org/session-message/view/20}}
\newline
\newline
\hyperref[sec:21]{< < < PREV SERMON < < <}
~
\hyperref[sec:toc]{[返目錄]}
~
\hyperref[sec:19]{> > > NEXT SERMON > > >}
\newline
\newline
約翰福音 1:1-1:14
\newline
\begin{longtable}{cl}
\hline
\hline
章節 & 經文 (和合本修訂版)\\
\hline
1:1 & \begin{tabularx}{0.7\textwidth}{X} 太初有道，道與神同在，道就是神。 \end{tabularx} \\ \\ \relax
1:2 & \begin{tabularx}{0.7\textwidth}{X} 這道太初與神同在。 \end{tabularx} \\ \\ \relax
1:3 & \begin{tabularx}{0.7\textwidth}{X} 萬物都是藉著他造的，沒有一樣不是藉著他造的。凡被造的， \end{tabularx} \\ \\ \relax
1:4 & \begin{tabularx}{0.7\textwidth}{X} 在他裡面有生命，這生命就是人的光。 \end{tabularx} \\ \\ \relax
1:5 & \begin{tabularx}{0.7\textwidth}{X} 光照在黑暗裡，黑暗卻沒有勝過光。 \end{tabularx} \\ \\ \relax
1:6 & \begin{tabularx}{0.7\textwidth}{X} 有一個人，是從神那裡差來的，名叫約翰。 \end{tabularx} \\ \\ \relax
1:7 & \begin{tabularx}{0.7\textwidth}{X} 這人來是為了作見證，是為那光作見證，要使眾人藉著他而信。 \end{tabularx} \\ \\ \relax
1:8 & \begin{tabularx}{0.7\textwidth}{X} 他不是那光，而是要為那光作見證。 \end{tabularx} \\ \\ \relax
1:9 & \begin{tabularx}{0.7\textwidth}{X} 那光是真光，來到世上，照亮所有的人。 \end{tabularx} \\ \\ \relax
1:10 & \begin{tabularx}{0.7\textwidth}{X} 他在世界，世界是藉著他造的，世界卻不認識他。 \end{tabularx} \\ \\ \relax
1:11 & \begin{tabularx}{0.7\textwidth}{X} 他來到自己的地方，自己的人並不接納他。 \end{tabularx} \\ \\ \relax
1:12 & \begin{tabularx}{0.7\textwidth}{X} 凡接納他的，就是信他名的人，他就賜他們權柄作神的兒女。 \end{tabularx} \\ \\ \relax
1:13 & \begin{tabularx}{0.7\textwidth}{X} 這些人不是從血生的，不是從情慾生的，也不是從人的意願生的，而是從神生的。 \end{tabularx} \\ \\ \relax
1:14 & \begin{tabularx}{0.7\textwidth}{X} 道成了肉身，住在我們中間，充充滿滿地有恩典有真理，我們也見過他的榮光，正是父獨一兒子的榮光。 \end{tabularx} \\ \\
[1ex]
\hline
\hline
\end{longtable}
每年的港九培靈研經會差不多都會下一次雨,但弟兄姊妹們風雨無阻,求神賜下祂如傾盆大雨般的恩典.
\newline
\newline
「那光是真光,照亮一切生在世上的人.」(約一9)那真光就是耶穌基督的光,主的光是真光.在《約》中,將耶穌基督的光的內容從不同的角度向我們講出來.
\newline
\newline
一．這真光是生命的光
\newline
\newline
「生命在祂裡頭,這生命就是人的光.」(一4)基督的生命是人的光,光與生命常連在一起.在大自然裡是這樣:陽光下,百花開放,彰顯出豐盛的生命.同樣,當基督的光在我們心中照耀,我們屬靈的美德之花開放,感謝天父,在基督裡有生命的光!這光結出果實,開放屬靈的花朵.《約》中所指的光是人內心的靈性之光,我們需要這生命的光照,讓我們的新生命發展,討主的喜悅,但願你,我的心是屬靈的向日葵,永遠跟著主,那是何等美麗的事!
\newline
\newline
二．這真光是引領我們人生道路的光
\newline
\newline
「我是世界的光.跟從我的,就不在黑暗裡走,必要得著生命的光.」(八12)主的光將人從黑暗裡引導出來.無可否認,世上有許多美好的事物,但總的來說,這是一個黑暗的世界.我們看到今天青少年多受黑暗潮流沖激,如果一個人心中有基督的光來引導,如果他跟從這位光明之主,就不在黑暗中行,就能脫離黑暗,勝過黑暗!「光照在黑暗裡,黑暗卻不接受光.」(5)感謝主!基督的光藉千萬人的得勝而勝過黑暗!撒旦是黑暗的王,但黑暗之君王不能勝過光明的君王!
\newline
\newline
我們跟從主,在祂引導下從黑暗中出來,不讓黑暗的權勢控制我們,我們常常面對基督光,被光照耀.主耶穌怎樣引導祂的羊?「既放出自己的羊來,就在前頭走,羊也跟著他,因為認得他的聲音.」(十4)我們看到主耶穌引導祂的羊:
\newline
\newline
(1)用榜樣,走在前面,羊看著牧人安然跟從.
\newline
\newline
(2)用祂的道,祂自己實行祂的道,所以祂的聲音又是榜樣!光明的主用祂光明的腳蹤來引導羊走在光明中,這引導的光何其寶貴!
\newline
\newline
三．機會的光
\newline
\newline
「光在你們中間還有不多的時候,應當趁著有光行走,免得黑暗臨到你們；那在黑暗裡行走的,不知道往何處去.你們應當趁著有光,信從這光,使你們成為光明之子.」(十二35-36)主耶穌說趁著......趁著有光,抓住有光的機會,免得在黑暗中行；在黑暗中不知方向,不曉得人生的歸宿是甚麼,不曉得永生是甚麼；要趁住有光來走永生之路,能夠進入那救恩之門.這光會過去,得救的感動會過去,有的人受了感動要信主,但是硬著心,一次又一次來抗拒這感動,結果這感動慢慢消失,這光慢慢消失,錯過了得救的機會!很多人經過救恩的門卻錯過了機會,希望在座的,沒有一位錯過這在光明中發現永生之門的機會.「那門是窄的,路是小的,找著的人也少.」(太七14)所以我們要注意,儆醒,抓住在光中得救的機會,抓住這機會,趁著有光,信從這光,就成為光明之子.光照在心靈中,心靈光明了,光照在人生中,有了光明的人生.最後有光明的歸宿,進到榮耀的天家,與我們的主永遠同在,成為光明之子,也是新的人生機會.「趁著白日,我們必須作作那差我來者的工；黑夜將到,就沒有人能作工了.我在世上的時候,是世上的光.」(約九4-5)蒙恩但沒有報恩,是一件令人後悔的事.
\newline
\newline
四．見證的光
\newline
\newline
「約翰是點著的明燈,你們情願暫時喜歡他的光.」(五35)施洗約翰見證主耶穌,是見證的明燈,他的光是見證主的光.「這人來,為要作見證,就是為光作見證,叫眾人因他可以信.」(一7)他來要為光作見證,為主耶穌作見證.「你們是世上的光.城造在山上是不能隱藏的.人點燈,不放在斗底下,是放在燈臺上,就照亮一家的人.」(太五14-15)我們每一個屬主的人都要發這光,基督徒在社會上要為主發光,叫人看見我們的好行為,將榮耀歸給我們在天上的父!平時我們在外面戴了一個面具,我們很注意給人一個好印象；回到家原來的面目出來了,不少人外面的面目與裡面的面目是兩回事.主說:「就照亮一家的人.」在家裡發光更不容易,因那是天長日久的事,一切都放鬆了,生活中自然的習慣不知不覺地流露出來了,那時候是真面目,是本相.一個人脾氣好不好,在外面是一回事,在家裡可能不同.這就是有些人結婚前很好,婚後問題愈來愈多的原因,所以每一個基督徒的家庭要學功課,接受主的提醒,把燈點起來放在燈台上,照亮一家的人.
\newline
\newline
五．避免跌倒的光
\newline
\newline
「耶穌回答說:『白日不是有十二小時麼?人在白日走路,就不至跌倒,因為看見這世上的光.若在黑夜走路,就必跌倒,因為他沒有光.』」(約十一9-10)光叫一個人不至於跌倒,主耶穌在此再一次說自己是世上的光,祂是世上的光,我們在這光裡行走不至於跌倒,在路上有甚麼絆腳的東西我們能看得清清楚楚,可以避開.主的光救我們脫離試探,讓我們看到試探在何方.社會上有很多陷阱,有很多試探人的地方和事物,平常的人不知不覺落在其中,但如一個人心中有基督的光照,在光明中走路,清楚看到陷阱在何處,試探在那裡,撒旦工作的方式在那裡,就會小心,謹慎防備,不至於落在魔鬼的詭計中.這時代,每一個基督徒都需要讓基督的光天天照耀自己的心,在基督的光中來看見.魔鬼,試探不會離開,只有行在基督光裡的人自己離開試探!要靠著主來追求,要作努力的嘗試,要智慧地遠離試探,抵擋魔鬼,逃避試探!這世界情慾的試探包圍每一個人；以前要去到色情場所,現在色情場所送進家門,無處可避,無處可逃,連青年人少年人都逃不掉；只有心裡有主的人能得勝!向撒旦,向罪惡誇勝!
\newline
\newline
六．光明與黑暗是敵對的
\newline
\newline
「光來到世間,世人因自己的行為是惡的,不愛光,倒愛黑暗,定他們的罪就是在此.」(三19)世界上很多人恨光,因為他們在黑暗中生活,在黑暗中作惡,他們不愛光倒愛黑暗.主耶穌責備那些人,顯出他們是黑暗的；他們就恨光,恨耶穌,一定要將他處死,就是那些恨光的領袖將耶穌交彼拉多釘在十字架上!十字架就是這樣的來由,基督徒背十字架也要付這樣的代價!
\newline
\newline
七．基督的榮光
\newline
\newline
「道成了肉身,住在我們中間,充充滿滿的有恩典有真理.我們也見過他的榮光,正是父獨生子的榮光.」(約一14)恩典和真理剛剛合成一對:只有恩典沒有真理是溺愛,單有真理沒有恩典只有處罰.恩典,真理配合在一起,是耶穌基督十字架的來源,祂擔當人的罪,為罪付代價；十字架是神慈愛,公義的交叉點,沒有神的愛就沒有十字架!沒有神的公義就沒有十字架!神的公義和慈愛在十字架上匯合了,這是神救人必須,必須之途徑.耶穌基督表現的榮光就是恩典,真理配合的榮光,(約八)中提到那行淫時被拿到的婦人,很多人要拿石頭打死她,因她違反了摩西的律法,他們把她帶到耶穌面前,問耶穌怎麼辦,明顯地要為難祂,但《聖經》上記載:發生了一件最感動人的事,主耶穌說:「我也不定你的罪.去吧,從此不要再犯罪了.」(11)應該定罪不定罪,第一句話是恩典,第二句話是真理.這是主的恩典,真理的榮光,這榮光把那些拿著石頭要打死這女人的人溶化散開了.當時主耶彎著腰在地上劃字,有人說是在劃:「我為她死,我為她死!」這是猜想,但是一個很可能的猜想!至少主耶穌的心是這樣想的.
\newline
\newline
「神就是光,在他毫無黑暗.這是我們從主所聽見,又報給你們的信息.我們若說是與神相交,卻仍在黑暗裡行,就是說謊話,不行真理了.我們若在光明中行,如同神在光明中,就彼此相交,他兒子耶穌的血也洗淨我們一切的罪.我們若說自己無罪,便是自欺,真理不在我們心裡了.我們若認自己的罪,神是信實的,是公義的,必要赦免我們的罪,洗淨我們一切的不義.」(約壹一5-9)這裡對神有一個重要的介紹:神就是光.神是光顯明一切的罪惡,神是愛要救罪人,神的這兩個屬性再一次在十字架上配合起來.我們如果與神交往,就不在黑暗中行,要認自己的罪,神必赦免我們的罪,洗淨我們一切的不義,這是光明的神與祂的愛連在一起!一個屬神的人是一個注重常常潔淨自己的人,(9)中的「認」字是一進行式,指繼續認罪,不是一次過認罪.不是信主時說「主啊,您赦免我現在的罪,將來的罪.」以後就可以犯罪了.不錯,信主時,基督的寶血都遮蓋了我們過去的罪,但信主之後,每一次做錯事都要認自己的罪,讓基督的寶血洗乾淨,靠著主不再犯罪.人有軟弱(但不是自己原諒自己的那種軟弱),有的時候有了過犯,心被光照,真正憂傷痛悔；如果犯了罪,又悔改,真正悔改,再認罪,基督的寶血再潔淨,遮蓋.(詩五一)中,我們看到大衛犯罪,但憂傷痛悔地悔改.所以一個基督徒犯了罪,要像大衛一樣,(1)要知罪,(2)要認罪,(3)要為罪憂傷,(4)要離開罪.
\newline
\newline
「黑暗漸漸過去,真光已經照耀.......愛弟兄的,就是住在光明中,在他並沒有絆跌的緣由.」(約壹二8,10)這是愛的光,愛人的就是在光明中.愛是光明,一個人有愛才能從黑暗中出來,在光明中行不至於跌倒.愛是真光,用真光來看才能得到真的認識,一個有愛的人有愛的眼光,從愛的眼光來看問題,就可以得到別人沒有的結論,得到真的,正確的結論.求神幫助我們用愛的眼光來看人,看事,看一切,這才不會跌倒,才沒有跌倒的緣由.帶有色眼鏡看人是不好的,只有用愛的眼鏡,用愛的顏色來看一切,立刻你就會覺得昇華了,你這人有了改變:有新的眼光,新的心態,新的感受,新的作法,那就是神的作法.
\newline
\newline


\newpage

\section{第71屆~港九培靈會~滕近輝~第六講}
\label{sec:19}
\textbf{真門徒}
\newline
\newline
連結: \href{https://www.hkbibleconference.org/session-message/view/19}{\texttt{https://www.hkbibleconference.org/session-message/view/19}}
\newline
\newline
\hyperref[sec:20]{< < < PREV SERMON < < <}
~
\hyperref[sec:toc]{[返目錄]}
~
\hyperref[sec:18]{> > > NEXT SERMON > > >}
\newline
\newline
約翰福音 8:31-8:31
\newline
\begin{longtable}{cl}
\hline
\hline
章節 & 經文 (和合本修訂版)\\
\hline
8:31 & \begin{tabularx}{0.7\textwidth}{X} 耶穌對信他的猶太人說：「你們若繼續遵守我的道，就真是我的門徒了。 \end{tabularx} \\ \\
[1ex]
\hline
\hline
\end{longtable}
「他們的數目有千千萬萬.」(啟五11)使徒約翰看見異象,有數不過來的人從各族各國各方環繞羔羊的寶座,在那裡唱歌,讚美,喜樂.眾所周知:《聖經》已經翻譯成二千多種文字,如果屬主的人用這樣多的文字,語言歌唱讚美神,這種喜樂的感受是很難描寫,表達出來的.
\newline
\newline
一．主的門徒是遵守主道的人
\newline
\newline
「你們若常常遵守我的道,就真是我的門徒.」(約八31)誰是真門徒呢?就是那常常遵守主道的人,主的道是我們人生道路的標準,規範.開始做基督徒時,要把我們的一言一行,一舉一動,思想感受,都放在基督的道的方格中,說來也真不容易!我們學習,再學習,學而時習之,最後成功；當我們能夠常常遵循主道時,心中何等喜樂!當我們的主看到一個一個遵守祂道的基督徒時,祂的心中也是何等喜樂,祂說:「你們真是我的門徒了.」
\newline
\newline
在《使》中,信耶穌的人不叫「信徒」,叫「門徒」.「信徒」是很好的名稱,「信」是最基本的與神的關係.為甚麼使徒時代要選用「門徒」一詞呢?「門徒」是學老師,一個信主的人要學主,從信主到學主,這是使徒時代教會一個重要的焦點.安提阿的「門徒」開始被稱為「基督徒」(使十一26),這名稱不是自己命名的,而是非基督徒命名的,因為他們在信耶穌的人身上看見一個表現:這些人信耶穌,敬拜耶穌,見證耶穌,效法耶穌.「基督徒」就是基督和人,意解人與基督聯在一起.《舊約》裡的「神人」是說這人信神,敬畏神,事奉神,見證神.《舊約》「神人」,《新約》「基督人」給我們看見教會歷史以此為重點的階段.一個人信了耶穌之後要學耶穌,效法耶穌,作門徒.如果今天一個人被稱為「門徒」的話,那麼這人一定很榮耀主,這真是一個很寶貴的名稱!
\newline
\newline
二．主的門徒是跟從主的人
\newline
\newline
(約一35-39)中,有兩個人因為聽了施洗約翰的見證,跟從了耶穌,作祂的門徒.這當中有一經歷的過程:
\newline
\newline
(1)看,那是觀察,觀察耶穌到底是個甚麼樣的人；
\newline
\newline
(2)同住,接受主的教導,訓練,與主住在一起；
\newline
\newline
(3)跟從主,成為主的門徒.
\newline
\newline
今天我們做主的門徒也不是糊里糊塗的；我們真是看出耶穌基督是主,經過各方面的詳細觀察,受了主的教導,清楚地認識了祂,經歷了祂,真心誠意地作祂的門徒.
\newline
\newline
巴不得每一個人都真正地願意跟從主耶穌!彼得跟從主耶穌,很可惜,有一次他是「遠遠地」跟從(可十四54),他當時很怕,因為主耶穌被逮捕了；結果彼得跌倒了,三次不認主!每一個「遠遠」跟從主的人免不了要跌倒!求主幫助我們緊緊地跟隨主,這樣我們可以站立得住!彼得後來復興了,在(約廿一22)中,主耶穌再一次呼召彼得跟從主；彼得緊跟主到底,最後為主被倒釘十字架.感謝神!一個曾經跌倒過的門徒,再起來繼續堅強穩定地跟從主,成為一個榮耀主的人.我們每一個人都可以這樣堅定地跟從主到底!如果我們曾經試過遠遠地跟從主,主赦免我們,在主的愛中我們復興!當彼得三次不認主時,《聖經》說主轉過身來看他,彼得看見主的眼睛；那不是定罪的眼睛,那是同情,了解的眼睛!彼得就出去痛哭.感謝神!那是他回轉的開始.我們雖然有時候軟弱,有時候遠遠跟從主,但是主的愛感動我們,溶化我們,終有一天我們改變了,我們決心永遠緊緊跟從主,走一條榮耀主的
\newline
\newline
(約十九)清楚地記載:使徒約翰跟從主,直到十字架下面；主耶穌在十字架上看見約翰,就將祂的母親交托給他.一個在十字架下面的人,一個緊緊跟從主直到十字架的人,主將託付給了他!主耶穌將託付交給每一個真正跟隨他的人.約翰看到被釘十字架的耶穌肋旁被槍刺之時,血和水流出來(34)；我曾讀一醫生的著作,他這樣寫:按照醫學來講,人的身體經歷很痛苦時,在十字架那樣的姿態中,他的心臟裡的血和水分開.約翰不是醫生,就算是,那時候也不知這醫學的事實,但他看見了,並寫下來,這是一個很寶貴的見證!而且這事有很重要的屬靈的意義:在(亞十三1)預言:「那日,必給大衛家和耶路撒冷的居民開一個泉源,洗除罪惡與污穢.」神開了一個救恩的泉源,洗淨人的罪；十字架上的主,肋旁流出血和水,血贖罪,水洗淨罪；這是一個跟從主到十字架下的人所看見的.
\newline
\newline
三．主的門徒是一個結果子的人
\newline
\newline
「你們多結果子,我父就因此得榮耀,你們也就是我的門徒了.」(約十五8)多結果子就是主的門徒了,我們從上下文看見要結靈性的果子.除此之外,(16)提到另外一種果子,主分派門徒去結果子,這是要「去」才能結的果子；這是傳福音,得靈魂的果子,去傳福音,為主得人.結出這兩種果子就是主的門徒!如果一個人信了主沒有任何果子結出來,那麼請問他是屬主的人嗎?信主後為甚麼完全沒有轉變,沒有果子結出來呢?這是一個很有道理的問號,一生一世,一個靈魂的果子都沒結出來,自己的心中有甚麼感覺呢?主的心中有甚麼感覺呢?!
\newline
\newline
很多年前,中國大陸還沒開放,有一個弟兄隻身來港,信主後很熱心傳福音,他用不停寫信的方法帶了自己國內的親友信主,後來又向周圍相識的人要他們在國內親友的地址,寫信向他們傳福音,很多人把未信神的親友告訴他,經過了大約六,七年,他成功用寫福音信的恩賜帶了一千多人信主!
\newline
\newline
在另一方面,我看到一個統計報告:一般的教會,帶人信主的會友大概百分之五,百分之九十五的人沒有做傳福音的工作!有一個人去中國內地去旅行,坐在火車上,看到對面有一人,看上去有學問,有地位,心中有感動向對方傳福音,但又怕他不會信,一直沒有開口,後來將自己口袋裡的小《聖經》拿出來讀,沒想到對面那人很想了解基督教,他很高興地剛想傳福音,但是到站了,那人要下車,他很後悔,神預備了一人給他,但他沒勇氣,沒開口,結果錯過機會.這見證提醒我們:許多時,聖靈預備了一些人的心,等候基督徒去向他們傳福音,但基督徒不開口,錯過時機!求主給我們信心,膽量和心志,來結靈魂和靈性的果子!主耶穌分派我們去結果子,多結果子就是祂的門徒了.
\newline
\newline
四．主的門徒是效法主的人
\newline
\newline
「我是你們的主,你們的夫子,尚且洗你們的腳,你們也當彼此洗腳.」(約十三14)主耶穌是門徒的主,給門徒洗腳,叫門徒要彼此洗腳,效法主作的事,學習像主.學生學書本,門徒學老師.「你們該效法我,像我效法基督一樣.」(林前十一1)保羅效法了基督,所以可以叫弟兄姊妹們效法他,一個效法主的人產生更多效法主的人.效法主的人有一種感力:別人和他接觸時感受他像基督,從中得到勉勵,跟著他的榜樣,起來和他一起效法基督.這樣門徒產生門徒,門徒去得人作門徒!主耶穌在大使命中說:「所以你們要去,使萬民作我的門徒......」《聖經》中有名詞:「各從其類」,甚麼類產生甚麼類,生物是這樣,靈性也如此,甚麼樣的人也影響別人像他,這是不可避免的.跟好人學好人,跟壞人學壞人.接近一個熱心的基督徒,你也會愈來愈熱心,接近一個冷淡的基督徒,你也會愈來愈冷淡.弟兄姊妹們,每一個人會受別人的影響,也可以影響別人,我們每一個都是可以幫助別人的基督徒,都是賜福的器皿.主要用你,主要用我,主要用我們每一個人作主的門徒!
\newline
\newline
五．主的門徒要彼此相愛的人
\newline
\newline
「你們若有彼此相愛的心,眾人因此就認出你們是我的門徒了.」(約十三35)彼此相愛,是主給門徒的新命令.為甚麼是新命令呢?因為有一個新榜樣,主說:「我怎樣愛你們,你們也要怎樣相愛.」(34)我們看見了主的愛,看見了主如何地愛我們,有新的感動,新的激勵,新的榜樣!這時候這愛的意義完全不同了；聽一句話是一回事,看見一個人實行愛是另一回事.《舊約》時代,以色列人都知道要愛鄰舍如己的命令,但過去那是命令,是教訓,是勉勵.但現在門徒看見主的愛,受了感動,有新的榜樣,新的動力,所以是新的命令!當我們彼此相愛的時候,人們就知道我們是主的門徒.
\newline
\newline
六．主的門徒是站穩在主立場上的人
\newline
\newline
「僕人不能大於主人.」(約十五20)主用「僕人」二字來稱呼門徒,表示「僕人」和「門徒」在這裡有同樣的意義.僕人不能大於主人,主人如果受逼迫,門徒也一定跟著受逼迫,如果主人站穩的立場,僕人也是那立場,如果主人為他的立場受逼迫,僕人也同樣受逼迫.主的門徒和主,站在同一個立場,這是很重要的；基督徒在社會上,面對潮流沖激之時,我們的立場是甚麼?我們不是騎牆派,牆頭草,基督徒有立場,不是你,我個人的立場,而是站在基督的立場上!我們不是看或跟潮流,而是看或跟耶穌!如果潮流中有合乎基督真理的,就應順潮流；如果是明顯違反基督真理的,就應反潮流.不是以潮流來決定我們,而是我們決定跟耶穌!一個有生命的基督徒,在世界的潮流裡,要跟主耶穌,走主的路,按照祂的道來決定我們的方向!超越潮流!這是主的門徒很重要的一件事.
\newline
\newline
(約十九38-40)中有兩個人原來是暗暗地作門徒,後來主耶穌釘了十字架,主在十字架上的表現大大地感動了他們的心,他們決定挺身而出,作公開的門徒；十字架改變了他們.有的人今天還是暗暗作門徒,覺得公開作門徒代價太大,但當他們更清楚認識主時,他們就壯膽,決心勇敢地,公開地承認耶穌,跟從耶穌,見證耶穌!當一個人真正認識十字架時,心的深處受感動,心就溶化了.保羅從逼迫教會到為教會受逼迫,我們看到:這樣的改變,後面都有很堅強的,很清楚的,很徹底的信仰,願主給我們這樣的心志:我們是基督的門徒,真門徒.
\newline
\newline


\newpage

\section{第71屆~港九培靈會~滕近輝~第七講}
\label{sec:18}
\textbf{真愛主}
\newline
\newline
連結: \href{https://www.hkbibleconference.org/session-message/view/18}{\texttt{https://www.hkbibleconference.org/session-message/view/18}}
\newline
\newline
\hyperref[sec:19]{< < < PREV SERMON < < <}
~
\hyperref[sec:toc]{[返目錄]}
~
\hyperref[sec:17]{> > > NEXT SERMON > > >}
\newline
\newline
約翰福音 12:1-12:8
\newline
\begin{longtable}{cl}
\hline
\hline
章節 & 經文 (和合本修訂版)\\
\hline
12:1 & \begin{tabularx}{0.7\textwidth}{X} 逾越節前六天，耶穌來到伯大尼，就是他使拉撒路從死人中復活的地方。 \end{tabularx} \\ \\ \relax
12:2 & \begin{tabularx}{0.7\textwidth}{X} 有人在那裡為耶穌預備宴席；馬大伺候，拉撒路也在同耶穌坐席的人中間。 \end{tabularx} \\ \\ \relax
12:3 & \begin{tabularx}{0.7\textwidth}{X} 馬利亞拿著一斤極貴的純哪噠香膏，抹耶穌的腳，又用自己頭髮去擦，屋裡充滿了膏的香氣。 \end{tabularx} \\ \\ \relax
12:4 & \begin{tabularx}{0.7\textwidth}{X} 有一個門徒，就是那將要出賣耶穌的加略人猶大，說： \end{tabularx} \\ \\ \relax
12:5 & \begin{tabularx}{0.7\textwidth}{X} 「為甚麼不把這香膏賣三百個銀幣去賙濟窮人呢？」 \end{tabularx} \\ \\ \relax
12:6 & \begin{tabularx}{0.7\textwidth}{X} 他說這話，並不是關心窮人，而是因為他是個賊，又管錢囊，常偷取錢囊中所存的。 \end{tabularx} \\ \\ \relax
12:7 & \begin{tabularx}{0.7\textwidth}{X} 耶穌說：「由她吧！她這香膏本是為我的安葬之日留著的。 \end{tabularx} \\ \\ \relax
12:8 & \begin{tabularx}{0.7\textwidth}{X} 因為常有窮人和你們在一起，但是你們不常有我。」 \end{tabularx} \\ \\
[1ex]
\hline
\hline
\end{longtable}
今天我們中間有很多愛主的人,愛主是最大的福份；剛才美蘭姊妹獻詩的時候,聯想到她的父親張有光牧師,我一直記得他低沉的歌聲,美蘭唱詩的恩賜自此而來；美蘭不但承繼了父親唱詩的恩賜,也跟著父親走上了事奉的道路!從「張有光牧師紀念集」中我得知,他辭世前心臟病已經很嚴重,但還在計劃出去傳揚主的福音,直到最後一刻,他的心還在為主.我們也紀念張牧師生前對港九培靈會貢獻了許多精神,力量和心血.感謝神!歷世歷代有很多敬虔愛主的人.
\newline
\newline
愛主是最大的福份,「我實在告訴你們,普天之下,無論在甚麼地方傳這福音,也要述說這女人所作的,以為記念.」(可十四9)主親自說:愛祂的,要得到永遠的紀念.我們的主很欣賞,悅納我們蒙恩的人的愛祂!祂的十字架的愛沒有落空,在你,我的身上沒有落空!我們今天之所以有愛主的心,完全是因為主先愛了我們!
\newline
\newline
你們有否聽到過三個字引起教會的大復興呢?在抗戰期間,我逃難到了中國的大後方,來到安徽的鶴州,聽到一個很感動人的事跡:有一次教會開培靈會,開始講道時,講員不能來到,臨時臨急請一位長老上台講話,這位長老沒有甚麼口才,但非常愛主,這是大家都知道的；到了台上,看見這麼多的人,更講不出話來,只結結巴巴地講:「我...愛...主!」大家受感動流淚,開聲禱告,很熱切地禱告!沒有人講道,只有三個字:「我愛主!」大家知道他很愛主,三個字的信息引起了復興!主的愛使我們復興!真願意在這培靈會中,對主的愛火被再一次挑旺起來!
\newline
\newline
一．馬利亞對主有真愛
\newline
\newline
大家都很熟悉(約十二1-8)這段經文,馬利亞將一斤極貴的真哪噠香膏抹耶穌的腳,跟著又拿自己的頭髮去擦,全屋滿了香氣,這是馬利亞對主真愛的表現.馬利亞「在耶穌腳前坐著聽祂的道」(路十39),她愛主,所以愛主的話,得到主的稱讚,說她選擇了那上好的福份,是不能奪去的.《約》中有三十二次用「愛」一字,其中二十二次是講主的愛,另十次是講愛主,這是本卷中顯著的重點.如果一個人讀《約》後沒有感受到愛的信息,那麼他一定沒有看通看透.在這一章裡,我們一同來思想,來紀念馬利亞對主的愛.
\newline
\newline
愛主和愛主的道是不能失去的,「有了我的命令又遵守的,這人就是愛我的；愛我的必蒙我父愛他,我也要愛他,並且要向他顯現.......人若愛我,就必遵守我的道；我父也必愛他,並且我們要到他那裡去,與他同住.」(十四21,23)愛主的人必得到四重回報:
\newline
\newline
(1)必蒙我父愛他,神愛愛主的人.
\newline
\newline
(2)主自己也要愛他.
\newline
\newline
(3)特別愛主,對主的認識深於其他人的人,主要向他顯現；主復活之後,首先向抹大拉的馬利亞顯現,抹大拉的馬利亞原來是個罪人,但信主悔改之後非常愛主,她愛主的心深切,主首先向她顯現!
\newline
\newline
(4)神和主要和他同住,「同住」是一種密切的關係,愛主的人與神,與主之間有與眾不同的密切關係.
\newline
\newline
二．馬利亞對主的愛是付代價的愛
\newline
\newline
她為了愛主,付了代價,極貴的香膏要用很多的錢才能買到,她不是一個有錢的人,但為了愛主付出了所有的!翻開史達德宣教士的傳記:第一頁是他在英國倫敦皇宮似的住宅,第二頁是他在非洲宣教住的小茅屋,下面的經文是:「他看為基督所受的凌辱比埃及的財物更寶貴.」(來十一26)這位宣教士真是把他的一切獻上了,他付了很大的代價；他的見證感動了許許多多人作宣教士,愛主的人感動,影響別人愛主.愛主:在主那方面,在人那方面都不落空,愛主得主永遠的紀念!大衛說:「我不肯用白得之物作燔祭獻耶和華我的神.」(撒下廿四24)他覺得用白白得來的獻給神是沒有價值,他要付代價,這是大衛愛神的心.
\newline
\newline
感謝神,歷世歷代以來有許多的人實在為主付了代價!神指著那些愛祂的人說:「我必使他們在我殿中,在我牆內,有記念,有名號,比有兒女的更美.我必賜他們永遠的名,不能剪除.」(賽五六5)這是神的應許!可見神是那麼重視,那麼寶貴愛祂的人!世上的事都是過眼煙雲,但是愛主的人有永遠的紀念,這是值得我們謹記的事!
\newline
\newline
三．馬利亞對主的愛是不計算的愛
\newline
\newline
馬利亞用一整斤真哪噠香膏抹主的腳；她原可以只用一點兒,一點兒香膏已經很香了,如果我們看到也會這樣想,但馬利亞沒有這想法,她將一斤的真哪噠香膏抹在主的腳上.保羅講到愛的時候,說:愛是不計算.真的愛不會斤斤計較,有的人對主處處打算盤,在奉獻上掙扎,為主很吝嗇,為自己卻很大方,不打算盤!馬利亞不是這樣,她沒有去想她不需要用一斤香膏這件事,她很自然的,很甘心的,完全用在主的身上!
\newline
\newline
在教會歷史上有一個真正屬靈的偉人,約翰衛斯理:他是英國一間最好的大學的畢業生；剛開始事奉時,薪水很低,但他拿出百分之八十來奉獻,後來升到年薪一千英磅,他仍然奉獻百分之八十!大約在二百五十年前人們發現了他的遺囑:內中只提到他的牧師袍,書及保留的懷錶,這就是他的遺產!但他一生中經過他手的錢不計其數,用在主的工作上,他對主有不計算的愛,沒有保留的愛!我們為著他而感謝神.
\newline
\newline
四．馬利亞對主的愛是謙卑的愛
\newline
\newline
她將極貴的真香膏抹在主的腳上,為甚麼不抹在主的頭上呢?相信她覺得自己不配,所以將它抹在主的腳上,這是何等謙卑的愛!亞伯拉罕事奉神也是持非常謙卑的態度,當耶和華再一次向他顯現時,他俯伏在地(創十八2-5,7-8),何等感恩,謙卑,他看事奉神是一件蒙恩的事!他把一隻牛犢看成一點點,他真的感覺這牛犢是「一點點」,這是他謙卑的心態,他覺得這樣事奉神是理當如此的!不是他有特別的好,而是他有一謙卑的心,得到主的悅納!至今幾千年了,我們還在讀,還在紀念.
\newline
\newline
我們為主作了一點點事,就很想人知道,很想人提一提,提得不夠,心裡就很不舒服,這是我們常有的心態.在這上面,我們從馬利亞身上可以學習功課.
\newline
\newline
五．馬利亞用她的頭髮擦乾主的腳
\newline
\newline
頭髮是女子最重視的,是女子的榮耀!她竟然用頭髮去擦主的腳,這真是將主放在第一位,這是高舉主的愛.在法國的君王中,路易十五送了一枚戒指給她的王后,上面刻著:神第一,法國第二,妳第三.這是像馬利亞一樣將主放在最高地位,這是難得的心態!六．馬利亞對主真的愛與猶大對主假的愛,成了強烈的對比,猶大說的很動聽:「這香膏為甚麼不賣三十兩銀子賙濟窮人呢?」(約十二5)把香膏賣掉,一大筆錢獻給主,放在錢袋裡不好嗎?他用假的愛想別人贊同他,表面上他很有愛心,記念窮人,關心窮人,但心底下,心中乃是出自私心,完全不是那回事.主說:「因為常有窮人和你們同在,只是你們不常有我.」行善很多機會,我們可以隨時行善,但有時候愛主的機會過去就沒有了,所以要抓住機會來愛主.
\newline
\newline
有一個弟兄,他有一個很強的遺憾；他太太過世之前,他沒有愛她,太太過身之後感到很對不起她,他感到他應該更多愛他的太太,有很深的痛苦一直壓在他心上,但是沒有機會了,沒有機會了!所以我們中間,丈夫對太太,太太對丈夫,要抓住機會愛對方,愛周圍的人,抓住機會愛主!因為我們不常有愛主的機會!
\newline
\newline


\newpage

\section{第71屆~港九培靈會~滕近輝~第八講}
\label{sec:17}
\textbf{真敬拜 -- 聖禮的內在化}
\newline
\newline
連結: \href{https://www.hkbibleconference.org/session-message/view/17}{\texttt{https://www.hkbibleconference.org/session-message/view/17}}
\newline
\newline
\hyperref[sec:18]{< < < PREV SERMON < < <}
~
\hyperref[sec:toc]{[返目錄]}
~
\hyperref[sec:16]{> > > NEXT SERMON > > >}
\newline
\newline
約翰福音 4:23-4:24
\newline
\begin{longtable}{cl}
\hline
\hline
章節 & 經文 (和合本修訂版)\\
\hline
4:23 & \begin{tabularx}{0.7\textwidth}{X} 時候將到，現在就是了，那真正敬拜父的，要用心靈和誠實敬拜他，因為父要這樣的人敬拜他。 \end{tabularx} \\ \\ \relax
4:24 & \begin{tabularx}{0.7\textwidth}{X} 神是靈，所以敬拜他的必須用心靈和誠實敬拜他。」 \end{tabularx} \\ \\
[1ex]
\hline
\hline
\end{longtable}
真正敬拜神要用心靈和誠實,「誠實」一詞原文可譯作「真理」,因為心靈和誠實在意義上是相似的,但心靈和真理是很清楚的兩件事,讓敬拜神的意義更加完全,更加深入.
\newline
\newline
一．敬拜真神
\newline
\newline
有人拜大山,大河,大樹,但我們不是甚麼都拜,是按照真理來敬拜神,是敬拜真神!《約》中三次提到真神,我們願意知道我們所拜的真神是怎樣的呢?為甚麼祂值得我們敬拜呢?「並且穿上新人；這新人是照著神的形像造的,有真理的仁義和聖潔.」(弗四24)神是真理,仁義,聖潔的,這是我們所敬拜的真神,如果一個人拜這位真神,他就會按照真理來生活,就會追求仁義的生活,也學習聖潔的生活!一個人崇拜的對象對他關係很大,可以決定他人生的價值觀,可以決定他人生追求的方面和路線,拜甚麼樣的神,這個人就是甚麼樣的人!有的人拜財神,將財放在第一,將他的價值觀表達出來.中國人很多拜關公,把關公當作財神來拜,其實關公很冤枉；關公是一個重義輕財的人,竟然成了財神!一個人將錢放在第一,他就會不擇手段來賺錢,很多人用不法手段,所以一個人崇拜的對象決定他的價值觀.如果一個人拜真理,仁義,聖潔的神,他的人生一定有一高尚目標!《聖經》叫我們敬拜真神,這樣的真神真是值得我們來敬拜!
\newline
\newline
怎樣來敬拜這位真神呢?首先要尊神為大,我們說崇拜,「崇」即崇高,將我們崇拜的對象放得很高,放在首位上,我們遵從祂,順服祂!《舊約》告訴我們:耶和華神是王；《新約》說,耶穌基督是王.大家記得彼拉多審問耶穌基督的時候,他說:「你是猶太人的王麼?」耶穌基督面對這問題很難回答；當時控告祂,要定祂罪的理由是說祂要造反,如果祂說祂是王,等於承認是造反的罪；我們的主回答:「你說我是王.我為此而生,也為此來到世間,特為真理作見證.凡屬真理的人就聽我的話.」(約十八37)祂承認是王,但是真理的王,不是政治的王；如果主耶穌是政治的王,祂早已過去,政治的權勢在歷史上興起衰敗,興起又衰敗,耶穌基督是真理的王,所以二千年來祂仍然是王,在許多人心目中祂是真理的王!主耶穌釘十字架時,十字架上有一牌子用希伯來文,希臘文,及拉丁文寫著:「猶太人的王.」他們當時這樣作的目的是要諷刺祂,但無形中,無意中為耶穌作了見證!亦成為一預言:希伯來文代表宗教範圍,希臘文代表文化範圍,拉丁文代表政治的範圍,三種文字代表三個範圍,耶穌基督在這三個範圍中都成了王!這個無意的見證實現了.宗教界信耶穌的人是最大的宗教,據耶魯大學一位很有名望,地位的學者,統計學家說:現在全世界六十億人口中有二十億是信耶穌的.宗教界耶穌是王!今日全世界的文化中,西方文化一般來說是有領導地位,而影響西方文化最大的因素是耶穌基督.在政治圈子內耶穌基督也是王嗎?是!因為主耶穌講到人的尊嚴,人的尊嚴是民主制度的基礎.「你們要完全,像你們的天父完全一樣.」(太五48)在《太》中,主耶穌用了「天父」十八次,即人類是天父上帝的兒女,這是人的尊嚴.今天在政治界最大的影響力是民主制度,那是很清楚地受了耶穌基督的影響,所以說:人的尊嚴是民主制度的基礎.十字架上的那牌子是說了預言,而且是已經應驗了的預言!
\newline
\newline
耶穌基督是王,是立憲和專制的君王!耶穌基督自己就是真理的王,真理就是祂的憲法,祂本性就是真理,仁義,聖潔,《聖經》說到神不能背乎自己,所以耶穌基督要按照祂的真理來作王,祂是立憲的君王!另一方面祂又是有絕對權利的君王,因為祂是創造之主,祂是主,祂有主權!感謝天父,我們敬拜的神是王,這位王是真理,仁義,聖潔的王不違反祂自己的本性,所以我們很放心.祂有絕對的權利!祂的權利是真理的權利,仁義的權利,聖潔的權利!這樣的權利愈強愈好,我們口服心服地在這樣的權利之下,將祂高舉起來.我們崇拜祂,將祂放在人生第一位,讓主居首位.有這樣的一位王來掌管自己的人生,自己就有救了,人生就不會墮落了,不會偏行己路了,不會誤入歧途了,就能走在人生的大道上!感謝神,我們可以這樣地來高舉祂,順服祂,讓祂居首位!
\newline
\newline
「我已經與基督同釘十字架,現在活著的不再是我,乃是基督在我裡面活著；並且我如今在肉身活著,是因信神的兒子而活；他是愛我,為我捨己.」(加二20)使徒保羅說他已與基督同釘十字架,現在是基督在他裡面活著.「不再是我,乃是基督」這是保羅的人生觀.「不再是我」,是否定自我中心,「乃是基督」是確定基督中心.誰願意否定自己呢?那就是知道否定了自我之後有更好的取代者,有更好的自己,知道讓耶穌基督作我人生的中心好過我自我中心!
\newline
\newline
從前,有一港大學生喜歡寫作,有一次寫了一篇自己很滿意的文章,去見他的教授,教授是一名作家,他請教授為他斧正.為甚麼自己滿意的文章要找人斧正呢?只有一個原因:知道教授比他程度更高,由教授來改一定愈改愈好,為了令自己的作品更好,所以請斧正:否定自己,確定教授!為甚麼基督徒說:「不再是我,乃是基督」?因為知道基督中心比自我中心有更好的人生表現!所以我們要將基督高舉作我們人生中心,在我們人生居首位.施洗約翰說:「他必興旺,我必衰微.」(約三30)基督必要興旺,「我」必衰微消失,即:要基督大,自己小.如果基督在一個人的身上顯大,一個人表現了基督,這個人的人生一定很美!以自我為中心的人生常常很醜惡,很多時候自己想大,要作主,但隨著靈性長進,發現了自我中心的壞處,發現了神的偉大在基督裡顯明出來,小小「我」藏在基督裡了!「自我」居首位就問題多多,世界上所有問題都是「自我」產生的.
\newline
\newline
二．榮耀神
\newline
\newline
「人憑著自己說,是求自己的榮耀；惟有求那差他來者的榮耀,這人是真的,在他心裡沒有不義.」(約七18)我們要求神的榮耀,追求榮耀神!「你們或喫或喝,無論作甚麼,都要為榮耀神而行.」(林前十31)榮耀神,這是人生最高原則.因為神配得榮耀,因為一個榮耀神的人一定榮神益人,一個將神放在第一的人一定關心別人,不做神不喜悅的事,一個愛神的人一定愛人!榮耀神形成了一個敬拜的人生；這不單單是參加一星期兩小時的敬拜聚會,乃是追求榮神,高舉祂!過一個有敬拜的人生.
\newline
\newline
「我又觀看,見羔羊站在錫安山,同他又有十四萬四千人,都有他的名和他父的名寫在額上.......他們在寶座前,並在四活物和眾長老前唱歌,彷彿是新歌；除了從地上買來的那十四萬四千人以外,沒有人能學這歌.」(啟十四1,3),這新歌沒有人能學會,如果是口唱的歌,再難唱都有人可以學會.那為甚麼呢,因為這不是口唱的歌,而是人生之歌,這是生命之歌!生活之歌,只有那樣人生的人才能唱這新歌!凡事追求榮耀祂,將神放在第一,我們的人生一定不斷地提昇,提昇,再提昇!真正的崇拜就有這作用.
\newline
\newline
三．讚美神
\newline
\newline
「和散那!奉主名來的以色列王是應當稱頌的!」(約十二13)這是對神的讚美,感謝.我們讚美神,到底讚美神的甚麼呢?「四活物各有六個翅膀,遍體內外都滿了眼睛.他們晝夜不住的說:聖哉!聖哉!聖哉!主神是昔在,今在,以後永在的全能者.......因為您創造了萬物,並且萬物因您的旨意被創造而有的.」(啟四8,11下)約翰提到在天上對神的敬拜是怎樣的?
\newline
\newline
(1)讚美神的美德,神是聖潔的.「惟有你們是被揀選的族類,是有君尊的祭司,是聖潔的國度,是屬神的子民,要叫你們宣揚那召你們出黑暗入奇妙光明者的美德.」(彼前二9)彼得宣揚祂的美德,宣揚神的美德是基督徒應有的責任；不但口裡宣揚,同時要生活上彰顯神的美德.
\newline
\newline
(2)讚美神無所不能的能力.
\newline
\newline
(3)讚美神的榮耀.
\newline
\newline
(4)讚美神的智慧,神創造萬物中有很深奧的智慧,大自然的奧秘實在驚人,裡面充滿著智慧.得救的人將來到了天上得賞賜,要與基督一同作王,管理這奇妙,偉大,無限的宇宙,科學家不停的研究,無止境地發現神的奇妙.
\newline
\newline
(5)讚美神的愛；神的寶座有虹圍繞(啟四3),《聖經》中虹象徵神的愛,虹有七種顏色,象徵神愛內容的豐富.虹是神與人立約的象徵,神指著空中的虹作一證據,表示與人所立恩典之約(創十一).虹圍繞神愛的寶座,恩典的寶座,何等美麗的象徵!神的美德,能力,榮耀,智慧,愛,構成我們讚美的內容.
\newline
\newline
用心靈敬拜神,《啟》中歌唱讚美時,有八次講「阿們」,讓我們立刻想到韓德爾寫的彌賽亞神曲中的「阿們頌」,這是敬拜神的人從心底發出:「是的,是的!」我們口唱心和地讚美神,保羅說:「以神的靈敬拜」(腓三3),我們的心靈被聖靈感動,從心靈深處發出的敬拜,絕對不是儀式上的敬拜.喜樂的敬拜；一個人從心裡發出喜樂的敬拜,絕不是不得已的；保羅和西拉在聖靈的帶領下越過亞洲的邊界進入歐洲,想不到抵達第一站(腓立比)就坐監,被棍打,戴手銬,戴腳鐐,被下到嚴重罪犯的內監,但他們在半夜最黑暗之時仍能唱詩,禱告,讚美神!神回應,結果神蹟事發生!(正如(賽六一1)所預言的:「主耶和華的靈在我身上；......報告被擄的得釋放,被囚的出監牢.」)福音到了歐洲,第一站就是讓那些囚犯聽見福音,悔改信主,成了歐洲第一間教會第一批會友,這是何等奇妙!這事證實了耶穌基督的使命,這兩位使者半夜喜樂地唱詩,禱告,讚美神,從心裡發出來的感謝,讚美產生神蹟,地大震動,神回應,使福音傳遍歐洲!以後,基督教成為羅馬帝國的國教,而且,羅馬帝國接受了基督教以後,做的第一件事是:決定廢除奴隸制度；羅馬帝國本身建立在奴隸制度基礎上,他們享受奴隸的服侍,這特權誰也不肯放棄!這真是被囚的得釋放,被擄的得自由,受壓制的得自由!這是主耶穌的工作,使保羅的感謝讚美成為大見證.我相信他們的感謝,讚美,就是最好,最好的敬拜!
\newline
\newline


\newpage

\section{第71屆~港九培靈會~滕近輝~第九講}
\label{sec:16}
\textbf{真事奉}
\newline
\newline
連結: \href{https://www.hkbibleconference.org/session-message/view/16}{\texttt{https://www.hkbibleconference.org/session-message/view/16}}
\newline
\newline
\hyperref[sec:17]{< < < PREV SERMON < < <}
~
\hyperref[sec:toc]{[返目錄]}
~
\hyperref[sec:9]{> > > NEXT SERMON > > >}
\newline
\newline
約翰福音 7:18-7:18
\newline
\begin{longtable}{cl}
\hline
\hline
章節 & 經文 (和合本修訂版)\\
\hline
7:18 & \begin{tabularx}{0.7\textwidth}{X} 憑著自己說的人是尋求自己的榮耀；但那尋求差他來那位的榮耀的人，他是真誠的，在他心裡沒有不義。 \end{tabularx} \\ \\
[1ex]
\hline
\hline
\end{longtable}
「人憑著自己說,是求自己的榮耀；惟有求那差他來者的榮耀,這人是真的,在他心裡沒有不義.」(約七18)這是真的事奉.這節《聖經》有四個重點:
\newline
\newline
(1)差者:神,神是差遣者.
\newline
\newline
(2)被差者:在這裡是指主耶穌自己而言,但這位被差的主耶穌也說:「父怎樣差遣了我,我也怎樣差遣你們.」(約廿21)所以主完成使命的心志是我們的榜樣.
\newline
\newline
(3)事奉者的心態:求差者的榮耀,不是求自己的好處,地位,利益,享受,滿足,乃是求神的榮耀!
\newline
\newline
(4)真事奉:心裡沒有不義.
\newline
\newline
我讀《約》時,發現主耶穌三十九次用「差」字,表現主耶穌一生念念不忘祂是父差來的,祂有很強的使命感:父差我,我要完成這使命!大家記得主十二歲時在聖殿講的話,小小年紀已經以父的事為念(路二49)；主在十字架上臨死時,大聲呼喊:「成了!」成了甚麼呢?成就了父的託付!主耶穌一生從開始到最後,一直表達了這強烈的使命感!一個人死在十字架上是多麼痛苦,但主耶穌想的是完成了父神的託付!
\newline
\newline
我們想想主耶穌怎樣事奉神呢?同時我們想一想主耶要我們怎事奉祂呢?「若有人服事我,就當跟從我；我在那裡,服事我的人也要在那裡；若有人服事我,我父必尊重他.」(約十二26)我們知道主耶穌在那裡,我們也要在那裡.那麼主耶穌在那裡呢?「我實實在在的告訴你們,一粒麥子不落在地裡死了,仍舊是一粒,若是死了,就結出許多子粒來.愛惜自己生命的,就失喪生命；在這世上恨惡自己生命的,就要保守生命到永生.」(24,25)主耶穌將自己的生命擺上,祂捨己放下了自己的生命,好像一粒麥子落在地裡,死了,祂獻上自己,以致於死!結果從祂的死結出許多子粒來,世界上千千萬萬得救的人都因為主的死而得救；這是主所走的路.一粒麥子所經的路,是主所在的地方.主在那裡,服事祂的人也在那裡；放下自己,捨己,不是為了自己而活著,而是為了完成主的使命而活著!如同主念念不忘父給祂的使命而完成一樣!我們與主同行,主的路是我們的路,主的所在是我們的所在最後,主的結果是我們的結果!當我們這樣與主走路的時候,我們也會結出許多子粒來!一九四九年,一位個子矮小的日本佈道家,被邀請到英國愛丁堡一間禮拜堂講道,許多許多的人來聽,因他是一粒麥子落在地裡死了的佈道家,他在東京貧民區,貧民窟傳道,與那些貧民打成一片,和他們一起住,一起生活.有一件事很感動人心:有一個貧民害了一種特別的病,每晚睡覺時需要一人抓住他的手才能入睡,這位傳道人每天抓住他的手睡,一連六個月,後來這人恢復了健康.當日本侵略中國的時候,他是唯一站出來反對日本政府,日本軍隊侵略中國的人,為此坐監直到戰事結束.影片《一粒麥子》就是描述他的事跡,所以英國人蜂擁聽他講道.一個真事奉的人是學習捨己的人,願意自己像那一粒麥子,這是主的呼召:「若有人服事我,就當跟從我；我在那裡,服事我的人也要在那裡.」
\newline
\newline
「我就是那在曠野有人聲喊著說:『修直主的道路』,正如先知以賽亞所說的.」(約一23)這是施洗約翰的榜樣.有人問他,你是《舊約》預言要來的那個以利亞嗎?你是要來的先知嗎?他回答:「不,我是曠野的聲音,只是一個見證者,一個服事者,我不是主角!」約翰的表現,顯示他有何等寶貴的心態.主耶穌和約翰,都給我們留下這樣寶貴的榜樣.「願你們平安!父怎樣差遣了我,我也照樣差遣你們.」(約廿21)我們看到主對門徒講的話,當門徒接受這差遣時,他們需要有甚麼條件呢?他們完成使命的因素是甚麼?
\newline
\newline
一．成功的保證
\newline
\newline
被主差遣,我們能完成祂的使命嗎?我們是這樣的軟弱,一想到困難,便會說:「算了,算了!與其不能成功,不如不開始!」所以心靈中的成功的保證感很重要,主耶穌向門徒顯現,將大使命給他們時,他們是很怕,很怕的,怕猶太人把門關了(關:鎖之意),怕猶太人抓他們,殺他們,是懼怕的心態!但是主向他們說:「願你們平安.」那時他們還沒有認出主來,當主第二次說:「願你們平安.」並把手和肋旁指給他們看時,感覺不同了；他們看見主手中的釘痕,看見主肋旁的槍傷,看見復活的主,他們滿了喜樂!主死而復活了,主勝過了死亡,祂與我們同在!這位主差遣他們,他們很放心,很有把握!因這位主能幫助他們完成使命.弟兄姊妹們,如果我們將自己獻給主來完成祂的使命,為祂而活,我們心中必須有這個認識:是復活的主,得勝的主差遣我!祂能給我力量,祂能幫助我,我的事奉不會落空!我自己不能作甚麼,但主能!這個保證感是每一個事奉主的人必須具備的條件.
\newline
\newline
二．被愛激勵
\newline
\newline
主耶穌把手和肋旁指給門徒看,他們再一次真實地感受到十字架的愛:這就是為他們被釘的雙手,這就是為他們被刺傷的肋旁!他們立刻心中充滿了感激:「這釘痕是為我,這槍傷是為我.」這是十字架的愛,這愛激勵他們,感動,推動他們!
\newline
\newline
三．聖靈的能力
\newline
\newline
在英國教會的歷史裡,有一位作宣教士的姊妹,她的事跡很感動人,很多人讀過她的傳記,這就是司瑪莉.她小時是一個很怕羞的女孩子,但長大後被基督的愛充滿,感動,激勵,奉獻自己到非洲去做宣教士,神大大使用她.她去世時,英國差會要派二十個弟兄去接替她的工作；一個被人看來軟弱,怕羞的女子,竟然成為這樣剛強,這樣有作為的宣教士.七十年代,我被邀去剛果講道,到達後首先被帶去看一個宣教士的墳墓；這是由美國宣道會差遣去的,同一次差去了三個人.到達後,第一個人由於那裡的水土較難適應,要求調到另一處了；第二個受不了,回美國去了；第三個留下,因水土不服,染病死在剛果.帶隊的說:「這裡建教會前,先建了一個宣教士的墳墓!」表面看,這三個宣教士是悲劇,但當這消息傳回美國時,有十一位弟兄姊妹站起來要步他們的後塵,繼承他們的遺志去傳福音!三個倒下去,十一個上前線,前仆後繼!我去時剛果宣道會已經有五萬多會友,一百多宣教士.去年我到南非參加一個華人差傳會議,聽到當地教會領袖講道:講道中有一報告,原來非洲的基督徒佔全非洲人口百分之十七,宣教士有二萬二千多人,真是令人難以置信!他們接受了福音,起來傳福音!一粒麥子落在地裡死了,沒有白白的死了,結出許多的子粒來.這些宣教的先鋒,都是被十字架愛的激勵而去.
\newline
\newline
「說了這話,就向他們吹一口氣,說:『你們受聖靈!』」(廿22)主耶穌向門徒吹了一口氣,象徵性地表明聖靈在五旬節時降臨,他們得到聖靈的能力,出去傳福音,大有功效.復活的主藉著聖靈的能力來作成工作,聖靈的能力令我們驚奇.五年前,我到越南的一個「冬令會」講道,有二百多青年聚集在一起,我向他們講奉獻的信息；他們帶我去看他們的教會,現在越南教會由一九五七年的五萬多增加到五十多萬,在共產政權之下基督徒人數居然增加了十倍,這是誰也想不到的,我心裡真是感受到神的恩典聖靈的能力.那些聚會的青年人很熱心,很愛主.有一位身體軟弱的姊妹帶領他們,因為工作的壓力,這位姊妹每晚只能睡四至五小時,但神保守她,祝福她的工作,神的能力處處彰顯!
\newline
\newline
五旬節聖靈降的時候,門徒被聖靈充滿,講起別國的話來.聖靈充滿與講起別國的話有甚麼關係呢?五旬節為甚麼有這現象呢?原來,這是事先的,預言性的保證,聖靈來了,藉著聖靈的工作和能力,福音要傳到全世界講各種不同話語的人中間去,當時這情形還沒有開始,但先有這一現象,這是神蹟,是預先表明的成功保證!今天我們看見這預表成全了,福音傳遍了全世界!我們知道聖靈的能力是何等浩大,有了聖靈的能力,我們擔起這使命時,心裡有平安,有把握.
\newline
\newline
四．主的榜樣
\newline
\newline
主耶穌說:「父怎樣差遣了我,我也照樣差遣你們.」(21)主的榜樣永遠是一個激勵,一想起主如何完成父的託付,我們真願意對主說:「主啊,我們願意事奉您,與您同行!」不久前的數據:宣教士人數比二十年前增加了三倍,今天全世界華人宣教士大概有八百人.
\newline
\newline
感謝神!這些福音的使者,被主興起來,他們得到這四個因素成為他們的力量,成為他們的保證,成為他們的激勵,得到一個永遠感動他們的榜樣!我們從《約》中看到「真事奉」,我們願意成為主賜福的器皿,我們願意接受主的託付,擔起主給我們的使命,願意跟從主的腳步!全世界千千萬萬的基督徒中不斷有人興起來,興起來,興起來!相信在座中有許許多多弟兄姊妹願意事奉主,你抓住機會,獻上你自己來為主作些工作,這必不落空,有永遠的價值,永遠的意義.
\newline
\newline


\newpage

\section{第71屆~港九培靈會~滕近輝~第十講}
\label{sec:9}
\textbf{真屬靈}
\newline
\newline
連結: \href{https://www.hkbibleconference.org/session-message/view/9}{\texttt{https://www.hkbibleconference.org/session-message/view/9}}
\newline
\newline
\hyperref[sec:16]{< < < PREV SERMON < < <}
~
\hyperref[sec:toc]{[返目錄]}
~
\hyperref[sec:749]{> > > NEXT SERMON > > >}
\newline
\newline
約翰福音 16:13-16:13
\newline
\begin{longtable}{cl}
\hline
\hline
章節 & 經文 (和合本修訂版)\\
\hline
16:13 & \begin{tabularx}{0.7\textwidth}{X} 但真理的靈來的時候，他要引導你們進入一切真理。因為他不是憑著自己說的，而是把他所聽見的都說出來，並且要把將要來的事向你們傳達。 \end{tabularx} \\ \\
[1ex]
\hline
\hline
\end{longtable}
一．聖靈是代求者
\newline
\newline
「我要求父,父就另外賜給你們一位保惠師(或作:訓慰師),叫他永遠與你們同在.」(十四16)主耶穌說到聖靈是保惠師,「保惠師」這名稱的意義之一:代求者.聖靈是一位代求者,祂為我們禱告,而且帶我們進入祂的禱告之中,教導我們如何和聖靈一樣的,有聖靈的心態,負擔的禱告.「況且我們的軟弱有聖靈幫助,我們本不曉得當怎樣禱告,只是聖靈親自用說不出來的歎息替我們禱告.」(羅八26)聖靈的禱告,祂親自用說不出的嘆息為我們禱告,「嘆息」是言詞不容易表達的一種負擔的感情,這是愛的負擔,內心的呻吟.聖靈關心我們,愛我們,知道我們的軟弱,祂知道我們的軟弱不是要定我們的罪,不是要拒絕我們,乃是同情我們,為我們的軟弱作呻吟的禱告.猶大說:「在聖靈裡禱告.」(猶20)意即像聖靈那樣禱告.
\newline
\newline
很久前讀到司布真牧師關於禱告的一個比喻:在神寶座旁有一個禱告的鈴,有繩子繫住的,繩子的一頭墜到地上.地上人禱告時好像抓住繩子來拉,有的人輕輕地拉,神寶座旁的鈴不會響；有的人兩手全力地拉,神寶座旁的鈴聲響了!這裡面有很重要的真理:有的人輕輕忽忽,禱告完了都不知自己求的甚麼,沒有入心,這樣的禱告難蒙垂聽.當一個將力量集中,好像用力拉繩,神就回應.從司布真牧師的傳記中,我們得知他的教會有兩層,上面一層是做禮拜用,下面一層與上面一樣大,上面做禮拜時,下面坐滿了人禱告,上面向人講,下面向神呼求!難怪他的教會那麼蒙福,他成了神大大使用的僕人!他講的這比喻是自己做到的.有聖靈負擔的禱告,這負擔從何而來?從聖靈愛的關心產生!聖靈把這愛的關心也嚐試放在我們心裡,我們不懂得這難學的功課,但一學到,功效就很大,很大!一九九五年,世界五十四個的國家的教會領袖在南韓開會；開會時有人提議:實際地去看一看韓國教會黎明的祈禱會,去那一間看呢?他們不想去看很出名的成了櫥窗式的教會,決定到最近所住酒店的一間禮拜堂去,那是一間衛理公會的教會.黎明前他們去到,見到一千多弟兄姊妹在那裡禱告.他們深深感受到:韓國教會在禱告上蒙恩,他們真是有這一負擔,難怪韓國教會興旺!
\newline
\newline
有負擔的禱告是教會歷史上很多大復興的共同特點,正如十八世紀莫拉維拉弟兄會的復興.原來他們的會友從不同地方,不同背景來聚會,有很多不同傳統,不同觀念,不同重點的人在一起,實很難合得來,得不到合一的心；後來他們決定禱告,十天來集中禱告,聖靈大大作工,他們的心被溶化,有了合而為一的心；跟著產生了大的,差傳的復興:頭十年差傳宣教士到四大洲傳福音,一七三二至一七四二年,他們差出很多宣教士,其中兩位宣教士到了中國西藏；在接著的一百五十年時間中,他們差遣出二千一百五十七位宣教士!他們的復興影響了約翰衛斯理,帶動英國的大復興運動,在約翰衛斯理的日記中:一七三九年一月一日,他和六十位同工一起通宵禱告,清晨三點鐘時,聖靈大大地澆灌下來,他們被聖靈充滿,充滿了敬拜,喜樂,讚美,被稱為「衛理公會的五旬節」.
\newline
\newline
我們又想到十八世紀在美國的大復興,那時有一位教會領袖名愛德華,他是一位學者,心中有感動,有負擔來禱告；記載中,他像保羅一樣,三天三夜禁食祈禱,聖靈大大的工作,跟著發生了北美洲的大復興!在中國教會歷史中,大家知道宋尚節被主使用,他初初領會時,常常早上四時起身禱告,他一一為許多請他代禱的人禱告.所以宋尚節領奮興會大會最顯著的特點就是禱告的熱忱,很多人在聚會中悔改,信主,熱心!這樣,禱告的熱潮展開了,復興的影響愈來愈廣大.我們看到歷代大復興的共同點是禱告,是在聖靈裡禱告!
\newline
\newline
二．聖靈是能力的靈
\newline
\newline
「我實實在在的告訴你們,我所作的事,信我的人也要作,並且要作比這更大的事,因為我往父那裡去.」(約十四12)這句話當時門徒很難了解,門徒怎可能做比主耶穌更大的事呢?但主就是這麼清楚地告訴門徒,為甚麼呢?「然而,我將真情告訴你們,我去是與你們有益的；我若不去,保惠師就不到你們這裡來；我若去,就差他來.」(十六7)主回到父那裡,父就差聖靈來!藉著這位聖靈的能力,屬主的人要做比主自己在世所做更大的事!主在世時受身體,時間的限制,只能同時在一個地方,但升天後聖靈降,使用所有的信徒,他們遍佈世界每一地方,聖靈藉他們所工作的總和比主耶穌在世上所作的更大.所以主耶穌所介紹的聖靈是能力的聖靈.在教會歷史中,我們真是看到聖靈的大能力:教會歷史最早階段受了十次大迫害,但每次迫害之後,信主的人更多,教會更興旺；到了第十次迫害,羅馬帝國接受了基督教作國教!在二百多年迫害的壓力之下,聖靈的能力大大彰顯,信主的人愈來愈多,殺害了的基督徒變成更多的基督!羅馬帝國被基督教的信仰征服了,這是一個極大的能力的表現.
\newline
\newline
在中國,庚子年義和團殺了許多基督徒,宣教士.結果呢?據聖經公會歷史報告:從庚子之後,《聖經》的銷量直線上升.這不是人的能力,是聖靈的能力!一九四九年前,中國只有一百多萬基督徒,而今天很多跡象令我們相信至少有六千萬基督徒,也可能到一億,比一百多萬增加了六十倍,這是在一種壓力的處境之下這樣發展的.
\newline
\newline
印尼大概有四千萬基督徒,是印尼全國人口的五分之一,一個回教國家,竟然有這麼多基督徒,是甚麼能力使這現象出現?!越南宣道會總會報告,從一九七五年的五萬多會友增加到現在的五十多萬會友,更令我們驚訝的是:在越南山地裡面,有許多部落民族,大概有七百萬,其中有五分之一人信了耶穌!聖靈真是的能力聖靈!
\newline
\newline
三．聖靈是結果子的聖靈
\newline
\newline
在(約十五)中,主耶穌說祂是真葡萄樹,我們屬祂的人,是樹上的枝子,枝子連在樹上就果子.葡萄樹的生命之漿是枝子和果子的來源,這生命之漿是否與聖靈有關呢?是否是指聖靈來說呢?我從前不以為然,至到多年前的一天,我在九龍塘宣道會夏天舉行的英語培靈會中,聽到講員湯瑪士說到(約十四,十五,十六)三章重點是講聖靈.但在(十五)中卻只有一次提到聖靈,為甚麼呢?因為葡萄樹的生命之漿是聖靈工作的象徵!「主在我裡面,我在主裡面.」是(十四)的特徵,是聖子指著聖靈的工作說的.基督藉著聖靈住在我們裡面,我們又藉聖靈住在主裡面,我們與主之間有密切的生命的關係,是聖靈將基督的生命流通在我們裡面,是聖靈帶我們進入基督,與祂同活.葡萄與枝子的聯繫是藉著生命之漿,因為有了生命之漿,所以枝子能結果子,而葡萄樹的生命之漿就是指著聖靈說的.(加五22)講聖靈結出九樣果子時,提到聖靈工作的特點,一個人信主之後,聖靈住他裡面,開始與人墮落的天性,與人原來的情慾相爭.人信主之前,想做甚麼就做甚麼,很自由；但重生之後,聖靈進來,人裡面有了爭戰,開始不自由了,心中受限制,受約束了!有的人信了主後覺得真麻煩,真不自由,退去了吧!其實這正是新生命,新生活的開始!那一個孩子長大不受約束,沒有一個可以為所欲為!父母一定要教導他,限制他,聖靈也是這樣.(加五19-22)講到「情慾」是甚麼?是罪惡.世界上各種的罪惡如果表露詳盡,可能不止一萬樣!聖靈和這些事相爭,當我們順服地站在聖靈那邊時,保羅說那時就開始結出聖靈的果子.聖靈的工作是限制,塑造我們,讓我們成為真正屬靈的兒女,結出聖靈的美麗的果子.
\newline
\newline
「只等真理的聖靈來了,他要引導你們明白一切的真理；因為他不是憑自己說的,乃是把他所聽見的都說出來,並要把將來的事告訴你們.」(約十六13)聖靈把我們帶入真理的實際,聖靈是實際的成就者,這信息在《新約》很多地方表明出來；如「我聽見從天上有聲音說:『你要寫下!』......聖靈說:『是的!』」(啟十四13)主耶穌寫給七個教會的書信,都有相同形式:每封書信開始時是主耶穌說話,結束時是聖靈說話.因主耶穌講的是客觀的真理,而聖靈把我們帶進客觀真理的實的勝利中.主耶穌說:「在主裡面而死的人有福了.」(啟十四13)聖靈在人的心裡說:「是的!」這人的心和聖靈同感,這「是的」成了當事人的心態!「他們息了自己的勞苦,作工的果效也隨著他們!」(13)每一個宣教士都有聖靈在他們心中說:「是的!你為主付的代價永不落空!」這就是真屬靈!
\newline
\newline


\newpage

\section{第71屆~港九培靈會~盧家駇~第一講}
\label{sec:749}
\textbf{貪愛世界}
\newline
\newline
連結: \href{https://www.hkbibleconference.org/session-message/view/749}{\texttt{https://www.hkbibleconference.org/session-message/view/749}}
\newline
\newline
\hyperref[sec:9]{< < < PREV SERMON < < <}
~
\hyperref[sec:toc]{[返目錄]}
~
\hyperref[sec:753]{> > > NEXT SERMON > > >}
\newline
\newline


\newpage

\section{第71屆~港九培靈會~盧家駇~第二講}
\label{sec:753}
\textbf{離棄愛心}
\newline
\newline
連結: \href{https://www.hkbibleconference.org/session-message/view/753}{\texttt{https://www.hkbibleconference.org/session-message/view/753}}
\newline
\newline
\hyperref[sec:749]{< < < PREV SERMON < < <}
~
\hyperref[sec:toc]{[返目錄]}
~
\hyperref[sec:756]{> > > NEXT SERMON > > >}
\newline
\newline


\newpage

\section{第71屆~港九培靈會~盧家駇~第三講}
\label{sec:756}
\textbf{取死身體}
\newline
\newline
連結: \href{https://www.hkbibleconference.org/session-message/view/756}{\texttt{https://www.hkbibleconference.org/session-message/view/756}}
\newline
\newline
\hyperref[sec:753]{< < < PREV SERMON < < <}
~
\hyperref[sec:toc]{[返目錄]}
~
\hyperref[sec:760]{> > > NEXT SERMON > > >}
\newline
\newline


\newpage

\section{第71屆~港九培靈會~盧家駇~第四講}
\label{sec:760}
\textbf{忘掉委身}
\newline
\newline
連結: \href{https://www.hkbibleconference.org/session-message/view/760}{\texttt{https://www.hkbibleconference.org/session-message/view/760}}
\newline
\newline
\hyperref[sec:756]{< < < PREV SERMON < < <}
~
\hyperref[sec:toc]{[返目錄]}
~
\hyperref[sec:763]{> > > NEXT SERMON > > >}
\newline
\newline


\newpage

\section{第71屆~港九培靈會~盧家駇~第五講}
\label{sec:763}
\textbf{隨流失去}
\newline
\newline
連結: \href{https://www.hkbibleconference.org/session-message/view/763}{\texttt{https://www.hkbibleconference.org/session-message/view/763}}
\newline
\newline
\hyperref[sec:760]{< < < PREV SERMON < < <}
~
\hyperref[sec:toc]{[返目錄]}
~
\hyperref[sec:764]{> > > NEXT SERMON > > >}
\newline
\newline


\newpage

\section{第71屆~港九培靈會~盧家駇~第六講}
\label{sec:764}
\textbf{缺乏目標}
\newline
\newline
連結: \href{https://www.hkbibleconference.org/session-message/view/764}{\texttt{https://www.hkbibleconference.org/session-message/view/764}}
\newline
\newline
\hyperref[sec:763]{< < < PREV SERMON < < <}
~
\hyperref[sec:toc]{[返目錄]}
~
\hyperref[sec:765]{> > > NEXT SERMON > > >}
\newline
\newline


\newpage

\section{第71屆~港九培靈會~盧家駇~第七講}
\label{sec:765}
\textbf{不傳福音}
\newline
\newline
連結: \href{https://www.hkbibleconference.org/session-message/view/765}{\texttt{https://www.hkbibleconference.org/session-message/view/765}}
\newline
\newline
\hyperref[sec:764]{< < < PREV SERMON < < <}
~
\hyperref[sec:toc]{[返目錄]}
~
\hyperref[sec:766]{> > > NEXT SERMON > > >}
\newline
\newline


\newpage

\section{第71屆~港九培靈會~盧家駇~第八講}
\label{sec:766}
\textbf{抗拒使命}
\newline
\newline
連結: \href{https://www.hkbibleconference.org/session-message/view/766}{\texttt{https://www.hkbibleconference.org/session-message/view/766}}
\newline
\newline
\hyperref[sec:765]{< < < PREV SERMON < < <}
~
\hyperref[sec:toc]{[返目錄]}
~
\hyperref[sec:643]{> > > NEXT SERMON > > >}
\newline
\newline


\newpage

\section{第72屆~港九培靈會~張慕皚~第一講}
\label{sec:643}
\textbf{內虛外暴的末世文化}
\newline
\newline
連結: \href{https://www.hkbibleconference.org/session-message/view/643}{\texttt{https://www.hkbibleconference.org/session-message/view/643}}
\newline
\newline
\hyperref[sec:766]{< < < PREV SERMON < < <}
~
\hyperref[sec:toc]{[返目錄]}
~
\hyperref[sec:644]{> > > NEXT SERMON > > >}
\newline
\newline
馬太福音 24:1-24:51
\newline
\begin{longtable}{cl}
\hline
\hline
章節 & 經文 (和合本修訂版)\\
\hline
24:1 & \begin{tabularx}{0.7\textwidth}{X} 耶穌出了聖殿，正離開的時候，門徒前來，把聖殿的建築指給他看。 \end{tabularx} \\ \\ \relax
24:2 & \begin{tabularx}{0.7\textwidth}{X} 耶穌回應他們說：「你們不是看見這一切嗎？我實在告訴你們，這裡將沒有一塊石頭會留在另一塊石頭上，而不被拆毀的。」 \end{tabularx} \\ \\ \relax
24:3 & \begin{tabularx}{0.7\textwidth}{X} 耶穌在橄欖山上坐著，門徒私下進前來問他：「請告訴我們，甚麼時候有這些事呢？你來臨和世代的終結有甚麼預兆呢？」 \end{tabularx} \\ \\ \relax
24:4 & \begin{tabularx}{0.7\textwidth}{X} 耶穌回答他們：「你們要謹慎，免得有人迷惑你們。 \end{tabularx} \\ \\ \relax
24:5 & \begin{tabularx}{0.7\textwidth}{X} 因為將有好些人冒我的名來，說『我是基督』，並且要迷惑許多人。 \end{tabularx} \\ \\ \relax
24:6 & \begin{tabularx}{0.7\textwidth}{X} 你們也將聽見打仗和打仗的風聲。注意，不要驚慌！因為這些事必須發生，但這還不是終結。 \end{tabularx} \\ \\ \relax
24:7 & \begin{tabularx}{0.7\textwidth}{X} 民要攻打民，國要攻打國，多處必有饑荒、地震。 \end{tabularx} \\ \\ \relax
24:8 & \begin{tabularx}{0.7\textwidth}{X} 這都是災難的起頭。 \end{tabularx} \\ \\ \relax
24:9 & \begin{tabularx}{0.7\textwidth}{X} 那時，人要使你們陷在患難裡，也要殺害你們；你們又要為我的名被萬民憎恨。 \end{tabularx} \\ \\ \relax
24:10 & \begin{tabularx}{0.7\textwidth}{X} 那時，會有許多人跌倒，也會彼此陷害，彼此憎恨； \end{tabularx} \\ \\ \relax
24:11 & \begin{tabularx}{0.7\textwidth}{X} 且有好些假先知起來，迷惑許多人。 \end{tabularx} \\ \\ \relax
24:12 & \begin{tabularx}{0.7\textwidth}{X} 因為不法的事增多，許多人的愛心漸漸冷淡了。 \end{tabularx} \\ \\ \relax
24:13 & \begin{tabularx}{0.7\textwidth}{X} 但堅忍到底的終必得救。 \end{tabularx} \\ \\ \relax
24:14 & \begin{tabularx}{0.7\textwidth}{X} 這天國的福音要傳遍天下，對萬民作見證，然後終結才來到。」 \end{tabularx} \\ \\ \relax
24:15 & \begin{tabularx}{0.7\textwidth}{X} 「當你們看見先知但以理所說的那『施行毀滅的褻瀆者』站在聖地（讀這經的人要會意）， \end{tabularx} \\ \\ \relax
24:16 & \begin{tabularx}{0.7\textwidth}{X} 那時，在猶太的，應當逃到山上； \end{tabularx} \\ \\ \relax
24:17 & \begin{tabularx}{0.7\textwidth}{X} 在屋頂上的，不要下來拿家裡的東西； \end{tabularx} \\ \\ \relax
24:18 & \begin{tabularx}{0.7\textwidth}{X} 在田裡的，不要回去取衣裳。 \end{tabularx} \\ \\ \relax
24:19 & \begin{tabularx}{0.7\textwidth}{X} 在那些日子，懷孕的和奶孩子的就苦了。 \end{tabularx} \\ \\ \relax
24:20 & \begin{tabularx}{0.7\textwidth}{X} 你們要祈求，好讓你們逃走的時候，不遇見冬天或安息日。 \end{tabularx} \\ \\ \relax
24:21 & \begin{tabularx}{0.7\textwidth}{X} 因為那時必有大災難，自從世界的起頭直到如今，從沒有這樣的災難，將來也不會有。 \end{tabularx} \\ \\ \relax
24:22 & \begin{tabularx}{0.7\textwidth}{X} 若不減少那些日子，凡血肉之軀的，就沒有一個能得救；可是為了選民，那些日子將減少。 \end{tabularx} \\ \\ \relax
24:23 & \begin{tabularx}{0.7\textwidth}{X} 那時，若有人對你們說：『看哪，基督在這裡！』或『在那裡！』你們不要信。 \end{tabularx} \\ \\ \relax
24:24 & \begin{tabularx}{0.7\textwidth}{X} 因為假基督和假先知將要起來，顯大神蹟、大奇事，如果可能，要把選民也迷惑了。 \end{tabularx} \\ \\ \relax
24:25 & \begin{tabularx}{0.7\textwidth}{X} 看哪，我已經預先告訴你們了。 \end{tabularx} \\ \\ \relax
24:26 & \begin{tabularx}{0.7\textwidth}{X} 若有人對你們說：『看哪，基督在曠野裡！』你們不要出去；或說：『看哪，基督在內室中！』你們不要信。 \end{tabularx} \\ \\ \relax
24:27 & \begin{tabularx}{0.7\textwidth}{X} 好像閃電從東邊發出，直照到西邊，人子來臨也要這樣。 \end{tabularx} \\ \\ \relax
24:28 & \begin{tabularx}{0.7\textwidth}{X} 屍首在哪裡，鷹也會聚在哪裡。」 \end{tabularx} \\ \\ \relax
24:29 & \begin{tabularx}{0.7\textwidth}{X} 「那些日子的災難一過去，太陽要變黑，月亮也不放光，眾星要從天上墜落，天上的萬象都要震動。 \end{tabularx} \\ \\ \relax
24:30 & \begin{tabularx}{0.7\textwidth}{X} 那時，人子的預兆要顯在天上，地上的萬族都要哀哭。他們要看見人子帶著能力和大榮耀，駕著天上的雲來臨。 \end{tabularx} \\ \\ \relax
24:31 & \begin{tabularx}{0.7\textwidth}{X} 他要差遣天使，用大聲的號筒，從四方，從天這邊直到天那邊，召集他的選民。」 \end{tabularx} \\ \\ \relax
24:32 & \begin{tabularx}{0.7\textwidth}{X} 「你們要從無花果樹學習功課：當樹枝發芽長葉的時候，你們就知道夏天近了。 \end{tabularx} \\ \\ \relax
24:33 & \begin{tabularx}{0.7\textwidth}{X} 同樣，當你們看見這一切，就知道那時候近了，就在門口了。 \end{tabularx} \\ \\ \relax
24:34 & \begin{tabularx}{0.7\textwidth}{X} 我實在告訴你們，這世代還沒有過去，這一切都要發生。 \end{tabularx} \\ \\ \relax
24:35 & \begin{tabularx}{0.7\textwidth}{X} 天地要廢去，我的話卻絕不廢去。」 \end{tabularx} \\ \\ \relax
24:36 & \begin{tabularx}{0.7\textwidth}{X} 「但那日子，那時辰，沒有人知道，連天上的天使也不知道，子也不知道，惟有父知道。 \end{tabularx} \\ \\ \relax
24:37 & \begin{tabularx}{0.7\textwidth}{X} 挪亞的日子怎樣，人子來臨也要怎樣。 \end{tabularx} \\ \\ \relax
24:38 & \begin{tabularx}{0.7\textwidth}{X} 在洪水以前的那些日子，人照常吃喝嫁娶，直到挪亞進方舟的那日， \end{tabularx} \\ \\ \relax
24:39 & \begin{tabularx}{0.7\textwidth}{X} 不知不覺洪水來了，把他們全都沖去。人子來臨也要這樣。 \end{tabularx} \\ \\ \relax
24:40 & \begin{tabularx}{0.7\textwidth}{X} 那時，兩個人在田裡，一個被接去，一個被撇下。 \end{tabularx} \\ \\ \relax
24:41 & \begin{tabularx}{0.7\textwidth}{X} 兩個女人推磨，一個被接去，一個被撇下。 \end{tabularx} \\ \\ \relax
24:42 & \begin{tabularx}{0.7\textwidth}{X} 所以，你們要警醒，因為不知道你們的主哪一天來到。 \end{tabularx} \\ \\ \relax
24:43 & \begin{tabularx}{0.7\textwidth}{X} 你們要知道，一家的主人若知道晚上甚麼時候有賊來，就必警醒，不讓賊挖穿房屋。 \end{tabularx} \\ \\ \relax
24:44 & \begin{tabularx}{0.7\textwidth}{X} 所以，你們也要預備，因為在你們想不到的時候，人子就來了。」 \end{tabularx} \\ \\ \relax
24:45 & \begin{tabularx}{0.7\textwidth}{X} 「那麼，誰是那忠心又精明的僕人，主人派他管理自己的家僕、按時分糧給他們的呢？ \end{tabularx} \\ \\ \relax
24:46 & \begin{tabularx}{0.7\textwidth}{X} 主人來到，看見僕人這樣做，那僕人就有福了。 \end{tabularx} \\ \\ \relax
24:47 & \begin{tabularx}{0.7\textwidth}{X} 我實在告訴你們，主人要派他管理所有的財產。 \end{tabularx} \\ \\ \relax
24:48 & \begin{tabularx}{0.7\textwidth}{X} 如果那惡僕心裡說：『我的主人會來得遲』， \end{tabularx} \\ \\ \relax
24:49 & \begin{tabularx}{0.7\textwidth}{X} 就動手打他的同伴，又和醉酒的人一同吃喝， \end{tabularx} \\ \\ \relax
24:50 & \begin{tabularx}{0.7\textwidth}{X} 在想不到的日子，不知道的時候，那僕人的主人要來， \end{tabularx} \\ \\ \relax
24:51 & \begin{tabularx}{0.7\textwidth}{X} 重重地懲罰他，定他和假冒為善的人同罪，在那裡他要哀哭切齒了。」 \end{tabularx} \\ \\
[1ex]
\hline
\hline
\end{longtable}


\newpage

\section{第72屆~港九培靈會~張慕皚~第二講}
\label{sec:644}
\textbf{在利己主義中邁向滅亡}
\newline
\newline
連結: \href{https://www.hkbibleconference.org/session-message/view/644}{\texttt{https://www.hkbibleconference.org/session-message/view/644}}
\newline
\newline
\hyperref[sec:643]{< < < PREV SERMON < < <}
~
\hyperref[sec:toc]{[返目錄]}
~
\hyperref[sec:645]{> > > NEXT SERMON > > >}
\newline
\newline
馬太福音 24:1-24:51
\newline
\begin{longtable}{cl}
\hline
\hline
章節 & 經文 (和合本修訂版)\\
\hline
24:1 & \begin{tabularx}{0.7\textwidth}{X} 耶穌出了聖殿，正離開的時候，門徒前來，把聖殿的建築指給他看。 \end{tabularx} \\ \\ \relax
24:2 & \begin{tabularx}{0.7\textwidth}{X} 耶穌回應他們說：「你們不是看見這一切嗎？我實在告訴你們，這裡將沒有一塊石頭會留在另一塊石頭上，而不被拆毀的。」 \end{tabularx} \\ \\ \relax
24:3 & \begin{tabularx}{0.7\textwidth}{X} 耶穌在橄欖山上坐著，門徒私下進前來問他：「請告訴我們，甚麼時候有這些事呢？你來臨和世代的終結有甚麼預兆呢？」 \end{tabularx} \\ \\ \relax
24:4 & \begin{tabularx}{0.7\textwidth}{X} 耶穌回答他們：「你們要謹慎，免得有人迷惑你們。 \end{tabularx} \\ \\ \relax
24:5 & \begin{tabularx}{0.7\textwidth}{X} 因為將有好些人冒我的名來，說『我是基督』，並且要迷惑許多人。 \end{tabularx} \\ \\ \relax
24:6 & \begin{tabularx}{0.7\textwidth}{X} 你們也將聽見打仗和打仗的風聲。注意，不要驚慌！因為這些事必須發生，但這還不是終結。 \end{tabularx} \\ \\ \relax
24:7 & \begin{tabularx}{0.7\textwidth}{X} 民要攻打民，國要攻打國，多處必有饑荒、地震。 \end{tabularx} \\ \\ \relax
24:8 & \begin{tabularx}{0.7\textwidth}{X} 這都是災難的起頭。 \end{tabularx} \\ \\ \relax
24:9 & \begin{tabularx}{0.7\textwidth}{X} 那時，人要使你們陷在患難裡，也要殺害你們；你們又要為我的名被萬民憎恨。 \end{tabularx} \\ \\ \relax
24:10 & \begin{tabularx}{0.7\textwidth}{X} 那時，會有許多人跌倒，也會彼此陷害，彼此憎恨； \end{tabularx} \\ \\ \relax
24:11 & \begin{tabularx}{0.7\textwidth}{X} 且有好些假先知起來，迷惑許多人。 \end{tabularx} \\ \\ \relax
24:12 & \begin{tabularx}{0.7\textwidth}{X} 因為不法的事增多，許多人的愛心漸漸冷淡了。 \end{tabularx} \\ \\ \relax
24:13 & \begin{tabularx}{0.7\textwidth}{X} 但堅忍到底的終必得救。 \end{tabularx} \\ \\ \relax
24:14 & \begin{tabularx}{0.7\textwidth}{X} 這天國的福音要傳遍天下，對萬民作見證，然後終結才來到。」 \end{tabularx} \\ \\ \relax
24:15 & \begin{tabularx}{0.7\textwidth}{X} 「當你們看見先知但以理所說的那『施行毀滅的褻瀆者』站在聖地（讀這經的人要會意）， \end{tabularx} \\ \\ \relax
24:16 & \begin{tabularx}{0.7\textwidth}{X} 那時，在猶太的，應當逃到山上； \end{tabularx} \\ \\ \relax
24:17 & \begin{tabularx}{0.7\textwidth}{X} 在屋頂上的，不要下來拿家裡的東西； \end{tabularx} \\ \\ \relax
24:18 & \begin{tabularx}{0.7\textwidth}{X} 在田裡的，不要回去取衣裳。 \end{tabularx} \\ \\ \relax
24:19 & \begin{tabularx}{0.7\textwidth}{X} 在那些日子，懷孕的和奶孩子的就苦了。 \end{tabularx} \\ \\ \relax
24:20 & \begin{tabularx}{0.7\textwidth}{X} 你們要祈求，好讓你們逃走的時候，不遇見冬天或安息日。 \end{tabularx} \\ \\ \relax
24:21 & \begin{tabularx}{0.7\textwidth}{X} 因為那時必有大災難，自從世界的起頭直到如今，從沒有這樣的災難，將來也不會有。 \end{tabularx} \\ \\ \relax
24:22 & \begin{tabularx}{0.7\textwidth}{X} 若不減少那些日子，凡血肉之軀的，就沒有一個能得救；可是為了選民，那些日子將減少。 \end{tabularx} \\ \\ \relax
24:23 & \begin{tabularx}{0.7\textwidth}{X} 那時，若有人對你們說：『看哪，基督在這裡！』或『在那裡！』你們不要信。 \end{tabularx} \\ \\ \relax
24:24 & \begin{tabularx}{0.7\textwidth}{X} 因為假基督和假先知將要起來，顯大神蹟、大奇事，如果可能，要把選民也迷惑了。 \end{tabularx} \\ \\ \relax
24:25 & \begin{tabularx}{0.7\textwidth}{X} 看哪，我已經預先告訴你們了。 \end{tabularx} \\ \\ \relax
24:26 & \begin{tabularx}{0.7\textwidth}{X} 若有人對你們說：『看哪，基督在曠野裡！』你們不要出去；或說：『看哪，基督在內室中！』你們不要信。 \end{tabularx} \\ \\ \relax
24:27 & \begin{tabularx}{0.7\textwidth}{X} 好像閃電從東邊發出，直照到西邊，人子來臨也要這樣。 \end{tabularx} \\ \\ \relax
24:28 & \begin{tabularx}{0.7\textwidth}{X} 屍首在哪裡，鷹也會聚在哪裡。」 \end{tabularx} \\ \\ \relax
24:29 & \begin{tabularx}{0.7\textwidth}{X} 「那些日子的災難一過去，太陽要變黑，月亮也不放光，眾星要從天上墜落，天上的萬象都要震動。 \end{tabularx} \\ \\ \relax
24:30 & \begin{tabularx}{0.7\textwidth}{X} 那時，人子的預兆要顯在天上，地上的萬族都要哀哭。他們要看見人子帶著能力和大榮耀，駕著天上的雲來臨。 \end{tabularx} \\ \\ \relax
24:31 & \begin{tabularx}{0.7\textwidth}{X} 他要差遣天使，用大聲的號筒，從四方，從天這邊直到天那邊，召集他的選民。」 \end{tabularx} \\ \\ \relax
24:32 & \begin{tabularx}{0.7\textwidth}{X} 「你們要從無花果樹學習功課：當樹枝發芽長葉的時候，你們就知道夏天近了。 \end{tabularx} \\ \\ \relax
24:33 & \begin{tabularx}{0.7\textwidth}{X} 同樣，當你們看見這一切，就知道那時候近了，就在門口了。 \end{tabularx} \\ \\ \relax
24:34 & \begin{tabularx}{0.7\textwidth}{X} 我實在告訴你們，這世代還沒有過去，這一切都要發生。 \end{tabularx} \\ \\ \relax
24:35 & \begin{tabularx}{0.7\textwidth}{X} 天地要廢去，我的話卻絕不廢去。」 \end{tabularx} \\ \\ \relax
24:36 & \begin{tabularx}{0.7\textwidth}{X} 「但那日子，那時辰，沒有人知道，連天上的天使也不知道，子也不知道，惟有父知道。 \end{tabularx} \\ \\ \relax
24:37 & \begin{tabularx}{0.7\textwidth}{X} 挪亞的日子怎樣，人子來臨也要怎樣。 \end{tabularx} \\ \\ \relax
24:38 & \begin{tabularx}{0.7\textwidth}{X} 在洪水以前的那些日子，人照常吃喝嫁娶，直到挪亞進方舟的那日， \end{tabularx} \\ \\ \relax
24:39 & \begin{tabularx}{0.7\textwidth}{X} 不知不覺洪水來了，把他們全都沖去。人子來臨也要這樣。 \end{tabularx} \\ \\ \relax
24:40 & \begin{tabularx}{0.7\textwidth}{X} 那時，兩個人在田裡，一個被接去，一個被撇下。 \end{tabularx} \\ \\ \relax
24:41 & \begin{tabularx}{0.7\textwidth}{X} 兩個女人推磨，一個被接去，一個被撇下。 \end{tabularx} \\ \\ \relax
24:42 & \begin{tabularx}{0.7\textwidth}{X} 所以，你們要警醒，因為不知道你們的主哪一天來到。 \end{tabularx} \\ \\ \relax
24:43 & \begin{tabularx}{0.7\textwidth}{X} 你們要知道，一家的主人若知道晚上甚麼時候有賊來，就必警醒，不讓賊挖穿房屋。 \end{tabularx} \\ \\ \relax
24:44 & \begin{tabularx}{0.7\textwidth}{X} 所以，你們也要預備，因為在你們想不到的時候，人子就來了。」 \end{tabularx} \\ \\ \relax
24:45 & \begin{tabularx}{0.7\textwidth}{X} 「那麼，誰是那忠心又精明的僕人，主人派他管理自己的家僕、按時分糧給他們的呢？ \end{tabularx} \\ \\ \relax
24:46 & \begin{tabularx}{0.7\textwidth}{X} 主人來到，看見僕人這樣做，那僕人就有福了。 \end{tabularx} \\ \\ \relax
24:47 & \begin{tabularx}{0.7\textwidth}{X} 我實在告訴你們，主人要派他管理所有的財產。 \end{tabularx} \\ \\ \relax
24:48 & \begin{tabularx}{0.7\textwidth}{X} 如果那惡僕心裡說：『我的主人會來得遲』， \end{tabularx} \\ \\ \relax
24:49 & \begin{tabularx}{0.7\textwidth}{X} 就動手打他的同伴，又和醉酒的人一同吃喝， \end{tabularx} \\ \\ \relax
24:50 & \begin{tabularx}{0.7\textwidth}{X} 在想不到的日子，不知道的時候，那僕人的主人要來， \end{tabularx} \\ \\ \relax
24:51 & \begin{tabularx}{0.7\textwidth}{X} 重重地懲罰他，定他和假冒為善的人同罪，在那裡他要哀哭切齒了。」 \end{tabularx} \\ \\
[1ex]
\hline
\hline
\end{longtable}


\newpage

\section{第72屆~港九培靈會~張慕皚~第三講}
\label{sec:645}
\textbf{根基毀壞,義人還能作甚麼?}
\newline
\newline
連結: \href{https://www.hkbibleconference.org/session-message/view/645}{\texttt{https://www.hkbibleconference.org/session-message/view/645}}
\newline
\newline
\hyperref[sec:644]{< < < PREV SERMON < < <}
~
\hyperref[sec:toc]{[返目錄]}
~
\hyperref[sec:646]{> > > NEXT SERMON > > >}
\newline
\newline
但以理書 1:1-1:21
\newline
\begin{longtable}{cl}
\hline
\hline
章節 & 經文 (和合本修訂版)\\
\hline
1:1 & \begin{tabularx}{0.7\textwidth}{X} 猶大王約雅敬在位第三年，巴比倫王尼布甲尼撒來到耶路撒冷，將城圍困。 \end{tabularx} \\ \\ \relax
1:2 & \begin{tabularx}{0.7\textwidth}{X} 主將猶大王約雅敬和神殿中的一些器皿交在他的手中。他就把他們帶到示拿地他神明的廟裡，將器皿收入他神明的庫房中。 \end{tabularx} \\ \\ \relax
1:3 & \begin{tabularx}{0.7\textwidth}{X} 王吩咐太監長亞施毗拿，從以色列人的王室後裔和貴族中帶進幾個人來， \end{tabularx} \\ \\ \relax
1:4 & \begin{tabularx}{0.7\textwidth}{X} 就是沒有殘疾、相貌俊美、通達各樣學問、知識聰明俱備、足能在王宮侍立的少年，要教他們迦勒底的文字和語言。 \end{tabularx} \\ \\ \relax
1:5 & \begin{tabularx}{0.7\textwidth}{X} 王從自己所用的膳和所飲的酒中，派給他們每日的分量，養育他們三年，好叫他們期滿以後侍立在王面前。 \end{tabularx} \\ \\ \relax
1:6 & \begin{tabularx}{0.7\textwidth}{X} 他們中間有猶大人但以理、哈拿尼雅、米沙利和亞撒利雅。 \end{tabularx} \\ \\ \relax
1:7 & \begin{tabularx}{0.7\textwidth}{X} 太監長給他們另外起名，稱但以理為伯提沙撒，稱哈拿尼雅為沙得拉，稱米沙利為米煞，稱亞撒利雅為亞伯尼。 \end{tabularx} \\ \\ \relax
1:8 & \begin{tabularx}{0.7\textwidth}{X} 但以理卻立志，不以王的膳和王所飲的酒玷污自己，於是懇求太監長容他不使自己玷污。 \end{tabularx} \\ \\ \relax
1:9 & \begin{tabularx}{0.7\textwidth}{X} 神使但以理在太監長眼前蒙恩，得憐憫。 \end{tabularx} \\ \\ \relax
1:10 & \begin{tabularx}{0.7\textwidth}{X} 太監長對但以理說：「我懼怕我主我王，他已經派給你們的飲食，何必讓他見你們的面貌比你們同年齡的少年憔悴呢？這樣，你們就使我的頭在王那裡不保了。」 \end{tabularx} \\ \\ \relax
1:11 & \begin{tabularx}{0.7\textwidth}{X} 但以理對太監長所派監管但以理、哈拿尼雅、米沙利、亞撒利雅的管理者說： \end{tabularx} \\ \\ \relax
1:12 & \begin{tabularx}{0.7\textwidth}{X} 「請你考驗僕人們十天，給我們素菜吃，清水喝， \end{tabularx} \\ \\ \relax
1:13 & \begin{tabularx}{0.7\textwidth}{X} 然後你親自觀察我們的面貌和那用王膳的少年的面貌；就照你所觀察的待你的僕人吧！」 \end{tabularx} \\ \\ \relax
1:14 & \begin{tabularx}{0.7\textwidth}{X} 管理者准許他們這件事，考驗他們十天。 \end{tabularx} \\ \\ \relax
1:15 & \begin{tabularx}{0.7\textwidth}{X} 過了十天，他們的身材看來比所有享用王膳的少年更加俊美健壯， \end{tabularx} \\ \\ \relax
1:16 & \begin{tabularx}{0.7\textwidth}{X} 於是管理者撤去王派給他們用的膳和所飲的酒，只給他們素菜。 \end{tabularx} \\ \\ \relax
1:17 & \begin{tabularx}{0.7\textwidth}{X} 這四個少年，神在各樣文字學問上賜給他們知識和聰明；但以理又明白各樣異象和夢兆。 \end{tabularx} \\ \\ \relax
1:18 & \begin{tabularx}{0.7\textwidth}{X} 王吩咐帶他們進宮的日子到了，太監長就把他們帶到尼布甲尼撒面前。 \end{tabularx} \\ \\ \relax
1:19 & \begin{tabularx}{0.7\textwidth}{X} 王與他們談論，在所有少年中找不到人能與但以理、哈拿尼雅、米沙利、亞撒利雅相比，於是他們就在王面前侍立。 \end{tabularx} \\ \\ \relax
1:20 & \begin{tabularx}{0.7\textwidth}{X} 王考問他們一切智慧和聰明的事，發現他們比全國所有的術士和巫師勝過十倍。 \end{tabularx} \\ \\ \relax
1:21 & \begin{tabularx}{0.7\textwidth}{X} 到居魯士王元年，但以理還健在。 \end{tabularx} \\ \\
[1ex]
\hline
\hline
\end{longtable}


\newpage

\section{第72屆~港九培靈會~張慕皚~第四講}
\label{sec:646}
\textbf{敵基督與全球化的災難}
\newline
\newline
連結: \href{https://www.hkbibleconference.org/session-message/view/646}{\texttt{https://www.hkbibleconference.org/session-message/view/646}}
\newline
\newline
\hyperref[sec:645]{< < < PREV SERMON < < <}
~
\hyperref[sec:toc]{[返目錄]}
~
\hyperref[sec:648]{> > > NEXT SERMON > > >}
\newline
\newline
馬太福音 24:1-24:51
\newline
\begin{longtable}{cl}
\hline
\hline
章節 & 經文 (和合本修訂版)\\
\hline
24:1 & \begin{tabularx}{0.7\textwidth}{X} 耶穌出了聖殿，正離開的時候，門徒前來，把聖殿的建築指給他看。 \end{tabularx} \\ \\ \relax
24:2 & \begin{tabularx}{0.7\textwidth}{X} 耶穌回應他們說：「你們不是看見這一切嗎？我實在告訴你們，這裡將沒有一塊石頭會留在另一塊石頭上，而不被拆毀的。」 \end{tabularx} \\ \\ \relax
24:3 & \begin{tabularx}{0.7\textwidth}{X} 耶穌在橄欖山上坐著，門徒私下進前來問他：「請告訴我們，甚麼時候有這些事呢？你來臨和世代的終結有甚麼預兆呢？」 \end{tabularx} \\ \\ \relax
24:4 & \begin{tabularx}{0.7\textwidth}{X} 耶穌回答他們：「你們要謹慎，免得有人迷惑你們。 \end{tabularx} \\ \\ \relax
24:5 & \begin{tabularx}{0.7\textwidth}{X} 因為將有好些人冒我的名來，說『我是基督』，並且要迷惑許多人。 \end{tabularx} \\ \\ \relax
24:6 & \begin{tabularx}{0.7\textwidth}{X} 你們也將聽見打仗和打仗的風聲。注意，不要驚慌！因為這些事必須發生，但這還不是終結。 \end{tabularx} \\ \\ \relax
24:7 & \begin{tabularx}{0.7\textwidth}{X} 民要攻打民，國要攻打國，多處必有饑荒、地震。 \end{tabularx} \\ \\ \relax
24:8 & \begin{tabularx}{0.7\textwidth}{X} 這都是災難的起頭。 \end{tabularx} \\ \\ \relax
24:9 & \begin{tabularx}{0.7\textwidth}{X} 那時，人要使你們陷在患難裡，也要殺害你們；你們又要為我的名被萬民憎恨。 \end{tabularx} \\ \\ \relax
24:10 & \begin{tabularx}{0.7\textwidth}{X} 那時，會有許多人跌倒，也會彼此陷害，彼此憎恨； \end{tabularx} \\ \\ \relax
24:11 & \begin{tabularx}{0.7\textwidth}{X} 且有好些假先知起來，迷惑許多人。 \end{tabularx} \\ \\ \relax
24:12 & \begin{tabularx}{0.7\textwidth}{X} 因為不法的事增多，許多人的愛心漸漸冷淡了。 \end{tabularx} \\ \\ \relax
24:13 & \begin{tabularx}{0.7\textwidth}{X} 但堅忍到底的終必得救。 \end{tabularx} \\ \\ \relax
24:14 & \begin{tabularx}{0.7\textwidth}{X} 這天國的福音要傳遍天下，對萬民作見證，然後終結才來到。」 \end{tabularx} \\ \\ \relax
24:15 & \begin{tabularx}{0.7\textwidth}{X} 「當你們看見先知但以理所說的那『施行毀滅的褻瀆者』站在聖地（讀這經的人要會意）， \end{tabularx} \\ \\ \relax
24:16 & \begin{tabularx}{0.7\textwidth}{X} 那時，在猶太的，應當逃到山上； \end{tabularx} \\ \\ \relax
24:17 & \begin{tabularx}{0.7\textwidth}{X} 在屋頂上的，不要下來拿家裡的東西； \end{tabularx} \\ \\ \relax
24:18 & \begin{tabularx}{0.7\textwidth}{X} 在田裡的，不要回去取衣裳。 \end{tabularx} \\ \\ \relax
24:19 & \begin{tabularx}{0.7\textwidth}{X} 在那些日子，懷孕的和奶孩子的就苦了。 \end{tabularx} \\ \\ \relax
24:20 & \begin{tabularx}{0.7\textwidth}{X} 你們要祈求，好讓你們逃走的時候，不遇見冬天或安息日。 \end{tabularx} \\ \\ \relax
24:21 & \begin{tabularx}{0.7\textwidth}{X} 因為那時必有大災難，自從世界的起頭直到如今，從沒有這樣的災難，將來也不會有。 \end{tabularx} \\ \\ \relax
24:22 & \begin{tabularx}{0.7\textwidth}{X} 若不減少那些日子，凡血肉之軀的，就沒有一個能得救；可是為了選民，那些日子將減少。 \end{tabularx} \\ \\ \relax
24:23 & \begin{tabularx}{0.7\textwidth}{X} 那時，若有人對你們說：『看哪，基督在這裡！』或『在那裡！』你們不要信。 \end{tabularx} \\ \\ \relax
24:24 & \begin{tabularx}{0.7\textwidth}{X} 因為假基督和假先知將要起來，顯大神蹟、大奇事，如果可能，要把選民也迷惑了。 \end{tabularx} \\ \\ \relax
24:25 & \begin{tabularx}{0.7\textwidth}{X} 看哪，我已經預先告訴你們了。 \end{tabularx} \\ \\ \relax
24:26 & \begin{tabularx}{0.7\textwidth}{X} 若有人對你們說：『看哪，基督在曠野裡！』你們不要出去；或說：『看哪，基督在內室中！』你們不要信。 \end{tabularx} \\ \\ \relax
24:27 & \begin{tabularx}{0.7\textwidth}{X} 好像閃電從東邊發出，直照到西邊，人子來臨也要這樣。 \end{tabularx} \\ \\ \relax
24:28 & \begin{tabularx}{0.7\textwidth}{X} 屍首在哪裡，鷹也會聚在哪裡。」 \end{tabularx} \\ \\ \relax
24:29 & \begin{tabularx}{0.7\textwidth}{X} 「那些日子的災難一過去，太陽要變黑，月亮也不放光，眾星要從天上墜落，天上的萬象都要震動。 \end{tabularx} \\ \\ \relax
24:30 & \begin{tabularx}{0.7\textwidth}{X} 那時，人子的預兆要顯在天上，地上的萬族都要哀哭。他們要看見人子帶著能力和大榮耀，駕著天上的雲來臨。 \end{tabularx} \\ \\ \relax
24:31 & \begin{tabularx}{0.7\textwidth}{X} 他要差遣天使，用大聲的號筒，從四方，從天這邊直到天那邊，召集他的選民。」 \end{tabularx} \\ \\ \relax
24:32 & \begin{tabularx}{0.7\textwidth}{X} 「你們要從無花果樹學習功課：當樹枝發芽長葉的時候，你們就知道夏天近了。 \end{tabularx} \\ \\ \relax
24:33 & \begin{tabularx}{0.7\textwidth}{X} 同樣，當你們看見這一切，就知道那時候近了，就在門口了。 \end{tabularx} \\ \\ \relax
24:34 & \begin{tabularx}{0.7\textwidth}{X} 我實在告訴你們，這世代還沒有過去，這一切都要發生。 \end{tabularx} \\ \\ \relax
24:35 & \begin{tabularx}{0.7\textwidth}{X} 天地要廢去，我的話卻絕不廢去。」 \end{tabularx} \\ \\ \relax
24:36 & \begin{tabularx}{0.7\textwidth}{X} 「但那日子，那時辰，沒有人知道，連天上的天使也不知道，子也不知道，惟有父知道。 \end{tabularx} \\ \\ \relax
24:37 & \begin{tabularx}{0.7\textwidth}{X} 挪亞的日子怎樣，人子來臨也要怎樣。 \end{tabularx} \\ \\ \relax
24:38 & \begin{tabularx}{0.7\textwidth}{X} 在洪水以前的那些日子，人照常吃喝嫁娶，直到挪亞進方舟的那日， \end{tabularx} \\ \\ \relax
24:39 & \begin{tabularx}{0.7\textwidth}{X} 不知不覺洪水來了，把他們全都沖去。人子來臨也要這樣。 \end{tabularx} \\ \\ \relax
24:40 & \begin{tabularx}{0.7\textwidth}{X} 那時，兩個人在田裡，一個被接去，一個被撇下。 \end{tabularx} \\ \\ \relax
24:41 & \begin{tabularx}{0.7\textwidth}{X} 兩個女人推磨，一個被接去，一個被撇下。 \end{tabularx} \\ \\ \relax
24:42 & \begin{tabularx}{0.7\textwidth}{X} 所以，你們要警醒，因為不知道你們的主哪一天來到。 \end{tabularx} \\ \\ \relax
24:43 & \begin{tabularx}{0.7\textwidth}{X} 你們要知道，一家的主人若知道晚上甚麼時候有賊來，就必警醒，不讓賊挖穿房屋。 \end{tabularx} \\ \\ \relax
24:44 & \begin{tabularx}{0.7\textwidth}{X} 所以，你們也要預備，因為在你們想不到的時候，人子就來了。」 \end{tabularx} \\ \\ \relax
24:45 & \begin{tabularx}{0.7\textwidth}{X} 「那麼，誰是那忠心又精明的僕人，主人派他管理自己的家僕、按時分糧給他們的呢？ \end{tabularx} \\ \\ \relax
24:46 & \begin{tabularx}{0.7\textwidth}{X} 主人來到，看見僕人這樣做，那僕人就有福了。 \end{tabularx} \\ \\ \relax
24:47 & \begin{tabularx}{0.7\textwidth}{X} 我實在告訴你們，主人要派他管理所有的財產。 \end{tabularx} \\ \\ \relax
24:48 & \begin{tabularx}{0.7\textwidth}{X} 如果那惡僕心裡說：『我的主人會來得遲』， \end{tabularx} \\ \\ \relax
24:49 & \begin{tabularx}{0.7\textwidth}{X} 就動手打他的同伴，又和醉酒的人一同吃喝， \end{tabularx} \\ \\ \relax
24:50 & \begin{tabularx}{0.7\textwidth}{X} 在想不到的日子，不知道的時候，那僕人的主人要來， \end{tabularx} \\ \\ \relax
24:51 & \begin{tabularx}{0.7\textwidth}{X} 重重地懲罰他，定他和假冒為善的人同罪，在那裡他要哀哭切齒了。」 \end{tabularx} \\ \\
[1ex]
\hline
\hline
\end{longtable}


\newpage

\section{第72屆~港九培靈會~張慕皚~第五講}
\label{sec:648}
\textbf{咆哮覓食的末世惡魔}
\newline
\newline
連結: \href{https://www.hkbibleconference.org/session-message/view/648}{\texttt{https://www.hkbibleconference.org/session-message/view/648}}
\newline
\newline
\hyperref[sec:646]{< < < PREV SERMON < < <}
~
\hyperref[sec:toc]{[返目錄]}
~
\hyperref[sec:650]{> > > NEXT SERMON > > >}
\newline
\newline
以弗所書 6:1-6:24
\newline
\begin{longtable}{cl}
\hline
\hline
章節 & 經文 (和合本修訂版)\\
\hline
6:1 & \begin{tabularx}{0.7\textwidth}{X} 作兒女的，你們要在主裡聽從父母，這是理所當然的。 \end{tabularx} \\ \\ \relax
6:2 & \begin{tabularx}{0.7\textwidth}{X} 當孝敬父母，使你得福，在世長壽。這是第一條帶應許的誡命。 \end{tabularx} \\ \\ \relax
6:3 & \begin{tabularx}{0.7\textwidth}{X} 見上節 \end{tabularx} \\ \\ \relax
6:4 & \begin{tabularx}{0.7\textwidth}{X} 作父親的，你們不要激怒兒女，但要照著主的教導和勸戒養育他們。 \end{tabularx} \\ \\ \relax
6:5 & \begin{tabularx}{0.7\textwidth}{X} 作僕人的，你們要懼怕戰兢，用誠實的心聽從你們肉身的主人，好像聽從基督一般； \end{tabularx} \\ \\ \relax
6:6 & \begin{tabularx}{0.7\textwidth}{X} 不要只在人的眼前這樣做，像僅是討人的喜歡，而是作基督的僕人，從心裡遵行神的旨意， \end{tabularx} \\ \\ \relax
6:7 & \begin{tabularx}{0.7\textwidth}{X} 甘心服侍，好像服侍主，不像服侍人， \end{tabularx} \\ \\ \relax
6:8 & \begin{tabularx}{0.7\textwidth}{X} 因為知道每個人所做的善事，不論是為奴的或是自主的，都必按所做的從主得到賞賜。 \end{tabularx} \\ \\ \relax
6:9 & \begin{tabularx}{0.7\textwidth}{X} 作主人的，你們待僕人也是一樣，不要威嚇他們，因為知道他們和你們在天上同有一位主，他並不偏待人。 \end{tabularx} \\ \\ \relax
6:10 & \begin{tabularx}{0.7\textwidth}{X} 最後，你們要靠著主，依賴他的大能大力作剛強的人。 \end{tabularx} \\ \\ \relax
6:11 & \begin{tabularx}{0.7\textwidth}{X} 要穿戴神所賜的全副軍裝，好抵擋魔鬼的詭計。 \end{tabularx} \\ \\ \relax
6:12 & \begin{tabularx}{0.7\textwidth}{X} 因為我們的爭戰並不是對抗有血有肉的人，而是對抗那些執政的、掌權的、管轄這幽暗世界的，以及天空靈界的惡魔。 \end{tabularx} \\ \\ \relax
6:13 & \begin{tabularx}{0.7\textwidth}{X} 所以，要拿起神所賜的全副軍裝，好在邪惡的日子能抵擋仇敵，並且完成了一切後還能站立得住。 \end{tabularx} \\ \\ \relax
6:14 & \begin{tabularx}{0.7\textwidth}{X} 所以，要站穩了，用真理當作帶子束腰，用公義當作護心鏡遮胸， \end{tabularx} \\ \\ \relax
6:15 & \begin{tabularx}{0.7\textwidth}{X} 又用和平的福音當作預備走路的鞋穿在腳上。 \end{tabularx} \\ \\ \relax
6:16 & \begin{tabularx}{0.7\textwidth}{X} 此外，要拿信德當作盾牌，用來撲滅那惡者一切燒著的箭。 \end{tabularx} \\ \\ \relax
6:17 & \begin{tabularx}{0.7\textwidth}{X} 要戴上救恩的頭盔，拿著聖靈的寶劍—就是神的道。 \end{tabularx} \\ \\ \relax
6:18 & \begin{tabularx}{0.7\textwidth}{X} 要靠著聖靈，隨時多方禱告祈求，並要為此警醒不倦，為眾聖徒祈求。 \end{tabularx} \\ \\ \relax
6:19 & \begin{tabularx}{0.7\textwidth}{X} 也要為我祈求，讓我有口才，能放膽開口講明福音的奧祕， \end{tabularx} \\ \\ \relax
6:20 & \begin{tabularx}{0.7\textwidth}{X} 我為這福音的奧祕作了帶鐵鏈的使者，讓我能照著當盡的本分放膽宣講。 \end{tabularx} \\ \\ \relax
6:21 & \begin{tabularx}{0.7\textwidth}{X} 今有親愛、忠心服事主的弟兄推基古，為了你們也明白我的事情和我的景況，他會讓你們知道一切的事。 \end{tabularx} \\ \\ \relax
6:22 & \begin{tabularx}{0.7\textwidth}{X} 我特意打發他到你們那裡去，好讓你們知道我們的情況，又讓他安慰你們的心。 \end{tabularx} \\ \\ \relax
6:23 & \begin{tabularx}{0.7\textwidth}{X} 願平安、慈愛、信心從父神和主耶穌基督歸給弟兄們。 \end{tabularx} \\ \\ \relax
6:24 & \begin{tabularx}{0.7\textwidth}{X} 願所有恆心愛我們主耶穌基督的人都蒙恩惠。 \end{tabularx} \\ \\
[1ex]
\hline
\hline
\end{longtable}


\newpage

\section{第72屆~港九培靈會~張慕皚~第六講}
\label{sec:650}
\textbf{面對惡魔,勝算在握}
\newline
\newline
連結: \href{https://www.hkbibleconference.org/session-message/view/650}{\texttt{https://www.hkbibleconference.org/session-message/view/650}}
\newline
\newline
\hyperref[sec:648]{< < < PREV SERMON < < <}
~
\hyperref[sec:toc]{[返目錄]}
~
\hyperref[sec:653]{> > > NEXT SERMON > > >}
\newline
\newline
以弗所書 6:1-6:24
\newline
\begin{longtable}{cl}
\hline
\hline
章節 & 經文 (和合本修訂版)\\
\hline
6:1 & \begin{tabularx}{0.7\textwidth}{X} 作兒女的，你們要在主裡聽從父母，這是理所當然的。 \end{tabularx} \\ \\ \relax
6:2 & \begin{tabularx}{0.7\textwidth}{X} 當孝敬父母，使你得福，在世長壽。這是第一條帶應許的誡命。 \end{tabularx} \\ \\ \relax
6:3 & \begin{tabularx}{0.7\textwidth}{X} 見上節 \end{tabularx} \\ \\ \relax
6:4 & \begin{tabularx}{0.7\textwidth}{X} 作父親的，你們不要激怒兒女，但要照著主的教導和勸戒養育他們。 \end{tabularx} \\ \\ \relax
6:5 & \begin{tabularx}{0.7\textwidth}{X} 作僕人的，你們要懼怕戰兢，用誠實的心聽從你們肉身的主人，好像聽從基督一般； \end{tabularx} \\ \\ \relax
6:6 & \begin{tabularx}{0.7\textwidth}{X} 不要只在人的眼前這樣做，像僅是討人的喜歡，而是作基督的僕人，從心裡遵行神的旨意， \end{tabularx} \\ \\ \relax
6:7 & \begin{tabularx}{0.7\textwidth}{X} 甘心服侍，好像服侍主，不像服侍人， \end{tabularx} \\ \\ \relax
6:8 & \begin{tabularx}{0.7\textwidth}{X} 因為知道每個人所做的善事，不論是為奴的或是自主的，都必按所做的從主得到賞賜。 \end{tabularx} \\ \\ \relax
6:9 & \begin{tabularx}{0.7\textwidth}{X} 作主人的，你們待僕人也是一樣，不要威嚇他們，因為知道他們和你們在天上同有一位主，他並不偏待人。 \end{tabularx} \\ \\ \relax
6:10 & \begin{tabularx}{0.7\textwidth}{X} 最後，你們要靠著主，依賴他的大能大力作剛強的人。 \end{tabularx} \\ \\ \relax
6:11 & \begin{tabularx}{0.7\textwidth}{X} 要穿戴神所賜的全副軍裝，好抵擋魔鬼的詭計。 \end{tabularx} \\ \\ \relax
6:12 & \begin{tabularx}{0.7\textwidth}{X} 因為我們的爭戰並不是對抗有血有肉的人，而是對抗那些執政的、掌權的、管轄這幽暗世界的，以及天空靈界的惡魔。 \end{tabularx} \\ \\ \relax
6:13 & \begin{tabularx}{0.7\textwidth}{X} 所以，要拿起神所賜的全副軍裝，好在邪惡的日子能抵擋仇敵，並且完成了一切後還能站立得住。 \end{tabularx} \\ \\ \relax
6:14 & \begin{tabularx}{0.7\textwidth}{X} 所以，要站穩了，用真理當作帶子束腰，用公義當作護心鏡遮胸， \end{tabularx} \\ \\ \relax
6:15 & \begin{tabularx}{0.7\textwidth}{X} 又用和平的福音當作預備走路的鞋穿在腳上。 \end{tabularx} \\ \\ \relax
6:16 & \begin{tabularx}{0.7\textwidth}{X} 此外，要拿信德當作盾牌，用來撲滅那惡者一切燒著的箭。 \end{tabularx} \\ \\ \relax
6:17 & \begin{tabularx}{0.7\textwidth}{X} 要戴上救恩的頭盔，拿著聖靈的寶劍—就是神的道。 \end{tabularx} \\ \\ \relax
6:18 & \begin{tabularx}{0.7\textwidth}{X} 要靠著聖靈，隨時多方禱告祈求，並要為此警醒不倦，為眾聖徒祈求。 \end{tabularx} \\ \\ \relax
6:19 & \begin{tabularx}{0.7\textwidth}{X} 也要為我祈求，讓我有口才，能放膽開口講明福音的奧祕， \end{tabularx} \\ \\ \relax
6:20 & \begin{tabularx}{0.7\textwidth}{X} 我為這福音的奧祕作了帶鐵鏈的使者，讓我能照著當盡的本分放膽宣講。 \end{tabularx} \\ \\ \relax
6:21 & \begin{tabularx}{0.7\textwidth}{X} 今有親愛、忠心服事主的弟兄推基古，為了你們也明白我的事情和我的景況，他會讓你們知道一切的事。 \end{tabularx} \\ \\ \relax
6:22 & \begin{tabularx}{0.7\textwidth}{X} 我特意打發他到你們那裡去，好讓你們知道我們的情況，又讓他安慰你們的心。 \end{tabularx} \\ \\ \relax
6:23 & \begin{tabularx}{0.7\textwidth}{X} 願平安、慈愛、信心從父神和主耶穌基督歸給弟兄們。 \end{tabularx} \\ \\ \relax
6:24 & \begin{tabularx}{0.7\textwidth}{X} 願所有恆心愛我們主耶穌基督的人都蒙恩惠。 \end{tabularx} \\ \\
[1ex]
\hline
\hline
\end{longtable}


\newpage

\section{第72屆~港九培靈會~張慕皚~第七講}
\label{sec:653}
\textbf{攜械披甲,嚴陣以待}
\newline
\newline
連結: \href{https://www.hkbibleconference.org/session-message/view/653}{\texttt{https://www.hkbibleconference.org/session-message/view/653}}
\newline
\newline
\hyperref[sec:650]{< < < PREV SERMON < < <}
~
\hyperref[sec:toc]{[返目錄]}
~
\hyperref[sec:655]{> > > NEXT SERMON > > >}
\newline
\newline
以弗所書 6:1-6:24
\newline
\begin{longtable}{cl}
\hline
\hline
章節 & 經文 (和合本修訂版)\\
\hline
6:1 & \begin{tabularx}{0.7\textwidth}{X} 作兒女的，你們要在主裡聽從父母，這是理所當然的。 \end{tabularx} \\ \\ \relax
6:2 & \begin{tabularx}{0.7\textwidth}{X} 當孝敬父母，使你得福，在世長壽。這是第一條帶應許的誡命。 \end{tabularx} \\ \\ \relax
6:3 & \begin{tabularx}{0.7\textwidth}{X} 見上節 \end{tabularx} \\ \\ \relax
6:4 & \begin{tabularx}{0.7\textwidth}{X} 作父親的，你們不要激怒兒女，但要照著主的教導和勸戒養育他們。 \end{tabularx} \\ \\ \relax
6:5 & \begin{tabularx}{0.7\textwidth}{X} 作僕人的，你們要懼怕戰兢，用誠實的心聽從你們肉身的主人，好像聽從基督一般； \end{tabularx} \\ \\ \relax
6:6 & \begin{tabularx}{0.7\textwidth}{X} 不要只在人的眼前這樣做，像僅是討人的喜歡，而是作基督的僕人，從心裡遵行神的旨意， \end{tabularx} \\ \\ \relax
6:7 & \begin{tabularx}{0.7\textwidth}{X} 甘心服侍，好像服侍主，不像服侍人， \end{tabularx} \\ \\ \relax
6:8 & \begin{tabularx}{0.7\textwidth}{X} 因為知道每個人所做的善事，不論是為奴的或是自主的，都必按所做的從主得到賞賜。 \end{tabularx} \\ \\ \relax
6:9 & \begin{tabularx}{0.7\textwidth}{X} 作主人的，你們待僕人也是一樣，不要威嚇他們，因為知道他們和你們在天上同有一位主，他並不偏待人。 \end{tabularx} \\ \\ \relax
6:10 & \begin{tabularx}{0.7\textwidth}{X} 最後，你們要靠著主，依賴他的大能大力作剛強的人。 \end{tabularx} \\ \\ \relax
6:11 & \begin{tabularx}{0.7\textwidth}{X} 要穿戴神所賜的全副軍裝，好抵擋魔鬼的詭計。 \end{tabularx} \\ \\ \relax
6:12 & \begin{tabularx}{0.7\textwidth}{X} 因為我們的爭戰並不是對抗有血有肉的人，而是對抗那些執政的、掌權的、管轄這幽暗世界的，以及天空靈界的惡魔。 \end{tabularx} \\ \\ \relax
6:13 & \begin{tabularx}{0.7\textwidth}{X} 所以，要拿起神所賜的全副軍裝，好在邪惡的日子能抵擋仇敵，並且完成了一切後還能站立得住。 \end{tabularx} \\ \\ \relax
6:14 & \begin{tabularx}{0.7\textwidth}{X} 所以，要站穩了，用真理當作帶子束腰，用公義當作護心鏡遮胸， \end{tabularx} \\ \\ \relax
6:15 & \begin{tabularx}{0.7\textwidth}{X} 又用和平的福音當作預備走路的鞋穿在腳上。 \end{tabularx} \\ \\ \relax
6:16 & \begin{tabularx}{0.7\textwidth}{X} 此外，要拿信德當作盾牌，用來撲滅那惡者一切燒著的箭。 \end{tabularx} \\ \\ \relax
6:17 & \begin{tabularx}{0.7\textwidth}{X} 要戴上救恩的頭盔，拿著聖靈的寶劍—就是神的道。 \end{tabularx} \\ \\ \relax
6:18 & \begin{tabularx}{0.7\textwidth}{X} 要靠著聖靈，隨時多方禱告祈求，並要為此警醒不倦，為眾聖徒祈求。 \end{tabularx} \\ \\ \relax
6:19 & \begin{tabularx}{0.7\textwidth}{X} 也要為我祈求，讓我有口才，能放膽開口講明福音的奧祕， \end{tabularx} \\ \\ \relax
6:20 & \begin{tabularx}{0.7\textwidth}{X} 我為這福音的奧祕作了帶鐵鏈的使者，讓我能照著當盡的本分放膽宣講。 \end{tabularx} \\ \\ \relax
6:21 & \begin{tabularx}{0.7\textwidth}{X} 今有親愛、忠心服事主的弟兄推基古，為了你們也明白我的事情和我的景況，他會讓你們知道一切的事。 \end{tabularx} \\ \\ \relax
6:22 & \begin{tabularx}{0.7\textwidth}{X} 我特意打發他到你們那裡去，好讓你們知道我們的情況，又讓他安慰你們的心。 \end{tabularx} \\ \\ \relax
6:23 & \begin{tabularx}{0.7\textwidth}{X} 願平安、慈愛、信心從父神和主耶穌基督歸給弟兄們。 \end{tabularx} \\ \\ \relax
6:24 & \begin{tabularx}{0.7\textwidth}{X} 願所有恆心愛我們主耶穌基督的人都蒙恩惠。 \end{tabularx} \\ \\
[1ex]
\hline
\hline
\end{longtable}


\newpage

\section{第72屆~港九培靈會~張慕皚~第八講}
\label{sec:655}
\textbf{深夜中的守望者}
\newline
\newline
連結: \href{https://www.hkbibleconference.org/session-message/view/655}{\texttt{https://www.hkbibleconference.org/session-message/view/655}}
\newline
\newline
\hyperref[sec:653]{< < < PREV SERMON < < <}
~
\hyperref[sec:toc]{[返目錄]}
~
\hyperref[sec:658]{> > > NEXT SERMON > > >}
\newline
\newline
馬太福音 25:1-25:46
\newline
\begin{longtable}{cl}
\hline
\hline
章節 & 經文 (和合本修訂版)\\
\hline
25:1 & \begin{tabularx}{0.7\textwidth}{X} 「那時，天國好比十個童女拿著燈出去迎接新郎。 \end{tabularx} \\ \\ \relax
25:2 & \begin{tabularx}{0.7\textwidth}{X} 其中有五個是愚拙的，五個是聰明的。 \end{tabularx} \\ \\ \relax
25:3 & \begin{tabularx}{0.7\textwidth}{X} 愚拙的拿著燈，卻沒有帶油； \end{tabularx} \\ \\ \relax
25:4 & \begin{tabularx}{0.7\textwidth}{X} 聰明的拿著燈，又盛了油在器皿裡。 \end{tabularx} \\ \\ \relax
25:5 & \begin{tabularx}{0.7\textwidth}{X} 新郎遲延的時候，她們都打盹，睡著了。 \end{tabularx} \\ \\ \relax
25:6 & \begin{tabularx}{0.7\textwidth}{X} 半夜有人喊：『看，新郎來了，你們出來迎接他。』 \end{tabularx} \\ \\ \relax
25:7 & \begin{tabularx}{0.7\textwidth}{X} 那些童女就都起來挑亮她們的燈。 \end{tabularx} \\ \\ \relax
25:8 & \begin{tabularx}{0.7\textwidth}{X} 愚拙的對聰明的說：『請分點油給我們，因為我們的燈要滅了。』 \end{tabularx} \\ \\ \relax
25:9 & \begin{tabularx}{0.7\textwidth}{X} 聰明的回答：『恐怕不夠你我用的；你們還是自己到賣油的那裡去買吧。』 \end{tabularx} \\ \\ \relax
25:10 & \begin{tabularx}{0.7\textwidth}{X} 她們去買的時候，新郎到了。那預備好了的，與他進去共赴婚宴，門就關了。 \end{tabularx} \\ \\ \relax
25:11 & \begin{tabularx}{0.7\textwidth}{X} 其餘的童女隨後也來了，說：『主啊，主啊，給我們開門！』 \end{tabularx} \\ \\ \relax
25:12 & \begin{tabularx}{0.7\textwidth}{X} 他卻回答：『我實在告訴你們，我不認識你們。』 \end{tabularx} \\ \\ \relax
25:13 & \begin{tabularx}{0.7\textwidth}{X} 所以，你們要警醒，因為那日子，那時辰，你們不知道。」 \end{tabularx} \\ \\ \relax
25:14 & \begin{tabularx}{0.7\textwidth}{X} 「天國又好比一個人要出外遠行，就叫了僕人來，把他的家業交給他們。 \end{tabularx} \\ \\ \relax
25:15 & \begin{tabularx}{0.7\textwidth}{X} 他按著各人的才幹，給他們銀子：一個給了五千，一個給了二千，一個給了一千，就出外遠行去了。 \end{tabularx} \\ \\ \relax
25:16 & \begin{tabularx}{0.7\textwidth}{X} 那領五千的立刻拿去做買賣，另外賺了五千。 \end{tabularx} \\ \\ \relax
25:17 & \begin{tabularx}{0.7\textwidth}{X} 那領二千的也照樣另賺了二千。 \end{tabularx} \\ \\ \relax
25:18 & \begin{tabularx}{0.7\textwidth}{X} 但那領一千的去掘開地，把主人的銀子埋藏了。 \end{tabularx} \\ \\ \relax
25:19 & \begin{tabularx}{0.7\textwidth}{X} 過了許久，那些僕人的主人來了，和他們算賬。 \end{tabularx} \\ \\ \relax
25:20 & \begin{tabularx}{0.7\textwidth}{X} 那領五千的又帶著另外的五千來，說：『主啊，你交給我五千。請看，我又賺了五千。』 \end{tabularx} \\ \\ \relax
25:21 & \begin{tabularx}{0.7\textwidth}{X} 主人說：『好，你這又善良又忠心的僕人，你在少許的事上忠心，我要派你管理許多的事，進來享受你主人的快樂吧！』 \end{tabularx} \\ \\ \relax
25:22 & \begin{tabularx}{0.7\textwidth}{X} 那領二千的也進前來，說：『主啊，你交給我二千。請看，我又賺了二千。』 \end{tabularx} \\ \\ \relax
25:23 & \begin{tabularx}{0.7\textwidth}{X} 主人說：『好，你這又善良又忠心的僕人，你在少許的事上忠心，我要派你管理許多的事，進來享受你主人的快樂吧！』 \end{tabularx} \\ \\ \relax
25:24 & \begin{tabularx}{0.7\textwidth}{X} 那領一千的也進前來，說：『主啊，我知道你，你是個嚴厲的人：沒有種的地方也要收割，沒有播的地方也要收穫， \end{tabularx} \\ \\ \relax
25:25 & \begin{tabularx}{0.7\textwidth}{X} 我就害怕，去把你的一千銀子埋藏在地裡。請看，你的銀子在這裡。』 \end{tabularx} \\ \\ \relax
25:26 & \begin{tabularx}{0.7\textwidth}{X} 他的主人回答他說：『你這又惡又懶的僕人，你既知道我沒有種的地方也要收割，沒有播的地方也要收穫， \end{tabularx} \\ \\ \relax
25:27 & \begin{tabularx}{0.7\textwidth}{X} 就該把我的銀子放給兌換銀錢的人，到我來的時候可以連本帶利收回。 \end{tabularx} \\ \\ \relax
25:28 & \begin{tabularx}{0.7\textwidth}{X} 把他這一千奪過來，給那有一萬的。 \end{tabularx} \\ \\ \relax
25:29 & \begin{tabularx}{0.7\textwidth}{X} 因為凡有的，還要加給他，叫他有餘；沒有的，連他所有的也要奪過來。 \end{tabularx} \\ \\ \relax
25:30 & \begin{tabularx}{0.7\textwidth}{X} 把這無用的僕人丟在外面黑暗裡，在那裡他要哀哭切齒了。』」 \end{tabularx} \\ \\ \relax
25:31 & \begin{tabularx}{0.7\textwidth}{X} 「當人子在他榮耀裡，同著眾天使來臨的時候，要坐在他榮耀的寶座上。 \end{tabularx} \\ \\ \relax
25:32 & \begin{tabularx}{0.7\textwidth}{X} 萬民都要聚集在他面前。他要把他們分別出來，好像牧人分別綿羊、山羊一般， \end{tabularx} \\ \\ \relax
25:33 & \begin{tabularx}{0.7\textwidth}{X} 把綿羊安置在右邊，山羊在左邊。 \end{tabularx} \\ \\ \relax
25:34 & \begin{tabularx}{0.7\textwidth}{X} 於是王要向他右邊的說：『你們這蒙我父賜福的，可來承受那創世以來為你們所預備的國。 \end{tabularx} \\ \\ \relax
25:35 & \begin{tabularx}{0.7\textwidth}{X} 因為我餓了，你們給我吃；渴了，你們給我喝；我流浪在外，你們留我住； \end{tabularx} \\ \\ \relax
25:36 & \begin{tabularx}{0.7\textwidth}{X} 我赤身露體，你們給我穿；我病了，你們看顧我；我在監獄裡，你們來看我。』 \end{tabularx} \\ \\ \relax
25:37 & \begin{tabularx}{0.7\textwidth}{X} 義人就回答：『主啊，我們甚麼時候見你餓了，給你吃；渴了，給你喝？ \end{tabularx} \\ \\ \relax
25:38 & \begin{tabularx}{0.7\textwidth}{X} 甚麼時候見你流浪在外，留你住；或是赤身露體，給你穿？ \end{tabularx} \\ \\ \relax
25:39 & \begin{tabularx}{0.7\textwidth}{X} 又甚麼時候見你病了，或是在監獄裡，來看你呢？』 \end{tabularx} \\ \\ \relax
25:40 & \begin{tabularx}{0.7\textwidth}{X} 王回答他們說：『我實在告訴你們，這些事你們做在我弟兄中一個最小的身上，就是做在我身上了。』 \end{tabularx} \\ \\ \relax
25:41 & \begin{tabularx}{0.7\textwidth}{X} 「王又要向那左邊的說：『你們這被詛咒的人，離開我！進入那為魔鬼和他的使者所預備的永火裡去！ \end{tabularx} \\ \\ \relax
25:42 & \begin{tabularx}{0.7\textwidth}{X} 因為我餓了，你們沒有給我吃；渴了，你們沒有給我喝； \end{tabularx} \\ \\ \relax
25:43 & \begin{tabularx}{0.7\textwidth}{X} 我流浪在外，你們沒有留我住；我赤身露體，你們沒有給我穿；我病了，我在監獄裡，你們沒有來看顧我。』 \end{tabularx} \\ \\ \relax
25:44 & \begin{tabularx}{0.7\textwidth}{X} 他們也要回答：『主啊，我們甚麼時候見你餓了，或渴了，或流浪在外，或赤身露體，或病了，或在監獄裡，沒有伺候你呢？』 \end{tabularx} \\ \\ \relax
25:45 & \begin{tabularx}{0.7\textwidth}{X} 王要回答：『我實在告訴你們，這些事你們沒有做在任何一個最小的弟兄身上，就是沒有做在我身上了。』 \end{tabularx} \\ \\ \relax
25:46 & \begin{tabularx}{0.7\textwidth}{X} 這些人要往永刑裡去；那些義人要往永生裡去。」 \end{tabularx} \\ \\
[1ex]
\hline
\hline
\end{longtable}


\newpage

\section{第72屆~港九培靈會~張慕皚~第九講}
\label{sec:658}
\textbf{離開巴比倫大城的呼召}
\newline
\newline
連結: \href{https://www.hkbibleconference.org/session-message/view/658}{\texttt{https://www.hkbibleconference.org/session-message/view/658}}
\newline
\newline
\hyperref[sec:655]{< < < PREV SERMON < < <}
~
\hyperref[sec:toc]{[返目錄]}
~
\hyperref[sec:660]{> > > NEXT SERMON > > >}
\newline
\newline
以賽亞書 52:1-52:15
\newline
\begin{longtable}{cl}
\hline
\hline
章節 & 經文 (和合本修訂版)\\
\hline
52:1 & \begin{tabularx}{0.7\textwidth}{X} 錫安哪，興起！興起！穿上你的能力！聖城耶路撒冷啊，穿上你華美的衣服！因為從今以後，未受割禮、不潔淨的必不再進入你中間。 \end{tabularx} \\ \\ \relax
52:2 & \begin{tabularx}{0.7\textwidth}{X} 耶路撒冷啊，抖去塵埃，起來坐在王位上！被擄的錫安哪，解開你頸上的鎖鏈！ \end{tabularx} \\ \\ \relax
52:3 & \begin{tabularx}{0.7\textwidth}{X} 耶和華如此說：「你們白白地被賣，也必不用銀子贖回。」 \end{tabularx} \\ \\ \relax
52:4 & \begin{tabularx}{0.7\textwidth}{X} 主耶和華如此說：「先前我的百姓下到埃及，在那裡寄居，末後又有亞述人欺壓他們。 \end{tabularx} \\ \\ \relax
52:5 & \begin{tabularx}{0.7\textwidth}{X} 我的百姓既是白白地被擄，如今我在這裡做甚麼呢？這是耶和華說的。轄制他們的人歡呼，我的名終日不斷受褻瀆，這是耶和華說的。 \end{tabularx} \\ \\ \relax
52:6 & \begin{tabularx}{0.7\textwidth}{X} 因此，我的百姓必認識我的名；在那日，他們必知道說這話的就是我。看哪，是我！」 \end{tabularx} \\ \\ \relax
52:7 & \begin{tabularx}{0.7\textwidth}{X} 在山上報佳音，傳平安，報好信息，傳揚救恩，那人的腳蹤何等佳美啊！他對錫安說：「你的神作王了！」 \end{tabularx} \\ \\ \relax
52:8 & \begin{tabularx}{0.7\textwidth}{X} 聽啊，你守望之人的聲音，他們揚聲一同歡唱；因為他們必親眼看見耶和華返回錫安。 \end{tabularx} \\ \\ \relax
52:9 & \begin{tabularx}{0.7\textwidth}{X} 耶路撒冷的廢墟啊，要出聲一同歡唱；因為耶和華安慰了他的百姓，救贖了耶路撒冷。 \end{tabularx} \\ \\ \relax
52:10 & \begin{tabularx}{0.7\textwidth}{X} 耶和華在萬國眼前露出聖臂，地的四極都要看見我們神的救恩。 \end{tabularx} \\ \\ \relax
52:11 & \begin{tabularx}{0.7\textwidth}{X} 離開吧！離開吧！你們要從巴比倫出來。你們扛抬耶和華器皿的人哪，不要沾染不潔淨之物，離去時務要保持潔淨。 \end{tabularx} \\ \\ \relax
52:12 & \begin{tabularx}{0.7\textwidth}{X} 你們出來必不致匆忙，也不致奔逃；因為耶和華要在你們前頭行，以色列的神必作你們的後盾。 \end{tabularx} \\ \\ \relax
52:13 & \begin{tabularx}{0.7\textwidth}{X} 看哪，我的僕人行事必有智慧，他必被高升，高舉，升到至高之處。 \end{tabularx} \\ \\ \relax
52:14 & \begin{tabularx}{0.7\textwidth}{X} 許多人因他驚奇---他的面貌比別人憔悴，他的外表比世人枯槁--- \end{tabularx} \\ \\ \relax
52:15 & \begin{tabularx}{0.7\textwidth}{X} 同樣，他也必使許多國家驚奇，君王要向他閉口。未曾傳給他們的，他們必看見；未曾聽見過的事，他們要明白。 \end{tabularx} \\ \\
[1ex]
\hline
\hline
\end{longtable}


\newpage

\section{第72屆~港九培靈會~張慕皚~第十講}
\label{sec:660}
\textbf{主耶穌何時再來?}
\newline
\newline
連結: \href{https://www.hkbibleconference.org/session-message/view/660}{\texttt{https://www.hkbibleconference.org/session-message/view/660}}
\newline
\newline
\hyperref[sec:658]{< < < PREV SERMON < < <}
~
\hyperref[sec:toc]{[返目錄]}
~
\hyperref[sec:606]{> > > NEXT SERMON > > >}
\newline
\newline
馬太福音 24:1-24:51
\newline
\begin{longtable}{cl}
\hline
\hline
章節 & 經文 (和合本修訂版)\\
\hline
24:1 & \begin{tabularx}{0.7\textwidth}{X} 耶穌出了聖殿，正離開的時候，門徒前來，把聖殿的建築指給他看。 \end{tabularx} \\ \\ \relax
24:2 & \begin{tabularx}{0.7\textwidth}{X} 耶穌回應他們說：「你們不是看見這一切嗎？我實在告訴你們，這裡將沒有一塊石頭會留在另一塊石頭上，而不被拆毀的。」 \end{tabularx} \\ \\ \relax
24:3 & \begin{tabularx}{0.7\textwidth}{X} 耶穌在橄欖山上坐著，門徒私下進前來問他：「請告訴我們，甚麼時候有這些事呢？你來臨和世代的終結有甚麼預兆呢？」 \end{tabularx} \\ \\ \relax
24:4 & \begin{tabularx}{0.7\textwidth}{X} 耶穌回答他們：「你們要謹慎，免得有人迷惑你們。 \end{tabularx} \\ \\ \relax
24:5 & \begin{tabularx}{0.7\textwidth}{X} 因為將有好些人冒我的名來，說『我是基督』，並且要迷惑許多人。 \end{tabularx} \\ \\ \relax
24:6 & \begin{tabularx}{0.7\textwidth}{X} 你們也將聽見打仗和打仗的風聲。注意，不要驚慌！因為這些事必須發生，但這還不是終結。 \end{tabularx} \\ \\ \relax
24:7 & \begin{tabularx}{0.7\textwidth}{X} 民要攻打民，國要攻打國，多處必有饑荒、地震。 \end{tabularx} \\ \\ \relax
24:8 & \begin{tabularx}{0.7\textwidth}{X} 這都是災難的起頭。 \end{tabularx} \\ \\ \relax
24:9 & \begin{tabularx}{0.7\textwidth}{X} 那時，人要使你們陷在患難裡，也要殺害你們；你們又要為我的名被萬民憎恨。 \end{tabularx} \\ \\ \relax
24:10 & \begin{tabularx}{0.7\textwidth}{X} 那時，會有許多人跌倒，也會彼此陷害，彼此憎恨； \end{tabularx} \\ \\ \relax
24:11 & \begin{tabularx}{0.7\textwidth}{X} 且有好些假先知起來，迷惑許多人。 \end{tabularx} \\ \\ \relax
24:12 & \begin{tabularx}{0.7\textwidth}{X} 因為不法的事增多，許多人的愛心漸漸冷淡了。 \end{tabularx} \\ \\ \relax
24:13 & \begin{tabularx}{0.7\textwidth}{X} 但堅忍到底的終必得救。 \end{tabularx} \\ \\ \relax
24:14 & \begin{tabularx}{0.7\textwidth}{X} 這天國的福音要傳遍天下，對萬民作見證，然後終結才來到。」 \end{tabularx} \\ \\ \relax
24:15 & \begin{tabularx}{0.7\textwidth}{X} 「當你們看見先知但以理所說的那『施行毀滅的褻瀆者』站在聖地（讀這經的人要會意）， \end{tabularx} \\ \\ \relax
24:16 & \begin{tabularx}{0.7\textwidth}{X} 那時，在猶太的，應當逃到山上； \end{tabularx} \\ \\ \relax
24:17 & \begin{tabularx}{0.7\textwidth}{X} 在屋頂上的，不要下來拿家裡的東西； \end{tabularx} \\ \\ \relax
24:18 & \begin{tabularx}{0.7\textwidth}{X} 在田裡的，不要回去取衣裳。 \end{tabularx} \\ \\ \relax
24:19 & \begin{tabularx}{0.7\textwidth}{X} 在那些日子，懷孕的和奶孩子的就苦了。 \end{tabularx} \\ \\ \relax
24:20 & \begin{tabularx}{0.7\textwidth}{X} 你們要祈求，好讓你們逃走的時候，不遇見冬天或安息日。 \end{tabularx} \\ \\ \relax
24:21 & \begin{tabularx}{0.7\textwidth}{X} 因為那時必有大災難，自從世界的起頭直到如今，從沒有這樣的災難，將來也不會有。 \end{tabularx} \\ \\ \relax
24:22 & \begin{tabularx}{0.7\textwidth}{X} 若不減少那些日子，凡血肉之軀的，就沒有一個能得救；可是為了選民，那些日子將減少。 \end{tabularx} \\ \\ \relax
24:23 & \begin{tabularx}{0.7\textwidth}{X} 那時，若有人對你們說：『看哪，基督在這裡！』或『在那裡！』你們不要信。 \end{tabularx} \\ \\ \relax
24:24 & \begin{tabularx}{0.7\textwidth}{X} 因為假基督和假先知將要起來，顯大神蹟、大奇事，如果可能，要把選民也迷惑了。 \end{tabularx} \\ \\ \relax
24:25 & \begin{tabularx}{0.7\textwidth}{X} 看哪，我已經預先告訴你們了。 \end{tabularx} \\ \\ \relax
24:26 & \begin{tabularx}{0.7\textwidth}{X} 若有人對你們說：『看哪，基督在曠野裡！』你們不要出去；或說：『看哪，基督在內室中！』你們不要信。 \end{tabularx} \\ \\ \relax
24:27 & \begin{tabularx}{0.7\textwidth}{X} 好像閃電從東邊發出，直照到西邊，人子來臨也要這樣。 \end{tabularx} \\ \\ \relax
24:28 & \begin{tabularx}{0.7\textwidth}{X} 屍首在哪裡，鷹也會聚在哪裡。」 \end{tabularx} \\ \\ \relax
24:29 & \begin{tabularx}{0.7\textwidth}{X} 「那些日子的災難一過去，太陽要變黑，月亮也不放光，眾星要從天上墜落，天上的萬象都要震動。 \end{tabularx} \\ \\ \relax
24:30 & \begin{tabularx}{0.7\textwidth}{X} 那時，人子的預兆要顯在天上，地上的萬族都要哀哭。他們要看見人子帶著能力和大榮耀，駕著天上的雲來臨。 \end{tabularx} \\ \\ \relax
24:31 & \begin{tabularx}{0.7\textwidth}{X} 他要差遣天使，用大聲的號筒，從四方，從天這邊直到天那邊，召集他的選民。」 \end{tabularx} \\ \\ \relax
24:32 & \begin{tabularx}{0.7\textwidth}{X} 「你們要從無花果樹學習功課：當樹枝發芽長葉的時候，你們就知道夏天近了。 \end{tabularx} \\ \\ \relax
24:33 & \begin{tabularx}{0.7\textwidth}{X} 同樣，當你們看見這一切，就知道那時候近了，就在門口了。 \end{tabularx} \\ \\ \relax
24:34 & \begin{tabularx}{0.7\textwidth}{X} 我實在告訴你們，這世代還沒有過去，這一切都要發生。 \end{tabularx} \\ \\ \relax
24:35 & \begin{tabularx}{0.7\textwidth}{X} 天地要廢去，我的話卻絕不廢去。」 \end{tabularx} \\ \\ \relax
24:36 & \begin{tabularx}{0.7\textwidth}{X} 「但那日子，那時辰，沒有人知道，連天上的天使也不知道，子也不知道，惟有父知道。 \end{tabularx} \\ \\ \relax
24:37 & \begin{tabularx}{0.7\textwidth}{X} 挪亞的日子怎樣，人子來臨也要怎樣。 \end{tabularx} \\ \\ \relax
24:38 & \begin{tabularx}{0.7\textwidth}{X} 在洪水以前的那些日子，人照常吃喝嫁娶，直到挪亞進方舟的那日， \end{tabularx} \\ \\ \relax
24:39 & \begin{tabularx}{0.7\textwidth}{X} 不知不覺洪水來了，把他們全都沖去。人子來臨也要這樣。 \end{tabularx} \\ \\ \relax
24:40 & \begin{tabularx}{0.7\textwidth}{X} 那時，兩個人在田裡，一個被接去，一個被撇下。 \end{tabularx} \\ \\ \relax
24:41 & \begin{tabularx}{0.7\textwidth}{X} 兩個女人推磨，一個被接去，一個被撇下。 \end{tabularx} \\ \\ \relax
24:42 & \begin{tabularx}{0.7\textwidth}{X} 所以，你們要警醒，因為不知道你們的主哪一天來到。 \end{tabularx} \\ \\ \relax
24:43 & \begin{tabularx}{0.7\textwidth}{X} 你們要知道，一家的主人若知道晚上甚麼時候有賊來，就必警醒，不讓賊挖穿房屋。 \end{tabularx} \\ \\ \relax
24:44 & \begin{tabularx}{0.7\textwidth}{X} 所以，你們也要預備，因為在你們想不到的時候，人子就來了。」 \end{tabularx} \\ \\ \relax
24:45 & \begin{tabularx}{0.7\textwidth}{X} 「那麼，誰是那忠心又精明的僕人，主人派他管理自己的家僕、按時分糧給他們的呢？ \end{tabularx} \\ \\ \relax
24:46 & \begin{tabularx}{0.7\textwidth}{X} 主人來到，看見僕人這樣做，那僕人就有福了。 \end{tabularx} \\ \\ \relax
24:47 & \begin{tabularx}{0.7\textwidth}{X} 我實在告訴你們，主人要派他管理所有的財產。 \end{tabularx} \\ \\ \relax
24:48 & \begin{tabularx}{0.7\textwidth}{X} 如果那惡僕心裡說：『我的主人會來得遲』， \end{tabularx} \\ \\ \relax
24:49 & \begin{tabularx}{0.7\textwidth}{X} 就動手打他的同伴，又和醉酒的人一同吃喝， \end{tabularx} \\ \\ \relax
24:50 & \begin{tabularx}{0.7\textwidth}{X} 在想不到的日子，不知道的時候，那僕人的主人要來， \end{tabularx} \\ \\ \relax
24:51 & \begin{tabularx}{0.7\textwidth}{X} 重重地懲罰他，定他和假冒為善的人同罪，在那裡他要哀哭切齒了。」 \end{tabularx} \\ \\
[1ex]
\hline
\hline
\end{longtable}


\newpage

\section{第72屆~港九培靈會~蘇穎睿~第一講}
\label{sec:606}
\textbf{新與變}
\newline
\newline
連結: \href{https://www.hkbibleconference.org/session-message/view/606}{\texttt{https://www.hkbibleconference.org/session-message/view/606}}
\newline
\newline
\hyperref[sec:660]{< < < PREV SERMON < < <}
~
\hyperref[sec:toc]{[返目錄]}
~
\hyperref[sec:607]{> > > NEXT SERMON > > >}
\newline
\newline


\newpage

\section{第72屆~港九培靈會~蘇穎睿~第二講}
\label{sec:607}
\textbf{見到就信?未必!}
\newline
\newline
連結: \href{https://www.hkbibleconference.org/session-message/view/607}{\texttt{https://www.hkbibleconference.org/session-message/view/607}}
\newline
\newline
\hyperref[sec:606]{< < < PREV SERMON < < <}
~
\hyperref[sec:toc]{[返目錄]}
~
\hyperref[sec:609]{> > > NEXT SERMON > > >}
\newline
\newline


\newpage

\section{第72屆~港九培靈會~蘇穎睿~第三講}
\label{sec:609}
\textbf{跳下去!等陣先!}
\newline
\newline
連結: \href{https://www.hkbibleconference.org/session-message/view/609}{\texttt{https://www.hkbibleconference.org/session-message/view/609}}
\newline
\newline
\hyperref[sec:607]{< < < PREV SERMON < < <}
~
\hyperref[sec:toc]{[返目錄]}
~
\hyperref[sec:612]{> > > NEXT SERMON > > >}
\newline
\newline


\newpage

\section{第72屆~港九培靈會~蘇穎睿~第四講}
\label{sec:612}
\textbf{封建枷鎖}
\newline
\newline
連結: \href{https://www.hkbibleconference.org/session-message/view/612}{\texttt{https://www.hkbibleconference.org/session-message/view/612}}
\newline
\newline
\hyperref[sec:609]{< < < PREV SERMON < < <}
~
\hyperref[sec:toc]{[返目錄]}
~
\hyperref[sec:613]{> > > NEXT SERMON > > >}
\newline
\newline


\newpage

\section{第72屆~港九培靈會~蘇穎睿~第五講}
\label{sec:613}
\textbf{民以食為先}
\newline
\newline
連結: \href{https://www.hkbibleconference.org/session-message/view/613}{\texttt{https://www.hkbibleconference.org/session-message/view/613}}
\newline
\newline
\hyperref[sec:612]{< < < PREV SERMON < < <}
~
\hyperref[sec:toc]{[返目錄]}
~
\hyperref[sec:615]{> > > NEXT SERMON > > >}
\newline
\newline


\newpage

\section{第72屆~港九培靈會~蘇穎睿~第六講}
\label{sec:615}
\textbf{是我,不要怕!}
\newline
\newline
連結: \href{https://www.hkbibleconference.org/session-message/view/615}{\texttt{https://www.hkbibleconference.org/session-message/view/615}}
\newline
\newline
\hyperref[sec:613]{< < < PREV SERMON < < <}
~
\hyperref[sec:toc]{[返目錄]}
~
\hyperref[sec:616]{> > > NEXT SERMON > > >}
\newline
\newline


\newpage

\section{第72屆~港九培靈會~蘇穎睿~第七講}
\label{sec:616}
\textbf{重見光明}
\newline
\newline
連結: \href{https://www.hkbibleconference.org/session-message/view/616}{\texttt{https://www.hkbibleconference.org/session-message/view/616}}
\newline
\newline
\hyperref[sec:615]{< < < PREV SERMON < < <}
~
\hyperref[sec:toc]{[返目錄]}
~
\hyperref[sec:617]{> > > NEXT SERMON > > >}
\newline
\newline


\newpage

\section{第72屆~港九培靈會~蘇穎睿~第八講}
\label{sec:617}
\textbf{等待難捱的日子}
\newline
\newline
連結: \href{https://www.hkbibleconference.org/session-message/view/617}{\texttt{https://www.hkbibleconference.org/session-message/view/617}}
\newline
\newline
\hyperref[sec:616]{< < < PREV SERMON < < <}
~
\hyperref[sec:toc]{[返目錄]}
~
\hyperref[sec:619]{> > > NEXT SERMON > > >}
\newline
\newline


\newpage

\section{第72屆~港九培靈會~蘇穎睿~第九講}
\label{sec:619}
\textbf{生離死別}
\newline
\newline
連結: \href{https://www.hkbibleconference.org/session-message/view/619}{\texttt{https://www.hkbibleconference.org/session-message/view/619}}
\newline
\newline
\hyperref[sec:617]{< < < PREV SERMON < < <}
~
\hyperref[sec:toc]{[返目錄]}
~
\hyperref[sec:622]{> > > NEXT SERMON > > >}
\newline
\newline


\newpage

\section{第72屆~港九培靈會~陳濟民~第一講}
\label{sec:622}
\textbf{耶穌的榮耀}
\newline
\newline
連結: \href{https://www.hkbibleconference.org/session-message/view/622}{\texttt{https://www.hkbibleconference.org/session-message/view/622}}
\newline
\newline
\hyperref[sec:619]{< < < PREV SERMON < < <}
~
\hyperref[sec:toc]{[返目錄]}
~
\hyperref[sec:625]{> > > NEXT SERMON > > >}
\newline
\newline


\newpage

\section{第72屆~港九培靈會~陳濟民~第二講}
\label{sec:625}
\textbf{生命的賜與者}
\newline
\newline
連結: \href{https://www.hkbibleconference.org/session-message/view/625}{\texttt{https://www.hkbibleconference.org/session-message/view/625}}
\newline
\newline
\hyperref[sec:622]{< < < PREV SERMON < < <}
~
\hyperref[sec:toc]{[返目錄]}
~
\hyperref[sec:627]{> > > NEXT SERMON > > >}
\newline
\newline


\newpage

\section{第72屆~港九培靈會~陳濟民~第三講}
\label{sec:627}
\textbf{從耶穌的「不孝」談起}
\newline
\newline
連結: \href{https://www.hkbibleconference.org/session-message/view/627}{\texttt{https://www.hkbibleconference.org/session-message/view/627}}
\newline
\newline
\hyperref[sec:625]{< < < PREV SERMON < < <}
~
\hyperref[sec:toc]{[返目錄]}
~
\hyperref[sec:629]{> > > NEXT SERMON > > >}
\newline
\newline


\newpage

\section{第72屆~港九培靈會~陳濟民~第四講}
\label{sec:629}
\textbf{恨死了}
\newline
\newline
連結: \href{https://www.hkbibleconference.org/session-message/view/629}{\texttt{https://www.hkbibleconference.org/session-message/view/629}}
\newline
\newline
\hyperref[sec:627]{< < < PREV SERMON < < <}
~
\hyperref[sec:toc]{[返目錄]}
~
\hyperref[sec:630]{> > > NEXT SERMON > > >}
\newline
\newline


\newpage

\section{第72屆~港九培靈會~陳濟民~第五講}
\label{sec:630}
\textbf{至死的愛}
\newline
\newline
連結: \href{https://www.hkbibleconference.org/session-message/view/630}{\texttt{https://www.hkbibleconference.org/session-message/view/630}}
\newline
\newline
\hyperref[sec:629]{< < < PREV SERMON < < <}
~
\hyperref[sec:toc]{[返目錄]}
~
\hyperref[sec:631]{> > > NEXT SERMON > > >}
\newline
\newline


\newpage

\section{第72屆~港九培靈會~陳濟民~第六講}
\label{sec:631}
\textbf{知與信}
\newline
\newline
連結: \href{https://www.hkbibleconference.org/session-message/view/631}{\texttt{https://www.hkbibleconference.org/session-message/view/631}}
\newline
\newline
\hyperref[sec:630]{< < < PREV SERMON < < <}
~
\hyperref[sec:toc]{[返目錄]}
~
\hyperref[sec:635]{> > > NEXT SERMON > > >}
\newline
\newline


\newpage

\section{第72屆~港九培靈會~陳濟民~第七講}
\label{sec:635}
\textbf{言與行}
\newline
\newline
連結: \href{https://www.hkbibleconference.org/session-message/view/635}{\texttt{https://www.hkbibleconference.org/session-message/view/635}}
\newline
\newline
\hyperref[sec:631]{< < < PREV SERMON < < <}
~
\hyperref[sec:toc]{[返目錄]}
~
\hyperref[sec:637]{> > > NEXT SERMON > > >}
\newline
\newline


\newpage

\section{第72屆~港九培靈會~陳濟民~第八講}
\label{sec:637}
\textbf{怎有可能超越}
\newline
\newline
連結: \href{https://www.hkbibleconference.org/session-message/view/637}{\texttt{https://www.hkbibleconference.org/session-message/view/637}}
\newline
\newline
\hyperref[sec:635]{< < < PREV SERMON < < <}
~
\hyperref[sec:toc]{[返目錄]}
~
\hyperref[sec:640]{> > > NEXT SERMON > > >}
\newline
\newline


\newpage

\section{第72屆~港九培靈會~陳濟民~第九講}
\label{sec:640}
\textbf{世人怎能信}
\newline
\newline
連結: \href{https://www.hkbibleconference.org/session-message/view/640}{\texttt{https://www.hkbibleconference.org/session-message/view/640}}
\newline
\newline
\hyperref[sec:637]{< < < PREV SERMON < < <}
~
\hyperref[sec:toc]{[返目錄]}
~
\hyperref[sec:587]{> > > NEXT SERMON > > >}
\newline
\newline


\newpage

\section{第73屆~港九培靈會~于力工~第一講}
\label{sec:587}
\textbf{由眾到一讀經法:用全部《聖經》來讀一段話}
\newline
\newline
連結: \href{https://www.hkbibleconference.org/session-message/view/587}{\texttt{https://www.hkbibleconference.org/session-message/view/587}}
\newline
\newline
\hyperref[sec:640]{< < < PREV SERMON < < <}
~
\hyperref[sec:toc]{[返目錄]}
~
\hyperref[sec:589]{> > > NEXT SERMON > > >}
\newline
\newline


\newpage

\section{第73屆~港九培靈會~于力工~第二講}
\label{sec:589}
\textbf{心理讀經法:用人的天性來讀主的話}
\newline
\newline
連結: \href{https://www.hkbibleconference.org/session-message/view/589}{\texttt{https://www.hkbibleconference.org/session-message/view/589}}
\newline
\newline
\hyperref[sec:587]{< < < PREV SERMON < < <}
~
\hyperref[sec:toc]{[返目錄]}
~
\hyperref[sec:592]{> > > NEXT SERMON > > >}
\newline
\newline


\newpage

\section{第73屆~港九培靈會~于力工~第三講}
\label{sec:592}
\textbf{性格讀經法:用人的性格來讀主的話}
\newline
\newline
連結: \href{https://www.hkbibleconference.org/session-message/view/592}{\texttt{https://www.hkbibleconference.org/session-message/view/592}}
\newline
\newline
\hyperref[sec:589]{< < < PREV SERMON < < <}
~
\hyperref[sec:toc]{[返目錄]}
~
\hyperref[sec:594]{> > > NEXT SERMON > > >}
\newline
\newline


\newpage

\section{第73屆~港九培靈會~于力工~第四講}
\label{sec:594}
\textbf{禱讀讀經法:用經文變成禱告的話語}
\newline
\newline
連結: \href{https://www.hkbibleconference.org/session-message/view/594}{\texttt{https://www.hkbibleconference.org/session-message/view/594}}
\newline
\newline
\hyperref[sec:592]{< < < PREV SERMON < < <}
~
\hyperref[sec:toc]{[返目錄]}
~
\hyperref[sec:597]{> > > NEXT SERMON > > >}
\newline
\newline


\newpage

\section{第73屆~港九培靈會~于力工~第五講}
\label{sec:597}
\textbf{創意讀經法:用已存之事物加上新的功用}
\newline
\newline
連結: \href{https://www.hkbibleconference.org/session-message/view/597}{\texttt{https://www.hkbibleconference.org/session-message/view/597}}
\newline
\newline
\hyperref[sec:594]{< < < PREV SERMON < < <}
~
\hyperref[sec:toc]{[返目錄]}
~
\hyperref[sec:598]{> > > NEXT SERMON > > >}
\newline
\newline


\newpage

\section{第73屆~港九培靈會~于力工~第六講}
\label{sec:598}
\textbf{由一到眾讀經法:用一處經文來看全部《聖經》的說明}
\newline
\newline
連結: \href{https://www.hkbibleconference.org/session-message/view/598}{\texttt{https://www.hkbibleconference.org/session-message/view/598}}
\newline
\newline
\hyperref[sec:597]{< < < PREV SERMON < < <}
~
\hyperref[sec:toc]{[返目錄]}
~
\hyperref[sec:599]{> > > NEXT SERMON > > >}
\newline
\newline


\newpage

\section{第73屆~港九培靈會~于力工~第七講}
\label{sec:599}
\textbf{靈修讀經法:用靈修的基本方法來讀主的話}
\newline
\newline
連結: \href{https://www.hkbibleconference.org/session-message/view/599}{\texttt{https://www.hkbibleconference.org/session-message/view/599}}
\newline
\newline
\hyperref[sec:598]{< < < PREV SERMON < < <}
~
\hyperref[sec:toc]{[返目錄]}
~
\hyperref[sec:600]{> > > NEXT SERMON > > >}
\newline
\newline


\newpage

\section{第73屆~港九培靈會~于力工~第八講}
\label{sec:600}
\textbf{邏輯讀經法:推理讀經,因果讀法}
\newline
\newline
連結: \href{https://www.hkbibleconference.org/session-message/view/600}{\texttt{https://www.hkbibleconference.org/session-message/view/600}}
\newline
\newline
\hyperref[sec:599]{< < < PREV SERMON < < <}
~
\hyperref[sec:toc]{[返目錄]}
~
\hyperref[sec:601]{> > > NEXT SERMON > > >}
\newline
\newline


\newpage

\section{第73屆~港九培靈會~于力工~第九講}
\label{sec:601}
\textbf{三一讀經法:讀出神的性情來}
\newline
\newline
連結: \href{https://www.hkbibleconference.org/session-message/view/601}{\texttt{https://www.hkbibleconference.org/session-message/view/601}}
\newline
\newline
\hyperref[sec:600]{< < < PREV SERMON < < <}
~
\hyperref[sec:toc]{[返目錄]}
~
\hyperref[sec:602]{> > > NEXT SERMON > > >}
\newline
\newline


\newpage

\section{第73屆~港九培靈會~于力工~第十講}
\label{sec:602}
\textbf{靈進讀經法:用靈命進深之系統來讀經}
\newline
\newline
連結: \href{https://www.hkbibleconference.org/session-message/view/602}{\texttt{https://www.hkbibleconference.org/session-message/view/602}}
\newline
\newline
\hyperref[sec:601]{< < < PREV SERMON < < <}
~
\hyperref[sec:toc]{[返目錄]}
~
\hyperref[sec:556]{> > > NEXT SERMON > > >}
\newline
\newline


\newpage

\section{第73屆~港九培靈會~李思敬~第一講}
\label{sec:556}
\textbf{我父在天}
\newline
\newline
連結: \href{https://www.hkbibleconference.org/session-message/view/556}{\texttt{https://www.hkbibleconference.org/session-message/view/556}}
\newline
\newline
\hyperref[sec:602]{< < < PREV SERMON < < <}
~
\hyperref[sec:toc]{[返目錄]}
~
\hyperref[sec:557]{> > > NEXT SERMON > > >}
\newline
\newline


\newpage

\section{第73屆~港九培靈會~李思敬~第二講}
\label{sec:557}
\textbf{願爾名聖,爾國臨格}
\newline
\newline
連結: \href{https://www.hkbibleconference.org/session-message/view/557}{\texttt{https://www.hkbibleconference.org/session-message/view/557}}
\newline
\newline
\hyperref[sec:556]{< < < PREV SERMON < < <}
~
\hyperref[sec:toc]{[返目錄]}
~
\hyperref[sec:560]{> > > NEXT SERMON > > >}
\newline
\newline


\newpage

\section{第73屆~港九培靈會~李思敬~第三講}
\label{sec:560}
\textbf{爾旨得成,在地若天}
\newline
\newline
連結: \href{https://www.hkbibleconference.org/session-message/view/560}{\texttt{https://www.hkbibleconference.org/session-message/view/560}}
\newline
\newline
\hyperref[sec:557]{< < < PREV SERMON < < <}
~
\hyperref[sec:toc]{[返目錄]}
~
\hyperref[sec:561]{> > > NEXT SERMON > > >}
\newline
\newline


\newpage

\section{第73屆~港九培靈會~李思敬~第四講}
\label{sec:561}
\textbf{所需之糧,今日賜我}
\newline
\newline
連結: \href{https://www.hkbibleconference.org/session-message/view/561}{\texttt{https://www.hkbibleconference.org/session-message/view/561}}
\newline
\newline
\hyperref[sec:560]{< < < PREV SERMON < < <}
~
\hyperref[sec:toc]{[返目錄]}
~
\hyperref[sec:562]{> > > NEXT SERMON > > >}
\newline
\newline


\newpage

\section{第73屆~港九培靈會~李思敬~第五講}
\label{sec:562}
\textbf{我免人負,求免我負}
\newline
\newline
連結: \href{https://www.hkbibleconference.org/session-message/view/562}{\texttt{https://www.hkbibleconference.org/session-message/view/562}}
\newline
\newline
\hyperref[sec:561]{< < < PREV SERMON < < <}
~
\hyperref[sec:toc]{[返目錄]}
~
\hyperref[sec:563]{> > > NEXT SERMON > > >}
\newline
\newline


\newpage

\section{第73屆~港九培靈會~李思敬~第六講}
\label{sec:563}
\textbf{俾勿我試,拯我出惡}
\newline
\newline
連結: \href{https://www.hkbibleconference.org/session-message/view/563}{\texttt{https://www.hkbibleconference.org/session-message/view/563}}
\newline
\newline
\hyperref[sec:562]{< < < PREV SERMON < < <}
~
\hyperref[sec:toc]{[返目錄]}
~
\hyperref[sec:564]{> > > NEXT SERMON > > >}
\newline
\newline


\newpage

\section{第73屆~港九培靈會~李思敬~第七講}
\label{sec:564}
\textbf{以國,權,榮,皆爾所有}
\newline
\newline
連結: \href{https://www.hkbibleconference.org/session-message/view/564}{\texttt{https://www.hkbibleconference.org/session-message/view/564}}
\newline
\newline
\hyperref[sec:563]{< < < PREV SERMON < < <}
~
\hyperref[sec:toc]{[返目錄]}
~
\hyperref[sec:566]{> > > NEXT SERMON > > >}
\newline
\newline


\newpage

\section{第73屆~港九培靈會~李思敬~第八講}
\label{sec:566}
\textbf{爰及世世,誠心所願}
\newline
\newline
連結: \href{https://www.hkbibleconference.org/session-message/view/566}{\texttt{https://www.hkbibleconference.org/session-message/view/566}}
\newline
\newline
\hyperref[sec:564]{< < < PREV SERMON < < <}
~
\hyperref[sec:toc]{[返目錄]}
~
\hyperref[sec:568]{> > > NEXT SERMON > > >}
\newline
\newline


\newpage

\section{第73屆~港九培靈會~李思敬~第九講}
\label{sec:568}
\textbf{凡求者得也}
\newline
\newline
連結: \href{https://www.hkbibleconference.org/session-message/view/568}{\texttt{https://www.hkbibleconference.org/session-message/view/568}}
\newline
\newline
\hyperref[sec:566]{< < < PREV SERMON < < <}
~
\hyperref[sec:toc]{[返目錄]}
~
\hyperref[sec:572]{> > > NEXT SERMON > > >}
\newline
\newline


\newpage

\section{第73屆~港九培靈會~樓恩德~第一講}
\label{sec:572}
\textbf{那人獨居不好}
\newline
\newline
連結: \href{https://www.hkbibleconference.org/session-message/view/572}{\texttt{https://www.hkbibleconference.org/session-message/view/572}}
\newline
\newline
\hyperref[sec:568]{< < < PREV SERMON < < <}
~
\hyperref[sec:toc]{[返目錄]}
~
\hyperref[sec:573]{> > > NEXT SERMON > > >}
\newline
\newline


\newpage

\section{第73屆~港九培靈會~樓恩德~第二講}
\label{sec:573}
\textbf{罪就伏在門前}
\newline
\newline
連結: \href{https://www.hkbibleconference.org/session-message/view/573}{\texttt{https://www.hkbibleconference.org/session-message/view/573}}
\newline
\newline
\hyperref[sec:572]{< < < PREV SERMON < < <}
~
\hyperref[sec:toc]{[返目錄]}
~
\hyperref[sec:575]{> > > NEXT SERMON > > >}
\newline
\newline


\newpage

\section{第73屆~港九培靈會~樓恩德~第三講}
\label{sec:575}
\textbf{漸漸挪移帳棚}
\newline
\newline
連結: \href{https://www.hkbibleconference.org/session-message/view/575}{\texttt{https://www.hkbibleconference.org/session-message/view/575}}
\newline
\newline
\hyperref[sec:573]{< < < PREV SERMON < < <}
~
\hyperref[sec:toc]{[返目錄]}
~
\hyperref[sec:577]{> > > NEXT SERMON > > >}
\newline
\newline


\newpage

\section{第73屆~港九培靈會~樓恩德~第四講}
\label{sec:577}
\textbf{從暗笑到喜笑}
\newline
\newline
連結: \href{https://www.hkbibleconference.org/session-message/view/577}{\texttt{https://www.hkbibleconference.org/session-message/view/577}}
\newline
\newline
\hyperref[sec:575]{< < < PREV SERMON < < <}
~
\hyperref[sec:toc]{[返目錄]}
~
\hyperref[sec:579]{> > > NEXT SERMON > > >}
\newline
\newline


\newpage

\section{第73屆~港九培靈會~樓恩德~第五講}
\label{sec:579}
\textbf{神必親自預備}
\newline
\newline
連結: \href{https://www.hkbibleconference.org/session-message/view/579}{\texttt{https://www.hkbibleconference.org/session-message/view/579}}
\newline
\newline
\hyperref[sec:577]{< < < PREV SERMON < < <}
~
\hyperref[sec:toc]{[返目錄]}
~
\hyperref[sec:580]{> > > NEXT SERMON > > >}
\newline
\newline


\newpage

\section{第73屆~港九培靈會~樓恩德~第六講}
\label{sec:580}
\textbf{必在這地昌盛}
\newline
\newline
連結: \href{https://www.hkbibleconference.org/session-message/view/580}{\texttt{https://www.hkbibleconference.org/session-message/view/580}}
\newline
\newline
\hyperref[sec:579]{< < < PREV SERMON < < <}
~
\hyperref[sec:toc]{[返目錄]}
~
\hyperref[sec:581]{> > > NEXT SERMON > > >}
\newline
\newline


\newpage

\section{第73屆~港九培靈會~樓恩德~第七講}
\label{sec:581}
\textbf{從不知到知道}
\newline
\newline
連結: \href{https://www.hkbibleconference.org/session-message/view/581}{\texttt{https://www.hkbibleconference.org/session-message/view/581}}
\newline
\newline
\hyperref[sec:580]{< < < PREV SERMON < < <}
~
\hyperref[sec:toc]{[返目錄]}
~
\hyperref[sec:582]{> > > NEXT SERMON > > >}
\newline
\newline


\newpage

\section{第73屆~港九培靈會~樓恩德~第八講}
\label{sec:582}
\textbf{不然我就死了}
\newline
\newline
連結: \href{https://www.hkbibleconference.org/session-message/view/582}{\texttt{https://www.hkbibleconference.org/session-message/view/582}}
\newline
\newline
\hyperref[sec:581]{< < < PREV SERMON < < <}
~
\hyperref[sec:toc]{[返目錄]}
~
\hyperref[sec:584]{> > > NEXT SERMON > > >}
\newline
\newline


\newpage

\section{第73屆~港九培靈會~樓恩德~第九講}
\label{sec:584}
\textbf{因信留下遺命}
\newline
\newline
連結: \href{https://www.hkbibleconference.org/session-message/view/584}{\texttt{https://www.hkbibleconference.org/session-message/view/584}}
\newline
\newline
\hyperref[sec:582]{< < < PREV SERMON < < <}
~
\hyperref[sec:toc]{[返目錄]}
~
\hyperref[sec:434]{> > > NEXT SERMON > > >}
\newline
\newline


\newpage

\section{第74屆~港九培靈會~戴紹曾~第一講}
\label{sec:434}
\textbf{蒙召成為見證}
\newline
\newline
連結: \href{https://www.hkbibleconference.org/session-message/view/434}{\texttt{https://www.hkbibleconference.org/session-message/view/434}}
\newline
\newline
\hyperref[sec:584]{< < < PREV SERMON < < <}
~
\hyperref[sec:toc]{[返目錄]}
~
\hyperref[sec:435]{> > > NEXT SERMON > > >}
\newline
\newline
以弗所書 4:1-4:1
\newline
\begin{longtable}{cl}
\hline
\hline
章節 & 經文 (和合本修訂版)\\
\hline
4:1 & \begin{tabularx}{0.7\textwidth}{X} 我為主作囚徒的勸你們，既然蒙召，行事為人就要與你們所蒙的呼召相稱。 \end{tabularx} \\ \\
[1ex]
\hline
\hline
\end{longtable}


\newpage

\section{第74屆~港九培靈會~戴紹曾~第二講}
\label{sec:435}
\textbf{蒙召承受福氣}
\newline
\newline
連結: \href{https://www.hkbibleconference.org/session-message/view/435}{\texttt{https://www.hkbibleconference.org/session-message/view/435}}
\newline
\newline
\hyperref[sec:434]{< < < PREV SERMON < < <}
~
\hyperref[sec:toc]{[返目錄]}
~
\hyperref[sec:439]{> > > NEXT SERMON > > >}
\newline
\newline
彼得前書 3:9-3:9
\newline
\begin{longtable}{cl}
\hline
\hline
章節 & 經文 (和合本修訂版)\\
\hline
3:9 & \begin{tabularx}{0.7\textwidth}{X} 不要以惡報惡，以辱罵還辱罵，倒要祝福，因為你們正是為此蒙召的，好使你們承受福氣。 \end{tabularx} \\ \\
[1ex]
\hline
\hline
\end{longtable}


\newpage

\section{第74屆~港九培靈會~戴紹曾~第三講}
\label{sec:439}
\textbf{蒙召得以自由}
\newline
\newline
連結: \href{https://www.hkbibleconference.org/session-message/view/439}{\texttt{https://www.hkbibleconference.org/session-message/view/439}}
\newline
\newline
\hyperref[sec:435]{< < < PREV SERMON < < <}
~
\hyperref[sec:toc]{[返目錄]}
~
\hyperref[sec:440]{> > > NEXT SERMON > > >}
\newline
\newline
路加福音 5:1-5:13
\newline
\begin{longtable}{cl}
\hline
\hline
章節 & 經文 (和合本修訂版)\\
\hline
5:1 & \begin{tabularx}{0.7\textwidth}{X} 耶穌站在革尼撒勒湖邊，眾人擁擠他，要聽神的道。 \end{tabularx} \\ \\ \relax
5:2 & \begin{tabularx}{0.7\textwidth}{X} 他見有兩隻船靠在湖邊，打魚的人卻離開船，洗網去了。 \end{tabularx} \\ \\ \relax
5:3 & \begin{tabularx}{0.7\textwidth}{X} 有一隻船是西門的，耶穌就上去，請他把船撐開，稍微離岸，就坐下，在船上教導眾人。 \end{tabularx} \\ \\ \relax
5:4 & \begin{tabularx}{0.7\textwidth}{X} 他講完了，對西門說：「把船開到水深的地方下網打魚。」 \end{tabularx} \\ \\ \relax
5:5 & \begin{tabularx}{0.7\textwidth}{X} 西門說：「老師，我們整夜勞累，並沒有打著甚麼。但依從你的話，我就下網。」 \end{tabularx} \\ \\ \relax
5:6 & \begin{tabularx}{0.7\textwidth}{X} 他們下了網，圈住許多魚，網險些裂開， \end{tabularx} \\ \\ \relax
5:7 & \begin{tabularx}{0.7\textwidth}{X} 就招手叫另一隻船上的同伴來幫助。他們就來，把魚裝滿了兩隻船，船甚至要沉下去。 \end{tabularx} \\ \\ \relax
5:8 & \begin{tabularx}{0.7\textwidth}{X} 西門‧彼得看見，就俯伏在耶穌膝前，說：「主啊，離開我，我是個罪人。」 \end{tabularx} \\ \\ \relax
5:9 & \begin{tabularx}{0.7\textwidth}{X} 他和一切跟他一起的人對打到了這一網的魚都很驚訝。 \end{tabularx} \\ \\ \relax
5:10 & \begin{tabularx}{0.7\textwidth}{X} 他的夥伴西庇太的兒子雅各、約翰，也是這樣。耶穌對西門說：「不要怕！從今以後，你要得人了。」 \end{tabularx} \\ \\ \relax
5:11 & \begin{tabularx}{0.7\textwidth}{X} 他們把兩隻船靠了岸，就撇下所有的，跟從了耶穌。 \end{tabularx} \\ \\ \relax
5:12 & \begin{tabularx}{0.7\textwidth}{X} 有一回，耶穌在一個城裡，有人滿身長了痲瘋，看見他，就俯伏在地，求他說：「主啊，你若肯，你能使我潔淨。」 \end{tabularx} \\ \\ \relax
5:13 & \begin{tabularx}{0.7\textwidth}{X} 耶穌伸手摸他，說：「我肯，你潔淨了吧！」痲瘋病立刻離開了他。 \end{tabularx} \\ \\
[1ex]
\hline
\hline
\end{longtable}


\newpage

\section{第74屆~港九培靈會~戴紹曾~第四講}
\label{sec:440}
\textbf{蒙召成為聖潔}
\newline
\newline
連結: \href{https://www.hkbibleconference.org/session-message/view/440}{\texttt{https://www.hkbibleconference.org/session-message/view/440}}
\newline
\newline
\hyperref[sec:439]{< < < PREV SERMON < < <}
~
\hyperref[sec:toc]{[返目錄]}
~
\hyperref[sec:441]{> > > NEXT SERMON > > >}
\newline
\newline
哥林多前書 1:15-1:16
\newline
\begin{longtable}{cl}
\hline
\hline
章節 & 經文 (和合本修訂版)\\
\hline
1:15 & \begin{tabularx}{0.7\textwidth}{X} 免得有人說你們是奉我的名受洗的。 \end{tabularx} \\ \\ \relax
1:16 & \begin{tabularx}{0.7\textwidth}{X} 我曾為司提法那家施過洗；此外我已記不清有沒有給別人施過洗。 \end{tabularx} \\ \\
[1ex]
\hline
\hline
\end{longtable}


\newpage

\section{第74屆~港九培靈會~戴紹曾~第五講}
\label{sec:441}
\textbf{蒙召與主聯合}
\newline
\newline
連結: \href{https://www.hkbibleconference.org/session-message/view/441}{\texttt{https://www.hkbibleconference.org/session-message/view/441}}
\newline
\newline
\hyperref[sec:440]{< < < PREV SERMON < < <}
~
\hyperref[sec:toc]{[返目錄]}
~
\hyperref[sec:442]{> > > NEXT SERMON > > >}
\newline
\newline
哥林多前書 1:9-1:9
\newline
\begin{longtable}{cl}
\hline
\hline
章節 & 經文 (和合本修訂版)\\
\hline
1:9 & \begin{tabularx}{0.7\textwidth}{X} 神是信實的，他呼召你們好與他兒子—我們的主耶穌基督—共享團契。 \end{tabularx} \\ \\
[1ex]
\hline
\hline
\end{longtable}


\newpage

\section{第74屆~港九培靈會~戴紹曾~第六講}
\label{sec:442}
\textbf{蒙召得享平安}
\newline
\newline
連結: \href{https://www.hkbibleconference.org/session-message/view/442}{\texttt{https://www.hkbibleconference.org/session-message/view/442}}
\newline
\newline
\hyperref[sec:441]{< < < PREV SERMON < < <}
~
\hyperref[sec:toc]{[返目錄]}
~
\hyperref[sec:443]{> > > NEXT SERMON > > >}
\newline
\newline
歌羅西書 3:15-3:15
\newline
\begin{longtable}{cl}
\hline
\hline
章節 & 經文 (和合本修訂版)\\
\hline
3:15 & \begin{tabularx}{0.7\textwidth}{X} 你們要讓基督所賜的和平在你們心裡作主，也為此蒙召，歸為一體。你們還要存感謝的心。 \end{tabularx} \\ \\
[1ex]
\hline
\hline
\end{longtable}


\newpage

\section{第74屆~港九培靈會~戴紹曾~第七講}
\label{sec:443}
\textbf{蒙召宣揚基督}
\newline
\newline
連結: \href{https://www.hkbibleconference.org/session-message/view/443}{\texttt{https://www.hkbibleconference.org/session-message/view/443}}
\newline
\newline
\hyperref[sec:442]{< < < PREV SERMON < < <}
~
\hyperref[sec:toc]{[返目錄]}
~
\hyperref[sec:444]{> > > NEXT SERMON > > >}
\newline
\newline
彼得前書 2:9-2:10
\newline
\begin{longtable}{cl}
\hline
\hline
章節 & 經文 (和合本修訂版)\\
\hline
2:9 & \begin{tabularx}{0.7\textwidth}{X} 不過，你們是被揀選的一族，是君尊的祭司，是神聖的國度，是屬神的子民，要使你們宣揚那召你們出黑暗入奇妙光明者的美德。 \end{tabularx} \\ \\ \relax
2:10 & \begin{tabularx}{0.7\textwidth}{X} 「你們從前不是子民，現在卻成了神的子民；從前未曾蒙憐憫，現在卻蒙了憐憫。」 \end{tabularx} \\ \\
[1ex]
\hline
\hline
\end{longtable}


\newpage

\section{第74屆~港九培靈會~戴紹曾~第八講}
\label{sec:444}
\textbf{蒙召為主受苦}
\newline
\newline
連結: \href{https://www.hkbibleconference.org/session-message/view/444}{\texttt{https://www.hkbibleconference.org/session-message/view/444}}
\newline
\newline
\hyperref[sec:443]{< < < PREV SERMON < < <}
~
\hyperref[sec:toc]{[返目錄]}
~
\hyperref[sec:445]{> > > NEXT SERMON > > >}
\newline
\newline
彼得前書 2:18-2:23
\newline
\begin{longtable}{cl}
\hline
\hline
章節 & 經文 (和合本修訂版)\\
\hline
2:18 & \begin{tabularx}{0.7\textwidth}{X} 你們作奴僕的，凡事要存敬畏的心順服主人；不但順服善良溫和的，就是乖僻的也要順服。 \end{tabularx} \\ \\ \relax
2:19 & \begin{tabularx}{0.7\textwidth}{X} 倘若你們為使良心對得起神，忍受冤屈的痛苦，這是可讚許的。 \end{tabularx} \\ \\ \relax
2:20 & \begin{tabularx}{0.7\textwidth}{X} 你們若因犯罪受責打而忍耐，有甚麼可稱讚的呢？但你們若因行善受苦而忍耐，這在神看來是可讚許的。 \end{tabularx} \\ \\ \relax
2:21 & \begin{tabularx}{0.7\textwidth}{X} 你們蒙召就是為此，因為基督也為你們受過苦，給你們留下榜樣，為要使你們跟隨他的腳蹤。 \end{tabularx} \\ \\ \relax
2:22 & \begin{tabularx}{0.7\textwidth}{X} 「他並沒有犯罪，口裡也沒有詭詐。」 \end{tabularx} \\ \\ \relax
2:23 & \begin{tabularx}{0.7\textwidth}{X} 他被辱罵不還口，受害也不說威嚇的話，只將自己交託給公義的審判者。 \end{tabularx} \\ \\
[1ex]
\hline
\hline
\end{longtable}


\newpage

\section{第74屆~港九培靈會~戴紹曾~第九講}
\label{sec:445}
\textbf{蒙召得享榮耀}
\newline
\newline
連結: \href{https://www.hkbibleconference.org/session-message/view/445}{\texttt{https://www.hkbibleconference.org/session-message/view/445}}
\newline
\newline
\hyperref[sec:444]{< < < PREV SERMON < < <}
~
\hyperref[sec:toc]{[返目錄]}
~
\hyperref[sec:446]{> > > NEXT SERMON > > >}
\newline
\newline
彼得前書 5:10-5:10
\newline
\begin{longtable}{cl}
\hline
\hline
章節 & 經文 (和合本修訂版)\\
\hline
5:10 & \begin{tabularx}{0.7\textwidth}{X} 那賜一切恩典的神曾在基督裡召了你們，得享他永遠的榮耀，在你們暫受苦難之後，必要親自成全你們，堅固你們，賜力量給你們，建立你們。 \end{tabularx} \\ \\
[1ex]
\hline
\hline
\end{longtable}


\newpage

\section{第74屆~港九培靈會~戴紹曾~第十講}
\label{sec:446}
\textbf{蒙召殷勤堅定}
\newline
\newline
連結: \href{https://www.hkbibleconference.org/session-message/view/446}{\texttt{https://www.hkbibleconference.org/session-message/view/446}}
\newline
\newline
\hyperref[sec:445]{< < < PREV SERMON < < <}
~
\hyperref[sec:toc]{[返目錄]}
~
\hyperref[sec:419]{> > > NEXT SERMON > > >}
\newline
\newline
彼得後書 1:10-1:10
\newline
\begin{longtable}{cl}
\hline
\hline
章節 & 經文 (和合本修訂版)\\
\hline
1:10 & \begin{tabularx}{0.7\textwidth}{X} 所以，弟兄們，要更加努力，使你們的蒙召和被選堅定不移。你們實行這幾樣，就永不失腳。 \end{tabularx} \\ \\
[1ex]
\hline
\hline
\end{longtable}


\newpage

\section{第74屆~港九培靈會~江耀全~第一講}
\label{sec:419}
\textbf{愛的福音}
\newline
\newline
連結: \href{https://www.hkbibleconference.org/session-message/view/419}{\texttt{https://www.hkbibleconference.org/session-message/view/419}}
\newline
\newline
\hyperref[sec:446]{< < < PREV SERMON < < <}
~
\hyperref[sec:toc]{[返目錄]}
~
\hyperref[sec:420]{> > > NEXT SERMON > > >}
\newline
\newline
路加福音 4:14-4:21
\newline
\begin{longtable}{cl}
\hline
\hline
章節 & 經文 (和合本修訂版)\\
\hline
4:14 & \begin{tabularx}{0.7\textwidth}{X} 耶穌帶著聖靈的能力回到加利利，他的名聲傳遍了四方。 \end{tabularx} \\ \\ \relax
4:15 & \begin{tabularx}{0.7\textwidth}{X} 他在各會堂裡教導人，眾人都稱讚他。 \end{tabularx} \\ \\ \relax
4:16 & \begin{tabularx}{0.7\textwidth}{X} 耶穌來到拿撒勒，就是他長大的地方。在安息日，照他素常的規矩進了會堂，站起來要念聖經。 \end{tabularx} \\ \\ \relax
4:17 & \begin{tabularx}{0.7\textwidth}{X} 有人把以賽亞先知的書交給他，他就打開，找到一處寫著： \end{tabularx} \\ \\ \relax
4:18 & \begin{tabularx}{0.7\textwidth}{X} 「主的靈在我身上，因為他用膏膏我，叫我傳福音給貧窮的人；差遣我宣告：被擄的得釋放，失明的得看見，受壓迫的得自由， \end{tabularx} \\ \\ \relax
4:19 & \begin{tabularx}{0.7\textwidth}{X} 宣告神悅納人的禧年。」 \end{tabularx} \\ \\ \relax
4:20 & \begin{tabularx}{0.7\textwidth}{X} 於是他把書捲起來，交還給管理人，就坐下。會堂裡的人都定睛看他。 \end{tabularx} \\ \\ \relax
4:21 & \begin{tabularx}{0.7\textwidth}{X} 耶穌對他們說：「你們聽見的這段經文，今天已經應驗了。」 \end{tabularx} \\ \\
[1ex]
\hline
\hline
\end{longtable}


\newpage

\section{第74屆~港九培靈會~江耀全~第二講}
\label{sec:420}
\textbf{愛的順服}
\newline
\newline
連結: \href{https://www.hkbibleconference.org/session-message/view/420}{\texttt{https://www.hkbibleconference.org/session-message/view/420}}
\newline
\newline
\hyperref[sec:419]{< < < PREV SERMON < < <}
~
\hyperref[sec:toc]{[返目錄]}
~
\hyperref[sec:421]{> > > NEXT SERMON > > >}
\newline
\newline
路加福音 1:46-1:56
\newline
\begin{longtable}{cl}
\hline
\hline
章節 & 經文 (和合本修訂版)\\
\hline
1:46 & \begin{tabularx}{0.7\textwidth}{X} 馬利亞說：「我心尊主為大； \end{tabularx} \\ \\ \relax
1:47 & \begin{tabularx}{0.7\textwidth}{X} 我靈以神我的救主為樂； \end{tabularx} \\ \\ \relax
1:48 & \begin{tabularx}{0.7\textwidth}{X} 因為他顧念他使女的卑微；從今以後，萬代要稱我有福。 \end{tabularx} \\ \\ \relax
1:49 & \begin{tabularx}{0.7\textwidth}{X} 因為那有權能的為我做了大事；他的名是聖的。 \end{tabularx} \\ \\ \relax
1:50 & \begin{tabularx}{0.7\textwidth}{X} 他憐憫敬畏他的人，直到世世代代。 \end{tabularx} \\ \\ \relax
1:51 & \begin{tabularx}{0.7\textwidth}{X} 他用膀臂施展大能；他趕散心裡妄想的狂傲人。 \end{tabularx} \\ \\ \relax
1:52 & \begin{tabularx}{0.7\textwidth}{X} 他叫有權柄的失位，叫卑賤的升高。 \end{tabularx} \\ \\ \relax
1:53 & \begin{tabularx}{0.7\textwidth}{X} 他叫飢餓的飽餐美食，叫富足的空手回去。 \end{tabularx} \\ \\ \relax
1:54 & \begin{tabularx}{0.7\textwidth}{X} 他扶助了他的僕人以色列，不忘記施憐憫， \end{tabularx} \\ \\ \relax
1:55 & \begin{tabularx}{0.7\textwidth}{X} 正如他對我們的列祖說過，『憐憫亞伯拉罕和他的後裔，直到永遠。』」 \end{tabularx} \\ \\ \relax
1:56 & \begin{tabularx}{0.7\textwidth}{X} 馬利亞和伊利莎白同住，約有三個月，然後回家去了。 \end{tabularx} \\ \\
[1ex]
\hline
\hline
\end{longtable}


\newpage

\section{第74屆~港九培靈會~江耀全~第三講}
\label{sec:421}
\textbf{愛的恩典}
\newline
\newline
連結: \href{https://www.hkbibleconference.org/session-message/view/421}{\texttt{https://www.hkbibleconference.org/session-message/view/421}}
\newline
\newline
\hyperref[sec:420]{< < < PREV SERMON < < <}
~
\hyperref[sec:toc]{[返目錄]}
~
\hyperref[sec:422]{> > > NEXT SERMON > > >}
\newline
\newline
路加福音 7:36-7:50
\newline
\begin{longtable}{cl}
\hline
\hline
章節 & 經文 (和合本修訂版)\\
\hline
7:36 & \begin{tabularx}{0.7\textwidth}{X} 有一個法利賽人請耶穌和他吃飯，耶穌就到那法利賽人家裡去坐席。 \end{tabularx} \\ \\ \relax
7:37 & \begin{tabularx}{0.7\textwidth}{X} 那城裡有一個女人，是個罪人，知道耶穌在法利賽人家裡坐席，就拿著盛滿香膏的玉瓶， \end{tabularx} \\ \\ \relax
7:38 & \begin{tabularx}{0.7\textwidth}{X} 站在耶穌背後，挨著他的腳哭，眼淚滴濕了耶穌的腳，就用自己的頭髮擦乾，又用嘴連連親他的腳，把香膏抹上。 \end{tabularx} \\ \\ \relax
7:39 & \begin{tabularx}{0.7\textwidth}{X} 請耶穌的法利賽人看見這事，心裡說：「這人若是先知，一定知道摸他的是誰，是個怎樣的女人；她是個罪人哪！」 \end{tabularx} \\ \\ \relax
7:40 & \begin{tabularx}{0.7\textwidth}{X} 耶穌回應他說：「西門，我有話要對你說。」西門說：「老師，請說。」 \end{tabularx} \\ \\ \relax
7:41 & \begin{tabularx}{0.7\textwidth}{X} 耶穌說：「有兩個人欠了某一個債主的錢，一個欠五百個銀幣，一個欠五十個銀幣。 \end{tabularx} \\ \\ \relax
7:42 & \begin{tabularx}{0.7\textwidth}{X} 因為他們無力償還，債主就開恩赦免了他們兩個人的債。那麼，這兩個人哪一個更愛他呢？」 \end{tabularx} \\ \\ \relax
7:43 & \begin{tabularx}{0.7\textwidth}{X} 西門回答：「我想是那多得赦免的人。」耶穌對他說：「你的判斷不錯。」 \end{tabularx} \\ \\ \relax
7:44 & \begin{tabularx}{0.7\textwidth}{X} 於是他轉過來向著那女人，對西門說：「你看見這女人嗎？我進了你的家，你沒有給我水洗腳，但這女人用眼淚滴濕了我的腳，又用頭髮擦乾。 \end{tabularx} \\ \\ \relax
7:45 & \begin{tabularx}{0.7\textwidth}{X} 你沒有親我，但這女人從我進來就不住地親我的腳。 \end{tabularx} \\ \\ \relax
7:46 & \begin{tabularx}{0.7\textwidth}{X} 你沒有用油抹我的頭，但這女人用香膏抹我的腳。 \end{tabularx} \\ \\ \relax
7:47 & \begin{tabularx}{0.7\textwidth}{X} 所以我告訴你，她許多的罪都赦免了，因為她愛的多；而那少得赦免的，愛的就少。」 \end{tabularx} \\ \\ \relax
7:48 & \begin{tabularx}{0.7\textwidth}{X} 於是耶穌對那女人說：「你的罪都赦免了。」 \end{tabularx} \\ \\ \relax
7:49 & \begin{tabularx}{0.7\textwidth}{X} 同席的人心裡說：「這是甚麼人，竟赦免人的罪呢？」 \end{tabularx} \\ \\ \relax
7:50 & \begin{tabularx}{0.7\textwidth}{X} 耶穌對那女人說：「你的信救了你，平安地回去吧！」 \end{tabularx} \\ \\
[1ex]
\hline
\hline
\end{longtable}


\newpage

\section{第74屆~港九培靈會~江耀全~第四講}
\label{sec:422}
\textbf{愛的承擔}
\newline
\newline
連結: \href{https://www.hkbibleconference.org/session-message/view/422}{\texttt{https://www.hkbibleconference.org/session-message/view/422}}
\newline
\newline
\hyperref[sec:421]{< < < PREV SERMON < < <}
~
\hyperref[sec:toc]{[返目錄]}
~
\hyperref[sec:423]{> > > NEXT SERMON > > >}
\newline
\newline
路加福音 10:25-10:37
\newline
\begin{longtable}{cl}
\hline
\hline
章節 & 經文 (和合本修訂版)\\
\hline
10:25 & \begin{tabularx}{0.7\textwidth}{X} 有一個律法師起來試探耶穌，說：「老師！我該做甚麼才可以承受永生？」 \end{tabularx} \\ \\ \relax
10:26 & \begin{tabularx}{0.7\textwidth}{X} 耶穌對他說：「律法上寫的是甚麼？你是怎樣念的呢？」 \end{tabularx} \\ \\ \relax
10:27 & \begin{tabularx}{0.7\textwidth}{X} 他回答說：「你要盡心、盡性、盡力、盡意愛主—你的神，又要愛鄰如己。」 \end{tabularx} \\ \\ \relax
10:28 & \begin{tabularx}{0.7\textwidth}{X} 耶穌對他說：「你回答得正確，你這樣做就會得永生。」 \end{tabularx} \\ \\ \relax
10:29 & \begin{tabularx}{0.7\textwidth}{X} 那人要證明自己有理，就對耶穌說：「誰是我的鄰舍呢？」 \end{tabularx} \\ \\ \relax
10:30 & \begin{tabularx}{0.7\textwidth}{X} 耶穌回答：「有一個人從耶路撒冷下耶利哥去，落在強盜手中。他們剝去他的衣裳，把他打個半死，丟下他走了。 \end{tabularx} \\ \\ \relax
10:31 & \begin{tabularx}{0.7\textwidth}{X} 偶然有一個祭司從那條路下來，看見他就從另一邊過去了。 \end{tabularx} \\ \\ \relax
10:32 & \begin{tabularx}{0.7\textwidth}{X} 又有一個利未人來到那裡，看見他，也照樣從另一邊過去了。 \end{tabularx} \\ \\ \relax
10:33 & \begin{tabularx}{0.7\textwidth}{X} 可是，有一個撒瑪利亞人路過那裡，看見他就動了慈心， \end{tabularx} \\ \\ \relax
10:34 & \begin{tabularx}{0.7\textwidth}{X} 上前用油和酒倒在他的傷處，包裹好了，扶他騎上自己的牲口，帶他到旅店裡去，照應他。 \end{tabularx} \\ \\ \relax
10:35 & \begin{tabularx}{0.7\textwidth}{X} 第二天，他拿出兩個銀幣來，交給店主，說：『請你照應他，額外的費用，我回來時會還你。』 \end{tabularx} \\ \\ \relax
10:36 & \begin{tabularx}{0.7\textwidth}{X} 你想，這三個人哪一個是落在強盜手中那人的鄰舍呢？」 \end{tabularx} \\ \\ \relax
10:37 & \begin{tabularx}{0.7\textwidth}{X} 他說：「是憐憫他的。」耶穌對他說：「你去，照樣做吧！」 \end{tabularx} \\ \\
[1ex]
\hline
\hline
\end{longtable}


\newpage

\section{第74屆~港九培靈會~江耀全~第五講}
\label{sec:423}
\textbf{愛的服侍}
\newline
\newline
連結: \href{https://www.hkbibleconference.org/session-message/view/423}{\texttt{https://www.hkbibleconference.org/session-message/view/423}}
\newline
\newline
\hyperref[sec:422]{< < < PREV SERMON < < <}
~
\hyperref[sec:toc]{[返目錄]}
~
\hyperref[sec:424]{> > > NEXT SERMON > > >}
\newline
\newline
路加福音 10:38-10:42
\newline
\begin{longtable}{cl}
\hline
\hline
章節 & 經文 (和合本修訂版)\\
\hline
10:38 & \begin{tabularx}{0.7\textwidth}{X} 他們繼續前行，耶穌進了一個村莊。有一個女人，名叫馬大，接他到自己家裡。 \end{tabularx} \\ \\ \relax
10:39 & \begin{tabularx}{0.7\textwidth}{X} 她有一個妹妹，名叫馬利亞，在主的腳前坐著聽他的道。 \end{tabularx} \\ \\ \relax
10:40 & \begin{tabularx}{0.7\textwidth}{X} 馬大伺候的事多，心裡忙亂，進前來，說：「主啊，我的妹妹留下我一個人伺候，你不在意嗎？請吩咐她來幫助我。」 \end{tabularx} \\ \\ \relax
10:41 & \begin{tabularx}{0.7\textwidth}{X} 主回答說：「馬大，馬大，你為許多的事操心煩惱， \end{tabularx} \\ \\ \relax
10:42 & \begin{tabularx}{0.7\textwidth}{X} 但是不可少的只有一件。馬利亞已經選擇了那上好的福分，是沒有人能從她奪去的。」 \end{tabularx} \\ \\
[1ex]
\hline
\hline
\end{longtable}


\newpage

\section{第74屆~港九培靈會~江耀全~第六講}
\label{sec:424}
\textbf{愛的寬宏}
\newline
\newline
連結: \href{https://www.hkbibleconference.org/session-message/view/424}{\texttt{https://www.hkbibleconference.org/session-message/view/424}}
\newline
\newline
\hyperref[sec:423]{< < < PREV SERMON < < <}
~
\hyperref[sec:toc]{[返目錄]}
~
\hyperref[sec:425]{> > > NEXT SERMON > > >}
\newline
\newline
路加福音 15:11-15:32
\newline
\begin{longtable}{cl}
\hline
\hline
章節 & 經文 (和合本修訂版)\\
\hline
15:11 & \begin{tabularx}{0.7\textwidth}{X} 耶穌又說：「一個人有兩個兒子。 \end{tabularx} \\ \\ \relax
15:12 & \begin{tabularx}{0.7\textwidth}{X} 小兒子對父親說：『父親，請你把我應得的家業分給我。』他父親就把財產分給他們。 \end{tabularx} \\ \\ \relax
15:13 & \begin{tabularx}{0.7\textwidth}{X} 過了不多幾天，小兒子把他一切所有的都收拾起來，往遠方去了。在那裡，他任意放蕩，浪費錢財。 \end{tabularx} \\ \\ \relax
15:14 & \begin{tabularx}{0.7\textwidth}{X} 他耗盡了一切所有的，又恰逢那地方有大饑荒，就窮困起來。 \end{tabularx} \\ \\ \relax
15:15 & \begin{tabularx}{0.7\textwidth}{X} 於是他去投靠當地的一個居民，那人打發他到田裡去放豬。 \end{tabularx} \\ \\ \relax
15:16 & \begin{tabularx}{0.7\textwidth}{X} 他恨不得拿豬所吃的豆莢充飢，也沒有人給他甚麼吃的。 \end{tabularx} \\ \\ \relax
15:17 & \begin{tabularx}{0.7\textwidth}{X} 他醒悟過來，就說：『我父親有多少雇工，糧食有餘，我倒在這裡餓死嗎？ \end{tabularx} \\ \\ \relax
15:18 & \begin{tabularx}{0.7\textwidth}{X} 我要起來，到我父親那裡去，對他說：父親！我得罪了天，又得罪了你， \end{tabularx} \\ \\ \relax
15:19 & \begin{tabularx}{0.7\textwidth}{X} 從今以後，我不配稱為你的兒子，把我當作一個雇工吧。』 \end{tabularx} \\ \\ \relax
15:20 & \begin{tabularx}{0.7\textwidth}{X} 於是他起來，往他父親那裡去。相離還遠，他父親看見，就動了慈心，跑去擁抱著他，連連親他。 \end{tabularx} \\ \\ \relax
15:21 & \begin{tabularx}{0.7\textwidth}{X} 兒子對他說：『父親！我得罪了天，又得罪了你，從今以後，我不配稱為你的兒子。』 \end{tabularx} \\ \\ \relax
15:22 & \begin{tabularx}{0.7\textwidth}{X} 父親卻吩咐僕人：『快把那上好的袍子拿出來給他穿，把戒指戴在他指頭上，把鞋穿在他腳上， \end{tabularx} \\ \\ \relax
15:23 & \begin{tabularx}{0.7\textwidth}{X} 把那肥牛犢牽來宰了，我們來吃喝慶祝； \end{tabularx} \\ \\ \relax
15:24 & \begin{tabularx}{0.7\textwidth}{X} 因為我這個兒子是死而復活，失而復得的。』他們就開始慶祝。 \end{tabularx} \\ \\ \relax
15:25 & \begin{tabularx}{0.7\textwidth}{X} 「那時，大兒子正在田裡。他回來，離家不遠時，聽見奏樂跳舞的聲音， \end{tabularx} \\ \\ \relax
15:26 & \begin{tabularx}{0.7\textwidth}{X} 就叫一個僮僕來，問是甚麼事。 \end{tabularx} \\ \\ \relax
15:27 & \begin{tabularx}{0.7\textwidth}{X} 僮僕對他說：『你弟弟回來了，你父親因為他無災無病地回來，把肥牛犢宰了。』 \end{tabularx} \\ \\ \relax
15:28 & \begin{tabularx}{0.7\textwidth}{X} 大兒子就生氣，不肯進去，他父親出來勸他。 \end{tabularx} \\ \\ \relax
15:29 & \begin{tabularx}{0.7\textwidth}{X} 他對父親說：『你看，我服侍你這麼多年，從來沒有違背過你的命令，而你從來沒有給我一隻小山羊，叫我和朋友們一同快樂。 \end{tabularx} \\ \\ \relax
15:30 & \begin{tabularx}{0.7\textwidth}{X} 但你這個兒子和娼妓吃光了你的財產，他一回來，你倒為他宰了肥牛犢。』 \end{tabularx} \\ \\ \relax
15:31 & \begin{tabularx}{0.7\textwidth}{X} 父親對他說：『兒啊！你常和我同在，我所有的一切都是你的； \end{tabularx} \\ \\ \relax
15:32 & \begin{tabularx}{0.7\textwidth}{X} 可是你這個弟弟是死而復活，失而復得的，所以我們理當歡喜慶祝。』」 \end{tabularx} \\ \\
[1ex]
\hline
\hline
\end{longtable}


\newpage

\section{第74屆~港九培靈會~江耀全~第七講}
\label{sec:425}
\textbf{愛的感恩}
\newline
\newline
連結: \href{https://www.hkbibleconference.org/session-message/view/425}{\texttt{https://www.hkbibleconference.org/session-message/view/425}}
\newline
\newline
\hyperref[sec:424]{< < < PREV SERMON < < <}
~
\hyperref[sec:toc]{[返目錄]}
~
\hyperref[sec:426]{> > > NEXT SERMON > > >}
\newline
\newline
路加福音 17:11-17:19
\newline
\begin{longtable}{cl}
\hline
\hline
章節 & 經文 (和合本修訂版)\\
\hline
17:11 & \begin{tabularx}{0.7\textwidth}{X} 耶穌往耶路撒冷去，經過撒瑪利亞和加利利中間的地區。 \end{tabularx} \\ \\ \relax
17:12 & \begin{tabularx}{0.7\textwidth}{X} 他進入一個村子，有十個痲瘋病人迎面而來，遠遠地站著， \end{tabularx} \\ \\ \relax
17:13 & \begin{tabularx}{0.7\textwidth}{X} 高聲說：「耶穌，老師啊，可憐我們吧！」 \end{tabularx} \\ \\ \relax
17:14 & \begin{tabularx}{0.7\textwidth}{X} 耶穌看見，就對他們說：「你們去，把身體給祭司檢查。」他們正去的時候就潔淨了。 \end{tabularx} \\ \\ \relax
17:15 & \begin{tabularx}{0.7\textwidth}{X} 其中有一個見自己已經好了，就回來大聲歸榮耀給神， \end{tabularx} \\ \\ \relax
17:16 & \begin{tabularx}{0.7\textwidth}{X} 又俯伏在耶穌腳前感謝他。這人是撒瑪利亞人。 \end{tabularx} \\ \\ \relax
17:17 & \begin{tabularx}{0.7\textwidth}{X} 耶穌回答說：「潔淨了的不是十個人嗎？那九個在哪裡呢？ \end{tabularx} \\ \\ \relax
17:18 & \begin{tabularx}{0.7\textwidth}{X} 除了這外族人，再沒有別人回來歸榮耀給神嗎？」 \end{tabularx} \\ \\ \relax
17:19 & \begin{tabularx}{0.7\textwidth}{X} 於是他對那人說：「起來，走吧，你的信救了你！」 \end{tabularx} \\ \\
[1ex]
\hline
\hline
\end{longtable}


\newpage

\section{第74屆~港九培靈會~江耀全~第八講}
\label{sec:426}
\textbf{愛的相遇}
\newline
\newline
連結: \href{https://www.hkbibleconference.org/session-message/view/426}{\texttt{https://www.hkbibleconference.org/session-message/view/426}}
\newline
\newline
\hyperref[sec:425]{< < < PREV SERMON < < <}
~
\hyperref[sec:toc]{[返目錄]}
~
\hyperref[sec:427]{> > > NEXT SERMON > > >}
\newline
\newline
路加福音 19:1-19:10
\newline
\begin{longtable}{cl}
\hline
\hline
章節 & 經文 (和合本修訂版)\\
\hline
19:1 & \begin{tabularx}{0.7\textwidth}{X} 耶穌進了耶利哥，要從那裡經過。 \end{tabularx} \\ \\ \relax
19:2 & \begin{tabularx}{0.7\textwidth}{X} 有一個人名叫撒該，作稅吏長，是個財主。 \end{tabularx} \\ \\ \relax
19:3 & \begin{tabularx}{0.7\textwidth}{X} 他要看看耶穌是怎樣的人，只因人多，他的身材又矮，所以看不見。 \end{tabularx} \\ \\ \relax
19:4 & \begin{tabularx}{0.7\textwidth}{X} 於是他跑到前頭，爬上桑樹，要看耶穌，因為耶穌要從那裡經過。 \end{tabularx} \\ \\ \relax
19:5 & \begin{tabularx}{0.7\textwidth}{X} 耶穌到了那裡，抬頭一看，對他說：「撒該，快下來！今天我必須住在你家裡。」 \end{tabularx} \\ \\ \relax
19:6 & \begin{tabularx}{0.7\textwidth}{X} 他就急忙下來，歡歡喜喜地接待耶穌。 \end{tabularx} \\ \\ \relax
19:7 & \begin{tabularx}{0.7\textwidth}{X} 眾人看見，都私下議論說：「他竟然到罪人家裡去住宿。」 \end{tabularx} \\ \\ \relax
19:8 & \begin{tabularx}{0.7\textwidth}{X} 撒該站著對主說：「主啊，我把所有的一半給窮人；我若勒索了誰，就還他四倍。」 \end{tabularx} \\ \\ \relax
19:9 & \begin{tabularx}{0.7\textwidth}{X} 耶穌對他說：「今天救恩到了這家，因為他也是亞伯拉罕的子孫。 \end{tabularx} \\ \\ \relax
19:10 & \begin{tabularx}{0.7\textwidth}{X} 人子來是要尋找和拯救失喪的人。」 \end{tabularx} \\ \\
[1ex]
\hline
\hline
\end{longtable}


\newpage

\section{第74屆~港九培靈會~江耀全~第九講}
\label{sec:427}
\textbf{愛的同在}
\newline
\newline
連結: \href{https://www.hkbibleconference.org/session-message/view/427}{\texttt{https://www.hkbibleconference.org/session-message/view/427}}
\newline
\newline
\hyperref[sec:426]{< < < PREV SERMON < < <}
~
\hyperref[sec:toc]{[返目錄]}
~
\hyperref[sec:428]{> > > NEXT SERMON > > >}
\newline
\newline
路加福音 24:13-24:35
\newline
\begin{longtable}{cl}
\hline
\hline
章節 & 經文 (和合本修訂版)\\
\hline
24:13 & \begin{tabularx}{0.7\textwidth}{X} 同一天，門徒中有兩個人往一個村子去；這村子名叫以馬忤斯，離耶路撒冷約有二十五里。 \end{tabularx} \\ \\ \relax
24:14 & \begin{tabularx}{0.7\textwidth}{X} 他們彼此談論所發生的這一切事。 \end{tabularx} \\ \\ \relax
24:15 & \begin{tabularx}{0.7\textwidth}{X} 正交談議論的時候，耶穌親自走近他們，和他們同行， \end{tabularx} \\ \\ \relax
24:16 & \begin{tabularx}{0.7\textwidth}{X} 可是他們的眼睛模糊了，沒認出他。 \end{tabularx} \\ \\ \relax
24:17 & \begin{tabularx}{0.7\textwidth}{X} 耶穌對他們說：「你們一邊走一邊談，彼此談論的是甚麼事呢？」他們就站住，臉上帶著愁容。 \end{tabularx} \\ \\ \relax
24:18 & \begin{tabularx}{0.7\textwidth}{X} 兩人中有一個名叫革流巴的回答：「你是在耶路撒冷的旅客中，惟一還不知道這幾天在那裡發生了甚麼事的人嗎？」 \end{tabularx} \\ \\ \relax
24:19 & \begin{tabularx}{0.7\textwidth}{X} 耶穌對他們說：「甚麼事呢？」他們對他說：「就是拿撒勒人耶穌的事。他是個先知，在神和眾百姓面前，說話行事都大有能力。 \end{tabularx} \\ \\ \relax
24:20 & \begin{tabularx}{0.7\textwidth}{X} 祭司長們和我們的官長竟把他解去，定了死罪，釘在十字架上。 \end{tabularx} \\ \\ \relax
24:21 & \begin{tabularx}{0.7\textwidth}{X} 但我們素來所盼望要救贖以色列民的就是他。不但如此，這些事發生到現在已經三天了。 \end{tabularx} \\ \\ \relax
24:22 & \begin{tabularx}{0.7\textwidth}{X} 還有，我們中間的幾個婦女使我們驚奇：她們清早去了墳墓， \end{tabularx} \\ \\ \relax
24:23 & \begin{tabularx}{0.7\textwidth}{X} 不見他的身體，就回來告訴我們，說她們看見了天使顯現，說他活了。 \end{tabularx} \\ \\ \relax
24:24 & \begin{tabularx}{0.7\textwidth}{X} 又有我們的幾個人往墳墓那裡去，所發現的正如婦女們所說的，只是沒有看見他。」 \end{tabularx} \\ \\ \relax
24:25 & \begin{tabularx}{0.7\textwidth}{X} 耶穌對他們說：「無知的人哪，先知所說的一切話，你們的心信得太遲鈍了。 \end{tabularx} \\ \\ \relax
24:26 & \begin{tabularx}{0.7\textwidth}{X} 基督不是必須受這些苦難，然後進入他的榮耀嗎？」 \end{tabularx} \\ \\ \relax
24:27 & \begin{tabularx}{0.7\textwidth}{X} 於是，他從摩西和眾先知起，凡經上所指著自己的話都給他們作了解釋。 \end{tabularx} \\ \\ \relax
24:28 & \begin{tabularx}{0.7\textwidth}{X} 他們走近所要去的村子，耶穌好像還要往前走， \end{tabularx} \\ \\ \relax
24:29 & \begin{tabularx}{0.7\textwidth}{X} 他們卻強留他說：「時候晚了，天快黑了，請你同我們住下吧。」耶穌就進去，要同他們住下。 \end{tabularx} \\ \\ \relax
24:30 & \begin{tabularx}{0.7\textwidth}{X} 坐下來和他們用餐的時候，耶穌拿起餅來，祝福了，擘開，遞給他們。 \end{tabularx} \\ \\ \relax
24:31 & \begin{tabularx}{0.7\textwidth}{X} 他們的眼睛開了，這才認出他來。耶穌卻從他們眼前消失了。 \end{tabularx} \\ \\ \relax
24:32 & \begin{tabularx}{0.7\textwidth}{X} 他們彼此說：「在路上他和我們說話，給我們講解聖經的時候，我們的心在我們裡面豈不是火熱的嗎？」 \end{tabularx} \\ \\ \relax
24:33 & \begin{tabularx}{0.7\textwidth}{X} 於是他們立刻起身，回耶路撒冷去，看見十一個使徒和與他們正在一起的人聚集在一處， \end{tabularx} \\ \\ \relax
24:34 & \begin{tabularx}{0.7\textwidth}{X} 說：「主果然復活了，已經顯現給西門看了。」 \end{tabularx} \\ \\ \relax
24:35 & \begin{tabularx}{0.7\textwidth}{X} 於是，兩個人把路上所遇到，和耶穌擘餅的時候怎麼被他們認出來的事，都述說了一遍。 \end{tabularx} \\ \\
[1ex]
\hline
\hline
\end{longtable}


\newpage

\section{第74屆~港九培靈會~溫偉耀~第一講}
\label{sec:428}
\textbf{軟弱與剛強 -- 寶貝與瓦器}
\newline
\newline
連結: \href{https://www.hkbibleconference.org/session-message/view/428}{\texttt{https://www.hkbibleconference.org/session-message/view/428}}
\newline
\newline
\hyperref[sec:427]{< < < PREV SERMON < < <}
~
\hyperref[sec:toc]{[返目錄]}
~
\hyperref[sec:429]{> > > NEXT SERMON > > >}
\newline
\newline
哥林多後書 4:7-4:11
\newline
\begin{longtable}{cl}
\hline
\hline
章節 & 經文 (和合本修訂版)\\
\hline
4:7 & \begin{tabularx}{0.7\textwidth}{X} 我們有這寶貝放在瓦器裡，為要顯明這莫大的能力是出於神，不是出於我們。 \end{tabularx} \\ \\ \relax
4:8 & \begin{tabularx}{0.7\textwidth}{X} 我們處處受困，卻不被捆住；內心困擾，卻沒有絕望； \end{tabularx} \\ \\ \relax
4:9 & \begin{tabularx}{0.7\textwidth}{X} 遭受迫害，卻不被撇棄；擊倒在地，卻不致滅亡。 \end{tabularx} \\ \\ \relax
4:10 & \begin{tabularx}{0.7\textwidth}{X} 我們身上常帶著耶穌的死，使耶穌的生也在我們身上顯明。 \end{tabularx} \\ \\ \relax
4:11 & \begin{tabularx}{0.7\textwidth}{X} 因為我們這活著的人常為耶穌被置於死地，使耶穌的生命在我們這必死的人身上顯明出來。 \end{tabularx} \\ \\
[1ex]
\hline
\hline
\end{longtable}


\newpage

\section{第74屆~港九培靈會~溫偉耀~第二講}
\label{sec:429}
\textbf{從高峰到深淵 -- 「三層天」與「一根刺」}
\newline
\newline
連結: \href{https://www.hkbibleconference.org/session-message/view/429}{\texttt{https://www.hkbibleconference.org/session-message/view/429}}
\newline
\newline
\hyperref[sec:428]{< < < PREV SERMON < < <}
~
\hyperref[sec:toc]{[返目錄]}
~
\hyperref[sec:430]{> > > NEXT SERMON > > >}
\newline
\newline
哥林多後書 12:1-12:10
\newline
\begin{longtable}{cl}
\hline
\hline
章節 & 經文 (和合本修訂版)\\
\hline
12:1 & \begin{tabularx}{0.7\textwidth}{X} 雖然自誇無益，我還是不得不誇。我現在要提到主的異象和啟示。 \end{tabularx} \\ \\ \relax
12:2 & \begin{tabularx}{0.7\textwidth}{X} 我認識一個在基督裡的人，他在十四年前被提到第三層天上去；或在身內，我不知道，或在身外，我也不知道，只有神知道。 \end{tabularx} \\ \\ \relax
12:3 & \begin{tabularx}{0.7\textwidth}{X} 我認識的這樣的一個人—或在身內，或在身外，我都不知道，只有神知道— \end{tabularx} \\ \\ \relax
12:4 & \begin{tabularx}{0.7\textwidth}{X} 他被提到樂園裡，聽見隱祕的言語，是人不可說的。 \end{tabularx} \\ \\ \relax
12:5 & \begin{tabularx}{0.7\textwidth}{X} 為這人，我要誇口；但是為我自己，除了我的軟弱以外，我並不誇口。 \end{tabularx} \\ \\ \relax
12:6 & \begin{tabularx}{0.7\textwidth}{X} 就是我願意誇口也不算狂，因為我會說實話；只是我絕口不談，恐怕有人把我看得太高了，過於他在我身上所看見所聽見的； \end{tabularx} \\ \\ \relax
12:7 & \begin{tabularx}{0.7\textwidth}{X} 又恐怕我因所得的啟示太高深，就過於高抬自己，所以有一根刺加在我身上，就是撒但的差役來折磨我，免得我過於高抬自己。 \end{tabularx} \\ \\ \relax
12:8 & \begin{tabularx}{0.7\textwidth}{X} 為了這事，我曾三次求主使這根刺離開我。 \end{tabularx} \\ \\ \relax
12:9 & \begin{tabularx}{0.7\textwidth}{X} 他對我說：「我的恩典是夠你用的，因為我的能力是在人的軟弱上顯得完全。」所以，我更喜歡誇耀自己的軟弱，好使基督的能力覆庇我。 \end{tabularx} \\ \\ \relax
12:10 & \begin{tabularx}{0.7\textwidth}{X} 為基督的緣故，我以軟弱、凌辱、艱難、迫害、困苦為可喜樂的事；因為我甚麼時候軟弱，甚麼時候就剛強了。 \end{tabularx} \\ \\
[1ex]
\hline
\hline
\end{longtable}


\newpage

\section{第74屆~港九培靈會~溫偉耀~第三講}
\label{sec:430}
\textbf{向人推薦自己 -- 怎樣給人留下深刻的印象?}
\newline
\newline
連結: \href{https://www.hkbibleconference.org/session-message/view/430}{\texttt{https://www.hkbibleconference.org/session-message/view/430}}
\newline
\newline
\hyperref[sec:429]{< < < PREV SERMON < < <}
~
\hyperref[sec:toc]{[返目錄]}
~
\hyperref[sec:431]{> > > NEXT SERMON > > >}
\newline
\newline
哥林多後書 2:14-3:6
\newline
\begin{longtable}{cl}
\hline
\hline
章節 & 經文 (和合本修訂版)\\
\hline
2:14 & \begin{tabularx}{0.7\textwidth}{X} 感謝神！他常率領我們在基督裡得勝，並藉著我們在各處顯揚那因認識基督而有的香氣。 \end{tabularx} \\ \\ \relax
2:15 & \begin{tabularx}{0.7\textwidth}{X} 因為無論在得救的人或在滅亡的人當中，我們都是基督馨香之氣，是獻給神的。 \end{tabularx} \\ \\ \relax
2:16 & \begin{tabularx}{0.7\textwidth}{X} 對滅亡的人，這是死而又死的氣味；對得救的人，這是生而又生的氣味。這些事誰能當得起呢？ \end{tabularx} \\ \\ \relax
2:17 & \begin{tabularx}{0.7\textwidth}{X} 我們不像許多人，把神的道當商品販賣，而是由於真誠，而是受命於神，在神面前憑著基督講道。 \end{tabularx} \\ \\ \relax
3:1 & \begin{tabularx}{0.7\textwidth}{X} 難道我們又開始推薦自己嗎？難道我們像某些人那樣要用人的推薦信介紹給你們，或用你們的推薦信給人嗎？ \end{tabularx} \\ \\ \relax
3:2 & \begin{tabularx}{0.7\textwidth}{X} 你們就是我們的推薦信，寫在我們心裡，被眾人所知道、所誦讀的， \end{tabularx} \\ \\ \relax
3:3 & \begin{tabularx}{0.7\textwidth}{X} 而你們顯明自己是基督的書信，藉著我們寫成的。不是用墨寫的，而是用永生神的靈寫的；不是寫在石版上，而是寫在心版上的。 \end{tabularx} \\ \\ \relax
3:4 & \begin{tabularx}{0.7\textwidth}{X} 我們藉著基督才對神有這樣的信心。 \end{tabularx} \\ \\ \relax
3:5 & \begin{tabularx}{0.7\textwidth}{X} 並不是我們憑自己配做甚麼事，我們之所以配做是出於神； \end{tabularx} \\ \\ \relax
3:6 & \begin{tabularx}{0.7\textwidth}{X} 他使我們能配作新約的執事，不是文字上的約，而是聖靈的約；因為文字使人死，聖靈能使人活。 \end{tabularx} \\ \\
[1ex]
\hline
\hline
\end{longtable}


\newpage

\section{第74屆~港九培靈會~溫偉耀~第四講}
\label{sec:431}
\textbf{誇口 -- 怎樣評價自己?}
\newline
\newline
連結: \href{https://www.hkbibleconference.org/session-message/view/431}{\texttt{https://www.hkbibleconference.org/session-message/view/431}}
\newline
\newline
\hyperref[sec:430]{< < < PREV SERMON < < <}
~
\hyperref[sec:toc]{[返目錄]}
~
\hyperref[sec:432]{> > > NEXT SERMON > > >}
\newline
\newline
哥林多後書 11:21-11:33
\newline
\begin{longtable}{cl}
\hline
\hline
章節 & 經文 (和合本修訂版)\\
\hline
11:21 & \begin{tabularx}{0.7\textwidth}{X} 說來慚愧，在這方面好像我們是太軟弱了。然而，我說句蠢話，人在甚麼事上敢誇口，我也敢誇口。 \end{tabularx} \\ \\ \relax
11:22 & \begin{tabularx}{0.7\textwidth}{X} 他們是希伯來人嗎？我也是。他們是以色列人嗎？我也是。他們是亞伯拉罕的後裔嗎？我也是。 \end{tabularx} \\ \\ \relax
11:23 & \begin{tabularx}{0.7\textwidth}{X} 他們是基督的用人嗎？我說句狂話，我更是。我比他們忍受更多勞苦，坐過更多次監牢，受過無數次的鞭打，常常冒死。 \end{tabularx} \\ \\ \relax
11:24 & \begin{tabularx}{0.7\textwidth}{X} 我被猶太人鞭打五次，每次四十減去一下； \end{tabularx} \\ \\ \relax
11:25 & \begin{tabularx}{0.7\textwidth}{X} 被棍打了三次，被石頭打了一次，遭海難三次，一晝一夜在深海裡掙扎。 \end{tabularx} \\ \\ \relax
11:26 & \begin{tabularx}{0.7\textwidth}{X} 我又屢次行遠路，遭江河的危險，盜賊的危險，同族人的危險，外族人的危險，城裡的危險，曠野的危險，海中的危險，假弟兄的危險。 \end{tabularx} \\ \\ \relax
11:27 & \begin{tabularx}{0.7\textwidth}{X} 我勞碌困苦，常常失眠，又飢又渴，忍飢耐寒，赤身露體。 \end{tabularx} \\ \\ \relax
11:28 & \begin{tabularx}{0.7\textwidth}{X} 除了這些外表的事以外，我還有為眾教會操心的事天天壓在我身上。 \end{tabularx} \\ \\ \relax
11:29 & \begin{tabularx}{0.7\textwidth}{X} 有誰軟弱，我不軟弱呢？有誰跌倒，我不焦急呢？ \end{tabularx} \\ \\ \relax
11:30 & \begin{tabularx}{0.7\textwidth}{X} 我若必須誇口，就誇我軟弱的事好了。 \end{tabularx} \\ \\ \relax
11:31 & \begin{tabularx}{0.7\textwidth}{X} 那永遠可稱頌之主耶穌的父神知道我不說謊。 \end{tabularx} \\ \\ \relax
11:32 & \begin{tabularx}{0.7\textwidth}{X} 在大馬士革的亞哩達王手下的提督把守大馬士革城，要捉拿我， \end{tabularx} \\ \\ \relax
11:33 & \begin{tabularx}{0.7\textwidth}{X} 我被人用筐子從城牆上的窗口縋下，逃脫了他的手。 \end{tabularx} \\ \\
[1ex]
\hline
\hline
\end{longtable}


\newpage

\section{第74屆~港九培靈會~溫偉耀~第五講}
\label{sec:432}
\textbf{火浴鳳凰 -- 跟苦難交個朋友}
\newline
\newline
連結: \href{https://www.hkbibleconference.org/session-message/view/432}{\texttt{https://www.hkbibleconference.org/session-message/view/432}}
\newline
\newline
\hyperref[sec:431]{< < < PREV SERMON < < <}
~
\hyperref[sec:toc]{[返目錄]}
~
\hyperref[sec:433]{> > > NEXT SERMON > > >}
\newline
\newline
哥林多後書 1:3-1:10
\newline
\begin{longtable}{cl}
\hline
\hline
章節 & 經文 (和合本修訂版)\\
\hline
1:3 & \begin{tabularx}{0.7\textwidth}{X} 願頌讚歸於神—我們主耶穌基督的父；他是發慈悲的父，賜各樣安慰的神。 \end{tabularx} \\ \\ \relax
1:4 & \begin{tabularx}{0.7\textwidth}{X} 我們在一切患難中，他安慰我們，使我們能用神所賜的安慰去安慰那些遭各樣患難的人。 \end{tabularx} \\ \\ \relax
1:5 & \begin{tabularx}{0.7\textwidth}{X} 正如我們跟基督同受許多苦楚，我們也靠基督得許多安慰。 \end{tabularx} \\ \\ \relax
1:6 & \begin{tabularx}{0.7\textwidth}{X} 如果我們受患難，那是為使你們得安慰，得拯救；如果我們得安慰，那也是為使你們得安慰，這安慰能使你們忍受我們所受同樣的苦楚。 \end{tabularx} \\ \\ \relax
1:7 & \begin{tabularx}{0.7\textwidth}{X} 我們為你們所存的盼望是確定的，因為知道你們分擔了我們的痛苦，也要分享我們的安慰。 \end{tabularx} \\ \\ \relax
1:8 & \begin{tabularx}{0.7\textwidth}{X} 弟兄們，我們不要你們不知道，我們從前在亞細亞遭遇苦難，因受到無法忍受的壓力，甚至連活命的指望都沒有了。 \end{tabularx} \\ \\ \relax
1:9 & \begin{tabularx}{0.7\textwidth}{X} 自己心裡也斷定是必死無疑，這是要使我們不依靠自己，只依靠使死人復活的神。 \end{tabularx} \\ \\ \relax
1:10 & \begin{tabularx}{0.7\textwidth}{X} 他曾救我們脫離那極大的死亡，他要繼續救我們，而且我們指望他將來還要救我們。 \end{tabularx} \\ \\
[1ex]
\hline
\hline
\end{longtable}


\newpage

\section{第74屆~港九培靈會~溫偉耀~第六講}
\label{sec:433}
\textbf{榮光的職事 -- 上帝要怎樣的人去事奉祂?}
\newline
\newline
連結: \href{https://www.hkbibleconference.org/session-message/view/433}{\texttt{https://www.hkbibleconference.org/session-message/view/433}}
\newline
\newline
\hyperref[sec:432]{< < < PREV SERMON < < <}
~
\hyperref[sec:toc]{[返目錄]}
~
\hyperref[sec:436]{> > > NEXT SERMON > > >}
\newline
\newline
哥林多後書 3:4-3:11
\newline
\begin{longtable}{cl}
\hline
\hline
章節 & 經文 (和合本修訂版)\\
\hline
3:4 & \begin{tabularx}{0.7\textwidth}{X} 我們藉著基督才對神有這樣的信心。 \end{tabularx} \\ \\ \relax
3:5 & \begin{tabularx}{0.7\textwidth}{X} 並不是我們憑自己配做甚麼事，我們之所以配做是出於神； \end{tabularx} \\ \\ \relax
3:6 & \begin{tabularx}{0.7\textwidth}{X} 他使我們能配作新約的執事，不是文字上的約，而是聖靈的約；因為文字使人死，聖靈能使人活。 \end{tabularx} \\ \\ \relax
3:7 & \begin{tabularx}{0.7\textwidth}{X} 那用字刻在石頭上屬死的事奉尚且有榮光，以致以色列人因摩西臉上那逐漸褪色的榮光不能定睛看他的臉， \end{tabularx} \\ \\ \relax
3:8 & \begin{tabularx}{0.7\textwidth}{X} 那屬聖靈的事奉不是更有榮光嗎？ \end{tabularx} \\ \\ \relax
3:9 & \begin{tabularx}{0.7\textwidth}{X} 若是那使人定罪的事奉有榮光，那使人稱義的事奉的榮光就越發大了。 \end{tabularx} \\ \\ \relax
3:10 & \begin{tabularx}{0.7\textwidth}{X} 那從前有榮光的，因這更大的榮光，就算不得有榮光了； \end{tabularx} \\ \\ \relax
3:11 & \begin{tabularx}{0.7\textwidth}{X} 若是那逐漸褪色的有榮光，這長存的就更有榮光了。 \end{tabularx} \\ \\
[1ex]
\hline
\hline
\end{longtable}


\newpage

\section{第74屆~港九培靈會~溫偉耀~第七講}
\label{sec:436}
\textbf{「零風險」投資 -- 捐輸的屬靈策略與回報}
\newline
\newline
連結: \href{https://www.hkbibleconference.org/session-message/view/436}{\texttt{https://www.hkbibleconference.org/session-message/view/436}}
\newline
\newline
\hyperref[sec:433]{< < < PREV SERMON < < <}
~
\hyperref[sec:toc]{[返目錄]}
~
\hyperref[sec:437]{> > > NEXT SERMON > > >}
\newline
\newline
哥林多後書 9:6-9:15
\newline
\begin{longtable}{cl}
\hline
\hline
章節 & 經文 (和合本修訂版)\\
\hline
9:6 & \begin{tabularx}{0.7\textwidth}{X} 還有一點：「少種的少收；多種的多收。」 \end{tabularx} \\ \\ \relax
9:7 & \begin{tabularx}{0.7\textwidth}{X} 各人要隨心所願，不要為難，不要勉強，因為神愛樂捐的人。 \end{tabularx} \\ \\ \relax
9:8 & \begin{tabularx}{0.7\textwidth}{X} 神能將各樣的恩惠多多加給你們，使你們凡事常常充足，能多做各樣善事。 \end{tabularx} \\ \\ \relax
9:9 & \begin{tabularx}{0.7\textwidth}{X} 如經上所記：「他施捨，賙濟貧窮；他的義行存到永遠。」 \end{tabularx} \\ \\ \relax
9:10 & \begin{tabularx}{0.7\textwidth}{X} 那賜種子給撒種的，賜糧食給人吃的，必多多加給你們種地的種子，又增添你們仁義的果子。 \end{tabularx} \\ \\ \relax
9:11 & \begin{tabularx}{0.7\textwidth}{X} 你們必凡事富足，能多多施捨，使人藉著我們而生感謝神的心。 \end{tabularx} \\ \\ \relax
9:12 & \begin{tabularx}{0.7\textwidth}{X} 因為辦這供給的事，不但補聖徒的缺乏，而且使許多人對神充滿更多的感謝。 \end{tabularx} \\ \\ \relax
9:13 & \begin{tabularx}{0.7\textwidth}{X} 他們從這供給的事上得了憑據，知道你們宣認基督，順服他的福音，慷慨捐助給他們和眾人，把榮耀歸給神。 \end{tabularx} \\ \\ \relax
9:14 & \begin{tabularx}{0.7\textwidth}{X} 他們也因神極大的恩賜顯在你們身上而切切想念你們，為你們祈禱。 \end{tabularx} \\ \\ \relax
9:15 & \begin{tabularx}{0.7\textwidth}{X} 感謝神，因他有說不盡的恩賜！ \end{tabularx} \\ \\
[1ex]
\hline
\hline
\end{longtable}


\newpage

\section{第74屆~港九培靈會~溫偉耀~第八講}
\label{sec:437}
\textbf{新造的人 -- 活出轉化了的屬靈生命}
\newline
\newline
連結: \href{https://www.hkbibleconference.org/session-message/view/437}{\texttt{https://www.hkbibleconference.org/session-message/view/437}}
\newline
\newline
\hyperref[sec:436]{< < < PREV SERMON < < <}
~
\hyperref[sec:toc]{[返目錄]}
~
\hyperref[sec:438]{> > > NEXT SERMON > > >}
\newline
\newline
哥林多後書 5:14-5:21
\newline
\begin{longtable}{cl}
\hline
\hline
章節 & 經文 (和合本修訂版)\\
\hline
5:14 & \begin{tabularx}{0.7\textwidth}{X} 原來基督的愛激勵我們；因我們這樣斷定，一人既替眾人死，眾人就都死了； \end{tabularx} \\ \\ \relax
5:15 & \begin{tabularx}{0.7\textwidth}{X} 並且他替眾人死，是叫那些活著的人不再為自己活，乃為替他們死而復活的主活。 \end{tabularx} \\ \\ \relax
5:16 & \begin{tabularx}{0.7\textwidth}{X} 所以，從今以後，我們不再按照人的看法來認識人，縱使我們曾經按照人的看法認識基督，如今卻不再這樣認識他了。 \end{tabularx} \\ \\ \relax
5:17 & \begin{tabularx}{0.7\textwidth}{X} 所以，若有人在基督裡，他就是新造的人，舊事已過，都變成新的了。 \end{tabularx} \\ \\ \relax
5:18 & \begin{tabularx}{0.7\textwidth}{X} 一切都是出於神；他藉著基督使我們與他和好，又將勸人與他和好的使命賜給我們。 \end{tabularx} \\ \\ \relax
5:19 & \begin{tabularx}{0.7\textwidth}{X} 這就是：神在基督裡使世人與自己和好，不將他們的過犯歸到他們身上，並且將這和好的信息託付了我們。 \end{tabularx} \\ \\ \relax
5:20 & \begin{tabularx}{0.7\textwidth}{X} 所以，我們作基督的特使，就好像神藉我們勸你們一般。我們替基督求你們，與神和好吧！ \end{tabularx} \\ \\ \relax
5:21 & \begin{tabularx}{0.7\textwidth}{X} 神使那無罪的，替我們成為罪，好使我們在他裡面成為神的義。 \end{tabularx} \\ \\
[1ex]
\hline
\hline
\end{longtable}


\newpage

\section{第74屆~港九培靈會~溫偉耀~第九講}
\label{sec:438}
\textbf{今生與永恆 -- 灑脫而進取的生命情操}
\newline
\newline
連結: \href{https://www.hkbibleconference.org/session-message/view/438}{\texttt{https://www.hkbibleconference.org/session-message/view/438}}
\newline
\newline
\hyperref[sec:437]{< < < PREV SERMON < < <}
~
\hyperref[sec:toc]{[返目錄]}
~
\hyperref[sec:399]{> > > NEXT SERMON > > >}
\newline
\newline
哥林多後書 5:1-5:10
\newline
\begin{longtable}{cl}
\hline
\hline
章節 & 經文 (和合本修訂版)\\
\hline
5:1 & \begin{tabularx}{0.7\textwidth}{X} 因為我們知道，我們這地上的帳篷若拆毀了，我們將有神所造的居所，不是人手所造的，而是在天上永存的。 \end{tabularx} \\ \\ \relax
5:2 & \begin{tabularx}{0.7\textwidth}{X} 我們在這帳篷裡嘆息，渴望得到那從天上來的居所，好像穿上衣服； \end{tabularx} \\ \\ \relax
5:3 & \begin{tabularx}{0.7\textwidth}{X} 倘若脫下也不至於赤身了。 \end{tabularx} \\ \\ \relax
5:4 & \begin{tabularx}{0.7\textwidth}{X} 其實，我們在這帳篷裡的人勞苦嘆息，並不是願意脫下地上的帳篷，而是願意穿上天上的居所，好使這必死的被生命吞滅了。 \end{tabularx} \\ \\ \relax
5:5 & \begin{tabularx}{0.7\textwidth}{X} 那為我們安排這事的是神，他賜給我們聖靈作憑據。 \end{tabularx} \\ \\ \relax
5:6 & \begin{tabularx}{0.7\textwidth}{X} 所以，我們總是勇敢的，並且知道，只要我們住在這身體內就是離開了主。 \end{tabularx} \\ \\ \relax
5:7 & \begin{tabularx}{0.7\textwidth}{X} 因為我們行事為人是憑著信心，不是憑著眼見。 \end{tabularx} \\ \\ \relax
5:8 & \begin{tabularx}{0.7\textwidth}{X} 我們勇敢，更情願離開身體，與主同住。 \end{tabularx} \\ \\ \relax
5:9 & \begin{tabularx}{0.7\textwidth}{X} 所以，無論是住在身內或住在身外，我們都立了志向要得主的喜悅。 \end{tabularx} \\ \\ \relax
5:10 & \begin{tabularx}{0.7\textwidth}{X} 因為我們眾人必須站在基督審判臺前受審，為使各人按著本身所行的，或善或惡受報。 \end{tabularx} \\ \\
[1ex]
\hline
\hline
\end{longtable}


\newpage

\section{第75屆~港九培靈會~焦源濂~第一分題:有福之人--第一講}
\label{sec:399}
\textbf{蒙福的源頭}
\newline
\newline
連結: \href{https://www.hkbibleconference.org/session-message/view/399}{\texttt{https://www.hkbibleconference.org/session-message/view/399}}
\newline
\newline
\hyperref[sec:438]{< < < PREV SERMON < < <}
~
\hyperref[sec:toc]{[返目錄]}
~
\hyperref[sec:400]{> > > NEXT SERMON > > >}
\newline
\newline
詩篇 1:1-1:6
\newline
\begin{longtable}{cl}
\hline
\hline
章節 & 經文 (和合本修訂版)\\
\hline
1:1 & \begin{tabularx}{0.7\textwidth}{X} 不從惡人的計謀，不站罪人的道路，不坐傲慢人的座位， \end{tabularx} \\ \\ \relax
1:2 & \begin{tabularx}{0.7\textwidth}{X} 惟喜愛耶和華的律法，晝夜思想他的律法；這人便為有福！ \end{tabularx} \\ \\ \relax
1:3 & \begin{tabularx}{0.7\textwidth}{X} 他要像一棵樹栽在溪水旁，按時候結果子，葉子也不枯乾。凡他所做的盡都順利。 \end{tabularx} \\ \\ \relax
1:4 & \begin{tabularx}{0.7\textwidth}{X} 惡人並不是這樣，卻像糠秕被風吹散。 \end{tabularx} \\ \\ \relax
1:5 & \begin{tabularx}{0.7\textwidth}{X} 因此，當審判的時候惡人必站立不住，罪人在義人的會眾中也是如此。 \end{tabularx} \\ \\ \relax
1:6 & \begin{tabularx}{0.7\textwidth}{X} 因為耶和華知道義人的道路，惡人的道路卻必滅亡。 \end{tabularx} \\ \\
[1ex]
\hline
\hline
\end{longtable}


\newpage

\section{第75屆~港九培靈會~焦源濂~第一分題:有福之人--第二講:蒙福的途徑}
\label{sec:400}
\textbf{第一分題:有福之人--第二講:蒙福的途徑}
\newline
\newline
連結: \href{https://www.hkbibleconference.org/session-message/view/400}{\texttt{https://www.hkbibleconference.org/session-message/view/400}}
\newline
\newline
\hyperref[sec:399]{< < < PREV SERMON < < <}
~
\hyperref[sec:toc]{[返目錄]}
~
\hyperref[sec:401]{> > > NEXT SERMON > > >}
\newline
\newline
詩篇 1:1-1:6
\newline
\begin{longtable}{cl}
\hline
\hline
章節 & 經文 (和合本修訂版)\\
\hline
1:1 & \begin{tabularx}{0.7\textwidth}{X} 不從惡人的計謀，不站罪人的道路，不坐傲慢人的座位， \end{tabularx} \\ \\ \relax
1:2 & \begin{tabularx}{0.7\textwidth}{X} 惟喜愛耶和華的律法，晝夜思想他的律法；這人便為有福！ \end{tabularx} \\ \\ \relax
1:3 & \begin{tabularx}{0.7\textwidth}{X} 他要像一棵樹栽在溪水旁，按時候結果子，葉子也不枯乾。凡他所做的盡都順利。 \end{tabularx} \\ \\ \relax
1:4 & \begin{tabularx}{0.7\textwidth}{X} 惡人並不是這樣，卻像糠秕被風吹散。 \end{tabularx} \\ \\ \relax
1:5 & \begin{tabularx}{0.7\textwidth}{X} 因此，當審判的時候惡人必站立不住，罪人在義人的會眾中也是如此。 \end{tabularx} \\ \\ \relax
1:6 & \begin{tabularx}{0.7\textwidth}{X} 因為耶和華知道義人的道路，惡人的道路卻必滅亡。 \end{tabularx} \\ \\
[1ex]
\hline
\hline
\end{longtable}


\newpage

\section{第75屆~港九培靈會~焦源濂~第一分題:有福之人--第三講}
\label{sec:401}
\textbf{蒙福的見證}
\newline
\newline
連結: \href{https://www.hkbibleconference.org/session-message/view/401}{\texttt{https://www.hkbibleconference.org/session-message/view/401}}
\newline
\newline
\hyperref[sec:400]{< < < PREV SERMON < < <}
~
\hyperref[sec:toc]{[返目錄]}
~
\hyperref[sec:403]{> > > NEXT SERMON > > >}
\newline
\newline
詩篇 1:1-1:6
\newline
\begin{longtable}{cl}
\hline
\hline
章節 & 經文 (和合本修訂版)\\
\hline
1:1 & \begin{tabularx}{0.7\textwidth}{X} 不從惡人的計謀，不站罪人的道路，不坐傲慢人的座位， \end{tabularx} \\ \\ \relax
1:2 & \begin{tabularx}{0.7\textwidth}{X} 惟喜愛耶和華的律法，晝夜思想他的律法；這人便為有福！ \end{tabularx} \\ \\ \relax
1:3 & \begin{tabularx}{0.7\textwidth}{X} 他要像一棵樹栽在溪水旁，按時候結果子，葉子也不枯乾。凡他所做的盡都順利。 \end{tabularx} \\ \\ \relax
1:4 & \begin{tabularx}{0.7\textwidth}{X} 惡人並不是這樣，卻像糠秕被風吹散。 \end{tabularx} \\ \\ \relax
1:5 & \begin{tabularx}{0.7\textwidth}{X} 因此，當審判的時候惡人必站立不住，罪人在義人的會眾中也是如此。 \end{tabularx} \\ \\ \relax
1:6 & \begin{tabularx}{0.7\textwidth}{X} 因為耶和華知道義人的道路，惡人的道路卻必滅亡。 \end{tabularx} \\ \\
[1ex]
\hline
\hline
\end{longtable}


\newpage

\section{第75屆~港九培靈會~焦源濂~第一分題:有福之人--第四講}
\label{sec:403}
\textbf{蒙福的為人}
\newline
\newline
連結: \href{https://www.hkbibleconference.org/session-message/view/403}{\texttt{https://www.hkbibleconference.org/session-message/view/403}}
\newline
\newline
\hyperref[sec:401]{< < < PREV SERMON < < <}
~
\hyperref[sec:toc]{[返目錄]}
~
\hyperref[sec:405]{> > > NEXT SERMON > > >}
\newline
\newline
詩篇 1:1-1:6
\newline
\begin{longtable}{cl}
\hline
\hline
章節 & 經文 (和合本修訂版)\\
\hline
1:1 & \begin{tabularx}{0.7\textwidth}{X} 不從惡人的計謀，不站罪人的道路，不坐傲慢人的座位， \end{tabularx} \\ \\ \relax
1:2 & \begin{tabularx}{0.7\textwidth}{X} 惟喜愛耶和華的律法，晝夜思想他的律法；這人便為有福！ \end{tabularx} \\ \\ \relax
1:3 & \begin{tabularx}{0.7\textwidth}{X} 他要像一棵樹栽在溪水旁，按時候結果子，葉子也不枯乾。凡他所做的盡都順利。 \end{tabularx} \\ \\ \relax
1:4 & \begin{tabularx}{0.7\textwidth}{X} 惡人並不是這樣，卻像糠秕被風吹散。 \end{tabularx} \\ \\ \relax
1:5 & \begin{tabularx}{0.7\textwidth}{X} 因此，當審判的時候惡人必站立不住，罪人在義人的會眾中也是如此。 \end{tabularx} \\ \\ \relax
1:6 & \begin{tabularx}{0.7\textwidth}{X} 因為耶和華知道義人的道路，惡人的道路卻必滅亡。 \end{tabularx} \\ \\
[1ex]
\hline
\hline
\end{longtable}


\newpage

\section{第75屆~港九培靈會~焦源濂~第一分題:有福之人--第五講}
\label{sec:405}
\textbf{蒙福的保証}
\newline
\newline
連結: \href{https://www.hkbibleconference.org/session-message/view/405}{\texttt{https://www.hkbibleconference.org/session-message/view/405}}
\newline
\newline
\hyperref[sec:403]{< < < PREV SERMON < < <}
~
\hyperref[sec:toc]{[返目錄]}
~
\hyperref[sec:408]{> > > NEXT SERMON > > >}
\newline
\newline
詩篇 1:1-1:6
\newline
\begin{longtable}{cl}
\hline
\hline
章節 & 經文 (和合本修訂版)\\
\hline
1:1 & \begin{tabularx}{0.7\textwidth}{X} 不從惡人的計謀，不站罪人的道路，不坐傲慢人的座位， \end{tabularx} \\ \\ \relax
1:2 & \begin{tabularx}{0.7\textwidth}{X} 惟喜愛耶和華的律法，晝夜思想他的律法；這人便為有福！ \end{tabularx} \\ \\ \relax
1:3 & \begin{tabularx}{0.7\textwidth}{X} 他要像一棵樹栽在溪水旁，按時候結果子，葉子也不枯乾。凡他所做的盡都順利。 \end{tabularx} \\ \\ \relax
1:4 & \begin{tabularx}{0.7\textwidth}{X} 惡人並不是這樣，卻像糠秕被風吹散。 \end{tabularx} \\ \\ \relax
1:5 & \begin{tabularx}{0.7\textwidth}{X} 因此，當審判的時候惡人必站立不住，罪人在義人的會眾中也是如此。 \end{tabularx} \\ \\ \relax
1:6 & \begin{tabularx}{0.7\textwidth}{X} 因為耶和華知道義人的道路，惡人的道路卻必滅亡。 \end{tabularx} \\ \\
[1ex]
\hline
\hline
\end{longtable}


\newpage

\section{第75屆~港九培靈會~焦源濂~第二分題:詩的人生--第六講}
\label{sec:408}
\textbf{人生之詩的秘訣}
\newline
\newline
連結: \href{https://www.hkbibleconference.org/session-message/view/408}{\texttt{https://www.hkbibleconference.org/session-message/view/408}}
\newline
\newline
\hyperref[sec:405]{< < < PREV SERMON < < <}
~
\hyperref[sec:toc]{[返目錄]}
~
\hyperref[sec:410]{> > > NEXT SERMON > > >}
\newline
\newline
詩篇 23:1-23:6
\newline
\begin{longtable}{cl}
\hline
\hline
章節 & 經文 (和合本修訂版)\\
\hline
23:1 & \begin{tabularx}{0.7\textwidth}{X} 耶和華是我的牧者，我必不致缺乏。 \end{tabularx} \\ \\ \relax
23:2 & \begin{tabularx}{0.7\textwidth}{X} 他使我躺臥在青草地上，領我在可安歇的水邊。 \end{tabularx} \\ \\ \relax
23:3 & \begin{tabularx}{0.7\textwidth}{X} 他使我的靈魂甦醒，為自己的名引導我走義路。 \end{tabularx} \\ \\ \relax
23:4 & \begin{tabularx}{0.7\textwidth}{X} 我雖然行過死蔭的幽谷，也不怕遭害，因為你與我同在；你的杖、你的竿，都安慰我。 \end{tabularx} \\ \\ \relax
23:5 & \begin{tabularx}{0.7\textwidth}{X} 在我敵人面前，你為我擺設筵席；你用油膏了我的頭，使我的福杯滿溢。 \end{tabularx} \\ \\ \relax
23:6 & \begin{tabularx}{0.7\textwidth}{X} 我一生一世必有恩惠慈愛隨著我；我且要住在耶和華的殿中，直到永遠。 \end{tabularx} \\ \\
[1ex]
\hline
\hline
\end{longtable}


\newpage

\section{第75屆~港九培靈會~焦源濂~第二分題:詩的人生--第七講}
\label{sec:410}
\textbf{無憂無慮的人生}
\newline
\newline
連結: \href{https://www.hkbibleconference.org/session-message/view/410}{\texttt{https://www.hkbibleconference.org/session-message/view/410}}
\newline
\newline
\hyperref[sec:408]{< < < PREV SERMON < < <}
~
\hyperref[sec:toc]{[返目錄]}
~
\hyperref[sec:414]{> > > NEXT SERMON > > >}
\newline
\newline
詩篇 23:1-23:6
\newline
\begin{longtable}{cl}
\hline
\hline
章節 & 經文 (和合本修訂版)\\
\hline
23:1 & \begin{tabularx}{0.7\textwidth}{X} 耶和華是我的牧者，我必不致缺乏。 \end{tabularx} \\ \\ \relax
23:2 & \begin{tabularx}{0.7\textwidth}{X} 他使我躺臥在青草地上，領我在可安歇的水邊。 \end{tabularx} \\ \\ \relax
23:3 & \begin{tabularx}{0.7\textwidth}{X} 他使我的靈魂甦醒，為自己的名引導我走義路。 \end{tabularx} \\ \\ \relax
23:4 & \begin{tabularx}{0.7\textwidth}{X} 我雖然行過死蔭的幽谷，也不怕遭害，因為你與我同在；你的杖、你的竿，都安慰我。 \end{tabularx} \\ \\ \relax
23:5 & \begin{tabularx}{0.7\textwidth}{X} 在我敵人面前，你為我擺設筵席；你用油膏了我的頭，使我的福杯滿溢。 \end{tabularx} \\ \\ \relax
23:6 & \begin{tabularx}{0.7\textwidth}{X} 我一生一世必有恩惠慈愛隨著我；我且要住在耶和華的殿中，直到永遠。 \end{tabularx} \\ \\
[1ex]
\hline
\hline
\end{longtable}


\newpage

\section{第75屆~港九培靈會~焦源濂~第二分題:詩的人生--第八講}
\label{sec:414}
\textbf{目標正確的人生}
\newline
\newline
連結: \href{https://www.hkbibleconference.org/session-message/view/414}{\texttt{https://www.hkbibleconference.org/session-message/view/414}}
\newline
\newline
\hyperref[sec:410]{< < < PREV SERMON < < <}
~
\hyperref[sec:toc]{[返目錄]}
~
\hyperref[sec:417]{> > > NEXT SERMON > > >}
\newline
\newline
詩篇 23:1-23:6
\newline
\begin{longtable}{cl}
\hline
\hline
章節 & 經文 (和合本修訂版)\\
\hline
23:1 & \begin{tabularx}{0.7\textwidth}{X} 耶和華是我的牧者，我必不致缺乏。 \end{tabularx} \\ \\ \relax
23:2 & \begin{tabularx}{0.7\textwidth}{X} 他使我躺臥在青草地上，領我在可安歇的水邊。 \end{tabularx} \\ \\ \relax
23:3 & \begin{tabularx}{0.7\textwidth}{X} 他使我的靈魂甦醒，為自己的名引導我走義路。 \end{tabularx} \\ \\ \relax
23:4 & \begin{tabularx}{0.7\textwidth}{X} 我雖然行過死蔭的幽谷，也不怕遭害，因為你與我同在；你的杖、你的竿，都安慰我。 \end{tabularx} \\ \\ \relax
23:5 & \begin{tabularx}{0.7\textwidth}{X} 在我敵人面前，你為我擺設筵席；你用油膏了我的頭，使我的福杯滿溢。 \end{tabularx} \\ \\ \relax
23:6 & \begin{tabularx}{0.7\textwidth}{X} 我一生一世必有恩惠慈愛隨著我；我且要住在耶和華的殿中，直到永遠。 \end{tabularx} \\ \\
[1ex]
\hline
\hline
\end{longtable}


\newpage

\section{第75屆~港九培靈會~焦源濂~第二分題:詩的人生--第九講}
\label{sec:417}
\textbf{不畏艱苦的人生}
\newline
\newline
連結: \href{https://www.hkbibleconference.org/session-message/view/417}{\texttt{https://www.hkbibleconference.org/session-message/view/417}}
\newline
\newline
\hyperref[sec:414]{< < < PREV SERMON < < <}
~
\hyperref[sec:toc]{[返目錄]}
~
\hyperref[sec:418]{> > > NEXT SERMON > > >}
\newline
\newline
詩篇 23:1-23:6
\newline
\begin{longtable}{cl}
\hline
\hline
章節 & 經文 (和合本修訂版)\\
\hline
23:1 & \begin{tabularx}{0.7\textwidth}{X} 耶和華是我的牧者，我必不致缺乏。 \end{tabularx} \\ \\ \relax
23:2 & \begin{tabularx}{0.7\textwidth}{X} 他使我躺臥在青草地上，領我在可安歇的水邊。 \end{tabularx} \\ \\ \relax
23:3 & \begin{tabularx}{0.7\textwidth}{X} 他使我的靈魂甦醒，為自己的名引導我走義路。 \end{tabularx} \\ \\ \relax
23:4 & \begin{tabularx}{0.7\textwidth}{X} 我雖然行過死蔭的幽谷，也不怕遭害，因為你與我同在；你的杖、你的竿，都安慰我。 \end{tabularx} \\ \\ \relax
23:5 & \begin{tabularx}{0.7\textwidth}{X} 在我敵人面前，你為我擺設筵席；你用油膏了我的頭，使我的福杯滿溢。 \end{tabularx} \\ \\ \relax
23:6 & \begin{tabularx}{0.7\textwidth}{X} 我一生一世必有恩惠慈愛隨著我；我且要住在耶和華的殿中，直到永遠。 \end{tabularx} \\ \\
[1ex]
\hline
\hline
\end{longtable}


\newpage

\section{第75屆~港九培靈會~焦源濂~第二分題:詩的人生--第十講}
\label{sec:418}
\textbf{經歷美好的人生}
\newline
\newline
連結: \href{https://www.hkbibleconference.org/session-message/view/418}{\texttt{https://www.hkbibleconference.org/session-message/view/418}}
\newline
\newline
\hyperref[sec:417]{< < < PREV SERMON < < <}
~
\hyperref[sec:toc]{[返目錄]}
~
\hyperref[sec:383]{> > > NEXT SERMON > > >}
\newline
\newline
詩篇 23:1-23:6
\newline
\begin{longtable}{cl}
\hline
\hline
章節 & 經文 (和合本修訂版)\\
\hline
23:1 & \begin{tabularx}{0.7\textwidth}{X} 耶和華是我的牧者，我必不致缺乏。 \end{tabularx} \\ \\ \relax
23:2 & \begin{tabularx}{0.7\textwidth}{X} 他使我躺臥在青草地上，領我在可安歇的水邊。 \end{tabularx} \\ \\ \relax
23:3 & \begin{tabularx}{0.7\textwidth}{X} 他使我的靈魂甦醒，為自己的名引導我走義路。 \end{tabularx} \\ \\ \relax
23:4 & \begin{tabularx}{0.7\textwidth}{X} 我雖然行過死蔭的幽谷，也不怕遭害，因為你與我同在；你的杖、你的竿，都安慰我。 \end{tabularx} \\ \\ \relax
23:5 & \begin{tabularx}{0.7\textwidth}{X} 在我敵人面前，你為我擺設筵席；你用油膏了我的頭，使我的福杯滿溢。 \end{tabularx} \\ \\ \relax
23:6 & \begin{tabularx}{0.7\textwidth}{X} 我一生一世必有恩惠慈愛隨著我；我且要住在耶和華的殿中，直到永遠。 \end{tabularx} \\ \\
[1ex]
\hline
\hline
\end{longtable}


\newpage

\section{第75屆~港九培靈會~蔡元雲~第一講}
\label{sec:383}
\textbf{同進曠野:廿一世紀的試探}
\newline
\newline
連結: \href{https://www.hkbibleconference.org/session-message/view/383}{\texttt{https://www.hkbibleconference.org/session-message/view/383}}
\newline
\newline
\hyperref[sec:418]{< < < PREV SERMON < < <}
~
\hyperref[sec:toc]{[返目錄]}
~
\hyperref[sec:385]{> > > NEXT SERMON > > >}
\newline
\newline
路加福音 4:1-4:44
\newline
\begin{longtable}{cl}
\hline
\hline
章節 & 經文 (和合本修訂版)\\
\hline
4:1 & \begin{tabularx}{0.7\textwidth}{X} 耶穌滿有聖靈，從約旦河回來，聖靈把他引到曠野， \end{tabularx} \\ \\ \relax
4:2 & \begin{tabularx}{0.7\textwidth}{X} 四十天受魔鬼的試探。在那些日子，他沒有吃甚麼，日子滿了，他餓了。 \end{tabularx} \\ \\ \relax
4:3 & \begin{tabularx}{0.7\textwidth}{X} 魔鬼對他說：「你若是神的兒子，叫這塊石頭變成食物吧。」 \end{tabularx} \\ \\ \relax
4:4 & \begin{tabularx}{0.7\textwidth}{X} 耶穌回答說：「經上記著：『人活著，不是單靠食物。』」 \end{tabularx} \\ \\ \relax
4:5 & \begin{tabularx}{0.7\textwidth}{X} 魔鬼又領他上了高山，霎時間把天下萬國都指給他看， \end{tabularx} \\ \\ \relax
4:6 & \begin{tabularx}{0.7\textwidth}{X} 對他說：「這一切權柄和榮華我都要給你，因為這原是交給我的，我願意給誰就給誰。 \end{tabularx} \\ \\ \relax
4:7 & \begin{tabularx}{0.7\textwidth}{X} 你若在我面前下拜，這一切都歸你。」 \end{tabularx} \\ \\ \relax
4:8 & \begin{tabularx}{0.7\textwidth}{X} 耶穌回答他說：「經上記著：『要拜主—你的神，惟獨事奉他。』」 \end{tabularx} \\ \\ \relax
4:9 & \begin{tabularx}{0.7\textwidth}{X} 魔鬼又領他到耶路撒冷去，叫他站在聖殿頂上，對他說：「你若是神的兒子，從這裡跳下去！ \end{tabularx} \\ \\ \relax
4:10 & \begin{tabularx}{0.7\textwidth}{X} 因為經上記著：『主要為你命令他的使者保護你； \end{tabularx} \\ \\ \relax
4:11 & \begin{tabularx}{0.7\textwidth}{X} 他們要用手托住你，免得你的腳碰在石頭上。』」 \end{tabularx} \\ \\ \relax
4:12 & \begin{tabularx}{0.7\textwidth}{X} 耶穌回答他說：「經上說：『不可試探主—你的神。』」 \end{tabularx} \\ \\ \relax
4:13 & \begin{tabularx}{0.7\textwidth}{X} 魔鬼用完了各樣的試探，就離開耶穌，再等時機。 \end{tabularx} \\ \\ \relax
4:14 & \begin{tabularx}{0.7\textwidth}{X} 耶穌帶著聖靈的能力回到加利利，他的名聲傳遍了四方。 \end{tabularx} \\ \\ \relax
4:15 & \begin{tabularx}{0.7\textwidth}{X} 他在各會堂裡教導人，眾人都稱讚他。 \end{tabularx} \\ \\ \relax
4:16 & \begin{tabularx}{0.7\textwidth}{X} 耶穌來到拿撒勒，就是他長大的地方。在安息日，照他素常的規矩進了會堂，站起來要念聖經。 \end{tabularx} \\ \\ \relax
4:17 & \begin{tabularx}{0.7\textwidth}{X} 有人把以賽亞先知的書交給他，他就打開，找到一處寫著： \end{tabularx} \\ \\ \relax
4:18 & \begin{tabularx}{0.7\textwidth}{X} 「主的靈在我身上，因為他用膏膏我，叫我傳福音給貧窮的人；差遣我宣告：被擄的得釋放，失明的得看見，受壓迫的得自由， \end{tabularx} \\ \\ \relax
4:19 & \begin{tabularx}{0.7\textwidth}{X} 宣告神悅納人的禧年。」 \end{tabularx} \\ \\ \relax
4:20 & \begin{tabularx}{0.7\textwidth}{X} 於是他把書捲起來，交還給管理人，就坐下。會堂裡的人都定睛看他。 \end{tabularx} \\ \\ \relax
4:21 & \begin{tabularx}{0.7\textwidth}{X} 耶穌對他們說：「你們聽見的這段經文，今天已經應驗了。」 \end{tabularx} \\ \\ \relax
4:22 & \begin{tabularx}{0.7\textwidth}{X} 眾人都稱讚他，並對他口中所出的恩言感到驚訝；他們說：「這不是約瑟的兒子嗎？」 \end{tabularx} \\ \\ \relax
4:23 & \begin{tabularx}{0.7\textwidth}{X} 耶穌對他們說：「你們一定會用這俗語向我說：『醫生，你醫治自己吧！我們聽見你在迦百農所做的事，也該在你自己的家鄉做吧。』」 \end{tabularx} \\ \\ \relax
4:24 & \begin{tabularx}{0.7\textwidth}{X} 他又說：「我實在告訴你們，沒有先知在自己家鄉被人接納的。 \end{tabularx} \\ \\ \relax
4:25 & \begin{tabularx}{0.7\textwidth}{X} 我對你們說實話，在以利亞的時候，天閉塞了三年六個月，遍地有大饑荒，那時，以色列中有許多寡婦， \end{tabularx} \\ \\ \relax
4:26 & \begin{tabularx}{0.7\textwidth}{X} 以利亞並沒有奉差往她們中任何一個人那裡去，只奉差往西頓的撒勒法一個寡婦那裡去。 \end{tabularx} \\ \\ \relax
4:27 & \begin{tabularx}{0.7\textwidth}{X} 在以利沙先知的時候，以色列中有許多痲瘋病人，但除了敘利亞的乃縵，沒有一個得潔淨的。」 \end{tabularx} \\ \\ \relax
4:28 & \begin{tabularx}{0.7\textwidth}{X} 會堂裡的人聽見這些話，都怒氣填胸， \end{tabularx} \\ \\ \relax
4:29 & \begin{tabularx}{0.7\textwidth}{X} 就起來趕他出城。他們的城造在山上；他們帶他到山崖，要把他推下去。 \end{tabularx} \\ \\ \relax
4:30 & \begin{tabularx}{0.7\textwidth}{X} 他卻從他們中間穿過去，走了。 \end{tabularx} \\ \\ \relax
4:31 & \begin{tabularx}{0.7\textwidth}{X} 耶穌下到迦百農，就是加利利的一座城，在安息日教導眾人。 \end{tabularx} \\ \\ \relax
4:32 & \begin{tabularx}{0.7\textwidth}{X} 他們對他的教導感到很驚奇，因為他的話裡有權柄。 \end{tabularx} \\ \\ \relax
4:33 & \begin{tabularx}{0.7\textwidth}{X} 在會堂裡有一個人，被污鬼的靈附著，大聲喊叫說： \end{tabularx} \\ \\ \relax
4:34 & \begin{tabularx}{0.7\textwidth}{X} 「唉！拿撒勒人耶穌，你為甚麼干擾我們？你來消滅我們嗎？我知道你是誰，你是神的聖者。」 \end{tabularx} \\ \\ \relax
4:35 & \begin{tabularx}{0.7\textwidth}{X} 耶穌斥責他說：「不要作聲，從這人身上出來吧！」鬼把那人摔倒在眾人中間，就出來了，卻沒有傷害他。 \end{tabularx} \\ \\ \relax
4:36 & \begin{tabularx}{0.7\textwidth}{X} 眾人都驚訝，彼此對問：「這是甚麼道理呢？因為他用權柄能力命令污靈，污靈就出來。」 \end{tabularx} \\ \\ \relax
4:37 & \begin{tabularx}{0.7\textwidth}{X} 於是耶穌的名聲傳遍了周圍各地。 \end{tabularx} \\ \\ \relax
4:38 & \begin{tabularx}{0.7\textwidth}{X} 耶穌出了會堂，進了西門的家。西門的岳母在發高燒，有些人為她求耶穌。 \end{tabularx} \\ \\ \relax
4:39 & \begin{tabularx}{0.7\textwidth}{X} 耶穌站在她旁邊，斥責那高燒，燒就退了。她立刻起來服事他們。 \end{tabularx} \\ \\ \relax
4:40 & \begin{tabularx}{0.7\textwidth}{X} 日落的時候，凡有病人的，不論害甚麼病，都帶到耶穌那裡。耶穌給他們每一個人按手，治好他們。 \end{tabularx} \\ \\ \relax
4:41 & \begin{tabularx}{0.7\textwidth}{X} 又有鬼從好些人身上出來，喊著說：「你是神的兒子！」耶穌斥責他們，不許他們說話，因為他們知道他是基督。 \end{tabularx} \\ \\ \relax
4:42 & \begin{tabularx}{0.7\textwidth}{X} 天亮的時候，耶穌出來，走到荒野的地方。眾人去找他，到了他那裡，要留住他，不讓他離開他們。 \end{tabularx} \\ \\ \relax
4:43 & \begin{tabularx}{0.7\textwidth}{X} 但耶穌對他們說：「我也必須在別的城傳神國的福音，因我奉差原是為此。」 \end{tabularx} \\ \\ \relax
4:44 & \begin{tabularx}{0.7\textwidth}{X} 於是耶穌在猶太的各會堂傳道。 \end{tabularx} \\ \\
[1ex]
\hline
\hline
\end{longtable}


\newpage

\section{第75屆~港九培靈會~蔡元雲~第二講}
\label{sec:385}
\textbf{同登高山:廿一世紀的文化救贖}
\newline
\newline
連結: \href{https://www.hkbibleconference.org/session-message/view/385}{\texttt{https://www.hkbibleconference.org/session-message/view/385}}
\newline
\newline
\hyperref[sec:383]{< < < PREV SERMON < < <}
~
\hyperref[sec:toc]{[返目錄]}
~
\hyperref[sec:387]{> > > NEXT SERMON > > >}
\newline
\newline
路加福音 6:1-6:49
\newline
\begin{longtable}{cl}
\hline
\hline
章節 & 經文 (和合本修訂版)\\
\hline
6:1 & \begin{tabularx}{0.7\textwidth}{X} 有一個安息日，耶穌從麥田經過。他的門徒摘了麥穗，用手搓著吃。 \end{tabularx} \\ \\ \relax
6:2 & \begin{tabularx}{0.7\textwidth}{X} 有幾個法利賽人說：「你們為甚麼做安息日不合法的事呢？」 \end{tabularx} \\ \\ \relax
6:3 & \begin{tabularx}{0.7\textwidth}{X} 耶穌回答他們：「大衛和跟從他的人飢餓時所做的事，你們沒有念過嗎？ \end{tabularx} \\ \\ \relax
6:4 & \begin{tabularx}{0.7\textwidth}{X} 他怎麼進了神的居所，拿供餅吃，又給跟從的人吃呢？這餅惟獨祭司可以吃，別人都不可以吃。」 \end{tabularx} \\ \\ \relax
6:5 & \begin{tabularx}{0.7\textwidth}{X} 他又對他們說：「人子是安息日的主。」 \end{tabularx} \\ \\ \relax
6:6 & \begin{tabularx}{0.7\textwidth}{X} 又有一個安息日，耶穌進了會堂教導人，在那裡有一個人，他的右手萎縮了。 \end{tabularx} \\ \\ \relax
6:7 & \begin{tabularx}{0.7\textwidth}{X} 文士和法利賽人窺探耶穌會不會在安息日治病，為要找把柄告他。 \end{tabularx} \\ \\ \relax
6:8 & \begin{tabularx}{0.7\textwidth}{X} 耶穌卻知道他們的意念，就對那萎縮了手的人說：「起來，站在當中！」那人就起來，站著。 \end{tabularx} \\ \\ \relax
6:9 & \begin{tabularx}{0.7\textwidth}{X} 耶穌對他們說：「我問你們，在安息日行善行惡，救命害命，哪樣是合法的呢？」 \end{tabularx} \\ \\ \relax
6:10 & \begin{tabularx}{0.7\textwidth}{X} 他就環視眾人，對那人說：「伸出手來！」他照著做，他的手就復原了。 \end{tabularx} \\ \\ \relax
6:11 & \begin{tabularx}{0.7\textwidth}{X} 他們怒氣填胸，彼此商議怎樣對付耶穌。 \end{tabularx} \\ \\ \relax
6:12 & \begin{tabularx}{0.7\textwidth}{X} 在那些日子，耶穌出去，上山祈禱，整夜向神禱告。 \end{tabularx} \\ \\ \relax
6:13 & \begin{tabularx}{0.7\textwidth}{X} 到了天亮，他叫門徒來，就從他們中間挑選十二個人，稱他們為使徒。 \end{tabularx} \\ \\ \relax
6:14 & \begin{tabularx}{0.7\textwidth}{X} 這十二個人有西門（耶穌又給他起名叫彼得），還有他弟弟安得烈，又有雅各和約翰，腓力和巴多羅買， \end{tabularx} \\ \\ \relax
6:15 & \begin{tabularx}{0.7\textwidth}{X} 馬太和多馬，亞勒腓的兒子雅各和激進黨的西門， \end{tabularx} \\ \\ \relax
6:16 & \begin{tabularx}{0.7\textwidth}{X} 雅各的兒子猶大和後來成為出賣者的加略人猶大。 \end{tabularx} \\ \\ \relax
6:17 & \begin{tabularx}{0.7\textwidth}{X} 耶穌和他們下了山，站在一塊平地上；在一起的有許多門徒，又有許多百姓從全猶太和耶路撒冷，並推羅、西頓的海邊來， \end{tabularx} \\ \\ \relax
6:18 & \begin{tabularx}{0.7\textwidth}{X} 都要聽他講道，又希望耶穌醫治他們的病；還有被污靈纏磨的，也得了醫治。 \end{tabularx} \\ \\ \relax
6:19 & \begin{tabularx}{0.7\textwidth}{X} 眾人都想要摸他，因為有能力從他身上發出來，治好了他們。 \end{tabularx} \\ \\ \relax
6:20 & \begin{tabularx}{0.7\textwidth}{X} 耶穌舉目看著門徒，說：「貧窮的人有福了！因為神的國是你們的。 \end{tabularx} \\ \\ \relax
6:21 & \begin{tabularx}{0.7\textwidth}{X} 現在飢餓的人有福了！因為你們將得飽足。現在哭泣的人有福了！因為你們將要歡笑。 \end{tabularx} \\ \\ \relax
6:22 & \begin{tabularx}{0.7\textwidth}{X} 人為人子的緣故憎恨你們，拒絕你們，辱罵你們，把你們當惡人除掉你們的名，你們就有福了！ \end{tabularx} \\ \\ \relax
6:23 & \begin{tabularx}{0.7\textwidth}{X} 在那日，你們要歡欣雀躍，因為你們在天上的賞賜是很多的；他們的祖宗也是這樣待先知的。 \end{tabularx} \\ \\ \relax
6:24 & \begin{tabularx}{0.7\textwidth}{X} 但你們富足的人有禍了！因為你們已經受過安慰。 \end{tabularx} \\ \\ \relax
6:25 & \begin{tabularx}{0.7\textwidth}{X} 你們現在飽足的人有禍了！因為你們將要飢餓。你們現在歡笑的人有禍了！因為你們將要哀慟哭泣。 \end{tabularx} \\ \\ \relax
6:26 & \begin{tabularx}{0.7\textwidth}{X} 人都說你們好的時候，你們有禍了！因為他們的祖宗也是這樣待假先知的。」 \end{tabularx} \\ \\ \relax
6:27 & \begin{tabularx}{0.7\textwidth}{X} 「可是我告訴你們這些聽的人，要愛你們的仇敵！要善待恨你們的人！ \end{tabularx} \\ \\ \relax
6:28 & \begin{tabularx}{0.7\textwidth}{X} 要祝福詛咒你們的人！要為凌辱你們的人禱告！ \end{tabularx} \\ \\ \relax
6:29 & \begin{tabularx}{0.7\textwidth}{X} 有人打你的臉，連另一邊也由他打。有人拿你的外衣，連內衣也由他拿去。 \end{tabularx} \\ \\ \relax
6:30 & \begin{tabularx}{0.7\textwidth}{X} 凡求你的，就給他；有人拿走你的東西，不要討回來。 \end{tabularx} \\ \\ \relax
6:31 & \begin{tabularx}{0.7\textwidth}{X} 「你們想要人怎樣待你們，你們也要怎樣待人。 \end{tabularx} \\ \\ \relax
6:32 & \begin{tabularx}{0.7\textwidth}{X} 你們若只愛那愛你們的人，有甚麼可感謝的呢？就是罪人也愛那愛他們的人。 \end{tabularx} \\ \\ \relax
6:33 & \begin{tabularx}{0.7\textwidth}{X} 你們若善待那善待你們的人，有甚麼可感謝的呢？就是罪人也是這樣做。 \end{tabularx} \\ \\ \relax
6:34 & \begin{tabularx}{0.7\textwidth}{X} 你們若借給人，希望從他收回，有甚麼可感謝的呢？就是罪人也借給罪人，再如數收回。 \end{tabularx} \\ \\ \relax
6:35 & \begin{tabularx}{0.7\textwidth}{X} 你們倒要愛仇敵，要善待他們，並要借給人不指望償還，你們的賞賜就很多了，你們必作至高者的兒子，因為他恩待那忘恩的和作惡的。 \end{tabularx} \\ \\ \relax
6:36 & \begin{tabularx}{0.7\textwidth}{X} 你們要仁慈，像你們的父是仁慈的。」 \end{tabularx} \\ \\ \relax
6:37 & \begin{tabularx}{0.7\textwidth}{X} 「你們不要評斷別人，就不被審判；你們不要定人的罪，就不被定罪；你們要饒恕人，就必蒙饒恕。 \end{tabularx} \\ \\ \relax
6:38 & \begin{tabularx}{0.7\textwidth}{X} 你們要給人，就必有給你們的，並且用十足的升斗，連搖帶按，上尖下流地倒在你們懷裡；因為你們用甚麼量器量給人，也必用甚麼量器量給你們。」 \end{tabularx} \\ \\ \relax
6:39 & \begin{tabularx}{0.7\textwidth}{X} 耶穌又用比喻對他們說：「瞎子豈能領瞎子，兩個人不是都要掉在坑裡嗎？ \end{tabularx} \\ \\ \relax
6:40 & \begin{tabularx}{0.7\textwidth}{X} 學生不高過老師，凡學成了的會和老師一樣。 \end{tabularx} \\ \\ \relax
6:41 & \begin{tabularx}{0.7\textwidth}{X} 為甚麼看見你弟兄眼中有刺，卻不想自己眼中有梁木呢？ \end{tabularx} \\ \\ \relax
6:42 & \begin{tabularx}{0.7\textwidth}{X} 你不見自己眼中有梁木，怎能對你弟兄說：『讓我去掉你眼中的刺』呢？你這假冒為善的人！先去掉自己眼中的梁木，然後才能看得清楚，好去掉你弟兄眼中的刺。」 \end{tabularx} \\ \\ \relax
6:43 & \begin{tabularx}{0.7\textwidth}{X} 「沒有好樹結壞果子，也沒有壞樹結好果子。 \end{tabularx} \\ \\ \relax
6:44 & \begin{tabularx}{0.7\textwidth}{X} 每一種樹木可以從其果子看出來。人不是從荊棘上摘無花果的，也不是從蒺藜裡摘葡萄的。 \end{tabularx} \\ \\ \relax
6:45 & \begin{tabularx}{0.7\textwidth}{X} 善人從他心裡所存的善發出善來，惡人從他所存的惡發出惡來；因為心裡所充滿的，口裡就說出來。」 \end{tabularx} \\ \\ \relax
6:46 & \begin{tabularx}{0.7\textwidth}{X} 「你們為甚麼稱呼我『主啊，主啊』，卻不照我的話做呢？ \end{tabularx} \\ \\ \relax
6:47 & \begin{tabularx}{0.7\textwidth}{X} 凡到我這裡來，聽了我的話又去做的，我要告訴你們他像甚麼人： \end{tabularx} \\ \\ \relax
6:48 & \begin{tabularx}{0.7\textwidth}{X} 他像一個人蓋房子，把地挖深，將根基立在磐石上，到發大水的時候，水沖那房子，房子總不動搖，因為蓋造得好。 \end{tabularx} \\ \\ \relax
6:49 & \begin{tabularx}{0.7\textwidth}{X} 但聽了不去做的，就像一個人在土地上蓋房子，沒有根基，水一沖，立刻倒塌了，並且那房子損壞得很厲害。」 \end{tabularx} \\ \\
[1ex]
\hline
\hline
\end{longtable}


\newpage

\section{第75屆~港九培靈會~蔡元雲~第三講}
\label{sec:387}
\textbf{同背十架:廿一世紀的苦難}
\newline
\newline
連結: \href{https://www.hkbibleconference.org/session-message/view/387}{\texttt{https://www.hkbibleconference.org/session-message/view/387}}
\newline
\newline
\hyperref[sec:385]{< < < PREV SERMON < < <}
~
\hyperref[sec:toc]{[返目錄]}
~
\hyperref[sec:390]{> > > NEXT SERMON > > >}
\newline
\newline
路加福音 9:1-9:62
\newline
\begin{longtable}{cl}
\hline
\hline
章節 & 經文 (和合本修訂版)\\
\hline
9:1 & \begin{tabularx}{0.7\textwidth}{X} 耶穌叫齊了十二使徒，給他們能力和權柄制伏一切的鬼，醫治疾病， \end{tabularx} \\ \\ \relax
9:2 & \begin{tabularx}{0.7\textwidth}{X} 又差遣他們宣講神的國，醫治病人， \end{tabularx} \\ \\ \relax
9:3 & \begin{tabularx}{0.7\textwidth}{X} 對他們說：「途中甚麼都不要帶；不要帶手杖和行囊，不要帶食物和銀錢，也不要帶兩件內衣。 \end{tabularx} \\ \\ \relax
9:4 & \begin{tabularx}{0.7\textwidth}{X} 你們無論進哪一家，就住在哪裡，也從那裡離開。 \end{tabularx} \\ \\ \relax
9:5 & \begin{tabularx}{0.7\textwidth}{X} 凡不接待你們的，你們離開那城的時候，要跺掉你們腳上的塵土，證明他們的不是。」 \end{tabularx} \\ \\ \relax
9:6 & \begin{tabularx}{0.7\textwidth}{X} 於是使徒出去，走遍各鄉傳福音，到處治病。 \end{tabularx} \\ \\ \relax
9:7 & \begin{tabularx}{0.7\textwidth}{X} 希律分封王聽見耶穌所做的一切事，就困惑起來，因為有人說：「約翰從死人中復活了。」 \end{tabularx} \\ \\ \relax
9:8 & \begin{tabularx}{0.7\textwidth}{X} 又有人說：「以利亞顯現了。」還有人說：「古時的一個先知又活了。」 \end{tabularx} \\ \\ \relax
9:9 & \begin{tabularx}{0.7\textwidth}{X} 希律說：「約翰我已經斬了，但這是甚麼人？關於他，我竟聽到這樣的事！」於是希律想要見他。 \end{tabularx} \\ \\ \relax
9:10 & \begin{tabularx}{0.7\textwidth}{X} 使徒們回來，把所做的事告訴耶穌，耶穌就私下帶他們離開那裡，往一座叫伯賽大的城去。 \end{tabularx} \\ \\ \relax
9:11 & \begin{tabularx}{0.7\textwidth}{X} 眾人知道了，就跟著他去；耶穌接待他們，對他們講論神國的事，治好那些需要醫治的人。 \end{tabularx} \\ \\ \relax
9:12 & \begin{tabularx}{0.7\textwidth}{X} 太陽快要下山，十二使徒進前來對他說：「請叫眾人散去，他們好往四面村莊鄉鎮裡去借宿和找吃的，因為我們這裡地方偏僻。」 \end{tabularx} \\ \\ \relax
9:13 & \begin{tabularx}{0.7\textwidth}{X} 耶穌對他們說：「你們給他們吃吧！」他們說：「我們不過有五個餅、兩條魚，若不去為這許多人買食物就不夠。」 \end{tabularx} \\ \\ \relax
9:14 & \begin{tabularx}{0.7\textwidth}{X} 那時，男人約有五千。耶穌對門徒說：「叫他們分組坐下，每組大約五十個人。」 \end{tabularx} \\ \\ \relax
9:15 & \begin{tabularx}{0.7\textwidth}{X} 門徒就這樣做了，叫眾人都坐下。 \end{tabularx} \\ \\ \relax
9:16 & \begin{tabularx}{0.7\textwidth}{X} 耶穌拿著這五個餅和兩條魚，望著天祝福，擘開，遞給門徒，擺在眾人面前。 \end{tabularx} \\ \\ \relax
9:17 & \begin{tabularx}{0.7\textwidth}{X} 所有的人都吃，並且吃飽了。他們把剩下的碎屑收拾起來，裝滿了十二個籃子。 \end{tabularx} \\ \\ \relax
9:18 & \begin{tabularx}{0.7\textwidth}{X} 耶穌獨自禱告的時候，門徒也同他在那裡。耶穌問他們：「眾人說我是誰？」 \end{tabularx} \\ \\ \relax
9:19 & \begin{tabularx}{0.7\textwidth}{X} 他們回答：「是施洗的約翰；有人說是以利亞；還有人說是古時的一個先知又活了。」 \end{tabularx} \\ \\ \relax
9:20 & \begin{tabularx}{0.7\textwidth}{X} 耶穌問他們：「你們說我是誰？」彼得回答：「是神所立的基督。」 \end{tabularx} \\ \\ \relax
9:21 & \begin{tabularx}{0.7\textwidth}{X} 耶穌切切吩咐他們，命令他們不可把這事告訴任何人； \end{tabularx} \\ \\ \relax
9:22 & \begin{tabularx}{0.7\textwidth}{X} 又說：「人子必須受許多的苦，被長老、祭司長和文士棄絕，並且被殺，第三天復活。」 \end{tabularx} \\ \\ \relax
9:23 & \begin{tabularx}{0.7\textwidth}{X} 耶穌又對眾人說：「若有人要跟從我，就當捨己，天天背起自己的十字架來跟從我。 \end{tabularx} \\ \\ \relax
9:24 & \begin{tabularx}{0.7\textwidth}{X} 因為凡要救自己生命的，必喪失生命；凡為我喪失生命的，他必救自己的生命。 \end{tabularx} \\ \\ \relax
9:25 & \begin{tabularx}{0.7\textwidth}{X} 人就是賺得全世界，卻喪失了自己，或賠上自己，有甚麼益處呢？ \end{tabularx} \\ \\ \relax
9:26 & \begin{tabularx}{0.7\textwidth}{X} 凡把我和我的道當作可恥的，人子在自己的榮耀裡，和天父與聖天使的榮耀裡來臨的時候，也要把那人當作可恥的。 \end{tabularx} \\ \\ \relax
9:27 & \begin{tabularx}{0.7\textwidth}{X} 我實在告訴你們，站在這裡的，有人在沒經歷死亡以前，必定看見神的國。」 \end{tabularx} \\ \\ \relax
9:28 & \begin{tabularx}{0.7\textwidth}{X} 說了這些話以後約有八天，耶穌帶著彼得、約翰、雅各上山去禱告。 \end{tabularx} \\ \\ \relax
9:29 & \begin{tabularx}{0.7\textwidth}{X} 正禱告的時候，他的面貌改變了，衣服潔白放光。 \end{tabularx} \\ \\ \relax
9:30 & \begin{tabularx}{0.7\textwidth}{X} 忽然有摩西和以利亞兩個人同耶穌說話； \end{tabularx} \\ \\ \relax
9:31 & \begin{tabularx}{0.7\textwidth}{X} 他們在榮光裡顯現，談論耶穌去世的事，就是他在耶路撒冷將要完成的事。 \end{tabularx} \\ \\ \relax
9:32 & \begin{tabularx}{0.7\textwidth}{X} 彼得和他的同伴都打盹，但一清醒，就看見耶穌的榮光和與他一起站著的那兩個人。 \end{tabularx} \\ \\ \relax
9:33 & \begin{tabularx}{0.7\textwidth}{X} 二人正要和耶穌分離的時候，彼得對耶穌說：「老師，我們在這裡真好！我們來搭三座棚，一座為你，一座為摩西，一座為以利亞。」他卻不知道自己在說些甚麼。 \end{tabularx} \\ \\ \relax
9:34 & \begin{tabularx}{0.7\textwidth}{X} 說這些話的時候，有一朵雲彩來遮蓋他們；他們一進入雲彩就很懼怕。 \end{tabularx} \\ \\ \relax
9:35 & \begin{tabularx}{0.7\textwidth}{X} 有聲音從雲彩裡出來，說：「這是我的兒子，我所揀選的。你們要聽從他！」 \end{tabularx} \\ \\ \relax
9:36 & \begin{tabularx}{0.7\textwidth}{X} 聲音停止後，只見耶穌獨自一人。當那些日子，門徒保持沉默，不把所看見的事告訴任何人。 \end{tabularx} \\ \\ \relax
9:37 & \begin{tabularx}{0.7\textwidth}{X} 第二天，他們下了山，有一大群人來迎見耶穌。 \end{tabularx} \\ \\ \relax
9:38 & \begin{tabularx}{0.7\textwidth}{X} 其中有一人喊著說：「老師！求你看看我的兒子，因為他是我的獨子。 \end{tabularx} \\ \\ \relax
9:39 & \begin{tabularx}{0.7\textwidth}{X} 他被靈拿住就突然喊叫，那靈又使他抽風，口吐白沫，並且重重地傷害他，不輕易放過他。 \end{tabularx} \\ \\ \relax
9:40 & \begin{tabularx}{0.7\textwidth}{X} 我求過你的門徒把那靈趕出去，他們卻不能。」 \end{tabularx} \\ \\ \relax
9:41 & \begin{tabularx}{0.7\textwidth}{X} 耶穌回答：「唉！這又不信又悖謬的世代啊，我和你們在一起，忍耐你們，要到幾時呢？把你的兒子帶到這裡來！」 \end{tabularx} \\ \\ \relax
9:42 & \begin{tabularx}{0.7\textwidth}{X} 他正來的時候，那鬼把他摔倒，使他重重地抽風。耶穌斥責那污靈，把孩子治好了，交給他父親。 \end{tabularx} \\ \\ \relax
9:43 & \begin{tabularx}{0.7\textwidth}{X} 眾人都詫異神的大能。 \end{tabularx} \\ \\ \relax
9:44 & \begin{tabularx}{0.7\textwidth}{X} 「你們要把這些話聽進去，因為人子將要被交在人手裡。」 \end{tabularx} \\ \\ \relax
9:45 & \begin{tabularx}{0.7\textwidth}{X} 門徒卻不明白這話，其中的意思對他們隱藏著，使他們不能明白，他們也不敢問這話的意思。 \end{tabularx} \\ \\ \relax
9:46 & \begin{tabularx}{0.7\textwidth}{X} 門徒互相議論，他們中間誰最大。 \end{tabularx} \\ \\ \relax
9:47 & \begin{tabularx}{0.7\textwidth}{X} 耶穌看出他們心中的議論，就領一個小孩子來，叫他站在自己旁邊， \end{tabularx} \\ \\ \relax
9:48 & \begin{tabularx}{0.7\textwidth}{X} 對他們說：「凡為我的名接納這小孩子的，就是接納我；凡接納我的，就是接納那差我來的。你們中間最小的，他就是最大的。」 \end{tabularx} \\ \\ \relax
9:49 & \begin{tabularx}{0.7\textwidth}{X} 約翰回應說：「老師，我們看見一個人奉你的名趕鬼，我們就阻止他，因為他不與我們一同跟從你。」 \end{tabularx} \\ \\ \relax
9:50 & \begin{tabularx}{0.7\textwidth}{X} 耶穌對他說：「不要阻止他，因為不抵擋你們的，就是幫助你們的。」 \end{tabularx} \\ \\ \relax
9:51 & \begin{tabularx}{0.7\textwidth}{X} 耶穌被接上升的日子將到，他決定面向耶路撒冷走去。 \end{tabularx} \\ \\ \relax
9:52 & \begin{tabularx}{0.7\textwidth}{X} 他打發使者在他前頭走；他們進了撒瑪利亞的一個村莊，要為他作準備。 \end{tabularx} \\ \\ \relax
9:53 & \begin{tabularx}{0.7\textwidth}{X} 那裡的人不接待他，因為他面向著耶路撒冷去。 \end{tabularx} \\ \\ \relax
9:54 & \begin{tabularx}{0.7\textwidth}{X} 他的門徒雅各和約翰看見了，就說：「主啊！你要我們吩咐火從天上降下來，燒滅他們嗎？」 \end{tabularx} \\ \\ \relax
9:55 & \begin{tabularx}{0.7\textwidth}{X} 耶穌轉身責備兩個門徒。 \end{tabularx} \\ \\ \relax
9:56 & \begin{tabularx}{0.7\textwidth}{X} 於是他們就往別的村莊去了。 \end{tabularx} \\ \\ \relax
9:57 & \begin{tabularx}{0.7\textwidth}{X} 他們在路上走的時候，有一個人對耶穌說：「你無論往哪裡去，我都要跟從你。」 \end{tabularx} \\ \\ \relax
9:58 & \begin{tabularx}{0.7\textwidth}{X} 耶穌對他說：「狐狸有洞，天空的飛鳥有窩，人子卻沒有枕頭的地方。」 \end{tabularx} \\ \\ \relax
9:59 & \begin{tabularx}{0.7\textwidth}{X} 他又對另一個人說：「來跟從我！」那人說：「主啊，容許我先回去埋葬我的父親。」 \end{tabularx} \\ \\ \relax
9:60 & \begin{tabularx}{0.7\textwidth}{X} 耶穌對他說：「讓死人埋葬他們的死人，你只管去傳講神的國。」 \end{tabularx} \\ \\ \relax
9:61 & \begin{tabularx}{0.7\textwidth}{X} 又有一人說：「主啊，我要跟從你，但容許我先去辭別我家裡的人。」 \end{tabularx} \\ \\ \relax
9:62 & \begin{tabularx}{0.7\textwidth}{X} 耶穌對他說：「手扶著犁向後看的人，不配進神的國。」 \end{tabularx} \\ \\
[1ex]
\hline
\hline
\end{longtable}


\newpage

\section{第75屆~港九培靈會~蔡元雲~第四講}
\label{sec:390}
\textbf{同闖禁區:廿一世紀的種族衝突}
\newline
\newline
連結: \href{https://www.hkbibleconference.org/session-message/view/390}{\texttt{https://www.hkbibleconference.org/session-message/view/390}}
\newline
\newline
\hyperref[sec:387]{< < < PREV SERMON < < <}
~
\hyperref[sec:toc]{[返目錄]}
~
\hyperref[sec:393]{> > > NEXT SERMON > > >}
\newline
\newline
路加福音 10:1-10:42
\newline
\begin{longtable}{cl}
\hline
\hline
章節 & 經文 (和合本修訂版)\\
\hline
10:1 & \begin{tabularx}{0.7\textwidth}{X} 這些事以後，主另外指定七十二個人，差遣他們兩個兩個地在他前面，往自己所要到的各城各地去。 \end{tabularx} \\ \\ \relax
10:2 & \begin{tabularx}{0.7\textwidth}{X} 他對他們說：「要收的莊稼多，做工的人少。所以，你們要求莊稼的主差遣做工的人出去收他的莊稼。 \end{tabularx} \\ \\ \relax
10:3 & \begin{tabularx}{0.7\textwidth}{X} 你們去吧！看！我差你們出去，如同羔羊進入狼群。 \end{tabularx} \\ \\ \relax
10:4 & \begin{tabularx}{0.7\textwidth}{X} 不要帶錢囊，不要帶行囊，不要帶鞋子；在路上也不要向人問安。 \end{tabularx} \\ \\ \relax
10:5 & \begin{tabularx}{0.7\textwidth}{X} 無論進哪一家，先要說：『願這一家平安。』 \end{tabularx} \\ \\ \relax
10:6 & \begin{tabularx}{0.7\textwidth}{X} 那裡若有當得平安的人，你們所求的平安就必臨到那家，不然，將歸還你們。 \end{tabularx} \\ \\ \relax
10:7 & \begin{tabularx}{0.7\textwidth}{X} 你們要住在那家，吃喝他們所供給的，因為工人得工錢是應當的；不要從這家搬到那家。 \end{tabularx} \\ \\ \relax
10:8 & \begin{tabularx}{0.7\textwidth}{X} 無論進哪一城，人若接待你們，給你們擺上甚麼食物，你們就吃甚麼。 \end{tabularx} \\ \\ \relax
10:9 & \begin{tabularx}{0.7\textwidth}{X} 要醫治那城裡的病人，對他們說：『神的國臨近你們了。』 \end{tabularx} \\ \\ \relax
10:10 & \begin{tabularx}{0.7\textwidth}{X} 無論進哪一城，人若不接待你們，你們就到大街上去，說： \end{tabularx} \\ \\ \relax
10:11 & \begin{tabularx}{0.7\textwidth}{X} 『就是你們城裡的塵土黏在我們的腳上，我們也當著你們擦去。但是，你們該知道神的國臨近了。』 \end{tabularx} \\ \\ \relax
10:12 & \begin{tabularx}{0.7\textwidth}{X} 我告訴你們，在那日子，所多瑪所受的，比那城還容易受呢！」 \end{tabularx} \\ \\ \relax
10:13 & \begin{tabularx}{0.7\textwidth}{X} 「哥拉汛哪，你有禍了！伯賽大啊，你有禍了！因為在你們中間所行的異能，若行在推羅、西頓，他們早已披麻蒙灰，坐在地上悔改了。 \end{tabularx} \\ \\ \relax
10:14 & \begin{tabularx}{0.7\textwidth}{X} 在審判的時候，推羅和西頓所受的，比你們還容易受呢！ \end{tabularx} \\ \\ \relax
10:15 & \begin{tabularx}{0.7\textwidth}{X} 迦百農啊，你以為要被舉到天上嗎？你要被推下陰間！」 \end{tabularx} \\ \\ \relax
10:16 & \begin{tabularx}{0.7\textwidth}{X} 耶穌又對門徒說：「聽從你們的就是聽從我；棄絕你們的就是棄絕我；棄絕我的就是棄絕差遣我來的那位。」 \end{tabularx} \\ \\ \relax
10:17 & \begin{tabularx}{0.7\textwidth}{X} 那七十二個人歡歡喜喜地回來，說：「主啊，因你的名，就是鬼也服了我們。」 \end{tabularx} \\ \\ \relax
10:18 & \begin{tabularx}{0.7\textwidth}{X} 耶穌對他們說：「我看見撒但從天上墜落，像閃電一樣。 \end{tabularx} \\ \\ \relax
10:19 & \begin{tabularx}{0.7\textwidth}{X} 我已經給你們權柄可以踐踏蛇和蠍子，又勝過仇敵一切的能力，絕沒有甚麼能害你們。 \end{tabularx} \\ \\ \relax
10:20 & \begin{tabularx}{0.7\textwidth}{X} 然而，不要因靈服了你們就歡喜，而要因你們的名記錄在天上歡喜。」 \end{tabularx} \\ \\ \relax
10:21 & \begin{tabularx}{0.7\textwidth}{X} 正當那時，耶穌被聖靈感動而歡喜快樂，說：「父啊，天地的主，我感謝你！因為你把這些事向聰明智慧的人隱藏起來，而向嬰孩啟示出來。父啊，是的，因為你的美意本是如此。 \end{tabularx} \\ \\ \relax
10:22 & \begin{tabularx}{0.7\textwidth}{X} 一切都是我父交給我的。除了父，沒有人知道子是誰；除了子和子所願意啟示的人，沒有人知道父是誰。」 \end{tabularx} \\ \\ \relax
10:23 & \begin{tabularx}{0.7\textwidth}{X} 耶穌轉身私下對門徒說：「看見你們所看見的，那眼睛有福了。 \end{tabularx} \\ \\ \relax
10:24 & \begin{tabularx}{0.7\textwidth}{X} 我告訴你們，從前有許多先知和君王要看你們所看的，卻沒有看見，要聽你們所聽的，卻沒有聽見。」 \end{tabularx} \\ \\ \relax
10:25 & \begin{tabularx}{0.7\textwidth}{X} 有一個律法師起來試探耶穌，說：「老師！我該做甚麼才可以承受永生？」 \end{tabularx} \\ \\ \relax
10:26 & \begin{tabularx}{0.7\textwidth}{X} 耶穌對他說：「律法上寫的是甚麼？你是怎樣念的呢？」 \end{tabularx} \\ \\ \relax
10:27 & \begin{tabularx}{0.7\textwidth}{X} 他回答說：「你要盡心、盡性、盡力、盡意愛主—你的神，又要愛鄰如己。」 \end{tabularx} \\ \\ \relax
10:28 & \begin{tabularx}{0.7\textwidth}{X} 耶穌對他說：「你回答得正確，你這樣做就會得永生。」 \end{tabularx} \\ \\ \relax
10:29 & \begin{tabularx}{0.7\textwidth}{X} 那人要證明自己有理，就對耶穌說：「誰是我的鄰舍呢？」 \end{tabularx} \\ \\ \relax
10:30 & \begin{tabularx}{0.7\textwidth}{X} 耶穌回答：「有一個人從耶路撒冷下耶利哥去，落在強盜手中。他們剝去他的衣裳，把他打個半死，丟下他走了。 \end{tabularx} \\ \\ \relax
10:31 & \begin{tabularx}{0.7\textwidth}{X} 偶然有一個祭司從那條路下來，看見他就從另一邊過去了。 \end{tabularx} \\ \\ \relax
10:32 & \begin{tabularx}{0.7\textwidth}{X} 又有一個利未人來到那裡，看見他，也照樣從另一邊過去了。 \end{tabularx} \\ \\ \relax
10:33 & \begin{tabularx}{0.7\textwidth}{X} 可是，有一個撒瑪利亞人路過那裡，看見他就動了慈心， \end{tabularx} \\ \\ \relax
10:34 & \begin{tabularx}{0.7\textwidth}{X} 上前用油和酒倒在他的傷處，包裹好了，扶他騎上自己的牲口，帶他到旅店裡去，照應他。 \end{tabularx} \\ \\ \relax
10:35 & \begin{tabularx}{0.7\textwidth}{X} 第二天，他拿出兩個銀幣來，交給店主，說：『請你照應他，額外的費用，我回來時會還你。』 \end{tabularx} \\ \\ \relax
10:36 & \begin{tabularx}{0.7\textwidth}{X} 你想，這三個人哪一個是落在強盜手中那人的鄰舍呢？」 \end{tabularx} \\ \\ \relax
10:37 & \begin{tabularx}{0.7\textwidth}{X} 他說：「是憐憫他的。」耶穌對他說：「你去，照樣做吧！」 \end{tabularx} \\ \\ \relax
10:38 & \begin{tabularx}{0.7\textwidth}{X} 他們繼續前行，耶穌進了一個村莊。有一個女人，名叫馬大，接他到自己家裡。 \end{tabularx} \\ \\ \relax
10:39 & \begin{tabularx}{0.7\textwidth}{X} 她有一個妹妹，名叫馬利亞，在主的腳前坐著聽他的道。 \end{tabularx} \\ \\ \relax
10:40 & \begin{tabularx}{0.7\textwidth}{X} 馬大伺候的事多，心裡忙亂，進前來，說：「主啊，我的妹妹留下我一個人伺候，你不在意嗎？請吩咐她來幫助我。」 \end{tabularx} \\ \\ \relax
10:41 & \begin{tabularx}{0.7\textwidth}{X} 主回答說：「馬大，馬大，你為許多的事操心煩惱， \end{tabularx} \\ \\ \relax
10:42 & \begin{tabularx}{0.7\textwidth}{X} 但是不可少的只有一件。馬利亞已經選擇了那上好的福分，是沒有人能從她奪去的。」 \end{tabularx} \\ \\
[1ex]
\hline
\hline
\end{longtable}


\newpage

\section{第75屆~港九培靈會~蔡元雲~第五講}
\label{sec:393}
\textbf{同動同靜:廿一世紀的屬靈操練}
\newline
\newline
連結: \href{https://www.hkbibleconference.org/session-message/view/393}{\texttt{https://www.hkbibleconference.org/session-message/view/393}}
\newline
\newline
\hyperref[sec:390]{< < < PREV SERMON < < <}
~
\hyperref[sec:toc]{[返目錄]}
~
\hyperref[sec:394]{> > > NEXT SERMON > > >}
\newline
\newline
路加福音 10:1-11:54
\newline
\begin{longtable}{cl}
\hline
\hline
章節 & 經文 (和合本修訂版)\\
\hline
10:1 & \begin{tabularx}{0.7\textwidth}{X} 這些事以後，主另外指定七十二個人，差遣他們兩個兩個地在他前面，往自己所要到的各城各地去。 \end{tabularx} \\ \\ \relax
10:2 & \begin{tabularx}{0.7\textwidth}{X} 他對他們說：「要收的莊稼多，做工的人少。所以，你們要求莊稼的主差遣做工的人出去收他的莊稼。 \end{tabularx} \\ \\ \relax
10:3 & \begin{tabularx}{0.7\textwidth}{X} 你們去吧！看！我差你們出去，如同羔羊進入狼群。 \end{tabularx} \\ \\ \relax
10:4 & \begin{tabularx}{0.7\textwidth}{X} 不要帶錢囊，不要帶行囊，不要帶鞋子；在路上也不要向人問安。 \end{tabularx} \\ \\ \relax
10:5 & \begin{tabularx}{0.7\textwidth}{X} 無論進哪一家，先要說：『願這一家平安。』 \end{tabularx} \\ \\ \relax
10:6 & \begin{tabularx}{0.7\textwidth}{X} 那裡若有當得平安的人，你們所求的平安就必臨到那家，不然，將歸還你們。 \end{tabularx} \\ \\ \relax
10:7 & \begin{tabularx}{0.7\textwidth}{X} 你們要住在那家，吃喝他們所供給的，因為工人得工錢是應當的；不要從這家搬到那家。 \end{tabularx} \\ \\ \relax
10:8 & \begin{tabularx}{0.7\textwidth}{X} 無論進哪一城，人若接待你們，給你們擺上甚麼食物，你們就吃甚麼。 \end{tabularx} \\ \\ \relax
10:9 & \begin{tabularx}{0.7\textwidth}{X} 要醫治那城裡的病人，對他們說：『神的國臨近你們了。』 \end{tabularx} \\ \\ \relax
10:10 & \begin{tabularx}{0.7\textwidth}{X} 無論進哪一城，人若不接待你們，你們就到大街上去，說： \end{tabularx} \\ \\ \relax
10:11 & \begin{tabularx}{0.7\textwidth}{X} 『就是你們城裡的塵土黏在我們的腳上，我們也當著你們擦去。但是，你們該知道神的國臨近了。』 \end{tabularx} \\ \\ \relax
10:12 & \begin{tabularx}{0.7\textwidth}{X} 我告訴你們，在那日子，所多瑪所受的，比那城還容易受呢！」 \end{tabularx} \\ \\ \relax
10:13 & \begin{tabularx}{0.7\textwidth}{X} 「哥拉汛哪，你有禍了！伯賽大啊，你有禍了！因為在你們中間所行的異能，若行在推羅、西頓，他們早已披麻蒙灰，坐在地上悔改了。 \end{tabularx} \\ \\ \relax
10:14 & \begin{tabularx}{0.7\textwidth}{X} 在審判的時候，推羅和西頓所受的，比你們還容易受呢！ \end{tabularx} \\ \\ \relax
10:15 & \begin{tabularx}{0.7\textwidth}{X} 迦百農啊，你以為要被舉到天上嗎？你要被推下陰間！」 \end{tabularx} \\ \\ \relax
10:16 & \begin{tabularx}{0.7\textwidth}{X} 耶穌又對門徒說：「聽從你們的就是聽從我；棄絕你們的就是棄絕我；棄絕我的就是棄絕差遣我來的那位。」 \end{tabularx} \\ \\ \relax
10:17 & \begin{tabularx}{0.7\textwidth}{X} 那七十二個人歡歡喜喜地回來，說：「主啊，因你的名，就是鬼也服了我們。」 \end{tabularx} \\ \\ \relax
10:18 & \begin{tabularx}{0.7\textwidth}{X} 耶穌對他們說：「我看見撒但從天上墜落，像閃電一樣。 \end{tabularx} \\ \\ \relax
10:19 & \begin{tabularx}{0.7\textwidth}{X} 我已經給你們權柄可以踐踏蛇和蠍子，又勝過仇敵一切的能力，絕沒有甚麼能害你們。 \end{tabularx} \\ \\ \relax
10:20 & \begin{tabularx}{0.7\textwidth}{X} 然而，不要因靈服了你們就歡喜，而要因你們的名記錄在天上歡喜。」 \end{tabularx} \\ \\ \relax
10:21 & \begin{tabularx}{0.7\textwidth}{X} 正當那時，耶穌被聖靈感動而歡喜快樂，說：「父啊，天地的主，我感謝你！因為你把這些事向聰明智慧的人隱藏起來，而向嬰孩啟示出來。父啊，是的，因為你的美意本是如此。 \end{tabularx} \\ \\ \relax
10:22 & \begin{tabularx}{0.7\textwidth}{X} 一切都是我父交給我的。除了父，沒有人知道子是誰；除了子和子所願意啟示的人，沒有人知道父是誰。」 \end{tabularx} \\ \\ \relax
10:23 & \begin{tabularx}{0.7\textwidth}{X} 耶穌轉身私下對門徒說：「看見你們所看見的，那眼睛有福了。 \end{tabularx} \\ \\ \relax
10:24 & \begin{tabularx}{0.7\textwidth}{X} 我告訴你們，從前有許多先知和君王要看你們所看的，卻沒有看見，要聽你們所聽的，卻沒有聽見。」 \end{tabularx} \\ \\ \relax
10:25 & \begin{tabularx}{0.7\textwidth}{X} 有一個律法師起來試探耶穌，說：「老師！我該做甚麼才可以承受永生？」 \end{tabularx} \\ \\ \relax
10:26 & \begin{tabularx}{0.7\textwidth}{X} 耶穌對他說：「律法上寫的是甚麼？你是怎樣念的呢？」 \end{tabularx} \\ \\ \relax
10:27 & \begin{tabularx}{0.7\textwidth}{X} 他回答說：「你要盡心、盡性、盡力、盡意愛主—你的神，又要愛鄰如己。」 \end{tabularx} \\ \\ \relax
10:28 & \begin{tabularx}{0.7\textwidth}{X} 耶穌對他說：「你回答得正確，你這樣做就會得永生。」 \end{tabularx} \\ \\ \relax
10:29 & \begin{tabularx}{0.7\textwidth}{X} 那人要證明自己有理，就對耶穌說：「誰是我的鄰舍呢？」 \end{tabularx} \\ \\ \relax
10:30 & \begin{tabularx}{0.7\textwidth}{X} 耶穌回答：「有一個人從耶路撒冷下耶利哥去，落在強盜手中。他們剝去他的衣裳，把他打個半死，丟下他走了。 \end{tabularx} \\ \\ \relax
10:31 & \begin{tabularx}{0.7\textwidth}{X} 偶然有一個祭司從那條路下來，看見他就從另一邊過去了。 \end{tabularx} \\ \\ \relax
10:32 & \begin{tabularx}{0.7\textwidth}{X} 又有一個利未人來到那裡，看見他，也照樣從另一邊過去了。 \end{tabularx} \\ \\ \relax
10:33 & \begin{tabularx}{0.7\textwidth}{X} 可是，有一個撒瑪利亞人路過那裡，看見他就動了慈心， \end{tabularx} \\ \\ \relax
10:34 & \begin{tabularx}{0.7\textwidth}{X} 上前用油和酒倒在他的傷處，包裹好了，扶他騎上自己的牲口，帶他到旅店裡去，照應他。 \end{tabularx} \\ \\ \relax
10:35 & \begin{tabularx}{0.7\textwidth}{X} 第二天，他拿出兩個銀幣來，交給店主，說：『請你照應他，額外的費用，我回來時會還你。』 \end{tabularx} \\ \\ \relax
10:36 & \begin{tabularx}{0.7\textwidth}{X} 你想，這三個人哪一個是落在強盜手中那人的鄰舍呢？」 \end{tabularx} \\ \\ \relax
10:37 & \begin{tabularx}{0.7\textwidth}{X} 他說：「是憐憫他的。」耶穌對他說：「你去，照樣做吧！」 \end{tabularx} \\ \\ \relax
10:38 & \begin{tabularx}{0.7\textwidth}{X} 他們繼續前行，耶穌進了一個村莊。有一個女人，名叫馬大，接他到自己家裡。 \end{tabularx} \\ \\ \relax
10:39 & \begin{tabularx}{0.7\textwidth}{X} 她有一個妹妹，名叫馬利亞，在主的腳前坐著聽他的道。 \end{tabularx} \\ \\ \relax
10:40 & \begin{tabularx}{0.7\textwidth}{X} 馬大伺候的事多，心裡忙亂，進前來，說：「主啊，我的妹妹留下我一個人伺候，你不在意嗎？請吩咐她來幫助我。」 \end{tabularx} \\ \\ \relax
10:41 & \begin{tabularx}{0.7\textwidth}{X} 主回答說：「馬大，馬大，你為許多的事操心煩惱， \end{tabularx} \\ \\ \relax
10:42 & \begin{tabularx}{0.7\textwidth}{X} 但是不可少的只有一件。馬利亞已經選擇了那上好的福分，是沒有人能從她奪去的。」 \end{tabularx} \\ \\ \relax
11:1 & \begin{tabularx}{0.7\textwidth}{X} 耶穌在一個地方禱告。禱告完了，有個門徒對他說：「主啊，求你教導我們禱告，像約翰教導他的門徒一樣。」 \end{tabularx} \\ \\ \relax
11:2 & \begin{tabularx}{0.7\textwidth}{X} 耶穌對他們說：「你們禱告的時候，要說：『父啊，願人都尊你的名為聖；願你的國降臨； \end{tabularx} \\ \\ \relax
11:3 & \begin{tabularx}{0.7\textwidth}{X} 我們日用的飲食，天天賜給我們。 \end{tabularx} \\ \\ \relax
11:4 & \begin{tabularx}{0.7\textwidth}{X} 赦免我們的罪，因為我們也赦免凡虧欠我們的人。不叫我們陷入試探。』」 \end{tabularx} \\ \\ \relax
11:5 & \begin{tabularx}{0.7\textwidth}{X} 耶穌又對他們說：「你們中間誰有一個朋友半夜到他那裡去，對他說：『朋友！請借給我三個餅； \end{tabularx} \\ \\ \relax
11:6 & \begin{tabularx}{0.7\textwidth}{X} 因為我有一個朋友旅途中來到我這裡，我沒有東西招待他。』 \end{tabularx} \\ \\ \relax
11:7 & \begin{tabularx}{0.7\textwidth}{X} 那人在裡面回答：『不要打擾我，門已經關了，孩子們也同我在床上了，我不能起來給你。』 \end{tabularx} \\ \\ \relax
11:8 & \begin{tabularx}{0.7\textwidth}{X} 我告訴你們，雖不因他是朋友起來給他，也會因他不顧面子地直求，起來照他所需要的給他。 \end{tabularx} \\ \\ \relax
11:9 & \begin{tabularx}{0.7\textwidth}{X} 我又告訴你們，祈求，就給你們；尋找，就找到；叩門，就給你們開門。 \end{tabularx} \\ \\ \relax
11:10 & \begin{tabularx}{0.7\textwidth}{X} 因為凡祈求的，就得著；尋找的，就找到；叩門的，就給他開門。 \end{tabularx} \\ \\ \relax
11:11 & \begin{tabularx}{0.7\textwidth}{X} 你們中間作父親的，誰有兒子求魚，反拿蛇當魚給他呢？ \end{tabularx} \\ \\ \relax
11:12 & \begin{tabularx}{0.7\textwidth}{X} 求雞蛋，反給他蠍子呢？ \end{tabularx} \\ \\ \relax
11:13 & \begin{tabularx}{0.7\textwidth}{X} 你們雖然不好，尚且知道拿好東西給兒女，何況天父，他豈不更要把聖靈賜給求他的人嗎？」 \end{tabularx} \\ \\ \relax
11:14 & \begin{tabularx}{0.7\textwidth}{X} 耶穌趕出一個使人成為啞巴的鬼，鬼出去了，啞巴就說出話來；眾人都很驚訝。 \end{tabularx} \\ \\ \relax
11:15 & \begin{tabularx}{0.7\textwidth}{X} 其中卻有人說：「他是靠著鬼王別西卜趕鬼。」 \end{tabularx} \\ \\ \relax
11:16 & \begin{tabularx}{0.7\textwidth}{X} 又有人試探耶穌，要他顯個來自天上的神蹟。 \end{tabularx} \\ \\ \relax
11:17 & \begin{tabularx}{0.7\textwidth}{X} 他知道他們的意念，就對他們說：「一國自相紛爭，必定荒蕪；一家自相紛爭，就必敗落。 \end{tabularx} \\ \\ \relax
11:18 & \begin{tabularx}{0.7\textwidth}{X} 撒但若自相紛爭，他的國怎能立得住呢？因為你們說我是靠著別西卜趕鬼。 \end{tabularx} \\ \\ \relax
11:19 & \begin{tabularx}{0.7\textwidth}{X} 我若靠著別西卜趕鬼，你們的子弟趕鬼又靠著誰呢？這樣，他們要作你們的判官。 \end{tabularx} \\ \\ \relax
11:20 & \begin{tabularx}{0.7\textwidth}{X} 我若靠著神的能力趕鬼，那麼，神的國就已臨到你們了。 \end{tabularx} \\ \\ \relax
11:21 & \begin{tabularx}{0.7\textwidth}{X} 壯士全副武裝，看守自己的住宅，他所有的都很安全； \end{tabularx} \\ \\ \relax
11:22 & \begin{tabularx}{0.7\textwidth}{X} 但有一個比他更強的來攻擊他，並且戰勝了他，就奪去他所倚靠的盔甲兵器，又分了他的掠物。 \end{tabularx} \\ \\ \relax
11:23 & \begin{tabularx}{0.7\textwidth}{X} 不跟我一起的，就是反對我；不與我一起收聚的，就是在拆散。」 \end{tabularx} \\ \\ \relax
11:24 & \begin{tabularx}{0.7\textwidth}{X} 「污靈離了人身，走遍無水之地尋找安歇之處，卻找不到。就說：『我要回到我原來的屋裡去。』 \end{tabularx} \\ \\ \relax
11:25 & \begin{tabularx}{0.7\textwidth}{X} 他到了，看見裡面打掃乾淨，修飾好了， \end{tabularx} \\ \\ \relax
11:26 & \begin{tabularx}{0.7\textwidth}{X} 就去另帶了七個比自己更惡的靈來，都進去住在那裡。那人後來的景況比先前更壞了。」 \end{tabularx} \\ \\ \relax
11:27 & \begin{tabularx}{0.7\textwidth}{X} 耶穌正說這些話的時候，眾人中間有一個女人高聲對他說：「懷你胎乳養你的有福了！」 \end{tabularx} \\ \\ \relax
11:28 & \begin{tabularx}{0.7\textwidth}{X} 耶穌卻說：「更有福的是聽神的道而遵守的人！」 \end{tabularx} \\ \\ \relax
11:29 & \begin{tabularx}{0.7\textwidth}{X} 當眾人越來越擁擠的時候，耶穌說：「這世代是一個邪惡的世代。他們求看神蹟，除了約拿的神蹟以外，再沒有神蹟給他們看了。 \end{tabularx} \\ \\ \relax
11:30 & \begin{tabularx}{0.7\textwidth}{X} 約拿怎樣為尼尼微人成了神蹟，人子也要照樣為這世代的人成為神蹟。 \end{tabularx} \\ \\ \relax
11:31 & \begin{tabularx}{0.7\textwidth}{X} 在審判的時候，南方的女王要起來定這世代的人的罪，因為她從地極而來，要聽所羅門智慧的話。看哪，比所羅門更大的在這裡！ \end{tabularx} \\ \\ \relax
11:32 & \begin{tabularx}{0.7\textwidth}{X} 在審判的時候，尼尼微人要起來定這世代的罪，因為尼尼微人聽了約拿所傳的就悔改了。看哪，比約拿更大的在這裡！」 \end{tabularx} \\ \\ \relax
11:33 & \begin{tabularx}{0.7\textwidth}{X} 「沒有人點燈放在地窖裡，或是斗底下，總是放在燈臺上，讓進來的人看見亮光。 \end{tabularx} \\ \\ \relax
11:34 & \begin{tabularx}{0.7\textwidth}{X} 你的眼睛就是身體的燈。當你的眼睛明亮，全身就光明，當眼睛昏花，全身就黑暗。 \end{tabularx} \\ \\ \relax
11:35 & \begin{tabularx}{0.7\textwidth}{X} 所以，你要注意，免得你裡面的光暗了。 \end{tabularx} \\ \\ \relax
11:36 & \begin{tabularx}{0.7\textwidth}{X} 若是你全身光明，毫無黑暗，就必全然光明，如同燈的明光照亮你。」 \end{tabularx} \\ \\ \relax
11:37 & \begin{tabularx}{0.7\textwidth}{X} 耶穌正說話的時候，有一個法利賽人請他吃飯，耶穌就進去坐席。 \end{tabularx} \\ \\ \relax
11:38 & \begin{tabularx}{0.7\textwidth}{X} 這法利賽人看見耶穌飯前不先洗手就很詫異。 \end{tabularx} \\ \\ \relax
11:39 & \begin{tabularx}{0.7\textwidth}{X} 主對他說：「如今你們法利賽人洗淨杯盤的外面，你們裡面卻滿了貪婪和邪惡。 \end{tabularx} \\ \\ \relax
11:40 & \begin{tabularx}{0.7\textwidth}{X} 無知的人哪！造外面的，不也造了裡面嗎？ \end{tabularx} \\ \\ \relax
11:41 & \begin{tabularx}{0.7\textwidth}{X} 只要把杯盤裡面的施捨給人，對你們來說一切就都潔淨了。 \end{tabularx} \\ \\ \relax
11:42 & \begin{tabularx}{0.7\textwidth}{X} 「但是你們法利賽人有禍了！因為你們將薄荷、芸香，和各樣蔬菜獻上十分之一，疏忽了公義和愛神的事；這原是你們該做的—至於其他也不可忽略。 \end{tabularx} \\ \\ \relax
11:43 & \begin{tabularx}{0.7\textwidth}{X} 你們法利賽人有禍了！因為你們喜愛會堂裡的高位，又喜歡人們在街市上向你們問安。 \end{tabularx} \\ \\ \relax
11:44 & \begin{tabularx}{0.7\textwidth}{X} 你們有禍了！因為你們如同不顯露的墳墓，走在上面的人並不知道。」 \end{tabularx} \\ \\ \relax
11:45 & \begin{tabularx}{0.7\textwidth}{X} 律法師中有一個回答耶穌，說：「老師，你這樣說也把我們侮辱了。」 \end{tabularx} \\ \\ \relax
11:46 & \begin{tabularx}{0.7\textwidth}{X} 耶穌說：「你們律法師也有禍了！因為你們把難挑的擔子放在別人身上，自己卻不肯動一個指頭去減輕這些擔子。 \end{tabularx} \\ \\ \relax
11:47 & \begin{tabularx}{0.7\textwidth}{X} 你們有禍了！因為你們建造先知的墳墓，那些先知正是你們的祖宗所殺的。 \end{tabularx} \\ \\ \relax
11:48 & \begin{tabularx}{0.7\textwidth}{X} 可見你們祖宗所做的事，你們是證人，你們也贊同，因為他們殺了先知，你們建造先知的墳墓。 \end{tabularx} \\ \\ \relax
11:49 & \begin{tabularx}{0.7\textwidth}{X} 所以，神的智慧也曾說：『我要差遣先知和使徒到他們那裡去，有的他們要殘殺，有的他們要迫害』， \end{tabularx} \\ \\ \relax
11:50 & \begin{tabularx}{0.7\textwidth}{X} 為使創世以來所流眾先知的血的罪都歸在這世代的人身上， \end{tabularx} \\ \\ \relax
11:51 & \begin{tabularx}{0.7\textwidth}{X} 就是從亞伯的血起，直到被殺在祭壇和聖所中間的撒迦利亞的血為止。是的，我告訴你們，這都要向這世代的人追討。 \end{tabularx} \\ \\ \relax
11:52 & \begin{tabularx}{0.7\textwidth}{X} 你們律法師有禍了！因為你們把知識的鑰匙奪了去，自己不進去，要進去的人，你們也阻擋他們。」 \end{tabularx} \\ \\ \relax
11:53 & \begin{tabularx}{0.7\textwidth}{X} 耶穌從那裡出來，文士和法利賽人就開始極力地催逼他，盤問他許多事， \end{tabularx} \\ \\ \relax
11:54 & \begin{tabularx}{0.7\textwidth}{X} 伺機要抓他的話柄。 \end{tabularx} \\ \\
[1ex]
\hline
\hline
\end{longtable}


\newpage

\section{第75屆~港九培靈會~蔡元雲~第六講}
\label{sec:394}
\textbf{同尋迷羊:廿一世紀的浪子}
\newline
\newline
連結: \href{https://www.hkbibleconference.org/session-message/view/394}{\texttt{https://www.hkbibleconference.org/session-message/view/394}}
\newline
\newline
\hyperref[sec:393]{< < < PREV SERMON < < <}
~
\hyperref[sec:toc]{[返目錄]}
~
\hyperref[sec:395]{> > > NEXT SERMON > > >}
\newline
\newline
路加福音 15:1-15:32
\newline
\begin{longtable}{cl}
\hline
\hline
章節 & 經文 (和合本修訂版)\\
\hline
15:1 & \begin{tabularx}{0.7\textwidth}{X} 許多稅吏和罪人都挨近耶穌，要聽他講道。 \end{tabularx} \\ \\ \relax
15:2 & \begin{tabularx}{0.7\textwidth}{X} 法利賽人和文士私下議論說：「這個人接納罪人，又同他們吃飯。」 \end{tabularx} \\ \\ \relax
15:3 & \begin{tabularx}{0.7\textwidth}{X} 耶穌就用比喻對他們說： \end{tabularx} \\ \\ \relax
15:4 & \begin{tabularx}{0.7\textwidth}{X} 「你們中間誰有一百隻羊，失去其中的一隻，不把這九十九隻留在曠野，去找那失去的羊，直到找著呢？ \end{tabularx} \\ \\ \relax
15:5 & \begin{tabularx}{0.7\textwidth}{X} 找到了，他就歡歡喜喜地把羊扛在肩上。 \end{tabularx} \\ \\ \relax
15:6 & \begin{tabularx}{0.7\textwidth}{X} 他回到家裡，請朋友和鄰舍來，對他們說：『你們和我一同歡喜吧，我失去的羊已經找到了！』 \end{tabularx} \\ \\ \relax
15:7 & \begin{tabularx}{0.7\textwidth}{X} 我告訴你們，一個罪人悔改，在天上也要這樣為他歡喜，比為九十九個不用悔改的義人歡喜還大呢！」 \end{tabularx} \\ \\ \relax
15:8 & \begin{tabularx}{0.7\textwidth}{X} 「同樣，哪一個婦人有十塊錢，若失落一塊，不點上燈，打掃屋子，細細地找，直到找著呢？ \end{tabularx} \\ \\ \relax
15:9 & \begin{tabularx}{0.7\textwidth}{X} 找到了，她就請朋友和鄰舍來，對她們說：『你們和我一同歡喜吧，我失落的那塊錢已經找到了！』 \end{tabularx} \\ \\ \relax
15:10 & \begin{tabularx}{0.7\textwidth}{X} 我告訴你們，一個罪人悔改，神的使者也是這樣為他歡喜。」 \end{tabularx} \\ \\ \relax
15:11 & \begin{tabularx}{0.7\textwidth}{X} 耶穌又說：「一個人有兩個兒子。 \end{tabularx} \\ \\ \relax
15:12 & \begin{tabularx}{0.7\textwidth}{X} 小兒子對父親說：『父親，請你把我應得的家業分給我。』他父親就把財產分給他們。 \end{tabularx} \\ \\ \relax
15:13 & \begin{tabularx}{0.7\textwidth}{X} 過了不多幾天，小兒子把他一切所有的都收拾起來，往遠方去了。在那裡，他任意放蕩，浪費錢財。 \end{tabularx} \\ \\ \relax
15:14 & \begin{tabularx}{0.7\textwidth}{X} 他耗盡了一切所有的，又恰逢那地方有大饑荒，就窮困起來。 \end{tabularx} \\ \\ \relax
15:15 & \begin{tabularx}{0.7\textwidth}{X} 於是他去投靠當地的一個居民，那人打發他到田裡去放豬。 \end{tabularx} \\ \\ \relax
15:16 & \begin{tabularx}{0.7\textwidth}{X} 他恨不得拿豬所吃的豆莢充飢，也沒有人給他甚麼吃的。 \end{tabularx} \\ \\ \relax
15:17 & \begin{tabularx}{0.7\textwidth}{X} 他醒悟過來，就說：『我父親有多少雇工，糧食有餘，我倒在這裡餓死嗎？ \end{tabularx} \\ \\ \relax
15:18 & \begin{tabularx}{0.7\textwidth}{X} 我要起來，到我父親那裡去，對他說：父親！我得罪了天，又得罪了你， \end{tabularx} \\ \\ \relax
15:19 & \begin{tabularx}{0.7\textwidth}{X} 從今以後，我不配稱為你的兒子，把我當作一個雇工吧。』 \end{tabularx} \\ \\ \relax
15:20 & \begin{tabularx}{0.7\textwidth}{X} 於是他起來，往他父親那裡去。相離還遠，他父親看見，就動了慈心，跑去擁抱著他，連連親他。 \end{tabularx} \\ \\ \relax
15:21 & \begin{tabularx}{0.7\textwidth}{X} 兒子對他說：『父親！我得罪了天，又得罪了你，從今以後，我不配稱為你的兒子。』 \end{tabularx} \\ \\ \relax
15:22 & \begin{tabularx}{0.7\textwidth}{X} 父親卻吩咐僕人：『快把那上好的袍子拿出來給他穿，把戒指戴在他指頭上，把鞋穿在他腳上， \end{tabularx} \\ \\ \relax
15:23 & \begin{tabularx}{0.7\textwidth}{X} 把那肥牛犢牽來宰了，我們來吃喝慶祝； \end{tabularx} \\ \\ \relax
15:24 & \begin{tabularx}{0.7\textwidth}{X} 因為我這個兒子是死而復活，失而復得的。』他們就開始慶祝。 \end{tabularx} \\ \\ \relax
15:25 & \begin{tabularx}{0.7\textwidth}{X} 「那時，大兒子正在田裡。他回來，離家不遠時，聽見奏樂跳舞的聲音， \end{tabularx} \\ \\ \relax
15:26 & \begin{tabularx}{0.7\textwidth}{X} 就叫一個僮僕來，問是甚麼事。 \end{tabularx} \\ \\ \relax
15:27 & \begin{tabularx}{0.7\textwidth}{X} 僮僕對他說：『你弟弟回來了，你父親因為他無災無病地回來，把肥牛犢宰了。』 \end{tabularx} \\ \\ \relax
15:28 & \begin{tabularx}{0.7\textwidth}{X} 大兒子就生氣，不肯進去，他父親出來勸他。 \end{tabularx} \\ \\ \relax
15:29 & \begin{tabularx}{0.7\textwidth}{X} 他對父親說：『你看，我服侍你這麼多年，從來沒有違背過你的命令，而你從來沒有給我一隻小山羊，叫我和朋友們一同快樂。 \end{tabularx} \\ \\ \relax
15:30 & \begin{tabularx}{0.7\textwidth}{X} 但你這個兒子和娼妓吃光了你的財產，他一回來，你倒為他宰了肥牛犢。』 \end{tabularx} \\ \\ \relax
15:31 & \begin{tabularx}{0.7\textwidth}{X} 父親對他說：『兒啊！你常和我同在，我所有的一切都是你的； \end{tabularx} \\ \\ \relax
15:32 & \begin{tabularx}{0.7\textwidth}{X} 可是你這個弟弟是死而復活，失而復得的，所以我們理當歡喜慶祝。』」 \end{tabularx} \\ \\
[1ex]
\hline
\hline
\end{longtable}


\newpage

\section{第75屆~港九培靈會~蔡元雲~第七講}
\label{sec:395}
\textbf{同死同活:廿一世紀的能力危機}
\newline
\newline
連結: \href{https://www.hkbibleconference.org/session-message/view/395}{\texttt{https://www.hkbibleconference.org/session-message/view/395}}
\newline
\newline
\hyperref[sec:394]{< < < PREV SERMON < < <}
~
\hyperref[sec:toc]{[返目錄]}
~
\hyperref[sec:397]{> > > NEXT SERMON > > >}
\newline
\newline
路加福音 23:1-24:53
\newline
\begin{longtable}{cl}
\hline
\hline
章節 & 經文 (和合本修訂版)\\
\hline
23:1 & \begin{tabularx}{0.7\textwidth}{X} 眾人都起來，把耶穌解到彼拉多面前。 \end{tabularx} \\ \\ \relax
23:2 & \begin{tabularx}{0.7\textwidth}{X} 他們開始控告他說：「我們見這人煽惑我們的國民，禁止我們納稅給凱撒，並說自己是基督，是王。」 \end{tabularx} \\ \\ \relax
23:3 & \begin{tabularx}{0.7\textwidth}{X} 彼拉多問耶穌：「你是猶太人的王嗎？」耶穌回答：「是你說的。」 \end{tabularx} \\ \\ \relax
23:4 & \begin{tabularx}{0.7\textwidth}{X} 彼拉多對祭司長們和眾人說：「我查不出這人有甚麼罪來。」 \end{tabularx} \\ \\ \relax
23:5 & \begin{tabularx}{0.7\textwidth}{X} 但他們越發竭力地說：「他煽動百姓，在猶太全地傳道，從加利利起，直到這裡了。」 \end{tabularx} \\ \\ \relax
23:6 & \begin{tabularx}{0.7\textwidth}{X} 彼拉多一聽見，就問：「這人是加利利人嗎？」 \end{tabularx} \\ \\ \relax
23:7 & \begin{tabularx}{0.7\textwidth}{X} 既知道耶穌屬希律所管，彼拉多就把他送到希律那裡去。那時希律正在耶路撒冷。 \end{tabularx} \\ \\ \relax
23:8 & \begin{tabularx}{0.7\textwidth}{X} 希律看見耶穌就非常高興；因為聽見過他的事，早就想要見他，並且指望看他行些神蹟， \end{tabularx} \\ \\ \relax
23:9 & \begin{tabularx}{0.7\textwidth}{X} 於是問他許多的話，耶穌卻一言不答。 \end{tabularx} \\ \\ \relax
23:10 & \begin{tabularx}{0.7\textwidth}{X} 那些祭司長和文士都站著，竭力控告他。 \end{tabularx} \\ \\ \relax
23:11 & \begin{tabularx}{0.7\textwidth}{X} 希律和他的士兵就藐視耶穌，戲弄他，給他穿上華麗的衣服，把他送回彼拉多那裡去。 \end{tabularx} \\ \\ \relax
23:12 & \begin{tabularx}{0.7\textwidth}{X} 從前希律和彼拉多彼此有仇，在那一天竟成了朋友。 \end{tabularx} \\ \\ \relax
23:13 & \begin{tabularx}{0.7\textwidth}{X} 彼拉多傳齊了眾祭司長、官長和百姓， \end{tabularx} \\ \\ \relax
23:14 & \begin{tabularx}{0.7\textwidth}{X} 對他們說：「你們解這人到我這裡，說他是煽惑百姓的。看哪，我也曾在你們面前審問他，並沒有查出這人犯過你們控告他的任何罪； \end{tabularx} \\ \\ \relax
23:15 & \begin{tabularx}{0.7\textwidth}{X} 就是希律也是如此，所以把他送回來。可見他沒有做甚麼該死的事。 \end{tabularx} \\ \\ \relax
23:16 & \begin{tabularx}{0.7\textwidth}{X} 所以，我要責打他，把他釋放。」 \end{tabularx} \\ \\ \relax
23:17 & \begin{tabularx}{0.7\textwidth}{X} 每逢這節期，總督必須釋放一個囚犯給他們。 \end{tabularx} \\ \\ \relax
23:18 & \begin{tabularx}{0.7\textwidth}{X} 眾人卻一齊喊著說：「除掉這個人！釋放巴拉巴給我們！」 \end{tabularx} \\ \\ \relax
23:19 & \begin{tabularx}{0.7\textwidth}{X} 這巴拉巴是因在城裡作亂和殺人而下在監裡的。 \end{tabularx} \\ \\ \relax
23:20 & \begin{tabularx}{0.7\textwidth}{X} 彼拉多願意釋放耶穌，就再次向他們講話。 \end{tabularx} \\ \\ \relax
23:21 & \begin{tabularx}{0.7\textwidth}{X} 無奈他們喊著說：「把他釘十字架！把他釘十字架！」 \end{tabularx} \\ \\ \relax
23:22 & \begin{tabularx}{0.7\textwidth}{X} 彼拉多第三次對他們說：「為甚麼呢？這人做了甚麼惡事呢？我並沒有查出他有甚麼該死的罪來。所以，我要責打他，把他釋放。」 \end{tabularx} \\ \\ \relax
23:23 & \begin{tabularx}{0.7\textwidth}{X} 他們大聲催逼彼拉多，要求他把耶穌釘十字架；他們的聲音終於得勝。 \end{tabularx} \\ \\ \relax
23:24 & \begin{tabularx}{0.7\textwidth}{X} 彼拉多這才照他們的要求定案； \end{tabularx} \\ \\ \relax
23:25 & \begin{tabularx}{0.7\textwidth}{X} 又把他們所要求的那因作亂和殺人而下在監裡的人釋放了，而把耶穌交給他們，隨他們的意思處置。 \end{tabularx} \\ \\ \relax
23:26 & \begin{tabularx}{0.7\textwidth}{X} 他們把耶穌帶去的時候，有一個古利奈人西門從鄉下來，他們就拿住他，把十字架擱在他身上，叫他背著跟在耶穌後面。 \end{tabularx} \\ \\ \relax
23:27 & \begin{tabularx}{0.7\textwidth}{X} 有許多百姓跟隨耶穌，其中有好些婦女為他號咷痛哭。 \end{tabularx} \\ \\ \relax
23:28 & \begin{tabularx}{0.7\textwidth}{X} 耶穌轉身對她們說：「耶路撒冷的女子，不要為我哭，要為你們自己和你們的兒女哭。 \end{tabularx} \\ \\ \relax
23:29 & \begin{tabularx}{0.7\textwidth}{X} 因為日子將到，人要說：『不生育的、未曾懷孕的，和未曾哺乳孩子的有福了！』 \end{tabularx} \\ \\ \relax
23:30 & \begin{tabularx}{0.7\textwidth}{X} 那時，人要向大山說：『倒在我們身上！』向小山說：『遮蓋我們！』 \end{tabularx} \\ \\ \relax
23:31 & \begin{tabularx}{0.7\textwidth}{X} 他們若在樹木青綠的時候做這些事，那麼在枯乾的時候將會怎麼樣呢？」 \end{tabularx} \\ \\ \relax
23:32 & \begin{tabularx}{0.7\textwidth}{X} 另外有兩個犯人也被帶來和耶穌一同處死。 \end{tabularx} \\ \\ \relax
23:33 & \begin{tabularx}{0.7\textwidth}{X} 到了一個地方，名叫髑髏地，他們就在那裡把耶穌釘在十字架上，又釘了兩個犯人：一個在右邊，一個在左邊。 \end{tabularx} \\ \\ \relax
23:34 & \begin{tabularx}{0.7\textwidth}{X} 〔 這時，耶穌說：「父啊！赦免他們，因為他們所做的，他們不知道。」〕士兵就抽籤分他的衣服。 \end{tabularx} \\ \\ \relax
23:35 & \begin{tabularx}{0.7\textwidth}{X} 百姓站在那裡觀看。官長也嘲笑他，說：「他救了別人，他若是基督，是神所揀選的，救救他自己吧！」 \end{tabularx} \\ \\ \relax
23:36 & \begin{tabularx}{0.7\textwidth}{X} 士兵也戲弄他，上前拿醋送給他喝， \end{tabularx} \\ \\ \relax
23:37 & \begin{tabularx}{0.7\textwidth}{X} 說：「你若是猶太人的王，救救你自己吧！」 \end{tabularx} \\ \\ \relax
23:38 & \begin{tabularx}{0.7\textwidth}{X} 在耶穌上方有一個牌子寫著：「這是猶太人的王。」 \end{tabularx} \\ \\ \relax
23:39 & \begin{tabularx}{0.7\textwidth}{X} 同釘的犯人中有一個譏笑他，說：「你不是基督嗎？救救你自己和我們吧！」 \end{tabularx} \\ \\ \relax
23:40 & \begin{tabularx}{0.7\textwidth}{X} 另一個就應聲責備他，說：「你是一樣受刑的，還不怕神嗎？ \end{tabularx} \\ \\ \relax
23:41 & \begin{tabularx}{0.7\textwidth}{X} 我們是應得的，因為我們是自作自受，但這個人沒有做過一件不對的事。」 \end{tabularx} \\ \\ \relax
23:42 & \begin{tabularx}{0.7\textwidth}{X} 他對耶穌說：「耶穌啊，你進入你國的時候，求你記念我。」 \end{tabularx} \\ \\ \relax
23:43 & \begin{tabularx}{0.7\textwidth}{X} 耶穌對他說：「我實在告訴你，今日你要同我在樂園裡了。」 \end{tabularx} \\ \\ \relax
23:44 & \begin{tabularx}{0.7\textwidth}{X} 那時大約是正午，全地都黑暗了，直到下午三點鐘， \end{tabularx} \\ \\ \relax
23:45 & \begin{tabularx}{0.7\textwidth}{X} 太陽變黑了，殿的幔子從當中裂為兩半。 \end{tabularx} \\ \\ \relax
23:46 & \begin{tabularx}{0.7\textwidth}{X} 耶穌大聲喊著說：「父啊，我將我的靈交在你手裡！」他說了這話，氣就斷了。 \end{tabularx} \\ \\ \relax
23:47 & \begin{tabularx}{0.7\textwidth}{X} 百夫長看見所發生的事，就歸榮耀給神，說：「這人真是個義人！」 \end{tabularx} \\ \\ \relax
23:48 & \begin{tabularx}{0.7\textwidth}{X} 聚集觀看這事的眾人，見了所發生的事，都捶著胸回去了。 \end{tabularx} \\ \\ \relax
23:49 & \begin{tabularx}{0.7\textwidth}{X} 所有與耶穌熟悉的人，和從加利利跟著他來的婦女們，都遠遠地站著，看這些事。 \end{tabularx} \\ \\ \relax
23:50 & \begin{tabularx}{0.7\textwidth}{X} 有一個人名叫約瑟，是個議員，為人善良正直， \end{tabularx} \\ \\ \relax
23:51 & \begin{tabularx}{0.7\textwidth}{X} 卻沒有附從別人的所謀所為。他是猶太的亞利馬太城人，素常盼望著神的國。 \end{tabularx} \\ \\ \relax
23:52 & \begin{tabularx}{0.7\textwidth}{X} 這人去見彼拉多，請求要耶穌的身體。 \end{tabularx} \\ \\ \relax
23:53 & \begin{tabularx}{0.7\textwidth}{X} 他把耶穌的身體取下來，用細麻布裹好，安放在鑿巖而成的墳墓裡；那墳墓從來沒有葬過人。 \end{tabularx} \\ \\ \relax
23:54 & \begin{tabularx}{0.7\textwidth}{X} 那日是預備日，安息日快到了。 \end{tabularx} \\ \\ \relax
23:55 & \begin{tabularx}{0.7\textwidth}{X} 那些從加利利和耶穌同來的婦女跟在後面，看見了墳墓和他的身體怎樣安放。 \end{tabularx} \\ \\ \relax
23:56 & \begin{tabularx}{0.7\textwidth}{X} 她們就回去，預備了香料香膏。在安息日，她們遵照誡命安息了。 \end{tabularx} \\ \\ \relax
24:1 & \begin{tabularx}{0.7\textwidth}{X} 七日的第一日，黎明的時候，那些婦女帶著所預備的香料來到墳墓那裡， \end{tabularx} \\ \\ \relax
24:2 & \begin{tabularx}{0.7\textwidth}{X} 發現石頭已經從墳墓滾開了， \end{tabularx} \\ \\ \relax
24:3 & \begin{tabularx}{0.7\textwidth}{X} 她們就進去，只是不見主耶穌的身體。 \end{tabularx} \\ \\ \relax
24:4 & \begin{tabularx}{0.7\textwidth}{X} 正在為這事困惑的時候，忽然有兩個人站在旁邊，衣服放光。 \end{tabularx} \\ \\ \relax
24:5 & \begin{tabularx}{0.7\textwidth}{X} 婦女們非常害怕，就俯伏在地上。那兩個人對她們說：「為甚麼在死人中找活人呢？ \end{tabularx} \\ \\ \relax
24:6 & \begin{tabularx}{0.7\textwidth}{X} 他不在這裡，已經復活了。要記得他還在加利利的時候怎樣告訴你們的， \end{tabularx} \\ \\ \relax
24:7 & \begin{tabularx}{0.7\textwidth}{X} 他說：『人子必須被交在罪人手裡，釘在十字架上，第三天復活。』」 \end{tabularx} \\ \\ \relax
24:8 & \begin{tabularx}{0.7\textwidth}{X} 她們就想起耶穌的話來。 \end{tabularx} \\ \\ \relax
24:9 & \begin{tabularx}{0.7\textwidth}{X} 於是她們從墳墓那裡回去，把這一切事告訴十一個使徒和其餘的人。 \end{tabularx} \\ \\ \relax
24:10 & \begin{tabularx}{0.7\textwidth}{X} 把這些事告訴使徒的有抹大拉的馬利亞、約亞拿，和雅各的母親馬利亞，還有跟她們在一起的婦女。 \end{tabularx} \\ \\ \relax
24:11 & \begin{tabularx}{0.7\textwidth}{X} 她們這些話，使徒以為是胡言，就不相信。 \end{tabularx} \\ \\ \relax
24:12 & \begin{tabularx}{0.7\textwidth}{X} 彼得起來，跑到墳墓前，俯身往裡看，只見細麻布，就回去了，因所發生的事而心裡驚訝。 \end{tabularx} \\ \\ \relax
24:13 & \begin{tabularx}{0.7\textwidth}{X} 同一天，門徒中有兩個人往一個村子去；這村子名叫以馬忤斯，離耶路撒冷約有二十五里。 \end{tabularx} \\ \\ \relax
24:14 & \begin{tabularx}{0.7\textwidth}{X} 他們彼此談論所發生的這一切事。 \end{tabularx} \\ \\ \relax
24:15 & \begin{tabularx}{0.7\textwidth}{X} 正交談議論的時候，耶穌親自走近他們，和他們同行， \end{tabularx} \\ \\ \relax
24:16 & \begin{tabularx}{0.7\textwidth}{X} 可是他們的眼睛模糊了，沒認出他。 \end{tabularx} \\ \\ \relax
24:17 & \begin{tabularx}{0.7\textwidth}{X} 耶穌對他們說：「你們一邊走一邊談，彼此談論的是甚麼事呢？」他們就站住，臉上帶著愁容。 \end{tabularx} \\ \\ \relax
24:18 & \begin{tabularx}{0.7\textwidth}{X} 兩人中有一個名叫革流巴的回答：「你是在耶路撒冷的旅客中，惟一還不知道這幾天在那裡發生了甚麼事的人嗎？」 \end{tabularx} \\ \\ \relax
24:19 & \begin{tabularx}{0.7\textwidth}{X} 耶穌對他們說：「甚麼事呢？」他們對他說：「就是拿撒勒人耶穌的事。他是個先知，在神和眾百姓面前，說話行事都大有能力。 \end{tabularx} \\ \\ \relax
24:20 & \begin{tabularx}{0.7\textwidth}{X} 祭司長們和我們的官長竟把他解去，定了死罪，釘在十字架上。 \end{tabularx} \\ \\ \relax
24:21 & \begin{tabularx}{0.7\textwidth}{X} 但我們素來所盼望要救贖以色列民的就是他。不但如此，這些事發生到現在已經三天了。 \end{tabularx} \\ \\ \relax
24:22 & \begin{tabularx}{0.7\textwidth}{X} 還有，我們中間的幾個婦女使我們驚奇：她們清早去了墳墓， \end{tabularx} \\ \\ \relax
24:23 & \begin{tabularx}{0.7\textwidth}{X} 不見他的身體，就回來告訴我們，說她們看見了天使顯現，說他活了。 \end{tabularx} \\ \\ \relax
24:24 & \begin{tabularx}{0.7\textwidth}{X} 又有我們的幾個人往墳墓那裡去，所發現的正如婦女們所說的，只是沒有看見他。」 \end{tabularx} \\ \\ \relax
24:25 & \begin{tabularx}{0.7\textwidth}{X} 耶穌對他們說：「無知的人哪，先知所說的一切話，你們的心信得太遲鈍了。 \end{tabularx} \\ \\ \relax
24:26 & \begin{tabularx}{0.7\textwidth}{X} 基督不是必須受這些苦難，然後進入他的榮耀嗎？」 \end{tabularx} \\ \\ \relax
24:27 & \begin{tabularx}{0.7\textwidth}{X} 於是，他從摩西和眾先知起，凡經上所指著自己的話都給他們作了解釋。 \end{tabularx} \\ \\ \relax
24:28 & \begin{tabularx}{0.7\textwidth}{X} 他們走近所要去的村子，耶穌好像還要往前走， \end{tabularx} \\ \\ \relax
24:29 & \begin{tabularx}{0.7\textwidth}{X} 他們卻強留他說：「時候晚了，天快黑了，請你同我們住下吧。」耶穌就進去，要同他們住下。 \end{tabularx} \\ \\ \relax
24:30 & \begin{tabularx}{0.7\textwidth}{X} 坐下來和他們用餐的時候，耶穌拿起餅來，祝福了，擘開，遞給他們。 \end{tabularx} \\ \\ \relax
24:31 & \begin{tabularx}{0.7\textwidth}{X} 他們的眼睛開了，這才認出他來。耶穌卻從他們眼前消失了。 \end{tabularx} \\ \\ \relax
24:32 & \begin{tabularx}{0.7\textwidth}{X} 他們彼此說：「在路上他和我們說話，給我們講解聖經的時候，我們的心在我們裡面豈不是火熱的嗎？」 \end{tabularx} \\ \\ \relax
24:33 & \begin{tabularx}{0.7\textwidth}{X} 於是他們立刻起身，回耶路撒冷去，看見十一個使徒和與他們正在一起的人聚集在一處， \end{tabularx} \\ \\ \relax
24:34 & \begin{tabularx}{0.7\textwidth}{X} 說：「主果然復活了，已經顯現給西門看了。」 \end{tabularx} \\ \\ \relax
24:35 & \begin{tabularx}{0.7\textwidth}{X} 於是，兩個人把路上所遇到，和耶穌擘餅的時候怎麼被他們認出來的事，都述說了一遍。 \end{tabularx} \\ \\ \relax
24:36 & \begin{tabularx}{0.7\textwidth}{X} 正說這些話的時候，耶穌親自站在他們當中，說：「願你們平安！」 \end{tabularx} \\ \\ \relax
24:37 & \begin{tabularx}{0.7\textwidth}{X} 他們卻驚慌害怕，以為所看見的是魂。 \end{tabularx} \\ \\ \relax
24:38 & \begin{tabularx}{0.7\textwidth}{X} 耶穌對他們說：「你們為甚麼驚恐不安？為甚麼心裡起疑惑呢？ \end{tabularx} \\ \\ \relax
24:39 & \begin{tabularx}{0.7\textwidth}{X} 你們看我的手和我的腳，就知道實在是我了。摸摸我，看，因為魂無骨無肉，你們看，我是有的。」 \end{tabularx} \\ \\ \relax
24:40 & \begin{tabularx}{0.7\textwidth}{X} 說了這話，他就把手和腳給他們看。 \end{tabularx} \\ \\ \relax
24:41 & \begin{tabularx}{0.7\textwidth}{X} 他們還在又驚又喜、不敢相信的時候，耶穌對他們說：「你們這裡有甚麼吃的沒有？」 \end{tabularx} \\ \\ \relax
24:42 & \begin{tabularx}{0.7\textwidth}{X} 他們給了他一片烤魚， \end{tabularx} \\ \\ \relax
24:43 & \begin{tabularx}{0.7\textwidth}{X} 他接過來，在他們面前吃了。 \end{tabularx} \\ \\ \relax
24:44 & \begin{tabularx}{0.7\textwidth}{X} 耶穌對他們說：「這就是我從前和你們同在時所告訴你們的話：摩西的律法、先知的書，和《詩篇》上所記一切指著我的話都必須應驗。」 \end{tabularx} \\ \\ \relax
24:45 & \begin{tabularx}{0.7\textwidth}{X} 於是耶穌開他們的心竅，使他們能明白聖經， \end{tabularx} \\ \\ \relax
24:46 & \begin{tabularx}{0.7\textwidth}{X} 又對他們說：「照經上所寫的，基督必受害，第三天從死人中復活， \end{tabularx} \\ \\ \relax
24:47 & \begin{tabularx}{0.7\textwidth}{X} 並且人們要奉他的名傳悔改、使罪得赦的道，從耶路撒冷起直傳到萬邦。 \end{tabularx} \\ \\ \relax
24:48 & \begin{tabularx}{0.7\textwidth}{X} 你們就是這些事的見證。 \end{tabularx} \\ \\ \relax
24:49 & \begin{tabularx}{0.7\textwidth}{X} 我要將我父所應許的降在你們身上，你們要在城裡等候，直到你們領受從上面來的能力。」 \end{tabularx} \\ \\ \relax
24:50 & \begin{tabularx}{0.7\textwidth}{X} 耶穌領他們出來，直到伯大尼附近，就舉手給他們祝福。 \end{tabularx} \\ \\ \relax
24:51 & \begin{tabularx}{0.7\textwidth}{X} 正祝福的時候，他離開他們，被帶到天上去了。 \end{tabularx} \\ \\ \relax
24:52 & \begin{tabularx}{0.7\textwidth}{X} 他們就拜他，帶著極大的喜樂回耶路撒冷去， \end{tabularx} \\ \\ \relax
24:53 & \begin{tabularx}{0.7\textwidth}{X} 常在聖殿裡稱頌神。 \end{tabularx} \\ \\
[1ex]
\hline
\hline
\end{longtable}


\newpage

\section{第75屆~港九培靈會~蔡元雲~第八講}
\label{sec:397}
\textbf{同被差遣:廿一世紀的普世宣教}
\newline
\newline
連結: \href{https://www.hkbibleconference.org/session-message/view/397}{\texttt{https://www.hkbibleconference.org/session-message/view/397}}
\newline
\newline
\hyperref[sec:395]{< < < PREV SERMON < < <}
~
\hyperref[sec:toc]{[返目錄]}
~
\hyperref[sec:365]{> > > NEXT SERMON > > >}
\newline
\newline
路加福音 24:1-24:53
\newline
\begin{longtable}{cl}
\hline
\hline
章節 & 經文 (和合本修訂版)\\
\hline
24:1 & \begin{tabularx}{0.7\textwidth}{X} 七日的第一日，黎明的時候，那些婦女帶著所預備的香料來到墳墓那裡， \end{tabularx} \\ \\ \relax
24:2 & \begin{tabularx}{0.7\textwidth}{X} 發現石頭已經從墳墓滾開了， \end{tabularx} \\ \\ \relax
24:3 & \begin{tabularx}{0.7\textwidth}{X} 她們就進去，只是不見主耶穌的身體。 \end{tabularx} \\ \\ \relax
24:4 & \begin{tabularx}{0.7\textwidth}{X} 正在為這事困惑的時候，忽然有兩個人站在旁邊，衣服放光。 \end{tabularx} \\ \\ \relax
24:5 & \begin{tabularx}{0.7\textwidth}{X} 婦女們非常害怕，就俯伏在地上。那兩個人對她們說：「為甚麼在死人中找活人呢？ \end{tabularx} \\ \\ \relax
24:6 & \begin{tabularx}{0.7\textwidth}{X} 他不在這裡，已經復活了。要記得他還在加利利的時候怎樣告訴你們的， \end{tabularx} \\ \\ \relax
24:7 & \begin{tabularx}{0.7\textwidth}{X} 他說：『人子必須被交在罪人手裡，釘在十字架上，第三天復活。』」 \end{tabularx} \\ \\ \relax
24:8 & \begin{tabularx}{0.7\textwidth}{X} 她們就想起耶穌的話來。 \end{tabularx} \\ \\ \relax
24:9 & \begin{tabularx}{0.7\textwidth}{X} 於是她們從墳墓那裡回去，把這一切事告訴十一個使徒和其餘的人。 \end{tabularx} \\ \\ \relax
24:10 & \begin{tabularx}{0.7\textwidth}{X} 把這些事告訴使徒的有抹大拉的馬利亞、約亞拿，和雅各的母親馬利亞，還有跟她們在一起的婦女。 \end{tabularx} \\ \\ \relax
24:11 & \begin{tabularx}{0.7\textwidth}{X} 她們這些話，使徒以為是胡言，就不相信。 \end{tabularx} \\ \\ \relax
24:12 & \begin{tabularx}{0.7\textwidth}{X} 彼得起來，跑到墳墓前，俯身往裡看，只見細麻布，就回去了，因所發生的事而心裡驚訝。 \end{tabularx} \\ \\ \relax
24:13 & \begin{tabularx}{0.7\textwidth}{X} 同一天，門徒中有兩個人往一個村子去；這村子名叫以馬忤斯，離耶路撒冷約有二十五里。 \end{tabularx} \\ \\ \relax
24:14 & \begin{tabularx}{0.7\textwidth}{X} 他們彼此談論所發生的這一切事。 \end{tabularx} \\ \\ \relax
24:15 & \begin{tabularx}{0.7\textwidth}{X} 正交談議論的時候，耶穌親自走近他們，和他們同行， \end{tabularx} \\ \\ \relax
24:16 & \begin{tabularx}{0.7\textwidth}{X} 可是他們的眼睛模糊了，沒認出他。 \end{tabularx} \\ \\ \relax
24:17 & \begin{tabularx}{0.7\textwidth}{X} 耶穌對他們說：「你們一邊走一邊談，彼此談論的是甚麼事呢？」他們就站住，臉上帶著愁容。 \end{tabularx} \\ \\ \relax
24:18 & \begin{tabularx}{0.7\textwidth}{X} 兩人中有一個名叫革流巴的回答：「你是在耶路撒冷的旅客中，惟一還不知道這幾天在那裡發生了甚麼事的人嗎？」 \end{tabularx} \\ \\ \relax
24:19 & \begin{tabularx}{0.7\textwidth}{X} 耶穌對他們說：「甚麼事呢？」他們對他說：「就是拿撒勒人耶穌的事。他是個先知，在神和眾百姓面前，說話行事都大有能力。 \end{tabularx} \\ \\ \relax
24:20 & \begin{tabularx}{0.7\textwidth}{X} 祭司長們和我們的官長竟把他解去，定了死罪，釘在十字架上。 \end{tabularx} \\ \\ \relax
24:21 & \begin{tabularx}{0.7\textwidth}{X} 但我們素來所盼望要救贖以色列民的就是他。不但如此，這些事發生到現在已經三天了。 \end{tabularx} \\ \\ \relax
24:22 & \begin{tabularx}{0.7\textwidth}{X} 還有，我們中間的幾個婦女使我們驚奇：她們清早去了墳墓， \end{tabularx} \\ \\ \relax
24:23 & \begin{tabularx}{0.7\textwidth}{X} 不見他的身體，就回來告訴我們，說她們看見了天使顯現，說他活了。 \end{tabularx} \\ \\ \relax
24:24 & \begin{tabularx}{0.7\textwidth}{X} 又有我們的幾個人往墳墓那裡去，所發現的正如婦女們所說的，只是沒有看見他。」 \end{tabularx} \\ \\ \relax
24:25 & \begin{tabularx}{0.7\textwidth}{X} 耶穌對他們說：「無知的人哪，先知所說的一切話，你們的心信得太遲鈍了。 \end{tabularx} \\ \\ \relax
24:26 & \begin{tabularx}{0.7\textwidth}{X} 基督不是必須受這些苦難，然後進入他的榮耀嗎？」 \end{tabularx} \\ \\ \relax
24:27 & \begin{tabularx}{0.7\textwidth}{X} 於是，他從摩西和眾先知起，凡經上所指著自己的話都給他們作了解釋。 \end{tabularx} \\ \\ \relax
24:28 & \begin{tabularx}{0.7\textwidth}{X} 他們走近所要去的村子，耶穌好像還要往前走， \end{tabularx} \\ \\ \relax
24:29 & \begin{tabularx}{0.7\textwidth}{X} 他們卻強留他說：「時候晚了，天快黑了，請你同我們住下吧。」耶穌就進去，要同他們住下。 \end{tabularx} \\ \\ \relax
24:30 & \begin{tabularx}{0.7\textwidth}{X} 坐下來和他們用餐的時候，耶穌拿起餅來，祝福了，擘開，遞給他們。 \end{tabularx} \\ \\ \relax
24:31 & \begin{tabularx}{0.7\textwidth}{X} 他們的眼睛開了，這才認出他來。耶穌卻從他們眼前消失了。 \end{tabularx} \\ \\ \relax
24:32 & \begin{tabularx}{0.7\textwidth}{X} 他們彼此說：「在路上他和我們說話，給我們講解聖經的時候，我們的心在我們裡面豈不是火熱的嗎？」 \end{tabularx} \\ \\ \relax
24:33 & \begin{tabularx}{0.7\textwidth}{X} 於是他們立刻起身，回耶路撒冷去，看見十一個使徒和與他們正在一起的人聚集在一處， \end{tabularx} \\ \\ \relax
24:34 & \begin{tabularx}{0.7\textwidth}{X} 說：「主果然復活了，已經顯現給西門看了。」 \end{tabularx} \\ \\ \relax
24:35 & \begin{tabularx}{0.7\textwidth}{X} 於是，兩個人把路上所遇到，和耶穌擘餅的時候怎麼被他們認出來的事，都述說了一遍。 \end{tabularx} \\ \\ \relax
24:36 & \begin{tabularx}{0.7\textwidth}{X} 正說這些話的時候，耶穌親自站在他們當中，說：「願你們平安！」 \end{tabularx} \\ \\ \relax
24:37 & \begin{tabularx}{0.7\textwidth}{X} 他們卻驚慌害怕，以為所看見的是魂。 \end{tabularx} \\ \\ \relax
24:38 & \begin{tabularx}{0.7\textwidth}{X} 耶穌對他們說：「你們為甚麼驚恐不安？為甚麼心裡起疑惑呢？ \end{tabularx} \\ \\ \relax
24:39 & \begin{tabularx}{0.7\textwidth}{X} 你們看我的手和我的腳，就知道實在是我了。摸摸我，看，因為魂無骨無肉，你們看，我是有的。」 \end{tabularx} \\ \\ \relax
24:40 & \begin{tabularx}{0.7\textwidth}{X} 說了這話，他就把手和腳給他們看。 \end{tabularx} \\ \\ \relax
24:41 & \begin{tabularx}{0.7\textwidth}{X} 他們還在又驚又喜、不敢相信的時候，耶穌對他們說：「你們這裡有甚麼吃的沒有？」 \end{tabularx} \\ \\ \relax
24:42 & \begin{tabularx}{0.7\textwidth}{X} 他們給了他一片烤魚， \end{tabularx} \\ \\ \relax
24:43 & \begin{tabularx}{0.7\textwidth}{X} 他接過來，在他們面前吃了。 \end{tabularx} \\ \\ \relax
24:44 & \begin{tabularx}{0.7\textwidth}{X} 耶穌對他們說：「這就是我從前和你們同在時所告訴你們的話：摩西的律法、先知的書，和《詩篇》上所記一切指著我的話都必須應驗。」 \end{tabularx} \\ \\ \relax
24:45 & \begin{tabularx}{0.7\textwidth}{X} 於是耶穌開他們的心竅，使他們能明白聖經， \end{tabularx} \\ \\ \relax
24:46 & \begin{tabularx}{0.7\textwidth}{X} 又對他們說：「照經上所寫的，基督必受害，第三天從死人中復活， \end{tabularx} \\ \\ \relax
24:47 & \begin{tabularx}{0.7\textwidth}{X} 並且人們要奉他的名傳悔改、使罪得赦的道，從耶路撒冷起直傳到萬邦。 \end{tabularx} \\ \\ \relax
24:48 & \begin{tabularx}{0.7\textwidth}{X} 你們就是這些事的見證。 \end{tabularx} \\ \\ \relax
24:49 & \begin{tabularx}{0.7\textwidth}{X} 我要將我父所應許的降在你們身上，你們要在城裡等候，直到你們領受從上面來的能力。」 \end{tabularx} \\ \\ \relax
24:50 & \begin{tabularx}{0.7\textwidth}{X} 耶穌領他們出來，直到伯大尼附近，就舉手給他們祝福。 \end{tabularx} \\ \\ \relax
24:51 & \begin{tabularx}{0.7\textwidth}{X} 正祝福的時候，他離開他們，被帶到天上去了。 \end{tabularx} \\ \\ \relax
24:52 & \begin{tabularx}{0.7\textwidth}{X} 他們就拜他，帶著極大的喜樂回耶路撒冷去， \end{tabularx} \\ \\ \relax
24:53 & \begin{tabularx}{0.7\textwidth}{X} 常在聖殿裡稱頌神。 \end{tabularx} \\ \\
[1ex]
\hline
\hline
\end{longtable}


\newpage

\section{第75屆~港九培靈會~蕭壽華~第一講}
\label{sec:365}
\textbf{只要剛強壯膽}
\newline
\newline
連結: \href{https://www.hkbibleconference.org/session-message/view/365}{\texttt{https://www.hkbibleconference.org/session-message/view/365}}
\newline
\newline
\hyperref[sec:397]{< < < PREV SERMON < < <}
~
\hyperref[sec:toc]{[返目錄]}
~
\hyperref[sec:366]{> > > NEXT SERMON > > >}
\newline
\newline
約書亞記 1:1-1:18
\newline
\begin{longtable}{cl}
\hline
\hline
章節 & 經文 (和合本修訂版)\\
\hline
1:1 & \begin{tabularx}{0.7\textwidth}{X} 耶和華的僕人摩西死了以後，耶和華對摩西的助手嫩的兒子約書亞說： \end{tabularx} \\ \\ \relax
1:2 & \begin{tabularx}{0.7\textwidth}{X} 「我的僕人摩西死了。現在你要起來，和眾百姓過這約旦河，往我所要賜給以色列人的地去。 \end{tabularx} \\ \\ \relax
1:3 & \begin{tabularx}{0.7\textwidth}{X} 凡你們腳掌所踏之地，我都照我所應許摩西的話賜給你們了。 \end{tabularx} \\ \\ \relax
1:4 & \begin{tabularx}{0.7\textwidth}{X} 從曠野和這黎巴嫩，直到大河，就是幼發拉底河，赫人的全地，又到大海日落的方向，都要作你們的疆土。 \end{tabularx} \\ \\ \relax
1:5 & \begin{tabularx}{0.7\textwidth}{X} 你一生的日子，必無人能在你面前站立得住。我怎樣與摩西同在，也必照樣與你同在；我必不撇下你，也不丟棄你。 \end{tabularx} \\ \\ \relax
1:6 & \begin{tabularx}{0.7\textwidth}{X} 你當剛強壯膽，因為你必使這百姓承受那地為業，就是我向他們列祖起誓要給他們的地。 \end{tabularx} \\ \\ \relax
1:7 & \begin{tabularx}{0.7\textwidth}{X} 只要剛強，大大壯膽，謹守遵行我僕人摩西所吩咐你的一切律法，不可偏離左右，使你無論往哪裡去，都可以順利。 \end{tabularx} \\ \\ \relax
1:8 & \begin{tabularx}{0.7\textwidth}{X} 這律法書不可離開你的口，總要晝夜思想，好使你謹守遵行這書上所寫的一切話。如此，你的道路就可以亨通，凡事順利。 \end{tabularx} \\ \\ \relax
1:9 & \begin{tabularx}{0.7\textwidth}{X} 我豈沒有吩咐你嗎？你當剛強壯膽，不要懼怕，也不要驚惶，因為你無論往哪裡去，耶和華你的神必與你同在。」 \end{tabularx} \\ \\ \relax
1:10 & \begin{tabularx}{0.7\textwidth}{X} 於是，約書亞吩咐百姓的官長說： \end{tabularx} \\ \\ \relax
1:11 & \begin{tabularx}{0.7\textwidth}{X} 「你們要走遍營中，吩咐百姓說：『當預備食物， 因為三日之內你們要過這約旦河，進去得耶和華－你們神賜給你們為業之地。』」 \end{tabularx} \\ \\ \relax
1:12 & \begin{tabularx}{0.7\textwidth}{X} 約書亞對呂便人、迦得人和瑪拿西半支派的人說： \end{tabularx} \\ \\ \relax
1:13 & \begin{tabularx}{0.7\textwidth}{X} 「你們要記得耶和華的僕人摩西所吩咐你們的話說：『耶和華－你們的神使你們得享安寧，必將這地賜給你們。』 \end{tabularx} \\ \\ \relax
1:14 & \begin{tabularx}{0.7\textwidth}{X} 你們的妻子、孩子和牲畜可以留在約旦河東、摩西所給你們的地。但你們中間所有大能的勇士都要帶著兵器，在你們的弟兄前面過去，你們要幫助他們。 \end{tabularx} \\ \\ \relax
1:15 & \begin{tabularx}{0.7\textwidth}{X} 等到耶和華使你們的弟兄和你們一樣得享平靜，並且得著耶和華－你們神所賜他們為業之地的時候，你們才可以回到你們所得之地，承受為業，就是耶和華的僕人摩西在約旦河東、向日出的方向所給你們的地。」 \end{tabularx} \\ \\ \relax
1:16 & \begin{tabularx}{0.7\textwidth}{X} 他們回答約書亞說：「凡你吩咐我們的，我們都必做；凡你差我們去的地方，我們都必去。 \end{tabularx} \\ \\ \relax
1:17 & \begin{tabularx}{0.7\textwidth}{X} 我們在一切事上怎樣聽從摩西，也必照樣聽從你。惟願耶和華－你的神與你同在，像與摩西同在一樣。 \end{tabularx} \\ \\ \relax
1:18 & \begin{tabularx}{0.7\textwidth}{X} 無論甚麼人違背你的命令，不聽從你所吩咐他的一切話，就必處死。你只要剛強壯膽！」 \end{tabularx} \\ \\
[1ex]
\hline
\hline
\end{longtable}


\newpage

\section{第75屆~港九培靈會~蕭壽華~第二講}
\label{sec:366}
\textbf{尋求真實的信心}
\newline
\newline
連結: \href{https://www.hkbibleconference.org/session-message/view/366}{\texttt{https://www.hkbibleconference.org/session-message/view/366}}
\newline
\newline
\hyperref[sec:365]{< < < PREV SERMON < < <}
~
\hyperref[sec:toc]{[返目錄]}
~
\hyperref[sec:367]{> > > NEXT SERMON > > >}
\newline
\newline
約書亞記 2:1-2:24
\newline
\begin{longtable}{cl}
\hline
\hline
章節 & 經文 (和合本修訂版)\\
\hline
2:1 & \begin{tabularx}{0.7\textwidth}{X} 嫩的兒子約書亞從什亭暗中派兩個人作探子，說：「你們去窺探那地和耶利哥。」於是二人去了，來到一個名叫喇合的妓女家裡，在那裡睡覺。 \end{tabularx} \\ \\ \relax
2:2 & \begin{tabularx}{0.7\textwidth}{X} 有人告訴耶利哥王說：「看哪，今夜有以色列人到這裡來窺探此地。」 \end{tabularx} \\ \\ \relax
2:3 & \begin{tabularx}{0.7\textwidth}{X} 耶利哥王派人到喇合那裡， 說：「你要交出那來到你這裡、進了你家的人，因為他們來是要窺探全地。」 \end{tabularx} \\ \\ \relax
2:4 & \begin{tabularx}{0.7\textwidth}{X} 但女人已把二人藏起來，卻說：「那兩個人確實到我這裡來過，他們從哪裡來，我卻不知道。 \end{tabularx} \\ \\ \relax
2:5 & \begin{tabularx}{0.7\textwidth}{X} 天黑、要關城門的時候，他們就出去了。他們往哪裡去我也不知道。你們趕快去追他們，就必追上。」 \end{tabularx} \\ \\ \relax
2:6 & \begin{tabularx}{0.7\textwidth}{X} 其實，這女人已經領二人上了屋頂，把他們藏在她擺列在屋頂的亞麻梗中。 \end{tabularx} \\ \\ \relax
2:7 & \begin{tabularx}{0.7\textwidth}{X} 那些人就往約旦河的路上追趕他們，直到渡口。追趕他們的人一出去，城門就關了。 \end{tabularx} \\ \\ \relax
2:8 & \begin{tabularx}{0.7\textwidth}{X} 二人還沒有睡之前，女人就上屋頂，到他們那裡， \end{tabularx} \\ \\ \relax
2:9 & \begin{tabularx}{0.7\textwidth}{X} 對他們說：「我知道耶和華已經把這地賜給你們了，並且我們也都懼怕你們。這地所有的居民在你們面前都融化了。 \end{tabularx} \\ \\ \relax
2:10 & \begin{tabularx}{0.7\textwidth}{X} 因為我們聽見你們出埃及的時候，耶和華怎樣在你們前面使紅海的水乾了，並且你們怎樣處置約旦河東的兩個亞摩利王，西宏和噩，把他們完全消滅。 \end{tabularx} \\ \\ \relax
2:11 & \begin{tabularx}{0.7\textwidth}{X} 我們一聽見就膽戰心驚，人人因你們的緣故勇氣全失。耶和華－你們的神是天上地下的神。 \end{tabularx} \\ \\ \relax
2:12 & \begin{tabularx}{0.7\textwidth}{X} 現在我既然恩待你們，求你們指著耶和華向我起誓，你們也要恩待我的父家。請你們給我一個確實的憑據， \end{tabularx} \\ \\ \relax
2:13 & \begin{tabularx}{0.7\textwidth}{X} 要救活我的父母、兄弟、姊妹，和所有屬他們的，拯救我們的性命脫離死亡。」 \end{tabularx} \\ \\ \relax
2:14 & \begin{tabularx}{0.7\textwidth}{X} 那二人對她說：「我們願意以性命來替你們死。你們若不洩漏我們這件事，當耶和華將這地賜給我們的時候，我們必以慈愛和誠信待你。」 \end{tabularx} \\ \\ \relax
2:15 & \begin{tabularx}{0.7\textwidth}{X} 於是女人用繩子把二人從窗戶縋下去，因為她的屋子是在城牆邊上，她也住在城牆上。 \end{tabularx} \\ \\ \relax
2:16 & \begin{tabularx}{0.7\textwidth}{X} 她對他們說：「你們暫且往山上去，免得追趕的人遇見你們。要在那裡躲藏三天，等追趕的人回來，你們才可以走自己的路。」 \end{tabularx} \\ \\ \relax
2:17 & \begin{tabularx}{0.7\textwidth}{X} 二人對她說：「你叫我們所起的誓與我們無關， \end{tabularx} \\ \\ \relax
2:18 & \begin{tabularx}{0.7\textwidth}{X} 除非，看哪，當我們來到這地的時候，你把這條朱紅線繩子繫在縋我們下去的窗戶上，並要叫你的父母、兄弟和你父的全家都聚集在你家中。 \end{tabularx} \\ \\ \relax
2:19 & \begin{tabularx}{0.7\textwidth}{X} 凡離開你家門往街上去的，他的血必歸到自己頭上，與我們無關；凡在你家裡的，若有人下手害他，他的血就歸到我們頭上。 \end{tabularx} \\ \\ \relax
2:20 & \begin{tabularx}{0.7\textwidth}{X} 你若洩漏我們這件事，你叫我們所起的誓就與我們無關了。」 \end{tabularx} \\ \\ \relax
2:21 & \begin{tabularx}{0.7\textwidth}{X} 女人說：「就照你們的話吧！」於是她送他們走了，就把朱紅繩子繫在窗戶上。 \end{tabularx} \\ \\ \relax
2:22 & \begin{tabularx}{0.7\textwidth}{X} 二人離開，到山上去，在那裡停留三天，直等到追趕的人回去。追趕的人一路尋找，卻找不著。 \end{tabularx} \\ \\ \relax
2:23 & \begin{tabularx}{0.7\textwidth}{X} 二人回來，下了山，過了河，來到嫩的兒子約書亞那裡，向他報告他們所遭遇的一切事。 \end{tabularx} \\ \\ \relax
2:24 & \begin{tabularx}{0.7\textwidth}{X} 他們對約書亞說：「耶和華果然將那全地交在我們手中了，並且那地所有的居民在我們面前都融化了。」 \end{tabularx} \\ \\
[1ex]
\hline
\hline
\end{longtable}


\newpage

\section{第75屆~港九培靈會~蕭壽華~第三講}
\label{sec:367}
\textbf{神在江河中開路}
\newline
\newline
連結: \href{https://www.hkbibleconference.org/session-message/view/367}{\texttt{https://www.hkbibleconference.org/session-message/view/367}}
\newline
\newline
\hyperref[sec:366]{< < < PREV SERMON < < <}
~
\hyperref[sec:toc]{[返目錄]}
~
\hyperref[sec:369]{> > > NEXT SERMON > > >}
\newline
\newline
約書亞記 3:1-3:17
\newline
\begin{longtable}{cl}
\hline
\hline
章節 & 經文 (和合本修訂版)\\
\hline
3:1 & \begin{tabularx}{0.7\textwidth}{X} 約書亞清早起來，和以色列眾人起行，離開什亭，來到約旦河，過河以前住在那裡。 \end{tabularx} \\ \\ \relax
3:2 & \begin{tabularx}{0.7\textwidth}{X} 過了三天，官長走遍營中， \end{tabularx} \\ \\ \relax
3:3 & \begin{tabularx}{0.7\textwidth}{X} 吩咐百姓說：「當你們看見利未家的祭司抬著耶和華－你們神的約櫃的時候，你們就要起行離開所住的地方，跟著約櫃走， \end{tabularx} \\ \\ \relax
3:4 & \begin{tabularx}{0.7\textwidth}{X} 使你們知道所當走的路，因為這條路是你們從來沒有走過的。只是你們要與約櫃相隔約二千肘，不可太靠近約櫃。」 \end{tabularx} \\ \\ \relax
3:5 & \begin{tabularx}{0.7\textwidth}{X} 約書亞吩咐百姓說：「你們要使自己分別為聖，因為明天耶和華必在你們中間行奇事。」 \end{tabularx} \\ \\ \relax
3:6 & \begin{tabularx}{0.7\textwidth}{X} 約書亞對祭司說：「你們抬起約櫃，在百姓的前面過去。」於是他們抬起約櫃，走在百姓前面。 \end{tabularx} \\ \\ \relax
3:7 & \begin{tabularx}{0.7\textwidth}{X} 耶和華對約書亞說：「從今日起，我必使你在以色列眾人眼前被尊為大，使他們知道我怎樣與摩西同在，也必照樣與你同在。 \end{tabularx} \\ \\ \relax
3:8 & \begin{tabularx}{0.7\textwidth}{X} 你要吩咐抬約櫃的祭司說：『你們到了約旦河的水邊，要在約旦河中站著。』」 \end{tabularx} \\ \\ \relax
3:9 & \begin{tabularx}{0.7\textwidth}{X} 約書亞對以色列人說：「你們近前，到這裡來，聽耶和華－你們神的話。」 \end{tabularx} \\ \\ \relax
3:10 & \begin{tabularx}{0.7\textwidth}{X} 約書亞說：「你們因這事會知道永生的神在你們中間，他必從你們面前趕出迦南人、赫人、希未人、比利洗人、革迦撒人、亞摩利人、耶布斯人。 \end{tabularx} \\ \\ \relax
3:11 & \begin{tabularx}{0.7\textwidth}{X} 看哪！全地之主的約櫃必在你們的前面過去，到約旦河裡。 \end{tabularx} \\ \\ \relax
3:12 & \begin{tabularx}{0.7\textwidth}{X} 現在， 你們要從以色列支派中選出十二個人，每支派一人。 \end{tabularx} \\ \\ \relax
3:13 & \begin{tabularx}{0.7\textwidth}{X} 當抬耶和華全地之主約櫃的祭司，腳掌踏入約旦河水裡的時候，約旦河的水，就是從上往下流的水，必然中斷，豎立成壘。」 \end{tabularx} \\ \\ \relax
3:14 & \begin{tabularx}{0.7\textwidth}{X} 百姓起行離開帳棚過約旦河的時候，抬約櫃的祭司在百姓的前面。 \end{tabularx} \\ \\ \relax
3:15 & \begin{tabularx}{0.7\textwidth}{X} 那時正是收割的日子，約旦河的水漲滿兩岸。抬約櫃的人到了約旦河，抬約櫃的祭司腳一入水邊， \end{tabularx} \\ \\ \relax
3:16 & \begin{tabularx}{0.7\textwidth}{X} 那從上往下流的水就在很遠的地方，在撒拉但旁邊的亞當城那裡停住，豎立成壘；那往亞拉巴海，就是鹽海下流的水全然中斷。於是，百姓在耶利哥的對面過了河。 \end{tabularx} \\ \\ \relax
3:17 & \begin{tabularx}{0.7\textwidth}{X} 抬耶和華約櫃的祭司在約旦河中的乾地上穩穩站著，以色列眾人都從乾地上過去，直到全國都過了約旦河。 \end{tabularx} \\ \\
[1ex]
\hline
\hline
\end{longtable}


\newpage

\section{第75屆~港九培靈會~蕭壽華~第四講}
\label{sec:369}
\textbf{調整與神的關係}
\newline
\newline
連結: \href{https://www.hkbibleconference.org/session-message/view/369}{\texttt{https://www.hkbibleconference.org/session-message/view/369}}
\newline
\newline
\hyperref[sec:367]{< < < PREV SERMON < < <}
~
\hyperref[sec:toc]{[返目錄]}
~
\hyperref[sec:371]{> > > NEXT SERMON > > >}
\newline
\newline
約書亞記 5:1-5:15
\newline
\begin{longtable}{cl}
\hline
\hline
章節 & 經文 (和合本修訂版)\\
\hline
5:1 & \begin{tabularx}{0.7\textwidth}{X} 約旦河西亞摩利人的眾王和靠海迦南人的眾王，聽見耶和華在以色列人前面使約旦河的水乾了，直到他們過了河，眾王因以色列人的緣故都膽戰心驚，勇氣全失。 \end{tabularx} \\ \\ \relax
5:2 & \begin{tabularx}{0.7\textwidth}{X} 那時，耶和華對約書亞說：「你要造火石刀，第二次為以色列人行割禮。」 \end{tabularx} \\ \\ \relax
5:3 & \begin{tabularx}{0.7\textwidth}{X} 約書亞就造了火石刀，在哈爾拉勒山為以色列人行割禮。 \end{tabularx} \\ \\ \relax
5:4 & \begin{tabularx}{0.7\textwidth}{X} 約書亞行割禮的原因是這樣：從埃及出來的眾百姓，所有能打仗的男丁，出了埃及以後，都死在曠野的路上。 \end{tabularx} \\ \\ \relax
5:5 & \begin{tabularx}{0.7\textwidth}{X} 這些從埃及出來的眾百姓都受過割禮；但是那些出埃及以後，在曠野的路上所生的眾百姓卻沒有受過割禮。 \end{tabularx} \\ \\ \relax
5:6 & \begin{tabularx}{0.7\textwidth}{X} 以色列人在曠野走了四十年，直到那從埃及出來，全國能打仗的人都消滅了，因為他們沒有聽從耶和華的話。耶和華曾向他們起誓，必不容許他們看見耶和華向他們列祖起誓要給我們的地，就是流奶與蜜之地。 \end{tabularx} \\ \\ \relax
5:7 & \begin{tabularx}{0.7\textwidth}{X} 他們的子孫，就是耶和華興起接續他們的，都沒有受過割禮；因為在路上他們沒有受割禮，約書亞就為他們行割禮。 \end{tabularx} \\ \\ \relax
5:8 & \begin{tabularx}{0.7\textwidth}{X} 全國的人都受了割禮，留在營中自己的地方，直到痊癒。 \end{tabularx} \\ \\ \relax
5:9 & \begin{tabularx}{0.7\textwidth}{X} 耶和華對約書亞說：「我今日將埃及的羞辱從你們身上除掉了。」因此，那地方名叫吉甲，直到今日。 \end{tabularx} \\ \\ \relax
5:10 & \begin{tabularx}{0.7\textwidth}{X} 以色列人在吉甲安營。正月十四日晚上，他們在耶利哥的平原守逾越節。 \end{tabularx} \\ \\ \relax
5:11 & \begin{tabularx}{0.7\textwidth}{X} 逾越節的第二日，他們吃了當地的出產，就在那一天，吃了無酵餅和烘過的穀物。 \end{tabularx} \\ \\ \relax
5:12 & \begin{tabularx}{0.7\textwidth}{X} 他們吃了當地出產的第二日，嗎哪就停止了。以色列人不再有嗎哪了。那一年，他們就吃迦南地的出產。 \end{tabularx} \\ \\ \relax
5:13 & \begin{tabularx}{0.7\textwidth}{X} 約書亞靠近耶利哥的時候，舉目觀看，看哪，有一個人站在他對面，手裡拿著拔出來的刀。約書亞到他那裡，對他說：「你是屬我們的，還是屬我們敵人的呢？」 \end{tabularx} \\ \\ \relax
5:14 & \begin{tabularx}{0.7\textwidth}{X} 他說：「不，我現在來是要作耶和華軍隊的元帥。」約書亞就臉伏於地下拜，說：「我主有甚麼話，請吩咐僕人吧！」 \end{tabularx} \\ \\ \relax
5:15 & \begin{tabularx}{0.7\textwidth}{X} 耶和華軍隊的元帥對約書亞說：「把你腳上的鞋脫下來，因為你所站的地方是聖的。」約書亞就照著做了。 \end{tabularx} \\ \\
[1ex]
\hline
\hline
\end{longtable}


\newpage

\section{第75屆~港九培靈會~蕭壽華~第五講}
\label{sec:371}
\textbf{信靠順服中的得勝}
\newline
\newline
連結: \href{https://www.hkbibleconference.org/session-message/view/371}{\texttt{https://www.hkbibleconference.org/session-message/view/371}}
\newline
\newline
\hyperref[sec:369]{< < < PREV SERMON < < <}
~
\hyperref[sec:toc]{[返目錄]}
~
\hyperref[sec:374]{> > > NEXT SERMON > > >}
\newline
\newline
約書亞記 6:1-6:21
\newline
\begin{longtable}{cl}
\hline
\hline
章節 & 經文 (和合本修訂版)\\
\hline
6:1 & \begin{tabularx}{0.7\textwidth}{X} 耶利哥的城門因以色列人的緣故，關得嚴緊，無人出入。 \end{tabularx} \\ \\ \relax
6:2 & \begin{tabularx}{0.7\textwidth}{X} 耶和華對約書亞說：「看，我已經把耶利哥城和耶利哥王，以及大能的勇士，都交在你手中。 \end{tabularx} \\ \\ \relax
6:3 & \begin{tabularx}{0.7\textwidth}{X} 你們要圍繞這城，所有的士兵繞城一次，六日你都要這樣做。 \end{tabularx} \\ \\ \relax
6:4 & \begin{tabularx}{0.7\textwidth}{X} 七個祭司要拿七個羊角走在約櫃前。到了第七日，你們要圍繞這城七次，祭司也要吹角。 \end{tabularx} \\ \\ \relax
6:5 & \begin{tabularx}{0.7\textwidth}{X} 羊角聲拖長的時候，你們一聽見角聲，眾百姓要大聲呼喊，城牆就必倒塌，各人要往前直上。」 \end{tabularx} \\ \\ \relax
6:6 & \begin{tabularx}{0.7\textwidth}{X} 嫩的兒子約書亞召了祭司來，對他們說：「你們抬起約櫃來，要有七個祭司拿七個羊角在耶和華的約櫃前。」 \end{tabularx} \\ \\ \relax
6:7 & \begin{tabularx}{0.7\textwidth}{X} 他又對百姓說：「你們向前去圍繞那城，帶兵器的要在耶和華的約櫃前過去。」 \end{tabularx} \\ \\ \relax
6:8 & \begin{tabularx}{0.7\textwidth}{X} 按照約書亞對百姓所說的，七個祭司拿了七個羊角在耶和華面前過去，他們吹著角，耶和華的約櫃在他們後面跟著。 \end{tabularx} \\ \\ \relax
6:9 & \begin{tabularx}{0.7\textwidth}{X} 帶兵器的走在吹角的祭司前面，後隊跟著約櫃走，號角繼續在吹。 \end{tabularx} \\ \\ \relax
6:10 & \begin{tabularx}{0.7\textwidth}{X} 約書亞吩咐百姓說：「你們不可呼喊，不可讓人聽見你們的聲音，連一句話也不可出你們的口，直到我對你們說『呼喊』的那日，你們才呼喊。」 \end{tabularx} \\ \\ \relax
6:11 & \begin{tabularx}{0.7\textwidth}{X} 這樣，約書亞使耶和華的約櫃圍繞那城，把城繞了一次。然後，眾人回到營裡，就在營裡住宿。 \end{tabularx} \\ \\ \relax
6:12 & \begin{tabularx}{0.7\textwidth}{X} 約書亞清早起來，祭司又抬起耶和華的約櫃。 \end{tabularx} \\ \\ \relax
6:13 & \begin{tabularx}{0.7\textwidth}{X} 七個祭司拿七個羊角，走在耶和華的約櫃前，他們吹著角；帶兵器的走在他們前面，後隊跟著耶和華的約櫃走，號角繼續在吹。 \end{tabularx} \\ \\ \relax
6:14 & \begin{tabularx}{0.7\textwidth}{X} 第二日，他們再把城圍繞一次，就回營裡去。六日都是這樣做。 \end{tabularx} \\ \\ \relax
6:15 & \begin{tabularx}{0.7\textwidth}{X} 第七日清早黎明時，他們起來，以同樣的方式圍繞城七次；惟獨這一日他們圍繞城七次。 \end{tabularx} \\ \\ \relax
6:16 & \begin{tabularx}{0.7\textwidth}{X} 到了第七次，祭司吹角的時候，約書亞對百姓說：「呼喊吧，因為耶和華已經把城交給你們了！ \end{tabularx} \\ \\ \relax
6:17 & \begin{tabularx}{0.7\textwidth}{X} 這城和其中所有的都要永獻給耶和華作當毀滅的，只有妓女喇合與她家中所有的可以存活，因為她隱藏了我們所派的使者。 \end{tabularx} \\ \\ \relax
6:18 & \begin{tabularx}{0.7\textwidth}{X} 但你們務必謹慎，不可取那當滅的物，免得你們受詛咒，取了那當滅的物，使以色列全營成為詛咒而遭受災禍。 \end{tabularx} \\ \\ \relax
6:19 & \begin{tabularx}{0.7\textwidth}{X} 只有金子、銀子和銅鐵的器皿都要歸耶和華為聖，放入耶和華的庫房中。」 \end{tabularx} \\ \\ \relax
6:20 & \begin{tabularx}{0.7\textwidth}{X} 於是百姓呼喊，祭司吹角。百姓一聽見角聲就大聲呼喊，城牆隨著倒塌。百姓上去進城，各人往前直上，把城奪取。 \end{tabularx} \\ \\ \relax
6:21 & \begin{tabularx}{0.7\textwidth}{X} 他們把城中所有的，無論男女老少，牛羊和驢，都用刀殺盡。 \end{tabularx} \\ \\
[1ex]
\hline
\hline
\end{longtable}


\newpage

\section{第75屆~港九培靈會~蕭壽華~第六講}
\label{sec:374}
\textbf{不容罪叫教會傾覆}
\newline
\newline
連結: \href{https://www.hkbibleconference.org/session-message/view/374}{\texttt{https://www.hkbibleconference.org/session-message/view/374}}
\newline
\newline
\hyperref[sec:371]{< < < PREV SERMON < < <}
~
\hyperref[sec:toc]{[返目錄]}
~
\hyperref[sec:377]{> > > NEXT SERMON > > >}
\newline
\newline
約書亞記 7:1-7:26
\newline
\begin{longtable}{cl}
\hline
\hline
章節 & 經文 (和合本修訂版)\\
\hline
7:1 & \begin{tabularx}{0.7\textwidth}{X} 以色列人在當滅之物上犯了罪。猶大支派中，謝拉的曾孫，撒底的孫子，迦米的兒子亞干取了當滅之物，耶和華的怒氣就向以色列人發作。 \end{tabularx} \\ \\ \relax
7:2 & \begin{tabularx}{0.7\textwidth}{X} 約書亞從耶利哥派人往伯特利東邊，靠近伯‧亞文的艾城去，對他們說：「你們上去窺探那地。」那些人就上去窺探艾城。 \end{tabularx} \\ \\ \relax
7:3 & \begin{tabularx}{0.7\textwidth}{X} 他們回到約書亞那裡，對他說：「眾百姓不必都上去，只要二、三千人上去就能攻取艾城；不必勞動眾百姓都上去，因為他們人少。」 \end{tabularx} \\ \\ \relax
7:4 & \begin{tabularx}{0.7\textwidth}{X} 於是百姓中約有三千人上那裡去，但他們竟在艾城的人面前逃跑。 \end{tabularx} \\ \\ \relax
7:5 & \begin{tabularx}{0.7\textwidth}{X} 艾城的人擊殺他們約三十六人，從城門前追趕他們，直到示巴琳，在下坡的地方擊敗他們。他們都膽戰心驚，融化如水。 \end{tabularx} \\ \\ \relax
7:6 & \begin{tabularx}{0.7\textwidth}{X} 約書亞和以色列的長老就撕裂衣服，在耶和華的約櫃前臉伏於地，直到晚上。他們把灰撒在頭上。 \end{tabularx} \\ \\ \relax
7:7 & \begin{tabularx}{0.7\textwidth}{X} 約書亞說：「唉！主耶和華啊，你為甚麼領這百姓過約旦河，把我們交在亞摩利人手中，使我們滅亡呢？我們不如住在約旦河的那邊！ \end{tabularx} \\ \\ \relax
7:8 & \begin{tabularx}{0.7\textwidth}{X} 主啊，求求你，以色列人既在仇敵面前轉身逃跑，我還有甚麼可說的呢？ \end{tabularx} \\ \\ \relax
7:9 & \begin{tabularx}{0.7\textwidth}{X} 迦南人和這地所有的居民聽見了就必圍困我們，把我們的名從地上除去。那時，你為你至大的名要怎樣做呢？」 \end{tabularx} \\ \\ \relax
7:10 & \begin{tabularx}{0.7\textwidth}{X} 耶和華對約書亞說：「起來！你的臉為何這樣俯伏呢？ \end{tabularx} \\ \\ \relax
7:11 & \begin{tabularx}{0.7\textwidth}{X} 以色列犯了罪，又違背了我所吩咐他們的約，又取了當滅之物。他們又偷竊，又行詭詐，又把那當滅的物與自己的器皿放在一起。 \end{tabularx} \\ \\ \relax
7:12 & \begin{tabularx}{0.7\textwidth}{X} 因此，以色列人在仇敵面前站立不住。他們在仇敵面前轉身逃跑，因為他們成了當滅的物。你們若不把當滅的物從你們中間除掉，我就不再與你們同在了。 \end{tabularx} \\ \\ \relax
7:13 & \begin{tabularx}{0.7\textwidth}{X} 你起來，去叫百姓分別為聖，說：『你們要為了明天使自己分別為聖，因為耶和華－以色列的神這樣說：以色列啊，在你中間有當滅的物；你們若不把你們中間當滅之物除掉，你在仇敵面前必站立不住！』 \end{tabularx} \\ \\ \relax
7:14 & \begin{tabularx}{0.7\textwidth}{X} 到了早晨，你們要按著支派近前來。耶和華所選的支派，要按著宗族近前來；耶和華所選的宗族，要按著家族近前來；耶和華所選的家族，要按著男丁，一個一個近前來。 \end{tabularx} \\ \\ \relax
7:15 & \begin{tabularx}{0.7\textwidth}{X} 被選的人有當滅之物在他那裡，他和他所有的必被火焚燒，因為他違背了耶和華的約，又因他在以色列中做了愚妄的事。」 \end{tabularx} \\ \\ \relax
7:16 & \begin{tabularx}{0.7\textwidth}{X} 於是，約書亞清早起來，召以色列按著支派近前來。選出來的是猶大支派。 \end{tabularx} \\ \\ \relax
7:17 & \begin{tabularx}{0.7\textwidth}{X} 他召猶大的宗族近前來，選出來的是謝拉宗族。他召謝拉宗族，按著男丁，一個一個近前來，選出來的是撒底。 \end{tabularx} \\ \\ \relax
7:18 & \begin{tabularx}{0.7\textwidth}{X} 他召撒底的家族，按著男丁，一個一個近前來，就選出猶大支派，謝拉的曾孫，撒底的孫子，迦米的兒子亞干。 \end{tabularx} \\ \\ \relax
7:19 & \begin{tabularx}{0.7\textwidth}{X} 約書亞對亞干說：「我兒，我勸你將榮耀歸給耶和華－以色列的神，在他面前認罪，把你所做的事告訴我，不可向我隱瞞。」 \end{tabularx} \\ \\ \relax
7:20 & \begin{tabularx}{0.7\textwidth}{X} 亞干回答約書亞說：「我實在得罪了耶和華－以色列的神。這是我所做的： \end{tabularx} \\ \\ \relax
7:21 & \begin{tabularx}{0.7\textwidth}{X} 我在所奪取的財物中看見一件美好的示拿外袍，二百舍客勒銀子，一條重五十舍客勒的金子。我貪愛這些物件，就拿去了。看哪，這些東西都埋在我帳棚內的地裡，銀子在外袍底下。」 \end{tabularx} \\ \\ \relax
7:22 & \begin{tabularx}{0.7\textwidth}{X} 約書亞就派使者跑到亞干的帳棚裡。看哪，那件外袍藏在他的帳棚裡，銀子在外袍底下。 \end{tabularx} \\ \\ \relax
7:23 & \begin{tabularx}{0.7\textwidth}{X} 他們從帳棚裡把這些東西取出來，拿到約書亞和以色列眾人那裡，倒在耶和華面前。 \end{tabularx} \\ \\ \relax
7:24 & \begin{tabularx}{0.7\textwidth}{X} 約書亞和以色列眾人把謝拉的曾孫亞干和那銀子、那件外袍、那條金子，以及亞干的兒女、牛、驢、羊、帳棚，和他所有的，都帶著上到亞割谷去。 \end{tabularx} \\ \\ \relax
7:25 & \begin{tabularx}{0.7\textwidth}{X} 約書亞說：「你為甚麼給我們招惹災禍呢？今日耶和華必使你遭受災禍。」於是以色列眾人用石頭打死他，用火焚燒他們，把石頭扔在其上。 \end{tabularx} \\ \\ \relax
7:26 & \begin{tabularx}{0.7\textwidth}{X} 眾人在亞干身上堆了一大堆石頭，直存到今日。於是耶和華轉意，不發他的烈怒。因此，那地方名叫亞割谷，直到今日。 \end{tabularx} \\ \\
[1ex]
\hline
\hline
\end{longtable}


\newpage

\section{第75屆~港九培靈會~蕭壽華~第七講}
\label{sec:377}
\textbf{修築祭壇奉獻自己}
\newline
\newline
連結: \href{https://www.hkbibleconference.org/session-message/view/377}{\texttt{https://www.hkbibleconference.org/session-message/view/377}}
\newline
\newline
\hyperref[sec:374]{< < < PREV SERMON < < <}
~
\hyperref[sec:toc]{[返目錄]}
~
\hyperref[sec:380]{> > > NEXT SERMON > > >}
\newline
\newline
約書亞記 8:30-8:35
\newline
\begin{longtable}{cl}
\hline
\hline
章節 & 經文 (和合本修訂版)\\
\hline
8:30 & \begin{tabularx}{0.7\textwidth}{X} 那時，約書亞在以巴路山上為耶和華－以色列的神築一座壇。 \end{tabularx} \\ \\ \relax
8:31 & \begin{tabularx}{0.7\textwidth}{X} 這壇是照耶和華的僕人摩西吩咐以色列人，用沒有動過鐵器的整塊石頭所築的，正如摩西律法書上所寫的。他們在這壇上給耶和華奉獻燔祭，又宰牲作為平安祭。 \end{tabularx} \\ \\ \relax
8:32 & \begin{tabularx}{0.7\textwidth}{X} 約書亞在那裡，當著以色列人面前，將摩西所寫的律法抄寫在石頭上。 \end{tabularx} \\ \\ \relax
8:33 & \begin{tabularx}{0.7\textwidth}{X} 以色列眾人，無論是本地人或寄居的，都和他們的長老、官長和審判官，站在約櫃兩旁，在抬耶和華約櫃的利未家的祭司面前，一半對著基利心山，一半對著以巴路山，照耶和華的僕人摩西先前所吩咐的，為以色列百姓祝福。 \end{tabularx} \\ \\ \relax
8:34 & \begin{tabularx}{0.7\textwidth}{X} 隨後，約書亞將律法上祝福和詛咒的話，照著律法書上一切所寫的，宣讀一遍。 \end{tabularx} \\ \\ \relax
8:35 & \begin{tabularx}{0.7\textwidth}{X} 摩西所吩咐的一切話，約書亞在以色列全會眾和婦女、孩童，以及住在他們中間的外人面前，沒有一句不宣讀的。 \end{tabularx} \\ \\
[1ex]
\hline
\hline
\end{longtable}


\newpage

\section{第75屆~港九培靈會~蕭壽華~第八講}
\label{sec:380}
\textbf{專心信靠跟隨神 - 迦勒的見證}
\newline
\newline
連結: \href{https://www.hkbibleconference.org/session-message/view/380}{\texttt{https://www.hkbibleconference.org/session-message/view/380}}
\newline
\newline
\hyperref[sec:377]{< < < PREV SERMON < < <}
~
\hyperref[sec:toc]{[返目錄]}
~
\hyperref[sec:338]{> > > NEXT SERMON > > >}
\newline
\newline
約書亞記 14:6-14:15
\newline
\begin{longtable}{cl}
\hline
\hline
章節 & 經文 (和合本修訂版)\\
\hline
14:6 & \begin{tabularx}{0.7\textwidth}{X} 猶大人來到吉甲，約書亞那裡，基尼洗族耶孚尼的兒子迦勒對約書亞說：「耶和華在加低斯‧巴尼亞指著我和你對神人摩西所說的話，你都知道。 \end{tabularx} \\ \\ \relax
14:7 & \begin{tabularx}{0.7\textwidth}{X} 耶和華的僕人摩西從加低斯‧巴尼亞差派我窺探這地的時候，我剛四十歲。我把心裡的話向他報告。 \end{tabularx} \\ \\ \relax
14:8 & \begin{tabularx}{0.7\textwidth}{X} 雖然同我上去的眾弟兄使百姓膽戰心驚，我仍然專心跟從耶和華－我的神。 \end{tabularx} \\ \\ \relax
14:9 & \begin{tabularx}{0.7\textwidth}{X} 那日，摩西起誓說：『你腳所踏之地必要歸你和你的子孫永遠為業，因為你專心跟從耶和華－我的神。』 \end{tabularx} \\ \\ \relax
14:10 & \begin{tabularx}{0.7\textwidth}{X} 現在，看哪，耶和華照他所說的使我活了這四十五年。當以色列人在曠野飄流的時候，耶和華曾對摩西說了這話。現在，看哪，我已經八十五歲了。 \end{tabularx} \\ \\ \relax
14:11 & \begin{tabularx}{0.7\textwidth}{X} 現今我還很健壯，像摩西差派我去的那天一樣；無論是戰爭，是出入，我現在的力量和那時的力量一樣。 \end{tabularx} \\ \\ \relax
14:12 & \begin{tabularx}{0.7\textwidth}{X} 請你將耶和華那日所說的這山區給我。那日你也曾聽說，這裡有亞衲族人，以及寬大堅固的城，或許耶和華會照他所說的與我同在，我就把他們趕出去。」 \end{tabularx} \\ \\ \relax
14:13 & \begin{tabularx}{0.7\textwidth}{X} 於是約書亞為耶孚尼的兒子迦勒祝福，把希伯崙給他為業。 \end{tabularx} \\ \\ \relax
14:14 & \begin{tabularx}{0.7\textwidth}{X} 所以希伯崙成了基尼洗族耶孚尼的兒子迦勒的產業，直到今日，因為他專心跟從耶和華－以色列的神。 \end{tabularx} \\ \\ \relax
14:15 & \begin{tabularx}{0.7\textwidth}{X} 希伯崙從前名叫基列‧亞巴；亞巴是亞衲族最尊貴的人。於是國中太平，沒有戰爭了。 \end{tabularx} \\ \\
[1ex]
\hline
\hline
\end{longtable}


\newpage

\section{第76屆~港九培靈會~李思敬~第一講}
\label{sec:338}
\textbf{入門與定向}
\newline
\newline
連結: \href{https://www.hkbibleconference.org/session-message/view/338}{\texttt{https://www.hkbibleconference.org/session-message/view/338}}
\newline
\newline
\hyperref[sec:380]{< < < PREV SERMON < < <}
~
\hyperref[sec:toc]{[返目錄]}
~
\hyperref[sec:339]{> > > NEXT SERMON > > >}
\newline
\newline
箴言 1:1-1:7
\newline
\begin{longtable}{cl}
\hline
\hline
章節 & 經文 (和合本修訂版)\\
\hline
1:1 & \begin{tabularx}{0.7\textwidth}{X} 大衛的兒子，以色列王所羅門的箴言： \end{tabularx} \\ \\ \relax
1:2 & \begin{tabularx}{0.7\textwidth}{X} 要使人懂得智慧和訓誨，明白通達的言語， \end{tabularx} \\ \\ \relax
1:3 & \begin{tabularx}{0.7\textwidth}{X} 使人領受明智的訓誨，就是公義、公平和正直， \end{tabularx} \\ \\ \relax
1:4 & \begin{tabularx}{0.7\textwidth}{X} 使愚蒙人靈巧，使年輕人有知識，有智謀。 \end{tabularx} \\ \\ \relax
1:5 & \begin{tabularx}{0.7\textwidth}{X} 智慧人聽見，增長學問，聰明人得著智謀， \end{tabularx} \\ \\ \relax
1:6 & \begin{tabularx}{0.7\textwidth}{X} 明白箴言和譬喻，懂得智慧人的言詞和謎語。 \end{tabularx} \\ \\ \relax
1:7 & \begin{tabularx}{0.7\textwidth}{X} 敬畏耶和華是知識的開端；愚妄人藐視智慧和訓誨。 \end{tabularx} \\ \\
[1ex]
\hline
\hline
\end{longtable}


\newpage

\section{第76屆~港九培靈會~李思敬~第二講}
\label{sec:339}
\textbf{動靜的抉擇}
\newline
\newline
連結: \href{https://www.hkbibleconference.org/session-message/view/339}{\texttt{https://www.hkbibleconference.org/session-message/view/339}}
\newline
\newline
\hyperref[sec:338]{< < < PREV SERMON < < <}
~
\hyperref[sec:toc]{[返目錄]}
~
\hyperref[sec:340]{> > > NEXT SERMON > > >}
\newline
\newline
箴言 1:8-1:33
\newline
\begin{longtable}{cl}
\hline
\hline
章節 & 經文 (和合本修訂版)\\
\hline
1:8 & \begin{tabularx}{0.7\textwidth}{X} 我兒啊，要聽你父親的訓誨，不可離棄你母親的教誨； \end{tabularx} \\ \\ \relax
1:9 & \begin{tabularx}{0.7\textwidth}{X} 因為這要作你頭上恩惠的華冠，作你頸上的項鏈。 \end{tabularx} \\ \\ \relax
1:10 & \begin{tabularx}{0.7\textwidth}{X} 我兒啊，罪人若引誘你，你不可隨從。 \end{tabularx} \\ \\ \relax
1:11 & \begin{tabularx}{0.7\textwidth}{X} 他們若說：「你與我們同去，我們要埋伏殺人流血，無故地潛藏，殺害無辜； \end{tabularx} \\ \\ \relax
1:12 & \begin{tabularx}{0.7\textwidth}{X} 我們好像陰間，把他們活活吞下，囫圇吞下，如吞下那下到地府的人； \end{tabularx} \\ \\ \relax
1:13 & \begin{tabularx}{0.7\textwidth}{X} 我們必得各樣寶物，將所奪來的裝滿房屋； \end{tabularx} \\ \\ \relax
1:14 & \begin{tabularx}{0.7\textwidth}{X} 你來與我們同夥，共用一個錢囊。」 \end{tabularx} \\ \\ \relax
1:15 & \begin{tabularx}{0.7\textwidth}{X} 我兒啊，不要與他們走同一道路，禁止你的腳走他們的路徑。 \end{tabularx} \\ \\ \relax
1:16 & \begin{tabularx}{0.7\textwidth}{X} 因為他們的腳奔跑行惡，他們急速殺人流血。 \end{tabularx} \\ \\ \relax
1:17 & \begin{tabularx}{0.7\textwidth}{X} 在飛鳥眼前張設網羅，一定會徒勞無功； \end{tabularx} \\ \\ \relax
1:18 & \begin{tabularx}{0.7\textwidth}{X} 同樣，他們埋伏，是自流己血，他們潛藏，是自害己命。 \end{tabularx} \\ \\ \relax
1:19 & \begin{tabularx}{0.7\textwidth}{X} 凡靠暴力歛財的，所行之路都是如此，這種念頭必奪去自己的生命。 \end{tabularx} \\ \\ \relax
1:20 & \begin{tabularx}{0.7\textwidth}{X} 智慧在街市上呼喊，在廣場上高聲吶喊， \end{tabularx} \\ \\ \relax
1:21 & \begin{tabularx}{0.7\textwidth}{X} 在熱鬧街頭呼叫，在城門口，在城中，發出言語，說： \end{tabularx} \\ \\ \relax
1:22 & \begin{tabularx}{0.7\textwidth}{X} 「你們無知的人喜愛無知，傲慢人喜歡傲慢，愚昧人恨惡知識，要到幾時呢？ \end{tabularx} \\ \\ \relax
1:23 & \begin{tabularx}{0.7\textwidth}{X} 你們當因我的責備回轉，我要將我的靈澆灌你們，將我的話指示你們。 \end{tabularx} \\ \\ \relax
1:24 & \begin{tabularx}{0.7\textwidth}{X} 因為我呼喚，你們不聽，我招手，無人理會。 \end{tabularx} \\ \\ \relax
1:25 & \begin{tabularx}{0.7\textwidth}{X} 你們忽視我一切的勸戒，拒聽我的責備。 \end{tabularx} \\ \\ \relax
1:26 & \begin{tabularx}{0.7\textwidth}{X} 你們遭難，我就發笑；驚恐臨到你們，驚恐如狂風來臨，災難好像暴風來到，急難痛苦臨到你們身上，我必嗤笑。 \end{tabularx} \\ \\ \relax
1:27 & \begin{tabularx}{0.7\textwidth}{X} 見上節 \end{tabularx} \\ \\ \relax
1:28 & \begin{tabularx}{0.7\textwidth}{X} 那時，他們就會呼求我，我卻不回答，懇切尋求我，卻尋不見。 \end{tabularx} \\ \\ \relax
1:29 & \begin{tabularx}{0.7\textwidth}{X} 因為他們恨惡知識，選擇不敬畏耶和華， \end{tabularx} \\ \\ \relax
1:30 & \begin{tabularx}{0.7\textwidth}{X} 不聽我的勸戒，藐視我一切的責備， \end{tabularx} \\ \\ \relax
1:31 & \begin{tabularx}{0.7\textwidth}{X} 所以他們要自食其果，飽脹在自己的計謀中。 \end{tabularx} \\ \\ \relax
1:32 & \begin{tabularx}{0.7\textwidth}{X} 愚蒙人背道，害死自己，愚昧人安逸，自取滅亡。 \end{tabularx} \\ \\ \relax
1:33 & \begin{tabularx}{0.7\textwidth}{X} 惟聽從我的，必安然居住，得享寧靜，不怕災禍。」 \end{tabularx} \\ \\
[1ex]
\hline
\hline
\end{longtable}


\newpage

\section{第76屆~港九培靈會~李思敬~第三講}
\label{sec:340}
\textbf{敬虔還須正直}
\newline
\newline
連結: \href{https://www.hkbibleconference.org/session-message/view/340}{\texttt{https://www.hkbibleconference.org/session-message/view/340}}
\newline
\newline
\hyperref[sec:339]{< < < PREV SERMON < < <}
~
\hyperref[sec:toc]{[返目錄]}
~
\hyperref[sec:341]{> > > NEXT SERMON > > >}
\newline
\newline
箴言 2:1-2:22
\newline
\begin{longtable}{cl}
\hline
\hline
章節 & 經文 (和合本修訂版)\\
\hline
2:1 & \begin{tabularx}{0.7\textwidth}{X} 我兒啊，你若領受我的言語，珍藏我的命令， \end{tabularx} \\ \\ \relax
2:2 & \begin{tabularx}{0.7\textwidth}{X} 留心聽智慧，專心求聰明； \end{tabularx} \\ \\ \relax
2:3 & \begin{tabularx}{0.7\textwidth}{X} 你若呼求明理，揚聲求聰明， \end{tabularx} \\ \\ \relax
2:4 & \begin{tabularx}{0.7\textwidth}{X} 尋找她，如尋找銀子，搜尋她，如搜尋寶藏， \end{tabularx} \\ \\ \relax
2:5 & \begin{tabularx}{0.7\textwidth}{X} 你就懂得敬畏耶和華，得以認識神。 \end{tabularx} \\ \\ \relax
2:6 & \begin{tabularx}{0.7\textwidth}{X} 因為耶和華賞賜智慧，知識和聰明都由他口而出。 \end{tabularx} \\ \\ \relax
2:7 & \begin{tabularx}{0.7\textwidth}{X} 他為正直人珍藏健全的知識，給行為純正的人作盾牌， \end{tabularx} \\ \\ \relax
2:8 & \begin{tabularx}{0.7\textwidth}{X} 為要保護公正的路，庇護虔誠人的道。 \end{tabularx} \\ \\ \relax
2:9 & \begin{tabularx}{0.7\textwidth}{X} 那時，你就明白公義、公平、正直，和一切完善的道路。 \end{tabularx} \\ \\ \relax
2:10 & \begin{tabularx}{0.7\textwidth}{X} 因為智慧要進入你的心，知識要使你內心歡愉。 \end{tabularx} \\ \\ \relax
2:11 & \begin{tabularx}{0.7\textwidth}{X} 智謀要庇護你，聰明必保護你， \end{tabularx} \\ \\ \relax
2:12 & \begin{tabularx}{0.7\textwidth}{X} 救你脫離惡人的道，脫離言談乖謬的人。 \end{tabularx} \\ \\ \relax
2:13 & \begin{tabularx}{0.7\textwidth}{X} 他們離棄正直的路，行走黑暗的道， \end{tabularx} \\ \\ \relax
2:14 & \begin{tabularx}{0.7\textwidth}{X} 喜歡作惡，喜愛惡人的錯謬。 \end{tabularx} \\ \\ \relax
2:15 & \begin{tabularx}{0.7\textwidth}{X} 他們的路歪曲，他們偏離中道。 \end{tabularx} \\ \\ \relax
2:16 & \begin{tabularx}{0.7\textwidth}{X} 智慧要救你遠離陌生女子，遠離那油嘴滑舌的外邦女子。 \end{tabularx} \\ \\ \relax
2:17 & \begin{tabularx}{0.7\textwidth}{X} 她離棄年輕時的配偶，忘了自己神聖的盟約。 \end{tabularx} \\ \\ \relax
2:18 & \begin{tabularx}{0.7\textwidth}{X} 她的家陷入死亡，她的路偏向陰魂。 \end{tabularx} \\ \\ \relax
2:19 & \begin{tabularx}{0.7\textwidth}{X} 凡到她那裡去的，不得回轉，也得不到生命的路。 \end{tabularx} \\ \\ \relax
2:20 & \begin{tabularx}{0.7\textwidth}{X} 智慧使你行善人的道，守義人的路。 \end{tabularx} \\ \\ \relax
2:21 & \begin{tabularx}{0.7\textwidth}{X} 正直人必在地上居住，完全人必在其上存留； \end{tabularx} \\ \\ \relax
2:22 & \begin{tabularx}{0.7\textwidth}{X} 惟惡人要從地上剪除，奸詐人要被拔出。 \end{tabularx} \\ \\
[1ex]
\hline
\hline
\end{longtable}


\newpage

\section{第76屆~港九培靈會~李思敬~第四講}
\label{sec:341}
\textbf{說「不」的勇氣}
\newline
\newline
連結: \href{https://www.hkbibleconference.org/session-message/view/341}{\texttt{https://www.hkbibleconference.org/session-message/view/341}}
\newline
\newline
\hyperref[sec:340]{< < < PREV SERMON < < <}
~
\hyperref[sec:toc]{[返目錄]}
~
\hyperref[sec:342]{> > > NEXT SERMON > > >}
\newline
\newline
箴言 3:1-3:35
\newline
\begin{longtable}{cl}
\hline
\hline
章節 & 經文 (和合本修訂版)\\
\hline
3:1 & \begin{tabularx}{0.7\textwidth}{X} 我兒啊，不要忘記我的教誨，你的心要謹守我的命令， \end{tabularx} \\ \\ \relax
3:2 & \begin{tabularx}{0.7\textwidth}{X} 因為它們必加給你長久的日子，生命的年數與平安。 \end{tabularx} \\ \\ \relax
3:3 & \begin{tabularx}{0.7\textwidth}{X} 不可使慈愛和誠信離開你，要繫在你頸項上，刻在你心版上。 \end{tabularx} \\ \\ \relax
3:4 & \begin{tabularx}{0.7\textwidth}{X} 這樣，你必在神和世人眼前蒙恩惠，有美好的見識。 \end{tabularx} \\ \\ \relax
3:5 & \begin{tabularx}{0.7\textwidth}{X} 你要專心仰賴耶和華，不可倚靠自己的聰明， \end{tabularx} \\ \\ \relax
3:6 & \begin{tabularx}{0.7\textwidth}{X} 在你一切所行的路上都要認定他，他必使你的道路平直。 \end{tabularx} \\ \\ \relax
3:7 & \begin{tabularx}{0.7\textwidth}{X} 不要自以為有智慧；要敬畏耶和華，遠離惡事。 \end{tabularx} \\ \\ \relax
3:8 & \begin{tabularx}{0.7\textwidth}{X} 這便醫治你的肉體，滋潤你的百骨。 \end{tabularx} \\ \\ \relax
3:9 & \begin{tabularx}{0.7\textwidth}{X} 你要以財物和一切初熟的土產尊崇耶和華， \end{tabularx} \\ \\ \relax
3:10 & \begin{tabularx}{0.7\textwidth}{X} 這樣，你的倉庫必充滿有餘，你的酒池有新酒盈溢。 \end{tabularx} \\ \\ \relax
3:11 & \begin{tabularx}{0.7\textwidth}{X} 我兒啊，不可輕看耶和華的管教，也不可厭煩他的責備， \end{tabularx} \\ \\ \relax
3:12 & \begin{tabularx}{0.7\textwidth}{X} 因為耶和華所愛的，他必責備，正如父親責備所喜愛的兒子。 \end{tabularx} \\ \\ \relax
3:13 & \begin{tabularx}{0.7\textwidth}{X} 得智慧，得聰明的，這人有福了。 \end{tabularx} \\ \\ \relax
3:14 & \begin{tabularx}{0.7\textwidth}{X} 因為智慧的獲利勝過銀子，所得的盈餘強如金子， \end{tabularx} \\ \\ \relax
3:15 & \begin{tabularx}{0.7\textwidth}{X} 比寶石更寶貴，你一切所喜愛的，都不足與其比較。 \end{tabularx} \\ \\ \relax
3:16 & \begin{tabularx}{0.7\textwidth}{X} 她的右手有長壽，左手有富貴。 \end{tabularx} \\ \\ \relax
3:17 & \begin{tabularx}{0.7\textwidth}{X} 她的道是安樂，她的路全是平安。 \end{tabularx} \\ \\ \relax
3:18 & \begin{tabularx}{0.7\textwidth}{X} 她給持守她的人作生命樹，謹守她的必定蒙福。 \end{tabularx} \\ \\ \relax
3:19 & \begin{tabularx}{0.7\textwidth}{X} 耶和華以智慧奠立地基，以聰明鋪設諸天， \end{tabularx} \\ \\ \relax
3:20 & \begin{tabularx}{0.7\textwidth}{X} 以知識使深淵裂開，使天空滴下甘露。 \end{tabularx} \\ \\ \relax
3:21 & \begin{tabularx}{0.7\textwidth}{X} 我兒啊，要謹守健全的知識和智謀，不可使它們偏離你的眼目。 \end{tabularx} \\ \\ \relax
3:22 & \begin{tabularx}{0.7\textwidth}{X} 這樣，它們必使你的生命有活力，又作你頸項的美飾。 \end{tabularx} \\ \\ \relax
3:23 & \begin{tabularx}{0.7\textwidth}{X} 那時，你就坦然行路，不致跌倒。 \end{tabularx} \\ \\ \relax
3:24 & \begin{tabularx}{0.7\textwidth}{X} 你躺下，必不懼怕；你躺臥，睡得香甜。 \end{tabularx} \\ \\ \relax
3:25 & \begin{tabularx}{0.7\textwidth}{X} 忽然來的驚恐，你不要害怕；惡人遭毀滅，也不要恐懼， \end{tabularx} \\ \\ \relax
3:26 & \begin{tabularx}{0.7\textwidth}{X} 因為耶和華是你的倚靠，他必保護你的腳不陷入羅網。 \end{tabularx} \\ \\ \relax
3:27 & \begin{tabularx}{0.7\textwidth}{X} 你的手若有行善的力量，不可推辭，要施與那應得的人。 \end{tabularx} \\ \\ \relax
3:28 & \begin{tabularx}{0.7\textwidth}{X} 你若手頭方便，不可對鄰舍說：「去吧，明天再來，我必給你。」 \end{tabularx} \\ \\ \relax
3:29 & \begin{tabularx}{0.7\textwidth}{X} 你的鄰舍既在你附近安居，不可設計害他。 \end{tabularx} \\ \\ \relax
3:30 & \begin{tabularx}{0.7\textwidth}{X} 人若未曾加害你，不可無故與他相爭。 \end{tabularx} \\ \\ \relax
3:31 & \begin{tabularx}{0.7\textwidth}{X} 不可嫉妒殘暴的人，不可選擇他的任何道路。 \end{tabularx} \\ \\ \relax
3:32 & \begin{tabularx}{0.7\textwidth}{X} 因為走偏方向的人是耶和華所憎惡的；正直人為他所親密。 \end{tabularx} \\ \\ \relax
3:33 & \begin{tabularx}{0.7\textwidth}{X} 耶和華詛咒惡人的家；義人的居所他卻賜福。 \end{tabularx} \\ \\ \relax
3:34 & \begin{tabularx}{0.7\textwidth}{X} 他譏誚那愛譏誚的人；但賜恩給謙卑的人。 \end{tabularx} \\ \\ \relax
3:35 & \begin{tabularx}{0.7\textwidth}{X} 智慧人必承受尊榮；愚昧人高升卻是羞辱。 \end{tabularx} \\ \\
[1ex]
\hline
\hline
\end{longtable}


\newpage

\section{第76屆~港九培靈會~李思敬~第五講}
\label{sec:342}
\textbf{小心思想防線}
\newline
\newline
連結: \href{https://www.hkbibleconference.org/session-message/view/342}{\texttt{https://www.hkbibleconference.org/session-message/view/342}}
\newline
\newline
\hyperref[sec:341]{< < < PREV SERMON < < <}
~
\hyperref[sec:toc]{[返目錄]}
~
\hyperref[sec:343]{> > > NEXT SERMON > > >}
\newline
\newline
箴言 4:1-5:6
\newline
\begin{longtable}{cl}
\hline
\hline
章節 & 經文 (和合本修訂版)\\
\hline
4:1 & \begin{tabularx}{0.7\textwidth}{X} 孩子們，要聽父親的訓誨，留心明白道理。 \end{tabularx} \\ \\ \relax
4:2 & \begin{tabularx}{0.7\textwidth}{X} 因我給你們好的教導，不可離棄我的教誨。 \end{tabularx} \\ \\ \relax
4:3 & \begin{tabularx}{0.7\textwidth}{X} 當我在父親面前還是小孩，是母親獨一嬌兒的時候， \end{tabularx} \\ \\ \relax
4:4 & \begin{tabularx}{0.7\textwidth}{X} 他教導我說：「你的心要持守我的話，遵守我的命令，你就會存活。 \end{tabularx} \\ \\ \relax
4:5 & \begin{tabularx}{0.7\textwidth}{X} 要獲得智慧，要獲得聰明，不可忘記， 也不可偏離我口中的言語。 \end{tabularx} \\ \\ \relax
4:6 & \begin{tabularx}{0.7\textwidth}{X} 不可離棄智慧，智慧就庇護你，要愛她，她就保護你。 \end{tabularx} \\ \\ \relax
4:7 & \begin{tabularx}{0.7\textwidth}{X} 智慧為首，所以要獲得智慧，要用你一切所有的換取聰明。 \end{tabularx} \\ \\ \relax
4:8 & \begin{tabularx}{0.7\textwidth}{X} 高舉智慧，她就使你升高，擁抱智慧，她就使你尊榮。 \end{tabularx} \\ \\ \relax
4:9 & \begin{tabularx}{0.7\textwidth}{X} 她必將恩惠的華冠加在你頭上，把榮冕賜給你。」 \end{tabularx} \\ \\ \relax
4:10 & \begin{tabularx}{0.7\textwidth}{X} 我兒啊，要聽，要領受我的言語，你就必延年益壽。 \end{tabularx} \\ \\ \relax
4:11 & \begin{tabularx}{0.7\textwidth}{X} 我已指教你走智慧的道，引導你行正直的路。 \end{tabularx} \\ \\ \relax
4:12 & \begin{tabularx}{0.7\textwidth}{X} 你行走，腳步沒有阻礙；你奔跑，也不致跌倒。 \end{tabularx} \\ \\ \relax
4:13 & \begin{tabularx}{0.7\textwidth}{X} 要持定訓誨，不可放鬆；要謹守它，因為它是你的生命。 \end{tabularx} \\ \\ \relax
4:14 & \begin{tabularx}{0.7\textwidth}{X} 不可行惡人的路，不要走壞人的道； \end{tabularx} \\ \\ \relax
4:15 & \begin{tabularx}{0.7\textwidth}{X} 要躲避，不可經過，要轉離而去。 \end{tabularx} \\ \\ \relax
4:16 & \begin{tabularx}{0.7\textwidth}{X} 他們若不行惡，難以成眠，不使人跌倒，就睡臥不安； \end{tabularx} \\ \\ \relax
4:17 & \begin{tabularx}{0.7\textwidth}{X} 因為他們以邪惡當餅吃，以暴力當酒喝。 \end{tabularx} \\ \\ \relax
4:18 & \begin{tabularx}{0.7\textwidth}{X} 但義人的路好像黎明的光，越照越明，直到正午。 \end{tabularx} \\ \\ \relax
4:19 & \begin{tabularx}{0.7\textwidth}{X} 惡人的道幽暗，自己不知因何跌倒。 \end{tabularx} \\ \\ \relax
4:20 & \begin{tabularx}{0.7\textwidth}{X} 我兒啊，要留心聽我的話，側耳聽我的言語， \end{tabularx} \\ \\ \relax
4:21 & \begin{tabularx}{0.7\textwidth}{X} 不可使它們偏離你的眼目，要存記在你心中。 \end{tabularx} \\ \\ \relax
4:22 & \begin{tabularx}{0.7\textwidth}{X} 因為找到它們的，就找到生命，得到全身的醫治。 \end{tabularx} \\ \\ \relax
4:23 & \begin{tabularx}{0.7\textwidth}{X} 你要保守你心，勝過保守一切，因為生命的泉源由心發出。 \end{tabularx} \\ \\ \relax
4:24 & \begin{tabularx}{0.7\textwidth}{X} 要離開歪曲的口，轉離偏邪的嘴唇。 \end{tabularx} \\ \\ \relax
4:25 & \begin{tabularx}{0.7\textwidth}{X} 你的兩眼要向前看，你的雙目直視前方。 \end{tabularx} \\ \\ \relax
4:26 & \begin{tabularx}{0.7\textwidth}{X} 要修平你腳下的路，你一切的道就必穩固。 \end{tabularx} \\ \\ \relax
4:27 & \begin{tabularx}{0.7\textwidth}{X} 不可偏左偏右，你的腳要離開邪惡。 \end{tabularx} \\ \\ \relax
5:1 & \begin{tabularx}{0.7\textwidth}{X} 我兒啊，要留心聽我的智慧，側耳聽我的聰明， \end{tabularx} \\ \\ \relax
5:2 & \begin{tabularx}{0.7\textwidth}{X} 為要使你謹守智謀，嘴唇保護知識。 \end{tabularx} \\ \\ \relax
5:3 & \begin{tabularx}{0.7\textwidth}{X} 因為陌生女子的嘴唇滴下蜂蜜，她的口比油更滑， \end{tabularx} \\ \\ \relax
5:4 & \begin{tabularx}{0.7\textwidth}{X} 後來卻苦似茵蔯，銳利如兩刃的劍。 \end{tabularx} \\ \\ \relax
5:5 & \begin{tabularx}{0.7\textwidth}{X} 她的腳墜落死亡，她的腳步踏入陰間， \end{tabularx} \\ \\ \relax
5:6 & \begin{tabularx}{0.7\textwidth}{X} 她無法找到生命的道路，她的路變遷不定，自己卻不知道。 \end{tabularx} \\ \\
[1ex]
\hline
\hline
\end{longtable}


\newpage

\section{第76屆~港九培靈會~李思敬~第六講}
\label{sec:343}
\textbf{離惡莫猶豫}
\newline
\newline
連結: \href{https://www.hkbibleconference.org/session-message/view/343}{\texttt{https://www.hkbibleconference.org/session-message/view/343}}
\newline
\newline
\hyperref[sec:342]{< < < PREV SERMON < < <}
~
\hyperref[sec:toc]{[返目錄]}
~
\hyperref[sec:344]{> > > NEXT SERMON > > >}
\newline
\newline
箴言 5:7-6:11
\newline
\begin{longtable}{cl}
\hline
\hline
章節 & 經文 (和合本修訂版)\\
\hline
5:7 & \begin{tabularx}{0.7\textwidth}{X} 孩子們，現在要聽從我，不可離棄我口中的言語。 \end{tabularx} \\ \\ \relax
5:8 & \begin{tabularx}{0.7\textwidth}{X} 你所行的道要遠離她，不可靠近她家的門口， \end{tabularx} \\ \\ \relax
5:9 & \begin{tabularx}{0.7\textwidth}{X} 免得將你的尊榮給別人，將你的歲月給殘忍的人； \end{tabularx} \\ \\ \relax
5:10 & \begin{tabularx}{0.7\textwidth}{X} 免得陌生人滿得你的財富，你勞苦所得的歸入外邦人的家。 \end{tabularx} \\ \\ \relax
5:11 & \begin{tabularx}{0.7\textwidth}{X} 在你人生終結，你皮肉和身體衰殘時，你必唉聲嘆氣， \end{tabularx} \\ \\ \relax
5:12 & \begin{tabularx}{0.7\textwidth}{X} 說：「我為何恨惡管教，心裡輕看責備呢？ \end{tabularx} \\ \\ \relax
5:13 & \begin{tabularx}{0.7\textwidth}{X} 我不聽從教師的話，也沒有側耳聽那教導我的。 \end{tabularx} \\ \\ \relax
5:14 & \begin{tabularx}{0.7\textwidth}{X} 在聚集的會眾中，我幾乎墜入深淵。」 \end{tabularx} \\ \\ \relax
5:15 & \begin{tabularx}{0.7\textwidth}{X} 你要喝自己池中的水，飲自己井裡的活水。 \end{tabularx} \\ \\ \relax
5:16 & \begin{tabularx}{0.7\textwidth}{X} 你的泉源豈可溢流在外？你的河水豈可流到街上？ \end{tabularx} \\ \\ \relax
5:17 & \begin{tabularx}{0.7\textwidth}{X} 讓它們惟獨歸你，不可與陌生人同享。 \end{tabularx} \\ \\ \relax
5:18 & \begin{tabularx}{0.7\textwidth}{X} 要使你的泉源蒙福，要喜愛你年輕時的妻子。 \end{tabularx} \\ \\ \relax
5:19 & \begin{tabularx}{0.7\textwidth}{X} 她如可愛的母鹿，如優美的母羊，願她的胸懷使你時時滿足，願你常常迷戀她的愛情。 \end{tabularx} \\ \\ \relax
5:20 & \begin{tabularx}{0.7\textwidth}{X} 我兒啊，你為何迷戀陌生女子？為何擁抱外邦女子的胸懷？ \end{tabularx} \\ \\ \relax
5:21 & \begin{tabularx}{0.7\textwidth}{X} 因為人所行的道都在耶和華眼前，他察驗人一切的路。 \end{tabularx} \\ \\ \relax
5:22 & \begin{tabularx}{0.7\textwidth}{X} 惡人被自己的罪孽抓住，被自己罪惡的繩索纏繞。 \end{tabularx} \\ \\ \relax
5:23 & \begin{tabularx}{0.7\textwidth}{X} 他因不受管教而死亡，因極度愚昧而走迷。 \end{tabularx} \\ \\ \relax
6:1 & \begin{tabularx}{0.7\textwidth}{X} 我兒啊，你若為朋友擔保，替陌生人擊掌， \end{tabularx} \\ \\ \relax
6:2 & \begin{tabularx}{0.7\textwidth}{X} 你就被口中的言語套住，被嘴裡的言語抓住。 \end{tabularx} \\ \\ \relax
6:3 & \begin{tabularx}{0.7\textwidth}{X} 我兒啊，你既落在朋友手中，當這樣行才可救自己：你要謙卑自己，去懇求你的朋友。 \end{tabularx} \\ \\ \relax
6:4 & \begin{tabularx}{0.7\textwidth}{X} 不要讓你的眼睛睡覺，不可容你的眼皮打盹。 \end{tabularx} \\ \\ \relax
6:5 & \begin{tabularx}{0.7\textwidth}{X} 要救自己，如羚羊脫離獵人的手，如鳥脫離捕鳥人的手。 \end{tabularx} \\ \\ \relax
6:6 & \begin{tabularx}{0.7\textwidth}{X} 懶惰人哪，你去察看螞蟻的動作，就可得智慧。 \end{tabularx} \\ \\ \relax
6:7 & \begin{tabularx}{0.7\textwidth}{X} 螞蟻沒有領袖，沒有官長，沒有君王， \end{tabularx} \\ \\ \relax
6:8 & \begin{tabularx}{0.7\textwidth}{X} 尚且在夏天預備食物，在收割時儲存糧食。 \end{tabularx} \\ \\ \relax
6:9 & \begin{tabularx}{0.7\textwidth}{X} 懶惰人哪，你要睡到幾時呢？你甚麼時候才睡醒呢？ \end{tabularx} \\ \\ \relax
6:10 & \begin{tabularx}{0.7\textwidth}{X} 再睡片時，打盹片時，抱著雙臂躺臥片時， \end{tabularx} \\ \\ \relax
6:11 & \begin{tabularx}{0.7\textwidth}{X} 你的貧窮就如盜賊來到，你的貧乏彷彿拿盾牌的人來臨。 \end{tabularx} \\ \\
[1ex]
\hline
\hline
\end{longtable}


\newpage

\section{第76屆~港九培靈會~李思敬~第七講}
\label{sec:344}
\textbf{還姦淫真面目}
\newline
\newline
連結: \href{https://www.hkbibleconference.org/session-message/view/344}{\texttt{https://www.hkbibleconference.org/session-message/view/344}}
\newline
\newline
\hyperref[sec:343]{< < < PREV SERMON < < <}
~
\hyperref[sec:toc]{[返目錄]}
~
\hyperref[sec:345]{> > > NEXT SERMON > > >}
\newline
\newline
箴言 6:12-7:23
\newline
\begin{longtable}{cl}
\hline
\hline
章節 & 經文 (和合本修訂版)\\
\hline
6:12 & \begin{tabularx}{0.7\textwidth}{X} 無賴的惡徒行事全憑歪曲的口， \end{tabularx} \\ \\ \relax
6:13 & \begin{tabularx}{0.7\textwidth}{X} 他眨眼傳神，以腳示意，用指點劃， \end{tabularx} \\ \\ \relax
6:14 & \begin{tabularx}{0.7\textwidth}{X} 存心乖謬，常設惡謀，散播紛爭。 \end{tabularx} \\ \\ \relax
6:15 & \begin{tabularx}{0.7\textwidth}{X} 所以，災難必突然臨到他，他必頃刻被毀，無從醫治。 \end{tabularx} \\ \\ \relax
6:16 & \begin{tabularx}{0.7\textwidth}{X} 耶和華所恨惡的有六樣，他心所憎惡的共有七樣： \end{tabularx} \\ \\ \relax
6:17 & \begin{tabularx}{0.7\textwidth}{X} 就是高傲的眼，撒謊的舌，殺害無辜的手， \end{tabularx} \\ \\ \relax
6:18 & \begin{tabularx}{0.7\textwidth}{X} 圖謀惡計的心，飛奔行惡的腳， \end{tabularx} \\ \\ \relax
6:19 & \begin{tabularx}{0.7\textwidth}{X} 口吐謊言的假證人，並在弟兄間散播紛爭的人。 \end{tabularx} \\ \\ \relax
6:20 & \begin{tabularx}{0.7\textwidth}{X} 我兒啊，要遵守你父親的命令，不可離棄你母親的教誨。 \end{tabularx} \\ \\ \relax
6:21 & \begin{tabularx}{0.7\textwidth}{X} 要常掛在你心上，繫在你頸項上。 \end{tabularx} \\ \\ \relax
6:22 & \begin{tabularx}{0.7\textwidth}{X} 你行走，她必引導你，你躺臥，她必保護你，你睡醒，她必與你談論。 \end{tabularx} \\ \\ \relax
6:23 & \begin{tabularx}{0.7\textwidth}{X} 因為誡命是燈，教誨是光，管教的責備是生命的道， \end{tabularx} \\ \\ \relax
6:24 & \begin{tabularx}{0.7\textwidth}{X} 要保護你遠離邪惡的婦女，遠離外邦女子諂媚的舌頭。 \end{tabularx} \\ \\ \relax
6:25 & \begin{tabularx}{0.7\textwidth}{X} 你不要因她的美色而動心，也不要被她的眼皮勾引。 \end{tabularx} \\ \\ \relax
6:26 & \begin{tabularx}{0.7\textwidth}{X} 因為連最後一塊餅都會被妓女拿走；有夫之婦會獵取寶貴的生命。 \end{tabularx} \\ \\ \relax
6:27 & \begin{tabularx}{0.7\textwidth}{X} 人若兜火在懷中，他的衣服豈能不燒著呢？ \end{tabularx} \\ \\ \relax
6:28 & \begin{tabularx}{0.7\textwidth}{X} 人若走在火炭上，他的腳豈能不燙傷呢？ \end{tabularx} \\ \\ \relax
6:29 & \begin{tabularx}{0.7\textwidth}{X} 與鄰舍之妻同寢的，也是如此，凡親近她的，難免受罰。 \end{tabularx} \\ \\ \relax
6:30 & \begin{tabularx}{0.7\textwidth}{X} 賊因飢餓偷竊充飢，人不藐視他， \end{tabularx} \\ \\ \relax
6:31 & \begin{tabularx}{0.7\textwidth}{X} 但若被抓到，要賠償七倍，他必賠上家中一切財物。 \end{tabularx} \\ \\ \relax
6:32 & \begin{tabularx}{0.7\textwidth}{X} 與婦人行姦淫的，便是無知，做這事的，必毀了自己。 \end{tabularx} \\ \\ \relax
6:33 & \begin{tabularx}{0.7\textwidth}{X} 他必受損傷和羞辱，他的羞恥不得消除。 \end{tabularx} \\ \\ \relax
6:34 & \begin{tabularx}{0.7\textwidth}{X} 丈夫因嫉恨發怒，報仇的時候絕不留情。 \end{tabularx} \\ \\ \relax
6:35 & \begin{tabularx}{0.7\textwidth}{X} 他不接受任何賠償，你送許多禮物，他也不肯和解。 \end{tabularx} \\ \\ \relax
7:1 & \begin{tabularx}{0.7\textwidth}{X} 我兒啊，要遵守我的言語，存記我的命令。 \end{tabularx} \\ \\ \relax
7:2 & \begin{tabularx}{0.7\textwidth}{X} 遵守我的命令就得存活，謹守我的教誨，好像保護眼中的瞳人。 \end{tabularx} \\ \\ \relax
7:3 & \begin{tabularx}{0.7\textwidth}{X} 要繫在你指頭上，刻在你心版上。 \end{tabularx} \\ \\ \relax
7:4 & \begin{tabularx}{0.7\textwidth}{X} 對智慧說「你是我的姊妹」，稱呼聰明為親人， \end{tabularx} \\ \\ \relax
7:5 & \begin{tabularx}{0.7\textwidth}{X} 她就保護你遠離陌生女子，遠離油嘴滑舌的外邦女子。 \end{tabularx} \\ \\ \relax
7:6 & \begin{tabularx}{0.7\textwidth}{X} 我曾在我房屋的窗戶內，透過窗格子往外觀看， \end{tabularx} \\ \\ \relax
7:7 & \begin{tabularx}{0.7\textwidth}{X} 看見在愚蒙人中，注意到孩兒中有一個無知的青年， \end{tabularx} \\ \\ \relax
7:8 & \begin{tabularx}{0.7\textwidth}{X} 從街上經過，靠近她的巷口，直往她家的路去， \end{tabularx} \\ \\ \relax
7:9 & \begin{tabularx}{0.7\textwidth}{X} 在黃昏，在傍晚，在半夜，黑暗之中。 \end{tabularx} \\ \\ \relax
7:10 & \begin{tabularx}{0.7\textwidth}{X} 看哪，有一個女子來迎接他，是妓女的打扮，有詭詐的心思。 \end{tabularx} \\ \\ \relax
7:11 & \begin{tabularx}{0.7\textwidth}{X} 她喧嚷，不守約束，她的腳在家裡留不住， \end{tabularx} \\ \\ \relax
7:12 & \begin{tabularx}{0.7\textwidth}{X} 有時在街市，有時在廣場，或在各巷口等候。 \end{tabularx} \\ \\ \relax
7:13 & \begin{tabularx}{0.7\textwidth}{X} 她拉住那青年吻他，厚著臉皮對他說： \end{tabularx} \\ \\ \relax
7:14 & \begin{tabularx}{0.7\textwidth}{X} 「我已獻了平安祭，今日我還了所許的願。 \end{tabularx} \\ \\ \relax
7:15 & \begin{tabularx}{0.7\textwidth}{X} 因此，我出來迎接你，渴望見你的面，我總算找到你了！ \end{tabularx} \\ \\ \relax
7:16 & \begin{tabularx}{0.7\textwidth}{X} 我已在床上鋪好被單，是埃及麻織的花紋布， \end{tabularx} \\ \\ \relax
7:17 & \begin{tabularx}{0.7\textwidth}{X} 又用沒藥、沉香、桂皮薰了我的床。 \end{tabularx} \\ \\ \relax
7:18 & \begin{tabularx}{0.7\textwidth}{X} 你來，讓我們飽享愛情，直到早晨，讓我們彼此親愛歡樂。 \end{tabularx} \\ \\ \relax
7:19 & \begin{tabularx}{0.7\textwidth}{X} 因為我丈夫不在家，出門遠行， \end{tabularx} \\ \\ \relax
7:20 & \begin{tabularx}{0.7\textwidth}{X} 他手帶錢囊，要到月圓才回家。」 \end{tabularx} \\ \\ \relax
7:21 & \begin{tabularx}{0.7\textwidth}{X} 這女子用許多巧言引誘他，用諂媚的嘴唇催逼他。 \end{tabularx} \\ \\ \relax
7:22 & \begin{tabularx}{0.7\textwidth}{X} 青年立刻跟隨她，好像牛去被宰殺，又像愚妄人帶著腳鐐去受刑， \end{tabularx} \\ \\ \relax
7:23 & \begin{tabularx}{0.7\textwidth}{X} 直到箭穿進他的肝，如同雀鳥急投羅網，卻不知會賠上自己的生命。 \end{tabularx} \\ \\
[1ex]
\hline
\hline
\end{longtable}


\newpage

\section{第76屆~港九培靈會~李思敬~第八講}
\label{sec:345}
\textbf{明辨生與死}
\newline
\newline
連結: \href{https://www.hkbibleconference.org/session-message/view/345}{\texttt{https://www.hkbibleconference.org/session-message/view/345}}
\newline
\newline
\hyperref[sec:344]{< < < PREV SERMON < < <}
~
\hyperref[sec:toc]{[返目錄]}
~
\hyperref[sec:346]{> > > NEXT SERMON > > >}
\newline
\newline
箴言 7:24-9:18
\newline
\begin{longtable}{cl}
\hline
\hline
章節 & 經文 (和合本修訂版)\\
\hline
7:24 & \begin{tabularx}{0.7\textwidth}{X} 孩子們，現在要聽從我，要留心聽我口中的言語。 \end{tabularx} \\ \\ \relax
7:25 & \begin{tabularx}{0.7\textwidth}{X} 你的心不可偏向她的道，不要誤入她的迷途。 \end{tabularx} \\ \\ \relax
7:26 & \begin{tabularx}{0.7\textwidth}{X} 因為她擊倒許多人，無數的人被她殺戮。 \end{tabularx} \\ \\ \relax
7:27 & \begin{tabularx}{0.7\textwidth}{X} 她的家是在陰間之路，下到死亡之宮。 \end{tabularx} \\ \\ \relax
8:1 & \begin{tabularx}{0.7\textwidth}{X} 智慧豈不呼喚？聰明豈不揚聲？ \end{tabularx} \\ \\ \relax
8:2 & \begin{tabularx}{0.7\textwidth}{X} 她站立在十字路口，在道路旁高處的頂上， \end{tabularx} \\ \\ \relax
8:3 & \begin{tabularx}{0.7\textwidth}{X} 在城門旁，城門口，入口處，她呼喊： \end{tabularx} \\ \\ \relax
8:4 & \begin{tabularx}{0.7\textwidth}{X} 「人哪，我呼喚你們，我向世人揚聲。 \end{tabularx} \\ \\ \relax
8:5 & \begin{tabularx}{0.7\textwidth}{X} 愚蒙人哪，你們要學習靈巧，愚昧人哪，你們的心要明辨。 \end{tabularx} \\ \\ \relax
8:6 & \begin{tabularx}{0.7\textwidth}{X} 你們當聽，因我要說尊貴的事，我要張開嘴唇講正直的事。 \end{tabularx} \\ \\ \relax
8:7 & \begin{tabularx}{0.7\textwidth}{X} 我的口要發出真理，我的嘴唇憎惡邪惡。 \end{tabularx} \\ \\ \relax
8:8 & \begin{tabularx}{0.7\textwidth}{X} 我口中的言語都是公義，並無奸詐和歪曲。 \end{tabularx} \\ \\ \relax
8:9 & \begin{tabularx}{0.7\textwidth}{X} 聰明人看為正確，有知識的，都以為正直。 \end{tabularx} \\ \\ \relax
8:10 & \begin{tabularx}{0.7\textwidth}{X} 你們當領受我的訓誨，勝過領受銀子，寧得知識，強如得上選的金子。 \end{tabularx} \\ \\ \relax
8:11 & \begin{tabularx}{0.7\textwidth}{X} 「因為智慧比寶石更美，一切可喜愛的都不足與其比較。 \end{tabularx} \\ \\ \relax
8:12 & \begin{tabularx}{0.7\textwidth}{X} 我－智慧以靈巧為居所，又尋得知識和智謀。 \end{tabularx} \\ \\ \relax
8:13 & \begin{tabularx}{0.7\textwidth}{X} 敬畏耶和華就是恨惡邪惡；我恨惡驕傲、狂妄、惡道，和乖謬的口。 \end{tabularx} \\ \\ \relax
8:14 & \begin{tabularx}{0.7\textwidth}{X} 我有策略和健全的知識，我聰明，又有能力。 \end{tabularx} \\ \\ \relax
8:15 & \begin{tabularx}{0.7\textwidth}{X} 君王藉我治國，王子藉我定公平， \end{tabularx} \\ \\ \relax
8:16 & \begin{tabularx}{0.7\textwidth}{X} 王公貴族，所有公義的審判官，都藉我掌權。 \end{tabularx} \\ \\ \relax
8:17 & \begin{tabularx}{0.7\textwidth}{X} 愛我的，我也愛他，懇切尋求我的，必尋見。 \end{tabularx} \\ \\ \relax
8:18 & \begin{tabularx}{0.7\textwidth}{X} 財富和尊榮在我，恆久的財寶和繁榮也在我。 \end{tabularx} \\ \\ \relax
8:19 & \begin{tabularx}{0.7\textwidth}{X} 我的果實勝過金子，強如純金，我的出產超乎上選的銀子。 \end{tabularx} \\ \\ \relax
8:20 & \begin{tabularx}{0.7\textwidth}{X} 我在公義的道上走，在公平的路中行， \end{tabularx} \\ \\ \relax
8:21 & \begin{tabularx}{0.7\textwidth}{X} 使愛我的承受財產，充滿他們的庫房。 \end{tabularx} \\ \\ \relax
8:22 & \begin{tabularx}{0.7\textwidth}{X} 「耶和華在造化的起頭，在太初創造萬物之先，就有了我。 \end{tabularx} \\ \\ \relax
8:23 & \begin{tabularx}{0.7\textwidth}{X} 從亙古，從太初，未有大地以前，我已被立。 \end{tabularx} \\ \\ \relax
8:24 & \begin{tabularx}{0.7\textwidth}{X} 沒有深淵，沒有大水的泉源，我已出生。 \end{tabularx} \\ \\ \relax
8:25 & \begin{tabularx}{0.7\textwidth}{X} 大山未曾奠定，小山未有之先，我已出生。 \end{tabularx} \\ \\ \relax
8:26 & \begin{tabularx}{0.7\textwidth}{X} 那時，他還沒有創造大地和田野，並世上頭一撮塵土。 \end{tabularx} \\ \\ \relax
8:27 & \begin{tabularx}{0.7\textwidth}{X} 他立高天，我在那裡，他在淵面的周圍劃出圓圈， \end{tabularx} \\ \\ \relax
8:28 & \begin{tabularx}{0.7\textwidth}{X} 上使穹蒼堅硬，下使淵源穩固， \end{tabularx} \\ \\ \relax
8:29 & \begin{tabularx}{0.7\textwidth}{X} 為滄海定出範圍，使水不越過界限，奠定大地的根基。 \end{tabularx} \\ \\ \relax
8:30 & \begin{tabularx}{0.7\textwidth}{X} 那時，我在他旁邊為工程師，天天充滿喜樂，時時在他面前歡笑， \end{tabularx} \\ \\ \relax
8:31 & \begin{tabularx}{0.7\textwidth}{X} 在他的全地歡笑，喜愛住在人世間。 \end{tabularx} \\ \\ \relax
8:32 & \begin{tabularx}{0.7\textwidth}{X} 「孩子們，現在要聽從我，謹守我道的有福了。 \end{tabularx} \\ \\ \relax
8:33 & \begin{tabularx}{0.7\textwidth}{X} 要聽訓誨，得智慧，不可棄絕。 \end{tabularx} \\ \\ \relax
8:34 & \begin{tabularx}{0.7\textwidth}{X} 聽從我，天天在我門口守望，在我門框旁等候的，那人有福了。 \end{tabularx} \\ \\ \relax
8:35 & \begin{tabularx}{0.7\textwidth}{X} 因為尋得我的，就尋得生命，他必蒙耶和華的恩惠。 \end{tabularx} \\ \\ \relax
8:36 & \begin{tabularx}{0.7\textwidth}{X} 得罪我的，害了自己的生命，凡恨惡我的，喜愛死亡。」 \end{tabularx} \\ \\ \relax
9:1 & \begin{tabularx}{0.7\textwidth}{X} 智慧建造房屋，鑿成七根柱子， \end{tabularx} \\ \\ \relax
9:2 & \begin{tabularx}{0.7\textwidth}{X} 宰殺牲畜，調好美酒，又擺設筵席， \end{tabularx} \\ \\ \relax
9:3 & \begin{tabularx}{0.7\textwidth}{X} 派遣女僕出去，自己在城中至高處呼喚： \end{tabularx} \\ \\ \relax
9:4 & \begin{tabularx}{0.7\textwidth}{X} 「誰是愚蒙的人，讓他轉到這裡來！」又對那無知的人說： \end{tabularx} \\ \\ \relax
9:5 & \begin{tabularx}{0.7\textwidth}{X} 「你們來，吃我的餅，喝我調的酒。 \end{tabularx} \\ \\ \relax
9:6 & \begin{tabularx}{0.7\textwidth}{X} 你們要離棄愚蒙，就得存活，並要走明智的道路。」 \end{tabularx} \\ \\ \relax
9:7 & \begin{tabularx}{0.7\textwidth}{X} 糾正傲慢人的，必招羞辱，責備惡人的，必被侮辱。 \end{tabularx} \\ \\ \relax
9:8 & \begin{tabularx}{0.7\textwidth}{X} 不要責備傲慢人，免得他恨你；要責備智慧人，他必愛你。 \end{tabularx} \\ \\ \relax
9:9 & \begin{tabularx}{0.7\textwidth}{X} 教導智慧人，他就越有智慧，指示義人，他就增長學問。 \end{tabularx} \\ \\ \relax
9:10 & \begin{tabularx}{0.7\textwidth}{X} 敬畏耶和華是智慧的開端，認識至聖者便是聰明。 \end{tabularx} \\ \\ \relax
9:11 & \begin{tabularx}{0.7\textwidth}{X} 藉著我，你的日子必增多，你生命的年數也必加添。 \end{tabularx} \\ \\ \relax
9:12 & \begin{tabularx}{0.7\textwidth}{X} 你若有智慧，是自己有智慧；你若傲慢，就自己承擔。 \end{tabularx} \\ \\ \relax
9:13 & \begin{tabularx}{0.7\textwidth}{X} 愚昧的女子喧嚷，她是愚蒙，一無所知。 \end{tabularx} \\ \\ \relax
9:14 & \begin{tabularx}{0.7\textwidth}{X} 她坐在自己家門口，在城中高處的座位上， \end{tabularx} \\ \\ \relax
9:15 & \begin{tabularx}{0.7\textwidth}{X} 呼喚過路的，向那些在路上直走的人說： \end{tabularx} \\ \\ \relax
9:16 & \begin{tabularx}{0.7\textwidth}{X} 「誰是愚蒙的人，讓他轉到這裡來！」又對那無知的人說： \end{tabularx} \\ \\ \relax
9:17 & \begin{tabularx}{0.7\textwidth}{X} 「偷來的水是甜的，暗藏的餅是美的。」 \end{tabularx} \\ \\ \relax
9:18 & \begin{tabularx}{0.7\textwidth}{X} 人卻不知有陰魂在她那裡，她召喚的人是在陰間的深處。 \end{tabularx} \\ \\
[1ex]
\hline
\hline
\end{longtable}


\newpage

\section{第76屆~港九培靈會~羅祖澄~第一講}
\label{sec:346}
\textbf{基督徒的新生命}
\newline
\newline
連結: \href{https://www.hkbibleconference.org/session-message/view/346}{\texttt{https://www.hkbibleconference.org/session-message/view/346}}
\newline
\newline
\hyperref[sec:345]{< < < PREV SERMON < < <}
~
\hyperref[sec:toc]{[返目錄]}
~
\hyperref[sec:347]{> > > NEXT SERMON > > >}
\newline
\newline
哥林多後書 5:11-5:21
\newline
\begin{longtable}{cl}
\hline
\hline
章節 & 經文 (和合本修訂版)\\
\hline
5:11 & \begin{tabularx}{0.7\textwidth}{X} 既然我們知道主是可畏的，就勸導人；但是神是認識我們的，我盼望你們的良心也認識我們。 \end{tabularx} \\ \\ \relax
5:12 & \begin{tabularx}{0.7\textwidth}{X} 我們不是向你們再推薦自己，而是要讓你們有誇耀我們的機會，使你們好面對那憑外貌、不憑內心誇耀的人。 \end{tabularx} \\ \\ \relax
5:13 & \begin{tabularx}{0.7\textwidth}{X} 如果我們癲狂，是為神；如果我們清醒，是為你們。 \end{tabularx} \\ \\ \relax
5:14 & \begin{tabularx}{0.7\textwidth}{X} 原來基督的愛激勵我們；因我們這樣斷定，一人既替眾人死，眾人就都死了； \end{tabularx} \\ \\ \relax
5:15 & \begin{tabularx}{0.7\textwidth}{X} 並且他替眾人死，是叫那些活著的人不再為自己活，乃為替他們死而復活的主活。 \end{tabularx} \\ \\ \relax
5:16 & \begin{tabularx}{0.7\textwidth}{X} 所以，從今以後，我們不再按照人的看法來認識人，縱使我們曾經按照人的看法認識基督，如今卻不再這樣認識他了。 \end{tabularx} \\ \\ \relax
5:17 & \begin{tabularx}{0.7\textwidth}{X} 所以，若有人在基督裡，他就是新造的人，舊事已過，都變成新的了。 \end{tabularx} \\ \\ \relax
5:18 & \begin{tabularx}{0.7\textwidth}{X} 一切都是出於神；他藉著基督使我們與他和好，又將勸人與他和好的使命賜給我們。 \end{tabularx} \\ \\ \relax
5:19 & \begin{tabularx}{0.7\textwidth}{X} 這就是：神在基督裡使世人與自己和好，不將他們的過犯歸到他們身上，並且將這和好的信息託付了我們。 \end{tabularx} \\ \\ \relax
5:20 & \begin{tabularx}{0.7\textwidth}{X} 所以，我們作基督的特使，就好像神藉我們勸你們一般。我們替基督求你們，與神和好吧！ \end{tabularx} \\ \\ \relax
5:21 & \begin{tabularx}{0.7\textwidth}{X} 神使那無罪的，替我們成為罪，好使我們在他裡面成為神的義。 \end{tabularx} \\ \\
[1ex]
\hline
\hline
\end{longtable}


\newpage

\section{第76屆~港九培靈會~羅祖澄~第二講}
\label{sec:347}
\textbf{新生命的標誌}
\newline
\newline
連結: \href{https://www.hkbibleconference.org/session-message/view/347}{\texttt{https://www.hkbibleconference.org/session-message/view/347}}
\newline
\newline
\hyperref[sec:346]{< < < PREV SERMON < < <}
~
\hyperref[sec:toc]{[返目錄]}
~
\hyperref[sec:348]{> > > NEXT SERMON > > >}
\newline
\newline
帖撒羅尼迦前書 1:1-1:3
\newline
\begin{longtable}{cl}
\hline
\hline
章節 & 經文 (和合本修訂版)\\
\hline
1:1 & \begin{tabularx}{0.7\textwidth}{X} 保羅、西拉、提摩太寫信給帖撒羅尼迦在父神和主耶穌基督裡的教會。願恩惠、平安歸給你們！ \end{tabularx} \\ \\ \relax
1:2 & \begin{tabularx}{0.7\textwidth}{X} 我們為你們眾人常常感謝神，禱告的時候提到你們， \end{tabularx} \\ \\ \relax
1:3 & \begin{tabularx}{0.7\textwidth}{X} 在我們的父神面前，不住地記念你們因信心所做的工作，因愛心所受的勞苦，因盼望我們主耶穌基督所存的堅忍。 \end{tabularx} \\ \\
[1ex]
\hline
\hline
\end{longtable}


\newpage

\section{第76屆~港九培靈會~羅祖澄~第三講}
\label{sec:348}
\textbf{生命的重整}
\newline
\newline
連結: \href{https://www.hkbibleconference.org/session-message/view/348}{\texttt{https://www.hkbibleconference.org/session-message/view/348}}
\newline
\newline
\hyperref[sec:347]{< < < PREV SERMON < < <}
~
\hyperref[sec:toc]{[返目錄]}
~
\hyperref[sec:349]{> > > NEXT SERMON > > >}
\newline
\newline
詩篇 127:1-127:5
\newline
\begin{longtable}{cl}
\hline
\hline
章節 & 經文 (和合本修訂版)\\
\hline
127:1 & \begin{tabularx}{0.7\textwidth}{X} 若不是耶和華建造房屋，建造的人就枉然勞力；若不是耶和華看守城池，看守的人枉然警醒。 \end{tabularx} \\ \\ \relax
127:2 & \begin{tabularx}{0.7\textwidth}{X} 你們清晨早起，夜晚安歇，吃勞碌得來的飯，本是枉然；惟有耶和華所親愛的，必叫他安然睡覺。 \end{tabularx} \\ \\ \relax
127:3 & \begin{tabularx}{0.7\textwidth}{X} 看哪，兒女是耶和華所賜的產業，所懷的胎是他所給的賞賜。 \end{tabularx} \\ \\ \relax
127:4 & \begin{tabularx}{0.7\textwidth}{X} 人在年輕時生的兒女好像勇士手中的箭。 \end{tabularx} \\ \\ \relax
127:5 & \begin{tabularx}{0.7\textwidth}{X} 箭袋充滿的人有福了！他們在城門口和仇敵爭論時必不蒙羞。 \end{tabularx} \\ \\
[1ex]
\hline
\hline
\end{longtable}


\newpage

\section{第76屆~港九培靈會~羅祖澄~第四講}
\label{sec:349}
\textbf{建立生命中的戰友}
\newline
\newline
連結: \href{https://www.hkbibleconference.org/session-message/view/349}{\texttt{https://www.hkbibleconference.org/session-message/view/349}}
\newline
\newline
\hyperref[sec:348]{< < < PREV SERMON < < <}
~
\hyperref[sec:toc]{[返目錄]}
~
\hyperref[sec:350]{> > > NEXT SERMON > > >}
\newline
\newline
腓立比書 1:1-1:11
\newline
\begin{longtable}{cl}
\hline
\hline
章節 & 經文 (和合本修訂版)\\
\hline
1:1 & \begin{tabularx}{0.7\textwidth}{X} 基督耶穌的僕人保羅和提摩太寫信給住腓立比、在基督耶穌裡的眾聖徒，以及諸位監督和執事。 \end{tabularx} \\ \\ \relax
1:2 & \begin{tabularx}{0.7\textwidth}{X} 願恩惠、平安從我們的父神和主耶穌基督歸給你們！ \end{tabularx} \\ \\ \relax
1:3 & \begin{tabularx}{0.7\textwidth}{X} 我每逢想念你們，就感謝我的神， \end{tabularx} \\ \\ \relax
1:4 & \begin{tabularx}{0.7\textwidth}{X} 每逢為你們眾人祈求的時候，總是歡歡喜喜地祈求， \end{tabularx} \\ \\ \relax
1:5 & \begin{tabularx}{0.7\textwidth}{X} 因為從第一天直到如今，你們都同心合意興旺福音。 \end{tabularx} \\ \\ \relax
1:6 & \begin{tabularx}{0.7\textwidth}{X} 我深信，那在你們心裡動了美好工作的，到了耶穌基督的日子必完成這工作。 \end{tabularx} \\ \\ \relax
1:7 & \begin{tabularx}{0.7\textwidth}{X} 我為你們眾人有這樣的想法原是應當的，因為你們常在我心裡；無論我是在捆鎖中，在辯明並證實福音的時候，你們都與我一同蒙恩。 \end{tabularx} \\ \\ \relax
1:8 & \begin{tabularx}{0.7\textwidth}{X} 我以基督耶穌的心腸切切想念你們眾人，這是神可以為我作證的。 \end{tabularx} \\ \\ \relax
1:9 & \begin{tabularx}{0.7\textwidth}{X} 我所禱告的就是：要你們的愛心，在知識和各樣見識上，不斷增長， \end{tabularx} \\ \\ \relax
1:10 & \begin{tabularx}{0.7\textwidth}{X} 使你們能分辨是非，在基督的日子作真誠無可指責的人， \end{tabularx} \\ \\ \relax
1:11 & \begin{tabularx}{0.7\textwidth}{X} 更靠著耶穌基督結滿仁義的果子，歸榮耀稱讚給神。 \end{tabularx} \\ \\
[1ex]
\hline
\hline
\end{longtable}


\newpage

\section{第76屆~港九培靈會~羅祖澄~第五講}
\label{sec:350}
\textbf{超越生命的虛空}
\newline
\newline
連結: \href{https://www.hkbibleconference.org/session-message/view/350}{\texttt{https://www.hkbibleconference.org/session-message/view/350}}
\newline
\newline
\hyperref[sec:349]{< < < PREV SERMON < < <}
~
\hyperref[sec:toc]{[返目錄]}
~
\hyperref[sec:351]{> > > NEXT SERMON > > >}
\newline
\newline
傳道書 2:1-2:11
\newline
\begin{longtable}{cl}
\hline
\hline
章節 & 經文 (和合本修訂版)\\
\hline
2:1 & \begin{tabularx}{0.7\textwidth}{X} 我心裡說：「來吧，讓我用喜樂試試你，使你享福！」看哪，這也是虛空。 \end{tabularx} \\ \\ \relax
2:2 & \begin{tabularx}{0.7\textwidth}{X} 論嬉笑，我說：「這是狂妄。」論享樂，「這有甚麼用呢？」 \end{tabularx} \\ \\ \relax
2:3 & \begin{tabularx}{0.7\textwidth}{X} 我心以智慧引導我，我心裡探究，如何用酒使身體舒暢，如何抓住愚昧，直等我看明世人在天下短暫一生中，當行何事為美。 \end{tabularx} \\ \\ \relax
2:4 & \begin{tabularx}{0.7\textwidth}{X} 我大興土木，為自己建造房屋，栽葡萄園， \end{tabularx} \\ \\ \relax
2:5 & \begin{tabularx}{0.7\textwidth}{X} 修造庭園和公園，在其中栽種各樣果樹， \end{tabularx} \\ \\ \relax
2:6 & \begin{tabularx}{0.7\textwidth}{X} 挖造水池，用以灌溉林中的幼樹。 \end{tabularx} \\ \\ \relax
2:7 & \begin{tabularx}{0.7\textwidth}{X} 我買了僕婢，也有生在家中的僕婢；又有許多牛群羊群，勝過我以前所有在耶路撒冷的人。 \end{tabularx} \\ \\ \relax
2:8 & \begin{tabularx}{0.7\textwidth}{X} 我為自己積蓄金銀，搜集各君王、各省份的財寶；又為自己得男女歌手和世人所喜愛的物，以及一個又一個的妃嬪。 \end{tabularx} \\ \\ \relax
2:9 & \begin{tabularx}{0.7\textwidth}{X} 這樣，我就日漸昌盛，勝過我以前所有在耶路撒冷的人。我的智慧仍然存留。 \end{tabularx} \\ \\ \relax
2:10 & \begin{tabularx}{0.7\textwidth}{X} 凡我眼所求的，我沒有克制它；我心所樂的，我沒有不享受。因我的心要為一切的勞碌快樂，這是我從一切勞碌中所得的報償。 \end{tabularx} \\ \\ \relax
2:11 & \begin{tabularx}{0.7\textwidth}{X} 後來，我回顧我手所經營的一切和我勞碌所做的工。看哪，全是虛空，全是捕風；在日光之下毫無益處。 \end{tabularx} \\ \\
[1ex]
\hline
\hline
\end{longtable}


\newpage

\section{第76屆~港九培靈會~羅祖澄~第六講}
\label{sec:351}
\textbf{勝過生命的困鎖}
\newline
\newline
連結: \href{https://www.hkbibleconference.org/session-message/view/351}{\texttt{https://www.hkbibleconference.org/session-message/view/351}}
\newline
\newline
\hyperref[sec:350]{< < < PREV SERMON < < <}
~
\hyperref[sec:toc]{[返目錄]}
~
\hyperref[sec:352]{> > > NEXT SERMON > > >}
\newline
\newline
約翰福音 20:19-20:23
\newline
\begin{longtable}{cl}
\hline
\hline
章節 & 經文 (和合本修訂版)\\
\hline
20:19 & \begin{tabularx}{0.7\textwidth}{X} 那日（就是七日的第一日）晚上，門徒因怕猶太人，所在的地方門都關了。耶穌來，站在當中，對他們說：「願你們平安！」 \end{tabularx} \\ \\ \relax
20:20 & \begin{tabularx}{0.7\textwidth}{X} 說了這話，他把手和肋旁給他們看。門徒一看見主就喜樂了。 \end{tabularx} \\ \\ \relax
20:21 & \begin{tabularx}{0.7\textwidth}{X} 於是耶穌又對他們說：「願你們平安！父怎樣差遣了我，我也照樣差遣你們。」 \end{tabularx} \\ \\ \relax
20:22 & \begin{tabularx}{0.7\textwidth}{X} 說了這話，他向他們吹一口氣，說：「領受聖靈吧！ \end{tabularx} \\ \\ \relax
20:23 & \begin{tabularx}{0.7\textwidth}{X} 你們赦免誰的罪，誰的罪就得赦免；你們不赦免誰的罪，誰的罪就不得赦免。」 \end{tabularx} \\ \\
[1ex]
\hline
\hline
\end{longtable}


\newpage

\section{第76屆~港九培靈會~羅祖澄~第七講}
\label{sec:352}
\textbf{得力的生命}
\newline
\newline
連結: \href{https://www.hkbibleconference.org/session-message/view/352}{\texttt{https://www.hkbibleconference.org/session-message/view/352}}
\newline
\newline
\hyperref[sec:351]{< < < PREV SERMON < < <}
~
\hyperref[sec:toc]{[返目錄]}
~
\hyperref[sec:353]{> > > NEXT SERMON > > >}
\newline
\newline
出埃及記 17:8-17:12
\newline
\begin{longtable}{cl}
\hline
\hline
章節 & 經文 (和合本修訂版)\\
\hline
17:8 & \begin{tabularx}{0.7\textwidth}{X} 那時，亞瑪力來到利非訂，和以色列爭戰。 \end{tabularx} \\ \\ \relax
17:9 & \begin{tabularx}{0.7\textwidth}{X} 摩西對約書亞說：「你為我們選出人來，出去和亞瑪力爭戰。明天我要站在山頂上，手裡拿著神的杖。」 \end{tabularx} \\ \\ \relax
17:10 & \begin{tabularx}{0.7\textwidth}{X} 於是，約書亞照著摩西對他所說的話去做，和亞瑪力爭戰。摩西、亞倫和戶珥都上了山頂。 \end{tabularx} \\ \\ \relax
17:11 & \begin{tabularx}{0.7\textwidth}{X} 摩西何時舉手，以色列就得勝；何時垂手，亞瑪力就得勝。 \end{tabularx} \\ \\ \relax
17:12 & \begin{tabularx}{0.7\textwidth}{X} 但摩西的雙手沉重，他們就搬一塊石頭來放在他下面，他就坐在上面。亞倫與戶珥扶著他的手，一個在這邊，一個在那邊，他的手就穩住，直到日落。 \end{tabularx} \\ \\
[1ex]
\hline
\hline
\end{longtable}


\newpage

\section{第76屆~港九培靈會~羅祖澄~第八講}
\label{sec:353}
\textbf{甘受差遣的生命}
\newline
\newline
連結: \href{https://www.hkbibleconference.org/session-message/view/353}{\texttt{https://www.hkbibleconference.org/session-message/view/353}}
\newline
\newline
\hyperref[sec:352]{< < < PREV SERMON < < <}
~
\hyperref[sec:toc]{[返目錄]}
~
\hyperref[sec:354]{> > > NEXT SERMON > > >}
\newline
\newline
以賽亞書 6:1-6:10
\newline
\begin{longtable}{cl}
\hline
\hline
章節 & 經文 (和合本修訂版)\\
\hline
6:1 & \begin{tabularx}{0.7\textwidth}{X} 當烏西雅王崩的那年，我看見主坐在高高的寶座上。他的衣裳下襬遮滿聖殿。 \end{tabularx} \\ \\ \relax
6:2 & \begin{tabularx}{0.7\textwidth}{X} 上有撒拉弗侍立，各有六個翅膀：兩個翅膀遮臉，兩個翅膀遮腳，兩個翅膀飛翔， \end{tabularx} \\ \\ \relax
6:3 & \begin{tabularx}{0.7\textwidth}{X} 彼此呼喊說：「聖哉！聖哉！聖哉！萬軍之耶和華；他的榮光遍滿全地！」 \end{tabularx} \\ \\ \relax
6:4 & \begin{tabularx}{0.7\textwidth}{X} 因呼喊者的聲音，門檻的根基震動，殿裡充滿了煙雲。 \end{tabularx} \\ \\ \relax
6:5 & \begin{tabularx}{0.7\textwidth}{X} 那時我說：「禍哉！我滅亡了！因為我是嘴唇不潔的人，住在嘴唇不潔的民中，又因我親眼看見大君王－萬軍之耶和華。」 \end{tabularx} \\ \\ \relax
6:6 & \begin{tabularx}{0.7\textwidth}{X} 有一撒拉弗向我飛來，手裡拿著燒紅的炭，是用火鉗從壇上取下來的， \end{tabularx} \\ \\ \relax
6:7 & \begin{tabularx}{0.7\textwidth}{X} 用炭沾我的口，說：「看哪，這炭沾了你的嘴唇，你的罪孽便除掉，你的罪惡就赦免了。」 \end{tabularx} \\ \\ \relax
6:8 & \begin{tabularx}{0.7\textwidth}{X} 我聽見主的聲音說：「我可以差遣誰呢？誰肯為我們去呢？」我說：「我在這裡，請差遣我！」 \end{tabularx} \\ \\ \relax
6:9 & \begin{tabularx}{0.7\textwidth}{X} 他說：「你去告訴這百姓說：『你們聽了又聽，卻不明白；看了又看，卻不曉得。』 \end{tabularx} \\ \\ \relax
6:10 & \begin{tabularx}{0.7\textwidth}{X} 要使這百姓心蒙油脂，耳朵發沉，眼睛昏花；恐怕他們眼睛看見，耳朵聽見，心裡明白，回轉過來，就得醫治。」 \end{tabularx} \\ \\
[1ex]
\hline
\hline
\end{longtable}


\newpage

\section{第76屆~港九培靈會~高文~第一講}
\label{sec:354}
\textbf{耶穌基督是主}
\newline
\newline
連結: \href{https://www.hkbibleconference.org/session-message/view/354}{\texttt{https://www.hkbibleconference.org/session-message/view/354}}
\newline
\newline
\hyperref[sec:353]{< < < PREV SERMON < < <}
~
\hyperref[sec:toc]{[返目錄]}
~
\hyperref[sec:355]{> > > NEXT SERMON > > >}
\newline
\newline
腓立比書 2:1-2:11
\newline
\begin{longtable}{cl}
\hline
\hline
章節 & 經文 (和合本修訂版)\\
\hline
2:1 & \begin{tabularx}{0.7\textwidth}{X} 所以，在基督裡若有任何勸勉，若有任何愛心的安慰，若有任何聖靈的團契，若有任何慈悲憐憫， \end{tabularx} \\ \\ \relax
2:2 & \begin{tabularx}{0.7\textwidth}{X} 你們就要意志相同，愛心相同，有一致的心思，一致的想法，使我的喜樂得以滿足。 \end{tabularx} \\ \\ \relax
2:3 & \begin{tabularx}{0.7\textwidth}{X} 凡事不可自私自利，不可貪圖虛榮；只要心存謙卑，各人看別人比自己強。 \end{tabularx} \\ \\ \relax
2:4 & \begin{tabularx}{0.7\textwidth}{X} 各人不要單顧自己的事，也要顧別人的事。 \end{tabularx} \\ \\ \relax
2:5 & \begin{tabularx}{0.7\textwidth}{X} 你們當以基督耶穌的心為心： \end{tabularx} \\ \\ \relax
2:6 & \begin{tabularx}{0.7\textwidth}{X} 他本有神的形像，卻不堅持自己與神同等； \end{tabularx} \\ \\ \relax
2:7 & \begin{tabularx}{0.7\textwidth}{X} 反倒虛己，取了奴僕的形像，成為人的樣式；既有人的樣子， \end{tabularx} \\ \\ \relax
2:8 & \begin{tabularx}{0.7\textwidth}{X} 就謙卑自己，存心順服，以至於死，且死在十字架上。 \end{tabularx} \\ \\ \relax
2:9 & \begin{tabularx}{0.7\textwidth}{X} 所以神把他升為至高，又賜給他超乎萬名之上的名， \end{tabularx} \\ \\ \relax
2:10 & \begin{tabularx}{0.7\textwidth}{X} 使一切在天上的、地上的和地底下的，因耶穌的名，眾膝都要跪下， \end{tabularx} \\ \\ \relax
2:11 & \begin{tabularx}{0.7\textwidth}{X} 眾口都要宣認：耶穌基督是主，歸榮耀給父神。 \end{tabularx} \\ \\
[1ex]
\hline
\hline
\end{longtable}


\newpage

\section{第76屆~港九培靈會~高文~第二講}
\label{sec:355}
\textbf{選擇耶穌的窄路}
\newline
\newline
連結: \href{https://www.hkbibleconference.org/session-message/view/355}{\texttt{https://www.hkbibleconference.org/session-message/view/355}}
\newline
\newline
\hyperref[sec:354]{< < < PREV SERMON < < <}
~
\hyperref[sec:toc]{[返目錄]}
~
\hyperref[sec:356]{> > > NEXT SERMON > > >}
\newline
\newline
馬太福音 7:13-7:14
\newline
\begin{longtable}{cl}
\hline
\hline
章節 & 經文 (和合本修訂版)\\
\hline
7:13 & \begin{tabularx}{0.7\textwidth}{X} 「你們要進窄門。因為通往滅亡的門是寬的，路是大的，進去的人也多； \end{tabularx} \\ \\ \relax
7:14 & \begin{tabularx}{0.7\textwidth}{X} 通往生命的門是窄的，路是小的，找到的人也少。」 \end{tabularx} \\ \\
[1ex]
\hline
\hline
\end{longtable}


\newpage

\section{第76屆~港九培靈會~高文~第三講}
\label{sec:356}
\textbf{耶穌的邀請}
\newline
\newline
連結: \href{https://www.hkbibleconference.org/session-message/view/356}{\texttt{https://www.hkbibleconference.org/session-message/view/356}}
\newline
\newline
\hyperref[sec:355]{< < < PREV SERMON < < <}
~
\hyperref[sec:toc]{[返目錄]}
~
\hyperref[sec:357]{> > > NEXT SERMON > > >}
\newline
\newline
路加福音 19:1-19:10
\newline
\begin{longtable}{cl}
\hline
\hline
章節 & 經文 (和合本修訂版)\\
\hline
19:1 & \begin{tabularx}{0.7\textwidth}{X} 耶穌進了耶利哥，要從那裡經過。 \end{tabularx} \\ \\ \relax
19:2 & \begin{tabularx}{0.7\textwidth}{X} 有一個人名叫撒該，作稅吏長，是個財主。 \end{tabularx} \\ \\ \relax
19:3 & \begin{tabularx}{0.7\textwidth}{X} 他要看看耶穌是怎樣的人，只因人多，他的身材又矮，所以看不見。 \end{tabularx} \\ \\ \relax
19:4 & \begin{tabularx}{0.7\textwidth}{X} 於是他跑到前頭，爬上桑樹，要看耶穌，因為耶穌要從那裡經過。 \end{tabularx} \\ \\ \relax
19:5 & \begin{tabularx}{0.7\textwidth}{X} 耶穌到了那裡，抬頭一看，對他說：「撒該，快下來！今天我必須住在你家裡。」 \end{tabularx} \\ \\ \relax
19:6 & \begin{tabularx}{0.7\textwidth}{X} 他就急忙下來，歡歡喜喜地接待耶穌。 \end{tabularx} \\ \\ \relax
19:7 & \begin{tabularx}{0.7\textwidth}{X} 眾人看見，都私下議論說：「他竟然到罪人家裡去住宿。」 \end{tabularx} \\ \\ \relax
19:8 & \begin{tabularx}{0.7\textwidth}{X} 撒該站著對主說：「主啊，我把所有的一半給窮人；我若勒索了誰，就還他四倍。」 \end{tabularx} \\ \\ \relax
19:9 & \begin{tabularx}{0.7\textwidth}{X} 耶穌對他說：「今天救恩到了這家，因為他也是亞伯拉罕的子孫。 \end{tabularx} \\ \\ \relax
19:10 & \begin{tabularx}{0.7\textwidth}{X} 人子來是要尋找和拯救失喪的人。」 \end{tabularx} \\ \\
[1ex]
\hline
\hline
\end{longtable}


\newpage

\section{第76屆~港九培靈會~高文~第四講}
\label{sec:357}
\textbf{為耶穌收割}
\newline
\newline
連結: \href{https://www.hkbibleconference.org/session-message/view/357}{\texttt{https://www.hkbibleconference.org/session-message/view/357}}
\newline
\newline
\hyperref[sec:356]{< < < PREV SERMON < < <}
~
\hyperref[sec:toc]{[返目錄]}
~
\hyperref[sec:358]{> > > NEXT SERMON > > >}
\newline
\newline
約翰福音 4:31-4:42
\newline
\begin{longtable}{cl}
\hline
\hline
章節 & 經文 (和合本修訂版)\\
\hline
4:31 & \begin{tabularx}{0.7\textwidth}{X} 就在這個時候，門徒求耶穌說：「拉比，請吃吧。」 \end{tabularx} \\ \\ \relax
4:32 & \begin{tabularx}{0.7\textwidth}{X} 耶穌對他們說：「我有食物吃，是你們不知道的。」 \end{tabularx} \\ \\ \relax
4:33 & \begin{tabularx}{0.7\textwidth}{X} 門徒就彼此說：「難道有人拿甚麼給他吃了嗎？」 \end{tabularx} \\ \\ \relax
4:34 & \begin{tabularx}{0.7\textwidth}{X} 耶穌對他們說：「我的食物就是要遵行差我來那位的旨意，完成他的工作。 \end{tabularx} \\ \\ \relax
4:35 & \begin{tabularx}{0.7\textwidth}{X} 你們不是說『到收割的時候還有四個月』嗎？我告訴你們，舉目向田觀看，莊稼熟了，可以收割了。 \end{tabularx} \\ \\ \relax
4:36 & \begin{tabularx}{0.7\textwidth}{X} 收割的人已經得工錢，為永生儲存五穀，使撒種的和收割的一同快樂。 \end{tabularx} \\ \\ \relax
4:37 & \begin{tabularx}{0.7\textwidth}{X} 『那人撒種，這人收割』，這話可見是真的。 \end{tabularx} \\ \\ \relax
4:38 & \begin{tabularx}{0.7\textwidth}{X} 我差你們去收你們所沒有辛勞的；別人辛勞，你們享受他們辛勞的成果。」 \end{tabularx} \\ \\ \relax
4:39 & \begin{tabularx}{0.7\textwidth}{X} 那城裡有好些撒瑪利亞人信了耶穌，因為那婦人作見證，說：「他把我素來所做的一切事都說了出來。」 \end{tabularx} \\ \\ \relax
4:40 & \begin{tabularx}{0.7\textwidth}{X} 於是撒瑪利亞人來見耶穌，求他在他們那裡住下，他就在那裡住了兩天。 \end{tabularx} \\ \\ \relax
4:41 & \begin{tabularx}{0.7\textwidth}{X} 因為耶穌的話，信的人就更多了。 \end{tabularx} \\ \\ \relax
4:42 & \begin{tabularx}{0.7\textwidth}{X} 他們對那婦人說：「現在我們信，不再是因為你的話，而是我們親自聽見了，知道這人真是世界的救主。」 \end{tabularx} \\ \\
[1ex]
\hline
\hline
\end{longtable}


\newpage

\section{第76屆~港九培靈會~高文~第五講}
\label{sec:358}
\textbf{耶穌的大使命}
\newline
\newline
連結: \href{https://www.hkbibleconference.org/session-message/view/358}{\texttt{https://www.hkbibleconference.org/session-message/view/358}}
\newline
\newline
\hyperref[sec:357]{< < < PREV SERMON < < <}
~
\hyperref[sec:toc]{[返目錄]}
~
\hyperref[sec:359]{> > > NEXT SERMON > > >}
\newline
\newline
馬太福音 28:16-28:20
\newline
\begin{longtable}{cl}
\hline
\hline
章節 & 經文 (和合本修訂版)\\
\hline
28:16 & \begin{tabularx}{0.7\textwidth}{X} 十一個門徒往加利利去，到了耶穌指定他們去的山上。 \end{tabularx} \\ \\ \relax
28:17 & \begin{tabularx}{0.7\textwidth}{X} 他們見了耶穌就拜他，然而還有人疑惑。 \end{tabularx} \\ \\ \relax
28:18 & \begin{tabularx}{0.7\textwidth}{X} 耶穌進前來，對他們說：「天上地下所有的權柄都賜給我了。 \end{tabularx} \\ \\ \relax
28:19 & \begin{tabularx}{0.7\textwidth}{X} 所以，你們要去，使萬民作我的門徒，奉父、子、聖靈的名給他們施洗， \end{tabularx} \\ \\ \relax
28:20 & \begin{tabularx}{0.7\textwidth}{X} 凡我所吩咐你們的，都教導他們遵守。看哪，我天天與你們同在，直到世代的終結。」 \end{tabularx} \\ \\
[1ex]
\hline
\hline
\end{longtable}


\newpage

\section{第76屆~港九培靈會~高文~第六講}
\label{sec:359}
\textbf{耶穌發出的決定性問題}
\newline
\newline
連結: \href{https://www.hkbibleconference.org/session-message/view/359}{\texttt{https://www.hkbibleconference.org/session-message/view/359}}
\newline
\newline
\hyperref[sec:358]{< < < PREV SERMON < < <}
~
\hyperref[sec:toc]{[返目錄]}
~
\hyperref[sec:360]{> > > NEXT SERMON > > >}
\newline
\newline
約翰福音 21:15-21:24
\newline
\begin{longtable}{cl}
\hline
\hline
章節 & 經文 (和合本修訂版)\\
\hline
21:15 & \begin{tabularx}{0.7\textwidth}{X} 他們吃完了早飯，耶穌對西門‧彼得說：「約翰的兒子西門，你愛我比這些更深嗎？」彼得對他說：「主啊，是的，你知道我愛你。」耶穌說：「你餵養我的小羊。」 \end{tabularx} \\ \\ \relax
21:16 & \begin{tabularx}{0.7\textwidth}{X} 耶穌第二次又對他說：「約翰的兒子西門，你愛我嗎？」彼得對他說：「主啊，是的，你知道我愛你。」耶穌說：「你牧養我的羊。」 \end{tabularx} \\ \\ \relax
21:17 & \begin{tabularx}{0.7\textwidth}{X} 耶穌第三次對他說：「約翰的兒子西門，你愛我嗎？」彼得因為耶穌第三次對他說「你愛我嗎」，就憂愁，對耶穌說：「主啊，你無所不知，你知道我愛你。」耶穌說：「你餵養我的羊。 \end{tabularx} \\ \\ \relax
21:18 & \begin{tabularx}{0.7\textwidth}{X} 我實實在在地告訴你，你年輕的時候，自己束上帶子，隨意往來；但年老的時候，你要伸出手來，別人要把你束上，帶你到不願意去的地方。」 \end{tabularx} \\ \\ \relax
21:19 & \begin{tabularx}{0.7\textwidth}{X} 耶穌說這話，是指彼得會怎樣死來榮耀神。說了這話，耶穌對他說：「你跟從我吧！」 \end{tabularx} \\ \\ \relax
21:20 & \begin{tabularx}{0.7\textwidth}{X} 彼得轉過身來，看見耶穌所愛的那門徒跟著，就是在晚餐時靠著耶穌胸膛說「主啊，出賣你的是誰」的那門徒。 \end{tabularx} \\ \\ \relax
21:21 & \begin{tabularx}{0.7\textwidth}{X} 彼得看見他，就問耶穌：「主啊，這個人怎樣呢？」 \end{tabularx} \\ \\ \relax
21:22 & \begin{tabularx}{0.7\textwidth}{X} 耶穌對他說：「假如我要他等到我來的時候還在，跟你有甚麼關係呢？你跟從我吧！」 \end{tabularx} \\ \\ \relax
21:23 & \begin{tabularx}{0.7\textwidth}{X} 於是這話在弟兄中間流傳，說那門徒不死。其實，耶穌不是說他不死，而是對彼得說：「假如我要他等到我來的時候還在，跟你有甚麼關係呢？」 \end{tabularx} \\ \\ \relax
21:24 & \begin{tabularx}{0.7\textwidth}{X} 這門徒就是為這些事作見證、並且記載這些事的，我們知道他的見證是真的。 \end{tabularx} \\ \\
[1ex]
\hline
\hline
\end{longtable}


\newpage

\section{第76屆~港九培靈會~高文~第七講}
\label{sec:360}
\textbf{耶穌的同在}
\newline
\newline
連結: \href{https://www.hkbibleconference.org/session-message/view/360}{\texttt{https://www.hkbibleconference.org/session-message/view/360}}
\newline
\newline
\hyperref[sec:359]{< < < PREV SERMON < < <}
~
\hyperref[sec:toc]{[返目錄]}
~
\hyperref[sec:362]{> > > NEXT SERMON > > >}
\newline
\newline
約翰福音 14:15-14:17
\newline
\begin{longtable}{cl}
\hline
\hline
章節 & 經文 (和合本修訂版)\\
\hline
14:15 & \begin{tabularx}{0.7\textwidth}{X} 「你們若愛我，就會遵守我的命令。 \end{tabularx} \\ \\ \relax
14:16 & \begin{tabularx}{0.7\textwidth}{X} 我要求父，父就賜給你們另外一位保惠師，使他永遠與你們同在。 \end{tabularx} \\ \\ \relax
14:17 & \begin{tabularx}{0.7\textwidth}{X} 他就是真理的靈，是世人不能接受的。因為他們既看不見他，也不認識他；你們卻認識他，因他常與你們同在，也要在你們裡面。 \end{tabularx} \\ \\
[1ex]
\hline
\hline
\end{longtable}


\newpage

\section{第76屆~港九培靈會~高文~第八講}
\label{sec:362}
\textbf{耶穌所賜的復興}
\newline
\newline
連結: \href{https://www.hkbibleconference.org/session-message/view/362}{\texttt{https://www.hkbibleconference.org/session-message/view/362}}
\newline
\newline
\hyperref[sec:360]{< < < PREV SERMON < < <}
~
\hyperref[sec:toc]{[返目錄]}
~
\hyperref[sec:363]{> > > NEXT SERMON > > >}
\newline
\newline
路加福音 24:44-24:49
\newline
\begin{longtable}{cl}
\hline
\hline
章節 & 經文 (和合本修訂版)\\
\hline
24:44 & \begin{tabularx}{0.7\textwidth}{X} 耶穌對他們說：「這就是我從前和你們同在時所告訴你們的話：摩西的律法、先知的書，和《詩篇》上所記一切指著我的話都必須應驗。」 \end{tabularx} \\ \\ \relax
24:45 & \begin{tabularx}{0.7\textwidth}{X} 於是耶穌開他們的心竅，使他們能明白聖經， \end{tabularx} \\ \\ \relax
24:46 & \begin{tabularx}{0.7\textwidth}{X} 又對他們說：「照經上所寫的，基督必受害，第三天從死人中復活， \end{tabularx} \\ \\ \relax
24:47 & \begin{tabularx}{0.7\textwidth}{X} 並且人們要奉他的名傳悔改、使罪得赦的道，從耶路撒冷起直傳到萬邦。 \end{tabularx} \\ \\ \relax
24:48 & \begin{tabularx}{0.7\textwidth}{X} 你們就是這些事的見證。 \end{tabularx} \\ \\ \relax
24:49 & \begin{tabularx}{0.7\textwidth}{X} 我要將我父所應許的降在你們身上，你們要在城裡等候，直到你們領受從上面來的能力。」 \end{tabularx} \\ \\
[1ex]
\hline
\hline
\end{longtable}


\newpage

\section{第76屆~港九培靈會~高文~第九講}
\label{sec:363}
\textbf{在耶穌的權柄裡得勝}
\newline
\newline
連結: \href{https://www.hkbibleconference.org/session-message/view/363}{\texttt{https://www.hkbibleconference.org/session-message/view/363}}
\newline
\newline
\hyperref[sec:362]{< < < PREV SERMON < < <}
~
\hyperref[sec:toc]{[返目錄]}
~
\hyperref[sec:364]{> > > NEXT SERMON > > >}
\newline
\newline
啟示錄 12:10-12:12
\newline
\begin{longtable}{cl}
\hline
\hline
章節 & 經文 (和合本修訂版)\\
\hline
12:10 & \begin{tabularx}{0.7\textwidth}{X} 我聽見在天上有大聲音說：「我神的救恩、能力、國度，和他所立的基督的權柄現在都來到了。因為那個在我們神面前、晝夜控告我們弟兄的，已經被摔下去了。 \end{tabularx} \\ \\ \relax
12:11 & \begin{tabularx}{0.7\textwidth}{X} 弟兄勝過那條龍是因羔羊的血，和因自己所見證的道。雖然至於死，他們也不惜自己的性命。 \end{tabularx} \\ \\ \relax
12:12 & \begin{tabularx}{0.7\textwidth}{X} 所以，諸天和住在其中的，你們都快樂吧！只是地和海有禍了！因為魔鬼知道自己的時候不多，就氣憤憤地下到你們那裡去了。」 \end{tabularx} \\ \\
[1ex]
\hline
\hline
\end{longtable}


\newpage

\section{第76屆~港九培靈會~高文~第十講}
\label{sec:364}
\textbf{與耶穌同赴婚筵}
\newline
\newline
連結: \href{https://www.hkbibleconference.org/session-message/view/364}{\texttt{https://www.hkbibleconference.org/session-message/view/364}}
\newline
\newline
\hyperref[sec:363]{< < < PREV SERMON < < <}
~
\hyperref[sec:toc]{[返目錄]}
~
\hyperref[sec:319]{> > > NEXT SERMON > > >}
\newline
\newline
啟示錄 19:1-19:9
\newline
\begin{longtable}{cl}
\hline
\hline
章節 & 經文 (和合本修訂版)\\
\hline
19:1 & \begin{tabularx}{0.7\textwidth}{X} 此後，我聽見好像有一大群人在天上大聲說：「哈利路亞！救恩、榮耀、權能都屬於我們的神。 \end{tabularx} \\ \\ \relax
19:2 & \begin{tabularx}{0.7\textwidth}{X} 他的判斷又真實又公義；因他判斷了那大淫婦，她用淫行敗壞了世界。神為他的僕人伸冤，向淫婦討流僕人血的罪。」 \end{tabularx} \\ \\ \relax
19:3 & \begin{tabularx}{0.7\textwidth}{X} 他們又一次說：「哈利路亞！燒淫婦的煙往上冒，直到永永遠遠。」 \end{tabularx} \\ \\ \relax
19:4 & \begin{tabularx}{0.7\textwidth}{X} 那二十四位長老和四活物就俯伏敬拜坐在寶座上的神，說：「阿們。哈利路亞！」 \end{tabularx} \\ \\ \relax
19:5 & \begin{tabularx}{0.7\textwidth}{X} 接著，有聲音從寶座出來說：「神的眾僕人哪，凡敬畏他的，無論大小，都要讚美我們的神！」 \end{tabularx} \\ \\ \relax
19:6 & \begin{tabularx}{0.7\textwidth}{X} 我聽見好像一大群人的聲音，像眾水的聲音，像大雷的聲音，說：「哈利路亞！因為主---我們的神、全能者，作王了。 \end{tabularx} \\ \\ \relax
19:7 & \begin{tabularx}{0.7\textwidth}{X} 我們要歡喜快樂，將榮耀歸給他；因為羔羊的婚期到了，他的新娘也自己預備好了， \end{tabularx} \\ \\ \relax
19:8 & \begin{tabularx}{0.7\textwidth}{X} 她蒙恩得穿明亮潔白的細麻衣：這細麻衣就是聖徒們的義行。」 \end{tabularx} \\ \\ \relax
19:9 & \begin{tabularx}{0.7\textwidth}{X} 天使對我說：「你要寫下來：凡被請赴羔羊婚宴的人有福了！」他又對我說：「這些都是神真實的話。」 \end{tabularx} \\ \\
[1ex]
\hline
\hline
\end{longtable}


\newpage

\section{第77屆~港九培靈會~吳獻章~第一講}
\label{sec:319}
\textbf{永恆的託付──摩西}
\newline
\newline
連結: \href{https://www.hkbibleconference.org/session-message/view/319}{\texttt{https://www.hkbibleconference.org/session-message/view/319}}
\newline
\newline
\hyperref[sec:364]{< < < PREV SERMON < < <}
~
\hyperref[sec:toc]{[返目錄]}
~
\hyperref[sec:320]{> > > NEXT SERMON > > >}
\newline
\newline
出埃及記 2:11-3:6
\newline
\begin{longtable}{cl}
\hline
\hline
章節 & 經文 (和合本修訂版)\\
\hline
2:11 & \begin{tabularx}{0.7\textwidth}{X} 過了一段日子，摩西長大了，他出去到他同胞那裡，看見他們的勞役。他看見一個埃及人打他的同胞，一個希伯來人。 \end{tabularx} \\ \\ \relax
2:12 & \begin{tabularx}{0.7\textwidth}{X} 他左右觀看，見沒有人，就把埃及人打死了，藏在沙土裡。 \end{tabularx} \\ \\ \relax
2:13 & \begin{tabularx}{0.7\textwidth}{X} 第二天他出去，看哪，有兩個希伯來人在打架，他就對那兇惡的人說：「你為甚麼打你同族的人呢？」 \end{tabularx} \\ \\ \relax
2:14 & \begin{tabularx}{0.7\textwidth}{X} 那人說：「誰立你作我們的領袖和審判官呢？難道你要殺我，像殺那埃及人一樣嗎？」摩西就懼怕，說：「這事一定是讓人知道了。」 \end{tabularx} \\ \\ \relax
2:15 & \begin{tabularx}{0.7\textwidth}{X} 法老聽見這事，就設法要殺摩西。於是摩西逃走，躲避法老，到了米甸地，坐在井旁。 \end{tabularx} \\ \\ \relax
2:16 & \begin{tabularx}{0.7\textwidth}{X} 米甸的祭司有七個女兒；她們來打水，打滿了槽，要給父親的羊群喝水。 \end{tabularx} \\ \\ \relax
2:17 & \begin{tabularx}{0.7\textwidth}{X} 有一些牧羊人來，把她們趕走，摩西卻起來幫助她們，取水給她們的羊群喝。 \end{tabularx} \\ \\ \relax
2:18 & \begin{tabularx}{0.7\textwidth}{X} 她們回到父親流珥那裡；他說：「今日你們為何這麼快就回來了呢？」 \end{tabularx} \\ \\ \relax
2:19 & \begin{tabularx}{0.7\textwidth}{X} 她們說：「有一個埃及人來救我們脫離牧羊人的手，他甚至打水給我們的羊群喝。」 \end{tabularx} \\ \\ \relax
2:20 & \begin{tabularx}{0.7\textwidth}{X} 他對女兒們說：「那人在哪裡？你們為甚麼撇下他呢？去請他來吃飯吧！」 \end{tabularx} \\ \\ \relax
2:21 & \begin{tabularx}{0.7\textwidth}{X} 摩西願意和那人同住， 那人就把女兒西坡拉給摩西為妻。 \end{tabularx} \\ \\ \relax
2:22 & \begin{tabularx}{0.7\textwidth}{X} 西坡拉生了一個兒子，摩西給他起名叫革舜，因他說：「我在外地作了寄居者。」 \end{tabularx} \\ \\ \relax
2:23 & \begin{tabularx}{0.7\textwidth}{X} 過了許多年，埃及王死了。以色列人因做苦工，就嘆息哀求；他們因苦工所發出的哀聲達於神。 \end{tabularx} \\ \\ \relax
2:24 & \begin{tabularx}{0.7\textwidth}{X} 神聽見他們的哀聲，就記念他與亞伯拉罕、以撒、雅各所立的約。 \end{tabularx} \\ \\ \relax
2:25 & \begin{tabularx}{0.7\textwidth}{X} 神看顧以色列人，神是知道的。 \end{tabularx} \\ \\ \relax
3:1 & \begin{tabularx}{0.7\textwidth}{X} 摩西牧放他岳父米甸祭司葉特羅的羊群，他領羊群往曠野的那一邊去，到了神的山，就是何烈山。 \end{tabularx} \\ \\ \relax
3:2 & \begin{tabularx}{0.7\textwidth}{X} 耶和華的使者在荊棘的火焰中向他顯現。摩西觀看，看哪，荊棘在火中焚燒，卻沒有燒燬。 \end{tabularx} \\ \\ \relax
3:3 & \begin{tabularx}{0.7\textwidth}{X} 摩西說：「我要轉過去看這大異象，這荊棘為何沒有燒燬呢？」 \end{tabularx} \\ \\ \relax
3:4 & \begin{tabularx}{0.7\textwidth}{X} 耶和華見摩西轉過去看，神就從荊棘裡呼叫他說：「摩西！摩西！」他說：「我在這裡。」 \end{tabularx} \\ \\ \relax
3:5 & \begin{tabularx}{0.7\textwidth}{X} 神說：「不要靠近這裡。把你腳上的鞋脫下來，因為你所站的地方是聖地」。 \end{tabularx} \\ \\ \relax
3:6 & \begin{tabularx}{0.7\textwidth}{X} 他又說：「我是你父親的神，是亞伯拉罕的神，以撒的神，雅各的神。」摩西蒙上臉，因為怕看神。 \end{tabularx} \\ \\
[1ex]
\hline
\hline
\end{longtable}


\newpage

\section{第77屆~港九培靈會~吳獻章~第二講}
\label{sec:320}
\textbf{臥虎藏龍}
\newline
\newline
連結: \href{https://www.hkbibleconference.org/session-message/view/320}{\texttt{https://www.hkbibleconference.org/session-message/view/320}}
\newline
\newline
\hyperref[sec:319]{< < < PREV SERMON < < <}
~
\hyperref[sec:toc]{[返目錄]}
~
\hyperref[sec:321]{> > > NEXT SERMON > > >}
\newline
\newline
撒母耳記上 16:1-16:13
\newline
\begin{longtable}{cl}
\hline
\hline
章節 & 經文 (和合本修訂版)\\
\hline
16:1 & \begin{tabularx}{0.7\textwidth}{X} 耶和華對撒母耳說：「我既厭棄掃羅作以色列的王，你為他悲傷要到幾時呢？你將膏油盛滿了角；來，我差遣你到伯利恆人耶西那裡去，因為我在他兒子中已看中了一個為我作王的。」 \end{tabularx} \\ \\ \relax
16:2 & \begin{tabularx}{0.7\textwidth}{X} 撒母耳說：「我怎麼能去呢？掃羅一聽見，就會殺我。」耶和華說：「你可以手裡牽一頭小母牛去，說：『我來是要向耶和華獻祭。』 \end{tabularx} \\ \\ \relax
16:3 & \begin{tabularx}{0.7\textwidth}{X} 你要請耶西來一同獻祭，我會指示你當做的事。我對你說的那個人，你要為我膏他。」 \end{tabularx} \\ \\ \relax
16:4 & \begin{tabularx}{0.7\textwidth}{X} 撒母耳遵照耶和華的話去做，來到伯利恆，城裡的長老都戰戰兢兢出來迎接他，有人問他說：「你是為平安來的嗎？」 \end{tabularx} \\ \\ \relax
16:5 & \begin{tabularx}{0.7\textwidth}{X} 他說：「為平安來的，我來是要向耶和華獻祭。你們要使自己分別為聖，來跟我一同獻祭。」撒母耳把耶西和他眾兒子分別為聖，請他們來一同獻祭。 \end{tabularx} \\ \\ \relax
16:6 & \begin{tabularx}{0.7\textwidth}{X} 他們來的時候，撒母耳看見以利押，就心裡說，耶和華的受膏者一定在耶和華面前了。 \end{tabularx} \\ \\ \relax
16:7 & \begin{tabularx}{0.7\textwidth}{X} 耶和華卻對撒母耳說：「不要只看他的外貌和他身材高大，我不揀選他。因為耶和華不像人看人，人是看外貌，耶和華是看內心。」 \end{tabularx} \\ \\ \relax
16:8 & \begin{tabularx}{0.7\textwidth}{X} 耶西叫亞比拿達從撒母耳面前經過，撒母耳說：「耶和華也不揀選他。」 \end{tabularx} \\ \\ \relax
16:9 & \begin{tabularx}{0.7\textwidth}{X} 耶西又叫沙瑪經過，撒母耳說：「耶和華也不揀選他。」 \end{tabularx} \\ \\ \relax
16:10 & \begin{tabularx}{0.7\textwidth}{X} 耶西叫他七個兒子都從撒母耳面前經過，撒母耳對耶西說：「這些都不是耶和華所揀選的。」 \end{tabularx} \\ \\ \relax
16:11 & \begin{tabularx}{0.7\textwidth}{X} 撒母耳對耶西說：「你的兒子都在這裡了嗎？」他說：「還有一個最小的，看哪，他正在放羊。」撒母耳對耶西說：「你派人去叫他來；他若不來這裡，我們必不坐席。」 \end{tabularx} \\ \\ \relax
16:12 & \begin{tabularx}{0.7\textwidth}{X} 耶西就派人去叫他來。他面色紅潤，雙目清秀，容貌俊美。耶和華說：「起來，膏他，因為這就是他了。」 \end{tabularx} \\ \\ \relax
16:13 & \begin{tabularx}{0.7\textwidth}{X} 撒母耳就用角裡的膏油，在他的兄長中膏了他。從這日起，耶和華的靈就大大感動大衛。撒母耳起身回拉瑪去了。 \end{tabularx} \\ \\
[1ex]
\hline
\hline
\end{longtable}


\newpage

\section{第77屆~港九培靈會~吳獻章~第三講}
\label{sec:321}
\textbf{權柄的失落──看父親的背影}
\newline
\newline
連結: \href{https://www.hkbibleconference.org/session-message/view/321}{\texttt{https://www.hkbibleconference.org/session-message/view/321}}
\newline
\newline
\hyperref[sec:320]{< < < PREV SERMON < < <}
~
\hyperref[sec:toc]{[返目錄]}
~
\hyperref[sec:322]{> > > NEXT SERMON > > >}
\newline
\newline


\newpage

\section{第77屆~港九培靈會~吳獻章~第四講}
\label{sec:322}
\textbf{被牧之羊}
\newline
\newline
連結: \href{https://www.hkbibleconference.org/session-message/view/322}{\texttt{https://www.hkbibleconference.org/session-message/view/322}}
\newline
\newline
\hyperref[sec:321]{< < < PREV SERMON < < <}
~
\hyperref[sec:toc]{[返目錄]}
~
\hyperref[sec:323]{> > > NEXT SERMON > > >}
\newline
\newline
詩篇 23:1-23:6
\newline
\begin{longtable}{cl}
\hline
\hline
章節 & 經文 (和合本修訂版)\\
\hline
23:1 & \begin{tabularx}{0.7\textwidth}{X} 耶和華是我的牧者，我必不致缺乏。 \end{tabularx} \\ \\ \relax
23:2 & \begin{tabularx}{0.7\textwidth}{X} 他使我躺臥在青草地上，領我在可安歇的水邊。 \end{tabularx} \\ \\ \relax
23:3 & \begin{tabularx}{0.7\textwidth}{X} 他使我的靈魂甦醒，為自己的名引導我走義路。 \end{tabularx} \\ \\ \relax
23:4 & \begin{tabularx}{0.7\textwidth}{X} 我雖然行過死蔭的幽谷，也不怕遭害，因為你與我同在；你的杖、你的竿，都安慰我。 \end{tabularx} \\ \\ \relax
23:5 & \begin{tabularx}{0.7\textwidth}{X} 在我敵人面前，你為我擺設筵席；你用油膏了我的頭，使我的福杯滿溢。 \end{tabularx} \\ \\ \relax
23:6 & \begin{tabularx}{0.7\textwidth}{X} 我一生一世必有恩惠慈愛隨著我；我且要住在耶和華的殿中，直到永遠。 \end{tabularx} \\ \\
[1ex]
\hline
\hline
\end{longtable}


\newpage

\section{第77屆~港九培靈會~吳獻章~第五講}
\label{sec:323}
\textbf{危險的任務}
\newline
\newline
連結: \href{https://www.hkbibleconference.org/session-message/view/323}{\texttt{https://www.hkbibleconference.org/session-message/view/323}}
\newline
\newline
\hyperref[sec:322]{< < < PREV SERMON < < <}
~
\hyperref[sec:toc]{[返目錄]}
~
\hyperref[sec:324]{> > > NEXT SERMON > > >}
\newline
\newline
以賽亞書 6:1-6:13
\newline
\begin{longtable}{cl}
\hline
\hline
章節 & 經文 (和合本修訂版)\\
\hline
6:1 & \begin{tabularx}{0.7\textwidth}{X} 當烏西雅王崩的那年，我看見主坐在高高的寶座上。他的衣裳下襬遮滿聖殿。 \end{tabularx} \\ \\ \relax
6:2 & \begin{tabularx}{0.7\textwidth}{X} 上有撒拉弗侍立，各有六個翅膀：兩個翅膀遮臉，兩個翅膀遮腳，兩個翅膀飛翔， \end{tabularx} \\ \\ \relax
6:3 & \begin{tabularx}{0.7\textwidth}{X} 彼此呼喊說：「聖哉！聖哉！聖哉！萬軍之耶和華；他的榮光遍滿全地！」 \end{tabularx} \\ \\ \relax
6:4 & \begin{tabularx}{0.7\textwidth}{X} 因呼喊者的聲音，門檻的根基震動，殿裡充滿了煙雲。 \end{tabularx} \\ \\ \relax
6:5 & \begin{tabularx}{0.7\textwidth}{X} 那時我說：「禍哉！我滅亡了！因為我是嘴唇不潔的人，住在嘴唇不潔的民中，又因我親眼看見大君王－萬軍之耶和華。」 \end{tabularx} \\ \\ \relax
6:6 & \begin{tabularx}{0.7\textwidth}{X} 有一撒拉弗向我飛來，手裡拿著燒紅的炭，是用火鉗從壇上取下來的， \end{tabularx} \\ \\ \relax
6:7 & \begin{tabularx}{0.7\textwidth}{X} 用炭沾我的口，說：「看哪，這炭沾了你的嘴唇，你的罪孽便除掉，你的罪惡就赦免了。」 \end{tabularx} \\ \\ \relax
6:8 & \begin{tabularx}{0.7\textwidth}{X} 我聽見主的聲音說：「我可以差遣誰呢？誰肯為我們去呢？」我說：「我在這裡，請差遣我！」 \end{tabularx} \\ \\ \relax
6:9 & \begin{tabularx}{0.7\textwidth}{X} 他說：「你去告訴這百姓說：『你們聽了又聽，卻不明白；看了又看，卻不曉得。』 \end{tabularx} \\ \\ \relax
6:10 & \begin{tabularx}{0.7\textwidth}{X} 要使這百姓心蒙油脂，耳朵發沉，眼睛昏花；恐怕他們眼睛看見，耳朵聽見，心裡明白，回轉過來，就得醫治。」 \end{tabularx} \\ \\ \relax
6:11 & \begin{tabularx}{0.7\textwidth}{X} 我就說：「主啊，這到幾時為止呢？」他說：「直到城鎮荒涼，無人居住，房屋空無一人，土地極其荒蕪； \end{tabularx} \\ \\ \relax
6:12 & \begin{tabularx}{0.7\textwidth}{X} 耶和華將人遷到遠方，國內被撇棄的土地很多。 \end{tabularx} \\ \\ \relax
6:13 & \begin{tabularx}{0.7\textwidth}{X} 國內剩下的人若還有十分之一，也必被吞滅。然而如同大樹與橡樹，雖被砍伐，殘幹卻仍存留，聖潔的苗裔是它的殘幹。」 \end{tabularx} \\ \\
[1ex]
\hline
\hline
\end{longtable}


\newpage

\section{第77屆~港九培靈會~吳獻章~第六講}
\label{sec:324}
\textbf{影響力的本質:名嘴背後}
\newline
\newline
連結: \href{https://www.hkbibleconference.org/session-message/view/324}{\texttt{https://www.hkbibleconference.org/session-message/view/324}}
\newline
\newline
\hyperref[sec:323]{< < < PREV SERMON < < <}
~
\hyperref[sec:toc]{[返目錄]}
~
\hyperref[sec:325]{> > > NEXT SERMON > > >}
\newline
\newline
約翰福音 1:6-1:7
\newline
\begin{longtable}{cl}
\hline
\hline
章節 & 經文 (和合本修訂版)\\
\hline
1:6 & \begin{tabularx}{0.7\textwidth}{X} 有一個人，是從神那裡差來的，名叫約翰。 \end{tabularx} \\ \\ \relax
1:7 & \begin{tabularx}{0.7\textwidth}{X} 這人來是為了作見證，是為那光作見證，要使眾人藉著他而信。 \end{tabularx} \\ \\
[1ex]
\hline
\hline
\end{longtable}


\newpage

\section{第77屆~港九培靈會~吳獻章~第七講}
\label{sec:325}
\textbf{危險的試探}
\newline
\newline
連結: \href{https://www.hkbibleconference.org/session-message/view/325}{\texttt{https://www.hkbibleconference.org/session-message/view/325}}
\newline
\newline
\hyperref[sec:324]{< < < PREV SERMON < < <}
~
\hyperref[sec:toc]{[返目錄]}
~
\hyperref[sec:326]{> > > NEXT SERMON > > >}
\newline
\newline
馬太福音 4:1-4:11
\newline
\begin{longtable}{cl}
\hline
\hline
章節 & 經文 (和合本修訂版)\\
\hline
4:1 & \begin{tabularx}{0.7\textwidth}{X} 當時，耶穌被聖靈引到曠野，受魔鬼的試探。 \end{tabularx} \\ \\ \relax
4:2 & \begin{tabularx}{0.7\textwidth}{X} 他禁食四十晝夜，後來就餓了。 \end{tabularx} \\ \\ \relax
4:3 & \begin{tabularx}{0.7\textwidth}{X} 那試探者進前來對他說：「你若是神的兒子，叫這些石頭變成食物吧。」 \end{tabularx} \\ \\ \relax
4:4 & \begin{tabularx}{0.7\textwidth}{X} 耶穌卻回答說：「經上記著：『人活著，不是單靠食物，乃是靠神口裡所出的一切話。』」 \end{tabularx} \\ \\ \relax
4:5 & \begin{tabularx}{0.7\textwidth}{X} 魔鬼就帶他進了聖城，叫他站在聖殿頂上， \end{tabularx} \\ \\ \relax
4:6 & \begin{tabularx}{0.7\textwidth}{X} 對他說：「你若是神的兒子，就跳下去！因為經上記著：『主要為你命令他的使者，用手托住你，免得你的腳碰在石頭上。』」 \end{tabularx} \\ \\ \relax
4:7 & \begin{tabularx}{0.7\textwidth}{X} 耶穌對他說：「經上又記著：『不可試探主—你的神。』」 \end{tabularx} \\ \\ \relax
4:8 & \begin{tabularx}{0.7\textwidth}{X} 魔鬼又帶他上了一座很高的山，將世上的萬國和萬國的榮華都指給他看， \end{tabularx} \\ \\ \relax
4:9 & \begin{tabularx}{0.7\textwidth}{X} 對他說：「你若俯伏拜我，我就把這一切賜給你。」 \end{tabularx} \\ \\ \relax
4:10 & \begin{tabularx}{0.7\textwidth}{X} 耶穌說：「撒但，退去！因為經上記著：『要拜主—你的神，惟獨事奉他。』」 \end{tabularx} \\ \\ \relax
4:11 & \begin{tabularx}{0.7\textwidth}{X} 於是，魔鬼離開了耶穌，立刻有天使來伺候他。 \end{tabularx} \\ \\
[1ex]
\hline
\hline
\end{longtable}


\newpage

\section{第77屆~港九培靈會~吳獻章~第八講}
\label{sec:326}
\textbf{基督的震撼教育}
\newline
\newline
連結: \href{https://www.hkbibleconference.org/session-message/view/326}{\texttt{https://www.hkbibleconference.org/session-message/view/326}}
\newline
\newline
\hyperref[sec:325]{< < < PREV SERMON < < <}
~
\hyperref[sec:toc]{[返目錄]}
~
\hyperref[sec:327]{> > > NEXT SERMON > > >}
\newline
\newline
馬可福音 15:37-16:8
\newline
\begin{longtable}{cl}
\hline
\hline
章節 & 經文 (和合本修訂版)\\
\hline
15:37 & \begin{tabularx}{0.7\textwidth}{X} 耶穌大喊一聲，氣就斷了。 \end{tabularx} \\ \\ \relax
15:38 & \begin{tabularx}{0.7\textwidth}{X} 殿的幔子從上到下裂為兩半。 \end{tabularx} \\ \\ \relax
15:39 & \begin{tabularx}{0.7\textwidth}{X} 對面站著的百夫長看見耶穌這樣斷氣，就說：「這人真是神的兒子！」 \end{tabularx} \\ \\ \relax
15:40 & \begin{tabularx}{0.7\textwidth}{X} 還有些婦女遠遠地觀看，其中有抹大拉的馬利亞，又有小雅各和約西的母親馬利亞，並有撒羅米， \end{tabularx} \\ \\ \relax
15:41 & \begin{tabularx}{0.7\textwidth}{X} 就是耶穌在加利利的時候，跟隨他、服事他的那些人，還有同耶穌上耶路撒冷的好些婦女。 \end{tabularx} \\ \\ \relax
15:42 & \begin{tabularx}{0.7\textwidth}{X} 到了晚上，因為這是預備日，就是安息日的前一日， \end{tabularx} \\ \\ \relax
15:43 & \begin{tabularx}{0.7\textwidth}{X} 有亞利馬太的約瑟前來，他是尊貴的議員，也是盼望著神國的，他放膽進去見彼拉多，請求要耶穌的身體。 \end{tabularx} \\ \\ \relax
15:44 & \begin{tabularx}{0.7\textwidth}{X} 彼拉多詫異耶穌已經死了，就叫百夫長來，問他耶穌是不是死了很久； \end{tabularx} \\ \\ \relax
15:45 & \begin{tabularx}{0.7\textwidth}{X} 既從百夫長得知實情，就把耶穌的身體賜給約瑟。 \end{tabularx} \\ \\ \relax
15:46 & \begin{tabularx}{0.7\textwidth}{X} 約瑟買了細麻布，把耶穌取下來，用細麻布裹好，安放在巖石中鑿出來的墓穴裡，又滾來一塊石頭擋住墓門。 \end{tabularx} \\ \\ \relax
15:47 & \begin{tabularx}{0.7\textwidth}{X} 抹大拉的馬利亞和約西的母親馬利亞都看見安放他的地方。 \end{tabularx} \\ \\ \relax
16:1 & \begin{tabularx}{0.7\textwidth}{X} 過了安息日，抹大拉的馬利亞、雅各的母親馬利亞，和撒羅米，買了香料，要去膏耶穌的身體。 \end{tabularx} \\ \\ \relax
16:2 & \begin{tabularx}{0.7\textwidth}{X} 七日的第一日清早，太陽出來後，她們來到墳墓那裡， \end{tabularx} \\ \\ \relax
16:3 & \begin{tabularx}{0.7\textwidth}{X} 彼此說：「誰要替我們把石頭從墓門滾開呢？」 \end{tabularx} \\ \\ \relax
16:4 & \begin{tabularx}{0.7\textwidth}{X} 她們抬頭一看，看見石頭已經滾開了，原來那石頭很大。 \end{tabularx} \\ \\ \relax
16:5 & \begin{tabularx}{0.7\textwidth}{X} 她們進了墳墓，看見一個年輕人坐在右邊，穿著白袍，就很驚奇。 \end{tabularx} \\ \\ \relax
16:6 & \begin{tabularx}{0.7\textwidth}{X} 那年輕人對她們說：「不要驚慌！你們尋找那釘十字架的拿撒勒人耶穌，他已經復活了，不在這裡。來看安放他的地方。 \end{tabularx} \\ \\ \relax
16:7 & \begin{tabularx}{0.7\textwidth}{X} 你們去，對他的門徒和彼得說：『他要比你們先到加利利去，在那裡你們會看見他，正如他從前所告訴你們的。』」 \end{tabularx} \\ \\ \relax
16:8 & \begin{tabularx}{0.7\textwidth}{X} 於是她們出來，從墳墓那裡逃走，又發抖又驚訝，甚麼也沒有告訴人，因為她們害怕。 \end{tabularx} \\ \\
[1ex]
\hline
\hline
\end{longtable}


\newpage

\section{第77屆~港九培靈會~吳獻章~第九講}
\label{sec:327}
\textbf{異象．團隊．委身──聖靈．人與宣教}
\newline
\newline
連結: \href{https://www.hkbibleconference.org/session-message/view/327}{\texttt{https://www.hkbibleconference.org/session-message/view/327}}
\newline
\newline
\hyperref[sec:326]{< < < PREV SERMON < < <}
~
\hyperref[sec:toc]{[返目錄]}
~
\hyperref[sec:328]{> > > NEXT SERMON > > >}
\newline
\newline
使徒行傳 16:1-16:10
\newline
\begin{longtable}{cl}
\hline
\hline
章節 & 經文 (和合本修訂版)\\
\hline
16:1 & \begin{tabularx}{0.7\textwidth}{X} 後來，保羅來到特庇，又到路司得。在那裡有一個門徒，名叫提摩太，是信主的猶太婦人的兒子，他父親卻是希臘人。 \end{tabularx} \\ \\ \relax
16:2 & \begin{tabularx}{0.7\textwidth}{X} 路司得和以哥念的弟兄都稱讚他。 \end{tabularx} \\ \\ \relax
16:3 & \begin{tabularx}{0.7\textwidth}{X} 保羅要帶他同去，只因那些地方的猶太人都知道他父親是希臘人，就給他行了割禮。 \end{tabularx} \\ \\ \relax
16:4 & \begin{tabularx}{0.7\textwidth}{X} 他們經過各城，把耶路撒冷使徒和長老所決定的規條交給門徒遵守。 \end{tabularx} \\ \\ \relax
16:5 & \begin{tabularx}{0.7\textwidth}{X} 於是眾教會信心越發堅固，人數天天增加。 \end{tabularx} \\ \\ \relax
16:6 & \begin{tabularx}{0.7\textwidth}{X} 因為聖靈禁止他們在亞細亞講道，他們就經過弗呂家、加拉太一帶地方。 \end{tabularx} \\ \\ \relax
16:7 & \begin{tabularx}{0.7\textwidth}{X} 到了每西亞的邊界，他們想要往庇推尼去，耶穌的靈卻不許。 \end{tabularx} \\ \\ \relax
16:8 & \begin{tabularx}{0.7\textwidth}{X} 他們就越過每西亞，下特羅亞去。 \end{tabularx} \\ \\ \relax
16:9 & \begin{tabularx}{0.7\textwidth}{X} 夜間，有異象向保羅顯現。有一個馬其頓人站著求他說：「請你過來，到馬其頓來幫助我們！」 \end{tabularx} \\ \\ \relax
16:10 & \begin{tabularx}{0.7\textwidth}{X} 保羅既看見這異象，我們就立即設法往馬其頓去，認為神呼召我們傳福音給那裡的人。 \end{tabularx} \\ \\
[1ex]
\hline
\hline
\end{longtable}


\newpage

\section{第77屆~港九培靈會~盧家駇~第一講}
\label{sec:328}
\textbf{起初的愛心}
\newline
\newline
連結: \href{https://www.hkbibleconference.org/session-message/view/328}{\texttt{https://www.hkbibleconference.org/session-message/view/328}}
\newline
\newline
\hyperref[sec:327]{< < < PREV SERMON < < <}
~
\hyperref[sec:toc]{[返目錄]}
~
\hyperref[sec:329]{> > > NEXT SERMON > > >}
\newline
\newline
啟示錄 2:1-2:7
\newline
\begin{longtable}{cl}
\hline
\hline
章節 & 經文 (和合本修訂版)\\
\hline
2:1 & \begin{tabularx}{0.7\textwidth}{X} 「你要寫信給以弗所教會的使者，說：『那右手拿著七顆星，在七個金燈臺中間行走的這樣說： \end{tabularx} \\ \\ \relax
2:2 & \begin{tabularx}{0.7\textwidth}{X} 我知道你的行為、勞碌、忍耐，也知道你不容忍惡人。你也曾察驗那自稱為使徒卻不是使徒的，看出他們是假的。 \end{tabularx} \\ \\ \relax
2:3 & \begin{tabularx}{0.7\textwidth}{X} 你能忍耐，曾為我的名勞苦而不困倦。 \end{tabularx} \\ \\ \relax
2:4 & \begin{tabularx}{0.7\textwidth}{X} 然而，有一件事我要責備你，就是你把起初的愛心拋棄了。 \end{tabularx} \\ \\ \relax
2:5 & \begin{tabularx}{0.7\textwidth}{X} 所以你要回想你是從哪裡墜落的，並且要悔改，做起初所做的工作。你若不悔改，我要到你那裡去，把你的燈臺從原處挪去。 \end{tabularx} \\ \\ \relax
2:6 & \begin{tabularx}{0.7\textwidth}{X} 然而你還有一件可取的事，就是你恨惡尼哥拉派的行為，這種行為也是我所恨惡的。 \end{tabularx} \\ \\ \relax
2:7 & \begin{tabularx}{0.7\textwidth}{X} 凡有耳朵的都應當聽聖靈向眾教會所說的話。得勝的，我必將神樂園中生命樹的果子賜給他吃。』」 \end{tabularx} \\ \\
[1ex]
\hline
\hline
\end{longtable}


\newpage

\section{第77屆~港九培靈會~盧家駇~第二講}
\label{sec:329}
\textbf{我愛我主教會}
\newline
\newline
連結: \href{https://www.hkbibleconference.org/session-message/view/329}{\texttt{https://www.hkbibleconference.org/session-message/view/329}}
\newline
\newline
\hyperref[sec:328]{< < < PREV SERMON < < <}
~
\hyperref[sec:toc]{[返目錄]}
~
\hyperref[sec:330]{> > > NEXT SERMON > > >}
\newline
\newline
以弗所書 2:11-2:13
\newline
\begin{longtable}{cl}
\hline
\hline
章節 & 經文 (和合本修訂版)\\
\hline
2:11 & \begin{tabularx}{0.7\textwidth}{X} 所以，你們要記得：從前你們按肉體是外邦人，是「沒受割禮的」；這名字是那些憑人手在肉身上「受割禮的人」所取的。 \end{tabularx} \\ \\ \relax
2:12 & \begin{tabularx}{0.7\textwidth}{X} 要記得那時候，你們與基督無關，與以色列選民團體隔絕，在所應許的約上是局外人，而且在世上沒有指望，沒有神。 \end{tabularx} \\ \\ \relax
2:13 & \begin{tabularx}{0.7\textwidth}{X} 從前你們是遠離神的人，如今卻在基督耶穌裡，靠著他的血，已經得以親近了。 \end{tabularx} \\ \\
[1ex]
\hline
\hline
\end{longtable}


\newpage

\section{第77屆~港九培靈會~盧家駇~第三講}
\label{sec:330}
\textbf{除去幽暗}
\newline
\newline
連結: \href{https://www.hkbibleconference.org/session-message/view/330}{\texttt{https://www.hkbibleconference.org/session-message/view/330}}
\newline
\newline
\hyperref[sec:329]{< < < PREV SERMON < < <}
~
\hyperref[sec:toc]{[返目錄]}
~
\hyperref[sec:331]{> > > NEXT SERMON > > >}
\newline
\newline
羅馬書 7:21-7:25
\newline
\begin{longtable}{cl}
\hline
\hline
章節 & 經文 (和合本修訂版)\\
\hline
7:21 & \begin{tabularx}{0.7\textwidth}{X} 我覺得有個律，就是我願意行善的時候，就有惡纏著我。 \end{tabularx} \\ \\ \relax
7:22 & \begin{tabularx}{0.7\textwidth}{X} 因為，按著我裡面的人，我喜歡神的律， \end{tabularx} \\ \\ \relax
7:23 & \begin{tabularx}{0.7\textwidth}{X} 但我看出肢體中另有個律和我內心的律交戰，把我擄去，使我附從那肢體中罪的律。 \end{tabularx} \\ \\ \relax
7:24 & \begin{tabularx}{0.7\textwidth}{X} 我真苦啊！誰能救我脫離這必死的身體呢？ \end{tabularx} \\ \\ \relax
7:25 & \begin{tabularx}{0.7\textwidth}{X} 感謝神，靠著我們的主耶穌基督就能！這樣看來，一方面，我內心順服神的律，另一方面，肉體卻順服罪的律了。 \end{tabularx} \\ \\
[1ex]
\hline
\hline
\end{longtable}


\newpage

\section{第77屆~港九培靈會~盧家駇~第四講}
\label{sec:331}
\textbf{心裡火熱}
\newline
\newline
連結: \href{https://www.hkbibleconference.org/session-message/view/331}{\texttt{https://www.hkbibleconference.org/session-message/view/331}}
\newline
\newline
\hyperref[sec:330]{< < < PREV SERMON < < <}
~
\hyperref[sec:toc]{[返目錄]}
~
\hyperref[sec:332]{> > > NEXT SERMON > > >}
\newline
\newline
羅馬書 12:1-12:11
\newline
\begin{longtable}{cl}
\hline
\hline
章節 & 經文 (和合本修訂版)\\
\hline
12:1 & \begin{tabularx}{0.7\textwidth}{X} 所以，弟兄們，我以神的慈悲勸你們，將身體獻上當作活祭，是聖潔的，是神所喜悅的，你們如此事奉乃是理所當然的。 \end{tabularx} \\ \\ \relax
12:2 & \begin{tabularx}{0.7\textwidth}{X} 不要效法這個世界，只要心意更新而變化，叫你們察驗何為神的善良、純全、可喜悅的旨意。 \end{tabularx} \\ \\ \relax
12:3 & \begin{tabularx}{0.7\textwidth}{X} 我憑著所賜我的恩對你們每一位說：不要把自己看得太高，要照著神所分給各人的信心來衡量，看得合乎中道。 \end{tabularx} \\ \\ \relax
12:4 & \begin{tabularx}{0.7\textwidth}{X} 正如我們一個身子上有好些肢體，肢體也不都有一樣的用處。 \end{tabularx} \\ \\ \relax
12:5 & \begin{tabularx}{0.7\textwidth}{X} 這樣，我們許多人在基督裡是一個身體，互相聯絡作肢體。 \end{tabularx} \\ \\ \relax
12:6 & \begin{tabularx}{0.7\textwidth}{X} 按著所得的恩典，我們各有不同的恩賜：或說預言，要按著信心的程度說預言； \end{tabularx} \\ \\ \relax
12:7 & \begin{tabularx}{0.7\textwidth}{X} 或服事的，要專一服事；或教導的，要專一教導； \end{tabularx} \\ \\ \relax
12:8 & \begin{tabularx}{0.7\textwidth}{X} 或勸勉的，要專一勸勉；施捨的，要誠實；治理的，要殷勤；憐憫人的，要樂意。 \end{tabularx} \\ \\ \relax
12:9 & \begin{tabularx}{0.7\textwidth}{X} 愛，不可虛假；惡，要厭惡；善，要親近。 \end{tabularx} \\ \\ \relax
12:10 & \begin{tabularx}{0.7\textwidth}{X} 愛弟兄，要相親相愛；恭敬人，要彼此推讓； \end{tabularx} \\ \\ \relax
12:11 & \begin{tabularx}{0.7\textwidth}{X} 殷勤，不可懶惰。要靈裡火熱；常常服侍主。 \end{tabularx} \\ \\
[1ex]
\hline
\hline
\end{longtable}


\newpage

\section{第77屆~港九培靈會~盧家駇~第五講}
\label{sec:332}
\textbf{完全奉獻}
\newline
\newline
連結: \href{https://www.hkbibleconference.org/session-message/view/332}{\texttt{https://www.hkbibleconference.org/session-message/view/332}}
\newline
\newline
\hyperref[sec:331]{< < < PREV SERMON < < <}
~
\hyperref[sec:toc]{[返目錄]}
~
\hyperref[sec:333]{> > > NEXT SERMON > > >}
\newline
\newline
使徒行傳 4:36-4:37
\newline
\begin{longtable}{cl}
\hline
\hline
章節 & 經文 (和合本修訂版)\\
\hline
4:36 & \begin{tabularx}{0.7\textwidth}{X} 有一個利未人，名叫約瑟，使徒稱他為巴拿巴（巴拿巴翻出來就是安慰之子），生在塞浦路斯。 \end{tabularx} \\ \\ \relax
4:37 & \begin{tabularx}{0.7\textwidth}{X} 他有田地，也賣了，把錢拿來，放在使徒腳前。 \end{tabularx} \\ \\
[1ex]
\hline
\hline
\end{longtable}


\newpage

\section{第77屆~港九培靈會~盧家駇~第六講}
\label{sec:333}
\textbf{新命令}
\newline
\newline
連結: \href{https://www.hkbibleconference.org/session-message/view/333}{\texttt{https://www.hkbibleconference.org/session-message/view/333}}
\newline
\newline
\hyperref[sec:332]{< < < PREV SERMON < < <}
~
\hyperref[sec:toc]{[返目錄]}
~
\hyperref[sec:334]{> > > NEXT SERMON > > >}
\newline
\newline
約翰福音 13:31-13:35
\newline
\begin{longtable}{cl}
\hline
\hline
章節 & 經文 (和合本修訂版)\\
\hline
13:31 & \begin{tabularx}{0.7\textwidth}{X} 猶大出去後，耶穌說：「如今人子得了榮耀，神在人子身上也得了榮耀。 \end{tabularx} \\ \\ \relax
13:32 & \begin{tabularx}{0.7\textwidth}{X} 如果神因人子得了榮耀，神也要因自己榮耀人子，並且要立刻榮耀他。 \end{tabularx} \\ \\ \relax
13:33 & \begin{tabularx}{0.7\textwidth}{X} 孩子們！我與你們同在的時候不多了；你們會找我，但我所去的地方，你們不能去。這話我曾對猶太人說過，現在也照樣對你們說。 \end{tabularx} \\ \\ \relax
13:34 & \begin{tabularx}{0.7\textwidth}{X} 我賜給你們一條新命令，乃是叫你們彼此相愛；我怎樣愛你們，你們也要怎樣彼此相愛。 \end{tabularx} \\ \\ \relax
13:35 & \begin{tabularx}{0.7\textwidth}{X} 你們若彼此相愛，眾人因此就認出你們是我的門徒了。」 \end{tabularx} \\ \\
[1ex]
\hline
\hline
\end{longtable}


\newpage

\section{第77屆~港九培靈會~盧家駇~第七講}
\label{sec:334}
\textbf{笑傲逆境}
\newline
\newline
連結: \href{https://www.hkbibleconference.org/session-message/view/334}{\texttt{https://www.hkbibleconference.org/session-message/view/334}}
\newline
\newline
\hyperref[sec:333]{< < < PREV SERMON < < <}
~
\hyperref[sec:toc]{[返目錄]}
~
\hyperref[sec:335]{> > > NEXT SERMON > > >}
\newline
\newline
腓立比書 4:10-4:14
\newline
\begin{longtable}{cl}
\hline
\hline
章節 & 經文 (和合本修訂版)\\
\hline
4:10 & \begin{tabularx}{0.7\textwidth}{X} 我靠主大大喜樂，因為你們關懷我的心如今又表現了出來；其實你們一直都關懷我，只是沒有機會罷了。 \end{tabularx} \\ \\ \relax
4:11 & \begin{tabularx}{0.7\textwidth}{X} 我並不是因缺乏而說這話，因為我已經學會無論在甚麼景況都可以知足。 \end{tabularx} \\ \\ \relax
4:12 & \begin{tabularx}{0.7\textwidth}{X} 我知道怎樣處卑賤，也知道怎樣處豐富；或飽足或飢餓，或有餘或缺乏，任何事情，任何景況，我都得了祕訣。 \end{tabularx} \\ \\ \relax
4:13 & \begin{tabularx}{0.7\textwidth}{X} 我靠著那加給我力量的，凡事都能做。 \end{tabularx} \\ \\ \relax
4:14 & \begin{tabularx}{0.7\textwidth}{X} 然而，你們能和我分擔憂患是一件好事。 \end{tabularx} \\ \\
[1ex]
\hline
\hline
\end{longtable}


\newpage

\section{第77屆~港九培靈會~盧家駇~第八講}
\label{sec:335}
\textbf{骨肉之親}
\newline
\newline
連結: \href{https://www.hkbibleconference.org/session-message/view/335}{\texttt{https://www.hkbibleconference.org/session-message/view/335}}
\newline
\newline
\hyperref[sec:334]{< < < PREV SERMON < < <}
~
\hyperref[sec:toc]{[返目錄]}
~
\hyperref[sec:336]{> > > NEXT SERMON > > >}
\newline
\newline
羅馬書 9:1-9:3
\newline
\begin{longtable}{cl}
\hline
\hline
章節 & 經文 (和合本修訂版)\\
\hline
9:1 & \begin{tabularx}{0.7\textwidth}{X} 我在基督裡說真話，不說謊話；我的良心被聖靈感動為我作證。 \end{tabularx} \\ \\ \relax
9:2 & \begin{tabularx}{0.7\textwidth}{X} 我非常憂愁，心裡時常傷痛。 \end{tabularx} \\ \\ \relax
9:3 & \begin{tabularx}{0.7\textwidth}{X} 為我弟兄，我骨肉之親，就是自己被詛咒，與基督分離，我也願意。 \end{tabularx} \\ \\
[1ex]
\hline
\hline
\end{longtable}


\newpage

\section{第77屆~港九培靈會~盧家駇~第九講}
\label{sec:336}
\textbf{遍傳地極}
\newline
\newline
連結: \href{https://www.hkbibleconference.org/session-message/view/336}{\texttt{https://www.hkbibleconference.org/session-message/view/336}}
\newline
\newline
\hyperref[sec:335]{< < < PREV SERMON < < <}
~
\hyperref[sec:toc]{[返目錄]}
~
\hyperref[sec:337]{> > > NEXT SERMON > > >}
\newline
\newline
馬太福音 28:19-28:20
\newline
\begin{longtable}{cl}
\hline
\hline
章節 & 經文 (和合本修訂版)\\
\hline
28:19 & \begin{tabularx}{0.7\textwidth}{X} 所以，你們要去，使萬民作我的門徒，奉父、子、聖靈的名給他們施洗， \end{tabularx} \\ \\ \relax
28:20 & \begin{tabularx}{0.7\textwidth}{X} 凡我所吩咐你們的，都教導他們遵守。看哪，我天天與你們同在，直到世代的終結。」 \end{tabularx} \\ \\
[1ex]
\hline
\hline
\end{longtable}


\newpage

\section{第77屆~港九培靈會~盧家駇~第十講}
\label{sec:337}
\textbf{候主再臨}
\newline
\newline
連結: \href{https://www.hkbibleconference.org/session-message/view/337}{\texttt{https://www.hkbibleconference.org/session-message/view/337}}
\newline
\newline
\hyperref[sec:336]{< < < PREV SERMON < < <}
~
\hyperref[sec:toc]{[返目錄]}
~
\hyperref[sec:310]{> > > NEXT SERMON > > >}
\newline
\newline
帖撒羅尼迦前書 4:13-4:18
\newline
\begin{longtable}{cl}
\hline
\hline
章節 & 經文 (和合本修訂版)\\
\hline
4:13 & \begin{tabularx}{0.7\textwidth}{X} 弟兄們，至於已睡了的人，我們不願意你們不知道，恐怕你們憂傷，像那些沒有指望的人一樣。 \end{tabularx} \\ \\ \relax
4:14 & \begin{tabularx}{0.7\textwidth}{X} 既然我們信耶穌死了，復活了，那些已經在耶穌裡睡了的人，神也必將他們與耶穌一同帶來。 \end{tabularx} \\ \\ \relax
4:15 & \begin{tabularx}{0.7\textwidth}{X} 我們照主的話告訴你們一件事：我們這活著還存留到主來臨的人，絕不會在那已經睡了的人之先。 \end{tabularx} \\ \\ \relax
4:16 & \begin{tabularx}{0.7\textwidth}{X} 因為，召集令一發，天使長的呼聲一叫，神的號角一吹，主必親自從天降臨；那在基督裡死了的人必先復活， \end{tabularx} \\ \\ \relax
4:17 & \begin{tabularx}{0.7\textwidth}{X} 然後我們這些活著還存留的人必和他們一同被提到雲裡，在空中與主相會。這樣，我們就要和主永遠同在。 \end{tabularx} \\ \\ \relax
4:18 & \begin{tabularx}{0.7\textwidth}{X} 所以，你們當用這些話彼此勸勉。 \end{tabularx} \\ \\
[1ex]
\hline
\hline
\end{longtable}


\newpage

\section{第77屆~港九培靈會~鮑維均~第一講}
\label{sec:310}
\textbf{神與我們同在}
\newline
\newline
連結: \href{https://www.hkbibleconference.org/session-message/view/310}{\texttt{https://www.hkbibleconference.org/session-message/view/310}}
\newline
\newline
\hyperref[sec:337]{< < < PREV SERMON < < <}
~
\hyperref[sec:toc]{[返目錄]}
~
\hyperref[sec:311]{> > > NEXT SERMON > > >}
\newline
\newline
馬太福音 1:18-1:25
\newline
\begin{longtable}{cl}
\hline
\hline
章節 & 經文 (和合本修訂版)\\
\hline
1:18 & \begin{tabularx}{0.7\textwidth}{X} 耶穌基督降生的事記在下面：他母親馬利亞已經許配給約瑟，還沒有迎娶，馬利亞就從聖靈懷了孕。 \end{tabularx} \\ \\ \relax
1:19 & \begin{tabularx}{0.7\textwidth}{X} 她丈夫約瑟是個義人，不願意當眾羞辱她，想要暗地裡把她休了。 \end{tabularx} \\ \\ \relax
1:20 & \begin{tabularx}{0.7\textwidth}{X} 正考慮這些事的時候，忽然主的使者在約瑟夢中向他顯現，說：「大衛的子孫約瑟，不要怕，把你的妻子馬利亞娶過來，因她所懷的孕是從聖靈來的。 \end{tabularx} \\ \\ \relax
1:21 & \begin{tabularx}{0.7\textwidth}{X} 她將要生一個兒子，你要給他起名叫耶穌，因他要將自己的百姓從罪惡裡救出來。」 \end{tabularx} \\ \\ \relax
1:22 & \begin{tabularx}{0.7\textwidth}{X} 這整件事的發生，是要應驗主藉先知所說的話： \end{tabularx} \\ \\ \relax
1:23 & \begin{tabularx}{0.7\textwidth}{X} 「必有童女懷孕生子；人要稱他的名為以馬內利。」（以馬內利翻出來就是「神與我們同在」。） \end{tabularx} \\ \\ \relax
1:24 & \begin{tabularx}{0.7\textwidth}{X} 約瑟醒來，就遵照主的使者的吩咐把妻子娶過來； \end{tabularx} \\ \\ \relax
1:25 & \begin{tabularx}{0.7\textwidth}{X} 但是沒有和她同房，直到她生了兒子，就給他起名叫耶穌。 \end{tabularx} \\ \\
[1ex]
\hline
\hline
\end{longtable}


\newpage

\section{第77屆~港九培靈會~鮑維均~第二講}
\label{sec:311}
\textbf{猶太人的王}
\newline
\newline
連結: \href{https://www.hkbibleconference.org/session-message/view/311}{\texttt{https://www.hkbibleconference.org/session-message/view/311}}
\newline
\newline
\hyperref[sec:310]{< < < PREV SERMON < < <}
~
\hyperref[sec:toc]{[返目錄]}
~
\hyperref[sec:312]{> > > NEXT SERMON > > >}
\newline
\newline
馬太福音 2:1-2:12
\newline
\begin{longtable}{cl}
\hline
\hline
章節 & 經文 (和合本修訂版)\\
\hline
2:1 & \begin{tabularx}{0.7\textwidth}{X} 在希律作王的時候，耶穌生在猶太的伯利恆。有幾個博學之士從東方來到耶路撒冷，說： \end{tabularx} \\ \\ \relax
2:2 & \begin{tabularx}{0.7\textwidth}{X} 「那生下來作猶太人之王的在哪裡？我們在東方看見他的星，特來拜他。」 \end{tabularx} \\ \\ \relax
2:3 & \begin{tabularx}{0.7\textwidth}{X} 希律王聽見了，就心裡不安；耶路撒冷全城的人也都不安。 \end{tabularx} \\ \\ \relax
2:4 & \begin{tabularx}{0.7\textwidth}{X} 他就召集了祭司長和民間的文士，問他們：「基督該生在哪裡？」 \end{tabularx} \\ \\ \relax
2:5 & \begin{tabularx}{0.7\textwidth}{X} 他們說：「在猶太的伯利恆。因為有先知記著： \end{tabularx} \\ \\ \relax
2:6 & \begin{tabularx}{0.7\textwidth}{X} 『猶大地的伯利恆啊，你在猶大諸城中並不是最小的；因為將來有一位統治者要從你那裡出來，牧養我以色列民。』」 \end{tabularx} \\ \\ \relax
2:7 & \begin{tabularx}{0.7\textwidth}{X} 於是，希律暗地裡召了博學之士來，查問那星是甚麼時候出現的， \end{tabularx} \\ \\ \relax
2:8 & \begin{tabularx}{0.7\textwidth}{X} 就派他們往伯利恆去，說：「你們去仔細尋訪那小孩子，找到了就來報信，我也好去拜他。」 \end{tabularx} \\ \\ \relax
2:9 & \begin{tabularx}{0.7\textwidth}{X} 他們聽了王的話就去了。忽然，在東方所看到的那顆星在前面引領他們，一直行到小孩子所在地方的上方就停住了。 \end{tabularx} \\ \\ \relax
2:10 & \begin{tabularx}{0.7\textwidth}{X} 他們看見那星，就非常歡喜； \end{tabularx} \\ \\ \relax
2:11 & \begin{tabularx}{0.7\textwidth}{X} 進了房子，看見小孩子和他母親馬利亞，就俯伏拜那小孩子，揭開寶盒，拿出黃金、乳香、沒藥，作為禮物獻給他。 \end{tabularx} \\ \\ \relax
2:12 & \begin{tabularx}{0.7\textwidth}{X} 因為在夢中得到主的指示，不要回去見希律，他們就從別的路回自己的家鄉去了。 \end{tabularx} \\ \\
[1ex]
\hline
\hline
\end{longtable}


\newpage

\section{第77屆~港九培靈會~鮑維均~第三講}
\label{sec:312}
\textbf{曠野之旅}
\newline
\newline
連結: \href{https://www.hkbibleconference.org/session-message/view/312}{\texttt{https://www.hkbibleconference.org/session-message/view/312}}
\newline
\newline
\hyperref[sec:311]{< < < PREV SERMON < < <}
~
\hyperref[sec:toc]{[返目錄]}
~
\hyperref[sec:313]{> > > NEXT SERMON > > >}
\newline
\newline
馬太福音 4:1-4:11
\newline
\begin{longtable}{cl}
\hline
\hline
章節 & 經文 (和合本修訂版)\\
\hline
4:1 & \begin{tabularx}{0.7\textwidth}{X} 當時，耶穌被聖靈引到曠野，受魔鬼的試探。 \end{tabularx} \\ \\ \relax
4:2 & \begin{tabularx}{0.7\textwidth}{X} 他禁食四十晝夜，後來就餓了。 \end{tabularx} \\ \\ \relax
4:3 & \begin{tabularx}{0.7\textwidth}{X} 那試探者進前來對他說：「你若是神的兒子，叫這些石頭變成食物吧。」 \end{tabularx} \\ \\ \relax
4:4 & \begin{tabularx}{0.7\textwidth}{X} 耶穌卻回答說：「經上記著：『人活著，不是單靠食物，乃是靠神口裡所出的一切話。』」 \end{tabularx} \\ \\ \relax
4:5 & \begin{tabularx}{0.7\textwidth}{X} 魔鬼就帶他進了聖城，叫他站在聖殿頂上， \end{tabularx} \\ \\ \relax
4:6 & \begin{tabularx}{0.7\textwidth}{X} 對他說：「你若是神的兒子，就跳下去！因為經上記著：『主要為你命令他的使者，用手托住你，免得你的腳碰在石頭上。』」 \end{tabularx} \\ \\ \relax
4:7 & \begin{tabularx}{0.7\textwidth}{X} 耶穌對他說：「經上又記著：『不可試探主—你的神。』」 \end{tabularx} \\ \\ \relax
4:8 & \begin{tabularx}{0.7\textwidth}{X} 魔鬼又帶他上了一座很高的山，將世上的萬國和萬國的榮華都指給他看， \end{tabularx} \\ \\ \relax
4:9 & \begin{tabularx}{0.7\textwidth}{X} 對他說：「你若俯伏拜我，我就把這一切賜給你。」 \end{tabularx} \\ \\ \relax
4:10 & \begin{tabularx}{0.7\textwidth}{X} 耶穌說：「撒但，退去！因為經上記著：『要拜主—你的神，惟獨事奉他。』」 \end{tabularx} \\ \\ \relax
4:11 & \begin{tabularx}{0.7\textwidth}{X} 於是，魔鬼離開了耶穌，立刻有天使來伺候他。 \end{tabularx} \\ \\
[1ex]
\hline
\hline
\end{longtable}


\newpage

\section{第77屆~港九培靈會~鮑維均~第四講}
\label{sec:313}
\textbf{願你的國降臨}
\newline
\newline
連結: \href{https://www.hkbibleconference.org/session-message/view/313}{\texttt{https://www.hkbibleconference.org/session-message/view/313}}
\newline
\newline
\hyperref[sec:312]{< < < PREV SERMON < < <}
~
\hyperref[sec:toc]{[返目錄]}
~
\hyperref[sec:314]{> > > NEXT SERMON > > >}
\newline
\newline
馬太福音 6:5-6:13
\newline
\begin{longtable}{cl}
\hline
\hline
章節 & 經文 (和合本修訂版)\\
\hline
6:5 & \begin{tabularx}{0.7\textwidth}{X} 「你們禱告的時候，不可像那假冒為善的人，愛站在會堂裡和十字路口禱告，故意讓人看見。我實在告訴你們，他們已經得了他們的賞賜。 \end{tabularx} \\ \\ \relax
6:6 & \begin{tabularx}{0.7\textwidth}{X} 你禱告的時候，要進入內室，關上門，向那在隱祕中的父禱告；你父在隱祕中察看，必將賞賜你。 \end{tabularx} \\ \\ \relax
6:7 & \begin{tabularx}{0.7\textwidth}{X} 你們禱告，不可像外邦人那樣重複一些空話，他們以為話多了必蒙垂聽。 \end{tabularx} \\ \\ \relax
6:8 & \begin{tabularx}{0.7\textwidth}{X} 你們不可效法他們。因為在你們祈求以前，你們所需要的，你們的父早已知道了。」 \end{tabularx} \\ \\ \relax
6:9 & \begin{tabularx}{0.7\textwidth}{X} 「所以，你們要這樣禱告：『我們在天上的父：願人都尊你的名為聖。 \end{tabularx} \\ \\ \relax
6:10 & \begin{tabularx}{0.7\textwidth}{X} 願你的國降臨；願你的旨意行在地上，如同行在天上。 \end{tabularx} \\ \\ \relax
6:11 & \begin{tabularx}{0.7\textwidth}{X} 我們日用的飲食，今日賜給我們。 \end{tabularx} \\ \\ \relax
6:12 & \begin{tabularx}{0.7\textwidth}{X} 免我們的債，如同我們免了人的債。 \end{tabularx} \\ \\ \relax
6:13 & \begin{tabularx}{0.7\textwidth}{X} 不叫我們陷入試探；救我們脫離那惡者。因為國度、權柄、榮耀，全是你的，直到永遠。阿們！』 \end{tabularx} \\ \\
[1ex]
\hline
\hline
\end{longtable}


\newpage

\section{第77屆~港九培靈會~鮑維均~第五講}
\label{sec:314}
\textbf{來跟從我}
\newline
\newline
連結: \href{https://www.hkbibleconference.org/session-message/view/314}{\texttt{https://www.hkbibleconference.org/session-message/view/314}}
\newline
\newline
\hyperref[sec:313]{< < < PREV SERMON < < <}
~
\hyperref[sec:toc]{[返目錄]}
~
\hyperref[sec:315]{> > > NEXT SERMON > > >}
\newline
\newline
馬太福音 9:18-9:26
\newline
\begin{longtable}{cl}
\hline
\hline
章節 & 經文 (和合本修訂版)\\
\hline
9:18 & \begin{tabularx}{0.7\textwidth}{X} 耶穌說這些話的時候，有一個會堂主管來，向他下跪，說：「我女兒剛死了，求你去按手在她身上，她就會活過來。」 \end{tabularx} \\ \\ \relax
9:19 & \begin{tabularx}{0.7\textwidth}{X} 耶穌就起來跟他去；門徒也跟了去。 \end{tabularx} \\ \\ \relax
9:20 & \begin{tabularx}{0.7\textwidth}{X} 這時，有一個女人，患了經血不止的病有十二年，來到耶穌背後，摸他的衣裳繸子； \end{tabularx} \\ \\ \relax
9:21 & \begin{tabularx}{0.7\textwidth}{X} 因為她心裡說：「我只要摸他的衣裳，就會痊癒。」 \end{tabularx} \\ \\ \relax
9:22 & \begin{tabularx}{0.7\textwidth}{X} 耶穌轉過來，看見她，就說：「女兒，放心！你的信救了你。」從那時起，這女人就痊癒了。 \end{tabularx} \\ \\ \relax
9:23 & \begin{tabularx}{0.7\textwidth}{X} 耶穌到了會堂主管的家裡，看見吹鼓手和亂哄哄的一群人， \end{tabularx} \\ \\ \relax
9:24 & \begin{tabularx}{0.7\textwidth}{X} 就說：「退去吧！這女孩不是死了，而是睡著了。」他們就嘲笑他。 \end{tabularx} \\ \\ \relax
9:25 & \begin{tabularx}{0.7\textwidth}{X} 眾人被趕出後，耶穌就進去，拉著她的手，女孩就起來了。 \end{tabularx} \\ \\ \relax
9:26 & \begin{tabularx}{0.7\textwidth}{X} 於是這消息傳遍了那地方。 \end{tabularx} \\ \\
[1ex]
\hline
\hline
\end{longtable}


\newpage

\section{第77屆~港九培靈會~鮑維均~第六講}
\label{sec:315}
\textbf{看卻看不見}
\newline
\newline
連結: \href{https://www.hkbibleconference.org/session-message/view/315}{\texttt{https://www.hkbibleconference.org/session-message/view/315}}
\newline
\newline
\hyperref[sec:314]{< < < PREV SERMON < < <}
~
\hyperref[sec:toc]{[返目錄]}
~
\hyperref[sec:316]{> > > NEXT SERMON > > >}
\newline
\newline
馬太福音 13:1-13:23
\newline
\begin{longtable}{cl}
\hline
\hline
章節 & 經文 (和合本修訂版)\\
\hline
13:1 & \begin{tabularx}{0.7\textwidth}{X} 就在那天，耶穌從房子裡出來，坐在海邊。 \end{tabularx} \\ \\ \relax
13:2 & \begin{tabularx}{0.7\textwidth}{X} 有一大群人到他那裡聚集，他只好上船坐下，眾人都站在岸上。 \end{tabularx} \\ \\ \relax
13:3 & \begin{tabularx}{0.7\textwidth}{X} 他用比喻對他們講了許多話。他說：「有一個撒種的出去撒種。 \end{tabularx} \\ \\ \relax
13:4 & \begin{tabularx}{0.7\textwidth}{X} 他撒的時候，有的落在路旁，飛鳥來把它們吃掉了。 \end{tabularx} \\ \\ \relax
13:5 & \begin{tabularx}{0.7\textwidth}{X} 有的落在土淺的石頭地上，因為土不深，很快就長出苗來， \end{tabularx} \\ \\ \relax
13:6 & \begin{tabularx}{0.7\textwidth}{X} 太陽出來一曬，因為沒有根就枯乾了。 \end{tabularx} \\ \\ \relax
13:7 & \begin{tabularx}{0.7\textwidth}{X} 有的落在荊棘裡，荊棘長起來，把它擠住了。 \end{tabularx} \\ \\ \relax
13:8 & \begin{tabularx}{0.7\textwidth}{X} 又有的落在好土裡，就結出果實，有一百倍的，有六十倍的，有三十倍的。 \end{tabularx} \\ \\ \relax
13:9 & \begin{tabularx}{0.7\textwidth}{X} 有耳的，就應當聽！」 \end{tabularx} \\ \\ \relax
13:10 & \begin{tabularx}{0.7\textwidth}{X} 門徒進前來問耶穌：「對眾人講話，為甚麼用比喻呢？」 \end{tabularx} \\ \\ \relax
13:11 & \begin{tabularx}{0.7\textwidth}{X} 耶穌回答他們說：「因為天國的奧祕只讓你們知道，不讓他們知道。 \end{tabularx} \\ \\ \relax
13:12 & \begin{tabularx}{0.7\textwidth}{X} 凡有的，還要給他，讓他有餘；凡沒有的，連他所有的也要奪去。 \end{tabularx} \\ \\ \relax
13:13 & \begin{tabularx}{0.7\textwidth}{X} 我之所以用比喻對他們講，是因為他們看卻看不清，聽卻聽不見，也不明白。 \end{tabularx} \\ \\ \relax
13:14 & \begin{tabularx}{0.7\textwidth}{X} 在他們身上，正應驗了以賽亞的預言：『你們聽了又聽，卻不明白，看了又看，卻看不清。 \end{tabularx} \\ \\ \relax
13:15 & \begin{tabularx}{0.7\textwidth}{X} 因為這百姓的心麻木，耳朵發沉，眼睛閉著，免得眼睛看見，耳朵聽見，心裡明白，回轉過來，我會醫治他們。』 \end{tabularx} \\ \\ \relax
13:16 & \begin{tabularx}{0.7\textwidth}{X} 但你們的眼睛是有福的，因為看得見；你們的耳朵也是有福的，因為聽得見。 \end{tabularx} \\ \\ \relax
13:17 & \begin{tabularx}{0.7\textwidth}{X} 我實在告訴你們，從前有許多先知和義人要看你們所看的，卻沒有看見；要聽你們所聽的，卻沒有聽見。」 \end{tabularx} \\ \\ \relax
13:18 & \begin{tabularx}{0.7\textwidth}{X} 「所以，你們要聽這撒種的比喻。 \end{tabularx} \\ \\ \relax
13:19 & \begin{tabularx}{0.7\textwidth}{X} 凡聽見天國的道而不明白的，那惡者就來，把撒在他心裡的奪了去；這就是撒在路旁的了。 \end{tabularx} \\ \\ \relax
13:20 & \begin{tabularx}{0.7\textwidth}{X} 撒在石頭地上的，就是人聽了道，立刻歡喜領受， \end{tabularx} \\ \\ \relax
13:21 & \begin{tabularx}{0.7\textwidth}{X} 只因心裡沒有根，不過是暫時的，一旦為道遭受患難或迫害，立刻就跌倒。 \end{tabularx} \\ \\ \relax
13:22 & \begin{tabularx}{0.7\textwidth}{X} 撒在荊棘裡的，就是人聽了道，後來有世上的憂慮、錢財的迷惑把道擠住了，結不出果實。 \end{tabularx} \\ \\ \relax
13:23 & \begin{tabularx}{0.7\textwidth}{X} 撒在好土裡的，就是人聽了道，明白了，後來結了果實，有一百倍的，有六十倍的，有三十倍的。」 \end{tabularx} \\ \\
[1ex]
\hline
\hline
\end{longtable}


\newpage

\section{第77屆~港九培靈會~鮑維均~第七講}
\label{sec:316}
\textbf{七十個七次}
\newline
\newline
連結: \href{https://www.hkbibleconference.org/session-message/view/316}{\texttt{https://www.hkbibleconference.org/session-message/view/316}}
\newline
\newline
\hyperref[sec:315]{< < < PREV SERMON < < <}
~
\hyperref[sec:toc]{[返目錄]}
~
\hyperref[sec:317]{> > > NEXT SERMON > > >}
\newline
\newline
馬太福音 18:21-18:35
\newline
\begin{longtable}{cl}
\hline
\hline
章節 & 經文 (和合本修訂版)\\
\hline
18:21 & \begin{tabularx}{0.7\textwidth}{X} 那時，彼得進前來，對耶穌說：「主啊，我弟兄得罪我，我當饒恕他幾次呢？到七次夠嗎？」 \end{tabularx} \\ \\ \relax
18:22 & \begin{tabularx}{0.7\textwidth}{X} 耶穌說：「我告訴你，不是到七次，而是到七十個七次。 \end{tabularx} \\ \\ \relax
18:23 & \begin{tabularx}{0.7\textwidth}{X} 因為天國好像一個王要和他僕人算賬。 \end{tabularx} \\ \\ \relax
18:24 & \begin{tabularx}{0.7\textwidth}{X} 他開始算的時候，有人帶了一個欠一萬他連得的僕人來。 \end{tabularx} \\ \\ \relax
18:25 & \begin{tabularx}{0.7\textwidth}{X} 因為他沒有甚麼償還之物，主人下令把他和他妻子兒女，以及一切所有的都賣了來償還。 \end{tabularx} \\ \\ \relax
18:26 & \begin{tabularx}{0.7\textwidth}{X} 那僕人就俯伏向他叩頭，說：『寬容我吧，我都會還你的。』 \end{tabularx} \\ \\ \relax
18:27 & \begin{tabularx}{0.7\textwidth}{X} 那僕人的主人就動了慈心，把他釋放了，並且免了他的債。 \end{tabularx} \\ \\ \relax
18:28 & \begin{tabularx}{0.7\textwidth}{X} 那僕人出來，遇見一個欠他一百個銀幣的同伴，就揪著他，扼住他的喉嚨，說：『把你所欠的還我！』 \end{tabularx} \\ \\ \relax
18:29 & \begin{tabularx}{0.7\textwidth}{X} 他的同伴就俯伏央求他，說：『寬容我吧，我會還你的。』 \end{tabularx} \\ \\ \relax
18:30 & \begin{tabularx}{0.7\textwidth}{X} 他不肯，卻把他下在監裡，直到他還了所欠的債。 \end{tabularx} \\ \\ \relax
18:31 & \begin{tabularx}{0.7\textwidth}{X} 同伴們看見他所做的事就很悲憤，把這一切的事都告訴了主人。 \end{tabularx} \\ \\ \relax
18:32 & \begin{tabularx}{0.7\textwidth}{X} 於是主人叫了他來，對他說：『你這惡奴才！你央求我，我就把你所欠的都免了； \end{tabularx} \\ \\ \relax
18:33 & \begin{tabularx}{0.7\textwidth}{X} 你不應該憐憫你的同伴，像我憐憫你嗎？』 \end{tabularx} \\ \\ \relax
18:34 & \begin{tabularx}{0.7\textwidth}{X} 主人就大怒，把他交給司刑的，直到他還清了所欠的債。 \end{tabularx} \\ \\ \relax
18:35 & \begin{tabularx}{0.7\textwidth}{X} 你們各人若不從心裡饒恕你的弟兄，我天父也要這樣待你們。」 \end{tabularx} \\ \\
[1ex]
\hline
\hline
\end{longtable}


\newpage

\section{第77屆~港九培靈會~鮑維均~第八講}
\label{sec:317}
\textbf{二人成為一體}
\newline
\newline
連結: \href{https://www.hkbibleconference.org/session-message/view/317}{\texttt{https://www.hkbibleconference.org/session-message/view/317}}
\newline
\newline
\hyperref[sec:316]{< < < PREV SERMON < < <}
~
\hyperref[sec:toc]{[返目錄]}
~
\hyperref[sec:318]{> > > NEXT SERMON > > >}
\newline
\newline
馬太福音 19:1-19:12
\newline
\begin{longtable}{cl}
\hline
\hline
章節 & 經文 (和合本修訂版)\\
\hline
19:1 & \begin{tabularx}{0.7\textwidth}{X} 耶穌說完了這些話，就離開加利利，來到猶太的境內、約旦河的東邊。 \end{tabularx} \\ \\ \relax
19:2 & \begin{tabularx}{0.7\textwidth}{X} 有一大群人跟著他，他就在那裡治好了他們。 \end{tabularx} \\ \\ \relax
19:3 & \begin{tabularx}{0.7\textwidth}{X} 有些法利賽人來試探耶穌說：「無論甚麼緣故，人休妻都合法嗎？」 \end{tabularx} \\ \\ \relax
19:4 & \begin{tabularx}{0.7\textwidth}{X} 耶穌回答：「那起初造人的，是造男造女， \end{tabularx} \\ \\ \relax
19:5 & \begin{tabularx}{0.7\textwidth}{X} 並且說：『因此，人要離開父母，與妻子結合，二人成為一體。』這經文你們沒有念過嗎？ \end{tabularx} \\ \\ \relax
19:6 & \begin{tabularx}{0.7\textwidth}{X} 既然如此，夫妻不再是兩個人，而是一體的了。所以，神配合的，人不可分開。」 \end{tabularx} \\ \\ \relax
19:7 & \begin{tabularx}{0.7\textwidth}{X} 法利賽人說：「這樣，摩西為甚麼吩咐給妻子休書就可以休她呢？」 \end{tabularx} \\ \\ \relax
19:8 & \begin{tabularx}{0.7\textwidth}{X} 耶穌說：「摩西因為你們的心硬，所以准許你們休妻，但起初並不是這樣。 \end{tabularx} \\ \\ \relax
19:9 & \begin{tabularx}{0.7\textwidth}{X} 我告訴你們，凡休妻另娶的，若不是為不貞的緣故，就是犯姦淫了。」 \end{tabularx} \\ \\ \relax
19:10 & \begin{tabularx}{0.7\textwidth}{X} 門徒對耶穌說：「丈夫和妻子的關係既是這樣，倒不如不娶。」 \end{tabularx} \\ \\ \relax
19:11 & \begin{tabularx}{0.7\textwidth}{X} 耶穌對他們說：「這話不是人人都能領受的，惟獨賜給誰，誰才能領受。 \end{tabularx} \\ \\ \relax
19:12 & \begin{tabularx}{0.7\textwidth}{X} 因為有人從母腹裡就是不宜結婚的，也有因人為的緣故不宜結婚的，並有為天國的緣故自己不結婚的。這話誰能領受，就領受吧。」 \end{tabularx} \\ \\
[1ex]
\hline
\hline
\end{longtable}


\newpage

\section{第77屆~港九培靈會~鮑維均~第九講}
\label{sec:318}
\textbf{無辜之人的血}
\newline
\newline
連結: \href{https://www.hkbibleconference.org/session-message/view/318}{\texttt{https://www.hkbibleconference.org/session-message/view/318}}
\newline
\newline
\hyperref[sec:317]{< < < PREV SERMON < < <}
~
\hyperref[sec:toc]{[返目錄]}
~
\hyperref[sec:707]{> > > NEXT SERMON > > >}
\newline
\newline
馬太福音 27:1-27:10
\newline
\begin{longtable}{cl}
\hline
\hline
章節 & 經文 (和合本修訂版)\\
\hline
27:1 & \begin{tabularx}{0.7\textwidth}{X} 到了早晨，眾祭司長和百姓的長老商議要處死耶穌， \end{tabularx} \\ \\ \relax
27:2 & \begin{tabularx}{0.7\textwidth}{X} 就把他綁著，解去，交給彼拉多總督。 \end{tabularx} \\ \\ \relax
27:3 & \begin{tabularx}{0.7\textwidth}{X} 這時，出賣耶穌的猶大看見耶穌已經定了罪，就後悔，把那三十塊銀錢拿回來給祭司長和長老， \end{tabularx} \\ \\ \relax
27:4 & \begin{tabularx}{0.7\textwidth}{X} 說：「我出賣了無辜人的血有罪了。」他們說：「那跟我們有甚麼相干？你自己承當吧！」 \end{tabularx} \\ \\ \relax
27:5 & \begin{tabularx}{0.7\textwidth}{X} 猶大就把那銀錢丟在殿裡，出去吊死了。 \end{tabularx} \\ \\ \relax
27:6 & \begin{tabularx}{0.7\textwidth}{X} 祭司長拾起銀錢來，說：「這是血價，不可放在聖殿的銀庫裡。」 \end{tabularx} \\ \\ \relax
27:7 & \begin{tabularx}{0.7\textwidth}{X} 他們商議，就用那銀錢買了窯戶的一塊田，用來埋葬外鄉人。 \end{tabularx} \\ \\ \relax
27:8 & \begin{tabularx}{0.7\textwidth}{X} 所以，那塊田直到今日還叫做「血田」。 \end{tabularx} \\ \\ \relax
27:9 & \begin{tabularx}{0.7\textwidth}{X} 這就應驗了先知耶利米所說的話：「他們用那三十塊銀錢，就是以色列人給那被估定的人所估定的價錢， \end{tabularx} \\ \\ \relax
27:10 & \begin{tabularx}{0.7\textwidth}{X} 買了窯戶的一塊田；這是照著主所吩咐我的。」 \end{tabularx} \\ \\
[1ex]
\hline
\hline
\end{longtable}


\newpage

\section{第78屆~港九培靈會~唐崇明~第一講}
\label{sec:707}
\textbf{甘為愛奴獻一生}
\newline
\newline
連結: \href{https://www.hkbibleconference.org/session-message/view/707}{\texttt{https://www.hkbibleconference.org/session-message/view/707}}
\newline
\newline
\hyperref[sec:318]{< < < PREV SERMON < < <}
~
\hyperref[sec:toc]{[返目錄]}
~
\hyperref[sec:708]{> > > NEXT SERMON > > >}
\newline
\newline
以弗所書 4:1-4:3
\newline
\begin{longtable}{cl}
\hline
\hline
章節 & 經文 (和合本修訂版)\\
\hline
4:1 & \begin{tabularx}{0.7\textwidth}{X} 我為主作囚徒的勸你們，既然蒙召，行事為人就要與你們所蒙的呼召相稱。 \end{tabularx} \\ \\ \relax
4:2 & \begin{tabularx}{0.7\textwidth}{X} 凡事要謙虛、溫柔、忍耐，用愛心互相寬容， \end{tabularx} \\ \\ \relax
4:3 & \begin{tabularx}{0.7\textwidth}{X} 以和平彼此聯繫，竭力保持聖靈所賜的合一。 \end{tabularx} \\ \\
[1ex]
\hline
\hline
\end{longtable}


\newpage

\section{第78屆~港九培靈會~唐崇明~第二講}
\label{sec:708}
\textbf{畢生所冀只為神}
\newline
\newline
連結: \href{https://www.hkbibleconference.org/session-message/view/708}{\texttt{https://www.hkbibleconference.org/session-message/view/708}}
\newline
\newline
\hyperref[sec:707]{< < < PREV SERMON < < <}
~
\hyperref[sec:toc]{[返目錄]}
~
\hyperref[sec:709]{> > > NEXT SERMON > > >}
\newline
\newline
腓立比書 3:9-3:12
\newline
\begin{longtable}{cl}
\hline
\hline
章節 & 經文 (和合本修訂版)\\
\hline
3:9 & \begin{tabularx}{0.7\textwidth}{X} 並且得以在他裡面，不是有自己因律法而得的義，而是有信基督的義，就是基於信，從神而來的義， \end{tabularx} \\ \\ \relax
3:10 & \begin{tabularx}{0.7\textwidth}{X} 使我認識基督，知道他復活的大能，並且知道和他一同受苦，效法他的死， \end{tabularx} \\ \\ \relax
3:11 & \begin{tabularx}{0.7\textwidth}{X} 或許我也得以從死人中復活。 \end{tabularx} \\ \\ \relax
3:12 & \begin{tabularx}{0.7\textwidth}{X} 這不是說我已經得著了，已經完全了；而是竭力追求，或許可以得著基督耶穌所要我得著的。 \end{tabularx} \\ \\
[1ex]
\hline
\hline
\end{longtable}


\newpage

\section{第78屆~港九培靈會~唐崇明~第三講}
\label{sec:709}
\textbf{父事為念定終生}
\newline
\newline
連結: \href{https://www.hkbibleconference.org/session-message/view/709}{\texttt{https://www.hkbibleconference.org/session-message/view/709}}
\newline
\newline
\hyperref[sec:708]{< < < PREV SERMON < < <}
~
\hyperref[sec:toc]{[返目錄]}
~
\hyperref[sec:710]{> > > NEXT SERMON > > >}
\newline
\newline
路加福音 2:41-2:52
\newline
\begin{longtable}{cl}
\hline
\hline
章節 & 經文 (和合本修訂版)\\
\hline
2:41 & \begin{tabularx}{0.7\textwidth}{X} 每年逾越節，他父母都上耶路撒冷去。 \end{tabularx} \\ \\ \relax
2:42 & \begin{tabularx}{0.7\textwidth}{X} 當他十二歲的時候，他們按著過節的規矩上去。 \end{tabularx} \\ \\ \relax
2:43 & \begin{tabularx}{0.7\textwidth}{X} 守滿了節期，他們回去，孩童耶穌仍舊在耶路撒冷。他的父母並不知道， \end{tabularx} \\ \\ \relax
2:44 & \begin{tabularx}{0.7\textwidth}{X} 以為他在同行的人中間，走了一天的路程才在親屬和熟悉的人中找他， \end{tabularx} \\ \\ \relax
2:45 & \begin{tabularx}{0.7\textwidth}{X} 既找不著，就回耶路撒冷去找他。 \end{tabularx} \\ \\ \relax
2:46 & \begin{tabularx}{0.7\textwidth}{X} 過了三天，他們發現他在聖殿裡，坐在教師中間，一面聽，一面問。 \end{tabularx} \\ \\ \relax
2:47 & \begin{tabularx}{0.7\textwidth}{X} 凡聽見他的人都對他的聰明和應對感到驚奇。 \end{tabularx} \\ \\ \relax
2:48 & \begin{tabularx}{0.7\textwidth}{X} 他父母看見就很驚奇。他母親對他說：「我兒啊，為甚麼對我們這樣做呢？看哪，你父親和我很焦急，到處找你！」 \end{tabularx} \\ \\ \relax
2:49 & \begin{tabularx}{0.7\textwidth}{X} 耶穌對他們說：「為甚麼找我呢？難道你們不知道我應當在我父的家裡嗎？」 \end{tabularx} \\ \\ \relax
2:50 & \begin{tabularx}{0.7\textwidth}{X} 他所說的這話，他們不明白。 \end{tabularx} \\ \\ \relax
2:51 & \begin{tabularx}{0.7\textwidth}{X} 他就同他們下去，回到拿撒勒，並且順從他們。他母親把這一切的事都存在心裡。 \end{tabularx} \\ \\ \relax
2:52 & \begin{tabularx}{0.7\textwidth}{X} 耶穌的智慧和身量，並神和人喜愛他的心，都一齊增長。 \end{tabularx} \\ \\
[1ex]
\hline
\hline
\end{longtable}


\newpage

\section{第78屆~港九培靈會~唐崇明~第四講}
\label{sec:710}
\textbf{愛得真誠愛得純}
\newline
\newline
連結: \href{https://www.hkbibleconference.org/session-message/view/710}{\texttt{https://www.hkbibleconference.org/session-message/view/710}}
\newline
\newline
\hyperref[sec:709]{< < < PREV SERMON < < <}
~
\hyperref[sec:toc]{[返目錄]}
~
\hyperref[sec:711]{> > > NEXT SERMON > > >}
\newline
\newline
哥林多前書 13:1-13:13
\newline
\begin{longtable}{cl}
\hline
\hline
章節 & 經文 (和合本修訂版)\\
\hline
13:1 & \begin{tabularx}{0.7\textwidth}{X} 我若能說人間的方言，甚至天使的語言，卻沒有愛，我就成為鳴的鑼、響的鈸一般。 \end{tabularx} \\ \\ \relax
13:2 & \begin{tabularx}{0.7\textwidth}{X} 我若有先知講道的能力，也明白各樣的奧祕，各樣的知識，而且有齊備的信心，使我能夠移山，卻沒有愛，我就算不了甚麼。 \end{tabularx} \\ \\ \relax
13:3 & \begin{tabularx}{0.7\textwidth}{X} 我若將所有的財產救濟窮人，又犧牲自己的身體讓人誇讚，卻沒有愛，仍然對我無益。 \end{tabularx} \\ \\ \relax
13:4 & \begin{tabularx}{0.7\textwidth}{X} 愛是恆久忍耐；又有恩慈；愛是不嫉妒；愛是不自誇，不張狂， \end{tabularx} \\ \\ \relax
13:5 & \begin{tabularx}{0.7\textwidth}{X} 不做害羞的事，不求自己的益處，不輕易發怒，不計算人的惡， \end{tabularx} \\ \\ \relax
13:6 & \begin{tabularx}{0.7\textwidth}{X} 不喜歡不義，只喜歡真理； \end{tabularx} \\ \\ \relax
13:7 & \begin{tabularx}{0.7\textwidth}{X} 凡事包容，凡事相信，凡事盼望，凡事忍耐。 \end{tabularx} \\ \\ \relax
13:8 & \begin{tabularx}{0.7\textwidth}{X} 愛是永不止息。先知講道之能終必歸於無有；說方言之能終必停止；知識也終必歸於無有。 \end{tabularx} \\ \\ \relax
13:9 & \begin{tabularx}{0.7\textwidth}{X} 我們現在所知道的有限，先知所講的也有限， \end{tabularx} \\ \\ \relax
13:10 & \begin{tabularx}{0.7\textwidth}{X} 等那完全的來到，這有限的必消逝。 \end{tabularx} \\ \\ \relax
13:11 & \begin{tabularx}{0.7\textwidth}{X} 我作孩子的時候，說話像孩子，心思像孩子，意念像孩子；既長大成人，就把孩子的事丟棄了。 \end{tabularx} \\ \\ \relax
13:12 & \begin{tabularx}{0.7\textwidth}{X} 我們現在是對著鏡子觀看，模糊不清；到那時，就要面對面了。我如今所認識的有限，到那時就全認識，如同主認識我一樣。 \end{tabularx} \\ \\ \relax
13:13 & \begin{tabularx}{0.7\textwidth}{X} 如今常存的有信，有望，有愛這三樣，其中最大的是愛。 \end{tabularx} \\ \\
[1ex]
\hline
\hline
\end{longtable}


\newpage

\section{第78屆~港九培靈會~唐崇明~第五講}
\label{sec:711}
\textbf{心灰意冷點心燈}
\newline
\newline
連結: \href{https://www.hkbibleconference.org/session-message/view/711}{\texttt{https://www.hkbibleconference.org/session-message/view/711}}
\newline
\newline
\hyperref[sec:710]{< < < PREV SERMON < < <}
~
\hyperref[sec:toc]{[返目錄]}
~
\hyperref[sec:712]{> > > NEXT SERMON > > >}
\newline
\newline
耶利米書 15:19-15:21
\newline
\begin{longtable}{cl}
\hline
\hline
章節 & 經文 (和合本修訂版)\\
\hline
15:19 & \begin{tabularx}{0.7\textwidth}{X} 所以耶和華如此說：「你若回轉，我就使你歸回，站在我面前。你若能將寶物和無用之物分別出來，你就可以當作我的口。他們必歸向你，你卻不可歸向他們。 \end{tabularx} \\ \\ \relax
15:20 & \begin{tabularx}{0.7\textwidth}{X} 我必使你向這百姓成為堅固的銅牆。他們必攻擊你，卻不能勝過你；因我與你同在，要拯救你，搭救你。這是耶和華說的。 \end{tabularx} \\ \\ \relax
15:21 & \begin{tabularx}{0.7\textwidth}{X} 我必搭救你脫離惡人的手，救贖你脫離殘暴之人的手。」 \end{tabularx} \\ \\
[1ex]
\hline
\hline
\end{longtable}


\newpage

\section{第78屆~港九培靈會~唐崇明~第六講}
\label{sec:712}
\textbf{信實廣大全是恩}
\newline
\newline
連結: \href{https://www.hkbibleconference.org/session-message/view/712}{\texttt{https://www.hkbibleconference.org/session-message/view/712}}
\newline
\newline
\hyperref[sec:711]{< < < PREV SERMON < < <}
~
\hyperref[sec:toc]{[返目錄]}
~
\hyperref[sec:713]{> > > NEXT SERMON > > >}
\newline
\newline
耶利米哀歌 3:19-3:27
\newline
\begin{longtable}{cl}
\hline
\hline
章節 & 經文 (和合本修訂版)\\
\hline
3:19 & \begin{tabularx}{0.7\textwidth}{X} 求你記得我的困苦和流離，它如茵蔯和苦膽一般； \end{tabularx} \\ \\ \relax
3:20 & \begin{tabularx}{0.7\textwidth}{X} 我心想念這些，就在我裡面憂悶。 \end{tabularx} \\ \\ \relax
3:21 & \begin{tabularx}{0.7\textwidth}{X} 但我的心回轉過來，因此就有指望； \end{tabularx} \\ \\ \relax
3:22 & \begin{tabularx}{0.7\textwidth}{X} 因耶和華的慈愛，我們不致滅絕，因他的憐憫永不斷絕， \end{tabularx} \\ \\ \relax
3:23 & \begin{tabularx}{0.7\textwidth}{X} 每早晨，這都是新的；你的信實極其廣大！ \end{tabularx} \\ \\ \relax
3:24 & \begin{tabularx}{0.7\textwidth}{X} 我心裡說：「耶和華是我的福分，因此，我要仰望他。」 \end{tabularx} \\ \\ \relax
3:25 & \begin{tabularx}{0.7\textwidth}{X} 凡等候耶和華，心裡尋求他的，耶和華必施恩給他。 \end{tabularx} \\ \\ \relax
3:26 & \begin{tabularx}{0.7\textwidth}{X} 人仰望耶和華，安靜等候他的救恩，這是好的。 \end{tabularx} \\ \\ \relax
3:27 & \begin{tabularx}{0.7\textwidth}{X} 人在年輕時負軛，這是好的。 \end{tabularx} \\ \\
[1ex]
\hline
\hline
\end{longtable}


\newpage

\section{第78屆~港九培靈會~唐崇明~第七講}
\label{sec:713}
\textbf{展翅上騰破雲層}
\newline
\newline
連結: \href{https://www.hkbibleconference.org/session-message/view/713}{\texttt{https://www.hkbibleconference.org/session-message/view/713}}
\newline
\newline
\hyperref[sec:712]{< < < PREV SERMON < < <}
~
\hyperref[sec:toc]{[返目錄]}
~
\hyperref[sec:719]{> > > NEXT SERMON > > >}
\newline
\newline
以賽亞書 40:27-40:31
\newline
\begin{longtable}{cl}
\hline
\hline
章節 & 經文 (和合本修訂版)\\
\hline
40:27 & \begin{tabularx}{0.7\textwidth}{X} 雅各啊，你為何說，以色列啊，你為何言，「我的道路向耶和華隱藏，我的冤屈神並不查問」？ \end{tabularx} \\ \\ \relax
40:28 & \begin{tabularx}{0.7\textwidth}{X} 你豈不曾知道嗎？你豈未曾聽見嗎？永在的神耶和華，創造地極的主，他不疲乏，也不困倦；他的智慧無法測度。 \end{tabularx} \\ \\ \relax
40:29 & \begin{tabularx}{0.7\textwidth}{X} 疲乏的，他賜能力；軟弱的，他加力量。 \end{tabularx} \\ \\ \relax
40:30 & \begin{tabularx}{0.7\textwidth}{X} 就是年輕人也要疲乏困倦，強壯的也必全然跌倒。 \end{tabularx} \\ \\ \relax
40:31 & \begin{tabularx}{0.7\textwidth}{X} 但那等候耶和華的必重新得力。他們必如鷹展翅上騰；他們奔跑卻不困倦，行走卻不疲乏。 \end{tabularx} \\ \\
[1ex]
\hline
\hline
\end{longtable}


\newpage

\section{第78屆~港九培靈會~唐崇明~第八講}
\label{sec:719}
\textbf{旗開得勝奔前程}
\newline
\newline
連結: \href{https://www.hkbibleconference.org/session-message/view/719}{\texttt{https://www.hkbibleconference.org/session-message/view/719}}
\newline
\newline
\hyperref[sec:713]{< < < PREV SERMON < < <}
~
\hyperref[sec:toc]{[返目錄]}
~
\hyperref[sec:722]{> > > NEXT SERMON > > >}
\newline
\newline
出埃及記 17:8-17:16
\newline
\begin{longtable}{cl}
\hline
\hline
章節 & 經文 (和合本修訂版)\\
\hline
17:8 & \begin{tabularx}{0.7\textwidth}{X} 那時，亞瑪力來到利非訂，和以色列爭戰。 \end{tabularx} \\ \\ \relax
17:9 & \begin{tabularx}{0.7\textwidth}{X} 摩西對約書亞說：「你為我們選出人來，出去和亞瑪力爭戰。明天我要站在山頂上，手裡拿著神的杖。」 \end{tabularx} \\ \\ \relax
17:10 & \begin{tabularx}{0.7\textwidth}{X} 於是，約書亞照著摩西對他所說的話去做，和亞瑪力爭戰。摩西、亞倫和戶珥都上了山頂。 \end{tabularx} \\ \\ \relax
17:11 & \begin{tabularx}{0.7\textwidth}{X} 摩西何時舉手，以色列就得勝；何時垂手，亞瑪力就得勝。 \end{tabularx} \\ \\ \relax
17:12 & \begin{tabularx}{0.7\textwidth}{X} 但摩西的雙手沉重，他們就搬一塊石頭來放在他下面，他就坐在上面。亞倫與戶珥扶著他的手，一個在這邊，一個在那邊，他的手就穩住，直到日落。 \end{tabularx} \\ \\ \relax
17:13 & \begin{tabularx}{0.7\textwidth}{X} 約書亞用刀打敗了亞瑪力和他的百姓。 \end{tabularx} \\ \\ \relax
17:14 & \begin{tabularx}{0.7\textwidth}{X} 耶和華對摩西說：「你要把這事記錄在書上作紀念，又念給約書亞聽：我要把亞瑪力的名字從天下全然塗去。」 \end{tabularx} \\ \\ \relax
17:15 & \begin{tabularx}{0.7\textwidth}{X} 摩西築了一座壇，起名叫「耶和華尼西」。 \end{tabularx} \\ \\ \relax
17:16 & \begin{tabularx}{0.7\textwidth}{X} 他說：「我指著耶和華的寶座發誓，耶和華必世世代代和亞瑪力爭戰。」 \end{tabularx} \\ \\
[1ex]
\hline
\hline
\end{longtable}


\newpage

\section{第78屆~港九培靈會~唐崇明~第九講}
\label{sec:722}
\textbf{無愧工人逾千軍}
\newline
\newline
連結: \href{https://www.hkbibleconference.org/session-message/view/722}{\texttt{https://www.hkbibleconference.org/session-message/view/722}}
\newline
\newline
\hyperref[sec:719]{< < < PREV SERMON < < <}
~
\hyperref[sec:toc]{[返目錄]}
~
\hyperref[sec:725]{> > > NEXT SERMON > > >}
\newline
\newline
提摩太後書 2:15-2:16
\newline
\begin{longtable}{cl}
\hline
\hline
章節 & 經文 (和合本修訂版)\\
\hline
2:15 & \begin{tabularx}{0.7\textwidth}{X} 你當竭力在神面前作一個經得起考驗、無愧的工人，按著正意講解真理的話。 \end{tabularx} \\ \\ \relax
2:16 & \begin{tabularx}{0.7\textwidth}{X} 要遠避世俗的空談，因為這等空談會使人進到更不敬虔的地步。 \end{tabularx} \\ \\
[1ex]
\hline
\hline
\end{longtable}


\newpage

\section{第78屆~港九培靈會~唐崇明~第十講}
\label{sec:725}
\textbf{重新立志顧羊群}
\newline
\newline
連結: \href{https://www.hkbibleconference.org/session-message/view/725}{\texttt{https://www.hkbibleconference.org/session-message/view/725}}
\newline
\newline
\hyperref[sec:722]{< < < PREV SERMON < < <}
~
\hyperref[sec:toc]{[返目錄]}
~
\hyperref[sec:664]{> > > NEXT SERMON > > >}
\newline
\newline
約翰福音 21:15-21:19
\newline
\begin{longtable}{cl}
\hline
\hline
章節 & 經文 (和合本修訂版)\\
\hline
21:15 & \begin{tabularx}{0.7\textwidth}{X} 他們吃完了早飯，耶穌對西門‧彼得說：「約翰的兒子西門，你愛我比這些更深嗎？」彼得對他說：「主啊，是的，你知道我愛你。」耶穌說：「你餵養我的小羊。」 \end{tabularx} \\ \\ \relax
21:16 & \begin{tabularx}{0.7\textwidth}{X} 耶穌第二次又對他說：「約翰的兒子西門，你愛我嗎？」彼得對他說：「主啊，是的，你知道我愛你。」耶穌說：「你牧養我的羊。」 \end{tabularx} \\ \\ \relax
21:17 & \begin{tabularx}{0.7\textwidth}{X} 耶穌第三次對他說：「約翰的兒子西門，你愛我嗎？」彼得因為耶穌第三次對他說「你愛我嗎」，就憂愁，對耶穌說：「主啊，你無所不知，你知道我愛你。」耶穌說：「你餵養我的羊。 \end{tabularx} \\ \\ \relax
21:18 & \begin{tabularx}{0.7\textwidth}{X} 我實實在在地告訴你，你年輕的時候，自己束上帶子，隨意往來；但年老的時候，你要伸出手來，別人要把你束上，帶你到不願意去的地方。」 \end{tabularx} \\ \\ \relax
21:19 & \begin{tabularx}{0.7\textwidth}{X} 耶穌說這話，是指彼得會怎樣死來榮耀神。說了這話，耶穌對他說：「你跟從我吧！」 \end{tabularx} \\ \\
[1ex]
\hline
\hline
\end{longtable}


\newpage

\section{第78屆~港九培靈會~楊錫鏘~第一講}
\label{sec:664}
\textbf{創造的主權}
\newline
\newline
連結: \href{https://www.hkbibleconference.org/session-message/view/664}{\texttt{https://www.hkbibleconference.org/session-message/view/664}}
\newline
\newline
\hyperref[sec:725]{< < < PREV SERMON < < <}
~
\hyperref[sec:toc]{[返目錄]}
~
\hyperref[sec:665]{> > > NEXT SERMON > > >}
\newline
\newline
創世記 1:1-1:31
\newline
\begin{longtable}{cl}
\hline
\hline
章節 & 經文 (和合本修訂版)\\
\hline
1:1 & \begin{tabularx}{0.7\textwidth}{X} 起初，神創造天地。 \end{tabularx} \\ \\ \relax
1:2 & \begin{tabularx}{0.7\textwidth}{X} 地是空虛混沌，深淵上面一片黑暗；神的靈運行在水面上。 \end{tabularx} \\ \\ \relax
1:3 & \begin{tabularx}{0.7\textwidth}{X} 神說：「要有光」，就有了光。 \end{tabularx} \\ \\ \relax
1:4 & \begin{tabularx}{0.7\textwidth}{X} 神看光是好的，於是神就把光和暗分開。 \end{tabularx} \\ \\ \relax
1:5 & \begin{tabularx}{0.7\textwidth}{X} 神稱光為「晝」，稱暗為「夜」。有晚上，有早晨，這是第一日。 \end{tabularx} \\ \\ \relax
1:6 & \begin{tabularx}{0.7\textwidth}{X} 神說：「眾水之間要有穹蒼，把水和水分開。」 \end{tabularx} \\ \\ \relax
1:7 & \begin{tabularx}{0.7\textwidth}{X} 神就造了穹蒼，把穹蒼以下的水和穹蒼以上的水分開。事就這樣成了。 \end{tabularx} \\ \\ \relax
1:8 & \begin{tabularx}{0.7\textwidth}{X} 神稱穹蒼為「天」。有晚上，有早晨，這是第二日。 \end{tabularx} \\ \\ \relax
1:9 & \begin{tabularx}{0.7\textwidth}{X} 神說：「天下面的水要聚集在一處，使乾地露出來。」事就這樣成了。 \end{tabularx} \\ \\ \relax
1:10 & \begin{tabularx}{0.7\textwidth}{X} 神稱乾地為「地」，稱聚集在一起的水為「海」。神看為好的。 \end{tabularx} \\ \\ \relax
1:11 & \begin{tabularx}{0.7\textwidth}{X} 神說：「地要長出植物，就是含種子的五穀菜蔬，和會結果子、果子裡有種子的樹，在地上各從其類。」事就這樣成了。 \end{tabularx} \\ \\ \relax
1:12 & \begin{tabularx}{0.7\textwidth}{X} 於是地長出了植物：含種子的五穀菜蔬，各從其類；會結果子、果子裡有種子的樹，各從其類。神看為好的。 \end{tabularx} \\ \\ \relax
1:13 & \begin{tabularx}{0.7\textwidth}{X} 有晚上，有早晨，這是第三日。 \end{tabularx} \\ \\ \relax
1:14 & \begin{tabularx}{0.7\textwidth}{X} 神說：「天上要有光體來分晝夜，讓它們作記號，定季節、日子、年份， \end{tabularx} \\ \\ \relax
1:15 & \begin{tabularx}{0.7\textwidth}{X} 它們要在天空發光，照在地上。」事就這樣成了。 \end{tabularx} \\ \\ \relax
1:16 & \begin{tabularx}{0.7\textwidth}{X} 於是神造了兩個大光體，大的管晝，小的管夜，又造了星辰。 \end{tabularx} \\ \\ \relax
1:17 & \begin{tabularx}{0.7\textwidth}{X} 神把這些光體擺列在天空，照在地上， \end{tabularx} \\ \\ \relax
1:18 & \begin{tabularx}{0.7\textwidth}{X} 管理晝夜，分別光暗。神看為好的。 \end{tabularx} \\ \\ \relax
1:19 & \begin{tabularx}{0.7\textwidth}{X} 有晚上，有早晨，這是第四日。 \end{tabularx} \\ \\ \relax
1:20 & \begin{tabularx}{0.7\textwidth}{X} 神說：「水要滋生眾多有生命之物；要有鳥飛在地面以上，天空之中。」 \end{tabularx} \\ \\ \relax
1:21 & \begin{tabularx}{0.7\textwidth}{X} 神就創造了大魚和在水裡滋生的各樣活動的生物，各從其類，以及各樣有翅膀的鳥，各從其類。神看為好的。 \end{tabularx} \\ \\ \relax
1:22 & \begin{tabularx}{0.7\textwidth}{X} 神就賜福給這一切，說：「要繁殖增多，充滿在海的水裡；飛鳥也要在地上增多。」 \end{tabularx} \\ \\ \relax
1:23 & \begin{tabularx}{0.7\textwidth}{X} 有晚上，有早晨，這是第五日。 \end{tabularx} \\ \\ \relax
1:24 & \begin{tabularx}{0.7\textwidth}{X} 神說：「地要生出有生命之物，各從其類，就是牲畜、爬行動物、地上的走獸，各從其類。」事就這樣成了。 \end{tabularx} \\ \\ \relax
1:25 & \begin{tabularx}{0.7\textwidth}{X} 於是神造了地上的走獸，各從其類；牲畜，各從其類；和地上一切的爬行動物，各從其類。神看為好的。 \end{tabularx} \\ \\ \relax
1:26 & \begin{tabularx}{0.7\textwidth}{X} 神說：「我們要照著我們的形像，按著我們的樣式造人，使他們管理海裡的魚、天空的鳥、地上的牲畜和全地，以及地上爬的一切爬行動物。」 \end{tabularx} \\ \\ \relax
1:27 & \begin{tabularx}{0.7\textwidth}{X} 神就照著他的形像創造人，照著神的形像創造他們；他創造了他們，有男有女。 \end{tabularx} \\ \\ \relax
1:28 & \begin{tabularx}{0.7\textwidth}{X} 神賜福給他們，神對他們說：「要生養眾多，遍滿這地，治理它；要管理海裡的魚、天空的鳥和地上各樣活動的生物。」 \end{tabularx} \\ \\ \relax
1:29 & \begin{tabularx}{0.7\textwidth}{X} 神說：「看哪，我把全地一切含種子的五穀菜蔬和一切會結果子、果子裡有種子的樹，都賜給你們；這些都可作食物。 \end{tabularx} \\ \\ \relax
1:30 & \begin{tabularx}{0.7\textwidth}{X} 至於地上一切的走獸、天空一切的飛鳥，並一切在地上爬行的，有生命的動物，我把綠色植物賜給牠們作食物。」事就這樣成了。 \end{tabularx} \\ \\ \relax
1:31 & \begin{tabularx}{0.7\textwidth}{X} 神看一切所造的，看哪，都非常好。有晚上，有早晨，這是第六日。 \end{tabularx} \\ \\
[1ex]
\hline
\hline
\end{longtable}


\newpage

\section{第78屆~港九培靈會~楊錫鏘~第二講}
\label{sec:665}
\textbf{創造的設計(一)}
\newline
\newline
連結: \href{https://www.hkbibleconference.org/session-message/view/665}{\texttt{https://www.hkbibleconference.org/session-message/view/665}}
\newline
\newline
\hyperref[sec:664]{< < < PREV SERMON < < <}
~
\hyperref[sec:toc]{[返目錄]}
~
\hyperref[sec:666]{> > > NEXT SERMON > > >}
\newline
\newline
創世記 1:1-1:10
\newline
\begin{longtable}{cl}
\hline
\hline
章節 & 經文 (和合本修訂版)\\
\hline
1:1 & \begin{tabularx}{0.7\textwidth}{X} 起初，神創造天地。 \end{tabularx} \\ \\ \relax
1:2 & \begin{tabularx}{0.7\textwidth}{X} 地是空虛混沌，深淵上面一片黑暗；神的靈運行在水面上。 \end{tabularx} \\ \\ \relax
1:3 & \begin{tabularx}{0.7\textwidth}{X} 神說：「要有光」，就有了光。 \end{tabularx} \\ \\ \relax
1:4 & \begin{tabularx}{0.7\textwidth}{X} 神看光是好的，於是神就把光和暗分開。 \end{tabularx} \\ \\ \relax
1:5 & \begin{tabularx}{0.7\textwidth}{X} 神稱光為「晝」，稱暗為「夜」。有晚上，有早晨，這是第一日。 \end{tabularx} \\ \\ \relax
1:6 & \begin{tabularx}{0.7\textwidth}{X} 神說：「眾水之間要有穹蒼，把水和水分開。」 \end{tabularx} \\ \\ \relax
1:7 & \begin{tabularx}{0.7\textwidth}{X} 神就造了穹蒼，把穹蒼以下的水和穹蒼以上的水分開。事就這樣成了。 \end{tabularx} \\ \\ \relax
1:8 & \begin{tabularx}{0.7\textwidth}{X} 神稱穹蒼為「天」。有晚上，有早晨，這是第二日。 \end{tabularx} \\ \\ \relax
1:9 & \begin{tabularx}{0.7\textwidth}{X} 神說：「天下面的水要聚集在一處，使乾地露出來。」事就這樣成了。 \end{tabularx} \\ \\ \relax
1:10 & \begin{tabularx}{0.7\textwidth}{X} 神稱乾地為「地」，稱聚集在一起的水為「海」。神看為好的。 \end{tabularx} \\ \\
[1ex]
\hline
\hline
\end{longtable}


\newpage

\section{第78屆~港九培靈會~楊錫鏘~第三講}
\label{sec:666}
\textbf{創造的設計(二)}
\newline
\newline
連結: \href{https://www.hkbibleconference.org/session-message/view/666}{\texttt{https://www.hkbibleconference.org/session-message/view/666}}
\newline
\newline
\hyperref[sec:665]{< < < PREV SERMON < < <}
~
\hyperref[sec:toc]{[返目錄]}
~
\hyperref[sec:667]{> > > NEXT SERMON > > >}
\newline
\newline
創世記 1:11-1:31
\newline
\begin{longtable}{cl}
\hline
\hline
章節 & 經文 (和合本修訂版)\\
\hline
1:11 & \begin{tabularx}{0.7\textwidth}{X} 神說：「地要長出植物，就是含種子的五穀菜蔬，和會結果子、果子裡有種子的樹，在地上各從其類。」事就這樣成了。 \end{tabularx} \\ \\ \relax
1:12 & \begin{tabularx}{0.7\textwidth}{X} 於是地長出了植物：含種子的五穀菜蔬，各從其類；會結果子、果子裡有種子的樹，各從其類。神看為好的。 \end{tabularx} \\ \\ \relax
1:13 & \begin{tabularx}{0.7\textwidth}{X} 有晚上，有早晨，這是第三日。 \end{tabularx} \\ \\ \relax
1:14 & \begin{tabularx}{0.7\textwidth}{X} 神說：「天上要有光體來分晝夜，讓它們作記號，定季節、日子、年份， \end{tabularx} \\ \\ \relax
1:15 & \begin{tabularx}{0.7\textwidth}{X} 它們要在天空發光，照在地上。」事就這樣成了。 \end{tabularx} \\ \\ \relax
1:16 & \begin{tabularx}{0.7\textwidth}{X} 於是神造了兩個大光體，大的管晝，小的管夜，又造了星辰。 \end{tabularx} \\ \\ \relax
1:17 & \begin{tabularx}{0.7\textwidth}{X} 神把這些光體擺列在天空，照在地上， \end{tabularx} \\ \\ \relax
1:18 & \begin{tabularx}{0.7\textwidth}{X} 管理晝夜，分別光暗。神看為好的。 \end{tabularx} \\ \\ \relax
1:19 & \begin{tabularx}{0.7\textwidth}{X} 有晚上，有早晨，這是第四日。 \end{tabularx} \\ \\ \relax
1:20 & \begin{tabularx}{0.7\textwidth}{X} 神說：「水要滋生眾多有生命之物；要有鳥飛在地面以上，天空之中。」 \end{tabularx} \\ \\ \relax
1:21 & \begin{tabularx}{0.7\textwidth}{X} 神就創造了大魚和在水裡滋生的各樣活動的生物，各從其類，以及各樣有翅膀的鳥，各從其類。神看為好的。 \end{tabularx} \\ \\ \relax
1:22 & \begin{tabularx}{0.7\textwidth}{X} 神就賜福給這一切，說：「要繁殖增多，充滿在海的水裡；飛鳥也要在地上增多。」 \end{tabularx} \\ \\ \relax
1:23 & \begin{tabularx}{0.7\textwidth}{X} 有晚上，有早晨，這是第五日。 \end{tabularx} \\ \\ \relax
1:24 & \begin{tabularx}{0.7\textwidth}{X} 神說：「地要生出有生命之物，各從其類，就是牲畜、爬行動物、地上的走獸，各從其類。」事就這樣成了。 \end{tabularx} \\ \\ \relax
1:25 & \begin{tabularx}{0.7\textwidth}{X} 於是神造了地上的走獸，各從其類；牲畜，各從其類；和地上一切的爬行動物，各從其類。神看為好的。 \end{tabularx} \\ \\ \relax
1:26 & \begin{tabularx}{0.7\textwidth}{X} 神說：「我們要照著我們的形像，按著我們的樣式造人，使他們管理海裡的魚、天空的鳥、地上的牲畜和全地，以及地上爬的一切爬行動物。」 \end{tabularx} \\ \\ \relax
1:27 & \begin{tabularx}{0.7\textwidth}{X} 神就照著他的形像創造人，照著神的形像創造他們；他創造了他們，有男有女。 \end{tabularx} \\ \\ \relax
1:28 & \begin{tabularx}{0.7\textwidth}{X} 神賜福給他們，神對他們說：「要生養眾多，遍滿這地，治理它；要管理海裡的魚、天空的鳥和地上各樣活動的生物。」 \end{tabularx} \\ \\ \relax
1:29 & \begin{tabularx}{0.7\textwidth}{X} 神說：「看哪，我把全地一切含種子的五穀菜蔬和一切會結果子、果子裡有種子的樹，都賜給你們；這些都可作食物。 \end{tabularx} \\ \\ \relax
1:30 & \begin{tabularx}{0.7\textwidth}{X} 至於地上一切的走獸、天空一切的飛鳥，並一切在地上爬行的，有生命的動物，我把綠色植物賜給牠們作食物。」事就這樣成了。 \end{tabularx} \\ \\ \relax
1:31 & \begin{tabularx}{0.7\textwidth}{X} 神看一切所造的，看哪，都非常好。有晚上，有早晨，這是第六日。 \end{tabularx} \\ \\
[1ex]
\hline
\hline
\end{longtable}


\newpage

\section{第78屆~港九培靈會~楊錫鏘~第四講}
\label{sec:667}
\textbf{創造的意義}
\newline
\newline
連結: \href{https://www.hkbibleconference.org/session-message/view/667}{\texttt{https://www.hkbibleconference.org/session-message/view/667}}
\newline
\newline
\hyperref[sec:666]{< < < PREV SERMON < < <}
~
\hyperref[sec:toc]{[返目錄]}
~
\hyperref[sec:668]{> > > NEXT SERMON > > >}
\newline
\newline
創世記 2:1-2:3
\newline
\begin{longtable}{cl}
\hline
\hline
章節 & 經文 (和合本修訂版)\\
\hline
2:1 & \begin{tabularx}{0.7\textwidth}{X} 天和地，以及萬象都完成了。 \end{tabularx} \\ \\ \relax
2:2 & \begin{tabularx}{0.7\textwidth}{X} 到第七日，神已經完成了造物之工，就在第七日安息了，歇了他所做一切的工。 \end{tabularx} \\ \\ \relax
2:3 & \begin{tabularx}{0.7\textwidth}{X} 神賜福給第七日，將它分別為聖，因為在這日，神安息了，歇了他所做一切創造的工。 \end{tabularx} \\ \\
[1ex]
\hline
\hline
\end{longtable}


\newpage

\section{第78屆~港九培靈會~楊錫鏘~第五講}
\label{sec:668}
\textbf{創造的禮物(一)}
\newline
\newline
連結: \href{https://www.hkbibleconference.org/session-message/view/668}{\texttt{https://www.hkbibleconference.org/session-message/view/668}}
\newline
\newline
\hyperref[sec:667]{< < < PREV SERMON < < <}
~
\hyperref[sec:toc]{[返目錄]}
~
\hyperref[sec:669]{> > > NEXT SERMON > > >}
\newline
\newline
創世記 2:4-2:17
\newline
\begin{longtable}{cl}
\hline
\hline
章節 & 經文 (和合本修訂版)\\
\hline
2:4 & \begin{tabularx}{0.7\textwidth}{X} 這就是天地創造的來歷。在耶和華神造地和天的時候， \end{tabularx} \\ \\ \relax
2:5 & \begin{tabularx}{0.7\textwidth}{X} 地上還沒有田野的草木，田間的菜蔬還沒有長出來，因為耶和華神還沒有降雨在地上，也沒有人耕種土地。 \end{tabularx} \\ \\ \relax
2:6 & \begin{tabularx}{0.7\textwidth}{X} 但是，有霧氣從地上騰，滋潤整個土地的表面。 \end{tabularx} \\ \\ \relax
2:7 & \begin{tabularx}{0.7\textwidth}{X} 耶和華神用地上的塵土造人，將生命之氣吹進他的鼻孔，這人就成了有靈的活人。 \end{tabularx} \\ \\ \relax
2:8 & \begin{tabularx}{0.7\textwidth}{X} 耶和華神在東方的伊甸栽了一個園子，把所造的人安置在那裡。 \end{tabularx} \\ \\ \relax
2:9 & \begin{tabularx}{0.7\textwidth}{X} 耶和華神使各樣的樹從土地裡長出來，可以悅人的眼目，好作食物。園子當中有生命樹和知善惡的樹。 \end{tabularx} \\ \\ \relax
2:10 & \begin{tabularx}{0.7\textwidth}{X} 有一條河從伊甸流出來，滋潤那園子，從那裡分成四個源頭： \end{tabularx} \\ \\ \relax
2:11 & \begin{tabularx}{0.7\textwidth}{X} 第一條名叫比遜，它環繞哈腓拉全地，在那裡有金子。 \end{tabularx} \\ \\ \relax
2:12 & \begin{tabularx}{0.7\textwidth}{X} 那地的金子很好，在那裡也有珍珠和紅瑪瑙。 \end{tabularx} \\ \\ \relax
2:13 & \begin{tabularx}{0.7\textwidth}{X} 第二條河名叫基訓，它環繞古實全地。 \end{tabularx} \\ \\ \relax
2:14 & \begin{tabularx}{0.7\textwidth}{X} 第三條河名叫底格里斯，它流到亞述的東邊。第四條河就是幼發拉底。 \end{tabularx} \\ \\ \relax
2:15 & \begin{tabularx}{0.7\textwidth}{X} 耶和華神把那人安置在伊甸園，讓他耕耘看管。 \end{tabularx} \\ \\ \relax
2:16 & \begin{tabularx}{0.7\textwidth}{X} 耶和華神吩咐那人說：「園中各樣樹上所出的，你可以隨意吃， \end{tabularx} \\ \\ \relax
2:17 & \begin{tabularx}{0.7\textwidth}{X} 只是知善惡的樹所出的，你不可吃，因為你吃它的日子必定死！」 \end{tabularx} \\ \\
[1ex]
\hline
\hline
\end{longtable}


\newpage

\section{第78屆~港九培靈會~楊錫鏘~第六講}
\label{sec:669}
\textbf{創造的禮物(二)}
\newline
\newline
連結: \href{https://www.hkbibleconference.org/session-message/view/669}{\texttt{https://www.hkbibleconference.org/session-message/view/669}}
\newline
\newline
\hyperref[sec:668]{< < < PREV SERMON < < <}
~
\hyperref[sec:toc]{[返目錄]}
~
\hyperref[sec:670]{> > > NEXT SERMON > > >}
\newline
\newline
創世記 2:18-2:25
\newline
\begin{longtable}{cl}
\hline
\hline
章節 & 經文 (和合本修訂版)\\
\hline
2:18 & \begin{tabularx}{0.7\textwidth}{X} 耶和華神說：「那人單獨一個不好，我要為他造一個配偶幫助他。」 \end{tabularx} \\ \\ \relax
2:19 & \begin{tabularx}{0.7\textwidth}{X} 耶和華神用泥土造了野地各樣的走獸和天空各樣的飛鳥，都帶到那人面前，看他叫甚麼。那人怎樣叫各樣的動物，那就是牠的名字。 \end{tabularx} \\ \\ \relax
2:20 & \begin{tabularx}{0.7\textwidth}{X} 那人就給一切牲畜、天空的飛鳥和野地各樣的走獸都起了名。只是亞當沒有找到配偶幫助他。 \end{tabularx} \\ \\ \relax
2:21 & \begin{tabularx}{0.7\textwidth}{X} 耶和華神使他沉睡，他就睡了；於是取下他的一根肋骨，又在原處把肉合起來。 \end{tabularx} \\ \\ \relax
2:22 & \begin{tabularx}{0.7\textwidth}{X} 耶和華神就用那人身上所取的肋骨造了一個女人，帶她到那人面前。 \end{tabularx} \\ \\ \relax
2:23 & \begin{tabularx}{0.7\textwidth}{X} 那人說：「這是我骨中的骨，肉中的肉，可以稱她為女人，因為她是從男人身上取出來的。」 \end{tabularx} \\ \\ \relax
2:24 & \begin{tabularx}{0.7\textwidth}{X} 因此，人要離開父母，與妻子結合，二人成為一體。 \end{tabularx} \\ \\ \relax
2:25 & \begin{tabularx}{0.7\textwidth}{X} 當時夫妻二人赤身露體，並不覺得羞恥。 \end{tabularx} \\ \\
[1ex]
\hline
\hline
\end{longtable}


\newpage

\section{第78屆~港九培靈會~楊錫鏘~第七講}
\label{sec:670}
\textbf{受造物的悖逆}
\newline
\newline
連結: \href{https://www.hkbibleconference.org/session-message/view/670}{\texttt{https://www.hkbibleconference.org/session-message/view/670}}
\newline
\newline
\hyperref[sec:669]{< < < PREV SERMON < < <}
~
\hyperref[sec:toc]{[返目錄]}
~
\hyperref[sec:671]{> > > NEXT SERMON > > >}
\newline
\newline
創世記 3:1-3:7
\newline
\begin{longtable}{cl}
\hline
\hline
章節 & 經文 (和合本修訂版)\\
\hline
3:1 & \begin{tabularx}{0.7\textwidth}{X} 耶和華神所造的，惟有蛇比田野一切的走獸更狡猾。蛇對女人說：「神豈是真說，你們不可吃園中任何樹上所出的嗎？」 \end{tabularx} \\ \\ \relax
3:2 & \begin{tabularx}{0.7\textwidth}{X} 女人對蛇說：「園中樹上的果子，我們都可以吃； \end{tabularx} \\ \\ \relax
3:3 & \begin{tabularx}{0.7\textwidth}{X} 只是園子中間那棵樹的果子，神曾說：『你們不可吃，也不可摸，免得你們死。』」 \end{tabularx} \\ \\ \relax
3:4 & \begin{tabularx}{0.7\textwidth}{X} 蛇對女人說：「你們不一定死； \end{tabularx} \\ \\ \relax
3:5 & \begin{tabularx}{0.7\textwidth}{X} 因為神知道，你們吃的日子眼睛就開了，你們就像神一樣知道善惡。」 \end{tabularx} \\ \\ \relax
3:6 & \begin{tabularx}{0.7\textwidth}{X} 於是女人見那棵樹好作食物，又悅人的眼目，那樹令人喜愛，能使人有智慧，她就摘下果子吃了，又給了與她一起的丈夫，他也吃了。 \end{tabularx} \\ \\ \relax
3:7 & \begin{tabularx}{0.7\textwidth}{X} 他們二人的眼睛就開了，知道自己赤身露體，就編織無花果樹的葉子，為自己做成裙子。 \end{tabularx} \\ \\
[1ex]
\hline
\hline
\end{longtable}


\newpage

\section{第78屆~港九培靈會~楊錫鏘~第八講}
\label{sec:671}
\textbf{受造物的審判}
\newline
\newline
連結: \href{https://www.hkbibleconference.org/session-message/view/671}{\texttt{https://www.hkbibleconference.org/session-message/view/671}}
\newline
\newline
\hyperref[sec:670]{< < < PREV SERMON < < <}
~
\hyperref[sec:toc]{[返目錄]}
~
\hyperref[sec:672]{> > > NEXT SERMON > > >}
\newline
\newline
創世記 3:7-3:19
\newline
\begin{longtable}{cl}
\hline
\hline
章節 & 經文 (和合本修訂版)\\
\hline
3:7 & \begin{tabularx}{0.7\textwidth}{X} 他們二人的眼睛就開了，知道自己赤身露體，就編織無花果樹的葉子，為自己做成裙子。 \end{tabularx} \\ \\ \relax
3:8 & \begin{tabularx}{0.7\textwidth}{X} 天起了涼風，那人和他妻子聽見耶和華神在園中來回行走的聲音，就藏在園裡的樹木中，躲避耶和華神的面。 \end{tabularx} \\ \\ \relax
3:9 & \begin{tabularx}{0.7\textwidth}{X} 耶和華神呼喚那人，對他說：「你在哪裡？」 \end{tabularx} \\ \\ \relax
3:10 & \begin{tabularx}{0.7\textwidth}{X} 他說：「我在園中聽見你的聲音，我就害怕；因為我赤身露體，我就藏了起來。」 \end{tabularx} \\ \\ \relax
3:11 & \begin{tabularx}{0.7\textwidth}{X} 耶和華神說：「誰告訴你，你是赤身露體呢？莫非你吃了那樹上所出的，就是我吩咐你不可吃的嗎？」 \end{tabularx} \\ \\ \relax
3:12 & \begin{tabularx}{0.7\textwidth}{X} 那人說：「你賜給我、與我一起的女人，是她把那樹上所出的給我，我就吃了。」 \end{tabularx} \\ \\ \relax
3:13 & \begin{tabularx}{0.7\textwidth}{X} 耶和華神對女人說：「你怎麼會做這種事呢？」女人說：「那蛇引誘我，我就吃了。」 \end{tabularx} \\ \\ \relax
3:14 & \begin{tabularx}{0.7\textwidth}{X} 耶和華神對蛇說：「你既做了這事，就必受詛咒，比一切的牲畜和野獸更重。你必用肚子行走，終生吃土。 \end{tabularx} \\ \\ \relax
3:15 & \begin{tabularx}{0.7\textwidth}{X} 我要使你和女人彼此為仇，你的後裔和女人的後裔也彼此為仇。他要傷你的頭，你要傷他的腳跟。 」 \end{tabularx} \\ \\ \relax
3:16 & \begin{tabularx}{0.7\textwidth}{X} 又對女人說：「我必多多加增你懷胎的痛苦，你生兒女時必多受痛苦。你必戀慕你丈夫，他必管轄你。」 \end{tabularx} \\ \\ \relax
3:17 & \begin{tabularx}{0.7\textwidth}{X} 又對亞當說：「你既聽從你妻子的話，吃了那樹上所出的，就是我吩咐你不可吃的，土地必因你的緣故受詛咒；你必終生勞苦才能從土地得吃的。 \end{tabularx} \\ \\ \relax
3:18 & \begin{tabularx}{0.7\textwidth}{X} 土地必給你長出荊棘和蒺藜來；你也要吃田間的五穀菜蔬。 \end{tabularx} \\ \\ \relax
3:19 & \begin{tabularx}{0.7\textwidth}{X} 你必汗流滿面才有食物可吃，直到你歸了土地，因為你是從土地而出的。你本是塵土，仍要歸回塵土。」 \end{tabularx} \\ \\
[1ex]
\hline
\hline
\end{longtable}


\newpage

\section{第78屆~港九培靈會~楊錫鏘~第九講}
\label{sec:672}
\textbf{受造物的蒙恩}
\newline
\newline
連結: \href{https://www.hkbibleconference.org/session-message/view/672}{\texttt{https://www.hkbibleconference.org/session-message/view/672}}
\newline
\newline
\hyperref[sec:671]{< < < PREV SERMON < < <}
~
\hyperref[sec:toc]{[返目錄]}
~
\hyperref[sec:673]{> > > NEXT SERMON > > >}
\newline
\newline
創世記 3:8-3:24
\newline
\begin{longtable}{cl}
\hline
\hline
章節 & 經文 (和合本修訂版)\\
\hline
3:8 & \begin{tabularx}{0.7\textwidth}{X} 天起了涼風，那人和他妻子聽見耶和華神在園中來回行走的聲音，就藏在園裡的樹木中，躲避耶和華神的面。 \end{tabularx} \\ \\ \relax
3:9 & \begin{tabularx}{0.7\textwidth}{X} 耶和華神呼喚那人，對他說：「你在哪裡？」 \end{tabularx} \\ \\ \relax
3:10 & \begin{tabularx}{0.7\textwidth}{X} 他說：「我在園中聽見你的聲音，我就害怕；因為我赤身露體，我就藏了起來。」 \end{tabularx} \\ \\ \relax
3:11 & \begin{tabularx}{0.7\textwidth}{X} 耶和華神說：「誰告訴你，你是赤身露體呢？莫非你吃了那樹上所出的，就是我吩咐你不可吃的嗎？」 \end{tabularx} \\ \\ \relax
3:12 & \begin{tabularx}{0.7\textwidth}{X} 那人說：「你賜給我、與我一起的女人，是她把那樹上所出的給我，我就吃了。」 \end{tabularx} \\ \\ \relax
3:13 & \begin{tabularx}{0.7\textwidth}{X} 耶和華神對女人說：「你怎麼會做這種事呢？」女人說：「那蛇引誘我，我就吃了。」 \end{tabularx} \\ \\ \relax
3:14 & \begin{tabularx}{0.7\textwidth}{X} 耶和華神對蛇說：「你既做了這事，就必受詛咒，比一切的牲畜和野獸更重。你必用肚子行走，終生吃土。 \end{tabularx} \\ \\ \relax
3:15 & \begin{tabularx}{0.7\textwidth}{X} 我要使你和女人彼此為仇，你的後裔和女人的後裔也彼此為仇。他要傷你的頭，你要傷他的腳跟。 」 \end{tabularx} \\ \\ \relax
3:16 & \begin{tabularx}{0.7\textwidth}{X} 又對女人說：「我必多多加增你懷胎的痛苦，你生兒女時必多受痛苦。你必戀慕你丈夫，他必管轄你。」 \end{tabularx} \\ \\ \relax
3:17 & \begin{tabularx}{0.7\textwidth}{X} 又對亞當說：「你既聽從你妻子的話，吃了那樹上所出的，就是我吩咐你不可吃的，土地必因你的緣故受詛咒；你必終生勞苦才能從土地得吃的。 \end{tabularx} \\ \\ \relax
3:18 & \begin{tabularx}{0.7\textwidth}{X} 土地必給你長出荊棘和蒺藜來；你也要吃田間的五穀菜蔬。 \end{tabularx} \\ \\ \relax
3:19 & \begin{tabularx}{0.7\textwidth}{X} 你必汗流滿面才有食物可吃，直到你歸了土地，因為你是從土地而出的。你本是塵土，仍要歸回塵土。」 \end{tabularx} \\ \\ \relax
3:20 & \begin{tabularx}{0.7\textwidth}{X} 那人給他妻子起名叫夏娃，因為她是眾生之母。 \end{tabularx} \\ \\ \relax
3:21 & \begin{tabularx}{0.7\textwidth}{X} 耶和華神用獸皮做衣服給亞當和他的妻子穿。 \end{tabularx} \\ \\ \relax
3:22 & \begin{tabularx}{0.7\textwidth}{X} 耶和華神說：「看哪，那人已經像我們中間的一個，知道善惡，現在恐怕他又伸手摘生命樹所出的來吃，就永遠活著。」 \end{tabularx} \\ \\ \relax
3:23 & \begin{tabularx}{0.7\textwidth}{X} 耶和華神就驅逐他出伊甸園，使他耕種土地，他原是從土地裡被取出來的。 \end{tabularx} \\ \\ \relax
3:24 & \begin{tabularx}{0.7\textwidth}{X} 耶和華神把那人趕出去，就在伊甸園東邊安設基路伯和發出火焰轉動的劍，把守生命樹的道路。 \end{tabularx} \\ \\
[1ex]
\hline
\hline
\end{longtable}


\newpage

\section{第78屆~港九培靈會~郭文池~第一講}
\label{sec:673}
\textbf{重尋相交生活的豐盛}
\newline
\newline
連結: \href{https://www.hkbibleconference.org/session-message/view/673}{\texttt{https://www.hkbibleconference.org/session-message/view/673}}
\newline
\newline
\hyperref[sec:672]{< < < PREV SERMON < < <}
~
\hyperref[sec:toc]{[返目錄]}
~
\hyperref[sec:681]{> > > NEXT SERMON > > >}
\newline
\newline
約翰一書 1:1-1:4
\newline
\begin{longtable}{cl}
\hline
\hline
章節 & 經文 (和合本修訂版)\\
\hline
1:1 & \begin{tabularx}{0.7\textwidth}{X} 論到從起初原有的生命之道，就是我們所聽見、所看見、親眼看過、親手摸過的－ \end{tabularx} \\ \\ \relax
1:2 & \begin{tabularx}{0.7\textwidth}{X} 這生命已經顯現出來，我們看見了，現在又作見證，把原與父同在，並且向我們顯現過的那永遠的生命傳揚給你們－ \end{tabularx} \\ \\ \relax
1:3 & \begin{tabularx}{0.7\textwidth}{X} 我們把所看見、所聽見的傳揚給你們，為要使你們也與我們有團契，而我們的團契是與父和他兒子耶穌基督所共有的。 \end{tabularx} \\ \\ \relax
1:4 & \begin{tabularx}{0.7\textwidth}{X} 我們把這些事寫給你們，使我們的喜樂得以滿足。 \end{tabularx} \\ \\
[1ex]
\hline
\hline
\end{longtable}


\newpage

\section{第78屆~港九培靈會~郭文池~第二講}
\label{sec:681}
\textbf{重尋聖潔生活的豐盛}
\newline
\newline
連結: \href{https://www.hkbibleconference.org/session-message/view/681}{\texttt{https://www.hkbibleconference.org/session-message/view/681}}
\newline
\newline
\hyperref[sec:673]{< < < PREV SERMON < < <}
~
\hyperref[sec:toc]{[返目錄]}
~
\hyperref[sec:684]{> > > NEXT SERMON > > >}
\newline
\newline
約翰一書 1:5-2:2
\newline
\begin{longtable}{cl}
\hline
\hline
章節 & 經文 (和合本修訂版)\\
\hline
1:5 & \begin{tabularx}{0.7\textwidth}{X} 神就是光，在他毫無黑暗；這是我們從主所聽見，又報給你們的信息。 \end{tabularx} \\ \\ \relax
1:6 & \begin{tabularx}{0.7\textwidth}{X} 我們若說，我們與神有團契，卻仍在黑暗裡行走，就是說謊話，不實行真理了。 \end{tabularx} \\ \\ \relax
1:7 & \begin{tabularx}{0.7\textwidth}{X} 我們若在光明中行走，如同神在光明中，就彼此有團契，他兒子耶穌的血就洗淨我們一切的罪。 \end{tabularx} \\ \\ \relax
1:8 & \begin{tabularx}{0.7\textwidth}{X} 我們若說自己沒有罪，就是欺騙自己，真理就不在我們裡面了。 \end{tabularx} \\ \\ \relax
1:9 & \begin{tabularx}{0.7\textwidth}{X} 我們若認自己的罪，神是信實的，是公義的，必要赦免我們的罪，洗淨我們一切的不義。 \end{tabularx} \\ \\ \relax
1:10 & \begin{tabularx}{0.7\textwidth}{X} 我們若說自己沒有犯過罪，就是把神當作說謊的，他的道就不在我們裡面了。 \end{tabularx} \\ \\ \relax
2:1 & \begin{tabularx}{0.7\textwidth}{X} 我的孩子們哪，我把這些話寫給你們，是要你們不犯罪。若有人犯罪，在父那裡我們有一位中保，就是那義者耶穌基督。 \end{tabularx} \\ \\ \relax
2:2 & \begin{tabularx}{0.7\textwidth}{X} 他為我們的罪作了贖罪祭，不單是為我們的罪，也是為普天下人的罪。 \end{tabularx} \\ \\
[1ex]
\hline
\hline
\end{longtable}


\newpage

\section{第78屆~港九培靈會~郭文池~第三講}
\label{sec:684}
\textbf{重尋光明生活的豐盛}
\newline
\newline
連結: \href{https://www.hkbibleconference.org/session-message/view/684}{\texttt{https://www.hkbibleconference.org/session-message/view/684}}
\newline
\newline
\hyperref[sec:681]{< < < PREV SERMON < < <}
~
\hyperref[sec:toc]{[返目錄]}
~
\hyperref[sec:687]{> > > NEXT SERMON > > >}
\newline
\newline
約翰一書 2:3-2:17
\newline
\begin{longtable}{cl}
\hline
\hline
章節 & 經文 (和合本修訂版)\\
\hline
2:3 & \begin{tabularx}{0.7\textwidth}{X} 我們若遵守神的命令，就知道我們確實認識他。 \end{tabularx} \\ \\ \relax
2:4 & \begin{tabularx}{0.7\textwidth}{X} 人若說「我認識他」，卻不遵守他的命令，就是說謊話的，真理就不在他裡面了。 \end{tabularx} \\ \\ \relax
2:5 & \begin{tabularx}{0.7\textwidth}{X} 凡遵守他的道的，愛神的心確實地在他裡面達到完全了。由此我們知道我們是在他裡面。 \end{tabularx} \\ \\ \relax
2:6 & \begin{tabularx}{0.7\textwidth}{X} 凡說自己住在他裡面的，就該照著他所行的去行。 \end{tabularx} \\ \\ \relax
2:7 & \begin{tabularx}{0.7\textwidth}{X} 親愛的，我寫給你們的不是一條新命令，而是你們從起初所受的舊命令；這舊命令就是你們所聽過的道。 \end{tabularx} \\ \\ \relax
2:8 & \begin{tabularx}{0.7\textwidth}{X} 然而，我寫給你們的是一條新命令，在基督裡是真實的，在你們也是真實的，因為黑暗漸漸消逝，真光已經在照耀。 \end{tabularx} \\ \\ \relax
2:9 & \begin{tabularx}{0.7\textwidth}{X} 人若說自己在光明中，卻恨他的弟兄，他到如今還是在黑暗裡。 \end{tabularx} \\ \\ \relax
2:10 & \begin{tabularx}{0.7\textwidth}{X} 那愛弟兄的，就是住在光明中，他不會使人失足犯罪。 \end{tabularx} \\ \\ \relax
2:11 & \begin{tabularx}{0.7\textwidth}{X} 惟獨那恨弟兄的，是在黑暗裡，也在黑暗裡行走，不知道往哪裡去，因為黑暗使他的眼睛瞎了。 \end{tabularx} \\ \\ \relax
2:12 & \begin{tabularx}{0.7\textwidth}{X} 孩子們哪，我寫信給你們，因為你們的罪藉著基督的名得了赦免。 \end{tabularx} \\ \\ \relax
2:13 & \begin{tabularx}{0.7\textwidth}{X} 父老們啊，我寫信給你們，因為你們認識從起初就有的那一位。青年們哪，我寫信給你們，因為你們勝過了那惡者。 \end{tabularx} \\ \\ \relax
2:14 & \begin{tabularx}{0.7\textwidth}{X} 孩子們哪，我曾寫信給你們，因為你們認識父。父老們啊，我曾寫信給你們，因為你們認識從起初就有的那一位。青年們哪，我曾寫信給你們，因為你們剛強，神的道常存在你們心裡，你們也勝過了那惡者。 \end{tabularx} \\ \\ \relax
2:15 & \begin{tabularx}{0.7\textwidth}{X} 不要愛世界和世界上的東西，若有人愛世界，愛父的心就不在他裡面了。 \end{tabularx} \\ \\ \relax
2:16 & \begin{tabularx}{0.7\textwidth}{X} 因為凡世界上的東西，好比肉體的情慾、眼目的情慾和今生的驕傲，都不是從父來的，而是從世界來的。 \end{tabularx} \\ \\ \relax
2:17 & \begin{tabularx}{0.7\textwidth}{X} 這世界和世上的情慾都要消逝，惟獨那遵行神旨意的人永遠常存。 \end{tabularx} \\ \\
[1ex]
\hline
\hline
\end{longtable}


\newpage

\section{第78屆~港九培靈會~郭文池~第四講}
\label{sec:687}
\textbf{重尋信心生活的豐盛}
\newline
\newline
連結: \href{https://www.hkbibleconference.org/session-message/view/687}{\texttt{https://www.hkbibleconference.org/session-message/view/687}}
\newline
\newline
\hyperref[sec:684]{< < < PREV SERMON < < <}
~
\hyperref[sec:toc]{[返目錄]}
~
\hyperref[sec:689]{> > > NEXT SERMON > > >}
\newline
\newline
約翰一書 2:18-2:29
\newline
\begin{longtable}{cl}
\hline
\hline
章節 & 經文 (和合本修訂版)\\
\hline
2:18 & \begin{tabularx}{0.7\textwidth}{X} 孩子們哪，如今是末世的時光了。你們曾聽過那敵基督者要來，現在有好些敵基督者已經出來了；由此我們就知道，如今是末世的時光了。 \end{tabularx} \\ \\ \relax
2:19 & \begin{tabularx}{0.7\textwidth}{X} 他們從我們中間出去，卻不是屬我們的，若是屬我們的，就必仍舊與我們同在。他們出去，這就顯明他們都不是屬我們的。 \end{tabularx} \\ \\ \relax
2:20 & \begin{tabularx}{0.7\textwidth}{X} 你們從那聖者受了恩膏，並且你們大家都知道。 \end{tabularx} \\ \\ \relax
2:21 & \begin{tabularx}{0.7\textwidth}{X} 我寫信給你們，不是因你們不認識真理，而是因你們認識，並且知道一切虛謊都不是從真理出來的。 \end{tabularx} \\ \\ \relax
2:22 & \begin{tabularx}{0.7\textwidth}{X} 誰是說謊話的呢？不就是那不認耶穌為基督的嗎？那不認父與子的，這個人就是敵基督的。 \end{tabularx} \\ \\ \relax
2:23 & \begin{tabularx}{0.7\textwidth}{X} 凡不認子的，就沒有父；宣認子的，連父也有了。 \end{tabularx} \\ \\ \relax
2:24 & \begin{tabularx}{0.7\textwidth}{X} 論到你們，務要將那從起初所聽見的常存在心裡；若將從起初所聽見的存在心裡，你們就會住在子裡面，也會住在父裡面。 \end{tabularx} \\ \\ \relax
2:25 & \begin{tabularx}{0.7\textwidth}{X} 基督所應許我們的就是永生。 \end{tabularx} \\ \\ \relax
2:26 & \begin{tabularx}{0.7\textwidth}{X} 我將這些話寫給你們，是論到那些迷惑你們的人說的。 \end{tabularx} \\ \\ \relax
2:27 & \begin{tabularx}{0.7\textwidth}{X} 至於你們，你們從基督所受的恩膏常存在你們心裡，並不用人教導你們，自有他的恩膏在凡事上教導你們。這恩膏是真的，不是假的，你們要按這恩膏的教導住在他裡面。 \end{tabularx} \\ \\ \relax
2:28 & \begin{tabularx}{0.7\textwidth}{X} 孩子們哪，你們要住在基督裡面。這樣，他若顯現，我們就可以坦然無懼；當他來臨的時候，在他面前不至於慚愧。 \end{tabularx} \\ \\ \relax
2:29 & \begin{tabularx}{0.7\textwidth}{X} 你們若知道他是公義的，就知道凡行公義的人都是他所生的。 \end{tabularx} \\ \\
[1ex]
\hline
\hline
\end{longtable}


\newpage

\section{第78屆~港九培靈會~郭文池~第五講}
\label{sec:689}
\textbf{重尋無罪生活的豐盛}
\newline
\newline
連結: \href{https://www.hkbibleconference.org/session-message/view/689}{\texttt{https://www.hkbibleconference.org/session-message/view/689}}
\newline
\newline
\hyperref[sec:687]{< < < PREV SERMON < < <}
~
\hyperref[sec:toc]{[返目錄]}
~
\hyperref[sec:694]{> > > NEXT SERMON > > >}
\newline
\newline
約翰一書 3:1-3:10
\newline
\begin{longtable}{cl}
\hline
\hline
章節 & 經文 (和合本修訂版)\\
\hline
3:1 & \begin{tabularx}{0.7\textwidth}{X} 你們看父賜給我們的是何等的慈愛，讓我們得以稱為神的兒女；我們也真是他的兒女。世人不認識我們的理由，是因他們未曾認識父。 \end{tabularx} \\ \\ \relax
3:2 & \begin{tabularx}{0.7\textwidth}{X} 親愛的，我們現在是神的兒女，將來如何還未顯明。我們所知道的是：基督顯現的時候，我們會像他，因為我們將見到他的本相。 \end{tabularx} \\ \\ \relax
3:3 & \begin{tabularx}{0.7\textwidth}{X} 凡對他有這指望的，就潔淨自己，像他是潔淨的一樣。 \end{tabularx} \\ \\ \relax
3:4 & \begin{tabularx}{0.7\textwidth}{X} 凡犯罪的，就是做違背律法的事；違背律法就是罪。 \end{tabularx} \\ \\ \relax
3:5 & \begin{tabularx}{0.7\textwidth}{X} 你們知道，基督曾顯現是要除掉罪；在他並沒有罪。 \end{tabularx} \\ \\ \relax
3:6 & \begin{tabularx}{0.7\textwidth}{X} 凡住在他裡面的，不犯罪；凡犯罪的，未曾看見他，也未曾認識他。 \end{tabularx} \\ \\ \relax
3:7 & \begin{tabularx}{0.7\textwidth}{X} 孩子們哪，不要讓人迷惑了你們；行義的才是義人，正如基督是義的。 \end{tabularx} \\ \\ \relax
3:8 & \begin{tabularx}{0.7\textwidth}{X} 犯罪的是出於魔鬼，因為魔鬼從起初就犯罪。神的兒子顯現出來，是為了要毀滅魔鬼的作為。 \end{tabularx} \\ \\ \relax
3:9 & \begin{tabularx}{0.7\textwidth}{X} 凡從神生的，不犯罪，因神的道存在他裡面，他也不能犯罪，因為他是由神所生的。 \end{tabularx} \\ \\ \relax
3:10 & \begin{tabularx}{0.7\textwidth}{X} 這就顯明誰是神的兒女，誰是魔鬼的兒女了。凡不行義的，不是出於神，不愛他弟兄的，也是如此。 \end{tabularx} \\ \\
[1ex]
\hline
\hline
\end{longtable}


\newpage

\section{第78屆~港九培靈會~郭文池~第六講}
\label{sec:694}
\textbf{重尋靈裡生活的豐盛}
\newline
\newline
連結: \href{https://www.hkbibleconference.org/session-message/view/694}{\texttt{https://www.hkbibleconference.org/session-message/view/694}}
\newline
\newline
\hyperref[sec:689]{< < < PREV SERMON < < <}
~
\hyperref[sec:toc]{[返目錄]}
~
\hyperref[sec:695]{> > > NEXT SERMON > > >}
\newline
\newline
約翰一書 3:11-4:6
\newline
\begin{longtable}{cl}
\hline
\hline
章節 & 經文 (和合本修訂版)\\
\hline
3:11 & \begin{tabularx}{0.7\textwidth}{X} 我們要彼此相愛。這就是你們從起初所聽到的信息。 \end{tabularx} \\ \\ \relax
3:12 & \begin{tabularx}{0.7\textwidth}{X} 不要像該隱；他是屬那邪惡者，殺了自己的弟弟。為甚麼殺了他呢？因為自己的行為是邪惡的，而弟弟的行為是正直的。 \end{tabularx} \\ \\ \relax
3:13 & \begin{tabularx}{0.7\textwidth}{X} 弟兄們，世人若恨你們，不要驚訝。 \end{tabularx} \\ \\ \relax
3:14 & \begin{tabularx}{0.7\textwidth}{X} 我們知道，我們已經出死入生了，因為我們愛弟兄。沒有愛心的，仍住在死中。 \end{tabularx} \\ \\ \relax
3:15 & \begin{tabularx}{0.7\textwidth}{X} 凡恨自己弟兄的，就是殺人的；你們知道，凡殺人的，沒有永生住在他裡面。 \end{tabularx} \\ \\ \relax
3:16 & \begin{tabularx}{0.7\textwidth}{X} 基督為我們捨命，我們從此就知道何為愛；我們也當為弟兄捨命。 \end{tabularx} \\ \\ \relax
3:17 & \begin{tabularx}{0.7\textwidth}{X} 凡有世上財物的，看見弟兄缺乏，卻關閉了惻隱的心，神的愛怎能住在他裡面呢？ \end{tabularx} \\ \\ \relax
3:18 & \begin{tabularx}{0.7\textwidth}{X} 孩子們哪，我們相愛，不要只在言語或舌頭上，總要以行為和真誠表現出來。 \end{tabularx} \\ \\ \relax
3:19 & \begin{tabularx}{0.7\textwidth}{X} 從這一點，我們會知道，我們是出於真理的，並且我們在神面前可以安心， \end{tabularx} \\ \\ \relax
3:20 & \begin{tabularx}{0.7\textwidth}{X} 即使我們的心責備自己，神比我們的心大，他知道一切。 \end{tabularx} \\ \\ \relax
3:21 & \begin{tabularx}{0.7\textwidth}{X} 親愛的，我們的心若不責備我們，在神面前就可以坦然無懼了。 \end{tabularx} \\ \\ \relax
3:22 & \begin{tabularx}{0.7\textwidth}{X} 我們一切所求的，就從他得著，因為我們遵守他的命令，行他所喜悅的事。 \end{tabularx} \\ \\ \relax
3:23 & \begin{tabularx}{0.7\textwidth}{X} 神的命令就是：我們要信他兒子耶穌基督的名，並且照他所賜給我們的命令彼此相愛。 \end{tabularx} \\ \\ \relax
3:24 & \begin{tabularx}{0.7\textwidth}{X} 遵守神命令的，住在神裡面，而神也住在他裡面。從這一點，我們知道神住在我們裡面，這是由於他所賜給我們的聖靈。 \end{tabularx} \\ \\ \relax
4:1 & \begin{tabularx}{0.7\textwidth}{X} 親愛的，一切的靈不可都信，總要察驗那些靈是否出於神，因為有許多假先知已經來到世上。 \end{tabularx} \\ \\ \relax
4:2 & \begin{tabularx}{0.7\textwidth}{X} 凡宣認耶穌基督是成了肉身而來的靈就是出於神的，由此你們可以認出神的靈來； \end{tabularx} \\ \\ \relax
4:3 & \begin{tabularx}{0.7\textwidth}{X} 凡不宣認耶穌的靈，不是出於神。這是那敵基督者的靈；你們從前聽見他要來，現在他已經在世上了。 \end{tabularx} \\ \\ \relax
4:4 & \begin{tabularx}{0.7\textwidth}{X} 孩子們哪，你們是屬神的，並且勝過了假先知，因為那在你們裡面的比那在世界上的更大。 \end{tabularx} \\ \\ \relax
4:5 & \begin{tabularx}{0.7\textwidth}{X} 他們是屬世界的，所以講論世界的事，而世人也聽從他們。 \end{tabularx} \\ \\ \relax
4:6 & \begin{tabularx}{0.7\textwidth}{X} 我們是屬神的，認識神的就聽從我們；不屬神的就不聽從我們。從此我們可以認出真理的靈和錯謬的靈來。 \end{tabularx} \\ \\
[1ex]
\hline
\hline
\end{longtable}


\newpage

\section{第78屆~港九培靈會~郭文池~第七講}
\label{sec:695}
\textbf{重尋相愛生活的豐盛}
\newline
\newline
連結: \href{https://www.hkbibleconference.org/session-message/view/695}{\texttt{https://www.hkbibleconference.org/session-message/view/695}}
\newline
\newline
\hyperref[sec:694]{< < < PREV SERMON < < <}
~
\hyperref[sec:toc]{[返目錄]}
~
\hyperref[sec:699]{> > > NEXT SERMON > > >}
\newline
\newline
約翰一書 4:7-5:4
\newline
\begin{longtable}{cl}
\hline
\hline
章節 & 經文 (和合本修訂版)\\
\hline
4:7 & \begin{tabularx}{0.7\textwidth}{X} 親愛的，我們要彼此相愛，因為愛是從神來的。凡有愛的都是由神而生，並且認識神。 \end{tabularx} \\ \\ \relax
4:8 & \begin{tabularx}{0.7\textwidth}{X} 沒有愛的就不認識神，因為神就是愛。 \end{tabularx} \\ \\ \relax
4:9 & \begin{tabularx}{0.7\textwidth}{X} 神差他獨一的兒子到世上來，使我們藉著他得生命；由此，神對我們的愛就顯明了。 \end{tabularx} \\ \\ \relax
4:10 & \begin{tabularx}{0.7\textwidth}{X} 不是我們愛神，而是神愛我們，差他的兒子為我們的罪作了贖罪祭；這就是愛。 \end{tabularx} \\ \\ \relax
4:11 & \begin{tabularx}{0.7\textwidth}{X} 親愛的，既然神這樣愛我們，我們也要彼此相愛。 \end{tabularx} \\ \\ \relax
4:12 & \begin{tabularx}{0.7\textwidth}{X} 從來沒有人見過神，我們若彼此相愛，神就住在我們裡面，他的愛在我們裡面得以完滿了。 \end{tabularx} \\ \\ \relax
4:13 & \begin{tabularx}{0.7\textwidth}{X} 因為神將他的靈賜給我們，由此我們知道我們是住在他裡面，而他也住在我們裡面。 \end{tabularx} \\ \\ \relax
4:14 & \begin{tabularx}{0.7\textwidth}{X} 父差子作世人的救主，這是我們所看見並且作見證的。 \end{tabularx} \\ \\ \relax
4:15 & \begin{tabularx}{0.7\textwidth}{X} 凡宣認耶穌為神兒子的，神就住在他裡面，而他也住在神裡面。 \end{tabularx} \\ \\ \relax
4:16 & \begin{tabularx}{0.7\textwidth}{X} 我們知道並且深信神是愛我們的。神就是愛，住在愛裡面的就是住在神裡面；神也住在他裡面。 \end{tabularx} \\ \\ \relax
4:17 & \begin{tabularx}{0.7\textwidth}{X} 由此，愛在我們裡面得以完滿：我們可以在審判的日子坦然無懼，因為基督如何，我們在這世上也如何。 \end{tabularx} \\ \\ \relax
4:18 & \begin{tabularx}{0.7\textwidth}{X} 在愛裡沒有懼怕；完滿的愛把懼怕驅逐出去，因為懼怕裡含著懲罰，懼怕的人在愛裡尚未得到完滿。 \end{tabularx} \\ \\ \relax
4:19 & \begin{tabularx}{0.7\textwidth}{X} 我們愛，因為神先愛我們。 \end{tabularx} \\ \\ \relax
4:20 & \begin{tabularx}{0.7\textwidth}{X} 人若說「我愛神」，卻恨他的弟兄，就是說謊了；不愛他看得見的弟兄，就不能愛看不見的神。 \end{tabularx} \\ \\ \relax
4:21 & \begin{tabularx}{0.7\textwidth}{X} 愛神的，也要愛弟兄；這是我們從神所受的命令。 \end{tabularx} \\ \\ \relax
5:1 & \begin{tabularx}{0.7\textwidth}{X} 凡信耶穌是基督的，都是從神生的；凡愛生他之神的，也必愛從神生的。 \end{tabularx} \\ \\ \relax
5:2 & \begin{tabularx}{0.7\textwidth}{X} 我們愛神，又實行他的命令，由此就知道我們愛神的兒女了。 \end{tabularx} \\ \\ \relax
5:3 & \begin{tabularx}{0.7\textwidth}{X} 我們遵守神的命令，這就是愛他了，而且他的命令並不是難守的。 \end{tabularx} \\ \\ \relax
5:4 & \begin{tabularx}{0.7\textwidth}{X} 因為凡從神生的就勝過世界；使我們勝過世界的就是我們的信心。 \end{tabularx} \\ \\
[1ex]
\hline
\hline
\end{longtable}


\newpage

\section{第78屆~港九培靈會~郭文池~第八講}
\label{sec:699}
\textbf{重尋救恩生活的豐盛}
\newline
\newline
連結: \href{https://www.hkbibleconference.org/session-message/view/699}{\texttt{https://www.hkbibleconference.org/session-message/view/699}}
\newline
\newline
\hyperref[sec:695]{< < < PREV SERMON < < <}
~
\hyperref[sec:toc]{[返目錄]}
~
\hyperref[sec:702]{> > > NEXT SERMON > > >}
\newline
\newline
約翰一書 5:5-5:12
\newline
\begin{longtable}{cl}
\hline
\hline
章節 & 經文 (和合本修訂版)\\
\hline
5:5 & \begin{tabularx}{0.7\textwidth}{X} 勝過世界的是誰呢？不就是那信耶穌是神兒子的嗎？ \end{tabularx} \\ \\ \relax
5:6 & \begin{tabularx}{0.7\textwidth}{X} 這藉著水和血而來的，就是耶穌基督，不是單用水，而是用水又用血，並且有聖靈作見證，因為聖靈就是真理。 \end{tabularx} \\ \\ \relax
5:7 & \begin{tabularx}{0.7\textwidth}{X} 作見證的有三： \end{tabularx} \\ \\ \relax
5:8 & \begin{tabularx}{0.7\textwidth}{X} 就是聖靈、水與血，這三樣也都是一致的。 \end{tabularx} \\ \\ \relax
5:9 & \begin{tabularx}{0.7\textwidth}{X} 既然我們領受人的見證，神的見證更該領受了，因為神的見證是為他兒子作的。 \end{tabularx} \\ \\ \relax
5:10 & \begin{tabularx}{0.7\textwidth}{X} 信神兒子的，就有這見證在他心裡；不信神的，就是把神當作說謊的，因為不信神為他兒子作的見證。 \end{tabularx} \\ \\ \relax
5:11 & \begin{tabularx}{0.7\textwidth}{X} 這見證就是：神賜給我們永生，而這永生是在他兒子裡面的。 \end{tabularx} \\ \\ \relax
5:12 & \begin{tabularx}{0.7\textwidth}{X} 那有神兒子的，就有生命；沒有神兒子的，就沒有生命。 \end{tabularx} \\ \\
[1ex]
\hline
\hline
\end{longtable}


\newpage

\section{第78屆~港九培靈會~郭文池~第九講}
\label{sec:702}
\textbf{重尋平安生活的豐盛}
\newline
\newline
連結: \href{https://www.hkbibleconference.org/session-message/view/702}{\texttt{https://www.hkbibleconference.org/session-message/view/702}}
\newline
\newline
\hyperref[sec:699]{< < < PREV SERMON < < <}
~
\hyperref[sec:toc]{[返目錄]}
~
\hyperref[sec:291]{> > > NEXT SERMON > > >}
\newline
\newline
約翰一書 5:13-5:21
\newline
\begin{longtable}{cl}
\hline
\hline
章節 & 經文 (和合本修訂版)\\
\hline
5:13 & \begin{tabularx}{0.7\textwidth}{X} 我把這些話寫給你們信奉神兒子之名的人，要讓你們知道自己有永生。 \end{tabularx} \\ \\ \relax
5:14 & \begin{tabularx}{0.7\textwidth}{X} 我們若照著神的旨意祈求，他就垂聽我們；這就是我們對他所存坦然無懼的心。 \end{tabularx} \\ \\ \relax
5:15 & \begin{tabularx}{0.7\textwidth}{X} 既然我們知道他聽我們一切所求的，就知道我們所求於他的，無不得著。 \end{tabularx} \\ \\ \relax
5:16 & \begin{tabularx}{0.7\textwidth}{X} 人若看見弟兄犯了不至於死的罪，就要為他祈求，神必將生命賜給他—有些人犯的罪是不至於死的；有的是至於死的罪，我不是說要為這罪祈求。 \end{tabularx} \\ \\ \relax
5:17 & \begin{tabularx}{0.7\textwidth}{X} 一切不義的事都是罪，但也有不至於死的罪。 \end{tabularx} \\ \\ \relax
5:18 & \begin{tabularx}{0.7\textwidth}{X} 我們知道，凡從神生的，必不犯罪；從神生的那一位，必保守他，那邪惡者無法加害於他。 \end{tabularx} \\ \\ \relax
5:19 & \begin{tabularx}{0.7\textwidth}{X} 我們知道，我們是屬神的，而全世界都伏在那邪惡者的權勢之下。 \end{tabularx} \\ \\ \relax
5:20 & \begin{tabularx}{0.7\textwidth}{X} 我們知道，神的兒子已經來到，並且將悟性賜給我們，使我們認識那位真實者，我們也在那位真實者裡面，就是在他兒子耶穌基督裡面。這是真神，也是永生。 \end{tabularx} \\ \\ \relax
5:21 & \begin{tabularx}{0.7\textwidth}{X} 孩子們哪，你們要遠避偶像。 \end{tabularx} \\ \\
[1ex]
\hline
\hline
\end{longtable}


\newpage

\section{第79屆~港九培靈會~奧思邦~第一講}
\label{sec:291}
\textbf{比喻目的:回應呼召}
\newline
\newline
連結: \href{https://www.hkbibleconference.org/session-message/view/291}{\texttt{https://www.hkbibleconference.org/session-message/view/291}}
\newline
\newline
\hyperref[sec:702]{< < < PREV SERMON < < <}
~
\hyperref[sec:toc]{[返目錄]}
~
\hyperref[sec:292]{> > > NEXT SERMON > > >}
\newline
\newline
馬可福音 4:1-4:20
\newline
\begin{longtable}{cl}
\hline
\hline
章節 & 經文 (和合本修訂版)\\
\hline
4:1 & \begin{tabularx}{0.7\textwidth}{X} 耶穌又在海邊教導人。有一大群人到他那裡聚集，他只好上船坐下。船在海裡，眾人都靠近海，站在岸上。 \end{tabularx} \\ \\ \relax
4:2 & \begin{tabularx}{0.7\textwidth}{X} 耶穌就用許多比喻教導他們。在教導的時候，他對他們說： \end{tabularx} \\ \\ \relax
4:3 & \begin{tabularx}{0.7\textwidth}{X} 「你們聽啊，有一個撒種的出去撒種。 \end{tabularx} \\ \\ \relax
4:4 & \begin{tabularx}{0.7\textwidth}{X} 他撒的時候，有的落在路旁，飛鳥來把它吃掉了。 \end{tabularx} \\ \\ \relax
4:5 & \begin{tabularx}{0.7\textwidth}{X} 有的落在土淺的石頭地上，因為土不深，很快就長出苗來， \end{tabularx} \\ \\ \relax
4:6 & \begin{tabularx}{0.7\textwidth}{X} 太陽出來一曬，因為沒有根就枯乾了。 \end{tabularx} \\ \\ \relax
4:7 & \begin{tabularx}{0.7\textwidth}{X} 有的落在荊棘裡，荊棘長起來，把它擠住了，就結不出果實。 \end{tabularx} \\ \\ \relax
4:8 & \begin{tabularx}{0.7\textwidth}{X} 又有的落在好土裡，就發芽長大，結出果實，有三十倍的，有六十倍的，有一百倍的。」 \end{tabularx} \\ \\ \relax
4:9 & \begin{tabularx}{0.7\textwidth}{X} 耶穌又說：「有耳可聽的，就應當聽！」 \end{tabularx} \\ \\ \relax
4:10 & \begin{tabularx}{0.7\textwidth}{X} 耶穌獨自一人的時候，跟隨他的人和十二使徒問他這些比喻的意思。 \end{tabularx} \\ \\ \relax
4:11 & \begin{tabularx}{0.7\textwidth}{X} 耶穌對他們說：「神國的奧祕只讓你們知道，若是對外人講，凡事就用比喻， \end{tabularx} \\ \\ \relax
4:12 & \begin{tabularx}{0.7\textwidth}{X} 要他們看了又看，卻看不清，聽了又聽，卻不明白，免得他們回轉過來，獲得赦免。」 \end{tabularx} \\ \\ \relax
4:13 & \begin{tabularx}{0.7\textwidth}{X} 耶穌又對他們說：「你們不明白這比喻嗎？這樣怎能明白一切的比喻呢？ \end{tabularx} \\ \\ \relax
4:14 & \begin{tabularx}{0.7\textwidth}{X} 撒種的人所撒的就是道。 \end{tabularx} \\ \\ \relax
4:15 & \begin{tabularx}{0.7\textwidth}{X} 那撒在路旁的種子，就是人聽了道，撒但立刻來，把撒在他們心裡的道奪了去。 \end{tabularx} \\ \\ \relax
4:16 & \begin{tabularx}{0.7\textwidth}{X} 那撒在石頭地上的，就是人聽了道，立刻歡喜領受， \end{tabularx} \\ \\ \relax
4:17 & \begin{tabularx}{0.7\textwidth}{X} 因心裡沒有根，不過是暫時的，一旦為道遭受患難或迫害，立刻就跌倒。 \end{tabularx} \\ \\ \relax
4:18 & \begin{tabularx}{0.7\textwidth}{X} 還有那撒在荊棘裡的，就是人聽了道， \end{tabularx} \\ \\ \relax
4:19 & \begin{tabularx}{0.7\textwidth}{X} 後來有世上的憂慮、錢財的迷惑，和別樣的私慾進來，把道擠住了，結不出果實。 \end{tabularx} \\ \\ \relax
4:20 & \begin{tabularx}{0.7\textwidth}{X} 那撒在好土裡的，就是人聽了道，領受了，並且結了果實，有三十倍的，有六十倍的，有一百倍的。」 \end{tabularx} \\ \\
[1ex]
\hline
\hline
\end{longtable}


\newpage

\section{第79屆~港九培靈會~奧思邦~第二講}
\label{sec:292}
\textbf{神恩浩瀚,不可思議}
\newline
\newline
連結: \href{https://www.hkbibleconference.org/session-message/view/292}{\texttt{https://www.hkbibleconference.org/session-message/view/292}}
\newline
\newline
\hyperref[sec:291]{< < < PREV SERMON < < <}
~
\hyperref[sec:toc]{[返目錄]}
~
\hyperref[sec:293]{> > > NEXT SERMON > > >}
\newline
\newline
路加福音 15:11-15:32
\newline
\begin{longtable}{cl}
\hline
\hline
章節 & 經文 (和合本修訂版)\\
\hline
15:11 & \begin{tabularx}{0.7\textwidth}{X} 耶穌又說：「一個人有兩個兒子。 \end{tabularx} \\ \\ \relax
15:12 & \begin{tabularx}{0.7\textwidth}{X} 小兒子對父親說：『父親，請你把我應得的家業分給我。』他父親就把財產分給他們。 \end{tabularx} \\ \\ \relax
15:13 & \begin{tabularx}{0.7\textwidth}{X} 過了不多幾天，小兒子把他一切所有的都收拾起來，往遠方去了。在那裡，他任意放蕩，浪費錢財。 \end{tabularx} \\ \\ \relax
15:14 & \begin{tabularx}{0.7\textwidth}{X} 他耗盡了一切所有的，又恰逢那地方有大饑荒，就窮困起來。 \end{tabularx} \\ \\ \relax
15:15 & \begin{tabularx}{0.7\textwidth}{X} 於是他去投靠當地的一個居民，那人打發他到田裡去放豬。 \end{tabularx} \\ \\ \relax
15:16 & \begin{tabularx}{0.7\textwidth}{X} 他恨不得拿豬所吃的豆莢充飢，也沒有人給他甚麼吃的。 \end{tabularx} \\ \\ \relax
15:17 & \begin{tabularx}{0.7\textwidth}{X} 他醒悟過來，就說：『我父親有多少雇工，糧食有餘，我倒在這裡餓死嗎？ \end{tabularx} \\ \\ \relax
15:18 & \begin{tabularx}{0.7\textwidth}{X} 我要起來，到我父親那裡去，對他說：父親！我得罪了天，又得罪了你， \end{tabularx} \\ \\ \relax
15:19 & \begin{tabularx}{0.7\textwidth}{X} 從今以後，我不配稱為你的兒子，把我當作一個雇工吧。』 \end{tabularx} \\ \\ \relax
15:20 & \begin{tabularx}{0.7\textwidth}{X} 於是他起來，往他父親那裡去。相離還遠，他父親看見，就動了慈心，跑去擁抱著他，連連親他。 \end{tabularx} \\ \\ \relax
15:21 & \begin{tabularx}{0.7\textwidth}{X} 兒子對他說：『父親！我得罪了天，又得罪了你，從今以後，我不配稱為你的兒子。』 \end{tabularx} \\ \\ \relax
15:22 & \begin{tabularx}{0.7\textwidth}{X} 父親卻吩咐僕人：『快把那上好的袍子拿出來給他穿，把戒指戴在他指頭上，把鞋穿在他腳上， \end{tabularx} \\ \\ \relax
15:23 & \begin{tabularx}{0.7\textwidth}{X} 把那肥牛犢牽來宰了，我們來吃喝慶祝； \end{tabularx} \\ \\ \relax
15:24 & \begin{tabularx}{0.7\textwidth}{X} 因為我這個兒子是死而復活，失而復得的。』他們就開始慶祝。 \end{tabularx} \\ \\ \relax
15:25 & \begin{tabularx}{0.7\textwidth}{X} 「那時，大兒子正在田裡。他回來，離家不遠時，聽見奏樂跳舞的聲音， \end{tabularx} \\ \\ \relax
15:26 & \begin{tabularx}{0.7\textwidth}{X} 就叫一個僮僕來，問是甚麼事。 \end{tabularx} \\ \\ \relax
15:27 & \begin{tabularx}{0.7\textwidth}{X} 僮僕對他說：『你弟弟回來了，你父親因為他無災無病地回來，把肥牛犢宰了。』 \end{tabularx} \\ \\ \relax
15:28 & \begin{tabularx}{0.7\textwidth}{X} 大兒子就生氣，不肯進去，他父親出來勸他。 \end{tabularx} \\ \\ \relax
15:29 & \begin{tabularx}{0.7\textwidth}{X} 他對父親說：『你看，我服侍你這麼多年，從來沒有違背過你的命令，而你從來沒有給我一隻小山羊，叫我和朋友們一同快樂。 \end{tabularx} \\ \\ \relax
15:30 & \begin{tabularx}{0.7\textwidth}{X} 但你這個兒子和娼妓吃光了你的財產，他一回來，你倒為他宰了肥牛犢。』 \end{tabularx} \\ \\ \relax
15:31 & \begin{tabularx}{0.7\textwidth}{X} 父親對他說：『兒啊！你常和我同在，我所有的一切都是你的； \end{tabularx} \\ \\ \relax
15:32 & \begin{tabularx}{0.7\textwidth}{X} 可是你這個弟弟是死而復活，失而復得的，所以我們理當歡喜慶祝。』」 \end{tabularx} \\ \\
[1ex]
\hline
\hline
\end{longtable}


\newpage

\section{第79屆~港九培靈會~奧思邦~第三講}
\label{sec:293}
\textbf{無知的人,罪有應得}
\newline
\newline
連結: \href{https://www.hkbibleconference.org/session-message/view/293}{\texttt{https://www.hkbibleconference.org/session-message/view/293}}
\newline
\newline
\hyperref[sec:292]{< < < PREV SERMON < < <}
~
\hyperref[sec:toc]{[返目錄]}
~
\hyperref[sec:294]{> > > NEXT SERMON > > >}
\newline
\newline
路加福音 12:13-12:21
\newline
\begin{longtable}{cl}
\hline
\hline
章節 & 經文 (和合本修訂版)\\
\hline
12:13 & \begin{tabularx}{0.7\textwidth}{X} 人群中有一個人對耶穌說：「老師！請你吩咐我的兄弟和我分家產。」 \end{tabularx} \\ \\ \relax
12:14 & \begin{tabularx}{0.7\textwidth}{X} 耶穌對他說：「你這個人！誰立我作你們的判官，或給你們分家產的呢？」 \end{tabularx} \\ \\ \relax
12:15 & \begin{tabularx}{0.7\textwidth}{X} 於是他對他們說：「你們要謹慎自守，躲避一切的貪心，因為人的生命不在於家道豐富。」 \end{tabularx} \\ \\ \relax
12:16 & \begin{tabularx}{0.7\textwidth}{X} 然後他用比喻對他們說：「有一個財主，田地出產豐富。 \end{tabularx} \\ \\ \relax
12:17 & \begin{tabularx}{0.7\textwidth}{X} 他自己心裡想：『我的出產沒有地方儲藏，怎麼辦呢？』 \end{tabularx} \\ \\ \relax
12:18 & \begin{tabularx}{0.7\textwidth}{X} 就說：『我要這麼辦：要把我的倉庫拆了，另蓋更大的，在那裡好儲藏我一切的糧食和財物， \end{tabularx} \\ \\ \relax
12:19 & \begin{tabularx}{0.7\textwidth}{X} 然後要對我自己說：你這個人哪，你有許多財物積存，可供多年享用，只管安安逸逸吃喝快樂吧！』 \end{tabularx} \\ \\ \relax
12:20 & \begin{tabularx}{0.7\textwidth}{X} 神卻對他說：『無知的人哪！今夜就要你的性命，你所預備的要歸誰呢？』 \end{tabularx} \\ \\ \relax
12:21 & \begin{tabularx}{0.7\textwidth}{X} 凡為自己積財，在神面前卻不富足的，也是這樣。」 \end{tabularx} \\ \\
[1ex]
\hline
\hline
\end{longtable}


\newpage

\section{第79屆~港九培靈會~奧思邦~第四講}
\label{sec:294}
\textbf{神發邀請,切勿推延}
\newline
\newline
連結: \href{https://www.hkbibleconference.org/session-message/view/294}{\texttt{https://www.hkbibleconference.org/session-message/view/294}}
\newline
\newline
\hyperref[sec:293]{< < < PREV SERMON < < <}
~
\hyperref[sec:toc]{[返目錄]}
~
\hyperref[sec:295]{> > > NEXT SERMON > > >}
\newline
\newline
馬太福音 22:1-22:14
\newline
\begin{longtable}{cl}
\hline
\hline
章節 & 經文 (和合本修訂版)\\
\hline
22:1 & \begin{tabularx}{0.7\textwidth}{X} 耶穌又用比喻對他們說： \end{tabularx} \\ \\ \relax
22:2 & \begin{tabularx}{0.7\textwidth}{X} 「天國好比一個王為他兒子擺設娶親的宴席。 \end{tabularx} \\ \\ \relax
22:3 & \begin{tabularx}{0.7\textwidth}{X} 他打發僕人去，請那些被邀的人來赴宴，他們卻不肯來。 \end{tabularx} \\ \\ \relax
22:4 & \begin{tabularx}{0.7\textwidth}{X} 王又打發別的僕人，說：『你們去告訴那被邀的人，我的宴席已經預備好了，牛和肥畜已經宰了，各樣都齊備，請你們來赴宴。』 \end{tabularx} \\ \\ \relax
22:5 & \begin{tabularx}{0.7\textwidth}{X} 那些人不理就走了，一個到自己田裡去，一個做買賣去。 \end{tabularx} \\ \\ \relax
22:6 & \begin{tabularx}{0.7\textwidth}{X} 其餘的抓住僕人，凌辱他們，把他們殺了。 \end{tabularx} \\ \\ \relax
22:7 & \begin{tabularx}{0.7\textwidth}{X} 王就大怒，發兵除滅那些兇手，燒燬他們的城。 \end{tabularx} \\ \\ \relax
22:8 & \begin{tabularx}{0.7\textwidth}{X} 於是王對僕人說：『喜宴已經齊備，只是所邀的人不配。 \end{tabularx} \\ \\ \relax
22:9 & \begin{tabularx}{0.7\textwidth}{X} 所以你們要往岔路口上去，凡遇見的，都邀來赴宴。』 \end{tabularx} \\ \\ \relax
22:10 & \begin{tabularx}{0.7\textwidth}{X} 那些僕人就出去，到大路上，凡遇見的，不論善惡都招聚了來，宴席上就坐滿了客人。 \end{tabularx} \\ \\ \relax
22:11 & \begin{tabularx}{0.7\textwidth}{X} 王進來見賓客，看到那裡有一個沒有穿禮服的， \end{tabularx} \\ \\ \relax
22:12 & \begin{tabularx}{0.7\textwidth}{X} 就對他說：『朋友，你到這裡來怎麼不穿禮服呢？』那人無言可答。 \end{tabularx} \\ \\ \relax
22:13 & \begin{tabularx}{0.7\textwidth}{X} 於是王對侍從說：『捆起他的手腳，把他扔在外邊的黑暗裡；在那裡他要哀哭切齒了。』 \end{tabularx} \\ \\ \relax
22:14 & \begin{tabularx}{0.7\textwidth}{X} 因為被召的人多，選上的人少。」 \end{tabularx} \\ \\
[1ex]
\hline
\hline
\end{longtable}


\newpage

\section{第79屆~港九培靈會~奧思邦~第五講}
\label{sec:295}
\textbf{善牧耶穌,唯一門路}
\newline
\newline
連結: \href{https://www.hkbibleconference.org/session-message/view/295}{\texttt{https://www.hkbibleconference.org/session-message/view/295}}
\newline
\newline
\hyperref[sec:294]{< < < PREV SERMON < < <}
~
\hyperref[sec:toc]{[返目錄]}
~
\hyperref[sec:296]{> > > NEXT SERMON > > >}
\newline
\newline
約翰福音 10:1-10:18
\newline
\begin{longtable}{cl}
\hline
\hline
章節 & 經文 (和合本修訂版)\\
\hline
10:1 & \begin{tabularx}{0.7\textwidth}{X} 「我實實在在地告訴你們，那不從門進羊圈，倒從別處爬進去的，就是賊，就是強盜。 \end{tabularx} \\ \\ \relax
10:2 & \begin{tabularx}{0.7\textwidth}{X} 那從門進去的才是羊的牧人。 \end{tabularx} \\ \\ \relax
10:3 & \begin{tabularx}{0.7\textwidth}{X} 看門的給他開門，羊也聽他的聲音。他按著名叫自己的羊，把羊領出來。 \end{tabularx} \\ \\ \relax
10:4 & \begin{tabularx}{0.7\textwidth}{X} 當他把自己的羊都放出來，就走在前面，羊也跟著他，因為牠們認得他的聲音。 \end{tabularx} \\ \\ \relax
10:5 & \begin{tabularx}{0.7\textwidth}{X} 羊絕不跟陌生人，反而會逃走，因為不認得陌生人的聲音。」 \end{tabularx} \\ \\ \relax
10:6 & \begin{tabularx}{0.7\textwidth}{X} 耶穌把這比方告訴他們，但他們不明白他所說的是甚麼。 \end{tabularx} \\ \\ \relax
10:7 & \begin{tabularx}{0.7\textwidth}{X} 所以，耶穌又對他們說：「我實實在在地告訴你們，我就是羊的門。 \end{tabularx} \\ \\ \relax
10:8 & \begin{tabularx}{0.7\textwidth}{X} 凡在我以前來的都是賊，是強盜；羊沒有聽從他們。 \end{tabularx} \\ \\ \relax
10:9 & \begin{tabularx}{0.7\textwidth}{X} 我就是門，凡從我進來的，必得安全，並且可進出，找到草吃。 \end{tabularx} \\ \\ \relax
10:10 & \begin{tabularx}{0.7\textwidth}{X} 盜賊來，無非要偷竊、殺害、毀壞；我來了，是要羊得生命，並且得的更豐盛。 \end{tabularx} \\ \\ \relax
10:11 & \begin{tabularx}{0.7\textwidth}{X} 「我是好牧人，好牧人為羊捨命。 \end{tabularx} \\ \\ \relax
10:12 & \begin{tabularx}{0.7\textwidth}{X} 雇工不是牧人，羊不是他自己的，他一看見狼來，就撇下羊群逃跑；狼抓住羊，把牠們趕散。 \end{tabularx} \\ \\ \relax
10:13 & \begin{tabularx}{0.7\textwidth}{X} 雇工逃走，因為他是雇工，對羊毫不關心。 \end{tabularx} \\ \\ \relax
10:14 & \begin{tabularx}{0.7\textwidth}{X} 我是好牧人；我認識我的羊，我的羊也認識我， \end{tabularx} \\ \\ \relax
10:15 & \begin{tabularx}{0.7\textwidth}{X} 正如父認識我，我也認識父一樣；並且我為羊捨命。 \end{tabularx} \\ \\ \relax
10:16 & \begin{tabularx}{0.7\textwidth}{X} 我另外有羊，不屬這圈裡的，我必須領牠們來，牠們也要聽我的聲音，並且要合成一群，歸一個牧人。 \end{tabularx} \\ \\ \relax
10:17 & \begin{tabularx}{0.7\textwidth}{X} 為此，我父愛我，因為我把命捨去，好再取回來。 \end{tabularx} \\ \\ \relax
10:18 & \begin{tabularx}{0.7\textwidth}{X} 沒有人奪去我的命，是我自己捨的；我有權捨棄，也有權再取回。這是我從我父所受的命令。」 \end{tabularx} \\ \\
[1ex]
\hline
\hline
\end{longtable}


\newpage

\section{第79屆~港九培靈會~奧思邦~第六講}
\label{sec:296}
\textbf{盡用資源,全為基督}
\newline
\newline
連結: \href{https://www.hkbibleconference.org/session-message/view/296}{\texttt{https://www.hkbibleconference.org/session-message/view/296}}
\newline
\newline
\hyperref[sec:295]{< < < PREV SERMON < < <}
~
\hyperref[sec:toc]{[返目錄]}
~
\hyperref[sec:297]{> > > NEXT SERMON > > >}
\newline
\newline
路加福音 16:1-16:13
\newline
\begin{longtable}{cl}
\hline
\hline
章節 & 經文 (和合本修訂版)\\
\hline
16:1 & \begin{tabularx}{0.7\textwidth}{X} 耶穌又對門徒說：「某財主有一個管家，有人向主人告管家浪費他的財物。 \end{tabularx} \\ \\ \relax
16:2 & \begin{tabularx}{0.7\textwidth}{X} 主人叫他來，對他說：『我聽到了，你做的是甚麼事？把你所經管的交代清楚，你不能再作我的管家了。』 \end{tabularx} \\ \\ \relax
16:3 & \begin{tabularx}{0.7\textwidth}{X} 那管家心裡說：『主人辭我，不用我再作管家，我將來做甚麼呢？鋤地嘛，沒有力氣；討飯嘛，怕羞。 \end{tabularx} \\ \\ \relax
16:4 & \begin{tabularx}{0.7\textwidth}{X} 我知道怎麼做，好叫人們在我不作管家之後，接我到他們家裡去。』 \end{tabularx} \\ \\ \relax
16:5 & \begin{tabularx}{0.7\textwidth}{X} 於是他把欠他主人債的，一個一個地叫了來，問頭一個說：『你欠我主人多少？』 \end{tabularx} \\ \\ \relax
16:6 & \begin{tabularx}{0.7\textwidth}{X} 他說：『一百簍油。』管家對他說：『拿你的賬，快坐下，寫五十。』 \end{tabularx} \\ \\ \relax
16:7 & \begin{tabularx}{0.7\textwidth}{X} 他問另一個說：『你欠多少？』他說：『一百石麥子。』管家對他說：『拿你的賬，寫八十。』 \end{tabularx} \\ \\ \relax
16:8 & \begin{tabularx}{0.7\textwidth}{X} 主人就誇獎這不義的管家做事精明，因為今世之子應付自己的世代比光明之子更加精明。 \end{tabularx} \\ \\ \relax
16:9 & \begin{tabularx}{0.7\textwidth}{X} 我又告訴你們，要藉著那不義的錢財結交朋友，到了錢財無用的時候，他們可以接你們到永遠的住處去。 \end{tabularx} \\ \\ \relax
16:10 & \begin{tabularx}{0.7\textwidth}{X} 人在最小的事上忠心，在大事上也忠心；在最小的事上不義，在大事上也不義。 \end{tabularx} \\ \\ \relax
16:11 & \begin{tabularx}{0.7\textwidth}{X} 若是你們在不義的錢財上不忠心，誰還把那真實的錢財託付你們呢？ \end{tabularx} \\ \\ \relax
16:12 & \begin{tabularx}{0.7\textwidth}{X} 如果你們在別人的東西上不忠心，誰還把你們自己的東西給你們呢？ \end{tabularx} \\ \\ \relax
16:13 & \begin{tabularx}{0.7\textwidth}{X} 一個僕人不能服侍兩個主；他不是恨這個愛那個，就是重這個輕那個。你們不能又服侍神，又服侍瑪門。」 \end{tabularx} \\ \\
[1ex]
\hline
\hline
\end{longtable}


\newpage

\section{第79屆~港九培靈會~奧思邦~第七講}
\label{sec:297}
\textbf{豈可或缺:彼此饒恕}
\newline
\newline
連結: \href{https://www.hkbibleconference.org/session-message/view/297}{\texttt{https://www.hkbibleconference.org/session-message/view/297}}
\newline
\newline
\hyperref[sec:296]{< < < PREV SERMON < < <}
~
\hyperref[sec:toc]{[返目錄]}
~
\hyperref[sec:298]{> > > NEXT SERMON > > >}
\newline
\newline
馬太福音 18:21-18:35
\newline
\begin{longtable}{cl}
\hline
\hline
章節 & 經文 (和合本修訂版)\\
\hline
18:21 & \begin{tabularx}{0.7\textwidth}{X} 那時，彼得進前來，對耶穌說：「主啊，我弟兄得罪我，我當饒恕他幾次呢？到七次夠嗎？」 \end{tabularx} \\ \\ \relax
18:22 & \begin{tabularx}{0.7\textwidth}{X} 耶穌說：「我告訴你，不是到七次，而是到七十個七次。 \end{tabularx} \\ \\ \relax
18:23 & \begin{tabularx}{0.7\textwidth}{X} 因為天國好像一個王要和他僕人算賬。 \end{tabularx} \\ \\ \relax
18:24 & \begin{tabularx}{0.7\textwidth}{X} 他開始算的時候，有人帶了一個欠一萬他連得的僕人來。 \end{tabularx} \\ \\ \relax
18:25 & \begin{tabularx}{0.7\textwidth}{X} 因為他沒有甚麼償還之物，主人下令把他和他妻子兒女，以及一切所有的都賣了來償還。 \end{tabularx} \\ \\ \relax
18:26 & \begin{tabularx}{0.7\textwidth}{X} 那僕人就俯伏向他叩頭，說：『寬容我吧，我都會還你的。』 \end{tabularx} \\ \\ \relax
18:27 & \begin{tabularx}{0.7\textwidth}{X} 那僕人的主人就動了慈心，把他釋放了，並且免了他的債。 \end{tabularx} \\ \\ \relax
18:28 & \begin{tabularx}{0.7\textwidth}{X} 那僕人出來，遇見一個欠他一百個銀幣的同伴，就揪著他，扼住他的喉嚨，說：『把你所欠的還我！』 \end{tabularx} \\ \\ \relax
18:29 & \begin{tabularx}{0.7\textwidth}{X} 他的同伴就俯伏央求他，說：『寬容我吧，我會還你的。』 \end{tabularx} \\ \\ \relax
18:30 & \begin{tabularx}{0.7\textwidth}{X} 他不肯，卻把他下在監裡，直到他還了所欠的債。 \end{tabularx} \\ \\ \relax
18:31 & \begin{tabularx}{0.7\textwidth}{X} 同伴們看見他所做的事就很悲憤，把這一切的事都告訴了主人。 \end{tabularx} \\ \\ \relax
18:32 & \begin{tabularx}{0.7\textwidth}{X} 於是主人叫了他來，對他說：『你這惡奴才！你央求我，我就把你所欠的都免了； \end{tabularx} \\ \\ \relax
18:33 & \begin{tabularx}{0.7\textwidth}{X} 你不應該憐憫你的同伴，像我憐憫你嗎？』 \end{tabularx} \\ \\ \relax
18:34 & \begin{tabularx}{0.7\textwidth}{X} 主人就大怒，把他交給司刑的，直到他還清了所欠的債。 \end{tabularx} \\ \\ \relax
18:35 & \begin{tabularx}{0.7\textwidth}{X} 你們各人若不從心裡饒恕你的弟兄，我天父也要這樣待你們。」 \end{tabularx} \\ \\
[1ex]
\hline
\hline
\end{longtable}


\newpage

\section{第79屆~港九培靈會~奧思邦~第八講}
\label{sec:298}
\textbf{信徒責任──才幹歸神}
\newline
\newline
連結: \href{https://www.hkbibleconference.org/session-message/view/298}{\texttt{https://www.hkbibleconference.org/session-message/view/298}}
\newline
\newline
\hyperref[sec:297]{< < < PREV SERMON < < <}
~
\hyperref[sec:toc]{[返目錄]}
~
\hyperref[sec:299]{> > > NEXT SERMON > > >}
\newline
\newline
馬太福音 25:14-25:30
\newline
\begin{longtable}{cl}
\hline
\hline
章節 & 經文 (和合本修訂版)\\
\hline
25:14 & \begin{tabularx}{0.7\textwidth}{X} 「天國又好比一個人要出外遠行，就叫了僕人來，把他的家業交給他們。 \end{tabularx} \\ \\ \relax
25:15 & \begin{tabularx}{0.7\textwidth}{X} 他按著各人的才幹，給他們銀子：一個給了五千，一個給了二千，一個給了一千，就出外遠行去了。 \end{tabularx} \\ \\ \relax
25:16 & \begin{tabularx}{0.7\textwidth}{X} 那領五千的立刻拿去做買賣，另外賺了五千。 \end{tabularx} \\ \\ \relax
25:17 & \begin{tabularx}{0.7\textwidth}{X} 那領二千的也照樣另賺了二千。 \end{tabularx} \\ \\ \relax
25:18 & \begin{tabularx}{0.7\textwidth}{X} 但那領一千的去掘開地，把主人的銀子埋藏了。 \end{tabularx} \\ \\ \relax
25:19 & \begin{tabularx}{0.7\textwidth}{X} 過了許久，那些僕人的主人來了，和他們算賬。 \end{tabularx} \\ \\ \relax
25:20 & \begin{tabularx}{0.7\textwidth}{X} 那領五千的又帶著另外的五千來，說：『主啊，你交給我五千。請看，我又賺了五千。』 \end{tabularx} \\ \\ \relax
25:21 & \begin{tabularx}{0.7\textwidth}{X} 主人說：『好，你這又善良又忠心的僕人，你在少許的事上忠心，我要派你管理許多的事，進來享受你主人的快樂吧！』 \end{tabularx} \\ \\ \relax
25:22 & \begin{tabularx}{0.7\textwidth}{X} 那領二千的也進前來，說：『主啊，你交給我二千。請看，我又賺了二千。』 \end{tabularx} \\ \\ \relax
25:23 & \begin{tabularx}{0.7\textwidth}{X} 主人說：『好，你這又善良又忠心的僕人，你在少許的事上忠心，我要派你管理許多的事，進來享受你主人的快樂吧！』 \end{tabularx} \\ \\ \relax
25:24 & \begin{tabularx}{0.7\textwidth}{X} 那領一千的也進前來，說：『主啊，我知道你，你是個嚴厲的人：沒有種的地方也要收割，沒有播的地方也要收穫， \end{tabularx} \\ \\ \relax
25:25 & \begin{tabularx}{0.7\textwidth}{X} 我就害怕，去把你的一千銀子埋藏在地裡。請看，你的銀子在這裡。』 \end{tabularx} \\ \\ \relax
25:26 & \begin{tabularx}{0.7\textwidth}{X} 他的主人回答他說：『你這又惡又懶的僕人，你既知道我沒有種的地方也要收割，沒有播的地方也要收穫， \end{tabularx} \\ \\ \relax
25:27 & \begin{tabularx}{0.7\textwidth}{X} 就該把我的銀子放給兌換銀錢的人，到我來的時候可以連本帶利收回。 \end{tabularx} \\ \\ \relax
25:28 & \begin{tabularx}{0.7\textwidth}{X} 把他這一千奪過來，給那有一萬的。 \end{tabularx} \\ \\ \relax
25:29 & \begin{tabularx}{0.7\textwidth}{X} 因為凡有的，還要加給他，叫他有餘；沒有的，連他所有的也要奪過來。 \end{tabularx} \\ \\ \relax
25:30 & \begin{tabularx}{0.7\textwidth}{X} 把這無用的僕人丟在外面黑暗裡，在那裡他要哀哭切齒了。』」 \end{tabularx} \\ \\
[1ex]
\hline
\hline
\end{longtable}


\newpage

\section{第79屆~港九培靈會~奧思邦~第九講}
\label{sec:299}
\textbf{信徒生活──禱告核心}
\newline
\newline
連結: \href{https://www.hkbibleconference.org/session-message/view/299}{\texttt{https://www.hkbibleconference.org/session-message/view/299}}
\newline
\newline
\hyperref[sec:298]{< < < PREV SERMON < < <}
~
\hyperref[sec:toc]{[返目錄]}
~
\hyperref[sec:300]{> > > NEXT SERMON > > >}
\newline
\newline
路加福音 11:1-11:13
\newline
\begin{longtable}{cl}
\hline
\hline
章節 & 經文 (和合本修訂版)\\
\hline
11:1 & \begin{tabularx}{0.7\textwidth}{X} 耶穌在一個地方禱告。禱告完了，有個門徒對他說：「主啊，求你教導我們禱告，像約翰教導他的門徒一樣。」 \end{tabularx} \\ \\ \relax
11:2 & \begin{tabularx}{0.7\textwidth}{X} 耶穌對他們說：「你們禱告的時候，要說：『父啊，願人都尊你的名為聖；願你的國降臨； \end{tabularx} \\ \\ \relax
11:3 & \begin{tabularx}{0.7\textwidth}{X} 我們日用的飲食，天天賜給我們。 \end{tabularx} \\ \\ \relax
11:4 & \begin{tabularx}{0.7\textwidth}{X} 赦免我們的罪，因為我們也赦免凡虧欠我們的人。不叫我們陷入試探。』」 \end{tabularx} \\ \\ \relax
11:5 & \begin{tabularx}{0.7\textwidth}{X} 耶穌又對他們說：「你們中間誰有一個朋友半夜到他那裡去，對他說：『朋友！請借給我三個餅； \end{tabularx} \\ \\ \relax
11:6 & \begin{tabularx}{0.7\textwidth}{X} 因為我有一個朋友旅途中來到我這裡，我沒有東西招待他。』 \end{tabularx} \\ \\ \relax
11:7 & \begin{tabularx}{0.7\textwidth}{X} 那人在裡面回答：『不要打擾我，門已經關了，孩子們也同我在床上了，我不能起來給你。』 \end{tabularx} \\ \\ \relax
11:8 & \begin{tabularx}{0.7\textwidth}{X} 我告訴你們，雖不因他是朋友起來給他，也會因他不顧面子地直求，起來照他所需要的給他。 \end{tabularx} \\ \\ \relax
11:9 & \begin{tabularx}{0.7\textwidth}{X} 我又告訴你們，祈求，就給你們；尋找，就找到；叩門，就給你們開門。 \end{tabularx} \\ \\ \relax
11:10 & \begin{tabularx}{0.7\textwidth}{X} 因為凡祈求的，就得著；尋找的，就找到；叩門的，就給他開門。 \end{tabularx} \\ \\ \relax
11:11 & \begin{tabularx}{0.7\textwidth}{X} 你們中間作父親的，誰有兒子求魚，反拿蛇當魚給他呢？ \end{tabularx} \\ \\ \relax
11:12 & \begin{tabularx}{0.7\textwidth}{X} 求雞蛋，反給他蠍子呢？ \end{tabularx} \\ \\ \relax
11:13 & \begin{tabularx}{0.7\textwidth}{X} 你們雖然不好，尚且知道拿好東西給兒女，何況天父，他豈不更要把聖靈賜給求他的人嗎？」 \end{tabularx} \\ \\
[1ex]
\hline
\hline
\end{longtable}


\newpage

\section{第79屆~港九培靈會~張慕皚~第一講}
\label{sec:300}
\textbf{回歸真愛的豐盛中}
\newline
\newline
連結: \href{https://www.hkbibleconference.org/session-message/view/300}{\texttt{https://www.hkbibleconference.org/session-message/view/300}}
\newline
\newline
\hyperref[sec:299]{< < < PREV SERMON < < <}
~
\hyperref[sec:toc]{[返目錄]}
~
\hyperref[sec:301]{> > > NEXT SERMON > > >}
\newline
\newline
瑪拉基書 1:1-1:5
\newline
\begin{longtable}{cl}
\hline
\hline
章節 & 經文 (和合本修訂版)\\
\hline
1:1 & \begin{tabularx}{0.7\textwidth}{X} 耶和華的話，藉瑪拉基傳給以色列的默示。 \end{tabularx} \\ \\ \relax
1:2 & \begin{tabularx}{0.7\textwidth}{X} 耶和華說：「我曾愛你們。」你們卻說：「你在何事上愛我們呢？」耶和華說：「以掃不是雅各的哥哥嗎？我卻愛雅各， \end{tabularx} \\ \\ \relax
1:3 & \begin{tabularx}{0.7\textwidth}{X} 惡以掃，使他的山嶺荒涼，把他的地業交給曠野的野狗。」 \end{tabularx} \\ \\ \relax
1:4 & \begin{tabularx}{0.7\textwidth}{X} 以東若說：「我們雖被毀壞，卻要重建荒廢之處。」萬軍之耶和華如此說：「任他們建造，我必拆毀；人必稱他們為『邪惡之境』，為『耶和華永遠惱怒之民』。」 \end{tabularx} \\ \\ \relax
1:5 & \begin{tabularx}{0.7\textwidth}{X} 你們必親眼看見，你們要說：「耶和華在以色列疆界之外必尊為大！」 \end{tabularx} \\ \\
[1ex]
\hline
\hline
\end{longtable}


\newpage

\section{第79屆~港九培靈會~張慕皚~第二講}
\label{sec:301}
\textbf{回歸真敬拜的豐盛中}
\newline
\newline
連結: \href{https://www.hkbibleconference.org/session-message/view/301}{\texttt{https://www.hkbibleconference.org/session-message/view/301}}
\newline
\newline
\hyperref[sec:300]{< < < PREV SERMON < < <}
~
\hyperref[sec:toc]{[返目錄]}
~
\hyperref[sec:302]{> > > NEXT SERMON > > >}
\newline
\newline
瑪拉基書 1:6-1:14
\newline
\begin{longtable}{cl}
\hline
\hline
章節 & 經文 (和合本修訂版)\\
\hline
1:6 & \begin{tabularx}{0.7\textwidth}{X} 萬軍之耶和華對你們說：「兒子孝敬父親，僕人敬畏主人；我既為父親，孝敬我的在哪裡呢？我既為主人，敬畏我的在哪裡呢？你們這些藐視我名的祭司啊！」你們卻說：「我們在何事上藐視你的名呢？」 \end{tabularx} \\ \\ \relax
1:7 & \begin{tabularx}{0.7\textwidth}{X} 「你們將不潔淨的食物獻在我的祭壇上，卻說：『我們在何事上使你不潔淨呢？』你們說，耶和華的供桌是可藐視的。 \end{tabularx} \\ \\ \relax
1:8 & \begin{tabularx}{0.7\textwidth}{X} 你們將瞎眼的獻為祭物，這不算為惡嗎？將瘸腿的、有病的獻上，這不算為惡嗎？那麼，請把這些獻給你的省長，他豈會悅納你，豈會抬舉你呢？這是萬軍之耶和華說的。」 \end{tabularx} \\ \\ \relax
1:9 & \begin{tabularx}{0.7\textwidth}{X} 現在我勸你們要懇求神，好讓他施恩給我們。這事既出於你們的手，他豈會抬舉你們任何人呢？這是萬軍之耶和華說的。 \end{tabularx} \\ \\ \relax
1:10 & \begin{tabularx}{0.7\textwidth}{X} 萬軍之耶和華說：「甚願你們中間有人把殿的門關上，免得你們徒然在我壇上燒火。我不喜歡你們，也不從你們手中悅納供物。」 \end{tabularx} \\ \\ \relax
1:11 & \begin{tabularx}{0.7\textwidth}{X} 萬軍之耶和華說：「從日出之地到日落之處，我的名在列國中必尊為大。在各處，人必奉我的名燒香，獻潔淨的供物，因為我的名在列國中必尊為大。 \end{tabularx} \\ \\ \relax
1:12 & \begin{tabularx}{0.7\textwidth}{X} 你們卻褻瀆我的名，說：『主的供桌是不潔淨的，供桌上的果子和食物是可藐視的。』 \end{tabularx} \\ \\ \relax
1:13 & \begin{tabularx}{0.7\textwidth}{X} 你們又說：『看哪，這些事何等煩瑣！』並嗤之以鼻。這是萬軍之耶和華說的。你們把搶來的、瘸腿的、有病的拿來獻上為祭，我豈能從你們手中悅納它呢？這是耶和華說的。 \end{tabularx} \\ \\ \relax
1:14 & \begin{tabularx}{0.7\textwidth}{X} 行詭詐的人是可詛咒的！他的群畜中雖有公羊，他許了願，卻將有殘疾的獻給主。因我是大君王，我的名在列國中是可畏的。這是萬軍之耶和華說的。」 \end{tabularx} \\ \\
[1ex]
\hline
\hline
\end{longtable}


\newpage

\section{第79屆~港九培靈會~張慕皚~第三講}
\label{sec:302}
\textbf{回歸行道的豐盛中}
\newline
\newline
連結: \href{https://www.hkbibleconference.org/session-message/view/302}{\texttt{https://www.hkbibleconference.org/session-message/view/302}}
\newline
\newline
\hyperref[sec:301]{< < < PREV SERMON < < <}
~
\hyperref[sec:toc]{[返目錄]}
~
\hyperref[sec:303]{> > > NEXT SERMON > > >}
\newline
\newline
瑪拉基書 2:1-2:9
\newline
\begin{longtable}{cl}
\hline
\hline
章節 & 經文 (和合本修訂版)\\
\hline
2:1 & \begin{tabularx}{0.7\textwidth}{X} 現在，眾祭司啊，這誡命是給你們的。 \end{tabularx} \\ \\ \relax
2:2 & \begin{tabularx}{0.7\textwidth}{X} 萬軍之耶和華說：「你們若不聽，若不放在心上，將榮耀歸給我的名，我就使詛咒臨到你們，使你們的福分變為詛咒；其實我已經詛咒了你們的福分，因你們不把誡命放在心上。 \end{tabularx} \\ \\ \relax
2:3 & \begin{tabularx}{0.7\textwidth}{X} 看哪，我要斥責你們的後裔，把糞抹在你們臉上，就是你們祭牲的糞。人要把你們和糞一起抬出去， \end{tabularx} \\ \\ \relax
2:4 & \begin{tabularx}{0.7\textwidth}{X} 你們就知道我頒這誡命給你們，使我與利未所立的約可以常存。這是萬軍之耶和華說的。 \end{tabularx} \\ \\ \relax
2:5 & \begin{tabularx}{0.7\textwidth}{X} 我曾與他立生命和平安的約。我將這兩樣賜給他，使他存敬畏的心；他就敬畏我，懼怕我的名。 \end{tabularx} \\ \\ \relax
2:6 & \begin{tabularx}{0.7\textwidth}{X} 真實的訓誨在他口中，他的嘴唇中沒有不義。他以平安和正直與我同行，使許多人回轉離開罪孽。 \end{tabularx} \\ \\ \relax
2:7 & \begin{tabularx}{0.7\textwidth}{X} 祭司的嘴唇當守護知識，人也當從他口中尋求訓誨，因為他是萬軍之耶和華的使者。 \end{tabularx} \\ \\ \relax
2:8 & \begin{tabularx}{0.7\textwidth}{X} 你們卻偏離正道，使許多人在這訓誨上絆跌。你們破壞了我與利未人所立的約。這是萬軍之耶和華說的。 \end{tabularx} \\ \\ \relax
2:9 & \begin{tabularx}{0.7\textwidth}{X} 所以我使你們被眾百姓藐視，看為卑賤；因你們不遵守我的道，在律法上看人的情面。」 \end{tabularx} \\ \\
[1ex]
\hline
\hline
\end{longtable}


\newpage

\section{第79屆~港九培靈會~張慕皚~第四講}
\label{sec:303}
\textbf{回歸貞潔生活的豐盛中(一)}
\newline
\newline
連結: \href{https://www.hkbibleconference.org/session-message/view/303}{\texttt{https://www.hkbibleconference.org/session-message/view/303}}
\newline
\newline
\hyperref[sec:302]{< < < PREV SERMON < < <}
~
\hyperref[sec:toc]{[返目錄]}
~
\hyperref[sec:304]{> > > NEXT SERMON > > >}
\newline
\newline
瑪拉基書 2:10-2:16
\newline
\begin{longtable}{cl}
\hline
\hline
章節 & 經文 (和合本修訂版)\\
\hline
2:10 & \begin{tabularx}{0.7\textwidth}{X} 我們豈不都有一位父嗎？豈不是一位神創造了我們嗎？為何互相行詭詐，褻瀆了神與我們列祖所立的約呢？ \end{tabularx} \\ \\ \relax
2:11 & \begin{tabularx}{0.7\textwidth}{X} 猶大行事詭詐，在以色列和耶路撒冷中行了可憎的事；因為猶大褻瀆耶和華所喜愛的聖殿，娶外邦神明的女子為妻。 \end{tabularx} \\ \\ \relax
2:12 & \begin{tabularx}{0.7\textwidth}{X} 凡做這事的，無論是清醒的或回應的，即使獻供物給萬軍之耶和華，耶和華也要將他從雅各的帳棚中剪除。 \end{tabularx} \\ \\ \relax
2:13 & \begin{tabularx}{0.7\textwidth}{X} 你們又再做這樣的事，使哭泣和嘆息的眼淚遮蓋耶和華的祭壇，以致耶和華不再理會那供物，也不喜歡從你們的手中收納。 \end{tabularx} \\ \\ \relax
2:14 & \begin{tabularx}{0.7\textwidth}{X} 你們還說：「這是為甚麼呢？」因為耶和華在你和你年輕時所娶的妻之間作證。她雖是你的配偶，你誓約的妻，你卻背棄她。 \end{tabularx} \\ \\ \relax
2:15 & \begin{tabularx}{0.7\textwidth}{X} 一個人如果還剩下一點靈性，他不會這麼做。這人在尋找甚麼呢？神的後裔！當謹守你們的靈性，誰也不可背棄年輕時所娶的妻。 \end{tabularx} \\ \\ \relax
2:16 & \begin{tabularx}{0.7\textwidth}{X} 耶和華－以色列的神說：「我恨惡休妻的事和衣服外面披上暴力的人。所以當謹守你們的心，不可行詭詐。這是萬軍之耶和華說的。」 \end{tabularx} \\ \\
[1ex]
\hline
\hline
\end{longtable}


\newpage

\section{第79屆~港九培靈會~張慕皚~第五講}
\label{sec:304}
\textbf{回歸貞潔生活的豐盛中(二)}
\newline
\newline
連結: \href{https://www.hkbibleconference.org/session-message/view/304}{\texttt{https://www.hkbibleconference.org/session-message/view/304}}
\newline
\newline
\hyperref[sec:303]{< < < PREV SERMON < < <}
~
\hyperref[sec:toc]{[返目錄]}
~
\hyperref[sec:305]{> > > NEXT SERMON > > >}
\newline
\newline
瑪拉基書 2:10-2:16
\newline
\begin{longtable}{cl}
\hline
\hline
章節 & 經文 (和合本修訂版)\\
\hline
2:10 & \begin{tabularx}{0.7\textwidth}{X} 我們豈不都有一位父嗎？豈不是一位神創造了我們嗎？為何互相行詭詐，褻瀆了神與我們列祖所立的約呢？ \end{tabularx} \\ \\ \relax
2:11 & \begin{tabularx}{0.7\textwidth}{X} 猶大行事詭詐，在以色列和耶路撒冷中行了可憎的事；因為猶大褻瀆耶和華所喜愛的聖殿，娶外邦神明的女子為妻。 \end{tabularx} \\ \\ \relax
2:12 & \begin{tabularx}{0.7\textwidth}{X} 凡做這事的，無論是清醒的或回應的，即使獻供物給萬軍之耶和華，耶和華也要將他從雅各的帳棚中剪除。 \end{tabularx} \\ \\ \relax
2:13 & \begin{tabularx}{0.7\textwidth}{X} 你們又再做這樣的事，使哭泣和嘆息的眼淚遮蓋耶和華的祭壇，以致耶和華不再理會那供物，也不喜歡從你們的手中收納。 \end{tabularx} \\ \\ \relax
2:14 & \begin{tabularx}{0.7\textwidth}{X} 你們還說：「這是為甚麼呢？」因為耶和華在你和你年輕時所娶的妻之間作證。她雖是你的配偶，你誓約的妻，你卻背棄她。 \end{tabularx} \\ \\ \relax
2:15 & \begin{tabularx}{0.7\textwidth}{X} 一個人如果還剩下一點靈性，他不會這麼做。這人在尋找甚麼呢？神的後裔！當謹守你們的靈性，誰也不可背棄年輕時所娶的妻。 \end{tabularx} \\ \\ \relax
2:16 & \begin{tabularx}{0.7\textwidth}{X} 耶和華－以色列的神說：「我恨惡休妻的事和衣服外面披上暴力的人。所以當謹守你們的心，不可行詭詐。這是萬軍之耶和華說的。」 \end{tabularx} \\ \\
[1ex]
\hline
\hline
\end{longtable}


\newpage

\section{第79屆~港九培靈會~張慕皚~第六講}
\label{sec:305}
\textbf{回歸真信心的豐盛中}
\newline
\newline
連結: \href{https://www.hkbibleconference.org/session-message/view/305}{\texttt{https://www.hkbibleconference.org/session-message/view/305}}
\newline
\newline
\hyperref[sec:304]{< < < PREV SERMON < < <}
~
\hyperref[sec:toc]{[返目錄]}
~
\hyperref[sec:306]{> > > NEXT SERMON > > >}
\newline
\newline
瑪拉基書 2:17-3:6
\newline
\begin{longtable}{cl}
\hline
\hline
章節 & 經文 (和合本修訂版)\\
\hline
2:17 & \begin{tabularx}{0.7\textwidth}{X} 你們用言語使耶和華厭煩，卻說：「我們在何事上使他厭煩呢？」因為你們說：「凡行惡的，耶和華看為善，並且喜愛他們；」又說：「公平的神在哪裡呢？」 \end{tabularx} \\ \\ \relax
3:1 & \begin{tabularx}{0.7\textwidth}{X} 萬軍之耶和華說：「看哪，我要差遣我的使者在我前面預備道路。你們所尋求的主必忽然來到他的殿；立約的使者，就是你們所仰慕的，看哪，快要來到。」 \end{tabularx} \\ \\ \relax
3:2 & \begin{tabularx}{0.7\textwidth}{X} 他來的日子，誰能當得起呢？他顯現的時候，誰能立得住呢？因為他如煉金匠的火，如漂洗者的鹼。 \end{tabularx} \\ \\ \relax
3:3 & \begin{tabularx}{0.7\textwidth}{X} 他必坐下如煉淨銀子的人，必潔淨利未人，熬煉他們像金銀一樣；他們就憑公義獻供物給耶和華。 \end{tabularx} \\ \\ \relax
3:4 & \begin{tabularx}{0.7\textwidth}{X} 那時，猶大和耶路撒冷所獻的供物必蒙耶和華悅納，彷彿古時之日、上古之年。 \end{tabularx} \\ \\ \relax
3:5 & \begin{tabularx}{0.7\textwidth}{X} 萬軍之耶和華說：「我必臨近你們，施行審判。我必速速作見證，警戒那些行邪術的、犯姦淫的、起假誓的、剝削雇工工錢的、欺壓孤兒寡婦的、屈枉寄居者的和不敬畏我的人。」 \end{tabularx} \\ \\ \relax
3:6 & \begin{tabularx}{0.7\textwidth}{X} 「我－耶和華是不改變的；所以，雅各的子孫啊，你們不致滅亡。 \end{tabularx} \\ \\
[1ex]
\hline
\hline
\end{longtable}


\newpage

\section{第79屆~港九培靈會~張慕皚~第七講}
\label{sec:306}
\textbf{回歸物質生活的豐盛中}
\newline
\newline
連結: \href{https://www.hkbibleconference.org/session-message/view/306}{\texttt{https://www.hkbibleconference.org/session-message/view/306}}
\newline
\newline
\hyperref[sec:305]{< < < PREV SERMON < < <}
~
\hyperref[sec:toc]{[返目錄]}
~
\hyperref[sec:307]{> > > NEXT SERMON > > >}
\newline
\newline
瑪拉基書 3:7-3:12
\newline
\begin{longtable}{cl}
\hline
\hline
章節 & 經文 (和合本修訂版)\\
\hline
3:7 & \begin{tabularx}{0.7\textwidth}{X} 從你們祖先的日子以來，你們就偏離我的律例而不遵守。現在你們要轉向我，我就轉向你們。這是萬軍之耶和華說的。你們卻說：『我們如何轉向呢？』 \end{tabularx} \\ \\ \relax
3:8 & \begin{tabularx}{0.7\textwidth}{X} 人豈可搶奪神呢？你們竟搶奪我！你們卻說：『我們在何事上搶奪你呢？』其實就是在你們當納的十分之一奉獻和當獻的供物上。 \end{tabularx} \\ \\ \relax
3:9 & \begin{tabularx}{0.7\textwidth}{X} 因你們全國上下都搶奪我的供物，詛咒就臨到你們身上。 \end{tabularx} \\ \\ \relax
3:10 & \begin{tabularx}{0.7\textwidth}{X} 你們要將當納的十分之一全然送入倉庫，使我家有糧，以此試試我，是否為你們敞開天上的窗戶，傾福與你們，甚至無處可容。這是萬軍之耶和華說的。 \end{tabularx} \\ \\ \relax
3:11 & \begin{tabularx}{0.7\textwidth}{X} 我必為你們斥責蝗蟲，不容牠毀壞你們的土產。你們田間的葡萄樹，果實未熟以先也不會掉落。這是萬軍之耶和華說的。 \end{tabularx} \\ \\ \relax
3:12 & \begin{tabularx}{0.7\textwidth}{X} 萬國必稱你們為有福的，因你們必成為喜樂之地。這是萬軍之耶和華說的。」 \end{tabularx} \\ \\
[1ex]
\hline
\hline
\end{longtable}


\newpage

\section{第79屆~港九培靈會~張慕皚~第八講}
\label{sec:307}
\textbf{回歸真屬靈的生命中(一)}
\newline
\newline
連結: \href{https://www.hkbibleconference.org/session-message/view/307}{\texttt{https://www.hkbibleconference.org/session-message/view/307}}
\newline
\newline
\hyperref[sec:306]{< < < PREV SERMON < < <}
~
\hyperref[sec:toc]{[返目錄]}
~
\hyperref[sec:308]{> > > NEXT SERMON > > >}
\newline
\newline
瑪拉基書 3:13-3:15
\newline
\begin{longtable}{cl}
\hline
\hline
章節 & 經文 (和合本修訂版)\\
\hline
3:13 & \begin{tabularx}{0.7\textwidth}{X} 耶和華說：「你們用話頂撞我。」你們卻說：「我們說了甚麼話頂撞你呢？」 \end{tabularx} \\ \\ \relax
3:14 & \begin{tabularx}{0.7\textwidth}{X} 你們說：「事奉神是枉然，我們遵守神所吩咐的，在萬軍之耶和華面前哀痛而行，有甚麼益處呢？ \end{tabularx} \\ \\ \relax
3:15 & \begin{tabularx}{0.7\textwidth}{X} 現在，我們稱狂傲的人為有福，並且行惡的人得以建立；他們雖然試探神，卻得以逃脫。」 \end{tabularx} \\ \\
[1ex]
\hline
\hline
\end{longtable}


\newpage

\section{第79屆~港九培靈會~張慕皚~第九講}
\label{sec:308}
\textbf{回歸真屬靈的豐盛中(二)}
\newline
\newline
連結: \href{https://www.hkbibleconference.org/session-message/view/308}{\texttt{https://www.hkbibleconference.org/session-message/view/308}}
\newline
\newline
\hyperref[sec:307]{< < < PREV SERMON < < <}
~
\hyperref[sec:toc]{[返目錄]}
~
\hyperref[sec:309]{> > > NEXT SERMON > > >}
\newline
\newline
瑪拉基書 3:16-3:18
\newline
\begin{longtable}{cl}
\hline
\hline
章節 & 經文 (和合本修訂版)\\
\hline
3:16 & \begin{tabularx}{0.7\textwidth}{X} 那時，敬畏耶和華的人彼此談論，耶和華側耳而聽，且有紀念冊在他面前，記錄那敬畏耶和華、思念他名的人。 \end{tabularx} \\ \\ \relax
3:17 & \begin{tabularx}{0.7\textwidth}{X} 萬軍之耶和華說：「在我所定的日子，他們必屬我，是我寶貴的產業。我必憐憫他們，如同人憐憫那服侍他的兒子。 \end{tabularx} \\ \\ \relax
3:18 & \begin{tabularx}{0.7\textwidth}{X} 那時你們必再一次看出義人和惡人，事奉神和不事奉神的人有何差別。」 \end{tabularx} \\ \\
[1ex]
\hline
\hline
\end{longtable}


\newpage

\section{第79屆~港九培靈會~張慕皚~第十講}
\label{sec:309}
\textbf{回歸真盼望的豐盛中}
\newline
\newline
連結: \href{https://www.hkbibleconference.org/session-message/view/309}{\texttt{https://www.hkbibleconference.org/session-message/view/309}}
\newline
\newline
\hyperref[sec:308]{< < < PREV SERMON < < <}
~
\hyperref[sec:toc]{[返目錄]}
~
\hyperref[sec:282]{> > > NEXT SERMON > > >}
\newline
\newline
瑪拉基書 4:1-4:6
\newline
\begin{longtable}{cl}
\hline
\hline
章節 & 經文 (和合本修訂版)\\
\hline
4:1 & \begin{tabularx}{0.7\textwidth}{X} 萬軍之耶和華說：「看哪，那日臨近，勢如燒著的火爐，凡狂傲的和行惡的都如碎秸，在那日被燒盡，根與枝條無一存留。 \end{tabularx} \\ \\ \relax
4:2 & \begin{tabularx}{0.7\textwidth}{X} 但是，對你們敬畏我名的人，必有公義的太陽出現，其光線有醫治的能力。你們必出來跳躍如圈裡的牛犢。 \end{tabularx} \\ \\ \relax
4:3 & \begin{tabularx}{0.7\textwidth}{X} 你們必踐踏惡人；在我所定的日子，他們必成為你們腳掌下的灰塵。這是萬軍之耶和華說的。 \end{tabularx} \\ \\ \relax
4:4 & \begin{tabularx}{0.7\textwidth}{X} 「你們當記念我僕人摩西的律法，就是我在何烈山為以色列眾人所吩咐他的律例典章。 \end{tabularx} \\ \\ \relax
4:5 & \begin{tabularx}{0.7\textwidth}{X} 「看哪，耶和華大而可畏之日未到以前，我要差遣以利亞先知到你們那裡去。 \end{tabularx} \\ \\ \relax
4:6 & \begin{tabularx}{0.7\textwidth}{X} 他必使父親的心轉向兒女，兒女的心轉向父親，免得我來詛咒這地。」 \end{tabularx} \\ \\
[1ex]
\hline
\hline
\end{longtable}


\newpage

\section{第79屆~港九培靈會~李寶珠~第一講}
\label{sec:282}
\textbf{主題介紹──前因後果}
\newline
\newline
連結: \href{https://www.hkbibleconference.org/session-message/view/282}{\texttt{https://www.hkbibleconference.org/session-message/view/282}}
\newline
\newline
\hyperref[sec:309]{< < < PREV SERMON < < <}
~
\hyperref[sec:toc]{[返目錄]}
~
\hyperref[sec:283]{> > > NEXT SERMON > > >}
\newline
\newline


\newpage

\section{第79屆~港九培靈會~李寶珠~第二講}
\label{sec:283}
\textbf{信的根據}
\newline
\newline
連結: \href{https://www.hkbibleconference.org/session-message/view/283}{\texttt{https://www.hkbibleconference.org/session-message/view/283}}
\newline
\newline
\hyperref[sec:282]{< < < PREV SERMON < < <}
~
\hyperref[sec:toc]{[返目錄]}
~
\hyperref[sec:284]{> > > NEXT SERMON > > >}
\newline
\newline


\newpage

\section{第79屆~港九培靈會~李寶珠~第三講}
\label{sec:284}
\textbf{信的証據}
\newline
\newline
連結: \href{https://www.hkbibleconference.org/session-message/view/284}{\texttt{https://www.hkbibleconference.org/session-message/view/284}}
\newline
\newline
\hyperref[sec:283]{< < < PREV SERMON < < <}
~
\hyperref[sec:toc]{[返目錄]}
~
\hyperref[sec:285]{> > > NEXT SERMON > > >}
\newline
\newline


\newpage

\section{第79屆~港九培靈會~李寶珠~第四講}
\label{sec:285}
\textbf{信的理據}
\newline
\newline
連結: \href{https://www.hkbibleconference.org/session-message/view/285}{\texttt{https://www.hkbibleconference.org/session-message/view/285}}
\newline
\newline
\hyperref[sec:284]{< < < PREV SERMON < < <}
~
\hyperref[sec:toc]{[返目錄]}
~
\hyperref[sec:286]{> > > NEXT SERMON > > >}
\newline
\newline


\newpage

\section{第79屆~港九培靈會~李寶珠~第五講}
\label{sec:286}
\textbf{愛的模式}
\newline
\newline
連結: \href{https://www.hkbibleconference.org/session-message/view/286}{\texttt{https://www.hkbibleconference.org/session-message/view/286}}
\newline
\newline
\hyperref[sec:285]{< < < PREV SERMON < < <}
~
\hyperref[sec:toc]{[返目錄]}
~
\hyperref[sec:287]{> > > NEXT SERMON > > >}
\newline
\newline


\newpage

\section{第79屆~港九培靈會~李寶珠~第六講}
\label{sec:287}
\textbf{愛的方式}
\newline
\newline
連結: \href{https://www.hkbibleconference.org/session-message/view/287}{\texttt{https://www.hkbibleconference.org/session-message/view/287}}
\newline
\newline
\hyperref[sec:286]{< < < PREV SERMON < < <}
~
\hyperref[sec:toc]{[返目錄]}
~
\hyperref[sec:288]{> > > NEXT SERMON > > >}
\newline
\newline


\newpage

\section{第79屆~港九培靈會~李寶珠~第七講}
\label{sec:288}
\textbf{望的焦點}
\newline
\newline
連結: \href{https://www.hkbibleconference.org/session-message/view/288}{\texttt{https://www.hkbibleconference.org/session-message/view/288}}
\newline
\newline
\hyperref[sec:287]{< < < PREV SERMON < < <}
~
\hyperref[sec:toc]{[返目錄]}
~
\hyperref[sec:289]{> > > NEXT SERMON > > >}
\newline
\newline


\newpage

\section{第79屆~港九培靈會~李寶珠~第八講}
\label{sec:289}
\textbf{望的眼光}
\newline
\newline
連結: \href{https://www.hkbibleconference.org/session-message/view/289}{\texttt{https://www.hkbibleconference.org/session-message/view/289}}
\newline
\newline
\hyperref[sec:288]{< < < PREV SERMON < < <}
~
\hyperref[sec:toc]{[返目錄]}
~
\hyperref[sec:290]{> > > NEXT SERMON > > >}
\newline
\newline


\newpage

\section{第79屆~港九培靈會~李寶珠~第九講}
\label{sec:290}
\textbf{望的視力}
\newline
\newline
連結: \href{https://www.hkbibleconference.org/session-message/view/290}{\texttt{https://www.hkbibleconference.org/session-message/view/290}}
\newline
\newline
\hyperref[sec:289]{< < < PREV SERMON < < <}
~
\hyperref[sec:toc]{[返目錄]}
~
\hyperref[sec:266]{> > > NEXT SERMON > > >}
\newline
\newline
\allsectionsfont{\centering}

\setlength\parindent{0pt}
\setlength{\columnsep}{1.25em}
\setlength{\parfillskip}{0pt}
\setlength{\tabcolsep}{1em}
\raggedbottom

\pagenumbering{gobble}


\newfontfamily\leftfont[Path=../fonts/fell_french_canon/, Ligatures=TeX, ItalicFont=IMFeFCit29C.otf, BoldFont=AveriaLibre-Bold.ttf]{IMFeFCrm29C.otf}
\newfontfamily\leftcitationfont[Path=../fonts/frankruehl/]{FrankRuehlCLM-Medium.ttf}
\newfontfamily\centerfont[Path=../fonts/garamond/, Ligatures=TeX, ItalicFont=EBGaramond-SemiBoldItalic.ttf]{EBGaramond-SemiBold.ttf}
\newfontfamily\rightfont[Path=../fonts/averia/, Ligatures=TeX, ItalicFont=AveriaLibre-RegularItalic.ttf, BoldFont=AveriaLibre-Bold.ttf, BoldItalicFont=AveriaLibre-BoldItalic.ttf]{AveriaLibre-Light.ttf}
\newfontfamily\rightcitationfont[Path=../fonts/rashi/]{Mekorot-Rashi.ttf}
\definecolor{hcolor}{HTML}{D3230C}
\definecolor{rcolor}{HTML}{D36F0C}
\newcommand{\chfont}[1]{\centerfont{\huge\textcolor{hcolor}{#1}}}
\newcommand{\leftcitation}[1]{\leftcitationfont{\Large\textcolor{hcolor}{#1}}}
\newcommand{\rightcitation}[1]{\rightcitationfont{\normalsize\textcolor{rcolor}{#1}}}
\newfontfamily\flowerfont[Path=../fonts/fell_flowers/]{IMFeFlow2.otf}

\begin{sloppypar}

\chapter*{\chfont{編按結語}}

\columnratio{0.5,0.5}\begin{paracol}{2}

\fontsize{11}{13}\leftfont \Large \leftcitation{א} \leftfont 余少好文．宏志博覽群書而不忘．善存藏經籍文獻備後時之用。\leftcitation{ב} \leftfont 歸主年時．受友所薦．聞道網海．\switchcolumn\fontsize{11}{13}\rightfont \Large \leftcitation{ח} \rightfont 有見粵道之危．國之封講道千言亦將就至．急之何則為？\leftcitation{ט} \rightfont 嘗聞猶太者之傳承．在其力守口述之

\end{paracol}


\columnratio{0.32,0.32,0.32}\begin{paracol}{3}

\fontsize{11}{13}\leftfont \Large 尤以吳約翰遜者 \switchcolumn[2]\fontsize{11}{13}\rightfont \Large 統．以煉千載不

\end{paracol}

\columnratio{0.32,0.32,0.32}
\begin{paracol}{3}\fontsize{11}{13}\leftfont \Large 為重．其載上之粵語講道緩緩入耳．收之藏其音頻．善妥整存．反復而嚼．受益無窮。\leftcitation{ג} \leftfont 我城我國既限．歷一四一九之不測．肺疫延年．信徒靈長屢受圍創．神州燈臺數盡指日可待．粵道之求與日俱增。\leftcitation{ד} \leftfont 觀乎社、經、法、媒、言、信、網之地．愈趨受鋤．自翔不果．授受壓力．粵道聖言亦愈漸艱難。\leftcitation{ה} \leftfont 況崇基例乎．學苑講道屢逆權勢者．其言末強受壓．舊章盡刪以存其身。講道釋數失傳．徒嘆奈何。

\switchcolumn

\fontsize{11}{13}\centerfont 
\begin{tikzpicture}
    \node (0,0) [xshift=-0.10cm, yshift=-1.0cm, opacity=0.10]{\includegraphics[width=0.30\textwidth]{../ot_frontcover.png}} ;
    \node (0,0) [xshift=+0.20cm, yshift=+2.0cm, opacity=0.10]{\includegraphics[width=0.20\textwidth]{../christ_on_cross.png}} ;
\end{tikzpicture}
\Large 

\leftcitation{ס} \centerfont 詩百又廿七載：
\leftcitation{ע} \centerfont 非耶和華建屋宇.則匠人之經營徒.
\leftcitation{פ} \centerfont 非耶和華衛城邑.則守者之儆醒徒.
\leftcitation{צ} \centerfont 余獻是卷予華人社區．願為福音流通之器．願獻斯微材為祭榮耀上帝．
\leftcitation{ק} \centerfont 阿門

\switchcolumn

\fontsize{11}{13}\rightfont \Large 滅．時越次聖殿期及當今。\leftcitation{י} \rightfont 猶太者力廣納之．筆錄以卷軸．便以傳、閱、頌、攜、守、鎖、抄、譯、釋、編，得書塔木德、密示拿等經傳．家喻戶曉．傳流若芳。\leftcitation{כ} \rightfont 猶太者文以載道．傳其口述．今我輩粵道之傳應當作如是．遂力行粵音識辨之法．載言載道．以盡忠傳粵道以待興。\leftcitation{ל} \rightfont 蒙下賜恩惠．無畏海量字音文書．既馭上帝之道．今廣及粵語講道．重駛編程之技．匯導粵音遂字稿．重塑講道現場．以傚猶太卷軸之舉便以傳流。\leftcitation{מ} \rightfont 是卷乃粵音口述傳之屬．莫通華文白話之語．

\end{paracol}

\columnratio{0.5,0.5}
\begin{paracol}{2}\fontsize{11}{13}\leftfont \Large \leftcitation{ו} \leftfont 斯殺一違儆百逆．既禁壓之．我輩聞風無奈．在所難免。\leftcitation{ז} \leftfont 另有異人例乎．以版權之名．脅網絡頻道之舉．同授礙予粵道之存流。

\switchcolumn

\fontsize{11}{13}\rightfont \Large 惟待後繼來者之傚．以譯釋傳之於神州華文地。\leftcitation{נ} \rightfont 今能排程驅馭圖靈以編彙文檔，其碼長共數千千亦無逢大礙．全蒙上帝保守。

\end{paracol}



\columnratio{1}\begin{paracol}{1}

\fontsize{11}{13}\rightfont \Large
~~~~~~~~~~~~~~~~~~~~~~~~~~~~~~~~~~~~~~~~~~~~~~~~~~~~~~~~~~~~~~~~~~~~~~~~~~~~~~~\leftcitation{ר} \rightfont 二零二三年二月一日

~~~~~~~~~~~~~~~~~~~~~~~~~~~~~~~~~~~~~~~~~~~~~~~~~~~~~~~~~~~~~~~~~~~~~~~~~~~~~~~\leftcitation{ש} \rightfont 米迦勒

~~~~~~~~~~~~~~~~~~~~~~~~~~~~~~~~~~~~~~~~~~~~~~~~~~~~~~~~~~~~~~~~~~~~~~~~~~~~~~~\leftcitation{ת} \rightfont 書於香港

\end{paracol}

\end{sloppypar}
\end{document}
